\documentclass[11pt,a4paper,oneside,openany]{stacks-project-book}

% Non-rendering source-note environments retained from the Stacks sources.
\usepackage{verbatim}
\newenvironment{reference}{\comment}{\endcomment}
\newenvironment{slogan}{\comment}{\endcomment}
\newenvironment{history}{\comment}{\endcomment}

% Diagrams and chapter-support packages.
\usepackage[all]{xy}
\xyoption{2cell}
\UseAllTwocells
\usepackage{multicol}
\usepackage{graphicx}

% Mainland Simplified-Chinese scientific typography.
\usepackage{fontspec}
\usepackage{xeCJK}
\defaultfontfeatures{Ligatures=TeX}
\setmainfont{Latin Modern Roman}
\setsansfont{Latin Modern Sans}
\setmonofont{Latin Modern Mono}
\setCJKmainfont[AutoFakeSlant=0.2]{Noto Serif CJK SC}
\setCJKsansfont{Noto Sans CJK SC}
\setCJKmonofont{Noto Sans CJK SC}
\newfontfamily\stackszhgreek{Noto Serif CJK SC}
\DeclareRobustCommand{\zhdelta}{{\stackszhgreek δ}}
\xeCJKsetup{PunctStyle=kaiming}
\XeTeXlinebreaklocale "zh"
\usepackage[a4paper,left=22mm,right=22mm,top=24mm,bottom=26mm,headheight=14pt,headsep=7mm,footskip=14mm]{geometry}
\linespread{1.18}
\setlength{\parindent}{2em}
\setlength{\parskip}{0pt}
\setlength{\emergencystretch}{1em}

% Localized structural names.
\renewcommand{\contentsname}{目录}
\renewcommand{\bibname}{参考文献}
\renewcommand{\proofname}{证明}
\renewcommand{\figurename}{图}
\renewcommand{\tablename}{表}
\renewcommand{\appendixname}{附录}

% Chinese numbered chapter heads while keeping numerical identifiers stable.
\makeatletter
\def\@makechapterhead#1{\global\topskip 7.5pc\relax
  \begingroup
  \fontsize{17}{22}\bfseries\centering
    \ifnum\c@secnumdepth>\m@ne
      \leavevmode \hskip-\leftskip
      \rlap{\vbox to\z@{\vss
          \ifx\chaptername\appendixname
            \centerline{\normalsize\mdseries 附录\enspace\thechapter}
          \else
            \centerline{\normalsize\mdseries 第\thechapter 章}
          \fi
          \vskip 3pc}}\hskip\leftskip\fi
     #1\par \endgroup
  \skip@34\p@ \advance\skip@-\normalbaselineskip
  \vskip\skip@ }
\makeatother

% Hyperlinks: current-volume labels resolve locally; absent chapters use tags.
\usepackage{hyperref}
\hypersetup{
  unicode=true,
  hidelinks,
  pdfborder={0 0 0},
  pdftitle={Stacks Project 简体中文版 - 116 章累积版},
  pdfauthor={The Stacks Project contributors; zh-Hans-CN translation production},
  pdfsubject={Unofficial AI-assisted Mainland Simplified Chinese translation of all 116 pinned Stacks Project chapters; not comprehensively human-reviewed}
}
\urlstyle{same}
% Generated from the frozen Stacks permanent-tag registry.
% Existing labels resolve inside this PDF; only absent labels fall back to tags.
\expandafter\def\csname stacksManual@foo-lemma-bar\endcsname{\href{https://github.com/stacks/stacks-project/blob/a04446e57ec1fbc252a871afcec7752fb2807b14/coding.tex}{\texttt{源码示例}}}
\DeclareRobustCommand{\StacksResolvedRef}[1]{%
  \ifcsname r@#1\endcsname
    \StacksOriginalRef{#1}%
  \else
    \ifcsname stacksTag@#1\endcsname
      \href{https://stacks.math.columbia.edu/tag/\csname stacksTag@#1\endcsname}{\texttt{\csname stacksTag@#1\endcsname}}%
    \else
      \ifcsname stacksManual@#1\endcsname
        \csname stacksManual@#1\endcsname%
      \else
        \StacksOriginalRef{#1}%
      \fi
    \fi
  \fi
}
\AtBeginDocument{%
  \let\StacksOriginalRef\ref
  \let\ref\StacksResolvedRef
}

% Theorem environments.
%
\theoremstyle{plain}
\newtheorem{theorem}[subsection]{定理}
\newtheorem{proposition}[subsection]{命题}
\newtheorem{lemma}[subsection]{引理}

\theoremstyle{definition}
\newtheorem{definition}[subsection]{定义}
\newtheorem{example}[subsection]{例}
\newtheorem{exercise}[subsection]{习题}
\newtheorem{situation}[subsection]{情形}

\theoremstyle{remark}
\newtheorem{remark}[subsection]{注}
\newtheorem{remarks}[subsection]{诸注}

\numberwithin{equation}{subsection}

% Macros
%
\def\lim{\mathop{\mathrm{lim}}\nolimits}
\def\colim{\mathop{\mathrm{colim}}\nolimits}
\def\Spec{\mathop{\mathrm{Spec}}}
\def\Hom{\mathop{\mathrm{Hom}}\nolimits}
\def\Ext{\mathop{\mathrm{Ext}}\nolimits}
\def\SheafHom{\mathop{\mathcal{H}\!\mathit{om}}\nolimits}
\def\SheafExt{\mathop{\mathcal{E}\!\mathit{xt}}\nolimits}
\def\Sch{\mathit{Sch}}
\def\Mor{\mathop{\mathrm{Mor}}\nolimits}
\def\Ob{\mathop{\mathrm{Ob}}\nolimits}
\def\Sh{\mathop{\mathit{Sh}}\nolimits}
\def\NL{\mathop{N\!L}\nolimits}
\def\CH{\mathop{\mathrm{CH}}\nolimits}
\def\proetale{{pro\text{-}\acute{e}tale}}
\def\etale{{\acute{e}tale}}
\def\QCoh{\mathit{QCoh}}
\def\Ker{\mathop{\mathrm{Ker}}}
\def\Im{\mathop{\mathrm{Im}}}
\def\Coker{\mathop{\mathrm{Coker}}}
\def\Coim{\mathop{\mathrm{Coim}}}

% Boxtimes
%
\DeclareMathSymbol{\boxtimes}{\mathbin}{AMSa}{"02}

%
% Macros for moduli stacks/spaces
%
\def\QCohstack{\mathcal{QC}\!\mathit{oh}}
\def\Cohstack{\mathcal{C}\!\mathit{oh}}
\def\Spacesstack{\mathcal{S}\!\mathit{paces}}
\def\Quotfunctor{\mathrm{Quot}}
\def\Hilbfunctor{\mathrm{Hilb}}
\def\Curvesstack{\mathcal{C}\!\mathit{urves}}
\def\Polarizedstack{\mathcal{P}\!\mathit{olarized}}
\def\Complexesstack{\mathcal{C}\!\mathit{omplexes}}
% \Pic is the operator that assigns to X its picard group, usage \Pic(X)
% \Picardstack_{X/B} denotes the Picard stack of X over B
% \Picardfunctor_{X/B} denotes the Picard functor of X over B
\def\Pic{\mathop{\mathrm{Pic}}\nolimits}
\def\Picardstack{\mathcal{P}\!\mathit{ic}}
\def\Picardfunctor{\mathrm{Pic}}
\def\Deformationcategory{\mathcal{D}\!\mathit{ef}}
\begin{document}
\frontmatter
\begin{titlepage}
\thispagestyle{empty}
\centering
\vspace*{25mm}
{\fontsize{30}{38}\selectfont\bfseries Stacks Project\par}
\vspace{11mm}
{\fontsize{27}{34}\selectfont\bfseries 简体中文版\par}
\vspace{8mm}
{\Large 116 章累积版\par}
\vspace{5mm}
{\large Mainland Simplified Chinese Cumulative Edition\par}
\vfill
{\large 当前收入 116 章\par}
\vspace{4mm}
{\normalsize 冻结英语 authority：\texttt{a04446e57ec1}\par}
\vspace{3mm}
{\normalsize 2026 年 9 月 19 日\par}
\vfill
{\small 人工智能辅助翻译；尚未经过全面的专家人工审校。\par}
\end{titlepage}

\chapter*{版本说明}
\addcontentsline{toc}{chapter}{版本说明}
本卷收入冻结英语版本的全部 116 章简体中文译文。章节覆盖完整并不意味着翻译质量已经得到全面认证。本译本由人工智能辅助生成，尚未经过全面的专家人工审校；术语、表达与翻译选择的核查仍在进行。欢迎相关领域专家提出修正意见。

中文 locale 明确为中国大陆规范的简体中文（\texttt{zh-Hans-CN}），不代表新加坡简体中文，也不代表台湾、香港或澳门的繁体中文本地化。正文采用 A4、11pt 正文字号、Noto Serif CJK SC、对称页边距和中文科技文献常用的首行缩进及行距。

英语 authority 固定为 Stacks Project 提交
\texttt{a04446e57ec1fbc252a871afcec7752fb2807b14}。翻译措辞的规范化与英语原文的数学勘误分开记录；本次修订不将数学勘误悄悄改入译文。翻译选择的依据、备选表达及尚存疑问见本版本 GitHub 仓库的\href{https://github.com/KokunoYumeto/stacks-zh-hans-cn/blob/main/review/translation-choices/README.md}{翻译选择记录}，可供专家审阅。卷内交叉引用在本 PDF 中解析，另保留通往 Stacks Project 永久标签的链接。

\begin{multicols}{2}
\begin{itemize}
\item 第 1 章：引言
\item 第 2 章：约定
\item 第 3 章：集合论
\item 第 4 章：范畴
\item 第 5 章：拓扑
\item 第 6 章：空间上的层
\item 第 7 章：位点与层
\item 第 8 章：栈
\item 第 9 章：域
\item 第 10 章：交换代数
\item 第 11 章：Brauer 群
\item 第 12 章：同调代数
\item 第 13 章：导出范畴
\item 第 14 章：单纯方法
\item 第 15 章：代数进阶
\item 第 16 章：环同态的光滑化
\item 第 17 章：模层
\item 第 18 章：站点上的模
\item 第 19 章：内射对象
\item 第 20 章：层上同调
\item 第 21 章：站点上的上同调
\item 第 22 章：微分分次代数
\item 第 23 章：除幂代数
\item 第 24 章：微分分次层
\item 第 25 章：超覆盖
\item 第 26 章：概形
\item 第 27 章：概形的构造
\item 第 28 章：概形的性质
\item 第 29 章：概形的态射
\item 第 30 章：概形的上同调
\item 第 31 章：除子
\item 第 32 章：概形的极限
\item 第 33 章：簇
\item 第 34 章：概形上的拓扑
\item 第 35 章：下降
\item 第 36 章：概形的导出范畴
\item 第 37 章：态射续篇
\item 第 38 章：平坦性续篇
\item 第 39 章：群胚概形
\item 第 40 章：群胚概形进阶
\item 第 41 章：概形的平展（\'Etale）态射
\item 第 42 章：Chow 同调与陈类
\item 第 43 章：相交理论
\item 第 44 章：曲线的 Picard 概形
\item 第 45 章：Weil 上同调理论
\item 第 46 章：适足模
\item 第 47 章：对偶化复形
\item 第 48 章：概形的对偶理论
\item 第 49 章：判别式与不同
\item 第 50 章：德拉姆上同调
\item 第 51 章：局部上同调
\item 第 52 章：代数几何与形式几何
\item 第 53 章：代数曲线
\item 第 54 章：曲面奇点的消解
\item 第 55 章：半稳定约化
\item 第 56 章：函子与态射
\item 第 57 章：代数簇的导出范畴
\item 第 58 章：概形的基本群
\item 第 59 章：\'Etale 上同调
\item 第 60 章：晶体上同调
\item 第 61 章：Pro-\'etale 上同调
\item 第 62 章：相对循环
\item 第 63 章：进一步的 \'Etale 上同调
\item 第 64 章：迹公式
\item 第 65 章：代数空间
\item 第 66 章：代数空间的性质
\item 第 67 章：代数空间的态射
\item 第 68 章：良好代数空间
\item 第 69 章：代数空间的上同调
\item 第 70 章：代数空间的极限
\item 第 71 章：代数空间上的除子
\item 第 72 章：域上的代数空间
\item 第 73 章：代数空间上的拓扑
\item 第 74 章：下降与代数空间
\item 第 75 章：空间的导出范畴
\item 第 76 章：空间态射进阶
\item 第 77 章：代数空间上的平坦性
\item 第 78 章：代数空间中的群胚
\item 第 79 章：代数空间中的群胚进阶
\item 第 80 章：自举
\item 第 81 章：代数空间的推出
\item 第 82 章：代数空间的周群
\item 第 83 章：群胚的商
\item 第 84 章：空间上同调续论
\item 第 85 章：单纯空间
\item 第 86 章：空间的对偶性
\item 第 87 章：形式代数空间
\item 第 88 章：形式空间的代数化
\item 第 89 章：再论曲面解消
\item 第 90 章：形式形变理论
\item 第 91 章：形变理论
\item 第 92 章：余切复形
\item 第 93 章：形变问题
\item 第 94 章：代数叠
\item 第 95 章：叠的例子
\item 第 96 章：代数叠上的层
\item 第 97 章：可表示性判据
\item 第 98 章：Artin 公理
\item 第 99 章：Quot 空间与 Hilbert 空间
\item 第 100 章：代数叠的性质
\item 第 101 章：代数叠的态射
\item 第 102 章：代数叠的极限
\item 第 103 章：代数叠的上同调
\item 第 104 章：叠的导出范畴
\item 第 105 章：代数叠导论
\item 第 106 章：叠态射的进一步讨论
\item 第 107 章：代数叠的几何
\item 第 108 章：模叠
\item 第 109 章：曲线模空间
\item 第 110 章：例
\item 第 111 章：练习
\item 第 112 章：文献指南
\item 第 113 章：待增内容
\item 第 114 章：编码风格
\item 第 115 章：已废弃内容
\item 第 116 章：GNU 自由文档许可证（GNU Free Documentation License）
\end{itemize}
\end{multicols}

原作版权：Copyright \copyright\ 2005--2025 Johan de Jong。原作与本衍生版依 GNU Free Documentation License 1.2 或其后版本发布，无不变章节、封面文字或封底文字。完整许可证见附录。原作者与贡献者、项目历史以及最新英语原文见 \url{https://stacks.math.columbia.edu/}。

\tableofcontents
\mainmatter
% Generated mechanically from the frozen introduction.tex target.
% Do not edit this generated file; edit the bound locale source instead.

% OK, start here.
%

\setcounter{chapter}{0}
\chapter{引言}



\phantomsection
\label{introduction-section-phantom}

\begin{verbatim}
Copyright (C) 2005 -- 2025 Johan de Jong
Permission is granted to copy, distribute and/or modify this
document under the terms of the GNU Free Documentation License,
Version 1.2 or any later version published by the Free Software
Foundation; with no Invariant Sections, no Front-Cover Texts,
and no Back-Cover Texts. A copy of the license is included in
the section entitled "GNU Free Documentation License".
\end{verbatim}



\section{概述}
\label{introduction-section-overview}

\noindent
除 Laumon 与 Moret-Bailly 的著作（见 \cite{LM-B}）以及 Fulton 等人尚在进行的工作之外，
我们认为，仍有必要编写一部关于代数栈及定义代数栈所需代数几何的开源教材。
Stacks Project 试图从交换代数出发奠定基础，继而经过概形与代数空间理论，
最终为代数栈理论建立一套全面的基础。

\medskip\noindent
我们设想本材料主要供在线阅读，因为它的一项关键特色是超链接：
这些链接能够迅速跳转至散布在许多不同页面上的内部引用。
如果浏览器使用嵌入式 PDF 或 DVI 查看器，跨文件链接应当可以正常工作。

\medskip\noindent
本项目由多人协作完成，我们诚邀读者参与。
阅读时如发现任何排印错误或内容错误，或有建议、补充材料或例子，
请发送电子邮件至
\href{mailto:stacks.project@gmail.com}{stacks.project@gmail.com}。
你可以下载包含全部源文件的 tar 归档包，解压后运行 make，
并在本地使用 DVI 或 PDF 查看器阅读。
也欢迎直接编辑这些 LaTeX 文件，并将改进稿通过电子邮件发给我们。


\section{贡献说明}
\label{introduction-section-attribution}

\noindent
本书涵盖范围极广；要准确而简洁地注明所有促成我们在此试图阐述之理论发展的数学家，
是一项艰巨的任务。我们希望最终能够引起足够的社群兴趣，找到愿意为每一章撰写历史评注
部分的贡献者。

\medskip\noindent
对本项目作出贡献者的姓名列于本书版本的扉页，并另有
\href{https://stacks.math.columbia.edu/tex/CONTRIBUTORS}{在线名单}。
这里列举若干主要贡献：
\begin{enumerate}
\item Jarod Alper 贡献了一章讨论代数栈相关文献的内容，见
《文献指南》第 \ref{guide-section-short-introductions} 节。
\item Bhargav Bhatt 撰写了概形的平展态射一章的初稿，见
《平展态射》第 \ref{etale-section-introduction} 节。
\item Bhargav Bhatt 撰写了
《代数进阶》第 \ref{more-algebra-section-formal-glueing} 节的初稿。
\item Kiran Kedlaya 贡献了
《下降》第 \ref{descent-section-descent-universally-injective} 节的初稿。
\item 下列内容的初稿：
\begin{enumerate}
\item 《代数》第 \ref{algebra-section-oka-families} 节，
\item 《内射对象》第 \ref{injectives-section-baer} 节，以及
\item 域这一章，见
《域》第 \ref{fields-section-introduction} 节，
\end{enumerate}
均取自
\href{https://math.uchicago.edu/~amathew/cr.html}{The CRing Project}，
由 Akhil Mathew 等人惠允使用。
\item Alex Perry 撰写了关于投射模和 Mittag--Leffler 模的材料，
其中包括
《代数》定理 \ref{algebra-theorem-ffdescent-projectivity} 的证明。
\item Alex Perry 撰写了 Schlessinger 与 Rim 意义下的形变理论一章，见
《形式形变理论》第 \ref{formal-defos-section-introduction} 节。
\item Thibaut Pugin、Zachary Maddock 与 Min Lee 为一门课程记录了讲义，
这些讲义构成了平展上同调一章与迹公式一章的基础。见
《平展上同调》第 \ref{etale-cohomology-section-introduction} 节及
《迹公式》第 \ref{trace-section-introduction} 节。
\item David Rydh 提出了许多有益意见，指出了若干错误，
在剩余 gerbe 相关材料上提供了关键帮助，并且是以下材料的最初提出者：
《空间中的群胚进阶》第
\ref{spaces-more-groupoids-section-finite} 节和
\ref{spaces-more-groupoids-section-etale-localize} 节。
\item Burt Totaro 贡献了《例》第
\ref{examples-section-non-descending-property-projective} 节、
\ref{examples-section-non-effective-descent-projective} 节，以及
《栈的性质》第 \ref{stacks-properties-section-dimension} 节。
\item pro-平展上同调一章（见
《pro-平展上同调》第 \ref{proetale-section-introduction} 节）
取自 Bhargav Bhatt 与 Peter Scholze 的一篇论文。
\item Bhargav Bhatt 贡献了《例》第
\ref{examples-section-non-algebraic-hom-stack} 节和
\ref{examples-section-flat-not-colimit-flat-finitely-presented} 节。
\item Ofer Gabber 发现了错误并提供了修正，还贡献了
《簇》引理 \ref{varieties-lemma-image-connected-component}、
《形式空间》引理 \ref{formal-spaces-lemma-completion-countably-indexed}、
《群胚进阶》第 \ref{more-groupoids-section-ind-quasi-affine} 节中的材料、
《空间的性质》第 \ref{spaces-properties-section-fpqc} 节的主要结果，
以及《平坦性进阶》命题
\ref{flat-proposition-finite-type-injective-into-flat-mod-m} 的证明。
\item J\'anos Koll\'ar 贡献了
《代数》引理 \ref{algebra-lemma-hart-serre-loc-thm} 和
《局部上同调》命题 \ref{local-cohomology-proposition-kollar}。
\item Kiran Kedlaya 撰写了
《代数进阶》第 \ref{more-algebra-section-beauville-laszlo} 节的初稿。
\item Matthew Emerton、Toby Gee 与 Brandon Levin 贡献了若干关于加厚的结果，
尤其包括《栈的态射进阶》中的引理
\ref{stacks-more-morphisms-lemma-reduced-diagonal}、
\ref{stacks-more-morphisms-lemma-thickening-diagonals} 和
\ref{stacks-more-morphisms-lemma-thickening-properties}。
\item Lena Min Ji 撰写了
《代数进阶》第 \ref{more-algebra-section-principal-radical-ideals} 节的初稿。
\item Matthew Emerton 与 Toby Gee 撰写了《栈的几何》第
\ref{stacks-geometry-section-multiplicities} 节和
\ref{stacks-geometry-section-dimension-of-algebraic-stacks} 节的初稿。
\end{enumerate}
% Generated mechanically from the frozen conventions.tex target.
% Do not edit this generated file; edit the bound locale source instead.

% OK, start here
%

\setcounter{chapter}{1}
\chapter{约定}



\phantomsection
\label{conventions-section-phantom}


\section{说明}
\label{conventions-section-comments}

\noindent
撰写这些文档时所采用约定背后的理念是：选择行得通的约定。

\section{集合论}
\label{conventions-section-sets}

\noindent
我们使用带选择公理的 Zermelo--Fraenkel 集合论。
参见 \cite{Kunen}。我们不使用集合论宇宙（这一点不同于 SGA4）。
我们不会着重讨论集合论问题，但会确保一切都（当然）严格无误，
因而也不会忽略这些问题。


\section{范畴}
\label{conventions-section-categories}

\noindent
一个范畴 $\mathcal{C}$ 由一个对象集以及每一对对象之间的一个态射集组成。
换言之，它就是其他文献中所谓的“小”范畴。
我们也会使用“大”范畴（其对象构成真类的范畴），
但仅限于《范畴》评注 \ref{categories-remark-big-categories} 中列出的那些。

\section{代数}
\label{conventions-section-algebra}

\noindent
在这些笔记中，环均指带有 $1$ 的交换环（即含幺交换环）。
因此，环的范畴以 $\mathbf{Z}$ 为始对象，
以 $\{0\}$ 为终对象（这是唯一满足 $1 = 0$ 的环）。
模均约定为含幺模。参见 \cite{Eisenbud}。

\section{记号}
\label{conventions-section-notation}

\noindent
这里的自然数是 $\mathbf{N} = \{1, 2, 3, \ldots\}$ 的元素。
整数是 $\mathbf{Z} = \{\ldots, -2, -1, 0, 1, 2, \ldots\}$ 的元素。
有理数域记作 $\mathbf{Q}$。
实数域记作 $\mathbf{R}$。
复数域记作 $\mathbf{C}$。
% Generated mechanically from the frozen sets.tex target.
% Do not edit this generated file; edit the bound locale source instead.

% OK, start here.
%

\setcounter{chapter}{2}
\chapter{集合论}



\phantomsection
\label{sets-section-phantom}


\section{引言}
\label{sets-section-introduction}

\noindent
我们不时需要用到一些集合论。我们采用 \cite{Kunen} 和 \cite{Jech}
所述的带选择公理的 Zermelo--Fraenkel 集合论（ZFC）。

\section{万物皆集合}
\label{sets-section-sets-everything}

\noindent
大多数数学家都把集合论视为数学的基本基础。那么，这在实际中究竟如何运作？
例如，怎样把“$X$ 是一个概形”这句话翻译成集合论语言？
我们只需逐层展开定义：概形是一个局部环化空间，并且每一点都有一个本身为仿射概形的开邻域。
局部环化空间是一个环化空间，其结构层的每个茎都是局部环。
环化空间是一个序对 $(X, \mathcal{O}_X)$，其中 $X$ 是拓扑空间，
$\mathcal{O}_X$ 是其上的环层。拓扑空间是一个序对 $(X, \tau)$，
其中 $X$ 是集合，$\tau \subset \mathcal{P}(X)$ 是一个由该集合的子集组成并满足
拓扑公理的集合。如此等等。

\medskip\noindent
那么，给定一个集合 $S$，我们怎样判断它是否为概形？
首先要看集合 $S$ 是否为序对。按照定义（见 \cite{Jech} 第 7 页），
这就是说存在集合 $a, b$，使得
$S$ 具有形式 $(a, b) := \{\{a\}, \{a, b\}\}$。
若是如此，再查看 $a$ 是否为某个序对 $(c, d)$。
如果是，就检验 $d \subset \mathcal{P}(c)$，并进一步检验 $d$
是否构成集合 $c$ 上的一个拓扑。如此等等。

\medskip\noindent
因此，虽然要在集合论中写出一个带单个自由变量 $x$ 的完整公式
$\phi_{scheme}(x)$ 来表达“$x$ 是概形”这一概念，需要相当多的工作，
但原则上确实可以做到。任何数学对象也都应当如此。

\section{类}
\label{sets-section-classes}

\noindent
我们非形式地使用{\it 类}这一概念。给定公式
$\phi(x, p_1, \ldots, p_n)$，我们把
$$
C = \{x : \phi(x, p_1, \ldots, p_n)\}
$$
称为一个{\it 类}。类比定义它的公式更便于操作，但严格说来，类并不是一个数学对象。
例如，若 $R$ 是环，我们可以考虑所有 $R$-模组成的类（毕竟，“$M$ 是一个
$R$-模”这句话可以翻译成集合论中的公式，而该公式便定义了一个类）。
不是集合的类称为{\it 真类}。

\medskip\noindent
这样，我们就可以考虑 $R$-模的范畴；它是一个“大”范畴，
换言之，其对象组成一个真类。类似地，我们也可以考虑概形的“大”范畴、
环的“大”范畴，等等。



\section{序数}
\label{sets-section-ordinals}

\noindent
若 $x\in T$ 蕴含 $x\subset T$，则集合 $T$ 称为{\it 传递的}。
若集合 $\alpha$ 是传递的，并且由 $\in$ 良序化，则称 $\alpha$ 为{\it 序数}。
在这种情形下，定义
$\alpha + 1 = \alpha \cup \{\alpha\}$；这仍是一个序数，称为 $\alpha$ 的{\it 后继}。
若存在序数 $\beta$ 使 $\alpha = \beta + 1$，则序数 $\alpha$
称为{\it 后继序数}。最小的序数是 $\emptyset$，也记作 $0$。
若 $\alpha$ 既不是 $0$，也不是后继序数，则称 $\alpha$ 为{\it 极限序数}，且有
$$
\alpha
=
\bigcup\nolimits_{\gamma \in \alpha} \gamma.
$$
第一个极限序数是 $\omega$，它也是第一个无限序数。
第一个不可数序数 $\omega_1$ 是所有可数序数组成的集合。
所有序数组成的族是一个真类。
它在如下意义下由 $\in$ 良序化：任意非空的序数集合（甚至序数类）都有最小元。
给定一个由序数组成的集合 $A$，定义 $A$ 的{\it 上确界}为
$\sup_{\alpha \in A} \alpha = \bigcup_{\alpha \in A} \alpha$；
它是不小于所有 $\alpha \in A$ 的最小序数。
给定任意良序集 $(S, <)$，存在唯一序数 $\alpha$ 使得
$(S, <) \cong (\alpha, \in)$；它称为该良序集的{\it 序型}。

\section{集合的层级}
\label{sets-section-sets-hierarchy}

\noindent
我们通过超限递归定义 $V_0 = \emptyset$、
$V_{\alpha + 1} = P(V_\alpha)$（幂集），并且对于极限序数 $\alpha$，定义
$$
V_\alpha = \bigcup\nolimits_{\beta < \alpha} V_\beta.
$$
注意，每个 $V_\alpha$ 都是传递集。

\begin{lemma}
\label{sets-lemma-axiom-regularity}
每个集合都属于某个序数 $\alpha$ 所对应的 $V_\alpha$。
\end{lemma}

\begin{proof}
见 \cite[Lemma 6.3]{Jech}。
\end{proof}

\noindent
\cite[Chapter III]{Kunen} 说明了该引理等价于基础公理。
集合 $S$ 的{\it 秩}是满足 $S \in V_{\alpha + 1}$ 的最小序数 $\alpha$。
所谓{\it 部分宇宙}，是指形如 $V_\alpha$、且大小适合当前语境的集合；
其具体含义将由上下文说明。

\section{基数}
\label{sets-section-cardinals}

\noindent
集合 $A$ 的{\it 势}是满足 $A$ 与 $\alpha$ 之间存在双射的最小序数 $\alpha$。
我们有时用记号 $\alpha = |A|$ 表示这一点。
序数 $\alpha$ 称为{\it 基数}，当且仅当它是某个集合 $A$ 的势；
换言之，当且仅当 $\alpha = |A|$。
我们用希腊字母 $\kappa$, $\lambda$ 表示基数。
第一个无限基数是 $\omega$；在基数的语境中，将它记作 $\aleph_0$。
若一个集合的势为 $\leq \aleph_0$，则称其为{\it 可数的}。
若 $\alpha$ 是序数，则用 $\alpha^+$ 表示满足 $> \alpha$ 的最小基数。
由此可以定义 $\aleph_1 = \aleph_0^+$、$\aleph_2 = \aleph_1^+$ 等；
事实上，对任意序数 $\alpha$，都可以通过超限递归定义 $\aleph_\alpha$。
注意有等式 $\aleph_1 = \omega_1$。

\medskip\noindent
基数 $\kappa, \lambda$ 的{\it 加法}记作 $\kappa \oplus \lambda$；
它是 $\kappa \amalg \lambda$ 的势。
基数 $\kappa, \lambda$ 的{\it 乘法}记作 $\kappa \otimes \lambda$；
它是 $\kappa \times \lambda$ 的势。
若 $\kappa$ 与 $\lambda$ 是无限基数，则
$\kappa \oplus \lambda = \kappa \otimes \lambda = \max(\kappa, \lambda)$。
基数 $\kappa, \lambda$ 的{\it 幂}记作 $\kappa^\lambda$；
它是从 $\lambda$ 到 $\kappa$ 的（集合）映射所成集合的势。
给定任意由基数组成的集合 $K$，$K$ 的{\it 上确界}为
$\sup_{\kappa \in K} \kappa = \bigcup_{\kappa \in K} \kappa$，
而这仍是一个基数。

\section{共尾性}
\label{sets-section-cofinality}

\noindent
良序集 $T$ 的一个{\it 共尾子集}是满足下述条件的子集 $S \subset T$：
$\forall t \in T \exists s\in S (t \leq s)$。
注意，良序集的子集在诱导序下仍是良序集。
给定序数 $\alpha$，其{\it 共尾性} $\text{cf}(\alpha)$
是这样一个最小序数 $\beta$：它是 $\alpha$ 的某个共尾子集的序型。
序数的共尾性总是基数。因此，也可以把 $\alpha$ 的共尾性定义为
$\alpha$ 的共尾子集所能具有的最小势。

\begin{lemma}
\label{sets-lemma-map-from-set-lifts}
设 $T = \colim_{\alpha < \beta} T_\alpha$ 是由小于给定序数 $\beta$
的序数指标的集合余极限。设 $\varphi : S \to T$ 是集合映射。
只要 $\beta$ 不是由 $S$ 指标的一族序数的极限，
即只要 $\beta$ 满足 $\text{cf}(\beta) > |S|$，
则 $\varphi$ 可提升为到某个 $T_\alpha$（其中 $\alpha < \beta$）的映射。
\end{lemma}

\begin{proof}
对每个元素 $s \in S$，选取 $\alpha_s < \beta$ 以及一个映到
$T$ 中 $\varphi(s)$ 的元素 $t_s \in T_{\alpha_s}$。
由假设，$\alpha = \sup_{s \in S} \alpha_s$ 严格小于 $\beta$。
于是，令 $\varphi_\alpha : S \to T_\alpha$ 把 $s$ 送到
$t_s$ 在 $T_\alpha$ 中的像，就得到所需解。
\end{proof}

\noindent
下面的论证实质上是 Grothendieck 关于存在任意大共尾性序数的论证；
他曾用它证明某些阿贝尔范畴中有足够多的内射对象，见 \cite{Tohoku}。

\begin{proposition}
\label{sets-proposition-exist-ordinals-large-cofinality}
设 $\kappa$ 为基数。则存在共尾性大于 $\kappa$ 的序数。
\end{proposition}

\begin{proof}
若 $\kappa$ 有限，则 $\omega = \text{cf}(\omega)$ 即可。
因此设 $\kappa$ 无限。
考虑势严格大于 $\kappa$ 的最小序数 $\alpha$。
我们声称 $\text{cf}(\alpha) > \kappa$。
注意 $\alpha$ 是极限序数：若 $\alpha = \beta + 1$，则
$|\alpha| = |\beta|$（因为 $\alpha$ 与 $\beta$ 都是无限的），
这与 $\alpha$ 的极小性矛盾。（当然，$\alpha$ 本身也是基数，但这里并不需要这一点。）
反设存在共尾子集 $S \subset \alpha$ 且 $|S| \leq \kappa$。
对 $\beta \in S$，即 $\beta < \alpha$，由 $\alpha$ 的极小性有
$|\beta| \leq \kappa$。
由于 $\alpha$ 是极限序数且 $S$ 在 $\alpha$ 中共尾，得到
$\alpha = \bigcup_{\beta \in S} \beta$。因此
$|\alpha| \leq |S| \otimes \kappa \leq \kappa \otimes \kappa \leq \kappa$，
这与 $\alpha$ 的选取矛盾。
\end{proof}



\section{反射原理}
\label{sets-section-reflection-principle}

\noindent
这部分材料有一些见于 \cite{Kunen} 中题为“Easy consistency proofs”的一章。

\medskip\noindent
设 $\phi(x_1, \ldots, x_n)$ 是集合论公式。
约定该记号表示 $\phi$ 的所有自由变量都包含在
$x_1, \ldots, x_n$ 之中。设 $M$ 为集合。
公式 $\phi^M(x_1, \ldots, x_n)$ 是把
$\phi(x_1, \ldots, x_n)$ 中所有 $\forall x$ 与 $\exists x$
分别替换为 $\forall x\in M$ 与 $\exists x\in M$ 所得到的公式。
因此，公式
$\phi(x_1, x_2) = \exists x (x\in x_1 \wedge x\in x_2)$
被变为
$\phi^M(x_1, x_2) = \exists x \in M (x\in x_1 \wedge x\in x_2)$。
公式 $\phi^M$ 称为 $\phi$ {\it 对 $M$ 的相对化}。

\begin{theorem}
\label{sets-theorem-reflection-principle}
给定集合论公式的一个{\bf 有限}族
$\phi_1(x_1, \ldots, x_n), \ldots, \phi_m(x_1, \ldots, x_n)$。
设 $M_0$ 为集合。则存在集合 $M$，使得
$M_0 \subset M$，并且
$\forall x_1, \ldots, x_n \in M$，都有
$$
\forall i = 1, \ldots, m, \ \phi_i^{M}(x_1, \ldots, x_n)
\Leftrightarrow
\forall i = 1, \ldots, m, \ \phi_i(x_1, \ldots, x_n).
$$
事实上，可以取 $M = V_\alpha$，其中 $\alpha$ 为某个极限序数。
\end{theorem}

\begin{proof}
见 \cite[Theorem 12.14]{Jech} 或 \cite[Theorem 7.4]{Kunen}。
\end{proof}

\noindent
我们把该定理解读如下：给定任意 $x_1, \ldots, x_n \in M$，
若让有界变量遍历所有集合，诸公式成立，当且仅当让这些有界变量遍历
$V_\alpha$ 的元素时诸公式成立。该定理直接处理集合论公式，因而是一个元定理。
它实际说明：给定至多以 $x_1, \ldots, x_n$ 为自由变量的有限公式表
$\phi_1, \ldots, \phi_m$，下列句子
$$
\begin{matrix}
\forall M_0\ \exists M, \ M_0 \subset M\ \forall x_1, \ldots, x_n \in M \\
\phi_1(x_1, \ldots, x_n) \wedge \ldots \wedge \phi_m(x_1, \ldots, x_n)
\leftrightarrow
\phi_1^M(x_1, \ldots, x_n) \wedge \ldots \wedge \phi_m^M(x_1, \ldots, x_n)
\end{matrix}
$$
在 ZFC 中可证。换言之，每当我们实际写下一个有限公式表 $\phi_i$ 时，
就得到一个定理。

\medskip\noindent
该定理在“普通数学”中有些难以直接使用，因为
$\phi_i^M(x_1, \ldots, x_n)$ 的含义并不十分清楚！
我们将改为利用反射原理证明中的思想，直接证明所需的存在性结果。

\section{构造概形的范畴}
\label{sets-section-categories-schemes}

\noindent
下面讨论如何运用上述思想：从一个初始概形集出发，构造一个在一系列自然运算下
封闭的“小”概形范畴。在此之前，先引入概形的大小。给定概形 $S$，定义
$$
\text{size}(S) = \max(\aleph_0, \kappa_1, \kappa_2),
$$
其中基数 $\kappa_1$ 和 $\kappa_2$ 定义如下：
\begin{enumerate}
\item 令 $\kappa_1$ 为 $S$ 的仿射开集所成集合的势。
\item 令 $\kappa_2$ 为所有 $\Gamma(U, \mathcal{O}_S)$ 的势的上确界，
其中 $U \subset S$ 遍历仿射开集。
\end{enumerate}

\begin{lemma}
\label{sets-lemma-bounded-size}
对每个基数 $\kappa$，存在集合 $A$，使得 $A$ 的每个元素都是概形，
并且对每个满足 $\text{size}(S) \leq \kappa$ 的概形 $S$，
都存在 $X \in A$ 使得 $X \cong S$（概形同构）。
\end{lemma}

\begin{proof}
略。提示：考虑任意概形如何同构于一个由仿射概形粘合而成的概形。
\end{proof}

\noindent
用 $Bound$ 表示把每个基数 $\kappa$ 映到下式的函数：
\begin{equation}
\label{sets-equation-bound}
Bound(\kappa) = \max\{\kappa^{\aleph_0}, \kappa^+\}.
\end{equation}
可以让这个函数增长得快得多，例如取
$Bound(\kappa) = \kappa^\kappa$，下面的结果仍然成立。
对任意序数 $\alpha$，用 $\Sch_\alpha$ 表示概形范畴中对象属于
$V_\alpha$ 的全子范畴。下面给出要证明的结果。

\begin{lemma}
\label{sets-lemma-construct-category}
沿用上述 $\text{size}$、$Bound$ 和 $\Sch_\alpha$ 的记号。
设 $S_0$ 是概形的集合。则存在具有下列性质的极限序数 $\alpha$：
\begin{enumerate}
\item
\label{sets-item-inclusion}
有 $S_0 \subset V_\alpha$；换言之，
$S_0 \subset \Ob(\Sch_\alpha)$。
\item
\label{sets-item-bounded}
对任意 $S \in \Ob(\Sch_\alpha)$ 以及满足
$\text{size}(T) \leq Bound(\text{size}(S))$ 的任意概形 $T$，
存在概形 $S' \in \Ob(\Sch_\alpha)$ 使得 $T \cong S'$。
\item
\label{sets-item-limit}
对任意可数\footnote{对象集与各态射集均可数。事实上，在 (3) 和 (4) 中，
可以把 $\aleph_0$ 换成任意基数，所得引理仍可证明。}图式范畴
$\mathcal{I}$ 及任意函子 $F : \mathcal{I} \to \Sch_\alpha$，
极限 $\lim_\mathcal{I} F$ 在 $\Sch_\alpha$ 中存在，当且仅当它在
$\Sch$ 中存在；而且在这种情形下，两者之间的自然态射是同构。
\item
\label{sets-item-colimit}
对任意可数指标范畴 $\mathcal{I}$ 及任意函子
$F : \mathcal{I} \to \Sch_\alpha$，
余极限 $\colim_\mathcal{I} F$ 在 $\Sch_\alpha$ 中存在，当且仅当它在
$\Sch$ 中存在；而且在这种情形下，两者之间的自然态射是同构。
\end{enumerate}
\end{lemma}

\begin{proof}
通过超限归纳定义一个把每个序数映到一个序数的函数 $f$。
令 $f(0) = 0$。给定 $f(\alpha)$，定义 $f(\alpha + 1)$ 为满足下列条件的
最小序数 $\beta$：
\begin{enumerate}
\item 有 $\alpha + 1 \leq \beta$ 且 $f(\alpha) \leq \beta$。
\item 对任意 $S \in \Ob(\Sch_{f(\alpha)})$ 以及满足
$\text{size}(T) \leq Bound(\text{size}(S))$ 的任意概形 $T$，
存在概形 $S' \in \Ob(\Sch_\beta)$ 使得 $T \cong S'$。
\item 对任意可数指标范畴 $\mathcal{I}$ 以及任意函子
$F : \mathcal{I} \to \Sch_{f(\alpha)}$，若极限
$\lim_\mathcal{I} F$ 或余极限 $\colim_\mathcal{I} F$ 在 $\Sch$ 中存在，
则它同构于 $\Sch_\beta$ 中的某个概形。
\end{enumerate}
为说明 $\beta$ 存在，论证如下。由于 $\Ob(\Sch_{f(\alpha)})$ 是集合，
可知
$\kappa =
\sup_{S \in \Ob(\Sch_{f(\alpha)})} Bound(\text{size}(S))$
存在且为基数。
依照引理 \ref{sets-lemma-bounded-size}，从 $\kappa$ 出发取得一个概形集 $A$。
存在一个由可数范畴组成的集合 $CountCat$，使每个可数范畴都同构于
$CountCat$ 的某个元素。因此在上面的 (3) 中，可以设
$\mathcal{I}$ 属于 $CountCat$。这意味着 (3) 中的序对
$(\mathcal{I}, F)$ 遍历一个集合。
于是存在一个以概形为元素的集合 $B$，使得对 (3) 中的每个
$(\mathcal{I}, F)$，若其极限或余极限存在，则该极限或余极限同构于
$B$ 中的一个元素。
因此，只要选取任意满足 $A \cup B \subset V_\beta$ 且
$\beta > \max\{\alpha + 1, f(\alpha)\}$ 的 $\beta$，条件 (1)--(3) 就成立。
由于任意非空序数族都有最小元，$f(\alpha + 1)$ 定义良好。
最后，若 $\alpha$ 为极限序数，则令
$f(\alpha) = \sup_{\alpha' < \alpha} f(\alpha')$。

\medskip\noindent
选取 $\beta_0$ 使得 $S_0 \subset V_{\beta_0}$。
由构造，$f(\beta) \geq \beta$，故也有
$S_0 \subset V_{f(\beta_0)}$。又因 $f$ 单调不减，
对任意 $\beta \geq \beta_0$ 都有 $S_0 \subset V_{f(\beta)}$。
接着选取共尾性满足
$\text{cf}(\beta_1) > \omega = \aleph_0$ 的任意序数
$\beta_1 > \beta_0$。这样的序数存在，因为序数的共尾性可以任意大；
见命题 \ref{sets-proposition-exist-ordinals-large-cofinality}。
我们声称 $\alpha = f(\beta_1)$ 是引理中问题的一个解。

\medskip\noindent
由上面对 $\beta_1 > \beta_0$ 的选取，引理的第一条性质成立。

\medskip\noindent
由于 $\beta_1$ 是极限序数（其共尾性无限），有
$f(\beta_1) = \sup_{\beta < \beta_1} f(\beta)$。
因此
$\{f(\beta) \mid \beta < \beta_1\} \subset f(\beta_1)$
是共尾子集，从而
$$
V_\alpha = V_{f(\beta_1)} = \bigcup\nolimits_{\beta < \beta_1} V_{f(\beta)}.
$$
现在令 $S \in \Ob(\Sch_\alpha)$。
定义 $\beta(S)$ 为满足
$S \in \Ob(\Sch_{f(\beta)})$ 的最小序数 $\beta$。
由上可知总有 $\beta(S) < \beta_1$。
由于
$\Ob(\Sch_{f(\beta + 1)}) \subset
\Ob(\Sch_\alpha)$，
由 $f$ 的上述构造可知引理的第二条性质成立。

\medskip\noindent
设 $\{S_1, S_2, \ldots\} \subset \Ob(\Sch_\alpha)$ 是一个可数族。
考虑函数 $\omega \to \beta_1$，$n \mapsto \beta(S_n)$。
由于 $\beta_1$ 的共尾性 $> \omega$，该函数的像不可能是共尾子集。
因此存在 $\beta < \beta_1$ 使得
$\{S_1, S_2, \ldots\} \subset \Ob(\Sch_{f(\beta)})$。
于是任意函子 $F : \mathcal{I} \to \Sch_\alpha$
都通过某个子范畴 $\Sch_{f(\beta)}$ 分解。
因此，若存在作为图式 $F$ 的余极限或极限的概形 $X$，
则由 $f$ 的构造，$X$ 同构于
$\Sch_{f(\beta + 1)}$ 中的对象，而后者是 $\Sch_\alpha$ 的子范畴。
这就证明了引理的最后两条断言。
\end{proof}

\begin{remark}
\label{sets-remark-how-to-use-reflection}
上面的引理也可以用反射原理证明，但必须谨慎。
假设句子 $\phi_{scheme}(X)$ 表达“$X$ 是概形”这一性质，那么公式
$\phi_{scheme}^{V_\alpha}(X)$ 究竟意味着什么？
反射原理确实说明可以找到 $\alpha$，使得对所有
$X \in V_\alpha$，都有
$\phi_{scheme}(X) \leftrightarrow \phi_{scheme}^{V_\alpha}(X)$，
但这本身毫无用处。只有把两个这样的陈述结合起来，才会得到有意义的结果。
例如，设 $\phi_{red}(X, Y)$ 表达“$X$、$Y$ 是概形，且 $Y$ 是
$X$ 的约化”这一性质（见《概形》定义
\ref{schemes-definition-reduced-induced-scheme}）。
对公式对
$\phi_1(X, Y) = \phi_{red}(X, Y)$、
$\phi_2(X) = \exists Y, \phi_1(X, Y)$
应用反射原理。容易看出，反射原理产生的任意 $\alpha$ 都有如下性质：
给定 $X \in \Ob(\Sch_\alpha)$，$X$ 的约化仍是 $\Sch_\alpha$ 的对象
（留作练习）。
\end{remark}

\begin{lemma}
\label{sets-lemma-bound-affine}
设 $S$ 为仿射概形，令 $R = \Gamma(S, \mathcal{O}_S)$。
则 $S$ 的大小等于 $\max\{ \aleph_0, |R|\}$。
\end{lemma}

\begin{proof}
$\Spec(R)$ 的仿射开集至多有 $\max\{|R|, \aleph_0\}$ 个。
这是显然的，因为任意仿射开集 $U \subset \Spec(R)$ 都是有限个主开集的并
$D(f_1) \cup \ldots \cup D(f_n)$，所以仿射开集的数目至多为
$\sup_n |R|^n = \max\{|R|, \aleph_0\}$；
见 \cite[Ch. I, 10.13]{Kunen}。
另一方面，
$\Gamma(U, \mathcal{O}) \subset R_{f_1} \times \ldots \times R_{f_n}$，
故
$|\Gamma(U, \mathcal{O})| \leq
\max\{\aleph_0, |R_{f_1}|, \ldots, |R_{f_n}|\}$。
因此只需证明 $|R_f| \leq \max\{\aleph_0, |R|\}$；证明略。
\end{proof}

\begin{lemma}
\label{sets-lemma-bound-size}
设 $S$ 为概形，$S = \bigcup_{i \in I} S_i$ 为一个开覆盖。则
$\text{size}(S) \leq \max\{|I|, \sup_i\{\text{size}(S_i)\}\}$。
\end{lemma}

\begin{proof}
令 $U \subset S$ 为任意仿射开集。由于 $U$ 拟紧，存在有限多个元素
$i_1, \ldots, i_n \in I$ 以及仿射开集
$U_i \subset U \cap S_i$，使得
$U = U_1 \cup U_2 \cup \ldots \cup U_n$。于是
$$
|\Gamma(U, \mathcal{O}_U)|
\leq
|\Gamma(U_1, \mathcal{O})|
\otimes
\ldots
\otimes
|\Gamma(U_n, \mathcal{O})|
\leq \sup\nolimits_i\{\text{size}(S_i)\}
$$
此外，这说明 $S$ 的仿射开集所成集合的势不超过下列集合的势：
$$
\coprod_{n \in \omega}
\coprod_{i_1, \ldots, i_n \in I}
\{\text{仿射开集，取自 }S_{i_1}\}
\times
\ldots
\times
\{\text{仿射开集，取自 }S_{i_n}\}.
$$
不交并内部每个集合的势至多为
$\sup_i\{\text{size}(S_i)\}$。
指标集的势至多为 $\max\{|I|, \aleph_0\}$；
见 \cite[Ch. I, 10.13]{Kunen}。
因此，由 \cite[Lemma 5.8]{Jech}，该余积的势至多为
$\max\{\aleph_0, |I|\}
\otimes \sup_i\{\text{size}(S_i)\}$。引理得证。
\end{proof}

\begin{lemma}
\label{sets-lemma-bound-size-fibre-product}
设 $f : X \to S$、$g : Y \to S$ 为概形态射。则
$\text{size}(X \times_S Y) \leq \max\{\text{size}(X), \text{size}(Y)\}$。
\end{lemma}

\begin{proof}
设 $S = \bigcup_{k \in K} S_k$ 为仿射开覆盖。
设 $X = \bigcup_{i \in I} U_i$、$Y = \bigcup_{j \in J} V_j$
为仿射开覆盖，并且 $I$、$J$ 的势分别
$\leq \text{size}(X), \text{size}(Y)$。
对每个 $i \in I$，存在由满足 $k \in K$ 的指标组成的有限集 $K_i$，使得
$f(U_i) \subset \bigcup_{k \in K_i} S_k$。
对每个 $j \in J$，存在由满足 $k \in K$ 的指标组成的有限集 $K_j$，使得
$g(V_j) \subset \bigcup_{k \in K_j} S_k$。
因此 $f(X), g(Y)$ 都包含在
$S' = \bigcup_{k \in K'} S_k$ 中，其中
$K' = \bigcup_{i \in I} K_i \cup \bigcup_{j \in J} K_j$。
注意 $K'$ 的势至多为 $\max\{\aleph_0, |I|, |J|\}$。
应用引理 \ref{sets-lemma-bound-size}，只需对 $k \in K'$ 证明
$\text{size}(f^{-1}(S_k) \times_{S_k} g^{-1}(S_k))
\leq \max\{\text{size}(X), \text{size}(Y)\}$。
换言之，可以设 $S$ 为仿射概形。

\medskip\noindent
设 $S$ 仿射。
令 $X = \bigcup_{i \in I} U_i$、$Y = \bigcup_{j \in J} V_j$
为仿射开覆盖，并且 $I$、$J$ 的势分别
$\leq \text{size}(X), \text{size}(Y)$。
再次由引理 \ref{sets-lemma-bound-size}，只需对积
$U_i \times_S V_j$ 证明本引理。
由引理 \ref{sets-lemma-bound-affine}，只需证明
$$
|A \otimes_C B| \leq \max\{\aleph_0, |A|, |B|\}.
$$
该不等式的证明略。
\end{proof}

\begin{lemma}
\label{sets-lemma-bound-finite-type}
设 $S$ 为概形。设 $f : X \to S$ 局部为有限型，且 $X$ 拟紧。
则 $\text{size}(X) \leq \text{size}(S)$。
\end{lemma}

\begin{proof}
可以找到有限仿射开覆盖
$X = \bigcup_{i = 1, \ldots n} U_i$，使每个 $U_i$
都映入 $S$ 的某个仿射开集 $S_i$。
因此由引理 \ref{sets-lemma-bound-size}，可约化到 $S$ 与 $X$ 都仿射的情形。
在此情形下，由引理 \ref{sets-lemma-bound-affine}，只需证明
$$
|A[x_1, \ldots, x_n]| \leq \max\{\aleph_0, |A|\}.
$$
该不等式的证明略。
\end{proof}

\noindent
在《代数》引理 \ref{algebra-lemma-epimorphism-cardinality} 中，
我们将证明：若 $A \to B$ 是环满态射，则
$|B| \leq \max(|A|, \aleph_0)$。
其概形版本是下面的引理。

\begin{lemma}
\label{sets-lemma-bound-monomorphism}
设 $f : X \to Y$ 为概形的单态射。
若下列性质中至少一项成立，则
$\text{size}(X) \leq \text{size}(Y)$：
\begin{enumerate}
\item $f$ 是拟紧的；
\item $f$ 局部为有限表现；
\item 可按需要在此增补更多条件。
\end{enumerate}
但对局部为有限型的单态射，该界一般不成立。
\end{lemma}

\begin{proof}
设 $Y = \bigcup_{j \in J} V_j$ 为 $Y$ 的仿射开覆盖，且
$|J| \leq \text{size}(Y)$。
由引理 \ref{sets-lemma-bound-size}，只需控制 $V_j$ 在 $X$ 中的逆像的大小。
因此可约化到 $Y$ 仿射的情形，设 $Y = \Spec(B)$。
对任意仿射开集 $\Spec(A) \subset X$，由上面的评注和
引理 \ref{sets-lemma-bound-affine} 有
$|A| \leq \max(|B|, \aleph_0) = \text{size}(Y)$。
因此只需证明 $X$ 的仿射开集至多有 $\text{size}(Y)$ 个。
若 $X$ 拟紧，这是显然的，故情形 (1) 成立。
在情形 (2) 中，可能出现的 $B$-代数 $A$ 的同构类数目由
$\text{size}(B)$ 控制，因为每个 $A$ 都是有限型 $B$-代数，
从而同构于某个
$B[x_1, \ldots, x_n]/(f_1, \ldots, f_m)$，
其中 $n, m$ 为某些整数，$f_j \in B[x_1, \ldots, x_n]$。
而 $X \to Y$ 是单态射，所以若存在定义在
$Y = \Spec(B)$ 上的态射 $\Spec(A) \to X$，它就是唯一的。
因此 $X$ 的仿射开集数目由上述同构类的数目控制。

\medskip\noindent
为证明引理的最后一句，考虑环
$B = \prod_{n \in \mathbf{N}} \mathbf{F}_2$，并令 $Y = \Spec(B)$。
对 $\mathbf{N}$ 上的每个超滤子 $\mathcal{U}$，得到一个剩余域为
$\mathbf{F}_2$ 的极大理想 $\mathfrak m_\mathcal{U}$；
映射 $B \to \mathbf{F}_2$ 把元素 $(x_n)$ 送到
$\lim_\mathcal{U} x_n$。细节略。
概形态射
$X = \coprod_\mathcal{U} \Spec(\mathbf{F}_2) \to Y$
是单态射，因为所有这些点彼此不同。
然而，$X$ 的仿射开子概形所成集合的势等于
$\mathbf{N}$ 上超滤子所成集合的势，即
$2^{2^{\aleph_0}}$。
由于 $|B| = 2^{\aleph_0} < 2^{2^{\aleph_0}}$，结论成立。
\end{proof}

\begin{lemma}
\label{sets-lemma-what-is-in-it}
设 $\alpha$ 为上述引理 \ref{sets-lemma-construct-category} 中的序数。
范畴 $\Sch_\alpha$ 满足下列性质：
\begin{enumerate}
\item 若 $X, Y, S \in \Ob(\Sch_\alpha)$，则对任意态射
$f : X \to S$、$g : Y \to S$，
纤维积 $X \times_S Y$ 在 $\Sch_\alpha$ 中存在，
且也是概形范畴中的纤维积。
\item 给定 $\Ob(\Sch_\alpha)$ 中任意至多可数的元素族
$S_1, S_2, \ldots$，余积 $\coprod_i S_i$ 在
$\Ob(\Sch_\alpha)$ 中存在，且也是概形范畴中的余积。
\item 对任意 $S \in \Ob(\Sch_\alpha)$ 以及任意开浸入
$U \to S$，存在 $V \in \Ob(\Sch_\alpha)$ 使 $V \cong U$。
\item 对任意 $S \in \Ob(\Sch_\alpha)$ 以及任意闭浸入
$T \to S$，存在 $S' \in \Ob(\Sch_\alpha)$ 使 $S' \cong T$。
\item 对任意 $S \in \Ob(\Sch_\alpha)$ 以及任意有限型态射
$T \to S$，存在 $S' \in \Ob(\Sch_\alpha)$ 使 $S' \cong T$。
\item 设概形 $S$ 有开覆盖
$S = \bigcup_{i \in I} S_i$，并且存在
$T \in \Ob(\Sch_\alpha)$，使得
(a) 对所有 $i \in I$ 都有
$\text{size}(S_i) \leq \text{size}(T)^{\aleph_0}$，
且 (b) $|I| \leq \text{size}(T)^{\aleph_0}$。
则 $S$ 同构于 $\Sch_\alpha$ 的一个对象。
\item 对任意 $S \in \Ob(\Sch_\alpha)$ 以及任意局部有限型态射
$f : T \to S$，若 $T$ 可由至多
$\text{size}(S)^{\aleph_0}$ 个仿射开集覆盖，
则存在 $S' \in \Ob(\Sch_\alpha)$ 使 $S' \cong T$。
例如，若 $T$ 可由至多
$|\mathbf{R}| = 2^{\aleph_0} = \aleph_0^{\aleph_0}$
个仿射开集覆盖，该结论成立。
\item 对任意 $S \in \Ob(\Sch_\alpha)$ 以及任意单态射
$T \to S$，若它局部为有限表现或拟紧，
则存在 $S' \in \Ob(\Sch_\alpha)$ 使 $S' \cong T$。
\item 设 $T \in \Ob(\Sch_\alpha)$ 为仿射概形，记
$R = \Gamma(T, \mathcal{O}_T)$。则下列任一概形都同构于
$\Sch_\alpha$ 中的某个概形：
\begin{enumerate}
\item 对任意理想 $I \subset R$，若完备化
$R^* = \lim_n R/I^n$，则概形 $\Spec(R^*)$；
\item 对任意有限型 $R$-代数 $R'$，概形 $\Spec(R')$；
\item 对任意局部化 $S^{-1}R$，概形 $\Spec(S^{-1}R)$；
\item 对任意素理想 $\mathfrak p \subset R$，概形
$\Spec(\overline{\kappa(\mathfrak p)})$；
\item 对任意子环 $R' \subset R$，概形 $\Spec(R')$；
\item 定义在势至多为 $|R|^{\aleph_0}$ 的环上的任意有限型概形；
\item 诸如此类。
\end{enumerate}
\end{enumerate}
\end{lemma}

\begin{proof}
断言 (1) 和 (2) 直接由定义得到。
断言 (3) 成立，因为 $S$ 的开子概形 $U$ 的大小显然不大于 $S$ 的大小。
断言 (4) 由 (5) 得到，断言 (5) 由 (7) 得到。
由引理 \ref{sets-lemma-bound-size}，
$S$ 的大小
$\leq \max\{|I|, \sup_i \text{size}(S_i)\}
\leq \text{size}(T)^{\aleph_0}$，故断言 (6) 成立。
断言 (7) 由 (6) 得到：对任意仿射开集 $V \subset T$，
由引理 \ref{sets-lemma-bound-finite-type} 有
$\text{size}(V) \leq \text{size}(S)$，
所以 (7) 的情形适用 (6)。
第 (8) 部分由引理 \ref{sets-lemma-bound-monomorphism} 得到。

\medskip\noindent
借助引理 \ref{sets-lemma-bound-affine}，断言 (9) 化为对环
$R^*$、$S^{-1}R$、$\overline{\kappa(\mathfrak p)}$、$R'$ 等的势给出上界。
其中最有意思的也许是环 $R^*$。
作为集合，它是满射 $R^{\mathbf{N}} \to R^*$ 的像。
由于 $|R^{\mathbf{N}}| = |R|^{\aleph_0}$，
我们选取的 $Bound(\kappa)$ 至少为 $\kappa^{\aleph_0}$，故结论成立。
总算如此！（域的代数闭包的势与该域的势相同，或者为 $\aleph_0$。）
\end{proof}

\begin{remark}
\label{sets-remark-what-is-not-in-it}
设 $R$ 为环，考虑环
$\prod_{\mathfrak p \in \Spec(R)} \kappa(\mathfrak p)$。
该环的势由 $|R|^{2^{|R|}}$ 控制，但一般不由
$|R|^{\aleph_0}$ 控制。
例如，若 $R = \mathbf{C}[x]$，则其势不由 $|R|^{\aleph_0}$ 控制；
若 $R = \prod_{n \in \mathbf{N}} \mathbf{F}_2$，
则其势不由 $|R|^{|R|}$ 控制。
因此，对上面引理 \ref{sets-lemma-what-is-in-it} 中的“诸如此类”不可不加辨别。
当然，如果在涉及 fppf/平展上同调的论证中确实需要考虑这些环，
可以把上面的函数 $Bound$ 改成
$\kappa \mapsto \kappa^{2^\kappa}$。
\end{remark}

\noindent
下面的引理使用 fpqc 覆盖的概念；该概念见
《概形上的拓扑》第 \ref{topologies-section-fpqc} 节。

\begin{lemma}
\label{sets-lemma-bound-by-covering}
设 $f : X \to Y$ 为概形态射。假设存在 fpqc 覆盖
$\{g_j : Y_j \to Y\}_{j \in J}$，使每个 $g_j$ 都通过 $f$ 分解。
则 $\text{size}(Y) \leq \text{size}(X)$。
\end{lemma}

\begin{proof}
令 $V \subset Y$ 为仿射开集。
按定义，存在 $n \geq 0$、映射
$a : \{1, \ldots, n\} \to J$ 以及仿射开集
$V_i \subset Y_{a(i)}$，使得
$V = g_{a(1)}(V_1) \cup \ldots \cup g_{a(n)}(V_n)$。
记 $h_j : Y_j \to X$ 为满足 $f \circ h_j = g_j$ 的态射。
则
$h_{a(1)}(V_1) \cup \ldots \cup h_{a(n)}(V_n)$
是 $f^{-1}(V)$ 的拟紧子集。
因此可找到拟紧开集 $W \subset f^{-1}(V)$，使其包含
$h_{a(i)}(V_i)$（$i = 1, \ldots, n$）。
特别地，$V = f(W)$。

\medskip\noindent
一方面，这说明 $Y$ 的仿射开集所成集合的势至多为
$X$ 的拟紧开集所成集合 $S$ 的势。
由于 $X$ 的每个拟紧开集都是有限个仿射开集的并，
该集合的势至多为
$\sup |S|^n = \max(\aleph_0, |S|)$。
另一方面，因为 $\{V_i \to V\}$ 是 fpqc 覆盖，有
$\mathcal{O}_Y(V) \subset
\prod_{i = 1, \ldots, n} \mathcal{O}_{Y_{a(i)}}(V_i)$。
又因 $V_i \to V$ 通过 $W$ 分解，故
$\mathcal{O}_Y(V) \subset \mathcal{O}_X(W)$。
再由 $W$ 有 $X$ 中仿射开集组成的有限覆盖，可知
$|\mathcal{O}_Y(V)|$ 由 $X$ 的大小控制。
现在由概形大小的定义即得本引理。
\end{proof}

\noindent
下面的引理使用 fppf 覆盖的概念；该概念见
《概形上的拓扑》第 \ref{topologies-section-fppf} 节。

\begin{lemma}
\label{sets-lemma-bound-fppf-covering}
设 $\{f_i : X_i \to X\}_{i \in I}$ 为一个概形的 fppf 覆盖。
则存在 fppf 覆盖 $\{W_j \to X\}_{j \in J}$，
它细化 $\{X_i \to X\}_{i \in I}$，并满足
$\text{size}(\coprod W_j) \leq \text{size}(X)$。
\end{lemma}

\begin{proof}
选取仿射开覆盖 $X = \bigcup_{a \in A} U_a$，使得
$|A| \leq \text{size}(X)$。
对每个 $a$，可以选取有限子集 $I_a \subset I$，
并对 $i \in I_a$ 选取拟紧开集
$W_{a, i} \subset X_i$，使得
$U_a = \bigcup_{i \in I_a} f_i(W_{a, i})$。
由引理 \ref{sets-lemma-bound-finite-type}，
$\text{size}(W_{a, i}) \leq \text{size}(X)$。
再由引理 \ref{sets-lemma-bound-size}，得到
$\text{size}(\coprod_a \coprod_{i \in I_a} W_{i, a})
\leq \text{size}(X)$。
\end{proof}




\section{带群作用的集合}
\label{sets-section-sets-with-group-action}

\noindent
设 $G$ 为群。用 $G\textit{-Sets}$ 表示 $G$-集合的“大”范畴。
对任意序数 $\alpha$，用 $G\textit{-Sets}_\alpha$ 表示
$G\textit{-Sets}$ 中对象属于 $V_\alpha$ 的全子范畴。
把 $G$-集合 $S$ 的大小定义为
$\text{size}(S) = \max\{\aleph_0, |G|, |S|\}$（这里 $|G|$ 和
$|S|$ 是底层集合的势）。同上，使用函数
$Bound(\kappa) = \kappa^{\aleph_0}$。

\begin{lemma}
\label{sets-lemma-sets-with-group-action}
沿用上述 $G$、$G\textit{-Sets}_\alpha$、$\text{size}$ 和
$Bound$ 的记号。设 $S_0$ 为 $G$-集合的集合。
则存在具有下列性质的极限序数 $\alpha$：
\begin{enumerate}
\item 有 $S_0 \cup \{{}_GG\} \subset \Ob(G\textit{-Sets}_\alpha)$。
\item 对任意 $S \in \Ob(G\textit{-Sets}_\alpha)$ 以及满足
$\text{size}(T) \leq Bound(\text{size}(S))$ 的任意 $G$-集合 $T$，
存在 $S' \in \Ob(G\textit{-Sets}_\alpha)$ 同构于 $T$。
\item 对任意可数指标范畴 $\mathcal{I}$ 以及任意函子
$F : \mathcal{I} \to G\textit{-Sets}_\alpha$，
极限 $\lim_\mathcal{I} F$ 与余极限
$\colim_\mathcal{I} F$ 都在 $G\textit{-Sets}_\alpha$ 中存在，
并且与 $G\textit{-Sets}$ 中的相同。
\end{enumerate}
\end{lemma}

\begin{proof}
略。证明与上述引理 \ref{sets-lemma-construct-category} 的证明相似，但更容易。
\end{proof}

\begin{lemma}
\label{sets-lemma-what-is-in-it-G-sets}
设 $\alpha$ 为上述引理 \ref{sets-lemma-sets-with-group-action} 中的序数。
范畴 $G\textit{-Sets}_\alpha$ 满足下列性质：
\begin{enumerate}
\item $G$-集合 ${}_GG$ 是 $G\textit{-Sets}_\alpha$ 的对象。
\item （余）积、纤维积和推出在 $G\textit{-Sets}_\alpha$ 中存在，
并且与 $G\textit{-Sets}$ 中相应的对象相同。
\item 给定 $G\textit{-Sets}_\alpha$ 的对象 $U$，
任意 $G$-稳定子集 $O \subset U$ 都同构于
$G\textit{-Sets}_\alpha$ 的某个对象。
\end{enumerate}
\end{lemma}

\begin{proof}
略。
\end{proof}

\section{位点的覆盖}
\label{sets-section-coverings-site}

\noindent
设 $\mathcal{C}$ 是一个范畴（如《范畴》定义
\ref{categories-definition-category} 中），并设
$\text{Cov}(\mathcal{C})$ 是由覆盖组成的真类，
满足《位点与层》定义 \ref{sites-definition-site} 中的性质
(1)、(2)、(3)。现将它们列出：
\begin{enumerate}
\item 若 $V \to U$ 为同构，则
$\{V \to U\} \in \text{Cov}(\mathcal{C})$。
\item 若
$\{U_i \to U\}_{i\in I} \in \text{Cov}(\mathcal{C})$，
并且对每个 $i$ 都有
$\{V_{ij} \to U_i\}_{j\in J_i} \in \text{Cov}(\mathcal{C})$，
则
$\{V_{ij} \to U\}_{i \in I, j\in J_i}
\in \text{Cov}(\mathcal{C})$。
\item 若
$\{U_i \to U\}_{i\in I}\in \text{Cov}(\mathcal{C})$，
且 $V \to U$ 是 $\mathcal{C}$ 的态射，则对所有 $i$，
$U_i \times_U V$ 都存在，并且
$\{U_i \times_U V \to V \}_{i\in I}
\in \text{Cov}(\mathcal{C})$。
\end{enumerate}
对序数 $\alpha$，令
$\text{Cov}(\mathcal{C})_\alpha =
\text{Cov}(\mathcal{C}) \cap V_\alpha$。
给定序数 $\alpha$ 和基数 $\kappa$，令
$\text{Cov}(\mathcal{C})_{\kappa, \alpha}$ 为下列元素组成的集合：
$\mathcal{U} =
\{\varphi_i : U_i \to U\}_{i\in I}
\in \text{Cov}(\mathcal{C})_\alpha$
且满足 $|I| \leq \kappa$。

\medskip\noindent
回顾下面的概念；见《位点与层》定义
\ref{sites-definition-combinatorial-tautological}。
设两族态射
$\{\varphi_i : U_i \to U\}_{i\in I}$ 与
$\{\psi_j : W_j \to U\}_{j\in J}$
在 $\mathcal{C}$ 中有同一靶。
若存在映射 $\alpha : I \to J$ 与 $\beta : J\to I$，
使得 $\varphi_i = \psi_{\alpha(i)}$ 且
$\psi_j = \varphi_{\beta(j)}$，则称这两族态射{\it 组合等价}。
这在具有固定靶的态射族上定义了一个等价关系。

\begin{lemma}
\label{sets-lemma-coverings-site}
沿用上述记号。设
$\text{Cov}_0 \subset \text{Cov}(\mathcal{C})$
是包含于 $\text{Cov}(\mathcal{C})$ 的一个集合。
则存在基数 $\kappa$ 和极限序数 $\alpha$，满足：
\begin{enumerate}
\item 有
$\text{Cov}_0 \subset
\text{Cov}(\mathcal{C})_{\kappa, \alpha}$。
\item 覆盖集
$\text{Cov}(\mathcal{C})_{\kappa, \alpha}$
满足《位点与层》定义 \ref{sites-definition-site} 中的
(1)、(2)、(3)（见上文）。换言之，
$(\mathcal{C}, \text{Cov}(\mathcal{C})_{\kappa, \alpha})$
是一个位点。
\item $\text{Cov}(\mathcal{C})$ 中的每个覆盖都组合等价于
$\text{Cov}(\mathcal{C})_{\kappa, \alpha}$ 中的某个覆盖。
\end{enumerate}
\end{lemma}

\begin{proof}
首先考虑具有固定靶的 $\mathcal{C}$-态射集合所成集合
$\mathcal{S}$。换言之，$\mathcal{S}$ 的元素是
$\text{Arrows}(\mathcal{C})$ 的子集 $T$，并且 $T$ 的所有元素有同一靶。
给定具有固定靶的态射族
$\mathcal{U} = \{\varphi_i : U_i \to U\}_{i\in I}$，定义
$Supp(\mathcal{U}) = \{ \varphi \in \text{Arrows}(\mathcal{C})
\mid \exists i\in I, \varphi = \varphi_i\}$。
注意，两族
$\mathcal{U} =  \{\varphi_i : U_i \to U\}_{i\in I}$ 与
$\mathcal{V} = \{V_j \to V\}_{j \in J}$
组合等价，当且仅当
$Supp(\mathcal{U}) = Supp(\mathcal{V})$。
接着定义子集
$\mathcal{S}_\tau \subset \mathcal{S}$：
$\mathcal{S}_\tau = \{ T \in \mathcal{S} \mid
\exists\ \mathcal{U} \in \text{Cov}(\mathcal{C}) \
T = Supp(\mathcal{U})\}$。
对每个 $T \in \mathcal{S}_\tau$，令 $\beta(T)$ 为满足下列条件的
最小序数 $\beta$：存在
$\mathcal{U} \in \text{Cov}(\mathcal{C})_\beta$，
使得 $T = \text{Supp}(\mathcal{U})$。
最后，令
$\beta_0 = \sup_{T \in S_\tau} \beta(T)$。
由此可知，每个 $\mathcal{U} \in \text{Cov}(\mathcal{C})$
都组合等价于
$\text{Cov}(\mathcal{C})_{\beta_0}$ 中的某个元素。

\medskip\noindent
令 $\kappa$ 为下列各基数的最大值：
$\aleph_0$、$|\text{Arrows}(\mathcal{C})|$，以及
$$
\sup\nolimits_{\{U_i \to U\}_{i\in I}
\in \text{Cov}(\mathcal{C})_{\beta_0}}
|I|,
\quad\text{与}\quad
\sup\nolimits_{\{U_i \to U\}_{i\in I} \in \text{Cov}_0} |I|.
$$
由于 $\kappa$ 是无限基数，有
$\kappa \otimes \kappa = \kappa$。
显然
$\text{Cov}(\mathcal{C})_{\beta_0} =
\text{Cov}(\mathcal{C})_{\kappa, \beta_0}$。

\medskip\noindent
通过超限归纳定义一个把每个序数映到序数的函数 $f$。
令 $f(0) = 0$。给定 $f(\alpha)$，定义 $f(\alpha + 1)$
为满足下列条件的最小序数 $\beta$：
\begin{enumerate}
\item 有 $\alpha + 1 \leq \beta$ 且
$f(\alpha) \leq \beta$。
\item 若
$\{U_i \to U\}_{i\in I}
\in \text{Cov}(\mathcal{C})_{\kappa, f(\alpha)}$，
并且对每个 $i$ 都有
$\{W_{ij} \to U_i\}_{j\in J_i}
\in \text{Cov}(\mathcal{C})_{\kappa, f(\alpha)}$，
则
$\{W_{ij} \to U\}_{i \in I, j\in J_i}
\in \text{Cov}(\mathcal{C})_{\kappa, \beta}$。
\item 若
$\{U_i \to U\}_{i\in I}
\in \text{Cov}(\mathcal{C})_{\kappa, \alpha}$，
且 $W \to U$ 是 $\mathcal{C}$ 的态射，则
$\{U_i \times_U W \to W \}_{i\in I}
\in \text{Cov}(\mathcal{C})_{\kappa, \beta}$。
\end{enumerate}
为说明 $\beta$ 存在，只需注意 (2) 和 (3) 中出现的所有覆盖
$\{W_{ij} \to U\}$ 与
$\{U_i \times_U W \to W \}$ 显然组成一个集合。
因此存在某个序数 $\beta$，使 $V_\beta$ 包含所有这些覆盖。
此外，覆盖 $\{W_{ij} \to U\}$ 的指标集的势满足
$\sum_{i \in I} |J_i|
\leq \kappa \otimes \kappa = \kappa$，
故这些覆盖属于
$\text{Cov}(\mathcal{C})_{\kappa, \beta}$。
由于任意非空序数族都有最小元，
$f(\alpha + 1)$ 定义良好。
最后，若 $\alpha$ 是极限序数，则令
$f(\alpha) = \sup_{\alpha' < \alpha} f(\alpha')$。

\medskip\noindent
选取序数 $\beta_1$，使得
$\text{Arrows}(\mathcal{C}) \subset V_{\beta_1}$、
$\text{Cov}_0 \subset V_{\beta_0}$，并且
$\beta_1 \geq \beta_0$。
由构造，$f(\beta_1) \geq \beta_1$，可知对
$V_{f(\beta_1)}$ 仍有相同性质。
又因 $f$ 单调不减，对任意 $\beta \geq \beta_1$ 仍然如此。
接着选取共尾性满足
$\text{cf}(\beta_2) > \kappa$ 的任意序数
$\beta_2 > \beta_1$。
这样的序数存在，因为序数的共尾性可以任意大；
见命题 \ref{sets-proposition-exist-ordinals-large-cofinality}。
我们声称序对 $\kappa$、
$\alpha = f(\beta_2)$ 是引理中问题的一个解。

\medskip\noindent
由上述对 $\kappa$、$\beta_2 > \beta_1 > \beta_0$ 的选取，
引理的第一条和第三条性质成立。

\medskip\noindent
由于 $\beta_2$ 是极限序数（其共尾性无限），有
$f(\beta_2) = \sup_{\beta < \beta_2} f(\beta)$。
因此
$\{f(\beta) \mid \beta < \beta_2\} \subset f(\beta_2)$
是共尾子集，从而
$$
V_\alpha = V_{f(\beta_2)}
= \bigcup\nolimits_{\beta < \beta_2} V_{f(\beta)}.
$$
现在令
$\mathcal{U} \in \text{Cov}_{\kappa, \alpha}$。
定义 $\beta(\mathcal{U})$ 为满足
$\mathcal{U} \in \text{Cov}_{\kappa, f(\beta)}$
的最小序数 $\beta$。由上可知总有
$\beta(\mathcal{U}) < \beta_2$。

\medskip\noindent
需要证明定义位点的性质 (1)、(2)、(3) 对序对
$(\mathcal{C}, \text{Cov}_{\kappa, \alpha})$ 成立。
第一条成立，因为按照对 $\beta_2$ 的选取，
$\mathcal{C}$ 的所有态射都属于 $V_{f(\beta_2)}$。
为证第三条，给定覆盖
$\mathcal{U} = \{U_i \to U\}_{i \in I}
\in \text{Cov}(\mathcal{C})_{\kappa, \alpha}$，
有 $\beta(\mathcal{U}) < \beta_2$；
因此由 $f$ 的构造，$\mathcal{U}$ 的任意基变换都属于
$\text{Cov}(\mathcal{C})_{\kappa, f(\beta + 1)}$，
进而属于
$\text{Cov}(\mathcal{C})_{\kappa, \alpha}$。

\medskip\noindent
最后证明第二条。设
$\{U_i \to U\}_{i\in I}
\in \text{Cov}(\mathcal{C})_{\kappa, f(\alpha)}$，
并且对每个 $i$ 都有
$\mathcal{W}_i = \{W_{ij} \to U_i\}_{j\in J_i}
\in \text{Cov}(\mathcal{C})_{\kappa, f(\alpha)}$。
考虑函数
$I \to \beta_2$，$i \mapsto \beta(\mathcal{W}_i)$。
由于 $\beta_2$ 的共尾性
$> \kappa \geq |I|$，该函数的像不可能是共尾子集。
因此存在 $\beta < \beta_1$，使对所有 $i \in I$ 都有
$\mathcal{W}_i \in \text{Cov}_{\kappa, f(\beta)}$。
由此，覆盖
$\{W_{ij} \to U\}_{i\in I, j \in J_i}$
属于
$\text{Cov}(\mathcal{C})_{\kappa, f(\beta + 1)}
\subset \text{Cov}(\mathcal{C})_{\kappa, \alpha}$，
如所要求。
\end{proof}

\begin{remark}
\label{sets-remark-better}
很可能存在某个极限序数 $\alpha$，使覆盖集
$\text{Cov}(\mathcal{C})_\alpha$ 满足本引理的条件。
毕竟，应用反射原理似乎正会给出这一结论
（但须顾及第 \ref{sets-section-reflection-principle} 节末尾以及
评注 \ref{sets-remark-how-to-use-reflection} 中所述的注意事项）。
\end{remark}

\section{阿贝尔范畴与内射对象}
\label{sets-section-abelian-categories-injectives}

\noindent
下面的引理适用于位点上的环层之模所成的范畴。

\begin{lemma}
\label{sets-lemma-abelian-injectives}
设给定一个大范畴 $\mathcal{A}$（见《范畴》评注
\ref{categories-remark-big-categories}）。
假设 $\mathcal{A}$ 是阿贝尔范畴，并且有足够多的内射对象。
见《同调代数》定义
\ref{homology-definition-abelian-category} 和
\ref{homology-definition-enough-injectives}。
则对 $\mathcal{A}$ 中任意给定的对象集
$\{A_s\}_{s\in S}$，存在阿贝尔子范畴
$\mathcal{A}' \subset \mathcal{A}$，满足：
\begin{enumerate}
\item 包含函子 $\mathcal{A}' \to \mathcal{A}$ 是正合的；
\item $\Ob(\mathcal{A}')$ 是集合；
\item 对每个 $s \in S$，$\Ob(\mathcal{A}')$ 都包含 $A_s$；
\item $\mathcal{A}'$ 有足够多的内射对象；
\item $\mathcal{A}'$ 的对象是内射对象，当且仅当它是
$\mathcal{A}$ 的内射对象。
\end{enumerate}
\end{lemma}

\begin{proof}
略。
\end{proof}
% Generated mechanically from the frozen categories.tex target.
% Do not edit this generated file; edit the bound locale source instead.


% OK, start here.
%

\setcounter{chapter}{3}
\chapter{范畴}



\phantomsection
\label{categories-section-phantom}


\section{引言}
\label{categories-section-introduction}

\noindent
范畴最早在 \cite{GenEqui} 中引入。
范畴所成的范畴（它是一个真类）是一个 $2$-范畴。
类似地，栈所成的范畴也构成一个 $2$-范畴。
如果你已经熟悉范畴，但尚不熟悉 $2$-范畴，
则应把第 \ref{categories-section-formal-cat-cat} 节作为后文形式定义的导论来读。

\section{定义}
\label{categories-section-definition-categories}

\noindent
我们回顾这些定义，部分目的是固定记号。

\begin{definition}
\label{categories-definition-category}
一个{\it 范畴} $\mathcal{C}$ 由下列数据组成：
\begin{enumerate}
\item 对象集 $\Ob(\mathcal{C})$。
\item 对每一对 $x, y \in \Ob(\mathcal{C})$，给定态射集
$\Mor_\mathcal{C}(x, y)$。
\item 对每个三元组 $x, y, z\in \Ob(\mathcal{C})$，给定复合映射
$ \Mor_\mathcal{C}(y, z) \times \Mor_\mathcal{C}(x, y)
\to \Mor_\mathcal{C}(x, z) $，记作 $(\phi, \psi) \mapsto
\phi \circ \psi$。
\end{enumerate}
这些数据须满足下列规则：
\begin{enumerate}
\item 对每个 $x\in \Ob(\mathcal{C})$，存在态射
$\text{id}_x\in \Mor_\mathcal{C}(x, x)$，使得只要相应复合有定义，
就有 $\text{id}_x \circ \phi = \phi$ 和
$\psi \circ \text{id}_x = \psi $。
\item 复合满足结合律；也就是说，只要相应复合有定义，就有
$(\phi \circ \psi) \circ \chi =
\phi \circ ( \psi \circ \chi)$。
\end{enumerate}
\end{definition}

\noindent
通常要求所有态射集 $\Mor_\mathcal{C}(x, y)$ 两两不交。
这样，态射 $\phi : x \to y$ 就有唯一的{\it 源} $x$
和唯一的{\it 靶} $y$。严格说来，这一要求并非必需；
但若不作此要求，在表述上述条件 (2) 时就必须格外谨慎。
这一约定很方便，我们往往假定它成立。
在这种情形下，若 $\phi$ 的源等于 $\psi$ 的靶，
则称 $\phi$ 与 $\psi$ {\it 可复合}，此时 $\phi \circ \psi$ 有定义。
一个等价的定义是把范畴定义为五元组
$(\text{Ob}, \text{Arrows}, s, t, \circ)$：
它由对象集、态射（箭头）集、源、靶和复合组成，并满足一长列公理。
我们偶尔会采用这一观点。

\begin{remark}
\label{categories-remark-big-categories}
大范畴。在有些文献中，范畴可以有一个真类的对象。
本书也允许这种情形，但仅限于下面列出的情况
（该表将随内容推进而更新）。
特别地，当我们说“设 $\mathcal{C}$ 为范畴”时，
默认 $\Ob(\mathcal{C})$ 是集合。
\begin{enumerate}
\item 集合的范畴 $\textit{Sets}$。
\item 阿贝尔群的范畴 $\textit{Ab}$。
\item 群的范畴 $\textit{Groups}$。
\item 给定群 $G$，带左 $G$-作用的集合所成范畴
$G\textit{-Sets}$。
\item 给定环 $R$，$R$-模的范畴 $\text{Mod}_R$。
\item 给定域 $k$，$k$ 上向量空间的范畴。
\item 环的范畴。
\item 除幂环的范畴，见《除幂代数》第
\ref{dpa-section-divided-power-rings} 节。
\item 概形的范畴。
\item 拓扑空间的范畴 $\textit{Top}$。
\item 给定拓扑空间 $X$，$X$ 上集合预层的范畴
$\textit{PSh}(X)$。
\item 给定拓扑空间 $X$，$X$ 上集合层的范畴
$\Sh(X)$。
\item 给定拓扑空间 $X$，$X$ 上阿贝尔群预层的范畴
$\textit{PAb}(X)$。
\item 给定拓扑空间 $X$，$X$ 上阿贝尔群层的范畴
$\textit{Ab}(X)$。
\item 给定小范畴 $\mathcal{C}$，从 $\mathcal{C}$ 到
$\textit{Sets}$ 的函子所成范畴。
\item 给定范畴 $\mathcal{C}$，$\mathcal{C}$ 上集合预层的范畴。
\item 给定位点 $\mathcal{C}$，$\mathcal{C}$ 上集合层的范畴。
\end{enumerate}
在这里逐项列出它们，原因之一是尽量避免处理类似
“大”范畴的“集合体”这样的东西；那相当于处理所有类的集合体，
而在我们看来，这无疑是元数学对象。
\end{remark}

\begin{remark}
\label{categories-remark-unique-identity}
由定义立即可知，范畴 $\mathcal{A}$ 的对象 $x$ 的任意两个恒等态射都相同。
因此，我们可以并且将会谈论 $x$ 的{\it 这个}恒等态射
$\text{id}_x$。
\end{remark}

\begin{definition}
\label{categories-definition-isomorphism}
若范畴 $\mathcal{C}$ 中的态射 $\phi : x \to y$ 满足：
存在态射 $\psi : y \to x$，使得
$\phi \circ \psi = \text{id}_y$ 且
$\psi \circ \phi = \text{id}_x$，
则称该态射为一个{\it 同构}。
\end{definition}

\noindent
同构 $\phi$ 有时也称为{\it 可逆}态射；定义中的态射 $\psi$
称为其{\it 逆}，记作 $\phi^{-1}$。只要存在，逆就是唯一的。
注意，给定范畴 $\mathcal{A}$ 的对象 $x$，
$\Mor_\mathcal{A}(x, x)$ 中的可逆元素集
$\text{Aut}_\mathcal{A}(x)$ 在复合下构成一个群。
这个群称为 $x$ 在 $\mathcal{A}$ 中的{\it 自同构}群。

\begin{definition}
\label{categories-definition-groupoid}
若一个范畴的每个态射都是同构，则称它为{\it 群胚}。
\end{definition}

\begin{example}
\label{categories-example-group-groupoid}
群 $G$ 给出一个只有单个对象 $x$ 的群胚，
其态射满足 $\Mor(x, x) = G$，复合规则由 $G$ 中的群运算给出。
每个只有单个对象的群胚都具有这种形式。
\end{example}

\begin{example}
\label{categories-example-set-groupoid}
集合 $C$ 给出如下定义的群胚 $\mathcal{C}$：
取 $\Ob(\mathcal{C}) := C$ 为对象集；对态射，
若 $x\neq y$，则令 $\Mor(x, y)$ 为空集，
若 $x = y$，则令它等于 $\{\text{id}_x\}$。
\end{example}

\begin{definition}
\label{categories-definition-functor}
两个范畴 $\mathcal{A}, \mathcal{B}$ 之间的{\it 函子}
$F : \mathcal{A} \to \mathcal{B}$ 由下列数据给出：
\begin{enumerate}
\item 映射 $F : \Ob(\mathcal{A}) \to \Ob(\mathcal{B})$。
\item 对每个 $x, y \in \Ob(\mathcal{A})$，给定映射
$F : \Mor_\mathcal{A}(x, y) \to \Mor_\mathcal{B}(F(x), F(y))$，
记作 $\phi \mapsto F(\phi)$。
\end{enumerate}
这些数据须按如下方式与复合及恒等态射相容：
对 $\mathcal{A}$ 中任意可复合态射对 $(\phi, \psi)$，有
$F(\phi \circ \psi) = F(\phi) \circ F(\psi)$；
并且 $F(\text{id}_x) = \text{id}_{F(x)}$。
\end{definition}

\noindent
注意，每个范畴 $\mathcal{A}$ 都有{\it 恒等}函子
$\text{id}_\mathcal{A}$。
此外，给定函子 $G : \mathcal{B} \to \mathcal{C}$ 和
$F : \mathcal{A} \to \mathcal{B}$，有一个按显然方式定义的
{\it 复合}函子 $G \circ F : \mathcal{A} \to \mathcal{C}$。

\begin{definition}
\label{categories-definition-faithful}
设 $F : \mathcal{A} \to \mathcal{B}$ 为函子。
\begin{enumerate}
\item 若对任意对象 $x, y \in \Ob(\mathcal{A})$，映射
$$
F : \Mor_\mathcal{A}(x, y) \to \Mor_\mathcal{B}(F(x), F(y))
$$
都是单射，则称 $F$ {\it 忠实}。
\item 若这些映射全都是双射，则称 $F$ {\it 全忠实}。
\item
若对任意对象 $y \in \Ob(\mathcal{B})$，都存在对象
$x \in \Ob(\mathcal{A})$，使得 $F(x)$ 在 $\mathcal{B}$ 中
同构于 $y$，则称函子 $F$ {\it 本质满射}。
\end{enumerate}
\end{definition}

\begin{definition}
\label{categories-definition-subcategory}
范畴 $\mathcal{B}$ 的一个{\it 子范畴}是范畴 $\mathcal{A}$，
其对象和箭头分别组成 $\mathcal{B}$ 的对象和箭头的子集，
并且 $\mathcal{A}$ 中的源、靶和复合与 $\mathcal{B}$ 中的相应运算一致，
而 $\mathcal{A}$ 中任一对象的恒等态射也与它在 $\mathcal{B}$ 中的恒等态射相同。
若对所有 $x, y \in \Ob(\mathcal{A})$ 都有
$\Mor_\mathcal{A}(x, y) = \Mor_\mathcal{B}(x, y)$，
则称 $\mathcal{A}$ 为 $\mathcal{B}$ 的{\it 满子范畴}。
若它是满子范畴，并且给定 $x \in \Ob(\mathcal{A})$ 时，
$\mathcal{B}$ 中每个与 $x$ 同构的对象也属于 $\mathcal{A}$，
则称 $\mathcal{A}$ 为 $\mathcal{B}$ 的{\it 严格满}子范畴。
\end{definition}

\noindent
若 $\mathcal{A} \subset \mathcal{B}$ 是子范畴，
则恒等映射给出从 $\mathcal{A}$ 到 $\mathcal{B}$ 的函子。
此外，子范畴 $\mathcal{A} \subset \mathcal{B}$ 是满的，
当且仅当包含函子是全忠实的。
注意，给定范畴 $\mathcal{B}$，$\mathcal{B}$ 的满子范畴所成集合
与 $\Ob(\mathcal{B})$ 的子集所成集合相同。

\begin{remark}
\label{categories-remark-functor-into-sets}
设 $\mathcal{A}$ 为范畴。
从 $\mathcal{A}$ 到 $\textit{Sets}$ 的函子 $F$ 是数学对象
（也就是说，它是集合，而不是类或集合论公式；见
《集合论》第 \ref{sets-section-sets-everything} 节），
尽管集合范畴是“大”的。
事实上，$F$ 在对象上的像 $F(\Ob(\mathcal{A}))$ 是集合；
于是，可以把 $F$ 看作 $\mathcal{A}$ 与集合范畴的一个满子范畴之间的函子，
其中该满子范畴的对象正是 $F(\Ob(\mathcal{A}))$ 的元素。
\end{remark}

\begin{example}
\label{categories-example-group-homomorphism-functor}
群同态 $p : G\to H$ 在例 \ref{categories-example-group-groupoid} 所述的相应群胚之间
给出一个函子。它忠实当且仅当 $p$ 为单射；
它全忠实当且仅当该同态为同构。
\end{example}

\begin{example}
\label{categories-example-category-over-X}
给定范畴 $\mathcal{C}$ 和对象 $X\in \Ob(\mathcal{C})$，
如下定义{\it $X$ 上对象的范畴}，记作 $\mathcal{C}/X$。
$\mathcal{C}/X$ 的对象是某个 $Y\in \Ob(\mathcal{C})$ 所给出的态射
$Y\to X$。对象 $Y\to X$ 与 $Y'\to X$ 之间的态射，
是在 $\mathcal{C}$ 中使显然图表交换的态射 $Y\to Y'$。
注意，有一个函子 $p_X : \mathcal{C}/X\to \mathcal{C}$，
它只是忘掉所给的态射。此外，给定 $\mathcal{C}$ 中的态射
$f : X'\to X$，与 $f$ 复合诱导函子
$F : \mathcal{C}/X' \to \mathcal{C}/X$，并且
$p_X\circ F = p_{X'}$。
\end{example}

\begin{example}
\label{categories-example-category-under-X}
给定范畴 $\mathcal{C}$ 和对象 $X\in \Ob(\mathcal{C})$，
如下定义{\it $X$ 下对象的范畴}，记作 $X/\mathcal{C}$。
$X/\mathcal{C}$ 的对象是某个 $Y\in \Ob(\mathcal{C})$ 所给出的态射
$X\to Y$。对象 $X\to Y$ 与 $X\to Y'$ 之间的态射，
是在 $\mathcal{C}$ 中使显然图表交换的态射 $Y\to Y'$。
注意，有一个函子 $p_X : X/\mathcal{C}\to \mathcal{C}$，
它只是忘掉所给的态射。此外，给定 $\mathcal{C}$ 中的态射
$f : X'\to X$，与 $f$ 复合诱导函子
$F : X/\mathcal{C} \to X'/\mathcal{C}$，并且
$p_{X'}\circ F = p_X$。
\end{example}




\begin{definition}
\label{categories-definition-transformation-functors}
设 $F, G : \mathcal{A} \to \mathcal{B}$ 为函子。
一个{\it 自然变换}，或一个{\it 函子态射}
$t : F \to G$，是一个族 $\{t_x\}_{x\in \Ob(\mathcal{A})}$，满足：
\begin{enumerate}
\item $t_x : F(x) \to G(x)$ 是范畴 $\mathcal{B}$ 中的态射；
\item 对 $\mathcal{A}$ 中每个态射 $\phi : x \to y$，下图交换：
$$
\xymatrix{
F(x) \ar[r]^{t_x} \ar[d]_{F(\phi)} & G(x) \ar[d]^{G(\phi)} \\
F(y) \ar[r]^{t_y} & G(y) }
$$
\end{enumerate}
\end{definition}

\noindent
有时我们用图表
$$
\xymatrix{
\mathcal{A}
\rtwocell^F_G{t}
&
\mathcal{B}
}
$$
表示 $t$ 是从 $F$ 到 $G$ 的态射。

\medskip\noindent
注意，每个函子 $F$ 都有{\it 恒等}变换
$\text{id}_F : F \to F$。
此外，给定函子态射 $t : F \to G$ 与函子态射 $s : E \to F$，
它们的{\it 复合} $t \circ s$ 由规则
$$
(t \circ s)_x = t_x \circ s_x : E(x) \to G(x)
$$
定义，其中 $x \in \Ob(\mathcal{A})$。
容易验证，这确实是从 $E$ 到 $G$ 的函子态射。
这样，给定范畴 $\mathcal{A}$ 与 $\mathcal{B}$，
就得到一个新范畴，即 $\mathcal{A}$ 与 $\mathcal{B}$ 之间的函子所成范畴。

\begin{remark}
\label{categories-remark-functors-sets-sets}
这是当 $\mathcal{A}$ 为“大”范畴时，同样结论不成立的一种情形。
例如，考虑函子 $\textit{Sets} \to \textit{Sets}$。
按照目前的定义，这样的函子是类而不是集合。
换言之，它由一个集合论公式（其中某些变量等于指定集合）给出！
试图把集合论的所有可能公式都纳入数学对象的定义，并不是好主意。
例如，考虑概形范畴上的层时，也会出现同样的问题。
我们稍后将回到这一点。
\end{remark}

\begin{definition}
\label{categories-definition-equivalence-categories}
一个{\it 范畴等价}
$F : \mathcal{A} \to \mathcal{B}$ 是满足下列条件的函子：
存在函子 $G : \mathcal{B} \to \mathcal{A}$，使复合
$F \circ G$ 与 $G \circ F$ 分别同构于恒等函子
$\text{id}_\mathcal{B}$ 与 $\text{id}_\mathcal{A}$。
在这种情形下，称 $G$ 为 $F$ 的一个{\it 拟逆}。
\end{definition}

\begin{lemma}
\label{categories-lemma-construct-quasi-inverse}
设 $F : \mathcal{A} \to \mathcal{B}$ 为全忠实函子。
假设对每个 $X \in \Ob(\mathcal{B})$，都给定
$\mathcal{A}$ 的对象 $j(X)$ 以及同构 $i_X : X \to F(j(X))$。
则存在唯一函子 $j : \mathcal{B} \to \mathcal{A}$，
$j$ 在对象上延拓上述规则，并且这些同构 $i_X$ 定义函子同构
$\text{id}_\mathcal{B} \to F \circ j$。
此外，$j$ 与 $F$ 是互为拟逆的范畴等价。
\end{lemma}

\begin{proof}
构造 $j : \mathcal{B} \to \mathcal{A}$ 分两步进行。
首先，定义映射 $j : \Ob(\mathcal{B}) \to \Ob(\mathcal{A})$，
它把 $X \in \mathcal{B}$ 对应到 $j(X)$。
其次，若 $X,Y \in \Ob(\mathcal{B})$ 且 $\phi : X \to Y$，
则考虑 $\phi' := i_Y \circ\phi\circ i_X^{-1}$。
由于 $F$ 全忠实，存在唯一的 $\varphi$ 满足 $F(\varphi) = \phi'$。
定义 $j(\phi) = \varphi$。我们略去 $j$ 是函子的验证。
由构造，图表
$$
\xymatrix{
X \ar[r]_-{i_X} \ar[d]_{\phi}
&
F(j(X)) \ar[d]^{F\circ j(\phi)}
\\
Y \ar[r]^-{i_Y}
&
F(j(Y))
}
$$
交换。因此，由于每个 $i_X$ 都是同构，
$\{i_X\}_X$ 是函子同构 $\text{id}_\mathcal{B} \to F\circ j$。
最后，还须证明 $j \circ F$ 同构于 $\text{id}_\mathcal{A}$。
然而，由于 $F$ 全忠实，要证明这一点，只需在后面与 $F$ 复合后证明；
也就是说，只需证明 $F \circ j \circ F$ 同构于
$F \circ \text{id}_\mathcal{A}$（略去一个小细节）。
由 $F \circ j \cong \text{id}_\mathcal{B}$，这显然成立。
\end{proof}

\begin{lemma}
\label{categories-lemma-equivalence-categories}
一个函子是范畴等价，当且仅当它既全忠实又本质满射。
\end{lemma}

\begin{proof}
设 $F : \mathcal{A} \to \mathcal{B}$ 本质满射且全忠实。
按照约定，所有范畴都是小的；又因 $F$ 本质满射，
利用选择公理，可以对每个 $X \in \Ob(\mathcal{B})$ 选取
$\mathcal{A}$ 的对象 $j(X)$ 以及同构 $i_X : X \to F(j(X))$。
于是，利用 $F$ 全忠实，应用引理 \ref{categories-lemma-construct-quasi-inverse}。
\end{proof}

\begin{definition}
\label{categories-definition-product-category}
设 $\mathcal{A}$、$\mathcal{B}$ 为范畴。
定义{\it 积范畴} $\mathcal{A} \times \mathcal{B}$：
其对象集为
$\Ob(\mathcal{A} \times \mathcal{B}) =
\Ob(\mathcal{A}) \times \Ob(\mathcal{B})$，并且
$$
\Mor_{\mathcal{A} \times \mathcal{B}}((x, y), (x', y'))
:=
\Mor_\mathcal{A}(x, x')\times
\Mor_\mathcal{B}(y, y').
$$
复合按分量定义。
\end{definition}

\section{对偶范畴与 Yoneda 引理}
\label{categories-section-opposite}

\begin{definition}
\label{categories-definition-opposite}
给定范畴 $\mathcal{C}$，其{\it 对偶范畴}
$\mathcal{C}^{opp}$ 与 $\mathcal{C}$ 有相同的对象，
但所有态射的方向都反转。
\end{definition}

\noindent
换言之，
$\Mor_{\mathcal{C}^{opp}}(x, y) = \Mor_\mathcal{C}(y, x)$。
$\mathcal{C}^{opp}$ 中的复合与 $\mathcal{C}$ 中相同，只是方向相反：
若 $\phi : y \to z$ 和 $\psi : x \to y$ 是
$\mathcal{C}^{opp}$ 中的态射，亦即 $\mathcal{C}$ 中的箭头
$z \to y$ 和 $y \to x$，则 $\phi \circ^{opp} \psi$ 是
$\mathcal{C}^{opp}$ 中的态射 $x \to z$，它对应于
$\mathcal{C}$ 中的复合 $z \to y \to x$。

\begin{definition}
\label{categories-definition-contravariant}
设 $\mathcal{C}$、$\mathcal{S}$ 为范畴。
从 $\mathcal{C}$ 到 $\mathcal{S}$ 的{\it 反变}函子 $F$
是函子 $\mathcal{C}^{opp}\to \mathcal{S}$。
\end{definition}

\noindent
具体地，反变函子 $F$ 由映射
$F : \Ob(\mathcal{C}) \to \Ob(\mathcal{S})$ 以及下列数据给出：
对 $\mathcal{C}$ 中每个态射 $\psi : x \to y$，
给定态射 $F(\psi) : F(y) \to F(x)$。
这些数据须满足：再给定态射 $\phi : y \to z$ 时，
作为态射 $F(z) \to F(x)$，有
$F(\phi \circ \psi) = F(\psi) \circ F(\phi)$。
（注意次序反转了。）

\begin{definition}
\label{categories-definition-presheaf}
设 $\mathcal{C}$ 为范畴。
\begin{enumerate}
\item $\mathcal{C}$ 上的一个{\it 集合预层}，简称{\it 预层}，
是从 $\mathcal{C}$ 到 $\textit{Sets}$ 的反变函子 $F$。
\item 预层范畴记作 $\textit{PSh}(\mathcal{C})$。
\end{enumerate}
\end{definition}

\noindent
当然，预层范畴是一个真类。

\begin{example}
\label{categories-example-hom-functor}
点函子。
对任意 $U\in \Ob(\mathcal{C})$，有反变函子
$$
\begin{matrix}
h_U & : & \mathcal{C}
&
\longrightarrow
&
\textit{Sets} \\
& &
X
&
\longmapsto
&
\Mor_\mathcal{C}(X, U)
\end{matrix}
$$
它把对象 $X$ 送到集合 $\Mor_\mathcal{C}(X, U)$。
换言之，$h_U$ 是一个预层。
给定态射 $f : X\to Y$，相应映射
$h_U(f) : \Mor_\mathcal{C}(Y, U)\to \Mor_\mathcal{C}(X, U)$
把 $\phi$ 送到 $\phi\circ f$。
我们总把这个预层记作 $h_U : \mathcal{C}^{opp} \to \textit{Sets}$。
它称为与 $U$ 相伴的{\it 可表预层}。
如果 $\mathcal{C}$ 是概形范畴，这个函子有时称为 $U$ 的
\emph{点函子}。
\end{example}

\noindent
注意，给定 $\mathcal{C}$ 中的态射 $\phi : U \to V$，
与态射 $U \to V$ 复合给出相应的函子自然变换
$h(\phi) : h_U \to h_V$。
这把 $\mathcal{C}$ 中态射的复合变成函子变换的复合。
换言之，得到函子
$$
h :
\mathcal{C}
\longrightarrow
\textit{PSh}(\mathcal{C})
$$
注意，靶是一个“大”范畴，见注 \ref{categories-remark-big-categories}。
另一方面，$h$ 是一个真正的数学对象（即一个集合），
参见注 \ref{categories-remark-functor-into-sets}。

\begin{lemma}[Yoneda 引理]
\label{categories-lemma-yoneda}
\begin{reference}
它以某种形式出现于 \cite{Yoneda-homology}；Grothendieck 在
\cite{Gr-II} 中使用了其推广形式。
\end{reference}
设 $U, V \in \Ob(\mathcal{C})$。
给定任意函子态射 $s : h_U \to h_V$，
存在唯一态射 $\phi : U \to V$，使得 $h(\phi) = s$。
换言之，函子 $h$ 是全忠实的。
更一般地，给定任意反变函子 $F$ 以及 $\mathcal{C}$ 的任意对象
$U$，有自然双射
$$
\Mor_{\textit{PSh}(\mathcal{C})}(h_U, F) \longrightarrow F(U),
\quad
s \longmapsto s_U(\text{id}_U).
$$
\end{lemma}

\begin{proof}
对第一个断言，只需取
$\phi = s_U(\text{id}_U) \in \Mor_\mathcal{C}(U, V)$。
对第二个断言，给定 $\xi \in F(U)$，如下定义 $s$：
$s_V : h_U(V) \to F(V)$ 把
$h_U(V) = \Mor_\mathcal{C}(V, U)$ 的元素 $f : V \to U$
送到 $F(f)(\xi)$。
\end{proof}

\begin{definition}
\label{categories-definition-representable-functor}
若反变函子 $F : \mathcal{C}\to \textit{Sets}$
同构于 $\mathcal{C}$ 的某个对象 $U$ 的点函子 $h_U$，
则称该函子{\it 可表}。
\end{definition}

\noindent
设 $\mathcal{C}$ 为范畴，并设
$F : \mathcal{C}^{opp} \to \textit{Sets}$ 为可表函子。
选取 $\mathcal{C}$ 的对象 $U$ 以及同构 $s : h_U \to F$。
Yoneda 引理保证序对 $(U, s)$ 在唯一同构意义下是唯一的。
对象 $U$ 称为{\it 表示} $F$ 的对象。
由 Yoneda 引理，变换 $s$ 对应于唯一元素 $\xi \in F(U)$。
该元素称为{\it 泛对象}。它具有下列性质：
对 $V \in \Ob(\mathcal{C})$，映射
$$
\Mor_\mathcal{C}(V, U) \longrightarrow F(V),\quad
(f : V \to U) \longmapsto F(f)(\xi)
$$
是双射。因此，$\xi$ 在如下意义下是泛的：
$F(V)$ 的每个元素都等于 $\xi$ 在 $F(f)$ 下的像，
其中 $f : V \to U$ 是 $\mathcal{C}$ 中唯一的态射。






\section{对象对的乘积}
\label{categories-section-products-pairs}

\begin{definition}
\label{categories-definition-products}
设 $x, y\in \Ob(\mathcal{C})$。
$x$ 与 $y$ 的一个{\it 乘积}是对象
$x \times y \in \Ob(\mathcal{C})$，连同态射
$p\in \Mor_{\mathcal C}(x \times y, x)$ 和
$q\in\Mor_{\mathcal C}(x \times y, y)$，
它们满足下列泛性质：对任意 $w\in \Ob(\mathcal{C})$
以及态射 $\alpha \in \Mor_{\mathcal C}(w, x)$ 和
$\beta \in \Mor_\mathcal{C}(w, y)$，
存在唯一的 $\gamma\in \Mor_{\mathcal C}(w, x \times y)$，
使图表
$$
\xymatrix{
w \ar[rrrd]^\beta \ar@{-->}[rrd]_\gamma \ar[rrdd]_\alpha & & \\
& & x \times y \ar[d]_p \ar[r]_q & y \\
& & x &
}
$$
交换。
\end{definition}

\noindent
若乘积存在，则它在唯一同构意义下是唯一的。
这由 Yoneda 引理得到，因为定义要求 $x \times y$ 是
$\mathcal{C}$ 的对象，并且函子性地依赖于 $w$ 满足
$$
h_{x \times y}(w) = h_x(w) \times h_y(w)
$$
换言之，乘积 $x \times y$ 是表示函子
$w \mapsto h_x(w) \times h_y(w)$ 的对象。

\begin{definition}
\label{categories-definition-has-products-of-pairs}
若对任意 $x, y \in \Ob(\mathcal{C})$，乘积 $x \times y$
都存在，则称范畴 $\mathcal{C}$ {\it 具有对象对的乘积}。
\end{definition}

\noindent
我们采用这一术语，是为了把该概念与“具有乘积”或“具有有限乘积”
区别开来；后二者通常另有所指（特别是，它们总蕴含终对象存在）。





\section{对象对的余积}
\label{categories-section-coproducts-pairs}

\begin{definition}
\label{categories-definition-coproducts}
设 $x, y \in \Ob(\mathcal{C})$。
$x$ 与 $y$ 的一个{\it 余积}，或{\it 并合和}，
是对象 $x \amalg y \in \Ob(\mathcal{C})$，连同态射
$i \in \Mor_{\mathcal C}(x, x \amalg y)$ 和
$j \in \Mor_{\mathcal C}(y, x \amalg y)$，
它们满足下列泛性质：对任意 $w \in \Ob(\mathcal{C})$
以及态射 $\alpha \in \Mor_{\mathcal C}(x, w)$ 和
$\beta \in \Mor_\mathcal{C}(y, w)$，
存在唯一的 $\gamma \in \Mor_{\mathcal C}(x \amalg y, w)$，
使图表
$$
\xymatrix{
& y \ar[d]^j \ar[rrdd]^\beta \\
x \ar[r]^i \ar[rrrd]_\alpha & x \amalg y \ar@{-->}[rrd]^\gamma \\
& & & w
}
$$
交换。
\end{definition}

\noindent
若余积存在，则它在唯一同构意义下是唯一的。
这由 Yoneda 引理（应用于对偶范畴）得到，因为定义要求
$x \amalg y$ 是 $\mathcal{C}$ 的对象，并且函子性地依赖于 $w$ 满足
$$
\Mor_\mathcal{C}(x \amalg y, w) =
\Mor_\mathcal{C}(x, w) \times \Mor_\mathcal{C}(y, w)
$$

\begin{definition}
\label{categories-definition-has-coproducts-of-pairs}
若对任意 $x, y \in \Ob(\mathcal{C})$，余积 $x \amalg y$
都存在，则称范畴 $\mathcal{C}$ {\it 具有对象对的余积}。
\end{definition}

\noindent
我们采用这一术语，是为了把该概念与“具有余积”或“具有有限余积”
区别开来；后二者通常另有所指（特别是，它们总蕴含
$\mathcal{C}$ 中存在始对象）。





\section{纤维积}
\label{categories-section-fibre-products}

\begin{definition}
\label{categories-definition-fibre-products}
设 $x, y, z\in \Ob(\mathcal{C})$，
$f\in \Mor_\mathcal{C}(x, y)$，并且
$g\in \Mor_{\mathcal C}(z, y)$。
$f$ 与 $g$ 的一个{\it 纤维积}是对象
$x \times_y z\in \Ob(\mathcal{C})$，连同态射
$p \in \Mor_{\mathcal C}(x \times_y z, x)$ 和
$q \in \Mor_{\mathcal C}(x \times_y z, z)$，使图表
$$
\xymatrix{
x \times_y z \ar[r]_q \ar[d]_p & z \ar[d]^g \\
x \ar[r]^f & y
}
$$
交换，并满足下列泛性质：对任意 $w\in \Ob(\mathcal{C})$
以及满足 $f \circ \alpha = g \circ \beta$ 的态射
$\alpha \in \Mor_{\mathcal C}(w, x)$ 和
$\beta \in \Mor_\mathcal{C}(w, z)$，
存在唯一的 $\gamma \in \Mor_{\mathcal C}(w, x \times_y z)$，
使图表
$$
\xymatrix{
w \ar[rrrd]^\beta \ar@{-->}[rrd]_\gamma \ar[rrdd]_\alpha & & \\
& & x \times_y z \ar[d]^p \ar[r]_q & z \ar[d]^g \\
& & x \ar[r]^f & y
}
$$
交换。
\end{definition}

\noindent
若纤维积存在，则它在唯一同构意义下是唯一的。
这由 Yoneda 引理得到，因为定义要求 $x \times_y z$ 是
$\mathcal{C}$ 的对象，并且函子性地依赖于 $w$ 满足
$$
h_{x \times_y z}(w) = h_x(w) \times_{h_y(w)} h_z(w)
$$
换言之，纤维积 $x \times_y z$ 是表示函子
$w \mapsto h_x(w) \times_{h_y(w)} h_z(w)$ 的对象。

\begin{definition}
\label{categories-definition-cartesian}
若范畴中的交换图表
$$
\xymatrix{
w \ar[r] \ar[d] &
z \ar[d] \\
x \ar[r] &
y
}
$$
满足：$w$ 以及态射 $w \to x$ 和 $w \to z$
构成态射 $x \to y$ 与 $z \to y$ 的纤维积，
则称该图表是{\it 笛卡尔的}。
\end{definition}

\begin{definition}
\label{categories-definition-has-fibre-products}
若对任意 $f\in \Mor_{\mathcal C}(x, y)$ 和
$g\in \Mor_{\mathcal C}(z, y)$，纤维积都存在，
则称范畴 $\mathcal{C}$ {\it 具有纤维积}。
\end{definition}

\begin{definition}
\label{categories-definition-representable-morphism}
若范畴 $\mathcal{C}$ 中的态射 $f : x \to y$ 满足：
对 $\mathcal{C}$ 中每个态射 $z \to y$，
纤维积 $x \times_y z$ 都存在，则称该态射{\it 可表}。
\end{definition}

\begin{lemma}
\label{categories-lemma-composition-representable}
设 $\mathcal{C}$ 为范畴。
设 $f : x \to y$ 和 $g : y \to z$ 可表。
则 $g \circ f : x \to z$ 可表。
\end{lemma}

\begin{proof}
设 $t \in \Ob(\mathcal C)$ 且
$ \varphi \in \Mor_{\mathcal C}(t,z) $。
由于 $g$ 与 $f$ 可表，得到交换图表
$$
\xymatrix{
y \times_z t \ar[r]_q \ar[d]_p &
t \ar[d]^{\varphi} \\
y \ar[r]^{g} & z
}
\quad\quad
\xymatrix{
x \times_y (y\! \times_z\! t) \ar[r]_{q'} \ar[d]_{p'} &
y \times_z t \ar[d]^p \\
x \ar[r]^f & y
}
$$
它们具有定义 \ref{categories-definition-fibre-products} 中的泛性质。
我们声称，令 $x \times_z t = x \times_y (y \times_z t)$，
并取态射 $q \circ q' : x \times_z t \to t$ 和
$p' : x \times_z t \to x$，就得到一个纤维积。
首先，由上述图表的交换性可知
$\varphi \circ q \circ q' = g \circ f \circ p'$。
为验证泛性质，设 $w \in \Ob(\mathcal C)$，
并设态射 $\alpha : w \to x$ 和 $\beta : w \to t$ 满足
$\varphi \circ \beta = g \circ f \circ \alpha$。
由纤维积的定义，存在唯一态射 $\delta$ 与 $\gamma$，使图表
$$
\xymatrix{
w \ar[rrrd]^\beta \ar@{-->}[rrd]_\delta \ar[rrdd]_{f\circ\alpha} & & \\
& & y \times_z t \ar[d]_p \ar[r]_q & t \ar[d]^{\varphi} \\
& & y \ar[r]^{g} & z
}
$$
以及
$$
\xymatrix{
w \ar[rrrd]^\delta \ar@{-->}[rrd]_\gamma \ar[rrdd]_{\alpha} & & \\
& & x \times_y (y\!\times_z\! t) \ar[d]_{p'} \ar[r]_{q'} &
y \times_z t \ar[d]^{p} \\
& & x \ar[r]^{f} & y
}
$$
交换。于是，$\gamma$ 使图表
$$
\xymatrix{
w \ar[rrrd]^\beta \ar@{-->}[rrd]_\gamma \ar[rrdd]_{\alpha} & & \\
& & x \times_z t \ar[d]_{p'} \ar[r]_{q\circ q'} & t \ar[d]^{\varphi} \\
& & x \ar[r]^{g\circ f} & z
}
$$
交换。为证明它的唯一性，设 $\gamma'$ 满足
$ q\circ q'\circ\gamma' = \beta $ 和
$ p'\circ \gamma' = \alpha $。
由于 $\gamma$ 唯一，只需证明
$ q'\circ\gamma' = \delta $ 和
$ p'\circ\gamma' = \alpha $ 即可。第二个等式正是我们的假设。
对第一个等式，还须使用上述态射的唯一性。
注意，$\delta$ 是满足 $ q\circ\delta = \beta $ 和
$ p\circ\delta = f\circ\alpha $ 的唯一态射。
我们已经假设 $ q\circ (q'\circ\gamma') = \beta $。
此外，由纤维积的定义，知道 $ f\circ p' = p\circ q' $。
因此：
$$
p\circ (q'\circ\gamma') =
(p\circ q')\circ\gamma' =
(f\circ p')\circ\gamma' =
f\circ (p'\circ\gamma') = f\circ\alpha.
$$
所以 $ q'\circ\gamma' = \delta $，证毕。
\end{proof}

\begin{lemma}
\label{categories-lemma-base-change-representable}
设 $\mathcal{C}$ 为范畴。
设 $f : x \to y$ 可表，并设 $y' \to y$ 是
$\mathcal{C}$ 中的态射。
则态射 $x' := x \times_y y' \to y'$ 也可表。
\end{lemma}

\begin{proof}
设 $z \to y'$ 为态射。纤维积 $x' \times_{y'} z$
应当表示函子
\begin{eqnarray*}
w & \mapsto & h_{x'}(w)\times_{h_{y'}(w)} h_z(w) \\
& = & (h_x(w) \times_{h_y(w)} h_{y'}(w)) \times_{h_{y'}(w)} h_z(w) \\
& = & h_x(w) \times_{h_y(w)} h_z(w)
\end{eqnarray*}
而后者由假设可表。
\end{proof}

\section{纤维积的例}
\label{categories-section-example-fibre-products}

\noindent
本节列举并描述一些纤维积的例子。

\medskip\noindent
首先是一个十分平凡的例子：集合范畴具有纤维积，
因而每个态射都可表。具体地，若 $f : X \to Y$ 和
$g : Z \to Y$ 是集合映射，则把 $X \times_Y Z$ 定义为
$X \times Z$ 中所有满足 $f(x) = g(z)$ 的序对 $(x, z)$ 所成的子集。
态射 $p : X \times_Y Z \to X$ 与 $q : X \times_Y Z \to Z$
是投影映射 $(x, z) \mapsto x$ 和 $(x, z) \mapsto z$。
最后，若态射 $\alpha : W \to X$ 与 $\beta : W \to Z$
满足 $f \circ \alpha = g \circ \beta$，则映射
$W \to X \times Z$，$w\mapsto (\alpha(w), \beta(w))$
显然取值于 $X \times_Y Z$，正合所需。

\medskip\noindent
许多范畴的对象是在集合上赋予某类代数结构；在这些范畴中，
底层集合的纤维积也给出该范畴中的纤维积。
例如，假设上述 $X$、$Y$ 与 $Z$ 是群，而 $f$、$g$ 是群同态。
那么，集合论纤维积 $X \times_Y Z$ 继承群结构；
只需用公式 $(x, z) \cdot (x', z') = (xx', zz')$
定义两个序对的乘积。下列范畴都可使用类似论证。
\begin{enumerate}
\item 群的范畴 $\textit{Groups}$。
\item 给定群 $G$ 时，带左 $G$-作用的集合所成范畴
$G\textit{-Sets}$。
\item 环的范畴。
\item 给定环 $R$ 时，$R$-模的范畴。
\end{enumerate}




\section{纤维积与可表性}
\label{categories-section-representable-map-presheaves}

\noindent
本节具体研究从一个范畴到集合范畴的反变函子所成范畴中的纤维积。
以后讨论位点、预层与层时，将用更一般的内容取代这里的讨论。
这类函子之间的“可表态射”概念颇有意义。

\begin{lemma}
\label{categories-lemma-fibre-product-presheaves}
设 $\mathcal{C}$ 为范畴。
设 $F, G, H : \mathcal{C}^{opp} \to \textit{Sets}$ 为函子，
并设 $a : F \to G$ 与 $b : H \to G$ 为函子变换。
则预层范畴 $\textit{PSh}(\mathcal{C})$ 中的纤维积
$F \times_{a, G, b} H$ 存在，而且对 $\mathcal{C}$ 的任意对象
$X$，它由公式
$$
(F \times_{a, G, b} H)(X) =
F(X) \times_{a_X, G(X), b_X} H(X)
$$
给出。
\end{lemma}

\begin{proof}
略。
\end{proof}

\noindent
作为特例，设有态射 $a : F \to G$、对象
$U \in \Ob(\mathcal{C})$ 以及元素 $\xi \in G(U)$。
根据 Yoneda 引理 \ref{categories-lemma-yoneda}，这给出变换
$\xi : h_U \to G$。此时的纤维积由规则
$$
(h_U \times_{\xi, G, a} F)(X) =
\{ (f, \xi') \mid f : X \to U, \ \xi' \in F(X), \ G(f)(\xi) = a_X(\xi')\}
$$
描述。若 $F$、$G$ 也可表，则只要纤维积存在，这就是表示它的函子；
见第 \ref{categories-section-fibre-products} 节。
与定义 \ref{categories-definition-representable-morphism} 的类比促使我们
定义可表变换的概念。

\begin{definition}
\label{categories-definition-representable-map-presheaves}
设 $\mathcal{C}$ 为范畴，并设
$F, G : \mathcal{C}^{opp} \to \textit{Sets}$ 为函子。
若对每个 $U \in \Ob(\mathcal{C})$ 及任意 $\xi \in G(U)$，
函子 $h_U \times_{\xi, G, a} F$ 都可表，
则称态射 $a : F \to G$ {\it 可表}，
或称{\it $F$ 在 $G$ 上相对可表}。
\end{definition}

\begin{lemma}
\label{categories-lemma-representable-over-representable}
设 $\mathcal{C}$ 为范畴。
设 $a : F \to G$ 是从 $\mathcal{C}$ 到 $\textit{Sets}$
的反变函子之间的态射。若 $a$ 可表，并且 $G$ 是可表函子，
则 $F$ 可表。
\end{lemma}

\begin{proof}
略。
\end{proof}

\begin{lemma}
\label{categories-lemma-representable-diagonal}
设 $\mathcal{C}$ 为范畴，并设
$F : \mathcal{C}^{opp} \to \textit{Sets}$ 为函子。
假设 $\mathcal{C}$ 具有对象对的乘积和纤维积。
下列条件等价：
\begin{enumerate}
\item 对角态射 $\Delta : F \to F \times F$ 可表；
\item 对 $\mathcal{C}$ 中每个 $U$ 及任意 $\xi \in F(U)$，
映射 $\xi : h_U \to F$ 可表；
\item 对 $\mathcal{C}$ 中每一对 $U, V$ 及任意
$\xi \in F(U)$、$\xi' \in F(V)$，纤维积
$h_U \times_{\xi, F, \xi'} h_V$ 可表。
\end{enumerate}
\end{lemma}

\begin{proof}
我们继续利用 Yoneda 引理，把 $F(U)$ 与从 $h_U \to F$
的函子变换等同起来。

\medskip\noindent
(2) 与 (3) 的等价性。设 $U, \xi, V, \xi'$ 如 (3) 中所述。
(2) 与 (3) 都恰好断言 $h_U \times_{\xi, F, \xi'} h_V$
可表；唯一差别是 (3) 关于 $U$ 与 $V$ 对称，而 (2) 并不对称。

\medskip\noindent
假设条件 (1) 成立。设 $U, \xi, V, \xi'$ 如 (3) 中所述。
注意，$h_U \times h_V = h_{U \times V}$ 可表。
用 $\eta : h_{U \times V} \to F \times F$ 表示对应于乘积
$\xi \times \xi' : h_U \times h_V \to F \times F$ 的映射。
由假设，纤维积
$F \times_{\Delta, F \times F, \eta} h_{U \times V}$ 可表。
这意味着存在 $W \in \Ob(\mathcal{C})$、态射
$W \to U$、$W \to V$ 以及 $h_W \to F$，使得
$$
\xymatrix{
h_W \ar[d] \ar[r] & h_U \times h_V \ar[d]^{\xi \times \xi'} \\
F \ar[r] & F \times F
}
$$
是笛卡尔图表。使用引理 \ref{categories-lemma-fibre-product-presheaves}
中纤维积的显式描述，读者可见这蕴含
$h_W = h_U \times_{\xi, F, \xi'} h_V$，正合所需。

\medskip\noindent
假设等价的条件 (2) 和 (3) 成立。
设 $U$ 为 $\mathcal{C}$ 的对象，并设
$(\xi, \xi') \in (F \times F)(U)$。
由 (3)，纤维积 $h_U \times_{\xi, F, \xi'} h_U$ 可表。
选取对象 $W$ 以及同构
$h_W \to h_U \times_{\xi, F, \xi'} h_U$。
两个投影 $\text{pr}_i : h_U \times_{\xi, F, \xi'} h_U \to h_U$
由 Yoneda 对应于态射 $p_i : W \to U$。
考虑 $W' = W \times_{(p_1, p_2), U \times U} U$。
因为
$$
h_{W'} = h_W \times_{h_U \times h_U} h_U
= (h_U \times_{\xi, F, \xi'} h_U) \times_{h_U \times h_U} h_U
= F \times_{F \times F} h_U.
$$
形式地可知 $W'$ 表示 $F \times_{\Delta, F \times F} h_U$。
因此 $\Delta$ 可表，证明完成。
\end{proof}









\section{推出}
\label{categories-section-pushouts}

\noindent
与纤维积对偶的概念是推出。

\begin{definition}
\label{categories-definition-pushouts}
设 $x, y, z\in \Ob(\mathcal{C})$，
$f\in \Mor_\mathcal{C}(y, x)$ 且
$g\in \Mor_{\mathcal C}(y, z)$。
$f$ 与 $g$ 的一个 {\it 推出}是对象
$x\amalg_y z\in \Ob(\mathcal{C})$，连同态射
$p\in \Mor_{\mathcal C}(x, x\amalg_y z)$ 和
$q\in\Mor_{\mathcal C}(z, x\amalg_y z)$，使图表
$$
\xymatrix{
y \ar[r]_g \ar[d]_f & z \ar[d]^q \\
x \ar[r]^p & x\amalg_y z
}
$$
交换，并满足下列泛性质：对任意 $w\in \Ob(\mathcal{C})$
及满足 $\alpha \circ f = \beta \circ g$ 的态射
$\alpha \in \Mor_{\mathcal C}(x, w)$ 与
$\beta \in \Mor_\mathcal{C}(z, w)$，存在唯一的
$\gamma\in \Mor_{\mathcal C}(x\amalg_y z, w)$，使图表
$$
\xymatrix{
y \ar[r]_g \ar[d]_f & z \ar[d]^q \ar[rrdd]^\beta & & \\
x \ar[r]^p \ar[rrrd]^\alpha & x \amalg_y z \ar@{-->}[rrd]^\gamma & & \\
& & & w
}
$$
交换。
\end{definition}

\noindent
若三元组 $(x\amalg_y z, p, q)$ 存在，可直接且容易地证明它
在唯一同构意义下是唯一的。另一种做法是把推出看作对偶范畴中的
纤维积，于是第 \ref{categories-section-fibre-products} 节的讨论立即给出这种唯一性。

\begin{definition}
\label{categories-definition-cocartesian}
若范畴中的交换图表
$$
\xymatrix{
y \ar[r] \ar[d] & z \ar[d] \\
x \ar[r] & w
}
$$
满足：$w$ 以及态射 $x \to w$、$z \to w$ 构成态射
$y \to x$ 与 $y \to z$ 的一个推出，则称该图表{\it 余笛卡尔}。
\end{definition}


\section{等化子}
\label{categories-section-equalizers}

\begin{definition}
\label{categories-definition-equalizers}
设 $X$、$Y$ 为范畴 $\mathcal{C}$ 的对象，
$a, b : X \to Y$ 为态射。若 $a \circ e = b \circ e$，且
$(Z, e)$ 满足下列泛性质，则称态射 $e : Z \to X$
为态射对 $(a, b)$ 的一个{\it 等化子}：对 $\mathcal{C}$ 中每个满足
$a \circ t = b \circ t$ 的态射 $t : W \to X$，存在唯一态射
$s : W \to Z$，使得 $t = e \circ s$。
\end{definition}

\noindent
与上面的纤维积情形一样，等化子一旦存在，便在唯一同构意义下唯一。
把此定义推广到多于 $2$ 个态射的情形也是直接的。

\section{余等化子}
\label{categories-section-coequalizers}

\begin{definition}
\label{categories-definition-coequalizers}
设 $X$、$Y$ 为范畴 $\mathcal{C}$ 的对象，
$a, b : X \to Y$ 为态射。若 $c \circ a = c \circ b$，且
$(Z, c)$ 满足下列泛性质，则称态射 $c : Y \to Z$
为态射对 $(a, b)$ 的一个{\it 余等化子}：对 $\mathcal{C}$ 中每个满足
$t \circ a = t \circ b$ 的态射 $t : Y \to W$，存在唯一态射
$s : Z \to W$，使得 $t = s \circ c$。
\end{definition}

\noindent
与上面的推出情形一样，余等化子一旦存在，便在唯一同构意义下唯一；
考虑对偶范畴，即可由等化子的唯一性推出这一点。
把此定义推广到多于 $2$ 个态射的情形也是直接的。

\section{始对象与终对象}
\label{categories-section-initial-final}

\begin{definition}
\label{categories-definition-initial-final}
设 $\mathcal{C}$ 为范畴。
\begin{enumerate}
\item 若对 $\mathcal{C}$ 的每个对象 $y$，恰有一个态射 $x \to y$，
则称 $\mathcal{C}$ 的对象 $x$ 为{\it 始对象}。
\item 若对 $\mathcal{C}$ 的每个对象 $y$，恰有一个态射 $y \to x$，
则称 $\mathcal{C}$ 的对象 $x$ 为{\it 终对象}。
\end{enumerate}
\end{definition}

\noindent
在集合范畴中，空集 $\emptyset$ 是始对象，而且事实上是唯一的始对象。
此外，任意{\it 单点集}，即只有一个元素的集合，都是终对象
（所以终对象并不唯一）。




\section{单态射与满态射}
\label{categories-section-mono-epi}

\begin{definition}
\label{categories-definition-mono-epi}
设 $\mathcal{C}$ 为范畴，$f : X \to Y$ 为 $\mathcal{C}$ 的态射。
\begin{enumerate}
\item 若对每个对象 $W$ 及每一对满足
$f \circ a = f \circ b$ 的态射 $a, b : W \to X$，均有
$a = b$，则称 $f$ 为{\it 单态射}。
\item 若对每个对象 $W$ 及每一对满足
$a \circ f = b \circ f$ 的态射 $a, b : Y \to W$，均有
$a = b$，则称 $f$ 为{\it 满态射}。
\end{enumerate}
\end{definition}

\begin{example}
\label{categories-example-mono-epi-sets}
在集合范畴中，单态射对应于单射映射，满态射对应于满射映射。
\end{example}

\begin{lemma}
\label{categories-lemma-characterize-mono-epi}
设 $\mathcal{C}$ 为范畴，$f : X \to Y$ 为 $\mathcal{C}$ 的态射。则
\begin{enumerate}
\item $f$ 是单态射，当且仅当 $X$ 是纤维积 $X \times_Y X$；
\item $f$ 是满态射，当且仅当 $Y$ 是推出 $Y \amalg_X Y$。
\end{enumerate}
\end{lemma}

\begin{proof}
假设 $f$ 为单态射。设 $W$ 为 $\mathcal C$ 的对象，并设
$\alpha, \beta \in \Mor_\mathcal C(W,X)$ 满足 $f\circ \alpha = f\circ
\beta$。由于 $f$ 是单态射，故 $\alpha = \beta$。此外，有交换图表
$$
\xymatrix{
X \ar[r]^{\text{id}_X} \ar[d]_{\text{id}_X} & X \ar[d]^{f} \\
X \ar[r]^{f} & Y
}
$$
它以 $\gamma := \alpha = \beta$ 满足泛性质。
因此 $X$ 的确是纤维积 $X\times_Y X$。

\medskip\noindent
假设 $X \times_Y X \cong X $。图表
$$
\xymatrix{
X \ar[r]^{\text{id}_X} \ar[d]_{\text{id}_X} & X \ar[d]^{f} \\
X \ar[r]^{f} & Y }
$$
交换；若 $W \in \Ob(\mathcal C)$ 且 $\alpha, \beta : W \to X$
满足 $f \circ \alpha = f \circ \beta$，则存在唯一的
$\gamma$ 满足
$$
\gamma = \text{id}_X\circ\gamma = \alpha = \beta
$$
从而证明 $\alpha = \beta$。

\medskip\noindent
第二点的证明完全相同，只需改用推出 $Y\amalg_X Y = Y$。
\end{proof}





\section{极限与余极限}
\label{categories-section-limits}

\noindent
设 $\mathcal{C}$ 为范畴。$\mathcal{C}$ 中的一个{\it 图表}，
就是一个函子 $M : \mathcal{I} \to \mathcal{C}$。
称 $\mathcal{I}$ 为{\it 指标范畴}，也称 $M$ 为
$\mathcal{I}$-{图表}。我们用记号 $M_i$ 表示
$\mathcal{I}$ 的对象 $i$ 的像。因此，对 $\mathcal{I}$ 中的态射
$\phi : i \to i'$，有 $M(\phi) : M_i \to M_{i'}$。

\begin{definition}
\label{categories-definition-limit}
范畴 $\mathcal{C}$ 中 $\mathcal{I}$-图表 $M$ 的一个{\it 极限}，
由 $\mathcal{C}$ 中的对象 $\lim_\mathcal{I} M$ 连同态射
$p_i : \lim_\mathcal{I} M \to M_i$ 给出，并满足：
\begin{enumerate}
\item 对 $\mathcal{I}$ 中的态射 $\phi : i \to i'$，有
$p_{i'} = M(\phi) \circ p_i$；
\item 对 $\mathcal{C}$ 的任意对象 $W$ 以及由
$i \in \Ob(\mathcal{I})$ 编号的任意态射族 $q_i : W \to M_i$，
若对 $\mathcal{I}$ 中所有 $\phi : i \to i'$ 都有
$q_{i'} = M(\phi) \circ q_i$，则存在唯一态射
$q : W \to \lim_\mathcal{I} M$，使得对 $\mathcal{I}$ 的每个对象
$i$ 均有 $q_i = p_i \circ q$。
\end{enumerate}
\end{definition}

\noindent
若极限 $(\lim_\mathcal{I} M, (p_i)_{i\in \Ob(\mathcal{I})})$ 存在，
则由定义中的唯一性要求，它在唯一同构意义下唯一。
对象对的乘积、纤维积与等化子都是极限的例子。
空图表的极限是 $\mathcal{C}$ 的终对象。
在集合范畴中，所有极限都存在。对偶概念是余极限。

\begin{definition}
\label{categories-definition-colimit}
范畴 $\mathcal{C}$ 中 $\mathcal{I}$-图表 $M$ 的一个{\it 余极限}，
由 $\mathcal{C}$ 中的对象 $\colim_\mathcal{I} M$ 连同态射
$s_i : M_i \to \colim_\mathcal{I} M$ 给出，并满足：
\begin{enumerate}
\item 对 $\mathcal{I}$ 中的态射 $\phi : i \to i'$，有
$s_i = s_{i'} \circ M(\phi)$；
\item 对 $\mathcal{C}$ 的任意对象 $W$ 以及由
$i \in \Ob(\mathcal{I})$ 编号的任意态射族 $t_i : M_i \to W$，
若对 $\mathcal{I}$ 中所有 $\phi : i \to i'$ 都有
$t_i = t_{i'} \circ M(\phi)$，则存在唯一态射
$t : \colim_\mathcal{I} M \to W$，使得对 $\mathcal{I}$ 的每个对象
$i$ 均有 $t_i = t \circ s_i$。
\end{enumerate}
\end{definition}

\noindent
若余极限 $(\colim_\mathcal{I} M, (s_i)_{i\in \Ob(\mathcal{I})})$ 存在，
则由定义中的唯一性要求，它在唯一同构意义下唯一。
对象对的余积、推出与余等化子都是余极限的例子。
空图表的余极限是 $\mathcal{C}$ 的始对象。
在集合范畴中，所有余极限都存在。

\begin{remark}
\label{categories-remark-diagram-small}
（余）极限的指标范畴绝不允许具有真类那么多的对象。
在本项目中，这意味着它不能是注 \ref{categories-remark-big-categories}
所列出的范畴之一。
\end{remark}

\begin{remark}
\label{categories-remark-limit-colim}
我们常用 $\lim_i M_i$、$\colim_i M_i$、
$\lim_{i\in \mathcal{I}} M_i$ 或 $\colim_{i\in \mathcal{I}} M_i$
代替以 $\mathcal{I}$ 为下标的写法。使用这种记号，并利用下文
第 \ref{categories-section-limit-sets} 节对集合的极限与余极限的描述，
可以作如下陈述。设 $M : \mathcal{I} \to \mathcal{C}$ 为图表。
\begin{enumerate}
\item 若对象 $\lim_i M_i$ 存在，则它满足
$$
\Mor_\mathcal{C}(W, \lim_i M_i)
=
\lim_i \Mor_\mathcal{C}(W, M_i)
$$
其中右侧的极限在集合范畴中取。
\item 若对象 $\colim_i M_i$ 存在，则它满足
$$
\Mor_\mathcal{C}(\colim_i M_i, W)
=
\lim_{i \in \mathcal{I}^{opp}} \Mor_\mathcal{C}(M_i, W)
$$
其中右侧取对偶范畴上的极限，并以集合范畴为值域。
\end{enumerate}
由 Yoneda 引理（及其对偶），这些公式分别完全确定极限与余极限。
\end{remark}

\begin{remark}
\label{categories-remark-cones-and-cocones}
设 $M : \mathcal{I} \to \mathcal{C}$ 为图表。
在此情形下，$M$ 的一个{\it 锥}由对象 $W$ 以及态射族
$q_i : W \to M_i$、$i \in \Ob(\mathcal{I})$ 给出，
并要求对 $\mathcal{I}$ 的所有态射 $\phi : i \to i'$，图表
$$
\xymatrix{
& W \ar[dl]_{q_i} \ar[dr]^{q_{i'}} \\
M_i \ar[rr]^{M(\phi)} & & M_{i'}
}
$$
交换。所有锥构成一个范畴，其态射概念是显然的。
显然，若 $M$ 的极限存在，它就是锥范畴中的终对象。
对偶地，$M$ 的一个{\it 余锥}由对象 $W$ 与态射族
$t_i : M_i \to W$ 给出，并要求对 $\mathcal{I}$ 中的所有态射
$\phi : i \to i'$，图表
$$
\xymatrix{
M_i \ar[rr]^{M(\phi)} \ar[dr]_{t_i} & & M_{i'} \ar[dl]^{t_{i'}} \\
& W
}
$$
交换。所有余锥构成一个范畴，其态射概念也是显然的。
与上面类似，$M$ 的余极限存在，当且仅当余锥范畴有始对象。
\end{remark}

\noindent
作为极限与余极限概念的应用，我们定义乘积与余积。

\begin{definition}
\label{categories-definition-product}
设 $I$ 为集合，并且对每个 $i \in I$ 给定范畴 $\mathcal{C}$
的对象 $M_i$。{\it 乘积} $\prod_{i\in I} M_i$ 按定义是
$\lim_\mathcal{I} M$（若其存在），其中 $\mathcal{I}$ 是以
$I$ 的元素为对象、且只以恒等态射为态射的范畴。
\end{definition}

\noindent
一个重要特例是 $I = \emptyset$；此时乘积是该范畴的终对象。
态射 $p_i : \prod M_i \to M_i$ 称为{\it 投影态射}。

\begin{definition}
\label{categories-definition-coproduct}
设 $I$ 为集合，并且对每个 $i \in I$ 给定范畴 $\mathcal{C}$
的对象 $M_i$。{\it 余积} $\coprod_{i\in I} M_i$ 按定义是
$\colim_\mathcal{I} M$（若其存在），其中 $\mathcal{I}$ 是以
$I$ 的元素为对象、且只以恒等态射为态射的范畴。
\end{definition}

\noindent
一个重要特例是 $I = \emptyset$；此时余积是该范畴的始对象。
注意，余积配有态射 $M_i \to \coprod M_i$。
这些态射有时称为{\it 余投影}。

\begin{lemma}
\label{categories-lemma-functorial-colimit}
设 $M : \mathcal{I} \to \mathcal{C}$ 与
$N : \mathcal{J} \to \mathcal{C}$ 为余极限均存在的图表。
设 $H : \mathcal{I} \to \mathcal{J}$ 为函子，并设
$t : M \to N \circ H$ 为函子变换。则存在唯一态射
$$
\theta :
\colim_\mathcal{I} M
\longrightarrow
\colim_\mathcal{J} N
$$
使所有图表
$$
\xymatrix{
M_i \ar[d]_{t_i} \ar[r]
&
\colim_\mathcal{I} M \ar[d]^{\theta}
\\
N_{H(i)} \ar[r]
&
\colim_\mathcal{J} N
}
$$
交换。
\end{lemma}

\begin{proof}
略。
\end{proof}

\begin{lemma}
\label{categories-lemma-functorial-limit}
设 $M : \mathcal{I} \to \mathcal{C}$ 与
$N : \mathcal{J} \to \mathcal{C}$ 为极限均存在的图表。
设 $H : \mathcal{I} \to \mathcal{J}$ 为函子，并设
$t : N \circ H \to M$ 为函子变换。则存在唯一态射
$$
\theta :
\lim_\mathcal{J} N
\longrightarrow
\lim_\mathcal{I} M
$$
使所有图表
$$
\xymatrix{
\lim_\mathcal{J} N \ar[d]^{\theta} \ar[r]
&
N_{H(i)} \ar[d]_{t_i}
\\
\lim_\mathcal{I} M \ar[r]
&
M_i
}
$$
交换。
\end{lemma}

\begin{proof}
略。
\end{proof}


\begin{lemma}
\label{categories-lemma-colimits-commute}
设 $\mathcal{I}$、$\mathcal{J}$ 为指标范畴，并设
$M : \mathcal{I} \times \mathcal{J} \to \mathcal{C}$ 为函子。
假设对所有 $i$，$M_{i, \infty} = \colim_j M_{i,j}$ 都存在。
则所得函子 $M_{-, \infty} : \mathcal{I} \to \mathcal{C}$
有余极限，当且仅当 $M$ 有余极限；此时二者的余极限相同。
特别地，只要所示各余极限都存在，就有
$$
\colim_i \colim_j M_{i, j}
=
\colim_{i, j} M_{i, j}
=
\colim_j \colim_i M_{i, j}
$$
极限的情形类似。
\end{lemma}

\begin{proof}
略。
\end{proof}

\begin{lemma}
\label{categories-lemma-limits-products-equalizers}
设 $M : \mathcal{I} \to \mathcal{C}$ 为图表。
记 $I = \Ob(\mathcal{I})$，$A = \text{Arrows}(\mathcal{I})$，
并以 $s, t : A \to I$ 表示源映射与靶映射。
假设 $\prod_{i \in I} M_i$ 与 $\prod_{a \in A} M_{t(a)}$ 存在，
并假设
$$
\xymatrix{
\prod_{i \in I} M_i
\ar@<1ex>[r]^\phi \ar@<-1ex>[r]_\psi
&
\prod_{a \in A} M_{t(a)}
}
$$
的等化子存在，其中各态射由其分量按如下方式确定：
$p_a \circ \psi = M(a) \circ p_{s(a)}$，且
$p_a \circ \phi = p_{t(a)}$。则这个等化子就是该图表的极限。
\end{lemma}

\begin{proof}
略。
\end{proof}


\begin{lemma}
\label{categories-lemma-colimits-coproducts-coequalizers}
\begin{slogan}
若所有余积与余等化子都存在，则所有余极限都存在。
\end{slogan}
设 $M : \mathcal{I} \to \mathcal{C}$ 为图表。
记 $I = \Ob(\mathcal{I})$，$A = \text{Arrows}(\mathcal{I})$，
并以 $s, t : A \to I$ 表示源映射与靶映射。
假设 $\coprod_{i \in I} M_i$ 与 $\coprod_{a \in A} M_{s(a)}$ 存在，
并假设
$$
\xymatrix{
\coprod_{a \in A} M_{s(a)}
\ar@<1ex>[r]^\phi \ar@<-1ex>[r]_\psi
&
\coprod_{i \in I} M_i
}
$$
的余等化子存在，其中各态射由其分量按如下方式确定：
分量 $M_{s(a)}$ 通过态射 $M(a)$，经 $\psi$ 映到分量
$M_{t(a)}$。分量 $M_{s(a)}$ 经 $\phi$，由恒等态射映到分量
$M_{s(a)}$。则这个余等化子就是该图表的余极限。
\end{lemma}

\begin{proof}
略。
\end{proof}






\section{集合范畴中的极限与余极限}
\label{categories-section-limit-sets}

\noindent
在 $\textit{Sets}$ 中，极限与余极限不仅存在，而且很容易描述。
具体地，设 $M : \mathcal{I} \to \textit{Sets}$、$i \mapsto M_i$
为集合图表，并记 $I = \Ob(\mathcal{I})$。其极限描述为
$$
\lim_\mathcal{I} M
=
\{
(m_i)_{i\in I} \in \prod\nolimits_{i\in I} M_i
\mid
\forall \phi : i \to i' \text{ 于 }\mathcal{I},
M(\phi)(m_i) = m_{i'}
\}.
$$
因此，可把极限的一个元素看作所有集合 $M_i$ 中元素的一个相容系。

\medskip\noindent
另一方面，余极限为
$$
\colim_\mathcal{I} M
=
(\coprod\nolimits_{i\in I} M_i)/\sim
$$
其中等价关系 $\sim$ 由下列关系生成：若 $m_i \in M_i$、
$m_{i'} \in M_{i'}$，且对某个 $\phi : i \to i'$ 有
$M(\phi)(m_i) = m_{i'}$，则令 $m_i \sim m_{i'}$。
换言之，$m_i \in M_i$ 与 $m_{i'} \in M_{i'}$ 等价，
当且仅当 $\mathcal{I}$ 中存在一条态射链
$$
\xymatrix{
&
i_1 \ar[ld] \ar[rd] & &
i_3 \ar[ld] & &
i_{2n-1} \ar[rd] & \\
i = i_0 & &
i_2 & &
\ldots & &
i_{2n} = i'
}
$$
以及元素 $m_{i_j} \in M_{i_j}$，使它们在上述 $\mathcal{I}$
中态射所诱导的映射 $M_{i_{2k-1}} \to M_{i_{2k-2}}$ 与
$M_{i_{2k-1}} \to M_{i_{2k}}$ 下彼此对应。

\medskip\noindent
这种对象用起来并不方便。但若图表是滤过的，其描述便简单得多。
我们将在第 \ref{categories-section-directed-colimits} 节说明这一点。



\section{连通极限}
\label{categories-section-connected-limits}

\noindent
若一个（余）极限的指标范畴是连通的，则称该（余）极限是连通的。

\begin{definition}
\label{categories-definition-category-connected}
若由
$x \sim y \Leftrightarrow \Mor_\mathcal{I}(x, y) \not = \emptyset$
生成的等价关系恰有一个等价类，则称范畴 $\mathcal{I}$
是{\it 连通的}。
\end{definition}

\noindent
这里遵循拓扑学定义 \ref{topology-definition-connected-components}
的约定：连通空间非空。下述结果在某种宽泛意义上刻画连通极限。

\begin{lemma}
\label{categories-lemma-connected-limit-over-X}
设 $\mathcal{C}$ 为范畴，$X$ 为 $\mathcal{C}$ 的对象，并设
$M : \mathcal{I} \to \mathcal{C}/X$ 为 $X$ 上对象范畴中的图表。
若指标范畴 $\mathcal{I}$ 连通，且 $M$ 在 $\mathcal{C}/X$
中的极限存在，则复合
$\mathcal{I} \to \mathcal{C}/X \to \mathcal{C}$
的极限也存在，且二者相同。
\end{lemma}

\begin{proof}
设 $L \to X$ 为表示 $\mathcal{C}/X$ 中极限的对象。考虑函子
$$
W \longmapsto \lim_i \Mor_\mathcal{C}(W, M_i).
$$
设 $(\varphi_i)$ 为右侧集合的元素。由于每个 $M_i$ 都配有态射
$s_i : M_i \to X$，得到态射
$f_i = s_i \circ \varphi_i : W \to X$。
但因为 $\mathcal{I}$ 连通，所有 $f_i$ 都相等；又因
$\mathcal{I}$ 非空，至少有一个 $f_i$。因此这个公共值
$W \to X$ 赋予 $W$ 以 $\mathcal{C}/X$ 中对象的结构，而
$(\varphi_i)$ 给出 $\lim_i \Mor_{\mathcal{C}/X}(W, M_i)$
的一个元素。由此得到唯一态射 $\phi : W \to L$，使每个
$\varphi_i$ 都是 $\phi$ 与 $L \to M_i$ 的复合，正合所需。
\end{proof}

\begin{lemma}
\label{categories-lemma-connected-colimit-under-X}
设 $\mathcal{C}$ 为范畴，$X$ 为 $\mathcal{C}$ 的对象，并设
$M : \mathcal{I} \to X/\mathcal{C}$ 为 $X$ 下对象范畴中的图表。
若指标范畴 $\mathcal{I}$ 连通，且 $M$ 在 $X/\mathcal{C}$
中的余极限存在，则复合
$\mathcal{I} \to X/\mathcal{C} \to \mathcal{C}$
的余极限也存在，且二者相同。
\end{lemma}

\begin{proof}
略。提示：本引理与引理 \ref{categories-lemma-connected-limit-over-X} 对偶。
\end{proof}




\section{共尾范畴与始范畴}
\label{categories-section-cofinal}

\noindent
文献中有时在下述定义里用“终”一词代替“共尾”。

\begin{definition}
\label{categories-definition-cofinal}
设 $H : \mathcal{I} \to \mathcal{J}$ 为范畴之间的函子。
若下列条件成立，则称{\it $\mathcal{I}$ 在 $\mathcal{J}$ 中共尾}，
或称 $H$ {\it 共尾}：
\begin{enumerate}
\item 对所有 $y \in \Ob(\mathcal{J})$，存在
$x \in \Ob(\mathcal{I})$ 及态射 $y \to H(x)$；
\item 给定 $y \in \Ob(\mathcal{J})$、
$x, x' \in \Ob(\mathcal{I})$ 及态射 $y \to H(x)$、
$y \to H(x')$，存在 $\mathcal{I}$ 中的态射序列
$$
x = x_0 \leftarrow x_1 \rightarrow x_2 \leftarrow x_3 \rightarrow \ldots
\rightarrow x_{2n} = x'
$$
以及 $\mathcal{J}$ 中的态射 $y \to H(x_i)$，其中
$y \to H(x_0)$ 与 $y \to H(x_{2n})$ 是给定态射，并且对
$k = 0, \ldots, n - 1$，图表
$$
\xymatrix{
& y \ar[ld] \ar[d] \ar[rd] \\
H(x_{2k}) & H(x_{2k + 1}) \ar[l] \ar[r] & H(x_{2k + 2})
}
$$
交换。
\end{enumerate}
换言之，固定 $\mathcal{J}$ 的对象 $y$，考虑所有序对
$(x, b)$ 构成的集合 $S$，其中 $x$ 是 $\mathcal{I}$ 的对象，
$b : y \to H(x)$ 是态射。在 $S$ 上考虑由下列关系生成的
等价关系：若存在态射 $a : x \to x'$ 满足
$b' = H(a) \circ b$，则令 $(x, b) \sim (x', b')$。
此时 $S$ 应恰好由一个等价类组成。
\end{definition}

\begin{lemma}
\label{categories-lemma-cofinal}
设 $H : \mathcal{I} \to \mathcal{J}$ 为范畴的函子，并假设
$\mathcal{I}$ 在 $\mathcal{J}$ 中共尾。则对每个图表
$M : \mathcal{J} \to \mathcal{C}$，只要任一边存在，就有典范同构
$$
\colim_\mathcal{I} M \circ H
=
\colim_\mathcal{J} M
$$
\end{lemma}

\begin{proof}
略。
\end{proof}

\begin{definition}
\label{categories-definition-initial}
设 $H : \mathcal{I} \to \mathcal{J}$ 为范畴之间的函子。
若下列条件成立，则称{\it $\mathcal{I}$ 在 $\mathcal{J}$ 中是始的}，
或称 $H$ {\it 是始的}：
\begin{enumerate}
\item 对所有 $y \in \Ob(\mathcal{J})$，存在
$x \in \Ob(\mathcal{I})$ 及态射 $H(x) \to y$；
\item 对任意 $y \in \Ob(\mathcal{J})$、
$x , x' \in \Ob(\mathcal{I})$ 以及 $\mathcal{J}$ 中的态射
$H(x) \to y$、$H(x') \to y$，存在 $\mathcal{I}$ 中的态射序列
$$
x = x_0 \leftarrow x_1 \rightarrow x_2 \leftarrow x_3 \rightarrow \ldots
\rightarrow x_{2n} = x'
$$
以及 $\mathcal{J}$ 中的态射 $H(x_i) \to y$，使得对
$k = 0, \ldots, n - 1$，图表
$$
\xymatrix{
H(x_{2k}) \ar[rd] &
H(x_{2k + 1}) \ar[l] \ar[r] \ar[d] &
H(x_{2k + 2}) \ar[ld] \\
& y
}
$$
交换。
\end{enumerate}
\end{definition}

\noindent
这正是“共尾”函子的对偶概念。

\begin{lemma}
\label{categories-lemma-initial}
设 $H : \mathcal{I} \to \mathcal{J}$ 为范畴的函子，并假设
$\mathcal{I}$ 在 $\mathcal{J}$ 中是始的。则对每个图表
$M : \mathcal{J} \to \mathcal{C}$，只要任一边存在，就有典范同构
$$
\lim_\mathcal{I} M \circ H = \lim_\mathcal{J} M
$$
\end{lemma}

\begin{proof}
略。
\end{proof}

\begin{lemma}
\label{categories-lemma-colimit-constant-connected-fibers}
设 $F : \mathcal{I} \to \mathcal{I}'$ 为函子。假设
\begin{enumerate}
\item $\mathcal{I}$ 在 $\mathcal{I}'$ 上的所有纤维范畴
（见定义 \ref{categories-definition-fibre-category}）均连通；
\item 对 $\mathcal{I}'$ 中每个态射 $\alpha' : x' \to y'$，
都存在 $\mathcal{I}$ 中的态射 $\alpha : x \to y$，使得
$F(\alpha) = \alpha'$。
\end{enumerate}
则对每个图表 $M : \mathcal{I}' \to \mathcal{C}$，余极限
$\colim_\mathcal{I} M \circ F$ 存在，当且仅当
$\colim_{\mathcal{I}'} M$ 存在；若存在，则二者相同。
\end{lemma}

\begin{proof}
可以先证明 $\mathcal{I}$ 在 $\mathcal{I}'$ 中共尾，
再应用引理 \ref{categories-lemma-cofinal}。也可如下直接证明。
只需证明，对 $\mathcal{C}$ 的任意对象 $T$，有
$$
\lim_{\mathcal{I}^{opp}} \Mor_\mathcal{C}(M_{F(i)}, T)
=
\lim_{(\mathcal{I}')^{opp}} \Mor_\mathcal{C}(M_{i'}, T)
$$
若 $(g_{i'})_{i' \in \Ob(\mathcal{I}')}$ 是右侧的元素，
令 $f_i = g_{F(i)}$，便得到左侧的元素
$(f_i)_{i \in \Ob(\mathcal{I})}$。反过来，设
$(f_i)_{i \in \Ob(\mathcal{I})}$ 为左侧的元素。
注意，在每个（连通）纤维范畴 $\mathcal{I}_{i'}$ 上，
函子 $M \circ F$ 都是取常值 $M_{i'}$ 的常值函子。
因此，所有满足 $F(i) = i'$ 的 $i \in \Ob(\mathcal{I})$
所对应的态射 $f_i$ 都相同，并确定一个良定义态射
$g_{i'} : M_{i'} \to T$。由假设 (2)，族
$(g_{i'})_{i' \in \Ob(\mathcal{I}')}$ 给出右侧的一个元素。
\end{proof}

\begin{lemma}
\label{categories-lemma-product-with-connected}
设 $\mathcal{I}$ 与 $\mathcal{J}$ 为范畴，并以
$p : \mathcal{I} \times \mathcal{J} \to \mathcal{J}$ 表示投影。
若 $\mathcal{I}$ 连通，则对图表 $M : \mathcal{J} \to \mathcal{C}$，
余极限 $\colim_\mathcal{J} M$ 存在，当且仅当
$\colim_{\mathcal{I} \times \mathcal{J}} M \circ p$ 存在；
若存在，则二者相等。
\end{lemma}

\begin{proof}
这是引理 \ref{categories-lemma-colimit-constant-connected-fibers} 的特例。
\end{proof}






\section{有限极限与余极限}
\label{categories-section-finite-limits}

\noindent
若一个（余）极限的指标范畴有限，即该指标范畴只有有限多个对象
和有限多个态射，则称它为{\it 有限}（余）极限。
若指标范畴非空，则称（余）极限{\it 非空}；
若指标范畴连通，则称（余）极限{\it 连通}，
见定义 \ref{categories-definition-category-connected}。
事实证明，有限指标范畴已经“足够多”。

\begin{lemma}
\label{categories-lemma-finite-diagram-category}
设 $\mathcal{I}$ 为满足下列条件的范畴：
\begin{enumerate}
\item $\Ob(\mathcal{I})$ 有限；
\item 存在有限多个态射
$f_1, \ldots, f_m \in \text{Arrows}(\mathcal{I})$，
使 $\mathcal{I}$ 的每个态射都能写成复合
$f_{j_1} \circ f_{j_2} \circ \ldots \circ f_{j_k}$。
\end{enumerate}
则存在函子 $F : \mathcal{J} \to \mathcal{I}$，使得
\begin{enumerate}
\item[(a)] $\mathcal{J}$ 是有限范畴；
\item[(b)] 对任意图表 $M : \mathcal{I} \to \mathcal{C}$，
$M$ 在 $\mathcal{I}$ 上的（余）极限存在，当且仅当
$M \circ F$ 在 $\mathcal{J}$ 上的（余）极限存在；在这种情形下，
这两个（余）极限典范同构。
\end{enumerate}
此外，$\mathcal{J}$ 连通（相应地，非空），当且仅当
$\mathcal{I}$ 连通（相应地，非空）。
\end{lemma}

\begin{proof}
设 $\Ob(\mathcal{I}) = \{x_1, \ldots, x_n\}$。
以 $s, t : \{1, \ldots, m\} \to \{1, \ldots, n\}$
表示满足 $f_j : x_{s(j)} \to x_{t(j)}$ 的函数。
令 $\Ob(\mathcal{J}) = \{y_1, \ldots, y_n, z_1, \ldots, z_n\}$。
除恒等态射外，再引入态射
$g_j : y_{s(j)} \to z_{t(j)}$、$j = 1, \ldots, m$，
以及态射 $h_i : y_i \to z_i$、$i = 1, \ldots, n$。
由于 $\mathcal{J}$ 中所有非恒等态射都从某个 $y$ 指向某个 $z$，
没有需要定义的复合，也没有需要检验的结合律。
令 $F(y_i) = F(z_i) = x_i$、$F(g_j) = f_j$ 且
$F(h_i) = \text{id}_{x_i}$。显然 $F$ 是函子，
$\mathcal{J}$ 是有限的，并且 $\mathcal{J}$ 连通（相应地，非空）
当且仅当 $\mathcal{I}$ 如此。

\medskip\noindent
设 $M : \mathcal{I} \to \mathcal{C}$ 为图表。
考虑 $\mathcal{C}$ 的对象 $W$ 以及定义 \ref{categories-definition-limit}
所述的态射 $q_i : W \to M(x_i)$。
把它们取作 $q_i : W \to M(F(y_i)) = M(F(z_i)) = M(x_i)$，
便得到定义 \ref{categories-definition-limit} 中图表 $M \circ F$ 所需的映射族。
反过来，设给定定义 \ref{categories-definition-limit} 中图表 $M \circ F$
所需的映射 $qy_i : W \to M(F(y_i))$ 与
$qz_i : W \to M(F(z_i))$。由于
$$
M(F(h_i)) = \text{id} : M(F(y_i)) = M(x_i) \longrightarrow M(x_i) = M(F(z_i))
$$
可知对所有 $i$ 都有 $qy_i = qz_i$。令 $q_i$ 等于这个公共值。
$q_{s(j)} = qy_{s(j)}$ 与 $q_{t(j)} = qz_{t(j)}$
对态射 $M(f_j)$ 的相容性，保证族 $q_i$ 对 $\mathcal{I}$
中的所有态射相容，因为根据假设，每个这样的态射都是态射 $f_j$
的复合。因此得到典范双射
$$
\lim_{B \in \Ob(\mathcal{J})} \Mor_\mathcal{C}(W, M(F(B)))
=
\lim_{A \in \Ob(\mathcal{I})} \Mor_\mathcal{C}(W, M(A))
$$
这蕴含引理关于极限的陈述。余极限的陈述可用同样方式证明
（证明略）。
\end{proof}

\begin{lemma}
\label{categories-lemma-fibre-products-equalizers-exist}
设 $\mathcal{C}$ 为范畴。下列条件等价：
\begin{enumerate}
\item $\mathcal{C}$ 中存在连通有限极限。
\item $\mathcal{C}$ 中存在等化子与纤维积。
\end{enumerate}
\end{lemma}

\begin{proof}
由于等化子与纤维积都是有限连通极限，故 (1) 蕴含 (2)。
反之，设 $\mathcal{I}$ 为有限连通指标范畴，并设
$F : \mathcal{J} \to \mathcal{I}$ 为引理
\ref{categories-lemma-finite-diagram-category} 的证明中构造的指标范畴函子。
于是可用 $\mathcal{J}$ 代替 $\mathcal{I}$。由此可假设
$\Ob(\mathcal{I}) = \{x_1, \ldots, x_n\} \amalg \{y_1, \ldots, y_m\}$，
其中 $n, m \geq 1$，并且 $\mathcal{I}$ 中所有非恒等态射都是
某些 $i$、$j$ 所对应的态射 $f : x_i \to y_j$。

\medskip\noindent
假设 $n > 1$。由于 $\mathcal{I}$ 连通，存在指标
$i_1, i_2$、$j_0$ 以及态射 $a : x_{i_1} \to y_{j_0}$、
$b : x_{i_2} \to y_{j_0}$。考虑范畴
$$
\mathcal{I}' =
\{x\} \amalg \{x_1, \ldots, \hat x_{i_1}, \ldots, \hat x_{i_2}, \ldots x_n\}
\amalg \{y_1, \ldots, y_m\}
$$
其中
$$
\Mor_{\mathcal{I}'}(x, y_j) = \Mor_\mathcal{I}(x_{i_1}, y_j)
\amalg \Mor_\mathcal{I}(x_{i_2}, y_j)
$$
而所有其他态射集都与 $\mathcal{I}$ 中相同。
对任意函子 $M : \mathcal{I} \to \mathcal{C}$，可构造函子
$M' : \mathcal{I}' \to \mathcal{C}$，令
$$
M'(x) = M(x_{i_1}) \times_{M(a), M(y_{j_0}), M(b)} M(x_{i_2})
$$
并且对与某个 $f : x_{i_1} \to y_j$ 对应的态射
$f' : x \to y_j$，令 $M'(f) = M(f) \circ \text{pr}_1$。
于是函子 $M$ 有极限，当且仅当函子 $M'$ 有极限（证明略）。
故由归纳法归约到 $n = 1$ 的情形。

\medskip\noindent
若 $n = 1$，则任意 $M : \mathcal{I} \to \mathcal{C}$ 的极限
是各对映射 $x_1 \to y_j$ 的逐次等化子，因而由假设存在。
\end{proof}

\begin{lemma}
\label{categories-lemma-almost-finite-limits-exist}
设 $\mathcal{C}$ 为范畴。下列条件等价：
\begin{enumerate}
\item $\mathcal{C}$ 中存在非空有限极限。
\item $\mathcal{C}$ 中存在对象对的乘积与等化子。
\item $\mathcal{C}$ 中存在对象对的乘积与纤维积。
\end{enumerate}
\end{lemma}

\begin{proof}
由于对象对的乘积、纤维积与等化子都是指标范畴非空的极限，
故 (1) 同时蕴含 (2) 与 (3)。假设 (2)。则有限非空乘积与
等化子存在。由引理 \ref{categories-lemma-limits-products-equalizers}，
有限非空极限存在，即 (1) 成立。假设 (3)。若
$a, b : A \to B$ 为 $\mathcal{C}$ 的态射，则 $a, b$ 的等化子为
$$
(A \times_{a, B, b} A)\times_{(\text{pr}_1, \text{pr}_2), A \times A, \Delta} A.
$$
因此 (3) 蕴含 (2)，引理得证。
\end{proof}

\begin{lemma}
\label{categories-lemma-finite-limits-exist}
设 $\mathcal{C}$ 为范畴。下列条件等价：
\begin{enumerate}
\item $\mathcal{C}$ 中存在有限极限。
\item 存在有限乘积与等化子。
\item 该范畴有终对象，且纤维积存在。
\end{enumerate}
\end{lemma}

\begin{proof}
由于有限乘积、纤维积、等化子与终对象都是有限指标范畴上的极限，
故 (1) 同时蕴含 (2) 与 (3)。由上面的引理
\ref{categories-lemma-limits-products-equalizers}，可知 (2) 蕴含 (1)。
假设 (3)。注意，乘积 $A \times B$ 是终对象上的纤维积。
若 $a, b : A \to B$ 为 $\mathcal{C}$ 的态射，则 $a, b$ 的等化子为
$$
(A \times_{a, B, b} A)\times_{(\text{pr}_1, \text{pr}_2), A \times A, \Delta} A.
$$
因此 (3) 蕴含 (2)，引理得证。
\end{proof}

\begin{lemma}
\label{categories-lemma-push-outs-coequalizers-exist}
设 $\mathcal{C}$ 为范畴。下列条件等价：
\begin{enumerate}
\item $\mathcal{C}$ 中存在连通有限余极限。
\item $\mathcal{C}$ 中存在余等化子与推出。
\end{enumerate}
\end{lemma}

\begin{proof}
略。提示：这与引理 \ref{categories-lemma-fibre-products-equalizers-exist} 对偶。
\end{proof}

\begin{lemma}
\label{categories-lemma-almost-finite-colimits-exist}
设 $\mathcal{C}$ 为范畴。下列条件等价：
\begin{enumerate}
\item $\mathcal{C}$ 中存在非空有限余极限。
\item $\mathcal{C}$ 中存在对象对的余积与余等化子。
\item $\mathcal{C}$ 中存在对象对的余积与推出。
\end{enumerate}
\end{lemma}

\begin{proof}
略。提示：这与引理 \ref{categories-lemma-almost-finite-limits-exist} 对偶。
\end{proof}

\begin{lemma}
\label{categories-lemma-colimits-exist}
设 $\mathcal{C}$ 为范畴。下列条件等价：
\begin{enumerate}
\item $\mathcal{C}$ 中存在有限余极限。
\item $\mathcal{C}$ 中存在有限余积与余等化子。
\item 该范畴有始对象，且推出存在。
\end{enumerate}
\end{lemma}

\begin{proof}
略。提示：这与引理 \ref{categories-lemma-finite-limits-exist} 对偶。
\end{proof}





\section{滤过余极限}
\label{categories-section-directed-colimits}

\noindent
在滤过图表上取余极限时，余极限较容易计算或描述。定义如下。

\begin{definition}
\label{categories-definition-directed}
若下列条件成立，则称图表 $M : \mathcal{I} \to \mathcal{C}$
是{\it 有向的}，或称{\it 滤过的}：
\begin{enumerate}
\item 范畴 $\mathcal{I}$ 至少有一个对象；
\item 对 $\mathcal{I}$ 的每一对对象 $x, y$，存在对象 $z$
及态射 $x \to z$、$y \to z$；
\item 对 $\mathcal{I}$ 的每一对对象 $x, y$ 以及 $\mathcal{I}$ 中每一对态射
$a, b : x \to y$，存在 $\mathcal{I}$ 中的态射 $c : y \to z$，
使得 $M(c \circ a) = M(c \circ b)$ 作为 $\mathcal{C}$ 中的态射相等。
\end{enumerate}
若 $\text{id} : \mathcal{I} \to \mathcal{I}$ 是滤过的，
则称指标范畴 $\mathcal{I}$ {\it 有向}或{\it 滤过}；
换言之，删去上面第 (3) 项中的 $M$。
\end{definition}

\noindent
可以看到，指标范畴滤过的任意图表都是滤过图表，而滤过余极限通常
正是以这种方式出现。事实上，若 $M : \mathcal{I} \to \mathcal{C}$
是滤过图表，则可把 $M$ 分解为
$\mathcal{I} \to \mathcal{I}' \to \mathcal{C}$，其中
$\mathcal{I}'$ 是滤过指标范畴\footnote{具体地，令
$\mathcal{I}'$ 与 $\mathcal{I}$ 有相同对象，但把
$\Mor_{\mathcal{I}'}(x, y)$ 定义为 $\Mor_\mathcal{I}(x, y)$
关于如下等价关系的商：若 $M(a) = M(b)$，则等同
$a, b : x \to y$。}，并且 $\colim_\mathcal{I} M$ 存在，
当且仅当 $\colim_{\mathcal{I}'} M'$ 存在；此时两个余极限典范同构。

\medskip\noindent
假设 $M : \mathcal{I} \to \textit{Sets}$ 为滤过图表。
此时，公式
$$
\colim_\mathcal{I} M
=
(\coprod\nolimits_{i\in I} M_i)/\sim
$$
中的等价关系可以简单描述为
$$
m_i \sim m_{i'}
\Leftrightarrow
\exists i'', \phi : i \to i'', \phi': i' \to i'',
M(\phi)(m_i) = M(\phi')(m_{i'}).
$$
换言之，两个元素在余极限中相等，当且仅当它们“最终变得相等”。

\begin{lemma}
\label{categories-lemma-directed-commutes}
设 $\mathcal{I}$ 与 $\mathcal{J}$ 为指标范畴。
假设 $\mathcal{I}$ 滤过，$\mathcal{J}$ 有限，并设
$M : \mathcal{I} \times \mathcal{J} \to \textit{Sets}$、
$(i, j) \mapsto M_{i, j}$ 为集合图表所成的图表。则
$$
\colim_i \lim_j M_{i, j}
=
\lim_j \colim_i M_{i, j}.
$$
特别地，$\mathcal{I}$ 上的余极限与集合的有限乘积、纤维积
及等化子可交换。
\end{lemma}

\begin{proof}
略。事实上，证明下述命题是一道有趣的练习：一个范畴滤过，
当且仅当该范畴上的余极限与有限极限（取值于集合范畴）可交换。
\end{proof}

\noindent
下面给出 $\mathcal{J}$ 无限时上述引理的反例。
令 $\mathcal{I}$ 的对象为 $\mathbf{N} = \{1, 2, 3, \ldots\}$，
且当 $i \leq i'$ 时恰有一个态射 $i \to i'$。
令 $\mathcal{J}$ 为离散范畴
$\mathbf{N} = \{1, 2, 3, \ldots\}$，即其态射只有恒等态射。
令 $M_{i, j} = \{1, 2, \ldots, i\}$；当 $i \leq i'$ 时，
态射 $M_{i, j} \to M_{i', j}$ 是显然的包含映射。
此时 $\colim_i M_{i, j} = \mathbf{N}$，因而
$$
\lim_j \colim_i M_{i, j}
=
\prod\nolimits_j \mathbf{N}
=
\mathbf{N}^\mathbf{N}
$$
另一方面，$\lim_j M_{i, j} = \prod\nolimits_j M_{i, j}$，故
$$
\colim_i \lim_j M_{i, j}
=
\bigcup\nolimits_i \{1, 2, \ldots, i\}^{\mathbf{N}}
$$
它比另一个极限更小。

\begin{lemma}
\label{categories-lemma-cofinal-in-filtered}
\begin{reference}
\cite[Proposition 3.2.4]{KS}
\end{reference}
设 $\mathcal{I}$ 为范畴，$\mathcal{J}$ 为其满子范畴。
假设 $\mathcal{I}$ 滤过，并且对 $\mathcal{I}$ 的任意对象 $i$，
都存在到 $\mathcal{J}$ 的某个对象 $j$ 的态射 $i \to j$。
则 $\mathcal{J}$ 滤过，且在 $\mathcal{I}$ 中共尾。
\end{lemma}

\begin{proof}
略。这是对有关概念的一道愉快练习。
\end{proof}

\noindent
事实证明，有时需要更精细地控制可以对指标范畴施加的条件。
因此，我们补充若干关于可能要求的性质的引理。

\begin{lemma}
\label{categories-lemma-preserve-products}
设 $\mathcal{I}$ 为指标范畴，即一个范畴。假设对
$\mathcal{I}$ 的每一对对象 $x, y$，都存在对象 $z$ 及态射
$x \to z$、$y \to z$。则
\begin{enumerate}
\item 若 $M$ 与 $N$ 是 $\mathcal{I}$ 上的集合图表，则
$\colim (M_i \times N_i) \to \colim M_i \times \colim N_i$ 为满射；
\item 一般而言，$\mathcal{I}$ 上集合图表的余极限不与有限非空乘积交换。
\end{enumerate}
\end{lemma}

\begin{proof}
证明 (1)。设 $(\overline{m}, \overline{n})$ 是
$\colim M_i \times \colim N_i$ 的元素。则对某些
$x, y \in \Ob(\mathcal{I})$，可找到 $m \in M_x$ 与 $n \in N_y$，
使 $m$ 映到 $\overline{m}$，而 $n$ 映到 $\overline{n}$。
见第 \ref{categories-section-limit-sets} 节。在 $\mathcal{I}$ 中选取
$a : x \to z$ 与 $b : y \to z$。于是
$(M(a)(m), N(b)(n))$ 是 $(M \times N)_z$ 的元素，
它在 $\colim (M_i \times N_i)$ 中的像映到
$(\overline{m}, \overline{n})$，正合所需。

\medskip\noindent
证明 (2)。设 $G$ 为非平凡群，并令 $\mathcal{I}$ 为只有一个对象、
自同态幺半群为 $G$ 的范畴。则 $\mathcal{I}$ 显然满足引理所述条件。
现在令 $G$ 以平移作用于自身，并把 $G$-集 $G$ 看作取值于集合的
$\mathcal{I}$-图表。则
$$
\colim_\mathcal{I} G \times \colim_\mathcal{I} G \cong G/G \times G/G
$$
不同构于
$$
\colim_\mathcal{I} (G \times G) \cong (G \times G)/G
$$
这个例子说明，即使只关心有限乘积，也不能直接删去引理
\ref{categories-lemma-directed-commutes} 中的附加条件。
\end{proof}

\begin{lemma}
\label{categories-lemma-colimits-abelian-as-sets}
设 $\mathcal{I}$ 为指标范畴，即一个范畴。假设对
$\mathcal{I}$ 的每一对对象 $x, y$，都存在对象 $z$ 及态射
$x \to z$、$y \to z$。设
$M : \mathcal{I} \to \textit{Ab}$ 为 $\mathcal{I}$ 上的阿贝尔群图表。
则 $M$ 在集合范畴中的余极限满射到 $M$ 在阿贝尔群范畴中的余极限。
\end{lemma}

\begin{proof}
回忆集合范畴中的余极限是无交并 $\coprod M_i$ 关于某关系的商，
见第 \ref{categories-section-limit-sets} 节。类似地，阿贝尔群范畴中的余极限
是直和 $\bigoplus M_i$ 的商。引理的假设意味着：给定
$i, j \in \Ob(\mathcal{I})$、$m \in M_i$ 与 $n \in M_j$，
可以找到对象 $k$ 以及态射 $a : i \to k$、$b : j \to k$。
因此，$m + n$ 在余极限中由 $M_k$ 的元素
$M(a)(m) + M(b)(n)$ 表示。故 $\coprod M_i$ 满射到该余极限。
\end{proof}

\begin{lemma}
\label{categories-lemma-split-into-connected}
设 $\mathcal{I}$ 为指标范畴，即一个范畴。假设对
$\mathcal{I}$ 中的每个实线图表
$$
\xymatrix{
x \ar[d] \ar[r] & y \ar@{..>}[d] \\
z \ar@{..>}[r] & w
}
$$
都存在对象 $w$ 和虚线箭头，使图表交换。
则 $\mathcal{I}$ 或为空，或为若干具有同一性质的连通范畴的
非空无交并。
\end{lemma}

\begin{proof}
若 $\mathcal{I}$ 是空范畴，则引理成立。否则，在
$\mathcal{I}$ 的对象上定义关系：若存在 $z$ 以及态射
$x \to z$、$y \to z$，则称 $x \sim y$。
由引理的假设，这是等价关系。因此 $\Ob(\mathcal{I})$
是若干等价类的无交并。令 $\mathcal{I}_j$ 为这些等价类所对应的
满子范畴。于是 $\mathcal{I} = \coprod \mathcal{I}_j$，
且每个 $\mathcal{I}_j$ 都非空，正合所需。
\end{proof}

\begin{lemma}
\label{categories-lemma-preserve-injective-maps}
设 $\mathcal{I}$ 为指标范畴，即一个范畴。假设对
$\mathcal{I}$ 中的每个实线图表
$$
\xymatrix{
x \ar[d] \ar[r] & y \ar@{..>}[d] \\
z \ar@{..>}[r] & w
}
$$
都存在对象 $w$ 和虚线箭头，使图表交换。则
\begin{enumerate}
\item $\mathcal{I}$ 上集合图表的单射态射 $M \to N$，
诱导集合的单射映射 $\colim M_i \to \colim N_i$；
\item 对阿贝尔群图表及其余极限，一般并非如此。
\end{enumerate}
\end{lemma}

\begin{proof}
若 $\mathcal{I}$ 是空范畴，则引理成立。因此可假设
$\mathcal{I}$ 非空。此时可写成
$\mathcal{I} = \coprod \mathcal{I}_j$，其中每个
$\mathcal{I}_j$ 都非空且满足同一性质，见引理
\ref{categories-lemma-split-into-connected}。由于
$\colim_\mathcal{I} M = \coprod_j \colim_{\mathcal{I}_j} M|_{\mathcal{I}_j}$，
这把 (1) 的证明归约到连通情形。

\medskip\noindent
假设 $\mathcal{I}$ 连通且 $M \to N$ 是单射，
即所有映射 $M_i \to N_i$ 都是单射。把 $M_i$ 与
$M_i \to N_i$ 的像等同，也就是把 $M_i$ 看作 $N_i$ 的子集。
下文直接使用第 \ref{categories-section-limit-sets} 节给出的余极限描述。
设 $s, s' \in \colim M_i$ 映到 $\colim N_i$ 的同一元素。
设 $s$ 来自 $M_i$ 的元素 $m$，$s'$ 来自 $M_{i'}$ 的元素 $m'$。
则可找到 $\mathcal{I}$ 的对象序列
$i = i_0, i_1, \ldots, i_n = i'$ 以及态射
$$
\xymatrix{
&
i_1 \ar[ld] \ar[rd] & &
i_3 \ar[ld] & &
i_{2n-1} \ar[rd] & \\
i = i_0 & &
i_2 & &
\ldots & &
i_{2n} = i'
}
$$
和元素 $n_{i_j} \in N_{i_j}$；在上述 $\mathcal{I}$ 中态射
所诱导的映射 $N_{i_{2k-1}} \to N_{i_{2k-2}}$ 与
$N_{i_{2k-1}} \to N_{i_{2k}}$ 下，这些元素彼此对应，
并且 $n_{i_0} = m$、$n_{i_{2n}} = m'$。
我们对 $n$ 归纳证明这蕴含 $s = s'$。初始情形 $n = 0$ 显然。
假设 $n \geq 1$。利用关于 $\mathcal{I}$ 的假设，得到交换图表
$$
\xymatrix{
& i_1 \ar[ld] \ar[rd] \\
i_0 \ar[rd] & & i_2 \ar[ld] \\
& w
}
$$
由于 $m$ 与 $n_{i_2}$ 都是元素 $n_{i_1}$ 的像，可知二者映到
$N_w$ 的同一元素。特别地，该元素是某个 $m'' \in M_w$，
它在 $\colim M_i$ 中给出的元素与 $s$ 相同。于是得到链
$$
\xymatrix{
&
i_3 \ar[ld] \ar[rd] & &
i_5 \ar[ld] & &
i_{2n-1} \ar[rd] & \\
w & &
i_4 & &
\ldots & &
i_{2n} = i'
}
$$
以及 $j \geq 3$ 时的元素 $n_{i_j}$；这条链比原链短。
这证明归纳步骤，从而完成 (1) 的证明。

\medskip\noindent
设 $G$ 为群，并令 $\mathcal{I}$ 为只有一个对象、自同态幺半群为
$G$ 的范畴。则 $\mathcal{I}$ 满足引理所述条件，因为给定
$g_1, g_2 \in G$，总可找到 $h_1, h_2 \in G$，使
$h_1 g_1 = h_2 g_2$。$\textit{Ab}$ 中 $\mathcal{I}$ 上的图表 $M$
等同于带 $G$-作用的阿贝尔群 $M$，而
$\colim_\mathcal{I} M$ 是 $M$ 的余不变量 $M_G$。
令 $G$ 为 $2$ 阶群，它平凡作用于
$M = \mathbf{Z}/2\mathbf{Z}$；把它映入
$N = \mathbf{Z}/2\mathbf{Z} \times \mathbf{Z}/2\mathbf{Z}$
的第一个直和项，并令 $G$ 的非平凡元素按
$(x, y) \mapsto (x + y, y)$ 作用。则
$M_G \to N_G$ 是零映射。
\end{proof}

\begin{lemma}
\label{categories-lemma-split-into-directed}
设 $\mathcal{I}$ 为指标范畴，即一个范畴。假设
\begin{enumerate}
\item 对 $\mathcal{I}$ 中每一对态射 $a : w \to x$ 与
$b : w \to y$，都存在对象 $z$ 及态射 $c : x \to z$、
$d : y \to z$，使得 $c \circ a = d \circ b$；
\item 对每一对态射 $a, b : x \to y$，都存在态射
$c : y \to z$，使得 $c \circ a = c \circ b$。
\end{enumerate}
则 $\mathcal{I}$ 是若干两两无交的滤过指标范畴
$\mathcal{I}_j$ 的并（可以为空）。
\end{lemma}

\begin{proof}
若 $\mathcal{I}$ 是空范畴，则引理成立。否则，在
$\mathcal{I}$ 的对象上定义关系：若存在 $z$ 以及态射
$x \to z$、$y \to z$，则称 $x \sim y$。
由引理的第一个假设，这是等价关系。因此
$\Ob(\mathcal{I})$ 是若干等价类的无交并。
令 $\mathcal{I}_j$ 为这些等价类所对应的满子范畴。
余下结论由定义立即得到。
\end{proof}

\begin{lemma}
\label{categories-lemma-almost-directed-commutes-equalizers}
设 $\mathcal{I}$ 为满足上面引理
\ref{categories-lemma-split-into-directed} 各项假设的指标范畴。
则 $\mathcal{I}$ 上的余极限与集合中的纤维积和等化子交换
（更一般地，与有限连通极限交换）。
\end{lemma}

\begin{proof}
由引理 \ref{categories-lemma-split-into-directed}，可写
$\mathcal{I} = \coprod \mathcal{I}_j$，其中每个
$\mathcal{I}_j$ 都滤过。由引理 \ref{categories-lemma-directed-commutes}，
$\mathcal{I}_j$ 上的余极限与等化子和纤维积交换。
因此，只需证明集合范畴中的等化子和纤维积与余积交换
（包括空余积）。换言之，给定集合 $J$、集合
$A_j, B_j, C_j$ 以及 $j \in J$ 时的集合映射
$A_j \to B_j$、$C_j \to B_j$，需证明
$$
(\coprod\nolimits_{j \in J} A_j)
\times_{(\coprod\nolimits_{j \in J} B_j)}
(\coprod\nolimits_{j \in J} C_j)
=
\coprod\nolimits_{j \in J} A_j \times_{B_j} C_j
$$
并且给定 $a_j, a'_j : A_j \to B_j$，还需证明
$$
\text{Equalizer}(
\coprod\nolimits_{j \in J} a_j,
\coprod\nolimits_{j \in J} a'_j)
=
\coprod\nolimits_{j \in J}
\text{Equalizer}(a_j, a'_j)
$$
即使 $J = \emptyset$，这些结论也成立。细节略。
\end{proof}



\section{余滤过极限}
\label{categories-section-codirected-limits}

\noindent
在余滤过图表上取极限时，极限较容易计算或描述。定义如下。

\begin{definition}
\label{categories-definition-codirected}
若下列条件成立，则称图表 $M : \mathcal{I} \to \mathcal{C}$
是{\it 反向的}或{\it 余滤过的}：
\begin{enumerate}
\item 范畴 $\mathcal{I}$ 至少有一个对象；
\item 对 $\mathcal{I}$ 的每一对对象 $x, y$，存在对象 $z$
及态射 $z \to x$、$z \to y$；
\item 对 $\mathcal{I}$ 的每一对对象 $x, y$ 以及
$\mathcal{I}$ 中的每一对态射 $a, b : x \to y$，
存在 $\mathcal{I}$ 中的态射 $c : w \to x$，使得
$M(a \circ c) = M(b \circ c)$ 作为 $\mathcal{C}$ 中的态射相等。
\end{enumerate}
若 $\text{id} : \mathcal{I} \to \mathcal{I}$ 是余滤过的，
则称指标范畴 $\mathcal{I}$ {\it 反向}或{\it 余滤过}；
换言之，删去上面第 (3) 项中的 $M$。
\end{definition}

\noindent
可以看到，指标范畴余滤过的任意图表都是余滤过图表，
而这种情形通常正是以这种方式出现。

\medskip\noindent
为说明集合的余滤过极限为何比一般极限“容易”，我们提及如下事实：
有限非空集合的余滤过图表具有非空极限
（引理 \ref{categories-lemma-nonempty-limit}）。
这一结论对有限非空集合的一般极限并不成立。










\section{预序集上的极限与余极限}
\label{categories-section-posets-limits}

\noindent
图表的一种特殊情形是预序集上的系。

\begin{definition}
\label{categories-definition-directed-set}
设 $I$ 为集合，$\leq$ 为 $I$ 上的二元关系。
\begin{enumerate}
\item 若 $\leq$ 具有传递性（若 $i \leq j$ 且 $j \leq k$，
则 $i \leq k$）和自反性（对所有 $i \in I$ 都有 $i \leq i$），
则称该关系为{\it 预序}。
\item 赋有预序的集合称为{\it 预序集}。
\item 若预序集 $(I, \leq)$ 满足 $I$ 非空，且
$\forall i, j \in I$，存在 $k \in I$ 使
$i \leq k, j \leq k$，则称其为{\it 有向集}。
\item 若 $\leq$ 是反对称的预序（若 $i \leq j$ 且
$j \leq i$，则 $i = j$），则称该关系为{\it 偏序}。
\item 赋有偏序的集合称为{\it 偏序集}。
\item 序关系为偏序的有向集称为{\it 有向偏序集}。
\end{enumerate}
\end{definition}

\noindent
谈论预序集时，通常从记号中略去 $\leq$；也就是说，
说预序集 $I$，而不说预序集 $(I, \leq)$。
给定预序集 $I$，符号 $\geq$ 由下述规则定义：对所有
$i, j \in I$，$i \geq j \Leftrightarrow j \leq i$。
“偏序集”有时简称为“poset”。

\medskip\noindent
给定预序集 $I$，可以构造一个范畴：对象是 $I$ 的元素；
若 $i \leq i'$，则恰有一个态射 $i \to i'$，否则没有这样的态射。
反过来，给定任意两个对象之间至多有一个箭头的范畴 $\mathcal{C}$，
集合 $\Ob(\mathcal{C})$ 上由规则
$x \leq y \Leftrightarrow \Mor_\mathcal{C}(x, y) \not = \emptyset$
定义了一个预序。

\begin{definition}
\label{categories-definition-system-over-poset}
设 $(I, \leq)$ 为预序集，$\mathcal{C}$ 为范畴。
\begin{enumerate}
\item $\mathcal{C}$ 中 $I$ 上的一个{\it 系}，有时也称
$\mathcal{C}$ 中 $I$ 上的{\it 归纳系}，由 $\mathcal{C}$ 的对象
$M_i$ 以及每当 $i \leq i'$ 时给定的态射
$f_{ii'} : M_i \to M_{i'}$ 构成，并满足 $f_{ii}
= \text{id}$，且每当 $i \leq i' \leq i''$ 时有
$f_{ii''} = f_{i'i''} \circ f_{i i'}$。
\item $\mathcal{C}$ 中 $I$ 上的一个{\it 逆系}，
有时也称 $\mathcal{C}$ 中 $I$ 上的{\it 射影系}，
由 $\mathcal{C}$ 的对象 $M_i$ 以及每当 $i' \leq i$ 时给定的态射
$f_{ii'} : M_i \to M_{i'}$ 构成，并满足 $f_{ii}
= \text{id}$，且每当 $i'' \leq i' \leq i$ 时有
$f_{ii''} = f_{i'i''} \circ f_{i i'}$。（注意不等式的方向反转。）
\end{enumerate}
我们用“$(M_i, f_{ii'})$ 是 $I$ 上的一个（逆）系”表示上述情形。
映射 $f_{ii'}$ 有时称为{\it 过渡映射}。
\end{definition}

\noindent
换言之，$I$ 上的一个系就是图表
$M : \mathcal{I} \to \mathcal{C}$，其中 $\mathcal{I}$ 是上面
与 $I$ 相联系的范畴：其对象是 $I$ 的元素，而且
$\mathcal{I}$ 中存在唯一箭头 $i \to i'$，当且仅当 $i \leq i'$。
逆系则是图表 $M : \mathcal{I}^{opp} \to \mathcal{C}$。
从这个观点看，可以对 $I$ 上的任意（逆）系取（余）极限。
不过，习惯上{\it 只对 $I$ 上的系取余极限}，
并且{\it 只对 $I$ 上的逆系取极限}。更精确地说：
给定 $I$ 上的系 $(M_i, f_{ii'})$，系 $(M_i, f_{ii'})$ 的余极限定义为
$$
\colim_{i \in I} M_i = \colim_\mathcal{I} M,
$$
即相应图表的余极限。给定 $I$ 上的逆系 $(M_i, f_{ii'})$，
逆系 $(M_i, f_{ii'})$ 的极限定义为
$$
\lim_{i \in I} M_i = \lim_{\mathcal{I}^{opp}} M,
$$
即相应图表的极限。

\begin{remark}
\label{categories-remark-preorder-versus-partial-order}
设 $I$ 为预序集。可以由 $I$ 构造一个典范偏序集
$\overline{I}$ 以及保序映射 $\pi : I \to \overline{I}$。
具体地，在 $I$ 上按规则
$$
i \sim j \Leftrightarrow (i \leq j\text{ 且 }j \leq i).
$$
定义等价关系 $\sim$。令 $\overline{I} = I/\sim$，并令
$\pi : I \to \overline{I}$ 为商映射。最后，$\overline{I}$
上存在唯一偏序，使得
$\pi(i) \leq \pi(j) \Leftrightarrow i \leq j$。
注意，若 $I$ 是有向集，则 $\overline{I}$ 是有向偏序集。
给定 $\overline{I}$ 上的（逆）系 $N$，令
$M_i = N_{\pi(i)}$，便得到 $I$ 上的（逆）系 $M$。
这一构造定义了 $I$ 与 $\overline{I}$ 上逆系范畴之间的函子。
事实上，这是一个范畴等价。原因是，若 $i \sim j$，则对
$I$ 上的任意系 $M$，映射 $M_i \to M_j$ 与 $M_j \to M_i$
互为逆同构。更精确地，选取 $\pi$ 的截面
$s : \overline{I} \to I$，上述函子的一个拟逆把 $M$ 送到
满足 $N_{\overline{i}} = M_{s(\overline{i})}$ 的 $N$。
最后，这种对应与系的余极限相容：若 $M$ 与 $N$ 如上相关，
并且 $\colim_{\overline{I}} N$ 或 $\colim_I M$ 任一存在，
则另一个也存在，且
$\colim_{\overline{I}} N = \colim_I M$。
对逆系及逆系的极限，有类似结论。
\end{remark}

\noindent
注 \ref{categories-remark-preorder-versus-partial-order} 的要点是：
计算一个系的余极限或一个逆系的极限时，总可假设该预序为偏序。

\begin{definition}
\label{categories-definition-directed-system}
设 $I$ 为预序集。若 $I$ 是有向集
（定义 \ref{categories-definition-directed-set}），即 $I$ 非空，
且对所有 $i_1, i_2 \in I$，存在 $i\in I$ 使
$i_1 \leq i$ 且 $i_2 \leq i$，则称系（相应地，逆系）
$(M_i, f_{ii'})$ 为{\it 有向系}（相应地，{\it 有向逆系}）。
\end{definition}

\noindent
在这种情形下，余极限有时（很不幸地）称为“直极限”。
我们不采用后一术语。事实证明，在下述意义上，
滤过范畴上的图表并不比有向系更一般。

\begin{lemma}
\label{categories-lemma-directed-category-system}
设 $\mathcal{I}$ 为滤过指标范畴。存在有向集 $I$ 以及
$\mathcal{I}$ 中 $I$ 上的系 $(x_i, \varphi_{ii'})$，
并具有下列性质：
\begin{enumerate}
\item 对任意范畴 $\mathcal{C}$ 以及任意取值于 $\mathcal{C}$ 的图表
$M : \mathcal{I} \to \mathcal{C}$，用
$(M(x_i), M(\varphi_{ii'}))$ 表示相应的 $I$ 上的系。
若 $\colim_{i \in I} M(x_i)$ 存在，则
$\colim_\mathcal{I} M$ 也存在，而且引理
\ref{categories-lemma-functorial-colimit} 中的变换
$$
\theta :
\colim_{i \in I} M(x_i)
\longrightarrow
\colim_\mathcal{I} M
$$
是同构。
\item 对任意范畴 $\mathcal{C}$ 以及 $\mathcal{C}$ 中任意图表
$M : \mathcal{I}^{opp} \to \mathcal{C}$，用
$(M(x_i), M(\varphi_{ii'}))$ 表示相应的 $I$ 上的逆系。
若 $\lim_{i \in I} M(x_i)$ 存在，则
$\lim_{\mathcal{I}^{opp}} M$ 也存在，而且引理
\ref{categories-lemma-functorial-limit} 中的变换
$$
\theta :
\lim_{\mathcal{I}^{opp}} M
\longrightarrow
\lim_{i \in I} M(x_i)
$$
是同构。
\end{enumerate}
\end{lemma}

\begin{proof}
如定义 \ref{categories-definition-system-over-poset} 后的正文所述，
可以把预序集看作范畴，把系看作函子。在整个证明中，
我们将自由地在这两个观点之间切换。为证明第一个陈述，
构造一个对应于有向集的范畴 $\mathcal{I}_0$\footnote{事实上，
我们的构造会得到一个有向偏序集。}以及共尾函子
$M_0 : \mathcal{I}_0 \to \mathcal{I}$。于是由引理
\ref{categories-lemma-cofinal}，图表 $M : \mathcal{I} \to \mathcal{C}$
的余极限与图表 $M \circ M_0 : \mathcal{I}_0 \to \mathcal{C}$
的余极限相同，从而得到所需陈述。第二个陈述与第一个对偶；
把 $\mathcal{C}$ 中的极限解释为 $\mathcal{C}^{opp}$ 中的余极限，
即可证明。细节略。

\medskip\noindent
若范畴 $\mathcal{F}$ 中存在有限箭头集 $F$，使
$\mathcal{F}$ 中每个箭头都能由 $F$ 中的箭头复合而得，
则称 $\mathcal{F}$ {\em 有限生成}。特别地，这蕴含
$\mathcal{F}$ 只有有限多个对象。我们先把证明归约到
$\mathcal{I}$ 具有下述性质的情形：
$\mathcal{I}$ 的每个有限生成子范畴都能扩张为具有唯一终对象的
有限生成子范畴。

\medskip\noindent
以 $\omega$ 表示有限序数所成的有向集，并把它看作滤过范畴。
容易验证，积范畴 $\mathcal{I}\times \omega$ 也滤过，且投影
$\Pi : \mathcal{I} \times \omega \to \mathcal{I}$ 共尾。

\medskip\noindent
现在设 $\mathcal{F}$ 为 $\mathcal{I}\times \omega$
的任意有限生成子范畴。利用滤过范畴的公理，并对
$\mathcal{F}$ 的有限生成元集作简单归纳，可为图表
$\mathcal{F} \to \mathcal{I} \times \omega$ 构造
$\mathcal{I} \times \omega$ 中的一个余锥
$(\{f_i\}, i_\infty)$。也就是说，对 $\mathcal{F}$ 中的每个对象
$i$，有态射 $f_i : i \to i_\infty$，并且对
$\mathcal{F}$ 中每个箭头 $f : i \to i'$，有
$f_i = f_{i'} \circ f$。还可选取 $i_\infty$，使得不存在从
$i_\infty$ 到 $\mathcal{F}$ 中对象的箭头。这是可能的，
因为总可把箭头 $f_i$ 后复合一个箭头；该箭头在
$\mathcal{I}$-分量上为恒等，在 $\omega$-分量上严格递增。
现在令 $\mathcal{F}^+$ 为如下范畴：它包含 $\mathcal{F}$
的所有对象与箭头，并加入对象 $i_\infty$、恒等箭头
$\text{id}_{i_\infty}$ 及箭头 $f_i$。由于在 $\mathcal{F}^+$
中不存在从 $i_\infty$ 到 $\mathcal{F}$ 任一对象的箭头，
$\mathcal{F}^+$ 的箭头集对复合封闭，所以
$\mathcal{F}^+$ 的确是范畴。按构造，它是 $\mathcal{I}$
的有限生成子范畴，并以 $i_\infty$ 为唯一终对象。
又由引理 \ref{categories-lemma-cofinal}，任意图表
$M : \mathcal{I} \to \mathcal{C}$ 的余极限都与
$M\circ\Pi$ 的余极限相同；这就给出了所需归约。

\medskip\noindent
$\mathcal{I}$ 中所有具有唯一终对象的有限生成子范畴所成的集合，
按包含关系自然排序。令 $\mathcal{I}_0$ 为与该集合对应的范畴。
还有函子 $M_0 : \mathcal{I}_0 \to \mathcal{I}$：它把
$\mathcal{I}_0$ 中的箭头 $\mathcal{F} \subset \mathcal{F'}$
送到从 $\mathcal{F}$ 的终对象到 $\mathcal{F}'$ 的终对象的唯一映射。
给定 $\mathcal{I}$ 的任意两个有限生成子范畴，由这两个范畴
生成的范畴仍有限生成。由关于 $\mathcal{I}$ 的假设，它还包含于
$\mathcal{I}$ 的某个具有唯一终对象的有限生成子范畴中。
这说明 $\mathcal{I}_0$ 有向。

\medskip\noindent
最后验证 $M_0$ 共尾。$\mathcal{I}$ 的每个对象，都是只含该对象
及其恒等箭头的子范畴中的终对象，故函子 $M_0$ 在对象上满射。
特别地，定义 \ref{categories-definition-cofinal} 的条件 (1) 成立。
给定 $\mathcal{I}$ 的对象 $i$、$\mathcal{I}_0$ 的对象
$\mathcal{F}_1, \mathcal{F}_2$，以及 $\mathcal{I}$ 中的映射
$\varphi_1 : i \to M_0(\mathcal{F}_1)$ 与
$\varphi_2 : i \to M_0(\mathcal{F}_2)$，可取
$\mathcal{F}_{12}$ 为具有唯一终对象的有限生成范畴，
其中包含 $\mathcal{F}_1$、$\mathcal{F}_2$ 及态射
$\varphi_1, \varphi_2$。所得图表
$$
\xymatrix{
& M_0(\mathcal{F}_{12}) & \\
M_0(\mathcal{F}_{1}) \ar[ru] & & M_0(\mathcal{F}_{2}) \ar[lu] \\
& i \ar[lu] \ar[ru]
}
$$
交换，因为它位于范畴 $\mathcal{F}_{12}$ 中，而
$M_0(\mathcal{F}_{12})$ 是该范畴的终对象。
因此条件 (2) 也成立，证明完成。
\end{proof}

\begin{remark}
\label{categories-remark-trick-needed}
注意，有限有向集 $(I, \geq)$ 总有最大对象 $i_\infty$。
因此，这类集合上的任意系 $(M_i, f_{ii'})$ 的余极限都是平凡的，
其含义是该余极限等于 $M_{i_\infty}$。相比之下，
由有限滤过范畴编号的余极限不一定平凡。例如，令
$\mathcal{I}$ 为只有一个对象 $i$ 和一个满足
$e = e \circ e$ 的非平凡态射 $e$ 的范畴。
图表 $M : \mathcal{I} \to Sets$ 的余极限是幂等映射
$M(e)$ 的像。这说明，在引理
\ref{categories-lemma-directed-category-system} 的证明中，采用类似于传到
$\mathcal{I}\times \omega$ 的技巧是必要的。
\end{remark}

\begin{lemma}
\label{categories-lemma-nonempty-limit}
若 $S : \mathcal{I} \to \textit{Sets}$ 是集合的余滤过图表，
且所有 $S_i$ 都有限非空，则 $\lim_i S_i$ 非空。
换言之，有限非空集合的有向逆系的极限非空。
\end{lemma}

\begin{proof}
由引理 \ref{categories-lemma-directed-category-system}，两个陈述等价。
设 $I$ 为有向集，并设 $(S_i)_{i \in I}$ 为 $I$ 上有限非空集合的逆系。
称一个{\it 子系} $T$ 是非空子集族
$T = (T_i)_{i \in I}$、$T_i \subset S_i$，并满足：
对所有 $i' \geq i$，过渡映射 $S_{i'} \to S_i$ 都把
$T_{i'}$ 映入 $T_i$。以 $\mathcal{T}$ 表示子系所成的集合。
我们按包含关系排列 $\mathcal{T}$。设 $T_\alpha$、$\alpha \in A$
是 $\mathcal{T}$ 中元素的一个全序族，并写
$T_\alpha = (T_{\alpha, i})_{i \in I}$。
令 $T_i = \bigcap_{\alpha \in A} T_{\alpha, i}$，便得到下界
$T = (T_i)_{i \in I}$；显然 $T_i$ 是 $S_i$ 的有限非空子集，
因为所有 $T_{\alpha, i}$ 都非空且 $T_\alpha$ 构成全序族。
因此可应用 Zorn 引理，得知 $\mathcal{T}$ 有极小元。

\medskip\noindent
下面分析极小元 $T \in \mathcal{T}$ 的形状。
首先注意，所有映射 $T_{i'} \to T_i$ 都是满射。
事实上，由于 $I$ 是有向集且 $T_i$ 有限，交
$T'_i = \bigcap_{i' \geq i} \Im(T_{i'} \to T_i)$ 非空。
因此 $T' = (T'_i)$ 是包含于 $T$ 的子系；由极小性，
$T' = T$。最后断言，对每个 $i$，该集合都是单点集。
事实上，若 $x \in T_i$，则可在 $i' \geq i$ 时定义
$T'_{i'} = (T_{i'} \to T_i)^{-1}(\{x\})$，并在
$j \not \geq i$ 时令 $T'_j = T_j$。这又是一个子系，
因为上面已看到子系 $T$ 的过渡映射都是满射。
由极小性，$T = T'$，这的确蕴含 $T_i$ 是单点集。
这对每个 $i \in I$ 都成立，所以对某些 $x_i \in S_i$，
有 $T_i = \{x_i\}$，并且对每个 $i' \geq i$，
在映射 $S_{i'} \to S_i$ 下有 $x_{i'} \mapsto x_i$。
换言之，$(x_i) \in \lim S_i$，引理得证。
\end{proof}






\section{本质常值系}
\label{categories-section-essentially-constant}

\noindent
设 $M : \mathcal{I} \to \mathcal{C}$ 为范畴 $\mathcal{C}$ 中的图表，
并假设指标范畴 $\mathcal{I}$ 滤过。此时有三个逐次增强的概念，
它们选出 $\mathcal{C}$ 的一个对象 $X$。第一个概念仅仅是
$$
X = \colim_{i \in \mathcal{I}} M_i.
$$
此时 $X$ 配有余投影 $M_i \to X$。更强的条件是：要求 $X$
为余极限，并存在 $i \in \mathcal{I}$ 及态射 $X \to M_i$，
使复合 $X \to M_i \to X$ 为 $\text{id}_X$。
更强的条件如下。

\begin{definition}
\label{categories-definition-essentially-constant-diagram}
设 $M : \mathcal{I} \to \mathcal{C}$ 为范畴 $\mathcal{C}$ 中的图表。
\begin{enumerate}
\item 假设指标范畴 $\mathcal{I}$ 滤过，并设
$(X, \{M_i \to X\}_i)$ 为 $M$ 的余锥，见注
\ref{categories-remark-cones-and-cocones}。若存在 $i \in \mathcal{I}$
及态射 $X \to M_i$，使得
\begin{enumerate}
\item $X \to M_i \to X$ 是 $\text{id}_X$；
\item 对所有 $j$，都存在 $k$ 以及态射 $i \to k$、$j \to k$，
使态射 $M_j \to M_k$ 等于复合
$M_j \to X \to M_i \to M_k$，
\end{enumerate}
则称 $M$ 是以 $X$ 为{\it 值}的{\it 本质常值}图表。
\item 假设指标范畴 $\mathcal{I}$ 余滤过，并设
$(X, \{X \to M_i\}_i)$ 为 $M$ 的锥，见注
\ref{categories-remark-cones-and-cocones}。若存在 $i \in \mathcal{I}$
及态射 $M_i \to X$，使得
\begin{enumerate}
\item $X \to M_i \to X$ 是 $\text{id}_X$；
\item 对所有 $j$，都存在 $k$ 以及态射 $k \to i$、$k \to j$，
使态射 $M_k \to M_j$ 等于复合
$M_k \to M_i \to X \to M_j$，
\end{enumerate}
则称 $M$ 是以 $X$ 为{\it 值}的{\it 本质常值}图表。
\end{enumerate}
使用本定义时，请牢记引理
\ref{categories-lemma-essentially-constant-is-limit-colimit}。
\end{definition}

\noindent
上下文会表明所指的是两个版本中的哪一个。若可能混淆，
我们称第一个版本中的 $M$ 作为 {\it ind-对象}本质常值，
称第二个版本中的它作为 {\it pro-对象}本质常值。
注 \ref{categories-remark-ind-category} 与 \ref{categories-remark-pro-category}
会进一步说明这一术语。事实上，我们常使用“本质常值系”这一术语，
而严格说来，它只对有向集上的系作了定义。

\begin{definition}
\label{categories-definition-essentially-constant-system}
设 $\mathcal{C}$ 为范畴。若把有向系 $(M_i, f_{ii'})$ 中的 $M$
视为函子 $I \to \mathcal{C}$ 时，它定义一个本质常值图表，
则称其为{\it 本质常值系}。若把有向逆系 $(M_i, f_{ii'})$ 中的 $M$
视为函子 $I^{opp} \to \mathcal{C}$ 时，它定义一个
本质常值逆图表，则称其为{\it 本质常值逆系}。
\end{definition}

\noindent
若 $(M_i, f_{ii'})$ 是本质常值系，且态射 $f_{ii'}$ 是单态射，
则对充分大的 $i \leq i'$，态射 $f_{ii'}$ 都是同构。
另一方面，考虑系
$$
\mathbf{Z}^2 \to \mathbf{Z}^2 \to \mathbf{Z}^2 \to \ldots
$$
其中映射由 $(a, b) \mapsto (a + b, 0)$ 给出。
该系以 $\mathbf{Z}$ 为值本质常值，但每个过渡映射都有核。

\medskip\noindent
下面给出一个非本质常值系的例子。令
$M = \bigoplus_{n \geq 0} \mathbf{Z}$，并令 $S : M \to M$
为移位算子 $(a_0, a_1, \ldots) \mapsto (a_1, a_2, \ldots)$。
此时，以 $S$ 为过渡映射的系 $M \to M \to M \to \ldots$
具有余极限 $0$，且复合 $0 \to M \to 0$ 是恒等态射，
但该系并非本质常值。

\medskip\noindent
下一个引理是一项合理性检验。

\begin{lemma}
\label{categories-lemma-essentially-constant-is-limit-colimit}
设 $M : \mathcal{I} \to \mathcal{C}$ 为图表。
若 $\mathcal{I}$ 滤过，且 $M$ 作为 ind-对象本质常值，
则 $X = \colim M_i$ 存在，并且 $M$ 以 $X$ 为值本质常值。
若 $\mathcal{I}$ 余滤过，且 $M$ 作为 pro-对象本质常值，
则 $X = \lim M_i$ 存在，并且 $M$ 以 $X$ 为值本质常值。
\end{lemma}

\begin{proof}
略。这是对各项定义的一道好练习。
\end{proof}

\begin{remark}
\label{categories-remark-ind-category}
设 $\mathcal{C}$ 为范畴。存在由 $\mathcal{C}$ 的
{\it ind-对象}构成的大范畴 $\text{Ind-}\mathcal{C}$。
具体地，若 $F : \mathcal{I} \to \mathcal{C}$ 与
$G : \mathcal{J} \to \mathcal{C}$ 是 $\mathcal{C}$ 中的滤过图表，
则可定义
$$
\Mor_{\text{Ind-}\mathcal{C}}(F, G) =
\lim_i \colim_j \Mor_\mathcal{C}(F(i), G(j)).
$$
存在典范函子 $\mathcal{C} \to \text{Ind-}\mathcal{C}$，
它把 $X$ 映到取常值 $X$ 的{\it 常值系}。
这是一个全忠实嵌入。用这种语言可知，图表 $F$ 本质常值，
当且仅当 $F$ 同构于一个常值系。若以后需要这些内容，
我们会在此将它表述成一个引理并加以证明。
\end{remark}

\begin{remark}
\label{categories-remark-pro-category}
设 $\mathcal{C}$ 为范畴。存在由 $\mathcal{C}$ 的
{\it pro-对象}构成的大范畴 $\text{Pro-}\mathcal{C}$。
具体地，若 $F : \mathcal{I} \to \mathcal{C}$ 与
$G : \mathcal{J} \to \mathcal{C}$ 是 $\mathcal{C}$ 中的余滤过图表，
则可定义
$$
\Mor_{\text{Pro-}\mathcal{C}}(F, G) =
\lim_j \colim_i \Mor_\mathcal{C}(F(i), G(j)).
$$
存在典范函子 $\mathcal{C} \to \text{Pro-}\mathcal{C}$，
它把 $X$ 映到取常值 $X$ 的{\it 常值系}。
这是一个全忠实嵌入。用这种语言可知，图表 $F$ 本质常值，
当且仅当 $F$ 同构于一个常值系。若以后需要这些内容，
我们会在此将它表述成一个引理并加以证明。
\end{remark}

\begin{example}
\label{categories-example-pro-morphism-inverse-systems}
设 $\mathcal{C}$ 为范畴，并设 $(X_n)$ 与 $(Y_n)$ 为
$\mathcal{C}$ 中在通常序关系下以 $\mathbf{N}$ 为指标的逆系。图示为
$$
\ldots \to X_3 \to X_2 \to X_1
\quad\text{且}\quad
\ldots \to Y_3 \to Y_2 \to Y_1
$$
设 $a : (X_n) \to (Y_n)$ 为 $\mathcal{C}$ 的 pro-对象的态射。
$a$ 包含什么数据？对每个 $n \in \mathbf{N}$，应存在一个
$m(n)$ 及态射 $a_n : X_{m(n)} \to Y_n$。这些态射应在下述意义下相容：
对所有 $n' \geq n$，存在
$m(n', n) \geq m(n'), m(n)$，使图表
$$
\xymatrix{
X_{m(n, n')} \ar[rr] \ar[d] & & X_{m(n)} \ar[d]^{a_n} \\
X_{m(n')} \ar[r]^{a_{n'}} & Y_{n'} \ar[r] & Y_n
}
$$
交换。把 $m(n)$ 替换为
$\max_{k, l \leq n}\{m(n, k), m(k, l)\}$ 后，可得到
$\ldots \geq m(3) \geq m(2) \geq m(1)$ 以及交换图表
$$
\xymatrix{
\ldots \ar[r] &
X_{m(3)} \ar[d]^{a_3} \ar[r] &
X_{m(2)} \ar[d]^{a_2} \ar[r] &
X_{m(1)} \ar[d]^{a_1} \\
\ldots \ar[r] &
Y_3 \ar[r] &
Y_2 \ar[r] &
Y_1
}
$$
给定满足 $m' \geq m$ 的递增映射
$m' : \mathbf{N} \to \mathbf{N}$，并令
$a'_i : X_{m'(i)} \to X_{m(i)} \to Y_i$，则序对
$(m', a')$ 定义同一个 pro-系态射。反过来，给定如上的两个序对
$(m_1, a_1)$ 与 $(m_1, a_2)$，它们定义同一个 pro-对象态射，
当且仅当存在 $m' \geq m_1, m_2$，使得 $a'_1 = a'_2$。
\end{example}

\begin{remark}
\label{categories-remark-pro-category-copresheaves}
设 $\mathcal{C}$ 为范畴，并设 $F : \mathcal{I} \to \mathcal{C}$、
$G : \mathcal{J} \to \mathcal{C}$ 为 $\mathcal{C}$ 中的余滤过图表。
考虑由下式定义的函子 $A, B : \mathcal{C} \to \textit{Sets}$：
$$
A(X) = \colim_i \Mor_\mathcal{C}(F(i), X)
\quad\text{且}\quad
B(X) = \colim_j \Mor_\mathcal{C}(G(j), X)
$$
我们断言，从 $F$ 到 $G$ 的 pro-系态射等同于函子变换
$t : B \to A$。具体地，给定 $t$，可把 $t$ 作用于
$B(G(j))$ 中 $\text{id}_{G(j)}$ 的类，得到元素的相容系
$\xi_j \in A(G(j)) = \colim_i \Mor_\mathcal{C}(F(i), G(j))$；
这恰是注 \ref{categories-remark-pro-category} 中
$\text{Pro-}\mathcal{C}$ 的态射定义。给定一个从 $F$ 到 $G$
的 pro-对象态射，如何构造变换 $B \to A$，这里从略。
\end{remark}

\begin{lemma}
\label{categories-lemma-image-essentially-constant}
设 $\mathcal{C}$ 为范畴，并设 $M : \mathcal{I} \to \mathcal{C}$
为指标范畴 $\mathcal{I}$ 滤过（相应地，余滤过）的图表。
设 $F : \mathcal{C} \to \mathcal{D}$ 为函子。
若 $M$ 作为 ind-对象（相应地，pro-对象）本质常值，
则 $F \circ M : \mathcal{I} \to \mathcal{D}$ 亦然。
\end{lemma}

\begin{proof}
若 $X$ 是 $M$ 的一个值，则由定义立即可知，$F(X)$ 是
$F \circ M$ 的一个值。
\end{proof}

\begin{lemma}
\label{categories-lemma-characterize-essentially-constant-ind}
设 $\mathcal{C}$ 为范畴，并设 $M : \mathcal{I} \to \mathcal{C}$
为指标范畴 $\mathcal{I}$ 滤过的图表。下列条件等价：
\begin{enumerate}
\item $M$ 是本质常值 ind-对象；
\item 存在余锥 $(X, \{M_i \to X\}_i)$，使得对
$\mathcal{C}$ 中任意 $W$，映射
$\colim_i \Mor_\mathcal{C}(W, M_i) \to \Mor_\mathcal{C}(W, X)$
为双射；
\item $X = \colim_i M_i$ 存在，并且对 $\mathcal{C}$ 中任意 $W$，
映射 $\colim_i \Mor_\mathcal{C}(W, M_i) \to \Mor_\mathcal{C}(W, X)$
为双射；
\item 存在 $\mathcal{I}$ 中的 $i$ 以及态射 $X \to M_i$，
使得对 $\mathcal{C}$ 中任意 $W$，映射 $\Mor_\mathcal{C}(W, X) \to
\colim_{j \in \mathcal{I}} \Mor_\mathcal{C}(W, M_j)$
为双射。
\end{enumerate}
在情形 (2)、(3)、(4) 中，本质常值系的值为 $X$。
\end{lemma}

\begin{proof}
显然 (3) 蕴含 (2)。假设 (2)。于是
$\text{id}_X \in \Mor_\mathcal{C}(X, X)$ 来自某个 $i$ 所对应的态射
$X \to M_i$，即 $X \to M_i \to X$ 是恒等态射。
于是对所有 $W$，两个映射
$$
\Mor_\mathcal{C}(W, X)
\longrightarrow
\colim_i \Mor_\mathcal{C}(W, M_i)
\longrightarrow
\Mor_\mathcal{C}(W, X)
$$
都是双射；其中第一个映射由上面找到的态射 $X \to M_i$ 诱导，
而二者的复合是恒等映射。这意味着复合
$$
\colim_i \Mor_\mathcal{C}(W, M_i)
\longrightarrow
\Mor_\mathcal{C}(W, X)
\longrightarrow
\colim_i \Mor_\mathcal{C}(W, M_i)
$$
也为恒等映射。令 $W = M_j$，并从余极限中的
$\text{id}_{M_j}$ 出发，可知对某个充分大的 $k$，
$M_j \to X \to M_i \to M_k$ 等于 $M_j \to M_k$。
这证明 (1) 成立。

\medskip\noindent
假设 (4)。设 $k$ 为 $\mathcal{I}$ 的对象。令 $W = M_k$，
可推出存在唯一态射 $M_k \to X$，并且存在 $j$ 及
$\mathcal{I}$ 中的态射 $k \to j$、$i \to j$，使
$M_k \to X \to M_i \to M_j$ 等于 $M_k \to M_j$。
唯一性保证得到余锥 $(X, \{M_k \to X\})$。
由此可见 (4) 蕴含 (2)；略去若干细节。

\medskip\noindent
我们略去 (1) 蕴含其他条件的证明。
\end{proof}

\begin{lemma}
\label{categories-lemma-characterize-essentially-constant-pro}
设 $\mathcal{C}$ 为范畴，并设 $M : \mathcal{I} \to \mathcal{C}$
为指标范畴 $\mathcal{I}$ 余滤过的图表。下列条件等价：
\begin{enumerate}
\item $M$ 是本质常值 pro-对象；
\item 存在锥 $(X, \{X \to M_i\})$，使得对
$\mathcal{C}$ 中任意 $W$，映射
$\colim_{i \in \mathcal{I}^{opp}} \Mor_\mathcal{C}(M_i, W) \to
\Mor_\mathcal{C}(X, W)$ 为双射；
\item $X = \lim_i M_i$ 存在，并且对 $\mathcal{C}$ 中任意 $W$，
映射
$\colim_{i \in \mathcal{I}^{opp}} \Mor_\mathcal{C}(M_i, W) \to
\Mor_\mathcal{C}(X, W)$ 为双射；
\item 存在 $\mathcal{I}$ 中的 $i$ 以及态射 $M_i \to X$，
使得对 $\mathcal{C}$ 中任意 $W$，映射 $\Mor_\mathcal{C}(X, W) \to
\colim_{j \in \mathcal{I}^{opp}} \Mor_\mathcal{C}(M_j, W)$ 为双射。
\end{enumerate}
在情形 (2)、(3)、(4) 中，本质常值系的值为 $X$。
\end{lemma}

\begin{proof}
本引理与引理 \ref{categories-lemma-characterize-essentially-constant-ind} 对偶。
\end{proof}

\begin{lemma}
\label{categories-lemma-cofinal-essentially-constant}
设 $\mathcal{C}$ 为范畴，$H : \mathcal{I} \to \mathcal{J}$
为滤过指标范畴之间的函子。若 $H$ 共尾，则任意图表
$M : \mathcal{J} \to \mathcal{C}$ 本质常值，当且仅当
$M \circ H$ 本质常值。
\end{lemma}

\begin{proof}
这由引理 \ref{categories-lemma-characterize-essentially-constant-ind}
与引理 \ref{categories-lemma-cofinal} 形式地推出。
\end{proof}

\begin{lemma}
\label{categories-lemma-essentially-constant-over-product}
设 $\mathcal{I}$ 与 $\mathcal{J}$ 为滤过范畴，并以
$p : \mathcal{I} \times \mathcal{J} \to \mathcal{J}$ 表示投影。
则 $\mathcal{I} \times \mathcal{J}$ 滤过，而且图表
$M : \mathcal{J} \to \mathcal{C}$ 本质常值，当且仅当
$M \circ p : \mathcal{I} \times \mathcal{J} \to \mathcal{C}$
本质常值。
\end{lemma}

\begin{proof}
我们略去 $\mathcal{I} \times \mathcal{J}$ 滤过的验证。
由于 $p$ 共尾（验证略），所述等价性由引理
\ref{categories-lemma-cofinal-essentially-constant} 得到。
\end{proof}

\begin{lemma}
\label{categories-lemma-initial-essentially-constant}
设 $\mathcal{C}$ 为范畴，$H : \mathcal{I} \to \mathcal{J}$
为余滤过指标范畴之间的函子。若 $H$ 是始的，则任意图表
$M : \mathcal{J} \to \mathcal{C}$ 本质常值，当且仅当
$M \circ H$ 本质常值。
\end{lemma}

\begin{proof}
这由引理 \ref{categories-lemma-characterize-essentially-constant-pro}、
\ref{categories-lemma-initial}、\ref{categories-lemma-cofinal} 以及下述事实形式地推出：
若 $\mathcal{I}$ 在 $\mathcal{J}$ 中是始的，则
$\mathcal{I}^{opp}$ 在 $\mathcal{J}^{opp}$ 中共尾。
\end{proof}





\section{正合函子}
\label{categories-section-exact-functor}

\noindent
本节定义正合函子。

\begin{definition}
\label{categories-definition-exact}
设 $F : \mathcal{A} \to \mathcal{B}$ 为函子。
\begin{enumerate}
\item 假设 $\mathcal{A}$ 中所有有限极限都存在。
若 $F$ 与所有有限极限交换，则称其是{\it 左正合}的。
\item 假设 $\mathcal{A}$ 中所有有限余极限都存在。
若 $F$ 与所有有限余极限交换，则称其是{\it 右正合}的。
\item 若 $F$ 同时左正合且右正合，则称其是{\it 正合}的。
\end{enumerate}
\end{definition}

\begin{lemma}
\label{categories-lemma-characterize-left-exact}
设 $F : \mathcal{A} \to \mathcal{B}$ 为函子。假设 $\mathcal{A}$ 中
所有有限极限都存在，见引理 \ref{categories-lemma-finite-limits-exist}。
下列条件等价：
\begin{enumerate}
\item $F$ 左正合；
\item $F$ 与有限积和等化子交换；
\item $F$ 把 $\mathcal{A}$ 的终对象变为 $\mathcal{B}$ 的终对象，
且与纤维积交换。
\end{enumerate}
\end{lemma}

\begin{proof}
引理 \ref{categories-lemma-limits-products-equalizers} 表明 (2) 蕴含 (1)。
假设 (3) 成立。终对象上的纤维积就是积。若
$a, b : A \to B$ 是 $\mathcal{A}$ 的态射，则 $a, b$ 的等化子为
$$
(A \times_{a, B, b} A)\times_{(\text{pr}_1, \text{pr}_2), A \times A, \Delta} A.
$$
故 (3) 蕴含 (2)。最后，(1) 蕴含 (3)，因为空极限是终对象，
而纤维积是极限。
\end{proof}

\begin{lemma}
\label{categories-lemma-characterize-right-exact}
设 $F : \mathcal{A} \to \mathcal{B}$ 为函子。假设 $\mathcal{A}$ 中
所有有限余极限都存在，见引理 \ref{categories-lemma-colimits-exist}。
下列条件等价：
\begin{enumerate}
\item $F$ 右正合；
\item $F$ 与有限余积和余等化子交换；
\item $F$ 把 $\mathcal{A}$ 的始对象变为 $\mathcal{B}$ 的始对象，
且与推出交换。
\end{enumerate}
\end{lemma}

\begin{proof}
这是引理 \ref{categories-lemma-characterize-left-exact} 的对偶。
\end{proof}




\section{伴随函子}
\label{categories-section-adjoint}

\begin{definition}
\label{categories-definition-adjoint}
设 $\mathcal{C}$、$\mathcal{D}$ 为范畴，并设
$u : \mathcal{C} \to \mathcal{D}$ 与
$v : \mathcal{D} \to \mathcal{C}$ 为函子。若存在双射
$$
\Mor_\mathcal{D}(u(X), Y)
\longrightarrow
\Mor_\mathcal{C}(X, v(Y))
$$
且它关于 $X \in \Ob(\mathcal{C})$ 与
$Y \in \Ob(\mathcal{D})$ 具有函子性，则称 $u$ 是 $v$ 的
{\it 左伴随}，或称 $v$ 是 $u$ 的{\it 右伴随}。
\end{definition}

\noindent
换言之，这表示给定了从 $\Mor_\mathcal{D}(u(-), -)$ 到
$\Mor_\mathcal{C}(-, v(-))$ 的函子同构；这些函子定义在
$\mathcal{C}^{opp} \times \mathcal{D} \to \textit{Sets}$ 上。
对 $\mathcal{C}$ 的任意对象 $X$，由 $\text{id}_{u(X)}$ 对应得到态射
$X \to v(u(X))$。类似地，对 $\mathcal{D}$ 的任意对象 $Y$，
由 $\text{id}_{v(Y)}$ 对应得到态射 $u(v(Y)) \to Y$。
这些映射称为{\it 伴随映射}。伴随映射关于 $X$ 与 $Y$ 具有函子性，
因而得到函子态射
$$
\eta : \text{id}_\mathcal{C} \to v \circ u\quad (\text{单位})
\quad\text{且}\quad
\epsilon : u \circ v \to \text{id}_\mathcal{D}\quad (\text{余单位}).
$$
此外，若 $\alpha : u(X) \to Y$ 与 $\beta : X \to v(Y)$ 为态射，
则下列条件等价：
\begin{enumerate}
\item $\alpha$ 与 $\beta$ 通过定义中的双射相互对应；
\item $\beta$ 是复合 $X \to v(u(X)) \xrightarrow{v(\alpha)} v(Y)$；
\item $\alpha$ 是复合 $u(X) \xrightarrow{u(\beta)} u(v(Y)) \to Y$。
\end{enumerate}
这样便可用伴随映射重新表述伴随函子的概念。

\begin{lemma}
\label{categories-lemma-adjoint-exists}
设 $u : \mathcal{C} \to \mathcal{D}$ 为范畴之间的函子。
若对每个 $y \in \Ob(\mathcal{D})$，函子
$x \mapsto \Mor_\mathcal{D}(u(x), y)$ 都可表，则 $u$ 有右伴随。
\end{lemma}

\begin{proof}
对每个 $y$，选择对象 $v(y)$ 以及函子同构
$\Mor_\mathcal{C}(-, v(y)) \to \Mor_\mathcal{D}(u(-), y)$。
由 Yoneda 引理（引理 \ref{categories-lemma-yoneda}），对任意态射
$g : y \to y'$，函子变换
$$
\Mor_\mathcal{C}(-, v(y)) \to \Mor_\mathcal{D}(u(-), y) \to
\Mor_\mathcal{D}(u(-), y') \to \Mor_\mathcal{C}(-, v(y'))
$$
对应于唯一态射 $v(g) : v(y) \to v(y')$。我们略去 $v$ 为函子且
为 $u$ 的右伴随的验证。
\end{proof}

\begin{lemma}
\label{categories-lemma-left-adjoint-composed-fully-faithful}
\begin{reference}
Bhargav Bhatt，私人通信。
\end{reference}
设 $u$ 是定义 \ref{categories-definition-adjoint} 中 $v$ 的左伴随。
\begin{enumerate}
\item 若 $v \circ u$ 全忠实，则 $u$ 全忠实。
\item 若 $u \circ v$ 全忠实，则 $v$ 全忠实。
\end{enumerate}
\end{lemma}

\begin{proof}
证明 (2)。假设 $u \circ v$ 全忠实。设 $X$、$Y$ 是
$\mathcal{D}$ 中的对象。则自然的复合映射
$$
\Mor(X,Y) \to \Mor(v(X),v(Y)) \to \Mor(u(v(X)), u(v(Y)))
$$
为双射，所以 $v$ 至少忠实。为证明全忠实性，只需证明上面的
第二个映射为单射。但 $u$ 与 $v$ 之间的伴随说明
$$
\Mor(v(X), v(Y)) \to \Mor(u(v(X)), u(v(Y))) \to \Mor(u(v(X)), Y)
$$
为双射；其中第一个映射是自然映射，第二个映射来自伴随的余单位
$u(v(Y)) \to Y$。因此
$\Mor(v(X), v(Y)) \to \Mor(u(v(X)), u(v(Y)))$
也为单射，正如所需。(1) 的证明与此对偶。
\end{proof}

\begin{lemma}
\label{categories-lemma-adjoint-fully-faithful}
设 $u$ 是定义 \ref{categories-definition-adjoint} 中 $v$ 的左伴随。则
\begin{enumerate}
\item $u$ 全忠实 $\Leftrightarrow$ $\text{id} \cong v \circ u$
$\Leftrightarrow$ $\eta : \text{id} \to v \circ u$ 是同构；
\item $v$ 全忠实 $\Leftrightarrow$
$u \circ v \cong \text{id}$ $\Leftrightarrow$
$\epsilon : u \circ v \to \text{id}$ 是同构。
\end{enumerate}
\end{lemma}

\begin{proof}
证明 (1)。假设 $u$ 全忠实。我们证明
$\eta_X : X \to v(u(X))$ 是同构。任取态射
$X' \to v(u(X))$。由伴随性，它对应于态射 $u(X') \to u(X)$；
由 $u$ 的全忠实性，它又对应于唯一态射 $X' \to X$。
因此，与 $\eta_X$ 后复合给出双射
$\Mor(X', X) \to \Mor(X', v(u(X)))$，故 $\eta_X$ 为同构。
若存在函子同构 $\text{id} \cong v \circ u$，则 $v \circ u$ 全忠实。
由引理 \ref{categories-lemma-left-adjoint-composed-fully-faithful}，$u$ 全忠实；
由上述论证，这蕴含 $\eta$ 为同构。因此 $3$ 个条件等价
（这些条件也等价于 $v \circ u$ 全忠实）。

\medskip\noindent
(2) 与 (1) 对偶。
\end{proof}

\begin{lemma}
\label{categories-lemma-adjoint-exact}
设 $u$ 是定义 \ref{categories-definition-adjoint} 中 $v$ 的左伴随。
\begin{enumerate}
\item 假设 $M : \mathcal{I} \to \mathcal{C}$ 为图表，且
$\colim_\mathcal{I} M$ 在 $\mathcal{C}$ 中存在。则
$u(\colim_\mathcal{I} M) = \colim_\mathcal{I} u \circ M$。
换言之，$u$ 与（可表的）余极限交换。
\item 假设 $M : \mathcal{I} \to \mathcal{D}$ 为图表，且
$\lim_\mathcal{I} M$ 在 $\mathcal{D}$ 中存在。则
$v(\lim_\mathcal{I} M) = \lim_\mathcal{I} v \circ M$。
换言之，$v$ 与可表极限交换。
\end{enumerate}
\end{lemma}

\begin{proof}
从余极限到一个对象的态射，等价于从该极限各组成对象到该对象的
相容态射系，见注 \ref{categories-remark-limit-colim}。因此
$$
\begin{matrix}
\Mor_\mathcal{D}(u(\colim_{i \in \mathcal{I}} M_i), Y) &
= & \Mor_\mathcal{C}(\colim_{i \in \mathcal{I}} M_i, v(Y)) \\
& = &
\lim_{i \in \mathcal{I}^{opp}}
\Mor_\mathcal{C}(M_i, v(Y)) \\
& = &
\lim_{i \in \mathcal{I}^{opp}}
\Mor_\mathcal{D}(u(M_i), Y)
\end{matrix}
$$
证明了 $u(\colim_{i \in \mathcal{I}} M_i)$ 就是所求余极限。
另一断言可类似证明。
\end{proof}

\begin{lemma}
\label{categories-lemma-exact-adjoint}
设 $u$ 是定义 \ref{categories-definition-adjoint} 中 $v$ 的左伴随。
\begin{enumerate}
\item 若 $\mathcal{C}$ 有有限余极限，则 $u$ 右正合。
\item 若 $\mathcal{D}$ 有有限极限，则 $v$ 左正合。
\end{enumerate}
\end{lemma}

\begin{proof}
这由定义与引理 \ref{categories-lemma-adjoint-exact} 立即得到。
\end{proof}

\begin{lemma}
\label{categories-lemma-unit-counit-relations}
设 $u : \mathcal{C} \to \mathcal{D}$ 是函子
$v : \mathcal{D} \to \mathcal{C}$ 的左伴随。设
$\eta_X : X \to v(u(X))$ 为单位，
$\epsilon_Y : u(v(Y)) \to Y$ 为余单位。则
$$
u(X) \xrightarrow{u(\eta_X)} u(v(u(X))
\xrightarrow{\epsilon_{u(X)}} u(X)
\quad\text{且}\quad
v(Y) \xrightarrow{\eta_{v(Y)}} v(u(v(Y))) \xrightarrow{v(\epsilon_Y)}
v(Y)
$$
都是恒等态射。
\end{lemma}

\begin{proof}
略。
\end{proof}

\begin{lemma}
\label{categories-lemma-transformation-between-functors-and-adjoints}
设 $u_1, u_2 : \mathcal{C} \to \mathcal{D}$ 为具有右伴随
$v_1, v_2 : \mathcal{D} \to \mathcal{C}$ 的函子。
设 $\beta : u_2 \to u_1$ 为函子变换，并设
$\beta^\vee : v_1 \to v_2$ 为相应的伴随函子变换。则
$$
\xymatrix{
u_2 \circ v_1 \ar[r]_\beta \ar[d]_{\beta^\vee} &
u_1 \circ v_1 \ar[d] \\
u_2 \circ v_2 \ar[r] & \text{id}
}
$$
可交换；其中未标记的箭头是余单位变换。
\end{lemma}

\begin{proof}
这是因为 $\beta^\vee_D : v_1D \to v_2D$ 是唯一满足下述条件的态射：
由它诱导的映射 $\Mor(C, v_1D) \to \Mor(C, v_2D)$，正是由
$\beta_C : u_2C \to u_1C$ 诱导的映射
$\Mor(u_1C, D) \to \Mor(u_2C, D)$。也就是说，由
$\beta_{v_1 D}$ 诱导的映射
$$
\Mor(u_1 v_1 D, D') \to \Mor(u_2 v_1 D, D')
$$
与由 $\beta^\vee_{D'}$ 诱导的映射
$$
\Mor(v_1 D, v_1 D') \to \Mor(v_1 D, v_2 D')
$$
相同。取 $D' = D$，可知余单位 $u_1 v_1 D \to D$ 与
$\beta_{v_1D}$ 前复合后，在伴随下对应于 $\beta^\vee_D$。
这恰好说明该图表在 $D$ 上取值后可交换。
\end{proof}

\begin{lemma}
\label{categories-lemma-compose-counits}
设 $\mathcal{A}$、$\mathcal{B}$、$\mathcal{C}$ 为范畴。设
$v : \mathcal{A} \to \mathcal{B}$ 与
$v' : \mathcal{B} \to \mathcal{C}$ 为函子，其左伴随分别为
$u$ 与 $u'$。则
\begin{enumerate}
\item 函子 $v'' = v' \circ v$ 的左伴随为
$u'' = u \circ u'$。
\item 给定 $\mathcal{A}$ 中的 $X$，有
\begin{equation}
\label{categories-equation-compose-counits}
\epsilon_X^v \circ u(\epsilon^{v'}_{v(X)}) = \epsilon^{v''}_X :
u''(v''(X)) \to X
\end{equation}
其中 $\epsilon$ 是各伴随的余单位。
\end{enumerate}
\end{lemma}

\begin{proof}
我们展开 (2) 中的公式，因为这也会立即证明 (1)。首先，伴随对
$(u, v)$ 与 $(u', v')$ 的余单位分别为映射
$\epsilon_X^v : u(v(X)) \to X$ 与
$\epsilon_Y^{v'} : u'(v'(Y)) \to Y$，见定义
\ref{categories-definition-adjoint} 后的讨论。以 (1) 中的 $u''$ 与 $v''$ 展开一切，
$$
u''(v''(X)) = u(u'(v'(v(X)))) \xrightarrow{u(\epsilon_{v(X)}^{v'})}
u(v(X)) \xrightarrow{\epsilon_X^v} X
$$
便得到 (\ref{categories-equation-compose-counits}) 左边的映射。
暂将其记为 $\epsilon_X^{v''}$。要说明它是伴随对
$(u'', v'')$ 的余单位，须证明：给定 $\mathcal{C}$ 中的 $Z$，
把态射 $\beta : Z \to v''(X)$ 映为
$\alpha = \epsilon_X^{v''} \circ u''(\beta) : u''(Z) \to X$
的规则在态射集合上为双射。这确实成立，因为该规则是下述两个规则的复合：
先把 $\beta$ 映为 $\epsilon_{v(X)}^{v'} \circ u'(\beta)$，
这由对 $(u', v')$ 的假设是双射；再把它映为
$\epsilon_X^v \circ u(\epsilon_{v(X)}^{v'} \circ u'(\beta))$，
这由对 $(u, v)$ 的假设也是双射。
\end{proof}






\section{一个可表性判据}
\label{categories-section-representable}

\noindent
下列引理常用于证明大范畴中泛对象的存在性；请参见注
\ref{categories-remark-how-to-use-it} 中的讨论。

\begin{lemma}
\label{categories-lemma-a-version-of-brown}
设 $\mathcal{C}$ 为有极限的大\footnote{见注
\ref{categories-remark-big-categories}。}范畴，并设
$F : \mathcal{C} \to \textit{Sets}$ 为函子。假设：
\begin{enumerate}
\item $F$ 与极限交换；
\item 存在 $\mathcal{C}$ 的一族对象 $\{x_i\}_{i \in I}$，且对每个
$i \in I$ 有元素 $f_i \in F(x_i)$，使得对
$y \in \Ob(\mathcal{C})$ 与 $g \in F(y)$，都存在 $i$ 及态射
$\varphi : x_i \to y$，满足 $F(\varphi)(f_i) = g$。
\end{enumerate}
则 $F$ 可表；也就是说，存在 $\mathcal{C}$ 的对象 $x$，使得
$$
F(y) = \Mor_\mathcal{C}(x, y)
$$
关于 $y$ 具有函子性。
\end{lemma}

\begin{proof}
令 $\mathcal{I}$ 为以下范畴：其对象是偶对 $(x_i, f_i)$，
而态射 $(x_i, f_i) \to (x_{i'}, f_{i'})$ 是
$\mathcal{C}$ 中满足 $F(\varphi)(f_i) = f_{i'}$ 的映射
$\varphi : x_i \to x_{i'}$。令
$$
x = \lim_{(x_i, f_i) \in \mathcal{I}} x_i
$$
（这并不是我们要找的 $x$，见下文）。由假设，该极限存在。
由于 $F$ 与极限交换，有
$$
F(x) = \lim_{(x_i, f_i) \in \mathcal{I}} F(x_i).
$$
因此存在泛元素 $f \in F(x)$，在对投影映射 $x \to x_i$
应用 $F$ 后，它映到 $f_i \in F(x_i)$。利用 $f$，得到函子变换
$$
\xi : \Mor_\mathcal{C}(x, - ) \longrightarrow F(-)
$$
（见第 \ref{categories-section-opposite} 节）。任取 $\mathcal{C}$ 的对象 $y$
及 $g \in F(y)$。由假设可选取 $x_i \to y$，使 $f_i$ 映到 $g$。
于是对映射
$$
x \longrightarrow x_i \longrightarrow y
$$
应用 $F$（第一个映射是定义 $x$ 的极限投影），会把 $f$ 映到 $g$。
故变换 $\xi$ 为满射。

\medskip\noindent
为了找出表示 $F$ 的对象，令 $e : x' \to x$ 为所有满足
$F(\varphi)(f) = f$ 的自映射 $\varphi : x \to x$ 的等化子。
由于 $F$ 与极限交换，它也与等化子交换，因而存在
$f' \in F(x')$ 映到 $F(x)$ 中的 $f$。由于 $\xi$ 为满射且
$f'$ 映到 $f$，可知 $\xi' : \Mor_\mathcal{C}(x', -) \to F(-)$
也为满射。最后，假设 $a, b : x' \to y$ 是两个满足
$F(a)(f') = F(b)(f')$ 的映射。我们须证明 $a = b$。
考虑等化子 $e' : x'' \to x'$。同样可找到映到 $f'$ 的
$f'' \in F(x'')$。选取映射 $\psi : x \to x''$，使
$F(\psi)(f) = f''$。于是
$e \circ e' \circ \psi : x \to x$ 是满足
$F(e \circ e' \circ \psi)(f) = f$ 的态射。因此
$e \circ e' \circ \psi \circ e = e$。由于 $e$ 是单态射，
这蕴含 $e'$ 是满态射，从而如所需有 $a = b$。
\end{proof}

\begin{remark}
\label{categories-remark-how-to-use-it}
上面的引理常用于构造某物上的自由某物，例如集合上的自由阿贝尔群、
集合上的自由群等等。以集合 $E$ 上的自由群为例，其想法是考虑函子
$$
F : \textit{Groups} \to \textit{Sets},\quad
G \longmapsto \text{Map}(E, G)
$$
该函子与极限交换。作为对象族，可以取由群 $G_i$ 及集合映射
$E \to G_i$ 组成的一族 $E \to G_i$，其中这些群的基数至多为
$\max(\aleph_0, |E|)$，并使这种结构的每个同构类至少出现一次。
事实上，若 $E \to G$ 是从 $E$ 到群 $G$ 的映射，则由其像生成的
子群 $G'$ 的基数至多为 $\max(\aleph_0, |E|)$。该引理说明函子可表，
故存在群 $F_E$，使
$\Mor_{\textit{Groups}}(F_E, G) = \text{Map}(E, G)$。
特别地，$F_E$ 的恒等态射对应于一个映射 $E \to F_E$，并且可证明
$F_E$ 由其像生成而不施加任何关系。

\medskip\noindent
另一种典型应用是：一旦已知极限存在，就可用该引理构造余极限。
我们以拓扑空间范畴为例；由拓扑，引理 \ref{topology-lemma-limits}，
该范畴有极限。具体地，假设
$\mathcal{I} \to \textit{Top}$，$i \mapsto X_i$ 是函子。
则可考虑
$$
F : \textit{Top} \longrightarrow \textit{Sets},\quad
Y \longmapsto \lim_\mathcal{I} \Mor_{\textit{Top}}(X_i, Y)
$$
该函子与极限交换。此外，给定任意拓扑空间 $Y$ 及
$F(Y)$ 的元素 $(\varphi_i : X_i \to Y)$，存在基数至多为
$|\coprod X_i|$ 的子空间 $Y' \subset Y$，使态射 $\varphi_i$
都映入 $Y'$。事实上，可在各 $\varphi_i$ 的像之并上取诱导拓扑。
所以显然满足引理的假设，从而得到表示函子 $F$ 的拓扑空间 $X$；
这恰好意味着 $X$ 是图表 $i \mapsto X_i$ 的余极限。
\end{remark}

\begin{theorem}[伴随函子定理]
\label{categories-theorem-adjoint-functor}
设 $G : \mathcal{C} \to \mathcal{D}$ 为大范畴之间的函子。
假设 $\mathcal{C}$ 有极限，$G$ 与极限交换，并且对
$\mathcal{D}$ 的每个对象 $y$，存在一组偶对
$(x_i, f_i)_{i \in I}$，其中 $x_i \in \Ob(\mathcal{C})$、
$f_i \in \Mor_\mathcal{D}(y, G(x_i))$；它满足：对任意偶对
$(x, f)$，其中 $x \in \Ob(\mathcal{C})$、
$f \in \Mor_\mathcal{D}(y, G(x))$，都存在 $i$ 及态射
$h : x_i \to x$，使 $f = G(h) \circ f_i$。则 $G$ 有左伴随 $F$。
\end{theorem}

\begin{proof}
这些假设蕴含：对 $\mathcal{D}$ 的每个对象 $y$，函子
$x \mapsto \Mor_\mathcal{D}(y, G(x))$ 都满足引理
\ref{categories-lemma-a-version-of-brown} 的假设。因此它由某个对象表示，
将该对象记为 $F(y)$。应用 Yoneda 引理（引理 \ref{categories-lemma-yoneda}），
规则 $y \mapsto F(y)$ 变成一个函子；由构造，它是 $G$ 的伴随。
细节略去。
\end{proof}





\section{范畴紧致对象}
\label{categories-section-compact}

\noindent
下面略述范畴中的“小”对象。

\begin{definition}
\label{categories-definition-compact-object}
设 $\mathcal{C}$ 为大\footnote{见注 \ref{categories-remark-big-categories}。}范畴。
若对每个使 $\colim_i M_i$ 存在的滤过图表
$M : \mathcal{I} \to \mathcal{C}$，都有
$$
\Mor_\mathcal{C}(X, \colim_i M_i) =
\colim_i \Mor_\mathcal{C}(X, M_i)
$$
则称 $\mathcal{C}$ 的对象 $X$ 是{\it 范畴紧致}的。
\end{definition}

\noindent
通常只在 $\mathcal{C}$ 有所有滤过余极限的假设下作此定义。

\begin{lemma}
\label{categories-lemma-extend-functor-by-colim}
设 $\mathcal{C}$ 与 $\mathcal{D}$ 为具有滤过余极限的大范畴。
设 $\mathcal{C}' \subset \mathcal{C}$ 为由 $\mathcal{C}$ 的范畴紧致对象
组成的小满子范畴，并且 $\mathcal{C}$ 的每个对象都是
$\mathcal{C}'$ 中对象的滤过余极限。则每个函子
$F' : \mathcal{C}' \to \mathcal{D}$ 都有唯一扩张
$F : \mathcal{C} \to \mathcal{D}$，它与滤过余极限交换。
\end{lemma}

\begin{proof}
对 $\mathcal{C}$ 的每个对象 $X$，可将 $X$ 写成滤过余极限
$X = \colim X_i$，其中 $X_i \in \Ob(\mathcal{C}')$。令
$$
F(X) = \colim F'(X_i)
$$
于 $\mathcal{D}$ 中。下文将证明此构造不依赖于 $X$ 的余极限表示。

\medskip\noindent
设给定 $\mathcal{C}$ 的态射 $\alpha : X \to Y$，且
$X = \colim_{i \in I} X_i$ 与 $Y = \colim_{j \in J} Y_i$
写成 $\mathcal{C}'$ 中对象的滤过余极限。对每个 $i \in I$，
由于 $X_i$ 是 $\mathcal{C}$ 的范畴紧致对象，可找到 $j \in J$
以及交换图表
$$
\xymatrix{
X_i \ar[r] \ar[d] & X \ar[d]^\alpha \\
Y_j \ar[r] & Y
}
$$
于是得到态射 $F'(X_i) \to F'(Y_j) \to F(Y)$，其中第二个态射是到
$F(Y) = \colim F'(Y_j)$ 的余投影。箭头
$\beta_i : F'(X_i) \to F(Y)$ 不依赖于 $j$ 的选择。
对 $i \leq i'$，复合
$$
F'(X_i) \to F'(X_{i'}) \xrightarrow{\beta_{i'}} F(Y)
$$
等于 $\beta_i$。因此由余极限的泛性质得到良定义的箭头
$$
F(\alpha) : F(X) = \colim F(X_i) \to F(Y)
$$
若 $\alpha' : Y \to Z$ 是 $\mathcal{C}$ 的另一个态射，且
$Z = \colim Z_k$ 也写成 $\mathcal{C}'$ 中对象的滤过余极限，
则可作为一项愉快的练习证明：诱导态射
$F(\alpha) : F(X) \to F(Y)$ 与 $F(\alpha') : F(Y) \to F(Z)$
复合为态射 $F(\alpha' \circ \alpha)$。细节略去。

\medskip\noindent
特别地，若给定两个表示
$X = \colim X_i$ 与 $X = \colim X'_{i'}$，它们都是
$\mathcal{C}'$ 中系统的滤过余极限，则得到互逆箭头
$\colim F'(X_i) \to \colim F'(X'_{i'})$ 与
$\colim F'(X'_{i'}) \to \colim F'(X_i)$。换言之，值 $F(X)$
不依赖于将 $X$ 表示成 $\mathcal{C}'$ 中对象的滤过余极限的选择。
结合上一段所讨论的 $F$ 的函子性，可知 $F$ 是函子。
此外，若 $X \in \Ob(\mathcal{C}')$，显然 $F(X) = F'(X)$。

\medskip\noindent
只要证明 $F$ 与滤过余极限交换，引理中的唯一性断言便很清楚
（否则该断言没有意义）。为此，假设
$X = \colim_{\lambda \in \Lambda} X_\lambda$
是 $\mathcal{C}$ 中的滤过余极限。由于 $F$ 是函子，当然得到映射
$$
\colim_\lambda F(X_\lambda) \longrightarrow F(X)
$$
另一方面，把 $X = \colim X_i$ 写成 $\mathcal{C}'$ 中对象的
滤过余极限。如上，对每个 $i \in I$，可选择
$\lambda \in \Lambda$ 及交换图表
$$
\xymatrix{
X_i \ar[rr] \ar[rd] & & X_\lambda \ar[ld] \\
& X
}
$$
同样，这给出与过渡态射相容的良定义态射
$F'(X_i) \to \colim_\lambda F(X_\lambda)$，因而给出态射
$$
F(X) = \colim_i F'(X_i) \longrightarrow \colim_\lambda F(X_\lambda)
$$
此态射与上面的态射互逆（细节略），并如所需证明
$F(X) = \colim_\lambda F(X_\lambda)$。
\end{proof}







\section{范畴中的局部化}
\label{categories-section-localization}

\noindent
本节的基本思想是：给定范畴 $\mathcal{C}$ 和一组箭头 $S$，
构造函子 $F : \mathcal{C} \to S^{-1}\mathcal{C}$，使 $S$ 的所有元素
在 $S^{-1}\mathcal{C}$ 中均可逆，并且 $F$ 在所有具有此性质的函子中
是泛的。本节的参考文献为 \cite[Chapter I, Section 2]{GZ}
与 \cite[Chapter II, Section 2]{Verdier}。

\begin{definition}
\label{categories-definition-multiplicative-system}
设 $\mathcal{C}$ 为范畴。若 $\mathcal{C}$ 的一组箭头 $S$ 具有下列性质，
则称它为{\it 左乘法系}：
\begin{enumerate}
\item[LMS1] $\mathcal{C}$ 每个对象的恒等态射都属于 $S$，且
$S$ 中两个可复合元素的复合仍属于 $S$。
\item[LMS2] 每个实线图表
$$
\xymatrix{
X \ar[d]_t \ar[r]_g & Y \ar@{..>}[d]^s \\
Z \ar@{..>}[r]^f & W
}
$$
若满足 $t \in S$，都可补成虚线交换方形，并使 $s \in S$。
\item[LMS3] 对每一对态射 $f, g : X \to Y$ 以及靶为 $X$ 的
$t \in S$，若 $f \circ t = g \circ t$，则存在源为 $Y$ 的
$s \in S$，使 $s \circ f = s \circ g$。
\end{enumerate}
若 $\mathcal{C}$ 的一组箭头 $S$ 具有下列性质，则称它为
{\it 右乘法系}：
\begin{enumerate}
\item[RMS1] $\mathcal{C}$ 每个对象的恒等态射都属于 $S$，且
$S$ 中两个可复合元素的复合仍属于 $S$。
\item[RMS2] 每个实线图表
$$
\xymatrix{
X \ar@{..>}[d]_t \ar@{..>}[r]_g & Y \ar[d]^s \\
Z \ar[r]^f & W
}
$$
若满足 $s \in S$，都可补成虚线交换方形，并使 $t \in S$。
\item[RMS3] 对每一对态射 $f, g : X \to Y$ 以及源为 $Y$ 的
$s \in S$，若 $s \circ f = s \circ g$，则存在靶为 $X$ 的
$t \in S$，使 $f \circ t = g \circ t$。
\end{enumerate}
若 $\mathcal{C}$ 的一组箭头 $S$ 同时是左乘法系和右乘法系，
则称它为{\it 乘法系}。换言之，这表示 MS1、MS2、MS3 成立，其中
MS1 $=$ LMS1 $+$ RMS1，MS2 $=$ LMS2 $+$ RMS2，且
MS3 $=$ LMS3 $+$ RMS3。（当然，LMS1 $=$ RMS1 $=$ MS1。）
\end{definition}

\noindent
这些条件可用于如下构造范畴 $S^{-1}\mathcal{C}$。

\medskip\noindent
{\bf 左分式演算。}
设 $\mathcal{C}$ 为范畴，$S$ 为左乘法系。如下定义一个新范畴
$S^{-1}\mathcal{C}$（我们将在引理 \ref{categories-lemma-left-localization}
的证明中验证此定义确实成立）：
\begin{enumerate}
\item 令 $\Ob(S^{-1}\mathcal{C}) = \Ob(\mathcal{C})$。
\item $S^{-1}\mathcal{C}$ 中的态射 $X \to Y$ 由偶对
$(f : X \to Y', s : Y \to Y')$ 给出，其中 $s \in S$，
并模去等价关系。（该等价关系定义如下。可以把偶对 $(f, s)$
的等价类看作 $s^{-1}f : X \to Y$。）
\item 称两个偶对 $(f_1 : X \to Y_1, s_1 : Y \to Y_1)$ 与
$(f_2 : X \to Y_2, s_2 : Y \to Y_2)$ 等价，如果存在第三个偶对
$(f_3 : X \to Y_3, s_3 : Y \to Y_3)$ 以及 $\mathcal{C}$ 的态射
$u : Y_1 \to Y_3$ 与 $v : Y_2 \to Y_3$，使它们嵌入交换图表
$$
\xymatrix{
 & Y_1 \ar[d]^u & \\
X \ar[ru]^{f_1} \ar[r]^{f_3} \ar[rd]_{f_2} &
Y_3 &
Y \ar[lu]_{s_1} \ar[l]_{s_3} \ar[ld]^{s_2} \\
& Y_2 \ar[u]_v &
}
$$
\item 偶对 $(f : X \to Y', s : Y \to Y')$ 与
$(g : Y \to Z', t : Z \to Z')$ 的等价类之复合，定义为偶对
$(h \circ f : X \to Z'', u \circ t : Z \to Z'')$ 的等价类；
其中选择 $h$ 与 $u \in S$ 使其嵌入交换图表
$$
\xymatrix{
Y \ar[d]_s \ar[r]_g & Z' \ar[d]^u \\
Y' \ar[r]^h & Z''
}
$$
这样的选择由假设存在。
\item $S^{-1} \mathcal{C}$ 中的恒等态射 $X \to X$ 是偶对
$(\text{id} : X \to X,
\text{id} : X \to X)$ 的等价类。
\end{enumerate}

\begin{lemma}
\label{categories-lemma-left-localization}
设 $\mathcal{C}$ 为范畴，$S$ 为左乘法系。
\begin{enumerate}
\item 上面在偶对上定义的关系是等价关系。
\item 上面给出的复合规则在等价类上良定义。
\item 复合满足结合律（且恒等态射满足恒等公理），
所以 $S^{-1}\mathcal{C}$ 是范畴。
\end{enumerate}
\end{lemma}

\begin{proof}
证明 (1)。称两个偶对
$p_1 = (f_1 : X \to Y_1, s_1 : Y \to Y_1)$ 与
$p_2 = (f_2 : X \to Y_2, s_2 : Y \to Y_2)$ 初等等价，
如果存在 $\mathcal{C}$ 的态射 $a : Y_1 \to Y_2$，使
$a \circ f_1 = f_2$ 且 $a \circ s_1 = s_2$。图示为
$$
\xymatrix{
X \ar@{=}[d] \ar[r]_{f_1} & Y_1 \ar[d]^a & Y \ar[l]^{s_1} \ar@{=}[d] \\
X \ar[r]^{f_2} & Y_2 & Y \ar[l]_{s_2}
}
$$
把此性质记为 $p_1Ep_2$。注意 $pEp$，且
$aEb, bEc \Rightarrow aEc$。（尽管如此命名，$E$ 并不是等价关系。）
断言 (1) 声称关系
$p \sim p' \Leftrightarrow \exists q: pEq \wedge p'Eq$
是等价关系（这里 $q$ 应为满足与 $p$ 和 $p'$ 相同条件的偶对）。
使用上面 $E$ 的性质作简单的形式论证可知，只须证明
$p_3Ep_1, p_3Ep_2 \Rightarrow p_1 \sim p_2$。
因此，假设给定交换图表
$$
\xymatrix{
 & Y_1 & \\
X \ar[ru]^{f_1} \ar[r]^{f_3} \ar[rd]_{f_2} &
Y_3 \ar[u]_{a_{31}} \ar[d]^{a_{32}} &
Y \ar[lu]_{s_1} \ar[l]_{s_3} \ar[ld]^{s_2} \\
& Y_2 &
}
$$
其中 $s_i \in S$。首先应用 LMS2 得到交换图表
$$
\xymatrix{
Y \ar[d]_{s_1} \ar[r]_{s_2} & Y_2 \ar@{..>}[d]^{a_{24}} \\
Y_1 \ar@{..>}[r]^{a_{14}} & Y_4
}
$$
且 $a_{24} \in S$。于是
$$
a_{14} \circ a_{31} \circ s_3 =
a_{14} \circ s_1 =
a_{24} \circ s_2 =
a_{24} \circ a_{32} \circ s_3.
$$
故由 LMS3，存在态射 $s_{44} : Y_4 \to Y'_4$，使
$s_{44} \in S$ 且
$s_{44} \circ a_{14} \circ a_{31}
= s_{44} \circ a_{24} \circ a_{32}$。
因此，把 $Y_4$、$a_{14}$、$a_{24}$ 分别替换为 $Y'_4$、
$s_{44} \circ a_{14}$、$s_{44} \circ a_{24}$ 后，可以假设
$a_{14} \circ a_{31} = a_{24} \circ a_{32}$
（并且仍有 $a_{24} \in S$ 以及
$a_{14} \circ s_1 = a_{24} \circ s_2$）。令
$$
f_4 =
a_{14} \circ f_1 =
a_{14} \circ a_{31} \circ f_3 =
a_{24} \circ a_{32} \circ f_3 =
a_{24} \circ f_2
$$
以及 $s_4 = a_{14} \circ s_1 = a_{24} \circ s_2$。则图表
$$
\xymatrix{
X \ar@{=}[d] \ar[r]_{f_1} & Y_1 \ar[d]^{a_{14}} & Y \ar[l]^{s_1} \ar@{=}[d] \\
X \ar[r]^{f_4} & Y_4 & Y \ar[l]_{s_4}
}
$$
可交换，并有 $s_4 \in S$（由 LMS1）。所以 $p_1 E p_4$，
其中 $p_4 = (f_4, s_4)$。类似地，$p_2 E p_4$。合并二者得到
$p_1 \sim p_2$。

\medskip\noindent
证明 (2)。设 $p = (f : X \to Y', s : Y \to Y')$ 与
$q = (g : Y \to Z', t : Z \to Z')$ 是上面复合定义中的偶对。
为了复合，选择交换图表
$$
\xymatrix{
Y \ar[d]_s \ar[r]_g & Z' \ar[d]^{u_2} \\
Y' \ar[r]^{h_2} & Z_2
}
$$
其中 $u_2 \in S$。先证明偶对
$r_2 = (h_2 \circ f : X \to Z_2, u_2 \circ t : Z \to Z_2)$
的等价类不依赖于 $(Z_2, h_2, u_2)$ 的选择。事实上，假设
$(Z_3, h_3, u_3)$ 是另一选择，相应复合为
$r_3 = (h_3 \circ f : X \to Z_3, u_3 \circ t : Z \to Z_3)$。
由 LMS2，可选择图表
$$
\xymatrix{
Z' \ar[d]_{u_2} \ar[r]_{u_3} & Z_3 \ar[d]^{u_{34}} \\
Z_2 \ar[r]^{h_{24}} & Z_4
}
$$
且 $u_{34} \in S$。有 $h_2 \circ s = u_2 \circ g$，并且类似地
$h_3 \circ s = u_3 \circ g$。现在
$$
u_{34} \circ h_3 \circ s
= u_{34} \circ u_3 \circ g
= h_{24} \circ u_2 \circ g
= h_{24} \circ h_2 \circ s.
$$
所以 LMS3 表明存在 $Z'_4$ 以及 $s_{44} : Z_4 \to Z'_4$，使
$s_{44} \circ u_{34} \circ h_3 = s_{44} \circ h_{24} \circ h_2$。
把 $Z_4$、$h_{24}$、$u_{34}$ 分别替换为 $Z'_4$、
$s_{44} \circ h_{24}$、$s_{44} \circ u_{34}$，可假设
$u_{34} \circ h_3 = h_{24} \circ h_2$。同时，关系
$u_{34} \circ u_3 = h_{24} \circ u_2$ 与 $u_{34} \in S$ 仍成立。
现可令 $h_4 = u_{34} \circ h_3 = h_{24} \circ h_2$ 且
$u_4 = u_{34} \circ u_3 = h_{24} \circ u_2$。于是得到交换图表
$$
\xymatrix{
X \ar@{=}[d] \ar[r]_{h_2\circ f} &
Z_2 \ar[d]^{h_{24}} &
Z \ar[l]^{u_2 \circ t} \ar@{=}[d] \\
X \ar@{=}[d] \ar[r]^{h_4\circ f} &
Z_4 &
Z \ar@{=}[d] \ar[l]_{u_4 \circ t} \\
X \ar[r]^{h_3 \circ f} &
Z_3 \ar[u]^{u_{34}} &
Z \ar[l]_{u_3 \circ t}
}
$$
因此得到偶对
$r_4 =
(h_4 \circ f : X \to Z_4, u_4 \circ t : Z \to Z_4)$；
上图表明 $r_2Er_4$ 且 $r_3Er_4$，从而如所需有
$r_2 \sim r_3$。所以现在可以把 $p \circ q$ 定义为以上方式得到的
所有可能偶对 $r$ 的等价类。

\medskip\noindent
为完成 (2) 的证明，须证明：给定偶对 $p_1, p_2, q$，若
$p_1Ep_2$，则每当复合有意义时，
$p_1 \circ q = p_2 \circ q$ 且 $q \circ p_1 = q \circ p_2$。
为此，写
$p_1 = (f_1 : X \to Y_1, s_1 : Y \to Y_1)$、
$p_2 = (f_2 : X \to Y_2, s_2 : Y \to Y_2)$，并令
$a : Y_1 \to Y_2$ 为 $\mathcal{C}$ 的态射，使
$f_2 = a \circ f_1$ 且 $s_2 = a \circ s_1$。
先假设 $q = (g : Y \to Z', t : Z \to Z')$。
此时选择左侧所示的交换图表
$$
\vcenter{
\xymatrix{
Y \ar[d]_{s_2} \ar[r]^g & Z' \ar[d]^u \\
Y_2 \ar[r]^h & Z''
}
}
\quad
\Rightarrow
\quad
\vcenter{
\xymatrix{
Y \ar[d]_{s_1} \ar[r]^g & Z' \ar[d]^u \\
Y_1 \ar[r]^{h \circ a} & Z''
}
}
$$
（其中 $u \in S$）；它蕴含右侧图表也可交换。利用这些图表可见，
两个复合 $q \circ p_1$ 与 $q \circ p_2$ 都是
$(h \circ a \circ f_1 : X \to Z'', u \circ t : Z \to Z'')$
的等价类。因此 $q \circ p_1 = q \circ p_2$。
另一种情形的证明略去；此时须证明
$p_1 \circ q = p_2 \circ q$。（它与已证明的情形类似。）

\medskip\noindent
证明 (3)。须证明复合的结合律。考虑实线图表
$$
\xymatrix{
& & & Z \ar[d] \\
& & Y \ar[d] \ar[r] & Z' \ar@{..>}[d] \\
& X \ar[d] \ar[r] & Y' \ar@{..>}[d] \ar@{..>}[r] & Z'' \ar@{..>}[d] \\
W \ar[r] & X' \ar@{..>}[r] & Y'' \ar@{..>}[r] & Z'''
}
$$
（其中竖直箭头都属于 $S$），它给出三个可复合的偶对。
利用 LMS2，可选择虚线箭头使各方形可交换，并使竖直箭头属于 $S$。
于是显然，三个偶对的复合是偶对
$(W \to Z''', Z \to Z''')$ 的等价类；该偶对由复合底行的横箭头
与右列的竖箭头得到。

\medskip\noindent
恒等公理的验证留给读者。
\end{proof}

\begin{remark}
\label{categories-remark-motivation-localization}
构造 $S^{-1} \mathcal{C}$ 的动机，是通过人为地为 $S$ 中的态射创造逆，
“迫使”它们可逆（代价是某些已有态射可能彼此等同）。这类似于
交换环关于乘法子集的局部化，更一般地，也类似于非交换环关于右分母集
的局部化（见 \cite[Section 10A]{Lam}）。这不只是一种类比：
$S^{-1} \mathcal{C}$ 的构造（更确切地说，它对加性范畴
$\mathcal{C}$ 的版本）实际上推广了后一种局部化。具体地，
一个非交换环可视为只有一个对象的预加性范畴（态射即环的元素）；
该环的一个乘法子集遂成为满足 LMS1（即 RMS1）的一组态射 $S$。
于是，这个范畴与这个子集 $S$ 的条件 RMS2、RMS3，分别转化为
右分母集的两个条件（“右可置换”与“右可逆”）；LMS 与左分母集的
情形类似。而 $S^{-1} \mathcal{C}$（赋予适当定义的加性结构）
就是对应于环之局部化的单对象范畴。
\end{remark}

\begin{definition}
\label{categories-definition-left-localization-as-fraction}
设 $\mathcal{C}$ 为范畴，$S$ 为 $\mathcal{C}$ 的态射构成的左乘法系。
给定 $\mathcal{C}$ 中任意态射 $f : X \to Y'$ 及 $S$ 中任意态射
$s : Y \to Y'$，用 {\it $s^{-1} f$} 表示偶对
$(f : X \to Y', s : Y \to Y')$ 的等价类。这是
$S^{-1} \mathcal{C}$ 中从 $X$ 到 $Y$ 的态射。
\end{definition}

\noindent
这个记号富有启发性，而且它所暗示的性质确实成立：给定
$\mathcal{C}$ 中任意态射 $f : X \to Y'$ 以及 $S$ 中任意两个态射
$s : Y \to Y'$ 与 $t : Y' \to Y''$，有
$\left(t \circ s\right)^{-1} \left(t \circ f\right) = s^{-1} f$。
此外，对 $\mathcal{C}$ 中任意 $f : X \to Y'$ 与
$g : Y' \to Z'$，以及所有 $S$ 中的 $s : Z \to Z'$，有
$s^{-1} \left(g \circ f\right) = \left(s^{-1} g\right) \circ
\left(\text{id}_{Y'}^{-1} f\right)$。最后，对 $\mathcal{C}$ 中任意
$f : X \to Y'$、所有 $S$ 中的 $s : Y \to Y'$ 及
$S$ 中的 $t : Z \to Y$，有
$\left(s \circ t\right)^{-1} f
= \left(t^{-1} \text{id}_Y\right)
\circ \left(s^{-1} f\right)$。这些都直接由定义可知。
我们可以“把具有相同靶的任意有限族态射写成有公分母的分式”。

\begin{lemma}
\label{categories-lemma-morphisms-left-localization}
设 $\mathcal{C}$ 为范畴，$S$ 为 $\mathcal{C}$ 的态射构成的左乘法系。
给定 $S^{-1}\mathcal{C}$ 中任意有限族态射 $g_i : X_i \to Y$
（以 $i$ 为指标），可以找到 $S$ 的元素 $s : Y \to Y'$ 以及
$\mathcal{C}$ 的一族态射 $f_i : X_i \to Y'$，使每个 $g_i$
都是偶对 $(f_i : X_i \to Y', s : Y \to Y')$ 的等价类。
\end{lemma}

\begin{proof}
对每个 $i$，选择 $g_i$ 的一个代表
$(X_i \to Y_i, s_i : Y \to Y_i)$。若能找到 $S$ 中的态射
$s : Y \to Y'$，使得对每个 $i$ 都存在态射
$a_i : Y_i \to Y'$ 满足 $a_i \circ s_i = s$，引理即得证。
若有两个指标 $i = 1, 2$，则可补全方形
$$
\xymatrix{
Y \ar[d]_{s_1} \ar[r]_{s_2} & Y_2 \ar[d]^{t_2} \\
Y_1 \ar[r]^{a_1} & Y'
}
$$
并使 $t_2 \in S$；定义 \ref{categories-definition-multiplicative-system}
保证这可以做到。此时 $s = t_2 \circ s_2 \in S$ 满足要求。
若有 $n > 2$ 个态射，就用上述技巧将问题归约到
$n - 1$ 个态射的情形，再由归纳法得证。
\end{proof}

\noindent
具有相同分母的态射之相等性有一个简单刻画。

\begin{lemma}
\label{categories-lemma-equality-morphisms-left-localization}
设 $\mathcal{C}$ 为范畴，$S$ 为 $\mathcal{C}$ 的态射构成的左乘法系。
设 $A, B : X \to Y$ 是 $S^{-1}\mathcal{C}$ 的态射，分别为
$(f : X \to Y', s : Y \to Y')$ 与
$(g : X \to Y', s : Y \to Y')$ 的等价类。下列条件等价：
\begin{enumerate}
\item $A = B$；
\item 存在 $S$ 中的态射 $t : Y' \to Y''$，使
$t \circ f = t \circ g$；
\item 存在态射 $a : Y' \to Y''$，使
$a \circ f = a \circ g$ 且 $a \circ s \in S$。
\end{enumerate}
\end{lemma}

\begin{proof}
我们将使用 $S^{-1}\mathcal{C}$ 是范畴这一事实
（引理 \ref{categories-lemma-left-localization}），并使用定义
\ref{categories-definition-left-localization-as-fraction} 的记号以及该定义后的讨论，
来等同 $S^{-1}\mathcal{C}$ 中的某些态射。因此写
$A = s^{-1}f$ 与 $B = s^{-1}g$。

\medskip\noindent
若 $A = B$，则
$(\text{id}_{Y'}^{-1}s) \circ A = (\text{id}_{Y'}^{-1}s) \circ B$。
有 $(\text{id}_{Y'}^{-1}s) \circ A = \text{id}_{Y'}^{-1}f$ 且
$(\text{id}_{Y'}^{-1}s) \circ B = \text{id}_{Y'}^{-1}g$。
$\text{id}_{Y'}^{-1}f$ 与 $\text{id}_{Y'}^{-1}g$ 相等，按定义意味着
存在交换图表
$$
\xymatrix{
 & Y' \ar[d]^u & \\
X \ar[ru]^f \ar[r]^h \ar[rd]_g &
Z &
Y' \ar[lu]_{\text{id}_{Y'}} \ar[l]_t \ar[ld]^{\text{id}_{Y'}} \\
& Y' \ar[u]_v &
}
$$
其中 $t \in S$。特别地，$u = v = t \in S$ 且
$t \circ f = t\circ g$。故 (1) 蕴含 (2)。

\medskip\noindent
(2) $\Rightarrow$ (3) 是立即的。假设 $a$ 如 (3) 中所述，记
$s' = a \circ s \in S$。则 $\text{id}_{Y''}^{-1}s'$ 是范畴
$S^{-1}\mathcal{C}$ 中的同构（逆为 $(s')^{-1}\text{id}_{Y''}$）。
所以为验证 $A = B$，只须验证
$\text{id}_{Y''}^{-1}s' \circ A = \text{id}_{Y''}^{-1}s' \circ B$。
利用定义 \ref{categories-definition-left-localization-as-fraction} 后文所讨论的规则，
计算得到
$\text{id}_{Y''}^{-1}s' \circ A =
\text{id}_{Y''}^{-1}(a \circ s) \circ s^{-1}f =
\text{id}_{Y''}^{-1}(a \circ f) =
\text{id}_{Y''}^{-1}(a \circ g) =
\text{id}_{Y''}^{-1}(a \circ s) \circ s^{-1}g =
\text{id}_{Y''}^{-1}s' \circ B$，故 (1) 成立。
\end{proof}

\begin{remark}
\label{categories-remark-left-localization-morphisms-colimit}
设 $\mathcal{C}$ 为范畴，$S$ 为左乘法系。给定 $\mathcal{C}$ 的对象
$Y$，以 $Y/S$ 表示如下范畴：其对象是满足 $s \in S$ 的
$s : Y \to Y'$，其态射是交换图表
$$
\xymatrix{
& Y \ar[ld]_s \ar[rd]^t & \\
Y' \ar[rr]^a & & Y''
}
$$
其中 $a : Y' \to Y''$ 任意。我们断言范畴 $Y/S$ 滤过
（见定义 \ref{categories-definition-directed}）。事实上，LMS1 蕴含
$\text{id}_Y : Y \to Y$ 属于 $Y/S$，故 $Y/S$ 非空。
LMS2 蕴含：给定 $s_1 : Y \to Y_1$ 与 $s_2 : Y \to Y_2$，
可找到图表
$$
\xymatrix{
Y \ar[d]_{s_1} \ar[r]_{s_2} & Y_2 \ar[d]^t \\
Y_1 \ar[r]^a & Y_3
}
$$
且 $t \in S$。所以 $s_1 : Y \to Y_1$ 与 $s_2 : Y \to Y_2$
在 $Y/S$ 中都有到 $t \circ s_2 : Y \to Y_3$ 的映射。
最后，给定 $Y/S$ 中从 $s_1 : Y \to Y_1$ 到
$s_2 : Y \to Y_2$ 的两个态射 $a, b$，有
$a \circ s_1 = b \circ s_1$；故由 LMS3，存在 $S$ 中的
$t : Y_2 \to Y_3$，使 $t \circ a = t \circ b$。
现在，引理 \ref{categories-lemma-morphisms-left-localization} 与
\ref{categories-lemma-equality-morphisms-left-localization} 的合并结果告诉我们
\begin{equation}
\label{categories-equation-left-localization-morphisms-colimit}
\Mor_{S^{-1}\mathcal{C}}(X, Y) =
\colim_{(s : Y \to Y') \in Y/S} \Mor_\mathcal{C}(X, Y')
\end{equation}
这个把 $S^{-1}\mathcal{C}$ 中的态射集表示为 $\mathcal{C}$ 中态射集的
滤过余极限的公式，有时很有用。
\end{remark}

\begin{lemma}
\label{categories-lemma-properties-left-localization}
设 $\mathcal{C}$ 为范畴，$S$ 为 $\mathcal{C}$ 的态射构成的左乘法系。
\begin{enumerate}
\item 规则 $X \mapsto X$ 以及
$(f : X \to Y) \mapsto (f : X \to Y, \text{id}_Y : Y \to Y)$
定义函子 $Q : \mathcal{C} \to S^{-1}\mathcal{C}$。
\item 对任意 $s \in S$，态射 $Q(s)$ 是
$S^{-1}\mathcal{C}$ 中的同构。
\item 若 $G : \mathcal{C} \to \mathcal{D}$ 为任意函子，且对每个
$s \in S$，$G(s)$ 都可逆，则存在唯一函子
$H : S^{-1}\mathcal{C} \to \mathcal{D}$，使 $H \circ Q = G$。
\end{enumerate}
\end{lemma}

\begin{proof}
(1) 与 (2) 很清楚。（在 (2) 中，$Q(s)$ 的逆是偶对
$(\text{id}_Y, s)$ 的等价类。）为说明 (3)，只须令
$H(X) = G(X)$，并令
$H((f : X \to Y', s : Y \to Y')) = G(s)^{-1} \circ G(f)$。
细节略去。
\end{proof}

\begin{lemma}
\label{categories-lemma-left-localization-limits}
设 $\mathcal{C}$ 为范畴，$S$ 为 $\mathcal{C}$ 的态射构成的左乘法系。
局部化函子 $Q : \mathcal{C} \to S^{-1}\mathcal{C}$ 与有限余极限交换。
\end{lemma}

\begin{proof}
设 $\mathcal{I}$ 为有限范畴，并设
$\mathcal{I} \to \mathcal{C}$、$i \mapsto X_i$ 为余极限存在的函子。
利用 (\ref{categories-equation-left-localization-morphisms-colimit})、$Y/S$ 滤过这一事实
以及引理 \ref{categories-lemma-directed-commutes}，有
\begin{align*}
\Mor_{S^{-1}\mathcal{C}}(Q(\colim X_i), Q(Y))
& =
\colim_{(s : Y \to Y') \in Y/S} \Mor_\mathcal{C}(\colim X_i, Y') \\
& =
\colim_{(s : Y \to Y') \in Y/S} \lim_i \Mor_\mathcal{C}(X_i, Y') \\
& =
\lim_i \colim_{(s : Y \to Y') \in Y/S} \Mor_\mathcal{C}(X_i, Y') \\
& =
\lim_i \Mor_{S^{-1}\mathcal{C}}(Q(X_i), Q(Y))
\end{align*}
而且，对每个 $j \in \Ob(\mathcal{I})$，此同构与两边到集合
$\Mor_{S^{-1}\mathcal{C}}(Q(X_j), Q(Y))$ 的投影交换。因此，
$Q(\colim X_i)$ 满足函子 $i \mapsto Q(X_i)$ 之余极限的泛性质，
故如所需，它就是该余极限。
\end{proof}

\begin{lemma}
\label{categories-lemma-left-localization-diagram}
设 $\mathcal{C}$ 为范畴，$S$ 为左乘法系。若
$f : X \to Y$、$f' : X' \to Y'$ 是 $\mathcal{C}$ 的两个态射，并且
$$
\xymatrix{
Q(X) \ar[d]_{Q(f)} \ar[r]_a & Q(X') \ar[d]^{Q(f')} \\
Q(Y) \ar[r]^b & Q(Y')
}
$$
是 $S^{-1}\mathcal{C}$ 中的交换图表，则存在 $\mathcal{C}$ 中的态射
$f'' : X'' \to Y''$ 以及 $\mathcal{C}$ 中的交换图表
$$
\xymatrix{
X \ar[d]_f \ar[r]_g & X'' \ar[d]^{f''} & X' \ar[d]^{f'} \ar[l]^s \\
Y \ar[r]^h & Y'' & Y' \ar[l]_t
}
$$
其中 $s, t \in S$，且 $a = s^{-1}g$、$b = t^{-1}h$。
\end{lemma}

\begin{proof}
如下选择映射与对象。首先，将 $a = s^{-1}g$ 写成某个
$S$ 中的 $s : X' \to X''$ 与某个 $g : X \to X''$。
由 LMS2，可以找到 $S$ 中的 $t : Y' \to Y''$ 以及
$f'' : X'' \to Y''$，使图表
$$
\xymatrix{
X \ar[d]_f \ar[r]_g & X'' \ar[d]^{f''} & X' \ar[d]^{f'} \ar[l]^s \\
Y & Y'' & Y' \ar[l]_t
}
$$
可交换。接下来，在该图表中，我们将反复改变对
$$
X'' \xrightarrow{f''} Y'' \xleftarrow{t} Y'
$$
的选择，方法是用一个满足 $d \circ t \in S$ 的态射
$d : Y'' \to Y'''$ 与 $t$ 和 $f''$ 后复合。根据注
\ref{categories-remark-left-localization-morphisms-colimit}，经过这种替换后，
可假设存在态射 $h : Y \to Y''$，使 $b = t^{-1}h$\footnote{还可用
更直接的方式看出这一点：把 $b = q^{-1}i$ 写成 $S$ 中的某个
$q : Y' \to Z$ 及某个 $i : Y \to Z$。由 LMS2，可找到
$S$ 中的 $r : Y'' \to Y'''$ 以及 $j : Z \to Y'''$，使
$j \circ q = r \circ t$。现在令 $d = r$ 且 $h = j \circ i$。}。
此时我们已具备引理中的所有条件，唯独尚不知道图表的左方形是否交换。
但 $S^{-1} \mathcal{C}$ 中复合的定义表明，
$b \circ Q\left(f\right)$ 是偶对
$(h \circ f : X \to Y'', t : Y' \to Y'')$ 的等价类
（因为 $b$ 是偶对 $(h : Y \to Y'', t : Y' \to Y'')$ 的等价类，
而 $Q\left(f\right)$ 是偶对
$(f : X \to Y, \text{id} : Y \to Y)$ 的等价类）；另一方面，
$Q\left(f'\right) \circ a$ 是偶对
$(f'' \circ g : X \to Y'', t : Y' \to Y'')$ 的等价类
（因为 $a$ 是偶对 $(g : X \to X'', s : X' \to X'')$ 的等价类，
而 $Q\left(f'\right)$ 是偶对
$(f' : X' \to Y', \text{id} : Y' \to Y')$ 的等价类）。
已知 $b \circ Q\left(f\right) = Q\left(f'\right) \circ a$，
故偶对 $(h \circ f : X \to Y'', t : Y' \to Y'')$ 与
$(f'' \circ g : X \to Y'', t : Y' \to Y'')$ 的等价类相等。
因此，由引理 \ref{categories-lemma-equality-morphisms-left-localization}，
可找到态射 $d : Y'' \to Y'''$，使 $d \circ t \in S$ 且
$d \circ h \circ f = d \circ f'' \circ g$。最后再作一次上述类型的
替换，便得到所需结论。
\end{proof}

\noindent
{\bf 右分式演算。}
设 $\mathcal{C}$ 为范畴，$S$ 为右乘法系。如下定义一个新范畴
$S^{-1}\mathcal{C}$（我们将在引理 \ref{categories-lemma-right-localization}
的证明中验证此定义确实成立）：
\begin{enumerate}
\item 令 $\Ob(S^{-1}\mathcal{C}) = \Ob(\mathcal{C})$。
\item $S^{-1}\mathcal{C}$ 中的态射 $X \to Y$ 由偶对
$(f : X' \to Y, s : X' \to X)$ 给出，其中 $s \in S$，
并模去等价关系。（该等价关系定义如下。可以把偶对 $(f, s)$
的等价类看作 $fs^{-1} : X \to Y$。）
\item 称两个偶对 $(f_1 : X_1 \to Y, s_1 : X_1 \to X)$ 与
$(f_2 : X_2 \to Y, s_2 : X_2 \to X)$ 等价，如果存在第三个偶对
$(f_3 : X_3 \to Y, s_3 : X_3 \to X)$ 以及 $\mathcal{C}$ 的态射
$u : X_3 \to X_1$ 与 $v : X_3 \to X_2$，使它们嵌入交换图表
$$
\xymatrix{
 & X_1 \ar[ld]_{s_1} \ar[rd]^{f_1} & \\
X &
X_3 \ar[l]_{s_3} \ar[u]_u \ar[d]^v \ar[r]^{f_3} &
Y \\
& X_2 \ar[lu]^{s_2} \ar[ru]_{f_2} &
}
$$
\item 偶对 $(f : X' \to Y, s : X' \to X)$ 与
$(g : Y' \to Z, t : Y' \to Y)$ 的等价类之复合，定义为偶对
$(g \circ h : X'' \to Z, s \circ u : X'' \to X)$ 的等价类；
其中选择 $h$ 与 $u \in S$ 使其嵌入交换图表
$$
\xymatrix{
X'' \ar[d]_u \ar[r]^h & Y' \ar[d]^t \\
X' \ar[r]^f & Y
}
$$
这样的选择由假设存在。
\item $S^{-1} \mathcal{C}$ 中的恒等态射 $X \to X$ 是偶对
$(\text{id} : X \to X,
\text{id} : X \to X)$ 的等价类。
\end{enumerate}

\begin{lemma}
\label{categories-lemma-right-localization}
设 $\mathcal{C}$ 为范畴，$S$ 为右乘法系。
\begin{enumerate}
\item 上面在偶对上定义的关系是等价关系。
\item 上面给出的复合规则在等价类上良定义。
\item 复合满足结合律（且恒等态射满足恒等公理），
所以 $S^{-1}\mathcal{C}$ 是范畴。
\end{enumerate}
\end{lemma}

\begin{proof}
本引理与引理 \ref{categories-lemma-left-localization} 对偶。把
$\mathcal{C}$ 换成其对偶范畴（其中 $S$ 是左乘法系），
即可由该引理形式地推出本引理。
\end{proof}

\begin{definition}
\label{categories-definition-right-localization-as-fraction}
设 $\mathcal{C}$ 为范畴，$S$ 为 $\mathcal{C}$ 的态射构成的右乘法系。
给定 $\mathcal{C}$ 中任意态射 $f : X' \to Y$ 及 $S$ 中任意态射
$s : X' \to X$，以 {\it $f s^{-1}$} 表示偶对
$(f : X' \to Y, s : X' \to X)$ 的等价类。这是
$S^{-1} \mathcal{C}$ 中从 $X$ 到 $Y$ 的态射。
\end{definition}

\noindent
与定义 \ref{categories-definition-left-localization-as-fraction} 中恒等式类似
（事实上对偶）的恒等式成立。我们可以“把具有相同源的任意有限族态射
写成有公分母的分式”。

\begin{lemma}
\label{categories-lemma-morphisms-right-localization}
设 $\mathcal{C}$ 为范畴，$S$ 为 $\mathcal{C}$ 的态射构成的右乘法系。
给定 $S^{-1}\mathcal{C}$ 中任意有限族态射 $g_i : X \to Y_i$
（以 $i$ 为指标），可以找到 $S$ 的元素 $s : X' \to X$
以及 $\mathcal{C}$ 的一族态射 $f_i : X' \to Y_i$，使 $g_i$
是偶对 $(f_i : X' \to Y_i, s : X' \to X)$ 的等价类。
\end{lemma}

\begin{proof}
本引理是引理 \ref{categories-lemma-morphisms-left-localization} 的对偶；
把所涉及的所有范畴换成其对偶范畴，即可由该引理形式地推出。
\end{proof}

\noindent
具有相同分母的态射之相等性有一个简单刻画。

\begin{lemma}
\label{categories-lemma-equality-morphisms-right-localization}
设 $\mathcal{C}$ 为范畴，$S$ 为 $\mathcal{C}$ 的态射构成的右乘法系。
设 $A, B : X \to Y$ 是 $S^{-1}\mathcal{C}$ 的态射，分别为
$(f : X' \to Y, s : X' \to X)$ 与
$(g : X' \to Y, s : X' \to X)$ 的等价类。下列条件等价：
\begin{enumerate}
\item $A = B$；
\item 存在 $S$ 中的态射 $t : X'' \to X'$，使
$f \circ t = g \circ t$；
\item 存在态射 $a : X'' \to X'$，使
$f \circ a = g \circ a$ 且 $s \circ a \in S$。
\end{enumerate}
\end{lemma}

\begin{proof}
这与引理 \ref{categories-lemma-equality-morphisms-left-localization} 对偶。
\end{proof}

\begin{remark}
\label{categories-remark-right-localization-morphisms-colimit}
设 $\mathcal{C}$ 为范畴，$S$ 为右乘法系。给定 $\mathcal{C}$ 的对象
$X$，以 $S/X$ 表示如下范畴：其对象是满足 $s \in S$ 的
$s : X' \to X$，其态射是交换图表
$$
\xymatrix{
X' \ar[rd]_s \ar[rr]_a & & X'' \ar[ld]^t \\
& X
}
$$
其中 $a : X' \to X''$ 任意。范畴 $S/X$ 余滤过
（见定义 \ref{categories-definition-codirected}）。
（这与注 \ref{categories-remark-left-localization-morphisms-colimit} 中的相应断言对偶。）
现在，引理 \ref{categories-lemma-morphisms-right-localization} 与
\ref{categories-lemma-equality-morphisms-right-localization} 的合并结果告诉我们
\begin{equation}
\label{categories-equation-right-localization-morphisms-colimit}
\Mor_{S^{-1}\mathcal{C}}(X, Y) =
\colim_{(s : X' \to X) \in (S/X)^{opp}} \Mor_\mathcal{C}(X', Y)
\end{equation}
这个把 $S^{-1}\mathcal{C}$ 中的态射表示为 $\mathcal{C}$ 中态射的
滤过余极限的公式，有时很有用。
\end{remark}

\begin{lemma}
\label{categories-lemma-properties-right-localization}
设 $\mathcal{C}$ 为范畴，$S$ 为 $\mathcal{C}$ 的态射构成的右乘法系。
\begin{enumerate}
\item 规则 $X \mapsto X$ 以及
$(f : X \to Y) \mapsto (f : X \to Y, \text{id}_X : X \to X)$
定义函子 $Q : \mathcal{C} \to S^{-1}\mathcal{C}$。
\item 对任意 $s \in S$，态射 $Q(s)$ 是
$S^{-1}\mathcal{C}$ 中的同构。
\item 若 $G : \mathcal{C} \to \mathcal{D}$ 为任意函子，且对每个
$s \in S$，$G(s)$ 都可逆，则存在唯一函子
$H : S^{-1}\mathcal{C} \to \mathcal{D}$，使 $H \circ Q = G$。
\end{enumerate}
\end{lemma}

\begin{proof}
本引理是引理 \ref{categories-lemma-properties-left-localization} 的对偶；
把所涉及的所有范畴换成其对偶范畴，即可由该引理形式地推出。
\end{proof}

\begin{lemma}
\label{categories-lemma-right-localization-limits}
设 $\mathcal{C}$ 为范畴，$S$ 为 $\mathcal{C}$ 的态射构成的右乘法系。
局部化函子 $Q : \mathcal{C} \to S^{-1}\mathcal{C}$ 与有限极限交换。
\end{lemma}

\begin{proof}
这与引理 \ref{categories-lemma-left-localization-limits} 对偶。
\end{proof}

\begin{lemma}
\label{categories-lemma-right-localization-diagram}
设 $\mathcal{C}$ 为范畴，$S$ 为右乘法系。若
$f : X \to Y$、$f' : X' \to Y'$ 是 $\mathcal{C}$ 的两个态射，并且
$$
\xymatrix{
Q(X) \ar[d]_{Q(f)} \ar[r]_a & Q(X') \ar[d]^{Q(f')} \\
Q(Y) \ar[r]^b & Q(Y')
}
$$
是 $S^{-1}\mathcal{C}$ 中的交换图表，则存在 $\mathcal{C}$ 中的态射
$f'' : X'' \to Y''$ 以及 $\mathcal{C}$ 中的交换图表
$$
\xymatrix{
X \ar[d]_f & X'' \ar[l]^s \ar[d]^{f''} \ar[r]_g & X' \ar[d]^{f'} \\
Y & Y'' \ar[l]_t \ar[r]^h & Y'
}
$$
其中 $s, t \in S$，且 $a = gs^{-1}$、$b = ht^{-1}$。
\end{lemma}

\begin{proof}
本引理与引理 \ref{categories-lemma-left-localization-diagram} 对偶。
\end{proof}

\noindent
{\bf 乘法系与双侧分式演算。}
若 $S$ 是乘法系，则左、右分式演算给出典范同构的范畴。

\begin{lemma}
\label{categories-lemma-multiplicative-system}
设 $\mathcal{C}$ 为范畴，$S$ 为乘法系。左分式范畴与右分式范畴
$S^{-1}\mathcal{C}$ 典范同构。
\end{lemma}

\begin{proof}
以 $\mathcal{C}_{left}$、$\mathcal{C}_{right}$ 表示两个分式范畴。
由引理 \ref{categories-lemma-properties-left-localization} 与
\ref{categories-lemma-properties-right-localization} 的泛性质，得到函子
$\mathcal{C}_{left} \to \mathcal{C}_{right}$ 与
$\mathcal{C}_{right} \to \mathcal{C}_{left}$。
由泛性质中的唯一性断言，这两个函子互为逆。
\end{proof}

\begin{definition}
\label{categories-definition-saturated-multiplicative-system}
设 $\mathcal{C}$ 为范畴，$S$ 为乘法系。若除了 MS1、MS2、MS3 之外，
还有
\begin{enumerate}
\item[MS4] 给定三个可复合态射 $f, g, h$，若
$fg, gh \in S$，则 $g \in S$，
\end{enumerate}
则称 $S$ 是{\it 饱和}的。
\end{definition}

\noindent
注意，饱和乘法系包含所有同构。此外，若 $f, g, h$ 是范畴中
可复合的态射，且 $fg, gh$ 是同构，则 $g$ 是同构
（因为此时 $g$ 同时有左逆和右逆，故可逆）。

\begin{lemma}
\label{categories-lemma-what-gets-inverted}
设 $\mathcal{C}$ 为范畴，$S$ 为乘法系。以
$Q : \mathcal{C} \to S^{-1}\mathcal{C}$ 表示局部化函子。集合
$$
\hat S = \{f \in \text{Arrows}(\mathcal{C}) \mid
Q(f) \text{ 为同构}\}
$$
等于
$$
S' = \{f \in \text{Arrows}(\mathcal{C}) \mid
\text{存在 }g, h\text{ 使得 }gf, fh \in S\}
$$
并且是包含 $S$ 的最小饱和乘法系。特别地，若 $S$ 饱和，
则 $\hat S = S$。
\end{lemma}

\begin{proof}
显然 $S \subset S' \subset \hat S$，因为 $S'$ 的元素映到
$S^{-1}\mathcal{C}$ 中同时具有左逆和右逆的态射。注意 $S'$ 满足
MS4，而 $\hat S$ 满足 MS1。接下来证明 $S' = \hat S$。

\medskip\noindent
设 $f \in \hat S$。令 $s^{-1}g = ht^{-1}$ 为
$S^{-1}\mathcal{C}$ 中的逆态射。（由引理
\ref{categories-lemma-multiplicative-system}，可以同时使用左分式和右分式描述
$S^{-1}\mathcal{C}$ 中的态射。）关系
$\text{id}_X = s^{-1}gf$ 在 $S^{-1}\mathcal{C}$ 中成立，意味着存在
交换图表
$$
\xymatrix{
 & X' \ar[d]^u & \\
X \ar[ru]^{gf} \ar[r]^{f'} \ar[rd]_{\text{id}_X} &
X'' &
X \ar[lu]_s \ar[l]_{s'} \ar[ld]^{\text{id}_X} \\
& X \ar[u]_v &
}
$$
其中有某些态射 $f', u, v$ 以及 $s' \in S$。因此
$ugf = s' \in S$。类似地，利用 $\text{id}_Y = fht^{-1}$，
可证明对某个 $w$ 有 $fhw \in S$。由此 $f \in S'$，故
$S' = \hat S$。只要证明 $S' = \hat S$ 是乘法系，就显然可推出
$S' = \hat S$ 是包含 $S$ 的最小饱和乘法系。

\medskip\noindent
上面的讨论已经处理了 MS1 与 MS4，所以为完成引理的证明，
须证明 $\hat S$ 满足 LMS2、RMS2、LMS3、RMS3。
先验证 $\hat S$ 的 LMS2。假设有实线图表
$$
\xymatrix{
X \ar[d]_t \ar[r]_g & Y \ar@{..>}[d]^s \\
Z \ar@{..>}[r]^f & W
}
$$
其中 $t \in \hat S$。选择态射 $a : Z \to Z'$，使 $at \in S$。
于是可对 $S$ 使用 LMS2，找到交换图表
$$
\xymatrix{
X \ar[d]_t \ar[r]_g & Y \ar[dd]^s \\
Z \ar[d]_a \\
Z' \ar[r]^{f'} & W
}
$$
令 $f = f' \circ a$ 即得结论。RMS2 的证明与此对偶。
最后，假设给定一对态射 $f, g : X \to Y$ 以及靶为 $X$ 的
$t \in \hat S$，并且 $ft = gt$。选择态射 $b$ 使 $tb \in S$。
于是 $ftb = gtb$，由 $S$ 的 LMS3，存在源为 $Y$ 的
$s \in S$，使 $sf = sg$，正如所需。RMS3 的证明与此对偶。
\end{proof}








\section{形式性质}
\label{categories-section-formal-cat-cat}

\noindent
本节讨论范畴的 $2$-范畴的一些形式性质。这将引出稍后对（严格）
$2$-范畴的定义。

\medskip\noindent
以 $\Ob(\textit{Cat})$ 表示所有范畴构成的类。对每一对范畴
$\mathcal{A}, \mathcal{B} \in \Ob(\textit{Cat})$，都有函子构成的
“小”范畴 $\text{Fun}(\mathcal{A}, \mathcal{B})$。像
$$
\xymatrix{
\mathcal{A}
\rruppertwocell^{F''}{t'}
\ar[rr]_(.3){F'}
\rrlowertwocell_F{t}
& &
\mathcal{B}
}
\text{ 复合为 }
\xymatrix{
\mathcal{A}
\rrtwocell^{F''}_F{\ \ t \circ t'}
& &
\mathcal{B}
}
$$
这样的函子变换之复合称为{\it 竖直}复合。对此使用通常的符号
$\circ$。接下来定义{\it 水平}复合。为此，先进一步说明现有结构。

\medskip\noindent
具体地，对每三个范畴 $\mathcal{A}$、$\mathcal{B}$、$\mathcal{C}$，
由函子复合得到复合法则
$$
\circ : \Ob(\text{Fun}(\mathcal{B}, \mathcal{C}))
\times
\Ob(\text{Fun}(\mathcal{A}, \mathcal{B}))
\longrightarrow
\Ob(\text{Fun}(\mathcal{A}, \mathcal{C}))
$$
该复合法则满足结合律，恒等函子充当单位。换言之，若忘掉函子变换，
可见 $\textit{Cat}$ 构成一个范畴。这个结构如何与函子之间的态射相互作用？

\medskip\noindent
给定函子 $F, F' : \mathcal{A} \to \mathcal{B}$ 的变换
$t : F \to F'$ 以及函子 $G : \mathcal{B} \to \mathcal{C}$，
可以定义函子变换 $G\circ F \to G \circ F'$。将此变换记为
${}_Gt$。对所有 $x \in \mathcal{A}$，它由公式
$({}_Gt)_x = G(t_x) : G(F(x)) \to G(F'(x))$ 给出。
这样，与 $G$ 复合成为函子
$$
\text{Fun}(\mathcal{A}, \mathcal{B})
\longrightarrow
\text{Fun}(\mathcal{A}, \mathcal{C}).
$$
只须验证 ${}_G(\text{id}_F) = \text{id}_{G \circ F}$ 以及
${}_G(t_1 \circ t_2) = {}_Gt_1 \circ {}_Gt_2$ 即可看出这一点。
当然还有 ${}_{\text{id}_\mathcal{B}}t = t$。

\medskip\noindent
类似地，给定函子 $G, G' : \mathcal{B} \to \mathcal{C}$ 的变换
$s : G \to G'$ 以及函子 $F : \mathcal{A} \to \mathcal{B}$，
可定义 $s_F$ 为函子变换 $G\circ F \to G' \circ F$，其中对所有
$x \in \mathcal{A}$，有
$(s_F)_x = s_{F(x)} : G(F(x)) \to G'(F(x))$。
这样，与 $F$ 复合成为函子
$$
\text{Fun}(\mathcal{B}, \mathcal{C})
\longrightarrow
\text{Fun}(\mathcal{A}, \mathcal{C}).
$$
只须验证 $(\text{id}_G)_F = \text{id}_{G\circ F}$ 以及
$(s_1 \circ s_2)_F = s_{1, F} \circ s_{2, F}$ 即可看出这一点。
当然还有 $s_{\text{id}_\mathcal{B}} = s$。

\medskip\noindent
只要各式有意义，这些构造还满足
$$
{}_{G_1}({}_{G_2}t) = {}_{G_1\circ G_2}t,
\ (s_{F_1})_{F_2} = s_{F_1 \circ F_2},
\text{ 且 }{}_H(s_F) = ({}_Hs)_F
$$
最后，给定函子 $F, F' : \mathcal{A} \to \mathcal{B}$、
$G, G' : \mathcal{B} \to \mathcal{C}$ 以及变换
$t : F \to F'$、$s : G \to G'$，下图可交换：
$$
\xymatrix{
G \circ F \ar[r]^{{}_Gt} \ar[d]_{s_F}
&
G \circ F' \ar[d]^{s_{F'}} \\
G' \circ F \ar[r]_{{}_{G'}t}
&
G' \circ F'
}
$$
换言之，${}_{G'}t \circ s_F = s_{F'}\circ {}_Gt$。
为证明这一点，只须考察任意对象 $x \in \Ob(\mathcal{A})$ 上的情形：
$$
\xymatrix{
G(F(x)) \ar[r]^{G(t_x)} \ar[d]_{s_{F(x)}}
&
G(F'(x)) \ar[d]^{s_{F'(x)}} \\
G'(F(x)) \ar[r]_{G'(t_x)}
&
G'(F'(x))
}
$$
该图可交换，因为 $s$ 是函子变换。这个相容关系使我们可以定义水平复合。

\begin{definition}
\label{categories-definition-horizontal-composition}
给定下式左侧所示图表：
$$
\xymatrix{
\mathcal{A}
\rtwocell^F_{F'}{t}
&
\mathcal{B}
\rtwocell^G_{G'}{s}
&
\mathcal{C}
}
\text{ 给出 }
\xymatrix{
\mathcal{A}
\rrtwocell^{G \circ F} _{G' \circ F'}{\ \ s \star t}
& &
\mathcal{C}
}
$$
定义{\it 水平}复合 $s \star t$ 为函子变换
${}_{G'}t \circ s_F = s_{F'}\circ {}_Gt$。
\end{definition}

\noindent
现在可将先前构造的变换 ${}_Gt$ 与 $s_F$ 恢复为
$ {}_Gt = \text{id}_G \star t $ 与 $ s_F = s \star \text{id}_F $。
此外，上面得到的所有规则都是下一个引理所述性质的推论。

\begin{lemma}
\label{categories-lemma-properties-2-cat-cats}
水平复合与竖直复合具有下列性质：
\begin{enumerate}
\item $\circ$ 与 $\star$ 满足结合律；
\item 恒等变换 $\text{id}_F$ 是 $\circ$ 的单位；
\item 恒等函子的恒等变换
$\text{id}_{\text{id}_\mathcal{A}}$ 是 $\star$ 与 $\circ$ 的单位；
\item 给定图表
$$
\xymatrix{
\mathcal{A}
\rruppertwocell^F{t}
\ar[rr]_(.3){F'}
\rrlowertwocell_{F''}{t'}
& &
\mathcal{B}
\rruppertwocell^G{s}
\ar[rr]_(.3){G'}
\rrlowertwocell_{G''}{s'}
& &
\mathcal{C}
}
$$
有 $ (s' \circ s) \star (t' \circ t) = (s' \star t') \circ (s \star t)$。
\end{enumerate}
\end{lemma}

\begin{proof}
用先前的记号，最后一个断言化为等式
$$
s'_{F''}
\circ
{}_{G'}t'
\circ
s_{F'}
\circ
{}_Gt
=
(s' \circ s)_{F''}
\circ
{}_G(t' \circ t).
$$
把上面所得结果用于中间的复合，可将左边改写为
$
s'_{F''}
\circ
s_{F''}
\circ
{}_Gt'
\circ
{}_Gt
$，它显然等于右边。
\end{proof}

\noindent
引理条件 (4) 的另一种表述是：函子复合与函子变换的水平复合给出函子
$$
(\circ, \star) :
\text{Fun}(\mathcal{B}, \mathcal{C})
\times
\text{Fun}(\mathcal{A}, \mathcal{B})
\longrightarrow
\text{Fun}(\mathcal{A}, \mathcal{C})
$$
其源为乘积范畴，见定义 \ref{categories-definition-product-category}。

\section{2-范畴}
\label{categories-section-2-categories}

\noindent
我们将定义栈的语境中出现的（严格）$2$-范畴。阅读前，请先看
第 \ref{categories-section-formal-cat-cat} 节与例
\ref{categories-example-2-1-category-of-categories}。基本上，就是取该例，并把
其中对象、$1$-态射、$2$-态射所满足的全部规则写出来。

\begin{definition}
\label{categories-definition-2-category}
一个（严格）{\it $2$-范畴} $\mathcal{C}$ 由下列数据组成：
\begin{enumerate}
\item 对象集 $\Ob(\mathcal{C})$。
\item 对每一对 $x, y \in \Ob(\mathcal{C})$，给定范畴
$\Mor_\mathcal{C}(x, y)$。$\Mor_\mathcal{C}(x, y)$ 的对象称为
{\it $1$-态射}，记为 $F : x \to y$；这些 $1$-态射之间的态射
称为{\it $2$-态射}，记为 $t : F' \to F$。
$\Mor_\mathcal{C}(x, y)$ 中 $2$-态射的复合称为{\it 竖直}复合；
对 $t : F' \to F$ 与 $t' : F'' \to F'$，记为 $t \circ t'$。
\item 对每三个 $x, y, z\in \Ob(\mathcal{C})$，给定函子
$$
(\circ, \star) :
\Mor_\mathcal{C}(y, z) \times \Mor_\mathcal{C}(x, y)
\longrightarrow
\Mor_\mathcal{C}(x, z).
$$
左侧一对 $1$-态射 $(F, G)$ 的像称为 $F$ 与 $G$ 的
{\it 复合}，记为 $F\circ G$。一对 $2$-态射 $(t, s)$ 的像称为
{\it 水平}复合，记为 $t \star s$。
\end{enumerate}
这些数据须满足下列规则：
\begin{enumerate}
\item 对象集与 $1$-态射集连同 $1$-态射的复合构成一个范畴。
\item $2$-态射的水平复合满足结合律。
\item 恒等 $1$-态射 $\text{id}_x$ 的恒等 $2$-态射
$\text{id}_{\text{id}_x}$ 是水平复合的单位。
\end{enumerate}
\end{definition}

\noindent
显然，这并不是一种很便于使用的对象。不过，有许多例子中如何使用它
相当清楚。迄今唯一的例子是如下 $2$-范畴：其对象是给定的一族范畴，
$1$-态射是这些范畴之间的函子，$2$-态射是函子的自然变换，见
第 \ref{categories-section-formal-cat-cat} 节。就本文而言，所有 $2$-范畴都将是
该例的子 $2$-范畴。以下说明何谓子 $2$-范畴。

\begin{definition}
\label{categories-definition-sub-2-category}
设 $\mathcal{C}$ 为 $2$-范畴。$\mathcal{C}$ 的一个
{\it 子 $2$-范畴} $\mathcal{C}'$ 由 $\Ob(\mathcal{C})$ 的子集
$\Ob(\mathcal{C}')$ 以及对所有 $x, y \in \Ob(\mathcal{C}')$，
范畴 $\Mor_\mathcal{C}(x, y)$ 的子范畴
$\Mor_{\mathcal{C}'}(x, y)$ 给出；要求这些数据连同运算
$\circ$（$1$-态射的复合）、$\circ$（$2$-态射的竖直复合）以及
$\star$（水平复合）构成一个 $2$-范畴。
\end{definition}

\begin{remark}
\label{categories-remark-big-2-categories}
大 $2$-范畴。在许多文献中，允许 $2$-范畴具有对象类
（但愿不允许“类的类”）。我们也允许这些“大”$2$-范畴，
不过仅限于下列情形（以后随进展更新）：
\begin{enumerate}
\item 范畴的 $2$-范畴 $\textit{Cat}$。
\item 范畴的 $(2, 1)$-范畴 $\textit{Cat}$。
\item 群胚的 $2$-范畴 $\textit{Groupoids}$；这是一个
$(2, 1)$-范畴。
\item 固定范畴上的纤维化范畴的 $2$-范畴。
\item 固定范畴上的纤维化范畴的 $(2, 1)$-范畴。
\item 固定范畴上的群胚纤维化范畴的 $2$-范畴；这是一个
$(2, 1)$-范畴。
\item 固定站点上的栈的 $2$-范畴。
\item 固定站点上的栈的 $(2, 1)$-范畴。
\item 固定站点上的群胚值栈的 $2$-范畴；这是一个
$(2, 1)$-范畴。
\item 固定站点上的广群值栈的 $2$-范畴；这是一个
$(2, 1)$-范畴。
\item 固定概形上的代数栈的 $2$-范畴；这是一个
$(2, 1)$-范畴。
\end{enumerate}
见定义 \ref{categories-definition-2-1-category}。注意，在每种情形中，
$2$-范畴 $\mathcal{C}$ 的对象类都是真类，但对所有对象
$x, y \in \Ob(C)$，范畴 $\Mor_\mathcal{C}(x, y)$ 都是“小”的
（按照我们的约定）。
\end{remark}

\noindent
第 \ref{categories-section-definition-categories} 节中定义的范畴等价概念，
可如下推广到更一般的 $2$-范畴情形。

\begin{definition}
\label{categories-definition-equivalence}
若存在 $1$-态射 $F : x \to y$ 与 $G : y \to x$，使
$F \circ G$ 与 $\text{id}_y$ 之间 $2$-同构，且
$G \circ F$ 与 $\text{id}_x$ 之间 $2$-同构，则称 $2$-范畴的两个对象
$x, y$ {\it 等价}。
\end{definition}

\noindent
有时需要说明从一个范畴到一个 $2$-范畴的函子是什么意思。

\begin{definition}
\label{categories-definition-functor-into-2-category}
设 $\mathcal{A}$ 为范畴，$\mathcal{C}$ 为 $2$-范畴。
\begin{enumerate}
\item 除非另有说明，从通常范畴到 $2$-范畴的{\it 函子}忽略
$2$-态射。换言之，它是一个“通常”函子，其靶是从 2-范畴忘掉所有
2-态射后形成的范畴。
\item 从 $\mathcal{A}$ 到 2-范畴 $\mathcal{C}$ 的
{\it 弱函子}，或{\it 伪函子} $\varphi$，由下列数据给出：
\begin{enumerate}
\item 映射 $\varphi : \Ob(\mathcal{A}) \to \Ob(\mathcal{C})$；
\item 对每一对 $x, y\in \Ob(\mathcal{A})$ 以及每个态射
$f : x \to y$，给定 $1$-态射
$\varphi(f) : \varphi(x) \to \varphi(y)$；
\item 对每个 $x\in \Ob(A)$，给定 $2$-态射
$\alpha_x : \text{id}_{\varphi(x)} \to \varphi(\text{id}_x)$；
\item 对 $\mathcal{A}$ 中每一对可复合态射 $f : x \to y$、
$g : y \to z$，给定 $2$-态射
$\alpha_{g, f} : \varphi(g \circ f) \to \varphi(g) \circ \varphi(f)$。
\end{enumerate}
这些数据须满足下列条件：
\begin{enumerate}
\item $2$-态射 $\alpha_x$ 与 $\alpha_{g, f}$ 都是同构；
\item 对 $\mathcal{A}$ 中任意态射 $f : x \to y$，有
$\alpha_{\text{id}_y, f} = \alpha_y \star \text{id}_{\varphi(f)}$：
$$
\xymatrix{
\varphi(x)
\rrtwocell^{\varphi(f)}_{\varphi(f)}{\ \ \ \ \text{id}_{\varphi(f)}}
& &
\varphi(y)
\rrtwocell^{\text{id}_{\varphi(y)}}_{\varphi(\text{id}_y)}{\alpha_y}
& &
\varphi(y)
}
=
\xymatrix{
\varphi(x)
\rrtwocell^{\varphi(f)}_{\varphi(\text{id}_y) \circ \varphi(f)}{\ \ \ \ \alpha_{\text{id}_y, f}}
& &
\varphi(y)
}
$$
\item 对 $\mathcal{A}$ 中任意态射 $f : x \to y$，有
$\alpha_{f, \text{id}_x} = \text{id}_{\varphi(f)} \star \alpha_x$；
\item 对 $\mathcal{A}$ 中任意三个可复合态射
$f : w \to x$、$g : x \to y$、$h : y \to z$，有
$$
(\text{id}_{\varphi(h)} \star \alpha_{g, f})
\circ
\alpha_{h, g \circ f}
=
(\alpha_{h, g} \star \text{id}_{\varphi(f)})
\circ
\alpha_{h \circ g, f}
$$
换言之，下列以 $1$-态射为对象、以 $2$-态射为箭头的图表可交换：
$$
\xymatrix{
\varphi(h \circ g \circ f)
\ar[d]_{\alpha_{h, g \circ f}}
\ar[rr]_{\alpha_{h \circ g, f}}
& &
\varphi(h \circ g) \circ \varphi(f)
\ar[d]^{\alpha_{h, g} \star \text{id}_{\varphi(f)}} \\
\varphi(h) \circ \varphi(g \circ f)
\ar[rr]^{\text{id}_{\varphi(h)} \star \alpha_{g, f}}
& &
\varphi(h) \circ \varphi(g) \circ \varphi(f)
}
$$
\end{enumerate}
\end{enumerate}
\end{definition}

\noindent
这仍然不是很便于使用的概念，但有时确实会出现。有一个定理断言，
任意伪函子都同构于一个函子。最后，还有{\it $2$-范畴之间的函子}
与{\it $2$-范畴之间的伪函子}的概念。后一概念会把我们带入
$3$-范畴的领域。我们希望几乎不惜一切代价避免定义它！

\section{(2, 1)-范畴}
\label{categories-section-2-1-categories}

\noindent
某些 $2$-范畴具有所有 $2$-态射都是同构的性质。它们将在下文发挥
重要作用，而且更容易使用。

\begin{definition}
\label{categories-definition-2-1-category}
一个（严格）{\it $(2, 1)$-范畴}，是所有 $2$-态射都是同构的
$2$-范畴。
\end{definition}

\begin{example}
\label{categories-example-2-1-category-of-categories}
注 \ref{categories-remark-big-2-categories} 中的 $2$-范畴 $\textit{Cat}$，
若只允许函子同构作为 $2$-态射，就可成为一个 $(2, 1)$-范畴。
\end{example}

\noindent
事实上，更一般地，任意 $2$-范畴 $\mathcal{C}$ 都产生一个
$(2, 1)$-范畴：取具有相同对象与 $1$-态射的子 $2$-范畴
$\mathcal{C}'$，但其 $2$-态射是 $\mathcal{C}$ 的可逆 $2$-态射。
在这种情形下，我们会说“{\it 令 $\mathcal{C}'$ 为与
$\mathcal{C}$ 关联的 $(2, 1)$-范畴}”或类似的话。例如，
群胚的 $(2, 1)$-范畴是指如下 $2$-范畴：对象是群胚，
$1$-态射是函子，$2$-态射是函子同构。不过这是个不好的例子，
因为群胚之间的函子之间的变换自动就是同构！

\begin{remark}
\label{categories-remark-other-2-categories}
因此，例 \ref{categories-example-2-1-category-of-categories} 的上述构造还有一些
变体，其中考虑群胚的 $2$-范畴、固定范畴上的群胚纤维化范畴，
或栈，等等。
\end{remark}






\section{2-纤维积}
\label{categories-section-2-fibre-products}

\noindent
本节介绍 $2$-纤维积。设 $\mathcal{C}$ 是一个 2-范畴。若图表
$$
\xymatrix{
w \ar[r] \ar[d] & y \ar[d] \\
x \ar[r] & z }
$$
中的两个 1-态射 $w \to y \to z$ 与 $w \to x \to z$ 是
2-同构的，则称该图表是 2-交换的。在 2-范畴中，要求图表
2-交换，比要求其严格交换更为自然。（事实上，有人或许会主张
根本不应使用严格 2-范畴；而在“弱”2-范畴中，1-态射图表
严格交换这一概念甚至没有意义。）相应地，纤维积的概念也必须调整。

\medskip\noindent
令 $\mathcal{C}$ 为 $2$-范畴。令 $x, y, z\in \Ob(\mathcal{C})$，并令
$f\in \Mor_\mathcal{C}(x, z)$ 与 $g\in \Mor_{\mathcal C}(y, z)$。
为了定义 $f$ 与 $g$ 的 2-纤维积，我们将考察 2-交换图表
$$
\xymatrix{
& w \ar[r]_a \ar[d]_b & x \ar[d]^{f} \\
& y \ar[r]^{g} & z. }
$$
在普通范畴的情形，纤维积是这类图表所成范畴中的终对象。相应地，
2-纤维积是某个 2-范畴中的终对象（见下述定义）。{\it 位于
$f$ 与 $g$ 上方的 $2$-交换图表的 $2$-范畴}是定义如下的 $2$-范畴：
\begin{enumerate}
\item 对象是如上的四元组 $(w, a, b, \phi)$，其中 $\phi$
是可逆 2-态射 $\phi : f \circ a \to g \circ b$；
\item 从 $(w', a', b', \phi')$ 到 $(w, a, b, \phi)$ 的 1-态射由
$(k : w' \to w, \alpha : a' \to a \circ k, \beta : b' \to b \circ k)$
给出，并使
$$
\xymatrix{
f \circ a'
\ar[rr]_{\text{id}_f \star \alpha}
\ar[d]_{\phi'}
& &
f \circ a \circ k
\ar[d]^{\phi \star \text{id}_k}
\\
g \circ b'
\ar[rr]^{\text{id}_g \star \beta}
& &
g \circ b \circ k
}
$$
交换；
\item 给定第二个 $1$-态射
$(k', \alpha', \beta') : (w'', a'', b'', \phi'') \to
(w', \alpha', \beta', \phi')$，$1$-态射的复合由下式给出：
$$
(k, \alpha, \beta) \circ (k', \alpha', \beta') =
(k \circ k',
(\alpha \star \text{id}_{k'}) \circ \alpha',
(\beta \star \text{id}_{k'}) \circ \beta'),
$$
\item 具有相同源与靶的两个 $1$-态射
$(k_i, \alpha_i, \beta_i)$（$i = 1, 2$）之间的一个 2-态射，
由一个 2-态射 $\delta : k_1 \to k_2$ 给出，并使
$$
\xymatrix{
a'
\ar[rd]_{\alpha_2}
\ar[r]_{\alpha_1} &
a \circ k_1
\ar[d]^{\text{id}_a \star \delta} &
&
b \circ k_1
\ar[d]_{\text{id}_b \star \delta} &
b'
\ar[l]^{\beta_1}
\ar[ld]^{\beta_2}
\\
&
a \circ k_2
&
&
b \circ k_2
&
}
$$
交换；
\item $2$-态射的竖直复合，由 $\mathcal{C}$ 中态射 $\delta$ 的
竖直复合给出；并且
\item 图表
$$
\xymatrix{
(w'', a'', b'', \phi'')
\rrtwocell^{(k'_1, \alpha'_1, \beta'_1)}_{(k'_2, \alpha'_2, \beta'_2)}{\delta'}
& &
(w', a', b', \phi')
\rrtwocell^{(k_1, \alpha_1, \beta_1)}_{(k_2, \alpha_2, \beta_2)}{\delta}
& &
(w, a, b, \phi)
}
$$
的水平复合由图表
$$
\xymatrix@C=12pc{
(w'', a'', b'', \phi'')
\rtwocell^{(k_1 \circ k'_1, (\alpha_1 \star \text{id}_{k'_1}) \circ \alpha'_1, (\beta_1 \star \text{id}_{k'_1}) \circ \beta'_1)}_{(k_2 \circ k'_2, (\alpha_2 \star \text{id}_{k'_2}) \circ \alpha'_2, (\beta_2 \star \text{id}_{k'_2}) \circ \beta'_2)}{\ \ \ \delta \star \delta'}
&
(w, a, b, \phi)
}
$$
给出。
\end{enumerate}
注意，如果 $\mathcal{C}$ 实际上是一个 $(2, 1)$-范畴，则上述 (2) 中的
态射 $\alpha$ 与 $\beta$ 自动也是同构\footnote{事实上，在 $2$-范畴的
情形，似乎还可以定义另一个 2-交换图表的 2-范畴，其中箭头
$\alpha$、$\beta$ 的方向反转，甚至只反转其中一个箭头的方向。这就是
我们稍后限制到 $(2, 1)$-范畴的原因。}。此外，如果 $\mathcal{C}$ 是
$(2, 1)$-范畴，那么 $2$-交换图表的 $2$-范畴也是 $(2, 1)$-范畴。

\begin{definition}
\label{categories-definition-final-object-2-category}
$(2, 1)$-范畴 $\mathcal{C}$ 的一个{\it 终对象}，是满足下列条件的对象
$x$：
\begin{enumerate}
\item 对每个 $y \in \Ob(\mathcal{C})$，都存在一个态射 $y \to x$；并且
\item 任意两个态射 $y \to x$ 都由唯一的 2-态射彼此同构。
\end{enumerate}
\end{definition}

\noindent
在更一般的 $2$-范畴情形，终对象很可能有不同的变体。我们不想在此
深入讨论，因而只在 $(2, 1)$-情形定义 $2$-纤维积。





\begin{definition}
\label{categories-definition-2-fibre-products}
令 $\mathcal{C}$ 为一个 $(2, 1)$-范畴。令
$x, y, z\in \Ob(\mathcal{C})$，并令
$f\in \Mor_\mathcal{C}(x, z)$
与 $g\in \Mor_{\mathcal C}(y, z)$。所谓{\it $f$ 与 $g$ 的
2-纤维积}，是上述 2-交换图表所成范畴中的一个终对象。
若 2-纤维积存在，我们将其记为
$x \times_z y\in \Ob(\mathcal{C})$，并将使图表
$$
\xymatrix{
& x \times_z y \ar[r]^{p} \ar[d]_q & x \ar[d]^{f} \\
& y \ar[r]^{g} & z }
$$
2-交换的所需态射记为
$p\in \Mor_{\mathcal C}(x \times_z y, x)$ 与
$q\in \Mor_{\mathcal C}(x \times_z y, y)$，将给出这一
2-交换性的可逆 2-态射记为 $\psi : f \circ p \to g \circ q$。
\end{definition}

\noindent
于是下列泛性质成立：对任意
$w\in \Ob(\mathcal{C})$ 以及态射
$a \in \Mor_{\mathcal C}(w, x)$ 和
$b \in \Mor_\mathcal{C}(w, y)$，若给定一个 2-同构
$\phi : f \circ a \to g\circ b$，
则存在 $\gamma \in \Mor_{\mathcal C}(w, x \times_z y)$，使图表
$$
\xymatrix{
w\ar[rrrd]^a \ar@{-->}[rrd]_\gamma \ar[rrdd]_b & & \\
& & x \times_z y \ar[r]_p \ar[d]_q & x \ar[d]^{f} \\
& & y \ar[r]^{g} & z }
$$
2-交换，并且在适当选择
$a \to p \circ \gamma$ 与 $b \to q \circ \gamma$ 后，图表
$$
\xymatrix{
f \circ a \ar[r] \ar[d]_\phi &
f \circ p \circ \gamma
\ar[d]^{\psi \star \text{id}_\gamma}
\\
g\circ b
\ar[r]
&
g \circ q \circ \gamma
}
$$
交换。此外，$\gamma$ 在同构意义下唯一。当然，精确的泛性质比这更
细致。我们以后所需的全部 2-纤维积情形，都来自下例：范畴的
2-范畴中的 2-纤维积。

\begin{example}
\label{categories-example-2-fibre-product-categories}
令 $\mathcal{A}$、$\mathcal{B}$ 和 $\mathcal{C}$ 为范畴。
令 $F : \mathcal{A} \to \mathcal{C}$ 与
$G : \mathcal{B} \to \mathcal{C}$ 为函子。如下定义范畴
$\mathcal{A} \times_\mathcal{C} \mathcal{B}$：
\begin{enumerate}
\item $\mathcal{A} \times_\mathcal{C} \mathcal{B}$ 的一个对象是三元组
$(A, B, f)$，其中 $A\in \Ob(\mathcal{A})$、
$B\in \Ob(\mathcal{B})$，而
$f : F(A) \to G(B)$ 是 $\mathcal{C}$ 中的一个同构；
\item 一个态射 $(A, B, f) \to (A', B', f')$ 由一对 $(a, b)$ 给出，
其中 $a : A \to A'$ 是 $\mathcal{A}$ 中的态射，
$b : B \to B'$ 是 $\mathcal{B}$ 中的态射，并使图表
$$
\xymatrix{
F(A) \ar[r]^f \ar[d]^{F(a)} & G(B) \ar[d]^{G(b)} \\
F(A') \ar[r]^{f'} & G(B')
}
$$
交换。
\end{enumerate}
此外，我们通过令
$$
p(A, B, f) = A, \quad q(A, B, f) = B,
$$
定义函子
$p : \mathcal{A} \times_\mathcal{C}\mathcal{B} \to \mathcal{A}$
和
$q : \mathcal{A} \times_\mathcal{C}\mathcal{B} \to \mathcal{B}$；
换言之，它们是遗忘函子。定义函子变换
$\psi : F \circ p \to G \circ q$。在对象
$\xi = (A, B, f)$ 上，它由
$\psi_\xi = f : F(p(\xi)) = F(A) \to G(B) = G(q(\xi))$ 给出。
\end{example}

\begin{lemma}
\label{categories-lemma-2-fibre-product-categories}
在范畴的 $(2, 1)$-范畴中，$2$-纤维积存在，并由
例 \ref{categories-example-2-fibre-product-categories} 的构造给出。
\end{lemma}

\begin{proof}
我们来验证泛性质。令 $\mathcal{W}$ 为范畴，令
$a : \mathcal{W} \to \mathcal{A}$ 与
$b : \mathcal{W} \to \mathcal{B}$ 为函子，并令
$t : F \circ a \to G \circ b$ 为函子同构。

\medskip\noindent
考虑由 $W \mapsto (a(W), b(W), t_W)$ 给出的函子
$\gamma : \mathcal{W} \to \mathcal{A} \times_\mathcal{C}\mathcal{B}$。
（省略对它是函子的验证。）再考虑由恒等式
$p \circ \gamma = a$ 与 $q \circ \gamma = b$ 得到的
$\alpha : a \to p \circ \gamma$ 和
$\beta : b \to q \circ \gamma$。于是显然，
$(\gamma, \alpha, \beta)$ 是从 $(W, a, b, t)$ 到
$(\mathcal{A} \times_\mathcal{C} \mathcal{B}, p, q, \psi)$ 的态射。

\medskip\noindent
令
$(k, \alpha', \beta') :
(W, a, b, t) \to (\mathcal{A} \times_\mathcal{C} \mathcal{B}, p, q, \psi)$
为第二个这样的态射。对 $\mathcal{W}$ 的一个对象 $W$，记
$k(W) = (a_k(W), b_k(W), t_{k, W})$。于是
$p(k(W)) = a_k(W)$，其余类同。映射 $\alpha'$ 对应于函子性映射
$\alpha' : a(W) \to a_k(W)$。由于我们是在范畴的
$(2, 1)$-范畴中工作，实际上每个映射
$a(W) \to a_k(W)$ 都是同构。利用这些同构（以及对应的
$b(W) \to b_k(W)$），可得到同构
$$
\delta_W :
\gamma(W) = (a(W), b(W), t_W)
\longrightarrow
(a_k(W), b_k(W), t_{k, W}) = k(W).
$$
直接可证 $\delta$ 定义了 $\gamma$ 与 $k$ 之间的
$2$-同构，它位于所需的 $2$-交换图表的 $2$-范畴中。
\end{proof}


\begin{remark}
\label{categories-remark-other-description-2-fibre-product}
令 $\mathcal{A}$、$\mathcal{B}$ 和 $\mathcal{C}$ 为范畴。
令 $F : \mathcal{A} \to \mathcal{C}$ 与
$G : \mathcal{B} \to \mathcal{C}$ 为函子。如下还可给出
$2$-纤维积 $\mathcal{A} \times_\mathcal{C} \mathcal{B}$ 的另一种、
稍为对称的构造。一个对象是五元组 $(A, B, C, a, b)$，其中
$A, B, C$ 分别是 $\mathcal{A}, \mathcal{B}, \mathcal{C}$ 的对象，
而 $a : F(A) \to C$ 与 $b : G(B) \to C$ 是同构。一个态射
$(A, B, C, a, b) \to (A', B', C', a', b')$ 由与态射
$a, b, a', b'$ 相容的三个态射
$A \to A', B \to B', C \to C'$ 给出。可以直接证明这会给出一个
$2$-纤维积。不过，更容易的做法是注意到，函子
$(A, B, C, a, b) \mapsto (A, B, b^{-1} \circ a)$ 给出了从五元组的
范畴到例 \ref{categories-example-2-fibre-product-categories} 所构造范畴的等价。
\end{remark}

\begin{lemma}
\label{categories-lemma-functoriality-2-fibre-product}
设
$$
\xymatrix{
& \mathcal{Y} \ar[d]_I \ar[rd]^K & \\
\mathcal{X} \ar[r]^H \ar[rd]^L &
\mathcal{Z} \ar[rd]^M & \mathcal{B} \ar[d]^G \\
& \mathcal{A} \ar[r]^F & \mathcal{C}
}
$$
是范畴的一个 $2$-交换图表。选择同构
$\alpha : G \circ K \to M \circ I$ 和
$\beta : M \circ H \to F \circ L$，就确定了与此情形相关的
$2$-纤维积之间的一个态射
$$
\mathcal{X} \times_\mathcal{Z} \mathcal{Y}
\longrightarrow
\mathcal{A} \times_\mathcal{C} \mathcal{B}
$$
\end{lemma}

\begin{proof}
在对象上使用函子
$$
(X, Y, \phi) \longmapsto (L(X), K(Y),
\alpha^{-1}_Y \circ M(\phi) \circ \beta^{-1}_X)
$$
并在态射上使用
$$
(a, b) \longmapsto (L(a), K(b))
$$
即可。
\end{proof}

\begin{lemma}
\label{categories-lemma-equivalence-2-fibre-product}
假设同引理 \ref{categories-lemma-functoriality-2-fibre-product}。
\begin{enumerate}
\item 若 $K$ 与 $L$ 忠实，则态射
$\mathcal{X} \times_\mathcal{Z} \mathcal{Y} \to
\mathcal{A} \times_\mathcal{C} \mathcal{B}$
忠实。
\item 若 $K$ 与 $L$ 全忠实而 $M$ 忠实，则态射
$\mathcal{X} \times_\mathcal{Z} \mathcal{Y} \to
\mathcal{A} \times_\mathcal{C} \mathcal{B}$
全忠实。
\item 若 $K$ 与 $L$ 是等价而 $M$ 全忠实，则态射
$\mathcal{X} \times_\mathcal{Z} \mathcal{Y} \to
\mathcal{A} \times_\mathcal{C} \mathcal{B}$
是等价。
\end{enumerate}
\end{lemma}

\begin{proof}
令 $(X, Y, \phi)$ 与 $(X', Y', \phi')$ 为
$\mathcal{X} \times_\mathcal{Z} \mathcal{Y}$ 的对象。置
$Z = H(X)$，并通过 $\phi$ 将其与 $I(Y)$ 等同。还通过
$\alpha_X$ 将 $M(Z)$ 与 $F(L(X))$ 等同，通过
$\beta_Y$ 将 $M(Z)$ 与 $G(K(Y))$ 等同。对
$Z' = H(X')$ 与 $M(Z')$ 也作类似处理。态射上的映射是
$$
\xymatrix{
\Mor_\mathcal{X}(X, X')
\times_{\Mor_\mathcal{Z}(Z, Z')}
\Mor_\mathcal{Y}(Y, Y')
\ar[d] \\
\Mor_\mathcal{A}(L(X), L(X'))
\times_{\Mor_\mathcal{C}(M(Z), M(Z'))}
\Mor_\mathcal{B}(K(Y), K(Y'))
}
$$
这个映射。因此 (1) 与 (2) 成立。进一步，若 $K$ 与 $L$
是等价而 $M$ 全忠实，则任意对象 $(A, B, \phi)$ 都在本质像中，
理由如下：选取 $X$、$Y$，使
$L(X) \cong A$ 且 $K(Y) \cong B$。于是 $M$ 的全忠实性保证
我们可以找到一个同构 $H(X) \cong I(Y)$。略去若干细节。
\end{proof}

\begin{lemma}
\label{categories-lemma-associativity-2-fibre-product}
设
$$
\xymatrix{
\mathcal{A} \ar[rd] & & \mathcal{C} \ar[ld] \ar[rd] & & \mathcal{E} \ar[ld] \\
& \mathcal{B} & & \mathcal{D}
}
$$
是范畴和函子的一个图表。那么存在范畴的典范同构
$$
(\mathcal{A} \times_\mathcal{B} \mathcal{C}) \times_\mathcal{D} \mathcal{E}
\cong
\mathcal{A} \times_\mathcal{B} (\mathcal{C} \times_\mathcal{D} \mathcal{E})
$$
\end{lemma}

\begin{proof}
如果你明白这里的意思，只需使用函子
$$
((A, C, \phi), E, \psi)
\longmapsto
(A, (C, E, \psi), \phi)
$$
即可。
\end{proof}

\noindent
今后处理两个以上范畴的纤维积时，我们不再书写括号。


\begin{lemma}
\label{categories-lemma-triple-2-fibre-product-pr02}
设
$$
\xymatrix{
\mathcal{A} \ar[rd] & & \mathcal{C} \ar[ld] \ar[rd] & & \mathcal{E} \ar[ld] \\
& \mathcal{B} \ar[rd]_F & & \mathcal{D} \ar[ld]^G \\
& & \mathcal{F} &
}
$$
是范畴和函子的一个交换图表。那么存在范畴的典范函子
$$
\text{pr}_{02} :
\mathcal{A} \times_\mathcal{B} \mathcal{C} \times_\mathcal{D} \mathcal{E}
\longrightarrow
\mathcal{A} \times_\mathcal{F} \mathcal{E}
$$
\end{lemma}

\begin{proof}
若将
$\mathcal{A} \times_\mathcal{B} \mathcal{C}
\times_\mathcal{D} \mathcal{E}$
写成
$(\mathcal{A} \times_\mathcal{B} \mathcal{C})
\times_\mathcal{D} \mathcal{E}$，
那么，如果你明白我的意思，只需使用函子
$$
((A, C, \phi), E, \psi)
\longmapsto
(A, E, G(\psi) \circ F(\phi))
$$
即可。
\end{proof}

\begin{lemma}
\label{categories-lemma-2-fibre-product-erase-factor}
设
$$
\mathcal{A} \to
\mathcal{B} \leftarrow \mathcal{C} \leftarrow \mathcal{D}
$$
是范畴和函子的一个图表。那么存在范畴的典范同构
$$
\mathcal{A} \times_\mathcal{B} \mathcal{C} \times_\mathcal{C} \mathcal{D}
\cong
\mathcal{A} \times_\mathcal{B} \mathcal{D}
$$
\end{lemma}

\begin{proof}
略。
\end{proof}

\noindent
我们断言，这意味着可以像使用普通纤维积一样使用这些 $2$-纤维积。
下面是稍后确实会用到的一些进一步引理。

\begin{lemma}
\label{categories-lemma-diagonal-1}
设
$$
\xymatrix{
\mathcal{C}_3 \ar[r] \ar[d] & \mathcal{S} \ar[d]^\Delta \\
\mathcal{C}_1 \times \mathcal{C}_2 \ar[r]^{G_1 \times G_2} &
\mathcal{S} \times \mathcal{S}
}
$$
是范畴的一个 $2$-纤维积。那么存在典范同构
$\mathcal{C}_3 \cong
\mathcal{C}_1 \times_{G_1, \mathcal{S}, G_2} \mathcal{C}_2$。
\end{lemma}

\begin{proof}
可以假设 $\mathcal{C}_3$ 是例
\ref{categories-example-2-fibre-product-categories} 中构造的范畴
$(\mathcal{C}_1 \times \mathcal{C}_2)\times_{\mathcal{S} \times \mathcal{S}}
\mathcal{S}$。于是，一个对象是三元组
$((X_1, X_2), S, \phi)$，其中
$\phi = (\phi_1, \phi_2) : (G_1(X_1), G_2(X_2)) \to (S, S)$
是同构。因此可以将三元组
$(X_1, X_2, \phi_2^{-1} \circ \phi_1)$ 与之对应。反过来，若
$(X_1, X_2, \psi)$ 是
$\mathcal{C}_1 \times_{G_1, \mathcal{S}, G_2} \mathcal{C}_2$ 的对象，
则可以将三元组
$((X_1, X_2), G_2(X_2), (\psi, \text{id}_{G_2(X_2)}))$
与之对应。我们断言，这两个构造给出互逆函子。略去对态射的处理，
以及证明它们互逆的过程。
\end{proof}

\begin{lemma}
\label{categories-lemma-diagonal-2}
设
$$
\xymatrix{
\mathcal{C}' \ar[r] \ar[d] & \mathcal{S} \ar[d]^\Delta \\
\mathcal{C} \ar[r]^{(G_1, G_2)} &
\mathcal{S} \times \mathcal{S}
}
$$
是范畴的一个 $2$-纤维积。那么存在典范同构
$$
\mathcal{C}' \cong
(\mathcal{C} \times_{G_1, \mathcal{S}, G_2} \mathcal{C})
\times_{(p, q), \mathcal{C} \times \mathcal{C}, \Delta}
\mathcal{C}.
$$
\end{lemma}

\begin{proof}
右端的一个对象由
$((C_1, C_2, \phi), C_3, \psi)$ 给出，其中
$\phi : G_1(C_1) \to G_2(C_2)$ 是同构，而
$\psi = (\psi_1, \psi_2) : (C_1, C_2) \to (C_3, C_3)$
是同构。因此可以将三元组
$(C_3, G_1(C_1), (G_1(\psi_1^{-1}), \phi^{-1} \circ G_2(\psi_2^{-1})))$
与之对应；它是 $\mathcal{C}'$ 的对象。略去细节。
\end{proof}

\begin{lemma}
\label{categories-lemma-fibre-product-after-map}
令 $\mathcal{A} \to \mathcal{C}$、$\mathcal{B} \to \mathcal{C}$
以及 $\mathcal{C} \to \mathcal{D}$ 为范畴之间的函子。那么图表
$$
\xymatrix{
\mathcal{A} \times_\mathcal{C} \mathcal{B} \ar[d] \ar[r] &
\mathcal{A} \times_\mathcal{D} \mathcal{B} \ar[d] \\
\mathcal{C} \ar[r]^-{\Delta_{\mathcal{C}/\mathcal{D}}} \ar[r] &
\mathcal{C} \times_\mathcal{D} \mathcal{C}
}
$$
是一个 $2$-纤维积图表。
\end{lemma}

\begin{proof}
略。
\end{proof}

\begin{lemma}
\label{categories-lemma-base-change-diagonal}
设
$$
\xymatrix{
\mathcal{U} \ar[d] \ar[r] & \mathcal{V} \ar[d] \\
\mathcal{X} \ar[r] & \mathcal{Y}
}
$$
是范畴的一个 $2$-纤维积。那么图表
$$
\xymatrix{
\mathcal{U} \ar[d] \ar[r] &
\mathcal{U} \times_\mathcal{V} \mathcal{U} \ar[d] \\
\mathcal{X} \ar[r] &
\mathcal{X} \times_\mathcal{Y} \mathcal{X}
}
$$
是 $2$-笛卡尔的。
\end{lemma}

\begin{proof}
这是一个纯粹的 $2$-范畴论陈述，在任何具有 $2$-纤维积的
$(2, 1)$-范畴中都成立。具体地说，它由下列等价链推出：
\begin{align*}
\mathcal{X} \times_{(\mathcal{X} \times_\mathcal{Y} \mathcal{X})}
(\mathcal{U} \times_\mathcal{V} \mathcal{U})
& =
\mathcal{X} \times_{(\mathcal{X} \times_\mathcal{Y} \mathcal{X})}
((\mathcal{X} \times_\mathcal{Y} \mathcal{V})
\times_\mathcal{V} (\mathcal{X} \times_\mathcal{Y} \mathcal{V})) \\
& =
\mathcal{X} \times_{(\mathcal{X} \times_\mathcal{Y} \mathcal{X})}
(\mathcal{X} \times_\mathcal{Y} \mathcal{X}
\times_\mathcal{Y} \mathcal{V}) \\
& =
\mathcal{X} \times_\mathcal{Y} \mathcal{V} = \mathcal{U}
\end{align*}
参见引理 \ref{categories-lemma-associativity-2-fibre-product} 和
\ref{categories-lemma-2-fibre-product-erase-factor}。
\end{proof}


\section{范畴上的范畴}
\label{categories-section-categories-over-categories}

\noindent
本节给定一个函子 $p : \mathcal{S} \to \mathcal{C}$。我们把
$\mathcal{S}$ 看作位于上方，把 $\mathcal{C}$ 看作位于下方。
为了明确所谈的对象，我们定义 $\mathcal{C}$ 上的范畴所成的
$2$-范畴。

\begin{definition}
\label{categories-definition-categories-over-C}
令 $\mathcal{C}$ 为范畴。所谓{\it $\mathcal{C}$ 上的范畴所成的
$2$-范畴}，是按如下方式定义的 $2$-范畴：
\begin{enumerate}
\item 它的对象是函子 $p : \mathcal{S} \to \mathcal{C}$。
\item 它的 $1$-态射
$(\mathcal{S}, p) \to (\mathcal{S}', p')$ 是满足
$p' \circ G = p$ 的函子 $G : \mathcal{S} \to \mathcal{S}'$。
\item 对
$G, H : (\mathcal{S}, p) \to (\mathcal{S}', p')$，它的 $2$-态射
$t : G \to H$ 是满足下述条件的函子态射：
对所有 $x \in \Ob(\mathcal{S})$，均有
$p'(t_x) = \text{id}_{p(x)}$。
\end{enumerate}
在这种情形，我们将 $(\mathcal{S}, p)$ 与
$(\mathcal{S}', p')$ 之间 $1$-态射所成的范畴记为
$$
\Mor_{\textit{Cat}/\mathcal{C}}(\mathcal{S}, \mathcal{S}')
$$
\end{definition}

\noindent
在这个 $2$-范畴中，我们完全按照
第 \ref{categories-section-formal-cat-cat} 节中对 $\textit{Cat}$ 的做法，
定义水平复合与竖直复合。$2$-范畴的公理成立，理由与它们在
$\textit{Cat}$ 中成立的理由相同。也可以利用这些公理在
$\textit{Cat}$ 中成立这一事实，再验证诸如“$\mathcal{C}$ 上的
$2$-态射的竖直复合仍给出 $\mathcal{C}$ 上的 $2$-态射”之类的
结论。此事显然。

\medskip\noindent
类似于空间之间映射的纤维，这里有纤维范畴的概念，以及与此情形
相关的一些提升概念。

\begin{definition}
\label{categories-definition-fibre-category}
令 $\mathcal{C}$ 为范畴。令
$p : \mathcal{S} \to \mathcal{C}$ 为 $\mathcal{C}$ 上的范畴。
\begin{enumerate}
\item 对象 $U\in \Ob(\mathcal{C})$ 上的{\it 纤维范畴}，是具有对象
$$
\Ob(\mathcal{S}_U) = \{x\in \Ob(\mathcal{S}) :
p(x) = U\}
$$
和态射
$$
\Mor_{\mathcal{S}_U}(x, y) = \{ \phi \in \Mor_\mathcal{S}(x, y) :
p(\phi) = \text{id}_U\}.
$$
的范畴 $\mathcal{S}_U$。
\item 对象 $U \in \Ob(\mathcal{C})$ 的一个{\it 提升}，是满足
$p(x) = U$ 的对象 $x\in \Ob(\mathcal{S})$；也就是说，
$x\in \Ob(\mathcal{S}_U)$。有时也说{\it $x$ 位于 $U$ 之上}。
\item 类似地，$\mathcal{C}$ 中态射 $f : V \to U$ 的一个
{\it 提升}，是 $\mathcal{S}$ 中满足 $p(\phi) = f$ 的态射
$\phi : y \to x$。有时也说{\it $\phi$ 位于 $f$ 之上}。
\end{enumerate}
\end{definition}

\noindent
这里可以作出一些观察。例如，若
$F : (\mathcal{S}, p) \to (\mathcal{S}', p')$ 是
$\mathcal{C}$ 上范畴的一个 $1$-态射，则 $F$ 诱导纤维范畴之间的
函子 $F : \mathcal{S}_U \to \mathcal{S}'_U$。对 $2$-态射亦同。

\medskip\noindent
下面给出必不可少的引理，它描述 $\mathcal{C}$ 上范畴所成的
$(2, 1)$-范畴中的 $2$-纤维积。

\begin{lemma}
\label{categories-lemma-2-product-categories-over-C}
令 $\mathcal{C}$ 为范畴。$\mathcal{C}$ 上的范畴所成的
$(2, 1)$-范畴具有 2-纤维积。设
$F : \mathcal{X} \to \mathcal{S}$ 与
$G : \mathcal{Y} \to \mathcal{S}$ 是 $\mathcal{C}$ 上范畴的态射。
一个显式的 2-纤维积
$\mathcal{X} \times_\mathcal{S}\mathcal{Y}$ 由下述描述给出：
\begin{enumerate}
\item $\mathcal{X} \times_\mathcal{S} \mathcal{Y}$ 的一个对象是四元组
$(U, x, y, f)$，其中 $U \in \Ob(\mathcal{C})$、
$x\in \Ob(\mathcal{X}_U)$、$y\in \Ob(\mathcal{Y}_U)$，且
$f : F(x) \to G(y)$ 是 $\mathcal{S}_U$ 中的同构；
\item 一个态射
$(U, x, y, f) \to (U', x', y', f')$ 由一对 $(a, b)$ 给出，
其中 $a : x \to x'$ 是 $\mathcal{X}$ 中的态射，
$b : y \to y'$ 是 $\mathcal{Y}$ 中的态射，并满足：
\begin{enumerate}
\item $a$ 与 $b$ 诱导同一个态射 $U \to U'$；并且
\item 图表
$$
\xymatrix{
F(x) \ar[r]^f \ar[d]^{F(a)} & G(y) \ar[d]^{G(b)} \\
F(x') \ar[r]^{f'} & G(y')
}
$$
交换。
\end{enumerate}
\end{enumerate}
在这种情形，函子
$p : \mathcal{X} \times_\mathcal{S}\mathcal{Y} \to \mathcal{X}$
和 $q : \mathcal{X} \times_\mathcal{S}\mathcal{Y} \to \mathcal{Y}$
是遗忘函子。变换 $\psi : F \circ p \to G \circ q$ 在对象
$\xi = (U, x, y, f)$ 上由
$\psi_\xi = f : F(p(\xi)) = F(x) \to G(y) = G(q(\xi))$ 给出。
\end{lemma}

\begin{proof}
我们来验证泛性质。令
$p_\mathcal{W} : \mathcal{W}\to \mathcal{C}$ 为
$\mathcal{C}$ 上的范畴，令
$X : \mathcal{W} \to \mathcal{X}$ 和
$Y : \mathcal{W} \to \mathcal{Y}$ 为 $\mathcal{C}$ 上的函子，
并令 $t : F \circ X \to G \circ Y$ 为 $\mathcal{C}$ 上的函子同构。
所需函子
$\gamma : \mathcal{W} \to \mathcal{X} \times_\mathcal{S} \mathcal{Y}$
由 $W \mapsto (p_\mathcal{W}(W), X(W), Y(W), t_W)$ 给出。
略去细节；比较引理 \ref{categories-lemma-2-fibre-product-categories}。
\end{proof}

\begin{example}
\label{categories-example-different-2-fibre-products}
引理 \ref{categories-lemma-2-product-categories-over-C} 中给出的范畴上的范畴之
$2$-纤维积构造，与引理 \ref{categories-lemma-2-fibre-product-categories}
中给出的范畴之 2-纤维积构造（如例
\ref{categories-example-2-fibre-product-categories}），一般会产生不等价的
结果。具体地，令 $\mathcal{S}$ 为仅有一个对象和两个箭头的群胚，
令 $\mathcal{X}$ 为仅有一个对象的离散范畴。将
$2$-纤维积 $\mathcal{X} \times_\mathcal{S} \mathcal{X}$
取作范畴，会得到含两个对象的离散范畴。然而，如果把它们全都视为
$\mathcal{S}$ 上的范畴，那么作为 $\mathcal{S}$ 上范畴的
$2$-纤维积 $\mathcal{X} \times_\mathcal{S} \mathcal{X}$
是仅有一个对象的离散范畴。区别在于（沿用引理
\ref{categories-lemma-2-product-categories-over-C} 的记号），在第一种情形
可以选择任意比较同构 $f$，而在第二种情形只能选择恒等箭头。
\end{example}

\begin{lemma}
\label{categories-lemma-fibre-2-fibre-product-categories-over-C}
令 $\mathcal{C}$ 为范畴。令
$f : \mathcal{X} \to \mathcal{S}$ 和
$g : \mathcal{Y} \to \mathcal{S}$ 为 $\mathcal{C}$ 上范畴的态射。
对 $\mathcal{C}$ 的任意对象 $U$，纤维范畴有下列恒等式：
$$
\left(\mathcal{X} \times_\mathcal{S}\mathcal{Y}\right)_U
=
\mathcal{X}_U \times_{\mathcal{S}_U} \mathcal{Y}_U
$$
\end{lemma}

\begin{proof}
略。
\end{proof}


\section{纤维化范畴}
\label{categories-section-fibred-categories}

\noindent
这里有必要非常简略地讨论纤维化范畴。

\medskip\noindent
令 $p : \mathcal{S} \to \mathcal{C}$ 为 $\mathcal{C}$ 上的范畴。
给定满足 $p(x) = U$ 的对象 $x \in \mathcal{S}$，以及态射
$f : V \to U$，可以尝试取某种“纤维积 $V \times_U x$”
（或者说，通过 $V \to U$ 对 $x$ 作一次{\it 基变换}）。具体地，
从对象 $z \in \mathcal{S}$ 到“$V \times_U x$”的一个态射，
应由一对 $(\varphi, g)$ 给出，其中
$\varphi : z \to x$、$g : p(z) \to V$，且
$p(\varphi) = f \circ g$。图示如下：
$$
\xymatrix{
z \ar@{~>}[d]^p \ar@{-}[r] &
? \ar[r] \ar@{~>}[d]^p &
x \ar@{~>}[d]^p \\
p(z) \ar[r] & V \ar[r]^f & U
}
$$
若这样的态射 $V \times_U x \to x$ 存在，则称其为强笛卡尔态射。

\begin{definition}
\label{categories-definition-cartesian-over-C}
令 $\mathcal{C}$ 为范畴。令
$p : \mathcal{S} \to \mathcal{C}$ 为 $\mathcal{C}$ 上的范畴。
一个{\it 强笛卡尔态射}，更精确地说，一个
{\it 强 $\mathcal{C}$-笛卡尔态射}，是 $\mathcal{S}$ 的态射
$\varphi : y \to x$，使得对每个 $z \in \Ob(\mathcal{S})$，
由 $\psi \longmapsto (\varphi \circ \psi, p(\psi))$ 给出的映射
$$
\Mor_\mathcal{S}(z, y)
\longrightarrow
\Mor_\mathcal{S}(z, x)
\times_{\Mor_\mathcal{C}(p(z), p(x))}
\Mor_\mathcal{C}(p(z), p(y)),
$$
是双射。
\end{definition}

\noindent
注意，由米田引理 \ref{categories-lemma-yoneda}，给定一个位于
$U \in \Ob(\mathcal{C})$ 之上的对象
$x \in \Ob(\mathcal{S})$，以及 $\mathcal{C}$ 的态射
$f : V \to U$，如果存在满足 $p(\varphi) = f$ 的强笛卡尔态射
$\varphi : y \to x$，那么 $(y, \varphi)$ 在唯一同构意义下唯一。
这从上述定义即显然，因为函子
$$
z
\longmapsto
\Mor_\mathcal{S}(z, x)
\times_{\Mor_\mathcal{C}(p(z), U)}
\Mor_\mathcal{C}(p(z), V)
$$
只依赖于数据 $(x, U, f : V \to U)$。因此，我们有时会用
$V \times_U x \to x$ 或 $f^*x \to x$ 表示作为 $f$ 的提升的
强笛卡尔态射。

\begin{lemma}
\label{categories-lemma-composition-cartesian}
令 $\mathcal{C}$ 为范畴。令
$p : \mathcal{S} \to \mathcal{C}$ 为 $\mathcal{C}$ 上的范畴。
\begin{enumerate}
\item 两个强笛卡尔态射的复合是强笛卡尔的。
\item $\mathcal{S}$ 的任意同构都是强笛卡尔的。
\item 若强笛卡尔态射 $\varphi$ 的 $p(\varphi)$ 是同构，则
该态射是同构。
\end{enumerate}
\end{lemma}

\begin{proof}
证明 (1)。令 $\varphi : y \to x$ 与 $\psi : z \to y$ 为
强笛卡尔态射。令 $t$ 为 $\mathcal{S}$ 的任意对象。于是
\begin{align*}
& \Mor_\mathcal{S}(t, z) \\
& =
\Mor_\mathcal{S}(t, y)
\times_{\Mor_\mathcal{C}(p(t), p(y))}
\Mor_\mathcal{C}(p(t), p(z)) \\
& =
\Mor_\mathcal{S}(t, x)
\times_{\Mor_\mathcal{C}(p(t), p(x))}
\Mor_\mathcal{C}(p(t), p(y))
\times_{\Mor_\mathcal{C}(p(t), p(y))}
\Mor_\mathcal{C}(p(t), p(z)) \\
& =
\Mor_\mathcal{S}(t, x)
\times_{\Mor_\mathcal{C}(p(t), p(x))}
\Mor_\mathcal{C}(p(t), p(z))
\end{align*}
故 $z \to x$ 是强笛卡尔的。

\medskip\noindent
证明 (2)。令 $y \to x$ 为同构。那么 $p(y) \to p(x)$ 也是同构。
因此
$\Mor_\mathcal{C}(p(z), p(y)) \to
\Mor_\mathcal{C}(p(z), p(x))$
是双射。于是
$\Mor_\mathcal{S}(z, x)
\times_{\Mor_\mathcal{C}(p(z), p(x))}
\Mor_\mathcal{C}(p(z), p(y))$
与 $\Mor_\mathcal{S}(z, x)$ 双射。因此，定义
\ref{categories-definition-cartesian-over-C} 中展示的映射是双射，因为
$y \to x$ 是同构；由此可知 $y \to x$ 是强笛卡尔的。

\medskip\noindent
证明 (3)。设 $\varphi : y \to x$ 是强笛卡尔的，且
$p(\varphi) : p(y) \to p(x)$ 是同构。将定义用于
$z = x$，可知 $(\text{id}_x, p(\varphi)^{-1})$ 来自唯一态射
$\chi : x \to y$。略去对 $\chi$ 是 $\varphi$ 之逆的验证。
\end{proof}

\begin{lemma}
\label{categories-lemma-cartesian-over-cartesian}
令 $F : \mathcal{A} \to \mathcal{B}$ 与
$G : \mathcal{B} \to \mathcal{C}$ 为范畴之间可复合的函子。
令 $x \to y$ 为 $\mathcal{A}$ 的态射。若 $x \to y$ 是强
$\mathcal{B}$-笛卡尔的，且 $F(x) \to F(y)$ 是强
$\mathcal{C}$-笛卡尔的，则 $x \to y$ 是强
$\mathcal{C}$-笛卡尔的。
\end{lemma}

\begin{proof}
这直接由定义推出。
\end{proof}

\begin{lemma}
\label{categories-lemma-strongly-cartesian-fibre-product}
令 $\mathcal{C}$ 为范畴。令
$p : \mathcal{S} \to \mathcal{C}$ 为 $\mathcal{C}$ 上的范畴。
令 $x \to y$ 与 $z \to y$ 为 $\mathcal{S}$ 的态射。假设
\begin{enumerate}
\item $x \to y$ 是强笛卡尔的；
\item $p(x) \times_{p(y)} p(z)$ 存在；并且
\item 在 $\mathcal{S}$ 中存在强笛卡尔态射 $a : w \to z$，满足
$p(w) = p(x) \times_{p(y)} p(z)$ 且
$p(a) = \text{pr}_2 : p(x) \times_{p(y)} p(z) \to p(z)$。
\end{enumerate}
那么纤维积 $x \times_y z$ 存在，并同构于 $w$。
\end{lemma}

\begin{proof}
由于 $x \to y$ 是强笛卡尔的，存在唯一态射
$b : w \to x$，使 $p(b) = \text{pr}_1$。为了说明 $w$ 是纤维积，
我们计算
\begin{align*}
& \Mor_\mathcal{S}(t, w) \\
& = \Mor_\mathcal{S}(t, z)
\times_{\Mor_\mathcal{C}(p(t), p(z))}
\Mor_\mathcal{C}(p(t), p(w)) \\
& = \Mor_\mathcal{S}(t, z)
\times_{\Mor_\mathcal{C}(p(t), p(z))}
(\Mor_\mathcal{C}(p(t), p(x))
\times_{\Mor_\mathcal{C}(p(t), p(y))}
\Mor_\mathcal{C}(p(t), p(z))) \\
& = \Mor_\mathcal{S}(t, z)
\times_{\Mor_\mathcal{C}(p(t), p(y))}
\Mor_\mathcal{C}(p(t), p(x)) \\
& = \Mor_\mathcal{S}(t, z)
\times_{\Mor_\mathcal{S}(t, y)}
\Mor_\mathcal{S}(t, y)
\times_{\Mor_\mathcal{C}(p(t), p(y))}
\Mor_\mathcal{C}(p(t), p(x)) \\
& = \Mor_\mathcal{S}(t, z)
\times_{\Mor_\mathcal{S}(t, y)}
\Mor_\mathcal{S}(t, x)
\end{align*}
正如所需。第一个等式成立，是因为 $a : w \to z$ 是强笛卡尔的；
最后一个等式成立，是因为 $x \to y$ 是强笛卡尔的。
\end{proof}

\begin{definition}
\label{categories-definition-fibred-category}
令 $\mathcal{C}$ 为范畴。令
$p : \mathcal{S} \to \mathcal{C}$ 为 $\mathcal{C}$ 上的范畴。
若对任意位于 $U \in \Ob(\mathcal{C})$ 之上的
$x \in \Ob(\mathcal{S})$，以及 $\mathcal{C}$ 的任意态射
$f : V \to U$，都存在位于 $f$ 之上的强笛卡尔态射
$f^*x \to x$，则称 $\mathcal{S}$ 是一个
{\it $\mathcal{C}$ 上的纤维化范畴}。
\end{definition}


\noindent
假设 $p : \mathcal{S} \to \mathcal{C}$ 是纤维化范畴。对定义中那样的
每个 $f : V \to U$ 和 $x\in \Ob(\mathcal{S}_U)$，可以选择一个
位于 $f$ 之上的强笛卡尔态射 $f^\ast x \to x$。由选择公理，
可以同时为所有 $f: V \to U = p(x)$ 选择
$f^*x \to x$。我们断言，对 $\mathcal{S}_U$ 中的每个态射
$\phi : x \to x'$ 以及每个 $f : V \to U$，都存在唯一的
$\mathcal{S}_V$ 中的态射
$f^\ast \phi : f^\ast x \to f^\ast x'$，使图表
$$
\xymatrix{
f^\ast x \ar[r]_{f^\ast \phi} \ar[d] & f^\ast x' \ar[d] \\
x \ar[r]^{\phi} & x' }
$$
交换。事实上，该箭头存在且唯一，是因为 $f^*x' \to x'$
是强笛卡尔的。这个箭头的唯一性保证 $f^\ast$（现在也已在态射上
定义）是函子 $ f^\ast : \mathcal{S}_U \to \mathcal{S}_V$。

\begin{definition}
\label{categories-definition-pullback-functor-fibred-category}
假设 $p : \mathcal{S} \to \mathcal{C}$ 是纤维化范畴。
\begin{enumerate}
\item 对 $p : \mathcal{S} \to \mathcal{C}$ 的一个
{\it 拉回选择}\footnote{这很可能不是标准术语。在一些文献中，这被
称为“cleavage”，但它会唤起错误的形象。也许“cleaving”会是更好的
词。一个相关概念是“splitting”，但在许多文献中，“splitting”是指
满足如下条件的拉回选择：对任意一对可复合态射，都有
$g^*f^* = (f \circ g)^*$。另见定义
\ref{categories-definition-split-fibred-category}。}
是指：对 $\mathcal{C}$ 的任意态射 $f: V \to U$ 和任意
$x \in \Ob(\mathcal{S}_U)$，选择一个位于 $f$ 之上的强笛卡尔态射
$f^\ast x \to x$。
\item 给定一个拉回选择，对 $\mathcal{C}$ 的任意态射
$f : V \to U$，上述函子
$f^* : \mathcal{S}_U \to \mathcal{S}_V$ 称为一个
{\it 拉回函子}（它与上面作出的选择 $f^*x \to x$ 相关）。
\end{enumerate}
\end{definition}

\noindent
当然，总可以假设我们的拉回选择具有
$\text{id}_U^*x = x$ 这一性质，不过在实践中，如果不对拉回施加
进一步假设，这个性质没有用处。

\begin{lemma}
\label{categories-lemma-fibred}
假设 $p : \mathcal{S} \to \mathcal{C}$ 是纤维化范畴，并且已经
给定 $p : \mathcal{S} \to \mathcal{C}$ 的一个拉回选择。
\begin{enumerate}
\item 对任意一对可复合态射 $f : V \to U$、
$g : W \to V$，存在唯一的函子同构
$$
\alpha_{g, f} :
(f \circ g)^\ast
\longrightarrow
g^\ast \circ f^\ast
$$
（作为函子 $\mathcal{S}_U \to \mathcal{S}_W$），使得对每个
$y\in \Ob(\mathcal{S}_U)$，下图交换：
$$
\xymatrix{
g^\ast f^\ast y \ar[r]
&
f^\ast y \ar[d] \\
(f \circ g)^\ast y \ar[r]
\ar[u]^{(\alpha_{g, f})_y}
&
y
}
$$
\item 若 $f = \text{id}_U$，则存在函子
$\mathcal{S}_U \to \mathcal{S}_U$ 的典范同构
$\alpha_U : \text{id} \to (\text{id}_U)^*$。
\item 四元组
$(U \mapsto \mathcal{S}_U, f \mapsto f^*, \alpha_{g, f}, \alpha_U)$
定义了从 $\mathcal{C}^{opp}$ 到范畴的 $(2, 1)$-范畴的一个
伪函子；见定义 \ref{categories-definition-functor-into-2-category}。
\end{enumerate}
\end{lemma}

\begin{proof}
事实上，显然 (1) 中的交换图表唯一确定纤维范畴
$\mathcal{S}_W$ 中的态射 $(\alpha_{g, f})_y$。它是同构，因为态射
$(f \circ g)^*y \to y$ 与复合
$g^*f^*y \to f^*y \to y$ 都是提升 $f \circ g$ 的强笛卡尔态射
（见定义 \ref{categories-definition-cartesian-over-C} 后的讨论和引理
\ref{categories-lemma-composition-cartesian}）。同理，因为
$\text{id}_x : x \to x$ 显然在 $\text{id}_U$ 上方是强笛卡尔的
（其中 $U = p(x)$），所以存在同构
$(\alpha_U)_x : x \to (\text{id}_U)^*x$。（当然，我们本来可以
预先假设，只要 $f$ 是恒等态射就有 $f^*x = x$；但为了保持一般性，
最好不作此假设。）略去对这样得到的 $\alpha_{g, f}$ 和
$\alpha_U$ 是函子变换的验证，也略去对 (3) 的验证。
\end{proof}

\begin{lemma}
\label{categories-lemma-fibred-equivalent}
令 $\mathcal{C}$ 为范畴。令 $\mathcal{S}_1$、$\mathcal{S}_2$
为 $\mathcal{C}$ 上的范畴。假设 $\mathcal{S}_1$ 与
$\mathcal{S}_2$ 作为 $\mathcal{C}$ 上的范畴等价。那么
$\mathcal{S}_1$ 在 $\mathcal{C}$ 上纤维化，当且仅当
$\mathcal{S}_2$ 在 $\mathcal{C}$ 上纤维化。
\end{lemma}

\begin{proof}
将给定函子记为 $p_i : \mathcal{S}_i \to \mathcal{C}$。令
$F : \mathcal{S}_1 \to \mathcal{S}_2$、
$G : \mathcal{S}_2 \to \mathcal{S}_1$ 为 $\mathcal{C}$ 上的函子，
并令
$i : F \circ G \to \text{id}_{\mathcal{S}_2}$、
$j : G \circ F \to \text{id}_{\mathcal{S}_1}$ 为
$\mathcal{C}$ 上的函子同构。我们断言，在这种情形，$F$ 将强
笛卡尔态射映为强笛卡尔态射。事实上，设
$\varphi : y \to x$ 是 $\mathcal{S}_1$ 中的强笛卡尔态射。
令 $f : V \to U$ 等于 $p_1(\varphi)$。假设
$z' \in \Ob(\mathcal{S}_2)$，其中 $W = p_2(z')$，并且给定
$g : W \to V$ 和 $\psi' : z' \to F(x)$，使
$p_2(\psi') = f \circ g$。那么
$$
\psi = j \circ G(\psi') : G(z') \to G(F(x)) \to x
$$
是 $\mathcal{S}_1$ 中满足 $p_1(\psi) = f \circ g$ 的态射。
因此，由假设，存在唯一的位于 $g$ 之上的态射
$\xi : G(z') \to y$，使 $\psi = \varphi \circ \xi$。这又给出态射
$$
\xi' = F(\xi) \circ i^{-1} : z' \to F(G(z')) \to F(y)
$$
它位于 $g$ 之上，并满足 $\psi' = F(\varphi) \circ \xi'$。
略去对 $\xi'$ 唯一性的验证。
\end{proof}

\noindent
由引理 \ref{categories-lemma-fibred-equivalent} 可知，等价将强笛卡尔态射映为
强笛卡尔态射。但对 $\mathcal{C}$ 上纤维化范畴之间的任意函子，
这未必成立。因此，我们如下定义纤维化范畴所成的 $2$-范畴。

\begin{definition}
\label{categories-definition-fibred-categories-over-C}
令 $\mathcal{C}$ 为范畴。所谓{\it $\mathcal{C}$ 上纤维化范畴
所成的 $2$-范畴}，是 $\mathcal{C}$ 上范畴所成的 $2$-范畴
（见定义 \ref{categories-definition-categories-over-C}）的子 $2$-范畴，
定义如下：
\begin{enumerate}
\item 它的对象是纤维化范畴
$p : \mathcal{S} \to \mathcal{C}$。
\item 它的 $1$-态射
$(\mathcal{S}, p) \to (\mathcal{S}', p')$ 是满足
$p' \circ G = p$ 的函子 $G : \mathcal{S} \to \mathcal{S}'$，
并且 $G$ 将强笛卡尔态射映为强笛卡尔态射。
\item 对
$G, H : (\mathcal{S}, p) \to (\mathcal{S}', p')$，它的 $2$-态射
$t : G \to H$ 是满足如下条件的函子态射：对所有
$x \in \Ob(\mathcal{S})$，均有
$p'(t_x) = \text{id}_{p(x)}$。
\end{enumerate}
在这种情形，我们将 $(\mathcal{S}, p)$ 与
$(\mathcal{S}', p')$ 之间 $1$-态射所成的范畴记为
$$
\Mor_{\textit{Fib}/\mathcal{C}}(\mathcal{S}, \mathcal{S}')
$$
\end{definition}

\noindent
注意 $1$-态射所满足的条件。还要注意，这是真正的 $2$-范畴，
而不是 $(2, 1)$-范畴。因此，在取 $2$-纤维积时，我们先转到
与之关联的 $(2, 1)$-范畴。


\begin{lemma}
\label{categories-lemma-2-product-fibred-categories-over-C}
令 $\mathcal{C}$ 为范畴。$\mathcal{C}$ 上纤维化范畴所成的
$(2, 1)$-范畴具有 2-纤维积，它们如引理
\ref{categories-lemma-2-product-categories-over-C} 所述。
\end{lemma}

\begin{proof}
这里基本上需要证明的是：给定 $\mathcal{C}$ 上纤维化范畴的态射
$F : \mathcal{X} \to \mathcal{S}$ 和
$G : \mathcal{Y} \to \mathcal{S}$，引理
\ref{categories-lemma-2-product-categories-over-C} 中描述的范畴
$\mathcal{X} \times_\mathcal{S} \mathcal{Y}$ 是纤维化的。
我们来证明 $\mathcal{X} \times_\mathcal{S} \mathcal{Y}$
有足够多的强笛卡尔态射。设 $(U, x, y, \phi)$ 是
$\mathcal{X} \times_\mathcal{S} \mathcal{Y}$ 的对象，并设
$f : V \to U$ 是 $\mathcal{C}$ 中的态射。选择 $\mathcal{X}$
中位于 $f$ 之上的强笛卡尔态射 $a : f^*x \to x$，以及
$\mathcal{Y}$ 中位于 $f$ 之上的强笛卡尔态射
$b : f^*y \to y$。由假设，$F(a)$ 与 $G(b)$ 是强笛卡尔的。
由于 $\phi : F(x) \to G(y)$ 是同构，强笛卡尔态射的唯一性给出
唯一同构 $f^*\phi : F(f^*x) \to G(f^*y)$，使
$G(b) \circ f^*\phi = \phi \circ F(a)$。换言之，
$(a, b) : (V, f^*x, f^*y, f^*\phi) \to (U, x, y, \phi)$
是 $\mathcal{X} \times_\mathcal{S} \mathcal{Y}$ 中的态射。
略去对它是强笛卡尔态射的验证（也略去对这些事实上是仅有的强
笛卡尔态射的验证）。
\end{proof}

\begin{lemma}
\label{categories-lemma-cute}
令 $\mathcal{C}$ 为范畴，令 $U \in \Ob(\mathcal{C})$。
若 $p : \mathcal{S} \to \mathcal{C}$ 是纤维化范畴，并且
$p$ 经由 $p' : \mathcal{S} \to \mathcal{C}/U$ 分解，那么
$p' : \mathcal{S} \to \mathcal{C}/U$ 是纤维化范畴。
\end{lemma}

\begin{proof}
设 $\varphi : x' \to x$ 关于 $p$ 是强笛卡尔的。我们断言，
$\varphi$ 关于 $p'$ 也是强笛卡尔的。置
$g = p'(\varphi)$，则对某些态射 $f : V \to U$ 和
$f' : V' \to U$，有 $g : V'/U \to V/U$。令
$z \in \Ob(\mathcal{S})$，并置 $p'(z) = (W \to U)$。
为了证明 $\varphi$ 对 $p'$ 是强笛卡尔的，需要证明由
$\psi' \longmapsto (\varphi \circ \psi', p'(\psi'))$ 给出的映射
$$
\Mor_\mathcal{S}(z, x')
\longrightarrow
\Mor_\mathcal{S}(z, x)
\times_{\Mor_{\mathcal{C}/U}(W/U, V/U)}
\Mor_{\mathcal{C}/U}(W/U, V'/U),
$$
是双射。设给定右端的元素 $(\psi, h)$。特别地，
$g \circ h = p(\psi)$；由 $\varphi$ 是强笛卡尔的这一条件，
得到唯一态射 $\psi' : z \to x'$，使
$\psi = \varphi \circ \psi'$ 且 $p(\psi') = h$。好，现在
$p'(\psi') : W/U \to V/U$ 是一个态射，其对应映射
$W \to V$ 是 $h$，因而它作为 $\mathcal{C}/U$ 中的态射等于
$h$。所以 $\psi'$ 是映到给定有序对 $(\psi, h)$ 的唯一态射
$z \to x'$。这证明了断言。

\medskip\noindent
最后，设给定 $g : V'/U \to V/U$ 以及满足
$p'(x) = V/U$ 的 $x$。由于
$p : \mathcal{S} \to \mathcal{C}$ 是纤维化范畴，存在满足
$p(\varphi) = g$ 的强笛卡尔态射 $\varphi : x' \to x$。
同样由上述论证可得
$p'(\varphi) = g : V'/U \to V/U$。而且如上所见，态射
$\varphi$ 是强笛卡尔的。因此定义
\ref{categories-definition-fibred-category} 的条件均得到满足，结论即得。
\end{proof}

\begin{lemma}
\label{categories-lemma-fibred-over-fibred}
令 $\mathcal{A} \to \mathcal{B} \to \mathcal{C}$ 为范畴之间的函子。
若 $\mathcal{A}$ 在 $\mathcal{B}$ 上纤维化，而
$\mathcal{B}$ 在 $\mathcal{C}$ 上纤维化，则
$\mathcal{A}$ 在 $\mathcal{C}$ 上纤维化。
\end{lemma}

\begin{proof}
这由定义和引理 \ref{categories-lemma-cartesian-over-cartesian} 推出。
\end{proof}

\begin{lemma}
\label{categories-lemma-fibred-category-representable-goes-up}
令 $p : \mathcal{S} \to \mathcal{C}$ 为纤维化范畴。
令 $x \to y$ 与 $z \to y$ 为 $\mathcal{S}$ 的态射，其中
$x \to y$ 是强笛卡尔的。若
$p(x) \times_{p(y)} p(z)$ 存在，则 $x \times_y z$ 存在，
$p(x \times_y z) = p(x) \times_{p(y)} p(z)$，并且
$x \times_y z \to z$ 是强笛卡尔的。
\end{lemma}

\begin{proof}
选取位于
$\text{pr}_2 : p(x) \times_{p(y)} p(z) \to p(z)$ 之上的强
笛卡尔态射 $\text{pr}_2^*z \to z$。于是由引理
\ref{categories-lemma-strongly-cartesian-fibre-product}，
$\text{pr}_2^*z = x \times_y z$。
\end{proof}

\begin{lemma}
\label{categories-lemma-ameliorate-morphism-fibred-categories}
令 $\mathcal{C}$ 为范畴。令 $F : \mathcal{X} \to \mathcal{Y}$
为 $\mathcal{C}$ 上纤维化范畴的一个 $1$-态射。存在
$\mathcal{C}$ 上纤维化范畴的 $1$-态射
$$
\xymatrix{
\mathcal{X} \ar@<1ex>[r]^u &
\mathcal{X}' \ar[r]^v \ar@<1ex>[l]^w & \mathcal{Y}
}
$$
使 $F = v \circ u$，并且
\begin{enumerate}
\item $u : \mathcal{X} \to \mathcal{X}'$ 全忠实；
\item $w$ 是 $u$ 的左伴随；并且
\item $v : \mathcal{X}' \to \mathcal{Y}$ 是纤维化范畴。
\end{enumerate}
\end{lemma}

\begin{proof}
将结构函子记为 $p : \mathcal{X} \to \mathcal{C}$ 和
$q : \mathcal{Y} \to \mathcal{C}$。我们如下显式构造
$\mathcal{X}'$。$\mathcal{X}'$ 的一个对象是四元组
$(U, x, y, f)$，其中 $x \in \Ob(\mathcal{X}_U)$、
$y \in \Ob(\mathcal{Y}_U)$，而
$f : y \to F(x)$ 是 $\mathcal{Y}_U$ 中的态射。一个态射
$(a, b) : (U, x, y, f) \to (U', x', y', f')$ 由
$a : x \to x'$ 和 $b : y \to y'$ 给出，其中
$p(a) = q(b) : U \to U'$ 且 $f' \circ b = F(a) \circ f$。

\medskip\noindent
同时为 $p$ 和 $q$ 选定拉回，并使用相同记号表示它们。令
$(U, x, y, f)$ 为对象，令 $h : V \to U$ 为态射。考虑由给定的
强笛卡尔映射 $h^*x \to x$ 与 $h^*y \to y$ 所产生的态射
$c : (V, h^*x, h^*y, h^*f) \to (U, x, y, f)$。
我们断言，$c$ 在 $\mathcal{X}'$ 中关于 $\mathcal{C}$ 是强
笛卡尔的。事实上，设给定 $\mathcal{X}'$ 的一个对象
$(W, x', y', f')$，一个位于 $W \to U$ 之上的态射
$(a, b) : (W, x', y', f') \to (U, x, y, f)$，以及
$W \to U$ 经由 $h$ 的分解 $W \to V \to U$。由于
$h^*x \to x$ 和 $h^*y \to y$ 是强笛卡尔的，得到位于给定态射
$W \to V$ 之上的态射
$a' : x' \to h^*x$ 与 $b' : y' \to h^*y$。考虑图表
$$
\xymatrix{
y' \ar[d]_{f'} \ar[r] & h^*y \ar[r] \ar[d]_{h^*f} & y \ar[d]_f \\
F(x') \ar[r] & F(h^*x) \ar[r] & F(x)
}
$$
外部矩形和右方块交换。由于 $F$ 是纤维化范畴的一个 $1$-态射，
态射 $F(h^*x) \to F(x)$ 是强笛卡尔的。于是由强笛卡尔态射的
泛性质，左方块交换。这证明 $\mathcal{X}'$ 在 $\mathcal{C}$
上纤维化。

\medskip\noindent
函子 $u : \mathcal{X} \to \mathcal{X}'$ 由
$x \mapsto (p(x), x, F(x), \text{id})$ 给出。它是全忠实的。
函子 $\mathcal{X}' \to \mathcal{Y}$ 由
$(U, x, y, f) \mapsto y$ 给出。函子
$w : \mathcal{X}' \to \mathcal{X}$ 由
$(U, x, y, f) \mapsto x$ 给出。根据我们对
$\mathcal{X}'$ 中关于 $\mathcal{C}$ 的强笛卡尔态射的描述，
这些函子都是 $\mathcal{C}$ 上纤维化范畴的 $1$-态射。
$w$ 与 $u$ 的伴随性意味着
$$
\Mor_\mathcal{X}(x, x') =
\Mor_{\mathcal{X}'}((U, x, y, f), (p(x'), x', F(x'), \text{id})),
$$
这直接由定义推出。

\medskip\noindent
最后，需要证明 $\mathcal{X}' \to \mathcal{Y}$ 是纤维化范畴。
令 $c : y' \to y$ 为 $\mathcal{Y}$ 中的态射，并令
$(U, x, y, f)$ 为 $\mathcal{X}'$ 中位于 $y$ 之上的对象。
置 $V = q(y')$，并令 $h = q(c) : V \to U$。令
$a : h^*x \to x$ 与 $b : h^*y \to y$ 为覆盖 $h$ 的强笛卡尔态射。
由于 $F$ 是纤维化范畴的一个 $1$-态射，我们可以将
$h^*F(x) = F(h^*x)$ 等同，并取强笛卡尔态射
$F(a) : F(h^*x) \to F(x)$。由 $b : h^*y \to y$ 的泛性质，
存在 $\mathcal{Y}_V$ 中的态射 $c' : y' \to h^*y$，使
$c = b \circ c'$。我们断言
$$
(a, c) : (V, h^*x, y', h^*f \circ c') \longrightarrow (U, x, y, f)
$$
在 $\mathcal{X}'$ 中关于 $\mathcal{Y}$ 是强笛卡尔的。为此，
令 $(W, x_1, y_1, f_1)$ 为 $\mathcal{X}'$ 的对象，令
$(a_1, b_1) : (W, x_1, y_1, f_1) \to (U, x, y, f)$ 为态射，
并设对某个态射 $b_1' : y_1 \to y'$ 有
$b_1 = c \circ b_1'$。那么
$$
(a_1', b_1') : (W, x_1, y_1, f_1) \longrightarrow (V, h^*x, y', h^*f \circ c')
$$
（其中 $a_1' : x_1 \to h^*x$ 是位于给定态射
$q(b_1') : W \to V$ 之上的唯一态射，使
$a_1 = a \circ a_1'$）就是所需态射。
\end{proof}


\section{惯性}
\label{categories-section-inertia}

\noindent
给定范畴 $\mathcal{C}$ 上的纤维化范畴
$p : \mathcal{S} \to \mathcal{C}$ 与
$p' : \mathcal{S}' \to \mathcal{C}$，以及一个 $1$-态射
$F : \mathcal{S} \to \mathcal{S}'$，在 $\mathcal{C}$ 上纤维化范畴
所成的 $(2, 1)$-范畴中有对角态射
$$
\Delta = \Delta_{\mathcal{S}/\mathcal{S}'} :
\mathcal{S} \longrightarrow \mathcal{S} \times_{\mathcal{S}'} \mathcal{S}
$$

\begin{lemma}
\label{categories-lemma-inertia-fibred-category}
令 $\mathcal{C}$ 为范畴。令
$p : \mathcal{S} \to \mathcal{C}$ 与
$p' : \mathcal{S}' \to \mathcal{C}$ 为纤维化范畴。令
$F : \mathcal{S} \to \mathcal{S}'$ 为 $\mathcal{C}$ 上纤维化范畴的
一个 $1$-态射。考虑 $\mathcal{C}$ 上的范畴
$\mathcal{I}_{\mathcal{S}/\mathcal{S}'}$，其
\begin{enumerate}
\item 对象是有序对 $(x, \alpha)$，其中
$x \in \Ob(\mathcal{S})$，且 $\alpha : x \to x$ 是满足
$F(\alpha) = \text{id}$ 的自同构；
\item 态射 $(x, \alpha) \to (y, \beta)$ 由使图表
$$
\xymatrix{
x\ar[r]_\phi\ar[d]_\alpha &
y\ar[d]^{\beta} \\
x\ar[r]^\phi &
y \\
}
$$
交换的态射 $\phi : x \to y$ 给出；并且
\item 函子 $\mathcal{I}_{\mathcal{S}/\mathcal{S}'} \to \mathcal{C}$
由 $(x, \alpha) \mapsto p(x)$ 给出。
\end{enumerate}
那么
\begin{enumerate}
\item 在 $\mathcal{C}$ 上范畴所成的 $(2, 1)$-范畴中存在等价
$$
\mathcal{I}_{\mathcal{S}/\mathcal{S}'} \longrightarrow
\mathcal{S}
\times_{\Delta, (\mathcal{S} \times_{\mathcal{S}'} \mathcal{S}), \Delta}
\mathcal{S}
$$
并且
\item $\mathcal{I}_{\mathcal{S}/\mathcal{S}'}$ 是
$\mathcal{C}$ 上的纤维化范畴。
\end{enumerate}
\end{lemma}

\begin{proof}
由引理 \ref{categories-lemma-2-product-fibred-categories-over-C} 和
\ref{categories-lemma-fibred-equivalent}，(2) 从 (1) 推出。因此只需证明 (1)。
以下不再特别说明地使用引理
\ref{categories-lemma-2-product-fibred-categories-over-C} 中的
$2$-纤维积构造。特别地，
$\mathcal{S}
\times_{\Delta, (\mathcal{S} \times_{\mathcal{S}'} \mathcal{S}), \Delta}
\mathcal{S}$
的一个对象是三元组 $(x, y, (\iota, \kappa))$，其中 $x$ 和 $y$
是 $\mathcal{S}$ 的对象，而
$(\iota, \kappa) : (x, x, \text{id}_{F(x)}) \to
(y, y, \text{id}_{F(y)})$
是 $\mathcal{S} \times_{\mathcal{S}'} \mathcal{S}$ 中的同构。
这恰好意味着 $\iota, \kappa : x \to y$ 是同构，且
$F(\iota) = F(\kappa)$。考虑函子
$$
I_{\mathcal{S}/\mathcal{S}'}
\longrightarrow
\mathcal{S}
\times_{\Delta, (\mathcal{S} \times_{\mathcal{S}'} \mathcal{S}), \Delta}
\mathcal{S}
$$
它将左端的对象 $(x, \alpha)$ 送到右端的对象
$(x, x, (\alpha, \text{id}_x))$，并将左端的态射 $\phi$
送到右端的态射 $(\phi, \phi)$。我们断言，该态射的一个拟逆由函子
$$
\mathcal{S}
\times_{\Delta, (\mathcal{S} \times_{\mathcal{S}'} \mathcal{S}), \Delta}
\mathcal{S}
\longrightarrow
I_{\mathcal{S}/\mathcal{S}'}
$$
给出；它将左端的对象 $(x, y, (\iota, \kappa))$ 送到右端的对象
$(x, \kappa^{-1} \circ \iota)$，并将左端的态射
$(\phi, \phi') : (x, y, (\iota, \kappa)) \to
(z, w, (\lambda, \mu))$
送到态射 $\phi$。事实上，复合上述两个函子所诱导的
$I_{\mathcal{S}/\mathcal{S}'}$ 的自函子严格等于恒等函子；而在
$\mathcal{S}
\times_{\Delta, (\mathcal{S} \times_{\mathcal{S}'} \mathcal{S}), \Delta}
\mathcal{S}$
上诱导的自函子，通过自然同构
$$
(\text{id}_x, \kappa) :
(x, x, (\kappa^{-1} \circ \iota, \text{id}_x))
\longrightarrow
(x, y, (\iota, \kappa)).
$$
同构于恒等函子。略去若干细节。
\end{proof}

\begin{definition}
\label{categories-definition-inertia-fibred-category}
令 $\mathcal{C}$ 为范畴。
\begin{enumerate}
\item 令 $F : \mathcal{S} \to \mathcal{S}'$ 为 $\mathcal{C}$
上纤维化范畴的一个 $1$-态射。$\mathcal{S}$ 在
$\mathcal{S}'$ 上的{\it 相对惯性}，是引理
\ref{categories-lemma-inertia-fibred-category} 的纤维化范畴
$\mathcal{I}_{\mathcal{S}/\mathcal{S}'} \to \mathcal{C}$。
\item 所谓 $\mathcal{S}$ 的{\it 惯性纤维化范畴
$\mathcal{I}_\mathcal{S}$}，是指
$\mathcal{I}_\mathcal{S} = \mathcal{I}_{\mathcal{S}/\mathcal{C}}$。
\end{enumerate}
\end{definition}

\noindent
注意，$\mathcal{C}$ 上的纤维化范畴之间有典范 $1$-态射
\begin{equation}
\label{categories-equation-inertia-structure-map}
\mathcal{I}_{\mathcal{S}/\mathcal{S}'} \longrightarrow \mathcal{S}
\quad\text{且}\quad
\mathcal{I}_\mathcal{S} \longrightarrow \mathcal{S}
\end{equation}
按照引理 \ref{categories-lemma-inertia-fibred-category} 的描述，它们只是将对象
$(x, \alpha)$ 映到对象 $x$，并将态射
$\phi : (x, \alpha) \to (y, \beta)$ 映到态射
$\phi : x \to y$。还有一个{\it 中性截面}
\begin{equation}
\label{categories-equation-neutral-section}
e : \mathcal{S} \to \mathcal{I}_{\mathcal{S}/\mathcal{S}'}
\quad\text{且}\quad
e : \mathcal{S} \to \mathcal{I}_\mathcal{S}
\end{equation}
由规则 $x \mapsto (x, \text{id}_x)$ 和
$(\phi : x \to y) \mapsto \phi$ 定义。它是
(\ref{categories-equation-inertia-structure-map}) 的右逆。给定一个
$2$-交换方块
$$
\xymatrix{
\mathcal{S}_1 \ar[d]_{F_1} \ar[r]_G & \mathcal{S}_2 \ar[d]^{F_2} \\
\mathcal{S}'_1 \ar[r]^{G'} & \mathcal{S}'_2
}
$$
有{\it 函子性映射}
\begin{equation}
\label{categories-equation-functorial}
\mathcal{I}_{\mathcal{S}_1/\mathcal{S}'_1}
\longrightarrow
\mathcal{I}_{\mathcal{S}_2/\mathcal{S}'_2}
\quad\text{且}\quad
\mathcal{I}_{\mathcal{S}_1}
\longrightarrow
\mathcal{I}_{\mathcal{S}_2}
\end{equation}
由规则 $(x, \alpha) \mapsto (G(x), G(\alpha))$ 和
$\phi \mapsto G(\phi)$ 定义。特别地，总有比较映射
\begin{equation}
\label{categories-equation-comparison}
\mathcal{I}_{\mathcal{S}/\mathcal{S}'}
\longrightarrow
\mathcal{I}_\mathcal{S}
\end{equation}
而且上述所有映射均与之相容。

\begin{lemma}
\label{categories-lemma-relative-inertia-as-fibre-product}
令 $F : \mathcal{S} \to \mathcal{S}'$ 为在范畴
$\mathcal{C}$ 上纤维化的范畴之间的一个 $1$-态射。那么图表
$$
\xymatrix{
\mathcal{I}_{\mathcal{S}/\mathcal{S}'}
\ar[d]_{F \circ (\ref{categories-equation-inertia-structure-map})}
\ar[rr]_{(\ref{categories-equation-comparison})} & &
\mathcal{I}_\mathcal{S} \ar[d]^{(\ref{categories-equation-functorial})} \\
\mathcal{S}' \ar[rr]^e & &
\mathcal{I}_{\mathcal{S}'}
}
$$
是一个 $2$-纤维积。
\end{lemma}

\begin{proof}
略。
\end{proof}


\section{群胚纤维化范畴}
\label{categories-section-fibred-groupoids}

\noindent
本节说明应如何理解群胚纤维化范畴，并看到它们基本上等同于取值于
群胚的 $(2, 1)$-范畴的函子。

\begin{definition}
\label{categories-definition-fibred-groupoids}
令 $p : \mathcal{S} \to \mathcal{C}$ 为函子。如果下列两个条件成立，
则称 $\mathcal{S}$ 在 $\mathcal{C}$ 上{\it 以群胚为纤维}：
\begin{enumerate}
\item 对 $\mathcal{C}$ 中的每个态射 $f : V \to U$ 以及
$U$ 的每个提升 $x$，都存在靶为 $x$ 的 $f$ 的提升
$\phi : y \to x$。
\item 对每一对态射 $\phi : y \to x$ 与 $ \psi : z \to x$，
以及任意满足 $p(\phi) \circ f = p(\psi)$ 的态射
$f : p(z) \to p(y)$，都存在唯一的 $f$ 的提升
$\chi : z \to y$，使 $\phi \circ \chi = \psi$。
\end{enumerate}
\end{definition}

\noindent
条件 (2) 换一种说法是：应用函子 $p$，会在下列交换图表中虚线箭头
所成的集合之间给出双射：
$$
\xymatrix{
y \ar[r] & x & p(y) \ar[r] & p(x) \\
z \ar@{-->}[u] \ar[ru] & & p(z) \ar@{-->}[u]\ar[ru] & \\
}
$$
还可以用下列方式理解第二个条件。设 $g : W \to V$ 和
$f : V \to U$ 是 $\mathcal{C}$ 中的态射。令
$x \in \Ob(\mathcal{S}_U)$。由第一个条件，可以将 $f$ 提升为
$ \phi : y \to x$，然后将 $g$ 提升为 $\psi : z \to y$。
也可以不作这个两步过程，而将 $g \circ f$ 直接提升为
$\gamma : z' \to x$。这给出图表
\begin{equation}
\label{categories-equation-fibred-groupoids}
\vcenter{
\xymatrix{
z' \ar@{-->}[d]\ar[rrd]^\gamma & & \\
z \ar@{-->}[u] \ar[r]^\psi \ar@{~>}[d]^p &
y \ar[r]^\phi \ar@{~>}[d]^p &
x \ar@{~>}[d]^p
\\
W \ar[r]^g & V \ar[r]^f & U \\
}
}
\end{equation}
中的实线箭头，其中波浪箭头表示的不是态射，而是函子 $p$。
把第二个条件应用于箭头 $\phi \circ \psi$、$\gamma$ 和
$\text{id}_W$，得到 $\mathcal{S}_W$ 中唯一的态射
$\chi : z \to z'$，使
$\gamma \circ \chi = \phi \circ \psi$。类似地，存在唯一态射
$z' \to z$。唯一性蕴含态射 $z' \to z$ 与 $z\to z'$ 互为逆，
换言之，它们是同构。

\medskip\noindent
从上述讨论应当已经清楚，群胚纤维化范畴与纤维化范畴关系非常
密切。下面给出相应结论。

\begin{lemma}
\label{categories-lemma-fibred-groupoids}
令 $p : \mathcal{S} \to \mathcal{C}$ 为函子。下列条件等价：
\begin{enumerate}
\item $p : \mathcal{S} \to \mathcal{C}$ 是群胚纤维化范畴；并且
\item 所有纤维范畴都是群胚，而且 $\mathcal{S}$ 是
$\mathcal{C}$ 上的纤维化范畴。
\end{enumerate}
进一步，在这种情形，$\mathcal{S}$ 的每个态射都是强笛卡尔的。
此外，如果对所有 $f: V \to U = p(x)$ 都给定一个位于 $f$
之上的 $f^\ast x \to x$，那么引理 \ref{categories-lemma-fibred} 中构造的数据
$(U \mapsto \mathcal{S}_U, f \mapsto f^*, \alpha_{f, g}, \alpha_U)$
定义了从 $\mathcal{C}^{opp}$ 到群胚的 $(2, 1)$-范畴的一个伪函子。
\end{lemma}

\begin{proof}
假设 $p : \mathcal{S} \to \mathcal{C}$ 是群胚纤维化范畴。
为了证明对 $U \in \Ob(\mathcal{C})$ 的所有纤维范畴
$\mathcal{S}_U$ 都是群胚，必须为 $\mathcal{S}_U$ 中每个
$f : y \to x$ 展示一个逆态射。左侧图表（在 $\mathcal{S}_U$ 中）
经 $p$ 映到右侧图表：
$$
\xymatrix{
y \ar[r]^f & x & U \ar[r]^{\text{id}_U} & U \\
x \ar@{-->}[u] \ar[ru]_{\text{id}_x} & &
U \ar@{-->}[u]\ar[ru]_{\text{id}_U} & \\
}
$$
因为只有 $\text{id}_U$ 能使右侧图表交换，所以存在唯一的
$g : x \to y$ 使左侧图表交换，故
$fg = \text{id}_x$。类似地，存在唯一的 $h : y \to x$，使
$gh = \text{id}_y$。于是 $fgh = f : y \to x$。由于
$fg = \text{id}_x$，得到 $h = f$。定义
\ref{categories-definition-fibred-groupoids} 的条件 (2) 恰好说明
$\mathcal{S}$ 的每个态射都是强笛卡尔的。因此，定义
\ref{categories-definition-fibred-groupoids} 的条件 (1) 蕴含
$\mathcal{S}$ 是 $\mathcal{C}$ 上的纤维化范畴。

\medskip\noindent
反过来，假设所有纤维范畴都是群胚，并且 $\mathcal{S}$ 是
$\mathcal{C}$ 上的纤维化范畴。需要验证定义
\ref{categories-definition-fibred-groupoids} 的条件 (1) 与 (2)。
第一个条件显然成立。令 $\phi : y \to x$、
$\psi : z \to x$ 和 $f : p(z) \to p(y)$ 满足
$p(\phi) \circ f = p(\psi)$，如定义
\ref{categories-definition-fibred-groupoids} 的条件 (2) 所述。记
$U = p(x)$、$V = p(y)$、$W = p(z)$、
$p(\phi) = g : V \to U$、$p(\psi) = h : W \to U$。
选择一个位于 $g$ 之上的强笛卡尔态射 $g^*x \to x$。
于是得到 $\mathcal{S}_V$ 中的态射 $i : y \to g^*x$，
因而它是同构。由于 $g^*x \to x$ 是强笛卡尔的，还得到对应于
有序对 $(\psi, f)$ 的态射 $j : z \to g^*x$。
容易验证 $\chi = i^{-1} \circ j$ 是一个解。

\medskip\noindent
我们已经在 (1) $\Rightarrow$ (2) 的证明中看到，
$\mathcal{S}$ 的每个态射都是强笛卡尔的。最后一个陈述直接由
引理 \ref{categories-lemma-fibred} 推出。
\end{proof}


\begin{lemma}
\label{categories-lemma-fibred-gives-fibred-groupoids}
令 $\mathcal{C}$ 为范畴。令
$p : \mathcal{S} \to \mathcal{C}$ 为纤维化范畴。如下定义
$\mathcal{S}$ 的子范畴 $\mathcal{S}'$：
\begin{enumerate}
\item $\Ob(\mathcal{S}') = \Ob(\mathcal{S})$；并且
\item 对 $x, y \in \Ob(\mathcal{S}')$，$x$ 与 $y$ 之间
$\mathcal{S}'$ 中的态射集合，是 $x$ 与 $y$ 之间
$\mathcal{S}$ 中的强笛卡尔态射集合。
\end{enumerate}
令 $p' : \mathcal{S}' \to \mathcal{C}$ 为 $p$ 在
$\mathcal{S}'$ 上的限制。那么
$p' : \mathcal{S}' \to \mathcal{C}$ 是群胚纤维化范畴。
\end{lemma}

\begin{proof}
注意，该构造确实有意义，因为由引理
\ref{categories-lemma-composition-cartesian}，$\mathcal{S}$ 任意对象的恒等
态射是强笛卡尔的，而且强笛卡尔态射的复合仍是强笛卡尔的。
定义 \ref{categories-definition-fibred-groupoids} 的第一个提升性质，来自
纤维化范畴的下述条件：给定任意态射 $f : V \to U$ 以及位于
$U$ 之上的 $x$，都存在位于 $f$ 之上的强笛卡尔态射
$\varphi : y \to x$。我们来验证 $\mathcal{C}$ 上的范畴
$p' : \mathcal{S}' \to \mathcal{C}$ 满足定义
\ref{categories-definition-fibred-groupoids} 的第二个提升性质。为此，仿照
定义 \ref{categories-definition-fibred-groupoids} 后的讨论论证。因此，在图表
\ref{categories-equation-fibred-groupoids} 中，态射 $\phi$、$\psi$ 和
$\gamma$ 是 $\mathcal{S}$ 的强笛卡尔态射。于是
$\gamma$ 与 $\phi \circ \psi$ 是 $\mathcal{S}$ 中位于
$\mathcal{C}$ 同一箭头之上的强笛卡尔态射，并且在
$\mathcal{S}$ 中有相同的靶。由定义
\ref{categories-definition-cartesian-over-C} 后的讨论，这意味着这两个箭头
按所需方式同构（这里还再次使用引理
\ref{categories-lemma-composition-cartesian}：$\mathcal{S}$ 中任意同构都是
强笛卡尔的）。
\end{proof}

\begin{example}
\label{categories-example-group-homomorphism-fibreedingroupoids}
群同态 $p : G \to H$ 产生例
\ref{categories-example-group-homomorphism-functor} 中那样的函子
$p : \mathcal{S}\to\mathcal{C}$。这个函子
$p : \mathcal{S}\to\mathcal{C}$ 是群胚纤维化范畴，当且仅当
$p$ 是满射。位于（唯一的）对象
$U\in \Ob(\mathcal{C})$ 之上的纤维范畴 $\mathcal{S}_U$，
是例 \ref{categories-example-group-groupoid} 中与 $p$ 的核关联的范畴。
\end{example}

\noindent
给定 $p : \mathcal{S} \to \mathcal{C}$，可以问：如果对所有
$U \in \Ob(\mathcal{C})$，纤维范畴 $\mathcal{S}_U$ 都是群胚，
那么 $\mathcal{S}$ 必然是 $\mathcal{C}$ 上的群胚纤维化范畴吗？
如下可见，答案是否定的。从一个群胚纤维化范畴
$p : \mathcal{S} \to \mathcal{C}$ 开始。改变 $\mathcal{S}$
中不映到某个对象的恒等态射的那些态射，并不会改变范畴
$\mathcal{S}_U$。所以可以破坏提升的存在性和唯一性条件。
一个例子是例 \ref{categories-example-group-homomorphism-fibreedingroupoids}
中的函子，其中 $G \to H$ 不是满射。下面是另一个例子。

\begin{example}
\label{categories-example-not-fibred-in-groupoids-but-fibre-cats-are}
令 $\Ob(\mathcal{C}) = \{A, B, T\}$，
$\Mor_\mathcal{C}(A, B) = \{f\}$、
$\Mor_\mathcal{C}(B, T) = \{g\}$、
$\Mor_\mathcal{C}(A, T) = \{h\} = \{gf\}, $，再加上每个对象的
恒等态射。下图画出了这个范畴。现在令
$\Ob(\mathcal{S}) = \{A', B', T'\}$，
$\Mor_\mathcal{S}(A', B') = \emptyset$、
$\Mor_\mathcal{S}(B', T') = \{g'\}$、
$\Mor_\mathcal{S}(A', T') = \{h'\}, $，再加上恒等态射。函子
$p : \mathcal{S} \to \mathcal{C}$ 是显然的。那么对每个
$U \in \Ob(\mathcal{C})$，$\mathcal{S}_U$ 都是只含一个对象及其
恒等态射的范畴，因而是群胚；但态射 $f: A \to B$ 无法提升。
类似地，若令
$\Mor_\mathcal{S}(A', B') = \{f'_1, f'_2\}$ 和
$ \Mor_\mathcal{S}(A', T') = \{h'\} = \{g'f'_1 \} = \{g'f'_2\}$，
那么纤维范畴不变，而下图中的 $f: A \to B$ 有两个提升。
$$
\xymatrix{
B' \ar[r]^{g'} & T' & & B \ar[r]^g & T & \\
A' \ar@{-->}[u]^{f'_1,\,f'_2} \ar[ru]_{h'} & & \ar@{}[u]^{\text{上图}} &
A \ar[u]^f \ar[ru]_{gf = h} & \\
}
$$
\end{example}

\noindent
稍后，我们想作出诸如“$\mathcal{C}$ 上的任意群胚纤维化范畴都等价
于一个分裂的范畴”，或“纤维范畴均为集合状的任意群胚纤维化范畴
都等价于一个以集合为纤维的范畴”这样的断言。等价的概念取决于
我们所使用的 $2$-范畴。

\begin{definition}
\label{categories-definition-categories-fibred-in-groupoids-over-C}
令 $\mathcal{C}$ 为范畴。所谓{\it $\mathcal{C}$ 上群胚纤维化范畴
所成的 $2$-范畴}，是 $\mathcal{C}$ 上纤维化范畴所成的
$2$-范畴（见定义 \ref{categories-definition-fibred-categories-over-C}）的
子 $2$-范畴，定义如下：
\begin{enumerate}
\item 它的对象是群胚纤维化范畴
$p : \mathcal{S} \to \mathcal{C}$。
\item 它的 $1$-态射
$(\mathcal{S}, p) \to (\mathcal{S}', p')$ 是满足
$p' \circ G = p$ 的函子 $G : \mathcal{S} \to \mathcal{S}'$
（由于每个态射都是强笛卡尔的，$G$ 自动保持它们）。
\item 对
$G, H : (\mathcal{S}, p) \to (\mathcal{S}', p')$，它的 $2$-态射
$t : G \to H$ 是满足下述条件的函子态射：对所有
$x \in \Ob(\mathcal{S})$，均有
$p'(t_x) = \text{id}_{p(x)}$。
\end{enumerate}
\end{definition}

\noindent
注意，每个 $2$-态射自动是同构！所以这实际上是
$(2, 1)$-范畴，而不仅是 $2$-范畴。下面给出关于 $2$-纤维积的
必备引理。

\begin{lemma}
\label{categories-lemma-2-product-fibred-categories}
令 $\mathcal{C}$ 为范畴。$\mathcal{C}$ 上群胚纤维化范畴所成的
$2$-范畴具有 2-纤维积，它们如引理
\ref{categories-lemma-2-product-categories-over-C} 所述。
\end{lemma}

\begin{proof}
由引理 \ref{categories-lemma-2-product-fibred-categories-over-C}，引理
\ref{categories-lemma-2-product-categories-over-C} 所描述的纤维积是
纤维化范畴。因此，只需证明纤维范畴都是群胚；见引理
\ref{categories-lemma-fibred-groupoids}。由引理
\ref{categories-lemma-fibre-2-fibre-product-categories-over-C}，只需证明
群胚的 $2$-纤维积是群胚；这显然成立（例如由引理
\ref{categories-lemma-2-fibre-product-categories} 中的构造可见）。
\end{proof}


\begin{remark}
\label{categories-remark-2-product-fibred-categories-same}
令 $\mathcal{C}$ 为范畴。令 $f : \mathcal{X} \to \mathcal{S}$
与 $g : \mathcal{Y} \to \mathcal{S}$ 为 $\mathcal{C}$ 上群胚纤维化
范畴的 $1$-态射。令 $p : \mathcal{S} \to \mathcal{C}$ 为给定函子。
我们断言，引理 \ref{categories-lemma-2-product-fibred-categories} 的
$2$-纤维积作为范畴典范等价于例
\ref{categories-example-2-fibre-product-categories} 中的纤维积。前者的对象是
满足 $p(\alpha) = \text{id}_U$ 的四元组
$(U, x, y, \alpha)$（见引理
\ref{categories-lemma-2-product-categories-over-C}），后者的对象是三元组
$(x, y, \alpha)$（见例
\ref{categories-example-2-fibre-product-categories}）。两个范畴之间的等价
由规则 $(U, x, y, \alpha) \mapsto (x, y, \alpha)$ 和
$(x, y, \alpha) \mapsto (p(f(x)), x, y', \alpha')$ 给出，其中
$\alpha' = g(\gamma)^{-1} \circ \alpha$，而
$\gamma : y' \to y$ 是箭头
$p(\alpha) : p(f(x)) \to p(g(y))$ 的一个提升。略去细节。
\end{remark}

\begin{lemma}
\label{categories-lemma-equivalence-fibred-categories}
令 $p : \mathcal{S}\to \mathcal{C}$ 与
$p' : \mathcal{S'}\to \mathcal{C}$ 为群胚纤维化范畴，并设
$G : \mathcal{S}\to \mathcal {S}'$ 是 $\mathcal{C}$ 上的函子。
\begin{enumerate}
\item $G$ 忠实（相应地，全忠实；相应地，是等价），当且仅当对每个
$U\in\Ob(\mathcal{C})$，诱导函子
$G_U : \mathcal{S}_U\to \mathcal{S}'_U$ 忠实
（相应地，全忠实；相应地，是等价）。
\item 若 $G$ 是等价，则 $G$ 是 $\mathcal{C}$ 上群胚纤维化范畴所成
$2$-范畴中的等价。
\end{enumerate}
\end{lemma}

\begin{proof}
令 $x, y$ 为 $\mathcal{S}$ 中位于同一对象 $U$ 之上的对象。
考虑交换图表
$$
\xymatrix{
\Mor_\mathcal{S}(x, y) \ar[rd]_p \ar[rr]_G & &
\Mor_{\mathcal{S}'}(G(x), G(y)) \ar[ld]^{p'} \\
& \Mor_\mathcal{C}(U, U) &
}
$$
从该图表可清楚看出，如果 $G$ 忠实（相应地，全忠实），那么每个
$G_U$ 也如此。

\medskip\noindent
假设 $G$ 是等价。对 $\mathcal{S}'$ 的每个对象 $x'$，都存在
$\mathcal{S}$ 的对象 $x$，使 $G(x)$ 同构于 $x'$。设 $x'$
位于 $U'$ 之上，而 $x$ 位于 $U$ 之上。那么 $\mathcal{C}$
中存在同构 $f : U' \to U$，即把 $p'$ 应用于同构
$x' \to G(x)$ 所得的同构。由群胚纤维化范畴的公理，
$\mathcal{S}$ 中存在位于 $f$ 之上的箭头 $f^*x \to x$。
因此，存在满足 $p'(\alpha) = \text{id}_{U'}$ 的同构
$\alpha : x' \to G(f^*x)$（这次使用 $\mathcal{S}'$ 的公理）。
综上，对 $\mathcal{S}'$ 的每个对象 $x'$，可以选择一个有序对
$(o_{x'}, \alpha_{x'})$，由 $\mathcal{S}$ 的对象 $o_{x'}$
以及同构 $\alpha_{x'} : x' \to G(o_{x'})$ 构成，并满足
$p'(\alpha_{x'}) = \text{id}_{p'(x')}$。从这里开始，按通常方式
进行（见引理 \ref{categories-lemma-equivalence-categories} 的证明），构造
逆函子 $F : \mathcal{S}' \to \mathcal{S}$：按
$x' \mapsto o_{x'}$，并将 $\varphi' : x' \to y'$ 映到唯一箭头
$\varphi_{\varphi'} : o_{x'} \to o_{y'}$，它满足
$\alpha_{y'}^{-1} \circ G(\varphi_{\varphi'}) \circ \alpha_{x'} = \varphi'$。
在这些选择下，$F$ 是 $\mathcal{C}$ 上的函子。略去对
$G \circ F$ 与 $F \circ G$ 分别和相应恒等函子 $2$-同构的验证
（在 $\mathcal{C}$ 上群胚纤维化范畴所成的 $2$-范畴中）。

\medskip\noindent
假设对所有 $U\in\Ob(\mathcal C)$，$G_U$ 都忠实
（相应地，全忠实）。为了证明 $G$ 忠实（相应地，全忠实），
需要证明，对任意对象 $x, y\in\Ob(\mathcal{S})$，$G$ 在
$\Mor_\mathcal{S}(x, y)$ 与
$\Mor_{\mathcal{S}'}(G(x), G(y))$ 之间诱导单射（相应地，双射）。
置 $U = p(x)$、$V = p(y)$。只需证明，对任意态射
$f : U \to V$，$G$ 在位于 $f$ 之上的态射 $x \to y$ 与位于
$f$ 之上的态射 $G(x) \to G(y)$ 之间诱导单射（相应地，双射）。
现在固定 $f : U \to V$，将一个拉回记为 $f^*y \to y$。
那么 $G(f^*y) \to G(y)$ 也是拉回。位于 $f$ 之上的从 $x$ 到
$y$ 的态射集合，与位于 $\text{id}_U$ 之上的从 $x$ 到
$f^*y$ 的态射集合双射。（由群胚纤维化范畴的第二个公理。）
类似地，位于 $f$ 之上的从 $G(x)$ 到 $G(y)$ 的态射集合，
与位于 $\text{id}_U$ 之上的从 $G(x)$ 到 $G(f^*y)$ 的态射集合
双射。因此，$G_U$ 忠实（相应地，全忠实）这一事实给出所需结论。

\medskip\noindent
最后，假设对所有 $U$，$G_U$ 都是等价，因而全忠实且本质满。
我们已经看到这蕴含 $G$ 全忠实，所以要证明它是等价，只需证明它
本质满。此事显然：若 $z'\in \Ob(\mathcal{S}')$，则
$z'\in \Ob(\mathcal{S}'_U)$，其中 $U = p'(z')$。由于 $G_U$
本质满，$z'$ 在 $\mathcal{S}'_U$ 中同构于形如 $G_U(z)$ 的对象，
其中 $z\in \Ob(\mathcal{S}_U)$。但 $\mathcal{S}'_U$ 中的态射
也是 $\mathcal{S}'$ 中的态射，所以 $z'$ 在 $\mathcal{S}'$ 中
同构于 $G(z)$。
\end{proof}


\begin{lemma}
\label{categories-lemma-fully-faithful-diagonal-equivalence}
令 $\mathcal{C}$ 为范畴。令 $p : \mathcal{S}\to \mathcal{C}$ 与
$p' : \mathcal{S'}\to \mathcal{C}$ 为群胚纤维化范畴。令
$G : \mathcal{S}\to \mathcal {S}'$ 为 $\mathcal{C}$ 上的函子。
那么 $G$ 全忠实，当且仅当对角函子
$$
\Delta_G :
\mathcal{S}
\longrightarrow
\mathcal{S} \times_{G, \mathcal{S}', G} \mathcal{S}
$$
是等价。
\end{lemma}

\begin{proof}
由引理 \ref{categories-lemma-equivalence-fibred-categories}，只需考察位于
$\mathcal{C}$ 的一个对象 $U$ 之上的纤维范畴。右端的一个对象
是三元组 $(x, x', \alpha)$，其中
$\alpha : G(x) \to G(x')$ 是 $\mathcal{S}'_U$ 中的态射。
函子 $\Delta_G$ 将 $\mathcal{S}_U$ 的对象 $x$ 映到三元组
$(x, x, \text{id}_{G(x)})$。注意，$(x, x', \alpha)$ 在
$\Delta_G$ 的本质像中，当且仅当对 $\mathcal{S}_U$ 中某个态射
$\beta : x \to x'$ 有 $\alpha = G(\beta)$（略去细节）。
因此，要使 $\Delta_G$ 成为等价，每个 $\alpha$ 都必须是某个态射
$\beta : x \to x'$ 的像，而且任意两个不同的态射
$\beta, \beta' : x \to x'$ 必须给出不同的态射
$G(\beta), G(\beta')$。这证明了引理。
\end{proof}

\begin{lemma}
\label{categories-lemma-morphisms-equivalent-fibred-groupoids}
令 $\mathcal{C}$ 为范畴。令 $\mathcal{S}_i$（$i = 1, 2, 3, 4$）
为 $\mathcal{C}$ 上的群胚纤维化范畴。设
$\varphi : \mathcal{S}_1 \to \mathcal{S}_2$ 与
$\psi : \mathcal{S}_3 \to \mathcal{S}_4$ 是
$\mathcal{C}$ 上的等价。那么
$$
\Mor_{\textit{Cat}/\mathcal{C}}(\mathcal{S}_2, \mathcal{S}_3)
\longrightarrow
\Mor_{\textit{Cat}/\mathcal{C}}(\mathcal{S}_1, \mathcal{S}_4),
\quad \alpha \longmapsto \psi \circ \alpha \circ \varphi
$$
是范畴等价。
\end{lemma}

\begin{proof}
这是一个一般事实，在任意 $2$-范畴中都成立。
\end{proof}

\begin{lemma}
\label{categories-lemma-inertia-fibred-groupoids}
令 $\mathcal{C}$ 为范畴。若
$p : \mathcal{S} \to \mathcal{C}$ 是群胚纤维化范畴，则惯性
纤维化范畴 $\mathcal{I}_\mathcal{S} \to \mathcal{C}$ 也是如此。
\end{lemma}

\begin{proof}
这从引理 \ref{categories-lemma-inertia-fibred-category} 中的构造即显然；
也可使用同一引理给出的等价
$I_\mathcal{S} \to \mathcal{S}
\times_{\Delta, \mathcal{S} \times_\mathcal{C} \mathcal{S}, \Delta}\mathcal{S}$，
再应用引理 \ref{categories-lemma-2-product-fibred-categories}。
\end{proof}

\begin{lemma}
\label{categories-lemma-cute-groupoids}
令 $\mathcal{C}$ 为范畴，令 $U \in \Ob(\mathcal{C})$。
若 $p : \mathcal{S} \to \mathcal{C}$ 是群胚纤维化范畴，并且
$p$ 经由 $p' : \mathcal{S} \to \mathcal{C}/U$ 分解，那么
$p' : \mathcal{S} \to \mathcal{C}/U$ 是群胚纤维化范畴。
\end{lemma}

\begin{proof}
我们已经在引理 \ref{categories-lemma-cute} 中看到，$p'$ 是纤维化范畴。
因此，只需证明纤维范畴都是群胚；见引理
\ref{categories-lemma-fibred-groupoids}。对 $V \in \Ob(\mathcal{C})$，有
$$
\mathcal{S}_V = \coprod\nolimits_{f : V \to U} \mathcal{S}_{(f : V \to U)}
$$
其中左端是 $p$ 的纤维范畴，右端是 $p'$ 的纤维范畴的不交并。
结论随即成立。
\end{proof}

\begin{lemma}
\label{categories-lemma-fibred-in-groupoids-over-fibred-in-groupoids}
令 $\mathcal{A} \to \mathcal{B} \to \mathcal{C}$ 为范畴之间的函子。
若 $\mathcal{A}$ 在 $\mathcal{B}$ 上以群胚为纤维，且
$\mathcal{B}$ 在 $\mathcal{C}$ 上以群胚为纤维，则
$\mathcal{A}$ 在 $\mathcal{C}$ 上以群胚为纤维。
\end{lemma}

\begin{proof}
可以直接由定义证明这一点。不过，我们将使用引理
\ref{categories-lemma-fibred-groupoids} 的判据。由引理
\ref{categories-lemma-fibred-over-fibred}，$\mathcal{A}$ 在
$\mathcal{C}$ 上纤维化。为了完成证明，只需说明对
$\mathcal{C}$ 中的 $U$，纤维范畴 $\mathcal{A}_U$ 是群胚。
事实上，若 $x \to y$ 是 $\mathcal{A}_U$ 的态射，则它在
$\mathcal{B}$ 中的像是同构，因为 $\mathcal{B}_U$ 是群胚。
但于是 $x \to y$ 是同构；例如，这可由引理
\ref{categories-lemma-composition-cartesian} 以及 $\mathcal{A}$ 的每个态射
都是强 $\mathcal{B}$-笛卡尔的这一事实推出（见引理
\ref{categories-lemma-fibred-groupoids}）。
\end{proof}

\begin{lemma}
\label{categories-lemma-fibred-groupoids-fibre-product-goes-up}
令 $p : \mathcal{S} \to \mathcal{C}$ 为群胚纤维化范畴。
令 $x \to y$ 与 $z \to y$ 为 $\mathcal{S}$ 的态射。若
$p(x) \times_{p(y)} p(z)$ 存在，则 $x \times_y z$ 存在，且
$p(x \times_y z) = p(x) \times_{p(y)} p(z)$。
\end{lemma}

\begin{proof}
由引理 \ref{categories-lemma-fibred-category-representable-goes-up} 推出。
\end{proof}


\begin{lemma}
\label{categories-lemma-ameliorate-morphism-categories-fibred-groupoids}
令 $\mathcal{C}$ 为范畴。令 $F : \mathcal{X} \to \mathcal{Y}$
为 $\mathcal{C}$ 上群胚纤维化范畴的一个 $1$-态射。存在由
$\mathcal{C}$ 上群胚纤维化范畴的 $1$-态射构成的分解
$\mathcal{X} \to \mathcal{X}' \to \mathcal{Y}$，使得
$\mathcal{X} \to \mathcal{X}'$ 是 $\mathcal{C}$ 上的等价，
并且 $\mathcal{X}'$ 是 $\mathcal{Y}$ 上的群胚纤维化范畴。
\end{lemma}

\begin{proof}
将结构函子记为 $p : \mathcal{X} \to \mathcal{C}$ 和
$q : \mathcal{Y} \to \mathcal{C}$。我们如下显式构造
$\mathcal{X}'$。$\mathcal{X}'$ 的一个对象是四元组
$(U, x, y, f)$，其中 $x \in \Ob(\mathcal{X}_U)$、
$y \in \Ob(\mathcal{Y}_U)$，而
$f : F(x) \to y$ 是 $\mathcal{Y}_U$ 中的同构。一个态射
$(a, b) : (U, x, y, f) \to (U', x', y', f')$ 由
$a : x \to x'$ 和 $b : y \to y'$ 给出，其中
$p(a) = q(b)$ 且 $f' \circ F(a) = b \circ f$。换言之，
$\mathcal{X}' = \mathcal{X} \times_{F, \mathcal{Y}, \text{id}} \mathcal{Y}$，
这里采用引理 \ref{categories-lemma-2-product-categories-over-C} 中的
$2$-纤维积构造。由引理 \ref{categories-lemma-2-product-fibred-categories}，
$\mathcal{X}'$ 是 $\mathcal{C}$ 上的群胚纤维化范畴，而且
$\mathcal{X}' \to \mathcal{Y}$ 是 $\mathcal{C}$ 上范畴的态射。
取函子 $\mathcal{X} \to \mathcal{X}'$，它在对象上由
$x \mapsto (p(x), x, F(x), \text{id}_{F(x)})$ 给出，在态射上由
$(a : x \to x') \mapsto (a, F(a))$ 给出。显然，复合
$\mathcal{X} \to \mathcal{X}' \to \mathcal{Y}$ 等于 $F$。
略去对 $\mathcal{X} \to \mathcal{X}'$ 是 $\mathcal{C}$ 上
纤维化范畴等价的验证。

\medskip\noindent
最后，需要证明 $\mathcal{X}' \to \mathcal{Y}$ 是群胚纤维化范畴。
令 $b : y' \to y$ 为 $\mathcal{Y}$ 中的态射，并令
$(U, x, y, f)$ 为 $\mathcal{X}'$ 中位于 $y$ 之上的对象。
因为 $\mathcal{X}$ 在 $\mathcal{C}$ 上以群胚为纤维，可以找到
位于 $U' = q(y') \to q(y) = U$ 之上的态射 $a : x' \to x$。
因为 $\mathcal{Y}$ 在 $\mathcal{C}$ 上以群胚为纤维，而且
$F(x') \to F(x)$ 与 $y' \to y$ 都位于同一个态射
$U' \to U$ 之上，可以找到位于 $\text{id}_{U'}$ 之上的
$f' : F(x') \to y'$，使 $f \circ F(a) = b \circ f'$。
由此得到
$(a, b) : (U', x', y', f') \to (U, x, y, f)$。
这验证了定义 \ref{categories-definition-fibred-groupoids} 的第一个条件 (1)。
为了验证 (2)，令
$(a, b) : (U', x', y', f') \to (U, x, y, f)$ 和
$(a', b') : (U'', x'', y'', f'') \to (U, x, y, f)$ 为
$\mathcal{X}'$ 的态射，并令 $b'' : y' \to y''$ 为
$\mathcal{Y}$ 中满足 $b' \circ b'' = b$ 的态射。需要证明存在唯一态射
$a'' : x' \to x''$，使
$f'' \circ F(a'') = b'' \circ f'$，并使
$(a', b') \circ (a'', b'') = (a, b)$。因为 $\mathcal{X}$
以群胚为纤维，存在唯一态射 $a'' : x' \to x''$，使
$a' \circ a'' = a$ 且 $p(a'') = q(b'')$。因为 $\mathcal{Y}$
以群胚为纤维，$F(a'')$ 是满足
$F(a') \circ F(a'') = F(a)$ 与 $q(F(a'')) = q(b'')$ 的唯一态射
$F(x') \to F(x'')$。关系
$f'' \circ F(a'') = b'' \circ f'$ 由此以及给定关系
$f \circ F(a) = b \circ f'$ 和
$f \circ F(a') = b' \circ f''$ 推出。
\end{proof}

\begin{lemma}
\label{categories-lemma-amelioration-unique}
令 $\mathcal{C}$ 为范畴。令 $F : \mathcal{X} \to \mathcal{Y}$
为 $\mathcal{C}$ 上群胚纤维化范畴的一个 $1$-态射。假设有
$2$-交换图表
$$
\xymatrix{
\mathcal{X}' \ar[rd]_f &
\mathcal{X} \ar[l]^a \ar[d]^F \ar[r]_b &
\mathcal{X}'' \ar[ld]^g \\
& \mathcal{Y}
}
$$
其中 $a$ 和 $b$ 是 $\mathcal{C}$ 上范畴的等价，而 $f$ 和 $g$
是群胚纤维化范畴。那么存在 $\mathcal{Y}$ 上范畴的等价
$h : \mathcal{X}'' \to \mathcal{X}'$，使 $h \circ b$ 与 $a$
作为 $\mathcal{C}$ 上范畴的 $1$-态射是 $2$-同构的。若上述图表
实际上交换，则可作出安排，使 $h \circ b$ 与 $a$ 作为
$\mathcal{Y}$ 上范畴的 $1$-态射是 $2$-同构的。
\end{lemma}

\begin{proof}
我们将证明，$\mathcal{Y}$ 上的 $\mathcal{X}'$ 和
$\mathcal{X}''$ 都等价于 $\mathcal{Y}$ 上的群胚纤维化范畴
$\mathcal{X} \times_{F, \mathcal{Y}, \text{id}} \mathcal{Y}$；
见引理
\ref{categories-lemma-ameliorate-morphism-categories-fibred-groupoids} 的证明。
在 $\mathcal{C}$ 上范畴所成的 $2$-范畴中，选择一个拟逆
$b^{-1} : \mathcal{X}'' \to \mathcal{X}$。由于图表的右三角形
$2$-交换，可见
$$
\xymatrix{
\mathcal{X} \ar[d]_F & \mathcal{X}'' \ar[l]^{b^{-1}} \ar[d]^g \\
\mathcal{Y} & \mathcal{Y} \ar[l]
}
$$
是 $2$-交换的。因此，由 $2$-纤维积的泛性质，得到一个 $1$-态射
$c : \mathcal{X}'' \to
\mathcal{X} \times_{F, \mathcal{Y}, \text{id}} \mathcal{Y}$。
此外，$c$ 是 $\mathcal{Y}$ 上范畴的态射（！），而且是等价
（由 $b$ 是等价这一假设；见引理
\ref{categories-lemma-equivalence-2-fibre-product}）。所以由引理
\ref{categories-lemma-equivalence-fibred-categories}，$c$ 是
$\mathcal{Y}$ 上群胚纤维化范畴所成的 $2$-范畴中的等价。

\medskip\noindent
还需构造 $c \circ b$ 与函子
$d : \mathcal{X} \to
\mathcal{X} \times_{F, \mathcal{Y}, \text{id}} \mathcal{Y}$ 之间的
$2$-同构；后者是在引理
\ref{categories-lemma-ameliorate-morphism-categories-fibred-groupoids} 的证明中
由 $x \mapsto (p(x), x, F(x), \text{id}_{F(x)})$ 构造的。
令 $\alpha : F \to g \circ b$ 与
$\beta : b^{-1} \circ b \to \text{id}$ 为 $\mathcal{C}$ 上范畴的
$1$-态射之间的 $2$-同构。注意，$c \circ b$ 在对象上由规则
$$
x \mapsto (p(x), b^{-1}(b(x)), g(b(x)), \alpha_x \circ F(\beta_x))
$$
给出。于是可见
$$
(\beta_x, \alpha_x) :
(p(x), x, F(x), \text{id}_{F(x)})
\longrightarrow
(p(x), b^{-1}(b(x)), g(b(x)), \alpha_x \circ F(\beta_x))
$$
是一个函子性同构，给出我们的 $2$-态射
$d \to b \circ c$。最后，如果图表交换，那么对所有 $x$，
$\alpha_x$ 都是恒等态射，从而这个 $2$-态射是在
$\mathcal{Y}$ 上范畴所成的 $2$-范畴中的 $2$-态射。
\end{proof}


\section{范畴预层}
\label{categories-section-presheaves-categories}

\noindent
本节将纤维化范畴的概念与关系密切的“范畴预层”概念作比较。
下面的例子说明基本构造。

\begin{example}
\label{categories-example-functor-categories}
令 $\mathcal{C}$ 为范畴。假设
$F : \mathcal{C}^{opp} \to \textit{Cat}$ 是到范畴的 $2$-范畴的函子；
见定义 \ref{categories-definition-functor-into-2-category}。对
$\mathcal{C}$ 中的 $f : V \to U$，我们启发性地将从 $F(U)$
到 $F(V)$ 的函子记为 $F(f) = f^\ast$。由此可以如下构造
$\mathcal{C}$ 上的纤维化范畴 $\mathcal{S}_F$。定义
$$
\Ob(\mathcal{S}_F) =
\{(U, x) \mid U\in \Ob(\mathcal{C}), x\in \Ob(F(U))\}.
$$
对 $(U, x), (V, y) \in \Ob(\mathcal{S}_F)$，定义
\begin{align*}
\Mor_{\mathcal{S}_F}((V, y), (U, x)) & =
\{ (f, \phi) \mid f \in \Mor_\mathcal{C}(V, U),
\phi \in \Mor_{F(V)}(y, f^\ast x)\} \\
& =
\coprod\nolimits_{f \in \Mor_\mathcal{C}(V, U)}
\Mor_{F(V)}(y, f^\ast x)
\end{align*}
为了定义复合，使用可复合的一对 $\mathcal{C}$ 中态射所满足的
$g^\ast \circ f^\ast = (f \circ g)^\ast$
（由取值于 $2$-范畴的函子的定义）。具体地，把
$\psi : z \to g^\ast y$ 与 $ \phi : y \to f^\ast x$ 的复合定义为
$ g^\ast(\phi) \circ \psi$。函子
$p_F : \mathcal{S}_F \to \mathcal{C}$ 由规则
$(U, x) \mapsto U$ 给出。我们来验证，这确实是纤维化范畴。
给定 $\mathcal{C}$ 中的 $f: V \to U$ 以及 $U$ 的提升
$(U, x)$，我们断言
$(f, \text{id}_{f^\ast x}): (V, {f^\ast x}) \to (U, x)$
是 $f$ 的强笛卡尔提升。需要证明左图中的 $h$ 确定右图中的
$(h, \nu)$：
$$
\xymatrix{
V \ar[r]^f &
U &
(V, f^*x) \ar[r]^{(f, \text{id}_{f^*x})} &
(U, x) \\
W \ar@{-->}[u]^h \ar[ru]_g & &
(W, z) \ar@{-->}[u]^{(h, \nu)} \ar[ru]_{(g, \psi)} &
}
$$
只需取 $\nu = \psi$；这是可行的，因为
$f \circ h = g$，从而 $g^*x = h^*f^*x$。此外，这是使右侧图表
交换的唯一提升。
\end{example}

\begin{definition}
\label{categories-definition-split-fibred-category}
令 $\mathcal{C}$ 为范畴。假设
$F : \mathcal{C}^{opp} \to \textit{Cat}$ 是到范畴的 $2$-范畴的
函子。把例 \ref{categories-example-functor-categories} 中构造的纤维化范畴
记为 $p_F : \mathcal{S}_F \to \mathcal{C}$。所谓
{\it 分裂纤维化范畴}，是一个在 $\mathcal{C}$ 上同构（！）于某个
这种范畴 {\it $\mathcal{S}_F$} 的纤维化范畴。
\end{definition}

\begin{lemma}
\label{categories-lemma-when-split}
令 $\mathcal{C}$ 为范畴。令 $\mathcal{S}$ 为
$\mathcal{C}$ 上的纤维化范畴。那么 $\mathcal{S}$ 是分裂的，
当且仅当对某个拉回选择（见定义
\ref{categories-definition-pullback-functor-fibred-category}），拉回函子
$(f \circ g)^*$ 与 $g^* \circ f^*$ 相等。
\end{lemma}

\begin{proof}
这直接由定义推出。
\end{proof}

\begin{lemma}
\label{categories-lemma-fibred-strict}
令 $ p : \mathcal{S} \to \mathcal{C}$ 为纤维化范畴。
存在一个反变函子 $F : \mathcal{C} \to \textit{Cat}$，使
$\mathcal{S}$ 与 $\mathcal{S}_F$ 在 $\mathcal{C}$ 上纤维化范畴所成
的 $2$-范畴中等价。换言之，每个纤维化范畴都等价于一个分裂的
纤维化范畴。
\end{lemma}

\begin{proof}
作一个拉回选择（见定义
\ref{categories-definition-pullback-functor-fibred-category}）。由引理
\ref{categories-lemma-fibred}，对 $\mathcal{C}$ 的每个态射 $f$，
得到拉回函子 $f^*$。

\medskip\noindent
如下构造一个新范畴 $\mathcal{S}'$。$\mathcal{S}'$ 的对象是有序对
$(x, f)$，它由 $\mathcal{C}$ 的态射 $f : V \to U$ 和位于 $U$
之上的 $\mathcal{S}$ 的对象 $x$ 组成；也就是说，
$x\in \Ob(\mathcal{S}_U)$。函子
$p' : \mathcal{S}' \to \mathcal{C}$ 将有序对 $(x, f)$ 映到态射
$f$ 的源；换言之，$p'(x, f : V\to U) = V$。一个态射
$\varphi : (x_1, f_1: V_1 \to U_1) \to
(x_2, f_2 : V_2 \to U_2)$
由有序对 $(\varphi, g)$ 给出，它由态射 $g : V_1 \to V_2$
以及满足 $p(\varphi) = g$ 的态射
$\varphi : f_1^\ast x_1 \to f_2^\ast x_2$ 组成。
定义复合法则毫无困难：对任意一对可复合态射，置
$(\varphi, g) \circ (\psi, h) =
(\varphi \circ \psi, g\circ h)$。有一个自然函子
$\mathcal{S} \to \mathcal{S}'$，它只是把位于 $U$ 之上的 $x$
映到有序对 $(x, \text{id}_U)$。

\medskip\noindent
此时需要验证 $p'$ 使 $\mathcal{S}'$ 成为 $\mathcal{C}$ 上的
纤维化范畴，并且需要验证 $\mathcal{S} \to \mathcal{S}'$
是 $\mathcal{C}$ 上范畴的一个等价，它将强笛卡尔态射映为
强笛卡尔态射。略去这些验证。

\medskip\noindent
最后，可以如下定义 $\mathcal{S}'$ 上的拉回函子：若
$g : V' \to V$ 且 $f : V \to U$，则在对象上置
$g^\ast(x, f) = (x, f \circ g)$。对于 $\mathcal{S}'_V$ 中态射之间
的态射
$(\varphi, \text{id}_V) : (x_1, f_1) \to (x_2, f_2)$，在态射上置
$g^\ast(\varphi, \text{id}_V) =
(g^\ast\varphi, \text{id}_{V'})$；这里使用引理
\ref{categories-lemma-fibred} 中的唯一等同
$g^\ast f_i^\ast x_i = (f_i \circ g)^\ast x_i$，将
$g^\ast\varphi$ 看作从 $(f_1 \circ g)^\ast x_1$ 到
$(f_2 \circ g)^\ast x_2$ 的态射。显然，这些拉回函子 $g^\ast$
满足 $g_1^\ast \circ g_2^\ast = (g_2\circ g_1)^\ast$；
换言之，$\mathcal{S}'$ 按所需方式分裂。
\end{proof}


\section{群胚预层}
\label{categories-section-presheaves-groupoids}

\noindent
本节将群胚纤维化范畴的概念与关系密切的“群胚预层”概念作比较。
下面的例子说明基本构造。

\begin{example}
\label{categories-example-functor-groupoids}
这是例 \ref{categories-example-functor-categories} 对“群胚预层”而非“范畴预层”
的类似物。输出将是群胚纤维化范畴，而非纤维化范畴。假设
$F : \mathcal{C}^{opp} \to \textit{Groupoids}$ 是到群胚范畴的函子；
见定义 \ref{categories-definition-functor-into-2-category}。对
$\mathcal{C}$ 中的 $f : V \to U$，我们启发性地将从 $F(U)$
到 $F(V)$ 的函子记为 $F(f) = f^\ast$。如下构造
$\mathcal{C}$ 上的群胚纤维化范畴 $\mathcal{S}_F$。定义
$$
\Ob(\mathcal{S}_F) =
\{(U, x) \mid U\in \Ob(\mathcal{C}), x\in \Ob(F(U))\}.
$$
对 $(U, x), (V, y) \in \Ob(\mathcal{S}_F)$，定义
\begin{align*}
\Mor_{\mathcal{S}_F}((V, y), (U, x))
& =
\{ (f, \phi) \mid f \in \Mor_\mathcal{C}(V, U),
\phi \in \Mor_{F(V)}(y, f^\ast x)\} \\
& =
\coprod\nolimits_{f \in \Mor_\mathcal{C}(V, U)}
\Mor_{F(V)}(y, f^\ast x)
\end{align*}
为了定义复合，使用可复合的一对 $\mathcal{C}$ 中态射所满足的
$g^\ast \circ f^\ast = (f \circ g)^\ast$
（由取值于 $2$-范畴的函子的定义）。具体地，把
$\psi : z \to g^\ast y$ 与 $ \phi : y \to f^\ast x$ 的复合定义为
$ g^\ast(\phi) \circ \psi$。函子
$p_F : \mathcal{S}_F \to \mathcal{C}$ 由规则
$(U, x) \mapsto U$ 给出。对每个 $U$，$F(U)$ 是群胚这一条件保证
$\mathcal{S}_F$ 在 $\mathcal{C}$ 上以群胚为纤维，因为我们已在
例 \ref{categories-example-functor-categories} 中看到 $\mathcal{S}_F$
是纤维化范畴；见引理 \ref{categories-lemma-fibred-groupoids}。
也可以如下直接证明定义 \ref{categories-definition-fibred-groupoids} 的条件
(1)、(2)：(1) 态射的提升存在，因为给定
$\mathcal{C}$ 中的 $f: V \to U$ 和 $\mathcal{S}_F$ 中位于 $U$
之上的对象 $(U, x)$，则
$(f, \text{id}_{f^\ast x}): (V, {f^\ast x}) \to (U, x)$
是 $f$ 的一个提升。(2) 假设给定下列实线图表：
$$
\xymatrix{
V \ar[r]^f & U & (V, y) \ar[r]^{(f, \phi)} & (U, x) \\
W \ar@{-->}[u]^h \ar[ru]_g & &
(W, z) \ar@{-->}[u]^{(h, \nu)} \ar[ru]_{(g, \psi)} & \\
}
$$
那么对虚线箭头有
$\nu = (h^\ast \phi)^{-1} \circ \psi$，所以给定 $h$ 后，
存在 $\nu$，并且由逆的唯一性，它是唯一的。
\end{example}

\begin{definition}
\label{categories-definition-split-category-fibred-in-groupoids}
令 $\mathcal{C}$ 为范畴。假设
$F : \mathcal{C}^{opp} \to \textit{Groupoids}$ 是到群胚的
$2$-范畴的函子。把例 \ref{categories-example-functor-groupoids} 中构造的
群胚纤维化范畴记为
$p_F : \mathcal{S}_F \to \mathcal{C}$。所谓
{\it 分裂群胚纤维化范畴}，是一个在 $\mathcal{C}$ 上同构（！）
于某个这种范畴 {\it $\mathcal{S}_F$} 的群胚纤维化范畴。
\end{definition}

\begin{lemma}
\label{categories-lemma-fibred-groupoids-strict}
令 $ p : \mathcal{S} \to \mathcal{C}$ 为群胚纤维化范畴。
存在一个反变函子 $F : \mathcal{C} \to \textit{Groupoids}$，
使 $\mathcal{S}$ 与 $\mathcal{S}_F$ 在 $\mathcal{C}$ 上等价。
换言之，每个群胚纤维化范畴都等价于一个分裂的群胚纤维化范畴。
\end{lemma}

\begin{proof}
作一个拉回选择（见定义
\ref{categories-definition-pullback-functor-fibred-category}）。由引理
\ref{categories-lemma-fibred} 和 \ref{categories-lemma-fibred-groupoids}，对
$\mathcal{C}$ 的每个态射 $f$，得到拉回函子 $f^*$。

\medskip\noindent
如下构造一个新范畴 $\mathcal{S}'$。$\mathcal{S}'$ 的对象是有序对
$(x, f)$，它由 $\mathcal{C}$ 的态射 $f : V \to U$ 和位于
$U$ 之上的 $\mathcal{S}$ 的对象 $x$ 组成；也就是说，
$x\in \Ob(\mathcal{S}_U)$。函子
$p' : \mathcal{S}' \to \mathcal{C}$ 将有序对 $(x, f)$ 映到态射
$f$ 的源；换言之，$p'(x, f : V\to U) = V$。一个态射
$\varphi : (x_1, f_1: V_1 \to U_1) \to
(x_2, f_2 : V_2 \to U_2)$
由有序对 $(\varphi, g)$ 给出，它由态射 $g : V_1 \to V_2$
以及满足 $p(\varphi) = g$ 的态射
$\varphi : f_1^\ast x_1 \to f_2^\ast x_2$ 组成。
定义复合法则毫无困难：对任意一对可复合态射，置
$(\varphi, g) \circ (\psi, h) =
(\varphi \circ \psi, g\circ h)$。有一个自然函子
$\mathcal{S} \to \mathcal{S}'$，它只是把位于 $U$ 之上的 $x$
映到有序对 $(x, \text{id}_U)$。

\medskip\noindent
此时需要验证 $p'$ 使 $\mathcal{S}'$ 成为 $\mathcal{C}$ 上的
群胚纤维化范畴，并且需要验证 $\mathcal{S} \to \mathcal{S}'$
是 $\mathcal{C}$ 上范畴的等价。略去这些验证。

\medskip\noindent
最后，可以如下定义 $\mathcal{S}'$ 上的拉回函子：若
$g : V' \to V$ 且 $f : V \to U$，则在对象上置
$g^\ast(x, f) = (x, f \circ g)$。对于 $\mathcal{S}'_V$ 中态射之间
的态射
$(\varphi, \text{id}_V) : (x_1, f_1) \to (x_2, f_2)$，在态射上置
$g^\ast(\varphi, \text{id}_V) =
(g^\ast\varphi, \text{id}_{V'})$；这里使用引理
\ref{categories-lemma-fibred-groupoids} 中的唯一等同
$g^\ast f_i^\ast x_i = (f_i \circ g)^\ast x_i$，将
$g^\ast\varphi$ 看作从 $(f_1 \circ g)^\ast x_1$ 到
$(f_2 \circ g)^\ast x_2$ 的态射。显然，这些拉回函子 $g^\ast$
满足 $g_1^\ast \circ g_2^\ast = (g_2\circ g_1)^\ast$；
换言之，$\mathcal{S}'$ 按所需方式分裂。
\end{proof}

\noindent
我们将在第 \ref{categories-section-representable-1-morphisms} 节看到这个引理的
另一种证明。


\section{集合纤维化范畴}
\label{categories-section-fibred-in-sets}

\begin{definition}
\label{categories-definition-discrete}
如果一个范畴仅有恒等态射，则称其为{\it 离散的}。
\end{definition}

\noindent
离散范畴只有一项有意义的信息：它的对象集合。因此，我们有时
混同离散范畴与集合。

\begin{definition}
\label{categories-definition-category-fibred-sets}
令 $\mathcal{C}$ 为范畴。所谓{\it 集合纤维化范畴}，或者
{\it 离散范畴纤维化范畴}，是所有纤维范畴均为离散范畴的
群胚纤维化范畴。
\end{definition}

\noindent
我们想阐明集合纤维化范畴与预层（见定义
\ref{categories-definition-presheaf}）之间的关系。为此，先作如下定义是
合理的。

\begin{definition}
\label{categories-definition-categories-fibred-in-sets-over-C}
令 $\mathcal{C}$ 为范畴。所谓{\it $\mathcal{C}$ 上集合纤维化范畴
所成的 $2$-范畴}，是 $\mathcal{C}$ 上群胚纤维化范畴所成范畴
（见定义 \ref{categories-definition-categories-fibred-in-groupoids-over-C}）
的子 $2$-范畴，定义如下：
\begin{enumerate}
\item 它的对象是集合纤维化范畴
$p : \mathcal{S} \to \mathcal{C}$。
\item 它的 $1$-态射
$(\mathcal{S}, p) \to (\mathcal{S}', p')$ 是满足
$p' \circ G = p$ 的函子 $G : \mathcal{S} \to \mathcal{S}'$
（由于每个态射都是强笛卡尔的，$G$ 自动保持它们）。
\item 对
$G, H : (\mathcal{S}, p) \to (\mathcal{S}', p')$，它的 $2$-态射
$t : G \to H$ 是满足下述条件的函子态射：对所有
$x \in \Ob(\mathcal{S})$，均有
$p'(t_x) = \text{id}_{p(x)}$。
\end{enumerate}
\end{definition}

\noindent
注意，每个 $2$-态射自动是同构。因此，这个 $2$-范畴实际上是
$(2, 1)$-范畴。下面给出关于 $2$-纤维积存在性的必备引理。

\begin{lemma}
\label{categories-lemma-2-product-categories-fibred-sets}
令 $\mathcal{C}$ 为范畴。$\mathcal{C}$ 上集合纤维化范畴所成的
2-范畴具有 2-纤维积。更精确地说，如果从集合纤维化范畴出发，
则引理 \ref{categories-lemma-2-product-categories-over-C} 中描述的
2-纤维积仍给出集合纤维化范畴。
\end{lemma}

\begin{proof}
略。
\end{proof}

\begin{example}
\label{categories-example-presheaf}
这是例 \ref{categories-example-functor-categories} 和
\ref{categories-example-functor-groupoids} 对预层而非“范畴预层”的类似物。
输出将是集合纤维化范畴，而非纤维化范畴。假设
$F : \mathcal{C}^{opp} \to \textit{Sets}$ 是预层。对
$\mathcal{C}$ 中的 $f : V \to U$，我们启发性地记
$F(f) = f^\ast : F(U) \to F(V)$。如下构造 $\mathcal{C}$ 上的
集合纤维化范畴 $\mathcal{S}_F$。定义
$$
\Ob(\mathcal{S}_F) =
\{(U, x) \mid U \in \Ob(\mathcal{C}), x \in \Ob(F(U))\}.
$$
对 $(U, x), (V, y) \in \Ob(\mathcal{S}_F)$，定义
\begin{align*}
\Mor_{\mathcal{S}_F}((V, y), (U, x))
& =
\{f \in \Mor_\mathcal{C}(V, U) \mid f^*x = y\}
\end{align*}
复合继承自 $\mathcal{C}$ 中的复合；这是可行的，因为对
$\mathcal{C}$ 中可复合的一对态射，有
$g^\ast \circ f^\ast = (f \circ g)^\ast$。函子
$p_F : \mathcal{S}_F \to \mathcal{C}$ 由规则
$(U, x) \mapsto U$ 给出。由于每个纤维范畴
$\mathcal{S}_{F, U}$ 都是以 $F(U)$ 为底层集合的离散范畴，
并且我们已在例 \ref{categories-example-functor-groupoids} 中看到
$\mathcal{S}_F$ 是群胚纤维化范畴，所以可知
$\mathcal{S}_F$ 以集合为纤维。
\end{example}

\begin{lemma}
\label{categories-lemma-2-category-fibred-sets}
\begin{slogan}
集合纤维化范畴恰好就是预层。
\end{slogan}
令 $\mathcal{C}$ 为范畴。集合纤维化范畴之间仅有的
$2$-态射是恒等态射。换言之，集合纤维化范畴所成的 $2$-范畴是
一个范畴。此外，存在范畴等价
$$
\left\{
\begin{matrix}
\text{预层范畴}\\
\text{集合范畴，位于 }\mathcal{C}
\end{matrix}
\right\}
\leftrightarrow
\left\{
\begin{matrix}
\text{范畴的范畴}\\
\text{在下列对象之上以集合为纤维： }\mathcal{C}
\end{matrix}
\right\}
$$
从左到右的函子是例 \ref{categories-example-presheaf} 中讨论的构造
$F \to \mathcal{S}_F$。从右到左的函子，把
$p : \mathcal{S} \to \mathcal{C}$ 映到对象预层
$U \mapsto \Ob(\mathcal{S}_U)$。
\end{lemma}

\begin{proof}
第一个断言显然，因为纤维范畴中仅有的态射是恒等态射。

\medskip\noindent
假设 $p :
\mathcal{S} \to \mathcal{C}$ 以集合为纤维。令
$f : V \to U$ 为 $\mathcal{C}$ 中的态射，令
$x \in \Ob(\mathcal{S}_U)$。那么对象 $f^\ast x$ 恰有一种选择。
因此，对引理 \ref{categories-lemma-fibred-groupoids} 中的 $f, g$，有
$(f \circ g)^\ast x = g^\ast(f^\ast x)$。所以可以把赋值
$U \mapsto \Ob(\mathcal{S}_U)$ 和 $f \mapsto f^\ast$
看作 $\mathcal{C}$ 上的预层。
\end{proof}

\noindent
下面给出集合纤维化范畴的一个重要例子。

\begin{example}
\label{categories-example-fibred-category-from-functor-of-points}
令 $\mathcal{C}$ 为范畴，令 $X \in \Ob(\mathcal{C})$。
考虑可表预层 $h_X = \Mor_\mathcal{C}(-, X)$
（见例 \ref{categories-example-hom-functor}）。另一方面，考虑例
\ref{categories-example-category-over-X} 中的范畴
$p : \mathcal{C}/X \to \mathcal{C}$。纤维范畴
$(\mathcal{C}/X)_U$ 的对象是态射 $h : U \to X$，而且只有
恒等态射。因此，在引理 \ref{categories-lemma-2-category-fibred-sets} 的
对应下，有
$$
h_X \longleftrightarrow \mathcal{C}/X.
$$
换言之，范畴 $\mathcal{C}/X$ 典范等价于例
\ref{categories-example-presheaf} 中与 $h_X$ 关联的范畴
$\mathcal{S}_{h_X}$。
\end{example}

\noindent
因此，很容易想把群胚纤维化范畴的 2-范畴中的“可表”对象定义为：
其关联预层可表的集合纤维化范畴。然而，这并不是适用的好定义，
因为我们更希望有一个在等价下不变的概念。为使这一点精确化，
我们将研究究竟哪些群胚纤维化范畴等价于集合纤维化范畴。


\section{集合胚纤维化范畴}
\label{categories-section-fibred-in-setoids}

\begin{definition}
\label{categories-definition-setoid}
如果一个范畴是群胚，且每个对象恰有一个自同构——恒等态射，
则称其为{\it 集合胚（setoid）}\footnote{吃了类固醇的集合！？}。
\end{definition}

\noindent
若 $C$ 是带等价关系 $\sim$ 的集合，则可以如下构造集合胚
$\mathcal{C}$：$\Ob(\mathcal{C}) = C$；除非 $x \sim y$，
否则 $\Mor_\mathcal{C}(x, y) = \emptyset$；而在该关系成立时，
置 $\Mor_\mathcal{C}(x, y) = \{1\}$。
$\sim$ 的传递性意味着可以复合态射。反过来，任意集合胚范畴在其
对象上定义一个等价关系（同构关系），而通过刚才描述的过程，
可以从这个关系恢复原范畴（在唯一同构意义下，而不只是在等价
意义下）。

\medskip\noindent
离散范畴是集合胚。对任意集合胚 $\mathcal{C}$，有一个典范过程
产生与之等价的离散范畴，即用同构类集合替代
$\Ob(\mathcal{C})$（并添加恒等态射）。用带等价关系的集合来说，
这对应于对等价关系取商。

\begin{definition}
\label{categories-definition-category-fibred-setoids}
令 $\mathcal{C}$ 为范畴。所谓{\it 集合胚纤维化范畴}，是所有纤维
范畴均为集合胚的群胚纤维化范畴。
\end{definition}

\noindent
下面将阐明集合胚纤维化范畴与集合纤维化范畴之间的关系。

\begin{definition}
\label{categories-definition-categories-fibred-in-setoids-over-C}
令 $\mathcal{C}$ 为范畴。所谓{\it $\mathcal{C}$ 上集合胚纤维化范畴
所成的 $2$-范畴}，是 $\mathcal{C}$ 上群胚纤维化范畴所成范畴
（见定义 \ref{categories-definition-categories-fibred-in-groupoids-over-C}）
的子 $2$-范畴，定义如下：
\begin{enumerate}
\item 它的对象是集合胚纤维化范畴
$p : \mathcal{S} \to \mathcal{C}$。
\item 它的 $1$-态射
$(\mathcal{S}, p) \to (\mathcal{S}', p')$ 是满足
$p' \circ G = p$ 的函子 $G : \mathcal{S} \to \mathcal{S}'$
（由于每个态射都是强笛卡尔的，$G$ 自动保持它们）。
\item 对
$G, H : (\mathcal{S}, p) \to (\mathcal{S}', p')$，它的 $2$-态射
$t : G \to H$ 是满足下述条件的函子态射：对所有
$x \in \Ob(\mathcal{S})$，均有
$p'(t_x) = \text{id}_{p(x)}$。
\end{enumerate}
\end{definition}

\noindent
注意，每个 $2$-态射自动是同构。因此，这个 $2$-范畴实际上是
$(2, 1)$-范畴。

\noindent
下面给出关于 $2$-纤维积存在性的必备引理。

\begin{lemma}
\label{categories-lemma-2-product-categories-fibred-setoids}
令 $\mathcal{C}$ 为范畴。$\mathcal{C}$ 上集合胚纤维化范畴所成的
2-范畴具有 2-纤维积。更精确地说，如果从集合胚纤维化范畴出发，
则引理 \ref{categories-lemma-2-product-categories-over-C} 中描述的
2-纤维积仍给出集合胚纤维化范畴。
\end{lemma}

\begin{proof}
略。
\end{proof}

\begin{lemma}
\label{categories-lemma-setoid-fibres}
令 $\mathcal{C}$ 为范畴。令 $\mathcal{S}$ 为
$\mathcal{C}$ 上的范畴。
\begin{enumerate}
\item 若 $\mathcal{S} \to \mathcal{S}'$ 是 $\mathcal{C}$ 上的等价，
而 $\mathcal{S}'$ 是 $\mathcal{C}$ 上的集合纤维化范畴，则
\begin{enumerate}
\item $\mathcal{S}$ 是 $\mathcal{C}$ 上的集合胚纤维化范畴；并且
\item 对每个 $U \in \Ob(\mathcal{C})$，映射
$\Ob(\mathcal{S}_U) \to \Ob(\mathcal{S}'_U)$ 将靶标识为源的
同构类集合。
\end{enumerate}
\item 若 $p : \mathcal{S} \to \mathcal{C}$ 是集合胚纤维化范畴，
则存在集合纤维化范畴
$p' : \mathcal{S}' \to \mathcal{C}$，以及 $\mathcal{C}$ 上的等价
$\text{can} : \mathcal{S} \to \mathcal{S}'$。
\end{enumerate}
\end{lemma}

\begin{proof}
我们证明 (2)。范畴 $\mathcal{S}'$ 的一个对象是有序对
$(U, \xi)$，其中 $U \in \Ob(\mathcal{C})$，而 $\xi$ 是
$\mathcal{S}_U$ 的对象的一个同构类。一个态射
$(U, \xi) \to (V , \psi)$ 由态射 $x \to y$ 给出，其中
$x \in \xi$ 且 $y \in \psi$。这里，如果两个态射
$x \to y$ 与 $x' \to y'$ 诱导相同的态射 $U \to V$，并且对
$\mathcal{S}_U$ 中同构 $x \to x'$ 和 $\mathcal{S}_V$ 中同构
$y \to y'$ 的某些选择，复合 $x \to x' \to y'$ 与
$x \to y \to y'$ 相同，就把这两个态射等同。由构造，
$\mathcal{S} \to \mathcal{S}'$ 在对象和态射上都有满映射。
把 $\mathcal{S}'$ 中的态射复合定义为使
$\mathcal{S} \to \mathcal{S}'$ 成为函子的唯一法则。
略去若干细节。
\end{proof}


\noindent
所以，集合胚纤维化范畴恰好就是等价于集合纤维化范畴的群胚
纤维化范畴。此外，由引理 \ref{categories-lemma-2-category-fibred-sets}，
集合纤维化范畴之间的等价是同构。

\begin{lemma}
\label{categories-lemma-2-category-fibred-setoids}
令 $\mathcal{C}$ 为范畴。引理 \ref{categories-lemma-setoid-fibres} 第 (2) 部分
的构造给出函子
$$
F :
\left\{
\begin{matrix}
\text{范畴的 2-范畴}\\
\text{在下列对象之上以集合胚为纤维： }\mathcal{C}
\end{matrix}
\right\}
\longrightarrow
\left\{
\begin{matrix}
\text{范畴的范畴}\\
\text{在下列对象之上以集合为纤维： }\mathcal{C}
\end{matrix}
\right\}
$$
（见定义 \ref{categories-definition-functor-into-2-category}）。这个函子在下述
意义下是等价：
\begin{enumerate}
\item 对满足 $F(f) = F(g)$ 的任意两个 1-态射
$f, g : \mathcal{S}_1 \to \mathcal{S}_2$，存在唯一的 2-同构
$f \to g$；
\item 对任意态射 $h : F(\mathcal{S}_1) \to F(\mathcal{S}_2)$，
存在满足 $F(f) = h$ 的 1-态射
$f : \mathcal{S}_1 \to \mathcal{S}_2$；并且
\item 任意集合纤维化范畴 $\mathcal{S}$ 都等于
$F(\mathcal{S})$。
\end{enumerate}
特别地，按规则
$F_i(U) = \Ob(\mathcal{S}_{i, U})/\cong$
定义 $F_i \in \textit{PSh}(\mathcal{C})$，则有
$$
\Mor_{\textit{Cat}/\mathcal{C}}(\mathcal{S}_1, \mathcal{S}_2)
\Big/
2\text{-同构}
=
\Mor_{\textit{PSh}(\mathcal{C})}(F_1, F_2)
$$
更精确地说，给定任意映射 $\phi : F_1 \to F_2$，存在一个
$1$-态射 $f : \mathcal{S}_1 \to \mathcal{S}_2$，它在对象的
同构类上诱导 $\phi$，并且在唯一 $2$-同构意义下唯一。
\end{lemma}

\begin{proof}
由引理 \ref{categories-lemma-2-category-fibred-sets}，$F$ 的靶是一个范畴，
所以这个断言有意义。引理 \ref{categories-lemma-setoid-fibres} 第 (2) 部分
的构造，把 $\mathcal{S}$ 映到一个集合纤维化范畴，其在 $U$
上方的值是 $\mathcal{S}_U$ 中同构类的集合。因此，显然它定义了
所述函子。令 $f, g : \mathcal{S}_1 \to \mathcal{S}_2$ 满足
$F(f) = F(g)$，如 (1) 所述。对 $\mathcal{C}$ 的每个对象 $U$
以及 $\mathcal{S}_{1, U}$ 的每个对象 $x$，由假设有
$f(x) \cong g(x)$。由于 $\mathcal{S}_2$ 以集合胚为纤维，
在 $\mathcal{S}_{2, U}$ 中存在唯一同构
$t_x : f(x) \to g(x)$。显然，规则 $x \mapsto t_x$ 给出所需的
$2$-同构 $f \to g$。略去 (2) 与 (3) 的证明。为了得到最后一个
断言，使用引理 \ref{categories-lemma-2-category-fibred-sets} 看出右端等于
$\Mor_{\textit{Cat}/\mathcal{C}}(F(\mathcal{S}_1), F(\mathcal{S}_2))$，
再应用上述 (1) 与 (2)。
\end{proof}

\noindent
下面给出集合胚纤维化范畴在所有群胚纤维化范畴中的另一种刻画。

\begin{lemma}
\label{categories-lemma-characterize-fibred-setoids-inertia}
令 $\mathcal{C}$ 为范畴。令
$p : \mathcal{S} \to \mathcal{C}$ 为群胚纤维化范畴。下列条件等价：
\begin{enumerate}
\item $p : \mathcal{S} \to \mathcal{C}$ 是集合胚纤维化范畴；并且
\item 典范 $1$-态射 $\mathcal{I}_\mathcal{S} \to \mathcal{S}$
（见 (\ref{categories-equation-inertia-structure-map})）是等价
（$\mathcal{C}$ 上范畴的等价）。
\end{enumerate}
\end{lemma}

\begin{proof}
假设 (2)。范畴 $\mathcal{I}_\mathcal{S}$ 的对象是
$(x, \alpha)$，其中 $x \in \mathcal{S}$，比如
$p(x) = U$，而 $\alpha : x \to x$ 是 $\mathcal{S}_U$ 中的态射。
因此，若
$\mathcal{I}_\mathcal{S} \to \mathcal{S}$ 是 $\mathcal{C}$
上的等价，则每一对对象 $(x, \alpha)$、$(x, \alpha')$ 都在
$\mathcal{I}_\mathcal{S}$ 位于 $U$ 上方的纤维范畴中同构。
考察 $\mathcal{I}_\mathcal{S}$ 中态射的定义，可知
$\alpha$、$\alpha'$ 在 $x$ 的自同构群中共轭。因此，取
$\alpha' = \text{id}_x$，便知 $x$ 的每个自同构都等于恒等态射。
由于 $\mathcal{S} \to \mathcal{C}$ 以群胚为纤维，这蕴含
$\mathcal{S} \to \mathcal{C}$ 以集合胚为纤维。略去
(1) $\Rightarrow$ (2) 的证明。
\end{proof}

\begin{lemma}
\label{categories-lemma-category-fibred-setoids-presheaves-products}
令 $\mathcal{C}$ 为范畴。引理
\ref{categories-lemma-2-category-fibred-setoids} 中将预层关联于集合胚
纤维化范畴的构造与积相容；其含义是，与 $2$-纤维积
$\mathcal{X} \times_\mathcal{Y} \mathcal{Z}$ 关联的预层，
是分别与 $\mathcal{X}, \mathcal{Y}, \mathcal{Z}$ 关联的预层的
纤维积。
\end{lemma}

\begin{proof}
令 $U \in \Ob(\mathcal{C})$。引理所说的恰好是
$$
\Ob((\mathcal{X} \times_\mathcal{Y} \mathcal{Z})_U)/\!\cong
\quad \text{等于} \quad
\Ob(\mathcal{X}_U)/\!\cong
\ \times_{\Ob(\mathcal{Y}_U)/\!\cong}
\ \Ob(\mathcal{Z}_U)/\!\cong
$$
略去其证明。（但注意，如果范畴 $\mathcal{Y}_U$ 不是集合胚，
这在一般情形并不成立。）
\end{proof}


\section{可表群胚纤维化范畴}
\label{categories-section-representable-fibred-groupoids}

\noindent
下面给出可表群胚纤维化范畴的定义。正如先前承诺的，这个定义在
等价下不变。

\begin{definition}
\label{categories-definition-representable-fibred-category}
令 $\mathcal{C}$ 为范畴。如果存在 $\mathcal{C}$ 的对象 $X$
以及一个等价 $j : \mathcal{S} \to \mathcal{C}/X$
（在 $\mathcal{C}$ 上群胚纤维化范畴所成的 $2$-范畴中），
则称群胚纤维化范畴
$p : \mathcal{S} \to \mathcal{C}$ 是{\it 可表的}。
\end{definition}

\noindent
通常会滥用记号说{\it $X$ 表示 $\mathcal{S}$}，而不提等价 $j$。
下面详述这意味着什么。

\begin{lemma}
\label{categories-lemma-characterize-representable-fibred-category}
令 $\mathcal{C}$ 为范畴。令
$p : \mathcal{S} \to \mathcal{C}$ 为群胚纤维化范畴。
\begin{enumerate}
\item $\mathcal{S}$ 可表，当且仅当下列条件成立：
\begin{enumerate}
\item $\mathcal{S}$ 以集合胚为纤维；并且
\item 预层 $U \mapsto \Ob(\mathcal{S}_U)/\cong$ 可表。
\end{enumerate}
\item 若 $\mathcal{S}$ 可表，则有序对 $(X, j)$ 在同构意义下唯一
确定，其中 $j$ 是等价 $j : \mathcal{S} \to \mathcal{C}/X$。
\end{enumerate}
\end{lemma}

\begin{proof}
第一个断言直接由引理 \ref{categories-lemma-setoid-fibres} 推出。
对第二个断言，设 $j' : \mathcal{S} \to \mathcal{C}/X'$
是第二个这样的有序对。选择一个 $1$-态射
$t' : \mathcal{C}/X' \to \mathcal{S}$，使
$j' \circ t' \cong \text{id}_{\mathcal{C}/X'}$ 且
$t' \circ j' \cong \text{id}_\mathcal{S}$。那么
$j \circ t' : \mathcal{C}/X' \to \mathcal{C}/X$ 是等价。
由引理 \ref{categories-lemma-2-category-fibred-sets}，它因而是同构。
所以由米田引理 \ref{categories-lemma-yoneda}（例如经由例
\ref{categories-example-fibred-category-from-functor-of-points}），
它由一个同构 $X' \to X$ 给出。
\end{proof}

\begin{lemma}
\label{categories-lemma-morphisms-representable-fibred-categories}
令 $\mathcal{C}$ 为范畴。令 $\mathcal{X}$、$\mathcal{Y}$
为 $\mathcal{C}$ 上的群胚纤维化范畴。假设
$\mathcal{X}$、$\mathcal{Y}$ 分别由 $\mathcal{C}$ 的对象
$X$、$Y$ 表示。那么
$$
\Mor_{\textit{Cat}/\mathcal{C}}(\mathcal{X}, \mathcal{Y})
\Big/
2\text{-同构}
=
\Mor_\mathcal{C}(X, Y)
$$
更精确地说，给定 $\phi : X \to Y$，存在一个 $1$-态射
$f : \mathcal{X} \to \mathcal{Y}$，它在对象的同构类上诱导
$\phi$，并且在唯一 $2$-同构意义下唯一。
\end{lemma}

\begin{proof}
由例 \ref{categories-example-fibred-category-from-functor-of-points}，有
$\mathcal{C}/X = \mathcal{S}_{h_X}$ 和
$\mathcal{C}/Y = \mathcal{S}_{h_Y}$。由引理
\ref{categories-lemma-2-category-fibred-setoids}，有
$$
\Mor_{\textit{Cat}/\mathcal{C}}(\mathcal{X}, \mathcal{Y})
\Big/
2\text{-同构}
=
\Mor_{\textit{PSh}(\mathcal{C})}(h_X, h_Y)
$$
由米田引理 \ref{categories-lemma-yoneda}，有
$\Mor_{\textit{PSh}(\mathcal{C})}(h_X, h_Y)
= \Mor_\mathcal{C}(X, Y)$。
\end{proof}


\section{2-米田引理}
\label{categories-section-2-yoneda}

\noindent
令 $\mathcal{C}$ 为范畴。$\mathcal{C}$ 上纤维范畴所成的 $2$-范畴
已在定义 \ref{categories-definition-fibred-categories-over-C} 中构造并定义。
若 $\mathcal{S}$、$\mathcal{S}'$ 是 $\mathcal{C}$ 上的纤维范畴，则
$$
\Mor_{\textit{Fib}/\mathcal{C}}(\mathcal{S}, \mathcal{S}')
$$
表示这个 $2$-范畴中 $1$-态射所成的范畴。下面给出纤维范畴情形下
米田引理的 $2$-范畴类似物。

\begin{lemma}[纤维范畴的 2-米田引理]
\label{categories-lemma-yoneda-2category-fibred}
令 $\mathcal{C}$ 为范畴。令 $\mathcal{S} \to \mathcal{C}$
为 $\mathcal{C}$ 上的纤维范畴。令 $U \in \Ob(\mathcal{C})$。
由 $G \mapsto G(\text{id}_U)$ 给出的函子
$$
\Mor_{\textit{Fib}/\mathcal{C}}(\mathcal{C}/U, \mathcal{S})
\longrightarrow
\mathcal{S}_U
$$
是等价。
\end{lemma}

\begin{proof}
为 $\mathcal{S}$ 选择拉回（见定义
\ref{categories-definition-pullback-functor-fibred-category}）。我们如下定义函子
$$
\mathcal{S}_U
\longrightarrow
\Mor_{\textit{Fib}/\mathcal{C}}(\mathcal{C}/U, \mathcal{S})
$$
给定 $x \in \Ob(\mathcal{S}_U)$，相应的函子满足：
\begin{enumerate}
\item 在对象上：$(f : V \to U) \mapsto f^*x$；
\item 在态射上：箭头 $(g : V'/U \to V/U)$ 映到复合
$$
(f \circ g)^*x \xrightarrow{(\alpha_{g, f})_x} g^*f^*x \rightarrow f^*x
$$
其中 $\alpha_{g, f}$ 如引理 \ref{categories-lemma-fibred} 中所定义。
\end{enumerate}
略去验证它是引理中函子的逆。
\end{proof}

\noindent
令 $\mathcal{C}$ 为范畴。$\mathcal{C}$ 上群胚纤维化范畴所成的
$2$-范畴，是 $\mathcal{C}$ 上范畴所成 $2$-范畴的一个“满”子
$2$-范畴（见定义
\ref{categories-definition-categories-fibred-in-groupoids-over-C}）。因此，若
$\mathcal{S}$、$\mathcal{S}'$ 是 $\mathcal{C}$ 上的群胚纤维化范畴，则
$$
\Mor_{\textit{Cat}/\mathcal{C}}(\mathcal{S}, \mathcal{S}')
$$
表示这个 $2$-范畴中 $1$-态射所成的范畴（见定义
\ref{categories-definition-categories-over-C}）。这些范畴全都是群胚，见定义
\ref{categories-definition-categories-fibred-in-groupoids-over-C} 后的说明。
下面给出米田引理的 $2$-范畴类似物。

\begin{lemma}[2-米田引理]
\label{categories-lemma-yoneda-2category}
令 $\mathcal{S}\to \mathcal{C}$ 为群胚纤维化范畴。
令 $U \in \Ob(\mathcal{C})$。由 $G \mapsto G(\text{id}_U)$ 给出的函子
$$
\Mor_{\textit{Cat}/\mathcal{C}}(\mathcal{C}/U, \mathcal{S})
\longrightarrow
\mathcal{S}_U
$$
是等价。
\end{lemma}

\begin{proof}
为 $\mathcal{S}$ 选择拉回（见定义
\ref{categories-definition-pullback-functor-fibred-category}）。我们如下定义函子
$$
\mathcal{S}_U
\longrightarrow
\Mor_{\textit{Cat}/\mathcal{C}}(\mathcal{C}/U, \mathcal{S})
$$
给定 $x \in \Ob(\mathcal{S}_U)$，相应的函子满足：
\begin{enumerate}
\item 在对象上：$(f : V \to U) \mapsto f^*x$；
\item 在态射上：箭头 $(g : V'/U \to V/U)$ 映到复合
$$
(f \circ g)^*x \xrightarrow{(\alpha_{g, f})_x} g^*f^*x \rightarrow f^*x
$$
其中 $\alpha_{g, f}$ 如引理 \ref{categories-lemma-fibred-groupoids} 中所定义。
\end{enumerate}
略去验证它是引理中函子的逆。
\end{proof}

\begin{remark}
\label{categories-remark-alternative-fibred-groupoids-strict}
我们可以用 $2$-米田引理给出引理
\ref{categories-lemma-fibred-groupoids-strict} 的另一证明。令
$p : \mathcal{S} \to \mathcal{C}$ 为群胚纤维化范畴。如下定义从
$\mathcal{C}$ 到群胚范畴的反变函子 $F$：对
$U\in \Ob(\mathcal{C})$，令
$$
F(U) = \Mor_{\textit{Cat}/\mathcal{C}}(\mathcal{C}/U, \mathcal{S}).
$$
若 $f : U \to V$，则诱导函子 $\mathcal{C}/U \to \mathcal{C}/V$
诱导态射 $F(f) : F(V) \to F(U)$。显然 $F$ 是函子。令
$\mathcal{S}'$ 为例 \ref{categories-example-functor-groupoids} 中与之相伴的
群胚纤维化范畴。有一个显然的 $\mathcal{C}$ 上函子
$G : \mathcal{S}' \to \mathcal{S}$，它把有序对 $(U, x)$——其中
$U \in \Ob(\mathcal{C})$ 且 $x \in F(U)$——映到
$x(\text{id}_U) \in \mathcal{S}$。现在引理
\ref{categories-lemma-yoneda-2category} 蕴含对每个 $U$，
$$
G_U : \mathcal{S}'_U = F(U)=
\Mor_{\textit{Cat}/\mathcal{C}}(\mathcal{C}/U, \mathcal{S})
\to
\mathcal{S}_U
$$
是等价，故由引理 \ref{categories-lemma-equivalence-fibred-categories}，
$G$ 是 $\mathcal{S}$ 与 $\mathcal{S}'$ 之间的等价。
\end{remark}





\section{可表 1-态射}
\label{categories-section-representable-1-morphisms}

\noindent
令 $\mathcal{C}$ 为范畴。本节说明 $\mathcal{C}$ 上两个群胚纤维化范畴之间
的 $1$-态射可表是什么意思。

\medskip\noindent
令 $\mathcal{C}$ 为范畴。令 $\mathcal{X}$、$\mathcal{Y}$ 为
$\mathcal{C}$ 上的群胚纤维化范畴。令 $U \in \Ob(\mathcal{C})$。
令 $F : \mathcal{X} \to \mathcal{Y}$ 和
$G : \mathcal{C}/U \to \mathcal{Y}$ 为 $\mathcal{C}$ 上群胚纤维化范畴
的 $1$-态射。我们想描述 $2$-纤维积
$$
\xymatrix{
(\mathcal{C}/U) \times_\mathcal{Y} \mathcal{X} \ar[r] \ar[d] &
\mathcal{X} \ar[d]^F \\
\mathcal{C}/U \ar[r]^G &
\mathcal{Y}
}
$$
令 $y = G(\text{id}_U) \in \mathcal{Y}_U$。为 $\mathcal{Y}$
选择拉回（见定义 \ref{categories-definition-pullback-functor-fibred-category}）。
那么 $G$ 同构于函子 $(f : V \to U) \mapsto f^*y$，见引理
\ref{categories-lemma-yoneda-2category} 及其证明。我们可以把
$(\mathcal{C}/U) \times_\mathcal{Y} \mathcal{X}$ 的对象看作四元组
$(V, f : V \to U, x, \phi)$，见引理
\ref{categories-lemma-2-product-categories-over-C}。利用上面对 $G$ 的描述，
可以把 $\phi$ 看作 $\mathcal{Y}_V$ 中的同构
$\phi : f^*y \to F(x)$。

\begin{lemma}
\label{categories-lemma-identify-fibre-product}
在上述情形中，
$(\mathcal{C}/U) \times_\mathcal{Y} \mathcal{X}$ 在
$\mathcal{C}/U$ 的对象 $f : V \to U$ 上的纤维范畴描述如下：
\begin{enumerate}
\item 对象是有序对 $(x, \phi)$，其中
$x \in \Ob(\mathcal{X}_V)$，且
$\phi : f^*y \to F(x)$ 是 $\mathcal{Y}_V$ 中的态射；
\item $(x, \phi)$ 与 $(x', \phi')$ 之间的态射集，是
$\mathcal{X}_V$ 中满足 $F(\psi) = \phi' \circ \phi^{-1}$ 的态射
$\psi : x \to x'$ 所成的集合。
\end{enumerate}
\end{lemma}

\begin{proof}
见上面的讨论。
\end{proof}

\begin{lemma}
\label{categories-lemma-prepare-representable-map-stack-in-groupoids}
令 $\mathcal{C}$ 为范畴。令 $\mathcal{X}$、$\mathcal{Y}$ 为
$\mathcal{C}$ 上的群胚纤维化范畴。令
$F : \mathcal{X} \to \mathcal{Y}$ 为 $1$-态射。令
$G : \mathcal{C}/U \to \mathcal{Y}$ 为 $1$-态射。那么
$$
(\mathcal{C}/U) \times_\mathcal{Y} \mathcal{X}
\longrightarrow
\mathcal{C}/U
$$
是群胚纤维化范畴。
\end{lemma}

\begin{proof}
我们已在引理 \ref{categories-lemma-2-product-fibred-categories} 中看到，复合
$$
(\mathcal{C}/U) \times_\mathcal{Y} \mathcal{X}
\longrightarrow
\mathcal{C}/U
\longrightarrow
\mathcal{C}
$$
是群胚纤维化范畴。于是该引理由引理
\ref{categories-lemma-cute-groupoids} 推出。
\end{proof}

\begin{definition}
\label{categories-definition-representable-map-categories-fibred-in-groupoids}
令 $\mathcal{C}$ 为范畴。令 $\mathcal{X}$、$\mathcal{Y}$ 为
$\mathcal{C}$ 上的群胚纤维化范畴。令
$F : \mathcal{X} \to \mathcal{Y}$ 为 $1$-态射。若对每个
$U \in \Ob(\mathcal{C})$ 及任意
$G : \mathcal{C}/U \to \mathcal{Y}$，群胚纤维化范畴
$$
(\mathcal{C}/U) \times_\mathcal{Y} \mathcal{X}
\longrightarrow
\mathcal{C}/U
$$
都是可表的，则称 $F$ 是{\it 可表的}，或称
{\it $\mathcal{X}$ 在 $\mathcal{Y}$ 上相对可表}。
\end{definition}

\begin{lemma}
\label{categories-lemma-spell-out-representable-map-stack-in-groupoids}
令 $\mathcal{C}$ 为范畴。令 $\mathcal{X}$、$\mathcal{Y}$ 为
$\mathcal{C}$ 上的群胚纤维化范畴。令
$F : \mathcal{X} \to \mathcal{Y}$ 为 $1$-态射。若 $F$ 可表，
则纤维范畴之间的每个函子
$$
F_U : \mathcal{X}_U \longrightarrow \mathcal{Y}_U
$$
都是忠实的。
\end{lemma}

\begin{proof}
由引理 \ref{categories-lemma-identify-fibre-product} 中对纤维范畴的描述，
以及引理 \ref{categories-lemma-characterize-representable-fibred-category}
中对可表纤维化范畴的刻画，这是显然的。
\end{proof}

\begin{lemma}
\label{categories-lemma-criterion-representable-map-stack-in-groupoids}
令 $\mathcal{C}$ 为范畴。令 $\mathcal{X}$、$\mathcal{Y}$ 为
$\mathcal{C}$ 上的群胚纤维化范畴。令
$F : \mathcal{X} \to \mathcal{Y}$ 为 $1$-态射。为
$\mathcal{Y}$ 选择拉回。假设
\begin{enumerate}
\item 纤维范畴之间的每个函子
$F_U : \mathcal{X}_U \longrightarrow \mathcal{Y}_U$ 都是忠实的；
\item 对每个 $U$ 及每个 $y \in \mathcal{Y}_U$，$\mathcal{C}/U$
上的预层
$$
(f : V \to U)
\longmapsto
\{(x, \phi) \mid x \in \mathcal{X}_V, \phi : f^*y \to F(x)\}/\cong
$$
是可表预层。
\end{enumerate}
那么 $F$ 可表。
\end{lemma}

\begin{proof}
由引理 \ref{categories-lemma-identify-fibre-product} 中对纤维范畴的描述，
以及引理 \ref{categories-lemma-characterize-representable-fibred-category}
中对可表纤维化范畴的刻画，这是显然的。
\end{proof}

\noindent
在陈述下一个引理之前，先指出群胚纤维化范畴所成的 $2$-范畴
是 $(2, 1)$-范畴，因而我们知道说它有终对象是什么意思（见定义
\ref{categories-definition-final-object-2-category}）。它确有终对象，即
$\text{id} : \mathcal{C} \to \mathcal{C}$。因此，我们把
$\mathcal{C}$ 上群胚纤维化范畴的 {\it $2$-积}定义为 $2$-纤维积
$$
\mathcal{X} \times \mathcal{Y} :=
\mathcal{X} \times_\mathcal{C} \mathcal{Y}.
$$
有了这个定义，下面的引理便有意义。

\begin{lemma}
\label{categories-lemma-representable-diagonal-groupoids}
令 $\mathcal{C}$ 为范畴。令 $\mathcal{S} \to \mathcal{C}$
为群胚纤维化范畴。假设 $\mathcal{C}$ 有对象对的积及纤维积。
下列条件等价：
\begin{enumerate}
\item 对角态射 $\mathcal{S} \to \mathcal{S} \times \mathcal{S}$ 可表；
\item 对 $\mathcal{C}$ 中每个 $U$，任意
$G : \mathcal{C}/U \to \mathcal{S}$ 都可表。
\end{enumerate}
\end{lemma}

\begin{proof}
假设对角态射可表，并给定 $U, G$。任取
$V \in \Ob(\mathcal{C})$ 和
$G' : \mathcal{C}/V \to \mathcal{S}$。注意
$\mathcal{C}/U \times \mathcal{C}/V = \mathcal{C}/U \times V$
可表。因此按假设，纤维积
$$
\xymatrix{
(\mathcal{C}/U \times V)
\times_{(\mathcal{S} \times \mathcal{S})}
\mathcal{S}
\ar[r] \ar[d] &
\mathcal{S} \ar[d] \\
\mathcal{C}/U \times V \ar[r]^{(G, G')} &
\mathcal{S} \times \mathcal{S}
}
$$
可表。这意味着存在 $\mathcal{C}$ 中的 $W \to U \times V$，使得
$$
\xymatrix{
\mathcal{C}/W \ar[d] \ar[r] & \mathcal{S} \ar[d] \\
\mathcal{C}/U \times \mathcal{C}/V \ar[r] & \mathcal{S} \times \mathcal{S}
}
$$
为笛卡尔图。由此得到所需的
$\mathcal{C}/W \cong \mathcal{C}/U \times_\mathcal{S} \mathcal{C}/V$
（见引理 \ref{categories-lemma-diagonal-1} 和注
\ref{categories-remark-2-product-fibred-categories-same}）。

\medskip\noindent
假设 (2) 成立。任取 $V \in \Ob(\mathcal{C})$ 及
$(G, G') : \mathcal{C}/V \to \mathcal{S} \times \mathcal{S}$。
我们必须证明
$\mathcal{C}/V \times_{\mathcal{S} \times \mathcal{S}} \mathcal{S}$
可表。已知
$\mathcal{C}/V \times_{G, \mathcal{S}, G'} \mathcal{C}/V$
可表，设它由 $\mathcal{C}/V$ 中的 $a : W \to V$ 表示。
等价
$$
\mathcal{C}/W \to \mathcal{C}/V \times_{G, \mathcal{S}, G'} \mathcal{C}/V
$$
再接第二投影到 $\mathcal{C}/V$，给出第二个态射
$a' : W \to V$。考虑
$W' = W \times_{(a, a'), V \times V} V$。存在等价
$$
\mathcal{C}/W' \cong
\mathcal{C}/V \times_{\mathcal{S} \times \mathcal{S}} \mathcal{S}
$$
即
\begin{eqnarray*}
\mathcal{C}/W' & \cong &
\mathcal{C}/W \times_{(\mathcal{C}/V \times \mathcal{C}/V)} \mathcal{C}/V \\
& \cong &
\left(\mathcal{C}/V \times_{(G, \mathcal{S}, G')} \mathcal{C}/V\right)
\times_{(\mathcal{C}/V \times \mathcal{C}/V)} \mathcal{C}/V \\
& \cong &
\mathcal{C}/V \times_{(\mathcal{S} \times \mathcal{S})} \mathcal{S}
\end{eqnarray*}
（最后一个同构见引理 \ref{categories-lemma-diagonal-2} 和注
\ref{categories-remark-2-product-fibred-categories-same}），这就证明了引理。
\end{proof}

\noindent
{\bf 文献说明：}
本节部分内容取自 Vistoli 的讲义 \cite{Vis2}。





\section{幺半范畴}
\label{categories-section-monoidal}

\noindent
令 $\mathcal{C}$ 为范畴。假设给定函子
$$
\otimes : \mathcal{C} \times \mathcal{C} \longrightarrow \mathcal{C}
$$
我们常常想知道 $\otimes$ 是否满足结合律，以及 $\otimes$ 是否有单位元。

\medskip\noindent
$(\mathcal{C}, \otimes)$ 的一个{\it 结合约束}，是函子性的同构
$$
\phi_{X, Y, Z} : X \otimes (Y \otimes Z) \to (X \otimes Y) \otimes Z
$$
使得对所有对象 $X, Y, Z, W$，图
\begin{equation}
\label{categories-equation-pentagon}
\vcenter{
\xymatrix{
X \otimes (Y \otimes ( Z \otimes W)) \ar[r] \ar[d] &
X \otimes ((Y \otimes Z) \otimes W) \ar[d] \\
(X \otimes Y) \otimes (Z \otimes W) \ar[d] &
(X \otimes (Y \otimes Z)) \otimes W \ar[ld] \\
((X \otimes Y) \otimes Z) \otimes W
}
}
\end{equation}
交换，其中每个箭头都由对 $\phi$ 的适当应用和 $\otimes$ 的函子性确定。

\medskip\noindent
文献 \cite[Theorem 3.1]{associativity} 说明了上述三元组
$(\mathcal{C}, \otimes, \phi)$ 如何具有协调性。换一种表述，这样的三元组
对所有 $n \geq 1$ 给出良定义的函子系
$$
\mathcal{C} \times \ldots \times \mathcal{C} \longrightarrow \mathcal{C},
\quad
(X_1, \ldots, X_n) \longmapsto X_1 \otimes \ldots \otimes X_n
$$
并对 $n, m \geq 1$ 给出函子性同构
$$
(X_1 \otimes \ldots \otimes X_n) \otimes (Y_1 \otimes \ldots \otimes Y_m)
\to
X_1 \otimes \ldots \otimes X_n \otimes Y_1 \otimes \ldots \otimes Y_m
$$
使得由这些同构组成的一切可能的图都交换（上面的五角星图是一个例子）。
当 $n = 5$ 时，这意味着各表达式
$$
\begin{matrix}
(((X_1 \otimes X_2) \otimes X_3) \otimes X_4) \otimes X_5 &
((X_1 \otimes (X_2 \otimes X_3)) \otimes X_4) \otimes X_5 \\
((X_1 \otimes X_2) \otimes (X_3 \otimes X_4)) \otimes X_5 &
(X_1 \otimes ((X_2 \otimes X_3) \otimes X_4)) \otimes X_5 \\
(X_1 \otimes (X_2 \otimes (X_3 \otimes X_4))) \otimes X_5 &
((X_1 \otimes X_2) \otimes X_3) \otimes (X_4 \otimes X_5) \\
(X_1 \otimes (X_2 \otimes X_3)) \otimes (X_4 \otimes X_5) &
(X_1 \otimes X_2) \otimes ((X_3 \otimes X_4) \otimes X_5) \\
(X_1 \otimes X_2) \otimes (X_3 \otimes (X_4 \otimes X_5)) &
X_1 \otimes (((X_2 \otimes X_3) \otimes X_4) \otimes X_5) \\
X_1 \otimes ((X_2 \otimes (X_3 \otimes X_4)) \otimes X_5) &
X_1 \otimes ((X_2 \otimes X_3) \otimes (X_4 \otimes X_5)) \\
X_1 \otimes (X_2 \otimes ((X_3 \otimes X_4) \otimes X_5)) &
X_1 \otimes (X_2 \otimes (X_3 \otimes (X_4 \otimes X_5)))
\end{matrix}
$$
都函子性同构于 $X_1 \otimes \ldots \otimes X_5$，而且这些同构与
利用 $\phi$ 在三个连续位置（若可能）得到的这些函子之间的一切可能同构
全都相容。以下我们将不加说明地使用这一点，并且在书写 $\otimes$ 的迭代时
不再写括号。

\medskip\noindent
上述三元组 $(\mathcal{C}, \otimes, \phi)$ 的一个{\it 单位元}，是
$\mathcal{C}$ 的对象 $\mathbf{1}$，连同函子性同构
$$
l : \mathbf{1} \otimes X \to X
\quad\text{且}\quad
r : X \otimes \mathbf{1} \to X
$$
使得对所有对象 $X, Y$，图
\begin{equation}
\label{categories-equation-triangle}
\vcenter{
\xymatrix{
X \otimes (\mathbf{1} \otimes Y) \ar[rr]_\phi \ar[rd]_{\text{id} \otimes l} & &
(X \otimes \mathbf{1}) \otimes Y \ar[ld]^{r \otimes \text{id}} \\
& X \otimes Y
}
}
\end{equation}
交换。我们常把单位元看作下面引理中的有序对 $(\mathbf{1}, 1)$。

\begin{lemma}
\label{categories-lemma-monoidal-unit}
令 $(\mathcal{C}, \otimes, \phi)$ 如上。$\mathcal{C}$ 中的单位元
$(\mathbf{1}, l, r)$ 与有序对 $(\mathbf{1}, 1)$ 一一对应，其中
$\mathbf{1}$ 是 $\mathcal{C}$ 的对象，且
$1 : \mathbf{1} \otimes \mathbf{1} \to \mathbf{1}$ 是同构，并且函子
$L : X \mapsto \mathbf{1} \otimes X$ 与
$R : X \mapsto X \otimes \mathbf{1}$ 都是等价。
\end{lemma}

\begin{proof}
给定单位元 $(\mathbf{1}, l, r)$，我们得到同构
$r : \mathbf{1} \otimes \mathbf{1} \to \mathbf{1}$；又因为 $L$ 与 $R$
都同构于恒等函子，所以它们是等价。反过来，假设给定
$(\mathbf{1}, 1)$，并且 $L$ 与 $R$ 是等价。我们得到函子性同构
$l_X : \mathbf{1} \otimes X \to X$ 和
$r_X : X \otimes \mathbf{1} \to X$，它们分别由
$L(l_X) = 1 \otimes \text{id}_X$ 和
$R(r_X) = \text{id}_X \otimes 1$ 所刻画。于是必须证明，对所有
$X$ 和 $Y$，从 $X \otimes \mathbf{1} \otimes Y$ 到
$X \otimes Y$ 的两个箭头 $\text{id}_X \otimes l_Y$ 与
$r_X \otimes \text{id}_Y$ 相同。这个性质只依赖于 $X$ 与 $Y$ 的
同构类。由于 $R$ 与 $L$ 是等价，只需对
$X = Z \otimes \mathbf{1}$ 和 $Y = \mathbf{1} \otimes W$ 证明它，
其中 $Z$ 和 $W$ 是某些对象。换言之，我们必须证明
$$
\text{id}_Z \otimes \text{id}_\mathbf{1} \otimes l_{\mathbf{1} \otimes W}
=
r_{Z \otimes \mathbf{1}} \otimes \text{id}_\mathbf{1} \otimes \text{id}_Y
$$
按构造，这两个映射分别等于
$\text{id}_Z \otimes 1 \otimes \text{id}_\mathbf{1} \otimes \text{id}_W$
与
$\text{id}_Z \otimes \text{id}_\mathbf{1} \otimes 1 \otimes \text{id}_W$。
因此只需证明
$1 \otimes \text{id}_\mathbf{1} = \text{id}_\mathbf{1} \otimes 1$。

\medskip\noindent
我们有
$\text{id}_\mathbf{1} \otimes 1 = r_\mathbf{1} \otimes \text{id}_\mathbf{1}$
以及
$l_\mathbf{1} \otimes \text{id}_\mathbf{1} = \text{id}_\mathbf{1} \otimes 1$。
可对 $\mathbf{1}$ 的某些自同构 $a, b$ 写成
$r_\mathbf{1} = a \circ 1$ 和 $l_\mathbf{1} = b \circ 1$。于是
$\text{id}_\mathbf{1} \otimes 1 =
(a \otimes \text{id}_\mathbf{1}) \circ (1 \otimes \text{id}_\mathbf{1})$
且
$1 \otimes \text{id}_\mathbf{1} =
(\text{id}_\mathbf{1} \otimes b) \circ (\text{id}_\mathbf{1} \otimes 1)$。
因此可写成
\begin{align*}
(1 \otimes \text{id}_\mathbf{1}) \circ
(\text{id}_\mathbf{1} \otimes 1 \otimes \text{id}_\mathbf{1})
& =
(1 \otimes \text{id}_\mathbf{1}) \circ
(\text{id}_\mathbf{1} \otimes \text{id}_\mathbf{1} \otimes b)
\circ
(\text{id}_\mathbf{1} \otimes \text{id}_\mathbf{1} \otimes 1) \\
& =
(\text{id}_\mathbf{1} \otimes b) \circ
(1 \otimes \text{id}_\mathbf{1}) \circ
(\text{id}_\mathbf{1} \otimes \text{id}_\mathbf{1} \otimes 1) \\
& =
(\text{id}_\mathbf{1} \otimes b) \circ
(1 \otimes 1)
\end{align*}
并且还有
\begin{align*}
(1 \otimes \text{id}_\mathbf{1}) \circ
(\text{id}_\mathbf{1} \otimes 1 \otimes \text{id}_\mathbf{1})
& =
(\text{id}_\mathbf{1} \otimes b) \circ
(\text{id}_\mathbf{1} \otimes 1) \circ
(a \otimes \text{id}_\mathbf{1} \otimes \text{id}_\mathbf{1}) \circ
(1 \otimes \text{id}_\mathbf{1} \otimes \text{id}_\mathbf{1}) \\
& =
(\text{id}_\mathbf{1} \otimes b) \circ
(a \otimes \text{id}_\mathbf{1}) \circ
(1 \otimes 1)
\end{align*}
这证明 $a \otimes \text{id}_\mathbf{1}$ 是恒等态射，从而如所需，
$a$ 是恒等态射。

\medskip\noindent
为完成证明，注意上述规则在单位元所成的范畴（适当定义）与有序对
$(\mathbf{1}, 1)$ 所成的范畴之间给出互逆等价。
\end{proof}

\begin{lemma}
\label{categories-lemma-monoidal-unit-unique}
令 $(\mathcal{C}, \otimes, \phi)$ 如上。令 $(\mathbf{1}, 1)$ 为单位元
（见引理 \ref{categories-lemma-monoidal-unit}）。那么
\begin{enumerate}
\item $1 \otimes \text{id}_\mathbf{1} = \text{id}_\mathbf{1} \otimes 1$
\item $\Gamma = \text{Mor}(\mathbf{1}, \mathbf{1})$ 是交换幺半群；
\item  $a = 1 \circ (a \otimes \text{id}_\mathbf{1}) \circ 1^{-1} =
1 \circ (\text{id}_\mathbf{1} \otimes a) \circ 1^{-1}$ 对所有
$a \in \Gamma$ 成立；
\item 任何其他单位元都经唯一同构同构于 $(\mathbf{1}, 1)$。
\end{enumerate}
\end{lemma}

\begin{proof}
第 (1) 部分已在引理 \ref{categories-lemma-monoidal-unit} 的证明中证得。
对 $a \in \Gamma$，有
\begin{align*}
(1 \circ (a \otimes \text{id}_\mathbf{1})) \otimes \text{id}_\mathbf{1}
& =
(1 \otimes \text{id}_\mathbf{1}) \circ
(a \otimes \text{id}_\mathbf{1} \otimes \text{id}_\mathbf{1}) \\
& =
(\text{id}_\mathbf{1} \otimes 1) \circ
(a \otimes \text{id}_\mathbf{1} \otimes \text{id}_\mathbf{1}) \\
& =
(a \otimes \text{id}_\mathbf{1}) \circ
(\text{id}_\mathbf{1} \otimes 1) \\
& =
(a \otimes \text{id}_\mathbf{1}) \circ
(1 \otimes \text{id}_\mathbf{1}) \\
& =
(a \circ 1) \otimes \text{id}_\mathbf{1}
\end{align*}
所以 $1 \circ (a \otimes \text{id}_\mathbf{1}) = a \circ 1$，
这给出 (3) 的一半。另一半同理。为证 (2)，注意对
$a, b \in \Gamma$，有
$$
a \circ b =
1 \circ (a \otimes \text{id}_\mathbf{1}) \circ 1^{-1}
\circ
1 \circ (\text{id}_\mathbf{1} \otimes b) \circ 1^{-1}
=
1 \circ (a \otimes b) \circ \circ 1^{-1}
$$
而 $b \circ a$ 也计算为同一表达式。因此 (2) 成立。为证 (4)，
设 $(\mathbf{1}', 1')$ 是第二个单位元。利用 $r$ 与 $l'$，有同构
$\mathbf{1}' \otimes \mathbf{1} \to \mathbf{1}'$ 和
$\mathbf{1}' \otimes \mathbf{1} \to \mathbf{1}$。因此存在同构
$t : \mathbf{1}' \to \mathbf{1}$。图
$$
\xymatrix{
\mathbf{1}' \otimes \mathbf{1}' \ar[r]_{t \otimes t} \ar[d]_{1'} & 
\mathbf{1} \otimes \mathbf{1} \ar[d]^1 \\
\mathbf{1}' \ar[r]^t & \mathbf{1}
}
$$
至多相差 $\mathbf{1}$ 的一个自同构 $a$ 而交换。把 $t$ 替换为
$a \circ t$ 后，该图将交换（提示：利用 (3) 看出
$1 \circ (a \otimes a) = a^2 \circ 1$）。
\end{proof}

\begin{definition}
\label{categories-definition-monoidal-category}
三元组 $(\mathcal{C}, \otimes, \phi)$ 中，$\mathcal{C}$ 是范畴，
$\otimes : \mathcal{C} \times \mathcal{C} \to \mathcal{C}$ 是函子，
$\phi$ 是结合约束；若存在单位元 $\mathbf{1}$，则称该三元组为
{\it 幺半范畴}。
\end{definition}

\noindent
我们总以 $\mathbf{1}$ 表示幺半范畴的一个单位元，并以
$1 : \mathbf{1} \otimes \mathbf{1} \to \mathbf{1}$ 表示一个选定同构；
由于有序对 $(\mathbf{1}, 1)$ 在唯一同构意义下确定（引理
\ref{categories-lemma-monoidal-unit-unique}），这样选择并无妨碍。

\begin{lemma}
\label{categories-lemma-monoidal-coherence}
在幺半范畴 $\mathcal{C}, \otimes, \phi, \mathbf{1}, 1$ 中，沿用引理
\ref{categories-lemma-monoidal-unit} 证明中的记号，我们有
\begin{enumerate}
\item 箭头 $1, r_\mathbf{1}, l_\mathbf{1} :
\mathbf{1} \otimes \mathbf{1} \to \mathbf{1}$ 相同；
\item 箭头
$l_X \otimes \text{id}_Y, l_{X \otimes Y} :
\mathbf{1} \otimes X \otimes Y \to X \otimes Y$ 相同；
\item 箭头
$\text{id}_X \otimes r_Y , r_{X \otimes Y} :
X \otimes Y \otimes \mathbf{1} \to X \otimes Y$ 相同。
\end{enumerate}
幺半范畴满足 \cite[Theorem 5.2]{associativity} 的假设。
\end{lemma}

\begin{proof}
引理 \ref{categories-lemma-monoidal-unit} 的证明已经给出 (1)。在引理
\ref{categories-lemma-monoidal-unit} 的证明中还看到，
$l_X$ 与 $l_{X \otimes Y}$ 分别是满足
$\text{id}_\mathbf{1} \otimes l_{X \otimes Y} =
1 \otimes \text{id}_{X \otimes Y}$ 和
$\text{id}_\mathbf{1} \otimes l_X = 1 \otimes \text{id}_X$ 的唯一态射。
第 (2) 部分立即推出。第 (3) 部分同理可证。结合
(\ref{categories-equation-pentagon}) 与 (\ref{categories-equation-triangle}) 的交换性，
这说明引理的最后一个断言成立。
\end{proof}

\noindent
文献 \cite[Theorem 5.2]{associativity} 说明了上述引理中的五元组
$(\mathcal{C}, \otimes, \phi, \mathbf{1}, 1)$ 如何具有协调性。
换一种表述，在我们所定义的幺半范畴中，对每个
$n \geq i \geq 0$ 都有函子性同构
$$
X_1 \otimes \ldots \otimes X_i \otimes \mathbf{1}
\otimes X_{i + 1} \otimes \ldots \otimes X_n
\to
X_1 \otimes \ldots \otimes X_n
$$
使得由这些同构组成的一切可能的图都交换。例如，从插入两个
$\mathbf{1}$ 开始，以不同次序使用这些同构，所得态射相同。除了书写
$\otimes$ 的迭代时省略括号这一约定外，从现在起还把
$X \otimes \mathbf{1}$ 和 $\mathbf{1} \otimes X$ 与 $X$ 等同，
不再另作说明。此外，我们会说“令 $\mathcal{C}$ 为幺半范畴”，
并默认给定 $\otimes, \phi, \mathbf{1}$。

\begin{definition}
\label{categories-definition-functor-monoidal-categories}
令 $\mathcal{C}$ 与 $\mathcal{C}'$ 为幺半范畴。幺半范畴的一个
{\it 函子} $F : \mathcal{C} \to \mathcal{C}'$，由所示函子 $F$
以及关于 $X$、$Y$ 函子性的同构
$$
F(X) \otimes F(Y) \to F(X \otimes Y)
$$
给出，使得对所有对象 $X$、$Y$、$Z$，图
$$
\xymatrix{
F(X) \otimes (F(Y) \otimes F(Z)) \ar[r] \ar[d] &
F(X) \otimes F(Y \otimes Z) \ar[r] &
F(X \otimes (Y \otimes Z)) \ar[d] \\
(F(X) \otimes F(Y)) \otimes F(Z) \ar[r] &
F(X \otimes Y) \otimes F(Z) \ar[r] &
F((X \otimes Y) \otimes Z)
}
$$
交换，并且 $F(\mathbf{1})$ 是 $\mathcal{C}'$ 中的单位元。
\end{definition}

\noindent
依照我们关于单位元的约定，若 $F$ 是幺半范畴的函子，总可假定
$F(\mathbf{1}) = \mathbf{1}$。例如，若 $A \to B$ 是环同态，则函子
$M \mapsto M \otimes_A B$ 是从 $\text{Mod}_A$ 到 $\text{Mod}_B$
的幺半范畴函子。

\begin{lemma}
\label{categories-lemma-invertible}
令 $\mathcal{C}$ 为幺半范畴。令 $X$ 为 $\mathcal{C}$ 的对象。
下列条件等价：
\begin{enumerate}
\item 函子 $L : Y \mapsto X \otimes Y$ 是等价；
\item 函子 $R : Y \mapsto Y \otimes X$ 是等价；
\item 存在对象 $X'$，使得
$X \otimes X' \cong X' \otimes X \cong \mathbf{1}$。
\end{enumerate}
\end{lemma}

\begin{proof}
假设 (1)。选择 $X'$ 使 $L(X') = \mathbf{1}$，即
$X \otimes X' \cong \mathbf{1}$。以 $L'$ 和 $R'$ 表示与 $X'$
对应的函子。等式 $X \otimes X' \cong \mathbf{1}$ 蕴含
$L \circ L' \cong \text{id}$。因此 $L'$ 必为 $L$ 的拟逆
（它依假设存在）。所以 $L' \circ L \cong \text{id}$，进而
$X' \otimes X \cong \mathbf{1}$。故 (3) 成立。

\medskip\noindent
(2) $\Rightarrow$ (3) 的证明与刚才的论证对偶。

\medskip\noindent
假设 (3)。显然 $L'$ 与 $L$ 互为拟逆，且 $R'$ 与 $R$ 互为拟逆。
故 (1) 与 (2) 成立。
\end{proof}

\begin{definition}
\label{categories-definition-invertible}
令 $\mathcal{C}$ 为幺半范畴。若引理 \ref{categories-lemma-invertible}
中的任一（从而全部）等价条件成立，则称 $\mathcal{C}$ 的对象 $X$
是{\it 可逆的}。
\end{definition}

\noindent
注意，若 $F : \mathcal{C} \to \mathcal{C}'$ 是幺半范畴的函子，
则 $F$ 把可逆对象映到可逆对象。

\begin{definition}
\label{categories-definition-dual}
给定幺半范畴 $(\mathcal{C}, \otimes, \phi)$ 及对象 $X$，一个
{\it 左对偶}是对象 $Y$，连同态射
$\eta : \mathbf{1} \to X \otimes Y$ 和
$\epsilon : Y \otimes X \to \mathbf{1}$，使得图
$$
\vcenter{
\xymatrix{
X \ar[rd]_1 \ar[r]_-{\eta \otimes 1} &
X \otimes Y \otimes X \ar[d]^{1 \otimes \epsilon} \\
& X
}
}
\quad\text{且}\quad
\vcenter{
\xymatrix{
Y \ar[rd]_1 \ar[r]_-{1 \otimes \eta} &
Y \otimes X \otimes Y \ar[d]^{\epsilon \otimes 1} \\
& Y
}
}
$$
交换。在这种情形下，我们称 $X$ 是 $Y$ 的一个{\it 右对偶}。
\end{definition}

\noindent
注意，若 $F : \mathcal{C} \to \mathcal{C}'$ 是幺半范畴的函子，
且 $Y$ 是 $X$ 的左对偶，则 $F(Y)$ 是 $F(X)$ 的左对偶。

\begin{lemma}
\label{categories-lemma-left-dual}
令 $\mathcal{C}$ 为幺半范畴。若 $Y$ 是 $X$ 的左对偶，则
$$
\Mor(Z' \otimes X, Z) = \Mor(Z', Z \otimes Y)
\quad\text{且}\quad
\Mor(Y \otimes Z', Z) = \Mor(Z', X \otimes Z)
$$
关于 $Z$ 与 $Z'$ 函子性地成立。
\end{lemma}

\begin{proof}
考虑映射
$$
\Mor(Z' \otimes X, Z) \to
\Mor(Z' \otimes X \otimes Y, Z \otimes Y) \to
\Mor(Z', Z \otimes Y)
$$
其中第二个箭头使用 $\eta$；以及映射序列
$$
\Mor(Z', Z \otimes Y) \to
\Mor(Z' \otimes X, Z \otimes Y \otimes X) \to
\Mor(Z' \otimes X, Z)
$$
其中第二个箭头使用 $\epsilon$。为证明这些箭头互逆，考虑映射
$a : Z' \to Z \otimes Y$。我们必须证明
$$
Z' \xrightarrow{\text{id}_{Z'} \otimes \eta}
Z' \otimes X \otimes Y \xrightarrow{a \otimes \text{id}_{X \otimes Y}}
Z \otimes Y \otimes X \otimes Y
\xrightarrow{\text{id}_Z \otimes \epsilon \otimes \text{id}_Y}
Z \otimes Y
$$
等于 $a$。前两个箭头的复合等于
$(\text{id}_{Z \otimes Y} \otimes \eta) \circ a :
Z' \to Z \otimes Y \to Z \otimes Y \otimes X \otimes Y$。
而 $\text{id}_{Z \otimes Y} \otimes \eta$ 与
$\text{id}_Z \otimes \epsilon \otimes \text{id}_Y$ 的复合，
依对偶的定义等于恒等态射。另一个复合同理。略去第二个等式的证明。
\end{proof}

\begin{remark}
\label{categories-remark-left-dual-adjoint}
引理 \ref{categories-lemma-left-dual} 特别说明，$Z \mapsto Z \otimes Y$ 是
$Z' \mapsto Z' \otimes X$ 的右伴随。特别地，伴随函子的唯一性保证，
若 $X$ 的左对偶存在，则它在唯一同构意义下唯一。反过来，假设函子
$Z \mapsto Z \otimes Y$ 是函子 $Z' \mapsto Z' \otimes X$ 的右伴随，
即给定关于 $Z$ 与 $Z'$ 均函子性的双射
$$
\Mor(Z' \otimes X, Z) \longrightarrow \Mor(Z', Z \otimes Y)
$$
伴随的单位给出关于 $Z$ 函子性的映射
$$
\eta_Z : Z \to Z \otimes X \otimes Y
$$
而伴随的余单位给出关于 $Z'$ 函子性的映射
$$
\epsilon_{Z'} : Z' \otimes Y \otimes X \to Z'
$$
特别地，我们得到
$\eta = \eta_\mathbf{1} : \mathbf{1} \to X \otimes Y$ 和
$\epsilon = \epsilon_\mathbf{1} : Y \otimes X \to \mathbf{1}$。
作为关于单位、余单位与伴随同构之关系的一个练习，读者可证明
$$
(\epsilon \otimes \text{id}_Y) \circ \eta_Y = \text{id}_Y
\quad\text{且}\quad
\epsilon_X \circ (\eta \otimes \text{id}_X) = \text{id}_X
$$
然而，这不足以证明
$(\epsilon \otimes \text{id}_Y) \circ (\text{id}_Y \otimes \eta) =
\text{id}_Y$ 以及
$(\text{id}_X \otimes \epsilon) \circ (\eta \otimes \text{id}_X) =
\text{id}_X$，因为一般而言，我们不知道
$\eta_Y = \text{id}_Y \otimes \eta$，也不知道
$\epsilon_X = \epsilon \otimes \text{id}_X$。为此，只要知道我们的伴随同构
具有下述性质便足够：对每个 $W, Z, Z'$，图
$$
\xymatrix{
\Mor(Z' \otimes X, Z) \ar[r] \ar[d]_{\text{id}_W \otimes -} &
\Mor(Z', Z \otimes Y) \ar[d]^{\text{id}_W \otimes -} \\
\Mor(W \otimes Z' \otimes X, W \otimes Z) \ar[r] &
\Mor(W \otimes Z', W \otimes Z \otimes Y)
}
$$
交换。若此条件成立，我们就说{\it 该伴随与给定张量结构相容}。
因此，要求 $Z \mapsto Z \otimes Y$ 是
$Z' \mapsto Z' \otimes X$ 的、与给定张量结构相容的右伴随，
等价于左对偶这一性质。
\end{remark}

\begin{lemma}
\label{categories-lemma-tensor-dual}
令 $\mathcal{C}$ 为幺半范畴。若 $Y_i$（$i = 1, 2$）分别是
$X_i$（$i = 1, 2$）的左对偶，则 $Y_2 \otimes Y_1$ 是
$X_1 \otimes X_2$ 的左对偶。
\end{lemma}

\begin{proof}
由伴随的唯一性和注 \ref{categories-remark-left-dual-adjoint} 得证。
\end{proof}

\noindent
$(\mathcal{C}, \otimes)$ 的一个{\it 交换约束}，是函子性同构
$$
\psi : X \otimes Y \longrightarrow Y \otimes X
$$
使得复合
$$
X \otimes Y \xrightarrow{\psi} Y \otimes X \xrightarrow{\psi} X \otimes Y
$$
是恒等态射。若对所有对象 $X, Y, Z$，图
\begin{equation}
\label{categories-equation-hexagon}
\vcenter{
\xymatrix{
X \otimes (Y \otimes Z) \ar[r]_\phi \ar[d]^\psi &
(X \otimes Y) \otimes Z \ar[r]_\psi &
Z \otimes (X \otimes Y) \ar[d]^\phi \\
X \otimes (Z \otimes Y) \ar[r]^\phi &
(X \otimes Z) \otimes Y \ar[r]^\psi &
(Z \otimes X) \otimes Y
}
}
\end{equation}
交换，则称 $\psi$ 与给定结合约束 $\phi$ {\it 相容}。

\begin{definition}
\label{categories-definition-symmetric-monoidal-category}
四元组 $(\mathcal{C}, \otimes, \phi, \psi)$ 中，$\mathcal{C}$ 是范畴，
$\otimes : \mathcal{C} \otimes \mathcal{C} \to \mathcal{C}$ 是函子，
$\phi$ 是结合约束，$\psi$ 是与 $\phi$ 相容的交换约束；若存在单位元，
则称该四元组为{\it 对称幺半范畴}。
\end{definition}

\noindent
确切地说，若 $(\mathcal{C}, \otimes, \phi, \psi)$ 是对称幺半范畴，
则 $(\mathcal{C}, \otimes, \phi)$ 是幺半范畴，因而可以使用上面讨论的
语言和记号。

\begin{lemma}
\label{categories-lemma-symmetric-monoidal-coherence}
在对称幺半范畴
$\mathcal{C}, \otimes, \phi, \psi, \mathbf{1}, 1$ 中，我们有
\begin{enumerate}
\item 箭头 $1 \circ \psi, 1 :
\mathbf{1} \otimes \mathbf{1} \to \mathbf{1}$ 相同；
\item 箭头 $\text{id}_X \otimes l_Y,
(l_X \otimes \text{id}) \circ (\psi \otimes \text{id}_Y):
X \otimes \mathbf{1} \otimes Y \to X \otimes Y$ 相同；
\end{enumerate}
对称幺半范畴满足 \cite[Theorem 5.1]{associativity} 的假设。
\end{lemma}

\begin{proof}
可写成 $\psi = a \otimes \text{id}_\mathbf{1}$，其中
$a : \mathbf{1} \to \mathbf{1}$ 是唯一同构。引理
\ref{categories-lemma-monoidal-unit-unique} 蕴含
$a \otimes \text{id}_\mathbf{1} = \text{id}_\mathbf{1} \otimes a$。
$\psi$ 的函子性说明图
$$
\xymatrix{
(\mathbf{1} \otimes \mathbf{1}) \otimes \mathbf{1}
\ar[r]_\psi
\ar[d]_{1 \otimes \text{id}_\mathbf{1}} &
\mathbf{1} \otimes (\mathbf{1} \otimes \mathbf{1})
\ar[d]^{\text{id}_\mathbf{1} \otimes 1} \\
\mathbf{1} \otimes \mathbf{1}  \ar[r]^\psi &
\mathbf{1} \otimes \mathbf{1}
}
$$
交换。因此上方箭头等于
$a \otimes \text{id}_\mathbf{1} \otimes \text{id}_\mathbf{1}$。
于是对 $X = Y = Z = \mathbf{1}$，(\ref{categories-equation-hexagon}) 表明
$a \otimes \text{id}_\mathbf{1} \otimes \text{id}_\mathbf{1}$ 等于自身的
平方。所以 $a = \text{id}_\mathbf{1}$。这就证明了 (1)。

\medskip\noindent
第 (2) 部分断言，如果在我们的幺半范畴中把源和靶都与 $X$ 等同，
则 $\psi : X \otimes \mathbf{1} \to \mathbf{1} \otimes X$ 是恒等态射。
这由 (\ref{categories-equation-hexagon}) 对 $\mathbf{1}, X, \mathbf{1}$ 的交换性推出，
即
$$
\xymatrix{
\mathbf{1} \otimes (X \otimes \mathbf{1}) \ar[r]_\phi \ar[d]^\psi &
(\mathbf{1} \otimes X) \otimes \mathbf{1} \ar[r]_\psi &
\mathbf{1} \otimes (\mathbf{1} \otimes X) \ar[d]^\phi \\
\mathbf{1} \otimes (\mathbf{1} \otimes X) \ar[r]^\phi &
(\mathbf{1} \otimes \mathbf{1}) \otimes X \ar[r]^\psi &
(\mathbf{1} \otimes \mathbf{1}) \otimes X
}
$$
交换，而我们知道此图中除一个以外的所有态射 $\psi$ 都等于恒等态射。

\medskip\noindent
除 (1) 与 (2) 外，图 (\ref{categories-equation-pentagon})、
(\ref{categories-equation-triangle})、(\ref{categories-equation-hexagon}) 的交换性以及引理
\ref{categories-lemma-monoidal-coherence} 的结果，蕴含引理的最后一个断言。
\end{proof}

\noindent
文献 \cite[Theorem 5.1]{associativity} 说明了上述引理中的六元组
$(\mathcal{C}, \otimes, \phi, \psi, \mathbf{1}, 1)$ 如何具有协调性。
换一种表述，在我们所定义的对称幺半范畴中，对每个 $n \geq 1$
以及 $\{1, \ldots, n\}$ 的每个置换 $\sigma$，都有函子性同构
$$
X_1 \otimes \ldots \otimes X_n
\to
X_{\sigma(1)} \otimes \ldots \otimes X_{\sigma(n)}
$$
使得由这些同构组成的一切可能的图都交换，并且这些同构与幺半范畴的
结构相容（例如，与在某一位置插入 $\mathbf{1}$ 时的同构相容）。

\begin{lemma}
\label{categories-lemma-dual-symmetric}
令 $(\mathcal{C}, \otimes, \phi, \psi)$ 为对称幺半范畴。
令 $X$ 为 $\mathcal{C}$ 的对象，并令 $Y$、
$\eta : \mathbf{1} \to X \otimes Y$、
$\epsilon : Y \otimes X \to \mathbf{1}$ 为定义
\ref{categories-definition-dual} 中 $X$ 的一个左对偶。那么
$\eta' = \psi \circ \eta : \mathbf{1} \to Y \otimes X$ 和
$\epsilon' = \epsilon \circ \psi : X \otimes Y \to \mathbf{1}$
使 $X$ 成为 $Y$ 的一个左对偶。
\end{lemma}

\begin{proof}
略。提示：这是一个愉快的定义练习。
\end{proof}

\begin{definition}
\label{categories-definition-functor-symmetric-monoidal-categories}
令 $\mathcal{C}$ 与 $\mathcal{C}'$ 为对称幺半范畴。对称幺半范畴的一个
{\it 函子} $F : \mathcal{C} \to \mathcal{C}'$，由所示函子 $F$ 以及
关于 $X$、$Y$ 函子性的同构
$$
F(X) \otimes F(Y) \to F(X \otimes Y)
$$
给出，使得 $F$ 是幺半范畴的函子，并且对所有对象 $X$ 与 $Y$，图
$$
\xymatrix{
F(X) \otimes F(Y) \ar[r] \ar[d] &
F(X \otimes Y) \ar[d] \\
F(Y) \otimes F(X) \ar[r] &
F(Y \otimes X)
}
$$
交换。
\end{definition}

\begin{remark}
\label{categories-remark-internal-hom-monoidal}
令 $\mathcal{C}$ 为幺半范畴。若对 $\mathcal{C}$ 的每对对象 $X, Y$，
都有 $\mathcal{C}$ 的对象 $hom(X, Y)$，使得
$$
\Mor(X, hom(Y, Z)) = \Mor(X \otimes Y, Z)
$$
关于 $X, Y, Z$ 函子性地成立，则称 $\mathcal{C}$ 有
{\it 内部 hom}。由米田引理，如果双函子
$(X, Y) \mapsto hom(X, Y)$ 存在，则它在唯一同构意义下确定。
给定内部 hom，我们得到典范映射
\begin{enumerate}
\item $hom(X, Y) \otimes X \to Y$；
\item $hom(Y, Z) \otimes hom(X, Y) \to hom(X, Z)$；
\item $Z \otimes hom(X, Y) \to hom(X, Z \otimes Y)$；
\item $Y \to hom(X, Y \otimes X)$；
\item 若 $\mathcal{C}$ 是对称幺半范畴，则有
$hom(Y, Z) \otimes X \to hom(hom(X, Y), Z)$。
\end{enumerate}
具体而言，(1) 中的映射是 $\text{id}_{hom(X, Y)}$ 在映射
$\Mor(hom(X, Y), hom(X, Y)) \to \Mor(hom(X, Y) \otimes X, Y)$
下的像。为用 $hom(X, Z)$ 的定义性质构造 (2) 中的映射，需要构造映射
$$
hom(Y, Z) \otimes hom(X, Y) \otimes X \longrightarrow Z
$$
而这样的映射存在，因为由 (1) 有映射
$hom(X, Y) \otimes X \to Y$ 与 $hom(Y, Z) \otimes Y \to Z$。
为用 $hom(X, Z \otimes Y)$ 的定义性质构造 (3) 中的映射，需要构造映射
$$
Z \otimes hom(X, Y) \otimes X \to Z \otimes Y
$$
为此使用 $\text{id}_Z \otimes a$，其中 $a$ 是 (1) 中的映射。
为构造 (4) 中的映射，注意由 (3) 已有映射
$Y \otimes hom(X, X) \to hom(X, Y \otimes X)$。所以只需构造映射
$\mathbf{1} \to hom(X, X)$；为此取
$\Mor(\mathbf{1}, hom(X, X))$ 中与
$\Mor(\mathbf{1} \otimes X, X)$ 中典范同构
$\mathbf{1} \otimes X \to X$ 对应的元素。最后来看 (5)。
由 $hom(hom(X, Y), Z)$ 的泛性质，只需构造映射
$$
hom(Y, Z) \otimes X \otimes hom(X, Y) \longrightarrow Z
$$
我们用交换约束交换最后两个张量积，然后使用映射
$hom(X, Y) \otimes X \to Y$ 与 $hom(Y, Z) \otimes Y \to Z$。
\end{remark}





\section{点线箭头范畴}
\label{categories-section-dotted-arrows}

\noindent 
我们讨论 $(2,1)$-范畴中的某些“点线箭头范畴”。在表述代数栈的各种
提升判据时会用到它们，例如见栈的态射，第
\ref{stacks-morphisms-section-valuative} 节，以及栈的态射进阶，第
\ref{stacks-more-morphisms-section-formally-smooth} 节。

\begin{definition}
\label{categories-definition-dotted-arrows}
令 $\mathcal{C}$ 为 $(2,1)$-范畴。考虑 $\mathcal{C}$ 中的
$2$-交换实线图
\begin{equation}
\label{categories-equation-dotted-arrows}
\vcenter{
\xymatrix{
S \ar[r]_-x \ar[d]_j & X \ar[d]^f \\
T \ar[r]^-y \ar@{..>}[ru] & Y
}
}
\end{equation}
固定一个见证此图 $2$-交换的 $2$-同构
$$
\gamma : y \circ j \rightarrow f \circ x
$$
给定 (\ref{categories-equation-dotted-arrows}) 与 $\gamma$，一个\emph{点线箭头}
是三元组 $(a, \alpha, \beta)$，其中包括态射 $a \colon T \to X$
及 $2$-同构
$\alpha : a \circ j \to x$、$\beta : y \to f \circ a$，并且
$\gamma = (\text{id}_f \star \alpha) \circ (\beta \star \text{id}_j)$；
换言之，图
$$
\xymatrix{
& f \circ a \circ j \ar[rd]^{\text{id}_f \star \alpha} \\
y \circ j \ar[ru]^{\beta \star \text{id}_j} \ar[rr]^\gamma & &
f \circ x
}
$$
交换。点线箭头的一个{\it 态射}
$(a, \alpha, \beta) \to (a', \alpha', \beta')$ 是满足
$\alpha = \alpha' \circ (\theta \star \text{id}_j)$ 且
$\beta' = (\text{id}_f \star \theta) \circ \beta$ 的
$2$-箭头 $\theta : a \to a'$。
\end{definition}

\noindent
在定义 \ref{categories-definition-dotted-arrows} 的情形中，有一个相伴的
\emph{点线箭头范畴}。这个范畴是群胚。一般而言，它可能依赖于
$\gamma$。下面两个引理说明，点线箭头范畴对基变换以及 $f$ 的复合
表现良好。

\begin{lemma}
\label{categories-lemma-cat-dotted-arrows-base-change}
令 $\mathcal{C}$ 为 $(2,1)$-范畴。假设给定 $\mathcal{C}$ 中的
$2$-交换图
$$
\xymatrix{
S \ar[r]_-{x'} \ar[d]_j &
X' \ar[d]^p \ar[r]_q &
X \ar[d]^f \\
T \ar[r]^-{y'} &
Y' \ar[r]^g &
Y
}
$$
其中右方块关于 $2$-同构
$\phi \colon g \circ p \to f \circ q$ 是 $2$-笛卡尔的。
选择 $2$-箭头
$\gamma' : y' \circ j \to p \circ x'$。令
$x = q \circ x'$、$y = g \circ y'$，并令
$\gamma : y \circ j \to f \circ x$ 为 $2$-同构
$\gamma = (\phi \star \text{id}_{x'}) \circ (\text{id}_g \star \gamma')$。
那么，左方块及 $\gamma'$ 的点线箭头范畴 $\mathcal{D}'$，等价于
外侧矩形及 $\gamma$ 的点线箭头范畴 $\mathcal{D}$。
\end{lemma}

\begin{proof}
有一个函子 $\mathcal{D}' \to \mathcal{D}$，它在对象上为
$(a, \alpha, \beta) \mapsto (q \circ a, \text{id}_q \star \alpha,
(\phi \star \text{id}_a) \circ (\text{id}_g \star \beta))$，
在箭头上为 $\theta \mapsto \text{id}_q \star \theta$。
由第 \ref{categories-section-2-fibre-products} 节所述 $2$-纤维积的泛性质，
形式地即可验证此函子 $\mathcal{D}' \to \mathcal{D}$ 是等价。
细节略去。
\end{proof}

\begin{lemma}
\label{categories-lemma-cat-dotted-arrows-composition}
令 $\mathcal{C}$ 为 $(2,1)$-范畴。假设给定 $\mathcal{C}$ 中的实线
$2$-交换图
$$
\xymatrix{
S \ar[r]_-x \ar[dd]_j & X \ar[d]^f \\
& Y \ar[d]^g \\
T \ar[r]^-z \ar@{..>}[ruu] & Z
}
$$
选择 $2$-同构 $\gamma \colon z \circ j \to g \circ f \circ x$。
令 $\mathcal{D}$ 为外侧矩形及 $\gamma$ 的点线箭头范畴。令
$\mathcal{D}'$ 为实线方块
$$
\xymatrix{
S \ar[r]_-{f \circ x} \ar[d]_j & Y \ar[d]^g \\
T \ar[r]^-z \ar@{..>}[ru] & Z
}
$$
及 $\gamma$ 的点线箭头范畴。那么 $\mathcal{D}$ 等价于一个范畴
$\mathcal{D}''$，后者具有下述性质：有函子
$\mathcal{D}'' \to \mathcal{D}'$，使 $\mathcal{D}''$ 成为
$\mathcal{D}'$ 上的群胚纤维化范畴；它的纤维范畴同构于形如
$$
\xymatrix{
S \ar[r]_-x \ar[d]_j & X \ar[d]^f \\
T \ar[r]^-y \ar@{..>}[ru] & Y
}
$$
的某些实线方块在选择了某个 $2$-同构
$y \circ j \to f \circ x$ 后所得的点线箭头范畴。
\end{lemma}

\begin{proof}
构造范畴 $\mathcal{D}''$，其对象是元组
$(a,\alpha,\beta,b,\eta)$，其中 $(a,\alpha,\beta)$ 是
$\mathcal{D}$ 的对象，$b \colon T \rightarrow Y$ 是 $1$-态射，
且 $\eta \colon b \rightarrow f \circ a$ 是 $2$-同构。
$\mathcal{D}''$ 中的态射
$(a,\alpha,\beta,b,\eta) \rightarrow (a',\alpha',\beta',b',\eta')$
是有序对 $(\theta_1,\theta_2)$，其中
$\theta_1 \colon a \rightarrow a'$ 在 $\mathcal{D}$ 中定义箭头
$(a, \alpha, \beta) \rightarrow (a', \alpha', \beta')$，且
$\theta_2 \colon b \rightarrow b'$ 是满足相容条件
$\eta' \circ \theta_2 = (\text{id}_f \star \theta_1) \circ \eta$
的 $2$-同构。

\medskip\noindent
有函子 $\mathcal{D}'' \rightarrow \mathcal{D}'$，它在对象上为
$(a, \alpha, \beta, b, \eta) \mapsto
(b, (\text{id}_f \star \alpha) \circ (\eta \star \text{id}_j),
(\text{id}_g \star \eta^{-1}) \circ \beta)$，在箭头上为
$(\theta_1,\theta_2) \mapsto \theta_2$。那么
$\mathcal{D}'' \rightarrow \mathcal{D}'$ 是群胚纤维化的。

\medskip\noindent
若 $(y, \delta, \epsilon)$ 是 $\mathcal{D}'$ 的对象，以
$\mathcal{D}_{y,\delta}$ 表示最后一个所示图在
$y \circ j \rightarrow f \circ x$ 取为 $\delta$ 时的点线箭头范畴。
有函子 $\mathcal{D}_{y,\delta} \rightarrow \mathcal{D}''$，它在对象上为
$(a, \alpha, \eta) \mapsto
(a, \alpha, (\text{id}_g \star \eta) \circ \epsilon, y, \eta)$，
在箭头上为 $\theta \mapsto (\theta, \text{id}_y)$。这给出从
$\mathcal{D}_{y,\delta}$ 到 $\mathcal{D}'' \rightarrow \mathcal{D}'$
在 $(y,\delta,\epsilon)$ 上的纤维范畴的同构。

\medskip\noindent
还有一个函子 $\mathcal{D} \rightarrow \mathcal{D}''$，它在对象上为
$(a,\alpha,\beta) \mapsto (a,\alpha,\beta,f \circ a, \text{id}_{f \circ a})$，
在箭头上为 $\theta \mapsto (\theta, \text{id}_f \star \theta)$。
此函子全忠实且本质满，因而是等价。细节略去。
\end{proof}












% P01 cursor: frozen categories.tex translated through source line 9855 (chapter complete).
% Generated mechanically from the frozen topology.tex target.
% Do not edit this generated file; edit the bound locale source instead.

% OK, start here.
%

\setcounter{chapter}{4}
\chapter{拓扑}



\phantomsection
\label{topology-section-phantom}


\section{引言}
\label{topology-section-introduction}

\noindent
本文将介绍基础拓扑学。可参见 \cite{Engelking}。

\section{基本概念}
\label{topology-section-topology-basic}

\noindent
下面列出拓扑学中的若干基本概念。其中一些将在后文作更详细的讨论，
另一些已在表中定义；其余则视为基础知识而不再定义。如果你对大多数斜体概念
尚不熟悉，我们建议先阅读一本拓扑学入门教材，再继续阅读本文。

\begin{enumerate}
\item
\label{topology-item-space}
$X$ 是一个{\it 拓扑空间}，
\item
\label{topology-item-point}
$x\in X$ 是一个{\it 点}，
\item
\label{topology-item-locally-closed}
$E \subset X$ 是一个{\it 局部闭}子集，
\item
\label{topology-item-closed-point}
$x\in X$ 是一个{\it 闭点}，
\item
\label{topology-item-dense}
$E \subset X$ 是一个{\it 稠密}子集，
\item
\label{topology-item-continuous}
$f : X_1 \to X_2$ 是{\it 连续的}，
\item
\label{topology-item-upper-semi-continuous}
扩展实值函数 $f : X \to \mathbf{R} \cup \{\infty, -\infty\}$ 称为
{\it 上半连续的}，如果对每个 $a \in \mathbf{R}$，
$\{x \in X \mid f(x) < a\}$ 都是开集，
\item
\label{topology-item-lower-semi-continuous}
扩展实值函数 $f : X \to \mathbf{R} \cup \{\infty, -\infty\}$ 称为
{\it 下半连续的}，如果对每个 $a \in \mathbf{R}$，
$\{x \in X \mid f(x) > a\}$ 都是开集，
\item 空间间的连续映射 $f : X \to Y$ 称为{\it 开映射}，如果对每个开集
$U \subset X$，$f(U)$ 在 $Y$ 中是开集，
\item 空间间的连续映射 $f : X \to Y$ 称为{\it 闭映射}，如果对每个闭集
$Z \subset X$，$f(Z)$ 在 $Y$ 中是闭集，
\item
\label{topology-item-neighbourhood}
点 $x \in X$ 的一个{\it 邻域}是任意子集 $E \subset X$，
它包含某个含有 $x$ 的开子集，
\item
\label{topology-item-induced-topology}
子集 $E \subset X$ 上的{\it 诱导拓扑}，
\item
\label{topology-item-covering}
$\mathcal{U} : U = \bigcup_{i \in I} U_i$ 是 $U$ 的一个
{\it 开覆盖}（注意：允许任意 $U_i$ 为空；当 $U$ 为空时，甚至允许 $I$ 为空集），
\item
\label{topology-item-subcover}
如 (\ref{topology-item-covering}) 中覆盖的一个{\it 子覆盖}，是满足 $I' \subset I$ 的
开覆盖 $\mathcal{U}' : U = \bigcup_{i \in I'} U_i$，
\item
\label{topology-item-refinement}
开覆盖 $\mathcal{V}$ 是开覆盖 $\mathcal{U}$ 的一个{\it 加细}
（若 $\mathcal{V} : U = \bigcup_{j \in J} V_j$ 且
$\mathcal{U} : U = \bigcup_{i \in I} U_i$，这意味着每个 $V_j$
都完全包含于某个 $U_i$ 中），
\item
\label{topology-item-fundamental-system}
{\it $\{ E_i \}_{i \in I}$ 是 $X$ 中点 $x$ 的一个邻域基本系}，
\item
\label{topology-item-Hausdorff}
拓扑空间 $X$ 称为{\it Hausdorff 空间}或{\it 分离空间}，当且仅当对任意两个
不同的点 $x, y \in X$，存在不交的开集 $U, V \subset X$，使得
$x \in U$、$y \in V$，
\item 两个拓扑空间的{\it 积}，
\label{topology-item-product}
\item
\label{topology-item-fibre-product}
一对连续映射 $f : X \to Y$ 和 $g : Z \to Y$ 的
{\it 纤维积 $X \times_Y Z$}，
\item
\label{topology-item-discrete-indiscrete}
集合上的{\it 离散拓扑}与{\it 平凡拓扑}，
\item 等等。
\end{enumerate}



\section{Hausdorff 空间}
\label{topology-section-Hausdorff}

\noindent
拓扑空间范畴具有有限积。

\begin{lemma}
\label{topology-lemma-Hausdorff}
设 $X$ 为拓扑空间。下列条件等价：
\begin{enumerate}
\item $X$ 是 Hausdorff 空间，
\item 对角线 $\Delta(X) \subset X \times X$ 是闭集。
\end{enumerate}
\end{lemma}

\begin{proof}
假设 $X$ 是 Hausdorff 空间。设 $(x, y) \not \in \Delta(X)$，即
$x \neq y$。存在 $X$ 中不交的开集 $U$ 与 $V$，使得
$x \in U$ 且 $y \in V$。于是 $U \times V \subset
(X \times X) \setminus \Delta(X) $。这说明 $ (X \times X) \setminus
\Delta(X)$ 是 $ X \times X$ 的开集，也就等价于说对角线
$\Delta(X) \subset X \times X$ 在 $X \times X$ 中是闭集。
反之也类似：补集 $(X \times X) \setminus \Delta(X)$ 恰由满足
$ x \neq y$ 的 $(x, y) \in X\times X$ 组成。因此，如果
$ \Delta(X) \subset X \times X$ 是闭集，那么由积拓扑的定义，
对每个这样的 $ (x, y)$，存在开集 $U, V \subset X$，满足
$(x, y) \in U \times V$ 且 $(U \times V)\cap \Delta(X) = \emptyset$。
换言之，$x \in U$、$y \in V$ 且 $U \cap V = \emptyset$。
\end{proof}

\begin{lemma}
\label{topology-lemma-graph-closed}
\begin{slogan}
到 Hausdorff 空间的映射之图是闭的。
\end{slogan}
设 $f : X \to Y$ 为拓扑空间间的连续映射。
若 $Y$ 是 Hausdorff 空间，则 $f$ 的图在 $X \times Y$ 中是闭的。
\end{lemma}

\begin{proof}
该图是映射 $X \times Y \to Y \times Y$ 下对角线的逆像。
因此本引理由引理 \ref{topology-lemma-Hausdorff} 得出。
\end{proof}

\begin{lemma}
\label{topology-lemma-section-closed}
设 $f : X \to Y$ 为拓扑空间间的连续映射。设 $s : Y \to X$
为满足 $f \circ s = \text{id}_Y$ 的连续映射。若 $X$ 是 Hausdorff 空间，
则 $s(Y)$ 是闭集。
\end{lemma}

\begin{proof}
由引理 \ref{topology-lemma-Hausdorff} 以及
$s(Y) = \{x \in X \mid x = s(f(x))\}$ 可得。
\end{proof}

\begin{lemma}
\label{topology-lemma-fibre-product-closed}
设 $X \to Z$ 和 $Y \to Z$ 为拓扑空间间的连续映射。
若 $Z$ 是 Hausdorff 空间，则 $X \times_Z Y$ 在 $X \times Y$ 中是闭的。
\end{lemma}

\begin{proof}
由引理 \ref{topology-lemma-Hausdorff} 可得，因为 $X \times_Z Y$ 是映射
$X \times Y \to Z \times Z$ 下 $\Delta(Z)$ 的逆像。
\end{proof}




\section{分离映射}
\label{topology-section-separated}

\noindent
本节只给出定义和若干简单引理。

\begin{definition}
\label{topology-definition-separated}
拓扑空间间的连续映射 $f : X \to Y$ 称为{\it 分离的}，当且仅当
对角映射 $\Delta : X \to X \times_Y X$ 是闭映射。
\end{definition}

\begin{lemma}
\label{topology-lemma-separated}
设 $f : X \to Y$ 为拓扑空间间的连续映射。下列条件等价：
\begin{enumerate}
\item $f$ 是分离的，
\item $\Delta(X) \subset X \times_Y X$ 是闭子集，
\item 对于映到 $Y$ 中同一点的任意不同两点 $x, x' \in X$，
存在分别为 $x$ 与 $x'$ 的不交开邻域。
\end{enumerate}
\end{lemma}

\begin{proof}
若 $f$ 分离，则由定义 \ref{topology-definition-separated}，$\Delta$ 是闭映射。
由于 $X$ 在 $X$ 中闭，$\Delta(X)$ 在 $ X\times_Y X $ 中闭。
因此 (1) 蕴含 (2)。

\medskip\noindent
假设 $\Delta(X) \subset X \times_Y X$ 是闭子集，并令 $U$ 表示其开补集。
这意味着存在开集 $W \subset X \times X$，使得
$W \cap (X \times_Y X) = U$。另一方面，由积拓扑的定义，若
$(x, x') \in W \cap (X \times_Y X)$，则存在 $X$ 的开集 $V$ 与 $V'$，
使得 $x \in V$、$x' \in V'$ 且 $V \times V' \subset W$。
若 $V \cap V' \neq \emptyset$，则存在 $z \in V \cap V'$。
然而 $(z, z) \in X \times_Y X$，所以
$(z,z) \in (V \times V') \cap (X\times_Y X) \subset U$，矛盾。
因此 $V \cap V' = \emptyset $，从而得到 $x$ 与 $x'$ 的两个不交开邻域。
这证明 (2) 蕴含 (3)。

\medskip\noindent
最后，假设对任意映到 $Y$ 中同一点的不同两点 $x, x' \in X$，
都存在 $x$ 的一个开邻域与 $x'$ 的一个开邻域，且二者不交。
设 $F$ 为 $X$ 的闭集。我们证明
$\Delta(F)$ 是 $X \times_Y X$ 的闭子集。取
$(x, x') \in X \times_Y X$ 且不属于 $\Delta(F)$。
于是或者 $x \not = x'$，或者 $x \not \in F$。
在第一种情形，选取 $x, x'$ 的不交开邻域 $V, V' \subset X$，于是
$(V \times V') \cap X \times_Y X$ 是 $(x, x')$ 的一个开邻域，
且不与 $\Delta(F)$ 相交。在第二种情形，
$((X \setminus F) \times X) \cap X \times_Y X$ 是 $(x, x')$ 的一个
开邻域，且不与 $\Delta(F)$ 相交。由此证明 (3) 蕴含 (1)。
\end{proof}

\begin{lemma}
\label{topology-lemma-from-hausdorff}
设 $f : X \to Y$ 为拓扑空间间的连续映射。若 $X$ 是 Hausdorff 空间，
则 $f$ 是分离的。
\end{lemma}

\begin{proof}
由引理 \ref{topology-lemma-separated} 与引理 \ref{topology-lemma-Hausdorff} 立即可得，
因为 $\Delta(X)$ 在 $X \times X$ 中闭蕴含 $\Delta(X)$ 在
$X \times_Y X$ 中闭。
\end{proof}

\begin{lemma}
\label{topology-lemma-base-change-separated}
设 $f : X \to Y$ 与 $Y' \to Y$ 为拓扑空间间的连续映射。
若 $f$ 分离，则 $f' : Y' \times_Y X \to Y'$ 分离。
\end{lemma}

\begin{proof}
由引理 \ref{topology-lemma-separated} 的刻画 (2) 可得。事实上，令
$X' = Y' \times_Y X$，则纤维积 $X' \times_{Y'} X'$ 中的对角线
$\Delta(X')$ 是 $X \times_Y X$ 中 $\Delta(X)$ 的逆像。
\end{proof}




\section{基}
\label{topology-section-bases}

\noindent
本节介绍拓扑空间之基的基本材料。

\begin{definition}
\label{topology-definition-base}
设 $X$ 为拓扑空间。$X$ 的一个子集族 $\mathcal{B}$ 称为
{\it $X$ 上拓扑的一个基}或{\it $X$ 上拓扑的一个基底}，如果满足：
\begin{enumerate}
\item 每个元素 $B \in \mathcal{B}$ 都在 $X$ 中开。
\item 对每个开集 $U \subset X$ 以及每个 $x \in U$，存在
$B \in \mathcal{B}$，使得 $x \in B \subset U$。
\end{enumerate}
\end{definition}

\noindent
下面的引理有时用来定义拓扑。

\begin{lemma}
\label{topology-lemma-make-base}
设 $X$ 为集合，$\mathcal{B}$ 为一个子集族。假设
$X = \bigcup_{B \in \mathcal{B}} B$，并且对任意
$x \in B_1 \cap B_2$，其中 $B_1, B_2 \in \mathcal{B}$，
都存在 $B_3 \in \mathcal{B}$，使得
$x \in B_3 \subset B_1 \cap B_2$。那么 $X$ 上存在唯一的拓扑，
使得 $\mathcal{B}$ 是该拓扑的一个基底。
\end{lemma}

\begin{proof}
令 $\sigma(\mathcal B)$ 表示 $X$ 的所有可写成 $\mathcal B$ 中元素之并的
子集所成的集合。我们断言 $\sigma(\mathcal{B})$ 是一个拓扑。
事实上，空集属于 $\sigma(\mathcal{B})$（作为空并），而 $X$ 属于
$\sigma(\mathcal{B})$（作为 $\mathcal{B}$ 所有元素之并）。显然
$\sigma(\mathcal{B})$ 对任意并封闭。最后，若
$U, V \in \sigma(\mathcal{B})$，则可写
$U = \bigcup_{i \in I} U_i$ 和 $V = \bigcup_{j \in J} V_j$，
其中对所有 $i \in I$ 与 $j \in J$，都有 $U_i, V_j \in \mathcal{B}$。
于是
$$
U \cap V = \bigcup\nolimits_{i \in I, j \in J}
U_i \cap V_j
$$
引理中的假设告诉我们 $U_i \cap V_j \in \sigma(\mathcal{B})$，
因而 $U \cap V$ 也属于其中。故 $\sigma(\mathcal{B})$ 是一个拓扑。
定义 \ref{topology-definition-base} 的性质 (1) 与 (2) 对此拓扑立即成立。
为证明该拓扑的唯一性，设 $\mathcal T$ 为 $X$ 上一个以
$\mathcal B$ 为基的拓扑。由定义 \ref{topology-definition-base} 的 (1)，
$\mathcal{B}$ 的每个元素当然都属于 $\mathcal{T}$，所以
$\sigma(\mathcal{B}) \subset \mathcal{T}$。反之，定义
\ref{topology-definition-base} 的 (2) 表明 $\mathcal{T}$ 的每个元素都是
$\mathcal{B}$ 中元素之并，即 $\mathcal{T} \subset \sigma(\mathcal{B})$。
证毕。
\end{proof}

\begin{lemma}
\label{topology-lemma-refine-covering-basis}
设 $X$ 为拓扑空间，$\mathcal{B}$ 为 $X$ 上拓扑的一个基底。
设 $\mathcal{U} : U = \bigcup_i U_i$ 为 $U \subset X$ 的一个开覆盖。
则存在开覆盖 $U = \bigcup V_j$，它是 $\mathcal{U}$ 的加细，
并且每个 $V_j$ 都是基底 $\mathcal{B}$ 的元素。
\end{lemma}

\begin{proof}
若 $ x \in U = \bigcup_{i\in I} U_i $，则存在 $ i_x \in I $，使得
$ x \in U_{i_x} $。因此存在 $ B_{i_x} \in \mathcal{B}$，满足
$ x \in B_{i_x} \subset U_{i_x}$。令
$J = \{i_x | x \in U\}$，并对 $j = i_x \in J$ 令 $V_j = B_{i_x}$。
这给出所需的 $U$ 的开覆盖 $\{V_j\}_{j \in J}$。
\end{proof}

\begin{definition}
\label{topology-definition-subbase}
设 $X$ 为拓扑空间。$X$ 的一个子集族 $\mathcal{B}$ 称为
{\it $X$ 上拓扑的一个子基}或{\it $X$ 上拓扑的一个子基底}，
如果 $\mathcal{B}$ 中元素的有限交构成 $X$ 上拓扑的一个基底。
\end{definition}

\noindent
特别地，$\mathcal{B}$ 的每个元素都是开集。

\begin{lemma}
\label{topology-lemma-subbase}
设 $X$ 为集合。给定 $X$ 的任意子集族 $\mathcal{B}$，$X$ 上存在唯一的拓扑，
使得 $\mathcal{B}$ 是该拓扑的一个子基。
\end{lemma}

\begin{proof}
按约定，$\bigcap_\emptyset B = X $。因此可将引理
\ref{topology-lemma-make-base} 应用于 $\mathcal B$ 中元素的有限交所成的集合。
\end{proof}

\begin{lemma}
\label{topology-lemma-create-map-from-subcollection}
设 $X$ 为拓扑空间，$\mathcal{B}$ 为 $X$ 的一个开集族。假设
$X = \bigcup_{U \in \mathcal{B}} U$，并且对 $U, V \in \mathcal{B}$ 有
$U \cap V = \bigcup_{W \in \mathcal{B}, W \subset U \cap V} W$。
那么存在拓扑空间间的连续映射 $f : X \to Y$，满足：
\begin{enumerate}
\item 对 $U \in \mathcal{B}$，像 $f(U)$ 是开集；
\item 对 $U \in \mathcal{B}$，有 $f^{-1}(f(U)) = U$；
\item 开集 $f(U)$（$U \in \mathcal{B}$）构成 $Y$ 上拓扑的一个基底。
\end{enumerate}
\end{lemma}

\begin{proof}
在 $X$ 的点上按下式定义等价关系 $\sim$：
$$
x \sim y \Leftrightarrow
(\forall U \in \mathcal{B} : x \in U \Leftrightarrow y \in U)
$$
令 $Y$ 为等价类的集合，$f : X \to Y$ 为自然映射。
按构造，(2) 成立。对 $\mathcal{B}$ 的假设恰好对应于引理
\ref{topology-lemma-make-base} 对子集 $f(U)$（$U \in \mathcal{B}$）之集合的假设。
因此 $Y$ 上存在唯一的拓扑使 (3) 成立，而 (1) 也随即显然成立。
\end{proof}

\section{浸没映射}
\label{topology-section-submersive}

\noindent
若 $X$ 是拓扑空间而 $E \subset X$ 是一个子集，我们通常在 $E$ 上赋予
{\it 诱导拓扑}。

\begin{lemma}
\label{topology-lemma-induced}
设 $X$ 为拓扑空间，$Y$ 为集合，并设 $f : Y \to X$ 为集合间的单射。
$Y$ 上的诱导拓扑可由下列任一陈述刻画：
\begin{enumerate}
\item 它是使 $f$ 连续的 $Y$ 上最弱拓扑，
\item $Y$ 的开子集恰为 $f^{-1}(U)$，其中 $U \subset X$ 为开集，
\item $Y$ 的闭子集恰为 $f^{-1}(Z)$，其中 $Z \subset X$ 为闭集。
\end{enumerate}
\end{lemma}

\begin{proof}
集合 $\mathcal{T} = \{f^{-1}(U) | U \subset X \text{ 开}\}$ 是 $Y$
上的一个拓扑。首先，$\emptyset = f^{-1}(\emptyset)$ 且 $f^{-1}(X) = Y$。
所以 $\mathcal{T}$ 包含 $\emptyset$ 与 $Y$。

\medskip\noindent
现设 $\{V_i\}_{i \in I}$ 为一族开子集，其中 $V_i \in \mathcal{T}$，
并写成 $V_i = f^{-1}(U_i)$，这里 $U_i$ 是 $X$ 的开子集。则
$$
\bigcup\nolimits_{i\in I} V_i =
\bigcup\nolimits_{i\in I} f^{-1}(U_i) =
f^{-1}\left(\bigcup\nolimits_{i\in I} U_i\right)
$$
由于 $\bigcup_{i\in I} U_i$ 在 $X$ 中开，故
$\bigcup_{i\in I} V_i \in \mathcal{T}$。
再设 $V_1, V_2 \in \mathcal{T}$。存在 $X$ 中的开集 $U_1, U_2$，
使得 $V_1 = f^{-1}(U_1)$ 且 $V_2 = f^{-1}(U_2)$。于是
$$
V_1 \cap V_2 = f^{-1}(U_1) \cap f^{-1}(U_2) = f^{-1}(U_1 \cap U_2)
$$
由于 $U_1 \cap U_2$ 在 $X$ 中开，故 $V_1 \cap V_2 \in \mathcal{T}$。

\medskip\noindent
根据连续映射的定义，$Y$ 上任何使 $f$ 连续的拓扑都包含
$\mathcal{T}$。所以 $\mathcal{T}$ 确实是使 $f$ 连续的 $Y$ 上最弱拓扑。
这证明 (1) 与 (2) 等价。

\medskip\noindent
对所有子集 $E \subset X$ 都有等式
$f^{-1}(X \setminus E) = Y \setminus f^{-1}(E)$，故 (2) 与 (3) 等价。
\end{proof}

\noindent
对偶地，若 $X$ 是拓扑空间且 $X \to Y$ 是集合间的满射，则可在 $Y$
上赋予{\it 商拓扑}。

\begin{lemma}
\label{topology-lemma-quotient}
设 $X$ 为拓扑空间，$Y$ 为集合，并设 $f : X \to Y$ 为集合间的满射。
$Y$ 上的商拓扑可由下列任一陈述刻画：
\begin{enumerate}
\item 它是使 $f$ 连续的 $Y$ 上最强拓扑，
\item $Y$ 的子集 $V$ 是开集，当且仅当 $f^{-1}(V)$ 是开集，
\item $Y$ 的子集 $Z$ 是闭集，当且仅当 $f^{-1}(Z)$ 是闭集。
\end{enumerate}
\end{lemma}

\begin{proof}
集合 $\mathcal{T} = \{V\subset Y | f^{-1}(V) \text{ 开}\}$ 是 $Y$
上的一个拓扑。首先，$\emptyset = f^{-1}(\emptyset)$ 且 $f^{-1}(Y) = X$。
所以 $\mathcal{T}$ 包含 $\emptyset$ 与 $Y$。

\medskip\noindent
设 $(V_i)_{i \in I}$ 为一族元素 $V_i \in \mathcal{T}$。则
$$
\bigcup\nolimits_{i\in I} f^{-1}(V_i) =
f^{-1}\left(\bigcup\nolimits_{i\in I} V_i\right)
$$
由于 $\bigcup_{i\in I}f^{-1}(V_i)$ 在 $X$ 中开，故
$\bigcup_{i\in I} V_i \in \mathcal{T}$。此外，若
$V_1, V_2 \in \mathcal{T}$，则
$$
f^{-1}(V_1) \cap f^{-1}(V_2) = f^{-1}(V_1 \cap V_2)
$$
由于 $f^{-1}(V_1) \cap f^{-1}(V_2)$ 在 $X$ 中开，故
$V_1 \cap V_2 \in \mathcal{T}$。

\medskip\noindent
最后，根据连续函数的定义，$Y$ 上任何使 $f$ 连续的拓扑都包含于
$\mathcal{T}$，所以 $\mathcal{T}$ 是使 $f$ 连续的 $Y$ 上最强拓扑。
这证明 (1) 与 (2) 等价。

\medskip\noindent
最后，对所有子集 $E \subset X$ 都有 $f^{-1}(X\setminus E) = Y
\setminus f^{-1}(E)$，故 (2) 与 (3) 等价。
\end{proof}

\noindent
设 $f : X \to Y$ 为拓扑空间间的连续映射。此时，作为集合间的映射，
我们得到分解 $X \to f(X) \to Y$。可以在 $f(X)$ 上赋予由满射
$X \to f(X)$ 给出的商拓扑，也可以赋予由单射 $f(X) \to Y$ 给出的
诱导拓扑。映射
$$
(f(X), \text{商拓扑})
\longrightarrow
(f(X), \text{诱导拓扑})
$$
是连续的。

\begin{definition}
\label{topology-definition-submersive}
设 $f : X \to Y$ 为拓扑空间间的连续映射。
\begin{enumerate}
\item 如果 $f(X)$ 上的诱导拓扑与商拓扑相同（见上面的讨论），
则称 $f$ 为一个{\it 拓扑空间的严格映射}。
\item 如果 $f$ 是满射且严格，则称 $f$ 是{\it 浸没的}
\footnote{这与微分流形之间浸没的概念大不相同！使用“严格且满射”来代替
“浸没”或许是一个好主意。}。
\end{enumerate}
\end{definition}

\noindent
因此，连续映射 $f : X \to Y$ 是浸没的，当且仅当 $f$ 是满射，
并且对任意 $T \subset Y$，$T$ 为开集或闭集当且仅当
$f^{-1}(T)$ 亦然。换言之，相对于满射 $X \to Y$，$Y$ 具有商拓扑。

\begin{lemma}
\label{topology-lemma-open-morphism-quotient-topology}
设 $f : X \to Y$ 为拓扑空间间满、开且连续的映射。
设 $T \subset Y$ 为子集。则
\begin{enumerate}
\item $f^{-1}(\overline{T}) = \overline{f^{-1}(T)}$，
\item $T \subset Y$ 是闭集，当且仅当 $f^{-1}(T)$ 是闭集，
\item $T \subset Y$ 是开集，当且仅当 $f^{-1}(T)$ 是开集，
\item $T \subset Y$ 是局部闭集，当且仅当 $f^{-1}(T)$ 是局部闭集。
\end{enumerate}
特别地，$f$ 是浸没的。
\end{lemma}

\begin{proof}
显然 $\overline{f^{-1}(T)} \subset f^{-1}(\overline{T})$。
若 $x \in X$ 且 $x \not \in \overline{f^{-1}(T)}$，则存在开邻域
$x \in U \subset X$，使得 $U \cap f^{-1}(T) = \emptyset$。
由于 $f$ 是开映射，$f(U)$ 是 $f(x)$ 的一个不与 $T$ 相交的开邻域。
所以 $x \not \in f^{-1}(\overline{T})$。这证明 (1)。
(2) 是 (1) 的简单推论。(3) 由 $f$ 开且满立即得出。
对于 (4)，若 $f^{-1}(T)$ 局部闭，则
$f^{-1}(T) \subset \overline{f^{-1}(T)} = f^{-1}(\overline{T})$
是开集。因此，将 (3) 应用于映射
$f^{-1}(\overline{T}) \to \overline{T}$，可知 $T$ 在
$\overline{T}$ 中开，即 $T$ 局部闭。
\end{proof}

\begin{lemma}
\label{topology-lemma-closed-morphism-quotient-topology}
设 $f : X \to Y$ 为拓扑空间间满、闭且连续的映射。
设 $T \subset Y$ 为子集。则
\begin{enumerate}
\item $\overline{T} = f(\overline{f^{-1}(T)})$，
\item $T \subset Y$ 是闭集，当且仅当 $f^{-1}(T)$ 是闭集，
\item $T \subset Y$ 是开集，当且仅当 $f^{-1}(T)$ 是开集，
\item $T \subset Y$ 是局部闭集，当且仅当
$f^{-1}(T)$ 是局部闭集。
\end{enumerate}
特别地，$f$ 是浸没的。
\end{lemma}

\begin{proof}
显然 $\overline{f^{-1}(T)} \subset f^{-1}(\overline{T})$。
于是 $T \subset f(\overline{f^{-1}(T)}) \subset \overline{T}$
是闭子集，从而得到 (1)。(2) 由 $f$ 是闭映射且满射立即可得。
(3) 由将 (2) 应用于 $T$ 的补集得到。
对于 (4)，若 $f^{-1}(T)$ 局部闭，则
$f^{-1}(T) \subset \overline{f^{-1}(T)}$ 是开集。
由 (1)，映射 $\overline{f^{-1}(T)} \to \overline{T}$ 是满射，
所以可将 (3) 应用于由 $f$ 诱导的映射
$\overline{f^{-1}(T)} \to \overline{T}$，从而得出 $T$ 在
$\overline{T}$ 中开，即 $T$ 局部闭。
\end{proof}

\section{连通分支}
\label{topology-section-connected-components}

\begin{definition}
\label{topology-definition-connected-components}
设 $X$ 为拓扑空间。
\begin{enumerate}
\item 如果 $X$ 非空，并且每当 $X = T_1 \amalg T_2$，其中
$T_i \subset X$ 既开又闭时，必有 $T_1 = \emptyset$ 或
$T_2 = \emptyset$，则称 $X$ 是{\it 连通的}。
\item 如果 $T \subset X$ 是 $X$ 的极大连通子集，则称 $T$ 是
$X$ 的一个{\it 连通分支}。
\end{enumerate}
\end{definition}

\noindent
空空间不是连通的。

\begin{lemma}
\label{topology-lemma-image-connected-space}
设 $f : X \to Y$ 为拓扑空间间的连续映射。若 $E \subset X$ 是连通子集，
则 $f(E) \subset Y$ 也是连通的。
\end{lemma}

\begin{proof}
设 $A \subset f(E)$ 为 $f(E)$ 的一个既开又闭的子集。由于 $f$ 连续，
$f^{-1}(A)$ 是 $E$ 的一个既开又闭的子集。由于 $E$ 连通，
$f^{-1}(A) = \emptyset$ 或 $f^{-1}(A) = E$。另一方面，$A\subset f(E)$
蕴含 $A = f(f^{-1}(A))$。事实上，若 $x\in f(f^{-1}(A))$，则存在
$y\in f^{-1}(A)$ 使 $f(y) = x$；又因为 $y\in f^{-1}(A)$，有
$f(y) \in A$，即 $x\in A$。反之，若 $x\in A$，由 $A\subset f(E)$，
存在 $y\in E$ 使 $f(y) = x$。因此 $y\in f^{-1}(A)$，从而
$x\in f(f^{-1}(A))$。于是 $A = \emptyset$ 或 $A = f(E)$，这证明
$f(E)$ 连通。
\end{proof}

\begin{lemma}
\label{topology-lemma-connected-components}
设 $X$ 为拓扑空间。
\begin{enumerate}
\item 若 $T \subset X$ 连通，则其闭包也连通。
\item $X$ 的任意连通分支都是闭的（但未必是开的）。
\item $X$ 的每个连通子集都包含于 $X$ 的唯一一个连通分支中。
\item $X$ 的每个点都包含于唯一一个连通分支中；换言之，$X$ 是其各连通分支
的不交并。
\end{enumerate}
\end{lemma}

\begin{proof}
设 $\overline{T}$ 为连通子集 $T$ 的闭包。假设
$\overline{T} = T_1 \amalg T_2$，其中 $T_i \subset \overline{T}$
既开又闭。则 $T = (T\cap T_1) \amalg (T \cap T_2)$。
所以 $T$ 等于二者之一，比如 $T = T_1 \cap T$。于是
$\overline{T} \subset T_1$。这蕴含 (1) 与 (2)。

\medskip\noindent
设 $A$ 为 $X$ 的一族非空连通子集，且
$\Omega = \bigcap_{T \in A} T$ 非空。我们断言
$E = \bigcup_{T \in A} T$ 连通。事实上，$E$ 包含 $\Omega$，故非空。
设 $E = E_1 \amalg E_2$，其中 $E_i$ 在 $E$ 中闭。重新编号后，可设
$E_1$ 与 $\Omega$ 相交。于是每个 $T \in A$ 都与 $E_1$ 相交；由于
$T$ 连通，它必包含于 $E_1$。故 $E \subset E_1$，断言得证。

\medskip\noindent
设 $W \subset X$ 为非空连通子集。将上一段的结果应用于 $X$ 中所有包含
$W$ 的连通子集所成的集合，便可看出 $E$ 是 $X$ 的一个连通分支。
这证明 (3) 中的存在性与唯一性。

\medskip\noindent
设 $x \in X$。在上一段中取 $W = \{x\}$，可知 $x$ 包含于 $X$ 的
唯一一个连通分支中。任意两个不同的连通分支必不交（由第二段的结果）。

\medskip\noindent
要得到连通分支不为开集的例子，只需取带有积拓扑的无穷积
$\prod_{n \in \mathbf{N}} \{0, 1\}$。它的连通分支都是单点集，而这些
单点集不是开集。
\end{proof}

\begin{remark}
\label{topology-remark-quasi-components}
\begin{reference}
\cite[Example 6.1.24]{Engelking}
\end{reference}
设 $X$ 为拓扑空间，$x \in X$。设 $Z \subset X$ 为 $X$ 中经过 $x$ 的
连通分支。考虑 $X$ 中所有包含 $x$ 的既开又闭子集之交 $E$。
显然 $Z \subset E$。一般而言 $Z \not = E$。例如，令
$X = \{x, y, z_1, z_2, \ldots\}$，并赋予以下开集族为基的拓扑：
对所有 $n$，取 $\{z_n\}$、$\{x, z_n, z_{n + 1}, \ldots\}$ 以及
$\{y, z_n, z_{n + 1}, \ldots\}$。则 $Z = \{x\}$ 且
$E = \{x, y\}$。细节从略。
\end{remark}

\begin{lemma}
\label{topology-lemma-connected-fibres-quotient-topology-connected-components}
设 $f : X \to Y$ 为拓扑空间间的连续映射。假设
\begin{enumerate}
\item $f$ 的所有纤维都连通，并且
\item 子集 $T \subset Y$ 是闭的，当且仅当 $f^{-1}(T)$ 是闭的。
\end{enumerate}
则 $f$ 诱导 $X$ 与 $Y$ 的连通分支集合之间的双射。
\end{lemma}

\begin{proof}
设 $T \subset Y$ 为一个连通分支。注意 $T$ 是闭的，参见引理
\ref{topology-lemma-connected-components}。只要证明 $f^{-1}(T)$ 连通，本引理即得证，
因为由引理 \ref{topology-lemma-image-connected-space}，$X$ 的任何连通子集都映入
$Y$ 的某个连通分支。假设 $f^{-1}(T) = Z_1 \amalg Z_2$，其中
$Z_1$、$Z_2$ 闭。对任意 $t \in T$，有
$f^{-1}(\{t\}) = Z_1 \cap f^{-1}(\{t\}) \amalg Z_2 \cap f^{-1}(\{t\})$。
由 (1)，$f^{-1}(\{t\})$ 连通，因而或者
$f^{-1}(\{t\}) \subset Z_1$，或者 $f^{-1}(\{t\}) \subset Z_2$。
换言之，$T = T_1 \amalg T_2$，其中 $f^{-1}(T_i) = Z_i$。
由 (2)，$T_i$ 在 $Y$ 中闭。因此按所需，$T_1 = \emptyset$ 或
$T_2 = \emptyset$。
\end{proof}

\begin{lemma}
\label{topology-lemma-connected-fibres-connected-components}
设 $f : X \to Y$ 为拓扑空间间的连续映射。假设
(a) $f$ 是开映射，
(b) $f$ 的所有纤维都连通。
则 $f$ 诱导 $X$ 与 $Y$ 的连通分支集合之间的双射。
\end{lemma}

\begin{proof}
这是引理
\ref{topology-lemma-connected-fibres-quotient-topology-connected-components} 的特殊情形。
\end{proof}

\begin{lemma}
\label{topology-lemma-finite-fibre-connected-components}
设 $f : X \to Y$ 为非空拓扑空间间的连续映射。假设
(a) $Y$ 连通，
(b) $f$ 是开映射且为闭映射，并且
(c) 存在点 $y\in Y$，使得纤维 $f^{-1}(y)$ 为有限集。
则 $X$ 至多有 $|f^{-1}(y)|$ 个连通分支。因此，$X$ 的任一连通分支
$T$ 都既开又闭，而 $f(T)$ 是 $Y$ 的非空且既开又闭的子集，故等于 $Y$。
\end{lemma}

\begin{proof}
若拓扑空间 $X$ 至少有 $N$ 个连通分支，其中 $N \in \mathbf{N}$，
则由归纳法可找到分解 $X = X_1 \amalg \ldots \amalg X_N$，将 $X$
写成 $N$ 个非空且既开又闭的子集 $X_1, \ldots , X_N$ 的不交并。
由于 $f$ 开且闭，每个 $f(X_i)$ 都是 $Y$ 的非空且既开又闭的子集，
从而等于 $Y$。特别地，对每个 $1 \leq i \leq N$，交集
$X_i \cap f^{-1}(y)$ 非空。因此 $f^{-1}(y)$ 至少有 $N$ 个元素。
\end{proof}

\begin{definition}
\label{topology-definition-totally-disconnected}
若拓扑空间的每个连通分支都是单点集，则称该空间{\it 全不连通}。
\end{definition}

\noindent
离散空间是全不连通的。全不连通空间未必离散；例如
$\mathbf{Q} \subset \mathbf{R}$ 全不连通但不离散。

\begin{lemma}
\label{topology-lemma-space-connected-components}
设 $X$ 为拓扑空间。令 $\pi_0(X)$ 为 $X$ 的连通分支集合。
令 $X \to \pi_0(X)$ 为将 $x \in X$ 映到 $X$ 中经过 $x$ 的连通分支
的映射。在 $\pi_0(X)$ 上赋予商拓扑。则 $\pi_0(X)$ 是全不连通空间，
并且从 $X$ 到任意全不连通空间 $Y$ 的连续映射 $X \to Y$ 都通过
$\pi_0(X)$ 分解。
\end{lemma}

\begin{proof}
由引理
\ref{topology-lemma-connected-fibres-quotient-topology-connected-components}，
$\pi_0(X)$ 的连通分支都是单点集。第二个陈述的证明从略。
\end{proof}

\begin{definition}
\label{topology-definition-locally-connected}
若拓扑空间 $X$ 的每个点 $x \in X$ 都有一个由连通邻域组成的邻域基本系，
则称 $X$ {\it 局部连通}。
\end{definition}

\begin{lemma}
\label{topology-lemma-locally-connected}
设 $X$ 为拓扑空间。若 $X$ 局部连通，则
\begin{enumerate}
\item $X$ 的任意开子集都局部连通，并且
\item $X$ 的连通分支都是开集。
\end{enumerate}
所以 $X$ 的开子集的连通分支也都是开集。特别地，每个点都有一个由开连通
邻域组成的邻域基本系。
\end{lemma}

\begin{proof}
对所有 $x\in X$，记 $\mathcal{N}(x)$ 为 $x$ 的连通邻域基本系，并设
$U\subset X$ 为 $X$ 的开子集。则对所有 $x\in U$，$U$ 是 $x$ 的邻域，
所以集合 $\{V \in \mathcal{N}(x) | V\subset U\}$ 非空，且为 $U$ 中
$x$ 的连通邻域基本系。因此 $U$ 局部连通，这证明 (1)。

\medskip\noindent
设 $x \in \mathcal{C} \subset X$，其中 $\mathcal{C}$ 为 $x$ 的连通分支。
由于 $X$ 局部连通，存在 $x$ 的一个连通邻域 $\mathcal{N}$。
因此由连通分支的定义，$\mathcal{N} \subset \mathcal{C}$，从而
$\mathcal{C}$ 是 $x$ 的一个邻域。这说明 $\mathcal{C}$ 是其每一点的邻域；
换言之，$\mathcal{C}$ 是开集，(2) 得证。
\end{proof}

\section{不可约分支}
\label{topology-section-irreducible-components}

\begin{definition}
\label{topology-definition-irreducible-components}
设 $X$ 为拓扑空间。
\begin{enumerate}
\item 如果 $X$ 非空，并且每当 $X = Z_1 \cup Z_2$，其中 $Z_i$ 为闭集时，
都有 $X = Z_1$ 或 $X = Z_2$，则称 $X$ 是{\it 不可约的}。
\item 如果 $Z \subset X$ 是 $X$ 的极大不可约子集，则称 $Z$ 是 $X$
的一个{\it 不可约分支}。
\end{enumerate}
\end{definition}

\noindent
不可约空间显然是连通的。

\begin{lemma}
\label{topology-lemma-image-irreducible-space}
设 $f : X \to Y$ 为拓扑空间间的连续映射。若 $E \subset X$ 是不可约子集，
则 $f(E) \subset Y$ 也不可约。
\end{lemma}

\begin{proof}
显然可以假设 $E = X$（即 $X$ 不可约）且 $f(E) = Y$（即 $f$ 满）。
首先，因为 $X \not = \emptyset$，所以 $Y \not = \emptyset$。
其次，假设 $Y = Y_1 \cup Y_2$，其中 $Y_1$、$Y_2$ 闭。
则 $X = X_1 \cup X_2$，其中 $X_i = f^{-1}(Y_i)$ 在 $X$ 中闭。
由对 $X$ 的假设，必有 $X = X_1$ 或 $X = X_2$；又因 $f$ 满，故
$Y = Y_1$ 或 $Y = Y_2$。
\end{proof}

\begin{lemma}
\label{topology-lemma-irreducible}
设 $X$ 为拓扑空间。
\begin{enumerate}
\item 若 $T \subset X$ 不可约，则它在 $X$ 中的闭包也不可约。
\item $X$ 的任一不可约分支都是闭的。
\item $X$ 的任一不可约子集都包含于 $X$ 的某个不可约分支中。
\item $X$ 的每个点都包含于 $X$ 的某个不可约分支中；换言之，$X$ 是其
不可约分支之并。
\end{enumerate}
\end{lemma}

\begin{proof}
设 $\overline{T}$ 为不可约子集 $T$ 的闭包。若
$\overline{T} = Z_1 \cup Z_2$，其中 $Z_i \subset \overline{T}$ 闭，
则 $T = (T\cap Z_1) \cup (T \cap Z_2)$，从而 $T$ 等于二者之一，比如
$T = Z_1 \cap T$。于是显然 $\overline{T} \subset Z_1$。这证明 (1)。
(2) 由 (1) 和不可约分支的定义立即得出。

\medskip\noindent
设 $T \subset X$ 不可约。考虑所有满足 $T \subset T' \subset X$ 的
不可约子集所成的集合 $A$。因为 $T \in A$，所以 $A$ 非空。
包含关系给出 $A$ 上的偏序：
$T' \leq T'' \Leftrightarrow T' \subset T''$。
选取极大全序子集 $A' \subset A$，并令
$E = \bigcup_{T' \in A'} T'$。我们断言 $E$ 不可约。事实上，假设
$E =  Z_1 \cup Z_2$ 是 $E$ 的两个闭子集之并。对每个 $T' \in A'$，
由 $T'$ 的不可约性，或者 $T' \subset Z_1$，或者 $T' \subset Z_2$。
假设对某个 $T'_0 \in A'$ 有 $T'_0 \not\subset Z_1$
（否则结论已经成立）。由于 $A'$ 是全序的，立即可见对所有
$T' \in A'$ 都有 $T' \subset Z_2$。故 $E = Z_2$。这证明 (3)。
因为单点空间不可约，(4) 是 (3) 的直接推论。
\end{proof}

\begin{lemma}
\label{topology-lemma-pick-irreducible-components}
设 $X$ 为拓扑空间，并假设 $X = \bigcup_{i = 1, \ldots, n} X_i$，
其中每个 $X_i$ 都是 $X$ 的不可约闭子集，并且没有任何 $X_i$ 包含于
其余各项之并。则每个 $X_i$ 都是 $X$ 的不可约分支，而 $X$ 的每个
不可约分支都是某个 $X_i$。
\end{lemma}

\begin{proof}
设 $Y$ 为 $X$ 的一个不可约分支。写成
$Y = \bigcup_{i = 1, \ldots, n} (Y \cap X_i)$。由于 $X_i$ 在 $X$ 中闭，
每个 $Y \cap X_i$ 在 $Y$ 中闭。由 $Y$ 的不可约性，$Y$ 等于某个
$Y \cap X_i$，即 $Y \subset X_i$。由不可约分支的极大性，得到
$Y = X_i$。

\medskip\noindent
反之，取任一 $X_i$。由引理 \ref{topology-lemma-irreducible}，它包含于 $X$ 的某个
不可约分支 $Y$ 中。由上一段，某个 $j$ 满足 $Y = X_j$。所以
$X_i \subset X_j$；由于原来的并中没有冗余项，$X_i = X_j$，故它是
一个不可约分支。
\end{proof}

\begin{lemma}
\label{topology-lemma-surjective-continuous-irreducible-components}
设 $f : X \to Y$ 为拓扑空间间满且连续的映射。若 $X$ 有有限个不可约分支，
比如有 $n$ 个，则 $Y$ 有 $\leq n$ 个不可约分支。
\end{lemma}

\begin{proof}
设 $X_1, \ldots, X_n$ 为 $X$ 的不可约分支。由引理
\ref{topology-lemma-image-irreducible-space} 与 \ref{topology-lemma-irreducible}，
$f(X_i)$ 的闭包 $Y_i \subset Y$ 不可约。因为 $f$ 满，
$Y$ 是各 $Y_i$ 之并。可以选取极小子集
$I \subset \{1, \ldots, n\}$，使得 $Y = \bigcup_{i \in I} Y_i$。
于是可应用引理 \ref{topology-lemma-pick-irreducible-components}，看出对
$i \in I$，$Y_i$ 恰为 $Y$ 的不可约分支。
\end{proof}

\noindent
单点集不可约。因此，若 $x \in X$ 是一点，则闭包
$\overline{\{x\}}$ 是 $X$ 的不可约闭子集。

\begin{definition}
\label{topology-definition-generic-point}
设 $X$ 为拓扑空间。
\begin{enumerate}
\item 设 $Z \subset X$ 为不可约闭子集。如果点 $\xi \in Z$ 满足
$Z = \overline{\{\xi\}}$，则称它为 $Z$ 的一个{\it 泛点}。
\item 如果对每个 $x, x' \in X$，只要 $x \not = x'$，就存在 $X$
的一个闭子集恰好包含这两点之一，则称空间 $X$ 为{\it Kolmogorov 空间}。
\item 如果每个不可约闭子集都有泛点，则称空间 $X$ {\it 拟清醒}。
\item 如果每个不可约闭子集都有唯一泛点，则称空间 $X$ {\it 清醒}。
\end{enumerate}
\end{definition}

\noindent
拓扑空间 $X$ 为 Kolmogorov、拟清醒或清醒，当且仅当从 $X$ 到 $X$
的不可约闭子集之集合的映射 $x\mapsto\overline{\{x\}}$ 分别为单射、
满射或双射。因此，一个拓扑空间清醒，当且仅当它拟清醒且为 Kolmogorov 空间。

\begin{lemma}
\label{topology-lemma-sober-subspace}
设 $X$ 为拓扑空间，且 $Y\subset X$。
\begin{enumerate}
\item 若 $X$ 为 Kolmogorov 空间，则 $Y$ 也是。
\item 假设 $Y$ 在 $X$ 中局部闭。若 $X$ 拟清醒，则 $Y$ 也拟清醒。
\item 假设 $Y$ 在 $X$ 中局部闭。若 $X$ 清醒，则 $Y$ 也清醒。
\end{enumerate}
\end{lemma}

\begin{proof}
(1) 的证明。假设 $X$ 为 Kolmogorov 空间。设 $x,y\in Y$ 且
$x\neq y$。则
$\overline{\overline{\{x\}}\cap Y}=\overline{\{x\}}\neq\overline{\{y\}}=
\overline{\overline{\{y\}}\cap Y}$。故
$\overline{\{x\}}\cap Y\neq\overline{\{y\}}\cap Y$。这说明 $Y$
是 Kolmogorov 空间。

\medskip\noindent
(2) 的证明。假设 $X$ 拟清醒。只需分别考虑 $Y$ 闭与 $Y$ 开的情形。
先假设 $Y$ 闭。设 $Z$ 为 $Y$ 的不可约闭子集。则 $Z$ 是 $X$ 的
不可约闭子集。因此存在 $x \in Z$，使得 $\overline{\{x\}} = Z$。
于是 $\overline{\{x\}} \cap Y = Z$。这说明 $Y$ 拟清醒。
再假设 $Y$ 开。设 $Z$ 为 $Y$ 的不可约闭子集。则 $\overline{Z}$ 是
$X$ 的不可约闭子集。因此存在 $x \in \overline{Z}$，使得
$\overline{\{x\}}=\overline{Z}$。若 $x\notin Y$，便得到矛盾
$Z=Z\cap Y\subset\overline{Z}\cap Y=\overline{\{x\}}\cap Y=\emptyset$。
所以 $x\in Y$。于是 $Z=\overline{Z}\cap Y=\overline{\{x\}}\cap Y$。
这说明 $Y$ 拟清醒。

\medskip\noindent
(3) 由 (1) 与 (2) 立即得出。
\end{proof}

\begin{lemma}
\label{topology-lemma-sober-local}
设 $X$ 为拓扑空间，$(X_i)_{i\in I}$ 为 $X$ 的一个覆盖。
\begin{enumerate}
\item 假设对每个 $i\in I$，$X_i$ 在 $X$ 中局部闭。则 $X$ 为
Kolmogorov 空间，当且仅当对每个 $i\in I$，$X_i$ 都是 Kolmogorov 空间。
\item 假设对每个 $i\in I$，$X_i$ 在 $X$ 中开。则 $X$ 拟清醒，
当且仅当对每个 $i\in I$，$X_i$ 都拟清醒。
\item 假设对每个 $i\in I$，$X_i$ 在 $X$ 中开。则 $X$ 清醒，
当且仅当对每个 $i\in I$，$X_i$ 都清醒。
\end{enumerate}
\end{lemma}

\begin{proof}
(1) 的证明。若 $X$ 为 Kolmogorov 空间，则由引理
\ref{topology-lemma-sober-subspace}，对每个 $i\in I$，$X_i$ 也是。
反之，假设对每个 $i\in I$，$X_i$ 都是 Kolmogorov 空间。
设 $x,y\in X$ 且 $\overline{\{x\}}=\overline{\{y\}}$。存在
$i\in I$ 使 $x\in X_i$。存在开子集 $U\subset X$，使得 $X_i$ 是
$U$ 的闭子集。若 $y\notin U$，则得到矛盾
$x\in\overline{\{x\}}\cap U=\overline{\{y\}}\cap U=\emptyset$。
故 $y\in U$。于是
$y\in\overline{\{y\}}\cap U=\overline{\{x\}}\cap U\subset X_i$，
这说明 $y\in X_i$。从而
$\overline{\{x\}}\cap X_i=\overline{\{y\}}\cap X_i$。
由于 $X_i$ 为 Kolmogorov 空间，得到 $x=y$。这说明 $X$ 为
Kolmogorov 空间。

\medskip\noindent
(2) 的证明。若 $X$ 拟清醒，则由引理 \ref{topology-lemma-sober-subspace}，
对每个 $i\in I$，$X_i$ 也拟清醒。反之，假设对每个 $i\in I$，
$X_i$ 都拟清醒。设 $Y$ 为 $X$ 的不可约闭子集。因为
$Y\neq\emptyset$，存在 $i\in I$ 使 $X_i\cap Y\neq\emptyset$。
由于 $X_i$ 在 $X$ 中开，$X_i\cap Y$ 在 $Y$ 中非空且开，因而不可约
并在 $Y$ 中稠密。于是 $X_i\cap Y$ 是 $X_i$ 的不可约闭子集。
由于 $X_i$ 拟清醒，存在 $x\in X_i\cap Y$，使得
$X_i\cap Y=\overline{\{x\}}\cap X_i\subset\overline{\{x\}}$。
因为 $X_i\cap Y$ 在 $Y$ 中稠密且 $Y$ 在 $X$ 中闭，所以
$Y=\overline{X_i\cap Y}\cap Y\subset\overline{X_i\cap Y}\subset
\overline{\{x\}}\subset Y$。因此 $Y=\overline{\{x\}}$。
这说明 $X$ 拟清醒。

\medskip\noindent
(3) 由 (1) 与 (2) 立即得出。
\end{proof}

\begin{example}
\label{topology-example-quasi-sober-not-kolmogorov}
设 $X$ 为基数至少为 $2$ 的平凡空间。则 $X$ 拟清醒但不是 Kolmogorov
空间。此外，它的单点集族是由离散空间、因而由 Kolmogorov 空间组成的
$X$ 的覆盖。
\end{example}

\begin{example}
\label{topology-example-kolmogorov-not-quasi-sober}
设 $Y$ 为无穷集，并赋予以 $Y$ 及 $Y$ 的有限子集为闭集的拓扑。
则 $Y$ 是 Kolmogorov 空间但不拟清醒。然而，它的单点集族是一个由离散空间、
因而由清醒空间组成的覆盖。
\end{example}

\begin{example}
\label{topology-example-not-kolmogorov-not-quasi-sober}
令 $X$ 与 $Y$ 如例 \ref{topology-example-quasi-sober-not-kolmogorov} 和例
\ref{topology-example-kolmogorov-not-quasi-sober}。则 $X\amalg Y$ 既不是
Kolmogorov 空间，也不拟清醒。
\end{example}

\begin{example}
\label{topology-example-sober-subspace}
设 $Z$ 为无穷集，且 $z\in Z$。在 $Z$ 上赋予以 $Z$ 及
$Z\setminus\{z\}$ 的有限子集为闭集的拓扑。则 $Z$ 清醒，但其子空间
$Z\setminus\{z\}$ 不拟清醒。
\end{example}

\begin{example}
\label{topology-example-Hausdorff}
回忆，拓扑空间 $X$ 为 Hausdorff 空间，当且仅当对任意不同的点
$x, y \in X$，存在不交的开集 $U, V \subset X$，使得
$x \in U$、$y \in V$。在这种情形，$X$ 不可约，当且仅当 $X$ 为单点集。
类似地，$X$ 的任一子集不可约，当且仅当它为单点集。因此 Hausdorff 空间
是清醒的。
\end{example}

\begin{lemma}
\label{topology-lemma-irreducible-on-top}
设 $f : X \to Y$ 为拓扑空间间的连续映射。假设
(a) $Y$ 不可约，
(b) $f$ 是开映射，并且
(c) 存在一族稠密的点 $y \in Y$，使得 $f^{-1}(y)$ 不可约。
则 $X$ 不可约。
\end{lemma}

\begin{proof}
假设 $X = Z_1 \cup Z_2$，其中 $Z_i$ 闭。考虑开集
$U_1 = Z_1 \setminus Z_2 = X \setminus Z_2$ 与
$U_2 = Z_2 \setminus Z_1 = X \setminus Z_1$。为得到矛盾，假设
$U_1$ 与 $U_2$ 都非空。由 (b)，$f(U_i)$ 是开集。由 (a)，$Y$
不可约，故 $f(U_1) \cap f(U_2) \not = \emptyset$。由 (c)，存在一个
对应于该交集中某一点的点 $y$，使得纤维 $X_y = f^{-1}(y)$ 不可约。
于是 $X_y \cap U_1$ 与 $X_y \cap U_2$ 是 $X_y$ 的两个非空不交开子集，
矛盾。
\end{proof}

\begin{lemma}
\label{topology-lemma-irreducible-fibres-irreducible-components}
设 $f : X \to Y$ 为拓扑空间间的连续映射。假设 (a) $f$ 是开映射，并且
(b) 对每个 $y \in Y$，纤维 $f^{-1}(y)$ 都不可约。
则 $f$ 诱导不可约分支之间的双射。
\end{lemma}

\begin{proof}
指出一点：假设 (b) 蕴含 $f$ 满（见定义
\ref{topology-definition-irreducible-components}）。设 $T \subset Y$ 为一个
不可约分支。注意 $T$ 是闭的，参见引理 \ref{topology-lemma-irreducible}。
只要证明 $f^{-1}(T)$ 不可约，本引理即得证，因为由引理
\ref{topology-lemma-image-irreducible-space}，$X$ 的任何不可约子集都映入 $Y$
的某个不可约分支。注意 $f^{-1}(T) \to T$ 满足引理
\ref{topology-lemma-irreducible-on-top} 的假设。因此结论成立。
\end{proof}

\noindent
下面引理中的构造有时称为“清醒化”。

\begin{lemma}
\label{topology-lemma-make-sober}
\begin{reference}
\cite[Page 68 ff]{EGA1-second}
\end{reference}
设 $X$ 为拓扑空间。存在典范连续映射
$$
c : X \longrightarrow X'
$$
从 $X$ 映到一个清醒拓扑空间 $X'$，它在从 $X$ 到清醒拓扑空间的连续映射中
具有泛性质。此外，赋值 $U' \mapsto c^{-1}(U')$ 给出 $X'$ 与 $X$
的开集之间的双射，并且与有限交和任意并相容。像 $c(X)$ 是 Kolmogorov
拓扑空间，而映射 $c : X \to c(X)$ 在从 $X$ 到 Kolmogorov 空间的映射中
具有泛性质。
\end{lemma}

\begin{proof}
令 $X'$ 为 $X$ 的不可约闭子集之集合，并令
$$
c : X \to X', \quad x \mapsto \overline{\{x\}}.
$$
对开集 $U \subset X$，令 $U' \subset X'$ 表示所有与 $U$ 相交的
$X$ 的不可约闭子集之集合。则 $c^{-1}(U') = U$。特别地，若
$U_1 \not = U_2$ 是 $X$ 中的开集，则 $U'_1 \not = U_2'$。
所以 $c$ 诱导 $X'$ 中形如 $U'$ 的子集与 $X$ 的开集之间的双射。

\medskip\noindent
设 $U_1, U_2$ 在 $X$ 中开。假设 $Z \in U'_1$ 且 $Z \in U'_2$。
则 $Z \cap U_1$ 与 $Z \cap U_2$ 是不可约空间 $Z$ 的非空开子集，
所以 $Z \cap U_1 \cap U_2$ 非空。因此
$(U_1 \cap U_2)' = U'_1 \cap U'_2$。规则 $U \mapsto U'$ 也与任意并相容
（细节从略）。于是显然，全体 $U'$ 构成 $X'$ 上的一个拓扑，并且得到
引理所述的双射。

\medskip\noindent
接着证明 $X'$ 清醒。设 $T \subset X'$ 为不可约闭子集。
设 $U \subset X$ 为满足 $X' \setminus T = U'$ 的开集。由引理中双射的性质，
$Z = X \setminus U$ 不可约。我们断言 $Z \in T$ 是唯一泛点。
事实上，任何不含 $Z$ 的形如 $V' \subset X'$ 的开集，都来自某个
不与 $Z$ 相交、即包含于 $U$ 的开集 $V \subset X$。

\medskip\noindent
最后验证泛性质。设 $f : X \to Y$ 为到清醒拓扑空间的连续映射。
令 $f' : X' \to Y$ 将不可约闭子集 $Z \subset X$ 映到
$\overline{f(Z)}$ 的唯一泛点。立即可见，作为集合间映射有
$f' \circ c = f$，而 $c$ 的性质蕴含 $f'$ 连续。连续映射 $f'$ 的唯一性的
验证从略。关于 Kolmogorov 空间的陈述之证明也从略。
\end{proof}

\begin{lemma}
\label{topology-lemma-remove-irreducible-connected}
设 $X$ 为连通拓扑空间，并且只有有限个不可约分支
$X_1, \ldots, X_n$。若 $n > 1$，则存在 $1 \leq j \leq n$，使得
$X' = \bigcup_{i \not = j} X_i$ 连通。
\end{lemma}

\begin{proof}
这是一个图论问题。令 $\Gamma$ 为以 $V = \{1, \ldots, n\}$ 为顶点集的图，
并且 $i$ 与 $j$ 之间有边，当且仅当 $X_i \cap X_j$ 非空。
$X$ 的连通性意味着 $\Gamma$ 连通。问题是找出 $1 \leq j \leq n$，
使 $\Gamma \setminus \{j\}$ 仍然连通。可选取距离最大的
$j, j' \in E$，此时 $j$ 即可（选一个叶子！）。细节从略。
\end{proof}

% P01-ERRATA-E0015: source line 1236 uses j,j' \in E although the graph's vertex set is V; canonical text retained.
\section{Noether 拓扑空间}
\label{topology-section-noetherian}

\begin{definition}
\label{topology-definition-noetherian}
若 $X$ 的闭子集满足降链条件，则称一个拓扑空间为{\it Noether 空间}。
若每一点都有一个 Noether 邻域，则称一个拓扑空间{\it 局部 Noether}。
\end{definition}

\begin{lemma}
\label{topology-lemma-Noetherian}
设 $X$ 为 Noether 拓扑空间。
\begin{enumerate}
\item $X$ 的任意子集带诱导拓扑后都是 Noether 空间。
\item 空间 $X$ 只有有限个不可约分支。
\item $X$ 的每个不可约分支都包含 $X$ 的一个非空开集。
\end{enumerate}
\end{lemma}

\begin{proof}
设 $T \subset X$ 为 $X$ 的一个子集。设
$T_1 \supset T_2 \supset \ldots$ 为 $T$ 的闭子集之降链。
写成 $T_i =  T \cap Z_i$，其中 $Z_i \subset X$ 闭。考虑闭子集的降链
$Z_1 \supset Z_1\cap Z_2 \supset Z_1 \cap Z_2 \cap Z_3 \ldots$。
由假设，此链稳定，因而原来的 $T_i$ 序列也稳定。所以 $T$ 为 Noether 空间。

\medskip\noindent
令 $A$ 为 $X$ 的所有不具有有限多个不可约分支的闭子集之集合。
为得到矛盾，假设 $A$ 非空。包含关系给出 $A$ 上的偏序：
$\alpha \leq \alpha' \Leftrightarrow Z_{\alpha} \subset Z_{\alpha'}$。
由降链条件，可以找到 $A$ 的一个最小元素，记为 $Z$。因为 $Z$ 不是
有限多个不可约分支之并，所以它不可约。故可以写成
$Z = Z' \cup Z''$，其中二者都是严格较小的闭子集。按构造，
$Z' = \bigcup Z'_i$ 与 $Z'' = \bigcup Z''_j$ 各自都是其不可约分支的有限并
（引理 \ref{topology-lemma-irreducible}）。因此
$Z = \bigcup Z'_i \cup \bigcup Z''_j$ 是不可约闭子集的有限并。
从这一表达式中去掉冗余项后，就得到 $Z$ 的不可约分支分解
（引理 \ref{topology-lemma-pick-irreducible-components}），矛盾。

\medskip\noindent
设 $Z \subset X$ 为 $X$ 的一个不可约分支。设
$Z_1, \ldots, Z_n$ 为 $X$ 的其他不可约分支。考虑
$U = Z \setminus (Z_1\cup\ldots\cup Z_n)$。它非空，否则不可约空间
$Z$ 将包含于某个其他 $Z_i$ 中。又因为
$X = Z \cup Z_1 \cup \ldots Z_n$（见引理 \ref{topology-lemma-irreducible}），
还有 $U = X \setminus (Z_1\cup\ldots\cup Z_n)$，故该集合在 $X$ 中开。
所以 $Z$ 包含 $X$ 的一个非空开集。
\end{proof}

\begin{lemma}
\label{topology-lemma-image-Noetherian}
设 $f : X \to Y$ 为拓扑空间间的连续映射。
\begin{enumerate}
\item 若 $X$ 为 Noether 空间，则 $f(X)$ 为 Noether 空间。
\item 若 $X$ 局部 Noether 且 $f$ 为开映射，则 $f(X)$ 局部 Noether。
\end{enumerate}
\end{lemma}

\begin{proof}
在情形 (1)，假设 $Z_1 \supset Z_2 \supset Z_3 \supset \ldots$ 是
$f(X)$ 的闭子集之降链（照例，它作为 $Y$ 的子集带诱导拓扑）。则
$f^{-1}(Z_1) \supset f^{-1}(Z_2) \supset f^{-1}(Z_3) \supset \ldots$
是 $X$ 的闭子集之降链，故该链稳定。由于
$f(f^{-1}(Z_i)) = Z_i$，可知
$Z_1 \supset Z_2 \supset Z_3 \supset \ldots$ 也稳定。
在情形 (2)，设 $y \in f(X)$。选取 $x \in X$ 使得 $f(x) = y$。
由假设，存在 $x$ 的一个 Noether 邻域 $E \subset X$。则由 (1)，
$f(E) \subset f(X)$ 是一个 Noether 邻域。
\end{proof}

\begin{lemma}
\label{topology-lemma-finite-union-Noetherian}
设 $X$ 为拓扑空间。设 $X_i \subset X$（$i = 1, \ldots, n$）为有限个
子集。若每个 $X_i$（带诱导拓扑）都是 Noether 空间，则
$\bigcup_{i = 1, \ldots, n}  X_i$（带诱导拓扑）也是 Noether 空间。
\end{lemma}

\begin{proof}
设 $\{F_m\}_{m \in \mathbf{N}}$ 为带诱导拓扑的
$X' = \bigcup_{i = 1, \ldots, n}  X_i$ 的闭子集之降链。
则可以找到 $X$ 的闭子集之降链 $\{G_m\}_{m \in \mathbf{N}}$，
使得对所有 $m$，$F_m = G_m \cap X'$（略去一个小细节）。
由于 $X_i$ 为 Noether 空间，而
$\{G_m \cap X_i\}_{m \in \mathbf{N}}$ 是 $X_i$ 的闭子集之降链，
故存在 $m_i \in \mathbf{N}$，使得对所有 $m \geq m_i$，都有
$G_m \cap X_i = G_{m_i} \cap X_i$。令
$m_0 = \max_{i = 1, \ldots, n} m_i$。则显然
$$
F_m = G_m \cap X' = G_m \cap (X_1 \cup \ldots \cup X_n) =
(G_m \cap X_1) \cup \ldots (G_m \cap X_n)
$$
在 $m \geq m_0$ 时稳定，证明完毕。
\end{proof}

\begin{example}
\label{topology-example-locally-Noetherian-no-closed-point}
任意非空的 Kolmogorov Noether 拓扑空间都有闭点（合并引理
\ref{topology-lemma-quasi-compact-closed-point} 与
\ref{topology-lemma-Noetherian-quasi-compact}）。令
$X = \{1, 2, 3, \ldots \}$。在 $X$ 上定义以 $\emptyset$、
$\{1, 2, \ldots, n\}$（$n \geq 1$）及 $X$ 为开集的拓扑。
于是 $X$ 是一个没有任何闭点的局部 Noether 拓扑空间。
该空间不可能是局部 Noether 概形的底拓扑空间，参见
Properties, Lemma \ref{properties-lemma-locally-Noetherian-closed-point}。
\end{example}

\begin{lemma}
\label{topology-lemma-locally-Noetherian-locally-connected}
设 $X$ 为局部 Noether 拓扑空间。则 $X$ 局部连通。
\end{lemma}

\begin{proof}
设 $x \in X$。设 $E$ 为 $x$ 的一个邻域。需要找到包含于 $E$ 的
$x$ 的一个连通邻域。由假设，存在 $x$ 的一个 Noether 邻域 $E'$。
于是 $E \cap E'$ 为 Noether 空间，见引理 \ref{topology-lemma-Noetherian}。
设 $E \cap E' = Y_1 \cup \ldots \cup Y_n$ 为其不可约分支分解，见引理
\ref{topology-lemma-Noetherian}。令 $E'' = \bigcup_{x \in Y_i} Y_i$。
这是 $E \cap E'$ 的一个含 $x$ 的连通子集。它包含 $E \cap E'$ 的开集
$E \cap E' \setminus (\bigcup_{x \not \in Y_i} Y_i)$，故它是 $X$ 中
$x$ 的一个邻域。本引理得证。
\end{proof}

\section{Krull 维数}
\label{topology-section-krull-dimension}

\begin{definition}
\label{topology-definition-Krull}
设 $X$ 为拓扑空间。
\begin{enumerate}
\item $X$ 的一个{\it 不可约闭子集链}是序列
$Z_0 \subset Z_1 \subset \ldots \subset Z_n \subset X$，其中 $Z_i$
闭且不可约，并且对 $i = 0, \ldots, n - 1$ 有
$Z_i \not = Z_{i + 1}$。
\item $X$ 的不可约闭子集链
$Z_0 \subset Z_1 \subset \ldots \subset Z_n \subset X$ 的
{\it 长度}是整数 $n$。
\item $X$ 的{\it 维数}，更准确地说是{\it Krull 维数} $\dim(X)$，
是由下式定义的 $\{-\infty, 0, 1, 2, 3, \ldots, \infty\}$ 中的元素：
$$
\dim(X) =
\sup \{\text{不可约闭子集链的长度}\}
$$
因此，$\dim(X) = -\infty$ 当且仅当 $X$ 为空空间。
\item 设 $x \in X$。$X$ 在 $x$ 处的{\it Krull 维数}定义为
$$
\dim_x(X) = \min \{\dim(U), x\in U\subset X\text{ 开}\}
$$
即当 $U$ 遍历 $X$ 中 $x$ 的开邻域时 $\dim(U)$ 的最小值。
\end{enumerate}
\end{definition}

\noindent
注意，若 $U' \subset U \subset X$ 都是开集，则
$\dim(U') \leq \dim(U)$。因此，若 $\dim_x(X) = d$，则 $x$ 有一个
由满足 $\dim(U) = \dim_x(X)$ 的开邻域 $U$ 组成的邻域基本系。

\begin{lemma}
\label{topology-lemma-dimension-supremum-local-dimensions}
设 $X$ 为拓扑空间。则 $\dim(X) = \sup \dim_x(X)$，其中上确界遍历
$X$ 的所有点 $x$。
\end{lemma}

\begin{proof}
显然，对所有 $x \in X$ 都有 $\dim(X) \geq \dim_x(X)$（见定义
\ref{topology-definition-Krull} 后的讨论）。这给出一个方向的不等式。
反之，设 $n \geq 0$，并假设 $\dim(X) \geq n$。则可找到不可约闭子集链
$Z_0 \subset Z_1 \subset \ldots \subset Z_n \subset X$。
选取 $x \in Z_0$。对 $x$ 的每个开邻域 $U$，得到 $U$ 中的
不可约闭子集链
$$
Z_0 \cap U \subset Z_1 \cap U \subset \ldots \subset Z_n \cap U
$$
事实上，$U \cap Z_i$ 在 $U$ 中闭且不可约，并且各包含关系都是严格的
（细节从略；提示：$U \cap Z_i$ 的闭包是 $Z_i$）。由此可见
$\dim_x(X) \geq n$，这证明另一个方向的不等式。
\end{proof}

\begin{example}
\label{topology-example-Krull-Rn}
通常的欧氏空间 $\mathbf{R}^n$ 的 Krull 维数是 $0$。
\end{example}

\begin{example}
\label{topology-example-krull-2set}
令 $X = \{s, \eta\}$，其开集为
$\{\emptyset, \{\eta\}, \{s, \eta\}\}$。在这种情形，一个极大的
不可约闭子集链是 $\{s\} \subset \{s, \eta\}$。因此
$\dim(X) = 1$。很容易推广此例，得到一个 Krull 维数为 $n$ 的
$(n + 1)$ 元拓扑空间。
\end{example}

\begin{definition}
\label{topology-definition-equidimensional}
设 $X$ 为拓扑空间。若 $X$ 的每个不可约分支都具有相同维数，
则称 $X$ {\it 等维}。
\end{definition}

\section{余维数与链状空间}
\label{topology-section-catenary-spaces}

\noindent
这里只定义不可约闭子集的余维数。

\begin{definition}
\label{topology-definition-codimension}
设 $X$ 为拓扑空间。设 $Y \subset X$ 为不可约闭子集。$Y$ 在 $X$ 中的
{\it 余维数}，是 $X$ 中以 $Y$ 起始的不可约闭子集链
$$
Y = Y_0 \subset Y_1 \subset \ldots \subset Y_e \subset X
$$
之长度 $e$ 的上确界。将其记为 $\text{codim}(Y, X)$。
\end{definition}

\noindent
余维数是 $\{0, 1, 2, \ldots\} \cup \{\infty\}$ 中的元素。
若 $\text{codim}(Y, X) < \infty$，则每条链都可以延拓为极大链
（但这些极大链未必具有相同长度）。

\begin{lemma}
\label{topology-lemma-codimension-at-generic-point}
设 $X$ 为拓扑空间。设 $Y \subset X$ 为不可约闭子集。
设 $U \subset X$ 为开子集，且 $Y \cap U$ 非空。则
$$
\text{codim}(Y, X) = \text{codim}(Y \cap U, U)
$$
\end{lemma}

\begin{proof}
规则 $T \mapsto \overline{T}$ 在 $U$ 的不可约闭子集与 $X$ 中和 $U$
相交的不可约闭子集之间定义一个保持包含关系的双射。由此容易得到本引理。
细节从略。
\end{proof}

\begin{example}
\label{topology-example-Noetherian-infinite-codimension}
令 $X = [0, 1]$ 为单位区间，并赋予以下拓扑：开集是 $[0, 1]$、
对 $n \in \mathbf{N}$ 的 $(1 - 1/n, 1]$ 以及 $\emptyset$。
所以闭集是 $\emptyset$、$\{0\}$、对 $n > 1$ 的
$[0, 1 - 1/n]$ 以及 $[0, 1]$。这显然是一个 Noether 拓扑空间。
但不可约闭子集 $Y = \{0\}$ 具有无穷余维数
$\text{codim}(Y, X) = \infty$。只需注意所有闭集
$[0, 1 - 1/n]$ 都不可约即可看出这一点。
\end{example}

\begin{definition}
\label{topology-definition-catenary}
设 $X$ 为拓扑空间。如果对每一对不可约闭子集 $T \subset T'$，都有
$\text{codim}(T, T') < \infty$，并且每条不可约闭子集极大链
$$
T = T_0 \subset T_1 \subset \ldots \subset T_e = T'
$$
都具有相同长度（等于余维数），则称 $X$ 是{\it 链状的}。
\end{definition}

\begin{lemma}
\label{topology-lemma-catenary}
设 $X$ 为拓扑空间。下列条件等价：
\begin{enumerate}
\item $X$ 是链状的，
\item $X$ 有一个由链状空间组成的开覆盖。
\end{enumerate}
此外，在这种情形，$X$ 的任意局部闭子空间都是链状的。
\end{lemma}

\begin{proof}
假设 $X$ 是链状的，并设 $U \subset X$ 为开子集。规则
$T \mapsto \overline{T}$ 在 $U$ 的不可约闭子集与 $X$ 中和 $U$
相交的不可约闭子集之间定义一个保持包含关系的双射。由此容易得到本引理。
细节从略。
\end{proof}

\begin{lemma}
\label{topology-lemma-catenary-in-codimension}
设 $X$ 为拓扑空间。下列条件等价：
\begin{enumerate}
\item $X$ 是链状的，并且
\item 对每一对不可约闭子集 $Y \subset Y'$，都有
$\text{codim}(Y, Y') < \infty$，且对每一组三个不可约闭子集
$Y \subset Y' \subset Y''$，都有
$$
\text{codim}(Y, Y'') = \text{codim}(Y, Y') + \text{codim}(Y', Y'').
$$
\end{enumerate}
\end{lemma}

\begin{proof}
先假设 $X$ 是链状的。由定义 \ref{topology-definition-catenary}，对每一对
不可约闭子集 $Y \subset Y'$ 都有 $\text{codim}(Y,Y') < \infty$。
设 $Y \subset Y' \subset Y''$ 为 $X$ 的三个不可约闭子集。设
$$
Y = Y_0 \subset Y_1 \subset ... \subset Y_{e_1} = Y'
$$
为 $Y$ 与 $Y'$ 之间的不可约闭子集极大链，其中
$e_1 = \text{codim}(Y,Y')$。又设
$$
Y' = Y_{e_1} \subset Y_{e_1 + 1}\subset ... \subset Y_{e_1 + e_2} = Y''
$$
为 $Y'$ 与 $Y''$ 之间的不可约闭子集极大链，其中
$e_2 = \text{codim}(Y',Y'')$。由于这两条链都是极大的，其连接
$$
Y = Y_0\subset Y_1 \subset ... \subset Y_{e_1} = Y' =
Y_{e_1} \subset Y_{e_1+1}\subset ... \subset Y_{e_1+e_2}=Y''
$$
也是 $Y$ 与 $Y''$ 之间的极大链，且其长度等于 $e_1 + e_2$。
因为 $X$ 是链状的，每条极大链都具有等于余维数的相同长度。
所以 (2) 中的等式
$\text{codim}(Y,Y'') = e_1 + e_2 = \text{codim}(Y,Y') + \text{codim}(Y',Y'')$
成立。

\medskip\noindent
反之，由归纳法证明：若 $Y = Y_1 \subset ... \subset Y_n = Y'$，则
$ \text{codim}(Y,Y') = \text{codim}(Y_1,Y_2) + ... +
\text{codim}(Y_{n-1},Y_n)$。因此，极大链被迫具有相同长度。
\end{proof}

\section{拟紧空间与拟紧映射}
\label{topology-section-quasi-compact}

\noindent
“紧”一词将保留给 Hausdorff 拓扑空间使用，而代数几何中出现的许多空间
并不是 Hausdorff 空间。

\begin{definition}
\label{topology-definition-quasi-compact}
拟紧性。
\begin{enumerate}
\item 若拓扑空间 $X$ 的每个开覆盖都有有限子覆盖，则称 $X$ {\it 拟紧}。
\item 若对每个拟紧开集 $V \subset Y$，连续映射 $f : X \to Y$
的逆像 $f^{-1}(V)$ 都拟紧，则称该映射{\it 拟紧}。
\item 若包含映射 $Z \to X$ 拟紧，则称子集 $Z \subset X$ {\it 逆紧}。
\end{enumerate}
\end{definition}

\noindent
许多拓扑学文本把拟紧且 Hausdorff 的空间称为{\it 紧空间}；另一些文本
则省略 Hausdorff 条件。为避免代数几何中的混淆，我们采用“拟紧”一词。
映射的拟紧性与“固有映射”的概念大不相同：在后一概念中，除了闭性与分离性，
还要求目标中任意拟紧子集的逆像拟紧；而在上面的定义中，我们仅考虑拟紧
{\it 开}集。

\begin{lemma}
\label{topology-lemma-composition-quasi-compact}
拟紧映射的复合是拟紧的。
\end{lemma}

\begin{proof}
由定义立即可得。
\end{proof}

\begin{lemma}
\label{topology-lemma-closed-in-quasi-compact}
拟紧拓扑空间的闭子集是拟紧的。
\end{lemma}

\begin{proof}
设 $E \subset X$ 为拟紧空间 $X$ 的闭子集。设 $E = \bigcup V_j$
为一个开覆盖。选取开集 $U_j \subset X$，使得 $V_j = E \cap U_j$。
则 $X = (X \setminus E) \cup \bigcup U_j$ 是 $X$ 的开覆盖。因此，对某个
$n$ 及某些指标 $j_i$，有
$X = (X \setminus E) \cup U_{j_1} \cup \ldots \cup U_{j_n}$。
于是按所需，$E = V_{j_1} \cup \ldots \cup V_{j_n}$。
\end{proof}

\begin{lemma}
\label{topology-lemma-quasi-compact-in-Hausdorff}
设 $X$ 为 Hausdorff 拓扑空间。
\begin{enumerate}
\item 若 $E \subset X$ 拟紧，则它是闭集。
\item 若 $E_1, E_2 \subset X$ 是不交的拟紧子集，则存在开集
$E_i \subset U_i$，使得 $U_1 \cap U_2 = \emptyset$。
\end{enumerate}
\end{lemma}

\begin{proof}
(1) 的证明。设 $x \in X$，且 $x \not \in E$。对每个 $e \in E$，
都可以找到不交的开集 $V_e$ 与 $U_e$，使得 $e \in V_e$ 且
$x \in U_e$。因为 $E \subset \bigcup V_e$，可以找到有限多个
$e_1, \ldots, e_n$，使得 $E \subset V_{e_1} \cup \ldots \cup V_{e_n}$。
于是 $U = U_{e_1} \cap \ldots \cap U_{e_n}$ 是 $x$ 的一个开邻域，
它不与 $V_{e_1} \cup \ldots \cup V_{e_n}$ 相交，特别地不与 $E$ 相交。
所以 $E$ 是闭的。

\medskip\noindent
(2) 的证明。在 (1) 的证明中已经看到，给定 $x \in E_1$，可以找到
开邻域 $x \in U_x$ 和一个开集 $E_2 \subset V_x$，使得
$U_x \cap V_x = \emptyset$。由于 $E_1$ 拟紧，可以找到有限多个
$x_i \in E_1$，使得 $E_1 \subset U = U_{x_1} \cup \ldots \cup U_{x_n}$。
取 $V = V_{x_1} \cap \ldots \cap V_{x_n}$，证明即告完成。
\end{proof}

\begin{lemma}
\label{topology-lemma-closed-in-compact}
设 $X$ 为拟紧 Hausdorff 空间，$E \subset X$。下列条件等价：
(a) $E$ 在 $X$ 中闭；(b) $E$ 拟紧。
\end{lemma}

\begin{proof}
蕴含关系 (a) $\Rightarrow$ (b) 是引理
\ref{topology-lemma-closed-in-quasi-compact}；蕴含关系 (b) $\Rightarrow$ (a)
是引理 \ref{topology-lemma-quasi-compact-in-Hausdorff}。
\end{proof}

\noindent
下面的结果其实只是拟紧性质的另一种表述。

\begin{lemma}
\label{topology-lemma-intersection-closed-in-quasi-compact}
设 $X$ 为拟紧拓扑空间。若 $\{Z_\alpha\}_{\alpha \in A}$ 是一族闭子集，
并且每个有限子族之交均非空，则
$\bigcap_{\alpha \in A} Z_\alpha$ 非空。
\end{lemma}

\begin{proof}
假设 $\bigcap_{\alpha \in A} Z_{\alpha} = \emptyset$。
取补集可得 $\bigcup_{\alpha \in A} (X\setminus Z_{\alpha})=X$。
因为各子集 $Z_\alpha$ 都闭，
$\bigcup_{\alpha \in A} (X \setminus Z_{\alpha})$ 是拟紧空间 $X$ 的
开覆盖。因此存在有限子集 $J \subset A$，使得
$X = \bigcup_{\alpha \in J} (X\setminus Z_{\alpha})$。
其补集于是为空，即 $\bigcap_{\alpha \in J} Z_{\alpha} = \emptyset$。
这证明存在有限子族 $\{Z_{\alpha}\}_{\alpha \in J}$ 满足
$\bigcap_{\alpha \in J} Z_{\alpha}=\emptyset$；由逆否命题得到结论。
\end{proof}

\begin{lemma}
\label{topology-lemma-image-quasi-compact}
设 $f : X \to Y$ 为拓扑空间间的连续映射。
\begin{enumerate}
\item 若 $X$ 拟紧，则 $f(X)$ 拟紧。
\item 若 $f$ 拟紧，则 $f(X)$ 逆紧。
\end{enumerate}
\end{lemma}

\begin{proof}
若 $f(X) = \bigcup V_i$ 是开覆盖，则
$X = \bigcup f^{-1}(V_i)$ 是开覆盖。因此，若 $X$ 拟紧，则对某些
$i_1, \ldots, i_n \in I$，有
$X = f^{-1}(V_{i_1}) \cup \ldots \cup f^{-1}(V_{i_n})$，从而
$f(X) = V_{i_1} \cup \ldots \cup V_{i_n}$。这证明 (1)。
假设 $f$ 拟紧，并设 $V \subset Y$ 为拟紧开集。则 $f^{-1}(V)$ 拟紧，
故由 (1)，$f(f^{-1}(V)) = f(X) \cap V$ 拟紧。因此 $f(X)$ 逆紧。
\end{proof}

\begin{lemma}
\label{topology-lemma-quasi-compact-closed-point}
设 $X$ 为拓扑空间。假设
\begin{enumerate}
\item $X$ 非空，
\item $X$ 拟紧，并且
\item $X$ 为 Kolmogorov 空间。
\end{enumerate}
则 $X$ 有闭点。
\end{lemma}

\begin{proof}
考虑 $X$ 中所有单点闭包所成的集合
$$
\mathcal{T} =
\{Z \subset X \mid Z = \overline{\{x\}} \text{ 对某个 }x \in X\}
$$
因为 $X$ 非空，它也非空。用包含关系将 $\mathcal{T}$ 视为偏序集。
设 $Z_\alpha$（$\alpha \in A$）是 $\mathcal{T}$ 的一个全序子集。
由引理 \ref{topology-lemma-intersection-closed-in-quasi-compact}，
$\bigcap_{\alpha \in A} Z_\alpha \not = \emptyset$。故存在某个
$x \in \bigcap_{\alpha \in A} Z_\alpha$，而且
$Z = \overline{\{x\}}\in \mathcal{T}$ 是该族的一个下界。
由 Zorn 引理，存在极小元素 $Z \in \mathcal{T}$。因为 $X$ 为
Kolmogorov 空间，可知对某个 $x$ 有 $Z = \{x\}$，且 $x \in X$
是闭点。
\end{proof}

\begin{lemma}
\label{topology-lemma-closed-points-quasi-compact}
设 $X$ 为拟紧 Kolmogorov 空间。则 $X$ 的闭点集合 $X_0$ 拟紧。
\end{lemma}

\begin{proof}
设 $X_0 = \bigcup U_{i, 0}$ 为开覆盖。对某个开集 $U_i \subset X$，
写成 $U_{i, 0} = X_0 \cap U_i$。考虑 $\bigcup U_i$ 的补集 $Z$。
它是 $X$ 的闭子集，故拟紧（引理 \ref{topology-lemma-closed-in-quasi-compact}），
并且是 Kolmogorov 空间。由引理 \ref{topology-lemma-quasi-compact-closed-point}，
若 $Z$ 非空，它就会有闭点，这与 $X_0 \subset \bigcup U_i$ 矛盾。
因此 $Z = \emptyset$ 且 $X = \bigcup U_i$。由于 $X$ 拟紧，该覆盖有
有限子覆盖，结论得证。
\end{proof}

\begin{lemma}
\label{topology-lemma-connected-component-intersection}
设 $X$ 为拓扑空间。假设
\begin{enumerate}
\item $X$ 拟紧，
\item $X$ 有一个由拟紧开集组成的拓扑基，并且
\item 两个拟紧开集之交拟紧。
\end{enumerate}
对任意 $x \in X$，$X$ 中包含 $x$ 的连通分支，等于 $X$ 中所有包含
$x$ 的既开又闭子集之交。
\end{lemma}

\begin{proof}
设 $T$ 为包含 $x$ 的连通分支。设
$S = \bigcap_{\alpha \in A} Z_\alpha$ 为 $X$ 中所有包含 $x$ 的既开又闭
子集 $Z_\alpha$ 之交。注意 $S$ 在 $X$ 中闭。还要注意，任意有限多个
$Z_\alpha$ 之交仍是某个 $Z_\alpha$。因为 $T$ 连通且 $x \in T$，有
$T \subset S$。只需证明 $S$ 连通。若不然，则存在不交并分解
$S = B \amalg C$，其中 $B$ 与 $C$ 在 $S$ 中既开又闭。特别地，
$B$ 与 $C$ 在 $X$ 中闭，所以由引理 \ref{topology-lemma-closed-in-quasi-compact}
及假设 (1)，二者拟紧。由假设 (2)，存在拟紧开集 $U, V \subset X$，
使得 $B = S \cap U$ 且 $C = S \cap V$（细节从略）。于是
$U \cap V \cap S = \emptyset$，故
$\bigcap_\alpha U \cap V \cap Z_\alpha = \emptyset$。
由假设 (3)，交 $U \cap V$ 拟紧。由引理
\ref{topology-lemma-intersection-closed-in-quasi-compact}，对某个
$\alpha' \in A$，有 $U \cap V \cap Z_{\alpha'} = \emptyset$。
因为 $X \setminus (U \cup V)$ 与 $S$ 不交，且在 $X$ 中闭、因而拟紧，
可以使用同一引理，看出对某个 $\alpha'' \in A$ 有
$Z_{\alpha''} \subset U \cup V$。于是
$Z_\alpha = Z_{\alpha'} \cap Z_{\alpha''}$ 包含于 $U \cup V$，
并与 $U \cap V$ 不交。所以
$Z_\alpha = U \cap Z_\alpha \amalg V \cap Z_\alpha$ 是两个开部分的不交并，
从而 $U \cap Z_\alpha$ 与 $V \cap Z_\alpha$ 在 $X$ 中既开又闭。
因此，若比如 $x \in B$，则可见 $S \subset U \cap Z_\alpha$，
从而 $C = \emptyset$。
\end{proof}

\begin{lemma}
\label{topology-lemma-connected-component-intersection-compact-Hausdorff}
设 $X$ 为拓扑空间。假设 $X$ 拟紧且为 Hausdorff 空间。
对任意 $x \in X$，$X$ 中包含 $x$ 的连通分支，等于 $X$ 中所有包含
$x$ 的既开又闭子集之交。
\end{lemma}

\begin{proof}
设 $T$ 为包含 $x$ 的连通分支。设
$S = \bigcap_{\alpha \in A} Z_\alpha$ 为 $X$ 中所有包含 $x$ 的既开又闭
子集 $Z_\alpha$ 之交。注意 $S$ 在 $X$ 中闭。还要注意，任意有限多个
$Z_\alpha$ 之交仍是某个 $Z_\alpha$。因为 $T$ 连通且 $x \in T$，有
$T \subset S$。只需证明 $S$ 连通。若不然，则存在不交并分解
$S = B \amalg C$，其中 $B$ 与 $C$ 在 $S$ 中既开又闭。特别地，
$B$ 与 $C$ 在 $X$ 中闭，所以由引理 \ref{topology-lemma-closed-in-quasi-compact}
可知二者拟紧。由引理 \ref{topology-lemma-quasi-compact-in-Hausdorff}，
存在不交的开集 $U, V \subset X$，使得 $B \subset U$ 且 $C \subset V$。
于是 $X \setminus U \cup V$ 在 $X$ 中闭，故拟紧
（引理 \ref{topology-lemma-closed-in-quasi-compact}）。由引理
\ref{topology-lemma-intersection-closed-in-quasi-compact}，对某个 $\alpha$ 有
$(X \setminus U \cup V) \cap Z_\alpha = \emptyset$。换言之，
$Z_\alpha \subset U \cup V$。因此
$Z_\alpha = Z_\alpha \cap V \amalg Z_\alpha \cap U$ 是两个开部分的不交并，
从而 $U \cap Z_\alpha$ 与 $V \cap Z_\alpha$ 在 $X$ 中既开又闭。
因此，若比如 $x \in B$，则可见 $S \subset U \cap Z_\alpha$，
从而 $C = \emptyset$。
\end{proof}

\begin{lemma}
\label{topology-lemma-closed-union-connected-components}
设 $X$ 为拓扑空间。假设
\begin{enumerate}
\item $X$ 拟紧，
\item $X$ 有一个由拟紧开集组成的拓扑基，并且
\item 两个拟紧开集之交拟紧。
\end{enumerate}
对一个子集 $T \subset X$，下列条件等价：
\begin{enumerate}
\item[(a)] $T$ 是 $X$ 的若干既开又闭子集之交，并且
\item[(b)] $T$ 在 $X$ 中闭，且是 $X$ 的若干连通分支之并。
\end{enumerate}
\end{lemma}

\begin{proof}
显然 (a) 蕴含 (b)。假设 (b)。设 $x \in X$，且 $x \not \in T$。
设 $x \in C \subset X$ 为 $X$ 中包含 $x$ 的连通分支。由引理
\ref{topology-lemma-connected-component-intersection}，可见
$C = \bigcap V_\alpha$ 是 $X$ 中所有包含 $C$ 的既开又闭子集
$V_\alpha$ 之交。特别地，任意两项之交 $V_\alpha \cap V_\beta$
仍作为某个 $V_\alpha$ 出现。因为 $T$ 是 $X$ 的若干连通分支之并，
$C \cap T = \emptyset$。所以
$T \cap \bigcap V_\alpha = \emptyset$。由于 $T$ 是拟紧空间的闭子集，
所以它拟紧（见引理 \ref{topology-lemma-closed-in-quasi-compact}）。由引理
\ref{topology-lemma-intersection-closed-in-quasi-compact}，可知对某个 $\alpha$ 有
$T \cap V_\alpha = \emptyset$。对这个 $\alpha$，
$U_\alpha = X \setminus V_\alpha$ 是 $X$ 的一个既开又闭的子集，
它包含 $T$ 但不含 $x$。本引理得证。
\end{proof}

\begin{lemma}
\label{topology-lemma-Noetherian-quasi-compact}
设 $X$ 为 Noether 拓扑空间。
\begin{enumerate}
\item 空间 $X$ 拟紧。
\item $X$ 的任意子集都逆紧。
\end{enumerate}
\end{lemma}

\begin{proof}
假设 $X = \bigcup U_i$ 是由集合 $I$ 索引的 $X$ 的开覆盖，并且不存在
有限开覆盖对它的加细。按下述方式归纳选取 $I$ 中的元素
$i_1, i_2, \ldots $：选取 $i_{n + 1}$，使得 $U_{i_{n + 1}}$
不包含于 $U_{i_1} \cup \ldots \cup U_{i_n}$。于是得到闭子集的严格
无穷降链
$X \supset (X \setminus U_{i_1}) \supset
(X \setminus U_{i_1} \cup U_{i_2}) \supset \ldots$。
这与 $X$ 为 Noether 空间矛盾，故第一个断言得证。第二个断言现在很清楚，
因为由引理 \ref{topology-lemma-Noetherian}，$X$ 的每个子集都是 Noether 空间。
\end{proof}

\begin{lemma}
\label{topology-lemma-quasi-compact-locally-Noetherian-Noetherian}
拟紧的局部 Noether 空间是 Noether 空间。
\end{lemma}

\begin{proof}
这些条件立即蕴含 $X$ 有一个由有限多个 Noether 子集组成的覆盖，故由引理
\ref{topology-lemma-finite-union-Noetherian}，该空间是 Noether 空间。
\end{proof}

\begin{lemma}[Alexander 子基定理]
\label{topology-lemma-subbase-theorem}
设 $X$ 为拓扑空间，$\mathcal{B}$ 为 $X$ 的一个子基。若每个由
$\mathcal{B}$ 中元素组成的 $X$ 的覆盖都有有限加细，则 $X$ 拟紧。
\end{lemma}

\begin{proof}
假设 $X$ 有一个不存在有限加细的开覆盖。利用 Zorn 引理，可以选取
一个不存在有限加细的极大开覆盖
$X = \bigcup_{i \in I} U_i$（细节从略）。换言之，若 $U \subset X$
是一个没有作为某个 $U_i$ 出现的开集，则覆盖
$X = U \cup \bigcup_{i \in I} U_i$ 确实有有限加细。
令 $I' \subset I$ 为满足 $U_i \in \mathcal{B}$ 的指标集合。
则 $\bigcup_{i \in I'} U_i \not = X$，否则由对 $\mathcal{B}$ 的假设，
会得到覆盖 $X$ 的有限加细。选取 $x \in X$，使得
$x \not \in \bigcup_{i \in I'} U_i$。选取 $i \in I$ 使 $x \in U_i$。
选取 $V_1, \ldots, V_n \in \mathcal{B}$，使得
$x \in V_1 \cap \ldots \cap V_n \subset U_i$。这是可行的，因为
$\mathcal{B}$ 是 $X$ 的一个子基。注意 $V_j$ 不作为某个 $U_i$ 出现。
由所选覆盖的极大性，对每个 $j$，存在
$i_{j, 1}, \ldots, i_{j, n_j} \in I$，使得
$X = V_j \cup U_{i_{j, 1}} \cup \ldots \cup U_{i_{j, n_j}}$。
因为 $V_1 \cap \ldots \cap V_n \subset U_i$，可得
$X = U_i \cup \bigcup U_{i_{j, l}}$，矛盾。
\end{proof}

% P01-ERRATA-E0016: source lines 1882 and 1884 omit parentheses around U \cup V in two set complements; canonical formulas retained.
% P01-ERRATA-E0017: source line 1947 writes X \setminus U_{i_1} \cup U_{i_2} instead of the complement of the growing union; canonical formula retained.
\section{局部拟紧空间}
\label{topology-section-locally-quasi-compact}

\noindent
回忆一点：一个点的邻域未必是开集。

\begin{definition}
\label{topology-definition-locally-quasi-compact}
若拓扑空间 $X$ 的每个点都有一个拟紧邻域基本系，则称该空间
{\it 局部拟紧}\footnote{这可能不是标准术语。文献中还使用下列定义：
(1) 每个点都有某个拟紧邻域；(2) 每个点都有一个闭拟紧邻域。
概形具有如下性质：每个点都有一个由开拟紧邻域组成的邻域基本系。}。
\end{definition}

\noindent
文献中的{\it 局部紧空间}一词通常指下列引理中的空间。

\begin{lemma}
\label{topology-lemma-locally-quasi-compact-Hausdorff}
Hausdorff 空间局部拟紧，当且仅当每个点都有一个拟紧邻域。
\end{lemma}

\begin{proof}
设 $X$ 为 Hausdorff 空间。设 $x \in X$，并设
$x \in E \subset X$ 为一个拟紧邻域。由引理
\ref{topology-lemma-quasi-compact-in-Hausdorff}，$E$ 闭。假设
$x \in U \subset X$ 是 $x$ 的一个开邻域。则 $Z = E \setminus U$
是 $E$ 的一个不含 $x$ 的闭子集。由引理
\ref{topology-lemma-quasi-compact-in-Hausdorff}，可以找到 $E$ 的一对不交开子集
$W, V \subset E$，使得 $x \in V$ 且 $Z \subset W$。
于是 $\overline{V} \subset E$ 是 $x$ 的一个包含于 $E \cap U$ 的闭邻域。
此外，$\overline{V}$ 作为 $E$ 的闭子集是拟紧的
（引理 \ref{topology-lemma-closed-in-quasi-compact}）。由此得到 $x$ 的一个
拟紧邻域基本系。
\end{proof}

\begin{lemma}[Baire 纲定理]
\label{topology-lemma-baire-category-locally-compact}
设 $X$ 为局部拟紧 Hausdorff 空间。设 $U_n \subset X$（$n \geq 1$）
为稠密开子集。则 $\bigcap_{n \geq 1} U_n$ 在 $X$ 中稠密。
\end{lemma}

\begin{proof}
将 $U_n$ 替换为 $\bigcap_{i = 1, \ldots, n} U_i$ 后，可以假设
$U_1 \supset U_2 \supset \ldots$。设 $x \in X$。我们将证明 $x$ 属于
$\bigcap_{n \geq 1} U_n$ 的闭包。为此，设 $E$ 为 $x$ 的一个邻域。
为证明 $E \cap \bigcap_{n \geq 1} U_n$ 非空，可以把 $E$ 换成一个
更小的邻域。将 $E$ 这样替换后，可以假设 $E$ 拟紧。

\medskip\noindent
置 $x_0 = x$ 且 $E_0 = E$。下面归纳选取点
$x_i \in E_{i - 1} \cap U_i$，以及 $x_i$ 的拟紧邻域 $E_i$，使得
$E_i \subset E_{i - 1} \cap U_i$。因为 $X$ 为 Hausdorff 空间，
各子集 $E_i \subset X$ 都是闭的
（引理 \ref{topology-lemma-quasi-compact-in-Hausdorff}）。又因各 $E_i$ 非空，
由引理 \ref{topology-lemma-intersection-closed-in-quasi-compact} 可知
$\bigcap_{i \geq 1} E_i$ 非空。由于
$\bigcap_{i \geq 1} E_i \subset E \cap \bigcap_{n \geq 1} U_n$，
本引理得证。

\medskip\noindent
基础情形 $i = 0$ 已在上面完成。现作归纳步骤。因为 $E_{i - 1}$ 是
$x_{i - 1}$ 的邻域，可以找到开集
$x_{i - 1} \in W \subset E_{i - 1}$。由于 $U_i$ 在 $X$ 中稠密，
$W \cap U_i$ 非空。任取 $x_i \in W \cap U_i$。由局部拟紧空间的定义，
可以找到 $x_i$ 的一个包含于 $W \cap U_i$ 的拟紧邻域 $E_i$。
于是按所需，$E_i \subset E_{i - 1} \cap U_i$。
\end{proof}

\begin{lemma}
\label{topology-lemma-relatively-compact-refinement}
设 $X$ 为 Hausdorff 且拟紧的空间。设
$X = \bigcup_{i \in I} U_i$ 为开覆盖。则存在开覆盖
$X = \bigcup_{i \in I} V_i$，使得对所有 $i$ 都有
$\overline{V_i} \subset U_i$。
\end{lemma}

\begin{proof}
设 $x \in X$。选取 $i(x) \in I$，使得 $x \in U_{i(x)}$。
因为 $X \setminus U_{i(x)}$ 与 $\{x\}$ 是 $X$ 的不交闭子集，
由引理 \ref{topology-lemma-closed-in-quasi-compact} 和
\ref{topology-lemma-quasi-compact-in-Hausdorff}，存在 $x$ 的开邻域 $U_x$，
其闭包与 $X \setminus U_{i(x)}$ 不交。所以
$\overline{U_x} \subset U_{i(x)}$。因为 $X$ 拟紧，存在有限多个点
$x_1, \ldots, x_m$，使得
$X = U_{x_1} \cup \ldots \cup U_{x_m}$。令
$V_i = \bigcup_{i = i(x_j)} U_{x_j}$，证明完成。
\end{proof}

\begin{lemma}
\label{topology-lemma-refine-covering}
设 $X$ 为 Hausdorff 且拟紧的空间。设
$X = \bigcup_{i \in I} U_i$ 为开覆盖。假设给定整数 $p \geq 0$，
并且对 $I$ 的每个 $(p + 1)$ 元组 $i_0, \ldots, i_p$，都给定开覆盖
$U_{i_0} \cap \ldots \cap U_{i_p} = \bigcup W_{i_0 \ldots i_p, k}$。
则存在开覆盖 $X = \bigcup_{j \in J} V_j$ 与映射
$\alpha : J \to I$，使得 $\overline{V_j} \subset U_{\alpha(j)}$，
并且对某个 $k$，每个 $V_{j_0} \cap \ldots \cap V_{j_p}$ 都包含于
$W_{\alpha(j_0) \ldots \alpha(j_p), k}$ 中。
\end{lemma}

\begin{proof}
因为 $X$ 拟紧，可以归约到 $I$ 有限的情形（细节从略）。
对有限的 $I$，关于 $p$ 归纳证明。基础情形 $p = 0$ 立即由取引理
\ref{topology-lemma-relatively-compact-refinement} 中的一个覆盖得到，该覆盖加细开覆盖
$X = \bigcup W_{i_0, k}$。

\medskip\noindent
归纳步骤。假设本引理对 $p - 1$ 已证。对 $I$ 的所有 $p$ 元组
$i'_0, \ldots, i'_{p - 1}$，令
$U_{i'_0} \cap \ldots \cap U_{i'_{p - 1}} =
\bigcup W_{i'_0 \ldots i'_{p - 1}, k}$ 为如下这些覆盖的一个共同加细：
$U_{i_0} \cap \ldots \cap U_{i_p} = \bigcup W_{i_0 \ldots i_p, k}$，
其中相应的 $(p + 1)$ 元组满足
$\{i'_0, \ldots, i'_{p - 1}\} = \{i_0, \ldots, i_p\}$（集合相等）。
（因为 $I$ 有限，这样的覆盖只有有限多个。）由归纳假设，对这些开集存在
一个解，记为 $X = \bigcup V_j$ 与 $\alpha : J \to I$。
此时，覆盖 $X = \bigcup_{j \in J} V_j$ 与 $\alpha$ 满足
$\overline{V_j} \subset U_{\alpha(j)}$；并且只要
$\alpha(j_0), \ldots, \alpha(j_p)$ 中有重复项，对某个 $k$，每个
$V_{j_0} \cap \ldots \cap V_{j_p}$ 都包含于
$W_{\alpha(j_0) \ldots \alpha(j_p), k}$。当然，可以且确实假设 $J$ 有限。

\medskip\noindent
固定两两不同的 $i_0, \ldots, i_p \in I$。考虑 $(p + 1)$ 元组
$j_0, \ldots, j_p \in J$，它满足
$i_0 = \alpha(j_0), \ldots, i_p = \alpha(j_p)$，并且
$V_{j_0} \cap \ldots \cap V_{j_p}$ 对任意 $k$ 都{\bf 不}包含于
$W_{\alpha(j_0) \ldots \alpha(j_p), k}$。令 $N$ 为这种
$(p + 1)$ 元组的个数。我们说明如何减小 $N$。由于
$$
\overline{V_{j_0}} \cap \ldots \cap \overline{V_{j_p}} \subset
U_{i_0} \cap \ldots \cap U_{i_p} = \bigcup W_{i_0 \ldots i_p, k}
$$
可以找到有限集 $K = \{k_1, \ldots, k_t\}$，使左端包含于
$\bigcup_{k \in K} W_{i_0 \ldots i_p, k}$。于是考虑开覆盖
$$
V_{j_0} =
(V_{j_0} \setminus (\overline{V_{j_1}} \cap \ldots \cap \overline{V_{j_p}}))
\cup (\bigcup\nolimits_{k \in K} V_{j_0} \cap W_{i_0 \ldots i_p, k})
$$
右端第一个开集与 $V_{j_1} \cap \ldots \cap V_{j_p}$ 之交为空，
而右端其他开集 $V_{j_0, k}$ 满足
$V_{j_0, k} \cap V_{j_1} \ldots \cap V_{j_p} \subset
W_{\alpha(j_0) \ldots \alpha(j_p), k}$。令 $J' = J \amalg K$。
对 $j \in J$，若 $j \not = j_0$，令 $V'_j = V_j$；并令
$V'_{j_0} =
V_{j_0} \setminus (\overline{V_{j_1}} \cap \ldots \cap \overline{V_{j_p}})$。
对 $k \in K$，令 $V'_k = V_{j_0, k}$。最后，映射
$\alpha' : J' \to I$ 在 $J$ 上由 $\alpha$ 给出，并将 $K$ 的每个元素
映到 $i_0$。简单检查可见，经此替换后 $N$ 减少一。
重复此过程 $N$ 次，就达到 $N = 0$ 的情形。

\medskip\noindent
为完成证明，对下述数目 $M$ 作归纳：具有两两不同分量的
$(p + 1)$ 元组 $i_0, \ldots, i_p \in I$ 的数目，其中存在
$(p + 1)$ 元组 $j_0, \ldots, j_p \in J$，满足
$i_0 = \alpha(j_0), \ldots, i_p = \alpha(j_p)$，且
$V_{j_0} \cap \ldots \cap V_{j_p}$ 对任意 $k$ 都{\bf 不}包含于
$W_{\alpha(j_0) \ldots \alpha(j_p), k}$。为此，断言上一段所作的操作
不会增大 $M$。这形式地来自如下事实：映射 $\alpha' : J' \to I$
通过某个映射 $\beta : J' \to J$ 分解，并且
$V'_{j'} \subset V_{\beta(j')}$。
\end{proof}

\begin{lemma}
\label{topology-lemma-lift-covering-of-a-closed}
设 $X$ 为 Hausdorff 且局部拟紧的空间。设 $Z \subset X$ 为拟紧
（因而闭的）子集。假设给定整数 $p \geq 0$、集合 $I$、对每个
$i \in I$ 给定开集 $U_i \subset X$，并且对 $I$ 的每个
$(p + 1)$ 元组 $i_0, \ldots, i_p$ 给定开集
$W_{i_0 \ldots i_p} \subset U_{i_0} \cap \ldots \cap U_{i_p}$，使得
\begin{enumerate}
\item $Z \subset \bigcup U_i$，并且
\item 对每个 $i_0, \ldots, i_p$，有
$W_{i_0 \ldots i_p} \cap Z = U_{i_0} \cap \ldots \cap U_{i_p} \cap Z$。
\end{enumerate}
则存在 $X$ 的开集 $V_i$，使得 $Z \subset \bigcup V_i$，
对所有 $i$ 都有 $\overline{V_i} \subset U_i$，并且对所有
$(p + 1)$ 元组 $i_0, \ldots, i_p$，都有
$V_{i_0} \cap \ldots \cap V_{i_p} \subset W_{i_0 \ldots i_p}$。
\end{lemma}

\begin{proof}
因为 $Z$ 拟紧，可以归约到 $I$ 有限的情形（细节从略）。
因为 $X$ 局部拟紧且 $Z$ 拟紧，可以找到拟紧邻域 $Z \subset E$，
即 $E$ 拟紧并且包含 $Z$ 在 $X$ 中的一个开邻域。若将 $X$ 替换为 $E$
后能证明该结果，则原结果随即成立。因此可假设 $X$ 拟紧。

\medskip\noindent
在 $I$ 有限且 $X$ 拟紧的情形，关于 $p$ 归纳证明。基础情形为
$p = 0$。此时 $X = (X \setminus Z) \cup \bigcup W_i$。由引理
\ref{topology-lemma-relatively-compact-refinement}，可以找到由开集组成的覆盖
$X = V \cup \bigcup V_i$，其中 $V_i \subset W_i$、
$V \subset X \setminus Z$，并且对所有 $i$，
$\overline{V_i} \subset W_i$。于是得到引理所提问题的一个解。

\medskip\noindent
归纳步骤。假设本引理对 $p - 1$ 已证。令
$W_{j_0 \ldots j_{p - 1}}$ 等于所有满足
$\{j_0, \ldots, j_{p - 1}\} = \{i_0, \ldots, i_p\}$（集合相等）的
$W_{i_0 \ldots i_p}$ 之交。由归纳假设，对这些开集存在一个解，记为
$V_i \subset U_i$。由 $W_{j_0 \ldots j_{p - 1}}$ 的选取方式可知，
对所有满足某个 $0 \leq a < b \leq p$ 有 $i_a = i_b$ 的
$(p + 1)$ 元组 $i_0, \ldots, i_p$，都有
$V_{i_0} \cap \ldots \cap V_{i_p} \subset W_{i_0 \ldots i_p}$。
所以，仅当对某个元素两两不同的 $(p + 1)$ 元组
$i_0, \ldots, i_p$ 有
$V_{i_0} \cap \ldots \cap V_{i_p} \not \subset W_{i_0 \ldots i_p}$ 时，
才需修改 $V_i$ 的选取。在这种情形，有
$$
T =
\overline{V_{i_0} \cap \ldots \cap V_{i_p} \setminus W_{i_0 \ldots i_p}}
\subset 
\overline{V_{i_0}} \cap \ldots \cap \overline{V_{i_p}} \setminus
W_{i_0 \ldots i_p}
$$
它是 $X$ 的闭子集，包含于 $U_{i_0} \cap \ldots \cap U_{i_p}$，
并且不与 $Z$ 相交。因此可将 $V_{i_0}$ 替换为
$V_{i_0} \setminus T$ 来“修正”该问题。对每个有问题的元组有限次重复此操作后，
本引理得证。
\end{proof}

\begin{lemma}
\label{topology-lemma-lift-covering-of-quasi-compact-hausdorff-subset}
设 $X$ 为拓扑空间。设 $Z \subset X$ 为拟紧子集，且 $Z$ 的任意两点在
$X$ 中都有不交开邻域。假设给定整数 $p \geq 0$、集合 $I$、对每个
$i \in I$ 给定开集 $U_i \subset X$，并且对 $I$ 的每个
$(p + 1)$ 元组 $i_0, \ldots, i_p$ 给定开集
$W_{i_0 \ldots i_p} \subset U_{i_0} \cap \ldots \cap U_{i_p}$，使得
\begin{enumerate}
\item $Z \subset \bigcup U_i$，并且
\item 对每个 $i_0, \ldots, i_p$，有
$W_{i_0 \ldots i_p} \cap Z = U_{i_0} \cap \ldots \cap U_{i_p} \cap Z$。
\end{enumerate}
则存在 $X$ 的开集 $V_i$，使得
\begin{enumerate}
\item $Z \subset \bigcup V_i$，
\item 对所有 $i$，$V_i \subset U_i$，
\item 对所有 $i$，$\overline{V_i} \cap Z \subset U_i$，并且
\item 对所有 $(p + 1)$ 元组 $i_0, \ldots, i_p$，有
$V_{i_0} \cap \ldots \cap V_{i_p} \subset W_{i_0 \ldots i_p}$。
\end{enumerate}
\end{lemma}

\begin{proof}
因为 $Z$ 拟紧，可以归约到 $I$ 有限的情形（细节从略）。
在 $I$ 有限的情形，关于 $p$ 归纳证明。

\medskip\noindent
基础情形为 $p = 0$。对 $z \in Z \cap U_i$ 与
$z' \in Z \setminus U_i$，存在 $X$ 中不交的开集
$z \in V_{z, z'}$ 与 $z' \in W_{z, z'}$。因为
$Z \setminus U_i$ 拟紧（引理 \ref{topology-lemma-closed-in-quasi-compact}），
可以选取有限多个 $z'_1, \ldots, z'_r$，使得
$Z \setminus U_i \subset W_{z, z'_1} \cup \ldots \cup W_{z, z'_r}$。
于是
$V_z = V_{z, z'_1} \cap \ldots \cap V_{z, z'_r} \cap U_i$
是 $z$ 的一个包含于 $U_i$ 的开邻域，并具有
$\overline{V_z} \cap Z \subset U_i$ 的性质。由于 $z$ 与 $i$ 任意，
并且 $Z$ 拟紧，可以找到有限列表 $z_1, i_1, \ldots, z_t, i_t$ 以及
开集 $V_{z_j} \subset U_{i_j}$，满足
$\overline{V_{z_j}} \cap Z \subset U_{i_j}$ 且
$Z \subset \bigcup V_{z_j}$。于是可令
$V_i = W_i \cap (\bigcup_{j : i = i_j} V_{z_j})$，解决 $p = 0$ 的情形。

\medskip\noindent
归纳步骤。假设本引理对 $p - 1$ 已证。令
$W_{j_0 \ldots j_{p - 1}}$ 等于所有满足
$\{j_0, \ldots, j_{p - 1}\} = \{i_0, \ldots, i_p\}$（集合相等）的
$W_{i_0 \ldots i_p}$ 之交。由归纳假设，对这些开集存在一个解，记为
$V_i \subset U_i$。由 $W_{j_0 \ldots j_{p - 1}}$ 的选取方式可知，
对所有满足某个 $0 \leq a < b \leq p$ 有 $i_a = i_b$ 的
$(p + 1)$ 元组 $i_0, \ldots, i_p$，都有
$V_{i_0} \cap \ldots \cap V_{i_p} \subset W_{i_0 \ldots i_p}$。
所以，仅当对某个元素两两不同的 $(p + 1)$ 元组
$i_0, \ldots, i_p$ 有
$V_{i_0} \cap \ldots \cap V_{i_p} \not \subset W_{i_0 \ldots i_p}$ 时，
才需修改 $V_i$ 的选取。在这种情形，有
$$
T =
\overline{V_{i_0} \cap \ldots \cap V_{i_p} \setminus W_{i_0 \ldots i_p}}
\subset
\overline{V_{i_0}} \cap \ldots \cap \overline{V_{i_p}} \setminus
W_{i_0 \ldots i_p}
$$
由开集 $V_i$ 的性质 (3)，它是 $X$ 的一个不与 $Z$ 相交的闭子集。
因此可将 $V_{i_0}$ 替换为 $V_{i_0} \setminus T$ 来“修正”该问题。
对每个有问题的元组有限次重复此操作后，本引理得证。
\end{proof}

\section{空间的极限}
\label{topology-section-limits}

\noindent
拓扑空间范畴具有积。具体地，若 $I$ 为集合，并且对 $i \in I$ 给定
拓扑空间 $X_i$，就在 $\prod_{i \in I} X_i$ 上赋予{\it 积拓扑}。
拓扑基取为形如 $\prod U_i$ 的集合，其中 $U_i \subset X_i$ 为开集，
且除有限多个 $i$ 外都有 $U_i = X_i$。

\medskip\noindent
拓扑空间范畴具有等化子。具体地，若 $a, b : X \to Y$ 是拓扑空间的态射，
则 $a$ 与 $b$ 的等化子是赋予诱导拓扑的子集
$\{x \in X \mid a(x) = b(x)\} \subset X$。

\begin{lemma}
\label{topology-lemma-limits}
拓扑空间范畴具有极限，而且到集合范畴的遗忘函子与极限交换。
\end{lemma}

\begin{proof}
这由上面的讨论和 Categories, Lemma
\ref{categories-lemma-limits-products-equalizers} 得出。
从上面的描述可知，遗忘函子与极限交换。另一种看法是使用 Categories,
Lemma \ref{categories-lemma-adjoint-exact}，并利用遗忘函子具有左伴随这一事实；
该左伴随就是把一个集合映到相应离散拓扑空间的函子。
\end{proof}

\begin{lemma}
\label{topology-lemma-describe-limits}
设 $\mathcal{I}$ 为余滤范畴。设 $i \mapsto X_i$ 为 $\mathcal{I}$ 上的
拓扑空间图。设 $X = \lim X_i$ 为极限，其投影映射为
$f_i : X \to X_i$。
\begin{enumerate}
\item $X$ 的任意开集都形如 $\bigcup_{j \in J} f_j^{-1}(U_j)$，
其中 $J \subset I$ 为某个子集，且 $U_j \subset X_j$ 为开集。
\item $X$ 的任意拟紧开集都形如 $f_i^{-1}(U_i)$，其中 $i$ 为某个指标，
且 $U_i \subset X_i$ 为某个开集。
\end{enumerate}
\end{lemma}

\begin{proof}
上面给出的极限构造表明，$X \subset \prod X_i$ 并带诱导拓扑。
$\prod X_i$ 的一个拓扑基由开集 $\prod U_i$ 构成，其中
$U_i \subset X_i$ 开，并且除有限多个 $i$ 外都有 $U_i = X_i$。
设 $i_1, \ldots, i_n \in \Ob(\mathcal{I})$ 是满足
$U_{i_j} \not = X_{i_j}$ 的对象。则
$$
X \cap \prod U_i = f_{i_1}^{-1}(U_{i_1}) \cap \ldots \cap f_{i_n}^{-1}(U_{i_n})
$$
对拓扑空间的一般极限而言，这些集合构成 $X$ 上拓扑的一个基。
不过，如果 $\mathcal{I}$ 如引理陈述中那样余滤，就可以选取
$j \in \Ob(\mathcal{I})$ 以及态射 $j \to i_l$（$l = 1, \ldots, n$）。令
$$
U_j =
(X_j \to X_{i_1})^{-1}(U_{i_1}) \cap \ldots \cap
(X_j \to X_{i_n})^{-1}(U_{i_n})
$$
于是显然 $X \cap \prod U_i = f_j^{-1}(U_j)$。因此，对 $X$ 的任意开集
$W$，存在集合 $A$、映射 $\alpha : A \to \Ob(\mathcal{I})$ 以及开集
$U_a \subset X_{\alpha(a)}$，使得
$W = \bigcup f_{\alpha(a)}^{-1}(U_a)$。置 $J = \Im(\alpha)$，并对
$j \in J$ 置 $U_j = \bigcup_{\alpha(a) = j} U_a$，即可看出
$W = \bigcup_{j \in J} f_j^{-1}(U_j)$。这证明 (1)。

\medskip\noindent
为证明 (2)，假设 $\bigcup_{j \in J} f_j^{-1}(U_j)$ 拟紧。
则对某些 $j_1, \ldots, j_m \in J$，它等于
$f_{j_1}^{-1}(U_{j_1}) \cup \ldots \cup f_{j_m}^{-1}(U_{j_m})$。
因为 $\mathcal{I}$ 余滤，可以选取 $i \in \Ob(\mathcal{I})$ 和态射
$i \to j_l$（$l = 1, \ldots, m$）。令
$$
U_i =
(X_i \to X_{j_1})^{-1}(U_{j_1}) \cup \ldots \cup
(X_i \to X_{j_m})^{-1}(U_{j_m})
$$
则按所需，该开集等于 $f_i^{-1}(U_i)$。
\end{proof}

\begin{lemma}
\label{topology-lemma-characterize-limit}
设 $\mathcal{I}$ 为余滤范畴。设 $i \mapsto X_i$ 为 $\mathcal{I}$ 上的
拓扑空间图。设 $X$ 为拓扑空间，并满足
\begin{enumerate}
\item 作为集合，$X = \lim X_i$（用 $f_i$ 表示投影映射），
\item 对 $i \in \Ob(\mathcal{I})$ 以及开集 $U_i \subset X_i$，
所有集合 $f_i^{-1}(U_i)$ 构成 $X$ 上拓扑的一个基。
\end{enumerate}
则作为拓扑空间，$X$ 是各 $X_i$ 的极限。
\end{lemma}

\begin{proof}
由引理 \ref{topology-lemma-describe-limits} 中对极限拓扑的描述得出。
\end{proof}

\begin{theorem}[Tychonov 定理]
\label{topology-theorem-tychonov}
拟紧空间的积拟紧。
\end{theorem}

\begin{proof}
设 $I$ 为集合，并对 $i \in I$ 设 $X_i$ 为拟紧拓扑空间。
置 $X = \prod X_i$。令 $\mathcal{B}$ 为 $X$ 的所有形如
$U_i \times \prod_{i' \in I, i' \not = i} X_{i'}$ 的子集之集合，
其中 $U_i \subset X_i$ 开。按构造，该族是 $X$ 上拓扑的一个子基。
由引理 \ref{topology-lemma-subbase-theorem}，只需证明任何由 $\mathcal{B}$ 中元素
$B_j$ 组成的覆盖 $X = \bigcup_{j \in J} B_j$ 都有有限加细。
可以分解 $J = \coprod J_i$，使得若 $j \in J_i$，则
$B_j = U_j \times \prod_{i' \not = i} X_{i'}$，其中
$U_j \subset X_i$ 开。若 $X_i = \bigcup_{j \in J_i} U_j$，
则存在有限加细，从而 $X = \bigcup_{j \in J} B_j$ 有有限加细。
若非如此，则对每个 $i$，可以选取一点 $x_i \in X_i$，使它不属于
$\bigcup_{j \in J_i} U_j$。但此时点 $x = (x_i)_{i \in I}$ 属于 $X$，
却不属于 $\bigcup_{j \in J} B_j$，矛盾。
\end{proof}

\noindent
若去掉空间 $X_i$ 为 Hausdorff 空间的假设，下面的引理不成立；参见
Examples, Section \ref{examples-section-lim-not-quasi-compact}。

\begin{lemma}
\label{topology-lemma-inverse-limit-quasi-compact}
设 $\mathcal{I}$ 为范畴，$i \mapsto X_i$ 为拓扑空间范畴中
$\mathcal{I}$ 上的图。若每个 $X_i$ 都拟紧且为 Hausdorff 空间，则
$\lim X_i$ 拟紧。
\end{lemma}

\begin{proof}
回忆，$\lim X_i$ 是 $\prod X_i$ 的子空间。由定理
\ref{topology-theorem-tychonov}，该积拟紧。因此只需证明 $\lim X_i$ 是
$\prod X_i$ 的闭子空间（引理 \ref{topology-lemma-closed-in-quasi-compact}）。
若 $\varphi : j \to k$ 是 $\mathcal{I}$ 的一个态射，令
$\Gamma_\varphi \subset X_j \times X_k$ 表示相应连续映射
$X_j \to X_k$ 的图。由引理 \ref{topology-lemma-graph-closed}，该图闭。
显然，$\lim X_i$ 是下列闭子集之交：
$$
\Gamma_\varphi \times \prod\nolimits_{l \not = j, k} X_l
\subset \prod X_i
$$
故结论成立。
\end{proof}

\noindent
下面的引理推广了 Categories, Lemma
\ref{categories-lemma-nonempty-limit}，并部分推广了引理
\ref{topology-lemma-intersection-closed-in-quasi-compact}。

\begin{lemma}
\label{topology-lemma-nonempty-limit}
设 $\mathcal{I}$ 为余滤范畴，$i \mapsto X_i$ 为拓扑空间范畴中
$\mathcal{I}$ 上的图。若每个 $X_i$ 都拟紧、Hausdorff 且非空，
则 $\lim X_i$ 非空。
\end{lemma}

\begin{proof}
在引理 \ref{topology-lemma-inverse-limit-quasi-compact} 的证明中已经看到，
$X = \lim X_i$ 是拟紧空间 $\prod X_i$ 内下列闭子集之交：
$$
Z_\varphi = \Gamma_\varphi \times \prod\nolimits_{l \not = j, k} X_l
$$
其中 $\varphi : j \to k$ 是 $\mathcal{I}$ 的一个态射，而
$\Gamma_\varphi \subset X_j \times X_k$ 是相应态射 $X_j \to X_k$ 的图。
所以由引理 \ref{topology-lemma-intersection-closed-in-quasi-compact}，只需证明这些
子集的任意有限交非空。假设
$\varphi_t : j_t \to k_t$（$t = 1, \ldots, n$）为 $\mathcal{I}$ 的
有限多个态射。因为 $\mathcal{I}$ 余滤，可以选取对象 $j$，并对每个
$t$ 选取态射 $\psi_t : j \to j_t$。对每一对 $t, t'$，如果
(a) $j_t = j_{t'}$，或 (b) $j_t = k_{t'}$，或 (c) $k_t = k_{t'}$，
就得到两个态射 $j \to l$；在情形 (a)、(b) 中 $l = j_t$，而在情形
(c) 中 $l = k_t$。因为 $\mathcal{I}$ 余滤，并且序对 $(t, t')$
只有有限多个，可以选取一个映射 $j' \to j$，使所有这些序对
$(t, t')$ 所对应的两个态射等化。选取元素 $x \in X_{j'}$，并对每个
$t$，令 $x_{j_t}$、相应地令 $x_{k_t}$，分别为 $x$ 在态射
$X_{j'} \to X_j \to X_{j_t}$、相应地在态射
$X_{j'} \to X_j \to X_{j_t} \to X_{k_t}$ 下的像。
对任一指标 $l \in \Ob(\mathcal{I})$，若对所有 $t$ 它都不等于
$j_t$ 或 $k_t$，则任选元素 $x_l \in X_l$（使用选择公理）。于是按构造，
$(x_i)_{i \in \Ob(\mathcal{I})}$ 属于交集
$$
Z_{\varphi_1} \cap \ldots \cap Z_{\varphi_n}
$$
证明完成。
\end{proof}

\section{可构造集}
\label{topology-section-constructible}

\begin{definition}
\label{topology-definition-constructible}
设 $X$ 为拓扑空间，$E \subset X$ 为 $X$ 的一个子集。
\begin{enumerate}
\item 若 $E$ 是有限多个形如 $U \cap V^c$ 的子集之并，其中
$U, V \subset X$ 在 $X$ 中开且逆紧，则称 $E$ 在 $X$ 中
{\it 可构造}\footnote{在 EGA I 第二版 \cite{EGA1-second} 中，
这称为“全局可构造”集，而“可构造”一词用于我们所称的局部可构造集。}。
\item 若存在开覆盖 $X = \bigcup V_i$，使得每个 $E \cap V_i$
在 $V_i$ 中可构造，则称 $E$ 在 $X$ 中{\it 局部可构造}。
\end{enumerate}
\end{definition}

\begin{lemma}
\label{topology-lemma-constructible}
可构造集族对有限交、有限并和补集封闭。
\end{lemma}

\begin{proof}
注意，若 $U_1$、$U_2$ 在 $X$ 中开且逆紧，则 $U_1 \cup U_2$ 也如此，
因为 $X$ 的两个拟紧子集之并拟紧。$U_1 \cap U_2$ 也逆紧。
事实上，假设 $U \subset X$ 是拟紧开集；因为 $U_2$ 在 $X$ 中逆紧，
$U_2 \cap U$ 拟紧；又因为 $U_1$ 在 $X$ 中逆紧，可知
$U_1 \cap (U_2 \cap U)$ 拟紧。由此形式地可证：可构造集的补集可构造，
有限多个可构造集之并可构造，而有限多个可构造集之交可构造。
\end{proof}

\begin{lemma}
\label{topology-lemma-inverse-images-constructibles}
设 $f : X \to Y$ 为拓扑空间间的连续映射。如果 $Y$ 的每个逆紧开子集的逆像
在 $X$ 中逆紧，则可构造集的逆像可构造。
\end{lemma}

\begin{proof}
这是因为 $f^{-1}(U \cap V^c) = f^{-1}(U) \cap f^{-1}(V)^c$，
再结合可构造集的定义即可。
\end{proof}

\begin{lemma}
\label{topology-lemma-open-immersion-constructible-inverse-image}
设 $U \subset X$ 为开集。对可构造集 $E \subset X$，交集
$E \cap U$ 在 $U$ 中可构造。
\end{lemma}

\begin{proof}
假设 $V \subset X$ 在 $X$ 中开且逆紧。由引理
\ref{topology-lemma-inverse-images-constructibles}，只需证明 $V \cap U$ 在 $U$
中逆紧。为此，设 $W \subset U$ 开且拟紧。则 $W$ 在 $X$ 中也开且拟紧。
因为 $V$ 在 $X$ 中逆紧，故 $V \cap W = V \cap U \cap W$ 拟紧。
\end{proof}

\begin{lemma}
\label{topology-lemma-quasi-compact-open-immersion-constructible-image}
设 $U \subset X$ 为逆紧开集，$E \subset U$。若 $E$ 在 $U$ 中可构造，
则 $E$ 在 $X$ 中可构造。
\end{lemma}

\begin{proof}
假设 $V, W \subset U$ 在 $U$ 中开且逆紧。则 $V, W$ 在 $X$ 中开且逆紧
（引理 \ref{topology-lemma-composition-quasi-compact}）。因此
$V \cap (U \setminus W) = V \cap (X \setminus W)$ 在 $X$ 中可构造。
又因 $U$ 的每个可构造子集都是有限多个形如
$V \cap (U \setminus W)$ 的子集之并，结论成立。
\end{proof}

\begin{lemma}
\label{topology-lemma-collate-constructible}
设 $X$ 为拓扑空间，$E \subset X$ 为子集。设
$X = V_1 \cup \ldots \cup V_m$ 为由逆紧开集组成的有限覆盖。
则 $E$ 在 $X$ 中可构造，当且仅当对每个 $j = 1, \ldots, m$，
$E \cap V_j$ 在 $V_j$ 中可构造。
\end{lemma}

\begin{proof}
若 $E$ 在 $X$ 中可构造，则由引理
\ref{topology-lemma-open-immersion-constructible-inverse-image}，对所有 $j$，
$E \cap V_j$ 在 $V_j$ 中可构造。反之，假设对每个
$j = 1, \ldots, m$，$E \cap V_j$ 在 $V_j$ 中可构造。由引理
\ref{topology-lemma-quasi-compact-open-immersion-constructible-image}，
$E = \bigcup E \cap V_j$ 是可构造集的有限并，因而可构造。
\end{proof}

\begin{lemma}
\label{topology-lemma-intersect-constructible-with-closed}
设 $X$ 为拓扑空间。设 $Z \subset X$ 为闭子集，并且
$X \setminus Z$ 拟紧。则对可构造集 $E \subset X$，交集
$E \cap Z$ 在 $Z$ 中可构造。
\end{lemma}

\begin{proof}
假设 $V \subset X$ 在 $X$ 中开且逆紧。由引理
\ref{topology-lemma-inverse-images-constructibles}，只需证明 $V \cap Z$ 在 $Z$
中逆紧。为此，设 $W \subset Z$ 开且拟紧。子集
$W' = W \cup (X \setminus Z)$ 拟紧、开，并且 $W = Z \cap W'$。
于是 $V \cap Z \cap W = V \cap Z \cap W'$ 是拟紧开集
$V \cap W'$ 的闭子集，因为 $V$ 在 $X$ 中逆紧。因此由引理
\ref{topology-lemma-closed-in-quasi-compact}，$V \cap Z \cap W$ 拟紧。
\end{proof}

\begin{lemma}
\label{topology-lemma-intersect-constructible-with-retrocompact}
设 $X$ 为拓扑空间，$T \subset X$ 为子集。假设
\begin{enumerate}
\item $T$ 在 $X$ 中逆紧，
\item 拟紧开集构成 $X$ 上拓扑的一个基。
\end{enumerate}
则对可构造集 $E \subset X$，交集 $E \cap T$ 在 $T$ 中可构造。
\end{lemma}

\begin{proof}
假设 $V \subset X$ 在 $X$ 中开且逆紧。由引理
\ref{topology-lemma-inverse-images-constructibles}，只需证明 $V \cap T$ 在 $T$
中逆紧。为此，设 $W \subset T$ 开且拟紧。由假设 (2)，可以找到
拟紧开集 $W' \subset X$，使得 $W = T \cap W'$（细节从略）。
于是 $V \cap T \cap W = V \cap T \cap W'$ 是 $T$ 与拟紧开集
$V \cap W'$ 之交，因为 $V$ 在 $X$ 中逆紧。因此
$V \cap T \cap W$ 拟紧。
\end{proof}

\begin{lemma}
\label{topology-lemma-closed-constructible-image}
设 $Z \subset X$ 为闭子集，其补集为逆紧开集。设 $E \subset Z$。
若 $E$ 在 $Z$ 中可构造，则 $E$ 在 $X$ 中可构造。
\end{lemma}

\begin{proof}
假设 $V \subset Z$ 在 $Z$ 中开且逆紧。考虑 $X$ 的开子集
$\tilde V = V \cup (X \setminus Z)$。设 $W \subset X$ 开且拟紧。则
$$
W \cap \tilde V =
\left(V \cap W\right) \cup \left((X \setminus Z) \cap W\right).
$$
第一部分拟紧，因为 $V \cap W = V \cap (Z \cap W)$，且
$(Z \cap W)$ 在 $Z$ 中开且拟紧
（引理 \ref{topology-lemma-closed-in-quasi-compact}），而 $V$ 在 $Z$ 中逆紧。
第二部分拟紧，因为 $(X \setminus Z)$ 在 $X$ 中逆紧。由此可知
$\tilde V$ 在 $X$ 中逆紧。因此，若 $V_1, V_2 \subset Z$ 开且逆紧，
则
$$
V_1 \cap (Z \setminus V_2) = \tilde V_1 \cap (X \setminus \tilde V_2)
$$
在 $X$ 中可构造。又因 $Z$ 的每个可构造子集都是有限多个形如
$V_1 \cap (Z \setminus V_2)$ 的子集之并，结论成立。
\end{proof}

\begin{lemma}
\label{topology-lemma-constructible-is-retrocompact}
设 $X$ 为拓扑空间。$X$ 的每个可构造子集都逆紧。
\end{lemma}

\begin{proof}
设 $E = \bigcup_{i = 1, \ldots, n} U_i \cap V_i^c$，其中 $U_i, V_i$
在 $X$ 中开且逆紧。设 $W \subset X$ 为拟紧开集。则
$E \cap W = \bigcup_{i = 1, \ldots, n} U_i \cap V_i^c \cap W$。
所以，只需证明：若 $U, V$ 开且逆紧，而 $W$ 开且拟紧，则
$U \cap V^c \cap W$ 拟紧。这是因为 $U \cap V^c \cap W$ 是拟紧集
$U \cap W$ 的闭子集，故可应用引理
\ref{topology-lemma-closed-in-quasi-compact}。
\end{proof}

\noindent
问题：若仅假设 $X$ 是拟紧拓扑空间，下面的引理是否仍然成立？
与引理 \ref{topology-lemma-intersect-constructible-with-closed} 比较。

\begin{lemma}
\label{topology-lemma-intersect-constructible-with-constructible}
设 $X$ 为拓扑空间。假设 $X$ 有一个由拟紧开集组成的基。
若 $E, E'$ 在 $X$ 中可构造，则交集 $E \cap E'$ 在 $E$ 中可构造。
\end{lemma}

\begin{proof}
合并引理 \ref{topology-lemma-intersect-constructible-with-retrocompact} 与
\ref{topology-lemma-constructible-is-retrocompact}。
\end{proof}

\begin{lemma}
\label{topology-lemma-constructible-in-constructible}
设 $X$ 为拓扑空间。假设 $X$ 有一个由拟紧开集组成的基。
设 $E$ 在 $X$ 中可构造，而 $F \subset E$ 在 $E$ 中可构造。
则 $F$ 在 $X$ 中可构造。
\end{lemma}

\begin{proof}
注意，$X$ 的任意逆紧子集 $T$ 的诱导拓扑都有一个由拟紧开集组成的基。
特别地，这对任意可构造子集成立
（引理 \ref{topology-lemma-constructible-is-retrocompact}）。写成
$E = E_1 \cup \ldots \cup E_n$，其中 $E_i = U_i \cap V_i^c$，而
$U_i, V_i \subset X$ 开且逆紧。注意，由引理
\ref{topology-lemma-intersect-constructible-with-constructible}，
$E_i = E \cap E_i$ 在 $E$ 中可构造。再由引理
\ref{topology-lemma-intersect-constructible-with-constructible}，
$F \cap E_i$ 在 $E_i$ 中可构造。所以，只需证明
$E = U \cap V^c$ 的情形，其中 $U, V \subset X$ 开且逆紧。
在这种情形，包含映射 $E \subset X$ 是复合
$$
E = U \cap V^c \to U \to X
$$
然后可将引理 \ref{topology-lemma-closed-constructible-image} 应用于第一个包含，
将引理 \ref{topology-lemma-quasi-compact-open-immersion-constructible-image}
应用于第二个包含。
\end{proof}

\begin{lemma}
\label{topology-lemma-locally-closed-constructible-image}
设 $X$ 为拟紧拓扑空间，它有一个由拟紧开集组成的基，并且任意两个拟紧开集
之交拟紧。设 $T \subset X$ 为局部闭子集，且 $T$ 拟紧，而 $T^c$
在 $X$ 中逆紧。则 $T$ 在 $X$ 中可构造。
\end{lemma}

\begin{proof}
注意，$T$ 拟紧且在 $\overline{T}$ 中开。利用由拟紧开集组成的基，
可以写成 $T = U \cap \overline{T}$，其中 $U$ 是 $X$ 中的拟紧开集。
则 $V = U \setminus T = U \cap T^c$ 在 $U$ 中逆紧，因为 $T^c$
在 $X$ 中逆紧。所以 $V$ 拟紧。由于任意两个拟紧开集之交拟紧，
$X$ 的任意拟紧开集都逆紧。因此，$T = U \cap V^c$，其中 $U$
与 $V = U \setminus T$ 都是 $X$ 的逆紧开集。更不必说，$T$ 在 $X$
中可构造。
\end{proof}

\begin{lemma}
\label{topology-lemma-collate-constructible-from-constructible}
设 $X$ 为拓扑空间，并且有一个由拟紧开集组成的拓扑基。
设 $E \subset X$ 为子集。设
$X = E_1 \cup \ldots \cup E_m$ 为由可构造子集组成的有限覆盖。
则 $E$ 在 $X$ 中可构造，当且仅当对每个 $j = 1, \ldots, m$，
$E \cap E_j$ 在 $E_j$ 中可构造。
\end{lemma}

\begin{proof}
合并引理 \ref{topology-lemma-intersect-constructible-with-constructible} 与
\ref{topology-lemma-constructible-in-constructible}。
\end{proof}

\begin{lemma}
\label{topology-lemma-generic-point-in-constructible}
设 $X$ 为拓扑空间。假设 $Z \subset X$ 不可约。设 $E \subset X$
为有限多个局部闭子集之并（例如，$E$ 可构造）。下列条件等价：
\begin{enumerate}
\item 交集 $E \cap Z$ 包含 $Z$ 的一个稠密开子集。
\item 交集 $E \cap Z$ 在 $Z$ 中稠密。
\end{enumerate}
若 $Z$ 有泛点 $\xi$，则这些条件还等价于
\begin{enumerate}
\item[(3)] 有 $\xi \in E$。
\end{enumerate}
\end{lemma}

\begin{proof}
蕴含关系 (1) $\Rightarrow$ (2) 显然。假设 (2)。注意，$E \cap Z$
是 $Z$ 的有限多个局部闭子集 $Z_i$ 之并。由于 $Z$ 不可约，必有某个
$Z_i$ 在 $Z$ 中稠密。因为该 $Z_i$ 在其闭包中开，所以它在 $Z$ 中
稠密且开。因此 (1) 成立。

\medskip\noindent
假设 $\xi \in Z$ 为泛点。若等价条件 (1)、(2) 成立，则
$\xi \in E$。反之，若 $\xi \in E$，则 $\xi \in E \cap Z$，
从而 $E \cap Z$ 在 $Z$ 中稠密。
\end{proof}

\section{可构造集与 Noether 空间}
\label{topology-section-constructible-Noetherian}

\begin{lemma}
\label{topology-lemma-constructible-Noetherian-space}
设 $X$ 为 Noether 拓扑空间。$X$ 中的可构造集恰为 $X$ 的有限多个
局部闭子集之并。
\end{lemma}

\begin{proof}
这立即由引理 \ref{topology-lemma-Noetherian-quasi-compact} 得出。
\end{proof}

\begin{lemma}
\label{topology-lemma-constructible-map-Noetherian}
设 $f : X \to Y$ 为 Noether 拓扑空间之间的连续映射。
若 $E \subset Y$ 在 $Y$ 中可构造，则 $f^{-1}(E)$ 在 $X$ 中可构造。
\end{lemma}

\begin{proof}
这立即由引理 \ref{topology-lemma-constructible-Noetherian-space} 和连续映射的定义得出。
\end{proof}

\begin{lemma}
\label{topology-lemma-characterize-constructible-Noetherian}
设 $X$ 为 Noether 拓扑空间，$E \subset X$ 为子集。下列条件等价：
\begin{enumerate}
\item $E$ 在 $X$ 中可构造；
\item 对每个不可约闭子集 $Z \subset X$，交集 $E \cap Z$ 或者包含
$Z$ 的一个非空开集，或者在 $Z$ 中不稠密。
\end{enumerate}
\end{lemma}

\begin{proof}
假设 $E$ 可构造，且 $Z \subset X$ 为不可约闭子集。由引理
\ref{topology-lemma-constructible-map-Noetherian}，$E \cap Z$ 在 $Z$ 中可构造。
因此，$E \cap Z$ 是 $Z$ 的有限多个非空局部闭子集 $T_i$ 之并。
显然，若所有 $T_i$ 都不在 $Z$ 中开，则 $E \cap Z$ 在 $Z$ 中不稠密。
由此可见 (1) 蕴含 (2)。

\medskip\noindent
反之，假设 (2) 成立。考虑由 $X$ 的所有闭子集 $Y$ 组成的集合
$\mathcal{S}$，其中 $E \cap Y$ 在 $Y$ 中不可构造。若
$\mathcal{S} \not = \emptyset$，则由于 $X$ 是 Noether 的，
它有一个极小元 $Y$。设
$Y = Y_1 \cup \ldots \cup Y_r$ 是 $Y$ 的不可约分支分解，参见引理
\ref{topology-lemma-Noetherian}。若 $r > 1$，则每个 $Y_i \cap E$ 在 $Y_i$
中可构造，因而是 $Y_i$ 的有限多个局部闭子集之并。于是
$E \cap Y$ 也是 $Y$ 的有限多个局部闭子集之并；由引理
\ref{topology-lemma-constructible-Noetherian-space}，$E \cap Y$ 在 $Y$ 中可构造。
这产生矛盾，故 $r = 1$。若 $r = 1$，则 $Y$ 不可约；由假设 (2)，
$E \cap Y$ 或者 (a) 包含 $Y$ 的一个开集 $V$，或者 (b) 在 $Y$ 中
不稠密。在情形 (a)，由 $Y$ 的极小性，$E \cap (Y \setminus V)$
是 $Y \setminus V$ 的有限多个局部闭子集之并。因此，$E \cap Y$
是 $Y$ 的有限多个局部闭子集之并，并由引理
\ref{topology-lemma-constructible-Noetherian-space} 可知其可构造。这产生矛盾，
所以必处于情形 (b)。在情形 (b)，存在真闭子集 $Y' \subset Y$，使得
$E \cap Y = E \cap Y'$。由 $Y$ 的极小性，$E \cap Y'$ 是 $Y'$
的有限多个局部闭子集之并；由此可见，$E \cap Y' = E \cap Y$
是 $Y$ 的有限多个局部闭子集之并，并由引理
\ref{topology-lemma-constructible-Noetherian-space} 可知其可构造。
这个矛盾完成了引理的证明。
\end{proof}

\begin{lemma}
\label{topology-lemma-constructible-neighbourhood-Noetherian}
设 $X$ 为 Noether 拓扑空间，$x \in X$，且 $E \subset X$ 在 $X$
中可构造。下列条件等价：
\begin{enumerate}
\item $E$ 是 $x$ 的邻域；
\item 对 $X$ 的每个包含 $x$ 的不可约闭子集 $Y$，交集 $E \cap Y$ 在 $Y$
中稠密。
\end{enumerate}
\end{lemma}

\begin{proof}
(1) 蕴含 (2) 是清楚的。假设 (2)。考虑由 $X$ 的所有包含 $x$
的闭子集 $Y$ 组成的集合 $\mathcal{S}$，其中 $E \cap Y$ 不是 $Y$
中 $x$ 的邻域。若 $\mathcal{S} \not = \emptyset$，则由于 $X$
是 Noether 的，它有一个极小元 $Y$。假设
$Y = Y_1 \cup Y_2$，其中 $Y_1$、$Y_2$ 是两个更小的非空闭子集。
若对 $i = 1, 2$ 都有 $x \in Y_i$，则 $Y_i \cap E$ 是 $Y_i$
中 $x$ 的邻域，从而 $Y \cap E$ 是 $Y$ 中 $x$ 的邻域，矛盾。
若 $x \in Y_1$ 但 $x \not\in Y_2$（不妨如此），则
$Y_1 \cap E$ 是 $Y_1$ 中 $x$ 的邻域，因而也是 $Y$ 中 $x$ 的邻域，
这同样产生矛盾。因此，$Y$ 是不可约闭集。由假设 (2)，$E \cap Y$
在 $Y$ 中稠密。于是 $E \cap Y$ 包含 $Y$ 的一个开集 $V$，参见引理
\ref{topology-lemma-characterize-constructible-Noetherian}。若 $x \in V$，则
$E \cap Y$ 是 $Y$ 中 $x$ 的邻域，矛盾。若 $x \not \in V$，则
$Y' = Y \setminus V$ 是包含 $x$ 的 $Y$ 的真闭子集。由 $Y$ 的极小性，
$E \cap Y'$ 包含 $Y'$ 中 $x$ 的一个开邻域 $V' \subset Y'$。
但此时 $V' \cup V$ 是 $Y$ 中包含于 $E$ 的 $x$ 的开邻域，矛盾。
这个矛盾完成了引理的证明。
\end{proof}

\begin{lemma}
\label{topology-lemma-characterize-open-Noetherian}
设 $X$ 为 Noether 拓扑空间，$E \subset X$ 为子集。下列条件等价：
\begin{enumerate}
\item $E$ 在 $X$ 中开；
\item 对每个不可约闭子集 $Y$，交集 $E \cap Y$ 或者为空，或者包含
$Y$ 的一个非空开集。
\end{enumerate}
\end{lemma}

\begin{proof}
这形式地由引理 \ref{topology-lemma-characterize-constructible-Noetherian} 与
\ref{topology-lemma-constructible-neighbourhood-Noetherian} 得出。
\end{proof}

\section{固有映射的刻画}
\label{topology-section-proper}

\noindent
本节讨论通常拓扑学中的固有映射概念。我们称拓扑空间的连续映射为
{\it 固有的}，如果它是普遍闭且分离的。虽然这与代数几何中固有态射的
定义很好地吻合，但它不同于 Bourbaki 的定义。采用我们对拓扑空间固有映射
的定义，固有基变换定理（Cohomology，定理
\ref{cohomology-theorem-proper-base-change}）无需任何进一步假设即成立。
此外，给定 $\mathbf{C}$ 上有限型概形的态射 $f : X \to Y$，有如下结论：
$f$ 作为概形态射是固有的，当且仅当在赋予经典拓扑的
$\mathbf{C}$-点上，连续映射 $f : X(\mathbf{C}) \to Y(\mathbf{C})$
是固有的。这在 \cite[Exp. XII, Prop.\ 3.2(v)]{SGA1} 中得到说明；
其中还有一个脚注指出，他们在拓扑学中采用的是附加了分离性的 Bourbaki
固有性概念。

\medskip\noindent
我们发现，为三个不同的概念分别命名很有用：分离、普遍闭，以及二者兼具
（即固有性）。对于局部紧 Hausdorff 空间的连续映射 $f : X \to Y$，
“固有”一词长期以来用于表示“$f^{-1}($紧集$) =$ 紧集”这一概念；
对于这样良好的空间，它等价于普遍闭。事实上，我们将会看到，若为清楚起见
用“拟紧”来表述原像条件，并加入映射为闭映射这一假设，则该条件一般地等价于
普遍闭。另见 \cite[Exercises 22-26 in Chapter II]{LangReal}，但须注意，
Lang 与 Bourbaki 一样，把“固有”用作“普遍闭”的同义词。

\begin{lemma}[管状引理]
\label{topology-lemma-tube}
设 $X$ 与 $Y$ 为拓扑空间，$A \subset X$ 与 $B \subset Y$ 为拟紧子集。
设 $A \times B \subset W \subset X \times Y$，其中 $W$ 在
$X \times Y$ 中开。则存在开集 $A \subset U \subset X$ 与
$B \subset V \subset Y$，使得 $U \times V \subset W$。
\end{lemma}

\begin{proof}
对每个 $a \in A$ 与 $b \in B$，存在 $X$ 的开集 $U_{(a, b)}$
和 $Y$ 的开集 $V_{(a, b)}$，使得
$(a, b) \in U_{(a, b)} \times V_{(a, b)} \subset W$。
固定 $b$，则存在有限多个 $a_1, \ldots, a_n$，使得
$A \subset U_{(a_1, b)} \cup \ldots \cup U_{(a_n, b)}$。因此
$$
A \times \{b\} \subset
(U_{(a_1, b)} \cup \ldots \cup U_{(a_n, b)}) \times
(V_{(a_1, b)} \cap \ldots \cap V_{(a_n, b)}) \subset W.
$$
所以，对每个 $b \in B$，存在开集 $U_b \subset X$ 与
$V_b \subset Y$，使得 $A \times \{b\} \subset U_b \times V_b \subset W$。
同上，存在有限多个 $b_1, \ldots, b_m$，使得
$B \subset V_{b_1} \cup \ldots \cup V_{b_m}$。此时结论成立，因为
$A \times B \subset
(U_{b_1} \cap \ldots \cap U_{b_m}) \times
(V_{b_1} \cup \ldots \cup V_{b_m})$。
\end{proof}

\noindent
以下定义中的术语可能与你习惯的略有不同。

\begin{definition}
\label{topology-definition-proper-map}
设 $f : X\to Y$ 为拓扑空间之间的连续映射。
\begin{enumerate}
\item 若每个闭子集的像都是闭的，则称映射 $f$ 为{\it 闭映射}。
\item 若对任意拓扑空间 $Z$，映射 $Z \times X\to Z \times Y$ 都是闭映射，
则称映射 $f$ 为 {\it Bourbaki-固有的}\footnote{这是 \cite{Bourbaki}
采用的术语。在文献中，这一性质有时也称为“普遍闭”。}。
\item 若每个拟紧子集 $V \subset Y$ 的原像 $f^{-1}(V)$ 都拟紧，
则称映射 $f$ 为{\it 拟固有的}。
\item 若对任意连续映射 $g: Z \to Y$，映射
$f': Z \times_Y X \to Z$ 都是闭映射，则称 $f$ 为{\it 普遍闭的}。
\item 若 $f$ 分离且普遍闭，则称 $f$ 为{\it 固有的}。
\end{enumerate}
\end{definition}

\noindent
下面的引理以后会用到。

\begin{lemma}
\label{topology-lemma-characterize-quasi-compact}
\begin{reference}
这是 \cite[I, p. 75, Lemme 1]{Bourbaki} 与
\cite[I, p. 76, Corrolaire 1]{Bourbaki} 的结合。
\end{reference}
拓扑空间 $X$ 拟紧，当且仅当对任意拓扑空间 $Z$，投影映射
$Z \times X \to Z$ 都是闭映射。
\end{lemma}

\begin{proof}
（另见下面的评注。）若 $X$ 不拟紧，则存在开覆盖
$X = \bigcup_{i \in I} U_i$，使得任意有限多个 $U_i$ 都不能覆盖 $X$。
令 $Z$ 为 $I$ 的幂集 $\mathcal{P}(I)$ 的子集，由 $I$ 本身以及 $I$
的所有非空有限子集组成。在 $Z$ 上定义拓扑，并以下列集合为拓扑基：
\begin{enumerate}
\item $Z\setminus\{I\}$ 的所有子集。
\item 对 $I$ 的每个有限子集 $K$，取集合
% P01-ERRATA-E0018: frozen source has an unmatched closing parenthesis in this set-builder formula; preserved here.
$U_K := \{J\subset I \mid J \in Z, \ K\subset J \})$。
\end{enumerate}
留给读者验证这确实是一个拓扑基。考虑由下式定义的 $Z \times X$ 的子集
$$
M = \{(J, x) \mid J \in Z, \ x \in \bigcap\nolimits_{i \in J} U_i^c)\}
$$
若 $(J, x) \not \in M$，则对某个 $i \in J$ 有 $x \in U_i$。
因此，$U_{\{i\}} \times U_i \subset Z \times X$ 是一个包含
$(J, x)$ 且不与 $M$ 相交的开子集。所以 $M$ 是闭的。
$M$ 到 $Z$ 的投影是 $Z-\{I\}$，它不是闭的。因此
$Z \times X \to Z$ 不是闭映射。

\medskip\noindent
假设 $X$ 拟紧。设 $Z$ 为拓扑空间，$M \subset  Z \times X$ 为闭子集。
设 $z \in Z$ 是一个不属于 $\text{pr}_1(M)$ 的点。由管状引理
\ref{topology-lemma-tube}，存在开集 $U \subset Z$，使得 $U \times X$ 包含于
$M$ 的补集。因此，$\text{pr}_1(M)$ 是闭的。
\end{proof}

\begin{remark}
\label{topology-remark-lemma-literature}
引理 \ref{topology-lemma-characterize-quasi-compact} 是
\cite[I, p. 75, Lemme 1]{Bourbaki} 与
\cite[I, p. 76, Corollaire 1]{Bourbaki} 的结合。
\end{remark}

\begin{theorem}
\label{topology-theorem-characterize-proper}
\begin{reference}
在 \cite[I, p. 75, Theorem 1]{Bourbaki} 中可找到：
(2) $\Leftrightarrow$ (4)。
在 \cite[I, p. 77, Proposition 6]{Bourbaki} 中可找到：
(2) $\Rightarrow$ (1)。
\end{reference}
设 $f: X\to Y$ 为拓扑空间之间的连续映射。下列条件等价：
\begin{enumerate}
\item 映射 $f$ 拟固有且为闭映射。
\item 映射 $f$ 是 Bourbaki-固有的。
\item 映射 $f$ 普遍闭。
\item 映射 $f$ 是闭映射，并且对任意 $y\in Y$，$f^{-1}(y)$ 拟紧。
\end{enumerate}
\end{theorem}

\begin{proof}
（另见下面的评注。）若映射 $f$ 满足 (1)，则它自动满足 (4)，因为单点
总是拟紧的。

\medskip\noindent
假设映射 $f$ 满足 (4)。我们将证明它普遍闭，即 (3) 成立。
设 $g : Z \to Y$ 为拓扑空间的连续映射，并考虑图表
$$
\xymatrix{
Z \times_Y X \ar[r]_{g'} \ar[d]_{f'} & X \ar[d]^f \\
Z \ar[r]^g & Y
}
$$
在证明中，我们会用到 $Z \times_Y X \to Z \times X$ 是到其像上的同胚；
也就是说，可将 $Z \times_Y X$ 与 $Z \times X$ 中相应的子集按诱导拓扑
等同起来。$f' : Z \times_Y X \to Z$ 的像是
$\Im(f') = \{z : g(z) \in f(X)\}$。因为 $f(X)$ 是闭的，所以
$\Im(f')$ 是 $Z$ 的闭子空间。考虑闭子集
$P \subset Z \times_Y X$。设 $z \in Z$，且 $z \not \in f'(P)$。
若 $z \not \in \Im(f')$，则 $Z \setminus \Im(f')$ 是避开 $f'(P)$
的一个开邻域。若 $z$ 属于 $\Im(f')$，则
$(f')^{-1}\{z\} = \{z\} \times f^{-1}\{g(z)\}$，并且由假设，
$f^{-1}\{g(z)\}$ 拟紧。因为 $P$ 是 $Z \times_Y X$ 的闭子集，
存在 $Z \times X$ 的闭子集 $P'$，使得
$P = P' \cap Z \times_Y X$。由于 $(f')^{-1}\{z\}$ 是
$P^c = P'^c \cup (Z \times_Y X)^c$ 的子集，并且
$(f')^{-1}\{z\}$ 与 $(Z \times_Y X)^c$ 不相交，故
$(f')^{-1}\{z\}$ 包含于 $P'^c$。可将管状引理 \ref{topology-lemma-tube}
应用于
$(f')^{-1}\{z\} = \{z\} \times f^{-1}\{g(z)\}
\subset (P')^c \subset Z \times X$。由此得到包含
$(f')^{-1}\{z\}$ 的 $V \times U$，其中 $U$ 和 $V$ 分别是 $X$ 与
$Z$ 中的开集，并且 $V \times U$ 与 $P'$ 的交为空。于是集合
$V \cap g^{-1}(Y-f(U^c))$ 在 $Z$ 中开，因为 $f$ 是闭映射；它包含
$z$，并且与 $P$ 的像不相交。因此 $f'(P)$ 是闭的。换言之，映射 $f$
普遍闭。

\medskip\noindent
蕴含关系 (3) $\Rightarrow$ (2) 是平凡的。事实上，给定任意拓扑空间
$Z$，考虑投影态射 $g : Z \times Y \to Y$。此时容易看出，$f'$ 就是映射
$Z \times X \to Z \times Y$；换言之，
$(Z \times Y) \times_Y X = Z \times X$。（这一等同纯粹是范畴性质，
本身与拓扑空间无关。）

\medskip\noindent
假设 $f$ 满足 (2)。我们将证明它满足 (1)。注意，$f$ 是闭映射，因为
$f$ 可与映射 $\{pt\} \times X \to \{pt\} \times Y$ 等同，而后者依假设
是闭映射。任取拟紧子集 $K \subset Y$。设 $Z$ 为任意拓扑空间。
因为 $Z \times X \to Z \times Y$ 是闭映射，所以映射
$Z \times f^{-1}(K) \to Z \times K$ 是闭映射（若 $T$ 在
$Z \times f^{-1}(K)$ 中闭，可将其写成
$T = Z \times f^{-1}(K) \cap T'$，其中 $T' \subset Z \times X$
闭）。因为 $K$ 拟紧，由引理 \ref{topology-lemma-characterize-quasi-compact}，
$K \times Z\to Z$ 是闭映射。因此复合
$Z \times f^{-1}(K)\to Z \times K \to Z$ 是闭映射；再次由引理
\ref{topology-lemma-characterize-quasi-compact}，$f^{-1}(K)$ 必拟紧。
\end{proof}

\begin{remark}
\label{topology-remark-proof-literature}
下面给出一些文献出处。在 \cite[I, p. 75, Theorem 1]{Bourbaki} 中可找到：
(2) $\Leftrightarrow$ (4)。在 \cite[I, p. 77, Proposition 6]{Bourbaki}
中可找到：(2) $\Rightarrow$ (1)。当然，平凡地有
(1) $\Rightarrow$ (4)。所以 (1)、(2) 与 (4) 等价。
(3) 与 (4) 的等价性见 \cite[Chapter II, Exercise 25]{LangReal}。
\end{remark}

\begin{lemma}
\label{topology-lemma-closed-map}
\begin{slogan}
从紧空间到 Hausdorff 空间的映射是普遍闭的。
\end{slogan}
设 $f : X \to Y$ 为拓扑空间的连续映射。若 $X$ 拟紧且 $Y$ 为
Hausdorff 空间，则 $f$ 普遍闭。
\end{lemma}

\begin{proof}
由于 $Y$ 的每个点都是闭点，由引理 \ref{topology-lemma-closed-in-quasi-compact}
可知，对所有 $y \in Y$，$X$ 的闭子集 $f^{-1}(y)$ 都拟紧。因此，
由定理 \ref{topology-theorem-characterize-proper}，只需证明 $f$ 是闭映射。
若 $E \subset X$ 闭，则它拟紧（引理 \ref{topology-lemma-closed-in-quasi-compact}），
所以 $f(E) \subset Y$ 拟紧（引理 \ref{topology-lemma-image-quasi-compact}），
进而 $f(E)$ 在 $Y$ 中闭（引理 \ref{topology-lemma-quasi-compact-in-Hausdorff}）。
\end{proof}

\begin{lemma}
\label{topology-lemma-bijective-map}
设 $f : X \to Y$ 为拓扑空间的连续映射。若 $f$ 是双射，$X$ 拟紧，
且 $Y$ 为 Hausdorff 空间，则 $f$ 是同胚。
\end{lemma}

\begin{proof}
只需证明 $f$ 是闭映射，因为这蕴含 $f^{-1}$ 连续。若 $T \subset X$
闭，则由引理 \ref{topology-lemma-closed-in-quasi-compact}，$T$ 拟紧；因此由引理
\ref{topology-lemma-image-quasi-compact}，$f(T)$ 拟紧；进而由引理
\ref{topology-lemma-quasi-compact-in-Hausdorff}，$f(T)$ 闭。
\end{proof}

\section{Jacobson 空间}
\label{topology-section-space-jacobson}

\begin{definition}
\label{topology-definition-space-jacobson}
设 $X$ 为拓扑空间，$X_0$ 为 $X$ 的闭点集。若每个闭子集
$Z \subset X$ 都是 $Z \cap X_0$ 的闭包，则称 $X$ 是
{\it Jacobson 空间}。
\end{definition}

\noindent
注意，拓扑空间 $X$ 是 Jacobson 空间，当且仅当 $X$ 的每个非空局部闭
子集都有一个在 $X$ 中闭的点。

\medskip\noindent
设 $X$ 为 Jacobson 空间，并设 $X_0$ 为赋予诱导拓扑的 $X$ 的闭点集。
显然，定义蕴含态射 $X_0 \to X$ 在 $X_0$ 的闭子集与 $X$ 的闭子集之间
诱导双射。因此，$X$ 的许多性质由 $X_0$ 继承。例如，$X$ 与 $X_0$
的 Krull 维数相同。

\begin{lemma}
\label{topology-lemma-jacobson-check-irreducible-closed}
设 $X$ 为拓扑空间，$X_0$ 为 $X$ 的闭点集。假设对每个点 $x\in X$，
交集 $X_0 \cap \overline{\{x\}}$ 在 $\overline{\{x\}}$ 中稠密。
则 $X$ 是 Jacobson 空间。
\end{lemma}

\begin{proof}
设 $Z$ 为 $X$ 的闭子集，$U$ 为 $X$ 的开子集，并且 $U\cap Z$ 非空。
于是，对 $x\in U\cap Z$，$\overline{\{x\}}\cap U$ 是 $Z\cap U$
的非空子集；由假设，它包含一个在 $X$ 中闭的点，正合所需。
\end{proof}

\begin{lemma}
\label{topology-lemma-non-jacobson-Noetherian-characterize}
设 $X$ 为具有拟紧开集基的 Kolmogorov 拓扑空间。若 $X$ 不是
Jacobson 空间，则存在非闭点 $x \in X$，使得 $\{x\}$ 局部闭。
\end{lemma}

\begin{proof}
因为 $X$ 不是 Jacobson 空间，存在 $X$ 中的闭集 $Z$ 与开集 $U$，
使得 $Z \cap U$ 非空且不包含在 $X$ 中闭的点。由于 $X$ 有拟紧开集基，
可将 $U$ 替换为 $Z\cap U$ 中某点的拟紧开邻域；故可假设 $U$ 是拟紧
开集。由引理 \ref{topology-lemma-quasi-compact-closed-point}，存在点
$x \in Z \cap U$，它在 $Z \cap U$ 中闭。因此 $\{x\}$ 局部闭，
但在 $X$ 中不闭。
\end{proof}

\begin{lemma}
\label{topology-lemma-jacobson-local}
设 $X$ 为拓扑空间，$X = \bigcup U_i$ 为一个开覆盖。则 $X$ 是
Jacobson 空间，当且仅当每个 $U_i$ 都是 Jacobson 空间。此外，在这种
情形有 $X_0 = \bigcup U_{i, 0}$。
\end{lemma}

\begin{proof}
设 $X$ 为拓扑空间，$X_0$ 为 $X$ 的闭点集，$U_{i, 0}$ 为 $U_i$
的闭点集。于是 $X_0 \cap U_i \subset U_{i, 0}$，但一般未必相等。

\medskip\noindent
先假设每个 $U_i$ 都是 Jacobson 空间。我们断言在这种情形
$X_0 \cap U_i = U_{i, 0}$。事实上，假设 $x \in U_{i, 0}$，即 $x$
在 $U_i$ 中闭。令 $\overline{\{x\}}$ 为其在 $X$ 中的闭包，并考虑
$\overline{\{x\}} \cap U_j$。若 $x \not \in U_j$，则
$\overline{\{x\}} \cap U_j = \emptyset$。若 $x \in U_j$，则
$U_i \cap U_j \subset U_j$ 是 $U_j$ 中包含 $x$ 的开子集。令
$T' = U_j \setminus U_i \cap U_j$ 且 $T = \{x\} \amalg T'$。
由于 $x$ 在 $U_i \cap U_j$ 中闭（因为 $x$ 在 $U_i$ 中闭），可知
$T$、$T'$ 都是 $U_j$ 的闭子集，且 $T$ 包含 $x$。因为 $U_j$ 是
Jacobson 空间，$U_j$ 的闭点在 $T$ 中稠密。鉴于
$T = \{x\} \amalg T'$，这只有在 $x$ 于 $U_j$ 中闭时才可能成立。
因此 $\overline{\{x\}} \cap U_j = \{x\}$。由此得出
$\overline{\{x\}} = \{ x \}$，正合所需。

\medskip\noindent
设 $Z \subset X$ 为闭子集（仍假设每个 $U_i$ 都是 Jacobson 空间）。
既然现在知道 $X_0 \cap Z  \cap U_i
= U_{i, 0} \cap Z$ 在 $Z \cap U_i$ 中稠密，便立即得出
$X_0 \cap Z$ 在 $Z$ 中稠密。

\medskip\noindent
反之，假设 $X$ 是 Jacobson 空间。设 $Z \subset U_i$ 为闭子集。
则 $X_0 \cap \overline{Z}$ 在 $\overline{Z}$ 中稠密。由于
$\overline{Z} \setminus Z$ 是闭的，$X_0 \cap Z$ 也在 $Z$ 中稠密。
由 $X_0 \cap U_i \subset U_{i, 0}$ 可知，$U_{i, 0} \cap Z$ 在 $Z$
中稠密。因此 $U_i$ 是 Jacobson 空间，正合所需。
\end{proof}

\begin{lemma}
\label{topology-lemma-jacobson-inherited}
设 $X$ 为 Jacobson 空间。下列类型的子集 $T \subset X$ 都是
Jacobson 空间：
\begin{enumerate}
\item 开子空间。
\item 闭子空间。
\item 局部闭子空间。
\item 局部闭子空间之并。
\item 可构造集。
\item 任意在 $X$ 上局部地为局部闭子集之并的子集 $T \subset X$。
\end{enumerate}
在每种情形，$T$ 的闭点都在 $X$ 中闭。
\end{lemma}

\begin{proof}
设 $X_0$ 为 $X$ 的闭点集。对任意子集 $T \subset X$，以 $(*)$
表示下列性质：
\begin{itemize}
\item[(*)] $T$ 的每个非空局部闭子集都有一个在 $X$ 中闭的点。
\end{itemize}
注意，总有 $X_0 \cap T \subset T_0$。因此性质 $(*)$ 蕴含 $T$
是 Jacobson 空间。此外，它显然蕴含 $T$ 的每个闭点都在 $X$ 中闭。

\medskip\noindent
假设 $T=\bigcup_i T_i$，其中 $T_i$ 在 $X$ 中局部闭。任取 $T$
中的非空局部闭子集 $A\subset T$，则存在某个 $T_i$，使得
$A\cap T_i$ 非空；它在 $T_i$ 中局部闭，从而也在 $X$ 中局部闭。
因为 $X$ 是 Jacobson 空间，$A$ 有一个在 $X$ 中闭的点。
\end{proof}

\begin{lemma}
\label{topology-lemma-finite-jacobson}
有限 Jacobson 空间是离散的。只有有限多个闭点的 Jacobson 空间是离散的。
\end{lemma}

\begin{proof}
若 $X$ 有限，则闭点集 $X_0 \subset X$ 有限。假设 $X_0$ 有限且
$X$ 是 Jacobson 空间。由 Jacobson 性质，$\overline{X_0} = X$。
现在
$X_0 =\{x_1, \ldots, x_n\} = \bigcup_{i = 1, \ldots, n} \{x_i\}$
是有限多个闭集之并，因而闭；所以 $X = \overline{X_0} = X_0$。
每个点都闭，并且由有限性，每个点都开。
\end{proof}

\begin{lemma}
\label{topology-lemma-jacobson-equivalent-locally-closed}
\begin{slogan}
对 Jacobson 空间而言，闭点反映拓扑的全部信息。
\end{slogan}
假设 $X$ 为 Jacobson 拓扑空间，并设 $X_0$ 为 $X$ 的闭点集。
存在一个保持包含关系的双射对应
$$
\{\text{局部闭子集的有限并，位于 } X\}
\leftrightarrow
\{\text{局部闭子集的有限并，位于 } X_0\}
$$
由 $E \mapsto E \cap X_0$ 给出。这个对应保持其中的局部闭子集、开子集
与闭子集。
\end{lemma}

\begin{proof}
只需证明对应 $E \mapsto E \cap X_0$ 是单射。事实上，若 $E\neq E'$，
则不妨设 $E\setminus E'$ 非空；它是有限多个局部闭集之并（略去细节）。
由于 $X$ 是 Jacobson 空间，可知
$(E \setminus E') \cap X_0 = E \cap X_0 \setminus E' \cap X_0$ 非空。
\end{proof}

\begin{lemma}
\label{topology-lemma-jacobson-equivalent-constructible}
假设 $X$ 为 Jacobson 拓扑空间，并设 $X_0$ 为 $X$ 的闭点集。
存在一个保持包含关系的双射对应
$$
\{\text{下列对象的可构造子集： } X\}
\leftrightarrow
\{\text{下列对象的可构造子集： } X_0\}
$$
由 $E \mapsto E \cap X_0$ 给出。这个对应保持其中的逆紧开子集以及
这些子集的补集。
\end{lemma}

\begin{proof}
由上面的引理 \ref{topology-lemma-jacobson-equivalent-locally-closed}，只需证明：
若 $U$ 在 $X$ 中开，则 $U\cap X_0$ 在 $X_0$ 中逆紧，当且仅当 $U$
在 $X$ 中逆紧。为此，只需证明：若 $U$ 在 $X$ 中开，则
$U\cap X_0$ 拟紧，当且仅当 $U$ 拟紧。由引理
\ref{topology-lemma-jacobson-inherited}，可以用 $U$ 替换 $X$，并假设 $U = X$。
最后注意，$X$ 的任意开集族 $\mathcal{U}$ 覆盖 $X$，当且仅当它覆盖
$X_0$。这是因为若闭集 $X\setminus \bigcup \mathcal{U}$ 非空，
利用 $X$ 的 Jacobson 性质可在其中找到一个属于 $X_0$ 的点。
\end{proof}

\section{特化}
\label{topology-section-specialization}

\begin{definition}
\label{topology-definition-specialization}
设 $X$ 为拓扑空间。
\begin{enumerate}
\item 若 $x, x' \in X$，且 $x \in \overline{\{x'\}}$，则称 $x$ 是
$x'$ 的一个{\it 特化}，或称 $x'$ 是 $x$ 的一个{\it 一般化}。
记作 $x' \leadsto x$。
\item 若对所有 $x' \in T$ 以及每个特化 $x' \leadsto x$，都有
$x \in T$，则称子集 $T \subset X$ {\it 对特化稳定}。
\item 若对所有 $x \in T$ 以及每个一般化 $x' \leadsto x$，都有
$x' \in T$，则称子集 $T \subset X$ {\it 对一般化稳定}。
\end{enumerate}
\end{definition}

\begin{lemma}
\label{topology-lemma-open-closed-specialization}
设 $X$ 为拓扑空间。
\begin{enumerate}
\item $X$ 的任意闭子集对特化稳定。
\item $X$ 的任意开子集对一般化稳定。
\item 子集 $T \subset X$ 对特化稳定，当且仅当其补集 $T^c$
对一般化稳定。
\end{enumerate}
\end{lemma}

\begin{proof}
设 $F$ 为 $X$ 的闭子集。若 $y\in F$，则 $\{y\} \subset F$；
由于 $F$ 闭，故 $\overline{\{y\}} \subset \overline{F} = F$。
所以，对 $y$ 的每个特化 $x$，都有 $x\in F$。

\medskip\noindent
设 $x, y\in X$ 且 $x\in \overline{\{y\}}$，并设 $T$ 为 $X$
的子集。说 $T$ 对特化稳定，就是说 $y\in T$ 蕴含 $x\in T$；
反过来，说 $T$ 对一般化稳定，就是说 $x\in T$ 蕴含 $y\in T$。
因此，取逆否命题即证得 (3)。

\medskip\noindent
对补集应用 (1) 与 (3)，即得第二条性质。
\end{proof}

\begin{lemma}
\label{topology-lemma-stable-specialization}
设 $T \subset X$ 为拓扑空间 $X$ 的子集。下列条件等价：
\begin{enumerate}
\item $T$ 对特化稳定；
\item $T$ 是 $X$ 的闭子集的（有向）并。
\end{enumerate}
\end{lemma}

\begin{proof}
假设 $T$ 对特化稳定，则对所有 $y\in T$ 都有
$\overline{\{y\}} \subset T$。所以
$T = \bigcup_{y\in T} \overline{\{y\}}$，这是 $X$ 的闭子集之并。
反之，假设 $T =
\bigcup_{i\in I}F_i$，其中 $F_i$ 是 $X$ 的闭子集。若 $y\in T$，
则存在 $i\in I$ 使得 $y\in F_i$。因为 $F_i$ 闭，有
$\overline{\{y\}} \subset F_i \subset T$，这证明了 $T$ 对特化稳定。
\end{proof}

\begin{definition}
\label{topology-definition-lift-specializations}
设 $f : X \to Y$ 为拓扑空间的连续映射。
\begin{enumerate}
\item 若给定 $Y$ 中的 $y' \leadsto y$ 以及任意满足
$f(x') = y'$ 的 $x'\in X$，都存在 $X$ 中 $x'$ 的一个特化
$x' \leadsto x$，使得 $f(x) = y$，则称{\it 特化可沿 $f$ 提升}，
或称 $f$ 是{\it 特化型的}。
\item 若给定 $Y$ 中的 $y' \leadsto y$ 以及任意满足
$f(x) = y$ 的 $x\in X$，都存在 $X$ 中 $x$ 的一个一般化
$x' \leadsto x$，使得 $f(x') = y'$，则称{\it 一般化可沿 $f$ 提升}，
或称 $f$ 是{\it 一般化型的}。
\end{enumerate}
\end{definition}

\begin{lemma}
\label{topology-lemma-lift-specialization-composition}
设 $f : X \to Y$ 与 $g : Y \to Z$ 为拓扑空间的连续映射。
若特化可沿 $f$ 和 $g$ 提升，则特化可沿 $g \circ f$ 提升。
“一般化可沿其提升”的情形也同样成立。
\end{lemma}

\begin{proof}
设 $z'\leadsto z$ 为 $Z$ 中的特化，并设 $x' \in X$ 满足
$g\circ f (x') = z'$。因为特化可沿 $g$ 提升，存在 $Y$ 中 $f(x')$
的一个特化 $f(x') \leadsto y$，使得 $g(y) = z$。同样，因为特化可沿
$f$ 提升，存在 $X$ 中 $x'$ 的一个特化 $x' \leadsto x$，使得
$f(x) = y$。这给出 $X$ 中 $x'$ 的一个特化 $x' \leadsto x$，满足
$g\circ f(x) = z$。换言之，特化可沿 $g\circ f$ 提升。
\end{proof}

\begin{lemma}
\label{topology-lemma-lift-specializations-images}
设 $f : X \to Y$ 为拓扑空间的连续映射。
\begin{enumerate}
\item 若特化可沿 $f$ 提升，并且 $T \subset X$ 对特化稳定，则
$f(T) \subset Y$ 对特化稳定。
\item 若一般化可沿 $f$ 提升，并且 $T \subset X$ 对一般化稳定，则
$f(T) \subset Y$ 对一般化稳定。
\end{enumerate}
\end{lemma}

\begin{proof}
设 $y' \leadsto y$ 为 $Y$ 中的特化，其中 $y'\in f(T)$；并取
$x'\in T$，使得 $f(x') = y'$。因为特化可沿 $f$ 提升，存在 $X$
中 $x'$ 的一个特化 $x'\leadsto x$，使得 $f(x) = y$。但 $T$
对特化稳定，所以 $x\in T$，进而 $y \in f(T)$。因此 $f(T)$
对特化稳定。

\medskip\noindent
(2) 的证明完全相同，只需使用一般化可沿 $f$ 提升。
\end{proof}

\begin{lemma}
\label{topology-lemma-closed-open-map-specialization}
设 $f : X \to Y$ 为拓扑空间的连续映射。
\begin{enumerate}
\item 若 $f$ 是闭映射，则特化可沿 $f$ 提升。
\item 若 $f$ 是开映射，$X$ 是 Noether 拓扑空间，$X$ 的每个不可约
闭子集都有泛点，并且 $Y$ 是 Kolmogorov 空间，则一般化可沿 $f$ 提升。
\end{enumerate}
\end{lemma}

\begin{proof}
假设 $f$ 是闭映射。给定 $Y$ 中的 $y' \leadsto y$ 以及任意满足
$f(x') = y'$ 的 $x'\in X$。考虑 $X$ 的闭子集
$T = \overline{\{x'\}}$。于是 $f(T) \subset Y$ 是闭子集，且
$y' \in f(T)$。所以也有 $y \in f(T)$。故 $y = f(x)$，其中
$x \in T$；也就是说，$x' \leadsto x$。

\medskip\noindent
假设 $f$ 是开映射，$X$ 是 Noether 的，$X$ 的每个不可约闭子集都有泛点，
并且 $Y$ 是 Kolmogorov 空间。给定 $Y$ 中的 $y' \leadsto y$ 以及任意
满足 $f(x) = y$ 的 $x \in X$。考虑
$T = f^{-1}(\{y'\}) \subset X$。取包含关系 $x \in U \subset X$
所给出的 $x$ 的开邻域。于是 $f(U) \subset Y$ 开且 $y \in f(U)$，
故也有 $y' \in f(U)$。换言之，$T \cap U \not = \emptyset$。
这证明 $x \in \overline{T}$。由于 $X$ 是 Noether 的，$T$ 也是
Noether 的（引理 \ref{topology-lemma-Noetherian}）。因此它有不可约分支分解
$T = T_1 \cup \ldots \cup T_n$。相应地，
$\overline{T} = \overline{T_1} \cup \ldots \cup \overline{T_n}$。
由上可知，对某个 $i$ 有 $x \in \overline{T_i}$。依假设，存在泛点
$x' \in \overline{T_i}$，从而 $x' \leadsto x$。由于
$x' \in \overline{T}$，可知 $f(x') \in \overline{\{y'\}}$。注意
$f(\overline{T_i}) = f(\overline{\{x'\}}) \subset \overline{\{f(x')\}}$。
若 $f(x') \not = y'$，则由于 $Y$ 是 Kolmogorov 空间，$f(x')$
不是不可约闭子集 $\overline{\{y'\}}$ 的泛点，并且包含关系
$\overline{\{f(x')\}} \subset \overline{\{y'\}}$ 是严格的；
也就是说，$y' \not \in f(\overline{T_i})$。这与
$f(T_i) = \{y'\}$ 矛盾。所以 $f(x') = y'$，结论得证。
\end{proof}

\begin{lemma}
\label{topology-lemma-quotient-kolmogorov}
假设 $s, t : R \to U$ 与 $\pi : U \to X$ 为拓扑空间的连续映射，并且
\begin{enumerate}
\item $\pi$ 是开映射；
\item $U$ 是清醒的；
\item $s, t$ 的纤维有限；
\item 一般化可沿 $s, t$ 提升；
\item $(t, s)(R) \subset U \times U$ 是 $U$ 上的等价关系，并且
$X$ 作为集合是 $U$ 关于此等价关系的商。
\end{enumerate}
则 $X$ 是 Kolmogorov 空间。
\end{lemma}

\begin{proof}
性质 (3) 与 (5) 蕴含点 $x$ 对应于该等价关系的一个有限等价类
$\{u_1, \ldots, u_n\} \subset U$。假设 $x' \in X$ 是另一个点，
对应于等价类 $\{u'_1, \ldots, u'_m\} \subset U$。假设对某个
$i, j$ 有 $u_i \leadsto u'_j$。那么，对任意满足
$s(r') = u'_j$ 的 $r' \in R$，由 (4) 可找到 $r \leadsto r'$，
使得 $s(r) = u_i$。所以 $t(r) \leadsto t(r')$。由于
$\{u'_1, \ldots, u'_m\} = t(s^{-1}(\{u'_j\}))$，可知
$\{u'_1, \ldots, u'_m\}$ 的每个元素都是 $\{u_1, \ldots, u_n\}$
中某个元素的特化。因此
$\overline{\{u_1\}} \cup \ldots \cup \overline{\{u_n\}}$
是等价类之并，从而形如 $\pi^{-1}(Z)$，其中 $Z \subset X$ 为某个子集。
由 (1)，$Z$ 在 $X$ 中闭；事实上
$Z = \overline{\{x\}}$，因为对每个 $i$ 都有
$\pi(\overline{\{u_i\}}) \subset \overline{\{x\}}$。换言之，
$x \leadsto x'$ 当且仅当 $x$ 在 $U$ 中的某个提升特化到 $x'$ 在 $U$
中的某个提升；这又当且仅当 $x'$ 在 $U$ 中的每个提升都是 $x$ 在 $U$
中的某个提升的特化。

\medskip\noindent
假设 $x \leadsto x'$ 与 $x' \leadsto x$ 都成立。仍如上，设 $x$ 对应于
$\{u_1, \ldots, u_n\}$，而 $x'$ 对应于
$\{u'_1, \ldots, u'_m\}$。由上一段的结果，可以找到序列
$$
\ldots \leadsto u'_{j_3} \leadsto u_{i_3} \leadsto u'_{j_2} \leadsto
u_{i_2} \leadsto u'_{j_1} \leadsto u_{i_1}
$$
它必定重复；所以由 (2) 可知
$\{u_1, \ldots, u_n\} = \{u'_1, \ldots, u'_m\}$，即 $x = x'$。
因此 $X$ 是 Kolmogorov 空间。
\end{proof}


\begin{lemma}
\label{topology-lemma-dimension-specializations-lift}
设 $f : X \to Y$ 为拓扑空间的态射。假设 $Y$ 是清醒拓扑空间，且
$f$ 满射。若特化或一般化二者之一可沿 $f$ 提升，则
$\dim(X) \geq \dim(Y)$。
\end{lemma}

\begin{proof}
假设特化可沿 $f$ 提升。
% P01-ERRATA-E0019: frozen source omits a subset relation before Z_e in the chain; preserved here.
% P01-ERRATA-E0020: frozen prose calls this a chain of closed subsets of X although the formula and proof place it in Y; preserved here.
设 $Z_0 \subset Z_1 \subset \ldots Z_e \subset Y$ 为 $X$ 的不可约闭子集链。
设 $\xi_e \in X$ 为一个映到 $Z_e$ 的泛点的点。依假设，存在 $X$
中的特化 $\xi_e \leadsto \xi_{e - 1}$，使得 $\xi_{e - 1}$ 映到
$Z_{e - 1}$ 的泛点。如此继续，得到特化序列
$$
\xi_e \leadsto \xi_{e - 1} \leadsto \ldots \leadsto \xi_0
$$
其中 $\xi_i$ 映到 $Z_i$ 的泛点。这显然蕴含不可约闭子集序列
$$
\overline{\{\xi_0\}} \subset
\overline{\{\xi_1\}} \subset \ldots
\overline{\{\xi_e\}}
$$
是 $X$ 中长度为 $e$ 的链。一般化可沿 $f$ 提升的情形类似。
\end{proof}

\begin{lemma}
\label{topology-lemma-characterize-closed-Noetherian}
设 $X$ 为 Noether 清醒拓扑空间，$E \subset X$ 为 $X$ 的子集。
\begin{enumerate}
\item 若 $E$ 可构造且对特化稳定，则 $E$ 闭。
\item 若 $E$ 可构造且对一般化稳定，则 $E$ 开。
\end{enumerate}
\end{lemma}

\begin{proof}
设 $E$ 可构造且对一般化稳定。设 $Y \subset X$ 为不可约闭子集，
其泛点为 $\xi \in Y$。若 $E \cap Y$ 非空，则由对一般化的稳定性，
它包含 $\xi$，因而在 $Y$ 中稠密；所以它包含 $Y$ 的一个非空开集，
参见引理 \ref{topology-lemma-characterize-constructible-Noetherian}。
由引理 \ref{topology-lemma-characterize-open-Noetherian}，$E$ 是开的。
这证明了 (2)。为证明 (1)，对 $E$ 在 $X$ 中的补集应用 (2)。
\end{proof}

\section{维数函数}
\label{topology-section-dimension-function}

\noindent
除非所考虑的空间是清醒的（定义 \ref{topology-definition-generic-point}），
否则考虑维数函数几乎没有意义。因此，可以用与 $X$ 关联的清醒拓扑空间
来改进下面的定义。由于概形的底拓扑空间是清醒的，我们不作这一改进。

\begin{definition}
\label{topology-definition-dimension-function}
设 $X$ 为拓扑空间。
\begin{enumerate}
\item 设 $x, y \in X$ 且 $x \not = y$。假设 $x \leadsto y$，即 $y$
是 $x$ 的一个特化。若不存在 $z \in X \setminus \{x, y\}$，使得
$x \leadsto z$ 且 $z \leadsto y$，则称 $y$ 是 $x$ 的一个
{\it 直接特化}。
\item 若映射 $\delta : X \to \mathbf{Z}$ 满足下列条件，则称它为
{\it 维数函数}\footnote{这很可能不是标准术语。这一概念通常只对
（局部）Noether 概形引入；在这种情形，条件 (a) 由 (b) 蕴含。}：
\begin{enumerate}
\item 每当 $x \leadsto y$ 且 $x \not = y$ 时，都有
$\delta(x) > \delta(y)$；
\item 对 $X$ 中的每个直接特化 $x \leadsto y$，都有
$\delta(x) = \delta(y) + 1$。
\end{enumerate}
\end{enumerate}
\end{definition}

\noindent
显然，若 $\delta$ 是维数函数，则对任意 $t \in \mathbf{Z}$，
$\delta + t$ 也是维数函数。下面是一个有趣的引理。

\begin{lemma}
\label{topology-lemma-dimension-function-catenary}
设 $X$ 为拓扑空间。若 $X$ 清醒且有维数函数，则 $X$ 是链状空间。
此外，对任意 $x \leadsto y$，有
$$
\delta(x) - \delta(y) =
\text{codim}\left(\overline{\{y\}}, \ \overline{\{x\}}\right).
$$
\end{lemma}

\begin{proof}
假设 $Y \subset Y' \subset X$ 为不可约闭子集，并设
$\xi \in Y$、$\xi' \in Y'$ 为它们的泛点。由定义立即可见
$\text{codim}(Y, Y') \leq \delta(\xi) - \delta(\xi') < \infty$。
事实上，第一个不等式是等式。设
$$
Y = Y_0 \subset Y_1 \subset \ldots \subset Y_e = Y'
$$
为任意不可约闭子集极大链，并以 $\xi_i \in Y_i$ 表示泛点。
那么 $\xi_i \leadsto \xi_{i + 1}$ 是直接特化。因此，
$e = \delta(\xi) - \delta(\xi')$，正合所需。这也证明了引理的最后一个
断言。
\end{proof}

\begin{lemma}
\label{topology-lemma-dimension-function-unique}
设 $X$ 为拓扑空间，并设 $\delta$、$\delta'$ 为 $X$ 上的两个维数函数。
若 $X$ 局部 Noether 且清醒，则 $\delta - \delta'$ 在 $X$ 上局部常值。
\end{lemma}

\begin{proof}
设 $x \in X$ 为一点。我们将证明 $\delta - \delta'$ 在 $x$ 的一个邻域
上常值。可以把 $X$ 替换为 $X$ 中 $x$ 的一个 Noether 开邻域。
因此可假设 $X$ Noether 且清醒。设 $Z_1, \ldots, Z_r$ 为 $X$
中经过 $x$ 的不可约分支。（因为 $X$ 是 Noether 的，它们只有有限多个，
参见引理 \ref{topology-lemma-Noetherian}。）设 $\xi_i \in Z_i$ 为泛点。
注意，$Z_1 \cup \ldots \cup Z_r$ 是 $X$ 中 $x$ 的一个邻域
（未必闭）。我们断言 $\delta - \delta'$ 在
$Z_1 \cup \ldots \cup Z_r$ 上常值。事实上，若 $y \in Z_i$，则
$$
\delta(x) - \delta(y) = \delta(x) - \delta(\xi_i) + \delta(\xi_i) - \delta(y)
= - \text{codim}(\overline{\{x\}}, Z_i)
+ \text{codim}(\overline{\{y\}}, Z_i)
$$
这是由引理 \ref{topology-lemma-dimension-function-catenary} 得出的。
对 $\delta'$ 也同样。因此结论成立。
\end{proof}

\begin{lemma}
\label{topology-lemma-locally-dimension-function}
设 $X$ 局部 Noether、清醒且链状。则每个点都有一个拥有维数函数的开邻域
$U \subset X$。
\end{lemma}

\begin{proof}
我们将反复使用：链状空间的开子空间仍为链状空间，参见引理
\ref{topology-lemma-catenary}；以及 Noether 拓扑空间只有有限多个不可约分支，
参见引理 \ref{topology-lemma-Noetherian}。在引理
\ref{topology-lemma-dimension-function-unique} 的证明中，我们已经看到如何构造
这样的函数。具体而言，先把 $X$ 替换为 $x$ 的一个 Noether 开邻域。
接着，设 $Z_1, \ldots, Z_r \subset X$ 为 $X$ 的不可约分支。设
$$
Z_i \cap Z_j = \bigcup Z_{ijk}
$$
为不可约分支分解。将 $X$ 替换为
$$
X \setminus \left(
\bigcup\nolimits_{x \not \in Z_i} Z_i
\cup
\bigcup\nolimits_{x \not \in Z_{ijk}} Z_{ijk}
\right)
$$
于是可以假设对所有 $i$ 都有 $x \in Z_i$，并且对所有 $i, j, k$
都有 $x \in Z_{ijk}$。对 $y \in X$，任取满足 $y \in Z_i$ 的 $i$，
并令
$$
\delta(y) = - \text{codim}(\overline{\{y\}}, Z_i)
+ \text{codim}(\overline{\{x\}}, Z_i).
$$
我们断言这是一个维数函数。先证明它是良定义的，即与 $i$ 的选择无关。
假设对某个 $i, j, k$ 有 $y \in Z_{ijk}$。则有（使用引理
\ref{topology-lemma-catenary-in-codimension}）
\begin{align*}
\delta(y) & =
- \text{codim}(\overline{\{y\}}, Z_i)
+ \text{codim}(\overline{\{x\}}, Z_i) \\
& =
- \text{codim}(\overline{\{y\}}, Z_{ijk})
- \text{codim}(Z_{ijk}, Z_i)
+ \text{codim}(\overline{\{x\}}, Z_{ijk})
+ \text{codim}(Z_{ijk}, Z_i) \\
& =
- \text{codim}(\overline{\{y\}}, Z_{ijk})
+ \text{codim}(\overline{\{x\}}, Z_{ijk})
\end{align*}
它关于 $i$ 与 $j$ 对称。略去验证它是维数函数的证明。
\end{proof}

\begin{remark}
\label{topology-remark-obstruction-to-dimension-function}
合并引理 \ref{topology-lemma-dimension-function-unique} 与
\ref{topology-lemma-locally-dimension-function} 可知，在链状、局部 Noether、清醒的
拓扑空间上，存在维数函数的障碍是 $H^1(X, \mathbf{Z})$ 的一个元素。
\end{remark}

\section{无处稠密集}
\label{topology-section-nowhere-dense}

\begin{definition}
\label{topology-definition-nowhere-dense}
设 $X$ 为拓扑空间。
\begin{enumerate}
\item 给定子集 $T \subset X$，$T$ 的{\it 内部}是 $X$ 中包含于 $T$
的最大开子集。
\item 若 $T$ 的闭包内部为空，则称子集 $T \subset X$ {\it 无处稠密}。
\end{enumerate}
\end{definition}

\begin{lemma}
\label{topology-lemma-nowhere-dense}
设 $X$ 为拓扑空间。有限多个无处稠密集之并仍为无处稠密集。
\end{lemma}

\begin{proof}
只需对两个无处稠密集证明本引理，一般结果即可由归纳法得出。
设 $A,B \subset X$ 为无处稠密子集。有
$\overline{A \cup B} = \overline{A} \cup \overline{B}$。因此，若
$U \subset \overline{A \cup B}$ 是 $X$ 的开子集，则
% P01-ERRATA-E0021: frozen source omits grouping around U intersect closure(B) and closure(A); formulas preserved here.
$U \setminus U \cap \overline{B}$ 是 $U$ 以及 $X$ 的开子集，并包含于
$\overline{A}$，因而为空。同理，$U \setminus U \cap \overline{A}$
为空。所以 $U = \emptyset$，正合所需。
\end{proof}

\begin{lemma}
\label{topology-lemma-image-nowhere-dense-open}
设 $X$ 为拓扑空间，$U \subset X$ 为开集，$T \subset U$ 为子集。
若 $T$ 在 $U$ 中无处稠密，则 $T$ 在 $X$ 中无处稠密。
\end{lemma}

\begin{proof}
假设 $T$ 在 $U$ 中无处稠密。假设 $x \in X$ 是 $T$ 在 $X$ 中的闭包
$\overline{T}$ 的一个内点。设 $x \in V \subset \overline{T}$，其中
$V \subset X$ 在 $X$ 中开。注意，$\overline{T} \cap U$ 是 $T$
在 $U$ 中的闭包。由于 $\overline{T} \cap U$ 的内部为空，必有
$V \cap U = \emptyset$。所以 $x$ 不可能属于 $U$ 的闭包，更不可能属于
$T$ 的闭包，矛盾。
\end{proof}

\begin{lemma}
\label{topology-lemma-nowhere-dense-local}
设 $X$ 为拓扑空间，$X = \bigcup U_i$ 为开覆盖，且 $T \subset X$
为子集。若对所有 $i$，$T \cap U_i$ 在 $U_i$ 中无处稠密，则 $T$
在 $X$ 中无处稠密。
\end{lemma}

\begin{proof}
以 $\overline{T}_i$ 表示 $T \cap U_i$ 在 $U_i$ 中的闭包。
有 $\overline{T} \cap U_i = \overline{T}_i$。取内部与同开集求交可交换，
所以
$$
(\text{下列对象的内部： }\overline{T}) \cap U_i =
\text{下列对象的内部： }(\overline{T} \cap U_i) =
\text{在下列对象中的内部： }U_i\text{ 属于 }\overline{T}_i
$$
由假设，最后一项为空。因此 $T$ 在 $X$ 中无处稠密。
\end{proof}

\begin{lemma}
\label{topology-lemma-closed-image-nowhere-dense}
设 $f : X \to Y$ 为拓扑空间的连续映射，$T \subset X$ 为子集。
若 $f$ 是 $X$ 到 $Y$ 的一个闭子集上的同胚，并且 $T$ 在 $X$
中无处稠密，则 $f(T)$ 在 $Y$ 中也无处稠密。
\end{lemma}

\begin{proof}
因为 $f(X)$ 在 $Y$ 中闭，并且 $f$ 是 $X$ 到 $f(X)$ 上的同胚，
所以 $f(T)$ 在 $Y$ 中的闭包等于 $f(\overline{T})$。因此，若
$V \subset Y$ 是包含于 $f(T)$ 闭包的开集，则 $U = f^{-1}(V)$
是包含于 $\overline{T}$ 的开集。所以 $U = \emptyset$，进而
$V = \emptyset$，正合所需。
\end{proof}

\begin{lemma}
\label{topology-lemma-open-inverse-image-closed-nowhere-dense}
设 $f : X \to Y$ 为拓扑空间的连续映射，$T \subset Y$ 为子集。
若 $f$ 是开映射，且 $T$ 是 $Y$ 的闭无处稠密子集，则 $f^{-1}(T)$
也是 $X$ 的闭无处稠密子集。若 $f$ 满射且为开映射，则 $T$ 是闭无处
稠密集，当且仅当 $f^{-1}(T)$ 是闭无处稠密集。
\end{lemma}

\begin{proof}
略。（提示：在第一种情形，$f^{-1}(T)$ 的内部映入 $T$ 的内部；
在第二种情形，$f^{-1}(T)$ 的内部满射到 $T$ 的内部。）
\end{proof}

\begin{lemma}
\label{topology-lemma-dense-in-constructible}
设 $X$ 为拓扑空间。设 $E \subset X$ 为有限多个局部闭子集之并
（例如，$E$ 可构造）。则 $E$ 包含其闭包的一个稠密开子集。
\end{lemma}

\begin{proof}
用 $E$ 的闭包替换 $X$ 后，可以假设 $E$ 在 $X$ 中稠密；我们需要证明
$E$ 包含一个稠密开集。于是 $E$ 的补集 $T$ 是有限多个局部闭子集之并，
设 $T = T_1 \cup \ldots \cup T_n$，并且 $T$ 在 $X$ 中无处稠密。
由于 $T_i$ 在其闭包 $Z_i$ 中开且稠密，可知 $Z_i$ 在 $X$ 中也无处
稠密。于是，由引理 \ref{topology-lemma-nowhere-dense}，$T$ 的闭包
$Z = Z_1 \cup \ldots \cup Z_n$ 在 $X$ 中无处稠密。由此得出，$E$
包含稠密开集 $X \setminus Z$。
\end{proof}

\section{仿有限空间}
\label{topology-section-profinite}

\noindent
定义如下。

\begin{definition}
\label{topology-definition-profinite}
若一个拓扑空间同胚于有限离散空间图表的极限，则称它是
{\it 仿有限的}。
\end{definition}

\noindent
这并不是仿有限空间最方便的刻画。

\begin{lemma}
\label{topology-lemma-profinite}
设 $X$ 为拓扑空间。下列条件等价：
\begin{enumerate}
\item $X$ 是仿有限空间；
\item $X$ 是 Hausdorff、拟紧且全不连通的。
\end{enumerate}
若这些条件成立，则 $X$ 是有限离散空间的余滤极限。
\end{lemma}

\begin{proof}
假设 (1)。选取有限离散空间的图表 $i \mapsto X_i$，使得
$X = \lim X_i$。因为每个 $X_i$ 都是 Hausdorff 且拟紧的，由引理
\ref{topology-lemma-inverse-limit-quasi-compact} 可知 $X$ 拟紧。若
$x, x' \in X$ 是不同的点，则 $x$ 与 $x'$ 在某个 $X_i$ 中的像不同。
所以 $x$ 与 $x'$ 有不交开邻域，即 $X$ 是 Hausdorff 的。用完全相同的
方法可见 $X$ 全不连通。

\medskip\noindent
假设 (2)。令 $\mathcal{I}$ 为所有有限不交并分解
$X = \coprod_{i \in I} U_i$ 构成的集合，其中对所有 $i \in I$，
$U_i$ 都是非空开（因而也闭）子集。对每个 $I \in \mathcal{I}$，
存在连续映射 $X \to I$，把 $U_i$ 的点送到 $i$。定义偏序如下：
对 $I, I' \in \mathcal{I}$，$I \leq I'$ 当且仅当 $I'$ 对应的覆盖
加细 $I$ 对应的覆盖。在这种情形，我们得到典范映射 $I' \to I$。
换言之，我们得到以 $\mathcal{I}$ 为指标的有限离散空间逆系统。
映射 $X \to I$ 相容地拼合，给出连续映射
$$
X \longrightarrow \lim_{I \in \mathcal{I}} I
$$
我们断言此映射是同胚，这便完成证明。（最后一个断言也随之成立，因为
偏序集 $\mathcal{I}$ 是有向的：给定 $X$ 的两个不交并分解，可以找到
一个同时加细二者的第三个分解。）事实上，该映射是单射，因为 $X$
全不连通，故 $\{x\}$ 是 $X$ 中所有包含 $x$ 的开闭子集之交
（引理 \ref{topology-lemma-connected-component-intersection-compact-Hausdorff}）；
该映射由引理 \ref{topology-lemma-intersection-closed-in-quasi-compact} 是满射。
由引理 \ref{topology-lemma-bijective-map}，该映射是同胚。
\end{proof}

\begin{lemma}
\label{topology-lemma-directed-inverse-limit-profinite}
仿有限空间的极限仍是仿有限空间。
\end{lemma}

\begin{proof}
设 $i \mapsto X_i$ 为以指标范畴 $\mathcal{I}$ 为指标的仿有限空间图表。
使用引理 \ref{topology-lemma-profinite} 对仿有限空间的刻画。特别地，每个 $X_i$
都是 Hausdorff、拟紧且全不连通的。由引理 \ref{topology-lemma-limits}，极限
$X = \lim X_i$ 存在。由引理 \ref{topology-lemma-inverse-limit-quasi-compact}，
极限 $X$ 拟紧。设 $x, x' \in X$ 为不同的点。则存在某个 $i$，使得
$x$ 与 $x'$ 在 $X_i$ 中的像 $x_i$ 与 $x'_i$ 在投影 $X \to X_i$
下彼此不同。
于是 $x_i$ 与 $x'_i$ 在 $X_i$ 中有不交开邻域。取这些开集的原像，
可知 $X$ 是 Hausdorff 的。同样，$x_i$ 与 $x'_i$ 位于 $X_i$ 的不同
连通分支中，从而 $x$ 与 $x'$ 必位于 $X$ 的不同连通分支中。因此
$X$ 全不连通。这完成了证明。
\end{proof}

\begin{lemma}
\label{topology-lemma-profinite-refine-open-covering}
设 $X$ 为仿有限空间。$X$ 的每个开覆盖都有一个形如
$X = \coprod U_i$ 的有限覆盖加细，其中 $U_i$ 开且闭。
\end{lemma}

\begin{proof}
将 $X = \lim X_i$ 写成以有向集 $I$ 为指标的有限离散空间逆系统的极限
（引理 \ref{topology-lemma-profinite}）。以 $f_i : X \to X_i$ 表示投影。
对每个点 $x = (x_i) \in X$，开邻域的一个基本系由
$f_i^{-1}(\{x_i\})$ 构成。因此，由于 $X$ 拟紧，可以假设有开覆盖
$$
X = f_{i_1}^{-1}(\{x_{i_1}\}) \cup \ldots \cup f_{i_n}^{-1}(\{x_{i_n}\})
$$
选取 $i \in I$，使得对 $j = 1, \ldots, n$ 有 $i \geq i_j$
（因为 $I$ 是有向集，这是可能的）。于是覆盖
$$
X = \coprod\nolimits_{t \in X_i} f_i^{-1}(\{t\})
$$
加细给定覆盖，并且具有所要求的形式。
\end{proof}

\begin{lemma}
\label{topology-lemma-pi0-profinite}
设 $X$ 为拓扑空间。若 $X$ 拟紧，并且 $X$ 的每个连通分支都是所有
包含它的开闭子集之交，则 $\pi_0(X)$ 是仿有限空间。
\end{lemma}

\begin{proof}
我们用引理 \ref{topology-lemma-profinite} 来证明。由于 $\pi_0(X)$ 是拟紧空间的像，
它拟紧（引理 \ref{topology-lemma-image-quasi-compact}）。按构造，它全不连通
（引理 \ref{topology-lemma-space-connected-components}）。设 $C, D \subset X$
为 $X$ 的两个不同连通分支。把 $C = \bigcap U_\alpha$ 写成 $X$ 中所有
包含 $C$ 的开闭子集之交。$U_\alpha$ 的任意有限交仍属此类。由于
$\bigcap U_\alpha \cap D = \emptyset$，可知对某个 $\alpha$ 有
$U_\alpha \cap D = \emptyset$（使用引理
\ref{topology-lemma-connected-components}、\ref{topology-lemma-closed-in-quasi-compact} 与
\ref{topology-lemma-intersection-closed-in-quasi-compact}）。因为 $U_\alpha$ 开且闭，
它是其所包含的连通分支之并；也就是说，$U_\alpha$ 是某个开闭子集
$V_\alpha \subset \pi_0(X)$ 的原像。这证明了对应于 $C$ 与 $D$ 的点
包含于不交开子集中，即 $\pi_0(X)$ 是 Hausdorff 的。
\end{proof}

\section{谱空间}
\label{topology-section-spectral}

\noindent
本节材料取自 \cite{Hochster} 与 \cite{Hochster-thesis}。Hochster 在其
博士论文中证明了（除其他结果外），谱空间恰好是能够作为某个环的谱出现的
拓扑空间。

\begin{definition}
\label{topology-definition-spectral-space}
若拓扑空间 $X$ 清醒且拟紧，任意两个拟紧开集之交拟紧，并且拟紧开集族
构成拓扑基，则称它为{\it 谱空间}。若谱空间的连续映射
$f : X \to Y$ 对每个拟紧开集的原像都拟紧，则称它为{\it 谱映射}。
\end{definition}

\noindent
换言之，谱空间的连续映射是谱映射，当且仅当它拟紧
（定义 \ref{topology-definition-quasi-compact}）。

\medskip\noindent
设 $X$ 为谱空间。$X$ 上的{\it 可构造拓扑}，是以集合 $U$ 与 $U^c$
为开集子基的拓扑，其中 $U$ 是 $X$ 的拟紧开集。注意，因为 $X$ 是谱空间，
开集 $U \subset X$ 逆紧，当且仅当 $U$ 拟紧。因此，可构造拓扑也可刻画为
使 $X$ 的每个可构造子集同时开且闭的最粗拓扑（可构造集的定义与性质见
第 \ref{topology-section-constructible} 节）。由此可知，$X$ 的子集在可构造拓扑
中开（相应地，闭），当且仅当它是可构造子集之并（相应地，交）。
由于拟紧开集族构成 $X$ 的拓扑基，可见可构造拓扑比 $X$ 上的给定拓扑更细。

\begin{lemma}
\label{topology-lemma-constructible-hausdorff-quasi-compact}
设 $X$ 为谱空间。可构造拓扑是 Hausdorff、全不连通且拟紧的。
\end{lemma}

\begin{proof}
设 $x, y \in X$ 且 $x \not = y$。因为 $X$ 清醒，存在恰好包含两点
$x, y$ 之一的开子集 $U$。不妨设 $x \in U$。可将 $U$ 替换为包含于
$U$ 的 $x$ 的拟紧开邻域。于是 $U$ 与 $U^c$ 在可构造拓扑中都开且闭。
因为 $x \in U$ 与 $y \in U^c$ 是可构造拓扑中的不交开集，所以 $X$
在可构造拓扑中是 Hausdorff 的。$U$ 的存在也蕴含 $x$ 与 $y$ 在可构造
拓扑中位于不同连通分支，故 $X$ 在可构造拓扑中全不连通。

\medskip\noindent
令 $\mathcal{B}$ 为所有子集 $B \subset X$ 构成的集合，其中 $B$
或者是拟紧开集，或者是具有拟紧补集的闭集。若 $B \in \mathcal{B}$，
则 $B^c \in \mathcal{B}$。由引理 \ref{topology-lemma-subbase-theorem}，只需证明
每个满足 $B_i \in \mathcal{B}$ 的覆盖
$X = \bigcup_{i \in I} B_i$ 都有有限加细。取补集可知，我们必须证明：
若 $\mathcal{B}$ 中的元素族 $\{B_i\}_{i \in I}$ 对所有 $n$ 以及所有
$i_1, \ldots, i_n \in I$ 都满足
$B_{i_1} \cap \ldots \cap B_{i_n} \not = \emptyset$，则它有公共交点。
可以并且确实假设：若 $i \not = i'$，则 $B_i \not = B_{i'}$。

\medskip\noindent
反设 $\{B_i\}_{i \in I}$ 是 $\mathcal{B}$ 中具有有限交性质但总交为空
的元素族。应用 Zorn 引理可假设该族极大（略去细节）。令
$I' \subset I$ 为使得 $B_i$ 闭的那些指标，并置
$Z = \bigcap_{i \in I'} B_i$。由引理
\ref{topology-lemma-intersection-closed-in-quasi-compact}，这是 $X$ 的非空闭子集。
若 $Z$ 可约，则可写成 $Z = Z' \cup Z''$，其中两个闭子集都
不等于 $Z$。特别地，这意味着可以找到与 $Z'$ 相交但不与
$Z''$ 相交的拟紧开集 $U' \subset X$。类似地，可以找到与 $Z''$
相交但不与 $Z'$ 相交的拟紧开集 $U'' \subset X$。置
$B' = X \setminus U'$ 与 $B'' = X \setminus U''$。注意
$Z'' \subset B'$ 且 $Z' \subset B''$。若存在有限多个指标
$i_1, \ldots, i_n \in I$，使得
$B' \cap B_{i_1} \cap \ldots \cap B_{i_n} = \emptyset$，并且还存在
有限多个指标 $j_1, \ldots, j_m \in I$，使得
$B'' \cap B_{j_1} \cap \ldots \cap B_{j_m} = \emptyset$，则可得
$Z \cap B_{i_1} \cap \ldots \cap B_{i_n} \cap B_{j_1} \cap \ldots \cap B_{j_m}
= \emptyset$。然而，集合
$B_{i_1} \cap \ldots \cap B_{i_n} \cap B_{j_1} \cap \ldots \cap B_{j_m}$
拟紧，因此可以找到有限多个指标 $i'_1, \ldots, i'_l \in I'$，使得
$B_{i_1} \cap \ldots \cap B_{i_n} \cap B_{j_1} \cap \ldots \cap
B_{j_m} \cap B_{i'_1} \cap \ldots \cap B_{i'_l} = \emptyset$，矛盾。
因此，可以把 $B'$ 或 $B''$ 之一加入给定族，这与极大性矛盾。
由此得出 $Z$ 不可约。但这也导致矛盾，因为此时对每个
$i \in I \setminus I'$，非空开集 $Z \cap B_i$（其非空性由上面同样的
论证得出）都包含 $Z$ 的唯一泛点。这个矛盾证明了引理。
\end{proof}

\begin{lemma}
\label{topology-lemma-fibres-spectral-map-quasi-compact}
设 $f : X \to Y$ 为谱空间之间的谱映射。则
\begin{enumerate}
\item $f$ 对可构造拓扑连续；
\item $f$ 的纤维拟紧；
\item 其像在可构造拓扑中闭。
\end{enumerate}
\end{lemma}

\begin{proof}
以 $X'$ 与 $Y'$ 表示分别赋予可构造拓扑的 $X$ 与 $Y$；由引理
\ref{topology-lemma-constructible-hausdorff-quasi-compact}，它们是拟紧 Hausdorff
空间。第 (1) 条是说 $X' \to Y'$ 连续，它立即由定义得出。第 (3) 条成立，
因为 $f(X')$ 是 Hausdorff 空间 $Y'$ 的拟紧子集，参见引理
\ref{topology-lemma-quasi-compact-in-Hausdorff}。我们有拓扑空间连续映射的交换图
$$
\xymatrix{
X' \ar[r] \ar[d] & X \ar[d] \\
Y' \ar[r] & Y
}
$$
因为 $Y'$ 是 Hausdorff 的，可知纤维 $X'_y$ 在 $X'$ 中闭。因为 $X'$
拟紧，可知 $X'_y$ 拟紧（引理 \ref{topology-lemma-closed-in-quasi-compact}）。
由于 $X'_y \to X_y$ 是满连续映射，由引理
\ref{topology-lemma-image-quasi-compact} 得出 $X_y$ 拟紧。
\end{proof}

\begin{lemma}
\label{topology-lemma-spectral-if-continuous-wrt-constructible-top}
设 $X$ 与 $Y$ 为谱空间，$f : X \to Y$ 为连续映射。则 $f$ 是谱映射，
当且仅当 $f$ 对可构造拓扑连续。
\end{lemma}

\begin{proof}
必要性是引理 \ref{topology-lemma-fibres-spectral-map-quasi-compact}。假设 $f$
对可构造拓扑连续。设 $V \subset Y$ 为拟紧开集。则 $V$ 在可构造拓扑
中开且闭。因此 $f^{-1}(V)$ 在可构造拓扑中开且闭。由引理
\ref{topology-lemma-constructible-hausdorff-quasi-compact}，$X$ 在可构造拓扑中
拟紧，所以 $f^{-1}(V)$ 在可构造拓扑中拟紧。由于恒等映射
$f^{-1}(V) \to f^{-1}(V)$ 从可构造拓扑到通常拓扑是满且连续的，
由引理 \ref{topology-lemma-image-quasi-compact}，$f^{-1}(V)$ 在 $X$ 的拓扑中拟紧。
这完成了证明。
\end{proof}

\begin{lemma}
\label{topology-lemma-spectral-sub}
设 $X$ 为谱空间。设 $E \subset X$ 在可构造拓扑中闭（例如可构造或闭）。
则赋予诱导拓扑的 $E$ 是谱空间。
\end{lemma}

\begin{proof}
设 $Z \subset E$ 为闭不可约子集，并设 $\eta$ 为 $Z$ 在 $X$ 中的闭包
$\overline{Z}$ 的泛点。为证明 $E$ 清醒，我们证明 $\eta \in E$。
若不然，因为 $E$ 在可构造拓扑中闭，存在可构造子集 $F \subset X$，
使得 $\eta \in F$ 且 $F \cap E = \emptyset$。由引理
\ref{topology-lemma-generic-point-in-constructible}，这蕴含
$F \cap \overline{Z}$ 包含 $\overline{Z}$ 的一个非空开子集。但这是不可能
的，因为 $\overline{Z}$ 是 $Z$ 的闭包，而 $Z \cap F = \emptyset$。

\medskip\noindent
因为 $E$ 在可构造拓扑中闭，所以它在可构造拓扑中拟紧
（引理 \ref{topology-lemma-closed-in-quasi-compact} 与
\ref{topology-lemma-constructible-hausdorff-quasi-compact}）。因此，它在来自 $X$
的拓扑中更是拟紧的。若 $U \subset X$ 是拟紧开集，则 $E \cap U$
在可构造拓扑中闭，因而拟紧（如上所见）。由此可知，$E$ 的拟紧开子集
是形如 $E \cap U$ 的交集，其中 $U$ 是 $X$ 的拟紧开集。这些集合构成
拓扑基。最后，给定两个拟紧开集 $U, U' \subset X$，交集
$(E \cap U) \cap (E \cap U') = E \cap (U \cap U')$，并且由于 $X$
是谱空间，$U \cap U'$ 拟紧。这完成了证明。
\end{proof}

\begin{lemma}
\label{topology-lemma-constructible-stable-specialization-closed}
设 $X$ 为谱空间。设 $E \subset X$ 为在可构造拓扑中闭的子集
（例如可构造子集）。
\begin{enumerate}
\item 若 $x \in \overline{E}$，则 $x$ 是 $E$ 中某点的特化。
\item 若 $E$ 对特化稳定，则 $E$ 闭。
\item 若 $E' \subset X$ 在可构造拓扑中开（例如可构造），并且对一般化
稳定，则 $E'$ 开。
\end{enumerate}
\end{lemma}

\begin{proof}
证明 (1)。设 $x \in \overline{E}$。令 $\{U_i\}$ 为 $x$ 的所有拟紧
开邻域构成的集合。有限多个 $U_i$ 之交仍属此类。对所有 $i$，交集
$U_i \cap E$ 非空。由于子集 $U_i \cap E$ 在可构造拓扑中闭，由引理
\ref{topology-lemma-constructible-hausdorff-quasi-compact} 与引理
\ref{topology-lemma-intersection-closed-in-quasi-compact} 可知
$\bigcap (U_i \cap E)$ 非空。因为 $\{U_i\}$ 是 $x$ 的开邻域基本系，
可知 $\bigcap U_i$ 是 $x$ 的一般化点集。因此，$x$ 是 $E$ 中某点的特化。

\medskip\noindent
第 (2) 条立即由 (1) 得出。

\medskip\noindent
证明 (3)。假设 $E'$ 如 (3) 所述。$E'$ 的补集在可构造拓扑中闭
（引理 \ref{topology-lemma-constructible}），并且对特化稳定
（引理 \ref{topology-lemma-open-closed-specialization}）。所以由 (2)，该补集闭；
也就是说，$E'$ 开。
\end{proof}

\begin{lemma}
\label{topology-lemma-two-points}
设 $X$ 为谱空间，$x, y \in X$。那么，或者存在一个同时特化到 $x$ 与
$y$ 的第三个点，或者存在分别包含 $x$ 与 $y$ 的不交开邻域。
\end{lemma}

\begin{proof}
令 $\{U_i\}$ 为 $x$ 的所有拟紧开邻域构成的集合。有限多个 $U_i$
之交仍属此类。令 $\{V_j\}$ 为 $y$ 的所有拟紧开邻域构成的集合。
有限多个 $V_j$ 之交仍属此类。若对某个 $i, j$，$U_i \cap V_j$ 为空，
则结论成立。否则，对所有 $i$ 与 $j$，交集 $U_i \cap V_j$ 都非空。
集合 $U_i \cap V_j$ 在 $X$ 的可构造拓扑中闭。由引理
\ref{topology-lemma-constructible-hausdorff-quasi-compact} 与引理
\ref{topology-lemma-intersection-closed-in-quasi-compact} 可知
$\bigcap (U_i \cap V_j)$ 非空。由于 $X$ 清醒，且 $\{U_i\}$ 是 $x$
的开邻域基本系，可知 $\bigcap U_i$ 是 $x$ 的一般化点集。类似地，
$\bigcap V_j$ 是 $y$ 的一般化点集。因此，
$\bigcap (U_i \cap V_j)$ 的任意元素都同时特化到 $x$ 与 $y$。
\end{proof}

\begin{lemma}
\label{topology-lemma-characterize-profinite-spectral}
设 $X$ 为谱空间。下列条件等价：
\begin{enumerate}
\item $X$ 仿有限；
\item $X$ 是 Hausdorff 的；
\item $X$ 全不连通；
\item 每个拟紧开集都闭；
\item 点之间不存在非平凡特化；
\item $X$ 的每个点都闭；
\item $X$ 的每个点都是 $X$ 的某个不可约分支的泛点；
\item 可构造拓扑等于 $X$ 上的给定拓扑；
% P01-ERRATA-E0022: frozen source contains the editorial placeholder “add more here”; translated literally.
\item 在此添加更多条件。
\end{enumerate}
\end{lemma}

\begin{proof}
引理 \ref{topology-lemma-profinite} 给出蕴含关系 (1) $\Rightarrow$ (3)。
不可约分支是闭的，所以若 $X$ 全不连通，则每个点都闭。因此 (3) 蕴含
(6)。(6) 与 (5) 的等价性是立即的；而因为 $X$ 清醒，
(6) $\Leftrightarrow$ (7) 成立。假设 (5)。则 $X$ 的所有可构造子集都闭
（引理 \ref{topology-lemma-constructible-stable-specialization-closed}），特别地，
所有拟紧开集都闭。所以 (5) 蕴含 (4)。因为 $X$ 清醒，对任意两点，
存在恰好包含其中一点的拟紧开集，故 (4) 蕴含 (2)。由可构造拓扑的定义，
(4) 与 (8) 等价。只需再证明 (2) 蕴含 (1)。假设 $X$ 是 Hausdorff 的。
每个拟紧开集也闭（引理 \ref{topology-lemma-quasi-compact-in-Hausdorff}）。
这蕴含 $X$ 全不连通。所以由引理 \ref{topology-lemma-profinite}，它仿有限。
\end{proof}

\begin{lemma}
\label{topology-lemma-spectral-pi0}
若 $X$ 为谱空间，则 $\pi_0(X)$ 为仿有限空间。
\end{lemma}

\begin{proof}
合并引理 \ref{topology-lemma-connected-component-intersection} 与
\ref{topology-lemma-pi0-profinite}。
\end{proof}

\begin{lemma}
\label{topology-lemma-product-spectral-spaces}
两个谱空间的积是谱空间。
\end{lemma}

\begin{proof}
设 $X$、$Y$ 为谱空间。以 $p : X \times Y \to X$ 与
$q : X \times Y \to Y$ 表示投影。设 $Z \subset X \times Y$ 为闭不可约
子集。那么 $p(Z) \subset X$ 不可约，且 $q(Z) \subset Y$ 不可约。
% P01-ERRATA-E0023: frozen source names p(X) and q(Y), while the proof immediately and necessarily uses p(Z) and q(Z); preserved here.
设 $x \in X$ 为 $p(X)$ 的闭包的泛点，并设 $y \in Y$ 为 $q(Y)$
的闭包的泛点。若 $(x, y) \not \in Z$，则存在开集
$x \in U \subset X$、$y \in V \subset Y$，使得
$Z \cap U \times V = \emptyset$。所以 $Z$ 包含于
$(X \setminus U) \times Y \cup X \times (Y \setminus V)$。由于 $Z$
不可约，可知或者 $Z \subset (X \setminus U) \times Y$，或者
$Z \subset X \times (Y \setminus V)$。在第一种情形，
$p(Z) \subset (X \setminus U)$；在第二种情形，
$q(Z) \subset (Y \setminus V)$。两种情形都荒谬，因为 $x$ 位于
$p(Z)$ 的闭包中，且 $y$ 位于 $q(Z)$ 的闭包中。因此
$(x, y) \in Z$，这意味着 $(x, y)$ 是 $Z$ 的泛点。

\medskip\noindent
$X \times Y$ 的拓扑基由形如 $U \times V$ 的开集组成，其中
$U \subset X$ 与 $V \subset Y$ 为拟紧开集（这里用到 $X$ 与 $Y$
都是谱空间）。那么，作为拟紧空间之积，$U \times V$ 拟紧。此外，
$X \times Y$ 的任意拟紧开集都是这种拟紧矩形 $U \times V$ 的有限并。
由此可知，任意两个这种开集之交仍拟紧（因为 $X$ 与 $Y$ 是谱空间）。
这完成了证明。
\end{proof}

\begin{lemma}
\label{topology-lemma-spectral-bijective}
设 $f : X \to Y$ 为拓扑空间的连续映射。若
\begin{enumerate}
\item $X$ 与 $Y$ 都是谱空间；
\item $f$ 是双射谱映射；
\item 一般化（相应地，特化）可沿 $f$ 提升，
\end{enumerate}
则 $f$ 是同胚。
\end{lemma}

\begin{proof}
因为 $f$ 是谱映射，它在 $X$ 与 $Y$ 的可构造拓扑之间定义连续映射。
由引理 \ref{topology-lemma-constructible-hausdorff-quasi-compact} 与
\ref{topology-lemma-bijective-map}，$X \to Y$ 在可构造拓扑中是同胚。
设 $U \subset X$ 为拟紧开集。则 $f(U)$ 在 $Y$ 中可构造。设
$y \in Y$ 特化到 $f(U)$ 中的某个点。由最后一个假设，$f^{-1}(y)$
特化到 $U$ 的某个点。所以 $f^{-1}(y) \in U$，从而 $y \in f(U)$。
由引理 \ref{topology-lemma-constructible-stable-specialization-closed}，可知
$f(U)$ 开。因此 $f$ 是同胚。为在特化可沿 $f$ 提升的情形证明本引理，
改为证明：若 $X \setminus Z$ 是 $X$ 的拟紧开集，则 $f(Z)$ 闭。
\end{proof}

\begin{lemma}
\label{topology-lemma-directed-inverse-limit-finite-sober-spectral-spaces}
有限清醒拓扑空间的有向逆系统之逆极限是谱拓扑空间。
\end{lemma}

\begin{proof}
设 $I$ 为有向集，$X_i$ 为以 $I$ 为指标的有限清醒空间逆系统。
令 $X = \lim X_i$，其存在性由引理 \ref{topology-lemma-limits} 得出。
作为集合，$X = \lim X_i$。以 $p_i : X \to X_i$ 表示投影。
因为 $I$ 有向，可以应用引理 \ref{topology-lemma-describe-limits}。拓扑基由开集
$p_i^{-1}(U_i)$ 给出，其中 $U_i \subset X_i$ 开。因为
$p_i^{-1}(U_i)$ 的开覆盖特别也是仿有限拓扑中的开覆盖，可知
$p_i^{-1}(U_i)$ 拟紧。给定 $U_i \subset X_i$ 与 $U_j \subset X_j$，
则对某个 $k \geq i, j$ 以及某个开集 $U_k \subset X_k$，有
$p_i^{-1}(U_i) \cap p_j^{-1}(U_j) = p_k^{-1}(U_k)$。最后，若
$Z \subset X$ 不可约且闭，则 $p_i(Z) \subset X_i$ 不可约，因而有唯一
泛点 $\xi_i$（因为 $X_i$ 是有限清醒拓扑空间）。于是
$\xi = \lim \xi_i$ 是 $Z$ 的泛点（它是 $Z$ 的点，因为 $Z$ 闭）。
这完成了证明。
\end{proof}

\begin{lemma}
\label{topology-lemma-spectral-closed-in-product-two-point-space}
设 $W$ 为恰有两个点的拓扑空间，其中一个点闭，另一个不闭。一个拓扑空间
是谱空间，当且仅当它同胚于若干个 $W$ 的积的一个子空间，且该子空间在
可构造拓扑中闭。
\end{lemma}

\begin{proof}
写成 $W = \{0, 1\}$，其中 $0$ 是 $1$ 的特化，但反之不然。设 $I$
为集合。由引理 \ref{topology-lemma-directed-inverse-limit-finite-sober-spectral-spaces}，
空间 $\prod_{i \in I} W$ 是谱空间。因此，由引理 \ref{topology-lemma-spectral-sub}，
$\prod_{i \in I} W$ 中在可构造拓扑下闭的子空间是谱空间。

\medskip\noindent
反之，设 $X$ 为谱空间。设 $U \subset X$ 为拟紧开集。考虑连续映射
$$
f_U : X \longrightarrow W
$$
它把 $U$ 中的每个点映到 $1$，把 $X \setminus U$ 中的每个点映到 $0$。
取这些映射的积，得到连续映射
$$
f = \prod f_U : X \longrightarrow \prod\nolimits_U W
$$
% P01-ERRATA-E0024: frozen source uses an undefined codomain Y at this point; preserved here until Y is defined below as f(X).
按构造，映射 $f : X \to Y$ 是谱映射。由引理
\ref{topology-lemma-fibres-spectral-map-quasi-compact}，$f$ 的像在可构造拓扑中闭。
若 $x', x \in X$ 不同，则因为 $X$ 清醒，或者 $x'$ 不是 $x$ 的特化，
或者反过来。无论哪种情形（因为拟紧开集构成 $X$ 的拓扑基），都存在
拟紧开集 $U \subset X$，使得 $f_U(x') \not = f_U(x)$。因此 $f$
是单射。令 $Y = f(X)$ 并赋予诱导拓扑。设 $y' \leadsto y$ 为 $Y$
中的特化，并设 $f(x') = y'$ 且 $f(x) = y$。如上论证可知
$x' \leadsto x$，因为否则存在某个 $U$，使得 $x \in U$ 且
$x' \not \in U$，这会蕴含 $f_U(x') \not \leadsto f_U(x)$。
由引理 \ref{topology-lemma-spectral-bijective}，得出 $f : X \to Y$ 是同胚。
\end{proof}

\begin{lemma}
\label{topology-lemma-spectral-inverse-limit-finite-sober-spaces}
一个拓扑空间是谱空间，当且仅当它是有限清醒拓扑空间的有向逆极限。
\end{lemma}

\begin{proof}
一个方向由引理
\ref{topology-lemma-directed-inverse-limit-finite-sober-spectral-spaces} 给出。
反之，假设 $X$ 是谱空间。由引理
\ref{topology-lemma-spectral-closed-in-product-two-point-space}，可以假设
$X \subset \prod_{i \in I} W$ 是在可构造拓扑中闭的子集，其中
$W = \{0, 1\}$。可以把
$$
\prod\nolimits_{i \in I} W =
\lim_{J \subset I\text{ 有限 }} \prod\nolimits_{j \in J} W
$$
写成余滤极限。对每个 $J$，令 $X_J \subset \prod_{j \in J} W$ 为
$X$ 的像。因为 $X$ 在赋予可构造拓扑的积中闭，可知作为集合
$X = \lim X_J$（略去细节）。关于极限的一个形式论证（略）表明，
作为拓扑空间也有 $X = \lim X_J$。
\end{proof}

\begin{lemma}
\label{topology-lemma-Noetherian-goes-to-spectral}
设 $X$ 为拓扑空间，并设 $c : X \to X'$ 为从 $X$ 到清醒拓扑空间的
泛映射，参见引理 \ref{topology-lemma-make-sober}。
\begin{enumerate}
\item 若 $X$ 拟紧，则 $X'$ 也拟紧。
\item 若 $X$ 拟紧、有拟紧开集基，并且任意两个拟紧开集之交拟紧，
则 $X'$ 是谱空间。
\item 若 $X$ 是 Noether 的，则 $X'$ 是 Noether 谱空间。
\end{enumerate}
\end{lemma}

\begin{proof}
设 $U \subset X$ 为开集，并设 $U' \subset X'$ 为对应开集，即满足
$c^{-1}(U') = U$ 的开集。于是 $U$ 拟紧，当且仅当 $U'$ 拟紧，因为沿
$c$ 拉回给出 $X$ 与 $X'$ 的开集之间保持并的双射。这特别证明了 (1)。

\medskip\noindent
证明 (2)。由上可知，$X'$ 有拟紧开集基。因为 $c^{-1}$ 也与成对开集之交
可交换，可知 $X'$ 的任意两个拟紧开集之交拟紧。最后，由 (1)，$X'$
拟紧；而按构造，它清醒。因此 $X'$ 是谱空间。

\medskip\noindent
证明 (3)。$X'$ 是 Noether 的，这是立即的，因为该性质由开子集的升链
条件定义，而这一条件对 $X$ 成立。我们已经在 (2) 中看到 $X'$ 是谱空间。
\end{proof}

\section{谱空间的极限}
\label{topology-section-limits-spectral}

\noindent
引理 \ref{topology-lemma-spectral-inverse-limit-finite-sober-spaces} 告诉我们，每个
谱空间都是有限清醒空间的余滤极限。每个有限清醒空间都是谱空间，并且
有限清醒空间的每个连续映射都是谱空间之间的谱映射。本节证明一些关于
由谱映射组成的谱拓扑空间系统之极限的引理。

\begin{lemma}
\label{topology-lemma-inverse-limit-spectral-spaces-quasi-compact}
设 $\mathcal{I}$ 为范畴，$i \mapsto X_i$ 为谱空间图表，并且对
$\mathcal{I}$ 中的 $a : j \to i$，相应映射 $f_a : X_j \to X_i$ 是谱映射。
\begin{enumerate}
\item 给定在可构造拓扑中闭的子集 $Z_i \subset X_i$，并且对
$\mathcal{I}$ 中所有 $a : j \to i$ 都有 $f_a(Z_j) \subset Z_i$，
则 $\lim Z_i$ 拟紧。
\item 空间 $X = \lim X_i$ 拟紧。
\end{enumerate}
\end{lemma}

\begin{proof}
由引理 \ref{topology-lemma-limits}，极限 $Z = \lim Z_i$ 存在。以 $X'_i$
表示赋予可构造拓扑的空间 $X_i$，以 $Z'_i$ 表示 $X'_i$ 的相应子空间。
设 $a : j \to i$ 为 $\mathcal{I}$ 中的态射。因为 $f_a$ 是谱映射，
它定义连续映射 $f_a : X'_j \to X'_i$。所以
$f_a|_{Z_j} : Z'_j \to Z'_i$ 是拟紧 Hausdorff 空间的连续映射
（由引理 \ref{topology-lemma-constructible-hausdorff-quasi-compact} 与
\ref{topology-lemma-closed-in-quasi-compact}）。因此，由引理
\ref{topology-lemma-inverse-limit-quasi-compact}，$Z' = \lim Z_i$ 拟紧。
映射 $Z'_i \to Z_i$ 连续，故 $Z' \to Z$ 连续且在底集合上为双射。
所以，作为满连续映射 $Z' \to Z$ 的像，$Z$ 拟紧
（引理 \ref{topology-lemma-image-quasi-compact}）。
\end{proof}

\begin{lemma}
\label{topology-lemma-inverse-limit-spectral-spaces-nonempty}
设 $\mathcal{I}$ 为余滤范畴，$i \mapsto X_i$ 为谱空间图表，并且对
$\mathcal{I}$ 中的 $a : j \to i$，相应映射 $f_a : X_j \to X_i$ 是谱映射。
\begin{enumerate}
\item 给定在可构造拓扑中闭的非空子集 $Z_i \subset X_i$，并且对
$\mathcal{I}$ 中所有 $a : j \to i$ 都有 $f_a(Z_j) \subset Z_i$，
则 $\lim Z_i$ 非空。
\item 若每个 $X_i$ 都非空，则 $X = \lim X_i$ 非空。
\end{enumerate}
\end{lemma}

\begin{proof}
以 $X'_i$ 表示赋予可构造拓扑的空间 $X_i$，以 $Z'_i$ 表示 $X'_i$
的相应子空间。设 $a : j \to i$ 为 $\mathcal{I}$ 中的态射。因为 $f_a$
是谱映射，它定义连续映射 $f_a : X'_j \to X'_i$。所以
$f_a|_{Z_j} : Z'_j \to Z'_i$ 是拟紧 Hausdorff 空间的连续映射
（由引理 \ref{topology-lemma-constructible-hausdorff-quasi-compact} 与
\ref{topology-lemma-closed-in-quasi-compact}）。由引理 \ref{topology-lemma-nonempty-limit}，
空间 $\lim Z'_i$ 非空。由于作为集合有 $\lim Z'_i = \lim Z_i$，结论成立。
\end{proof}

\begin{lemma}
\label{topology-lemma-inverse-limit-spectral-spaces-equal}
设 $\mathcal{I}$ 为余滤范畴，$i \mapsto X_i$ 为谱空间图表，并且对
$\mathcal{I}$ 中的 $a : j \to i$，相应映射 $f_a : X_j \to X_i$ 是谱映射。
令 $X = \lim X_i$，投影为 $p_i : X \to X_i$。设
$i \in \Ob(\mathcal{I})$，并设 $E, F \subset X_i$ 为子集，其中 $E$
在可构造拓扑中闭，$F$ 在可构造拓扑中开。那么
$p_i^{-1}(E) \subset p_i^{-1}(F)$，当且仅当存在 $\mathcal{I}$ 中的态射
$a : j \to i$，使得 $f_a^{-1}(E) \subset f_a^{-1}(F)$。
\end{lemma}

\begin{proof}
注意
$$
p_i^{-1}(E) \setminus p_i^{-1}(F) =
\lim_{a : j \to i} f_a^{-1}(E) \setminus f_a^{-1}(F)
$$
因为 $f_a$ 是谱映射，它对可构造拓扑连续，故集合
$f_a^{-1}(E) \setminus f_a^{-1}(F)$ 在可构造拓扑中闭。因此可应用引理
\ref{topology-lemma-inverse-limit-spectral-spaces-nonempty}，得到左边非空，当且仅当
右边每个空间都非空。
\end{proof}

\begin{lemma}
\label{topology-lemma-inverse-limit-spectral-spaces-constructible}
设 $\mathcal{I}$ 为余滤范畴，$i \mapsto X_i$ 为谱空间图表，并且对
$\mathcal{I}$ 中的 $a : j \to i$，相应映射 $f_a : X_j \to X_i$ 是谱映射。
令 $X = \lim X_i$，投影为 $p_i : X \to X_i$。设 $E \subset X$
为可构造子集。则存在 $i \in \Ob(\mathcal{I})$ 以及可构造子集
$E_i \subset X_i$，使得 $p_i^{-1}(E_i) = E$。若 $E$ 开（相应地，闭），
则可选取开（相应地，闭）的 $E_i$。
\end{lemma}

\begin{proof}
假设 $E$ 是 $X$ 的拟紧开集。由引理 \ref{topology-lemma-describe-limits}，
可以对某个 $i$ 与某个开集 $U_i \subset X_i$ 写成
$E = p_i^{-1}(U_i)$。将 $U_i = \bigcup U_{i, \alpha}$ 写成拟紧开集之并。
因为 $E$ 拟紧，可以找到 $\alpha_1, \ldots, \alpha_n$，使得
$E = p_i^{-1}(U_{i, \alpha_1} \cup \ldots \cup U_{i, \alpha_n})$。
所以 $E_i = U_{i, \alpha_1} \cup \ldots \cup U_{i, \alpha_n}$ 合用。

\medskip\noindent
假设 $E$ 是闭可构造子集。则 $E^c$ 是拟紧开集。由上一段的结果，
对某个 $i$ 以及拟紧开集 $F_i \subset X_i$，有
$E^c = p_i^{-1}(F_i)$。于是 $E = p_i^{-1}(F_i^c)$，正合所需。

\medskip\noindent
若 $E$ 一般，可以写成 $E = \bigcup_{l = 1, \ldots, n} U_l \cap Z_l$，
其中 $U_l$ 是可构造开集，$Z_l$ 是可构造闭集。由前几段的结果，可写成
$U_l = p_{i_l}^{-1}(U_{l, i_l})$ 与
$Z_l = p_{j_l}^{-1}(Z_{l, j_l})$，其中
$U_{l, i_l} \subset X_{i_l}$ 为可构造开集，
$Z_{l, j_l} \subset X_{j_l}$ 为可构造闭集。由于 $\mathcal{I}$ 余滤，
可以选取 $\mathcal{I}$ 的对象 $k$ 以及态射 $a_l : k \to i_l$ 与
$b_l : k \to j_l$。于是取
$E_k = \bigcup_{l = 1, \ldots, n}
f_{a_l}^{-1}(U_{l, i_l}) \cap f_{b_l}^{-1}(Z_{l, j_l})$，得到 $X_k$
的一个可构造子集，其在 $X$ 中的原像为 $E$。
\end{proof}

\begin{lemma}
\label{topology-lemma-directed-inverse-limit-spectral-spaces}
设 $\mathcal{I}$ 为余滤指标范畴，$i \mapsto X_i$ 为谱空间图表，并且对
$\mathcal{I}$ 中的 $a : j \to i$，相应映射 $f_a : X_j \to X_i$ 是谱映射。
则逆极限 $X = \lim X_i$ 是谱拓扑空间，并且投影
$p_i : X \to X_i$ 是谱映射。
\end{lemma}

\begin{proof}
极限 $X = \lim X_i$ 存在（引理 \ref{topology-lemma-limits}），并由引理
\ref{topology-lemma-inverse-limit-spectral-spaces-quasi-compact} 拟紧。

\medskip\noindent
以 $p_i : X \to X_i$ 表示投影。因为 $\mathcal{I}$ 余滤，可以应用引理
\ref{topology-lemma-describe-limits}。所以 $X$ 的拓扑基由开集 $p_i^{-1}(U_i)$
给出，其中 $U_i \subset X_i$ 开。由于 $X_i$ 的拓扑基由拟紧开集给出，
可知 $X$ 的一个拓扑基由 $p_i^{-1}(U_i)$ 给出，其中
$U_i \subset X_i$ 是拟紧开集。形式论证表明，作为拓扑空间有
$$
p_i^{-1}(U_i) = \lim_{a : j \to i} f_a^{-1}(U_i)
$$
因为每个 $f_a$ 都是谱映射，集合 $f_a^{-1}(U_i)$ 在 $X_j$ 的可构造
拓扑中闭，所以由引理 \ref{topology-lemma-inverse-limit-spectral-spaces-quasi-compact}，
$p_i^{-1}(U_i)$ 拟紧。因此，$X$ 有一个由拟紧开集组成的拓扑基。

\medskip\noindent
$X$ 的任意拟紧开集 $U$ 都形如 $U = p_i^{-1}(U_i)$，其中 $i$ 为某个
指标，$U_i \subset X_i$ 为某个拟紧开集（见引理
\ref{topology-lemma-inverse-limit-spectral-spaces-constructible}）。给定拟紧开集
$U_i \subset X_i$ 与 $U_j \subset X_j$，则对某个 $k$ 以及拟紧开集
$U_k \subset X_k$，有
$p_i^{-1}(U_i) \cap p_j^{-1}(U_j) = p_k^{-1}(U_k)$。事实上，选取 $k$
以及态射 $k \to i$ 与 $k \to j$，并令 $U_k$ 为 $U_i$ 与 $U_j$
在 $X_k$ 中的拉回之交。由此可见，$X$ 的任意两个拟紧开集之交仍为
拟紧开集。

\medskip\noindent
最后，设 $Z \subset X$ 不可约且闭。则 $p_i(Z) \subset X_i$ 不可约，
因而 $Z_i = \overline{p_i(Z)}$ 有唯一泛点 $\xi_i$（因为 $X_i$ 是谱空间）。
对 $\mathcal{I}$ 中的 $a : j \to i$，有 $f_a(\xi_j) = \xi_i$，因为
$\overline{f_a(Z_j)} = Z_i$。所以 $\xi = \lim \xi_i$ 是 $X$ 的点。
断言：$\xi \in Z$。若不然，可以找到包含 $\xi$ 且与 $Z$ 不相交的
拟紧开集。它应当形如 $p_i^{-1}(U_i)$，其中 $i$ 为某个指标，
$U_i \subset X_i$ 为拟紧开集。于是 $\xi_i \in U_i$，但
$p_i(Z) \cap U_i = \emptyset$，这与 $\xi_i \in \overline{p_i(Z)}$ 矛盾。
所以 $\xi \in Z$，从而 $\overline{\{\xi\}} \subset Z$。反过来，
每个 $z \in Z$ 都在 $\xi$ 的闭包中。事实上，给定 $z$ 的拟紧开邻域
$U$，将 $U = p_i^{-1}(U_i)$ 写成上述形式，其中 $i$ 为某个指标，
$U_i \subset X_i$ 为拟紧开集。由 $p_i(z) \in U_i$ 得
$\xi_i \in U_i$，进而 $\xi \in U$。因此 $\xi$ 是 $Z$ 的泛点。
略去 $\xi$ 是 $Z$ 的唯一泛点的证明（提示：证明第二个泛点必须等于
$\xi$；为此证明它在 $X_i$ 中必映到 $\xi_i$，因为由 $X_i$ 的谱性，
不可约集 $Z_i$ 有唯一泛点）。这完成了证明。
\end{proof}

\begin{lemma}
\label{topology-lemma-descend-opens}
设 $\mathcal{I}$ 为余滤指标范畴，$i \mapsto X_i$ 为谱空间图表，并且对
$\mathcal{I}$ 中的 $a : j \to i$，相应映射 $f_a : X_j \to X_i$ 是谱映射。
置 $X = \lim X_i$，并以 $p_i : X \to X_i$ 表示投影。
\begin{enumerate}
\item 给定任意拟紧开集 $U \subset X$，存在 $i \in \Ob(\mathcal{I})$
以及拟紧开集 $U_i \subset X_i$，使得 $p_i^{-1}(U_i) = U$。
\item 给定拟紧开集 $U_i \subset X_i$ 与 $U_j \subset X_j$，并假设
$p_i^{-1}(U_i) \subset p_j^{-1}(U_j)$，则存在
$k \in \Ob(\mathcal{I})$ 以及态射 $a : k \to i$ 与 $b : k \to j$，
使得 $f_a^{-1}(U_i) \subset f_b^{-1}(U_j)$。
\item 若 $U_i, U_{1, i}, \ldots, U_{n, i} \subset X_i$ 是拟紧开集，且
$p_i^{-1}(U_i) = p_i^{-1}(U_{1, i}) \cup \ldots \cup p_i^{-1}(U_{n, i})$，
则对 $\mathcal{I}$ 中的某个态射 $a : j \to i$，有
$f_a^{-1}(U_i) = f_a^{-1}(U_{1, i}) \cup \ldots \cup f_a^{-1}(U_{n, i})$。
\item 与 (3) 相同的断言，但把并换成交。
\end{enumerate}
\end{lemma}

\begin{proof}
第 (1) 条是引理 \ref{topology-lemma-inverse-limit-spectral-spaces-constructible}
的特殊情形。第 (2) 条是引理
\ref{topology-lemma-inverse-limit-spectral-spaces-equal} 的特殊情形，因为拟紧开集
在可构造拓扑中既开又闭。第 (3)、(4) 条形式地由 (1)、(2) 以及子集取原像
与取并、取交可交换这一事实得出。
\end{proof}

\begin{lemma}
\label{topology-lemma-make-spectral-space}
设 $W$ 为谱空间 $X$ 的子集。下列条件等价：
\begin{enumerate}
\item $W$ 是可构造集之交，并且对一般化封闭；
\item $W$ 拟紧且对一般化封闭；
\item 存在拟紧子集 $E \subset X$，使得 $W$ 是特化到 $E$ 的点集；
\item $W$ 是拟紧开子集之交；
\item
\label{topology-item-intersection-quasi-compact-open}
存在非空集合 $I$ 以及拟紧开集 $U_i \subset X$（$i \in I$），使得
$W = \bigcap U_i$，并且对所有 $i, j \in I$，存在 $k \in I$ 满足
$U_k \subset U_i \cap U_j$。
\end{enumerate}
在这种情形，(a) $W$ 是谱空间，(b) 作为拓扑空间有 $W = \lim U_i$，
并且 (c) 对任意包含 $W$ 的开集 $U$，存在某个 $i$ 使得 $U_i \subset U$。
\end{lemma}

\begin{proof}
设 $W \subset X$ 满足 (1)。则 $W$ 在可构造拓扑中闭，故在可构造拓扑
中拟紧（由引理 \ref{topology-lemma-constructible-hausdorff-quasi-compact} 与
\ref{topology-lemma-closed-in-quasi-compact}），进而在 $X$ 的拓扑中拟紧
（因为 $X$ 的开集在可构造拓扑中开）。所以 (2) 成立。

\medskip\noindent
取 $E = W$，显然可见 (2) 蕴含 (3)。

\medskip\noindent
设 $X$ 为谱空间，且 $E \subset W$ 如 (3) 所述。因为 $W$ 的每个点
都特化到 $E$ 的某个点，所以 $W$ 的包含 $E$ 的开集等于 $W$。
因此，因为 $E$ 拟紧，$W$ 也拟紧。若 $x \in X$ 且 $x \not \in W$，
则 $Z = \overline{\{x\}}$ 与 $W$ 不相交。因为 $W$ 拟紧，可以找到拟紧
开集 $U$，使得 $W \subset U$ 且 $U \cap Z = \emptyset$。由此得出 (4)。

\medskip\noindent
若 $W = \bigcap_{j \in J} U_j$，则令 $I$ 等于 $J$ 的有限子集集，并对
$i = \{j_1, \ldots, j_r\}$ 令
$U_i = U_{j_1} \cap \ldots \cap U_{j_r}$，即可看出 (4) 蕴含 (5)。
(5) 蕴含 (1) 是立即的。

\medskip\noindent
设 $I$ 与 $U_i$ 如 (5) 所述。因为 $W = \bigcap U_i$，由极限的泛性质，
有 $W = \lim U_i$。于是由引理
\ref{topology-lemma-directed-inverse-limit-spectral-spaces}，$W$ 是谱空间。
设 $U \subset X$ 为 $W$ 的开邻域。于是
$E_i = U_i \cap (X \setminus U)$ 是谱空间 $Z = X \setminus U$
中的可构造子集族，其总交为空。利用 $Z$ 上的谱拓扑拟紧
% P01-ERRATA-E0025: frozen source says “spectral topology” but invokes the lemma proving quasi-compactness of the constructible topology; preserved here.
（引理 \ref{topology-lemma-constructible-hausdorff-quasi-compact}），再由引理
\ref{topology-lemma-intersection-closed-in-quasi-compact} 可知，对某个 $i$ 有
$E_i = \emptyset$。
\end{proof}

\begin{lemma}
\label{topology-lemma-make-spectral-space-minus}
设 $X$ 为谱空间，$E \subset X$ 为可构造子集。设 $W \subset X$ 为
$X$ 中特化到 $E$ 的某个点的点集。则 $W \setminus E$ 是谱空间。
若 $W = \bigcap U_i$，其中 $U_i$ 如引理 \ref{topology-lemma-make-spectral-space}
的 (\ref{topology-item-intersection-quasi-compact-open}) 所述，则
$W \setminus E = \lim (U_i \setminus E)$。
\end{lemma}

\begin{proof}
因为 $E$ 可构造，所以它拟紧，因而引理 \ref{topology-lemma-make-spectral-space}
可应用于 $W$。若 $E$ 可构造，则对所有 $i \in I$，$E$ 在 $U_i$
中可构造。因此，由引理 \ref{topology-lemma-spectral-sub}，$U_i \setminus E$
是谱空间。因为 $W \setminus E = \bigcap (U_i \setminus E)$，由极限的
泛性质有
% P01-ERRATA-E0026: frozen source omits grouping around the objects U_i minus E after lim; preserved here.
$W \setminus E = \lim U_i \setminus E$。于是由引理
\ref{topology-lemma-directed-inverse-limit-spectral-spaces}，$W \setminus E$ 是谱空间。
\end{proof}

\section{Stone-{\v C}ech 紧化}
\label{topology-section-stone-cech}

\noindent
拓扑空间 $X$ 的 Stone-{\v C}ech 紧化，是从 $X$ 到 Hausdorff 拟紧空间
$\beta(X)$ 的映射 $X \to \beta(X)$，它对这类映射具有泛性质。
我们用一个标准论证证明其存在；该论证使用下面的简单引理。

\begin{lemma}
\label{topology-lemma-dense-image}
设 $f : X \to Y$ 为拓扑空间的连续映射。假设 $f(X)$ 在 $Y$ 中稠密，
且 $Y$ 是 Hausdorff 的。则 $Y$ 的基数不超过 $P(P(X))$ 的基数，
其中 $P$ 表示幂集运算。
\end{lemma}

\begin{proof}
令 $S = f(X) \subset Y$。令 $\mathcal{D}$ 为 $Y$ 的所有闭域组成的集合，
即等于自身内部之闭包的子集 $D \subset Y$。注意，$Y$ 的开子集之闭包
是闭域。对 $y \in Y$，考虑集合
$$
I_y = \{T \subset S \mid \text{ 存在 }
D \in \mathcal{D}\text{ 其中 }T = S \cap D\text{ 且 }y \in D\}.
$$
因为 $S$ 在 $Y$ 中稠密，所以对每个闭域 $D$，$S \cap D$ 在 $D$
中稠密。因此，若对 $D, D' \in \mathcal{D}$ 有
$D \cap S = D' \cap S$，则 $D = D'$。于是，$I_y = I_{y'}$ 蕴含
$y = y'$，因为 Hausdorff 条件保证可以找到包含 $y$ 而不包含 $y'$
的闭域。结论随即成立。
\end{proof}

\noindent
设 $X$ 为拓扑空间。由引理 \ref{topology-lemma-dense-image}，存在集合 $I$，
其元素为具有稠密像的连续映射 $f : X \to Y$ 的同构类，其中 $Y$
为 Hausdorff 且拟紧。对 $i \in I$，选取代表 $f_i : X \to Y_i$。
考虑映射
$$
\prod f_i : X \longrightarrow \prod\nolimits_{i \in I} Y_i
$$
并以 $\beta(X)$ 表示其像的闭包。因为每个 $Y_i$ 都是 Hausdorff 的，
$\beta(X)$ 也是。因为每个 $Y_i$ 都拟紧，$\beta(X)$ 也拟紧
（使用定理 \ref{topology-theorem-tychonov} 与引理
\ref{topology-lemma-closed-in-quasi-compact}）。

\medskip\noindent
我们来证明典范映射 $X \to \beta(X)$ 对到 Hausdorff 拟紧空间的映射
满足泛性质。具体地，设 $f : X \to Y$ 为这种态射。令 $Z \subset Y$
为 $f(X)$ 的闭包。那么 $X \to Z$ 同构于映射 $f_i : X \to Y_i$
之一，设为 $f_{i_0} : X \to Y_{i_0}$。所以 $f$ 按所需方式分解为
$X \to \beta(X) \to \prod Y_i \to Y_{i_0} \cong Z \to Y$。

\begin{lemma}
\label{topology-lemma-one-point-compactification}
设 $X$ 为 Hausdorff 局部拟紧空间。存在映射 $X \to X^*$，它把 $X$
等同于拟紧 Hausdorff 空间 $X^*$ 的开子空间，并且
$X^* \setminus X$ 是单点集（单点紧化）。特别地，映射
$X \to \beta(X)$ 把 $X$ 等同于 $\beta(X)$ 的开子空间。
\end{lemma}

\begin{proof}
置 $X^* = X \amalg \{\infty\}$。规定 $X^*$ 的子集 $V$ 为开集，
如果或者 $V \subset X$ 在 $X$ 中开，或者 $\infty \in V$，并且
$U = V \cap X$ 是 $X$ 的开集，满足 $X \setminus U$ 拟紧。
略去验证这定义了一个拓扑。显然，
% P01-ERRATA-E0027: frozen source says the map identifies X with an open subspace of X, not X-star; preserved here.
$X \to X^*$ 把 $X$ 等同于 $X$ 的开子空间。

\medskip\noindent
因为 $X$ 局部拟紧，每个点 $x \in X$ 都有拟紧邻域
$x \in E \subset X$。于是 $E$ 闭（引理
\ref{topology-lemma-quasi-compact-in-Hausdorff} 的第 (1) 条），并且
$V = (X \setminus E) \amalg \{\infty\}$ 是 $\infty$ 的开邻域，
与 $E$ 的内部不交。因此 $X^*$ 是 Hausdorff 的。

\medskip\noindent
设 $X^* = \bigcup V_i$ 为开覆盖。则对某个 $i$（设为 $i_0$），有
$\infty \in V_{i_0}$。按构造，$Z = X^* \setminus V_{i_0}$ 拟紧。
因此覆盖 $Z \subset \bigcup_{i \not = i_0} Z \cap V_i$ 有有限加细，
这蕴含 $X^*$ 的给定覆盖有有限加细。所以 $X^*$ 拟紧。

\medskip\noindent
由 Stone-{\v C}ech 紧化的泛性质，映射 $X \to X^*$ 分解为
$X \to \beta(X) \to X^*$。令 $\varphi : \beta(X) \to X^*$ 为该分解。
那么 $X \to \varphi^{-1}(X)$ 是 $\varphi^{-1}(X) \to X$ 的截面，
所以其像闭（引理 \ref{topology-lemma-section-closed}）。由于 $X \to \beta(X)$
的像稠密，得出 $X = \varphi^{-1}(X)$。
\end{proof}

\section{极端不连通空间}
\label{topology-section-extremally-disconnected}

\noindent
本节材料取自 \cite{Gleason}（并按 \cite{Rainwater} 作了轻微修改）。
Gleason 的论文证明，在拟紧 Hausdorff 空间范畴中，“投射对象”恰好是
极端不连通空间。

\begin{definition}
\label{topology-definition-extremally-disconnected}
若拓扑空间 $X$ 的每个开子集之闭包仍开，则称 $X$ 为
{\it 极端不连通的}。
\end{definition}

\noindent
若 $X$ 是 Hausdorff 且极端不连通的，则 $X$ 全不连通
（一般情形下这不成立）。若 $X$ 拟紧、Hausdorff 且极端不连通，则由引理
\ref{topology-lemma-profinite}，$X$ 仿有限；但逆命题一般不成立。例如，$p$-进整数
$\mathbf{Z}_p = \lim \mathbf{Z}/p^n\mathbf{Z}$ 是并非极端不连通的
仿有限空间。事实上，若 $U \subset \mathbf{Z}_p$ 为所有赋值为偶数的
非零元素构成的集合，则 $U$ 开，但其闭包 $U \cup \{0\}$ 不开。

\begin{lemma}
\label{topology-lemma-image-open-technical}
设 $f : X \to Y$ 为拓扑空间的连续映射。假设 $f$ 满射，并且对每个真闭
子集 $E \subset X$ 都有 $f(E) \not = Y$。那么，对开集 $U \subset X$，
子集 $f(U)$ 包含于 $Y \setminus f(X \setminus U)$ 的闭包。
\end{lemma}

\begin{proof}
取 $y \in f(U)$，并设 $V \subset Y$ 为 $y$ 的任意开邻域。我们证明
$V$ 与 $Y \setminus f(X \setminus U)$ 相交。注意，
$W = U \cap f^{-1}(V)$ 是 $X$ 的非空开子集，所以
$f(X \setminus W) \not = Y$。取 $y' \in Y$，使得
$y' \not \in f(X \setminus W)$。容易证明 $y' \in V$ 且
$y' \in Y \setminus f(X \setminus U)$。
\end{proof}

\begin{lemma}
\label{topology-lemma-intersection-empty}
设 $X$ 为极端不连通空间。若 $U, V \subset X$ 是不交开子集，则
$\overline{U}$ 与 $\overline{V}$ 也不交。
\end{lemma}

\begin{proof}
由假设，$\overline{U}$ 开，所以 $V \cap \overline{U}$ 是与 $U$ 不交的
开集；它必为空，因为 $\overline{U}$ 是 $X$ 中所有包含 $U$ 的闭子集之交。
这意味着开集 $\overline{V} \cap \overline{U}$ 避开 $V$，故由同一论证
它为空。
\end{proof}

\begin{lemma}
\label{topology-lemma-isomorphism}
设 $f : X \to Y$ 为 Hausdorff 拟紧拓扑空间的连续映射。若 $Y$
极端不连通，$f$ 满射，并且对 $X$ 的每个真闭子集 $Z$ 都有
$f(Z) \not = Y$，则 $f$ 是同胚。
\end{lemma}

\begin{proof}
由引理 \ref{topology-lemma-bijective-map}，只需证明 $f$ 单射。假设
$x, x' \in X$ 为不同的点，且 $y = f(x) = f(x')$。选取 $x, x'$
的不交开邻域 $U, U' \subset X$。注意，$f$ 是闭映射
（引理 \ref{topology-lemma-closed-map}），故
$T = f(X \setminus U)$ 与 $T' = f(X \setminus U')$ 在 $Y$ 中闭。
因为 $X$ 是 $X \setminus U$ 与 $X \setminus U'$ 之并，可知
$Y = T \cup T'$。由引理 \ref{topology-lemma-image-open-technical}，$y$ 同时属于
$Y \setminus T$ 的闭包与 $Y \setminus T'$ 的闭包。另一方面，由引理
\ref{topology-lemma-intersection-empty}，二者之交为空。这样便得到所需矛盾。
\end{proof}

\begin{lemma}
\label{topology-lemma-find-compact-subset}
设 $f : X \to Y$ 为 Hausdorff 拟紧拓扑空间之间的满连续映射。
则存在拟紧子集 $E \subset X$，使得 $f(E) = Y$，但对所有真闭子集
$E' \subset E$ 都有 $f(E') \not = Y$。
\end{lemma}

\begin{proof}
我们将不再特别说明地使用：$X$ 的拟紧子集恰好是其闭子集
（引理 \ref{topology-lemma-closed-in-compact}）。考虑所有满足 $f(E) = Y$
的拟紧子集 $E \subset X$ 按包含关系组成的集合 $\mathcal{E}$。
我们用 Zorn 引理证明 $\mathcal{E}$ 有极小元。为此，只需证明：
给定 $\mathcal{E}$ 中的全序族 $E_\lambda$，交集
$\bigcap E_\lambda$ 属于 $\mathcal{E}$。它是闭的，因而拟紧。
对每个 $y \in Y$，集合 $E_\lambda \cap f^{-1}(\{y\})$ 非空且闭，
故由引理 \ref{topology-lemma-intersection-closed-in-quasi-compact}，交集
$\bigcap E_\lambda \cap f^{-1}(\{y\}) = \bigcap (E_\lambda \cap f^{-1}(\{y\}))$
非空。这完成了证明。
\end{proof}

\begin{proposition}
\label{topology-proposition-projective-in-category-hausdorff-qc}
设 $X$ 为 Hausdorff 拟紧拓扑空间。下列条件等价：
\begin{enumerate}
\item $X$ 极端不连通；
\item 对任意满连续映射 $f : Y \to X$，若 $Y$ 是 Hausdorff 拟紧的，
则存在连续截面；
\item 对拟紧 Hausdorff 空间连续映射的任意实线交换图
$$
\xymatrix{
& Y \ar[d] \\
X \ar@{..>}[ru] \ar[r] & Z
}
$$
若 $Y \to Z$ 满射，则在拓扑空间范畴中存在使图表交换的点线箭头。
\end{enumerate}
\end{proposition}

\begin{proof}
显然，(3) 蕴含 (2)。另一方面，若 (2) 成立，并且 $X \to Z$ 与
$Y \to Z$ 如 (3) 所述，则 (2) 保证投影 $X \times_Z Y \to X$ 有截面，
这蕴含存在合适的点线箭头（略去细节）。所以 (3) 与 (2) 等价。

\medskip\noindent
假设 $X$ 极端不连通，并设 $f : Y \to X$ 如 (2) 所述。由引理
\ref{topology-lemma-find-compact-subset}，存在拟紧子集 $E \subset Y$，使得
$f(E) = X$，但对所有真闭子集 $E' \subset E$ 都有 $f(E') \not = X$。
由引理 \ref{topology-lemma-isomorphism}，$f|_E : E \to X$ 是同胚，其逆给出
所需截面。

\medskip\noindent
假设 (2)。设 $U \subset X$ 为开集，补集为 $Z$。考虑满连续映射
$f : \overline{U} \amalg Z \to X$。设 $\sigma$ 为一个截面。于是
$\overline{U} = \sigma^{-1}(\overline{U})$ 开。因此 $X$ 极端不连通。
\end{proof}

\begin{lemma}
\label{topology-lemma-rainwater}
设 $f : X \to X$ 为 Hausdorff 拓扑空间的满连续自映射。若 $f$ 不是
$\text{id}_X$，则存在真闭子集 $E \subset X$，使得 $X = E \cup f(E)$。
\end{lemma}

\begin{proof}
取 $p \in X$，使得 $f(p) \not = p$。选取不交开邻域
$p \in U$、$f(p) \in V$，并置
% P01-ERRATA-E0028: frozen source omits parentheses around U intersect f^{-1}(V); the subsequent complement argument requires them, but the formula is preserved here.
$E = X \setminus U \cap f^{-1}(V)$。于是 $p \not \in E$，故 $E$ 是真闭
子集。若 $x \in X$，则或者 $x \in E$；否则
$x \in U \cap f^{-1}(V)$。写成 $x = f(y)$（因为 $f$ 满射，这是可能的）。
若 $y \in U \cap f^{-1}(V)$，则会有 $x = f(y) \in V$，这与
$x \in U$ 矛盾。所以 $y \in E$，并且 $x \in f(E)$。
\end{proof}

\begin{example}
\label{topology-example-stone-Cech-discrete}
可以用命题 \ref{topology-proposition-projective-in-category-hausdorff-qc} 看出，
离散空间 $X$ 的 Stone-{\v C}ech 紧化 $\beta(X)$ 极端不连通。
事实上，设 $f : Y \to \beta(X)$ 为连续满射，其中 $Y$ 拟紧且 Hausdorff。
因为 $X$ 离散，可以把映射 $X \to \beta(X)$ 提升为连续（！）映射
$X \to Y$。由 Stone-{\v C}ech 紧化的泛性质，得到分解
$X \to \beta(X) \to Y$。因为 $\beta(X) \to Y \to \beta(X)$ 在稠密
子集 $X$ 上等于恒等映射，可知我们得到了一个截面。特别地，离散空间的
Stone-{\v C}ech 紧化全不连通，从而仿有限（见定义
\ref{topology-definition-extremally-disconnected} 后的讨论与引理 \ref{topology-lemma-profinite}）。
\end{example}

\noindent
利用例 \ref{topology-example-stone-Cech-discrete} 所提供的大量极端不连通空间，
可以证明每个拟紧 Hausdorff 空间在拟紧 Hausdorff 空间范畴中都有
“投射覆盖”。

\begin{lemma}
\label{topology-lemma-existence-projective-cover}
\begin{slogan}
每个拟紧 Hausdorff 空间都有典范的极端不连通覆盖
\end{slogan}
设 $X$ 为拟紧 Hausdorff 空间。存在连续满射 $X' \to X$，其中 $X'$
拟紧、Hausdorff 且极端不连通。若要求 $X'$ 的每个真闭子集都不映满
$X$，则 $X'$ 在同构意义下唯一。
\end{lemma}

\begin{proof}
令 $Y = X$，但赋予离散拓扑。令 $X' = \beta(Y)$。连续映射
$Y \to X$ 分解为 $Y \to X' \to X$。由例
\ref{topology-example-stone-Cech-discrete}，这证明了引理的第一个断言。

\medskip\noindent
由引理 \ref{topology-lemma-find-compact-subset}，可以找到拟紧子集 $E \subset X'$，
它满射到 $X$，且 $E$ 的任何真闭子集都不满射到 $X$。因为 $X'$
极端不连通，存在 $X$ 上的连续映射 $f : X' \to E$
（命题 \ref{topology-proposition-projective-in-category-hausdorff-qc}）。把 $f$ 与映射
$E \to X'$ 复合，得到连续自映射 $f|_E : E \to E$。注意，$f|_E$
必定满射，否则其像会是真闭子集并满射到 $X$。所以 $f|_E$ 必须是
$\text{id}_E$，否则引理 \ref{topology-lemma-rainwater} 表明 $E$ 并非极小。
因此，$\text{id}_E$ 经过极端不连通空间 $X'$ 分解。形式的范畴论论证
（使用命题 \ref{topology-proposition-projective-in-category-hausdorff-qc} 的刻画）
表明 $E$ 极端不连通。

\medskip\noindent
% P01-ERRATA-E0029: frozen proof does not explicitly replace X' by the minimal E before invoking minimality of X' in the uniqueness argument.
为证明唯一性，假设有第二个极小覆盖 $X'' \to X$。由命题
\ref{topology-proposition-projective-in-category-hausdorff-qc} 中证明的提升性质，
可以找到 $X$ 上的连续映射 $g : X' \to X''$。注意，$g$ 是闭映射
（引理 \ref{topology-lemma-closed-map}）。因此 $g(X') \subset X''$ 是满射到
$X$ 的闭子集；由 $X''$ 的极小性，得出 $g(X') = X''$。另一方面，
若 $E \subset X'$ 是真闭子集，则 $g(E) \not = X''$，因为由 $X'$
的极小性，$E$ 不满射到 $X$。由引理 \ref{topology-lemma-isomorphism}，$g$ 是同构。
\end{proof}

\begin{remark}
\label{topology-remark-size-projective-cover}
设 $X$ 为拟紧 Hausdorff 空间，并设 $\kappa$ 为不小于 $X$ 基数的无限
基数。那么，极小拟紧 Hausdorff 极端不连通覆盖
$X' \to X$（引理 \ref{topology-lemma-existence-projective-cover}）的基数至多为
$2^{2^\kappa}$。事实上，选取子集 $S \subset X'$，使其双射地映到 $X$。
由 $X'$ 的极小性，集合 $S$ 在 $X'$ 中稠密。因此，由引理
\ref{topology-lemma-dense-image}，$|X'| \leq 2^{2^\kappa}$。
\end{remark}

\section{杂项}
\label{topology-section-miscellany}


\noindent
下面的引理适用于与拟分离概形关联的底拓扑空间。

\begin{lemma}
\label{topology-lemma-topology-quasi-separated-scheme}
设 $X$ 为满足下列条件的拓扑空间：
\begin{enumerate}
\item 有一个由拟紧开集组成的拓扑基；
\item 任意两个拟紧开集之交拟紧。
\end{enumerate}
则
\begin{enumerate}
\item $X$ 局部拟紧；
\item 拟紧开集 $U \subset X$ 逆紧；
\item 任意拟紧开集 $U \subset X$ 都有一个开覆盖的共尾系
$\mathcal{U} : U = \bigcup_{j\in J} U_j$，其中 $J$ 有限，并且所有
$U_j$ 与 $U_j \cap U_{j'}$ 都拟紧；
% P01-ERRATA-E0030: frozen source contains the editorial placeholder “add more here”; translated literally.
\item 在此添加更多条件。
\end{enumerate}
\end{lemma}

\begin{proof}
略。
\end{proof}

\begin{definition}
\label{topology-definition-isolated-point}
设 $X$ 为拓扑空间。若 $\{x\}$ 在 $X$ 中开，则称 $x \in X$ 是 $X$
的{\it 孤立点}。
\end{definition}



\section{分划与分层}
\label{topology-section-stratifications}

\noindent
分层可以用许多不同方式定义。欢迎对本节的定义选择提出意见。

\begin{definition}
\label{topology-definition-paritition}
设 $X$ 为拓扑空间。$X$ 的一个{\it 分划}是分解
$X = \coprod X_i$，其中 $X_i$ 是局部闭子集。$X_i$ 称为该分划的
{\it 分块}。给定 $X$ 的两个分划，若其中一个的各分块都是另一个的
若干分块之并，则称后者{\it 加细}前者。
\end{definition}

\noindent
任意拓扑空间 $X$ 都有按连通分支给出的分划。若 $X$ 有有限多个不可约
分支 $Z_1, \ldots, Z_r$，则有一个分块为
$X_I = \bigcap_{i \in I} Z_i \setminus (\bigcup_{i \not \in I} Z_i)$
的分划，其指标是子集 $I \subset \{1, \ldots, r\}$；该分划加细按连通
分支给出的分划。

\begin{definition}
\label{topology-definition-good-stratification}
设 $X$ 为拓扑空间。$X$ 的一个{\it 良好分层}是分划
$X = \coprod X_i$，满足对所有 $i, j \in I$ 都有
$$
X_i \cap \overline{X_j} \not = \emptyset
\Rightarrow
X_i \subset \overline{X_j}.
$$
\end{definition}

\noindent
给定良好分层 $X = \coprod_{i \in I} X_i$，通过规定 $i \leq j$
当且仅当 $X_i \subset \overline{X_j}$，在 $I$ 上得到偏序。于是
$$
\overline{X_j} = \bigcup\nolimits_{i \leq j} X_i
$$
然而，代数几何中经常出现的情形是：在最后一个展示公式中，左边只是右边
的子集。这引出下面的定义。

\begin{definition}
\label{topology-definition-stratification}
设 $X$ 为拓扑空间。$X$ 的一个{\it 分层}由分划
$X = \coprod_{i \in I} X_i$ 以及 $I$ 上的一个偏序给出，使得对每个
$j \in I$ 都有
$$
\overline{X_j} \subset \bigcup\nolimits_{i \leq j} X_i
$$
分块 $X_i$ 称为该分层的{\it 层}。
\end{definition}

\noindent
我们常对分层附加条件。例如，若分层{\it 局部有限}，即每个点都有只与
有限多个层相交的邻域，则它尤其良好。更一般地，引入下面的定义。

\begin{definition}
\label{topology-definition-locally-finite}
设 $X$ 为拓扑空间，$I$ 为集合，并且对 $i \in I$ 设
$E_i \subset X$ 为子集。若对所有 $x \in X$，都存在 $x$ 的开邻域 $U$，
使得 $\{i \in I | E_i \cap U \not = \emptyset\}$ 有限，则称集合族
$\{E_i\}_{i \in I}$ {\it 局部有限}。
\end{definition}


\begin{remark}
\label{topology-remark-locally-finite-stratification}
给定拓扑空间 $X$ 的局部有限分层 $X = \coprod X_i$，得到由 $I$ 指标化的
$X$ 的闭子集族 $Z_i = \bigcup_{j \leq i} X_j$，并且
$$
Z_i \cap Z_j = \bigcup\nolimits_{k \leq i, j} Z_k
$$
反之，给定由偏序集 $I$ 指标化的闭子集 $Z_i \subset X$，假设
$X = \bigcup Z_i$，每个点都有只与有限多个 $Z_i$ 相交的邻域，并且上面的
展示公式成立；那么，令 $X_i = Z_i \setminus \bigcup_{j < i} Z_j$，
得到 $X$ 的一个局部有限分层。
\end{remark}

\begin{lemma}
\label{topology-lemma-partition-refined-by-stratification}
设 $X$ 为拓扑空间，$X = \coprod X_i$ 为 $X$ 的有限分划。
则存在加细它的 $X$ 的有限分层。
\end{lemma}

\begin{proof}
令 $T_i = \overline{X_i}$ 且 $\Delta_i = T_i \setminus X_i$。令 $S$
为所有 $T_i$ 与 $\Delta_i$ 的交所组成的集合。（例如
$T_1 \cap T_2 \cap \Delta_4$ 是 $S$ 的一个元素。）那么
$S = \{Z_s\}$ 是 $X$ 的有限闭子集族，并且对所有 $s, s' \in S$，
有 $Z_s \cap Z_{s'} \in S$。在 $S$ 上定义包含偏序。然后令
$Y_s = Z_s \setminus \bigcup_{s' < s} Z_{s'}$，便得到所需分层。
\end{proof}

\begin{lemma}
\label{topology-lemma-constructible-partition-refined-by-stratification}
设 $X$ 为拓扑空间。假设 $X = T_1 \cup \ldots \cup T_n$ 被写成可构造
子集之并。则存在有限分层 $X = \coprod X_i$，其中每个 $X_i$ 可构造，
并且每个 $T_k$ 都是若干层之并。
\end{lemma}

\begin{proof}
由可构造子集的定义，可将每个 $T_i$ 写成有限多个 $U \cap V^c$ 之并，
其中 $U, V \subset X$ 是逆紧开集。因此可假设
$T_i = U_i \cap V_i^c$，其中 $U_i, V_i \subset X$ 是逆紧开集。
令 $S$ 为 $X$ 的有限闭子集集，由
$\emptyset, X, U_i^c, V_i^c$ 以及它们的有限交组成。若 $Z \in S$，
则 $Z$ 在 $X$ 中可构造（引理 \ref{topology-lemma-constructible}）。此外，对所有
$Z, Z' \in S$，有 $Z \cap Z' \in S$。在 $S$ 上定义包含偏序。
对 $Z \in S$，令
$X_Z = Z \setminus \bigcup_{Z' < Z,\ Z' \in S} Z'$，得到满足引理所述
性质的分层 $X = \coprod_{Z \in S} X_Z$。
\end{proof}

\begin{lemma}
\label{topology-lemma-noetherian-partition-refined-by-stratification}
设 $X$ 为 Noether 拓扑空间。$X$ 的任意有限分划都可由有限良好分层加细。
\end{lemma}

\begin{proof}
设 $X = \coprod X_i$ 为 $X$ 的有限分划，并设 $Z$ 为 $X$ 的一个不可约
分支。由于 $X = \bigcup \overline{X_i}$，且指标集有限，存在某个 $i$
使得 $Z \subset \overline{X_i}$。因为 $X_i$ 局部闭，这蕴含
$Z \cap X_i$ 包含 $Z$ 的一个开集。所以 $Z \cap X_i$ 包含 $X$ 的一个
开集 $U$（引理 \ref{topology-lemma-Noetherian}）。写成
$X_i = U \amalg X_i^1 \amalg X_i^2$，其中
$X_i^1 = (X_i \setminus U) \cap \overline{U}$ 且
$X_i^2 = (X_i \setminus U) \cap \overline{U}^c$。对 $i' \not = i$，令
$X_{i'}^1 = X_{i'} \cap \overline{U}$ 且
$X_{i'}^2 = X_{i'} \cap \overline{U}^c$。于是
$$
X \setminus U = \coprod X^k_l
$$
是一个满足 $\overline{U} \setminus U = \bigcup X_l^1$ 的分划。
注意，$X \setminus U$ 闭且严格小于 $X$。由 Noether 归纳法，可以把这个
分划加细为有限良好分层
$X \setminus U = \coprod_{\alpha \in A} T_\alpha$。于是
$X = U \amalg \coprod_{\alpha \in A} T_\alpha$ 是加细原分划的 $X$
的有限良好分层。
\end{proof}







\section{空间的余极限}
\label{topology-section-colimits}

\noindent
拓扑空间范畴有余积。具体地，若 $I$ 是集合，并且对 $i \in I$ 给定拓扑
空间 $X_i$，则在集合 $\coprod_{i \in I} X_i$ 上赋予{\it 余积拓扑}。
我们用形如 $U_i$ 的集合构成这一拓扑的基，其中 $U_i \subset X_i$ 开。

\medskip\noindent
拓扑空间范畴有余等化子。具体地，若 $a, b : X \to Y$ 为拓扑空间的态射，
则 $a$ 与 $b$ 的余等化子，是集合范畴中的余等化子 $Y/\sim$ 赋予商拓扑
（第 \ref{topology-section-submersive} 节）。

\begin{lemma}
\label{topology-lemma-colimits}
拓扑空间范畴有余极限，并且到集合范畴的遗忘函子与余极限可交换。
\end{lemma}

\begin{proof}
这由上面的讨论以及 Categories，引理
\ref{categories-lemma-colimits-coproducts-coequalizers} 得出。
Categories，评注 \ref{categories-remark-how-to-use-it} 中勾勒了余极限存在性的
另一个证明。由上可知，遗忘函子与余极限可交换。另一种看法是使用
Categories，引理 \ref{categories-lemma-adjoint-exact}，以及遗忘函子有右伴随
这一事实；该右伴随把集合赋予相应的混沌（或无差别）拓扑空间。
\end{proof}

\section{拓扑群、环与模}
\label{topology-section-topological-groups}

\noindent
本节仅简要给出一些定义和初等性质。

\begin{definition}
\label{topology-definition-topological-group}
所谓一个{\it 拓扑群}，是指赋予了拓扑的群 $G$，使得乘法
$G \times G \to G$，$(x, y) \mapsto xy$ 与取逆映射
$G \to G$，$x \mapsto x^{-1}$ 都连续。
所谓一个{\it 拓扑群同态}，是指连续的群同态。
\end{definition}

\noindent
若 $G$ 是拓扑群且 $H \subset G$ 是子群，则 $H$ 赋予诱导拓扑后是
拓扑群。若 $G$ 是拓扑群且 $G \to H$ 是群的满射，则 $H$ 赋予商拓扑后
是拓扑群。

\begin{example}
\label{topology-example-automorphisms-of-a-set}
设 $E$ 为集合。我们可以在自映射集 $\text{Map}(E, E)$ 上赋予紧开拓扑；
也就是说，对给定的 $f : E \to E$，由集合
$U_S(f) = \{f' : E \to E \mid f'|_S = f|_S\}$ 构成 $f$ 的一个邻域基本系，
其中 $S \subset E$ 有限。当 $E$ 赋予离散拓扑时，作用
$$
\text{Map}(E, E) \times E \longrightarrow E,\quad
(f, e) \longmapsto f(e)
$$
连续。若 $X$ 是拓扑空间且 $X \times E \to E$ 是连续映射，则映射
$X \to \text{Map}(E, E)$ 连续。换言之，紧开拓扑是使上面所示的“作用”
映射连续的最粗拓扑。复合映射
$$
\text{Map}(E, E) \times \text{Map}(E, E) \to \text{Map}(E, E)
$$
也连续（利用上面对邻域的描述很容易验证）。最后，若
$\text{Aut}(E) \subset \text{Map}(E, E)$ 是可逆映射的子集，则取逆映射
$i : \text{Aut}(E) \to \text{Aut}(E)$，$f \mapsto f^{-1}$ 也连续。事实上，
设 $S \subset E$ 有限，则
$i^{-1}(U_S(f^{-1})) = U_{f^{-1}(S)}(f)$。因此
$\text{Aut}(E)$ 是定义 \ref{topology-definition-topological-group} 所述的拓扑群。
\end{example}

\begin{lemma}
\label{topology-lemma-topological-group-limits}
拓扑群范畴有极限，并且极限与到以下范畴的遗忘函子可交换：
(a) 拓扑空间范畴；(b) 群范畴。
\end{lemma}

\begin{proof}
只需证明积和等化子的存在性及交换性，见 Categories，引理
\ref{categories-lemma-limits-products-equalizers}。设 $G_i$，$i \in I$ 是一族
拓扑群。取通常的积 $G = \prod G_i$，并赋予积拓扑。由于作为拓扑空间有
$G \times G = \prod (G_i \times G_i)$（因为在任意范畴中，积与积可交换），
所以 $G$ 上的乘法连续。取逆映射也同样连续。设
$a, b : G \to H$ 是两个拓扑群同态。作为等化子，我们可以直接取 $a$ 与
$b$ 作为拓扑空间映射的等化子；这与把它们视为群同态时的等化子赋予诱导
拓扑所得对象相同。
\end{proof}

\begin{lemma}
\label{topology-lemma-profinite-group}
设 $G$ 为拓扑群。下列条件等价：
\begin{enumerate}
\item 作为拓扑空间，$G$ 是仿有限的；
\item $G$ 是有限离散拓扑群所成图表的极限；
\item $G$ 是有限离散拓扑群的余滤极限。
\end{enumerate}
\end{lemma}

\begin{proof}
对于拓扑空间，我们已有相应结果，见引理 \ref{topology-lemma-profinite}。结合引理
\ref{topology-lemma-topological-group-limits}，可知只需证明 (1) 蕴含 (3)。

\medskip\noindent
我们先证明，单位元 $e$ 的每个邻域 $E$ 都包含一个开子群。事实上，由于
$G$ 是有限离散拓扑空间的余滤极限（引理 \ref{topology-lemma-profinite}），可以选择
到有限离散空间 $T$ 的连续映射 $f : G \to T$，使得
$f^{-1}(f(\{e\})) \subset E$。考虑
$$
H = \{g \in G \mid f(gg') = f(g')\text{ 对所有 }g' \in G\}
$$
这是 $G$ 的一个子群，并且包含于 $E$。因此只需证明 $H$ 是开的。
% P01-ERRATA: P01-E0032 -- the frozen proof does not handle t outside f(G), when W is empty.
取
$t \in T$，并令 $W = f^{-1}(\{t\})$。注意 $W \subset G$ 是开闭的，
尤其是拟紧的。对每个 $w \in W$，存在开邻域 $e \in U_w \subset G$ 和
$w \in U'_w \subset W$，使得 $U_wU'_w \subset W$。由拟紧性，可以找到
$w_1, \ldots, w_n$，使得 $W = \bigcup U'_{w_i}$。于是
% P01-ERRATA: P01-E0033 -- the frozen assertion leaves g free; it needs g in U_t.
$U_t = U_{w_1} \cap \ldots \cap U_{w_n}$ 是 $e$ 的开邻域，并且对所有
$w \in W$ 都有 $f(gw) = t$。由于 $T$ 有限，可知
$\bigcap_{t \in T} U_t \subset H$ 是 $e$ 的开邻域。又因为
$H \subset G$ 是子群，所以 $H$ 是开的。

\medskip\noindent
假设 $H \subset G$ 是开子群。由于 $G$ 拟紧，可知 $H$ 在 $G$ 中的指数
有限。设 $G = Hg_1 \cup \ldots \cup Hg_n$。则
% P01-ERRATA: P01-E0034 -- these are left-coset representatives, but the written conjugates use the right-coset orientation.
$N = \bigcap_{i = 1, \ldots, n} g_iHg_i^{-1}$ 是包含于 $H$ 的开正规子群。
由于 $N$ 的指数也有限，可知 $G \to G/N$ 是到一个有限离散拓扑群的满射。

\medskip\noindent
考虑映射
$$
G \longrightarrow \lim_{N \subset G\text{ 开且正规}} G/N
$$
我们断言该映射是拓扑群的同构。这就完成了引理的证明，因为右侧的极限是
余滤的（两个开正规子群之交仍是开正规子群）。该映射连续，因为每个
$G \to G/N$ 都连续。该映射为单射，因为 $G$ 是 Hausdorff 的，并且由
上面的论证，$e$ 的每个邻域都包含某个 $N$。由引理
\ref{topology-lemma-intersection-closed-in-quasi-compact}，该映射为满射。由引理
\ref{topology-lemma-bijective-map}，该映射是同胚。
\end{proof}

\begin{definition}
\label{topology-definition-profinite-group}
若一个拓扑群满足引理 \ref{topology-lemma-profinite-group} 中的等价条件，则称它为
{\it 仿有限群}。
\end{definition}

\noindent
若 $G_1 \to G_2 \to G_3 \to \ldots$ 是一个拓扑群系统，则它在拓扑群中
的余极限 $G = \colim G_n$（引理 \ref{topology-lemma-topological-group-colimits}）
一般不同于它在拓扑空间中的余极限（引理 \ref{topology-lemma-colimits}），尽管二者
具有相同的底层集合。见 Examples，第 \ref{examples-section-colimit-topology} 节。

\begin{lemma}
\label{topology-lemma-topological-group-colimits}
拓扑群范畴有余极限，并且余极限与到群范畴的遗忘函子可交换。
\end{lemma}

\begin{proof}
我们将用 Categories，评注 \ref{categories-remark-how-to-use-it} 中的论证来
证明余极限存在。具体地，假设
% P01-ERRATA: P01-E0031 -- frozen source writes Top although the stated codomain is TopGroup.
$\mathcal{I} \to \textit{Top}$，$i \mapsto G_i$ 是一个取值于拓扑群范畴
$\textit{TopGroup}$ 的函子。于是可以考虑
$$
F : \textit{TopGroup} \longrightarrow \textit{Sets},\quad
H \longmapsto \lim_\mathcal{I} \Mor_{\textit{TopGroup}}(G_i, H)
$$
该函子与极限可交换。此外，给定任意拓扑群 $H$ 以及 $F(H)$ 的一个元素
$(\varphi_i : G_i \to H)$，存在基数不超过 $|\coprod G_i|$ 的子群
$H' \subset H$（这里的余积取在群范畴中，即诸 $G_i$ 的自由积），使得
同态 $\varphi_i$ 的像都落在 $H'$ 中。事实上，可以取由诸
$\varphi_i$ 的像生成的子群，并赋予诱导拓扑。因此，Categories，引理
\ref{categories-lemma-a-version-of-brown} 的假设显然满足，从而得到表示函子
$F$ 的拓扑群 $G$；这恰好意味着 $G$ 是图表 $i \mapsto G_i$ 的余极限。

\medskip\noindent
为证明它与到群范畴的遗忘函子可交换，我们使用 Categories，引理
\ref{categories-lemma-adjoint-exact}。事实上，该遗忘函子有右伴随，即把一个群
赋予相应的混沌（或无差别）拓扑群的函子。
\end{proof}

\begin{definition}
\label{topology-definition-topological-ring}
所谓一个{\it 拓扑环}，是指赋予了拓扑的环 $R$，使得加法
$R \times R \to R$，$(x, y) \mapsto x + y$ 和乘法
$R \times R \to R$，$(x, y) \mapsto xy$ 都连续。
所谓一个{\it 拓扑环同态}，是指连续的环同态。
\end{definition}

\noindent
在 Stacks Project 中，环是有 $1$ 的交换环。若 $R$ 是拓扑环，则
$(R, +)$ 是拓扑群，因为 $x \mapsto -x$ 连续。若 $R$ 是拓扑环且
$R' \subset R$ 是子环，则 $R'$ 赋予诱导拓扑后是拓扑环。若 $R$ 是拓扑环
且 $R \to R'$ 是环的满射，则 $R'$ 赋予商拓扑后是拓扑环。

\begin{lemma}
\label{topology-lemma-topological-ring-limits}
拓扑环范畴有极限，并且极限与到以下范畴的遗忘函子可交换：
(a) 拓扑空间范畴；(b) 环范畴。
\end{lemma}

\begin{proof}
只需证明积和等化子的存在性及交换性，见 Categories，引理
\ref{categories-lemma-limits-products-equalizers}。设 $R_i$，$i \in I$ 是一族
拓扑环。取通常的积 $R = \prod R_i$，并赋予积拓扑。由于作为拓扑空间有
$R \times R = \prod (R_i \times R_i)$（因为在任意范畴中，积与积可交换），
所以 $R$ 上的加法和乘法连续。设 $a, b : R \to R'$ 是两个拓扑环同态。
作为等化子，我们可以直接取 $a$ 与 $b$ 作为拓扑空间映射的等化子；这与
把它们视为环同态时的等化子赋予诱导拓扑所得对象相同。
\end{proof}

\begin{lemma}
\label{topology-lemma-topological-ring-colimits}
拓扑环范畴有余极限，并且余极限与到环范畴的遗忘函子可交换。
\end{lemma}

\begin{proof}
与引理 \ref{topology-lemma-topological-group-colimits} 的证明中所用完全相同的论证表明
余极限存在。为证明它与到环范畴的遗忘函子可交换，我们使用 Categories，
引理 \ref{categories-lemma-adjoint-exact}。事实上，该遗忘函子有右伴随，即把
一个环赋予相应的混沌（或无差别）拓扑环的函子。
\end{proof}

\begin{definition}
\label{topology-definition-topological-module}
设 $R$ 为拓扑环。所谓一个{\it 拓扑模}，是指赋予了拓扑的 $R$-模 $M$，
使得加法 $M \times M \to M$ 与纯量乘法 $R \times M \to M$ 都连续。
所谓一个{\it 拓扑模同态}，是指连续的模同态。
\end{definition}

\noindent
若 $R$ 是拓扑环且 $M$ 是拓扑模，则 $(M, +)$ 是拓扑群，因为
$x \mapsto -x$ 连续。若 $R$ 是拓扑环，$M$ 是拓扑模且
$M' \subset M$ 是子模，则 $M'$ 赋予诱导拓扑后是拓扑模。若 $R$ 是拓扑环，
$M$ 是拓扑模且 $M \to M'$ 是模的满射，则 $M'$ 赋予商拓扑后是拓扑模。

\begin{lemma}
\label{topology-lemma-topological-module-limits}
设 $R$ 为拓扑环。$R$ 上的拓扑模范畴有极限，并且极限与到以下范畴的遗忘
函子可交换：(a) 拓扑空间范畴；(b) $R$-模范畴。
\end{lemma}

\begin{proof}
只需证明积和等化子的存在性及交换性，见 Categories，引理
\ref{categories-lemma-limits-products-equalizers}。设 $M_i$，$i \in I$ 是一族
$R$ 上的拓扑模。取通常的积 $M = \prod M_i$，并赋予积拓扑。由于作为
拓扑空间有 $M \times M = \prod (M_i \times M_i)$（因为在任意范畴中，
积与积可交换），所以 $M$ 上的加法连续。乘法 $R \times M \to M$ 也同样
连续。设 $a, b : M \to M'$ 是两个 $R$ 上的拓扑模同态。作为等化子，我们
可以直接取 $a$ 与 $b$ 作为拓扑空间映射的等化子；这与把它们视为模同态时
的等化子赋予诱导拓扑所得对象相同。
\end{proof}

\begin{lemma}
\label{topology-lemma-topological-module-colimits}
设 $R$ 为拓扑环。$R$ 上的拓扑模范畴有余极限，并且余极限与到 $R$-模范畴
的遗忘函子可交换。
\end{lemma}

\begin{proof}
与引理 \ref{topology-lemma-topological-group-colimits} 的证明中所用完全相同的论证表明
余极限存在。为证明它与到 $R$-模范畴的遗忘函子可交换，我们使用
Categories，引理 \ref{categories-lemma-adjoint-exact}。事实上，该遗忘函子有
右伴随，即把一个模赋予相应的混沌（或无差别）拓扑模的函子。
\end{proof}








% P01 cursor: frozen topology.tex translated through source line 6020; chapter complete pending structural/build QA.
% Generated mechanically from the frozen sheaves.tex target.
% Do not edit this generated file; edit the bound locale source instead.

% OK, start here.
%

\setcounter{chapter}{5}
\chapter{空间上的层}



\phantomsection
\label{sheaves-section-phantom}


\section{引言}
\label{sheaves-section-introduction}

\noindent
本文将说明拓扑空间上层的基本性质。可参见 \cite{Godement}。

\medskip\noindent
后续文档中关于站点上层的讨论将取代这里的内容。不过，在此简要定义其中一些
概念或许仍有意义。








\section{基本概念}
\label{sheaves-section-sheaves-basic}

\noindent
下面列出拓扑学中的若干基本概念。

\begin{enumerate}
\item 设 $X$ 为拓扑空间。“设 $U = \bigcup_{i \in I} U_i$ 为一个开覆盖”
是指：$I$ 是集合，并且对每个 $i \in I$，给定开子集 $U_i \subset X$，
使得 $U$ 是诸 $U_i$ 之并。允许 $I = \emptyset$；此时不存在 $U_i$，并且
$U = \emptyset$。当 $I \not = \emptyset$ 时，也允许任意一些乃至所有
$U_i$ 为空。
% P01-ERRATA: P01-E0035 -- the frozen reader-facing list contains an editorial placeholder.
\item 等等，等等。
\end{enumerate}









\section{预层}
\label{sheaves-section-presheaves}


\begin{definition}
\label{sheaves-definition-presheaf}
设 $X$ 为拓扑空间。
\begin{enumerate}
\item $X$ 上的一个{\it 集合预层 $\mathcal{F}$}，是这样一条规则：对每个开集
$U \subset X$ 指定一个集合 $\mathcal{F}(U)$，并对每个包含关系
$V \subset U$ 指定一个映射
$\rho^U_V : \mathcal{F}(U) \to \mathcal{F}(V)$，使得
$\rho^U_U = \text{id}_{\mathcal{F}(U)}$，且每当
$W \subset V \subset U$ 时，都有
$\rho^U_W = \rho^V_W \circ \rho ^U_V$。
\item $X$ 上集合预层的一个{\it 态射
$\varphi : \mathcal{F} \to \mathcal{G}$}，是这样一条规则：对每个开集
$U \subset X$ 指定一个与限制映射相容的集合映射
$\varphi : \mathcal{F}(U) \to \mathcal{G}(U)$；也就是说，每当
$V \subset U \subset X$ 为开集时，图表
$$
\xymatrix{
\mathcal{F}(U) \ar[r]^\varphi \ar[d]^{\rho^U_V} &
\mathcal{G}(U) \ar[d]^{\rho^U_V} \\
\mathcal{F}(V) \ar[r]^\varphi & \mathcal{G}(V)
}
$$
交换。
\item $X$ 上集合预层的范畴记作 $\textit{PSh}(X)$。
\end{enumerate}
\end{definition}

\noindent
集合 $\mathcal{F}(U)$ 的元素称为 $\mathcal{F}$ 在 $U$ 上的{\it 截面}。
对每个 $V \subset U$，映射
$\rho^U_V : \mathcal{F}(U) \to \mathcal{F}(V)$ 称为{\it 限制映射}。
若 $s\in \mathcal{F}(U)$，我们采用记号 $s|_V := \rho^U_V(s)$。这一记号
与拓扑学中函数限制的概念相容，因为若 $W \subset V \subset U$，且 $s$
是 $\mathcal{F}$ 在 $U$ 上的截面，则由上述定义中限制映射的性质，有
$s|_W = (s|_V)|_W$。

\medskip\noindent
另一种常用记号是用符号 $\Gamma(U, -)$ 或 $H^0(U, -)$ 表示开集 $U$
上的截面。换言之，下列等式是恒真的：
$$
\Gamma(U, \mathcal{F}) = \mathcal{F}(U) = H^0(U, \mathcal{F}).
$$
本章不使用这一记号，但其他章节会使用。

\begin{definition}
\label{sheaves-definition-constant-presheaf}
设 $X$ 为拓扑空间，$A$ 为集合。{\it 取值为 $A$ 的常值预层}，是把集合
$A$ 指派给每个开集 $U \subset X$，并且所有限制映射均为
$\text{id}_A$ 的预层。
\end{definition}

\section{阿贝尔预层}
\label{sheaves-section-abelian-presheaves}

\noindent
本节简要指出预层范畴的一些特征，它们使我们能够定义阿贝尔群预层。

\begin{example}
\label{sheaves-example-singleton-presheaf}
设 $X$ 为拓扑空间。考虑一条规则 $\mathcal{F}$，它把一个单元素集与 $X$
的每个开子集相联系。由于从任意集合到单元素集存在唯一映射，故存在唯一的
限制映射 $\rho^U_V$。所得结构是 $X$ 上的集合预层。由刚才所述的单元素集
性质，它是 $X$ 上集合预层范畴的终对象。因此，它在唯一同构的意义下也是
唯一的。我们有时用 $*$ 表示该预层。
\end{example}

\begin{lemma}
\label{sheaves-lemma-product-presheaves}
设 $X$ 为拓扑空间。$X$ 上集合预层的范畴有积（见 Categories，定义
\ref{categories-definition-product}）。此外，积
$\mathcal{F} \times \mathcal{G}$ 在开集 $U$ 上的截面集，等于
$\mathcal{F}$ 和 $\mathcal{G}$ 在 $U$ 上的截面集之积。
\end{lemma}

\begin{proof}
具体地，设 $\mathcal{F}$ 和 $\mathcal{G}$ 是拓扑空间 $X$ 上的集合预层。
考虑规则 $U \mapsto \mathcal{F}(U) \times \mathcal{G}(U)$，并将它记作
$\mathcal{F} \times \mathcal{G}$。若 $V \subset U \subset X$ 为开集，
则定义限制映射
$$
(\mathcal{F} \times \mathcal{G})(U)
\longrightarrow
(\mathcal{F} \times \mathcal{G})(V)
$$
为 $(s, t) \mapsto (s|_V, t|_V)$。于是立即可见，
$\mathcal{F} \times \mathcal{G}$ 是一个预层。此外，还有投影映射
$p : \mathcal{F} \times \mathcal{G} \to \mathcal{F}$ 和
$q : \mathcal{F} \times \mathcal{G} \to \mathcal{G}$。留给读者证明：
对任意第三个预层 $\mathcal{H}$，都有
$\Mor(\mathcal{H}, \mathcal{F} \times \mathcal{G})
= \Mor(\mathcal{H}, \mathcal{F}) \times
\Mor(\mathcal{H}, \mathcal{G})$。
\end{proof}

\noindent
回忆一下，若 $(A, + : A \times A \to A, - : A \to A, 0\in A)$ 是阿贝尔群，
则零元映射和取负映射由加法律唯一确定。换言之，说“设 $(A, +)$ 为阿贝尔群”
是有意义的。

\begin{lemma}
\label{sheaves-lemma-abelian-presheaves}
设 $X$ 为拓扑空间，$\mathcal{F}$ 为集合预层。考虑 $\mathcal{F}$ 上的以下
几类结构：
\begin{enumerate}
\item 对每个开集 $U$，在 $\mathcal{F}(U)$ 上给定阿贝尔群结构，使所有限制
映射均为阿贝尔群同态。
\item 给定预层映射 $+ : \mathcal{F} \times \mathcal{F} \to \mathcal{F}$、
预层映射 $- : \mathcal{F} \to \mathcal{F}$ 以及映射
$0 : * \to \mathcal{F}$（见例 \ref{sheaves-example-singleton-presheaf}），并使它们
满足通常阿贝尔群中 $+, -, 0$ 的所有公理。
\item 给定预层映射 $+ : \mathcal{F} \times \mathcal{F} \to \mathcal{F}$、
预层映射 $- : \mathcal{F} \to \mathcal{F}$ 以及映射
$0 : * \to \mathcal{F}$，使得对每个开集 $U \subset X$，四元组
$(\mathcal{F}(U), +, -, 0)$ 是阿贝尔群，
\item 给定预层映射 $+ : \mathcal{F} \times \mathcal{F}
\to \mathcal{F}$，使得对每个开集 $U \subset X$，映射
$+ : \mathcal{F}(U) \times \mathcal{F}(U) \to \mathcal{F}(U)$
定义一个阿贝尔群结构。
\end{enumerate}
上述 (1)--(4) 各类数据的集合之间存在自然双射。
\end{lemma}

\begin{proof}
略。
\end{proof}

\noindent
该引理说明，在预层范畴中给定一个阿贝尔群对象 $\mathcal{F}$，等价于给定
一个集合预层 $\mathcal{F}$，使所有集合 $\mathcal{F}(U)$ 都被赋予阿贝尔群
结构，并使所有限制映射均为群同态。对于大多数代数结构，我们都将用这一
方式处理此类对象的（预）层；也就是说，我们把此类对象的一个（预）层定义为
集合（预）层 $\mathcal{F}$，其所有截面集 $\mathcal{F}(U)$ 均以与限制映射
相容的方式被赋予这一结构。

\begin{definition}
\label{sheaves-definition-abelian-presheaves}
设 $X$ 为拓扑空间。
\begin{enumerate}
\item $X$ 上的一个{\it 阿贝尔群预层}，或称 $X$ 上的一个
{\it 阿贝尔预层}，是集合预层 $\mathcal{F}$，使得对每个开集
$U \subset X$，集合 $\mathcal{F}(U)$ 被赋予阿贝尔群结构，并且所有限制
映射 $\rho^U_V$ 都是阿贝尔群同态；见上面的引理
\ref{sheaves-lemma-abelian-presheaves}。
\item $X$ 上阿贝尔预层的一个{\it 态射}
$\varphi : \mathcal{F} \to \mathcal{G}$，是集合预层的一个态射，并且对每个
开集 $U \subset X$，它诱导阿贝尔群同态
$\mathcal{F}(U) \to \mathcal{G}(U)$。
\item $X$ 上阿贝尔群预层的范畴记作 $\textit{PAb}(X)$。
\end{enumerate}
\end{definition}

\begin{example}
\label{sheaves-example-direct-sum-points}
设 $X$ 为拓扑空间。假设对每个 $x \in X$ 给定阿贝尔群 $M_x$。对开集
$U \subset X$，令
$$
\mathcal{F}(U) = \bigoplus\nolimits_{x \in U} M_x.
$$
我们把该阿贝尔群的典型元素记作 $\sum_{i = 1}^n m_{x_i}$，其中
$x_i \in U$ 且 $m_{x_i} \in M_{x_i}$。（当然，总可以选择一种表示，使
$x_1, \ldots, x_n$ 两两不同。）对开集 $V \subset U \subset X$，定义限制
映射 $\mathcal{F}(U) \to \mathcal{F}(V)$，把元素
$s = \sum_{i = 1}^n m_{x_i}$ 映到元素
$s|_V = \sum_{x_i \in V} m_{x_i}$。留给读者验证，这是一个阿贝尔群预层。
\end{example}



\section{代数结构预层}
\label{sheaves-section-presheaves-structures}

\noindent
我们来澄清代数结构预层的定义。假设 $\mathcal{C}$ 是一个范畴，且
$F : \mathcal{C} \to \textit{Sets}$ 是忠实函子。通常，$F$ 是一个“遗忘”
函子。对对象 $M \in \Ob(\mathcal{C})$，我们常把 $F(M)$ 称为对象 $M$
的{\it 底层集合}。若 $M \to M'$ 是 $\mathcal{C}$ 中的态射，则称
$F(M) \to F(M')$ 为{\it 底层集合映射}。事实上，我们常常不区分一个对象
与其底层集合，对态射也同样如此。因此，若集合映射 $F(M) \to F(M')$
等于 $\mathcal{C}$ 中某个态射 $f : M \to M'$ 的 $F(f)$，则称它为一个
{\it 代数结构态射}。

\medskip\noindent
仿照上面的定义 \ref{sheaves-definition-abelian-presheaves}，一个“$\mathcal{C}$ 中对象
的预层”可以由以下数据定义：
\begin{enumerate}
\item 一个集合预层 $\mathcal{F}$，以及
\item 对每个开集 $U \subset X$，选择一个对象
$A(U) \in \Ob(\mathcal{C})$。
\end{enumerate}
这些数据须满足以下条件（采用上面的说法）：
\begin{enumerate}
\item 对每个开集 $U \subset X$，集合 $\mathcal{F}(U)$ 是 $A(U)$ 的底层
集合，并且
\item 对每个开集 $V \subset U \subset X$，集合映射
$\rho_V^U: \mathcal{F}(U) \to \mathcal{F}(V)$ 是代数结构态射。
\end{enumerate}
换言之，对 $X$ 中每个开集 $V \subset U$，限制映射 $\rho^U_V$ 是范畴
$\mathcal{C}$ 中某个唯一态射 $\alpha^U_V : A(U) \to A(V)$ 的像
$F(\alpha^U_V)$。唯一性由 $F$ 忠实这一条件所保证；它还蕴含：每当
$W \subset V \subset U$ 为 $X$ 中的开集时，有
$\alpha^U_W = \alpha^V_W \circ \alpha^U_V$。系统
$(A(-), \alpha^U_V)$ 就是我们将定义的 $X$ 上取值于 $\mathcal{C}$ 的预层；
比较 Sites，定义 \ref{sites-definition-presheaf}。通过规则
$\mathcal{F}(U) = F(A(U))$ 和 $\rho_V^U = F(\alpha_V^U)$，可以恢复我们的
集合预层 $(\mathcal{F}, \rho_V^U)$。

\begin{definition}
\label{sheaves-definition-presheaf-values-in-category}
设 $X$ 为拓扑空间，$\mathcal{C}$ 为范畴。
\begin{enumerate}
% P01-ERRATA: P01-E0036 -- unlike the set-valued definition, the frozen definition omits the identity restriction axiom.
\item $X$ 上一个{\it 取值于 $\mathcal{C}$ 的预层 $\mathcal{F}$}，由这样一条
规则给出：对每个开集 $U \subset X$ 指定 $\mathcal{C}$ 的一个对象
$\mathcal{F}(U)$，并对每个包含关系 $V \subset U$ 指定 $\mathcal{C}$ 中的
一个态射 $\rho_V^U : \mathcal{F}(U) \to \mathcal{F}(V)$，使得每当
$W \subset V \subset U$ 时，都有
$\rho_W^U = \rho_W^V \circ \rho_V^U$。
\item 取值于 $\mathcal{C}$ 的预层的一个{\it 态射
$\varphi : \mathcal{F} \to \mathcal{G}$}，由与限制态射相容的
$\mathcal{C}$ 中的态射 $\varphi : \mathcal{F}(U) \to \mathcal{G}(U)$ 给出。
\end{enumerate}
\end{definition}

\begin{definition}
\label{sheaves-definition-underlying-presheaf-sets}
设 $X$ 为拓扑空间，$\mathcal{C}$ 为范畴，并设
$F : \mathcal{C} \to \textit{Sets}$ 为忠实函子。设 $\mathcal{F}$ 是 $X$
上取值于 $\mathcal{C}$ 的预层。集合预层 $U \mapsto F(\mathcal{F}(U))$
称为 $\mathcal{F}$ 的{\it 底层集合预层}。
\end{definition}

\noindent
依照定义 \ref{sheaves-definition-presheaf-values-in-category} 前的讨论，通常仍用同一字母
$\mathcal{F}$ 表示其底层集合预层。特别地，“设 $s \in \mathcal{F}(U)$”或
“设 $s$ 是 $\mathcal{F}$ 在 $U$ 上的一个截面”是指
$s \in F(\mathcal{F}(U))$。

\medskip\noindent
这些记号和定义尤其适用于：{\it （不一定阿贝尔的）群、环、固定环上的模、
固定域上的向量空间等的预层}，以及{\it 它们之间的态射}。

\section{模预层}
\label{sheaves-section-presheaves-modules}

\noindent
假设 $\mathcal{O}$ 是 $X$ 上的环预层。我们想定义 $X$ 上
$\mathcal{O}$-模预层的概念。仿照定义 \ref{sheaves-definition-abelian-presheaves}，
很自然地会把它定义为集合预层 $\mathcal{F}$，使得对每个开集
$U \subset X$，集合 $\mathcal{F}(U)$ 被赋予与（$\mathcal{F}$ 和
$\mathcal{O}$ 的）限制映射相容的 $\mathcal{O}(U)$-模结构。不过，习惯上
采用下面这个等价的定义。

\begin{definition}
\label{sheaves-definition-presheaf-modules}
设 $X$ 为拓扑空间，$\mathcal{O}$ 为 $X$ 上的环预层。
\begin{enumerate}
\item 一个{\it $\mathcal{O}$-模预层}，由一个阿贝尔预层 $\mathcal{F}$ 和一个
集合预层映射
$$
\mathcal{O} \times \mathcal{F} \longrightarrow \mathcal{F}
$$
给出，使得对每个开集 $U \subset X$，映射
$\mathcal{O}(U) \times \mathcal{F}(U) \to \mathcal{F}(U)$
在阿贝尔群 $\mathcal{F}(U)$ 上定义一个 $\mathcal{O}(U)$-模结构。
\item $\mathcal{O}$-模预层的一个{\it 态射
$\varphi : \mathcal{F} \to \mathcal{G}$}，是阿贝尔预层的一个态射
$\varphi : \mathcal{F} \to \mathcal{G}$，使图表
$$
\xymatrix{
\mathcal{O} \times \mathcal{F} \ar[r] \ar[d]_{\text{id} \times \varphi} &
\mathcal{F} \ar[d]^{\varphi} \\
\mathcal{O} \times \mathcal{G} \ar[r] &
\mathcal{G}
}
$$
交换。
\item 上述 $\mathcal{O}$-模态射的集合记作
$\Hom_\mathcal{O}(\mathcal{F}, \mathcal{G})$。
\item $\mathcal{O}$-模预层的范畴记作 $\textit{PMod}(\mathcal{O})$。
\end{enumerate}
\end{definition}

\noindent
假设 $\mathcal{O}_1 \to \mathcal{O}_2$ 是 $X$ 上环预层的一个态射。在此
情形下，若 $\mathcal{F}$ 是 $\mathcal{O}_2$-模预层，则可以借助复合映射
$$
\mathcal{O}_1 \times \mathcal{F}
\to
\mathcal{O}_2 \times \mathcal{F}
\to
\mathcal{F}.
$$
把 $\mathcal{F}$ 看作 $\mathcal{O}_1$-模预层。为表明环的限制，有时把它记作
$\mathcal{F}_{\mathcal{O}_1}$。我们称之为 $\mathcal{F}$ 的{\it 限制}。
由此得到限制函子
$$
\textit{PMod}(\mathcal{O}_2)
\longrightarrow
\textit{PMod}(\mathcal{O}_1)
$$

\medskip\noindent
另一方面，给定 $\mathcal{O}_1$-模预层 $\mathcal{G}$，可以按规则
$$
\left(\mathcal{O}_2 \otimes_{p, \mathcal{O}_1} \mathcal{G}\right)(U)
=
\mathcal{O}_2(U) \otimes_{\mathcal{O}_1(U)} \mathcal{G}(U)
$$
构造 $\mathcal{O}_2$-模预层
$\mathcal{O}_2 \otimes_{p, \mathcal{O}_1} \mathcal{G}$。下标 $p$ 表示
“预层”而不是“点”。该预层称为{\it 张量积预层}。由此得到{\it 变环}函子
$$
\textit{PMod}(\mathcal{O}_1)
\longrightarrow
\textit{PMod}(\mathcal{O}_2)
$$

\begin{lemma}
\label{sheaves-lemma-adjointness-tensor-restrict-presheaves}
在上述 $X$、$\mathcal{O}_1$、$\mathcal{O}_2$、$\mathcal{F}$ 和
$\mathcal{G}$ 的条件下，存在典范双射
$$
\Hom_{\mathcal{O}_1}(\mathcal{G}, \mathcal{F}_{\mathcal{O}_1})
=
\Hom_{\mathcal{O}_2}(
\mathcal{O}_2 \otimes_{p, \mathcal{O}_1} \mathcal{G},
\mathcal{F}
)
$$
换言之，限制函子与变环函子互为伴随。
\end{lemma}

\begin{proof}
这由以下事实得出：对环同态 $A \to B$，限制函子与变环函子互为伴随。
\end{proof}





\section{层}
\label{sheaves-section-sheaves}

\noindent
本节说明层条件。

\begin{definition}
\label{sheaves-definition-sheaf}
设 $X$ 为拓扑空间。
\begin{enumerate}
\item $X$ 上的一个{\it 集合层 $\mathcal{F}$}，是满足以下附加性质的集合
预层：给定任意开覆盖 $U = \bigcup_{i \in I} U_i$，以及任意一族截面
$s_i \in \mathcal{F}(U_i)$，$i \in I$，使得 $\forall i, j\in I$ 都有
$$
s_i|_{U_i \cap U_j} = s_j|_{U_i \cap U_j}
$$
则存在唯一截面 $s \in \mathcal{F}(U)$，使得对所有 $i \in I$ 都有
$s_i = s|_{U_i}$。
\item 一个{\it 集合层态射}就是一个集合预层态射。
\item $X$ 上集合层的范畴记作 $\Sh(X)$。
\end{enumerate}
\end{definition}

\begin{remark}
\label{sheaves-remark-confusion}
关于是否有必要说明层在空集 $\emptyset \subset X$ 上的截面集，总会有一点
混淆。确有必要；如果正确理解定义，我们其实已经作了说明。事实上，空集由
空开覆盖覆盖，因此上述定义中的“截面族 $s_i$”实际构成空积的一个元素，而
空积是层所取值范畴的终对象。换言之，正确理解定义便会自动推出
$\mathcal{F}(\emptyset) = \textit{终对象}$；对集合层而言，这就是一个
单元素集。如果不喜欢这一论证，也可以直接要求
$\mathcal{F}(\emptyset) = \{*\}$。

\medskip\noindent
特别地，该条件于是保证：若 $U, V \subset X$ 为开集且{\it 不交}，则
$$
\mathcal{F}(U \cup V) = \mathcal{F}(U) \times \mathcal{F}(V).
$$
（因为终对象上的纤维积就是积。）
\end{remark}

\begin{example}
\label{sheaves-example-basic-continuous-maps}
设 $X$、$Y$ 为拓扑空间。考虑规则 $\mathcal{F}$，它把集合
$$
\mathcal{F}(U) = \{ f : U \to Y \mid f \text{ 连续}\}
$$
以及显然的限制映射与开集 $U \subset X$ 相联系。我们断言
$\mathcal{F}$ 是一个层。为此，假设 $U = \bigcup_{i\in I} U_i$ 是开覆盖，
并且 $f_i \in \mathcal{F}(U_i)$，$i\in I$，满足对所有 $i, j \in I$ 都有
$f_i |_{U_i \cap U_j} = f_j|_{U_i \cap U_j}$。在此情形下，定义
$f : U \to Y$ 如下：对任意满足 $u \in U_i$ 的 $i \in I$，令 $f(u)$ 等于
$f_i(u)$ 的值。由假设，这是良定义的。此外，映射 $f : U \to Y$ 在
$U_i$ 上的限制与连续映射 $f_i$ 相同。因此 $f$ 显然连续！
\end{example}

\noindent
可以利用该例的结果定义常值层。具体地，设 $A$ 为集合，并在 $A$ 上赋予
离散拓扑。设 $U \subset X$ 为开子集。于是
$$
\{ f : U \to A \mid f\text{ 连续}\}
=
\{ f : U \to A \mid f\text{ 局部常值}\}.
$$
因此，把所有到 $A$ 的局部常值映射指派给一个开集的规则是一个层。

\begin{definition}
\label{sheaves-definition-constant-sheaf}
设 $X$ 为拓扑空间，$A$ 为集合。取值为 $A$ 的{\it 常值层}记作
$\underline{A}$ 或 $\underline{A}_X$；它把所有局部常值映射 $U \to A$
组成的集合指派给开集 $U \subset X$，其限制映射由函数的限制给出。
\end{definition}

\begin{example}
\label{sheaves-example-sheaf-product-pointwise}
设 $X$ 为拓扑空间，并设 $(A_x)_{x \in X}$ 是由点 $x \in X$ 标号的一族
集合 $A_x$。我们将由这些数据构造一个集合层 $\Pi$。对开集
$U \subset X$，令
$$
\Pi(U) = \prod\nolimits_{x \in U} A_x.
$$
对开集 $V \subset U \subset X$，按下列规则定义限制映射：元素
$s = (a_x)_{x\in U} \in \Pi(U)$ 限制为
$s|_V = (a_x)_{x \in V}$。这显然定义了一个集合预层。我们断言它是一个层。
事实上，设 $U = \bigcup U_i$ 为开覆盖。假设 $s_i \in \Pi(U_i)$，并且
$s_i$ 与 $s_j$ 在 $U_i \cap U_j$ 上相同。写
$s_i = (a_{i, x})_{x\in U_i}$。相容性条件蕴含：每当
$x \in U_i \cap U_j$ 时，集合 $A_x$ 中有 $a_{i, x} = a_{j, x}$。因此，
存在唯一元素 $s = (a_x)_{x\in U}$ 属于
$\Pi(U) = \prod_{x\in U} A_x$，使得只要对某个 $i$ 有 $x \in U_i$，就有
$a_x = a_{i, x}$。当然，对所有 $i$，该元素 $s$ 都满足
$s|_{U_i} = s_i$。
\end{example}

\begin{example}
\label{sheaves-example-direct-sum-points-not-sheaf}
设 $X$ 为拓扑空间。假设对每个 $x\in X$ 给定阿贝尔群 $M_x$。考虑例
\ref{sheaves-example-direct-sum-points} 中定义的预层
$\mathcal{F} : U \mapsto \bigoplus_{x \in U} M_x$。一般而言，它不是一个层。
例如，若 $X$ 是赋予离散拓扑的无限集合，则层条件将蕴含
$\mathcal{F}(X) = \prod_{x\in X} \mathcal{F}(\{x\})$，但按照定义，
$\mathcal{F}(X) = \bigoplus_{x \in X} M_x
= \bigoplus_{x \in X} \mathcal{F}(\{x\})$。一般而言，无限直和不同于
无限直积。

\medskip\noindent
不过，若 $X$ 是一个拓扑空间，且 $X$ 的每个开集都拟紧，则
$\mathcal{F}$ {\it 的确}是一个层。留给读者作为练习。
\end{example}



\section{阿贝尔层}
\label{sheaves-section-abelian-sheaves}

\begin{definition}
\label{sheaves-definition-abelian-sheaf}
设 $X$ 为拓扑空间。
\begin{enumerate}
\item $X$ 上的一个{\it 阿贝尔层}，或称 $X$ 上的一个{\it 阿贝尔群层}，
是 $X$ 上的阿贝尔预层，并且其底层集合预层是一个层。
\item 阿贝尔群层的范畴记作 $\textit{Ab}(X)$。
\end{enumerate}
\end{definition}

\noindent
设 $X$ 为拓扑空间。对阿贝尔预层 $\mathcal{F}$，关于开覆盖
$U = \bigcup U_i$ 的层条件常表述为：阿贝尔群复形
$$
0 \to \mathcal{F}(U)
\to \prod\nolimits_i \mathcal{F}(U_i)
\to \prod\nolimits_{(i_0, i_1)} \mathcal{F}(U_{i_0} \cap U_{i_1})
$$
正合。第一个映射是通常的映射，而第二个映射把元素 $(s_i)_{i \in I}$
映到元素
$$
(
s_{i_0}|_{U_{i_0} \cap U_{i_1}} -
s_{i_1}|_{U_{i_0} \cap U_{i_1}}
)_{(i_0, i_1)}
\in \prod\nolimits_{(i_0, i_1)} \mathcal{F}(U_{i_0} \cap U_{i_1})
$$

\section{代数结构层}
\label{sheaves-section-sheaves-structures}

\noindent
我们来澄清某些类型结构之层的定义。首先重新表述层条件。设
$\mathcal{F}$ 是拓扑空间 $X$ 上的集合预层。层条件可以表述如下。设
$U = \bigcup_{i\in I} U_i$ 为开覆盖。考虑图表
$$
\xymatrix{
\mathcal{F}(U) \ar[r]
&
\prod\nolimits_{i\in I}
\mathcal{F}(U_i)
\ar@<1ex>[r] \ar@<-1ex>[r]
&
\prod\nolimits_{(i_0, i_1) \in I \times I}
\mathcal{F}(U_{i_0} \cap U_{i_1})
}
$$
这里，左侧映射由规则 $s \mapsto \prod_{i \in I} s|_{U_i}$ 定义。右侧的
两个映射分别是
$$
\prod\nolimits_i s_i
\mapsto
\prod\nolimits_{(i_0, i_1)} s_{i_0}|_{U_{i_0} \cap U_{i_1}}
\text{ 分别 }
\prod\nolimits_i s_i
\mapsto
\prod\nolimits_{(i_0, i_1)} s_{i_1}|_{U_{i_0} \cap U_{i_1}}.
$$
层条件恰好说明，左侧箭头是右侧两个箭头的等化子。只要一个范畴有积，这一
表述就立即推广到取值于该范畴的预层。

\begin{definition}
\label{sheaves-definition-sheaf-values-in-category}
设 $X$ 为拓扑空间，$\mathcal{C}$ 为有积的范畴。若 $X$ 上取值于
$\mathcal{C}$ 的预层 $\mathcal{F}$ 对每个开覆盖都有图表
$$
\xymatrix{
\mathcal{F}(U) \ar[r]
&
\prod\nolimits_{i\in I}
\mathcal{F}(U_i)
\ar@<1ex>[r] \ar@<-1ex>[r]
&
\prod\nolimits_{(i_0, i_1) \in I \times I}
\mathcal{F}(U_{i_0} \cap U_{i_1})
}
$$
在范畴 $\mathcal{C}$ 中是等化子图表，则称它为一个{\it 层}。
\end{definition}

\noindent
假设 $\mathcal{C}$ 是范畴，且 $F : \mathcal{C} \to \textit{Sets}$ 是忠实
函子。一个值得牢记的例子是：$\mathcal{C}$ 是阿贝尔群范畴，而 $F$ 是遗忘
函子。考虑 $X$ 上取值于 $\mathcal{C}$ 的预层 $\mathcal{F}$。我们希望用其
底层集合预层（定义 \ref{sheaves-definition-underlying-presheaf-sets}）重新表述上述
条件。注意，底层集合预层是集合层，当且仅当施加遗忘函子 $F$ 后，所有集合
图表
$$
\xymatrix{
F(\mathcal{F}(U)) \ar[r]
&
\prod\nolimits_{i\in I}
F(\mathcal{F}(U_i))
\ar@<1ex>[r] \ar@<-1ex>[r]
&
\prod\nolimits_{(i_0, i_1) \in I \times I}
F(\mathcal{F}(U_{i_0} \cap U_{i_1}))
}
$$
都是等化子图表！因此，我们希望 $\mathcal{C}$ 有积和等化子，并希望 $F$
与它们可交换。这等价于 $\mathcal{C}$ 有极限且 $F$ 与极限可交换；见
Categories，引理 \ref{categories-lemma-limits-products-equalizers}。但这仍不
充分（见例 \ref{sheaves-example-sheaves-topological-spaces}）；还需要 $F$
{\it 反映同构}。这一性质是指：给定 $\mathcal{C}$ 中的态射
$f : A \to A'$，若（且仅若）$F(f)$ 是双射，则 $f$ 是同构。

\begin{lemma}
\label{sheaves-lemma-sheaves-structure}
假设范畴 $\mathcal{C}$ 和函子 $F : \mathcal{C} \to \textit{Sets}$ 满足：
\begin{enumerate}
\item $F$ 忠实，
\item $\mathcal{C}$ 有极限，并且 $F$ 与极限可交换，以及
\item 函子 $F$ 反映同构。
\end{enumerate}
设 $X$ 为拓扑空间，并设 $\mathcal{F}$ 是取值于 $\mathcal{C}$ 的预层。
则 $\mathcal{F}$ 是层，当且仅当其底层集合预层是层。
\end{lemma}

\begin{proof}
假设 $\mathcal{F}$ 是层。则 $\mathcal{F}(U)$ 是上述图表的等化子；由假设，
$F(\mathcal{F}(U))$ 是相应集合图表的等化子。因此 $F(\mathcal{F})$ 是集合层。

\medskip\noindent
反之，假设 $F(\mathcal{F})$ 是层。令 $E \in \Ob(\mathcal{C})$ 为定义
\ref{sheaves-definition-sheaf-values-in-category} 中两个平行箭头的等化子。仅由
$\mathcal{F}$ 是预层，就得到典范态射 $\mathcal{F}(U) \to E$。由假设，
诱导映射 $F(\mathcal{F}(U)) \to F(E)$ 是同构，因为 $F(E)$ 是相应集合图表
的等化子。因此，由引理的条件 (3)，可知 $\mathcal{F}(U) \to E$ 是同构。
\end{proof}

\noindent
该引理特别适用于{\it 群、环、固定环上的代数、固定环上的模、固定域上的
向量空间等的层}。换言之，它们是群、环、固定环上的模、固定域上的向量
空间等的预层，并且其底层集合预层是层。

\begin{example}
\label{sheaves-example-C0-sheaf-rings}
设 $X$ 为拓扑空间。对每个开集 $U \subset X$，考虑
$\mathbf{R}$-代数
$\mathcal{C}^{0}(U) = \{ f : U \to \mathbf{R} \mid f\text{ 连续}\}$。
存在显然的限制映射，使其成为 $X$ 上的 $\mathbf{R}$-代数预层。由例
\ref{sheaves-example-basic-continuous-maps}，它是集合层。因此，由引理
\ref{sheaves-lemma-sheaves-structure}，它是 $X$ 上的 $\mathbf{R}$-代数层。
\end{example}

\begin{example}
\label{sheaves-example-sheaves-topological-spaces}
考虑拓扑空间范畴 $\textit{Top}$。存在自然的忠实函子
$\textit{Top} \to \textit{Sets}$，它与积和等化子可交换，但不反映同构。
事实上，结果表明引理 \ref{sheaves-lemma-sheaves-structure} 的类似命题是错误的。
具体地，设 $X = \mathbf{N}$ 并赋予离散拓扑。设 $A_i$，
$i \in \mathbf{N}$，为离散拓扑空间。对任意子集
$U \subset \mathbf{N}$，定义 $\mathcal{F}(U) = \prod_{i\in U} A_i$，
并赋予离散拓扑。于是这是一个拓扑空间预层，其底层集合预层是层，见例
\ref{sheaves-example-sheaf-product-pointwise}。然而，如果每个 $A_i$ 至少有两个元素，
那么按照定义 \ref{sheaves-definition-sheaf-values-in-category}，它并不是拓扑空间层。
读者可以验证，若在每个 $\mathcal{F}(U) = \prod_{i\in U} A_i$ 上赋予
{\it 积拓扑}，则确实得到 $X$ 上的拓扑空间层。
\end{example}


\section{模层}
\label{sheaves-section-sheaves-modules}

\begin{definition}
\label{sheaves-definition-sheaf-modules}
设 $X$ 为拓扑空间，$\mathcal{O}$ 为 $X$ 上的环层。
\begin{enumerate}
\item 一个{\it $\mathcal{O}$-模层}，是一个 $\mathcal{O}$-模预层
$\mathcal{F}$（见定义 \ref{sheaves-definition-presheaf-modules}），并且其底层阿贝尔群
预层 $\mathcal{F}$ 是层。
\item 一个{\it $\mathcal{O}$-模层态射}，是一个 $\mathcal{O}$-模预层态射。
% P01-ERRATA: P01-E0037 -- frozen English has singular "set of morphism".
\item 给定 $\mathcal{O}$-模层 $\mathcal{F}$ 和 $\mathcal{G}$，用
$\Hom_\mathcal{O}(\mathcal{F}, \mathcal{G})$ 表示 $\mathcal{O}$-模层态射
的集合。
\item $\mathcal{O}$-模层的范畴记作 $\textit{Mod}(\mathcal{O})$。
\end{enumerate}
\end{definition}

\noindent
即使 $\mathcal{O}$ 只是一个环预层，这一定义也有某种意义；不过，我们不知
它有什么有用的例子，并且当 $\mathcal{O}$ 不是环层时，将避免使用
“$\mathcal{O}$-模层”这一术语。







\section{茎}
\label{sheaves-section-stalks}

\noindent
设 $X$ 为拓扑空间，$x \in X$ 为一点，并设 $\mathcal{F}$ 是 $X$ 上的集合
预层。$\mathcal{F}$ 在 $x$ 处的{\it 茎}是集合
$$
\mathcal{F}_x
=
\colim_{x\in U} \mathcal{F}(U)
$$
其中余极限遍历 $X$ 中 $x$ 的开邻域 $U$ 所成集合。开邻域集合按（反向）
包含关系偏序：规定 $U \geq U' \Leftrightarrow U \subset U'$。系统中的转移
映射由 $\mathcal{F}$ 的限制映射给出。关于系统上（余）极限的记号和术语，
见 Categories，第 \ref{categories-section-posets-limits} 节。注意，这一余极限
是有向余极限。因此，很容易描述 $\mathcal{F}_x$。具体地，
$$
\mathcal{F}_x
=
\{
(U, s)
\mid
x\in U, s\in \mathcal{F}(U)
\}/\sim
$$
其中等价关系定义为：$(U, s) \sim (U', s')$，当且仅当存在开集
$U'' \subset U \cap U'$，满足 $x \in U''$ 且
$s|_{U''} = s'|_{U''}$。给定一对 $(U, s)$，有时用 $s_x$ 表示
$\mathcal{F}_x$ 中与 $(U, s)$ 的等价类相应的元素。有时也说“$s$ 在
$\mathcal{F}_x$ 中的像”来指 $s_x$。例如，给定两对 $(U, s)$ 和
$(U', s')$，有时说“$s$ 与 $s'$ 在 $\mathcal{F}_x$ 中相等”，表示
$s_x = s'_x$。其他作者使用“$s$ 在 $x$ 处的芽”这一术语。

\medskip\noindent
该定义的一个显然推论是：对任意开集 $U \subset X$，存在典范映射
$$
\mathcal{F}(U)
\longrightarrow
\prod\nolimits_{x \in U} \mathcal{F}_x
$$
由 $s \mapsto \prod_{x \in U} (U, s)$ 定义。请想一想！

\begin{lemma}
\label{sheaves-lemma-sheaf-subset-stalks}
设 $\mathcal{F}$ 是拓扑空间 $X$ 上的集合层。对每个开集 $U \subset X$，
映射
$$
\mathcal{F}(U)
\longrightarrow
\prod\nolimits_{x \in U} \mathcal{F}_x
$$
是单射。
\end{lemma}

\begin{proof}
假设 $s, s' \in \mathcal{F}(U)$ 对所有 $x \in U$ 都映到每个茎
$\mathcal{F}_x$ 中的同一元素。这意味着，对每个 $x \in U$，存在开集
$V^x \subset U$，$x \in V^x$，使得 $s|_{V^x} = s'|_{V^x}$。但
$U = \bigcup_{x \in U} V^x$ 是一个开覆盖。因此，由层条件中的唯一性，
可知 $s = s'$。
\end{proof}

\begin{definition}
\label{sheaves-definition-separated}
设 $X$ 为拓扑空间。若对每个开集 $U \subset X$，映射
$\mathcal{F}(U) \to \prod_{x \in U} \mathcal{F}_x$ 都是单射，则称 $X$
上的集合预层 $\mathcal{F}$ 是{\it 分离的}。
\end{definition}

\noindent
另一个观察是，茎 $\mathcal{F}_x$ 的构造关于预层 $\mathcal{F}$ 是函子性的。
换言之，它给出函子
$$
\textit{PSh}(X) \longrightarrow \textit{Sets},
\ \mathcal{F} \longmapsto \mathcal{F}_x.
$$
该函子称为{\it 茎函子}。具体地，若
$\varphi : \mathcal{F} \to \mathcal{G}$ 是预层态射，则按规则
$(U, s) \mapsto (U, \varphi(s))$ 定义
$\varphi_x : \mathcal{F}_x \to \mathcal{G}_x$。为验证这确实可行，需要检查：
若 $(U, s) = (U', s')$ 在 $\mathcal{F}_x$ 中成立，则
$(U, \varphi(s)) = (U', \varphi(s'))$ 在 $\mathcal{G}_x$ 中也成立。
这由 $\varphi$ 与限制映射相容显然可得。

\begin{example}
\label{sheaves-example-stalk-constant-presheaf}
设 $X$ 为拓扑空间，$A$ 为集合。暂用 $A_p$ 表示取值为 $A$ 的常值预层
（$p$ 表示预层，而不是点）。存在到取值为 $A$ 的常值层的典范预层映射
$A_p \to \underline{A}$。对每个点，都有典范双射
$A = (A_p)_x = \underline{A}_x$；其中第二个映射由函子性从映射
$A_p \to \underline{A}$ 诱导。
\end{example}



\begin{example}
\label{sheaves-example-germs-functions}
假设 $X = \mathbf{R}^n$ 并赋予 Euclid 拓扑。考虑 $X$ 上的
$\mathcal{C}^\infty$ 函数预层，记作
$\mathcal{C}^\infty_{\mathbf{R}^n}$。换言之，
$\mathcal{C}^\infty_{\mathbf{R}^n}(U)$ 是 $\mathcal{C}^\infty$ 函数
$f : U \to \mathbf{R}$ 所成集合。与例
\ref{sheaves-example-basic-continuous-maps} 一样，容易证明这是一个层。事实上，它是
$\mathbf{R}$-向量空间层。

\medskip\noindent
接着，设 $x \in X = \mathbf{R}^n$ 为一点。应如何理解茎
$\mathcal{C}^\infty_{\mathbf{R}^n, x}$ 中的元素？这样的元素由定义域包含
$x$ 的一个 $\mathcal{C}^\infty$ 函数 $f$ 给出。两个这样的函数 $f$、$g$
确定同一个茎元素，当且仅当它们在 $x$ 的某个邻域上一致。换言之，
$\mathcal{C}^\infty_{\mathbf{R}^n, x}$ 的一个元素，就是有时所谓的
{\it $\mathcal{C}^\infty$ 函数在 $x$ 处的芽}。
\end{example}

\begin{example}
\label{sheaves-example-sheaf-product-pointwise-stalk}
设 $X$ 为拓扑空间，并对每个 $x \in X$ 设 $A_x$ 为一个集合。考虑例
\ref{sheaves-example-sheaf-product-pointwise} 中的层
$\mathcal{F} : U \mapsto \prod_{x\in U} A_x$。这里只想指出，
$\mathcal{F}$ 在 $x$ 处的茎 $\mathcal{F}_x$ 一般{\it 不}等于集合 $A_x$。
当然，存在映射 $\mathcal{F}_x \to A_x$，但一般能说的仅此而已。例如，假设
$x = \lim x_n$，且对所有 $n \not = m$ 都有 $x_n \not = x_m$，并假设对所有
$y \in X$ 都有 $A_y = \{0, 1\}$。则 $\mathcal{F}_x$ 映满由 $0$ 和 $1$
组成的序列尾所成的（无限）集合。事实上，$x$ 的每个开邻域都包含几乎所有
$x_n$。另一方面，若 $x$ 的每个邻域都包含一点 $y$，使得
$A_y = \emptyset$，则 $\mathcal{F}_x = \emptyset$。
\end{example}



\section{阿贝尔预层的茎}
\label{sheaves-section-stalks-abelian-presheaves}

\noindent
我们先处理阿贝尔群的情形，以此作为一般情形的模型。

\begin{lemma}
\label{sheaves-lemma-stalk-abelian-presheaf}
设 $X$ 为拓扑空间，$\mathcal{F}$ 为 $X$ 上的阿贝尔群预层。在
$\mathcal{F}_x$ 上存在唯一的阿贝尔群结构，使得对每个开集
$U \subset X$，$x\in U$，映射
$\mathcal{F}(U) \to \mathcal{F}_x$ 是群同态。此外，
$$
\mathcal{F}_x
=
\colim_{x\in U} \mathcal{F}(U)
$$
在阿贝尔群范畴中成立。
\end{lemma}

\begin{proof}
把两个元素 $(U, s)$ 和 $(V, t)$ 之和定义为
$(U \cap V, s|_{U\cap V} + t|_{U \cap V})$。其余性质容易验证。
\end{proof}

\noindent
上述证明的关键在于，开邻域的偏序集是有向集（Categories，定义
\ref{categories-definition-directed-set}）。事实上，两个阿贝尔群 $A, B$
的余积是直和 $A \oplus B$，而集合范畴中的余积是不交并
$A \amalg B$；这表明一般而言，阿贝尔群范畴中的余极限不同于集合范畴中的
余极限。


\section{代数结构预层的茎}
\label{sheaves-section-stalks-presheaves-structures}

\noindent
对任何一种使得有向余极限与遗忘函子可交换的代数结构，引理
\ref{sheaves-lemma-stalk-abelian-presheaf} 的证明都同样适用。

\begin{lemma}
\label{sheaves-lemma-stalk-presheaf-values-in-category}
设 $\mathcal{C}$ 为范畴，$F : \mathcal{C} \to \textit{Sets}$ 为函子。
假设：
\begin{enumerate}
\item $F$ 忠实，并且
\item $\mathcal{C}$ 中存在有向余极限，且 $F$ 与它们可交换。
\end{enumerate}
设 $X$ 为拓扑空间，$x \in X$，并设 $\mathcal{F}$ 是取值于
$\mathcal{C}$ 的预层。则
$$
\mathcal{F}_x = \colim_{x\in U} \mathcal{F}(U)
$$
在 $\mathcal{C}$ 中存在。它的底层集合等于 $\mathcal{F}$ 的底层集合预层的
茎。此外，构造 $\mathcal{F} \mapsto \mathcal{F}_x$ 是从取值于
$\mathcal{C}$ 的预层范畴到 $\mathcal{C}$ 的函子。
\end{lemma}

\begin{proof}
略。
\end{proof}

\noindent
按照定义，所有态射 $\mathcal{F}(U) \to \mathcal{F}_x$ 都是范畴
$\mathcal{C}$ 中的态射；施加遗忘函子 $F$ 后，它们成为底层集合层的相应
% P01-ERRATA: P01-E0038 -- the frozen text says underlying sheaf, but only an underlying presheaf is assumed.
映射。与往常一样，我们不区分 $\mathcal{C}$ 中的态射与其底层集合映射；
由于 $F$ 忠实，这是允许的。

\medskip\noindent
该引理特别适用于：{\it （不一定阿贝尔的）群、环、固定环上的模、固定域上的
向量空间的预层}。

\section{模预层的茎}
\label{sheaves-section-stalk-presheaves-modules}

\begin{lemma}
\label{sheaves-lemma-stalk-module}
设 $X$ 为拓扑空间，$\mathcal{O}$ 为 $X$ 上的环预层，$\mathcal{F}$ 为
$\mathcal{O}$-模预层，并设 $x \in X$。由乘法映射
$\mathcal{O} \times \mathcal{F} \to \mathcal{F}$ 得到的典范映射
$\mathcal{O}_x \times \mathcal{F}_x \to \mathcal{F}_x$，在阿贝尔群
$\mathcal{F}_x$ 上定义一个 $\mathcal{O}_x$-模结构。
\end{lemma}

\begin{proof}
略。
\end{proof}

\begin{lemma}
\label{sheaves-lemma-stalk-tensor-presheaf-modules}
设 $X$ 为拓扑空间，$\mathcal{O} \to \mathcal{O}'$ 是 $X$ 上环预层的一个
态射，$\mathcal{F}$ 是 $\mathcal{O}$-模预层，并设 $x \in X$。作为
$\mathcal{O}'_x$-模，有
$$
\mathcal{F}_x \otimes_{\mathcal{O}_x} \mathcal{O}'_x
=
(\mathcal{F} \otimes_{p, \mathcal{O}} \mathcal{O}')_x
$$
\end{lemma}

\begin{proof}
略。
\end{proof}


\section{代数结构}
\label{sheaves-section-algebraic-structures}

\noindent
本节对上述各节中遇到的概念作轻度形式化。

\begin{definition}
\label{sheaves-definition-algebraic-structure}
一种{\it 代数结构类型}由一个范畴 $\mathcal{C}$ 和一个函子
$F : \mathcal{C} \to \textit{Sets}$ 给出，并满足以下性质：
\begin{enumerate}
\item $F$ 忠实，
\item $\mathcal{C}$ 有极限，且 $F$ 与极限可交换，
\item $\mathcal{C}$ 有滤过余极限，且 $F$ 与它们可交换，以及
\item $F$ 反映同构。
\end{enumerate}
\end{definition}

\noindent
我们作此定义，是为了明确指出下文若干论证中将用到的性质。不过，我们实际
上不会深入研究这一概念，因为按照约定，除 Categories，评注
\ref{categories-remark-big-categories} 所列者外，我们不得研究“大”范畴。
其中，以下范畴具有所需性质。

\begin{lemma}
\label{sheaves-lemma-list-algebraic-structures}
下列范畴在配备显然的遗忘函子后，定义代数结构类型：
\begin{enumerate}
\item 带基点集合的范畴。
\item 阿贝尔群范畴。
\item 群范畴。
\item 幺半群范畴。
\item 环范畴。
\item 固定环 $R$ 上的 $R$-模范畴。
\item 固定域上的 Lie 代数范畴。
\end{enumerate}
\end{lemma}

\begin{proof}
略。
\end{proof}

\noindent
从现在起，我们将借助底层集合（预）层来理解代数结构的（预）层及其茎。
由引理 \ref{sheaves-lemma-sheaves-structure} 和
\ref{sheaves-lemma-stalk-presheaf-values-in-category}，这是允许的。

\medskip\noindent
本节余下部分指出一些以后会用到的代数结构结果。

\begin{lemma}
\label{sheaves-lemma-properties-algebraic-structures}
设 $(\mathcal{C}, F)$ 为一种代数结构类型。
\begin{enumerate}
\item $\mathcal{C}$ 有终对象 $0$，且 $F(0) = \{ * \}$。
\item $\mathcal{C}$ 有积，且 $F(\prod A_i) = \prod F(A_i)$。
\item $\mathcal{C}$ 有纤维积，且
$F(A \times_B C) = F(A)\times_{F(B)}F(C)$。
\item $\mathcal{C}$ 有等化子；若 $E \to A$ 是
$a, b : A \to B$ 的等化子，则 $F(E) \to F(A)$ 是
$F(a), F(b) : F(A) \to F(B)$ 的等化子。
\item $A \to B$ 是单态射，当且仅当 $F(A) \to F(B)$ 是单射。
\item 若 $F(a) : F(A) \to F(B)$ 是满射，则 $a$ 是满态射。
\item 给定 $A_1 \to A_2 \to A_3 \to \ldots$，则
$\colim A_i$ 存在，且 $F(\colim A_i) = \colim F(A_i)$；对任意滤过余极限
更一般地也成立。
\end{enumerate}
\end{lemma}

\begin{proof}
略。唯一值得说明的是 (5)：$A \to B$ 是单态射，当且仅当
$A \to A \times_B A$ 是同构；再利用 $F$ 反映同构即可。
\end{proof}

\begin{lemma}
\label{sheaves-lemma-image-contained-in}
设 $(\mathcal{C}, F)$ 为一种代数结构类型。假设
$A, B, C \in \Ob(\mathcal{C})$，并设 $f : A \to B$ 和
$g : C \to B$ 是 $\mathcal{C}$ 的态射。若 $F(g)$ 是单射，且
$\Im(F(f)) \subset \Im(F(g))$，则 $f$ 可作如下分解：存在态射
$t : A \to C$，使得
$f = g \circ t$。
\end{lemma}

\begin{proof}
考虑 $A \times_B C$。假设蕴含
$F(A \times_B C) = F(A) \times_{F(B)} F(C) = F(A)$。由于 $F$ 反映同构，
所以 $A = A \times_B C$。结论随即得出。
\end{proof}

\begin{example}
\label{sheaves-example-application-lemma-image-contained-in}
该引理将经常应用于以下情形。假设在 $\mathcal{C}$ 中有图表
$$
\xymatrix{
A \ar[r] & B \ar[d] \\
C \ar[r] & D
}
$$
假设 $C \to D$ 在底层集合上为单射，并且复合映射 $A \to B \to D$ 在
底层集合上的像包含于 $C \to D$ 的像中。于是，在 $\mathcal{C}$ 中得到
交换图表
$$
\xymatrix{
A \ar[r] \ar[d] & B \ar[d] \\
C \ar[r] & D
}
$$
\end{example}

\begin{example}
\label{sheaves-example-sheaf-product-pointwise-algebraic-structure}
% P01-ERRATA: P01-E0039 -- the frozen sentence calls the functor F itself a type, whereas Definition 14.1 makes the pair (C,F) the type.
设 $F : \mathcal{C} \to \textit{Sets}$ 为一种代数结构。设 $X$ 为拓扑空间。
假设对每个 $x \in X$ 给定对象 $A_x \in \Ob(\mathcal{C})$。考虑 $X$ 上
取值于 $\mathcal{C}$ 的预层 $\Pi$，它由规则
$\Pi(U) = \prod_{x \in U} A_x$ 定义（限制映射是显然的）。注意，相应的
集合预层 $U \mapsto F(\Pi(U)) = \prod_{x \in U} F(A_x)$ 由例
\ref{sheaves-example-sheaf-product-pointwise} 是一个层。因此，$\Pi$ 是类型
$(\mathcal{C} , F)$ 的代数结构层。这给出了许多阿贝尔群层、群层、环层等
例子。
\end{example}












\section{正合性与点}
\label{sheaves-section-exactness-points}

\noindent
在任意范畴中，都有满态射、单态射、同构等概念。

\begin{lemma}
\label{sheaves-lemma-points-exactness}
设 $X$ 为拓扑空间，并设 $\varphi : \mathcal{F} \to \mathcal{G}$
为 $X$ 上集合层的态射。
\begin{enumerate}
\item 映射 $\varphi$ 是层范畴中的单态射，当且仅当对所有 $x \in X$，
映射 $\varphi_x : \mathcal{F}_x \to \mathcal{G}_x$ 都是单射。
\item 映射 $\varphi$ 是层范畴中的满态射，当且仅当对所有 $x \in X$，
映射 $\varphi_x : \mathcal{F}_x \to \mathcal{G}_x$ 都是满射。
\item 映射 $\varphi$ 是层范畴中的同构，当且仅当对所有 $x \in X$，
映射 $\varphi_x : \mathcal{F}_x \to \mathcal{G}_x$ 都是双射。
\end{enumerate}
\end{lemma}

\begin{proof}
略。
\end{proof}

\noindent
由此可见，在集合层范畴中，满态射与单态射可描述如下。

\begin{definition}
\label{sheaves-definition-injective-surjective}
设 $X$ 为拓扑空间。
\begin{enumerate}
\item 若对每个开集 $U \subset X$ 都有
$\mathcal{F}(U) \subset \mathcal{G}(U)$，且 $\mathcal{G}$ 的限制映射
诱导出 $\mathcal{F}$ 的限制映射，则称预层 $\mathcal{F}$ 是预层
$\mathcal{G}$ 的一个{\it 子预层}。若 $\mathcal{F}$ 和 $\mathcal{G}$
都是层，则称 $\mathcal{F}$ 为 $\mathcal{G}$ 的一个{\it 子层}。
有时用记号 $\mathcal{F} \subset \mathcal{G}$ 表示这一关系。
\item $X$ 上集合预层的态射
$\varphi : \mathcal{F} \to \mathcal{G}$ 称为{\it 单射的}，当且仅当
对 $X$ 中每个开集 $U$，映射
$\mathcal{F}(U) \to \mathcal{G}(U)$ 都是单射。
\item $X$ 上集合预层的态射
$\varphi : \mathcal{F} \to \mathcal{G}$ 称为{\it 满射的}，当且仅当
对 $X$ 中每个开集 $U$，映射
$\mathcal{F}(U) \to \mathcal{G}(U)$ 都是满射。
\item $X$ 上集合层的态射
$\varphi : \mathcal{F} \to \mathcal{G}$ 称为{\it 单射的}，当且仅当
对 $X$ 中每个开集 $U$，映射
$\mathcal{F}(U) \to \mathcal{G}(U)$ 都是单射。
\item $X$ 上集合层的态射
$\varphi : \mathcal{F} \to \mathcal{G}$ 称为{\it 满射的}，当且仅当
对 $X$ 的每个开集 $U$ 及 $\mathcal{G}(U)$ 的每个截面 $s$，都存在开覆盖
$U = \bigcup U_i$，使得对所有 $i$，$s|_{U_i}$ 都属于
$\mathcal{F}(U_i) \to \mathcal{G}(U_i)$ 的像。
\end{enumerate}
\end{definition}


\begin{lemma}
\label{sheaves-lemma-characterize-epi-mono}
设 $X$ 为拓扑空间。
\begin{enumerate}
\item 预层范畴中的满态射（相应地，单态射）恰为预层的满射映射
（相应地，单射映射）。
\item 层范畴中的满态射（相应地，单态射）恰为层的满射映射
（相应地，单射映射），也恰为在所有茎上均满射（相应地，单射）的映射。
\end{enumerate}
\end{lemma}

\begin{proof}
略。
\end{proof}


\begin{lemma}
\label{sheaves-lemma-check-homomorphism-stalks}
% P01-ERRATA: P01-E0040 -- the frozen sentence begins with lowercase "let" after the lemma label.
设 $X$ 为拓扑空间。设 $(\mathcal{C}, F)$ 为一种代数结构类型。
假设 $\mathcal{F}$、$\mathcal{G}$ 是 $X$ 上取值于 $\mathcal{C}$ 的层。
设 $\varphi : \mathcal{F} \to \mathcal{G}$ 为底层集合层之间的映射。
若对所有点 $x \in X$，映射
$\mathcal{F}_x \to \mathcal{G}_x$ 都是代数结构的态射，则 $\varphi$
是代数结构层的态射。
\end{lemma}

\begin{proof}
设 $U$ 为 $X$ 的开子集。考虑下列（底层）集合图表
$$
\xymatrix{
\mathcal{F}(U) \ar[r] \ar[d] &
\prod_{x \in U} \mathcal{F}_x \ar[d] \\
\mathcal{G}(U) \ar[r] &
\prod_{x \in U} \mathcal{G}_x
}
$$
由假设及前面的结果，除左侧竖直箭头外，其余箭头都是代数结构的态射。
此外，底部水平箭头是单射；见引理
\ref{sheaves-lemma-sheaf-subset-stalks}。因此，由引理
\ref{sheaves-lemma-image-contained-in} 可得结论；另见例
\ref{sheaves-example-application-lemma-image-contained-in}。
\end{proof}

\noindent
阿贝尔层的短正合列等内容将在模层一章中讨论。见 Modules，第
\ref{modules-section-kernels} 节。




\section{层化}
\label{sheaves-section-sheafification}

\noindent
本节说明如何对拓扑空间上的预层作层化。在这种情形下，我们将用茎来描述
层化。这不同于 Sites，第 \ref{sites-section-sheafification} 节所述的一般
过程，并且或许更易理解。

\medskip\noindent
基本构造如下。设 $\mathcal{F}$ 为拓扑空间 $X$ 上的集合预层。
对每个开集 $U \subset X$，定义
$$
\mathcal{F}^{\#}(U)
=
\{
(s_u) \in \prod\nolimits_{u \in U} \mathcal{F}_u
\text{ 使得 }(*)
\}
$$
其中 $(*)$ 是如下性质：
\begin{itemize}
\item[(*)] 对每个 $u \in U$，存在开邻域 $u \in V \subset U$ 以及截面
$\sigma \in \mathcal{F}(V)$，使得对所有 $v \in V$，在
$\mathcal{F}_v$ 中都有 $s_v = (V, \sigma)$。
\end{itemize}
注意，$(*)$ 是对每个 $u \in U$ 的条件；给定 $u \in U$ 后，该条件是否
成立，只依赖于各值 $s_v$，其中 $v$ 位于 $u$ 的任一开邻域中。因此显然，若
$V \subset U \subset X$ 为开集，则投影映射
$$
\prod\nolimits_{u \in U} \mathcal{F}_u
\longrightarrow
\prod\nolimits_{v \in V} \mathcal{F}_v
$$
把 $\mathcal{F}^{\#}(U)$ 的元素映入 $\mathcal{F}^{\#}(V)$。
以这些映射作为限制映射，就使 $\mathcal{F}^\#$ 成为 $X$ 上的集合预层。

\medskip\noindent
此外，第 \ref{sheaves-section-stalks} 节中描述的映射
$\mathcal{F}(U) \to \prod_{u \in U} \mathcal{F}_u$ 显然以
$\mathcal{F}^{\#}(U)$ 为其像的容纳集合。又若
$V \subset U \subset X$ 为开集，则有交换图表
$$
\xymatrix{
\mathcal{F}(U) \ar[r] \ar[d] &
\mathcal{F}^{\#}(U) \ar[r] \ar[d] &
\prod_{u\in U} \mathcal{F}_u \ar[d] \\
\mathcal{F}(V) \ar[r] &
\mathcal{F}^{\#}(V) \ar[r] &
\prod_{v\in V} \mathcal{F}_v
}
$$
其中竖直映射均由限制映射诱导。因此得到一个典范预层态射
$\mathcal{F} \to \mathcal{F}^{\#}$。

\medskip\noindent
在例 \ref{sheaves-example-sheaf-product-pointwise} 中，我们已经看到，规则
$\Pi(\mathcal{F}) : U \mapsto \prod_{u\in U} \mathcal{F}_u$
配备显然的限制映射后是一个层。按构造，$\mathcal{F}^{\#}$ 是它的一个
子预层。换言之，有预层态射
$$
\mathcal{F} \to \mathcal{F}^\# \to \Pi(\mathcal{F}).
$$
此外，把上述序列赋给 $\mathcal{F}$ 的规则显然对预层 $\mathcal{F}$
具有函子性。下面各引理的证明中将使用这一记号。

\begin{lemma}
\label{sheaves-lemma-sheafification-sheaf}
预层 $\mathcal{F}^{\#}$ 是一个层。
\end{lemma}

\begin{proof}
与其阅读这里的证明，读者或许更适合自己找出解释。事实上，该引理成立的
理由与连续函数预层是层的理由相同；见例
\ref{sheaves-example-basic-continuous-maps}（利用“espace étalé”可将这一类比精确化）。

\medskip\noindent
无论如何，设 $U = \bigcup U_i$ 为开覆盖。假设
$s_i = (s_{i, u})_{u \in U_i} \in \mathcal{F}^{\#}(U_i)$，且
$s_i$ 与 $s_j$ 在 $U_i \cap U_j$ 上一致。由于 $\Pi(\mathcal{F})$ 是层，
可在 $\prod_{u\in U} \mathcal{F}_u$ 中找到元素
$s = (s_u)_{u\in U}$，其在 $U_i$ 上的限制为 $s_i$。还需验证性质 $(*)$。
取 $u \in U$。则对某个 $i$ 有 $u \in U_i$。由 $s_i$ 的性质 $(*)$，
存在开集 $V$、$u \in V \subset U_i$ 以及
$\sigma \in \mathcal{F}(V)$，使得对所有 $v \in V$，在
$\mathcal{F}_v$ 中有 $s_{i, v} = (V, \sigma)$。由于
$s_{i, v} = s_v$，便得到 $s$ 的性质 $(*)$。
\end{proof}

\begin{lemma}
\label{sheaves-lemma-stalk-sheafification}
设 $X$ 为拓扑空间，$\mathcal{F}$ 为 $X$ 上的集合预层，并设
$x \in X$。则 $\mathcal{F}_x = \mathcal{F}^\#_x$。
\end{lemma}

\begin{proof}
映射 $\mathcal{F}_x \to \mathcal{F}^\#_x$ 是单射，因为映射
$\mathcal{F}_x \to \Pi(\mathcal{F})_x$ 已经是单射。事实上，有典范映射
$\Pi(\mathcal{F})_x \to \mathcal{F}_x$，它是映射
$\mathcal{F}_x \to \Pi(\mathcal{F})_x$ 的左逆；见例
\ref{sheaves-example-sheaf-product-pointwise-stalk}。
为证明满射性，设 $\overline{s} \in \mathcal{F}^\#_x$。
可找到 $x$ 的开邻域 $U$，使得 $\overline{s}$ 是 $(U, s)$ 的等价类，
其中 $s \in \mathcal{F}^\#(U)$。按定义，这意味着存在 $x$ 的开邻域
$V \subset U$ 以及截面 $\sigma \in \mathcal{F}(V)$，使得 $s|_V$
是 $\sigma$ 在 $\Pi(\mathcal{F})(V)$ 中的像。显然，$(V, \sigma)$ 的类
给出 $\mathcal{F}_x$ 的一个元素，并映到 $\overline{s}$。
\end{proof}

\begin{lemma}
\label{sheaves-lemma-sheafify-universal}
设 $\mathcal{F}$ 为 $X$ 上的集合预层。任何映入集合层的映射
$\mathcal{F} \to \mathcal{G}$ 都唯一分解为
$\mathcal{F} \to \mathcal{F}^\# \to \mathcal{G}$。
\end{lemma}

\begin{proof}
显然有交换图表
$$
\xymatrix{
\mathcal{F} \ar[r] \ar[d] &
\mathcal{F}^\# \ar[r] \ar[d] &
\Pi(\mathcal{F}) \ar[d] \\
\mathcal{G} \ar[r] &
\mathcal{G}^\# \ar[r] &
\Pi(\mathcal{G}) \\
}
$$
因此，只需证明 $\mathcal{G} = \mathcal{G}^\#$。由引理
\ref{sheaves-lemma-points-exactness}，只需证明对每个点 $x \in X$，映射
$\mathcal{G}_x \to \mathcal{G}^\#_x$ 都是双射。这正是上面的引理
\ref{sheaves-lemma-stalk-sheafification}。
\end{proof}

\noindent
这个引理实际上说明有一对伴随函子：
$i : \Sh(X) \to \textit{PSh}(X)$（包含）和
$\# : \textit{PSh}(X) \to \Sh(X)$（层化）。公式为
$$
\Mor_{\textit{PSh}(X)}(\mathcal{F}, i(\mathcal{G}))
=
\Mor_{\Sh(X)}(\mathcal{F}^\#, \mathcal{G})
$$
这说明层化是包含函子的左伴随。见 Categories，第
\ref{categories-section-adjoint} 节。

\begin{example}
\label{sheaves-example-sheafify-constant}
记号见例 \ref{sheaves-example-stalk-constant-presheaf}。映射
$A_p \to \underline{A}$ 诱导出映射
$A_p^\# \to \underline{A}$。容易看出后者是同构。换言之：取值为 $A$
的常值预层的层化，就是取值为 $A$ 的常值层。
\end{example}

\begin{lemma}
\label{sheaves-lemma-separated-presheaf-into-sheaf}
设 $X$ 为拓扑空间。预层 $\mathcal{F}$ 是分离的（见定义
\ref{sheaves-definition-separated}），
当且仅当典范映射 $\mathcal{F} \to \mathcal{F}^\#$ 是单射。
\end{lemma}

\begin{proof}
这由本节中 $\mathcal{F}^\#$ 的构造立即可见。
\end{proof}

\begin{lemma}
\label{sheaves-lemma-sheafify-epi-mono}
设 $X$ 为拓扑空间。集合预层的满射（相应地，单射）态射经层化后仍为
满射（相应地，单射）。
\end{lemma}

\begin{proof}
略。
\end{proof}


\section{阿贝尔预层的层化}
\label{sheaves-section-sheafify-abelian-presheaves}

\noindent
下面这个看似奇怪的引理很可能并非必需，但在处理代数结构预层的层化时
非常方便。

\begin{lemma}
\label{sheaves-lemma-diagram-fibre-product}
设 $X$ 为拓扑空间，$\mathcal{F}$ 为 $X$ 上的集合预层，并设
$U \subset X$ 为开集。存在典范纤维积图表
$$
\xymatrix{
\mathcal{F}^\#(U) \ar[d] \ar[r] &
\Pi(\mathcal{F})(U) \ar[d] \\
\prod_{x \in U} \mathcal{F}_x
\ar[r] &
\prod_{x \in U} \Pi(\mathcal{F})_x
}
$$
其中各映射如下：
\begin{enumerate}
\item 左侧竖直映射的各分量为
$\mathcal{F}^\#(U) \to \mathcal{F}^\#_x = \mathcal{F}_x$，
其中等式来自引理 \ref{sheaves-lemma-stalk-sheafification}。
\item 顶部水平映射来自第 \ref{sheaves-section-sheafification} 节所述的预层映射
$\mathcal{F} \to \Pi(\mathcal{F})$。
\item 右侧竖直映射有显然的分量映射
$\Pi(\mathcal{F})(U) \to \Pi(\mathcal{F})_x$。
\item 底部水平映射的各分量为
$\mathcal{F}_x \to \Pi(\mathcal{F})_x$；它们来自第
\ref{sheaves-section-sheafification} 节所述的预层映射
$\mathcal{F} \to \Pi(\mathcal{F})$。
\end{enumerate}
\end{lemma}

\begin{proof}
该图表显然交换。需要证明它是纤维积图表。底部水平箭头是单射，因为所有
映射 $\mathcal{F}_x \to \Pi(\mathcal{F})_x$ 都是单射（见引理
\ref{sheaves-lemma-stalk-sheafification} 的证明开头）。截面
$s \in \Pi(\mathcal{F})(U)$ 属于 $\mathcal{F}^\#$，当且仅当 $(*)$
成立。而 $(*)$ 表示，在每一点附近，截面 $s$ 都来自 $\mathcal{F}$ 的一个
截面。按茎函子的定义，这等价于：$s$ 在每个茎
$\Pi(\mathcal{F})_x$ 中的值都来自茎 $\mathcal{F}_x$ 的某个元素。
引理得证。
\end{proof}


\begin{lemma}
\label{sheaves-lemma-sheafify-abelian-presheaf}
设 $X$ 为拓扑空间，$\mathcal{F}$ 为 $X$ 上的阿贝尔预层。
则 $\mathcal{F}^\#$ 上存在唯一的阿贝尔层结构，使得
$\mathcal{F} \to \mathcal{F}^\#$ 是阿贝尔预层的态射。此外，有如下
伴随性：
$$
\Mor_{\textit{PAb}(X)}(\mathcal{F}, i(\mathcal{G}))
=
\Mor_{\textit{Ab}(X)}(\mathcal{F}^\#, \mathcal{G}).
$$
\end{lemma}

\begin{proof}
回忆第 \ref{sheaves-section-sheafification} 节定义的集合层
$\Pi(\mathcal{F})$。所有茎 $\mathcal{F}_x$ 都是阿贝尔群；见引理
\ref{sheaves-lemma-stalk-abelian-presheaf}。因此，由例
\ref{sheaves-example-sheaf-product-pointwise-algebraic-structure}，
$\Pi(\mathcal{F})$ 是阿贝尔群层。映射
$\mathcal{F} \to \Pi(\mathcal{F})$ 显然也是阿贝尔预层的态射。
若能证明第 \ref{sheaves-section-sheafification} 节中的条件 $(*)$ 对每个开子集
$U \subset X$ 都定义 $\Pi(\mathcal{F})(U)$ 的一个子群，则
$\mathcal{F}^\#$ 便典范地继承阿贝尔层结构。直接证明这一点相当容易，
我们把寻找一个简洁论证留给读者。这里采用的论证可推广到代数结构预层，
如下所述：
% P01-ERRATA: P01-E0041 -- the frozen source has subject-verb disagreement: "Lemma ... show".
引理 \ref{sheaves-lemma-diagram-fibre-product} 表明 $\mathcal{F}^\#(U)$ 是一个
阿贝尔群图表的纤维积。因此，$\mathcal{F}^\#$ 如所需地成为阿贝尔子群。

\medskip\noindent
注意，此时由引理 \ref{sheaves-lemma-stalk-abelian-presheaf}，
$\mathcal{F}^\#_x$ 是阿贝尔群；而
$\mathcal{F}_x \to \mathcal{F}^\#_x$ 是双射（引理
\ref{sheaves-lemma-stalk-sheafification}），也是阿贝尔群同态。因此
$\mathcal{F}_x \to \mathcal{F}^\#_x$ 是阿贝尔群同构。下文将直接使用
这一点，不再另行说明。

\medskip\noindent
为证明伴随性，使用集合预层层化的伴随性。例如，若
$\psi : \mathcal{F} \to i(\mathcal{G})$ 是预层态射，则得到层态射
$\psi' : \mathcal{F}^\# \to \mathcal{G}$。需要验证它是阿贝尔层的态射。
例如，可以先由引理 \ref{sheaves-lemma-stalk-sheafification} 注意到这在各茎上
成立，再使用上面的引理 \ref{sheaves-lemma-check-homomorphism-stalks}。
\end{proof}



\section{代数结构预层的层化}
\label{sheaves-section-sheafification-presheaves-structures}


\begin{lemma}
\label{sheaves-lemma-sheafify-presheaf-structures}
设 $X$ 为拓扑空间，$(\mathcal{C}, F)$ 为一种代数结构类型，并设
$\mathcal{F}$ 为 $X$ 上取值于 $\mathcal{C}$ 的预层。则存在取值于
$\mathcal{C}$ 的层 $\mathcal{F}^\#$ 以及取值于 $\mathcal{C}$ 的预层态射
$\mathcal{F} \to \mathcal{F}^\#$，并具有以下性质：
\begin{enumerate}
\item 映射 $\mathcal{F} \to \mathcal{F}^\#$ 将 $\mathcal{F}^\#$ 的底层
集合层等同于 $\mathcal{F}$ 的底层集合预层的层化。
\item 对任意态射 $\mathcal{F} \to \mathcal{G}$，其中 $\mathcal{G}$
为取值于 $\mathcal{C}$ 的层，都存在唯一分解
$\mathcal{F} \to \mathcal{F}^\# \to \mathcal{G}$。
\end{enumerate}
\end{lemma}

\begin{proof}
证明与引理 \ref{sheaves-lemma-sheafify-abelian-presheaf} 的证明相同，只需反复应用
引理 \ref{sheaves-lemma-image-contained-in}（另见例
\ref{sheaves-example-application-lemma-image-contained-in}）。不过，主要思想是把
$\mathcal{F}^\#(U)$ 定义为 $\mathcal{C}$ 中下列图表的纤维积
$$
\xymatrix{
 &
\Pi(\mathcal{F})(U) \ar[d] \\
\prod_{x \in U} \mathcal{F}_x
\ar[r] &
\prod_{x \in U} \Pi(\mathcal{F})_x
}
$$
对照引理 \ref{sheaves-lemma-diagram-fibre-product}。
\end{proof}

\section{模预层的层化}
\label{sheaves-section-sheafification-presheaves-modules}

\begin{lemma}
\label{sheaves-lemma-sheafification-presheaf-modules}
设 $X$ 为拓扑空间，$\mathcal{O}$ 为 $X$ 上的环预层，$\mathcal{F}$ 为
$\mathcal{O}$-模预层。设 $\mathcal{O}^\#$ 为 $\mathcal{O}$ 的层化，
并设 $\mathcal{F}^\#$ 为 $\mathcal{F}$ 作为阿贝尔群预层的层化。
存在集合层的映射
$$
\mathcal{O}^\# \times \mathcal{F}^\#
\longrightarrow
\mathcal{F}^\#
$$
使得图表
$$
\xymatrix{
\mathcal{O} \times \mathcal{F} \ar[r] \ar[d] &
\mathcal{F} \ar[d] \\
\mathcal{O}^\# \times \mathcal{F}^\# \ar[r] &
\mathcal{F}^\#
}
$$
交换，并使 $\mathcal{F}^\#$ 成为 $\mathcal{O}^\#$-模层。此外，若
$\mathcal{G}$ 是 $\mathcal{O}^\#$-模层，则任何 $\mathcal{O}$-模预层态射
$\mathcal{F} \to \mathcal{G}$（其靶是限制为 $\mathcal{O}$-模的
$\mathcal{G}$）都唯一分解为
$\mathcal{F} \to \mathcal{F}^\# \to \mathcal{G}$，其中
$\mathcal{F}^\# \to \mathcal{G}$ 是 $\mathcal{O}^\#$-模态射。
\end{lemma}

\begin{proof}
略。
\end{proof}

\noindent
这实际上意味着函子
$i : \textit{Mod}(\mathcal{O}^\#) \to \textit{PMod}(\mathcal{O})$
（它把限制与将层包含进预层结合起来）和引理中的层化函子
${}^\# : \textit{PMod}(\mathcal{O}) \to \textit{Mod}(\mathcal{O}^\#)$
互为伴随。公式为
$$
\Mor_{\textit{PMod}(\mathcal{O})}(\mathcal{F}, i\mathcal{G})
=
\Mor_{\textit{Mod}(\mathcal{O}^\#)}(\mathcal{F}^\#, \mathcal{G})
$$

\medskip\noindent
设 $X$ 为拓扑空间，并设
$\mathcal{O}_1 \to \mathcal{O}_2$ 为 $X$ 上环层的态射。在第
\ref{sheaves-section-presheaves-modules} 节中，我们已经定义了与此情形相伴的
模预层限制函子和变环函子。

\medskip\noindent
若 $\mathcal{F}$ 是 $\mathcal{O}_2$-模层，则 $\mathcal{F}$ 的限制
$\mathcal{F}_{\mathcal{O}_1}$ 显然是 $\mathcal{O}_1$-模层。由此得到
限制函子
$$
\textit{Mod}(\mathcal{O}_2)
\longrightarrow
\textit{Mod}(\mathcal{O}_1)
$$

\medskip\noindent
另一方面，给定 $\mathcal{O}_1$-模层 $\mathcal{G}$，
$\mathcal{O}_2$-模预层
$\mathcal{O}_2 \otimes_{p, \mathcal{O}_1} \mathcal{G}$ 一般并不是层。
因此，定义{\it 张量积层}
$\mathcal{O}_2 \otimes_{\mathcal{O}_1} \mathcal{G}$ 为
$$
\mathcal{O}_2 \otimes_{\mathcal{O}_1} \mathcal{G}
=
(\mathcal{O}_2 \otimes_{p, \mathcal{O}_1} \mathcal{G})^\#
$$
即预层构造的层化。由此得到{\it 变环}函子
$$
\textit{Mod}(\mathcal{O}_1)
\longrightarrow
\textit{Mod}(\mathcal{O}_2)
$$

\begin{lemma}
\label{sheaves-lemma-adjointness-tensor-restrict}
在上述 $X$、$\mathcal{O}_1$、$\mathcal{O}_2$、$\mathcal{F}$ 和
$\mathcal{G}$ 下，存在典范双射
$$
\Hom_{\mathcal{O}_1}(\mathcal{G}, \mathcal{F}_{\mathcal{O}_1})
=
\Hom_{\mathcal{O}_2}(
\mathcal{O}_2 \otimes_{\mathcal{O}_1} \mathcal{G},
\mathcal{F}
)
$$
换言之，限制函子与变环函子互为伴随。
\end{lemma}

\begin{proof}
这由引理 \ref{sheaves-lemma-adjointness-tensor-restrict-presheaves} 以及如下事实得出：
$\Hom_{\mathcal{O}_2}(
\mathcal{O}_2 \otimes_{\mathcal{O}_1} \mathcal{G},
\mathcal{F}
)
=
\Hom_{\mathcal{O}_2}(
\mathcal{O}_2 \otimes_{p, \mathcal{O}_1} \mathcal{G},
\mathcal{F}
)$，因为 $\mathcal{F}$ 是层。
\end{proof}

\begin{lemma}
\label{sheaves-lemma-stalk-tensor-sheaf-modules}
% P01-ERRATA: P01-E0042 -- the frozen source omits "of" in "a sheaf O-modules".
设 $X$ 为拓扑空间。设 $\mathcal{O} \to \mathcal{O}'$ 为 $X$ 上环层的
态射，$\mathcal{F}$ 为 $\mathcal{O}$-模层，并设 $x \in X$。则作为
$\mathcal{O}'_x$-模，有
$$
\mathcal{F}_x \otimes_{\mathcal{O}_x} \mathcal{O}'_x
=
(\mathcal{F} \otimes_\mathcal{O} \mathcal{O}')_x
$$
\end{lemma}

\begin{proof}
这直接由引理 \ref{sheaves-lemma-stalk-tensor-presheaf-modules} 以及取茎与层化可交换
这一事实得到。
\end{proof}

















\section{连续映射与层}
\label{sheaves-section-presheaves-functorial}

\noindent
设 $f : X \to Y$ 为拓扑空间的连续映射。我们将定义预层与层的直像函子和
逆像函子。

\medskip\noindent
设 $\mathcal{F}$ 为 $X$ 上的集合预层。定义 $\mathcal{F}$ 的{\it 直像}
为如下规则
$$
f_*\mathcal{F}(V) = \mathcal{F}(f^{-1}(V))
$$
其中 $V \subset Y$ 为任意开集。给定开集
$V_1 \subset V_2 \subset Y$，限制映射由下图的交换性给出
$$
\xymatrix{
f_*\mathcal{F}(V_2) \ar[d] \ar@{=}[r] &
\mathcal{F}(f^{-1}(V_2)) \ar[d]^{\text{下列对象的限制： }\mathcal{F}} \\
f_*\mathcal{F}(V_1) \ar@{=}[r] &
\mathcal{F}(f^{-1}(V_1))
}
$$
显然，这定义了一个集合预层。该构造显然对预层 $\mathcal{F}$ 具有
函子性，因而得到函子
$$
f_* : \textit{PSh}(X) \longrightarrow \textit{PSh}(Y).
$$

\begin{lemma}
\label{sheaves-lemma-pushforward-sheaf}
设 $f : X \to Y$ 为连续映射，$\mathcal{F}$ 为 $X$ 上的集合层。
则 $f_*\mathcal{F}$ 是 $Y$ 上的层。
\end{lemma}

\begin{proof}
这立即来自如下事实：若 $V = \bigcup V_j$ 是 $Y$ 中的开覆盖，则
$f^{-1}(V) = \bigcup f^{-1}(V_j)$ 是 $X$ 中的开覆盖。
\end{proof}

\noindent
因此得到函子
$$
f_* : \Sh(X) \longrightarrow \Sh(Y).
$$
它与复合以如下强意义相容。

\begin{lemma}
\label{sheaves-lemma-pushforward-composition}
设 $f : X \to Y$ 和 $g : Y \to Z$ 为拓扑空间的连续映射。
函子 $(g \circ f)_*$ 与 $g_* \circ f_*$ 相等（对集合预层和集合层均如此）。
\end{lemma}

\begin{proof}
这是因为 $(g \circ f)_*\mathcal{F}(W) =
\mathcal{F}((g \circ f)^{-1}W)$，
$(g_* \circ f_*)\mathcal{F}(W) = \mathcal{F}(f^{-1} g^{-1} W)$，并且
$(g \circ f)^{-1}W = f^{-1} g^{-1} W$。
\end{proof}

\noindent
设 $\mathcal{G}$ 为 $Y$ 上的集合预层。给定预层 $\mathcal{G}$ 的
{\it 逆像预层} $f_p\mathcal{G}$，定义为预层上直像 $f_*$ 的左伴随。
换言之，它应当是 $X$ 上的预层 $f_p \mathcal{G}$，满足
$$
\Mor_{\textit{PSh}(X)}(f_p\mathcal{G}, \mathcal{F})
=
\Mor_{\textit{PSh}(Y)}(\mathcal{G}, f_*\mathcal{F}).
$$
由米田引理，这唯一确定了逆像。事实上，它确实存在。

\begin{lemma}
\label{sheaves-lemma-pullback-presheaves}
设 $f : X \to Y$ 为连续映射。存在函子
$f_p : \textit{PSh}(Y) \to \textit{PSh}(X)$，它是 $f_*$ 的左伴随。
对预层 $\mathcal{G}$，该函子由规则
$$
f_p\mathcal{G}(U) = \colim_{f(U) \subset V} \mathcal{G}(V)
$$
确定，其中余极限取遍 $Y$ 中 $f(U)$ 的所有开邻域 $V$。
这些余极限取遍有向偏序集。（$f_p\mathcal{G}$ 的限制映射将在证明中说明。）
\end{lemma}

\begin{proof}
余极限取遍所有包含 $f(U)$ 的开子集 $V \subset Y$ 所成的偏序集，
次序取反向包含。这是有向偏序集，因为若 $V, V'$ 属于其中，则
$V \cap V'$ 也属于其中。此外，若 $U_1 \subset U_2$，则 $f(U_2)$ 的
每个开邻域也是 $f(U_1)$ 的开邻域。因此，定义
$f_p\mathcal{G}(U_2)$ 的系统是定义 $f_p\mathcal{G}(U_1)$ 的系统的
子系统，从而得到限制映射（例如应用 Categories，引理
\ref{categories-lemma-functorial-colimit} 中的一般结论）。

\medskip\noindent
注意，余极限的构造显然对 $\mathcal{G}$ 具有函子性，限制映射亦然。
因此，我们已经把 $f_p$ 定义为函子。

\medskip\noindent
一个有用的小观察是，存在典范映射
$\mathcal{G}(U) \to f_p\mathcal{G}(f^{-1}(U))$，因为
$f(f^{-1}(U))$ 的开邻域系统包含元素 $U$。它与限制映射相容。
换言之，存在典范映射
$i_\mathcal{G} : \mathcal{G} \to f_* f_p \mathcal{G}$。

\medskip\noindent
设 $\mathcal{F}$ 为 $X$ 上的集合预层。假设
$\psi : f_p\mathcal{G} \to \mathcal{F}$ 是集合预层的映射。
相应的映射 $\mathcal{G} \to f_*\mathcal{F}$ 为
$f_*\psi \circ i_\mathcal{G} :
\mathcal{G} \to f_* f_p \mathcal{G} \to f_* \mathcal{F}$。

\medskip\noindent
另一个有用的小观察是，存在典范映射
$c_\mathcal{F} : f_p f_* \mathcal{F} \to \mathcal{F}$。
事实上，设 $U \subset X$ 为开集。对 $Y$ 中每个开邻域
$V \supset f(U)$，都存在映射
$f_*\mathcal{F}(V) = \mathcal{F}(f^{-1}(V))\to \mathcal{F}(U)$，
即 $\mathcal{F}$ 上的限制映射。并且，这与 $\mathcal{F}$ 在包含
$f(U)$ 的各个开集之 $f^{-1}$ 上取值之间的限制映射相容。因此得到典范映射
$f_p f_* \mathcal{F}(U) \to \mathcal{F}(U)$。另一个平凡的验证表明，
这些映射与限制映射相容，并定义了集合预层的映射 $c_\mathcal{F}$。

\medskip\noindent
假设 $\varphi : \mathcal{G} \to f_*\mathcal{F}$ 是集合预层的映射。
考虑 $f_p\varphi :
f_p \mathcal{G} \to f_p f_* \mathcal{F}$。再与 $c_\mathcal{F}$ 后复合，
便得到所需映射
$c_\mathcal{F} \circ f_p\varphi : f_p\mathcal{G} \to \mathcal{F}$。
我们略去验证此构造与上面给出的反向构造互为逆过程。
\end{proof}

\begin{lemma}
\label{sheaves-lemma-stalk-pullback-presheaf}
设 $f : X \to Y$ 为连续映射，$x \in X$，并设 $\mathcal{G}$ 为 $Y$
上的集合预层。存在茎的典范双射
$(f_p\mathcal{G})_x = \mathcal{G}_{f(x)}$。
\end{lemma}

\begin{proof}
可如下看出这一点
\begin{eqnarray*}
(f_p\mathcal{G})_x
& = &
\colim_{x \in U} f_p\mathcal{G}(U) \\
& = &
\colim_{x \in U} \colim_{f(U) \subset V} \mathcal{G}(V) \\
& = &
\colim_{f(x) \in V} \mathcal{G}(V) \\
& = &
\mathcal{G}_{f(x)}
\end{eqnarray*}
这里使用了 Categories，引理 \ref{categories-lemma-colimits-commute}，
以及如下事实：$Y$ 中任意包含 $f(x)$ 的开集 $V$ 都出现在上面的第三种
描述中。细节略去。
\end{proof}

\noindent
设 $\mathcal{G}$ 为 $Y$ 上的集合层。定义{\it 逆像层}
$f^{-1}\mathcal{G}$ 为
$$
f^{-1}\mathcal{G} = (f_p\mathcal{G})^\# .
$$
逆像 $f^{-1}$ 是层上直像的左伴随。换言之，
$$
\Mor_{\Sh(X)}(f^{-1}\mathcal{G}, \mathcal{F})
=
\Mor_{\Sh(Y)}(\mathcal{G}, f_*\mathcal{F}).
$$
事实上，有
\begin{eqnarray*}
\Mor_{\Sh(X)}(f^{-1}\mathcal{G}, \mathcal{F})
& = &
\Mor_{\textit{PSh}(X)}(f_p\mathcal{G}, \mathcal{F}) \\
& = &
\Mor_{\textit{PSh}(Y)}(\mathcal{G}, f_*\mathcal{F}) \\
& = &
\Mor_{\Sh(Y)}(\mathcal{G}, f_*\mathcal{F})
\end{eqnarray*}
第一个等式使用了层化是层到预层之包含的左伴随；第二个等式使用了在预层
上 $f_p$ 是 $f_*$ 的左伴随。我们将在引理
\ref{sheaves-lemma-f-map} 的证明中回到这一陈述。

\begin{lemma}
\label{sheaves-lemma-stalk-pullback}
设 $x \in X$，并设 $\mathcal{G}$ 为 $Y$ 上的集合层。存在茎的典范双射
$(f^{-1}\mathcal{G})_x = \mathcal{G}_{f(x)}$。
\end{lemma}

\begin{proof}
这是引理 \ref{sheaves-lemma-stalk-sheafification} 与
\ref{sheaves-lemma-stalk-pullback-presheaf} 的结合。
\end{proof}

\begin{lemma}
\label{sheaves-lemma-pullback-composition}
设 $f : X \to Y$ 和 $g : Y \to Z$ 为拓扑空间的连续映射。函子
$(g \circ f)^{-1}$ 与 $f^{-1} \circ g^{-1}$ 典范同构。类似地，在预层上
有 $(g \circ f)_p \cong f_p \circ g_p$。
\end{lemma}

\begin{proof}
使用伴随函子在唯一同构意义下唯一这一事实，以及引理
\ref{sheaves-lemma-pushforward-composition} 即可看出。
\end{proof}

\begin{definition}
\label{sheaves-definition-f-map}
设 $f : X \to Y$ 为连续映射，$\mathcal{F}$ 为 $X$ 上的集合层，并设
$\mathcal{G}$ 为 $Y$ 上的集合层。一个{\it $f$-映射
$\xi : \mathcal{G} \to \mathcal{F}$}，是由开子集 $V \subset Y$
索引的一族映射
$\xi_V : \mathcal{G}(V) \to \mathcal{F}(f^{-1}(V))$，使得
$$
\xymatrix{
\mathcal{G}(V) \ar[r]_{\xi_V} \ar[d]_{\text{下列对象的限制： }\mathcal{G}} &
\mathcal{F}(f^{-1}V) \ar[d]^{\text{下列对象的限制： }\mathcal{F}} \\
\mathcal{G}(V') \ar[r]^{\xi_{V'}} &
\mathcal{F}(f^{-1}V')
}
$$
对所有开集 $V' \subset V \subset Y$ 均交换。
\end{definition}

\noindent
文献中有时采用引理 \ref{sheaves-lemma-f-map} 第 (4) 项那样的另一种定义，
但该引理表明二者实际上没有区别。

\begin{lemma}
\label{sheaves-lemma-f-map}
设 $f : X \to Y$ 为连续映射。以下四个集合之间存在双射：
\begin{enumerate}
\item 映射 $\mathcal{G} \to f_*\mathcal{F}$ 的集合；
\item 映射 $f^{-1}\mathcal{G} \to \mathcal{F}$ 的集合；
\item $f$-映射 $\xi : \mathcal{G} \to \mathcal{F}$ 的集合；以及
\item 对所有满足 $f(U) \subset V$ 的开集 $U \subset X$ 和
$V \subset Y$，映射
$\xi_{U, V} : \mathcal{G}(V) \to \mathcal{F}(U)$ 的所有与各限制映射
相容的族所成的集合。
\end{enumerate}
这些双射对 $\mathcal{F} \in \Sh(X)$ 和 $\mathcal{G} \in \Sh(Y)$
具有函子性。
\end{lemma}

\begin{proof}
按定义，层映射 $a : \mathcal{G} \to f_*\mathcal{F}$ 是如下规则：
对 $Y$ 的每个开集 $V$，指定映射
$a_V : \mathcal{G}(V) \to f_*\mathcal{F}(V)$，而且有
$f_*\mathcal{F}(V) = \mathcal{F}(f^{-1}(V))$。因此，至少其“数据”恰好
对应于从 $\mathcal{G}$ 到 $\mathcal{F}$ 的 $f$-映射 $\xi$ 所需的数据。
为证明 (1) 与 (3) 中的集合双射，只需观察：$a$ 是层映射，当且仅当相应的
映射族 $a_V$ 满足定义 \ref{sheaves-definition-f-map} 中的条件。

\medskip\noindent
回忆 $f^{-1}\mathcal{G}$ 是 $f_p\mathcal{G}$ 的层化。由层化的泛性质，
层映射 $b : f^{-1}\mathcal{G} \to \mathcal{F}$ 等价于预层映射
$b_p : f_p\mathcal{G} \to \mathcal{F}$，其中 $f_p$ 是本节前面定义的
函子。给出这样的映射 $b_p$，就需要对 $X$ 的每个开集 $U$ 指定映射
$$
b_{p, U} :
\colim_{f(U) \subset V} \mathcal{G}(V)
\longrightarrow
\mathcal{F}(U)
$$
并使其与限制映射相容。我们可以并且将把 $b_{p, U}$ 看成映射族
$b_{p, U, V} : \mathcal{G}(V) \to \mathcal{F}(U)$，其中 $V$ 取遍 $Y$
中满足 $f(U) \subset V$ 的开集。这些映射必须与所有可能的限制映射
相容。换言之，$b_p$ 对应于 (4) 中的一个映射族。当然，反过来，这样的
映射族定义映射 $b_p$，继而定义映射
$b : f^{-1}\mathcal{G} \to \mathcal{F}$。

\medskip\noindent
为完成引理的证明，必须说明如下“遗忘结构”规则是双射：它把 (4) 中的族
$\xi_{U, V}$ 送到满足 $\xi_V = \xi_{f^{-1}(V), V}$ 的 $f$-映射 $\xi$。
为此，若 $\xi$ 是通常意义下的 $f$-映射，只需把
$\tilde \xi_{U, V}$ 定义为 $\xi_V : \mathcal{G}(V) \to
\mathcal{F}(f^{-1}(V))$ 与限制映射
% P01-ERRATA: P01-E0043 -- the frozen source misspells "restriction" as "restruction".
$\mathcal{F}(f^{-1}(V)) \to \mathcal{F}(U)$ 的复合。它恰因
$f(U) \subset V$，亦即 $U \subset f^{-1}(V)$，而有定义。证明完成。
\end{proof}

\noindent
有时，考虑 $f$-映射比考虑 $X$ 或 $Y$ 上各层之间的映射更为方便。
如下定义 $f$-映射的复合。

\begin{definition}
\label{sheaves-definition-composition-f-maps}
假设 $f : X \to Y$ 和 $g : Y \to Z$ 是拓扑空间的连续映射。假设
$\mathcal{F}$ 是 $X$ 上的层，$\mathcal{G}$ 是 $Y$ 上的层，且
$\mathcal{H}$ 是 $Z$ 上的层。设
$\varphi : \mathcal{G} \to \mathcal{F}$ 为 $f$-映射。
% P01-ERRATA: P01-E0044 -- two frozen loci use the article "an" before "$g$-map".
设 $\psi : \mathcal{H} \to \mathcal{G}$ 为 $g$-映射。
{\it $\varphi$ 与 $\psi$ 的复合}是 $(g \circ f)$-映射
$\varphi \circ \psi$，由下列图表的交换性定义
$$
\xymatrix{
\mathcal{H}(W) \ar[rr]_{(\varphi \circ \psi)_W}
\ar[rd]_{\psi_W} & &
\mathcal{F}(f^{-1}g^{-1}W) \\
&
\mathcal{G}(g^{-1}W)
\ar[ru]_{\varphi_{g^{-1}W}}
}
$$
\end{definition}

\noindent
我们把验证此定义确实可行留给读者。另一种理解方式是把
$\varphi \circ \psi$ 看成复合
$$
\mathcal{H}
\xrightarrow{\psi}
g_*\mathcal{G}
\xrightarrow{g_*\varphi}
g_* f_* \mathcal{F} = (g \circ f)_* \mathcal{F}
$$
现在，是否觉得使用 $f$-映射思考多少会更容易一些？

\medskip\noindent
最后，给定连续映射 $f : X \to Y$ 和 $f$-映射
$\varphi : \mathcal{G} \to \mathcal{F}$，对所有 $x \in X$ 都有自然的
茎映射
$$
\varphi_x : \mathcal{G}_{f(x)} \longrightarrow \mathcal{F}_x
$$
$\mathcal{G}_{f(x)}$ 中元素的代表 $(V, s)$ 被映到 $\mathcal{F}_x$
中以 $(f^{-1}V, \varphi_V(s))$ 为代表的元素。我们把验证它良定义留给
读者。另一种说法是：它是使所有图表
$$
\xymatrix{
\mathcal{F}(f^{-1}V) \ar[r] &
\mathcal{F}_x \\
\mathcal{G}(V) \ar[r] \ar[u]^{\varphi_V} &
\mathcal{G}_{f(x)} \ar[u]^{\varphi_x}
}
$$
（其中 $f(x) \in V \subset Y$ 为开集）交换的唯一映射。

\begin{lemma}
\label{sheaves-lemma-compose-f-maps-stalks}
假设 $f : X \to Y$ 和 $g : Y \to Z$ 是拓扑空间的连续映射。假设
$\mathcal{F}$ 是 $X$ 上的层，$\mathcal{G}$ 是 $Y$ 上的层，且
$\mathcal{H}$ 是 $Z$ 上的层。设
$\varphi : \mathcal{G} \to \mathcal{F}$ 为 $f$-映射。
% P01-ERRATA: P01-E0044 -- second locus of the same article error.
设 $\psi : \mathcal{H} \to \mathcal{G}$ 为 $g$-映射。设 $x \in X$ 为一点。
茎映射
$(\varphi \circ \psi)_x : \mathcal{H}_{g(f(x))}
\to \mathcal{F}_x$ 是如下复合
$$
\mathcal{H}_{g(f(x))}
\xrightarrow{\psi_{f(x)}}
\mathcal{G}_{f(x)}
\xrightarrow{\varphi_x}
\mathcal{F}_x
$$
\end{lemma}

\begin{proof}
这立即来自定义 \ref{sheaves-definition-composition-f-maps} 及上面对茎映射的定义。
\end{proof}



\section{连续映射与阿贝尔层}
\label{sheaves-section-abelian-presheaves-functorial}

\noindent
设 $f : X \to Y$ 为连续映射。我们断言存在函子
\begin{eqnarray*}
f_* : \textit{PAb}(X) & \longrightarrow & \textit{PAb}(Y) \\
f_* : \textit{Ab}(X) & \longrightarrow & \textit{Ab}(Y) \\
f_p : \textit{PAb}(Y) & \longrightarrow & \textit{PAb}(X) \\
f^{-1} : \textit{Ab}(Y) & \longrightarrow & \textit{Ab}(X)
\end{eqnarray*}
它们具有与第 \ref{sheaves-section-presheaves-functorial} 节中相应函子类似的性质。
为说明这一点，作如下论证。

\medskip\noindent
每个函子的构造方式都与第 \ref{sheaves-section-presheaves-functorial} 节中相应
函子的构造方式相同。之所以可行，是因为该节中的所有余极限都是有向余极限
（不过下面仍会详细说明）。

\medskip\noindent
首先，给定 $X$ 上的阿贝尔预层 $\mathcal{F}$ 和 $Y$ 上的阿贝尔预层
$\mathcal{G}$，将
\begin{eqnarray*}
f_*\mathcal{F}(V) & = & \mathcal{F}(f^{-1}(V)) \\
f_p\mathcal{G}(U) & = & \colim_{f(U) \subset V} \mathcal{G}(V)
\end{eqnarray*}
定义为阿贝尔群。限制映射与集合预层的限制映射相同（而且它们全是阿贝尔群
同态）。

\medskip\noindent
赋值 $\mathcal{F} \mapsto f_*\mathcal{F}$ 和
% P01-ERRATA: P01-E0046 -- the frozen source uses \to rather than \mapsto for its second assignment.
$\mathcal{G} \to f_p\mathcal{G}$ 是阿贝尔群预层范畴上的函子。
这是显然的：例如，阿贝尔预层的映射
$\mathcal{G}_1 \to \mathcal{G}_2$ 给出有向系统之间的映射
$\{\mathcal{G}_1(V)\}_{f(U) \subset V} \to
\{\mathcal{G}_2(V)\}_{f(U) \subset V}$，其中所有映射都是同态，因而又
给出阿贝尔群同态
$f_p\mathcal{G}_1(U) \to f_p\mathcal{G}_2(U)$。

\medskip\noindent
函子 $f_*$ 与 $f_p$ 在阿贝尔群预层范畴上互为伴随，即
$$
\Mor_{\textit{PAb}(X)}(f_p\mathcal{G}, \mathcal{F})
=
\Mor_{\textit{PAb}(Y)}(\mathcal{G}, f_*\mathcal{F}).
$$
为证明这一点，注意引理 \ref{sheaves-lemma-pullback-presheaves} 的证明中的映射
$i_\mathcal{G} : \mathcal{G} \to f_* f_p\mathcal{G}$ 是阿贝尔预层的映射。
因此，若 $\psi : f_p\mathcal{G} \to \mathcal{F}$ 是阿贝尔预层的映射，
则相应的映射 $\mathcal{G} \to f_*\mathcal{F}$，即映射
$f_*\psi \circ i_\mathcal{G} :
\mathcal{G} \to f_* f_p \mathcal{G} \to f_* \mathcal{F}$，
也是阿贝尔预层的映射。反过来，指出引理
\ref{sheaves-lemma-pullback-presheaves} 的证明中的映射
$c_\mathcal{F} : f_p f_* \mathcal{F} \to \mathcal{F}$ 也是阿贝尔预层的
映射（因为它由 $\mathcal{F}$ 的限制映射构成，而这些限制映射全是同态）。
因此，给定阿贝尔预层的映射
$\varphi : \mathcal{G} \to f_*\mathcal{F}$，映射
$c_\mathcal{F} \circ f_p\varphi : f_p\mathcal{G} \to \mathcal{F}$
同样是阿贝尔预层的映射。由于作为集合预层映射上的构造，
$\psi \mapsto f_*\psi$ 与
$\varphi \mapsto c_\mathcal{F} \circ f_p\varphi$ 互为逆构造，
它们在阿贝尔预层的映射上也互为逆构造。

\medskip\noindent
% P01-ERRATA: P01-E0045 -- the frozen source reverses X and Y here; it places F on Y and f_*F on X despite f:X->Y.
若 $\mathcal{F}$ 是 $Y$ 上的阿贝尔层，则 $f_*\mathcal{F}$ 是 $X$ 上的
阿贝尔层。这由阿贝尔层的定义以及集合层情形下的相应事实得出；见引理
\ref{sheaves-lemma-pushforward-sheaf}。这定义了阿贝尔层范畴上的函子 $f_*$。

\medskip\noindent
仍像以前一样，定义 $f^{-1}\mathcal{G} = (f_p\mathcal{G})^\#$。
$f_*$ 与 $f^{-1}$ 的伴随性如集合预层情形一样形式地得出。论证如下：
\begin{eqnarray*}
\Mor_{\textit{Ab}(X)}(f^{-1}\mathcal{G}, \mathcal{F})
& = &
\Mor_{\textit{PAb}(X)}(f_p\mathcal{G}, \mathcal{F}) \\
& = &
\Mor_{\textit{PAb}(Y)}(\mathcal{G}, f_*\mathcal{F}) \\
& = &
\Mor_{\textit{Ab}(Y)}(\mathcal{G}, f_*\mathcal{F})
\end{eqnarray*}

\begin{lemma}
\label{sheaves-lemma-pullback-abelian-stalk}
设 $f : X \to Y$ 为连续映射。
\begin{enumerate}
\item 设 $\mathcal{G}$ 为 $Y$ 上的阿贝尔预层，并设 $x \in X$。
引理 \ref{sheaves-lemma-stalk-pullback-presheaf} 的双射
$\mathcal{G}_{f(x)} \to (f_p\mathcal{G})_x$ 是阿贝尔群同构。
\item 设 $\mathcal{G}$ 为 $Y$ 上的阿贝尔层，并设 $x \in X$。
引理 \ref{sheaves-lemma-stalk-pullback} 的双射
$\mathcal{G}_{f(x)} \to (f^{-1}\mathcal{G})_x$ 是阿贝尔群同构。
\end{enumerate}
\end{lemma}

\begin{proof}
略。
\end{proof}

\noindent
给定连续映射 $f : X \to Y$，以及 $X$ 上的阿贝尔群层
$\mathcal{F}$ 和 $Y$ 上的阿贝尔群层 $\mathcal{G}$，阿贝尔群层的
{\it $f$-映射 $\mathcal{G} \to \mathcal{F}$}这一概念有意义。
可以完全仿照定义 \ref{sheaves-definition-f-map} 来定义（把集合映射换为阿贝尔群
同态），也可以直接说它等同于阿贝尔层的映射
$\mathcal{G} \to f_*\mathcal{F}$。下文将自由使用这一概念。
$\mathcal{G}$ 与 $\mathcal{F}$ 之间的 $f$-映射所成的群，与群
$\Mor_{\textit{Ab}(X)}(f^{-1}\mathcal{G}, \mathcal{F})$ 及
$\Mor_{\textit{Ab}(Y)}(\mathcal{G}, f_*\mathcal{F})$ 典范双射。

\medskip\noindent
$f$-映射的复合完全仿照集合层的 $f$-映射情形定义。此外，给定如上的
$f$-映射 $\mathcal{G} \to \mathcal{F}$，其诱导的茎映射
$$
\varphi_x : \mathcal{G}_{f(x)} \longrightarrow \mathcal{F}_x
$$
是阿贝尔群同态。


\section{连续映射与代数结构层}
\label{sheaves-section-presheaves-structures-functorial}

\noindent
设 $(\mathcal{C}, F)$ 为一种代数结构类型。对拓扑空间 $X$，引入记号：
\begin{enumerate}
\item $\textit{PSh}(X, \mathcal{C})$ 表示取值于 $\mathcal{C}$ 的预层范畴。
\item $\Sh(X, \mathcal{C})$ 表示取值于 $\mathcal{C}$ 的层范畴。
\end{enumerate}
设 $f : X \to Y$ 为拓扑空间的连续映射。与上一节相同的论证表明，
存在函子
\begin{eqnarray*}
f_* : \textit{PSh}(X, \mathcal{C}) &
\longrightarrow &
\textit{PSh}(Y, \mathcal{C}) \\
f_* : \Sh(X, \mathcal{C}) &
\longrightarrow &
\Sh(Y, \mathcal{C}) \\
f_p : \textit{PSh}(Y, \mathcal{C}) &
\longrightarrow &
\textit{PSh}(X, \mathcal{C}) \\
f^{-1} : \Sh(Y, \mathcal{C}) &
\longrightarrow &
\Sh(X, \mathcal{C})
\end{eqnarray*}
其构造方式及性质均与为阿贝尔（预）层构造的函子相同。特别地，有交换图表
$$
\xymatrix{
\textit{PSh}(X, \mathcal{C}) \ar[r]^{f_*} \ar[d]^F &
\textit{PSh}(Y, \mathcal{C}) \ar[d]^F &
\Sh(X, \mathcal{C}) \ar[r]^{f_*} \ar[d]^F &
\Sh(Y, \mathcal{C}) \ar[d]^F
\\
\textit{PSh}(X) \ar[r]^{f_*} &
\textit{PSh}(Y) &
\Sh(X) \ar[r]^{f_*} &
\Sh(Y)
\\
\textit{PSh}(Y, \mathcal{C}) \ar[r]^{f_p} \ar[d]^F &
\textit{PSh}(X, \mathcal{C}) \ar[d]^F &
\Sh(Y, \mathcal{C}) \ar[r]^{f^{-1}} \ar[d]^F &
\Sh(X, \mathcal{C}) \ar[d]^F
\\
\textit{PSh}(Y) \ar[r]^{f_p} &
\textit{PSh}(X) &
\Sh(Y) \ar[r]^{f^{-1}} &
\Sh(X)
}
$$

\medskip\noindent
应当牢记的主要公式如下
\begin{eqnarray*}
f_*\mathcal{F}(V) & = & \mathcal{F}(f^{-1}(V)) \\
f_p\mathcal{G}(U) & = & \colim_{f(U) \subset V} \mathcal{G}(V) \\
f^{-1}\mathcal{G} & = & (f_p\mathcal{G})^\# \\
(f_p\mathcal{G})_x & = & \mathcal{G}_{f(x)} \\
(f^{-1}\mathcal{G})_x & = & \mathcal{G}_{f(x)}
\end{eqnarray*}
每个公式都在范畴 $\mathcal{C}$ 中成立；取底层集合后，便得到集合预层的
相应公式。此外，还有伴随性
\begin{eqnarray*}
\Mor_{\textit{PSh}(X, \mathcal{C})}(f_p\mathcal{G}, \mathcal{F})
& = &
\Mor_{\textit{PSh}(Y, \mathcal{C})}(\mathcal{G}, f_*\mathcal{F}) \\
\Mor_{\Sh(X, \mathcal{C})}(f^{-1}\mathcal{G}, \mathcal{F})
& = &
\Mor_{\Sh(Y, \mathcal{C})}(\mathcal{G}, f_*\mathcal{F}).
\end{eqnarray*}
为证明这些性质，主要步骤是把引理 \ref{sheaves-lemma-pullback-presheaves}
的证明中出现的映射
$$
i_\mathcal{G} : \mathcal{G} \longrightarrow f_*f_p\mathcal{G}
$$
和
$$
c_\mathcal{F} : f_p f_* \mathcal{F} \longrightarrow \mathcal{F}
$$
构造为取值于 $\mathcal{C}$ 的预层态射。由于这些构造与集合预层情形完全
相同，可放心地留给读者。

\medskip\noindent
给定连续映射 $f : X \to Y$，以及 $X$ 上的代数结构层
$\mathcal{F}$ 和 $Y$ 上的代数结构层 $\mathcal{G}$，代数结构层的
{\it $f$-映射 $\mathcal{G} \to \mathcal{F}$}这一概念有意义。
可以完全仿照定义 \ref{sheaves-definition-f-map} 来定义（把集合映射换为
$\mathcal{C}$ 中的态射），也可以直接说它等同于代数结构层的映射
$\mathcal{G} \to f_*\mathcal{F}$。下文将自由使用这一概念。
$\mathcal{G}$ 与 $\mathcal{F}$ 之间的 $f$-映射所成的集合，与集合
$\Mor_{\Sh(X, \mathcal{C})}(f^{-1}\mathcal{G}, \mathcal{F})$ 及
$\Mor_{\Sh(Y, \mathcal{C})}(\mathcal{G}, f_*\mathcal{F})$ 典范双射。

\medskip\noindent
$f$-映射的复合完全仿照集合层的 $f$-映射情形定义。此外，给定如上的
$f$-映射 $\mathcal{G} \to \mathcal{F}$，其诱导的茎映射
$$
\varphi_x : \mathcal{G}_{f(x)} \longrightarrow \mathcal{F}_x
$$
是代数结构的同态。


\begin{lemma}
\label{sheaves-lemma-f-map-sets-algebraic-structures}
设 $f : X \to Y$ 为拓扑空间的连续映射。假设给定 $X$ 上的代数结构层
$\mathcal{F}$ 和 $Y$ 上的代数结构层 $\mathcal{G}$。设
$\varphi : \mathcal{G} \to \mathcal{F}$ 为底层集合层的 $f$-映射。
若对每个开集 $V \subset Y$，集合映射
$\varphi_V : \mathcal{G}(V) \to \mathcal{F}(f^{-1}V)$ 都是
$\mathcal{C}$ 中某个态射在底层集合上的作用，则 $\varphi$ 来自代数结构层
之间唯一的 $f$-态射。
\end{lemma}

\begin{proof}
略。
\end{proof}




\section{连续映射与模层}
\label{sheaves-section-presheaves-modules-functorial}

\noindent
模层的情形更为复杂。原因在于，定义逆像函子与直像函子的自然背景是
赋环空间，我们将在下面定义它。先陈述几个显然的引理。

\begin{lemma}
\label{sheaves-lemma-pushforward-presheaf-module}
设 $f : X \to Y$ 为拓扑空间的连续映射。设 $\mathcal{O}$ 为 $X$ 上的
环预层，$\mathcal{F}$ 为 $\mathcal{O}$-模预层。存在底层集合预层的
自然映射
$$
f_*\mathcal{O} \times f_*\mathcal{F}
\longrightarrow
f_*\mathcal{F}
$$
使 $f_*\mathcal{F}$ 成为 $f_*\mathcal{O}$-模预层。该构造对
$\mathcal{F}$ 具有函子性。
\end{lemma}

\begin{proof}
% P01-ERRATA: P01-E0047 -- two frozen proofs use "Let ... is open" instead of "be open".
设 $V \subset Y$ 为开集。将引理中的映射定义为
$$
f_*\mathcal{O}(V) \times f_*\mathcal{F}(V)
=
\mathcal{O}(f^{-1}V) \times \mathcal{F}(f^{-1}V)
\to
\mathcal{F}(f^{-1}V)
=
f_*\mathcal{F}(V).
$$
其中间的箭头是 $X$ 上的乘法映射。我们把验证它与限制映射相容并在
$f_*\mathcal{F}$ 上定义 $f_*\mathcal{O}$-模结构留给读者。
\end{proof}

\begin{lemma}
\label{sheaves-lemma-pullback-presheaf-module}
设 $f : X \to Y$ 为拓扑空间的连续映射。设 $\mathcal{O}$ 为 $Y$ 上的
环预层，$\mathcal{G}$ 为 $\mathcal{O}$-模预层。存在底层集合预层的
自然映射
$$
f_p\mathcal{O} \times f_p\mathcal{G}
\longrightarrow
f_p\mathcal{G}
$$
使 $f_p\mathcal{G}$ 成为 $f_p\mathcal{O}$-模预层。该构造对
$\mathcal{G}$ 具有函子性。
\end{lemma}

\begin{proof}
% P01-ERRATA: P01-E0047 -- second locus of the same finite-verb error.
设 $U \subset X$ 为开集。将引理中的映射定义为
\begin{eqnarray*}
f_p\mathcal{O}(U) \times f_p\mathcal{G}(U)
& = &
\colim_{f(U) \subset V} \mathcal{O}(V)
\times
\colim_{f(U) \subset V} \mathcal{G}(V) \\
& = &
\colim_{f(U) \subset V} (\mathcal{O}(V)\times \mathcal{G}(V)) \\
& \to &
\colim_{f(U) \subset V} \mathcal{G}(V) \\
& = &
f_p\mathcal{G}(U).
\end{eqnarray*}
其中间的箭头是 $Y$ 上的乘法映射。第二个等式成立，是因为有向余极限与
有限极限可交换；见 Categories，引理
\ref{categories-lemma-directed-commutes}。我们把验证它与限制映射相容并在
$f_p\mathcal{G}$ 上定义 $f_p\mathcal{O}$-模结构留给读者。
\end{proof}

\noindent
设 $f : X \to Y$ 为连续映射。设 $\mathcal{O}_X$ 为 $X$ 上的环预层，
$\mathcal{O}_Y$ 为 $Y$ 上的环预层。于是目前已经定义了函子
\begin{eqnarray*}
f_* : \textit{PMod}(\mathcal{O}_X) &
\longrightarrow &
\textit{PMod}(f_*\mathcal{O}_X) \\
f_p : \textit{PMod}(\mathcal{O}_Y) &
\longrightarrow &
\textit{PMod}(f_p\mathcal{O}_Y)
\end{eqnarray*}
它们满足如下相容性。

\begin{lemma}
\label{sheaves-lemma-adjoint-push-pull-presheaves-modules}
设 $f : X \to Y$ 为拓扑空间的连续映射。设 $\mathcal{O}$ 为 $Y$ 上的
环预层，$\mathcal{G}$ 为 $\mathcal{O}$-模预层，并设 $\mathcal{F}$ 为
$f_p\mathcal{O}$-模预层。则
$$
\Mor_{\textit{PMod}(f_p\mathcal{O})}(f_p\mathcal{G}, \mathcal{F})
=
\Mor_{\textit{PMod}(\mathcal{O})}(\mathcal{G}, f_*\mathcal{F}).
$$
这里使用引理 \ref{sheaves-lemma-pullback-presheaf-module} 和
\ref{sheaves-lemma-pushforward-presheaf-module}，并通过映射
$i_\mathcal{O} : \mathcal{O} \to f_*f_p\mathcal{O}$ 把
$f_*\mathcal{F}$ 看作 $\mathcal{O}$-模（该映射最初在引理
\ref{sheaves-lemma-pullback-presheaves} 的证明中定义）。
\end{lemma}

\begin{proof}
注意，由第 \ref{sheaves-section-abelian-presheaves-functorial} 节，有
$$
\Mor_{\textit{PAb}(X)}(f_p\mathcal{G}, \mathcal{F})
=
\Mor_{\textit{PAb}(Y)}(\mathcal{G}, f_*\mathcal{F}).
$$
因此，需要证明在此对应下，模映射所成的子集彼此对应。此外，该对应由规则
$$
(\psi : f_p\mathcal{G} \to \mathcal{F})
\longmapsto
(f_*\psi \circ i_\mathcal{G} :
\mathcal{G} \to f_* \mathcal{F})
$$
确定，反方向则由规则
$$
(\varphi : \mathcal{G} \to f_* \mathcal{F})
\longmapsto
(c_\mathcal{F} \circ f_p\varphi : f_p\mathcal{G} \to \mathcal{F})
$$
确定，其中 $i_\mathcal{G}$ 和 $c_\mathcal{F}$ 如第
\ref{sheaves-section-abelian-presheaves-functorial} 节所述。因此，利用 $f_*$ 与
$f_p$ 的函子性，只需验证映射
$i_\mathcal{G} : \mathcal{G} \to f_* f_p \mathcal{G}$ 和
$c_\mathcal{F} : f_p f_* \mathcal{F} \to \mathcal{F}$ 与模结构相容；
这留给读者。
\end{proof}

\begin{lemma}
\label{sheaves-lemma-adjoint-pull-push-presheaves-modules}
设 $f : X \to Y$ 为拓扑空间的连续映射。设 $\mathcal{O}$ 为 $X$ 上的
环预层，$\mathcal{F}$ 为 $\mathcal{O}$-模预层，并设 $\mathcal{G}$ 为
$f_*\mathcal{O}$-模预层。则
$$
\Mor_{\textit{PMod}(\mathcal{O})}(
\mathcal{O} \otimes_{p, f_pf_*\mathcal{O}} f_p\mathcal{G}, \mathcal{F})
=
\Mor_{\textit{PMod}(f_*\mathcal{O})}(\mathcal{G}, f_*\mathcal{F}).
$$
这里使用引理 \ref{sheaves-lemma-pullback-presheaf-module} 和
\ref{sheaves-lemma-pushforward-presheaf-module}，并在张量积的定义中使用映射
$c_\mathcal{O} : f_pf_*\mathcal{O} \to \mathcal{O}$。
\end{lemma}

\begin{proof}
这由下列等式得到
\begin{eqnarray*}
\Mor_{\textit{PMod}(\mathcal{O})}(
\mathcal{O} \otimes_{p, f_pf_*\mathcal{O}} f_p\mathcal{G}, \mathcal{F})
& = &
\Mor_{\textit{PMod}(f_pf_*\mathcal{O})}(
f_p\mathcal{G}, \mathcal{F}_{f_pf_*\mathcal{O}}) \\
& = &
\Mor_{\textit{PMod}(f_*\mathcal{O})}(\mathcal{G},
f_*(\mathcal{F}_{f_pf_*\mathcal{O}})) \\
& = &
\Mor_{\textit{PMod}(f_*\mathcal{O})}(\mathcal{G}, f_*\mathcal{F}).
\end{eqnarray*}
第一个等式是引理 \ref{sheaves-lemma-adjointness-tensor-restrict-presheaves}。
第二个等式是引理 \ref{sheaves-lemma-adjoint-push-pull-presheaves-modules}。
第三个等式由阿贝尔层的等式
$f_*(\mathcal{F}_{f_pf_*\mathcal{O}}) = f_*\mathcal{F}$ 给出，且该等式是
$f_*\mathcal{O}$-线性的。事实上，在引理
\ref{sheaves-lemma-pullback-presheaves} 的证明所述伴随下，
$\text{id}_{f_*\mathcal{O}}$ 对应于 $c_\mathcal{O}$，因而
$\text{id}_{f_*\mathcal{O}} = f_*c_\mathcal{O} \circ i_{f_*\mathcal{O}}$。
\end{proof}

\begin{lemma}
\label{sheaves-lemma-pushforward-module}
设 $f : X \to Y$ 为拓扑空间的连续映射。设 $\mathcal{O}$ 为 $X$ 上的
环层，$\mathcal{F}$ 为 $\mathcal{O}$-模层。引理
\ref{sheaves-lemma-pushforward-presheaf-module} 中定义的直像
$f_*\mathcal{F}$ 是 $f_*\mathcal{O}$-模层。
\end{lemma}

\begin{proof}
由定义和引理 \ref{sheaves-lemma-pushforward-sheaf} 显然可得。
\end{proof}

\begin{lemma}
\label{sheaves-lemma-pullback-module}
设 $f : X \to Y$ 为拓扑空间的连续映射。设 $\mathcal{O}$ 为 $Y$ 上的
环层，$\mathcal{G}$ 为 $\mathcal{O}$-模层。存在底层集合预层的自然映射
$$
f^{-1}\mathcal{O} \times f^{-1}\mathcal{G}
\longrightarrow
f^{-1}\mathcal{G}
$$
使 $f^{-1}\mathcal{G}$ 成为 $f^{-1}\mathcal{O}$-模层。
\end{lemma}

\begin{proof}
回忆 $f^{-1}$ 定义为函子 $f_p$ 与层化的复合。因此，该引理是引理
\ref{sheaves-lemma-pullback-presheaf-module} 和引理
\ref{sheaves-lemma-sheafification-presheaf-modules} 的结合。
\end{proof}

\noindent
设 $f : X \to Y$ 为连续映射。设 $\mathcal{O}_X$ 为 $X$ 上的环层，
$\mathcal{O}_Y$ 为 $Y$ 上的环层。于是现在已经定义了函子
\begin{eqnarray*}
f_* : \textit{Mod}(\mathcal{O}_X) &
\longrightarrow &
\textit{Mod}(f_*\mathcal{O}_X) \\
f^{-1} : \textit{Mod}(\mathcal{O}_Y) &
\longrightarrow &
\textit{Mod}(f^{-1}\mathcal{O}_Y)
\end{eqnarray*}
它们满足如下相容性。

\begin{lemma}
\label{sheaves-lemma-adjoint-push-pull-modules}
设 $f : X \to Y$ 为拓扑空间的连续映射。设 $\mathcal{O}$ 为 $Y$ 上的
环层，$\mathcal{G}$ 为 $\mathcal{O}$-模层，并设 $\mathcal{F}$ 为
$f^{-1}\mathcal{O}$-模层。则
$$
\Mor_{\textit{Mod}(f^{-1}\mathcal{O})}(f^{-1}\mathcal{G}, \mathcal{F})
=
\Mor_{\textit{Mod}(\mathcal{O})}(\mathcal{G}, f_*\mathcal{F}).
$$
这里使用引理 \ref{sheaves-lemma-pullback-module} 和
\ref{sheaves-lemma-pushforward-module}，并通过
$\mathcal{O} \to f_*f^{-1}\mathcal{O}$ 限制标量，把
$f_*\mathcal{F}$ 看作 $\mathcal{O}$-模。
\end{lemma}

\begin{proof}
使用下列等式论证
\begin{eqnarray*}
\Mor_{\textit{Mod}(f^{-1}\mathcal{O})}(f^{-1}\mathcal{G}, \mathcal{F})
& = &
\Mor_{\textit{Mod}(f_p\mathcal{O})}(f_p\mathcal{G}, \mathcal{F}) \\
& = &
\Mor_{\textit{Mod}(\mathcal{O})}(\mathcal{G}, f_*\mathcal{F}).
\end{eqnarray*}
其中第二个等式来自引理
\ref{sheaves-lemma-adjoint-push-pull-presheaves-modules}，第一个等式来自引理
\ref{sheaves-lemma-sheafification-presheaf-modules}。
\end{proof}

\begin{lemma}
\label{sheaves-lemma-adjoint-pull-push-modules}
设 $f : X \to Y$ 为拓扑空间的连续映射。设 $\mathcal{O}$ 为 $X$ 上的
环层，$\mathcal{F}$ 为 $\mathcal{O}$-模层，并设 $\mathcal{G}$ 为
$f_*\mathcal{O}$-模层。则
$$
\Mor_{\textit{Mod}(\mathcal{O})}(
\mathcal{O} \otimes_{f^{-1}f_*\mathcal{O}} f^{-1}\mathcal{G}, \mathcal{F})
=
\Mor_{\textit{Mod}(f_*\mathcal{O})}(\mathcal{G}, f_*\mathcal{F}).
$$
这里使用引理 \ref{sheaves-lemma-pullback-module} 和
\ref{sheaves-lemma-pushforward-module}，并在张量积的定义中使用典范映射
$f^{-1}f_*\mathcal{O} \to \mathcal{O}$。
\end{lemma}

\begin{proof}
这由下列等式得到
\begin{eqnarray*}
\Mor_{\textit{Mod}(\mathcal{O})}(
\mathcal{O} \otimes_{f^{-1}f_*\mathcal{O}} f^{-1}\mathcal{G}, \mathcal{F})
& = &
\Mor_{\textit{Mod}(f^{-1}f_*\mathcal{O})}(
f^{-1}\mathcal{G}, \mathcal{F}_{f^{-1}f_*\mathcal{O}}) \\
& = &
\Mor_{\textit{Mod}(f_*\mathcal{O})}(\mathcal{G}, f_*\mathcal{F}).
\end{eqnarray*}
它们是引理 \ref{sheaves-lemma-adjointness-tensor-restrict} 与
\ref{sheaves-lemma-adjoint-push-pull-modules} 的结合。
\end{proof}

\noindent
设 $f : X \to Y$ 为连续映射。设 $\mathcal{O}_X$ 为 $X$ 上的环（预）层，
$\mathcal{O}_Y$ 为 $Y$ 上的环（预）层。于是目前已经定义了函子
\begin{eqnarray*}
f_* : \textit{PMod}(\mathcal{O}_X) &
\longrightarrow &
\textit{PMod}(f_*\mathcal{O}_X) \\
f_* : \textit{Mod}(\mathcal{O}_X) &
\longrightarrow &
\textit{Mod}(f_*\mathcal{O}_X) \\
f_p : \textit{PMod}(\mathcal{O}_Y) &
\longrightarrow &
\textit{PMod}(f_p\mathcal{O}_Y) \\
f^{-1} : \textit{Mod}(\mathcal{O}_Y) &
\longrightarrow &
\textit{Mod}(f^{-1}\mathcal{O}_Y)
\end{eqnarray*}
显然，模层上的函子对 $(f_*, f^{-1})$ 通常并不伴随，因为它们的靶范畴
不相同。确切地说，如上所见，只有在环层 $\mathcal{O}_X,
\mathcal{O}_Y$ 奇迹般地满足关系
$\mathcal{O}_X = f^{-1}\mathcal{O}_Y$ 和
$\mathcal{O}_Y = f_*\mathcal{O}_X$ 时才成立，而实践中几乎从不如此。
我们暂时中断讨论，定义恰当的态射概念；对于这种态射，模层上存在适当的
伴随函子对。

\section{赋环空间}
\label{sheaves-section-ringed-spaces}

\noindent
设 $X$ 为拓扑空间，$\mathcal{O}_X$ 为 $X$ 上的环层。应当把环层
$\mathcal{O}_X$ 看作 $X$ 上的函数层。若 $f : X \to Y$ 是“适当的”
映射，则通过复合，$Y$ 上的函数变成 $X$ 上的函数。因此，应当存在从
$\mathcal{O}_Y$ 到 $\mathcal{O}_X$ 的自然 $f$-映射；见定义
\ref{sheaves-definition-f-map} 和引理 \ref{sheaves-lemma-f-map}。一个具体例子见下面的
例 \ref{sheaves-example-continuous-map-ringed}。相关的抽象定义如下。

\begin{definition}
\label{sheaves-definition-ringed-space}
一个{\it 赋环空间}是一个二元组 $(X, \mathcal{O}_X)$，其中 $X$ 为
拓扑空间，$\mathcal{O}_X$ 为 $X$ 上的环层。一个{\it 赋环空间态射}
$(X, \mathcal{O}_X) \to (Y, \mathcal{O}_Y)$ 是一个二元组，由连续映射
$f : X \to Y$ 和环层的 $f$-映射
$f^\sharp : \mathcal{O}_Y \to \mathcal{O}_X$ 组成。
\end{definition}

\begin{example}
\label{sheaves-example-continuous-map-ringed}
设 $f : X \to Y$ 为拓扑空间的连续映射。考虑 $X$ 上的连续实值函数层
$\mathcal{C}^0_X$ 和 $Y$ 上的连续实值函数层 $\mathcal{C}^0_Y$；见例
\ref{sheaves-example-C0-sheaf-rings}。我们断言，存在与 $f$ 相伴的自然 $f$-映射
$f^\sharp : \mathcal{C}^0_Y \to \mathcal{C}^0_X$。事实上，只需用规则
\begin{eqnarray*}
\mathcal{C}^0_Y(V) & \longrightarrow & \mathcal{C}^0_X(f^{-1}V) \\
h & \longmapsto & h \circ f
\end{eqnarray*}
来定义。严格说来，应写作
$f^\sharp(h) = h \circ f|_{f^{-1}(V)}$。显然，这是定义
\ref{sheaves-definition-f-map} 中那样的映射族，并与 $\mathbf{R}$-代数结构相容。
因此，它是 $\mathbf{R}$-代数层的 $f$-映射；见引理
\ref{sheaves-lemma-f-map-sets-algebraic-structures}。

\medskip\noindent
当然，在许多其他情形下，某种几何态射也有典范的赋环空间态射与之相伴。
例如，若 $M$、$N$ 是 $\mathcal{C}^\infty$-流形，且
$f : M \to N$ 是无穷次可微映射，则 $f$ 诱导典范赋环空间态射
$(M, \mathcal{C}_M^\infty) \to (N, \mathcal{C}^\infty_N)$。
其构造（与上述构造完全相同）留给读者。
\end{example}

\noindent
赋环空间态射如何复合也许并非完全显然，因此在此明确写出。

\begin{definition}
\label{sheaves-definition-composition-maps-ringed-spaces}
设
$(f, f^\sharp) : (X, \mathcal{O}_X) \to (Y, \mathcal{O}_Y)$ 和
$(g, g^\sharp) : (Y, \mathcal{O}_Y) \to (Z, \mathcal{O}_Z)$
为赋环空间态射。用规则
$$
(g, g^\sharp) \circ (f, f^\sharp) = (g \circ f, f^\sharp \circ g^\sharp).
$$
定义{\it 赋环空间态射的复合}。这里使用定义
\ref{sheaves-definition-composition-f-maps} 中定义的 $f$-映射复合。
\end{definition}


\section{赋环空间态射与模}
\label{sheaves-section-ringed-spaces-functoriality-modules}

\noindent
现在已经引入足够的记号，可以定义沿赋环空间态射的模之逆像与直像。

\begin{definition}
\label{sheaves-definition-pushforward}
设 $(f, f^\sharp) : (X, \mathcal{O}_X) \to (Y, \mathcal{O}_Y)$
为赋环空间态射。
\begin{enumerate}
\item 设 $\mathcal{F}$ 为 $\mathcal{O}_X$-模层。将 $\mathcal{F}$ 的
{\it 直像}定义为如下 $\mathcal{O}_Y$-模层：作为阿贝尔群层，它等于
$f_*\mathcal{F}$；其模结构是把引理 \ref{sheaves-lemma-pushforward-module}
给出的模结构沿
$f^\sharp : \mathcal{O}_Y \to f_*\mathcal{O}_X$ 限制所得的结构。
\item 设 $\mathcal{G}$ 为 $\mathcal{O}_Y$-模层。将{\it 逆像}
$f^*\mathcal{G}$ 定义为由公式
$$
f^*\mathcal{G}
=
\mathcal{O}_X \otimes_{f^{-1}\mathcal{O}_Y} f^{-1}\mathcal{G}
$$
给出的 $\mathcal{O}_X$-模层，其中环映射
$f^{-1}\mathcal{O}_Y \to \mathcal{O}_X$ 是对应于 $f^\sharp$ 的映射，
模结构则由引理 \ref{sheaves-lemma-pullback-module} 给出。
\end{enumerate}
\end{definition}

\noindent
于是已经定义函子
\begin{eqnarray*}
f_* : \textit{Mod}(\mathcal{O}_X)
& \longrightarrow &
\textit{Mod}(\mathcal{O}_Y) \\
f^* : \textit{Mod}(\mathcal{O}_Y)
& \longrightarrow &
\textit{Mod}(\mathcal{O}_X)
\end{eqnarray*}
关于这些函子的最后一个结果是：它们确如预期那样互为伴随。

\begin{lemma}
\label{sheaves-lemma-adjoint-pullback-pushforward-modules}
设 $(f, f^\sharp) : (X, \mathcal{O}_X) \to (Y, \mathcal{O}_Y)$
为赋环空间态射。设 $\mathcal{F}$ 为 $\mathcal{O}_X$-模层，
$\mathcal{G}$ 为 $\mathcal{O}_Y$-模层。存在典范双射
$$
\Hom_{\mathcal{O}_X}(f^*\mathcal{G}, \mathcal{F})
=
\Hom_{\mathcal{O}_Y}(\mathcal{G}, f_*\mathcal{F}).
$$
换言之，函子 $f^*$ 是 $f_*$ 的左伴随。
\end{lemma}

\begin{proof}
这由前面的工作得到：
\begin{eqnarray*}
\Hom_{\mathcal{O}_X}(f^*\mathcal{G}, \mathcal{F})
& = &
\Mor_{\textit{Mod}(\mathcal{O}_X)}(
\mathcal{O}_X \otimes_{f^{-1}\mathcal{O}_Y} f^{-1}\mathcal{G}, \mathcal{F}) \\
& = &
\Mor_{\textit{Mod}(f^{-1}\mathcal{O}_Y)}(
f^{-1}\mathcal{G}, \mathcal{F}_{f^{-1}\mathcal{O}_Y}) \\
& = &
\Hom_{\mathcal{O}_Y}(\mathcal{G}, f_*\mathcal{F}).
\end{eqnarray*}
这里使用引理 \ref{sheaves-lemma-adjointness-tensor-restrict} 和
\ref{sheaves-lemma-adjoint-push-pull-modules}。
\end{proof}

\begin{lemma}
\label{sheaves-lemma-push-pull-composition-modules}
设 $f : X \to Y$ 和 $g : Y \to Z$ 为赋环空间态射。函子
$(g \circ f)_*$ 与 $g_* \circ f_*$ 相等，并有函子的典范同构
$(g \circ f)^* \cong f^* \circ g^*$。
\end{lemma}

\begin{proof}
关于直像的结论来自引理 \ref{sheaves-lemma-pushforward-composition} 及我们的定义。
关于逆像的结论则由与引理 \ref{sheaves-lemma-pullback-composition} 的证明相同的
论证得到。
\end{proof}


\noindent
给定赋环空间态射
$(f, f^\sharp) : (X, \mathcal{O}_X) \to (Y, \mathcal{O}_Y)$，
$\mathcal{O}_X$-模层 $\mathcal{F}$，以及 $Y$ 上的
$\mathcal{O}_Y$-模层 $\mathcal{G}$，模层的
{\it $f$-映射 $\varphi : \mathcal{G} \to \mathcal{F}$}这一概念有意义。
可以把它定义为阿贝尔层的 $f$-映射
$\varphi : \mathcal{G} \to \mathcal{F}$（见定义
\ref{sheaves-definition-f-map} 和引理 \ref{sheaves-lemma-f-map}），并要求对所有开集
$V \subset Y$，映射
$$
\mathcal{G}(V) \longrightarrow \mathcal{F}(f^{-1}V)
$$
都是 $\mathcal{O}_Y(V)$-模映射。这里通过映射
$f^\sharp_V : \mathcal{O}_Y(V) \to \mathcal{O}_X(f^{-1}V)$，把
$\mathcal{F}(f^{-1}V)$ 看作 $\mathcal{O}_Y(V)$-模。
$\mathcal{G}$ 与 $\mathcal{F}$ 之间的 $f$-映射所成的集合，与集合
$\Mor_{\textit{Mod}(\mathcal{O}_X)}(f^*\mathcal{G}, \mathcal{F})$ 及
$\Mor_{\textit{Mod}(\mathcal{O}_Y)}(\mathcal{G}, f_*\mathcal{F})$
典范双射。见上文。

\medskip\noindent
$f$-映射的复合完全仿照集合层的 $f$-映射情形定义。此外，给定如上的
$f$-映射 $\mathcal{G} \to \mathcal{F}$ 和 $x \in X$，诱导的茎映射
$$
\varphi_x : \mathcal{G}_{f(x)} \longrightarrow \mathcal{F}_x
$$
是 $\mathcal{O}_{Y, f(x)}$-模映射，其中 $\mathcal{F}_x$ 上的
$\mathcal{O}_{Y, f(x)}$-模结构来自 $\mathcal{O}_{X, x}$-模结构以及映射
$f^\sharp_x : \mathcal{O}_{Y, f(x)} \to \mathcal{O}_{X, x}$。
下面是一个相关引理。

\begin{lemma}
\label{sheaves-lemma-stalk-pullback-modules}
设 $(f, f^\sharp) : (X, \mathcal{O}_X) \to (Y, \mathcal{O}_Y)$
为赋环空间态射。设 $\mathcal{G}$ 为 $\mathcal{O}_Y$-模层，并设
$x \in X$。则作为 $\mathcal{O}_{X, x}$-模，有
$$
(f^*\mathcal{G})_x =
\mathcal{G}_{f(x)}
\otimes_{\mathcal{O}_{Y, f(x)}}
\mathcal{O}_{X, x}
$$
其中右侧的张量积使用
$f^\sharp_x : \mathcal{O}_{Y, f(x)} \to \mathcal{O}_{X, x}$。
\end{lemma}

\begin{proof}
这由引理 \ref{sheaves-lemma-stalk-tensor-sheaf-modules} 以及 $x$ 处逆像层的茎与
$f(x)$ 处相应茎的等同得到。例如，见第
\ref{sheaves-section-presheaves-structures-functorial} 节中的公式。
\end{proof}



\section{摩天楼层与茎}
\label{sheaves-section-skyscraper-sheaves}

\begin{definition}
\label{sheaves-definition-skyscraper-sheaf}
设 $X$ 为拓扑空间。
\begin{enumerate}
\item 设 $x \in X$ 为一点。记 $i_x : \{x\} \to X$ 为包含映射。
设 $A$ 为集合，并把 $A$ 看作单点空间 $\{x\}$ 上的层。称
$i_{x, *}A$ 为{\it 在 $x$ 处取值为 $A$ 的摩天楼层}。
\item 若上述 (1) 中的 $A$ 是阿贝尔群，则把 $i_{x, *}A$ 看作 $X$
上的阿贝尔群层。
\item 若上述 (1) 中的 $A$ 是代数结构，则把 $i_{x, *}A$ 看作代数结构层。
\item 若 $(X, \mathcal{O}_X)$ 是赋环空间，则把
$i_x : \{x\} \to X$ 看作赋环空间态射
$(\{x\}, \mathcal{O}_{X, x}) \to (X, \mathcal{O}_X)$；若 $A$ 是
$\mathcal{O}_{X, x}$-模，则把 $i_{x, *}A$ 看作
$\mathcal{O}_X$-模层。
\item 若存在 $X$ 的一点 $x$ 和集合 $A$，使得
$\mathcal{F} \cong i_{x, *}A$，则称集合层 $\mathcal{F}$ 为
{\it 摩天楼层}。
\item 若存在 $X$ 的一点 $x$ 和阿贝尔群 $A$，使得
$\mathcal{F} \cong i_{x, *}A$ 在阿贝尔群层意义下成立，则称阿贝尔群层
$\mathcal{F}$ 为{\it 摩天楼层}。
\item 若存在 $X$ 的一点 $x$ 和代数结构 $A$，使得
$\mathcal{F} \cong i_{x, *}A$ 在代数结构层意义下成立，则称代数结构层
$\mathcal{F}$ 为{\it 摩天楼层}。
\item 若 $(X, \mathcal{O}_X)$ 是赋环空间，$\mathcal{F}$ 是
$\mathcal{O}_X$-模层，并且存在一点 $x \in X$ 和
$\mathcal{O}_{X, x}$-模 $A$，使得
$\mathcal{F} \cong i_{x, *}A$ 在 $\mathcal{O}_X$-模层意义下成立，
则称 $\mathcal{F}$ 为{\it 摩天楼层}。
\end{enumerate}
\end{definition}

\begin{lemma}
\label{sheaves-lemma-skyscraper-stalks}
设 $X$ 为拓扑空间，$x \in X$ 为一点，$A$ 为集合。对任意点
$x' \in X$，在 $x$ 处取值为 $A$ 的摩天楼层在 $x'$ 处的茎为
$$
(i_{x, *}A)_{x'} =
\left\{
\begin{matrix}
A & \text{若} & x' \in \overline{\{x\}} \\
\{*\} & \text{若} & x' \not\in \overline{\{x\}}
\end{matrix}
\right.
$$
阿贝尔群、代数结构及模层情形下有类似描述。
\end{lemma}

\begin{proof}
略。
\end{proof}

\begin{lemma}
\label{sheaves-lemma-stalk-skyscraper-adjoint}
设 $X$ 为拓扑空间，$x \in X$ 为一点。函子
$\mathcal{F} \mapsto \mathcal{F}_x$ 与 $A \mapsto i_{x, *}A$ 互为伴随。
公式为
$$
\Mor_{\textit{Sets}}(\mathcal{F}_x, A)
=
\Mor_{\Sh(X)}(\mathcal{F}, i_{x, *}A).
$$
阿贝尔群和代数结构情形下也有类似陈述。对模层则有
$$
\Hom_{\mathcal{O}_{X, x}}(\mathcal{F}_x, A)
=
\Hom_{\mathcal{O}_X}(\mathcal{F}, i_{x, *}A).
$$
\end{lemma}

\begin{proof}
略。提示：茎函子可视为态射 $i_x : \{x\} \to X$ 的逆像函子。
于是，伴随性来自 $i_x^{-1}$ 与 $i_{x, *}$ 的伴随性（对模层则来自
$i_x^*$ 与 $i_{x, *}$ 的伴随性）。
\end{proof}








\section{预层的极限与余极限}
\label{sheaves-section-limits-presheaves}

\noindent
设 $X$ 为拓扑空间，并设
$\mathcal{I} \to \textit{PSh}(X)$，$i \mapsto \mathcal{F}_i$ 为图表。
\begin{enumerate}
\item $\lim_i \mathcal{F}_i$ 与 $\colim_i \mathcal{F}_i$ 都存在。
\item 对任意开集 $U \subset X$，有
$$
(\lim_i \mathcal{F}_i)(U) =
\lim_i \mathcal{F}_i(U)
$$
以及
$$
(\colim_i \mathcal{F}_i)(U) =
\colim_i \mathcal{F}_i(U).
$$
\item 设 $x \in X$ 为一点。一般而言，$\lim_i \mathcal{F}_i$ 在 $x$
处的茎不等于各茎的极限。但若指标范畴有限，则二者相等。换言之，茎函子
是左正合的（见 Categories，定义 \ref{categories-definition-exact}）。
\item 设 $x \in X$。总有
$$
(\colim_i \mathcal{F}_i)_x =
\colim_i \mathcal{F}_{i, x}.
$$
\end{enumerate}
所有证明都很容易。

\section{层的极限与余极限}
\label{sheaves-section-limits-sheaves}

\noindent
设 $X$ 为拓扑空间，并设
$\mathcal{I} \to \Sh(X)$，$i \mapsto \mathcal{F}_i$ 为图表。
\begin{enumerate}
\item $\lim_i \mathcal{F}_i$ 与 $\colim_i \mathcal{F}_i$ 都存在。
\item 包含函子 $i : \Sh(X) \to \textit{PSh}(X)$ 与极限可交换。
换言之，可以把层范畴中的极限作为预层范畴中的极限来计算。特别地，对任意
开集 $U \subset X$，有
$$
(\lim_i \mathcal{F}_i)(U) = \lim_i \mathcal{F}_i(U).
$$
\item 包含函子 $i : \Sh(X) \to \textit{PSh}(X)$ 一般不与余极限
可交换（甚至不与有限余极限可交换——想想满射）。余极限的计算方法是：
先在预层范畴中取余极限，再层化，即
$$
\colim_i \mathcal{F}_i =
\Big(U \mapsto \colim_i \mathcal{F}_i(U)\Big)^\#.
$$
\item 设 $x \in X$ 为一点。一般而言，$\lim_i \mathcal{F}_i$ 在 $x$
处的茎不等于各茎的极限。但若指标范畴有限，则二者相等。换言之，茎函子
是左正合的。
\item 设 $x \in X$。总有
$$
(\colim_i \mathcal{F}_i)_x = \colim_i \mathcal{F}_{i, x}.
$$
\item 层化函子 ${}^\# : \textit{PSh}(X) \to \Sh(X)$ 与所有余极限
及有限极限可交换，但不与所有极限可交换。
\end{enumerate}
所有证明都很容易。下面给出一个关于层的有向余极限的结论。

\begin{lemma}
\label{sheaves-lemma-directed-colimits-sections}
设 $X$ 为拓扑空间，$I$ 为有向集。设
$(\mathcal{F}_i, \varphi_{ii'})$ 为 $I$ 上的集合层系统；见 Categories，
第 \ref{categories-section-posets-limits} 节。设 $U \subset X$ 为开子集。
考虑典范映射
$$
\Psi :
\colim_i \mathcal{F}_i(U)
\longrightarrow
\left(\colim_i \mathcal{F}_i\right)(U)
$$
\begin{enumerate}
\item 若所有转移映射都是单射，则对任意开集 $U$，$\Psi$ 都是单射。
\item 若 $U$ 拟紧，则 $\Psi$ 是单射。
\item 若 $U$ 拟紧且所有转移映射都是单射，则 $\Psi$ 是同构。
\item 若 $U$ 有一个开覆盖的共尾系
$\mathcal{U} : U = \bigcup_{j\in J} U_j$，其中 $J$ 有限，且对所有
$j, j' \in J$，$U_j \cap U_{j'}$ 均拟紧，则 $\Psi$ 是双射。
\end{enumerate}
\end{lemma}

\begin{proof}
假设所有转移映射都是单射。此时，预层
$\mathcal{F}' : V \mapsto \colim_i \mathcal{F}_i(V)$ 是分离的（见定义
\ref{sheaves-definition-separated}）。由上述讨论，有
$(\mathcal{F}')^\# = \colim_i \mathcal{F}_i$。由引理
\ref{sheaves-lemma-separated-presheaf-into-sheaf}，映射
$\mathcal{F}' \to (\mathcal{F}')^\#$ 是单射。这证明了 (1)。

\medskip\noindent
假设 $U$ 拟紧。设 $s \in \mathcal{F}_i(U)$ 和
$s' \in \mathcal{F}_{i'}(U)$ 在左侧给出的元素经 $\Psi$ 后有相同的像。
由于 $U$ 拟紧，这意味着存在有限开覆盖
$U = \bigcup_{j = 1, \ldots, m} U_j$，并且对每个 $j$，存在指标
$i_j \in I$，$i_j \geq i$，$i_j \geq i'$，使得
$\varphi_{ii_j}(s)$ 与 $\varphi_{i'i_j}(s')$ 在 $U_j$ 上的限制相同。
取 $i''\in I$，使其相对于所有 $i_j$ 均为 $\geq$。于是，对所有 $j$，
$\varphi_{ii''}(s)$ 与 $\varphi_{i'i''}(s')$ 在开集 $U_j$ 上一致，
从而 $\varphi_{ii''}(s) = \varphi_{i'i''}(s')$。这证明了 (2)。

\medskip\noindent
假设 $U$ 拟紧且所有转移映射都是单射。设 $s$ 为 $\Psi$ 的靶中的元素。
由于 $U$ 拟紧，存在有限开覆盖
$U = \bigcup_{j = 1, \ldots, m} U_j$；对每个 $j$，存在指标
$i_j \in I$ 及 $s_j \in \mathcal{F}_{i_j}(U_j)$，使得对所有 $j$，
$s|_{U_j}$ 都来自 $s_j$。取 $i \in I$，使其相对于所有 $i_j$ 均为
$\geq$。由 (1)，
各截面 $\varphi_{i_ji}(s_j)$ 在交集 $U_j \cap U_{j'}$ 上一致。
因此，它们粘合成截面 $s' \in \mathcal{F}_i(U)$，且该截面经 $\Psi$
映到 $s$。这证明了 (3)。

\medskip\noindent
假设 (4) 的条件成立。特别地，$U$ 拟紧，因而由 (2)，$\Psi$ 是单射。
设 $s$ 为 $\Psi$ 的靶中的元素。由假设，存在有限开覆盖
$U = \bigcup_{j = 1, \ldots, m} U_j$，其中对所有 $j, j' \in J$，
$U_j \cap U_{j'}$ 均拟紧；并且对每个 $j$，存在指标 $i_j \in I$ 及
$s_j \in \mathcal{F}_{i_j}(U_j)$，使得对所有 $j$，$s|_{U_j}$ 是
$s_j$ 的像。由于 $U_j \cap U_{j'}$ 拟紧，可以应用 (2)，从而存在
$i_{jj'} \in I$，满足 $i_{jj'} \geq i_j$、$i_{jj'} \geq i_{j'}$，
并使 $\varphi_{i_ji_{jj'}}(s_j)$ 与
$\varphi_{i_{j'}i_{jj'}}(s_{j'})$ 在 $U_j \cap U_{j'}$ 上一致。
选择指标 $i \in I$ 大于等于所有 $i_{jj'}$。于是，
$\mathcal{F}_i$ 的各截面 $\varphi_{i_ji}(s_j)$ 粘合成 $U$ 上的一个
$\mathcal{F}_i$ 截面。该截面如所需地映到元素 $s$。
\end{proof}

\begin{example}
\label{sheaves-example-conditions-needed-colimit}
令 $X = \{s_1, s_2, \xi_1, \xi_2, \xi_3, \ldots\}$ 为集合。规定子集
$U \subset X$ 为开集，当且仅当 $s_1 \in U$ 或 $s_2 \in U$ 蕴含
$U$ 包含所有 $\xi_i$。令 $U_n = \{\xi_n, \xi_{n + 1}, \ldots\}$，
并令 $j_n : U_n \to X$ 为包含映射。置
$\mathcal{F}_n = j_{n, *}\underline{\mathbf{Z}}$。存在转移映射
$\mathcal{F}_n \to \mathcal{F}_{n + 1}$。令
$\mathcal{F} = \colim \mathcal{F}_n$。注意，若 $m < n$，则
$\mathcal{F}_{n, \xi_m} = 0$，因为 $\{\xi_m\}$ 是 $X$ 的开子集，且与
$U_n$ 不相交。因此，对所有 $n$，都有 $\mathcal{F}_{\xi_n} = 0$。
另一方面，茎 $\mathcal{F}_{s_i}$，$i = 1, 2$，是余极限
$$
M = \colim_n \prod\nolimits_{m \geq n} \mathbf{Z}
$$
它并不为零。由此得出，层 $\mathcal{F}$ 是在闭点 $s_1$ 与 $s_2$ 处
取值为 $M$ 的摩天楼层之直和。因此
$\Gamma(X, \mathcal{F}) = M \oplus M$。另一方面，读者可以验证
$\colim_n \Gamma(X, \mathcal{F}_n) = M$。所以，上面的引理
\ref{sheaves-lemma-directed-colimits-sections} 第 (4) 项确实需要某种条件。
\end{example}

\noindent
前一引理还有一个处理谱空间图表上各层的版本。为陈述它，引入一些记号。
设 $\mathcal{I}$ 为余滤过指标范畴。设
$i \mapsto X_i$ 为 $\mathcal{I}$ 上的谱空间图表，使得对
$\mathcal{I}$ 中的 $a : j \to i$，相应映射
$f_a : X_j \to X_i$ 是谱映射。置 $X = \lim X_i$，并记
$p_i : X \to X_i$ 为投影。

\begin{lemma}
\label{sheaves-lemma-compute-pullback-to-limit}
在上述情形下，设 $i \in \Ob(\mathcal{I})$，$\mathcal{G}$ 为 $X_i$
上的层。对拟紧开集 $U_i \subset X_i$，有
$$
p_i^{-1}\mathcal{G}(p_i^{-1}(U_i)) =
\colim_{a : j \to i} f_a^{-1}\mathcal{G}(f_a^{-1}(U_i))
$$
\end{lemma}

\begin{proof}
先证明典范映射
$\colim_{a : j \to i} f_a^{-1}\mathcal{G}(f_a^{-1}(U_i)) \to
p_i^{-1}\mathcal{G}(p_i^{-1}(U_i))$ 是单射。设 $s, s'$ 是某个
$a : j \to i$ 下 $f_a^{-1}(U_i)$ 上 $f_a^{-1}\mathcal{G}$ 的截面。
对 $b : k \to j$，令 $Z_k \subset f_{a \circ b}^{-1}(U_i)$ 为如下点
$x$ 所成的闭子集：$s$ 与 $s'$ 在茎
$(f_{a \circ b}^{-1}\mathcal{G})_x$ 中的像不同。若对所有
$b : k \to j$，$Z_k$ 都非空，则由 Topology，引理
\ref{topology-lemma-inverse-limit-spectral-spaces-nonempty}，
$\lim_{b : k \to j} Z_k$ 也非空。于是，对
$x \in \lim_{b : k \to j} Z_k \subset X$（注意
$\mathcal{I}/j \to \mathcal{I}$ 是始函子），$s$ 与 $s'$ 在
$p_i^{-1}\mathcal{G}$ 的 $x$ 处茎中的像也不同，因为对上述所有
$b : k \to j$，都有
% P01-ERRATA: P01-E0048 -- the frozen source reverses the composable arrows here as b\circ a; surrounding loci use a\circ b.
$(p_i^{-1}\mathcal{G})_x = (f_{b \circ a}^{-1}\mathcal{G})_{p_k(x)}$。
因此，若 $s$ 与 $s'$ 在
$p_i^{-1}\mathcal{G}(p_i^{-1}(U_i))$ 中的像相同，则对某个
$b : k \to j$，$Z_k$ 为空。这证明了单射性。

\medskip\noindent
再证满射性。设 $s$ 为 $p_i^{-1}(U_i)$ 上 $p_i^{-1}\mathcal{G}$ 的截面。
由 Topology，引理
\ref{topology-lemma-directed-inverse-limit-spectral-spaces}，集合
$p_i^{-1}(U_i)$ 是谱空间 $X$ 的拟紧开集。按逆像层的构造，可以找到开覆盖
$p_i^{-1}(U_i) = \bigcup_{l \in L} W_l$、开集
$V_{l, i} \subset X_i$ 以及截面
$s_{l, i} \in \mathcal{G}(V_{l, i})$，使得
$p_i(W_l) \subset V_{l, i}$ 且
$p_i^{-1}s_{l, i}|_{W_l} = s|_{W_l}$。由于 $X$ 和 $X_i$ 都是谱空间，
且 $p_i^{-1}(U_i)$ 是拟紧开集，可以假设 $L$ 有限，并且对所有 $l$，
$W_l$ 与 $V_{l, i}$ 都是拟紧开集。于是可应用 Topology，引理
\ref{topology-lemma-descend-opens}，找到 $a : j \to i$ 和由拟紧开集组成的
开覆盖 $f_a^{-1}(U_i) = \bigcup_{l \in L} W_{l, j}$；它拉回到 $X$
后是开覆盖 $p_i^{-1}(U_i) = \bigcup_{l \in L} W_l$，并且还有
$W_{l, j} \subset f_a^{-1}(V_{l, i})$。把 $s_{l, i}$ 经 $f_a$ 拉回后
限制到 $W_{l, j}$，记为 $s_{l, j}$。于是，$s_{l, j}$ 与
$s_{l', j}$ 在 $(f_a^{-1}\mathcal{G})(W_{l, j} \cap W_{l', j})$
中限制所得的元素，拉回后成为
$(p_i^{-1}\mathcal{G})(W_l \cap W_{l'})$ 中同一个元素，即 $s$ 的限制。
因此，由单射性，可找到 $b : k \to j$，使各截面
$f_b^{-1}s_{l, j}$ 如所需地粘合成 $f_{a \circ b}^{-1}(U_i)$ 上的截面。
\end{proof}

\noindent
接下来，除谱空间的余滤过系统 $X_i$ 外，再假设给定
\begin{enumerate}
\item 对每个 $i \in \Ob(\mathcal{I})$，给定 $X_i$ 上的层
$\mathcal{F}_i$；
\item 对 $a : j \to i$，给定 $f_a$-映射
$\varphi_a : \mathcal{F}_i \to \mathcal{F}_j$。
\end{enumerate}
并且每当 $c = a \circ b$ 时，满足
$\varphi_c = \varphi_b \circ \varphi_a$。在 $X$ 上置
$\mathcal{F} = \colim p_i^{-1}\mathcal{F}_i$。

\begin{lemma}
\label{sheaves-lemma-descend-opens}
在上述情形下，设 $i \in \Ob(\mathcal{I})$，并设
$U_i \subset X_i$ 为拟紧开集。则
$$
\colim_{a : j \to i} \mathcal{F}_j(f_a^{-1}(U_i)) = \mathcal{F}(p_i^{-1}(U_i))
$$
\end{lemma}

\begin{proof}
回忆 $p_i^{-1}(U_i)$ 是谱空间 $X$ 的拟紧开集；见 Topology，引理
\ref{topology-lemma-directed-inverse-limit-spectral-spaces}。因此，引理
\ref{sheaves-lemma-directed-colimits-sections} 适用，并有
$$
\mathcal{F}(p_i^{-1}(U_i)) =
\colim_{a : j \to i} p_j^{-1}\mathcal{F}_j(p_i^{-1}(U_i)).
$$
形式论证表明
$$
\colim_{a : j \to i} \mathcal{F}_j(f_a^{-1}(U_i)) =
\colim_{a : j \to i} \colim_{b : k \to j}
f_b^{-1}\mathcal{F}_j(f_{a \circ b}^{-1}(U_i))
$$
因此，只需证明
$$
p_j^{-1}\mathcal{F}_j(p_i^{-1}(U_i)) =
\colim_{b : k \to j} f_b^{-1}\mathcal{F}_j(f_{a \circ b}^{-1}(U_i))
$$
这正是把引理 \ref{sheaves-lemma-compute-pullback-to-limit} 应用于
$\mathcal{F}_j$ 及拟紧开集 $f_a^{-1}(U_i)$。
\end{proof}

















\section{基与层}
\label{sheaves-section-bases}

\noindent
有时，拓扑存在一个由比一般开集更易处理的开集组成的基。为便于在这种
情形下处理（预）层，这里给出一些定义和简单引理。

\begin{definition}
\label{sheaves-definition-presheaf-basis}
设 $X$ 为拓扑空间，$\mathcal{B}$ 为 $X$ 的拓扑基。
\begin{enumerate}
\item {\it $\mathcal{B}$ 上的集合预层 $\mathcal{F}$}是如下规则：
对每个 $U \in \mathcal{B}$ 指定集合 $\mathcal{F}(U)$；对
$\mathcal{B}$ 中元素的每个包含关系 $V \subset U$ 指定映射
$\rho^U_V : \mathcal{F}(U) \to \mathcal{F}(V)$；并且对所有
$U \in \mathcal{B}$ 都有
$\rho^U_U = \text{id}_{\mathcal{F}(U)}$，每当
$\mathcal{B}$ 中有 $W \subset V \subset U$ 时，都有
$\rho^U_W = \rho^V_W \circ \rho ^U_V$。
\item {\it $\mathcal{B}$ 上集合预层的态射
$\varphi : \mathcal{F} \to \mathcal{G}$}是如下规则：对每个元素
$U \in \mathcal{B}$ 指定与限制映射相容的集合映射
% P01-ERRATA: P01-E0049 -- two basis-presheaf definitions write the component as phi rather than phi_U.
$\varphi : \mathcal{F}(U) \to \mathcal{G}(U)$。
\end{enumerate}
\end{definition}

\noindent
与通常的预层情形一样，我们使用截面、截面的限制等术语。特别地，可把
$\mathcal{F}$ 在一点 $x \in X$ 处的{\it 茎}定义为余极限
$$
\mathcal{F}_x = \colim_{U\in \mathcal{B}, x\in U} \mathcal{F}(U).
$$
与 $X$ 上预层的茎一样，这个极限是有向的。原因在于，所有满足
$U\in \mathcal{B}$、$x \in U$ 的元素构成 $x$ 的开邻域基本系。

\medskip\noindent
很容易构造例子，说明 $X$ 上预层的概念与 $X$ 的拓扑基上预层的概念
相差很大。层的情形则不会发生这种现象。因此，下述概念更为有用。

\begin{definition}
\label{sheaves-definition-sheaf-basis}
设 $X$ 为拓扑空间，$\mathcal{B}$ 为 $X$ 的拓扑基。
\begin{enumerate}
\item {\it $\mathcal{B}$ 上的集合层 $\mathcal{F}$}是
$\mathcal{B}$ 上满足下述附加性质的集合预层：给定任意
$U \in \mathcal{B}$、任意覆盖
$U = \bigcup_{i \in I} U_i$（其中 $U_i \in \mathcal{B}$），以及任意覆盖
$U_i \cap U_j = \bigcup_{k \in I_{ij}} U_{ijk}$（其中
$U_{ijk} \in \mathcal{B}$），层条件成立：
\begin{itemize}
\item[(**)] 对任意截面族 $s_i \in \mathcal{F}(U_i)$，$i \in I$，
若对 $\forall i, j\in I$、$\forall k\in I_{ij}$ 均有
$$
s_i|_{U_{ijk}} = s_j|_{U_{ijk}}
$$
则存在唯一截面 $s \in \mathcal{F}(U)$，使得对所有 $i \in I$ 都有
$s_i = s|_{U_i}$。
\end{itemize}
\item {\it $\mathcal{B}$ 上集合层的态射}就是集合预层的态射。
\end{enumerate}
\end{definition}

\noindent
首先说明，只需在一个覆盖的共尾系上验证层条件 $(**)$。在上述定义的
情形下，设 $U \in \mathcal{B}$。暂记
$\text{Cov}_\mathcal{B}(U)$ 为所有由 $\mathcal{B}$ 中元素组成的 $U$
之覆盖所成的集合。注意，$\text{Cov}_\mathcal{B}(U)$ 按加细关系成为
预序集。若对每个 $\mathcal{U} \in \text{Cov}_\mathcal{B}(U)$，都存在
一个加细 $\mathcal{U}$ 的覆盖 $\mathcal{V} \in C$，则称子集
$C \subset \text{Cov}_\mathcal{B}(U)$ 为共尾系。

\begin{lemma}
\label{sheaves-lemma-cofinal-systems-coverings}
沿用上述记号。对每个 $U \in \mathcal{B}$，设
$C(U) \subset \text{Cov}_\mathcal{B}(U)$ 为共尾系。对每个
$U \in \mathcal{B}$ 以及 $C(U)$ 中的每个覆盖
$\mathcal{U} : U = \bigcup U_i$，给定覆盖
$\mathcal{U}_{ij} : U_i \cap U_j = \bigcup U_{ijk}$，其中
$U_{ijk} \in \mathcal{B}$。设 $\mathcal{F}$ 为 $\mathcal{B}$ 上的
集合预层。则以下条件等价：
\begin{enumerate}
\item 预层 $\mathcal{F}$ 是 $\mathcal{B}$ 上的层。
\item 对每个 $U \in \mathcal{B}$ 以及 $C(U)$ 中的每个覆盖
$\mathcal{U} : U = \bigcup U_i$，层条件 $(**)$ 都成立（相对于给定的
覆盖 $\mathcal{U}_{ij}$）。
\end{enumerate}
\end{lemma}

\begin{proof}
需要证明 (2) 蕴含 (1)。假设 $U \in \mathcal{B}$，且
$\mathcal{U} : U = \bigcup_{i\in I} U_i$ 是由 $\mathcal{B}$ 中元素
组成的任意覆盖。由于系统 $C(U)$ 共尾，可找到 $C(U)$ 中加细
$\mathcal{U}$ 的元素
$\mathcal{V} : U = \bigcup_{j \in J} V_j$。这表示存在映射
$\alpha : J \to I$，使得 $V_j \subset U_{\alpha(j)}$。

\medskip\noindent
注意，若截面 $s, s' \in \mathcal{F}(U)$ 满足
$s|_{U_i} = s'|_{U_i}$，则
$$
s|_{V_j}
= (s|_{U_{\alpha(j)}})|_{V_j}
= (s'|_{U_{\alpha(j)}})|_{V_j}
= s'|_{V_j}
$$
对所有 $j$ 成立。因此，由覆盖 $\mathcal{V}$ 的 $(**)$ 中的唯一性，
得到 $s = s'$。这样便证明了任意覆盖 $\mathcal{U}$ 的 $(**)$ 中的
唯一性部分。

\medskip\noindent
再假设
$U_i \cap U_{i'} = \bigcup_{k \in I_{ii'}} U_{ii'k}$ 是由
$U_{ii'k} \in \mathcal{B}$ 组成的任意覆盖。尝试证明系统
$(\mathcal{U}, \mathcal{U}_{ij})$ 的 $(**)$ 中的存在性部分。为此，设
$s_i \in \mathcal{F}(U_i)$，并假设对所有 $i, i', k$ 均有
$$
s_i|_{U_{ii'k}} = s_{i'}|_{U_{ii'k}}
$$
置 $t_j = s_{\alpha(j)}|_{V_j}$，其中 $\mathcal{V}$ 和 $\alpha$ 如上。

\medskip\noindent
此处论证有一个小障碍。设
$\mathcal{V}_{jj'} : V_j \cap V_{j'} = \bigcup_{l \in J_{jj'}} V_{jj'l}$
为引理陈述中给定的覆盖。事先并不清楚对所有 $j, j', l$ 是否有
$$
t_j|_{V_{jj'l}} = t_{j'}|_{V_{jj'l}}
$$
为说明这一点，注意，由对元素族 $s_i$ 的假设，对所有
$k \in I_{\alpha(j)\alpha(j')}$，确有
$$
t_j|_W = t_{j'}|_W \text{ 对所有 } W \in \mathcal{B},
W \subset V_{jj'l} \cap U_{\alpha(j)\alpha(j')k}
$$
又因为
$V_j \cap V_{j'} \subset U_{\alpha(j)} \cap U_{\alpha(j')}$，可见
$t_j|_{V_{jj'l}}$ 与 $t_{j'}|_{V_{jj'l}}$ 在由 $\mathcal{B}$ 中元素组成的
$V_{jj'l}$ 的某个覆盖的各成员上相同。因此，由上面已经证明的唯一性部分，
最终得到所需的 $t_j|_{V_{jj'l}}$ 与 $t_{j'}|_{V_{jj'l}}$ 的等式。
随后，由 $(\mathcal{V}, \mathcal{V}_{jj'})$ 的性质 $(**)$，得到元素
$t \in \mathcal{F}(U)$ 的存在性。

\medskip\noindent
这里又有一个小问题。已知 $t$ 在 $V_j$ 上限制为 $t_j$，但尚不知道
$t$ 在 $U_i$ 上限制为 $s_i$。为得出结论，注意各集合
$U_i \cap V_j$，$j \in J$，覆盖 $U_i$。因而，各集合
$U_{i \alpha(j) k} \cap V_j$，$j\in J$，$k \in I_{i\alpha(j)}$，
也覆盖 $U_i$。我们把如下验证留给读者：在包含于某个开集
$U_{i \alpha(j) k} \cap V_j$（其中 $j\in J$、
$k \in I_{i\alpha(j)}$）的任意 $W \in \mathcal{B}$ 上，$t$ 与 $s_i$
限制为 $\mathcal{F}$ 的同一个截面。因此，由上述唯一性部分即得结论。
\end{proof}

\begin{lemma}
\label{sheaves-lemma-cofinal-systems-coverings-standard-case}
设 $X$ 为拓扑空间，$\mathcal{B}$ 为 $X$ 的拓扑基。假设对每个三元组
$U, U', U'' \in \mathcal{B}$，只要 $U' \subset U$ 且
$U'' \subset U$，便有 $U' \cap U'' \in \mathcal{B}$。对每个
$U \in \mathcal{B}$，设
$C(U) \subset \text{Cov}_\mathcal{B}(U)$ 为共尾系。设 $\mathcal{F}$
为 $\mathcal{B}$ 上的集合预层。则以下条件等价：
\begin{enumerate}
\item 预层 $\mathcal{F}$ 是 $\mathcal{B}$ 上的层。
\item 对每个 $U \in \mathcal{B}$、$C(U)$ 中的每个覆盖
$\mathcal{U} : U = \bigcup U_i$，以及每个满足
$s_i|_{U_i \cap U_j} = s_j|_{U_i \cap U_j}$ 的截面族
$s_i \in \mathcal{F}(U_i)$，都存在唯一截面
$s \in \mathcal{F}(U)$，其在 $U_i$ 上限制为 $s_i$。
\end{enumerate}
\end{lemma}

\begin{proof}
这是上面引理 \ref{sheaves-lemma-cofinal-systems-coverings} 在每个覆盖
$\mathcal{U}_{ij}$ 只含一个元素这一特殊情形下的改述。不过，这一情形
也简单得多，直接证明是一个容易的练习。
\end{proof}

\begin{lemma}
\label{sheaves-lemma-condition-star-sections}
设 $X$ 为拓扑空间，$\mathcal{B}$ 为 $X$ 的拓扑基，
$U \in \mathcal{B}$，并设 $\mathcal{F}$ 为 $\mathcal{B}$ 上的集合层。
映射
$$
\mathcal{F}(U) \to \prod\nolimits_{x \in U} \mathcal{F}_x
$$
将 $\mathcal{F}(U)$ 与满足下述性质的元素 $(s_x)_{x\in U}$ 等同：
\begin{itemize}
\item[(*)] 对任意 $x \in U$，存在 $V \in \mathcal{B}$，满足
$x \in V \subset U$，以及截面 $\sigma \in \mathcal{F}(V)$，使得对所有
$y \in V$，在 $\mathcal{F}_y$ 中都有 $s_y = (V, \sigma)$。
\end{itemize}
\end{lemma}

\begin{proof}
首先，由定义 \ref{sheaves-definition-sheaf-basis} 的层条件中的唯一性，映射
$\mathcal{F}(U) \to \prod\nolimits_{x \in U} \mathcal{F}_x$ 是单射。
设 $(s_x)$ 为右侧满足 $(*)$ 的任意元素。显然，这意味着可以找到覆盖
$U = \bigcup U_i$，$U_i \in \mathcal{B}$，使得
$(s_x)_{x \in U_i}$ 来自某个 $\sigma_i \in \mathcal{F}(U_i)$。
对每个 $y \in U_i \cap U_j$，截面 $\sigma_i$ 与 $\sigma_j$ 在茎
$\mathcal{F}_y$ 中一致。因此，存在元素
$V_{ijy} \in \mathcal{B}$，$y \in V_{ijy}$，使得
$\sigma_i|_{V_{ijy}} = \sigma_j|_{V_{ijy}}$。于是，定义
\ref{sheaves-definition-sheaf-basis} 的层条件 $(**)$ 可应用于系统
$\sigma_i$，从而得到具有所需性质的截面 $s \in \mathcal{F}(U)$。
\end{proof}

\noindent
设 $X$ 为拓扑空间，$\mathcal{B}$ 为 $X$ 的拓扑基。从 $X$ 上的集合层
范畴到 $\mathcal{B}$ 上的集合层范畴，有自然的限制函子。事实上，它是
范畴等价。用具体的语言来说，这意味着如下结果。

\begin{lemma}
\label{sheaves-lemma-extend-off-basis}
设 $X$ 为拓扑空间，$\mathcal{B}$ 为 $X$ 的拓扑基，并设
$\mathcal{F}$ 为 $\mathcal{B}$ 上的集合层。存在唯一的 $X$ 上集合层
$\mathcal{F}^{ext}$，使得对所有 $U \in \mathcal{B}$，有
$\mathcal{F}^{ext}(U) = \mathcal{F}(U)$，且与限制映射相容。
\end{lemma}

\begin{proof}
先构造具有所需性质的预层 $\mathcal{F}^{ext}$。具体地，对任意开集
$U \subset X$，把 $\mathcal{F}^{ext}(U)$ 定义为满足引理
\ref{sheaves-lemma-condition-star-sections} 的 $(*)$ 的元素
$(s_x)_{x \in U}$ 所成的集合。显然，存在限制映射使
$\mathcal{F}^{ext}$ 成为集合预层。又由引理
\ref{sheaves-lemma-condition-star-sections}，每当 $U$ 是基 $\mathcal{B}$ 的元素时，
都有 $\mathcal{F}(U) = \mathcal{F}^{ext}(U)$。要说明
$\mathcal{F}^{ext}$ 是层，可仿照引理
\ref{sheaves-lemma-sheafification-sheaf} 的证明论证。
\end{proof}

\noindent
注意，在该引理的情形下有
$$
\mathcal{F}_x = \mathcal{F}_x^{ext}
$$
这是因为 $\mathcal{B}$ 中包含 $x$ 的元素构成 $x$ 的开邻域基本系。

\begin{lemma}
\label{sheaves-lemma-restrict-basis-equivalence}
设 $X$ 为拓扑空间，$\mathcal{B}$ 为 $X$ 的拓扑基。记
$\Sh(\mathcal{B})$ 为 $\mathcal{B}$ 上的层范畴。存在范畴等价
$$
\Sh(X) \longrightarrow \Sh(\mathcal{B})
$$
它把 $X$ 上的层送到其在 $\mathcal{B}$ 各成员上的限制。
\end{lemma}

\begin{proof}
逆函子由上面的引理 \ref{sheaves-lemma-extend-off-basis} 给出。验证显然的函子性
留给读者。
\end{proof}

\noindent
至此结束对基 $\mathcal{B}$ 上集合层的讨论。设 $(\mathcal{C}, F)$
为一种代数结构类型。本节末尾将指出，上述构造对取值于 $\mathcal{C}$
的层也成立。先简要定义相关概念。

\begin{definition}
\label{sheaves-definition-sheaf-structures-basis}
设 $X$ 为拓扑空间，$\mathcal{B}$ 为 $X$ 的拓扑基，并设
$(\mathcal{C}, F)$ 为一种代数结构类型。
\begin{enumerate}
\item {\it $\mathcal{B}$ 上取值于 $\mathcal{C}$ 的预层
$\mathcal{F}$}是如下规则：对每个 $U \in \mathcal{B}$ 指定
$\mathcal{C}$ 的对象 $\mathcal{F}(U)$；对 $\mathcal{B}$ 中元素的每个
包含关系 $V \subset U$，指定 $\mathcal{C}$ 中的态射
$\rho^U_V : \mathcal{F}(U) \to \mathcal{F}(V)$；并且对所有
$U \in \mathcal{B}$ 都有
$\rho^U_U = \text{id}_{\mathcal{F}(U)}$，每当 $\mathcal{B}$ 中有
$W \subset V \subset U$ 时，都有
$\rho^U_W = \rho^V_W \circ \rho ^U_V$。
\item {\it $\mathcal{B}$ 上取值于 $\mathcal{C}$ 的预层态射
$\varphi : \mathcal{F} \to \mathcal{G}$}是如下规则：对每个元素
$U \in \mathcal{B}$ 指定一个与限制映射相容的代数结构态射
% P01-ERRATA: P01-E0049 -- second locus of the component-subscript omission.
$\varphi : \mathcal{F}(U) \to \mathcal{G}(U)$。
\item 给定 $\mathcal{B}$ 上取值于 $\mathcal{C}$ 的预层
$\mathcal{F}$，称 $U \mapsto F(\mathcal{F}(U))$ 为底层集合预层。
\item {\it $\mathcal{B}$ 上取值于 $\mathcal{C}$ 的层
$\mathcal{F}$}是 $\mathcal{B}$ 上取值于 $\mathcal{C}$、且其底层
集合预层为层的预层。
\end{enumerate}
\end{definition}

\noindent
此时，可把 $\mathcal{B}$ 上取值于 $\mathcal{C}$ 的预层在
$x \in X$ 处的{\it 茎}定义为有向余极限
$$
\mathcal{F}_x = \colim_{U\in \mathcal{B}, x\in U} \mathcal{F}(U).
$$
由对 $\mathcal{C}$ 的假设，它作为 $\mathcal{C}$ 的对象存在。
此外，$\mathcal{F}_x$ 的底层集合正是 $\mathcal{B}$ 上底层集合预层的茎。

\medskip\noindent
注意，引理 \ref{sheaves-lemma-cofinal-systems-coverings}、
\ref{sheaves-lemma-cofinal-systems-coverings-standard-case} 和
\ref{sheaves-lemma-condition-star-sections} 涉及的层性质，是用相伴集合预层定义的。
因此，它们原封不动地推广到取值于 $\mathcal{C}$ 的预层。引理
% P01-ERRATA: P01-E0050 -- the frozen source has subject-verb disagreement: "analogue ... need".
\ref{sheaves-lemma-extend-off-basis} 的类似结果需要稍加谨慎。陈述如下。

\begin{lemma}
\label{sheaves-lemma-extend-off-basis-structures}
设 $X$ 为拓扑空间，$(\mathcal{C}, F)$ 为一种代数结构类型，
$\mathcal{B}$ 为 $X$ 的拓扑基，并设 $\mathcal{F}$ 为
$\mathcal{B}$ 上取值于 $\mathcal{C}$ 的层。存在唯一的 $X$ 上取值于
$\mathcal{C}$ 的层 $\mathcal{F}^{ext}$，使得对所有
$U \in \mathcal{B}$，有 $\mathcal{F}^{ext}(U) = \mathcal{F}(U)$，
且与限制映射相容。
\end{lemma}

\begin{proof}
由对二元组 $(\mathcal{C}, F)$ 施加的条件，只需找出一个预层
$\mathcal{F}^{ext}$，它在底层集合预层的层面上完成正确的构造。因此，
首要任务是对每个开集 $U \subset X$ 构造合适的对象
$\mathcal{F}^{ext}(U)$。可以在 $\mathcal{B}$ 上预层的背景下仿照引理
\ref{sheaves-lemma-diagram-fibre-product} 来做。不过，下面这种略有不同
（但基本等价）的方法也可以：把它定义为有向余极限
$$
\mathcal{F}^{ext}(U)
:=
\colim_\mathcal{U} FIB(\mathcal{U})
$$
其中余极限取遍所有由 $U_i \in \mathcal{B}$ 组成的覆盖
$\mathcal{U} : U = \bigcup_{i\in I} U_i$，相应对象由下列纤维积图表定义
$$
\xymatrix{
FIB(\mathcal{U}) \ar[r] \ar[d] &
\prod\nolimits_{x\in U} \mathcal{F}_x \ar[d] \\
\prod\nolimits_{i\in I} \mathcal{F}(U_i) \ar[r] &
\prod\nolimits_{i \in I} \prod\nolimits_{x\in U_i} \mathcal{F}_x
}
$$
由通常的论证，见引理 \ref{sheaves-lemma-image-contained-in} 和例
\ref{sheaves-example-application-lemma-image-contained-in}，只需说明此构造在底层
集合上的结果与上面用 $(**)$ 所作的定义相同。细节留给读者。
\end{proof}

\noindent
注意，在该引理的情形下，作为 $\mathcal{C}$ 中的对象，有
$$
\mathcal{F}_x = \mathcal{F}_x^{ext}
$$
这是因为 $\mathcal{B}$ 中包含 $x$ 的元素构成 $x$ 的开邻域基本系。

\begin{lemma}
\label{sheaves-lemma-restrict-basis-equivalence-structures}
设 $X$ 为拓扑空间，$\mathcal{B}$ 为 $X$ 的拓扑基，并设
$(\mathcal{C}, F)$ 为一种代数结构类型。记
$\Sh(\mathcal{B}, \mathcal{C})$ 为 $\mathcal{B}$ 上取值于
$\mathcal{C}$ 的层范畴。存在范畴等价
$$
\Sh(X, \mathcal{C})
\longrightarrow
\Sh(\mathcal{B}, \mathcal{C})
$$
它把 $X$ 上的层送到其在 $\mathcal{B}$ 各成员上的限制。
\end{lemma}

\begin{proof}
% P01-ERRATA: P01-E0051 -- three frozen inverse-functor proofs say "in given" instead of "is given".
逆函子由上面的引理 \ref{sheaves-lemma-extend-off-basis-structures} 给出。
验证显然的函子性留给读者。
\end{proof}

\noindent
最后讨论基上的模（预）层。我们将使用如下容易的事实：基上的集合预层
范畴有积，而且积通过在基的各元素上逐项取值之积来描述。

\begin{definition}
\label{sheaves-definition-sheaf-modules-basis}
设 $X$ 为拓扑空间，$\mathcal{B}$ 为 $X$ 的拓扑基，并设
$\mathcal{O}$ 为 $\mathcal{B}$ 上的环预层。
\begin{enumerate}
\item {\it $\mathcal{B}$ 上的 $\mathcal{O}$-模预层
$\mathcal{F}$}是 $\mathcal{B}$ 上的阿贝尔群预层，连同一个集合预层态射
$\mathcal{O} \times \mathcal{F} \to \mathcal{F}$，使得对所有
$U \in \mathcal{B}$，映射
$\mathcal{O}(U) \times \mathcal{F}(U) \to \mathcal{F}(U)$ 都把群
$\mathcal{F}(U)$ 变成 $\mathcal{O}(U)$-模。
\item {\it $\mathcal{B}$ 上 $\mathcal{O}$-模预层的态射
$\varphi : \mathcal{F} \to \mathcal{G}$}是 $\mathcal{B}$ 上的阿贝尔
预层态射，它对每个 $U \in \mathcal{B}$ 都诱导
$\mathcal{O}(U)$-模同态
$\mathcal{F}(U) \to \mathcal{G}(U)$。
\item 假设 $\mathcal{O}$ 是 $\mathcal{B}$ 上的环层。
{\it $\mathcal{B}$ 上的 $\mathcal{O}$-模层 $\mathcal{F}$}是
$\mathcal{B}$ 上的 $\mathcal{O}$-模预层，且其底层阿贝尔群预层是层。
\end{enumerate}
\end{definition}

\noindent
可把 $\mathcal{B}$ 上的 $\mathcal{O}$-模预层在 $x \in X$ 处的
{\it 茎}定义为有向余极限
$$
\mathcal{F}_x = \colim_{U\in \mathcal{B}, x\in U} \mathcal{F}(U).
$$
它是 $\mathcal{O}_x$-模。

\medskip\noindent
注意，引理 \ref{sheaves-lemma-cofinal-systems-coverings}、
\ref{sheaves-lemma-cofinal-systems-coverings-standard-case} 和
\ref{sheaves-lemma-condition-star-sections} 涉及的层性质，是用相伴集合预层定义的。
因此，它们原封不动地推广到 $\mathcal{O}$-模预层。引理
\ref{sheaves-lemma-extend-off-basis} 的类似结果如下。

\begin{lemma}
\label{sheaves-lemma-extend-off-basis-module}
设 $X$ 为拓扑空间，$\mathcal{B}$ 为 $X$ 的拓扑基，
$\mathcal{O}$ 为 $\mathcal{B}$ 上的环层，并设 $\mathcal{F}$ 为
$\mathcal{B}$ 上的 $\mathcal{O}$-模层。设 $\mathcal{O}^{ext}$ 为
扩张 $\mathcal{O}$ 的 $X$ 上环层，并设 $\mathcal{F}^{ext}$ 为扩张
$\mathcal{F}$ 的 $X$ 上阿贝尔层；见引理
\ref{sheaves-lemma-extend-off-basis-structures}。存在典范映射
$$
\mathcal{O}^{ext} \times \mathcal{F}^{ext}
\longrightarrow
\mathcal{F}^{ext}
$$
它在 $\mathcal{B}$ 的各元素上与给定映射相同，并赋予
$\mathcal{F}^{ext}$ 一个 $\mathcal{O}^{ext}$-模结构。
\end{lemma}

\begin{proof}
只需在集合预层层面构造乘法映射。看出这一点也许最容易的方式，是直接证明：
若 $(f_x)_{x \in U}$、$f_x \in \mathcal{O}_x$ 与
$(m_x)_{x \in U}$、$m_x \in \mathcal{F}_x$ 都满足 $(*)$，则元素
$(f_xm_x)_{x \in U}$ 也满足 $(*)$。于是得到所需结论，因为在引理
\ref{sheaves-lemma-extend-off-basis} 的证明中，我们用满足 $(*)$ 的茎元素族来
构造扩张。
\end{proof}

\noindent
注意，在该引理的情形下，作为 $\mathcal{O}_x$-模，有
$$
\mathcal{F}_x = \mathcal{F}_x^{ext}
$$
这是因为 $\mathcal{B}$ 中包含 $x$ 的元素构成 $x$ 的开邻域基本系；
也可以简单地说，这是因为底层集合上的结论成立。

\begin{lemma}
\label{sheaves-lemma-restrict-basis-equivalence-modules}
设 $X$ 为拓扑空间，$\mathcal{B}$ 为 $X$ 的拓扑基，并设
$\mathcal{O}$ 为 $X$ 上的环层。记
$\textit{Mod}(\mathcal{O}|_\mathcal{B})$ 为 $\mathcal{B}$ 上
$\mathcal{O}|_\mathcal{B}$-模层的范畴。存在范畴等价
$$
\textit{Mod}(\mathcal{O})
\longrightarrow
\textit{Mod}(\mathcal{O}|_\mathcal{B})
$$
它把 $X$ 上的 $\mathcal{O}$-模层送到其在 $\mathcal{B}$ 各成员上的限制。
\end{lemma}

\begin{proof}
% P01-ERRATA: P01-E0051 -- another locus of "in given".
逆函子由上面的引理 \ref{sheaves-lemma-extend-off-basis-module} 给出。
验证显然的函子性留给读者。
\end{proof}

\noindent
最后讨论这一理论与连续映射的关系。由于上面的工作，这现在很容易。
先处理在靶上给定一个基的情形。

\begin{lemma}
\label{sheaves-lemma-f-map-basis-below-structures}
设 $f : X \to Y$ 为拓扑空间的连续映射，$(\mathcal{C}, F)$ 为一种
代数结构类型，$\mathcal{F}$ 为 $X$ 上取值于 $\mathcal{C}$ 的层，
$\mathcal{G}$ 为 $Y$ 上取值于 $\mathcal{C}$ 的层，并设
$\mathcal{B}$ 为 $Y$ 的拓扑基。假设对每个 $V \in \mathcal{B}$ 给定
$\mathcal{C}$ 中的态射
$$
\varphi_V :
\mathcal{G}(V)
\longrightarrow
\mathcal{F}(f^{-1}V)
$$
并且它们与限制映射相容。则存在唯一的 $f$-映射（见定义
\ref{sheaves-definition-f-map} 及第
\ref{sheaves-section-presheaves-structures-functorial} 节中对 $f$-映射的讨论）
$\varphi : \mathcal{G} \to \mathcal{F}$，它对
$V \in \mathcal{B}$ 恢复 $\varphi_V$。
\end{lemma}

\begin{proof}
这是平凡的，因为这族映射等同于 $\mathcal{G}$ 与
$f_*\mathcal{F}$ 在 $\mathcal{B}$ 上的限制之间的态射。由引理
\ref{sheaves-lemma-restrict-basis-equivalence-structures}，这等同于给定从
$\mathcal{G}$ 到 $f_*\mathcal{F}$ 的态射；由引理 \ref{sheaves-lemma-f-map}，
又等同于 $f$-映射。关于如何处理代数结构层的情形，另见引理
\ref{sheaves-lemma-f-map-sets-algebraic-structures} 及其前面的讨论。
\end{proof}

\noindent
下面是赋环空间的类似结果。

\begin{lemma}
\label{sheaves-lemma-f-map-basis-below-modules}
设 $(f, f^\sharp) : (X, \mathcal{O}_X) \to (Y, \mathcal{O}_Y)$
为赋环空间态射，$\mathcal{F}$ 为 $\mathcal{O}_X$-模层，
$\mathcal{G}$ 为 $\mathcal{O}_Y$-模层，并设 $\mathcal{B}$ 为 $Y$
的拓扑基。假设对每个 $V \in \mathcal{B}$ 给定
$\mathcal{O}_Y(V)$-模映射
$$
\varphi_V :
\mathcal{G}(V)
\longrightarrow
\mathcal{F}(f^{-1}V)
$$
（其中 $\mathcal{F}(f^{-1}V)$ 通过
$f^\sharp_V : \mathcal{O}_Y(V) \to \mathcal{O}_X(f^{-1}V)$ 获得模结构），
并且它们与限制映射相容。则存在唯一的 $f$-映射（见第
\ref{sheaves-section-ringed-spaces-functoriality-modules} 节中对 $f$-映射的讨论）
$\varphi : \mathcal{G} \to \mathcal{F}$，它对
$V \in \mathcal{B}$ 恢复 $\varphi_V$。
\end{lemma}

\begin{proof}
与上面代数结构层相应引理的证明相同。
\end{proof}

\begin{lemma}
\label{sheaves-lemma-f-map-basis-above-and-below-structures}
设 $f : X \to Y$ 为拓扑空间的连续映射，$(\mathcal{C}, F)$ 为一种
代数结构类型，$\mathcal{F}$ 为 $X$ 上取值于 $\mathcal{C}$ 的层，
$\mathcal{G}$ 为 $Y$ 上取值于 $\mathcal{C}$ 的层。设
$\mathcal{B}_Y$ 为 $Y$ 的拓扑基，$\mathcal{B}_X$ 为 $X$ 的拓扑基。
假设对每个 $V \in \mathcal{B}_Y$ 和每个满足
$f(U) \subset V$ 的 $U \in \mathcal{B}_X$，给定 $\mathcal{C}$ 中的态射
$$
\varphi_V^U :
\mathcal{G}(V)
\longrightarrow
\mathcal{F}(U)
$$
并且它们与限制映射相容。则存在唯一的 $f$-映射（见定义
\ref{sheaves-definition-f-map} 及第
\ref{sheaves-section-presheaves-structures-functorial} 节中对 $f$-映射的讨论）
$\varphi : \mathcal{G} \to \mathcal{F}$，使得对上述每一对 $(U, V)$，
$\varphi_V^U$ 恢复为复合
$$
\mathcal{G}(V) \xrightarrow{\varphi_V}
\mathcal{F}(f^{-1}(V)) \xrightarrow{\text{restr.}}
\mathcal{F}(U)
$$
\end{lemma}

\begin{proof}
% P01-ERRATA: P01-E0052 -- the frozen proof begins "Let us first proves" and later says "the maps on stalks is".
先对集合层证明这一点。固定开集 $V \subset Y$，并取
$s \in \mathcal{G}(V)$。我们将构造元素
$\varphi_V(s) \in \mathcal{F}(f^{-1}V)$。对每个
$x \in f^{-1}V$，可在茎 $\mathcal{F}_x$ 中定义一个值
$\varphi(s)_x$：选取满足 $x \in U \subset f^{-1}V$ 的
$U \in \mathcal{B}_X$，并令 $\varphi(s)_x$ 等于茎中
$(U, \varphi_V^U(s))$ 的等价类。显然，族
$(\varphi(s)_x)_{x \in f^{-1}V}$ 满足条件 $(*)$，因为当 $U$ 变化时，
各映射 $\varphi_V^U$ 与层 $\mathcal{F}$ 中的限制相容。因此，由引理
\ref{sheaves-lemma-extend-off-basis} 的证明，
$(\varphi(s)_x)_{x \in f^{-1}V}$ 对应于
$\mathcal{F}(f^{-1}V)$ 的唯一元素 $\varphi_V(s)$。这样便定义了集合映射
$\varphi_V : \mathcal{G}(V) \to \mathcal{F}(f^{-1}V)$。
$\varphi_V$ 与 $\varphi_V^U$ 之间的相容性由引理
\ref{sheaves-lemma-condition-star-sections} 得出。

\medskip\noindent
我们把如下验证留给读者：当 $V \in \mathcal{B}_Y$ 变化时，
$\varphi_V$ 的构造与限制映射相容。因此，可应用上面的引理
\ref{sheaves-lemma-f-map-basis-below-structures}，将它们“粘合”为所需的
$f$-映射。

\medskip\noindent
最后，注意这样构造的集合层映射具有如下性质：茎映射
$$
\mathcal{G}_{f(x)} \longrightarrow \mathcal{F}_x
$$
是如下映射系统的余极限：$V \in \mathcal{B}_Y$ 取遍包含 $f(x)$ 的
元素，$U \in \mathcal{B}_X$ 取遍包含 $x$ 的元素，而映射为
$\varphi_V^U$。特别地，若 $\mathcal{G}$ 与 $\mathcal{F}$ 是代数结构层
的底层集合层，则可见各茎上的映射是代数结构的态射。因此，相伴的底层
集合层映射 $f^{-1}\mathcal{G} \to \mathcal{F}$ 满足引理
\ref{sheaves-lemma-f-map-sets-algebraic-structures} 的假设。由此得出
$f^{-1}\mathcal{G} \to \mathcal{F}$ 是取值于 $\mathcal{C}$ 的层态射。
而由伴随性，这意味着 $\varphi$ 是代数结构层的 $f$-映射。
\end{proof}

\begin{lemma}
\label{sheaves-lemma-f-map-basis-above-and-below-modules}
设 $(f, f^\sharp) : (X, \mathcal{O}_X) \to (Y, \mathcal{O}_Y)$
为赋环空间态射，$\mathcal{F}$ 为 $\mathcal{O}_X$-模层，
$\mathcal{G}$ 为 $\mathcal{O}_Y$-模层。设 $\mathcal{B}_Y$ 为 $Y$
的拓扑基，$\mathcal{B}_X$ 为 $X$ 的拓扑基。假设对每个
$V \in \mathcal{B}_Y$ 和每个满足 $f(U) \subset V$ 的
$U \in \mathcal{B}_X$，给定 $\mathcal{O}_Y(V)$-模映射
$$
\varphi_V^U :
\mathcal{G}(V)
\longrightarrow
\mathcal{F}(U)
$$
并且它们与限制映射相容。这里，$\mathcal{F}(U)$ 上的
$\mathcal{O}_Y(V)$-模结构来自 $\mathcal{O}_X(U)$-模结构以及映射
$f^\sharp_V : \mathcal{O}_Y(V)
\to \mathcal{O}_X(f^{-1}V) \to \mathcal{O}_X(U)$。
则存在唯一的模层 $f$-映射（见定义 \ref{sheaves-definition-f-map} 及第
\ref{sheaves-section-ringed-spaces-functoriality-modules} 节中对 $f$-映射的讨论）
$\varphi : \mathcal{G} \to \mathcal{F}$，使得对上述每一对 $(U, V)$，
$\varphi_V^U$ 恢复为复合
$$
\mathcal{G}(V) \xrightarrow{\varphi_V}
\mathcal{F}(f^{-1}(V)) \xrightarrow{\text{restr.}}
\mathcal{F}(U)
$$
\end{lemma}

\begin{proof}
与上面的证明类似，略去。
\end{proof}

\section{开浸入与（预）层}
\label{sheaves-section-open-immersions}

\noindent
设 $X$ 为拓扑空间，并设 $j : U \to X$ 为开子集 $U$ 到 $X$ 的包含。
在第 \ref{sheaves-section-presheaves-functorial} 节中，我们已经定义函子
$j_*$ 与 $j^{-1}$，使得 $j_*$ 是 $j^{-1}$ 的右伴随。对开浸入，
$j^{-1}$ 还有一个左伴随，记作 $j_!$。先指出，在开浸入情形下，
$j^{-1}$ 有特别简单的描述。

\begin{lemma}
\label{sheaves-lemma-j-pullback}
设 $X$ 为拓扑空间，并设 $j : U \to X$ 为开子集 $U$ 到 $X$ 的包含。
\begin{enumerate}
\item 设 $\mathcal{G}$ 为 $X$ 上的集合预层。预层
$j_p\mathcal{G}$（见第 \ref{sheaves-section-presheaves-functorial} 节）由规则
$V \mapsto \mathcal{G}(V)$ 给出，其中 $V \subset U$ 为开集。
\item 设 $\mathcal{G}$ 为 $X$ 上的集合层。层 $j^{-1}\mathcal{G}$
由规则 $V \mapsto \mathcal{G}(V)$ 给出，其中 $V \subset U$ 为开集。
\item 对任意点 $u \in U$ 及 $X$ 上的任意层 $\mathcal{G}$，有茎的
典范等同
$$
j^{-1}\mathcal{G}_u = (\mathcal{G}|_U)_u = \mathcal{G}_u.
$$
\item 在 $U$ 的预层范畴上，有 $j_pj_* = \text{id}$。
\item 在 $U$ 的层范畴上，有 $j^{-1}j_* = \text{id}$。
\end{enumerate}
同样的描述对阿贝尔群（预）层、代数结构（预）层及模（预）层均成立。
\end{lemma}

\begin{proof}
$j_p\mathcal{G}(V)$ 定义中的余极限取遍所有满足 $V \subset W$ 的
$X$ 中开集 $W \subset X$，次序取反向包含。因此，它有最大元，即 $V$。
这证明了 (1)。而 (2) 成立，是因为若 $\mathcal{G}$ 是层，则对
$V \subset U$ 为开集的赋值 $V \mapsto \mathcal{G}(V)$ 显然也是层。
断言 (3) 由 (2) 得出，因为包含于 $U$ 的 $u$ 的开邻域在 $X$ 中
$u$ 的所有开邻域中共尾。(4) 与 (5) 由计算
$j^{-1}j_*\mathcal{F}(V) = j_*\mathcal{F}(V) = \mathcal{F}(V)$ 得出。

\medskip\noindent
完全相同的论证适用于阿贝尔群（预）层和代数结构（预）层。
\end{proof}

\begin{definition}
\label{sheaves-definition-restriction}
设 $X$ 为拓扑空间，并设 $j : U \to X$ 为开子集的包含。
\begin{enumerate}
\item 设 $\mathcal{G}$ 为 $X$ 上的集合、阿贝尔群或代数结构预层。
引理 \ref{sheaves-lemma-j-pullback} 中描述的预层 $j_p\mathcal{G}$ 称为
{$\mathcal{G}$ 在 $U$ 上的\it 限制}，记作 $\mathcal{G}|_U$。
\item 设 $\mathcal{G}$ 为 $X$ 上的集合层、阿贝尔群层或代数结构层。
层 $j^{-1}\mathcal{G}$ 称为{$\mathcal{G}$ 在 $U$ 上的\it 限制}，
记作 $\mathcal{G}|_U$。
\item 若 $(X, \mathcal{O})$ 是赋环空间，则二元组
$(U, \mathcal{O}|_U)$ 称为{\it 与 $U$ 相伴的
$(X, \mathcal{O})$ 的开子空间}。
\item 若 $\mathcal{G}$ 是 $\mathcal{O}$-模预层，则
$\mathcal{G}|_U$ 连同乘法映射
$\mathcal{O}|_U \times \mathcal{G}|_U \to \mathcal{G}|_U$
（见引理 \ref{sheaves-lemma-pullback-module}）称为
{$\mathcal{G}$ 在 $U$ 上的\it 限制}。
\end{enumerate}
\end{definition}

\noindent
% P01-ERRATA: P01-E0053 -- item (4) defines restriction for module presheaves, then the next sentence says that definition is left to the reader.
我们把模预层的限制之定义留给读者。好了，本节将讨论限制函子的一个左伴随。
先给出集合（预）层情形下的定义。

\begin{definition}
\label{sheaves-definition-j-shriek}
设 $X$ 为拓扑空间，并设 $j : U \to X$ 为开子集的包含。
\begin{enumerate}
\item 设 $\mathcal{F}$ 为 $U$ 上的集合预层。定义
{\it $\mathcal{F}$ 的以空集延拓 $j_{p!}\mathcal{F}$}为 $X$ 上由规则
$$
j_{p!}\mathcal{F}(V) =
\left\{
\begin{matrix}
\emptyset & \text{若} & V \not \subset U \\
\mathcal{F}(V) & \text{若} & V \subset U
\end{matrix}
\right.
$$
定义的集合预层，其限制映射是显然的。
\item 设 $\mathcal{F}$ 为 $U$ 上的集合层。定义
{\it $\mathcal{F}$ 的以空集延拓 $j_!\mathcal{F}$}为预层
$j_{p!}\mathcal{F}$ 的层化。
\end{enumerate}
\end{definition}

\begin{lemma}
\label{sheaves-lemma-j-shriek}
设 $X$ 为拓扑空间，并设 $j : U \to X$ 为开子集的包含。
\begin{enumerate}
\item 函子 $j_{p!}$ 是限制函子 $j_p$ 的左伴随（见引理
\ref{sheaves-lemma-j-pullback}）。
\item 函子 $j_!$ 是限制的左伴随；公式为
$$
\Mor_{\Sh(X)}(j_!\mathcal{F}, \mathcal{G})
=
\Mor_{\Sh(U)}(\mathcal{F}, j^{-1}\mathcal{G})
=
\Mor_{\Sh(U)}(\mathcal{F}, \mathcal{G}|_U)
$$
它对 $\mathcal{F}$ 和 $\mathcal{G}$ 双函子地成立。
\item 设 $\mathcal{F}$ 为 $U$ 上的集合层。层 $j_!\mathcal{F}$ 的茎
描述如下
$$
j_{!}\mathcal{F}_x =
\left\{
\begin{matrix}
\emptyset & \text{若} & x \not \in U \\
\mathcal{F}_x & \text{若} & x \in U
\end{matrix}
\right.
$$
\item 在 $U$ 的预层范畴上，有 $j_pj_{p!} = \text{id}$。
\item 在 $U$ 的层范畴上，有 $j^{-1}j_! = \text{id}$。
\end{enumerate}
\end{lemma}

\begin{proof}
要把 $j_{p!}\mathcal{F}$ 映入 $\mathcal{G}$，只需在每当
$V \subset U$ 时给出与限制映射相容的映射
$\mathcal{F}(V) \to \mathcal{G}(V)$。而由引理
\ref{sheaves-lemma-j-pullback}，映射 $\mathcal{F} \to \mathcal{G}|_U$
也有同样的描述。$j_!$ 与限制的伴随性由此及层化的性质得出。
茎的等同由以空集延拓的定义和茎的定义显然可得。(4) 与 (5) 由计算层在
$U$ 的任意开集上的取值得出。
\end{proof}

\noindent
注意，若 $\mathcal{F}$ 是 $U$ 上的阿贝尔群层，则一般而言，按上述定义的
$j_!\mathcal{F}$ 不是阿贝尔群层，例如因为它的一些茎为空（因而当然不是
阿贝尔群）。所以，需要依据所考虑层的类型修改 $j_!$ 的定义。在以空集
延拓的定义中选择空集，是因为它是集合范畴的始对象。因此，在阿贝尔群
情形下使用 $0$（更一般地，对取值于任意阿贝尔范畴的层也如此）。

\begin{definition}
\label{sheaves-definition-j-shriek-structures}
设 $X$ 为拓扑空间，并设 $j : U \to X$ 为开子集的包含。
\begin{enumerate}
\item 设 $\mathcal{F}$ 为 $U$ 上的阿贝尔预层。定义
{\it $\mathcal{F}$ 的以 $0$ 延拓 $j_{p!}\mathcal{F}$}为 $X$ 上由规则
$$
j_{p!}\mathcal{F}(V) =
\left\{
\begin{matrix}
0 & \text{若} & V \not \subset U \\
\mathcal{F}(V) & \text{若} & V \subset U
\end{matrix}
\right.
$$
定义的阿贝尔预层，其限制映射是显然的。
\item 设 $\mathcal{F}$ 为 $U$ 上的阿贝尔层。定义
{\it $\mathcal{F}$ 的以 $0$ 延拓 $j_!\mathcal{F}$}为阿贝尔预层
$j_{p!}\mathcal{F}$ 的层化。
\item 设 $\mathcal{C}$ 为具有始对象 $e$ 的范畴，$\mathcal{F}$ 为
$U$ 上取值于 $\mathcal{C}$ 的预层。定义
{\it $\mathcal{F}$ 的以 $e$ 延拓 $j_{p!}\mathcal{F}$}为 $X$ 上取值于
$\mathcal{C}$、且由规则
$$
j_{p!}\mathcal{F}(V) =
\left\{
\begin{matrix}
e & \text{若} & V \not \subset U \\
\mathcal{F}(V) & \text{若} & V \subset U
\end{matrix}
\right.
$$
定义的预层，其限制映射是显然的。
\item 设 $(\mathcal{C}, F)$ 为一种代数结构类型，且 $\mathcal{C}$
有始对象 $e$。设 $\mathcal{F}$ 为 $U$ 上（给定类型的）代数结构层。
定义{\it $\mathcal{F}$ 的以 $e$ 延拓 $j_!\mathcal{F}$}为上述预层
$j_{p!}\mathcal{F}$ 的层化。
\item 设 $\mathcal{O}$ 为 $X$ 上的环预层，$\mathcal{F}$ 为
$\mathcal{O}|_U$-模预层。在这种情形下，把{\it 以 $0$ 延拓}定义为如下
$\mathcal{O}$-模预层：作为阿贝尔预层，它等于 $j_{p!}\mathcal{F}$，
并配备乘法映射
$\mathcal{O} \times j_{p!}\mathcal{F} \to j_{p!}\mathcal{F}$。
\item 设 $\mathcal{O}$ 为 $X$ 上的环层，$\mathcal{F}$ 为
$\mathcal{O}|_U$-模层。在这种情形下，把{\it 以 $0$ 延拓}定义为如下
$\mathcal{O}$-模：作为阿贝尔层，它等于 $j_!\mathcal{F}$，并配备乘法
映射 $\mathcal{O} \times j_!\mathcal{F} \to j_!\mathcal{F}$。
\end{enumerate}
\end{definition}

\noindent
确实可以在代数结构层的背景下定义 $j_!$（见下文）。不过，所得对象取决于
所涉及的代数结构类型。例如，若 $\mathcal{O}$ 是 $U$ 上的环层，则
$j_{!, rings}\mathcal{O} \not = j_{!, abelian}\mathcal{O}$，因为环范畴的
始对象是 $\mathbf{Z}$，而阿贝尔群范畴的始对象是 $0$。特别地，与迄今
所见的情形不同，函子 $j_!$ {\it 不与取底层集合层可交换}！像往常一样，
分别处理阿贝尔群（预）层、代数结构（预）层和模（预）层的情形。

\begin{lemma}
\label{sheaves-lemma-j-shriek-abelian}
设 $X$ 为拓扑空间，并设 $j : U \to X$ 为开子集的包含。考虑阿贝尔
（预）层的限制函子与以 $0$ 延拓函子。
\begin{enumerate}
\item 函子 $j_{p!}$ 是限制函子 $j_p$ 的左伴随（见引理
\ref{sheaves-lemma-j-pullback}）。
\item 函子 $j_!$ 是限制的左伴随；公式为
$$
\Mor_{\textit{Ab}(X)}(j_!\mathcal{F}, \mathcal{G})
=
\Mor_{\textit{Ab}(U)}(\mathcal{F}, j^{-1}\mathcal{G})
=
\Mor_{\textit{Ab}(U)}(\mathcal{F}, \mathcal{G}|_U)
$$
它对 $\mathcal{F}$ 和 $\mathcal{G}$ 双函子地成立。
\item 设 $\mathcal{F}$ 为 $U$ 上的阿贝尔层。层 $j_!\mathcal{F}$ 的茎
描述如下
$$
j_{!}\mathcal{F}_x =
\left\{
\begin{matrix}
0 & \text{若} & x \not \in U \\
\mathcal{F}_x & \text{若} & x \in U
\end{matrix}
\right.
$$
\item 在 $U$ 的阿贝尔预层范畴上，有 $j_pj_{p!} = \text{id}$。
\item 在 $U$ 的阿贝尔层范畴上，有 $j^{-1}j_! = \text{id}$。
\end{enumerate}
\end{lemma}

\begin{proof}
略。
\end{proof}

\begin{lemma}
\label{sheaves-lemma-j-shriek-structures}
设 $X$ 为拓扑空间，并设 $j : U \to X$ 为开子集的包含。设
$(\mathcal{C}, F)$ 为一种代数结构类型，且 $\mathcal{C}$ 有始对象 $e$。
考虑上面定义的代数结构（预）层的限制函子与以 $e$ 延拓函子。
\begin{enumerate}
\item 函子 $j_{p!}$ 是限制函子 $j_p$ 的左伴随（见引理
\ref{sheaves-lemma-j-pullback}）。
\item 函子 $j_!$ 是限制的左伴随；公式为
$$
\Mor_{\Sh(X, \mathcal{C})}(j_!\mathcal{F}, \mathcal{G})
=
\Mor_{\Sh(U, \mathcal{C})}(\mathcal{F}, j^{-1}\mathcal{G})
=
\Mor_{\Sh(U, \mathcal{C})}(\mathcal{F}, \mathcal{G}|_U)
$$
它对 $\mathcal{F}$ 和 $\mathcal{G}$ 双函子地成立。
\item 设 $\mathcal{F}$ 为 $U$ 上的层。层 $j_!\mathcal{F}$ 的茎描述如下
$$
j_{!}\mathcal{F}_x =
\left\{
\begin{matrix}
e & \text{若} & x \not \in U \\
\mathcal{F}_x & \text{若} & x \in U
\end{matrix}
\right.
$$
\item 在 $U$ 上的代数结构预层范畴中，有
$j_pj_{p!} = \text{id}$。
\item 在 $U$ 上的代数结构层范畴中，有 $j^{-1}j_! = \text{id}$。
\end{enumerate}
\end{lemma}

\begin{proof}
略。
\end{proof}

\begin{lemma}
\label{sheaves-lemma-j-shriek-modules}
设 $(X, \mathcal{O})$ 为赋环空间，并设
$j : (U, \mathcal{O}|_U) \to (X, \mathcal{O})$ 为开子空间。
考虑上面定义的模（预）层的限制函子与以 $0$ 延拓函子。
\begin{enumerate}
\item 函子 $j_{p!}$ 是限制的左伴随；公式为
$$
\Mor_{\textit{PMod}(\mathcal{O})}(j_{p!}\mathcal{F}, \mathcal{G})
=
\Mor_{\textit{PMod}(\mathcal{O}|_U)}(\mathcal{F}, \mathcal{G}|_U)
$$
它对 $\mathcal{F}$ 和 $\mathcal{G}$ 双函子地成立。
\item 函子 $j_!$ 是限制的左伴随；公式为
$$
\Mor_{\textit{Mod}(\mathcal{O})}(j_!\mathcal{F}, \mathcal{G})
=
\Mor_{\textit{Mod}(\mathcal{O}|_U)}(\mathcal{F}, \mathcal{G}|_U)
$$
它对 $\mathcal{F}$ 和 $\mathcal{G}$ 双函子地成立。
\item
% P01-ERRATA: P01-E0054 -- the frozen source calls F an O-module sheaf on U rather than an O|_U-module sheaf.
设 $\mathcal{F}$ 为 $U$ 上的 $\mathcal{O}$-模层。层
$j_!\mathcal{F}$ 的茎描述如下
$$
j_{!}\mathcal{F}_x =
\left\{
\begin{matrix}
0 & \text{若} & x \not \in U \\
\mathcal{F}_x & \text{若} & x \in U
\end{matrix}
\right.
$$
\item 在 $U$ 上的 $\mathcal{O}|_U$-模层范畴中，有
$j^{-1}j_! = \text{id}$。
\end{enumerate}
\end{lemma}

\begin{proof}
略。
\end{proof}

\noindent
注意，由上述各引理，函子 $j_*$ 和 $j_!$ 都是从 $U$ 上的层范畴到
$X$ 上的层范畴的全忠实嵌入。只有对函子 $j_!$，才能容易地描述其本质像。

\begin{lemma}
\label{sheaves-lemma-equivalence-categories-open}
设 $X$ 为拓扑空间，并设 $j : U \to X$ 为开子集的包含。函子
$$
j_! : \Sh(U) \longrightarrow \Sh(X)
$$
是全忠实的。其本质像恰由如下层 $\mathcal{G}$ 组成：对所有
$x \in X \setminus U$，都有 $\mathcal{G}_x = \emptyset$。
\end{lemma}

\begin{proof}
全忠实性由 $j^{-1} j_! = \text{id}$ 形式地得出。已经看到，函子像中的
任意层都具有引理所述的茎性质。反过来，假设 $\mathcal{G}$ 具有指定性质。
那么很容易验证
$$
j_! j^{-1} \mathcal{G} \to \mathcal{G}
$$
在所有茎上都是同构，因而是同构。
\end{proof}

\begin{lemma}
\label{sheaves-lemma-equivalence-categories-open-abelian}
设 $X$ 为拓扑空间，并设 $j : U \to X$ 为开子集的包含。函子
$$
j_! : \textit{Ab}(U) \longrightarrow \textit{Ab}(X)
$$
是全忠实的。其本质像恰由如下层 $\mathcal{G}$ 组成：对所有
$x \in X \setminus U$，都有 $\mathcal{G}_x = 0$。
\end{lemma}

\begin{proof}
略。
\end{proof}

\begin{lemma}
\label{sheaves-lemma-equivalence-categories-open-structures}
设 $X$ 为拓扑空间，并设 $j : U \to X$ 为开子集的包含。设
$(\mathcal{C}, F)$ 为一种代数结构类型，且 $\mathcal{C}$ 有始对象 $e$。
函子
$$
j_! : \Sh(U, \mathcal{C}) \longrightarrow \Sh(X, \mathcal{C})
$$
是全忠实的。其本质像恰由如下层 $\mathcal{G}$ 组成：对所有
$x \in X \setminus U$，都有 $\mathcal{G}_x = e$。
\end{lemma}

\begin{proof}
略。
\end{proof}


\begin{lemma}
\label{sheaves-lemma-equivalence-categories-open-modules}
设 $(X, \mathcal{O})$ 为赋环空间，并设
$j : (U, \mathcal{O}|_U) \to (X, \mathcal{O})$ 为开子空间。函子
$$
j_! : \textit{Mod}(\mathcal{O}|_U) \longrightarrow \textit{Mod}(\mathcal{O})
$$
是全忠实的。其本质像恰由如下层 $\mathcal{G}$ 组成：对所有
$x \in X \setminus U$，都有 $\mathcal{G}_x = 0$。
\end{lemma}

\begin{proof}
略。
\end{proof}

\begin{remark}
\label{sheaves-remark-j-shriek-not-exact}
设 $j : U \to X$ 如上为拓扑空间的开浸入。设 $x \in X$、
$x \not \in U$，并设 $\mathcal{F}$ 为 $U$ 上的集合层。由引理
\ref{sheaves-lemma-j-shriek}，有 $j_!\mathcal{F}_x = \emptyset$。因此，除非
$U = X$，$j_!$ 不会把 $\Sh(U)$ 的终对象变为 $\Sh(X)$ 的终对象。
按照 Categories，第 \ref{categories-section-exact-functor} 节中的约定，
这意味着作为集合层范畴之间的函子，$j_!$ 不是左正合的。稍后将证明
阿贝尔层上的 $j_!$ 是正合的；见 Modules，引理
\ref{modules-lemma-j-shriek-exact}。
\end{remark}







\section{闭浸入与（预）层}
\label{sheaves-section-closed-immersions}

\noindent
设 $X$ 为拓扑空间，并设 $i : Z \to X$ 为闭子集 $Z$ 到 $X$ 的包含。
在第 \ref{sheaves-section-presheaves-functorial} 节中，我们已经定义函子
$i_*$ 与 $i^{-1}$，使得 $i_*$ 是 $i^{-1}$ 的右伴随。

\begin{lemma}
\label{sheaves-lemma-stalks-closed-pushforward}
设 $X$ 为拓扑空间，$i : Z \to X$ 为闭子集 $Z$ 到 $X$ 的包含，并设
$\mathcal{F}$ 为 $Z$ 上的集合层。$i_*\mathcal{F}$ 的茎描述如下
$$
i_*\mathcal{F}_x =
\left\{
\begin{matrix}
\{*\} & \text{若} & x \not \in Z \\
\mathcal{F}_x & \text{若} & x \in Z
\end{matrix}
\right.
$$
其中 $\{*\}$ 表示单元素集合。此外，在 $Z$ 上的集合层范畴中，
$i^{-1}i_* = \text{id}$。对 $Z$ 上的阿贝尔层、相应地对 $Z$ 上的
代数结构层，同一结论也成立；此时必须把 $\{*\}$ 分别换为 $0$、
相应地换为代数结构范畴的终对象。
\end{lemma}

\begin{proof}
若 $x \not \in Z$，则存在任意小且与 $Z$ 不相交的 $x$ 的开邻域 $U$。
由于 $\mathcal{F}$ 是层，对任意这样的 $U$，都有
$\mathcal{F}(i^{-1}(U)) = \{*\}$；见评注 \ref{sheaves-remark-confusion}。
这证明了第一种情形。第二种情形来自如下事实：对 $z \in Z$，$z$ 的任意
开邻域都形如 $Z \cap U$，其中 $U$ 是 $X$ 的某个开集。为证明
$i^{-1}i_* = \text{id}$，考虑典范映射
$i^{-1}i_*\mathcal{F} \to \mathcal{F}$。由上可知，它在各茎上是同构，
因而是同构。

\medskip\noindent
对阿贝尔群层和代数结构层，作同样的论证。
\end{proof}


\begin{lemma}
\label{sheaves-lemma-equivalence-categories-closed}
设 $X$ 为拓扑空间，并设 $i : Z \to X$ 为闭子集的包含。函子
$$
i_* : \Sh(Z) \longrightarrow \Sh(X)
$$
是全忠实的。其本质像恰由如下层 $\mathcal{G}$ 组成：对所有
$x \in X \setminus Z$，都有 $\mathcal{G}_x = \{*\}$。
\end{lemma}

\begin{proof}
全忠实性由 $i^{-1} i_* = \text{id}$ 形式地得出。已经看到，函子像中的
任意层都具有引理所述的茎性质。反过来，假设 $\mathcal{G}$ 具有指定性质。
那么很容易验证
$$
\mathcal{G} \to i_* i^{-1} \mathcal{G}
$$
在所有茎上都是同构，因而是同构。
\end{proof}

\begin{lemma}
\label{sheaves-lemma-equivalence-categories-closed-abelian}
设 $X$ 为拓扑空间，并设 $i : Z \to X$ 为闭子集的包含。函子
$$
i_* : \textit{Ab}(Z) \longrightarrow \textit{Ab}(X)
$$
是全忠实的。其本质像恰由如下层 $\mathcal{G}$ 组成：对所有
$x \in X \setminus Z$，都有 $\mathcal{G}_x = 0$。
\end{lemma}

\begin{proof}
略。
\end{proof}

\begin{lemma}
\label{sheaves-lemma-equivalence-categories-closed-structures}
设 $X$ 为拓扑空间，并设 $i : Z \to X$ 为闭子集的包含。设
$(\mathcal{C}, F)$ 为具有终对象 $0$ 的一种代数结构类型。函子
$$
i_* : \Sh(Z, \mathcal{C}) \longrightarrow \Sh(X, \mathcal{C})
$$
是全忠实的。其本质像恰由如下层 $\mathcal{G}$ 组成：对所有
$x \in X \setminus Z$，都有 $\mathcal{G}_x = 0$。
\end{lemma}

\begin{proof}
略。
\end{proof}

\begin{remark}
\label{sheaves-remark-i-star-not-exact}
设 $i : Z \to X$ 如上为拓扑空间的闭浸入。设 $x \in X$、
$x \not \in Z$，并设 $\mathcal{F}$ 为 $Z$ 上的集合层。由引理
\ref{sheaves-lemma-stalks-closed-pushforward}，有
$(i_*\mathcal{F})_x = \{ * \}$。因此，若
$\mathcal{F} = * \amalg *$，其中 $*$ 是单元素层，则
% P01-ERRATA: P01-E0055 -- the frozen formula drops parentheses around the stalk of i_*F, unlike the preceding formula.
$i_*\mathcal{F}_x = \{*\} \not = i_*(*)_x \amalg i_*(*)_x$，
因为后者是二元素集合。按照 Categories，第
\ref{categories-section-exact-functor} 节中的约定，这意味着作为集合层
范畴之间的函子，$i_*$ 不是右正合的。特别地，它不可能有右伴随；见
Categories，引理 \ref{categories-lemma-exact-adjoint}。

\medskip\noindent
另一方面，稍后将看到（见 Modules，引理
\ref{modules-lemma-i-star-right-adjoint}），阿贝尔层上的 $i_*$ 是正合的，
而且确实有右伴随，即把 $X$ 上的阿贝尔层送到其支撑于 $Z$ 的截面层的函子。
\end{remark}

\begin{remark}
\label{sheaves-remark-closed-immersion-spaces}
我们尚未讨论闭浸入与赋环空间之间的关系。这是因为，赋环空间闭浸入的概念
最好在拟凝聚层的背景下讨论；见 Modules，第
\ref{modules-section-closed-immersion} 节。
\end{remark}

\section{层的粘合}
\label{sheaves-section-glueing-sheaves}

\noindent
本节将粘合定义在 $X$ 的一个覆盖的各成员上的层。先处理映射。

\begin{lemma}
\label{sheaves-lemma-glue-maps}
设 $X$ 为拓扑空间，$X = \bigcup U_i$ 为开覆盖，并设
$\mathcal{F}$、$\mathcal{G}$ 为 $X$ 上的集合层。给定一族层映射
$$
\varphi_i :
\mathcal{F}|_{U_i}
\longrightarrow
\mathcal{G}|_{U_i}
$$
使得对所有 $i, j \in I$，映射 $\varphi_i, \varphi_j$ 在
$\mathcal{F}|_{U_i \cap U_j} \to \mathcal{G}|_{U_i \cap U_j}$ 上限制为
同一个映射，则存在唯一的层映射
$$
\varphi :
\mathcal{F}
\longrightarrow
\mathcal{G}
$$
其在每个 $U_i$ 上的限制与 $\varphi_i$ 相同。
\end{lemma}

\begin{proof}
对每个开子集 $U \subset X$，定义
$$
\varphi_U : \mathcal{F}(U) \to \mathcal{G}(U), \quad
s \mapsto \varphi_U(s)
$$
其中 $\varphi_U(s)$ 是满足
$$
(\varphi_U(s))|_{U \cap U_i}
= (\varphi_i)_{U \cap U_i}(s|_{U \cap U_i}).
$$
的唯一截面。由层公理以及如下事实，这样的截面存在且唯一：
\begin{align*}
((\varphi_i)_{U \cap U_i}(s|_{U \cap U_i}))|_{U \cap U_i \cap U_j}
&= (\varphi_i)_{U \cap U_i \cap U_j}(s|_{U \cap U_i \cap U_j})\\
&= (\varphi_j)_{U \cap U_i \cap U_j}(s|_{U \cap U_i \cap U_j})\\
&= ((\varphi_j)_{U \cap U_j}(s|_{U \cap U_j}))|_{U \cap U_i \cap U_j}.
\end{align*}
这族映射确实给出层映射：若 $V \subset U \subset X$ 为开子集，则
$$
(\varphi_U(s))|_V
= \varphi_V(s|_V)
$$
因为对每个 $i \in I$ 都有
\begin{align*}
(\varphi_U(s))|_{V \cap U_i}
&= ((\varphi_U(s))|_{U \cap U_i})|_{V \cap U_i}\\
&= ((\varphi_i)_{U \cap U_i}(s|_{U \cap U_i}))|_{V \cap U_i}\\
&= (\varphi_i)_{V \cap U_i}(s|_{V \cap U_i})\\
% P01-ERRATA: P01-E0056 -- the frozen final line uses s_V instead of the restriction s|_V.
&= \varphi_V(s_{V})|_{V \cap U_i}.
\end{align*}
此外，它在每个 $U_i$ 上的限制都与 $\varphi_i$ 相同，因为给定开子集
$U \subset X$ 和 $s \in \mathcal{F}(U \cap U_i)$，有
\begin{align*}
\varphi_{U \cap U_i}(s)
&= \varphi_{U \cap U_i}(s)|_{U \cap U_i}\\
&= (\varphi_i)_{U \cap U_i}(s|_{U \cap U_i})\\
&= (\varphi_i)_{U \cap U_i}(s).
\end{align*}
\end{proof}

\noindent
前一引理说明，给定拓扑空间 $X$ 上的两个层 $\mathcal{F}$、
$\mathcal{G}$，规则
$$
U \longmapsto
\Mor_{\Sh(U)}(
\mathcal{F}|_U, \mathcal{G}|_U)
$$
定义一个层。这是一种{\it 内部 Hom 层}。它在集合层的背景下很少使用，
在模层的背景下更常用；见 Modules，第
\ref{modules-section-internal-hom} 节。

\medskip\noindent
设 $X$ 为拓扑空间，$X = \bigcup_{i\in I} U_i$ 为开覆盖。对每个
$i \in I$，设 $\mathcal{F}_i$ 为 $U_i$ 上的集合层。对每一对
$i, j \in I$，设
$$
\varphi_{ij} :
\mathcal{F}_i|_{U_i \cap U_j}
\longrightarrow
\mathcal{F}_j|_{U_i \cap U_j}
$$
为集合层的同构。再假设对每个指标三元组 $i, j, k \in I$，下图交换
$$
\xymatrix{
\mathcal{F}_i|_{U_i \cap U_j \cap U_k}
\ar[rr]_{\varphi_{ik}}
\ar[rd]_{\varphi_{ij}} & &
\mathcal{F}_k|_{U_i \cap U_j \cap U_k} \\
&
\mathcal{F}_j|_{U_i \cap U_j \cap U_k}
\ar[ru]_{\varphi_{jk}}
}
$$
称这样的一族数据 $(\mathcal{F}_i, \varphi_{ij})$ 为
{\it 关于覆盖 $X = \bigcup U_i$ 的集合层粘合数据}。

\begin{lemma}
\label{sheaves-lemma-glue-sheaves}
设 $X$ 为拓扑空间，$X = \bigcup_{i\in I} U_i$ 为开覆盖。给定任意
关于覆盖 $X = \bigcup U_i$ 的集合层粘合数据
$(\mathcal{F}_i, \varphi_{ij})$，则存在 $X$ 上的集合层
$\mathcal{F}$，连同同构
$$
\varphi_i : \mathcal{F}|_{U_i} \to \mathcal{F}_i
$$
使图表
$$
\xymatrix{
\mathcal{F}|_{U_i \cap U_j} \ar[r]_{\varphi_i} \ar[d]_{\text{id}} &
\mathcal{F}_i|_{U_i \cap U_j} \ar[d]^{\varphi_{ij}} \\
\mathcal{F}|_{U_i \cap U_j} \ar[r]^{\varphi_j} &
\mathcal{F}_j|_{U_i \cap U_j}
}
$$
交换。
\end{lemma}

\begin{proof}
第一种证明。在此证明中，给出 $\mathcal{F}$ 在开集 $W \subset X$ 上
截面集合的公式。具体地，定义
$$
\mathcal{F}(W) =
\{
(s_i)_{i \in I} \mid
s_i \in \mathcal{F}_i(W \cap U_i),
\varphi_{ij}(s_i|_{W \cap U_i \cap U_j}) = s_j|_{W \cap U_i \cap U_j}
\}.
$$
对 $W' \subset W$，限制映射通过把每个 $s_i$ 限制到
$W' \cap U_i$ 来定义。$\mathcal{F}$ 的层条件立即由每个
$\mathcal{F}_i$ 的层条件得出。

\medskip\noindent
还需证明 $\mathcal{F}|_{U_i}$ 同构地映到 $\mathcal{F}_i$。设
$W \subset U_i$。此时，$\mathcal{F}(W)$ 定义中的条件蕴含
$s_j = \varphi_{ij}(s_i|_{W \cap U_j})$。而粘合数据定义中的各图表交换，
保证可以从{\it 任意}截面 $s \in \mathcal{F}_i(W)$ 出发，通过令
$s_i = s$ 和 $s_j = \varphi_{ij}(s_i|_{W \cap U_j})$ 得到相容族。

\medskip\noindent
第二种证明（概述）。设 $\mathcal{B}$ 为所有满足如下条件的开集
$U \subset X$ 所成的集合：对某个 $i \in I$，有 $U \subset U_i$。
则 $\mathcal{B}$ 是 $X$ 的拓扑基。对 $U \in \mathcal{B}$，选取满足
$U \subset U_i$ 的 $i \in I$，并置
$\mathcal{F}(U) = \mathcal{F}_i(U)$。利用同构 $\varphi_{ij}$，可见这一
规定“不依赖于 $i$ 的选择”。利用 $\mathcal{F}_i$ 的限制映射，可知
$\mathcal{F}$ 是 $\mathcal{B}$ 上的层。最后，使用引理
\ref{sheaves-lemma-extend-off-basis}，把 $\mathcal{F}$ 扩张为 $X$ 上的唯一层
$\mathcal{F}$。
\end{proof}

\begin{lemma}
\label{sheaves-lemma-glue-sheaves-structures}
设 $X$ 为拓扑空间，$X = \bigcup U_i$ 为开覆盖。设
$(\mathcal{F}_i, \varphi_{ij})$ 为阿贝尔群层、相应地代数结构层、
相应地某个 $X$ 上环层 $\mathcal{O}$ 的 $\mathcal{O}$-模层的粘合数据。
则上面引理 \ref{sheaves-lemma-glue-sheaves} 的证明中的构造，分别给出阿贝尔群层、
代数结构层、$\mathcal{O}$-模层。
\end{lemma}

\begin{proof}
这是因为在该构造中，开集 $W$ 上的截面集合 $\mathcal{F}(W)$ 是下列映射的
等化子
$$
\xymatrix{
\prod_{i \in I} \mathcal{F}_i(W \cap U_i)
\ar@<1ex>[r]
\ar@<-1ex>[r]
&
\prod_{i, j\in I} \mathcal{F}_i(W \cap U_i \cap U_j)
}
$$
在所考虑的每种情形下，该等化子都给出相关范畴中的对象，其底层集合就是
所引用引理中考虑的对象。
\end{proof}

\begin{lemma}
\label{sheaves-lemma-mapping-property-glue}
设 $X$ 为拓扑空间，$X = \bigcup_{i\in I} U_i$ 为开覆盖。把集合层
$\mathcal{F}$ 送到关于覆盖 $X = \bigcup U_i$ 的如下粘合数据的函子
$$
(\mathcal{F}|_{U_i},
(\mathcal{F}|_{U_i})|_{U_i \cap U_j}
\to
(\mathcal{F}|_{U_j})|_{U_i \cap U_j}
)
$$
定义了 $\Sh(X)$ 与粘合数据范畴之间的范畴等价。对阿贝尔层、相应地
代数结构层、相应地 $\mathcal{O}$-模层，也有类似陈述。
\end{lemma}

\begin{proof}
由引理 \ref{sheaves-lemma-glue-maps}，该函子全忠实；由引理
\ref{sheaves-lemma-glue-sheaves}，它本质满（通过一个显式给出的拟逆函子）。
\end{proof}

\noindent
这个引理意味着：若层 $\mathcal{F}$ 由粘合数据
$(\mathcal{F}_i, \varphi_{ij})$ 构造，且 $\mathcal{G}$ 是 $X$ 上的层，
则态射 $f : \mathcal{F} \to \mathcal{G}$ 由一族与粘合映射
$\varphi_{ij}$ 相容的层态射
$$
f_i : \mathcal{F}_i \longrightarrow \mathcal{G}|_{U_i}
$$
给出。类似地，给定层态射 $g : \mathcal{G} \to \mathcal{F}$，等价于
给定一族与粘合映射 $\varphi_{ij}$ 相容的层态射
$$
g_i : \mathcal{G}|_{U_i} \longrightarrow \mathcal{F}_i
$$
% Generated mechanically from the frozen sites.tex target.
% Do not edit this generated file; edit the bound locale source instead.

% OK, start here.
%

\setcounter{chapter}{6}
\chapter{位点与层}



\phantomsection
\label{sites-section-phantom}


\section{引言}
\label{sites-section-introduction}

\noindent
Grothendieck 引入位点这一概念，是为了研究概形的 \'etale 拓扑中的层。
这一概念的基本参考文献或许是 \cite{SGA4}。我们采用的位点概念与
\cite{SGA4} 中的不同；我们所谓的位点，在 \cite[Expos\'e II, D\'efinition 1.3]{SGA4}
中称为赋有预拓扑的范畴。之所以这样处理，是因为
在代数几何中，使用一个给定的覆盖类往往很方便，例如定义概形的某个性质何时
在给定拓扑中是局部性质时即是如此；见 Descent，第
\ref{descent-section-descending-properties} 节。本文的论述将紧随
\cite{ArtinTopologies}。我们不使用宇宙。












\section{预层}
\label{sites-section-presheaves}

\noindent
设 $\mathcal{C}$ 为范畴。一个{\it 集合预层}是从 $\mathcal{C}$ 到
$\textit{Sets}$ 的反变函子 $\mathcal{F}$（见 Categories，备注
\ref{categories-remark-functor-into-sets}）。因此，对 $\mathcal{C}$ 的每个
对象 $U$，都有一个集合 $\mathcal{F}(U)$。这个集合的元素称为
$\mathcal{F}$ 在 $U$ 上的{\it 截面}。对每个态射 $f : V \to U$，映射
$\mathcal{F}(f) : \mathcal{F}(U) \to \mathcal{F}(V)$ 称为{\it 限制映射}，
并且常记作 $f^\ast : \mathcal{F}(U) \to \mathcal{F}(V)$。也可以说，
$f^*(s)$ 是 $s$ 沿 $f$ 的{\it 拉回}。函子性意味着
$g^* f^* (s) = (f \circ g)^*(s)$。有时我们采用记号
$s|_V := f^\ast(s)$。这一记号与拓扑学中的函数限制概念相容，因为若
$W \to V \to U$ 是 $\mathcal{C}$ 中的态射，而 $s$ 是 $\mathcal{F}$ 在
$U$ 上的截面，则由 $\mathcal{F}$ 的函子性可知
$s|_W = (s|_V)|_W$。当然，我们必须谨慎，因为态射 $V \to U$ 很可能不止一个，
而相应的拉回映射肯定未必相等。

\begin{definition}
\label{sites-definition-presheaves-sets}
$\mathcal{C}$ 上的一个{\it 集合预层}，是从 $\mathcal{C}$ 到
$\textit{Sets}$ 的反变函子。{\it 预层态射}是函子之间的变换。集合预层的
范畴记作 $\textit{PSh}(\mathcal{C})$。
\end{definition}

\noindent
注意，对 $\mathcal{C}$ 的任意对象 $U$，点函子 $h_U$（见 Categories，例
\ref{categories-example-hom-functor}）都是一个预层。这些预层称为
{\it 可表预层}。它们具有如下良好性质：对任意预层 $\mathcal{F}$，都有
\begin{equation}
\label{sites-equation-map-representable-into-presheaf}
\Mor_{\textit{PSh}(\mathcal{C})}(h_U, \mathcal{F})
=
\mathcal{F}(U).
\end{equation}
这就是 Yoneda 引理（Categories，引理 \ref{categories-lemma-yoneda}）。

\medskip\noindent
类似地，我们可以定义阿贝尔群预层、环预层等概念。更一般地，还可以定义取值于
某个范畴的预层。

\begin{definition}
\label{sites-definition-presheaf}
设 $\mathcal{C}$、$\mathcal{A}$ 为范畴。$\mathcal{C}$ 上一个取值于
$\mathcal{A}$ 的{\it 预层} $\mathcal{F}$，是从 $\mathcal{C}$ 到
$\mathcal{A}$ 的反变函子，即
$\mathcal{F} : \mathcal{C}^{opp} \to \mathcal{A}$。
$\mathcal{C}$ 上取值于 $\mathcal{A}$ 的预层 $\mathcal{F} \to \mathcal{G}$
的一个{\it 态射}，是从函子 $\mathcal{F}$ 到 $\mathcal{G}$ 的一个变换。
\end{definition}

\noindent
这些分别构成 $\mathcal{C}$ 上取值于 $\mathcal{A}$ 的预层范畴的对象和态射。

\begin{remark}
\label{sites-remark-big-presheaves}
如前所述，我们可以考虑取值于 Categories，备注
\ref{categories-remark-big-categories} 所列任意一个“大全畴”的预层范畴。
这些同样会是“大全畴”；随着论述推进，它们也会列入上述备注中。
\end{remark}






















\section{预层的单射与满射}
\label{sites-section-injective-surjective}

\begin{definition}
\label{sites-definition-presheaves-injective-surjective}
设 $\mathcal{C}$ 为范畴，并设 $\varphi : \mathcal{F}
\to \mathcal{G}$ 是集合预层的映射。
\begin{enumerate}
\item 若对 $\mathcal{C}$ 的每个对象 $U$，映射
$\varphi_U : \mathcal{F}(U) \to \mathcal{G}(U)$ 都是单射，则称
$\varphi$ 是{\it 单射的}。
\item 若对 $\mathcal{C}$ 的每个对象 $U$，映射
$\varphi_U : \mathcal{F}(U) \to \mathcal{G}(U)$ 都是满射，则称
$\varphi$ 是{\it 满射的}。
\end{enumerate}
\end{definition}

\begin{lemma}
\label{sites-lemma-mono-epi}
上面定义的单射（相应地，满射）恰好是
$\textit{PSh}(\mathcal{C})$ 中的单态射（相应地，满态射）。一个映射是同构，
当且仅当它既是单射又是满射。
\end{lemma}

\begin{proof}
我们先证明 $\varphi : \mathcal{F} \to \mathcal{G}$ 是单射，当且仅当它是
$\textit{PSh}(\mathcal{C})$ 中的单态射。“仅当”方向是直接的，故只需证明
“当”方向。假设 $\varphi$ 是单态射。取
$U \in \Ob(\mathcal{C})$；我们需要证明 $\varphi_U$ 是单射。于是取满足
$\varphi_U (a) = \varphi_U (b)$ 的
$a, b \in \mathcal{F}(U)$；需要验证 $a = b$。在同构
(\ref{sites-equation-map-representable-into-presheaf}) 下，
$\mathcal{F}(U)$ 的元素 $a$ 和 $b$ 对应于两个自然变换
$a', b' \in \Mor_{\textit{PSh}(\mathcal{C})}(h_U, \mathcal{F})$。
类似地，在相应的同构
$\Mor_{\textit{PSh}(\mathcal{C})}(h_U, \mathcal{G})
= \mathcal{G}(U)$ 下，$\mathcal{G}(U)$ 中两个相等的元素
$\varphi_U (a)$ 和 $\varphi_U (b)$ 对应于两个自然变换
$\varphi \circ a', \varphi \circ b'
\in \Mor_{\textit{PSh}(\mathcal{C})}(h_U, \mathcal{G})$，
因此这两个自然变换也必相等。故
$\varphi \circ a' = \varphi \circ b'$；由于 $\varphi$ 是单态射，便有
$a' = b'$，进而 $a = b$。这就证明了 (1)。

\medskip\noindent
下面证明 $\varphi : \mathcal{F} \to \mathcal{G}$ 是满射，当且仅当它是
$\textit{PSh}(\mathcal{C})$ 中的满态射。“仅当”方向是直接的，故只需证明
“当”方向。假设 $\varphi$ 是满态射。

\medskip\noindent
对 $\textit{Sets}$ 范畴中的任意两个态射 $f : A \to B$ 和
$g : A \to C$，用 $\text{inl}_{f,g}$ 和 $\text{inr}_{f,g}$ 表示从
$B$ 和 $C$ 到 $B \coprod_A C$ 的两个典范映射。（这里，推出在
$\textit{Sets}$ 中计算。）

\medskip\noindent
现在定义 $\mathcal{C}$ 上的集合预层 $\mathcal{H}$：对每个
$U \in \Ob(\mathcal{C})$，令
$\mathcal{H}(U) = \mathcal{G}(U) \coprod_{\mathcal{F}(U)} \mathcal{G}(U)$
（其中推出在 $\textit{Sets}$ 中计算，并由映射
$\varphi_U : \mathcal{F}(U) \to \mathcal{G}(U)$ 诱导）；它在态射上的作用
以显然的方式定义（利用推出的函子性）。于是存在两个自然变换
$i_1 : \mathcal{G} \to \mathcal{H}$ 和
$i_2 : \mathcal{G} \to \mathcal{H}$，它们在对象
$U \in \Ob(\mathcal{C})$ 处的分量分别由映射
$\text{inl}_{\varphi_U, \varphi_U}$ 和
$\text{inr}_{\varphi_U, \varphi_U}$ 给出。推出的定义表明
$i_1 \circ \varphi = i_2 \circ \varphi$，因而 $i_1 = i_2$（因为
$\varphi$ 是满态射）。所以对每个 $U \in \Ob(\mathcal{C})$，都有
$\text{inl}_{\varphi_U, \varphi_U}
= \text{inr}_{\varphi_U, \varphi_U}$。于是 $\varphi_U$ 必为满射（一个简单的
组合论论证表明，若 $f : A \to B$ 是 $\textit{Sets}$ 中的态射，则
$\text{inl}_{f,f} = \text{inr}_{f,f}$ 当且仅当 $f$ 是满射）。换言之，
$\varphi$ 是满射，(2) 得证。

\medskip\noindent
最后证明 $\varphi : \mathcal{F} \to \mathcal{G}$ 既是单射又是满射，
当且仅当它是 $\textit{PSh}(\mathcal{C})$ 中的同构。这一次，“当”方向是
直接的。为证明“仅当”方向，只须注意：若 $\varphi$ 既单又满，则对每个
$U \in \Ob(\mathcal{C})$，映射 $\varphi_U$ 都可逆；对所有 $U$，这些映射的逆
可组合成自然变换 $\mathcal{G} \to \mathcal{F}$，它正是 $\varphi$ 的逆。
\end{proof}
\begin{definition}
\label{sites-definition-sub-presheaf}
若对每个对象 $U \in \Ob(\mathcal{C})$，集合 $\mathcal{F}(U)$ 都是
$\mathcal{G}(U)$ 的子集，并且这一关系与限制映射相容，则称 $\mathcal{F}$
是 $\mathcal{G}$ 的一个{\it 子预层}。
\end{definition}

\noindent
换言之，包含映射 $\mathcal{F}(U) \to \mathcal{G}(U)$ 可以粘合成一个
（单射的）预层态射 $\mathcal{F} \to \mathcal{G}$。

\begin{lemma}
\label{sites-lemma-image}
设 $\mathcal{C}$ 为范畴。假设
$\varphi : \mathcal{F} \to \mathcal{G}$ 是 $\mathcal{C}$ 上集合预层的态射。
则存在唯一的子预层 $\mathcal{G}' \subset \mathcal{G}$，使得 $\varphi$ 分解为
$\mathcal{F} \to \mathcal{G}' \to \mathcal{G}$，并且第一个映射是满射。
\end{lemma}

\begin{proof}
为证明存在性，只须对每个 $U \in \Ob(C)$ 令
% P01-ERRATA: P01-E0057 -- the frozen source omits \mathcal around C.
$\mathcal{G}'(U) = \varphi_U \left(\mathcal{F}(U)\right)$
（并从 $\mathcal{G}$ 继承其在态射上的作用），再证明这定义了
$\mathcal{G}$ 的一个子预层，且 $\varphi$ 分解为
$\mathcal{F} \to \mathcal{G}' \to \mathcal{G}$，其中第一个映射是满射。
唯一性是直接的。
\end{proof}

\begin{definition}
\label{sites-definition-image}
采用引理 \ref{sites-lemma-image} 的记号。我们称 $\mathcal{G}'$ 为
{\it $\varphi$ 的像}。
\end{definition}















\section{预层的极限与余极限}
\label{sites-section-limits-colimits-PSh}

\noindent
设 $\mathcal{C}$ 为范畴。范畴 $\textit{PSh}(\mathcal{C})$ 中存在极限和
余极限。此外，对任意 $U \in \Ob(\mathcal{C})$，函子
$$
\textit{PSh}(\mathcal{C})
\longrightarrow
\textit{Sets}, \quad
\mathcal{F}
\longmapsto
\mathcal{F}(U)
$$
与极限和余极限交换。或许证明这些陈述的最简便方法如下。给定一个图表
$
\mathcal{F} :
\mathcal{I}
\to
\textit{PSh}(\mathcal{C})
$，定义预层
$$
\mathcal{F}_{\lim} :
U
\longmapsto
\lim_{i \in \mathcal{I}} \mathcal{F}_i(U)
\text{  且  }
\mathcal{F}_{\colim} :
U
\longmapsto
\colim_{i \in \mathcal{I}} \mathcal{F}_i(U)
$$
显然存在投影映射 $\mathcal{F}_{\lim} \to \mathcal{F}_i$ 和典范映射
$\mathcal{F}_i \to \mathcal{F}_{\colim}$。这些映射满足 Categories，定义
\ref{categories-definition-limit}（相应地，Categories，定义
\ref{categories-definition-colimit}）中极限（相应地，余极限）映射的要求。
事实上，它们显然在 $\mathcal{F}$ 上构成一个锥，相应地构成一个余锥。
此外，若 $(\mathcal{G}, q_i : \mathcal{G} \to \mathcal{F}_i)$ 是另一组
这样的数据（如极限定义中那样），则对每个 $U$，我们得到一族满足适当函子性
条件的映射 $\mathcal{G}(U) \to \mathcal{F}_i(U)$，因而得到唯一映射
$\mathcal{G}(U) \to \mathcal{F}_{\lim}(U)$。容易验证，随着 $U$ 变化，
这些映射彼此相容，并来自所需的映射
$\mathcal{G} \to \mathcal{F}_{\lim}$。余极限情形的论证类似。





















\section{预层范畴的函子性}
\label{sites-section-functoriality-PSh}

\noindent
设 $u : \mathcal{C} \to \mathcal{D}$ 为范畴之间的函子。在这种情况下，记
$$
u^p :
\textit{PSh}(\mathcal{D})
\longrightarrow
\textit{PSh}(\mathcal{C})
$$
为这样的函子：它把 $\mathcal{D}$ 上的预层 $\mathcal{G}$ 送到预层
$u^p\mathcal{G} = \mathcal{G} \circ u$。注意，由上一节可知，这个函子与所有
极限和余极限交换。

\medskip\noindent
对 $V \in \Ob(\mathcal{D})$，用 $\mathcal{I}^u_V$ 表示如下范畴：
\begin{equation}
\label{sites-equation-colim-category}
\begin{matrix}
\Ob(\mathcal{I}^u_V)
&
=
&
\{
(U, \phi)
\mid
U \in \Ob(\mathcal{C}),
\phi : V \to u(U)
\}
\\
\Mor_{\mathcal{I}^u_V}((U, \phi), (U', \phi'))
&
=
&
\{
f : U \to U' \text{ 于 }\mathcal{C}
\mid
u(f) \circ \phi = \phi'
\}
\end{matrix}
\end{equation}
有时我们省略记号中的上标 ${}^u$，只写 $\mathcal{I}_V$。我们将利用这些范畴
定义函子 $u^p$ 的左伴随。在此之前，先证明几个技术性引理。

\begin{lemma}
\label{sites-lemma-almost-directed}
设 $u : \mathcal{C} \to \mathcal{D}$ 为范畴之间的函子。假设
$\mathcal{C}$ 有纤维积和等化子，且 $u$ 与它们交换。那么范畴
$(\mathcal{I}_V)^{opp}$ 满足 Categories，引理
\ref{categories-lemma-split-into-directed} 的假设。
\end{lemma}

\begin{proof}
需要验证两个条件。

\medskip\noindent
首先，假设给定三个对象
$\phi : V \to u(U)$、$\phi' : V \to u(U')$ 和
$\phi'' : V \to u(U'')$，以及态射 $a : U' \to U$、
$b : U'' \to U$，使得
$u(a) \circ \phi' = \phi$ 且 $u(b) \circ \phi'' = \phi$。
我们需要证明存在另一个对象 $\phi''' : V \to u(U''')$ 以及态射
$c : U''' \to U'$ 和 $d : U''' \to U''$，使得
$u(c) \circ \phi''' = \phi'$、$u(d) \circ \phi''' = \phi''$ 且
$a \circ c = b \circ d$。取 $U''' = U' \times_U U''$，并以 $c$ 和 $d$
为投影态射。由于 $u$ 与纤维积交换，这一构造有效；验证从略。

\medskip\noindent
其次，假设给定两个对象 $\phi : V \to u(U)$ 和
$\phi' : V \to u(U')$，以及态射
$a, b : (U, \phi) \to (U', \phi')$。我们需要找到一个等化 $a$ 和 $b$
的态射 $c : (U'', \phi'') \to (U, \phi)$。令
$c : U'' \to U$ 为 $\mathcal{C}$ 中 $a$ 和 $b$ 的等化子。由于 $u$ 与
等化子交换，并且
$u(a) \circ \phi = u(b) \circ \phi = \phi'$，我们得到一个态射
$\phi'' : V \to u(U'')$。
\end{proof}

\begin{lemma}
\label{sites-lemma-directed}
设 $u : \mathcal{C} \to \mathcal{D}$ 为范畴之间的函子。假设
\begin{enumerate}
\item 范畴 $\mathcal{C}$ 有终对象 $X$，且 $u(X)$ 是 $\mathcal{D}$ 的终对象；
\item 范畴 $\mathcal{C}$ 有纤维积，且 $u$ 与纤维积交换。
\end{enumerate}
那么指标范畴 $(\mathcal{I}^u_V)^{opp}$ 是滤过的（见 Categories，定义
\ref{categories-definition-directed}）。
\end{lemma}

\begin{proof}
这些假设蕴含引理 \ref{sites-lemma-almost-directed} 的假设成立（见 Categories，
第 \ref{categories-section-finite-limits} 节中的讨论）。由 Categories，引理
\ref{categories-lemma-split-into-directed}，可知 $\mathcal{I}_V$ 是若干定向范畴
的不交并（可能为空）。因此只需证明 $\mathcal{I}_V$ 连通。

\medskip\noindent
首先证明 $\mathcal{I}_V$ 非空。事实上，令 $X$ 为由假设存在的
$\mathcal{C}$ 的终对象。又因假设 $u(X)$ 是 $\mathcal{D}$ 的终对象，便有
态射 $V \to u(X)$。它给出 $\mathcal{I}_V$ 的一个对象。

\medskip\noindent
其次证明 $\mathcal{I}_V$ 连通。设
$\phi_1 : V \to u(U_1)$ 和 $\phi_2 : V \to u(U_2)$ 属于
$\Ob(\mathcal{I}_V)$。由假设，$U_1\times U_2$ 存在，且
$u(U_1\times U_2) = u(U_1)\times u(U_2)$。考虑由积的泛性质中
$(\phi_1, \phi_2)$ 所对应的态射
$\phi : V \to u(U_1\times U_2)$。显然，对象
$\phi : V \to u(U_1\times U_2)$ 同时映到
$\phi_1 : V \to u(U_1)$ 和 $\phi_2 : V \to u(U_2)$。
\end{proof}

\noindent
给定 $\mathcal{D}$ 中的 $g : V' \to V$，在对象上令
$\overline{g}(U, \phi) = (U, \phi \circ g)$，便得到函子
$\overline{g} : \mathcal{I}_V \to \mathcal{I}_{V'}$。给定
$\mathcal{C}$ 上的预层 $\mathcal{F}$，得到函子
$$
\mathcal{F}_V :
\mathcal{I}_V^{opp}
\longrightarrow
\textit{Sets}, \quad
(U, \phi)
\longmapsto
\mathcal{F}(U).
$$
换言之，$\mathcal{F}_V$ 是 $\mathcal{I}_V$ 上的集合预层。注意，
$\mathcal{F}_{V'} \circ \overline{g} = \mathcal{F}_V$。定义
$$
u_p\mathcal{F}(V) =
\colim_{\mathcal{I}_V^{opp}} \mathcal{F}_V
$$
由余极限的性质，对每个 $(U, \phi) \in \Ob(\mathcal{I}_V)$，得到典范映射
$\mathcal{F}(U)\xrightarrow{c(\phi)}u_p\mathcal{F}(V)$。对上述
$g : V' \to V$，由 Categories，引理
\ref{categories-lemma-functorial-colimit}，存在一个与
$\mathcal{F}_{V'} \circ \overline{g} = \mathcal{F}_V$ 相容的典范限制映射
$g^* : u_p\mathcal{F}(V) \to u_p\mathcal{F}(V')$。它是使得对所有
$(U, \phi) \in \Ob(\mathcal{I}_V)$，图表
$$
\xymatrix{
\mathcal{F}(U) \ar[r]^{c(\phi)} \ar[d]_{\text{id}}
&
u_p\mathcal{F}(V) \ar[d]^{g^*}
\\
\mathcal{F}(U) \ar[r]^{c(\phi \circ g)}
&
u_p\mathcal{F}(V')
}
$$
交换的唯一映射。这些映射的唯一性表明，我们由此得到一个预层。这个预层记作
$u_p\mathcal{F}$。

\begin{lemma}
\label{sites-lemma-recover}
存在一个与限制映射（$\mathcal{F}$ 上和 $u_p\mathcal{F}$ 上的限制映射）相容的
典范映射 $\mathcal{F}(U) \to u_p\mathcal{F}(u(U))$。
\end{lemma}

\begin{proof}
这正是上面引入的映射 $c(\text{id}_{u(U)})$。
\end{proof}

\noindent
注意，预层的任意映射 $\mathcal{F} \to \mathcal{F}'$ 都给出函子之间的相容
映射族 $\mathcal{F}_V \to \mathcal{F}'_V$，从而给出预层映射
$u_p\mathcal{F} \to u_p\mathcal{F}'$。换言之，我们定义了函子
$$
u_p :
\textit{PSh}(\mathcal{C})
\longrightarrow
\textit{PSh}(\mathcal{D})
$$

\begin{lemma}
\label{sites-lemma-adjoints-u}
函子 $u_p$ 是函子 $u^p$ 的左伴随。换言之，公式
$$
\Mor_{\textit{PSh}(\mathcal{C})}(\mathcal{F}, u^p\mathcal{G})
=
\Mor_{\textit{PSh}(\mathcal{D})}(u_p\mathcal{F}, \mathcal{G})
$$
对 $\mathcal{F}$ 和 $\mathcal{G}$ 双函子自然地成立。
\end{lemma}

\begin{proof}
设 $\mathcal{G}$ 是 $\mathcal{D}$ 上的预层，$\mathcal{F}$ 是
$\mathcal{C}$ 上的预层。我们将通过构造两个方向的映射来证明所示公式；留给
读者验证它们互为逆映射。

\medskip\noindent
给定映射 $\alpha : u_p \mathcal{F} \to \mathcal{G}$，得到
$u^p\alpha : u^p u_p \mathcal{F} \to u^p \mathcal{G}$。引理
\ref{sites-lemma-recover} 表明存在映射
$\mathcal{F} \to u^p u_p \mathcal{F}$。两者复合即给出所需映射。（这一构造的
好处是，它对眼前的一切显然都具有函子性。）

\medskip\noindent
反过来，给定映射 $\beta : \mathcal{F} \to u^p\mathcal{G}$，得到映射
$u_p\beta : u_p\mathcal{F} \to u_p u^p\mathcal{G}$。取
$V \in \Ob(\mathcal{D})$。我们断言，$\mathcal{I}_V$ 上的函子
$u^p\mathcal{G}_V$ 到取常值 $\mathcal{G}(V)$ 的常值函子有一个典范映射。
事实上，对 $\mathcal{I}_V$ 的每个对象 $(U, \phi)$，
$u^p\mathcal{G}_V$ 在该对象上的值是 $\mathcal{G}(u(U))$，它通过
$\mathcal{G}(\phi) = \phi^*$ 映到 $\mathcal{G}(V)$。由于
$\mathcal{G}$ 本身是函子，这构成函子之间的变换。由此得到映射
$u_p u^p \mathcal{G}(V) \to \mathcal{G}(V)$。另一个直接验证表明，该映射
对 $V$ 具有函子性，从而给出预层映射
$u_p u^p \mathcal{G} \to \mathcal{G}$。复合
$u_p\mathcal{F} \to u_p u^p\mathcal{G} \to \mathcal{G}$
即为所需映射。
\end{proof}

\begin{remark}
\label{sites-remark-functoriality-presheaves-values}
假设 $\mathcal{A}$ 是这样一个范畴：任意图表
$\mathcal{I}_Y \to \mathcal{A}$ 在 $\mathcal{A}$ 中都有余极限。那么显然，
在取值于 $\mathcal{A}$ 的预层范畴上，也存在与上面完全相同方式定义的函子
$u^p$ 和 $u_p$。此外，伴随对 $u^p$ 与 $u_p$ 的伴随性在此情形下仍然成立。
\end{remark}

\begin{lemma}
\label{sites-lemma-pullback-representable-presheaf}
设 $u : \mathcal{C} \to \mathcal{D}$ 为范畴之间的函子。对
$\mathcal{C}$ 的任意对象 $U$，都有 $u_ph_U = h_{u(U)}$。
\end{lemma}

\begin{proof}
由 $u_p$ 与 $u^p$ 的伴随性，有
$$
\Mor_{\textit{PSh}(\mathcal{D})}(u_ph_U, \mathcal{G})
=
\Mor_{\textit{PSh}(\mathcal{C})}(h_U, u^p\mathcal{G})
=
u^p\mathcal{G}(U) =
\mathcal{G}(u(U))
$$
因而由 Yoneda 引理可知，作为预层，$u_ph_U = h_{u(U)}$。
\end{proof}
















\section{位点}
\label{sites-section-sites-definitions}

\noindent
我们的位点概念使用下面这类结构。

\begin{definition}
\label{sites-definition-family-morphisms-fixed-target}
设 $\mathcal{C}$ 为范畴，见 Conventions，第
\ref{conventions-section-categories} 节。$\mathcal{C}$ 中一个
{\it 具有固定靶的态射族}，由一个对象 $U \in \Ob(\mathcal{C})$、一个集合
$I$，以及对每个 $i\in I$ 给出的一个靶为 $U$ 的 $\mathcal{C}$ 中态射
$U_i \to U$ 构成。我们用记号 $\{U_i \to U\}_{i\in I}$ 表示它。
\end{definition}

\noindent
集合 $I$ 可能为空！这一记号意在提示拓扑学中的开覆盖。

\begin{definition}
\label{sites-definition-site}
一个{\it 位点}\footnote{如引言所述，这一记号与 \cite{SGA4} 中的不同。}
由一个范畴 $\mathcal{C}$ 和一族具有固定靶的态射族
$\{U_i \to U\}_{i \in I}$ 所成的集合 $\text{Cov}(\mathcal{C})$ 给出；这些
态射族称为 $\mathcal{C}$ 的{\it 覆盖}，并满足下列公理：
\begin{enumerate}
\item 若 $V \to U$ 是同构，则
$\{V \to U\} \in \text{Cov}(\mathcal{C})$。
\item 若 $\{U_i \to U\}_{i\in I} \in \text{Cov}(\mathcal{C})$，且对每个
$i$ 都有 $\{V_{ij} \to U_i\}_{j\in J_i} \in \text{Cov}(\mathcal{C})$，则
$\{V_{ij} \to U\}_{i \in I, j\in J_i} \in \text{Cov}(\mathcal{C})$。
\item 若 $\{U_i \to U\}_{i\in I}\in \text{Cov}(\mathcal{C})$，且
$V \to U$ 是 $\mathcal{C}$ 的态射，则对所有 $i$，纤维积
$U_i \times_U V$ 都存在，并且
$\{U_i \times_U V \to V \}_{i\in I} \in \text{Cov}(\mathcal{C})$。
\end{enumerate}
\end{definition}

\noindent
作几点说明。在公理 (1) 中，我们要求存在 $\mathcal{C}$ 的一个覆盖
$\{U_i \to U\}_{i \in I}$，使得 $I = \{i\}$ 是单元素集，且态射
$U_i \to U$ 等于 (1) 中给定的态射 $V \to U$。下文中，我们经常用
$\{V \to U\}$ 表示指标集为单元素集的具有固定靶的态射族。在公理 (3) 中，
我们要求：对 $i \in I$ 的纤维积 $U_i \times_U V$ 的某种选择，所述覆盖存在。

\begin{remark}
\label{sites-remark-no-big-sites}
（关于集合论问题——初读时可跳过。）引入位点的主要原因，是研究位点上的层
范畴；它推广了在代数几何中极其重要的拓扑空间上的层范畴。为避免考虑“类的类”
之类的对象，与我们对范畴和 $2$-范畴的做法不同，我们不允许位点是“大全畴”。

\medskip\noindent
假设 $\mathcal{C}$ 是范畴，且 $\text{Cov}(\mathcal{C})$ 是满足上述
(1)、(2)、(3) 的覆盖真类。我们也不允许它作为位点，主要因为接下来要在覆盖
上取极限。不过，有几种自然方法可以把 $\text{Cov}(\mathcal{C})$ 替换成覆盖
的一个集合或稍有不同的结构，而所得层范畴不变。例如：
\begin{enumerate}
\item 在 Sets，第 \ref{sets-section-coverings-site} 节中，我们说明如何选取一个
适当的覆盖集，使所得层范畴相同。
\item 也可以取相伴拓扑（见定义 \ref{sites-definition-topology-associated-site}）。
$\mathcal{C}$ 上所得拓扑具有相同的层范畴。两个拓扑的层范畴相同，当且仅当
它们相等；见定理 \ref{sites-theorem-topology-and-topos}。范畴上的拓扑由对象上的筛
的选择给出。$\mathcal{C}$ 上所有可能的筛，乃至所有可能的拓扑，所成的集合
都是集合。
\item 还可以略微修改位点的概念，见下面的备注
\ref{sites-remark-shrink-coverings}，从而得到一个典范的覆盖集。
\end{enumerate}
每种方案都有一些小缺点。对第一种方案，需要验证后续构造不依赖于覆盖集的选择。
对第二种方案，需要学习拓扑，并为位点重新完成许多论证。对第三种方案，见备注
\ref{sites-remark-shrink-coverings} 的最后一句。

\medskip\noindent
我们将采用上述定义 \ref{sites-definition-site} 中的位点。给定一个范畴
$\mathcal{C}$ 以及如上的覆盖真类，我们用 Sets，引理
\ref{sets-lemma-coverings-site} 将其替换为一个覆盖集，从而产生一个位点。
下面的引理 \ref{sites-lemma-choice-set-coverings-immaterial} 将证明，所得层范畴
（托扑斯）与这一选择无关。若读者愿意，可自行采用另外两种策略之一处理这些
问题。
\end{remark}

\begin{example}
\label{sites-example-site-topological}
设 $X$ 为拓扑空间。令 $X_{Zar}$ 为这样的范畴：其对象是 $X$ 中的所有开集
$U$，态射仅为包含映射。也就是说，$X_{Zar}$ 的任意两个对象之间至多有一个
态射。现在定义
$\{U_i \to U\}_{i \in I}\in \text{Cov}(X_{Zar})$，当且仅当
$\bigcup U_i = U$。上述条件 (1) 和 (2) 显然成立；一旦注意到在
$X_{Zar}$ 中 $U \times V = U \cap V$，条件 (3) 也显然成立。特别要注意，
空集必须是 $X_{Zar}$ 的元素，否则一般而言这一构造无法成立。此外，正如稍后
将看到的，同样重要的是允许{\it 空集的空覆盖作为一个覆盖}！我们按照 Sets，
引理 \ref{sets-lemma-coverings-site} 选取适当的覆盖集
$\text{Cov}(X_{Zar})_{\kappa, \alpha}$，从而把 $X_{Zar}$ 变成位点。
% P01-ERRATA: P01-E0058 -- the frozen source has the ungrammatical plural phrase "a presheaves".
位点 $X_{Zar}$ 上的预层和层（定义见下文），恰好与 Sheaves，第
\ref{sheaves-section-introduction} 节所定义的拓扑空间上的通常预层和层概念
一致。
\end{example}

\begin{example}
\label{sites-example-site-on-group}
设 $G$ 为群。考虑范畴 $G\textit{-Sets}$，其对象是带左 $G$-作用的集合
$X$，态射是 $G$-等变映射。一个重要例子是 ${}_GG$：它是一个 $G$-集合，
其底层集合是 $G$，作用由左乘给出。这个范畴有纤维积，见 Categories，第
\ref{categories-section-example-fibre-products} 节。若
$\bigcup_{i\in I} \varphi_i(U_i) = U$，则称
$\{\varphi_i : U_i \to U\}_{i\in I}$ 为覆盖。这在
$G\textit{-Sets}$ 上给出一个覆盖类，容易看出它满足定义
\ref{sites-definition-site} 的条件 (1)、(2)、(3)。所得结构还不是位点，因为底层
范畴的对象汇集与覆盖汇集都是真类。
% P01-ERRATA: P01-E0059 -- the frozen source says "replace by ... by".
首先按照 Sets，引理 \ref{sets-lemma-sets-with-group-action}，把
$G\textit{-Sets}$ 替换为一个全子范畴 $G\textit{-Sets}_\alpha$。然后如
Sets，引理 \ref{sets-lemma-coverings-site} 那样适当限制覆盖汇集，所得位点
$(G\textit{-Sets}_\alpha,
\text{Cov}_{\kappa, \alpha'}(G\textit{-Sets}_\alpha))$
记作 $\mathcal{T}_G$。

\medskip\noindent
特别地，若群 $G$ 可数，则可令 $\mathcal{T}_G$ 为可数 $G$-集合的范畴，并
把所有指标集 $I$ 可数的联合满态射族
$\{\varphi_i : U_i \to U\}_{i \in I}$ 取作覆盖。
\end{example}

\begin{example}
\label{sites-example-indiscrete}
设 $\mathcal{C}$ 为范畴。有一种典范方式把它变成位点，其中 $U$ 的覆盖正是
$\{f : V \to U \mid f\text{ 为同构}\}$。这个位点上的层就是
$\mathcal{C}$ 上的预层。
% P01-ERRATA: P01-E0060 -- the frozen source says "This corresponding topology".
相应的拓扑称为{\it 混沌拓扑}或{\it 不可分拓扑}。
\end{example}
















\section{层}
\label{sites-section-sheaves}

\noindent
设 $\mathcal{C}$ 为位点。在引入取值于一个范畴的层这一概念之前，我们先说明
一个集合预层成为层是什么意思。设 $\mathcal{F}$ 是 $\mathcal{C}$ 上的集合
预层，并设
$\{U_i \to U\}_{i\in I}$ 是 $\text{Cov}(\mathcal{C})$ 的元素。根据假设，
所有纤维积 $U_i \times_U U_j$ 都在 $\mathcal{C}$ 中存在。有两个自然映射
$$
\xymatrix{
\prod\nolimits_{i\in I}
\mathcal{F}(U_i)
\ar@<1ex>[r]^-{\text{pr}_0^*} \ar@<-1ex>[r]_-{\text{pr}_1^*}
&
\prod\nolimits_{(i_0, i_1) \in I \times I}
\mathcal{F}(U_{i_0} \times_U U_{i_1})
}
$$
如所示，记为 $\text{pr}^*_i$，$i = 0, 1$。事实上，左端的一个元素对应于
一族 $(s_i)_{i\in I}$，其中每个 $s_i$ 都是 $\mathcal{F}$ 在 $U_i$ 上的
截面。对每一对 $(i_0, i_1) \in I \times I$，都有投影态射
$$
\text{pr}^{(i_0, i_1)}_{i_0} :
U_{i_0} \times_U U_{i_1}
\longrightarrow
U_{i_0}
\text{ 且 }
\text{pr}^{(i_0, i_1)}_{i_1} :
U_{i_0} \times_U U_{i_1}
\longrightarrow
U_{i_1}.
$$
因此，可以通过第一个映射拉回截面 $s_{i_0}$，也可以通过第二个映射拉回截面
$s_{i_1}$。明确地说，上述映射为
$$
\text{pr}_0^* :
(s_i)_{i\in I}
\longmapsto
\Big(
\text{pr}^{(i_0, i_1), *}_{i_0}(s_{i_0})
\Big)_{(i_0, i_1) \in I \times I}
$$
和
$$
\text{pr}_1^* :
(s_i)_{i\in I}
\longmapsto
\Big(
\text{pr}^{(i_0, i_1), *}_{i_1}(s_{i_1})
\Big)_{(i_0, i_1) \in I \times I}.
$$
最后考虑自然映射
$$
\mathcal{F}(U)
\longrightarrow
\prod\nolimits_{i\in I}
\mathcal{F}(U_i), \quad
s
\longmapsto
(s|_{U_i})_{i \in I}
$$
这里用 $s|_{U_i}$ 表示 $s$ 沿映射 $U_i \to U$ 的拉回。由
$\mathcal{F}$ 的函子性及纤维积图表的交换性，显然有
$\text{pr}_0^*( (s|_{U_i})_{i \in I} ) =
\text{pr}_1^*( (s|_{U_i})_{i \in I} )$。

\begin{definition}
\label{sites-definition-sheaf-sets}
设 $\mathcal{C}$ 为位点，$\mathcal{F}$ 为 $\mathcal{C}$ 上的集合预层。若对
每个覆盖 $\{U_i \to U\}_{i \in I} \in \text{Cov}(\mathcal{C})$，图表
\begin{equation}
\label{sites-equation-sheaf-condition}
\xymatrix{
\mathcal{F}(U) \ar[r]
&
\prod\nolimits_{i\in I}
\mathcal{F}(U_i)
\ar@<1ex>[r]^-{\text{pr}_0^*} \ar@<-1ex>[r]_-{\text{pr}_1^*}
&
\prod\nolimits_{(i_0, i_1) \in I \times I}
\mathcal{F}(U_{i_0} \times_U U_{i_1})
}
\end{equation}
把第一个箭头表示为 $\text{pr}_0^*$ 与 $\text{pr}_1^*$ 的等化子，则称
$\mathcal{F}$ 是一个{\it 层}。
\end{definition}

\noindent
粗略地说，这意味着：若给定截面 $s_i \in \mathcal{F}(U_i)$，使得对所有
$(i, j) \in I \times I$，在 $\mathcal{F}(U_i \times_U U_j)$ 中都有
$$
s_i|_{U_i \times_U U_j} = s_j|_{U_i \times_U U_j}
$$
则存在唯一的 $s \in \mathcal{F}(U)$，使得 $s_i = s|_{U_i}$。

\begin{remark}
\label{sites-remark-sheaf-condition-empty-covering}
若覆盖 $\{U_i \to U\}_{i \in I}$ 是空族（即 $I = \emptyset$），则层条件
表示 $\mathcal{F}(U) = \{*\}$ 是单元素集。这是因为在
(\ref{sites-equation-sheaf-condition}) 中，第二、第三个集合都是集合范畴中的空积；
它们是集合范畴的终对象，因而是单元素集。
\end{remark}

\begin{example}
\label{sites-example-sheaves-topological}
设 $X$ 为拓扑空间，令 $X_{Zar}$ 为例
\ref{sites-example-site-topological} 中构造的位点。$X_{Zar}$ 上的层概念与
Sheaves，定义 \ref{sheaves-definition-sheaf} 中引入的 $X$ 上层的概念一致。
\end{example}

\begin{example}
\label{sites-example-topological-wrong}
设 $X$ 为拓扑空间。考虑位点 $X'_{Zar}$：除了不允许空集的空覆盖之外，它与例
\ref{sites-example-site-topological} 中的位点 $X_{Zar}$ 相同。换言之，我们允许覆盖
$\{\emptyset \to \emptyset\}$，但不允许指标集为空的覆盖。容易证明，这仍然
定义一个位点。不过，我们断言 $X'_{Zar}$ 上的层与 $X_{Zar}$ 上的层不同。
例如，作为一个极端情形，考虑 $X = \{p\}$ 为单元素集的情况。此时
$X'_{Zar}$ 的对象是 $\emptyset, X$；$\emptyset$ 的每个覆盖都可由
$\{\emptyset \to \emptyset\}$ 加细，而 $X$ 的每个覆盖都可由
$\{X \to X\}$ 加细。显然，这里的一个层可由集合
$\mathcal{F}(\emptyset)$ 的任意选择、集合 $\mathcal{F}(X)$ 的任意选择，
以及限制映射 $\mathcal{F}(X) \to \mathcal{F}(\emptyset)$ 的任意选择给出。
因此，$X'_{Zar}$ 上的层等同于二点空间 $\{\eta, p\}$ 上的通常层，其中开集为
$\{\emptyset, \{\eta\}, \{p, \eta\}\}$。一般而言，$X'_{Zar}$ 上的层
等同于空间 $X \amalg \{\eta\}$ 上的层；该空间的开集是空集以及形如
$U \cup \{\eta\}$ 的任意集合，其中 $U \subset X$ 为开集。
\end{example}


\begin{definition}
\label{sites-definition-category-sheaves-sets}
集合层范畴 {\it $\Sh(\mathcal{C})$} 是
$\textit{PSh}(\mathcal{C})$ 的全子范畴，其对象为集合层。
\end{definition}

\noindent
设 $\mathcal{A}$ 为范畴。若对每个作为覆盖族指标集出现的 $I$，
$\mathcal{A}$ 中都存在以 $I$ 和 $I \times I$ 为指标的积，则上述定义
\ref{sites-definition-sheaf-sets} 有意义，并定义了 $\mathcal{C}$ 上取值于
$\mathcal{A}$ 的层这一概念。注意，$\mathcal{A}$ 中的图表
$$
\xymatrix{
\mathcal{F}(U) \ar[r]
&
\prod\nolimits_{i\in I}
\mathcal{F}(U_i)
\ar@<1ex>[r]^-{\text{pr}_0^*} \ar@<-1ex>[r]_-{\text{pr}_1^*}
&
\prod\nolimits_{(i_0, i_1) \in I \times I}
\mathcal{F}(U_{i_0} \times_U U_{i_1})
}
$$
是等化子图表，当且仅当对 $\mathcal{A}$ 的每个对象 $X$，集合图表
$$
\xymatrix{
\Mor_\mathcal{A}(X, \mathcal{F}(U)) \ar[r]
&
\prod
\Mor_\mathcal{A}(X, \mathcal{F}(U_i))
\ar@<1ex>[r]^-{\text{pr}_0^*} \ar@<-1ex>[r]_-{\text{pr}_1^*}
&
\prod
\Mor_\mathcal{A}(X, \mathcal{F}(U_{i_0} \times_U U_{i_1}))
}
$$
是等化子图表。

\medskip\noindent
现在假设 $\mathcal{A}$ 任意。设 $\mathcal{F}$ 为取值于 $\mathcal{A}$ 的
预层。任取对象 $X\in \Ob(\mathcal{A})$。由规则
$$
\mathcal{F}_X(U) = \Mor_\mathcal{A}(X, \mathcal{F}(U)).
$$
定义一个集合预层 $\mathcal{F}_X$。由上可知，要求所有预层
$\mathcal{F}_X$ 都是集合层，便得到一个良好定义。

\begin{definition}
\label{sites-definition-sheaf}
设 $\mathcal{C}$ 为位点，$\mathcal{A}$ 为范畴，$\mathcal{F}$ 是
$\mathcal{C}$ 上取值于 $\mathcal{A}$ 的预层。若对 $\mathcal{A}$ 的所有对象
$X$，上面定义的集合预层 $\mathcal{F}_X$ 都是层，则称 $\mathcal{F}$ 是一个
{\it 层}。
\end{definition}












\section{具有固定靶的态射族}
\label{sites-section-refinements}

\noindent
本节引入若干与具有相同靶的态射族有关的概念。

\begin{definition}
\label{sites-definition-morphism-coverings}
设 $\mathcal{C}$ 为范畴。设
$\mathcal{U} = \{U_i \to U\}_{i\in I}$ 是 $\mathcal{C}$ 中一个具有固定靶的
态射族，$\mathcal{V} = \{V_j \to V\}_{j\in J}$ 是另一个这样的态射族。
\begin{enumerate}
\item 从 $\mathcal{U}$ 到 $\mathcal{V}$ 的一个
{\it $\mathcal{C}$ 中具有固定靶的映射族的态射}，简称
{\it 从 $\mathcal{U}$ 到 $\mathcal{V}$ 的态射}，由一个态射 $U \to V$、
一个集合映射 $\alpha : I \to J$，以及对每个 $i\in I$ 给出的一个态射
$U_i \to V_{\alpha(i)}$ 构成，使得图表
$$
\xymatrix{
U_i \ar[r] \ar[d]
&
V_{\alpha(i)} \ar[d]
\\
U \ar[r]
&
V
}
$$
交换。
\item 在 $U = V$ 且 $U \to V$ 是恒等态射的特殊情形下，我们称
$\mathcal{U}$ 是态射族 $\mathcal{V}$ 的一个{\it 加细}。
\end{enumerate}
\end{definition}

\noindent
一个平凡但重要的观察是：若
$\mathcal{V} = \{V_j \to V\}_{j \in J}$ 是{\it 空映射族}，即
$J = \emptyset$，则任何满足 $I \not = \emptyset$ 的族
$\mathcal{U} = \{U_i \to V\}_{i \in I}$ 都不能加细 $\mathcal{V}$！

\begin{definition}
\label{sites-definition-combinatorial-tautological}
设 $\mathcal{C}$ 为范畴。设
$\mathcal{U} = \{\varphi_i : U_i \to U\}_{i\in I}$ 和
$\mathcal{V} = \{\psi_j : V_j \to U\}_{j\in J}$ 是两个具有固定靶的态射族。
\begin{enumerate}
\item 若存在映射 $\alpha : I \to J$ 和 $\beta : J\to I$，使得
$\varphi_i = \psi_{\alpha(i)}$ 且 $\psi_j = \varphi_{\beta(j)}$，则称
$\mathcal{U}$ 与 $\mathcal{V}$ {\it 组合等价}。
\item 若存在映射 $\alpha : I \to J$ 和 $\beta : J\to I$，并且对所有
$i\in I$ 和 $j \in J$ 都有交换图表
$$
\xymatrix{
U_i \ar[rd] \ar[rr] & &
V_{\alpha(i)} \ar[ld] & &
V_j \ar[rd] \ar[rr] & &
U_{\beta(j)} \ar[ld] \\
&
U & & & &
U &
}
$$
其中水平箭头都是同构，则称 $\mathcal{U}$ 与 $\mathcal{V}$
{\it 自明等价}。
\end{enumerate}
\end{definition}

\begin{lemma}
\label{sites-lemma-tautological-combinatorial}
设 $\mathcal{C}$ 为范畴。设
$\mathcal{U} = \{\varphi_i : U_i \to U\}_{i\in I}$ 和
$\mathcal{V} = \{\psi_j : V_j \to U\}_{j\in J}$ 是两个具有相同固定靶的
态射族。
\begin{enumerate}
\item 若 $\mathcal{U}$ 与 $\mathcal{V}$ 组合等价，则它们自明等价。
\item 若 $\mathcal{U}$ 与 $\mathcal{V}$ 自明等价，则 $\mathcal{U}$ 是
$\mathcal{V}$ 的加细，且 $\mathcal{V}$ 是 $\mathcal{U}$ 的加细。
\item “组合等价”关系是所有具有固定靶的态射族上的等价关系。
\item “自明等价”关系是所有具有固定靶的态射族上的等价关系。
\item “$\mathcal{U}$ 加细 $\mathcal{V}$ 且 $\mathcal{V}$ 加细
$\mathcal{U}$”关系是所有具有固定靶的态射族上的等价关系。
\end{enumerate}
\end{lemma}

\begin{proof}
从略。
\end{proof}

\noindent
在下面的引理中，给定一个范畴 $\mathcal{C}$、$\mathcal{C}$ 上的预层
$\mathcal{F}$，以及满足所有纤维积 $U_i \times_U U_{i'}$ 都存在的态射族
$\mathcal{U} = \{U_i \to U\}_{i\in I}$，若图表
(\ref{sites-equation-sheaf-condition}) 是等化子图表，我们就说
{\it $\mathcal{F}$ 关于 $\mathcal{U}$ 的层条件}成立。

\begin{lemma}
\label{sites-lemma-tautological-same-sheaf}
设 $\mathcal{C}$ 为范畴。设
$\mathcal{U} = \{\varphi_i : U_i \to U\}_{i\in I}$ 和
$\mathcal{V} = \{\psi_j : V_j \to U\}_{j\in J}$ 是两个具有相同固定靶的
态射族。假设纤维积 $U_i \times_U U_{i'}$ 和
$V_j \times_U V_{j'}$ 都存在。若 $\mathcal{U}$ 与 $\mathcal{V}$ 自明等价，
则对 $\mathcal{C}$ 上任意预层 $\mathcal{F}$，$\mathcal{F}$ 关于
$\mathcal{U}$ 的层条件等价于 $\mathcal{F}$ 关于 $\mathcal{V}$ 的层条件。
\end{lemma}

\begin{proof}
首先注意，若 $\varphi : A \to B$ 是 $\mathcal{C}$ 中的同构，则
$\varphi^* : \mathcal{F}(B) \to \mathcal{F}(A)$ 是同构。设
$\beta : J \to I$ 为映射，并设
$\chi_j : V_j \to U_{\beta(j)}$ 是假设所保证的 $U$ 上同构。我们证明，
$\mathcal{V}$ 的层条件蕴含 $\mathcal{U}$ 的层条件。假设给定截面
$s_i \in \mathcal{F}(U_i)$，使得对所有
$(i, i') \in I \times I$，在 $\mathcal{F}(U_i \times_U U_{i'})$ 中有
$$
s_i|_{U_i \times_U U_{i'}} = s_{i'}|_{U_i \times_U U_{i'}}
$$
定义 $s_j = \chi_j^*s_{\beta(j)}$。对任意一对
$(j, j') \in J \times J$，态射
$\chi_j \times_{\text{id}_U} \chi_{j'} : V_j \times_U V_{j'} \to
U_{\beta(j)} \times_U U_{\beta(j')}$ 也是同构。因此，由结构传递可知
$$
s_j|_{V_j \times_U V_{j'}} = s_{j'}|_{V_j \times_U V_{j'}}
$$
$\mathcal{V}$ 的层条件蕴含存在唯一的 $s$，使得对所有 $j \in J$，都有
$s|_{V_j} = s_j$。由证明开头的观察，这也蕴含对所有
$i \in \Im(\beta)$，都有 $s|_{U_i} = s_i$。现设 $i \in I$ 且
$i \not \in \Im(\beta)$。对此 $i$，在 $U$ 上有同构
$U_i \to V_{\alpha(i)} \to U_{\beta(\alpha(i))}$。它给出态射
$U_i \to U_i \times_U U_{\beta(\alpha(i))}$，且此态射是投影的截面。
由于 $s_i$ 和 $s_{\beta(\alpha(i))}$ 在纤维积上限制为同一元素，可知
$s_{\beta(\alpha(i))}$ 沿
$U_i \to U_{\beta(\alpha(i))}$ 拉回为 $s_i$。因此也有
$s_i = s|_{U_i}$，正合所需。
\end{proof}

\begin{lemma}
\label{sites-lemma-compare-separated-presheaf-condition}
设 $\mathcal{C}$ 为范畴。设
$\mathcal{V} = \{V_j \to U\}_{j \in J} \to
\mathcal{U} = \{U_i \to U\}_{i \in I}$ 是 $\mathcal{C}$ 中具有固定靶的映射
族的一个态射，由 $\text{id} : U \to U$、$\alpha : J \to I$ 和
$f_j : V_j \to U_{\alpha(j)}$ 给出。设 $\mathcal{F}$ 是
$\mathcal{C}$ 上的预层。若
$\mathcal{F}(U) \to \prod_{j \in J} \mathcal{F}(V_j)$ 是单射，则
$\mathcal{F}(U) \to \prod_{i \in I} \mathcal{F}(U_i)$ 也是单射。
\end{lemma}

\begin{proof}
从略。
\end{proof}

\begin{lemma}
\label{sites-lemma-compare-sheaf-condition}
设 $\mathcal{C}$ 为范畴。设
$\mathcal{V} = \{V_j \to U\}_{j \in J} \to
\mathcal{U} = \{U_i \to U\}_{i \in I}$ 是 $\mathcal{C}$ 中具有固定靶的映射
族的一个态射，由 $\text{id} : U \to U$、$\alpha : J \to I$ 和
$f_j : V_j \to U_{\alpha(j)}$ 给出。设 $\mathcal{F}$ 是
$\mathcal{C}$ 上的预层。假设
\begin{enumerate}
\item 纤维积 $U_i \times_U U_{i'}$、$U_i \times_U V_j$、
$V_j \times_U V_{j'}$ 均存在；
\item $\mathcal{F}$ 满足关于 $\mathcal{V}$ 的层条件；
\item 对每个 $i \in I$，映射
$\mathcal{F}(U_i) \to \prod_{j \in J} \mathcal{F}(V_j \times_U U_i)$
都是单射。
\end{enumerate}
那么 $\mathcal{F}$ 满足关于 $\mathcal{U}$ 的层条件。
\end{lemma}

\begin{proof}
由引理 \ref{sites-lemma-compare-separated-presheaf-condition}，映射
$\mathcal{F}(U) \to \prod \mathcal{F}(U_i)$ 是单射。假设给定
$s_i \in \mathcal{F}(U_i)$，使得对所有 $i, i' \in I$，都有
$s_i|_{U_i \times_U U_{i'}} = s_{i'}|_{U_i \times_U U_{i'}}$。令
$s_j = f_j^*(s_{\alpha(j)}) \in \mathcal{F}(V_j)$。由于态射 $f_j$ 都是
$U$ 上的态射，我们得到与投影映射下的 $f_j, f_{j'}$ 相容的诱导态射
% P01-ERRATA: P01-E0061 -- the frozen target uses undefined i,i' instead of j,j'.
$f_{jj'} : V_j \times_U V_{j'} \to
U_{\alpha(i)} \times_U U_{\alpha(i')}$。由此可得
$$
s_j|_{V_j \times_U V_{j'}}
= f_{jj'}^*(s_{\alpha(j)}|_{U_{\alpha(j)} \times_U U_{\alpha(j')}})
= f_{jj'}^*(s_{\alpha(j')}|_{U_{\alpha(j)} \times_U U_{\alpha(j')}})
= s_{j'}|_{V_j \times_U V_{j'}}
$$
对所有 $j, j' \in J$ 成立。因此，由 $\mathcal{F}$ 关于
$\mathcal{V}$ 的层条件，得到截面 $s \in \mathcal{F}(U)$，其在每个 $V_j$
上的限制为 $s_j$。若能证明对任意 $i \in I$，$s$ 在 $U_i$ 上的限制为
$s_i$，则证明完成。由于 $\mathcal{F}$ 满足 (3)，只须证明对所有
$j \in J$，$s$ 和 $s_i$ 在 $U_i \times_U V_j$ 上限制为同一元素。为此使用
$$
s|_{U_i \times_U V_j} = s_j|_{U_i \times_U V_j} =
(\text{id} \times f_j)^*s_{\alpha(j)}|_{U_i \times_U U_{\alpha(j)}} =
(\text{id} \times f_j)^*s_i|_{U_i \times_U U_{\alpha(j)}} =
s_i|_{U_i \times_U V_j}
$$
正合所需。
\end{proof}

\begin{lemma}
\label{sites-lemma-refine-same-topology}
设 $\mathcal{C}$ 为范畴。设 $\text{Cov}_i$，$i = 1, 2$，是两组具有固定靶的
态射族集合，并且各自在 $\mathcal{C}$ 上定义一个位点结构。
\begin{enumerate}
\item 若每个 $\mathcal{U} \in \text{Cov}_1$ 都与某个
$\mathcal{V} \in \text{Cov}_2$ 自明等价，则
$\Sh(\mathcal{C}, \text{Cov}_2) \subset
\Sh(\mathcal{C}, \text{Cov}_1)$。
若还对每个 $\mathcal{U} \in \text{Cov}_2$，都存在某个与之自明等价的
$\mathcal{V} \in \text{Cov}_1$，则两个层范畴相等。
% P01-ERRATA: P01-E0062 -- the frozen source has subject-verb disagreement in "category ... are".
\item 假设对每个 $\mathcal{U} \in \text{Cov}_1$，都存在
$\mathcal{V} \in \text{Cov}_2$，使得 $\mathcal{V}$ 加细 $\mathcal{U}$。
在这种情况下
$\Sh(\mathcal{C}, \text{Cov}_2) \subset
\Sh(\mathcal{C}, \text{Cov}_1)$。
若还对每个 $\mathcal{U} \in \text{Cov}_2$，都存在
$\mathcal{V} \in \text{Cov}_1$，使得 $\mathcal{V}$ 加细 $\mathcal{U}$，
则两个层范畴相等。
\end{enumerate}
\end{lemma}

\begin{proof}
(1) 直接由引理 \ref{sites-lemma-tautological-same-sheaf} 和定义得到。

\medskip\noindent
证明 (2)。设 $\mathcal{F}$ 是位点
$(\mathcal{C}, \text{Cov}_2)$ 上的集合层。取
$\mathcal{U} \in \text{Cov}_1$，写成
$\mathcal{U} = \{U_i \to U\}_{i \in I}$。由假设，可以选择
$\mathcal{U}$ 的一个加细 $\mathcal{V} \in \text{Cov}_2$；写成
$\mathcal{V} = \{V_j \to U\}_{j \in J}$，其加细数据由
$\alpha : J \to I$ 和 $f_j : V_j \to U_{\alpha(j)}$ 给出。注意，
$\mathcal{F}$ 满足 $\mathcal{V}$ 的层条件，也满足覆盖
$\{V_j \times_U U_i \to U_i\}_{j \in J}$ 的层条件，因为这些覆盖属于
$\text{Cov}_2$。故由引理 \ref{sites-lemma-compare-sheaf-condition}，
$\mathcal{F}$ 满足 $\mathcal{U}$ 的层条件。
\end{proof}

\begin{lemma}
\label{sites-lemma-choice-set-coverings-immaterial}
设 $\mathcal{C}$ 为范畴。设 $\text{Cov}(\mathcal{C})$ 是满足定义
\ref{sites-definition-site} 中条件 (1)、(2)、(3) 的覆盖真类。设
$\text{Cov}_1, \text{Cov}_2 \subset \text{Cov}(\mathcal{C})$ 是
$\text{Cov}(\mathcal{C})$ 的两个子集，它们都赋予 $\mathcal{C}$ 位点结构。
若每个覆盖 $\mathcal{U} \in \text{Cov}(\mathcal{C})$ 都与
$\text{Cov}_1$ 中的某个覆盖组合等价，并且与 $\text{Cov}_2$ 中的某个覆盖
组合等价，则
$\Sh(\mathcal{C}, \text{Cov}_1) =
\Sh(\mathcal{C}, \text{Cov}_2)$。
\end{lemma}

\begin{proof}
由上述引理 \ref{sites-lemma-refine-same-topology} 和
\ref{sites-lemma-tautological-combinatorial} 即得，因为假设蕴含每个覆盖
$\mathcal{U} \in \text{Cov}_1 \subset \text{Cov}(\mathcal{C})$
都与 $\text{Cov}_2$ 的一个元素组合等价；交换 $\text{Cov}_1$ 与
$\text{Cov}_2$ 的角色亦然。
\end{proof}

















\section{G-集合的例子}
\label{sites-section-example-sheaf-G-sets}

\noindent
作为例子，考虑例 \ref{sites-example-site-on-group} 中的位点 $\mathcal{T}_G$。
我们将描述 $\mathcal{T}_G$ 上的层范畴。结果将与定义 $\mathcal{T}_G$ 时所作
的选择无关。事实上，在证明中只需要位点 $\mathcal{T}_G$ 的下列性质：
\begin{enumerate}
\item[(a)] $\mathcal{T}_G$ 是 $G\textit{-Sets}$ 的全子范畴；
\item[(b)] $\mathcal{T}_G$ 包含 $G$-集合 ${}_GG$；
\item[(c)] $\mathcal{T}_G$ 有纤维积，且它们与 $G\textit{-Sets}$ 中的相同；
\item[(d)] 给定 $U \in \Ob(\mathcal{T}_G)$ 和 $G$-不变子集
$O \subset U$，存在 $\mathcal{T}_G$ 的一个与 $O$ 同构的对象；
\item[(e)] 任何满射映射族 $\{U_i \to U\}_{i \in I}$，其中
$U, U_i \in \Ob(\mathcal{T}_G)$，都与 $\mathcal{T}_G$ 的某个覆盖组合等价。
\end{enumerate}
这些性质由 Sets，引理 \ref{sets-lemma-what-is-in-it-G-sets} 和
\ref{sets-lemma-coverings-site} 成立。

\medskip\noindent
注意，映射
$$
\Hom_G({}_GG, {}_GG)
\longrightarrow
G^{opp},
\varphi
\longmapsto
\varphi(1)
$$
是群同构。其逆映射把 $g \in G$ 送到映射
$R_g : s \mapsto sg$（即右乘）。注意
$R_{g_1g_2} = R_{g_2} \circ R_{g_1}$，故必须取对群。

\medskip\noindent
这意味着，对 $\mathcal{T}_G$ 上的每个预层 $\mathcal{F}$，集合
$\mathcal{F}({}_GG)$ 按如下方式继承一个 $G$-集合结构：对
$g \in G$ 和 $s \in \mathcal{F}({}_GG)$，把 $g \cdot s$ 定义为
$\mathcal{F}(R_g)(s)$。这是左作用，因为
$$
(g_1g_2) \cdot s  = \mathcal{F}(R_{g_1g_2})(s) =
\mathcal{F}(R_{g_2}\circ R_{g_1})(s) =
\mathcal{F}(R_{g_1})( \mathcal{F}(R_{g_2})(s)) =
g_1 \cdot (g_2 \cdot s).
$$
这里使用了 $\mathcal{F}$ 的反变性。注意，若
$\mathcal{F} \to \mathcal{G}$ 是 $\mathcal{T}_G$ 上集合预层的态射，则得到
映射 $\mathcal{F}({}_GG) \to \mathcal{G}({}_GG)$，它与刚定义的
$G$-作用相容。综上，我们构造了函子
$$
\textit{PSh}(\mathcal{T}_G)
\longrightarrow
G\textit{-Sets}, \quad
\mathcal{F}
\longmapsto
\mathcal{F}({}_GG).
$$
留给读者验证，这一构造具有如下良好性质：可表预层 $h_U$ 被映到一个与 $U$
典范同构的对象。用公式表示，即
$h_U({}_GG) = \Hom_G({}_GG, U) \cong U$。

\medskip\noindent
假设 $S$ 是一个 $G$-集合。用公式\footnote{它看起来可能是由 $S$ 定义的可表
预层，但未必如此，因为 $S$ 可能不是 $\mathcal{T}_G$ 的对象；该位点被选成
一个足够大的 $G$-集合之集。}
$$
\mathcal{F}_S(U)
=
\Mor_{G\textit{-Sets}}(U, S).
$$
定义预层 $\mathcal{F}_S$。这显然是一个预层。另一方面，假设
$\{U_i \to U\}_{i\in I}$ 是 $\mathcal{T}_G$ 中的覆盖。这蕴含
$\coprod_i U_i \to U$ 是满射。因此映射
$$
\mathcal{F}_S(U)
=
\Mor_{G\textit{-Sets}}(U, S)
\longrightarrow
\prod \mathcal{F}_S(U_i)
=
\prod \Mor_{G\textit{-Sets}}(U_i, S)
$$
显然是单射。再给定一族 $G$-等变映射 $s_i : U_i \to S$，使得所有图表
$$
\xymatrix{
U_i \times_U U_j \ar[d] \ar[r]
&
U_j \ar[d]^{s_j}
\\
U_i \ar[r]^{s_i}
&
S
}
$$
都交换，则存在唯一的 $G$-等变映射 $s : U \to S$，使得 $s_i$ 是复合
$U_i \to U \to S$。事实上，定义 $s(u) = s_i(u_i)$，其中 $i\in I$ 是任一
满足存在某个 $u_i \in U_i$ 经映射 $U_i \to U$ 映到 $u$ 的指标。上述图表
的交换性恰好蕴含这一构造是良定义的。综上，我们构造了函子
$$
G\textit{-Sets}
\longrightarrow
\Sh(\mathcal{T}_G), \quad
S
\longmapsto
\mathcal{F}_S
.
$$

\medskip\noindent
现在得到如下范畴与函子的图表：
$$
\xymatrix{
\textit{PSh}(\mathcal{T}_G) \ar[rr]^{\mathcal{F} \mapsto \mathcal{F}({}_GG)}
&
&
G\textit{-Sets} \ar[ld]_{S \mapsto \mathcal{F}_S}
\\
&
\Sh(\mathcal{T}_G) \ar[lu]
&
}
$$
由定义立即可知
$\mathcal{F}_S({}_GG) = \Mor_G({}_GG, S) = S$，最后一个等号来自在 $1$
处取值。这几乎证明了下面的命题。

\begin{proposition}
\label{sites-proposition-sheaves-on-group}
函子 $\mathcal{F} \mapsto \mathcal{F}({}_GG)$ 和
$S \mapsto \mathcal{F}_S$ 给出 $\Sh(\mathcal{T}_G)$ 与
$G\textit{-Sets}$ 之间的一对拟逆等价。
\end{proposition}

\begin{proof}
我们已经看到，按一个方向复合这两个函子得到与恒等函子同构的函子。在另一个
方向上，对任意层 $\mathcal{H}$，存在自然的层映射
$$
can :
\mathcal{H}
\longrightarrow
\mathcal{F}_{\mathcal{H}({}_GG)}.
$$
事实上，对 $\mathcal{T}_G$ 的任意对象 $U$，令 $can_U$ 为映射
$$
\begin{matrix}
\mathcal{H}(U)
&
\longrightarrow
&
\mathcal{F}_{\mathcal{H}({}_GG)}(U)
=
\Mor_G(U, \mathcal{H}({}_GG))
\\
s
&
\longmapsto
&
(u \mapsto \alpha_u^*s).
\end{matrix}
$$
这里 $\alpha_u : {}_GG \to U$ 是映射 $\alpha_u(g) = gu$，而
$\alpha_u^* : \mathcal{H}(U) \to \mathcal{H}({}_GG)$ 是拉回映射。
一个平凡但容易混淆的验证表明，这确实是预层映射。我们需要证明 $can$ 是同构。
为此，证明对所有 $U \in \Ob(\mathcal{T}_G)$，$can_U$ 都是同构。
把 $U = {}_GG$ 这一重要但容易的情形留给读者。$\mathcal{T}_G$ 的一般对象
$U$ 是若干 $G$-轨道的不交并：$U = \coprod_{i\in I} O_i$。映射族
$\{O_i \to U\}_{i \in I}$ 与 $\mathcal{T}_G$ 中的某个覆盖自明等价（由本节
开头列出的 $\mathcal{T}_G$ 的性质）。故由引理
\ref{sites-lemma-tautological-same-sheaf}，层 $\mathcal{H}$ 满足关于
$\{O_i \to U\}_{i \in I}$ 的层性质。这个覆盖的层性质蕴含
$\mathcal{H}(U) = \prod_i \mathcal{H}(O_i)$。因此，只须证明当 $U$ 由单个
$G$-轨道构成时，$can_U$ 是同构。取 $u \in U$，并令 $H \subset G$ 为其
稳定子。显然，$\Mor_G(U, \mathcal{H}({}_GG)) = \mathcal{H}({}_GG)^H$
等于 $H$-不变元素组成的子集。另一方面，考虑由 $g \mapsto gu$ 给出的覆盖
$\{{}_GG \to U\}$（它同样只是与 $\mathcal{T}_G$ 中的某个覆盖组合等价，
而这同样无关紧要）。注意，纤维积 $({}_GG)\times_U ({}_GG)$ 等于
$\{(g, gh), g\in G, h\in H\} \cong \coprod_{h\in H} {}_GG$。因此，这个覆盖
的层性质表述为
$$
\xymatrix{
\mathcal{H}(U) \ar[r]
&
\mathcal{H}({}_GG)
\ar@<1ex>[r]^-{\text{pr}_0^*} \ar@<-1ex>[r]_-{\text{pr}_1^*}
&
\prod_{h \in H}
\mathcal{H}({}_GG).
}
$$
到因子 $\mathcal{H}({}_GG)$ 的两个映射 $\text{pr}_i^*$ 相差乘以 $h$。
现在，结论由此以及 $can$ 在 $U = {}_GG$ 时是同构这一事实得到。
\end{proof}






















\section{层化}
\label{sites-section-sheafification}

\noindent
为了定义层化，我们研究一个覆盖的第零 {\v C}ech 上同调群及其函子性。
% P01-ERRATA: P01-E0063 -- set-valued H^0 need not be a group.

\medskip\noindent
设 $\mathcal{F}$ 是 $\mathcal{C}$ 上的集合预层，并设
$\mathcal{U} = \{U_i \to U\}_{i \in I}$ 是 $\mathcal{C}$ 的一个覆盖。
用记号 $\mathcal{F}(\mathcal{U})$ 表示等化子
$$
H^0(\mathcal{U}, \mathcal{F})
=
\{
(s_i)_{i\in I} \in \prod\nolimits_i \mathcal{F}(U_i)
\mid
s_i|_{U_i \times_U U_j} = s_j|_{U_i \times_U U_j}
\ \forall i, j \in I
\}.
$$
稍后将看到，这是 $\mathcal{F}$ 在 $U$ 上关于覆盖 $\mathcal{U}$ 的第零
{\v C}ech 上同调。顺便指出，只要所有态射 $U_i \to U$ 都可表，便可以定义
$H^0(\mathcal{U}, \mathcal{F})$；也就是说，$\mathcal{U}$ 不必是该位点的
覆盖。存在典范映射
$\mathcal{F}(U) \to H^0(\mathcal{U}, \mathcal{F})$。显然，覆盖的一个态射
$\mathcal{U} \to \mathcal{V}$ 诱导交换图表
$$
\xymatrix{
& U_i \ar[rr] & & V_{\alpha(i)} \\
U_i \times_U U_j \ar[rr] \ar[ur] \ar[dr] & &
V_{\alpha(i)} \times_V V_{\alpha(j)} \ar[ur] \ar[dr] & \\
& U_j \ar[rr] & & V_{\alpha(j)}
}.
$$
它进而产生映射 $H^0(\mathcal{V}, \mathcal{F}) \to
H^0(\mathcal{U}, \mathcal{F})$，并与映射
$\mathcal{F}(V) \to \mathcal{F}(U)$ 相容。

\medskip\noindent
由构造，预层 $\mathcal{F}$ 是层，当且仅当对 $\mathcal{C}$ 的每个覆盖
$\mathcal{U}$，自然映射
$\mathcal{F}(U) \to H^0(\mathcal{U}, \mathcal{F})$ 都是双射。我们将用这一
概念证明下面关于层的极限的简单引理。

\begin{lemma}
\label{sites-lemma-limit-sheaf}
\begin{slogan}
层图表的极限存在，并且与作为预层的极限一致
\end{slogan}
设 $\mathcal{F} : \mathcal{I} \to \Sh(\mathcal{C})$ 是一个图表。那么
$\lim_\mathcal{I} \mathcal{F}$ 存在，并且等于预层范畴中的极限。
\end{lemma}

\begin{proof}
设 $\lim_i \mathcal{F}_i$ 为作为预层的极限。我们证明它是层；由此立即可知，
它也是层范畴中的极限。为证明层性质，设
$\mathcal{V} = \{V_j \to V\}_{j\in J}$ 为覆盖。取
$(s_j)_{j\in J}$ 为 $H^0(\mathcal{V}, \lim_i \mathcal{F}_i)$ 的元素。
利用投影映射，得到 $H^0(\mathcal{V}, \mathcal{F}_i)$ 中的元素
$(s_{j, i})_{j\in J}$。由 $\mathcal{F}_i$ 的层性质，存在唯一的
$s_i \in \mathcal{F}_i(V)$，使得 $s_{j, i} = s_i|_{V_j}$。设
$\phi : i \to i'$ 是指标范畴的态射。我们要证明
$\mathcal{F}(\phi) : \mathcal{F}_i \to \mathcal{F}_{i'}$ 把 $s_i$ 映到
$s_{i'}$。我们知道，对所有 $j$，这对截面 $s_{i, j}$ 和 $s_{i', j}$ 成立；
故由 $\mathcal{F}_{i'}$ 的层性质，所需结论成立。此时得到
$(\lim_i \mathcal{F}_i)(V)$ 的一个元素
$s = (s_i)_{i \in \Ob(\mathcal{I})}$。留给读者验证，该元素具有所需性质
$s_j = s|_{V_j}$。
\end{proof}

\begin{example}
\label{sites-example-singleton-sheaf}
一个特殊例子是空图表上的极限。它给出（预）层范畴中的终对象。该预层把一个
单元素集赋给 $\mathcal{C}$ 的每个对象 $U$，其限制映射均为唯一映射，并且
这个预层还是层。我们常用 $*$ 表示这个层。
\end{example}

\noindent
令 $\mathcal{J}_U$ 为 $U$ 的所有覆盖组成的范畴。换言之，
$\mathcal{J}_U$ 的对象是 $\mathcal{C}$ 中 $U$ 的覆盖，态射是加细。根据我们
对位点的约定，这确实是一个范畴，也就是说，其对象和态射的汇集都是集合。
注意，$\Ob(\mathcal{J}_U)$ 非空，因为 $\{\text{id}_U\}$ 是其中一个对象。
由上述讨论，构造
$\mathcal{U} \mapsto H^0(\mathcal{U}, \mathcal{F})$ 是
$\mathcal{J}_U$ 上的反变函子。定义
$$
\mathcal{F}^{+}(U)
=
\colim_{\mathcal{J}_U^{opp}}
H^0(\mathcal{U}, \mathcal{F})
$$
关于极限和余极限的讨论，见 Categories，第
\ref{categories-section-limits} 节。我们指出，稍后将看到
$\mathcal{F}^{+}(U)$ 是 $\mathcal{F}$ 在 $U$ 上的第零 {\v C}ech 上同调。

\medskip\noindent
在进一步讨论这个余极限的结构之前，先把集合族
$\mathcal{F}^{+}(U)$，$U \in \Ob(\mathcal{C})$，变成一个预层。具体地，
设 $V \to U$ 是 $\mathcal{C}$ 的态射。由位点公理，存在函子\footnote{这一
构造实际上涉及对纤维积 $U_i \times_U V$ 的选择，因而涉及选择公理。所得映射
与所作选择无关；见下文。}
$$
\mathcal{J}_U
\longrightarrow
\mathcal{J}_V, \quad
\{U_i \to U\}
\longmapsto
\{U_i \times_U V \to V\}.
$$
注意，投影映射给出函子性的覆盖态射
$\{U_i \times_U V \to V\} \to \{U_i \to U\}$，从而由上述构造给出
函子性的集合映射
$H^0(\{U_i \to U\}, \mathcal{F}) \to
H^0(\{U_i \times_U V \to V\}, \mathcal{F})$。换言之，存在从函子
$H^0(-, \mathcal{F}) : \mathcal{J}_U^{opp} \to \textit{Sets}$ 到复合
$\mathcal{J}_U^{opp} \to \mathcal{J}_V^{opp}
\xrightarrow{H^0(-, \mathcal{F})} \textit{Sets}$ 的函子变换。因此，由余极限
的一般性质，得到典范映射 $\mathcal{F}^+(U) \to \mathcal{F}^+(V)$。
用上面对集合 $\mathcal{F}^+(U)$ 的描述来说，它把由
$s = (s_i) \in H^0(\{U_i \to U\}, \mathcal{F})$ 所代表的元素送到由
$(s_i|_{V \times_U U_i})
\in H^0(\{U_i \times_U V \to V\}, \mathcal{F})$ 所代表的元素。

\begin{lemma}
\label{sites-lemma-plus-presheaf}
上述构造定义一个预层 $\mathcal{F}^+$，并带有一个典范的预层映射
$\mathcal{F} \to \mathcal{F}^+$。
\end{lemma}

\begin{proof}
只须证明：给定态射 $W \to V \to U$，复合
$\mathcal{F}^+(U) \to \mathcal{F}^+(V) \to \mathcal{F}^+(W)$ 等于映射
$\mathcal{F}^+(U) \to \mathcal{F}^+(W)$。可以直接证明这一点：验证对于给定
的覆盖 $\{U_i \to U\}$ 和
$s = (s_i) \in H^0(\{U_i \to U\}, \mathcal{F})$，存在典范同构
$W \times_U U_i \cong W \times_V (V \times_U U_i)$，并且经此同构，
$s_i|_{W \times_U U_i}$ 对应于
$(s_i|_{V \times_U U_i})|_{W \times_V (V \times_U U_i)}$。
\end{proof}

\noindent
更间接地，引理 \ref{sites-lemma-independent-refinement} 的结果表明，可以沿从
$W$ 的任一覆盖到 $\{U_i \to U\}$ 的任意态射拉回上述元素 $s$，而所得总是
$\mathcal{F}^+(W)$ 中的同一元素。

\begin{lemma}
\label{sites-lemma-plus-functorial}
对应 $\mathcal{F} \mapsto (\mathcal{F} \to \mathcal{F}^+)$ 是一个函子。
\end{lemma}

\begin{proof}
我们不作证明，只准确陈述需要证明的内容。设
$\mathcal{F} \to \mathcal{G}$ 是预层映射。证明下图交换：
$$
\xymatrix{
\mathcal{F} \ar[r] \ar[d]
&
\mathcal{F}^{+} \ar[d]
\\
\mathcal{G} \ar[r]
&
\mathcal{G}^{+}
}
$$
\end{proof}

\noindent
下面两个引理表明，上述余极限是定向集上的余极限。

\begin{lemma}
\label{sites-lemma-common-refinement}
给定位点 $\mathcal{C}$ 的某个对象 $U$ 的一对覆盖
$\{U_i \to U\}$ 和 $\{V_j \to U\}$，存在一个共同加细的覆盖。
\end{lemma}

\begin{proof}
由于 $\mathcal{C}$ 是位点，对每个 $i$，族
$\{V_j \times_U U_i \to U_i\}_j$ 都是覆盖。于是，另一个公理蕴含
$\{V_j \times_U U_i \to U\}_{i, j}$ 是 $U$ 的覆盖。显然，这个覆盖加细
两个给定覆盖。
\end{proof}

\begin{lemma}
\label{sites-lemma-independent-refinement}
覆盖的任意两个态射 $f, g: \mathcal{U} \to \mathcal{V}$，若诱导同一个态射
$U \to V$，则诱导同一个映射
$H^0(\mathcal{V}, \mathcal{F}) \to  H^0(\mathcal{U}, \mathcal{F})$。
\end{lemma}

\begin{proof}
设 $\mathcal{U} = \{U_i \to U\}_{i\in I}$，
$\mathcal{V} = \{V_j \to V\}_{j\in J}$。态射 $f$ 由映射
$U\to V$、映射 $\alpha : I \to J$ 以及映射
$f_i : U_i \to V_{\alpha(i)}$ 构成。类似地，$g$ 决定映射
$\beta : I \to J$ 和映射 $g_i : U_i \to V_{\beta(i)}$。由于 $f$ 与 $g$
诱导同一个映射 $U\to V$，图表
$$
\xymatrix{
&
V_{\alpha(i)} \ar[dr]
\\
U_i \ar[ur]^{f_i} \ar[dr]_{g_i}
&
&
V
\\
&
V_{\beta(i)} \ar[ur]
}
$$
对每个 $i\in I$ 都交换。因此，$f$ 与 $g$ 通过纤维积分解：
$$
\xymatrix{
&
V_{\alpha(i)}
\\
U_i \ar[r]^-\varphi \ar[ur]^{f_i} \ar[dr]_{g_i}
&
V_{\alpha(i)} \times_V V_{\beta(i)} \ar[u]_{\text{pr}_1} \ar[d]^{\text{pr}_2}
\\
&
V_{\beta(i)}.
}
$$
现在取 $s = (s_j)_j \in H^0(\mathcal{V}, \mathcal{F})$。则对所有
$i\in I$：
$$
(f^*s)_i
=
f_i^*(s_{\alpha(i)})
=
\varphi^*\text{pr}_1^*(s_{\alpha(i)})
=
\varphi^*\text{pr}_2^*(s_{\beta(i)})
=
g_i^*(s_{\beta(i)})
=
(g^*s)_i,
$$
其中中间的等号来自 $H^0(\mathcal{V}, \mathcal{F})$ 的定义。这表明由
$f$ 与 $g$ 诱导的映射
$H^0(\mathcal{V}, \mathcal{F}) \to H^0(\mathcal{U}, \mathcal{F})$ 相等。
\end{proof}

\begin{remark}
\label{sites-remark-both-refine-same-H0}
特别地，这个引理表明，若 $\mathcal{U}$ 是 $\mathcal{V}$ 的加细，并且
$\mathcal{V}$ 是 $\mathcal{U}$ 的加细，则存在典范等同
$H^0(\mathcal{U}, \mathcal{F}) =
H^0(\mathcal{V}, \mathcal{F})$。
\end{remark}

\noindent
由这两个引理以及 $\mathcal{J}_U$ 非空这一事实，可知图表
$H^0(-, \mathcal{F}) : \mathcal{J}_U^{opp}
\to \textit{Sets}$ 是滤过的；见 Categories，定义
\ref{categories-definition-directed}。因此，由 Categories，第
\ref{categories-section-directed-colimits} 节，余极限
$\mathcal{F}^{+}(U)$ 可以按如下直接方式描述。集合 $\mathcal{F}^{+}(U)$
中的每个元素，都来自某个 $U$ 的覆盖 $\mathcal{U}$ 的一个元素
$s \in H^0(\mathcal{U}, \mathcal{F})$。再给定一个元素
$s' \in H^0(\mathcal{U}', \mathcal{F})$，则 $s$ 与 $s'$ 决定余极限中的
同一元素，当且仅当存在 $U$ 的一个覆盖 $\mathcal{V}$，以及加细
$f : \mathcal{V} \to \mathcal{U}$ 和
$f' : \mathcal{V} \to \mathcal{U}'$，使得在
$H^0(\mathcal{V}, \mathcal{F})$ 中有 $f^*s = (f')^*s'$。由于平凡覆盖
$\{\text{id}_U\}$ 是 $\mathcal{J}_U$ 的对象，得到典范映射
$\mathcal{F}(U) \to \mathcal{F}^+(U)$。

\begin{lemma}
\label{sites-lemma-plus-surjective}
映射 $\theta : \mathcal{F} \to \mathcal{F}^+$ 具有如下性质：对
$\mathcal{C}$ 的每个对象 $U$ 及每个截面 $s \in \mathcal{F}^+(U)$，存在
覆盖 $\{U_i \to U\}$，使得 $s|_{U_i}$ 属于
$\theta : \mathcal{F}(U_i) \to \mathcal{F}^{+}(U_i)$ 的像。
\end{lemma}

\begin{proof}
事实上，取一个覆盖 $\{U_i \to U\}$，使 $s$ 来自元素
$(s_i) \in H^0(\{U_i \to U\}, \mathcal{F})$。由引理
\ref{sites-lemma-independent-refinement}，为计算 $s$ 拉回到
$\mathcal{F}^+(U_i)$ 的元素，可以使用覆盖 $\{U_i \to U_i\}$ 及显然的覆盖
态射 $\{U_i \to U_i\} \to \{U_i \to U\}$。而使用这个覆盖时，$s$ 在
$U_i$ 上的限制恰好是 $\theta(s_i)$。
\end{proof}

\begin{definition}
\label{sites-definition-separated}
若对所有覆盖 $\{U_i \rightarrow U\}$，映射
$\mathcal{F}(U) \to \prod \mathcal{F}(U_i)$ 都是单射，则称位点
$\mathcal{C}$ 上的集合预层 $\mathcal{F}$ 是{\it 分离的}。
% P01-ERRATA: P01-E0064 -- the frozen source says "for all coverings of {family}".
\end{definition}

\begin{theorem}
\label{sites-theorem-plus}
对上述 $\mathcal{F}$，有：
\begin{enumerate}
\item
\label{sites-item-sep}
预层 $\mathcal{F}^+$ 是分离的。
\item
\label{sites-item-sheaf}
若 $\mathcal{F}$ 是分离的，则 $\mathcal{F}^+$ 是层，并且预层映射
$\mathcal{F} \to \mathcal{F}^+$ 是单射。
\item
\label{sites-item-plus-iso}
若 $\mathcal{F}$ 是层，则 $\mathcal{F} \to \mathcal{F}^+$ 是同构。
\item
\label{sites-item-plusplus}
预层 $\mathcal{F}^{++}$ 总是层。
\end{enumerate}
\end{theorem}

\begin{proof}
证明 (\ref{sites-item-sep})。假设 $s, s' \in \mathcal{F}^+(U)$，并假设存在某个
覆盖 $\{U_i \to U\}$，使得对所有 $i$，都有
$s|_{U_i} = s'|_{U_i}$。现在有 $U$ 的三个覆盖：上述覆盖
$\{U_i \to U\}$，如引理 \ref{sites-lemma-plus-surjective} 中用于 $s$ 的覆盖
$\mathcal{U}$，以及用于 $s'$ 的类似覆盖 $\mathcal{U}'$。由引理
\ref{sites-lemma-common-refinement}，可以找到一个共同加细，记作
$\{W_j \to U\}$。这意味着存在
$s_j, s'_j \in \mathcal{F}(W_j)$，使得
$s|_{W_j} = \theta(s_j)$，对 $s'|_{W_j}$ 亦然，并且
$\theta(s_j) = \theta(s'_j)$。最后一个等式意味着存在某个覆盖
$\{W_{jk} \to W_j\}$，使得
$s_j|_{W_{jk}} = s'_j|_{W_{jk}}$。由于 $\{W_{jk} \to U\}$ 是覆盖，
可知 $s, s'$ 映到 $H^0(\{W_{jk} \to U\}, \mathcal{F})$ 的同一元素，
正合所需。

\medskip\noindent
证明 (\ref{sites-item-sheaf})。$\mathcal{F} \to \mathcal{F}^+$ 显然是单射，因为
所有映射 $\mathcal{F}(U) \to H^0(\mathcal{U}, \mathcal{F})$ 都是单射。
同样显然，若 $\mathcal{U} \to \mathcal{U}'$ 是加细，则
$H^0(\mathcal{U}', \mathcal{F})
\to H^0(\mathcal{U}, \mathcal{F})$ 是单射。现在假设
$\{U_i \to U\}$ 是覆盖，并设 $(s_i)$ 是一族
$\mathcal{F}^+(U_i)$ 的元素，满足对所有 $i, i' \in I$ 的层条件
$s_i|_{U_i \times_U U_{i'}} = s_{i'}|_{U_i \times_U U_{i'}}$。
选择如引理 \ref{sites-lemma-plus-surjective} 中的覆盖
$\{U_{ij} \to U_i\}$，使得 $s_i|_{U_{ij}}$ 是唯一元素
$s_{ij} \in \mathcal{F}(U_{ij})$ 的像。层条件蕴含 $s_{ij}$ 与
$s_{i'j'}$ 在 $U_{ij} \times_U U_{i'j'}$ 上相等，因为该对象映到
$U_i \times_U U_{i'}$，而我们在那里已有等式。因此
$(s_{ij}) \in H^0(\{U_{ij} \to U\}, \mathcal{F})$ 产生一个元素
$s \in \mathcal{F}^+(U)$。留给读者验证 $s|_{U_i} = s_i$。

\medskip\noindent
证明 (\ref{sites-item-plus-iso})。这由定义立即得到，因为层性质恰好说明，对
$U$ 的每个覆盖 $\mathcal{U}$，每个映射
$\mathcal{F} \to H^0(\mathcal{U}, \mathcal{F})$ 都是双射。
% P01-ERRATA: P01-E0065 -- the frozen source omits (U) from the domain of this map.

\medskip\noindent
陈述 (\ref{sites-item-plusplus}) 现在是显然的。
\end{proof}

\begin{definition}
\label{sites-definition-associated-sheaf}
设 $\mathcal{C}$ 为位点，$\mathcal{F}$ 是 $\mathcal{C}$ 上的集合预层。
层 $\mathcal{F}^\# := \mathcal{F}^{++}$ 连同典范映射
$\mathcal{F} \to \mathcal{F}^\#$，称为{\it 与 $\mathcal{F}$ 相伴的层}。
\end{definition}

\begin{proposition}
\label{sites-proposition-sheafification-adjoint}
典范映射 $\mathcal{F} \to \mathcal{F}^\#$ 具有如下泛性质：对任意映射
$\mathcal{F} \to \mathcal{G}$，其中 $\mathcal{G}$ 是集合层，存在唯一映射
$\mathcal{F}^\# \to \mathcal{G}$，使得复合
$\mathcal{F} \to \mathcal{F}^\# \to \mathcal{G}$ 等于给定映射。
\end{proposition}

\begin{proof}
由引理 \ref{sites-lemma-plus-functorial}，得到交换图表
$$
\xymatrix{
\mathcal{F} \ar[r] \ar[d]
&
\mathcal{F}^{+} \ar[r] \ar[d]
&
\mathcal{F}^{++} \ar[d]
\\
\mathcal{G} \ar[r]
&
\mathcal{G}^{+} \ar[r]
&
\mathcal{G}^{++}
}
$$
而由定理 \ref{sites-theorem-plus}，下方的水平映射都是同构。唯一性来自引理
\ref{sites-lemma-plus-surjective}；该引理说明，$\mathcal{F}^\#$ 的每个截面在局部
都来自 $\mathcal{F}$ 的截面。
\end{proof}

\noindent
由此结果容易看出，函子 $\mathcal{F}
\mapsto (\mathcal{F} \to \mathcal{F}^\#)$ 在函子的唯一同构意义下唯一。
事实上，暂用 $i : \Sh(\mathcal{C}) \to \textit{PSh}(\mathcal{C})$ 表示
包含函子。上述结果实际上说明
$$
\Mor_{\textit{PSh}(\mathcal{C})}(\mathcal{F}, i(\mathcal{G}))
=
\Mor_{\Sh(\mathcal{C})}(\mathcal{F}^\#, \mathcal{G}).
$$
换言之，层化函子是包含函子 $i$ 的左伴随。

\begin{example}
\label{sites-example-constant-sheaf}
设 $\mathcal{C}$ 为位点，$E$ 为集合。{\it 取常值 $E$ 的常值层
$\underline{E}$}，是取常值 $E$ 的常值函子所给预层的层化。它由下列规则刻画：
对 $\mathcal{C}$ 上的每个层 $\mathcal{F}$，
$\Mor_{\Sh(\mathcal{C})}(\underline{E}, \mathcal{F}) =
\Mor_{\textit{Sets}}(E, \Gamma(\mathcal{C}, \mathcal{F}))$；这里的记号来自定义
\ref{sites-definition-global-sections}。
\end{example}

\noindent
以几个引理结束本节。

\begin{lemma}
\label{sites-lemma-colimit-sheaves}
\begin{slogan}
层范畴中的余极限等于预层范畴中余极限的层化。
\end{slogan}
设 $\mathcal{F} : \mathcal{I} \to \Sh(\mathcal{C})$ 是一个图表。那么
$\colim_\mathcal{I} \mathcal{F}$ 存在，并且是预层范畴中余极限的层化。
\end{lemma}

\begin{proof}
由于层化函子是左伴随，它与所有余极限交换；见 Categories，引理
\ref{categories-lemma-adjoint-exact}。所以，由于
$\textit{PSh}(\mathcal{C})$ 有余极限，可知 $\Sh(\mathcal{C})$ 也有余极限
（它们是预层中余极限的层化）。
\end{proof}

\begin{lemma}
\label{sites-lemma-sheafification-exact}
函子 $\textit{PSh}(\mathcal{C}) \to \Sh(\mathcal{C})$，
$\mathcal{F} \mapsto \mathcal{F}^\#$，是正合的。
\end{lemma}

\begin{proof}
由于它是左伴随，故右正合；见 Categories，引理
\ref{categories-lemma-exact-adjoint}。另一方面，由引理
\ref{sites-lemma-common-refinement} 和引理 \ref{sites-lemma-independent-refinement}，
$\mathcal{F}^+$ 的构造中的余极限实际上取在定向集
$\Ob(\mathcal{J}_U)$ 上，其中
$\mathcal{U} \geq \mathcal{U}'$ 当且仅当 $\mathcal{U}$ 是
$\mathcal{U}'$ 的加细。因此，由 Categories，引理
\ref{categories-lemma-directed-commutes}，可知
$\mathcal{F} \to \mathcal{F}^+$ 与有限极限交换（视为从预层到预层的函子）。
再用引理 \ref{sites-lemma-limit-sheaf} 即得结论。
\end{proof}

\begin{lemma}
\label{sites-lemma-sections-sheafification}
设 $\mathcal{C}$ 为位点，$\mathcal{F}$ 为 $\mathcal{C}$ 上的集合预层。记
$\theta^2 : \mathcal{F} \to \mathcal{F}^\#$ 为从 $\mathcal{F}$ 到其层化的
典范映射。设 $U$ 为 $\mathcal{C}$ 的对象，并设
$s \in \mathcal{F}^\#(U)$。存在覆盖 $\{U_i \to U\}$ 和截面
$s_i \in \mathcal{F}(U_i)$，使得
\begin{enumerate}
\item $s|_{U_i} = \theta^2(s_i)$；
\item 对每个 $i, j$，存在 $\mathcal{C}$ 的一个覆盖
$\{U_{ijk} \to U_i \times_U U_j\}$，使得 $s_i$ 与 $s_j$ 拉回到每个
$U_{ijk}$ 后相等。
\end{enumerate}
反过来，给定任意覆盖 $\{U_i \to U\}$ 和满足 (2) 的元素
$s_i \in \mathcal{F}(U_i)$，则存在唯一的截面
$s \in \mathcal{F}^\#(U)$，使得 (1) 成立。
\end{lemma}

\begin{proof}
从略。
\end{proof}

\begin{lemma}
\label{sites-lemma-isomorphism-sheafifications}
设 $\mathcal{C}$ 为位点。设 $\mathcal{F} \to \mathcal{G}$ 是
$\mathcal{C}$ 上集合预层的映射。记 $\mathcal{B}$ 为所有满足
$\mathcal{F}(U) \to \mathcal{G}(U)$ 是双射的
$U \in \Ob(\mathcal{C})$ 所成的集合。若 $\mathcal{C}$ 的每个对象都有一个
由 $\mathcal{B}$ 的元素组成的覆盖，则
$\mathcal{F}^\# \to \mathcal{G}^\#$ 是同构。
\end{lemma}

\begin{proof}
取 $U \in \Ob(\mathcal{C})$。我们证明
$\mathcal{F}^\#(U) \to \mathcal{G}^\#(U)$ 是满射。为此，我们将使用引理
\ref{sites-lemma-sections-sheafification}，不再另行说明。对任意
$s \in \mathcal{G}^\#(U)$，存在覆盖 $\{U_i \to U\}$ 和截面
$s_i \in \mathcal{G}(U_i)$，使得
\begin{enumerate}
\item $s|_{U_i}$ 是 $s_i$ 经 $\mathcal{G} \to \mathcal{G}^\#$ 的像；
\item 对每个 $i, j$，存在 $\mathcal{D}$ 的一个覆盖
$\{U_{ijk} \to U_i \times_U U_j\}$，使得 $s_i$ 与 $s_j$ 拉回到每个
$U_{ijk}$ 后相等。
% P01-ERRATA: P01-E0066 -- the frozen source names undefined site D instead of C.
\end{enumerate}
由假设，对每个 $i$，可以选择一个覆盖 $\{U_{i, a} \to U_i\}$，其中
$U_{i, a} \in \mathcal{B}$。于是 $\{U_{i, a} \to U\}$ 是覆盖。用
$s_{i, a}$ 表示 $s_i$ 在 $\mathcal{G}(U_{i, a})$ 中的像，可知
$s_{i, a}$ 与 $s_{j, b}$ 拉回到覆盖
$$
\{(U_{i, a} \times_U U_{j, b}) \times_{U_i \times_U U_j} U_{ijk}
\to U_{i, a} \times_U U_{j, b}\}
$$
的成员上后相等。因此，可以假设 $U_i \in \mathcal{B}$。重复此论证，还可
以假设对所有 $i, j, k$ 都有 $U_{ijk} \in \mathcal{B}$（细节从略）。于是，
由于引理陈述中对 $\mathcal{B}$ 的定义保证
$\mathcal{F}(U_i) \to \mathcal{G}(U_i)$ 是双射，得到唯一的
$t_i \in \mathcal{F}(U_i)$ 映到 $s_i$。$t_i$ 与 $t_j$ 拉回到每个
$U_{ijk}$ 后在 $\mathcal{F}(U_{ijk})$ 中相等，因为
$U_{ijk} \in \mathcal{B}$，并且 $t_i$ 与 $t_j$ 已有相等性。
% P01-ERRATA: P01-E0067 -- the frozen rationale is circular; the known agreement is for s_i,s_j.
于是该引理表明，存在唯一的 $t \in \mathcal{F}^\#(U)$，使得
$t|_{U_i}$ 是 $t_i$ 的像。由此 $t$ 映到 $s$。我们略去
$\mathcal{F}^\#(U) \to \mathcal{G}^\#(U)$ 为单射的证明。
\end{proof}









\section{层的单射与满射}
\label{sites-section-sheaves-injective}

\begin{definition}
\label{sites-definition-sheaves-injective-surjective}
设 $\mathcal{C}$ 为位点，并设 $\varphi : \mathcal{F}
\to \mathcal{G}$ 是集合层的映射。
\begin{enumerate}
\item 若对 $\mathcal{C}$ 的每个对象 $U$，映射
$\varphi : \mathcal{F}(U) \to \mathcal{G}(U)$ 都是单射，则称
$\varphi$ 是{\it 单射的}。
\item 若对 $\mathcal{C}$ 的每个对象 $U$ 和每个截面
$s\in \mathcal{G}(U)$，都存在一个覆盖 $\{U_i \to U\}$，使得对所有 $i$，
限制 $s|_{U_i}$ 都属于
$\varphi : \mathcal{F}(U_i) \to \mathcal{G}(U_i)$ 的像，则称
$\varphi$ 是{\it 满射的}。
\end{enumerate}
\end{definition}

\begin{lemma}
\label{sites-lemma-mono-epi-sheaves}
上面定义的单射（相应地，满射）恰好是范畴 $\Sh(\mathcal{C})$ 中的单态射
（相应地，满态射）。层的映射是同构，当且仅当它既是单射又是满射。
\end{lemma}

\begin{proof}
从略。
\end{proof}

\begin{lemma}
\label{sites-lemma-coequalizer-surjection}
设 $\mathcal{C}$ 为位点。设 $\mathcal{F} \to \mathcal{G}$ 是集合层的满射。
则图表
$$
\xymatrix{
\mathcal{F} \times_\mathcal{G} \mathcal{F}
\ar@<1ex>[r] \ar@<-1ex>[r]
&
\mathcal{F} \ar[r]
&
\mathcal{G}}
$$
把 $\mathcal{G}$ 表示为余等化子。
\end{lemma}

\begin{proof}
设 $\mathcal{H}$ 为集合层，并设 $\varphi : \mathcal{F} \to \mathcal{H}$
是等化两个映射
$\mathcal{F} \times_\mathcal{G} \mathcal{F} \to \mathcal{F}$ 的层映射。
设 $\mathcal{G}' \subset \mathcal{G}$ 为映射
$\mathcal{F} \to \mathcal{G}$ 的预层像。由于积
$\mathcal{F} \times_\mathcal{G} \mathcal{F}$ 可以在预层范畴中计算，可知它
等于预层积 $\mathcal{F} \times_{\mathcal{G}'} \mathcal{F}$。因此
$\varphi$ 诱导唯一的预层映射
$\psi' : \mathcal{G}' \to \mathcal{H}$。由引理
\ref{sites-lemma-mono-epi-sheaves}，$\mathcal{G}$ 是 $\mathcal{G}'$ 的层化，故
$\psi'$ 唯一延拓为层映射 $\psi : \mathcal{G} \to \mathcal{H}$。我们略去
验证 $\varphi$ 等于 $\psi$ 与给定映射的复合。
\end{proof}














\section{可表层}
\label{sites-section-representable-sheaves}

\noindent
设 $\mathcal{C}$ 为范畴。典范拓扑是使所有可表预层都成为层的最细拓扑
（它在定义 \ref{sites-definition-canonical-topology} 中有形式定义，但这里不需要）。
该拓扑未必总是 $\mathcal{C}$ 上某个位点结构的相伴拓扑。若
$\mathcal{C}$ 有纤维积，我们将给出生成此拓扑的一类覆盖。先给出下面的一般
定义。

\begin{definition}
\label{sites-definition-universal-effective-epimorphisms}
设 $\mathcal{C}$ 为范畴。若族 $\{U_i \to U\}_{i \in I}$ 中所有态射
$U_i \to U$ 都可表（见 Categories，定义
\ref{categories-definition-representable-morphism}），并且对任意
$X\in \Ob(\mathcal{C})$，序列
$$
\xymatrix{
\Mor_\mathcal{C}(U, X) \ar[r]
&
\prod\nolimits_{i \in I} \Mor_\mathcal{C}(U_i, X)
\ar@<1ex>[r] \ar@<-1ex>[r]
&
\prod\nolimits_{(i, j) \in I^2} \Mor_\mathcal{C}(U_i \times_U U_j, X)
}
$$
是等化子图表，则称该族是一个{\it 有效满态射}。若对任意态射
$V \to U$，基变换 $\{U_i \times_U V \to V\}$ 都是有效满态射，则称族
$\{U_i \to U\}$ 是一个{\it 泛有效满态射}。
\end{definition}

\noindent
泛有效满态射族所成的类满足定义 \ref{sites-definition-site} 的公理。若
$\mathcal{C}$ 有纤维积，则相伴拓扑就是典范拓扑。（在这种情况下，要得到一个
位点，可按 Sets，引理 \ref{sets-lemma-coverings-site} 论证。）

\medskip\noindent
反过来，假设 $\mathcal{C}$ 是一个位点，且所有可表预层都是层。于是显然，
所有覆盖都是泛有效满态射。因此，在位点语境中，下面的定义才是“正确”的。

\begin{definition}
\label{sites-definition-weaker-than-canonical}
若 $\mathcal{C}$ 的所有覆盖都是泛有效满态射，则称位点 $\mathcal{C}$ 上的
拓扑{\it 弱于典范拓扑}，或称该拓扑是{\it 次典范的}。
\end{definition}

\noindent
可表层是同时为层的可表预层。由于当拓扑不是次典范时最好避免这一术语，我们
只在次典范情形下给出其形式定义。

\begin{definition}
\label{sites-definition-representable-sheaf}
设 $\mathcal{C}$ 为拓扑次典范的位点。Yoneda 嵌入 $h$（见 Categories，
第 \ref{categories-section-opposite} 节）把 $\mathcal{C}$ 表示为
$\mathcal{C}$ 的层范畴的全子范畴。在这种情况下，我们称形如 $h_U$ 的层为
$\mathcal{C}$ 上的{\it 可表层}，其中 $U \in \Ob(\mathcal{C})$。
记号：有时，与 $U$ 相伴的可表层 $h_U$ 记作 {\it $\underline{U}$}。
\end{definition}

\noindent
注意，在该定义的情形下，对每个层 $\mathcal{F}$，都有
$$
\Mor_{\Sh(\mathcal{C})}(h_U, \mathcal{F}) = \mathcal{F}(U)
$$
因为该等式对预层成立；见
(\ref{sites-equation-map-representable-into-presheaf})。一般而言，预层 $h_U$ 不是
层；要得到层，必须对它们作层化。在这种情况下，仍有
\begin{equation}
\label{sites-equation-map-representable-into-sheaf}
\Mor_{\Sh(\mathcal{C})}(h_U^\#, \mathcal{F}) =
\Mor_{\textit{PSh}(\mathcal{C})}(h_U, \mathcal{F}) =
\mathcal{F}(U)
\end{equation}
对每个层 $\mathcal{F}$ 成立。事实上，第一个等式来自 $\#$ 的伴随性质，
第二个等式就是
(\ref{sites-equation-map-representable-into-presheaf})。

\begin{lemma}
\label{sites-lemma-covering-surjective-after-sheafification}
\begin{slogan}
覆盖经层化后成为满射。
\end{slogan}
设 $\mathcal{C}$ 为位点。若 $\{U_i \to U\}_{i \in I}$ 是位点
$\mathcal{C}$ 的一个覆盖，则集合预层的态射
$$
\coprod\nolimits_{i \in I} h_{U_i} \to h_U
$$
经层化后成为满射。
\end{lemma}

\begin{proof}
由上述引理 \ref{sites-lemma-mono-epi-sheaves}，需要证明
$\coprod\nolimits_{i \in I} h_{U_i}^\# \to h_U^\#$ 是满态射。设
$\mathcal{F}$ 为集合层。态射 $h_U^\# \to \mathcal{F}$ 对应于截面
$s \in \mathcal{F}(U)$。因此，映射
$\Mor(h_U^\#, \mathcal{F})
\to \prod_i \Mor(h_{U_i}^\#, \mathcal{F})$ 的单射性直接来自
$\mathcal{F}$ 的层性质。
\end{proof}

\noindent
下一个引理说明：在拓扑弱于典范拓扑的情况下，每个层在某种意义下都由可表层
组成。

\begin{lemma}
\label{sites-lemma-sheaf-coequalizer-representable}
设 $\mathcal{C}$ 为位点。设 $E \subset \Ob(\mathcal{C})$ 是一个子集，使得
$\mathcal{C}$ 的每个对象都有由 $E$ 的元素组成的覆盖。设 $\mathcal{F}$ 为
集合层。存在集合层的图表
$$
\xymatrix{
\mathcal{F}_1 \ar@<1ex>[r] \ar@<-1ex>[r] &
\mathcal{F}_0 \ar[r] &
\mathcal{F}
}
$$
把 $\mathcal{F}$ 表示为余等化子，并且对 $i = 0, 1$，$\mathcal{F}_i$ 都是
形如 $h_U^\#$（其中 $U \in E$）的层的余积。
\end{lemma}

\begin{proof}
首先证明存在所需类型的满态射 $\mathcal{F}_0 \to \mathcal{F}$。事实上，只须取
$$
\mathcal{F}_0 =
\coprod\nolimits_{U \in E, s \in \mathcal{F}(U)}
(h_U)^\# \longrightarrow \mathcal{F}
$$
这里，箭头限制在与 $(U, s)$ 对应的分量上时，把元素
$\text{id}_U \in h_U^\#(U)$ 映到截面 $s \in \mathcal{F}(U)$。由引理
\ref{sites-lemma-mono-epi-sheaves} 以及对 $E$ 的条件，这是一条满态射。为构造
$\mathcal{F}_1$，先令
$\mathcal{G} = \mathcal{F}_0 \times_\mathcal{F} \mathcal{F}_0$，再像上面那样
构造满态射 $\mathcal{F}_1 \to \mathcal{G}$。见引理
\ref{sites-lemma-coequalizer-surjection}。
\end{proof}

\begin{lemma}
\label{sites-lemma-colimit-representable}
设 $\mathcal{C}$ 为位点，$\mathcal{F}$ 为 $\mathcal{C}$ 上的集合层。则存在
图表 $\mathcal{I} \to \mathcal{C}$，$i \mapsto U_i$，使得
$$
\mathcal{F} = \colim_{i \in \mathcal{I}} h_{U_i}^\#
$$
此外，若 $E \subset \Ob(\mathcal{C})$ 是一个子集，使 $\mathcal{C}$ 的每个
对象都有由 $E$ 的元素组成的覆盖，则可假设对所有
$i \in \Ob(\mathcal{I})$，$U_i$ 都是 $E$ 的元素。
\end{lemma}

\begin{proof}
令 $\mathcal{I}$ 为如下范畴：其对象是二元组 $(U, s)$，其中
$U \in \Ob(\mathcal{C})$ 且 $s \in \mathcal{F}(U)$；态射
$(U, s) \to (U', s')$ 是 $\mathcal{C}$ 中满足 $f^*s' = s$ 的
态射 $f : U \to U'$。对 $\mathcal{I}$ 的每个对象 $(U, s)$，由 Yoneda 引理，
元素 $s$ 定义预层映射 $s : h_U \to \mathcal{F}$，从而由层化的泛性质得到
映射 $h_U^\# \to \mathcal{F}$。立即可见，这些映射与范畴
$\mathcal{I}$ 中的态射相容。因此，得到预层映射
$\colim_{(U, s)} h_U \to \mathcal{F}$（其中余极限在预层范畴中取），以及层
映射 $\colim_{(U, s)} (h_U)^\# \to \mathcal{F}$（其中余极限在层范畴中取）。
由于层化是包含函子
$\Sh(\mathcal{C}) \to \textit{PSh}(\mathcal{C})$ 的左伴随（命题
\ref{sites-proposition-sheafification-adjoint}），由 Categories，引理
\ref{categories-lemma-adjoint-exact}，有
$\colim (h_U)^\# = (\colim h_U)^\#$。因此只须证明
$\colim_{(U, s)} h_U \to \mathcal{F}$ 是预层同构。为此，证明对
$\mathcal{C}$ 的每个对象 $X$，映射
$\colim_{(U, s)} h_U(X) \to \mathcal{F}(X)$ 是双射。事实上，它有如下逆映射：
把元素 $t \in \mathcal{F}(X)$ 送到集合
$\colim_{(U, s)} h_U(X)$ 中由 $(X, t)$ 和
$\text{id}_X \in h_X(X)$ 对应的元素。

\medskip\noindent
最后一个陈述的证明从略。
\end{proof}





\section{连续函子}
\label{sites-section-continuous-functors}

\begin{definition}
\label{sites-definition-continuous}
设 $\mathcal{C}$ 和 $\mathcal{D}$ 为位点。若对每个
$\{V_i \to V\}_{i\in I} \in \text{Cov}(\mathcal{C})$，函子
$u : \mathcal{C} \to \mathcal{D}$ 都满足下列条件，则称它是
{\it 连续的}：
\begin{enumerate}
\item $\{u(V_i) \to u(V)\}_{i\in I}$ 属于 $\text{Cov}(\mathcal{D})$；
\item 对 $\mathcal{C}$ 中任意态射 $T \to V$，态射
$u(T \times_V V_i) \to u(T) \times_{u(V)} u(V_i)$ 是同构。
\end{enumerate}
\end{definition}

\noindent
回忆一下，给定上述函子 $u$ 和 $\mathcal{D}$ 上的集合预层
$\mathcal{F}$，我们已把 $u^p\mathcal{F}$ 定义为预层
$\mathcal{F} \circ u$，换言之，对 $\mathcal{C}$ 的每个对象 $V$，有
$$
u^p\mathcal{F} (V) = \mathcal{F}(u(V))
$$

\begin{lemma}
\label{sites-lemma-pushforward-sheaf}
设 $\mathcal{C}$ 和 $\mathcal{D}$ 为位点。设
$u : \mathcal{C} \to \mathcal{D}$ 为连续函子。若 $\mathcal{F}$ 是
$\mathcal{D}$ 上的层，则 $u^p\mathcal{F}$ 也是层。
\end{lemma}

\begin{proof}
设 $\{V_i \to V\}$ 为覆盖。由假设，$\{u(V_i) \to u(V)\}$ 是
$\mathcal{D}$ 中的覆盖，并且
$u(V_i \times_V V_j) = u(V_i)\times_{u(V)}u(V_j)$。所以，
$u^p\mathcal{F}$ 关于覆盖 $\{V_i \to V\}$ 的层条件，恰好就是
$\mathcal{F}$ 关于覆盖 $\{u(V_i) \to u(V)\}$ 的层条件。
\end{proof}

\noindent
为避免混淆，有时把 $u^p$ 限制在集合层子范畴上所得的函子记作
$$
u^s :
\Sh(\mathcal{D})
\longrightarrow
\Sh(\mathcal{C})
$$
回忆一下，$u^p$ 有左伴随
$u_p : \textit{PSh}(\mathcal{C}) \to \textit{PSh}(\mathcal{D})$；见第
\ref{sites-section-functoriality-PSh} 节。

\begin{lemma}
\label{sites-lemma-adjoint-sheaves}
在引理 \ref{sites-lemma-pushforward-sheaf} 的情形下，函子
$u_s : \mathcal{G} \mapsto (u_p \mathcal{G})^\#$ 是 $u^s$ 的左伴随。
\end{lemma}

\begin{proof}
直接由引理 \ref{sites-lemma-adjoints-u} 和命题
\ref{sites-proposition-sheafification-adjoint} 得到。
\end{proof}

\noindent
下面是一个技术性引理。

\begin{lemma}
\label{sites-lemma-technical-up}
在引理 \ref{sites-lemma-pushforward-sheaf} 的情形下，对 $\mathcal{C}$ 上的任意预层
$\mathcal{G}$，都有
$(u_p\mathcal{G})^\# = (u_p(\mathcal{G}^\#))^\#$。
\end{lemma}

\begin{proof}
对 $\mathcal{D}$ 上的任意层 $\mathcal{F}$，有
\begin{eqnarray*}
\Mor_{\Sh(\mathcal{D})}(u_s(\mathcal{G}^\#), \mathcal{F})
& = &
\Mor_{\Sh(\mathcal{C})}(\mathcal{G}^\#, u^s\mathcal{F}) \\
& = &
\Mor_{\textit{PSh}(\mathcal{C})}(\mathcal{G}^\#, u^p\mathcal{F}) \\
& = &
\Mor_{\textit{PSh}(\mathcal{C})}(\mathcal{G}, u^p\mathcal{F}) \\
& = &
\Mor_{\textit{PSh}(\mathcal{D})}(u_p\mathcal{G}, \mathcal{F}) \\
& = &
\Mor_{\Sh(\mathcal{D})}((u_p\mathcal{G})^\#, \mathcal{F})
\end{eqnarray*}
结论由 Yoneda 引理得到。
\end{proof}

\begin{lemma}
\label{sites-lemma-pullback-representable-sheaf}
设 $u : \mathcal{C} \to \mathcal{D}$ 为位点之间的连续函子。对
$\mathcal{C}$ 的任意对象 $U$，都有 $u_sh_U^\# = h_{u(U)}^\#$。
\end{lemma}

\begin{proof}
由引理 \ref{sites-lemma-pullback-representable-presheaf} 和
\ref{sites-lemma-technical-up} 得到。
\end{proof}


\begin{remark}
\label{sites-remark-quasi-continuous}
（初读时可跳过。）设 $\mathcal{C}$ 和 $\mathcal{D}$ 为位点。我们利用自明
等价映射族的定义（见定义 \ref{sites-definition-combinatorial-tautological}），
（略微）削弱连续性的定义条件。设
$u : \mathcal{C} \to \mathcal{D}$ 为函子。若对每个
$\mathcal{V} = \{V_i \to V\}_{i\in I} \in \text{Cov}(\mathcal{C})$，
都有下列性质，则称 $u$ 是{\it 拟连续的}：
\begin{enumerate}
\item[(1')] 映射族 $\{u(V_i) \to u(V)\}_{i\in I}$ 与
$\text{Cov}(\mathcal{D})$ 的一个元素自明等价；
\item[(2)] 对 $\mathcal{C}$ 中任意态射 $T \to V$，态射
$u(T \times_V V_i) \to u(T) \times_{u(V)} u(V_i)$ 是同构。
\end{enumerate}
我们将看到，当 $u$ 拟连续时，引理 \ref{sites-lemma-pushforward-sheaf} 和
\ref{sites-lemma-adjoint-sheaves} 同样成立。

\medskip\noindent
首先注意，态射 $u(V_i) \to u(V)$ 是可表的，因为由第一个条件，它们与可表
态射同构。特别地，族
$u(\mathcal{V}) = \{u(V_i) \to u(V)\}_{i\in I}$ 对 $\mathcal{D}$ 上的任意
预层 $\mathcal{F}$ 都给出一个第零 {\v C}ech 上同调群
$H^0(u(\mathcal{V}), \mathcal{F})$。
% P01-ERRATA: P01-E0063 -- this is the same set-valued H^0/group mismatch.
设 $\mathcal{U} = \{U_j \to u(V)\}_{j \in J}$ 是
$\text{Cov}(\mathcal{D})$ 的一个元素，并且与
$\{u(V_i) \to u(V)\}_{i \in I}$ 自明等价。注意，$u(\mathcal{V})$ 是
$\mathcal{U}$ 的加细，反之亦然。因此，由备注
\ref{sites-remark-both-refine-same-H0}，有
$H^0(u(\mathcal{V}), \mathcal{F}) = H^0(\mathcal{U}, \mathcal{F})$。
特别地，若 $\mathcal{F}$ 是层，则由用第零 {\v C}ech 上同调群表述的层性质，
有 $\mathcal{F}(u(V)) = H^0(u(\mathcal{V}), \mathcal{F})$。由此可知，若
$\mathcal{F}$ 是层，则 $u^p\mathcal{F}$ 是层，因为
$H^0(\mathcal{V}, u^p\mathcal{F}) =
H^0(u(\mathcal{V}), \mathcal{F})$，而刚才已经看到，它等于
$\mathcal{F}(u(V)) = u^p\mathcal{F}(V)$。这样便证明了引理
\ref{sites-lemma-pushforward-sheaf}。引理 \ref{sites-lemma-adjoint-sheaves} 随即成立。
\end{remark}







\section{位点态射}
\label{sites-section-morphism-sites}

\begin{definition}
\label{sites-definition-morphism-sites}
设 $\mathcal{C}$ 和 $\mathcal{D}$ 为位点。一个{\it 位点态射}
$f : \mathcal{D} \to \mathcal{C}$，由连续函子
$u : \mathcal{C} \to \mathcal{D}$ 给出，并要求函子 $u_s$ 正合。
\end{definition}

\noindent
注意，函子 $u$ 的方向与态射 $f$ 的方向{\it 相反}。若
$f \leftrightarrow u$ 是位点态射，则采用记号
$f^{-1} = u_s$ 和 $f_* = u^s$。函子 $f^{-1}$ 称为{\it 逆像函子}，函子
$f_*$ 称为{\it 直像函子}。与拓扑学中一样，有下列伴随性质：
$$
\Mor_{\Sh(\mathcal{D})}(f^{-1}\mathcal{G}, \mathcal{F})
=
\Mor_{\Sh(\mathcal{C})}(\mathcal{G}, f_*\mathcal{F})
$$
这一定义的动机来自下面的例子。

\begin{example}
\label{sites-example-continuous-map}
设 $f : X  \to Y$ 是拓扑空间的连续映射。回忆一下，我们有位点
$X_{Zar}$ 和 $Y_{Zar}$；见例 \ref{sites-example-site-topological}。考虑函子
$u : Y_{Zar} \to X_{Zar}$，$V \mapsto f^{-1}(V)$。这个函子显然连续，因为
开覆盖的逆像仍是开覆盖。（实际上，这取决于如何为 $X_{Zar}$ 与 $Y_{Zar}$
选择覆盖集。但无论如何，该函子是拟连续的；见备注
\ref{sites-remark-quasi-continuous}。）容易验证，函子 $u^s$ 等于拓扑学中通常的
直像函子 $f_*$。因此，由于 $u_s$ 是伴随函子，而通常的拓扑逆像函子
$f^{-1}$ 也是伴随函子，得到典范同构 $f^{-1} \cong u_s$。由于 $f^{-1}$
正合，可知 $u_s$ 正合。故 $u$ 定义位点态射
$f : X_{Zar} \to Y_{Zar}$；我们也可以把它记作 $f$，因为已经看到函子
$u_s, u^s$ 与它们通常的概念一致。
\end{example}

\begin{example}
\label{sites-example-finer-topology}
设 $\mathcal{C}$ 为范畴。设
$$
\text{Cov}(\mathcal{C}) \supset \text{Cov}'(\mathcal{C})
$$
是两组具有固定靶的态射族集合，并且都把 $\mathcal{C}$ 变成位点。记
$\mathcal{C}_\tau$ 为对应于 $\text{Cov}(\mathcal{C})$ 的位点，记
$\mathcal{C}_{\tau'}$ 为对应于 $\text{Cov}'(\mathcal{C})$ 的位点。我们断言，
$\mathcal{C}$ 上的恒等函子定义位点态射
$$
\epsilon : \mathcal{C}_\tau \longrightarrow \mathcal{C}_{\tau'}
$$
事实上，注意 $\text{id} : \mathcal{C}_{\tau'} \to \mathcal{C}_\tau$ 连续，
因为每个 $\tau'$-覆盖都是 $\tau$-覆盖。因此，函子
$\epsilon_* = \text{id}^s$ 是底层预层上的恒等函子。故 $\epsilon_*$ 的左伴随
$\epsilon^{-1}$ 把 $\tau'$-层 $\mathcal{F}$ 送到 $\mathcal{F}$ 的
$\tau$-层化（由层化的泛性质）。$\tau'$-层的有限极限与预层的有限极限一致
（引理 \ref{sites-lemma-limit-sheaf}），而 $\tau$-层化与有限极限交换（引理
\ref{sites-lemma-sheafification-exact}）。因此，$\epsilon^{-1}$ 左正合。又因为
$\epsilon^{-1}$ 是左伴随，它也右正合（Categories，引理
\ref{categories-lemma-exact-adjoint}）。所以 $\epsilon^{-1}$ 正合，我们已经
验证了定义 \ref{sites-definition-morphism-sites} 的全部条件。
\end{example}

\begin{lemma}
\label{sites-lemma-composition-morphisms-sites}
\begin{slogan}
位点函子的复合保持位点函子的连续性。
\end{slogan}
设 $\mathcal{C}_i$，$i = 1, 2, 3$，为位点。设
$u : \mathcal{C}_2 \to \mathcal{C}_1$ 和
$v : \mathcal{C}_3 \to \mathcal{C}_2$ 是诱导位点态射的连续函子。那么函子
$u \circ v : \mathcal{C}_3 \to \mathcal{C}_1$ 连续，并定义位点态射
$\mathcal{C}_1 \to \mathcal{C}_3$。
\end{lemma}

\begin{proof}
由定义立即可知 $u \circ v$ 是连续函子。此外，显然有
$(u \circ v)^p = v^p \circ u^p$，从而
$(u \circ v)^s = v^s \circ u^s$。因此，函子 $(u \circ v)_s$ 与
$u_s \circ v_s$ 都是 $(u \circ v)^s$ 的左伴随。故
$(u \circ v)_s \cong u_s \circ v_s$，于是 $(u \circ v)_s$ 作为正合函子的
复合是正合的。
\end{proof}

\begin{definition}
\label{sites-definition-composition-morphisms-sites}
设 $\mathcal{C}_i$，$i = 1, 2, 3$，为位点。设
$f : \mathcal{C}_1 \to \mathcal{C}_2$ 和
$g : \mathcal{C}_2 \to \mathcal{C}_3$ 是位点态射，分别由连续函子
$u : \mathcal{C}_2 \to \mathcal{C}_1$ 和
$v : \mathcal{C}_3 \to \mathcal{C}_2$ 给出。{\it 复合} $g \circ f$ 是对应于
函子 $u \circ v$ 的位点态射。
\end{definition}

\noindent
在这种情况下，有 $(g \circ f)_* = g_* \circ f_*$ 及
$(g \circ f)^{-1} = f^{-1} \circ g^{-1}$（见引理
\ref{sites-lemma-composition-morphisms-sites} 的证明）。

\begin{lemma}
\label{sites-lemma-directed-morphism}
设 $\mathcal{C}$ 和 $\mathcal{D}$ 为位点，且
$u : \mathcal{C} \to \mathcal{D}$ 连续。假设第
\ref{sites-section-functoriality-PSh} 节中的所有范畴
$(\mathcal{I}_V^u)^{opp}$ 都是滤过的。那么 $u$ 定义位点态射
$\mathcal{D} \to \mathcal{C}$；换言之，$u_s$ 正合。
\end{lemma}

\begin{proof}
由于 $u_s$ 是 $u^s$ 的左伴随，$u_s$ 右正合；见 Categories，引理
\ref{categories-lemma-exact-adjoint}。因此只须证明 $u_s$ 左正合，换言之，
只须证明 $u_s$ 与有限极限交换。由于范畴
$\mathcal{I}_Y^{opp}$ 是滤过的，由 Categories，引理
\ref{categories-lemma-directed-commutes}，可知 $u_p$ 与有限极限交换（这里也
使用了 $\textit{PSh}$ 中极限的描述，见第
\ref{sites-section-limits-colimits-PSh} 节）。又因层化也与有限极限交换（引理
\ref{sites-lemma-sheafification-exact}），结合 $u_s = \# \circ u_p$ 即得结论。
\end{proof}

\begin{proposition}
\label{sites-proposition-get-morphism}
设 $\mathcal{C}$ 和 $\mathcal{D}$ 为位点。设
$u : \mathcal{C} \to \mathcal{D}$ 连续，并进一步假设：
\begin{enumerate}
\item 范畴 $\mathcal{C}$ 有终对象 $X$，且 $u(X)$ 是 $\mathcal{D}$ 的终对象；
\item 范畴 $\mathcal{C}$ 有纤维积，且 $u$ 与纤维积交换。
\end{enumerate}
那么 $u$ 定义位点态射 $\mathcal{D} \to \mathcal{C}$；换言之，$u_s$ 正合。
\end{proposition}

\begin{proof}
这由引理 \ref{sites-lemma-directed} 和 \ref{sites-lemma-directed-morphism} 得到。
\end{proof}

\begin{remark}
\label{sites-remark-explain-left-exact}
上述命题 \ref{sites-proposition-get-morphism} 的条件等价于说 $u$ 左正合，即与有限
极限交换。见 Categories，引理
\ref{categories-lemma-finite-limits-exist} 和
\ref{categories-lemma-characterize-left-exact}。用终对象和纤维积来表述似乎更
自然，因为在例子中，这种表述似乎有更多几何意义。
\end{remark}

\noindent
引理 \ref{sites-lemma-continuous-with-continuous-left-adjoint} 将提供另一种方法，
证明连续函子产生位点态射。

\begin{remark}
\label{sites-remark-quasi-continuous-morphism-sites}
（初读时可跳过。）设 $\mathcal{C}$ 和 $\mathcal{D}$ 为位点。仿照定义
\ref{sites-definition-morphism-sites}，我们说一个{\it 拟位点态射
$f : \mathcal{D} \to \mathcal{C}$}，由一个拟连续函子
$u : \mathcal{C} \to \mathcal{D}$（见备注
\ref{sites-remark-quasi-continuous}）给出，并要求 $u_s$ 正合。命题
\ref{sites-proposition-get-morphism} 在此情形下的类似命题，只须把“连续”替换为
“拟连续”，并把“态射”替换为“拟态射”。证明完全相同。
\end{remark}

\noindent
定义 \ref{sites-definition-morphism-sites} 中 $u_s$ 正合这一条件不能省略。例如，
若只假设 $u$ 连续，下面引理的结论未必成立。

\begin{lemma}
\label{sites-lemma-morphism-of-sites-covering}
设 $f : \mathcal{D} \to \mathcal{C}$ 是由函子
$u : \mathcal{C} \to \mathcal{D}$ 给出的位点态射。给定
$\mathcal{D}$ 的任意对象 $V$，存在覆盖 $\{V_j \to V\}$，使得对每个 $j$，
都存在态射 $V_j \to u(U_j)$，其中 $U_j$ 是 $\mathcal{C}$ 的某个对象。
\end{lemma}

\begin{proof}
由于 $f^{-1} = u_s$ 正合，有 $f^{-1}* = *$，其中 $*$ 表示层范畴的终对象
（例 \ref{sites-example-singleton-sheaf}）。由于 $f^{-1}* = u_s*$ 是 $u_p*$ 的
层化，可知存在覆盖 $\{V_j \to V\}$，使得 $(u_p*)(V_j)$ 非空。又因为
$(u_p*)(V_j)$ 是范畴 $\mathcal{I}^u_{V_j}$ 上的余极限，而该范畴的对象是态射
$V_j \to u(U)$，引理得证。
\end{proof}


























\section{托扑斯}
\label{sites-section-topoi}

\noindent
下面给出适合本文用途的托扑斯定义。托扑斯就是一个位点上的层范畴；要指定一个
托扑斯，只须指定该位点。托扑斯与位点之间真正的区别在于态射的定义。事实上，
有许多托扑斯态射并非来自底层位点的态射。

\begin{definition}[托扑斯]
\label{sites-definition-topos}
一个{\it 托扑斯}是某个位点 $\mathcal{C}$ 上的层范畴
$\Sh(\mathcal{C})$。
\begin{enumerate}
\item 设 $\mathcal{C}$、$\mathcal{D}$ 为位点。从
$\Sh(\mathcal{D})$ 到 $\Sh(\mathcal{C})$ 的一个{\it 托扑斯态射} $f$，
由一对函子
$f_* : \Sh(\mathcal{D}) \to \Sh(\mathcal{C})$ 和
$f^{-1} : \Sh(\mathcal{C}) \to \Sh(\mathcal{D})$ 给出，并满足：
\begin{enumerate}
\item 双函子自然地有
$$
\Mor_{\Sh(\mathcal{D})}(f^{-1}\mathcal{G}, \mathcal{F})
=
\Mor_{\Sh(\mathcal{C})}(\mathcal{G}, f_*\mathcal{F})
$$
\item 函子 $f^{-1}$ 与有限极限交换，即它左正合。
\end{enumerate}
\item 设 $\mathcal{C}$、$\mathcal{D}$、$\mathcal{E}$ 为位点。给定托扑斯态射
$f :\Sh(\mathcal{D}) \to \Sh(\mathcal{C})$ 和
$g :\Sh(\mathcal{E}) \to \Sh(\mathcal{D})$，{\it 复合
$f\circ g$} 是由函子
$(f \circ g)_* = f_* \circ g_*$ 和
$(f \circ g)^{-1} = g^{-1} \circ f^{-1}$ 定义的托扑斯态射。
\end{enumerate}
\end{definition}

\noindent
假设 $\alpha : \mathcal{S}_1 \to \mathcal{S}_2$ 是（可能为“大的”）范畴的
等价。若 $\mathcal{S}_1$、$\mathcal{S}_2$ 是托扑斯，则令
$f_* = \alpha$，并令 $f^{-1}$ 等于 $\alpha$ 的一个拟逆，便得到托扑斯态射
$f : \mathcal{S}_1 \to \mathcal{S}_2$。此外，这个态射是托扑斯
$2$-范畴中的等价（见第 \ref{sites-section-2-category} 节）。所以，若
$\mathcal{S}$ 等价于某个位点上的层范畴（而不一定等于某个位点上的层范畴），
说“$\mathcal{S}$ 是一个托扑斯”是有意义的。我们偶尔采用这种记号滥用。

\medskip\noindent
{\it 空托扑斯}是位点 $\mathcal{C}$ 上的层所成的托扑斯，其中
$\mathcal{C}$ 是空范畴。我们有时用 $\emptyset$ 表示这个位点。这个位点只有
唯一一个层（因为 $\emptyset$
没有对象）。所以 $\Sh(\emptyset)$ 等价于只有一个对象和一个态射的范畴。

\medskip\noindent
{\it 点托扑斯}是如下位点 $\mathcal{C}$ 上的层所成的托扑斯：它只有一个对象
$pt$ 和一个态射 $\text{id}_{pt}$，唯一的覆盖是
$\{\text{id}_{pt}\}$。我们简记这个位点为 $pt$。显然，层范畴
$ = $ 预层范畴 $ = $ 集合范畴。用公式表示：
$\Sh(pt) = \textit{Sets}$。

\medskip\noindent
设 $\mathcal{C}$ 和 $\mathcal{D}$ 为位点。设
$f : \Sh(\mathcal{D}) \to \Sh(\mathcal{C})$ 为托扑斯态射。注意，$f_*$
与所有极限交换，而 $f^{-1}$ 与所有余极限交换；见 Categories，引理
\ref{categories-lemma-adjoint-exact}。特别地，上述定义中对 $f^{-1}$ 的条件
保证 $f^{-1}$ 正合。托扑斯态射通常使用引理
\ref{sites-lemma-cocontinuous-morphism-topoi} 或下一个引理来构造。

\begin{lemma}
\label{sites-lemma-morphism-sites-topoi}
给定与函子 $u : \mathcal{C} \to \mathcal{D}$ 对应的位点态射
$f : \mathcal{D} \to \mathcal{C}$，函子对
$(f^{-1} = u_s, f_* = u^s)$ 是托扑斯态射
$\Sh(\mathcal{D}) \to \Sh(\mathcal{C})$。
\end{lemma}

\begin{proof}
这由定义 \ref{sites-definition-morphism-sites} 显然成立。
\end{proof}

\begin{remark}
\label{sites-remark-pt-topos}
有许多位点产生托扑斯 $\Sh(pt)$。下面是一个有用的例子。假设 $S$ 是一个
（由集合组成的）集合，并且至少含有一个非空元素。令 $\mathcal{S}$ 为这样的
范畴：对象是 $S$ 的元素，态射是任意集合映射。假设 $\mathcal{S}$ 有纤维积。
例如，若 $S = \mathcal{P}(\text{无限集})$ 是任意无限集合的幂集，则
此条件成立（集合论练习）。把满射映射族声明为覆盖，从而使
$\mathcal{S}$ 成为位点（并如 Sets，第 \ref{sets-section-coverings-site} 节
那样选择一个适当且充分大的覆盖族集合）。我们断言
$\Sh(\mathcal{S})$ 等价于集合范畴。

\medskip\noindent
先证明 $S$ 含有一个单元素集 $e \in S$ 的情形。在此情形下，存在托扑斯等价
$i : \Sh(pt) \to \Sh(\mathcal{S})$，由函子
\begin{equation}
\label{sites-equation-sheaves-pt-sets}
i^{-1}\mathcal{F} = \mathcal{F}(e), \quad
i_*E = (U \mapsto \Mor_{\textit{Sets}}(U, E))
\end{equation}
给出。事实上，假设 $\mathcal{F}$ 是 $\mathcal{S}$ 上的层。对任意
$U \in \Ob(\mathcal{S}) = S$，都可以找到覆盖
$\{\varphi_u : e \to U\}_{u \in U}$，其中 $\varphi_u$ 把 $e$ 的唯一元素
映到 $u \in U$。此时层条件蕴含
$\mathcal{F}(U) = \prod_{u \in U} \mathcal{F}(e)$。换言之，
$\mathcal{F}(U) = \Mor_{\textit{Sets}}(U, \mathcal{F}(e))$。此外，这条规则
与限制映射相容。因此函子
$$
i_* :
\textit{Sets} = \Sh(pt)
\longrightarrow
\Sh(\mathcal{S}), \quad
E \longmapsto (U \mapsto \Mor_{\textit{Sets}}(U, E))
$$
是范畴等价，其逆就是上面给出的函子 $i^{-1}$。

\medskip\noindent
若 $\mathcal{S}$ 不含单元素集，则上述定义的函子 $i_*$ 仍有意义。为证明它在
此情形下仍为等价，任取一个非空的 $\tilde e \in S$，并取一个像为单元素集的
映射 $\varphi : \tilde e \to \tilde e$。对任意层 $\mathcal{F}$，令
$$
\mathcal{F}(e) :=
\Im(
\mathcal{F}(\varphi) :
\mathcal{F}(\tilde e)
\longrightarrow
\mathcal{F}(\tilde e)
)
$$
并证明这给出 $i_*$ 的拟逆。细节从略。
\end{remark}

\begin{remark}
\label{sites-remark-morphism-topoi-big}
（与托扑斯态射有关的集合论问题。初读时可跳过。）上述定义的一个托扑斯态射
不是集合，而是类。换言之，它由数学公式而非数学对象给出。虽然可以考虑两个
给定托扑斯之间所有态射的汇集，但把它作为数学对象引入并非好主意。另一方面，
假设 $\mathcal{C}$ 和 $\mathcal{D}$ 是给定的位点。考虑函子
$\Phi : \mathcal{C} \to \Sh(\mathcal{D})$。这样的函子是集合，换言之，它
是数学对象。可以依次对 $\Phi$ 提出下列问题：
\begin{enumerate}
\item 给定 $\mathcal{D}$ 上的层 $\mathcal{F}$，规则
$U \mapsto \Mor_{\Sh(\mathcal{D})}(\Phi(U), \mathcal{F})$ 是否定义
$\mathcal{C}$ 上的层？若是，它定义函子
$\Phi_* : \Sh(\mathcal{D}) \to \Sh(\mathcal{C})$。
\item $\Phi_*$ 是否有左伴随？若有，把这个左伴随记作 $\Phi^{-1}$。
\item $\Phi^{-1}$ 是否正合？
\end{enumerate}
若最后一个问题的答案仍为“是”，便得到托扑斯态射
$(\Phi_*, \Phi^{-1})$。此外，给定任意托扑斯态射
$(f_*, f^{-1})$，可令 $\Phi(U) = f^{-1}(h_U^\#)$，并得到上述函子
$\Phi$，使得 $f_* \cong \Phi_*$ 且 $f^{-1} \cong \Phi^{-1}$（与伴随性质
相容）。结论是：通过使用 $\Phi$ 的汇集而非托扑斯态射，我们 (a) 把托扑斯
态射概念替换为数学对象，并且 (b) $\Phi$ 的汇集构成一个类（而不是类的
汇集）。当然，还可以说得更多，例如可以更准确地阐明上述条件 (2) 和 (3) 的
意义；我们在第 \ref{sites-section-points} 节对托扑斯的点讨论这一点。
\end{remark}

\begin{remark}
\label{sites-remark-quasi-continuous-morphism-topoi}
（初读时可跳过。）设 $\mathcal{C}$ 和 $\mathcal{D}$ 为位点。拟位点态射
$f : \mathcal{D} \to \mathcal{C}$（见备注
\ref{sites-remark-quasi-continuous-morphism-sites}）完全如引理
\ref{sites-lemma-morphism-sites-topoi} 那样，产生从
$\Sh(\mathcal{D})$ 到 $\Sh(\mathcal{C})$ 的托扑斯态射 $f$。
\end{remark}














\section{G-集合与态射}
\label{sites-section-G-sets-morphisms}

\noindent
设 $\varphi : G \to H$ 为群同态。按照例 \ref{sites-example-site-on-group} 和第
\ref{sites-section-example-sheaf-G-sets} 节，选择（适当的）位点
$\mathcal{T}_G$ 和 $\mathcal{T}_H$。令
$u : \mathcal{T}_H \to \mathcal{T}_G$ 为这样的函子：它把 $H$-集合 $U$
送到具有相同底层集合的 $G$-集合 $U_\varphi$，其 $G$-作用由
$g \cdot \xi = \varphi(g)\xi$ 定义，其中 $g \in G$ 且 $\xi \in U$。
显然，$u$ 与有限极限交换并且连续\footnote{集合论说明：先选择
$\mathcal{T}_H$，再选择 $\mathcal{T}_G$，使其包含 $u(\mathcal{T}_H)$，
并使 $\mathcal{T}_H$ 中的每个覆盖都对应于 $\mathcal{T}_G$ 中的覆盖。这可由
Sets，引理 \ref{sets-lemma-sets-with-group-action}、
\ref{sets-lemma-what-is-in-it-G-sets} 和
\ref{sets-lemma-coverings-site} 实现。}。应用命题
\ref{sites-proposition-get-morphism} 和引理 \ref{sites-lemma-morphism-sites-topoi}，得到
与 $\varphi$ 相伴的托扑斯态射
$$
f : \Sh(\mathcal{T}_G) \longrightarrow \Sh(\mathcal{T}_H)
$$
由命题 \ref{sites-proposition-sheaves-on-group}，得到一对伴随函子
$$
f_* : G\textit{-Sets} \longrightarrow H\textit{-Sets}, \quad
f^{-1} : H\textit{-Sets} \longrightarrow G\textit{-Sets}.
$$
下面具体求出这两个函子。

\medskip\noindent
先求 $f_*$ 的公式。回忆一下，给定 $G$-集合 $S$，$\mathcal{T}_G$ 上的相应
层 $\mathcal{F}_S$ 由规则
$\mathcal{F}_S(U) = \Mor_G(U, S)$ 给出。另一方面，给定
$\mathcal{T}_H$ 上的层 $\mathcal{G}$，相应的 $H$-集合由规则
$\mathcal{G}({}_HH)$ 给出。因此
$$
f_*S = \Mor_{G\textit{-Sets}}(({}_HH)_\varphi, S).
$$
进一步展开，得到
$$
f_*S = \{ a : H \to S \mid a(gh) = ga(h) \}
$$
其左 $H$-作用由
$(h \cdot a)(h') = a(h'h)$ 给出，其中 $a \in f_*S$ 任意。
% P01-ERRATA: P01-E0068 -- g is a G-element but a's argument lies in H; the formula omits \varphi(g).

\medskip\noindent
接着显式计算 $f^{-1}$。注意，由于 $\mathcal{T}_G$ 和
$\mathcal{T}_H$ 上的拓扑都是次典范的，所有可表预层都是层。此外，给定
$\mathcal{T}_H$ 的对象 $V$，可知 $f^{-1}h_V$ 等于 $h_{u(V)}$（见引理
\ref{sites-lemma-pullback-representable-sheaf}）。所以，对可表层有
$f^{-1}S = S_\varphi$。由于 $\mathcal{T}_H$ 上的每个层都是可表层的余积，
可知一般情形下也成立。因此，对任意 $H$-集合 $T$，有
$$
f^{-1}T = T_\varphi.
$$
$f^{-1}$ 与 $f_*$ 之间的伴随性由公式
$$
\Mor_{G\textit{-Sets}}(T_\varphi, S) =
\Mor_{H\textit{-Sets}}(T, f_*S)
$$
体现，其中 $f_*S$ 如上。这也可以直接证明。此外，由此显然可知
$(f^{-1}, f_*)$ 构成伴随对，且 $f^{-1}$ 正合。因此，也可以不用上述构造，
而直接以这种方式定义托扑斯态射
$f : G\textit{-Sets} \to H\textit{-Sets}$。














\section{拟紧对象与余极限}
\label{sites-section-quasi-compact}

\noindent
为能使用与拓扑空间情形相同的语言，我们引入如下术语。

\begin{definition}
\label{sites-definition-quasi-compact}
设 $\mathcal{C}$ 为位点。若对 $\mathcal{C}$ 中任一覆盖
$\mathcal{U} = \{U_i \to U\}_{i \in I}$，都存在另一个覆盖
$\mathcal{V} = \{V_j \to U\}_{j \in J}$，以及一个由
$\text{id} : U \to U$、$\alpha : J \to I$ 和
$V_j \to U_{\alpha(j)}$ 给出的具有固定靶的态射族之态射
$\mathcal{V} \to \mathcal{U}$（见定义
\ref{sites-definition-morphism-coverings}），使得 $\alpha$ 的像是 $I$ 的有限
子集，则称 $\mathcal{C}$ 的对象 $U$ 是{\it 拟紧的}。
\end{definition}

\noindent
当然，通常的概念足以推出 $U$ 拟紧。

\begin{lemma}
\label{sites-lemma-conclude-quasi-compact}
设 $\mathcal{C}$ 为位点，$U$ 为 $\mathcal{C}$ 的对象。考虑下列条件：
\begin{enumerate}
\item $U$ 拟紧；
\item 对 $\mathcal{C}$ 中每个覆盖 $\{U_i \to U\}_{i \in I}$，都存在
$\mathcal{C}$ 中加细 $\mathcal{U}$ 的有限覆盖
$\{V_j \to U\}_{j = 1, \ldots, m}$；
\item 对 $\mathcal{C}$ 中每个覆盖 $\{U_i \to U\}_{i \in I}$，都存在
有限子集 $I' \subset I$，使得
$\{U_i \to U\}_{i \in I'}$ 是 $\mathcal{C}$ 中的覆盖。
\end{enumerate}
则总有 (3) $\Rightarrow$ (2) $\Rightarrow$ (1)，但反向蕴含一般不成立。
\end{lemma}

\begin{proof}
这些蕴含由定义立即可得。令 $X = [0, 1] \subset \mathbf{R}$ 为拓扑空间
（带通常的 $\epsilon$-$\delta$ 拓扑）。令 $\mathcal{C}$ 为 $X$ 的开子空间
所成的范畴，以包含映射为态射，并取通常的开覆盖（与例
\ref{sites-example-site-topological} 比较）。不过，我们随后改变 $\mathcal{C}$ 中
覆盖的概念：除形如 $\{U \to U\}$ 的覆盖外，排除所有有限覆盖。容易看出，
这会成为定义 \ref{sites-definition-site} 所述的位点。在这个位点中，对象
$X = U = [0, 1]$ 在定义 \ref{sites-definition-quasi-compact} 的意义下拟紧，但
$U$ 不满足 (2)。留给读者构造满足 (2) 而不满足 (3) 的例子。
\end{proof}

\noindent
下面给出拟紧对象的托扑斯论含义。

\begin{lemma}
\label{sites-lemma-quasi-compact}
设 $\mathcal{C}$ 为位点，$U$ 为 $\mathcal{C}$ 的对象。下列条件等价：
\begin{enumerate}
\item $U$ 拟紧；
\item 对每个层的满射
$\coprod_{i \in I} \mathcal{F}_i \to h_U^\#$，都存在有限子集
$J \subset I$，使得
$\coprod_{i \in J} \mathcal{F}_i \to h_U^\#$ 是满射。
\end{enumerate}
\end{lemma}

\begin{proof}
假设 (1) 成立，并设 $\coprod_{i \in I} \mathcal{F}_i \to h_U^\#$ 为满射。
则 $\text{id}_U$ 是 $h_U^\#$ 在 $U$ 上的截面。因此存在覆盖
$\{U_a \to U\}_{a \in A}$，并且对每个 $a \in A$，存在
$\coprod_{i \in I} \mathcal{F}_i$ 在 $U_a$ 上的截面 $s_a$，它映到
$\text{id}_U$。根据把预层余积层化以构造余积的方法（引理
\ref{sites-lemma-colimit-sheaves}），对每个 $a$ 都存在覆盖
$\{U_{ab} \to U_a\}_{b \in B_a}$；对所有 $b \in B_a$，都存在
$\iota(b) \in I$ 以及 $\mathcal{F}_{\iota(b)}$ 在 $U_{ab}$ 上的截面
$s_{b}$，它映到 $\text{id}_U|_{U_{ab}}$。因此，用
$\{U_{ab} \to U\}_{a \in A, b \in B_a}$ 代替覆盖
$\{U_a \to U\}_{a \in A}$ 后，可以假设有映射 $\iota : A \to I$，
且对每个 $a \in A$，有 $\mathcal{F}_{\iota(a)}$ 在 $U_a$ 上的截面
$s_a$，它映到 $\text{id}_U$。
由于 $U$ 拟紧，存在覆盖 $\{V_c \to U\}_{c \in C}$、像为有限集的映射
$\alpha : C \to A$，以及 $U$ 上的态射
$V_c \to U_{\alpha(c)}$。于是 $J = \Im(\iota \circ \alpha) \subset I$
满足要求，因为
$\coprod_{c \in C} h_{V_c}^\# \to h_U^\#$ 是满射（引理
\ref{sites-lemma-covering-surjective-after-sheafification}），并且经由
$\coprod_{i \in J} \mathcal{F}_i \to h_U^\#$ 分解。（这里用到复合
$h_{V_c}^\# \to h_{U_{\alpha(c)}}
\xrightarrow{s_{\alpha(c)}} \mathcal{F}_{\iota(\alpha(c))} \to h_U^\#$
就是由态射 $V_c \to U$ 诱导的映射 $h_{V_c}^\# \to h_U^\#$，因为
$s_{\alpha(c)}$ 映到 $\text{id}_U|_{U_{\alpha(c)}}$。）
% P01-ERRATA: P01-E0069 -- the intermediate representable is missing sheafification marker \#.

\medskip\noindent
假设 (2) 成立。令 $\{U_i \to U\}_{i \in I}$ 为覆盖。由引理
\ref{sites-lemma-covering-surjective-after-sheafification} 可知
$\coprod_{i \in I} h_{U_i}^\# \to h_U^\#$ 是满射。因此可找到有限子集
$J \subset I$，使得
$\coprod_{j \in J} h_{U_j}^\# \to h_U^\#$ 是满射。再按上述论证，得到
$\mathcal{C}$ 中 $U$ 的覆盖 $\{V_c \to U\}_{c \in C}$ 和映射
$\iota : C \to J$，使得 $\text{id}_U$ 可提升为
$h_{U_{\iota(c)}}^\#$ 在 $V_c$ 上的截面 $s_c$。
% P01-ERRATA: P01-E0070 -- source says "a section of s_c of"; translated by its only type-correct reading.
进一步加细该覆盖后，可以假设
$s_c \in h_{U_{\iota(c)}}(V_c)$ 映到 $\text{id}_U$。于是
$s_c : V_c \to U_{\iota(c)}$ 是 $U$ 上的态射，结论成立。
\end{proof}

\noindent
上述引理启发如下定义。

\begin{definition}
\label{sites-definition-quasi-compact-topos}
若对 $\Sh(\mathcal{C})$ 中任一满射
$\coprod_{i \in I} \mathcal{F}_i \to \mathcal{F}$，都存在有限子集
$J \subset I$，使得
$\coprod_{i \in J} \mathcal{F}_i \to \mathcal{F}$ 是满射，则称对象
$\mathcal{F}$ 是{\it 拟紧的}。若托扑斯
$\Sh(\mathcal{C})$ 的终对象 $*$ 是拟紧对象，则称
$\Sh(\mathcal{C})$ 是{\it 拟紧的}。
\end{definition}

\noindent
由引理 \ref{sites-lemma-quasi-compact}，若位点 $\mathcal{C}$ 有终对象 $X$，则
$\Sh(\mathcal{C})$ 拟紧当且仅当 $X$ 拟紧。

\begin{lemma}
\label{sites-lemma-image-of-quasi-compact}
\begin{slogan}
层的满射传递拟紧性。
\end{slogan}
设 $\mathcal{C}$ 为位点。
\begin{enumerate}
\item 若 $U \to V$ 是 $\mathcal{C}$ 的态射，且
$h_U^\# \to h_V^\#$ 是满射，而 $U$ 拟紧，则 $V$ 拟紧。
\item 若 $\mathcal{F} \to \mathcal{G}$ 是集合层的满射，且
$\mathcal{F}$ 拟紧，则 $\mathcal{G}$ 拟紧。
\end{enumerate}
\end{lemma}

\begin{proof}
略。
\end{proof}

\begin{lemma}
\label{sites-lemma-coproduct-quasi-compact}
\begin{slogan}
层的有限余积保持拟紧性。
\end{slogan}
% P01-ERRATA: P01-E0071 -- source slogan misspells "coproducts" as "copoducts".
设 $\mathcal{C}$ 为位点。若 $n \geq 1$，且
$\mathcal{F}_1, \ldots, \mathcal{F}_n$ 是 $\mathcal{C}$ 上的拟紧层，则
$\coprod_{i = 1, \ldots, n} \mathcal{F}_i$ 拟紧。
\end{lemma}

\begin{proof}
略。
\end{proof}

\noindent
下列两个引理是 Sheaves，引理
\ref{sheaves-lemma-directed-colimits-sections} 对位点的类似结果。

\begin{lemma}
\label{sites-lemma-directed-colimits-sections}
设 $\mathcal{C}$ 为位点，且
$\mathcal{I} \to \Sh(\mathcal{C})$，$i \mapsto \mathcal{F}_i$
是集合层的滤过图表。设 $U \in \Ob(\mathcal{C})$。考虑典范映射
$$
\Psi :
\colim_i \mathcal{F}_i(U)
\longrightarrow
\left(\colim_i \mathcal{F}_i\right)(U)
$$
采用上面引入的术语：
\begin{enumerate}
\item 若所有转移映射都是单射，则对任意 $U$，$\Psi$ 都是单射。
\item 若 $U$ 拟紧，则 $\Psi$ 是单射。
\item 若 $U$ 拟紧，且所有转移映射都是单射，则 $\Psi$ 是同构。
\item 若 $U$ 有一个由覆盖 $\{U_j \to U\}_{j \in J}$ 组成的共尾系，
其中 $J$ 有限，且对所有 $j, j' \in J$，
$U_j \times_U U_{j'}$ 都拟紧，则 $\Psi$ 是双射。
\end{enumerate}
\end{lemma}

\begin{proof}
假设所有转移映射都是单射。此时预层
$\mathcal{F}' : V \mapsto \colim_i \mathcal{F}_i(V)$ 是分离的（见定义
\ref{sites-definition-separated}）。由引理 \ref{sites-lemma-colimit-sheaves} 有
$(\mathcal{F}')^\# = \colim_i \mathcal{F}_i$。由定理
\ref{sites-theorem-plus} 可知
$\mathcal{F}' \to (\mathcal{F}')^\#$ 是单射。这证明了 (1)。

\medskip\noindent
假设 $U$ 拟紧。设 $s \in \mathcal{F}_i(U)$ 和
$s' \in \mathcal{F}_{i'}(U)$ 在左边给出的元素经 $\Psi$ 后具有相同的像。
这意味着可以选择覆盖 $\{U_a \to U\}_{a \in A}$，并且对每个
$a \in A$ 选择指标 $i_a \in I$，其中 $i_a \geq i$、$i_a \geq i'$，
使得 $\varphi_{ii_a}(s) = \varphi_{i'i_a}(s')$。因为 $U$ 拟紧，可以选择
覆盖 $\{V_b \to U\}_{b \in B}$、像为有限集的映射
$\alpha : B \to A$，以及 $U$ 上的态射
$V_b \to U_{\alpha(b)}$。选择 $i''\in I$，使它不小于（即 $\geq$）所有
$i_{\alpha(b)}$；这是可能的，因为 $\alpha$ 的像有限。于是对所有
$b \in B$，$\varphi_{ii''}(s)$ 和 $\varphi_{i'i''}(s)$ 在 $V_b$ 上相等，从而
$\varphi_{ii''}(s) = \varphi_{i'i''}(s)$。这证明了 (2)。
% P01-ERRATA: P01-E0072 -- both second terms use s but the setup requires s'; formula preserved.

\medskip\noindent
假设 $U$ 拟紧，且所有转移映射都是单射。令 $s$ 为 $\Psi$ 的靶中的元素。
存在覆盖 $\{U_a \to U\}_{a \in A}$；对每个 $a \in A$，存在指标
$i_a \in I$ 和截面 $s_a \in \mathcal{F}_{i_a}(U_a)$，使得对所有
$a \in A$，$s|_{U_a}$ 都来自 $s_a$。因为 $U$ 拟紧，可以选择覆盖
$\{V_b \to U\}_{b \in B}$、像为有限集的映射
$\alpha : B \to A$，以及 $U$ 上的态射
$V_b \to U_{\alpha(b)}$。选择 $i \in I$，使它不小于（即 $\geq$）所有
$i_{\alpha(b)}$；这是可能的，因为 $\alpha$ 的像有限。由 (1)，截面
$s_b = \varphi_{i_{\alpha(b)} i}(s_{\alpha(b)})|_{V_b}$
在 $V_b \times_U V_{b'}$ 上相符。因此它们粘合为截面
$s' \in \mathcal{F}_i(U)$，该截面经 $\Psi$ 映到 $s$。这证明了 (3)。

\medskip\noindent
假设 (4) 的假设成立。由引理 \ref{sites-lemma-conclude-quasi-compact}，对象
$U$ 拟紧，故由 (2)，$\Psi$ 是单射。为证明满射性，令 $s$ 为 $\Psi$ 的靶中
的元素。根据假设，存在有限覆盖
$\{U_j \to U\}_{j = 1, \ldots, m}$，使得对所有
$1 \leq j, j' \leq m$，$U_j \times_U U_{j'}$ 都拟紧；并且对每个 $j$
存在指标 $i_j \in I$ 和 $s_j \in \mathcal{F}_{i_j}(U_j)$，使得
$s|_{U_j}$ 是 $s_j$ 的像，对所有 $j$ 都成立。由于
$U_j \times_U U_{j'}$ 拟紧，可以应用
(2)，从而存在 $i_{jj'} \in I$，满足 $i_{jj'} \geq i_j$、
$i_{jj'} \geq i_{j'}$，并使
$\varphi_{i_ji_{jj'}}(s_j)$ 与
$\varphi_{i_{j'}i_{jj'}}(s_{j'})$ 在
$U_j \times_U U_{j'}$ 上相等。选择指标 $i \in I$，使其不小于所有
$i_{jj'}$。于是 $\mathcal{F}_i$ 的截面
$\varphi_{i_ji}(s_j)$ 粘合成 $\mathcal{F}_i$ 在 $U$ 上的截面。
该截面按要求映到元素 $s$。
\end{proof}

\begin{lemma}
\label{sites-lemma-directed-colimits-global-sections}
设 $\mathcal{C}$ 为位点，且
$\mathcal{I} \to \Sh(\mathcal{C})$，$i \mapsto \mathcal{F}_i$
是集合层的滤过图表。考虑典范映射
$$
\Psi :
\colim_i \Gamma(\mathcal{C}, \mathcal{F}_i)
\longrightarrow
\Gamma(\mathcal{C}, \colim_i \mathcal{F}_i)
$$
有下列结论：
\begin{enumerate}
\item 若所有转移映射都是单射，则 $\Psi$ 是单射。
\item 若 $\Sh(\mathcal{C})$ 拟紧，则 $\Psi$ 是单射。
\item 若 $\Sh(\mathcal{C})$ 拟紧，且所有转移映射都是单射，则
$\Psi$ 是同构。
\item 假设存在集合 $S \subset \Ob(\Sh(\mathcal{C}))$，具有如下性质：
\begin{enumerate}
\item 对每个满射 $\mathcal{F} \to *$，都存在
$\mathcal{K} \in S$ 和映射 $\mathcal{K} \to \mathcal{F}$，使得
$\mathcal{K} \to *$ 是满射；
\item 对 $\mathcal{K} \in S$，积
$\mathcal{K} \times \mathcal{K}$ 拟紧。
\end{enumerate}
则 $\Psi$ 是双射。
\end{enumerate}
\end{lemma}

\begin{proof}
(1) 的证明。假设所有转移映射都是单射。此时预层
$\mathcal{F}' : V \mapsto \colim_i \mathcal{F}_i(V)$ 是分离的（见定义
\ref{sites-definition-separated}）。由引理 \ref{sites-lemma-colimit-sheaves} 有
$(\mathcal{F}')^\# = \colim_i \mathcal{F}_i$。由定理
\ref{sites-theorem-plus} 可知
$\mathcal{F}' \to (\mathcal{F}')^\#$ 是单射。这证明了 (1)。

\medskip\noindent
(2) 的证明。假设 $\Sh(\mathcal{C})$ 拟紧。回忆一下，对
$\Sh(\mathcal{C})$ 中的所有 $\mathcal{F}$，都有
$\Gamma(\mathcal{C}, \mathcal{F}) = \Mor(*, \mathcal{F})$。
令 $a_i, b_i : * \to \mathcal{F}_i$；对 $i' \geq i$，以
$a_{i'}, b_{i'} : * \to \mathcal{F}_{i'}$ 表示它们与系统转移映射的复合。
令 $a = \colim_{i' \geq i} a_{i'}$，并对 $b$ 作类似定义。对
$i' \geq i$，记
$$
E_{i'} = \text{Equalizer}(a_{i'}, b_{i'}) \subset *
\quad\text{且}\quad
E = \text{Equalizer}(a, b) \subset *
$$
由 Categories，引理 \ref{categories-lemma-directed-commutes}，有
$E = \colim_{i' \geq i} E_{i'}$。从而
$\coprod_{i' \geq i} E_{i'} \to E$ 是层的满射。因此，若 $E = *$，
即 $a = b$，则因为 $*$ 拟紧，可知对某个 $i' \geq i$ 有
$E_{i'} = *$，进而对某个 $i' \geq i$ 有 $a_{i'} = b_{i'}$。
这证明了 (2)。

\medskip\noindent
(3) 的证明。假设 $\Sh(\mathcal{C})$ 拟紧，且所有转移映射都是单射。
令 $a : * \to \colim \mathcal{F}_i$ 为映射。则
$E_i = a^{-1}(\mathcal{F}_i) \subset *$ 是子层，并且
$\colim E_i = *$（由上面的引用）。故对某个 $i$ 有 $E_i = *$，从而
$a$ 的像按要求包含在 $\mathcal{F}_i$ 中。

\medskip\noindent
(4) 的证明。设 $S \subset \Ob(\Sh(\mathcal{C}))$ 满足 (4)(a)、(b)。
把 (4)(a) 应用于 $\text{id} : * \to *$，可找到
$\mathcal{K} \in S$，使得 $\mathcal{K} \to *$ 是满射。映射
$\mathcal{K} \times \mathcal{K} \to \mathcal{K} \to *$ 都是满射。
由 (4)(b) 和引理 \ref{sites-lemma-image-of-quasi-compact}，可知
$\mathcal{K}$ 与 $\Sh(\mathcal{C})$ 都拟紧。因此由 (2)，$\Psi$ 是单射。
令 $\mathcal{F} = \colim \mathcal{F}_i$，并令
$s : * \to \mathcal{F}$ 为该余极限的整体截面。由于
$\coprod \mathcal{F}_i \to \mathcal{F}$ 是满射，投影
$$
\coprod\nolimits_{i \in I} * \times_{s, \mathcal{F}} \mathcal{F}_i \to *
$$
是满射。由 (4)(a)，得到 $\mathcal{K} \in S$ 和映射
$\mathcal{K} \to \coprod_{i \in I} * \times_{s, \mathcal{F}} \mathcal{F}_i$，
其中 $\mathcal{K} \to *$ 是满射。因为 $\mathcal{K}$ 拟紧，对某个有限子集
$I' \subset I$，可得到分解
$\mathcal{K} \to \coprod_{i' \in I'} * \times_{s, \mathcal{F}}
\mathcal{F}_{i'}$。令 $i \in I$ 为有限子集 $I'$ 的一个上界。转移映射定义
映射 $\coprod_{i' \in I'} \mathcal{F}_{i'} \to \mathcal{F}_i$，后者又产生
映射 $\mathcal{K} \to * \times_{s, \mathcal{F}} \mathcal{F}_i$。换言之，
得到 $\mathcal{K} \in S$，其中 $\mathcal{K} \to *$ 是满射，并有交换图
$$
\xymatrix{
\mathcal{K} \times \mathcal{K} \ar@<1ex>[r] \ar@<-1ex>[r] &
\mathcal{K} \ar[d] \ar[r] & {*} \ar[d]^s \\
& \mathcal{F}_i \ar[r] & \mathcal{F} \ar@{=}[r] & \colim \mathcal{F}_i
}
$$
注意该图的顶行是余等化子。因此，只需证明增大 $i$ 后，两个诱导映射
$a_i, b_i : \mathcal{K} \times \mathcal{K} \to \mathcal{F}_i$
相等即可。下一段使用与 (2) 的证明完全相同的论证来说明这一点；建议读者
跳过证明的剩余部分。
% P01-ERRATA: P01-E0073 -- source has the duplicated phrase "This is done shown".

\medskip\noindent
对 $i' \geq i$，以
$a_{i'}, b_{i'} : \mathcal{K} \times \mathcal{K} \to \mathcal{F}_{i'}$
表示 $a_i, b_i$ 与系统转移映射的复合。令
$a = \colim_{i' \geq i} a_{i'} :
\mathcal{K} \times \mathcal{K} \to \mathcal{F}$，并对 $b$ 作类似定义。
由上图的交换性，有 $a = b$。对 $i' \geq i$，记
$$
E_{i'} = \text{Equalizer}(a_{i'}, b_{i'}) \subset
\mathcal{K} \times \mathcal{K}
\quad\text{且}\quad
E = \text{Equalizer}(a, b) \subset
\mathcal{K} \times \mathcal{K}
$$
由 Categories，引理 \ref{categories-lemma-directed-commutes}，有
$E = \colim_{i' \geq i} E_{i'}$。从而
$\coprod_{i' \geq i} E_{i'} \to E$ 是层的满射。因为 $a = b$，有
$E = \mathcal{K} \times \mathcal{K}$。又因为由 (4)(b)，
$\mathcal{K} \times \mathcal{K}$ 拟紧，所以对某个 $i' \geq i$ 有
$E_{i'} = \mathcal{K} \times \mathcal{K}$，进而对某个 $i' \geq i$ 有
$a_{i'} = b_{i'}$。
\end{proof}

\begin{remark}
\label{sites-remark-stronger-conditions}
设 $\mathcal{C}$ 为位点。可以用多种方式保证引理
\ref{sites-lemma-directed-colimits-global-sections} 第 (4) 部分的假设成立。下面
列出几种。
\begin{enumerate}
\item 假设存在集合 $\mathcal{B} \subset \Ob(\mathcal{C})$，具有如下性质：
\begin{enumerate}
\item 对每个满射 $\mathcal{F} \to *$，都存在 $m \geq 0$ 和
$U_1, \ldots, U_m \in \mathcal{B}$，使得 $\mathcal{F}(U_j)$ 非空，且
$\coprod h_{U_j}^\# \to *$ 是满射；
\item 对 $U, U' \in \mathcal{B}$，层
$h_U^\# \times h_{U'}^\#$ 拟紧。
\end{enumerate}
\item 假设存在集合 $\mathcal{B} \subset \Ob(\mathcal{C})$，具有如下性质：
\begin{enumerate}
\item 存在 $m \geq 0$ 和 $U_1, \ldots, U_m \in \mathcal{B}$，使得
$\coprod h_{U_j}^\# \to *$ 是满射；
\item 对 $U \in \mathcal{B}$，$U$ 的任意覆盖都可由有限覆盖
$\{U_j \to U\}_{j = 1, \ldots, m}$ 加细，其中
$U_j \in \mathcal{B}$；
\item 对 $U, U' \in \mathcal{B}$，存在 $m \geq 0$、
$U_1, \ldots, U_m \in \mathcal{B}$，以及态射
$U_j \to U$ 和 $U_j \to U'$，使得
$\coprod h_{U_j}^\# \to h_U^\# \times h_{U'}^\#$ 是满射。
\end{enumerate}
\item 假设：
\begin{enumerate}
\item $\Sh(\mathcal{C})$ 拟紧；
\item $\mathcal{C}$ 的每个对象都有一个覆盖，其成员均为拟紧对象；
\item 若 $U$ 和 $U'$ 拟紧，则层
$h_U^\# \times h_{U'}^\#$ 拟紧。
\end{enumerate}
\end{enumerate}
在情形 (1) 和 (2) 中，令 $S \subset \Ob(\Sh(\mathcal{C}))$ 为由
$U \in \mathcal{B}$ 的层 $h_U^\#$ 的有限余积组成的集合。在情形 (3) 中，
令 $S \subset \Ob(\Sh(\mathcal{C}))$ 为由拟紧对象
$U \in \Ob(\mathcal{C})$ 的层 $h_U^\#$ 的有限余积组成的集合。
\end{remark}

\noindent
以后将需要如下关于余极限可能情形的界。

\begin{lemma}
\label{sites-lemma-colimit-over-ordinal-sections}
设 $\mathcal{C}$ 为位点，$\beta$ 为序数。设
$\beta \to \Sh(\mathcal{C})$，$\alpha \mapsto \mathcal{F}_\alpha$
为 $\beta$ 上的层系。对 $U \in \Ob(\mathcal{C})$，考虑典范映射
$$
\colim_{\alpha < \beta} \mathcal{F}_\alpha(U)
\longrightarrow
\left(\colim_{\alpha < \beta} \mathcal{F}_\alpha \right)(U)
$$
若 $\beta$ 的共尾性足够大，则该映射对所有 $U$ 都是双射。
\end{lemma}

\begin{proof}
左边是预层范畴中所取余极限 $\mathcal{F}_{\colim}$ 在 $U$ 上的值，见第
\ref{sites-section-limits-colimits-PSh} 节。回忆一下，
$\colim_{\alpha < \beta} \mathcal{F}_\alpha$ 是
$\mathcal{F}_{\colim}$ 的层化 $\mathcal{F}_{\colim}^\#$，见引理
\ref{sites-lemma-colimit-sheaves}。令
$\mathcal{U} = \{U_i \to U\}_{i \in I}$ 为 $\mathcal{C}$ 的覆盖所成集合
$\text{Cov}(\mathcal{C})$ 的元素。若 $\beta$ 的共尾性大于 $I$ 的基数，
则断言
$$
H^0(\mathcal{U}, \mathcal{F}_{\colim}) =
\colim H^0(\mathcal{U}, \mathcal{F}_\alpha) =
\colim \mathcal{F}_\alpha(U) = \mathcal{F}_{\colim}(U)
$$
第二、第三个等号显然。对第一个等号，设
$s = (s_i) \in H^0(\mathcal{U}, \mathcal{F}_{\colim})$。则对每个 $i$，
元素 $s_i$ 来自某个 $\alpha_i < \beta$ 的元素
$s_{i, \alpha_i} \in \mathcal{F}_{\alpha_i}(U_i)$。由共尾性假设，可选择
与 $i$ 无关的 $\alpha_i = \alpha$。
% P01-ERRATA: P01-E0074 -- the source elides replacing all alpha_i by a common upper bound alpha.
于是对某个 $\alpha_{i, j} < \beta$，$s_i$ 和 $s_j$ 映到
$\mathcal{F}_{\alpha_{i, j}}(U_i \times_U U_j)$ 中的同一元素。由于
$I \times I$ 的基数也小于 $\beta$ 的共尾性，增大 $\alpha$ 后，可以假设
对所有 $i, j$ 都有 $\alpha_{i, j} = \alpha$。这证明自然映射
$\colim H^0(\mathcal{U}, \mathcal{F}_\alpha)
\to H^0(\mathcal{U}, \mathcal{F}_{\colim})$ 是满射。非常类似的论证表明
它是单射。特别地，$\mathcal{F}_{\colim}$ 对 $\mathcal{U}$ 满足层条件。
因此，若 $\beta$ 的共尾性大于所有覆盖的指标集 $I$ 的基数所成集合的上确界，
便得到结论。
% P01-ERRATA: P01-E0075 -- source has "cardinality if I x I" instead of "of".
\end{proof}







\section{位点的余极限}
\label{sites-section-colimit-sites}

\noindent
当位点是一个位点逆系统的极限时，我们需要引理
\ref{sites-lemma-directed-colimits-sections} 的类似结果。为简单起见，只说明覆盖的
指标集有限时的构造。

\begin{situation}
\label{sites-situation-inverse-limit-sites}
这里给定：
\begin{enumerate}
\item 余滤过指标范畴 $\mathcal{I}$；
\item 对 $i \in \Ob(\mathcal{I})$，给定位点 $\mathcal{C}_i$，使得
$\mathcal{C}_i$ 中每个覆盖都有有限指标集；
\item 对 $\mathcal{I}$ 中的态射 $a : i \to j$，给定由连续函子
$u_a : \mathcal{C}_j \to \mathcal{C}_i$ 所确定的位点态射
$f_a : \mathcal{C}_i \to \mathcal{C}_j$。
\end{enumerate}
并且只要 $\mathcal{I}$ 中 $c = a \circ b$，就有
$f_a \circ f_b = f_c$。
\end{situation}

\begin{lemma}
\label{sites-lemma-colimit-sites}
在情形 \ref{sites-situation-inverse-limit-sites} 中，可以如下构造位点
$(\mathcal{C}, \text{Cov}(\mathcal{C}))$：
\begin{enumerate}
\item 作为范畴，$\mathcal{C} = \colim \mathcal{C}_i$；
\item $\text{Cov}(\mathcal{C})$ 是各个
$\text{Cov}(\mathcal{C}_i)$ 在
$u_i : \mathcal{C}_i \to \mathcal{C}$ 下之像的并集。
\end{enumerate}
\end{lemma}

\begin{proof}
我们对位点态射复合的定义蕴含：只要 $\mathcal{I}$ 中
$c = a \circ b$，就有 $u_b \circ u_a = u_c$。公式
$\mathcal{C} = \colim \mathcal{C}_i$ 意味着
$\Ob(\mathcal{C}) = \colim \Ob(\mathcal{C}_i)$ 且
$\text{Arrows}(\mathcal{C}) = \colim \text{Arrows}(\mathcal{C}_i)$。
于是源、靶和复合都从各个 $\text{Arrows}(\mathcal{C}_i)$ 上的源、靶和
复合继承而来。这样便得到一个范畴。以
$u_i : \mathcal{C}_i \to \mathcal{C}$ 表示显然的函子。注意，给定
$\mathcal{C}$ 中任意有限图表，都存在某个 $i$，使该图表是
$\mathcal{C}_i$ 中一个图表的像。

\medskip\noindent
令 $\{U^t \to U\}$ 为 $\mathcal{C}$ 的覆盖。先证明：若
$V \to U$ 是 $\mathcal{C}$ 的态射，则 $U^t \times_U V$ 存在。由上述说明
和覆盖的定义，可以找到 $i$、$\mathcal{C}_i$ 的覆盖
$\{U_i^t \to U_i\}$ 以及态射 $V_i \to U_i$，它们在 $u_i$ 下的像就是
给定数据。断言 $U^t \times_U V$ 是
$U^t_i \times_{U_i} V_i$ 在 $u_i$ 下的像。事实上，对
$\mathcal{I}$ 中每个 $a : j \to i$，函子 $u_a$ 连续，所以
$u_a(U^t_i \times_{U_i} V_i) = u_a(U^t_i) \times_{u_a(U_i)} u_a(V_i)$。
特别地，如有需要，可以用 $j$ 代替 $i$。因此，若 $W$ 是
$\mathcal{C}$ 的另一个对象，则可以假设 $W = u_i(W_i)$，并得到
\begin{align*}
& \Mor_\mathcal{C}(W, u_i(U^t_i \times_{U_i} V_i)) \\
& =
\colim_{a : j \to i}
\Mor_{\mathcal{C}_j}(u_a(W_i), u_a(U^t_i \times_{U_i} V_i)) \\
& =
\colim_{a : j \to i}
\Mor_{\mathcal{C}_j}(u_a(W_i), u_a(U^t_i))
\times_{\Mor_{\mathcal{C}_j}(u_a(W_i), u_a(U_i))}
\Mor_{\mathcal{C}_j}(u_a(W_i), u_a(V_i)) \\
& =
\Mor_\mathcal{C}(W, U^t)
\times_{\Mor_\mathcal{C}(W, U)}
\Mor_\mathcal{C}(W, V)
\end{align*}
因为滤过余极限与有限极限交换（Categories，引理
\ref{categories-lemma-directed-commutes}）。同时可知
$\{U^t \times_U V \to V\}$ 是 $\mathcal{C}$ 中的覆盖。这样就验证了
定义 \ref{sites-definition-site} 的公理 (3)。

\medskip\noindent
为验证定义 \ref{sites-definition-site} 的公理 (2)，令
$\{U^t \to U\}_{t \in T}$ 为 $\mathcal{C}$ 的覆盖，并且对每个 $t$，令
$\{U^{ts} \to U^t\}$ 为 $\mathcal{C}$ 的覆盖。可以找到 $i$ 和
$\mathcal{C}_i$ 的覆盖 $\{U^t_i \to U_i\}_{t \in T}$，其在 $u_i$ 下的像
是 $\{U^t \to U\}$。因为 $T$ 是{\bf 有限的}，可以选择
$\mathcal{I}$ 中的 $a : j \to i$ 和 $\mathcal{C}_j$ 中的覆盖
$\{U^{ts}_j \to u_a(U^t_i)\}$，使它们在 $u_j$ 下的像给出
$\{U^{ts} \to U^t\}$。于是，把公理 (2) 应用于位点 $\mathcal{C}_j$，
可知 $\{U^{ts} \to U\}$ 是 $\mathcal{C}$ 的覆盖。

\medskip\noindent
定义 \ref{sites-definition-site} 的公理 (1) 的证明从略。
\end{proof}

\begin{lemma}
\label{sites-lemma-compute-pullback-to-limit}
在情形 \ref{sites-situation-inverse-limit-sites} 中，令
$u_i : \mathcal{C}_i \to \mathcal{C}$ 如引理
\ref{sites-lemma-colimit-sites} 中所构造。则 $u_i$ 定义位点态射
$f_i : \mathcal{C} \to \mathcal{C}_i$。对
$U_i \in \Ob(\mathcal{C}_i)$ 以及 $\mathcal{C}_i$ 上的层
$\mathcal{F}$，有
\begin{equation}
\label{sites-equation-compute-pullback-to-limit}
f_i^{-1}\mathcal{F}(u_i(U_i)) =
\colim_{a : j \to i} f_a^{-1}\mathcal{F}(u_a(U_i))
\end{equation}
\end{lemma}

\begin{proof}
由引理 \ref{sites-lemma-colimit-sites} 证明中的论证，立即可知各函子 $u_i$ 连续。
为完成证明，必须说明 $f_i^{-1} := u_{i, s}$ 是正合函子
$\Sh(\mathcal{C}_i) \to \Sh(\mathcal{C})$。事实上，只需证明
$f_i^{-1}$ 左正合，因为它作为左伴随是右正合的（Categories，引理
\ref{categories-lemma-exact-adjoint}）。先证明
(\ref{sites-equation-compute-pullback-to-limit})，然后由此推出正合性。

\medskip\noindent
对 $\mathcal{C}$ 的任意对象 $V$，可以选取 $a : j \to i$ 和对象
$V_j \in \Ob(\mathcal{C})$，使得 $V = u_j(V_j)$。然后可令
% P01-ERRATA: P01-E0076 -- source says "pick a a" before the arrow a : j -> i.
% P01-ERRATA: P01-E0077 -- V_j is typed as an object of C, but must lie in C_j for u_j(V_j).
$$
\mathcal{G}(V) = \colim_{b : k \to j} f_{a \circ b}^{-1}\mathcal{F}(u_b(V_j))
$$
该余极限的值 $\mathcal{G}(V)$ 与 $b : j \to i$ 的选择以及满足
$u_j(V_j) = V$ 的对象 $V_j$ 的选择无关；验证从略。
% P01-ERRATA: P01-E0078 -- the choice was denoted a : j -> i; source switches to b and reuses b incompatibly.
此外，若 $\alpha : V \to V'$ 是 $\mathcal{C}$ 的态射，则可以选择
$b : j \to i$ 和态射 $\alpha_j : V_j \to V'_j$，使得
$u_j(\alpha_j) = \alpha$。利用沿态射
$u_b(\alpha_j) : u_b(V_j) \to u_b(V'_j)$ 的限制映射，可诱导映射
$\mathcal{G}(V') \to \mathcal{G}(V)$。验证表明 $\mathcal{G}$ 是预层
（从略）。事实上，$\mathcal{G}$ 满足层条件。任一
$\mathcal{C}$ 中的覆盖 $\mathcal{U} = \{U^t \to U\}$ 都来自某个有限层级。
设 $\mathcal{U}_j = \{U^t_j \to U_j\}$ 在 $\mathcal{I}$ 中某个
$a : j \to i$ 下由 $u_j$ 映到 $\mathcal{U}$。则
$$
H^0(\mathcal{U}, \mathcal{G}) =
\colim_{b : k \to j} H^0(u_b(\mathcal{U}_j), f_{b \circ a}^{-1}\mathcal{F}) =
\colim_{b : k \to j} f_{b \circ a}^{-1}\mathcal{F}(u_b(U_j)) =
\mathcal{G}(U)
$$
如所需。第一个等号成立是因为滤过余极限与有限极限交换（Categories，引理
\ref{categories-lemma-directed-commutes}）。
% P01-ERRATA: P01-E0079 -- b \circ a is ill-typed here; the setup and earlier formula require a \circ b.
根据构造，$\mathcal{G}(U)$ 由
(\ref{sites-equation-compute-pullback-to-limit}) 的右边给出。因此，如果能证明
$\mathcal{G}$ 等于 $f_i^{-1}\mathcal{F}$，则
(\ref{sites-equation-compute-pullback-to-limit}) 成立。

\medskip\noindent
本段验证 $\mathcal{G}$ 与 $f_i^{-1}\mathcal{F}$ 典范同构。强烈建议读者
跳过本段。为验证这一点，必须证明有双射
$\Mor_{\Sh(\mathcal{C})}(\mathcal{G}, \mathcal{H}) =
\Mor_{\Sh(\mathcal{C}_i)}(\mathcal{F}, f_{i, *}\mathcal{H})$，并且它对
$\mathcal{C}$ 上的层 $\mathcal{H}$ 是函子性的，其中
$f_{i, *} = u_i^p$。映射 $\mathcal{G} \to \mathcal{H}$ 等价于相容的映射系
$$
\varphi_{a, b, V_j} :
f_{a \circ b}^{-1}\mathcal{F}(u_b(V_j))
\longrightarrow
\mathcal{H}(u_j(V_j))
$$
这里 $a : j \to i$、$b : k \to j$ 且
$V_j \in \Ob(\mathcal{C}_j)$ 任意。相容性迫使映射
$\varphi_{a, b, V_j}$ 等于
$\varphi_{a \circ b, \text{id}, u_b(V_j)}$。给定
$a : j \to i$，映射族 $\varphi_{a, \text{id}, V_j}$ 对应于层的映射
$\varphi_a : f_a^{-1}\mathcal{F} \to f_{j, *}\mathcal{H}$。
$\varphi_{a, \text{id}, u_a(V_i)}$ 与
$\varphi_{\text{id}, \text{id}, V_i}$ 之间的相容性蕴含：
$\varphi_a$ 是映射 $\varphi_{id}$ 关于下列等式给出的伴随：
$$
\Mor_{\Sh(\mathcal{C}_j)}(f_a^{-1}\mathcal{F}, f_{j, *}\mathcal{H}) =
\Mor_{\Sh(\mathcal{C}_i)}(\mathcal{F}, f_{a, *}f_{j, *}\mathcal{H}) =
\Mor_{\Sh(\mathcal{C}_i)}(\mathcal{F}, f_{i, *}\mathcal{H})
$$
% P01-ERRATA: P01-E0080 -- source has subject-verb disagreement: "compatibilities ... implies".
最终可知，整个映射系 $\varphi_{a, b, V_j}$ 由映射
$\varphi_{id} : \mathcal{F} \to f_{i, *}\mathcal{H}$ 决定。反过来，给定
这样的映射 $\psi : \mathcal{F} \to f_{i, *}\mathcal{H}$，反向阅读刚才的
论证即可构造映射族 $\varphi_{a, b, V_j}$。这便完成了
$\mathcal{G} = f_i^{-1}\mathcal{F}$ 的证明。

\medskip\noindent
假设 (\ref{sites-equation-compute-pullback-to-limit}) 成立。则函子
$\mathcal{F} \mapsto f_i^{-1}\mathcal{F}(U)$ 与有限极限交换，因为层的
有限极限在预层范畴中计算（引理 \ref{sites-lemma-limit-sheaf}），各函子
$f_a^{-1}$ 与有限极限交换，并且滤过余极限与有限极限交换。为说明对
$\mathcal{C}$ 的一般对象 $V$，函子
$\mathcal{F} \mapsto f_i^{-1}\mathcal{F}(V)$ 与有限极限交换，可以对上面
给出的 $f_i^{-1}\mathcal{F}(V) = \mathcal{G}(V)$ 公式使用同一论证。因此
$f_i^{-1}$ 左正合，引理得证。
\end{proof}

\begin{lemma}
\label{sites-lemma-colimit}
在情形 \ref{sites-situation-inverse-limit-sites} 中，假设给定：
\begin{enumerate}
\item 对所有 $i \in \Ob(\mathcal{I})$，给定 $\mathcal{C}_i$ 上的层
$\mathcal{F}_i$；
\item 对 $a : j \to i$，给定 $\mathcal{C}_j$ 上的层映射
$\varphi_a : f_a^{-1}\mathcal{F}_i \to \mathcal{F}_j$。
\end{enumerate}
并且只要 $c = a \circ b$，就有
$\varphi_c = \varphi_b \circ f_b^{-1}\varphi_a$。在引理
\ref{sites-lemma-colimit-sites} 的位点 $\mathcal{C}$ 上，令
$\mathcal{F} = \colim f_i^{-1}\mathcal{F}_i$。设
$i \in \Ob(\mathcal{I})$ 且 $X_i \in \text{Ob}(\mathcal{C}_i)$。则
$$
\colim_{a : j \to i} \mathcal{F}_j(u_a(X_i)) = \mathcal{F}(u_i(X_i))
$$
\end{lemma}

\begin{proof}
形式论证表明
$$
\colim_{a : j \to i} \mathcal{F}_i(u_a(X_i)) =
\colim_{a : j \to i} \colim_{b : k \to j}
f_b^{-1}\mathcal{F}_j(u_{a \circ b}(X_i))
$$
% P01-ERRATA: P01-E0081 -- the left side uses F_i on an object of C_j; statement and typing require F_j.
由 (\ref{sites-equation-compute-pullback-to-limit})，内层余极限等于
$f_j^{-1}\mathcal{F}_j(u_i(X_i))$，故由引理
\ref{sites-lemma-directed-colimits-sections} 得到结论。
\end{proof}

\begin{lemma}
\label{sites-lemma-colimit-push-pull}
在情形 \ref{sites-situation-inverse-limit-sites} 中，假设有
$\mathcal{C}$ 上的层 $\mathcal{F}$。则
$$
\mathcal{F} = \colim f_i^{-1}f_{i, *}\mathcal{F}
$$
其中对 $a : j \to i$，转移映射是 $f_j^{-1}\varphi_a$；这里
$\varphi_a : f_a^{-1}f_{i, *}\mathcal{F} \to f_{j, *}\mathcal{F}$
是典范映射，满足与引理 \ref{sites-lemma-colimit} 中相同的余圈条件。
\end{lemma}

\begin{proof}
对态射
$$
\varphi_a : f_a^{-1}f_{i, *}\mathcal{F} \to f_{j, *}\mathcal{F}
$$
我们取恒等映射
$$
f_{i, *}\mathcal{F} \to f_{a, *}f_{j, *}\mathcal{F}
$$
的伴随。因此 $\varphi_a$ 是伴随对
$(f_a^{-1}, f_{a, *})$ 的余单位。必须证明，对所有
$a : j \to i$ 和 $b : k \to i$，若复合 $c = a \circ b$，则
$\varphi_c = \varphi_b \circ f_b^{-1}\varphi_a$。
% P01-ERRATA: P01-E0082 -- b is typed k -> i, so c = a \circ b is undefined; b must target j.
这由 Categories，引理 \ref{categories-lemma-compose-counits} 得出。最后，必须
证明由伴随给出的映射
$$
\colim_{i \in I} f_i^{-1}f_{i, *}\mathcal{F}
\longrightarrow
\mathcal{F}
$$
是同构。对 $\mathcal{C}$ 的对象 $U$，需要证明映射
$$
(\colim_{i \in I} f_i^{-1}\mathcal{F}_i)(U) \to \mathcal{F}(U)
$$
是双射。选择 $i$ 和 $\mathcal{C}_i$ 的对象 $U_i$，使得
$u_i(U_i) = U$。则由引理 \ref{sites-lemma-colimit}，左边等于
% P01-ERRATA: P01-E0083 -- F_i is undefined in this lemma; both displayed occurrences should be f_{i,*}F.
$$
(\colim_{i \in I} f_i^{-1}\mathcal{F}_i)(U) =
\colim_{a : j \to i} f_{j, *}\mathcal{F}(u_a(U_i))
$$
因为 $u_j(u_a(U_i)) = U$，根据定义，对所有
$a : j \to i$ 都有
$f_{j, *}\mathcal{F}(u_a(U_i)) = \mathcal{F}(U)$。因此该余极限的值是
$\mathcal{F}(U)$，证明完成。
\end{proof}












\section{预层的进一步函子性}
\label{sites-section-more-functoriality-PSh}

\noindent
本节重新讨论第 \ref{sites-section-functoriality-PSh} 节的内容。设
$u : \mathcal{C} \to \mathcal{D}$ 为范畴之间的函子。回忆
$$
u^p :
\textit{PSh}(\mathcal{D})
\longrightarrow
\textit{PSh}(\mathcal{C})
$$
是这样的函子：它把 $\mathcal{D}$ 上的 $\mathcal{G}$ 送到预层
$u^p\mathcal{G} = \mathcal{G} \circ u$。事实表明，该函子不仅有左伴随
（即 $u_p$），而且还有右伴随。

\medskip\noindent
具体地，对任意 $V \in \Ob(\mathcal{D})$，定义范畴
${}_V\mathcal{I} = {}_V^u\mathcal{I}$。其对象是序对
$(U, \psi : u(U) \to V)$。注意，这里的箭头方向与第
\ref{sites-section-functoriality-PSh} 节定义范畴 $\mathcal{I}_V^u$ 时所用的
箭头方向相反。态射 $(U, \psi) \to (U', \psi')$ 由满足
$\psi = \psi' \circ u(\alpha)$ 的态射 $\alpha : U \to U'$ 给出。
此外，给定 $\mathcal{C}$ 上任意集合预层 $\mathcal{F}$，引入函子
${}_V\mathcal{F} : {}_V\mathcal{I}^{opp} \to \textit{Sets}$，由规则
${}_V\mathcal{F}(U, \psi) = \mathcal{F}(U)$ 定义。令
$$
{}_pu(\mathcal{F})(V) := \lim_{{}_V\mathcal{I}^{opp}} {}_V\mathcal{F}
$$
作为极限，对 ${}_V\mathcal{I}$ 的每个对象 $(U, \psi)$，都有投影映射
$c(\psi) : {}_pu(\mathcal{F})(V) \to \mathcal{F}(U)$。事实上，
$$
{}_pu(\mathcal{F})(V)
=
\left\{
\begin{matrix}
\text{族 }
s_{(U, \psi)} \in \mathcal{F}(U) \\
\forall \beta : (U_1, \psi_1) \to (U_2, \psi_2)
\text{ 于 }{}_V\mathcal{I} \\
\text{ 有 } \beta^*s_{(U_2, \psi_2)} = s_{(U_1, \psi_1)}
\end{matrix}
\right\}
$$
其中对应由 $s \mapsto s_{(U, \psi)} = c(\psi)(s)$ 给出。留给读者定义与
$\mathcal{D}$ 的任意态射 $V' \to V$ 相伴的限制映射
${}_pu(\mathcal{F})(V) \to {}_pu(\mathcal{F})(V')$。所得预层记为
${}_pu\mathcal{F}$。

\begin{lemma}
\label{sites-lemma-recover-pu}
存在与限制映射相容的典范映射
${}_pu\mathcal{F}(u(U)) \to \mathcal{F}(U)$。
\end{lemma}

\begin{proof}
这就是上面的投影映射 $c(\text{id}_{u(U)})$。
\end{proof}

\noindent
注意，预层的任意映射 $\mathcal{F} \to \mathcal{F}'$ 都产生函子间相容的
映射系 ${}_V\mathcal{F} \to {}_V\mathcal{F}'$，从而产生预层的映射
${}_pu\mathcal{F} \to {}_pu\mathcal{F}'$。换言之，我们定义了函子
$$
{}_pu :
\textit{PSh}(\mathcal{C})
\longrightarrow
\textit{PSh}(\mathcal{D})
$$

\begin{lemma}
\label{sites-lemma-adjoints-pu}
函子 ${}_pu$ 是函子 $u^p$ 的右伴随。换言之，公式
$$
\Mor_{\textit{PSh}(\mathcal{C})}(u^p\mathcal{G}, \mathcal{F})
=
\Mor_{\textit{PSh}(\mathcal{D})}(\mathcal{G}, {}_pu\mathcal{F})
$$
关于 $\mathcal{F}$ 和 $\mathcal{G}$ 双函子地成立。
\end{lemma}

\begin{proof}
证明与引理 \ref{sites-lemma-adjoints-u} 的证明完全相同。注意，引理
\ref{sites-lemma-recover-pu} 中的映射
$u^p{}_pu \mathcal{F} \to \mathcal{F}$ 正是用于从右边得到左边的映射。

\medskip\noindent
另一种证法是：把 $\mathcal{C}$ 上的集合预层 $\mathcal{F}$ 看作
$\mathcal{C}^{opp}$ 上取值于 $\textit{Sets}^{opp}$ 的预层
$\mathcal{F}'$，并对 $\mathcal{D}$ 作类似处理。验证
$({}_pu \mathcal{F})' = u_p(\mathcal{F}')$ 且
$(u^p\mathcal{G})' = u^p(\mathcal{G}')$。由备注
\ref{sites-remark-functoriality-presheaves-values}，对取值于
$\textit{Sets}^{opp}$ 的预层，$u_p$ 与 $u^p$ 互为伴随。结论由此形式地
得到。
\end{proof}

\noindent
因此，给定范畴的函子 $u : \mathcal{C} \to \mathcal{D}$，可在预层范畴
之间得到函子序列
$$
u_p, u^p, {}_pu
$$
其中每一对相邻函子中，前者是后者的左伴随。

\begin{lemma}
\label{sites-lemma-adjoint-functors}
设 $u : \mathcal{C} \to \mathcal{D}$ 和
$v : \mathcal{D} \to \mathcal{C}$ 为范畴的函子。假设 $v$ 是 $u$ 的
右伴随。则：
\begin{enumerate}
\item 对 $\mathcal{D}$ 中任意 $V$，有 $u^ph_V = h_{v(V)}$；
\item 范畴 $\mathcal{I}^v_U$ 有始对象；
\item 范畴 ${}_V^u\mathcal{I}$ 有终对象；
\item ${}_pu = v^p$；
\item $u^p = v_p$。
\end{enumerate}
\end{lemma}

\begin{proof}
(1) 的证明。设 $V$ 为 $\mathcal{D}$ 的对象。由假设，
$u^ph_V = h_{v(V)}$，因为
$u^ph_V(U) = \Mor_\mathcal{D}(u(U), V) = \Mor_\mathcal{C}(U, v(V))$。

\medskip\noindent
(2) 的证明。设 $U$ 为 $\mathcal{C}$ 的对象。令
$\eta : U \to v(u(U))$ 为映射 $\text{id} : u(U) \to u(U)$ 的伴随映射。
断言 $(u(U), \eta)$ 是 $\mathcal{I}_U^v$ 的始对象。事实上，给定
$\mathcal{I}_U^v$ 的对象 $(V, \phi : U \to v(V))$，态射 $\phi$ 的伴随是
映射 $\psi : u(U) \to V$，后者定义态射
$(u(U), \eta) \to (V, \phi)$。

\medskip\noindent
(3) 的证明。设 $V$ 为 $\mathcal{D}$ 的对象。令
$\xi : u(v(V)) \to V$ 为映射 $\text{id} : v(V) \to v(V)$ 的伴随映射。
断言 $(v(V), \xi)$ 是 ${}_V^u\mathcal{I}$ 的终对象。事实上，给定
${}_V^u\mathcal{I}$ 的对象 $(U, \psi : u(U) \to V)$，态射 $\psi$ 的
伴随是映射 $\phi : U \to v(V)$，后者定义态射
$(U, \psi) \to (v(V), \xi)$。

\medskip\noindent
因此，对 $\mathcal{C}$ 上任意预层 $\mathcal{F}$，有
\begin{eqnarray*}
v^p\mathcal{F}(V)
& = &
\mathcal{F}(v(V)) \\
& = &
\Mor_{\textit{PSh}(\mathcal{C})}(h_{v(V)}, \mathcal{F}) \\
& = &
\Mor_{\textit{PSh}(\mathcal{C})}(u^ph_V, \mathcal{F}) \\
& = &
\Mor_{\textit{PSh}(\mathcal{D})}(h_V, {}_pu\mathcal{F}) \\
& = &
{}_pu\mathcal{F}(V)
\end{eqnarray*}
这证明了 (4)。由伴随函子的唯一性，(5) 随之成立。
\end{proof}





\section{余连续函子}
\label{sites-section-cocontinuous-functors}

\noindent
还有另一种构造托扑斯态射的方法。它使用位点之间的余连续函子，定义如下。

\begin{definition}
\label{sites-definition-cocontinuous}
设 $\mathcal{C}$ 和 $\mathcal{D}$ 为位点，并设
$u : \mathcal{C} \to \mathcal{D}$ 为函子。若对每个
$U \in \Ob(\mathcal{C})$ 以及 $\mathcal{D}$ 的每个覆盖
$\{V_j \to u(U)\}_{j \in J}$，都存在 $\mathcal{C}$ 的覆盖
$\{U_i \to U\}_{i\in I}$，使得态射族
$\{u(U_i) \to u(U)\}_{i \in I}$ 加细覆盖
$\{V_j \to u(U)\}_{j \in J}$，则称函子 $u$ 是{\it 余连续的}。
\end{definition}

\noindent
注意，$\{u(U_i) \to u(U)\}_{i \in I}$ 一般{\it 不是}位点
$\mathcal{D}$ 的覆盖。

\begin{lemma}
\label{sites-lemma-pu-sheaf}
设 $\mathcal{C}$ 和 $\mathcal{D}$ 为位点，
$u : \mathcal{C} \to \mathcal{D}$ 余连续，并设 $\mathcal{F}$ 为
$\mathcal{C}$ 上的层。则 ${}_pu\mathcal{F}$ 是 $\mathcal{D}$ 上的层，
将其记为 ${}_su\mathcal{F}$。
\end{lemma}

\begin{proof}
令 $\{V_j \to V\}_{j \in J}$ 为位点 $\mathcal{D}$ 的覆盖。必须证明
$$
\xymatrix{
&\ \phantom{}_pu\mathcal{F}(V) \ar[r] &
\prod {}_pu\mathcal{F}(V_j) \ar@<1ex>[r] \ar@<-1ex>[r] &
\prod {}_pu\mathcal{F}(V_j \times_V V_{j'}) &
}
$$
是等化子图。因为 ${}_pu$ 是 $u^p$ 的右伴随，有
$$
{}_pu\mathcal{F}(V) =
\Mor_{\textit{PSh}(\mathcal{D})}(h_V, {}_pu\mathcal{F}) =
\Mor_{\textit{PSh}(\mathcal{C})}(u^ph_V, \mathcal{F}) =
\Mor_{\Sh(\mathcal{C})}((u^ph_V)^\#, \mathcal{F})
$$
因此，只需证明
\begin{equation}
\label{sites-equation-coequalizer}
\xymatrix{
\coprod u^p h_{V_j \times_V V_{j'}} \ar@<1ex>[r] \ar@<-1ex>[r] &
\coprod u^p h_{V_j} \ar[r] &
u^p h_V
}
\end{equation}
在层化后成为余等化子图。（回忆一下，层范畴中的余积是预层范畴中余积的
层化，见引理 \ref{sites-lemma-colimit-sheaves}。）

\medskip\noindent
先证明 (\ref{sites-equation-coequalizer}) 的第二个箭头在层化后成为满射。为此使用
引理 \ref{sites-lemma-mono-epi-sheaves}。因此，只需证明：$u^ph_V$ 在 $U$ 上的
截面 $s$ 可在 $U$ 的某个覆盖的各成员上提升为
$\coprod u^p h_{V_j}$ 的截面。注意，$s$ 是态射
$s : u(U) \to V$。于是
$\{V_j \times_{V, s} u(U) \to u(U)\}$ 是 $\mathcal{D}$ 的覆盖。因为 $u$
余连续，存在覆盖 $\{U_i \to U\}$，使得
$\{u(U_i) \to u(U)\}$ 加细
$\{V_j \times_{V, s} u(U) \to u(U)\}$。这意味着每个限制
$s|_{U_i} : u(U_i) \to V$ 都经由某个 $j$ 的态射
$s_i : u(U_i) \to V_j$ 分解；换言之，$s|_{U_i}$ 按要求在
$u^ph_{V_j}(U_i) \to u^ph_V(U_i)$ 的像中。

\medskip\noindent
设 $s, s' \in (\coprod u^ph_{V_j})^\#(U)$ 映到
$(u^ph_V)^\#(U)$ 中同一元素。为完成引理的证明，说明把 $U$ 替换为某个
覆盖的各成员后，$s, s'$ 是 $\coprod u^p h_{V_j \times_V V_{j'}}$ 的同一
截面在 (\ref{sites-equation-coequalizer}) 的两个映射下的像。先把 $U$ 替换为
某个覆盖的各成员，从而可以假设 $s \in u^ph_{V_j}(U)$ 且
$s' \in u^ph_{V_{j'}}(U)$。再次作这样的替换，可以保证 $s$ 和 $s'$ 在
$u^ph_V(U)$ 中，而不只是在层化中，具有相同的像。因此
$s : u(U) \to V_j$ 和 $s' : u(U) \to V_{j'}$ 是 $\mathcal{D}$ 的态射，
使得
$$
\xymatrix{
u(U) \ar[r]_{s'} \ar[d]_s & V_{j'} \ar[d] \\
V_j \ar[r] & V
}
$$
交换。于是得到 $t = (s, s') : u(U) \to V_j \times_V V_{j'}$，即截面
$t \in u^ph_{V_j \times_V V_{j'}}(U)$，它按要求映到 $s, s'$。
\end{proof}

\begin{lemma}
\label{sites-lemma-exact-cocontinuous}
设 $\mathcal{C}$ 和 $\mathcal{D}$ 为位点，
$u : \mathcal{C} \to \mathcal{D}$ 余连续。函子
$\Sh(\mathcal{D}) \to \Sh(\mathcal{C})$，
$\mathcal{G} \mapsto (u^p\mathcal{G})^\#$ 是上面引理
\ref{sites-lemma-pu-sheaf} 中引入的函子 ${}_su$ 的左伴随。此外，它是正合的。
\end{lemma}

\begin{proof}
如下证明伴随性质：
\begin{eqnarray*}
\Mor_{\Sh(\mathcal{C})}
((u^p\mathcal{G})^\#, \mathcal{F})
& = &
\Mor_{\textit{PSh}(\mathcal{C})}
(u^p\mathcal{G}, \mathcal{F}) \\
& = &
\Mor_{\textit{PSh}(\mathcal{D})}
(\mathcal{G}, {}_pu\mathcal{F}) \\
& = &
\Mor_{\Sh(\mathcal{D})}
(\mathcal{G}, {}_su\mathcal{F}).
\end{eqnarray*}
因此它是左伴随，从而右正合，见 Categories，引理
\ref{categories-lemma-exact-adjoint}。已经看到层化是左正合的，见引理
\ref{sites-lemma-sheafification-exact}。此外，由引理
\ref{sites-lemma-limit-sheaf}，包含函子
$i : \Sh(\mathcal{D}) \to \textit{PSh}(\mathcal{D})$ 左正合。最后，函子
$u^p$ 左正合，因为它是右伴随（即 $u_p$ 的右伴随）。因此，该函子是左正合
函子之复合 ${}^\# \circ u^p \circ i$，所以左正合。
\end{proof}

\noindent
以一个技术性引理结束本节。

\begin{lemma}
\label{sites-lemma-technical-pu}
在引理 \ref{sites-lemma-exact-cocontinuous} 的情形中。对 $\mathcal{D}$ 上任意
预层 $\mathcal{G}$，有
$(u^p\mathcal{G})^\# = (u^p(\mathcal{G}^\#))^\#$。
\end{lemma}
% P01-ERRATA: P01-E0084 -- the first sentence of the lemma is a fragment ending with a period.

\begin{proof}
对 $\mathcal{C}$ 上任意层 $\mathcal{F}$，有
\begin{eqnarray*}
\Mor_{\Sh(\mathcal{C})}((u^p(\mathcal{G}^\#))^\#, \mathcal{F})
& = &
\Mor_{\Sh(\mathcal{D})}(\mathcal{G}^\#, {}_su\mathcal{F}) \\
& = &
\Mor_{\Sh(\mathcal{D})}(\mathcal{G}^\#, {}_pu\mathcal{F}) \\
& = &
\Mor_{\textit{PSh}(\mathcal{D})}(\mathcal{G}, {}_pu\mathcal{F}) \\
& = &
\Mor_{\textit{PSh}(\mathcal{C})}(u^p\mathcal{G}, \mathcal{F}) \\
& = &
\Mor_{\Sh(\mathcal{C})}((u^p\mathcal{G})^\#, \mathcal{F})
\end{eqnarray*}
结论由米田引理得出。
\end{proof}

\begin{remark}
\label{sites-remark-cartesian-cocontinuous}
设 $u : \mathcal{C} \to \mathcal{D}$ 为范畴之间的函子。给定
$\mathcal{D}$ 中的态射 $g : u(U) \to V$ 和 $f : W \to V$，可以考虑函子
$$
\mathcal{C}^{opp} \longrightarrow \textit{Sets},\quad
T \longmapsto
\Mor_\mathcal{C}(T, U)
\times_{\Mor_\mathcal{D}(u(T), V)}
\Mor_\mathcal{D}(u(T), W)
$$
若该函子可表，以 $U \times_{g, V, f} W$ 表示 $\mathcal{C}$ 中相应的对象。
假设 $\mathcal{C}$ 和 $\mathcal{D}$ 是位点。考虑性质 $P$：对
$\mathcal{D}$ 的每个覆盖 $\{f_j : V_j \to V\}$ 和任意态射
$g : u(U) \to V$，有：
\begin{enumerate}
\item 对所有 $i$，$U \times_{g, V, f_i} V_i$ 存在；
\item $\{U \times_{g, V, f_i} V_i \to U\}$ 是 $\mathcal{C}$ 的覆盖。
\end{enumerate}
% P01-ERRATA: P01-E0085 -- the covering is indexed by j, but both property clauses switch to undefined i.
请注意这与连续函子定义的相似之处。若 $u$ 具有 $P$，则 $u$ 余连续
（细节从略）。我们将遇到的许多余连续函子都满足 $P$。
\end{remark}



\section{余连续函子与托扑斯态射}
\label{sites-section-cocontinuous-morphism-topoi}

\noindent
由上文清楚可知，余连续函子 $u$ 给出一个与 $u$ 同方向的托扑斯态射。
因此，这与定义 \ref{sites-definition-morphism-sites}、命题
\ref{sites-proposition-get-morphism} 和引理
\ref{sites-lemma-morphism-sites-topoi} 中在某些条件下与连续 $u$ 相伴的托扑斯态射
方向相反。

\begin{lemma}
\label{sites-lemma-cocontinuous-morphism-topoi}
设 $\mathcal{C}$ 和 $\mathcal{D}$ 为位点，并设
$u : \mathcal{C} \to \mathcal{D}$ 余连续。函子
$g_* = {}_su$ 和 $g^{-1} = (u^p\ )^\#$ 定义从
$\Sh(\mathcal{C})$ 到 $\Sh(\mathcal{D})$ 的托扑斯态射 $g$。
\end{lemma}

\begin{proof}
这正是引理 \ref{sites-lemma-exact-cocontinuous} 的内容。
\end{proof}

\begin{lemma}
\label{sites-lemma-composition-cocontinuous}
\begin{slogan}
位点函子的复合保持位点函子的余连续性。
\end{slogan}
设 $u : \mathcal{C} \to \mathcal{D}$ 和
$v : \mathcal{D} \to \mathcal{E}$ 为余连续函子。则
$v \circ u$ 余连续，并且 $h = g \circ f$，其中
$f : \Sh(\mathcal{C}) \to \Sh(\mathcal{D})$、相应地
$g : \Sh(\mathcal{D}) \to \Sh(\mathcal{E})$、相应地
$h : \Sh(\mathcal{C}) \to \Sh(\mathcal{E})$，分别是与 $u$、相应地
$v$、相应地 $v \circ u$ 相伴的托扑斯态射。
\end{lemma}

\begin{proof}
令 $U \in \Ob(\mathcal{C})$。设 $\{E_i \to v(u(U))\}$ 是
$\mathcal{E}$ 中 $U$ 的一个覆盖。根据假设，存在 $\mathcal{D}$ 中的覆盖
$\{D_j \to u(U)\}$，使得 $\{v(D_j) \to v(u(U))\}$ 加细
$\{E_i \to v(u(U))\}$。同样根据假设，存在 $\mathcal{C}$ 中的覆盖
$\{C_l \to U\}$，使得 $\{u(C_l) \to u(U)\}$ 加细
$\{D_j \to u(U)\}$。于是
$\{v(u(C_l)) \to v(u(U))\}$ 确实加细覆盖
$\{E_i \to v(u(U))\}$。这证明 $v \circ u$ 余连续。
% P01-ERRATA: P01-E0086 -- the source calls {E_i -> v(u(U))} a covering of U in E; U is not an E-object.
为证明最后的断言，只需说明
${}_sv \circ {}_su = {}_s(v \circ u)$。由引理 \ref{sites-lemma-pu-sheaf}，只需
证明 ${}_pv \circ {}_pu = {}_p(v \circ u)$。因为 ${}_pu$、相应地
${}_pv$、相应地 ${}_p(v \circ u)$，分别是 $u^p$、相应地 $v^p$、
相应地 $(v \circ u)^p$ 的右伴随，所以只需证明
$u^p \circ v^p = (v \circ u)^p$。这直接由定义得出。
\end{proof}

\begin{example}
\label{sites-example-open-immersion-cocontinuous}
设 $X$ 为拓扑空间，$j : U  \to X$ 为开子空间的包含映射。回忆一下，
我们有位点 $X_{Zar}$ 和 $U_{Zar}$，见例
\ref{sites-example-site-topological}。还记得，与 $j$ 相伴的函子
$u : X_{Zar} \to U_{Zar}$ 连续，并产生位点态射
$U_{Zar} \to X_{Zar}$，见例 \ref{sites-example-continuous-map}。它也给出托扑斯
态射 $(j_*, j^{-1})$。接着考虑函子
$v : U_{Zar} \to X_{Zar}$，$V \mapsto v(V) = V$（仍是同一个开集，但现在
把它看作 $X_{Zar}$ 的对象）。该函子余连续。事实上，若
$v(V) = \bigcup_{j \in J} W_j$ 是 $X$ 中的开覆盖，则每个 $W_j$ 必为
$U$ 的子集，因而形如 $v(V_j)$，而且显然
$V = \bigcup_{j \in J} V_j$ 是 $U$ 中的开覆盖。由上面的引理
\ref{sites-lemma-cocontinuous-morphism-topoi}，可知有与 $v$ 相伴的托扑斯态射
$$
\Sh(U) \longrightarrow \Sh(X)
$$
由 ${}_sv$ 和 $(v^p\ )^\#$ 给出。断言事实上
$(v^p\ )^\# = j^{-1}$ 且 ${}_sv = j_*$；换言之，这与上面给出的托扑斯
态射相同。看出这一点的最简便方法或许是注意：对 $X$ 上任意层
$\mathcal{G}$，有 $v^p\mathcal{G}(V) = \mathcal{G}(V)$；根据 Sheaves，引理
\ref{sheaves-lemma-j-pullback}，这描述了 $j^{-1}\mathcal{G}$（因此在此情形
无需层化）。${}_sv$ 与 $j_*$ 的相等由伴随函子的唯一性得出（也可直接
计算）。
\end{example}

\begin{example}
\label{sites-example-open-map-cocontinuous}
这是例 \ref{sites-example-open-immersion-cocontinuous} 的一个小推广。设
$f : X \to Y$ 为拓扑空间的连续映射，并假设 $f$ 是开映射。回忆一下，
我们有位点 $X_{Zar}$ 和 $Y_{Zar}$，见例
\ref{sites-example-site-topological}。还记得，与 $f$ 相伴的函子
$u : Y_{Zar} \to X_{Zar}$ 连续，并产生位点态射
$X_{Zar} \to Y_{Zar}$，见例 \ref{sites-example-continuous-map}。它也给出托扑斯
态射 $(f_*, f^{-1})$。接着考虑函子
$v : X_{Zar} \to Y_{Zar}$，$U \mapsto v(U) = f(U)$。该函子余连续。
事实上，若 $f(U) = \bigcup_{j \in J} V_j$ 是 $Y$ 中的开覆盖，则令
$U_j = f^{-1}(V_j) \cap U$，可得开覆盖 $U = \bigcup U_j$，使得
$f(U) = \bigcup f(U_j)$ 加细 $f(U) = \bigcup V_j$。由上面的引理
\ref{sites-lemma-cocontinuous-morphism-topoi}，可知有与 $v$ 相伴的托扑斯态射
$$
\Sh(X) \longrightarrow \Sh(Y)
$$
由 ${}_sv$ 和 $(v^p\ )^\#$ 给出。断言事实上
$(v^p\ )^\# = f^{-1}$ 且 ${}_sv = f_*$；换言之，这与上面给出的托扑斯
态射相同。对 $Y$ 上任意层 $\mathcal{G}$，有
$v^p\mathcal{G}(U) = \mathcal{G}(f(U))$。另一方面，可以计算
$u_p\mathcal{G}(U) = \colim_{f(U) \subset V} \mathcal{G}(V)
= \mathcal{G}(f(U))$，因为显然
$(f(U), U \subset f^{-1}(f(U)))$ 是第
\ref{sites-section-functoriality-PSh} 节中范畴 $\mathcal{I}_U^u$ 的始对象。
因此 $u_p = v^p$，从而 $f^{-1} = u_s = (v^p\ )^\#$。${}_sv$ 与 $f_*$
的相等由伴随函子的唯一性得出（也可直接计算）。
\end{example}

\noindent
在第一个例 \ref{sites-example-open-immersion-cocontinuous} 中，函子 $v$ 也连续。
但在第二个例 \ref{sites-example-open-map-cocontinuous} 中，它一般不连续，因为
定义 \ref{sites-definition-continuous} 的条件 (2) 可能不成立。因此下述引理适用于
第一个例，但不适用于第二个例。

\begin{lemma}
\label{sites-lemma-when-shriek}
\begin{slogan}
对同时连续且余连续的位点函子，相伴托扑斯态射中的层化是多余的。
\end{slogan}
设 $\mathcal{C}$ 和 $\mathcal{D}$ 为位点，并设
$u : \mathcal{C} \to \mathcal{D}$ 为函子。假设：
\begin{enumerate}
\item[(a)] $u$ 余连续；
\item[(b)] $u$ 连续。
\end{enumerate}
令 $g : \Sh(\mathcal{C}) \to \Sh(\mathcal{D})$ 为相伴的托扑斯态射。则：
\begin{enumerate}
\item 公式 $g^{-1} = (u^p\ )^\#$ 中无需层化，换言之，
$g^{-1}(\mathcal{G})(U) = \mathcal{G}(u(U))$；
\item $g^{-1}$ 有左伴随 $g_{!} = (u_p\ )^\#$；
\item $g^{-1}$ 与任意极限和余极限交换。
\end{enumerate}
\end{lemma}

\begin{proof}
由引理 \ref{sites-lemma-pushforward-sheaf}，对 $\mathcal{D}$ 上任意层
$\mathcal{G}$，预层 $u^p\mathcal{G}$ 是 $\mathcal{C}$ 上的层（记为
$u^s\mathcal{G}$）。由引理 \ref{sites-lemma-adjoint-sheaves}，函子
$\mathcal{F} \mapsto (u_p\mathcal{F})^\#$ 是 $u^s$ 的左伴随。关于极限和
余极限的断言由 Categories，第 \ref{categories-section-adjoint} 节中的讨论
得出。
\end{proof}

\noindent
在上述引理 \ref{sites-lemma-when-shriek} 的情形中，可知有伴随函子序列
$$
g_{!}, \ g^{-1}, \ g_*.
$$
函子 $g_!$ 一般{\it 不}正合，因为一般而言，它不把
$\Sh(\mathcal{C})$ 的终对象变为 $\Sh(\mathcal{D})$ 的终对象。见
Sheaves，备注 \ref{sheaves-remark-j-shriek-not-exact}。另一方面，在例
\ref{sites-example-open-immersion-cocontinuous} 的拓扑情形中，函子 $j_!$ 在
阿贝尔层上正合，见 Modules，引理 \ref{modules-lemma-j-shriek-exact}。
下述引理把它推广到位点情形。

\begin{lemma}
\label{sites-lemma-preserve-equalizers}
设 $\mathcal{C}$ 和 $\mathcal{D}$ 为位点，并设
$u : \mathcal{C} \to \mathcal{D}$ 为函子。假设：
\begin{enumerate}
\item[(a)] $u$ 余连续；
\item[(b)] $u$ 连续；
\item[(c)] $\mathcal{C}$ 中存在纤维积和等化子，并且 $u$ 与它们交换。
\end{enumerate}
此时上述函子 $g_!$ 与纤维积和等化子交换（更一般地，与有限连通极限
交换）。
\end{lemma}

\begin{proof}
假设 (a)、(b)、(c)。有 $g_! = (u_p\ )^\#$。回忆一下（引理
\ref{sites-lemma-limit-sheaf}），层的极限等于相应的预层极限，而层化与有限极限
交换（引理 \ref{sites-lemma-sheafification-exact}）。因此，只需说明 $u_p$ 与
纤维积和等化子交换。为此，只需第 \ref{sites-section-functoriality-PSh} 节中的
范畴 $(\mathcal{I}_V^u)^{opp}$ 上的余极限与纤维积和等化子交换。这由引理
\ref{sites-lemma-almost-directed} 和 Categories，引理
\ref{categories-lemma-almost-directed-commutes-equalizers} 得出。
\end{proof}

\noindent
下述引理处理一种甚至更像与拓扑空间开浸入相伴之态射的情形。

\begin{lemma}
\label{sites-lemma-back-and-forth}
设 $\mathcal{C}$ 和 $\mathcal{D}$ 为位点，并设
$u : \mathcal{C} \to \mathcal{D}$ 为函子。假设：
\begin{enumerate}
\item[(a)] $u$ 余连续；
\item[(b)] $u$ 连续；
\item[(c)] $u$ 全忠实。
\end{enumerate}
对上述 $g_!, g^{-1}, g_*$，任意 $\mathcal{C}$ 上的层 $\mathcal{F}$ 的
典范映射 $\mathcal{F} \to g^{-1}g_!\mathcal{F}$ 和
$g^{-1}g_*\mathcal{F} \to \mathcal{F}$ 都是同构。
\end{lemma}

\begin{proof}
设 $X$ 为 $\mathcal{C}$ 的对象。引理 \ref{sites-lemma-pu-sheaf} 和
\ref{sites-lemma-when-shriek} 已经表明，对函子
$g^{-1} = (u^p\ )^\#$ 和 $g_{*} = ({}_pu\ )^\#$，层化不是必需的。
可以计算
$(g^{-1}g_{*}\mathcal{F})(X) = g_{*}\mathcal{F}(u(X))
= \lim \mathcal{F}(Y)$。这里的极限遍历序对
$(Y, u(Y) \to u(X))$ 所成的范畴，其中不要求态射
$u(Y) \to u(X)$ 形如 $u(\alpha)$，$\alpha$ 为 $\mathcal{C}$ 的态射。
由假设 (c)，它们自动来自 $\mathcal{C}$ 的态射，因而该极限是在
$(X, u(\text{id}_X))$ 上的值，即 $\mathcal{F}(X)$。这证明
$g^{-1}g_{*}\mathcal{F} = \mathcal{F}$。

\medskip\noindent
另一方面，$(g^{-1}g_{!}\mathcal{F})(X) =
g_{!}\mathcal{F}(u(X)) = (u_p\mathcal{F})^\#(u(X))$，并且
$u_p\mathcal{F}(u(X)) = \colim \mathcal{F}(Y)$。这里的余极限遍历序对
$(Y, u(X) \to u(Y))$ 所成的范畴，其中不要求态射
$u(X) \to u(Y)$ 形如 $u(\alpha)$，$\alpha$ 为 $\mathcal{C}$ 的态射。
由假设 (c)，它们自动来自 $\mathcal{C}$ 的态射，因而该余极限是在
$(X, u(\text{id}_X))$ 上的值，即 $\mathcal{F}(X)$。因此，对每个
$X \in \Ob(\mathcal{C})$，有
$u_p\mathcal{F}(u(X)) = \mathcal{F}(X)$。由于 $u$ 余连续且连续，
$\mathcal{D}$ 中 $u(X)$ 的任意覆盖都可由 $\mathcal{D}$ 的一个覆盖
(!) $\{u(X_i) \to u(X)\}$ 加细，其中
$\{X_i \to X\}$ 是 $\mathcal{C}$ 中的覆盖。这也蕴含
$(u_p\mathcal{F})^+(u(X)) = \mathcal{F}(X)$，因为在定义
$(u_p\mathcal{F})^+$ 在 $u(X)$ 上之值的余极限中，可以限制到上述覆盖
$\{u(X_i) \to u(X)\}$ 所成的共尾系。因此，对 $\mathcal{C}$ 的所有对象
$X$，同样有 $(u_p\mathcal{F})^+(u(X)) = \mathcal{F}(X)$。再重复一次该
论证，便对 $\mathcal{C}$ 的所有对象 $X$ 得到
$(u_p\mathcal{F})^\#(u(X)) = \mathcal{F}(X)$。这产生所需等式
$g^{-1}g_!\mathcal{F} = \mathcal{F}$。
\end{proof}

\noindent
最后给出一种没有相应拓扑例子的情形。我们将用该引理考察在保持相同拓扑的
同时扩大概形的“部分宇宙”时会发生什么。在引理的情形中，托扑斯态射
$g : \Sh(\mathcal{C}) \to \Sh(\mathcal{D})$ 把
$\Sh(\mathcal{C})$ 识别为 $\Sh(\mathcal{D})$ 的子托扑斯（第
\ref{sites-section-subtopoi} 节），而且给定嵌入还有一个收缩。

\begin{lemma}
\label{sites-lemma-bigger-site}
设 $\mathcal{C}$ 和 $\mathcal{D}$ 为位点，并设
$u : \mathcal{C} \to \mathcal{D}$ 为函子。假设：
\begin{enumerate}
\item[(a)] $u$ 余连续；
\item[(b)] $u$ 连续；
\item[(c)] $u$ 全忠实；
\item[(d)] $\mathcal{C}$ 中存在纤维积，且 $u$ 与纤维积交换；
\item[(e)] 存在终对象
$e_\mathcal{C} \in \Ob(\mathcal{C})$、
$e_\mathcal{D} \in \Ob(\mathcal{D})$，使得
$u(e_\mathcal{C}) = e_\mathcal{D}$。
\end{enumerate}
令 $g_!, g^{-1}, g_*$ 如上。则 $u$ 定义位点态射
$f : \mathcal{D} \to \mathcal{C}$，并且 $f_* = g^{-1}$、
$f^{-1} = g_!$。复合
$$
\xymatrix{
\Sh(\mathcal{C}) \ar[r]^g &
\Sh(\mathcal{D}) \ar[r]^f &
\Sh(\mathcal{C})
}
$$
同构于托扑斯 $\Sh(\mathcal{C})$ 的恒等态射。此外，函子 $f^{-1}$
全忠实。
\end{lemma}

\begin{proof}
由假设，函子 $u$ 满足命题 \ref{sites-proposition-get-morphism} 的假设。因此，
$u$ 定义位点态射，从而定义引理 \ref{sites-lemma-morphism-sites-topoi} 中的
托扑斯态射 $f$。公式 $f_* = g^{-1}$ 和 $f^{-1} = g_!$ 由所引引理及引理
\ref{sites-lemma-when-shriek} 清楚可得。由引理
\ref{sites-lemma-back-and-forth}，有
$f_* \circ g_* = g^{-1} \circ g_* \cong \text{id}$，以及
$g^{-1} \circ f^{-1} = g^{-1} \circ g_! \cong \text{id}$。

\medskip\noindent
还须证明 $f^{-1}$ 全忠实。令
$\mathcal{F}, \mathcal{G} \in \Ob(\Sh(\mathcal{C}))$。必须证明映射
$$
\Mor_{\Sh(\mathcal{C})}(\mathcal{F}, \mathcal{G})
\longrightarrow
\Mor_{\Sh(\mathcal{D})}(f^{-1}\mathcal{F}, f^{-1}\mathcal{G})
$$
是双射。但右边等于
\begin{align*}
\Mor_{\Sh(\mathcal{D})}(f^{-1}\mathcal{F}, f^{-1}\mathcal{G})
& =
\Mor_{\Sh(\mathcal{C})}(\mathcal{F}, f_*f^{-1}\mathcal{G}) \\
& =
\Mor_{\Sh(\mathcal{C})}(\mathcal{F}, g^{-1}f^{-1}\mathcal{G}) \\
& =
\Mor_{\Sh(\mathcal{C})}(\mathcal{F}, \mathcal{G})
\end{align*}
（第一个等号由伴随得到），这证明了所需结论。
\end{proof}

\begin{example}
\label{sites-example-closed-map-cocontinuous-false}
设 $X$ 为拓扑空间，$i : Z  \to X$ 为一个子集（赋以诱导拓扑）的包含映射。
考虑函子 $u : X_{Zar} \to Z_{Zar}$，
$U \mapsto u(U) = Z \cap U$。乍看之下，该函子似乎也余连续。毕竟，既然
$Z$ 带有诱导拓扑，$U\cap Z$ 的任意覆盖难道不应来自 $X$ 中 $U$ 的某个
覆盖吗？并非如此！例如，若 $U \cap Z = \emptyset$ 会怎样？此时空覆盖是
$U \cap Z$ 的覆盖，而空覆盖只能由空覆盖加细。因此得到
$u$ 余连续 $\Rightarrow$ $X$ 的每个非空开集 $U$ 都与 $Z$ 有非空交。
% P01-ERRATA: P01-E0087 -- source has the malformed phrase "shouldn't any covering ... it come".
% P01-ERRATA: P01-E0088 -- source omits the verb in "$u$ cocontinuous".
但这并不充分。例如，若 $X = \mathbf{R}$ 是带通常拓扑的实数直线，且
$Z = \mathbf{R} \setminus \{0\}$，则 $Z$ 有开覆盖
$Z = \{x < 0\} \cup \bigcup_n \{1/n < x\}$，它不可能由 $X$ 的任何开覆盖
之限制来加细。
\end{example}







\section{具有右伴随的余连续函子}
\label{sites-section-cocontinuous-adjoint}

\noindent
设 $\mathcal{C}$ 和 $\mathcal{D}$ 为位点。设
$u : \mathcal{C} \to \mathcal{D}$ 和
$v : \mathcal{D} \to \mathcal{C}$ 为底层范畴的函子，并且 $v$ 是 $u$ 的
右伴随。在此情形中，若 $v$ 连续，则 $u$ 余连续（引理
\ref{sites-lemma-continuous-with-continuous-left-adjoint}）。若 $u$ 余连续，则常常
（但并非总是如此，见例 \ref{sites-example-not-equivalent}）$v$ 连续；若确实如此，
则 $v$ 定义一个位点态射，其相伴托扑斯态射与 $u$ 所定义的相同。

\begin{lemma}
\label{sites-lemma-have-functor-other-way}
设 $\mathcal{C}$ 和 $\mathcal{D}$ 为位点，设
$u : \mathcal{C} \to \mathcal{D}$ 和
$v : \mathcal{D} \to \mathcal{C}$ 为函子。假设 $u$ 余连续，且 $v$ 是
$u$ 的右伴随。令 $g : \Sh(\mathcal{C}) \to \Sh(\mathcal{D})$ 为与 $u$
相伴的托扑斯态射，见引理 \ref{sites-lemma-cocontinuous-morphism-topoi}。则：
\begin{enumerate}
\item 对 $\mathcal{C}$ 上的层 $\mathcal{F}$，层 $g_*\mathcal{F}$ 等于
预层 $v^p\mathcal{F}$；换言之，
$(g_*\mathcal{F})(V) = \mathcal{F}(v(V))$；
\item 对 $\mathcal{D}$ 上的层 $\mathcal{G}$，有
$g^{-1}\mathcal{G} = (v_p\mathcal{G})^\#$。
\end{enumerate}
\end{lemma}

\begin{proof}
对 (1) 中的 $\mathcal{F}$，有
$$
g_*\mathcal{F} =
{}_su\mathcal{F} =
{}_pu\mathcal{F} =
v^p\mathcal{F} =
\mathcal{F} \circ v
$$
第一个等号是引理 \ref{sites-lemma-cocontinuous-morphism-topoi}。第二个等号是
引理 \ref{sites-lemma-pu-sheaf}。第三个等号是引理
\ref{sites-lemma-adjoint-functors}。最后一个等号是第
\ref{sites-section-functoriality-PSh} 节中 $v^p$ 的定义。这证明 (1)。
对 (2) 中的 $\mathcal{G}$，有
$$
g^{-1}\mathcal{G} = (u^p\mathcal{G})^\# = (v_p\mathcal{G})^\#
$$
第一个等号是引理 \ref{sites-lemma-cocontinuous-morphism-topoi}。第二个等号是
引理 \ref{sites-lemma-adjoint-functors}。
\end{proof}

\begin{lemma}
\label{sites-lemma-have-functor-other-way-morphism}
记号与假设同引理 \ref{sites-lemma-have-functor-other-way}。若另外 $v$ 连续，则
$v$ 定义位点态射 $f : \mathcal{C} \to \mathcal{D}$，其相伴托扑斯态射
等于 $g$。
\end{lemma}

\begin{proof}
使用引理 \ref{sites-lemma-have-functor-other-way} 的结果而不再另行说明。为证明
$v$ 定义引理所述的位点态射 $f$，必须说明 $v_s$ 是正合函子（见定义
\ref{sites-definition-morphism-sites}）。由于
$v_s\mathcal{G} = (v_p\mathcal{G})^\# = g^{-1}\mathcal{G}$，这由 $g$ 是
托扑斯态射这一事实得出。于是 $f^{-1} = v_s = g^{-1}$，从而发现 $f = g$
作为托扑斯态射成立。
\end{proof}

\begin{example}
\label{sites-example-finer-topology-bis}
本例延续例 \ref{sites-example-finer-topology} 的讨论，并沿用其中的记号
$\mathcal{C}, \tau, \tau', \epsilon$。注意，恒等函子
$v : \mathcal{C}_{\tau'} \to \mathcal{C}_\tau$ 连续，而恒等函子
$u : \mathcal{C}_\tau \to \mathcal{C}_{\tau'}$ 余连续。此外，$u$ 是 $v$ 的
左伴随。因此，引理 \ref{sites-lemma-have-functor-other-way} 和
\ref{sites-lemma-have-functor-other-way-morphism} 的结果适用，从而 $v$ 定义位点
态射，即
$$
\epsilon : \mathcal{C}_\tau \longrightarrow \mathcal{C}_{\tau'}
$$
其相应托扑斯态射与余连续函子 $u$ 所伴随的托扑斯态射相同。
\end{example}

\begin{lemma}
\label{sites-lemma-continuous-with-continuous-left-adjoint}
设 $\mathcal{C}$ 和 $\mathcal{D}$ 为位点，并设
$v : \mathcal{D} \to \mathcal{C}$ 为连续函子。假设 $v$ 有左伴随
$u : \mathcal{C} \to \mathcal{D}$。则：
\begin{enumerate}
\item $u$ 余连续；
\item 引理 \ref{sites-lemma-have-functor-other-way} 和
\ref{sites-lemma-have-functor-other-way-morphism} 的结论成立。
\end{enumerate}
特别地，$v$ 定义位点态射 $f : \mathcal{C} \to \mathcal{D}$。
\end{lemma}

\begin{proof}
令 $U$ 为 $\mathcal{C}$ 的对象，并令
$\{V_j \to u(U)\}_{j \in J}$ 为 $\mathcal{D}$ 中的覆盖。则
$\{v(V_j) \to v(u(U))\}_{i \in I}$ 是 $\mathcal{C}$ 中的覆盖。通过伴随
映射 $U \to v(u(U))$，可以基变换为覆盖
$\{W_j \to U\}_{j \in J}$，其中
$W_j = v(V_j) \times_{v(u(U))} U$。以 $p_j : W_j \to v(V_j)$ 表示第一
投影，得到映射
% P01-ERRATA: P01-E0089 -- the pulled-forward covering switches from j in J to undefined i in I.
$$
u(W_j) \xrightarrow{u(p_j)} u(v(V_j)) \longrightarrow V_j
$$
其中第二个箭头是伴随映射。这确定具有固定靶的态射族之态射
$\{u(W_j) \to u(U)\} \to \{V_j \to u(U)\}$，表明 $u$ 的确余连续。
其余断言由引理 \ref{sites-lemma-have-functor-other-way} 和
\ref{sites-lemma-have-functor-other-way-morphism} 立即得出。
\end{proof}

\begin{example}
\label{sites-example-not-equivalent}
设 $\mathcal{C}$ 和 $\mathcal{D}$ 为位点。设
$u : \mathcal{C} \to \mathcal{D}$ 和
$v : \mathcal{D} \to \mathcal{C}$ 为底层范畴的函子，并且 $v$ 是 $u$ 的
右伴随。引理 \ref{sites-lemma-continuous-with-continuous-left-adjoint} 表明，若
$v$ 连续，则 $u$ 余连续。反过来，若 $u$ 余连续，不能由此断定 $v$ 连续。
% P01-ERRATA: P01-E0090 -- source misspells "cocontinuous" as "cocontinous".
将用概形的大 \'etale 位点和平滑位点给出这一现象的例子，不过大概也有
初等例子。具体而言，考虑概形 $S$ 以及位点
$(\Sch/S)_\etale$ 和 $(\Sch/S)_{smooth}$。可以假设这些位点具有同一底层
范畴，见 Topologies，备注 \ref{topologies-remark-choice-sites}。令
$u = v = \text{id}$。则作为从 $(\Sch/S)_\etale$ 到
$(\Sch/S)_{smooth}$ 的函子，$u$ 余连续，因为概形的每个平滑覆盖都可由
\'etale 覆盖加细，见 More on Morphisms，引理
\ref{more-morphisms-lemma-etale-dominates-smooth}。反之，从
$(\Sch/S)_{smooth}$ 到 $(\Sch/S)_\etale$ 的函子 $v$ 不连续，因为平滑
覆盖一般不是 \'etale 覆盖。
\end{example}





\section{具有左伴随的余连续函子}
\label{sites-section-cocontinuous-left-adjoint}

\noindent
余连续函子 $u$ 可能有左伴随 $w$。

\begin{lemma}
\label{sites-lemma-have-left-adjoint}
设 $\mathcal{C}$ 和 $\mathcal{D}$ 为位点。令
$g : \Sh(\mathcal{C}) \to \Sh(\mathcal{D})$ 为与连续且余连续的函子
$u : \mathcal{C} \to \mathcal{D}$ 相伴的托扑斯态射，见引理
\ref{sites-lemma-cocontinuous-morphism-topoi} 和 \ref{sites-lemma-when-shriek}。
\begin{enumerate}
\item 若 $w : \mathcal{D} \to \mathcal{C}$ 是 $u$ 的左伴随，则：
\begin{enumerate}
\item $g_!\mathcal{F}$ 是与预层 $w^p\mathcal{F}$ 相伴的层；
\item $g_!$ 正合。
\end{enumerate}
\item 若 $w$ 是连续的左伴随，则 $g_!$ 有左伴随。
\item 若 $w$ 是余连续的左伴随，则 $g_! = h^{-1}$ 且
$g^{-1} = h_*$，其中
$h : \Sh(\mathcal{D}) \to \Sh(\mathcal{C})$ 是与 $w$ 相伴的托扑斯态射。
\end{enumerate}
\end{lemma}

\begin{proof}
回忆一下，$g_!\mathcal{F}$ 是 $u_p\mathcal{F}$ 的层化。因此，(1)(a) 由
引理 \ref{sites-lemma-adjoint-functors} 中 $u_p = w^p$ 这一事实得出。

\medskip\noindent
为说明 (1)(b)，注意 $g_!$ 作为左伴随与所有余极限交换（Categories，引理
\ref{categories-lemma-adjoint-exact}）。令
$i \mapsto \mathcal{F}_i$ 为 $\Sh(\mathcal{C})$ 中的有限图表。则
$\lim \mathcal{F}_i$ 在预层范畴中计算（引理
\ref{sites-lemma-limit-sheaf}）。因为 $w^p$ 是右伴随（引理
\ref{sites-lemma-adjoints-u}），有
$w^p \lim \mathcal{F}_i = \lim w^p\mathcal{F}_i$。层化正合（引理
\ref{sites-lemma-sheafification-exact}），故结论由 (1)(a) 得出。

\medskip\noindent
假设 $w$ 连续。则 $g_! = (w^p\ )^\# = w^s$，但层化不是必需的，并且
$g_!$ 有左伴随 $w_s$，见引理 \ref{sites-lemma-pushforward-sheaf} 和
\ref{sites-lemma-adjoint-sheaves}。

\medskip\noindent
假设 $w$ 余连续。等式 $g_! = h^{-1}$ 由 (1)(a) 和定义得出。等式
$g^{-1} = h_*$ 由 $g_! = h^{-1}$ 和伴随函子的唯一性得出。或者，也可由
引理 \ref{sites-lemma-have-functor-other-way} 推出。
\end{proof}









\section{下惊叹号的存在性}
\label{sites-section-lower-shriek}

\noindent
本节讨论托扑斯态射 $f$ 的一些情形；在这些情形中，$f^{-1}$ 有左伴随
$f_!$。

\begin{lemma}
\label{sites-lemma-existence-lower-shriek}
设 $\mathcal{C}$、$\mathcal{D}$ 为两个位点。设
$f : \Sh(\mathcal{D}) \to \Sh(\mathcal{C})$ 为托扑斯态射。设
$E \subset \Ob(\mathcal{D})$ 为子集，满足：
\begin{enumerate}
\item 对 $V \in E$，存在 $\mathcal{C}$ 上的层 $\mathcal{G}$，使得
$f^{-1}\mathcal{F}(V) =
\Mor_{\Sh(\mathcal{C})}(\mathcal{G}, \mathcal{F})$ 对
$\Sh(\mathcal{C})$ 中的 $\mathcal{F}$ 函子地成立；
\item $\mathcal{D}$ 的每个对象都有一个由 $E$ 中对象组成的覆盖。
\end{enumerate}
则 $f^{-1}$ 有左伴随 $f_!$。
\end{lemma}

\begin{proof}
由米田引理（Categories，引理 \ref{categories-lemma-yoneda}），与
$V \in E$ 对应的层 $\mathcal{G}_V$ 由公式
$f^{-1}\mathcal{F}(V) =
\Mor_{\Sh(\mathcal{C})}(\mathcal{G}_V, \mathcal{F})$ 在唯一同构意义下
确定。回忆一下，
$f^{-1}\mathcal{F}(V) =
\Mor_{\Sh(\mathcal{D})}(h_V^\#, f^{-1}\mathcal{F})$。以
$i_V : h_V^\# \to f^{-1}\mathcal{G}_V$ 表示与
$\Mor(\mathcal{G}_V, \mathcal{G}_V)$ 中的 $\text{id}$ 对应的映射。(1) 中的
函子性蕴含该双射由下式给出：
$$
\Mor_{\Sh(\mathcal{C})}(\mathcal{G}_V, \mathcal{F}) \to
\Mor_{\Sh(\mathcal{D})}(h_V^\#, f^{-1}\mathcal{F}),\quad
\varphi \mapsto f^{-1}\varphi \circ i_V
$$
对任意 $V_1, V_2 \in E$，存在典范映射
$$
\Mor_{\Sh(\mathcal{D})}(h^\#_{V_2}, h^\#_{V_1})
\to
\Hom_{\Sh(\mathcal{C})}(\mathcal{G}_{V_2}, \mathcal{G}_{V_1}),\quad
\varphi \mapsto f_!(\varphi)
$$
其特征为
$f^{-1}(f_!(\varphi)) \circ i_{V_2} = i_{V_1} \circ \varphi$。注意，
$\varphi \mapsto f_!(\varphi)$ 与复合相容；这可直接由上述刻画看出。因此，
$h_V^\# \mapsto \mathcal{G}_V$ 和
$\varphi \mapsto f_!\varphi$ 构成一个以 $\Sh(\mathcal{D})$ 的全子范畴
为定义域的函子，其中该全子范畴的对象是各个 $h_V^\#$。

\medskip\noindent
设 $J$ 为集合，并设 $J \to E$，$j \mapsto V_j$ 为映射。利用上述双射的
乘积，得到函子性双射
$$
\Mor_{\Sh(\mathcal{C})}(\coprod \mathcal{G}_{V_j}, \mathcal{F})
\longrightarrow
\Mor_{\Sh(\mathcal{D})}(\coprod h_{V_j}^\#, f^{-1}\mathcal{F})
$$
因此，可以把函子 $f_!$ 扩张到 $\Sh(\mathcal{D})$ 的全子范畴，该全子
范畴的对象是 $V \in E$ 的 $h_V^\#$ 的余积。

\medskip\noindent
给定 $\mathcal{D}$ 上任意层 $\mathcal{H}$，选择余等化子图
$$
\xymatrix{
\mathcal{H}_1 \ar@<1ex>[r] \ar@<-1ex>[r] &
\mathcal{H}_0 \ar[r] &
\mathcal{H}
}
$$
其中 $\mathcal{H}_i = \coprod h_{V_{i, j}}^\#$ 是一个余积，且
$V_{i, j} \in E$。由假设 (2)，这是可能的，见引理
\ref{sites-lemma-sheaf-coequalizer-representable}（对于担心集合论问题的读者，
注意引理 \ref{sites-lemma-sheaf-coequalizer-representable} 中给出的构造是典范的）。
把 $f_!(\mathcal{H})$ 定义为 $\mathcal{C}$ 上使
$$
\xymatrix{
f_!\mathcal{H}_1 \ar@<1ex>[r] \ar@<-1ex>[r] &
f_!\mathcal{H}_0 \ar[r] &
f_!\mathcal{H}
}
$$
成为余等化子图的层。于是
\begin{align*}
\Mor(f_!\mathcal{H}, \mathcal{F})
& =
\text{Equalizer}(
\xymatrix{
\Mor(f_!\mathcal{H}_0, \mathcal{F}) \ar@<1ex>[r] \ar@<-1ex>[r] &
\Mor(f_!\mathcal{H}_1, \mathcal{F})
}
) \\
& =
\text{Equalizer}(
\xymatrix{
\Mor(\mathcal{H}_0, f^{-1}\mathcal{F}) \ar@<1ex>[r] \ar@<-1ex>[r] &
\Mor(\mathcal{H}_1, f^{-1}\mathcal{F})
}
) \\
& =
\Hom(\mathcal{H}, f^{-1}\mathcal{F})
\end{align*}
因此可把 $f_!$ 扩张到 $\mathcal{D}$ 上整个层范畴。
\end{proof}








\section{局部化}
\label{sites-section-localize}

\noindent
设 $\mathcal{C}$ 为位点，并设 $U \in \Ob(\mathcal{C})$。关于 $U$ 上对象
所成范畴 $\mathcal{C}/U$ 的定义，见 Categories，例
\ref{categories-example-category-over-X}。如下把 $\mathcal{C}/U$ 变为位点：
当且仅当 $U$ 上对象的态射族 $\{V_j \to V\}$ 是 $\mathcal{C}$ 中的覆盖时，
才宣称它是 $\mathcal{C}/U$ 的覆盖。考虑遗忘函子
$$
j_U : \mathcal{C}/U \longrightarrow \mathcal{C}.
$$
它显然余连续且连续。因此，由引理 \ref{sites-lemma-when-shriek}，得到托扑斯态射
$$
j_U : \Sh(\mathcal{C}/U) \longrightarrow \Sh(\mathcal{C})
$$
由 $j_U^{-1}$ 和 $j_{U*}$ 给出，并有 $j_U^{-1}$ 的左伴随函子 $j_{U!}$。

\begin{definition}
\label{sites-definition-localize}
设 $\mathcal{C}$ 为位点，并设 $U \in \Ob(\mathcal{C})$。
\begin{enumerate}
\item 位点 $\mathcal{C}/U$ 称为{\it 位点 $\mathcal{C}$ 在对象 $U$ 处的
局部化}。
\item 托扑斯态射
$j_U : \Sh(\mathcal{C}/U) \to \Sh(\mathcal{C})$ 称为{\it 局部化态射}。
\item 函子 $j_{U*}$ 称为{\it 直像函子}。
\item 对 $\mathcal{C}$ 上的层 $\mathcal{F}$，层
$j_U^{-1}\mathcal{F}$ 称为 $\mathcal{F}$ {\it 到 $\mathcal{C}/U$ 的限制}。
\item 对 $\mathcal{C}/U$ 上的层 $\mathcal{G}$，层
$j_{U!}\mathcal{G}$ 称为 $\mathcal{G}$ {\it 以空集的延拓}。
\end{enumerate}
\end{definition}

\noindent
如所预期，限制 $j_U^{-1}\mathcal{F}$ 是由规则
$j_U^{-1}\mathcal{F}(X/U) = \mathcal{F}(X)$ 定义的层。在此情形中，以空集
的延拓也有非常简单的描述，如下。

\begin{lemma}
\label{sites-lemma-describe-j-shriek}
设 $\mathcal{C}$ 为位点，$U \in \Ob(\mathcal{C})$，并设
$\mathcal{G}$ 为 $\mathcal{C}/U$ 上的预层。则
$j_{U!}(\mathcal{G}^\#)$ 是与下述预层相伴的层：
$$
V
\longmapsto
\coprod\nolimits_{\varphi \in \Mor_\mathcal{C}(V, U)}
\mathcal{G}(V \xrightarrow{\varphi} U)
$$
其限制映射显然。
\end{lemma}

\begin{proof}
由引理 \ref{sites-lemma-when-shriek}，有
$j_{U!}(\mathcal{G}^\#) = ((j_U)_p\mathcal{G}^\#)^\#$。由引理
\ref{sites-lemma-technical-up}，它等于 $((j_U)_p\mathcal{G})^\#$。因此，只需证明
对 $\mathcal{C}/U$ 上任意预层 $\mathcal{G}$，$(j_U)_p$ 由上述公式给出。
现在，根据第 \ref{sites-section-functoriality-PSh} 节中的定义，有
% P01-ERRATA: P01-E0091 -- source begins the sentence with the editing remnant "OK, and".
$$
(j_U)_p\mathcal{G}(V)
=
\colim_{(W/U, V \to W)} \mathcal{G}(W)
$$
显然，对每个 $\varphi : V \to U$，序对 $(W/U, V \to W)$ 所成范畴都有
对象 $O_\varphi = (\varphi : V \to U, \text{id} : V \to V)$；而且从某个
$O_\varphi$ 到任意对象恰有一个态射。结论随之成立。
\end{proof}

\begin{lemma}
\label{sites-lemma-describe-j-shriek-representable}
设 $\mathcal{C}$ 为位点，$U \in \Ob(\mathcal{C})$，并设 $X/U$ 为
$\mathcal{C}/U$ 的对象。则 $j_{U!}(h_{X/U}^\#) = h_X^\#$。
\end{lemma}

\begin{proof}
以 $p : X \to U$ 表示 $X$ 的结构态射。由引理
\ref{sites-lemma-describe-j-shriek}，$j_{U!}(h_{X/U}^\#)$ 是与下述预层相伴的层：
$$
V
\longmapsto
\coprod\nolimits_{\varphi \in \Mor_\mathcal{C}(V, U)}
\{\psi : V \to X \mid p \circ \psi = \varphi\}
$$
这显然就是 $\Mor_\mathcal{C}(V, X)$，故引理成立。
\end{proof}

\noindent
由上述任一引理，有 $j_{U!}(*) = h_U^\#$。因此，对 $\mathcal{C}/U$ 上的
每个层 $\mathcal{G}$，都有层的典范映射
$j_{U!}\mathcal{G} \to h_U^\#$。这刻画了 $j_{U!}$ 本质像中的层。

\begin{lemma}
\label{sites-lemma-essential-image-j-shriek}
设 $\mathcal{C}$ 为位点，$U \in \Ob(\mathcal{C})$。函子 $j_{U!}$ 给出
范畴等价
$$
\Sh(\mathcal{C}/U)
\longrightarrow
\Sh(\mathcal{C})/h_U^\#
$$
\end{lemma}

\begin{proof}
把 $\mathcal{C}/U$ 的对象记作序对 $(X, a)$，其中 $X$ 是
$\mathcal{C}$ 的对象，$a : X \to U$ 是 $\mathcal{C}$ 的态射。类似地，
$\Sh(\mathcal{C})/h_U^\#$ 的对象是序对 $(\mathcal{F}, \varphi)$。函子
$\Sh(\mathcal{C}/U) \to \Sh(\mathcal{C})/h_U^\#$ 把 $\mathcal{G}$ 送到
序对 $(j_{U!}\mathcal{G}, \gamma)$，其中 $\gamma$ 是
$j_{U!}\mathcal{G} \to j_{U!}*$ 与等同
$j_{U!}* = h_U^\#$ 的复合。

\medskip\noindent
构造从 $\Sh(\mathcal{C})/h_U^\#$ 到 $\Sh(\mathcal{C}/U)$ 的函子。假设
给定 $(\mathcal{F}, \varphi)$。对 $\mathcal{C}/U$ 的对象 $(X, a)$，考虑
集合 $\mathcal{F}_\varphi(X, a)$；它由 $s \in \mathcal{F}(X)$ 中经
$\varphi$ 映到 $a \in \Mor_\mathcal{C}(X, U) = h_U(X)$ 在
$h_U^\#(X)$ 中之像的元素组成。容易看出，
$(X, a) \mapsto \mathcal{F}_\varphi(X, a)$ 是 $\mathcal{C}/U$ 上的层。
显然，规则 $(\mathcal{F}, \varphi) \mapsto \mathcal{F}_\varphi$ 定义函子
$\Sh(\mathcal{C})/h_U^\# \to \Sh(\mathcal{C}/U)$。

\medskip\noindent
还考虑函子
$\textit{PSh}(\mathcal{C})/h_U \to \textit{PSh}(\mathcal{C}/U)$，
$(\mathcal{F}, \varphi) \mapsto \mathcal{F}_\varphi$，其中
$\mathcal{F}_\varphi(X, a)$ 定义为 $\mathcal{F}(X)$ 中映到
$a \in h_U(X)$ 的元素所成的集合。断言图
$$
\xymatrix{
\textit{PSh}(\mathcal{C})/h_U \ar[r] \ar[d] &
\textit{PSh}(\mathcal{C}/U) \ar[d] \\
\Sh(\mathcal{C})/h_U^\# \ar[r] &
\Sh(\mathcal{C}/U)
}
$$
交换，其中竖直箭头由层化给出。为看出这一点\footnote{另一种做法是，
对预层用笛卡尔图
$$
\vcenter{
\xymatrix{
\mathcal{F}_\varphi \ar[r] \ar[d] & {*} \ar[d] \\
\mathcal{F}|_{\mathcal{C}/U} \ar[r] & h_U|_{\mathcal{C}/U}
}
}
\quad\text{对预层以及}\quad
\vcenter{
\xymatrix{
\mathcal{F}_\varphi \ar[r] \ar[d] & {*} \ar[d] \\
\mathcal{F}|_{\mathcal{C}/U} \ar[r] & h_U^\#|_{\mathcal{C}/U}
}
}
$$
描述 $\mathcal{F}_\varphi$，并对层使用后一图，再利用到
$\mathcal{C}/U$ 的限制与层化交换这一事实。}，只需证明该构造与引理
\ref{sites-lemma-plus-presheaf}、\ref{sites-lemma-plus-functorial} 及定理
\ref{sites-theorem-plus} 中的函子
$\mathcal{F} \mapsto \mathcal{F}^+$ 交换。与
$\mathcal{F} \mapsto \mathcal{F}^+$ 的交换性来自如下事实：给定
$(X, a)$，$\mathcal{C}/U$ 中 $(X, a)$ 的覆盖范畴与 $\mathcal{C}$ 中 $X$
的覆盖范畴典范等同。

\medskip\noindent
接着，令 $\textit{PSh}(\mathcal{C}/U) \to
\textit{PSh}(\mathcal{C})/h_U$ 把 $\mathcal{G}$ 送到序对
$(j_{U!}^{PSh}\mathcal{G}, \gamma)$，其中
$j_{U!}^{PSh}\mathcal{G}$ 是由引理 \ref{sites-lemma-describe-j-shriek} 中公式
定义的预层，而 $\gamma$ 是
$j_{U!}^{PSh}\mathcal{G} \to j_{U!}*$ 与等同
$j_{U!}^{PSh}* = h_U$ 的复合（由公式显然可知）。
% P01-ERRATA: P01-E0092 -- source omits "is" in "where j_{U!}^{PSh}G the presheaf".
% P01-ERRATA: P01-E0093 -- the target of the first map drops the PSh superscript, inconsistent with the stated identification.
于是立即清楚可知图
$$
\xymatrix{
\textit{PSh}(\mathcal{C}/U) \ar[r] \ar[d] &
\textit{PSh}(\mathcal{C})/h_U \ar[d] \\
\Sh(\mathcal{C}/U) \ar[r] &
\Sh(\mathcal{C})/h_U^\#
}
$$
交换，其中竖直箭头是层化。综合上述各点，只需证明：对
$\textit{PSh}(\mathcal{C}/U)$ 中的 $\mathcal{G}$，有函子性同构
$(j_{U!}^{PSh}\mathcal{G})_\gamma = \mathcal{G}$；并且对
$\textit{PSh}(\mathcal{C})/h_U$ 中的 $(\mathcal{F}, \varphi)$，有
$j_{U!}^{PSh}\mathcal{F}_\varphi = \mathcal{F}$。预层
$(j_{U!}^{PSh}\mathcal{G})_\gamma$ 在 $(X, a)$ 上的值是映射
$$
\coprod\nolimits_{a' : X \to U} \mathcal{G}(X, a') \to \Mor_\mathcal{C}(X, U)
$$
在 $a$ 上的纤维，即 $\mathcal{G}(X, a)$。这证明第一个等式。预层
$j_{U!}^{PSh}\mathcal{F}_\varphi$ 在 $X$ 上的值是
% P01-ERRATA: P01-E0094 -- source has the duplicated phrase "is on X is".
$$
\coprod\nolimits_{a : X \to U} \mathcal{F}_\varphi(X, a) =
\mathcal{F}(X)
$$
因为给定集合映射 $S \to S'$，集合 $S$ 是其各纤维的无交并。
\end{proof}

\noindent
引理 \ref{sites-lemma-essential-image-j-shriek} 说明，函子 $j_{U!}$ 是复合
$$
\Sh(\mathcal{C}/U) \rightarrow
\Sh(\mathcal{C})/h_U^\# \rightarrow
\Sh(\mathcal{C})
$$
其中第一个箭头是等价。

\begin{lemma}
\label{sites-lemma-j-shriek-commutes-equalizers-fibre-products}
设 $\mathcal{C}$ 为位点，$U \in \Ob(\mathcal{C})$。函子 $j_{U!}$ 与纤维积和
等化子交换（更一般地，与有限连通极限交换）。特别地，若
$\mathcal{F} \subset \mathcal{F}'$ 于 $\Sh(\mathcal{C}/U)$ 中，则
$j_{U!}\mathcal{F} \subset j_{U!}\mathcal{F}'$。
\end{lemma}

\begin{proof}
利用引理 \ref{sites-lemma-essential-image-j-shriek} 以及范畴等价与所有极限交换这一
事实，问题归结为函子
$\Sh(\mathcal{C})/h_U^\# \rightarrow \Sh(\mathcal{C})$ 与纤维积和等化子
交换。也可以使用引理 \ref{sites-lemma-describe-j-shriek} 中对 $j_{U!}$ 的描述，
并利用层化是正合的，直接证明这一点。（此外，当 $\mathcal{C}$ 有纤维积和
等化子时，结论也由引理 \ref{sites-lemma-preserve-equalizers} 得出。）
\end{proof}

\begin{lemma}
\label{sites-lemma-j-shriek-reflects-injections-surjections}
设 $\mathcal{C}$ 为位点，$U \in \Ob(\mathcal{C})$。函子 $j_{U!}$ 反映单射与
满射。
\end{lemma}

\begin{proof}
我们必须证明 $j_{U!}$ 反映单态射与满态射，见引理
\ref{sites-lemma-mono-epi-sheaves}。利用引理
\ref{sites-lemma-essential-image-j-shriek}，这归结为函子
$\Sh(\mathcal{C})/h_U^\# \to \Sh(\mathcal{C})$ 反映单态射与满态射这一
事实。
\end{proof}

\begin{lemma}
\label{sites-lemma-compute-j-shriek-restrict}
设 $\mathcal{C}$ 为位点，$U \in \Ob(\mathcal{C})$。对 $\mathcal{C}$ 上的任意
层 $\mathcal{F}$，有
$j_{U!}j_U^{-1}\mathcal{F} = \mathcal{F} \times h_U^\#$。
\end{lemma}

\begin{proof}
这由引理 \ref{sites-lemma-describe-j-shriek} 中对 $j_{U!}$ 的描述显然可知。
\end{proof}

\begin{lemma}
\label{sites-lemma-relocalize}
设 $\mathcal{C}$ 为位点。设 $f : V \to U$ 是 $\mathcal{C}$ 的态射。则存在
连续且余连续函子构成的交换图
$$
\xymatrix{
\mathcal{C}/V \ar[rd]_{j_V} \ar[rr]_j & &
\mathcal{C}/U \ar[ld]^{j_U} \\
& \mathcal{C} &
}
$$
函子 $j : \mathcal{C}/V \to \mathcal{C}/U$，
$(a : W \to V) \mapsto (f \circ a : W \to U)$，通过等同
$(\mathcal{C}/U)/(V/U) = \mathcal{C}/V$ 与函子
$j_{V/U} : (\mathcal{C}/U)/(V/U) \to \mathcal{C}/U$ 等同。此外，有
$j_{V!} = j_{U!} \circ j_!$、
$j_V^{-1} = j^{-1} \circ j_U^{-1}$ 以及
$j_{V*} = j_{U*} \circ j_*$。
\end{lemma}

\begin{proof}
图的交换性是立即的。$j$ 与 $j_{V/U}$ 相符由定义得出。由引理
\ref{sites-lemma-composition-cocontinuous}，可见下列拓扑斯态射图
\begin{equation}
\label{sites-equation-relocalize}
\vcenter{
\xymatrix{
\Sh(\mathcal{C}/V) \ar[rd]_{j_V} \ar[rr]_j & &
\Sh(\mathcal{C}/U) \ar[ld]^{j_U} \\
& \Sh(\mathcal{C}) &
}
}
\end{equation}
交换。这证明
$j_V^{-1} = j^{-1} \circ j_U^{-1}$ 和 $j_{V*} = j_{U*} \circ j_*$。
等式 $j_{V!} = j_{U!} \circ j_!$ 由伴随性质形式地得出。
\end{proof}

\begin{lemma}
\label{sites-lemma-relocalize-explicit}
记号 $\mathcal{C}$、$f : V \to U$、$j_U$、$j_V$ 和 $j$ 如引理
\ref{sites-lemma-relocalize}。通过引理
\ref{sites-lemma-essential-image-j-shriek} 的等同
$\Sh(\mathcal{C}/V) = \Sh(\mathcal{C})/h_V^\#$ 与
$\Sh(\mathcal{C}/U) = \Sh(\mathcal{C})/h_U^\#$，有
\begin{enumerate}
\item 函子 $j^{-1}$ 有如下描述：
$$
j^{-1}(\mathcal{H} \xrightarrow{\varphi} h_U^\#)
=
(\mathcal{H} \times_{\varphi, h_U^\#, f} h_V^\# \to h_V^\#).
$$
\item 函子 $j_!$ 有如下描述：
$$
j_!(\mathcal{H} \xrightarrow{\varphi} h_V^\#) =
(\mathcal{H} \xrightarrow{h_f \circ \varphi} h_U^\#)
$$
\end{enumerate}
\end{lemma}

\begin{proof}
先证明 (2)。回忆等同
$\Sh(\mathcal{C}/V) \to \Sh(\mathcal{C})/h_V^\#$ 把 $\mathcal{G}$ 送到
$j_{V!}\mathcal{G} \to j_{V!}(*) = h_V^\#$，而
$\Sh(\mathcal{C}/U) \to \Sh(\mathcal{C})/h_U^\#$ 的情形也类似。因此，
$j_!\mathcal{G}$ 被送到
$j_{U!}(j_!\mathcal{G}) \to j_{U!}(*) = h_U^\#$；又由引理
\ref{sites-lemma-relocalize} 有 $j_{U!}j_! = j_{V!}$，故 (2) 成立。

\medskip\noindent
读者现在可以利用 $j^{-1}$ 是 $j_!$ 的右伴随，并利用 (1) 中的规则是 (2)
中规则的右伴随，来证明 (1)。下面给出直接证明。假设
$\varphi : \mathcal{H} \to h_U^\#$ 是 $\Sh(\mathcal{C})/h_U^\#$ 的对象。
由引理 \ref{sites-lemma-essential-image-j-shriek} 的证明，它对应于
$\mathcal{C}/U$ 上按规则
$$
(a : W \to U)
\longmapsto
\{ s \in \mathcal{H}(W) \mid \varphi(s) = a\}
$$
定义的层 $\mathcal{H}_\varphi$。这个 $\mathcal{C}/U$ 上的层到
$\mathcal{C}/V$ 的拉回
$j^{-1}\mathcal{H}_\varphi$ 由规则
$$
(a : W \to V)
\longmapsto
\{ s \in \mathcal{H}(W) \mid \varphi(s) = f \circ a\}
$$
给出，这是因为把 $j^{-1} = j_{U/V}^{-1}$ 描述为
$\mathcal{H}_\varphi$ 到 $\mathcal{C}/V$ 的限制。
% P01-ERRATA: P01-E0095 -- source writes j_{U/V}^{-1}, while the defined localization functor is j=j_{V/U}.
另一方面，把该规则应用于 $\Sh(\mathcal{C})/h_V^\#$ 的对象
$$
\xymatrix{
\mathcal{H}' = \mathcal{H} \times_{\varphi, h_U^\#, f} h_V^\#
\ar[rr]^-{\varphi'} & & h_V^\#
}
$$
便得到 $\mathcal{H}'_{\varphi'}$，它由
\begin{align*}
(a : W \to V)
\longmapsto
&  \{ s' \in \mathcal{H}'(W) \mid \varphi'(s') = a\} \\
= &
\{ (s, a') \in \mathcal{H}(W) \times h_V^\#(W) \mid
a' = a \text{ 且 } \varphi(s) = f \circ a'\}
\end{align*}
给出。这恰与上面对 $j^{-1}\mathcal{H}_\varphi$ 的描述是同一规则。
\end{proof}

\begin{remark}
\label{sites-remark-localize-presheaves}
局部化与预层。设 $\mathcal{C}$ 为范畴，$U$ 是 $\mathcal{C}$ 的对象。严格地说，
函子 $j_U^{-1}$、$j_{U*}$ 与 $j_{U!}$ 尚未对预层定义。不过，我们当然可以把
预层视为 $\mathcal{C}$ 上混沌拓扑的层（见例 \ref{sites-example-indiscrete}）。因此也
得到函子
$$
j_U^{-1} :
\textit{PSh}(\mathcal{C})
\longrightarrow
\textit{PSh}(\mathcal{C}/U)
$$
以及函子
$$
j_{U*}, j_{U!} :
\textit{PSh}(\mathcal{C}/U)
\longrightarrow
\textit{PSh}(\mathcal{C})
$$
它们分别是 $j_U^{-1}$ 的右、左伴随。由引理
\ref{sites-lemma-describe-j-shriek} 可见，$j_{U!}\mathcal{G}$ 是预层
$$
V \longmapsto
\coprod\nolimits_{\varphi \in \Mor_\mathcal{C}(V, U)}
\mathcal{G}(V \xrightarrow{\varphi} U)
$$
此外，函子 $j_{U!}$ 与纤维积和等化子交换。
\end{remark}

\begin{remark}
\label{sites-remark-localization-cartesian-cocontinuous}
设 $\mathcal{C}$ 为位点，$U \to V$ 是 $\mathcal{C}$ 的态射。余连续函子
$\mathcal{C}/U \to \mathcal{C}$ 与
$j : \mathcal{C}/U \to \mathcal{C}/V$（引理 \ref{sites-lemma-relocalize}）满足
备注 \ref{sites-remark-cartesian-cocontinuous} 的性质 $P$。例如，若有对象
$(X/U)$、$(W/V)$、态射 $g : j(X/U) \to (W/V)$ 以及覆盖
$\{f_i  : (W_i/V) \to (W/V)\}$，则
$(X \times_W W_i/U)$ 是 $(X/U) \times_{g, (W/V), f_i} (W_i/V)$ 的一个化身，
而族 $\{(X \times_W W_i/U) \to (X/U)\}$ 是 $\mathcal{C}/U$ 的覆盖。
\end{remark}






\section{层的粘合}
\label{sites-section-glueing-sheaves}

\noindent
本节类似于《层》，第 \ref{sheaves-section-glueing-sheaves} 节。

\begin{lemma}
\label{sites-lemma-glue-maps}
\begin{slogan}
层的映射可以粘合。
\end{slogan}
设 $\mathcal{C}$ 为位点，$\{U_i \to U\}$ 是 $\mathcal{C}$ 的覆盖，
$\mathcal{F}$、$\mathcal{G}$ 是 $\mathcal{C}$ 上的层。给定一族层映射
$$
\varphi_i :
\mathcal{F}|_{\mathcal{C}/U_i}
\longrightarrow
\mathcal{G}|_{\mathcal{C}/U_i}
$$
使得对所有 $i, j \in I$，映射 $\varphi_i, \varphi_j$ 在相应交叠上的
限制为同一映射
$\varphi_{ij} : \mathcal{F}|_{\mathcal{C}/U_i \times_U U_j} \to
\mathcal{G}|_{\mathcal{C}/U_i \times_U U_j}$，则存在唯一的层映射
$$
\varphi :
\mathcal{F}|_{\mathcal{C}/U}
\longrightarrow
\mathcal{G}|_{\mathcal{C}/U}
$$
其到每个 $\mathcal{C}/U_i$ 的限制都与 $\varphi_i$ 相符。
\end{lemma}

\begin{proof}
引理中所用的限制是引理 \ref{sites-lemma-relocalize} 中的限制。设 $V/U$ 是
$\mathcal{C}/U$ 的对象。令 $V_i = U_i \times_U V$，并记
$\mathcal{V} = \{V_i \to V\}$。注意
$(U_i \times_U U_j) \times_U V = V_i \times_V V_j$。于是有
$\mathcal{F}|_{\mathcal{C}/U_i}(V_i/U_i) = \mathcal{F}(V_i)$，以及
$\mathcal{F}|_{\mathcal{C}/U_i \times_U U_j}(V_i \times_V V_j/U_i \times_U U_j)
= \mathcal{F}(V_i \times_V V_j)$，对 $\mathcal{G}$ 也类似。因此，可以把
$V$ 上截面间的 $\varphi$ 定义为图中的虚线箭头：
$$
\xymatrix{
\mathcal{F}(V) \ar@{=}[r] &
H^0(\mathcal{V}, \mathcal{F}) \ar@{..>}[d] \ar[r] &
\prod \mathcal{F}(V_i)
\ar[d]_{\prod \varphi_i}
\ar@<1ex>[r] \ar@<-1ex>[r] &
\prod \mathcal{F}(V_i \times_V V_j) \ar[d]_{\prod \varphi_{ij}} \\
\mathcal{G}(V) \ar@{=}[r] &
H^0(\mathcal{V}, \mathcal{G}) \ar[r] &
\prod \mathcal{G}(V_i)
\ar@<1ex>[r] \ar@<-1ex>[r] &
\prod \mathcal{G}(V_i \times_V V_j)
}
$$
等号来自层条件；记号 $H^0(\mathcal{V}, -)$ 见第
\ref{sites-section-sheafification} 节。我们略去这些映射与限制映射相容的验证。
\end{proof}

\noindent
前一引理蕴含：给定位点 $\mathcal{C}$ 上的两个层 $\mathcal{F}$、
$\mathcal{G}$，规则
$$
U \longmapsto
\Mor_{\Sh(\mathcal{C}/U)}(
\mathcal{F}|_{\mathcal{C}/U}, \mathcal{G}|_{\mathcal{C}/U})
$$
定义一个层 $\SheafHom(\mathcal{F},\mathcal{G})$。这是一种
{\it 内 Hom 层}。在集合层的语境中很少使用它，而在模层的语境中更常用；
见《位点上的模》，第 \ref{sites-modules-section-internal-hom} 节。

\begin{lemma}
\label{sites-lemma-internal-hom-sheaf}
\begin{slogan}
位点上的层范畴是笛卡尔闭的
\end{slogan}
设 $\mathcal{C}$ 为位点，$\mathcal{F}$、$\mathcal{G}$ 与 $\mathcal{H}$ 是
$\mathcal{C}$ 上的层。存在典范双射
$$
\Mor_{\Sh(\mathcal{C})}(\mathcal{F}\times\mathcal{G},\mathcal{H}) =
\Mor_{\Sh(\mathcal{C})}(\mathcal{F},\SheafHom(\mathcal{G},\mathcal{H}))
$$
并且它关于三个变元都是函子性的。
\end{lemma}

\begin{proof}
引理断言函子 $-\times\mathcal{G}$ 与 $\SheafHom(\mathcal{G},-)$ 互为伴随。
为证明这一点，我们使用单位和余单位的概念，见《范畴》，第
\ref{categories-section-adjoint} 节。单位
$$
\eta_\mathcal{F} :
\mathcal{F}
\longrightarrow
\SheafHom(\mathcal{G},\mathcal{F}\times\mathcal{G})
$$
把 $s \in \mathcal{F}(U)$ 送到映射
$\mathcal{G}|_{\mathcal{C}/U} \to
\mathcal{F}|_{\mathcal{C}/U}\times\mathcal{G}|_{\mathcal{C}/U}$；该映射在
$V/U$ 上由
$$
\mathcal{G}(V) \longrightarrow \mathcal{F}(V)\times \mathcal{G}(V), \quad
t \longmapsto (s|_{V},t).
$$
给出。余单位
$$
\epsilon_{\mathcal{H}} :
\SheafHom(\mathcal{G}, \mathcal{H}) \times \mathcal{G}
\longrightarrow
\mathcal{H}
$$
是求值映射。它由规则
$$
\Mor_{\Sh(\mathcal{C}/U)}(
\mathcal{G}|_{\mathcal{C}/U}, \mathcal{H}|_{\mathcal{C}/U})
\times \mathcal{G}(U)
\longrightarrow
\mathcal{H}(U),\quad
(\varphi, s) \longmapsto \varphi(s).
$$
给出。于是，对每个 $\varphi : \mathcal{F} \times \mathcal{G} \to \mathcal{H}$，
相应的态射 $\mathcal{F} \to \SheafHom(\mathcal{G},\mathcal{H})$ 把每个截面
$s \in \mathcal{F}(U)$ 送到 $\mathcal{C}/U$ 上的层态射；该态射在 $V/U$
上的截面间由
$$
\mathcal{G}(V) \longrightarrow \mathcal{H}(V),\quad
t \longmapsto \varphi(s|_V, t).
$$
给出。反过来，对每个
$\psi : \mathcal{F} \to \SheafHom(\mathcal{G}, \mathcal{H})$，相应的态射
$\mathcal{F} \times \mathcal{G} \to \mathcal{H}$ 在对象 $U$ 上的截面间由
$$
\mathcal{F}(U) \times \mathcal{G}(U) \longrightarrow \mathcal{H}(U),\quad
(s, t) \longmapsto \psi(s)(t)
$$
给出。我们略去证明这些构造互逆的细节。
\end{proof}

\begin{lemma}
\label{sites-lemma-hom-sheaf-hU}
设 $\mathcal{C}$ 为位点，$U \in \Ob(\mathcal{C})$。则对
$\Sh(\mathcal{C})$ 中的 $\mathcal{F}$，有
$\SheafHom(h_U^\#, \mathcal{F}) = j_*(\mathcal{F}|_{\mathcal{C}/U})$。
\end{lemma}

\begin{proof}
可以直接构造层的同构来证明。这里改用如下论证。设 $\mathcal{G}$ 是
$\mathcal{C}$ 上的层。则
\begin{align*}
\Mor(\mathcal{G}, j_*(\mathcal{F}|_{\mathcal{C}/U}))
& =
\Mor(\mathcal{G}|_{\mathcal{C}/U}, \mathcal{F}|_{\mathcal{C}/U}) \\
& =
\Mor(j_!(\mathcal{G}|_{\mathcal{C}/U}), \mathcal{F}) \\
& =
\Mor(\mathcal{G} \times h_U^\#, \mathcal{F}) \\
& =
\Mor(\mathcal{G}, \SheafHom(h_U^\#, \mathcal{F}))
\end{align*}
由米田引理即可得出结论。这里使用了引理
\ref{sites-lemma-internal-hom-sheaf} 和
\ref{sites-lemma-compute-j-shriek-restrict}。
\end{proof}

\noindent
设 $\mathcal{C}$ 为位点，$\{U_i \to U\}_{i \in I}$ 是
$\mathcal{C}$ 的覆盖。对每个 $i \in I$，设 $\mathcal{F}_i$ 是
$\mathcal{C}/U_i$ 上的集合层。对每一对 $i, j \in I$，设
$$
\varphi_{ij} :
\mathcal{F}_i|_{\mathcal{C}/U_i \times_U U_j}
\longrightarrow
\mathcal{F}_j|_{\mathcal{C}/U_i \times_U U_j}
$$
是集合层的同构。再假设，对指标的每个三元组 $i, j, k \in I$，下图交换：
$$
\xymatrix{
\mathcal{F}_i|_{\mathcal{C}/U_i \times_U U_j \times_U U_k}
\ar[rr]_{\varphi_{ik}}
\ar[rd]_{\varphi_{ij}} & &
\mathcal{F}_k|_{\mathcal{C}/U_i \times_U U_j \times_U U_k} \\
&
\mathcal{F}_j|_{\mathcal{C}/U_i \times_U U_j \times_U U_k}
\ar[ru]_{\varphi_{jk}}
}
$$
我们把这样一族数据 $(\mathcal{F}_i, \varphi_{ij})$ 称为
{\it 关于覆盖 $\{U_i \to U\}_{i \in I}$ 的集合层粘合数据}。

\begin{lemma}
\label{sites-lemma-glue-sheaves}
设 $\mathcal{C}$ 为位点，$\{U_i \to U\}_{i \in I}$ 是
$\mathcal{C}$ 的覆盖。给定关于覆盖 $\{U_i \to U\}_{i \in I}$ 的任意
集合层粘合数据 $(\mathcal{F}_i, \varphi_{ij})$，存在
$\mathcal{C}/U$ 上的集合层 $\mathcal{F}$ 以及同构
$$
\varphi_i : \mathcal{F}|_{\mathcal{C}/U_i} \to \mathcal{F}_i
$$
使得下列各图交换：
$$
\xymatrix{
\mathcal{F}|_{\mathcal{C}/U_i \times_U U_j}
\ar[d]_{\text{id}} \ar[r]_{\varphi_i} &
\mathcal{F}_i|_{\mathcal{C}/U_i \times_U U_j} \ar[d]^{\varphi_{ij}} \\
\mathcal{F}|_{\mathcal{C}/U_i \times_U U_j} \ar[r]^{\varphi_j} &
\mathcal{F}_j|_{\mathcal{C}/U_i \times_U U_j}
}
$$
\end{lemma}

\begin{proof}
说明如何构造 $\mathcal{C}/U$ 上的层 $\mathcal{F}$。设
$a : V \to U$ 是 $\mathcal{C}/U$ 的对象。则
$$
\mathcal{F}(V/U) = \{
(s_i)_{i \in I} \in \prod_{i \in I} \mathcal{F}_i(U_i \times_U V/U_i)
\mid
\varphi_{ij}(s_i|_{U_i \times_U U_j \times_U V})
=
s_j|_{U_i \times_U U_j \times_U V}
\}
$$
我们略去限制映射的构造，略去这是一个层的验证，也略去同构
$\varphi_i$ 的构造以及引理中各图交换性的证明。
\end{proof}

\noindent
设 $\mathcal{C}$ 为位点，$\{U_i \to U\}_{i \in I}$ 是
$\mathcal{C}$ 的覆盖，$\mathcal{F}$ 是 $\mathcal{C}/U$ 上的层。与
$\mathcal{F}$ 相伴的{\it 典范粘合数据}由限制
$\mathcal{F}|_{\mathcal{C}/U_i}$ 以及典范同构
$$
\left(\mathcal{F}|_{\mathcal{C}/U_i}\right)|_{\mathcal{C}/U_i \times_U U_j}
=
\left(\mathcal{F}|_{\mathcal{C}/U_j}\right)|_{\mathcal{C}/U_i \times_U U_j}
$$
给出；这些同构来自如下事实：复合函子
$\mathcal{C}/U_i \times_U U_j \to \mathcal{C}/U_i \to \mathcal{C}/U$
与
$\mathcal{C}/U_i \times_U U_j \to \mathcal{C}/U_j \to \mathcal{C}/U$
相等。

\begin{lemma}
\label{sites-lemma-mapping-property-glue}
设 $\mathcal{C}$ 为位点，$\{U_i \to U\}_{i \in I}$ 是
$\mathcal{C}$ 的覆盖。把 $\mathcal{C}/U$ 上的 $\mathcal{F}$ 送到其典范
粘合数据的函子，给出范畴 $\Sh(\mathcal{C}/U)$ 与粘合数据范畴之间的等价。
\end{lemma}

\begin{proof}
在引理 \ref{sites-lemma-glue-maps} 中，我们看到该函子是全忠实的；在引理
\ref{sites-lemma-glue-sheaves} 中，我们证明它是本质满的（方法是显式构造一个
拟逆函子）。
\end{proof}

\noindent
设 $\mathcal{C}$ 为位点。下面讨论在整个位点 $\mathcal{C}$ 上粘合层的一个
版本。对 $\mathcal{C}$ 的每个对象 $U$，设 $\mathcal{F}_U$ 是
$\mathcal{C}/U$ 上的层。回忆：对 $\mathcal{C}$ 中的每个态射
$f : V \rightarrow U$，都有相伴函子
$j_f : \mathcal{C}/V \rightarrow \mathcal{C}/U$，它由
$(a : W \to V) \mapsto (f \circ a : W \to U)$ 给出。对每个这样的 $f$，
设
$$
c_f : j_f^{-1} \mathcal{F}_U \to \mathcal{F}_V
$$
是层的同构。再假设，对 $\mathcal{C}$ 中任意两个箭头
$f : V \to U$ 与 $g : W \to V$，复合
$c_g \circ j_{g}^{-1} c_f$ 等于 $c_{f \circ g}$。我们把这样一族数据
$(\mathcal{F}_{U}, c_f)$ 称为{\it $\mathcal{C}$ 上集合层的绝对粘合数据}。
绝对粘合数据间的态射
$(\mathcal{F}_{U}, c_f) \to (\mathcal{G}_{U}, c'_f)$ 由一族层态射
$(\varphi_U)$ 给出，其中 $\varphi_U : \mathcal{F}_U \to \mathcal{G}_U$，
并且
$$
\xymatrix{
j_f^{-1} \mathcal{F}_U \ar[r]_{c_f} \ar[d]_{j_f^{-1} \varphi_U} &
\mathcal{F}_V \ar[d]^{\varphi_V} \\
j_f^{-1} \mathcal{G}_U \ar[r]^{c'_f} &
\mathcal{G}_V
}
$$
对 $\mathcal{C}$ 中每个态射 $f : V \to U$ 都交换。

\medskip\noindent
与 $\mathcal{C}$ 上任意层 $\mathcal{F}$ 相伴的，是其{\it 典范绝对粘合数据}
$(\mathcal{F}|_{\mathcal{C}/U},c_f)$；其中，对 $f : V \to U$，典范同构
$c_f : j_f^{-1} \mathcal{F}|_{\mathcal{C}/U} \to
\mathcal{F}|_{\mathcal{C}/V}$ 来自引理 \ref{sites-lemma-relocalize} 中的关系
$j_V = j_U \circ j_f$。$\mathcal{C}$ 的层之间的任意态射
$\varphi : \mathcal{F} \to \mathcal{G}$ 都诱导典范绝对粘合数据间的态射
$(\varphi|_{\mathcal{C}/U})$。

\begin{lemma}
\label{sites-lemma-glue-sheaves-absolute}
设 $\mathcal{C}$ 为位点。把 $\mathcal{C}$ 上的 $\mathcal{F}$ 送到其典范
绝对粘合数据的函子，给出范畴 $\Sh(\mathcal{C})$ 与绝对粘合数据范畴之间的
等价。
\end{lemma}

\begin{proof}
给定绝对粘合数据 $(\mathcal{F}_U, c_f)$，令
$\mathcal{F}(U) = \mathcal{F}_U(U)$，从而构造 $\mathcal{C}$ 上的层
$\mathcal{F}$；沿 $f : V \to U$ 的限制由交换图给出：
% P01-ERRATA: P01-E0096 -- source sentence omits "is" before "given by the commutative diagram".
$$
\xymatrix{
\mathcal{F}_U(U) \ar[r] \ar@{=}[d] &
\mathcal{F}_U(V) \ar[r]^{c_f} &
\mathcal{F}_V(V) \ar@{=}[d] \\
\mathcal{F}(U) \ar[rr] &  &
\mathcal{F}(V)
}
$$
相容条件 $c_g \circ j_{g}^{-1} c_f = c_{f \circ g}$ 保证
$\mathcal{F}$ 是预层，也保证映射
$c_f : \mathcal{F}_U(V) \to \mathcal{F}(V)$ 对每个 $U$ 都定义同构
$\mathcal{F}_U \to \mathcal{F}|_{\mathcal{C}/U}$。由于每个
$\mathcal{F}_U$ 都是层，这也蕴含 $\mathcal{F}$ 是层。刚才构造的函子
$(\mathcal{F}_U, c_f) \mapsto \mathcal{F}$ 是把 $\mathcal{C}$ 上的层送到
其典范粘合数据的函子的拟逆。其余细节略去。
\end{proof}

\begin{remark}
\label{sites-remark-variant-glue-sheaves-absolute}
引理 \ref{sites-lemma-glue-sheaves-absolute} 有一个会在代数几何中出现的变体。
具体地，假设 $\mathcal{C}$ 是具有所有纤维积的位点，并且对每个
$U \in \Ob(\mathcal{C})$，都给定满足下列性质的满子范畴
$U_\tau \subset \mathcal{C}/U$：
\begin{enumerate}
\item $U/U$ 属于 $U_\tau$；
\item 对 $U_\tau$ 中的 $V/U$ 以及 $\mathcal{C}$ 的覆盖
$\{V_j \to V\}$，有 $V_j/U$ 属于 $U_\tau$；且
\item 对 $\mathcal{C}$ 的态射 $U' \to U$ 以及 $U_\tau$ 中的 $V/U$，
基变换 $V' = V \times_U U'$ 属于 $U'_\tau$。
\end{enumerate}
在此情形，对 $\mathcal{C}$ 中所有 $U$，$U_\tau$ 都是位点；并且对
$\mathcal{C}$ 的每个态射 $f : U' \to U$，基变换函子
$U_\tau \to U'_\tau$ 定义位点态射 $f_\tau : U_\tau \to U'_\tau$。
此时所得粘合陈述如下：$\mathcal{C}$ 上的层 $\mathcal{F}$ 由下列数据给出：
\begin{enumerate}
\item 对每个 $U \in \Ob(\mathcal{C})$，给定 $U_\tau$ 上的层
$\mathcal{F}_U$；
\item 对 $\mathcal{C}$ 中每个 $f : U' \to U$，给定映射
$c_f : f_\tau^{-1}\mathcal{F}_U \to \mathcal{F}_{U'}$。
\end{enumerate}
这些数据须满足下列条件：
\begin{enumerate}
\item[(a)] 给定 $\mathcal{C}$ 中的 $f : U' \to U$ 和
$g : U'' \to U'$，复合 $c_g \circ g_\tau^{-1}c_f$ 等于
$c_{f \circ g}$；且
\item[(b)] 若 $f : U' \to U$ 属于 $U_\tau$，则 $c_f$ 是同构。
\end{enumerate}
如果以后需要这一结果，我们会在此精确陈述并证明。（注意，这一结果与上述
陈述略有不同，因为这里并不要求所有映射 $c_f$ 都是同构！）
\end{remark}



\section{进一步的局部化}
\label{sites-section-more-localize}

\noindent
本节证明几个关于局部化的引理，其中会对位点或作局部化的对象施加一些附加
假设。
% P01-ERRATA: P01-E0097 -- source has "hypotheses on the site on or the object" instead of "on the site or on the object".

\begin{lemma}
\label{sites-lemma-describe-j-shriek-good-site}
设 $\mathcal{C}$ 为位点，$U \in \Ob(\mathcal{C})$。若
$\mathcal{C}$ 上的拓扑是次典范的，见定义
\ref{sites-definition-weaker-than-canonical}，并且 $\mathcal{G}$ 是
$\mathcal{C}/U$ 上的层，则
$$
j_{U!}(\mathcal{G})(V)
=
\coprod\nolimits_{\varphi \in \Mor_\mathcal{C}(V, U)}
\mathcal{G}(V \xrightarrow{\varphi} U),
$$
换言之，引理 \ref{sites-lemma-describe-j-shriek} 中无需层化。
\end{lemma}

\begin{proof}
设 $\mathcal{V} = \{V_i \to V\}_{i \in I}$ 是位点 $\mathcal{C}$ 中
$V$ 的覆盖。我们直接验证引理 \ref{sites-lemma-describe-j-shriek} 中预层
$\mathcal{H}$ 的层条件。设
$(s_i, \varphi_i)_{i \in I} \in \prod_i \mathcal{H}(V_i)$。这意味着
$\varphi_i : V_i \to U$ 是 $\mathcal{C}$ 中的态射，并且
$s_i \in \mathcal{G}(V_i \xrightarrow{\varphi_i} U)$。
% P01-ERRATA: P01-E0098 -- source joins these sentences with a comma before capitalized "This".
序对 $(s_i, \varphi_i)$ 到 $V_i \times_V V_j$ 的限制是
$(s_i|_{V_i \times_V V_j/U}, \text{pr}_1 \circ \varphi_i)$；类似地，序对
$(s_j, \varphi_j)$ 到 $V_i \times_V V_j$ 的限制是
$(s_j|_{V_i \times_V V_j/U}, \text{pr}_2 \circ \varphi_j)$。因此，若族
$(s_i, \varphi_i)$ 属于 $\check{H}^0(\mathcal{V}, \mathcal{H})$，则可见
$\text{pr}_1 \circ \varphi_i = \text{pr}_2 \circ \varphi_j$。
$\mathcal{C}$ 上的拓扑弱于典范拓扑这一条件于是蕴含：存在唯一态射
$\varphi : V \to U$，使 $\varphi_i$ 是 $V_i \to V$ 与 $\varphi$ 的复合。
此时，$\mathcal{G}$ 的层条件保证各截面 $s_i$ 粘合为唯一截面
$s \in \mathcal{G}(V \xrightarrow{\varphi} U)$。故如所需，
$(s, \varphi) \in \mathcal{H}(V)$。
\end{proof}

\begin{lemma}
\label{sites-lemma-localize-given-products}
设 $\mathcal{C}$ 为位点，$U \in \Ob(\mathcal{C})$。假设
$\mathcal{C}$ 有任意两个对象的积。则
\begin{enumerate}
\item 函子 $j_U$ 有连续的右伴随，即函子 $v(X) = X \times U / U$；
\item 函子 $v$ 定义位点态射 $\mathcal{C}/U \to \mathcal{C}$，其相伴的
拓扑斯态射等于 $j_U : \Sh(\mathcal{C}/U) \to \Sh(\mathcal{C})$；且
\item 有 $j_{U*}\mathcal{F}(X) = \mathcal{F}(X \times U/U)$。
\end{enumerate}
\end{lemma}

\begin{proof}
函子 $v$ 是 $j_U$ 的右伴随，意即给定 $Y/U$ 和 $X$ 时有
$$
\Mor_\mathcal{C}(Y, X)
=
\Mor_{\mathcal{C}/U}(Y/U, X \times U/U)
$$
这是显然的。为验证 $v$ 连续，设 $\{X_i \to X\}$ 是
$\mathcal{C}$ 的覆盖。由位点的第三条公理（定义 \ref{sites-definition-site}），
可见
$$
\{X_i \times_X (X \times U) \to X \times_X (X \times U)\}
=
\{X_i \times U \to X \times U\}
$$
也是 $\mathcal{C}$ 的覆盖。因此 $v$ 连续。引理的其余陈述由引理
\ref{sites-lemma-have-functor-other-way} 和
\ref{sites-lemma-have-functor-other-way-morphism} 得出。
\end{proof}

\begin{lemma}
\label{sites-lemma-relocalize-given-fibre-products}
设 $\mathcal{C}$ 为位点，$U \to V$ 是 $\mathcal{C}$ 的态射。假设
$\mathcal{C}$ 有纤维积。设 $j$ 如引理 \ref{sites-lemma-relocalize}。则
\begin{enumerate}
\item 函子 $j : \mathcal{C}/U \to \mathcal{C}/V$ 有连续的右伴随，即函子
$v : (X/V) \mapsto (X \times_V U/U)$；
\item 函子 $v$ 定义位点态射 $\mathcal{C}/U \to \mathcal{C}/V$，其相伴的
拓扑斯态射等于 $j$；且
\item 有 $j_*\mathcal{F}(X/V) = \mathcal{F}(X \times_V U/U)$。
\end{enumerate}
\end{lemma}

\begin{proof}
由引理 \ref{sites-lemma-localize-given-products} 得出，因为根据引理
\ref{sites-lemma-relocalize}，$j$ 可以视为局部化函子。
\end{proof}

\noindent
开浸入的一条基本性质是：先限制直像，或者先作空集延拓再限制，都会还原原来
的层。对上面与 $j_U$ 相伴的函子，这并不总成立。当 $U$ 是“终对象的子对象”
时则成立。
% P01-ERRATA: P01-E0099 -- source compound subject "the restriction ... and the restriction ..." takes singular "produces".

\begin{lemma}
\label{sites-lemma-restrict-back}
设 $\mathcal{C}$ 为位点，$U \in \Ob(\mathcal{C})$。假设
$\mathcal{C}$ 中每个 $X$ 至多有一个到 $U$ 的态射。设 $\mathcal{F}$ 是
$\mathcal{C}/U$ 上的层。典范映射
$\mathcal{F} \to j_U^{-1}j_{U!}\mathcal{F}$ 与
$j_U^{-1}j_{U*}\mathcal{F} \to \mathcal{F}$ 都是同构。
\end{lemma}

\begin{proof}
这是引理 \ref{sites-lemma-back-and-forth} 的特殊情形，因为对 $U$ 的假设等价于
局部化函子 $\mathcal{C}/U \to \mathcal{C}$ 的全忠实性。
% P01-ERRATA: P01-E0100 -- source uses the malformed noun phrase "the fully faithfulness".
\end{proof}

\begin{lemma}
\label{sites-lemma-localize-cartesian-square}
设 $\mathcal{C}$ 为位点。设
$$
\xymatrix{
U' \ar[d] \ar[r] & U \ar[d] \\
V' \ar[r] & V
}
$$
是 $\mathcal{C}$ 的交换图。引理 \ref{sites-lemma-relocalize} 的态射给出连续且
余连续函子的交换图以及拓扑斯的交换图
$$
\vcenter{
\xymatrix{
\mathcal{C}/U' \ar[d]_{j_{U'/V'}} \ar[r]_{j_{U'/U}} &
\mathcal{C}/U \ar[d]^{j_{U/V}} \\
\mathcal{C}/V' \ar[r]^{j_{V'/V}} & \mathcal{C}/V
}
}
\quad\text{且}\quad
\vcenter{
\xymatrix{
\Sh(\mathcal{C}/U') \ar[d]_{j_{U'/V'}} \ar[r]_{j_{U'/U}} &
\Sh(\mathcal{C}/U) \ar[d]^{j_{U/V}} \\
\Sh(\mathcal{C}/V') \ar[r]^{j_{V'/V}} &
\Sh(\mathcal{C}/V)
}
}
$$
此外，若 $\mathcal{C}$ 的初始图是笛卡尔的，则
$j_{V'/V}^{-1} \circ j_{U/V, *} = j_{U'/V', *} \circ j_{U'/U}^{-1}$。
\end{lemma}

\begin{proof}
引理第一条陈述中左方形的交换性由定义立即得出。由引理
\ref{sites-lemma-composition-cocontinuous}，它蕴含拓扑斯图的交换性。假设该图
是笛卡尔的。由伴随函子的唯一性，要证明
$j_{V'/V}^{-1} \circ j_{U/V, *} = j_{U'/V', *} \circ j_{U'/U}^{-1}$，
等价于证明
$j_{U/V}^{-1} \circ j_{V'/V!} = j_{U'/U!} \circ j_{U'/V'}^{-1}$。
借助引理 \ref{sites-lemma-essential-image-j-shriek} 的等同，可以把拓扑斯图看作
$$
\xymatrix{
\Sh(\mathcal{C})/h_{U'}^\# \ar[d] \ar[r] &
\Sh(\mathcal{C})/h_U^\# \ar[d] \\
\Sh(\mathcal{C})/h_{V'}^\# \ar[r] &
\Sh(\mathcal{C})/h_V^\#
}
$$
而由引理 \ref{sites-lemma-relocalize-explicit}，我们知道如何解释函子
$j^{-1}$ 与 $j_!$。因此必须证明：给定 $\mathcal{F} \to h_{V'}^\#$，有
$$
\mathcal{F} \times_{h_{V'}^\#} h_{U'}^\# =
\mathcal{F} \times_{h_V^\#} h_U^\#
$$
作为带有到 $h_U^\#$ 的映射的层。这是真的，因为
$h_{U'} = h_{V'} \times_{h_V} h_U$；而且由于层化是正合的，也有
$$
h_{U'}^\# = h_{V'}^\# \times_{h_V^\#} h_U^\#
$$
\end{proof}












\section{局部化与态射}
\label{sites-section-localize-morphisms}

\noindent
为理解局部化与位点及拓扑斯的态射之间的关系，下列引理很重要。
% P01-ERRATA: P01-E0101 -- source omits "the" in "understand relation between".

\begin{lemma}
\label{sites-lemma-localize-morphism}
设 $f : \mathcal{C} \to \mathcal{D}$ 是对应于连续函子
$u : \mathcal{D} \to \mathcal{C}$ 的位点态射。设
$V \in \Ob(\mathcal{D})$，并令 $U = u(V)$。则函子
$u' : \mathcal{D}/V \to \mathcal{C}/U$，
$V'/V \mapsto u(V')/U$，确定一个位点态射
$f' : \mathcal{C}/U \to \mathcal{D}/V$。态射 $f'$ 嵌入拓扑斯的交换图
$$
\xymatrix{
\Sh(\mathcal{C}/U) \ar[r]_{j_U} \ar[d]_{f'} &
\Sh(\mathcal{C}) \ar[d]^f \\
\Sh(\mathcal{D}/V) \ar[r]^{j_V} &
\Sh(\mathcal{D}).
}
$$
利用引理 \ref{sites-lemma-essential-image-j-shriek} 的等同
$\Sh(\mathcal{C}/U) = \Sh(\mathcal{C})/h_U^\#$ 与
$\Sh(\mathcal{D}/V) = \Sh(\mathcal{D})/h_V^\#$，函子
$(f')^{-1}$ 可由规则
$$
(f')^{-1}(\mathcal{H} \xrightarrow{\varphi} h_V^\#)
=
(f^{-1}\mathcal{H} \xrightarrow{f^{-1}\varphi} h_U^\#).
$$
描述。最后，有 $f'_*j_U^{-1} = j_V^{-1}f_*$。
\end{lemma}

\begin{proof}
显然 $u'$ 连续，因而得到函子
$f'_* = (u')^s = (u')^p$（见第 \ref{sites-section-functoriality-PSh} 节与第
\ref{sites-section-continuous-functors} 节）以及伴随
$(f')^{-1} = (u')_s = ((u')_p\ )^\#$。断言
$f'_*j_U^{-1} = j_V^{-1}f_*$ 来自
$$
(j_V^{-1}f_*\mathcal{F})(V'/V)
= f_*\mathcal{F}(V') = \mathcal{F}(u(V'))
= (j_U^{-1}\mathcal{F})(u(V')/U)
= (f'_*j_U^{-1}\mathcal{F})(V'/V)
$$
这甚至对预层也成立。先验并不清楚的是：$(f')^{-1}$ 是正合的、该图交换，
并且对 $(f')^{-1}$ 的描述成立。

\medskip\noindent
设 $\mathcal{H}$ 是 $\mathcal{D}/V$ 上的层。计算
$j_{U!}(f')^{-1}\mathcal{H}$。有
\begin{eqnarray*}
j_{U!}(f')^{-1}\mathcal{H}
& =
((j_U)_p(u'_p\mathcal{H})^\#)^\# \\
& =
((j_U)_pu'_p\mathcal{H})^\# \\
& =
(u_p(j_V)_p\mathcal{H})^\# \\
& =
f^{-1}j_{V!}\mathcal{H}
\end{eqnarray*}
第一个等式由展开定义得出。第二个等式由引理
\ref{sites-lemma-technical-up} 得出。第三个等式来自
$u \circ j_V = j_U \circ u'$。第四个等式再次由引理
\ref{sites-lemma-technical-up} 得出。
% P01-ERRATA: P01-E0102 -- source presents all four explanations as sentence fragments "The ... equality by/because ...".
上述等式全都是函子同构，因此可以把它解释为下列范畴与函子图交换：
$$
\xymatrix{
\Sh(\mathcal{C}/U) \ar[r]_{j_{U!}} &
\Sh(\mathcal{C})/h_U^\# \ar[r] &
\Sh(\mathcal{C}) \\
\Sh(\mathcal{D}/V) \ar[r]^{j_{V!}} \ar[u]^{(f')^{-1}} &
\Sh(\mathcal{D})/h_V^\# \ar[r] \ar[u]^{f^{-1}} &
\Sh(\mathcal{D}) \ar[u]^{f^{-1}}
}
$$
中间箭头有意义，因为
$f^{-1}h_V^\# = (h_{u(V)})^\# = h_U^\#$；见引理
\ref{sites-lemma-pullback-representable-sheaf}。特别地，这证明了引理陈述中对
$(f')^{-1}$ 的描述。根据引理 \ref{sites-lemma-essential-image-j-shriek}，左侧
水平箭头是等价；又因为根据假设 $f^{-1}$ 是正合的，故得出
$(f')^{-1} = u'_s$ 是正合的。具体地，因为它是左伴随，所以已经右正合
（《范畴》，引理 \ref{categories-lemma-adjoint-exact}）。因此只需证明它把
终对象变为终对象，并且与纤维积交换（《范畴》，引理
\ref{categories-lemma-characterize-left-exact}）。对诱导函子
$f^{-1} : \Sh(\mathcal{D})/h_V^\# \to \Sh(\mathcal{C})/h_U^\#$，两点都显然。
这证明 $f'$ 是位点态射。

\medskip\noindent
还须验证 $(f')^{-1}j_V^{-1} = j_U^{-1}f^{-1}$。为此，使用上面的公式以及
引理 \ref{sites-lemma-compute-j-shriek-restrict} 中的描述。具体地，二者结合给出：
对 $\mathcal{D}$ 上的任意层 $\mathcal{G}$，有
$$
j_{U!}(f')^{-1}j_V^{-1}\mathcal{G}
=
f^{-1}j_{V!}j_V^{-1}\mathcal{G}
=
f^{-1}(\mathcal{G} \times h_V^\#)
=
f^{-1}\mathcal{G} \times h_U^\#
=
j_{U!}j_U^{-1}f^{-1}\mathcal{G}.
$$
由于函子 $j_{U!}$ 诱导等价
$\Sh(\mathcal{C}/U) \to \Sh(\mathcal{C})/h_U^\#$，结论成立。
\end{proof}

\noindent
下列引理是上面更一般的引理 \ref{sites-lemma-localize-morphism} 的特殊情形。

\begin{lemma}
\label{sites-lemma-localize-morphism-strong}
设 $\mathcal{C}$、$\mathcal{D}$ 为位点，
$u : \mathcal{D} \to \mathcal{C}$ 为函子。设
$V \in \Ob(\mathcal{D})$，令 $U = u(V)$。假设
\begin{enumerate}
\item $\mathcal{C}$ 和 $\mathcal{D}$ 有所有有限极限；
\item $u$ 连续；且
\item $u$ 与有限极限交换。
\end{enumerate}
存在位点态射的交换图
$$
\xymatrix{
\mathcal{C}/U \ar[r]_{j_U} \ar[d]_{f'} & \mathcal{C} \ar[d]^f \\
\mathcal{D}/V \ar[r]^{j_V} & \mathcal{D}
}
$$
其中右竖直箭头对应于 $u$，左竖直箭头对应于函子
$u' : \mathcal{D}/V \to \mathcal{C}/U$，
$V'/V \mapsto u(V')/u(V)$；水平箭头则对应于引理
\ref{sites-lemma-localize-given-products} 中的函子
$\mathcal{C} \to \mathcal{C}/U$，$X \mapsto X \times U$，以及
$\mathcal{D} \to \mathcal{D}/V$，$Y \mapsto Y \times V$。
% P01-ERRATA: P01-E0103 -- source omits the slice structure notation /U and /V from these two functor values, unlike the cited lemma.
此外，相伴的拓扑斯态射图等于引理 \ref{sites-lemma-localize-morphism} 的图。
特别地，有 $f'_*j_U^{-1} = j_V^{-1}f_*$。
\end{lemma}

\begin{proof}
注意，$u$ 满足命题 \ref{sites-proposition-get-morphism} 的假设，因此由该命题诱导
位点态射 $f : \mathcal{C} \to \mathcal{D}$。显然，$u$ 诱导所述函子
$u'$。显然，这一函子也满足命题 \ref{sites-proposition-get-morphism} 的假设。
因此得到位点态射 $f' : \mathcal{C}/U \to \mathcal{D}/V$。根据我们对位点
态射复合的定义（见定义 \ref{sites-definition-composition-morphisms-sites}），并且
由于
$$
u(Y \times V) = u(Y) \times u(V) = u(Y) \times U
$$
表明与引理之图相反的范畴与函子图交换，故原图交换。
\end{proof}

\noindent
至此，我们可以局部化一个位点，知道如何再次局部化，也可以在下方位点的一个
对象处局部化位点态射。把这些结合起来，就得到下列类型的图。

\begin{lemma}
\label{sites-lemma-relocalize-morphism}
设 $f : \mathcal{C} \to \mathcal{D}$ 是对应于连续函子
$u : \mathcal{D} \to \mathcal{C}$ 的位点态射。设
$V \in \Ob(\mathcal{D})$、$U \in \Ob(\mathcal{C})$，且
$c : U \to u(V)$ 是 $\mathcal{C}$ 的态射。存在拓扑斯的交换图
$$
\xymatrix{
\Sh(\mathcal{C}/U) \ar[r]_{j_U} \ar[d]_{f_c} &
\Sh(\mathcal{C}) \ar[d]^f \\
\Sh(\mathcal{D}/V) \ar[r]^{j_V} &
\Sh(\mathcal{D}).
}
$$
有 $f_c = f' \circ j_{U/u(V)}$，其中
$f' : \Sh(\mathcal{C}/u(V)) \to \Sh(\mathcal{D}/V)$ 如引理
\ref{sites-lemma-localize-morphism}，而
$j_{U/u(V)} : \Sh(\mathcal{C}/U) \to \Sh(\mathcal{C}/u(V))$ 如引理
\ref{sites-lemma-relocalize}。利用引理
\ref{sites-lemma-essential-image-j-shriek} 的等同
$\Sh(\mathcal{C}/U) = \Sh(\mathcal{C})/h_U^\#$ 与
$\Sh(\mathcal{D}/V) = \Sh(\mathcal{D})/h_V^\#$，函子
$(f_c)^{-1}$ 可由规则
$$
(f_c)^{-1}(\mathcal{H} \xrightarrow{\varphi} h_V^\#)
=
(f^{-1}\mathcal{H} \times_{f^{-1}\varphi, h_{u(V)}^\#, c} h_U^\#
\rightarrow h_U^\#).
$$
描述。最后，给定任意态射 $b : V' \to V$、$a : U' \to U$ 以及
$c' : U' \to u(V')$，使得
$$
\xymatrix{
U' \ar[r]_-{c'} \ar[d]_a & u(V') \ar[d]^{u(b)} \\
U \ar[r]^-c & u(V)
}
$$
交换，则图
$$
\xymatrix{
\Sh(\mathcal{C}/U') \ar[r]_{j_{U'/U}} \ar[d]_{f_{c'}} &
\Sh(\mathcal{C}/U) \ar[d]^{f_c} \\
\Sh(\mathcal{D}/V') \ar[r]^{j_{V'/V}} &
\Sh(\mathcal{D}/V).
}
$$
交换。
\end{lemma}

\begin{proof}
这一引理可谓自证，它更多是把理论发展到此处已经知道的事项汇集起来。例如，
第一个方形的交换性来自图 (\ref{sites-equation-relocalize}) 的交换性以及引理
\ref{sites-lemma-localize-morphism} 中图的交换性。对 $f_c^{-1}$ 的描述来自合并
引理 \ref{sites-lemma-relocalize-explicit} 与引理
\ref{sites-lemma-localize-morphism}。最后一个方形的交换性随即由等式
$$
f^{-1}\mathcal{H} \times_{h_{u(V)}^\#, c} h_U^\# \times_{h_U^\#} h_{U'}^\#
=
f^{-1}(\mathcal{H} \times_{h_V^\#} h_{V'}^\#)
\times_{h_{u(V'), c'}^\#} h_{U'}^\#
$$
得出；这在利用
$f^{-1}h_V^\# = h_{u(V)}^\#$ 和
$f^{-1}h_{V'}^\# = h_{u(V')}^\#$ 后是形式的，见引理
\ref{sites-lemma-pullback-representable-sheaf}。
% P01-ERRATA: P01-E0104 -- source puts c' inside h_{u(V'),c'}^#; expected fibre-product indexing separates the base h_{u(V')}^# and map c'.
\end{proof}

\noindent
下列引理给出另一种局部化的函子性，适用于拓扑斯态射来自余连续函子的情形。
这里的图与引理 \ref{sites-lemma-localize-morphism} 中的图不同；特别地，引理
\ref{sites-lemma-localize-morphism} 中的等式
$f'_*j_U^{-1} = j_V^{-1}f_*$ 在下列引理的情形中一般不成立。

\begin{lemma}
\label{sites-lemma-localize-cocontinuous}
设 $\mathcal{C}$、$\mathcal{D}$ 为位点，
$u : \mathcal{C} \to \mathcal{D}$ 是余连续函子。设 $U$ 是
$\mathcal{C}$ 的对象，并令 $V = u(U)$。有交换图
$$
\xymatrix{
\mathcal{C}/U \ar[r]_{j_U} \ar[d]_{u'} & \mathcal{C} \ar[d]^u \\
\mathcal{D}/V \ar[r]^-{j_V} & \mathcal{D}
}
$$
其中左竖直箭头为
$u' : \mathcal{C}/U \to \mathcal{D}/V$，$U'/U \mapsto V'/V$。
% P01-ERRATA: P01-E0105 -- V' is undefined; the functor value used in the proof is u(U')/V.
则 $u'$ 也余连续，并且得到拓扑斯的交换图
$$
\xymatrix{
\Sh(\mathcal{C}/U) \ar[r]_{j_U} \ar[d]_{f'} &
\Sh(\mathcal{C}) \ar[d]^f \\
\Sh(\mathcal{D}/V) \ar[r]^-{j_V} &
\Sh(\mathcal{D})
}
$$
其中 $f$（相应地 $f'$）对应于 $u$（相应地 $u'$）。
\end{lemma}

\begin{proof}
第一个图的交换性是显然的。只要证明 $u'$ 余连续，它就蕴含第二个图交换。

\medskip\noindent
设 $U'/U$ 是 $\mathcal{C}/U$ 的对象。设
$\{V_j/V \to u(U')/V\}_{j \in J}$ 是 $\mathcal{D}/V$ 中
$u(U')/V$ 的覆盖。由于 $u$ 余连续，存在覆盖
$\{U_i' \to U'\}_{i \in I}$，使族
$\{u(U_i') \to u(U')\}$ 加细 $\mathcal{D}$ 中的覆盖
$\{V_j \to u(U')\}$。换言之，存在指标集映射
$\alpha : I \to J$ 以及在 $U'$ 上的态射
$\phi_i : u(U_i') \to V_{\alpha(i)}$。
% P01-ERRATA: P01-E0106 -- these are morphisms in D and must be over u(U'), not over U' (an object of C).
通过复合 $U'_i \to U' \to U$，把 $U_i'$ 看作 $U$ 上的对象。则
$\{U'_i/U \to U'/U\}$ 是 $\mathcal{C}/U$ 的覆盖，使
$\{u(U_i')/V \to u(U')/V\}$ 加细
$\{V_j/V \to u(U')/V\}$（使用同一个 $\alpha$ 和同一些映射
$\phi_i$）。因此 $u' : \mathcal{C}/U \to \mathcal{D}/V$ 余连续。
\end{proof}

\begin{lemma}
\label{sites-lemma-localize-cocontinuous-downstairs}
设 $\mathcal{C}$、$\mathcal{D}$ 为位点，
$u : \mathcal{C} \to \mathcal{D}$ 是余连续函子。设 $V$ 是
$\mathcal{D}$ 的对象。令 ${}^u_V\mathcal{I}$ 为第
\ref{sites-section-more-functoriality-PSh} 节引入的范畴。有交换图
$$
\vcenter{
\xymatrix{
\,_V^u\mathcal{I} \ar[r]_j \ar[d]_{u'} &
\mathcal{C} \ar[d]^u \\
\mathcal{D}/V \ar[r]^-{j_V} &
\mathcal{D}
}
}
\quad\text{其中}\quad
\begin{matrix}
j : (U, \psi) \mapsto U \\
u' : (U, \psi) \mapsto (\psi : u(U) \to V)
\end{matrix}
$$
规定 ${}^u_V\mathcal{I}$ 的一族态射
$\{(U_i, \psi_i) \to (U, \psi)\}$ 是覆盖，当且仅当
$\{U_i \to U\}$ 是 $\mathcal{C}$ 中的覆盖。则
\begin{enumerate}
\item ${}^u_V\mathcal{I}$ 是位点；
\item $j$ 连续且余连续；
\item $u'$ 余连续；
\item 得到拓扑斯的交换图
$$
\xymatrix{
\Sh({}^u_V\mathcal{I}) \ar[r]_j \ar[d]_{f'} &
\Sh(\mathcal{C}) \ar[d]^f \\
\Sh(\mathcal{D}/V) \ar[r]^-{j_V} &
\Sh(\mathcal{D})
}
$$
其中 $f$（相应地 $f'$）对应于 $u$（相应地 $u'$）；且
\item 有 $f'_*j^{-1} = j_V^{-1}f_*$。
\end{enumerate}
\end{lemma}

\begin{proof}
(1)、(2)、(3) 和 (4) 是定义以及函子 $j$ 与纤维积交换这一事实的直接推论。
细节略去。为证明 (5)，回忆 $f_*$ 由 ${}_su = {}_pu$ 给出。因此，
$j_V^{-1}f_*\mathcal{F}$ 在 $V'/V$ 上的值，是
${}_pu\mathcal{F}$ 在 $V'$ 上的值，也就是 $\mathcal{F}$ 在范畴
${}^u_{V'}\mathcal{I}$ 上各值的极限。显然有范畴等价
$$
{}^u_{V'}\mathcal{I} \to {}^{u'}_{V'/V}\mathcal{I}
$$
另一方面，$f'_*j^{-1}\mathcal{F}$ 在 $V'/V$ 上的值由
$j^{-1}\mathcal{F}$ 在范畴 ${}^{u'}_{V'/V}\mathcal{I}$ 上各值的极限给出；
而且 $j^{-1}\mathcal{F}$ 在 ${}^u_V\mathcal{I}$ 的对象上的值正是
$\mathcal{F}$ 的值（由引理 \ref{sites-lemma-when-shriek}，因为 $j$ 连续且
余连续），故 (5) 成立。
\end{proof}

\noindent
下列两个结果的性质略有不同。

\begin{lemma}
\label{sites-lemma-special-square-cocontinuous}
假设给定位点 $\mathcal{C}', \mathcal{C}, \mathcal{D}', \mathcal{D}$ 以及函子
$$
\xymatrix{
\mathcal{C}' \ar[r]_{v'} \ar[d]_{u'} &
\mathcal{C} \ar[d]^u \\
\mathcal{D}' \ar[r]^v &
\mathcal{D}
}
$$
假设
\begin{enumerate}
\item $u$、$u'$、$v$ 与 $v'$ 余连续，并由引理
\ref{sites-lemma-cocontinuous-morphism-topoi} 分别产生拓扑斯态射
$f$、$f'$、$g$ 与 $g'$；
\item $v \circ u' = u \circ v'$；
\item $v$ 与 $v'$ 既连续又余连续；且
\item 对 $\mathcal{D}'$ 的任意对象 $V'$，由 $v$ 给出的函子
${}^{u'}_{V'}\mathcal{I} \to {}^{\ \ \ u}_{v(V')}\mathcal{I}$ 是余终的。
\end{enumerate}
则 $f'_* \circ (g')^{-1} = g^{-1} \circ f_*$，并且
$g'_! \circ (f')^{-1} = f^{-1} \circ g_!$。
\end{lemma}

\begin{proof}
范畴 ${}^{u'}_{V'}\mathcal{I}$ 与
${}^{\ \ \ u}_{v(V')}\mathcal{I}$ 定义于第
\ref{sites-section-more-functoriality-PSh} 节。条件 (4) 中的函子把
${}^{u'}_{V'}\mathcal{I}$ 的对象 $\psi : u'(U') \to V'$ 送到
${}^{\ \ \ u}_{v(V')}\mathcal{I}$ 的对象
$v(\psi) : uv'(U') = vu'(U') \to v(V')$。回忆，$g^{-1}$ 由
$v^p$ 给出（引理 \ref{sites-lemma-when-shriek}），而 $f_*$ 由
${}_su = {}_pu$ 给出。因此，$g^{-1}f_*\mathcal{F}$ 在 $V'$ 上的值，是
${}_pu\mathcal{F}$ 在 $v(V')$ 上的值，即极限
$$
\lim_{u(U) \to v(V') \in \Ob({}^{\ \ \ u}_{v(V')}\mathcal{I}^{opp})}
\mathcal{F}(U)
$$
同理，$f'_*(g')^{-1}\mathcal{F}$ 在 $V'$ 上的值由极限
$$
\lim_{u'(U') \to V' \in \Ob({}^{u'}_{V'}\mathcal{I}^{opp})} \mathcal{F}(v'(U'))
$$
给出。因此，假设 (4) 及《范畴》，引理 \ref{categories-lemma-initial}
表明二者相同，从而证明引理的第一个等式。第二个等式由第一个等式及伴随函子
的唯一性得出。
\end{proof}

\begin{lemma}
\label{sites-lemma-special-square-continuous}
假设给定位点 $\mathcal{C}', \mathcal{C}, \mathcal{D}', \mathcal{D}$ 以及函子
$$
\xymatrix{
\mathcal{C}' \ar[r]_{v'} &
\mathcal{C} \\
\mathcal{D}' \ar[r]^v \ar[u]^{u'} &
\mathcal{D} \ar[u]_u
}
$$
沿用第 \ref{sites-section-morphism-sites} 节与第
\ref{sites-section-cocontinuous-morphism-topoi} 节的记号，假设
\begin{enumerate}
\item $u$ 与 $u'$ 连续，并产生位点态射 $f$ 与 $f'$；
\item $v$ 与 $v'$ 余连续，并产生拓扑斯态射 $g$ 与 $g'$；
\item $u \circ v = v' \circ u'$；且
\item $v$ 与 $v'$ 既连续又余连续。
\end{enumerate}
则\footnote{在这一一般性下，我们不知道 $f \circ g'$ 与
$g \circ f'$ 作为拓扑斯态射是否相等（从前者到后者有典范 $2$-箭头，但它未必
是同构）。}
$f'_* \circ (g')^{-1} = g^{-1} \circ f_*$，并且
$g'_! \circ (f')^{-1} = f^{-1} \circ g_!$。
\end{lemma}

\begin{proof}
具体地，有
$$
f'_*(g')^{-1}\mathcal{F} = (u')^p((v')^p\mathcal{F})^\# =
(u')^p(v')^p\mathcal{F}
$$
第一个等式由定义得出，第二个等式由引理 \ref{sites-lemma-when-shriek} 得出。又有
$$
g^{-1}f_*\mathcal{F} = (v^pu^p\mathcal{F})^\# =
((u')^p(v')^p\mathcal{F})^\# = (u')^p(v')^p\mathcal{F}
$$
第一个等式由定义得出，第二个等式来自
$u \circ v = v' \circ u'$，第三个等式则因为已经看到
$(u')^p(v')^p\mathcal{F}$ 是层。这证明
$f'_* \circ (g')^{-1} = g^{-1} \circ f_*$；而等式
$g'_! \circ (f')^{-1} = f^{-1} \circ g_!$ 由左伴随的唯一性得出。
% P01-ERRATA: P01-E0107 -- source proof again uses fragmentary "The first equality by ..." clauses for both displays.
\end{proof}
















\section{拓扑斯的态射}
\label{sites-section-morphisms-topoi}

\noindent
本节证明，每个拓扑斯态射都等价于一个来自位点态射的拓扑斯态射。请与
\cite[Expos\'e III, Proposition 4.9.4]{SGA4} 比较。

\begin{lemma}
\label{sites-lemma-equivalence}
设 $\mathcal{C}$、$\mathcal{D}$ 为位点，
$u : \mathcal{C} \to \mathcal{D}$ 为函子。假设
\begin{enumerate}
\item $u$ 余连续；
\item $u$ 连续；
\item 给定 $\mathcal{C}$ 中的 $a, b : U' \to U$，若
$u(a) = u(b)$，则 $\mathcal{C}$ 中存在覆盖
$\{f_i : U'_i \to U'\}$，使 $a \circ f_i = b \circ f_i$；
\item 给定 $U', U \in \Ob(\mathcal{C})$ 以及 $\mathcal{D}$ 中的态射
$c : u(U') \to u(U)$，存在 $\mathcal{C}$ 中的覆盖
$\{f_i : U_i' \to U'\}$ 以及态射 $c_i : U_i' \to U$，使
$u(c_i) = c \circ u(f_i)$；且
\item 给定 $V \in \Ob(\mathcal{D})$，$\mathcal{D}$ 中存在一个形如
$\{u(U_i) \to V\}_{i \in I}$ 的 $V$ 的覆盖。
\end{enumerate}
则由引理 \ref{sites-lemma-cocontinuous-morphism-topoi} 与余连续函子 $u$ 相伴的
拓扑斯态射
$$
g : \Sh(\mathcal{C}) \longrightarrow \Sh(\mathcal{D})
$$
是等价。
\end{lemma}

\begin{proof}
假设 $u$ 满足性质 (1) -- (5)。将证明伴随映射
$$
\mathcal{G} \longrightarrow g_*g^{-1}\mathcal{G}
\quad\text{且}\quad
g^{-1}g_*\mathcal{F} \longrightarrow \mathcal{F}
$$
是同构。

\medskip\noindent
注意，引理 \ref{sites-lemma-when-shriek} 适用，并且对 $\mathcal{D}$ 上任意层
$\mathcal{G}$，有 $g^{-1}\mathcal{G}(U) = \mathcal{G}(u(U))$。接着，设
$\mathcal{F}$ 是 $\mathcal{C}$ 上的层，$V$ 是 $\mathcal{D}$ 的对象。由定义，
有 $g_*\mathcal{F}(V) = \lim_{u(U) \to V} \mathcal{F}(U)$。因此
$$
g^{-1}g_*\mathcal{F}(U) = \lim_{U', u(U') \to u(U)} \mathcal{F}(U')
$$
其中态射 $\psi : u(U') \to u(U)$ 不一定具有形式 $u(\alpha)$。这样一些序对
$(U', \psi)$ 所成的范畴有终对象，即 $(U, \text{id})$，它产生从该极限到
$\mathcal{F}(U)$ 的映射。
% P01-ERRATA: P01-E0109 -- (U,id) need not be final: a morphism from (U',psi) to it would require psi=u(alpha), which the preceding sentence explicitly says need not hold.
设 $(s_{(U', \psi)})$ 是该极限的一个元素。要证明
$s_{(U', \psi)}$ 由值
$s_{(U, \text{id})} \in \mathcal{F}(U)$ 唯一确定。由性质 (4)，给定任意
$(U', \psi)$，都存在覆盖 $\{U'_i \to U'\}$，使复合
$u(U'_i) \to u(U') \to u(U)$ 具有形式 $u(c_i)$，其中
$c_i : U'_i \to U$ 是 $\mathcal{C}$ 中的态射。因此
$$
s_{(U', \psi)}|_{U'_i} = c_i^*(s_{(U, \text{id})}).
$$
因为 $\mathcal{F}$ 是层，确实可得 $s_{(U', \psi)}$ 由
$s_{(U, \text{id})}$ 确定。这证明唯一性。

为证明存在性，任取
$s \in \mathcal{F}(U)$、$\psi : u(U') \to u(U)$、
$\{f_i : U_i' \to U'\}$ 以及如上的 $c_i : U_i' \to U$，使
$\psi \circ u(f_i) = u(c_i)$。我们断言存在（唯一的）元素
$s_{(U', \psi)} \in \mathcal{F}(U')$，使
$$
s_{(U', \psi)}|_{U'_i} = c_i^*(s).
$$
事实上，先验并不清楚元素
$c_i^*(s)|_{U_i' \times_{U'} U_j'}$ 与
$c_j^*(s)|_{U_i' \times_{U'} U_j'}$ 是否相等，因为图
$$
\xymatrix{
U_i' \times_{U'} U_j' \ar[r]_-{\text{pr}_2} \ar[d]_{\text{pr}_1} &
U_j' \ar[d]^{c_j} \\
U_i' \ar[r]^{c_i} & U}
$$
不一定交换。但引理的条件 (3) 保证存在覆盖
$\{f_{ijk} : U'_{ijk} \to U_i' \times_{U'} U_j'\}_{k \in K_{ij}}$，使
$c_i \circ \text{pr}_1 \circ f_{ijk} = c_j \circ \text{pr}_2 \circ f_{ijk}$。
因此
$$
f_{ijk}^* \left(c_i^*s|_{U_i' \times_{U'} U_j'}\right)
=
f_{ijk}^* \left(c_j^*s|_{U_i' \times_{U'} U_j'}\right)
$$
所以，由 $\mathcal{F}$ 的层条件，
$c_i^*(s)|_{U_i' \times_{U'} U_j'} = c_j^*(s)|_{U_i' \times_{U'} U_j'}$；
再由 $\mathcal{F}$ 的层条件，$s_{(U', \psi)}$ 存在。唯一性保证如此得到的族
$(s_{(U', \psi)})$ 是该极限的元素，并且 $s_{(U, \psi)} = s$。这证明
$g^{-1}g_*\mathcal{F} \to \mathcal{F}$ 是同构。
% P01-ERRATA: P01-E0108 -- source concludes s_(U,psi)=s although the distinguished component throughout is s_(U,id).

\medskip\noindent
设 $\mathcal{G}$ 是 $\mathcal{D}$ 上的层，$V$ 是 $\mathcal{D}$ 的对象。于是
$$
g_*g^{-1}\mathcal{G}(V) =
\lim_{U, \psi : u(U) \to V} \mathcal{G}(u(U))
$$
由前一段可见，层 $g_*g^{-1}\mathcal{G}$ 在形如 $V = u(U)$ 的对象
$V$ 上的值等于 $\mathcal{G}(u(U))$。（形式地，这是因为
$g^{-1}g_*g^{-1} \cong g^{-1}$ 以及证明开头给出的 $g^{-1}$ 的描述；非形式地，
只需比较这里和上面的极限。）因此，伴随映射
$\mathcal{G} \to g_*g^{-1}\mathcal{G}$ 在每个形如 $u(U)$ 的对象上的截面间
都是双射。由公理 (5)，存在由形如 $u(U)$ 的对象构成的 $V$ 的覆盖，因而容易
看出该伴随映射是同构。
\end{proof}

\noindent
给引理 \ref{sites-lemma-equivalence} 中那样的余连续函子一个名称会很方便。

\begin{definition}
\label{sites-definition-special-cocontinuous-functor}
设 $\mathcal{C}$、$\mathcal{D}$ 为位点。{\it 从 $\mathcal{C}$ 到
$\mathcal{D}$ 的特殊余连续函子 $u$}，是满足引理
\ref{sites-lemma-equivalence} 的假设与结论的余连续函子
$u : \mathcal{C} \to \mathcal{D}$。
\end{definition}

\begin{lemma}
\label{sites-lemma-localize-special-cocontinuous}
设 $\mathcal{C}$、$\mathcal{D}$ 为位点，
$u : \mathcal{C} \to \mathcal{D}$ 是特殊余连续函子。对
$\mathcal{C}$ 的每个对象 $U$，都有引理
\ref{sites-lemma-localize-cocontinuous} 中那样的交换图
$$
\xymatrix{
\mathcal{C}/U \ar[r]_{j_U} \ar[d] & \mathcal{C} \ar[d]^u \\
\mathcal{D}/u(U) \ar[r]^-{j_{u(U)}} & \mathcal{D}
}
$$
左竖直箭头是特殊余连续函子。因此，在拓扑斯的交换图
$$
\xymatrix{
\Sh(\mathcal{C}/U) \ar[r]_{j_U} \ar[d] &
\Sh(\mathcal{C}) \ar[d]^u \\
\Sh(\mathcal{D}/u(U)) \ar[r]^-{j_{u(U)}} &
\Sh(\mathcal{D})
}
$$
中，竖直箭头都是等价。
\end{lemma}

\begin{proof}
引理 \ref{sites-lemma-localize-cocontinuous} 已经给出这些图的存在性与交换性。
我们必须对诱导函子 $u : \mathcal{C}/U \to \mathcal{D}/u(U)$ 验证引理
\ref{sites-lemma-equivalence} 的假设 (1) -- (5)。这完全是机械的。

\medskip\noindent
性质 (1)。这是引理 \ref{sites-lemma-localize-cocontinuous}。

\medskip\noindent
性质 (2)。设 $\{U_i'/U \to U'/U\}_{i \in I}$ 是
$\mathcal{C}/U$ 中 $U'/U$ 的覆盖。由于 $u$ 连续，
$\{u(U_i')/u(U) \to u(U')/u(U)\}_{i \in I}$ 是
$\mathcal{D}/u(U)$ 中 $u(U')/u(U)$ 的覆盖。因此，(2) 对
$u : \mathcal{C}/U \to \mathcal{D}/u(U)$ 成立。

\medskip\noindent
性质 (3)。设 $a, b : U''/U \to U'/U$ 是 $\mathcal{C}/U$ 中满足
$u(a) = u(b)$ 于 $\mathcal{D}/u(U)$ 中的态射。由于 $u$ 满足 (3)，存在
$\mathcal{C}$ 中的覆盖 $\{f_i : U''_i \to U''\}$，使
$a \circ f_i = b \circ f_i$。这给出 $\mathcal{C}/U$ 中的覆盖
$\{f_i : U''_i/U \to U''/U\}$，使 $a \circ f_i = b \circ f_i$。
因此，(3) 对 $u : \mathcal{C}/U \to \mathcal{D}/u(U)$ 成立。

\medskip\noindent
性质 (4)。给定 $U''/U, U'/U \in \Ob(\mathcal{C}/U)$ 以及
$\mathcal{D}/u(U)$ 中的态射
$c : u(U'')/u(U) \to u(U')/u(U)$。由于 $u$ 满足性质 (4)，存在
$\mathcal{C}$ 中的覆盖 $\{f_i : U_i'' \to U''\}$ 以及态射
$c_i : U_i'' \to U'$，使 $u(c_i) = c \circ u(f_i)$。通过复合
$U_i'' \to U'' \to U$，把 $U_i''$ 看作 $U$ 上的对象。$c_i$ 未必是
$U$ 上的态射！但由于 $u(c_i)$ 是 $u(U)$ 上的态射，可以应用 $u$ 的性质
(3)，找到覆盖 $\{f_{ik} : U''_{ik} \to U''_i\}$，使
$c_{ik} = c_i \circ f_{ik} : U''_{ik} \to U'$ 是 $U$ 上的态射。因此，
$\{f_i \circ f_{ik} : U''_{ik}/U \to U''/U\}$ 是
$\mathcal{C}/U$ 中的覆盖，使 $u(c_{ik}) = c \circ u(f_{ik})$。
% P01-ERRATA: P01-E0110 -- source conclusion omits f_i: it should use u(f_i circ f_ik), matching the displayed covering and preceding equality.
因此，(4) 对 $u : \mathcal{C}/U \to \mathcal{D}/u(U)$ 成立。

\medskip\noindent
性质 (5)。设 $h : V \to u(U)$ 是 $\mathcal{D}/u(U)$ 的对象。由于
$u$ 满足性质 (5)，$\mathcal{D}$ 中存在覆盖
$\{c_i : u(U_i) \to V\}$。由性质 (4)，可以找到覆盖
$\{f_{ij} : U_{ij} \to U_i\}$ 和态射 $c_{ij} : U_{ij} \to U$，使
$u(c_{ij}) = h \circ c_i \circ u(f_{ij})$。因此，
$\{u(U_{ij})/u(U) \to V/u(U)\}$ 是 $\mathcal{D}/u(U)$ 中具有所需形式的
覆盖，从而得出 (5) 对 $u : \mathcal{C}/U \to \mathcal{D}/u(U)$ 成立。
\end{proof}

\begin{lemma}
\label{sites-lemma-special-equivalence}
设 $\mathcal{C}$ 为位点。设
$\mathcal{C}' \subset \Sh(\mathcal{C})$ 是满足下列条件的满子范畴（其对象构成
一个集合）：
\begin{enumerate}
\item 对所有 $U \in \Ob(\mathcal{C})$，有
$h_U^\# \in \Ob(\mathcal{C}')$；且
\item $\mathcal{C}'$ 在 $\Sh(\mathcal{C})$ 中的纤维积下保持不变。
\end{enumerate}
规定 $\mathcal{C}'$ 的覆盖是任意一族映射
$\{\mathcal{F}_i \to \mathcal{F}\}_{i \in I}$，使
$\coprod_{i \in I} \mathcal{F}_i \to \mathcal{F}$ 是层的满射。则
\begin{enumerate}
\item $\mathcal{C}'$ 是位点（选择一个覆盖集之后，见《集合》，引理
\ref{sets-lemma-coverings-site}）；
\item $\mathcal{C}'$ 上的可表预层是层（即 $\mathcal{C}'$ 上的拓扑是
次典范的，见定义 \ref{sites-definition-weaker-than-canonical}）；
\item 函子 $v : \mathcal{C} \to \mathcal{C}'$，$U \mapsto h_U^\#$ 是特殊
余连续函子，因而诱导等价
$g : \Sh(\mathcal{C}) \to \Sh(\mathcal{C}')$；
\item 对任意 $\mathcal{F} \in \Ob(\mathcal{C}')$，有
$g^{-1}h_\mathcal{F} = \mathcal{F}$；且
\item 对任意 $U \in \Ob(\mathcal{C})$，有
$g_*h_U^\# = h_{v(U)} = h_{h_U^\#}$。
\end{enumerate}
\end{lemma}

\begin{proof}
提醒：上面的某些陈述乍看可能有些令人困惑；这是因为 $\mathcal{C}'$ 的对象也
可以看作 $\mathcal{C}$ 上的层！
% P01-ERRATA: P01-E0111 -- source has the ungrammatical phrase "may look be a bit confusing".
我们略去证明引理所述的 $\mathcal{C}'$ 的覆盖满足定义
\ref{sites-definition-site} 的条件。

\medskip\noindent
假设 $\{\mathcal{F}_i \to \mathcal{F}\}$ 是层态射的满族。设
$\mathcal{G}$ 是另一个层。引理的 (2) 断言下图
$$
\xymatrix{
\Mor_{\Sh(\mathcal{C})}(
\coprod_{i \in I} \mathcal{F}_i, \mathcal{G})
\ar@<1ex>[r] \ar@<-1ex>[r]
&
\Mor_{\Sh(\mathcal{C})}(
\coprod_{(i_0, i_1) \in I \times I}
\mathcal{F}_{i_0} \times_\mathcal{F} \mathcal{F}_{i_1}, \mathcal{G})
}
$$
的等化子是 $\Mor_{\Sh(\mathcal{C})}(\mathcal{F}, \mathcal{G}).$
这是显然的（例如使用引理 \ref{sites-lemma-coequalizer-surjection}）。

\medskip\noindent
为证明 (3)，必须验证引理 \ref{sites-lemma-equivalence} 的条件 (1) -- (5)。
$v$ 余连续这一事实等价于引理 \ref{sites-lemma-mono-epi-sheaves} 中对层满射的
描述。函子 $v$ 连续，因为 $U \mapsto h_U^\#$ 与纤维积交换，并且把覆盖
变为覆盖（见引理 \ref{sites-lemma-sheafification-exact} 和引理
\ref{sites-lemma-covering-surjective-after-sheafification}）。引理
\ref{sites-lemma-equivalence} 的性质 (3)、(4) 是关于态射
$f : h_{U'}^\# \to h_U^\#$ 的陈述。这样的态射等同于
$h_U^\#(U')$ 的元素。因此，(3) 和 (4) 由层化的构造立即得出。引理
\ref{sites-lemma-equivalence} 的性质 (5) 是引理
\ref{sites-lemma-sheaf-coequalizer-representable}。以
$g : \Sh(\mathcal{C}) \to \Sh(\mathcal{C}')$ 表示由引理
\ref{sites-lemma-equivalence} 与 $v$ 相伴的拓扑斯等价。
% P01-ERRATA: P01-E0112 -- source "Denote g ... the equivalence" omits "by" or otherwise needs recasting.

\medskip\noindent
设 $\mathcal{F}$ 如引理的 (4)。对任意 $U \in \Ob(\mathcal{C})$，有
$$
g^{-1}h_\mathcal{F}(U) = h_\mathcal{F}(v(U))
= \Mor_{\Sh(\mathcal{C})}(h_U^\#, \mathcal{F})
= \mathcal{F}(U)
$$
第一个等式由引理 \ref{sites-lemma-when-shriek} 得出。因此 (4) 成立。
% P01-ERRATA: P01-E0113 -- source explanation is a sentence fragment: "The first equality by ...".

\medskip\noindent
设 $\mathcal{F} \in \Ob(\mathcal{C}')$，并设
$U \in \Ob(\mathcal{C})$。则
\begin{align*}
g_*h_U^\#(\mathcal{F})
& =
\Mor_{\Sh(\mathcal{C}')}(h_\mathcal{F}, g_*h_U^\#) \\
& =
\Mor_{\Sh(\mathcal{C})}(g^{-1}h_\mathcal{F}, h_U^\#) \\
& =
\Mor_{\Sh(\mathcal{C})}(\mathcal{F}, h_U^\#) \\
& =
\Mor_{\mathcal{C}'}(\mathcal{F}, h_U^\#)
\end{align*}
如所需（其中第三个等式已在上面证明）。
\end{proof}

\noindent
利用这一点，可以把任意拓扑斯整理为定义在具有所有有限极限的位点上。

\begin{lemma}
\label{sites-lemma-topos-good-site}
设 $\Sh(\mathcal{C})$ 为拓扑斯，$\{\mathcal{F}_i\}_{i \in I}$ 是
$\mathcal{C}$ 上的一个层集。存在拓扑斯等价
$g : \Sh(\mathcal{C}) \to \Sh(\mathcal{C}')$；它由位点间的特殊余连续函子
$u : \mathcal{C} \to \mathcal{C}'$ 诱导，并且
\begin{enumerate}
\item $\mathcal{C}'$ 具有次典范拓扑；
\item $\mathcal{C}'$ 的一族态射 $\{V_j \to V\}$ 是（组合等价于）
$\mathcal{C}'$ 的覆盖，当且仅当 $\coprod h_{V_j} \to h_V$ 是满射；
\item $\mathcal{C}'$ 有纤维积和终对象（即 $\mathcal{C}'$ 有所有有限极限）；
\item $\mathcal{C}'$ 上可表层的每个子层都是可表的；且
\item 每个 $g_*\mathcal{F}_i$ 都是可表层。
\end{enumerate}
\end{lemma}

\begin{proof}
考虑满子范畴 $\mathcal{C}_1 \subset \Sh(\mathcal{C})$，其对象包括所有
$U \in \Ob(\mathcal{C})$ 对应的 $h_U^\#$、给定的层
$\mathcal{F}_i$，以及终层 $*$（见例 \ref{sites-example-singleton-sheaf}）。
我们将归纳定义满子范畴
$$
\mathcal{C}_1 \subset \mathcal{C}_2 \subset \mathcal{C}_2 \subset \ldots
\subset \Sh(\mathcal{C})
$$
% P01-ERRATA: P01-E0114 -- source chain repeats C_2; the third term should be C_3.
具体地，给定 $\mathcal{C}_n$，令 $\mathcal{C}_{n + 1}$ 为由
$\mathcal{C}_n$ 的对象的一切纤维积和子层组成的满子范畴。（注意，
$\mathcal{C}_{n + 1}$ 有一个对象集。）令
$\mathcal{C}' = \bigcup_{n \geq 1} \mathcal{C}_n$。
$\mathcal{C}'$ 中的覆盖是 $\mathcal{C}'$ 的对象间的任意一族态射
$\{\mathcal{G}_j \to \mathcal{G}\}_{j \in J}$，使
$\coprod \mathcal{G}_j \to \mathcal{G}$ 作为 $\mathcal{C}$ 上的层映射是
满射。函子 $v : \mathcal{C} \to \mathcal{C'}$ 由
$U \mapsto h_U^\#$ 给出。应用引理 \ref{sites-lemma-special-equivalence}。
\end{proof}

\noindent
下面是本节的目标。

\begin{lemma}
\label{sites-lemma-morphism-topoi-comes-from-morphism-sites}
\begin{reference}
这一陈述与 \cite[Proposition 4.9.4. Expos\'e IV]{SGA4} 密切相关。为得到完整
结果，还应使用 \cite[Remarque 4.7.4, Expos\'e IV]{SGA4}。
\end{reference}
设 $\mathcal{C}$、$\mathcal{D}$ 为位点。设
$f : \Sh(\mathcal{C}) \to \Sh(\mathcal{D})$ 是拓扑斯态射。则存在位点
$\mathcal{C}'$ 以及函子图
$$
\xymatrix{
\mathcal{C} \ar[r]_v & \mathcal{C}' & \mathcal{D} \ar[l]^u
}
$$
使得
\begin{enumerate}
\item 函子 $v$ 是特殊余连续函子；
\item 函子 $u$ 与纤维积交换、连续，并定义位点态射
$\mathcal{C}' \to \mathcal{D}$；且
\item 拓扑斯态射 $f$ 与拓扑斯态射的复合
$$
\Sh(\mathcal{C}) \longrightarrow
\Sh(\mathcal{C}') \longrightarrow
\Sh(\mathcal{D})
$$
相符，其中第一个箭头由引理 \ref{sites-lemma-equivalence} 从 $v$ 得到，第二个
箭头由引理 \ref{sites-lemma-morphism-sites-topoi} 从 $u$ 得到。
\end{enumerate}
\end{lemma}

\begin{proof}
考虑满子范畴 $\mathcal{C}_1 \subset \Sh(\mathcal{C})$，其对象包括对所有
$U \in \Ob(\mathcal{C})$ 和所有 $V \in \Ob(\mathcal{D})$ 的
$h_U^\#$ 与 $f^{-1}h_V^\#$。令 $\mathcal{C}_{n + 1}$ 为由
$\mathcal{C}_n$ 的对象的一切纤维积组成的满子范畴。令
$\mathcal{C}' = \bigcup_{n \geq 1} \mathcal{C}_n$。
$\mathcal{C}'$ 中的覆盖是任意一族
$\{\mathcal{F}_i \to \mathcal{F}\}_{i \in I}$，使
$\coprod_{i \in I} \mathcal{F}_i \to \mathcal{F}$ 作为
$\mathcal{C}$ 上的层映射是满射。函子
$v : \mathcal{C} \to \mathcal{C'}$ 由 $U \mapsto h_U^\#$ 给出。函子
$u : \mathcal{D} \to \mathcal{C'}$ 由 $V \mapsto f^{-1}h_V^\#$ 给出。

\medskip\noindent
(1) 由引理 \ref{sites-lemma-special-equivalence} 得出。

\medskip\noindent
证明引理的 (2) 与 (3)。函子 $u$ 与纤维积交换，因为
$V \mapsto h_V^\#$ 和 $f^{-1}$ 都如此。此外，由于 $f^{-1}$ 正合并且与任意
余极限交换，可见它把覆盖变为层态射的满族。因此 $u$ 连续。为证明它定义
位点态射，还须证明 $u_s$ 正合。为此，将证明
$g^{-1} \circ u_s = f^{-1}$。事实上，届时因为 $g^{-1}$ 是等价并且
$f^{-1}$ 正合，即可得出结论。因为 $g^{-1}$ 伴随于 $g_*$，$u_s$ 伴随于
$u^s$，而 $f^{-1}$ 伴随于 $f_*$，所以证明
$u^s \circ g_* = f_*$ 也足够。设 $\mathcal{F}$ 是 $\mathcal{C}$ 上的层，
$V$ 是 $\mathcal{D}$ 的对象。则
\begin{align*}
(u^s g_{\ast}\mathcal{F})(V)
& =
(g_{\ast}\mathcal{F})(f^{-1}h_V^\#) \\
& =
\Mor_{Sh(\mathcal{C}')}(h_{f^{-1}h_V^\#},g_{\ast}\mathcal{F}) \\
& =
\Mor_{Sh(\mathcal{C})}(g^{-1}h_{f^{-1}h_V^\#},\mathcal{F}) \\
& =
\Mor_{Sh(\mathcal{C})}(f^{-1}h_V^\#,\mathcal{F}) \\
& =
\Mor_{Sh(\mathcal{D})}(h_V^\#,f_{\ast}\mathcal{F}) \\
& =
f_{\ast}\mathcal{F}(V)
\end{align*}
第一个等式是因为 $u^s = u^p$。第二个等式是米田引理。第三个等式由
$g^{-1}$ 与 $g_*$ 的伴随性得出。第四个等式由引理
\ref{sites-lemma-special-equivalence} (4) 得出。第五个等式由 $f^{-1}$ 与 $f_*$ 的
伴随性得出。第六个等式由米田引理得出。因此
$u^s g_*\mathcal{F} = f_*\mathcal{F}$，从而完成引理的证明。
% P01-ERRATA: P01-E0115 -- source gives several equality explanations as sentence fragments ("The first equality because", "The third equality by", etc.).
\end{proof}

\begin{remark}
\label{sites-remark-morphism-topoi-comes-from-morphism-sites}
记号和假设如引理 \ref{sites-lemma-morphism-topoi-comes-from-morphism-sites}。若位点
$\mathcal{D}$ 有终对象和纤维积，则函子
$u : \mathcal{D} \to \mathcal{C}'$ 满足命题
\ref{sites-proposition-get-morphism} 的所有假设。事实上，除引理中所述的性质外，
$u$ 还把 $\mathcal{D}$ 的终对象变为 $\mathcal{C}'$ 的终对象。这由 $u$ 的
构造显然可知。因此，若先把引理 \ref{sites-lemma-topos-good-site} 应用于
$\mathcal{D}$，再把引理
\ref{sites-lemma-morphism-topoi-comes-from-morphism-sites} 应用于所得拓扑斯态射
$\Sh(\mathcal{C}) \to \Sh(\mathcal{D}')$，便得到下列陈述：
% P01-ERRATA: P01-E0116 -- source says plural "Lemmas" before a single reference to lemma-topos-good-site.
任意拓扑斯态射 $f : \Sh(\mathcal{C}) \to \Sh(\mathcal{D})$ 都嵌入交换图
$$
\xymatrix{
\Sh(\mathcal{C}) \ar[d]_g \ar[r]_f &
\Sh(\mathcal{D}) \ar[d]^e \\
\Sh(\mathcal{C}') \ar[r]^{f'} &
\Sh(\mathcal{D}')
}
$$
并满足下列性质：
\begin{enumerate}
\item 态射 $e$ 与 $g$ 是由特殊余连续函子
$\mathcal{C} \to \mathcal{C}'$ 和 $\mathcal{D} \to \mathcal{D}'$ 给出的等价；
\item 位点 $\mathcal{C}'$ 与 $\mathcal{D}'$ 有纤维积、终对象，并具有
次典范拓扑；
\item 态射 $f' : \mathcal{C}' \to \mathcal{D}'$ 来自一个位点态射，它对应于
命题 \ref{sites-proposition-get-morphism} 适用的函子
$u : \mathcal{D}' \to \mathcal{C}'$；且
\item 给定 $\mathcal{C}$（相应地 $\mathcal{D}$）上的任意层集
$\mathcal{F}_i$（相应地 $\mathcal{G}_j$），可以假设其中每个层在
$\mathcal{C}'$（相应地 $\mathcal{D}'$）上都是可表层。
\end{enumerate}
把 $\mathcal{C}$ 和 $\mathcal{D}$ 替换为 $\mathcal{C}'$ 和
$\mathcal{D}'$ 往往很有用。
\end{remark}

\begin{remark}
\label{sites-remark-equivalence-topoi-comes-from-morphism-sites}
记号和假设如引理 \ref{sites-lemma-morphism-topoi-comes-from-morphism-sites}。进一步
假设原拓扑斯态射 $\Sh(\mathcal{C}) \to \Sh(\mathcal{D})$ 是等价。则引理
\ref{sites-lemma-morphism-topoi-comes-from-morphism-sites} 的证明中的构造给出两个
函子
$$
\mathcal{C} \rightarrow \mathcal{C}' \leftarrow \mathcal{D}
$$
它们都是特殊余连续函子。因此，在这种情况下，实际上可以把拓扑斯态射分解为
复合
$$
\Sh(\mathcal{C}) \rightarrow
\Sh(\mathcal{C}') =
\Sh(\mathcal{D}') \leftarrow
\Sh(\mathcal{D})
$$
如备注 \ref{sites-remark-morphism-topoi-comes-from-morphism-sites}，但中间态射为恒等
态射。
\end{remark}









\section{拓扑斯的局部化}
\label{sites-section-localize-topoi}

\noindent
我们把局部化的部分材料重复到表面上更一般的拓扑斯情形。实际上，这并不更
一般，因为总可以扩张底层位点，从而假设在位点的对象处作局部化。

\begin{lemma}
\label{sites-lemma-localize-topos}
设 $\mathcal{C}$ 为位点，$\mathcal{F}$ 是 $\mathcal{C}$ 上的层。则范畴
$\Sh(\mathcal{C})/\mathcal{F}$ 是拓扑斯。存在典范拓扑斯态射
$$
j_\mathcal{F} :
\Sh(\mathcal{C})/\mathcal{F}
\longrightarrow
\Sh(\mathcal{C})
$$
它是第 \ref{sites-section-localize} 节意义下的局部化，并且
\begin{enumerate}
\item 函子 $j_\mathcal{F}^{-1}$ 是函子
$\mathcal{H} \mapsto \mathcal{H} \times \mathcal{F}/\mathcal{F}$；且
\item 函子 $j_{\mathcal{F}!}$ 是遗忘函子
$\mathcal{G}/\mathcal{F} \mapsto \mathcal{G}$。
\end{enumerate}
\end{lemma}

\begin{proof}
应用引理 \ref{sites-lemma-topos-good-site}。这意味着可以假设 $\mathcal{C}$ 是
具有次典范拓扑的位点，并且对某个 $U \in \Ob(\mathcal{C})$ 有
$\mathcal{F} = h_U = h_U^\#$。因此，第 \ref{sites-section-localize} 节的材料适用。
特别地，存在等价
$\Sh(\mathcal{C}/U) = \Sh(\mathcal{C})/h_U^\#$，使复合
$$
\Sh(\mathcal{C}/U)
\to
\Sh(\mathcal{C})/h_U^\#
\to \Sh(\mathcal{C})
$$
等于 $j_{U!}$，见引理 \ref{sites-lemma-essential-image-j-shriek}。以
$a : \Sh(\mathcal{C})/h_U^\# \to \Sh(\mathcal{C}/U)$ 表示逆函子，于是
$j_{\mathcal{F}!} = j_{U!} \circ a$、
$j_\mathcal{F}^{-1} = a^{-1} \circ j_U^{-1}$，以及
$j_{\mathcal{F}, *} = j_{U, *} \circ a$。
% P01-ERRATA: P01-E0117 -- source "Denote a ... the inverse functor" omits "by" or otherwise needs recasting.
对 $j_{\mathcal{F}!}$ 的描述由上述内容得出。对
$j_\mathcal{F}^{-1}$ 的描述由引理
\ref{sites-lemma-compute-j-shriek-restrict} 得出。
\end{proof}

\begin{lemma}
\label{sites-lemma-localize-pushforward}
在引理 \ref{sites-lemma-localize-topos} 的情形中，函子
$j_{\mathcal{F}, *}$ 把 $\varphi : \mathcal{G} \to \mathcal{F}$ 送到层
$$
U
\longmapsto
\{\alpha : \mathcal{F}|_U \to \mathcal{G}|_U
\text{ 使得 } \alpha \text{ 是下列映射的右逆： }\varphi|_U \}.
$$
% P01-ERRATA: P01-E0118 -- source says "is the one associates to" rather than "is the one that associates to".
\end{lemma}

\begin{proof}
对任意 $\varphi : \mathcal{G} \to \mathcal{F}$，用记号
$\mathcal{G}/\mathcal{F}$ 表示 $\Sh(\mathcal{C})/\mathcal{F}$ 的相应对象。有
$$
(j_{\mathcal{F}, *}(\mathcal{G}/\mathcal{F}))(U) =
\Mor_{\Sh(\mathcal{C})}(h_U^\#, j_{\mathcal{F}, *}(\mathcal{G}/\mathcal{F})) =
\Mor_{\Sh(\mathcal{C})/\mathcal{F}}(j_{\mathcal{F}}^{-1}h_U^{\#},
(\mathcal{G}/\mathcal{F})).
$$
由引理 \ref{sites-lemma-localize-topos}，该集合是在集合映射
$$
\Mor_{\Sh(\mathcal{C})}(h_U^\# \times \mathcal{F}, \mathcal{G})
\xrightarrow{\varphi \circ}
\Mor_{\Sh(\mathcal{C})}(h_U^\# \times \mathcal{F}, \mathcal{F}).
$$
下元素 $h_U^\# \times \mathcal{F} \to \mathcal{F}$ 上的纤维。由引理
\ref{sites-lemma-internal-hom-sheaf} 中的伴随，有
\begin{align*}
\Mor_{\Sh(\mathcal{C})}(h_U^{\#}\times\mathcal{F}, \mathcal{G})
& =
\Mor_{\Sh(\mathcal{C})}(h_U^{\#},\SheafHom(\mathcal{F}, \mathcal{G})) \\
& =
\Mor_{\Sh(\mathcal{C}/U)}(\mathcal{F}|_{\mathcal{C}/U},
\mathcal{G}|_{\mathcal{C}/U}), \\
\Mor_{\Sh(\mathcal{C})}(h_U^{\#} \times \mathcal{F}, \mathcal{F})
& =
\Mor_{\Sh(\mathcal{C})}(h_U^{\#},\SheafHom(\mathcal{F},\mathcal{F})) \\
& =
\Mor_{\Sh(\mathcal{C}/U)}(\mathcal{F}|_{\mathcal{C}/U},
\mathcal{F}|_{\mathcal{C}/U}),
\end{align*}
在此伴随下，映射 $\varphi\circ$ 变为映射
$$
\Mor_{\Sh(\mathcal{C}/U)}(\mathcal{F}|_{\mathcal{C}/U},
\mathcal{G}|_{\mathcal{C}/U})
\longrightarrow
\Mor_{\Sh(\mathcal{C}/U)}(\mathcal{F}|_{\mathcal{C}/U},
\mathcal{F}|_{\mathcal{C}/U}),\quad
\psi \longmapsto \varphi|_{\mathcal{C}/U} \circ \psi,
$$
而元素 $h_U^\# \times \mathcal{F} \to \mathcal{F}$ 变为
$\text{id}_{\mathcal{F}|_{\mathcal{C}/U}}$。因此
$(j_{\mathcal{F}, *}\mathcal{G}/\mathcal{F})(U)$ 同构于映射
$$
\Mor_{\Sh(\mathcal{C}/U)}(\mathcal{F}|_{\mathcal{C}/U},
\mathcal{G}|_{\mathcal{C}/U})
\xrightarrow{\varphi|_{\mathcal{C}/U}\circ}
\Mor_{\Sh(\mathcal{C}/U)}(\mathcal{F}|_{\mathcal{C}/U},
\mathcal{F}|_{\mathcal{C}/U}),
$$
在 $\text{id}_{\mathcal{F}|_{\mathcal{C}/U}}$ 上的纤维，即
$\{\alpha : \mathcal{F}|_U \to \mathcal{G}|_U
\text{ 使得 } \alpha \text{ 是下列映射的右逆： }\varphi|_U \}$，
如所需。
\end{proof}




\begin{lemma}
\label{sites-lemma-localize-topos-site}
设 $\mathcal{C}$ 为位点，$\mathcal{F}$ 是 $\mathcal{C}$ 上的层。令
$\mathcal{C}/\mathcal{F}$ 为序对 $(U, s)$ 所成的范畴，其中
$U \in \Ob(\mathcal{C})$ 且 $s \in \mathcal{F}(U)$。规定
$\mathcal{C}/\mathcal{F}$ 的覆盖是一族
$\{(U_i, s_i) \to (U, s)\}$，使 $\{U_i \to U\}$ 是
$\mathcal{C}$ 的覆盖。则
$j : \mathcal{C}/\mathcal{F} \to \mathcal{C}$ 是位点的连续且余连续函子，
它诱导拓扑斯态射
$j : \Sh(\mathcal{C}/\mathcal{F}) \to \Sh(\mathcal{C})$。事实上，存在
等价 $\Sh(\mathcal{C}/\mathcal{F}) =
\Sh(\mathcal{C})/\mathcal{F}$，它把 $j$ 变为 $j_\mathcal{F}$。
\end{lemma}

\begin{proof}
我们略去验证 $\mathcal{C}/\mathcal{F}$ 是位点以及 $j$ 连续且余连续。由引理
\ref{sites-lemma-when-shriek}，存在所述拓扑斯态射 $j$，满足
$j^{-1}\mathcal{G}(U, s) = \mathcal{G}(U)$，并且 $j^{-1}$ 有左伴随
$j_!$。$\mathcal{C}/\mathcal{F}$ 上的态射
$\varphi : * \to j^{-1}\mathcal{G}$，等同于一条规则：给每个序对
$(U, s)$ 指定截面 $\varphi(s) \in \mathcal{G}(U)$，并与限制映射相容。
因此，这又等同于 $\mathcal{C}$ 上的态射
$\varphi : \mathcal{F} \to \mathcal{G}$。得出 $j_!* = \mathcal{F}$。特别地，
对每个 $\mathcal{H} \in \Sh(\mathcal{C}/\mathcal{F})$，都有典范映射
$$
j_!\mathcal{H} \to j_!* = \mathcal{F}
$$
也就是说，得到函子
$j'_! : \Sh(\mathcal{C}/\mathcal{F}) \to
\Sh(\mathcal{C})/\mathcal{F}$。该函子的逆，是把
$\Sh(\mathcal{C})/\mathcal{F}$ 的对象
$\varphi : \mathcal{G} \to \mathcal{F}$ 送到层
$$
a(\mathcal{G}/\mathcal{F}) :
(U, s) \longmapsto \{t \in \mathcal{G}(U) \mid \varphi(t) = s\}
$$
的规则。我们略去验证 $a(\mathcal{G}/\mathcal{F})$ 是层以及 $a$ 是
$j'_!$ 的逆。
\end{proof}

\begin{definition}
\label{sites-definition-localize-topos}
设 $\mathcal{C}$ 为位点，$\mathcal{F}$ 是 $\mathcal{C}$ 上的层。
\begin{enumerate}
\item 拓扑斯 $\Sh(\mathcal{C})/\mathcal{F}$ 称为
{\it 拓扑斯 $\Sh(\mathcal{C})$ 在 $\mathcal{F}$ 处的局部化}。
\item 引理 \ref{sites-lemma-localize-topos} 的拓扑斯态射
$j_\mathcal{F} :
\Sh(\mathcal{C})/\mathcal{F}
\to
\Sh(\mathcal{C})$ 称为{\it 局部化态射}。
\end{enumerate}
\end{definition}

\noindent
我们将证明：每当层 $\mathcal{F}$ 对位点的某个对象 $U$ 等于
$h_U^\#$ 时，拓扑斯的局部化就等于位点在 $U$ 处的局部化上的层范畴。此外，
还将验证所有函子性都与这一等同相容。

\begin{lemma}
\label{sites-lemma-localize-compare}
设 $\mathcal{C}$ 为位点，对 $\mathcal{C}$ 的某个对象 $U$，令
$\mathcal{F} = h_U^\#$。则引理 \ref{sites-lemma-localize-topos} 构造的
$j_\mathcal{F} : \Sh(\mathcal{C})/\mathcal{F}
\to \Sh(\mathcal{C})$ 与第 \ref{sites-section-localize} 节构造的拓扑斯态射
$j_U : \Sh(\mathcal{C}/U) \to \Sh(\mathcal{C})$ 相符；这里使用引理
\ref{sites-lemma-essential-image-j-shriek} 的等同
$\Sh(\mathcal{C}/U) = \Sh(\mathcal{C})/h_U^\#$。
\end{lemma}

\begin{proof}
我们已在引理 \ref{sites-lemma-essential-image-j-shriek} 中看到，复合
$\Sh(\mathcal{C}/U) \to \Sh(\mathcal{C})/h_U^\#
\to \Sh(\mathcal{C})$ 是 $j_{U!}$。由引理
\ref{sites-lemma-localize-topos}，函子
$\Sh(\mathcal{C})/h_U^\# \to \Sh(\mathcal{C})$ 是
$j_{\mathcal{F}!}$。因此，在这一等同下
$j_{\mathcal{F}!} = j_{U!}$。所以
$j_\mathcal{F}^{-1} = j_U^{-1}$（由伴随性），进而
$j_{\mathcal{F}, *} = j_{U, *}$（再次由伴随性）。
\end{proof}

\begin{lemma}
\label{sites-lemma-relocalize-topos}
设 $\mathcal{C}$ 为位点。若
$s : \mathcal{G} \to \mathcal{F}$ 是 $\mathcal{C}$ 上的层态射，则存在自然交换的
拓扑斯态射图式
$$
\xymatrix{
\Sh(\mathcal{C})/\mathcal{G} \ar[rd]_{j_\mathcal{G}} \ar[rr]_j & &
\Sh(\mathcal{C})/\mathcal{F} \ar[ld]^{j_\mathcal{F}} \\
& \Sh(\mathcal{C}) &
}
$$
其中 $j = j_{\mathcal{G}/\mathcal{F}}$ 是拓扑斯
$\Sh(\mathcal{C})/\mathcal{F}$ 在对象 $\mathcal{G}/\mathcal{F}$ 处的局部化。
特别地，有
$$
j^{-1}(\mathcal{H} \to \mathcal{F}) =
(\mathcal{H} \times_\mathcal{F} \mathcal{G} \to \mathcal{G})
$$
以及
$$
j_!(\mathcal{E} \xrightarrow{e} \mathcal{G}) =
(\mathcal{E} \xrightarrow{s \circ e} \mathcal{F}).
$$
\end{lemma}

\begin{proof}
对 $j^{-1}$ 与 $j_!$ 的描述来自引理
\ref{sites-lemma-localize-topos} 中对这些函子的描述。由这些函子作为遗忘函子的描述，
函子等式 $j_{\mathcal{G}!} = j_{\mathcal{F}!} \circ j_!$ 是清楚的。由伴随性还得到等式
$j_\mathcal{G}^{-1} = j^{-1} \circ j_\mathcal{F}^{-1}$ 以及
$j_{\mathcal{G}, *} = j_{\mathcal{F}, *} \circ j_*$。
\end{proof}

\begin{lemma}
\label{sites-lemma-relocalize-compare}
设 $\mathcal{C}$ 与 $s : \mathcal{G} \to \mathcal{F}$ 如引理
\ref{sites-lemma-relocalize-topos} 中所述。若
$\mathcal{G} = h_V^\#$、$\mathcal{F} = h_U^\#$，且
$s : \mathcal{G} \to \mathcal{F}$ 来自 $\mathcal{C}$ 的态射 $V \to U$，
则引理 \ref{sites-lemma-relocalize-topos} 中的图式等同于图式
(\ref{sites-equation-relocalize})；这里使用等同
$\Sh(\mathcal{C}/V) = \Sh(\mathcal{C})/h_V^\#$
与
$\Sh(\mathcal{C}/U) = \Sh(\mathcal{C})/h_U^\#$，
它们来自引理 \ref{sites-lemma-essential-image-j-shriek}。
\end{lemma}

\begin{proof}
这是因为对 $j^{-1}$ 的描述相符。见引理
\ref{sites-lemma-relocalize-explicit} 与引理
\ref{sites-lemma-relocalize-topos}。
\end{proof}






\section{局部化与拓扑斯态射}
\label{sites-section-localize-morphisms-topoi}

\noindent
本节是第 \ref{sites-section-localize-morphisms} 节对拓扑斯态射的类似版本。

\begin{lemma}
\label{sites-lemma-localize-morphism-topoi}
设 $f : \Sh(\mathcal{C}) \to \Sh(\mathcal{D})$
为拓扑斯态射，$\mathcal{G}$ 是 $\mathcal{D}$ 上的层。令
$\mathcal{F} = f^{-1}\mathcal{G}$。则存在拓扑斯的交换图式
$$
\xymatrix{
\Sh(\mathcal{C})/\mathcal{F} \ar[r]_{j_\mathcal{F}} \ar[d]_{f'} &
\Sh(\mathcal{C}) \ar[d]^f \\
\Sh(\mathcal{D})/\mathcal{G} \ar[r]^{j_\mathcal{G}} &
\Sh(\mathcal{D}).
}
$$
态射 $f'$ 由性质
$$
(f')^{-1}(\mathcal{H} \xrightarrow{\varphi} \mathcal{G})
=
(f^{-1}\mathcal{H} \xrightarrow{f^{-1}\varphi} \mathcal{F})
$$
刻画，并且有 $f'_*j_\mathcal{F}^{-1} = j_\mathcal{G}^{-1}f_*$。
\end{lemma}

\begin{proof}
由于命题是关于拓扑斯的，并不涉及底层位点，我们可以任意更换位点。因此，
由评注 \ref{sites-remark-morphism-topoi-comes-from-morphism-sites} 中的讨论，可以假设
$f$ 由连续函子 $u : \mathcal{D} \to \mathcal{C}$ 给出；该函子满足命题
\ref{sites-proposition-get-morphism} 的假设，两个位点都具有所有有限极限与次典范拓扑，
并且对 $\mathcal{D}$ 的某个对象 $V$ 有 $\mathcal{G} = h_V$。于是由引理
\ref{sites-lemma-pullback-representable-sheaf}，
$\mathcal{F} = f^{-1}h_V = h_{u(V)}$。由引理
\ref{sites-lemma-localize-morphism}，得到拓扑斯态射的交换图式
% P01-ERRATA: P01-E0119 -- the proof uses U below without defining U; it should set U = u(V) or use u(V) explicitly.
$$
\xymatrix{
\Sh(\mathcal{C}/U) \ar[r]_{j_U} \ar[d]_{f'} &
\Sh(\mathcal{C}) \ar[d]^f \\
\Sh(\mathcal{D}/V) \ar[r]^{j_V} &
\Sh(\mathcal{D}),
}
$$
并有 $f'_*j_U^{-1} = j_V^{-1}f_*$。由引理
\ref{sites-lemma-localize-compare}，可以把 $j_\mathcal{F}$ 与 $j_U$ 等同，
把 $j_\mathcal{G}$ 与 $j_V$ 等同。对 $(f')^{-1}$ 的描述见引理
\ref{sites-lemma-localize-morphism}。
\end{proof}

\begin{lemma}
\label{sites-lemma-localize-morphism-compare}
设 $f : \mathcal{C} \to \mathcal{D}$ 是由连续函子
$u : \mathcal{D} \to \mathcal{C}$ 给出的位点态射。设 $V$ 是
$\mathcal{D}$ 的对象，令 $U = u(V)$。令
$\mathcal{G} = h_V^\#$，并令
$\mathcal{F} = h_U^\# = f^{-1}h_V^\#$（见引理
\ref{sites-lemma-pullback-representable-sheaf}）。则引理
\ref{sites-lemma-localize-morphism-topoi} 的拓扑斯态射图式与引理
\ref{sites-lemma-localize-morphism} 的拓扑斯态射图式相符；这里使用引理
\ref{sites-lemma-localize-compare} 的等同
$j_\mathcal{F}= j_U$ 与 $j_\mathcal{G} = j_V$。
\end{lemma}

\begin{proof}
这并非全然平凡，因为在引理
\ref{sites-lemma-localize-morphism-topoi} 的证明中，为给出 $f$ 而选取的位点态射，
可能不同于本引理给定的位点态射。但在两种情形下，函子 $(f')^{-1}$ 都由同一规则
描述。因此它们相符，所关联的拓扑斯态射也相同。略去一些细节。
\end{proof}

\begin{lemma}
\label{sites-lemma-relocalize-morphism-topoi}
设 $f : \Sh(\mathcal{C}) \to \Sh(\mathcal{D})$
为拓扑斯态射。设 $\mathcal{G} \in \Sh(\mathcal{D})$、
$\mathcal{F} \in \Sh(\mathcal{C})$，且
$s : \mathcal{F} \to f^{-1}\mathcal{G}$ 是层态射。存在拓扑斯的交换图式
$$
\xymatrix{
\Sh(\mathcal{C})/\mathcal{F} \ar[r]_{j_\mathcal{F}} \ar[d]_{f_s} &
\Sh(\mathcal{C}) \ar[d]^f \\
\Sh(\mathcal{D})/\mathcal{G} \ar[r]^{j_\mathcal{G}} &
\Sh(\mathcal{D}).
}
$$
有 $f_s = f' \circ j_{\mathcal{F}/f^{-1}\mathcal{G}}$，其中
$f' :
\Sh(\mathcal{C})/f^{-1}\mathcal{G}
\to
\Sh(\mathcal{D})/\mathcal{F}$
% P01-ERRATA: P01-E0120 -- the codomain above should be Sh(D)/G; F is a sheaf on C and is ill-typed here.
如引理 \ref{sites-lemma-localize-morphism-topoi} 中所述，而
$j_{\mathcal{F}/f^{-1}\mathcal{G}} :
\Sh(\mathcal{C})/\mathcal{F}
\to
\Sh(\mathcal{C})/f^{-1}\mathcal{G}$
如引理 \ref{sites-lemma-relocalize-topos} 中所述。函子 $(f_s)^{-1}$ 由规则
$$
(f_s)^{-1}(\mathcal{H} \xrightarrow{\varphi} \mathcal{G})
=
(f^{-1}\mathcal{H} \times_{f^{-1}\varphi, f^{-1}\mathcal{G}, s} \mathcal{F}
\rightarrow \mathcal{F}).
$$
刻画。最后，给定任意态射 $b : \mathcal{G}' \to \mathcal{G}$、
$a : \mathcal{F}' \to \mathcal{F}$ 与
$s' : \mathcal{F}' \to f^{-1}\mathcal{G}'$，使得
$$
\xymatrix{
\mathcal{F}' \ar[r]_-{s'} \ar[d]_a & f^{-1}\mathcal{G}' \ar[d]^{f^{-1}b} \\
\mathcal{F} \ar[r]^-s & f^{-1}\mathcal{G}
}
$$
交换，则图式
$$
\xymatrix{
\Sh(\mathcal{C})/\mathcal{F}'
\ar[r]_{j_{\mathcal{F}'/\mathcal{F}}} \ar[d]_{f_{s'}} &
\Sh(\mathcal{C})/\mathcal{F} \ar[d]^{f_s} \\
\Sh(\mathcal{D})/\mathcal{G}' \ar[r]^{j_{\mathcal{G}'/\mathcal{G}}} &
\Sh(\mathcal{D})/\mathcal{G}.
}
$$
交换。
\end{lemma}

\begin{proof}
第一个方块的交换性由引理
\ref{sites-lemma-relocalize-topos} 中图式的交换性与引理
\ref{sites-lemma-localize-morphism-topoi} 中图式的交换性得出。对
$f_s^{-1}$ 的描述，则由引理
\ref{sites-lemma-localize-morphism-topoi} 中对 $(f')^{-1}$ 的描述与引理
\ref{sites-lemma-relocalize-topos} 中对
$(j_{\mathcal{F}/f^{-1}\mathcal{G}})^{-1}$ 的描述合并得出。最后一个方块的
交换性于是由等式
$$
f^{-1}\mathcal{H} \times_{f^{-1}\mathcal{G}, s} \mathcal{F}
\times_\mathcal{F} \mathcal{F}'
=
f^{-1}(\mathcal{H} \times_\mathcal{G} \mathcal{G}')
\times_{f^{-1}\mathcal{G}', s'} \mathcal{F}'
$$
得出，该等式是形式的。
\end{proof}

\begin{lemma}
\label{sites-lemma-relocalize-morphism-compare}
设 $f : \mathcal{C} \to \mathcal{D}$ 是由连续函子
$u : \mathcal{D} \to \mathcal{C}$ 给出的位点态射。设 $V$ 是
$\mathcal{D}$ 的对象，$c : U \to u(V)$ 是态射。令
$\mathcal{G} = h_V^\#$ 且 $\mathcal{F} = h_U^\# = f^{-1}h_V^\#$。
% P01-ERRATA: P01-E0121 -- h_U# need not equal f^{-1}h_V# = h_{u(V)}#; the statement should only set G = h_V# and F = h_U#.
令 $s : \mathcal{F} \to f^{-1}\mathcal{G}$ 为 $c$ 诱导的映射。则引理
\ref{sites-lemma-relocalize-morphism} 的拓扑斯态射图式与引理
\ref{sites-lemma-relocalize-morphism-topoi} 的拓扑斯态射图式相符；这里使用引理
\ref{sites-lemma-localize-compare} 的等同
$j_\mathcal{F} = j_U$ 与 $j_\mathcal{G} = j_V$。
\end{lemma}

\begin{proof}
这由合并引理 \ref{sites-lemma-relocalize-compare} 与
\ref{sites-lemma-localize-morphism-compare} 得出。
\end{proof}



















\section{点}
\label{sites-section-points}

\begin{definition}
\label{sites-definition-point-topos}
设 $\mathcal{C}$ 为位点。{it 拓扑斯 $\Sh(\mathcal{C})$ 的点}，是从
$\Sh(pt)$ 到 $\Sh(\mathcal{C})$ 的拓扑斯态射 $p$。
\end{definition}

\noindent
我们将用函子 $u : \mathcal{C} \to \textit{Sets}$ 定义位点的点。稍后会看到，
$u$ 将定义一个位点态射，它诱导与 $\mathcal{C}$ 关联的拓扑斯的一个点，见引理
\ref{sites-lemma-site-point-morphism}。

\medskip\noindent
设 $\mathcal{C}$ 为位点。令 $p = u$ 为函子
$u : \mathcal{C} \to \textit{Sets}$。引入这种看似奇特的说法，是因为谈论
函子的邻域似乎有些古怪；因此，我们把“点” $p$ 看作由一个范畴数据——即函子
$u$——给出的几何对象。$p$ 实际上等于 $u$ 这一事实并不重要。$p$ 的一个
{it 邻域}是序对 $(U, x)$，其中
$U \in \Ob(\mathcal{C})$ 且 $x \in u(U)$。一个{it 邻域态射}
$(V, y) \to (U, x)$，由 $\mathcal{C}$ 的态射
$\alpha :V \to U$ 给出，并满足 $u(\alpha)(y) = x$。注意，邻域范畴不是
一个“大”范畴。

\medskip\noindent
我们把预层 $\mathcal{F}$ 在 $p$ 处的{it 茎}定义为
\begin{equation}
\label{sites-equation-stalk}
\mathcal{F}_p = \colim_{\{(U, x)\}^{opp}} \mathcal{F}(U).
\end{equation}
该余极限取遍 $p$ 的邻域范畴的反范畴。换言之，$\mathcal{F}_p$ 的一个元素由
三元组 $(U, x, s)$ 给出，其中 $(U, x)$ 是 $p$ 的邻域，且
$s \in \mathcal{F}(U)$。三元组的相等关系，是由
$(U, x, s) \sim (V, y, \alpha^*s)$ 生成的等价关系，其中 $\alpha$ 如上。

\medskip\noindent
注意，若 $\varphi : \mathcal{F} \to \mathcal{G}$ 是集合预层的态射，
则得到典范茎映射
$\varphi_p : \mathcal{F}_p \to \mathcal{G}_p$。因而得到一个{it 茎函子}
$$
\textit{PSh}(\mathcal{C}) \longrightarrow \textit{Sets}, \quad
\mathcal{F} \longmapsto \mathcal{F}_p.
$$
我们用任意函子 $p = u : \mathcal{C} \to \textit{Sets}$ 定义了茎函子；
此定义无需任何条件便能成立\footnote{应当尽量避免对所有 $U$ 都有
$u(U) = \emptyset$ 的情形。}。另一方面，除非 $p$ 确实是一个点（见下述定义），
否则最好不要使用这一概念，因为一般来说茎函子不具有良好性质。

\begin{definition}
\label{sites-definition-point}
设 $\mathcal{C}$ 为位点。位点 $\mathcal{C}$ 的一个{it 点 $p$}，由函子
$u : \mathcal{C} \to \textit{Sets}$ 给出，并且满足
\begin{enumerate}
\item 对 $\mathcal{C}$ 的每个覆盖 $\{U_i \to U\}$，映射
$\coprod u(U_i) \to u(U)$ 是满射。
\item 对 $\mathcal{C}$ 的每个覆盖 $\{U_i \to U\}$ 以及每个态射
$V \to U$，映射
$u(U_i \times_U V) \to u(U_i) \times_{u(U)} u(V)$ 都是双射。
\item 茎函子 $\Sh(\mathcal{C}) \to \textit{Sets}$，
$\mathcal{F} \mapsto \mathcal{F}_p$，是左正合的。
\end{enumerate}
\end{definition}

\noindent
这些条件应当很熟悉，因为它们仿照定义
\ref{sites-definition-continuous} 与
\ref{sites-definition-morphism-sites} 中的条件。注意，(3) 蕴含
$*_p = \{*\}$，见例 \ref{sites-example-singleton-sheaf}。因此至少对某个 $U$ 有
$u(U) \not = \emptyset$（因为空余极限给出空集）。我们将在下文证明（引理
\ref{sites-lemma-point-site-topos}），这确实会诱导拓扑斯
$\Sh(\mathcal{C})$ 的一个点。在此之前，先对一般函子 $u$ 证明若干引理。

\begin{lemma}
\label{sites-lemma-points-recover}
设 $\mathcal{C}$ 为位点，$p = u : \mathcal{C} \to \textit{Sets}$ 是函子。
对 $U \in \Ob(\mathcal{C})$，存在函子性的同构
$(h_U)_p = u(U)$。
\end{lemma}

\begin{proof}
$(h_U)_p$ 的一个元素由三元组 $(V, y, f)$ 给出，其中
$V \in \Ob(\mathcal{C})$、$y\in u(V)$ 且
$f \in h_U(V) = \Mor_\mathcal{C}(V, U)$。若存在态射
$\phi : V \to V'$，满足 $u(\phi)(y) = y'$ 与
$f' \circ \phi = f$，则两个三元组 $(V, y, f)$、$(V', y', f')$
确定同一元素；一般还必须取这类等同的链，才能得到正确的等价关系。无论如何，
每个 $(V, y, f)$ 都等价于元素
$(U, u(f)(y), \text{id}_U)$。若存在上述 $\phi$，则三元组
$(V, y, f)$、$(V', y', f')$ 确定同一个元素
$(U, u(f)(y), \text{id}_U) = (U, u(f')(y'), \text{id}_U)$。
% P01-ERRATA: P01-E0122 -- the source calls these equivalent stalk representatives an "object" and later the "same triple"; both should be "element".
这证明了映射
$u(U) \to (h_U)_p$，$x \mapsto \text{下列对象的类： }(U, x, \text{id}_U)$
是双射。
\end{proof}

\noindent
设 $\mathcal{C}$ 为位点，$p = u : \mathcal{C} \to \textit{Sets}$ 是函子。
类比第 \ref{sites-section-functoriality-PSh} 节中的构造，对给定集合 $E$，按规则
\begin{equation}
\label{sites-equation-skyscraper}
U
\longmapsto
u^pE(U) = \Mor_{\textit{Sets}}(u(U), E) = \text{Map}(u(U), E).
\end{equation}
定义预层 $u^pE$。这定义了函子
$u^p : \textit{Sets} \to \textit{PSh}(\mathcal{C})$，$E \mapsto u^pE$。

\begin{lemma}
\label{sites-lemma-adjoint-point-push-stalk}
对任意函子 $u : \mathcal{C} \to \textit{Sets}$，函子 $u^p$ 是预层上的
茎函子的右伴随。
% P01-ERRATA: P01-E0123 -- the source's first sentence ends after "For any functor ...", leaving a sentence fragment.
\end{lemma}

\begin{proof}
设 $\mathcal{F}$ 是 $\mathcal{C}$ 上的预层，$E$ 是集合。态射
$\mathcal{F} \to u^pE$ 由一族相容映射
$\mathcal{F}(U) \to \text{Map}(u(U), E)$ 给出，也就是由一族相容映射
$\mathcal{F}(U) \times u(U) \to E$ 给出。而按 $\mathcal{F}_p$ 的定义，映射
$\mathcal{F}_p \to E$ 由一条规则给出：它给每个三元组 $(U, x, \sigma)$
指定 $E$ 中的一个元素，并使等价三元组映到同一元素；见围绕方程
(\ref{sites-equation-stalk}) 的讨论。这也正是一族相容映射
$\mathcal{F}(U) \times u(U) \to E$。
\end{proof}

\noindent
类比第 \ref{sites-section-continuous-functors} 节，有如下引理。

\begin{lemma}
\label{sites-lemma-point-pushforward-sheaf}
设 $\mathcal{C}$ 为位点，$p = u : \mathcal{C} \to \textit{Sets}$ 是函子。
假设对 $\mathcal{C}$ 的每个覆盖 $\{U_i \to U\}$，
\begin{enumerate}
\item 映射 $\coprod u(U_i) \to u(U)$ 是满射；且
\item 映射
$u(U_i \times_U U_j) \to u(U_i) \times_{u(U)} u(U_j)$ 是满射。
\end{enumerate}
则有
\begin{enumerate}
\item 对所有集合 $E$，预层 $u^pE$ 都是层，将它记为 $u^sE$；
\item 茎函子 $\Sh(\mathcal{C}) \to \textit{Sets}$ 与函子
$u^s: \textit{Sets} \to \Sh(\mathcal{C})$ 互为伴随；且
\item 对每个集合预层 $\mathcal{F}$，都有
$\mathcal{F}_p = \mathcal{F}^\#_p$。
\end{enumerate}
\end{lemma}

\begin{proof}
第一条断言直接来自层的定义、假设 (1) 和 (2) 以及 $u^pE$ 的定义。第二条是
$u^p$ 与茎函子之间的伴随性（引理
\ref{sites-lemma-adjoint-point-push-stalk}）限制到层之后的重述。第三条断言是因为，
对任意集合 $E$，由层化的伴随性质有
$$
\text{Map}(\mathcal{F}_p, E) =
\Mor_{\textit{PSh}(\mathcal{C})}(\mathcal{F}, u^pE) =
\Mor_{\Sh(\mathcal{C})}(\mathcal{F}^\#, u^sE) =
\text{Map}(\mathcal{F}^\#_p, E)
$$
\end{proof}

\noindent
特别地，当 $p = u$ 是点时，引理
\ref{sites-lemma-point-pushforward-sheaf} 成立。在这种情形下，我们把层 $u^sE$
看作在 $p$ 处取值 $E$ 的“摩天楼”层。

\begin{definition}
\label{sites-definition-pushforward-point}
设 $p$ 是由函子 $u$ 给出的位点 $\mathcal{C}$ 的点。对集合 $E$，定义
$p_*E = u^sE$，即上述引理 \ref{sites-lemma-point-pushforward-sheaf} 描述的层。
我们有时称它为{it 摩天楼层}。
% P01-ERRATA: P01-E0124 -- the source needs punctuation or "to be" in "define p_*E = u^sE the sheaf".
\end{definition}

\noindent
特别地，摩天楼层与茎具有如下伴随性质：
$$
\Mor_{\Sh(\mathcal{C})}(\mathcal{F}, p_*E)
=
\text{Map}(\mathcal{F}_p, E)
$$
这说明记号 $p^{-1}\mathcal{F} = \mathcal{F}_p$ 是自然的，我们有时会使用它。

\begin{lemma}
\label{sites-lemma-point-site-topos}
设 $\mathcal{C}$ 为位点。
\begin{enumerate}
\item 设 $p$ 是位点 $\mathcal{C}$ 的点。则上面引入的函子对
$(p_*, p^{-1})$ 定义拓扑斯态射
$\Sh(pt) \to \Sh(\mathcal{C})$。
\item 设 $p = (p_*, p^{-1})$ 是拓扑斯 $\Sh(\mathcal{C})$ 的点。
则函子 $u : U \mapsto p^{-1}(h_U^\#)$ 诱导位点 $\mathcal{C}$ 的点
$p'$，其关联的拓扑斯态射 $(p'_*, (p')^{-1})$ 等于 $p$。
\end{enumerate}
\end{lemma}

\begin{proof}
(1) 的证明。由上述结果，函子 $p_*$ 与 $p^{-1}$ 互为伴随。根据定义
\ref{sites-definition-point}，函子 $p^{-1}$ 必须是正合的。因此定义
\ref{sites-definition-topos} 中施加的条件全都满足，故 (1) 成立。

\medskip\noindent
(2) 的证明。设 $\{U_i \to U\}$ 是 $\mathcal{C}$ 的覆盖。则
$\coprod (h_{U_i})^\# \to h_U^\#$ 是满射，见引理
\ref{sites-lemma-covering-surjective-after-sheafification}。由于 $p^{-1}$ 是正合的
（由拓扑斯态射的定义），得出 $\coprod u(U_i) \to u(U)$ 是满射。这证明了定义
\ref{sites-definition-point} 的第 (1) 条。层化是正合的，见引理
\ref{sites-lemma-sheafification-exact}。因此，若 $U \times_V W$ 在
$\mathcal{C}$ 中存在，则
$$
h_{U \times_V W}^\# =  h_U^\# \times_{h_V^\#} h_W^\#
$$
又因 $p^{-1}$ 正合，看到
$u(U \times_V W) = u(U) \times_{u(V)} u(W)$。这证明了定义
\ref{sites-definition-point} 的第 (2) 条。令 $p' = u$，并令
$\mathcal{F}_{p'}$ 为用 $u$ 按方程
(\ref{sites-equation-stalk}) 定义的茎函子。存在典范比较映射
$c : \mathcal{F}_{p'} \to \mathcal{F}_p = p^{-1}\mathcal{F}$。具体而言，
给定表示 $\mathcal{F}_{p'}$ 的元素 $\xi$ 的三元组 $(U, x, \sigma)$，
把 $\sigma$ 看作映射 $\sigma : h_U^\# \to \mathcal{F}$；由于
$x \in u(U) = p^{-1}(h_U^\#)$，可以令
$c(\xi) = p^{-1}(\sigma)(x)$。由引理
\ref{sites-lemma-points-recover}，有 $(h_U)_{p'} = u(U)$。因为定义
\ref{sites-definition-point} 的条件 (1) 与 (2) 对 $p'$ 成立，由引理
\ref{sites-lemma-point-pushforward-sheaf}，还有
$(h_U^\#)_{p'} = (h_U)_{p'}$。因此
$$
(h_U^\#)_{p'} = (h_U)_{p'} = u(U) = p^{-1}(h_U^\#)
$$
我们断言这个双射等于比较映射
$c : (h_U^\#)_{p'} \to p^{-1}(h_U^\#)$（略去验证）。$\mathcal{C}$ 上任意层
都是形如 $h_U^\#$ 的层的余积之间映射的余等化子，见引理
\ref{sites-lemma-sheaf-coequalizer-representable}。茎函子
$\mathcal{F} \mapsto \mathcal{F}_{p'}$ 与函子 $p^{-1}$ 都与任意余极限交换
（因为二者都是左伴随）。得出对每个层 $\mathcal{F}$，$c$ 都是同构。因此茎函子
$\mathcal{F} \mapsto \mathcal{F}_{p'}$ 同构于 $p^{-1}$，特别地，它是正合的。
这证明了定义 \ref{sites-definition-point} 的条件 (3) 成立，且 $p'$ 是点。最后一条断言
已在上面证明，因为我们已经看到 $p^{-1} = (p')^{-1}$。
\end{proof}

\noindent
事实上，一个点总对应于某个位点态射，如下一个引理所示。

\begin{lemma}
\label{sites-lemma-site-point-morphism}
设 $\mathcal{C}$ 为位点，$p$ 是由
$u : \mathcal{C} \to \textit{Sets}$ 给出的 $\mathcal{C}$ 的点。设 $S_0$
为无限集，使得对所有 $U \in \Ob(\mathcal{C})$ 都有
$u(U) \subset S_0$。令 $\mathcal{S}$ 为评注
\ref{sites-remark-pt-topos} 中由幂集 $S = \mathcal{P}(S_0)$ 构造的位点。则
\begin{enumerate}
\item 存在等价 $i : \Sh(pt) \to \Sh(\mathcal{S})$；
\item 函子 $u : \mathcal{C} \to \mathcal{S}$ 诱导位点态射
$f : \mathcal{S} \to \mathcal{C}$；且
\item 复合
$$
\Sh(pt) \to
\Sh(\mathcal{S}) \to
\Sh(\mathcal{C})
$$
是引理 \ref{sites-lemma-point-site-topos} 的拓扑斯态射 $(p_*, p^{-1})$。
\end{enumerate}
\end{lemma}

\begin{proof}
第 (1) 条已在评注
\ref{sites-remark-pt-topos} 中看到。此外，回忆该等价把集合 $E$ 关联到
$\mathcal{S}$ 上的层 $i_*E$，它由规则
$V \mapsto \Mor_{\textit{Sets}}(V, E)$ 定义。第 (2) 条由
$\mathcal{C}$ 的点的定义（定义 \ref{sites-definition-point}）与位点态射的定义
（定义 \ref{sites-definition-morphism-sites}）直接得出。最后，考虑 $f_*i_*E$。
按构造有
$$
f_*i_*E(U) = i_*E(u(U)) = \Mor_{\textit{Sets}}(u(U), E)
$$
这等于 $p_*E(U)$，见方程
(\ref{sites-equation-skyscraper})。这证明了 (3)。
\end{proof}

\noindent
与拓扑情形不同，取值 $E$ 的摩天楼层的茎不总是 $E$。一般成立的是如下结论。

\begin{lemma}
\label{sites-lemma-stalk-skyscraper}
设 $\mathcal{C}$ 为位点，
$p : \Sh(pt) \to \Sh(\mathcal{C})$ 是与 $\mathcal{C}$ 关联的拓扑斯的点。
对任意集合 $E$，存在典范映射
$$
E \longrightarrow (p_*E)_p \longrightarrow E
$$
其复合为 $\text{id}_E$。
\end{lemma}

\begin{proof}
总有伴随映射 $(p_*E)_p = p^{-1}p_*E \to E$。当
$E = \{*\}$ 时，这一映射是同构，因为 $p_*$ 与 $p^{-1}$ 都是左正合的，
因而把终对象变为终对象。因此，给定 $e \in E$，可以考虑映射
$i_e : \{*\} \to E$，它给出
$$
\xymatrix{
p^{-1}p_*\{*\} \ar[rr]_{p^{-1}p_*i_e} \ar[d]_{\cong} & & p^{-1}p_*E \ar[d] \\
\{*\} \ar[rr]^{i_e} & & E
}
$$
由此得到所需映射 $E \to (p_*E)_p = p^{-1}p_*E$。
\end{proof}

\begin{lemma}
\label{sites-lemma-skyscraper-functor-exact}
设 $\mathcal{C}$ 为位点，
$p : \Sh(pt) \to \Sh(\mathcal{C})$ 是与 $\mathcal{C}$ 关联的拓扑斯的点。
函子 $p_* : \textit{Sets} \to \Sh(\mathcal{C})$ 具有如下性质：它与任意极限
交换；它是左正合的；它是忠实的；它把满射变为满射；它与余等化子交换；它反映
单射、满射和同构。
\end{lemma}

\begin{proof}
因为 $p_*$ 是右伴随，所以它与任意极限交换且左正合。$p^{-1}p_*E \to E$
是典范分裂满射这一事实，蕴含 $p_*$ 是忠实的，并反映单射、满射与同构。由引理
\ref{sites-lemma-point-site-topos}，可以假设 $p$ 来自底层位点 $\mathcal{C}$ 的一个点
$u : \mathcal{C} \to \textit{Sets}$。在此情形下，层 $p_*E$ 由
$$
p_*E(U) = \Mor_{\textit{Sets}}(u(U), E)
$$
给出，见方程 \ref{sites-equation-skyscraper} 与定义
\ref{sites-definition-pushforward-point}。从该公式立即得出，$p_*$ 把满射变为满射，
并把余等化子变为余等化子。
\end{proof}






\section{构造点}
\label{sites-section-construct-points}

\noindent
本节给出判据，判断从位点到集合范畴的函子何时定义该位点的点。

\begin{lemma}
\label{sites-lemma-neighbourhoods-cofiltered}
设 $\mathcal{C}$ 为位点，$p = u : \mathcal{C} \to \textit{Sets}$ 是函子。
若 $p$ 的邻域范畴是余滤过的，则茎函子
(\ref{sites-equation-stalk}) 是左正合的。
\end{lemma}

\begin{proof}
设 $\mathcal{I} \to \Sh(\mathcal{C})$，
$i \mapsto \mathcal{F}_i$ 是层的有限图式。需要证明该系统极限的茎与茎的极限
相符。令 $\mathcal{F}$ 为该系统作为{it 预层}的极限。根据引理
\ref{sites-lemma-limit-sheaf}，它是层，也是层范畴中的极限。因此只须证明
$\mathcal{F}_p = \lim_\mathcal{I} \mathcal{F}_{i, p}$。还要回忆
$\mathcal{F}$ 具有简单描述，见第
\ref{sites-section-limits-colimits-PSh} 节。故只须证明
$$
\lim_i \colim_{\{(U, x)\}^{opp}} \mathcal{F}_i(U)
=
\colim_{\{(U, x)\}^{opp}} \lim_i \mathcal{F}_i(U).
$$
这由 Categories, Lemma
\ref{categories-lemma-directed-commutes} 成立，因为按假设，邻域范畴的反范畴
是滤过的。
\end{proof}

\begin{lemma}
\label{sites-lemma-neighbourhoods-directed}
设 $\mathcal{C}$ 为位点。假设 $\mathcal{C}$ 有终对象 $X$ 和纤维积。设
$p = u : \mathcal{C} \to \textit{Sets}$ 是函子，满足
\begin{enumerate}
\item $u(X)$ 是单点集；且
\item 对每一对具有同一余域的态射 $U \to W$ 与 $V \to W$，映射
$u(U \times_W V) \to u(U) \times_{u(W)} u(V)$ 是双射。
\end{enumerate}
则 $p$ 的邻域范畴是余滤过的，因而茎函子
$\Sh(\mathcal{C}) \to \textit{Sets}$，
$\mathcal{F} \to \mathcal{F}_p$，与有限极限交换。
% P01-ERRATA: P01-E0125 -- the source has the duplicated phrase "Then the the category".
\end{lemma}

\begin{proof}
请注意它与引理 \ref{sites-lemma-directed} 的相似性。关于 $\mathcal{C}$ 的假设蕴含
$\mathcal{C}$ 有有限极限。见 Categories, Lemma
\ref{categories-lemma-finite-limits-exist}。假设 (1) 蕴含邻域范畴非空。
设 $(U, x)$ 与 $(V, y)$ 是邻域。则由 (2)，
$u(U \times V) = u(U \times_X V) =
u(U) \times_{u(X)} u(V) = u(U) \times u(V)$。因此存在同时映到
$(U, x)$ 与 $(V, y)$ 的邻域 $(U \times_X V, z)$。
设 $a, b : (V, y) \to (U, x)$ 是邻域范畴中的两个态射。令 $W$ 为
$a, b : V \to U$ 的等化子。与 Categories, Lemma
\ref{categories-lemma-finite-limits-exist} 的证明相同，可以用纤维积写出 $W$：
$$
W = (V \times_{a, U, b} V) \times_{(pr_1, pr_2), V \times V, \Delta} V
$$
(2) 中的双射性保证存在 $z \in u(W)$ 映到 $((y, y), y)$。于是
$(W, z) \to (V, y)$ 如所需地等化 $a, b$。“因而”部分就是引理
\ref{sites-lemma-neighbourhoods-cofiltered}。
\end{proof}

\begin{proposition}
\label{sites-proposition-point-limits}
设 $\mathcal{C}$ 为位点，并假设 $\mathcal{C}$ 中存在有限极限（即
$\mathcal{C}$ 有纤维积和终对象）。这种位点 $\mathcal{C}$ 的一个点 $p$，由函子
$u : \mathcal{C} \to \textit{Sets}$ 给出，并满足
\begin{enumerate}
\item $u$ 与有限极限交换；且
\item 若 $\{U_i \to U\}$ 是覆盖，则
$\coprod_i u(U_i) \to u(U)$ 是满射。
\end{enumerate}
\end{proposition}

\begin{proof}
先设 $p$ 是由函子 $u$ 给出的点（定义
\ref{sites-definition-point}）。条件 (2) 直接由点的定义满足。由引理
\ref{sites-lemma-points-recover}，有 $(h_U)_p = u(U)$。由引理
\ref{sites-lemma-point-pushforward-sheaf}，有
$(h_U^\#)_p = (h_U)_p$。因此看到 $u$ 等于函子的复合
$$
\mathcal{C} \xrightarrow{h}
\textit{PSh}(\mathcal{C}) \xrightarrow{{}^\#}
\Sh(\mathcal{C}) \xrightarrow{()_p}
\textit{Sets}
$$
这些函子中的每一个都是左正合的，故 $u$ 满足 (1)。

\medskip\noindent
反之，设 $u$ 满足 (1) 和 (2)。此时立即看出，$u$ 满足定义
\ref{sites-definition-point} 的前两个条件。它的茎函子是正合的，因为由引理
\ref{sites-lemma-point-pushforward-sheaf} 它是左伴随，而由引理
\ref{sites-lemma-neighbourhoods-directed} 它与有限极限交换。
\end{proof}

\begin{remark}
\label{sites-remark-improve-proposition-points-limits}
事实上，设 $\mathcal{C}$ 为位点，假设 $\mathcal{C}$ 有终对象 $X$ 与纤维积。
设 $p = u: \mathcal{C} \to \textit{Sets}$ 为函子，满足
\begin{enumerate}
\item $u(X) = \{*\}$ 是单点集；
% P01-ERRATA: P01-E0126 -- the source omits "is" in "$u(X)=\{*\}$ a singleton".
\item 对每一对具有同一余域的态射 $U \to W$ 与 $V \to W$，映射
$u(U \times_W V) \to u(U) \times_{u(W)} u(V)$ 是满射。
\item 对每个覆盖 $\{U_i \to U\}$，映射
$\coprod u(U_i) \to u(U)$ 是满射。
\end{enumerate}
一般来说，$p$ {\bf 不是} $\mathcal{C}$ 的点。一个例子是具有两个对象
$\{U, X\}$ 的范畴 $\mathcal{C}$，其中恰有一个非恒等箭头，即
$U \to X$。在 $\mathcal{C}$ 上赋予平凡拓扑，即仅有的覆盖是
$\{U \to U\}$ 与 $\{X \to X\}$。层 $\mathcal{F}$ 与预层是一回事，并由
三元组 $(A, B, A \to B)$ 组成：即 $A = \mathcal{F}(X)$、
$B = \mathcal{F}(U)$，而 $A \to B$ 是对应于 $U \to X$ 的限制映射。
注意 $U \times_X U = U$，故纤维积存在。考虑函子 $u = p$，其中
$u(X) = \{*\}$ 且 $u(U) = \{*_1, *_2\}$。它满足 (1)、(2) 和 (3)，
但相应的茎函子 (\ref{sites-equation-stalk}) 是函子
$$
(A, B, A \to B) \longmapsto B \amalg_A B
$$
它不是正合的。具体来说，考虑层的单射
$(\emptyset, \{1\}, \emptyset \to \{1\}) \to (\{1\}, \{1\}, \{1\} \to \{1\})$，
茎函子把它变为集合间的非单射
$$
\{1\} \amalg \{1\} \longrightarrow \{1\} \amalg_{\{1\}} \{1\}
$$
\end{remark}

\begin{example}
\label{sites-example-point-topological}
设 $X$ 为拓扑空间，$X_{Zar}$ 是例
\ref{sites-example-site-topological} 的位点。设 $x \in X$ 是一点。考虑函子
$$
u : X_{Zar} \longrightarrow \textit{Sets}, \quad
U \mapsto
\left\{
\begin{matrix}
\emptyset & \text{若} & x \not \in U \\
\{*\} & \text{若} & x \in U
\end{matrix}
\right.
$$
该函子与积和纤维积交换，并把覆盖变为满射族。因此得到位点
$X_{Zar}$ 的一个点 $p$。立即可验证，茎函子与 Sheaves, Section
\ref{sheaves-section-stalks} 中定义的 $x$ 处的茎相符。
\end{example}

\begin{example}
\label{sites-example-point-topology}
设 $X$ 为拓扑空间。拓扑斯 $\Sh(X)$ 的点是什么？为看清这一点，令
$X_{Zar}$ 为例 \ref{sites-example-site-topological} 的位点。由引理
\ref{sites-lemma-point-site-topos}，$\Sh(X)$ 的点对应于该位点的点。设 $p$ 是
由函子 $u : X_{Zar} \to \textit{Sets}$ 给出的位点 $X_{Zar}$ 的点。
我们将使用命题 \ref{sites-proposition-point-limits} 中对此类 $u$ 的刻画。它立即蕴含
$u(\emptyset) = \emptyset$ 与 $u(U \cap V) = u(U) \times u(V)$。特别地，
通过对角映射有 $u(U) = u(U) \times u(U)$，这蕴含 $u(U)$ 或为单点集，
或为空集。此外，若 $U = \bigcup U_i$ 是开覆盖，则
$$
u(U) = \emptyset \Rightarrow \forall i, \ u(U_i) = \emptyset
\quad\text{且}\quad
u(U) \not = \emptyset \Rightarrow \exists i, \ u(U_i) \not = \emptyset.
$$
由此得出，存在唯一的最大开集 $W \subset X$ 满足 $u(W) = \emptyset$，即所有满足
$u(U) = \emptyset$ 的开集 $U$ 的并。令 $Z = X \setminus W$。若
$Z = Z_1 \cup Z_2$，其中 $Z_i \subset Z$ 闭，则
$W = (X \setminus Z_1) \cap (X \setminus Z_2)$，故
$\emptyset = u(W) = u(X \setminus Z_1) \times u(X \setminus Z_2)$；于是
$u(X \setminus Z_1) = \emptyset$ 或
$u(X \setminus Z_2) = \emptyset$。这意味着
$X \setminus Z_1 = W$ 或 $X \setminus Z_2 = W$。换言之，$Z$ 不可约。
现在看到 $u$ 由规则
$$
u : X_{Zar} \longrightarrow \textit{Sets}, \quad
U \mapsto
\left\{
\begin{matrix}
\emptyset & \text{若} & Z \cap U = \emptyset \\
\{*\} & \text{若} & Z \cap U \not = \emptyset
\end{matrix}
\right.
$$
描述。注意，对任意不可约闭子集 $Z \subset X$，该函子都满足命题
\ref{sites-proposition-point-limits} 的假设 (1)、(2)，因而定义一个点。换言之，
位点 $X_{Zar}$ 的点与 $X$ 的不可约闭子集一一对应。特别地，若 $X$ 是清醒拓扑
空间，则 $X_{Zar}$ 的点与 $X$ 的点一一对应，见例
\ref{sites-example-point-topological}。
\end{example}

\begin{example}
\label{sites-example-point-G-sets}
考虑例 \ref{sites-example-site-on-group} 与第
\ref{sites-section-example-sheaf-G-sets} 节描述的位点 $\mathcal{T}_G$。遗忘函子
$u : \mathcal{T}_G \to \textit{Sets}$ 与积和纤维积交换，并把覆盖变为满射族。
因此它定义 $\mathcal{T}_G$ 的一个点。我们把
$\Sh(\mathcal{T}_G)$ 与 $G\textit{-Sets}$ 等同。茎函子
$$
p^{-1} :
\Sh(\mathcal{T}_G) = G\textit{-Sets}
\longrightarrow
\textit{Sets}
$$
就是遗忘函子。正向像 $p_*$ 是函子
$$
\textit{Sets}
\longrightarrow
\Sh(\mathcal{T}_G) = G\textit{-Sets}
$$
它把集合 $S$ 映到 $G$-集合 $\text{Map}(G, S)$，作用为
$g \cdot \psi = \psi \circ R_g$，其中 $R_g$ 是右乘。特别地，作为集合有
$p^{-1}p_*S = \text{Map}(G, S)$，而引理
\ref{sites-lemma-stalk-skyscraper} 的映射
$S \to \text{Map}(G, S) \to S$ 就是显然的映射。
\end{example}

\begin{example}
\label{sites-example-indiscrete-points}
设 $\mathcal{C}$ 是赋予混沌拓扑的范畴（例
\ref{sites-example-indiscrete}）。对 $\mathcal{C}$ 的每个对象 $U_0$，函子
$u : U \mapsto \Mor_\mathcal{C}(U_0, U)$ 都定义 $\mathcal{C}$ 的一个点
$p$。事实上，定义 \ref{sites-definition-point} 的条件 (1) 与 (2) 是直接的，
因为仅有的覆盖由恒等映射给出。条件 (3) 成立，因为
$\mathcal{F}_p = \mathcal{F}(U_0)$，并且由于该拓扑是混沌拓扑，在 $U_0$
上取截面是正合函子。
% P01-ERRATA: P01-E0127 -- the source repeats "Condition (2)" where exactness proves condition (3).
% P01-ERRATA: P01-E0128 -- the source calls the topology "discrete" although Example example-indiscrete names it chaotic/indiscrete.
\end{example}





\section{点与拓扑斯态射}
\label{sites-section-functorial-points}

\noindent
本节对点与拓扑斯态射作若干评注。

\begin{lemma}
\label{sites-lemma-point-functor}
设 $u : \mathcal{C} \to \mathcal{D}$ 是函子，
$v : \mathcal{D} \to \textit{Sets}$ 是函子，并令
$w = v \circ u$。分别以 $q$、$p$ 表示与 $v$、$w$ 关联的茎函子
(\ref{sites-equation-stalk})。则对 $\mathcal{C}$ 上的预层 $\mathcal{F}$，函子性地有
$(u_p\mathcal{F})_q = \mathcal{F}_p$。
\end{lemma}

\begin{proof}
这是一个简单的范畴事实。我们有
\begin{align*}
(u_p\mathcal{F})_q
& =
\colim_{(V, y)} \colim_{U, \phi : V \to u(U)} \mathcal{F}(U) \\
& = \colim_{(V, y, U, \phi : V \to u(U))} \mathcal{F}(U) \\
& = \colim_{(U, x)} \mathcal{F}(U) \\
& = \mathcal{F}_p
\end{align*}
第一个等式由 $u_p$ 的定义与茎函子的定义成立。注意 $y \in v(V)$。在第二个
等式中，只是合并了余极限。为说明第三个等式，把 Categories, Lemma
\ref{categories-lemma-colimit-constant-connected-fibers} 应用于按规则
$$
F((V, y, U, \phi : V \to u(U))) = (U, v(\phi)(y)).
$$
定义的图式范畴之间的函子 $F$。这是有意义的，因为
$w(U) = v(u(U))$。来验证 Categories, Lemma
\ref{categories-lemma-colimit-constant-connected-fibers} 的假设。注意 $F$ 有右逆
$(U, x) \mapsto (u(U), x, U, \text{id} : u(U) \to u(U))$。这同样有意义，
因为 $x \in w(U) = v(u(U))$。另一方面，总有 $(U, x)$ 上的纤维范畴中的态射
$$
(V, y, U, \phi : V \to u(U))
\longrightarrow
(u(U), v(\phi)(y), U, \text{id} : u(U) \to u(U))
$$
这表明各纤维范畴连通。第四个也是最后一个等式是清楚的。
\end{proof}

\begin{lemma}
\label{sites-lemma-point-morphism-sites}
\begin{slogan}
位点的映射定义点上的映射，而且拉回保持茎。
\end{slogan}
设 $f : \mathcal{D} \to \mathcal{C}$ 是由连续函子
$u : \mathcal{C} \to \mathcal{D}$ 给出的位点态射。设 $q$ 是由函子
$v : \mathcal{D} \to \textit{Sets}$ 给出的 $\mathcal{D}$ 的点，见定义
\ref{sites-definition-point}。则函子
$v \circ u : \mathcal{C} \to \textit{Sets}$ 定义 $\mathcal{C}$ 的一个点
$p$；此外，对 $\mathcal{C}$ 上的任意层 $\mathcal{F}$，存在典范等同
$$
(f^{-1}\mathcal{F})_q = \mathcal{F}_p
$$
\end{lemma}

\begin{proof}[引理 \ref{sites-lemma-point-morphism-sites} 的第一种证明]
注意，由于 $u$ 连续且 $v$ 定义一个点，$v \circ u$ 立即满足定义
\ref{sites-definition-point} 的条件 (1) 与 (2)。来证明所显示的等式。设
$\mathcal{F}$ 是 $\mathcal{C}$ 上的层。则
$$
(f^{-1}\mathcal{F})_q = (u_s\mathcal{F})_q =
(u_p \mathcal{F})_q = \mathcal{F}_p
$$
第一个等式因为 $f^{-1} = u_s$，第二个等式由引理
\ref{sites-lemma-point-pushforward-sheaf}，第三个等式由引理
\ref{sites-lemma-point-functor}。现在又看到，$p$ 也满足定义
\ref{sites-definition-point} 的条件 (3)，因为它是正合函子的复合。证明完毕。
\end{proof}

\begin{proof}[引理 \ref{sites-lemma-point-morphism-sites} 的第二种证明]
由引理
\ref{sites-lemma-site-point-morphism}，可把 $(q_*, q^{-1})$ 分解为
$$
\Sh(pt) \xrightarrow{i} \Sh(\mathcal{S})
\xrightarrow{h} \Sh(\mathcal{D})
$$
其中第二个拓扑斯态射来自由函子
$v : \mathcal{D} \to \mathcal{S}$ 诱导的位点态射
$h : \mathcal{S} \to \mathcal{D}$（这是有意义的，因为
$\mathcal{S} \subset \textit{Sets}$ 是包含 $v$ 的像中每个对象的满子范畴）。由引理
\ref{sites-lemma-composition-morphisms-sites}，复合
$v \circ u : \mathcal{C} \to \mathcal{S}$ 定义位点态射
$g : \mathcal{S} \to \mathcal{C}$。特别地，函子
$v \circ u : \mathcal{C} \to \mathcal{S}$ 连续；根据
$\mathcal{S}$ 中覆盖的定义（见评注
\ref{sites-remark-pt-topos}），这意味着 $v \circ u$ 满足定义
\ref{sites-definition-point} 的条件 (1) 与 (2)。另一方面，看到
$$
g_*i_*E(U) = i_*E(v(u(U)) = \Mor_{\textit{Sets}}(v(u(U)), E)
$$
% P01-ERRATA: P01-E0129 -- the frozen display is missing a closing parenthesis after i_*E(v(u(U))).
这是由评注 \ref{sites-remark-pt-topos} 中对 $i$ 的构造得出的。注意，这与如下函子的
公式相同，而它等于 $(v \circ u)^pE$，见方程
(\ref{sites-equation-skyscraper})。
% P01-ERRATA: P01-E0130 -- the source says "the same as the formula for which is equal to" and omits the object of "for".
由引理
\ref{sites-lemma-point-pushforward-sheaf}，函子
$g_*i_* = (v \circ u)^p = (v \circ u)^s$ 是茎函子
$\mathcal{F} \mapsto \mathcal{F}_q$ 的右伴随。
% P01-ERRATA: P01-E0131 -- this stalk functor is p^{-1} on Sh(C), so the subscript should be p, not q.
因此看到茎函子 $p^{-1}$ 典范同构于 $i^{-1} \circ g^{-1}$。所以它是正合的，
并得出 $p$ 是点。最后，由构造有 $g = f \circ h$，故
$p^{-1} = i^{-1} \circ h^{-1} \circ f^{-1} = q^{-1} \circ f^{-1}$；
也就是说，得到本引理所显示的公式。
\end{proof}

\begin{lemma}
\label{sites-lemma-point-morphism-topoi}
设 $f : \Sh(\mathcal{D}) \to \Sh(\mathcal{C})$ 是拓扑斯态射，
$q : \Sh(pt) \to \Sh(\mathcal{D})$ 是一个点。则
$p = f \circ q$ 是拓扑斯 $\Sh(\mathcal{C})$ 的点，而且对
$\mathcal{C}$ 上的任意层 $\mathcal{F}$，存在典范等同
$$
(f^{-1}\mathcal{F})_q = \mathcal{F}_p
$$
\end{lemma}

\begin{proof}
这由定义以及拓扑斯态射可以复合这一事实立即得出。
\end{proof}





\section{局部化与点}
\label{sites-section-localize-points}

\noindent
本节证明，局部化 $\mathcal{C}/U$ 的点可以用简单方式从 $\mathcal{C}$ 的点构造。

\begin{lemma}
\label{sites-lemma-point-localize}
设 $\mathcal{C}$ 为位点，$p$ 是由
$u : \mathcal{C} \to \textit{Sets}$ 给出的 $\mathcal{C}$ 的点。设 $U$ 是
$\mathcal{C}$ 的对象，且 $x \in u(U)$。函子
$$
v : \mathcal{C}/U \longrightarrow \textit{Sets}, \quad
(\varphi : V \to U) \longmapsto \{y \in u(V) \mid u(\varphi)(y) = x\}
$$
定义位点 $\mathcal{C}/U$ 的一个点 $q$，使图式
$$
\xymatrix{
& \Sh(pt) \ar[d]^p \ar[ld]_q \\
\Sh(\mathcal{C}/U) \ar[r]^{j_U} &
\Sh(\mathcal{C})
}
$$
交换。换言之，对 $\mathcal{C}$ 上的任意层有
$\mathcal{F}_p = (j_U^{-1}\mathcal{F})_q$。
% P01-ERRATA: P01-E0132 -- the source says "for any sheaf on C" but omits the sheaf symbol F.
\end{lemma}

\begin{proof}
如引理 \ref{sites-lemma-site-point-morphism} 中那样选取 $S$ 与 $\mathcal{S}$。
可以如该引理中一样等同 $\Sh(pt) = \Sh(\mathcal{S})$，并把由 $u$ 诱导的
拓扑斯态射写成
$p = f : \Sh(\mathcal{S}) \to \Sh(\mathcal{C})$。由引理
\ref{sites-lemma-localize-morphism}，得到拓扑斯的交换图式
$$
\xymatrix{
\Sh(\mathcal{S}/u(U)) \ar[r]_-{j_{u(U)}} \ar[d]_{p'} &
\Sh(\mathcal{S}) \ar[d]^p \\
\Sh(\mathcal{C}/U) \ar[r]^{j_U} &
\Sh(\mathcal{C}),
}
$$
其中 $p'$ 由函子 $u' : \mathcal{C}/U \to \mathcal{S}/u(U)$ 给出，
$V/U \mapsto u(V)/u(U)$。考虑函子
$j_x : \mathcal{S} \cong \mathcal{S}/x$，它把集合 $E$ 关联到集合 $E$，
并赋予取常值 $x$ 的常值映射 $E \to u(U)$。则 $j_x$ 是全忠实的余连续函子，
并有连续右伴随
$v_x : (\psi : E \to u(U)) \mapsto \psi^{-1}(\{x\})$。注意
$j_{u(U)} \circ j_x = \text{id}_\mathcal{S}$ 且
$v_x \circ u' = v$。这些观察蕴含如下拓扑斯交换图式
$$
\xymatrix{
\Sh(\mathcal{S}) \ar[rd]^a \ar[rdd]_q
\ar `r[rrr] `d[dd]^p [rrdd] & & & \\
& \Sh(\mathcal{S}/u(U)) \ar[r]_-{j_{u(U)}} \ar[d]^{p'} &
\Sh(\mathcal{S}) \ar[d]^p & \\
& \Sh(\mathcal{C}/U) \ar[r]^{j_U} &
\Sh(\mathcal{C}) &
}
$$
具体如下：
\begin{enumerate}
\item 态射
$a : \Sh(\mathcal{S}) \to \Sh(\mathcal{S}/u(U))$ 是与余连续函子
$j_x$ 关联的拓扑斯态射；它等于与连续函子 $v_x$ 关联的态射，见引理
\ref{sites-lemma-cocontinuous-morphism-topoi} 与第
\ref{sites-section-cocontinuous-adjoint} 节。
\item 复合 $p \circ j_{u(U)} \circ a = p$，因为
$j_{u(U)} \circ j_x = \text{id}_\mathcal{S}$。
\item 复合 $p' \circ a$ 给出拓扑斯态射。此外，它是与连续函子
$v_x \circ u' = v$ 关联的拓扑斯态射。因此 $v$ 确实定义
$\mathcal{C}/U$ 的一个点 $q$，并且由构造它嵌入上述图式。
\end{enumerate}
本引理的证明结束。
\end{proof}

\begin{lemma}
\label{sites-lemma-points-above-point}
设 $\mathcal{C}$、$p$、$u$、$U$ 如引理
\ref{sites-lemma-point-localize} 中所述。引理
\ref{sites-lemma-point-localize} 的构造，在位于 $p$ 上方的
$\mathcal{C}/U$ 的点 $q$ 与 $u(U)$ 的元素 $x$ 之间给出一一对应。
\end{lemma}

\begin{proof}
设 $q$ 是由函子 $v : \mathcal{C}/U \to \textit{Sets}$ 给出的
$\mathcal{C}/U$ 的点，并且作为拓扑斯态射有 $j_U \circ q = p$。回忆对
$\mathcal{C}$ 的任意对象 $V$，有
$u(V) = p^{-1}(h_V^\#)$，见引理
\ref{sites-lemma-point-site-topos}。类似地，对 $\mathcal{C}/U$ 的任意对象
$V/U$，有 $v(V/U) = q^{-1}(h_{V/U}^\#)$。考虑如下两个图式
$$
\vcenter{
\xymatrix{
\Mor_{\mathcal{C}/U}(W/U, V/U) \ar[r] \ar[d] &
\Mor_\mathcal{C}(W, V) \ar[d] \\
\Mor_{\mathcal{C}/U}(W/U, U/U) \ar[r] &
\Mor_\mathcal{C}(W, U)
}
}
\quad
\vcenter{
\xymatrix{
h_{V/U}^\# \ar[r] \ar[d] & j_U^{-1}(h_V^\#) \ar[d] \\
h_{U/U}^\# \ar[r] & j_U^{-1}(h_U^\#)
}
}
$$
右侧图式是 $\mathcal{C}/U$ 上预层图式的层化，该预层图式把 $W/U$ 映到
左侧的集合图式。（这里有一个小技术点，即
$(j_U^{-1}h_V)^\# = j_U^{-1}(h_V^\#)$，对 $h_U$ 亦然；见引理
\ref{sites-lemma-technical-pu}。）注意左侧的集合图式是笛卡尔的。由于层化正合
（引理 \ref{sites-lemma-sheafification-exact}），得出右侧图式也是笛卡尔的。

\medskip\noindent
把正合函子 $q^{-1}$ 应用于右侧图式，得到集合的笛卡尔图式
$$
\xymatrix{
v(V/U) \ar[r] \ar[d] & u(V) \ar[d] \\
v(U/U) \ar[r] & u(U)
}
$$
这里使用了 $q^{-1} \circ j^{-1} = p^{-1}$。
% P01-ERRATA: P01-E0133 -- the source uses undefined j^{-1}; the localization functor here is j_U^{-1}.
由于 $U/U$ 是 $\mathcal{C}/U$ 的终对象，$v(U/U)$ 是单点集。因此
$v(U/U)$ 在 $u(U)$ 中的像是一个元素 $x$，而顶部水平映射给出双射
$v(V/U) \to  \{y \in u(V) \mid y \mapsto x\text{ 于 }u(U)\}$，如所需。
\end{proof}

\begin{lemma}
\label{sites-lemma-stalk-j-shriek}
设 $\mathcal{C}$ 为位点，$p$ 是由
$u : \mathcal{C} \to \textit{Sets}$ 给出的 $\mathcal{C}$ 的点，$U$ 是
$\mathcal{C}$ 的对象。对 $\mathcal{C}/U$ 上任意层 $\mathcal{G}$，有
$$
(j_{U!}\mathcal{G})_p =
\coprod\nolimits_q \mathcal{G}_q
$$
其中余积取遍按引理
\ref{sites-lemma-point-localize} 与元素 $x \in u(U)$ 关联的 $\mathcal{C}/U$ 的点 $q$。
\end{lemma}

\begin{proof}
使用引理 \ref{sites-lemma-describe-j-shriek} 中的描述：
$j_{U!}\mathcal{G}$ 是预层
$V \mapsto
\coprod\nolimits_{\varphi \in \Mor_\mathcal{C}(V, U)}
\mathcal{G}(V/_\varphi U)$
所关联的层。此外，由引理
\ref{sites-lemma-point-pushforward-sheaf}，$j_{U!}\mathcal{G}$ 在 $p$ 处的茎等于
该预层的茎。因此
$$
(j_{U!}\mathcal{G})_p =
\colim_{(V, y)} \coprod\nolimits_{\varphi : V \to U} \mathcal{G}(V/_\varphi U)
$$
对该余极限的每个元素 $(V, y, \varphi, s)$，可以指定
$x = u(\varphi)(y) \in u(U)$。因此得到
$$
(j_{U!}\mathcal{G})_p =
\coprod\nolimits_{x \in u(U)}
\colim_{(\varphi : V \to U, y), \ u(\varphi)(y) = x} \mathcal{G}(V/_\varphi U).
$$
根据我们对点 $q$ 的构造，这等于本引理中的表达式。
\end{proof}

\begin{remark}
\label{sites-remark-not-pushforward}
警告：引理 \ref{sites-lemma-stalk-j-shriek} 的结果对 $j_{U, *}$ 没有类似结论。
\end{remark}





\section{拓扑斯的 2-态射}
\label{sites-section-2-category}

\noindent
本节简要讨论拓扑斯的 $2$-态射这一概念。

\begin{definition}
\label{sites-definition-2-morphism-topoi}
设 $f, g : \Sh(\mathcal{C}) \to \Sh(\mathcal{D})$
是两个拓扑斯态射。一个{it 从 $f$ 到 $g$ 的 2-态射}由函子间的变换
$t : f_* \to g_*$ 给出。
\end{definition}

\noindent
在图示中，我们有时把 $t$ 表示如下：
$$
\xymatrix{
\Sh(\mathcal{C})
\rrtwocell^f_g{t}
&
&
\Sh(\mathcal{D})
}
$$
注意，由于 $f^{-1}$ 伴随于 $f_*$，$g^{-1}$ 伴随于 $g_*$，$t$ 还诱导函子间的
变换 $t : g^{-1} \to f^{-1}$（通常用同一符号表示）；它由下列图式交换这一条件
唯一刻画：
$$
\xymatrix{
\Mor_{\Sh(\mathcal{D})}(\mathcal{G}, f_*\mathcal{F})
\ar[d]_{t \circ -} \ar@{=}[r] &
\Mor_{\Sh(\mathcal{C})}(f^{-1}\mathcal{G}, \mathcal{F})
\ar[d]^{- \circ t}
\\
\Mor_{\Sh(\mathcal{D})}(\mathcal{G}, g_*\mathcal{F})
\ar@{=}[r] &
\Mor_{\Sh(\mathcal{C})}(g^{-1}\mathcal{G}, \mathcal{F})
}
$$
由于集合论上的困难（见评注
\ref{sites-remark-morphism-topoi-big}），我们得不到拓扑斯的 2-范畴。不过仍可定义
水平复合与垂直复合，并证明 Categories, Section
\ref{categories-section-2-categories} 所列的严格 2-范畴公理成立。具体来说，
2-态射的垂直复合是清楚的（只须复合函子间的变换）；1-态射的复合已在定义
\ref{sites-definition-topos} 中定义；而
$$
\xymatrix{
\Sh(\mathcal{C})
\rtwocell^f_g{t}
&
\Sh(\mathcal{D})
\rtwocell^{f'}_{g'}{s}
&
\Sh(\mathcal{E})
}
$$
的水平复合，由 Categories, Definition
\ref{categories-definition-horizontal-composition} 中引入的函子间变换
$s \star t$ 定义。明确地说，$s \star t$ 由
$$
\xymatrix{
f'_*f_*\mathcal{F}
\ar[r]^{f'_*t} &
f'_*g_*\mathcal{F}
\ar[r]^s &
g'_*g_*\mathcal{F}
}
\quad\text{或}\quad
\xymatrix{
f'_*f_*\mathcal{F}
\ar[r]^s &
g'_*f_*\mathcal{F}
\ar[r]^{g'_*t} &
g'_*g_*\mathcal{F}
}
$$
给出（这些映射相等）。由于这些定义与 Categories, Section
\ref{categories-section-formal-cat-cat} 中的定义相符，由 Categories, Lemma
\ref{categories-lemma-properties-2-cat-cats} 得出：在这些定义下，严格 2-范畴的
公理成立。




\section{点之间的态射}
\label{sites-section-morphisms-points}

\begin{lemma}
\label{sites-lemma-maps-u-points}
设 $\mathcal{C}$ 为位点，
$u, u' : \mathcal{C} \to \textit{Sets}$ 是两个函子，而
$t : u' \to u$ 是函子间的变换。则得到典范茎函子变换
$t_{stalk} : \mathcal{F}_{p'} \to \mathcal{F}_p$；通过引理
\ref{sites-lemma-points-recover} 中的等同，它与 $t$ 相符。
\end{lemma}

\begin{proof}
略。
\end{proof}

\begin{definition}
\label{sites-definition-morphism-points}
设 $\mathcal{C}$ 为位点。设 $p, p'$ 是由函子
$u, u' : \mathcal{C} \to \textit{Sets}$ 给出的 $\mathcal{C}$ 的点。一个
{it 态射 $f : p \to p'$}，由函子间的变换
$$
f_u : u' \to u.
$$
给出。
\end{definition}

\noindent
注意函子间的变换沿相反方向。正如稍后会看到的，这可以通过把态射 $f$ 看作
一种 $2$-箭头来理解；图示如下：
$$
\xymatrix{
\textit{Sets}
=
\Sh(pt)
\rrtwocell^p_{p'}{f}
&
&
\Sh(\mathcal{C})
}
$$
具体来说，稍后会看到，$f_u$ 在摩天楼层构造之间诱导典范函子变换
$p_* \to p'_*$。

\medskip\noindent
这是一个相当重要的概念，值得在此作更完整的处理。期望事项如下：
\begin{enumerate}
\item 描述例 \ref{sites-example-point-G-sets} 中所述 $\mathcal{T}_G$ 的点的自同构。
\item 用 $\Mor(p_*, p'_*)$ 描述 $\Mor(p, p')$。
\item 拓扑空间中点的特殊化。证明：若在拓扑空间 $X$ 中
$x' \in \overline{\{x\}}$，则存在态射 $p \to p'$，其中 $p$（相应地，$p'$）
是 $X_{Zar}$ 中与 $x$（相应地，$x'$）关联的点。
\end{enumerate}




\section{有足够多点的位点}
\label{sites-section-sites-enough-points}

\begin{definition}
\label{sites-definition-enough-points}
设 $\mathcal{C}$ 为位点。
\begin{enumerate}
\item 若每个层映射 $\phi : \mathcal{F} \to \mathcal{G}$，只要在所有茎
$\mathcal{F}_{p_i} \to \mathcal{G}_{p_i}$ 上都是同构，就必为同构，则称点族
$\{p_i\}_{i\in I}$ 是{it 保守的}。
% P01-ERRATA: P01-E0134 -- the source calls F_{p_i}->G_{p_i} "fibres" although they are stalks.
\item 若存在一个保守点族，则称 $\mathcal{C}$ {it 有足够多的点}。
\end{enumerate}
\end{definition}

\noindent
事实证明，这时可以在茎上检验“正合性”。

\begin{lemma}
\label{sites-lemma-exactness-stalks}
设 $\mathcal{C}$ 为位点，$\{p_i\}_{i\in I}$ 是保守点族。则
\begin{enumerate}
\item 给定任意层映射 $\varphi : \mathcal{F} \to \mathcal{G}$，
$\forall i, \varphi_{p_i}$ 为单射蕴含 $\varphi$ 为单射。
\item 给定任意层映射 $\varphi : \mathcal{F} \to \mathcal{G}$，
$\forall i, \varphi_{p_i}$ 为满射蕴含 $\varphi$ 为满射。
\item 给定任意一对层映射
$\varphi_1, \varphi_2 : \mathcal{F} \to \mathcal{G}$，
$\forall i, \varphi_{1, p_i} = \varphi_{2, p_i}$ 蕴含
$\varphi_1 = \varphi_2$。
\item 给定有限图式 $\mathcal{G} : \mathcal{J}
\to \Sh(\mathcal{C})$、层 $\mathcal{F}$ 以及态射
$q_j : \mathcal{F} \to \mathcal{G}_j$，则 $(\mathcal{F}, q_j)$ 是该图式的
极限，当且仅当对每个 $i$，茎
$(\mathcal{F}_{p_i}, (q_j)_{p_i})$ 是相应极限。
\item 给定有限图式 $\mathcal{F} : \mathcal{J}
\to \Sh(\mathcal{C})$、层 $\mathcal{G}$ 以及态射
$e_j : \mathcal{F}_j \to \mathcal{G}$，则 $(\mathcal{G}, e_j)$ 是该图式的
余极限，当且仅当对每个 $i$，茎
$(\mathcal{G}_{p_i}, (e_j)_{p_i})$ 是相应余极限。
\end{enumerate}
\end{lemma}

\begin{proof}
将反复使用如下事实：所有茎函子都与任意有限极限和余极限交换，因而与积、纤维积
等交换。还将使用单射就是单态射、满射就是满态射这一事实。层映射
$\varphi : \mathcal{F} \to \mathcal{G}$ 为单射，当且仅当
$\mathcal{F} \to \mathcal{F} \times_\mathcal{G} \mathcal{F}$
是同构。因此 (1) 成立。类似地，
$\varphi : \mathcal{F} \to \mathcal{G}$ 为满射，当且仅当
$\mathcal{G} \amalg_\mathcal{F} \mathcal{G} \to \mathcal{G}$
是同构。因此 (2) 成立。映射 $a, b : \mathcal{F} \to \mathcal{G}$ 相等，
当且仅当
$\text{Equalizer}(a, b : \mathcal{F} \to \mathcal{G}) \to \mathcal{F}$
是同构\footnote{等化子可以用积与纤维积表示；见 Categories, Lemma
\ref{categories-lemma-finite-limits-exist} 的证明。}。因此 (3) 成立。
断言 (4) 与 (5) 直接由定义及本证明开头的评注得出。
\end{proof}

\begin{lemma}
\label{sites-lemma-enough}
设 $\mathcal{C}$ 为位点，$\{(p_i, u_i)\}_{i\in I}$ 是点族。该族保守，当且仅当
对每个层 $\mathcal{F}$、每个 $U\in \Ob(\mathcal{C})$ 以及每一对互异截面
$s, s' \in \mathcal{F}(U)$，$s \not = s'$，都存在 $i$ 与
$x\in u_i(U)$，使三元组 $(U, x, s)$ 与 $(U, x, s')$ 定义
$\mathcal{F}_{p_i}$ 的不同元素。
\end{lemma}

\begin{proof}
设该族保守，且 $\mathcal{F}$、$U$ 与 $s, s'$ 如引理中所述。截面 $s$、$s'$
定义两个不同映射 $a, a' : (h_U)^\# \to \mathcal{F}$。因此由引理
\ref{sites-lemma-exactness-stalks}，存在 $i$ 使
$a_{p_i} \not = a'_{p_i}$。回忆由引理
\ref{sites-lemma-points-recover} 与
\ref{sites-lemma-point-pushforward-sheaf}，有
$(h_U)^\#_{p_i} = u_i(U)$。因此存在 $x \in u_i(U)$，使
$\mathcal{F}_{p_i}$ 中 $a_{p_i}(x) \not = a'_{p_i}(x)$。展开定义便看到
$(U, x, s)$ 与 $(U, x, s')$ 如本引理的陈述所述。

\medskip\noindent
为证明逆命题，假设本引理关于点的存在性的条件。设
$\phi : \mathcal{F} \to \mathcal{G}$ 是在所有茎上都为同构的层映射。需要证明
$\phi$ 既单又满，见引理
\ref{sites-lemma-mono-epi-sheaves}。单射性是该假设的直接推论。为证明满射性，需要证明
$\mathcal{G} \amalg_\mathcal{F} \mathcal{G} \to \mathcal{G}$ 是同构
（Categories, Lemma
\ref{categories-lemma-characterize-mono-epi}）。由于该映射显然满，只须检验其单射性；
这又由假设中
$\mathcal{G} \amalg_\mathcal{F} \mathcal{G} \to \mathcal{G}$
在所有茎上都是单射得出。
\end{proof}

\noindent
在下一个引理中，根据引理
\ref{sites-lemma-points-above-point}，点 $q_{i, x}$ 恰好是位于点 $p_i$ 上方的
$\mathcal{C}/U$ 的所有点。

\begin{lemma}
\label{sites-lemma-localize-enough}
设 $\mathcal{C}$ 为位点，$U$ 是 $\mathcal{C}$ 的对象。设
$\{(p_i, u_i)\}_{i\in I}$ 是 $\mathcal{C}$ 的点族。对
$x \in u_i(U)$，令 $q_{i, x}$ 为引理
\ref{sites-lemma-point-localize} 中构造的 $\mathcal{C}/U$ 的点。若
$\{p_i\}$ 是保守点族，则
$\{q_{i, x}\}_{i \in I, x \in u_i(U)}$ 是 $\mathcal{C}/U$ 的保守点族。
特别地，若 $\mathcal{C}$ 有足够多的点，则每个局部化 $\mathcal{C}/U$ 也如此。
% P01-ERRATA: P01-E0135 -- the source begins the third sentence with lowercase "let" after a period.
\end{lemma}

\begin{proof}
我们知道 $j_{U!}$ 诱导等价
$j_{U!} : \Sh(\mathcal{C}/U) \to \Sh(\mathcal{C})/h_U^\#$，见引理
\ref{sites-lemma-essential-image-j-shriek}。此外，还知道
$(j_{U!}\mathcal{G})_{p_i} = \coprod_x \mathcal{G}_{q_{i, x}}$，见引理
\ref{sites-lemma-stalk-j-shriek}。因此结论形式地成立。
\end{proof}

\noindent
下一个引理告诉我们，可以在位点上局部地检验点的存在性。

\begin{lemma}
\label{sites-lemma-enough-points-local}
设 $\mathcal{C}$ 为位点，$\{U_i\}_{i \in I}$ 是 $\mathcal{C}$ 的对象族。假设
\begin{enumerate}
\item $\coprod h_{U_i}^\# \to *$ 是层的满射；且
\item 每个局部化 $\mathcal{C}/U_i$ 都有足够多的点。
\end{enumerate}
则 $\mathcal{C}$ 有足够多的点。
\end{lemma}

\begin{proof}
对每个 $i \in I$，设 $\{p_j\}_{j \in J_i}$ 是 $\mathcal{C}/U_i$ 的保守点族。
对 $j \in J_i$，以
$q_j : \Sh(pt) \to \Sh(\mathcal{C})$ 表示 $p_j$ 与局部化态射
$\Sh(\mathcal{C}/U_i) \to \Sh(\mathcal{C})$ 的复合。
% P01-ERRATA: P01-E0136 -- the source's "denote q_j ... the composition" omits "by" or requires recasting.
则 $q_j$ 是点，见引理
\ref{sites-lemma-point-morphism-topoi}。我们断言点族
$\{q_j\}_{j \in \coprod J_i}$ 保守。事实上，设
$\mathcal{F} \to \mathcal{G}$ 是 $\mathcal{C}$ 上的层映射，使得对所有
$j \in \coprod J_i$，映射
$\mathcal{F}_{q_j} \to \mathcal{G}_{q_j}$ 都是同构。设 $W$ 是
$\mathcal{C}$ 的对象。由假设 (1)，存在覆盖 $\{W_a \to W\}$ 与态射
$W_a \to U_{i(a)}$。由于由引理
\ref{sites-lemma-point-morphism-topoi} 有
$(\mathcal{F}|_{\mathcal{C}/U_{i(a)}})_{p_j} = \mathcal{F}_{q_j}$ 与
$(\mathcal{G}|_{\mathcal{C}/U_{i(a)}})_{p_j} = \mathcal{G}_{q_j}$，
而点族 $\{p_j\}_{j \in J_{i(a)}}$ 保守，所以
$\mathcal{F}|_{U_{i(a)}} \to \mathcal{G}|_{U_{i(a)}}$ 是同构。因此对每个 $a$，
$\mathcal{F}(W_a) \to \mathcal{G}(W_a)$ 是双射。类似地，对每个 $a, b$，
$\mathcal{F}(W_a \times_W W_b) \to \mathcal{G}(W_a \times_W W_b)$
是双射。由层条件，这表明
$\mathcal{F}(W) \to \mathcal{G}(W)$ 是双射，即
$\mathcal{F} \to \mathcal{G}$ 是同构。
\end{proof}

\begin{lemma}
\label{sites-lemma-check-morphism-sites}
设 $u : \mathcal{C} \to \mathcal{D}$ 是位点的连续函子。设
$\{(q_i, v_i)\}_{i\in I}$ 是 $\mathcal{D}$ 的保守点族。若每个函子
$u_i = v_i \circ u$ 都定义 $\mathcal{C}$ 的一个点，则 $u$ 定义位点态射
$f : \mathcal{D} \to \mathcal{C}$。
\end{lemma}

\begin{proof}
以 $p_i$ 表示 $\textit{PSh}(\mathcal{C})$ 上与函子 $u_i$ 对应的茎函子
(\ref{sites-equation-stalk})。有
$$
(f^{-1}\mathcal{F})_{q_i} =
(u_s\mathcal{F})_{q_i} =
(u_p\mathcal{F})_{q_i} =
\mathcal{F}_{p_i}
$$
% P01-ERRATA: P01-E0137 -- f is used before the proof establishes that u defines a morphism of sites; u_s should be used directly until exactness is proved.
第一个等式因为 $f^{-1} = u_s$，第二个等式由引理
\ref{sites-lemma-point-pushforward-sheaf}，第三个等式由引理
\ref{sites-lemma-point-functor}。因此，若 $p_i$ 是点，则先沿 $f$ 拉回、再在
$q_i$ 处取茎是正合函子。由于点族 $\{q_i\}$ 保守，这蕴含 $f^{-1}$ 是正合函子；
于是由定义 \ref{sites-definition-morphism-sites}，看到 $f$ 是位点态射。
\end{proof}





\section{点的存在性判据}
\label{sites-section-criterion-points}

\noindent
本节对应于 Deligne 为 \cite[Expos\'e VI]{SGA4} 所写的附录。事实上，
两者几乎逐字相同。

\medskip\noindent
设 $\mathcal{C}$ 为位点。假设 $(I, \geq)$ 是有向集，并且
$(U_i, f_{ii'})$ 是 $I$ 上的逆系统，见 Categories, Definition
\ref{categories-definition-system-over-poset}。给定数据
$(I, \geq, U_i, f_{ii'})$，定义
$$
u : \mathcal{C} \longrightarrow \textit{Sets}, \quad
u(V) = \colim_i \Mor_\mathcal{C}(U_i , V)
$$
令 $\mathcal{F} \mapsto \mathcal{F}_p$ 为第
\ref{sites-section-points} 节中与 $u$ 关联的茎函子。由定义直接可见，在这种特殊情形下，
实际上有
$$
\mathcal{F}_p = \colim_i \mathcal{F}(U_i)
$$
注意，$u$ 与所有有限极限（指 $\mathcal{C}$ 中可表示的那些）交换，因为每个函子
$V \mapsto \Mor_\mathcal{C}(U_i , V)$ 都如此；见 Categories, Lemma
\ref{categories-lemma-directed-commutes}。
% P01-ERRATA: P01-E0138 -- the source uses plural "do" after singular "each of the functors"; "does" is required.

\medskip\noindent
若 $J \subset I$，$J$ 上的序由 $I$ 上的序诱导，并且
$V_j = U_j$、$g_{jj'} = f_{jj'}$（换言之，$J$ 上的逆系统由 $I$ 上的逆系统
诱导），则称系统 $(I, \geq, U_i, f_{ii'})$ 是
$(J, \geq, V_j, g_{jj'})$ 的一个{it 加细}。
% P01-ERRATA: P01-E0139 -- the source omits "is" in "the ordering on J induced from that of I".
令 $u$ 为与 $(I, \geq, U_i, f_{ii'})$ 关联的函子，$u'$ 为与
$(J, \geq, V_j, g_{jj'})$ 关联的函子。这诱导函子间的变换
$$
u' \longrightarrow u
$$
原因仅仅是 $u'$ 的余极限取遍 $u$ 的余极限中各系统的一个子系统。特别地，
得到关联的茎函子变换 $\mathcal{F}_{p'} \to \mathcal{F}_p$，见引理
\ref{sites-lemma-maps-u-points}。

\begin{lemma}
\label{sites-lemma-refine}
设 $\mathcal{C}$ 为位点。设 $(J, \geq, V_j, g_{jj'})$ 是上述系统，关联的函子对
为 $(u', p')$。设 $\mathcal{F}$ 是 $\mathcal{C}$ 上的层，
$s, s' \in \mathcal{F}_{p'}$ 是不同元素。设
$\{W_k \to W\}$ 是 $\mathcal{C}$ 的有限覆盖，且 $f \in u'(W)$。
则存在 $(J, \geq, V_j, g_{jj'})$ 的一个加细
$(I, \geq, U_i, f_{ii'})$，使 $s, s'$ 映到
$\mathcal{F}_p$ 的不同元素，而且 $f$ 在 $u(W)$ 中的像落在某个
$u(W_k)$ 的像中。
\end{lemma}

\begin{proof}
存在 $j_0 \in J$，使 $f$ 由 $f' : V_{j_0} \to W$ 定义。对
$j \geq j_0$，令
$V_{j, k} = V_j \times_{f'\circ f_{j j_0}, W} W_k$。
% P01-ERRATA: P01-E0140 -- the transition map in the J-system is g_{j j_0}; f_{j j_0} is undefined here.
则 $\{V_{j, k} \to V_j\}$ 是位点 $\mathcal{C}$ 中的有限覆盖。因此
$\mathcal{F}(V_j) \subset \prod_k \mathcal{F}(V_{j, k})$。再次由
Categories, Lemma \ref{categories-lemma-directed-commutes}，看到
$$
\mathcal{F}_{p'} =
\colim_j \mathcal{F}(V_j)
\longrightarrow
\prod\nolimits_k \colim_j \mathcal{F}(V_{j, k})
$$
是单射。因此存在 $k$，使 $s$ 与 $s'$ 在
$\colim_j \mathcal{F}(V_{j, k})$ 中的像不同。令
$J_0 = \{j \in J, j \geq j_0\}$ 且 $I = J \amalg J_0$。在 $I$ 上规定序，
使第二个并项的任何元素都不小于第一个并项的任何元素，而其余次序使用 $J$ 上的序。
若 $j \in I$ 位于第一个并项，就使用 $V_j$；若 $j \in I$ 位于第二个并项，
就使用 $V_{j, k}$。略去该逆系统的转移映射的定义。由上述结果，$s, s'$ 在
$\mathcal{F}_p$ 中的像不同。此外，由构造，$f'$ 在 $V_{j, k}$ 上的限制分解经过
$W_k$。
\end{proof}

\begin{lemma}
\label{sites-lemma-refine-all-at-once}
设 $\mathcal{C}$ 为位点。设 $(J, \geq, V_j, g_{jj'})$ 是上述系统，关联的函子对
为 $(u', p')$。设 $\mathcal{F}$ 是 $\mathcal{C}$ 上的层，
$s, s' \in \mathcal{F}_{p'}$ 是不同元素。则存在
$(J, \geq, V_j, g_{jj'})$ 的一个加细
$(I, \geq, U_i, f_{ii'})$，使 $s, s'$ 映到
$\mathcal{F}_p$ 的不同元素，并且对位点 $\mathcal{C}$ 的每个有限覆盖
$\{W_k \to W\}$ 以及任意 $f \in u'(W)$，$f$ 在 $u(W)$ 中的像都落在某个
$u(W_k)$ 的像中。
\end{lemma}

\begin{proof}
令 $E$ 为序对 $(\{W_k \to W\}, f\in u'(W))$ 所成的集合。考虑序对
$(E' \subset E, (I, \geq, U_i, f_{ii'}))$，满足
\begin{enumerate}
\item $(I, \geq, U_i, g_{ii'})$ 是
$(J, \geq, V_j, g_{jj'})$ 的加细；
% P01-ERRATA: P01-E0141 -- the tuple is defined with transition maps f_{ii'}, but item (1) switches to undefined g_{ii'}.
\item $s, s'$ 映到 $\mathcal{F}_p$ 的不同元素；且
\item 对每个序对 $(\{W_k \to W\}, f\in u'(W)) \in E'$，$f$ 在
$u(W)$ 中的像都落在某个 $u(W_k)$ 的像中。
\end{enumerate}
在第一因子上按包含关系、第二因子上按加细关系对这些序对排序。以
$\mathcal{S}$ 表示所有上述序对
$(E' \subset E, (I, \geq, U_i, f_{ii'}))$ 所成的类。
% P01-ERRATA: P01-E0142 -- the source's "Denote S the class" omits "by" or requires recasting.
我们断言 $\mathcal{S}$ 满足 Zorn 引理的假设。事实上，设
$(E'_a, (I_a, \geq, U_i, f_{ii'}))_{a \in A}$ 是 $\mathcal{S}$ 的全序子集。
则可以定义 $E' = \bigcup_{a \in A} E'_a$，并令
$I = \bigcup_{a \in A} I_a$。我们断言相应序对
$(E' , (I, \geq, U_i, f_{ii'}))$ 是 $\mathcal{S}$ 的元素。条件 (1) 与 (3)
是清楚的。对条件 (2)，注意
$$
u = \colim_{a \in A} u_a
\text{ 并相应地 }
\mathcal{F}_p = \colim_{a \in A} \mathcal{F}_{p_a}
$$
$s, s'$ 的像在该茎中互异，由集合的有向余极限的描述得出；见
Categories, Section \ref{categories-section-directed-colimits}。在这种情形下，
我们将简写为
$(E', (I, \ldots)) = \bigcup_{a \in A}(E'_a, (I_a, \ldots))$。

\medskip\noindent
好；若 $\mathcal{S}$ 是集合，则 Zorn 引理适用，并与上面的引理
\ref{sites-lemma-refine} 结合而容易证明本引理。但它并不适用，因为
$\mathcal{S}$ 是类。为绕开这一点，在 $E$ 上选取一个良序。对 $e \in E$，令
$E'_e = \{e' \in E \mid e' \leq e\}$。用超限递归构造序对
$(E'_e, (I_e, \ldots)) \in \mathcal{S}$，使得
$e_1 \leq e_2 \Rightarrow (E'_{e_1}, (I_{e_1}, \ldots))
\leq (E'_{e_2}, (I_{e_2}, \ldots))$。设 $e \in E$，比如
$e = (\{W_k \to W\}, f\in u'(W))$。若 $e$ 有前驱 $e - 1$，就按引理
\ref{sites-lemma-refine} 对系统
$e = (\{W_k \to W\}, f\in u'(W))$ 取
$(I_{e - 1}, \ldots)$ 的一个加细 $(I_e, \ldots)$。若 $e$ 没有前驱，就对系统
$e = (\{W_k \to W\}, f\in u'(W))$ 取
$\bigcup_{e' < e} (I_{e'}, \ldots)$ 的一个加细 $(I_e, \ldots)$。最后，
并 $\bigcup_{e \in E} I_e$ 就是本引理所提出问题的一个解。
\end{proof}

\begin{proposition}
\label{sites-proposition-criterion-points}
\begin{reference}
\cite[Expos\'e VI, Appendix by Deligne, Proposition 9.0]{SGA4}
\end{reference}
设 $\mathcal{C}$ 为位点。假设
\begin{enumerate}
\item $\mathcal{C}$ 中存在有限极限；且
\item 每个覆盖 $\{U_i \to U\}_{i \in I}$ 都有一个由
$\mathcal{C}$ 的有限覆盖给出的加细。
\end{enumerate}
则 $\mathcal{C}$ 有足够多的点。
\end{proposition}

\begin{proof}
需要证明：给定 $\mathcal{C}$ 上任意层 $\mathcal{F}$、任意
$U \in \Ob(\mathcal{C})$ 以及任意两个不同截面
$s, s' \in \mathcal{F}(U)$，存在一点 $p$，使 $s, s'$ 在
$\mathcal{F}_p$ 中的像不同。见引理
\ref{sites-lemma-enough}。考虑系统 $(J, \geq, V_j, g_{jj'})$，其中
$J = \{1\}$、$V_1 = U$、$g_{11} = \text{id}_U$。应用引理
\ref{sites-lemma-refine-all-at-once}。由该引理的结果，得到加细原系统的系统
$(I, \geq, U_i, f_{ii'})$，使
$s_p \not = s'_p$，而且对位点 $\mathcal{C}$ 的每个有限覆盖
$\{W_k \to W\}$，映射 $\coprod_k u(W_k) \to u(W)$ 都是满射。由于
$\mathcal{C}$ 的每个覆盖都可以由有限覆盖加细，得出对位点
$\mathcal{C}$ 的{\it 任意}覆盖 $\{W_k \to W\}$，
$\coprod_k u(W_k) \to u(W)$ 都是满射。这蕴含 $u = p$ 是点，见命题
\ref{sites-proposition-point-limits}（以及本节开头保证 $u$ 与有限极限交换的讨论）。
\end{proof}

\begin{lemma}
\label{sites-lemma-criterion-points}
设 $\mathcal{C}$ 为位点。设 $I$ 是集合，并且对每个 $i \in I$，设
$U_i$ 是 $\mathcal{C}$ 的对象，满足
\begin{enumerate}
\item $\coprod h_{U_i}$ 满射到 $\Sh(\mathcal{C})$ 的终对象；且
% P01-ERRATA: P01-E0143 -- h_{U_i} is generally only a presheaf; the object of Sh(C) should be h_{U_i}^# as in lemma-enough-points-local.
\item $\mathcal{C}/U_i$ 满足命题
\ref{sites-proposition-criterion-points} 的假设。
\end{enumerate}
则 $\mathcal{C}$ 有足够多的点。
\end{lemma}

\begin{proof}
由假设 (2) 与上述命题，$\mathcal{C}/U_i$ 有足够多的点。通过引理
\ref{sites-lemma-point-morphism-sites} 的过程，$\mathcal{C}/U_i$ 的点给出
$\mathcal{C}$ 的点。因此只须证明：若
$\phi : \mathcal{F} \to \mathcal{G}$ 是 $\mathcal{C}$ 上的层映射，且对所有
$i$，$\phi|_{\mathcal{C}/U_i}$ 都是同构，则 $\phi$ 是同构。由假设 (1)，
对 $\mathcal{C}$ 的每个对象 $W$，存在覆盖
$\{W_j \to W\}_{j \in J}$，使对每个 $j \in J$，存在 $i \in I$ 与态射
$f_j : W_j \to U_i$。于是映射
$\mathcal{F}(W_j) \to \mathcal{G}(W_j)$ 都是双射；类似地，
$\mathcal{F}(W_j \times_W W_{j'}) \to \mathcal{G}(W_j \times_W W_{j'})$
也都是双射。层条件告诉我们，$\mathcal{F}(W) \to \mathcal{G}(W)$ 是双射，
如所需。
\end{proof}



\section{弱可缩对象}
\label{sites-section-w-contractible}

\noindent
位点的一个{\it 弱可缩}对象，是指满足下述引理中各等价条件的对象。

\begin{lemma}
\label{sites-lemma-w-contractible}
设 $\mathcal{C}$ 为位点，$U$ 是 $\mathcal{C}$ 的对象。下列条件等价：
\begin{enumerate}
\item 对每个覆盖 $\{U_i \to U\}$，存在层映射
$h_U^\# \to \coprod h_{U_i}^\#$，它是
$\coprod h_{U_i} \to h_U$ 的层化的右逆。
\item 对每个集合层的满射 $\mathcal{F} \to \mathcal{G}$，映射
$\mathcal{F}(U) \to \mathcal{G}(U)$ 都是满射。
\end{enumerate}
\end{lemma}

\begin{proof}
假设 (1)，并设 $\mathcal{F} \to \mathcal{G}$ 是集合层的满射。对
$s \in \mathcal{G}(U)$，存在覆盖 $\{U_i \to U\}$ 以及
$t_i \in \mathcal{F}(U_i)$，其像为 $s|_{U_i}$；见定义
\ref{sites-definition-sheaves-injective-surjective}。通过
(\ref{sites-equation-map-representable-into-sheaf})，把 $t_i$ 看成映射
$t_i : h_{U_i}^\# \to \mathcal{F}$。将
$\coprod t_i : \coprod h_{U_i}^\# \to \mathcal{F}$ 与 (1) 给出的映射
$h_U^\# \to \coprod h_{U_i}^\#$ 复合，就得到映到 $s$ 的截面
$t \in \mathcal{F}(U)$。因此 (2) 成立。
% P01-ERRATA: P01-E0144 -- the source duplicates the predicate in "we get from (1) we obtain".

\medskip\noindent
假设 (2) 成立。设 $\{U_i \to U\}$ 为覆盖。则
$\coprod h_{U_i}^\# \to h_U^\#$ 是满射（引理
\ref{sites-lemma-covering-surjective-after-sheafification}）。因此由 (2)，存在
$\coprod h_{U_i}^\#$ 的截面 $s$，映到 $h_U^\#$ 的截面
$\text{id}_U$。该截面对应于映射
$h_U^\# \to \coprod h_{U_i}^\#$，它是
$\coprod h_{U_i} \to h_U$ 的层化的右逆，从而证明 (1)。
\end{proof}

\begin{definition}
\label{sites-definition-w-contractible}
设 $\mathcal{C}$ 为位点。
\begin{enumerate}
\item 若引理 \ref{sites-lemma-w-contractible} 的等价条件成立，则称
$\mathcal{C}$ 的对象 $U$ 是{\it 弱可缩的}。
\item 若 $\mathcal{C}$ 的每个对象 $U$ 都有一个覆盖
$\{U_i \to U\}$，且对所有 $i$，$U_i$ 都是弱可缩的，则称该位点有
{\it 足够多的弱可缩对象}。
\item 更一般地，若 $P$ 是 $\mathcal{C}$ 的对象的一种性质，并且
$\mathcal{C}$ 的每个对象 $U$ 都有覆盖 $\{U_i \to U\}$，使所有
对所有 $i$，$U_i$ 都具有 $P$，则称 $\mathcal{C}$ 有{\it 足够多的 $P$ 对象}。
\end{enumerate}
\end{definition}

\noindent
$\mathbf{A}^1_\mathbf{C}$ 的小 \'etale 位点没有任何弱可缩对象。另一方面，
任意概形的小 pro-\'etale 位点都有足够多的可缩对象。







\section{正像的正合性}
\label{sites-section-pushforward}

\noindent
设 $f$ 为拓扑斯态射。一般而言，函子 $f_*$ 仅左正合。可以对此函子附加许多条件，
以挑出特定类别的拓扑斯态射。我们在此汇集这些条件，并指出其中若干逻辑依赖关系。
下述引理有些部分纯属范畴论（也就是说，它们不依赖于拓扑斯态射；有一对伴随函子
就足够了）。

\begin{lemma}
\label{sites-lemma-exactness-properties}
设 $f : \Sh(\mathcal{C}) \to \Sh(\mathcal{D})$ 是拓扑斯态射。考虑下列性质
（针对集合层）：
\begin{enumerate}
\item $f_*$ 是忠实函子；
\item $f_*$ 是全忠实函子；
\item 对 $\Sh(\mathcal{C})$ 中的所有 $\mathcal{F}$，
$f^{-1}f_*\mathcal{F} \to \mathcal{F}$ 都是满射；
\item $f_*$ 将满射映为满射；
\item $f_*$ 与余等化子交换；
\item $f_*$ 与推出交换；
\item 对 $\Sh(\mathcal{C})$ 中的所有 $\mathcal{F}$，
$f^{-1}f_*\mathcal{F} \to \mathcal{F}$ 都是同构；
\item $f_*$ 反映单射；
\item $f_*$ 反映满射；
\item $f_*$ 反映双射；且
\item 对任意满射 $\mathcal{F} \to f^{-1}\mathcal{G}$，存在满射
$\mathcal{G}' \to \mathcal{G}$，使
$f^{-1}\mathcal{G}' \to f^{-1}\mathcal{G}$ 通过
$\mathcal{F} \to f^{-1}\mathcal{G}$ 分解。
\end{enumerate}
则有下列蕴含关系：
\begin{enumerate}
\item[(a)] (2) $\Rightarrow$ (1),
\item[(b)] (3) $\Rightarrow$ (1),
\item[(c)] (7) $\Rightarrow$ (1), (2), (3), (8), (9), (10).
\item[(d)] (3) $\Leftrightarrow$ (9),
\item[(e)] (6) $\Rightarrow$ (4) and (5) $\Rightarrow$ (4),
\item[(f)] (4) $\Leftrightarrow$ (11),
\item[(g)] (9) $\Rightarrow$ (8), (10), and
\item[(h)] (2) $\Leftrightarrow$ (7).
\end{enumerate}
图示如下：
$$
\xymatrix{
(6) \ar@{=>}[rd] & & & & & (9) \ar@{=>}[r] \ar@{=>}[rd] & (8) \\
& (4) \ar@{<=>}[r] & (11) &
(2) \ar@{<=>}[r] &
(7) \ar@{=>}[ru] \ar@{=>}[rd] & & (10) \\
(5) \ar@{=>}[ur] & & & & & (3) \ar@{=>}[r] & (1)
}
$$
\end{lemma}

\begin{proof}
(a) 的证明：由定义立即得出。

\medskip\noindent
(b) 的证明。设 $a, b : \mathcal{F} \to \mathcal{F}'$ 是
$\mathcal{C}$ 上的层映射。若 $f_*a = f_*b$，则
$f^{-1}f_*a = f^{-1}f_*b$。考虑交换图
$$
\xymatrix{
\mathcal{F} \ar@<-1ex>[r] \ar@<1ex>[r] & \mathcal{F}' \\
f^{-1}f_*\mathcal{F} \ar@<-1ex>[r] \ar@<1ex>[r] \ar[u] &
f^{-1}f_*\mathcal{F}' \ar[u]
}
$$
若下面两条箭头相等而竖直箭头为满射，则上面两条箭头相等。因此 (b) 得证。

\medskip\noindent
(c) 的证明。设 $a : \mathcal{F} \to \mathcal{F}'$ 是
$\mathcal{C}$ 上的层映射。考虑交换图
$$
\xymatrix{
\mathcal{F} \ar[r] & \mathcal{F}' \\
f^{-1}f_*\mathcal{F} \ar[r] \ar[u] &
f^{-1}f_*\mathcal{F}' \ar[u]
}
$$
若 (7) 成立，则竖直箭头均为同构。因此，若 $f_*a$ 为单射（相应地，满射或双射），
则下面的箭头为单射（相应地，满射或双射），从而上面的箭头也为单射（相应地，满射或
双射）。于是 (7) 蕴含 (8)、(9)、(10)。显然 (7) 蕴含 (3)。蕴含关系
(7) $\Rightarrow$ (2), (1) 由下面将证明的 (a) 和 (h) 得出。

\medskip\noindent
(d) 的证明。假设 (3)。设 $a : \mathcal{F} \to \mathcal{F}'$ 是
$\mathcal{C}$ 上的层映射，并且 $f_*a$ 是满射。由于 $f^{-1}$ 正合，这蕴含
$f^{-1}f_*a : f^{-1}f_*\mathcal{F} \to f^{-1}f_*\mathcal{F}'$
为满射。结合 (3)，得出 $a$ 是满射。这就是说 (9) 成立。
假设 (9)。设 $\mathcal{F}$ 是 $\mathcal{C}$ 上的层。需要证明映射
$f^{-1}f_*\mathcal{F} \to \mathcal{F}$ 是满射。只须证明
$f_*f^{-1}f_*\mathcal{F} \to f_*\mathcal{F}$ 是满射。而这是成立的，因为存在典范映射
$f_*\mathcal{F} \to f_*f^{-1}f_*\mathcal{F}$，它是该映射的一侧逆。

\medskip\noindent
(e) 的证明。以下不再说明地使用 Categories, Lemma
\ref{categories-lemma-characterize-mono-epi}。若
$\mathcal{F} \to \mathcal{F}'$ 是满射，则
$\mathcal{F}' \amalg_\mathcal{F} \mathcal{F}' \to \mathcal{F}'$
是同构。因此 (6) 蕴含
$$
f_*\mathcal{F}' \amalg_{f_*\mathcal{F}} f_*\mathcal{F}' =
f_*(\mathcal{F}' \amalg_\mathcal{F} \mathcal{F}')
\longrightarrow
f_*\mathcal{F}'
$$
也是同构。进而这蕴含 $f_*\mathcal{F} \to f_*\mathcal{F}'$ 是满射。因此 (6)
蕴含 (4)。若 $\mathcal{F} \to \mathcal{F}'$ 是满射，则由引理
\ref{sites-lemma-coequalizer-surjection}，$\mathcal{F}'$ 是两个投影
$\mathcal{F} \times_{\mathcal{F}'} \mathcal{F} \to \mathcal{F}$
的余等化子。因此若 (5) 成立，则 $f_*\mathcal{F}'$ 是两个投影
$$
f_*(\mathcal{F} \times_{\mathcal{F}'} \mathcal{F}) =
f_*\mathcal{F} \times_{f_*\mathcal{F}'} f_*\mathcal{F}
\longrightarrow
f_*\mathcal{F}
$$
的余等化子，而这显然意味着 $f_*\mathcal{F} \to f_*\mathcal{F}'$ 是满射。因此 (5)
也蕴含 (4)。

\medskip\noindent
(f) 的证明。假设 (4)。设 $\mathcal{F} \to f^{-1}\mathcal{G}$ 是
$\mathcal{C}$ 上的层的满射。由 (4)，
$f_*\mathcal{F} \to f_*f^{-1}\mathcal{G}$ 是满射。令
$\mathcal{G}'$ 为纤维积
$$
\xymatrix{
f_*\mathcal{F} \ar[r] & f_*f^{-1}\mathcal{G} \\
\mathcal{G}' \ar[u] \ar[r] & \mathcal{G} \ar[u]
}
$$
于是 $\mathcal{G}' \to \mathcal{G}$ 也是满射。考虑交换图
$$
\xymatrix{
\mathcal{F} \ar[r] & f^{-1}\mathcal{G} \\
f^{-1}f_*\mathcal{F} \ar[r] \ar[u] & f^{-1}f_*f^{-1}\mathcal{G} \ar[u] \\
f^{-1}\mathcal{G}' \ar[u] \ar[r] & f^{-1}\mathcal{G} \ar[u]
}
$$
由此得到所需结论。反过来，假设 (11)。设
$a : \mathcal{F} \to \mathcal{F}'$ 是 $\mathcal{C}$ 上的层的满射。
% P01-ERRATA: P01-E0145 -- the source omits the article in "be surjective map of sheaves".
考虑纤维积图
$$
\xymatrix{
\mathcal{F} \ar[r] & \mathcal{F}' \\
\mathcal{F}'' \ar[u] \ar[r] & f^{-1}f_*\mathcal{F}' \ar[u]
}
$$
由于下面的水平箭头是满射，由 (11) 可找到满射
$\gamma : \mathcal{G}' \to f_*\mathcal{F}'$，使 $f^{-1}\gamma$ 通过
$\mathcal{F}'' \to f^{-1}f_*\mathcal{F}'$ 分解：
$$
\xymatrix{
& \mathcal{F} \ar[r] & \mathcal{F}' \\
f^{-1}\mathcal{G}' \ar[r] & \mathcal{F}'' \ar[u] \ar[r] &
f^{-1}f_*\mathcal{F}' \ar[u]
}
$$
用 $f_*$ 将其向下推，得到交换图
$$
\xymatrix{
& f_*\mathcal{F} \ar[r] & f_*\mathcal{F}' \\
f_*f^{-1}\mathcal{G}' \ar[r] & f_*\mathcal{F}'' \ar[u] \ar[r] &
f_*f^{-1}f_*\mathcal{F}' \ar[u] \\
\mathcal{G}' \ar[u] \ar[rr] & & f_*\mathcal{F}' \ar[u]
}
$$
这证明了 (4)。

\medskip\noindent
(g) 的证明。假设 (9)。以下不再说明地使用 Categories, Lemma
\ref{categories-lemma-characterize-mono-epi}。设
$a : \mathcal{F} \to \mathcal{F}'$ 是 $\mathcal{C}$ 上的层映射，且
$f_*a$ 是单射。这意味着
$f_*\mathcal{F} \to f_*\mathcal{F} \times_{f_*\mathcal{F}'} f_*\mathcal{F} =
f_*(\mathcal{F} \times_{\mathcal{F}'} \mathcal{F})$ 是同构。于是由 (9)，
$\mathcal{F} \to \mathcal{F} \times_{\mathcal{F}'} \mathcal{F}$ 是满射，
即为同构。因此 $a$ 是单射，也就是说 (8) 成立。由于 (10) 显然等价于
(8) $+$ (9)，(g) 得证。

\medskip\noindent
(h) 的证明。这就是 Categories, Lemma
\ref{categories-lemma-adjoint-fully-faithful}。
\end{proof}

\noindent
下面给出位点态射的一个条件，它保证函子 $f_*$ 将满射映为满射。

\begin{lemma}
\label{sites-lemma-weaker}
设 $f : \mathcal{D} \to \mathcal{C}$ 是由连续函子
$u : \mathcal{C} \to \mathcal{D}$ 所关联的位点态射。假设对
$\mathcal{C}$ 的任意对象 $U$ 以及 $\mathcal{D}$ 中的任意覆盖
$\{V_j \to u(U)\}$，都存在 $\mathcal{C}$ 中的覆盖 $\{U_i \to U\}$，使层映射
$$
\coprod h_{u(U_i)}^\# \to h_{u(U)}^\#
$$
通过层映射
$$
\coprod h_{V_j}^\# \to h_{u(U)}^\#.
$$
分解。则 $f_*$ 将层的满射映为层的满射。
\end{lemma}

\begin{proof}
设 $a : \mathcal{F} \to \mathcal{G}$ 是 $\mathcal{D}$ 上层的满射。设 $U$ 是
$\mathcal{C}$ 的对象，并设
$s \in f_*\mathcal{G}(U) = \mathcal{G}(u(U))$。由假设，存在覆盖
$\{V_j \to u(U)\}$ 以及截面 $s_j \in \mathcal{F}(V_j)$，使
$a(s_j) = s|_{V_j}$。现在可以把截面 $s$、$s_j$ 与 $a$ 看成给出下列层映射交换图：
$$
\xymatrix{
\coprod h_{V_j}^\# \ar[r]_-{\coprod s_j} \ar[d] & \mathcal{F} \ar[d]^a \\
h_{u(U)}^\# \ar[r]^s & \mathcal{G}
}
$$
由假设，存在覆盖 $\{U_i \to U\}$，使上面的交换图可以扩充如下：
$$
\xymatrix{
& \coprod h_{V_j}^\# \ar[r]_-{\coprod s_j} \ar[d] & \mathcal{F} \ar[d]^a \\
\coprod h_{u(U_i)}^\# \ar[r] \ar[ur] &
h_{u(U)}^\# \ar[r]^s & \mathcal{G}
}
$$
因为 $\mathcal{F}$ 是层，从左下角到右上角的映射对应于一族截面
$s_i \in \mathcal{F}(u(U_i))$，即截面 $s_i \in f_*\mathcal{F}(U_i)$。
图的交换性蕴含 $a(s_i)$ 等于 $s$ 在 $U_i$ 上的限制。换言之，我们已经证明
$f_*a$ 是层的满射。
\end{proof}

\begin{example}
\label{sites-example-weaker-not-coequalizer}
假设 $f : \mathcal{D} \to \mathcal{C}$ 满足引理 \ref{sites-lemma-weaker}
的假设。一般而言，$f_*$ 并不与余等化子或推出交换。事实上，设 $f$ 是拓扑空间态射
$X = \{1, 2\} \to Y = \{*\}$ 所关联的位点态射（见例
\ref{sites-example-continuous-map}），其中 $Y$ 是单点空间，而
$X = \{1, 2\}$ 是具有两个点的离散空间。$X$ 上的层 $\mathcal{F}$ 由一对集合
$(A_1, A_2)$ 给出。于是 $f_*\mathcal{F}$ 对应于集合 $A_1 \times A_2$。
因此，若
$a = (a_1, a_2), b = (b_1, b_2) : (A_1, A_2) \to (B_1, B_2)$
是 $X$ 上的层映射，则 $a, b$ 的余等化子为 $(C_1, C_2)$，其中 $C_i$ 是
$a_i, b_i$ 的余等化子；而 $f_*a, f_*b$ 的余等化子是
$$
a_1 \times a_2, b_1 \times b_2 :
A_1 \times A_2 \longrightarrow B_1 \times B_2
$$
的余等化子，一般不同于 $C_1 \times C_2$。事实上，若
$A_2 = \emptyset$，则 $A_1 \times A_2 = \emptyset$，故所示两箭头的余等化子为
$B_1 \times B_2$；但一般而言 $C_1 \not = B_1$。对推出可作类似例子。
\end{example}

\noindent
下述引理给出一个判据，用以判断何时位点态射的函子 $f_*$ 反映单射与满射。
注意，这也蕴含 $f_*$ 是忠实函子，并且映射
$f^{-1}f_*\mathcal{F} \to \mathcal{F}$ 总是满射。

\begin{lemma}
\label{sites-lemma-cover-from-below}
设 $f : \mathcal{D} \to \mathcal{C}$ 是由函子
$u : \mathcal{C} \to \mathcal{D}$ 给出的位点态射。假设对
$\mathcal{D}$ 的每个对象 $V$，都存在 $\mathcal{C}$ 的对象 $U_i$ 以及态射
$u(U_i) \to V$，使 $\{u(U_i) \to V\}$ 是 $\mathcal{D}$ 的覆盖。在这种情形下，
函子 $f_* : \Sh(\mathcal{D}) \to \Sh(\mathcal{C})$ 反映单射与满射。
\end{lemma}

\begin{proof}
设 $\alpha : \mathcal{F} \to \mathcal{G}$ 是 $\mathcal{D}$ 上的层映射。
% P01-ERRATA: P01-E0146 -- the source says "be maps of sheaves" for the single map alpha.
由假设，对 $\mathcal{D}$ 的每个对象 $V$，层条件给出
$\mathcal{F}(V) \subset \prod \mathcal{F}(u(U_i)) =
\prod f_*\mathcal{F}(U_i)$，其中 $U_i \in \Ob(\mathcal{C})$ 适当；对
$\mathcal{G}$ 亦然。因此，若 $f_*\alpha$ 是单射，则显然 $\alpha$ 是单射。
换言之，$f_*$ 反映单射。

\medskip\noindent
假设 $f_*\alpha$ 是满射。对上述 $V, U_i, u(U_i) \to V$ 以及截面
$s \in \mathcal{G}(V)$，存在覆盖 $\{U_{ij} \to U_i\}$，使
$s|_{u(U_{ij})}$ 落在 $\mathcal{F}(u(U_{ij}))$ 的像中。由于
$\{u(U_{ij}) \to V\}$ 是覆盖（因为 $u$ 连续，并由位点公理），我们得出
$s$ 局部落在像中。因此 $\alpha$ 是满射。换言之，$f_*$ 反映满射。
\end{proof}

\begin{example}
\label{sites-example-cover-from-below}
我们构造一个满足引理 \ref{sites-lemma-cover-from-below} 假设的态射
$f : \mathcal{D} \to \mathcal{C}$。具体地，设
$\varphi : G \to H$ 是有限群的态射。考虑可数 $G$-集合与 $H$-集合的位点
$\mathcal{D} = \mathcal{T}_G$ 与 $\mathcal{C} = \mathcal{T}_H$，其覆盖为联合满射
的可数态射族（例 \ref{sites-example-site-on-group}）。设
$u : \mathcal{T}_H \to \mathcal{T}_G$ 是节
\ref{sites-section-G-sets-morphisms} 中描述的函子，而
$f : \mathcal{T}_G \to \mathcal{T}_H$ 是相应的位点态射。若 $\varphi$ 是单射，
则每个可数 $G$-集合，作为 $G$-集合，都是某个可数 $H$-集合的商（若
$\varphi$ 不是单射，这一点不成立）。因此 $f$ 满足引理
\ref{sites-lemma-cover-from-below} 的假设。
若 $\mathcal{T}_G$ 上的层 $\mathcal{F}$ 对应于 $G$-集合 $S$，则典范映射
$$
f^{-1}f_*\mathcal{F} \longrightarrow \mathcal{F}
$$
对应于映射
$$
\text{Map}_G(H, S) \longrightarrow S,\quad
a \longmapsto a(1_H)
$$
若 $\varphi$ 单射但不满射，则该映射是满射（正如引理
\ref{sites-lemma-cover-from-below} 所应有的结论），但一般不是单射（例如取
$G = \{1\}$、$H = \{1, \sigma\}$ 以及 $S = \{1, 2\}$）。此外，函子
$f_*$ 不与余等化子或推出交换（取 $G = \{1\}$ 且
$H = \{1, \sigma\}$）。
\end{example}

\section{几乎余连续的函子}
\label{sites-section-almost-cocontinuous}

\noindent
设 $\mathcal{C}$ 为位点。范畴 $\textit{PSh}(\mathcal{C})$ 有始对象，即把空集
赋给 $\mathcal{C}$ 的每个对象的预层。将该预层记为 $\emptyset$。由层化的性质，
$\emptyset$ 的层化 $\emptyset^\#$ 是 $\mathcal{C}$ 上层的范畴
$\Sh(\mathcal{C})$ 的始对象。

\begin{definition}
\label{sites-definition-empty}
设 $\mathcal{C}$ 为位点。若 $\emptyset^\# \to h_U^\#$ 是层的同构，则称
$\mathcal{C}$ 的对象 $U$ 在{\it 层论意义下为空}。
\end{definition}

\noindent
下述引理把这个概念表述得更明确。

\begin{lemma}
\label{sites-lemma-characterize-empty}
设 $\mathcal{C}$ 为位点，$U$ 是 $\mathcal{C}$ 的对象。下列条件等价：
\begin{enumerate}
\item $U$ 在层论意义下为空；
\item 对每个层 $\mathcal{F}$，$\mathcal{F}(U)$ 都是单点集；
\item $\emptyset^\#(U)$ 是单点集；
\item $\emptyset^\#(U)$ 非空；且
\item 空族是 $\mathcal{C}$ 中 $U$ 的一个覆盖。
\end{enumerate}
此外，若 $U$ 在层论意义下为空，则对 $\mathcal{C}$ 的任意态射
$U' \to U$，对象 $U'$ 也在层论意义下为空。
\end{lemma}

\begin{proof}
对任意层 $\mathcal{F}$，有
$\mathcal{F}(U) = \Mor_{\Sh(\mathcal{C})}(h_U^\#, \mathcal{F})$。
因此 (1) 与 (2) 等价。显然 (2) 蕴含 (3)，而 (3) 蕴含 (4)。
若 $U$ 的每个覆盖都由非空族给出，则由加号构造的定义，
$\emptyset^+(U)$ 为空。注意 $\emptyset^+ = \emptyset^\#$，因为
$\emptyset$ 是分离预层；见定理 \ref{sites-theorem-plus}。所以 (4) 蕴含 (5)。
若 (5) 成立，则由 $\mathcal{F}$ 的层条件，对每个层 $\mathcal{F}$，
$\mathcal{F}(U)$ 都是单点集；见注 \ref{sites-remark-sheaf-condition-empty-covering}。
因此 (5) 蕴含 (2)，从而 (1) -- (5) 等价。引理的最后一个断言由定义
\ref{sites-definition-site} 的公理 (3) 对 $U$ 的空覆盖应用而得。
% P01-ERRATA: P01-E0147 -- the source omits "to" in "applied the empty covering of U".
\end{proof}

\begin{definition}
\label{sites-definition-almost-cocontinuous}
设 $\mathcal{C}$、$\mathcal{D}$ 为位点，并设
$u : \mathcal{C} \to \mathcal{D}$ 为函子。若对 $\mathcal{C}$ 的每个对象
$U$ 以及每个覆盖 $\{V_j \to u(U)\}_{j \in J}$，都存在 $\mathcal{C}$ 中的覆盖
$\{U_i \to U\}_{i \in I}$，使对 $I$ 中每个 $i$，下列两个条件至少有一个成立：
\begin{enumerate}
\item $u(U_i)$ 在层论意义下为空；或
\item 对某个 $j \in J$，态射 $u(U_i) \to u(U)$ 通过 $V_j$ 分解，
\end{enumerate}
则称 $u$ 是{\it 几乎余连续的}。
\end{definition}

\noindent
该定义的动机来自拓扑空间的闭浸入 $i : Z \to X$。如例
\ref{sites-example-closed-map-cocontinuous-false} 所讨论的，连续函子
$X_{Zar} \to Z_{Zar}$，$U \mapsto Z \cap U$ 并不余连续，但它在上述意义下
几乎余连续。我们知道，$i_*$ 对集合层虽不正合，对阿贝尔群层却正合；见
Sheaves, Remark \ref{sheaves-remark-i-star-not-exact}。
% P01-ERRATA: P01-E0148 -- the source needs commas around the parenthetical "while not exact on sheaves of sets".
这对连续且几乎余连续的函子普遍成立。

\begin{lemma}
\label{sites-lemma-almost-cocontinuous-sheafification}
设 $\mathcal{C}$、$\mathcal{D}$ 为位点，并设
$u : \mathcal{C} \to \mathcal{D}$ 为函子。假设 $u$ 连续且几乎余连续。设
$\mathcal{G}$ 是 $\mathcal{D}$ 上的预层，并且每当 $V$ 在层论意义下为空时，
$\mathcal{G}(V)$ 都是单点集。则
$(u^p\mathcal{G})^\# = u^p(\mathcal{G}^\#)$。
\end{lemma}

\begin{proof}
设 $U \in \Ob(\mathcal{C})$。需要证明
$(u^p\mathcal{G})^\#(U) = u^p(\mathcal{G}^\#)(U)$。只须证明
$(u^p\mathcal{G})^+(U) = u^p(\mathcal{G}^+)(U)$，因为
$\mathcal{G}^+$ 仍是满足本引理假设的预层。我们有
$$
u^p(\mathcal{G}^+)(U) =
\mathcal{G}^+(u(U)) =
\colim_\mathcal{V} \check H^0(\mathcal{V}, \mathcal{G})
$$
其中余极限遍历 $\mathcal{D}$ 中 $u(U)$ 的所有覆盖 $\mathcal{V}$。另一方面，
$$
u^p(\mathcal{G})^+(U) =
\colim_\mathcal{U} \check H^0(u(\mathcal{U}), \mathcal{G})
$$
其中余极限遍历 $\mathcal{C}$ 中 $U$ 的覆盖
$\mathcal{U} = \{U_i \to U\}_{i \in I}$ 所成的范畴，并且
$u(\mathcal{U}) = \{u(U_i) \to u(U)\}_{i \in I}$。$u$ 连续这一条件意味着每个
$u(\mathcal{U})$ 都是覆盖。写成 $I = I_1 \amalg I_2$，其中
$$
I_2 = \{i \in I \mid u(U_i)\text{ 在层论意义下为空}\}
$$
则 $u(\mathcal{U})' = \{u(U_i) \to u(U)\}_{i \in I_1}$ 仍是该对象的覆盖，
% P01-ERRATA: P01-E0149 -- the source says "still a covering of because" and omits the covered object u(U).
因为其余每一块都可由空族覆盖，所以可以由定义 \ref{sites-definition-site} 的公理 (2)
删去。此外，由我们对 $\mathcal{G}$ 的假设，
$\check H^0(u(\mathcal{U}), \mathcal{G}) =
\check H^0(u(\mathcal{U})', \mathcal{G})$。最后，$u$ 几乎余连续这一条件蕴含：
对 $u(U)$ 的每个覆盖 $\mathcal{V}$，都存在 $U$ 的覆盖 $\mathcal{U}$，使
$u(\mathcal{U})'$ 加细 $\mathcal{V}$。所以，上面两个余极限具有相同的值，
如所需。
\end{proof}

\begin{lemma}
\label{sites-lemma-continuous-almost-cocontinuous}
设 $\mathcal{C}$、$\mathcal{D}$ 为位点，并设
$u : \mathcal{C} \to \mathcal{D}$ 为函子。假设 $u$ 连续且几乎余连续。则
$u^s = u^p : \Sh(\mathcal{D}) \to \Sh(\mathcal{C})$ 与推出和余等化子交换
（更一般地，与有限连通余极限交换）。
\end{lemma}

\begin{proof}
设 $\mathcal{I}$ 为有限连通指标范畴。设
$\mathcal{I} \to \Sh(\mathcal{D})$，$i \mapsto \mathcal{G}_i$ 为一个图。
% P01-ERRATA: P01-E0150 -- the source says "Let ... by a diagram" instead of "be a diagram".
我们知道，该图的余极限是预层范畴中余极限的层化；见引理
\ref{sites-lemma-colimit-sheaves}。以 $\colim^{Psh}$ 表示预层范畴中的余极限。
由于 $\mathcal{I}$ 有限且连通，预层
$\colim^{Psh}_i \mathcal{G}_i$ 满足引理
\ref{sites-lemma-almost-cocontinuous-sheafification} 的假设（因为单点集的有限连通余极限
是单点集）。所以该引理给出
\begin{align*}
u^s(\colim_i \mathcal{G}_i) & =
u^s((\colim^{Psh}_i \mathcal{G}_i)^\#) \\
& = (u^p(\colim^{Psh}_i \mathcal{G}_i))^\# \\
& = (\colim^{PSh}_i u^p(\mathcal{G}_i))^\# \\
& = \colim_i u^s(\mathcal{G}_i)
\end{align*}
如所需。
% P01-ERRATA: P01-E0151 -- the source inconsistently writes Psh and PSh for the same presheaf colimit superscript.
\end{proof}

\begin{lemma}
\label{sites-lemma-morphism-of-sites-almost-cocontinuous}
设 $f : \mathcal{D} \to \mathcal{C}$ 是与连续函子
$u : \mathcal{C} \to \mathcal{D}$ 关联的位点态射。若 $u$ 几乎余连续，则
$f_*$ 与推出和余等化子交换（更一般地，与有限连通余极限交换）。
\end{lemma}

\begin{proof}
这是引理 \ref{sites-lemma-continuous-almost-cocontinuous} 的特例。
\end{proof}

\section{子拓扑斯}
\label{sites-section-subtopoi}

\noindent
定义如下。

\begin{definition}
\label{sites-definition-embedding}
设 $\mathcal{C}$ 和 $\mathcal{D}$ 为位点。若 $f_*$ 全忠实，则称拓扑斯态射
$f : \Sh(\mathcal{D}) \to \Sh(\mathcal{C})$ 为一个{\it 嵌入}。
\end{definition}

\noindent
由引理 \ref{sites-lemma-exactness-properties}，这等价于要求：对
$\mathcal{D}$ 上的每个层 $\mathcal{F}$，伴随映射
$f^{-1}f_*\mathcal{F} \to \mathcal{F}$ 都是同构。

\begin{definition}
\label{sites-definition-subtopos}
设 $\mathcal{C}$ 为位点。若存在拓扑斯嵌入
$f : \Sh(\mathcal{D}) \to \Sh(\mathcal{C})$，使严格满子范畴
$E \subset \Sh(\mathcal{C})$ 等于函子 $f_*$ 的本质像，则称 $E$ 为一个
{\it 子拓扑斯}。
\end{definition}

\noindent
下述引理构造的子拓扑斯将在后面的定义中称为“开”的。

\begin{lemma}
\label{sites-lemma-open-subtopos}
设 $\mathcal{C}$ 为位点，$\mathcal{F}$ 是 $\mathcal{C}$ 上的层。下列条件等价：
\begin{enumerate}
\item $\mathcal{F}$ 是 $\Sh(\mathcal{C})$ 的终对象的子对象；且
\item 拓扑斯 $\Sh(\mathcal{C})/\mathcal{F}$ 是
$\Sh(\mathcal{C})$ 的子拓扑斯。
\end{enumerate}
\end{lemma}

\begin{proof}
我们已在引理 \ref{sites-lemma-localize-topos} 中看到，
$\Sh(\mathcal{C})/\mathcal{F}$ 是拓扑斯。这里回顾其证明。先应用引理
\ref{sites-lemma-topos-good-site}，可假设 $\mathcal{C}$ 是具有次典范拓扑、纤维积、
终对象 $X$ 以及满足 $\mathcal{F} = h_U$ 的对象 $U$ 的位点。引理
\ref{sites-lemma-localize-topos} 的证明表明，拓扑斯态射
$j_\mathcal{F} : \Sh(\mathcal{C})/\mathcal{F} \to \Sh(\mathcal{C})$
（在某些识别下）等于局部化态射
$j_U : \Sh(\mathcal{C}/U) \to \Sh(\mathcal{C})$。

\medskip\noindent
假设 (2)。这意味着对 $\mathcal{C}/U$ 上的所有层 $\mathcal{G}$，
$j_U^{-1}j_{U, *}\mathcal{G} \to \mathcal{G}$ 都是同构。对
$\mathcal{C}/U$ 的任意对象 $Z/U$，由引理
\ref{sites-lemma-localize-given-products}，有
$$
(j_{U, *}h_{Z/U})(U) = \Mor_{\mathcal{C}/U}(U \times_X U/U, Z/U)
$$
在上面的等式中令 $\mathcal{G} = h_{Z/U}$，便对所有 $Z/U$ 得到
$$
\Mor_{\mathcal{C}/U}(U \times_X U/U, Z/U) = \Mor_{\mathcal{C}/U}(U, Z/U)
$$
由 Yoneda 引理（Categories, Lemma \ref{categories-lemma-yoneda}），这蕴含
$U \times_X U = U$。由 Categories, Lemma
\ref{categories-lemma-characterize-mono-epi}，$U \to X$ 是单态射；换言之，
(1) 成立。

\medskip\noindent
假设 (1)。则由引理 \ref{sites-lemma-restrict-back}，
$j_U^{-1} j_{U, *} = \text{id}$。
\end{proof}

\begin{definition}
\label{sites-definition-open-subtopos}
设 $\mathcal{C}$ 为位点。若存在 $\Sh(\mathcal{C})$ 的终对象的子层
$\mathcal{F}$，使严格满子范畴 $E \subset \Sh(\mathcal{C})$ 是引理
\ref{sites-lemma-open-subtopos} 所述的子拓扑斯
$\Sh(\mathcal{C})/\mathcal{F}$，则称 $E$ 为一个{\it 开子拓扑斯}。
\end{definition}

\noindent
这意味着 $\Sh(\mathcal{C})$ 的开子拓扑斯之集，与
$\Sh(\mathcal{C})$ 的终对象的子对象之集之间存在双射。给定一个开子拓扑斯，
存在一个“闭”补。

\begin{lemma}
\label{sites-lemma-closed-subtopos}
设 $\mathcal{C}$ 为位点，$\mathcal{F}$ 是 $\Sh(\mathcal{C})$ 的终对象
$*$ 的子层。由所有满足
$\mathcal{F} \times \mathcal{G} \to \mathcal{F}$ 为同构的层
$\mathcal{G}$ 组成的满子范畴，是 $\Sh(\mathcal{C})$ 的子拓扑斯。
\end{lemma}

\begin{proof}
应用引理 \ref{sites-lemma-topos-good-site}，可假设 $\mathcal{C}$ 是具有该引理所列性质
的位点。特别地，$\mathcal{C}$ 有终对象 $X$（因而 $* = h_X$），也有满足
$\mathcal{F} = h_U$ 的对象 $U$。

\medskip\noindent
令 $\mathcal{D} = \mathcal{C}$ 作为范畴，但规定一族态射
$\{V_j \to V\}$ 是覆盖，当且仅当
$\{V_i \to V\} \cup \{U \times_X V \to V\}$ 是覆盖。
% P01-ERRATA: P01-E0152 -- the source switches the covering-family index from j to i without declaration.
由我们对 $\mathcal{C}$ 的选择，这恰好意味着
$$
h_{U \times_X V} \amalg \coprod h_{V_i} \longrightarrow h_V
$$
是满射。我们断言 $\mathcal{D}$ 是位点，即覆盖满足定义
\ref{sites-definition-site} 的条件 (1)、(2)、(3)。条件 (1) 成立。对条件 (2)，设
$\{V_i \to V\}$ 与 $\{V_{ij} \to V_i\}$ 是 $\mathcal{D}$ 的覆盖。则复合
$$
\coprod \left(
h_{U \times_X V_i} \amalg \coprod h_{V_{ij}}
\right) \longrightarrow
h_{U \times_X V} \amalg \coprod h_{V_i} \longrightarrow h_V
$$
是满射。由于每个态射 $U \times_X V_i \to V$ 都通过
$U \times_X V$ 分解，得出
$$
h_{U \times_X V} \amalg \coprod h_{V_{ij}} \longrightarrow h_V
$$
是满射，即 $\{V_{ij} \to V\}$ 是 $\mathcal{D}$ 中 $V$ 的覆盖。条件 (3)
类似地得出，因为层的满射在基变换下仍是满射。

\medskip\noindent
注意，（恒等）函子 $u : \mathcal{C} \to \mathcal{D}$ 连续，并与纤维积及终对象
交换。因此得到位点态射 $f : \mathcal{D} \to \mathcal{C}$（命题
\ref{sites-proposition-get-morphism}）。注意 $f_*$ 在底层预层上是恒等函子，因而全忠实。
为完成证明，需要说明 $f_*$ 的本质像正是本引理指定的满子范畴
$E \subset \Sh(\mathcal{C})$。为此，注意
$\mathcal{G} \in \Ob(\Sh(\mathcal{C}))$ 属于 $E$，当且仅当对
$\mathcal{C}$ 的所有对象 $V$，$\mathcal{G}(U \times_X V)$ 都是单点集。因此，
这样的层对 $\mathcal{D}$ 的所有覆盖都满足层性质（论证略）。反过来，若
$\mathcal{G}$ 对 $\mathcal{D}$ 的所有覆盖都满足层性质，则
$\mathcal{G}(U \times_X V)$ 是单点集，因为在 $\mathcal{D}$ 中，对象
$U \times_X V$ 由空覆盖覆盖。
\end{proof}

\begin{definition}
\label{sites-definition-closed-subtopos}
设 $\mathcal{C}$ 为位点。若存在 $\Sh(\mathcal{C})$ 的终对象的子层
$\mathcal{F}$，使严格满子范畴 $E \subset \Sh(\mathcal{C})$ 是引理
\ref{sites-lemma-closed-subtopos} 所述的子拓扑斯，则称 $E$ 为一个
{\it 闭子拓扑斯}。
\end{definition}

\noindent
现在可以定义拓扑斯的闭浸入与开浸入。

\begin{definition}
\label{sites-definition-immersion-topoi}
设 $f : \Sh(\mathcal{D}) \to \Sh(\mathcal{C})$ 为拓扑斯态射。
\begin{enumerate}
\item 若 $f$ 是嵌入，且 $f_*$ 的本质像是开子拓扑斯，则称 $f$ 是
{\it 开浸入}。
\item 若 $f$ 是嵌入，且 $f_*$ 的本质像是闭子拓扑斯，则称 $f$ 是
{\it 闭浸入}。
\end{enumerate}
\end{definition}

\begin{lemma}
\label{sites-lemma-closed-immersion}
设 $i : \Sh(\mathcal{D}) \to \Sh(\mathcal{C})$ 为拓扑斯的闭浸入。则
$i_*$ 全忠实，将满射映为满射，与余等化子交换，与推出交换，并反映单射、满射与双射。
\end{lemma}

\begin{proof}
设 $\mathcal{F}$ 是 $\Sh(\mathcal{C})$ 的终对象 $*$ 的子层，并设
$E \subset \Sh(\mathcal{C})$ 是由所有满足
$\mathcal{F} \times \mathcal{G} \to \mathcal{F}$ 为同构的
$\mathcal{G}$ 组成的满子范畴。由引理 \ref{sites-lemma-closed-subtopos}，函子
$i_*$ 同构于包含函子 $\iota : E \to \Sh(\mathcal{C})$。

\medskip\noindent
设 $j_{\mathcal{F}} : \Sh(\mathcal{C})/\mathcal{F} \to \Sh(\mathcal{C})$
为局部化函子（引理 \ref{sites-lemma-localize-topos}）。注意，$E$ 也可以描述为所有满足
$j_\mathcal{F}^{-1}\mathcal{G} = *$ 的层 $\mathcal{G}$ 之集。

\medskip\noindent
设 $a, b : \mathcal{G}_1 \to \mathcal{G}_2$ 是 $E$ 的两个态射。
% P01-ERRATA: P01-E0153 -- the source says "two morphism" instead of "two morphisms".
要证明 $\iota$ 与余等化子交换，只须证明 $a$、$b$ 在
$\Sh(\mathcal{C})$ 中的余等化子属于 $E$。这是清楚的，因为两个态射
$* \to *$ 的余等化子是 $*$，并且 $j_\mathcal{F}^{-1}$ 正合。推出的情形类似。

\medskip\noindent
因此 $i_*$ 满足引理 \ref{sites-lemma-exactness-properties} 的性质 (5)、(6)、(7)，
从而态射 $i$ 满足该引理中提到的所有性质，特别是本引理提到的那些性质。
\end{proof}

\section{代数结构层}
\label{sites-section-sheaves-algebraic-structures}

\noindent
在 Sheaves, Section \ref{sheaves-section-algebraic-structures} 中，我们把一种
代数结构类型定义为序对 $(\mathcal{A}, s)$，其中 $\mathcal{A}$ 是范畴，而
$s : \mathcal{A} \to \textit{Sets}$ 是满足下列条件的函子：
\begin{enumerate}
\item $s$ 忠实；
\item $\mathcal{A}$ 有极限，且 $s$ 与极限交换；
\item $\mathcal{A}$ 有滤过余极限，且 $s$ 与它们交换；且
\item $s$ 反映同构。
\end{enumerate}
对这样的代数结构类型，我们已经看到：空间 $X$ 上取值于 $\mathcal{A}$ 的预层
$\mathcal{F}$ 是层，当且仅当与它关联的集合预层是层。此外，我们阐明了茎的概念；
并且给定连续映射 $f : X \to Y$，定义了代数结构层上的伴随函子正像与逆像，
它们与底层集合层上的正像与逆像一致。再者，把基上的代数结构层扩张到空间的所有
开集，也如预期那样运作。
% P01-ERRATA: P01-E0154 -- the source's compound subject "pushforward and pullback" takes "agree", not "agrees".

\medskip\noindent
这部分材料在位点与层的框架中仍然成立。设 $(\mathcal{A}, s)$ 为一种代数结构类型，
$\mathcal{C}$ 为位点。分别以
$\textit{PSh}(\mathcal{C}, \mathcal{A})$ 与
$\Sh(\mathcal{C}, \mathcal{A})$ 表示 $\mathcal{C}$ 上取值于
$\mathcal{A}$ 的预层范畴与层范畴。
\begin{itemize}
\item[] ($\alpha$) 取值于 $\mathcal{A}$ 的预层是层，当且仅当其底层集合预层是层。
见 Sheaves, Lemma \ref{sheaves-lemma-sheaves-structure} 的证明。
\item[] ($\beta$) 给定取值于 $\mathcal{A}$ 的预层 $\mathcal{F}$，预层
${\mathcal{F}}^\# = (\mathcal{F}^+)^+$ 是层。这是因为层化过程中出现的余极限
是滤过的，甚至是遍历有向集的余极限（见节 \ref{sites-section-sheafification}，特别是
引理 \ref{sites-lemma-sheafification-exact} 的证明），并且 $s$ 与滤过余极限交换。
\item[] ($\gamma$) 得到交换图
$$
\xymatrix{
\Sh(\mathcal{C}, \mathcal{A}) \ar@<1ex>[r] \ar[d]_s &
\textit{PSh}(\mathcal{C}, \mathcal{A}) \ar@<1ex>[l]^{{}^\#} \ar[d]^s\\
\Sh(\mathcal{C}) \ar@<1ex>[r] &
\textit{PSh}(\mathcal{C}) \ar@<1ex>[l]
}
$$
\item[] ($\delta$) 当且仅当 $\mathcal{F}$ 是代数结构层时，
$\mathcal{F} = \mathcal{F}^\#$。
\item[] ($\epsilon$) 函子 ${}^\#$ 左伴随于包含函子：
$$
\Mor_{\textit{PSh}(\mathcal{C}, \mathcal{A})}(\mathcal{G}, \mathcal{F})
=
\Mor_{\Sh(\mathcal{C}, \mathcal{A})}(\mathcal{G}^\#, \mathcal{F})
$$
其证明与命题 \ref{sites-proposition-sheafification-adjoint} 的证明相同。
\item[] ($\zeta$) 函子 $\mathcal{F} \mapsto \mathcal{F}^\#$ 左正合。
其证明与引理 \ref{sites-lemma-sheafification-exact} 的证明相同。
\end{itemize}

\begin{definition}
\label{sites-definition-pushforward-algebraic-structures}
设 $f : \mathcal{D} \to \mathcal{C}$ 是由函子
$u : \mathcal{C} \to \mathcal{D}$ 给出的位点态射。对代数结构预层，按规则
$u^p\mathcal{F}(U) = \mathcal{F}(uU)$ 定义{\it 正像}函子；对代数结构层，
以同一规则定义，即 $f_*\mathcal{F}(U) = \mathcal{F}(uU)$。
\end{definition}

\noindent
问题出现在试图定义逆像时。原因是节 \ref{sites-section-functoriality-PSh} 中定义函子
$u_p$ 的余极限未必滤过。因此，一般而言，上面的公理不足以定义代数结构（预）层的
逆像。尽管如此，几乎在所有情形中，下述引理都足以定义代数结构（预）层的正像与逆像。
% P01-ERRATA: P01-E0155 -- the source says "trying the define" instead of "trying to define".

\begin{lemma}
\label{sites-lemma-push-pull-good-case}
假设函子 $u : \mathcal{C} \to \mathcal{D}$ 满足命题
\ref{sites-proposition-get-morphism} 的假设，因而产生位点态射
$f : \mathcal{D} \to \mathcal{C}$。在这种情形下，逆像函子
$f^{-1}$（相应地，$u_p$）与正像函子 $f_*$（相应地，$u^p$）扩张为代数结构层
（相应地，预层）范畴上的一对伴随函子。此外，这些函子与取底层集合层
（相应地，预层）的操作交换。
\end{lemma}

\begin{proof}
上面已经定义了 $f_* = u^p$。在命题 \ref{sites-proposition-get-morphism} 的证明过程中，
我们看到，在该命题的假设下，定义 $u_p$ 所用的所有余极限都是滤过的。因此由代数结构
类型的定义，可以用完全相同的余极限在代数结构预层上定义 $u_p$。$u_p$ 与 $u^p$
的伴随性，其证明与引理 \ref{sites-lemma-adjoints-u} 的证明完全相同。上面对代数结构预层
之层化的讨论，于是蕴含可以定义
$f^{-1}(\mathcal{F}) = (u_p\mathcal{F})^\#$。
\end{proof}

\noindent
下面简要讨论如何处理一般位点态射、更一般地任意拓扑斯态射的逆像与正像。

\medskip\noindent
设 $\mathcal{C}$ 为位点。在 $\mathcal{A} = \textit{Ab}$ 的情形，可以把
$\mathcal{C}$ 上的阿贝尔（预）层看成四元组
$(\mathcal{F}, +, 0, i)$。其中数据为
\begin{enumerate}
\item[(D1)] $\mathcal{F}$ 是集合层；
\item[(D2)] $+ : \mathcal{F} \times \mathcal{F} \to \mathcal{F}$ 是集合层的态射；
\item[(D3)] $0 : * \to \mathcal{F}$ 是从单点层（见例
\ref{sites-example-singleton-sheaf}）到 $\mathcal{F}$ 的态射；且
\item[(D4)] $i : \mathcal{F} \to \mathcal{F}$ 是集合层的态射。
\end{enumerate}
这些数据必须满足下列公理：
\begin{enumerate}
\item[(A1)] $+$ 结合且交换；
\item[(A2)] $0$ 是 $+$ 的单位元；且
\item[(A3)] $+ \circ (1, i) = 0 \circ (\mathcal{F} \to *)$。
\end{enumerate}
对照 Sheaves, Lemma \ref{sheaves-lemma-abelian-presheaves}。设
$f : \mathcal{D} \to \mathcal{C}$ 为位点态射。注意，由于 $f^{-1}$ 正合，有
$f^{-1}* = *$ 以及
$f^{-1}(\mathcal{F} \times \mathcal{F}) =
f^{-1}\mathcal{F} \times f^{-1}\mathcal{F}$。因此可以直接把
$f^{-1}\mathcal{F}$ 定义为四元组
$(f^{-1}\mathcal{F}, f^{-1}+, f^{-1}0, f^{-1}i)$。由于 $f^{-1}$ 是与有限极限
交换的函子，这些公理得以保持。最后，不难验证 $f_*$ 与 $f^{-1}$ 如通常那样伴随。

\medskip\noindent
在 \cite{SGA4} 中使用了这种方法。他们引入所谓
``{\it esp\`ece the structure alg\'ebrique $\ll$d\'efinie
par limites projectives finie$\gg$}''。对这样的 esp\`ece，可以用上述方法定义一对伴随
函子 $f^{-1}$ 与 $f_*$。这显然适用于实践中遇到的大多数代数结构。我们不把该构造
形式化，而只列出这种方法适用的代数结构（逐一验证）。事实上，该方法适用于任意
拓扑斯态射。
% P01-ERRATA: P01-E0156 -- the quoted French has "the" for "de" and singular "finie" where "limites" requires plural "finies".

\begin{proposition}
\label{sites-proposition-functoriality-algebraic-structures-topoi}
\begin{slogan}
拓扑斯态射保持代数结构。
\end{slogan}
设 $\mathcal{C}$、$\mathcal{D}$ 为位点。设 $f = (f^{-1}, f_*)$ 为从
$\Sh(\mathcal{D}) \to \Sh(\mathcal{C})$ 的拓扑斯态射。上述方法对下列代数结构
类型给出代数结构层上的一对伴随函子 $(f^{-1}, f_*)$，且与取底层集合层相容：
\begin{enumerate}
\item 带基点集合；
\item 阿贝尔群；
\item 群；
\item 幺半群；
\item 环；
\item 固定环上的模；且
\item 固定域上的李代数。
\end{enumerate}
此外，在每种情形中，上面标为 ($\alpha$)、($\beta$)、($\gamma$)、($\delta$)、
($\epsilon$)、($\zeta$) 的结果都成立。
\end{proposition}

\begin{proof}
命题的最后一个断言成立，是因为每个所列范畴配备显然的遗忘函子后，确实都是本节
开头所解释意义下的一种代数结构类型。见 Sheaves, Lemma
\ref{sheaves-lemma-list-algebraic-structures}。

\medskip\noindent
(2) 的证明。按照命题前的讨论，把阿贝尔群层看成四元组
$(\mathcal{F}, +, 0, i)$。若 $(\mathcal{F}, +, 0, i)$ 定义在
$\mathcal{C}$ 上，则其逆像定义为
$(f^{-1}\mathcal{F}, f^{-1}+, f^{-1}0, f^{-1}i)$。若
$(\mathcal{G}, +, 0, i)$ 定义在 $\mathcal{D}$ 上，则其正像定义为
$(f_*\mathcal{G}, f_*+, f_*0, f_*i)$。这是可行的，因为
$f_*(\mathcal{G} \times \mathcal{G}) = f_*\mathcal{G} \times f_*\mathcal{G}$。
伴随性由集合层函子的伴随性得出，因为
$$
\Mor_{\textit{Ab}(\mathcal{C})}
((\mathcal{F}_1, +, 0, i), (\mathcal{F}_2, +, 0, i))
=
\left\{
\begin{matrix}
\varphi \in \Mor_{\Sh(\mathcal{C})}
(\mathcal{F}_1, \mathcal{F}_2) \\
\varphi \text{ 与下列映射相容： }+, 0, i
\end{matrix}
\right\}
$$
细节留给读者。

\medskip\noindent
这种方法也适用于环层：把一个（有单位元的）环层看成满足初等代数教材中所列公理的
六元组 $(\mathcal{O}, + , 0, i, \cdot, 1)$ 即可。

\medskip\noindent
带基点集合层是序对 $(\mathcal{F}, p)$，其中 $\mathcal{F}$ 是集合层，而
$p : * \to \mathcal{F}$ 是集合层的映射。

\medskip\noindent
群层由带有适当公理的四元组 $(\mathcal{F}, \cdot, 1, i)$ 给出。

\medskip\noindent
幺半群层由带有适当公理的序对 $(\mathcal{F}, \cdot)$ 给出。
% P01-ERRATA: P01-E0157 -- the source has singular "axiom" after "suitable" although monoid structure requires multiple laws.

\medskip\noindent
设 $R$ 为环。$R$-模层由五元组
$(\mathcal{F}, +, 0, i, \{\lambda_r\}_{r \in R})$ 给出，其中四元组
$(\mathcal{F}, +, 0, i)$ 是上述阿贝尔群层，而
$\lambda_r : \mathcal{F} \to \mathcal{F}$ 是满足下列条件的一族集合层态射：
% P01-ERRATA: P01-E0158 -- the source says "An sheaf" instead of "A sheaf".
$\lambda_r \circ 0 = 0$，
$\lambda_r \circ + = + \circ (\lambda_r, \lambda_r)$，
$\lambda_{r + r'} =
+ \circ \lambda_r \times \lambda_{r'} \circ (\text{id}, \text{id})$，
$\lambda_{rr'} = \lambda_r \circ \lambda_{r'}$，
$\lambda_1 = \text{id}$，$\lambda_0 = 0 \circ (\mathcal{F} \to *)$。
% P01-ERRATA: P01-E0159 -- the distributivity formula needs parentheses around (lambda_r times lambda_r') to make the composition typed and unambiguous.
\end{proof}

\noindent
我们将在 Modules on Sites, Section
\ref{sites-modules-section-sheaves-modules} 中讨论环层上的模层范畴。

\begin{remark}
\label{sites-remark-no-pullback-presheaves}
设 $\mathcal{C}$、$\mathcal{D}$ 为位点。设
$u : \mathcal{D} \to \mathcal{C}$ 是连续函子，它产生位点态射
$\mathcal{C} \to \mathcal{D}$。注意，即使在阿贝尔群的情形，我们也尚未定义阿贝尔群
预层的逆像函子。由于阿贝尔群范畴中所有余极限均可表出，当然可以用定义集合预层上
$u_p$ 时所用的相同余极限，在阿贝尔预层上定义函子 $u_p^{ab}$。在阿贝尔预层范畴上，
$u_p^{ab}$ 也将左伴随于 $u^p$。然而，在底层集合预层上，$u_p^{ab}$ 未必总与
$u_p$ 一致。
\end{remark}

\section{拉回映射}
\label{sites-section-pullback}

\noindent
有时位点 $\mathcal{C}$ 没有终对象。在这种情形下，按如下方式定义整体截面函子。

\begin{definition}
\label{sites-definition-global-sections}
位点 $\mathcal{C}$ 上的集合预层 $\mathcal{F}$ 的{\it 整体截面}是集合
$$
\Gamma(\mathcal{C}, \mathcal{F}) =
\Mor_{\textit{PSh}(\mathcal{C})}(*, \mathcal{F})
$$
其中 $*$ 是 $\mathcal{C}$ 上预层范畴的终对象，即把单点集赋给每个对象的预层。
% P01-ERRATA: P01-E0160 -- the source's plural subject "global sections" is followed by singular "is".
\end{definition}

\noindent
当然，同一定义也适用于层。下面给出一种计算整体截面的方式。

\begin{lemma}
\label{sites-lemma-compute-global-sections}
设 $\mathcal{C}$ 为位点。设 $a, b : V \to U$ 是 $\mathcal{C}$ 的态射，使
$$
\xymatrix{
h_V^\# \ar@<1ex>[r] \ar@<-1ex>[r] & h_U^\# \ar[r] & {*}
}
$$
是 $\Sh(\mathcal{C})$ 中的余等化子。则
$\Gamma(\mathcal{C}, \mathcal{F})$ 是
$a^*, b^* : \mathcal{F}(U) \to \mathcal{F}(V)$ 的等化子。
% P01-ERRATA: P01-E0161 -- the source calls a,b:V->U "objects" of C; they are morphisms.
\end{lemma}

\begin{proof}
因为 $\Mor_{\Sh(\mathcal{C})}(h_U^\#, \mathcal{F}) = \mathcal{F}(U)$，
这由定义即得。
\end{proof}

\noindent
现在设 $f : \Sh(\mathcal{D}) \to \Sh(\mathcal{C})$ 为拓扑斯态射。则对
$\mathcal{C}$ 上的任意层 $\mathcal{F}$，存在拉回映射
$$
f^{-1} :
\Gamma(\mathcal{C}, \mathcal{F})
\longrightarrow
\Gamma(\mathcal{D}, f^{-1}\mathcal{F})
$$
事实上，由于 $f^{-1}$ 正合，它把 $*$ 映为 $*$。所以 $\mathcal{F}$ 在
$\mathcal{C}$ 上的整体截面 $s$，即层映射 $s : * \to \mathcal{F}$，可以拉回为
$f^{-1}s : * = f^{-1}* \to f^{-1}\mathcal{F}$。

\medskip\noindent
可以稍作推广：考虑分别定义在 $\mathcal{C}$、$\mathcal{D}$ 上的一对层
$\mathcal{F}$、$\mathcal{G}$，连同一个映射
$f^{-1}\mathcal{F} \to \mathcal{G}$。将上述构造与显然的映射
$\Gamma(\mathcal{D}, f^{-1}\mathcal{F}) \to \Gamma(\mathcal{D}, \mathcal{G})$
复合，便得到映射
$$
\Gamma(\mathcal{C}, \mathcal{F})
\longrightarrow
\Gamma(\mathcal{D}, \mathcal{G})
$$
该映射有时也称为拉回映射。

\medskip\noindent
自然出现的一个稍微更一般的构造如下。假设有拓扑斯态射的交换图
$$
\xymatrix{
\Sh(\mathcal{D}) \ar[rd]_h \ar[rr]_f & &
\Sh(\mathcal{C}) \ar[ld]^g \\
& \Sh(\mathcal{B})
}
$$
再假设 $\mathcal{F}$ 是 $\mathcal{C}$ 上的层。则存在{\it 拉回映射}
$$
f^{-1} : g_*\mathcal{F} \longrightarrow h_*f^{-1}\mathcal{F}
$$
它正是由识别
$g_*f_*f^{-1}\mathcal{F} = h_*f^{-1}\mathcal{F}$，以及对典范映射
$\mathcal{F} \to f_*f^{-1}\mathcal{F}$ 应用 $g_*$ 所得到的映射。若 $g$
是恒等态射，则该映射在整体截面上与上面的拉回映射一致。

\medskip\noindent
在前一段的情形下，再设有分别定义在 $\mathcal{C}$、$\mathcal{D}$ 上的一对层
$\mathcal{F}$、$\mathcal{G}$，连同一个映射
$f^{-1}\mathcal{F} \to \mathcal{G}$。把上述拉回映射与对
$f^{-1}\mathcal{F} \to \mathcal{G}$ 应用 $h_*$ 所得的映射复合，得到
$$
g_*\mathcal{F} \longrightarrow h_*\mathcal{G}
$$
限制到 $\mathcal{B}$ 的某个对象上的截面，就恢复上面讨论的整体截面“拉回映射”
（选取适当位点）。

\medskip\noindent
更一般地，设有拓扑斯的交换图
$$
\xymatrix{
\Sh(\mathcal{D}) \ar[d]_h \ar[r]_f &
\Sh(\mathcal{C}) \ar[d]^g \\
\Sh(\mathcal{B}) \ar[r]^e & \Sh(\mathcal{A})
}
$$
以及 $\mathcal{C}$ 上的层 $\mathcal{G}$。则存在{\it 基变换映射}
$$
e^{-1}g_*\mathcal{G} \longrightarrow h_*f^{-1}\mathcal{G}.
$$
该映射伴随于映射
$g_*\mathcal{G} \to e_*h_*f^{-1}\mathcal{G} = (e \circ h)_*f^{-1}\mathcal{G}$，
后者就是刚才描述的拉回映射。

\begin{remark}
\label{sites-remark-compose-base-change}
考虑交换图
$$
\xymatrix{
\Sh(\mathcal{B}') \ar[r]_k \ar[d]_{f'} &
\Sh(\mathcal{B}) \ar[d]^f \\
\Sh(\mathcal{C}') \ar[r]^l \ar[d]_{g'} &
\Sh(\mathcal{C}) \ar[d]^g \\
\Sh(\mathcal{D}') \ar[r]^m &
\Sh(\mathcal{D})
}
$$
则两个方块的基变换映射复合后给出外部矩形的基变换映射。更精确地，复合
\begin{align*}
m^{-1} \circ (g \circ f)_*
& =
m^{-1} \circ g_* \circ f_* \\
& \to g'_* \circ l^{-1} \circ f_* \\
& \to g'_* \circ f'_* \circ k^{-1} \\
& = (g' \circ f')_* \circ k^{-1}
\end{align*}
是该矩形的基变换映射。验证略去。
\end{remark}

\begin{remark}
\label{sites-remark-compose-base-change-horizontal}
考虑交换图
$$
\xymatrix{
\Sh(\mathcal{C}'') \ar[r]_{g'} \ar[d]_{f''} &
\Sh(\mathcal{C}') \ar[r]_g \ar[d]_{f'} &
\Sh(\mathcal{C}) \ar[d]^f \\
\Sh(\mathcal{D}'') \ar[r]^{h'} &
\Sh(\mathcal{D}') \ar[r]^h &
\Sh(\mathcal{D})
}
$$
的环化拓扑斯。则两个方块的基变换映射复合后给出外部矩形的基变换映射。更精确地，
复合
\begin{align*}
(h \circ h')^{-1} \circ f_*
& =
(h')^{-1} \circ h^{-1} \circ f_* \\
& \to (h')^{-1} \circ f'_* \circ g^{-1} \\
& \to f''_* \circ (g')^{-1} \circ g^{-1} \\
& = f''_* \circ (g \circ g')^{-1}
\end{align*}
是该矩形的基变换映射。验证略去。
% P01-ERRATA: P01-E0162 -- the source calls this a diagram of ringed topoi although only bare sheaf topoi are shown and no ring data enter the construction.
\end{remark}

\section{与 SGA4 的比较}
\label{sites-section-notation-comparison}

\noindent
节 \ref{sites-section-functoriality-PSh} 中函子 $u^p$、$u_p$ 的记号，以及节
\ref{sites-section-continuous-functors} 中 $u^s$、$u_s$ 的记号，取自
\cite[pages 14 and 42]{ArtinTopologies}。作出这些选择之后，节
\ref{sites-section-more-functoriality-PSh} 中函子 ${}_pu$ 与节
\ref{sites-section-cocontinuous-functors} 中函子 ${}_su$ 的记号看起来也很合理。
本节将我们的记号与 SGA4 的记号比较。

\medskip\noindent
{\bf 预层：}设 $u : \mathcal{C} \to \mathcal{D}$ 为范畴间的函子。
在 \cite[Exposee I, Section 5]{SGA4} 中，函子 $u^p$ 记为 $u^*$。
在 \cite[Exposee I, Proposition 5.1]{SGA4} 中，函子 $u_p$ 记为 $u_!$；
在 \cite[Exposee I, Proposition 5.1]{SGA4} 中，函子 ${}_pu$ 记为 $u_*$。
换言之，
$$
u_p, u^p, {}_pu\quad(SP)
\quad\text{与}\quad
u_!, u^*, u_*\quad(SGA4)
$$
需要提醒读者，SGA4 的不同部分对这些函子使用不同记号。

\medskip\noindent
{\bf 层与连续函子：}假设 $\mathcal{C}$ 与 $\mathcal{D}$ 为位点，而
$u : \mathcal{C} \to \mathcal{D}$ 是连续函子（定义
\ref{sites-definition-continuous}）。在 \cite[Exposee III, 1.11]{SGA4} 中，
函子 $u^s$ 记为 $u_s$；在 \cite[Exposee III, Proposition 1.2]{SGA4} 中，
函子 $u_s$ 记为 $u^s$。换言之，
$$
u_s, u^s\quad(SP)
\quad\text{与}\quad
u^s, u_s\quad(SGA4)
$$
当 $u$ 定义位点态射 $f : \mathcal{D} \to \mathcal{C}$（定义
\ref{sites-definition-morphism-sites}）时，相应的拓扑斯态射（引理
\ref{sites-lemma-morphism-sites-topoi}）与 \cite[Exposee IV, (4.9.1.1)]{SGA4}
中的相同。

\medskip\noindent
{\bf 层与余连续函子：}假设 $\mathcal{C}$ 与 $\mathcal{D}$ 为位点，而
$u : \mathcal{C} \to \mathcal{D}$ 是余连续函子（定义
\ref{sites-definition-cocontinuous}）。在
\cite[Exposee III, Proposition 2.3]{SGA4} 中，函子 ${}_su$（引理
\ref{sites-lemma-pu-sheaf}）记为 $u_*$；在
\cite[Exposee III, Proposition 2.3]{SGA4} 中，函子 $(u^p\ )^\#$ 记为 $u^*$。
换言之，
$$
(u^p\ )^\#, {}_su\quad(SP)
\quad\text{与}\quad
u^*, u_*\quad(SGA4)
$$
因此，引理 \ref{sites-lemma-cocontinuous-morphism-topoi} 中与 $u$ 关联的拓扑斯态射，
与 \cite[Exposee IV, 4.7]{SGA4} 中的相同。

\medskip\noindent
{\bf 拓扑斯态射：}若 $f$ 是由函子 $(f^{-1}, f_*)$ 给出的拓扑斯态射，则在
\cite[Exposee IV, Definition 3.1]{SGA4} 中，函子 $f^{-1}$ 记为 $f^*$。
我们将用 $f^{-1}$ 表示集合层、或更一般地代数结构层的逆像（节
\ref{sites-section-sheaves-algebraic-structures}）。对于环化拓扑斯的态射，我们将用
$f^*$ 表示模层的逆像（Modules on Sites, Definition
\ref{sites-modules-definition-pushforward}）。

\section{拓扑}
\label{sites-section-topologies}

\noindent
本节按照 \cite{SGA4} 定义范畴上的拓扑。可以用这种语言发展层与拓扑斯的全部理论。
关于这些材料的现代论述见 \cite{KS}。不过，对代数几何最重要的情形，是位点在其
底层范畴上定义的拓扑。因此，我们强烈建议初次阅读的读者
{\bf 跳过本节以及本章后面的所有其他章节}！

\begin{definition}
\label{sites-definition-sieve}
设 $\mathcal{C}$ 为范畴，$U \in \Ob(\mathcal{C})$。$U$ 上的一个
{\it 筛 $S$} 是子预层 $S \subset h_U$。
\end{definition}

\noindent
换言之，$U$ 上的筛对每个对象 $T \in \Ob(\mathcal{C})$，从所有态射
$T \to U$ 所成的集合中选出一个子集 $S(T)$。事实上，对子集族
$S(T) \subset h_U(T) = \Mor_\mathcal{C}(T, U)$ 的唯一要求是下列规则：
\begin{equation}
\label{sites-equation-property-sieve}
\left.
\begin{matrix}
(\alpha : T \to U) \in S(T) \\
g : T' \to T
\end{matrix}
\right\} \Rightarrow
(\alpha \circ g : T' \to U) \in S(T')
\end{equation}
一个有用的直观图景，是把映射 $S \to h_U$ 看成“从 $S$ 到 $U$ 的态射”。

\begin{lemma}
\label{sites-lemma-sieves-set}
设 $\mathcal{C}$ 为范畴，$U \in \Ob(\mathcal{C})$。
\begin{enumerate}
\item $U$ 上所有筛构成集合。
\item 包含关系在该集合上定义偏序。
\item 筛的并与交仍是筛。
\item
\label{sites-item-sieve-generated}
给定 $\mathcal{C}$ 中以 $U$ 为目标的态射族
$\{U_i \to U\}_{i\in I}$，存在唯一最小的 $U$ 上的筛 $S$，使每个
$U_i \to U$ 都属于 $S(U_i)$。
\item 筛 $S = h_U$ 是最大筛。
\item 空子预层是最小筛。
\end{enumerate}
\end{lemma}

\begin{proof}
由子预层的定义，预层 $\mathcal{F}$ 的所有子预层所成的集族，是
$\prod_{U \in \Ob(\mathcal{C})} \mathcal{P}(\mathcal{F}(U))$ 的子集。
后者是集合。（这里 $\mathcal{P}(A)$ 表示 $A$ 的幂集。）因此 $U$ 上所有筛构成集合。

\medskip\noindent
偏序定义如下：$S \leq S'$ 当且仅当对所有 $T \to U$，有
$S(T) \subset S'(T)$。记作 $S \subset S'$。

\medskip\noindent
给定 $U$ 上的一族筛 $S_i$（$i \in I$），可以把 $\bigcup S_i$ 定义为满足
$(\bigcup S_i)(T) = \bigcup S_i(T)$（对所有 $T \in \Ob(\mathcal{C})$）的筛。
交 $\bigcap S_i$ 以同样方式定义。

\medskip\noindent
给定命题中的 $\{U_i \to U\}_{i\in I}$，考虑预层态射
$h_{U_i} \to h_U$。直接把 $S$ 定义为这些预层映射之像（定义
\ref{sites-definition-image}）的并。

\medskip\noindent
引理的最后两个断言显然。
\end{proof}

\begin{definition}
\label{sites-definition-sieve-generated}
设 $\mathcal{C}$ 为范畴。给定 $\mathcal{C}$ 中以 $U$ 为目标的态射族
$\{f_i : U_i \to U\}_{i\in I}$，称引理 \ref{sites-lemma-sieves-set} 第
(\ref{sites-item-sieve-generated}) 部分所述的 $U$ 上的筛 $S$ 为
{\it 由态射 $f_i$ 生成的 $U$ 上的筛}。
\end{definition}

\begin{definition}
\label{sites-definition-pullback-sieve}
设 $\mathcal{C}$ 为范畴，$f : V \to U$ 是 $\mathcal{C}$ 的态射，而
$S \subset h_U$ 是筛。定义 $S$ 沿 $f$ 的{\it 拉回}为 $V$ 上的筛
$S \times_U V$，规则为
$$
(\alpha : T \to V) \in (S \times_U V)(T)
\Leftrightarrow
(f \circ \alpha : T \to U) \in S(T)
$$
\end{definition}

\noindent
留给读者验证这确实是筛（提示：使用等式 \ref{sites-equation-property-sieve}）。
有时也称 $S \times_U V$ 为 $S$ 沿 $f : V \to U$ 的{\it 基变换}。

\begin{lemma}
\label{sites-lemma-pullback-sieve-section}
设 $\mathcal{C}$ 为范畴，$U \in \Ob(\mathcal{C})$，而 $S$ 是 $U$ 上的筛。
若 $f : V \to U$ 属于 $S$，则 $S \times_U V = h_V$ 是最大筛。
\end{lemma}

\begin{proof}
由定义显然。
\end{proof}

\begin{definition}
\label{sites-definition-topology}
设 $\mathcal{C}$ 为范畴。$\mathcal{C}$ 上的一个{\it 拓扑}，是把 $U$ 上所有筛
所成集合的一个子集 $J(U)$ 赋给每个 $U \in \Ob(\mathcal{C})$ 的规则，并满足：
\begin{enumerate}
\item 对 $\mathcal{C}$ 中每个态射 $f : V \to U$ 以及每个
$S \in J(U)$，拉回 $S \times_U V$ 属于 $J(V)$。
\item 若 $S$、$S'$ 是 $U \in \Ob(\mathcal{C})$ 上的筛，$S \in J(U)$，
并且对所有 $f \in S(V)$，拉回 $S' \times_U V$ 都属于 $J(V)$，则
$S'$ 属于 $J(U)$。
\item 对每个 $U \in \Ob(\mathcal{C})$，最大筛 $S = h_U$ 属于 $J(U)$。
\end{enumerate}
\end{definition}

\noindent
在这种情形下，属于 $J(U)$ 的筛称为{\it 覆盖筛}。

\begin{lemma}
\label{sites-lemma-topology-basic}
设 $\mathcal{C}$ 为范畴，$J$ 是 $\mathcal{C}$ 上的拓扑，而
$U \in \Ob(\mathcal{C})$。
\begin{enumerate}
\item $J(U)$ 中元素的有限交仍属于 $J(U)$。
\item 若 $S \in J(U)$ 且 $S' \supset S$，则 $S' \in J(U)$。
\end{enumerate}
\end{lemma}

\begin{proof}
(1) 的证明。空族之交的情形由定义 \ref{sites-definition-topology} 的条件 (3) 得出。
设 $S, S' \in J(U)$，并考虑 $S'' = S \cap S'$。对 $S(V)$ 中的每个
$V \to U$，有
$$
S' \times_U V = S'' \times_U V
$$
这是因为 $V \to U$ 已经属于 $S$。所以由定义的第二条公理，
$S'' \in J(U)$。

\medskip\noindent
(2) 的证明。设 $S \in J(U)$ 且 $S' \supset S$。对 $S(V)$ 中的每个
$V \to U$，由引理 \ref{sites-lemma-pullback-sieve-section}，有
$S' \times_U V = h_V$。所以由第三条公理，$S' \times_U V \in J(V)$；
再由第二条公理，$S' \in J(U)$。
\end{proof}

\begin{definition}
\label{sites-definition-finer}
设 $\mathcal{C}$ 为范畴，$J$、$J'$ 是 $\mathcal{C}$ 上的两个拓扑。若对
$\mathcal{C}$ 的每个对象 $U$ 都有 $J'(U) \subset J(U)$，则称 $J$ 比
$J'$ {\it 更细}或{\it 更强}。在这种情形下，也称 $J'$ 比 $J$
{\it 更粗}或{\it 更弱}。
\end{definition}

\noindent
换言之，$J'$ 的任意覆盖筛都是 $J$ 的覆盖筛。$\mathcal{C}$ 上存在最细拓扑，
即每个筛都是覆盖筛的拓扑，称为 $\mathcal{C}$ 的{\it 离散拓扑}。也存在最粗拓扑，
即对所有对象 $U$ 都有 $J(U) = \{h_U\}$ 的拓扑，称为{\it 混沌拓扑}或
{\it 平凡拓扑}。

\begin{lemma}
\label{sites-lemma-play-with-topologies}
设 $\mathcal{C}$ 为范畴，而 $\{J_i\}_{i\in I}$ 是一族拓扑。
\begin{enumerate}
\item 规则 $J(U) = \bigcap J_i(U)$ 定义 $\mathcal{C}$ 上的拓扑。
\item 存在比所有拓扑 $J_i$ 都更细的最粗拓扑。
\end{enumerate}
\end{lemma}

\begin{proof}
第一部分由定义直接得出。第二部分由取所有比各 $J_i$ 都更细的拓扑之交得出。
\end{proof}

\noindent
至此，可以不加任何动机地定义层。

\begin{definition}
\label{sites-definition-sheaf-sets-topology}
设范畴 $\mathcal{C}$ 配备拓扑 $J$，而 $\mathcal{F}$ 是 $\mathcal{C}$ 上的集合预层。
若对每个 $U \in \Ob(\mathcal{C})$ 以及 $U$ 的每个覆盖筛 $S$，典范映射
$$
\Mor_{\textit{PSh}(\mathcal{C})}(h_U, \mathcal{F})
\longrightarrow
\Mor_{\textit{PSh}(\mathcal{C})}(S, \mathcal{F})
$$
都是双射，则称 $\mathcal{F}$ 是 $\mathcal{C}$ 上的{\it 层}。
\end{definition}

\noindent
回忆所示公式的左端等于 $\mathcal{F}(U)$。换言之，$\mathcal{F}$ 是层，当且仅当
$\mathcal{F}$ 在 $U$ 上的一个截面，等同于由
$(\alpha : T \to U) \in S(T)$ 参数化的一族相容截面
$s_{T, \alpha} \in \mathcal{F}(T)$；并且这对 $U$ 上每个覆盖筛 $S$ 都成立。

\begin{lemma}
\label{sites-lemma-topology-presheaves-sheaves}
设 $\mathcal{C}$ 为范畴，$\{ \mathcal{F}_i \}_{i\in I}$ 是
$\mathcal{C}$ 上的一族集合预层。对每个 $U \in \Ob(\mathcal{C})$，以
$J(U)$ 表示 $U$ 上满足下列性质的筛 $S$ 所成的集合：对每个态射
$V \to U$，映射
% P01-ERRATA: P01-E0163 -- the source's "denote J(U) the set" omits "by" or needs recasting.
$$
\Mor_{\textit{PSh}(\mathcal{C})}(h_V, \mathcal{F}_i)
\longrightarrow
\Mor_{\textit{PSh}(\mathcal{C})}(S \times_U V, \mathcal{F}_i)
$$
对所有 $i \in I$ 都是双射。则 $J$ 定义 $\mathcal{C}$ 上的拓扑。它是使所有
$\mathcal{F}_i$ 都成为层的最细拓扑。
\end{lemma}

\begin{proof}
首先注意，若 $J$ 是拓扑，则引理的最后一个断言立即成立。

\medskip\noindent
对每个 $i \in I$ 与 $U \in \Ob(\mathcal{C})$，令 $J_i(U)$ 为 $U$ 上满足
下列性质的筛 $S$ 所成的集合：对每个态射 $V \to U$，映射
$$
\Mor_{\textit{PSh}(\mathcal{C})}(h_V, \mathcal{F}_i)
\longrightarrow
\Mor_{\textit{PSh}(\mathcal{C})}(S \times_U V, \mathcal{F}_i)
$$
是双射。注意 $J(U) = \bigcap J_i(U)$。因此，若证明 $J_i$ 定义拓扑，就由引理
\ref{sites-lemma-play-with-topologies} 得到结论。于是归约到下一段讨论的情形。

\medskip\noindent
证明当预层族仅含一个预层 $\mathcal{F}$ 时，$J$ 是拓扑。拓扑的第一条和第三条公理
立即得到验证。设有对象 $U$ 以及 $U$ 上的筛 $S, S'$，使
$S \in J(U)$，并且对 $S(V)$ 中的所有 $V \to U$，都有
$S' \times_U V \in J(V)$。需要证明 $S' \in J(U)$。换言之，需要证明对任意
$f : W \to U$，映射
$$
\mathcal{F}(W) =
\Mor_{\textit{PSh}(\mathcal{C})}(h_W, \mathcal{F})
\longrightarrow
\Mor_{\textit{PSh}(\mathcal{C})}(S' \times_U W, \mathcal{F})
$$
都是双射。取元素
$\varphi \in \Mor_{\textit{PSh}(\mathcal{C})}(S' \times_U W, \mathcal{F})$。
我们将构造映到 $\varphi$ 的截面 $s \in \mathcal{F}(W)$。

\medskip\noindent
设 $\alpha : V \to W$ 是 $S \times_U W$ 的一个元素。由拉回的定义，复合
$f \circ \alpha : V \to W  \to U$ 属于 $S$。因此，由对筛对 $S, S'$ 的假设，
$S' \times_U V$ 属于 $J(V)$。现在有预层的交换图
$$
\xymatrix{
S' \times_U V \ar[r] \ar[d] & h_V \ar[d] \\
S' \times_U W \ar[r] & h_W
}
$$
$\varphi$ 在 $S' \times_U V$ 上的限制对应于元素
$s_{V, \alpha} \in \mathcal{F}(V)$。这是由 $J$ 的定义以及
$S' \times_U V$ 属于 $J(V)$ 得出的。留给读者验证：规则
$(V, \alpha) \mapsto s_{V, \alpha}$ 定义元素
$\psi \in \Mor_{\textit{PSh}(\mathcal{C})}(S \times_U W, \mathcal{F})$。
由于 $S \in J(U)$，从 $J$ 的定义立即看出，$\psi$ 对应于
$\mathcal{F}(W)$ 的一个元素 $s$。

\medskip\noindent
留给读者验证：构造 $\varphi \mapsto s$ 是上面所示自然映射的逆。
\end{proof}

\begin{definition}
\label{sites-definition-canonical-topology}
设 $\mathcal{C}$ 为范畴。使所有可表预层均为层的 $\mathcal{C}$ 上最细拓扑（见引理
\ref{sites-lemma-topology-presheaves-sheaves}），称为 $\mathcal{C}$ 的
{\it 典范拓扑}。
\end{definition}

\section{位点定义的拓扑}
\label{sites-section-topology-site}

\noindent
假设 $\mathcal{C}$ 是范畴，而 $\text{Cov}_1(\mathcal{C})$ 与
$\text{Cov}_2(\mathcal{C})$ 是在 $\mathcal{C}$ 上定义位点结构的两组覆盖。
在这种情形下，$\text{Cov}_1(\mathcal{C})$ 与
$\text{Cov}_2(\mathcal{C})$ 的（集合）层范畴可能相同；例如见引理
\ref{sites-lemma-refine-same-topology}。

\medskip\noindent
事实上，$\mathcal{C}$ 关于某个拓扑 $J$ 的层范畴与 $J$ 彼此决定。
这是一个非平凡的断言，稍后将在定理 \ref{sites-theorem-topology-and-topos}
中讨论。

\medskip\noindent
暂且接受这一点，研究位点所决定的拓扑就是合理的。

\begin{lemma}
\label{sites-lemma-site-gives-topology}
设 $\mathcal{C}$ 是覆盖为 $\text{Cov}(\mathcal{C})$ 的位点。对
$\mathcal{C}$ 的每个对象 $U$，令 $J(U)$ 为 $U$ 上具有下述性质的筛 $S$
所成的集合：存在覆盖
$\{f_i : U_i \to U\}_{i\in I} \in \text{Cov}(\mathcal{C})$，使由 $f_i$
生成的筛 $S'$（见定义 \ref{sites-definition-sieve-generated}）包含于 $S$。
\begin{enumerate}
\item 该 $J$ 是 $\mathcal{C}$ 上的拓扑。
\item 预层 $\mathcal{F}$ 是该拓扑的层（见定义
\ref{sites-definition-sheaf-sets-topology}），当且仅当它是该位点上的层（见定义
\ref{sites-definition-sheaf-sets}）。
\end{enumerate}
\end{lemma}

\begin{proof}
为证明第一个断言，只需注意位点定义（定义 \ref{sites-definition-site}）的公理
(1)、(2)、(3) 分别直接蕴含拓扑定义（定义 \ref{sites-definition-topology}）的公理
(3)、(2)、(1)。作为例子，证明 $J$ 具有性质 (2)。设 $U$ 是
$\mathcal{C}$ 的对象，而 $S, S'$ 是 $U$ 上的筛，满足 $S \in J(U)$，并且对
$S(V)$ 中的每个 $V \to U$ 都有 $S' \times_U V \in J(V)$。由 $J(U)$
的定义，可找到该位点的覆盖 $\{f_i : U_i \to U\}$，使 $S$，即
$h_{U_i} \to h_U$ 的像所成的筛，包含于 $S$。
% P01-ERRATA: P01-E0164 -- the source reads "such that S the image ... is contained in S", duplicating S and omitting the grammatical construction.
由于每个 $S'\times_U U_i$ 属于 $J(U_i)$，可见存在该位点的覆盖
$\{U_{ij} \to U_i\}$，使 $h_{U_{ij}} \to h_{U_i}$ 包含于
$S' \times_U U_i$。由基变换的定义，这意味着
$h_{U_{ij}} \to h_U$ 包含于子预层 $S' \subset h_U$。由位点的公理 (2)，
$\{U_{ij} \to U\}$ 是 $U$ 的覆盖；再由 $J$ 的定义，得出
$S' \in J(U)$。

\medskip\noindent
设 $\mathcal{F}$ 为预层。假设 $\mathcal{F}$ 是拓扑 $J$ 的层。我们将证明
$\mathcal{F}$ 也是该位点上的层。设
$\{f_i : U_i \to U\}_{i\in I}$ 是该位点的覆盖，而
$s_i \in \mathcal{F}(U_i)$ 是满足
$s_i|_{U_i \times_U U_j} = s_j|_{U_i \times_U U_j}$（对所有 $i, j$）的一族截面。
需要证明存在唯一的截面 $s \in \mathcal{F}(U)$，其在各 $U_i$ 上的限制为
$s_i$。设 $S \subset h_U$ 为由 $f_i$ 生成的筛。按定义，$S \in J(U)$。
与其直接构造 $s$，由拓扑中的层条件，只须构造元素
% P01-ERRATA: P01-E0165 -- the source spells "Instead" as two words, "In stead".
$$
\varphi \in \Mor_{\textit{PSh}(\mathcal{C})}(S, \mathcal{F}).
$$
对某个对象 $T \in \mathcal{U}$，取 $\alpha \in S(T)$。
% P01-ERRATA: P01-E0166 -- the category U is undefined here; T should be an object of C.
这恰好意味着 $\alpha : T \to U$ 是一个通过某个 $i\in I$ 的 $f_i$ 分解的态射
（可能通过多于 $1$ 个）。选取这样的指标 $i$ 与分解
$\alpha = f_i \circ \alpha_i$，并定义 $\varphi(\alpha) = \alpha_i^* s_i$。
若 $i'$、$\alpha = f_i \circ \alpha_{i'}'$ 是第二种选择，则
% P01-ERRATA: P01-E0167 -- the second factorization uses f_i again; it should use f_{i'}.
$\alpha_i^* s_i = (\alpha_{i'}')^* s_{i'}$，这恰由对所有 $i, j$ 的条件
$s_i|_{U_i \times_U U_j} = s_j|_{U_i \times_U U_j}$ 得出。因此
$\varphi(\alpha)$ 良定义。留给读者验证：$\varphi$（它进而决定 $s$）是正确的，
即 $s$ 的相应限制为 $s_i$。

\medskip\noindent
设 $\mathcal{F}$ 为预层。假设 $\mathcal{F}$ 是位点
$(\mathcal{C}, \text{Cov}(\mathcal{C}))$ 上的层。我们将证明
$\mathcal{F}$ 也是拓扑 $J$ 的层。设 $U$ 是 $\mathcal{C}$ 的对象，而 $S$ 是
$U$ 上关于拓扑 $J$ 的覆盖筛。设
$$
\varphi \in \Mor_{\textit{PSh}(\mathcal{C})}(S, \mathcal{F}).
$$
需要证明 $\mathcal{F}(U) =
\Mor_{\textit{PSh}(\mathcal{C})}(h_U, \mathcal{F})$ 中存在唯一元素，限制为
$\varphi$。由定义，存在覆盖
$\{f_i : U_i \to U\}_{i\in I} \in \text{Cov}(\mathcal{C})$，使
$f_i : U_i \to U$ 属于 $S(U_i)$。因此可以令
$s_i = \varphi(f_i) \in \mathcal{F}(U_i)$。容易验证，对所有 $i, j$，有
$s_i|_{U_i \times_U U_j} = s_j|_{U_i \times_U U_j}$。于是由
$\mathcal{F}$ 在位点 $(\mathcal{C}, \text{Cov}(\mathcal{C}))$ 上的层条件，
得到所需截面 $s$。细节留给读者。
\end{proof}

\begin{definition}
\label{sites-definition-topology-associated-site}
设 $\mathcal{C}$ 是覆盖为 $\text{Cov}(\mathcal{C})$ 的位点。上面引理
\ref{sites-lemma-site-gives-topology} 中构造的拓扑 $J$，称为
{\it 与 $\mathcal{C}$ 关联的拓扑}。
\end{definition}

\noindent
设 $\mathcal{C}$ 为范畴。设 $\text{Cov}_1(\mathcal{C})$ 与
$\text{Cov}_2(\mathcal{C})$ 是在 $\mathcal{C}$ 上定义位点结构的两组覆盖。
它们定义的拓扑完全可能相同。若如此，则称
$\text{Cov}_1(\mathcal{C})$ 与 $\text{Cov}_2(\mathcal{C})$
在 $\mathcal{C}$ 上{\it 定义相同的拓扑}。并且在这种情形下，由引理
\ref{sites-lemma-site-gives-topology}，层范畴相同。
% P01-ERRATA: P01-E0168 -- the source calls Cov_1 and Cov_2 "two coverings" although each denotes a set/class of coverings defining a site.

\medskip\noindent
通常，我们只关心一组覆盖所定义的拓扑，并把选择不同覆盖集的可能性看作研究该拓扑
的工具。

\begin{remark}
\label{sites-remark-enlarge-coverings}
扩大覆盖类。显然，若 $\text{Cov}(\mathcal{C})$ 在 $\mathcal{C}$ 上定义位点结构，
则可以向 $\mathcal{C}$ 加入任意一组具有固定目标的态射族，只要它们与
$\text{Cov}(\mathcal{C})$ 的元素恒真等价（见定义
\ref{sites-definition-combinatorial-tautological}），而不改变拓扑。
% P01-ERRATA: P01-E0169 -- the source says to add covering families "to C"; they should be added to Cov(C), not to the category itself.
\end{remark}

\begin{remark}
\label{sites-remark-shrink-coverings}
缩小覆盖类。设 $\mathcal{C}$ 为位点。考虑集合
$$
\mathcal{S} = P(\text{Arrows}(\mathcal{C})) \times \Ob(\mathcal{C})
$$
其中 $P(\text{Arrows}(\mathcal{C}))$ 是态射集的幂集，即由所有态射集组成的集合。
令 $\mathcal{S}_\tau \subset \mathcal{S}$ 为所有满足下列条件的
$(T, U) \in \mathcal{S}$ 所成的子集：
(a) 所有 $\varphi \in T$ 的目标都是 $U$；
(b) 族 $\{\varphi\}_{\varphi \in T}$ 与 $\text{Cov}(\mathcal{C})$
中的某个覆盖恒真等价（见定义 \ref{sites-definition-combinatorial-tautological}）。
显然，若把 $\mathcal{S}_\tau$ 的元素看作覆盖，并不会恰好得到定义
\ref{sites-definition-site} 中所定义的位点概念。所得结构
$(\mathcal{C}, \mathcal{S}_\tau)$ 满足略经修改的条件。修改后的条件为：
\begin{enumerate}
\item[(0')] $\text{Cov}(\mathcal{C}) \subset
P(\text{Arrows}(\mathcal{C})) \times \Ob(\mathcal{C})$,
\item[(1')] 若 $V \to U$ 是同构，则
$(\{V \to U\}, U) \in \text{Cov}(\mathcal{C})$。
\item[(2')] 若 $(T, U) \in \text{Cov}(\mathcal{C})$，并且对
$T$ 中的 $f : U' \to U$ 给定
$(T_f, U') \in \text{Cov}(\mathcal{C})$，则令
$T' = \{f \circ f' \mid f \in T,\ f' \in T_f\}$，有
$(T', U) \in \text{Cov}(\mathcal{C})$。
\item[(3')] 若 $(T, U) \in \text{Cov}(\mathcal{C})$，且
$g : V \to U$ 是 $\mathcal{C}$ 的态射，则
\begin{enumerate}
\item 对 $T$ 中的 $f : U' \to U$，$U' \times_{f, U, g} V$ 存在；且
\item 对纤维积的某种选择令
$T' = \{\text{pr}_2 : U' \times_{f, U, g} V \to V \mid f : U' \to U \in T\}$，
则 $(T', V) \in \text{Cov}(\mathcal{C})$。
\end{enumerate}
\end{enumerate}
容易验证：给定满足上述 (0') -- (3') 的结构，适当扩大
$\text{Cov}(\mathcal{C})$ 后（对照 Sets, Section
\ref{sets-section-coverings-site}），得到一个位点。至少从拓扑角度看，该概念与实际的
位点概念显然差别不大。它有两个好处：由上面的条件 (0')，覆盖自动构成集合；同样由
(0')，所有这类结构的总体也构成集合。代价是必须到处写“恒真等价”。
\end{remark}

\section{拓扑中的层化}
\label{sites-section-sheafification-topology}

\noindent
本节说明拓扑中与层化构造对应的构造。

\medskip\noindent
设 $\mathcal{C}$ 为范畴，$J$ 是 $\mathcal{C}$ 上的拓扑，而
$\mathcal{F}$ 是集合预层。对每个 $U \in \Ob(\mathcal{C})$，把
$$
L\mathcal{F}(U)
=
\colim_{S \in J(U)^{opp}}
\Mor_{\textit{PSh}(\mathcal{C})}(S, \mathcal{F})
$$
定义为一个余极限。这里把 $J(U)$ 看成按包含关系排序的偏序集；见引理
\ref{sites-lemma-sieves-set}。该系统的转移映射定义如下。若
$S \subset S'$ 属于 $J(U)$，则 $S \to S'$ 是预层态射，因而存在自然的限制映射
$$
\Mor_{\textit{PSh}(\mathcal{C})}(S', \mathcal{F})
\longrightarrow
\Mor_{\textit{PSh}(\mathcal{C})}(S, \mathcal{F}).
$$
所以，只要反转 $J(U)$ 上的次序（这正是上标 ${}^{opp}$ 所表示的意思），
$S \mapsto \Mor_{\textit{PSh}(\mathcal{C})}(S, \mathcal{F})$
就是 Categories, Definition \ref{categories-definition-system-over-poset}
意义下的有向系统。特别地，因为 $h_U \in J(U)$，存在来自识别
$\mathcal{F}(U) = \Mor_{\textit{PSh}(\mathcal{C})}(h_U, \mathcal{F})$
的典范映射
$$
\ell : \mathcal{F}(U) \longrightarrow L\mathcal{F}(U)
$$
此外，定义 $L\mathcal{F}(U)$ 的余极限是有向的，因为对 $U$ 上任意一对覆盖筛
$S, S'$，筛 $S \cap S'$ 仍是覆盖筛；见引理 \ref{sites-lemma-sieves-set}。
% P01-ERRATA: P01-E0170 -- closure of covering sieves under finite intersections is Lemma lemma-topology-basic, not Lemma lemma-sieves-set.

\medskip\noindent
设 $f : V \to U$ 是 $\mathcal{C}$ 中的态射，并设 $S \in J(U)$。有交换图
$$
\xymatrix{
S \times_U V \ar[r] \ar[d] & h_V \ar[d] \\
S \ar[r] & h_U
}
$$
可以用左侧竖直映射得到典范限制映射
$$
\Mor_{\textit{PSh}(\mathcal{C})}(S, \mathcal{F})
\to \Mor_{\textit{PSh}(\mathcal{C})}(S \times_U V, \mathcal{F}).
$$
基变换 $S \mapsto S \times_U V$ 诱导保序映射 $J(U) \to J(V)$，而限制映射
定义 Categories, Lemma {categories-lemma-functorial-colimit} 所述的函子变换。
% P01-ERRATA: P01-E0171 -- the source omits the \ref command around categories-lemma-functorial-colimit.
所以得到自然的限制映射
$$
L\mathcal{F}(U) \longrightarrow L\mathcal{F}(V).
$$

\begin{lemma}
\label{sites-lemma-L-presheaf}
在上述情形中：
% P01-ERRATA: P01-E0172 -- the source's "In the situation above." is a sentence fragment before the list.
\begin{enumerate}
\item 赋值 $U \mapsto L\mathcal{F}(U)$ 连同上面定义的限制映射构成预层。
\item 各映射 $\ell$ 粘合为预层态射
$\ell : \mathcal{F} \to L\mathcal{F}$。
\item 规则 $\mathcal{F} \mapsto (\mathcal{F} \xrightarrow{\ell}
L\mathcal{F})$ 是函子。
\item 若 $\mathcal{F}$ 是 $\mathcal{G}$ 的子预层，则
$L\mathcal{F}$ 是 $L\mathcal{G}$ 的子预层。
\item 映射 $\ell : \mathcal{F} \to L\mathcal{F}$ 具有下列性质：对每个截面
$s \in L\mathcal{F}(U)$，存在 $U$ 上的覆盖筛 $S$ 以及元素
$\varphi \in \Mor_{\textit{PSh}(\mathcal{C})}(S, \mathcal{F})$，使
$\ell(\varphi)$ 等于 $s$ 在 $S$ 上的限制。
\end{enumerate}
\end{lemma}

\begin{proof}
略。
\end{proof}

\begin{definition}
\label{sites-definition-presheaf-separated-topology}
设 $\mathcal{C}$ 为范畴，而 $J$ 是 $\mathcal{C}$ 上的拓扑。若对每个对象 $U$
以及 $U$ 上的每个覆盖筛 $S$，典范映射
$\mathcal{F}(U) \to \Mor_{\textit{PSh}(\mathcal{C})}(S, \mathcal{F})$
都是单射，则称集合预层 $\mathcal{F}$ 是{\it 分离的}。
\end{definition}

\begin{theorem}
\label{sites-theorem-L-topology}
设 $\mathcal{C}$ 为范畴，$J$ 是 $\mathcal{C}$ 上的拓扑，而
$\mathcal{F}$ 是集合预层。
\begin{enumerate}
\item 预层 $L\mathcal{F}$ 是分离的。
\item 若 $\mathcal{F}$ 分离，则 $L\mathcal{F}$ 是层，并且预层映射
$\mathcal{F} \to L\mathcal{F}$ 是单射。
\item 若 $\mathcal{F}$ 是层，则 $\mathcal{F} \to L\mathcal{F}$ 是同构。
\item 预层 $LL\mathcal{F}$ 总是层。
\end{enumerate}
\end{theorem}

\begin{proof}
(3) 由 $L$ 的定义与层的定义（定义 \ref{sites-definition-sheaf-sets-topology}）
立即得出。(4) 形式地由其余部分得出。

\medskip\noindent
概述 (1) 的证明。设 $S$ 是对象 $U$ 的覆盖筛。设
$\varphi_i \in L\mathcal{F}(U)$（$i = 1, 2$）在
$\Mor_{\textit{PSh}(\mathcal{C})}(S, L\mathcal{F})$ 中映到同一元素。
可以找到 $U$ 上的一个覆盖筛 $S'$，使两个 $\varphi_i$ 都由元素
$\varphi_i \in \Mor_{\textit{PSh}(\mathcal{C})}(S', \mathcal{F})$ 表示。
把 $S$ 与 $S'$ 都替换为仍是覆盖筛的 $S' \cap S$（见引理
\ref{sites-lemma-sieves-set}），可假设 $S' = S$。设
$V\in \Ob(\mathcal{C})$，而 $\alpha : V \to U$ 属于 $S(V)$。则
$S \times_U V = h_V$；见引理 \ref{sites-lemma-pullback-sieve-section}。所以
$\varphi_i$ 沿 $V \to U$ 的限制对应于 $\mathcal{F}$ 在 $V$ 上的截面
$s_{i, V, \alpha}$。假设意味着存在 $V$ 的覆盖筛 $S_{V, \alpha}$，使
$s_{i, V, \alpha}$ 在
$\Mor_{\textit{PSh}(\mathcal{C})}(S_{V, \alpha}, \mathcal{F})$
中限制为同一元素。考虑 $U$ 上由下列规则定义的筛 $S''$：
\begin{eqnarray}
\label{sites-equation-S-prime-prime}
(f : T \to U) \in S''(T)
& \Leftrightarrow &
\exists\ V , \ \alpha : V \to U, \ \alpha \in S(V), \nonumber \\
& &
\exists\ g : T \to V, \ g \in S_{V, \alpha}(T), \\
& &
f = \alpha \circ g \nonumber
\end{eqnarray}
由拓扑的公理 (2)，$S''$ 是 $U$ 上的覆盖筛。由构造，
$\varphi_1$ 与 $\varphi_2$ 在
$\Mor_{\textit{PSh}(\mathcal{C})}(S'', L\mathcal{F})$ 中限制为同一元素，
如所需。

\medskip\noindent
概述 (2) 的证明。假设 $\mathcal{F}$ 是 $\mathcal{C}$ 上关于拓扑 $J$
分离的集合预层。设 $S$ 是 $\mathcal{C}$ 的对象 $U$ 的覆盖筛。假设
$\varphi \in \Mor_\mathcal{C}(S, L\mathcal{F})$。需要找到限制为 $\varphi$
的元素 $s \in L\mathcal{F}(U)$。
% P01-ERRATA: P01-E0173 -- S and LF are presheaves, so the Hom must be in PSh(C), not in C.
设 $V\in \Ob(\mathcal{C})$，而 $\alpha : V \to U$ 属于 $S(V)$。值
$\varphi(\alpha) \in L\mathcal{F}(V)$ 由 $V$ 的一个覆盖筛
$S_{V, \alpha}$ 以及预层态射
$\varphi_{V, \alpha} : S_{V, \alpha} \to \mathcal{F}$ 给出。与上面的证明一样，
通过等式 (\ref{sites-equation-S-prime-prime}) 定义 $U$ 上的覆盖筛 $S''$。用下列简单
规则定义
$$
\varphi'' : S'' \longrightarrow \mathcal{F}
$$
：对每个 $f : T \to U$、$f \in S''(T)$，按照等式
(\ref{sites-equation-S-prime-prime}) 选取 $V, \alpha, g$，然后令
$$
\varphi''(f) = \varphi_{V, \alpha}(g).
$$
我们断言这与 $V, \alpha, g$ 的选择无关。考虑另一个这样的选择$ V', \alpha', g'$。
% P01-ERRATA: P01-E0174 -- the source omits a space before the math span in "choice$ V'".
$\varphi_{V, \alpha}$ 与 $\varphi_{V', \alpha'}$ 在 $T$ 上下列两个覆盖筛之交上的
限制相同：
$$
(S_{V, \alpha} \times_{V, g} T) \cap (S_{V', \alpha'} \times_{V', g'} T)
$$
事实上，这两个限制都对应于 $\varphi$ 通过 $f$ 在 $T$ 上的限制，而所需等式由
$\mathcal{F}$ 分离得出。把公共限制记为 $\psi$。
% P01-ERRATA: P01-E0175 -- the source's "Denote the common restriction psi" omits "by".
选择无关性由
$\varphi_{V, \alpha}(g) = \psi(\text{id}_T) =
\varphi_{V', \alpha'}(g')$ 得出。于是 $\varphi''$ 给出元素
$s \in L\mathcal{F}(U)$。留给读者检查 $s$ 限制为 $\varphi$。
\end{proof}

\begin{definition}
\label{sites-definition-associated-sheaf-topology}
设范畴 $\mathcal{C}$ 配备拓扑 $J$，而 $\mathcal{F}$ 是
$\mathcal{C}$ 上的集合预层。层 $\mathcal{F}^\# := LL\mathcal{F}$ 连同典范映射
$\mathcal{F} \to \mathcal{F}^\#$，称为{\it 与 $\mathcal{F}$ 关联的层}。
\end{definition}

\begin{proposition}
\label{sites-proposition-sheafification-adjoint-topology}
设范畴 $\mathcal{C}$ 配备一个拓扑，而 $\mathcal{F}$ 是
$\mathcal{C}$ 上的集合预层。典范映射 $\mathcal{F} \to \mathcal{F}^\#$ 具有下列
泛性质：对任意映射 $\mathcal{F} \to \mathcal{G}$，其中
$\mathcal{G}$ 是集合层，存在唯一映射 $\mathcal{F}^\# \to \mathcal{G}$，使
$\mathcal{F} \to \mathcal{F}^\# \to \mathcal{G}$ 等于给定映射。
\end{proposition}

\begin{proof}
与命题 \ref{sites-proposition-sheafification-adjoint} 的证明相同。
\end{proof}

\section{拓扑与层}
\label{sites-section-topology-and-sheaves}

\begin{lemma}
\label{sites-lemma-sieve-sheafification}
设范畴 $\mathcal{C}$ 配备拓扑 $J$，$U$ 是 $\mathcal{C}$ 的对象，而
$S$ 是 $U$ 上的筛。下列条件等价：
\begin{enumerate}
\item 筛 $S$ 是覆盖筛。
\item 映射 $S \to h_U$ 的层化 $S^\# \to h_U^\#$ 是同构。
\end{enumerate}
\end{lemma}

\begin{proof}
先作两点一般说明。将使用 $S^\# = LLS$ 与 $h_U^\# = LLh_U$。特别地，
由引理 \ref{sites-lemma-L-presheaf}，$S^\# \to h_U^\#$ 是单射。注意
$\text{id}_U \in h_U(U)$。因此，它产生 $Lh_U$ 以及
$h_U^\# = LLh_U$ 在 $U$ 上的截面，我们也将这些截面记为 $\text{id}_U$。

\medskip\noindent
假设 $S$ 是覆盖筛。显然，只须找到态射 $h_U \to S^\#$，使其与前述单射的复合
$h_U \to h_U^\#$ 是典范映射。为找到这种映射，
只须找到限制为 $\text{id}_U$ 的截面 $s \in S^\#(U)$。但由于 $S$ 是覆盖筛，
元素 $\text{id}_S \in \Mor_{\textit{PSh}(\mathcal{C})}(S, S)$
产生 $LS$ 在 $U$ 上的一个截面，它在 $Lh_U$ 中限制为 $\text{id}_U$。
所以结论成立。

\medskip\noindent
假设 $S^\# \to h_U^\#$ 是同构。设 $1 \in S^\#(U)$ 是与
$h_U^\#(U)$ 中的 $\text{id}_U$ 对应的元素。由于 $S^\# = LLS$，存在
$U$ 上的覆盖筛 $S'$，使 $1$ 来自
$$
\varphi \in \Mor_{\textit{PSh}(\mathcal{C})}(S', LS).
$$
这又意味着，对每个 $\alpha : V \to U$、$\alpha\in S'(V)$，存在
$V$ 上的覆盖筛 $S_{V, \alpha}$，使 $\varphi(\alpha)$ 对应于预层态射
$S_{V, \alpha} \to S$。换言之，$S_{V, \alpha}$ 包含于
$S \times_U V$。由拓扑的第二条公理，$S$ 是覆盖筛。
\end{proof}

\begin{theorem}
\label{sites-theorem-topology-and-topos}
设 $\mathcal{C}$ 为范畴，而 $J$、$J'$ 是 $\mathcal{C}$ 上的拓扑。下列条件等价：
\begin{enumerate}
\item $J = J'$；
\item 拓扑 $J$ 的层与拓扑 $J'$ 的层相同。
\end{enumerate}
\end{theorem}

\begin{proof}
若 $J = J'$，则层的概念相同，这是恒真式。反过来，引理
\ref{sites-lemma-sieve-sheafification} 用层化函子刻画覆盖筛。而层化函子
$\textit{PSh}(\mathcal{C}) \to \Sh(\mathcal{C}, J)$ 左伴随于包含函子
$\Sh(\mathcal{C}, J) \to \textit{PSh}(\mathcal{C})$。因此，若子范畴
$\Sh(\mathcal{C}, J)$ 与 $\Sh(\mathcal{C}, J')$ 相同，则层化函子相同，
因而覆盖筛族相同。
\end{proof}

\begin{lemma}
\label{sites-lemma-finer-topology}
假设与记号如定理 \ref{sites-theorem-topology-and-topos}。则
$J \subset J'$ 当且仅当拓扑 $J'$ 的每个层都是拓扑 $J$ 的层。
% P01-ERRATA: P01-E0176 -- the source's fragment "Assumption and notation as in" needs a finite verb and plural "Assumptions".
\end{lemma}

\begin{proof}
一个方向是清楚的。对另一个方向，假设
$\Sh(\mathcal{C}, J') \subset \Sh(\mathcal{C}, J)$。由形式论证，这蕴含：
若 $\mathcal{F}$ 是集合预层，而
$\mathcal{F} \to \mathcal{F}^\#$、相应地
$\mathcal{F} \to \mathcal{F}^{\#, \prime}$ 是关于 $J$、相应地 $J'$ 的层化，
则存在典范映射 $\mathcal{F}^\# \to \mathcal{F}^{\#, \prime}$，使
$\mathcal{F} \to \mathcal{F}^\# \to \mathcal{F}^{\#, \prime}$
等于典范映射 $\mathcal{F} \to \mathcal{F}^{\#, \prime}$。当然，
$\mathcal{F}^\# \to \mathcal{F}^{\#, \prime}$ 把第二个层识别为第一个层关于拓扑
$J'$ 的层化。把它应用于引理 \ref{sites-lemma-sieve-sheafification} 的映射
$S \to h_U$，得到交换图
$$
\xymatrix{
S \ar[r] \ar[d] &
S^\# \ar[r] \ar[d] &
S^{\#, \prime} \ar[d] \\
h_U \ar[r] &
h_U^\# \ar[r] &
h_U^{\#, \prime}
}
$$
显然，若 $S$ 是拓扑 $J$ 的覆盖筛，则中间的竖直映射是同构（由该引理），
从而右侧竖直映射也是同构，因为它是中间竖直映射关于 $J'$ 的层化。再次由该引理，
$S$ 也是 $J'$ 的覆盖筛。
\end{proof}

\section{拓扑与连续函子}
\label{sites-section-topologies-continuous-functors}

\noindent
说明连续函子如何在层上给出一对伴随函子。





\section{点与拓扑}
\label{sites-section-points-topologies}

\noindent
回忆节 \ref{sites-section-points}：给定函子
$p = u : \mathcal{C} \to \textit{Sets}$，可以定义茎函子
$$
\textit{PSh}(\mathcal{C}) \longrightarrow \textit{Sets},
\mathcal{F} \longmapsto \mathcal{F}_p.
$$

\begin{definition}
\label{sites-definition-point-topology}
设 $\mathcal{C}$ 为范畴，而 $J$ 是 $\mathcal{C}$ 上的拓扑。该拓扑的一个
{\it 点 $p$} 由满足下列条件的函子
$u : \mathcal{C} \to \textit{Sets}$ 给出：
\begin{enumerate}
\item 对 $U$ 上的每个覆盖筛 $S$，映射
$S_p \to (h_U)_p$ 是满射。
\item 茎函子 $\Sh(\mathcal{C}) \to \textit{Sets}$，
$\mathcal{F} \to \mathcal{F}_p$，是正合的。
\end{enumerate}
% P01-ERRATA: P01-E0177 -- the source uses \to rather than \mapsto for the action F to F_p, unlike the same stalk functor displayed immediately above.
\end{definition}
% Generated mechanically from the frozen stacks.tex target.
% Do not edit this generated file; edit the bound locale source instead.

% OK, start here.
%

\setcounter{chapter}{7}
\chapter{栈}



\phantomsection
\label{stacks-section-phantom}


\section{引言}
\label{stacks-section-introduction}

\noindent
在这一很短的章中，我们引入栈与群胚栈。见 \cite{DM} 与 \cite{Vis2}。


\section{与纤维化范畴相关的态射预层}
\label{stacks-section-morphisms}

\noindent
令 $\mathcal{C}$ 为范畴。令 $p : \mathcal{S} \to \mathcal{C}$ 为纤维化范畴，
见 Categories, Section \ref{categories-section-fibred-categories}。
设 $x, y\in \Ob(\mathcal{S}_U)$ 是 $U$ 上方纤维范畴中的对象。
我们将定义函子
$$
\mathit{Mor}(x, y) : (\mathcal{C}/U)^{opp} \longrightarrow \textit{Sets}.
$$
换言之，这是 $\mathcal{C}/U$ 上的预层，见
Sites, Definition \ref{sites-definition-presheaf}。
依照 Categories,
Definition \ref{categories-definition-pullback-functor-fibred-category}
选定拉回。对 $f : V \to U$，置
$$
\mathit{Mor}(x, y)(f : V \to U) =
\Mor_{\mathcal{S}_V}(f^\ast x, f^\ast y).
$$
令 $f' : V' \to U$ 为 $\mathcal{C}/U$ 的第二个对象。
还需定义与 $\mathcal{C}/U$ 中态射
$g : V'/U  \to V/U$ 对应的限制映射；换言之，
$g : V' \to V$ 且 $f' = f \circ g$。该映射为
$$
\Mor_{\mathcal{S}_V}(f^\ast x, f^\ast y)
\longrightarrow
\Mor_{\mathcal{S}_{V'}}({f'}^\ast x, {f'}^\ast y), \quad
\phi \longmapsto \phi|_{V'}
$$
它基本上就是 $g^\ast$，但这会把左边的元素 $\phi$
变为
$\Mor_{\mathcal{S}_{V'}}(g^\ast f^\ast x, g^\ast f^\ast y)$
中的元素 $g^\ast \phi$。此时使用 Categories,
Lemma \ref{categories-lemma-fibred} 中的变换 $\alpha_{g, f}$。
用公式表示，该限制映射为
$$
\phi|_{V'} =
(\alpha_{g, f})_y^{-1} \circ
g^\ast \phi \circ
(\alpha_{g, f})_x.
$$
当然，实际使用时没有人会以这种方式思考该限制映射。
我们只在此处这样做一次，以验证下列引理。

\begin{lemma}
\label{stacks-lemma-painful}
上述构造确实给出一个预层。
\end{lemma}

\begin{proof}
令 $g : V'/U \to V/U$ 如上，并令
$g' : V''/U \to V'/U$ 同样为 $\mathcal{C}/U$ 中的态射。
于是 $f' = f \circ g$，且
$f'' = f' \circ g' = f \circ g \circ g'$。
令 $\phi \in \Mor_{\mathcal{S}_V}(f^\ast x, f^\ast y)$。则
\begin{eqnarray*}
& &
(\alpha_{g \circ g', f})_y^{-1} \circ
(g \circ g')^\ast \phi \circ
(\alpha_{g \circ g', f})_x
\\
& = &
(\alpha_{g \circ g', f})_y^{-1} \circ
(\alpha_{g', g})_{f^*y}^{-1} \circ
(g')^*g^\ast \phi \circ
(\alpha_{g', g})_{f^*x} \circ
(\alpha_{g \circ g', f})_x
\\
& = &
(\alpha_{g', f'})_y^{-1} \circ
(g')^*(\alpha_{g, f})_y^{-1} \circ
(g')^* g^\ast \phi \circ
(g')^*(\alpha_{g, f})_x
\circ
(\alpha_{g', f'})_x
\\
& = &
(\alpha_{g', f'})_y^{-1} \circ
(g')^*\Big(
(\alpha_{g, f})_y^{-1} \circ
g^\ast \phi \circ
(\alpha_{g, f})_x
\Big) \circ
(\alpha_{g', f'})_x
\end{eqnarray*}
这正是所需的等式
$\phi|_{V''} = (\phi|_{V'})|_{V''}$。
第一个等式成立，是因为 $\alpha_{g', g}$ 是函子之间的变换，因而
$$
\xymatrix{
(g \circ g')^*f^*x
\ar[rr]_{(g \circ g')^\ast \phi}
\ar[d]_{(\alpha_{g', g})_{f^*x}} & &
(g \circ g')^*f^*y
\ar[d]^{(\alpha_{g', g})_{f^*y}} \\
(g')^*g^*f^*x
\ar[rr]^{(g')^*g^\ast \phi} & &
(g')^*g^*f^*y
}
$$
交换。第二个等式由伪函子的性质 (d) 得出，因为
$f' = f \circ g$（见 Categories,
Definition \ref{categories-definition-functor-into-2-category}）。
最后一个等式来自 $(g')^*$ 是函子这一事实。
\end{proof}

\noindent
从现在起，我们常常不再明说变换 $\alpha_{g, f}$，而直接等同函子
$g^* \circ f^*$ 与 $(f \circ g)^*$。特别地，给定
$g : V'/U \to V/U$，预层 $\mathit{Mor}(x, y)$ 的限制映射
有时记作 $\phi \mapsto g^*\phi$。我们用一个定义形式化该构造。

\begin{definition}
\label{stacks-definition-mor-presheaf}
令 $\mathcal{C}$ 为范畴。令
$p : \mathcal{S} \to \mathcal{C}$ 为纤维化范畴，见
Categories, Section \ref{categories-section-fibred-categories}。
给定 $\mathcal{C}$ 的对象 $U$ 以及纤维范畴的对象
$x$、$y$，
所谓{\it 从 $x$ 到 $y$ 的态射预层}，是上面描述的预层
$$
(f : V \to U) \longmapsto \Mor_{\mathcal{S}_V}(f^*x, f^*y)
$$
记为 $\mathit{Mor}(x, y)$。其子预层 $\mathit{Isom}(x, y)$
在 $V$ 上的值是纤维范畴 $\mathcal{S}_V$ 中所有同构
$f^*x \to f^*y$ 所成的集合，称为{\it 从 $x$ 到 $y$ 的同构预层}。
\end{definition}

\noindent
若 $\mathcal{S}$ 以群胚为纤维，则当然有
$\mathit{Isom}(x, y) = \mathit{Mor}(x, y)$，并且通常使用
$\mathit{Isom}$ 这一记号。

\begin{lemma}
\label{stacks-lemma-presheaf-mor-map-fibred-categories}
令 $F : \mathcal{S}_1 \to \mathcal{S}_2$ 为范畴 $\mathcal{C}$
上纤维化范畴的 $1$-态射。令
$U \in \Ob(\mathcal{C})$ 且
$x, y\in \Ob((\mathcal{S}_1)_U)$。则 $F$ 在
$\mathcal{C}/U$ 上定义预层的典范态射
$$
\mathit{Mor}_{\mathcal{S}_1}(x, y)
\longrightarrow
\mathit{Mor}_{\mathcal{S}_2}(F(x), F(y))
$$
\end{lemma}

\begin{proof}
由 Categories,
Definition \ref{categories-definition-fibred-categories-over-C}，
函子 $F$ 把强笛卡尔态射映为强笛卡尔态射。因此，若
$f : V \to U$ 是 $\mathcal{C}$ 中的态射，则存在典范同构
$\alpha_V : f^*F(x) \to F(f^*x)$ 与
$\beta_V : f^*F(y) \to F(f^*y)$，使得
$f^*F(x) \to F(f^*x) \to F(x)$ 是典范态射
$f^*F(x) \to F(x)$；对 $\beta_V$ 亦然。因此可用
$\phi \mapsto \beta_V^{-1} \circ F(\phi) \circ \alpha_V$
定义
$$
\xymatrix{
\mathit{Mor}_{\mathcal{S}_1}(x, y)(f : V \to U) \ar@{=}[r] &
\Mor_{\mathcal{S}_{1, V}}(f^\ast x, f^\ast y) \ar[d] \\
\mathit{Mor}_{\mathcal{S}_2}(F(x), F(y))(f : V \to U) \ar@{=}[r] &
\Mor_{\mathcal{S}_{2, V}}(f^\ast F(x), f^\ast F(y))
}
$$
略去其与限制映射相容的验证。
\end{proof}

\begin{remark}
\label{stacks-remark-alternative}
设 $p : \mathcal{S} \to \mathcal{C}$ 以群胚为纤维。
此时可证明 Lemma \ref{stacks-lemma-painful}；所用的是 Categories,
Lemma \ref{categories-lemma-fibred-strict}，后者说明，
$\mathcal{S} \to \mathcal{C}$ 等价于反变函子
$F : \mathcal{C} \to \textit{Groupoids}$ 所对应的范畴。
对与 $F$ 对应的纤维化范畴，严格地有
$g^* \circ f^* = (f \circ g)^*$，无需使用映射
$\alpha_{g, f}$。此时该引理（更为）平凡。
当然，这还用到了：换成等价的纤维化范畴时，
$\mathit{Mor}(x, y)$ 预层不变；这由
Lemma \ref{stacks-lemma-presheaf-mor-map-fibred-categories} 得出。
\end{remark}

\begin{lemma}
\label{stacks-lemma-isom-as-2-fibre-product}
令 $\mathcal{C}$ 为范畴。令
$p : \mathcal{S} \to \mathcal{C}$ 为纤维化范畴，见
Categories, Section \ref{categories-section-fibred-categories}。
令 $U \in \Ob(\mathcal{C})$ 且
$x, y \in \Ob(\mathcal{S}_U)$。
% P01-ERRATA: P01-E0178 -- "Denote x, y ... also the corresponding 1-morphisms" needs "also by" or recasting.
仍以 $x, y : \mathcal{C}/U \to \mathcal{S}$ 记相应的 $1$-态射，
见 Categories, Lemma \ref{categories-lemma-yoneda-2category}。则
\begin{enumerate}
\item $2$-纤维积
$\mathcal{S} \times_{\mathcal{S} \times \mathcal{S}, (x, y)} \mathcal{C}/U$
在 $\mathcal{C}/U$ 上以集合胚为纤维；并且
\item $\mathit{Isom}(x, y)$ 是与该集合胚纤维化范畴对应的集合预层，
见 Categories, Lemma \ref{categories-lemma-2-category-fibred-setoids}。
\end{enumerate}
\end{lemma}

\begin{proof}
略。提示：该 $2$-纤维积的对象是
$(a : V \to U, z, (\alpha, \beta))$，其中
$\alpha : z \to a^*x$ 与 $\beta : z \to a^*y$
是 $\mathcal{S}_V$ 中的同构。因此，把这样的对象指派给同构
$\beta \circ \alpha^{-1}$，便得到它与
$\mathit{Isom}(x, y)$ 之间的关系。
\end{proof}





\section{纤维化范畴中的下降数据}
\label{stacks-section-descent-data}

\noindent
本节在纤维化范畴的抽象框架中定义下降数据的概念。
在此之前先指出，这与拟凝聚层的下降数据
（Descent, Section \ref{descent-section-equivalence}）
以及概形之上的概形的下降数据
（Descent, Section \ref{descent-section-descent-datum}）
完全类似。

\medskip\noindent
我们采用如下约定：投影映射
$\text{pr}_i : X \times \ldots \times X \to X$
的下标从 $i = 0$ 开始。因此有
$\text{pr}_0, \text{pr}_1 : X \times X  \to X$、
$\text{pr}_0, \text{pr}_1, \text{pr}_2 : X \times X \times X  \to X$
等等。

\begin{definition}
\label{stacks-definition-descent-data}
令 $\mathcal{C}$ 为范畴。令
$p : \mathcal{S} \to \mathcal{C}$ 为纤维化范畴。
依照 Categories,
Definition \ref{categories-definition-pullback-functor-fibred-category}
选定拉回。令
$\mathcal{U} = \{f_i : U_i \to U\}_{i \in I}$
为 $\mathcal{C}$ 中的态射族。假设所有纤维积
$U_i \times_U U_j$ 与
$U_i \times_U U_j \times_U U_k$ 均存在。
\begin{enumerate}
\item 所谓 $\mathcal{S}$ 中{\it 相对于族
$\{f_i : U_i \to U\}$ 的下降数据
$(X_i, \varphi_{ij})$}，由下列资料给出：对每个 $i \in I$，
给定 $\mathcal{S}_{U_i}$ 的对象 $X_i$；对每对
$(i, j) \in I^2$，给定 $\mathcal{S}_{U_i \times_U U_j}$
中的同构
$\varphi_{ij} : \text{pr}_0^*X_i \to \text{pr}_1^*X_j$；
并要求对每组三个下标 $(i, j, k) \in I^3$，图表
$$
\xymatrix{
\text{pr}_0^*X_i \ar[rd]_{\text{pr}_{01}^*\varphi_{ij}}
\ar[rr]_{\text{pr}_{02}^*\varphi_{ik}} & &
\text{pr}_2^*X_k \\
& \text{pr}_1^*X_j \ar[ru]_{\text{pr}_{12}^*\varphi_{jk}} &
}
$$
在范畴
$\mathcal{S}_{U_i \times_U U_j \times_U U_k}$ 中交换。
这称为{\it 余循环条件}。
\item 所谓{\it 下降数据的态射
$\psi : (X_i, \varphi_{ij}) \to
(X'_i, \varphi'_{ij})$}，由态射族
$\psi = (\psi_i)_{i\in I}$ 给出，其中
$\psi_i : X_i \to X'_i$ 是 $\mathcal{S}_{U_i}$ 中的态射，
并且所有图表
$$
\xymatrix{
\text{pr}_0^*X_i \ar[r]_{\varphi_{ij}} \ar[d]_{\text{pr}_0^*\psi_i}
& \text{pr}_1^*X_j \ar[d]^{\text{pr}_1^*\psi_j} \\
\text{pr}_0^*X'_i \ar[r]^{\varphi'_{ij}} &
\text{pr}_1^*X'_j \\
}
$$
在范畴 $\mathcal{S}_{U_i \times_U U_j}$ 中交换。
\item 相对于 $\mathcal{U}$ 的下降数据所成的范畴记作
$DD(\mathcal{U})$。
\end{enumerate}
\end{definition}

\noindent
如果每个态射 $f_i : U_i \to U$ 都是{\it 可表的}，则纤维积
$U_i \times_U U_j$ 与
$U_i \times_U U_j \times_U U_k$ 存在；见 Categories,
Definition \ref{categories-definition-representable-morphism}。
回忆，在位点中，覆盖 $\{U_i \to U\}$ 的条件之一是每个态射均可表，
见 Sites, Definition \ref{sites-definition-site} 第 (3) 部分。
事实上，上述定义主要用于 $\mathcal{C}$ 是位点且
$\{U_i \to U\}$ 是 $\mathcal{C}$ 的覆盖的情形。
不过，下降数据只是可以如上定义的抽象构造。这一点很有用：
例如，给定 $\mathcal{C}$ 上的纤维化范畴，可以考察下降数据对哪些族有效，
并尝试把这些族用作位点的覆盖族。

\begin{remarks}
\label{stacks-remarks-definition-descent-datum}
关于 Definition \ref{stacks-definition-descent-data} 有两点说明。
令 $p : \mathcal{S} \to \mathcal{C}$ 为纤维化范畴，并令
$\{f_i : U_i \to U\}_{i \in I}$ 与
$(X_i, \varphi_{ij})$ 如
Definition \ref{stacks-definition-descent-data} 中所述。
\begin{enumerate}
\item 存在对角态射
$\Delta : U_i \to U_i \times_U U_i$。
沿该态射拉回 $\varphi_{ii}$，得到自同构
$\Delta^\ast \varphi_{ii} \in \text{Aut}_{U_i}(X_i)$。
把三元组 $(i, i, i)$ 的余循环条件沿
$\Delta_{123} : U_i \to U_i \times_U U_i \times_U U_i$
拉回，可推出
$\Delta^\ast \varphi_{ii} \circ \Delta^\ast \varphi_{ii} =
\Delta^\ast \varphi_{ii}$；因而
$\Delta^\ast \varphi_{ii} = \text{id}_{X_i}$。
\item 存在态射
$\Delta_{13}: U_i \times_U U_j \to U_i \times_U U_j \times_U U_i$。
把三元组 $(i, j, i)$ 的余循环条件沿该态射拉回，得到恒等式
$(\sigma^\ast \varphi_{ji}) \circ \varphi_{ij} =
\text{id}_{\text{pr}_0^\ast X_i}$，其中
$\sigma : U_i \times_U U_j \to U_j \times_U U_i$
是交换态射。
\end{enumerate}
\end{remarks}

\begin{lemma}
\label{stacks-lemma-pullback}
（下降数据的拉回。）
令 $\mathcal{C}$ 为范畴。令
$p : \mathcal{S} \to \mathcal{C}$ 为纤维化范畴。
% P01-ERRATA: P01-E0179 -- "Make a choice pullbacks" is missing "of".
依照 Categories,
Definition \ref{categories-definition-pullback-functor-fibred-category}
选定拉回。令
$\mathcal{U} = \{f_i : U_i \to U\}_{i \in I}$ 与
$\mathcal{V} = \{V_j \to V\}_{j \in J}$
为 $\mathcal{C}$ 中具有固定目标的态射族。
% P01-ERRATA: P01-E0180 -- the source says "be a families"; the plural subject requires "be families".
假设所有纤维积
$U_i \times_U U_{i'}$、
$U_i \times_U U_{i'} \times_U U_{i''}$、
$V_j \times_V V_{j'}$ 与
$V_j \times_V V_{j'} \times_V V_{j''}$ 均存在。
令 $\alpha : I \to J$、$h : U \to V$ 以及
$g_i : U_i \to V_{\alpha(i)}$ 给出具有固定目标的映射族之间的态射，
见 Sites, Definition \ref{sites-definition-morphism-coverings}。
\begin{enumerate}
\item 令 $(Y_j, \varphi_{jj'})$ 为相对于族
$\{V_j \to V\}$ 的下降数据。系统
$$
\left(
g_i^*Y_{\alpha(i)},
(g_i \times g_{i'})^*\varphi_{\alpha(i)\alpha(i')}
\right)
$$
是相对于 $\mathcal{U}$ 的下降数据。
\item 该构造定义从相对于 $\mathcal{V}$ 的下降数据到相对于
$\mathcal{U}$ 的下降数据的函子。
\item 再给定 $\alpha' : I \to J$、$h' : U \to V$ 以及
$g'_i : U_i \to V_{\alpha'(i)}$ 所给出的具有固定目标的映射族之间的态射；
% P01-ERRATA: P01-E0181 -- the source omits an article before "morphism of families" in the second choice.
若 $h = h'$，则由两组资料所得的下降数据函子典范同构。
\end{enumerate}
\end{lemma}

\begin{proof}
略。
\end{proof}

\begin{definition}
\label{stacks-definition-pullback-functor}
令 $\mathcal{U} = \{U_i \to U\}_{i \in I}$、
$\mathcal{V} = \{V_j \to V\}_{j \in J}$、
$\alpha : I \to J$、$h : U \to V$ 以及
$g_i : U_i \to V_{\alpha(i)}$ 如
Lemma \ref{stacks-lemma-pullback} 中所述。该引理中构造的函子
$$
(Y_j, \varphi_{jj'}) \longmapsto
(g_i^*Y_{\alpha(i)}, (g_i \times g_{i'})^*\varphi_{\alpha(i)\alpha(i')})
$$
称为下降数据上的{\it 拉回函子}。
\end{definition}

\noindent
给定 $h : U \to V$，若存在覆盖 $h$ 的态射
$\tilde h : \mathcal{U} \to \mathcal{V}$，则由
Lemma \ref{stacks-lemma-pullback}，$\tilde h^*$ 与
$\tilde h$ 的选择无关。因此，我们有时直接以 $h^*$ 表示拉回函子。

\begin{definition}
\label{stacks-definition-effective-descent-datum}
令 $\mathcal{C}$ 为范畴。令
$p : \mathcal{S} \to \mathcal{C}$ 为纤维化范畴。
依照 Categories,
Definition \ref{categories-definition-pullback-functor-fibred-category}
选定拉回。令
$\mathcal{U} = \{f_i : U_i \to U\}_{i \in I}$
为目标为 $U$ 的态射族。假设所有纤维积
$U_i \times_U U_j$ 与
$U_i \times_U U_j \times_U U_k$ 均存在。
\begin{enumerate}
\item 给定 $\mathcal{S}_U$ 的对象 $X$，所谓{\it 平凡下降数据}，
是相对于族 $\{\text{id}_U : U \to U\}$ 的下降数据
$(X, \text{id}_X)$。
\item 给定 $\mathcal{S}_U$ 的对象 $X$，把平凡下降数据
$(X, \text{id}_X)$ 沿显然的映射
$\{f_i : U_i \to U\} \to \{\text{id}_U : U \to U\}$
拉回，就在对象族 $f_i^*X$ 上得到一个{\it 典范下降数据}。
记该下降数据为 $(f_i^*X, can)$。
\item 相对于 $\{f_i : U_i \to U\}$ 的下降数据
$(X_i, \varphi_{ij})$ 称为{\it 有效的}，如果存在
$\mathcal{S}_U$ 的对象 $X$，使
$(X_i, \varphi_{ij})$ 同构于 $(f_i^*X, can)$。
\end{enumerate}
\end{definition}

\noindent
注意，把 $X \in \mathcal{S}_U$ 指派到它相对于
$\mathcal{U}$ 的典范下降数据，这一规则定义函子
$$
\mathcal{S}_U \longrightarrow DD(\mathcal{U}).
$$
下降数据有效，当且仅当它属于该函子的本质像。
下面具体写出典范下降数据。

\begin{lemma}
\label{stacks-lemma-trivial-cocycle}
在 Definition \ref{stacks-definition-effective-descent-datum} 第 (2) 部分的情形下，
映射 $can_{ij} : \text{pr}_0^*f_i^*X \to \text{pr}_1^*f_j^*X$
等于
$(\alpha_{\text{pr}_1, f_j})_X \circ
(\alpha_{\text{pr}_0, f_i})_X^{-1}$，其中
$\alpha_{\cdot, \cdot}$ 如 Categories,
Lemma \ref{categories-lemma-fibred} 中所述，并且我们使用映射
$U_i \times_U U_j \to U$ 之间的等式
$f_i \circ \text{pr}_0 = f_j \circ \text{pr}_1$。
\end{lemma}

\begin{proof}
略。
\end{proof}

\begin{lemma}
\label{stacks-lemma-compare-descent-condition}
令 $\mathcal{C}$ 为范畴。令
$
\mathcal{V} = \{V_j \to U\}_{j \in J} \to
\mathcal{U} = \{U_i \to U\}_{i \in I}
$
为 $\mathcal{C}$ 中具有固定目标的映射族之间的态射，它由
$\text{id} : U \to U$、
$\alpha : J \to I$ 以及
$f_j : V_j \to U_{\alpha(j)}$ 给出。令
$p : \mathcal{S} \to \mathcal{C}$ 为纤维化范畴。若
\begin{enumerate}
\item 对 $0 \leq p \leq 3$ 与 $0 \leq q \leq 3$，当
$p + q \geq 2$，并且
$i_1, \ldots, i_p \in I$、$j_1, \ldots, j_q \in J$ 时，纤维积
$U_{i_1} \times_U \ldots \times_U U_{i_p} \times_U
V_{j_1} \times_U \ldots \times_U V_{j_q}$ 存在；
\item 函子 $\mathcal{S}_U \to DD(\mathcal{V})$ 是等价；
\item 对每个 $i \in I$，函子
$\mathcal{S}_{U_i} \to DD(\mathcal{V}_i)$ 是全忠实的；且
\item 对每个 $i, i' \in I$，函子
$\mathcal{S}_{U_i \times_U U_{i'}} \to DD(\mathcal{V}_{ii'})$
是忠实的。
\end{enumerate}
这里
$\mathcal{V}_i = \{U_i \times_U V_j \to U_i\}_{j \in J}$，而
$\mathcal{V}_{ii'} =
\{U_i \times_U U_{i'} \times_U V_j \to
U_i \times_U U_{i'}\}_{j \in J}$。
则 $\mathcal{S}_U \to DD(\mathcal{U})$ 是等价。
\end{lemma}

\begin{proof}
条件 (1) 保证有足够的纤维积可供使用，从而上述陈述有意义。
我们先证明函子 $\mathcal{S}_U \to DD(\mathcal{U})$
本质满射。设给定相对于 $\mathcal{U}$ 的下降数据
$(X_i, \varphi_{ii'})$。
由引理 \ref{stacks-lemma-pullback}，可将其拉回为相对于
$\mathcal{V}$ 的下降数据 $(X_j, \varphi_{jj'})$。
由假设 (2)，该下降数据有效，因而存在
$\mathcal{S}_U$ 中的对象 $X$，使得平凡下降数据
$(X, \text{id}_X)$ 沿态射
$\mathcal{V} \to \{U \to U\}$ 的拉回同构于
$(X_j, \varphi_{jj'})$。其次，注意有图式
$$
\xymatrix{
\mathcal{V}_i \ar[r] \ar[d] &
\mathcal{V} \ar[r] &
\mathcal{U} \ar[d] \\
\{U_i \to U_i\} \ar[rr] \ar[rru] & &
\{U \to U\}
}
$$
其箭头是 $\mathcal{C}$ 中具有固定目标的映射族之间的态射。
这个图式并不交换，但由引理 \ref{stacks-lemma-pullback}，
由此得到的下降数据拉回函子典范同构。因此，
$(X, \text{id}_X)$ 与 $(X_i, \text{id}_{X_i})$
在 $DD(\mathcal{V}_i)$ 中的拉回是同构对象。
由假设 (3)，遂得到范畴 $\mathcal{S}_{U_i}$ 中的同构
$(U_i \to U)^*X \to X_i$。
我们略去验证这些箭头与态射 $\varphi_{ii'}$ 相容；
提示：使用条件 (4) 中函子的忠实性。
我们也略去验证函子
$\mathcal{S}_U \to DD(\mathcal{U})$ 全忠实。
\end{proof}













\section{栈}
\label{stacks-section-definition}

\noindent
下面给出栈的定义。它把纤维化范畴的概念与下降的概念结合起来。

\begin{definition}
\label{stacks-definition-stack}
令 $\mathcal{C}$ 为位点。$\mathcal{C}$ 上的一个{\it 栈}
是 $\mathcal{C}$ 上的范畴 $p : \mathcal{S} \to \mathcal{C}$，
并满足下列条件：
\begin{enumerate}
\item $p : \mathcal{S} \to \mathcal{C}$ 是纤维化范畴，见
Categories, Definition \ref{categories-definition-fibred-category}；
\item 对任意 $U \in \Ob(\mathcal{C})$ 及任意
$x, y \in \mathcal{S}_U$，预层 $\mathit{Mor}(x, y)$（见
Definition \ref{stacks-definition-mor-presheaf}）是位点
$\mathcal{C}/U$ 上的层；且
\item 对位点 $\mathcal{C}$ 的任意覆盖
$\mathcal{U} = \{f_i : U_i \to U\}_{i \in I}$，
$\mathcal{S}$ 中相对于 $\mathcal{U}$ 的任意下降数据都是有效的。
\end{enumerate}
\end{definition}

\noindent
我们认为，上述表述是理解栈最方便的方式。具体而言，给定
$\mathcal{C}$ 上的一个范畴，要验证它是栈，就依次检验性质
(1)、(2) 与 (3)。如果该范畴不是纤维化的，性质 (2) 和 (3)
当然没有意义。没有 (2)，就无法证明 (3) 中的下降在唯一同构意义下唯一，
也无法证明其函子性。

\medskip\noindent
下一个引理给出一个等价定义。

\begin{lemma}
\label{stacks-lemma-stack-equivalences}
令 $\mathcal{C}$ 为位点。
令 $p : \mathcal{S} \to \mathcal{C}$ 为
$\mathcal{C}$ 上的纤维化范畴。下列条件等价：
\begin{enumerate}
\item $\mathcal{S}$ 是 $\mathcal{C}$ 上的栈；且
\item 对位点 $\mathcal{C}$ 的任意覆盖
$\mathcal{U} = \{f_i : U_i \to U\}_{i \in I}$，将对象送到其
典范下降数据的函子
$$
\mathcal{S}_U \longrightarrow DD(\mathcal{U})
$$
是等价。
\end{enumerate}
\end{lemma}

\begin{proof}
略。
\end{proof}

\begin{lemma}
\label{stacks-lemma-substack}
令 $p : \mathcal{S} \to \mathcal{C}$ 为位点
$\mathcal{C}$ 上的栈。令 $\mathcal{S}'$ 为
$\mathcal{S}$ 的子范畴。假设
\begin{enumerate}
\item 若 $\varphi : y \to x$ 是 $\mathcal{S}$ 的强笛卡尔态射，
且 $x$ 是 $\mathcal{S}'$ 的对象，则 $y$ 同构于
$\mathcal{S}'$ 的某个对象；
\item $\mathcal{S}'$ 是 $\mathcal{S}$ 的满子范畴；且
\item 若 $\{f_i : U_i \to U\}$ 是 $\mathcal{C}$ 的覆盖，
$x$ 是 $\mathcal{S}$ 中位于 $U$ 上的对象，并且对每个 $i$，
$f_i^*x$ 都同构于 $\mathcal{S}'$ 的某个对象，则
$x$ 同构于 $\mathcal{S}'$ 的某个对象。
\end{enumerate}
那么 $\mathcal{S}' \to \mathcal{C}$ 是栈。
\end{lemma}

\begin{proof}
略。提示：
第一个条件保证 $\mathcal{S}'$ 是纤维化范畴。
% P01-E0182: source-side wording candidate; canonical meaning is preserved.
第二个条件保证 $\mathcal{S}'$ 的 $\mathit{Isom}$-预层都是层
（因为它们与 $\mathcal{S}$ 中相应的预层相同）。
第三个条件保证 $\mathcal{S}'$ 中满足下降条件：我们可先在
$\mathcal{S}$ 中下降，再由 (3) 得知所得对象同构于
$\mathcal{S}'$ 的某个对象。
\end{proof}

\begin{lemma}
\label{stacks-lemma-stack-equivalent}
令 $\mathcal{C}$ 为位点。
令 $\mathcal{S}_1$、$\mathcal{S}_2$ 为
$\mathcal{C}$ 上的范畴。假设 $\mathcal{S}_1$ 与 $\mathcal{S}_2$ 作为
$\mathcal{C}$ 上的范畴彼此等价。那么
$\mathcal{S}_1$ 是 $\mathcal{C}$ 上的栈，当且仅当
$\mathcal{S}_2$ 是 $\mathcal{C}$ 上的栈。
\end{lemma}

\begin{proof}
令 $F : \mathcal{S}_1 \to \mathcal{S}_2$、
$G : \mathcal{S}_2 \to \mathcal{S}_1$ 为
$\mathcal{C}$ 上的函子，并令
$i : F \circ G \to \text{id}_{\mathcal{S}_2}$、
$j : G \circ F \to \text{id}_{\mathcal{S}_1}$ 为
$\mathcal{C}$ 上函子的同构。由
Categories, Lemma \ref{categories-lemma-fibred-equivalent}，
$\mathcal{S}_1$ 纤维化，当且仅当
$\mathcal{S}_2$ 在 $\mathcal{C}$ 上纤维化。因此可假设
$\mathcal{S}_1$ 与 $\mathcal{S}_2$ 均纤维化。
此外，Categories, Lemma
\ref{categories-lemma-fibred-equivalent} 的证明表明，
$F$ 与 $G$ 把强笛卡尔态射映为强笛卡尔态射；换言之，
$F$ 与 $G$ 是 $\mathcal{C}$ 上纤维化范畴的 $1$-态射。
这意味着，给定 $U \in \Ob(\mathcal{C})$ 及
$x, y \in \mathcal{S}_{1, U}$，下列预层可相互识别：
% P01-E0183: source-side category-subscript candidate; frozen display preserved.
$$
\mathit{Mor}_{\mathcal{S}_1}(x, y),
\mathit{Mor}_{\mathcal{S}_1}(F(x), F(y)) :
(\mathcal{C}/U)^{opp} \longrightarrow \textit{Sets}.
$$
见引理 \ref{stacks-lemma-presheaf-mor-map-fibred-categories}。因此，
第一个是层，当且仅当第二个是层。
最后还需证明：若 $\mathcal{S}_1$ 中的每个下降数据都有效，
则 $\mathcal{S}_2$ 中的每个下降数据也都有效。
为此，令 $(X_i, \varphi_{ii'})$ 为 $\mathcal{S}_2$ 中
相对于位点 $\mathcal{C}$ 的覆盖 $\{U_i \to U\}$ 的下降数据。
% P01-E0184: source-side missing-preposition candidate; meaning preserved.
那么 $(G(X_i), G(\varphi_{ii'}))$ 是 $\mathcal{S}_1$ 中
相对于覆盖 $\{U_i \to U\}$ 的下降数据。
令 $X$ 为 $\mathcal{S}_{1, U}$ 的对象，使下降数据
$(f_i^*X, can)$ 同构于 $(G(X_i), G(\varphi_{ii'}))$。
于是 $F(X)$ 是 $\mathcal{S}_{2, U}$ 的对象，使下降数据
$(f_i^*F(X), can)$ 同构于
$(F(G(X_i)), F(G(\varphi_{ii'})))$；后者又利用 $i$
同构于原下降数据 $(X_i, \varphi_{ii'})$。
\end{proof}

\noindent
$\mathcal{C}$ 上栈的 $2$-范畴定义如下。

\begin{definition}
\label{stacks-definition-stacks-over-C}
令 $\mathcal{C}$ 为位点。{\it $\mathcal{C}$ 上栈的
$2$-范畴}是 $\mathcal{C}$ 上纤维化范畴的
$2$-范畴（见 Categories, Definition
\ref{categories-definition-fibred-categories-over-C}）的如下子
$2$-范畴：
\begin{enumerate}
\item 其对象为栈 $p : \mathcal{S} \to \mathcal{C}$。
\item 其 $1$-态射
$(\mathcal{S}, p) \to (\mathcal{S}', p')$ 为函子
$G : \mathcal{S} \to \mathcal{S}'$，满足
$p' \circ G = p$，且 $G$ 把强笛卡尔态射映为强笛卡尔态射。
\item 对
$G, H : (\mathcal{S}, p) \to (\mathcal{S}', p')$，
其 $2$-态射 $t : G \to H$ 为函子之间的态射，并且对所有
$x \in \Ob(\mathcal{S})$ 都有
$p'(t_x) = \text{id}_{p(x)}$。
\end{enumerate}
\end{definition}

\begin{lemma}
\label{stacks-lemma-2-product-stacks}
令 $\mathcal{C}$ 为位点。$\mathcal{C}$ 上栈的
$(2, 1)$-范畴具有 2-纤维积，其描述与
Categories, Lemma
\ref{categories-lemma-2-product-categories-over-C} 中相同。
\end{lemma}

\begin{proof}
令 $f : \mathcal{X} \to \mathcal{S}$ 与
$g : \mathcal{Y} \to \mathcal{S}$ 为上述意义下
$\mathcal{C}$ 上栈之间的 $1$-态射。由
Categories, Lemma
\ref{categories-lemma-2-product-categories-over-C} 描述的范畴
$\mathcal{X} \times_\mathcal{S} \mathcal{Y}$，依
Categories, Lemma
\ref{categories-lemma-2-product-fibred-categories-over-C} 是纤维化范畴。
（这里用到了 $f$ 与 $g$ 保持强笛卡尔态射。）
还需证明态射预层是层，并且相对于 $\mathcal{C}$ 的覆盖的下降是有效的。

\medskip\noindent
回忆，$\mathcal{X} \times_\mathcal{S} \mathcal{Y}$ 的对象
由四元组 $(U, x, y, \phi)$ 给出，它位于
$\mathcal{C}$ 的对象 $U$ 上。再令
% P01-E0185: source-side missing-article candidate; meaning preserved.
$(U, x', y', \phi')$ 为位于 $U$ 上的第二个对象。
回忆，$\phi : f(x) \to g(y)$ 与
$\phi' : f(x') \to g(y')$ 是范畴 $\mathcal{S}_U$ 中的同构。
用这些同构来识别
$z = f(x) = g(y)$ 与 $z' = f(x') = g(y')$。
% P01-E0186: source-side agreement candidate; meaning preserved.
在这些识别下，显然作为预层有
$$
\mathit{Mor}((U, x, y, \phi), (U, x', y', \phi'))
=
\mathit{Mor}(x, x')
\times_{\mathit{Mor}(z, z')}
\mathit{Mor}(y, y')
$$
然而，预层范畴中的纤维积保持层（Sites, Lemma
\ref{sites-lemma-limit-sheaf}），故这是一个层。

\medskip\noindent
令 $\mathcal{U} = \{f_i : U_i \to U\}_{i \in I}$ 为位点
$\mathcal{C}$ 的覆盖。令 $(X_i, \chi_{ij})$ 为
$\mathcal{X} \times_\mathcal{S} \mathcal{Y}$ 中相对于
$\mathcal{U}$ 的下降数据。按上述方式写
$X_i = (U_i, x_i, y_i, \phi_i)$，并依范畴
$\mathcal{X} \times_\mathcal{S} \mathcal{Y}$ 的定义写
$\chi_{ij} = (\varphi_{ij}, \psi_{ij})$（见
Categories, Lemma
\ref{categories-lemma-2-product-categories-over-C}）。
显然，$(x_i, \varphi_{ij})$ 是 $\mathcal{X}$ 中的下降数据，
而 $(y_i, \psi_{ij})$ 是 $\mathcal{Y}$ 中的下降数据。
由于 $\mathcal{X}$ 与 $\mathcal{Y}$ 都是栈，这些下降数据有效。
于是得到 $x \in \Ob(\mathcal{X}_U)$ 与
$y \in \Ob(\mathcal{Y}_U)$，它们与下降数据相容，并满足
$x_i = x|_{U_i}$ 及 $y_i = y|_{U_i}$。
令 $z = f(x)$、$z' = g(y)$；二者都是 $\mathcal{S}_U$ 的对象。
态射 $\phi_i$ 是 $\mathit{Isom}(z, z')(U_i)$ 的元素，并满足
$\phi_i|_{U_i \times_U U_j} = \phi_j|_{U_i \times_U U_j}$。
因此，由 $\mathit{Isom}(z, z')$ 的层性质，得到同构
$\phi : z = f(x) \to z' = g(y)$。
我们略去验证，与
$(\mathcal{X} \times_\mathcal{S} \mathcal{Y})_U$ 的对象
$(U, x, y, \phi)$ 相联系的典范下降数据同构于起初给定的下降数据。
\end{proof}

\begin{lemma}
\label{stacks-lemma-characterize-ff}
令 $\mathcal{C}$ 为位点。
令 $\mathcal{S}_1$、$\mathcal{S}_2$ 为
$\mathcal{C}$ 上的栈。令
$F : \mathcal{S}_1 \to \mathcal{S}_2$ 为 $1$-态射。
下列条件等价：
\begin{enumerate}
\item $F$ 全忠实；
\item 对每个 $U \in \Ob(\mathcal{C})$ 以及每个
$x, y \in \Ob(\mathcal{S}_{1, U})$，映射
$$
F :
\mathit{Mor}_{\mathcal{S}_1}(x, y)
\longrightarrow
\mathit{Mor}_{\mathcal{S}_2}(F(x), F(y))
$$
是位点 $\mathcal{C}/U$ 上层之间的同构。
\end{enumerate}
\end{lemma}

\begin{proof}
假设 (1)。对 (2) 中的 $U, x, y$，上述映射 $F$
在位于 $U$ 上的 $\mathcal{C}$ 的对象 $V$ 处取值为如下映射：
$$
\begin{aligned}
F : \Mor_{\mathcal{S}_{1, V}}(x|_V, y|_V) &\to \\
& \Mor_{\mathcal{S}_{2, V}}(F(x|_V), F(y|_V))
\end{aligned}
$$
由于 $F$ 全忠实，相应映射
$\Mor_{\mathcal{S}_1}(x|_V, y|_V) \to
\Mor_{\mathcal{S}_2}(F(x|_V), F(y|_V))$
为双射。纤维范畴 $\mathcal{S}_{1, V}$ 中的态射，恰为
$\mathcal{S}_1$ 中 $x|_V$ 与 $y|_V$ 之间位于
$\text{id}_V$ 上的态射。类似地，纤维范畴
$\mathcal{S}_{2, V}$ 中的态射，恰为 $\mathcal{S}_2$ 中
$F(x|_V)$ 与 $F(y|_V)$ 之间位于 $\text{id}_V$ 上的态射。
因此，$F$ 也在这两者之间诱导双射，故 (2) 成立。

\medskip\noindent
假设 (2)。设给定 $\mathcal{C}$ 的对象 $U$、$V$，
以及 $x \in \Ob(\mathcal{S}_{1, U})$ 和
$y \in \Ob(\mathcal{S}_{1, V})$。要证明 $F$ 全忠实，
只需证明它在位于某个固定 $f : U \to V$ 上的态射之间诱导双射。
在 $\mathcal{S}_1$ 中选取位于 $f$ 上的强笛卡尔态射
$f^*y \to y$。于是，$\mathcal{S}_1$ 中位于 $f$ 上的
态射 $x \to y$ 的集合与
$\Mor_{\mathcal{S}_{1, U}}(x, f^*y)$ 之间有双射。
由于 $F$ 作为 $\mathcal{C}$ 上栈的 $2$-范畴中的
$1$-态射保持强笛卡尔态射，$\mathcal{S}_2$ 中位于
$f$ 上的态射 $F(x) \to F(y)$ 的集合与
$\Mor_{\mathcal{S}_{2, U}}(F(x), F(f^*y))$ 之间也有双射。
又因 $F$ 诱导双射
$\Mor_{\mathcal{S}_{1, U}}(x, f^*y) \to
\Mor_{\mathcal{S}_{2, U}}(F(x), F(f^*y))$，
可知 (1) 成立。
\end{proof}

\begin{lemma}
\label{stacks-lemma-characterize-essentially-surjective-when-ff}
令 $\mathcal{C}$ 为位点。
令 $\mathcal{S}_1$、$\mathcal{S}_2$ 为
$\mathcal{C}$ 上的栈。令
$F : \mathcal{S}_1 \to \mathcal{S}_2$ 为全忠实的
$1$-态射。下列条件等价：
\begin{enumerate}
\item $F$ 是等价；
\item 对每个 $U \in \Ob(\mathcal{C})$ 及每个
$x \in \Ob(\mathcal{S}_{2, U})$，存在覆盖
$\{f_i : U_i \to U\}$，使得 $f_i^*x$ 属于函子
$F : \mathcal{S}_{1, U_i} \to \mathcal{S}_{2, U_i}$
的本质像。
\end{enumerate}
\end{lemma}

\begin{proof}
(1) $\Rightarrow$ (2) 是立即的。
为证明 (2) 蕴含 (1)，需证明 (2) 中的每个
$x$ 均属于函子 $F$ 的本质像。
为此，选取 (2) 中的覆盖、对象
$x_i \in \Ob(\mathcal{S}_{1, U_i})$ 以及同构
$\varphi_i : F(x_i) \to f_i^*x$。由下式取箭头
$$
\varphi_{ij} :
x_i|_{U_i \times_U U_j}
\longrightarrow
x_j|_{U_i \times_U U_j}
$$
使 $F(\varphi_{ij}) = \varphi_j^{-1} \circ \varphi_i$，
便得到 $\mathcal{S}_1$ 中相对于
$\{f_i : U_i \to U\}$ 的下降数据。
由栈的公理，该下降数据有效，故得到
$\mathcal{S}_1$ 中位于 $U$ 上的对象 $x_1$。
我们略去验证 $F(x_1)$ 在 $U$ 上同构于 $x$。
\end{proof}

\begin{remark}
\label{stacks-remark-stack-make-small}
（缩减一个“大的”栈以得到栈。）
令 $\mathcal{C}$ 为位点。假设
% P01-E0187: source-side missing-article candidate; meaning preserved.
$p : \mathcal{S} \to \mathcal{C}$ 是从一个“大的”范畴到
$\mathcal{C}$ 的函子，也就是说，假设
$\mathcal{S}$ 的对象构成真类。最后，假设
$p : \mathcal{S} \to \mathcal{C}$ 满足
Definition \ref{stacks-definition-stack} 的条件 (1)、(2)、(3)。
一般而言，无法用一个等价范畴替换
% P01-E0188: source-side article candidate; meaning preserved.
$p : \mathcal{S} \to \mathcal{C}$ 而得到一个栈。原因是，
% P01-E0189: source-side number-agreement candidate; meaning preserved.
纤维范畴 $\mathcal{S}_U$ 的对象同构类可能构成真类。
另一方面，假设
\begin{enumerate}
\item[(4)] 对每个 $U \in \Ob(\mathcal{C})$，存在集合
$S_U \subset \Ob(\mathcal{S}_U)$，使得
$\mathcal{S}_U$ 的每个对象都在 $\mathcal{S}_U$ 中同构于
$S_U$ 的某个元素。
\end{enumerate}
在这种情况下，可找到 $\mathcal{S}$ 的满子范畴
$\mathcal{S}_{small}$，并令
$p_{small} = p|_{\mathcal{S}_{small}}$，使得
\begin{enumerate}
\item[(a)] 函子
$p_{small} : \mathcal{S}_{small} \to \mathcal{C}$ 定义一个栈；且
\item[(b)] 嵌入 $\mathcal{S}_{small} \to \mathcal{S}$
全忠实且本质满射。
\end{enumerate}
（提示：对每个 $U \in \Ob(\mathcal{C})$，令
$\alpha(U)$ 表示满足
$\Ob(\mathcal{S}_U) \cap V_{\alpha(U)}$ 满射到
$\mathcal{S}_U$ 的对象同构类集合的最小序数，并令
$\alpha = \sup_{U \in \Ob(\mathcal{C})} \alpha(U)$。
然后取
$\Ob(\mathcal{S}_{small}) = \Ob(\mathcal{S}) \cap V_\alpha$。
所用记号见 Sets, Section \ref{sets-section-sets-hierarchy}。）
\end{remark}







\section{群胚中的栈}
\label{stacks-section-stacks-in-groupoids}

\noindent
在各类栈中，纤维化于群胚的栈稍易理解。我们如下重新定义它们。

\begin{definition}
\label{stacks-definition-stack-in-groupoids}
位点 $\mathcal{C}$ 上的一个{\it 群胚中的栈}，是
$\mathcal{C}$ 上的范畴 $p : \mathcal{S} \to \mathcal{C}$，
并满足：
\begin{enumerate}
\item $p : \mathcal{S} \to \mathcal{C}$ 在
$\mathcal{C}$ 上纤维化于群胚（见
Categories, Definition \ref{categories-definition-fibred-groupoids}）；
\item 对所有 $U \in \Ob(\mathcal{C})$ 及所有
$x, y\in \Ob(\mathcal{S}_U)$，预层
$\mathit{Isom}(x, y)$ 是位点 $\mathcal{C}/U$ 上的层；且
\item 对 $\mathcal{C}$ 中的所有覆盖
$\mathcal{U} = \{U_i \to U\}$，相对于
$\mathcal{U}$ 的所有下降数据 $(x_i, \phi_{ij})$ 都有效。
\end{enumerate}
\end{definition}

\noindent
通常最难检验的是第三个条件。下面的引理将此定义与栈的概念作比较。

\begin{lemma}
\label{stacks-lemma-stack-in-groupoids-stack}
令 $\mathcal{C}$ 为位点。
令 $p : \mathcal{S} \to \mathcal{C}$ 为
$\mathcal{C}$ 上的范畴。下列条件等价：
\begin{enumerate}
\item $\mathcal{S}$ 是 $\mathcal{C}$ 上的群胚中的栈；
\item $\mathcal{S}$ 是 $\mathcal{C}$ 上的栈，且所有纤维范畴均为群胚；且
\item $\mathcal{S}$ 在 $\mathcal{C}$ 上纤维化于群胚，
并且是 $\mathcal{C}$ 上的栈。
\end{enumerate}
\end{lemma}

\begin{proof}
略；不过可参见 Categories, Lemma
\ref{categories-lemma-fibred-groupoids}。
\end{proof}

\begin{lemma}
\label{stacks-lemma-stack-gives-stack-groupoids}
令 $\mathcal{C}$ 为位点。
令 $p : \mathcal{S} \to \mathcal{C}$ 为栈。
令 $p' : \mathcal{S}' \to \mathcal{C}$ 为与
$\mathcal{S}$ 相联系、由
Categories, Lemma
\ref{categories-lemma-fibred-gives-fibred-groupoids}
构造的纤维化于群胚的范畴。那么
$p' : \mathcal{S}' \to \mathcal{C}$ 是群胚中的栈。
\end{lemma}

\begin{proof}
回忆，$\mathcal{S}'$ 中的态射恰为 $\mathcal{S}$ 的强笛卡尔态射，
而 $\mathcal{S}$ 的任何同构都是这样的态射。因此，
$\mathcal{S}'$ 中的下降数据与 $\mathcal{S}$ 中的下降数据
完全相同。现在应用引理 \ref{stacks-lemma-stack-equivalences}。
若干细节从略。
\end{proof}

\begin{lemma}
\label{stacks-lemma-stack-in-groupoids-equivalent}
令 $\mathcal{C}$ 为位点。
令 $\mathcal{S}_1$、$\mathcal{S}_2$ 为
$\mathcal{C}$ 上的范畴。假设 $\mathcal{S}_1$ 与
$\mathcal{S}_2$ 作为 $\mathcal{C}$ 上的范畴彼此等价。
那么 $\mathcal{S}_1$ 是 $\mathcal{C}$ 上的群胚中的栈，
当且仅当 $\mathcal{S}_2$ 是 $\mathcal{C}$ 上的群胚中的栈。
\end{lemma}

\begin{proof}
由引理 \ref{stacks-lemma-stack-in-groupoids-stack} 与
\ref{stacks-lemma-stack-equivalent} 合并即得。
\end{proof}

\noindent
$\mathcal{C}$ 上群胚中的栈的 $2$-范畴定义如下。

\begin{definition}
\label{stacks-definition-stacks-in-groupoids-over-C}
令 $\mathcal{C}$ 为位点。
{\it $\mathcal{C}$ 上群胚中的栈的 $2$-范畴}
是 $\mathcal{C}$ 上栈的 $2$-范畴（见
Definition \ref{stacks-definition-stacks-over-C}）的如下子
$2$-范畴：
\begin{enumerate}
\item 其对象为群胚中的栈
$p : \mathcal{S} \to \mathcal{C}$。
\item 其 $1$-态射
$(\mathcal{S}, p) \to (\mathcal{S}', p')$ 为满足
$p' \circ G = p$ 的函子
$G : \mathcal{S} \to \mathcal{S}'$。
（由于每个态射都是强笛卡尔的，每个函子都保持它们。）
\item 对
$G, H : (\mathcal{S}, p) \to (\mathcal{S}', p')$，
其 $2$-态射 $t : G \to H$ 为函子之间的态射，并且对所有
$x \in \Ob(\mathcal{S})$ 都有
$p'(t_x) = \text{id}_{p(x)}$。
\end{enumerate}
\end{definition}

\noindent
注意，任意 $2$-态射自动为同构，因而
$\mathcal{C}$ 上群胚中的栈的 $2$-范畴实际上是一个
（严格）$(2, 1)$-范畴。

\begin{lemma}
\label{stacks-lemma-2-product-stacks-in-groupoids}
% P01-E0190: source says category although the object was defined over a site.
令 $\mathcal{C}$ 为范畴。
$\mathcal{C}$ 上群胚中的栈的 $2$-范畴具有 2-纤维积，
其描述与 Categories, Lemma
\ref{categories-lemma-2-product-categories-over-C} 中相同。
\end{lemma}

\begin{proof}
由 Categories, Lemma
\ref{categories-lemma-2-product-fibred-categories}
以及引理 \ref{stacks-lemma-stack-in-groupoids-stack}、
\ref{stacks-lemma-2-product-stacks} 立即可得。
\end{proof}








\section{类集合中的栈}
\label{stacks-section-stacks-in-setoids}

\noindent
本节仅简要说明：集合中的栈与集合层是同一回事。
记号请参见 Categories, Section
\ref{categories-section-fibred-in-setoids}。

\begin{definition}
\label{stacks-definition-stack-in-sets}
令 $\mathcal{C}$ 为位点。
\begin{enumerate}
\item $\mathcal{C}$ 上的一个{\it 类集合（setoid）中的栈}，
是 $\mathcal{C}$ 上所有纤维范畴均为类集合的栈。
\item 一个{\it 集合中的栈}，或{\it 离散范畴中的栈}，
是 $\mathcal{C}$ 上所有纤维范畴均为离散范畴的栈。
\end{enumerate}
\end{definition}

\noindent
由 Section \ref{stacks-section-stacks-in-groupoids} 中的讨论，
这等同于纤维范畴为类集合（相应地，为离散范畴）的群胚中的栈。
此外，它也等同于一个纤维化于类集合（相应地，集合）且本身为栈的范畴。

\begin{lemma}
\label{stacks-lemma-when-stack-in-sets}
令 $\mathcal{C}$ 为位点。在
$$
\left\{
\begin{matrix}
\text{预层范畴}\\
\text{集合范畴，位于 }\mathcal{C}
\end{matrix}
\right\}
\leftrightarrow
\left\{
\begin{matrix}
\text{范畴的范畴}\\
\text{在下列对象之上以集合为纤维： }\mathcal{C}
\end{matrix}
\right\}
$$
这一 Categories, Lemma
\ref{categories-lemma-2-category-fibred-sets} 的等价之下，
集合中的栈恰好对应于层。
\end{lemma}

\begin{proof}
略。提示：证明下降的有效性恰好对应于层条件。
\end{proof}

\begin{lemma}
\label{stacks-lemma-stack-in-setoids-characterize}
令 $\mathcal{C}$ 为位点。
令 $\mathcal{S}$ 为 $\mathcal{C}$ 上纤维化于类集合的范畴。
那么 $\mathcal{S}$ 是类集合中的栈，当且仅当唯一的、
纤维化于集合的等价范畴 $\mathcal{S}'$（见
Categories, Lemma \ref{categories-lemma-setoid-fibres}）
是集合中的栈。换言之，当且仅当预层
$$
U \longmapsto \Ob(\mathcal{S}_U)/\!\!\cong
$$
是层。
\end{lemma}

\begin{proof}
略。
\end{proof}

\begin{lemma}
\label{stacks-lemma-stack-in-setoids-equivalent}
令 $\mathcal{C}$ 为位点。
令 $\mathcal{S}_1$、$\mathcal{S}_2$ 为
$\mathcal{C}$ 上的范畴。假设 $\mathcal{S}_1$ 与
$\mathcal{S}_2$ 作为 $\mathcal{C}$ 上的范畴彼此等价。
那么 $\mathcal{S}_1$ 是 $\mathcal{C}$ 上的类集合中的栈，
当且仅当 $\mathcal{S}_2$ 是 $\mathcal{C}$ 上的类集合中的栈。
\end{lemma}

\begin{proof}
由 Categories, Lemma
\ref{categories-lemma-setoid-fibres} 可知，
$\mathcal{C}$ 上的范畴 $\mathcal{S}$ 在
$\mathcal{C}$ 上纤维化于类集合，当且仅当它在
$\mathcal{C}$ 上等价于某个纤维化于集合的范畴。
因此，$\mathcal{S}_1$ 在 $\mathcal{C}$ 上纤维化于类集合，
当且仅当 $\mathcal{S}_2$ 在 $\mathcal{C}$ 上纤维化于类集合。
现在由引理 \ref{stacks-lemma-stack-in-setoids-characterize} 即得结论。
\end{proof}

\noindent
$\mathcal{C}$ 上类集合中的栈的 $2$-范畴定义如下。

\begin{definition}
\label{stacks-definition-stacks-in-setoids-over-C}
令 $\mathcal{C}$ 为位点。
{\it $\mathcal{C}$ 上类集合中的栈的 $2$-范畴}
是 $\mathcal{C}$ 上栈的 $2$-范畴（见
Definition \ref{stacks-definition-stacks-over-C}）的如下子
$2$-范畴：
\begin{enumerate}
\item 其对象为类集合中的栈
$p : \mathcal{S} \to \mathcal{C}$。
\item 其 $1$-态射
$(\mathcal{S}, p) \to (\mathcal{S}', p')$ 为满足
$p' \circ G = p$ 的函子
$G : \mathcal{S} \to \mathcal{S}'$。
（由于每个态射都是强笛卡尔的，每个函子都保持它们。）
\item 对
$G, H : (\mathcal{S}, p) \to (\mathcal{S}', p')$，
其 $2$-态射 $t : G \to H$ 为函子之间的态射，并且对所有
$x \in \Ob(\mathcal{S})$ 都有
$p'(t_x) = \text{id}_{p(x)}$。
\end{enumerate}
\end{definition}

\noindent
注意，任意 $2$-态射自动为同构，因而
$\mathcal{C}$ 上类集合中的栈的 $2$-范畴实际上是一个
（严格）$(2, 1)$-范畴。

\begin{lemma}
\label{stacks-lemma-2-product-stacks-in-setoids}
令 $\mathcal{C}$ 为位点。
$\mathcal{C}$ 上类集合中的栈的 $2$-范畴具有 2-纤维积，
其描述与 Categories, Lemma
\ref{categories-lemma-2-product-categories-over-C} 中相同。
\end{lemma}

\begin{proof}
由 Categories, Lemmas
\ref{categories-lemma-2-product-fibred-categories} 与
\ref{categories-lemma-2-product-categories-fibred-setoids}
以及引理 \ref{stacks-lemma-stack-in-groupoids-stack} 与
\ref{stacks-lemma-2-product-stacks} 立即可得。
\end{proof}

\begin{lemma}
\label{stacks-lemma-2-fibre-product-gives-stack-in-setoids}
令 $\mathcal{C}$ 为位点。
令 $\mathcal{S}, \mathcal{T}$ 为 $\mathcal{C}$ 上
群胚中的栈，并令 $\mathcal{R}$ 为 $\mathcal{C}$ 上
类集合中的栈。令
$f : \mathcal{T} \to \mathcal{S}$ 与
$g : \mathcal{R} \to \mathcal{S}$ 为 $1$-态射。
若 $f$ 忠实，则 $2$-纤维积
$$
\mathcal{T} \times_{f, \mathcal{S}, g} \mathcal{R}
$$
是 $\mathcal{C}$ 上类集合中的栈。
\end{lemma}

\begin{proof}
由 Categories, Lemma
\ref{categories-lemma-2-product-categories-over-C}
中对 $2$-纤维积的显式描述立即可得。
\end{proof}

\begin{lemma}
\label{stacks-lemma-2-fibre-product-stacks-in-setoids-over-stack-in-groupoids}
令 $\mathcal{C}$ 为位点。
令 $\mathcal{S}$ 为 $\mathcal{C}$ 上群胚中的栈，并令
$\mathcal{S}_i$，$i = 1, 2$ 为 $\mathcal{C}$ 上
类集合中的栈。令 $f_i : \mathcal{S}_i \to \mathcal{S}$
为 $1$-态射。那么 $2$-纤维积
$$
\mathcal{S}_1 \times_{f_1, \mathcal{S}, f_2} \mathcal{S}_2
$$
是 $\mathcal{C}$ 上类集合中的栈。
\end{lemma}

\begin{proof}
这是引理
\ref{stacks-lemma-2-fibre-product-gives-stack-in-setoids}
在 $f_2$ 忠实时的特殊情形。
\end{proof}

\begin{lemma}
\label{stacks-lemma-faithful-descent}
令 $\mathcal{C}$ 为位点。令
$$
\xymatrix{
\mathcal{T}_2 \ar[r] \ar[d]_{G'} & \mathcal{T}_1 \ar[d]^G \\
\mathcal{S}_2 \ar[r]^F & \mathcal{S}_1
}
$$
为 $\mathcal{C}$ 上群胚中的栈所组成的 $2$-笛卡尔图式。假设
\begin{enumerate}
\item 对每个 $U \in \Ob(\mathcal{C})$ 和
$x \in \Ob((\mathcal{S}_1)_U)$，存在覆盖
$\{U_i \to U\}$，使得 $x|_{U_i}$ 属于
$F : (\mathcal{S}_2)_{U_i} \to (\mathcal{S}_1)_{U_i}$
的本质像；且
\item $G'$ 忠实。
\end{enumerate}
那么 $G$ 忠实。
\end{lemma}

\begin{proof}
可假设 $\mathcal{T}_2$ 就是由
Categories, Lemma
\ref{categories-lemma-2-product-categories-over-C}
描述的范畴
$\mathcal{S}_2 \times_{\mathcal{S}_1} \mathcal{T}_1$。
由 Categories, Lemma
\ref{categories-lemma-equivalence-fibred-categories}，
$G, G'$ 的忠实性可在纤维范畴上检验。
设 $y, y'$ 为 $\mathcal{T}_1$ 中位于
$\mathcal{C}$ 的对象 $U$ 上的对象。令
$\alpha, \beta : y \to y'$ 为 $(\mathcal{T}_1)_U$ 中满足
$G(\alpha) = G(\beta)$ 的态射。我们的目标是证明
$\alpha = \beta$。转而考虑
$\gamma = \alpha^{-1} \circ \beta$，可知
$G(\gamma) = \text{id}_{G(y)}$，而我们必须证明
$\gamma = \text{id}_y$。由假设，可找到覆盖
$\{U_i \to U\}$，使 $G(y)|_{U_i}$ 属于
$F :(\mathcal{S}_2)_{U_i} \to (\mathcal{S}_1)_{U_i}$
的本质像。因为只需对每个 $i$ 证明
$\gamma|_{U_i} = \text{id}$，故可假设对
$\mathcal{S}_2$ 中位于 $U$ 上的某个对象 $x$，有
% P01-E0191: source calls the already displayed f “morphisms”; it is the
% isomorphism supplied by essential-image membership.
$f : F(x) \to G(y)$，且 $f$ 是
$(\mathcal{S}_1)_U$ 中的同构。在这种情况下，得到纤维范畴
$\mathcal{S}_2 \times_{\mathcal{S}_1} \mathcal{T}_1$
中位于 $U$ 上的态射
$$
(1, \gamma) : (U, x, y, f) \longrightarrow (U, x, y, f)
$$
它在 $G'$ 下于 $\mathcal{S}_1$ 中的像是
$\text{id}_x$。由于 $G'$ 忠实，遂有
$\gamma = \text{id}_y$，结论得证。
\end{proof}

\begin{lemma}
\label{stacks-lemma-stack-in-setoids-descent}
令 $\mathcal{C}$ 为位点。令
$$
\xymatrix{
\mathcal{T}_2 \ar[r] \ar[d] & \mathcal{T}_1 \ar[d]^G \\
\mathcal{S}_2 \ar[r]^F & \mathcal{S}_1
}
$$
为 $\mathcal{C}$ 上群胚中的栈所组成的 $2$-笛卡尔图式。若
\begin{enumerate}
\item $F : \mathcal{S}_2 \to \mathcal{S}_1$ 全忠实；
\item 对每个 $U \in \Ob(\mathcal{C})$ 和
$x \in \Ob((\mathcal{S}_1)_U)$，存在覆盖
$\{U_i \to U\}$，使得 $x|_{U_i}$ 属于
$F : (\mathcal{S}_2)_{U_i} \to (\mathcal{S}_1)_{U_i}$
的本质像；且
\item $\mathcal{T}_2$ 是类集合中的栈，
\end{enumerate}
那么 $\mathcal{T}_1$ 是类集合中的栈。
\end{lemma}

\begin{proof}
可假设 $\mathcal{T}_2$ 就是由
Categories, Lemma
\ref{categories-lemma-2-product-categories-over-C}
描述的范畴
$\mathcal{S}_2 \times_{\mathcal{S}_1} \mathcal{T}_1$。
选取 $U \in \Ob(\mathcal{C})$ 与
$y \in \Ob((\mathcal{T}_1)_U)$。
我们必须证明 $\mathcal{C}/U$ 上的层
$\mathit{Aut}(y)$ 是平凡的。
% P01-E0192: source has duplicated “To to”.
为此可以用 $U$ 的一个覆盖的各成员替换 $U$。
因此，由假设 (2)，可假设存在对象
$x \in \Ob((\mathcal{S}_2)_U)$ 与同构
$f : F(x) \to G(y)$。
于是 $y' = (U, x, y, f)$ 是 $\mathcal{T}_2$
中位于 $U$ 上的对象，它在投影
$\mathcal{T}_2 \to \mathcal{T}_1$ 下映到 $y$。
由 (1) 中 $F$ 全忠实，映射
$\mathit{Aut}(y') \to \mathit{Aut}(y)$ 是满射；
这里使用 Categories, Lemma
\ref{categories-lemma-2-product-categories-over-C}
中对 $\mathcal{T}_2$ 的态射的显式描述。
由 (3)，层 $\mathit{Aut}(y')$ 是平凡的，故得引理结论。
\end{proof}

\begin{lemma}
\label{stacks-lemma-relative-sheaf-over-stack-is-stack}
令 $\mathcal{C}$ 为位点。令
$F : \mathcal{S} \to \mathcal{T}$ 为
$\mathcal{C}$ 上纤维化于群胚的范畴之间的 $1$-态射。假设：
\begin{enumerate}
\item $\mathcal{T}$ 是 $\mathcal{C}$ 上群胚中的栈；
\item 对每个 $U \in \Ob(\mathcal{C})$，纤维范畴之间的函子
$\mathcal{S}_U \to \mathcal{T}_U$ 忠实；
\item 对每个 $U$ 与每个 $y \in \Ob(\mathcal{T}_U)$，预层
$$
(h : V \to U)
\longmapsto
\{(x, f) \mid x \in \Ob(\mathcal{S}_V), f : F(x) \to f^*y\text{ 相对于 }V\}/\cong
$$
是 $\mathcal{C}/U$ 上的层。
\end{enumerate}
那么 $\mathcal{S}$ 是 $\mathcal{C}$ 上群胚中的栈。
\end{lemma}

\begin{proof}
我们必须证明态射的下降与对象的下降。

\medskip\noindent
态射的下降。令 $\{U_i \to U\}$ 为 $\mathcal{C}$ 的覆盖。
令 $x, x'$ 为 $\mathcal{S}$ 中位于 $U$ 上的对象。
对每个 $i$，令
$\alpha_i : x|_{U_i} \to x'|_{U_i}$ 为位于 $U_i$ 上的态射，
使得 $\alpha_i$ 与 $\alpha_j$ 限制为同一个态射
$x|_{U_i \times_U U_j} \to x'|_{U_i \times_U U_j}$。
由于 $\mathcal{T}$ 是群胚中的栈，存在位于 $U$ 上的态射
$\beta : F(x) \to F(x')$，其在 $U_i$ 上的限制为
$F(\alpha_i)$。于是可把
$\xi = (x, \beta)$ 与
$\xi' = (x', \text{id}_{F(x')})$
视为假设 (3) 中与 $y = F(x')$ 相联系的预层在
$U$ 上的截面。另一方面，$\xi$ 与 $\xi'$ 在
$U_i$ 上的限制分别为
$(x|_{U_i}, F(\alpha_i))$ 与
$(x'|_{U_i}, \text{id}_{F(x'|_{U_i})})$。
态射 $\alpha_i$ 给出这两者之间的同构。
因此，由假设 (3)，$\xi$ 与 $\xi'$ 同构。
这意味着存在位于 $U$ 上的态射
$\alpha : x \to x'$，使 $F(\alpha) = \beta$。
由于 $F$ 在纤维范畴上忠实，得到
$\alpha|_{U_i} = \alpha_i$。

\medskip\noindent
对象的下降。令 $\{U_i \to U\}$ 为 $\mathcal{C}$ 的覆盖。
令 $(x_i, \varphi_{ij})$ 为 $\mathcal{S}$ 中相对于该覆盖的
下降数据。由于 $\mathcal{T}$ 是群胚中的栈，存在
$\mathcal{T}_U$ 中的对象 $y$ 以及同构
$\beta_i : F(x_i) \to y|_{U_i}$，使得
$F(\varphi_{ij}) = \beta_j|_{U_i \times_U U_j} \circ
(\beta_i|_{U_i \times_U U_j})^{-1}$。
于是 $(x_i, \beta_i)$ 是假设 (3) 中定义的、与
$y$ 相联系的预层在 $U$ 上的截面。
此外，$\varphi_{ij}$ 定义从
$(x_i, \beta_i)|_{U_i \times_U U_j}$ 到
$(x_j, \beta_j)|_{U_i \times_U U_j}$ 的同构。
因此，由假设 (3)，存在 $U$ 上的偶对 $(x, \beta)$，
其在 $U_i$ 上的限制同构于 $(x_i, \beta_i)$。
这意味着存在态射
$\alpha_i : x_i \to x|_{U_i}$，满足
$\beta_i = \beta|_{U_i} \circ F(\alpha_i)$。
由于 $F$ 在纤维范畴上忠实，计算表明
$\varphi_{ij} = \alpha_j|_{U_i \times_U U_j} \circ
(\alpha_i|_{U_i \times_U U_j})^{-1}$。证明完毕。
\end{proof}










\section{惯性栈}
\label{stacks-section-the-inertia-stack}

\noindent
令 $p : \mathcal{S} \to \mathcal{C}$ 与
$p' : \mathcal{S}' \to \mathcal{C}$ 为范畴
$\mathcal{C}$ 上的纤维化范畴。
令 $F : \mathcal{S} \to \mathcal{S}'$ 为
$\mathcal{C}$ 上纤维化范畴之间的 $1$-态射。
回忆，我们在 Categories, Definition
\ref{categories-definition-inertia-fibred-category} 中定义了一个
{\it 相对惯性纤维化范畴}
$\mathcal{I}_{\mathcal{S}/\mathcal{S}'} \to \mathcal{C}$：
其对象为偶对 $(x , \alpha)$，其中
$x \in \Ob(\mathcal{S})$ 且 $\alpha : x \to x$ 满足
$F(\alpha) = \text{id}_{F(x)}$。还有一个绝对版本，即
$\mathcal{S}$ 的{\it 惯性} $\mathcal{I}_\mathcal{S}$。
若 $\mathcal{S}$ 与 $\mathcal{S}'$ 都是栈，则这些惯性范畴
实际上是 $\mathcal{C}$ 上的栈。

\begin{lemma}
\label{stacks-lemma-inertia}
令 $\mathcal{C}$ 为位点。令
$p : \mathcal{S} \to \mathcal{C}$ 与
$p' : \mathcal{S}' \to \mathcal{C}$ 为位点
$\mathcal{C}$ 上的栈。令
$F : \mathcal{S} \to \mathcal{S}'$ 为
$\mathcal{C}$ 上栈之间的 $1$-态射。
\begin{enumerate}
\item 惯性 $\mathcal{I}_{\mathcal{S}/\mathcal{S}'}$ 与
$\mathcal{I}_\mathcal{S}$ 是 $\mathcal{C}$ 上的栈。
\item 若 $\mathcal{S}, \mathcal{S}'$ 是
$\mathcal{C}$ 上群胚中的栈，则
$\mathcal{I}_{\mathcal{S}/\mathcal{S}'}$ 与
$\mathcal{I}_\mathcal{S}$ 亦然。
\item 若 $\mathcal{S}, \mathcal{S}'$ 是
$\mathcal{C}$ 上类集合中的栈，则
$\mathcal{I}_{\mathcal{S}/\mathcal{S}'}$ 与
$\mathcal{I}_\mathcal{S}$ 亦然。
\end{enumerate}
\end{lemma}

\begin{proof}
前三个断言由引理 \ref{stacks-lemma-2-product-stacks}、
\ref{stacks-lemma-2-product-stacks-in-groupoids}、
\ref{stacks-lemma-2-product-stacks-in-setoids} 以及
Categories, Lemma
\ref{categories-lemma-inertia-fibred-category} (1)
中的等价得出。
\end{proof}

\begin{lemma}
\label{stacks-lemma-characterize-stack-in-setoids}
令 $\mathcal{C}$ 为位点。
若 $\mathcal{S}$ 是群胚中的栈，则典范 $1$-态射
$\mathcal{I}_\mathcal{S} \to \mathcal{S}$ 是等价，
当且仅当 $\mathcal{S}$ 是类集合中的栈。
\end{lemma}

\begin{proof}
直接由 Categories, Lemma
\ref{categories-lemma-characterize-fibred-setoids-inertia} 得出。
\end{proof}




\section{纤维化范畴的栈化}
\label{stacks-section-stackify}

\noindent
下一个引理中过程的结果称为位点上纤维化范畴的{\it 栈化}。

\begin{lemma}
\label{stacks-lemma-stackify}
令 $\mathcal{C}$ 为位点。
令 $p : \mathcal{S} \to \mathcal{C}$ 为
$\mathcal{C}$ 上的纤维化范畴。存在栈
$p' : \mathcal{S}' \to \mathcal{C}$ 以及
$\mathcal{C}$ 上纤维化范畴之间的 $1$-态射
$G : \mathcal{S} \to \mathcal{S}'$（见
Categories, Definition
\ref{categories-definition-fibred-categories-over-C}），满足：
\begin{enumerate}
\item 对每个 $U \in \Ob(\mathcal{C})$ 以及任意
$x, y \in \Ob(\mathcal{S}_U)$，由 $G$ 诱导的映射
$$
\mathit{Mor}(x, y) \longrightarrow \mathit{Mor}(G(x), G(y))
$$
把右端识别为左端的层化；且
\item 对每个 $U \in \Ob(\mathcal{C})$ 及任意
$x' \in \Ob(\mathcal{S}'_U)$，存在覆盖
$\{U_i \to U\}_{i \in I}$，使得对每个
$i \in I$，对象 $x'|_{U_i}$ 属于函子
$G : \mathcal{S}_{U_i} \to \mathcal{S}'_{U_i}$ 的本质像。
\end{enumerate}
此外，这些条件在唯一 $2$-同构意义下确定栈 $\mathcal{S}'$。
\end{lemma}

\begin{proof}[朴素方法的证明]
在这种证明方法中，我们分阶段进行。

\medskip\noindent
首先，给定 $U$ 上的 $x$ 与 $\mathcal{S}$ 的任意对象
$y$，若 $\mathcal{S}$ 的两个态射
$a, b : x \to y$ 位于 $\mathcal{C}$ 的同一个箭头上，
并且存在 $\mathcal{C}$ 的覆盖
$\{f_i : U_i \to U\}$，使得复合
$$
f_i^*x \to x \xrightarrow{a} y,
\quad
f_i^*x \to x \xrightarrow{b} y
$$
相等，则称它们{\it 局部相等}。这在 $\mathcal{S}$ 的箭头上
给出等价关系 $\sim$。若 $b \sim b'$，则
$a \circ b \circ c \sim a \circ b' \circ c$（验证从略）。
因此可对该等价关系取商，得到 $\mathcal{C}$ 上的新范畴
$\mathcal{S}^1$，以及态射
$G^1 : \mathcal{S} \to \mathcal{S}^1$。

\medskip\noindent
可验证 $G^1$ 保持强笛卡尔态射，并且
$\mathcal{S}^1$ 是 $\mathcal{C}$ 上的纤维化范畴。
验证从略。于是问题归约到局部相等的态射彼此相等的情形。

\medskip\noindent
其次，我们如下加入态射。给定 $U$ 上的 $x$ 与
% P01-E0193: source has “any object y of lying over V”.
位于 $V$ 上的任意对象 $y$，一个{\it 从 $x$ 到 $y$
的局部定义态射}由下列资料给出：
\begin{enumerate}
\item 态射 $f : U \to V$；
\item $U$ 的覆盖 $\{f_i : U_i \to U\}$；以及
\item 满足 $p(a_i) = f \circ f_i$ 的态射
$a_i : f_i^*x \to y$，
\end{enumerate}
并要求复合
$$
(f_i \times f_j)^*x \to f_i^*x \xrightarrow{a_i} y,
\quad
(f_i \times f_j)^*x \to f_j^*x \xrightarrow{a_j} y
$$
相等。注意，通常的态射 $a : x \to y$ 给出局部定义态射
$(p(a) : U \to V, \{\text{id}_U\}, a)$。
称两个局部定义态射
$(f, \{f_i : U_i \to U\}, a_i)$ 与
$(g, \{g_j : U'_j \to U\}, b_j)$ {\it 相等}，如果
$f = g$，且复合
$$
(f_i \times g_j)^*x \to f_i^*x \xrightarrow{a_i} y,
\quad
(f_i \times g_j)^*x \to g_j^*x \xrightarrow{b_j} y
$$
相等（这是正确的条件，因为我们处在局部相等态射彼此相等的情形）。
要复合从 $x$ 到 $y$ 的局部定义态射
$(f, \{f_i : U_i \to U\}, a_i)$ 与从 $y$ 到位于
$W$ 上的 $z$ 的局部定义态射
$(g, \{g_j : V_j \to V\}, b_j)$，只需取
$g \circ f : U \to W$、覆盖
$\{U_i \times_V V_j \to U\}$，并以复合
$$
x|_{U_i \times_V V_j}
\xrightarrow{\text{pr}_0^*a_i}
y|_{V_j}
\xrightarrow{b_j}
z
$$
为映射。我们略去验证这是局部定义态射。

\medskip\noindent
可验证，以与 $\mathcal{S}$ 相同的对象为对象、以局部定义态射为态射的
$\mathcal{S}^2$ 是 $\mathcal{C}$ 上的范畴，存在
$\mathcal{C}$ 上的函子
$G^2 : \mathcal{S} \to \mathcal{S}^2$，
% P01-E0194: source says “strongly cartesian objects”; frozen wording retained.
该函子保持强笛卡尔对象，并且 $\mathcal{S}^2$ 是
$\mathcal{C}$ 上的纤维化范畴。验证从略。
通过检验使用局部定义态射的效果是对（分离的）态射预层取层化，
问题便归约到 $\mathcal{S}$ 的态射预层全为层的情形。

\medskip\noindent
最后，在态射预层全为层的情形，为保证最终结果中下降条件有效，
必须加入对象。最简单的做法是考虑范畴 $\mathcal{S}'$，
其对象为偶对 $(\mathcal{U}, \xi)$，其中
$\mathcal{U} = \{U_i \to U\}$ 是 $\mathcal{C}$ 的覆盖，而
% P01-E0195: source omits “to” in “relative U”.
$\xi = (X_i, \varphi_{ii'})$ 是相对于 $\mathcal{U}$ 的下降数据。
设给定两组这样的资料
$(\mathcal{U}, \xi) = (\{f_i : U_i \to U\},  x_i, \varphi_{ii'})$ 与
$(\mathcal{V}, \eta) = (\{g_j : V_j \to V\},  y_j, \psi_{jj'})$。
我们把
$$
\Mor_{\mathcal{S}'}((\mathcal{U}, \xi), (\mathcal{V}, \eta))
$$
定义为所有 $(f, a_{ij})$ 组成的集合，其中
$f : U \to V$，且
$$
a_{ij} :
x_i|_{U_i \times_V V_j}
\longrightarrow
y_j
$$
是 $\mathcal{S}$ 中位于 $U_i \times_V V_j \to V_j$
上的态射。它们须满足下列条件：对任意
$i, i' \in I$ 与 $j, j' \in J$，令
$W = (U_i \times_U U_{i'}) \times_V (V_j \times_V V_{j'})$。
则图式
$$
\xymatrix{
x_i|_W \ar[r]_{a_{ij}|_W} \ar[d]_{\varphi_{ii'}|_W} &
y_j|_W \ar[d]^{\psi_{jj'}|_W} \\
x_{i'}|_W \ar[r]^{a_{i'j'}|_W} &
y_{j'}|_W
}
$$
交换。此时必须验证：
\begin{enumerate}
\item 上述态射有定义良好的复合；
\item 这使 $\mathcal{S}'$ 成为 $\mathcal{C}$ 上的范畴；
\item 存在 $\mathcal{C}$ 上的函子
$G : \mathcal{S} \to \mathcal{S}'$；
\item 对 $\mathcal{S}$ 的对象 $x, y$，有
$\Mor_\mathcal{S}(x, y) = \Mor_{\mathcal{S}'}(G(x), G(y))$；
\item $\mathcal{S}'$ 的任意对象局部地来自
$\mathcal{S}$ 的对象，即引理的 (2) 成立；
\item $G$ 保持强笛卡尔态射；
\item $\mathcal{S}'$ 是 $\mathcal{C}$ 上的纤维化范畴；且
\item $\mathcal{S}'$ 是 $\mathcal{C}$ 上的栈。
\end{enumerate}
这些都不困难，只是事项很多。细节从略。
\end{proof}

\begin{proof}[较不朴素的证明]
下面给出一个较不朴素的证明。
由 Categories, Lemma
\ref{categories-lemma-fibred-strict}，存在纤维化范畴之间的等价
$\mathcal{S} \to \mathcal{S}'$，其中
$\mathcal{S}'$ 是分裂纤维化范畴，即其拉回函子严格地复合。
显然，对 $\mathcal{S}'$ 成立的引理蕴含对
$\mathcal{S}$ 成立的引理。因此，可把 $\mathcal{S}$
看作取值于范畴的预层。

\medskip\noindent
暂时忘掉 $2$-态射，把 $2$-范畴 $\textit{Cat}$
看成一个范畴。按照 Categories, Section
\ref{categories-section-definition-categories}，把范畴看成五元组
$(\text{Ob}, \text{Arrows}, s, t, \circ)$。
考虑遗忘函子
$$
forget : \textit{Cat} \to \textit{Sets} \times \textit{Sets}, \quad
(\text{Ob}, \text{Arrows}, s, t, \circ)
\mapsto
(\text{Ob}, \text{Arrows}).
$$
那么 $forget$ 忠实，$\textit{Cat}$ 有极限且
$forget$ 与极限交换，$\textit{Cat}$ 有有向余极限且
$forget$ 与之交换，并且 $forget$ 反映同构。
可以对取值于 $\textit{Cat}$ 的预层作层化；通过与
Sites, Section \ref{sites-section-sheaves-algebraic-structures}
第一部分中类似的论证，所得构造与 $forget$ 交换。
将此应用于 $\mathcal{S}$，得到层化
$\mathcal{S}^\#$；它有对象层与态射层，二者分别为
$\mathcal{S}$ 相应预层的层化。在这种情况下，很容易看出映射
$\mathcal{S} \to \mathcal{S}^\#$ 具有引理中的性质 (1) 与 (2)。

\medskip\noindent
然而，范畴 $\mathcal{S}^\#$ 可能还不是栈，因为尽管对象预层是层，
下降条件仍可能不成立。为补救这一点，必须加入更多对象。
不过，上述论证确实把问题归约到
$\mathcal{S} = \mathcal{S}_F$ 的情形，其中
$F : \mathcal{C}^{opp} \to \textit{Cat}$ 是取值于范畴的某个层（！）。
在此情形，考虑如下定义的函子
$F' : \mathcal{C}^{opp} \to \textit{Cat}$：
\begin{enumerate}
\item 集合 $\Ob(F'(U))$ 是所有偶对
$(\mathcal{U}, \xi)$ 组成的集合，其中
$\mathcal{U} = \{U_i \to U\}$ 是 $U$ 的覆盖，而
$\xi = (x_i, \varphi_{ii'})$ 是相对于
$\mathcal{U}$ 的下降数据。
\item $F'(U)$ 中从 $(\mathcal{U}, \xi)$ 到
$(\mathcal{V}, \eta)$ 的态射，是
$$
\colim \Mor_{DD(\mathcal{W})}(a^*\xi, b^*\eta)
$$
的一个元素；这里余极限遍历所有共同加细
$a : \mathcal{W} \to \mathcal{U}$、
$b : \mathcal{W} \to \mathcal{V}$。
该余极限是滤过的（验证从略）。
% P01-E0196: source says F(U); the category being defined is F'(U).
因此，$F(U)$ 中态射的复合可通过寻找共同加细并在
$DD(\mathcal{W})$ 中复合来定义。
\item 给定 $h : V \to U$ 以及 $F'(U)$ 的对象
$(\mathcal{U}, \xi)$，令 $F'(h)(\mathcal{U}, \xi)$
等于 $(V \times_U \mathcal{U}, \text{pr}_1^*\xi)$。
更确切地，若 $\mathcal{U} = \{U_i \to U\}$ 且
$\xi = (x_i, \varphi_{ii'})$，则
$V \times_U \mathcal{U} = \{V \times_U U_i \to V\}$；
它带有典范态射
$\text{pr}_1 : V \times_U \mathcal{U} \to \mathcal{U}$，而
$\text{pr}_1^*\xi$ 是 $\xi$ 沿该态射的拉回
（见 Definition \ref{stacks-definition-pullback-functor}）。
\item 给定 $h : V \to U$、对象
$(\mathcal{U}, \xi)$ 与 $(\mathcal{V}, \eta)$，
以及它们之间由
$a : \mathcal{W} \to \mathcal{U}$、
$b : \mathcal{W} \to \mathcal{V}$ 和
$\alpha : a^*\xi \to b^*\eta$ 表示的态射，则
$F'(h)(\alpha)$ 由
$a' : V \times_U\mathcal{W} \to V \times_U\mathcal{U}$、
$b' : V \times_U\mathcal{W} \to V \times_U\mathcal{V}$，
以及态射 $\alpha$ 沿映射
$V \times_U \mathcal{W} \to \mathcal{W}$ 的拉回
$\alpha'$ 表示。由于 $\mathcal{S}_F$ 中的拉回严格交换，
这一构造成立。
\end{enumerate}
存在映射 $F \to F'$，它把 $F(U)$ 的对象 $x$
送到 $F'(U)$ 的对象
$(\{U \to U\}, (x, triv))$。
此时必须检验相应函子
$\mathcal{S}_F \to \mathcal{S}_{F'}$
具有引理中的性质 (1)、(2)，并最终检验
$\mathcal{S}_{F'}$ 是栈。细节从略。
\end{proof}

\begin{lemma}
\label{stacks-lemma-stackify-universal-property}
令 $\mathcal{C}$ 为位点。
令 $p : \mathcal{S} \to \mathcal{C}$ 为
$\mathcal{C}$ 上的纤维化范畴。令
% P01-E0197: source omits “be” before “the stack and 1-morphism”.
$p' : \mathcal{S}' \to \mathcal{C}$ 与
$G : \mathcal{S} \to \mathcal{S}'$ 为引理
\ref{stacks-lemma-stackify} 中构造的栈与 $1$-态射。
该构造具有如下泛性质：给定栈
$q : \mathcal{X} \to \mathcal{C}$ 以及
$\mathcal{C}$ 上纤维化范畴之间的 $1$-态射
$F : \mathcal{S} \to \mathcal{X}$，存在
$1$-态射 $H : \mathcal{S}' \to \mathcal{X}$，使图式
$$
\xymatrix{
\mathcal{S} \ar[rr]_F \ar[rd]_G & & \mathcal{X} \\
& \mathcal{S}' \ar[ru]_H
}
$$
为 $2$-交换的。
\end{lemma}

\begin{proof}
略。提示：假设 $x' \in \Ob(\mathcal{S}'_U)$。
由引理 \ref{stacks-lemma-stackify} 的结论，存在覆盖
$\{U_i \to U\}_{i \in I}$，使得对某些
$x_i \in \Ob(\mathcal{S}_{U_i})$ 有
$x'|_{U_i} = G(x_i)$。此外，存在覆盖
$\{U_{ijk} \to U_i \times_U U_j\}$ 与同构
$\alpha_{ijk} : x_i|_{U_{ijk}} \to x_j|_{U_{ijk}}$，
满足 $G(\alpha_{ijk}) = \text{id}_{x'|_{U_{ijk}}}$。
令 $y_i = F(x_i)$。可以检验
$$
F(\alpha_{ijk}) : y_i|_{U_{ijk}} \to y_j|_{U_{ijk}}
$$
在交叠处一致，因而（由于 $\mathcal{X}$ 是栈）定义态射
$\beta_{ij} : y_i|_{U_i \times_U U_j} \to
y_j|_{U_i \times_U U_j}$。
其次，检验 $\beta_{ij}$ 定义下降数据。
由于 $\mathcal{X}$ 是栈，这些下降数据有效，因而找到
% P01-E0198: source says G(x_i), while this proof set y_i = F(x_i).
$\mathcal{X}_U$ 的对象 $y$，它在 $U_i$ 上与
$G(x_i)$ 一致。提示是令 $H(x') = y$。
\end{proof}

\begin{lemma}
\label{stacks-lemma-stackify-universal-property-more}
记号与假设同引理
\ref{stacks-lemma-stackify-universal-property}。
由上述引理证明中的构造，得到范畴之间的典范等价
$$
\Mor_{\textit{Fib}/\mathcal{C}}(\mathcal{S}, \mathcal{X})
=
\Mor_{\textit{Stacks}/\mathcal{C}}(\mathcal{S}', \mathcal{X})
$$
\end{lemma}

\begin{proof}
略。
\end{proof}

\begin{lemma}
\label{stacks-lemma-stackification-fibre-product-fibred-categories}
令 $\mathcal{C}$ 为位点。
令 $f : \mathcal{X} \to \mathcal{Y}$ 与
$g : \mathcal{Z} \to \mathcal{Y}$ 为
$\mathcal{C}$ 上纤维化范畴之间的态射。
在这种情况下，$2$-纤维积的栈化是各栈化的
$2$-纤维积。
\end{lemma}

\begin{proof}
以 $\mathcal{X}', \mathcal{Y}', \mathcal{Z}'$
表示相应栈化，并以 $\mathcal{W}$ 表示
$\mathcal{X} \times_\mathcal{Y} \mathcal{Z}$ 的栈化。
由 $2$-纤维积的构造，存在典范 $1$-态射
$\mathcal{X} \times_\mathcal{Y} \mathcal{Z} \to
\mathcal{X}' \times_{\mathcal{Y}'} \mathcal{Z}'$。
由于第二个 $2$-纤维积是栈（见引理
\ref{stacks-lemma-2-product-stacks}），由栈化的泛性质
（见引理 \ref{stacks-lemma-stackify-universal-property}），
这个 $1$-态射诱导 $1$-态射
$h : \mathcal{W} \to
\mathcal{X}' \times_{\mathcal{Y}'} \mathcal{Z}'$。
现在 $h$ 是栈之间的态射；应用引理
\ref{stacks-lemma-characterize-ff} 与
\ref{stacks-lemma-characterize-essentially-surjective-when-ff}，
可检验它是等价。

\medskip\noindent
因此，先证明 $h$ 在 $\mathit{Mor}$-层上诱导同构。
令 $\xi, \xi'$ 为 $\mathcal{W}$ 中位于
$U \in \Ob(\mathcal{C})$ 上的对象。我们要证明
$$
h : \mathit{Mor}(\xi, \xi') \longrightarrow \mathit{Mor}(h(\xi), h(\xi'))
$$
是同构。为此，可在 $U$ 上局部地工作（见
Sites, Section \ref{sites-section-glueing-sheaves}）。
因此，由 $\mathcal{W}$ 的构造（见引理
\ref{stacks-lemma-stackify}），可假设 $\xi, \xi'$
确实来自 $\mathcal{X} \times_\mathcal{Y} \mathcal{Z}$
中位于 $U$ 上的对象
$(x, z, \alpha)$ 与 $(x', z', \alpha')$。
再次由同一引理可知，在此情形，
$\mathit{Mor}(\xi, \xi')$ 是下列预层的层化：
$$
V/U \longmapsto
\Mor_{\mathcal{X}_V}(x|_V, x'|_V)
\times_{\Mor_{\mathcal{Y}_V}(f(x)|_V, f(x')|_V)}
\Mor_{\mathcal{Z}_V}(z|_V, z'|_V)
$$
而 $\mathit{Mor}(h(\xi), h(\xi'))$ 等于纤维积
% P01-E0199: frozen source display lacks the closing parenthesis after j(f(x')).
$$
\mathit{Mor}(i(x), i(x'))
\times_{\mathit{Mor}(j(f(x)), j(f(x'))}
\mathit{Mor}(k(z), k(z'))
$$
其中 $i : \mathcal{X} \to \mathcal{X}'$、
$j : \mathcal{Y} \to \mathcal{Y}'$ 与
$k : \mathcal{Z} \to \mathcal{Z}'$ 是典范函子。
因此，本段第一个显示映射是同构，因为层化是正合的
（从而预层的纤维积的层化，就是各层化的纤维积）。

\medskip\noindent
最后，必须检验
$\mathcal{X}' \times_{\mathcal{Y}'} \mathcal{Z}'$
中位于 $U$ 上的任意对象，在 $U$ 上局部地属于 $h$
的本质像。把这样的对象写成三元组
$(x', z', \alpha)$。那么 $x'$ 局部地来自
$\mathcal{X}$ 的对象，$z'$ 局部地来自
$\mathcal{Z}$ 的对象；在对 $x'$、$z'$ 作适当替换后，
$\mathcal{Y}'_U$ 的态射 $\alpha$ 局部地来自
$\mathcal{Y}$ 的态射。换言之，我们已经证明，
$\mathcal{X}' \times_{\mathcal{Y}'} \mathcal{Z}'$
中位于 $U$ 上的任意对象，在 $U$ 上局部地属于
$\mathcal{X} \times_\mathcal{Y} \mathcal{Z} \to
\mathcal{X}' \times_{\mathcal{Y}'} \mathcal{Z}'$
的本质像，因此更属于 $h$ 的局部本质像。
\end{proof}

\begin{lemma}
\label{stacks-lemma-stackification-inertia}
令 $\mathcal{C}$ 为位点。
令 $\mathcal{X}$ 为 $\mathcal{C}$ 上的纤维化范畴。
惯性纤维化范畴 $\mathcal{I}_\mathcal{X}$ 的栈化，
% P01-E0200: source omits the article in “is inertia of”.
是 $\mathcal{X}$ 的栈化的惯性。
\end{lemma}

\begin{proof}
这是因为，由引理
\ref{stacks-lemma-stackification-fibre-product-fibred-categories}，
栈化与 $2$-纤维积相容；并且范畴在
$\mathcal{C}$ 上的惯性可用 $2$-纤维积表示，见
Categories, Lemma \ref{categories-lemma-inertia-fibred-category}。
\end{proof}





\section{纤维化于群胚的范畴的栈化}
\label{stacks-section-stackify-groupoids}

\noindent
结论如下。

\begin{lemma}
\label{stacks-lemma-stackify-groupoids}
令 $\mathcal{C}$ 为位点。
令 $p : \mathcal{S} \to \mathcal{C}$ 为
$\mathcal{C}$ 上纤维化于群胚的范畴。
存在群胚中的栈
$p' : \mathcal{S}' \to \mathcal{C}$，以及
$\mathcal{C}$ 上纤维化于群胚的范畴之间的 $1$-态射
$G : \mathcal{S} \to \mathcal{S}'$（见
Categories, Definition
\ref{categories-definition-categories-fibred-in-groupoids-over-C}），满足：
\begin{enumerate}
\item 对每个 $U \in \Ob(\mathcal{C})$ 以及任意
$x, y \in \Ob(\mathcal{S}_U)$，由 $G$ 诱导的映射
$$
\mathit{Mor}(x, y) \longrightarrow \mathit{Mor}(G(x), G(y))
$$
把右端识别为左端的层化；且
\item 对每个 $U \in \Ob(\mathcal{C})$ 及任意
$x' \in \Ob(\mathcal{S}'_U)$，存在覆盖
$\{U_i \to U\}_{i \in I}$，使得对每个
$i \in I$，对象 $x'|_{U_i}$ 属于函子
$G : \mathcal{S}_{U_i} \to \mathcal{S}'_{U_i}$ 的本质像。
\end{enumerate}
此外，这些条件在唯一 $2$-同构意义下确定群胚中的栈
$\mathcal{S}'$。
\end{lemma}

\begin{proof}
应用引理 \ref{stacks-lemma-stackify}。再应用引理
\ref{stacks-lemma-stack-in-groupoids-stack}，所得结果是群胚中的栈。
\end{proof}

\begin{lemma}
\label{stacks-lemma-stackify-groupoids-universal-property}
令 $\mathcal{C}$ 为位点。
令 $p : \mathcal{S} \to \mathcal{C}$ 为
$\mathcal{C}$ 上纤维化于群胚的范畴。令
% P01-E0201: source omits “be” before the predicate.
$p' : \mathcal{S}' \to \mathcal{C}$ 与
$G : \mathcal{S} \to \mathcal{S}'$ 为引理
\ref{stacks-lemma-stackify-groupoids} 中构造的群胚中的栈与 $1$-态射。
该构造具有如下泛性质：给定群胚中的栈
$q : \mathcal{X} \to \mathcal{C}$，以及
$\mathcal{C}$ 上范畴之间的 $1$-态射
$F : \mathcal{S} \to \mathcal{X}$，存在
$1$-态射 $H : \mathcal{S}' \to \mathcal{X}$，使图式
$$
\xymatrix{
\mathcal{S} \ar[rr]_F \ar[rd]_G & & \mathcal{X} \\
& \mathcal{S}' \ar[ru]_H
}
$$
为 $2$-交换的。
\end{lemma}

\begin{proof}
这是引理 \ref{stacks-lemma-stackify-universal-property} 的特殊情形。
\end{proof}

\begin{lemma}
\label{stacks-lemma-stackification-fibre-product-categories-fibred-in-groupoids}
令 $\mathcal{C}$ 为位点。
令 $f : \mathcal{X} \to \mathcal{Y}$ 与
$g : \mathcal{Z} \to \mathcal{Y}$ 为
$\mathcal{C}$ 上纤维化于群胚的范畴之间的态射。
在这种情况下，$2$-纤维积的栈化是各栈化的
$2$-纤维积。
\end{lemma}

\begin{proof}
这是引理
\ref{stacks-lemma-stackification-fibre-product-fibred-categories}
的特殊情形。
\end{proof}










\section{继承拓扑}
\label{stacks-section-topology}

\noindent
事实证明，位点上的纤维化范畴从底位点继承一个典范拓扑。

\begin{lemma}
\label{stacks-lemma-topology-inherited}
令 $\mathcal{C}$ 为位点。令
$p : \mathcal{S} \to \mathcal{C}$ 为纤维化范畴。
令 $\text{Cov}(\mathcal{S})$ 为 $\mathcal{S}$ 中具有固定目标的
态射族 $\{x_i \to x\}_{i \in I}$ 的集合，并要求：
(a) 每个 $x_i \to x$ 都是强笛卡尔的；且
(b) $\{p(x_i) \to p(x)\}_{i \in I}$ 是
$\mathcal{C}$ 的覆盖。那么
$(\mathcal{S}, \text{Cov}(\mathcal{S}))$ 是位点。
\end{lemma}

\begin{proof}
需检验 Sites, Definition \ref{sites-definition-site} 的三个条件。
\begin{enumerate}
\item 若 $x \to y$ 是 $\mathcal{S}$ 的同构，则由
Categories, Lemma \ref{categories-lemma-composition-cartesian}，
它是强笛卡尔的，而 $p(x) \to p(y)$ 是
$\mathcal{C}$ 的同构。因此
$\{p(x) \to p(y)\}$ 是 $\mathcal{C}$ 的覆盖，从而
$\{x \to y\} \in \text{Cov}(\mathcal{S})$。
\item 若
$\{x_i \to x\}_{i\in I} \in \text{Cov}(\mathcal{S})$，
并且对每个 $i$ 都有
$\{y_{ij} \to x_i\}_{j\in J_i} \in \text{Cov}(\mathcal{S})$，
则由 Categories, Lemma
\ref{categories-lemma-composition-cartesian}，
每个复合 $y_{ij} \to x$ 都是强笛卡尔的，而且
$\{p(y_{ij}) \to p(x)\}_{i \in I, j\in J_i}
\in \text{Cov}(\mathcal{C})$。故也有
$\{y_{ij} \to x\}_{i \in I, j\in J_i}
\in \text{Cov}(\mathcal{S})$。
\item 假设
$\{x_i \to x\}_{i\in I}\in \text{Cov}(\mathcal{S})$，
且 $y \to x$ 是 $\mathcal{S}$ 的态射。
由于 $\{p(x_i) \to p(x)\}$ 是 $\mathcal{C}$ 的覆盖，
$p(x_i) \times_{p(x)} p(y)$ 存在。因此
Categories, Lemma
\ref{categories-lemma-fibred-category-representable-goes-up}
蕴含 $x_i \times_x y$ 存在、
$p(x_i \times_x y) = p(x_i) \times_{p(x)} p(y)$，
且 $x_i \times_x y \to y$ 是强笛卡尔的。
又因为
$\{p(x_i) \times_{p(x)} p(y) \to p(y) \}_{i\in I}
\in \text{Cov}(\mathcal{C})$，所以
$\{x_i \times_x y \to y \}_{i\in I}
\in \text{Cov}(\mathcal{S})$。
\end{enumerate}
证明完毕。
\end{proof}

\noindent
注意，若 $p : \mathcal{S} \to \mathcal{C}$ 纤维化于群胚，
则引理 \ref{stacks-lemma-topology-inherited} 中位点
$\mathcal{S}$ 的覆盖可刻画为
$$
\{x_i \to x\} \in \text{Cov}(\mathcal{S})
\Leftrightarrow
\{p(x_i) \to p(x)\} \in \text{Cov}(\mathcal{C})
$$
因为 $\mathcal{S}$ 的每个态射都是强笛卡尔的。

\begin{definition}
\label{stacks-definition-topology-inherited}
令 $\mathcal{C}$ 为位点。令
$p : \mathcal{S} \to \mathcal{C}$ 为纤维化范畴。
% P01-E0202: source omits “that” in “We say ... is”.
称引理 \ref{stacks-lemma-topology-inherited} 中的
$(\mathcal{S}, \text{Cov}(\mathcal{S}))$ 为
{\it $\mathcal{S}$ 上从 $\mathcal{C}$ 继承的位点结构}。
有时也把这表述为：{\it $\mathcal{S}$ 配备从
$\mathcal{C}$ 继承的拓扑}。
\end{definition}

\noindent
特别地，在这种情形得到层托扑斯 $\Sh(\mathcal{S})$。
事实证明，这个托扑斯关于纤维化范畴的 $1$-态射具有函子性。

\begin{lemma}
\label{stacks-lemma-topology-inherited-functorial}
令 $\mathcal{C}$ 为位点。令
$F : \mathcal{X} \to \mathcal{Y}$ 为
$\mathcal{C}$ 上纤维化范畴之间的 $1$-态射。
那么，在从 $\mathcal{C}$ 继承的位点结构之间，
$F$ 是连续且余连续的函子。因此，$F$ 诱导托扑斯之间的态射
$f : \Sh(\mathcal{X}) \to \Sh(\mathcal{Y})$，满足
$f_* = {}_sF = {}_pF$ 与 $f^{-1} = F^s = F^p$。
特别地，对 $\mathcal{Y}$ 上的层 $\mathcal{G}$
与 $\mathcal{X}$ 的对象 $x$，有
$f^{-1}(\mathcal{G})(x) = \mathcal{G}(F(x))$。
\end{lemma}

\begin{proof}
先证明 $F$ 连续。令
$\{x_i \to x\}_{i \in I}$ 为 $\mathcal{X}$ 的覆盖。
由 Categories, Definition
\ref{categories-definition-fibred-categories-over-C}，
函子 $F$ 把强笛卡尔态射映为强笛卡尔态射，故
$\{F(x_i) \to F(x)\}_{i \in I}$ 是 $\mathcal{Y}$ 的覆盖。
这证明了 Sites, Definition
\ref{sites-definition-continuous} 的 (1)。
此外，令 $x' \to x$ 为 $\mathcal{X}$ 的态射。
由 Categories, Lemma
\ref{categories-lemma-fibred-category-representable-goes-up}，
纤维积 $x_i \times_x x'$ 存在，且
$x_i \times_x x' \to x'$ 是强笛卡尔的。
因此 $F(x_i \times_x x') \to F(x')$ 是强笛卡尔的。
对 $\mathcal{Y}$ 应用 Categories, Lemma
\ref{categories-lemma-fibred-category-representable-goes-up}，
这意味着
$F(x_i \times_x x') = F(x_i) \times_{F(x)} F(x')$。
这证明了 Sites, Definition
\ref{sites-definition-continuous} 的 (2)，故 $F$ 连续。

\medskip\noindent
其次证明 $F$ 余连续。令 $x \in \Ob(\mathcal{X})$，并令
$\{y_i \to F(x)\}_{i \in I}$ 为 $\mathcal{Y}$ 中的覆盖。
以 $\{U_i \to U\}_{i \in I}$ 表示 $\mathcal{C}$
中相应的覆盖。对每个 $i$，选取 $\mathcal{X}$ 中位于
$U_i \to U$ 上的强笛卡尔态射 $x_i \to x$。
% P01-E0203: source has “both a strongly cartesian morphisms”.
那么 $F(x_i) \to F(x)$ 与 $y_i \to F(x)$ 都是
$\mathcal{Y}$ 中位于 $U_i \to U$ 上的强笛卡尔态射。
因此，在 $\mathcal{Y}_{U_i}$ 中存在唯一的同构
$F(x_i) \to y_i$，并与到 $F(x)$ 的映射相容。
所以 $\{x_i \to x\}_{i \in I}$ 是 $\mathcal{X}$ 的覆盖，
而 $\{F(x_i) \to F(x)\}_{i \in I}$ 同构于
$\{y_i \to F(x)\}_{i \in I}$。故 $F$ 余连续，见
Sites, Definition \ref{sites-definition-cocontinuous}。

\medskip\noindent
最后一个断言由前两个断言得出，见 Sites, Lemmas
\ref{sites-lemma-cocontinuous-morphism-topoi}、
\ref{sites-lemma-pu-sheaf} 与
\ref{sites-lemma-when-shriek}。
\end{proof}

\begin{lemma}
\label{stacks-lemma-localizing}
令 $\mathcal{C}$ 为位点。令
$p : \mathcal{X} \to \mathcal{C}$ 为纤维化于群胚的范畴。
令 $x \in \Ob(\mathcal{X})$ 位于
$U = p(x)$ 上。函子 $p$ 诱导位点之间的等价
$\mathcal{X}/x \to \mathcal{C}/U$，其中
$\mathcal{X}$ 配备从 $\mathcal{C}$ 继承的拓扑。
\end{lemma}

\begin{proof}
这里，$\mathcal{C}/U$ 是位点 $\mathcal{C}$ 在对象
$U$ 处的局部化，$\mathcal{X}/x$ 亦然。
由 Categories, Definition
\ref{categories-definition-fibred-groupoids} 可知，规则
$x'/x \mapsto p(x')/p(x)$ 定义范畴之间的等价
$\mathcal{X}/x \to \mathcal{C}/U$。于是，由
Definition \ref{stacks-definition-topology-inherited}，
$\mathcal{X}/x$ 中 $x'$ 的覆盖与
$\mathcal{C}/U$ 中 $p(x')$ 的覆盖一一对应。
\end{proof}

\begin{lemma}
\label{stacks-lemma-stack-in-groupoids-over-stack-in-groupoids}
令 $\mathcal{C}$ 为位点。令
$p : \mathcal{X} \to \mathcal{C}$ 与
$q : \mathcal{Y} \to \mathcal{C}$ 为群胚中的栈。
令 $F : \mathcal{X} \to \mathcal{Y}$ 为
$\mathcal{C}$ 上范畴之间的 $1$-态射。
若 $F$ 使 $\mathcal{X}$ 成为
$\mathcal{Y}$ 上纤维化于群胚的范畴，则
$\mathcal{X}$ 是 $\mathcal{Y}$ 上群胚中的栈
（配备从 $\mathcal{C}$ 继承的拓扑）。
\end{lemma}

\begin{proof}
来证明对象的下降。令 $\{y_i \to y\}$ 为
$\mathcal{Y}$ 的覆盖。令 $(x_i, \varphi_{ij})$ 为
$\mathcal{X}$ 中相对于该覆盖的下降数据。
那么 $(x_i, \varphi_{ij})$ 也是相对于
$\mathcal{C}$ 的覆盖 $\{q(y_i) \to q(y)\}$ 的下降数据。
由于 $\mathcal{X}$ 是群胚中的栈，得到位于
$q(y)$ 上的对象 $x$ 以及位于 $q(y_i)$ 上、与
$\varphi_{ij}$ 相容的同构
$\psi_i : x|_{q(y_i)} \to x_i$；也就是说，
$$
\varphi_{ij} = \psi_j|_{q(y_i) \times_{q(y)} q(y_j)}
\circ \psi_i^{-1}|_{q(y_i) \times_{q(y)} q(y_j)}.
$$
考虑 $\mathcal{C}/p(x)$ 上的层
$\mathit{I} = \mathit{Isom}_\mathcal{Y}(F(x), y)$。
% P01-E0204: q is not defined on x_i; p(x_i)=q(y_i) is the typed base.
注意，$s_i = F(\psi_i) \in \mathit{I}(q(x_i))$，
因为 $F(x_i) = y_i$。由于
$F(\varphi_{ij}) = \text{id}$（因为我们从
$\{y_i \to y\}$ 上的下降数据出发），上述公式表明
$s_i|_{q(y_i) \times_{q(y)} q(y_j)}
= s_j|_{q(y_i) \times_{q(y)} q(y_j)}$。
因此局部截面 $s_i$ 黏合为 $s : F(x) \to y$。
由于 $F$ 纤维化于群胚，$x$ 同构于某个满足
$F(x') = y$ 的对象 $x'$。
我们略去验证，$\mathcal{X}$ 在 $y$ 上的纤维范畴中的
$x'$ 是由下降数据 $(x_i, \varphi_{ij})$ 提出的下降问题的解。
还略去证明 $\mathcal{X}/\mathcal{Y}$ 的
$\mathit{Isom}$-预层具有层性质。
\end{proof}

\begin{lemma}
\label{stacks-lemma-stack-over-stack}
令 $\mathcal{C}$ 为位点。令
$p : \mathcal{X} \to \mathcal{C}$ 为栈。
为 $\mathcal{X}$ 配备从 $\mathcal{C}$ 继承的拓扑，
并令 $q : \mathcal{Y} \to \mathcal{X}$ 为栈。
那么 $\mathcal{Y}$ 是 $\mathcal{C}$ 上的栈。
若 $p$ 与 $q$ 定义群胚中的栈，则
$\mathcal{Y}$ 是 $\mathcal{C}$ 上群胚中的栈。
\end{lemma}

\begin{proof}
检验 Definition \ref{stacks-definition-stack} 的三个条件，以证明
$\mathcal{Y}$ 是 $\mathcal{C}$ 上的栈。
由 Categories, Lemma
\ref{categories-lemma-fibred-over-fibred} 可知，
$\mathcal{Y}$ 是 $\mathcal{C}$ 上的纤维化范畴。
所以条件 (1) 成立。

\medskip\noindent
令 $U$ 为 $\mathcal{C}$ 的对象，并令
$y_1, y_2$ 为 $\mathcal{Y}$ 中位于 $U$ 上的对象。
在 $\mathcal{X}$ 中记 $x_i = q(y_i)$。
考虑 $\mathcal{C}/U$ 上的预层映射
$$
q : \mathit{Mor}_{\mathcal{Y}/\mathcal{C}}(y_1, y_2)
\longrightarrow
\mathit{Mor}_{\mathcal{X}/\mathcal{C}}(x_1, x_2)
$$
见引理 \ref{stacks-lemma-presheaf-mor-map-fibred-categories}。
令 $\{U_i \to U\}$ 为覆盖，并令 $\varphi_i$
为左侧预层在 $U_i$ 上的截面，使 $\varphi_i$ 与
$\varphi_j$ 在 $U_i \times_U U_j$ 上限制为同一截面。
需找到限制为 $\varphi_i$ 的态射
$\varphi : y_1 \to y_2$。
注意，由于第二个预层是层（根据假设），对某个
位于 $U$ 上的态射 $\psi : x_1 \to x_2$，有
$q(\varphi_i) = \psi|_{U_i}$。
令 $y_{12} \to y_2$ 为 $\mathcal{Y}$ 中位于
$\psi$ 上的强 $\mathcal{X}$-笛卡尔态射。
那么 $\varphi_i$ 对应于位于 $x_1|_{U_i}$ 上的态射
$\varphi'_i : y_1|_{U_i} \to y_{12}|_{U_i}$。
换言之，$\varphi'_i$ 现在定义预层
$$
\mathit{Mor}_{\mathcal{Y}/\mathcal{X}}(y_1, y_{12})
$$
在覆盖 $\{x_1|_{U_i} \to x_1\}$ 各成员上的局部截面。
由假设，它们黏合为唯一的态射
$y_1 \to y_{12}$；将其与给定态射
$y_{12} \to y_2$ 复合，便得到所需态射
$y_1 \to y_2$。

\medskip\noindent
最后，证明下降数据有效。令
$\{f_i : U_i \to U\}$ 为 $\mathcal{C}$ 的覆盖，并令
$(y_i, \varphi_{ij})$ 为相对于该覆盖的下降数据
（Definition \ref{stacks-definition-descent-data}）。
令 $x_i = q(y_i)$、$\psi_{ij} = q(\varphi_{ij})$，
便得到 $\mathcal{X}$ 中相对于该覆盖的下降数据
$(x_i, \psi_{ij})$。由对 $\mathcal{X}$ 的假设，
可假设 $x_i = x|_{U_i}$，而 $\psi_{ij}$ 等于典范下降数据
（Definition \ref{stacks-definition-effective-descent-datum}）。
在此情形，$\{x|_{U_i} \to x\}$ 是覆盖，并可把
$(y_i, \varphi_{ij})$ 看作相对于该覆盖的下降数据。
% P01-E0205: source says Y is a stack over C, which is the conclusion;
% the standing hypothesis is that q:Y→X is a stack over X.
由 $\mathcal{Y}$ 是 $\mathcal{C}$ 上栈这一假设，
可知该下降数据有效，从而完成条件 (3) 的证明。

\medskip\noindent
最后一个断言成立，是因为 $\mathcal{Y}$ 是
$\mathcal{C}$ 上的栈，并由 Categories, Lemma
\ref{categories-lemma-fibred-in-groupoids-over-fibred-in-groupoids}
纤维化于群胚。
\end{proof}





\section{Gerbe（胚）}
\label{stacks-section-gerbes}

\noindent
Gerbe（胚）是群胚中的栈的一种特殊类型。

\begin{definition}
\label{stacks-definition-gerbe}
位点 $\mathcal{C}$ 上的一个{\it Gerbe（胚）}，是
$\mathcal{C}$ 上的范畴
$p : \mathcal{S} \to \mathcal{C}$，满足：
\begin{enumerate}
\item $p : \mathcal{S} \to \mathcal{C}$ 是
$\mathcal{C}$ 上群胚中的栈（见
Definition \ref{stacks-definition-stack-in-groupoids}）；
\item 对 $U \in \Ob(\mathcal{C})$，存在
$\mathcal{C}$ 中的覆盖 $\{U_i \to U\}$，使
$\mathcal{S}_{U_i}$ 非空；且
\item 对 $U \in \Ob(\mathcal{C})$ 与
$x, y \in \Ob(\mathcal{S}_U)$，存在
$\mathcal{C}$ 中的覆盖 $\{U_i \to U\}$，使得在
$\mathcal{S}_{U_i}$ 中有
$x|_{U_i} \cong y|_{U_i}$。
\end{enumerate}
\end{definition}

\noindent
换言之，胚是这样的群胚中的栈：任意两个对象局部同构，
并且对象局部存在。

\begin{lemma}
\label{stacks-lemma-gerbe-equivalent}
令 $\mathcal{C}$ 为位点。
令 $\mathcal{S}_1$、$\mathcal{S}_2$ 为
$\mathcal{C}$ 上的范畴。假设 $\mathcal{S}_1$ 与
$\mathcal{S}_2$ 作为 $\mathcal{C}$ 上的范畴彼此等价。
那么 $\mathcal{S}_1$ 是 $\mathcal{C}$ 上的胚，
当且仅当 $\mathcal{S}_2$ 是 $\mathcal{C}$ 上的胚。
\end{lemma}

\begin{proof}
假设 $\mathcal{S}_1$ 是 $\mathcal{C}$ 上的胚。
由引理 \ref{stacks-lemma-stack-in-groupoids-equivalent} 可知，
$\mathcal{S}_2$ 是 $\mathcal{C}$ 上群胚中的栈。
令 $F : \mathcal{S}_1 \to \mathcal{S}_2$、
$G : \mathcal{S}_2 \to \mathcal{S}_1$ 为
$\mathcal{C}$ 上范畴之间的等价。
给定 $U \in \Ob(\mathcal{C})$，存在覆盖
$\{U_i \to U\}$，使 $(\mathcal{S}_1)_{U_i}$ 非空。
应用 $F$ 可知 $(\mathcal{S}_2)_{U_i}$ 非空。
给定 $U \in \Ob(\mathcal{C})$ 与
$x, y \in \Ob((\mathcal{S}_2)_U)$，存在
$\mathcal{C}$ 中的覆盖 $\{U_i \to U\}$，使得在
$(\mathcal{S}_1)_{U_i}$ 中有
$G(x)|_{U_i} \cong G(y)|_{U_i}$。
由 Categories, Lemma
\ref{categories-lemma-equivalence-fibred-categories}，
这蕴含在 $(\mathcal{S}_2)_{U_i}$ 中有
$x|_{U_i} \cong y|_{U_i}$。
\end{proof}

\noindent
我们希望稍微推广胚的定义。具体地，令
$F : \mathcal{X} \to \mathcal{Y}$ 为位点
$\mathcal{C}$ 上群胚中的栈之间的 $1$-态射。
我们想说明 $\mathcal{X}$ 是 $\mathcal{Y}$ 上的胚
意味着什么。由 Section \ref{stacks-section-topology}，
范畴 $\mathcal{Y}$ 从 $\mathcal{C}$ 继承位点结构。
一个朴素的猜测是：只要求
$\mathcal{X} \to \mathcal{Y}$ {\it 是}上述意义下的胚。
可是，如此得到的概念在用 $\mathcal{C}$ 上群胚中的等价栈替换
$\mathcal{X}$ 时并不保持不变；即使只是
$\mathcal{Y}$ 上纤维化于群胚这一性质，也会如此。
不过，事实证明，可以用一个在 $\mathcal{Y}$ 上纤维化于群胚的、
$\mathcal{C}$ 上群胚中的等价栈替换 $\mathcal{X}$；
此时，在 $\mathcal{Y}$ 上为胚这一性质与替换的选择无关。
精确表述如下。

\begin{lemma}
\label{stacks-lemma-when-gerbe}
令 $\mathcal{C}$ 为位点。令
$p : \mathcal{X} \to \mathcal{C}$ 与
$q : \mathcal{Y} \to \mathcal{C}$ 为群胚中的栈。
令 $F : \mathcal{X} \to \mathcal{Y}$ 为
$\mathcal{C}$ 上范畴之间的 $1$-态射。下列条件等价：
\begin{enumerate}
\item 对某个（等价地，任意）分解
$F = F' \circ a$，其中
$a : \mathcal{X} \to \mathcal{X}'$ 是
$\mathcal{C}$ 上范畴之间的等价，而 $F'$ 纤维化于群胚，
映射 $F' : \mathcal{X}' \to \mathcal{Y}$ 是胚
（$\mathcal{Y}$ 配备从 $\mathcal{C}$ 继承的拓扑）。
\item 满足下列两个条件：
\begin{enumerate}
\item 对位于 $U \in \Ob(\mathcal{C})$ 上的
$y \in \Ob(\mathcal{Y})$，存在
$\mathcal{C}$ 中的覆盖 $\{U_i \to U\}$ 以及
$\mathcal{X}$ 中位于 $U_i$ 上的对象 $x_i$，
使得在 $\mathcal{Y}_{U_i}$ 中有
$F(x_i) \cong y|_{U_i}$；且
\item 对 $U \in \Ob(\mathcal{C})$、
$x, x' \in \Ob(\mathcal{X}_U)$ 以及
$\mathcal{Y}_U$ 中的 $b : F(x) \to F(x')$，
存在 $\mathcal{C}$ 中的覆盖 $\{U_i \to U\}$
以及 $\mathcal{X}_{U_i}$ 中的态射
$a_i : x|_{U_i} \to x'|_{U_i}$，满足
$F(a_i) = b|_{U_i}$。
\end{enumerate}
\end{enumerate}
\end{lemma}

\begin{proof}
由 Categories, Lemma
\ref{categories-lemma-ameliorate-morphism-categories-fibred-groupoids}，
存在分解 $F = F' \circ a$，其中
$a : \mathcal{X} \to \mathcal{X}'$ 是
$\mathcal{C}$ 上范畴之间的等价，而 $F'$ 纤维化于群胚。
由 Categories, Lemma
\ref{categories-lemma-amelioration-unique}，
给定任意两个这样的分解
$F = F' \circ a = F'' \circ b$，
$\mathcal{X}'$ 与 $\mathcal{X}''$ 作为
$\mathcal{Y}$ 上的范畴彼此等价。
因此，引理 \ref{stacks-lemma-gerbe-equivalent} 保证条件 (1)
与分解的选择无关。此外，这意味着可以假设
$\mathcal{X}' =
\mathcal{X} \times_{F, \mathcal{Y}, \text{id}} \mathcal{Y}$，
如 Categories, Lemma
\ref{categories-lemma-ameliorate-morphism-categories-fibred-groupoids}
的证明中所述。

\medskip\noindent
来证明 (a)、(b) 蕴含
$\mathcal{X}' \to \mathcal{Y}$ 是胚。
首先，由引理
\ref{stacks-lemma-stack-in-groupoids-over-stack-in-groupoids} 可知，
$\mathcal{X}' \to \mathcal{Y}$ 是群胚中的栈。
其次，令 $y$ 为 $\mathcal{Y}$ 中位于
$U \in \Ob(\mathcal{C})$ 上的对象。
由 (a)，可找到 $\mathcal{C}$ 中的覆盖
$\{U_i \to U\}$、$\mathcal{X}$ 中位于 $U_i$
上的对象 $x_i$，以及 $\mathcal{Y}_{U_i}$ 中的同构
$f_i : F(x_i) \to y|_{U_i}$。
那么 $(U_i, x_i, y|_{U_i}, f_i)$ 是
$\mathcal{X}'_{U_i}$ 的对象；即 Definition
\ref{stacks-definition-gerbe} 的第二个条件成立。
最后，令 $(U, x, y, f)$ 与 $(U, x', y, f')$
为 $\mathcal{X}'$ 中位于同一对象
$y \in \Ob(\mathcal{Y})$ 上的对象。
令 $b = (f')^{-1} \circ f$。
由条件 (b)，可找到覆盖 $\{U_i \to U\}$ 以及
$\mathcal{X}_{U_i}$ 中满足
$F(a_i) = b|_{U_i}$ 的同构
$a_i : x|_{U_i} \to x'|_{U_i}$。
那么
$$
(a_i, \text{id}) : (U, x, y, f)|_{U_i} \to (U, x', y, f')|_{U_i}
$$
是 $\mathcal{X}'_{U_i}$ 中所需的态射。
这证明 (2) 蕴含 (1)。

\medskip\noindent
为证明 (1) 蕴含 (2)，反向阅读上一段的论证即可。细节从略。
\end{proof}

\begin{definition}
\label{stacks-definition-gerbe-over-stack-in-groupoids}
令 $\mathcal{C}$ 为位点。令 $\mathcal{X}$ 与
$\mathcal{Y}$ 为 $\mathcal{C}$ 上群胚中的栈。
令 $F : \mathcal{X} \to \mathcal{Y}$ 为
$\mathcal{C}$ 上范畴之间的 $1$-态射。
若满足引理 \ref{stacks-lemma-when-gerbe} 中的等价条件，
则称 $\mathcal{X}$ 是 $\mathcal{Y}$ {\it 上的胚}。
\end{definition}

\noindent
当 $\mathcal{Y} = \mathcal{C}$ 时，这一定义与
Definition \ref{stacks-definition-gerbe} 不冲突，因为在
Lemma \ref{stacks-lemma-when-gerbe} 的 (1) 中可取
$\mathcal{X}' = \mathcal{X}$。
注意，Lemma \ref{stacks-lemma-when-gerbe} 的条件 (2)(a)、(2)(b)
在精神上非常接近 Definition \ref{stacks-definition-gerbe} 的
条件 (2)、(3)。具体而言，(2)(a) 说对象同构类的预层之间的映射
在层化后成为满射。此外，(2)(b) 说
$$
\mathit{Isom}_\mathcal{X}(x, x')
\longrightarrow
\mathit{Isom}_\mathcal{Y}(F(x), F(x'))
$$
对任意 $U$ 与 $x, x' \in \Ob(\mathcal{X}_U)$，
都是 $\mathcal{C}/U$ 上层之间的满射。

\begin{lemma}
\label{stacks-lemma-base-change-gerbe}
令 $\mathcal{C}$ 为位点。令
$$
\xymatrix{
\mathcal{X}' \ar[r]_{G'} \ar[d]_{F'} & \mathcal{X} \ar[d]^F \\
\mathcal{Y}' \ar[r]^G & \mathcal{Y}
}
$$
为 $\mathcal{C}$ 上群胚中的栈所组成的 $2$-纤维积。
若 $\mathcal{X}$ 是 $\mathcal{Y}$ 上的胚，则
$\mathcal{X}'$ 是 $\mathcal{Y}'$ 上的胚。
\end{lemma}

\begin{proof}
% P01-E0206: source omits “we” before “may assume”.
由 $2$-纤维积的唯一性性质，可假设
$\mathcal{X}' = \mathcal{Y}' \times_\mathcal{Y} \mathcal{X}$，
如 Categories, Lemma
\ref{categories-lemma-2-product-categories-over-C} 中所述。
来证明引理 \ref{stacks-lemma-when-gerbe} 的性质 (2)(a)、(2)(b)
对 $\mathcal{Y}' \times_\mathcal{Y} \mathcal{X}
\to \mathcal{Y}'$ 成立。

\medskip\noindent
令 $y'$ 为 $\mathcal{Y}'$ 中位于
$\mathcal{C}$ 的对象 $U$ 上的对象。
由假设，存在 $U$ 的覆盖 $\{U_i \to U\}$
以及对象 $x_i \in \mathcal{X}_{U_i}$，并有同构
$\alpha_i : G(y')|_{U_i} \to F(x_i)$。
那么 $(U_i, y'|_{U_i}, x_i, \alpha_i)$ 是
$\mathcal{Y}' \times_\mathcal{Y} \mathcal{X}$ 中位于
$U_i$ 上的对象，其在 $\mathcal{Y}'$ 中的像为
$y'|_{U_i}$。因此 (2)(a) 成立。

\medskip\noindent
令 $U \in \Ob(\mathcal{C})$，令
$x'_1, x'_2$ 为
$\mathcal{Y}' \times_\mathcal{Y} \mathcal{X}$
中位于 $U$ 上的对象，并令
$b' : F'(x'_1) \to F'(x'_2)$ 为
$\mathcal{Y}'_U$ 中的态射。写
$x'_i = (U, y'_i, x_i, \alpha_i)$。
% P01-E0208: the source reverses the two projections F' and G'.
注意，$F'(x'_i) = x_i$ 且 $G'(x'_i) = y'_i$。
由假设，存在 $\mathcal{C}$ 中的覆盖
$\{U_i \to U\}$，以及 $\mathcal{X}_{U_i}$ 中满足
$F(a_i) = G(b')|_{U_i}$ 的态射
$a_i : x_1|_{U_i} \to x_2|_{U_i}$。
那么 $(b'|_{U_i}, a_i)$ 是 (2)(b) 所要求的态射
$x'_1|_{U_i} \to x'_2|_{U_i}$。
\end{proof}

\begin{lemma}
\label{stacks-lemma-composition-gerbe}
令 $\mathcal{C}$ 为位点。令
$F : \mathcal{X} \to \mathcal{Y}$ 与
$G : \mathcal{Y} \to \mathcal{Z}$ 为
$\mathcal{C}$ 上群胚中的栈之间的 $1$-态射。
若 $\mathcal{X}$ 是 $\mathcal{Y}$ 上的胚，且
$\mathcal{Y}$ 是 $\mathcal{Z}$ 上的胚，则
$\mathcal{X}$ 是 $\mathcal{Z}$ 上的胚。
\end{lemma}

\begin{proof}
来证明引理 \ref{stacks-lemma-when-gerbe} 的性质 (2)(a)、(2)(b)
对 $\mathcal{X} \to \mathcal{Z}$ 成立。

\medskip\noindent
令 $z$ 为 $\mathcal{Z}$ 中位于
$\mathcal{C}$ 的对象 $U$ 上的对象。
由对 $G$ 的假设，存在 $U$ 的覆盖
$\{U_i \to U\}$ 与对象
$y_i \in \mathcal{Y}_{U_i}$，使
$G(y_i) \cong z|_{U_i}$。
由对 $F$ 的假设，存在覆盖
$\{U_{ij} \to U_i\}$ 与对象
$x_{ij} \in \mathcal{X}_{U_{ij}}$，使
$F(x_{ij}) \cong y_i|_{U_{ij}}$。
那么 $\{U_{ij} \to U\}$ 是 $\mathcal{C}$ 的覆盖，且
$(G \circ F)(x_{ij}) \cong z|_{U_{ij}}$。
因此 (2)(a) 成立。

\medskip\noindent
令 $U \in \Ob(\mathcal{C})$，令 $x_1, x_2$
为 $\mathcal{X}$ 中位于 $U$ 上的对象，并令
$c : (G \circ F)(x_1) \to (G \circ F)(x_2)$ 为
$\mathcal{Z}_U$ 中的态射。由对 $G$ 的假设，
存在 $U$ 的覆盖 $\{U_i \to U\}$ 以及
$\mathcal{Y}_{U_i}$ 中满足
$G(b_i) = c|_{U_i}$ 的态射
$b_i : F(x_1)|_{U_i} \to F(x_2)|_{U_i}$。
由对 $F$ 的假设，存在覆盖
$\{U_{ij} \to U_i\}$ 以及
$\mathcal{X}_{U_{ij}}$ 中满足
$F(a_{ij}) = b_i|_{U_{ij}}$ 的态射
$a_{ij} : x_1|_{U_{ij}} \to x_2|_{U_{ij}}$。
那么 $\{U_{ij} \to U\}$ 是 $\mathcal{C}$ 的覆盖，且
$(G \circ F)(a_{ij}) = c|_{U_{ij}}$，正如 (2)(b) 所要求。
\end{proof}

\begin{lemma}
\label{stacks-lemma-gerbe-descent}
令 $\mathcal{C}$ 为位点。令
$$
\xymatrix{
\mathcal{X}' \ar[r]_{G'} \ar[d]_{F'} & \mathcal{X} \ar[d]^F \\
\mathcal{Y}' \ar[r]^G & \mathcal{Y}
}
$$
为 $\mathcal{C}$ 上群胚中的栈所组成的 $2$-笛卡尔图式。
若对每个 $U \in \Ob(\mathcal{C})$ 与
$x \in \Ob(\mathcal{Y}_U)$，都存在覆盖
$\{U_i \to U\}$，使得 $x|_{U_i}$ 属于
$G : \mathcal{Y}'_{U_i} \to \mathcal{Y}_{U_i}$
的本质像，并且 $\mathcal{X}'$ 是
$\mathcal{Y}'$ 上的胚，则 $\mathcal{X}$ 是
$\mathcal{Y}$ 上的胚。
\end{lemma}

\begin{proof}
% P01-E0207: source omits “we” before “may assume”.
由 $2$-纤维积的唯一性性质，可假设
$\mathcal{X}' = \mathcal{Y}' \times_\mathcal{Y} \mathcal{X}$，
如 Categories, Lemma
\ref{categories-lemma-2-product-categories-over-C} 中所述。
来证明引理 \ref{stacks-lemma-when-gerbe} 的性质 (2)(a)、(2)(b)
对 $\mathcal{X} \to \mathcal{Y}$ 成立。

\medskip\noindent
令 $y$ 为 $\mathcal{Y}$ 中位于
$\mathcal{C}$ 的对象 $U$ 上的对象。
由假设，存在 $U$ 的覆盖 $\{U_i \to U\}$ 与对象
$y'_i \in \mathcal{Y}'_{U_i}$，使
$G(y'_i) \cong y|_{U_i}$。
由 $\mathcal{X}' \to \mathcal{Y}'$ 的 (2)(a)，
存在覆盖 $\{U_{ij} \to U_i\}$ 以及
$\mathcal{X}'$ 中位于 $U_{ij}$ 上的对象 $x'_{ij}$，
% P01-E0209: source uses the ungrammatical “with ... isomorphic”.
使 $F'(x'_{ij})$ 同构于 $y'_i$ 在
$U_{ij}$ 上的限制。那么
$\{U_{ij} \to U\}$ 是 $\mathcal{C}$ 的覆盖，
而 $G'(x'_{ij})$ 是 $\mathcal{X}$ 中位于
$U_{ij}$ 上的对象，其在 $\mathcal{Y}$ 中的像同构于
$y|_{U_{ij}}$。这证明了
$\mathcal{X} \to \mathcal{Y}$ 的 (2)(a)。

\medskip\noindent
令 $U \in \Ob(\mathcal{C})$，令 $x_1, x_2$
为 $\mathcal{X}$ 中位于 $U$ 上的对象，并令
$b : F(x_1) \to F(x_2)$ 为 $\mathcal{Y}_U$ 中的态射。
由假设，可选取覆盖 $\{U_i \to U\}$ 以及
$\mathcal{Y}'$ 中位于 $U_i$ 上的对象 $y'_i$，
使存在同构
$\alpha_i : G(y'_i) \to F(x_1)|_{U_i}$。
于是得到 $\mathcal{X}'$ 中位于 $U_i$ 上的对象
$$
x'_{1i} = (U_i, y'_i, x_1|_{U_i}, \alpha_i)
\quad\text{且}\quad
x'_{2i} = (U_i, y'_i, x_2|_{U_i}, b|_{U_i} \circ \alpha_i)
$$
$y'_i$ 上的恒等态射是态射
$F'(x'_{1i}) \to F'(x'_{2i})$。
由 $\mathcal{X}' \to \mathcal{Y}'$ 的 (2)(b)，
存在覆盖 $\{U_{ij} \to U_i\}$ 与满足
$F'(a'_{ij}) = \text{id}_{y'_i}|_{U_{ij}}$ 的态射
$a'_{ij} : x'_{1i}|_{U_{ij}} \to x'_{2i}|_{U_{ij}}$。
展开 $\mathcal{Y}' \times_\mathcal{Y} \mathcal{X}$
中态射的定义，可知
$G'(a'_{ij}) : x_1|_{U_{ij}} \to x_2|_{U_{ij}}$
正是所求态射；即
$\mathcal{X} \to \mathcal{Y}$ 的 (2)(b) 成立。
\end{proof}

\noindent
所有自同构层均为阿贝尔层的胚在代数几何中发挥重要作用。

\begin{lemma}
\label{stacks-lemma-gerbe-abelian-auts}
令 $p : \mathcal{S} \to \mathcal{C}$ 为位点
$\mathcal{C}$ 上的胚。假设对所有
$U \in \Ob(\mathcal{C})$ 与
$x \in \Ob(\mathcal{S}_U)$，$\mathcal{C}/U$
上的群层
$\mathit{Aut}(x) = \mathit{Isom}(x, x)$ 都是阿贝尔的。
那么存在
\begin{enumerate}
\item $\mathcal{C}$ 上的阿贝尔群层 $\mathcal{G}$；
\item 对每个 $U \in \Ob(\mathcal{C})$ 与每个
$x \in \Ob(\mathcal{S}_U)$，一个同构
$\mathcal{G}|_U \to \mathit{Aut}(x)$，
\end{enumerate}
使得对每个 $U$ 以及 $\mathcal{S}_U$ 中的每个态射
$\varphi : x \to y$，图式
$$
\xymatrix{
\mathcal{G}|_U \ar[d] \ar@{=}[rr] & & \mathcal{G}|_U \ar[d] \\
\mathit{Aut}(x)
\ar[rr]^{\alpha \mapsto \varphi \circ \alpha \circ \varphi^{-1}} & &
\mathit{Aut}(y)
}
$$
交换。
\end{lemma}

\begin{proof}
令 $x, y$ 为 $\mathcal{S}$ 的两个对象，且
$U = p(x) = p(y)$。

\medskip\noindent
若存在位于 $U$ 上的态射 $\varphi : x \to y$，
则它是同构，因而确实得到同构
$\mathit{Aut}(x) \to \mathit{Aut}(y)$，它把
$\alpha$ 送到
$\varphi \circ \alpha \circ \varphi^{-1}$。
此外，因为假设 $\mathit{Aut}(x)$ 可交换，
一个简单计算表明该同构与 $\varphi$ 的选择无关：
若 $\psi$ 是第二个这样的映射，则
$$
\varphi \circ \alpha \circ \varphi^{-1} =
\psi \circ \psi^{-1} \circ \varphi \circ \alpha \circ \varphi^{-1} =
\psi \circ \alpha \circ \psi^{-1} \circ \varphi \circ \varphi^{-1} =
\psi \circ \alpha \circ \psi^{-1}
$$
其结果是层之间的典范同构
$\mathit{Aut}(x) \to \mathit{Aut}(y)$。
再者，若有第三个对象 $z$ 以及态射 $y \to z$
（因而也有态射 $x \to z$），则典范同构
$\mathit{Aut}(x) \to \mathit{Aut}(y)$、
$\mathit{Aut}(y) \to \mathit{Aut}(z)$ 与
$\mathit{Aut}(x) \to \mathit{Aut}(z)$ 相容，其意义是图式
$$
\xymatrix{
\mathit{Aut}(x) \ar[rd] \ar[rr] & &  \mathit{Aut}(z) \\
& \mathit{Aut}(y) \ar[ru]
}
$$
交换。

\medskip\noindent
若不存在位于 $U$ 上从 $x$ 到 $y$ 的态射，
则可选取覆盖 $\{U_i \to U\}$，使得存在态射
$x|_{U_i} \to y|_{U_i}$。这给出典范同构
$$
\mathit{Aut}(x)|_{U_i} \longrightarrow \mathit{Aut}(y)|_{U_i}
$$
它们在 $U_i \times_U U_j$ 上一致（由典范性）。
由层的黏合（Sites, Lemma \ref{sites-lemma-glue-maps}），
得到唯一同构
$\mathit{Aut}(x) \to \mathit{Aut}(y)$，其在任意
$U_i$ 上的限制均为上一段的典范同构。
与上面类似，若有位于 $U$ 上的第三个对象，
这些典范同构满足相容性。

\medskip\noindent
如果 $\mathcal{S}$ 在 $U$ 上的纤维范畴为空，怎么办？
在这种情况下，可以找到覆盖 $\{U_i \to U\}$
以及 $\mathcal{S}$ 中位于 $U_i$ 上的对象 $x_i$。
令 $\mathcal{G}_i = \mathit{Aut}(x_i)$。
由上述论证，得到典范同构
$$
\varphi_{ij} :
\mathcal{G}_i|_{U_i \times_U U_j}
\longrightarrow
\mathcal{G}_j|_{U_i \times_U U_j}
$$
它们在 $U_i \times_U U_j \times_U U_k$ 上的限制满足
Sites, Section \ref{sites-section-glueing-sheaves}
中说明的余圈条件。
由 Sites, Lemma \ref{sites-lemma-glue-sheaves}，
得到 $U$ 上的层 $\mathcal{G}$；它在 $U_i$ 上的限制
以与黏合映射 $\varphi_{ij}$ 相容的方式给出
$\mathcal{G}_i$。

\medskip\noindent
若 $\mathcal{C}$ 有终对象 $U$，证明便告完成，
因为可以取 $\mathcal{G}$ 为刚才构造的层。
在一般情形，还需验证在不同 $U$ 上构造的层
$\mathcal{G}$ 以典范方式相容。此验证从略。
\end{proof}








\section{栈的函子性}
\label{stacks-section-inverse-image}

\noindent
本节研究改变栈的基位点时会发生什么。初次阅读时可以跳过本节。

\medskip\noindent
令 $u : \mathcal{C} \to \mathcal{D}$ 为范畴之间的函子。
令 $p : \mathcal{S} \to \mathcal{D}$ 为
$\mathcal{D}$ 上的范畴。在这种情况下，以
$u^p\mathcal{S}$ 表示如下定义的 $\mathcal{C}$ 上的范畴：
\begin{enumerate}
\item $u^p\mathcal{S}$ 的对象是偶对 $(U, y)$，其中
$U$ 是 $\mathcal{C}$ 的对象，而
$y$ 是 $\mathcal{S}_{u(U)}$ 的对象。
\item 态射 $(a, \beta) : (U, y) \to (U', y')$ 由
$\mathcal{C}$ 的态射 $a : U \to U'$ 与
$\mathcal{S}$ 的态射 $\beta : y \to y'$ 给出，并满足
$p(\beta) = u(a)$。
\end{enumerate}
注意，按这些定义，$u^p\mathcal{S}$ 在 $U$ 上的纤维范畴
等于 $\mathcal{S}$ 在 $u(U)$ 上的纤维范畴。

\begin{lemma}
\label{stacks-lemma-fibred-category-pushforward}
在上述情形中，若 $\mathcal{S}$ 是 $\mathcal{D}$ 上的纤维化范畴，
则 $u^p\mathcal{S}$ 是 $\mathcal{C}$ 上的纤维化范畴。
\end{lemma}

\begin{proof}
在阅读此证明前，请先参阅 Categories, Definitions
\ref{categories-definition-cartesian-over-C} 与
\ref{categories-definition-fibred-category} 周围的讨论。
令 $(a, \beta) : (U, y) \to (U', y')$ 为
$u^p\mathcal{S}$ 的态射。我们断言，
$(a, \beta)$ 是强笛卡尔的，当且仅当 $\beta$ 是强笛卡尔的。
先假设 $\beta$ 强笛卡尔。考虑
$u^p\mathcal{S}$ 的任意第二个态射
$(a_1, \beta_1) : (U_1, y_1) \to (U', y')$。
那么
\begin{align*}
& \Mor_{u^p\mathcal{S}}((U_1, y_1), (U, y)) \\
& =
\Mor_\mathcal{C}(U_1, U)
\times_{\Mor_\mathcal{D}(u(U_1), u(U))}
\Mor_\mathcal{S}(y_1, y) \\
& =
\Mor_\mathcal{C}(U_1, U)
\times_{\Mor_\mathcal{D}(u(U_1), u(U))}
\Mor_\mathcal{S}(y_1, y')
\times_{\Mor_\mathcal{D}(u(U_1), u(U'))}
\Mor_\mathcal{D}(u(U_1), u(U)) \\
& =
\Mor_\mathcal{S}(y_1, y')
\times_{\Mor_\mathcal{D}(u(U_1), u(U'))}
\Mor_\mathcal{C}(U_1, U) \\
& =
\Mor_{u^p\mathcal{S}}((U_1, y_1), (U', y'))
\times_{\Mor_\mathcal{C}(U_1, U')}
\Mor_\mathcal{C}(U_1, U)
\end{align*}
其中第二个等号来自 $\beta$ 强笛卡尔。因此，确实可知
$(a, \beta)$ 强笛卡尔。反之，假设 $(a, \beta)$ 强笛卡尔。
在 $\mathcal{S}$ 中选取满足
$p(\beta') = u(a)$ 的强笛卡尔态射
$\beta' : y'' \to y'$。
% P01-E0210: source has the typo “bot” for “both”.
% P01-E0211: source's two displayed source/target pairs are internally
% inconsistent with the original morphism and the common lift of a.
那么 $(a, \beta) : (U, y) \to (U, y')$ 与
$(a, \beta') : (U, y'') \to (U, y)$ 均为强笛卡尔态射，
并提升 $a$。因此，由强笛卡尔态射的唯一性（见
Categories, Section \ref{categories-section-fibred-categories}
中的讨论），存在 $\mathcal{S}_{u(U)}$ 中的同构
$\iota : y \to y''$，满足
$\beta = \beta' \circ \iota$；由 Categories, Lemma
\ref{categories-lemma-composition-cartesian}，这蕴含
$\beta$ 在 $\mathcal{S}$ 中强笛卡尔。

\medskip\noindent
最后还须证明：给定 $(U', y')$ 与 $U \to U'$，
可以在 $u^p\mathcal{S}$ 中找到提升态射
$U \to U'$ 的强笛卡尔态射
$(U, y) \to (U', y')$。由上述结论及假设，这可由选取提升
$u(U) \to u(U')$ 的强笛卡尔态射 $y \to y'$ 得到。
\end{proof}

\begin{lemma}
\label{stacks-lemma-stack-pushforward}
令 $u : \mathcal{C} \to \mathcal{D}$ 为位点之间的连续函子。
令 $p : \mathcal{S} \to \mathcal{D}$ 为
$\mathcal{D}$ 上的栈。那么 $u^p\mathcal{S}$
是 $\mathcal{C}$ 上的栈。
\end{lemma}

\begin{proof}
引理 \ref{stacks-lemma-fibred-category-pushforward} 已说明
$u^p\mathcal{S}$ 是 $\mathcal{C}$ 上的纤维化范畴。
此外，该引理的证明说明：
$u^p\mathcal{S}$ 的态射 $(a, \beta)$ 强笛卡尔，
当且仅当 $\beta$ 在 $\mathcal{S}$ 中强笛卡尔。
因此，给定 $\mathcal{C}$ 的态射 $a : U \to U'$，
我们不仅有等式
% P01-E0212: source omits u on the fibre indices; compare source line 2695.
$(u^p\mathcal{S})_U = \mathcal{S}_U$ 与
$(u^p\mathcal{S})_{U'} = \mathcal{S}_{U'}$，
而且在这些等式之下拉回函子一致；写成公式即
$a^*(U', y') = (U, u(a)^*y')$。

\medskip\noindent
说明这些后，令
$\mathcal{U} = \{U_i \to U\}$ 为 $\mathcal{C}$ 的覆盖。
由于 $u$ 连续，
$\mathcal{V} = \{u(U_i) \to u(U)\}$ 是
$\mathcal{D}$ 的覆盖，并且
$u(U_i \times_U U_j) = u(U_i) \times_{u(U)} u(U_j)$；
三重纤维积
$U_i \times_U U_j \times_U U_k$ 也同样。
由于我们有纤维范畴与拉回的这些识别，可知
% P01-E0213: source twice says “descend data” instead of “descent data”.
相对于 $\mathcal{U}$ 的下降数据与相对于
$\mathcal{V}$ 的下降数据完全相同。
根据假设，$\mathcal{S}$ 中下降有效，故
$u^p\mathcal{S}$ 中也如此。
\end{proof}

\begin{lemma}
\label{stacks-lemma-stack-in-groupoids-pushforward}
令 $u : \mathcal{C} \to \mathcal{D}$ 为位点之间的连续函子。
令 $p : \mathcal{S} \to \mathcal{D}$ 为
$\mathcal{D}$ 上群胚中的栈。那么
$u^p\mathcal{S}$ 是 $\mathcal{C}$ 上群胚中的栈。
\end{lemma}

\begin{proof}
这立即由引理 \ref{stacks-lemma-stack-pushforward}
以及所有纤维范畴均为群胚这一事实得出。
\end{proof}

\begin{definition}
\label{stacks-definition-pushforward-stack}
令 $f : \mathcal{D} \to \mathcal{C}$ 为由连续函子
$u : \mathcal{C} \to \mathcal{D}$ 给出的位点态射。
令 $\mathcal{S}$ 为 $\mathcal{D}$ 上的纤维化范畴。
在这种情况下，以{\it $f_*\mathcal{S}$} 表示上面定义的
纤维化范畴 $u^p\mathcal{S}$。称
$f_*\mathcal{S}$ 为 $\mathcal{S}$ {\it 沿 $f$ 的推前}。
\end{definition}

\noindent
由上述结论可知，若 $\mathcal{S}$ 是栈（群胚中的栈），则
$f_*\mathcal{S}$ 是栈（群胚中的栈）。
栈的拉回较难定义（我们的这个特定构造还需要额外假设——
欢迎写出并提交更一般的构造）。我们分几步完成。

\medskip\noindent
令 $u : \mathcal{C} \to \mathcal{D}$ 为范畴之间的函子。
令 $p : \mathcal{S} \to \mathcal{C}$ 为
$\mathcal{C}$ 上的范畴。在这种情况下，定义范畴
$u_{pp}\mathcal{S}$ 如下：
\begin{enumerate}
\item $u_{pp}\mathcal{S}$ 的对象是三元组
$(U, \phi : V \to u(U), x)$，其中
$U \in \Ob(\mathcal{C})$，映射
$\phi : V \to u(U)$ 是 $\mathcal{D}$ 的态射，而
$x \in \Ob(\mathcal{S}_U)$。
\item $u_{pp}\mathcal{S}$ 的态射
$$
(U_1, \phi_1 : V_1 \to u(U_1), x_1)
\longrightarrow
(U_2, \phi_2 : V_2 \to u(U_2), x_2)
$$
由 $(a, b, \alpha)$ 给出，其中
$a : U_1 \to U_2$ 是 $\mathcal{C}$ 的态射，
$b : V_1 \to V_2$ 是 $\mathcal{D}$ 的态射，而
% P01-E0214: source omits the article in “is morphism”.
$\alpha : x_1 \to x_2$ 是 $\mathcal{S}$ 的态射，
满足 $p(\alpha) = a$，并且图式
$$
\xymatrix{
V_1 \ar[d]_{\phi_1} \ar[r]_b & V_2 \ar[d]^{\phi_2} \\
u(U_1) \ar[r]^{u(a)} & u(U_2)
}
$$
在 $\mathcal{D}$ 中交换。
\end{enumerate}
通过
$$
p_{pp} : u_{pp}\mathcal{S} \longrightarrow \mathcal{D},
\quad
(U, \phi : V \to u(U), x) \longmapsto V.
$$
把 $u_{pp}\mathcal{S}$ 看作 $\mathcal{D}$ 上的范畴。
$u_{pp}\mathcal{S}$ 在 $\mathcal{D}$ 的对象 $V$
上的纤维范畴没有简单描述。

\begin{lemma}
\label{stacks-lemma-right-multiplicative-system}
在上述情形中，假设：
\begin{enumerate}
\item $p : \mathcal{S} \to \mathcal{C}$ 是纤维化范畴；
\item $\mathcal{C}$ 有非空有限极限；且
\item $u : \mathcal{C} \to \mathcal{D}$ 与非空有限极限交换。
\end{enumerate}
考虑由形如
$$
(a, \text{id}_V, \alpha) :
(U', \phi' : V \to u(U'), x')
\longrightarrow
(U, \phi : V \to u(U), x)
$$
且 $\alpha$ 强笛卡尔的态射组成的集合
$R \subset \text{Arrows}(u_{pp}\mathcal{S})$。
那么 $R$ 是右乘法系。
\end{lemma}

\begin{proof}
按照 Categories, Definition
\ref{categories-definition-multiplicative-system}，
需检验 RMS1、RMS2、RMS3。
条件 RMS1 成立，因为强笛卡尔态射的复合仍为强笛卡尔态射，
见 Categories, Lemma
\ref{categories-lemma-composition-cartesian}。

\medskip\noindent
为检验 RMS2，假设有 $u_{pp}\mathcal{S}$ 的态射
$$
(a, b, \alpha) :
(U_1, \phi_1 : V_1 \to u(U_1), x_1)
\longrightarrow
(U, \phi : V \to u(U), x)
$$
以及 $R$ 中的态射
$$
(c, \text{id}_V, \gamma) :
(U', \phi' : V \to u(U'), x')
\longrightarrow
(U, \phi : V \to u(U), x)
$$
其中 $\gamma$ 强笛卡尔。在这种情况下，令
$U'_1 = U_1 \times_U U'$，并以
$a' : U'_1 \to U'$ 与
$c' : U'_1 \to U_1$ 表示投影。
因为 $u(U'_1) = u(U_1) \times_{u(U)} u(U')$，
$\phi'_1 = (\phi_1, \phi') : V_1 \to u(U'_1)$
是 $\mathcal{D}$ 中的态射。
令 $\gamma_1 : x_1' \to x_1$ 为 $\mathcal{S}$
中满足
% P01-E0215: p(gamma_1) is C-valued but source sets it equal to
% the D-morphism phi'_1; the typed projection is c'.
$p(\gamma_1) = \phi'_1$ 的强笛卡尔态射
（它存在，因为 $\mathcal{S}$ 是 $\mathcal{C}$ 上的纤维化范畴）。
由于 $\gamma : x' \to x$ 强笛卡尔，存在唯一态射
$\alpha' : x'_1 \to x'$，满足 $p(\alpha') = a'$。
此时可知
% P01-E0216: the source display has inconsistent U/U' and phi/phi' entries.
$$
(a', b, \alpha') :
(U_1, \phi_1 : V_1 \to u(U'_1), x'_1)
\longrightarrow
(U, \phi : V \to u(U'), x')
$$
是态射，并且
% P01-E0217: source writes phi instead of phi_1 in the codomain tuple.
$$
(c', \text{id}_{V_1}, \gamma_1) :
(U'_1, \phi'_1 : V_1 \to u(U'_1), x'_1)
\longrightarrow
(U_1, \phi : V_1 \to u(U_1), x_1)
$$
是 $R$ 的元素；它
% P01-E0218: source has singular “which form”.
构成 RMS2 所提出存在性问题的一个解。

\medskip\noindent
最后，假设
$$
(a, b, \alpha), (a', b', \alpha') :
(U_1, \phi_1 : V_1 \to u(U_1), x_1)
\longrightarrow
(U, \phi : V \to u(U), x)
$$
是 $u_{pp}\mathcal{S}$ 的两个态射，并假设
$$
(c, \text{id}_V, \gamma) :
(U, \phi : V \to u(U), x)
\longrightarrow
(U', \phi : V \to u(U'), x')
$$
是 $R$ 中使态射 $(a, b, \alpha)$ 与
$(a', b', \alpha')$ 相等的元素。
特别地，这蕴含 $b = b'$。令 $d : U_2 \to U_1$
为 $a, a'$ 的等化子；它存在（见
Categories, Lemma
\ref{categories-lemma-almost-finite-limits-exist}）。
此外，$u(d) : u(U_2) \to u(U_1)$ 是
$u(a), u(a')$ 的等化子；故（因为 $b = b'$）
存在态射 $\phi_2 : V_1 \to u(U_2)$，使得
% P01-E0219: source repeats phi_1 on both sides; the factorization
% should use phi_2 on the right.
$\phi_1 = u(d) \circ \phi_1$。
令 $\delta : x_2 \to x_1$ 为 $\mathcal{S}$ 中满足
% P01-E0220: p(delta) is C-valued, while u(d) lies in D; d is typed.
$p(\delta) = u(d)$ 的强笛卡尔态射。
现在断言 $\alpha \circ \delta = \alpha' \circ \delta$。
这是因为 $\gamma$ 强笛卡尔、
$\gamma \circ \alpha \circ \delta
= \gamma \circ \alpha' \circ \delta$，且
$p(\alpha \circ \delta) = p(\alpha' \circ \delta)$。
因此箭头
$$
(d, \text{id}_{V_1}, \delta) :
(U_2, \phi_2 : V_1 \to u(U_2), x_2)
\longrightarrow
(U_1, \phi_1 : V_1 \to u(U_1), x_1)
$$
是 $R$ 的元素，并使
$(a, b, \alpha)$ 与 $(a', b', \alpha')$ 相等。
所以 $R$ 也满足 RMS3。
\end{proof}

\begin{lemma}
\label{stacks-lemma-fibred-category-pullback}
记号与假设同引理
\ref{stacks-lemma-right-multiplicative-system}。
令 $u_p\mathcal{S} = R^{-1}u_{pp}\mathcal{S}$，见
Categories, Section \ref{categories-section-localization}。
那么 $u_p\mathcal{S}$ 是 $\mathcal{D}$ 上的纤维化范畴。
\end{lemma}

\begin{proof}
使用 Categories, Lemma
\ref{categories-lemma-right-localization} 正上方给出的
$u_p\mathcal{S}$ 的描述。注意，函子
$p_{pp} : u_{pp}\mathcal{S} \to \mathcal{D}$
把 $R$ 的每个元素都映为恒等态射。
因此，由 Categories, Lemma
\ref{categories-lemma-properties-right-localization}，
得到延拓给定函子的典范函子
$p_p : u_p\mathcal{S} \to \mathcal{D}$。
我们就是这样把 $u_p\mathcal{S}$ 看作
$\mathcal{D}$ 上的范畴。

\medskip\noindent
首先刻画 $u_p\mathcal{S}$ 中的
$\mathcal{D}$-强笛卡尔态射。
$u_p\mathcal{S}$ 的态射 $f : X \to Y$ 是偶对
$(f' : X' \to Y, r : X' \to X)$ 的等价类，其中
$r \in R$。事实上，在 $u_p\mathcal{S}$ 中，采用显然的记号，
有 $f = (f', 1) \circ (r, 1)^{-1}$。
注意，同构总是强笛卡尔的，强笛卡尔态射的复合也如此，见
Categories, Lemma \ref{categories-lemma-composition-cartesian}。
故 $f$ 强笛卡尔，当且仅当 $(f', 1)$ 强笛卡尔。
因此，下面的断言完全刻画强笛卡尔态射。断言：
$u_{pp}\mathcal{S}$ 的态射
$$
(a, b, \alpha) :
X_1 = (U_1, \phi_1 : V_1 \to u(U_1), x_1)
\longrightarrow
(U_2, \phi_2 : V_2 \to u(U_2), x_2) = X_2
$$
的像 $f = ((a, b, \alpha), 1)$ 在
$u_p\mathcal{S}$ 中强笛卡尔，当且仅当
$\alpha$ 是 $\mathcal{S}$ 的强笛卡尔态射。

\medskip\noindent
假设 $\alpha$ 强笛卡尔。
令 $X = (U, \phi : V \to u(U), x)$ 为另一个对象，并令
$f_2 : X \to X_2$ 为 $u_p\mathcal{S}$ 的态射，满足
% P01-E0221: b_1 is D-valued here; source types it as U→U_1 in C.
$p_p(f_2) = b \circ b_1$，其中
$b_1 : U \to U_1$。
要证明 $f$ 强笛卡尔，需证明在 $u_p\mathcal{S}$ 中存在唯一态射
$f_1 : X \to X_1$，使得
$p_p(f_1) = b_1$，并且在 $u_p\mathcal{S}$ 中有
$f_2 = f \circ f_1$。
写 $f_2 = (f'_2 : X' \to X_2, r : X' \to X)$。
仍可在 $u_p\mathcal{S}$ 中写成
$f_2 = (f'_2, 1) \circ (r, 1)^{-1}$。
由于 $(r, 1)$ 是在 $\mathcal{D}$ 中的像为恒等态射的同构，
寻找具有所需性质的态射 $f_1 : X \to X_1$，
等同于寻找 $u_p\mathcal{S}$ 中的态射
$f'_1 : X' \to X_1$，满足
% P01-E0222: source uses p instead of the localization's base functor p_p.
$p(f'_1) = b_1$ 与 $f'_2 = f \circ f'_1$。
因此，可假设 $f_2$ 形如
$f_2 = ((a_2, b_2, \alpha_2), 1)$，其中
$b_2 = b \circ b_1$。图示如下：
$$
\xymatrix{
& & (U_1, V_1 \to u(U_1), x_1) \ar[d]^{(a, b, \alpha)} \\
(U, V \to u(U), x) \ar[rr]^{(a_2, b_2, \alpha_2)} & &
(U_2, V_2 \to u(U_2), x_2)
}
$$
现在容易构造态射 $f_1$。
具体地，令 $U' = U \times_{U_2} U_1$，投影为
$c : U' \to U$ 与 $a_1 : U' \to U_1$。
选取提升态射 $c$ 的强笛卡尔态射
$\gamma : x' \to x$。由于 $b_2 = b \circ b_1$，并且
$u(U') = u(U) \times_{u(U_2)} u(U_1)$，可知
$\phi' = (\phi, \phi_1 \circ b_1) : V \to u(U')$。
由于 $\alpha$ 强笛卡尔，且
$a \circ a_1 = a_2 \circ c = p(\alpha_2 \circ \gamma)$，
存在提升 $a_1$ 的态射
$\alpha_1 : x' \to x_1$，满足
$\alpha \circ \alpha_1 = \alpha_2 \circ \gamma$。
令 $X' = (U', \phi' : V \to u(U'), x')$。
于是可知
$$
f_1 =
((a_1, b_1, \alpha_1) : X' \to X_1, (c, \text{id}_V, \gamma) : X' \to X) :
X \longrightarrow X_1
$$
可行；事实上，图式
$$
\xymatrix{
(U', \phi' : V \to u(U'), x')
\ar[d]_{(c, \text{id}_V, \gamma)}
\ar[rr]_{(a_1, b_1, \alpha_1)}
& & (U_1, V_1 \to u(U_1), x_1) \ar[d]^{(a, b, \alpha)} \\
(U, V \to u(U), x) \ar[rr]^{(a_2, b_2, \alpha_2)} & &
(U_2, V_2 \to u(U_2), x_2)
}
$$
由构造交换。这证明了存在性。

\medskip\noindent
接下来证明唯一性；仍处于
$f = ((a, b, \alpha), 1)$ 与
$f_2 = ((a_2, b_2, \alpha_2), 1)$ 的特殊情形。
我们强烈建议读者跳过这一部分。
假设 $g_1, g'_1 : X \to X_1$ 是
$u_p\mathcal{S}$ 的两个态射，满足
$p_p(g_1) = p_p(g'_1) = b_1$ 与
$f_2 = f \circ g_1 = f \circ g'_1$。
目标是证明 $g_1 = g'_1$。由 Categories, Lemma
\ref{categories-lemma-morphisms-right-localization}，
可把 $g_1$ 与 $g'_1$ 分别表示为
$(f_1 : X' \to X_1, r : X' \to X)$ 与
$(f'_1 : X' \to X_1, r : X' \to X)$ 的等价类，
其中 $r \in R$。由 Categories, Lemma
\ref{categories-lemma-equality-morphisms-right-localization}，
$f_2 = f \circ g_1 = f \circ g'_1$ 意味着存在
$u_{pp}\mathcal{S}$ 中的态射 $r' : X'' \to X'$，
使 $r' \circ r \in R$，并且
$$
(a, b, \alpha) \circ f_1 \circ r' =
(a, b, \alpha) \circ f'_1 \circ r' =
(a_2, b_2, \alpha_2) \circ r'
$$
在 $u_{pp}\mathcal{S}$ 中成立。
注意，此时 $g_1$ 由 $(f_1 \circ r', r \circ r')$ 表示，
$g'_1$ 也类似。因此，可假设
$$
(a, b, \alpha) \circ f_1 =
(a, b, \alpha) \circ f'_1 =
(a_2, b_2, \alpha_2).
$$
写
$r = (c, \text{id}_V, \gamma) :
(U', \phi' : V \to u(U'), x')$、
$f_1 = (a_1, b_1, \alpha_1)$ 与
$f'_1 = (a'_1, b_1, \alpha'_1)$。
这里使用了条件 $p_p(g_1) = p_p(g'_1)$。
上述等式现在等价于
$a \circ a_1 = a \circ a'_1 = a_2 \circ c$ 与
$\alpha \circ \alpha_1
= \alpha \circ \alpha'_1 = \alpha_2 \circ \gamma$。
在这种情形，未必有 $a_1 = a'_1$。
所以还须再与一个来自 $R$ 的态射预复合。
具体地，令 $U'' = \text{Eq}(a_1, a'_1)$ 为
$a_1$ 与 $a'_1$ 的等化子；它是 $U'$ 的子对象。
以 $c' : U'' \to U'$ 表示典范单态射。
由上述态射之间的关系可知，$V \to u(U')$ 分解经过
$u(U'') = u(\text{Eq}(a_1, a'_1))
= \text{Eq}(u(a_1), u(a'_1))$。
因而得到新对象
$(U'', \phi'' : V \to u(U''), x'')$，其中
% P01-E0223: source says gamma' lifts gamma; gamma is not a base arrow.
$\gamma' : x'' \to x'$ 是提升 $\gamma$ 的强笛卡尔态射。
于是，可用 $R$ 的元素
$(c', \text{id}_V, \gamma')$ 与 $f_1$、$f'_1$ 预复合。
这样做后，即用
$(U'', \phi'' : V \to u(U''), x'')$ 替换
$(U', \phi' : V \to u(U'), x')$ 后，回到先前的情形，
并且现在还有 $a_1 = a'_1$。
此时，由 $\alpha$ 强笛卡尔这一事实（！）形式地推出
$\alpha_1 = \alpha'_1$。因此，正如所需，
$g_1 = g'_1$。

\medskip\noindent
我们略去证明如下事实：对
$u_p\mathcal{S}$ 中形如 $((a, b, \alpha), 1)$
的任意强笛卡尔态射，态射 $\alpha$ 在
$\mathcal{S}$ 中强笛卡尔。
（在证明的其余部分不需要强笛卡尔态射的这一刻画，
尽管本节后面会使用它。）

\medskip\noindent
令 $(U, \phi : V \to u(U), x)$ 为
$u_p\mathcal{S}$ 的对象。令 $b : V' \to V$
为 $\mathcal{D}$ 的态射。那么态射
$$
(\text{id}_U, b, \text{id}_x) :
(U, \phi \circ b : V' \to u(U), x)
\longrightarrow
(U, \phi : V \to u(U), x)
$$
由前面各段的结论强笛卡尔，证明完成。
\end{proof}

\begin{lemma}
\label{stacks-lemma-fibred-groupoids-category-pullback}
记号与假设同引理
\ref{stacks-lemma-fibred-category-pullback}。
若 $\mathcal{S}$ 纤维化于群胚，则
$u_p\mathcal{S}$ 纤维化于群胚。
\end{lemma}

\begin{proof}
由引理 \ref{stacks-lemma-fibred-category-pullback} 可知，
$u_p\mathcal{S}$ 是纤维化范畴。
令 $f : X \to Y$ 为 $u_p\mathcal{S}$ 的态射，并满足
$p_p(f) = \text{id}_V$。若能证明 $f$ 可逆，便告完成；见
Categories, Lemma \ref{categories-lemma-fibred-groupoids}。
把 $f$ 写成偶对 $((a, b, \alpha), r)$ 的等价类，
其中 $r \in R$。那么 $p_p(r) = \text{id}_V$，
故 $p_{pp}((a, b, \alpha)) = \text{id}_V$，从而
$b = \text{id}_V$。但 $\mathcal{S}$ 的任意态射均为强笛卡尔态射，
见 Categories, Lemma
\ref{categories-lemma-fibred-groupoids}；因此
$(a, b, \alpha) \in R$ 在 $u_p\mathcal{S}$ 中可逆，正如所需。
\end{proof}

\begin{lemma}
\label{stacks-lemma-adjointness-pullback-pushforward}
令 $u : \mathcal{C} \to \mathcal{D}$ 为函子。
令 $p : \mathcal{S} \to \mathcal{C}$ 与
$q : \mathcal{T} \to \mathcal{D}$ 分别为
$\mathcal{C}$ 与 $\mathcal{D}$ 上的范畴。假设：
\begin{enumerate}
\item $p : \mathcal{S} \to \mathcal{C}$ 是纤维化范畴；
\item $q : \mathcal{T} \to \mathcal{D}$ 是纤维化范畴；
\item $\mathcal{C}$ 有非空有限极限；且
\item $u : \mathcal{C} \to \mathcal{D}$ 与非空有限极限交换。
\end{enumerate}
那么，态射范畴之间有典范等价
$$
\Mor_{\textit{Fib}/\mathcal{C}}(\mathcal{S}, u^p\mathcal{T})
=
\Mor_{\textit{Fib}/\mathcal{D}}(u_p\mathcal{S}, \mathcal{T})
$$
\end{lemma}

\begin{proof}
在此证明中，以记号 $x/U$ 表示
$\mathcal{S}$ 中位于 $\mathcal{C}$ 的 $U$ 上的对象 $x$。
类似地，$y/V$ 表示 $\mathcal{T}$ 中位于
$\mathcal{D}$ 的 $V$ 上的对象 $y$。
同理，$\alpha/a : x/U \to x'/U'$ 表示态射
$\alpha : x \to x'$，它在 $\mathcal{C}$ 中的像为
$a : U \to U'$。

\medskip\noindent
令 $G : u_p\mathcal{S} \to \mathcal{T}$ 为
$\mathcal{D}$ 上纤维化范畴之间的 $1$-态射。
% P01-E0224: source omits “by” in “Denote G' ... the composition”.
以 $G' : u_{pp}\mathcal{S} \to \mathcal{T}$ 表示
$G$ 与典范（局部化）函子
$u_{pp}\mathcal{S} \to u_p\mathcal{S}$ 的复合。
考虑由如下规则给出的函子
$H : \mathcal{S} \to u^p\mathcal{T}$：
在对象上，
$$
H(x/U) = (U, G'(U, \text{id}_{u(U)} : u(U) \to u(U), x))
$$
而在态射上，
% P01-E0225: source's displayed morphism uses the tuple order (alpha,a),
% whereas objects of u^pT were defined with morphisms (a,beta).
$$
H((\alpha, a) : x/U \to x'/U') = G'(a, u(a), \alpha)
$$
由于 $G$ 把强笛卡尔态射变为强笛卡尔态射，可知若
$\alpha$ 强笛卡尔，则 $H(\alpha)$ 强笛卡尔。
事实上，引理 \ref{stacks-lemma-fibred-category-pullback} 的证明已说明，
在这种情况下，映射 $(a, u(a), \alpha)$ 在
$u_p\mathcal{S}$ 中成为强笛卡尔态射。
显然，该构造关于 $G$ 具有函子性，于是得到函子
$$
A :
\Mor_{\textit{Fib}/\mathcal{D}}(u_p\mathcal{S}, \mathcal{T})
\longrightarrow
\Mor_{\textit{Fib}/\mathcal{C}}(\mathcal{S}, u^p\mathcal{T})
$$

\medskip\noindent
反之，令 $H : \mathcal{S} \to u^p\mathcal{T}$
为纤维化范畴之间的 $1$-态射。
回忆，$u^p\mathcal{T}$ 的对象是偶对 $(U, y)$，其中
$y \in \Ob(\mathcal{T}_{u(U)})$。
% P01-E0226: source omits “by” in “We denote pr ... the functor”.
以 $\text{pr} : u^p\mathcal{T} \to \mathcal{T}$
表示函子 $(U, y) \mapsto y$。
在这种情况下，按如下规则定义函子
$G' : u_{pp}\mathcal{S} \to \mathcal{T}$。
在对象上，
$$
G'(U, \phi : V \to u(U), x) = \phi^*\text{pr}(H(x))
$$
而令
$$
G'((a, b, \alpha) : (U, \phi : V \to u(U), x)
\to (U', \phi' : V' \to u(U'), x'))
=
\beta
$$
为满足 $q(\beta) = b$ 且使图式
$$
\xymatrix{
\phi^*\text{pr}(H(x)) \ar[d] \ar[r]_-{\beta} &
(\phi')^*\text{pr}(H(x')) \ar[d] \\
\text{pr}(H(x)) \ar[r]^{\text{pr}(H(a, \alpha))} & \text{pr}(H(x'))
}
$$
交换的唯一态射
$\beta : \phi^*\text{pr}(H(x))
\to (\phi')^*\text{pr}(H(x'))$。
这样的态射存在且唯一，因为 $\mathcal{T}$ 是纤维化范畴。

\medskip\noindent
检验若 $r \in R$，则 $G'(r)$ 是同构。
具体地，若
$$
(a, \text{id}_V, \alpha) :
(U', \phi' : V \to u(U'), x')
\longrightarrow
(U, \phi : V \to u(U), x)
$$
是引理 \ref{stacks-lemma-right-multiplicative-system} 的右乘法系
$R$ 中的元素，且 $\alpha$ 强笛卡尔，则
$H(\alpha)$ 强笛卡尔，而
$\text{pr}(H(\alpha))$ 强笛卡尔；见引理
\ref{stacks-lemma-fibred-category-pushforward} 的证明。
因此，此时态射 $\beta$ 满足
$q(\beta) = \text{id}_V$ 且强笛卡尔。
所以，由 Categories, Lemma
\ref{categories-lemma-composition-cartesian}，
$\beta$ 是同构。于是，由 Categories, Lemma
\ref{categories-lemma-properties-right-localization}，
得到典范延拓 $G : u_p\mathcal{S} \to \mathcal{T}$。

\medskip\noindent
接下来证明 $G$ 把强笛卡尔态射变为强笛卡尔态射。
% P01-E0227: source has the extraneous article in “is a strongly cartesian”.
假设 $f : X \to Y$ 是强笛卡尔态射。
由 $u_p\mathcal{S}$ 中强笛卡尔态射的刻画，可把
$f$ 写成
% P01-E0228: source makes the roof arrow r land in Y; it should land in X.
$((a, b, \alpha) : X' \to Y, r : X' \to Y)$，
其中 $r \in R$，且 $\alpha$ 在 $\mathcal{S}$ 中强笛卡尔。
由上述结论，只需证明
% P01-E0229: source omits the comma between b and alpha.
$G'(a, b \alpha)$ 强笛卡尔。
和之前一样，$\alpha$ 强笛卡尔这一条件蕴含
$\text{pr}(H(a, \alpha)) :
\text{pr}(H(x)) \to \text{pr}(H(x'))$
在 $\mathcal{T}$ 中强笛卡尔。
% P01-E0230: source has number disagreement in “all arrows ... is”.
由于上面的交换方块中除 $\beta$ 外的所有箭头现在都强笛卡尔，
$\beta$ 也强笛卡尔，正如所需。
显然，构造 $H \mapsto G$ 关于 $H$ 具有函子性，
于是得到函子
$$
B :
\Mor_{\textit{Fib}/\mathcal{C}}(\mathcal{S}, u^p\mathcal{T})
\longrightarrow
\Mor_{\textit{Fib}/\mathcal{D}}(u_p\mathcal{S}, \mathcal{T})
$$
为完成引理证明，还须证明函子 $A$ 与 $B$ 互为拟逆。
验证从略。
\end{proof}

\begin{definition}
\label{stacks-definition-pullback-stack}
令 $f : \mathcal{D} \to \mathcal{C}$ 为由连续函子
$u : \mathcal{C} \to \mathcal{D}$ 给出的位点态射，
并假设它满足 Sites, Proposition
\ref{sites-proposition-get-morphism} 的假设与结论。
令 $\mathcal{S}$ 为 $\mathcal{C}$ 上的栈。
在这种情况下，以{\it $f^{-1}\mathcal{S}$} 表示
上面构造的 $\mathcal{D}$ 上纤维化范畴
$u_p\mathcal{S}$ 的栈化。称
$f^{-1}\mathcal{S}$ 为 $\mathcal{S}$ {\it 沿 $f$ 的拉回}。
\end{definition}

\noindent
当然，若 $\mathcal{S}$ 是群胚中的栈，则由引理
\ref{stacks-lemma-stackify-groupoids} 与
\ref{stacks-lemma-fibred-groupoids-category-pullback}，
$f^{-1}\mathcal{S}$ 是群胚中的栈。

\begin{lemma}
\label{stacks-lemma-adjointness-pullback-pushforward-stacks}
令 $f : \mathcal{D} \to \mathcal{C}$ 为由连续函子
$u : \mathcal{C} \to \mathcal{D}$ 给出的位点态射，
并假设它满足 Sites, Proposition
\ref{sites-proposition-get-morphism} 的假设与结论。
令 $p : \mathcal{S} \to \mathcal{C}$ 与
$q : \mathcal{T} \to \mathcal{D}$ 为栈。
那么，态射范畴之间有典范等价
$$
\Mor_{\textit{Stacks}/\mathcal{C}}(\mathcal{S}, f_*\mathcal{T})
=
\Mor_{\textit{Stacks}/\mathcal{D}}(f^{-1}\mathcal{S}, \mathcal{T})
$$
\end{lemma}

\begin{proof}
对 $i = 1, 2$，栈的 $i$-态射与纤维化范畴的
$i$-态射相同；见 Definition \ref{stacks-definition-stacks-over-C}。
由引理 \ref{stacks-lemma-adjointness-pullback-pushforward} 已有
$$
\Mor_{\textit{Fib}/\mathcal{C}}(\mathcal{S}, u^p\mathcal{T})
=
\Mor_{\textit{Fib}/\mathcal{D}}(u_p\mathcal{S}, \mathcal{T})
$$
又因 $u^p\mathcal{T} = f_*\mathcal{T}$，而
$f^{-1}\mathcal{S}$ 是 $u_p\mathcal{S}$ 的栈化，
结论由引理
\ref{stacks-lemma-stackify-universal-property-more} 得出。
\end{proof}

\begin{lemma}
\label{stacks-lemma-technical-up}
令 $f : \mathcal{D} \to \mathcal{C}$ 为由连续函子
$u : \mathcal{C} \to \mathcal{D}$ 给出的位点态射，
并假设它满足 Sites, Proposition
\ref{sites-proposition-get-morphism} 的假设与结论。
令 $\mathcal{S} \to \mathcal{C}$ 为纤维化范畴，并令
$\mathcal{S} \to \mathcal{S}'$ 为 $\mathcal{S}$ 的栈化。
那么 $f^{-1}\mathcal{S}'$ 是
$u_p\mathcal{S}$ 的栈化。
\end{lemma}

\begin{proof}
略。提示：这是 Sites, Lemma
\ref{sites-lemma-technical-up} 的类似结论。
\end{proof}

\noindent
下面的引理告诉我们：只要
$\Sch'_{fppf}$ 包含 $\Sch_{fppf}$（见
Topologies, Section \ref{topologies-section-change-alpha}），
$\Sch_{fppf}$ 上栈的 $2$-范畴就是
$\Sch'_{fppf}$ 上栈的 $2$-范畴的一个“满 2-子范畴”。

\begin{lemma}
\label{stacks-lemma-bigger-site}
令 $\mathcal{C}$ 与 $\mathcal{D}$ 为位点。
令 $u : \mathcal{C} \to \mathcal{D}$ 为满足
Sites, Lemma \ref{sites-lemma-bigger-site} 的假设的函子。
令 $f : \mathcal{D} \to \mathcal{C}$ 为相应的位点态射。
那么：
\begin{enumerate}
\item 对每个栈
$p : \mathcal{S} \to \mathcal{C}$，典范函子
$\mathcal{S} \to f_*f^{-1}\mathcal{S}$ 是栈之间的等价；
\item 给定 $\mathcal{C}$ 上的栈
$\mathcal{S}, \mathcal{S}'$，构造 $f^{-1}$ 诱导态射范畴之间的等价
$$
\Mor_{\textit{Stacks}/\mathcal{C}}(\mathcal{S}, \mathcal{S}')
\longrightarrow
\Mor_{\textit{Stacks}/\mathcal{D}}(f^{-1}\mathcal{S}, f^{-1}\mathcal{S}')
$$
\end{enumerate}
\end{lemma}

\begin{proof}
注意，由引理
\ref{stacks-lemma-adjointness-pullback-pushforward-stacks}，
有范畴之间的等价
$$
\Mor_{\textit{Stacks}/\mathcal{D}}(f^{-1}\mathcal{S}, f^{-1}\mathcal{S}')
=
\Mor_{\textit{Stacks}/\mathcal{C}}(\mathcal{S}, f_*f^{-1}\mathcal{S}')
$$
所以 (2) 由 (1) 得出。

\medskip\noindent
为证明 (1)，将使用引理
\ref{stacks-lemma-characterize-essentially-surjective-when-ff}。
该引理说明，必须证明
$can : \mathcal{S} \to f_*f^{-1}\mathcal{S}$
全忠实，并且 $f_*f^{-1}\mathcal{S}$ 的所有对象
局部地属于本质像。

\medskip\noindent
我们快速描述函子 $can$；见引理
\ref{stacks-lemma-adjointness-pullback-pushforward} 的证明。
为此，引入函子
$c'' : \mathcal{S} \to u_{pp}\mathcal{S}$，定义为
$c''(x/U) = (U, \text{id} : u(U) \to u(U), x)$ 与
$c''(\alpha/a) = (a, u(a), \alpha)$。
令 $c' : \mathcal{S} \to u_p\mathcal{S}$ 为
$c''$ 与典范函子
$u_{pp}\mathcal{S} \to u_p\mathcal{S}$ 的复合。
令 $c : \mathcal{S} \to f^{-1}\mathcal{S}$ 为
$c'$ 与典范函子
$u_p\mathcal{S} \to f^{-1}\mathcal{S}$ 的复合。
那么 $can : \mathcal{S} \to f_*f^{-1}\mathcal{S}$
是这样的函子：它把 $x/U$ 送到偶对
$(U, c(x))$，把 $\alpha/a$ 送到态射
$(a, c(\alpha))$。

\medskip\noindent
% P01-E0231: source uses “Fully faithfulness” instead of “Full faithfulness”.
全忠实性。为证明这一点，将使用引理
\ref{stacks-lemma-characterize-ff}。
令 $U \in \Ob(\mathcal{C})$。
令 $x, y \in \mathcal{S}_U$。
首先，由于 $u$ 全忠实，直接从 $f_*$ 的定义得到
$$
\Mor_{(f_*f^{-1}\mathcal{S})_U}(can(x), can(y))
=
\Mor_{(f^{-1}\mathcal{S})_{u(U)}}(c(x), c(y))
$$
% P01-E0232: source has “Similar holds” rather than “Similarly”.
拉回到任意 $U'/U$ 后，类似结论也成立。
由于 $f^{-1}\mathcal{S}$ 是 $u_p\mathcal{S}$ 的栈化，
且 $u$ 连续并余连续，预层
$$
U'/U \longmapsto
\Mor_{(f^{-1}\mathcal{S})_{u(U')}}(c(x|_{U'}), c(y|_{U'}))
$$
是预层
$$
U'/U \longmapsto
\Mor_{(u_p\mathcal{S})_{u(U')}}(c'(x|_{U'}), c'(y|_{U'}))
$$
的层化。因此，为完成全忠实性的证明，只需证明对任意
$U$ 与 $x, y$，映射
% P01-E0233: source indexes the D-fibre of u_pS by U instead of u(U).
$$
\Mor_{\mathcal{S}_U}(x, y)
\longrightarrow
\Mor_{(u_p\mathcal{S})_U}(c'(x), c'(y))
$$
是双射。$u_p\mathcal{S}$ 中位于 $u(U)$ 上的态射
$f : x \to y$ 由如下图式的一个等价类给出：
$$
\xymatrix{
(U', \phi : u(U) \to u(U'), x')
\ar[d]_{(c, \text{id}_{u(U)}, \gamma)}
\ar[r]_{(a, b, \alpha)} &
(U, \text{id} : u(U) \to u(U), y)
\\
(U, \text{id} : u(U) \to u(U), x)
}
$$
其中 $\gamma$ 强笛卡尔，且
$b = \text{id}_{u(U)}$。
由于 $u$ 全忠实，可对某个态射
$c' : U \to U'$ 写成 $\phi = u(c')$；
此时可知 $a \circ c' = \text{id}_U$，且
% P01-E0234: c:U'→U, so c∘c' is id_U, not id_{U'}.
$c \circ c' = \text{id}_{U'}$。
由于 $\gamma$ 强笛卡尔，可以找到提升 $c'$ 的态射
$\gamma' : x \to x'$，满足
$\gamma \circ \gamma' = \text{id}_x$。
由 $u_p\mathcal{S}$ 中态射等价类的定义可知，
$u_{pp}\mathcal{S}$ 的态射
$$
\xymatrix{
(U, \text{id} : u(U) \to u(U), x)
\ar[rr]_{(\text{id}, \text{id}, \alpha \circ \gamma')} & &
(U, \text{id} : u(U) \to u(U), y)
}
$$
在 $u_p\mathcal{S}$ 中诱导态射 $f$。
这证明该映射满射。单射性的证明从略。

\medskip\noindent
最后，必须证明 $f_*f^{-1}\mathcal{S}$ 的任意对象
局部地来自 $\mathcal{S}$ 的对象。这由各构造显然成立
（细节从略）。
\end{proof}




\section{栈与局部化}
\label{stacks-section-localize}

\noindent
令 $\mathcal{C}$ 为位点。
令 $U$ 为 $\mathcal{C}$ 的对象。
我们想把 $\mathcal{C}/U$ 上的栈理解为
$\mathcal{C}$ 上的栈连同一个到 $U$ 的态射。
下一个引理说明，当预层 $h_U$ 是层时，为何这更容易做到。

\begin{lemma}
\label{stacks-lemma-when-localization-stack}
令 $\mathcal{C}$ 为位点。令
$U \in \Ob(\mathcal{C})$。
那么 $j_U : \mathcal{C}/U \to \mathcal{C}$
是 $\mathcal{C}$ 上的栈，当且仅当 $h_U$ 是层。
\end{lemma}

\begin{proof}
结合引理 \ref{stacks-lemma-stack-in-setoids-characterize} 与
Categories, Example
\ref{categories-example-fibred-category-from-functor-of-points}。
\end{proof}

\noindent
假设 $\mathcal{C}$ 为位点，$U$ 是 $\mathcal{C}$ 的对象，
且与其相联系的可表预层是层。以
$j : \mathcal{C}/U \to \mathcal{C}$ 表示局部化函子。

\medskip\noindent
{\bf 构造 A。} 令
$p : \mathcal{S} \to \mathcal{C}/U$ 为位点
$\mathcal{C}/U$ 上的栈。如下定义栈
$j_!p : j_!\mathcal{S} \to \mathcal{C}$：
\begin{enumerate}
\item 作为范畴，$j_!\mathcal{S} = \mathcal{S}$；且
\item 函子 $j_!p : j_!\mathcal{S} \to \mathcal{C}$
就是复合 $j \circ p$。
\end{enumerate}
我们略去验证这是栈（提示：使用 $h_U$ 是层来黏合到
$U$ 的态射）。存在典范函子
$$
j_!\mathcal{S} \longrightarrow \mathcal{C}/U
$$
即函子 $p$；它是 $\mathcal{C}$ 上栈之间的 $1$-态射。

\medskip\noindent
{\bf 构造 B。}
令 $q : \mathcal{T} \to \mathcal{C}$ 为
$\mathcal{C}$ 上的栈，并配备 $\mathcal{C}$ 上栈之间的态射
$p : \mathcal{T} \to \mathcal{C}/U$。
在这种情况下，$p : \mathcal{T} \to \mathcal{C}/U$
自动为 $\mathcal{C}/U$ 上的栈。

\begin{lemma}
\label{stacks-lemma-localize-stacks}
假设 $\mathcal{C}$ 为位点，$U$ 是 $\mathcal{C}$ 的对象，
且与其相联系的可表预层是层。上述构造 A、B 定义
$2$-范畴之间互逆（！）的函子
$$
\left\{
\begin{matrix}
2\text{-范畴，其对象为}\\
\text{下列对象上的叠： }\mathcal{C}/U
\end{matrix}
\right\}
\leftrightarrow
\left\{
\begin{matrix}
2\text{-二元组范畴 }(\mathcal{T}, p)
\text{ 由下列对象组成：} \\
\text{叠 }\mathcal{T}\text{ 相对于 }\mathcal{C}\text{ 以及一个态射} \\
p : \mathcal{T} \to \mathcal{C}/U\text{ 叠范畴，位于 }\mathcal{C}
\end{matrix}
\right\}
$$
\end{lemma}

\begin{proof}
这是显然的。
\end{proof}
% Generated mechanically from the frozen fields.tex target.
% Do not edit this generated file; edit the bound locale source instead.

% OK, start here.
%

\setcounter{chapter}{8}
\chapter{域}



\phantomsection
\label{fields-section-phantom}



\section{引言}
\label{fields-section-introduction}

\noindent
在本章中，我们将讨论域论。回忆，一个{\it 域}是非零环，且其中
每个非零元素均可逆。等价地，域仅有两个理想 $(0)$ 与 $(1)$，
因为任意非零元素都是单位。因此，域将是后文所发展理论中许多问题的
最简单情形。

\medskip\noindent
域扩张理论与标准交换代数的观感有所不同，例如，任意域态射都是单射。
尽管如此，涉及环的问题往往可以归约为关于域的问题。例如，任意整环都可
嵌入一个域（它的分式域）；任意{\it 局部环}（即具有唯一极大理想的环；
我们尚未定义这个术语）都有与之相伴的剩余域（即模其极大理想所得的商）。
因此，掌握域扩张将很有用。




\section{基本定义}
\label{fields-section-definitions}

\noindent
由于我们把本章置于讨论交换代数的章节之前，必须先在这里引入一些基本定义，
然后再在代数各章中更详细地讨论它们。

\begin{definition}
\label{fields-definition-field}
一个{\it 域}是每个非零元素都可逆的非零环。给定一个域，
其{\it 子域}是本身亦为域的子环。
\end{definition}

\noindent
对域 $k$，我们以 $k^*$ 表示子集 $k \setminus \{0\}$。
这推广了通常的记号 $R^*$，后者表示环 $R$ 中可逆元素所成的群。

\begin{definition}
\label{fields-definition-domain}
一个{\it 整环}是以 $0$ 为唯一零因子的非零环。
\end{definition}



\section{域的例}
\label{fields-section-examples}

\noindent
首先给出若干域的例。读者应当回忆：若 $R$ 是环且 $I \subset R$ 是理想，
则 $R/I$ 是域，当且仅当 $I$ 是极大理想。

\begin{example}[有理数]
\label{fields-example-rational-numbers}
有理数构成一个域，称为{\it 有理数域}，记作 $\mathbf{Q}$。
\end{example}

\begin{example}[素域]
\label{fields-example-prime-field}
若 $p$ 是素数，则 $\mathbf{Z}/(p)$ 是域，记作
$\mathbf{F}_p$。事实上，$(p)$ 是 $\mathbf{Z}$ 的
极大理想。因此，域可以是有限的：$\mathbf{F}_p$ 含有 $p$ 个元素。
\end{example}

\begin{example}
% P01-E0235: the frozen label misspells “polynomial” as “polymial”.
\label{fields-example-quotient-polymial-ring}
在主理想整环中，由不可约元生成的理想是极大理想。现在，若 $k$ 是域，
则多项式环 $k[x]$ 是主理想整环。因此，若 $P \in k[x]$ 是不可约
多项式（即不能分解为较低次数因子乘积的非常数多项式），则
$k[x]/(P)$ 是域。它以自然方式包含 $k$ 的一个副本。
这是构造域的一种非常一般的方法。例如，复数域 $\mathbf{C}$
可构造为 $\mathbf{R}[x]/(x^2 + 1)$。
\end{example}

\begin{example}[分式域]
\label{fields-example-quotient-field}
回忆，给定整环 $A$，存在从 $A$ 到某个域 $F$ 的嵌入 $A \to F$；
其构造方式与由 $\mathbf{Z}$ 构造 $\mathbf{Q}$ 完全相同。
形式上，$F$ 的元素是分式 $a/b$ 的等价类，其中
$a, b \in A$ 且 $b \not = 0$。照例，$a/b = a'/b'$ 当且仅当
$ab' = ba'$。域 $F$ 称为 $A$ 的{\it 分式域}。
分式域具有如下普遍性质：给定到域 $K$ 的单射环映射
$\varphi : A \to K$，存在唯一映射 $\psi : F \to K$ 使图表
$$
\xymatrix{
F \ar[r]_\psi & K \\
A \ar[u] \ar[ru]_\varphi
}
$$
交换。事实上，这一映射的定义很清楚：置
$\psi(a/b) = \varphi(a)\varphi(b)^{-1}$；$\varphi$ 的单射性保证，
若 $ b \not = 0$，则 $\varphi(b) \not = 0$。
\end{example}

\begin{example}[有理函数域]
\label{fields-example-field-of-rational-functions}
若 $k$ 是域，则可以考虑 $k$ 上的有理函数域 $k(x)$。
它是多项式环 $k[x]$ 的分式域。换言之，它由商 $F/G$ 的等价类组成，
其中 $F, G \in k[x]$ 且 $G \not = 0$。
\end{example}

\begin{example}
\label{fields-example-field-of-meromorphic-functions}
令 $X$ 为黎曼曲面。以 $\mathbf{C}(X)$ 表示 $X$ 上亚纯函数的集合。
在函数的乘法与加法下，$\mathbf{C}(X)$ 是环。事实上，
$\mathbf{C}(X)$ 是域：若非零函数 $f(z)$ 亚纯，则 $1/f(z)$ 亦亚纯。
例如，令 $S^2$ 为黎曼球面；复分析告诉我们，亚纯函数环
$\mathbf{C}(S^2)$ 就是有理函数域 $\mathbf{C}(z)$。
\end{example}



\section{向量空间}
\label{fields-section-vector-spaces}

\noindent
域之所以性质良好，一个原因是域上的模（即向量空间）理论非常简单。

\begin{lemma}
\label{fields-lemma-vector-space-is-free}
若 $k$ 是域，则每个 $k$-模都是自由模。
\end{lemma}

\begin{proof}
事实上，线性代数告诉我们，$k$-模（即向量空间）$V$ 有一个
{\it 基} $\mathcal{B} \subset V$；它给出以 $\mathcal{B}$
为基的自由向量空间到 $V$ 的同构。
\end{proof}

\begin{lemma}
\label{fields-lemma-field-semi-simple}
域上的每个模正合列都分裂。
\end{lemma}

\begin{proof}
由引理 \ref{fields-lemma-vector-space-is-free} 立即得到，因为每个向量空间
都是投射模。
\end{proof}

\noindent
这也是后续各章中的许多理论不会对域给出太多新内容的另一个原因：
域上的模行为如此简单。注意，引理 \ref{fields-lemma-field-semi-simple}
是关于 $k$-模的{\it 范畴}的陈述（这里 $k$ 是域），因为正合性的概念
本质上是关于箭头的，即只使用范畴论概念；事实上，它可以在所谓的
{\it 阿贝尔范畴}中表述。

\medskip\noindent
此后，由于域上的模理论就是线性代数，而域的理想理论并不十分有趣，
我们将研究本章真正的主题：域的{\it 扩张}。


\section{域的特征}
\label{fields-section-more-fields}

\noindent
在环范畴中，$\mathbf{Z}$ 是一个{\it 始对象}：对任意环 $R$，
恰有一个从 $\mathbf{Z}$ 到它的映射。域范畴中不存在这样的始对象。
不过，存在一族对象，使每个域都恰好从其中一个对象接受唯一映射，
而不从其余对象接受任何映射。
% P01-E0236: the frozen English sentence at lines 196-197 has the malformed
% phrase “every field can be mapped into ... by exactly one of them”.

\medskip\noindent
令 $F$ 为域。把 $F$ 看作环，得到环映射
$f : \mathbf{Z} \to F$。该环映射的像是整环（因为它是域的子环），
故 $f$ 的核是 $\mathbf{Z}$ 中的素理想。因此，$f$ 的核或者是
$(0)$，或者是某个素数 $p$ 所生成的 $(p)$。

\medskip\noindent
在第一种情形下，$f$ 是单射；此时我们把 $\mathbf{Z}$ 看作
$F$ 的子环。又因 $F$ 的每个非零元素都可逆，所以对
$p, q \in \mathbf{Z}$ 且 $q \not = 0$，谈论 $p/q \in F$
是有意义的。因此，在这种情形下，我们可以并且确实把
$\mathbf{Q}$ 看作 $F$ 的子环。容易看出，此时它是 $F$ 的最小子域。

\medskip\noindent
在第二种情形下，即 $\Ker(f) = (p)$ 时，
$\mathbf{Z}/(p) = \mathbf{F}_p$ 是 $F$ 的子环。显然，它是
$F$ 的最小子域。

\medskip\noindent
由此可见，每个域都含有一个最小子域；它或者是 $\mathbf{Q}$，
或者是对某个素数 $p$ 而言的有限域 $\mathbf{F}_p$。
% P01-E0237: frozen lines 220-222 say “either Q or finite equal to F_p”.

\begin{definition}
\label{fields-definition-characteristic}
域 $F$ 的{\it 特征}定义如下：若 $\mathbf{Z} \subset F$，则特征为
$0$；若 $F$ 中 $p = 0$，则特征为素数 $p$。$F$ 的{\it 素子域}
是 $F$ 的最小子域：特征为零时，它是
$\mathbf{Q} \subset F$；特征为 $p > 0$ 时，它是
$\mathbf{F}_p \subset F$。
\end{definition}

\noindent
容易看出，若 $E \subset F$ 是子域，则 $E$ 与 $F$ 的特征相同。

\begin{example}
\label{fields-example-characteristic}
$\mathbf{F}_p$ 的特征是 $p$，而 $\mathbf{Q}$ 的特征是 $0$。
\end{example}


\section{域扩张}
\label{fields-section-extensions}

\noindent
一般而言，我们所关心的与其说是孤立的域，不如说是域的{\it 扩张}。
这或许类似于不研究孤立的环，而研究固定环上的{\it 代数}。
域的便利之处在于，下一个结果表明，“一个域上的另一个域”这一概念
恰好就是域扩张。

\begin{lemma}
\label{fields-lemma-field-maps-injective}
若 $F$ 是域而 $R$ 是非零环，则任意环同态
$\varphi : F \to R$ 都是单射。
\end{lemma}

\begin{proof}
事实上，令 $a \in \Ker(\varphi)$ 为非零元素。则
$\varphi(1) = \varphi(a^{-1} a) = \varphi(a^{-1}) \varphi(a) = 0$。
因此 $1 = \varphi(1) = 0$，从而 $R$ 是零环。
\end{proof}

\begin{definition}
\label{fields-definition-extension}
若域 $F$ 包含于域 $E$，则称 $E$ 是 $F$ 的一个{\it 域扩张}。
我们写 $E/F$ 表示 $E$ 是 $F$ 的扩张。
\end{definition}

\noindent
因此，若 $F, F'$ 是域，且 $F \to F'$ 是任意环同态，
则由引理 \ref{fields-lemma-field-maps-injective} 可知它是单射；略微滥用术语，
可把 $F'$ 看作 $F$ 的扩张。或者说，$F$ 的域扩张就是恰好为域的
$F$-代数。这与一般环的情形完全不同，因为环同态未必是单射。

\medskip\noindent
令 $k$ 为域。存在一个由 $k$ 的域扩张组成的{\it 范畴}。
该范畴的对象是扩张 $E/k$，即域之间的（必然单射的）态射
$$
k \to E,
$$
而扩张 $E/k$ 与 $E'/k$ 之间的态射是 $k$-代数态射
$E \to E'$；等价地，它是交换图表
$$
\xymatrix{
E \ar[rr] & & E' \\
& k \ar[ru] \ar[lu] &
}
$$
在 $k$ 的扩张范畴中，态射 $E \to E'$ 的集合记作 $\Mor_k(E, E')$。
% P01-E0238: frozen line 295 says “morphisms from E \to E'”.

\begin{definition}
\label{fields-definition-tower}
一个域塔 $E_n/E_{n - 1}/\ldots/E_0$ 由域扩张序列
$E_n/E_{n - 1}$、$E_{n - 1}/E_{n - 2}$、$\ldots$、$E_1/E_0$
组成。
\end{definition}

\noindent
下面给出几个域扩张的例。

\begin{example}
\label{fields-example-monogenic-extension}
令 $k$ 为域，$P \in k[x]$ 为不可约多项式。我们已经看到
$k[x]/(P)$ 是域（例 \ref{fields-example-quotient-polymial-ring}）。
它又以显然的方式成为 $k$-代数，故是 $k$ 的扩张。
\end{example}

\begin{example}
\label{fields-example-field-of-meromorphic-functions-extension-C}
若 $X$ 是黎曼曲面，则其亚纯函数域
$\mathbf{C}(X)$（例 \ref{fields-example-field-of-meromorphic-functions}）
是 $\mathbf{C}$ 的扩域，因为 $\mathbf{C}$ 的任意元素都在 $X$
上给出一个亚纯的、事实上全纯的常值函数。
\end{example}

\noindent
令 $F/k$ 为域扩张，$S \subset F$ 为任意子集。则 $F$ 中存在一个
包含 $S$ 的{\it 最小}子扩张（即包含 $k$ 的 $F$ 的子域）。
为此，考虑所有同时包含 $S$ 与 $k$ 的 $F$ 的子域，并取其交；
容易验证该交仍是域。标准论证进一步表明，这恰是由 $k$ 与 $S$ 中的
元素经有限次基本代数运算（加、乘、减、除）所得的 $F$ 中元素的集合。

\begin{definition}
\label{fields-definition-generated-by}
令 $k$ 为域。若 $F/k$ 是域扩张且 $S \subset F$，则以
$k(S)$ 表示 $F$ 中包含 $k$ 与 $S$ 的最小子域。我们称
$S$ {\it 生成域扩张} $k(S)/k$。若 $S = \{\alpha\}$ 是单元素集，
则写 $k(\alpha)$ 而不写 $k(\{\alpha\})$。若存在有限子集
$S \subset F$ 使 $F = k(S)$，则称 $F/k$ 为{\it 有限生成域扩张}。
\end{definition}

\noindent
例如，$\mathbf{C}$ 在 $\mathbf{R}$ 上由 $i$ 生成。

\begin{exercise}
\label{fields-exercise-C-not-countably-generated}
证明 $\mathbf{C}$ 在 $\mathbf{Q}$ 上不存在可数生成集。
\end{exercise}

\noindent
现在分类由一个元素生成的扩张。

\begin{lemma}[单扩张的分类]
\label{fields-lemma-field-extension-generated-by-one-element}
若域扩张 $F/k$ 由一个元素生成，则它在 $k$ 上同构于有理函数域
$k(t)/k$，或者同构于某个扩张 $k[t]/(P)$，其中
$P \in k[t]$ 不可约。
\end{lemma}

\noindent
我们将看到，许多最重要的域扩张都由一个元素生成，因而这个结果确实有用。

\begin{proof}
令 $\alpha \in F$ 满足 $F = k(\alpha)$；按假设，这样的 $\alpha$
存在。存在环态射
$$
k[t] \to F
$$
把不定元 $t$ 送到 $\alpha$。其像是整环，故其核是素理想。
因此，该核或者是 $(0)$，或者是由某个不可约 $P \in k[t]$
生成的 $(P)$。

\medskip\noindent
若核为不可约多项式 $P \in k[t]$ 所生成的 $(P)$，则该映射经
$k[t]/(P)$ 分解，并诱导域态射 $k[t]/(P) \to F$。
由于其像包含 $\alpha$，容易看出该映射满射，因而是同构。
在这种情形下，$k[t]/(P) \simeq F$。

\medskip\noindent
若核平凡，则有单射 $k[t] \to F$。于是可定义分式域 $k(t)$
到 $F$ 的态射：给定 $R(t), Q(t) \in k[t]$ 的商
$R(t)/Q(t)$，把它送到 $R(\alpha)/Q(\alpha)$。
$k[t] \to F$ 的单射性意味着，除非 $Q$ 是零多项式，否则
$Q(\alpha) \neq 0$。$k[t]$ 的分式域就是有理函数域 $k(t)$，
故得到态射 $k(t) \to F$，其像包含 $\alpha$。
于是它是满射，因而是同构。
\end{proof}

\begin{lemma}
\label{fields-lemma-common-extension-field}
令 $k$ 为域，$E/k$ 与 $F/k$ 为域扩张。则存在共同扩域 $M/k$；
亦即，存在扩域以及 $k$-扩张态射 $E \to M$ 与 $F \to M$。
\end{lemma}

\begin{proof}
我们只证明 $E$ 是 $k$ 的有限生成域扩张的情形；一般情形由此通过
佐恩引理型论证得到（细节从略）。

\medskip\noindent
首先，设 $E$ 是 $k$ 的单扩张。由
引理 \ref{fields-lemma-field-extension-generated-by-one-element}，
这意味着或者 $E = k(t)$ 是有理函数域，或者
$E = k[t]/(P)$，其中 $P \in k[t]$ 是某个不可约多项式。
在第一种情形下，取 $M = F(t)$ 为有理函数域，并有显然的映射
$E \to M$ 与 $F \to M$。在第二种情形下，在 $F[t]$ 中取
$P$ 的像的一个不可约因子 $Q$，并令 $M = F[t]/(Q)$；
同样有显然的映射 $E \to M$ 与 $F \to M$。

\medskip\noindent
若 $E = k(\alpha_1, \ldots, \alpha_n)$，则对 $n$ 归纳可找到
扩张 $M/k$ 以及映射 $F \to M$ 与
$k(\alpha_1, \ldots, \alpha_{n - 1}) \to M$。
由上一段讨论的单扩张情形，可找到扩张
$M'/k(\alpha_1, \ldots, \alpha_{n - 1})$ 以及映射
$M \to M'$ 与 $k(\alpha_1, \ldots, \alpha_n) \to M'$。
把 $M'$ 看作 $k$ 的扩张即可。
\end{proof}



\section{有限扩张}
\label{fields-section-finite-extensions}

\noindent
若 $F/E$ 是域扩张，则显然 $F$ 也是 $E$ 上的向量空间
（标量作用就是 $F$ 中的乘法）。

\begin{definition}
\label{fields-definition-degree}
令 $F/E$ 为域扩张。把 $F$ 看作 $E$-向量空间时的维数称为该扩张的
{\it 次数}，记作 $[F : E]$。若 $[F : E] < \infty$，则称 $F$
是 $E$ 的{\it 有限}扩张。
\end{definition}

\begin{example}
\label{fields-example-C-over-R}
域 $\mathbf{C}$ 是 $\mathbf{R}$ 上以 $1, i$ 为基的二维向量空间。
因此，$\mathbf{C}$ 是 $\mathbf{R}$ 的二次有限扩张。
\end{example}

\begin{lemma}
\label{fields-lemma-finite-goes-up}
令 $K/E/F$ 为代数域扩张塔。
若 $K$ 在 $F$ 上有限，则 $K$ 在 $E$ 上有限。
\end{lemma}

\begin{proof}
由定义直接得到。
\end{proof}

\noindent
现在考虑最重要的特殊例，即引理
\ref{fields-lemma-field-extension-generated-by-one-element} 所给出的情形；
接下来的两个例将计算其次数。

\begin{example}[有理函数域的次数]
\label{fields-example-degree-rational-function-field}
若 $k$ 是任意域，则有理函数域 $k(t)$ {\it 不是}有限扩张。
例如，元素
$\left\{t^n, n \in \mathbf{Z}\right\}$ 在 $k$ 上线性无关。

\medskip\noindent
事实上，若 $k$ 不可数，则 $k(t)$ 作为 $k$-向量空间具有
{\it 不可数}维。为说明这一点，我们断言元素族
$\{1/(t- \alpha), \alpha \in k\} \subset k(t)$ 在
$k$ 上线性无关。它们之间的任意非平凡关系都会导致矛盾：
例如，在 $\mathbf{C}$ 上考虑时，这是因为把
$\frac{1}{t-\alpha}$ 看作 $\mathbf{C}$ 上的亚纯函数，
它在 $\alpha$ 处有极点，而在别处没有极点。
因此，对 $c_i \in k^*$ 以及互异的 $\alpha_i \in k$，
任意和 $\sum c_i \frac{1}{t - \alpha_i}$ 都会在每个
$\alpha_i$ 处有极点，因而不可能为零。

\medskip\noindent
有趣的是，这给出了复数域上希尔伯特零点定理的一个简短证明。
稍一般的结果见代数，定理
\ref{algebra-theorem-uncountable-nullstellensatz}。
\end{example}

\begin{lemma}
\label{fields-lemma-finite-finitely-generated}
有限域扩张是有限生成域扩张。逆命题不成立。
\end{lemma}

\begin{proof}
令 $F/E$ 为有限域扩张，并令 $\alpha_1, \ldots, \alpha_n$
为 $F$ 作为 $E$-向量空间的一组基。于是
$F = E(\alpha_1, \ldots, \alpha_n)$，故 $F/E$ 是有限生成域扩张。
例 \ref{fields-example-degree-rational-function-field} 表明逆命题不成立。
\end{proof}

\begin{example}[单代数扩张的次数]
\label{fields-example-degree-simple-algebraic-extension}
考虑例 \ref{fields-example-monogenic-extension} 中讨论的单生成域扩张
$E/k$。换言之，$E = k[t]/(P)$，其中
$P \in k[t]$ 是不可约多项式。则次数 $[E : k]$ 就是多项式
$P$ 的次数 $d = \deg(P)$。事实上，设
\begin{equation}
\label{fields-equation-P}
P = a_d t^d + a_{d - 1} t^{d - 1} + \ldots + a_0.
\end{equation}
其中 $a_d \not = 0$。于是 $1, t, \ldots, t^{d - 1}$ 在
$k[t]/(P)$ 中的像在 $k$ 上线性无关，因为涉及这些像的任意关系
次数都严格小于 $P$ 的次数，而 $P$ 是理想 $(P)$ 中次数最小的元素。

\medskip\noindent
反过来，集合 $S = \{1, t, \ldots, t^{d - 1}\}$（更准确地说，
这些元素的像）张成向量空间 $k[t]/(P)$。事实上，由
(\ref{fields-equation-P}) 可知 $a_d t^d$ 属于 $S$ 的张成空间。
由于 $a_d$ 可逆，$t^d$ 也属于 $S$ 的张成空间。
类似地，关系 $t P(t) = 0$ 表明 $t^{d + 1}$ 的像属于
$\{1, t, \ldots, t^d\}$ 的张成空间；由刚才所证，它也属于
$S$ 的张成空间。向上归纳可得，对 $n \geq d$，$t^n$ 的像
属于 $S$ 的张成空间。
\end{example}

\noindent
例如，这证实了 $[\mathbf{C}: \mathbf{R}] = 2$。
更一般地，若 $k$ 是域且 $\alpha \in k$ 不是平方元，则不可约多项式
$x^2 - \alpha \in k[x]$ 可用来构造二次扩张
$k[x]/(x^2 - \alpha)$。我们把它写作 $k(\sqrt{\alpha})$。
显然，这类扩张称为{\it 二次扩张}。

\medskip\noindent
次数的基本事实是它在{\it 域塔中具有乘法性}。

\begin{lemma}[乘法性]
\label{fields-lemma-multiplicativity-degrees}
给定域塔 $F/E/k$，则
$$
[F:k] = [F:E][E:k]
$$
\end{lemma}

\begin{proof}
令 $\alpha_1, \ldots, \alpha_n \in F$ 为 $F$ 的一组 $E$-基，
令 $\beta_1, \ldots, \beta_m \in E$ 为 $E$ 的一组 $k$-基。
所要证明的是乘积集合
$\{\alpha_i \beta_j, 1 \leq i \leq n, 1 \leq j \leq m\}$
是 $F$ 的一组 $k$-基。先验证它们在 $k$ 上张成 $F$。

\medskip\noindent
按假设，$\{\alpha_i\}$ 在 $E$ 上张成 $F$。因此若
$f \in F$，则存在 $a_i \in E$ 使
$$
f = \sum\nolimits_i a_i \alpha_i,
$$
并且对每个 $i$，可对某些 $b_{ij} \in k$ 写成
$a_i = \sum b_{ij} \beta_j$。合并这些等式得到
$$
f = \sum\nolimits_{i,j} b_{ij} \alpha_i \beta_j,
$$
这证明 $\{\alpha_i \beta_j\}$ 在 $k$ 上张成 $F$。

\medskip\noindent
现在假设存在一个非平凡关系
$$
\sum\nolimits_{i,j} c_{ij} \alpha_i \beta_j = 0
$$
其中 $c_{ij} \in k$。于是
$$
\sum\nolimits_i \alpha_i \left( \sum\nolimits_j c_{ij} \beta_j \right) = 0,
$$
且括号内各项属于 $E$，因为 $\beta_j$ 也属于其中。
$\{\alpha_i\}$ 在 $E$ 上线性无关，故各个内层和都为零；
再由 $\{\beta_j\}$ 在 $k$ 上线性无关，所有 $c_{ij}$ 都为零。
\end{proof}

\noindent
我们暂且岔开一步，引入一个略为旁支的定义。

\begin{definition}
\label{fields-definition-number-field}
若域 $K$ 的特征为 $0$，且扩张 $K/\mathbf{Q}$ 有限，
则称之为{\it 数域}。
\end{definition}

\noindent
数域是代数数论中的基本对象。后文将看到，对数域中整数
$\mathbf{Z}$ 的类似物，某种类似唯一分解的性质仍然成立
（尽管严格的唯一分解通常不成立！）。


\section{代数扩张}
\label{fields-section-algebraic-extensions}

\noindent
一类重要的扩张是其中每个元素都生成有限扩张的扩张。
% P01-E0239: source lines 605-606 have singular “class” with plural “are”.

\begin{definition}
\label{fields-definition-algebraic}
考虑域扩张 $F/E$。若元素 $\alpha \in F$ 是某个系数在
$E$ 中的非零多项式的根，则称 $\alpha$ 在 $E$ 上是
{\it 代数的}。若 $F$ 的所有元素都如此，则称
$F$ 是 $E$ 的{\it 代数扩张}。
\end{definition}

\noindent
由引理 \ref{fields-lemma-field-extension-generated-by-one-element}，
子扩张 $E(\alpha)$ 或者同构于有理函数域 $E(t)$，或者同构于
商环 $E[t]/(P)$，其中 $P \in E[t]$ 是不可约多项式。
在后一种情形中，$\alpha$ 在 $E$ 上代数（事实上，
引理 \ref{fields-lemma-field-extension-generated-by-one-element} 的证明
表明可取 $P$ 使 $\alpha$ 是 $P$ 的根）；在前一种情形中则不代数。

\begin{example}
\label{fields-example-C-algebraic-over-R}
域 $\mathbf{C}$ 是 $\mathbf{R}$ 的代数扩张。具体而言，若
$\alpha = a + ib$ 属于 $\mathbf{C}$，则
$\alpha^2 - 2a\alpha + a^2 + b^2 = 0$ 是 $\alpha$ 在
$\mathbf{R}$ 上满足的多项式方程。
\end{example}

\begin{example}
\label{fields-example-compact-riemann-surface-is-finite-over-P1}
令 $X$ 为紧黎曼曲面，并令
$f \in \mathbf{C}(X) - \mathbf{C}$ 为 $X$ 上任意非常值亚纯函数
（见例 \ref{fields-example-field-of-meromorphic-functions}）。
已知 $\mathbf{C}(X)$ 在由 $f$ 生成的子扩张
$\mathbf{C}(f)$ 上代数。这里不作证明。
\end{example}

\begin{lemma}
\label{fields-lemma-algebraic-goes-up}
令 $K/E/F$ 为域扩张塔。
\begin{enumerate}
\item 若 $\alpha \in K$ 在 $F$ 上代数，则 $\alpha$ 在
$E$ 上代数。
\item 若 $K$ 在 $F$ 上代数，则 $K$ 在 $E$ 上代数。
\end{enumerate}
\end{lemma}

\begin{proof}
由定义立即得到。
\end{proof}

\noindent
现在说明有限性与代数性之间的深刻联系。

\begin{lemma}
\label{fields-lemma-finite-is-algebraic}
有限扩张是代数扩张。事实上，扩张 $E/k$ 代数，当且仅当
对每个 $\alpha \in E$，由它生成的子扩张 $k(\alpha)/k$ 都有限。
\end{lemma}

\noindent
一般而言，代数扩张远不必有限。

\begin{proof}
设 $E/k$ 有限，次数为 $n$。取 $\alpha \in E$。则元素
$1, \alpha, \ldots, \alpha^n$ 在 $k$ 上线性相关，否则必有
$[E : k] > n$。一个线性相关关系就给出了 $\alpha$ 必须满足的
所需多项式。

\medskip\noindent
对于最后一个断言，注意单生成扩张 $k(\alpha)/k$ 有限，
当且仅当 $\alpha$ 在 $k$ 上代数；这是例
\ref{fields-example-degree-rational-function-field} 与
\ref{fields-example-degree-simple-algebraic-extension} 的结论。
因此，若 $E/k$ 代数，则每个 $k(\alpha)/k$（$\alpha \in E$）
都是有限扩张；反之亦然。
\end{proof}

\noindent
可以从上一证明（实质上是从例
\ref{fields-example-degree-rational-function-field} 与
\ref{fields-example-degree-simple-algebraic-extension}）抽取出一个结论：
单生成扩张有限，当且仅当它是代数扩张。下一结果将使用这一观察。

\begin{lemma}
\label{fields-lemma-algebraic-finitely-generated}
\begin{slogan}
有限生成代数扩张是有限扩张。
\end{slogan}
令 $k$ 为域，并令 $\alpha_1, \alpha_2, \ldots, \alpha_n$
为某个扩域中的元素，且每个 $\alpha_i$ 都在 $k$ 上代数。
则扩张 $k(\alpha_1, \ldots, \alpha_n)/k$ 有限。
也就是说，有限生成代数扩张是有限扩张。
\end{lemma}

\begin{proof}
事实上，每个扩张
$k(\alpha_{1}, \ldots, \alpha_{i+1})/k(\alpha_1, \ldots, \alpha_{i})$
都由一个元素生成且为代数扩张，因而有限。
由次数的乘法性（引理 \ref{fields-lemma-multiplicativity-degrees}）即得结论。
\end{proof}

\noindent
在 $\mathbf{Q}$ 上代数的复数统称为{\it 代数数}。
例如，$\sqrt{2}$ 是代数数，$i$ 是代数数，而 $\pi$ 不是。
代数数构成一个域，这是一个基本事实；然而，仅从“一个数恰在它满足
一个有理系数非零多项式方程时是代数数”这一定义出发，
如何证明这一点并不显然（例如，通过多项式方程来证明）。
% P01-E0240: source line 712 has singular “set” with plural “are”.

\begin{lemma}
\label{fields-lemma-algebraic-elements}
令 $E/k$ 为域扩张。$E$ 中在 $k$ 上代数的元素构成
$E/k$ 的一个子扩张。
\end{lemma}

\begin{proof}
令 $\alpha, \beta \in E$ 在 $k$ 上代数。由引理
\ref{fields-lemma-algebraic-finitely-generated}，
$k(\alpha, \beta)/k$ 是有限扩张。因此
$k(\alpha + \beta) \subset k(\alpha, \beta)$ 是有限扩张，
由引理 \ref{fields-lemma-finite-is-algebraic} 可知 $\alpha + \beta$
代数。$\alpha$ 与 $\beta$ 的差、积和商也同理。
\end{proof}

\noindent
域扩张的许多良好性质与环的情形一样，会沿域塔以及取复合域而保持。

\begin{lemma}
\label{fields-lemma-algebraic-permanence}
令 $E/k$ 与 $F/E$ 为代数域扩张。则 $F/k$ 是代数域扩张。
\end{lemma}

\begin{proof}
取 $\alpha \in F$。则 $\alpha$ 在 $E$ 上代数。
关键观察是，$\alpha$ 在 $k$ 的某个有限生成子扩张上代数。
% P01-E0241: source line 748 says “subextension of k”, although the
% constructed field contains k and lies in E.
也就是说，存在有限集 $S \subset E$，使得 $\alpha$ 在
$k(S)$ 上代数：这是清楚的，因为代数性意味着存在某个由
$\alpha$ 满足的多项式 $E[x]$，而可取 $S$ 为该多项式的系数集合。
因此，$\alpha$ 在 $k(S)$ 上代数。特别地，扩张
$k(S, \alpha)/ k(S)$ 有限。由于 $S$ 是有限集且 $k(S)/k$
代数，引理 \ref{fields-lemma-algebraic-finitely-generated} 表明
$k(S)/k$ 有限。利用乘法性（引理
\ref{fields-lemma-multiplicativity-degrees}），可知
$k(S,\alpha)/k$ 有限，所以 $\alpha$ 在 $k$ 上代数。
\end{proof}

\noindent
上一论证中的证明方法——在 $E$ 上代数这一性质
{\it 下降}到 $E$ 的一个有限生成子扩张——是一个在代数中反复出现的思想。
例如，它常使我们能把一般的交换代数问题归约到诺特情形。

\begin{lemma}
\label{fields-lemma-size-algebraic-extension}
令 $E/F$ 为代数域扩张。则 $E$ 的基数 $|E|$ 至多为
$\max(\aleph_0, |F|)$。
\end{lemma}

\begin{proof}
令 $S$ 为系数在 $F$ 中的非常数多项式集合。
对每个 $P \in S$，根的集合
$r(P, E) = \{\alpha \in E \mid P(\alpha) = 0\}$
是有限的（细节从略）。此外，$E$ 在 $F$ 上代数这一事实意味着
$E = \bigcup_{P \in S} r(P, E)$。显然，$S$ 的基数不超过
$\max(\aleph_0, |F|)$，因为它是 $F$ 的有限乘积的可数并。
因此，$E$ 的基数也有同样的上界。
\end{proof}

\begin{lemma}
\label{fields-lemma-subalgebra-algebraic-extension-field}
令 $E/F$ 为有限域扩张，或更一般地为代数域扩张。
任意子环 $F \subset R \subset E$ 都是域。
\end{lemma}

\begin{proof}
令 $\alpha \in R$ 非零。则 $1, \alpha, \alpha^2, \ldots$
都属于 $R$。由引理 \ref{fields-lemma-finite-is-algebraic}，存在非平凡关系
$a_0 + a_1 \alpha + \ldots + a_d \alpha^d = 0$，其中
$a_i \in F$。可以假设 $a_0 \not = 0$，因为否则可用
$\alpha$ 除该关系以降低 $d$。于是
$$
a_0 = \alpha (- a_1  - \ldots - a_d \alpha^{d - 1})
$$
这证明 $\alpha$ 的逆是 $R$ 中的元素
$a_0^{-1} (- a_1  - \ldots - a_d \alpha^{d - 1})$。
\end{proof}

\begin{lemma}
\label{fields-lemma-algebraic-extension-self-map}
令 $E/F$ 为代数域扩张。任意 $F$-代数映射
$f : E \to E$ 都是自同构。
% P01-E0242: source line 808 omits “be” after “Let E/F”.
\end{lemma}

\begin{proof}
若 $E/F$ 有限，则 $f : E \to E$ 是有限维向量空间之间的
$F$-线性单射（引理 \ref{fields-lemma-field-maps-injective}），因而是双射。
一般情形下仍可看出 $f$ 是单射。令 $\alpha \in E$，并令
$P \in F[x]$ 为满足 $P(\alpha) = 0$ 的多项式。
令 $E' \subset E$ 为由 $P$ 在 $E$ 中的各根
$\alpha = \alpha_1, \ldots, \alpha_n$ 生成的 $E$ 的子域。
由引理 \ref{fields-lemma-algebraic-finitely-generated}，$E'$ 在 $F$ 上有限。
由于 $f$ 保持根的集合，得到 $f|_{E'} : E' \to E'$。
由证明的第一部分，$f|_{E'}$ 是同构，所以 $\alpha$ 属于 $f$ 的像。
\end{proof}






\section{最小多项式}
\label{fields-section-minimal-polynomials}

\noindent
令 $E/k$ 为域扩张，并令 $\alpha \in E$ 在 $k$ 上代数。
则 $\alpha$ 满足 $k[x]$ 中的一个（非平凡）多项式方程。
考虑满足 $P(\alpha) = 0$ 的多项式 $P \in k[x]$ 所成的集合；
按假设，该集合不只含零多项式。容易看出它是一个{\it 理想}。
事实上，它是映射
$$
k[x] \to E, \quad x \mapsto \alpha
$$
的核。由于 $k[x]$ 是主理想整环，该理想存在一个{\it 生成元}
$P \in k[x]$。不失一般性地假设 $P$ 首一，则 $P$ 唯一确定。

\begin{definition}
\label{fields-definition-minimal-polynomial}
上述多项式 $P$ 称为 $\alpha$ 在 $k$ 上的{\it 最小多项式}。
\end{definition}

\noindent
最小多项式有如下刻画：它是使 $\alpha$ 为零的次数最小的首一多项式。
$P$ 的任意非常数倍数都有更高次数，并且只有 $P$ 的倍数才能使
$\alpha$ 为零。这解释了“{\it 最小}”这一名称。

\medskip\noindent
显然，最小多项式{\it 不可约}。这等价于断言，$k[x]$ 中由使
$\alpha$ 为零的多项式组成的理想是素理想。这又来自如下事实：
映射 $k[x] \to E, x \mapsto \alpha$ 的目标是整环
（甚至是域），故其核是素理想。

\begin{lemma}
\label{fields-lemma-degree-minimal-polynomial}
最小多项式的次数为 $[k(\alpha) : k]$。
\end{lemma}

\begin{proof}
这只是引理 \ref{fields-lemma-field-extension-generated-by-one-element}
中论证的重新表述：若 $P$ 是 $\alpha$ 的最小多项式，则映射
$$
k[x]/(P) \to k(\alpha), \quad x \mapsto \alpha
$$
如该证明所述是同构，而我们已经计算过这类扩张的次数
（见例 \ref{fields-example-degree-simple-algebraic-extension}）。
\end{proof}

\noindent
因此，上述证明的观察是：若 $\alpha \in E$ 代数，则
$k(\alpha) \subset E$ 同构于 $k[x]/(P)$。


\section{代数闭包}
\label{fields-section-algebraic-closure}

\noindent
“代数学基本定理”断言 $\mathbf{C}$ 是代数闭域。
这一结果有一个使用复分析中刘维尔定理的漂亮证明；我们将给出另一证明
（见引理 \ref{fields-lemma-C-algebraically-closed}）。
% P01-E0243: source lines 895-898 join two independent clauses with a comma.

\begin{definition}
\label{fields-definition-algebraically-closed}
若域 $F$ 的每个代数扩张 $E/F$ 都平凡，即 $E = F$，
则称该域为{\it 代数闭域}。
\end{definition}

\noindent
这未必是每本教材采用的定义。下面的引理将它与另一常用定义比较。

\begin{lemma}
\label{fields-lemma-algebraically-closed}
令 $F$ 为域。下列条件等价：
\begin{enumerate}
\item $F$ 是代数闭域；
\item $F$ 上的每个不可约多项式都是一次多项式；
\item $F$ 上的每个非常数多项式都有根；
\item $F$ 上的每个非常数多项式都是一次因子的乘积。
\end{enumerate}
\end{lemma}

\begin{proof}
若 $F$ 是代数闭域，则每个不可约多项式都是一次的。
事实上，若存在次数 $> 1$ 的不可约多项式，它就生成一个非平凡的有限
（因而代数）域扩张；见例
\ref{fields-example-degree-simple-algebraic-extension}。
所以 (1) 蕴含 (2)。若每个不可约多项式都是一次的，
则每个不可约多项式都有根，从而每个非常数多项式都有根。
故 (2) 蕴含 (3)。

\medskip\noindent
假设每个非常数多项式都有根。令 $P \in F[x]$ 非常数。
若 $P(\alpha) = 0$ 且 $\alpha \in F$，则由带余除法可知，
对某个 $Q \in F[x]$ 有 $P = (x - \alpha)Q$。
因此，对次数归纳可知，任意非常数多项式都可写成
$c \prod (x - \alpha_i)$。

\medskip\noindent
最后，假设 $F$ 上每个非常数多项式都是一次因子的乘积。
令 $E/F$ 为代数扩张。则 $E$ 的所有单子扩张
$F(\alpha)/F$ 必然平凡（因为按假设，唯一的不可约多项式是一次的）。
因此 $E = F$。这表明 (4) 蕴含 (1)，证明完成。
\end{proof}

\noindent
现在要定义域的一个“泛”代数扩张。实际上必须谨慎：
代数闭包{\it 不是}泛对象。也就是说，代数闭包并不在
{\it 唯一}同构意义下唯一；它只在同构意义下唯一。
不过，即使它不具有函子性，仍会非常方便。

\begin{definition}
\label{fields-definition-algebraic-closure}
令 $F$ 为域。$F$ 的一个{\it 代数闭包}是包含 $F$ 的域
$\overline{F}$，并且满足：
\begin{enumerate}
\item $\overline{F}$ 在 $F$ 上代数。
\item $\overline{F}$ 是代数闭域。
\end{enumerate}
\end{definition}

\noindent
若 $F$ 是代数闭域，则 $F$ 就是它自身的代数闭包。
现在证明基本的存在性结果。

\begin{theorem}
\label{fields-theorem-existence-algebraic-closure}
每个域都有代数闭包。
\end{theorem}

\noindent
这个证明与本章其余内容大多没有直接关系。不过，我们需要知道，
把一个域嵌入代数闭域是{\it 可能的}，并且以后常假定已经作了这种嵌入。

\begin{proof}
令 $F$ 为域。由引理 \ref{fields-lemma-size-algebraic-extension}，
$F$ 的任意代数扩张的基数不超过 $\max(\aleph_0, |F|)$。
选择一个包含 $F$ 的集合 $S$，使 $|S| > \max(\aleph_0, |F|)$。
考虑三元组 $(E, \sigma_E, \mu_E)$，其中
\begin{enumerate}
\item $E$ 是满足 $F \subset E \subset S$ 的集合；
\item $\sigma_E : E \times E \to E$ 与
$\mu_E : E \times E \to E$ 是集合映射，使得
$(E, \sigma_E, \mu_E)$ 定义 $F$ 的域扩张结构
（特别地，对 $a, b \in F$ 有 $\sigma_E(a, b) = a +_F b$，
$\mu_E$ 也类似）；
\item $E/F$ 是代数域扩张。
\end{enumerate}
所有三元组 $(E, \sigma_E, \mu_E)$ 构成一个集合 $I$。
对 $i \in I$，把相应的 $F$-域扩张记作
$E_i = (E_i, \sigma_i, \mu_i)$。
% P01-E0244: source lines 991-992 omit “by” in “denote ... the corresponding”.
在 $I$ 上定义偏序：规定 $i \leq i'$，当且仅当
$E_i \subset E_{i'}$（这是有意义的，因为 $E_i$ 与 $E_{i'}$
都是同一集合 $S$ 的子集），并且
$\sigma_i = \sigma_{i'}|_{E_i \times E_i}$ 与
$\mu_i = \mu_{i'}|_{E_i \times E_i}$；换言之，$E_{i'}$
是 $E_i$ 的域扩张。

\medskip\noindent
令 $T \subset I$ 为全序子集。显然，
$E_T = \bigcup_{i \in T} E_i$ 连同诱导映射
$\sigma_T = \bigcup \sigma_i$ 与 $\mu_T = \bigcup \mu_i$
仍是 $I$ 的元素。换言之，$I$ 的每个全序子集都在 $I$ 中有上界。
% P01-E0245: source line 1003 says “totally order” and “a upper bound”.
由佐恩引理，$I$ 中存在极大元素 $(E, \sigma_E, \mu_E)$。
我们断言 $E$ 是代数闭包。按 $I$ 的定义，扩张 $E/F$ 代数，
故只需证明 $E$ 是代数闭域。

\medskip\noindent
反证。假设 $E$ 不是代数闭域。则存在 $E$ 上次数 $> 1$ 的
不可约多项式 $P$；见引理 \ref{fields-lemma-algebraically-closed}。
由引理 \ref{fields-lemma-finite-is-algebraic}，得到非平凡有限扩张
$E' = E[x]/(P)$。注意，由引理 \ref{fields-lemma-algebraic-permanence}，
$E'/F$ 代数。因此 $E'$ 的基数满足
$\leq \max(\aleph_0, |F|)$。由初等集合论，可以把已有单射
$E \subset S$ 扩张为单射 $E' \to S$。换言之，可把
$E'$ 看作集合 $I$ 的元素，这与 $E$ 的极大性矛盾。
矛盾完成证明。
\end{proof}

\begin{lemma}
\label{fields-lemma-map-into-algebraic-closure}
令 $F$ 为域，令 $\overline{F}$ 为 $F$ 的代数闭包，并令
$M/F$ 为代数扩张。则存在 $F$-扩张态射
$M \to \overline{F}$。
\end{lemma}

\begin{proof}
考虑由偶对 $(E, \varphi)$ 构成的集合 $I$，其中
$F \subset E \subset M$ 是子扩张，而
$\varphi : E \to \overline{F}$ 是 $F$-扩张态射。
在 $I$ 上定义偏序：$(E, \varphi) \leq (E', \varphi')$，
当且仅当 $E \subset E'$ 且 $\varphi'|_E = \varphi$。
若 $T = \{(E_t,  \varphi_t)\} \subset I$ 是全序子集，则
$\bigcup \varphi_t : \bigcup E_t \to \overline{F}$ 是 $I$ 的元素。
因此，$I$ 的每个全序子集都有上界。由佐恩引理，$I$ 中存在
极大元素 $(E, \varphi)$。我们断言 $E = M$，这将完成证明。
否则，取 $\alpha \in M$ 且 $\alpha \not \in E$。
则 $\alpha$ 在 $E$ 上代数；见引理 \ref{fields-lemma-algebraic-goes-up}。
% P01-E0246: source line 1040 says “The alpha is algebraic” rather than “Then”.
令 $P$ 为 $\alpha$ 在 $E$ 上的最小多项式，令
$P^\varphi$ 为 $P$ 经 $\varphi$ 在 $\overline{F}[x]$ 中的像。
由于 $\overline{F}$ 是代数闭域，$P^\varphi$ 在
$\overline{F}$ 中有根 $\beta$。于是可把 $\varphi$ 扩张为
$\varphi' : E(\alpha) = E[x]/(P) \to \overline{F}$，
方法是把 $x$ 送到 $\beta$。这与 $(E, \varphi)$ 的极大性矛盾，
正合所需。
\end{proof}

\begin{lemma}
\label{fields-lemma-algebraic-closures-isomorphic}
一个域的任意两个代数闭包同构。
\end{lemma}

\begin{proof}
令 $F$ 为域。若 $M$ 与 $\overline{F}$ 都是 $F$ 的代数闭包，
则由引理 \ref{fields-lemma-map-into-algebraic-closure}，存在
$F$-扩张态射 $\varphi : M \to \overline{F}$。
其像 $\varphi(M)$ 是代数闭域。另一方面，由引理
\ref{fields-lemma-algebraic-goes-up}，扩张
$\varphi(M) \subset \overline{F}$ 代数。
因此 $\varphi(M) = \overline{F}$。
\end{proof}





\section{互素多项式}
\label{fields-section-relatively-prime}

\noindent
令 $K$ 为代数闭域。正如引理 \ref{fields-lemma-algebraically-closed}
中所见，环 $K[x]$ 的理想结构非常简单。特别地，每个多项式
$P \in K[x]$ 都可写成
$$
P = c(x - \alpha_1) \ldots (x - \alpha_n)
$$
其中 $c$ 是常数项，而 $\alpha_1, \ldots, \alpha_n \in k$
是 $P$ 的根（按重数计）。显然，$K[x]$ 中唯一的不可约多项式是
一次多项式 $c(x - \alpha)$，其中 $c, \alpha \in K$
且 $c \neq 0$。
% P01-E0247: in source line 1081, c is the leading coefficient, not the
% constant term, in the displayed factorization.
% P01-E0248: source line 1081 places the roots in lowercase k although the
% entire paragraph fixes the algebraically closed field K.

\begin{definition}
\label{fields-definition-relatively-prime}
若 $k$ 是任意域，则称 $k[x]$ 中两个多项式{\it 互素}，
若它们在 $k[x]$ 中生成单位理想。
\end{definition}

\noindent
继续上面的讨论：若 $K$ 是代数闭域，则 $K[x]$ 中两个多项式互素，
当且仅当它们没有公共根。这是因为 $K[x]$ 的极大理想都形如
$(x - \alpha)$，其中 $\alpha \in K$。因此，若
$F, G \in K[x]$ 没有公共根，则 $(F, G)$ 不可能包含在任何
$(x - \alpha)$ 中（否则它们会在 $\alpha$ 处有公共根）。

\medskip\noindent
若 $k$ {\it 不是}代数闭域，上述讨论仍可给出关于 $k[x]$
中两个多项式何时生成单位理想的信息。

\begin{lemma}
\label{fields-lemma-relatively-prime-polynomials}
$k[x]$ 中两个多项式互素，恰好当它们在 $k$ 的某个代数闭包
$\overline{k}$ 中没有公共根。
\end{lemma}

\begin{proof}
所要证明的是：任意两个多项式 $P, Q$ 在 $k[x]$ 中生成 $(1)$，
当且仅当它们在 $\overline{k}[x]$ 中生成 $(1)$。
这是线性代数的一个结论：系数在 $k$ 中的线性方程组有解，
当且仅当它在 $k$ 的任意扩域中有解。
因此，可以归约到代数闭域的情形；此时结论由已经证明的内容显然成立。
\end{proof}





\section{可分代数扩张}
\label{fields-section-separable-extensions}

\noindent
在特征 $p$ 的域上，不可约多项式会出现一种特殊现象。
下面的引理对此作出说明。

\begin{lemma}
\label{fields-lemma-irreducible-polynomials}
令 $F$ 为域，令 $P \in F[x]$ 为 $F$ 上的不可约多项式。
令 $P' = \text{d}P/\text{d}x$ 为 $P$ 关于 $x$ 的导数。
则会出现下列两种情形之一：
\begin{enumerate}
\item $P$ 与 $P'$ 互素；或者
\item $P'$ 是零多项式。
\end{enumerate}
第二种情形只可能在 $F$ 的特征 $p > 0$ 时发生。
此时 $P(x) = Q(x^q)$，其中 $q = p^f$ 是 $p$ 的幂，
而 $Q \in F[x]$ 是不可约多项式，且 $Q$ 与 $Q'$ 互素。
\end{lemma}

\begin{proof}
注意 $P'$ 的次数 $< \deg(P)$。因此，若 $P$ 与 $P'$ 不互素，
则 $(P, P') = (R)$，其中 $R$ 的次数 $< \deg(P)$，
这与 $P$ 的不可约性矛盾。这证明了 (1) 与 (2) 的二分情形。

\medskip\noindent
假设处于情形 (2)，且 $P = a_d x^d + \ldots + a_0$。
则 $P' = da_d x^{d - 1} + \ldots + a_1$。在特征 $0$ 时，
这迫使 $a_d, \ldots, a_1 = 0$，从而 $P$ 为常数，矛盾。
因此，特征 $p$ 必为正。此时，条件 $P' = 0$ 迫使
凡 $p$ 不整除 $i$ 时都有 $a_i = 0$。
换言之，对某个非常数多项式 $P_1$，有
$P(x) = P_1(x^p)$。显然，$P_1$ 也不可约。
对次数归纳可知，$P_1(x) = Q(x^q)$，如引理所述；
因此 $P(x) = Q(x^{pq})$，引理得证。
\end{proof}

\begin{definition}
\label{fields-definition-separable}
令 $F$ 为域，令 $K/F$ 为域扩张。
\begin{enumerate}
\item 若 $F$ 上的不可约多项式 $P$ 与其导数互素，
则称它{\it 可分}。
\item 给定在 $F$ 上代数的 $\alpha \in K$，若它的最小多项式在
$F$ 上可分，则称 $\alpha$ 在 $F$ 上{\it 可分}。
\item 若 $K$ 是 $F$ 的代数扩张，并且 $K$ 的每个元素都在
$F$ 上可分，则称 $K$ 在 $F$ 上
{\it 可分}\footnote{对于非代数扩张，这一定义没有意义，
也不是正确的定义。请参见代数，第
\ref{algebra-section-separability} 节与第
\ref{algebra-section-separability-continued} 节。}
\end{enumerate}
\end{definition}

\noindent
由引理 \ref{fields-lemma-irreducible-polynomials}，在特征 $0$ 时，
每个不可约多项式都可分，扩域中的每个代数元素都可分，
并且每个代数扩张都可分。

\begin{lemma}
\label{fields-lemma-separable-goes-up}
令 $K/E/F$ 为代数域扩张塔。
\begin{enumerate}
\item 若 $\alpha \in K$ 在 $F$ 上可分，则 $\alpha$ 在
$E$ 上可分。
\item 若 $K$ 在 $F$ 上可分，则 $K$ 在 $E$ 上可分。
\end{enumerate}
\end{lemma}

\begin{proof}
以下不再特别说明地使用引理 \ref{fields-lemma-irreducible-polynomials}。
令 $P$ 为 $\alpha$ 在 $F$ 上的最小多项式，令 $Q$ 为
$\alpha$ 在 $E$ 上的最小多项式。则在多项式环 $E[x]$ 中
$Q$ 整除 $P$；设 $P = QR$。于是 $P' = Q'R + QR'$。
因此，若 $Q' = 0$，则 $Q$ 同时整除 $P$ 与 $P'$，
故由该引理有 $P' = 0$。这证明了 (1)。
(2) 立即由 (1) 和定义得到。
\end{proof}

\begin{lemma}
\label{fields-lemma-recognize-separable}
令 $F$ 为域。$F$ 上的不可约多项式 $P$ 可分，当且仅当
$P$ 在 $F$ 的某个代数闭包中具有两两不同的根。
\end{lemma}

\begin{proof}
假设 $\alpha \in \overline{F}$ 同时是 $P$ 与 $P'$ 的根。
则对某个多项式 $Q$ 有 $P = (x - \alpha)Q$。
求导得到 $P' = Q + (x - \alpha)Q'$。因此 $\alpha$ 是 $Q$ 的根。
由此可见，若 $P$ 与 $P'$ 有公共根，则 $P$ 的根并非两两不同。
反过来，若 $P$ 有重根，即 $(x - \alpha)^2$ 整除 $P$，
则 $\alpha$ 同时是 $P$ 与 $P'$ 的根。结合引理
\ref{fields-lemma-relatively-prime-polynomials} 即得结论。
\end{proof}

\begin{lemma}
\label{fields-lemma-nr-roots-unchanged}
令 $F$ 为域，令 $\overline{F}$ 为 $F$ 的代数闭包。
令 $p > 0$ 为 $F$ 的特征，并令 $P$ 为 $F$ 上的多项式。
则 $P$ 与 $P(x^p)$ 在 $\overline{F}$ 中的根集具有相同基数
（不计重数）。
% P01-E0249: source lines 1228-1229 use plural “have” with singular “set”.
\end{lemma}

\begin{proof}
显然，$\alpha$ 是 $P(x^p)$ 的根，当且仅当 $\alpha^p$ 是
$P$ 的根。换言之，$P(x^p)$ 的根就是各个 $x^p - \beta$
的根，其中 $\beta$ 是 $P$ 的根。因此，只需证明映射
$\overline{F} \to \overline{F}$，$\alpha \mapsto \alpha^p$
是双射。它是满射，因为 $\overline{F}$ 是代数闭域，
即每个元素都有 $p$ 次根。它也是单射，因为
$\alpha^p = \beta^p$ 蕴含 $(\alpha - \beta)^p = 0$，
这里使用了特征为 $p$。当然，在域中，$x^p = 0$ 蕴含 $x = 0$。
\end{proof}

\noindent
令 $F$ 为域，令 $P$ 为 $F$ 上的不可约多项式。
已知 $P = Q(x^q)$，其中 $Q$ 是某个可分不可约多项式
（引理 \ref{fields-lemma-irreducible-polynomials}），而 $q$ 是特征
$p$ 的一个幂（若特征为零，则
$q = 1$\footnote{本章中一个方便的约定是置 $0^0 = 1$。}
且 $Q = P$）。由引理 \ref{fields-lemma-nr-roots-unchanged}，
$P$ 与 $Q$ 在 $F$ 的任意代数闭包中的根数相同。
由引理 \ref{fields-lemma-recognize-separable}，这个数等于 $Q$ 的次数。

\begin{definition}
\label{fields-definition-separable-degree}
令 $F$ 为域，令 $P$ 为 $F$ 上的不可约多项式。
$P$ 的{\it 可分次数}是 $P$ 在 $F$ 的任意代数闭包中的根集基数
（见上面的讨论），记作 $\deg_s(P)$。
\end{definition}

\noindent
$P$ 的可分次数总是整除其次数，商是特征的一个幂。
若特征为零，则 $\deg_s(P) = \deg(P)$。

\begin{situation}
\label{fields-situation-finitely-generated}
这里 $F$ 为域，$K/F$ 是由元素
$\alpha_1, \ldots, \alpha_n \in K$ 生成的有限扩张。
% P01-E0250: source line 1270 says “Here F be a field”.
置 $K_0 = F$，以及
$$
K_i = F(\alpha_1, \ldots, \alpha_i)
$$
从而得到有限扩张塔
$K = K_n / K_{n - 1} / \ldots / K_0 = F$。
以 $P_i$ 表示 $\alpha_i$ 在 $K_{i - 1}$ 上的最小多项式。
最后，固定 $F$ 的一个代数闭包 $\overline{F}$。
\end{situation}

\noindent
令 $F$、$K$、$\alpha_i$ 与 $\overline{F}$ 如情形
\ref{fields-situation-finitely-generated} 所述。
假设 $\varphi : K \to \overline{F}$ 是 $F$-扩张态射。
于是得到映射 $\varphi_i : K_i \to \overline{F}$。
$P_i \in K_{i - 1}[x]$ 经 $\varphi_{i - 1}$ 的像是多项式
$P_i^{\varphi_{i - 1}} \in \overline{F}[x]$，并且
$\varphi(\alpha_i)$ 是该多项式的根。

\begin{lemma}
\label{fields-lemma-count-embeddings}
在情形 \ref{fields-situation-finitely-generated} 中，对应
$$
\Mor_F(K, \overline{F})
\longrightarrow
\{(\beta_1, \ldots, \beta_n)\text{ 如下}\},
\quad
\varphi \longmapsto (\varphi(\alpha_1), \ldots, \varphi(\alpha_n))
$$
是双射。这里右端是 $\overline{F}$ 中元素的 $n$-元组
$(\beta_1, \ldots, \beta_n)$ 所成的集合，并满足：
\begin{enumerate}
\item $\beta_1 \in \overline{F}$ 是 $P_1$ 的根；
令 $\varphi_1 : K_1 \to \overline{F}$ 为把 $\alpha_1$
送到 $\beta_1$ 的 $F$-同态；
\item $\beta_2 \in \overline{F}$ 是 $P_2^{\varphi_1}$ 的根；
令 $\varphi_2 : K_2 \to \overline{F}$ 为把 $\alpha_2$
送到 $\beta_2$ 且扩张 $\varphi_1$ 的同态；
\item 如此继续，直到
\item $\beta_n \in \overline{F}$ 是
$P_n^{\varphi_{n - 1}}$ 的根。
\end{enumerate}
每一步中，同态 $\varphi_i$ 都存在且唯一，因为
$K_i = K_{i - 1}[x]/(P_i)$，而 $\beta_i$ 是
$P_i^{\varphi_{i - 1}}$ 的根。
% P01-E0251: source lines 1313 and 1320 misspell “homomorphism” as
% “homorphism”.
\end{lemma}

\begin{proof}
上文已经讨论了从左到右的映射。
令 $(\beta_1, \ldots, \beta_n)$ 为右端的元素。
令 $\varphi : K = K_n \to \overline{F}$ 为把 $\alpha_n$
送到 $\beta_n$ 且扩张 $\varphi_{n - 1}$ 的唯一同态。
唯一性表明两个构造互为逆。
\end{proof}

\begin{lemma}
\label{fields-lemma-count-embeddings-explicitly}
在情形 \ref{fields-situation-finitely-generated} 中，有
$|\Mor_F(K, \overline{F})| = \prod_{i = 1}^n \deg_s(P_i)$。
\end{lemma}

\begin{proof}
这立即由引理 \ref{fields-lemma-count-embeddings} 得到。
注意，这里默默使用的一个关键事实是不可约多项式可分次数的良定义性；
这一点已在定义 \ref{fields-definition-separable-degree} 之前说明。
\end{proof}

\noindent
现在用上述结果刻画可分域扩张。

\begin{lemma}
\label{fields-lemma-separably-generated-separable}
假设与记号如情形 \ref{fields-situation-finitely-generated}。
若每个 $P_i$ 都可分，即每个 $\alpha_i$ 都在
$K_{i - 1}$ 上可分，则
$$
|\Mor_F(K, \overline{F})| = [K : F]
$$
且域扩张 $K/F$ 可分。若某个 $\alpha_i$ 在 $K_{i - 1}$ 上
不可分，则
$|\Mor_F(K, \overline{F})| < [K : F]$。
\end{lemma}

\begin{proof}
若 $\alpha_i$ 在 $K_{i - 1}$ 上可分，则
$\deg_s(P_i) = \deg(P_i) = [K_i : K_{i - 1}]$
（最后一个等式由引理 \ref{fields-lemma-degree-minimal-polynomial} 得到）。
由乘法性（引理 \ref{fields-lemma-multiplicativity-degrees}），有
$$
[K : F] = \prod [K_i : K_{i - 1}] = \prod \deg(P_i) =
\prod \deg_s(P_i) = |\Mor_F(K, \overline{F})|
$$
其中最后一个等式是引理
\ref{fields-lemma-count-embeddings-explicitly}。完全相同的论证表明，
若某个 $\alpha_i$ 在 $K_{i - 1}$ 上不可分，则有严格不等式
$|\Mor_F(K, \overline{F})| < [K : F]$。

\medskip\noindent
最后，再次假设每个 $\alpha_i$ 都在 $K_{i - 1}$ 上可分。
我们将证明 $K/F$ 可分。任取 $\gamma = \gamma_1 \in K$。
可以找到另外一些元素 $\gamma_2, \ldots, \gamma_m$，使
$K = F(\gamma_1, \ldots, \gamma_m)$（例如，可取
$\gamma_2 = \alpha_1, \ldots, \gamma_{n + 1} = \alpha_n$）。
由引理已经证明的后一部分可知，若 $\gamma$ 在 $F$ 上不可分，
则有严格不等式 $|\Mor_F(K, \overline{F})| < [K : F]$，
这与引理最先证明的部分矛盾。
% P01-E0252: source lines 1379-1380 say “already prove above”.
\end{proof}

\begin{lemma}
\label{fields-lemma-separable-equality}
令 $K/F$ 为有限域扩张，令 $\overline{F}$ 为 $F$ 的代数闭包。
则
$$
|\Mor_F(K, \overline{F})| \leq [K : F]
$$
且等号成立，当且仅当 $K$ 在 $F$ 上可分。
\end{lemma}

\begin{proof}
这是引理 \ref{fields-lemma-separably-generated-separable} 的推论。
事实上，由于 $K/F$ 有限，可找到有限多个元素
$\alpha_1, \ldots, \alpha_n \in K$，它们在 $F$ 上生成 $K$
（例如，可取 $\alpha_i$ 为 $K$ 在 $F$ 上的一组基）。
若 $K/F$ 可分，则由引理 \ref{fields-lemma-separable-goes-up}，
每个 $\alpha_i$ 都在 $F(\alpha_1, \ldots, \alpha_{i - 1})$
上可分；再由引理 \ref{fields-lemma-separably-generated-separable} 得到等号。
另一方面，若等号成立，则不论怎样选择
$\alpha_1, \ldots, \alpha_n$，引理
\ref{fields-lemma-separably-generated-separable} 都表明 $\alpha_1$
在 $F$ 上可分。由于可用 $K$ 的任意元素开始这一序列，
故 $K$ 在 $F$ 上可分。
\end{proof}

\begin{lemma}
\label{fields-lemma-separable-permanence}
令 $E/k$ 与 $F/E$ 为可分代数域扩张。则 $F/k$ 是可分域扩张。
\end{lemma}

\begin{proof}
取 $\alpha \in F$。则 $\alpha$ 在 $E$ 上可分且代数。
令 $P = x^d + \sum_{i < d} a_i x^i$ 为 $\alpha$ 在 $E$ 上的
最小多项式。每个 $a_i$ 都在 $k$ 上可分且代数。考虑域塔
$$
k \subset k(a_0) \subset k(a_0, a_1) \subset \ldots \subset
k(a_0, \ldots, a_{d - 1}) \subset k(a_0, \ldots, a_{d - 1}, \alpha)
$$
由于 $a_i$ 在 $k$ 上可分且代数，由引理
\ref{fields-lemma-separable-goes-up}，它在
$k(a_0, \ldots, a_{i - 1})$ 上也可分且代数。
最后，$\alpha$ 在 $k(a_0, \ldots, a_{d - 1})$ 上可分且代数，
因为它是 $P$ 的根，而该多项式不可约
（因为它在可能更大的域 $E$ 上不可约）且可分
（因为它在 $E$ 上可分）。因此，由引理
\ref{fields-lemma-separably-generated-separable}，
$k(a_0, \ldots, a_{d - 1}, \alpha)$ 在 $k$ 上可分，
从而 $\alpha$ 在 $k$ 上可分，如所需。
\end{proof}

\begin{lemma}
\label{fields-lemma-separable-elements}
令 $E/k$ 为域扩张。则 $E$ 中在 $k$ 上可分的元素构成
$E/k$ 的一个子扩张。
\end{lemma}

\begin{proof}
令 $\alpha, \beta \in E$ 在 $k$ 上可分。由引理
\ref{fields-lemma-separable-goes-up}，$\beta$ 在 $k(\alpha)$ 上可分。
由引理 \ref{fields-lemma-separably-generated-separable}
（取 $n = 2$、$\alpha_1 = \alpha$ 与 $\alpha_2 = \beta$），
$k(\alpha, \beta)$ 在 $k$ 上可分。
\end{proof}





\section{特征标的线性无关性}
\label{fields-section-independence-characters}

\noindent
以下是所需陈述。

\begin{lemma}
\label{fields-lemma-independence-characters}
令 $L$ 为域，令 $G$ 为幺半群，例如群。设
$\chi_1, \ldots, \chi_n : G \to L$ 是两两不同的幺半群同态，
其中把 $L$ 视为乘法幺半群。则
$\chi_1, \ldots, \chi_n$ 在 $L$ 上线性无关：若
$\lambda_1, \ldots, \lambda_n \in L$ 不全为零，
则对某个 $g \in G$ 有 $\sum \lambda_i\chi_i(g) \not = 0$。
\end{lemma}

\begin{proof}
若 $n = 1$，结论成立，因为当 $e \in G$ 是单位元时，
$\chi_1(e) = 1$。对 $n > 1$ 归纳证明。
设 $\lambda_1, \ldots, \lambda_n \in L$ 不全为零。
若某个 $\lambda_i = 0$，则由关于 $n$ 的归纳假设即得结论。
我们要证明对某个 $g \in G$ 有
$\sum \lambda_i\chi_i(g) \not = 0$，故除以 $-\lambda_n$ 后，
可假设 $\lambda_n = -1$。唯一可能出问题的情形是
$$
\chi_n(g) = \sum\nolimits_{i = 1, \ldots, n - 1} \lambda_i\chi_i(g)
$$
对所有 $g \in G$ 都成立。固定 $h \in G$，于是还会有
\begin{align*}
\chi_n(h)\chi_n(g) & = \chi_n(hg) \\
& = \sum\nolimits_{i = 1, \ldots, n - 1} \lambda_i\chi_i(hg) \\
& = \sum\nolimits_{i = 1, \ldots, n - 1} \lambda_i\chi_i(h) \chi_i(g)
\end{align*}
把前一关系乘以 $\chi_n(h)$ 再相减，得到
$$
0 = \sum\nolimits_{i = 1, \ldots, n - 1}
\lambda_i (\chi_n(h) - \chi_i(h)) \chi_i(g)
$$
对所有 $g \in G$ 成立。由于 $\lambda_i \not = 0$，
归纳假设表明对所有 $i$ 都有 $\chi_n(h) = \chi_i(h)$。
上面对 $h$ 的选择任意，所以对 $i \leq n - 1$ 有
$\chi_i = \chi_n$，这与特征标 $\chi_i$ 两两不同的假设矛盾。
\end{proof}

\begin{lemma}
\label{fields-lemma-sums-of-powers}
令 $L$ 为域。设 $n \geq 1$，且
$\alpha_1, \ldots, \alpha_n \in L$ 是 $L$ 中两两不同的元素。
则存在 $e \geq 0$，使
$\sum_{i = 1, \ldots, n} \alpha_i^e \not = 0$。
\end{lemma}

\begin{proof}
把特征标的线性无关性（引理
\ref{fields-lemma-independence-characters}）应用于幺半群同态
$\mathbf{Z}_{\geq 0} \to L$，$e \mapsto \alpha_i^e$。
\end{proof}

\begin{lemma}
\label{fields-lemma-independence-embeddings}
令 $K/F$ 与 $L/F$ 为域扩张，令
$\sigma_1, \ldots, \sigma_n : K \to L$ 为两两不同的
$F$-扩张态射。则 $\sigma_1, \ldots, \sigma_n$
在 $L$ 上线性无关：若 $\lambda_1, \ldots, \lambda_n \in L$
不全为零，则对某个 $\alpha \in K$ 有
$\sum \lambda_i\sigma_i(\alpha) \not = 0$。
\end{lemma}

\begin{proof}
把引理 \ref{fields-lemma-independence-characters}
应用于 $\sigma_i$ 在单位群上的限制。
\end{proof}

\begin{lemma}
\label{fields-lemma-finite-separable-tensor-alg-closed}
令 $K/F$ 与 $L/F$ 为域扩张，其中 $K/F$ 有限可分，
而 $L$ 代数闭。则映射
$$
K \otimes_F L
\longrightarrow
\prod\nolimits_{\sigma \in \Hom_F(K, L)} L,\quad
\alpha \otimes \beta \mapsto (\sigma(\alpha)\beta)_\sigma
$$
是 $L$-代数同构。
\end{lemma}

\begin{proof}
选择 $K$ 作为 $F$-向量空间的一组基
$\alpha_1, \ldots, \alpha_n$。由引理
\ref{fields-lemma-separable-equality}（再加一个略去的小论证），集合
$\Hom_F(K, L)$ 有 $n$ 个元素，记作
$\sigma_1, \ldots, \sigma_n$。特别地，两端作为
$L$-向量空间具有相同维数 $n$。因此，若该映射不是同构，
它就有核。换言之，存在不全为零的
$\mu_j \in L$，$j = 1, \ldots, n$，使
$\sum \alpha_j \otimes \mu_j$ 属于核。也就是说，
对所有 $i$ 都有 $\sum \sigma_i(\alpha_j)\mu_j = 0$。
这意味着以 $\sigma_i(\alpha_j)$ 为元素的
$n \times n$ 矩阵不可逆。因此，可找到不全为零的
$\lambda_1, \ldots, \lambda_n \in L$，使
对所有 $j$ 都有 $\sum \lambda_i\sigma_i(\alpha_j) = 0$。
现在，任意元素 $\alpha \in K$ 都可写成
$\alpha = \sum \beta_j \alpha_j$，其中 $\beta_j \in F$，
于是得到
$$
\sum \lambda_i\sigma_i(\alpha) =
\sum \lambda_i\sigma_i(\sum \beta_j \alpha_j) =
\sum \beta_j \sum \lambda_i\sigma_i(\alpha_j) = 0
$$
这与引理 \ref{fields-lemma-independence-embeddings} 矛盾。
\end{proof}







\section{纯不可分扩张}
\label{fields-section-purely-inseparable}

\noindent
纯不可分扩张与上一节定义的可分扩张相对。这类扩张只在正特征中出现。

\begin{definition}
\label{fields-definition-purely-inseparable}
令 $F$ 为特征 $p > 0$ 的域，令 $K/F$ 为扩张。
\begin{enumerate}
\item 若存在 $p$ 的某个幂 $q$，使得 $\alpha^q \in F$，
则称元素 $\alpha \in K$ 在 $F$ 上{\it 纯不可分}。
\item 扩张 $K/F$ 称为{\it 纯不可分}，当且仅当 $K$ 的每个元素
都在 $F$ 上纯不可分。
\end{enumerate}
若有域扩张 $L/M$（不对 $M$ 的特征作任何条件），则称该扩张
{\it 纯不可分}，若 $L = M$，或者 $M$ 的特征是某个素数 $p$，
且 $L/M$ 在上述意义下纯不可分。
\end{definition}

\noindent
注意，纯不可分扩张必然是代数扩张。
令 $F$ 为特征 $p > 0$ 的域。纯不可分扩张的一个例是：
给元素 $t \in F$ 添加一个它原先没有的 $p$ 次根。
具体而言，下面的引理表明 $P = x^p - t$ 不可约，因而
$$
K = F[x]/(P) = F[t^{1/p}]
$$
是域。而 $K$ 在 $F$ 上纯不可分，因为 $K$ 的每个元素
$$
a_0 + a_1t^{1/p} + \ldots + a_{p - 1}t^{(p - 1)/p},\quad a_i \in F
$$
的 $p$ 次幂等于
$$
(a_0 + a_1t^{1/p} + \ldots + a_{p - 1}t^{(p - 1)/p})^p =
a_0^p + a_1^p t + \ldots + a_{p - 1}^pt^{p - 1} \in F
$$
这种情形出现在 $\mathbf{F}_p$ 上的有理函数域
$\mathbf{F}_p(t)$ 中。

\begin{lemma}
\label{fields-lemma-take-pth-root}
令 $p$ 为素数，令 $F$ 为特征 $p$ 的域。
令 $t \in F$ 为在 $F$ 中没有 $p$ 次根的元素。
则多项式 $x^p - t$ 在 $F$ 上不可约。
\end{lemma}

\begin{proof}
假设存在分解 $x^p - t = f g$。求导得到
$f' g + f g' = 0$。由于 $f$ 与 $g$ 的次数都小于 $p$，
$f' = 0$ 与 $g' = 0$ 都不成立。此外，
$\deg(f') < \deg(f)$ 且 $\deg(g') < \deg(g)$。
由此可知 $f$ 与 $g$ 有公共因子。因此，若 $x^p - t$ 可约，
则它形如 $x^p - t = c f^n$，其中 $f$ 不可约，
$c \in F^*$ 且 $n > 1$。由于 $p$ 是素数，这迫使
$n = p$ 且 $f$ 为一次多项式，进而意味着 $x^p - t$ 在
$F$ 中有根，矛盾。
\end{proof}

\noindent
我们将看到，在特征 $p$ 中取 $p$ 次根是一种非常重要的运算。

\begin{lemma}
\label{fields-lemma-purely-inseparable-permanence}
令 $E/k$ 与 $F/E$ 为纯不可分域扩张。则 $F/k$ 是纯不可分域扩张。
\end{lemma}

\begin{proof}
设 $k$ 的特征为 $p$。取 $\alpha \in F$。则对某个
$p$ 的幂 $q$ 有 $\alpha^q \in E$。继而，对某个
$p$ 的幂 $q'$ 有 $(\alpha^q)^{q'} \in k$。
所以 $\alpha^{qq'} \in k$。
\end{proof}

\begin{lemma}
\label{fields-lemma-purely-inseparable-elements}
令 $E/k$ 为域扩张。则 $E$ 中在 $k$ 上纯不可分的元素构成
$E/k$ 的一个子扩张。
\end{lemma}

\begin{proof}
令 $p$ 为 $k$ 的特征。设 $\alpha, \beta \in E$ 在 $k$ 上
纯不可分。对某些 $p$ 的幂 $q, q'$，有
$\alpha^q \in k$ 且 $\beta^{q'} \in k$。
若 $q''$ 是 $p$ 的幂，则
$(\alpha + \beta)^{q''} = \alpha^{q''} + \beta^{q''}$。
因此，若 $q'' \geq q, q'$，则 $\alpha + \beta$
在 $k$ 上纯不可分。$\alpha$ 与 $\beta$ 的差、积与商也同理。
\end{proof}

\begin{lemma}
\label{fields-lemma-finite-purely-inseparable}
令 $E/F$ 为特征 $p > 0$ 的有限纯不可分域扩张。
则存在元素序列 $\alpha_1, \ldots, \alpha_n \in E$，
使我们得到域塔
$$
E = F(\alpha_1, \ldots, \alpha_n) \supset
F(\alpha_1, \ldots, \alpha_{n - 1}) \supset
\ldots
\supset F(\alpha_1) \supset F
$$
其中每个中间扩张的次数都是 $p$，且由添加一个 $p$ 次根得到。
具体而言，
$\alpha_i^p \in F(\alpha_1, \ldots, \alpha_{i - 1})$
是在 $F(\alpha_1, \ldots, \alpha_{i - 1})$ 中没有
$p$ 次根的元素，这里 $i = 1, \ldots, n$。
\end{lemma}

\begin{proof}
对 $E/F$ 的次数归纳。若扩张次数为 $1$，结论显然
（取 $n = 0$）。否则，选择 $\alpha \in E$ 且
$\alpha \not \in F$。设对某个 $r > 0$ 有
$\alpha^{p^r} \in F$。取最小的 $r$，并以
$\alpha^{p^{r - 1}}$ 替换 $\alpha$。于是
$\alpha \not \in F$，但 $\alpha^p \in F$。
此时 $t = \alpha^p$ 不是 $F$ 中的 $p$ 次幂
（否则 $\alpha \in F$；见引理
\ref{fields-lemma-nr-roots-unchanged} 或其证明）。
因此，$F \subset F(\alpha)$ 是次数为 $p$ 的子扩张
（引理 \ref{fields-lemma-take-pth-root}）。由归纳假设，可找到
$\alpha_1, \ldots, \alpha_n \in E$，它们生成
$E/F(\alpha)$ 并满足引理的结论。序列
$\alpha, \alpha_1, \ldots, \alpha_n$ 对扩张 $E/F$ 合乎要求。
\end{proof}

\begin{lemma}
\label{fields-lemma-separable-first}
\begin{slogan}
任意代数域扩张都唯一地分解为先可分扩张、后纯不可分扩张。
\end{slogan}
令 $E/F$ 为代数域扩张。存在唯一的子扩张 $E/E_{sep}/F$，
使 $E_{sep}/F$ 可分，而 $E/E_{sep}$ 纯不可分。
\end{lemma}

\begin{proof}
若特征为零，置 $E_{sep} = E$。假设特征为 $p > 0$。
令 $E_{sep}$ 为 $E$ 中在 $F$ 上可分的元素集合。
由引理 \ref{fields-lemma-separable-elements}，它是子扩张；
当然，$E_{sep}$ 在 $F$ 上可分。
给定 $E$ 中的 $\alpha$，存在 $p$ 的某个幂 $q$，
使 $\alpha^q$ 在 $F$ 上可分。具体地，$q$ 就是使
$\alpha$ 的最小多项式形如 $P(x^q)$ 的那个 $p$ 的幂，
其中 $P$ 可分且代数；见引理
\ref{fields-lemma-irreducible-polynomials}。所以 $E/E_{sep}$ 纯不可分。
唯一性显然。
\end{proof}

\begin{definition}
\label{fields-definition-insep-degree}
令 $E/F$ 为代数域扩张，令 $E_{sep}$ 为引理
\ref{fields-lemma-separable-first} 中得到的子扩张。
\begin{enumerate}
\item 整数 $[E_{sep} : F]$ 称为该扩张的{\it 可分次数}，
记作 $[E : F]_s$。
\item 整数 $[E : E_{sep}]$ 称为该扩张的{\it 不可分次数}，
或{\it 不可分度}，记作 $[E : F]_i$。
\end{enumerate}
\end{definition}

\noindent
当然，在特征 $0$ 时，有 $[E : F] = [E : F]_s$ 以及
$[E : F]_i = 1$。由乘法性（引理
\ref{fields-lemma-multiplicativity-degrees}），有
$$
[E : F] = [E : F]_s [E : F]_i
$$
即使这些次数中有些是无限的也成立。事实上，可分次数与不可分次数
也分别具有乘法性（见引理
\ref{fields-lemma-multiplicativity-all-degrees}）。

\begin{lemma}
\label{fields-lemma-separable-degree}
令 $K/F$ 为有限扩张，令 $\overline{F}$ 为 $F$ 的代数闭包。
则 $[K : F]_s = |\Mor_F(K, \overline{F})|$。
\end{lemma}

\begin{proof}
先证明 $K/F$ 纯不可分的情形。我们断言，此时存在唯一映射
$K \to \overline{F}$。可选择引理
\ref{fields-lemma-finite-purely-inseparable} 中的元素序列
$\alpha_1, \ldots, \alpha_n \in K$ 来看出这一点。
$\alpha_i$ 在 $F(\alpha_1, \ldots, \alpha_{i - 1})$ 上的
不可约多项式是 $x^p - \alpha_i^p$。应用引理
\ref{fields-lemma-count-embeddings-explicitly}，得到
$|\Mor_F(K, \overline{F})| = 1$。另一方面，此时
$[K : F]_s = 1$，所以等式成立。

\medskip\noindent
回到一般有限扩张 $K/F$。依照引理
\ref{fields-lemma-separable-first}，选择 $F \subset K_s \subset K$。
由引理 \ref{fields-lemma-separable-equality}，有
$|\Mor_F(K_s, \overline{F})| = [K_s : F] = [K : F]_s$。
另一方面，根据上一段的结果，每个域映射
$\sigma' : K_s \to \overline{F}$ 都唯一扩张为域映射
$\sigma : K \to \overline{F}$。换言之，
$|\Mor_F(K, \overline{F})| = |\Mor_F(K_s, \overline{F})|$，
证明完成。
\end{proof}

\begin{lemma}[乘法性]
\label{fields-lemma-multiplicativity-all-degrees}
给定代数域扩张塔 $K/E/F$，则
$$
[K : F]_s = [K : E]_s [E : F]_s
\quad\text{且}\quad
[K : F]_i = [K : E]_i [E : F]_i
$$
\end{lemma}

\begin{proof}
先证明 $K$ 在 $F$ 上有限的情形。由于通常次数具有乘法性
（引理 \ref{fields-lemma-multiplicativity-degrees}），
只需证明两个公式中的一个。由引理
\ref{fields-lemma-separable-degree}，有
$[K : F]_s = |\Mor_F(K, \overline{F})|$。
由同一引理，对任意
$\sigma \in \Mor_F(E, \overline{F})$，$\sigma$
扩张为映射 $\tau : K \to \overline{F}$ 的方式数为
$[K : E]_s$。事实上，通过
$E \cong \sigma(E) \subset \overline{F}$，
可把 $\overline{F}$ 看作 $E$ 的代数闭包。
再结合 $\sigma$ 有
$[E : F]_s = |\Mor_F(E, \overline{F})|$ 种选择，即得结论。

\medskip\noindent
扩张为无限扩张时的证明从略。
\end{proof}









\section{正规扩张}
\label{fields-section-normal}

\noindent
令 $P \in F[x]$ 为域 $F$ 上的非常数多项式。若存在
$c \in F^*$、$n \geq 1$、$\alpha_1, \ldots, \alpha_n \in F$，
使得在 $F[x]$ 中有
$$
P = c(x - \alpha_1) \ldots (x - \alpha_n),
$$
则称 $P$ {\it 在 $F$ 上完全分解为一次因子}，或简称
{\it 在 $F$ 上完全分解}。正规扩张定义如下。

\begin{definition}
\label{fields-definition-normal}
令 $E/F$ 为代数域扩张。若对所有 $\alpha \in E$，
$\alpha$ 在 $F$ 上的最小多项式 $P$ 都在 $E$ 上完全分解为
一次因子，则称 $E$ 在 $F$ 上{\it 正规}。
\end{definition}

\noindent
与可分扩张的情形一样，建立这一概念的基本性质需要一些工作。

\begin{lemma}
\label{fields-lemma-normal-goes-up}
令 $K/E/F$ 为代数域扩张塔。若 $K$ 在 $F$ 上正规，
则 $K$ 在 $E$ 上正规。
\end{lemma}

\begin{proof}
取 $\alpha \in K$。令 $P$ 为 $\alpha$ 在 $F$ 上的最小多项式，
令 $Q$ 为 $\alpha$ 在 $E$ 上的最小多项式。
则在多项式环 $E[x]$ 中，$Q$ 整除 $P$；设 $P = QR$。
所以，若 $P$ 在 $K$ 上完全分解，则 $Q$ 也如此。
\end{proof}

\begin{lemma}
\label{fields-lemma-intersect-normal}
令 $F$ 为域，令 $M/F$ 为代数扩张。设
$M/E_i/F$（$i \in I$）为子扩张，且 $E_i/F$ 正规。
则 $\bigcap E_i$ 在 $F$ 上正规。
\end{lemma}

\begin{proof}
由定义直接得到。
\end{proof}

\begin{lemma}
\label{fields-lemma-separable-first-normal}
令 $E/F$ 为正规代数域扩张。则引理
\ref{fields-lemma-separable-first} 的子扩张 $E/E_{sep}/F$ 正规。
\end{lemma}

\begin{proof}
若特征为零，则 $E_{sep} = E$，结论显然。
若特征为 $p > 0$，则 $E_{sep}$ 是 $E$ 中在 $F$ 上可分的
元素集合。若 $\alpha \in E_{sep}$ 的最小多项式为 $P$，写成
$P = c(x - \alpha)(x - \alpha_2) \ldots (x - \alpha_d)$，
其中 $\alpha_2, \ldots, \alpha_d \in E$。
由于 $P$ 是可分多项式，且 $\alpha_i$ 是 $P$ 的根，
可知 $\alpha_i \in E_{sep}$，如所需。
\end{proof}

\begin{lemma}
\label{fields-lemma-characterize-normal}
令 $E/F$ 为代数域扩张，令 $\overline{F}$ 为 $F$ 的代数闭包。
下列条件等价：
\begin{enumerate}
\item $E$ 在 $F$ 上正规；
\item 对每一对 $\sigma, \sigma' \in \Mor_F(E, \overline{F})$，
都有 $\sigma(E) = \sigma'(E)$。
\end{enumerate}
\end{lemma}

\begin{proof}
令 $\mathcal{P}$ 为 $E$ 的所有元素在 $F$ 上的全部最小多项式所成的
集合。置
$$
T =
\{\beta \in \overline{F} \mid P(\beta) = 0\text{ 对某个 }P \in \mathcal{P}\}
$$
显然，若 $E$ 在 $F$ 上正规，则对所有
$\sigma \in \Mor_F(E, \overline{F})$ 有 $\sigma(E) = T$。
因此 (1) 蕴含 (2)。

\medskip\noindent
反过来，假设 (2)。取 $\beta \in T$。
可找到相应的 $\alpha \in E$，其最小多项式
$P \in \mathcal{P}$ 使 $\beta$ 为零。由于
$F(\alpha) = F[x]/(P)$，可找到元素
$\sigma_0 \in \Mor_F(F(\alpha), \overline{F})$，
把 $\alpha$ 送到 $\beta$。由引理
\ref{fields-lemma-map-into-algebraic-closure}，可把 $\sigma_0$
扩张为某个 $\sigma \in \Mor_F(E, \overline{F})$。
于是，$\beta$ 属于所有嵌入
$\sigma : E \to \overline{F}$ 的共同像。
这说明对任意 $\sigma$ 都有 $\sigma(E) = T$。
固定一个 $\sigma$。现在令 $P \in \mathcal{P}$。
由引理 \ref{fields-lemma-algebraically-closed}，对某个 $n$ 以及
$\beta_i \in \overline{F}$，可写成
$$
P = (x - \beta_1) \ldots (x - \beta_n)
$$
注意 $\beta_i \in T$。因此，对某个 $\alpha_i \in E$，
有 $\beta_i = \sigma(\alpha_i)$。于是
$P = (x - \alpha_1) \ldots (x - \alpha_n)$ 在 $E$ 上完全分解。
证明完成。
\end{proof}

\begin{lemma}
\label{fields-lemma-normally-generated}
令 $E/F$ 为代数域扩张。若 $E$ 在 $F$ 上由
$\alpha_i \in E$（$i \in I$）生成，并且对每个 $i$，
$\alpha_i$ 在 $F$ 上的最小多项式都在 $E$ 中完全分解，
则 $E/F$ 正规。
\end{lemma}

\begin{proof}
令 $P_i$ 为 $\alpha_i$ 在 $F$ 上的最小多项式。
令 $\alpha_i = \alpha_{i, 1}, \alpha_{i, 2}, \ldots, \alpha_{i, d_i}$
为 $P_i$ 在 $E$ 上的根。给定两个 $F$ 上的嵌入
$\sigma, \sigma' : E \to \overline{F}$，有
$$
\{\sigma(\alpha_{i, 1}), \ldots, \sigma(\alpha_{i, d_i})\} =
\{\sigma'(\alpha_{i, 1}), \ldots, \sigma'(\alpha_{i, d_i})\}
$$
因为两端都等于 $P_i$ 在 $\overline{F}$ 中的根集。
元素 $\alpha_{i, j}$ 在 $F$ 上生成 $E$，故
$\sigma(E) = \sigma'(E)$。由引理
\ref{fields-lemma-characterize-normal}，$E/F$ 正规。
\end{proof}

\begin{lemma}
\label{fields-lemma-lift-maps}
令 $L/M/K$ 为代数扩张塔。
\begin{enumerate}
\item 若 $M/K$ 正规，则 $L/K$ 的任意自同构 $\tau$
诱导自同构 $\tau|_M : M \to M$。
\item 若 $L/K$ 正规，则任意 $K$-代数映射
$\sigma : M \to L$ 都可扩张为 $L$ 的自同构。
\end{enumerate}
\end{lemma}

\begin{proof}
选择 $L$ 的代数闭包 $\overline{L}$
（定理 \ref{fields-theorem-existence-algebraic-closure}）。

\medskip\noindent
令 $\tau$ 如 (1) 所述。由引理 \ref{fields-lemma-characterize-normal}，
作为 $\overline{L}$ 的子域有 $\tau(M) = M$，故
$\tau|_M : M \to M$ 是自同构。

\medskip\noindent
令 $\sigma : M \to L$ 如 (2) 所述。
由引理 \ref{fields-lemma-map-into-algebraic-closure}，可把 $\sigma$
扩张为映射 $\tau : L \to \overline{L}$，即图表
$$
\xymatrix{
L \ar[r]_\tau & \overline{L} \\
M \ar[u] \ar[ru]_\sigma & K \ar[l] \ar[u]
}
$$
交换。由引理 \ref{fields-lemma-characterize-normal}，
$\tau(L) = L$。故 $\tau : L \to L$ 是扩张 $\sigma$ 的自同构。
\end{proof}

\begin{definition}
\label{fields-definition-automorphisms}
令 $E/F$ 为域扩张。以 $\text{Aut}(E/F)$ 或
$\text{Aut}_F(E)$ 表示 $E$ 作为 $F$-扩张范畴中对象的自同构群。
$\text{Aut}(E/F)$ 的元素称为{\it $E$ 在 $F$ 上的自同构}，
或{\it $E/F$ 的自同构}。
\end{definition}

\noindent
下面用自同构刻画正规扩张。

\begin{lemma}
\label{fields-lemma-normal-and-automorphisms}
令 $E/F$ 为有限扩张。则
$$
|\text{Aut}(E/F)| \leq [E : F]_s
$$
且等号成立，当且仅当 $E$ 在 $F$ 上正规。
\end{lemma}

\begin{proof}
选择 $F$ 的代数闭包 $\overline{F}$。回忆
$[E : F]_s = |\Mor_F(E, \overline{F})|$。
取元素 $\sigma_0 \in \Mor_F(E, \overline{F})$。则映射
$$
\text{Aut}(E/F) \longrightarrow \Mor_F(E, \overline{F}),\quad
\tau \longmapsto \sigma_0 \circ \tau
$$
是单射，从而得到不等式。若等号成立，则每个
$\sigma \in \Mor_F(E, \overline{F})$ 都由 $\sigma_0$
预合成某个自同构得到。因此 $\sigma(E) = \sigma_0(E)$。
由引理 \ref{fields-lemma-characterize-normal}，$E$ 在 $F$ 上正规。

\medskip\noindent
反过来，假设 $E/F$ 正规。由引理
\ref{fields-lemma-characterize-normal}，对所有
$\sigma \in \Mor_F(E, \overline{F})$ 都有
$\sigma(E) = \sigma_0(E)$。于是置
$\tau = \sigma_0^{-1} \circ \sigma$，得到 $E$ 在 $F$ 上的
自同构。因此上面展示的映射是满射。
\end{proof}

\begin{lemma}
\label{fields-lemma-normal-embeddings-differ-by-aut}
令 $L/K$ 为代数正规域扩张，令 $E/K$ 为域扩张。
则或者不存在从 $L$ 到 $E$ 的 $K$-嵌入，或者存在一个
$\tau : L \to E$，且其他每一个都形如
$\tau \circ \sigma$，其中 $\sigma \in \text{Aut}(L/K)$。
\end{lemma}

\begin{proof}
给定 $\tau$，用 $\tau(L) \subset E$ 替换 $L$，
再应用引理 \ref{fields-lemma-lift-maps}。
\end{proof}





\section{分裂域}
% P01-E0253: frozen source label misspells "fields" as "fieds"; retained unchanged.
\label{fields-section-splitting-fieds}

\noindent
下面的引理是构造正规域扩张的一个有用工具。

\begin{lemma}
\label{fields-lemma-splitting-field}
令 $F$ 为域，令 $P \in F[x]$ 为非常数多项式。
存在最小的域扩张 $E/F$，使得 $P$ 在 $E$ 上完全分解。
此外，域扩张 $E/F$ 正规，并且在（不唯一的）同构意义下唯一。
\end{lemma}

\begin{proof}
选择一个代数闭包 $\overline{F}$。于是可在 $\overline{F}[x]$ 中写成
$P = c (x - \beta_1) \ldots (x - \beta_n)$，见
引理 \ref{fields-lemma-algebraically-closed}。注意 $c \in F^*$。置
$E = F(\beta_1, \ldots, \beta_n)$。显然，在要求
$P$ 在 $E$ 上完全分解的条件下，$E$ 是最小的。

\medskip\noindent
接下来，令 $E'$ 是 $F$ 的另一个最小域扩张，使得
$P$ 在 $E'$ 上完全分解。写成
$P = c (x - \alpha_1) \ldots (x - \alpha_n)$，其中 $c \in F$ 且
$\alpha_i \in E'$。由最小性再次得到
$E' = F(\alpha_1, \ldots, \alpha_n)$。此外，若任取
$\sigma : E' \to \overline{F}$
（引理 \ref{fields-lemma-map-into-algebraic-closure}），
便立刻看出，对某个置换
$\tau : \{1, \ldots, n\} \to \{1, \ldots, n\}$，有
$\sigma(\alpha_i) = \beta_{\tau(i)}$。
因此 $\sigma(E') = E$。由引理 \ref{fields-lemma-characterize-normal}，
这说明 $E'$ 是 $F$ 的正规扩张，并且
$E \cong E'$ 作为 $F$ 的扩张同构，从而证明完毕。
\end{proof}

\begin{definition}
\label{fields-definition-splitting-field}
令 $F$ 为域，令 $P \in F[x]$ 为非常数多项式。
引理 \ref{fields-lemma-splitting-field} 中构造的域扩张 $E/F$
称为{\it $P$ 在 $F$ 上的分裂域}。
\end{definition}

\begin{lemma}
\label{fields-lemma-normal-closure}
\begin{slogan}
有限域扩张的正规闭包存在。
\end{slogan}
令 $E/F$ 为有限域扩张。存在唯一的最小有限扩张 $K/E$，
使得 $K$ 在 $F$ 上正规。
\end{lemma}

\begin{proof}
选择 $E$ 在 $F$ 上的生成元 $\alpha_1, \ldots, \alpha_n$。
令 $P_1, \ldots, P_n$ 分别为
$\alpha_1, \ldots, \alpha_n$ 在 $F$ 上的最小多项式。
置 $P = P_1 \ldots P_n$。注意
$(x - \alpha_1) \ldots (x - \alpha_n)$ 整除 $P$，因为
每个 $(x - \alpha_i)$ 都整除 $P_i$。设
$P = (x - \alpha_1) \ldots (x - \alpha_n)Q$。
令 $K/E$ 为 $P$ 在 $E$ 上的分裂域。
我们声称 $K$ 也是 $P$ 在 $F$ 上的分裂域
（这蕴含 $K$ 在 $F$ 上正规）。
这一点是清楚的，因为 $K/E$ 由 $Q$ 在 $E$ 上的根生成，
而 $E$ 由 $(x - \alpha_1) \ldots (x - \alpha_n)$
在 $F$ 上的根生成；因此 $K$ 由 $P$ 在 $F$ 上的根生成。

\medskip\noindent
唯一性。设 $K'/E$ 是第二个最小扩张，使得
$K'/F$ 正规。选择一个代数闭包 $\overline{F}$
和一个嵌入 $\sigma_0 : E \to \overline{F}$。由
引理 \ref{fields-lemma-map-into-algebraic-closure}，
可将 $\sigma_0$ 扩张为 $\sigma : K \to \overline{F}$ 和
$\sigma' : K' \to \overline{F}$。
由引理 \ref{fields-lemma-intersect-normal}，
$\sigma(K) \cap \sigma'(K')$ 在 $F$ 上正规。
根据最小性得到 $\sigma(K) = \sigma'(K')$。
于是 $\sigma^{-1} \circ \sigma' : K' \to K$
给出 $E$ 的扩张之间的一个同构。
\end{proof}

\begin{definition}
\label{fields-definition-normal-closure}
令 $E/F$ 为有限域扩张。引理 \ref{fields-lemma-normal-closure}
中构造的域扩张 $K/E$
% P01-E0254: frozen English omits "of" in "normal closure E over F".
称为{\it $E$ 在 $F$ 上的正规闭包}。
\end{definition}

\noindent
可以在任意给定的正规扩张内部构造正规闭包。

\begin{lemma}
\label{fields-lemma-normal-closure-inside-normal}
令 $L/K$ 为代数正规扩张。
\begin{enumerate}
\item 若 $L/M/K$ 是子扩张塔且 $M/K$ 有限，则存在扩张塔
$L/M'/M/K$，其中 $M'/K$ 有限且正规。
\item 若 $L/M'/M/K$ 是扩张塔，且 $M/K$ 正规、$M'/M$ 有限，
则存在扩张塔 $L/M''/M'/M/K$，其中 $M''/M$
有限且 $M''/K$ 正规。
\end{enumerate}
\end{lemma}

\begin{proof}
(1) 的证明。令 $M'$ 为 $L/K$ 的包含 $M$ 且在 $K$ 上正规的
最小子扩张。由引理 \ref{fields-lemma-normal-closure}，
这是 $M/K$ 的正规闭包，并且在 $K$ 上有限。

\medskip\noindent
(2) 的证明。令 $\alpha_1, \ldots, \alpha_n \in M'$
在 $M$ 上生成 $M'$。令 $P_1, \ldots, P_n$ 分别为
$\alpha_1, \ldots, \alpha_n$ 在 $K$ 上的最小多项式。
令 $\alpha_{i, j}$ 为 $P_i$ 在 $L$ 中的根。令
$M''  = M(\alpha_{i, j})$。由引理
\ref{fields-lemma-normally-generated}
（应用于生成元集合 $M \cup \{\alpha_{i, j}\}$），
可得 $M''$ 在 $K$ 上正规。
\end{proof}

\noindent
下面的引理有时可用于证明正规闭包的性质。

\begin{lemma}
\label{fields-lemma-normal-closure-tensor-product}
令 $L/K$ 为有限扩张，令 $M/L$ 为 $L$ 在 $K$ 上的正规闭包。
则存在 $K$-代数的满射
$$
L \otimes_K L \otimes_K \ldots \otimes_K L \longrightarrow M
$$
其中张量因子的数目可取为
$[L : K]_s \leq [L : K]$。
\end{lemma}

\begin{proof}
选择 $K$ 的一个代数闭包 $\overline{K}$。
置 $n = [L : K]_s = |\Mor_K(L, \overline{K})|$；
等号来自引理 \ref{fields-lemma-separable-degree}。
设 $\Mor_K(L, \overline{K}) = \{\sigma_1, \ldots, \sigma_n\}$。
令 $M' \subset \overline{K}$ 为由
$\sigma_i(L)$（$i = 1, \ldots, n$）生成的 $K$-子代数。
由于 $\overline{K}$ 的任意 $K$-子代数都是域，$M'$ 是域。
从 $M'$ 到 $\overline{K}$ 的任意 $K$-代数映射都会置换
$\sigma_i$，因而把 $M'$ 映入且映满 $M'$。按构造，域 $M'$
由 $\sigma_1(L)$ 中元素的共轭元生成。由此再结合
引理 \ref{fields-lemma-characterize-normal}，可得 $M'$ 在 $K$ 上正规，
并且它是 $\overline{K}$ 中包含 $\sigma_1(L)$ 的最小正规子扩张。
由正规闭包的唯一性，$M \cong M'$。
最后，存在满射
$$
L \otimes_K L \otimes_K \ldots \otimes_K L \longrightarrow M',
\quad
\lambda_1 \otimes \ldots \otimes \lambda_n \longmapsto
\sigma_1(\lambda_1) \ldots \sigma_n(\lambda_n)
$$
并注意按定义有 $n \leq [L : K]$。
\end{proof}












\section{单位根}
\label{fields-section-roots-of-1}

\noindent
令 $F$ 为域。对整数 $n \geq 1$，置
$$
\mu_n(F) = \{\zeta \in F \mid \zeta^n = 1\}
$$
这称为{\it $n$ 次单位根群}，或称为{\it $n$ 次方根为 $1$ 的元素群}。
它在乘法下是阿贝尔群，单位元为 $1$。
注意，在域中，次数为 $d$ 的多项式至多有 $d$ 个根。
因此，由于该群由一个次数为 $n$ 的多项式方程定义，
可见 $|\mu_n(F)| \leq n$。
当然，$\mu_n(F)$ 的每个元素的阶都整除 $n$。
此外，各子群
$$
\mu_d(F) \subset \mu_n(F),\quad d | n
$$
至多有 $d$ 个元素。这蕴含 $\mu_n(F)$ 是循环群。

\begin{lemma}
\label{fields-lemma-cyclic}
令 $A$ 为指数整除 $n$ 的阿贝尔群，并且对所有 $d | n$，
$\{x \in A \mid dx = 0\}$ 的基数至多为 $d$。
则 $A$ 是阶整除 $n$ 的循环群。
\end{lemma}

\begin{proof}
这些条件蕴含 $|A| \leq n$，特别地，$A$ 是有限群。
有限阿贝尔群的结构表明
$A = \mathbf{Z}/e_1\mathbf{Z} \oplus \ldots \oplus \mathbf{Z}/e_r\mathbf{Z}$，
其中某些整数满足 $1 < e_1 | e_2 | \ldots | e_r$。
这将蕴含 $\{x \in A \mid e_1 x = 0\}$ 的基数为 $e_1^r$。
因此 $r = 1$。
\end{proof}

\noindent
将此应用于域 $\mathbf{F}_p$，便得到著名结论：
群 $(\mathbf{Z}/p\mathbf{Z})^*$ 是循环群。
有限域一节还会对此作更多说明。

\medskip\noindent
还有一个常用的观察：若 $F$ 的特征为 $p > 0$，则
$\mu_{p^n}(F) = \{1\}$。这是因为，正如我们在
引理 \ref{fields-lemma-nr-roots-unchanged} 的证明中看到的那样，
在特征为 $p$ 的域上，取 $p$ 次幂是单射。
（当然，这也直接由该引理的陈述得出。）






\section{有限域}
\label{fields-section-finite}

\noindent
令 $F$ 为有限域。显然 $F$ 具有正特征，因为不可能有单射
$\mathbf{Q} \to F$。设 $F$ 的特征为 $p$。
扩张 $\mathbf{F}_p \subset F$ 是有限的。
因此可见，$F$ 有 $q = p^f$ 个元素，其中 $f \geq 1$。

\medskip\noindent
考察单位群 $F^*$。这是有限阿贝尔群，因而具有某个指数 $e$。
于是 $F^* = \mu_e(F)$，由第 \ref{fields-section-roots-of-1} 节的讨论可知，
$F^*$ 是阶为 $q - 1$ 的循环群。（事后也可得 $e = q - 1$。）
特别地，若 $\alpha \in F^*$ 是生成元，则显然
$$
F = \mathbf{F}_p(\alpha)
$$
换言之，扩张 $F/\mathbf{F}_p$ 由单个元素生成。
当然，对有限域的任意扩张 $E/F$ 也是如此
（因为 $E$ 已经由单个元素在其素域上生成）。





\section{本原元}
\label{fields-section-primitive-element}

\noindent
令 $E/F$ 为有限域扩张。若 $\alpha \in E$ 且 $E = F(\alpha)$，
则称该元素为{\it $E$ 在 $F$ 上的本原元}。

\begin{lemma}[本原元]
\label{fields-lemma-primitive-element}
令 $E/F$ 为有限域扩张。下列条件等价：
\begin{enumerate}
\item $E$ 在 $F$ 上存在本原元；
\item 只有有限多个子扩张 $E/K/F$。
\end{enumerate}
此外，若 $E/F$ 可分，则 (1) 和 (2) 成立。
\end{lemma}

\begin{proof}
令 $\alpha \in E$ 为本原元。令 $P$ 为 $\alpha$ 在 $F$ 上的
最小多项式。接下来，令 $E/K/F$ 为一个子扩张。
令 $Q$ 为 $\alpha$ 在 $K$ 上的最小多项式。注意
$\deg(Q) = [E : K]$。写成
$Q = x^d + \sum_{i < d} a_i x^i$，我们声称
$K$ 等于 $L = F(a_0, \ldots, a_{d - 1})$。
确实，$\alpha$ 在 $L$ 上的次数为 $d$，且 $L \subset K$。
因此 $[E : L] = [E : K]$，从而 $[K : L] = 1$，即 $K = L$。
故只需证明多项式 $Q$ 的可能性至多有限。
这是清楚的，因为在 $K[x]$ 中有分解 $P = QR$，因而尤其在
$E[x]$ 中也有此分解。由于 $E[x]$ 中具有唯一分解，
$P$ 在 $E[x]$ 中至多有有限多个首一因子。

\medskip\noindent
若 $F$ 是有限域（等价地，$E$ 是有限域），则由
第 \ref{fields-section-finite} 节的讨论，$E/F$ 有本原元。
接下来，假设 $F$ 无限，并且只有至多有限多个真子域
$E/K/F$。将它们列为 $K_1, \ldots, K_N$。
于是每个 $K_i \subset E$ 都是一个真 $F$-向量子空间。
由于 $F$ 无限，可找到向量 $\alpha \in E$，使得对所有 $i$
都有 $\alpha \not \in K_i$
（一个向量空间绝不可能等于有限多个真向量子空间的并；
细节从略）。于是 $\alpha$ 是 $E$ 在 $F$ 上的本原元。

\medskip\noindent
已证明 (1) 与 (2) 等价，现在来证明引理的最后一个陈述。
选择 $F$ 的一个代数闭包 $\overline{F}$。
枚举其中的元素
$\sigma_1, \ldots, \sigma_n \in \Mor_F(E, \overline{F})$。
由于 $E/F$ 可分，由引理 \ref{fields-lemma-separable-equality}，
有 $n = [E : F]$。注意若 $i \not = j$，则
$$
V_{ij} = \Ker(\sigma_i - \sigma_j : E \longrightarrow \overline{F})
$$
不等于 $E$。因此，按上一段的论证，可找到
$\alpha \in E$，使得对所有 $i \not = j$ 都有
$\alpha \not \in V_{ij}$。从而
$|\Mor_F(F(\alpha), \overline{F})| \geq n$。
另一方面，$[F(\alpha) : F] \leq [E : F]$。
故由引理 \ref{fields-lemma-separable-equality} 得到等号，
进而得出 $E = F(\alpha)$。
\end{proof}






\section{迹与范数}
\label{fields-section-trace-pairing}

\noindent
令 $L/K$ 为有限域扩张。由引理
\ref{fields-lemma-vector-space-is-free}，
可选择一个 $K$-模同构 $L \cong K^{\oplus n}$。
当然，$n = [L : K]$ 就是该域扩张的次数。
利用这一同构，可得到一个 $K$-代数映射
$$
L \longrightarrow \text{Mat}(n \times n, K),\quad
\alpha \longmapsto \text{乘以 }\alpha\text{ 的线性变换矩阵}
$$
因此，给定 $\alpha \in L$，可取相应矩阵的迹与行列式。
这些量当然与上面所选的基无关。更典范地说，只需把
$L$ 看作有限维 $K$-向量空间，便有
$\text{Trace}_K(\alpha : L \to L)$ 以及行列式
$\det_K(\alpha : L \to L)$。

\begin{definition}
\label{fields-definition-trace-norm}
令 $L/K$ 为有限域扩张。对 $\alpha \in L$，定义其{\it 迹}
为 $\text{Trace}_{L/K}(\alpha) = \text{Trace}_K(\alpha : L \to L)$，
定义其{\it 范数}
为 $\text{Norm}_{L/K}(\alpha) = \det_K(\alpha : L \to L)$。
\end{definition}

\noindent
由定义显然可知，$\text{Trace}_{L/K}$ 是 $K$-线性的，并且对
$\alpha \in K$ 满足
$\text{Trace}_{L/K}(\alpha) = [L : K]\alpha$。
类似地，$\text{Norm}_{L/K}$ 是乘性的，并且对 $\alpha \in K$ 满足
$\text{Norm}_{L/K}(\alpha) = \alpha^{[L : K]}$。
这是练习 \ref{exercises-exercise-trace-det} 和
\ref{exercises-exercise-trace-det-rings} 中所讨论的更一般构造的特例。

\begin{lemma}
\label{fields-lemma-characteristic-vs-minimal-polynomial}
令 $L/K$ 为有限域扩张。令 $\alpha \in L$，并令 $P$
为 $\alpha$ 在 $K$ 上的最小多项式。则 $K$-线性映射
$\alpha : L \to L$ 的特征多项式等于 $P^e$，其中
$e \deg(P) = [L : K]$。
\end{lemma}

\begin{proof}
选择 $L$ 在 $K(\alpha)$ 上的一组基
$\beta_1, \ldots, \beta_e$。由引理
\ref{fields-lemma-degree-minimal-polynomial} 和
\ref{fields-lemma-multiplicativity-degrees}，
$e$ 满足 $e \deg(P) = [L : K]$。
于是可见，$L = \bigoplus K(\alpha) \beta_i$ 是分解为
$\alpha$-不变子空间的直和。因此，$\alpha : L \to L$
的特征多项式等于 $\alpha : K(\alpha) \to K(\alpha)$
的特征多项式的 $e$ 次幂。

\medskip\noindent
为完成证明，可假设 $L = K(\alpha)$。
在这种情况下，由 Cayley--Hamilton 定理，$\alpha$
是特征多项式的根。由于特征多项式与最小多项式次数相同，
故二者相等。
\end{proof}

\begin{lemma}
\label{fields-lemma-trace-and-norm-from-minimal-polynomial}
令 $L/K$ 为有限域扩张。令 $\alpha \in L$，并令
$P = x^d + a_1 x^{d - 1} + \ldots + a_d$
为 $\alpha$ 在 $K$ 上的最小多项式。则
$$
\text{Norm}_{L/K}(\alpha) = (-1)^{[L : K]} a_d^e
\quad\text{且}\quad
\text{Trace}_{L/K}(\alpha) = - e a_1
$$
其中 $e d = [L : K]$。
\end{lemma}

\begin{proof}
这立即由引理 \ref{fields-lemma-characteristic-vs-minimal-polynomial}
和定义得到。
\end{proof}

\begin{lemma}
\label{fields-lemma-trace-and-norm-linear}
令 $L/K$ 为有限域扩张。令 $V$ 为有限维 $L$-向量空间，
令 $\varphi : V \to V$ 为 $L$-线性映射。则
$$
\text{Trace}_K(\varphi : V \to V) =
\text{Trace}_{L/K}(\text{Trace}_L(\varphi : V \to V))
$$
并且
$$
\det\nolimits_K(\varphi : V \to V) =
\text{Norm}_{L/K}(\det\nolimits_L(\varphi : V \to V))
$$
\end{lemma}

\begin{proof}
选择同构 $V = L^{\oplus n}$，使得 $\varphi$ 对应于一个
$n \times n$ 矩阵。在迹的情形，公式两边关于 $\varphi$
都是可加的。因此可假设 $\varphi$ 对应于仅在 $(i, j)$ 位置
有一个非零元素的矩阵。此时直接计算表明两边相等。

\medskip\noindent
在范数的情形，若 $\varphi$ 有非零核，则两边都为零。
因此可假设 $\varphi$ 对应于 $\text{GL}_n(L)$ 的一个元素。
公式两边关于 $\varphi$ 都是乘性的。由于
$\text{GL}_n(L)$ 的每个元素都是初等矩阵的乘积，
可假设 $\varphi$ 具有下列两种形式之一：
$$
E_{12}(\lambda) =
\left(
\begin{matrix}
1 & \lambda & \ldots \\
0 & 1 & \ldots \\
\ldots & \ldots & \ldots
\end{matrix}
\right)
\quad\text{或}\quad
E_1(a) =
\left(
\begin{matrix}
a & 0 & \ldots \\
0 & 1 & \ldots \\
\ldots & \ldots & \ldots
\end{matrix}
\right)
$$
（因为如有需要，也可置换基元素）。
在两种情形下，直接计算都容易验证该公式。
\end{proof}

\begin{lemma}
\label{fields-lemma-trace-and-norm-tower}
令 $M/L/K$ 为有限域扩张塔。则
$$
\text{Trace}_{M/K} = \text{Trace}_{L/K} \circ \text{Trace}_{M/L}
\quad\text{且}\quad
\text{Norm}_{M/K} = \text{Norm}_{L/K} \circ \text{Norm}_{M/L}
$$
\end{lemma}

\begin{proof}
把 $M$ 看作 $L$ 上的向量空间，并应用引理
\ref{fields-lemma-trace-and-norm-linear}。
\end{proof}

\noindent
迹配对利用迹来定义。

\begin{definition}
\label{fields-definition-trace-pairing}
令 $L/K$ 为有限域扩张。$L/K$ 的{\it 迹配对}是对称
$K$-双线性形式
$$
Q_{L/K} : L \times L \longrightarrow K,\quad
(\alpha, \beta) \longmapsto \text{Trace}_{L/K}(\alpha\beta)
$$
\end{definition}

\noindent
事实表明，有限域扩张可分，当且仅当其迹配对非退化。

\begin{lemma}
\label{fields-lemma-separable-trace-pairing}
令 $L/K$ 为有限域扩张。下列条件等价：
\begin{enumerate}
\item $L/K$ 可分；
\item $\text{Trace}_{L/K}$ 不恒为零；
\item 迹配对 $Q_{L/K}$ 非退化。
\end{enumerate}
\end{lemma}

\begin{proof}
显然 (3) 蕴含 (2)。若 (2) 成立，则取 $\gamma \in L$，
使得 $\text{Trace}_{L/K}(\gamma) \not = 0$。
于是，对任意非零 $\alpha \in L$，有
$Q_{L/K}(\alpha, \gamma/\alpha) \not = 0$。
因此 $Q_{L/K}$ 非退化。这证明了 (2) 与 (3) 等价。

\medskip\noindent
假设 $K$ 的特征为 $p$，且 $L = K(\alpha)$，其中
$\alpha \not \in K$ 且 $\alpha^p \in K$。则
$\text{Trace}_{L/K}(1) = p = 0$。
对 $i = 1, \ldots, p - 1$，可见
$x^p - \alpha^{pi}$ 是 $\alpha^i$ 在 $K$ 上的最小多项式，
并由引理 \ref{fields-lemma-trace-and-norm-from-minimal-polynomial} 得到
$\text{Trace}_{L/K}(\alpha^i) = 0$。
因此，对这种次数为 $p$ 的纯不可分扩张，
$\text{Trace}_{L/K}$ 恒为零。

\medskip\noindent
假设 $L/K$ 不可分。则存在子域塔 $L/K'/K$，使得
$L/K'$ 是上一段所述次数为 $p$ 的纯不可分扩张，见引理
\ref{fields-lemma-separable-first} 和
\ref{fields-lemma-finite-purely-inseparable}。
因此，由引理 \ref{fields-lemma-trace-and-norm-tower}，
$\text{Trace}_{L/K}$ 恒为零。

\medskip\noindent
反过来，假设 $L/K$ 可分。对次数作归纳，将证明
$\text{Trace}_{L/K}$ 不恒为零。
因此，由引理 \ref{fields-lemma-trace-and-norm-tower}，可假设
$L/K$ 由单个元素 $\alpha$ 生成
（利用迹非零即为满射这一事实）。
我们必须证明，对某个 $e \geq 0$，
$\text{Trace}_{L/K}(\alpha^e)$ 非零。
令 $P = x^d + a_1 x^{d - 1} + \ldots + a_d$
为 $\alpha$ 在 $K$ 上的最小多项式。
于是 $P$ 也是线性映射 $\alpha : L \to L$ 的特征多项式，
见引理 \ref{fields-lemma-characteristic-vs-minimal-polynomial}。
% P01-E0255: frozen source uses undefined lowercase k in L/k; retained.
由于 $L/k$ 可分，由引理 \ref{fields-lemma-recognize-separable}，
可见 $P$ 在 $K$ 的一个代数闭包 $\overline{K}$ 中有
$d$ 个两两不同的根 $\alpha_1, \ldots, \alpha_d$。
因此，它们是 $\alpha : L \to L$ 的特征值。
由线性代数，$\alpha^e$ 的迹等于
$\alpha_1^e + \ldots + \alpha_d^e$。
于是结论由引理 \ref{fields-lemma-sums-of-powers} 得出。
\end{proof}

\noindent
令 $K$ 为域，令 $Q : V \times V \to K$ 为有限维
$K$-向量空间上的双线性形式。设 $\dim_K(V) = n$。
于是 $Q$ 定义线性映射 $Q : V \to V^*$，$v \mapsto Q(v, -)$，
其中 $V^* = \Hom_K(V, K)$ 是对偶向量空间。因而得到线性映射
$$
\det(Q) : \wedge^n(V) \longrightarrow \wedge^n(V)^*
$$
若选择基元素 $\omega \in \wedge^n(V)$，则可写成
$\det(Q)(\omega) = \lambda \omega^*$，其中 $\omega^*$
是 $\wedge^n(V)^*$ 中的对偶基元素。若把所选的
$\omega$ 改为 $c \omega$，其中 $c \in K^*$，则
$\omega^*$ 改为 $c^{-1} \omega^*$，从而
$\lambda$ 改为 $c^2 \lambda$。因此，$\lambda$ 在
$K/(K^*)^2$ 中的类是良定义的，称为
{\it $Q$ 的判别式}。展开定义可见，若
$\{v_1, \ldots, v_n\}$ 是 $V$ 在 $K$ 上的一组基，则
$$
\lambda = \det(Q(v_i, v_j)_{1 \leq i, j \leq n})
$$
注意，判别式非零，当且仅当 $Q$ 非退化。

\begin{definition}
\label{fields-definition-discriminant}
令 $L/K$ 为有限域扩张。{\it $L/K$ 的判别式}
是迹配对 $Q_{L/K}$ 的判别式。
\end{definition}

\noindent
由以上讨论和引理 \ref{fields-lemma-separable-trace-pairing}，
可见判别式非零，当且仅当 $L/K$ 可分。
对 $a \in K$，我们常说“判别式是 $a$”，
尽管更准确的说法是判别式为 $a$ 在 $K/(K^*)^2$ 中的类。

\begin{exercise}
\label{fields-exercise-quadratic-discriminant}
令 $L/K$ 为次数 $2$ 的扩张。证明恰有下列一种情形发生：
\begin{enumerate}
\item 判别式为 $0$，$K$ 的特征为 $2$，并且
$L/K$ 是通过对 $K$ 的某个元素开平方得到的纯不可分扩张；
\item 判别式为 $1$，$K$ 的特征为 $2$，并且
$L/K$ 是次数为 $2$ 的可分扩张；
\item 判别式不是平方，$K$ 的特征不为 $2$，并且
$L$ 由 $K$ 添入判别式的平方根得到。
\end{enumerate}
\end{exercise}












\section{Galois 理论}
\label{fields-section-galois-theory}

\noindent
定义如下。

\begin{definition}
\label{fields-definition-galois}
若域扩张 $E/F$ 是代数的、可分的并且正规的，
则称它为{\it Galois 扩张}。
\end{definition}

\noindent
事实表明，有限扩张是 Galois 扩张，当且仅当它具有“正确”数量的自同构。

\begin{lemma}
\label{fields-lemma-finite-Galois}
令 $E/F$ 为有限域扩张。则 $E$ 在 $F$ 上为 Galois 扩张，
当且仅当 $|\text{Aut}(E/F)| = [E : F]$。
\end{lemma}

\begin{proof}
假设 $|\text{Aut}(E/F)| = [E : F]$。由引理
\ref{fields-lemma-normal-and-automorphisms}，这蕴含
$E/F$ 可分且正规，因而是 Galois 扩张。
反过来，若 $E/F$ 可分，则 $[E : F] = [E : F]_s$；
若 $E/F$ 还正规，则引理
\ref{fields-lemma-normal-and-automorphisms} 蕴含
$|\text{Aut}(E/F)| = [E : F]$。
\end{proof}

\noindent
受上面引理的启发，按如下方式引入 Galois 群。

\begin{definition}
\label{fields-definition-galois-group}
若 $E/F$ 是 Galois 扩张，则群 $\text{Aut}(E/F)$
称为{\it Galois 群}，记为 $\text{Gal}(E/F)$。
\end{definition}

\noindent
若 $L/K$ 是无限 Galois 扩张，则应把 Galois 群视为拓扑群。
我们将在第 \ref{fields-section-infinite-galois} 节回到这一点。

\begin{lemma}
\label{fields-lemma-galois-goes-up}
令 $K/E/F$ 为代数域扩张塔。
若 $K$ 在 $F$ 上是 Galois 扩张，则 $K$ 在 $E$ 上也是 Galois 扩张。
\end{lemma}

\begin{proof}
合并引理 \ref{fields-lemma-normal-goes-up} 与
\ref{fields-lemma-separable-goes-up} 即得。
\end{proof}

\begin{lemma}
\label{fields-lemma-normal-closure-galois}
令 $L/K$ 为有限可分域扩张。令 $M$ 为 $L$ 在 $K$ 上的正规闭包
（定义 \ref{fields-definition-normal-closure}）。则 $M/K$ 是 Galois 扩张。
\end{lemma}

\begin{proof}
引理 \ref{fields-lemma-separable-first} 的子扩张塔
$M/M_{sep}/K$ 由引理 \ref{fields-lemma-separable-first-normal} 可知是正规的。
由于 $L/K$ 可分，有 $L \subset M_{sep}$。
由最小性，$M = M_{sep}$，证明完成。
\end{proof}

\noindent
令 $G$ 为作用在域 $K$ 上的群（作用由域自同构给出）。
我们常使用记号
$$
K^G = \{x \in K \mid \sigma(x) = x \ \forall \sigma \in G\}
$$
并称它为 $G$ 在 $K$ 上作用的{\it 不动域}。

\begin{lemma}
\label{fields-lemma-galois-over-fixed-field}
令 $K$ 为域，令 $G$ 为忠实作用在 $K$ 上的有限群。
则扩张 $K/K^G$ 是 Galois 扩张，且
$[K : K^G] = |G|$，此扩张的 Galois 群为 $G$。
\end{lemma}

\begin{proof}
给定 $\alpha \in K$，考察 $\alpha$ 在群作用下的轨道
$G \cdot \alpha \subset K$。考察多项式
$$
P = \prod\nolimits_{\beta \in G \cdot \alpha} (x - \beta) \in K[x]
$$
整个引理的关键在于，此多项式在 $G$ 的作用下不变，
因而系数属于 $K^G$。确切地说，对 $\tau \in G$，有
$$
P^\tau = \prod\nolimits_{\beta \in G \cdot \alpha} (x - \tau(\beta)) =
\prod\nolimits_{\beta \in G \cdot \alpha} (x - \beta) = P
$$
因为映射 $\beta \mapsto \tau(\beta)$ 是轨道
$G \cdot \alpha$ 的一个置换。因此 $P \in K^G[x]$。
又因为 $\alpha$ 是其轨道中的元素，有 $P(\alpha) = 0$，
所以扩张 $K/K^G$ 是代数的。此外，$\alpha$ 在 $K^G$ 上的
最小多项式 $Q$ 整除刚构造的多项式 $P$。
因此 $Q$ 可分（例如由引理 \ref{fields-lemma-recognize-separable}），
从而 $K/K^G$ 可分。故 $K/K^G$ 是 Galois 扩张。
为完成证明，只需证明 $[K : K^G] = |G|$，因为此时由引理
\ref{fields-lemma-finite-Galois}，$G$ 就是 Galois 群。

\medskip\noindent
选择有限多个元素 $\alpha_i \in K$（$i = 1, \ldots, n$），
使得对 $i = 1, \ldots, n$ 有
$\sigma(\alpha_i) = \alpha_i$ 即蕴含
$\sigma$ 是 $G$ 的单位元。置
$$
L = K^G(\{\sigma(\alpha_i); 1 \leq i \leq n, \sigma \in G\}) \subset K
$$
并注意，$G$ 在 $K$ 上的作用诱导 $G$ 在 $L$ 上的作用。
我们将证明 $L$ 在 $K^G$ 上的次数为 $|G|$。
这将完成证明，因为若 $L \subset K$ 是真包含，
则可将一个元素 $\alpha \in K$（$\alpha \not \in L$）
加入元素列表 $\alpha_1, \ldots, \alpha_n$，却不增大 $L$，
这显然矛盾。于是问题归结为 $K/K^G$ 有限的情形，
下一段将处理这一情形。

\medskip\noindent
假设 $K/K^G$ 有限。由引理 \ref{fields-lemma-primitive-element}，
可找到 $\alpha \in K$，使得 $K = K^G(\alpha)$。
由本证明第一段的构造可见，$\alpha$ 在 $K$ 上的次数至多为
$|G|$。
% P01-E0256: frozen source says "over K"; the intended base appears to be K^G.
另一方面，该次数不可能小于 $|G|$，因为按构造，$G$
忠实作用在 $K^G(\alpha) = L$ 上，再应用引理
\ref{fields-lemma-normal-and-automorphisms} 的不等式即可。
\end{proof}

\begin{theorem}[Galois 理论基本定理]
\label{fields-theorem-galois-theory}
令 $L/K$ 为有限 Galois 扩张，其 Galois 群为 $G$。
则 $K = L^G$，且映射
$$
\{G\text{ 的子群}\}
\longrightarrow
\{\text{中间扩张 }L/M/K\},\quad
H \longmapsto L^H
$$
是双射，其逆映射把 $M$ 送到 $\text{Gal}(L/M)$。
$G$ 的正规子群 $H$ 恰好对应于满足 $M/K$ 为 Galois 扩张的子扩张
$M$。
\end{theorem}

\begin{proof}
由引理 \ref{fields-lemma-galois-goes-up}，给定子扩张 $L/M/K$，
扩张 $L/M$ 是 Galois 的。当然，$L/M$ 也是有限的
（引理 \ref{fields-lemma-finite-goes-up}）。因此，由引理
\ref{fields-lemma-finite-Galois}，
$|\text{Gal}(L/M)| = [L : M]$。
反过来，若 $H \subset G$ 为有限子群，则由引理
\ref{fields-lemma-galois-over-fixed-field}，
$[L : L^H] = |H|$。由这两个观察即可形式地得到定理中的双射对应。

\medskip\noindent
若 $H \subset G$ 正规，则 $L^H$ 在 $G$ 的作用下不变，
从而得到典范映射 $G/H \to \text{Aut}(L^H/K)$。
由于 $\text{Gal}(L/L^H) = H$，此映射必为单射。
因此 $|G/H| = [L^H : K]$，由引理
\ref{fields-lemma-finite-Galois}，$L^H$ 是 Galois 扩张。

\medskip\noindent
反过来，假设 $K \subset M \subset L$ 且 $M/K$ 是 Galois 扩张。
由引理 \ref{fields-lemma-lift-maps}，每个元素
$\tau \in \text{Gal}(L/K)$ 都诱导元素
$\tau|_M \in \text{Gal}(M/K)$。这诱导 Galois 群的同态
$\text{Gal}(L/K) \to \text{Gal}(M/K)$，其核为 $H$。
因此 $H$ 是正规子群。
\end{proof}

\begin{lemma}
\label{fields-lemma-ses-galois}
令 $L/M/K$ 为域塔。假设 $L/K$ 与 $M/K$ 都是有限 Galois 扩张。
则得到有限群的短正合列
$$
1 \to \text{Gal}(L/M) \to \text{Gal}(L/K) \to \text{Gal}(M/K) \to 1
$$
\end{lemma}

\begin{proof}
确实，由引理 \ref{fields-lemma-lift-maps}，每个元素
$\tau \in \text{Gal}(L/K)$ 都诱导元素
$\tau|_M \in \text{Gal}(M/K)$，这给出右侧的同态。
左侧映射把左边的群等同于右侧箭头的核。
该序列正合，因为由有限扩张塔中次数的乘法性，
各群的阶恰好相符（引理 \ref{fields-lemma-multiplicativity-degrees}）。
也可直接使用引理 \ref{fields-lemma-lift-maps}
看出右侧映射是满射。
\end{proof}






\section{无限 Galois 理论}
\label{fields-section-infinite-galois}

\noindent
Galois 群带有典范拓扑。

\begin{lemma}
\label{fields-lemma-galois-profinite}
令 $E/F$ 为 Galois 扩张。在 $\text{Gal}(E/F)$ 上赋予使得
$$
\text{Gal}(E/F) \times E \longrightarrow E
$$
连续的最粗拓扑，其中 $E$ 赋予离散拓扑。则：
\begin{enumerate}
\item 对任意拓扑空间 $X$ 以及映射 $X \to \text{Gal}(E/F)$，
若作用 $X \times E \to E$ 连续，则诱导映射
$X \to \text{Gal}(E/F)$ 连续；
\item 此拓扑使 $\text{Gal}(E/F)$ 成为有限型拓扑群。
\end{enumerate}
\end{lemma}

\begin{proof}
在整个证明中，把 $E$ 视为离散拓扑空间。
回忆，自映射集合 $\text{Map}(E, E)$ 上的紧开拓扑是使作用
$\text{Map}(E, E) \times E \to E$ 连续的泛性拓扑。
精确陈述见拓扑学，例 \ref{topology-example-automorphisms-of-a-set}。
引理中 $\text{Gal}(E/F)$ 上的拓扑，是由单射
$\text{Gal}(E/F) \to \text{Map}(E, E)$ 诱导的拓扑。
因此，泛性质 (1) 由紧开拓扑相应的泛性质得出。
由于可逆自映射集合 $\text{Aut}(E)$ 在紧开拓扑下构成拓扑群，
见拓扑学，例 \ref{topology-example-automorphisms-of-a-set}，
并且 $\text{Gal}(E/F) = \text{Aut}(E/F) \to \text{Map}(E, E)$
经由 $\text{Aut}(E)$ 分解，所以得到一个拓扑群。
换言之，我们利用单射
$$
\text{Gal}(E/F) \subset \text{Aut}(E)
$$
在 $\text{Gal}(E/F)$ 上赋予诱导的拓扑群结构
（见拓扑学，第 \ref{topology-section-topological-groups} 节）。
按构造，这是使作用
$\text{Gal}(E/F) \times E \to E$ 连续的最粗拓扑群结构。

\medskip\noindent
为证明 $\text{Gal}(E/F)$ 是有限型群，作如下论证
（由于我们在证明 Galois 群是有限群逆极限之前就定义了其拓扑，
这一论证必然不是通常的论证）。
由拓扑学，引理 \ref{topology-lemma-profinite-group}，
只需证明 $\text{Gal}(E/F)$ 的底拓扑空间是有限型空间。
对任意子集 $S \subset E$，考虑集合
$$
G(S) = \{ f : S \to E \mid
\begin{matrix}
\text{对所有 }\alpha \in S,\ f(\alpha)\text{ 是}\\
\alpha\text{ 在 }F\text{ 上的极小多项式的根}
\end{matrix}
\}
$$
由于多项式只有有限多个根，对每个有限的 $S \subset E$，
$G(S)$ 都是有限集。若 $S \subset S'$，则限制给出映射
$G(S') \to G(S)$。还要注意，若
$\alpha \in S \cap F$ 且 $f \in G(S)$，则
$f(\alpha) = \alpha$，因为此时最小多项式是线性的。
考虑有限型拓扑空间
$$
G = \lim_{S \subset E\text{ 为有限集}} G(S)
$$
以及典范映射
$$
c : \text{Gal}(E/F) \longrightarrow G,\quad
\sigma \longmapsto (\sigma|_S : S \to E)_S
$$
此映射是单射。展开定义可见，上面定义的
$\text{Gal}(E/F)$ 上的拓扑正是从 $G$ 诱导的拓扑。
元素 $(f_S) \in G$ 属于 $c$ 的像，当且仅当在有意义时满足
(A) $f_S(\alpha) + f_S(\beta) = f_S(\alpha + \beta)$ 和
(M) $f_S(\alpha)f_S(\beta) = f_S(\alpha\beta)$
（即 $\alpha, \beta, \alpha + \beta, \alpha\beta \in S$）。
确切地说，这意味着
$\lim f_S : E \to E$ 是 $F$-代数映射，因而由引理
\ref{fields-lemma-algebraic-extension-self-map} 可知是自同构。
对给定三元组 $(S, \alpha, \beta)$，条件 (A) 与 (M)
定义 $G$ 的闭子集。因此 $\text{Gal}(E/F)$ 同胚于一个有限型空间的
闭子集，从而其本身也是有限型空间。
\end{proof}

\begin{lemma}
\label{fields-lemma-galois-infinite}
令 $L/M/K$ 为域塔。假设 $L/K$ 与 $M/K$ 都是 Galois 扩张。
则存在典范的连续满同态
$c : \text{Gal}(L/K) \to \text{Gal}(M/K)$。
\end{lemma}

\begin{proof}
由引理 \ref{fields-lemma-lift-maps}，给定
$\text{Gal}(L/K)$ 中的 $\tau : L \to L$，其限制
$\tau|_M : M \to M$ 是 $\text{Gal}(M/K)$ 的元素。
这定义同态 $c$。连续性来自该拓扑的泛性质：作用
$$
\text{Gal}(L/K) \times M \longrightarrow M,\quad
(\tau, x) \longmapsto \tau(x) = c(\tau)(x)
$$
连续，因为 $M \subset L$，且作用
$\text{Gal}(L/K) \times L \to L$ 连续。
于是由引理 \ref{fields-lemma-galois-profinite} 的 (1) 得到
$c$ 连续。引理 \ref{fields-lemma-lift-maps} 还说明此映射是满射。
\end{proof}

\noindent
下面给出理解无限 Galois 扩张之 Galois 群的一种更标准方式。

\begin{lemma}
\label{fields-lemma-infinite-galois-limit}
令 $L/K$ 为 Galois 扩张，其 Galois 群为 $G$。
令 $\Lambda$ 为有限 Galois 子扩张的集合，即
$\lambda \in \Lambda$ 对应于 $L/L_\lambda/K$，
其中 $L_\lambda/K$ 是有限 Galois 扩张，Galois 群为
$G_\lambda$。在 $\Lambda$ 上按如下规则定义偏序：
$\lambda \geq \lambda'$ 当且仅当
$L_\lambda \supset L_{\lambda'}$。则：
\begin{enumerate}
\item $\Lambda$ 是有向偏序集；
\item $L_\lambda$ 是 $\Lambda$ 上的 $K$-扩张系统，并且
$L = \colim L_\lambda$；
\item $G_\lambda$ 是 $\Lambda$ 上有限群的逆系统，转移映射均为满射，
并且
$$
G = \lim_{\lambda \in \Lambda} G_\lambda
$$
作为有限型群成立；
\item 每个投影 $G \to G_\lambda$ 都连续且为满射。
\end{enumerate}
\end{lemma}

\begin{proof}
$L$ 的每个包含 $K$ 的子域在 $K$ 上都可分
（这直接由定义得出）。令 $S \subset L$ 为有限子集。
则 $K(S)/K$ 有限，并且存在扩张塔
$L/E/K(S)/K$，使得 $E/K$ 是有限 Galois 扩张，见引理
\ref{fields-lemma-normal-closure-inside-normal}。
因此，对某个 $\lambda \in \Lambda$，有 $E = L_\lambda$。
这当然说明集合 $\Lambda$ 非空。
此外，给定 $\lambda_1, \lambda_2 \in \Lambda$，由引理
\ref{fields-lemma-finite-finitely-generated}，可对有限集
$S_1, S_2 \subset L$ 写成
$L_{\lambda_i} = K(S_i)$。
于是存在 $\lambda \in \Lambda$，使得
$K(S_1 \cup S_2) \subset L_\lambda$。因此
$\lambda \geq \lambda_1, \lambda_2$，且
$\Lambda$ 有向（范畴论，定义
\ref{categories-definition-directed-system}）。
最后，由于 $L$ 中每个元素都包含在某个
$\lambda \in \Lambda$ 对应的 $L_\lambda$ 中，
由范畴论，第 \ref{categories-section-directed-colimits} 节
对滤过余极限的描述，得到 $\colim L_\lambda = L$。

\medskip\noindent
若 $\Lambda$ 中 $\lambda \geq \lambda'$，则由引理
\ref{fields-lemma-ses-galois} 得到典范满射
$G_\lambda \to G_{\lambda'}$，
$\sigma \mapsto \sigma|_{L_{\lambda'}}$。
因此得到有限群的逆系统，其转移映射均为满射。

\medskip\noindent
回忆 $G = \text{Aut}(L/K)$。由引理
\ref{fields-lemma-galois-infinite}，对 $\sigma \in G$，
其在 $L_\lambda$ 上的限制
$\sigma|_{L_\lambda}$ 是 $G_\lambda$ 的元素。
此外，此过程给出连续满射 $G \to G_\lambda$。
由于 $G_\lambda$ 逆系统中的转移映射也由限制给出，
显然得到典范连续映射
$$
G \longrightarrow \lim_{\lambda \in \Lambda} G_\lambda
$$
连续性来自拓扑群范畴中极限的定义；回忆，由拓扑学，引理
\ref{topology-lemma-topological-group-limits}，
这些极限与到集合范畴和拓扑空间范畴的遗忘函子交换。
另一方面，由于 $L = \colim L_\lambda$，显然逆极限的任一元素
（视作集合意义下的元素）都定义 $L$ 的一个自同构。
因此该映射为双射。两边的拓扑都是有限型的，而有限型空间之间的
连续双射是同胚（拓扑学，引理
\ref{topology-lemma-bijective-map}），故证明完成。
\end{proof}

\begin{theorem}[无限 Galois 理论基本定理]
% P01-E0257: frozen source label misspells "infinite" as "inifinite"; retained.
\label{fields-theorem-inifinite-galois-theory}
令 $L/K$ 为 Galois 扩张。令 $G = \text{Gal}(L/K)$
为视作有限型拓扑群的 Galois 群
（引理 \ref{fields-lemma-galois-profinite}）。则 $K = L^G$，且映射
$$
\{G\text{ 的闭子群}\}
\longrightarrow
\{\text{中间扩张 }L/M/K\},\quad
H \longmapsto L^H
$$
是双射，其逆映射把 $M$ 送到 $\text{Gal}(L/M)$。
有限子扩张 $M$ 恰好对应于开子群 $H \subset G$。
$G$ 的正规闭子群 $H$ 恰好对应于在 $K$ 上为 Galois 扩张的子扩张
$M$。
\end{theorem}

\begin{proof}
以下将不再特别说明地使用有限 Galois 理论的结论
（定理 \ref{fields-theorem-galois-theory}）。
令 $S \subset L$ 为有限子集。存在扩张塔
$L/E/K$，使得 $K(S) \subset E$ 且
$E/K$ 是有限 Galois 扩张，见引理
\ref{fields-lemma-normal-closure-inside-normal}。
换言之，$L/K$ 是其有限 Galois 子扩张之并。
对这样的 $E$，由引理 \ref{fields-lemma-galois-infinite}，
映射 $\text{Gal}(L/K) \to \text{Gal}(E/K)$ 连续且为满射；
即其核是开的，因为 $\text{Gal}(E/K)$ 上的拓扑是离散的。
特别地，$L \setminus K$ 中没有元素被
$\text{Gal}(L/K)$ 固定，因为
$E^{\text{Gal}(E/K)} = K$。这证明 $L^G = K$。

\medskip\noindent
由引理 \ref{fields-lemma-galois-goes-up}，给定子扩张 $L/M/K$，
扩张 $L/M$ 是 Galois 扩张。从 $G$ 上拓扑的定义立即可知，
子群 $\text{Gal}(L/M)$ 是闭的。
将以上结论应用于 $L/M$，得到
$L^{\text{Gal}(L/M)} = M$。

\medskip\noindent
反过来，令 $H \subset G$ 为闭子群。我们声称
$H = \text{Gal}(L/L^H)$。包含关系
$H \subset \text{Gal}(L/L^H)$ 是显然的。
假设 $g \in \text{Gal}(L/L^H)$。令 $S \subset L$ 为有限子集。
我们将证明 $g$ 的开邻域
$U_S(g) = \{g' \in G \mid g'(s) = g(s)\}$ 与 $H$ 相交。
由于 $H$ 闭，这蕴含 $g \in H$。
令 $L/E/K$ 为本证明第一段中那样的包含 $K(S)$ 的有限 Galois 子扩张，
并考虑同态
$c : \text{Gal}(L/K) \to \text{Gal}(E/K)$。
则 $L^H \cap E = E^{c(H)}$。由于 $g$ 固定 $L^H$，
它也固定 $E^{c(H)}$，故由有限 Galois 理论，
$c(g) \in c(H)$。取 $h \in H$ 使得
$c(h) = c(g)$。于是如所需，$h \in U_S(g)$。

\medskip\noindent
至此，我们已经建立闭子群与子扩张之间的对应。

\medskip\noindent
假设 $H \subset G$ 是开的。按上面的论证可知，对某个充分大的
有限 Galois 子扩张 $E$，$H$ 包含
$\text{Gal}(L/E)$，并且 $L^H$ 包含在 $E$ 中，
因而在 $K$ 上有限。反过来，若 $M$ 是有限子扩张，
则 $M$ 由某个有限子集 $S$ 生成，相应子群为开集
$U_S(e)$，其中 $e \in G$ 是单位元。

\medskip\noindent
假设 $K \subset M \subset L$ 且 $M/K$ 是 Galois 扩张。
由引理 \ref{fields-lemma-galois-infinite}，存在 Galois 群之间的连续满同态
$\text{Gal}(L/K) \to \text{Gal}(M/K)$，
其核为 $\text{Gal}(L/M)$。因此
$\text{Gal}(L/M)$ 是正规闭子群。

\medskip\noindent
最后，假设 $N \subset G$ 正规且闭。对本证明第一段中的任意
$L/E/K$，像 $c(N) \subset \text{Gal}(E/K)$ 是正规子群。
因此 $L^N = \bigcup E^{c(N)}$ 是 $K$ 的 Galois 扩张之并
（由有限 Galois 理论），从而在 $K$ 上为 Galois 扩张。
\end{proof}

\begin{lemma}
\label{fields-lemma-ses-infinite-galois}
令 $L/M/K$ 为域塔。假设 $L/K$ 与 $M/K$ 都是 Galois 扩张。
则得到有限型拓扑群的短正合列
$$
1 \to \text{Gal}(L/M) \to \text{Gal}(L/K) \to \text{Gal}(M/K) \to 1
$$
\end{lemma}

\begin{proof}
这是引理 \ref{fields-lemma-galois-infinite} 的重新表述。
\end{proof}






\section{复数}
\label{fields-section-complex-numbers}

\noindent
代数学基本定理断言，复数域是代数闭域。本节对此作简要讨论。

\medskip\noindent
首先要说明，证明这一点需要用到少量微积分知识。
我们将使用如下直观事实：实数域上的每个奇数次多项式都有实根。
具体地，对某个 $k \geq 0$ 和 $a_{2k + 1} \not = 0$，令
$P(x) = a_{2k + 1} x^{2k + 1} + \ldots + a_0 \in \mathbf{R}[x]$。
不妨假设 $a_{2k + 1} > 0$。当 $x \in \mathbf{R}$
为充分大的正数时，由于项
$a_{2k + 1} x^{2k + 1}$ 支配所有其他项，有 $P(x) > 0$。
类似地，若 $x \ll 0$，则由于同一原因有 $P(x) < 0$
（这里用到了次数为奇数）。因此，由介值定理，存在
$x \in \mathbf{R}$ 使得 $P(x) = 0$。

\medskip\noindent
由上面可以推出，$\mathbf{R}$ 没有非平凡的奇数次域扩张，
因为这种扩张中的元素将具有奇数次最小多项式。

\medskip\noindent
接下来，令 $K/\mathbf{R}$ 为有限 Galois 扩张，Galois 群为 $G$。
令 $P \subset G$ 为一个 $2$-Sylow 子群。则
$K^P/\mathbf{R}$ 是奇数次扩张，故由上面可知
$K^P = \mathbf{R}$，这又蕴含 $G = P$。
（这些论证当然都依赖于 Galois 理论。）
因此 $G$ 是 $2$-群。若 $G$ 非平凡，则
$\mathbf{C} \subset K$，因为 $\mathbf{C}$
（在同构意义下）是 $\mathbf{R}$ 唯一的次数 $2$ 扩张。
若 $G$ 有多于 $2$ 个元素，就会得到 $\mathbf{C}$ 的二次扩张。
这是荒谬的，因为每个复数都有平方根。

\medskip\noindent
结论是：$\mathbf{C}$ 代数闭。事实上，若不然，就会得到
一个非平凡有限扩张 $K/\mathbf{C}$；由引理
\ref{fields-lemma-normal-closure}，可假设它在 $\mathbf{R}$ 上正规
（因而是 Galois 的）。但上面已经说明，此时
$K = \mathbf{C}$。

\begin{lemma}[代数学基本定理]
\label{fields-lemma-C-algebraically-closed}
域 $\mathbf{C}$ 是代数闭域。
\end{lemma}

\begin{proof}
见上面的讨论。
\end{proof}





\section{Kummer 扩张}
\label{fields-section-Kummer}

\noindent
令 $K$ 为域。令 $n \geq 2$ 为整数，并假设 $K$ 包含一个本原
$n$ 次方根为 $1$ 的元素。令 $a \in K$。令 $L$ 为通过添入方程
$x^n = a$ 的一个根 $b$ 所得的 $K$ 的扩张。
则 $L/K$ 是 Galois 扩张。若
$G = \text{Gal}(L/K)$ 为 Galois 群，则映射
$$
G \longrightarrow \mu_n(K),\quad \sigma \longmapsto \sigma(b)/b
$$
是群的单同态。特别地，$G$ 作为循环群 $\mu_n(K)$ 的子群，
是阶整除 $n$ 的循环群。Kummer 理论给出其逆命题。

\begin{lemma}[Kummer 扩张]
\label{fields-lemma-Kummer}
令 $L/K$ 为域的 Galois 扩张，其 Galois 群为
$\mathbf{Z}/n\mathbf{Z}$。进一步假设 $K$ 的特征与 $n$ 互素，
并且 $K$ 包含一个本原 $n$ 次方根为 $1$ 的元素。
则 $L = K[z]$，其中 $z^n \in K$。
\end{lemma}

\begin{proof}
令 $\zeta \in K$ 为一个本原 $n$ 次方根为 $1$ 的元素。
令 $\sigma$ 为 $\text{Gal}(L/K)$ 的一个生成元。
把 $\sigma : L \to L$ 看作 $K$-线性算子。
注意，作为线性算子，有 $\sigma^n - 1 = 0$。
应用特征标的线性无关性
（引理 \ref{fields-lemma-independence-characters}），可知不存在
次数 $< n$ 的 $K$ 上多项式零化 $\sigma$。
因此，$\sigma$ 作为线性算子的最小多项式为 $x^n - 1$。
由于 $\zeta$ 是 $x^n - 1$ 的根，线性代数表明存在
$0 \neq z \in L$，使得 $\sigma(z) = \zeta z$。
此 $z$ 满足 $z^n \in K$，因为
$\sigma(z^n) = (\zeta z)^n = z^n$。此外，
$z, \sigma(z), \ldots, \sigma^{n - 1}(z) =
z, \zeta z, \ldots \zeta^{n - 1} z$ 两两不同，
这保证 $z$ 在 $K$ 上生成 $L$。
因此如所需，$L = K[z]$。
\end{proof}

\begin{lemma}
\label{fields-lemma-adjoint-pth-root-unity}
令 $K$ 为域，代数闭包为 $\overline{K}$。
令 $p$ 为不同于 $K$ 的特征的素数。
令 $\zeta \in \overline{K}$ 为一个本原 $p$ 次方根为 $1$ 的元素。
则 $K(\zeta)/K$ 是次数整除 $p - 1$ 的 Galois 扩张。
\end{lemma}

\begin{proof}
多项式 $x^p - 1$ 在 $K(\zeta)$ 上完全分解，因为它的根为
$1, \zeta, \zeta^2, \ldots, \zeta^{p - 1}$。
因此 $K(\zeta)/K$ 是分裂域，从而正规。
由于 $x^p - 1$ 是可分多项式，该扩张可分。
所以该扩张是 Galois 扩张。$K(\zeta)$ 在 $K$ 上的任意自同构
都把 $\zeta$ 送到某个 $\zeta^i$，其中
$1 \leq i \leq p - 1$。因此 Galois 群是
$(\mathbf{Z}/p\mathbf{Z})^*$ 的子群。
\end{proof}

\begin{lemma}
\label{fields-lemma-subfields-kummer}
\begin{reference}
\cite[Theorem 5.2]{Radical}
\end{reference}
令 $K$ 为域。令 $L/K$ 为次数 $e$ 的有限扩张，并假设它由元素
$\alpha$ 生成，其中 $a = \alpha^e \in K$。
若 $L$ 中每个 $e$ 次单位根都包含在 $K$ 中，则任意子扩张
% P01-E0258: frozen source writes "sub extension" as two words.
$L/L'/K$ 都由某个 $\alpha^d$ 生成，其中 $d | e$。
\end{lemma}

\begin{proof}
注意，对 $d | e$，子域 $K(\alpha^d)$ 满足
$[K(\alpha^d) : K] = e/d$ 以及
$[L : K(\alpha^d)] = d$。
令 $L/L'/K$ 为子扩张。设 $d = [L : L']$。
若 $\alpha^d \in L'$，则由次数可知
$L' = K(\alpha^d)$。
令 $P \in L'[x]$ 为 $\alpha$ 在 $L'$ 上的最小多项式。
则 $P$ 整除 $x^e - a$，且 $P$ 的次数为 $d$。
在 $x^e - a$ 在 $L$ 上的一个分裂域中写成
$$
x^e - a = \prod\nolimits_{i = 1, \ldots, e} (x - \zeta_i \alpha)
$$
$\zeta_i$ 是 $e$ 次单位根，重新编号后有
$$
P = \prod\nolimits_{i = 1, \ldots, d} (x - \zeta_i \alpha)
$$
$P$ 的常数项等于
% P01-E0259: frozen formula omits the factor (-1)^d in the constant term.
$$
c = (\prod\nolimits_{i = 1, \ldots, d} \zeta_i) \alpha^d
$$
并且属于 $L' \subset L$。由于 $\alpha \in L$，这蕴含
$\zeta = \prod_{i = 1, \ldots, d} \zeta_i$ 属于 $L$，
从而由假设属于 $K$。因此
$\alpha^d = \zeta^{-1}c \in L'$，结论得证。
\end{proof}





\section{Artin--Schreier 扩张}
\label{fields-section-Artin-Schreier}

\noindent
令 $K$ 为特征 $p > 0$ 的域。令 $a \in K$。
令 $L$ 为通过添入方程 $x^p - x = a$ 的一个根 $b$
所得的 $K$ 的扩张。则 $L/K$ 是 Galois 扩张。
若 $G = \text{Gal}(L/K)$ 为 Galois 群，则映射
$$
G \longrightarrow \mathbf{Z}/p\mathbf{Z},\quad
\sigma \longmapsto \sigma(b) - b
$$
是群的单同态。特别地，$G$ 作为
$\mathbf{Z}/p\mathbf{Z}$ 的子群，是阶整除 $p$ 的循环群。
Artin--Schreier 扩张理论给出其逆命题。

\begin{lemma}[Artin--Schreier 扩张]
\label{fields-lemma-Artin-Schreier}
令 $L/K$ 为特征 $p > 0$ 的域的 Galois 扩张，
其 Galois 群为 $\mathbf{Z}/p\mathbf{Z}$。
则 $L = K[z]$，其中 $z^p - z \in K$。
\end{lemma}

\begin{proof}
令 $\sigma$ 为 $\text{Gal}(L/K)$ 的一个生成元。
把 $\sigma : L \to L$ 看作 $K$-线性算子。
注意，作为线性算子，有 $\sigma^p - 1 = 0$。
应用特征标的线性无关性
（引理 \ref{fields-lemma-independence-characters}），不存在次数
$< p$ 的多项式零化 $\sigma$。
因此，$\sigma$ 的最小多项式为
$x^p - 1 = (x - 1)^p$。
这蕴含存在 $w \in L$，使得
$(\sigma - 1)^{p - 1}(w) = y$ 非零。
于是 $\sigma(y) = y$，即 $y \in K$。因此
$z = y^{-1}(\sigma - 1)^{p - 2}(w)$ 满足
$\sigma(z) = z + 1$。由于 $z \not \in K$，有 $L = K[z]$。
又因为
$\sigma(z^p - z) = (z + 1)^p - (z + 1) = z^p - z$，
可见 $z^p - z \in K$，证明完成。
\end{proof}








\section{超越性}
\label{fields-section-transcendence}

\noindent
回忆标准定义。

\begin{definition}
\label{fields-definition-transcendence}
令 $K/k$ 为域扩张。
\begin{enumerate}
\item 若把 $X_i$ 送到 $x_i$ 的映射
$$
k[X_i; i\in I] \longrightarrow K
$$
是单射，则称 $K$ 中的元素族 $\{x_i\}_{i \in I}$
在 $k$ 上{\it 代数无关}。
\item 多项式环 $k[x_i; i \in I]$ 的分式域记为
$k(x_i; i\in I)$。
\item $k$ 的{\it 纯超越扩张}，是与 $k$ 上某个多项式环的
分式域同构的任意域扩张 $K/k$。
\item $K/k$ 的{\it 超越基}，是一个在 $k$ 上代数无关的元素族
$\{x_i\}_{i \in I}$，并且扩张
$K/k(x_i; i\in I)$ 是代数的。
\end{enumerate}
\end{definition}

\begin{example}
\label{fields-example-pi-transcendental}
域 $\mathbf{Q}(\pi)$ 是纯超越的，因为 $\pi$
不是任何有理系数非零多项式的根。特别地，
$\mathbf{Q}(\pi) \cong \mathbf{Q}(x)$。
\end{example}

\begin{lemma}
\label{fields-lemma-transcendence-degree}
令 $E/F$ 为域扩张。$E$ 在 $F$ 上存在超越基。
任意两个超越基具有相同基数。
\end{lemma}

\begin{proof}
令 $A$ 为 $E$ 的代数无关子集。令 $G$ 为 $E$ 的一个包含
$A$ 且生成 $E/F$ 的子集。我们声称，可以找到超越基 $B$，
使得 $A \subset B \subset G$。
为证明这一点，考虑代数无关子集组成的族 $\mathcal{B}$；
其成员是 $G$ 中包含 $A$ 的子集。用包含关系在
$\mathcal{B}$ 上定义偏序。于是 $\mathcal{B}$
至少包含元素 $A$。$\mathcal{B}$ 的全序子集 $T$
中各元素的并，在 $F$ 上是 $E$ 的代数无关子集，
因为任何代数依赖关系都会已经出现在 $T$ 的某个元素中
（多项式只涉及有限多个变量）。该并也包含 $A$ 且包含在 $G$ 中。
由 Zorn 引理，存在极大元素 $B \in \mathcal{B}$。
现在声称 $E$ 在 $F(B)$ 上是代数的。这是因为，若不然，
由于 $F(G) = E$，就会存在元素
$f \in G$ 在 $F(B)$ 上超越。于是 $B \cup\{f\}$
将代数无关，与 $B$ 的极大性矛盾。因此 $B$ 是所需超越基。

\medskip\noindent
令 $B$ 与 $B'$ 为两个超越基。不失一般性，可假设
$|B'| \leq |B|$。现在把证明分成两种情形。
第一种情形是 $B$ 为无限集。对每个 $\alpha \in B'$，
存在有限集 $B_{\alpha} \subset B$，使得 $\alpha$ 在
$F(B_{\alpha})$ 上是代数的，因为任何代数依赖关系只使用
有限多个不定元。定义
$B^* = \bigcup_{\alpha\in B'} B_{\alpha}$。
按构造，$B^* \subset B$，但我们声称这两个集合其实相等。
若不相等，设存在元素 $\beta \in B \setminus B^*$。
已知 $\beta$ 在 $F(B')$ 上是代数的，而后者在
$F(B^*)$ 上是代数的。因此 $\beta$ 在 $F(B^*)$ 上是代数的，
矛盾。所以
$|B| \leq |\bigcup_{\alpha \in B'} B_{\alpha}|$。
若 $B'$ 有限，则 $B$ 也有限，故可假设 $B'$ 无限；此时
$$
|B| \leq |\bigcup\nolimits_{\alpha \in B'} B_{\alpha}| = |B'|
$$
因为每个 $B_\alpha$ 有限而 $B'$ 无限。
因此在无限情形，$|B| = |B'|$。

\medskip\noindent
现在考察 $B$ 有限的情形。此时 $B'$ 也有限，设
$B = \{\alpha_1, \ldots, \alpha_n\}$ 且
$B' = \{\beta_1, \ldots, \beta_m\}$，其中 $m \leq n$。
对 $m$ 作归纳：若 $m = 0$，则 $E/F$ 是代数的，故
$B = \emptyset$，从而 $n = 0$。若 $m > 0$，则存在不可约多项式
$f \in F[x, y_1, \ldots, y_n]$，使得
$f(\beta_1, \alpha_1, \ldots, \alpha_n) = 0$，并且 $x$
在 $f$ 中出现。由于 $\beta_1$ 在 $F$ 上不是代数的，
$f$ 必须涉及某个 $y_i$；不失一般性，假设 $f$ 使用 $y_1$。
令 $B^* = \{\beta_1, \alpha_2, \ldots, \alpha_n\}$。
我们声称 $B^*$ 是 $E/F$ 的一个基。为证明此断言，
注意到有代数扩张塔
$$
E/ F(B^*, \alpha_1) / F(B^*)
$$
因为 $\alpha_1$ 在 $F(B^*)$ 上是代数的。
现在声称 $B^*$（按元素的重数计）在 $F$ 上代数无关；
否则，就会存在不可约多项式
$g\in F[x, y_2, \ldots, y_n]$，使得
$g(\beta_1, \alpha_2, \ldots, \alpha_n) = 0$。
该多项式必然涉及 $x$，从而使 $\beta_1$ 在
$F(\alpha_2, \ldots, \alpha_n)$ 上代数；
这又会使 $\alpha_1$ 在
$F(\alpha_2, \ldots, \alpha_n)$ 上代数，这是不可能的。
因此，$\{\alpha_2, \ldots, \alpha_n\}$ 和
$\{\beta_2, \ldots, \beta_m\}$ 是 $E$ 在
$F(\beta_1)$ 上的基，由归纳假设得到 $m = n$。
\end{proof}

\begin{definition}
\label{fields-definition-transcendence-degree}
令 $K/k$ 为域扩张。$K$ 在 $k$ 上的{\it 超越次数}
是 $K$ 在 $k$ 上任一超越基的基数，记为
$\text{trdeg}_k(K)$。
\end{definition}

\begin{lemma}
\label{fields-lemma-transcendence-degree-tower}
令 $L/K/k$ 为域扩张。则
$$
\text{trdeg}_k(L) =
\text{trdeg}_K(L) +
\text{trdeg}_k(K).
$$
\end{lemma}

\begin{proof}
选择 $K$ 在 $k$ 上的超越基 $A \subset K$。
选择 $L$ 在 $K$ 上的超越基 $B \subset L$。
容易看出，$A \cup B$ 是 $L$ 在 $k$ 上的超越基。
\end{proof}

\begin{example}
\label{fields-example-pi-e-transcendental}
考虑由添入数 $e$ 和 $\pi$ 得到的域扩张
$\mathbf{Q}(e, \pi)$。由于 $e$ 与 $\pi$ 在有理数域上都超越，
此域扩张的超越次数至少为 $1$。然而，若 $e$ 与 $\pi$
代数无关，则此域扩张的超越次数可能为 $2$。
这是否成立尚不知晓，因而确定
$\text{trdeg}(\mathbf{Q}(e, \pi))$ 的问题仍是开放的。
\end{example}

\begin{example}
\label{fields-example-function-field-not-unique-transcendence-basis}
令 $F$ 为域，且 $E = F(t)$。由于 $E = F(t)$，
$\{t\}$ 是一个超越基。然而，$\{t^2\}$ 也是超越基，
因为 $F(t)/F(t^2)$ 是代数扩张。这说明，虽然总可把扩张
$E/F$ 分解为代数扩张 $E/F'$ 与纯超越扩张 $F'/F$，
但这种分解不唯一，并且依赖于超越基的选择。
\end{example}

\begin{example}
\label{fields-example-riemann-surface-transcendence}
令 $X$ 为紧 Riemann 曲面。则其函数域
$\mathbf{C}(X)$（见例 \ref{fields-example-field-of-meromorphic-functions}）
在 $\mathbf{C}$ 上的超越次数为一。事实上，$\mathbf{C}$
的{\it 任意}有限生成、超越次数为一的扩张都来自某个 Riemann 曲面。
紧 Riemann 曲面与（非常值）全纯映射所成范畴，甚至和
$\mathbf{C}$ 的超越次数为 $1$ 的有限生成扩张及
$\mathbf{C}$-代数态射所成范畴的反范畴等价。见 \cite{Forster}。

\medskip\noindent
上述陈述也有代数版本。给定代数闭域 $k$（例如复数域）上的
射影空间中的一条（不可约）代数曲线，可以考察其“有理函数域”：
大致说来，就是形如多项式之商的函数，其中分母不在曲线上恒为零。
光滑射影曲线与曲线的非常值态射所成范畴，和 $k$ 的超越次数为一的
有限生成扩张所成范畴之间存在类似的反等价
（代数曲线，定理 \ref{curves-theorem-curves-rational-maps}）。
见 \cite{H}。
\end{example}

\begin{definition}
\label{fields-definition-algebraically-closed-in}
令 $K/k$ 为域扩张。
\begin{enumerate}
\item {\it $k$ 在 $K$ 中的代数闭包}是 $K$ 的子域
$k'$，由 $K$ 中在 $k$ 上代数的元素组成。
\item 若 $K$ 中每个在 $k$ 上代数的元素都包含在 $k$ 中，
则称 $k$ 在{\it $K$ 中代数闭}。
\end{enumerate}
\end{definition}

\begin{lemma}
\label{fields-lemma-purely-transcendental-degree}
令 $k'/k$ 为有限域扩张。令
$k'(x_1, \ldots, x_r)/k(x_1, \ldots, x_r)$
为所诱导的纯超越扩张之间的扩张。则
$[k'(x_1, \ldots, x_r) : k(x_1, \ldots, x_r)] = [k' : k] < \infty$。
\end{lemma}

\begin{proof}
由扩张次数的乘法性
（引理 \ref{fields-lemma-multiplicativity-degrees}），
只需证明 $k'$ 由单个元素 $\alpha \in k'$ 在 $k$ 上生成的情形。
令 $f \in k[T]$ 为 $\alpha$ 在 $k$ 上的最小多项式。
于是 $k'(x_1, \ldots, x_r)$ 由
$\alpha, x_1, \ldots, x_r$ 在 $k$ 上生成，因而
$k'(x_1, \ldots, x_r)$ 由 $\alpha$ 在
$k(x_1, \ldots, x_r)$ 上生成。
故只需证明 $f$ 作为 $k(x_1, \ldots, x_r)[T]$ 的元素仍不可约。
我们只概述证明。显然，$f$ 作为
$k[x_1, \ldots, x_r, T]$ 的元素不可约；例如，这是因为
$f$ 作为 $T$ 的多项式是首一的，并且
$k[x_1, \ldots, x_r, T]$ 中任何可能的分解，在令
$x_i$ 等于 $0$ 后都会导出 $k[T]$ 中的分解。
由 Gauss 引理得到结论。
\end{proof}

\begin{lemma}
\label{fields-lemma-algebraic-closure-in-finitely-generated}
令 $K/k$ 为有限生成域扩张。$k$ 在 $K$ 中的代数闭包
在 $k$ 上有限。
\end{lemma}

\begin{proof}
令 $x_1, \ldots, x_r \in K$ 为 $K$ 在 $k$ 上的一个超越基。
则 $n = [K : k(x_1, \ldots, x_r)] < \infty$。
假设 $k \subset k' \subset K$，其中 $k'/k$ 有限。
在这种情况下，
$[k'(x_1, \ldots, x_r) : k(x_1, \ldots, x_r)] = [k' : k] < \infty$，
见引理 \ref{fields-lemma-purely-transcendental-degree}。因此
$$
[k' : k] = [k'(x_1, \ldots, x_r) : k(x_1, \ldots, x_r)] \leq
[K : k(x_1, \ldots, x_r)] = n.
$$
换言之，有限子扩张的次数有界，引理即得证。
\end{proof}







\section{线性无交扩张}
\label{fields-section-linearly-disjoint}

\noindent
令 $k$ 为域，$K$ 与 $L$ 为 $k$ 的域扩张。
再假设 $K$ 与 $L$ 都嵌入某个更大的域 $\Omega$ 中。

\begin{definition}
\label{fields-definition-compositum}
考虑域扩张图
\begin{equation}
\label{fields-equation-inside-omega}
\vcenter{
\xymatrix{
L \ar[r] & \Omega \\
k \ar[r] \ar[u] & K \ar[u]
}
}
\end{equation}
{\it $K$ 与 $L$ 在 $\Omega$ 中的合成域}，记为 $KL$，
是 $\Omega$ 中同时包含 $L$ 与 $K$ 的最小子域。
\end{definition}

\noindent
显然，$KL$ 由集合 $K \cup L$ 在 $k$ 上生成，
由集合 $K$ 在 $L$ 上生成，也由集合 $L$ 在 $K$ 上生成。

\medskip\noindent
{\bf 警告：}合成域的（同构类）依赖于 $K$ 与 $L$
嵌入 $\Omega$ 的选择。
% P01-E0260: frozen source says "composition" where the defined term is "compositum".
例如，考虑数域
$K = \mathbf{Q}(2^{1/8}) \subset \mathbf{R}$ 与
$L = \mathbf{Q}(2^{1/12}) \subset \mathbf{R}$。
在 $\mathbf{R}$ 中的合成域是
$\mathbf{Q}(2^{1/24})$，它在 $\mathbf{Q}$ 上的次数为 $24$。
然而，若把 $K = \mathbf{Q}[x]/(x^8 - 2)$ 嵌入
$\mathbf{C}$，其中把 $x$ 送到 $2^{1/8}e^{2\pi i/8}$，
则合成域 $\mathbf{Q}(2^{1/12}, 2^{1/8}e^{2\pi i/8})$
包含 $i = e^{2\pi i/4}$，并且在 $\mathbf{Q}$ 上的次数为 $48$
（我们略去次数为 $48$ 的证明，但 $i$ 的存在已经证明
两个合成域不同构）。

\begin{definition}
\label{fields-definition-linearly-disjoint}
考虑如式 (\ref{fields-equation-inside-omega}) 的域图。
若映射
$$
K \otimes_k L \longrightarrow KL,\quad
\sum x_i \otimes y_i \longmapsto \sum x_i y_i
$$
是单射，则称 $K$ 与 $L$ {\it 在 $\Omega$ 中于 $k$ 上线性无交}。
\end{definition}

\noindent
下面的引理似乎不适合放在其他地方。

\begin{lemma}
\label{fields-lemma-normal-case}
令 $E/F$ 为正规代数域扩张。存在子扩张
$E / E_{sep} /F$ 与 $E / E_{insep} / F$，使得：
\begin{enumerate}
\item $F \subset E_{sep}$ 是 Galois 扩张，且
$E_{sep} \subset E$ 是纯不可分扩张；
\item $F \subset E_{insep}$ 是纯不可分扩张，且
$E_{insep} \subset E$ 是 Galois 扩张；
\item $E = E_{sep} \otimes_F E_{insep}$。
\end{enumerate}
\end{lemma}

\begin{proof}
我们在引理 \ref{fields-lemma-separable-first} 中找到了子域 $E_{sep}$。
置 $E_{insep} = E^{\text{Aut}(E/F)}$。细节从略。
\end{proof}










\section{回顾}
\label{fields-section-algebraic}

\noindent
本节快速回顾以上内容。

\medskip\noindent
令 $K/k$ 为域扩张，令 $\alpha \in K$。则有下列可能：
\begin{enumerate}
\item 元素 $\alpha$ 在 $k$ 上超越。
\item 元素 $\alpha$ 在 $k$ 上代数。以
$P(T) \in k[T]$ 表示它的{\it 最小多项式}。
% P01-E0261: frozen English omits "by" after "Denote".
这是系数属于 $k$ 的首一多项式
$P(T) = T^d + a_1 T^{d - 1} + \ldots + a_d$。
它不可约且满足 $P(\alpha) = 0$。这些性质唯一确定 $P$，
整数 $d$ 称为{\it $\alpha$ 在 $k$ 上的次数}。
又分两种情形：
\begin{enumerate}
\item 多项式 $\text{d}P/\text{d}T$ 不恒为零。
这等价于如下条件：
$P(T) = \prod_{i = 1, \ldots, d} (T - \alpha_i)$，
其中 $\alpha_1, \ldots, \alpha_d$ 是 $k$ 的代数闭包中
两两不同的元素。在这种情况下，称 $\alpha$ 在 $k$ 上{\it 可分}。
\item $\text{d}P/\text{d}T$ 恒为零。在这种情况下，
$k$ 的特征 $p$ 为
% P01-E0262: frozen source has a malformed greater-than-zero math span without p.
$ > 0$，并且 $P$ 实际上是 $T^p$ 的多项式。
显然，存在最大的幂 $q = p^e$，使得 $P$ 是 $T^q$ 的多项式。
于是元素 $\alpha^q$ 在 $k$ 上可分。
\end{enumerate}
\end{enumerate}

\begin{definition}
\label{fields-definition-separable-algebraic}
代数域扩张。
\begin{enumerate}
\item 若 $K$ 的每个元素都在 $k$ 上代数，则称域扩张
$K/k$ 为{\it 代数扩张}。
\item 若每个 $\alpha \in k'$ 都在 $k$ 上可分，
则称代数扩张 $k'/k$ 为{\it 可分扩张}。
\item 若 $k$ 的特征为 $p > 0$，并且对每个元素
$\alpha \in k'$，都存在 $p$ 的一个幂 $q$，使得
$\alpha^q \in k$，则称代数扩张 $k'/k$ 为{\it 纯不可分扩张}。
\item 若对每个 $\alpha \in k'$，$\alpha$ 在 $k$ 上的
最小多项式 $P(T) \in k[T]$ 都在 $k'$ 上完全分解为线性因子，
则称代数扩张 $k'/k$ 为{\it 正规扩张}。
\item 若代数扩张 $k'/k$ 可分且正规，则称它为
{\it Galois 扩张}。
\end{enumerate}
\end{definition}

\noindent
下面的引理似乎不适合放在其他地方。

\begin{lemma}
\label{fields-lemma-pth-root}
令 $K$ 为特征 $p > 0$ 的域。令 $L/K$ 为可分代数扩张。
令 $\alpha \in L$。
\begin{enumerate}
\item 若 $\alpha$ 在 $K$ 上的最小多项式的系数都是
$K$ 中的 $p$ 次幂，则 $\alpha$ 是 $L$ 中的 $p$ 次幂。
\item 更一般地，若 $P \in K[T]$ 是满足下列条件的多项式：
(a) $\alpha$ 是 $P$ 的根；(b) $P$ 在一个代数闭包中有两两不同的根；
(c) $P$ 的所有系数都是 $p$ 次幂；则 $\alpha$
是 $L$ 中的 $p$ 次幂。
\end{enumerate}
\end{lemma}

\begin{proof}
由定义，(2) 蕴含 (1)。假设 $P$ 如 (2) 所述。
写成
$P(T) = \sum\nolimits_{i = 0}^d a_i T^{d - i}$ 且 $a_i = b_i^p$。
多项式
$Q(T) = \sum\nolimits_{i = 0}^d b_i T^{d - i}$
在代数闭包中也有两两不同的根，因为 $Q$ 的根是 $P$ 的根的
$p$ 次根。若 $\alpha$ 不是 $p$ 次幂，则
$T^p - \alpha$ 是 $L$ 上的不可约多项式
（引理 \ref{fields-lemma-take-pth-root}）。
此外，$Q$ 与 $T^p - \alpha$ 在代数闭包
$\overline{L}$ 中有共同根。
因此 $Q$ 与 $T^p - \alpha$ 不互素，这蕴含在 $L[T]$ 中
$T^p - \alpha | Q$。这与 $Q$ 的根两两不同矛盾。
\end{proof}
% Generated mechanically from the frozen algebra.tex target.
% Do not edit this generated file; edit the bound locale source instead.

% OK, start here.
%

\setcounter{chapter}{9}
\chapter{交换代数}



\phantomsection
\label{algebra-section-phantom}





\section{引言}
\label{algebra-section-introduction}

\noindent
本文将介绍交换代数的基础知识。可参阅 \cite{MatCA}。






\section{约定}
\label{algebra-section-conventions}

\noindent
环均指带 $1$ 的交换环。零环也是环；事实上，它是唯一没有素理想的环。
我们将使用 Kronecker 符号 $\delta_{ij}$。若 $R \to S$ 是环同态，且
$\mathfrak q$ 是 $S$ 的一个素理想，则用记号
``$\mathfrak p = R \cap \mathfrak q$''
表示 $\mathfrak q$ 在 $R \to S$ 下的逆像素理想，即使 $R$ 不是 $S$ 的子环，
甚至即使 $R \to S$ 不是单射，也采用这一记号。






\section{基本概念}
\label{algebra-section-rings-basic}

\noindent
下面列出交换代数中的一些基本概念。其中有些概念将在后文作更详细的讨论，
有些在此表中定义，其余则视为基础概念而不再定义。若你不熟悉以下大多数斜体概念，
建议先阅读一本代数入门教材，再继续阅读。

\begin{enumerate}
\item $R$ 是一个{\it 环}，
\label{algebra-item-ring}
\item $x\in R$ 是{\it 幂零元}，
\label{algebra-item-ring-element-nilpotent}
\item $x\in R$ 是{\it 零因子}，
\label{algebra-item-ring-element-zerodivisor}
\item $x\in R$ 是{\it 可逆元}，
\label{algebra-item-ring-element-unit}
\item $e \in R$ 是{\it 幂等元}，
\label{algebra-item-ring-element-idempotent}
\item 若幂等元 $e \in R$ 满足 $e = 1$ 或 $e = 0$，则称其为{\it 平凡的}，
\label{algebra-item-idempotent-trivial}
\item $\varphi : R_1 \to R_2$ 是{\it 环同态}，
\label{algebra-item-ring-homomorphism}
\item
\label{algebra-item-ring-homomorphism-finite-presentation}
$\varphi : R_1 \to R_2$ 是{\it 有限表示的}，亦即
{\it $R_2$ 是有限表示的 $R_1$-代数}，
见定义 \ref{algebra-definition-finite-type}，
\item
\label{algebra-item-ring-homomorphism-finite-type}
$\varphi : R_1 \to R_2$ 是{\it 有限型的}，亦即
{\it $R_2$ 是有限型 $R_1$-代数}，
见定义 \ref{algebra-definition-finite-type}，
\item
\label{algebra-item-ring-homomorphism-finite}
$\varphi : R_1 \to R_2$ 是{\it 有限的}，亦即
{\it $R_2$ 是有限 $R_1$-代数}，
\item $R$ 是{\it （整）环}，
\label{algebra-item-ring-domain}
\item $R$ 是{\it 约化的}，
\label{algebra-item-ring-reduced}
\item $R$ 是{\it Noether 环}，
\label{algebra-item-ring-Noetherian}
\item $R$ 是{\it 主理想整环}，简称 {\it PID}，
\label{algebra-item-ring-PID}
\item $R$ 是{\it 欧几里得整环}，
\label{algebra-item-ring-Euclidean}
\item $R$ 是{\it 唯一分解整环}，简称 {\it UFD}，
\label{algebra-item-ring-UFD}
\item $R$ 是{\it 离散赋值环}，简称 {\it dvr}，
\label{algebra-item-ring-dvr}
\item $K$ 是{\it 域}，
\label{algebra-item-field}
\item $L/K$ 是{\it 域扩张}，
\label{algebra-item-field-extension}
\item $L/K$ 是{\it 代数域扩张}，
\label{algebra-item-field-extension-algebraic}
\item $\{t_i\}_{i\in I}$ 是 $L$ 在 $K$ 上的{\it 超越基}，
\label{algebra-item-transcendence-basis}
\item $L$ 在 $K$ 上的{\it 超越次数} $\text{trdeg}(L/K)$，
\label{algebra-item-transcendence-degree}
\item 域 $k$ 是{\it 代数闭的}，
\label{algebra-item-algebraically-closed}
\item
\label{algebra-item-extend-into-algebraically-closed}
若 $L/K$ 是代数扩张，且 $\Omega/K$ 是扩张而 $\Omega$
为代数闭域，则存在环同态 $L \to \Omega$，它延拓 $K$ 上给定的同态，
\item $I \subset R$ 是{\it 理想}，
\label{algebra-item-ideal}
\item $I \subset R$ 是{\it 根理想}，
\label{algebra-item-ideal-radical}
\item 若 $I$ 是理想，则有其{\it 根} $\sqrt{I}$，
\label{algebra-item-radical-ideal}
\item
\label{algebra-item-ideal-nilpotent}
$I \subset R$ 为{\it 幂零理想}，意指对某个 $n \in \mathbf{N}$ 有 $I^n = 0$，
\item
\label{algebra-item-ideal-locally-nilpotent}
$I \subset R$ 为{\it 局部幂零理想}，意指 $I$ 的每个元素均为幂零元，
\item $\mathfrak p \subset R$ 是{\it 素理想}，
\label{algebra-item-prime-ideal}
\item
\label{algebra-item-prime-product-ideals}
若 $\mathfrak p \subset R$ 是素理想，$I, J \subset R$ 是理想，且
$IJ\subset \mathfrak p$，则 $I \subset \mathfrak p$ 或 $J \subset \mathfrak p$。
\item $\mathfrak m \subset R$ 是{\it 极大理想}，
\label{algebra-item-maximal-ideal}
\item 任一非零环都有极大理想，
\label{algebra-item-exists-maximal-ideal}
\item
\label{algebra-item-jacobson-radical}
$R$ 的{\it Jacobson 根}是 $R$ 的所有极大理想之交
$\text{rad}(R) = \bigcap_{\mathfrak m \subset R} \mathfrak m$，
\item 由子集 $T \subset R$ {\it 生成}的理想 $(T)$，
\label{algebra-item-ideal-generated-by}
\item {\it 商环} $R/I$，
\label{algebra-item-quotient-ring}
\item 环 $R$ 中的理想 $I$ 是素理想，当且仅当 $R/I$ 是整环，
\label{algebra-item-characterize-prime-ideal}
\item
\label{algebra-item-characterize-maximal-ideal}
环 $R$ 中的理想 $I$ 是极大理想，当且仅当环 $R/I$ 是域，
\item
\label{algebra-item-inverse-image-ideal}
若 $\varphi : R_1 \to R_2$ 是环同态，且 $I \subset R_2$ 是理想，
则 $\varphi^{-1}(I)$ 是 $R_1$ 的理想，
\item
\label{algebra-item-image-ideal}
若 $\varphi : R_1 \to R_2$ 是环同态，且 $I \subset R_1$ 是理想，
则 $\varphi(I) \cdot R_2$（有时记作 $I \cdot R_2$ 或 $IR_2$）
是由 $\varphi(I)$ 生成的 $R_2$ 的理想，
\item
\label{algebra-item-inverse-image-prime}
若 $\varphi : R_1 \to R_2$ 是环同态，且 $\mathfrak p \subset R_2$
是素理想，则 $\varphi^{-1}(\mathfrak p)$ 是 $R_1$ 的素理想，
\item $M$ 是{\it $R$-模}，
\label{algebra-item-module}
\item
\label{algebra-item-annihilator}
对 $m \in M$，元素 $m$ 在 $R$ 中的{\it 零化子}为
$I = \{f \in R \mid fm = 0\}$，
\item $N \subset M$ 是{\it $R$-子模}，
\label{algebra-item-submodule}
\item $M$ 是{\it Noether $R$-模}，
\label{algebra-item-Noetherian-module}
\item $M$ 是{\it 有限 $R$-模}，
\label{algebra-item-finite-module}
\item $M$ 是{\it 有限生成 $R$-模}，
\label{algebra-item-finitely-generated-module}
\item $M$ 是{\it 有限表示 $R$-模}，
\label{algebra-item-finitely-presented-module}
\item $M$ 是{\it 自由 $R$-模}，
\label{algebra-item-free-module}
\item
\label{algebra-item-extension-free}
若 $0 \to K \to L \to M \to 0$ 是 $R$-模的短正合列，且 $K$、$M$
均为自由模，则 $L$ 也是自由模，
\item 若 $N \subset M \subset L$ 是 $R$-模，则 $L/M = (L/N)/(M/N)$，
\label{algebra-item-isomorphism-theorem}
\item $S$ 是 {\it $R$ 的乘法子集}，
\label{algebra-item-multiplicative-subset}
\item $R$ 的{\it 局部化} $R \to S^{-1}R$，
\label{algebra-item-localization-ring}
\item
\label{algebra-item-localization-zero}
若 $R$ 是环，$S$ 是 $R$ 的乘法子集，则 $S^{-1}R$ 是零环，
当且仅当 $S$ 含有 $0$，
\item
\label{algebra-item-localize-nonzerodivisors}
若 $R$ 是环，且乘法子集 $S$ 完全由非零因子组成，
则 $R \to S^{-1}R$ 是单射，
\item 若 $\varphi : R_1 \to R_2$ 是环同态，且 $S$ 是 $R_1$ 的乘法子集，
则 $\varphi(S)$ 是 $R_2$ 的乘法子集，
\item
\label{algebra-item-products-multiplicative-subsets}
若 $S$、$S'$ 是 $R$ 的乘法子集，并以 $SS'$ 表示乘积集合
$SS' = \{r \in R \mid \exists s\in S, \exists s' \in S', r = ss'\}$，
则 $SS'$ 是 $R$ 的乘法子集，
\item
\label{algebra-item-localization-localization}
若 $S$、$S'$ 是 $R$ 的乘法子集，且 $\overline{S}$ 表示 $S$ 在
$(S')^{-1}R$ 中的像，则
$(SS')^{-1}R = \overline{S}^{-1}((S')^{-1}R)$，
\item $R$-模 $M$ 的{\it 局部化} $S^{-1}M$，
\label{algebra-item-localization-module}
\item
\label{algebra-item-localization-exact}
函子 $M \mapsto S^{-1}M$ 保持单射、满射和正合性，
\item
\label{algebra-item-localization-localization-module}
若 $S$、$S'$ 是 $R$ 的乘法子集，且 $M$ 是 $R$-模，则
$(SS')^{-1}M = S^{-1}((S')^{-1}M)$，
\item
\label{algebra-item-localize-ideal}
若 $R$ 是环，$I$ 是 $R$ 的理想，$S$ 是 $R$ 的乘法子集，
则 $S^{-1}I$ 是 $S^{-1}R$ 的理想，并且
$S^{-1}R/S^{-1}I = \overline{S}^{-1}(R/I)$，其中 $\overline{S}$
是 $S$ 在 $R/I$ 中的像，
\item
\label{algebra-item-ideal-in-localization}
若 $R$ 是环，$S$ 是 $R$ 的乘法子集，则 $S^{-1}R$ 的任一理想 $I'$
均形如 $S^{-1}I$，其中可取 $I$ 为 $I'$ 在 $R$ 中的逆像，
\item
\label{algebra-item-submodule-in-localization}
若 $R$ 是环，$M$ 是 $R$-模，$S$ 是 $R$ 的乘法子集，则
$S^{-1}M$ 的任一子模 $N'$ 均形如 $S^{-1}N$，其中 $N \subset M$
是某个子模；并且可取 $N$ 为 $N'$ 在 $M$ 中的逆像，
\item 若 $S = \{1, f, f^2, \ldots\}$，则 $R_f = S^{-1}R$ 且 $M_f = S^{-1}M$，
\label{algebra-item-localize-f}
\item
\label{algebra-item-localize-p}
若对某个素理想 $\mathfrak p$ 有
$S = R \setminus \mathfrak p = \{x\in R \mid x\not\in \mathfrak p\}$，
则通常记 $R_{\mathfrak p} = S^{-1}R$ 且
$M_{\mathfrak p} = S^{-1}M$，
\item {\it 局部环}是恰有一个极大理想的环，
\label{algebra-item-local-ring}
\item {\it 半局部环}是只有有限多个极大理想的环，
\label{algebra-item-semi-local-ring}
\item
\label{algebra-item-localize-p-local-ring}
若 $\mathfrak p$ 是 $R$ 中的素理想，则 $R_{\mathfrak p}$ 是局部环，
其极大理想为 $\mathfrak p R_{\mathfrak p}$，
\item
\label{algebra-item-residue-field}
环 $R$ 中素理想 $\mathfrak p$ 的{\it 剩余域}记作 $\kappa(\mathfrak p)$，
它是整环 $R/\mathfrak p$ 的分式域；并且等于
$R_\mathfrak p/\mathfrak pR_\mathfrak p
= (R \setminus \mathfrak p)^{-1}R/\mathfrak p$，
\item 给定 $R$ 以及 $M_1$、$M_2$，有{\it 张量积} $M_1 \otimes_R M_2$，
\label{algebra-item-tensor-product}
\item
\label{algebra-item-cauchy-binet}
给定环 $R$ 中分别为 $m \times n$ 和 $n \times m$ 的矩阵 $A$、$B$，
则在 $R$ 中有下式：
$$
\det(AB) = \sum \det(A_S)\det({}_SB)
$$
此处求和遍历
所有大小为 $m$ 的子集 $S \subset \{1, \ldots, n\}$，$A_S$ 是由 $A$
中对应于 $S$ 的列构成的 $m \times m$ 子矩阵，而 ${}_SB$ 是由 $B$
中对应于 $S$ 的行构成的 $m \times m$ 子矩阵，
\item 等等。
\end{enumerate}




\section{蛇形引理}
\label{algebra-section-snake}

\noindent
蛇形引理及其各种变体在同调一章的 Abel 范畴框架中讨论，
见同调，第 \ref{homology-section-abelian-categories} 节。

\begin{lemma}
\label{algebra-lemma-snake}
\begin{reference}
\cite[III, Lemma 3.3]{Cartan-Eilenberg}
\end{reference}
给定行正合的 Abel 群交换图
$$
\xymatrix{
& X \ar[r] \ar[d]^\alpha &
Y \ar[r] \ar[d]^\beta &
Z \ar[r] \ar[d]^\gamma &
0 \\
0 \ar[r] & U \ar[r] & V \ar[r] & W
}
$$
则存在典范正合列
$$
\Ker(\alpha) \to \Ker(\beta) \to \Ker(\gamma)
\to
\Coker(\alpha) \to \Coker(\beta) \to \Coker(\gamma)
$$
此外，若 $X \to Y$ 是单射，则第一个映射为单射；若 $V \to W$
是满射，则最后一个映射为满射。
\end{lemma}

\begin{proof}
映射 $\partial : \Ker(\gamma) \to \Coker(\alpha)$ 定义如下。
取 $z \in \Ker(\gamma)$，选择一个映到 $z$ 的 $y \in Y$。
于是 $\beta(y) \in V$ 在 $W$ 中映为零，故 $\beta(y)$ 是某个
$u \in U$ 的像。令 $\partial z = \overline{u}$，即 $u$ 在 $\alpha$
的余核中的类。正合性的证明从略。
\end{proof}

\section{有限模与有限表示模}
\label{algebra-section-module-finite-type}

\noindent
下面给出一些基本记号和引理。

\begin{definition}
\label{algebra-definition-module-finite-type}
设 $R$ 为环，$M$ 为 $R$-模。
\begin{enumerate}
\item 若存在 $n \in \mathbf{N}$ 及 $x_1, \ldots, x_n \in M$，
使 $M$ 的每个元素都是诸 $x_i$ 的 $R$-线性组合，则称 $M$ 为
{\it 有限 $R$-模}，亦称{\it 有限生成 $R$-模}。等价地，这意味着对某个
$n \in \mathbf{N}$ 存在满射 $R^{\oplus n} \to M$。
\item 若存在整数 $n, m \in \mathbf{N}$ 以及正合列
$$
R^{\oplus m} \longrightarrow R^{\oplus n} \longrightarrow M \longrightarrow 0
$$
则称 $M$ 为{\it 有限表示 $R$-模}，亦称{\it 有限表示的 $R$-模}。
\end{enumerate}
\end{definition}

\noindent
非正式地说，$M$ 是有限表示 $R$-模，当且仅当它有限生成，且这些生成元之间的
关系所成的模也有限生成。定义中那样一个正合列的选取称为 $M$ 的一个{\it 表示}。

\begin{lemma}
\label{algebra-lemma-lift-map}
设 $R$ 为环。设 $\alpha : R^{\oplus n} \to M$ 和 $\beta : N \to M$
为模同态。若 $\Im(\alpha) \subset \Im(\beta)$，则存在 $R$-模同态
$\gamma : R^{\oplus n} \to N$，使得 $\alpha = \beta \circ \gamma$。
\end{lemma}

\begin{proof}
令 $e_i = (0, \ldots, 0, 1, 0, \ldots, 0)$ 为 $R^{\oplus n}$ 的第 $i$ 个基向量。
由假设可取 $x_i \in N$，使得 $\alpha(e_i) = \beta(x_i)$。令
$\gamma(a_1, \ldots, a_n) = \sum a_i x_i$。由构造即有
$\alpha = \beta \circ \gamma$。
\end{proof}

\begin{lemma}
\label{algebra-lemma-extension}
设 $R$ 为环，并设
$$
0 \to M_1 \to M_2 \to M_3 \to 0
$$
为 $R$-模的短正合列。
\begin{enumerate}
\item 若 $M_1$ 和 $M_3$ 是有限 $R$-模，则 $M_2$ 是有限 $R$-模。
\item 若 $M_1$ 和 $M_3$ 是有限表示 $R$-模，则 $M_2$ 是有限表示 $R$-模。
\item 若 $M_2$ 是有限 $R$-模，则 $M_3$ 是有限 $R$-模。
\item 若 $M_2$ 是有限表示 $R$-模，且 $M_1$ 是有限 $R$-模，则
$M_3$ 是有限表示 $R$-模。
\item 若 $M_3$ 是有限表示 $R$-模，且 $M_2$ 是有限 $R$-模，则
$M_1$ 是有限 $R$-模。
\end{enumerate}
\end{lemma}

\begin{proof}
先证 (1)。若 $x_1, \ldots, x_n$ 是 $M_1$ 的生成元，而
$y_1, \ldots, y_m \in M_2$ 在 $M_3$ 中的像生成 $M_3$，则
$x_1, \ldots, x_n, y_1, \ldots, y_m$ 生成 $M_2$。

\medskip\noindent
(3) 由定义立即得到。

\medskip\noindent
证明 (5)。设 $M_3$ 有限表示且 $M_2$ 有限。选取一个表示
$$
R^{\oplus m} \to R^{\oplus n} \to M_3 \to 0
$$
由引理 \ref{algebra-lemma-lift-map}，存在映射 $R^{\oplus n} \to M_2$，
使下图中的实线部分
$$
\xymatrix{
& R^{\oplus m} \ar[r] \ar@{..>}[d] & R^{\oplus n} \ar[r] \ar[d] &
M_3 \ar[r] \ar[d]^{\text{id}} & 0 \\
0 \ar[r] & M_1 \ar[r] & M_2 \ar[r] & M_3 \ar[r] & 0
}
$$
交换，因而得到虚线箭头。由蛇形引理（引理 \ref{algebra-lemma-snake}）可知存在同构
$$
\Coker(R^{\oplus m} \to M_1)
\cong
\Coker(R^{\oplus n} \to M_2)
$$
特别地，$\Coker(R^{\oplus m} \to M_1)$ 是有限 $R$-模。又因
$\Im(R^{\oplus m} \to M_1)$ 由 (3) 为有限模，所以由 (1) 得 $M_1$ 有限。

\medskip\noindent
证明 (4)。设 $M_2$ 有限表示且 $M_1$ 有限。选取表示
$R^{\oplus m} \to R^{\oplus n} \to M_2 \to 0$，并选取满射
$R^{\oplus k} \to M_1$。由引理 \ref{algebra-lemma-lift-map}，复合映射
$R^{\oplus k} \to M_1 \to M_2$ 可分解为
$R^{\oplus k} \to R^{\oplus n} \to M_2$。于是
$R^{\oplus k + m} \to R^{\oplus n} \to M_3 \to 0$
是一个表示。

\medskip\noindent
证明 (2)。设 $M_1$ 和 $M_3$ 均为有限表示模。(1) 的证明中的论证给出交换图
$$
\xymatrix{
0 \ar[r] & R^{\oplus n} \ar[d] \ar[r] & R^{\oplus n + m} \ar[d] \ar[r] &
R^{\oplus m} \ar[d] \ar[r] & 0 \\
0 \ar[r] & M_1 \ar[r] & M_2 \ar[r] & M_3 \ar[r] & 0
}
$$
其中竖直箭头均为满射。由蛇形引理得到短正合列
$$
0 \to \Ker(R^{\oplus n} \to M_1) \to
\Ker(R^{\oplus n + m} \to M_2) \to
\Ker(R^{\oplus m} \to M_3) \to 0
$$
由 (5)，外侧两个模均有限，故中间的模也有限。再由 (4)，$M_2$ 为有限表示模。
\end{proof}

\begin{lemma}
\label{algebra-lemma-trivial-filter-finite-module}
\begin{slogan}
有限模存在滤过，使其逐次商均为循环模。
\end{slogan}
设 $R$ 为环，$M$ 为有限 $R$-模。存在由有限 $R$-子模组成的滤过
$$
0 = M_0 \subset M_1 \subset \ldots \subset M_n = M
$$
使每个商 $M_i/M_{i - 1}$ 都同构于 $R/I_i$，其中 $I_i$ 是 $R$ 的某个理想。
\end{lemma}

\begin{proof}
对 $M$ 的生成元数作归纳。设 $x_1, \ldots, x_r \in M$ 为一组生成元。
令 $M' = Rx_1 \subset M$。则 $M/M'$ 有 $r - 1$ 个生成元，故可应用归纳假设。
并且显然 $M' \cong R/I_1$，其中 $I_1 = \{f \in R \mid fx_1 = 0\}$。
\end{proof}

\begin{lemma}
\label{algebra-lemma-finite-over-subring}
设 $R \to S$ 为环同态，$M$ 为 $S$-模。若 $M$ 作为 $R$-模是有限的，
则 $M$ 作为 $S$-模也是有限的。
\end{lemma}

\begin{proof}
事实上，$M$ 的任一 $R$-生成集也是 $M$ 的 $S$-生成集，
因为其 $R$-模结构由 $R$ 在 $S$ 中的像诱导。
\end{proof}


\section{有限型与有限表示的环同态}
\label{algebra-section-finite-type}

\begin{definition}
\label{algebra-definition-finite-type}
设 $R \to S$ 为环同态。
\begin{enumerate}
\item 若存在 $n \in \mathbf{N}$ 以及 $R$-代数满射
$R[x_1, \ldots, x_n] \to S$，则称 $R \to S$ 为{\it 有限型的}，
亦称{\it $S$ 是有限型 $R$-代数}。
% P01-E0263: source wording candidate recorded in ERRATA_CANDIDATES.jsonl.
\item 若存在整数 $n, m \in \mathbf{N}$、多项式
$f_1, \ldots, f_m \in R[x_1, \ldots, x_n]$ 以及下列 $R$-代数同构：
$$
R[x_1, \ldots, x_n]/(f_1, \ldots, f_m) \cong S
$$
则称 $R \to S$ 为{\it 有限表示的}。
\end{enumerate}
\end{definition}

\noindent
非正式地说，$R \to S$ 为有限表示，当且仅当 $S$ 作为 $R$-代数有限生成，
且这些生成元之间的关系所成的理想有限生成。定义中的一个满射
$R[x_1, \ldots, x_n] \to S$ 有时称为 $S$ 的一个{\it 表示}。

\begin{lemma}
\label{algebra-lemma-compose-finite-type}
有限型与有限表示具有下列稳定性质。
\begin{enumerate}
\item 有限型环同态的复合仍为有限型。
\item 有限表示环同态的复合仍为有限表示。
\item 给定 $R \to S' \to S$，若 $R \to S$ 为有限型，则
$S' \to S$ 为有限型。
\item 给定 $R \to S' \to S$，若 $R \to S$ 为有限表示，且
$R \to S'$ 为有限型，则 $S' \to S$ 为有限表示。
\end{enumerate}
\end{lemma}

\begin{proof}
只证最后一个断言。写成
$S = R[x_1, \ldots, x_n]/(f_1, \ldots, f_m)$ 及
$S' = R[y_1, \ldots, y_a]/I$。设 $y_i$ 的类 $\bar y_i$ 在 $S$
中映为 $h_i \bmod (f_1, \ldots, f_m)$。于是显然
$S = S'[x_1, \ldots, x_n]/(f_1, \ldots, f_m,
h_1 - \bar y_1, \ldots, h_a - \bar y_a)$。
\end{proof}

\begin{lemma}
\label{algebra-lemma-finite-presentation-independent}
设 $R \to S$ 为有限表示环同态。对任一满射
$\alpha : R[x_1, \ldots, x_n] \to S$，$\alpha$ 的核都是
$R[x_1, \ldots, x_n]$ 中的有限生成理想。
\end{lemma}

\begin{proof}
写成 $S = R[y_1, \ldots, y_m]/(f_1, \ldots, f_k)$。
选取 $g_i \in R[y_1, \ldots, y_m]$ 作为 $\alpha(x_i)$ 的提升。
于是 $S = R[x_i, y_j]/(f_l, x_i - g_i)$。选取
$h_j \in R[x_1, \ldots, x_n]$，使 $\alpha(h_j)$ 对应于
$y_j \bmod (f_1, \ldots, f_k)$。考虑映射
$\psi : R[x_i, y_j] \to R[x_i]$，$x_i \mapsto x_i$，
$y_j \mapsto h_j$。则 $\alpha$ 的核是 $(f_l, x_i - g_i)$
在 $\psi$ 下的像，结论成立。
\end{proof}

\begin{lemma}
\label{algebra-lemma-finitely-presented-over-subring}
设 $R \to S$ 为环同态，$M$ 为 $S$-模。假设 $R \to S$ 为有限型，
且 $M$ 作为 $R$-模有限表示，则 $M$ 作为 $S$-模也有限表示。
\end{lemma}

\begin{proof}
这与引理 \ref{algebra-lemma-compose-finite-type} 中 (4) 的证明类似。
可设 $S = R[x_1, \ldots, x_n]/J$。选取
$y_1, \ldots, y_m \in M$，使其作为 $R$-模生成 $M$；并选取关系
$\sum a_{ij} y_j = 0$，$i = 1, \ldots, t$，使其生成
$R^{\oplus m} \to M$ 的核。对任意 $i = 1, \ldots, n$ 和
$j = 1, \ldots, m$，写成
$$
x_i y_j = \sum a_{ijk} y_k
$$
其中 $a_{ijk} \in R$。考虑由 $y_1, \ldots, y_m$ 生成的 $S$-模 $N$，
其关系为 $\sum a_{ij} y_j = 0$，$i = 1, \ldots, t$，以及
$x_i y_j = \sum a_{ijk} y_k$，$i = 1, \ldots, n$、
$j = 1, \ldots, m$。于是 $N$ 有表示
$$
S^{\oplus nm + t} \longrightarrow S^{\oplus m} \longrightarrow N
\longrightarrow 0
$$
由构造存在满射 $\varphi : N \to M$。为完成证明，只需证明 $\varphi$
为单射。设对某些 $b_j \in S$ 有 $z = \sum b_j y_j \in N$。
可将 $b_j$ 视为以 $R$ 中元素为系数、关于 $x_1, \ldots, x_n$ 的多项式。
应用形如 $x_i y_j = \sum a_{ijk} y_k$ 的关系，可归纳降低这些多项式的次数。
故对某些 $c_j \in R$ 有 $z = \sum c_j y_j$。于是若
$\varphi(z) = 0$，则向量 $(c_1, \ldots, c_m)$ 是诸向量
$(a_{i1}, \ldots, a_{im})$ 的 $R$-线性组合，从而如所需地得到 $z = 0$。
\end{proof}





\section{有限环同态}
\label{algebra-section-finite}

\noindent
定义如下。

\begin{definition}
\label{algebra-definition-finite-ring-map}
设 $\varphi : R \to S$ 为环同态。若 $S$ 作为 $R$-模有限，
则称 $\varphi : R \to S$ {\it 有限}。
\end{definition}

\begin{lemma}
\label{algebra-lemma-finite-module-over-finite-extension}
设 $R \to S$ 为有限环同态，$M$ 为 $S$-模。则 $M$ 作为 $R$-模有限，
当且仅当 $M$ 作为 $S$-模有限。
\end{lemma}

\begin{proof}
其中一个方向由引理 \ref{algebra-lemma-finite-over-subring} 得到。
为证另一个方向，设 $M$ 作为 $S$-模有限。取
$x_1, \ldots, x_n \in S$，使其作为 $R$-模生成 $S$；取
$y_1, \ldots, y_m \in M$，使其作为 $S$-模生成 $M$。
则诸 $x_i y_j$ 作为 $R$-模生成 $M$。
\end{proof}

\begin{lemma}
\label{algebra-lemma-finite-transitive}
设 $R \to S$ 与 $S \to T$ 均为有限环同态，则 $R \to T$ 有限。
\end{lemma}

\begin{proof}
若诸 $t_i$ 作为 $S$-模生成 $T$，诸 $s_j$ 作为 $R$-模生成 $S$，
则诸 $t_i s_j$ 作为 $R$-模生成 $T$。（亦可由引理
\ref{algebra-lemma-finite-module-over-finite-extension} 得到。）
\end{proof}

\begin{lemma}
\label{algebra-lemma-finite-finite-type}
设 $\varphi : R \to S$ 为环同态。
\begin{enumerate}
\item 若 $\varphi$ 有限，则 $\varphi$ 为有限型。
\item 若 $S$ 作为 $R$-模为有限表示，则 $\varphi$ 为有限表示。
\end{enumerate}
\end{lemma}

\begin{proof}
对 (1)，若 $x_1, \ldots, x_n \in S$ 作为 $R$-模生成 $S$，
则 $x_1, \ldots, x_n$ 也作为 $R$-代数生成 $S$。对 (2)，设关系
$\sum r_j^ix_i = 0$，$j = 1, \ldots, m$ 作为 $S$ 的 $R$-模生成元 $x_i$
之间的关系的一组生成元。此外，对某些 $r_i \in R$ 写成
$1 = \sum r_ix_i$，并对某些 $r_{ij}^k \in R$ 写成
$x_ix_j = \sum r_{ij}^k x_k$。于是作为 $R$-代数有
$$
S = R[t_1, \ldots, t_n]/
(\sum r_j^it_i,\ 1 - \sum r_it_i,\ t_it_j - \sum r_{ij}^k t_k)
$$
这证明了 (2)。
\end{proof}

\noindent
关于有限环同态的更多内容，见第 \ref{algebra-section-finite-ring-extensions} 节。



\section{余极限}
\label{algebra-section-colimits}

\noindent
本节部分内容与范畴论中关于余极限的一般讨论重叠，见范畴论第
\ref{categories-section-limits}--\ref{categories-section-posets-limits} 节。
预序集的概念定义于范畴论，定义 \ref{categories-definition-directed-set}；
它比偏序集的概念稍弱。

\begin{definition}
\label{algebra-definition-directed-system}
设 $(I, \leq)$ 为预序集。一个{\it 以 $I$ 为指标的 $R$-模系统
$(M_i, \mu_{ij})$}，由以 $I$ 为指标的一族 $R$-模
$\{M_i\}_{i\in I}$ 以及一族 $R$-模同态
$\{\mu_{ij} : M_i \to M_j\}_{i \leq j}$ 构成，并且对所有
$i \leq j \leq k$ 有
$$
\mu_{ii} = \text{id}_{M_i}\quad
\mu_{ik} = \mu_{jk}\circ \mu_{ij}
$$
若 $I$ 是有向集，则称 $(M_i, \mu_{ij})$ 为{\it 有向系统}。
\end{definition}

\noindent
这与范畴论，定义 \ref{categories-definition-system-over-poset}
及第 \ref{categories-section-posets-limits} 节中定义的概念相同。
任意范畴中一个图式或系统的余极限之定义，见范畴论，定义
\ref{categories-definition-colimit}。

\begin{lemma}
\label{algebra-lemma-colimit}
设 $(M_i, \mu_{ij})$ 是预序集 $I$ 上的 $R$-模系统。
该系统 $(M_i, \mu_{ij})$ 的余极限是商 $R$-模
$(\bigoplus_{i\in I} M_i) /Q$，其中 $Q$ 是由所有形如
$$
\iota_i(x_i) - \iota_j(\mu_{ij}(x_i))
$$
的元素生成的 $R$-子模，而
$\iota_i : M_i \to \bigoplus_{i\in I} M_i$ 是自然嵌入。
将余极限记作 $M = \colim_i M_i$，投影记作
$\pi : \bigoplus_{i\in I} M_i \to M$，并令
$\phi_i = \pi \circ \iota_i : M_i \to M$。
\end{lemma}

\begin{proof}
此引理是范畴论，引理
\ref{categories-lemma-colimits-coproducts-coequalizers} 的特例，
不过我们也直接证明这一情形。注意在上述构造中
$\phi_i = \phi_j\circ \mu_{ij}$。要证明偶对 $(M, \phi_i)$ 是余极限，
须验证其泛性质：对任一其他这样的偶对 $(Y, \psi_i)$，其中
$\psi_i : M_i \to Y$ 且 $\psi_i = \psi_j\circ \mu_{ij}$，
存在唯一的 $R$-模同态 $g : M \to Y$，使下图交换：
$$
\xymatrix{
M_i \ar[rr]^{\mu_{ij}} \ar[dr]^{\phi_i} \ar[ddr]_{\psi_i} & &
M_j\ar[dl]_{\phi_j} \ar[ddl]^{\psi_j} \\
& M \ar[d]^{g}\\
& Y
}
$$
这是显然的，因为可在直和 $\bigoplus M_i$ 的直和项 $M_i$ 上
取映射 $\psi_i$，由此定义 $g$。
\end{proof}

\begin{lemma}
\label{algebra-lemma-directed-colimit}
设 $(M_i, \mu_{ij})$ 是预序集 $I$ 上的 $R$-模系统，并假设 $I$ 有向。
该系统 $(M_i, \mu_{ij})$ 的余极限典范同构于如下定义的模 $M$：
\begin{enumerate}
\item 作为集合，令
$$
M = \left(\coprod\nolimits_{i \in I} M_i\right)/\sim
$$
其中，对 $m \in M_i$ 和 $m' \in M_{i'}$，规定
$$
m \sim m' \Leftrightarrow
\mu_{ij}(m) = \mu_{i'j}(m')\text{ 对某个 }j \geq i, i'
$$
\item 作为 Abel 群，对 $m \in M_i$ 和 $m' \in M_{i'}$，
定义 $m$ 与 $m'$ 在 $M$ 中的类之和为
$\mu_{ij}(m) + \mu_{i'j}(m')$ 的类，其中 $j \in I$ 是任一满足
$i \leq j$ 且 $i' \leq j$ 的指标；
\item 作为 $R$-模，对 $m \in M_i$ 和 $x \in R$，定义 $x$ 与 $m$
在 $M$ 中的类之积为 $xm$ 在 $M$ 中的类。
\end{enumerate}
典范映射 $\phi_i : M_i \to M$ 由典范映射
$M_i \to \coprod_{i \in I} M_i$ 诱导。
\end{lemma}

\begin{proof}
从略。可与范畴论，第 \ref{categories-section-directed-colimits} 节比较。
\end{proof}

\begin{lemma}
\label{algebra-lemma-zero-directed-limit}
设 $(M_i, \mu_{ij})$ 为有向系统，并设 $M = \colim M_i$，其映射为
$\mu_i : M_i \to M$。则对 $x_i \in M_i$，$\mu_i(x_i) = 0$，
当且仅当存在 $j \geq i$ 使得 $\mu_{ij}(x_i) = 0$。
\end{lemma}

\begin{proof}
这由引理 \ref{algebra-lemma-directed-colimit} 中有向余极限的描述立即得到。
\end{proof}

\begin{example}
\label{algebra-example-zero-colimit-different}
考虑偏序集 $I = \{a, b, c\}$，其中 $a < b$、$a < c$，且无其他严格不等式。
$I$ 上的一个系统 $(M_a, M_b, M_c, \mu_{ab}, \mu_{ac})$ 由三个
$R$-模 $M_a, M_b, M_c$ 以及两个 $R$-模同态
$\mu_{ab} : M_a \to M_b$ 和 $\mu_{ac} : M_a \to M_c$ 构成。
该系统的余极限就是
$$
M := \colim_{i \in I} M_i = \Coker(M_a \to M_b \oplus M_c)
$$
其中映射为 $\mu_{ab} \oplus -\mu_{ac}$。因此，典范映射
$M_a \to M$ 的核为 $\Ker(\mu_{ab}) + \Ker(\mu_{ac})$；
而典范映射 $M_b \to M$ 的核是 $\Ker(\mu_{ac})$ 在映射
$\mu_{ab}$ 下的像。故引理 \ref{algebra-lemma-zero-directed-limit}
的结论对一般系统显然不成立。
\end{example}


\begin{definition}
\label{algebra-definition-homomorphism-directed-systems}
设 $(M_i, \mu_{ij})$、$(N_i, \nu_{ij})$ 是同一预序集 $I$ 上的
$R$-模系统。从 $(M_i, \mu_{ij})$ 到 $(N_i, \nu_{ij})$ 的
一个{\it 系统同态} $\Phi$，按定义是由 $R$-模同态
$\phi_i : M_i \to N_i$ 组成的一族，并且对所有 $i \leq j$ 有
$\phi_j \circ \mu_{ij} = \nu_{ij} \circ \phi_i$。
\end{definition}

\noindent
用范畴语言说，这与相应图式 $M : I \to \text{Mod}_R$ 和
$N : I \to \text{Mod}_R$ 之间的函子变换是同一个概念。
下一个引理是范畴论，引理 \ref{categories-lemma-functorial-colimit} 的特例。

\begin{lemma}
\label{algebra-lemma-homomorphism-limit}
设 $(M_i, \mu_{ij})$、$(N_i, \nu_{ij})$ 是同一预序集上的 $R$-模系统。
从 $(M_i, \mu_{ij})$ 到 $(N_i, \nu_{ij})$ 的系统同态
$\Phi = (\phi_i)$ 诱导唯一同态
$$
\colim \phi_i : \colim M_i \longrightarrow \colim N_i
$$
使得对所有 $i \in I$，下图交换：
$$
\xymatrix{
M_i \ar[r] \ar[d]_{\phi_i} & \colim M_i \ar[d]^{\colim \phi_i} \\
N_i \ar[r] & \colim N_i
}
$$
\end{lemma}

\begin{proof}
写成 $M = \colim M_i$、$N = \colim N_i$ 及 $\phi = \colim \phi_i$
（后者尚待构造）。以下将不再另作说明地使用引理 \ref{algebra-lemma-colimit}
中对 $M$ 和 $N$ 的显式描述。引理的条件等价于下图交换：
$$
\xymatrix{
\bigoplus_{i\in I} M_i \ar[r] \ar[d]_{\bigoplus\phi_i} & M \ar[d]^\phi \\
\bigoplus_{i\in I} N_i \ar[r] & N
}
$$
因此，若 $\phi$ 存在，则其唯一性显然。为证明 $\phi$ 存在，
只须说明上方水平箭头的核在 $\bigoplus \phi_i$ 下映入下方水平箭头的核。
为此，取 $j \leq k$ 及 $x_j \in M_j$，则
$$
(\bigoplus \phi_i)(x_j - \mu_{jk}(x_j))
=
\phi_j(x_j) - \phi_k(\mu_{jk}(x_j))
=
\phi_j(x_j) - \nu_{jk}(\phi_j(x_j))
$$
如所需地属于下方水平箭头的核。
\end{proof}

\begin{lemma}
\label{algebra-lemma-directed-colimit-exact}
\begin{slogan}
滤过余极限是正合的。有向余极限是正合的。
\end{slogan}
设 $I$ 为有向集。设 $(L_i, \lambda_{ij})$、$(M_i, \mu_{ij})$ 和
$(N_i, \nu_{ij})$ 为 $I$ 上的 $R$-模系统。设
$\varphi_i : L_i \to M_i$ 与 $\psi_i : M_i \to N_i$ 为 $I$ 上的系统同态。
假设对每个 $i \in I$，$R$-模序列
$$
\xymatrix{
L_i \ar[r]^{\varphi_i} &
M_i \ar[r]^{\psi_i} &
N_i
}
$$
是同调为 $H_i$ 的复形。则诸 $R$-模 $H_i$ 构成 $I$ 上的系统；
$R$-模序列
$$
\xymatrix{
\colim_i L_i \ar[r]^\varphi &
\colim_i M_i \ar[r]^\psi &
\colim_i N_i
}
$$
也是复形；若以 $H$ 表示其同调，则
$$
H = \colim_i H_i.
$$
\end{lemma}

\begin{proof}
显然
$
\xymatrix{
\colim_i L_i \ar[r]^\varphi &
\colim_i M_i \ar[r]^\psi &
\colim_i N_i
}
$
是复形。对每个 $i \in I$，存在典范 $R$-模同态 $H_i \to H$：
它把每个 $[m] \in H_i = \Ker(\psi_i) / \Im(\varphi_i)$ 映到
$m$ 在 $\colim_i M_i$ 中的像于
$H = \Ker(\psi) / \Im(\varphi)$ 中的剩余类。这些同态给出同态
$\colim_i H_i \to H$。只须证明它既满又单。

\medskip\noindent
以下将反复使用引理 \ref{algebra-lemma-directed-colimit} 中关于 $I$ 上余极限的描述，
而不再另作说明。取 $h \in H$。由于
$H = \Ker(\psi)/\Im(\varphi)$，$h$ 是模 $\Im(\varphi)$ 意义下
某个元素 $[m]$ 的类，而该元素属于 $\Ker(\psi) \subset \colim_i M_i$。
选取 $i$，使 $[m]$
来自某个 $m \in M_i$。再选取 $j \geq i$，使
$\nu_{ij}(\psi_i(m)) = 0$；由于 $[m] \in \Ker(\psi)$，这样的指标
确实存在。以 $j$ 代替 $i$，并以 $\mu_{ij}(m)$ 代替 $m$，
即可设 $m \in \Ker(\psi_i)$。这证明 $\colim_i H_i \to H$ 为满射。

\medskip\noindent
设 $h_i \in H_i$ 在 $H$ 中的像为零。由于
$H_i = \Ker(\psi_i)/\Im(\varphi_i)$，可用元素
$m \in \Ker(\psi_i) \subset M_i$ 表示 $h_i$。
$h_i$ 在 $H$ 中消失的假设意味着 $m$ 在 $\colim_i M_i$ 中的类
属于 $\varphi$ 的像。因此，存在 $j \geq i$ 及 $l \in L_j$，使得
$\varphi_j(l) = \mu_{ij}(m)$。这显然表明 $h_i$ 在 $H_j$ 中的像为零，
从而证明 $\colim_i H_i \to H$ 为单射。
\end{proof}

\begin{example}
\label{algebra-example-colimit-not-exact}
一般而言，取余极限并不正合。考虑偏序集 $I = \{a, b, c\}$，
其中 $a < b$、$a < c$，且无其他严格不等式，正如例
\ref{algebra-example-zero-colimit-different} 中那样。考虑系统同态
$(0, \mathbf{Z}, \mathbf{Z}, 0, 0) \to
(\mathbf{Z}, \mathbf{Z}, \mathbf{Z}, 1, 1)$。
由例 \ref{algebra-example-zero-colimit-different} 中对余极限的描述可知，
相应的余极限同态不是单射，尽管该系统同态在每个对象上都是单射。
故引理 \ref{algebra-lemma-directed-colimit-exact} 的结论对一般系统不成立。
\end{example}

\begin{lemma}
\label{algebra-lemma-almost-directed-colimit-exact}
设 $\mathcal{I}$ 为满足范畴论，引理
\ref{categories-lemma-split-into-directed} 假设的指标范畴。
则在 $\mathcal{I}$ 上取 Abel 群图式的余极限是正合的（即引理
\ref{algebra-lemma-directed-colimit-exact} 的类似结论在此情形成立）。
\end{lemma}

\begin{proof}
由范畴论，引理 \ref{categories-lemma-split-into-directed}，
可写 $\mathcal{I} = \coprod_{j \in J} \mathcal{I}_j$，其中每个
$\mathcal{I}_j$ 都是滤过范畴，而 $J$ 可以为空。由范畴论，引理
\ref{categories-lemma-directed-category-system}，在指标范畴
$\mathcal{I}_j$ 上取余极限，等同于在某个有向集上取余极限。
因此，引理 \ref{algebra-lemma-directed-colimit-exact} 可用于这些余极限。
问题由此归结为证明：在 $R$-模范畴中，对集合 $J$ 取余积是正合的。
换言之，给定 $R$ 模的正合列 $L_j \to M_j \to N_j$，须证明
% P01-E0264: source wording candidate recorded in ERRATA_CANDIDATES.jsonl.
$$
\bigoplus\nolimits_{j \in J} L_j
\longrightarrow
\bigoplus\nolimits_{j \in J} M_j
\longrightarrow
\bigoplus\nolimits_{j \in J} N_j
$$
正合。这可直接验证，即使 $J$ 为空也成立。
\end{proof}
\section{局部化}
\label{algebra-section-localization}

\begin{definition}
\label{algebra-definition-multiplicative-subset}
设 $R$ 为环，$S$ 是 $R$ 的子集。若
$1\in S$，且 $S$ 对乘法封闭，即
$s, s' \in S \Rightarrow ss' \in S$，
则称 $S$ 为{\it $R$ 的乘法子集}。
\end{definition}

\noindent
给定环 $A$ 及其乘法子集 $S$，在 $A \times S$ 上定义如下关系：
$$
(x, s) \sim (y, t)
\Leftrightarrow
\exists u \in S \text{ 使得 } (xt-ys)u = 0
$$
容易验证这是等价关系。用 $x/s$（或 $\frac{x}{s}$）表示
$(x, s)$ 的等价类，并用 $S^{-1}A$ 表示所有等价类组成的集合。
在 $S^{-1}A$ 中定义加法与乘法如下：
$$
x/s + y/t = (xt + ys)/st, \quad
x/s \cdot y/t = xy/st
$$
可以验证，在这些运算下 $S^{-1}A$ 成为环。

\begin{definition}
\label{algebra-definition-localization}
此环称为{\it $A$ 关于 $S$ 的局部化}。
\end{definition}

\noindent
从 $A$ 到其局部化 $S^{-1}A$ 有自然环同态
$$
A \longrightarrow S^{-1}A, \quad x \longmapsto x/1
$$
有时称为{\it 局部化同态}。一般而言，局部化同态不一定是单射；
若 $S$ 不含零因子，则它是单射。事实上，若 $x/1 = 0$，
则存在 $u\in S$，使得 $A$ 中 $xu = 0$；由于 $S$ 中没有零因子，
必有 $x = 0$。环的局部化具有下列泛性质。

\begin{proposition}
\label{algebra-proposition-universal-property-localization}
设 $f : A \to B$ 为环同态，并把 $S$ 的每个元素映为 $B$ 中的单位。
则存在唯一同态 $g : S^{-1}A \to B$，使下图交换：
$$
\xymatrix{
A \ar[rr]^{f} \ar[dr] & & B \\
& S^{-1}A \ar[ur]_g
}
$$
\end{proposition}

\begin{proof}
存在性。如下定义映射 $g$。对 $x/s\in S^{-1}A$，令
$g(x/s) = f(x)f(s)^{-1}\in B$。由定义容易验证这是良定义的环同态，
并且显然使上图交换。

\medskip\noindent
唯一性。现设 $g' : S^{-1}A \to B$ 满足 $g'(x/1) = f(x)$，
证明 $g = g'$。由图的交换性，对 $s \in S$ 有
$f(s) = g'(s/1)$。而 $B$ 中 $g'(1/s)f(s) = 1$，
故 $g'(1/s) = f(s)^{-1}$，从而
$g'(x/s) = g'(x/1)g'(1/s) = f(x)f(s)^{-1} = g(x/s)$。
\end{proof}

\begin{lemma}
\label{algebra-lemma-localization-zero}
局部化 $S^{-1}A$ 是零环，当且仅当 $0\in S$。
\end{lemma}

\begin{proof}
若 $0\in S$，则按定义任意有序对 $(a, s)\sim (0, 1)$。
若 $0\not \in S$，则显然在 $S^{-1}A$ 中 $1/1 \neq 0/1$。
\end{proof}

\begin{lemma}
\label{algebra-lemma-localization-and-modules}
设 $R$ 为环，$S \subset R$ 为乘法子集。$S^{-1}R$-模范畴等价于
满足下列性质的 $R$-模 $N$ 所成范畴：每个 $s \in S$ 在 $N$ 上的作用
都是自同构。
\end{lemma}

\begin{proof}
给出该等价的函子把 $S^{-1}R$-模 $M$ 关联到同一个模，但通过局部化同态
$R \to S^{-1}R$ 把它视为 $R$-模。反之，若 $N$ 是 $R$-模，
且每个 $s \in S$ 都通过自同构 $s_N$ 作用，则可令 $x/s$ 通过
$x_N  \circ s_N^{-1}$ 作用，从而把 $N$ 视为 $S^{-1}R$-模。
这两个函子互为拟逆的验证从略。
\end{proof}

\noindent
环的局部化概念可推广为模的局部化。设 $A$ 为环，$S$ 是 $A$ 的乘法
子集，$M$ 是 $A$-模。在 $M \times S$ 上定义如下关系：
$$
(m, s) \sim (n, t)
\Leftrightarrow
\exists u\in S \text{ 使得 } (mt-ns)u = 0
$$
这显然是等价关系。
% P01-E0265: source wording candidate recorded in ERRATA_CANDIDATES.jsonl.
用 $m/s$（或 $\frac{m}{s}$）表示 $(m, s)$ 的等价类，并用
$S^{-1}M$ 表示所有等价类组成的集合。定义加法与标量乘法如下：
% P01-E0266: source formula candidate recorded in ERRATA_CANDIDATES.jsonl.
$$
m/s + n/t = (mt + ns)/st,\quad
m/s\cdot n/t = mn/st
$$
显然，这使 $S^{-1}M$ 成为 $S^{-1}A$-模。

\begin{definition}
\label{algebra-definition-localization-module}
$S^{-1}A$-模 $S^{-1}M$ 称为 $M$ 关于 $S$ 的{\it 局部化}。
\end{definition}

\noindent
注意，有时把 $A$-模同态 $M \to S^{-1}M$，$m \mapsto m/1$
称为{\it 局部化同态}。它满足下列泛性质。

\begin{lemma}
\label{algebra-lemma-universal-property-localization-module}
设 $R$ 为环。
% P01-E0267: source wording candidate recorded in ERRATA_CANDIDATES.jsonl.
设 $S \subset R$ 为乘法子集，$M$、$N$ 为 $R$-模。
假设 $S$ 的所有元素在 $N$ 上的作用都是自同构。则由局部化同态诱导的
典范映射
$$
\Hom_R(S^{-1}M, N) \longrightarrow \Hom_R(M, N)
$$
是同构。
\end{lemma}

\begin{proof}
该映射良定义且为 $R$-线性的事实是显然的。
单射性：设 $\alpha \in \Hom_R(S^{-1}M, N)$，并任取
$m/s \in S^{-1}M$。由于
$s \cdot \alpha(m/s) = \alpha(m/1)$，故
$ \alpha(m/s) =s^{-1}(\alpha (m/1))$；所以 $\alpha$ 完全由它在
$M$ 于 $S^{-1}M$ 中的像上的取值决定。
满射性：给定一个
% P01-E0268: source typography candidate recorded in ERRATA_CANDIDATES.jsonl.
R-线性映射 $\beta : M \rightarrow N$。须证明它可以“延拓”到
$S^{-1}M$。定义集合映射
$$
M \times S \rightarrow N,\quad
(m,s) \mapsto s^{-1}\beta(m)
$$
显然，该映射保持上述等价关系，故下降为良定义映射
$\alpha : S^{-1}M \rightarrow N$。还须证明它为 $R$-线性的。
取 $r, r' \in R$、$s, s' \in S$ 以及 $m, m' \in M$，则
\begin{align*}
\alpha(r \cdot m/s + r' \cdot m' /s')
& =
\alpha((r \cdot s' \cdot m + r' \cdot s \cdot m') /(ss')) \\
& =
(ss')^{-1}\beta(r \cdot s' \cdot m + r' \cdot s \cdot m') \\
& =
(ss')^{-1} (r \cdot s' \beta (m) + r' \cdot s \beta (m')) \\
& =
r \alpha (m/s) + r' \alpha (m' /s')
\end{align*}
所需结论成立。
\end{proof}

\begin{example}
\label{algebra-example-localize-at-prime}
设 $A$ 为环，$M$ 为 $A$-模。以下是几个重要的局部化例子。
\begin{enumerate}
\item 给定 $A$ 的素理想 $\mathfrak p$，考虑
$S = A\setminus\mathfrak p$。立即可知 $S$ 是乘法子集。在此情形，
分别用 $A_\mathfrak p$ 和 $M_\mathfrak p$ 表示 $A$ 和 $M$ 关于 $S$
的局部化，称为{\it $A$、相应地 $M$ 在 $\mathfrak p$ 处的局部化}。
\item 设 $f\in A$。考虑 $S = \{1, f, f^2, \ldots\}$。
这显然是 $A$ 的乘法子集。此时分别用 $A_f$（相应地 $M_f$）表示局部化
$S^{-1}A$（相应地 $S^{-1}M$），称为
{\it $A$、相应地 $M$ 关于 $f$ 的局部化}。
注意，$A_f = 0$ 当且仅当 $f$ 在 $A$ 中是幂零元。
\item 令 $S = \{f \in A \mid f \text{ is not a zerodivisor in }A\}$。
这是 $A$ 的乘法子集。此时，环 $Q(A) = S^{-1}A$ 称为 $A$ 的
{\it 全商环}，也称为{\it 全分式环}。
\item 若 $A$ 是整环，则全商环 $Q(A)$ 就是 $A$ 的分式域。请参见
域，例 \ref{fields-example-quotient-field}。
\end{enumerate}
\end{example}

\begin{lemma}
\label{algebra-lemma-localization-colimit}
设 $R$ 为环，$S \subset R$ 为乘法子集，$M$ 为 $R$-模。则
$$
S^{-1}M = \colim_{f \in S} M_f
$$
其中 $S$ 上的预序由下式给出：
$f \geq f' \Leftrightarrow f = f'f''$，其中某个 $f'' \in R$；
此时映射 $M_{f'} \to M_f$ 由
$m/(f')^e \mapsto m(f'')^e/f^e$ 给出。
\end{lemma}

\begin{proof}
从略。提示：使用引理
\ref{algebra-lemma-universal-property-localization-module} 的泛性质。
\end{proof}

\noindent
以下一段中，设 $A$ 为环，$M, N$ 为 $A$ 上的模。

\medskip\noindent
若 $S$、$S'$ 是 $A$ 的乘法子集，则显然
$$
SS' = \{ss' : s\in S, \ s'\in S'\}
$$
也是 $A$ 的乘法子集。于是有下列结论。

\begin{proposition}
\label{algebra-proposition-localize-twice}
设 $\overline{S}$ 为 $S$ 在 $S'^{-1}A$ 中的像，则
$(SS')^{-1}A$ 同构于 $\overline{S}^{-1}(S'^{-1}A)$。
\end{proposition}

\begin{proof}
把 $x\in A$ 映到 $x/1\in (SS')^{-1}A$ 的映射，由泛性质诱导出
把 $x/s\in S'^{-1}A$ 映到 $x/s \in (SS')^{-1}A$ 的映射。
$\overline{S}$ 中元素的像在 $(SS')^{-1}A$ 中可逆。再次由泛性质，
得到映射
$f : \overline{S}^{-1}(S'^{-1}A) \to (SS')^{-1}A$，
它把 $(x/s')/(s/1)$ 映到 $x/ss'$。

\medskip\noindent
另一方面，把 $x\in A$ 映到 $(x/1)/(1/1)$ 的、从 $A$ 到
$\overline{S}^{-1}(S'^{-1}A)$ 的映射，再次由泛性质诱导出映射
$g : (SS')^{-1}A \to \overline{S}^{-1}(S'^{-1}A)$，
它把 $x/ss'$ 映到 $(x/s')/(s/1)$。
立即可验证 $f$ 与 $g$ 互为逆映射，故二者都是同构。
\end{proof}

\noindent
对于模 $M$，有下列结论。

\begin{proposition}
\label{algebra-proposition-localize-twice-module}
把 $S'^{-1}M$ 视为 $A$-模，则 $S^{-1}(S'^{-1}M)$ 同构于
$(SS')^{-1}M$。
\end{proposition}

\begin{proof}
% P01-E0269: source wording candidate recorded in ERRATA_CANDIDATES.jsonl.
注意，对于给定的 $A$-模 M，我们尚未证明 $S^{-1}M$ 的任何泛性质。
因此不能像前一个证明那样论证，而必须显式构造此同构。

\medskip\noindent
如下定义映射：
\begin{align*}
& f : S^{-1}(S'^{-1}M) \longrightarrow (SS')^{-1}M, \quad \frac{x/s'}{s}\mapsto
x/ss'\\
& g : (SS')^{-1}M \longrightarrow S^{-1}(S'^{-1}M), \quad x/t\mapsto
\frac{x/s'}{s}\ \text{对某个 }s\in S, s'\in S', \text{ 且 }
t = ss'
\end{align*}
必须验证这些同态良定义，即不依赖于分式的选择。
% P01-E0270: source wording candidate recorded in ERRATA_CANDIDATES.jsonl.
这容易验证；证明它们互为逆映射也很直接。
\end{proof}

\noindent
% P01-E0271: source precision candidate recorded in ERRATA_CANDIDATES.jsonl.
若 $u : M \to N$ 是 $A$ 同态，则局部化确实诱导出良定义的
$S^{-1}A$ 同态 $S^{-1}u : S^{-1}M \to S^{-1}N$，把 $x/s$ 映到
$u(x)/s$。立即可验证此构造具有函子性，故 $S^{-1}$ 实际上是从
$A$-模范畴到 $S^{-1}A$-模范畴的函子。此外，该函子是正合的，
如下命题所示。

\begin{proposition}
\label{algebra-proposition-localization-exact}
\begin{slogan}
局部化是正合的。
\end{slogan}
% P01-E0272: source base-ring candidate recorded in ERRATA_CANDIDATES.jsonl.
设 $L\xrightarrow{u} M\xrightarrow{v} N$ 为 $R$-模的正合列。则
$S^{-1}L \to S^{-1}M \to S^{-1}N$ 也正合。
\end{proposition}

\begin{proof}
首先，由于局部化是函子，显然
$S^{-1}L \to S^{-1}M \to S^{-1}N$ 是复形。其次，设 $x/s$
在 $S^{-1}N$ 中的像为零，其中 $x/s \in S^{-1}M$。则按定义，存在
$t\in S$，使得 $N$ 中 $v(xt) = v(x)t = 0$，亦即
$xt \in \Ker(v)$。由 $L \to M \to N$ 的正合性，存在 $L$ 中的某个
$y$，使 $xt = u(y)$。于是 $x/s$ 是 $y/st$ 的像。这就证明了正合性。
\end{proof}

\begin{lemma}
\label{algebra-lemma-localize-quotient-modules}
局部化保持商；即若 $N$ 是 $M$ 的子模，则
$S^{-1}(M/N)\simeq (S^{-1}M)/(S^{-1}N)$。
\end{lemma}

\begin{proof}
由正合列
$$
0 \longrightarrow N \longrightarrow M \longrightarrow M/N \longrightarrow 0
$$
得到
$$
0 \longrightarrow S^{-1}N \longrightarrow S^{-1}M
\longrightarrow S^{-1}(M/N) \longrightarrow 0
$$
% P01-E0273: source result-type candidate recorded in ERRATA_CANDIDATES.jsonl.
引理随即成立。
\end{proof}

\noindent
% P01-E0274: source result-type candidate recorded in ERRATA_CANDIDATES.jsonl.
在前一引理中，对 $A$ 的理想 $I$ 取 $N = I$、$M = A$，
便知作为 $A$-模有
$S^{-1}A/S^{-1}I \simeq S^{-1}(A/I)$。下一个命题表明它们作为环也同构。

\begin{proposition}
\label{algebra-proposition-localize-quotient}
设 $I$ 是 $A$ 的理想，$S$ 是 $A$ 的乘法子集。则
$S^{-1}I$ 是 $S^{-1}A$ 的理想，且
$\overline{S}^{-1}(A/I)$ 同构于 $S^{-1}A/S^{-1}I$，其中
$\overline{S}$ 是 $S$ 在 $A/I$ 中的像。
\end{proposition}

\begin{proof}
由于 $I$ 本身是理想，$S^{-1}I$ 为理想的事实是显然的。定义
$$
f : S^{-1}A\longrightarrow \overline{S}^{-1}(A/I), \quad x/s\mapsto
\overline{x}/\overline{s}
$$
其中 $\overline{x}$ 与 $\overline{s}$ 分别是 $x$ 与 $s$ 在
$A/I$ 中的像。本证明中继续使用类似记号。由 $S^{-1}A$ 的泛性质，
该映射良定义；其核包含 $S^{-1}I$，故它诱导映射
$$
\overline{f} : S^{-1}A/S^{-1}I \longrightarrow \overline{S}^{-1}(A/I), \quad
\overline{x/s}\mapsto \overline{x}/\overline{s}
$$

\medskip\noindent
另一方面，把 $x$ 映到 $\overline{x/1}$ 的映射
$A \to S^{-1}A/S^{-1}I$ 诱导出把 $\overline{x}$ 映到
$\overline{x/1}$ 的映射 $A/I \to S^{-1}A/S^{-1}I$。
$\overline{S}$ 的像在 $S^{-1}A/S^{-1}I$ 中可逆，故由泛性质诱导映射
$$
g : \overline{S}^{-1}(A/I) \longrightarrow S^{-1}A/S^{-1}I, \quad
\frac{\overline{x}}{\overline{s}}\mapsto \overline{x/s}
$$
显然，$\overline{f}$ 与 $g$ 互为逆映射，故二者都是同构。
\end{proof}

\noindent
现在考察子模在局部化下的行为。

\begin{lemma}
\label{algebra-lemma-submodule-localization}
$S^{-1}M$ 的任一子模 $N'$ 都具有形式 $S^{-1}N$，其中
$N\subset M$；事实上，可以取 $N$ 为 $N'$ 在 $M$ 中的逆像。
\end{lemma}

\begin{proof}
设 $N$ 为 $N'$ 在 $M$ 中的逆像。于是可见
% P01-E0275: source inclusion-direction candidate recorded in ERRATA_CANDIDATES.jsonl.
$S^{-1}N\supset N'$。为证明二者相等，取 $S^{-1}N$ 中的
$x/s$，其中 $s\in S$ 且 $x\in N$。于是 $x/1\in N'$。
% P01-E0276: source base-ring candidate recorded in ERRATA_CANDIDATES.jsonl.
由于 $N'$ 是 $S^{-1}R$-子模，有
$x/s = x/1\cdot 1/s\in N'$。证明完毕。
\end{proof}

\noindent
令 $M = A$，并令 $N = I$ 为 $A$ 的理想，得到下列推论；
它可以视为命题 \ref{algebra-proposition-localize-quotient} 第一部分的逆命题。

\begin{lemma}
\label{algebra-lemma-ideal-in-localization}
\begin{slogan}
环之局部化中的理想都是理想的局部化。
\end{slogan}
$S^{-1}A$ 的每个理想 $I'$ 都具有形式 $S^{-1}I$；
可以取 $I$ 为 $I'$ 在 $A$ 中的逆像。
\end{lemma}

\begin{proof}
立即由引理 \ref{algebra-lemma-submodule-localization} 得到。
\end{proof}

\section{内 Hom}
\label{algebra-section-hom}

\noindent
若 $R$ 是环，$M$、$N$ 是 $R$-模，则
$$
\Hom_R(M, N) = \{ \varphi : M \to N\}
$$
是从 $M$ 到 $N$ 的 $R$-线性映射所成的集合。照常令
$(\varphi + \psi)(m) = \varphi(m) + \psi(m)$，便赋予此集合 Abel 群
结构。事实上，$\Hom_R(M, N)$ 还按规则
$(x \varphi)(m) = x \varphi(m) = \varphi(xm)$ 成为 $R$-模。

\medskip\noindent
给定 $R$-模同态 $a : M \to M'$ 与 $b : N \to N'$，
可分别用 $a$ 与 $b$ 对同态作前复合及后复合。由此得到交换图
$$
\xymatrix{
\Hom_R(M', N) \ar[d]_{- \circ a} \ar[r]_{b \circ -} &
\Hom_R(M', N') \ar[d]^{- \circ a} \\
\Hom_R(M, N) \ar[r]^{b \circ -} &
\Hom_R(M, N')
}
$$
事实上，图中的映射都是 $R$-模同态。因此 $\Hom_R$ 定义了加性函子
$$
\text{Mod}_R^{opp} \times \text{Mod}_R \longrightarrow \text{Mod}_R, \quad
(M, N) \longmapsto \Hom_R(M, N)
$$

\begin{lemma}
\label{algebra-lemma-hom-exact}
正合性与 $\Hom_R$。设 $R$ 为环，$M_1$、$M_2$、$M_3$ 为
$R$-模，并设 $M_1 \to M_2$ 与 $M_2 \to M_3$ 为 $R$-模同态。
\begin{enumerate}
\item $M_1 \to M_2 \to M_3 \to 0$ 正合，当且仅当对所有 $R$-模 $N$，
$0 \to \Hom_R(M_3, N) \to \Hom_R(M_2, N) \to \Hom_R(M_1, N)$
正合。
\item $0 \to M_1 \to M_2 \to M_3$ 正合，当且仅当对所有 $R$-模 $N$，
$$
0 \to \Hom_R(N, M_1) \to \Hom_R(N, M_2) \to \Hom_R(N, M_3)
$$
正合。
\end{enumerate}
\end{lemma}

\begin{proof}
从略。
\end{proof}

\begin{lemma}
\label{algebra-lemma-hom-from-finitely-presented}
设 $R$ 为环，$M$ 为有限表示 $R$-模，$N$ 为 $R$-模。
\begin{enumerate}
\item 对 $f \in R$，有
$\Hom_R(M, N)_f = \Hom_{R_f}(M_f, N_f) = \Hom_R(M_f, N_f)$；
\item 对 $R$ 的乘法子集 $S$，有
$$
S^{-1}\Hom_R(M, N) = \Hom_{S^{-1}R}(S^{-1}M, S^{-1}N) =
\Hom_R(S^{-1}M, S^{-1}N).
$$
\end{enumerate}
\end{lemma}

\begin{proof}
(1) 是 (2) 的特殊情形。(2) 中第二个等号由引理
\ref{algebra-lemma-localization-and-modules} 得到。选取一个表示
$$
\bigoplus\nolimits_{j = 1, \ldots, m} R
\longrightarrow
\bigoplus\nolimits_{i = 1, \ldots, n} R
\to M \to 0.
$$
由引理 \ref{algebra-lemma-hom-exact}，得到正合列
$$
0 \to
\Hom_R(M, N) \to
\bigoplus\nolimits_{i = 1, \ldots, n} N
\longrightarrow
\bigoplus\nolimits_{j = 1, \ldots, m} N.
$$
关于 $S$ 作局部化并使用命题 \ref{algebra-proposition-localization-exact}，
得到正合列
$$
0 \to
S^{-1}\Hom_R(M, N) \to
\bigoplus\nolimits_{i = 1, \ldots, n} S^{-1}N
\longrightarrow
\bigoplus\nolimits_{j = 1, \ldots, m} S^{-1}N
$$
又因 $S^{-1}M$ 位于正合列
$$
\bigoplus\nolimits_{j = 1, \ldots, m} S^{-1}R
\longrightarrow
\bigoplus\nolimits_{i = 1, \ldots, n} S^{-1}R \to S^{-1}M \to 0
$$
中，该列由引理 \ref{algebra-lemma-hom-exact} 诱导出正合列
$$
0 \to
\Hom_{S^{-1}R}(S^{-1}M, S^{-1}N) \to
\bigoplus\nolimits_{i = 1, \ldots, n} S^{-1}N
\longrightarrow
\bigoplus\nolimits_{j = 1, \ldots, m} S^{-1}N
$$
而这正是上面的正合列，故结论成立。
\end{proof}

\section{有限模与有限表示模的刻画}
\label{algebra-section-colim-and-hom}

\noindent
给定环 $R$ 上的模 $N$，可以用函子 $\Hom_R(N, -)$ 来刻画
$N$ 是否为有限模或有限表示模。

\begin{lemma}
\label{algebra-lemma-characterize-finite-module-hom}
设 $R$ 为环，$N$ 为 $R$-模。下列条件等价：
\begin{enumerate}
\item $N$ 是有限 $R$-模；
\item 对 $R$-模的任意滤过余极限 $M = \colim M_i$，映射
$\colim \Hom_R(N, M_i) \to \Hom_R(N, M)$ 是单射。
\end{enumerate}
\end{lemma}

\begin{proof}
设 (1) 成立，并为 $N$ 选取生成元 $x_1, \ldots, x_m$。
若 $N \to M_i$ 是模同态，且复合
$N \to M_i \to M$ 为零，则由于
$M = \colim_{i' \geq i} M_{i'}$，对每个
$j \in \{1, \ldots, m\}$，都能找到
% P01-E0277: source article candidate recorded in ERRATA_CANDIDATES.jsonl.
某个 $i' \geq i$，使 $x_j$ 在 $M_{i'}$ 中的像为零。
由于 $x_j$ 只有有限多个，可以找到一个同时适用于所有生成元的 $i'$。
于是复合 $N \to M_i \to M_{i'}$ 为零，从而该映射是单射，即 (2) 成立。

\medskip\noindent
设 (2) 成立。对有限子集 $E \subset N$，用 $N_E \subset N$
表示由 $E$ 中元素生成的 $R$-子模。则
$0 = \colim N/N_E$ 是滤过余极限。因此，
% P01-E0278: source factorization wording candidate recorded in ERRATA_CANDIDATES.jsonl.
$\text{id} : N \to N$ 对某个 $E$ 通过 $N_E$ 分解；换言之，$N$
是有限生成的。
\end{proof}

\noindent
为便于引用，现定义模中元素之间的关系。

\begin{definition}
\label{algebra-definition-relation}
设 $R$ 为环，$M$ 为 $R$-模。设 $n \geq 0$，并对
$i = 1, \ldots, n$ 给定 $x_i \in M$。$M$ 中
$x_1, \ldots, x_n$ 之间的一个{\it 关系}，是满足
$\sum_{i = 1, \ldots, n} f_i x_i = 0$ 的元素列
$f_1, \ldots, f_n \in R$。
\end{definition}

\begin{lemma}
\label{algebra-lemma-module-colimit-fp}
设 $R$ 为环，$M$ 为 $R$-模。则 $M$ 是某个 $R$-模有向系统
$(M_i, \mu_{ij})$ 的余极限，其中每个 $M_i$ 都是有限表示 $R$-模。
\end{lemma}

\begin{proof}
考虑任意有限子集 $S \subset M$，以及 $S$ 中元素之间的任意有限关系集
$E$。于是，每个 $s \in S$ 对应于 $x_s \in M$，而每个
$e \in E$ 由元素向量 $f_{e, s} \in R$ 组成，满足
$\sum f_{e, s} x_s = 0$。令 $M_{S, E}$ 为映射
$$
R^{\# E} \longrightarrow R^{\# S}, \quad
(g_e)_{e\in E} \longmapsto (\sum g_e f_{e, s})_{s\in S}.
$$
的余核。有典范映射 $M_{S, E} \to M$。若 $S \subset S'$，
且 $E$ 中元素通过此映射对应于 $E'$ 中的关系，则有显然的映射
$M_{S, E} \to M_{S', E'}$，并且它与到 $M$ 的映射交换。
令 $I$ 为所有有序对 $(S, E)$ 组成的集合，并按上述包含关系排序。
显然，这个有向系统的余极限是 $M$。
\end{proof}

\begin{lemma}
\label{algebra-lemma-characterize-finitely-presented-module-hom}
设 $R$ 为环，$N$ 为 $R$-模。下列条件等价：
\begin{enumerate}
\item $N$ 是有限表示 $R$-模；
\item 对 $R$-模的任意滤过余极限 $M = \colim M_i$，映射
$\colim \Hom_R(N, M_i) \to \Hom_R(N, M)$ 是双射。
\end{enumerate}
\end{lemma}

\begin{proof}
设 (1) 成立，并选取正合列 $F_{-1} \to F_0 \to N \to 0$，
其中 $F_i$ 是有限自由模。于是有关于 $R$-模 $M$ 函子性的正合列
$$
0 \to \Hom_R(N, M) \to \Hom_R(F_0, M) \to \Hom_R(F_{-1}, M)
$$
。函子 $\Hom_R(F_i, M)$ 与滤过余极限交换，因为
$\Hom_R(R^{\oplus n}, M) = M^{\oplus n}$。
由于滤过余极限是正合的（引理 \ref{algebra-lemma-directed-colimit-exact}），
可知 (2) 成立。

\medskip\noindent
设 (2) 成立。由引理 \ref{algebra-lemma-module-colimit-fp}，
可把 $N = \colim N_i$ 写成滤过余极限，其中对所有 $i$，$N_i$
都是有限表示模。因此，对某个 $i$，$\text{id}_N$ 通过 $N_i$ 分解。
这意味着 $N$ 是某个有限表示 $R$-模（即 $N_i$）的直和因子，
故其也是有限表示模。
\end{proof}

\section{张量积}
\label{algebra-section-tensor-product}

\begin{definition}
\label{algebra-definition-bilinear}
设 $R$ 为环，$M, N, P$ 为三个 $R$-模。称映射
$f : M \times N \to P$（其中 $M \times N$ 仅视为两个 $R$-模的笛卡尔
积）是{\it $R$-双线性的}，如果对每个 $x \in M$，从 $N$ 到 $P$ 的映射
$y\mapsto f(x, y)$ 都是 $R$-线性的，并且对每个 $y\in N$，映射
$x\mapsto f(x, y)$ 也都是 $R$-线性的。
\end{definition}

\begin{lemma}
\label{algebra-lemma-tensor-product}
设 $M, N$ 为 $R$-模。则存在一对 $(T, g)$，其中 $T$ 是 $R$-模，
$g : M \times N \to T$ 是 $R$-双线性映射，并具有下列泛性质：
对任意 $R$-模 $P$ 及任意 $R$-双线性映射
$f : M \times N \to P$，存在唯一的 $R$-线性映射
$\tilde{f} : T \to P$，使得 $f = \tilde{f} \circ g$。
换言之，下图交换：
$$
\xymatrix{
M \times N \ar[rr]^f \ar[dr]_g & & P\\
& T \ar[ur]_{\tilde f}
}
$$
此外，若 $(T, g)$ 与 $(T', g')$ 是两对具有此性质的对象，
则存在唯一同构 $j : T \to T'$，使 $j\circ g = g'$。
\end{lemma}

\noindent
满足上述泛性质的 $R$-模 $T$ 称为 $R$-模 $M$ 与 $N$ 的
{\it 张量积}，记为 $M \otimes_R N$。

\begin{proof}
先证明这样的 $R$-模 $T$ 存在。设 $M, N$ 为 $R$-模。
令 $T$ 为商模 $P/Q$，其中 $P$ 是自由 $R$-模
$R^{(M \times N)}$，而 $Q$ 是由下列各类元素生成的 $R$-模
（$x\in M, y\in N$）：
\begin{align*}
(x + x', y) - (x, y) - (x', y), \\
(x, y + y') - (x, y) - (x, y'), \\
(ax, y) - a(x, y), \\
(x, ay) - a(x, y)
\end{align*}
用 $\pi : M \times N \to T$ 表示自然映射。验证双线性条件时，
上述关系即说明该映射为 $R$-双线性的。
% P01-E0280: source wording candidate recorded in ERRATA_CANDIDATES.jsonl.
把像记为 $\pi(x, y) = x \otimes y$；这些元素生成 $T$。
现设 $f : M \times N \to P$ 为 $R$-双线性映射，则可把映射
$f'(x \otimes y) = f(x, y)$ 线性延拓，定义
$f' : T \to P$。显然 $f = f'\circ \pi$。此外，$f'$ 由其在生成集
$\{x \otimes y : x\in M, y\in N\}$ 上的取值唯一确定。
设另有一对 $(T', g')$ 满足同样性质。则存在唯一的
$j : T \to T'$ 以及 $j' : T' \to T$，使
$g' = j\circ g$、$g = j'\circ g'$。
% P01-E0279: source composition-order defect recorded in ERRATA_CANDIDATES.jsonl.
但映射 $(j\circ j') \circ g$ 与 $g$ 都满足泛性质，故由唯一性二者相等，
从而 $j'\circ j$ 是 $T$ 上的恒等映射。类似地，
$(j'\circ j) \circ g' = g'$，且 $j\circ j'$ 是 $T'$ 上的恒等映射。
所以 $j$ 是同构。
\end{proof}

\begin{lemma}
\label{algebra-lemma-flip-tensor-product}
设 $M, N, P$ 为 $R$-模，则双线性映射
\begin{align*}
(x, y) & \mapsto y \otimes x\\
(x + y, z) & \mapsto x \otimes z + y \otimes z\\
(r, x) & \mapsto rx
\end{align*}
分别诱导唯一同构
\begin{align*}
M \otimes_R N & \to N \otimes_R M, \\
(M\oplus N)\otimes_R P & \to (M \otimes_R P)\oplus(N \otimes_R P),  \\
R \otimes_R M & \to M
\end{align*}
\end{lemma}

\begin{proof}
从略。
\end{proof}

\noindent
可以把两个 $R$-模的张量积推广到有限多个 $R$-模，并在多重张量积与
多线性映射之间建立对应。用几乎相同的构造，可以证明下列结论。

\begin{lemma}
\label{algebra-lemma-multilinear}
设 $M_1, \ldots, M_r$ 为 $R$-模。则存在一对 $(T, g)$，
由一个 $R$-模 T 及一个具有如下泛性质的 $R$-多线性映射
$g : M_1\times \ldots \times M_r \to T$ 组成：
对任意 $R$-多线性映射 $f : M_1\times \ldots \times M_r \to P$，
存在唯一的 $R$-模同态 $f' : T \to P$，使 $f'\circ g = f$。
这样的模 $T$ 在唯一同构意义下唯一。把它记为
$M_1\otimes_R \ldots \otimes_R M_r$，并把泛多线性映射记为
$(m_1, \ldots, m_r) \mapsto m_1 \otimes \ldots \otimes m_r$。
\end{lemma}

\begin{proof}
从略。
\end{proof}

\begin{lemma}
\label{algebra-lemma-transitive}
同态
$$
(M \otimes_R N)\otimes_R P \to
M \otimes_R N \otimes_R P \to
M \otimes_R (N \otimes_R P)
$$
由下式给出：
$f((x \otimes y)\otimes z) = x \otimes y \otimes z$ 以及
$g(x \otimes y \otimes z) = x \otimes (y \otimes z)$，
其中 $x\in M, y\in N, z\in P$。这些同态良定义且为同构。
\end{lemma}

\begin{proof}
证明 $f$ 良定义且为同构；对 $g$ 的证明完全类似。固定任意
$z\in P$，则映射 $(x, y)\mapsto x \otimes y \otimes z$，
$x\in M, y\in N$，关于 $x$、$y$ 是 $R$-双线性的，因而诱导同态
$f_z : M \otimes N \to M \otimes N \otimes P$，满足
$f_z(x \otimes y) = x \otimes y \otimes z$。
再考虑由 $(w, z)\mapsto f_z(w)$ 给出的映射
$(M \otimes N)\times P \to M \otimes N \otimes P$。
此映射为 $R$-双线性的，故诱导
$f : (M \otimes_R N)\otimes_R P \to M \otimes_R N \otimes_R P$，
且 $f((x \otimes y)\otimes z) = x \otimes y \otimes z$。
为构造其逆映射，注意映射
$\pi : M \times N \times P \to (M \otimes N)\otimes P$
是 $R$-三线性的。因此，它诱导一个符合泛性质的 $R$-线性映射
$h : M \otimes N \otimes P \to (M \otimes N)\otimes P$。
在此可见，$h(x \otimes y \otimes z) = (x \otimes y)\otimes z$。
由 $f$ 与 $h$ 的显式表达式，$f\circ h$ 与 $h\circ f$ 分别是
$M \otimes N \otimes P$ 与 $(M \otimes N)\otimes P$ 上的恒等映射，
故 $f$ 即所求同构。
\end{proof}

\noindent
归纳可见，这一结论推广到多重张量积。再结合引理
\ref{algebra-lemma-flip-tensor-product}，可知 $R$-模范畴上的张量积运算
具有结合性、交换性与分配性。

\begin{definition}
\label{algebra-definition-bimodule}
若 Abel 群 $N$ 既是 $A$-模又是 $B$-模，且对所有 $a \in A$、
$b \in B$，乘以 $a$ 与乘以 $b$ 的作用交换，即对所有
$n \in N$ 有 $b(an) = a(bn)$，则称之为
{\it $(A, B)$-双模}。在这种情形下，通常把 $B$-作用写在右边：
对 $b \in B$ 与 $n \in N$，用 $nb$ 表示 $n$ 乘以 $b$ 的结果。
采用此约定，上述相容性即为：对所有 $a\in A, b\in B, x\in N$，
有 $(ax)b = a(xb)$。用简记 $_AN_B$ 表示 $(A, B)$-双模 $N$。
\end{definition}

\begin{lemma}
\label{algebra-lemma-tensor-with-bimodule}
设 $M$ 为 $A$-模，$P$ 为 $B$-模，$N$ 为 $(A, B)$-双模。
则模 $(M \otimes_A N)\otimes_B P$ 与
$M \otimes_A(N \otimes_B P)$ 都可以赋予 $(A, B)$-双模结构，而且
$$
(M \otimes_A N)\otimes_B P \cong M \otimes_A(N \otimes_B P).
$$
\end{lemma}

\begin{proof}
先验地，$M \otimes_A N$ 是 $A$-模；但可以用规则
$$
(x \otimes y)b = x \otimes yb, \quad x\in M, y\in N, b\in B
$$
赋予它 $B$-模结构。于是 $M \otimes_A N$ 成为 $(A, B)$-双模。
$N \otimes_B P$ 的情形类似，进而
$(M \otimes_A N)\otimes_B P$ 与 $M \otimes_A(N \otimes_B P)$
也同样如此。由引理 \ref{algebra-lemma-transitive}，这两个模通过同一个映射，
% P01-E0281: source duplicated-word candidate recorded in ERRATA_CANDIDATES.jsonl.
既作为 $A$-模也作为 $B$-模同构。
\end{proof}

\begin{lemma}
\label{algebra-lemma-hom-from-tensor-product}
对任意三个 $R$-模 $M, N, P$，有
$$
\Hom_R(M \otimes_R N, P) \cong \Hom_R(M, \Hom_R(N, P))
$$
\end{lemma}

\begin{proof}
$R$-线性映射 $\hat{f}\in \Hom_R(M \otimes_R N, P)$ 对应于
$R$-双线性映射 $f : M \times N \to P$。对每个 $x\in M$，
由泛性质，映射 $y\mapsto f(x, y)$ 是 $R$-线性的。因此，$f$
对应于映射 $\phi_f : M \to \Hom_R(N, P)$。该映射为 $R$-线性的，
因为
$$
\phi_f(ax + y)(z) =
f(ax + y, z) = af(x, z)+f(y, z) =
(a\phi_f(x)+\phi_f(y))(z),
$$
其中任意 $a \in R$、$x \in M$、$y \in M$ 及 $z \in N$。
反之，任意 $f \in \Hom_R(M, \Hom_R(N, P))$ 定义一个 $R$-双线性映射
$M \times N \to P$，即 $(x, y)\mapsto f(x)(y)$。
故这给出两个模 $\Hom_R(M \otimes_R N, P)$ 与
$\Hom_R(M, \Hom_R(N, P))$ 之间的自然一一对应。
\end{proof}

\begin{lemma}[张量积与余极限交换]
\label{algebra-lemma-tensor-products-commute-with-limits}
设 $(M_i, \mu_{ij})$ 是预序集 $I$ 上的系统，$N$ 是 $R$-模。则
$$
\colim (M_i \otimes N) \cong (\colim M_i)\otimes N.
$$
此外，此同构由同态
$\mu_i \otimes 1: M_i \otimes N \to M \otimes N$
诱导，其中 $M = \colim_i M_i$，自然映射为 $\mu_i : M_i \to M$。
\end{lemma}

\begin{proof}
第一种证明。由引理 \ref{algebra-lemma-hom-from-tensor-product}，
函子 $M' \mapsto M' \otimes_R N$ 是函子
$N' \mapsto \Hom_R(N, N')$ 的左伴随。因此，
$M' \mapsto M' \otimes_R N$ 与所有余极限交换；见范畴论，引理
\ref{categories-lemma-adjoint-exact}。

\medskip\noindent
第二种直接证明。令 $P = \colim (M_i \otimes N)$，其余投影为
$\lambda_i : M_i \otimes N \to P$；令 $M = \colim M_i$，其余投影为
$\mu_i : M_i \to M$。则对所有 $i\leq j$，下图交换：
$$
\xymatrix{
M_i \otimes N \ar[r]_{\mu_i \otimes 1} \ar[d]_{\mu_{ij} \otimes 1} &
M \otimes N \ar[d]^{\text{id}} \\
M_j \otimes N \ar[r]^{\mu_j \otimes 1} &
M \otimes N
}
$$
由引理 \ref{algebra-lemma-homomorphism-limit}，这些映射诱导唯一同态
$\psi : P \to M \otimes N$，使
$\mu_i \otimes 1 = \psi \circ \lambda_i$。

\medskip\noindent
为构造逆映射，对每个 $i\in I$，都有典范的 $R$-双线性映射
$g_i : M_i \times N \to M_i \otimes N$。这诱导唯一映射
$\widehat{\phi} : M \times N \to P$，使
$\widehat{\phi} \circ (\mu_i \times 1) = \lambda_i \circ g_i$。
它为 $R$-双线性的，故诱导 $R$-线性映射
$\phi : M \otimes N \to P$。由下列交换图
$$
\xymatrix{
M_i \times N \ar[r]^{g_i} \ar[d]^{\mu_i \times \text{id}} &
M_i \otimes N\ar[r]_{\text{id}} \ar[d]_{\lambda_i} &
M_i \otimes N \ar[d]_{\mu_i \otimes \text{id}} \ar[rd]^{\lambda_i} \\
M \times N \ar[r]^{\widehat{\phi}} &
P \ar[r]^{\psi} & M \otimes N \ar[r]^{\phi} & P
}
$$
可知 $\psi\circ\widehat{\phi} = g$，其中 g 是典范 $R$-双线性映射
$g : M \times N \to M \otimes N$。故 $\psi\circ\phi$ 是
$M \otimes N$ 上的恒等映射。由右侧方形与三角形，
$\phi\circ\psi$ 也是 $P$ 上的恒等映射。
\end{proof}

\begin{lemma}
\label{algebra-lemma-tensor-product-exact}
设
\begin{align*}
M_1\xrightarrow{f} M_2\xrightarrow{g} M_3 \to 0
\end{align*}
为 $R$-模及同态的正合列，并设 $N$ 为任意 $R$-模。则序列
\begin{equation}
\label{algebra-equation-2ndex}
M_1\otimes N\xrightarrow{f \otimes 1} M_2\otimes N \xrightarrow{g \otimes 1}
M_3\otimes N \to 0
\end{equation}
正合。换言之，函子 $- \otimes_R N$ 是{\it 右正合的}；也就是说，
把原右正合列中的每一项与之作张量积会保持正合性。
\end{lemma}

\begin{proof}
对每个 $R$-模 $P$，向第一个正合列应用函子
$\Hom(-, \Hom(N, P))$，得到
$$
0 \to
\Hom(M_3, \Hom(N, P)) \to
\Hom(M_2, \Hom(N, P)) \to
\Hom(M_1, \Hom(N, P))
$$
它由引理 \ref{algebra-lemma-hom-exact} (1) 正合。
由引理 \ref{algebra-lemma-hom-from-tensor-product}，这变为序列
$$
0 \to \Hom(M_3 \otimes N, P) \to
\Hom(M_2 \otimes N, P) \to \Hom(M_1 \otimes N, P)
$$
因此也正合。再次使用引理 \ref{algebra-lemma-hom-exact} (1)，
即得到所求正合列。
\end{proof}

\begin{remark}
\label{algebra-remark-tensor-product-not-exact}
然而，张量积一般并不保持正合列。换言之，若
$M_1 \to M_2 \to M_3$ 正合，则对任意 $R$-模 $N$，
$M_1 \otimes N \to M_2 \otimes N \to M_3 \otimes N$
未必正合。
\end{remark}



\begin{example}
\label{algebra-example-tensor-product-not-exact}
考虑单射 $2 : \mathbf{Z}\to \mathbf{Z}$，把它视为
$\mathbf{Z}$-模同态。令 $N = \mathbf{Z}/2$。则诱导映射
$\mathbf{Z} \otimes \mathbf{Z}/2 \to \mathbf{Z} \otimes \mathbf{Z}/2$
不是单射。事实上，对
$x \otimes y\in \mathbf{Z} \otimes \mathbf{Z}/2$，
$$
(2 \otimes 1)(x \otimes y) = 2x \otimes y = x \otimes 2y = x \otimes 0 = 0
$$
因此，诱导映射是零映射，而 $\mathbf{Z} \otimes N\neq 0$。
\end{example}

\begin{remark}
\label{algebra-remark-flat-module}
对 $R$-模 $N$，若函子 $-\otimes_R N$ 正合，即与 $N$ 作张量积保持所有
正合列，则称 $N$ 为
% P01-E0282: source article candidate recorded in ERRATA_CANDIDATES.jsonl.
{\it 平坦} $R$-模。我们将在第 \ref{algebra-section-flat} 节进一步讨论这一概念。
\end{remark}

\begin{lemma}
\label{algebra-lemma-tensor-finiteness}
设 $R$ 为环，$M$、$N$ 为 $R$-模。
\begin{enumerate}
\item 若 $N$ 与 $M$ 都是有限模，则 $M \otimes_R N$ 也是有限模。
\item 若 $N$ 与 $M$ 都是有限表示模，则 $M \otimes_R N$ 也是有限表示模。
\end{enumerate}
\end{lemma}

\begin{proof}
设 $M$ 为有限模。选取表示
$0 \to K \to R^{\oplus n} \to M \to 0$。由引理
\ref{algebra-lemma-tensor-product-exact}，这给出正合列
$K \otimes_R N \to N^{\oplus n} \to M \otimes_R N \to 0$。
所以，若 $N$ 也有限，则 $M \otimes_R N$ 是有限模的商，因而有限；
见引理 \ref{algebra-lemma-extension}。类似地，若 $N$ 与 $M$ 都是有限表示模，
则 $K$ 是有限模，而 $M \otimes_R N$ 是有限表示模
$N^{\oplus n}$ 关于有限模 $K \otimes_R N$ 的商，因而是有限表示模；
见引理 \ref{algebra-lemma-extension}。
\end{proof}

\begin{lemma}
\label{algebra-lemma-tensor-localization}
设 $M$ 为 $R$-模。则 $S^{-1}R$-模 $S^{-1}M$ 与
$S^{-1}R \otimes_R M$ 典范同构，且典范同构
$f : S^{-1}R \otimes_R M \to S^{-1}M$ 由下式给出：
$$
f((a/s) \otimes m) = am/s, \forall a \in R, m \in M, s \in S
$$
\end{lemma}

\begin{proof}
显然，由 $f'(a/s, m) = am/s$ 给出的映射
$f' : S^{-1}R \times M \to S^{-1}M$ 是双线性的；
故由泛性质，该映射诱导引理陈述中的唯一 $S^{-1}R$-模同态
$f : S^{-1}R \otimes_R M \to S^{-1}M$。
事实上，$S^{-1}M$ 的每个元素都具有形式 $m/s$，其中
$m\in M, s\in S$；而 $S^{-1}R \otimes_R M$ 的每个元素都具有形式
$1/s \otimes m$。为看出后一事实，把 $S^{-1}R \otimes_R M$ 中的元素写成
$$
\sum_k \frac{a_k}{s_k} \otimes m_k =
\sum_k \frac{a_k t_k}{s} \otimes m_k =
\frac{1}{s} \otimes \sum_k {a_k t_k}m_k = \frac{1}{s} \otimes m
$$
% P01-E0283: source capitalization candidate recorded in ERRATA_CANDIDATES.jsonl.
其中 $m = \sum_k {a_k t_k}m_k$。于是 $f$ 显然是满射；
若 $f(\frac{1}{s} \otimes m) = m/s = 0$，则存在 $t \in S$，
使得 $M$ 中 $tm = 0$。于是
$$
\frac{1}{s} \otimes m = \frac{1}{st} \otimes tm = \frac{1}{st} \otimes 0 = 0
$$
因此 $f$ 是单射。
\end{proof}

\begin{lemma}
\label{algebra-lemma-tensor-product-localization}
设 $M, N$ 为 $R$-模，则有典范 $S^{-1}R$-模同构
$f : S^{-1}M \otimes_{S^{-1}R}S^{-1}N \to S^{-1}(M \otimes_R N)$，
由下式给出：
$$
f((m/s)\otimes(n/t)) = (m \otimes n)/st
$$
\end{lemma}

\begin{proof}
反复使用引理 \ref{algebra-lemma-tensor-with-bimodule} 与引理
\ref{algebra-lemma-tensor-localization}，并注意 $S^{-1}R$ 是
$(R, S^{-1}R)$-双模，可知这两个 $S^{-1}R$-模同构：
\begin{align*}
S^{-1}(M \otimes_R N) & \cong S^{-1}R \otimes_R (M \otimes_R N)\\
 & \cong S^{-1}M \otimes_R N\\
 & \cong (S^{-1}M \otimes_{S^{-1}R}S^{-1}R)\otimes_R N\\
 & \cong S^{-1}M \otimes_{S^{-1}R}(S^{-1}R \otimes_R N)\\
 & \cong S^{-1}M \otimes_{S^{-1}R}S^{-1}N
\end{align*}
容易看出，这一同构正是引理所述同构。
\end{proof}
\section{张量代数}
\label{algebra-section-tensor-algebra}

\noindent
设 $R$ 为环，$M$ 为 $R$-模。定义{\it $M$ 在 $R$ 上的张量代数}为
非交换 $R$-代数
$$
\text{T}(M) = \text{T}_R(M) =
\bigoplus\nolimits_{n \geq 0} \text{T}^n(M)
$$
其中
$\text{T}^0(M) = R$、
$\text{T}^1(M) = M$、
$\text{T}^2(M) = M \otimes_R M$、
$\text{T}^3(M) = M \otimes_R M \otimes_R M$，依此类推。
乘法由纯张量上的下列规则定义：
$$
(x_1 \otimes x_2 \otimes \ldots \otimes x_n)
\cdot
(y_1 \otimes y_2 \otimes \ldots \otimes y_m)
=
x_1 \otimes x_2 \otimes \ldots \otimes x_n \otimes
y_1 \otimes y_2 \otimes \ldots \otimes y_m
$$
再按线性延拓。

\medskip\noindent
定义{\it $M$ 在 $R$ 上的外代数 $\wedge(M)$}为 $\text{T}(M)$
关于由元素 $x \otimes x \in \text{T}^2(M)$ 生成的双边理想的商。
纯张量 $x_1 \otimes \ldots \otimes x_n$ 在 $\wedge^n(M)$ 中的像记为
$x_1 \wedge \ldots \wedge x_n$。这些元素生成 $\wedge^n(M)$；
它们关于每个 $x_i$ 都是 $R$-线性的，并且当两个 $x_i$ 相等时为零
（即作为 $x_1, x_2, \ldots, x_n$ 的函数是交替的）。
$\wedge(M)$ 上的乘法是分次交换的；即每个 $x \in M$ 与 $y \in M$
都满足 $x \wedge y = - y \wedge x$。

\medskip\noindent
例如，设 $M = Rx_1 \oplus \ldots \oplus Rx_n$ 是有限自由模。
此时 $\wedge(M)$ 是自由 $R$-模，其一组基为元素
$$
x_{i_1} \wedge \ldots \wedge x_{i_r}
$$
其中 $0 \leq r \leq n$ 且 $1 \leq i_1 < i_2 < \ldots < i_r \leq n$。

\medskip\noindent
定义{\it $M$ 在 $R$ 上的对称代数 $\text{Sym}(M)$}为 $\text{T}(M)$
关于由元素 $x \otimes y - y \otimes x \in \text{T}^2(M)$
生成的双边理想的商。纯张量 $x_1 \otimes \ldots \otimes x_n$
在 $\text{Sym}^n(M)$ 中的像简记为 $x_1 \ldots x_n$。
这些元素生成 $\text{Sym}^n(M)$，并关于每个 $x_i$ 都是 $R$-线性的。
若元素列 $x_1, \ldots, x_n$ 是元素列 $x_1', \ldots, x_n'$ 的一个排列，
则 $x_1 \ldots x_n = x_1' \ldots x_n'$。因此，$\text{Sym}(M)$ 是交换的。

\medskip\noindent
仍以有限自由模 $M = Rx_1 \oplus \ldots \oplus Rx_n$ 为例。
此时 $\text{Sym}(M) = R[x_1, \ldots, x_n]$ 是多项式代数。

\begin{lemma}
\label{algebra-lemma-free-tensor-algebra}
设 $R$ 为环，$M$ 为 $R$-模。若 $M$ 是自由 $R$-模，
则其每个对称幂与外幂也是自由模。
\end{lemma}

\begin{proof}
从略；有限自由的情形见上文。
\end{proof}

\begin{lemma}
\label{algebra-lemma-presentation-sym-exterior}
设 $R$ 为环，并设
$M_2 \to M_1 \to M \to 0$ 为 $R$-模的正合列。则有正合列
$$
M_2 \otimes_R \text{Sym}^{n - 1}(M_1)
\to
\text{Sym}^n(M_1)
\to
\text{Sym}^n(M)
\to
0
$$
以及类似的正合列
$$
M_2 \otimes_R \wedge^{n - 1}(M_1)
\to
\wedge^n(M_1)
\to
\wedge^n(M)
\to
0
$$
\end{lemma}

\begin{proof}
从略。
\end{proof}

\begin{lemma}
\label{algebra-lemma-present-sym-wedge}
设 $R$ 为环，$M$ 为 $R$-模。设 $x_i$（$i \in I$）是 $M$
作为 $R$-模的一组给定生成元，并设 $n \geq 2$。存在典范正合列
$$
\bigoplus_{1 \leq j_1 < j_2 \leq n}
\bigoplus_{i_1, i_2 \in I}
\text{T}^{n - 2}(M)
\oplus
\bigoplus_{1 \leq j_1 < j_2 \leq n}
\bigoplus_{i \in I}
\text{T}^{n - 2}(M)
\to
\text{T}^n(M)
\to
\wedge^n(M)
\to
0
$$
其中，第一个直和项中的纯张量
$m_1 \otimes \ldots \otimes m_{n - 2}$ 映为
\begin{align*}
\underbrace{
m_1 \otimes \ldots \otimes x_{i_1} \otimes \ldots
\otimes x_{i_2} \otimes \ldots \otimes m_{n - 2}
}_{\text{其中 } x_{i_1} \text{ 且 } x_{i_2}
\text{ occupying slots } j_1 \text{ 且 } j_2
\text{ in the tensor}} \\
+
\underbrace{
m_1 \otimes \ldots \otimes x_{i_2} \otimes \ldots
\otimes x_{i_1} \otimes \ldots \otimes m_{n - 2}
}_{\text{其中 } x_{i_2} \text{ 且 } x_{i_1}
\text{ occupying slots } j_1 \text{ 且 } j_2
\text{ in the tensor}}
\end{align*}
而第二个直和项中的 $m_1 \otimes \ldots \otimes m_{n - 2}$ 映为
$$
\underbrace{
m_1 \otimes \ldots \otimes x_i \otimes \ldots
\otimes x_i \otimes \ldots \otimes m_{n - 2}
}_{\text{其中 } x_{i} \text{ 且 } x_{i}
\text{ occupying slots } j_1 \text{ 且 } j_2
\text{ in the tensor}}
$$
还有典范正合列
$$
\bigoplus_{1 \leq j_1 < j_2 \leq n}
\bigoplus_{i_1, i_2 \in I}
\text{T}^{n - 2}(M)
\to
\text{T}^n(M)
\to
\text{Sym}^n(M)
\to
0
$$
其中，纯张量 $m_1 \otimes \ldots \otimes m_{n - 2}$ 映为
\begin{align*}
\underbrace{
m_1 \otimes \ldots \otimes x_{i_1} \otimes \ldots
\otimes x_{i_2} \otimes \ldots \otimes m_{n - 2}
}_{\text{其中 } x_{i_1} \text{ 且 } x_{i_2}
\text{ occupying slots } j_1 \text{ 且 } j_2
\text{ in the tensor}} \\
-
\underbrace{
m_1 \otimes \ldots \otimes x_{i_2} \otimes \ldots
\otimes x_{i_1} \otimes \ldots \otimes m_{n - 2}
}_{\text{其中 } x_{i_2} \text{ 且 } x_{i_1}
\text{ occupying slots } j_1 \text{ 且 } j_2
\text{ in the tensor}}
\end{align*}
\end{lemma}

\begin{proof}
从略。
\end{proof}

\begin{lemma}
\label{algebra-lemma-present-wedge}
设 $A \to B$ 为环同态，$M$ 为 $B$-模，并设 $n > 1$。
$A$-线性映射
$M \otimes_A \ldots \otimes_A M \to \wedge^n_B(M)$ 的核，
作为 $A$-模由下列元素生成：满足 $m_i = m_j$ 的
$m_1 \otimes \ldots \otimes m_n$，其中
$i \not = j$、$m_1, \ldots, m_n \in M$；以及元素
$m_1 \otimes \ldots \otimes bm_i \otimes \ldots \otimes m_n -
m_1 \otimes \ldots \otimes bm_j \otimes \ldots \otimes m_n$，
其中 $i \not = j$、$m_1, \ldots, m_n \in M$ 且 $b \in B$。
\end{lemma}

\begin{proof}
从略。
\end{proof}

\begin{lemma}
\label{algebra-lemma-colimit-tensor-algebra}
\begin{slogan}
取张量代数与滤过余极限交换。
\end{slogan}
设 $R$ 为环，$M_i$ 为 $R$-模的有向系统。则
$\colim_i \text{T}(M_i) = \text{T}(\colim_i M_i)$，
对称代数与外代数也有类似结论。
\end{lemma}

\begin{proof}
从略。提示：应用引理 \ref{algebra-lemma-tensor-products-commute-with-limits}。
\end{proof}

\begin{lemma}
\label{algebra-lemma-tensor-algebra-localization}
设 $R$ 为环，$S \subset R$ 为乘法子集。则对任意 $R$-模 $M$，
% P01-E0284: source notation candidate recorded in ERRATA_CANDIDATES.jsonl.
$S^{-1}T_R(M) = T_{S^{-1}R}(S^{-1}M)$。对称代数与外代数也类似。
\end{lemma}

\begin{proof}
从略。提示：应用引理 \ref{algebra-lemma-tensor-product-localization}。
\end{proof}

\section{基变换}
\label{algebra-section-base-change}

\noindent
如下正式引入代数中的基变换。

\begin{definition}
\label{algebra-definition-base-change}
设 $\varphi : R \to S$ 为环同态，$M$ 为 $S$-模，并设
$R \to R'$ 为任意环同态。$\varphi$ 沿 $R \to R'$ 的
{\it 基变换}是环同态 $R' \to S \otimes_R R'$。在这种情形下，
常记 $S' = S \otimes_R R'$。$S$-模 $M$ 的{\it 基变换}是
$S'$-模 $M \otimes_R R'$。
\end{definition}

\noindent
若对某族变量 $x_i$（$i \in I$）及某族多项式
$f_j \in R[x_i]$（$j \in J$），有 $S = R[x_i]/(f_j)$，则
$S \otimes_R R' = R'[x_i]/(f'_j)$，其中 $f'_j \in R'[x_i]$
是 $f_j$ 在由 $R \to R'$ 诱导的映射
$R[x_i] \to R'[x_i]$ 下的像。这个简单观察是理解基变换的关键。

\begin{lemma}
\label{algebra-lemma-base-change-finiteness}
设 $R \to S$ 为环同态，$M$ 为 $S$-模。设 $R \to R'$ 为环同态，
并令 $S' = S \otimes_R R'$ 与 $M' = M \otimes_R R'$ 为相应基变换。
\begin{enumerate}
\item 若 $M$ 是有限 $S$-模，则基变换 $M'$ 是有限 $S'$-模。
\item 若 $M$ 是有限表示 $S$-模，则基变换 $M'$ 是有限表示 $S'$-模。
\item 若 $R \to S$ 为有限型，则基变换 $R' \to S'$ 也为有限型。
\item 若 $R \to S$ 为有限表示，则基变换 $R' \to S'$ 也为有限表示。
\end{enumerate}
\end{lemma}

\begin{proof}
(1) 的证明。取满的 $S$-线性映射
$S^{\oplus n} \to M \to 0$。由引理
\ref{algebra-lemma-flip-tensor-product} 与 \ref{algebra-lemma-tensor-product-exact}，
与 $R^\prime$ 作张量积后得到满射
${S^\prime}^{\oplus n} \to M^\prime \rightarrow 0$，
故 $M^\prime$ 是有限生成 $S^\prime$-模。
(2) 的证明。取表示
$S^{\oplus m} \to S^{\oplus n} \to M \to 0$。
由引理 \ref{algebra-lemma-flip-tensor-product} 与
\ref{algebra-lemma-tensor-product-exact}，与 $R^\prime$ 作张量积后得到
$S^\prime$-模 $M^\prime$ 的有限表示
${S^\prime}^{\oplus m} \to {S^\prime}^{\oplus n} \to M^\prime \to 0$。
(3) 的证明。由引理前的观察即可；按假设可取 $I$ 为有限集。
(4) 的证明。同理由引理前的观察即可；按假设可取 $I$ 与 $J$ 为有限集。
\end{proof}

\noindent
设 $\varphi : R \to S$ 为环同态。给定 $S$-模 $N$，可按规则
$r \cdot n = \varphi(r)n$ 得到 $R$-模 $N_R$。
这有时称为把 $N$ {\it 限制}到 $R$。

\begin{lemma}
\label{algebra-lemma-adjoint-tensor-restrict}
设 $R \to S$ 为环同态。函子
$\text{Mod}_S \to \text{Mod}_R$，$N \mapsto N_R$（限制），
与函子 $\text{Mod}_R \to \text{Mod}_S$，
$M \mapsto M \otimes_R S$（基变换）互为伴随。公式为
$$
\Hom_R(M, N_R) = \Hom_S(M \otimes_R S, N)
$$
\end{lemma}

\begin{proof}
若 $\alpha : M \to N_R$ 是 $R$-模同态，则按规则
$\alpha'(m \otimes s) = s\alpha(m)$ 定义
$\alpha' : M \otimes_R S \to N$。若
$\beta : M \otimes_R S \to N$ 是 $S$-模同态，则按规则
$\beta'(m) = \beta(m \otimes 1)$ 定义 $\beta' : M \to N_R$。
这两个构造互逆的验证从略。
\end{proof}

\noindent
上述引理说明限制函子有左伴随，即基变换。它也有右伴随。

\begin{lemma}
\label{algebra-lemma-adjoint-hom-restrict}
设 $R \to S$ 为环同态。函子
$\text{Mod}_S \to \text{Mod}_R$，$N \mapsto N_R$（限制），
与函子 $\text{Mod}_R \to \text{Mod}_S$，
$M \mapsto \Hom_R(S, M)$ 互为伴随。公式为
$$
\Hom_R(N_R, M) = \Hom_S(N, \Hom_R(S, M))
$$
\end{lemma}

\begin{proof}
若 $\alpha : N_R \to M$ 是 $R$-模同态，则按规则
$\alpha'(n) = (s \mapsto \alpha(sn))$ 定义
$\alpha' : N \to \Hom_R(S, M)$。若
$\beta : N \to \Hom_R(S, M)$ 是 $S$-模同态，则按规则
$\beta'(n) = \beta(n)(1)$ 定义 $\beta' : N_R \to M$。
这两个构造互逆的验证从略。
\end{proof}

\begin{lemma}
\label{algebra-lemma-hom-from-tensor-product-variant}
设 $R \to S$ 为环同态。给定 $S$-模 $M, N$ 及 $R$-模 $P$，有
$$
\Hom_R(M \otimes_S N, P) = \Hom_S(M, \Hom_R(N, P))
$$
\end{lemma}

\begin{proof}
这可以直接证明，也可由引理 \ref{algebra-lemma-adjoint-hom-restrict}
与 \ref{algebra-lemma-hom-from-tensor-product} 推出。事实上，
\begin{align*}
\Hom_R(M \otimes_S N, P)
& =
\Hom_S(M \otimes_S N, \Hom_R(S, P)) \\
& =
\Hom_S(M, \Hom_S(N, \Hom_R(S, P))) \\
& =
\Hom_S(M, \Hom_R(N, P))
\end{align*}
正合所求。
\end{proof}

\section{杂项}
\label{algebra-section-miscellany}

\noindent
本节的证明不应引用基本概念一节，即第 \ref{algebra-section-rings-basic} 节之外
的任何结果。

\begin{lemma}
\label{algebra-lemma-product-ideals-in-prime}
设 $R$ 为环，$I$、$J$ 为两个理想，$\mathfrak p$ 是包含乘积 $IJ$
的素理想。则 $\mathfrak{p}$ 包含 $I$ 或 $J$。
\end{lemma}

\begin{proof}
反设不然，取 $x \in I \setminus \mathfrak p$ 与
$y \in J \setminus \mathfrak p$。其乘积属于
$IJ \subset \mathfrak p$，这与 $\mathfrak p$ 为素理想矛盾。
\end{proof}

\begin{lemma}[素理想回避]
\label{algebra-lemma-silly}
\begin{slogan}
1. 在仿射概形中，若有限多个点包含于一个开子集中，则它们包含于一个
更小的主开子集中。
2. 仿射开集在仿射概形的给定有限集的邻域中共尾。
\end{slogan}
设 $R$ 为环。设 $I_i \subset R$（$i = 1, \ldots, r$）与
$J \subset R$ 为理想。假设
\begin{enumerate}
\item 对 $i = 1, \ldots, r$，有 $J \not\subset I_i$；且
\item 除至多两个外，所有 $I_i$ 都是素理想。
\end{enumerate}
则存在 $x \in J$，使得对所有 $i$ 都有 $x\not\in I_i$。
\end{lemma}

\begin{proof}
当 $r = 1$ 时结论成立。若 $r = 2$，取 $x, y \in J$，
使 $x \not \in I_1$ 且 $y \not \in I_2$。除非
$x \in I_2$ 且 $y \in I_1$，否则结论已得。在后一情形，
$x + y$ 不可能属于 $I_1$（否则 $x + y - y \in I_1$），
也不可能属于 $I_2$。

\medskip\noindent
设 $r \geq 3$，并假定结论对 $r - 1$ 成立。重新编号后，
可设 $I_r$ 为素理想。还可设各 $I_i$ 之间不存在包含关系。
取 $x \in J$，使对所有 $i = 1, \ldots, r - 1$ 都有
$x \not \in I_i$。若 $x \not\in I_r$，结论已得；故设
$x \in I_r$。若 $J I_1 \ldots I_{r - 1} \subset I_r$，
则由引理 \ref{algebra-lemma-product-ideals-in-prime} 有
$J \subset I_r$，矛盾。取
$y \in J I_1 \ldots I_{r - 1}$ 且 $y \not \in I_r$。
则 $x + y$ 满足要求。
\end{proof}

\begin{lemma}
\label{algebra-lemma-silly-silly}
设 $R$ 为环，$x \in R$，$I \subset R$ 为理想，并设
$\mathfrak p_i$（$i = 1, \ldots, r$）为素理想。
假设对 $i = 1, \ldots, r$，有
$x + I \not \subset \mathfrak p_i$。则存在 $y \in I$，
使得对所有 $i$ 都有 $x + y \not \in \mathfrak p_i$。
\end{lemma}

\begin{proof}
可设各 $\mathfrak p_i$ 之间不存在包含关系。重新排序后，可设
当 $i < s$ 时 $x \not \in \mathfrak p_i$，而当 $i \geq s$ 时
$x \in \mathfrak p_i$。若 $s = r + 1$，结论已得。否则，可找到
$y \in I$，使 $y \not \in \mathfrak p_s$。选取
$f \in \bigcap_{i < s} \mathfrak p_i$，使
$f \not \in \mathfrak p_s$。
% P01-E0285: source element-membership candidate recorded in ERRATA_CANDIDATES.jsonl.
于是 $x + fy$ 不属于 $\mathfrak p_1, \ldots, \mathfrak p_s$。
对 $s$ 归纳即得结论。
\end{proof}

\begin{lemma}[中国剩余定理]
\label{algebra-lemma-chinese-remainder}
设 $R$ 为环。
\begin{enumerate}
\item 若理想 $I_1, \ldots, I_r$ 满足：当 $a \not = b$ 时
$I_a + I_b = R$，则
$I_1 \cap \ldots \cap I_r = I_1I_2\ldots I_r$，且
$R/(I_1I_2\ldots I_r) \cong R/I_1 \times \ldots \times R/I_r$。
\item 若 $\mathfrak m_1, \ldots, \mathfrak m_r$ 是两两不同的极大理想，
则当 $a \not = b$ 时 $\mathfrak m_a + \mathfrak m_b = R$，
故可应用上述结论。
\end{enumerate}
\end{lemma}

\begin{proof}
先证明 $I_1 \cap \ldots \cap I_r = I_1 \ldots I_r$；这也会推出诱导环同态
$R/(I_1 \ldots I_r) \rightarrow R/I_1 \times \ldots \times R/I_r$
的单射性。包含关系
$I_1 \cap \ldots \cap I_r \supset I_1 \ldots I_r$ 总成立，
因为理想对乘以任意环元素封闭。为证明反向包含，断言理想
$$
I_1 \ldots \hat I_i \ldots I_r,\quad i = 1, \ldots, r
$$
生成环 $R$。对 $r$ 归纳。当 $r = 2$ 时成立。若 $r > 2$，
则 $R$ 是理想
$I_1 \ldots \hat I_i \ldots I_{r - 1}$（$i = 1, \ldots, r - 1$）
之和。因此，$I_r$ 是理想
$I_1 \ldots \hat I_i \ldots I_r$（$i = 1, \ldots, r - 1$）
之和。对这些理想按相反次序应用同样论证，可知 $I_1$ 是理想
$I_1 \ldots \hat I_i \ldots I_r$（$i = 2, \ldots, r$）之和。
由假设 $R = I_1 + I_r$，可知 $R$ 是上述所列理想之和。
因而可找到元素
$a_i \in I_1 \ldots \hat I_i \ldots I_r$，使其总和为一。
把此等式乘以 $I_1 \cap \ldots \cap I_r$ 的任一元素，即得反向包含。

还须证明典范映射
$R/(I_1 \ldots I_r) \rightarrow R/I_1 \times \ldots \times R/I_r$
为满射。为此，考察它在上面所得等式
$1 = \sum_{i=1}^r a_i$ 上的作用。一方面，环态射把 1 映到 1；
另一方面，当 $j \neq i$ 时，任意 $a_i$ 在 $R/I_j$ 中的像为零。
因此，$a_i$ 在 $R/I_i$ 中的像是单位元。故给定任意元素
$(\bar{b_1}, \ldots, \bar{b_r}) \in R/I_1 \times \ldots \times R/I_r$，
元素 $\sum_{i=1}^r a_i \cdot b_i$ 是它在 $R$ 中的一个原像。

\medskip\noindent
为证明 (2)，由极大理想两两不同的定义，对 $a \neq b$ 有
$\mathfrak{m}_a + \mathfrak{m}_b = R$，故上述结论适用。
\end{proof}

\begin{lemma}
\label{algebra-lemma-matrix-left-inverse}
设 $R$ 为环，$n \geq m$。设 $A$ 是系数在 $R$ 中的
$n \times m$ 矩阵，并设 $J \subset R$ 是由 $A$ 的
$m \times m$ 子式生成的理想。
\begin{enumerate}
\item 对任意 $f \in J$，存在 $m \times n$ 矩阵 $B$，使
$BA = f 1_{m \times m}$。
\item 若 $f \in R$，且对某个 $m \times n$ 矩阵 $B$ 有
$BA = f 1_{m \times m}$，则 $f^m \in J$。
\end{enumerate}
\end{lemma}

\begin{proof}
对满足 $I \subset \{1, \ldots, n\}$ 与 $|I| = m$ 的这类子集，
用 $E_I$ 表示投影
$$
R^{\oplus n} = \bigoplus\nolimits_{i \in \{1, \ldots, n\}} R
\longrightarrow \bigoplus\nolimits_{i \in I} R
$$
的 $m \times n$ 矩阵，并令 $A_I = E_I A$；即 $A_I$ 是由 $A$
中下标属于 $I$ 的行构成的 $m \times m$ 矩阵。
设 $B_I$ 为 $A_I$ 的伴随矩阵（代数余子式矩阵的转置），即满足
$A_I B_I = B_I A_I = \det(A_I) 1_{m \times m}$。
$A$ 的 $m \times m$ 子式正是所有满足
$I \subset \{1, \ldots, n\}$、$|I| = m$ 的这类子集所对应的行列式
$\det A_I$。若 $f \in J$，则对某些 $c_I \in R$，可写成
$f = \sum c_I \det(A_I)$。令 $B = \sum c_I B_I E_I$，即得 (1)。

\medskip\noindent
若 $f 1_{m \times m} = BA$，则由柯西--比内公式
(\ref{algebra-item-cauchy-binet})，
$f^m = \sum b_I \det(A_I)$，其中 $b_I$ 是由 $B$ 中下标属于 $I$
的列构成的 $m \times m$ 矩阵的行列式。
\end{proof}

\begin{lemma}
\label{algebra-lemma-matrix-right-inverse}
设 $R$ 为环，$n \geq m$。设 $A = (a_{ij})$ 是系数在 $R$ 中的
$n \times m$ 矩阵，按分块形式写为
$$
A =
\left(
\begin{matrix}
A_1 \\
A_2
\end{matrix}
\right)
$$
其中 $A_1$ 的大小为 $m \times m$。设 $B$ 为 $A_1$ 的伴随矩阵
（代数余子式矩阵的转置）。则
$$
AB = 
\left(
\begin{matrix}
f 1_{m \times m} \\
C
\end{matrix}
\right)
$$
其中 $f = \det(A_1)$，而 $c_{ij}$ 在相差一个符号的意义下，
是 $A$ 中由行 $1, \ldots, \hat j, \ldots, m, i$ 确定的
$m \times m$ 子式。
\end{lemma}

\begin{proof}
% P01-E0286: source missing-identity candidate recorded in ERRATA_CANDIDATES.jsonl.
由于伴随矩阵满足 $A_1B = B A_1 = f$，故 $AB$ 表达式的第一个分块正确。
注意
$$
c_{ij} = \sum\nolimits_k a_{ik}b_{kj} = \sum (-1)^{j + k}a_{ik} \det(A_1^{jk})
$$
% P01-E0287: source index-name candidate recorded in ERRATA_CANDIDATES.jsonl.
其中 $A_1^{ij}$ 表示从 $A_1$ 中删去第 $j$ 行与第 $k$ 列所得矩阵。
最后这个表达式正是引理陈述中矩阵行列式的按行展开。
\end{proof}

\begin{lemma}
\label{algebra-lemma-map-cannot-be-injective}
\begin{slogan}
若有限自由模的源的秩大于目标的秩，则该映射不可能是单射。
\end{slogan}
设 $R$ 为非零环，$n \geq 1$，$M$ 为由 $< n$ 个元素生成的 $R$-模。
则任意 $R$-模同态 $f : R^{\oplus n} \to M$ 的核都非零。
\end{lemma}

\begin{proof}
选取满射 $R^{\oplus n - 1} \to M$。可把映射 $f$ 提升为映射
$f' : R^{\oplus n} \to R^{\oplus n - 1}$（引理
\ref{algebra-lemma-lift-map}）。只须证明 $f'$ 的核非零。
映射 $f' : R^{\oplus n} \to R^{\oplus n - 1}$ 由矩阵
$A = (a_{ij})$ 给出。若某个 $a_{ij}$ 不是幂零元，比如
$a = a_{ij}$ 不是，则可把 $R$ 替换为局部化 $R_a$，从而可设
$a_{ij}$ 是单位。事实上，若局部化后找到非零核，则原来就有非零核，
因为局部化是正合的；见命题 \ref{algebra-proposition-localization-exact}。
在这种情形下，可以在 $R^{\oplus n}$ 与 $R^{\oplus n - 1}$
上都作
% P01-E0288: source terminology candidate recorded in ERRATA_CANDIDATES.jsonl.
基的变换，归约到
$$
A =
\left(
\begin{matrix}
1 & 0 & 0 & \ldots \\
0 & a_{22} & a_{23} & \ldots \\
0 & a_{32} & \ldots \\
\ldots & \ldots
\end{matrix}
\right)
$$
的情形。于是，此时对 $n$ 归纳即可。否则，每个 $a_{ij}$ 都是幂零元。
令 $I = (a_{ij}) \subset R$。注意，对某个 $m \geq 0$，
有 $I^{m + 1} = 0$。令 $m$ 为使 $I^m \not = 0$ 的最大整数。
于是 $(I^m)^{\oplus n}$ 包含于该映射的核中，结论得证。
\end{proof}

\begin{lemma}
\label{algebra-lemma-rank}
\begin{slogan}
有限自由模的秩是良定义的。
\end{slogan}
设 $R$ 为非零环，$n, m \geq 0$ 为整数。若作为 $R$-模，
$R^{\oplus n}$ 同构于 $R^{\oplus m}$，则 $n = m$。
\end{lemma}

\begin{proof}
立即由引理 \ref{algebra-lemma-map-cannot-be-injective} 得到。
\end{proof}

\section{凯莱--哈密顿}
\label{algebra-section-cayley-hamilton}

\begin{lemma}
\label{algebra-lemma-charpoly}
设 $R$ 为环，$A = (a_{ij})$ 是系数在 $R$ 中的 $n \times n$ 矩阵。
设 $P(x) \in R[x]$ 为 $A$ 的特征多项式（定义为
$\det(x\text{id}_{n \times n} - A)$）。则在
$\text{Mat}(n \times n, R)$ 中有 $P(A) = 0$。
\end{lemma}

\begin{proof}
分几步把问题归约到线性代数中熟知的凯莱--哈密顿定理：
\begin{enumerate}
\item 若 $\phi :S \rightarrow R$ 是环态射，而 $b_{ij}$ 是
$a_{ij}$ 在此映射下的原像，则只须对 $S$ 与 $(b_{ij})$ 证明该陈述，
因为 $\phi$ 是环态射。
\item 若 $\psi :R \hookrightarrow S$ 是单射环态射，则显然只须对
$S$ 以及视为 $S$ 中元素的 $a_{ij}$ 证明结论。
\item 因此，可先归约到多项式环情形
$R = \mathbf{Z}[X_{ij}]$、$a_{ij} = X_{ij}$，再进一步归约到
$R = \mathbf{Q}(X_{ij})$；在后一情形可直接应用凯莱--哈密顿定理。
\end{enumerate}
\end{proof}

\begin{lemma}
\label{algebra-lemma-charpoly-module}
设 $R$ 为环，$M$ 为有限 $R$-模，$\varphi : M \to M$ 为自同态。
则存在首一多项式 $P \in R[T]$，使得作为 $M$ 的自同态有
$P(\varphi) = 0$。
\end{lemma}

\begin{proof}
选取满的 $R$-模同态 $R^{\oplus n} \to M$，它由
$(a_1, \ldots, a_n) \mapsto \sum a_ix_i$ 给出，其中
$x_i \in M$ 是某些生成元。选取
$(a_{i1}, \ldots, a_{in}) \in R^{\oplus n}$，使
$\varphi(x_i) = \sum a_{ij} x_j$。换言之，下图交换：
$$
\xymatrix{
R^{\oplus n} \ar[d]_A \ar[r] & M \ar[d]^\varphi \\
R^{\oplus n} \ar[r] & M
}
$$
其中 $A = (a_{ij})$。由引理 \ref{algebra-lemma-charpoly}，
存在首一多项式 $P$，使 $P(A) = 0$。于是 $P(\varphi) = 0$。
\end{proof}

\begin{lemma}
\label{algebra-lemma-charpoly-module-ideal}
设 $R$ 为环，$I \subset R$ 为理想，$M$ 为有限 $R$-模。
设 $\varphi : M \to M$ 为满足 $\varphi(M) \subset IM$ 的自同态。
则存在首一多项式
% P01-E0289: source polynomial-variable candidate recorded in ERRATA_CANDIDATES.jsonl.
$P = t^n + a_1 t^{n - 1} + \ldots + a_n \in R[T]$，
使得 $a_j \in I^j$，并且作为 $M$ 的自同态有 $P(\varphi) = 0$。
\end{lemma}

\begin{proof}
选取满的 $R$-模同态 $R^{\oplus n} \to M$，它由
$(a_1, \ldots, a_n) \mapsto \sum a_ix_i$ 给出，其中
$x_i \in M$ 是某些生成元。选取
$(a_{i1}, \ldots, a_{in}) \in I^{\oplus n}$，使
$\varphi(x_i) = \sum a_{ij} x_j$。换言之，下图交换：
$$
\xymatrix{
R^{\oplus n} \ar[d]_A \ar[r] & M \ar[d]^\varphi \\
I^{\oplus n} \ar[r] & M
}
$$
其中 $A = (a_{ij})$。由引理 \ref{algebra-lemma-charpoly}，多项式
$P(t) = \det(t\text{id}_{n \times n} - A)$ 具有所有所需性质。
\end{proof}

\noindent
作为一个有趣的应用，证明下列令人惊讶的引理。

\begin{lemma}
\label{algebra-lemma-fun}
设 $R$ 为环，$M$ 为有限 $R$-模，并设
$\varphi : M \to M$ 为满的 $R$-模同态。则 $\varphi$ 是同构。
\end{lemma}

\begin{proof}[第一种证明]
记 $R' = R[x]$，并把 $M$ 视为有限 $R'$-模，其中 $x$ 通过
$\varphi$ 作用。令 $I = (x) \subset R'$。由 $\varphi$ 为满射的假设，
有 $IM = M$。因此，可对作为 $R'$-模的 $M$、理想 $I$ 以及自同态
$\text{id}_M$ 应用引理 \ref{algebra-lemma-charpoly-module-ideal}。由此得到
$(1 + a_1 + \ldots + a_n)\text{id}_M = 0$，其中 $a_j \in I$。
对某个 $b_j(x) \in R[x]$，写 $a_j = b_j(x)x$。换回 $\varphi$，
得到
$\text{id}_M = -(\sum_{j = 1, \ldots, n} b_j(\varphi)) \varphi$，
故 $\varphi$ 可逆。
\end{proof}

\begin{proof}[第二种证明]
对 $M$ 在 $R$ 上的生成元个数作归纳。若 $M$ 由一个元素生成，
则对某个理想 $I \subset R$，有 $M \cong R/I$。
此时可把 $R$ 替换为 $R/I$，从而 $M = R$。
此时 $\varphi : R \to R$ 由 $M$ 上乘以某个元素 $r \in R$ 给出。
$\varphi$ 的满射性迫使 $r$ 可逆，因为 $\varphi$ 必须命中 $1$；
所以 $\varphi$ 可逆。

\medskip\noindent
现假设引理对由 $n - 1$ 个元素生成的模已得证明，并考察由 $n$
个元素生成的模 $M$。令 $A$ 表示环 $R[t]$，并令 $t$ 通过
$\varphi$ 作用，从而把模 $M$ 视为 $A$-模。由于 $M$ 在 $R$ 上有限，
它在 $R[t]$ 上也有限；而我们要证明的 $\varphi$ 的单射性是集合论性质，
故不妨证明 $A$ 上的自同态 $t : M \to M$ 是单射。
问题已归约到自同态由底环中元素乘法给出的情形。
令 $M' \subset M$ 表示由 $M$ 的前 $n - 1$ 个生成元生成的
子-$A$-模，并考虑交换图
$$
\xymatrix{
0 \ar[r] & M' \ar[r]\ar[d]^{\varphi\mid_{M'}} & M\ar[d]^\varphi \ar[r] &
M/M' \ar[d]^{\varphi \bmod M'} \ar[r] & 0 \\
0 \ar[r] & M' \ar[r]                                  & M \ar[r]
           & M/M' \ar[r]                                  & 0,
}
$$
$\varphi$ 在 $M'$ 上的限制以及 $\varphi$ 在商 $M/M'$ 上诱导的映射
均为良定义，因为 $\varphi$ 是乘以底环中的一个元素，而 $M'$ 与
$M/M'$ 各自都是 $A$-模。由 $n = 1$ 的情形，映射
$M/M' \to M/M'$ 是同构。追图表明 $\varphi|_{M'}$ 为满射，
故由归纳假设 $\varphi|_{M'}$ 是同构。再由蛇形引理，
中间的竖直映射必为同构。
\end{proof}
\section{环的谱}
\label{algebra-section-spectrum-ring}

\noindent
我们约定把作为拓扑空间的环谱纳入代数章，而把仿射概形纳入概形章；
这一划分只是任意的编排选择。

\begin{definition}
\label{algebra-definition-spectrum-ring}
设 $R$ 为环。
\begin{enumerate}
\item $R$ 的{\it 谱}是 $R$ 的全体素理想组成的集合，通常记作
$\Spec(R)$。
\item 给定子集 $T \subset R$，令 $V(T) \subset \Spec(R)$ 为包含
$T$ 的素理想之集，即 $V(T) = \{ \mathfrak p \in
\Spec(R) \mid \forall f\in T, f\in \mathfrak p\}$。
\item 给定元素 $f \in R$，令 $D(f) \subset \Spec(R)$ 为不包含
$f$ 的素理想之集。
\end{enumerate}
\end{definition}

\begin{lemma}
\label{algebra-lemma-Zariski-topology}
设 $R$ 为环。
\begin{enumerate}
\item 环 $R$ 的谱为空，当且仅当 $R$ 为零环。
\item 每个非零环都有极大理想。
\item 每个非零环都有极小素理想。
\item 给定理想 $I \subset R$ 以及素理想
$I \subset \mathfrak p$，存在素理想
$I \subset \mathfrak q \subset \mathfrak p$，使得
$\mathfrak q$ 在 $I$ 上极小。
\item 若 $T \subset R$，且 $(T)$ 是 $T$ 在 $R$ 中生成的理想，
则 $V((T)) = V(T)$。
\item 若 $I$ 为理想而 $\sqrt{I}$ 为其根理想，参见基本概念
(\ref{algebra-item-radical-ideal})，则 $V(I) = V(\sqrt{I})$。
\item 给定 $R$ 的理想 $I$，有 $\sqrt{I} =
\bigcap_{I \subset \mathfrak p} \mathfrak p$。
\item 若 $I$ 为理想，则 $V(I) = \emptyset$ 当且仅当 $I$ 为单位理想。
\item 若 $I$, $J$ 是 $R$ 的理想，则 $V(I) \cup V(J) =
V(I \cap J)$。
\item 若 $(I_a)_{a\in A}$ 是 $R$ 的一族理想，则
$\bigcap_{a\in A} V(I_a) = V(\bigcup_{a\in A} I_a)$。
\item 若 $f \in R$，则 $D(f) \amalg V(f) = \Spec(R)$。
\item 若 $f \in R$，则 $D(f) = \emptyset$ 当且仅当 $f$ 幂零。
\item 若对某个单位 $u \in R$ 有 $f = u f'$，则 $D(f) = D(f')$。
\item 若 $I \subset R$ 为理想，而 $\mathfrak p$ 是 $R$ 的素理想且
$\mathfrak p \not\in V(I)$，则存在 $f \in R$，使得
$\mathfrak p \in D(f)$ 且 $D(f) \cap V(I) = \emptyset$。
\item 若 $f, g \in R$，则 $D(fg) = D(f) \cap D(g)$。
\item 若对 $i \in I$ 有 $f_i \in R$，则
$\bigcup_{i\in I} D(f_i)$ 是 $V(\{f_i \}_{i\in I})$ 在
$\Spec(R)$ 中的补集。
\item 若 $f \in R$ 且 $D(f) = \Spec(R)$，则 $f$ 为单位。
\end{enumerate}
\end{lemma}

\begin{proof}
依次证明以上各项。
\begin{enumerate}
\item 这是 (2) 或 (3) 的直接推论。
\item 令 $\mathfrak{A}$ 为 $R$ 的全体真理想组成的集合。它按包含关系
排序，而且非空，因为 $(0) \in \mathfrak{A}$ 是真理想。设 $A$ 是
$\mathfrak A$ 的全序子集。于是 $\bigcup_{I \in A} I$ 确为理想。
由于对所有 $I \in A$ 都有 1 $\notin I$，这个并不包含 1，因而仍为
真理想。因此 $\bigcup_{I \in A} I$ 属于 $\mathfrak{A}$，并且是
$A$ 的一个上界。由佐恩引理，$\mathfrak{A}$ 有极大元，而它正是所求的
极大理想。
\item 因为 $R$ 非零，它含有一个极大理想，而极大理想是素理想。因此
$R$ 的全体素理想组成的集合 $\mathfrak{A}$ 非空。按反向包含关系排列
$\mathfrak{A}$。设 $A$ 是 $\mathfrak{A}$ 的全序子集。显然
$J = \bigcap_{I \in A} I$ 是理想；但它为素理想并不那么显然。设
$xy \in J$，则对所有 $I \in A$ 都有 $xy \in I$。现令
$B = \{I \in A | y \in I\}$，并令 $K = \bigcap_{I \in B} I$。
由于 $A$ 全序，或者 $K = J$（此时 $y \in J$，结论成立），或者
$K \supset J$，并且对每个真包含于 $K$ 的 $I \in A$，都有
$y \notin I$。但这些 $I$ 都是素理想，所以对所有这样的
$I, x \in I$。因此 $x \in J$。两种情形下 $J$ 都是素理想。
故由佐恩引理得到一个极大元；在这里，它就是极小素理想。
\item 论证与 (3) 完全相同，只需仅考虑包含于 $\mathfrak{p}$ 且包含
$I$ 的素理想。
\item $(T)$ 是包含 $T$ 的最小理想。因此若 $T \subset I$，其中右端
为某个理想，则也有 $(T) \subset I$。所以若 $I \in V(T)$，
则也有 $I \in V((T))$。反向包含显然。
\item 由 $I \subset \sqrt{I}, V(\sqrt{I}) \subset V(I)$。现设
$\mathfrak{p} \in V(I)$，并取 $x \in \sqrt{I}$。则对某个 $n$ 有
$x^n \in I$，从而 $x^n \in \mathfrak{p}$。由于 $\mathfrak{p}$
为素理想，一个直接的归纳论证给出 $x \in \mathfrak{p}$。故
$\sqrt{I} \subset \mathfrak{p}$，于是
$\mathfrak{p} \in V(\sqrt{I})$。
\item 设 $f \in R \setminus \sqrt{I}$。则对所有 $n$ 都有
$f^n \notin I$。因此 $S = \{1, f, f^2, \ldots\}$ 是不含 $0$ 的
乘法子集。取包含 $S^{-1}I$ 的素理想
$\bar{\mathfrak{p}} \subset S^{-1}R$。令 $\mathfrak{p}$ 为
$\bar{\mathfrak{p}}$ 在 $R$ 中的原像；它是包含 $I$ 且不与 $S$
相交的素理想。这说明
$\bigcap_{I \subset \mathfrak p} \mathfrak p \subset \sqrt{I}$。
反之，若 $a \in \sqrt{I}$，则对某个 $n$ 有 $a^n \in I$。所以当
$I \subset \mathfrak{p}$ 时，$a^n \in \mathfrak{p}$。由于
$\mathfrak{p}$ 为素理想，故 $a \in \mathfrak{p}$，从而得证等式。
\item $I$ 不是单位理想，当且仅当 $I$ 包含于某个极大理想（将 (2)
用于环 $R/I$ 即可看出），而该极大理想因而是素理想。
% P01-E0290: malformed ideal-containment disjunction recorded in ERRATA_CANDIDATES.jsonl.
\item 若 $\mathfrak{p} \in V(I) \cup V(J)$，则 $I \subset \mathfrak{p}$
或 $J \subset \mathfrak{p}$，从而 $I \cap J \subset \mathfrak{p}$。
反之，若 $I \cap J \subset \mathfrak{p}$，则
$IJ \subset \mathfrak{p}$；由于 $\mathfrak{p}$ 为素理想，故
$I$ 或 $J$ 中有一个包含于 $\mathfrak{p}$。
\item $\mathfrak{p} \in \bigcap_{a \in A} V(I_a) \Leftrightarrow
I_a \subset \mathfrak{p}, \forall a \in A \Leftrightarrow
\mathfrak{p} \in V(\bigcup_{a\in A} I_a)$。
\item 若 $\mathfrak{p}$ 为素理想且 $f \in R$，则严格地说，或者
$f \in \mathfrak{p}$，或者 $f \notin \mathfrak{p}$；这正是该不交并
所表达的内容。
\item 若 $a \in R$ 幂零，则对某个 $n$ 有 $a^n = 0$。于是对每个
素理想都有 $a^n \in \mathfrak{p}$。归纳可得
$a \in \mathfrak{p}$，故 $D(a) = \emptyset$。反之，由 (7)，若
$a \in R$ 不幂零，则存在不包含它的素理想。
\item $f \in \mathfrak{p} \Leftrightarrow uf \in \mathfrak{p}$，
因为 $u$ 可逆。
\item 若 $\mathfrak{p} \notin V(I)$，则
$\exists f \in I \setminus \mathfrak{p}$。于是
$f \notin \mathfrak{p}$，所以 $\mathfrak{p} \in D(f)$。此外，若
$\mathfrak{q} \in D(f)$，则 $f \notin \mathfrak{q}$，从而 $I$
不包含于 $\mathfrak{q}$。故 $D(f) \cap V(I) = \emptyset$。
\item 若 $fg \in \mathfrak{p}$，则 $f \in \mathfrak{p}$ 或
$g \in \mathfrak{p}$。因此，若 $f \notin \mathfrak{p}$ 且
$g \notin \mathfrak{p}$，则 $fg \notin \mathfrak{p}$。又因为
$\mathfrak{p}$ 是理想，若 $fg \notin \mathfrak{p}$，则
$f \notin \mathfrak{p}$ 且 $g \notin \mathfrak{p}$。
\item $\mathfrak{p} \in \bigcup_{i \in I} D(f_i) \Leftrightarrow \exists i \in
I, f_i \notin \mathfrak{p} \Leftrightarrow \mathfrak{p} \in \Spec(R)
\setminus V(\{f_i\}_{i \in I})$。
\item 若 $D(f) = \Spec(R)$，则 $V(f) = \emptyset$，因而
$fR = R$，所以 $f$ 是单位。
\end{enumerate}
\end{proof}

\noindent
该引理说明，定义 \ref{algebra-definition-spectrum-ring} 中的子集 $V(T)$
构成 $\Spec(R)$ 上某个拓扑的闭集。它还说明集合 $D(f)$ 是开集，
并构成该拓扑的一组基。

\begin{definition}
\label{algebra-definition-Zariski-topology}
设 $R$ 为环。以各集合 $V(T)$ 为闭集的 $\Spec(R)$ 上的拓扑称为
{\it 扎里斯基}拓扑。开子集 $D(f)$ 称为 $\Spec(R)$ 的
{\it 标准开集}。
\end{definition}

\noindent
由上下文应当能够看清我们是仅把 $\Spec(R)$ 视为集合，还是把它视为
拓扑空间。

\begin{lemma}
\label{algebra-lemma-spec-functorial}
\begin{slogan}
谱的函子性
\end{slogan}
设 $\varphi : R \to R'$ 为环同态。诱导映射
$$
\Spec(\varphi) : \Spec(R') \longrightarrow \Spec(R),
\quad
\mathfrak p' \longmapsto \varphi^{-1}(\mathfrak p')
$$
关于扎里斯基拓扑连续。事实上，对任意元素 $f \in R$，有
$\Spec(\varphi)^{-1}(D(f)) = D(\varphi(f))$。
\end{lemma}

\begin{proof}
由基本概念 (\ref{algebra-item-inverse-image-prime})，
$\mathfrak p := \varphi^{-1}(\mathfrak p')$ 确为 $R$ 的素理想。
引理的最后一个断言直接来自定义，并蕴含第一个断言。
\end{proof}

\noindent
若 $\varphi' : R' \to R''$ 是第二个环同态，则复合
$$
\Spec(R'')
\longrightarrow
\Spec(R')
\longrightarrow
\Spec(R)
$$
等于 $\Spec(\varphi' \circ \varphi)$。换言之，$\Spec$ 是从环范畴
到拓扑空间范畴的反变函子。

\begin{lemma}
\label{algebra-lemma-spec-localization}
设 $R$ 为环，$S \subset R$ 为乘法子集。由 $\Spec$ 的函子性，
映射 $R \to S^{-1}R$ 诱导同胚
$$
\Spec(S^{-1}R)
\longrightarrow
\{\mathfrak p \in \Spec(R) \mid S \cap \mathfrak p = \emptyset \}
$$
其中右端赋予从 $\Spec(R)$ 的扎里斯基拓扑诱导的拓扑。逆映射由
$\mathfrak p \mapsto S^{-1}\mathfrak p = \mathfrak p(S^{-1}R)$ 给出。
\end{lemma}

\begin{proof}
记引理中箭头的右端为 $D$。取素理想
% P01-E0291: missing copula in source construction recorded in ERRATA_CANDIDATES.jsonl.
$\mathfrak p' \subset S^{-1}R$，并令 $\mathfrak p$ 为
$\mathfrak p'$ 在 $R$ 中的原像。由于 $\mathfrak p'$ 不包含 $1$，
可知 $\mathfrak p$ 不包含 $S$ 的任何元素。因此 $\mathfrak p \in D$，
从而像包含于 $D$。现设 $\mathfrak p \in D$。依假设，像
$\overline{S}$ 不包含 $0$。由基本概念
(\ref{algebra-item-localization-zero})，$\overline{S}^{-1}(R/\mathfrak p)$
不是零环。由基本概念 (\ref{algebra-item-localize-ideal})，可见
$S^{-1}R / S^{-1}\mathfrak p = \overline{S}^{-1}(R/\mathfrak p)$
为整环，故 $S^{-1}\mathfrak p$ 为素理想。这个环等式还表明
$S^{-1}\mathfrak p$ 在 $R$ 中的原像等于 $\mathfrak p$，因为由基本概念
(\ref{algebra-item-localize-nonzerodivisors})，映射
$R/\mathfrak p \to \overline{S}^{-1}(R/\mathfrak p)$ 为单射。
这就证明了映射 $\Spec(S^{-1}R) \to \Spec(R)$ 双射到 $D$，并且其
逆映射如上所述。由引理 \ref{algebra-lemma-spec-functorial}，该映射连续。
最后，设 $D(g) \subset \Spec(S^{-1}R)$ 为标准开集。将其写成
$g = h/s$，其中 $h\in R$ 且 $s\in S$。由于 $g$ 与
$h/1$ 相差一个单位，在 $\Spec(S^{-1}R)$ 中有
$D(g) = D(h/1)$。于是由引理 \ref{algebra-lemma-spec-functorial} 及上述
双射性，$D(g) = D(h/1)$ 的像为 $D \cap D(h)$。这也证明了该映射
为开映射。
\end{proof}

\begin{lemma}
\label{algebra-lemma-standard-open}
设 $R$ 为环，$f \in R$。由 $\Spec$ 的函子性，映射
$R \to R_f$ 诱导同胚
$$
\Spec(R_f) \longrightarrow D(f) \subset \Spec(R).
$$
其逆映射由 $\mathfrak p \mapsto \mathfrak p \cdot R_f$ 给出。
\end{lemma}

\begin{proof}
这是引理 \ref{algebra-lemma-spec-localization} 的特殊情形。
\end{proof}

\noindent
谱的每个“仿射开集”并不都是标准开集。参见例
\ref{algebra-example-affine-open-not-standard}。

\begin{lemma}
\label{algebra-lemma-spec-closed}
设 $R$ 为环，$I \subset R$ 为理想。由 $\Spec$ 的函子性，映射
$R \to R/I$ 诱导同胚
$$
\Spec(R/I) \longrightarrow V(I) \subset \Spec(R).
$$
其逆映射由 $\mathfrak p \mapsto \mathfrak p / I$ 给出。
\end{lemma}

\begin{proof}
像显然包含于 $V(I)$。另一方面，若 $\mathfrak p \in V(I)$，则
$\mathfrak p \supset I$，从而可以考虑理想
$\mathfrak p /I \subset R/I$。由基本概念
(\ref{algebra-item-isomorphism-theorem})，可知
$(R/I)/(\mathfrak p/I) = R/\mathfrak p$ 为整环，所以
$\mathfrak p/I$ 是素理想。由此立即可见 $D(f + I)$ 的像为
$D(f) \cap V(I)$，因而该映射是同胚。
\end{proof}



\begin{lemma}
\label{algebra-lemma-quasi-compact}
\begin{slogan}
环的谱是拟紧的
\end{slogan}
设 $R$ 为环。空间 $\Spec(R)$ 是拟紧的。
\end{lemma}

\begin{proof}
% P01-E0292: refinement/subcover terminology candidate recorded in ERRATA_CANDIDATES.jsonl.
只需证明 $\Spec(R)$ 的任何标准开集覆盖都可由一个有限覆盖加细。设
$\Spec(R) = \cup D(f_i)$，其中 $\{f_i\}_{i\in I}$ 是 $R$ 中的一族
元素。这意味着 $\cap V(f_i) = \emptyset$。由引理
\ref{algebra-lemma-Zariski-topology}，这意味着
$V(\{f_i \}) = \emptyset$。再由同一引理，这意味着各 $f_i$ 生成
$R$ 的单位理想。也就是说，可以把 $1$ 写成一个{\it 有限}和：
$1 = \sum_{i \in J} r_i f_i$，其中 $J \subset I$ 有限。于是
$\Spec(R) = \cup_{i \in J} D(f_i)$。
\end{proof}

\begin{lemma}
\label{algebra-lemma-topology-spec}
设 $R$ 为环。$X = \Spec(R)$ 上的拓扑具有下列性质：
\begin{enumerate}
\item $X$ 拟紧；
\item $X$ 有一组由拟紧开集构成的拓扑基；
\item 任意两个拟紧开集的交仍拟紧。
\end{enumerate}
\end{lemma}

\begin{proof}
环的谱拟紧，参见引理 \ref{algebra-lemma-quasi-compact}。它有一组由标准开集
$D(f) = \Spec(R_f)$ 构成的拓扑基（引理
\ref{algebra-lemma-standard-open}），而由第一条可知这些开集拟紧。
两个标准开集的交拟紧，因为 $D(f) \cap D(g) = D(fg)$。
给定任意两个拟紧开集 $U, V \subset X$，可写成
$U = D(f_1) \cup \ldots \cup D(f_n)$ 以及
$V = D(g_1) \cup \ldots \cup D(g_m)$。于是
% P01-E0293: unindexed union candidate recorded in ERRATA_CANDIDATES.jsonl.
$U \cap V = \bigcup D(f_ig_j)$，它是拟紧的。
\end{proof}
\section{局部环}
\label{algebra-section-local-rings}

\noindent
局部环是代数几何中最基本而常用的工具。

\begin{definition}
\label{algebra-definition-local-ring}
恰有一个极大理想的环称为{\it 局部环}。若 $R$ 为局部环，通常把它的
极大理想记作 $\mathfrak m_R$，并把域 $R/\mathfrak m_R$ 称为局部环
$R$ 的{\it 剩余域}。我们常说“设 $(R, \mathfrak m)$ 为局部环”或
“设 $(R, \mathfrak m, \kappa)$ 为局部环”，以表示 $R$ 是局部环、
$\mathfrak m$ 是其唯一极大理想，而
$\kappa = R/\mathfrak m$ 是其剩余域。局部环之间的{\it 局部同态}
是环映射 $\varphi : R \to S$，其中 $R$ 与 $S$ 均为局部环，并且
$\varphi(\mathfrak m_R) \subset \mathfrak m_S$。若已知 $R$ 与 $S$
都是局部环，则“{\it 局部环映射 $\varphi : R \to S$}”意指
$\varphi$ 是局部环之间的局部同态。
\end{definition}

\noindent
域是局部环。域之间的任意环映射都是局部环之间的局部同态。

\medskip\noindent
环 $R$ 在素理想 $\mathfrak p$ 处的局部化 $R_\mathfrak p$ 是局部环，
其极大理想为 $\mathfrak p R_\mathfrak p$。事实上，由引理
\ref{algebra-lemma-spec-localization}，$R_\mathfrak p$ 的每个素理想都包含于
素理想 $\mathfrak p R_\mathfrak p$（所以后者是极大理想，也是
$R_\mathfrak p$ 唯一的极大理想）。$R_\mathfrak p$ 的剩余域记作
$\kappa(\mathfrak p)$；我们称它为{\it $\mathfrak p$ 的剩余域}。
由命题 \ref{algebra-proposition-localize-quotient}，可将 $\kappa(\mathfrak p)$
等同于整环 $R/\mathfrak p$ 的分式域。经由复合
$$
\Spec(\kappa(\mathfrak p)) \to \Spec(R_\mathfrak p) \to \Spec(R)
$$
源中的唯一一点映到靶中的点 $\mathfrak p$。

\medskip\noindent
设 $\varphi : R \to S$ 为环映射。令 $\mathfrak q \subset S$ 为
素理想，并考虑 $R$ 的素理想
$\mathfrak p = \varphi^{-1}(\mathfrak q)$。由于
$\varphi(\mathfrak p) \subset \mathfrak q$，诱导环映射
$$
R_\mathfrak p \to S_\mathfrak q,\quad r/g \mapsto \varphi(r)/\varphi(g)
$$
是局部环映射，并且得到剩余域间的诱导映射
$\kappa(\mathfrak p) \to \kappa(\mathfrak q)$。


\begin{example}
\label{algebra-example-not-local}
若 $R$ 为局部环，而 $\mathfrak p \subset R$ 是非极大的素理想，
则 $R \to R_\mathfrak p$ 不是局部同态。
\end{example}

\begin{lemma}
\label{algebra-lemma-characterize-local-ring}
设 $R$ 为环。下列条件等价：
\begin{enumerate}
\item $R$ 是局部环；
\item $\Spec(R)$ 恰有一个闭点；
\item $R$ 有一个极大理想 $\mathfrak m$，并且
$R \setminus \mathfrak m$ 的每个元素都是单位；
\item $R$ 不是零环，并且对每个 $x \in R$，$x$ 与 $1 - x$ 中至少
一个可逆，也可能二者均可逆。
\end{enumerate}
\end{lemma}

\begin{proof}
设 $R$ 为环，$\mathfrak m$ 为极大理想。若
$x \in R \setminus \mathfrak m$ 且 $x$ 不是单位，则存在包含
$x$ 的极大理想 $\mathfrak m'$。因此 $R$ 至少有两个极大理想。
反过来，若 $\mathfrak m'$ 是另一个极大理想，则可取
$x \in \mathfrak m'$ 且 $x \not \in \mathfrak m$。显然 $x$
不是单位。这证明了 (1) 与 (3) 等价。(1) 与 (2) 的等价是同义反复。
若 $R$ 局部，则 (4) 成立，因为 $x$ 或者属于 $\mathfrak m$，或者
不属于它。若 (4) 成立，而 $\mathfrak m$, $\mathfrak m'$ 是不同的
极大理想，则由中国剩余定理（引理
\ref{algebra-lemma-chinese-remainder}），可选取 $x \in R$ 使得
$x \bmod \mathfrak m' = 0$ 且 $x \bmod \mathfrak m = 1$。
此元素 $x$ 不可逆，而 $1 - x$ 也不可逆，矛盾。因此 (4) 与 (1)
等价。
\end{proof}

\begin{lemma}
\label{algebra-lemma-characterize-local-ring-map}
设 $\varphi : R \to S$ 为环映射，并假设 $R$ 与 $S$ 是局部环。
下列条件等价：
\begin{enumerate}
\item $\varphi$ 是局部环映射；
\item $\varphi(\mathfrak m_R) \subset \mathfrak m_S$；
\item $\varphi^{-1}(\mathfrak m_S) = \mathfrak m_R$；
\item 对任意 $x \in R$，若 $\varphi(x)$ 在 $S$ 中可逆，则 $x$
在 $R$ 中可逆。
\end{enumerate}
\end{lemma}

\begin{proof}
由定义，条件 (1) 与 (2) 等价。若 (3) 成立，则 (2) 成立。
反过来，若 (2) 成立，则 $\varphi^{-1}(\mathfrak m_S)$ 是包含极大理想
$\mathfrak m_R$ 的素理想，故
$\varphi^{-1}(\mathfrak m_S) = \mathfrak m_R$。最后，由引理
\ref{algebra-lemma-characterize-local-ring}，(4) 是 (2) 的逆否命题。
\end{proof}

\begin{remark}
\label{algebra-remark-fundamental-diagram}
与环映射 $\varphi : R \to S$ 及素理想 $\mathfrak p \subset R$
相关联的一个基本交换图如下：
$$
\xymatrix{
\kappa(\mathfrak p) \otimes_R S =
S_{\mathfrak p}/{\mathfrak p}S_{\mathfrak p}
&
S_{\mathfrak p} \ar[l]
&
S \ar[r] \ar[l]
&
S/\mathfrak pS \ar[r]
&
(R \setminus \mathfrak p)^{-1}S/\mathfrak pS
\\
\kappa(\mathfrak p) =
R_{\mathfrak p}/{\mathfrak p}R_{\mathfrak p} \ar[u]
&
R_{\mathfrak p} \ar[u] \ar[l]
&
R \ar[u] \ar[r] \ar[l]
&
R/\mathfrak p \ar[u] \ar[r]
&
\kappa(\mathfrak p) \ar[u]
}
$$
在此图中，最左列与最右列相同。在谱上，各水平映射诱导到其像的
同胚，而各方块诱导拓扑空间的纤维积方块（参见引理
\ref{algebra-lemma-spec-localization} 与 \ref{algebra-lemma-spec-closed}）。这说明
$\mathfrak p$ 属于谱上该映射的像，当且仅当
$S \otimes_R \kappa(\mathfrak p)$ 不是零环。若确实存在位于
$\mathfrak p$ 上方的素理想 $\mathfrak q \subset S$，即满足
$\mathfrak p = \varphi^{-1}(\mathfrak q)$，则可把该图扩充为下图：
$$
\xymatrix{
\kappa(\mathfrak q) = S_{\mathfrak q}/{\mathfrak q}S_{\mathfrak q}
&
S_{\mathfrak q} \ar[l]
&
S \ar[r] \ar[l]
&
S/\mathfrak q \ar[r]
&
\kappa(\mathfrak q)
\\
\kappa(\mathfrak p) \otimes_R S =
S_{\mathfrak p}/{\mathfrak p}S_{\mathfrak p} \ar[u]
&
S_{\mathfrak p} \ar[u] \ar[l]
&
S \ar[u] \ar[r] \ar[l]
&
S/\mathfrak pS \ar[u] \ar[r]
&
(R \setminus \mathfrak p)^{-1}S/\mathfrak pS \ar[u]
\\
\kappa(\mathfrak p) =
R_{\mathfrak p}/{\mathfrak p}R_{\mathfrak p} \ar[u]
&
R_{\mathfrak p} \ar[u] \ar[l]
&
R \ar[u] \ar[r] \ar[l]
&
R/\mathfrak p \ar[u] \ar[r]
&
\kappa(\mathfrak p) \ar[u]
}
$$
在此图中，最左列与最右列仍然相同；在谱上，各水平映射仍诱导到其像
的同胚。
\end{remark}

\begin{lemma}
\label{algebra-lemma-in-image}
设 $\varphi : R \to S$ 为环映射，$\mathfrak p$ 为 $R$ 的素理想。
下列条件等价：
\begin{enumerate}
\item $\mathfrak p$ 属于映射 $\Spec(S) \to \Spec(R)$ 的像；
\item $S \otimes_R \kappa(\mathfrak p) \not = 0$；
\item $S_{\mathfrak p}/\mathfrak p S_{\mathfrak p} \not = 0$；
\item $(S/\mathfrak pS)_{\mathfrak p} \not = 0$；
\item $\mathfrak p = \varphi^{-1}(\mathfrak pS)$。
\end{enumerate}
\end{lemma}

\begin{proof}
注 \ref{algebra-remark-fundamental-diagram} 已经给出了前两个条件的等价性；
其余条件只是它的不同表述。
\end{proof}

\begin{remark}
\label{algebra-remark-local-ring-fibre}
设 $R \to S$ 为环映射。令 $\mathfrak q$ 为 $S$ 的素理想，并位于
$R$ 的素理想 $\mathfrak p$ 上方。由注
\ref{algebra-remark-fundamental-diagram}，素理想 $\mathfrak q$ 对应于纤维环
$F = S \otimes_R \kappa(\mathfrak p)$ 的唯一素理想
$\overline{\mathfrak q}$。于是
$$
F_{\overline{\mathfrak q}}
\cong S_\mathfrak q \otimes_{R_\mathfrak p} \kappa(\mathfrak p)
\cong S_\mathfrak q/\mathfrak p S_\mathfrak q
$$
事实上，有一个显然的环映射
$F \to S_\mathfrak q \otimes_{R_\mathfrak p} \kappa(\mathfrak p)$，
容易看出它同构于 $F \to F_{\overline{\mathfrak q}}$。第二个同构来自
$\kappa(\mathfrak p)$ 是 $R_\mathfrak p$ 对
$\mathfrak pR_\mathfrak p$ 的商这一事实。
\end{remark}
\section{环的雅可布森根}
\label{algebra-section-radical}

\noindent
回顾环 $R$ 的{\it 雅可布森根} $\text{rad}(R)$ 是 $R$ 的全体极大理想
之交。若 $R$ 局部，则 $\text{rad}(R)$ 就是 $R$ 的极大理想。

\begin{lemma}
\label{algebra-lemma-contained-in-radical}
设 $R$ 为环，其雅可布森根为 $\text{rad}(R)$。令 $I \subset R$
为理想。下列条件等价：
\begin{enumerate}
\item $I \subset \text{rad}(R)$；
\item $1 + I$ 的每个元素都是 $R$ 中的单位。
\end{enumerate}
在这种情形下，$R$ 中每个映为 $R/I$ 中单位的元素本身都是单位。
\end{lemma}

\begin{proof}
若 $f \in \text{rad}(R)$，则对 $R$ 的每个极大理想
$\mathfrak m$ 都有 $f \in \mathfrak m$。所以对 $R$ 的每个极大理想
$\mathfrak m$ 都有 $1 + f \not \in \mathfrak m$。因此
$\Spec(R)$ 的闭子集 $V(1 + f)$ 为空。这说明 $1 + f$ 是单位，
参见引理 \ref{algebra-lemma-Zariski-topology}。

\medskip\noindent
反过来，假设对所有 $f \in I$，$1 + f$ 都是单位。若
$\mathfrak m$ 是极大理想且 $I \not \subset \mathfrak m$，则
$I + \mathfrak m = R$。故对某个 $g \in \mathfrak m$ 与
$f \in I$ 有 $1 = f + g$。于是 $g = 1 + (-f)$ 不是单位，矛盾。

\medskip\noindent
对最后一个断言，设 $f \in R$ 映为 $R/I$ 中的单位。可以找到
$g \in R$，它映为 $f \bmod I$ 的乘法逆元。于是
$fg = 1 \bmod I$。由 (2)，$fg$ 是 $R$ 中的单位，从而 $f$ 是单位。
\end{proof}

\begin{lemma}
\label{algebra-lemma-surjective-on-spec-units}
设 $\varphi : R \to S$ 为环映射，并且诱导映射
$\Spec(S) \to \Spec(R)$ 为满射。则元素 $x \in R$ 是单位，当且仅当
$\varphi(x) \in S$ 是单位。
\end{lemma}

\begin{proof}
若 $x$ 是单位，则 $\varphi(x)$ 也是单位。反过来，若
$\varphi(x)$ 是单位，则对所有 $\mathfrak q \in \Spec(S)$ 都有
$\varphi(x) \not \in \mathfrak q$。因此对所有
$\mathfrak q \in \Spec(S)$ 都有
$x \not \in \varphi^{-1}(\mathfrak q) = \Spec(\varphi)(\mathfrak q)$。
由于 $\Spec(\varphi)$ 为满射，由引理
\ref{algebra-lemma-Zariski-topology} 的第 (17) 项可知 $x$ 是单位。
\end{proof}




\section{中山引理}
\label{algebra-section-nakayama}

\noindent
我们引述 \cite{MatCA}：“这个简单而重要的引理归功于
T.~Nakayama、G.~Azumaya 和 W.~Krull。其优先权并不明确；虽然它通常
称为中山引理，已故的 Nakayama 教授却不喜欢这个名称。”

\begin{lemma}[中山引理]
\label{algebra-lemma-NAK}
\begin{reference}
\cite[1.M Lemma (NAK) page 11]{MatCA}
\end{reference}
\begin{history}
我们引述 \cite{MatCA}：“这个简单而重要的引理归功于
T.~Nakayama、G.~Azumaya 和 W.~Krull。其优先权并不明确；虽然它通常
称为中山引理，已故的 Nakayama 教授却不喜欢这个名称。”
\end{history}
设 $R$ 为环，其雅可布森根为 $\text{rad}(R)$。令 $M$ 为 $R$-模，
$I \subset R$ 为理想。
\begin{enumerate}
\item
\label{algebra-item-nakayama}
若 $IM = M$ 且 $M$ 有限，则存在 $f \in 1 + I$ 使得 $fM = 0$。
\item 若 $IM = M$、$M$ 有限且 $I \subset \text{rad}(R)$，则 $M = 0$。
\item 若 $N, N' \subset M$、$M = N + IN'$ 且 $N'$ 有限，则存在
$f \in 1 + I$，使得 $fM \subset N$ 且 $M_f = N_f$。
\item 若 $N, N' \subset M$、$M = N + IN'$、$N'$ 有限且
$I \subset \text{rad}(R)$，则 $M = N$。
\item 若 $N \to M$ 为模映射、$N/IN \to M/IM$ 为满射且 $M$ 有限，
则存在 $f \in 1 + I$，使得 $N_f \to M_f$ 为满射。
\item 若 $N \to M$ 为模映射、$N/IN \to M/IM$ 为满射、$M$ 有限且
$I \subset \text{rad}(R)$，则 $N \to M$ 为满射。
\item 若 $x_1, \ldots, x_n \in M$ 生成 $M/IM$ 且 $M$ 有限，则存在
$f \in 1 + I$，使得 $x_1, \ldots, x_n$ 在 $R_f$ 上生成 $M_f$。
\item 若 $x_1, \ldots, x_n \in M$ 生成 $M/IM$、$M$ 有限且
$I \subset \text{rad}(R)$，则 $M$ 由 $x_1, \ldots, x_n$ 生成。
\item 若 $IM = M$ 且 $I$ 幂零，则 $M = 0$。
\item 若 $N, N' \subset M$、$M = N + IN'$ 且 $I$ 幂零，则 $M = N$。
\item 若 $N \to M$ 为模映射、$I$ 幂零且
$N/IN \to M/IM$ 为满射，则 $N \to M$ 为满射。
\item 若 $\{x_\alpha\}_{\alpha \in A}$ 是 $M$ 中生成 $M/IM$ 的一族
元素且 $I$ 幂零，则 $M$ 由各 $x_\alpha$ 生成。
\end{enumerate}
\end{lemma}

\begin{proof}
证明 (\ref{algebra-item-nakayama})。选取 $M$ 在 $R$ 上的一组生成元
$y_1, \ldots, y_m$。对每个 $i$，可写成
$y_i = \sum z_{ij} y_j$，其中 $z_{ij} \in I$（因为 $M = IM$）。
换言之，$\sum_j (\delta_{ij} - z_{ij})y_j = 0$。令 $f$ 为
$m \times m$ 矩阵 $A = (\delta_{ij} - z_{ij})$ 的行列式。注意
$f \in 1 + I$（因为矩阵 $A$ 的每个分量模 $I$ 同余于
$m \times m$ 单位矩阵的相应分量）。由引理
\ref{algebra-lemma-matrix-left-inverse} (1)，存在 $m \times m$ 矩阵 $B$，
使得 $BA = f 1_{m \times m}$。展开可见，对所有 $h$ 与 $j$，有
$\sum_{i} b_{hi} a_{ij} = f \delta_{hj}$；因而对每个 $h$，
$\sum_{i, j} b_{hi} a_{ij} y_j
= \sum_{j} f \delta_{hj} y_j = f y_h$。换言之，对每个 $h$ 都有
$0 = f y_h$（因为每个 $i$ 都满足 $\sum_j a_{ij} y_j = 0$）。
这说明 $f$ 零化 $M$。

\medskip\noindent
% P01-E0294: missing article in source unit statement recorded in ERRATA_CANDIDATES.jsonl.
由引理 \ref{algebra-lemma-contained-in-radical}，$1 + \text{rad}(R)$ 的元素是
$R$ 的可逆元。因此 (\ref{algebra-item-nakayama}) 蕴含 (2)。把 (1) 用于
$M/N$ 得到 (3)，这里由于 $N'$ 有限，该商有限。把 (2) 用于
$M/N$ 得到 (4)，这里同样由于 $N'$ 有限，该商有限。把 (3) 用于
$M$ 及其子模 $\Im(N \to M)$ 和 $M$，得到 (5)。把 (4) 用于 $M$
及其子模 $\Im(N \to M)$ 和 $M$，得到 (6)。把 (5) 用于映射
$R^{\oplus n} \to M$，$(a_1, \ldots, a_n) \mapsto
a_1x_1 + \ldots + a_nx_n$，得到 (7)。把 (6) 用于映射
$R^{\oplus n} \to M$，$(a_1, \ldots, a_n) \mapsto
a_1x_1 + \ldots + a_nx_n$，得到 (8)。

\medskip\noindent
第 (9) 项成立，因为若 $M = IM$，则对所有 $n \geq 0$ 都有
$M = I^nM$，而 $I$ 幂零意味着对某个 $n \gg 0$ 有 $I^n = 0$。
第 (10)、(11)、(12) 项由 (9) 及上述论证推出。
\end{proof}

\begin{lemma}
\label{algebra-lemma-NAK-localization}
设 $R$ 为环，$S \subset R$ 为乘法子集，$I \subset R$ 为理想，
$M$ 为有限 $R$-模。若 $x_1, \ldots, x_r \in M$ 作为
$S^{-1}(R/I)$-模生成 $S^{-1}(M/IM)$，则存在 $f \in S + I$，使得
$x_1, \ldots, x_r$ 作为 $R_f$-模生成 $M_f$。\footnote{特殊情形：
(I) $I = 0$。引理说明，若 $x_1, \ldots, x_r$ 生成 $S^{-1}M$，
则对某个 $f \in S$，$x_1, \ldots, x_r$ 生成 $M_f$。
(II) $I = \mathfrak p$ 是素理想且 $S = R \setminus \mathfrak p$。
引理说明，若 $x_1, \ldots, x_r$ 生成
$M \otimes_R \kappa(\mathfrak p)$，则对某个满足 $f \in R$、
$f \not \in \mathfrak p$ 的元素，$x_1, \ldots, x_r$ 生成 $M_f$。}
\end{lemma}

\begin{proof}
先看特殊情形 $I = 0$。设 $y_1, \ldots, y_s$ 是 $M$ 在 $R$ 上的
生成元。由于 $S^{-1}M$ 由 $x_1, \ldots, x_r$ 生成，对每个 $i$，
可在 $S^{-1}M$ 中写成
$y_i = \sum (a_{ij}/s_{ij})x_j$，其中 $a_{ij} \in R$ 且
$s_{ij} \in S$。乘以各 $s_{ij}$ 的乘积 $s \in S$，可见在
$S^{-1}M$ 中 $sy_i = \sum a'_{ij}x_j$，其中
$a'_{ij} \in R$。这又意味着存在 $t_i \in S$，使得在 $M$ 中有
$t_isy_i = \sum t_ia'_{ij}x_j$。因此，若 $t \in S$ 是各 $t_i$ 的
乘积，则在 $M_{st}$ 中，$y_i$ 属于 $x_1, \ldots, x_r$ 生成的
$R_{st}$-子模。故 $x_1, \ldots, x_r$ 生成 $M_{st}$。
% P01-E0295: subject--verb agreement defect recorded in ERRATA_CANDIDATES.jsonl.

\medskip\noindent
一般情形。由特殊情形，可找到 $s \in S$，使得
$x_1, \ldots, x_r$ 在 $(R/I)_s$ 上生成 $(M/IM)_s$。由引理
\ref{algebra-lemma-NAK}，可找到 $g \in 1 + I_s \subset R_s$，使得
$x_1, \ldots, x_r$ 在 $(R_s)_g$ 上生成 $(M_s)_g$。写成
$g = 1 + i/s'$。则 $f = ss' + is$ 满足要求；细节从略。
\end{proof}

\begin{lemma}
\label{algebra-lemma-when-surjective-local}
设 $A \to B$ 为局部环之间的局部同态。假设：
\begin{enumerate}
\item $B$ 作为 $A$-模有限；
\item $\mathfrak m_B$ 是有限生成理想；
\item $A \to B$ 诱导剩余域之间的同构；
\item $\mathfrak m_A/\mathfrak m_A^2 \to
\mathfrak m_B/\mathfrak m_B^2$ 为满射。
\end{enumerate}
则 $A \to B$ 为满射。
\end{lemma}

\begin{proof}
为证明 $A \to B$ 满射，把它视为 $A$-模映射并应用引理
\ref{algebra-lemma-NAK} (6)。由此只需证明
$A/\mathfrak m_A \to B/\mathfrak m_AB$ 为满射。由于
$A/\mathfrak m_A = B/\mathfrak m_B$，只需证明
$\mathfrak m_AB \to \mathfrak m_B$ 为满射。把
$\mathfrak m_AB \to \mathfrak m_B$ 视为 $B$-模映射并应用引理
\ref{algebra-lemma-NAK} (6)。于是只需看出
$\mathfrak m_AB/\mathfrak m_A\mathfrak m_B \to
\mathfrak m_B/\mathfrak m_B^2$ 为满射。这由假设 (4) 推出。
\end{proof}
\section{谱的开闭子集}
\label{algebra-section-open-and-closed}

\noindent
事实表明，谱中既开又闭的子集与环的幂等元相对应。

\begin{lemma}
\label{algebra-lemma-idempotent-spec}
设 $R$ 为环，$e \in R$ 为幂等元。此时
$$
\Spec(R) = D(e) \amalg D(1-e).
$$
\end{lemma}

\begin{proof}
注意，整环中的幂等元 $e$ 只能是 $1$ 或 $0$。因此
\begin{eqnarray*}
D(e)
& = &
\{ \mathfrak p \in \Spec(R)
\mid
e \not\in \mathfrak p \} \\
& = &
\{ \mathfrak p \in \Spec(R)
\mid
e \not = 0\text{ 于 }\kappa(\mathfrak p) \} \\
& = &
\{ \mathfrak p \in \Spec(R)
\mid
e = 1\text{ 于 }\kappa(\mathfrak p) \}
\end{eqnarray*}
类似地，
\begin{eqnarray*}
D(1-e)
& = &
\{ \mathfrak p \in \Spec(R)
\mid
1 - e \not\in \mathfrak p \} \\
& = &
\{ \mathfrak p \in \Spec(R)
\mid
e \not = 1\text{ 于 }\kappa(\mathfrak p) \} \\
& = &
\{ \mathfrak p \in \Spec(R)
\mid
e = 0\text{ 于 }\kappa(\mathfrak p) \}
\end{eqnarray*}
由于 $e$ 在任意剩余域中的像只能是 $1$ 或 $0$，可知 $D(e)$ 与
$D(1-e)$ 覆盖整个 $\Spec(R)$。
\end{proof}

\begin{lemma}
\label{algebra-lemma-spec-product}
设 $R_1$ 与 $R_2$ 为环，令 $R = R_1 \times R_2$。映射
$R \to R_1$、$(x, y) \mapsto x$ 以及 $R \to R_2$、
$(x, y) \mapsto y$ 分别诱导连续映射
$\Spec(R_1) \to \Spec(R)$ 与 $\Spec(R_2) \to \Spec(R)$。
诱导映射
$$
\Spec(R_1) \amalg \Spec(R_2)
\longrightarrow
\Spec(R)
$$
是同胚。换言之，$R = R_1\times R_2$ 的谱是 $R_1$ 的谱与
$R_2$ 的谱的不交并。
\end{lemma}

\begin{proof}
写成 $1 = e_1 + e_2$，其中 $e_1 = (1, 0)$ 且
$e_2 = (0, 1)$。注意，$e_1$ 与 $e_2 = 1 - e_1$ 都是幂等元。
留给读者证明 $R_1 = R_{e_1}$ 是 $R$ 在 $e_1$ 处的局部化。
$e_2$ 的情形类似。因此，引理的断言由引理
\ref{algebra-lemma-idempotent-spec} 与引理 \ref{algebra-lemma-standard-open} 合并推出。
\end{proof}

\noindent
引入函数的粘合引理后，我们还会再次证明下一个引理。参见节
\ref{algebra-section-tilde-module-sheaf}。

\begin{lemma}
\label{algebra-lemma-disjoint-decomposition}
设 $R$ 为环。对每个既开又闭的 $U \subset \Spec(R)$，存在唯一的
幂等元 $e \in R$，使得 $U = D(e)$。这给出既开又闭的子集
$U \subset \Spec(R)$ 与幂等元 $e \in R$ 之间的一一对应。
\end{lemma}

\begin{proof}
设 $U \subset \Spec(R)$ 既开又闭。因为 $U$ 闭，由引理
\ref{algebra-lemma-quasi-compact}，它拟紧；其补集同样拟紧。把它写成
标准开集的有限并 $U = \bigcup_{i = 1}^n D(f_i)$。类似地，写成
$\Spec(R) \setminus U = \bigcup_{j = 1}^m D(g_j)$。由于
$\emptyset = D(f_i) \cap D(g_j) = D(f_i g_j)$，由引理
\ref{algebra-lemma-Zariski-topology}，$f_i g_j$ 幂零。令
$I = (f_1, \ldots, f_n) \subset R$，并令
$J = (g_1, \ldots, g_m) \subset R$。注意，$V(J)$ 等于 $U$，
$V(I)$ 等于 $U$ 的补集，所以
$\Spec(R) = V(I) \amalg V(J)$。由上述幂零性可知，对某个充分大的
整数 $N$，有 $(IJ)^N = (0)$。由于
$\bigcup D(f_i) \cup \bigcup D(g_j) = \Spec(R)$，由引理
\ref{algebra-lemma-Zariski-topology} 可知 $I + J = R$。把这个等式取
$2N$ 次幂，得到 $I^N + J^N = R$。写成 $1 = x + y$，其中
$x \in I^N$ 且 $y \in J^N$。因为 $I^N J^N = (0)$，有
$0 = xy = x(1 - x)$。故 $x = x^2$ 是幂等元，并且属于
$I^N \subset I$。幂等元 $y = 1 - x$ 属于 $J^N \subset J$。
这说明，对每个 $\mathfrak p \in V(J)$，幂等元 $x$ 在剩余域
$\kappa(\mathfrak p)$ 中映为 $1$；而对每个
$\mathfrak p \in V(I)$，$x$ 在 $\kappa(\mathfrak p)$ 中映为 $0$。

\medskip\noindent
为证明唯一性，设 $e_1, e_2$ 是 $R$ 中不同的幂等元。必须证明存在
素理想 $\mathfrak p$，使得 $e_1 \in \mathfrak p$ 而
$e_2 \not \in \mathfrak p$，或反过来。写
$e_i' = 1 - e_i$。若 $e_1 \not = e_2$，则
$0 \not = e_1 - e_2  = e_1(e_2 + e_2') - (e_1 + e_1')e_2
= e_1 e_2' - e_1' e_2$。所以幂等元 $e_1 e_2' \not = 0$ 或
$e_1' e_2 \not = 0$ 至少有一个成立。
% P01-E0296: nonzero qualification on idempotent recorded in ERRATA_CANDIDATES.jsonl.
幂等元不幂零，因而由引理 \ref{algebra-lemma-Zariski-topology} 可找到素理想
$\mathfrak p$，使得 $e_1e_2' \not \in \mathfrak p$ 或
$e_1'e_2 \not \in \mathfrak p$。容易看出这给出了所需的素理想。
\end{proof}

\begin{lemma}
\label{algebra-lemma-characterize-spec-connected}
设 $R$ 为非零环。则 $\Spec(R)$ 连通，当且仅当 $R$ 没有非平凡
幂等元。
\end{lemma}

\begin{proof}
这由引理 \ref{algebra-lemma-disjoint-decomposition} 及连通拓扑空间的定义
立即得到。
\end{proof}

\begin{lemma}
\label{algebra-lemma-ideal-is-squared-union-connected}
设 $I \subset R$ 为环 $R$ 的有限生成理想，并且 $I = I^2$。则：
\begin{enumerate}
\item 存在幂等元 $e \in R$，使得 $I = (e)$；
\item 对幂等元 $e' = 1 - e \in R$，有 $R/I \cong R_{e'}$；
\item $V(I)$ 在 $\Spec(R)$ 中既开又闭。
\end{enumerate}
\end{lemma}

\begin{proof}
由中山引理 \ref{algebra-lemma-NAK}，存在元素 $f = 1 + i$、$i \in I$，
使得 $fI = 0$。于是 $f^2 = f + fi = f$，故它为幂等元。
考虑幂等元 $e = 1 - f = -i \in I$。对 $j \in I$，有
$ej = j - fj = j$，所以 $I = (e)$。这证明了 (1)。

\medskip\noindent
(2) 与 (3) 由 (1) 推出。事实上，
$V(I) = V(e) = \Spec(R) \setminus D(e)$；由引理
\ref{algebra-lemma-idempotent-spec} 或引理
\ref{algebra-lemma-disjoint-decomposition}，它既开又闭。这证明了 (3)。
为证明 (2)，注意映射 $R \to R_{e'}$ 为满射，因为在 $R_{e'}$ 中
$x/(e')^n = x/e' = xe'/(e')^2 = xe'/e' = x/1$。映射
$R \to R_{e'}$ 的核是 $R$ 中被 $e'$ 的某个正整数次幂零化的元素之集。
由于 $e'$ 幂等，这就是被 $e'$ 零化的元素所成的理想；又因为
$e + e' = 1$ 是一对正交幂等元，该理想就是 $I = (e)$。这证明了 (2)。
\end{proof}
\section{谱的连通分支}
\label{algebra-section-connected-components}

\noindent
谱的连通分支并不像初看时可能以为的那样容易理解。这是因为我们习惯于
局部连通空间的拓扑，而环的谱一般并不局部连通。

\begin{lemma}
\label{algebra-lemma-closed-union-connected-components}
设 $R$ 为环，$T \subset \Spec(R)$ 为谱的子集。下列条件等价：
\begin{enumerate}
\item $T$ 闭，并且是 $\Spec(R)$ 的若干连通分支之并；
\item $T$ 是 $\Spec(R)$ 中若干既开又闭子集之交；
\item $T = V(I)$，其中 $I \subset R$ 是由幂等元生成的理想。
\end{enumerate}
此外，若 (3) 中的理想存在，则它是唯一的。
\end{lemma}

\begin{proof}
由引理 \ref{algebra-lemma-topology-spec} 及拓扑，引理
\ref{topology-lemma-closed-union-connected-components}，可知 (1) 与
(2) 等价。假设 (2)，并写成 $T = \bigcap U_\alpha$，其中
$U_\alpha \subset \Spec(R)$ 既开又闭。由引理
\ref{algebra-lemma-disjoint-decomposition}，对某个幂等元
$e_\alpha \in R$ 有 $U_\alpha = D(e_\alpha)$。令
$I = (1 - e_\alpha)$，便有 $T = V(I)$，即 (3) 成立。最后假设 (3)。
写成 $T = V(I)$，并令 $I = (e_\alpha)$，其中各 $e_\alpha$ 为幂等元。
于是显然
$T = \bigcap V(e_\alpha) = \bigcap D(1 - e_\alpha)$。

\medskip\noindent
设 $I$ 是由幂等元生成的理想。令 $e \in R$ 为满足
$V(I) \subset V(e)$ 的幂等元。由引理 \ref{algebra-lemma-Zariski-topology}，
对某个 $n \geq 1$ 有 $e^n \in I$。由于 $e$ 幂等，这意味着
$e \in I$。因此，$I$ 恰由满足 $T \subset V(e)$ 的那些幂等元 $e$
生成。换言之，闭子集 $T$ 完全确定理想 $I$，这证明了唯一性。
\end{proof}

\begin{lemma}
\label{algebra-lemma-connected-component}
设 $R$ 为环。$\Spec(R)$ 的一个连通分支形如 $V(I)$，其中 $I$
是由幂等元生成的理想，并且 $R$ 的每个幂等元在 $R/I$ 中都映为
$0$ 或 $1$。
\end{lemma}

\begin{proof}
设 $\mathfrak p$ 为 $R$ 的素理想。由引理
\ref{algebra-lemma-topology-spec}，可知拓扑，引理
\ref{topology-lemma-connected-component-intersection} 的假设对拓扑空间
% P01-E0297: source verb-form defect recorded in ERRATA_CANDIDATES.jsonl.
$\Spec(R)$ 成立。因此，$\mathfrak p$ 在 $\Spec(R)$ 中的连通分支是
$\Spec(R)$ 中包含 $\mathfrak p$ 的所有既开又闭子集之交。所以它等于
$V(I)$，其中 $I$ 由幂等元 $e \in R$ 生成，而 $e$ 在
$\kappa(\mathfrak p)$ 中映为 $0$，参见引理
\ref{algebra-lemma-disjoint-decomposition}。
显然，不属于这一族的任意幂等元 $e$ 在 $R/I$ 中映为 $1$。
\end{proof}
\section{粘合性质}
\label{algebra-section-more-glueing}

\noindent
本节汇集若干如下形式的标准结果：若某个性质在一个标准开覆盖的每个
成员上成立，则它整体成立。事实上，往往只需在局部环层面检验，正如下
一个引理所示。

\begin{lemma}
\label{algebra-lemma-characterize-zero-local}
设 $R$ 为环。
\begin{enumerate}
\item 对 $R$-模 $M$ 的元素 $x$，下列条件等价：
\begin{enumerate}
\item $x = 0$；
\item 对所有 $\mathfrak p \in \Spec(R)$，$x$ 在 $M_\mathfrak p$ 中
的像为零；
\item 对 $R$ 的所有极大理想 $\mathfrak m$，$x$ 在
$M_{\mathfrak m}$ 中的像为零。
\end{enumerate}
换言之，映射 $M \to \prod_{\mathfrak m} M_{\mathfrak m}$ 为单射。
\item 给定 $R$-模 $M$，下列条件等价：
\begin{enumerate}
\item $M$ 为零；
\item 对所有 $\mathfrak p \in \Spec(R)$，$M_{\mathfrak p}$ 为零；
\item 对 $R$ 的所有极大理想 $\mathfrak m$，$M_{\mathfrak m}$ 为零。
\end{enumerate}
\item 给定 $R$-模的复形 $M_1 \to M_2 \to M_3$，下列条件等价：
\begin{enumerate}
\item $M_1 \to M_2 \to M_3$ 正合；
\item 对 $R$ 的每个素理想 $\mathfrak p$，局部化
$M_{1, \mathfrak p} \to M_{2, \mathfrak p} \to M_{3, \mathfrak p}$
正合；
\item 对 $R$ 的每个极大理想 $\mathfrak m$，局部化
$M_{1, \mathfrak m} \to M_{2, \mathfrak m} \to M_{3, \mathfrak m}$
正合。
\end{enumerate}
\item 给定 $R$-模映射 $f : M \to M'$，下列条件等价：
\begin{enumerate}
\item $f$ 为单射；
\item 对 $R$ 的所有素理想 $\mathfrak p$，
$f_{\mathfrak p} : M_\mathfrak p \to M'_\mathfrak p$ 为单射；
\item 对 $R$ 的所有极大理想 $\mathfrak m$，
$f_{\mathfrak m} : M_\mathfrak m \to M'_\mathfrak m$ 为单射。
\end{enumerate}
\item 给定 $R$-模映射 $f : M \to M'$，下列条件等价：
\begin{enumerate}
\item $f$ 为满射；
\item 对 $R$ 的所有素理想 $\mathfrak p$，
$f_{\mathfrak p} : M_\mathfrak p \to M'_\mathfrak p$ 为满射；
\item 对 $R$ 的所有极大理想 $\mathfrak m$，
$f_{\mathfrak m} : M_\mathfrak m \to M'_\mathfrak m$ 为满射。
\end{enumerate}
\item 给定 $R$-模映射 $f : M \to M'$，下列条件等价：
\begin{enumerate}
\item $f$ 为双射；
\item 对 $R$ 的所有素理想 $\mathfrak p$，
$f_{\mathfrak p} : M_\mathfrak p \to M'_\mathfrak p$ 为双射；
\item 对 $R$ 的所有极大理想 $\mathfrak m$，
$f_{\mathfrak m} : M_\mathfrak m \to M'_\mathfrak m$ 为双射。
\end{enumerate}
\end{enumerate}
\end{lemma}

\begin{proof}
设 $x \in M$ 如 (1) 所述。令 $I = \{f \in R \mid fx = 0\}$。
容易看出 $I$ 是理想（它是 $x$ 的零化子）。条件 (1)(c) 意味着，
对每个极大理想 $\mathfrak m$，存在 $f \in R \setminus \mathfrak m$
使得 $fx =0$。换言之，$V(I)$ 不包含闭点。由引理
\ref{algebra-lemma-Zariski-topology} 可知 $I$ 是单位理想。因此 $x$ 为零，
即 (1)(a) 成立。这证明了 (1)。

\medskip\noindent
把 (1) 同时应用于 $M$ 的所有元素，即得 (2)。

\medskip\noindent
证明 (3)。令 $H$ 为该序列的同调，即
$H = \Ker(M_2 \to M_3)/\Im(M_1 \to M_2)$。由命题
\ref{algebra-proposition-localization-exact}，$H_\mathfrak p$ 是序列
$M_{1, \mathfrak p} \to M_{2, \mathfrak p} \to M_{3, \mathfrak p}$
的同调。因此 (3) 是 (2) 的推论。

\medskip\noindent
(4) 与 (5) 是 (3) 的特殊情形。把 (4) 与 (5) 合并，形式地得到 (6)。
\end{proof}

\begin{lemma}
\label{algebra-lemma-cover}
\begin{slogan}
模与代数的扎里斯基局部性质
\end{slogan}
设 $R$ 为环，$M$ 为 $R$-模，$S$ 为 $R$-代数。假设
$f_1, \ldots, f_n$ 是 $R$ 中的有限元素列，并满足
$\bigcup D(f_i) = \Spec(R)$，换言之，$(f_1, \ldots, f_n) = R$。
\begin{enumerate}
\item 若每个 $M_{f_i} = 0$，则 $M = 0$。
\item 若每个 $M_{f_i}$ 都是有限 $R_{f_i}$-模，则 $M$ 是有限
$R$-模。
\item 若每个 $M_{f_i}$ 都是有限表示 $R_{f_i}$-模，则 $M$ 是有限表示
$R$-模。
\item 设 $M \to N$ 为 $R$-模映射。若对每个 $i$，
$M_{f_i} \to N_{f_i}$ 都是同构，则 $M \to N$ 是同构。
\item 设 $0 \to M'' \to M \to M' \to 0$ 为 $R$-模的复形。
若对每个 $i$，$0 \to M''_{f_i} \to M_{f_i} \to M'_{f_i} \to 0$
都正合，则 $0 \to M'' \to M \to M' \to 0$ 正合。
\item 若每个 $R_{f_i}$ 都是诺特环，则 $R$ 是诺特环。
\item 若每个 $S_{f_i}$ 都是有限型 $R_{f_i}$-代数，则 $S$ 是有限型
$R$-代数。
\item 若每个 $S_{f_i}$ 都在 $R_{f_i}$ 上有限表示，则 $S$ 是有限表示
$R$-代数。
\end{enumerate}
\end{lemma}

\begin{proof}
依次证明各项。
\begin{enumerate}
\item 由命题 \ref{algebra-proposition-localize-twice}，该假设蕴含对所有
$\mathfrak p \in \Spec(R)$ 都有 $M_\mathfrak p = 0$，故由引理
\ref{algebra-lemma-characterize-zero-local} 得出结论。
\item 对每个 $i$，取 $M_{f_i}$ 的有限生成集 $X_i$。不失一般性，
可假设 $X_i$ 的元素都在局部化映射 $M \rightarrow M_{f_i}$ 的像中，
于是取 $M$ 中由 $X_i$ 各元素的原像组成的有限集合 $Y_i$。令 $Y$
为这些集合之并，它仍是有限集。考虑显然的 $R$-线性映射
$R^Y \rightarrow M$，把基元素 $e_y$ 映到 $y$。依假设，在 $R$
的任意素理想 $\mathfrak p$ 处局部化后，此映射都是满射；故由引理
\ref{algebra-lemma-characterize-zero-local}，它本身满射，从而 $M$ 有限生成。
\item 由 (2)，有短正合列
$$
0 \rightarrow K \rightarrow R^n \rightarrow M \rightarrow 0
$$
由于局部化是正合函子，且 $M_{f_i}$ 有限表示，由引理
\ref{algebra-lemma-extension}，对所有 $1 \leq i \leq n$，$K_{f_i}$ 都有限生成。
% P01-E0298: overloaded cover/generator index recorded in ERRATA_CANDIDATES.jsonl.
由 (2)，这说明 $K$ 是有限 $R$-模，因而 $M$ 有限表示。
\item 由命题 \ref{algebra-proposition-localize-twice}，该假设蕴含在所有素理想
处的诱导局部化态射都是同构，所以由引理
\ref{algebra-lemma-characterize-zero-local} 得出结论。
\item 由命题 \ref{algebra-proposition-localize-twice}，该假设蕴含在所有素理想
处的诱导局部化序列都短正合，所以由引理
\ref{algebra-lemma-characterize-zero-local} 得出结论。
\item 我们证明 $R$ 的每个理想都有有限生成集。为此，令
$I \subset R$ 为任意理想。由命题
\ref{algebra-proposition-localization-exact}，每个
$I_{f_i} \subset R_{f_i}$ 都是理想。依假设，它们都有限生成，故由
(2) 得出结论。
\item 对每个 $i$，取 $S_{f_i}$ 的有限生成集 $X_i$。不失一般性，
可假设 $X_i$ 的元素都在局部化映射 $S \rightarrow S_{f_i}$ 的像中，
于是取 $S$ 中由 $X_i$ 各元素的原像组成的有限集合 $Y_i$。令 $Y$
为这些集合之并，它仍是有限集。考虑由 $Y$ 诱导的代数同态
$R[X_y]_{y \in Y} \rightarrow S$。因为它是代数同态，其像 $T$ 是
$R$-模 $S$ 的 $R$-子模，所以可以考虑商模 $S/T$。依假设，在各
$f_i$ 处局部化后它为零，故由 (1) 它为零，因而 $S$ 是有限型
$R$-代数。
\item 由上一项，存在满的 $R$-代数同态
$R[X_1, \ldots, X_n] \rightarrow S$。令 $K$ 为该映射的核。
这是 $R[X_1, \ldots, X_n]$ 中的理想，并且在每个 $f_i$ 处局部化后
有限生成。由于各 $f_i$ 在 $R$ 中生成单位理想，它们在
$R[X_1, \ldots, X_n]$ 中也生成单位理想，所以应用 (2) 即完成证明。
\end{enumerate}
\end{proof}

\begin{lemma}
\label{algebra-lemma-cover-upstairs}
设 $R \to S$ 为环映射。假设 $g_1, \ldots, g_n$ 是 $S$ 中的有限
元素列，并满足 $\bigcup D(g_i) = \Spec(S)$，换言之，
$(g_1, \ldots, g_n) = S$。
\begin{enumerate}
\item 若每个 $S_{g_i}$ 都是有限型 $R$-代数，则 $S$ 是有限型
$R$-代数。
\item 若每个 $S_{g_i}$ 都在 $R$ 上有限表示，则 $S$ 在 $R$ 上
有限表示。
\end{enumerate}
\end{lemma}

\begin{proof}
选取 $h_1, \ldots, h_n \in S$，使得 $\sum h_i g_i = 1$。

\medskip\noindent
证明 (1)。对每个 $i$，选取元素的有限列
$x_{i, j} \in S_{g_i}$、$j = 1, \ldots, m_i$，使其作为
$R$-代数生成 $S_{g_i}$。对某个 $y_{i, j} \in S$ 及某个
$n_{i, j} \ge 0$，写成
$x_{i, j} = y_{i, j}/g_i^{n_{i, j}}$。考虑由
$g_1, \ldots, g_n$、$h_1, \ldots, h_n$ 以及
$y_{i, j}$（$i = 1, \ldots, n$、$j = 1, \ldots, m_i$）生成的
$R$-子代数 $S' \subset S$。由于局部化正合（命题
\ref{algebra-proposition-localization-exact}），可知
$S'_{g_i} \to S_{g_i}$ 为单射。另一方面，由 $y_{i, j}$ 的选取，
它是满射。由于 $h_1, \ldots, h_n \in S'$，各元素
$g_1, \ldots, g_n$ 在 $S'$ 中生成单位理想。因此，把
$S' \to S$ 视为 $S'$-模映射，由引理 \ref{algebra-lemma-cover}，它是同构。

\medskip\noindent
证明 (2)。已经知道 $S$ 是有限型的。对某个理想 $J$，写成
$S = R[x_1, \ldots, x_m]/J$。对每个 $i$，选取 $g_i$ 的提升
$g'_i \in R[x_1, \ldots, x_m]$，并选取 $h_i$ 的提升
$h'_i \in R[x_1, \ldots, x_m]$。于是
$$
S_{g_i} = R[x_1, \ldots, x_m, y_i]/(J_i + (1 - y_ig'_i))
$$
其中 $J_i$ 是 $R[x_1, \ldots, x_m, y_i]$ 中由 $J$ 生成的理想。
略去一个小细节。由引理 \ref{algebra-lemma-finite-presentation-independent}，
可以选取有限元素列 $f_{i, j} \in J$、$j = 1, \ldots, m_i$，
使得 $f_{i, j}$ 在 $J_i$ 中的像与 $1 - y_ig'_i$ 一起生成理想
$J_i + (1 - y_ig'_i)$。令
$$
S' = R[x_1, \ldots, x_m]/\left(\sum h'_ig'_i - 1, f_{i, j}; 
i = 1, \ldots, n, j = 1, \ldots, m_i\right)
$$
存在满的 $R$-代数映射 $S' \to S$。各元素
$g'_1, \ldots, g'_n$ 在 $S'$ 中的类生成单位理想，并且由构造，映射
$S'_{g'_i} \to S_{g_i}$ 都是单射。因此可以像 (1) 中那样得出结论。
\end{proof}
\section{函数的粘合}
\label{algebra-section-tilde-module-sheaf}

\noindent
本节将证明：给定由标准开集组成的开覆盖
$$
\Spec(R) = \bigcup\nolimits_{i = 1}^n D(f_i)
$$
以及对每个 $i$ 给定元素 $h_i \in R_{f_i}$，若作为
$R_{f_i f_j}$ 的元素有 $h_i = h_j$，则存在唯一的 $h \in R$，
使得 $h$ 在 $R_{f_i}$ 中的像为 $h_i$。这个结果可以作两种解释：
\begin{enumerate}
\item 规则 $D(f) \mapsto R_f$ 在标准开集上给出一个环层，参见层，
节 \ref{sheaves-section-bases}。
\item 若把 $R_f$ 的元素看成 $D(f)$ 上的“代数”函数或“正则”函数，
那么它们可以像拓扑流形上的连续函数、相应地像可微流形上的可微函数
那样粘合。
\end{enumerate}

\begin{lemma}
\label{algebra-lemma-cover-module}
设 $R$ 为环，$f_1, \ldots, f_n$ 是 $R$ 中生成单位理想的元素，
$M$ 为 $R$-模。序列
$$
0 \to
M \xrightarrow{\alpha}
\bigoplus\nolimits_{i = 1}^n M_{f_i} \xrightarrow{\beta}
\bigoplus\nolimits_{i, j = 1}^n M_{f_i f_j}
$$
正合，其中 $\alpha(m) = (m/1, \ldots, m/1)$，并且
$\beta(m_1/f_1^{e_1}, \ldots, m_n/f_n^{e_n})
= (m_i/f_i^{e_i} - m_j/f_j^{e_j})_{(i, j)}$。
\end{lemma}

\begin{proof}
只需证明该序列在任意极大理想 $\mathfrak m$ 处的局部化正合，参见
引理 \ref{algebra-lemma-characterize-zero-local}。由于
$f_1, \ldots, f_n$ 生成单位理想，存在 $i$ 使得
$f_i \not \in \mathfrak m$。重新编号后可假设 $i = 1$。注意
$(M_{f_i})_\mathfrak m = (M_\mathfrak m)_{f_i}$ 以及
$(M_{f_if_j})_\mathfrak m = (M_\mathfrak m)_{f_if_j}$，参见命题
\ref{algebra-proposition-localize-twice-module}。特别地，因为 $f_1$ 是单位，
有 $(M_{f_1})_\mathfrak m = M_\mathfrak m$ 以及
$(M_{f_1 f_i})_\mathfrak m = (M_\mathfrak m)_{f_i}$。注意，序列中的
映射是由引理 \ref{algebra-lemma-universal-property-localization-module} 及
$M$ 上的恒等映射给出的典范映射。综上，把 $R$ 替换为
$R_\mathfrak m$、把 $M$ 替换为 $M_\mathfrak m$、把 $f_i$ 替换为其
在 $R_\mathfrak m$ 中的像，并把 $f_1$ 替换为
$1 \in R_\mathfrak m$，便归约到 $f_1 = 1$ 的情形。

\medskip\noindent
假设 $f_1 = 1$。此时 $\alpha$ 的单射性显然。设
$m = (m_i) \in \bigoplus_{i = 1}^n M_{f_i}$ 属于 $\beta$ 的核。
则 $m_1 \in M_{f_1} = M$。此外，$\beta(m) = 0$ 蕴含 $m_1$ 与
$m_i$ 在 $M_{f_1f_i} = M_{f_i}$ 中映为同一元素。因此
$\alpha(m_1) = m$，证明完成。
\end{proof}

\begin{lemma}
\label{algebra-lemma-standard-covering}
设 $R$ 为环，并设 $f_1, f_2, \ldots f_n\in R$ 在 $R$ 中生成
单位理想。则下列序列正合：
% P01-E0299: missing separator in source element list recorded in ERRATA_CANDIDATES.jsonl.
$$
0 \longrightarrow
R \longrightarrow
\bigoplus\nolimits_i R_{f_i} \longrightarrow
\bigoplus\nolimits_{i, j}R_{f_if_j}
$$
其中映射 $\alpha : R \longrightarrow \bigoplus_i R_{f_i}$ 与
$\beta : \bigoplus_i R_{f_i} \longrightarrow \bigoplus_{i, j} R_{f_if_j}$
定义为
$$
\alpha(x) = \left(\frac{x}{1}, \ldots, \frac{x}{1}\right)
\text{ 且 }
\beta\left(\frac{x_1}{f_1^{r_1}}, \ldots, \frac{x_n}{f_n^{r_n}}\right)
=
\left(\frac{x_i}{f_i^{r_i}}-\frac{x_j}{f_j^{r_j}}~\text{于}~R_{f_if_j}\right).
$$
\end{lemma}

\begin{proof}
这是引理 \ref{algebra-lemma-cover-module} 的特殊情形。
\end{proof}

\noindent
下面的结果我们在上文已经见过，但为了方便，现将它明确陈述出来。

\begin{lemma}
\label{algebra-lemma-disjoint-implies-product}
设 $R$ 为环。若 $\Spec(R) = U \amalg V$，且 $U$ 与 $V$ 都开，
则 $R \cong R_1 \times R_2$，并且经由引理
\ref{algebra-lemma-spec-product} 中的映射，有 $U \cong \Spec(R_1)$ 与
$V \cong \Spec(R_2)$。此外，$R_1$ 与 $R_2$ 都既是环 $R$ 的局部化，
又是它的商。
\end{lemma}

\begin{proof}
由引理 \ref{algebra-lemma-disjoint-decomposition}，对某个幂等元 $e$，有
$U = D(e)$ 及 $V = D(1-e)$。由引理
\ref{algebra-lemma-standard-covering} 可知
$R \cong R_e \times R_{1 - e}$（因为显然 $R_{e(1-e)} = 0$，所以
粘合条件是平凡的；当然，在这种情形下直接证明该乘积分解也很容易）。
引理得证。
\end{proof}

\begin{lemma}
\label{algebra-lemma-when-injective-covering}
设 $R$ 为环。设 $f_1, \ldots, f_n \in R$。
设 $M$ 为 $R$-模。则 $M \to \bigoplus M_{f_i}$ 是单射，当且仅当
$$
M \longrightarrow \bigoplus\nolimits_{i = 1, \ldots, n} M, \quad
m \longmapsto (f_1m, \ldots, f_nm)
$$
是单射。
\end{lemma}

\begin{proof}
映射 $M \to \bigoplus M_{f_i}$ 是单射，当且仅当对于所有
$m \in M$ 及 $e_1, \ldots, e_n \geq 1$，若
$f_i^{e_i}m = 0$，$i = 1, \ldots, n$，则 $m = 0$。
这显然推出所显示的映射是单射。反之，假设所显示的映射是单射，且
$m \in M$、$e_1, \ldots, e_n \geq 1$ 满足
$f_i^{e_i}m = 0$，$i = 1, \ldots, n$。若所有 $i$ 都有
$e_i = 1$，则由所显示映射的单射性立即得到 $m = 0$。
接下来对 $e = \sum e_i$ 归纳，证明任意这样的数据都有同一结论。
基始情形是 $e = n$，已在上面处理。若某个 $e_i > 1$，令
$m' = f_im$。由归纳假设可知 $m' = 0$。因此
$f_i m = 0$，即可以取 $e_i = 1$，这使 $e$ 减小，结论得证。
\end{proof}

\noindent
下面的引理最好在平坦下降的更一般语境中陈述和证明。不过它与上文十分契合，
所以在此陈述也很自然。

\begin{lemma}
\label{algebra-lemma-glue-modules}
设 $R$ 为环，$f_1, \ldots, f_n \in R$。假设给定以下数据：
\begin{enumerate}
\item 对每个 $i$，给定一个 $R_{f_i}$-模 $M_i$。
\item 对每一对 $i, j$，给定一个 $R_{f_if_j}$-模同构
$\psi_{ij} : (M_i)_{f_j} \to (M_j)_{f_i}$。
\end{enumerate}
并且这些数据满足“上圈条件”，即所有图表
$$
\xymatrix{
(M_i)_{f_jf_k}
\ar[rd]_{\psi_{ij}}
\ar[rr]^{\psi_{ik}}
& &
(M_k)_{f_if_j} \\
&
(M_j)_{f_if_k} \ar[ru]_{\psi_{jk}}
}
$$
均交换（对所有三元组 $i, j, k$）。由这些数据定义
$$
M = \Ker\left(
\bigoplus\nolimits_{1 \leq i \leq n} M_i
\longrightarrow
\bigoplus\nolimits_{1 \leq i, j \leq n} (M_i)_{f_j}
\right)
$$
其中 $(m_1, \ldots, m_n)$ 映到一个元素，其第 $(i, j)$ 个分量为
$m_i/1 - \psi_{ji}(m_j/1)$。则自然映射 $M \to M_i$ 诱导同构
$M_{f_i} \to M_i$。此外，对所有 $m \in M$，有
$\psi_{ij}(m/1) = m/1$（采用显然的记号）。
\end{lemma}

\begin{proof}
要证明 $M_{f_1} \to M_1$ 是同构，只须证明它在
$R_{f_1}$ 的每个素理想 $\mathfrak p'$ 处的局部化都是同构，见引理
\ref{algebra-lemma-characterize-zero-local}。由引理 \ref{algebra-lemma-standard-open}，
可写成 $\mathfrak p' = \mathfrak p R_{f_1}$，其中某个素理想
$\mathfrak p \subset R$ 满足 $f_1 \not \in \mathfrak p$。
由于局部化是正合的（命题 \ref{algebra-proposition-localization-exact}），有
\begin{align*}
(M_{f_1})_{\mathfrak p'} & =
M_\mathfrak p \\
& =
\Ker\left(
\bigoplus\nolimits_{1 \leq i \leq n} M_{i, \mathfrak p}
\longrightarrow
\bigoplus\nolimits_{1 \leq i, j \leq n} ((M_i)_{f_j})_\mathfrak p
\right) \\
& =
\Ker\left(
\bigoplus\nolimits_{1 \leq i \leq n} M_{i, \mathfrak p}
\longrightarrow
\bigoplus\nolimits_{1 \leq i, j \leq n} (M_{i, \mathfrak p})_{f_j}
\right)
\end{align*}
这里还使用了命题 \ref{algebra-proposition-localize-twice-module}。
由于 $f_1$ 在 $R_\mathfrak p$ 中是单位，将 $R$ 替换为
$R_\mathfrak p$、将 $f_i$ 替换为其在 $R_\mathfrak p$ 中的像 $f_i$、
将 $M$ 替换为 $M_\mathfrak p$，并将 $f_1$ 替换为 $1$，
便约化到 $f_1 = 1$ 的情形。

\medskip\noindent
假设 $f_1 = 1$。则对 $j = 2, \ldots, n$，
$\psi_{1j} : (M_1)_{f_j} \to M_j$ 是同构。
若用这些同构作恒等识别 $M_j = (M_1)_{f_j}$，则
$\psi_{ij} : (M_1)_{f_if_j} \to (M_1)_{f_if_j}$
就是典范恒等识别。因此复形
$$
0 \to M_1 \to
\bigoplus\nolimits_{1 \leq i \leq n} (M_1)_{f_i}
\longrightarrow
\bigoplus\nolimits_{1 \leq i, j \leq n}
(M_1)_{f_if_j}
$$
由引理 \ref{algebra-lemma-cover-module} 是正合的。因此在这种情形下，第一个映射将
$M_1$ 与 $M$ 等同，其余结论也都清楚。
\end{proof}

\section{零因子与全分式环}
\label{algebra-section-total-quotient-ring}

\noindent
极小素理想处的局部环具有下列性质。

\begin{lemma}
\label{algebra-lemma-minimal-prime-reduced-ring}
设 $\mathfrak p$ 为环 $R$ 的极小素理想。$R_{\mathfrak p}$ 的极大理想
中的每个元素都幂零。若 $R$ 既约，则 $R_{\mathfrak p}$ 是域。
\end{lemma}

\begin{proof}
若 ${\mathfrak p}R_{\mathfrak p}$ 的某个元素 $x$ 不幂零，则
$D(x) \not = \emptyset$，参见引理 \ref{algebra-lemma-Zariski-topology}。
这与 $\mathfrak p$ 的极小性矛盾。若 $R$ 既约，则
${\mathfrak p}R_{\mathfrak p} = 0$，因而该局部环是域。
\end{proof}

\begin{lemma}
\label{algebra-lemma-reduced-ring-sub-product-fields}
设 $R$ 为既约环。则：
\begin{enumerate}
\item $R$ 是某个域乘积的子环；
\item $R \to \prod_{\mathfrak p\text{ minimal}} R_{\mathfrak p}$
把它嵌入一个域的乘积；
\item $\bigcup_{\mathfrak p\text{ minimal}} \mathfrak p$ 是 $R$ 的
零因子之集。
\end{enumerate}
\end{lemma}

\begin{proof}
由引理 \ref{algebra-lemma-minimal-prime-reduced-ring}，每个环
$R_\mathfrak p$ 都是域。特别地，环映射 $R \to R_\mathfrak p$
的核是 $\mathfrak p$。由引理 \ref{algebra-lemma-Zariski-topology}，有
$\bigcap_{\mathfrak p} \mathfrak p = (0)$。所以 (2) 与 (1) 成立。
若 $x y = 0$ 且 $y \not = 0$，则存在极小素理想 $\mathfrak p$，
使得 $y \not \in \mathfrak p$。因此 $x \in \mathfrak p$。所以
$R$ 的每个零因子都属于
$\bigcup_{\mathfrak p\text{ minimal}} \mathfrak p$。反过来，假设对某个
极小素理想 $\mathfrak p$ 有 $x \in \mathfrak p$。于是 $x$ 在
$R_\mathfrak p$ 中映为零，故存在 $y \in R$、
$y \not \in \mathfrak p$，使得 $xy = 0$。换言之，$x$ 是零因子。
这完成了 (3) 及整个引理的证明。
\end{proof}

\noindent
环 $R$ 的全分式环 $Q(R)$ 已在例
\ref{algebra-example-localize-at-prime} 中引入。

\begin{lemma}
\label{algebra-lemma-total-ring-fractions}
设 $R$ 为环，$S \subset R$ 是由非零因子组成的乘法子集。则
$Q(R) \cong Q(S^{-1}R)$。特别地，$Q(R) \cong Q(Q(R))$。
\end{lemma}

\begin{proof}
若 $x \in S^{-1}R$ 是非零因子，并且对某个 $r \in R$、$f \in S$
有 $x = r/f$，则 $r$ 是 $R$ 中的非零因子。引理随即得证。
\end{proof}

\noindent
可以应用粘合结果证明关于全分式环 $Q(R)$ 的一个结论；该环已在例
\ref{algebra-example-localize-at-prime} 中引入。

\begin{lemma}
\label{algebra-lemma-total-ring-fractions-no-embedded-points}
设 $R$ 为环。假设 $R$ 有有限多个极小素理想
$\mathfrak q_1, \ldots, \mathfrak q_t$，并且
$\mathfrak q_1 \cup \ldots \cup \mathfrak q_t$ 是 $R$ 的零因子之集。
则全分式环 $Q(R)$ 等于
$R_{\mathfrak q_1} \times \ldots \times R_{\mathfrak q_t}$。
\end{lemma}

\begin{proof}
因为每个非零因子都属于 $R \setminus \mathfrak q_i$，存在自然映射
% P01-E0300: element/subset wording defect recorded in ERRATA_CANDIDATES.jsonl.
$Q(R) \to R_{\mathfrak q_i}$。因而有自然映射
$Q(R) \to R_{\mathfrak q_1} \times \ldots \times R_{\mathfrak q_t}$。
对任意非极小素理想 $\mathfrak p \subset R$，由引理
\ref{algebra-lemma-silly} 可知
$\mathfrak p \not \subset \mathfrak q_1 \cup \ldots \cup \mathfrak q_t$。
因此 $\Spec(Q(R)) = \{\mathfrak q_1, \ldots, \mathfrak q_t\}$
（将二者视为 $\Spec(R)$ 的子集，参见引理
\ref{algebra-lemma-spec-localization}）。所以 $\Spec(Q(R))$ 是有限离散集，
从而 $Q(R) = A_1 \times \ldots \times A_t$，其中
$\Spec(A_i) = \{\mathfrak{q}_i\}$，参见引理
\ref{algebra-lemma-disjoint-implies-product}。此外，$A_i$ 是 $R$ 的一个局部化，
并且是局部环。因此 $A_i \cong R_{\mathfrak q_i}$。
\end{proof}
\section{谱的不可约分支}
\label{algebra-section-irreducible}

\noindent
我们将证明，环谱的不可约分支对应于环中的极小素理想。

\begin{lemma}
\label{algebra-lemma-irreducible}
设 $R$ 为环。
\begin{enumerate}
\item 对素理想 $\mathfrak p \subset R$，单点集 $\{\mathfrak p\}$
在扎里斯基拓扑中的闭包为 $V(\mathfrak p)$。用公式表示即
$\overline{\{\mathfrak p\}} = V(\mathfrak p)$。
\item $\Spec(R)$ 的不可约闭子集恰为各子集 $V(\mathfrak p)$，其中
$\mathfrak p \subset R$ 为素理想。
\item $\Spec(R)$ 的不可约分支（参见拓扑，定义
\ref{topology-definition-irreducible-components}）恰为各子集
$V(\mathfrak p)$，其中 $\mathfrak p \subset R$ 为极小素理想。
\end{enumerate}
\end{lemma}

\begin{proof}
注意，若 $ \mathfrak p \in V(I)$，则 $I \subset \mathfrak p$。
所以显然 $\overline{\{\mathfrak p\}} = V(\mathfrak p)$。特别地，
$V(\mathfrak p)$ 是单点集的闭包，因而不可约。第二个断言蕴含第三个。
为证明第二个断言，设 $V(I) \subset \Spec(R)$，其中 $I$ 是根理想。
若 $I$ 不是素理想，则选取 $a, b\in R$、$a, b\not \in I$，使得
$ab\in I$。此时 $V(I, a) \cup V(I, b) = V(I)$；但由引理
\ref{algebra-lemma-Zariski-topology}，$V(I, b) = V(I)$ 与
$V(I, a) = V(I)$ 都不成立。因此 $V(I)$ 不可约。
\end{proof}

\noindent
换言之，该引理表明，$\Spec(R)$ 的每个不可约闭子集都形如
$V(\mathfrak p)$，其中 $\mathfrak p$ 为某个素理想。由于
$V(\mathfrak p) = \overline{\{\mathfrak p\}}$，每个不可约闭子集都有
唯一泛点，参见拓扑，定义 \ref{topology-definition-generic-point}。
特别地，$\Spec(R)$ 是清醒拓扑空间。把这个事实记录为下列引理。

\begin{lemma}
\label{algebra-lemma-spec-spectral}
环的谱是谱空间；参见拓扑，定义
\ref{topology-definition-spectral-space}。
\end{lemma}

\begin{proof}
形式上，这由引理 \ref{algebra-lemma-irreducible} 与引理
\ref{algebra-lemma-topology-spec} 推出。另见上文的讨论。
\end{proof}

\begin{lemma}
\label{algebra-lemma-irreducible-components-containing-x}
设 $R$ 为环，$\mathfrak p \subset R$ 为素理想。
\begin{enumerate}
\item $\Spec(R)$ 中经过 $\mathfrak p$ 的不可约闭子集之集，与素理想
$\mathfrak q \subset R_{\mathfrak p}$ 一一对应。
\item $\Spec(R)$ 中经过 $\mathfrak p$ 的不可约分支之集，与极小素理想
$\mathfrak q \subset R_{\mathfrak p}$ 一一对应。
\end{enumerate}
\end{lemma}

\begin{proof}
这由引理 \ref{algebra-lemma-irreducible} 及引理
\ref{algebra-lemma-spec-localization} 中对 $\Spec(R_\mathfrak p)$ 的描述推出；
后者说明 $\Spec(R_\mathfrak p)$ 对应于 $R$ 中满足
$\mathfrak q \subset \mathfrak p$ 的素理想 $\mathfrak q$。
\end{proof}

\begin{lemma}
\label{algebra-lemma-standard-open-containing-maximal-point}
设 $R$ 为环，$\mathfrak p$ 为 $R$ 的极小素理想。令
$W \subset \Spec(R)$ 为不包含点 $\mathfrak p$ 的拟紧开集。则存在
$f \in R$、$f \not \in \mathfrak p$，使得
$D(f) \cap W = \emptyset$。
\end{lemma}

\begin{proof}
因为 $W$ 拟紧，可以把它写成标准仿射开集 $D(g_i)$ 的有限并，其中
$i = 1, \ldots, n$。由于 $\mathfrak p \not \in W$，对所有 $i$ 都有
$g_i \in \mathfrak p$。由引理
\ref{algebra-lemma-minimal-prime-reduced-ring}，每个 $g_i$ 在
$R_{\mathfrak p}$ 中都幂零。因此可以找到
$f \in R$、$f \not \in \mathfrak p$，使得对所有 $i$，某个
$n_i > 0$ 满足 $f g_i^{n_i} = 0$。这个 $D(f)$ 即满足要求。
\end{proof}

\begin{lemma}
\label{algebra-lemma-ring-with-only-minimal-primes}
设 $R$ 为环，并把 $X = \Spec(R)$ 视为拓扑空间。下列条件等价：
\begin{enumerate}
\item $X$ 仿有限；
\item $X$ 是 Hausdorff 空间；
\item $X$ 全不连通；
\item $X$ 的每个拟紧开集都闭；
\item 其素理想之间不存在非平凡包含关系；
\item 每个素理想都是极大理想；
\item 每个素理想都是极小素理想；
\item 每个标准开集 $D(f) \subset X$ 都闭；
% P01-E0301: unfinished source placeholder recorded in ERRATA_CANDIDATES.jsonl.
\item 此处再添加条件。
\end{enumerate}
\end{lemma}

\begin{proof}
第一种证明。显然，(5)、(6)、(7) 等价。由于每个拟紧开集都是标准
开集的有限并，(4) 与 (8) 显然等价。蕴含关系 (7) $\Rightarrow$ (4)
由引理 \ref{algebra-lemma-standard-open-containing-maximal-point} 得到。
假设 (4) 成立。令 $\mathfrak p, \mathfrak p'$ 为 $R$ 的两个不同
素理想。选取 $f \in \mathfrak p'$、$f \not \in \mathfrak p$
（必要时交换 $\mathfrak p$ 与 $\mathfrak p'$）。于是
$\mathfrak p' \not \in D(f)$ 而 $\mathfrak p \in D(f)$。由 (4)，
开集 $D(f)$ 也闭。因此 $\mathfrak p$ 与 $\mathfrak p'$ 分别位于两个
不交开邻域中，且这两个邻域之并为 $X$。所以 $X$ 是 Hausdorff 且
全不连通的。因此 (4) $\Rightarrow$ (2) 和 (3)。若 (3) 成立，则
$\Spec(R)$ 的点之间不可能有特化关系，从而 (5) 成立。若 $X$ 是
Hausdorff 空间，则每个点都闭，所以 (2) 蕴含 (6)。因此 (2)、(3)、
(4)、(5)、(6)、(7)、(8) 等价。每个仿有限空间都是 Hausdorff 的，
所以 (1) 蕴含 (2)。若 $X$ 满足 (2) 与 (3)，则由引理
\ref{algebra-lemma-quasi-compact}，$X$ 拟紧；再由拓扑，引理
\ref{topology-lemma-profinite}，可知该空间仿有限。

\medskip\noindent
第二种证明。除 (4) 与 (8) 的等价外，其余结论由引理
\ref{algebra-lemma-spec-spectral} 及纯拓扑事实推出；参见拓扑，引理
\ref{topology-lemma-characterize-profinite-spectral}。
\end{proof}
\section{环谱的例}
\label{algebra-section-examples-spectra}

\noindent
本节给出若干谱的例子。

\begin{example}
\label{algebra-example-spec-Zxmodx2minus4}
本例描述 $X = \Spec(\mathbf{Z}[x]/(x^2 - 4))$。令 $\mathfrak{p}$ 为
$X$ 中的任意素理想，并令
$\phi : \mathbf{Z} \to \mathbf{Z}[x]/(x^2 - 4)$ 为自然环映射。于是
$ \phi^{-1}(\mathfrak p)$ 是 $\mathbf{Z}$ 中的素理想。若
$ \phi^{-1}(\mathfrak p) = (2)$，则因为 $\mathfrak p$ 包含 $2$，
它经由映射
$ \mathbf{Z}[x]/(x^2 - 4) \to  \mathbf{Z}[x]/(x^2 - 4, 2)$ 对应于
$$
\mathbf{Z}[x]/(x^2 - 4, 2) \cong (\mathbf{Z}/2\mathbf{Z})[x]/(x^2)
$$
中的素理想。$(\mathbf{Z}/2\mathbf{Z})[x]/(x^2)$ 中的任意素理想
对应于 $(\mathbf{Z}/2\mathbf{Z})[x]$ 中包含 $(x^2)$ 的素理想，因而
这样的素理想包含 $x$。因为
$(\mathbf{Z}/2\mathbf{Z}) \cong (\mathbf{Z}/2\mathbf{Z})[x]/(x)$ 是域，
$(x)$ 是极大理想。任意素理想都包含 $(x)$，而 $(x)$ 极大，故该环
仅有一个素理想 $(x)$。所以在这种情形下
$\mathfrak p = (2, x)$。现若 $ \phi^{-1}(\mathfrak p) = (q)$，其中
$q > 2$，则因为 $\mathfrak p$ 包含 $q$，它经由映射
$ \mathbf{Z}[x]/(x^2 - 4) \to  \mathbf{Z}[x]/(x^2 - 4, q)$ 对应于
$$
\mathbf{Z}[x]/(x^2 - 4, q) \cong (\mathbf{Z}/q\mathbf{Z})[x]/(x^2 - 4)
$$
中的素理想。$(\mathbf{Z}/q\mathbf{Z})[x]/(x^2 - 4)$ 中的任意素理想
对应于 $(\mathbf{Z}/q\mathbf{Z})[x]$ 中包含
$(x^2 - 4) = (x -2)(x + 2)$ 的素理想。因此，这些素理想必须包含
$x -2$ 或 $x + 2$。因为 $(\mathbf{Z}/q\mathbf{Z})[x]$ 是主理想整环，
所有非零素理想都极大，故 $(\mathbf{Z}/q\mathbf{Z})[x]$ 中恰有两个包含
$(x-2)(x + 2)$ 的素理想，
即 $(x-2)$ 与 $(x + 2)$。总之，存在两个素理想
$(q, x-2)$ 与 $(q, x + 2)$，因为
$2 \neq -2 \in \mathbf{Z}/(q)$。最后处理
$\phi^{-1}(\mathfrak p) = (0)$ 的情形。注意，$\mathfrak p$ 对应于
$\mathbf{Z}[x]$ 中包含
$(x^2 - 4) = (x -2)(x + 2)$ 的素理想。因此 $\mathfrak p$ 包含
$(x-2)$ 或 $(x + 2)$。所以依假设，$\mathfrak p$ 对应于
$\mathbf{Z}[x]/(x - 2)$ 中的一个素理想，或
$\mathbf{Z}[x]/(x + 2)$ 中的一个素理想，并且它与 $\mathbf{Z}$ 仅交于
$0$。由于 $\mathbf{Z}[x]/(x - 2) \cong \mathbf{Z}$ 且
$\mathbf{Z}[x]/(x + 2) \cong \mathbf{Z}$，这意味着 $\mathfrak p$
必须对应于其中一个环中的 $0$。因此在原环中，
$\mathfrak p = (x - 2)$ 或 $\mathfrak p = (x + 2)$。
\end{example}

\begin{example}
\label{algebra-example-spec-Zx}
本例描述 $X = \Spec(\mathbf{Z}[x])$。固定 $\mathfrak p \in X$。
令 $\phi : \mathbf{Z} \to \mathbf{Z}[x]$，并注意
$\phi^{-1}(\mathfrak p) \in \Spec(\mathbf{Z})$。若
$\phi^{-1}(\mathfrak p) = (q)$，其中 $q$ 是满足 $q > 0$ 的素数，
则 $\mathfrak p$ 对应于 $(\mathbf{Z}/(q))[x]$ 中的素理想，而该理想
必须由一个在 $(\mathbf{Z}/(q))[x]$ 中不可约的多项式生成。
% P01-E0302: omitted zero prime in positive-characteristic fibre recorded in ERRATA_CANDIDATES.jsonl.
若选择这个多项式的最小次数代表元，则它在 $\mathbf{Z}[x]$ 中也不可约。
% P01-E0303: primitive/monic lift qualification recorded in ERRATA_CANDIDATES.jsonl.
因此在这种情形下，$\mathfrak p = (q, f_q)$，其中 $f_q$ 是
$\mathbf{Z}[x]$ 中的不可约多项式，并且将其视为
$(\mathbf{Z}/(q) [x])$ 中的多项式时仍不可约。现假设
$\phi^{-1}(\mathfrak p) = (0)$。在这种情形下，$\mathfrak p$ 必须由
非常数多项式生成；由于 $\mathfrak p$ 为素理想，可假设这些多项式在
$\mathbf{Z}[x]$ 中不可约。
% P01-E0304: omitted zero ideal in generic fibre recorded in ERRATA_CANDIDATES.jsonl.
由高斯引理，这些多项式在 $\mathbf{Q}[x]$ 中也不可约。因为
$\mathbf{Q}[x]$ 是欧几里得整环，若至少有两个不同的不可约元 $f, g$
生成 $\mathfrak p$，则对某些 $a, b \in \mathbf{Q}[x]$，有
$1 = af + bg$。乘以一个公分母，可见对
$\bar{a}, \bar{b} \in \mathbf{Z}[x]$ 及某个非零的
$m \in \mathbf{Z}$，有 $m = \bar{a}f + \bar{b} g$，矛盾。因此
$\mathfrak p$ 由 $\mathbf{Z}[x]$ 中的一个不可约多项式生成。
\end{example}

\begin{example}
\label{algebra-example-spec-kxy}
设 $k$ 为任意域。本例描述 $X = \Spec(k[x, y])$。显然 $(0)$ 为
素理想；由于 $k[x, y]$ 是唯一分解整环，由不可约多项式生成的任何
主理想也都是素理想。现假设 $\mathfrak p$ 是 $X$ 中的非主理想。
因为 $k[x, y]$ 是诺特唯一分解整环，素理想 $\mathfrak p$ 可由有限个
不可约多项式 $(f_1, \ldots, f_n)$ 生成。现断言：若 $f, g$ 是
$k[x, y]$ 中不相伴的不可约多项式，则
$(f, g) \cap k[x] \neq 0$。为此，只需证明把 $f$ 与 $g$ 视为
$k(x)[y]$ 中的多项式时互素。在这种情形下，$k(x)[y]$ 是欧几里得整环；
应用欧几里得算法并清除分母，得到
$p = af + bg$，其中 $p, a, b \in k[x]$。
% P01-E0305: coefficient-ring defect after clearing denominators recorded in ERRATA_CANDIDATES.jsonl.
所以假设情况并非如此，即存在某个非单位 $h \in k(x)[y]$ 同时整除
$f$ 与 $g$。于是由高斯引理，对某些 $a, b \in k(x)$，有
$ah | f$ 与 $bh | g$，其中 $ah, bh \in k[x]$。
% P01-E0306: ambient polynomial-ring defect recorded in ERRATA_CANDIDATES.jsonl.
由不可约性，$ah = f$ 且 $bh = g$（因为 $h \notin k(x)$）。所以回到
$k(x)[y]$ 中，$f, g $ 相伴，因为 $\frac{a}{b} g = f$。由于 $k(x)$
是 $k[x]$ 的分式域，可对没有公因子的元素 $r , s \in k[x]$ 写成
$g = \frac{r}{s} f $。这蕴含在 $k[x, y]$ 中有 $sg = rf$；由于
$k[x, y]$ 是唯一分解整环，$s$ 必须整除 $f$。故 $s = 1$ 或
$s = f$。若 $s = f$，则 $r = g$，这意味着 $f, g \in k[x]$，因而
它们在 $k(x)$ 中必为单位，并在 $k(x)[y]$ 中互素，与假设矛盾。
若 $s = 1$，则 $g = rf$，仍为矛盾。因此 $f, g$ 必须在欧几里得整环
$k(x)[y]$ 中互素。由此归约到 $\mathfrak p$ 包含某个不可约多项式
$p \in k[x] \subset k[x, y]$ 的情形。由上文，$\mathfrak p$ 对应于环
$k[x, y]/(p) = k(\alpha)[y]$ 中的素理想，其中 $\alpha$ 是 $k$ 上的
代数元，其极小多项式为 $p$。这是主理想整环，故任一素理想对应于
$(0)$ 或 $k(\alpha)[y]$ 中的一个不可约多项式。因此
$\mathfrak p$ 形如 $(p)$ 或 $(p, f)$，其中 $f$ 是 $k[x, y]$ 中的
多项式，并且在商 $k[x, y]/(p)$ 中不可约。
\end{example}

\begin{example}
\label{algebra-example-affine-open-not-standard}
考虑环
$$
R = \{ f \in \mathbf{Q}[z]\text{ 其中 }f(0) = f(1) \}.
$$
考虑映射
$$
\varphi : \mathbf{Q}[A, B] \to R
$$
其中 $\varphi(A) = z^2-z$ 且 $\varphi(B) = z^3-z^2$。容易验证
$(A^3 - B^2 + AB) \subset \Ker(\varphi)$，并且
$A^3 - B^2 + AB$ 不可约。假设 $\varphi$ 满射；那么，因为 $R$ 是
整环（它是某个整环的子环），$\Ker(\varphi)$ 必须是
$\mathbf{Q}[A, B]$ 的素理想。包含 $(A^3-B^2 + AB)$ 的素理想是
$(A^3-B^2 + AB)$ 本身，以及任意满足如下条件的极大理想 $(f, g)$：
$f, g\in\mathbf{Q}[A, B]$，且 $f$ 模 $g$ 不可约。
% P01-E0307: unsupported maximal-ideal classification recorded in ERRATA_CANDIDATES.jsonl.
但 $R$ 不是域，所以核必为 $(A^3-B^2 + AB)$；因此 $\varphi$ 给出同构
$R \to \mathbf{Q}[A, B]/(A^3-B^2 + AB)$。
% P01-E0308: induced-isomorphism direction recorded in ERRATA_CANDIDATES.jsonl.

\medskip\noindent
为看出 $\varphi$ 满射，必须把任意 $f\in R$ 表成 $A(z) = z^2-z$
与 $B(z) = z^3-z^2$ 的 $\mathbf{Q}$-系数多项式。注意关系
$zA(z) = B(z)$。令 $a = f(0) = f(1)$。于是 $z(z-1)$ 必须整除
$f(z)-a$，所以可写成
$f(z) = z(z-1)g(z)+a = A(z)g(z)+a$。若 $\deg(g) < 2$，则
$h(z) = c_1z + c_0$，且
$f(z) = A(z)(c_1z + c_0)+a = c_1B(z)+c_0A(z)+a$，故结论成立。
% P01-E0309: h/g variable mismatch recorded in ERRATA_CANDIDATES.jsonl.
若 $\deg(g)\geq 2$，则由多项式除法算法，可写成
$g(z) = A(z)h(z)+b_1z + b_0$（$\deg(h)\leq\deg(g)-2$），所以
$f(z) = A(z)^2h(z)+b_1B(z)+b_0A(z)$。
% P01-E0310: omitted additive constant recorded in ERRATA_CANDIDATES.jsonl.
继续对 $h(z)$ 应用除法并迭代，可将 $f(z)$ 表成 $A(z)$ 与 $B(z)$
的多项式；故 $\varphi$ 满射。

\medskip\noindent
现令 $a \in \mathbf{Q}$、$a \neq 0, \frac{1}{2}, 1$，并考虑
$$
R_a = \{ f \in \mathbf{Q}[z, \frac{1}{z-a}]\text{ 其中 }f(0) = f(1)
\}.
$$
这同样是有限生成 $\mathbf{Q}$-代数：容易验证函数
$z^2-z$、$z^3-z$ 与 $\frac{a^2-a}{z-a}+z$ 作为
$\mathbf{Q}$-代数生成 $R_a$。有下列包含关系：
$$
R\subset R_a\subset\mathbf{Q}[z, \frac{1}{z-a}], \quad
R\subset\mathbf{Q}[z]\subset\mathbf{Q}[z, \frac{1}{z-a}].
$$
回顾（引理 \ref{algebra-lemma-spec-localization}），对环 T 与乘法子集
$S\subset T$，环映射 $T \to S^{-1}T$ 诱导谱上的映射
$\Spec(S^{-1}T) \to \Spec(T)$，它同胚到子集
$$
\{\mathfrak p \in \Spec(T) \mid S \cap \mathfrak p = \emptyset\}
\subset \Spec(T).
$$
若对某个 $f\in T$ 有 $S = \{ 1, f, f^2, \ldots\}$，则这是开集
$D(f)\subset T$。
% P01-E0311: standard-open ambient-space type error recorded in ERRATA_CANDIDATES.jsonl.
现验证环映射 $R \to R_a$ 的相应性质：我们将通过验证 $\theta$ 是单射
局部同胚，证明由包含 $R\subset R_a$ 诱导的映射
$\theta : \Spec(R_a) \to \Spec(R)$ 同胚到 $\Spec(R)$ 的一个开子集。
为此，使用下面将描述的两个基本开集所组成的 $\Spec(R_a)$ 开覆盖。
对任意 $r\in\mathbf{Q}$，令 $\text{ev}_r : R \to \mathbf{Q}$ 为在
$r$ 处求值给出的同态。注意，当 $r = 0$ 及 $r = 1-a$ 时，它可以
延拓为同态 $\text{ev}_r' : R_a \to \mathbf{Q}$（后者是因为
$\frac{1}{z-a}$ 在 $z = 1-a$ 处有定义，这是由于
$a\neq\frac{1}{2}$）。然而，$\text{ev}_a$ 不能延拓到 $R_a$。
写 $\mathfrak{m}_r = \Ker(\text{ev}_r)$。有
$$
\mathfrak{m}_0 = (z^2-z, z^3-z),
$$
$$
\mathfrak{m}_a = ((z-1 + a)(z-a), (z^2-1 + a)(z-a)), \text{ 且}
$$
$$
\mathfrak{m}_{1-a} = ((z-1 + a)(z-a), (z-1 + a)(z^2-a)).
$$
为验证这些等式，注意各右端显然包含于左端。然后把生成元用 $A$ 与
$B$ 表示，并将 $R$ 视为 $\mathbf{Q}[A, B]/(A^3-B^2 + AB)$，由此
检查各右端为极大理想。注意，$\mathfrak{m}_a$ 不在 $\theta$ 的像中：
$$
(z^2 - z)^2(z - a)\left(\frac{a^2 - a}{z - a} + z\right) =
(z^2 - z)^2(a^2 - a) + (z^2 - z)^2(z - a)z
$$
左端属于 $\mathfrak m_a R_a$，因为 $(z^2 - z)(z - a)$ 属于
$\mathfrak m_a$，并且
$(z^2 - z)(\frac{a^2 - a}{z - a} + z)$ 属于 $R_a$。类似地，元素
$(z^2 - z)^2(z - a)z$ 属于 $\mathfrak m_a R_a$，因为
$(z^2 - z)$ 属于 $R_a$，而 $(z^2 - z)(z - a)$ 属于
$\mathfrak m_a$。由于 $a \not \in \{0, 1\}$，可得
$(z^2 - z)^2 \in \mathfrak m_a R_a$。因此 $R_a$ 的任何理想 $I$
都不可能满足 $I \cap R = \mathfrak m_a$，因为这样的 $I$ 必须包含
$(z^2 - z)^2$，而它属于 $R$ 却不属于 $\mathfrak m_a$。基本开集
$D((z-1 + a)(z-a))\subset\Spec(R)$ 等于闭集
$\{\mathfrak{m}_a, \mathfrak{m}_{1-a}\}$ 的补集。然后验证
$R_{(z-1 + a)(z-a)} = (R_a)_{(z-1 + a)(z-a)}$；把这个局部化环记作
$R'$，便可知映射 $R \to R'$ 分解为 $R \to R_a \to R'$。于是由
引理 \ref{algebra-lemma-spec-localization}，这些映射把
$\Spec(R') \subset \Spec(R_a)$ 与 $\Spec(R') \subset \Spec(R)$
表示为开子集；所以 $\theta : \Spec(R_a) \to \Spec(R)$ 限制到
$D((z-1 + a)(z-a))$ 时，同胚到一个开子集。类似地，$\theta$ 限制到
$D((z^2 + z + 2a-2)(z-a)) \subset \Spec(R_a)$ 时，同胚到开子集
$D((z^2 + z + 2a-2)(z-a)) \subset \Spec(R)$。依
$z^2 + z + 2a-2$ 在 $\mathbf{Q}$ 上是否不可约，前一个基本开集的
补集由一或两个闭点以及闭点 $\mathfrak{m}_a$ 组成。此外，在 $R_a$
中由元素 $(z^2 + z + 2a-a)(z-a)$ 与 $(z-1 + a)(z-a)$ 生成的理想
是整个 $R_a$，所以这两个基本开集覆盖 $\Spec(R_a)$。
% P01-E0312: inconsistent coefficient in covering generator recorded in ERRATA_CANDIDATES.jsonl.
因此，为证明 $\theta$ 同胚到
$\Spec(R)-\{\mathfrak{m}_a\}$，只需证明这一或两个点绝不可能等于
$\mathfrak{m}_{1-a}$。事实确实如此，因为 $1-a$ 是
$z^2 + z + 2a-2$ 的根，当且仅当 $a = 0$ 或 $a = 1$；二者均不出现。

\medskip\noindent
尽管这个同胚仿效了在 $R$ 的某个元素处局部化的行为，而且
$\mathbf{Q}[z, \frac{1}{z-a}]$ 是 $\mathbf{Q}[z]$ 在极大理想
$(z-a)$ 处的局部化，环 $R_a$ 却{\it 不是} $R$ 的局部化：任何局部化
$S^{-1}R$ 都会产生比原环 $R$ 更多的单位。
% P01-E0313: localization-at-maximal-ideal error recorded in ERRATA_CANDIDATES.jsonl.
% P01-E0314: overbroad units claim recorded in ERRATA_CANDIDATES.jsonl.
$R$ 的单位是 $\mathbf{Q}^\times$，也就是 $\mathbf{Q}$ 的单位。事实上，
容易看出 $R_a$ 的单位为 $\mathbf{Q}^*$。具体而言，
$\mathbf{Q}[z, \frac{1}{z - a}]$ 的单位形如 $c (z - a)^n$，其中
$c \in \mathbf{Q}^*$ 且 $n \in \mathbf{Z}$；显然，只有当 $n = 0$
时，它们才属于 $R_a$。所以 $R_a$ 的单位并不比 $R$ 多，因而不可能是
$R$ 的局部化。

\medskip\noindent
我们使用 $a\neq 0, 1$ 来保证 $\frac{1}{z-a}$ 在 $z = 0, 1$ 处
有意义。还在若干处使用了 $a\neq 1/2$：(1) 为了能够讨论
$R_a$ 上 $\text{ev}_{1-a}$ 的核，这保证 $\mathfrak{m}_{1-a}$ 是
$R_a$ 的一个点（即 $R_a$ 恰缺少 $R$ 的一个点）。(2) 在最后，为了
推出 $(z-a)^{k + \ell}$ 只有在 $k = \ell = 0$ 时才能属于 $R$；
% P01-E0315: unexplained exponent variables recorded in ERRATA_CANDIDATES.jsonl.
事实上，若 $a = 1/2$，则只要 $k + \ell$ 为偶数，它就属于 $R$。
所以 $R_a$ 的单位确实会比 $R$ 更多，并且 $R_a$ 可能是 $R$ 的局部化。
\end{example}

\section{关于素理想的一则元观察}
\label{algebra-section-oka-families}

\noindent
本节取自 CRing 项目。设 $R$ 为环，$S \subset R$ 为乘法子集。
引理 \ref{algebra-lemma-spec-localization} 的一个推论是：关于“不与 $S$ 相交”
这一性质极大的理想 $I \subset R$ 是素理想。理由是，对局部化环
$S^{-1}R$ 的某个极大理想 $\mathfrak m$，有 $I = R \cap \mathfrak m$。
事实表明，对于理想的许多性质，不具有该性质的极大理想都是素理想。
\cite{Lam-Reyes} 发展了一种看出这一点的一般方法。本节暂离主线，
解释这一现象。

\medskip\noindent
设 $R$ 为环。若 $I$ 是 $R$ 的理想且 $a \in R$，定义
$$
(I : a) = \left\{ x \in R \mid xa \in I\right\}.
$$
更一般地，若 $J \subset R$ 是理想，定义
$$
(I : J) = \left\{ x \in R \mid xJ \subset I\right\}.
$$

\begin{lemma}
\label{algebra-lemma-colon}
设 $R$ 为环。若 $J \subset R$ 为主理想，则对任意理想
$I \subset J$，都有 $I = J (I : J)$。
\end{lemma}

\begin{proof}
设 $J = (a)$。则 $(I : J) = (I : a)$。由于 $I \subset J$，
任意 $y \in I$ 都可写成 $y = xa$，其中 $x \in (I : a)$。
故 $I \subset J (I : J)$。反之，若 $x \in (I : a)$，则
$xJ = (xa) \subset I$，这证明了另一包含关系。
\end{proof}

\noindent
设 $\mathcal{F}$ 为 $R$ 的一些理想组成的集合。我们关心怎样的条件
能够保证 $\mathcal{F}$ 的补集中的极大元为素理想。

\begin{definition}
\label{algebra-definition-oka-family}
设 $R$ 为环，$\mathcal{F}$ 为 $R$ 的理想组成的集合。若
$R \in \mathcal{F}$，并且只要 $I \subset R$ 是理想，且对某个
$a \in R$ 有 $(I : a), (I, a) \in \mathcal{F}$，便有
$I \in \mathcal{F}$，则称 $\mathcal{F}$ 为一个 {\it Oka 族}。
\end{definition}

\noindent
下面给出若干 Oka 族的例子。第一个就是本节引言中讨论的基本例子。

\begin{example}
\label{algebra-example-oka-family-not-meet-multiplicative-set}
设 $R$ 为环，$S$ 为 $R$ 的乘法子集。我们断言
$\mathcal{F} = \{I \subset R \mid I \cap S \not = \emptyset\}$
是 Oka 族。事实上，设对某个 $a \in R$ 有
$(I : a), (I, a) \in \mathcal{F}$。取
$s \in (I, a) \cap S$ 与 $s' \in (I : a) \cap S$。则
$ss' \in I \cap S$，因而 $I \in \mathcal{F}$。所以
$\mathcal{F}$ 是 Oka 族。
\end{example}

\begin{example}
\label{algebra-example-oka-family-finitely-generated}
设 $R$ 为环，$I \subset R$ 为理想，$a \in R$。若
$(I : a)$ 由 $a_1, \ldots, a_n$ 生成，而 $(I, a)$ 由
$a, b_1, \ldots, b_m$ 生成，其中 $b_1, \ldots, b_m \in I$，
则 $I$ 由
$aa_1, \ldots, aa_n, b_1, \ldots, b_m$ 生成。为说明这一点，
注意若 $x \in I$，则 $x \in (I, a)$ 是
$a, b_1, \ldots, b_m$ 的线性组合，但 $a$ 的系数必属于
$(I:a)$。由此可知，有限生成理想组成的族是 Oka 族。
\end{example}

\begin{example}
\label{algebra-example-oka-family-principal}
我们证明环 $R$ 的主理想组成的族是 Oka 族。事实上，设
$I \subset R$ 为理想，$a \in R$，并且 $(I, a)$ 与 $(I : a)$
都是主理想。注意 $(I : a) = (I : (I, a))$。令 $J = (I, a)$，
则 $J$ 与 $(I : J)$ 都是主理想。由引理 \ref{algebra-lemma-colon}，
有 $I = J (I : J)$。因此在这里，由于 $J = (I, a)$ 与
$(I : J)$ 都是主理想，$I$ 也是主理想。
\end{example}

\begin{example}
\label{algebra-example-oka-family-bound-cardinality}
设 $R$ 为环，$\kappa$ 为无限基数。至多可由 $\kappa$ 个元素生成的
理想组成的族是 Oka 族。论证与
例 \ref{algebra-example-oka-family-finitely-generated} 中的论证类似，故略去。
\end{example}

\begin{example}
\label{algebra-example-oka-family-property-modules}
设 $A$ 为环，$I \subset A$ 为理想，$a \in A$ 为元素。有短正合列
$0 \to A/(I : a) \to A/I \to A/(I, a) \to 0$，其中第一个箭头由
乘以 $a$ 给出。因此，若 $P$ 是一个对扩张稳定且对 $0$ 成立的
$A$-模性质，则使 $A/I$ 具有 $P$ 的那些理想 $I$ 组成 Oka 族。
\end{example}

\begin{proposition}
\label{algebra-proposition-oka}
若 $\mathcal{F}$ 是理想的 Oka 族，则 $\mathcal{F}$ 的补集中的
任意极大元都是素理想。
\end{proposition}

\begin{proof}
设 $I \not \in \mathcal{F}$ 关于不属于 $\mathcal{F}$ 极大，
但 $I$ 不是素理想。注意 $I \not = R$，因为
$R \in \mathcal{F}$。由于 $I$ 不素，可找到 $a, b \in R - I$，
使 $ab \in I$。于是 $(I, a) \neq I$，且 $(I : a)$ 包含
$b \not \in I$，故也有 $(I : a) \neq I$。所以
$(I : a), (I, a)$ 都严格包含 $I$，从而必属于 $\mathcal{F}$。
由 Oka 条件，$I \in \mathcal{F}$，矛盾。
\end{proof}

\noindent
至此，我们可以把上面大多数例子转化为关于环中素理想的引理。

\begin{lemma}
\label{algebra-lemma-simple}
设 $R$ 为环，$S$ 为 $R$ 的乘法子集。关于性质
$I \cap S = \emptyset$ 极大的理想 $I \subset R$ 是素理想。
\end{lemma}

\begin{proof}
这就是本节引言中讨论的例子。另一种证明是把
例 \ref{algebra-example-oka-family-not-meet-multiplicative-set}
与命题 \ref{algebra-proposition-oka} 结合起来。
\end{proof}

\begin{lemma}
\label{algebra-lemma-cohen}
设 $R$ 为环。
\begin{enumerate}
\item 关于“不有限生成”极大的理想 $I \subset R$ 是素理想。
\item 若 $R$ 的每个素理想都有限生成，则 $R$ 的每个理想都有限生成
\footnote{稍后我们将称 $R$ 为诺特环。}。
\end{enumerate}
\end{lemma}

\begin{proof}
第一条是例 \ref{algebra-example-oka-family-finitely-generated} 与
命题 \ref{algebra-proposition-oka} 的直接推论。对于第二条，设存在一个
不有限生成的理想 $I \subset R$。由不有限生成理想组成的全序链
$\left\{I_\alpha\right\}$ 的并仍不有限生成；事实上，若
$I = \bigcup I_\alpha$ 由 $a_1, \ldots, a_n$ 生成，则所有生成元
都属于某个 $I_\alpha $，因而会生成它。由佐恩引理，存在一个关于
不有限生成极大的理想。由第一条，该理想是素理想。
\end{proof}

\begin{lemma}
\label{algebra-lemma-primes-principal}
设 $R$ 为环。
\begin{enumerate}
\item 关于“不是主理想”极大的理想 $I \subset R$ 是素理想。
\item 若 $R$ 的每个素理想都是主理想，则 $R$ 的每个理想都是主理想。
\end{enumerate}
\end{lemma}

\begin{proof}
第一条由例 \ref{algebra-example-oka-family-principal} 与
命题 \ref{algebra-proposition-oka} 得出。对于第二条，设存在一个不是主理想的
理想 $I \subset R$。由不是主理想的理想组成的全序链
$\left\{I_\alpha\right\}$ 的并不是主理想；
% P01-E0316: missing copula in frozen source recorded in ERRATA_CANDIDATES.jsonl.
事实上，若 $I = \bigcup I_\alpha$ 由 $a$ 生成，则 $a$ 属于某个
$I_\alpha $，并且 $a$ 会生成该理想。由佐恩引理，存在一个关于不是主理想
极大的理想。由第一条，该理想必为素理想。
\end{proof}

\begin{lemma}
\label{algebra-lemma-characterize-domain}
设 $R$ 为环。
\begin{enumerate}
\item 在不含非零因子的理想中极大的理想是素理想。
\item 若 $R$ 非零，且 $R$ 中每个非零素理想都包含一个非零因子，
则 $R$ 是整环。
\end{enumerate}
\end{lemma}

\begin{proof}
考虑由非零因子组成的集合 $S$。它是 $R$ 的乘法子集。因此，
关于不与 $S$ 相交极大的任何理想都是素理想，见
引理 \ref{algebra-lemma-simple}。所以，若每个非零素理想都包含一个非零因子，
则 $(0)$ 是素理想，即 $R$ 是整环。
\end{proof}

\begin{remark}
\label{algebra-remark-cohen-bound-cardinality}
设 $R$ 为环，$\kappa$ 为无限基数。应用
例 \ref{algebra-example-oka-family-bound-cardinality} 与
命题 \ref{algebra-proposition-oka} 可知，关于不能由 $\kappa$ 个元素生成
这一性质极大的任何理想都是素理想。这个结果不太有用，因为存在一个环
$R$，它的每个素理想都可由 $\aleph_0$ 个元素生成，但某个理想却不可以。
具体地，设 $k$ 为域，$T$ 为基数大于 $\aleph_0$ 的集合，并令
$$
R = k[\{x_n\}_{n \geq 1}, \{z_{t, n}\}_{t \in T, n \geq 0}]/
(x_n^2, z_{t, n}^2, x_n z_{t, n} - z_{t, n - 1})
$$
这是一个局部环，其唯一素理想为 $\mathfrak m = (x_n)$。但是理想
$(z_{t, n})$ 不能由可数多个元素生成。
\end{remark}

\begin{example}
\label{algebra-example-noetherian-topology-on-spec}
\begin{reference}
Lukas Heger 于 2020 年 11 月 12 日的评论。
\end{reference}
设 $R$ 为环，$X = \Spec(R)$。由于 $X$ 的闭子集对应于 $R$ 的根理想
（引理 \ref{algebra-lemma-Zariski-topology}），故 $X$ 是诺特拓扑空间当且仅当
根理想满足升链条件。
% P01-E0317: “radical ideas” source typo recorded in ERRATA_CANDIDATES.jsonl.
这又当且仅当每个根理想都是一个有限生成理想的根（细节略去）。令
$$
\mathcal{F} = \{I \subset R \mid
\sqrt{I} = \sqrt{(f_1, \ldots, f_n)}\text{ 对某个 }n
\text{ 且 }f_1, \ldots, f_n \in R\}.
$$
读者可利用恒等式
$$
\sqrt{I} = \sqrt{(I, a)(I : a)}
$$
证明 $\mathcal{F}$ 是 Oka 族；该恒等式对任意理想 $I \subset R$
及任意元素 $a \in R$ 成立。另一方面，若有一条理想全序链
$\{I_\alpha\}$，其中没有理想属于 $\mathcal{F}$，则其并
$I = \bigcup I_\alpha$ 也不可能属于 $\mathcal{F}$。否则，若
$\sqrt{I} = \sqrt{(f_1, \ldots, f_n)}$，则对某个 $e$ 有
$f_i^e \in I$；于是对某个不依赖于 $i$ 的 $\alpha$，有
$f_i^e \in I_\alpha$；继而
$\sqrt{I_\alpha} = \sqrt{(f_1, \ldots, f_n)}$，矛盾。因此，
若不属于 $\mathcal{F}$ 的理想集合非空，它便有极大元；完全仿照
引理 \ref{algebra-lemma-cohen} 可知，$X$ 是诺特拓扑空间当且仅当 $R$ 的
每个素理想都等于某个 $\sqrt{(f_1, \ldots, f_n)}$，其中
$f_1, \ldots, f_n \in R$。若以后需要这个结果，我们会在此处
严谨地陈述并证明它。
\end{example}

\section{有限表示环映射的像}
\label{algebra-section-images-finite-presentation}

\noindent
本节证明由环映射 $R \to S$ 诱导的映射
$\Spec(S) \to \Spec(R)$ 的拓扑性质，主要是切瓦莱定理。为此，
我们将使用可构造集、拟紧集、逆紧集等概念；这些概念定义于
拓扑学，第 \ref{topology-section-constructible} 节。

\begin{lemma}
\label{algebra-lemma-qc-open}
设 $U \subset \Spec(R)$ 为开集。以下条件等价：
\begin{enumerate}
\item $U$ 在 $\Spec(R)$ 中逆紧；
\item $U$ 拟紧；
\item $U$ 是有限个标准开集之并；
\item 存在有限生成理想 $I \subset R$，使得
$X \setminus V(I) = U$。
% P01-E0318: undefined X in frozen source recorded in ERRATA_CANDIDATES.jsonl.
\end{enumerate}
\end{lemma}

\begin{proof}
因为 $\Spec(R)$ 拟紧，所以 (1) $\Rightarrow$ (2)，见
引理 \ref{algebra-lemma-quasi-compact}。因为标准开集构成拓扑的一组基，
所以 (2) $\Rightarrow$ (3)。证明 (3) $\Rightarrow$ (1)。令
$U = \bigcup_{i = 1\ldots n} D(f_i)$。要证明 $U$ 在 $\Spec(R)$ 中
逆紧，只须证明对 $\Spec(R)$ 的任意拟紧开集 $V$，$U \cap V$ 拟紧。
由 (2) $\Rightarrow$ (3)，可写
$V = \bigcup_{j = 1\ldots m} D(g_j)$。每个标准开集都同胚于某个环的谱，
故为拟紧的，见引理 \ref{algebra-lemma-standard-open} 与
\ref{algebra-lemma-quasi-compact}。因此
$U \cap V =
(\bigcup_{i = 1\ldots n} D(f_i)) \cap (\bigcup_{j = 1\ldots m} D(g_j))
= \bigcup_{i, j} D(f_i g_j)$ 是有限个拟紧开集之并，因而拟紧。
最后，由引理 \ref{algebra-lemma-Zariski-topology}，(4) 与 (3) 等价。
\end{proof}

\begin{lemma}
\label{algebra-lemma-affine-map-quasi-compact}
设 $\varphi : R \to S$ 为环映射。诱导的连续映射
$f : \Spec(S) \to \Spec(R)$ 是拟紧的。对任意可构造集
$E \subset \Spec(R)$，逆像 $f^{-1}(E)$ 在 $\Spec(S)$ 中可构造。
\end{lemma}

\begin{proof}
先证明任意拟紧开集 $U \subset \Spec(R)$ 的逆像拟紧。由
引理 \ref{algebra-lemma-qc-open}，可将 $U$ 写成有限个标准开集之并。
% P01-E0319: “finite open” source typo recorded in ERRATA_CANDIDATES.jsonl.
于是由引理 \ref{algebra-lemma-spec-functorial}，$f^{-1}(U)$ 是有限个标准开集
之并。再次应用引理 \ref{algebra-lemma-qc-open}，可知 $f^{-1}(U)$ 拟紧。
第二个断言由拓扑学，引理
\ref{topology-lemma-inverse-images-constructibles} 得出。
\end{proof}

\begin{lemma}
\label{algebra-lemma-constructible}
设 $R$ 为环。$\Spec(R)$ 的一个子集可构造，当且仅当它可写成有限个
形如 $D(f) \cap V(g_1, \ldots, g_m)$ 的子集之并，其中
$f, g_1, \ldots, g_m \in R$。
\end{lemma}

\begin{proof}
由引理 \ref{algebra-lemma-qc-open}，子集 $D(f)$ 与
$V(g_1, \ldots, g_m)$ 的补集都是逆紧开集。因此
$D(f) \cap V(g_1, \ldots, g_m)$ 是可构造子集，其任意有限并亦然。
反之，设 $T \subset \Spec(R)$ 可构造。由拓扑学，定义
\ref{topology-definition-constructible}，可设
$T = U \cap V^c$，其中 $U, V \subset \Spec(R)$ 是逆紧开集。
由引理 \ref{algebra-lemma-qc-open}，可写
$U = \bigcup_{i = 1, \ldots, n} D(f_i)$ 及
$V = \bigcup_{j = 1, \ldots, m} D(g_j)$。于是
$T = \bigcup_{i = 1, \ldots, n} \big(D(f_i) \cap V(g_1, \ldots, g_m)\big)$。
\end{proof}

\begin{lemma}
\label{algebra-lemma-constructible-is-image}
设 $R$ 为环，$T \subset \Spec(R)$ 可构造。则存在有限表示环映射
$R \to S$，使 $T$ 是 $\Spec(S)$ 在 $\Spec(R)$ 中的像。
\end{lemma}

\begin{proof}
有限个环之积的谱是这些谱的不交并，见
引理 \ref{algebra-lemma-spec-product}。因此，若 $T = T_1 \cup T_2$，且结论对
$T_1$ 与 $T_2$ 成立，则结论对 $T$ 也成立。由
引理 \ref{algebra-lemma-constructible}，可设
$T = D(f) \cap V(g_1, \ldots, g_m)$。在这种情形下，$T$ 是映射
$\Spec((R/(g_1, \ldots, g_m))_f) \to \Spec(R)$ 的像，见引理
\ref{algebra-lemma-standard-open} 与 \ref{algebra-lemma-spec-closed}。
\end{proof}

\begin{lemma}
\label{algebra-lemma-open-fp}
设 $R$ 为环，$f$ 为 $R$ 的元素，并令 $S = R_f$。则
$\Spec(S)$ 的可构造子集的像在 $\Spec(R)$ 中可构造。
\end{lemma}

\begin{proof}
以下不再说明地反复使用引理 \ref{algebra-lemma-qc-open}。设 $U, V$ 是
$\Spec(S)$ 中的拟紧开集。我们证明 $U \cap V^c$ 的像可构造。
在引理 \ref{algebra-lemma-standard-open} 给出的识别
$\Spec(S) = D(f)$ 下，集合 $U, V$ 对应于 $\Spec(R)$ 的拟紧开集
$U', V'$。因而只须证明 $U' \cap (V')^c$ 在 $\Spec(R)$ 中可构造，
这是显然的。
\end{proof}

\begin{lemma}
\label{algebra-lemma-closed-fp}
设 $R$ 为环，$I$ 为 $R$ 的有限生成理想，并令 $S = R/I$。
则 $\Spec(S)$ 的可构造子集的像在 $\Spec(R)$ 中可构造。
\end{lemma}

\begin{proof}
若 $I = (f_1, \ldots, f_m)$，则 $V(I)$ 是
$\bigcup D(f_i)$ 的补集，见引理 \ref{algebra-lemma-Zariski-topology}。
因此，由引理 \ref{algebra-lemma-qc-open}，它是可构造的。将映射
$R \to S$ 记作 $f \mapsto \overline{f}$。我们须证明：若
$\overline{U}, \overline{V}$ 是 $\Spec(S)$ 的逆紧开集，则
$\overline{U} \cap \overline{V}^c$ 在 $\Spec(R)$ 中的像可构造。
由引理 \ref{algebra-lemma-qc-open}，可写
$\overline{U} = \bigcup D(\overline{g_i})$。令
$U = \bigcup D({g_i})$，则 $\overline{U}$ 的像是
$U \cap V(I)$，它在 $\Spec(R)$ 中可构造。类似地，对
$\Spec(R)$ 的某个逆紧开集 $V$，$\overline{V}$ 的像等于
$V \cap V(I)$。故 $\overline{U} \cap \overline{V}^c$ 的像等于
$U \cap V(I) \cap V^c$，如所需。
\end{proof}

\begin{lemma}
\label{algebra-lemma-affineline-open}
设 $R$ 为环。映射 $\Spec(R[x]) \to \Spec(R)$ 是开映射，且任意
标准开集的像都是拟紧开集。
\end{lemma}

\begin{proof}
只须证明标准开集 $D(f)$（其中 $f\in R[x]$）的像是拟紧开集。
$D(f)$ 的像就是 $\Spec(R[x]_f) \to \Spec(R)$ 的像。设
$\mathfrak p \subset R$ 为素理想，并令 $\overline{f}$ 为 $f$ 在
$\kappa(\mathfrak p)[x]$ 中的像。回忆引理 \ref{algebra-lemma-in-image}：
$\mathfrak p$ 属于该像，当且仅当
$R[x]_f \otimes_R \kappa(\mathfrak p) =
\kappa(\mathfrak p)[x]_{\overline{f}}$ 不是零环。这恰好等价于 $f$
在 $\kappa(\mathfrak p)[x]$ 中的像非零；换言之，$f$ 的某个系数
不属于 $\mathfrak p$。因此，若
$f = a_d x^d + \ldots + a_0$，则 $D(f)$ 的像为
$D(a_d) \cup \ldots \cup D(a_0)$。
\end{proof}

\noindent
下面证明稍后将使用的特征多项式性质。

\begin{lemma}
\label{algebra-lemma-characteristic-polynomial-prime}
设 $R \to A$ 为环同态，并假设作为 $R$-模有
$A \cong R^{\oplus n}$。设 $f \in A$。乘法映射 $m_f: A
\to A$ 是 $R$-线性的，因而有特征多项式
$P(T) = T^n + r_{n-1}T^{n-1} + \ldots + r_0 \in R[T]$。
对任意素理想 $\mathfrak{p} \in \Spec(R)$，$f$ 在
$A \otimes_R \kappa(\mathfrak{p})$ 上幂零地作用，当且仅当
$\mathfrak p \in V(r_0, \ldots, r_{n-1})$。
\end{lemma}

\begin{proof}
只须证明：乘法映射
$m_{\bar f}: A \otimes_R \kappa(\mathfrak p) \to
A \otimes_R \kappa(\mathfrak p)$ 的特征多项式
$\bar P(T) \in \kappa(\mathfrak p)[T]$，恰为 $P(T)$ 在
$\kappa(\mathfrak p)[T]$ 中的像，结论便容易得出。这里该映射以
$\bar f$ 乘 $A \otimes_R \kappa(\mathfrak p)$ 的元素，而它是视为
$\kappa(\mathfrak p)$ 中元素的 $f$ 的像。
% P01-E0320: incorrect ambient ring for bar f recorded in ERRATA_CANDIDATES.jsonl.
以 $R$-模 $A$ 的一组基 $e_1, \ldots, e_n$ 表示乘法映射 $m_f$，
设其矩阵为 $(a_{ij})$，矩阵元属于 $R$。那么
$A \otimes_R \kappa(\mathfrak p) \cong (R \otimes_R
\kappa(\mathfrak p))^{\oplus n} = \kappa(\mathfrak p)^n$，这是
$\kappa(\mathfrak p)$ 上以
$e_1 \otimes 1, \ldots, e_n \otimes 1$ 为基的 $n$ 维向量空间。
像为 $\bar f = f \otimes 1$，故乘法映射 $m_{\bar f}$ 的矩阵是
$(a_{ij} \otimes 1)$。因此，特征多项式恰为 $P(T)$ 的像。

\medskip\noindent
由线性代数，一个线性变换在 $n$ 维向量空间上幂零地作用，
当且仅当其特征多项式是 $T^n$（因为特征多项式整除极小多项式的
某个幂）。所以，$f$ 在 $A \otimes_R \kappa(\mathfrak p)$ 上
幂零地作用，当且仅当 $\bar P(T) = T^n$。这又当且仅当对所有
$0 \leq i \leq n - 1$ 都有 $r_i \in \mathfrak p$，即
$\mathfrak p \in V(r_0, \ldots, r_{n - 1}).$
\end{proof}

\begin{lemma}
\label{algebra-lemma-affineline-special}
设 $R$ 为环，$f, g \in R[x]$ 为多项式。假设 $g$ 的首项系数是
$R$ 的单位。存在元素 $r_i\in R$（$i = 1\ldots, n$），使
$D(f) \cap V(g)$ 在 $\Spec(R)$ 中的像为
$\bigcup_{i = 1, \ldots, n} D(r_i)$。
% P01-E0321: subject-verb agreement defect recorded in ERRATA_CANDIDATES.jsonl.
\end{lemma}

\begin{proof}
写成 $g = ux^d + a_{d-1}x^{d-1} + \ldots + a_0$，其中 $d$ 是 $g$
的次数，故 $u \in R^*$。考虑环 $A = R[x]/(g)$。作为 $R$-模，
它是有限自由的，以 $1, x, \ldots, x^{d-1}$ 的像为基。考虑 $A$ 上
乘以 $f$ 的像。这是一个 $R$-模映射。因此可令 $P(T) \in R[T]$
为该映射的特征多项式。写成
$P(T) = T^d + r_{d-1} T^{d-1} + \ldots + r_0$。我们断言
$r_0, \ldots, r_{d-1}$ 具有所需性质。下面将使用特征多项式的性质
$$
\mathfrak p \in V(r_0, \ldots, r_{d-1})
\Leftrightarrow
\text{multiplication by }f\text{ is nilpotent on }
A \otimes_R \kappa(\mathfrak p).
$$
该性质已在引理 \ref{algebra-lemma-characteristic-polynomial-prime} 中证明。

\medskip\noindent
设 $\mathfrak q\in D(f) \cap V(g)$，并令
$\mathfrak p = \mathfrak q \cap R$。则有一个与乘以 $f$ 相容的非零映射
$A \otimes_R \kappa(\mathfrak p) \to \kappa(\mathfrak q)$。
而 $f$ 在 $\kappa(\mathfrak q)$ 上作为单位作用。因此可知
$\mathfrak p \not \in  V(r_0, \ldots, r_{d-1})$。

\medskip\noindent
另一方面，设对 $R$ 的某个素理想 $\mathfrak p$ 及某个
$0 \leq i \leq d - 1$，有 $r_i \not\in \mathfrak p$。则乘以 $f$
在代数 $A \otimes_R \kappa(\mathfrak p)$ 上不是幂零的。因此存在
不包含 $f$ 的像的素理想 $\overline{\mathfrak q} \subset
A \otimes_R \kappa(\mathfrak p)$。$\overline{\mathfrak q}$ 在
$R[x]$ 中的逆像是 $D(f) \cap V(g)$ 的一个元素，并映到
$\mathfrak p$。
\end{proof}

\begin{theorem}[切瓦莱定理]
\label{algebra-theorem-chevalley}
设 $R \to S$ 是有限表示的。则 $\Spec(S)$ 的可构造子集在
$\Spec(R)$ 中的像可构造。
\end{theorem}

\begin{proof}
写成 $S = R[x_1, \ldots, x_n]/(f_1, \ldots, f_m)$。可将
$R \to S$ 分解为
$R \to R[x_1] \to R[x_1, x_2]
\to \ldots \to R[x_1, \ldots, x_{n-1}] \to S$。所以可设
$S = R[x]/(f_1, \ldots, f_m)$。此时再将映射分解为
$R \to R[x] \to S$；由引理 \ref{algebra-lemma-closed-fp}，问题归结到
$S = R[x]$ 的情形。由引理 \ref{algebra-lemma-qc-open}，只须证明若
% P01-E0322: missing “it” in frozen source recorded in ERRATA_CANDIDATES.jsonl.
$T = (\bigcup_{i = 1\ldots n} D(f_i)) \cap V(g_1, \ldots, g_m)$，
其中 $f_i , g_j \in R[x]$，则其在 $\Spec(R)$ 中的像可构造。
由于可构造集的有限并可构造，只须处理 $n = 1$，即
$T = D(f) \cap V(g_1, \ldots, g_m)$ 的情形。

\medskip\noindent
注意若 $c \in R$，则
$$
\Spec(R) =
V(c) \amalg D(c) =
\Spec(R/(c)) \amalg \Spec(R_c),
$$
相应地，$\Spec(R[x]) =
V(c) \amalg D(c) = \Spec(R/(c)[x]) \amalg
\Spec(R_c[x])$。$T = D(f) \cap V(g_1, \ldots, g_m)$ 与每一部分的交
仍有相同形式，其中 $f$、$g_i$ 分别由其在 $R/(c)[x]$、$R_c[x]$
中的像代替。注意，$T$ 在 $\Spec(R)$ 中的像是
$T \cap V(c)$ 与 $T \cap D(c)$ 的像之并。利用引理
\ref{algebra-lemma-open-fp} 和 \ref{algebra-lemma-closed-fp}，只须分别证明这两部分
的像在 $\Spec(R/(c))$ 与 $\Spec(R_c)$ 中可构造。

\medskip\noindent
设如上有 $T = D(f) \cap V(g_1, \ldots, g_m)$，且
$\deg(g_1) \leq \deg(g_2) \leq \ldots \leq \deg(g_m)$。
我们将对 $m$ 以及诸 $g_i$ 的次数作归纳。令
$d_1 = \deg(g_1)$，即 $g_1 = c x^{d_1} + l.o.t$，其中非零的
$c \in R$。把 $R$ 分成 $R/(c)$ 与 $R_c$ 两部分后，要么降低
$g_1$ 的次数（由归纳覆盖），要么归结到 $c$ 可逆的情形。若 $c$ 可逆
且 $m > 1$，写成 $g_2 = c' x^{d_2} + l.o.t$。此时考虑
$g_2' = g_2 - (c'/c) x^{d_2 - d_1} g_1$。因为理想
$(g_1, g_2, \ldots, g_m)$ 与 $(g_1, g_2', g_3, \ldots, g_m)$
相等，所以 $T = D(f) \cap V(g_1, g_2', g_3\ldots, g_m)$。
% P01-E0323: missing separator after g_3 recorded in ERRATA_CANDIDATES.jsonl.
但这里 $g_2'$ 的次数严格小于 $g_2$ 的次数，故此情形由归纳覆盖。

\medskip\noindent
上述归纳的基本情形为：(a) $T = D(f) \cap V(g)$，其中 $g$ 的首项系数
可逆；以及 (b) $T = D(f)$。这两种情形分别由引理
\ref{algebra-lemma-affineline-special} 与 \ref{algebra-lemma-affineline-open} 处理。
% P01-E0324: “bases case” source typo recorded in ERRATA_CANDIDATES.jsonl.
\end{proof}

\section{关于像的进一步结果}
\label{algebra-section-more-images}

\noindent
本节汇集几个关于环映射在 $\Spec$ 上之像的补充引理。例如，
另见第 \ref{algebra-section-going-up} 节。

\begin{lemma}
\label{algebra-lemma-generic-finite-presentation}
设 $R \subset S$ 是整环的包含，并假设 $R \to S$ 是有限型的。
则存在非零的 $f \in R$ 与非零的 $g \in S$，使
$R_f \to S_{fg}$ 是有限表示的。
\end{lemma}

\begin{proof}
对 $S$ 在 $R$ 上的生成元个数作归纳。证明过程中，可将 $R$ 换成
$R_f$，将 $S$ 换成 $S_f$，其中 $f \in R$ 非零。

\medskip\noindent
设 $S$ 在 $R$ 上由一个元素生成。则对某个素理想
$\mathfrak q \subset R[x]$，有 $S = R[x]/\mathfrak q$。若
$\mathfrak q = (0)$，无须证明。若 $\mathfrak q \not = (0)$，
取 $h \in \mathfrak q$ 为关于 $x$ 次数最小的非零元素。写成
$h = f x^d + a_{d - 1} x^{d - 1} + \ldots + a_0$，其中
$a_i \in R$ 且 $f \not = 0$。在 $R$ 与 $S$ 中将 $f$ 可逆化后，
可设 $h$ 首一。于是得到满的 $R$-代数映射
$R[x]/(h) \to S$。作为 $R$-模，有
$R[x]/(h) = R \oplus Rx \oplus \ldots \oplus Rx^{d - 1}$；由 $d$
的极小性，$R[x]/(h)$ 到 $S$ 的映射是单射。因此
$R[x]/(h) \cong S$ 在 $R$ 上有限表示。

\medskip\noindent
设 $S$ 在 $R$ 上由 $n > 1$ 个元素生成。设
$x_1, \ldots, x_n \in S$ 生成 $S$，并记 $S' \subset S$ 为由
$x_1, \ldots, x_{n-1}$ 生成的 $R$-子代数。由归纳假设，存在非零的
$f\in R$ 与 $g \in S'$，使 $R_f \to S'_{fg}$ 有限表示。
% P01-E0325: missing article before “induction hypothesis” recorded in ERRATA_CANDIDATES.jsonl.
接着对 $S'_{fg} \to S_{fg}$ 应用归纳假设，可知存在
$f' \in S'_{fg}$ 与 $g' \in S_{fg}$，使
$S'_{fgf'} \to S_{fgf'g'}$ 有限表示。余下结论留给读者。
\end{proof}

\begin{lemma}
\label{algebra-lemma-characterize-image-finite-type}
设 $R \to S$ 为有限型环映射。记 $X = \Spec(R)$、
$Y = \Spec(S)$，并将诱导的谱映射写成 $f : Y \to X$。
% P01-E0326: malformed “Write ... the induced map” recorded in ERRATA_CANDIDATES.jsonl.
设 $E \subset Y = \Spec(S)$ 为可构造集。若点 $\xi \in X$ 属于
$f(E)$，则 $\overline{\{\xi\}} \cap f(E)$ 包含
$\overline{\{\xi\}}$ 的一个稠密开子集。
\end{lemma}

\begin{proof}
设 $\xi \in X$ 是 $f(E)$ 的一点。选取映到 $\xi$ 的点
$\eta \in E$。令 $\mathfrak p \subset R$ 为对应于 $\xi$ 的素理想，
$\mathfrak q \subset S$ 为对应于 $\eta$ 的素理想。考虑图
$$
\xymatrix{
\eta \ar[r] \ar@{|->}[d] & E \cap Y' \ar[r] \ar[d] &
Y' = \Spec(S/\mathfrak q) \ar[r] \ar[d] &
Y \ar[d] \\
\xi \ar[r] & f(E) \cap X' \ar[r] &
X' = \Spec(R/\mathfrak p) \ar[r] &
X
}
$$
由引理 \ref{algebra-lemma-affine-map-quasi-compact}，$E \cap Y'$ 在 $Y'$ 中
可构造。于是可以用 $X'$ 代替 $X$，用 $Y'$ 代替 $Y$。因此可设
$R \subset S$ 是整环的包含，$\xi$ 是 $X$ 的泛点，而 $\eta$ 是
$Y$ 的泛点。结合引理 \ref{algebra-lemma-generic-finite-presentation} 与
切瓦莱定理（定理 \ref{algebra-theorem-chevalley}），可知存在稠密开集
$U \subset X$、$V \subset Y$，使 $f(V) \subset U$，并且
$f : V \to U$ 将可构造集映为可构造集。注意 $E \cap V$ 在 $V$
中可构造，见拓扑学，引理
\ref{topology-lemma-open-immersion-constructible-inverse-image}。
所以 $f(E \cap V)$ 在 $U$ 中可构造并包含 $\xi$。由拓扑学，引理
\ref{topology-lemma-generic-point-in-constructible}，
$f(E \cap V)$ 包含一个稠密开集 $U' \subset U$。
\end{proof}

\noindent
本节最后给出几个与可构造集无关的谱映射之像的结果。

\begin{lemma}
\label{algebra-lemma-surjective-spec-radical-ideal}
设 $\varphi : R \to S$ 为环映射。以下条件等价：
\begin{enumerate}
\item 映射 $\Spec(S) \to \Spec(R)$ 是满射。
\item 对任意理想 $I \subset R$，$\sqrt{IS}$ 在 $R$ 中的逆像等于
$\sqrt{I}$。
\item 对任意根理想 $I \subset R$，$IS$ 在 $R$ 中的逆像等于 $I$。
\item 对 $R$ 的每个素理想 $\mathfrak p$，$\mathfrak p S$ 在 $R$
中的逆像是 $\mathfrak p$。
\end{enumerate}
在这种情形下，任意基变换后仍有同样的结论：给定环映射
$R \to R'$，环映射 $R' \to R' \otimes_R S$ 也具有等价性质
(1)、(2)、(3)。
\end{lemma}

\begin{proof}
若 $J \subset S$ 为理想，则
$\sqrt{\varphi^{-1}(J)} = \varphi^{-1}(\sqrt{J})$。这说明 (2) 与
(3) 等价。(3) $\Rightarrow$ (4) 是直接的。若 $I \subset R$ 是根理想，
则引理 \ref{algebra-lemma-Zariski-topology} 保证
$I = \bigcap_{I \subset \mathfrak p} \mathfrak p$。因此
(4) $\Rightarrow$ (2)。由引理 \ref{algebra-lemma-in-image}，有
$\mathfrak p = \varphi^{-1}(\mathfrak p S)$ 当且仅当
$\mathfrak p$ 属于该像。所以 (1) $\Leftrightarrow$ (4)。因此
(1)、(2)、(3)、(4) 等价。

\medskip\noindent
假设 (1) 成立。设 $R \to R'$ 为环映射，并设
$\mathfrak p' \subset R'$ 是位于 $R$ 的素理想 $\mathfrak p$ 之上的
素理想。要证明 $\mathfrak p'$ 属于
$\Spec(R' \otimes_R S) \to \Spec(R')$ 的像，由引理
\ref{algebra-lemma-in-image}，须证明
$(R' \otimes_R S) \otimes_{R'} \kappa(\mathfrak p')$ 非零。但有
$$
(R' \otimes_R S) \otimes_{R'} \kappa(\mathfrak p') =
S \otimes_R \kappa(\mathfrak p)
\otimes_{\kappa(\mathfrak p)} \kappa(\mathfrak p')
$$
它非零，因为由假设 $S \otimes_R \kappa(\mathfrak p)$ 非零，且
$\kappa(\mathfrak p) \to \kappa(\mathfrak p')$ 是域扩张。
\end{proof}

\begin{lemma}
\label{algebra-lemma-domain-image-dense-set-points-generic-point}
设 $R$ 为整环，$\varphi : R \to S$ 为环映射。以下条件等价：
\begin{enumerate}
\item 环映射 $R \to S$ 是单射。
\item $\Spec(S) \to \Spec(R)$ 的像包含一个稠密点集。
\item 存在素理想 $\mathfrak q \subset S$，它在 $R$ 中的逆像是
$(0)$。
\end{enumerate}
\end{lemma}

\begin{proof}
设 $K$ 为整环 $R$ 的分式域。假设 $R \to S$ 为单射。由于局部化正合，
可知 $K \to S \otimes_R K$ 是单射。因此，由引理
\ref{algebra-lemma-in-image}，存在一个映到 $(0)$ 的素理想。

\medskip\noindent
注意 $(0)$ 在 $\Spec(R)$ 中稠密，所以最后一个条件蕴含第二个条件。

\medskip\noindent
假设第二个条件成立。设 $f \in R$、$f \not = 0$。由于 $R$ 是整环，
可知 $V(f)$ 是 $R$ 的真闭子集。
% P01-E0327: V(f) has the wrong ambient space in frozen source; recorded in ERRATA_CANDIDATES.jsonl.
由假设，存在 $S$ 的素理想 $\mathfrak q$，使
$\varphi(f) \not \in \mathfrak q$。故 $\varphi(f) \not = 0$，
从而 $R \to S$ 是单射。
\end{proof}

\begin{lemma}
\label{algebra-lemma-injective-minimal-primes-in-image}
设 $R \subset S$ 为单射环映射。则 $\Spec(S) \to \Spec(R)$ 命中
所有极小素理想。
\end{lemma}

\begin{proof}
设 $\mathfrak p \subset R$ 为极小素理想。此时 $R_{\mathfrak p}$
有唯一素理想。因此只须证明 $S_{\mathfrak p}$ 非零。这由局部化正合
得出，见命题 \ref{algebra-proposition-localization-exact}。
\end{proof}

\begin{lemma}
\label{algebra-lemma-image-dense-generic-points}
设 $R \to S$ 为环映射。以下条件等价：
\begin{enumerate}
\item $R \to S$ 的核由幂零元组成。
\item $R$ 的极小素理想属于 $\Spec(S) \to \Spec(R)$ 的像。
\item $\Spec(S) \to \Spec(R)$ 的像在 $\Spec(R)$ 中稠密。
\end{enumerate}
\end{lemma}

\begin{proof}
令 $I = \Ker(R \to S)$。注意
$\sqrt{(0)} = \bigcap_{\mathfrak q \subset S} \mathfrak q$，见引理
\ref{algebra-lemma-Zariski-topology}。所以
$\sqrt{I} = \bigcap_{\mathfrak q \subset S} R \cap \mathfrak q$。
于是 $V(I) = V(\sqrt{I})$ 是 $\Spec(S) \to \Spec(R)$ 之像的闭包。
这说明 (1) 与 (3) 等价。显然 (2) 蕴含 (3)。最后假设 (1)。
可以用 $R/I$ 代替 $R$，用 $S/IS$ 代替 $S$，而不影响谱的拓扑和
映射。因此，(1) $\Rightarrow$ (2) 由引理
\ref{algebra-lemma-injective-minimal-primes-in-image} 得出。
\end{proof}

\begin{lemma}
\label{algebra-lemma-minimal-prime-image-minimal-prime}
设 $R \to S$ 为环映射。若极小素理想 $\mathfrak p \subset R$ 属于
$\Spec(S) \to \Spec(R)$ 的像，则它是某个极小素理想的像。
\end{lemma}

\begin{proof}
设 $\mathfrak p = \mathfrak q \cap R$。选取满足
$\mathfrak r \subset \mathfrak q$ 的极小素理想
$\mathfrak r \subset S$，见引理 \ref{algebra-lemma-Zariski-topology}。
由 $\mathfrak p$ 的极小性可知
$\mathfrak p = \mathfrak r \cap R$。
\end{proof}

\begin{lemma}
\label{algebra-lemma-intersection-not-zero}
设 $A \subset B$ 是整环的包含，且诱导分式域的代数扩张。若
$J \subset B$ 为非零理想，则 $A \cap J$ 也非零。因此，
$\Spec(B)$ 的真闭子集的像在 $\Spec(A)$ 中不稠密。
\end{lemma}

\begin{proof}
取非零元素 $x \in J$。由于 $x$ 在 $A$ 的分式域上代数，存在
$d \geq 1$ 及 $a_0, \ldots, a_d \in A$，其中
$a_0, a_d \not = 0$，使得在 $B$ 中
$a_d x^d + a_{d - 1} x^{d - 1} + \ldots + a_0 = 0$。
于是 $a_0 \in A \cap J$。
\end{proof}

\section{诺特环}
\label{algebra-section-Noetherian}

\noindent
若环 $R$ 的任意理想都有限生成，则称 $R$ 是{\it 诺特的}。
这显然等价于 $R$ 的理想满足升链条件。由引理 \ref{algebra-lemma-cohen}，
只须检验 $R$ 的每个素理想都有限生成。

\begin{lemma}
\label{algebra-lemma-Noetherian-permanence}
\begin{slogan}
诺特性质在取有限型扩张与局部化时保持。
\end{slogan}
诺特环上的任意有限生成环都是诺特环。诺特环的任意局部化都是诺特环。
\end{lemma}

\begin{proof}
关于局部化的断言来自如下事实：任意理想 $J \subset S^{-1}R$ 都形如
$I \cdot S^{-1}R$。诺特环 $R$ 的任意商 $R/I$ 都是诺特的，
因为任意理想 $\overline{J} \subset R/I$ 都形如 $J/I$，其中
$I \subset J \subset R$ 为理想。因此只须证明：若 $R$ 诺特，
则 $R[X]$ 也诺特。设
$J_1 \subset J_2 \subset \ldots$ 是 $R[X]$ 中的理想升链。
考虑理想 $I_{i, d}$，其定义为 $J_i$ 中次数为 $d$ 的多项式的
首项系数在 $R$ 中组成的理想。显然，只要 $i \leq i'$ 且
$d \leq d'$，就有 $I_{i, d} \subset I_{i', d'}$。由 $R$ 中的
升链条件，所有 $I_{i, d}$ 中至多只有有限个互异理想。
（提示：$\mathbf{N} \times \mathbf{N}$ 的任意无限子集都包含
一条递增无限序列。）取充分大的 $i_0$，使对所有 $i \geq i_0$ 及
所有 $d$，都有 $I_{i, d} = I_{i_0, d}$。设对某个
$i \geq i_0$ 有 $f \in J_i$。对次数 $d = \deg(f)$ 作归纳，
证明 $f \in J_{i_0}$。事实上，存在 $g\in J_{i_0}$，其次数为 $d$，
且首项系数与 $f$ 相同。由归纳，$f - g \in J_{i_0}$，结论得证。
\end{proof}

\begin{lemma}
\label{algebra-lemma-Noetherian-power-series}
若 $R$ 是诺特环，则形式幂级数环 $R[[x_1, \ldots, x_n]]$ 也是诺特环。
\end{lemma}

\begin{proof}
由于
$R[[x_1, \ldots, x_{n + 1}]] \cong R[[x_1, \ldots, x_n]][[x_{n + 1}]]$，
只须证明：若 $R$ 诺特，则 $R[[x]]$ 诺特。设 $I \subset R[[x]]$
为理想。须证明 $I$ 有限生成。对每个整数 $d$，记
$I_d = \{a \in R \mid ax^d + \text{h.o.t.} \in I\}$。由于 $R$ 诺特，
可知 $I_0 \subset I_1 \subset \ldots$ 稳定。选取 $d_0$，使
$I_{d_0} = I_{d_0 + 1} = \ldots$。对每个 $d \leq d_0$，选取元素
$f_{d, j} \in I \cap (x^d)$（$j = 1, \ldots, n_d$），使得若写成
$f_{d, j} = a_{d, j}x^d + \text{h.o.t}$，则
$I_d = (a_{d, j})$。记
$I' = (\{f_{d, j}\}_{d = 0, \ldots, d_0, j = 1, \ldots, n_d})$。
显然 $I' \subset I$。取 $f \in I$。首先可选取 $c_{d, i} \in R$，使
$$
f - \sum c_{d, i} f_{d, i} \in (x^{d_0 + 1}) \cap I.
$$
接着可选取 $c_{i, 1} \in R$（$i = 1, \ldots, n_{d_0}$），使
$$
f - \sum c_{d, i} f_{d, i} - \sum c_{i, 1}xf_{d_0, i} \in (x^{d_0 + 2}) \cap I.
$$
再选取 $c_{i, 2} \in R$（$i = 1, \ldots, n_{d_0}$），使
$$
f - \sum c_{d, i} f_{d, i} - \sum c_{i, 1}xf_{d_0, i}
- \sum c_{i, 2}x^2f_{d_0, i}
\in (x^{d_0 + 3}) \cap I.
$$
依此类推。最后可知
$$
f = \sum c_{d, i} f_{d, i} +
\sum\nolimits_i (\sum\nolimits_e c_{i, e} x^e)f_{d_0, i}
$$
属于 $I'$，如所需。
% P01-E0328: element/ideal “contained in” wording recorded in ERRATA_CANDIDATES.jsonl.
\end{proof}

\noindent
下面的引理虽然容易，却很有用，因为有限型 $\mathbf{Z}$-代数经常出现于
一种称为“绝对诺特约化”的技巧中。

\begin{lemma}
\label{algebra-lemma-obvious-Noetherian}
域上的任意有限型代数都是诺特的。$\mathbf{Z}$ 上的任意有限型代数
都是诺特的。
\end{lemma}

\begin{proof}
这直接来自引理 \ref{algebra-lemma-Noetherian-permanence}，以及域都是诺特环、
$\mathbf{Z}$ 是诺特环（因为它是主理想整环）这两个事实。
% P01-E0329: missing article in “is Noetherian ring” recorded in ERRATA_CANDIDATES.jsonl.
\end{proof}

\begin{lemma}
\label{algebra-lemma-Noetherian-finite-type-is-finite-presentation}
设 $R$ 为诺特环。
\begin{enumerate}
\item 任意有限 $R$-模都是有限表示的。
\item 有限 $R$-模的任意子模都是有限的。
\item 任意有限型 $R$-代数在 $R$ 上有限表示。
\end{enumerate}
\end{lemma}

\begin{proof}
设 $M$ 为有限 $R$-模。由引理
\ref{algebra-lemma-trivial-filter-finite-module}，可找到 $M$ 的有限滤过，
其逐次商形如 $R/I$。由于任意理想都有限生成，每个商 $R/I$
都有限表示。因此由引理 \ref{algebra-lemma-extension}，$M$ 有限表示。
这证明了 (1)。

\medskip\noindent
设 $N \subset M$ 为子模。由于 $M$ 有限，商 $M/N$ 有限。
由 (1)，$M/N$ 有限表示。因此，由引理 \ref{algebra-lemma-extension} 的 (5)，
$N$ 有限。这证明了 (2)。

\medskip\noindent
为证明 (3)，注意由引理 \ref{algebra-lemma-Noetherian-permanence}，
$R[x_1, \ldots, x_n]$ 的任意理想都有限生成。
\end{proof}

\begin{lemma}
\label{algebra-lemma-Noetherian-topology}
若 $R$ 是诺特环，则 $\Spec(R)$ 是诺特拓扑空间，见拓扑学，定义
\ref{topology-definition-noetherian}。
\end{lemma}

\begin{proof}
这是因为 $\Spec(R)$ 的任意闭子集都唯一地形如 $V(I)$，其中 $I$
是根理想，见引理 \ref{algebra-lemma-Zariski-topology}。并且这一对应反转包含关系。
所以结论由定义得出。
\end{proof}

\begin{lemma}
\label{algebra-lemma-Noetherian-irreducible-components}
\begin{slogan}
诺特仿射概形只有有限多个泛点。
\end{slogan}
若 $R$ 是诺特环，则 $\Spec(R)$ 只有有限多个不可约分支。换言之，
$R$ 只有有限多个极小素理想。
\end{lemma}

\begin{proof}
由引理 \ref{algebra-lemma-Noetherian-topology} 和拓扑学，引理
\ref{topology-lemma-Noetherian}，可知不可约分支只有有限多个。
由引理 \ref{algebra-lemma-irreducible}，它们对应于 $R$ 的极小素理想。
\end{proof}

\begin{lemma}
\label{algebra-lemma-Noetherian-base-change-finite-type}
设 $R \to S$ 为环映射，$R \to R'$ 为有限型映射。若 $S$ 诺特，
则基变换 $S' = R' \otimes_R S$ 诺特。
\end{lemma}

\begin{proof}
由引理 \ref{algebra-lemma-base-change-finiteness}，有限型在基变换下保持。
因此 $S \to S'$ 是有限型的。由于 $S$ 诺特，可以应用引理
\ref{algebra-lemma-Noetherian-permanence}。
\end{proof}

\begin{lemma}
\label{algebra-lemma-Noetherian-field-extension}
设 $k$ 为域，$R$ 为诺特 $k$-代数。若 $K/k$ 是有限生成域扩张，
则 $K \otimes_k R$ 诺特。
\end{lemma}

\begin{proof}
由于 $K/k$ 是有限生成域扩张，存在有限生成 $k$-代数
$B \subset K$，使 $K$ 是 $B$ 的分式域。换言之，
$K = S^{-1}B$，其中 $S = B \setminus \{0\}$。于是
$K \otimes_k R = S^{-1}(B \otimes_k R)$。由引理
\ref{algebra-lemma-Noetherian-base-change-finite-type}，$B \otimes_k R$
诺特。最后，由引理 \ref{algebra-lemma-Noetherian-permanence}，
$K \otimes_k R = S^{-1}(B \otimes_k R)$ 诺特。
\end{proof}

\noindent
下面是几个有时很有用的有趣引理。

\begin{lemma}
\label{algebra-lemma-subring-of-local-ring}
设 $R$ 为环，$\mathfrak p \subset R$ 为素理想。在以下每种情形下，
都存在 $f \in R$、$f \not \in \mathfrak p$，使
$R_f \to R_\mathfrak p$ 是单射：
\begin{enumerate}
\item $R$ 是整环；
\item $R$ 诺特；或
\item $R$ 约化且只有有限多个极小素理想。
\end{enumerate}
\end{lemma}

\begin{proof}
若 $R$ 是整环，则 $R \subset R_\mathfrak p$，故取 $f = 1$ 即可。
若 $R$ 诺特，则 $R \to R_\mathfrak p$ 的核 $I$ 是有限生成理想，
因而可找到 $f \in R$、$f \not \in \mathfrak p$，使 $IR_f = 0$。
对此 $f$，映射 $R_f \to R_\mathfrak p$ 是单射，故 $f$ 可用。
若 $R$ 约化，且极小素理想只有有限多个
$\mathfrak p_1, \ldots, \mathfrak p_n$，则可选择
$f \in \bigcap_{\mathfrak p_i \not \subset \mathfrak p} \mathfrak p_i$、
$f \not \in \mathfrak p$。事实上，若 $\mathfrak{p}_i\not\subset
\mathfrak{p}$，则存在 $f_i \in \mathfrak{p}_i$、
$f_i \not\in \mathfrak{p}$，并且 $f = \prod f_i$ 可用。对此 $f$，
有 $R_f \subset R_\mathfrak p$，因为 $R_f$ 的极小素理想对应于
$R_\mathfrak p$ 的极小素理想，并且可以应用引理
\ref{algebra-lemma-reduced-ring-sub-product-fields}（略去若干细节）。
\end{proof}

\begin{lemma}
\label{algebra-lemma-surjective-endo-noetherian-ring-is-iso}
诺特环的任意满自同态都是同构。
\end{lemma}

\begin{proof}
若 $f : R \to R$ 是这样的自同态但不是单射，则
$$
\Ker(f) \subset \Ker(f \circ f) \subset
\Ker(f \circ f \circ f) \subset \ldots
$$
将是一条严格递增的理想链。
\end{proof}

\section{局部幂零理想}
\label{algebra-section-locally-nilpotent}

\noindent
定义如下。

\begin{definition}
\label{algebra-definition-locally-nilpotent-ideal}
设 $R$ 为环，$I \subset R$ 为理想。若对每个 $x \in I$，都存在
$n \in \mathbf{N}$ 使 $x^n = 0$，则称 $I$ 是{\it 局部幂零的}。
若存在 $n \in \mathbf{N}$ 使 $I^n = 0$，则称 $I$ 是{\it 幂零的}。
\end{definition}

\begin{example}
\label{algebra-example-locally-nilpotent-not-nilpotent}
设 $R = k[x_n | n \in \mathbf{N}]$ 为域 $k$ 上含无穷多个变量的
多项式环。设 $I$ 为由 $x_n^n$（$n \in \mathbf{N}$）生成的理想，
并令 $S = R/I$。则由诸 $x_n$ 的像（$n \in \mathbf{N}$）生成的理想
$J \subset S$ 局部幂零，但并不幂零。事实上，由于幂零元的
$S$-线性组合仍为幂零元，要证明 $J$ 局部幂零，只须注意其所有生成元
都是幂零的（这显然成立）。另一方面，对每个 $n \in \mathbf{N}$，
都有 $x_{n + 1}^n \not \in I$，故 $J^n \not = 0$。所以 $J$ 不幂零。
\end{example}

\begin{lemma}
\label{algebra-lemma-locally-nilpotent}
设 $R \to R'$ 为环映射，$I \subset R$ 为局部幂零理想。
则 $IR'$ 是 $R'$ 的局部幂零理想。
\end{lemma}

\begin{proof}
这来自如下事实：若 $x, y \in R'$ 都幂零，则 $x + y$ 也幂零。
具体地，若 $x^n = 0$ 且 $y^m = 0$，则
$(x + y)^{n + m - 1} = 0$。
\end{proof}

\begin{lemma}
\label{algebra-lemma-locally-nilpotent-unit}
设 $R$ 为环，$I \subset R$ 为局部幂零理想。元素 $x$ 是 $R$ 的单位，
当且仅当 $x$ 在 $R/I$ 中的像是单位。
\end{lemma}

\begin{proof}
若 $x$ 是 $R$ 中的单位，则其像显然是 $R/I$ 中的单位。只须证明反向。
假设 $y \in R$ 的像是 $x$ 的像在 $R/I$ 中的逆元。于是对某个
$z \in I$，有 $xy = 1 - z$。这意味着模 $xR$ 有
$1\equiv z$。由于 $z$ 属于局部幂零理想 $I$，对某个充分大的 $N$，
有 $z^N = 0$。故模 $xR$ 有 $1 = 1^N \equiv z^N = 0$。
换言之，$x$ 整除 $1$，所以是单位。
\end{proof}

\begin{lemma}
\label{algebra-lemma-Noetherian-power}
\begin{slogan}
诺特环的一个理想若逐元素幂零，则该理想幂零。
\end{slogan}
设 $R$ 为诺特环，$I, J$ 为 $R$ 的理想。假设
$J \subset \sqrt{I}$。则对某个 $n$ 有 $J^n \subset I$。
特别地，在诺特环中，“局部幂零理想”与“幂零理想”这两个概念相同。
\end{lemma}

\begin{proof}
写成 $J = (f_1, \ldots, f_s)$。由假设 $f_i^{d_i} \in I$。
取 $n = d_1 + d_2 + \ldots + d_s + 1$。
\end{proof}

\begin{lemma}
\label{algebra-lemma-lift-idempotents}
设 $R$ 为环，$I \subset R$ 为局部幂零理想。则
$R \to R/I$ 在幂等元上诱导双射。
\end{lemma}

\begin{proof}[引理 \ref{algebra-lemma-lift-idempotents} 的第一种证明]
由于 $I$ 局部幂零，它包含于每个素理想。因此
$\Spec(R/I) = V(I) = \Spec(R)$。所以该引理由
引理 \ref{algebra-lemma-disjoint-decomposition} 得出。
\end{proof}

\begin{proof}[引理 \ref{algebra-lemma-lift-idempotents} 的第二种证明]
设 $\overline{e} \in R/I$ 是幂等元。须将 $\overline{e}$ 提升为 $R$
的幂等元。

\medskip\noindent
首先任选 $\overline{e}$ 的一个提升 $f \in R$，并令
$x = f^2 - f$。则 $x \in I$，所以 $x$ 幂零（因为 $I$ 局部幂零）。
现在令 $J$ 为 $R$ 中由 $x$ 生成的理想。由于它由幂零元 $x$ 生成，
$J$ 是幂零的，而不只是局部幂零的。

\medskip\noindent
现在假设已经找到 $\overline{e}$ 的一个提升 $e \in R$，使对某个
$k \geq 1$ 有 $e^2 - e \in J^k$。令
$e' = e - (2e - 1)(e^2 - e) = 3e^2 - 2e^3$；它是
$\overline{e}$ 的另一个提升（因为 $\overline{e}$ 的幂等性给出
$e^2 - e \in I$）。简单计算可得
$$
(e')^2 - e' = (4e^2 - 4e - 3)(e^2 - e)^2 \in J^{2k}
$$

\medskip\noindent
所以，我们从 $\overline{e}$ 的一个满足 $e^2 - e \in J^k$ 的提升
$e$ 出发，得到了 $\overline{e}$ 的一个满足
$(e')^2 - e' \in J^{2k}$ 的提升 $e'$。以此方式可逐次改进近似
（从 $e = f$ 开始，它对 $k = 1$ 符合要求）。最终到达
$J^k = 0$ 的阶段；此时有 $\overline{e}$ 的一个提升 $e$ 满足
$e^2 - e \in J^k = 0$，即该 $e$ 是幂等元。

\medskip\noindent
由此可见，若 $\overline{e} \in R/I$ 是任意幂等元，则存在一个
$\overline{e}$ 的提升，并且该提升是 $R$ 的幂等元。还须证明此提升唯一。设 $e_1$ 与 $e_2$
为两个这样的提升。须证明 $e_1 = e_2$。

\medskip\noindent
由 $e_1$ 与 $e_2$ 的定义，有 $e_1 \equiv e_2
\mod I$，并且 $e_1$ 与 $e_2$ 都幂等。由
$e_1 \equiv e_2 \mod I$ 可知 $e_1 - e_2 \in I$，故
$e_1 - e_2$ 幂零（因为 $I$ 局部幂零）。直接计算（使用 $e_1$ 与
$e_2$ 的幂等性）得到 $(e_1 - e_2)^3 = e_1 - e_2$。由此归纳可知，
对任意正奇数 $k$，都有 $(e_1 - e_2)^k = e_1 - e_2$。由于所有充分大的
$k$ 都满足 $(e_1 - e_2)^k = 0$（因为 $e_1 - e_2$ 幂零），
可知 $e_1 - e_2 = 0$，即 $e_1 = e_2$，证明完成。
\end{proof}

\begin{lemma}
\label{algebra-lemma-lift-idempotents-noncommutative}
设 $A$ 为可能不交换的代数，$e \in A$ 为满足
$x = e^2 - e$ 幂零的元素。则存在形如
$e' = e + x(\sum a_{i, j}e^ix^j) \in A$ 的幂等元，其中
$a_{i, j} \in \mathbf{Z}$。
\end{lemma}

\begin{proof}
考虑环 $R_n = \mathbf{Z}[e]/((e^2 - e)^n)$。显然，若能对每个 $R_n$
证明结论，则引理成立。在 $R_n$ 中考虑理想 $I = (e^2 - e)$，
并应用引理 \ref{algebra-lemma-lift-idempotents}。
\end{proof}

\begin{lemma}
\label{algebra-lemma-lift-nth-roots}
设 $R$ 为环，$I \subset R$ 为局部幂零理想。设 $n \geq 1$ 为在
$R/I$ 中可逆的整数。则
\begin{enumerate}
\item $n$ 次幂映射 $1 + I \to 1 + I$，$1 + x \mapsto (1 + x)^n$
是双射；
\item $R$ 的一个单位是 $n$ 次幂，当且仅当其在 $R/I$ 中的像是
$n$ 次幂。
% P01-E0330: “a nth power” article defect recorded in ERRATA_CANDIDATES.jsonl.
\end{enumerate}
\end{lemma}

\begin{proof}
设 $a \in R$ 为单位，且其在 $R/I$ 中的像与某个 $b \in R$ 的
$b^n$ 的像相同。则 $b$ 是单位（引理
\ref{algebra-lemma-locally-nilpotent-unit}），且对某个 $x \in I$ 有
$ab^{-n} = 1 + x$。由 (1)，$ab^{-n} = c^n$。所以 (2) 由 (1) 得出。

\medskip\noindent
证明 (1)。这是因为映射 $1 + x \mapsto (1 + x)^n$ 有逆映射。
具体地，可以考虑把 $1 + x$ 映为下式的映射：
\begin{align*}
(1 + x)^{1/n}
& =
1 + {1/n \choose 1}x +
{1/n \choose 2}x^2 +
{1/n \choose 3}x^3 + \ldots \\
& =
1 + \frac{1}{n} x + \frac{1 - n}{2n^2}x^2 +
\frac{(1 - n)(1 - 2n)}{6n^3}x^3 + \ldots
\end{align*}
这与初等微积分中一样。该式有意义，因为当 $k \gg 0$ 时
$x^k = 0$，所以级数是有限的；而且每个系数
${1/n \choose k} \in \mathbf{Z}[1/n]$（细节略去；注意由引理
\ref{algebra-lemma-locally-nilpotent-unit}，$n$ 在 $R$ 中可逆）。
\end{proof}



\section{一个有趣的问题}
\label{algebra-section-curiosity}

\noindent
引理 \ref{algebra-lemma-disjoint-implies-product} 说明了当某个理想
$I \subset R$ 的 $V(I)$ 为开集时会发生什么。但若
$\Spec(S^{-1}R)$ 在 $\Spec(R)$ 中是闭的，又会怎样？以下两个引理
给出部分答案。更多信息见第 \ref{algebra-section-pure-ideals} 节。

\begin{lemma}
\label{algebra-lemma-invert-closed-quotient}
设 $R$ 为环，$S \subset R$ 为乘法子集。假设映射
$\Spec(S^{-1}R) \to \Spec(R)$ 的像是闭的。则对某个理想
$I \subset R$，有 $S^{-1}R \cong R/I$。
\end{lemma}

\begin{proof}
令 $I = \Ker(R \to S^{-1}R)$，于是 $V(I)$ 包含该像。设该像为闭子集
$V(I') \subset \Spec(R)$，其中 $I' \subset R$ 为某个理想。因此
$V(I') \subset V(I)$。对 $f \in I'$，元 $f/1 \in S^{-1}R$ 属于每个
素理想。因此对某个 $n \geq 1$，$f^n$ 在 $S^{-1}R$ 中映为零
（引理 \ref{algebra-lemma-Zariski-topology}）。
% P01-E0331: source uses “is contained in every prime ideal” for an element; recorded in ERRATA_CANDIDATES.jsonl.
故 $V(I') = V(I)$。由此可知，每个 $g \in S$ 模 $I$ 都可逆。因此得到
环映射 $R/I \to S^{-1}R$ 与 $S^{-1}R \to R/I$。由 $I$ 的选取，
第一个映射是单射。第二个映射是
$S^{-1}R \to S^{-1}(R/I) = R/I$；由于局部化是正合的，其核为
$S^{-1}I$。因为 $S^{-1}I = 0$，第二个映射也是单射。因此
$S^{-1}R \cong R/I$。
\end{proof}

\begin{lemma}
\label{algebra-lemma-invert-closed-split}
设 $R$ 为环，$S \subset R$ 为乘法子集。假设映射
$\Spec(S^{-1}R) \to \Spec(R)$ 的像是闭的。若 $R$ 是诺特环，或
$\Spec(R)$ 是诺特拓扑空间，或 $S$ 作为幺半群是有限生成的，则对某个环
$R'$，有 $R \cong S^{-1}R \times R'$。
\end{lemma}

\begin{proof}
由引理 \ref{algebra-lemma-invert-closed-quotient}，对某个理想 $I \subset R$，
有 $S^{-1}R \cong R/I$。由引理 \ref{algebra-lemma-disjoint-implies-product}，
只须证明 $V(I)$ 是开集。若 $R$ 是诺特环，则 $\Spec(R)$ 是诺特拓扑空间，
见引理 \ref{algebra-lemma-Noetherian-topology}。若 $\Spec(R)$ 是诺特拓扑空间，
则补集 $\Spec(R) \setminus V(I)$ 是拟紧的，见拓扑学，引理
\ref{topology-lemma-Noetherian-quasi-compact}。因此存在有限多个
$f_1, \ldots, f_n \in I$，使得 $V(I) = V(f_1, \ldots, f_n)$。
由于每个 $f_i$ 在 $S^{-1}R$ 中都映为零，故存在 $g \in S$，使对
$i = 1, \ldots, n$ 都有 $gf_i = 0$。因而按所需有
$D(g) = V(I)$。若 $S$ 作为幺半群是有限生成的，设 $S$ 由
$g_1, \ldots, g_m$ 生成，则
$S^{-1}R \cong R_{g_1 \ldots g_m}$，从而
$V(I) = D(g_1 \ldots g_m)$。
\end{proof}



\section{希尔伯特零点定理}
\label{algebra-section-nullstellensatz}

\begin{theorem}[希尔伯特零点定理]
\label{algebra-theorem-nullstellensatz}
设 $k$ 为域。
\begin{enumerate}
\item
\label{algebra-item-finite-kappa}
对任意极大理想 $\mathfrak m \subset k[x_1, \ldots, x_n]$，域扩张
$\kappa(\mathfrak m)/k$ 是有限的。
\item
\label{algebra-item-polynomial-ring-Jacobson}
任意根理想 $I \subset k[x_1, \ldots, x_n]$ 都是包含它的极大理想之交。
\end{enumerate}
任意有限型 $k$-代数也满足同样的结论。
\end{theorem}

\begin{proof}
只须对多项式代数 $k[x_1, \ldots, x_n]$ 证明定理的
(\ref{algebra-item-finite-kappa})，因为任意有限生成 $k$-代数都是这种多项式代数
的商。对 $n$ 作归纳。$n = 0$ 的情形显然。设 $\mathfrak m$ 是
$k[x_1, \ldots, x_n]$ 中的极大理想。令
$\mathfrak p \subset k[x_n]$ 为 $\mathfrak m$ 与 $k[x_n]$ 的交。

\medskip\noindent
若 $\mathfrak p \not = (0)$，则 $\mathfrak p$ 是极大理想，并由一个
首一不可约多项式 $P$ 生成（因为 $k[x_n]$ 中有欧几里得算法）。于是
$k' = k[x_n]/\mathfrak p$ 是 $k$ 的有限域扩张，并包含于
$\kappa(\mathfrak m)$。在此情形，得到满射
$$
k'[x_1, \ldots, x_{n-1}]
\to
k'[x_1, \ldots, x_n] =
k' \otimes_k k[x_1, \ldots, x_n]
\longrightarrow
\kappa(\mathfrak m)
$$
因此由归纳假设，$\kappa(\mathfrak m)$ 是 $k'$ 的有限扩张。所以
$\kappa(\mathfrak m)$ 也是 $k$ 的有限扩张。

\medskip\noindent
若 $\mathfrak p = (0)$，考虑环扩张
$k[x_n] \subset k[x_1, \ldots, x_n]/\mathfrak m$。这是有限生成环扩张，
故由引理 \ref{algebra-lemma-obvious-Noetherian} 与
\ref{algebra-lemma-Noetherian-finite-type-is-finite-presentation}，它是有限表示的。
因而由定理 \ref{algebra-theorem-chevalley}，
$\Spec(k[x_1, \ldots, x_n]/\mathfrak m)$ 在 $\Spec(k[x_n])$ 中的像
是可构造的。由于该像包含 $(0)$，可知它包含某个非零
$f\in k[x_n]$ 的标准开集 $D(f)$。显然 $D(f)$ 是无限的，这与
$k[x_1, \ldots, x_n]/\mathfrak m$ 为域（因而其谱只有一个点）的
假设矛盾。

\medskip\noindent
证明 (\ref{algebra-item-polynomial-ring-Jacobson})。设
$I \subset R$ 为根理想，其中 $R$ 是 $k$ 上有限型的。取
$f \in R$ 且 $f \not \in I$。须找一个极大理想
$\mathfrak m \subset R$，使得 $I \subset \mathfrak m$ 且
$f \not \in \mathfrak m$。环 $(R/I)_f$ 非零，因为若其中 $1 = 0$，
则 $f^n \in I$；由于 $I$ 是根理想，这将推出 $f \in I$，与对 $f$
的假设矛盾。因此可在 $(R/I)_f$ 中选取一个极大理想
$\mathfrak m'$，见引理 \ref{algebra-lemma-Zariski-topology}。令
$\mathfrak m \subset R$ 为 $\mathfrak m'$ 在 $R$ 中的逆像。可见
$I \subset \mathfrak m$ 且 $f \not \in \mathfrak m$。若能证明
$\mathfrak m$ 是 $R$ 的极大理想，结论即得。显然有
$$
k \subset R/\mathfrak m \subset \kappa(\mathfrak m').
$$
由 (\ref{algebra-item-finite-kappa})，域扩张 $\kappa(\mathfrak m')/k$ 是有限的。
因而由域论，引理
\ref{fields-lemma-subalgebra-algebraic-extension-field}，$R/\mathfrak m$
是域。所以 $\mathfrak m$ 是极大理想，证明完成。
\end{proof}

\begin{lemma}
\label{algebra-lemma-field-finite-type-over-domain}
设 $R$ 为环，$K$ 为域。若 $R \subset K$ 且 $K$ 在 $R$ 上为有限型，
则存在 $f \in R$，使 $R_f$ 为域，并且 $K/R_f$ 是有限域扩张。
\end{lemma}

\begin{proof}
由引理 \ref{algebra-lemma-characterize-image-finite-type}，存在非空开集
$U \subset \Spec(R)$，它包含于 $\Spec(K) \to \Spec(R)$ 的像
$\{(0)\}$ 中。选取 $f \in R$、$f \not = 0$，使
$D(f) \subset U$，即 $D(f) = \{(0)\}$。于是 $R_f$ 是谱恰有一个点
的整环，故 $R_f$ 是域。此时 $K$ 是域 $R_f$ 上的有限生成代数，
因而由希尔伯特零点定理（定理 \ref{algebra-theorem-nullstellensatz}），
它是 $R_f$ 的有限域扩张。
\end{proof}



\section{雅可布森环}
\label{algebra-section-ring-jacobson}

\noindent
设 $R$ 为环。$\Spec(R)$ 的闭点就是 $R$ 的极大理想。代数几何中自然
出现的环往往有许多极大理想，例如域上或 $\mathbf{Z}$ 上的有限型代数。
我们将证明，这些都是雅可布森环的例子。

\begin{definition}
\label{algebra-definition-ring-jacobson}
设 $R$ 为环。若每个根理想 $I$ 都是包含它的极大理想之交，则称 $R$ 为
\textit{雅可布森环}。
\end{definition}

\begin{lemma}
\label{algebra-lemma-finite-type-field-Jacobson}
域上的任意有限型代数都是雅可布森环。
\end{lemma}

\begin{proof}
这由定理 \ref{algebra-theorem-nullstellensatz} 与定义
\ref{algebra-definition-ring-jacobson} 得出。
\end{proof}

\begin{lemma}
\label{algebra-lemma-jacobson-prime}
设 $R$ 为环。若 $R$ 的每个素理想都是包含它的极大理想之交，则 $R$ 是
雅可布森环。
\end{lemma}

\begin{proof}
这立即来自如下事实：每个根理想 $I \subset R$ 都是包含它的素理想之交。
见引理 \ref{algebra-lemma-Zariski-topology}。
\end{proof}

\begin{lemma}
\label{algebra-lemma-jacobson}
环 $R$ 是雅可布森环，当且仅当 $\Spec(R)$ 是雅可布森空间；见拓扑学，
定义 \ref{topology-definition-space-jacobson}。
\end{lemma}

\begin{proof}
先设 $R$ 是雅可布森环。令 $Z \subset \Spec(R)$ 为闭子集。须证明
$Z$ 中的闭点集在 $Z$ 中稠密。令 $U \subset \Spec(R)$ 为开集，且
$U \cap Z$ 非空。须证明 $Z \cap U$ 含有 $\Spec(R)$ 的一个闭点。
由于标准开集构成 $\Spec(R)$ 拓扑的一组基，可以设 $U = D(f)$。
由引理 \ref{algebra-lemma-Zariski-topology}，可以设 $Z = V(I)$，其中 $I$ 是
根理想。又有 $f \not \in I$。按假设，存在极大理想
$\mathfrak m \subset R$，使 $I \subset \mathfrak m$ 而
$f \not\in \mathfrak m$。因此按所需，
$\mathfrak m \in D(f) \cap V(I) = U \cap Z$。

\medskip\noindent
反之，设 $\Spec(R)$ 是雅可布森空间。令 $I \subset R$ 为根理想，并令
$J = \cap_{I \subset \mathfrak m} \mathfrak m$ 为包含 $I$ 的极大理想之交。
显然 $J$ 是根理想，$V(J) \subset V(I)$，并且 $V(J)$ 是 $V(I)$ 中
包含 $V(I)$ 全部闭点的最小闭子集。由假设可知 $V(J) = V(I)$。
但引理 \ref{algebra-lemma-Zariski-topology} 表明，扎里斯基闭集与根理想之间存在
双射，故按所需有 $I = J$。
\end{proof}

\begin{lemma}
\label{algebra-lemma-characterize-jacobson}
设 $R$ 为环。若 $R$ 不是雅可布森环，则存在素理想
$\mathfrak p \subset R$ 及元素 $f \in R$，满足：
\begin{enumerate}
\item $\mathfrak p$ 不是极大理想；
\item $f \not \in \mathfrak p$；
\item $V(\mathfrak p) \cap D(f) = \{\mathfrak p\}$；
\item $(R/\mathfrak p)_f$ 是域。
\end{enumerate}
另一方面，若 $R$ 是雅可布森环，则对任意满足 (1) 与 (2) 的
$(\mathfrak p, f)$，集合 $V(\mathfrak p) \cap D(f)$ 是无限的。
\end{lemma}

\begin{proof}
设 $R$ 不是雅可布森环。由引理 \ref{algebra-lemma-jacobson}，这意味着存在闭子集
$T \subset \Spec(R)$，其闭点集 $T_0 \subset T$ 在 $T$ 中不稠密。
选取 $f \in R$，使 $T_0 \subset V(f)$ 而 $T \not \subset V(f)$。
注意，若 $T = V(I)$，则 $T \cap D(f)$ 同胚于
$\Spec((R/I)_f)$；见引理 \ref{algebra-lemma-spec-closed} 与
\ref{algebra-lemma-standard-open}。由于任意环都有极大理想（引理
\ref{algebra-lemma-Zariski-topology}），可以选取空间 $T \cap D(f)$ 的一个闭点
$t$。于是 $t$ 对应于素理想 $\mathfrak p \subset R$；它不是极大理想
（因为 $t \not \in T_0$）。所以 (1) 成立。由构造，
$f \not \in \mathfrak p$，故 (2) 成立。由于 $t$ 是
$T \cap D(f)$ 的闭点，可知
$V(\mathfrak p) \cap D(f) = \{\mathfrak p\}$，即 (3) 成立。因此
$(R/\mathfrak p)_f$ 是谱只有一个点的整环，故 (4) 成立（例如合用引理
\ref{algebra-lemma-characterize-local-ring} 与
\ref{algebra-lemma-minimal-prime-reduced-ring}）。

\medskip\noindent
反之，设 $R$ 是雅可布森环，且 $(\mathfrak p, f)$ 满足 (1) 与 (2)。若
$V(\mathfrak p) \cap D(f) =
\{\mathfrak p, \mathfrak q_1, \ldots, \mathfrak q_t\}$，则由
$\mathfrak p \not = \mathfrak q_i$ 可知，存在元素 $g \in R$，使得
$g \not \in \mathfrak p$，但对所有 $i$ 都有 $g \in \mathfrak q_i$。
因而 $V(\mathfrak p) \cap D(fg) = \{\mathfrak p\}$。这是不可能的，
因为 $\Spec(R)$ 是雅可布森拓扑空间，所以 $\Spec(R)$ 的每个局部闭子集
至少含有一个闭点。
\end{proof}

\begin{lemma}
\label{algebra-lemma-pid-jacobson}
环 $\mathbf{Z}$ 是雅可布森环。更一般地，设环 $R$ 满足：
\begin{enumerate}
\item $R$ 是整环；
\item $R$ 是诺特环；
\item 任意非零素理想都是极大理想；
\item $R$ 有无限多个极大理想。
\end{enumerate}
则 $R$ 是雅可布森环。
\end{lemma}

\begin{proof}
设 $R$ 满足 (1)、(2)、(3) 与 (4)。该结论意指
$(0) = \bigcap_{\mathfrak m \subset R} \mathfrak m$。由于 $R$ 有无限多个
极大理想，只须证明任意非零 $x \in R$ 至多属于有限多个极大理想，换言之，
$V(x)$ 是有限的。
% P01-E0332: source uses “x is contained in ... maximal ideals” for an element; recorded in ERRATA_CANDIDATES.jsonl.
由引理 \ref{algebra-lemma-spec-closed}，$V(x)$ 同胚于 $\Spec(R/xR)$。由假设 (3)，
$R/xR$ 的每个素理想都是极小的，因而对应于 $\Spec(R/xR)$ 的一个不可约
分支（引理 \ref{algebra-lemma-irreducible}）。由于 $R/xR$ 是诺特环，拓扑空间
$\Spec(R/xR)$ 是诺特的（引理 \ref{algebra-lemma-Noetherian-topology}），且只有
有限多个不可约分支（拓扑学，引理 \ref{topology-lemma-Noetherian}）。
所以按所需，$V(x)$ 是有限的。
\end{proof}

\begin{example}
\label{algebra-example-infinite-product-fields-jacobson}
设 $A$ 为无限集。对每个 $\alpha \in A$，令 $k_\alpha$ 为域。我们断言
$R = \prod_{\alpha\in A} k_\alpha$ 是雅可布森环。首先注意，任意元素
$f \in R$ 都有形式 $f = ue$，其中 $u \in R$ 为单位，$e\in R$ 为
幂等元（留给读者）。因而 $D(f) = D(e)$，并且
$R_f = R_e = R/(1-e)$ 是 $R$ 的商。事实上，任何具有此性质的环都是
雅可布森环。设 $\mathfrak p \subset R$ 为素理想，且
$f \in R$、$f \not \in \mathfrak p$。须找出 $R$ 的一个极大理想
$\mathfrak m$，使 $\mathfrak p \subset \mathfrak m$ 且
$f \not\in \mathfrak m$。因为 $R_f$ 是 $R$ 的商，$R_f$ 的任意极大
理想都对应于 $R$ 的一个不含 $f$ 的极大理想。因此，选取 $R_f$ 的一个
包含 $\mathfrak p R_f$ 的极大理想即可得到结论。
\end{example}

\begin{example}
\label{algebra-example-not-jacobson}
设整环 $R$ 只有有限多个极大理想 $\mathfrak m_i$，
$i = 1, \ldots, n$。除非它是域，否则它不是雅可布森环。事实上，
在此情形，$(0)$ 不是极大理想之交，因为
$(0) \not =
\mathfrak m_1 \cap \mathfrak m_2 \cap \ldots \cap \mathfrak m_n
\supset \mathfrak m_1 \cdot \mathfrak m_2 \cdot \ldots
\cdot \mathfrak m_n \not = 0$。
特别地，离散赋值环，或任何至少有两个素理想的局部环，都不是雅可布森环。
\end{example}

\begin{lemma}
\label{algebra-lemma-finite-residue-extension-closed}
设 $R \to S$ 为环映射。令 $\mathfrak m \subset R$ 为极大理想，
$\mathfrak q \subset S$ 为位于 $\mathfrak m$ 之上的素理想，并且
$\kappa(\mathfrak q)/\kappa(\mathfrak m)$ 是代数域扩张。则
$\mathfrak q$ 是 $S$ 的极大理想。
\end{lemma}

\begin{proof}
考虑图表
$$
\xymatrix{
S \ar[r] & S/\mathfrak q \ar[r] & \kappa(\mathfrak q) \\
R \ar[r] \ar[u] & R/\mathfrak m \ar[u]
}
$$
可见 $\kappa(\mathfrak m) \subset S/\mathfrak q \subset
\kappa(\mathfrak q)$。由于域扩张
$\kappa(\mathfrak m) \subset \kappa(\mathfrak q)$ 是代数的，介于
$\kappa(\mathfrak m)$ 与 $\kappa(\mathfrak q)$ 之间的任意环都是域
（域论，引理 \ref{fields-lemma-subalgebra-algebraic-extension-field}）。
因此 $S/\mathfrak q$ 是域，并且事实上等于 $\kappa(\mathfrak q)$。
\end{proof}

\begin{lemma}
\label{algebra-lemma-dimension}
设 $k$ 为域，$V$ 为 $k$ 上的非零向量空间。假设 $V$ 的维数（这是一个
基数）小于 $k$ 的基数。则对任意线性算子 $T : V \to V$，都存在某个
首一多项式 $P(t) \in k[t]$，使 $P(T)$ 不可逆。
\end{lemma}

\begin{proof}
否则，$V$ 继承域 $k(t)$ 上的向量空间结构。但 $k(t)$ 在 $k$ 上的维数
至少是 $k$ 的基数；例如，这是因为诸元 $\frac{1}{t - \lambda}$ 在
$k$ 上线性无关。
\end{proof}

\noindent
下面是希尔伯特零点定理的另一版本。

\begin{theorem}
\label{algebra-theorem-uncountable-nullstellensatz}
设 $k$ 为域。令 $S$ 为由诸元 $\{x_i\}_{i \in I}$ 在 $k$ 上生成的
$k$-代数。假设 $I$ 的基数小于 $k$ 的基数。则
\begin{enumerate}
\item 对所有极大理想 $\mathfrak m \subset S$，域扩张
$\kappa(\mathfrak m)/k$ 是代数的；
\item $S$ 是雅可布森环。
\end{enumerate}
\end{theorem}

\begin{proof}
若 $I$ 有限，则结论由希尔伯特零点定理，即定理
\ref{algebra-theorem-nullstellensatz} 得出。以下假设 $I$ 无限。只须对以变量
$x_i$ 生成的多项式环中的极大理想
$\mathfrak m \subset k[\{x_i\}_{i \in I}]$ 证明结论，因为 $S$ 是它的
一个商。由于 $I$ 无限，所有单项式
$x_{i_1}^{e_1} \ldots x_{i_r}^{e_r}$ 的集合的基数至多等于 $I$ 的基数，
其中 $i_1, \ldots, i_r \in I$ 且 $e_1, \ldots, e_r \geq 0$。
这是因为 $I \times \ldots \times I$ 的基数等于 $I$ 的基数，并且
$\bigcup_{n \geq 0} I^n$ 的基数也相同。（若 $I$ 有限，这不成立；在该
情形，只有当 $k$ 不可数时，此证明才适用。）

\medskip\noindent
为得到矛盾，取 $T \in \kappa(\mathfrak m)$，使其在 $k$ 上超越。
注意，由乘以 $T$ 给出的 $k$-线性映射
$T : \kappa(\mathfrak m) \to \kappa(\mathfrak m)$ 具有如下性质：对所有
首一多项式 $P(t) \in k[t]$，$P(T)$ 都可逆。此外，
$\kappa(\mathfrak m)$ 在 $k$ 上的维数至多为 $I$ 的基数，因为它是
$k$ 上向量空间 $k[\{x_i\}_{i \in I}]$ 的商（如上所见，后者的维数为
$\# I$）。这与引理 \ref{algebra-lemma-dimension} 矛盾。

\medskip\noindent
为证明 $S$ 是雅可布森环，作如下论证。否则，由引理
\ref{algebra-lemma-characterize-jacobson}，存在素理想
$\mathfrak q \subset S$ 与元素 $f \in S$、$f \not \in \mathfrak q$，
使 $\mathfrak q$ 不是极大理想，而 $(S/\mathfrak q)_f$ 是域。但注意，
$(S/\mathfrak q)_f$ 可由至多 $\# I + 1$ 个元素生成。因此域扩张
$(S/\mathfrak q)_f/k$ 是代数的（由证明的第一部分）。这推出
$\kappa(\mathfrak q)$ 是 $k$ 的代数扩张，故由引理
\ref{algebra-lemma-finite-residue-extension-closed}，$\mathfrak q$ 是极大理想。
此矛盾完成证明。
\end{proof}

\begin{lemma}
\label{algebra-lemma-base-change-Jacobson}
设 $k$ 为域，$S$ 为 $k$-代数。对任意基数大于 $S$ 的基数的域扩张
$K/k$，有：
\begin{enumerate}
\item 对 $S_K$ 的每个极大理想 $\mathfrak m$，域
$\kappa(\mathfrak m)$ 在 $K$ 上是代数的；
\item $S_K$ 是雅可布森环。
\end{enumerate}
\end{lemma}

\begin{proof}
选取 $k \subset K$，使 $K$ 的基数大于 $S$ 的基数。由于 $S$ 的元素
生成 $K$-代数 $S_K$，可应用定理
\ref{algebra-theorem-uncountable-nullstellensatz}。
\end{proof}

\begin{example}
\label{algebra-example-countable-trick-does-not-work}
若 $k$ 是可数域且 $I$ 也可数，定理
\ref{algebra-theorem-uncountable-nullstellensatz} 证明中的技巧确实不适用。
设 $k$ 为可数域，$x$ 为变量，并令 $k(x)$ 为 $x$ 的有理函数域。考虑
多项式代数 $R = k[x, \{x_f\}_{f \in k[x]-\{0\}}]$。令
$I = (\{fx_f - 1\}_{f\in k[x] - \{0\}})$。注意，$I$ 是 $R$ 中的
真理想。选取极大理想 $I \subset \mathfrak m$。则
$k \subset R/\mathfrak m$ 同构于 $k(x)$，并且不在 $k$ 上代数。
\end{example}

\begin{lemma}
\label{algebra-lemma-Jacobson-invert-element}
设 $R$ 为雅可布森环，$f \in R$。环 $R_f$ 是雅可布森环，并且 $R_f$
的极大理想对应于 $R$ 的不含 $f$ 的极大理想。
\end{lemma}

\begin{proof}
由拓扑学，引理 \ref{topology-lemma-jacobson-inherited}，可知
$D(f) = \Spec(R_f)$ 是雅可布森空间，并且 $D(f)$ 的闭点对应于
$\Spec(R)$ 中恰好落在 $D(f)$ 内的闭点。因此由引理
\ref{algebra-lemma-jacobson}，$R_f$ 是雅可布森环。
\end{proof}

\begin{example}
\label{algebra-example-localize-not-preserve-closed-points}
下面这个简单例子表明，当 $R$ 不是雅可布森环时，引理
\ref{algebra-lemma-Jacobson-invert-element} 的结论不成立。考虑环
$R = \mathbf{Z}_{(2)}$，即 $\mathbf{Z}$ 在素理想 $(2)$ 处的局部化。
$R$ 在元素 $2$ 处的局部化同构于 $\mathbf{Q}$；写成公式即
$R_2 \cong \mathbf{Q}$。显然，映射 $R \to R_2$ 把
$\Spec(\mathbf{Q})$ 的闭点映到 $\Spec(R)$ 的泛点。
\end{example}

\begin{example}
\label{algebra-example-infinite-localize-not-preserve-closed-points}
下面这个简单例子表明，即使 $R$ 是雅可布森环，若对无限多个元素作局部化，
引理 \ref{algebra-lemma-Jacobson-invert-element} 的结论也不成立。令
$R = \mathbf{Z}$，并考虑 $R$ 在所有非零元素组成的集合处的局部化
$(R \setminus \{0\})^{-1}R \cong \mathbf{Q}$。显然，映射
$\mathbf{Z} \to \mathbf{Q}$ 把 $\Spec(\mathbf{Q})$ 的闭点映到
$\Spec(\mathbf{Z})$ 的泛点。
\end{example}

\begin{lemma}
\label{algebra-lemma-Jacobson-mod-ideal}
设 $R$ 为雅可布森环，$I \subset R$ 为理想。环 $R/I$ 是雅可布森环，
并且 $R/I$ 的极大理想对应于 $R$ 的包含 $I$ 的极大理想。
\end{lemma}

\begin{proof}
证明与引理 \ref{algebra-lemma-Jacobson-invert-element} 的证明相同。
\end{proof}

\begin{lemma}
\label{algebra-lemma-silly-jacobson}
设 $R$ 为雅可布森环，$K$ 为域。设 $R \subset K$，且 $K$ 在 $R$ 上
为有限型。则 $R$ 是域，并且 $K/R$ 是有限域扩张。
% P01-E0333: source joins “Let R subset K and K is ...” ungrammatically; recorded in ERRATA_CANDIDATES.jsonl.
\end{lemma}

\begin{proof}
首先注意，$R$ 是整环。由引理
\ref{algebra-lemma-field-finite-type-over-domain}，对某个非零 $f \in R$，
$R_f$ 是域且 $K/R_f$ 是有限域扩张。因此 $(0)$ 是 $R_f$ 的极大理想，
再由引理 \ref{algebra-lemma-Jacobson-invert-element} 可知，$(0)$ 是 $R$ 的
极大理想。
\end{proof}

\begin{proposition}
\label{algebra-proposition-Jacobson-permanence}
设 $R$ 为雅可布森环，$R \to S$ 为有限型环映射。则：
\begin{enumerate}
\item 环 $S$ 是雅可布森环；
\item 映射 $\Spec(S) \to \Spec(R)$ 把闭点映为闭点；
\item 对位于 $\mathfrak m \subset R$ 之上的极大理想
$\mathfrak m' \subset S$，域扩张
$\kappa(\mathfrak m')/\kappa(\mathfrak m)$ 是有限的。
\end{enumerate}
\end{proposition}

\begin{proof}
令 $\mathfrak m' \subset S$ 为极大理想，并令
$R \cap \mathfrak m' = \mathfrak m$。映射
$R/\mathfrak m \to S/\mathfrak m'$ 满足引理
\ref{algebra-lemma-silly-jacobson} 的条件；这是由引理
\ref{algebra-lemma-Jacobson-mod-ideal} 得到的。因此 $R/\mathfrak m$ 是域，
$\mathfrak m$ 是极大理想，并且诱导的剩余域扩张是有限的。这证明 (2)
与 (3)。

\medskip\noindent
若 $S$ 不是雅可布森环，则由引理 \ref{algebra-lemma-characterize-jacobson}，
存在 $S$ 的非极大素理想 $\mathfrak q$ 及
$g \in S$、$g \not\in \mathfrak q$，使 $(S/\mathfrak q)_g$ 是域。
为得到矛盾，证明 $\mathfrak q$ 是极大理想。令
$\mathfrak p = \mathfrak q \cap R$。映射
$R/\mathfrak p \to (S/\mathfrak q)_g$ 满足引理
\ref{algebra-lemma-silly-jacobson} 的条件；这是由引理
\ref{algebra-lemma-Jacobson-mod-ideal} 得到的。因此 $R/\mathfrak p$ 是域，且域扩张
$\kappa(\mathfrak p) \to (S/\mathfrak q)_g = \kappa(\mathfrak q)$
是有限的，从而是代数的。于是由引理
\ref{algebra-lemma-finite-residue-extension-closed}，$\mathfrak q$ 是 $S$ 的
极大理想。矛盾。
\end{proof}

\begin{lemma}
\label{algebra-lemma-corollary-jacobson}
$\mathbf{Z}$ 上的任意有限型代数都是雅可布森环。
\end{lemma}

\begin{proof}
合用引理 \ref{algebra-lemma-pid-jacobson} 与命题
\ref{algebra-proposition-Jacobson-permanence}。
\end{proof}

\begin{lemma}
\label{algebra-lemma-image-finite-type-map-Jacobson-rings}
设 $R \to S$ 为雅可布森环之间的有限型环映射。记
$X = \Spec(R)$、$Y = \Spec(S)$，并将诱导的谱映射记为
$f : Y \to X$。令 $E \subset Y = \Spec(S)$ 为可构造集。用下标
${}_0$ 表示拓扑空间的闭点集。
\begin{enumerate}
\item 有 $f(E)_0 = f(E_0) = X_0 \cap f(E)$。
\item 点 $\xi \in X$ 属于 $f(E)$，当且仅当
$\overline{\{\xi\}} \cap f(E_0)$ 在 $\overline{\{\xi\}}$ 中稠密。
\end{enumerate}
\end{lemma}

\begin{proof}
有连续映射的交换图
$$
\xymatrix{
E \ar[r] \ar[d] & Y \ar[d] \\
f(E) \ar[r] & X
}
$$
设 $x \in f(E)$ 在 $f(E)$ 中是闭的。则
$f^{-1}(\{x\})\cap E$ 非空且在 $E$ 中闭。对两个包含关系
$$
f^{-1}(\{x\}) \cap E \subset E \subset Y
$$
应用拓扑学，引理 \ref{topology-lemma-jacobson-inherited}，可找到点
$y \in f^{-1}(\{x\}) \cap E$，它在 $Y$ 中是闭的。换言之，存在
$y \in Y_0$ 且 $y \in E_0$，并映到 $x$。因此 $x \in f(E_0)$。
这证明 $f(E)_0 \subset f(E_0)$。命题
\ref{algebra-proposition-Jacobson-permanence} 蕴含
$f(E_0) \subset X_0 \cap f(E)$。包含关系
$X_0 \cap f(E) \subset f(E)_0$ 显然。这证明第一个断言。

\medskip\noindent
设 $\xi \in f(E)$。由引理 \ref{algebra-lemma-characterize-image-finite-type}，
集合 $f(E) \cap \overline{\{\xi\}}$ 包含
$\overline{\{\xi\}}$ 的一个稠密开子集。由于 $X$ 是雅可布森空间，
可知 $f(E) \cap \overline{\{\xi\}}$ 含有一个稠密闭点集；见拓扑学，
引理 \ref{topology-lemma-jacobson-inherited}。由本引理的 (1) 得出结论。

\medskip\noindent
另一方面，设 $\overline{\{\xi\}} \cap f(E_0)$ 在
$\overline{\{\xi\}}$ 中稠密。由
引理 \ref{algebra-lemma-constructible-is-image}，存在有限表示环映射
$S \to S'$，使 $E$ 是 $Y' := \Spec(S') \to Y$ 的像。对环映射
$S \to S'$ 应用本引理的第一部分，$E_0$ 是 $Y'_0$ 的像。因此以
$S'$ 代替 $S$，可以假设 $E = Y$。设 $\xi$ 对应于
$\mathfrak p \subset R$。考虑图表
$$
\xymatrix{
S \ar[r] & S/\mathfrak p S \\
R \ar[r] \ar[u] & R/\mathfrak p \ar[u]
}
$$
该图表以及 $f(Y_0) \cap V(\mathfrak p)$ 在 $V(\mathfrak p)$ 中的
稠密性表明，态射 $R/\mathfrak p \to S/\mathfrak p S$ 满足引理
\ref{algebra-lemma-domain-image-dense-set-points-generic-point} 的条件 (2)。
因此存在映到 $(0)$ 的素理想
$\overline{\mathfrak q} \subset S/\mathfrak pS$。换言之，
$\overline{\mathfrak q}$ 在 $S$ 中的逆像 $\mathfrak q$ 按所需映到
$\mathfrak p$。
\end{proof}

\noindent
上述引理的结论是：可以从像的闭点集读出 $f$ 的像。当映射是有限表示时，
这一点稍为更好，因为此时知道可构造集的像仍为可构造集。在陈述之前，先引入
一些记号。以 $\text{Constr}(X)$ 表示可构造集的集合。设
$R \to S$ 为环映射。记 $X = \Spec(R)$、$Y = \Spec(S)$，并将诱导的
谱映射记为 $f : Y \to X$。用下标 ${}_0$ 表示拓扑空间的闭点集。


\begin{lemma}
\label{algebra-lemma-conclude-jacobson-Noetherian}
沿用上述记号。假设 $R$ 是诺特雅可布森环，并进一步假设 $R \to S$
为有限型。存在交换图
$$
\xymatrix{
\text{Constr}(Y) \ar[r]^{E \mapsto E_0} \ar[d]^{E \mapsto f(E)} &
\text{Constr}(Y_0) \ar[d]^{E \mapsto f(E)} \\
\text{Constr}(X) \ar[r]^{E \mapsto E_0} &
\text{Constr}(X_0)
}
$$
其中水平箭头是拓扑学，引理
\ref{topology-lemma-jacobson-equivalent-constructible} 中的双射。
\end{lemma}

\begin{proof}
由于 $R \to S$ 为有限型，由引理
\ref{algebra-lemma-Noetherian-finite-type-is-finite-presentation}，它是有限表示的。
因此由切瓦莱定理（定理 \ref{algebra-theorem-chevalley}），$X$ 中可构造集的像
在 $Y$ 中是可构造的。结合引理
\ref{algebra-lemma-image-finite-type-map-Jacobson-rings} 即得本引理。
% P01-E0334: source reverses the domain/codomain of the constructible-image statement; recorded in ERRATA_CANDIDATES.jsonl.
\end{proof}

\noindent
为说明雅可布森环的用法，给出以下两个例子。

\begin{example}
\label{algebra-example-product-matrices-zero}
设 $k$ 为域。空间 $\Spec(k[x, y]/(xy))$ 有两个不可约分支，即 $x$ 轴
与 $y$ 轴。作为推广，令
$$
R = k[x_{11}, x_{12}, x_{21}, x_{22}, y_{11}, y_{12}, y_{21}, y_{22}]/
\mathfrak a,
$$
其中 $\mathfrak a$ 是
$k[x_{11}, x_{12}, x_{21}, x_{22}, y_{11}, y_{12}, y_{21}, y_{22}]$
中的理想，由下列 $2 \times 2$ 乘积矩阵的各个条目生成：
$$
\left(
\begin{matrix}
x_{11} & x_{12}\\
x_{21} & x_{22}
\end{matrix}
\right)
 \left(
\begin{matrix}
y_{11} & y_{12}\\
y_{21} & y_{22}
\end{matrix}
\right).
$$
本例将描述 $\Spec(R)$。

\medskip\noindent
为证明关于 $\Spec(k[x, y]/(xy))$ 的断言，论证如下。若
$\mathfrak p \subset k[x, y]$ 是任意包含 $xy$ 的理想，则 $x$ 或 $y$
应属于 $\mathfrak p$。因此这样的极小素理想恰为 $(x)$ 与 $(y)$。
% P01-E0335: source says arbitrary ideal, but the product implication requires a prime ideal; recorded in ERRATA_CANDIDATES.jsonl.
当 $k$ 代数闭时，这些分支的 $\text{max-Spec}$ 可分别视为 $y$ 轴与
$x$ 轴的点集。

\medskip\noindent
对上述推广，注意，至少在 $k$ 代数闭时，可以把
$k[x_{11}, x_{12}, x_{21}, x_{22}, y_{11}, y_{12}, y_{21}, y_{22}])$
的谱的闭点与如下矩阵空间等同：
% P01-E0336: source polynomial-ring expression has an unmatched closing parenthesis; recorded in ERRATA_CANDIDATES.jsonl.
$$
\left\{ (X, Y) \in \text{Mat}(2, k)\times \text{Mat}(2, k) \mid
X = \left(
\begin{matrix}
x_{11} & x_{12}\\
x_{21} & x_{22}
\end{matrix}
\right),
Y= \left(
\begin{matrix}
y_{11} & y_{12}\\
y_{21} & y_{22}
\end{matrix}
\right)
\right\}
$$
现在定义
$\text{GL}(2, k)\times \text{GL}(2, k)\times \text{GL}(2, k)$
在矩阵空间 $\{(X, Y)\}$ 上的群作用：
$$
(g_1, g_2, g_3) \times (X, Y) \mapsto ((g_1Xg_2^{-1}, g_2Yg_3^{-1})).
$$
这里还要注意，代数集
$$
\text{GL}(2, k)\times \text{GL}(2, k)\times \text{GL}(2, k) \subset
\text{Mat}(2, k)\times \text{Mat}(2, k) \times \text{Mat}(2, k)
$$
不可约，因为它是下列整环的极大谱：
$$
k[x_{11}, x_{12}, \ldots, z_{21}, z_{22}, (x_{11}x_{22}-x_{12}x_{21})^{-1}
, (y_{11}y_{22}-y_{12}y_{21})^{-1}, (z_{11}z_{22}-z_{12}z_{21})^{-1}].
$$
由于不可约代数集的像仍不可约，只须分类集合
$\{(X, Y)\in \text{Mat}(2, k)\times \text{Mat}(2, k)|XY = 0\}$ 的轨道，
再取其闭包。由标准线性代数，归结为以下三种情形：
\begin{enumerate}
\item $\exists (g_1, g_2)$，使 $g_1Xg_2^{-1} = I_{2\times 2}$。
此时 $Y$ 必为 $0$，而这作为代数集在群作用下不变。因此，此轨道包含于
由素理想 $(y_{11}, y_{12}, y_{21}, y_{22})$ 定义的不可约代数集中。
取闭包可知，$(y_{11}, y_{12}, y_{21}, y_{22})$ 确实是一个分支。
\item $\exists (g_1, g_2)$，使
$$
g_1Xg_2^{-1} = \left(
\begin{matrix}
1 & 0 \\
0 & 0
\end{matrix}
\right).
$$
这一情形当且仅当 $X$ 为秩 1 的矩阵；进一步，这样的 $X$ 消去 $Y$
当且仅当
$$
x_{11}y_{11}+x_{12}y_{21} = 0; \quad x_{11}y_{12}+x_{12}y_{22} = 0;
$$
$$
x_{21}y_{11}+x_{22}y_{21} = 0; \quad x_{21}y_{12}+x_{22}y_{22} = 0.
$$
固定一个秩 1 的 $X$。满足上述方程的非零 $Y$ 构成不可约代数集，
理由如下（$Y = 0$ 属于前一种情形）：由
$0 = g_1Xg_2^{-1}g_2Y$ 可知
$$
g_2Y = \left(
\begin{matrix}
0 & 0 \\
y_{21}' & y_{22}'
\end{matrix}
\right).
$$
再在右侧施以 $g_3$ 给出的 $\text{GL}(2, k)$ 作用，可将 $g_2Y$ 化为
$$
g_2Yg_3^{-1} = \left(
\begin{matrix}
0 & 0 \\
0 & 1
\end{matrix}
\right),
$$
所以这样的 $Y$ 构成不可约代数集，同构于 $\text{GL}(2, k)$ 在此作用下
的像。最后注意，$X$ 的“秩 1”条件构成不可约代数集
$\det X = x_{11}x_{22} - x_{12}x_{21} = 0$ 的稠密开子集。于是可知，
全部五个方程在非零矩阵对空间的开子集中定义一个不可约分支
$(x_{11}y_{11}+x_{12}y_{21}, x_{11}y_{12}+x_{12}y_{22}, x_{21}y_{11}
+x_{22}y_{21}, x_{21}y_{12}+x_{22}y_{22}, x_{11}x_{22}-x_{12}x_{21})$。
可以证明，方程对 $\det X = 0$、$\det Y = 0$ 在 $\Spec(R)$ 中截出
一个不可约分支，且上述轨迹是它的稠密开子集。
\item $\exists (g_1, g_2)$，使 $g_1Xg_2^{-1} = 0$；等价地，$X = 0$。
此时 $Y$ 可以任取，因而该分支由
$(x_{11}, x_{12}, x_{21}, x_{22})$ 定义。
\end{enumerate}
\end{example}


\begin{example}
\label{algebra-example-idempotent-matrices}
再举一例，考虑 $R = k[\{t_{ij}\}_{i, j = 1}^{n}]/\mathfrak a$，其中
$\mathfrak a$ 是由乘积矩阵 $T^2-T$ 的各条目生成的理想，
$T = (t_{ij})$。由线性代数可知，在由
$g, T \mapsto gTg^{-1}$ 定义的 $GL(n, k)$-作用下，$T$ 按其秩分类，
且每个 $T$ 都共轭于某个
$\text{diag}(1, \ldots, 1, 0, \ldots, 0)$，后者有 $r$ 个 1 与
$n-r$ 个 0。因此，这种 $\text{diag}(1, \ldots, 1, 0, \ldots, 0)$
在群作用下的每个轨道都构成一个不可约分支，而每个幂等矩阵都包含在某个
这样的轨道中。接下来证明，任意两个不同轨道必不相交。为此只需构造在不同
轨道上取不同值的多项式函数。在特征为 0 的情形，可以取
$f(t_{ij}) = trace(T) = \sum_{i = 1}^nt_{ii}$。在正特征情形，情况稍微
棘手，因为即使 $T \neq 0$，也可能有 $trace(T) = 0$。例如，当
$char = 3$ 时，
% P01-E0337: source has “classified by the its rank”; recorded in ERRATA_CANDIDATES.jsonl.
$$
trace\left(
\begin{matrix}
1 & & \\
& 1 & \\
& & 1
\end{matrix}
\right) = 3 = 0
$$
无论如何，可以用其他函数分离这些分支。例如在特征 3 的情形，
$tr(\wedge^3T)$ 在对应于 $diag(1, 1, 1)$ 的分支上取值 1，而在其他
分支上取值 0。
% P01-E0338: the last claim fails for general n (e.g. rank 4 in characteristic 3); recorded in ERRATA_CANDIDATES.jsonl.
\end{example}

\section{有限环扩张与整环扩张}
\label{algebra-section-finite-ring-extensions}

\noindent
关于有限环映射与整环映射的平凡引理。先回顾定义。

\begin{definition}
\label{algebra-definition-integral-ring-map}
设 $\varphi : R \to S$ 为环映射。
\begin{enumerate}
\item 元素 $s \in S$ 称为{\it 整于 $R$}，如果存在首一多项式
$P(x) \in R[x]$，使得 $P^\varphi(s) = 0$；其中
$P^\varphi(x) \in S[x]$ 是 $P$ 在
$\varphi : R[x] \to S[x]$ 下的像。
\item  若每个 $s \in S$ 都整于 $R$，则称环映射 $\varphi$ 是{\it 整的}。
\end{enumerate}
\end{definition}

\begin{lemma}
\label{algebra-lemma-characterize-integral-element}
设 $\varphi : R \to S$ 为环映射，$y \in S$。若存在 $S$ 的有限生成
$R$-子模 $M$，使 $1 \in M$ 且 $yM \subset M$，则 $y$ 整于 $R$。
\end{lemma}

\begin{proof}
考虑映射 $\varphi : M \to M$，$x \mapsto y \cdot x$。由引理
\ref{algebra-lemma-charpoly-module}，存在首一多项式 $P \in R[T]$，使
$P(\varphi) = 0$。在环 $S$ 中有
$P(y) = P(y) \cdot 1 = P(\varphi)(1) = 0$。
\end{proof}

\begin{lemma}
\label{algebra-lemma-finite-is-integral}
有限环映射是整的。
\end{lemma}

\begin{proof}
设 $R \to S$ 为有限环映射，$y \in S$。对 $M = S$ 应用引理
\ref{algebra-lemma-characterize-integral-element}，可知 $y$ 整于 $R$。
\end{proof}

\begin{lemma}
\label{algebra-lemma-characterize-integral}
设 $\varphi : R \to S$ 为环映射，且 $s_1, \ldots, s_n$ 是 $S$ 的有限个
元素。则下列条件等价：所有 $s_i$（$i = 1, \ldots, n$）都整于 $R$；
存在一个有限于 $R$ 的 $R$-子代数 $S' \subset S$，包含所有 $s_i$。
\end{lemma}

\begin{proof}
若每个 $s_i$ 都整，则由 $\varphi(R)$ 和各 $s_i$ 生成的子代数
在 $R$ 上有限。事实上，若 $s_i$ 在 $R$ 上满足次数为 $d_i$ 的首一方程，
则该子代数作为 $R$-模由满足 $0 \leq e_i \leq d_i - 1$ 的元素
$s_1^{e_1} \ldots s_n^{e_n}$ 生成。反之，若给定一个包含所有 $s_i$、
且为有限的 $R$-子代数 $S'$，则由引理
\ref{algebra-lemma-finite-is-integral}，所有 $s_i$ 都整。
\end{proof}

\begin{lemma}
\label{algebra-lemma-characterize-finite-in-terms-of-integral}
设 $R \to S$ 为环映射。以下条件等价：
\begin{enumerate}
\item $R \to S$ 是有限的；
\item $R \to S$ 是整的且为有限型的；
\item 存在 $x_1, \ldots, x_n \in S$，它们作为 $R$-代数生成 $S$，
且每个 $x_i$ 都整于 $R$。
\end{enumerate}
\end{lemma}

\begin{proof}
由引理 \ref{algebra-lemma-characterize-integral} 立即得到。
\end{proof}

\begin{lemma}
\label{algebra-lemma-integral-transitive}
\begin{slogan}
整环映射的复合仍为整环映射
\end{slogan}
设 $R \to S$ 与 $S \to T$ 是整环映射，则 $R \to T$ 是整环映射。
\end{lemma}

\begin{proof}
取 $t \in T$。令 $P(x) \in S[x]$ 为满足 $P(t) = 0$ 的首一多项式。
对 $P$ 的有限个系数应用引理
\ref{algebra-lemma-characterize-integral}。于是 $t$ 整于某个在 $R$ 上有限的
子代数 $S' \subset S$。再次应用引理
\ref{algebra-lemma-characterize-integral}，得到一个在 $S'$ 上有限且包含 $t$ 的
子代数 $T' \subset T$。对 $R \to S' \to T'$ 应用引理
\ref{algebra-lemma-finite-transitive}，可知 $T'$ 在 $R$ 上有限。于是由引理
\ref{algebra-lemma-finite-is-integral}，$t$ 整于 $R$。
\end{proof}

\begin{lemma}
\label{algebra-lemma-integral-closure-is-ring}
设 $R \to S$ 为环同态。集合
$$
S' = \{s \in S \mid s\text{ is integral over }R\}
$$
是 $S$ 的一个 $R$-子代数。
\end{lemma}

\begin{proof}
由引理 \ref{algebra-lemma-characterize-integral} 与
\ref{algebra-lemma-finite-is-integral}，结论显然。
\end{proof}

\begin{lemma}
\label{algebra-lemma-finite-product-integral}
设 $R_i\to S_i$（$i = 1, \ldots, n$）为环映射。令 $R$ 与 $S$ 分别为
所有 $R_i$ 与所有 $S_i$ 的乘积。则元素
$s = (s_1, \ldots, s_n) \in S$ 整于 $R$，当且仅当每个 $s_i$ 整于
$R_i$。
\end{lemma}

\begin{proof}
略。
\end{proof}

\begin{definition}
\label{algebra-definition-integral-closure}
设 $R \to S$ 为环映射。S 中整于 $R$ 的元素所成的环
$S' \subset S$（见引理 \ref{algebra-lemma-integral-closure-is-ring}）称为
$R$ 在 $S$ 中的{\it 整闭包}。若 $R \subset S$，且 $R = S'$，则称
$R$ 在 $S$ 中{\it 整闭}。
\end{definition}

\noindent
特别地，$R \to S$ 是整的，当且仅当 $R$ 在 $S$ 中的整闭包就是 $S$。

\begin{lemma}
\label{algebra-lemma-finite-product-integral-closure}
设 $R_i\to S_i$（$i = 1, \ldots, n$）为环映射。记 $R_i$ 在 $S_i$ 中的
整闭包为 $S'_i$。进一步令 $R$ 与 $S$ 分别为所有 $R_i$ 与所有 $S_i$
的乘积。则 $R$ 在 $S$ 中的整闭包是各 $S'_i$ 的乘积。特别地，
$R \to S$ 整闭，当且仅当每个 $R_i \to S_i$ 都整闭。
\end{lemma}

\begin{proof}
由引理 \ref{algebra-lemma-finite-product-integral} 立即得到。
\end{proof}

\begin{lemma}
\label{algebra-lemma-integral-closure-localize}
整闭包与局部化交换：若 $A \to B$ 为环映射，且 $S \subset A$ 为乘法子集，
则 $S^{-1}A$ 在 $S^{-1}B$ 中的整闭包是 $S^{-1}B'$，其中
$B' \subset B$ 是 $A$ 在 $B$ 中的整闭包。
\end{lemma}

\begin{proof}
由于局部化是正合的，有 $S^{-1}B' \subset S^{-1}B$。设 $x \in B'$ 且
$f \in S$。于是对某些 $a_i \in A$，在 $B$ 中有
$$
x^d + \sum_{i = 1, \ldots, d} a_i x^{d - i} = 0
$$
因此在 $S^{-1}B$ 中也有
$$
(x/f)^d + \sum\nolimits_{i = 1, \ldots, d} a_i/f^i (x/f)^{d - i} = 0
$$
由此可知 $S^{-1}B'$ 包含于 $S^{-1}A$ 在 $S^{-1}B$ 中的整闭包。
反之，设 $x/f \in S^{-1}B$ 整于 $S^{-1}A$。于是对某些
$a_i \in A$ 与 $f_i \in S$，在 $S^{-1}B$ 中有
$$
(x/f)^d + \sum\nolimits_{i = 1, \ldots, d} (a_i/f_i) (x/f)^{d - i} = 0
$$
这意味着，对某个适当的 $f' \in S$，有
$$
(f'f_1 \ldots f_d x)^d +
\sum\nolimits_{i = 1, \ldots, d}
f^i(f')^if_1^i \ldots f_i^{i - 1} \ldots f_d^i a_i
(f'f_1 \ldots f_dx)^{d - i} = 0
$$
因此 $f'f_1\ldots f_dx \in B'$，从而按所需
$x/f \in S^{-1}B'$。
\end{proof}

\begin{lemma}
\label{algebra-lemma-integral-closure-stalks}
\begin{slogan}
环上代数的元素整于该环，当且仅当它在该环的每个素理想处都局部整。
\end{slogan}
设 $\varphi : R \to S$ 为环映射，$x \in S$。以下条件等价：
\begin{enumerate}
\item $x$ 整于 $R$；
\item 对每个素理想 $\mathfrak p \subset R$，元素
$x \in S_{\mathfrak p}$ 整于 $R_{\mathfrak p}$。
\end{enumerate}
\end{lemma}

\begin{proof}
(1) 蕴含 (2) 显然。设 (2) 成立。考虑由 $\varphi(R)$ 与 $x$ 生成的
$R$-代数 $S' \subset S$。令 $\mathfrak p$ 为 $R$ 的一个素理想。则在
$S_{\mathfrak p}$ 中，对某些 $a_i \in R_{\mathfrak p}$ 有
$$
x^d + \sum_{i = 1, \ldots, d} \varphi(a_i) x^{d - i} = 0
$$
考察其中出现的分母可知，存在 $f \in R$、$f \not \in \mathfrak p$，使
$a_i \in R_f$ 且在 $S_f$ 中有
$$
x^d + \sum_{i = 1, \ldots, d} \varphi(a_i) x^{d - i} = 0
$$
这说明 $S'_f$ 在 $R_f$ 上有限。由于 $\mathfrak p$ 任意，且
$\Spec(R)$ 拟紧（引理 \ref{algebra-lemma-quasi-compact}），可找到有限多个
$f_1, \ldots, f_n \in R$，它们生成 $R$ 的单位理想，并使每个
$S'_{f_i}$ 在 $R_{f_i}$ 上有限。因此由引理 \ref{algebra-lemma-cover}，
$S'$ 在 $R$ 上有限。于是由引理 \ref{algebra-lemma-characterize-integral}，
$x$ 整于 $R$。
\end{proof}

\begin{lemma}
\label{algebra-lemma-base-change-integral}
\begin{slogan}
整性与有限性在基变换下保持
\end{slogan}
设 $R \to S$ 与 $R \to R'$ 为环映射，令 $S' = R' \otimes_R S$。
\begin{enumerate}
\item 若 $R \to S$ 整，则 $R' \to S'$ 整。
\item 若 $R \to S$ 有限，则 $R' \to S'$ 有限。
\end{enumerate}
\end{lemma}

\begin{proof}
证明 (1)。令 $s_i \in S$ 为生成 $S$ 的一组元素。这些元素各自都满足
一个在 $R$ 上的首一多项式方程 $P_i$，其系数属于 $R$。因此元素
$1 \otimes s_i \in S'$ 在 $R'$ 上生成 $S'$，并满足相应的
$R'$ 上的多项式 $P_i'$。由于这些元素在 $R'$ 上生成 $S'$，可知 $S'$ 整于 $R'$。
证明 (2) 略。
\end{proof}

\begin{lemma}
\label{algebra-lemma-integral-local}
设 $R \to S$ 为环映射，且 $f_1, \ldots, f_n \in R$ 生成单位理想。
\begin{enumerate}
\item 若每个 $R_{f_i} \to S_{f_i}$ 都整，则 $R \to S$ 整。
\item 若每个 $R_{f_i} \to S_{f_i}$ 都有限，则 $R \to S$ 有限。
\end{enumerate}
\end{lemma}

\begin{proof}
证明 (1)。

设 $s \in S$。考虑 R[x] 中满足 $P(s) = 0$ 的多项式 $P$ 所成的理想
$I \subset R[x]$。令 $J \subset R$ 表示 $I$ 中元素的首项系数所成的
理想（原文特意标注为 (!)）。由假设并消去分母可知，对所有 $i$，某些
$n_i \geq 0$，所有 $f_i^{n_i} \in J$。因此 $J$ 包含 $1$，从而 $s$
整于 $R$。证明 (2) 略。
\end{proof}

\begin{lemma}
\label{algebra-lemma-integral-permanence}
设 $A \to B \to C$ 为环映射。
\begin{enumerate}
\item 若 $A \to C$ 是整的，则 $B \to C$ 也是整的。
\item 若 $A \to C$ 是有限的，则 $B \to C$ 也是有限的。
\end{enumerate}
\end{lemma}

\begin{proof}
略。
\end{proof}

\begin{lemma}
\label{algebra-lemma-integral-closure-transitive}
设 $A \to B \to C$ 为环映射。令 $B'$ 为 $A$ 在 $B$ 中的整闭包，令
$C'$ 为 $B'$ 在 $C$ 中的整闭包。则 $C'$ 是 $A$ 在 $C$ 中的整闭包。
\end{lemma}

\begin{proof}
略。
\end{proof}

\begin{lemma}
\label{algebra-lemma-integral-overring-surjective}
设 $R \to S$ 为整环扩张且 $R \subset S$。则
$\varphi : \Spec(S) \to \Spec(R)$ 是满射。
\end{lemma}

\begin{proof}
令 $\mathfrak p \subset R$ 为素理想。由引理
\ref{algebra-lemma-in-image}，须证明 $\mathfrak pS_{\mathfrak p} \not = S_{\mathfrak p}$。
局部化 $R_{\mathfrak p} \to S_{\mathfrak p}$ 是单射（因为局部化正合），
并且由引理 \ref{algebra-lemma-integral-closure-localize} 或
\ref{algebra-lemma-base-change-integral} 是整的。因此可用
$R_{\mathfrak p}$、$S_{\mathfrak p}$ 分别代替 $R$、$S$，并设 $R$ 是
带极大理想 $\mathfrak m$ 的局部环；只须证明 $\mathfrak mS \not = S$。
假设 $1 = \sum f_i s_i$，其中 $f_i \in \mathfrak m$、$s_i \in S$，以导出
矛盾。取 $R \subset S' \subset S$，使 $R \to S'$ 有限且
$s_i \in S'$；见引理 \ref{algebra-lemma-characterize-integral}。等式
$1 = \sum f_i s_i$ 意味着有限 $R$-模 $S'$ 满足 $S' = \mathfrak m S'$。
由中山引理 \ref{algebra-lemma-NAK}，得到 $S' = 0$，矛盾。
\end{proof}

\begin{lemma}
\label{algebra-lemma-integral-under-field}
设 $R$ 为环，$K$ 为域。若 $R \subset K$ 且 $K$ 整于 $R$，则 $R$ 是域，
且 $K$ 是代数扩张。若 $R \subset K$ 且 $K$ 在 $R$ 上有限，则 $R$ 是域，
且 $K$ 是有限代数扩张。
\end{lemma}

\begin{proof}
假设 $R \subset K$ 是整扩张。由引理
\ref{algebra-lemma-integral-overring-surjective}，$\Spec(R)$ 只有 $1$ 个点。
显然 $R$ 是整环，故 $R = R_{(0)}$ 是域（引理
\ref{algebra-lemma-minimal-prime-reduced-ring}）。其余断言立即得到。
\end{proof}

\begin{lemma}
\label{algebra-lemma-integral-over-field}
设 $k$ 为域，$S$ 为 $k$ 上的 $k$-代数。
\begin{enumerate}
\item 若 $S$ 是整环且在 $k$ 上有限维，则 $S$ 是域。
\item 若 $S$ 整于 $k$ 且为整环，则 $S$ 是域。
\item 若 $S$ 整于 $k$，则 $S$ 的每个素理想都是极大理想（更多结论见
引理 \ref{algebra-lemma-ring-with-only-minimal-primes}）。
\end{enumerate}
\end{lemma}

\begin{proof}
关于素理想的断言来自“integral $+$ domain $\Rightarrow$ field”这一断言。设 $S$ 整于 $k$
且 $S$ 为整环。取 $s \in S$。由引理
\ref{algebra-lemma-characterize-integral}，可找到包含 $s$ 的有限维
$k$-子代数 $k \subset S' \subset S$。若证明第一条断言，则可知 $S$ 是域。
现在假设 $S$ 在 $k$ 上有限维且为整环。取 $s\in S$。由于 $S$ 是整环，
乘法映射 $s : S \to S$ 由维数相等可知是满射。因此存在 $s_1 \in S$，
使 $ss_1 = 1$。所以 $S$ 是域。
\end{proof}

\begin{lemma}
\label{algebra-lemma-integral-no-inclusion}
设 $R \to S$ 为整环映射。令
$\mathfrak q, \mathfrak q' \in \Spec(S)$ 为不同的素理想，且它们在
$\Spec(R)$ 中的像相同。则既不存在 $\mathfrak q \subset \mathfrak q'$，
也不存在 $\mathfrak q' \subset \mathfrak q$。
\end{lemma}

\begin{proof}
令 $\mathfrak p \subset R$ 为像。由注记
\ref{algebra-remark-fundamental-diagram}，素理想 $\mathfrak q, \mathfrak q'$
对应于 $S \otimes_R \kappa(\mathfrak p)$ 中的理想。因此由引理
\ref{algebra-lemma-integral-over-field} 得到结论。
\end{proof}

\begin{lemma}
\label{algebra-lemma-finite-finite-fibres}
设 $R \to S$ 为有限映射。则
$\Spec(S) \to \Spec(R)$ 的纤维是有限的。
\end{lemma}

\begin{proof}
由注记 \ref{algebra-remark-fundamental-diagram} 的讨论，纤维是环
$S \otimes_R \kappa(\mathfrak p)$ 的谱。由于 $R \to S$ 有限，这些纤维环
在 $\kappa(\mathfrak p)$ 上有限，因而由引理
\ref{algebra-lemma-Noetherian-permanence} 是诺特的。由引理
\ref{algebra-lemma-integral-no-inclusion}，$S \otimes_R \kappa(\mathfrak p)$
的每个素理想都是极小素理想。因此由引理
\ref{algebra-lemma-Noetherian-irreducible-components}，这样的素理想至多有有限多个。
\end{proof}

\begin{lemma}
\label{algebra-lemma-integral-going-up}
设 $R \to S$ 为环映射，且 $S$ 整于 $R$。令
$\mathfrak p \subset \mathfrak p' \subset R$ 为素理想。设
$\mathfrak q$ 是 $S$ 中映到 $\mathfrak p$ 的素理想。则存在素理想
$\mathfrak q'$，满足 $\mathfrak q \subset \mathfrak q'$，并且映到
$\mathfrak p'$。
\end{lemma}

\begin{proof}
可以用 $R$ 替换为 $R/\mathfrak p$，用 $S$ 替换为 $S/\mathfrak q$。
于是归约为整环扩张 $R \subset S$ 以及素理想
$\mathfrak p' \subset R$ 的情形。由引理
\ref{algebra-lemma-integral-overring-surjective} 即得。
\end{proof}

\noindent
上述引理所表达的性质称为环映射 $R \to S$ 的“上升性质”（going up
property），见定义 \ref{algebra-definition-going-up-down}。

\begin{lemma}
\label{algebra-lemma-finite-finitely-presented-extension}
设 $R \to S$ 为有限且有限表示的环映射，$M$ 为 $S$-模。则 $M$ 作为
$R$-模有限表示，当且仅当 $M$ 作为 $S$-模有限表示。
\end{lemma}

\begin{proof}
其中一个方向由引理 \ref{algebra-lemma-finitely-presented-over-subring} 得出。
反之，设 $M$ 作为 $S$-模有限表示。取表示
$$
S^{\oplus m} \longrightarrow
S^{\oplus n} \longrightarrow
M \longrightarrow 0
$$
由于 $S$ 作为 $R$-模有限，$S^{\oplus n} \to M$ 的核是有限
$R$-模。因此由引理 \ref{algebra-lemma-extension}，只须证明 $S$ 作为
$R$-模是有限表示的。

\medskip\noindent
取 $y_1, \ldots, y_n \in S$，使 $y_1, \ldots, y_n$ 作为 $R$-模生成 $S$。由引理
\ref{algebra-lemma-characterize-integral-element}，每个 $y_i$ 都整于 $R$。
选取首一多项式 $P_i(x) \in R[x]$，使 $P_i(y_i) = 0$。考虑环
$$
S' = R[x_1, \ldots, x_n]/(P_1(x_1), \ldots, P_n(x_n))
$$
由引理 \ref{algebra-lemma-compose-finite-type}，$S$ 作为 $S'$-代数是有限表示的。
由于 $S' \to S$ 是满射，由引理
\ref{algebra-lemma-finite-presentation-independent}，核
$J = \Ker(S' \to S)$ 作为理想是有限生成的。因此 $J$ 是有限
$S'$-模（直接由定义得到）。于是 $S = \Coker(J \to S')$ 作为
$S'$-模是有限表示的，这是引理 \ref{algebra-lemma-extension} 的结论。
因此，如第一段所述，只须证明 $S'$ 作为 $R$-模是有限表示的。事实上，
$S'$ 是自由 $R$-模，其基为满足 $0 \leq e_i < \deg(P_i)$ 的单项式
$x_1^{e_1} \ldots x_n^{e_n}$。具体地，将 $R \to S'$ 写成复合
$$
R \to R[x_1]/(P_1(x_1)) \to
R[x_1, x_2]/(P_1(x_1), P_2(x_2)) \to
\ldots \to S'
$$
这表明，此序列中的第 $i$ 个环作为第 $(i - 1)$ 个环的模是自由的，其基为
$1, x_i, \ldots, x_i^{\deg(P_i) - 1}$。由归纳法容易得到结论。细节略。
\end{proof}

\begin{lemma}
\label{algebra-lemma-silly-normal}
设 $R$ 为环，$x, y \in R$ 为非零因子。令
$R[x/y] \subset R_{xy}$ 为由 $x/y$ 生成的 $R$-子代数，并类似地定义
$R[y/x]$ 与 $R[x/y, y/x]$。若 $R$ 在 $R_x$ 或 $R_y$ 中整闭，则序列
$$
0 \to R \xrightarrow{(-1, 1)} R[x/y] \oplus R[y/x] \xrightarrow{(1, 1)}
R[x/y, y/x] \to 0
$$
是 $R$-模的短正合序列。
\end{lemma}

\begin{proof}
由于 $x/y \cdot y/x = 1$，映射
$R[x/y] \oplus R[y/x] \to R[x/y, y/x]$ 是满射。令
$\alpha \in R[x/y] \cap R[y/x]$。为证明中间处正合，须证明
$\alpha \in R$。由假设，可写成
$$
\alpha = a_0 + a_1 x/y + \ldots + a_n (x/y)^n =
b_0 + b_1 y/x + \ldots + b_m(y/x)^m
$$
其中 $n, m \geq 0$ 且 $a_i, b_j \in R$。取 $N > \max(n, m)$。
考虑 $R_{xy}$ 中由下列元素生成的有限 $R$-子模 $M$：
$$
(x/y)^N, (x/y)^{N - 1}, \ldots, x/y, 1, y/x, \ldots,
(y/x)^{N - 1}, (y/x)^N
$$
断言 $\alpha M \subset M$。事实上，对 $0 \leq i \leq N$，显然有
$(x/y)^i (b_0 + b_1 y/x + \ldots + b_m(y/x)^m) \in M$；并且
$(y/x)^i (a_0 + a_1 x/y + \ldots + a_n(x/y)^n) \in M$；对 $0 \leq i \leq N$。因此由引理
\ref{algebra-lemma-characterize-integral-element}，$\alpha$ 整于 $R$。注意
$\alpha \in R_x$，所以若 $R$ 在 $R_x$ 中整闭，则 $\alpha \in R$，结论得证。
\end{proof}

\section{正规环}
\label{algebra-section-normal-rings}

\noindent
我们先介绍正规整环的概念，然后介绍（非常一般的）正规环概念。

\begin{definition}
\label{algebra-definition-domain-normal}
若整环 $R$ 在其分式域中整闭，则称其为{\it 正规}。
\end{definition}

\begin{lemma}
\label{algebra-lemma-integral-closure-in-normal}
设 $R \to S$ 为环映射。若 $S$ 是正规整环，则 $R$
在 $S$ 中的整闭包是正规整环。
\end{lemma}

\begin{proof}
略。
\end{proof}

\noindent
下述概念在研究正规性时偶尔有用。

\begin{definition}
\label{algebra-definition-almost-integral}
设 $R$ 为整环。
\begin{enumerate}
\item 分式域中的元素 $g$ 称为{\it 几乎整于 $R$}，若存在 $r \in R$，
$r\not = 0$，使得对所有 $n \geq 0$ 都有 $rg^n \in R$。
\item 若 $R$ 的分式域中的每个几乎整元素都属于 $R$，则称整环 $R$
为{\it 完全正规}。
\end{enumerate}
\end{definition}

\noindent
下述引理表明，诺特整环是正规整环，当且仅当它完全正规。

\begin{lemma}
\label{algebra-lemma-almost-integral}
设 $R$ 为分式域 $K$ 中的整环。若 $u, v \in K$ 几乎整于 $R$，
则 $u + v$ 与 $uv$ 也几乎整于 $R$。任何整于 $R$ 的元素 $g \in K$
都几乎整于 $R$。若 $R$ 是诺特环，则反之亦然。
\end{lemma}

\begin{proof}
若对所有 $n \geq 0$ 都有 $ru^n \in R$，且
$v^nr' \in R$，则对所有 $n \geq 0$，
$(uv)^nrr'$ 与 $(u + v)^nrr'$ 都属于 $R$，对所有 $n \geq 0$。于是得到第一个断言。
设 $g \in K$ 整于 $R$。此时存在 $d > 0$，使环 $R[g]$ 作为
$R$-模由 $1, g, \ldots, g^d$ 生成。令 $r \in R$ 为
$1, g, \ldots, g^d \in K$ 的公共分母。于是
$rR[g] \subset R$，所以 $g$ 几乎整于 $R$。

\medskip\noindent
设 $R$ 是诺特环，且 $g \in K$ 几乎整于 $R$。取定义中的
$r \in R$，满足 $r\not = 0$。则作为 $R$-模，
$R[g] \subset \frac{1}{r}R$。由于 $R$ 是诺特环，这意味着
$R[g]$ 在 $R$ 上有限。因此 $g$ 整于 $R$，见引理
\ref{algebra-lemma-finite-is-integral}。
\end{proof}

\begin{lemma}
\label{algebra-lemma-localize-normal-domain}
正规整环的任意局部化仍是正规整环。
\end{lemma}

\begin{proof}
设 $R$ 为正规整环，$S \subset R$ 为乘法子集。设 $g$ 是
$R$ 的分式域中的元素，且整于 $S^{-1}R$。令
$P = x^d + \sum_{j < d} a_j x^j$ 为满足 $P(g) = 0$ 的多项式，
其中 $a_i \in S^{-1}R$。选取 $s \in S$，使得对所有 $i$ 都有
$sa_i \in R$。于是 $sg$ 满足首一多项式
$x^d + \sum_{j < d} s^{d-j}a_j x^j$，其系数
$s^{d-j}a_j$ 属于 $R$。由于 $R$ 正规，故 $sg \in R$。
于是 $g \in S^{-1}R$。
\end{proof}

\begin{lemma}
\label{algebra-lemma-PID-normal}
主理想整环是正规环。
\end{lemma}

\begin{proof}
设 $R$ 为主理想整环。设 $g = a/b$ 是 $R$ 的分式域中
整于 $R$ 的元素。由于 $R$ 是主理想整环，我们可以约去 $a$ 与
$b$ 的公因子，并设 $(a, b) = R$。此时，$r_i \in R$ 的任意方程
$(a/b)^n + r_{n-1} (a/b)^{n-1} + \ldots + r_0 = 0$
都会推出 $a^n \in (b)$。除非 $b$ 是 $R$ 中的单位，否则这与
$(a, b) = R$ 矛盾。
\end{proof}

\begin{lemma}
\label{algebra-lemma-prepare-polynomial-ring-normal}
设 $R$ 为分式域为 $K$ 的整环。设
$f = \sum \alpha_i x^i$ 是 $K[x]$ 中的元素。
\begin{enumerate}
\item 若 $f$ 整于 $R[x]$，则所有的 $\alpha_i$ 都整于 $R$；并且
\item 若 $f$ 几乎整于 $R[x]$，则所有的 $\alpha_i$ 都几乎整于 $R$。
\end{enumerate}
\end{lemma}

\begin{proof}
我们先证明第二个断言。写成
$f = \alpha_0 + \alpha_1 x + \ldots + \alpha_r x^r$，其中
$\alpha_r \not = 0$。根据假设，存在
$h = b_0 + b_1 x + \ldots + b_s x^s \in R[x]$，$b_s \not = 0$，
使得对所有 $n \geq 0$ 都有 $f^n h \in R[x]$。这意味着对所有
$n \geq 0$ 都有 $b_s \alpha_r^n \in R$。因此 $\alpha_r$ 几乎整于
$R$。由于几乎整元素的集合构成一个子环（引理
\ref{algebra-lemma-almost-integral}），我们得到
$f - \alpha_r x^r = \alpha_0 + \alpha_1 x + \ldots + \alpha_{r - 1} x^{r - 1}$
几乎整于 $R[x]$。对 $r$ 归纳即得。

\medskip\noindent
为证明第一个断言，我们使用绝对诺特化约。具体地，写
$\alpha_i = a_i / b_i$，并令
$P(t) = t^d + \sum_{j < d} f_j t^j$ 为系数 $f_j \in R[x]$ 且满足
$P(f) = 0$ 的多项式。令 $f_j = \sum f_{ji}x^i$。考虑由
$R$ 中有限个元素 $a_i, b_i, f_{ji}$ 生成的子环
$R_0 \subset R$。它是整环；令 $K_0$ 为其分式域。由于 $R_0$ 是有限型
$\mathbf{Z}$-代数，它是诺特的，见引理
\ref{algebra-lemma-obvious-Noetherian}。此时 $f \in K_0[x]$ 整于 $R_0[x]$ 仍然成立，
因为 $R$ 中关于元素 $a_i, b_i, f_{ji}$ 的所有恒等式在 $R_0$ 中也成立。
由引理 \ref{algebra-lemma-almost-integral}，元素 $f$ 几乎整于 $R_0[x]$。
由该引理的第二个断言，元素 $\alpha_i$ 几乎整于 $R_0$。又因为 $R_0$ 是诺特的，
它们整于 $R_0$，见引理 \ref{algebra-lemma-almost-integral}。当然，它们于是也整于 $R$。
\end{proof}

\begin{lemma}
\label{algebra-lemma-polynomial-domain-normal}
设 $R$ 为正规整环。则 $R[x]$ 是正规整环。
\end{lemma}

\begin{proof}
当 $R$ 是域 $K$ 时结论成立，因为 $K[x]$ 是欧几里得整环，因而是主理想整环，
再由引理 \ref{algebra-lemma-PID-normal} 得到它是正规整环。设 $g$ 是 $R[x]$ 的分式域中
整于 $R[x]$ 的元素。由于 $g$ 整于 $K[x]$（其中 $K$ 是 $R$ 的分式域），我们可以写成
$g = \alpha_d x^d + \alpha_{d-1}x^{d-1} + \ldots + \alpha_0$，其中 $\alpha_i \in K$。
由引理 \ref{algebra-lemma-prepare-polynomial-ring-normal}，元素 $\alpha_i$ 整于 $R$，
因而属于 $R$。
\end{proof}

\begin{lemma}
\label{algebra-lemma-power-series-over-Noetherian-normal-domain}
设 $R$ 为诺特正规整环。则 $R[[x]]$ 是诺特正规整环。
\end{lemma}

\begin{proof}
由引理 \ref{algebra-lemma-Noetherian-power-series}，幂级数环是诺特的。设
$f, g \in R[[x]]$ 为非零元素，且 $w = f/g$ 整于 $R[[x]]$。令 $K$ 为 $R$ 的分式域。
由于劳朗级数环
$K((x)) = K[[x]][1/x]$ 是域，我们可以写成
$w = a_n x^n + a_{n + 1} x^{n + 1} + \ldots$，其中 $n \in \mathbf{Z}$，
$a_i \in K$，且 $a_n \not = 0$。由引理 \ref{algebra-lemma-almost-integral}，存在
$R[[x]]$ 中的非零元素 $h = b_m x^m + b_{m + 1} x^{m + 1} + \ldots$，其中
$b_m \not = 0$，使得对所有 $e \geq 1$ 都有 $w^e h \in R[[x]]$。我们得到
$n \geq 0$，且对所有 $e \geq 1$ 都有 $b_m a_n^e \in R$。
由于 $R$ 是诺特环，由同一引理可知 $a_n \in R$。现在，若
$a_n, a_{n + 1}, \ldots, a_{N - 1} \in R$，则可对
$w - a_n x^n - \ldots - a_{N - 1} x^{N - 1} = a_N x^N + \ldots$
应用同样的论证。如此可知所有 $a_i \in R$，引理得证。
\end{proof}

\begin{lemma}
\label{algebra-lemma-normality-is-local}
设 $R$ 为整环。以下命题等价：
\begin{enumerate}
\item 整环 $R$ 是正规整环；
\item 对每个素理想 $\mathfrak p \subset R$，局部环
$R_{\mathfrak p}$ 是正规整环；并且
\item 对每个极大理想 $\mathfrak m$，环 $R_{\mathfrak m}$
是正规整环。
\end{enumerate}
\end{lemma}

\begin{proof}
由引理 \ref{algebra-lemma-localize-normal-domain}，得到 (1) $\Rightarrow$ (2)。
(2) $\Rightarrow$ (3) 显然。(3) $\Rightarrow$ (1) 则由以下事实得到：对于任意整环 $R$，
有
$$
R = \bigcap\nolimits_{\mathfrak m} R_{\mathfrak m}
$$
（在 $R$ 的分式域内）。事实上，若 $g$ 属于右端，则理想
$I = \{x \in R \mid xg \in R\}$ 不包含于任何极大理想 $\mathfrak m$ 中，
从而 $I = R$。
\end{proof}

\noindent
引理 \ref{algebra-lemma-normality-is-local} 表明，下面的定义与定义
\ref{algebra-definition-domain-normal} 相容。（这是 EGA 中的定义——见
\cite[IV, 5.13.5 and 0, 4.1.4]{EGA}。）

\begin{definition}
\label{algebra-definition-ring-normal}
若对每个素理想 $\mathfrak p \subset R$，局部化 $R_{\mathfrak p}$
是正规整环（见定义 \ref{algebra-definition-domain-normal}），则称环 $R$ 为{\it 正规}。
\end{definition}

\noindent
注意，正规环是约化环，因为 $R$ 是其在所有素理想处局部化所得环的乘积的子环
（例如见引理 \ref{algebra-lemma-characterize-zero-local}）。

\begin{lemma}
\label{algebra-lemma-normal-ring-integrally-closed}
正规环在其全分式环中整闭。
\end{lemma}

\begin{proof}
设 $R$ 为正规环。设 $x \in Q(R)$ 是 $R$ 的全分式环中整于 $R$ 的元素。
令 $I = \{f \in R, fx \in R\}$。设 $\mathfrak p \subset R$ 为素理想。
由于 $R \to R_{\mathfrak p}$ 是平坦的，有
$R_{\mathfrak p} \subset Q(R) \otimes_R R_{\mathfrak p}$。由于
$R_{\mathfrak p}$ 是正规整环，故 $x \otimes 1$ 是 $R_{\mathfrak p}$ 中的元素。
因此可找到 $a, f \in R$，$f \not \in \mathfrak p$，使得
$x \otimes 1 = a \otimes 1/f$。这意味着 $fx - a$ 在
$Q(R) \otimes_R R_{\mathfrak p} = Q(R)_{\mathfrak p}$ 中映为零，而这又意味着
存在 $f' \in R$，$f' \not \in \mathfrak p$，使得在 $R$ 中有
$f'fx = f'a$。换言之，$ff' \in I$。因此 $I$ 是一个不包含于 $R$ 的任何素理想
的理想，即 $I = R$，且 $x \in R$。
\end{proof}

\begin{lemma}
\label{algebra-lemma-localization-normal-ring}
正规环的局部化仍是正规环。
\end{lemma}

\begin{proof}
略。
\end{proof}

\begin{lemma}
\label{algebra-lemma-polynomial-ring-normal}
设 $R$ 为正规环。则 $R[x]$ 是正规环。
\end{lemma}

\begin{proof}
设 $\mathfrak q$ 是 $R[x]$ 中的素理想。令
$\mathfrak p = R \cap \mathfrak q$。由引理
\ref{algebra-lemma-polynomial-domain-normal}，可知 $R_{\mathfrak p}[x]$ 是正规整环。
于是由引理 \ref{algebra-lemma-localize-normal-domain}，$(R[x])_{\mathfrak q}$ 是正规整环。
\end{proof}

\begin{lemma}
\label{algebra-lemma-finite-product-normal}
正规环的有限乘积是正规环。
\end{lemma}

\begin{proof}
只须证明两个正规环（记为 $R$ 和 $S$）的乘积是正规环。由引理
\ref{algebra-lemma-disjoint-decomposition}，$R\times S$ 的素理想具有
$\mathfrak{p}\times S$ 和 $R\times \mathfrak{q}$ 的形式，其中
$\mathfrak{p}$ 和 $\mathfrak{q}$ 分别是 $R$ 和 $S$ 的素理想。局部化给出
$(R\times S)_{\mathfrak{p}\times S}=R_{\mathfrak{p}}$，它由假设是正规整环。
对 $S$ 的情形同理。
\end{proof}

\begin{lemma}
\label{algebra-lemma-characterize-reduced-ring-normal}
设 $R$ 为环。假设 $R$ 是约化的并且只有有限多个极小素理想。
则以下命题等价：
\begin{enumerate}
\item $R$ 是正规环；
\item $R$ 在其全分式环中整闭；并且
\item $R$ 是正规整环的有限乘积。
\end{enumerate}
\end{lemma}

\begin{proof}
(1) $\Rightarrow$ (2) 以及 (3) $\Rightarrow$ (1) 的蕴含一般都成立，
见引理 \ref{algebra-lemma-normal-ring-integrally-closed} 和
\ref{algebra-lemma-finite-product-normal}。

\medskip\noindent
设 $\mathfrak p_1, \ldots, \mathfrak p_n$ 是 $R$ 的极小素理想。
由引理 \ref{algebra-lemma-reduced-ring-sub-product-fields} 和
\ref{algebra-lemma-total-ring-fractions-no-embedded-points}，有
$Q(R) = R_{\mathfrak p_1} \times \ldots \times R_{\mathfrak p_n}$；
并且由引理 \ref{algebra-lemma-minimal-prime-reduced-ring}，每个因子都是域。
记 $e_i = (0, \ldots, 0, 1, 0, \ldots, 0)$ 为 $Q(R)$ 的第 $i$ 个幂等元。

\medskip\noindent
若 $R$ 在 $Q(R)$ 中整闭，则特别地包含幂等元 $e_i$，从而可知 $R$ 是
$n$ 个整环的乘积（见第 \ref{algebra-section-connected-components} 节和
\ref{algebra-section-tilde-module-sheaf} 节）。每个因子具有
$R/\mathfrak p_i$ 的形式，其分式域为 $R_{\mathfrak p_i}$。
由引理 \ref{algebra-lemma-finite-product-integral-closure}，每个映射
$R/\mathfrak p_i \to R_{\mathfrak p_i}$ 都是整闭的。
因此 $R$ 是正规整环的有限乘积。
\end{proof}

\begin{lemma}
\label{algebra-lemma-colimit-normal-ring}
设 $(R_i, \varphi_{ii'})$ 是环的有向系统
（范畴论定义 \ref{algebra-definition-directed-system}）。若每个 $R_i$ 都是正规环，
则 $R = \colim_i R_i$ 也是正规环。
\end{lemma}

\begin{proof}
设 $\mathfrak p \subset R$ 为素理想。令
$\mathfrak p_i = R_i \cap \mathfrak p$（通常的记号滥用）。于是有
$R_{\mathfrak p} = \colim_i (R_i)_{\mathfrak p_i}$。
由于每个 $(R_i)_{\mathfrak p_i}$ 都是正规整环，我们归约为证明正规整环的情形。
若 $a, b \in R$ 且 $a/b$ 满足 $P(T) \in R[T]$ 这一首一多项式，
则可找到充分大的 $i \in I$，使得 $a, b$ 分别来自 $R_i$ 上的元素
$a_i, b_i$，而 $P$ 来自首一多项式 $P_i\in R_i[T]$，并且
$P_i(a_i/b_i)=0$。由于 $R_i$ 正规，有 $a_i/b_i \in R_i$，从而
$a/b \in R$。
\end{proof}

\section{正规环上整扩张的下降}
\label{algebra-section-going-down-integral-over-normal}

\noindent
我们先稍微研究“整于理想”的元素这一概念，然后证明本节标题所指的定理。

\begin{definition}
\label{algebra-definition-integral-over-ideal}
设 $\varphi : R \to S$ 为环映射。设 $I \subset R$ 为理想。
若元素 $g \in S$ 满足下述条件，则称该元素
{\it 整于 $I$}：存在首一多项式
$P = x^d + \sum_{j < d} a_j x^j$，其系数满足
$a_j \in I^{d-j}$，并且
$P^\varphi(g) = 0$ 在 $S$ 中成立。
\end{definition}

\noindent
这一概念主要用于 $\varphi = \text{id}_R : R \to R$ 的情形。
此时，整于 $I$ 的元素所成的集合 $I'$ 称为{\it $I$ 的整闭包}。
我们将看到 $I'$ 是 $R$ 的一个理想（当然 $I \subset I'$）。

\begin{lemma}
\label{algebra-lemma-characterize-integral-ideal}
设 $\varphi : R \to S$ 为环映射。设 $I \subset R$ 为理想。
令 $A = \sum I^nt^n \subset R[t]$ 为多项式环中由
$R \oplus It \subset R[t]$ 生成的子环。元素 $s \in S$ 整于 $I$，
当且仅当元素 $st \in S[t]$ 整于 $A$。
\end{lemma}

\begin{proof}
假设 $st$ 整于 $A$。令
$P = x^d + \sum_{j < d} a_j x^j$ 为系数在 $A$ 中的首一多项式，
满足 $P^\varphi(st) = 0$。令 $a_j' \in A$ 为 $a_j$ 的
$d-j$ 次部分，换言之，$a_j' = a_j'' t^{d-j}$，其中
$a_j'' \in I^{d-j}$。根据次数的原因，仍有
$(st)^d + \sum_{j < d} \varphi(a_j'') t^{d-j} (st)^j = 0$。
因此
$s^d + \sum_{j < d} \varphi(a_j'') s^j = 0$，
从而 $s$ 整于 $I$。

\medskip\noindent
假设 $s$ 整于 $I$。设
$P = x^d + \sum_{j < d} a_j x^j$，其中 $a_j \in I^{d-j}$。
则立即得到多项式
$Q = x^d + \sum_{j < d} (a_j t^{d-j}) x^j$，其系数在 $A$ 中，
这证明了 $st$ 整于 $A$。
\end{proof}

\begin{lemma}
\label{algebra-lemma-integral-over-ideal-is-submodule}
设 $\varphi : R \to S$ 为环映射。设 $I \subset R$ 为理想。
$S$ 中整于 $I$ 的元素所成的集合是 $S$ 的一个 $R$-子模。
此外，若 $s \in S$ 整于 $R$，而 $s'$ 整于 $I$，则
$ss'$ 整于 $I$。
\end{lemma}

\begin{proof}
我们将使用引理 \ref{algebra-lemma-integral-closure-is-ring} 而不再另行说明。
由引理 \ref{algebra-lemma-characterize-integral-ideal} 的刻画，集合对加法封闭是显然的，
下文沿用该引理的记号。任何整于 $R$ 的元素 $s \in S$，都对应于 $S[t]$ 中的
次数 $0$ 元素 $s$，并且该元素整于 $A$（因为 $R \subset A$）。
因此，$S[t]$ 上乘以 $s$ 的运算保持整于 $A$ 的性质，
\end{proof}

\begin{lemma}
\label{algebra-lemma-integral-integral-over-ideal}
设 $\varphi : R \to S$ 为整环映射。设 $I \subset R$ 为理想。
则 $IS$ 中的每个元素都整于 $I$。
\end{lemma}

\begin{proof}
由引理 \ref{algebra-lemma-integral-over-ideal-is-submodule} 立即得到。
\end{proof}

\begin{lemma}
\label{algebra-lemma-polynomials-divide}
设 $K$ 为域。设 $n, m \in \mathbf{N}$，且
$a_0, \ldots, a_{n - 1}, b_0, \ldots, b_{m - 1} \in K$。
若多项式 $x^n + a_{n - 1}x^{n - 1} + \ldots + a_0$
在 $K[x]$ 中整除多项式 $x^m + b_{m - 1} x^{m - 1} + \ldots + b_0$，
则
\begin{enumerate}
\item 对于包含元素 $b_0, \ldots, b_{m - 1}$ 的 $K$ 的任意子环
$R_0$，$a_0, \ldots, a_{n - 1}$ 都整于该子环；并且
\item 对于包含元素
$a_0, \ldots, a_{n - 1}, b_0, \ldots, b_{m - 1}$ 的任意子环
$R \subset K$，每个 $a_i$ 都属于下列根理想：
$$
\sqrt{(b_0, \ldots, b_{m-1})R}
$$
\end{enumerate}
\end{lemma}

\begin{proof}
设 $L/K$ 为域扩张，使得可以写成
$$
x^m + b_{m - 1} x^{m - 1} + \ldots + b_0 =
\prod_{i = 1}^m (x - \beta_i)
$$
其中 $\beta_i \in L$。见域一章的
\ref{fields-section-splitting-fieds} 节。
每个 $\beta_i$ 都整于 $R_0$。由于每个 $a_i$ 都是
$\beta_1, \ldots, \beta_m$ 的齐次多项式，我们得到 $a_i$ 也具有同样性质
（使用引理 \ref{algebra-lemma-integral-closure-is-ring}）。这证明了 (1)。

\medskip\noindent
设 $R$ 如 (2) 中所述。选取 $c_0, \ldots, c_{m - n - 1} \in K$，使得
$$
\begin{matrix}
x^m + b_{m - 1} x^{m - 1} + \ldots + b_0 =  \\
(x^n + a_{n - 1}x^{n - 1} + \ldots + a_0)
(x^{m - n} + c_{m - n - 1}x^{m - n - 1}+ \ldots + c_0).
\end{matrix}
$$
该等式推出
\begin{align*}
c_{m - n - 1} & = b_{m - 1} - a_{n - 1}, \\
c_{m - n - 2} & = b_{m - 2} - a_{n - 2} - a_{n - 1}c_{m - n - 1}, \\
\ldots
\end{align*}
因此对所有 $j$ 都有 $c_j \in R$。除去根式
$\sqrt{(b_0, \ldots, b_{m - 1})}$ 后得到约化环
$\overline{R}$。我们必须证明像
$\overline{a}_i \in \overline{R}$ 都为零。在
$\overline{R}[x]$ 中有关系
$$
\begin{matrix}
x^m = x^m + \overline{b}_{m - 1} x^{m - 1} + \ldots + \overline{b}_0 = \\
(x^n + \overline{a}_{n - 1}x^{n - 1} + \ldots + \overline{a}_0)
(x^{m - n} + \overline{c}_{m - n - 1}x^{m - n - 1}+ \ldots + \overline{c}_0).
\end{matrix}
$$
容易看出这意味着对所有 $i$ 都有 $\overline{a}_i = 0$。事实上，由引理
\ref{algebra-lemma-minimal-prime-reduced-ring}，$\overline{R}$ 在极小素理想
$\mathfrak{p}$ 处的局部化是域，而 $\overline{R}_{\mathfrak p}[x]$ 是唯一分解整环。
因此
$f = x^n + \sum \overline{a}_i x^i$
与 $x^n$ 相伴；又由于 $f$ 是首一的，在
$\overline{R}_{\mathfrak p}[x]$ 中有 $f = x^n$。
于是存在 $s \in \overline{R}$，$s \not\in \mathfrak p$，使得
$s(f - x^n) = 0$。因此所有 $\overline{a}_i$ 都属于
$\mathfrak p$，再由引理 \ref{algebra-lemma-reduced-ring-sub-product-fields} 得到结论。
\end{proof}

\begin{lemma}
\label{algebra-lemma-minimal-polynomial-normal-domain}
设 $R \subset S$ 为整环的包含。假设 $R$ 正规。设 $g \in S$ 整于
$R$。则 $g$ 的最小多项式的系数属于 $R$。
\end{lemma}

\begin{proof}
设 $P = x^m + b_{m-1} x^{m-1} + \ldots + b_0$
为系数在 $R$ 中且满足 $P(g) = 0$ 的多项式。设
$Q = x^n + a_{n-1}x^{n-1} + \ldots + a_0$
为 $g$ 在 $R$ 的分式域 $K$ 上的最小多项式。于是 $Q$ 在 $K[x]$ 中整除
$P$。由引理 \ref{algebra-lemma-polynomials-divide} 可知 $a_i$ 整于 $R$。
由于 $R$ 正规，这意味着它们属于 $R$。
\end{proof}

\begin{proposition}
\label{algebra-proposition-going-down-normal-integral}
设 $R \subset S$ 为整环的包含。假设 $R$ 正规且 $S$ 整于 $R$。
设 $\mathfrak p \subset \mathfrak p' \subset R$ 为素理想。设
$\mathfrak q'$ 是 $S$ 中满足 $\mathfrak p' = R \cap \mathfrak q'$ 的素理想。
则存在素理想 $\mathfrak q$，满足
$\mathfrak q \subset \mathfrak q'$，并且 $\mathfrak p = R \cap \mathfrak q$。
换言之，$R \to S$ 满足下降性质，见定义
\ref{algebra-definition-going-up-down}。
\end{proposition}

\begin{proof}
令 $\mathfrak p$、$\mathfrak p'$ 和 $\mathfrak q'$ 如命题所述。我们要证明存在
素理想 $\mathfrak q$，使得 $\mathfrak q \subset \mathfrak q'$ 且
$R \cap \mathfrak q = \mathfrak p$。这等价于寻找一个映到 $\mathfrak p$ 的
$S_{\mathfrak q'}$ 的素理想。
根据引理 \ref{algebra-lemma-in-image}，我们需要证明
$\mathfrak p S_{\mathfrak q'} \cap R
= \mathfrak p$。取 $z \in \mathfrak p S_{\mathfrak q'} \cap R$。
可写成 $z = y/g$，其中 $y \in \mathfrak pS$ 且
$g \in S$，$g \not\in \mathfrak q'$。换一种写法就是 $zg = y$。

\medskip\noindent
由引理 \ref{algebra-lemma-integral-integral-over-ideal}，存在首一多项式
$P = x^m + b_{m-1} x^{m-1} + \ldots + b_0$，
其中 $b_i \in \mathfrak p$，并且 $P(y) = 0$。

\medskip\noindent
由引理 \ref{algebra-lemma-minimal-polynomial-normal-domain}，$g$ 在 $K$ 上的
最小多项式的系数属于 $R$。将其写成
$Q = x^n + a_{n-1} x^{n-1} + \ldots
+ a_0$。注意，并非所有 $a_i$（$i = n-1, \ldots, 0$）
都属于 $\mathfrak p$，因为这会推出
$g^n = \sum_{j < n} a_j g^j \in \mathfrak pS
\subset \mathfrak p'S
\subset \mathfrak q'$，
从而产生矛盾。

\medskip\noindent
由于 $y = zg$，由上面的论证立即可知
$Q' = x^n + za_{n-1} x^{n-1} + \ldots + z^{n}a_0$
是 $y$ 的最小多项式。因此 $Q'$ 整除 $P$，由引理
\ref{algebra-lemma-polynomials-divide} 可知
$z^ja_{n - j} \in \sqrt{(b_0, \ldots, b_{m-1})}
\subset \mathfrak p$，$j =  1, \ldots, n$。
由于并非所有 $a_i$（$i = n-1, \ldots, 0$）都属于 $\mathfrak p$，
我们得到 $z \in \mathfrak p$，命题得证。
\end{proof}

\section{平坦模与平坦环映射}
\label{algebra-section-flat}

\noindent
一个常用结果是：若 $M = \colim_{i\in \mathcal{I}} M_i$
是 $R$-模的余极限，且 $N$ 是 $R$-模，则
$$
M \otimes N
=
\colim_{i\in \mathcal{I}} M_i \otimes_R N,
$$
见引理 \ref{algebra-lemma-tensor-products-commute-with-limits}。
通常将这一性质表述为{\it $\otimes$ 与余极限可交换}。
另一个常用结果是：若 $0 \to N_1 \to N_2 \to N_3 \to 0$
是正合序列，且 $M$ 是任意 $R$-模，则
$$
M \otimes_R N_1
\to
M \otimes_R N_2
\to
M \otimes_R N_3
\to
0
$$
仍然正合，见引理 \ref{algebra-lemma-tensor-product-exact}。
这两个性质都说明函子
$N \mapsto M \otimes_R N$ {\it 是右正合的}。
见范畴论的 \ref{categories-section-exact-functor} 节和
同调代数的 \ref{homology-section-functors} 节。
若 $N \mapsto N \otimes_R M$ 也是左正合的，即它是正合的，
则称 $R$-模 $M$ 为平坦模。下面给出精确定义。

\begin{definition}
\label{algebra-definition-flat}
设 $R$ 为环。
\begin{enumerate}
\item 若每当
$N_1 \to N_2 \to N_3$ 是 $R$-模的正合序列时，
序列 $M \otimes_R N_1 \to M \otimes_R N_2 \to M \otimes_R N_3$
也正合，则称 $R$-模 $M$ {\it 平坦}。
\item 若 $R$-模的复形
$N_1 \to N_2 \to N_3$ 正合，当且仅当
序列 $M \otimes_R N_1 \to M \otimes_R N_2 \to M \otimes_R N_3$
正合，则称 $R$-模 $M$ {\it 忠实平坦}。
\item 若环映射 $R \to S$ 使得 $S$ 作为 $R$-模平坦，
则称该环映射{\it 平坦}。
\item 若环映射 $R \to S$ 使得 $S$ 作为 $R$-模忠实平坦，
则称该环映射{\it 忠实平坦}。
\end{enumerate}
\end{definition}

\noindent
下面给出一个使用平坦性条件的例子。

\begin{lemma}
\label{algebra-lemma-flat-intersect-ideals}
设 $R$ 为环。设 $I, J \subset R$ 为理想。设 $M$ 为平坦的
$R$-模。则 $IM \cap JM = (I \cap J)M$。
\end{lemma}

\begin{proof}
考虑正合序列 $0 \to I \cap J \to R \to R/I \oplus R/J$。
与平坦模 $M$ 作张量，得到正合序列
$$
0 \to (I \cap J) \otimes_R M \to M \to M/IM \oplus M/JM
$$
由于 $M \to M/IM \oplus M/JM$ 的核等于
$IM \cap JM$，结论成立。
\end{proof}

\begin{lemma}
\label{algebra-lemma-colimit-flat}
设 $R$ 为环。设 $\{M_i, \varphi_{ii'}\}$ 为平坦 $R$-模所成的有向系统。
则 $\colim_i M_i$ 是平坦的 $R$-模。
\end{lemma}

\begin{proof}
这是因为 $\otimes$ 与余极限可交换，且有向余极限是正合的，见引理
\ref{algebra-lemma-directed-colimit-exact}。
\end{proof}

\begin{lemma}
\label{algebra-lemma-composition-flat}
（忠实）平坦环映射的复合仍是（忠实）平坦的。
若 $R \to R'$ 是（忠实）平坦的，且 $M'$ 是
（忠实）平坦的 $R'$-模，则 $M'$ 是（忠实）平坦的 $R$-模。
\end{lemma}

\begin{proof}
引理的第一个断言是第二个断言的特例，因此只须证明后者。设
$R \to R'$ 为平坦环映射，$M'$ 为平坦的 $R'$-模。
我们需要证明 $M'$ 是平坦的 $R$-模。设
$N_1 \to N_2 \to N_3$ 是 $R$-模的正合复形。
于是复形 $R' \otimes_R N_1 \to
R' \otimes_R N_2 \to R' \otimes_R N_3$ 正合（因为 $R'$
作为 $R$-模平坦），从而复形
$M' \otimes_{R'} \left(R' \otimes_R N_1\right)
\to M' \otimes_{R'} \left(R' \otimes_R N_2\right)
\to M' \otimes_{R'} \left(R' \otimes_R N_3\right)$
正合（因为 $M'$ 是平坦的 $R'$-模）。由于对任意 $R$-模 $N$ 都有函子性同构
$M' \otimes_{R'} \left(R' \otimes_R N\right)
\cong \left(M' \otimes_{R'} R'\right) \otimes_R N
\cong M' \otimes_R N$
（由引理 \ref{algebra-lemma-tensor-with-bimodule} 和
\ref{algebra-lemma-flip-tensor-product}），上述复形同构于复形
$M' \otimes_R N_1 \to M' \otimes_R N_2 \to M' \otimes_R N_3$，
因而后者也正合。这说明 $M'$ 是平坦的 $R$-模。沿相反方向追溯这一论证可知：
若 $R \to R'$ 忠实平坦，且 $M'$ 作为 $R'$-模忠实平坦，
则 $M'$ 作为 $R$-模也忠实平坦。
\end{proof}

\begin{lemma}
\label{algebra-lemma-flat}
设 $M$ 为 $R$-模。以下命题等价：
\begin{enumerate}
\item
\label{algebra-item-flat}
$M$ 在 $R$ 上平坦。
\item
\label{algebra-item-injective}
对于每个 $R$-模的单射 $N \subset N'$，
映射 $N \otimes_R M \to N'\otimes_R M$ 是单射。
\item
\label{algebra-item-f-ideal}
对于每个理想 $I \subset R$，映射
$I \otimes_R M \to R \otimes_R M = M$ 是单射。
\item
\label{algebra-item-ffg-ideal}
对于每个有限生成理想 $I \subset R$，
映射 $I \otimes_R M \to R \otimes_R M = M$ 是单射。
\end{enumerate}
\end{lemma}

\begin{proof}
蕴含（\ref{algebra-item-flat}）蕴含（\ref{algebra-item-injective}）
蕴含（\ref{algebra-item-f-ideal}）蕴含（\ref{algebra-item-ffg-ideal}）
都是平凡的。因此我们证明（\ref{algebra-item-ffg-ideal}）蕴含（\ref{algebra-item-flat}）。
设 $N_1 \to N_2 \to N_3$ 正合。令
$K = \Ker(N_2 \to N_3)$，$Q = \Im(N_2 \to N_3)$。
于是得到映射
$$
N_1 \otimes_R M \to
K \otimes_R M \to
N_2 \otimes_R M \to
Q \otimes_R M \to
N_3 \otimes_R M
$$
注意第一个和第三个箭头都是满射。因此只要证明第二个和第四个箭头是单射即可，
从而完成证明\footnote{更详细的论证如下：
假设第二个和第四个箭头是单射。将引理
\ref{algebra-lemma-tensor-product-exact} 应用于正合序列
$K \to N_2 \to Q \to 0$，得到序列
$K \otimes_R M \to N_2 \otimes_R M \to
Q \otimes_R M \to 0$ 正合。因此
$\Ker \left(N_2 \otimes_R M \to Q \otimes_R M\right)
= \Im \left(K \otimes_R M \to N_2 \otimes_R M\right)$。
由于
$\Im \left(K \otimes_R M \to N_2 \otimes_R M\right)
= \Im \left(N_1 \otimes_R M \to N_2 \otimes_R M\right)$
（这是因为 $N_1 \otimes_R M \to
K \otimes_R M$ 满射），并且
$\Ker \left(N_2 \otimes_R M \to Q \otimes_R M\right)
= \Ker \left(N_2 \otimes_R M \to N_3 \otimes_R M\right)$
（这是因为 $Q \otimes_R M \to
N_3 \otimes_R M$ 单射），于是得到
$\Ker \left(N_2 \otimes_R M \to N_3 \otimes_R M\right)
= \Im \left(N_1 \otimes_R M \to N_2 \otimes_R M\right)$，
这说明函子 $- \otimes_R M$ 是正合的，因而 $M$ 平坦。}。
所以只须证明 $- \otimes_R M$ 将单射的 $R$-模映射变为单射的
$R$-模映射。

\medskip\noindent
假设 $K \to N$ 是单射的 $R$-模映射，并令
$x \in \Ker(K \otimes_R M \to N \otimes_R M)$。
我们要证明 $x$ 为零。$R$-模 $K$ 是其有限
$R$-子模的并，因此 $K \otimes_R M$ 是形如
$K_i \otimes_R M$ 的 $R$-模的余极限，其中 $K_i$ 遍历 $K$ 的所有有限
$R$-子模（因为张量积与余极限可交换）。
因此，对某个 $i$，$x$ 来自元素
$x_i \in K_i \otimes_R M$。于是可假设 $K$
是有限 $R$-模。以下均作此假设。我们把单射
$K \to N$ 视为包含，从而 $K \subset N$。

\medskip\noindent
$R$-模 $N$ 是所有包含 $K$ 的有限
$R$-子模的并。因此 $N \otimes_R M$
是形如 $N_i \otimes_R M$ 的 $R$-模的余极限，其中 $N_i$ 遍历 $N$ 中所有
包含 $K$ 的有限 $R$-子模（同样因为张量积与余极限可交换）。
注意这是一个有向系统上的余极限（因为 $N$ 的两个有限子模之和仍然有限）。
所以由引理 \ref{algebra-lemma-zero-directed-limit}，元素
$x \in K \otimes_R M$ 至少在其中一个这样的 $R$-模
$N_i \otimes_R M$ 中映为零（因为 $x$ 在 $N \otimes_R M$ 中映为零）。
因此可假设 $N$ 是有限 $R$-模。

\medskip\noindent
假设 $N$ 是有限 $R$-模。写成 $N = R^{\oplus n}/L$，并令
$K = L'/L$，其中 $L \subset L' \subset R^{\oplus n}$。
对于任意 $R$-子模 $G \subset R^{\oplus n}$，都有一个典范映射
$G \otimes_R M \to M^{\oplus n}$，它由复合映射
$G \otimes_R M \to R^n \otimes_R M = M^{\oplus n}$ 得到。
只须证明 $L \otimes_R M \to M^{\oplus n}$
和 $L' \otimes_R M \to M^{\oplus n}$ 都是单射。
事实上，此时
$K \otimes_R M = L' \otimes_R M/L \otimes_R M \to M^{\oplus n}/L \otimes_R M$
也是单射\footnote{将 $L' \otimes_R M$ 和 $L \otimes_R M$
分别识别为 $M^{\oplus n}$ 的子模后，这一点显然（这是合法的，因为映射
$L \otimes_R M \to M^{\oplus n}$ 和
$L' \otimes_R M \to M^{\oplus n}$ 都是单射，并且与显然的映射
$L' \otimes_R M \to L \otimes_R M$ 交换）。}。

\medskip\noindent
所以当 $L \subset R^{\oplus n}$ 是 $R$-子模时，只须证明
$L \otimes_R M \to M^{\oplus n}$ 是单射。
我们对 $n$ 归纳证明。下面处理归纳基础情形 $n = 1$。
归纳步骤中设 $n > 1$，并令
$L' = L \cap R \oplus 0^{\oplus n - 1}$。于是 $L'' = L/L'$ 是
$R^{\oplus n - 1}$ 的子模。我们得到图表
$$
\xymatrix{
&
L' \otimes_R M \ar[r] \ar[d] &
L \otimes_R M \ar[r] \ar[d] &
L'' \otimes_R M \ar[r] \ar[d] &
0 \\
0 \ar[r] &
M \ar[r] &
M^{\oplus n} \ar[r] &
M^{\oplus n - 1} \ar[r] & 0
}
$$
由归纳假设和基础情形，左右两条竖直箭头都是单射。两行都是正合的，
因此中间的竖直箭头也是单射。

\medskip\noindent
上述归纳的基础情形是 $L \subset R$ 为理想的情形。
换言之，对 $R$ 的任意理想 $I$，我们要证明
$I \otimes_R M \to M$ 是单射。当 $I$ 有限生成时我们已经知道结论成立。
然而，$I = \bigcup I_\alpha$ 是包含于其中的有限生成理想 $I_\alpha$ 的并。
换言之，$I = \colim I_\alpha$。由于 $\otimes$ 与余极限可交换，有
$I \otimes_R M = \colim I_\alpha \otimes_R M$；又因为由假设所有
$I_\alpha \otimes_R M \to M$ 的映射都是单射，
$I \otimes_R M \to M$ 也具有同样性质。
\end{proof}

\begin{lemma}
\label{algebra-lemma-colimit-rings-flat}
设 $\{R_i, \varphi_{ii'}\}$ 为有向集 $I$ 上的环系统。
令 $R = \colim_i R_i$。
\begin{enumerate}
\item 若 $M$ 是 $R$-模，并且对所有 $i$，$M$ 作为 $R_i$-模都是平坦的，
则 $M$ 是平坦的 $R$-模。
\item 对 $i \in I$，令 $M_i$ 为平坦的 $R_i$-模；对 $i' \geq i$，
令 $f_{ii'} : M_i \to M_{i'}$ 为 $\varphi_{ii'}$-线性映射，满足
$f_{i' i''} \circ f_{i i'} = f_{i i''}$。则
$M = \colim_{i \in I} M_i$ 是平坦的 $R$-模。
\end{enumerate}
\end{lemma}

\begin{proof}
第 (1) 部分是第 (2) 部分在对所有 $i$ 取 $M_i = M$ 且
$f_{i i'} = \text{id}_M$ 时的特例。下面证明 (2)。
设 $\mathfrak a \subset R$ 为有限生成理想。由
引理 \ref{algebra-lemma-flat}，只须证明 $\mathfrak a \otimes_R M \to M$ 是单射。
可找到 $i \in I$ 以及有限生成理想 $\mathfrak a' \subset R_i$，使得
$\mathfrak a = \mathfrak a'R$。于是
$\mathfrak a = \colim_{i' \geq i} \mathfrak a'R_{i'}$。
由于 $\otimes$ 与余极限可交换，映射 $\mathfrak a \otimes_R M \to M$
是下列映射的余极限：
$$
\mathfrak a'R_{i'} \otimes_{R_{i'}} M_{i'} \longrightarrow M_{i'}
$$
由假设这些映射都是单射。由于在 $I$ 上的余极限由引理
\ref{algebra-lemma-directed-colimit-exact} 可知是正合的，结论成立。
\end{proof}

\begin{lemma}
\label{algebra-lemma-flat-base-change}
设 $M$ 在 $R$ 上（忠实）平坦，且 $R \to R'$ 为环映射。
则 $M \otimes_R R'$ 在 $R'$ 上（忠实）平坦。
\end{lemma}

\begin{proof}
对于任意 $R'$-模 $N$，都有典范同构
$N \otimes_{R'} (R'\otimes_R M)
= N \otimes_R M$。因此函子
$-\otimes_{R'}(R'\otimes_R M)$ 所需的正合性性质，
由函子 $-\otimes_R M$ 的相应正合性性质推出。
\end{proof}

\begin{lemma}
\label{algebra-lemma-flatness-descends}
设 $R \to R'$ 为忠实平坦的环映射。设 $M$ 为 $R$-模，并令
$M' = R' \otimes_R M$。则 $M$ 在 $R$ 上平坦，当且仅当
$M'$ 在 $R'$ 上平坦。
\end{lemma}

\begin{proof}
由引理 \ref{algebra-lemma-flat-base-change}，若 $M$ 平坦，则 $M'$ 平坦。
反之，假设 $M'$ 平坦。设
$N_1 \to N_2 \to N_3$ 为 $R$-模的正合序列。
我们要证明 $N_1 \otimes_R M \to N_2 \otimes_R M \to N_3 \otimes_R M$
正合。由于 $R \to R'$ 平坦，我们知道
$N_1 \otimes_R R' \to N_2 \otimes_R R' \to N_3 \otimes_R R'$ 正合。
由 $M'$ 的平坦性，
$N_1 \otimes_R R' \otimes_{R'} M'
\to N_2 \otimes_R R' \otimes_{R'} M'
\to N_3 \otimes_R R' \otimes_{R'} M'$ 正合。
可以将其写成
$N_1 \otimes_R M \otimes_R R'
\to N_2 \otimes_R M \otimes_R R'
\to N_3 \otimes_R M \otimes_R R'$。
最后，由忠实平坦性可知
$N_1 \otimes_R M \to N_2 \otimes_R M \to N_3 \otimes_R M$
正合。
\end{proof}

\begin{lemma}
\label{algebra-lemma-flatness-descends-more-general}
设 $R$ 为环。设 $S \to S'$ 为 $R$-代数的平坦映射。
设 $M$ 为 $S$-模，并令 $M' = S' \otimes_S M$。
\begin{enumerate}
\item 若 $M$ 在 $R$ 上平坦，则 $M'$ 在 $R$ 上平坦。
\item 若 $S \to S'$ 忠实平坦，则 $M$ 在 $R$ 上平坦，当且仅当
$M'$ 在 $R$ 上平坦。
\end{enumerate}
\end{lemma}

\begin{proof}
设 $N \to N'$ 为 $R$-模的单射。由 $S \to S'$ 的平坦性，有
$$
\Ker(N \otimes_R M \to N' \otimes_R M) \otimes_S S'
=
\Ker(N \otimes_R M' \to N' \otimes_R M')
$$
若 $M$ 在 $R$ 上平坦，则左端为零，由平坦性的第二个刻画（引理
\ref{algebra-lemma-flat}）可知 $M'$ 在 $R$ 上平坦。
若 $M'$ 在 $R$ 上平坦，则右端为零；若再假设 $S \to S'$ 忠实平坦，
则这推出 $\Ker(N \otimes_R M \to N' \otimes_R M)$ 为零，
从而 $M$ 在 $R$ 上平坦。
\end{proof}

\begin{lemma}
\label{algebra-lemma-flat-permanence}
设 $R \to S$ 为环映射。设 $M$ 为 $S$-模。
若 $M$ 作为 $R$-模平坦且作为 $S$-模忠实平坦，
则 $R \to S$ 平坦。
\end{lemma}

\begin{proof}
设 $N_1 \to N_2 \to N_3$ 为 $R$-模的正合序列。
由假设，$N_1 \otimes_R M \to N_2 \otimes_R M \to N_3 \otimes_R M$
正合。可将其写成
$$
N_1 \otimes_R S \otimes_S M
\to
N_2 \otimes_R S \otimes_S M
\to
N_3 \otimes_R S \otimes_S M.
$$
由 $M$ 在 $S$ 上的忠实平坦性，得到
$N_1 \otimes_R S \to N_2 \otimes_R S \to N_3 \otimes_R S$ 正合。
因此 $R \to S$ 平坦。
\end{proof}

\noindent
设 $R$ 为环。设 $M$ 为 $R$-模。设
$\sum f_i x_i = 0$ 为 $M$ 中的一个关系。
若存在整数 $m \geq 0$、元素 $y_j \in M$（$j = 1, \ldots, m$）以及元素
$a_{ij} \in R$（$i = 1, \ldots, n$，$j = 1, \ldots, m$），使得
$$
x_i = \sum\nolimits_j a_{ij} y_j, \forall i,
\quad\text{且}\quad
0 = \sum\nolimits_i f_ia_{ij}, \forall j.
$$
则称关系 $\sum f_i x_i$ 为{\it 平凡的}。

\begin{lemma}[平坦性的方程判据]
\label{algebra-lemma-flat-eq}
$R$ 上的模 $M$ 平坦，当且仅当 $M$ 中的每个关系都是平凡的。
\end{lemma}

\begin{proof}
假设 $M$ 平坦，设 $\sum f_i x_i = 0$ 为 $M$ 中的关系。
令 $I = (f_1, \ldots, f_n)$，并令
$K = \Ker(R^n \to I, (a_1, \ldots, a_n) \mapsto \sum_i a_i f_i)$。
于是有短正合序列
$0 \to K \to R^n \to I \to 0$。此时
$\sum f_i \otimes x_i$ 是 $I \otimes_R M$ 中的元素，
并映到 $R \otimes_R M = M$ 中的零。由平坦性，
$\sum f_i \otimes x_i$ 在 $I \otimes_R M$ 中为零。
因此存在 $K \otimes_R M$ 中的元素映到
$\sum e_i \otimes x_i \in R^n \otimes_R M$，其中 $e_i$
是 $R^n$ 的第 $i$ 个基元素。将该元素写成
$\sum k_j \otimes y_j$，再将 $k_j$ 在 $R^n$ 中的像写成
$\sum a_{ij} e_i$，即得结论。

\medskip\noindent
假设每个关系都是平凡的。设 $I$ 为有限生成理想，设
$x = \sum f_i \otimes x_i$ 是 $I \otimes_R M$ 中映到
$R \otimes_R M = M$ 中零的元素。这恰好意味着
$\sum f_i x_i$ 是 $M$ 中的关系。由它平凡这一事实容易得到
$x$ 为零，因为
$$
x
=
\sum f_i \otimes x_i
=
\sum f_i \otimes \left(\sum a_{ij}y_j\right)
=
\sum \left(\sum f_i a_{ij}\right) \otimes y_j
=
0
$$
\end{proof}

\begin{lemma}
\label{algebra-lemma-flat-tor-zero}
设 $R$ 为环，$0 \to M'' \to M' \to M \to 0$
为短正合序列，且 $N$ 为 $R$-模。若 $M$ 平坦，则
$N \otimes_R M'' \to N \otimes_R M'$ 是单射，即序列
$$
0 \to N \otimes_R M'' \to N \otimes_R M' \to N \otimes_R M \to 0
$$
是短正合序列。
\end{lemma}

\begin{proof}
设 $R^{(I)} \to N$ 是从自由模到 $N$ 的满射，核为 $K$。
将蛇形引理应用于下图即可得到结论：
$$
\begin{matrix}
 & & 0 & & 0 & & 0 & & \\
 & & \uparrow & & \uparrow & & \uparrow & & \\
 & & M''\otimes_R N & \to & M' \otimes_R N & \to & M \otimes_R N & \to & 0 \\
 & & \uparrow & & \uparrow & & \uparrow & & \\
0 & \to & (M'')^{(I)} & \to & (M')^{(I)} & \to & M^{(I)} & \to & 0 \\
 & & \uparrow & & \uparrow & & \uparrow & & \\
 & & M''\otimes_R K & \to & M' \otimes_R K & \to & M \otimes_R K & \to & 0 \\
 & & & & & & \uparrow & & \\
 & & & & & & 0 & &
\end{matrix}
$$
其中各行各列均正合。中间一行正合，因为与自由模 $R^{(I)}$ 作张量的运算是正合的。
\end{proof}

\begin{lemma}
\label{algebra-lemma-flat-ses}
设 $0 \to M' \to M \to M'' \to 0$ 为 $R$-模的短正合序列。
若 $M'$ 和 $M''$ 平坦，则 $M$ 也平坦。
若 $M$ 和 $M''$ 平坦，则 $M'$ 也平坦。
\end{lemma}

\begin{proof}
我们使用如下判据：若对于每个理想 $I \subset R$，映射
$N \otimes_R I \to N$ 都是单射，则模 $N$ 平坦，
见引理 \ref{algebra-lemma-flat}。取理想 $I \subset R$。
考虑图表
$$
\begin{matrix}
0 & \to & M' & \to & M & \to & M'' & \to & 0 \\
& & \uparrow & & \uparrow & & \uparrow & & \\
& & M'\otimes_R I & \to & M \otimes_R I & \to & M''\otimes_R I & \to & 0
\end{matrix}
$$
其中两行正合。这立即证明第一个断言。
第二个断言则是因为若 $M''$ 平坦，则由引理
\ref{algebra-lemma-flat-tor-zero}，左下方的水平箭头是单射。
\end{proof}

\begin{lemma}
\label{algebra-lemma-easy-ff}
设 $R$ 为环。设 $M$ 为 $R$-模。以下命题等价：
\begin{enumerate}
\item $M$ 忠实平坦；并且
\item $M$ 平坦，且对所有 $R$-模同态 $\alpha : N \to N'$，
有 $\alpha = 0$ 当且仅当 $\alpha \otimes \text{id}_M = 0$。
\end{enumerate}
\end{lemma}

\begin{proof}
若 $M$ 忠实平坦，则
$0 \to \Ker(\alpha) \to N \to N'$ 正合，当且仅当与 $M$ 作张量后仍然正合。
这证明 (1) 蕴含 (2)。
反之，假设 (2)。设 $N_1 \to N_2 \to N_3$ 为复形，并假设
$N_1 \otimes_R M \to N_2 \otimes_R M \to N_3\otimes_R M$
正合。取 $x \in \Ker(N_2 \to N_3)$，考虑映射
$\alpha : R \to N_2/\Im(N_1)$，
$r \mapsto rx + \Im(N_1)$。由复形 $-\otimes_R M$ 的正合性，
$\alpha \otimes \text{id}_M$ 为零。根据假设 $\alpha$ 为零。
因此 $x$ 属于 $N_1 \to N_2$ 的像。
\end{proof}

\begin{lemma}
\label{algebra-lemma-ff}
\begin{slogan}
平坦模忠实平坦，当且仅当它的纤维非零。
\end{slogan}
设 $M$ 为平坦的 $R$-模。以下命题等价：
\begin{enumerate}
\item $M$ 忠实平坦；
\item 对每个非零 $R$-模 $N$，张量积 $M \otimes_R N$
非零；
\item 对所有 $\mathfrak p \in \Spec(R)$，张量积
$M \otimes_R \kappa(\mathfrak p)$ 非零；并且
\item 对所有 $R$ 的极大理想 $\mathfrak m$，张量积
$M \otimes_R \kappa(\mathfrak m) = M/{\mathfrak m}M$
非零。
\end{enumerate}
\end{lemma}

\begin{proof}
假设 $M$ 忠实平坦且 $N \not = 0$。由引理 \ref{algebra-lemma-easy-ff}，
非零映射 $1 : N \to N$ 诱导非零映射
$M \otimes_R N \to M \otimes_R N$，因此 $M \otimes_R N \not = 0$。
故 (1) 蕴含 (2)。(2) $\Rightarrow$ (3) $\Rightarrow$ (4) 立即成立。

\medskip\noindent
假设 (4)。设 $N_1 \to N_2 \to N_3$ 为复形，且
$N_1 \otimes_R M \to N_2\otimes_R M \to
N_3\otimes_R M$ 正合。令 $H$ 为该复形的上同调，即
$H = \Ker(N_2 \to N_3)/\Im(N_1 \to N_2)$。
为完成证明，我们证明 $H = 0$。由平坦性，$H \otimes_R M = 0$。
取 $x \in H$，令 $I = \{f \in R \mid fx = 0 \}$
为其湮没子。由于 $R/I \subset H$，由 $M$ 的平坦性得到
$M/IM \subset H \otimes_R M = 0$。若 $I \not =  R$，可选取
极大理想 $I \subset \mathfrak m \subset R$。这立即产生矛盾。
\end{proof}

\begin{lemma}
\label{algebra-lemma-ff-rings}
设 $R \to S$ 为平坦环映射。以下命题等价：
\begin{enumerate}
\item $R \to S$ 忠实平坦；
\item 诱导的 $\Spec$ 映射是满射；并且
\item 每个闭点 $x \in \Spec(R)$
都在映射 $\Spec(S) \to \Spec(R)$ 的像中。
\end{enumerate}
\end{lemma}

\begin{proof}
由引理 \ref{algebra-lemma-ff} 很快得到结论，因为我们在注记
\ref{algebra-remark-fundamental-diagram} 中看到，$\mathfrak p$ 在像中，当且仅当环
$S \otimes_R \kappa(\mathfrak p)$ 非零。
\end{proof}

\begin{lemma}
\label{algebra-lemma-local-flat-ff}
局部环之间的平坦局部环同态是忠实平坦的。
\end{lemma}

\begin{proof}
由引理 \ref{algebra-lemma-ff-rings} 立即得到。
\end{proof}

\noindent
平坦性与局部化相容良好。

\begin{lemma}
\label{algebra-lemma-flat-localization}
设 $R$ 为环。设 $S \subset R$ 为乘法子集。
\begin{enumerate}
\item 局部化 $S^{-1}R$ 是平坦的 $R$-代数。
\item 若 $M$ 是 $S^{-1}R$-模，则 $M$ 是平坦的 $R$-模，
当且仅当 $M$ 是平坦的 $S^{-1}R$-模。
\item 设 $M$ 为 $R$-模。则 $M$ 是平坦的 $R$-模，当且仅当对 $R$ 的所有素理想
$\mathfrak p$，$M_{\mathfrak p}$ 是平坦的 $R_{\mathfrak p}$-模。
\item 设 $M$ 为 $R$-模。则 $M$ 是平坦的 $R$-模，当且仅当对 $R$ 的所有极大理想
$\mathfrak m$，$M_{\mathfrak m}$ 是平坦的 $R_{\mathfrak m}$-模。
\item 设 $R \to A$ 为环映射，$M$ 为 $A$-模，且
$g_1, \ldots, g_m \in A$ 为生成 $A$ 的单位理想的元素。则 $M$ 在 $R$ 上平坦，
当且仅当每个局部化 $M_{g_i}$ 在 $R$ 上平坦。
\item 设 $R \to A$ 为环映射，且 $M$ 为 $A$-模。则 $M$ 是平坦的 $R$-模，
当且仅当对 $A$ 的所有素理想 $\mathfrak q$，局部化
$M_{\mathfrak q}$ 是平坦的 $R_{\mathfrak p}$-模
（其中 $\mathfrak p$ 是位于 $\mathfrak q$ 下方的 $R$ 的素理想）。
\item 设 $R \to A$ 为环映射，且 $M$ 为 $A$-模。则 $M$ 是平坦的 $R$-模，
当且仅当对 $A$ 的所有极大理想 $\mathfrak m$，局部化
$M_{\mathfrak m}$ 是平坦的 $R_{\mathfrak p}$-模
（其中 $\mathfrak p = R \cap \mathfrak m$）。
\end{enumerate}
\end{lemma}

\begin{proof}
我们证明引理的最后一个断言。证明中将反复使用局部化是正合的并且与张量积可交换这一事实，
见 \ref{algebra-section-localization} 节和 \ref{algebra-section-tensor-product} 节。

\medskip\noindent
设 $R \to A$ 为环映射，$M$ 为 $A$-模。假设对 $A$ 的所有极大理想
$\mathfrak m$，$M_{\mathfrak m}$ 是平坦的 $R_{\mathfrak p}$-模，其中
$\mathfrak p = R \cap \mathfrak m$。设 $I \subset R$ 为理想。
我们要证明映射 $I \otimes_R M \to M$ 是单射。
可以将其视为 $A$-模的映射。由假设，局部化
$(I \otimes_R M)_{\mathfrak m} \to M_{\mathfrak m}$ 是单射，因为
$(I \otimes_R M)_{\mathfrak m} =
I_{\mathfrak p} \otimes_{R_{\mathfrak p}} M_{\mathfrak m}$。
因此由引理 \ref{algebra-lemma-characterize-zero-local}，$I \otimes_R M \to M$ 的核为零。
于是 $M$ 在 $R$ 上平坦。

\medskip\noindent
反之，假设 $M$ 在 $R$ 上平坦。取 $A$ 中位于 $R$ 的素理想
$\mathfrak p$ 上方的素理想 $\mathfrak q$。设
$I \subset R_{\mathfrak p}$ 为理想。我们要证明
$I \otimes_{R_{\mathfrak p}} M_{\mathfrak q} \to M_{\mathfrak q}$
是单射。可写成 $I = J_{\mathfrak p}$，其中 $J \subset R$ 为理想。
于是映射 $I \otimes_{R_{\mathfrak p}} M_{\mathfrak q} \to M_{\mathfrak q}$
正是映射 $J \otimes_R M \to M$ 在 $\mathfrak q$ 处的局部化，而后者是单射。
由于局部化是正合的，可知 $M_{\mathfrak q}$ 是平坦的
$R_{\mathfrak p}$-模。

\medskip\noindent
这证明了 (7) 和 (6)。其他断言可由最后一个断言直接得到（证明略）。
\end{proof}

\begin{lemma}
\label{algebra-lemma-flat-going-down}
设 $R \to S$ 为平坦映射。设 $\mathfrak p \subset \mathfrak p'$ 为 $R$ 的素理想。
设 $\mathfrak q' \subset S$ 为映到 $\mathfrak p'$ 的 $S$ 的素理想。
则存在素理想 $\mathfrak q \subset \mathfrak q'$，映到 $\mathfrak p$。
\end{lemma}

\begin{proof}
由引理 \ref{algebra-lemma-flat-localization}，局部环映射
$R_{\mathfrak p'} \to S_{\mathfrak q'}$ 平坦。
由引理 \ref{algebra-lemma-local-flat-ff}，该局部环映射忠实平坦。
由引理 \ref{algebra-lemma-ff-rings}，存在一个映到
$\mathfrak p R_{\mathfrak p'}$ 的素理想。该素理想在 $S$ 中的逆像即为所求。
\end{proof}

\noindent
引理中描述的 $R \to S$ 的性质称为“下降性质”。
见定义 \ref{algebra-definition-going-up-down}。

\begin{lemma}
\label{algebra-lemma-colimit-faithfully-flat}
设 $R$ 为环。设 $\{S_i, \varphi_{ii'}\}$ 为忠实平坦的 $R$-代数所成的有向系统。
则 $S = \colim_i S_i$ 是忠实平坦的 $R$-代数。
\end{lemma}

\begin{proof}
由引理 \ref{algebra-lemma-colimit-flat}，可知 $S$ 平坦。
设 $\mathfrak m \subset R$ 为极大理想。由引理
\ref{algebra-lemma-ff-rings}，环 $S_i/\mathfrak m S_i$ 均非零。
因此 $S/\mathfrak mS = \colim S_i/\mathfrak mS_i$ 也非零，
因为 $1$ 不等于零。于是映射 $\Spec(S) \to \Spec(R)$ 的像包含
$\mathfrak m$，由引理 \ref{algebra-lemma-ff-rings} 可知 $R \to S$
忠实平坦。
\end{proof}

\section{支集与湮没子}
\label{algebra-section-supp-and-ann}

\noindent
一些非常基本的定义和引理。

\begin{definition}
\label{algebra-definition-support-module}
设 $R$ 为环，$M$ 为 $R$-模。$M$ 的{\it 支集}是集合
$$
\text{Supp}(M)
=
\{
\mathfrak p \in \Spec(R)
\mid
M_{\mathfrak p} \not = 0
\}
$$
\end{definition}

\begin{lemma}
\label{algebra-lemma-support-zero}
\begin{slogan}
环上的模支集为空，当且仅当它是平凡模。
\end{slogan}
设 $R$ 为环，$M$ 为 $R$-模。则
$$
M = (0) \Leftrightarrow \text{Supp}(M) = \emptyset.
$$
\end{lemma}

\begin{proof}
事实上，引理 \ref{algebra-lemma-characterize-zero-local} 甚至表明：
若 $M$ 非零，则 $\text{Supp}(M)$ 总是包含一个极大理想。
\end{proof}

\begin{definition}
\label{algebra-definition-annihilator}
设 $R$ 为环，$M$ 为 $R$-模。
\begin{enumerate}
\item 给定元素 $m \in M$，$m$ 的{\it 湮没子}是理想
$$
\text{Ann}_R(m) = \text{Ann}(m) = \{f \in R \mid fm = 0\}.
$$
\item $M$ 的{\it 湮没子}是理想
$$
\text{Ann}_R(M) = \text{Ann}(M) = \{f \in R \mid fm = 0\ \forall m \in M\}.
$$
\end{enumerate}
\end{definition}

\begin{lemma}
\label{algebra-lemma-annihilator-flat-base-change}
设 $R \to S$ 为平坦环映射。设 $M$ 为 $R$-模且
$m \in M$。则 $\text{Ann}_R(m) S = \text{Ann}_S(m \otimes 1)$。
若 $M$ 是有限 $R$-模，则
$\text{Ann}_R(M) S = \text{Ann}_S(M \otimes_R S)$。
\end{lemma}

\begin{proof}
令 $I = \text{Ann}_R(m)$。按定义有正合序列
$0 \to I \to R \to M$，其中映射 $R \to M$ 将 $f$ 送到 $fm$。
利用平坦性得到正合序列
$0 \to I \otimes_R S \to S \to M \otimes_R S$，这证明了第一个断言。
若 $m_1, \ldots, m_n$ 是 $M$ 的一组生成元，则
$\text{Ann}_R(M) = \bigcap \text{Ann}_R(m_i)$。类似地，
$\text{Ann}_S(M \otimes_R S) = \bigcap \text{Ann}_S(m_i \otimes 1)$。
令 $I_i = \text{Ann}_R(m_i)$。于是只须证明
$\bigcap_{i = 1, \ldots, n} (I_i S) = (\bigcap_{i = 1, \ldots, n} I_i)S$。
这就是引理 \ref{algebra-lemma-flat-intersect-ideals}。
\end{proof}

\begin{lemma}
\label{algebra-lemma-support-closed}
设 $R$ 为环，$M$ 为 $R$-模。若 $M$ 有限，则
$\text{Supp}(M)$ 是闭集。更准确地，若 $I = \text{Ann}(M)$ 是 $M$ 的湮没子，
则 $V(I) = \text{Supp}(M)$。
\end{lemma}

\begin{proof}
我们证明 $V(I) = \text{Supp}(M)$。

\medskip\noindent
假设 $\mathfrak p \in \text{Supp}(M)$。则 $M_{\mathfrak p} \not = 0$。
取 $m \in M$，使其在 $M_\mathfrak p$ 中的像非零。
由局部化 $M_\mathfrak p$ 的构造，$m$ 的湮没子包含于
$\mathfrak p$。因此当然 $I = \text{Ann}(M)$ 也必须包含于 $\mathfrak p$。

\medskip\noindent
反之，假设 $\mathfrak p \not \in \text{Supp}(M)$。则
$M_{\mathfrak p} = 0$。令 $x_1, \ldots, x_r \in M$ 为生成元。
由引理 \ref{algebra-lemma-localization-colimit}，存在
$f \in R$，$f\not\in \mathfrak p$，使得 $x_i/1 = 0$ 在 $M_f$ 中成立。
因此对某个 $n_i \geq 1$ 有 $f^{n_i} x_i = 0$。
于是当 $n = \max\{n_i\}$ 时有 $f^nM = 0$，结论成立。
\end{proof}

\begin{lemma}
\label{algebra-lemma-support-base-change}
设 $R \to R'$ 为环映射，$M$ 为有限 $R$-模。
则 $\text{Supp}(M \otimes_R R')$ 是 $\text{Supp}(M)$ 的逆像。
\end{lemma}

\begin{proof}
设 $\mathfrak p \in \text{Supp}(M)$。由中山引理
（引理 \ref{algebra-lemma-NAK}）可知
$$
M \otimes_R \kappa(\mathfrak p) = M_\mathfrak p/\mathfrak p M_\mathfrak p
$$
是非零的 $\kappa(\mathfrak p)$-向量空间。
因此对每个位于 $\mathfrak p$ 上方的素理想 $\mathfrak p' \subset R'$，有
$$
(M \otimes_R R')_{\mathfrak p'}/\mathfrak p' (M \otimes_R R')_{\mathfrak p'} =
(M \otimes_R R') \otimes_{R'} \kappa(\mathfrak p') =
M \otimes_R \kappa(\mathfrak p) \otimes_{\kappa(\mathfrak p)}
\kappa(\mathfrak p')
$$
非零。这意味着 $\mathfrak p' \in \text{Supp}(M \otimes_R R')$。
反之，若 $\mathfrak p' \subset R'$ 是位于任意素理想 $\mathfrak p \subset R$ 上方的素理想，
则
$$
(M \otimes_R R')_{\mathfrak p'} =
M_\mathfrak p \otimes_{R_\mathfrak p} R'_{\mathfrak p'}.
$$
因此，若 $\mathfrak p' \in \text{Supp}(M \otimes_R R')$
位于素理想 $\mathfrak p \subset R$ 上方，则
$\mathfrak p \in \text{Supp}(M)$。
\end{proof}

\begin{lemma}
\label{algebra-lemma-support-element}
设 $R$ 为环，$M$ 为 $R$-模，且 $m \in M$。
则 $\mathfrak p \in V(\text{Ann}(m))$，当且仅当
$m$ 不映到 $M_\mathfrak p$ 中的零。
\end{lemma}

\begin{proof}
可用 $Rm \subset M$ 替换 $M$。此时 (1)
$\text{Ann}(m) = \text{Ann}(M)$；(2) $m$ 不映到 $M_\mathfrak p$ 中的零，
当且仅当 $\mathfrak p \in \text{Supp}(M)$。
结论现在由引理 \ref{algebra-lemma-support-closed} 得到。
\end{proof}

\begin{lemma}
\label{algebra-lemma-support-finite-presentation-constructible}
设 $R$ 为环，$M$ 为 $R$-模。若 $M$ 是有限表示的 $R$-模，则
$\text{Supp}(M)$ 是 $\Spec(R)$ 的闭子集，且其补集拟紧。
\end{lemma}

\begin{proof}
取表示
$$
R^{\oplus m} \longrightarrow R^{\oplus n} \longrightarrow M \to 0
$$
令 $A \in \text{Mat}(n \times m, R)$ 为第一个映射的矩阵。
由中山引理 \ref{algebra-lemma-NAK}，有
$$
M_{\mathfrak p} \not = 0 \Leftrightarrow
M \otimes \kappa(\mathfrak p) \not = 0 \Leftrightarrow
\text{rank}(A \bmod \mathfrak p) < n.
$$
因此，若 $I$ 是由 $A$ 的 $n \times n$ 子式生成的 $R$ 的理想，
则 $\text{Supp}(M) = V(I)$。由于 $I$
有限生成，设 $I = (f_1, \ldots, f_t)$，则
$\Spec(R) \setminus V(I)$ 是标准开集 $D(f_i)$ 的有限并，因而拟紧。
\end{proof}

\begin{lemma}
\label{algebra-lemma-support-quotient}
设 $R$ 为环，$M$ 为 $R$-模。
\begin{enumerate}
\item 若 $M$ 有限，则 $M/IM$ 的支集为
$\text{Supp}(M) \cap V(I)$。
\item 若 $N \subset M$，则 $\text{Supp}(N) \subset
\text{Supp}(M)$。
\item 若 $Q$ 是 $M$ 的商模，则 $\text{Supp}(Q) \subset
\text{Supp}(M)$。
\item 若 $0 \to N \to M \to Q \to 0$ 是短正合序列，
则 $\text{Supp}(M) = \text{Supp}(Q) \cup \text{Supp}(N)$。
\end{enumerate}
\end{lemma}

\begin{proof}
函子 $M \mapsto M_{\mathfrak p}$ 是正合的。这立即推出除第一个断言外的所有断言。
对于第一个断言，我们需要证明：若 $M_\mathfrak p \not = 0$ 且
$I \subset \mathfrak p$，则 $(M/IM)_{\mathfrak p}
= M_\mathfrak p/IM_\mathfrak p \not = 0$。这由中山引理
\ref{algebra-lemma-NAK} 得到。
\end{proof}

\section{上升与下降}
\label{algebra-section-going-up}

\noindent
设 $\mathfrak p$、$\mathfrak p'$ 为环 $R$ 的素理想。
令 $X = \Spec(R)$，赋予扎里斯基拓扑。记 $x \in X$ 为对应于
$\mathfrak p$ 的点，记 $x' \in X$ 为对应于
$\mathfrak p'$ 的点。则有：
$$
x' \leadsto x \Leftrightarrow \mathfrak p' \subset \mathfrak p.
$$
换言之，$x$ 是 $x'$ 的一个特殊化，当且仅当
$\mathfrak p' \subset \mathfrak p$。
术语和记号见拓扑一章的 \ref{topology-section-specialization} 节。

\begin{definition}
\label{algebra-definition-going-up-down}
设 $\varphi : R \to S$ 为环映射。
\begin{enumerate}
\item 若给定 $R$ 中的素理想 $\mathfrak p \subset \mathfrak p'$，
以及 $S$ 中位于 $\mathfrak p$ 上方的素理想 $\mathfrak q$，存在 $S$ 中的素理想
$\mathfrak q'$，满足 (a) $\mathfrak q \subset \mathfrak q'$，且 (b)
$\mathfrak q'$ 位于 $\mathfrak p'$ 上方，则称
$\varphi : R \to S$ 满足{\it 上升}。
\item 若给定 $R$ 中的素理想 $\mathfrak p \subset \mathfrak p'$，
以及 $S$ 中位于 $\mathfrak p'$ 上方的素理想 $\mathfrak q'$，存在 $S$ 中的素理想
$\mathfrak q$，满足 (a) $\mathfrak q \subset \mathfrak q'$，且 (b)
$\mathfrak q$ 位于 $\mathfrak p$ 上方，则称
$\varphi : R \to S$ 满足{\it 下降}。
\end{enumerate}
\end{definition}

\noindent
到目前为止，我们已经见到以下情形：
\begin{enumerate}
\item 整环映射满足上升，见引理 \ref{algebra-lemma-integral-going-up}。
\item 作为特例，有限环映射满足上升。
\item 作为特例，商映射 $R \to R/I$ 满足上升。
\item 平坦环映射满足下降，见引理 \ref{algebra-lemma-flat-going-down}。
\item 作为特例，任意局部化满足下降。
\item 若 $R \subset S$ 为整环的扩张，$R$ 正规且 $S$ 整于 $R$，
则满足下降，见命题 \ref{algebra-proposition-going-down-normal-integral}。
\end{enumerate}
下面是下降成立的另一种情形。

\begin{lemma}
\label{algebra-lemma-open-going-down}
设 $R \to S$ 为环映射。若诱导映射
$\varphi : \Spec(S) \to \Spec(R)$ 是开映射，则
$R \to S$ 满足下降。
\end{lemma}

\begin{proof}
设 $\mathfrak p \subset \mathfrak p' \subset R$，且
$\mathfrak q' \subset S$ 位于 $\mathfrak p'$ 上方。由于 $\varphi$ 是开映射，
对每个 $g \in S$，$g \not \in \mathfrak q'$，都有
$\mathfrak p$ 属于 $D(g) \subset \Spec(S)$ 的像。换言之，
$S_g \otimes_R \kappa(\mathfrak p)$ 非零。由于 $S_{\mathfrak q'}$
是这些 $S_g$ 的有向余极限，这意味着
$S_{\mathfrak q'} \otimes_R \kappa(\mathfrak p)$ 非零，见
引理 \ref{algebra-lemma-localization-colimit} 和
\ref{algebra-lemma-tensor-products-commute-with-limits}。
因此 $\mathfrak p$ 属于
$\Spec(S_{\mathfrak q'}) \to \Spec(R)$ 的像，正是所需结论。
\end{proof}

\begin{lemma}
\label{algebra-lemma-going-up-down-specialization}
设 $R \to S$ 为环映射。
\begin{enumerate}
\item $R \to S$ 满足下降，当且仅当一般化沿映射
$\Spec(S) \to \Spec(R)$ 提升，见拓扑一章定义
\ref{topology-definition-lift-specializations}。
\item $R \to S$ 满足上升，当且仅当特殊化沿映射
$\Spec(S) \to \Spec(R)$ 提升，见拓扑一章定义
\ref{topology-definition-lift-specializations}。
\end{enumerate}
\end{lemma}

\begin{proof}
略。
\end{proof}

\begin{lemma}
\label{algebra-lemma-going-up-down-composition}
设 $R \to S$ 和 $S \to T$ 为满足下降的环映射。
则 $R \to T$ 也满足下降。上升的情形同样成立。
\end{lemma}

\begin{proof}
根据引理 \ref{algebra-lemma-going-up-down-specialization}，结论由拓扑一章的引理
\ref{topology-lemma-lift-specialization-composition} 得到。
\end{proof}

\begin{lemma}
\label{algebra-lemma-image-stable-specialization-closed}
设 $R \to S$ 为环映射。令 $T \subset \Spec(R)$
为 $\Spec(S)$ 的像。若 $T$ 在特殊化下稳定，则 $T$ 是闭集。
\end{lemma}

\begin{proof}
我们给出两个证明。

\medskip\noindent
第一个证明。设 $\mathfrak p \subset R$ 为素理想，且
$\Spec(R)$ 中对应的点属于 $T$ 的闭包。这意味着对于每个
$f \in R$，$f \not \in \mathfrak p$，都有
$D(f) \cap T \not = \emptyset$。注意 $D(f) \cap T$
是 $\Spec(S_f)$ 在 $\Spec(R)$ 中的像。因此 $S_f \not = 0$。
换言之，在环 $S_f$ 中 $1 \not = 0$。由于 $S_{\mathfrak p}$ 是环 $S_f$ 的有向余极限，
可知在 $S_{\mathfrak p}$ 中 $1 \not = 0$。
换言之，$S_{\mathfrak p} \not = 0$；考察
$\Spec(S_{\mathfrak p}) \to \Spec(S) \to \Spec(R)$ 的像，
可知存在 $\mathfrak p' \in T$，满足 $\mathfrak p' \subset \mathfrak p$。
由于假设 $T$ 在特殊化下闭，得到 $\mathfrak p$ 是 $T$ 中的点，命题得证。

\medskip\noindent
第二个证明。令 $I = \Ker(R \to S)$。可用 $R/I$ 替换 $R$。
此时环映射 $R \to S$ 是单射。由引理
\ref{algebra-lemma-injective-minimal-primes-in-image}，$R$ 的所有极小素理想
都包含于像 $T$ 中。因此若 $T$ 在特殊化下稳定，则它包含所有素理想。
\end{proof}

\begin{lemma}
\label{algebra-lemma-going-up-closed}
设 $R \to S$ 为环映射。以下命题等价：
\begin{enumerate}
\item $R \to S$ 满足上升；并且
\item 映射 $\Spec(S) \to \Spec(R)$ 是闭映射。
\end{enumerate}
\end{lemma}

\begin{proof}
一般而言，拓扑空间的闭映射会使特殊化沿映射提升，见拓扑一章的引理
\ref{topology-lemma-closed-open-map-specialization}。
因此第二个条件推出第一个条件。

\medskip\noindent
假设 $R \to S$ 满足上升。设 $V(I) \subset \Spec(S)$ 为闭集。
我们要证明 $V(I)$ 在 $\Spec(R)$ 中的像是闭集。
环映射 $S \to S/I$ 显然满足上升。因此由引理
\ref{algebra-lemma-going-up-down-composition}，$R \to S \to S/I$ 满足上升。
将 $S$ 替换为 $S/I$ 后，只须证明 $\Spec(S)$ 在 $\Spec(R)$ 中的像
$T$ 是闭集。由拓扑一章引理
\ref{topology-lemma-open-closed-specialization} 和
\ref{topology-lemma-lift-specializations-images}，该像在特殊化下稳定。
于是结论由引理 \ref{algebra-lemma-image-stable-specialization-closed} 得到。
\end{proof}

\begin{lemma}
\label{algebra-lemma-constructible-stable-specialization-closed}
设 $R$ 为环，设 $E \subset \Spec(R)$ 为可构造子集。
\begin{enumerate}
\item 若 $E$ 在特殊化下稳定，则 $E$ 是闭集。
\item 若 $E$ 在一般化下稳定，则 $E$ 是开集。
\end{enumerate}
\end{lemma}

\begin{proof}
第一个证明。第一个断言由引理
\ref{algebra-lemma-image-stable-specialization-closed} 与引理
\ref{algebra-lemma-constructible-is-image} 结合得到。第二个断言成立是因为可构造集的补集
仍然可构造（见拓扑一章引理 \ref{topology-lemma-constructible}），
再应用本引理的第一部分以及拓扑一章引理
\ref{topology-lemma-open-closed-specialization}。

\medskip\noindent
第二个证明。

由于根据引理 \ref{algebra-lemma-spec-spectral}，$\Spec(R)$ 是谱空间，
这只是拓扑一章引理
\ref{topology-lemma-constructible-stable-specialization-closed} 的一个特例。
\end{proof}

\begin{proposition}
\label{algebra-proposition-fppf-open}
设 $R \to S$ 是平坦且有限表示的环映射。
则 $\Spec(S) \to \Spec(R)$ 是开映射。
更一般地，任意有限表示且满足下降性质的环映射
$R \to S$ 都具有此性质。
\end{proposition}

\begin{proof}
若 $R \to S$ 平坦，则由引理
\ref{algebra-lemma-flat-going-down}，$R \to S$ 满足下降性质。
因此，为证明命题，只须假设 $R \to S$ 有限表示且满足下降性质。

\medskip\noindent
由于 $g \in S$ 的标准开集 $D(g) \subset \Spec(S)$
构成拓扑的一个基，只须证明 $D(g)$ 的像是开集。
回忆（引理 \ref{algebra-lemma-standard-open}），
$\Spec(S_g) \to \Spec(S)$ 将 $\Spec(S_g)$ 同胚映到 $D(g)$。
由于 $S \to S_g$ 满足下降性质（见上文），由引理
\ref{algebra-lemma-going-up-down-composition} 可知 $R \to S_g$
满足下降性质。因此将 $S$ 替换为 $S_g$ 后，只须证明该像是开集。
由切瓦莱定理（定理 \ref{algebra-theorem-chevalley}），该像是可构造集 $E$。
又由于 $R \to S$ 满足下降性质，$E$ 在一般化下稳定，见拓扑一章引理
\ref{topology-lemma-open-closed-specialization} 和
\ref{topology-lemma-lift-specializations-images}。
因此由引理 \ref{algebra-lemma-constructible-stable-specialization-closed}，
$E$ 是开集。
\end{proof}

\begin{lemma}
\label{algebra-lemma-same-image}
设 $k$ 为域，$R$、$S$ 为 $k$-代数。
设 $S' \subset S$ 为 $k$-子代数，且 $f \in S' \otimes_k R$。
在交换图表
$$
\xymatrix{
\Spec((S \otimes_k R)_f) \ar[rd] \ar[rr] & &
\Spec((S' \otimes_k R)_f) \ar[ld] \\
& \Spec(R) &
}
$$
中，两条对角箭头的像相同。
\end{lemma}

\begin{proof}
设 $\mathfrak p \subset R$ 属于西南箭头的像。这意味着（引理
\ref{algebra-lemma-in-image}）
$$
(S' \otimes_k R)_f \otimes_R \kappa(\mathfrak p)
=
(S' \otimes_k \kappa(\mathfrak p))_f
$$
不是零环，即 $S' \otimes_k \kappa(\mathfrak p)$ 不是零环，且 $f$
在其中的像不是幂零元。环映射
$$
S' \otimes_k \kappa(\mathfrak p) \to S \otimes_k \kappa(\mathfrak p)
$$
是单射。因此 $S \otimes_k \kappa(\mathfrak p)$ 也不是零环，且 $f$
在其中的像不是幂零元。于是
$$
(S \otimes_k R)_f \otimes_R \kappa(\mathfrak p)
$$
不是零环。因此（引理 \ref{algebra-lemma-in-image}），$\mathfrak p$ 属于东南箭头的像，
这正是所需结论。
\end{proof}

\begin{lemma}
\label{algebra-lemma-map-into-tensor-algebra-open}
设 $k$ 为域。
设 $R$ 和 $S$ 为 $k$-代数。
映射 $\Spec(S \otimes_k R) \to \Spec(R)$
是开映射。
\end{lemma}

\begin{proof}
设 $f \in S \otimes_k R$。只须证明标准开集 $D(f)$ 的像是开集。
取一个有限型 $k$-子代数 $S' \subset S$，使得
$f \in S' \otimes_k R$。映射 $R \to S' \otimes_k R$ 平坦且有限表示，
所以由命题 \ref{algebra-proposition-fppf-open}，
$$
\Spec((S' \otimes_k R)_f) \to \Spec(R)
$$
的像 $U$ 是开集。由引理 \ref{algebra-lemma-same-image}，这也正是 $D(f)$ 的像，
命题得证。
\end{proof}

\noindent
下面是一个有时很有用的棘手引理。

\begin{lemma}
\label{algebra-lemma-unique-prime-over-localize-below}
设 $R \to S$ 为环映射。
设 $\mathfrak p \subset R$ 为素理想。假设
\begin{enumerate}
\item 在 $\mathfrak p$ 上方存在唯一的素理想
$\mathfrak q \subset S$，并且
\item 下列两者之一成立：
\begin{enumerate}
\item $R \to S$ 满足上升性质，或
\item $R \to S$ 满足下降性质，且 $R$ 的每个素理想上方至多有一个
素理想属于 $S$。
\end{enumerate}
\end{enumerate}
则 $S_{\mathfrak p} = S_{\mathfrak q}$。
\end{lemma}

\begin{proof}
考虑任意一个素理想 $\mathfrak q' \subset S$，它对应于
$\Spec(S_{\mathfrak p})$ 中的一个点。这意味着
$\mathfrak p' = R \cap \mathfrak q'$ 包含于 $\mathfrak p$。
图示如下：
$$
\xymatrix{
\mathfrak q' \ar@{-}[d] \ar@{-}[r] & ? \ar@{-}[r] \ar@{-}[d] & S \ar@{-}[d] \\
\mathfrak p' \ar@{-}[r] & \mathfrak p \ar@{-}[r] & R
}
$$
假设 (1) 和 (2)(a)。由上升性质，存在素理想 $\mathfrak q'' \subset S$，
使得 $\mathfrak q' \subset \mathfrak q''$ 且 $\mathfrak q''$ 位于
$\mathfrak p$ 上方。由 $\mathfrak q$ 的唯一性，得到
$\mathfrak q'' = \mathfrak q$。换言之，$\mathfrak q'$ 定义了
$\Spec(S_{\mathfrak q})$ 中的一个点。

\medskip\noindent
假设 (1) 和 (2)(b)。由下降性质，存在素理想
$\mathfrak q'' \subset \mathfrak q$，位于 $\mathfrak p'$ 上方。
由位于 $\mathfrak p'$ 上方的素理想的唯一性，得到
$\mathfrak q' = \mathfrak q''$。换言之，$\mathfrak q'$ 定义了
$\Spec(S_{\mathfrak q})$ 中的一个点。

\medskip\noindent
在两种情形中，我们都得到双射
$$
\Spec(S_{\mathfrak q}) \to \Spec(S_{\mathfrak p})
$$
。
显然，这意味着 $S - \mathfrak q$ 中的所有元素都在 $S_{\mathfrak p}$ 中可逆，
换言之 $S_{\mathfrak p} = S_{\mathfrak q}$。
\end{proof}

\noindent
下面的引理是平坦环映射的下降性质的一个推广。

\begin{lemma}
\label{algebra-lemma-going-down-flat-module}
设 $R \to S$ 为环映射。设 $N$ 为有限生成的 $S$-模，且作为 $R$-模平坦。
在 $\text{Supp}(N) \subset \Spec(S)$ 上赋予诱导拓扑。
则一般化沿映射 $\text{Supp}(N) \to \Spec(R)$ 提升。
\end{lemma}

\begin{proof}
该陈述的含义如下。设 $\mathfrak p \subset \mathfrak p' \subset R$ 为素理想。
设 $\mathfrak q' \subset S$ 为素理想，且 $\mathfrak q' \in \text{Supp}(N)$。
则存在素理想 $\mathfrak q \subset \mathfrak q'$，满足
$\mathfrak q \in \text{Supp}(N)$ 且位于 $\mathfrak p$ 上方。
由于 $N$ 在 $R$ 上平坦，由引理 \ref{algebra-lemma-flat-localization} 可知
$N_{\mathfrak q'}$ 在 $R_{\mathfrak p'}$ 上平坦。
由于 $N_{\mathfrak q'}$ 是 $S_{\mathfrak q'}$ 上的有限模，且因
$\mathfrak q' \in \text{Supp}(N)$ 而非零，由 Nakayama 引理
\ref{algebra-lemma-NAK} 可知
$$
N_{\mathfrak q'} \otimes_{S_{\mathfrak q'}} \kappa(\mathfrak q')
$$
非零。因此
$$
N_{\mathfrak q'} \otimes_{R_{\mathfrak p'}} \kappa(\mathfrak p')
$$
也非零。由引理 \ref{algebra-lemma-ff} 得到
$$
N_{\mathfrak q'} \otimes_{R_{\mathfrak p'}} \kappa(\mathfrak p)
$$
非零。令
$$
J \subset S_{\mathfrak q'} \otimes_{R_{\mathfrak p'}} \kappa(\mathfrak p)
$$
为有限非零模
$$
N_{\mathfrak q'} \otimes_{R_{\mathfrak p'}} \kappa(\mathfrak p)
$$
的湮没子。由于 $J$ 是真理想，我们可以选取素理想 $\mathfrak q \subset S$，
使其对应于
$$
S_{\mathfrak q'} \otimes_{R_{\mathfrak p'}} \kappa(\mathfrak p)/J
$$
的一个素理想。该素理想属于 $N$ 的支集，位于 $\mathfrak p$ 上方，且包含于
$\mathfrak q'$，正是所需结论。
\end{proof}

\section{可分扩张}
\label{algebra-section-separability}

\noindent
本节讨论非代数域扩张的可分性。这与几何约化代数的概念密切相关，见定义
\ref{algebra-definition-geometrically-reduced}。

\begin{definition}
\label{algebra-definition-separable-field-extension}
设 $K/k$ 为域扩张。
\begin{enumerate}
\item 若存在 $K/k$ 的超越基 $\{x_i; i \in I\}$，使得扩张
$K/k(x_i; i \in I)$ 是可分代数扩张，则称 $K$ 在 $k$ 上{\it 可分生成}。
\item 若对于每个子扩张 $k \subset K' \subset K$，其中 $K'$ 在 $k$ 上有限生成，
扩张 $K'/k$ 都是可分生成的，则称 $K$ 在 $k$ 上{\it 可分}。
\end{enumerate}
\end{definition}

\noindent
由这个略显笨拙的定义，尚不清楚一个可分生成的域扩张本身是否可分。
下面将证明确实如此，见引理 \ref{algebra-lemma-separably-generated-separable}。

\begin{lemma}
\label{algebra-lemma-subextensions-are-separable}
设 $K/k$ 为可分域扩张。对于任意子扩张 $K/K'/k$，域扩张
$K'/k$ 都是可分的。
\end{lemma}

\begin{proof}
这直接由定义得到。
\end{proof}

\begin{lemma}
\label{algebra-lemma-generating-finitely-generated-separable-field-extensions}
设 $K/k$ 是一个可分生成且有限生成的域扩张。
令 $r = \text{trdeg}_k(K)$。则存在 $K$ 中的元素
$x_1, \ldots, x_{r + 1}$，满足
\begin{enumerate}
\item $x_1, \ldots, x_r$ 是 $K$ 在 $k$ 上的超越基；
\item $K = k(x_1, \ldots, x_{r + 1})$；
\item $x_{r + 1}$ 在 $k(x_1, \ldots, x_r)$ 上可分。
\end{enumerate}
\end{lemma}

\begin{proof}
将定义与域一章引理 \ref{fields-lemma-primitive-element} 结合即可。
\end{proof}

\begin{lemma}
\label{algebra-lemma-make-separably-generated}
设 $K/k$ 为有限生成域扩张。存在图表
$$
\xymatrix{
K \ar[r] & K' \\
k \ar[u] \ar[r] & k' \ar[u]
}
$$
其中 $k'/k$、$K'/K$ 是有限纯不可分域扩张，且 $K'/k'$ 是可分生成域扩张。
\end{lemma}

\begin{proof}
只有在 $k$ 的特征为 $p > 0$ 时该引理才有意义。
取 $x_1, \ldots, x_r$ 为 $K$ 在 $k$ 上的超越基。
由于 $K$ 在 $k$ 上有限生成，扩张
$k(x_1, \ldots, x_r) \subset K$ 是有限扩张。
令 $K/K_{sep}/k(x_1, \ldots, x_r)$ 为域一章引理
\ref{fields-lemma-separable-first} 中给出的子扩张。
若 $K = K_{sep}$，则结论成立。我们对
$d = [K : K_{sep}]$ 使用归纳法。

\medskip\noindent
假设 $d > 1$。取 $\beta \in K$，使得
$\alpha = \beta^p \in K_{sep}$ 且 $\beta \not \in K_{sep}$。
令 $P = T^n + a_1T^{n - 1} + \ldots + a_n$
为 $\alpha$ 在 $k(x_1, \ldots, x_r)$ 上的最小多项式。
令 $k'/k$ 为通过添加 $p$ 次根得到的有限纯不可分扩张，
使得每个 $a_i$ 都是
$k'(x_1^{1/p}, \ldots, x_r^{1/p})$ 中的 $p$ 次幂。
这样的扩张存在；细节略去。
令 $L$ 为适合下图的域
$$
\xymatrix{
K \ar[r] & L \\
k(x_1, \ldots, x_r) \ar[u] \ar[r] & k'(x_1^{1/p}, \ldots, x_r^{1/p}) \ar[u]
}
$$
我们可以并且确实假设 $L$ 是 $K$ 与
$k'(x_1^{1/p}, \ldots, x_r^{1/p})$ 的复合域。
令 $L/L_{sep}/k'(x_1^{1/p}, \ldots, x_r^{1/p})$
为域一章引理 \ref{fields-lemma-separable-first} 中给出的子扩张。
则 $L_{sep}$ 是 $K_{sep}$ 与
$k'(x_1^{1/p}, \ldots, x_r^{1/p})$ 的复合域。
元素 $\alpha \in L_{sep}$ 是多项式 $P$ 的一个根；该多项式的所有系数
都是 $k'(x_1^{1/p}, \ldots, x_r^{1/p})$ 中的 $p$ 次幂，且其根两两不同。
由域一章引理 \ref{fields-lemma-pth-root}，可知存在某个
$\alpha' \in L_{sep}$，使得 $\alpha = (\alpha')^p$。
显然，这意味着 $\beta$ 映到 $\alpha' \in L_{sep}$。
换言之，我们得到域塔
$$
\xymatrix{
K \ar[r] & L \\
K_{sep}(\beta) \ar[r] \ar[u] & L_{sep} \ar[u] \\
K_{sep} \ar[r] \ar[u] & L_{sep} \ar@{=}[u] \\
k(x_1, \ldots, x_r) \ar[u] \ar[r] & k'(x_1^{1/p}, \ldots, x_r^{1/p}) \ar[u] \\
k \ar[r] \ar[u] & k' \ar[u]
}
$$
因此该构造导致新的情形
$[L : L_{sep}] < [K : K_{sep}]$。由归纳假设，存在
$k' \subset k''$ 和 $L \subset L'$，使其对于扩张 $L/k'$ 满足本引理的结论。
那么，对于扩张 $K/k$，扩张 $k''/k$ 和 $L'/K$ 就满足要求。
引理得证。
\end{proof}

\section{几何约化代数}
\label{algebra-section-geometrically-reduced}

\noindent
关于几何约化代数的主要结果是引理
\ref{algebra-lemma-geometrically-reduced-finite-purely-inseparable-extension}。
建议读者读完定义后跳到该引理。

\begin{definition}
\label{algebra-definition-geometrically-reduced}
设 $k$ 为域，$S$ 为 $k$-代数。
若对于每个域扩张 $K/k$，$K$-代数
$K \otimes_k S$ 都是约化的，则称 $S$ 在 $k$ 上{\it 几何约化}。
\end{definition}

\noindent
设 $k$ 为域，且 $S$ 为约化的 $k$-代数。要检验 $S$ 是几何约化的，
只须检验 $\overline{k} \otimes_k S$ 是约化的（其中 $\overline{k}$ 表示
$k$ 的代数闭包）。事实上，只须对有限纯不可分域扩张 $k'/k$ 检验即可。见
引理 \ref{algebra-lemma-geometrically-reduced-finite-purely-inseparable-extension}。

\begin{lemma}
\label{algebra-lemma-subalgebra-separable}
几何约化代数的初等性质。设 $k$ 为域，$S$ 为 $k$-代数。
\begin{enumerate}
\item 若 $S$ 在 $k$ 上几何约化，则其每个 $k$-子代数也几何约化。
\item 若 $S$ 的所有有限生成 $k$-子代数都是几何约化的，则 $S$ 也是几何约化的。
\item 几何约化 $k$-代数的有向余极限也是几何约化的。
\item 若 $S$ 在 $k$ 上几何约化，则 $S$ 的任意局部化也在 $k$ 上几何约化。
\end{enumerate}
\end{lemma}

\begin{proof}
略。第二和第三个性质由张量积与余极限交换这一事实得到。
\end{proof}

\begin{lemma}
\label{algebra-lemma-geometrically-reduced-permanence}
设 $k$ 为域。若 $R$ 在 $k$ 上几何约化，且 $S \subset R$ 为乘法子集，
则局部化 $S^{-1}R$ 在 $k$ 上几何约化。
若 $R$ 在 $k$ 上几何约化，则 $R[x]$ 在 $k$ 上几何约化。
\end{lemma}

\begin{proof}
略。提示：约化环的局部化仍是约化的，且局部化与张量积交换。
\end{proof}

\noindent
在下面引理的证明中，我们会反复使用以下观察：若
$R' \subset R$ 且 $S' \subset S$ 是 $k$-代数的包含，则映射
$R' \otimes_k S' \to R \otimes_k S$ 是单射。

\begin{lemma}
\label{algebra-lemma-limit-argument}
设 $k$ 为域，$R$、$S$ 为 $k$-代数。
\begin{enumerate}
\item 若 $R \otimes_k S$ 非约化，则存在有限生成子代数 $R' \subset R$、
$S' \subset S$，使得 $R' \otimes_k S'$ 非约化。
\item 若 $R \otimes_k S$ 含有非零零因子，则存在有限生成子代数
$R' \subset R$、$S' \subset S$，使得 $R' \otimes_k S'$ 含有非零零因子。
\item 若 $R \otimes_k S$ 含有非平凡幂等元，则存在有限生成子代数
$R' \subset R$、$S' \subset S$，使得 $R' \otimes_k S'$ 含有非平凡幂等元。
\end{enumerate}
\end{lemma}

\begin{proof}
设 $z \in R \otimes_k S$ 为幂零元。可写成
$z = \sum_{i = 1, \ldots, n} x_i \otimes y_i$。
因此可以取 $R'$ 为 $x_i$ 生成的 $k$-子代数，取 $S'$ 为 $y_i$
生成的 $k$-子代数。第二和第三个断言的证明相同。
\end{proof}

\begin{lemma}
\label{algebra-lemma-geometrically-reduced-any-reduced-base-change}
令 $k$ 为域。
令 $S$ 为几何约化的 $k$-代数。
令 $R$ 为任意约化的 $k$-代数。
则 $R \otimes_k S$ 是约化的。
\end{lemma}

\begin{proof}
由引理 \ref{algebra-lemma-limit-argument}，我们可以假设 $R$ 是 $k$ 上的有限型代数。
于是，作为约化的诺特环，$R$ 嵌入域的有限乘积中
（见引理 \ref{algebra-lemma-total-ring-fractions-no-embedded-points}、
\ref{algebra-lemma-Noetherian-irreducible-components} 和
\ref{algebra-lemma-minimal-prime-reduced-ring}）。
因此我们可以假设 $R$ 是域的有限乘积。在这种情形下，由定义
\ref{algebra-definition-geometrically-reduced} 可知 $R \otimes_k S$ 是约化的。
\end{proof}

\begin{lemma}
\label{algebra-lemma-separable-extension-preserves-reducedness}
令 $k$ 为域。
令 $S$ 为约化的 $k$-代数。
令 $K/k$ 或为可分域扩张，或为可分生成的域扩张。
则 $K \otimes_k S$ 是约化的。
\end{lemma}

\begin{proof}
假设 $k \subset K$ 为可分扩张。
由引理 \ref{algebra-lemma-limit-argument}，我们可以假设 $S$ 是 $k$ 上的有限型代数，
且 $K$ 在 $k$ 上有限生成。
于是 $S$ 嵌入域的有限乘积中，即其全分式环（见引理
\ref{algebra-lemma-minimal-prime-reduced-ring} 和
\ref{algebra-lemma-total-ring-fractions-no-embedded-points}）。
因此实际上可以假设 $S$ 是整环。
按照引理 \ref{algebra-lemma-generating-finitely-generated-separable-field-extensions}，
取 $x_1, \ldots, x_{r + 1} \in K$。
令 $P \in k(x_1, \ldots, x_r)[T]$ 为 $x_{r + 1}$ 的最小多项式。它是可分多项式。
容易看出
$k[x_1, \ldots, x_r] \otimes_k S = S[x_1, \ldots, x_r]$ 是整环。
这推出 $k(x_1, \ldots, x_r) \otimes_k S$ 是整环，
因为它是 $S[x_1, \ldots, x_r]$ 的局部化。
环扩张 $k(x_1, \ldots, x_r) \otimes_k S \subset K \otimes_k S$
由单个元素 $x_{r + 1}$ 以及一个方程（即 $P$）生成。因此 $K \otimes_k S$ 嵌入
$F[T]/(P)$，其中 $F$ 是 $k(x_1, \ldots, x_r) \otimes_k S$ 的分式域。
由于 $P$ 可分，这个环是域的有限乘积，结论成立。

\medskip\noindent
此时我们尚不知道可分生成的域扩张是可分的，
所以还必须在这种情形下证明该引理。
为此，设 $\{x_i\}_{i \in I}$ 是 $K$ 在 $k$ 上的分离超越基。
对于任意有限集的元素 $\lambda_j \in K$，存在 $T \subset I$ 的有限子集，使得
$k(\{x_i\}_{i\in T}) \subset k(\{x_i\}_{i \in T} \cup \{\lambda_j\})$
是有限可分扩张。因此 $K$ 是 $k$ 上有限生成且可分生成扩张的有向余极限。
于是前一段的论证同样适用于此情形。
\end{proof}

\begin{lemma}
\label{algebra-lemma-generic-points-geometrically-reduced}
令 $k$ 为域，$S$ 为 $k$-代数。假设 $S$ 是约化的，且 $S_{\mathfrak p}$ 是几何
约化的，其中 $\mathfrak p$ 遍历 $S$ 的每个极小素理想。
则 $S$ 是几何约化的。
\end{lemma}

\begin{proof}
由于 $S$ 是约化的，映射
$S \to \prod_{\mathfrak p\text{ minimal}} S_{\mathfrak p}$
是单射，见引理 \ref{algebra-lemma-reduced-ring-sub-product-fields}。
若 $K/k$ 是域扩张，则映射
$$
S \otimes_k K \to (\prod S_\mathfrak p) \otimes_k K \to
\prod S_\mathfrak p \otimes_k K
$$
是单射：第一个映射因为 $k \to K$ 是平坦的，第二个直接可见，
因为 $K$ 是自由的 $k$-模。由于 $S_\mathfrak p$ 是几何约化的，
右侧的环是约化的。因此 $S \otimes_k K$ 是约化环的子环，故为约化的。
\end{proof}

\begin{lemma}
\label{algebra-lemma-separable-algebraic-diagonal}
令 $k'/k$ 为可分代数扩张。
则存在乘法子集 $S \subset k' \otimes_k k'$，使得乘法映射
$k' \otimes_k k' \to k'$ 被识别为
$k' \otimes_k k' \to S^{-1}(k' \otimes_k k')$。
\end{lemma}

\begin{proof}
先假设 $k'/k$ 是有限可分扩张。则 $k' = k(\alpha)$，
见《域》，引理 \ref{fields-lemma-primitive-element}。
令 $P \in k[x]$ 为 $\alpha$ 在 $k$ 上的最小多项式。
则 $P$ 是不可约、可分且首一的多项式，见《域》，
小节 \ref{fields-section-separable-extensions}。
于是 $k'[x]/(P) \to k' \otimes_k k'$，
$\sum \alpha_i x^i \mapsto \alpha_i \otimes \alpha^i$ 是同构。
我们可以在 $k'[x]$ 中分解 $P = (x - \alpha) Q$，由于 $P$ 可分，
可知 $Q(\alpha) \not = 0$。
于是显然，乘法子集 $S'$ 由
$Q$ 在 $k'[x]/(P)$ 中生成并起作用，即
$k' = (S')^{-1}(k'[x]/(P))$。
通过结构传递，$S'$ 在 $k' \otimes_k k'$ 中的像 $S$ 也起作用。

\medskip\noindent
在一般情形下，写 $k' = \bigcup k_i$ 为其有限子域扩张（在 $k$ 上）的并。
对每个 $i$，存在乘法子集 $S_i \subset k_i \otimes_k k_i$，使得
$k_i = S_i^{-1}(k_i \otimes_k k_i)$。令
$S$ 为 $\bigcup S_i \subset k' \otimes_k k'$ 的乘法闭包。
显然，乘法映射将 $S$ 的每个元素映为 $k'$ 中的可逆元素。
利用 $k' \otimes_k k'$ 是环 $k_i \otimes_k k_i$ 的并，可知乘法映射核中的
每个元素都被映为 $S^{-1}(k' \otimes_k k')$ 中的零。
利用局部化的正合性，结论得证。
\end{proof}

\begin{lemma}
\label{algebra-lemma-geometrically-reduced-over-separable-algebraic}
令 $k'/k$ 为可分代数域扩张。
令 $A$ 为 $k'$-代数。则 $A$ 在
$k$ 上几何约化，当且仅当它在 $k'$ 上几何约化。
\end{lemma}

\begin{proof}
假设 $A$ 在 $k'$ 上几何约化。
令 $K/k$ 为域扩张。由引理
\ref{algebra-lemma-separable-extension-preserves-reducedness}，$K \otimes_k k'$ 是约化环。
因此由引理 \ref{algebra-lemma-geometrically-reduced-any-reduced-base-change}，可知
$K \otimes_k A = (K \otimes_k k') \otimes_{k'} A$ 是约化的。

\medskip\noindent
假设 $A$ 在 $k$ 上几何约化。令 $K/k'$ 为域扩张。则
$$
K \otimes_{k'} A = (K \otimes_k A) \otimes_{(k' \otimes_k k')} k'
$$
由于由引理 \ref{algebra-lemma-separable-algebraic-diagonal}，
$k' \otimes_k k' \to k'$ 是局部化，故可知 $K \otimes_{k'} A$
是约化代数的局部化，因此是约化的。
\end{proof}

\section{可分扩张（续）}
\label{algebra-section-separability-continued}

\noindent
本节继续讨论小节 \ref{algebra-section-separability} 中开始的内容。

\begin{lemma}
\label{algebra-lemma-mini-separability}
令 $k$ 为特征为 $p > 1$ 的域。令 $K/k$ 为由
$x_1, \ldots, x_{n + 1} \in K$ 生成的域扩张，满足
\begin{enumerate}
\item $\{x_1, \ldots, x_n\}$ 是 $K/k$ 的超越基，
\item 对于每个 $k$-线性无关的子集
$\{a_1, \ldots, a_m\}$（其中 $K$），集合
$\{a^p_1, \ldots, a_m^p\}$ 是 $k$-线性无关的。
\end{enumerate}
则存在 $1 \leq j \leq n+1$，使得
$\{ x_1, \ldots, \widehat{x}_j, \ldots, x_{n+1}\}$
是 $K / k$ 的分离超越基。
\end{lemma}

\begin{proof}
由假设，$x_{n + 1}$ 在 $k(x_1, \ldots, x_n)$ 上代数，
因此存在非零多项式 $F \in k[X_1, \ldots, X_{n + 1}]$，使得
$F(x_1, \ldots, x_{n+1}) = 0$。取总次数最小的 $F$。
则 $F$ 是不可约的，因为至少有一个不可约因子也必须具有相同的性质。

\medskip\noindent
我们断言，对某个 $i$，$X_i$ 在 $F$ 中出现的幂不全是 $p$ 的倍数。
反设 $F$ 中出现的所有指数都是 $p$ 的倍数，则集合
$$
\{x_1^{\alpha_1} \ldots x^{\alpha_{n+1}}_{n+1} \mid \lambda_\alpha \neq 0\}
\subset
K
$$
是 $k$-线性相关的，其中 $\lambda_\alpha$ 是 $F$ 的系数。
由假设 (2)，我们得到集合
$$
\{x_1^{\alpha_1 / p} \ldots x^{\alpha_{n+1} / p}_{n+1}
\mid \lambda_\alpha \neq 0 \}
$$
也是 $k$-线性相关的，这与 $\deg(F)$ 的最小性矛盾。

\medskip\noindent
取某个 $i$，使 $p$ 次幂以外的 $X_i$ 幂出现在 $F$ 中。于是可见 $x_i$ 在下列域上代数：
$$
L = k(x_1, \ldots, x_{i - 1}, x_{i + 1}, \ldots, x_{n+1})
$$
由《域》，引理 \ref{fields-lemma-transcendence-degree}，可知下列元素：
$$
x_1, \ldots, x_{i - 1}, x_{i + 1}, \ldots, x_{n+1}
$$
是 $K/k$ 的超越基。因此 $L$ 是由下列元素：
$$
x_1, \ldots, x_{i - 1}, x_{i + 1}, \ldots, x_{n + 1}
$$
生成的、在 $k$ 上的多项式环的分式域。
由 Gauss 引理，我们得到
$$
P(T) =
F(x_1, \ldots, x_{i - 1}, T, x_{i + 1}, \ldots, x_{n + 1}) \in L[T]
$$
是不可约的。按构造，$P(T)$ 不包含于 $L[T^p]$ 中。
因此 $K/L$ 是可分的，正如所要求的。
\end{proof}

\noindent
令 $p$ 为素数，令 $k$ 为特征为 $p$ 的域。
在此情形下，我们用 $k^{1/p}$ 表示通过向 $k$ 添入
$p$ 次根（对应于 $k$ 中的所有元素）并添入 $k$ 所得的扩张。
（换言之，它是代数闭包中由 $k$ 生成的、由 $p$ 次根（取自 $k$ 的元素）生成的子域。）

\begin{lemma}
\label{algebra-lemma-characterize-separable-field-extensions}
令 $k$ 为特征为 $p > 0$ 的域。
令 $K/k$ 为域扩张。
以下命题等价：
\begin{enumerate}
\item $K$ 在 $k$ 上可分，
\item 对于每个 $k$-线性无关子集 $\{a_1, \ldots, a_m\}$，
在 $K$ 中，其对应的集合 $\{a^p_1, \ldots, a_m^p\}$ 是 $k$-线性无关的，
\item 环 $K \otimes_k k^{1/p}$ 是约化的，并且
\item $K$ 在 $k$ 上几何约化。
\end{enumerate}
\end{lemma}
\begin{proof}
蕴含关系 (1) $\Rightarrow$ (4) 由引理
\ref{algebra-lemma-separable-extension-preserves-reducedness} 得出。
蕴含关系 (4) $\Rightarrow$ (3) 是直接的。

\medskip\noindent
假设 (3)。考虑由
$m : K \otimes_k k^{1/p} \rightarrow K$ 给出的环同态
$$
\lambda \otimes \mu \rightarrow \lambda^p \mu^p
$$
注意，对于所有 $x \in K \otimes_k k^{1/p}$，都有
$x^p = m(x) \otimes 1$。
由于 $K \otimes_k k^{1/p}$ 是约化的，
可见 $m$ 是单射。若 $\{a_1, \ldots, a_m\} \subset K$
是 $k$-线性无关的，则 $\{a_1 \otimes 1, \ldots, a_m \otimes 1\}$
是 $k^{1/p}$-线性无关的。由 $m$ 的单射性可知，
$a_1^p, \ldots, a_m^p$ 的任何非平凡 $k$-线性组合都不为零。
因此 (3) 推出 (2)。

\medskip\noindent
假设 (2)。为证明 (1)，可假设 $K$ 在 $k$ 上有限生成，
并且需要证明 $K$ 在 $k$ 上可分生成。
令 $\{x_1, \ldots, x_d\}$ 为 $K/k$ 的超越基。由
《域》，引理 \ref{fields-lemma-algebraic-finitely-generated}，
有 $[K : K'] < \infty$，其中 $K' = k(x_1, \ldots, x_d)$。
选取超越基，使不可分次数 $[K : K']_i$ 最小。
若 $K / K'$ 可分，则结论成立。
反设并导出矛盾。则存在 $x_{d + 1} \in K$，它在 $K'$ 上不可分，
特别地，$[K'(x_{d+1}) : K']_i > 1$。于是由
引理 \ref{algebra-lemma-mini-separability}，存在 $1 \leq j \leq n + 1$，
使得 $K'' = k(x_1, \ldots, \widehat{x}_j, \ldots, x_{d+1})$
满足 $[K'(x_{d+1}) : K'']_i = 1$。由次数的乘法性，
$[K : K'']_i < [K : K']_i$，矛盾得证。
\end{proof}

\begin{lemma}
\label{algebra-lemma-separably-generated-separable}
可分生成的域扩张是可分的。
\end{lemma}

\begin{proof}
结合引理 \ref{algebra-lemma-separable-extension-preserves-reducedness} 与
引理 \ref{algebra-lemma-characterize-separable-field-extensions}。
\end{proof}

\noindent
下面的引理将使用完美闭包的概念，
其定义见定义 \ref{algebra-definition-perfection}。

\begin{lemma}
\label{algebra-lemma-geometrically-reduced-finite-purely-inseparable-extension}
令 $k$ 为域。令 $S$ 为 $k$-代数。
以下命题等价：
\begin{enumerate}
\item 对每个有限的纯不可分扩张，$k' \otimes_k S$ 是约化的，其中 $k'$ 是 $k$ 的扩张，
\item $k^{1/p} \otimes_k S$ 是约化的，
\item $k^{perf} \otimes_k S$ 是约化的，其中 $k^{perf}$ 是
$k$ 的完美闭包，
\item $\overline{k} \otimes_k S$ 是约化的，其中 $\overline{k}$ 是
$k$ 的代数闭包，并且
\item $S$ 在 $k$ 上几何约化。
\end{enumerate}
\end{lemma}

\begin{proof}
注意，任何有限纯不可分扩张 $k'/k$ 都嵌入完美闭包中。
此外，在 $k^{perf}$ 中，$k^{1/p}$ 嵌入 $k^{perf}$，而后者嵌入
$\overline{k}$。因此显然有
(5) $\Rightarrow$ (4) $\Rightarrow$ (3) $\Rightarrow$ (2)，
并且 (3) $\Rightarrow$ (1)。

\medskip\noindent
我们证明 (1) $\Rightarrow$ (5)。
假设 $k' \otimes_k S$ 对每个有限
纯不可分扩张 $k'$（在 $k$ 上）都是约化的。令 $K/k$ 为域扩张。我们需要证明 $K \otimes_k S$
是约化的。由引理 \ref{algebra-lemma-limit-argument}，可归约到
$K/k$ 是有限生成域扩张的情形。取如下图表：
$$
\xymatrix{
K \ar[r] & K' \\
k \ar[u] \ar[r] & k' \ar[u]
}
$$
如引理 \ref{algebra-lemma-make-separably-generated} 中所述。
由假设，$k' \otimes_k S$ 是约化的。
由引理 \ref{algebra-lemma-separable-extension-preserves-reducedness}，
可知 $K' \otimes_k S$ 是约化的。
因此 $K \otimes_k S$ 是约化的，正如所要求的。

\medskip\noindent
最后证明 (2) $\Rightarrow$ (5)。
假设 $k^{1/p} \otimes_k S$ 是约化的。则 $S$ 是约化的。
此外，对极小素理想处的每个局部化 $S_{\mathfrak p}$，
素理想 $\mathfrak p$，环 $k^{1/p}\otimes_k S_{\mathfrak p}$
是 $k^{1/p} \otimes_k S$ 的局部化，因而是约化的。
但是由引理 \ref{algebra-lemma-minimal-prime-reduced-ring}，$S_{\mathfrak p}$ 是域，
所以由引理 \ref{algebra-lemma-characterize-separable-field-extensions}，
$S_{\mathfrak p}$ 是几何约化的。
由引理 \ref{algebra-lemma-generic-points-geometrically-reduced}，
可知 $S$ 是几何约化的。
\end{proof}

\section{完美域}
\label{algebra-section-perfect-fields}

\noindent
下面给出定义。

\begin{definition}
\label{algebra-definition-perfect}
令 $k$ 为域。若 $k$ 的每个域扩张都在 $k$ 上可分，
则称 $k$ 为{\it 完美的}。
\end{definition}

\begin{lemma}
\label{algebra-lemma-perfect}
域 $k$ 完美，当且仅当它是特征为 $0$ 的域，或是特征为 $p > 0$ 的域，
且每个元素都有一个 $p$ 次根。
\end{lemma}

\begin{proof}
特征为零的情形显然。
假设 $k$ 的特征为 $p > 0$。
若 $k$ 完美，则所有向 $k$ 添入某个元素的
$p$ 次根所得到的域扩张都必须是平凡的，因此 $k$ 的每个元素都有 $p$ 次根。
反之，若每个元素都有 $p$ 次根，则 $k = k^{1/p}$，且 $k$ 的每个域扩张都
由引理 \ref{algebra-lemma-characterize-separable-field-extensions} 可知是可分的。
\end{proof}

\begin{lemma}
\label{algebra-lemma-make-separable}
令 $K/k$ 为有限生成的域扩张。
存在图表
$$
\xymatrix{
K \ar[r] & K' \\
k \ar[u] \ar[r] & k' \ar[u]
}
$$
其中 $k'/k$、$K'/K$ 是有限纯不可分域扩张，且 $K'/k'$ 是可分域扩张。
在此情形下，可以假设 $K' = k'K$ 是复合域，
并且还可以假设 $K' = (k' \otimes_k K)_{red}$。
\end{lemma}

\begin{proof}
由引理 \ref{algebra-lemma-make-separably-generated}，可找到一个图表，使
$K'/k'$ 可分生成。
由引理 \ref{algebra-lemma-separably-generated-separable}，这推出 $K'$ 在 $k'$ 上可分。
复合域 $k'K$ 是 $K'/k'$ 的子扩张，因此
$k' \subset k'K$ 由引理 \ref{algebra-lemma-subextensions-are-separable} 可知是可分的。
环 $(k' \otimes_k K)_{red}$ 是整环，因为对某个
$n \gg 0$，映射 $x \mapsto x^{p^n}$ 将其映入 $K$。
因此由引理 \ref{algebra-lemma-integral-over-field}，它是域。
于是 $(k' \otimes_k K)_{red} \to K'$ 将其同构地映到 $k'K$ 上。
\end{proof}

\begin{lemma}
\label{algebra-lemma-perfection}
\begin{slogan}
每个域都有唯一的完美闭包。
\end{slogan}
对每个域 $k$，存在纯不可分扩张 $k'/k$，使得 $k'$ 是完美的。
域扩张 $k'/k$ 在唯一同构意义下是唯一的。
\end{lemma}

\begin{proof}
若 $k$ 的特征为零，则 $k' = k$ 是唯一选择。假设 $k$ 的特征为 $p > 0$。
对每个 $n > 0$，存在唯一的代数扩张
$k \subset k^{1/p^n}$，使得 (a) 每个元素 $\lambda \in k$
在 $k^{1/p^n}$ 中都有 $p^n$ 次根，并且 (b) 对每个元素
$\mu \in k^{1/p^n}$，都有 $\mu^{p^n} \in k$。
具体地，考虑环映射 $k \to k^{1/p^n} = k$，$x \mapsto x^{p^n}$。
该映射是单射，并满足 (a) 和 (b)。显然，作为经由映射
$y \mapsto y^p$ 的 $k$-扩张，有
$k^{1/p^n} \subset k^{1/p^{n + 1}}$。于是可取
$k' = \bigcup k^{1/p^n}$。
省略一些细节。
\end{proof}

\begin{definition}
\label{algebra-definition-perfection}
令 $k$ 为域。引理 \ref{algebra-lemma-perfection} 中的域扩张 $k'/k$
称为 $k$ 的{\it 完美闭包}。记作 $k^{perf}/k$。
\end{definition}

\noindent
注意，若 $k'/k$ 是任意代数纯不可分扩张，则
$k'$ 是 $k^{perf}$ 的子扩张，即 $k^{perf}/k'/k$。事实上，
由引理 \ref{algebra-lemma-perfection} 的唯一性，$(k')^{perf}$ 与 $k^{perf}$ 同构。

\begin{lemma}
\label{algebra-lemma-perfect-reduced}
令 $k$ 为完美域。
任意约化 $k$-代数在 $k$ 上都是几何约化的。
令 $R$、$S$ 为 $k$-代数。
假设 $R$ 和 $S$ 都是约化的。
则 $k$-代数 $R \otimes_k S$ 是约化的。
\end{lemma}

\begin{proof}
第一个陈述由引理
\ref{algebra-lemma-geometrically-reduced-finite-purely-inseparable-extension} 得出。
第二个陈述应用第一个陈述以及引理
\ref{algebra-lemma-geometrically-reduced-any-reduced-base-change} 即得。
\end{proof}









\section{万有同胚}
\label{algebra-section-universal-homeomorphism}

\noindent
设 $k'/k$ 为代数纯不可分域扩张。那么，对任意 $k$-代数 $R$，
环映射 $R \to k' \otimes_k R$ 都诱导谱之间的同胚。
其原因是下面稍为一般的引理 \ref{algebra-lemma-p-ring-map}。

\begin{lemma}
\label{algebra-lemma-surjective-locally-nilpotent-kernel}
设 $\varphi : R \to S$ 为满映射，且其核局部幂零。
则 $\varphi$ 诱导谱之间的同胚以及剩余域之间的同构。
对任意环映射 $R \to R'$，环映射
$R' \to R' \otimes_R S$ 也是核局部幂零的满映射。
\end{lemma}

\begin{proof}
由引理 \ref{algebra-lemma-spec-closed}，映射 $\Spec(S) \to \Spec(R)$
是到闭子集 $V(\Ker(\varphi))$ 上的同胚。当然，
$V(\Ker(\varphi)) = \Spec(R)$，因为 $R$ 的每个素理想都包含
$R$ 的每个幂零元。这也推出关于剩余域的陈述。
由张量积的右正合性可知，$\Ker(\varphi)R'$ 是满映射
$R' \to R' \otimes_R S$ 的核。最后的陈述于是由
引理 \ref{algebra-lemma-locally-nilpotent} 得出。
\end{proof}

\begin{lemma}
\label{algebra-lemma-powers-field}
\begin{reference}
\cite[Lemma 3.1.6]{Alper-adequate}
\end{reference}
设 $k'/k$ 为域扩张。下列条件等价：
\begin{enumerate}
\item 对每个 $x \in k'$，存在 $n > 0$，使得 $x^n \in k$；
\item $k' = k$，或者 $k$ 和 $k'$ 的特征为 $p > 0$，并且
$k'/k$ 是纯不可分扩张，或 $k$ 和 $k'$ 都是 $\mathbf{F}_p$ 的代数扩张。
\end{enumerate}
\end{lemma}

\begin{proof}
注意，(2) 中列出的每一种情形都满足 (1)。
因此，假设 $k'/k$ 满足 (1)，并证明我们处于 (2) 的某种情形。
排除 $k = k'$ 的情形后，可以假设 $k' \not = k$。
显然，$k'/k$ 是代数扩张。因此可以假设 $k'/k$ 是非平凡有限扩张。
令 $k'/k'_{sep}/k$ 为域，引理 \ref{fields-lemma-separable-first}
中得到的可分子扩张。须证明 $k = k'_{sep}$，或 $k$ 是
$\mathbf{F}_p$ 的代数扩张。因此可以假设 $k'/k$ 是非平凡有限可分扩张，
而我们须证明 $k$ 是 $\mathbf{F}_p$ 的代数扩张。

\medskip\noindent
取 $x \in k'$，$x \not \in k$。取 $n, m > 0$，使得
$x^n \in k$ 且 $(x + 1)^m \in k$。令 $\overline{k}$ 为 $k$ 的代数闭包。
可以选取嵌入 $\sigma, \tau : k' \to \overline{k}$，使得
$\sigma(x) \not = \tau(x)$。这由域，第
\ref{fields-section-separable-extensions} 节中的讨论得出
（更确切地说，将 $k'$ 换成由 $x$ 生成的 $k$-扩张后，
它由域，引理 \ref{fields-lemma-count-embeddings} 得出）。
于是，对 $\overline{k}$ 中某个 $n$ 次单位根 $\zeta$，有
$\sigma(x) = \zeta \tau(x)$。类似地，对某个 $m$ 次单位根
$\zeta' \in \overline{k}$，有
$\sigma(x + 1) = \zeta' \tau(x + 1)$。
由于 $\sigma(x + 1) \not = \tau(x + 1)$，可知 $\zeta' \not = 1$。于是
$$
\zeta' (\tau(x) + 1) =
\zeta' \tau(x + 1) =
\sigma(x + 1) =
\sigma(x) + 1 =
\zeta \tau(x) + 1
$$
推出
$$
\tau(x) (\zeta' - \zeta) = 1 - \zeta'
$$
因此 $\zeta' \not = \zeta$，且
$$
\tau(x) = (1 - \zeta')/(\zeta' - \zeta)
$$
故 $k'$ 中不属于 $k$ 的每个元素都在素子域上代数。
由于 $k'$ 由 $k'$ 中不属于 $k$ 的元素在素子域上生成，
可知 $k'$（从而 $k$）在素子域上代数。

\medskip\noindent
最后，若 $k$ 的特征为 $0$，则上述结论如下导致矛盾
（鼓励读者自行寻找证明）。对每个有理数 $y$，类似地得到单位根
$\zeta_y$，使得 $\sigma(x + y) = \zeta_y\tau(x + y)$。于是有
$$
\zeta \tau(x) + y = \zeta_y(\tau(x) + y)
$$
并且由上面关于 $\tau(x)$ 的公式可知
$\zeta_y \in \mathbf{Q}(\zeta, \zeta')$。由于数域
$\mathbf{Q}(\zeta, \zeta')$ 仅含有限多个单位根，
可找到两个不同的有理数 $y, y'$，使得
$\zeta_y = \zeta_{y'}$。于是
$$
y - y' =
\sigma(x + y) - \sigma(x + y') =
\zeta_y(\tau(x + y)) - \zeta_{y'}\tau(x + y') = \zeta_y(y - y')
$$
这推出 $\zeta_y = 1$，矛盾。
\end{proof}

\begin{lemma}
\label{algebra-lemma-powers}
设 $\varphi : R \to S$ 为环映射。若
\begin{enumerate}
\item 对任意 $x \in S$，存在 $n > 0$，使得
$x^n$ 属于 $\varphi$ 的像；
\item $\Ker(\varphi)$ 局部幂零，
\end{enumerate}
则 $\varphi$ 诱导谱之间的同胚，并诱导满足
引理 \ref{algebra-lemma-powers-field} 中等价条件的剩余域扩张。
\end{lemma}

\begin{proof}
假设 (1) 和 (2)。设 $\mathfrak q, \mathfrak q'$ 为 $S$ 的素理想，
它们位于 $R$ 的同一素理想 $\mathfrak p$ 上。假设 $x \in S$，且
$x \in \mathfrak q$、$x \not \in \mathfrak q'$。则对所有 $n > 0$，
$x^n \in \mathfrak q$ 且 $x^n \not \in \mathfrak q'$。
若对某个 $n > 0$，有 $x^n = \varphi(y)$，其中 $y \in R$，则
$$
x^n \in \mathfrak q \Rightarrow y \in \mathfrak p \Rightarrow
x^n \in \mathfrak q'
$$
这导致矛盾。因此不存在上述 $x$，从而
$\mathfrak q = \mathfrak q'$；也就是说，谱上的映射是单射。
由假设 (2)，核 $I = \Ker(\varphi)$ 包含于每个素理想中，故作为拓扑空间，
$\Spec(R) = \Spec(R/I)$。由假设 (1)，诱导映射 $R/I \to S$ 是整的，
所以引理 \ref{algebra-lemma-integral-overring-surjective} 表明
$\Spec(S) \to \Spec(R/I)$ 是满射。结合上述结论，
$\Spec(S) \to \Spec(R)$ 是双射。
若任取 $x \in S$，并选取 $y \in R$，使得对某个 $n > 0$，
$\varphi(y) = x^n$，则在上述双射下，开集
$D(x) \subset \Spec(S)$ 对应于开集 $D(y) \subset \Spec(R)$。
因此映射 $\Spec(S) \to \Spec(R)$ 是同胚。

\medskip\noindent
为证明关于剩余域的陈述，设 $\mathfrak q \subset S$ 为位于素理想
$\mathfrak p \subset R$ 上的素理想。令 $x \in \kappa(\mathfrak q)$。
将 $\kappa(\mathfrak q)$ 视为局部环 $S_\mathfrak q$ 的剩余域，
则 $x$ 是某个 $y/z \in S_\mathfrak q$ 的像，其中
$y \in S$、$z \in S$、$z \not \in \mathfrak q$。
选取 $n, m > 0$，使得 $y^n, z^m$ 属于 $\varphi$ 的像。
那么 $x^{nm}$ 是 $(y/z)^{nm} = (y^n)^m/(z^m)^n$ 的剩余类，
而后者属于 $R_\mathfrak p \to S_\mathfrak q$ 的像。
因此 $x^{nm}$ 属于 $\kappa(\mathfrak p) \to \kappa(\mathfrak q)$ 的像。
\end{proof}

\begin{lemma}
\label{algebra-lemma-2-3-ring-map}
设 $\varphi : R \to S$ 为环映射。假设
\begin{enumerate}
\item[(a)] 作为 $R$-代数，$S$ 由满足
$x^2, x^3 \in \varphi(R)$ 的元素 $x$ 生成；
\item[(b)] $\Ker(\varphi)$ 局部幂零。
\end{enumerate}
则 $\varphi$ 诱导剩余域之间的同构以及谱之间的同胚。
对任意环映射 $R \to R'$，环映射 $R' \to R' \otimes_R S$
也满足 (a) 和 (b)。
\end{lemma}

\begin{proof}
假设 (a) 和 (b)。由于 $S$ 在 $R$ 上整，谱上的映射是闭映射；
参见引理 \ref{algebra-lemma-going-up-closed} 和
\ref{algebra-lemma-integral-going-up}。由引理
\ref{algebra-lemma-image-dense-generic-points}，其像稠密。
因此 $\Spec(S) \to \Spec(R)$ 是满射。
若 $\mathfrak q \subset S$ 是位于 $\mathfrak p \subset R$ 上的素理想，
则域扩张 $\kappa(\mathfrak q)/\kappa(\mathfrak p)$ 由元素
$\alpha \in \kappa(\mathfrak q)$ 生成，这些元素的平方和立方
均在 $\kappa(\mathfrak p)$ 中。因此显然
$\alpha \in \kappa(\mathfrak p)$，从而
$\kappa(\mathfrak q) = \kappa(\mathfrak p)$。
若 $\mathfrak q, \mathfrak q'$ 是位于 $\mathfrak p$ 上的两个不同素理想，
则 (a) 中作为 $S$ 的生成元的某个 $x$，在
$\kappa(\mathfrak q) = \kappa(\mathfrak p)$ 和
$\kappa(\mathfrak q') = \kappa(\mathfrak p)$ 中的像必然不同。
这与 $x^2$ 和 $x^3$ 的像都相同这一事实矛盾。
这就证明 $\Spec(S) \to \Spec(R)$ 是单射，因而是同胚
（由已经证明的结论）。

\medskip\noindent
由于 $\varphi$ 诱导谱之间的同胚，它在谱上尤其是满射；
此性质在任意基变换下保持，参见引理
\ref{algebra-lemma-surjective-spec-radical-ideal}。
因此，对任意 $R \to R'$，环映射 $R' \to R' \otimes_R S$ 的核
由幂零元组成；参见引理 \ref{algebra-lemma-image-dense-generic-points}。
换言之，$R' \to R' \otimes_R S$ 满足 (b)。
显然，(a) 在基变换下保持。
\end{proof}

\begin{lemma}
\label{algebra-lemma-help-with-powers}
设 $p$ 为素数，$n, m > 0$ 为两个整数。存在整数 $a$，使得
$(x + y)^{p^a}, p^a(x + y) \in \mathbf{Z}[x^{p^n}, p^nx, y^{p^m}, p^my]$。
\end{lemma}

\begin{proof}
只要 $a \geq n, m$，关于 $p^a(x + y)$ 的结论就是显然的。
事实上，取 $a \gg n, m$。写成
$$
(x + y)^{p^a}  = \sum\nolimits_{i, j \geq 0, i + j = p^a}
{p^a \choose i, j} x^iy^j
$$
对每个满足 $i + j = p^a$ 的 $i, j \geq 0$，写
$i = q p^n + r$，其中 $r \in \{0, \ldots, p^n - 1\}$；并写
$j = q' p^m + r'$，其中 $r' \in \{0, \ldots, p^m - 1\}$。
若
$$
p^{nr + mr'} \text{ divides } {p^a \choose i, j}
$$
成立，则条件
$(x + y)^{p^a} \in \mathbf{Z}[x^{p^n}, p^nx, y^{p^m}, p^my]$
成立。若 $r = r' = 0$，则此整除性成立。若 $r \not = 0$，则写成
$$
{p^a \choose i, j} = \frac{p^a}{i} {p^a - 1 \choose i - 1, j}
$$
由于 $r \not = 0$，有理数 $p^a/i$ 的 $p$-进赋值至少为
$a - (n - 1)$（因为 $i$ 不被 $p^n$ 整除）。
所以在此情形下，${p^a \choose i, j}$ 被 $p^{a - n + 1}$ 整除。
类似地，若 $r' \not = 0$，则 ${p^a \choose i, j}$ 被
$p^{a - m + 1}$ 整除。取 $a = np^n + mp^m + n + m$ 即可。
\end{proof}

\begin{lemma}
\label{algebra-lemma-p-ring-map-field}
设 $k'/k$ 为域扩张，$p$ 为素数。下列条件等价：
\begin{enumerate}
\item 作为 $k$ 的域扩张，$k'$ 由满足下述条件的元素 $x$ 生成：
存在 $n > 0$，使得 $x^{p^n} \in k$ 且 $p^nx \in k$；
\item $k = k'$，或者 $k$ 和 $k'$ 的特征均为 $p$，且
$k'/k$ 是纯不可分扩张。
\end{enumerate}
\end{lemma}

\begin{proof}
令 $x \in k'$。若存在 $n > 0$，使得 $x^{p^n} \in k$ 且
$p^nx \in k$，并且特征不是 $p$，则 $x \in k$。
若特征为 $p$，则 $x^{p^n} \in k$，所以 $x$ 在 $k$ 上纯不可分。
\end{proof}

\begin{lemma}
\label{algebra-lemma-p-ring-map}
设 $\varphi : R \to S$ 为环映射，$p$ 为素数。假设
\begin{enumerate}
\item[(a)] 作为 $R$-代数，$S$ 由满足下述条件的元素 $x$ 生成：
存在 $n > 0$，使得 $x^{p^n} \in \varphi(R)$ 且
$p^nx \in \varphi(R)$；
\item[(b)] $\Ker(\varphi)$ 局部幂零。
\end{enumerate}
则 $\varphi$ 诱导谱之间的同胚，并诱导满足引理
\ref{algebra-lemma-p-ring-map-field} 中等价条件的剩余域扩张。
对任意环映射 $R \to R'$，环映射
$R' \to R' \otimes_R S$ 也满足 (a) 和 (b)。
\end{lemma}

\begin{proof}
假设 (a) 和 (b)。注意，(b) 等价于引理 \ref{algebra-lemma-powers} 的条件 (2)。
令 $T \subset S$ 为由下述元素组成的集合：$x \in S$，且存在整数
$n > 0$，使得 $x^{p^n} , p^n x \in \varphi(R)$。
断言 $T = S$。这将证明引理 \ref{algebra-lemma-powers} 的条件 (1) 成立，
从而 $\varphi$ 诱导谱之间的同胚。
由假设 (a)，只须证明 $T \subset S$ 是 $R$-子代数。
若 $x \in T$ 且 $y \in R$，则显然 $yx \in T$。
假设 $x, y \in T$，并且 $n, m > 0$ 满足
$x^{p^n}, y^{p^m}, p^n x, p^m y \in \varphi(R)$。
则 $(xy)^{p^{n + m}}, p^{n + m}xy \in \varphi(R)$，
因而 $xy \in T$。由引理 \ref{algebra-lemma-help-with-powers}，有 $x + y \in T$，
断言得证。

\medskip\noindent
由于 $\varphi$ 诱导谱之间的同胚，它在谱上尤其是满射；
此性质在任意基变换下保持，参见引理
\ref{algebra-lemma-surjective-spec-radical-ideal}。
因此，对任意 $R \to R'$，环映射 $R' \to R' \otimes_R S$ 的核
由幂零元组成；参见引理 \ref{algebra-lemma-image-dense-generic-points}。
换言之，$R' \to R' \otimes_R S$ 满足 (b)。
显然，(a) 在基变换下保持。
最后，关于剩余域的条件由 (a) 得出，因为 $S$ 作为 $R$-代数的生成元
映到各剩余域扩张的生成元。
\end{proof}

\begin{lemma}
\label{algebra-lemma-radicial}
设 $\varphi : R \to S$ 为环映射。假设
\begin{enumerate}
\item $\varphi$ 诱导谱上的单射；
\item $\varphi$ 诱导纯不可分剩余域扩张。
\end{enumerate}
则对任意环映射 $R \to R'$，映射 $R' \to R' \otimes_R S$
也满足性质 (1) 和 (2)。
\end{lemma}

\begin{proof}
置 $S' = R' \otimes_R S$，于是有环谱的连续映射的交换图
$$
\xymatrix{
\Spec(S') \ar[r] \ar[d] & \Spec(S) \ar[d] \\
\Spec(R') \ar[r] & \Spec(R)
}
$$
设 $\mathfrak p' \subset R'$ 为位于 $\mathfrak p \subset R$ 上的素理想。
若 $S$ 中没有位于 $\mathfrak p$ 上的素理想，则 $S'$ 中也没有位于
$\mathfrak p'$ 上的素理想。否则，由注记 \ref{algebra-remark-fundamental-diagram}，
$F = S \otimes_R \kappa(\mathfrak p)$ 有唯一的素理想 $\mathfrak r$，
且其剩余域在 $\kappa(\mathfrak p)$ 上纯不可分。考虑环映射
$$
\kappa(\mathfrak p) \to F \to \kappa(\mathfrak r)
$$
由引理 \ref{algebra-lemma-minimal-prime-reduced-ring}，理想
$\mathfrak r \subset F$ 局部幂零，所以可将引理
\ref{algebra-lemma-surjective-locally-nilpotent-kernel} 应用于环映射
$F \to \kappa(\mathfrak r)$。还可将引理 \ref{algebra-lemma-p-ring-map}
应用于环映射 $\kappa(\mathfrak p) \to \kappa(\mathfrak r)$。
因此，在映射
$$
\kappa(\mathfrak p') \to
\kappa(\mathfrak p') \otimes_{\kappa(\mathfrak p)} F \to
\kappa(\mathfrak p') \otimes_{\kappa(\mathfrak p)} \kappa(\mathfrak r)
$$
中，复合映射和第二个箭头均诱导谱上的双射和纯不可分剩余域扩张。
这推出第一个映射也具有同样的性质。由于
$$
\kappa(\mathfrak p') \otimes_{\kappa(\mathfrak p)} F =
\kappa(\mathfrak p') \otimes_{\kappa(\mathfrak p)}
\kappa(\mathfrak p) \otimes_R S =
\kappa(\mathfrak p') \otimes_R S =
\kappa(\mathfrak p') \otimes_{R'} R' \otimes_R S
$$
由注记 \ref{algebra-remark-fundamental-diagram} 中的讨论可得结论。
\end{proof}

\begin{lemma}
\label{algebra-lemma-radicial-integral}
设 $\varphi : R \to S$ 为环映射。假设
\begin{enumerate}
\item $\varphi$ 是整的；
\item $\varphi$ 诱导谱上的单射；
\item $\varphi$ 诱导纯不可分剩余域扩张。
\end{enumerate}
则 $\varphi$ 诱导从 $\Spec(S)$ 到 $\Spec(R)$ 的某个闭子集上的同胚，
并且对任意环映射 $R \to R'$，映射 $R' \to R' \otimes_R S$
满足性质 (1)、(2)、(3)。
\end{lemma}

\begin{proof}
由引理 \ref{algebra-lemma-going-up-closed} 和
\ref{algebra-lemma-integral-going-up}，谱上的映射是闭映射。
由引理 \ref{algebra-lemma-radicial} 和 \ref{algebra-lemma-base-change-integral}，
这些性质在基变换下保持。
\end{proof}

\begin{lemma}
\label{algebra-lemma-radicial-integral-bijective}
设 $\varphi : R \to S$ 为环映射。假设
\begin{enumerate}
\item $\varphi$ 是整的；
\item $\varphi$ 诱导谱上的双射；
\item $\varphi$ 诱导纯不可分剩余域扩张。
\end{enumerate}
则 $\varphi$ 诱导谱之间的同胚，并且对任意环映射
$R \to R'$，映射 $R' \to R' \otimes_R S$ 满足性质 (1)、(2)、(3)。
\end{lemma}

\begin{proof}
由引理 \ref{algebra-lemma-radicial-integral} 和
\ref{algebra-lemma-surjective-spec-radical-ideal} 得出。
\end{proof}

\begin{lemma}
\label{algebra-lemma-universally-bijective}
设 $\varphi : R \to S$ 为满足下列条件的环映射：
\begin{enumerate}
\item $\varphi$ 的核局部幂零；
\item 作为 $R$-代数，$S$ 由满足下述条件的元素 $x$ 生成：
存在 $n > 0$ 及多项式 $P(T) \in R[T]$，使其在 $S[T]$ 中的像为
$(T - x)^n$。
\end{enumerate}
则 $\Spec(S) \to \Spec(R)$ 是同胚，并且 $R \to S$
诱导纯不可分剩余域扩张。此外，条件 (1) 和 (2) 在任意基变换下保持。
\end{lemma}

\begin{proof}
可以将 $R$ 替换为 $R/\Ker(\varphi)$；参见引理
\ref{algebra-lemma-surjective-locally-nilpotent-kernel}。
假设 (2) 表明，$S$ 由在 $R$ 上整的元素作为 $R$-代数生成。
因此 $R \subset S$ 是整的
（引理 \ref{algebra-lemma-integral-closure-is-ring}）。特别地，
$\Spec(S) \to \Spec(R)$ 是满的闭映射
（引理 \ref{algebra-lemma-integral-overring-surjective}、
\ref{algebra-lemma-going-up-closed} 和 \ref{algebra-lemma-integral-going-up}）。

\medskip\noindent
设 $x \in S$ 为 (2) 中的某个生成元；也就是说，存在
$n > 0$，使得 $(T - x)^n \in R[T]$。设 $\mathfrak p \subset R$ 为素理想。
由上述结论及引理 \ref{algebra-lemma-in-image}，环
$\kappa(\mathfrak p) \otimes_R S$ 非零。
若 $\kappa(\mathfrak p)$ 的特征为零，则 $nx \in R$ 表明
$1 \otimes x$ 属于 $\kappa(\mathfrak p) \to
\kappa(\mathfrak p) \otimes_R S$ 的像。因此
$\kappa(\mathfrak p) \to \kappa(\mathfrak p) \otimes_R S$ 是同构。
若 $\kappa(\mathfrak p)$ 的特征为 $p > 0$，则写
$n = p^k m$，其中 $m$ 与 $p$ 互素。在
$\kappa(\mathfrak p) \otimes_R S[T]$ 中有
$$
(T - 1 \otimes x)^n = ((T - 1 \otimes x)^{p^k})^m =
(T^{p^k} - 1 \otimes x^{p^k})^m
$$
因而可知 $mx^{p^k} \in R$。这表明 $1 \otimes x^{p^k}$ 属于
$\kappa(\mathfrak p) \to \kappa(\mathfrak p) \otimes_R S$ 的像。
所以可将引理 \ref{algebra-lemma-p-ring-map} 应用于
$\kappa(\mathfrak p) \to \kappa(\mathfrak p) \otimes_R S$。
在两种情形下，都可断定 $\kappa(\mathfrak p) \otimes_R S$
有唯一的素理想，且其剩余域在 $\kappa(\mathfrak p)$ 上纯不可分。
由注记 \ref{algebra-remark-fundamental-diagram} 可知 $\varphi$ 在谱上为双射。

\medskip\noindent
关于基变换的陈述是立即的。
\end{proof}











\section{几何不可约代数}
\label{algebra-section-algebras-over-fields}

\noindent
若对每个域扩张 $k'/k$，代数 $S \otimes_k k'$ 都有唯一的极小素理想，
则称域 $k$ 上的代数 $S$ 几何不可约。本节将发展与此概念有关的一些理论。

\begin{lemma}
\label{algebra-lemma-flat-fibres-irreducible}
设 $R \to S$ 为环映射。假设
\begin{enumerate}
\item[(a)] $\Spec(R)$ 不可约；
\item[(b)] $R \to S$ 平坦；
\item[(c)] $R \to S$ 是有限表示的；
\item[(d)] 对 $R$ 中稠密的一族素理想 $\mathfrak p$，纤维环
$S \otimes_R \kappa(\mathfrak p)$ 的谱不可约。
\end{enumerate}
则 $\Spec(S)$ 不可约。更一般地，将 (b) $+$ (c) 替换为
“映射 $\Spec(S) \to \Spec(R)$ 是开映射”，结论仍成立。
\end{lemma}

\begin{proof}
假设 (b) 和 (c) 表明谱上的映射是开映射；参见命题
\ref{algebra-proposition-fppf-open}。因此，本引理由拓扑，引理
\ref{topology-lemma-irreducible-on-top} 得出。
\end{proof}

\begin{lemma}
\label{algebra-lemma-separably-closed-irreducible}
设 $k$ 为可分闭域，$R$、$S$ 为 $k$-代数。
若 $R$、$S$ 各有唯一的极小素理想，则 $R \otimes_k S$ 也如此。
\end{lemma}

\begin{proof}
设 $k \subset \overline{k}$ 为完美闭包；参见定义
\ref{algebra-definition-perfection}。由假设，$\overline{k}$ 是代数闭域。
环映射 $R \to R \otimes_k \overline{k}$、
$S \to S \otimes_k \overline{k}$ 以及
$R \otimes_k S \to (R \otimes_k S) \otimes_k \overline{k}
= (R \otimes_k \overline{k}) \otimes_{\overline{k}} (S \otimes_k \overline{k})$
满足引理 \ref{algebra-lemma-p-ring-map} 的假设。因此可以假设 $k$ 是代数闭域。

\medskip\noindent
可以把 $R$ 和 $S$ 替换为它们的约化。因此可以假设 $R$ 和 $S$ 是整环。
由引理 \ref{algebra-lemma-perfect-reduced}，$R \otimes_k S$ 是约化的。
所以其谱可约，当且仅当它含有非零零因子。
由引理 \ref{algebra-lemma-limit-argument}，可约化到如下情形：
$R$ 和 $S$ 是代数闭域 $k$ 上的有限型整环。

\medskip\noindent
注意，环映射 $R \to R \otimes_k S$ 是有限表示的且平坦。
此外，对 $R$ 的每个极大理想 $\mathfrak m$，有
$(R \otimes_k S) \otimes_R R/\mathfrak m \cong S$，因为由希尔伯特零点定理
\ref{algebra-theorem-nullstellensatz}，$k \cong R/\mathfrak m$。
而且，由于 $\Spec(R)$ 是雅可布森空间，极大理想之集在 $R$ 的谱中稠密；
参见引理 \ref{algebra-lemma-finite-type-field-Jacobson}。
因此，可将引理 \ref{algebra-lemma-flat-fibres-irreducible} 应用于环映射
$R \to R \otimes_k S$，从而如所需，$R \otimes_k S$ 的谱不可约。
\end{proof}

\begin{lemma}
\label{algebra-lemma-geometrically-irreducible}
设 $k$ 为域，$R$ 为 $k$-代数。下列条件等价：
\begin{enumerate}
\item 对每个域扩张 $k'/k$，$R \otimes_k k'$ 的谱不可约；
\item 对每个有限可分域扩张 $k'/k$，$R \otimes_k k'$ 的谱不可约；
\item $R \otimes_k \overline{k}$ 的谱不可约，其中 $\overline{k}$
是 $k$ 的可分代数闭包；
\item $R \otimes_k \overline{k}$ 的谱不可约，其中 $\overline{k}$
是 $k$ 的代数闭包。
\end{enumerate}
\end{lemma}

\begin{proof}
显然，(1) 推出 (2)。

\medskip\noindent
假设 (2)，并令 $\overline{k}$ 为 $k$ 的可分代数闭包。
假设 $\mathfrak q_i \subset R \otimes_k \overline{k}$，$i = 1, 2$
是两个极小素理想。对每个有限子扩张 $\overline{k}/k'/k$，
扩张 $k'/k$ 可分，并且环映射
$R \otimes_k k' \to R \otimes_k \overline{k}$ 平坦。
所以 $\mathfrak p_i = (R \otimes_k k') \cap \mathfrak q_i$
是极小素理想（因为由引理 \ref{algebra-lemma-flat-going-down}，
平坦环映射满足下降）。于是由假设 (2)，
$\mathfrak p_1 = \mathfrak p_2$。由于 $\overline{k} = \bigcup k'$，
可知 $\mathfrak q_1 = \mathfrak q_2$。因此
$\Spec(R \otimes_k \overline{k})$ 不可约。

\medskip\noindent
假设 (3)，并令 $\overline{k}$ 为 $k$ 的代数闭包。
设 $\overline{k}/\overline{k}'/k$ 为相应的 $k$ 的可分代数闭包。
则 $\overline{k}/\overline{k}'$ 是纯不可分扩张（在正特征情形），
或是平凡扩张。因此，例如由引理 \ref{algebra-lemma-p-ring-map}，映射
$R \otimes_k \overline{k}' \to R \otimes_k \overline{k}$
诱导谱之间的同胚。故 (4) 成立。

\medskip\noindent
假设 (4)。设 $k'/k$ 为任意域扩张，并令 $\overline{k}$ 为 $k$ 的代数闭包。
可以选取一个域 $F$，使得 $k'$ 和 $\overline{k}$ 都同构于 $F$ 的子域。于是
$$
R \otimes_k F = (R \otimes_k \overline{k}) \otimes_{\overline{k}} F
$$
故由引理 \ref{algebra-lemma-separably-closed-irreducible}，
$R \otimes_k F$ 有唯一的极小素理想。最后，环映射
$R \otimes_k k' \to R \otimes_k F$ 平坦且单射，因此
$R \otimes_k k'$ 的任一极小素理想都是 $R \otimes_k F$ 的
某个极小素理想的像（由引理
\ref{algebra-lemma-injective-minimal-primes-in-image} 和下降）。
所以这样的极小素理想只有一个，证明完毕。
\end{proof}

\begin{definition}
\label{algebra-definition-geometrically-irreducible}
设 $k$ 为域，$S$ 为 $k$-代数。若对每个域扩张 $k'/k$，
$S \otimes_k k'$ 的谱不可约\footnote{不可约空间非空。}，
则称 $S$ {it 在 $k$ 上几何不可约}。
\end{definition}

\noindent
由引理 \ref{algebra-lemma-geometrically-irreducible}，只须对有限可分域扩张
$k'/k$ 检验这一点，或者取 $k'$ 等于 $k$ 的可分代数闭包加以检验。

\begin{lemma}
\label{algebra-lemma-separably-closed-irreducible-implies-geometric}
设 $k$ 为域，$R$ 为 $k$-代数。若 $k$ 是可分闭域，则
$R$ 在 $k$ 上几何不可约，当且仅当 $R$ 的谱不可约。
\end{lemma}

\begin{proof}
由定义 \ref{algebra-definition-geometrically-irreducible} 后的注记立即得到。
\end{proof}

\begin{lemma}
\label{algebra-lemma-subalgebra-geometrically-irreducible}
设 $k$ 为域，$S$ 为 $k$-代数。
\begin{enumerate}
\item 若 $S$ 在 $k$ 上几何不可约，则每个 $k$-子代数也几何不可约。
\item 若 $S$ 的所有有限生成 $k$-子代数都几何不可约，
则 $S$ 几何不可约。
\item 几何不可约 $k$-代数的有向余极限几何不可约。
\end{enumerate}
\end{lemma}

\begin{proof}
设 $S' \subset S$ 为子代数。那么，对任意扩张 $k'/k$，环映射
$S' \otimes_k k' \to S \otimes_k k'$ 仍是单射。
故 (1) 由引理 \ref{algebra-lemma-injective-minimal-primes-in-image} 得出
（并使用连续映射下不可约空间的像不可约这一事实）。
第二个和第三个性质由张量积与余极限交换这一事实得出。
\end{proof}

\begin{lemma}
\label{algebra-lemma-geometrically-irreducible-any-base-change}
设 $k$ 为域，$S$ 为几何不可约的 $k$-代数，$R$ 为任意 $k$-代数。映射
$$
\Spec(R \otimes_k S) \longrightarrow \Spec(R)
$$
诱导不可约分支之间的双射。
\end{lemma}

\begin{proof}
回忆，不可约分支对应于极小素理想
（引理 \ref{algebra-lemma-irreducible}）。由于 $R \to R \otimes_k S$ 平坦，
由下降（引理 \ref{algebra-lemma-flat-going-down}），
$R \otimes_k S$ 的任一极小素理想都位于 $R$ 的某个极小素理想上。
反之，若 $\mathfrak p \subset R$ 为（极小）素理想，则由
$R \to R \otimes_k S$ 的平坦性，有
$$
R \otimes_k S/\mathfrak p(R \otimes_k S)
=
(R/\mathfrak p) \otimes_k S
\subset
\kappa(\mathfrak p) \otimes_k S
$$
由假设，环 $\kappa(\mathfrak p) \otimes_k S$ 的谱不可约。
因此，$R \otimes_k S/\mathfrak p(R \otimes_k S)$ 有唯一的极小素理想
（引理 \ref{algebra-lemma-injective-minimal-primes-in-image}）。
换言之，不可约集 $V(\mathfrak p)$ 的逆像不可约。本引理因而得证。
\end{proof}

\noindent
下面对几何不可约域扩张的概念作一些注记。

\begin{lemma}
\label{algebra-lemma-field-extension-geometrically-irreducible}
设 $K/k$ 为域扩张。若 $k$ 在 $K$ 中代数闭，则 $K$ 在 $k$ 上几何不可约。
\end{lemma}

\begin{proof}
假设 $k$ 在 $K$ 中代数闭。由定义
\ref{algebra-definition-geometrically-irreducible} 和引理
\ref{algebra-lemma-geometrically-irreducible}，只须证明：
对每个有限可分扩张 $k'/k$，$K \otimes_k k'$ 的谱不可约。
设 $k'$ 由 $\alpha \in k'$ 在 $k$ 上生成；参见域，引理
\ref{fields-lemma-primitive-element}。设
$P = T^d + a_1 T^{d - 1} + \ldots + a_d \in k[T]$ 为 $\alpha$ 的极小多项式。
则 $K \otimes_k k' \cong K[T]/(P)$。
$K[T]/(P)$ 的谱可约的唯一可能，是 $P$ 在 $K[T]$ 中可约。
假设 $P = P_1 P_2$ 是 $K[T]$ 中的非平凡因式分解，以导出矛盾。
由引理 \ref{algebra-lemma-polynomials-divide}，$P_1$ 和 $P_2$ 的系数在 $k$ 上代数。
我们的假设表明 $P_1$ 和 $P_2$ 的系数属于 $k$，
这与 $P$ 在 $k$ 上不可约矛盾。
\end{proof}

\begin{lemma}
\label{algebra-lemma-geometrically-irreducible-transitive}
设 $K/k$ 为几何不可约域扩张，$S$ 为几何不可约的 $K$-代数。
则 $S$ 在 $k$ 上几何不可约。
\end{lemma}

\begin{proof}
由定义 \ref{algebra-definition-geometrically-irreducible} 和引理
\ref{algebra-lemma-geometrically-irreducible}，只须证明：
对每个有限可分扩张 $k'/k$，$S \otimes_k k'$ 的谱不可约。
由于 $K$ 在 $k$ 上几何不可约，可知 $K' = K \otimes_k k'$
是 $K$ 的有限可分域扩张。因为假设 $S$ 在 $K$ 上几何不可约，
所以 $S \otimes_k k' = S \otimes_K K'$ 的谱不可约。
\end{proof}

\begin{lemma}
\label{algebra-lemma-geometrically-irreducible-base-change-transcendental}
设 $K/k$ 为域扩张。下列条件等价：
\begin{enumerate}
\item $K$ 在 $k$ 上几何不可约；
\item 纯超越扩张所诱导的扩张 $K(t)/k(t)$ 几何不可约。
\end{enumerate}
\end{lemma}

\begin{proof}
假设 (1)。以 $\Omega$ 表示 $k(t)$ 的一个代数闭包。
由定义 \ref{algebra-definition-geometrically-irreducible}，可知
$$
K \otimes_k \Omega = K \otimes_k k(t) \otimes_{k(t)} \Omega
$$
的谱不可约。由于 $K(t)$ 是 $K \otimes_k k(T)$ 的局部化，
可知 $K(t) \otimes_{k(t)} \Omega$ 的谱不可约。
所以由引理 \ref{algebra-lemma-geometrically-irreducible}，
$K(t)/k(t)$ 几何不可约。

\medskip\noindent
假设 (2)。设 $k'/k$ 为域扩张。须证明 $K \otimes_k k'$ 有唯一的极小素理想。
已知
$$
K(t) \otimes_{k(t)} k'(t)
$$
的谱不可约，也就是说，它有唯一的极小素理想。
由于存在单射
$K \otimes_k k' \to K(t) \otimes_{k(t)} k'(t)$（省略细节），
结论由引理 \ref{algebra-lemma-injective-minimal-primes-in-image} 和
\ref{algebra-lemma-minimal-prime-image-minimal-prime} 得出。
\end{proof}

\begin{lemma}
\label{algebra-lemma-geometrically-irreducible-add-transcendental}
设 $K/L/M$ 为域塔，且 $L/M$ 几何不可约。
设 $x \in K$ 在 $L$ 上超越。则 $L(x)/M(x)$ 几何不可约。
\end{lemma}

\begin{proof}
这由引理 \ref{algebra-lemma-geometrically-irreducible-base-change-transcendental}
得出，因为域 $L(x)$ 和 $M(x)$ 分别是 $L$ 和 $M$ 的纯超越扩张。
\end{proof}

\begin{lemma}
\label{algebra-lemma-geometrically-irreducible-separable-elements}
设 $K/k$ 为域扩张。下列条件等价：
\begin{enumerate}
\item $K/k$ 几何不可约；
\item 每个在 $k$ 上可分代数的元素 $\alpha \in K$ 都属于 $k$。
\end{enumerate}
\end{lemma}

\begin{proof}
假设 (1)，并设 $\alpha \in K$ 在 $k$ 上可分代数。
则 $k' = k(\alpha)$ 是 $k$ 的有限可分扩张，并且包含在 $K$ 中。
由引理 \ref{algebra-lemma-subalgebra-geometrically-irreducible}，
扩张 $k'/k$ 几何不可约。特别地，$k' \otimes_k \overline{k}$
的谱不可约（因而，若它是域的积，则恰有一个因子）。
由域，引理 \ref{fields-lemma-finite-separable-tensor-alg-closed}，
可知 $\Hom_k(k', \overline{k})$ 有一个元素；进而由域，引理
\ref{fields-lemma-separable-equality}，可知 $k' = k$。故 (2) 成立。

\medskip\noindent
假设 (2)。令 $k' \subset K$ 为由在 $k$ 上代数的元素组成的子域。
由引理 \ref{algebra-lemma-field-extension-geometrically-irreducible}，
扩张 $K/k'$ 几何不可约。由假设，$k'/k$ 是纯不可分扩张。
由引理 \ref{algebra-lemma-p-ring-map}，扩张 $k'/k$ 几何不可约。
所以由引理 \ref{algebra-lemma-geometrically-irreducible-transitive}，
$K/k$ 几何不可约。
\end{proof}

\begin{lemma}
\label{algebra-lemma-make-geometrically-irreducible}
设 $K/k$ 为域扩张。考虑由在 $k$ 上可分代数的元素组成的子扩张
$K/k'/k$。则 $K$ 在 $k'$ 上几何不可约。若 $K/k$ 为有限生成域扩张，
则 $[k' : k] < \infty$。
\end{lemma}

\begin{proof}
第一个陈述立即由引理
\ref{algebra-lemma-geometrically-irreducible-separable-elements} 以及下述事实得出：
由可分代数扩张的传递性，在 $k'$ 上可分代数的元素都属于 $k'$；
参见域，引理 \ref{fields-lemma-separable-permanence}。
若 $K/k$ 有限生成，则由域，引理
\ref{fields-lemma-algebraic-closure-in-finitely-generated}，
$k'$ 在 $k$ 上有限。
\end{proof}

\begin{lemma}
\label{algebra-lemma-Galois-orbit}
设 $K/k$ 为域扩张，$\overline{k}/k$ 为可分代数闭包。
则 $\text{Gal}(\overline{k}/k)$ 在
$\overline{k} \otimes_k K$ 的素理想上可迁地作用。
\end{lemma}

\begin{proof}
设 $K/k'/k$ 为引理 \ref{algebra-lemma-make-geometrically-irreducible}
中得到的子扩张。注意，由于 $k \subset \overline{k}$ 是整的，
$\overline{k} \otimes_k K$ 和 $\overline{k} \otimes_k k'$ 的所有素理想
都是极大理想；参见引理 \ref{algebra-lemma-integral-no-inclusion}。
由引理 \ref{algebra-lemma-geometrically-irreducible-any-base-change}，映射
$$
\Spec(\overline{k} \otimes_k K) \to \Spec(\overline{k} \otimes_k k')
$$
是双射，因为 (1) 所有素理想都是极小素理想，(2)
$\overline{k} \otimes_k K = (\overline{k} \otimes_k k') \otimes_{k'} K$，
并且 (3) $K$ 在 $k'$ 上几何不可约。
因此，只须证明 $\text{Gal}(\overline{k}/k)$ 在
$\overline{k} \otimes_k k'$ 的素理想上的作用可迁。

\medskip\noindent
由于 $\overline{k} \otimes_k k'$ 的每个素理想都是极大理想，
其剩余域同构于 $\overline{k}$。因此，
$\overline{k} \otimes_k k'$ 的素理想与
$\Hom_k(k', \overline{k})$ 的元素一一对应，其中
$\sigma \in \Hom_k(k', \overline{k})$ 对应于映射
$1 \otimes \sigma : \overline{k} \otimes_k k' \to \overline{k}$
的核 $\mathfrak p_\sigma$。特别地，
$\text{Gal}(\overline{k}/k)$ 如所需地在此集合上可迁作用。
\end{proof}





\section{几何连通代数}
\label{algebra-section-geometrically-connected}

\begin{lemma}
\label{algebra-lemma-separably-closed-connected}
设 $k$ 为可分闭域，$R$、$S$ 为 $k$-代数。若
$\Spec(R)$ 和 $\Spec(S)$ 连通，则
$\Spec(R \otimes_k S)$ 也连通。
\end{lemma}

\begin{proof}
回忆，$\Spec(R)$ 连通，当且仅当 $R$ 没有非平凡幂等元；
参见引理 \ref{algebra-lemma-characterize-spec-connected}。
所以由引理 \ref{algebra-lemma-limit-argument}，可以假设 $R$ 和 $S$ 都是
$k$ 上的有限型代数。此时 $R$ 和 $S$ 是诺特环，并且只有有限多个极小素理想；
参见引理 \ref{algebra-lemma-Noetherian-irreducible-components}。
因此，可对 $n + m$ 作归纳，其中 $n$、$m$ 分别是
$\Spec(R)$、$\Spec(S)$ 的不可约分支数。
当然，$n$ 或 $m$ 为零的情形是平凡的。若 $n = m = 1$，
即 $\Spec(R)$ 和 $\Spec(S)$ 都只有一个不可约分支，
则结论由引理 \ref{algebra-lemma-separably-closed-irreducible} 得出。
假设 $n > 1$。设 $\mathfrak p \subset R$ 为对应于不可约闭子集
$T \subset \Spec(R)$ 的极小素理想。令 $T' \subset \Spec(R)$ 为其余
$n - 1$ 个不可约分支之并。选取理想 $I \subset R$，使得
$T' = V(I) = \Spec(R/I)$（引理 \ref{algebra-lemma-spec-closed}）。
仔细选取极小素理想，还可以使 $T'$ 连通；参见拓扑，引理
\ref{topology-lemma-remove-irreducible-connected}。那么
$T \cup T' = \Spec(R)$，并且
$T \cap T' = V(\mathfrak p + I) = \Spec(R/(\mathfrak p + I))$
非空，因为假设 $\Spec(R)$ 连通。
$T$ 在 $\Spec(R \otimes_k S)$ 中的逆像是
$\Spec(R/\mathfrak p \otimes_k S)$，而 $T'$ 在
$\Spec(R \otimes_k S)$ 中的逆像是 $\Spec(R/I \otimes_k S)$。
由归纳假设，这两个空间都连通。$T \cap T'$ 的逆像是
$\Spec(R/(\mathfrak p + I) \otimes_k S)$，它非空。
因此 $\Spec(R \otimes_k S)$ 连通。
\end{proof}

\begin{lemma}
\label{algebra-lemma-geometrically-connected}
设 $k$ 为域，$R$ 为 $k$-代数。下列条件等价：
\begin{enumerate}
\item 对每个域扩张 $k'/k$，$R \otimes_k k'$ 的谱连通；
\item 对每个有限可分域扩张 $k'/k$，$R \otimes_k k'$ 的谱连通。
\end{enumerate}
\end{lemma}

\begin{proof}
对任意域扩张 $k'/k$，$R \otimes_k k'$ 的谱连通，
等价于 $R \otimes_k k'$ 没有非平凡幂等元；参见引理
\ref{algebra-lemma-characterize-spec-connected}。假设 (2)。
设 $k \subset \overline{k}$ 为 $k$ 的可分代数闭包。
由引理 \ref{algebra-lemma-limit-argument}，可知 (2) 等价于
$R \otimes_k \overline{k}$ 没有非平凡幂等元。
对任意域扩张 $k'/k$，存在域扩张 $\overline{k}'/\overline{k}$，
使得 $k' \subset \overline{k}'$。由引理
\ref{algebra-lemma-separably-closed-connected}，
$R \otimes_k \overline{k}'$ 没有非平凡幂等元。
若 $R \otimes_k k'$ 有非平凡幂等元，则
$R \otimes_k \overline{k}'$ 也有，矛盾。
\end{proof}

\begin{definition}
\label{algebra-definition-geometrically-connected}
设 $k$ 为域，$S$ 为 $k$-代数。若对每个域扩张 $k'/k$，
$S \otimes_k k'$ 的谱连通，则称 $S$ {it 在 $k$ 上几何连通}。
\end{definition}

\noindent
由引理 \ref{algebra-lemma-geometrically-connected}，只须对有限可分域扩张
$k'/k$ 检验这一点。

\begin{lemma}
\label{algebra-lemma-separably-closed-connected-implies-geometric}
设 $k$ 为域，$R$ 为 $k$-代数。若 $k$ 是可分闭域，则
$R$ 在 $k$ 上几何连通，当且仅当 $R$ 的谱连通。
\end{lemma}

\begin{proof}
由定义 \ref{algebra-definition-geometrically-connected} 后的注记立即得到。
\end{proof}

\begin{lemma}
\label{algebra-lemma-subalgebra-geometrically-connected}
设 $k$ 为域，$S$ 为 $k$-代数。
\begin{enumerate}
\item 若 $S$ 在 $k$ 上几何连通，则每个 $k$-子代数也几何连通。
\item 若 $S$ 的所有有限生成 $k$-子代数都几何连通，则 $S$ 几何连通。
\item 几何连通 $k$-代数的有向余极限几何连通。
\end{enumerate}
\end{lemma}

\begin{proof}
这由连通性的如下刻画得出：不存在非平凡幂等元。
第二个和第三个性质由张量积与余极限交换这一事实得出。
\end{proof}

\noindent
下面的引理将被更一般的簇，引理
\ref{varieties-lemma-bijection-connected-components} 所取代。

\begin{lemma}
\label{algebra-lemma-geometrically-connected-any-base-change}
设 $k$ 为域，$S$ 为几何连通的 $k$-代数，$R$ 为任意 $k$-代数。映射
$$
R \longrightarrow R \otimes_k S
$$
诱导幂等元之间的双射，而映射
$$
\Spec(R \otimes_k S) \longrightarrow \Spec(R)
$$
诱导连通分支之间的双射。
\end{lemma}

\begin{proof}
第二个断言由第一个断言结合引理 \ref{algebra-lemma-connected-component} 得出。
由引理 \ref{algebra-lemma-subalgebra-geometrically-connected} 和
\ref{algebra-lemma-limit-argument}，可以假设 $R$ 和 $S$ 都是 $k$ 上的有限型代数。
于是 $R \otimes_k S$ 也是 $k$ 上的有限型代数。
注意，在此情形下，所有这些环都是诺特环，因而它们的谱只有有限多个连通分支
（因为只有有限多个不可约分支；参见引理
\ref{algebra-lemma-Noetherian-irreducible-components}）。
特别地，这里涉及的所有连通分支都是开集！因此，由引理
\ref{algebra-lemma-disjoint-implies-product} 可知，在此情形下，
本引理的第一个陈述等价于第二个陈述。下面证明第二个陈述。
由于代数 $S$ 几何连通且非零，映射
$X = \Spec(R \otimes_k S) \to \Spec(R) = Y$ 的所有纤维都连通且非空。
此外，由于 $R \to R \otimes_k S$ 是有限表示的平坦映射，
映射 $X \to Y$ 是开映射
（命题 \ref{algebra-proposition-fppf-open}）。拓扑，引理
\ref{topology-lemma-connected-fibres-connected-components} 表明
$X \to Y$ 诱导连通分支之间的双射。
\end{proof}




\section{几何整代数}
\label{algebra-section-geometrically-integral}

\noindent
定义如下。

\begin{definition}
\label{algebra-definition-geometrically-integral}
设 $k$ 为域，$S$ 为 $k$-代数。若对每个域扩张 $k'/k$，
环 $S \otimes_k k'$ 都是整环，则称 $S$ {it 在 $k$ 上几何整}。
\end{definition}

\noindent
关于几何整代数的任何问题都可以转化为关于几何约化且不可约代数的问题。

\begin{lemma}
\label{algebra-lemma-geometrically-integral}
设 $k$ 为域，$S$ 为 $k$-代数。此时，$S$ 在 $k$ 上几何整，
当且仅当 $S$ 在 $k$ 上既几何不可约又几何约化。
\end{lemma}

\begin{proof}
省略。
\end{proof}

\begin{lemma}
\label{algebra-lemma-characterize-geometrically-integral}
设 $k$ 为域，$S$ 为 $k$-代数。下列条件等价：
\begin{enumerate}
\item $S$ 在 $k$ 上几何整；
\item 对每个有限域扩张 $k'/k$，环 $S \otimes_k k'$ 是整环；
\item $S \otimes_k \overline{k}$ 是整环，其中 $\overline{k}$
是 $k$ 的代数闭包。
\end{enumerate}
\end{lemma}

\begin{proof}
由引理 \ref{algebra-lemma-geometrically-integral}、
\ref{algebra-lemma-geometrically-reduced-finite-purely-inseparable-extension} 和
\ref{algebra-lemma-geometrically-irreducible} 得出。
\end{proof}

\begin{lemma}
\label{algebra-lemma-geometrically-integral-any-integral-base-change}
设 $k$ 为域，$S$ 为几何整的 $k$-代数。
设 $R$ 为 $k$-代数且为整环。则 $R \otimes_k S$ 是整环。
\end{lemma}

\begin{proof}
由引理 \ref{algebra-lemma-geometrically-reduced-any-reduced-base-change}，
环 $R \otimes_k S$ 是约化的；由引理
\ref{algebra-lemma-geometrically-irreducible-any-base-change}，
环 $R \otimes_k S$ 不可约（其谱只有一个不可约分支），
所以 $R \otimes_k S$ 是整环。
\end{proof}





\section{赋值环}
\label{algebra-section-valuation-rings}

\noindent
先给出一些定义。

\begin{definition}
\label{algebra-definition-valuation-ring}
赋值环：
\begin{enumerate}
\item 设 $K$ 为域，$A$、$B$ 为包含于 $K$ 中的局部环。
若 $A \subset B$ 且 $\mathfrak m_A = A \cap \mathfrak m_B$，
则称 $B$ {it 支配} $A$。
\item 设 $A$ 为环。若 $A$ 是局部整环，并且在 $A$ 的分式域所含局部环中，
$A$ 关于支配关系是极大的，则称 $A$ 为{it 赋值环}。
\item 设 $A$ 为分式域为 $K$ 的赋值环。若 $R \subset K$ 是 $K$ 的子环，
且 $R \subset A$，则称 $A$ {it 以 $R$ 为中心}。
\end{enumerate}
\end{definition}

\noindent
依此定义，域是赋值环。

\begin{lemma}
\label{algebra-lemma-dominate}
设 $K$ 为域，$A \subset K$ 为局部子环。则存在分式域为 $K$ 且支配 $A$ 的赋值环。
\end{lemma}

\begin{proof}
用支配关系将 $K$ 的局部子环之集视为偏序集。
假设 $\{A_i\}_{i \in I}$ 是 $K$ 的一族全序局部子环。
则 $B = \bigcup A_i$ 是支配所有 $A_i$ 的局部子环。
因此，由佐恩引理，只须证明：若 $A \subset K$ 是分式域不等于 $K$ 的局部环，
则存在局部环 $B \subset K$，$B \not = A$，且它支配 $A$。

\medskip\noindent
取不属于 $A$ 的分式域的 $t \in K$。
若 $t$ 在 $A$ 上超越，则 $A[t] \subset K$，从而
$A[t]_{(t, \mathfrak m)} \subset K$ 是不同于 $A$ 且支配 $A$ 的局部环。
假设 $t$ 在 $A$ 上代数。则对某个非零 $a \in A$，元素 $at$ 在 $A$ 上整。
此时，由 $A$ 和 $ta$ 生成的子环 $A' \subset K$ 在 $A$ 上有限。
由引理 \ref{algebra-lemma-integral-overring-surjective}，存在位于
$\mathfrak m$ 上的素理想 $\mathfrak m' \subset A'$。
于是 $A'_{\mathfrak m'}$ 支配 $A$。若 $A = A'_{\mathfrak m'}$，
则 $t$ 属于 $A$ 的分式域，这与假设矛盾。因此如所需，
$A \not = A'_{\mathfrak m'}$。
\end{proof}

\begin{lemma}
\label{algebra-lemma-valuation-ring-normal}
设 $A$ 为赋值环。则 $A$ 是正规整环。
\end{lemma}

\begin{proof}
假设 $x$ 属于 $A$ 的分式域，且在 $A$ 上整。
以 $A'$ 表示由 $A$ 和 $x$ 生成的 $K$ 的子环。
由于 $A\subset A'$ 是整扩张，由引理
\ref{algebra-lemma-integral-overring-surjective}，存在位于
$\mathfrak m$ 上的素理想 $\mathfrak m' \subset A'$。
于是 $A'_{\mathfrak m'}$ 支配 $A$。由于 $A$ 是赋值环，
可知 $A=A'_{\mathfrak m'}$，从而 $x\in A$。
\end{proof}


\begin{lemma}
\label{algebra-lemma-valuation-ring-x-or-x-inverse}
设 $A$ 为极大理想是 $\mathfrak m$、分式域是 $K$ 的赋值环。
令 $x \in K$。则 $x \in A$ 或 $x^{-1} \in A$，或二者皆成立。
\end{lemma}

\begin{proof}
假设 $x$ 不属于 $A$。以 $A'$ 表示由 $A$ 和 $x$ 生成的 $K$ 的子环。
由于 $A$ 是赋值环，$A'$ 中没有位于 $\mathfrak m$ 上的素理想。
因为 $\mathfrak m$ 是极大理想，可知 $V(\mathfrak m A') = \emptyset$。
于是由引理 \ref{algebra-lemma-Zariski-topology}，$\mathfrak m A' = A'$。
故可写成 $1 = \sum_{i = 0}^d t_i x^i$，其中 $t_i \in \mathfrak m$。
这推出 $(1 - t_0) (x^{-1})^d - \sum t_i (x^{-1})^{d - i} = 0$。
特别地，$x^{-1}$ 在 $A$ 上整，所以由引理
\ref{algebra-lemma-valuation-ring-normal}，$x^{-1} \in A$。
\end{proof}

\begin{lemma}
\label{algebra-lemma-x-or-x-inverse-valuation-ring}
设 $A \subset K$ 为域 $K$ 的子环，并且对所有 $x \in K$，
$x \in A$ 或 $x^{-1} \in A$，或二者皆成立。
则 $A$ 是分式域为 $K$ 的赋值环。
\end{lemma}

\begin{proof}
若 $A$ 不等于 $K$，则 $A$ 不是域，并且存在非零极大理想
$\mathfrak m$。若 $\mathfrak m'$ 是另一个极大理想，则选取
$x, y \in A$，使得 $x \in \mathfrak m$、$y \not \in \mathfrak m$、
$x \not \in \mathfrak m'$ 且 $y \in \mathfrak m'$。
那么 $x/y \in A$ 和 $y/x \in A$ 都不成立，这与本引理的假设矛盾。
所以 $A$ 是局部环。假设 $A'$ 是包含于 $K$ 中且支配 $A$ 的局部环。
令 $x \in A'$。须证明 $x \in A$。若不然，则 $x^{-1} \in A$，
而且当然 $x^{-1} \in \mathfrak m_A$。但是这样就有
$x^{-1} \in \mathfrak m_{A'}$，与 $x \in A'$ 矛盾。
\end{proof}

\begin{lemma}
\label{algebra-lemma-colimit-valuation-rings}
\begin{slogan}
赋值环在滤过直极限下稳定
\end{slogan}
设 $I$ 为有向集，$(A_i, \varphi_{ij})$ 为 $I$ 上的赋值环系统。
则 $A = \colim A_i$ 是赋值环。
\end{lemma}

\begin{proof}
显然 $A$ 是整环。设 $a, b \in A$。
引理 \ref{algebra-lemma-x-or-x-inverse-valuation-ring} 表明，须证明在 $A$ 中
$a | b$ 或 $b | a$。选取充分大的 $i$，使得存在映到 $a, b$ 的
$a_i, b_i \in A_i$。将引理 \ref{algebra-lemma-valuation-ring-x-or-x-inverse}
应用于 $A_i$ 中的 $a_i, b_i$，即得 $A$ 中 $a, b$ 的结论。
\end{proof}

\begin{lemma}
\label{algebra-lemma-valuation-ring-cap-field}
设 $L/K$ 为域扩张。若 $B \subset L$ 是赋值环，
则 $A = K \cap B$ 是赋值环。
\end{lemma}

\begin{proof}
可以把 $L$ 替换为 $B$ 的分式域 $F$，并把 $K$ 替换为
$K \cap F$。然后本引理由引理
\ref{algebra-lemma-valuation-ring-x-or-x-inverse} 和
\ref{algebra-lemma-x-or-x-inverse-valuation-ring} 结合得出。
\end{proof}

\begin{lemma}
\label{algebra-lemma-valuation-ring-cap-field-finite}
设 $L/K$ 为代数域扩张。若 $B \subset L$ 是分式域为 $L$
且不是域的赋值环，则 $A = K \cap B$ 是赋值环且不是域。
\end{lemma}

\begin{proof}
由引理 \ref{algebra-lemma-valuation-ring-cap-field}，环 $A$ 是赋值环。
若 $A$ 是域，则 $A = K$。于是 $A = K \subset B$ 是整扩张，
故 $B$ 的素理想之间不存在真包含
（引理 \ref{algebra-lemma-integral-no-inclusion}）。
这与 $B$ 是局部整环而不是域的假设矛盾。
\end{proof}

\begin{lemma}
\label{algebra-lemma-make-valuation-rings}
设 $A$ 为赋值环。对任意素理想 $\mathfrak p \subset A$，
商环 $A/\mathfrak p$ 是赋值环。局部化 $A_\mathfrak p$ 也如此；
事实上，$A$ 的任意局部化都是赋值环。
\end{lemma}

\begin{proof}
使用引理 \ref{algebra-lemma-x-or-x-inverse-valuation-ring}
所给出的赋值环刻画。
\end{proof}

\begin{lemma}
\label{algebra-lemma-stack-valuation-rings}
设 $A'$ 为剩余域是 $K$ 的赋值环，$A$ 为分式域是 $K$ 的赋值环。则
$C = \{\lambda \in A' \mid \lambda \bmod \mathfrak m_{A'} \in A\}$
是赋值环。
\end{lemma}

\begin{proof}
注意，$\mathfrak m_{A'} \subset C$ 且 $C/\mathfrak m_{A'} = A$。
特别地，$C$ 的分式域等于 $A'$ 的分式域。
将使用引理 \ref{algebra-lemma-x-or-x-inverse-valuation-ring} 的判据来证明本引理。
令 $x$ 为 $C$ 的分式域中的元素。由该引理，可以假设 $x \in A'$。
若 $x \in \mathfrak m_{A'}$，则 $x \in C$。若非如此，则 $x$ 是 $A'$ 的单位，
并且还有 $x^{-1} \in A'$。因此，再次由该引理，$x$ 或 $x^{-1}$
映到 $A$ 的某个元素。
\end{proof}

\begin{lemma}
\label{algebra-lemma-find-valuation-rings}
设 $A$ 为分式域是 $K$ 的正规整环。
\begin{enumerate}
\item 对每个 $x \in K$，$x \not \in A$，存在分式域是 $K$ 的赋值环
$A \subset V \subset K$，使得 $x \not \in V$。
\item 若 $A$ 是局部环，还可以选取支配 $A$ 的 $V$。
\end{enumerate}
换言之，$A$ 是 $K$ 中所有包含 $A$ 的赋值环之交；若 $A$ 是局部环，
则 $A$ 是 $K$ 中所有支配 $A$ 的赋值环之交。
\end{lemma}

\begin{proof}
假设 $x \in K$，$x \not \in A$。考虑 $B = A[x^{-1}]$。
则 $x \not \in B$。事实上，若
$x = a_0 + a_1x^{-1} + \ldots + a_d x^{-d}$，
则 $x^{d + 1} - a_0x^d - \ldots - a_d = 0$，于是 $x$ 在 $A$ 上整，
这与 $A$ 正规矛盾。因此 $x^{-1}$ 不是 $B$ 的单位。
所以 $V(x^{-1}) \subset \Spec(B)$ 非空
（引理 \ref{algebra-lemma-Zariski-topology}），并且可以选取满足
$x^{-1} \in \mathfrak p$ 的素理想 $\mathfrak p \subset B$。
选取支配 $B_\mathfrak p$ 的赋值环 $V \subset K$
（引理 \ref{algebra-lemma-dominate}）。由于 $x^{-1} \in \mathfrak m_V$，
有 $x \not \in V$。

\medskip\noindent
若 $A$ 是局部环，则断言
$x^{-1} B + \mathfrak m_A B \not = B$。事实上，若
$1 = (a_0 + a_1x^{-1} + \ldots + a_d x^{-d})x^{-1} +
a'_0 + \ldots + a'_d x^{-d}$，其中 $a_i \in A$ 且
$a'_i \in \mathfrak m_A$，则会得到
$$
(1 - a'_0) x^{d + 1} - (a_0 + a'_1) x^d - \ldots - a_d = 0
$$
由于 $a'_0 \in \mathfrak m_A$，可知 $1 - a'_0$ 是 $A$ 的单位，
从而像先前一样，$x$ 会在 $A$ 上整，导致矛盾。
再选取素理想 $\mathfrak p \supset x^{-1} B + \mathfrak m_A B$，
即可得到支配 $A$ 的 $V$。
\end{proof}

\noindent
{it 全序阿贝尔群}是二元组 $(\Gamma, \geq)$，其中 $\Gamma$ 是赋有全序
$\geq$ 的阿贝尔群，并且对所有 $\gamma, \gamma', \gamma'' \in \Gamma$，
$\gamma \geq \gamma' \Rightarrow
\gamma + \gamma'' \geq \gamma' + \gamma''$。

\begin{lemma}
\label{algebra-lemma-valuation-group}
设 $A$ 为分式域是 $K$ 的赋值环。置
$\Gamma = K^*/A^*$（群运算按加法记）。对 $\gamma, \gamma' \in \Gamma$，
定义 $\gamma \geq \gamma'$，当且仅当 $\gamma - \gamma'$ 属于
$A - \{0\} \to \Gamma$ 的像。则 $(\Gamma, \geq)$ 是全序阿贝尔群。
\end{lemma}

\begin{proof}
省略；但它容易由引理 \ref{algebra-lemma-valuation-ring-x-or-x-inverse} 得出。
注意，当 $A = K$ 时，得到赋有唯一全序的零群 $\Gamma = \{0\}$。
\end{proof}

\begin{definition}
\label{algebra-definition-value-group}
设 $A$ 为赋值环。
\begin{enumerate}
\item 引理 \ref{algebra-lemma-valuation-group} 的全序阿贝尔群
$(\Gamma, \geq)$ 称为赋值环 $A$ 的{it 值群}。
\item 映射 $v : A - \{0\} \to \Gamma$ 以及
$v : K^* \to \Gamma$ 都称为与 $A$ 相伴的{it 赋值}。
\item 若 $\Gamma \cong \mathbf{Z}$，则称赋值环 $A$ 为{it 离散赋值环}。
\end{enumerate}
\end{definition}

\noindent
注意，若 $\Gamma \cong \mathbf{Z}$，则满足 $1 \geq 0$ 的这种同构是唯一的。
按此方式选取同构后，该次序即为整数的通常次序。

\begin{lemma}
\label{algebra-lemma-properties-valuation}
设 $A$ 为赋值环。赋值 $v : A -\{0\} \to \Gamma_{\geq 0}$
具有下列性质：
\begin{enumerate}
\item $v(a) = 0 \Leftrightarrow a \in A^*$；
\item $v(ab) = v(a) + v(b)$；
\item 当 $a + b \not = 0$ 时，$v(a + b) \geq \min(v(a), v(b))$。
\end{enumerate}
\end{lemma}

\begin{proof}
省略。
\end{proof}

\begin{lemma}
\label{algebra-lemma-characterize-valuation-ring}
设 $A$ 为环。下列条件等价：
\begin{enumerate}
\item $A$ 是赋值环；
\item $A$ 是局部整环，并且 $A$ 的每个有限生成理想都是主理想。
\end{enumerate}
\end{lemma}

\begin{proof}
假设 $A$ 是赋值环，并设 $f_1, \ldots, f_n \in A$。
选取 $i$，使得 $v(f_i)$ 在所有 $v(f_j)$ 中最小。
则 $(f_i) = (f_1, \ldots, f_n)$。反之，假设 $A$ 是局部整环，
并且 $A$ 的每个有限生成理想都是主理想。
取 $f, g \in A$，并写成 $(f, g) = (h)$。则对某些
$a, b, c, d \in A$，有 $f = ah$、$g = bh$ 且 $h = cf + dg$。
所以 $ac + bd = 1$，从而 $a$ 或 $b$ 是单位；也就是说，
$g/f$ 或 $f/g$ 属于 $A$。由引理
\ref{algebra-lemma-x-or-x-inverse-valuation-ring}，这表明 $A$ 是赋值环。
\end{proof}

\begin{lemma}
\label{algebra-lemma-valuation-valuation-ring}
设 $(\Gamma, \geq)$ 为全序阿贝尔群，$K$ 为域，
$v : K^* \to \Gamma$ 为阿贝尔群同态，并且对满足
$a, b, a + b$ 均非零的 $a, b \in K$，有
$v(a + b) \geq \min(v(a), v(b))$。则
$$
A =
\{
x \in K \mid x = 0 \text{ 或 } v(x) \geq 0
\}
$$
是值群为 $\Im(v) \subset \Gamma$ 的赋值环，其极大理想为
$$
\mathfrak m =
\{
x \in K \mid x = 0 \text{ 或 } v(x) > 0
\}
$$
而其单位群为
$$
A^* =
\{
x \in K^* \mid v(x) = 0
\}.
$$
\end{lemma}

\begin{proof}
省略。
\end{proof}

\noindent
设 $(\Gamma, \geq)$ 为全序阿贝尔群。$\Gamma$ 的{it 理想}是子集
$I \subset \Gamma$，使得 $I$ 的所有元素都 $\geq 0$，并且
$\gamma \in I$、$\gamma' \geq \gamma$ 推出 $\gamma' \in I$。
若这种理想满足 $0 \not \in I$ 以及
$\gamma + \gamma' \in I, \gamma, \gamma' \geq 0
\Rightarrow \gamma \in I \text{ 或 } \gamma' \in I$，
则称它为{it 素理想}。

\begin{lemma}
\label{algebra-lemma-ideals-valuation-ring}
设 $A$ 为赋值环。$A$ 中的理想与 $\Gamma$ 的理想按 $1 - 1$ 对应。
此双射保持包含关系，并将素理想映到素理想。
\end{lemma}

\begin{proof}
省略。
\end{proof}

\begin{lemma}
\label{algebra-lemma-valuation-ring-Noetherian-discrete}
赋值环是诺特环，当且仅当它是离散赋值环或域。
\end{lemma}

\begin{proof}
假设 $A$ 是离散赋值环，其赋值
$v : A \setminus \{0\} \to \mathbf{Z}$ 已正规化，使得
$\Im(v) = \mathbf{Z}_{\geq 0}$。由引理
\ref{algebra-lemma-ideals-valuation-ring}，$A$ 的理想是各子集
$I_n = \{0\} \cup v^{-1}(\mathbf{Z}_{\geq n})$。
显然，满足 $v(x) = n$ 的任意元素 $x \in A$ 都生成 $I_n$。
因此 $A$ 是主理想整环，当然是诺特环。

\medskip\noindent
假设 $A$ 是值群为 $\Gamma$ 的诺特赋值环。
由引理 \ref{algebra-lemma-ideals-valuation-ring}，可知 $\Gamma$ 的理想满足升链条件。
可以假设 $A$ 不是域；也就是说，存在 $\gamma \in \Gamma$，使得
$\gamma > 0$。将升链条件应用于各子集
$\gamma + \Gamma_{\geq 0}$，其中 $\gamma > 0$，可知存在大于 $0$
的最小元素 $\gamma_0$。设 $\gamma \in \Gamma$ 为满足
$\gamma > 0$ 的元素。考虑元素序列 $\gamma$、
$\gamma - \gamma_0$、$\gamma - 2\gamma_0$ 等。
由升链条件，它们不可能全都 $> 0$。设
$\gamma - n \gamma_0$ 是其中最后一个 $\geq 0$ 的元素。
由 $\gamma_0$ 的极小性，$0 = \gamma - n \gamma_0$。
因此如所需，$\Gamma$ 是循环群。
\end{proof}























\section{再论诺特环}
\label{algebra-section-Noetherian-again}


\begin{lemma}
\label{algebra-lemma-Noetherian-basic}
设 $R$ 为诺特环。任意有限 $R$-模都是有限表示的。
有限 $R$-模的任意子模都是有限的。
有限 $R$-模的 $R$-子模满足升链条件。
\end{lemma}

\begin{proof}
先证明有限 $R$-模 $M$ 的任意子模 $N$ 都是有限的。
对 $M$ 的生成元个数作归纳。若该数为 $1$，则对某些理想
$I \subset J \subset R$，有 $N = J/I \subset M = R/I$。
所以诺特性的定义给出结论。若 $M$ 的生成元多于 $1$ 个，
则可找到短正合列 $0 \to M' \to M \to M'' \to 0$，
其中 $M'$ 和 $M''$ 的生成元更少。注意，置
$N' = M' \cap N$ 及 $N'' = \Im(N \to M'')$，可得到 $N$ 的类似短正合列。
因此，结论由归纳假设得出，因为 $N$ 的生成元个数至多为
$N'$ 的生成元个数与 $N''$ 的生成元个数之和。

\medskip\noindent
为证明 $M$ 是有限表示的，只须把前面的结论应用于一个表示
$R^n \to M$ 的核。

\medskip\noindent
众所周知且容易证明，$M$ 的 $R$-子模满足升链条件，
等价于 $M$ 的每个子模都是有限 $R$-模。省略证明。
\end{proof}

\begin{lemma}[阿廷--里斯]
\label{algebra-lemma-Artin-Rees}
假设 $R$ 是诺特环，$I \subset R$ 是理想。
设 $N \subset M$ 为有限 $R$-模。则存在常数 $c > 0$，使得对所有
$n \geq c$，有 $I^n M \cap N  =  I^{n-c}(I^cM \cap N)$。
\end{lemma}

\begin{proof}
考虑环 $S = R \oplus I \oplus I^2 \oplus \ldots
= \bigoplus_{n \geq 0} I^n$。约定 $I^0 = R$。
乘法映射 $I^n \times I^m$ 通过 $R$ 中的乘法映到 $I^{n + m}$。
注意，若 $I = (f_1, \ldots, f_t)$，则 $S$ 是诺特环
$R[X_1, \ldots, X_t]$ 的商。该映射把单项式
$X_1^{e_1}\ldots X_t^{e_t}$ 映到 $f_1^{e_1}\ldots f_t^{e_t}$。
所以 $S$ 是诺特环。类似地，考虑模
$M \oplus IM \oplus I^2M \oplus \ldots
= \bigoplus_{n \geq 0} I^nM$。这是有限生成 $S$-模。
事实上，若 $x_1, \ldots, x_r$ 在 $R$ 上生成 $M$，
则它们也在 $S$ 上生成 $\bigoplus_{n \geq 0} I^nM$。
接着考虑子模 $\bigoplus_{n \geq 0} I^nM \cap N$。
容易验证，这是 $S$-子模。由引理 \ref{algebra-lemma-Noetherian-basic}，
它是有限生成 $S$-模；设其生成元为
$\xi_j \in \bigoplus_{n \geq 0} I^nM \cap N$，$j = 1, \ldots, s$。
将每个 $\xi_j$ 分解为齐次部分后，可以假设对某个 $d_j$，
每个 $\xi_j \in I^{d_j}M \cap N$。置 $c = \max\{d_j\}$。
则对所有 $n \geq c$，$I^nM \cap N$ 中的每个元素都形如
$\sum h_j \xi_j$，其中 $h_j \in I^{n - d_j}$。
本引理于是由此以及平凡观察
$I^{n-d_j}(I^{d_j}M \cap N) \subset I^{n-c}(I^cM \cap N)$ 得出。
\end{proof}

\begin{lemma}
\label{algebra-lemma-map-AR}
假设 $0 \to K \to M \xrightarrow{f} N$ 是诺特环 $R$ 上
有限生成模的正合列。设 $I \subset R$ 为理想。则存在 $c$，使得对所有
$n \geq c$，有
$$
f^{-1}(I^nN) = K + I^{n-c}f^{-1}(I^cN)
\quad\text{且}\quad
f(M) \cap I^nN \subset f(I^{n - c}M)
$$
\end{lemma}

\begin{proof}
将引理 \ref{algebra-lemma-Artin-Rees} 应用于 $\Im(f) \subset N$，并注意映射
$f : I^{n-c}M \to I^{n-c}f(M)$ 是满射。
\end{proof}

\begin{lemma}[克鲁尔交定理]
\label{algebra-lemma-intersect-powers-ideal-module-zero}
设 $R$ 为诺特局部环，$I \subset R$ 为真理想，$M$ 为有限 $R$-模。
则 $\bigcap_{n \geq 0} I^nM = 0$。
\end{lemma}

\begin{proof}
置 $N = \bigcap_{n \geq 0} I^nM$。则对所有 $n \geq 0$，
$N = I^nM \cap N$。由阿廷--里斯引理 \ref{algebra-lemma-Artin-Rees}，
对某个充分大的 $n$，有 $N = I^nM \cap N \subset IN$。
由中山引理 \ref{algebra-lemma-NAK}，$N = 0$。
\end{proof}

\begin{lemma}
\label{algebra-lemma-intersection-powers-ideal-module}
设 $R$ 为诺特环，$I \subset R$ 为理想，$M$ 为有限 $R$-模。
令 $N = \bigcap_n I^n M$。
\begin{enumerate}
\item 对每个满足 $I \subset \mathfrak p$ 的素理想 $\mathfrak p$，
存在 $f \in R$，$f \not \in \mathfrak p$，使得 $N_f = 0$。
\item 若 $I$ 包含于 $R$ 的雅可布森根，则 $N = 0$。
\end{enumerate}
\end{lemma}

\begin{proof}
证明 (1)。设 $x_1, \ldots, x_n$ 为模 $N$ 的生成元；参见引理
\ref{algebra-lemma-Noetherian-basic}。对每个满足 $I \subset \mathfrak p$
的素理想 $\mathfrak p$，由引理
\ref{algebra-lemma-intersect-powers-ideal-module-zero}，$N$ 在局部化
$M_{\mathfrak p}$ 中的像为零。因此，可找到
$g_i \in R$，$g_i \not \in \mathfrak p$，使得 $x_i$ 在 $N_{g_i}$ 中映到零。
于是 $N_{g_1g_2\ldots g_n} = 0$。

\medskip\noindent
(2) 由 (1) 和引理 \ref{algebra-lemma-characterize-zero-local} 得出。
\end{proof}

\begin{remark}
\label{algebra-remark-intersection-powers-ideal}
引理 \ref{algebra-lemma-intersect-powers-ideal-module-zero} 特别表明：
若 $I \subset R$ 是诺特局部环 $R$ 中的非单位理想，则
$\bigcap_n I^n = (0)$。更一般地，设 $R$ 为诺特环，$I \subset R$ 为理想。
假设 $f \in \bigcap_{n \in \mathbf{N}} I^n$。
则引理 \ref{algebra-lemma-intersection-powers-ideal-module} 表明，
对每个素理想 $I \subset \mathfrak p$，存在
$g \in R$，$g \not \in \mathfrak p$，使得 $f$ 在 $R_g$ 中映到零。
在代数几何中，我们把它表述为：“$f$ 在 $\Spec(R)$ 的闭集
$V(I)$ 的某个开邻域中为零”。
\end{remark}

\begin{lemma}[阿廷--泰特]
\label{algebra-lemma-Artin-Tate}
设 $R$ 为诺特环，$S$ 为有限生成 $R$-代数。
若 $T \subset S$ 是 $R$-子代数，且 $S$ 作为 $T$-模有限生成，
则 $T$ 在 $R$ 上是有限型的。
\end{lemma}

\begin{proof}
选取作为 $R$-代数生成 $S$ 的元素 $x_1, \ldots, x_n \in S$。
在 $S$ 中选取作为 $T$-模生成 $S$ 的 $y_1, \ldots, y_m$。
于是存在 $a_{ij} \in T$，使得 $x_i = \sum a_{ij} y_j$。
还存在 $b_{ijk} \in T$，使得 $y_i y_j = \sum b_{ijk} y_k$。
令 $T' \subset T$ 为由 $a_{ij}$ 和 $b_{ijk}$ 生成的 $R$-子代数。
它是有限生成 $R$-代数，因而是诺特环。考虑代数
$$
S' = T'[Y_1, \ldots, Y_m]/(Y_i Y_j - \sum b_{ijk} Y_k).
$$
注意，$S'$ 在 $T'$ 上有限；事实上，作为 $T'$-模，它由
$1, Y_1, \ldots, Y_m$ 的类生成。考虑把 $Y_i$ 映到 $y_i$ 的
$T'$-代数同态 $S' \to S$。由于 $a_{ij} \in T'$，可知 $x_j$
属于此映射的像。因此 $S' \to S$ 是满射。
故 $S$ 在 $T'$ 上也是有限的。由于 $T'$ 是诺特环，
可知 $T \subset S$ 在 $T'$ 上有限，证明完成。
\end{proof}






















\section{长度}
\label{algebra-section-length}

\begin{definition}
\label{algebra-definition-length}
设 $R$ 为环。对任意 $R$-模 $M$，定义 $M$ 在 $R$ 上的{it 长度}为
$$
\text{length}_R(M)
=
\sup
\{
n
\mid
\exists\ 0 = M_0 \subset M_1 \subset \ldots \subset M_n = M,
\text{ }M_i \not = M_{i + 1}
\}.
$$
\end{definition}

\noindent
换言之，它是子模链长度的上确界。一个子模链何时是另一个子模链的加细，
有一个显然的概念。这在所有子模链之集上给出偏序；若 $M$ 非零，
最小链的形状是 $0 = M_0 \subset M_1 = M$。
注意一个显然事实：若 $M$ 的长度有限，则每条链都可加细为极大链。
但是，所有极大链长度相同并不如此显然（稍后将看到）。

\begin{lemma}
\label{algebra-lemma-finite-length-finite}
\begin{slogan}
有限长度模是有限模。
\end{slogan}
设 $R$ 为环，$M$ 为 $R$-模。
若 $\text{length}_R(M) < \infty$，则 $M$ 是有限 $R$-模。
\end{lemma}

\begin{proof}
省略。
\end{proof}

\begin{lemma}
\label{algebra-lemma-length-additive}
\begin{slogan}
长度在短正合列中可加。
\end{slogan}
若 $0 \to M' \to M \to M'' \to 0$ 是 $R$ 上模的短正合列，
则 $M$ 的长度等于 $M'$ 与 $M''$ 的长度之和。
\end{lemma}

\begin{proof}
给定 $M'$ 和 $M''$ 的长度分别为 $n', n''$ 的滤过，
容易构造 $M$ 的相应滤过，其长度为 $n' + n''$。于是
$\text{length}_R M \geq \text{length}_R M' + \text{length}_R M''$。
反之，给定 $M$ 的滤过
$M_0 \subset M_1 \subset \ldots \subset M_n$，考虑诱导滤过
$M_i' = M_i \cap M'$ 和 $M_i'' = \Im(M_i \to M'')$。
令 $n'$（相应地，$n''$）为滤过 $\{M'_i\}$（相应地，$\{M''_i\}$）
的步数。若 $M_i' = M_{i + 1}'$ 且 $M_i'' = M_{i + 1}''$，
则 $M_i = M_{i + 1}$。所以 $n' + n'' \geq n$。
结合先前结论即得证明。
\end{proof}

\begin{lemma}
\label{algebra-lemma-length-infinite}
设 $R$ 为极大理想是 $\mathfrak m$ 的局部环。
若 $M$ 是 $R$-模，且对所有 $n \geq 0$，
$\mathfrak m^n M \not = 0$，则 $\text{length}_R(M) = \infty$。
换言之，若 $M$ 长度有限，则对某个 $n$，$\mathfrak m^nM = 0$。
\end{lemma}

\begin{proof}
假设对所有 $n\geq 0$，$\mathfrak m^n M \not = 0$。
选取 $x \in M$ 及 $f_1, \ldots, f_n \in \mathfrak m$，使得
$f_1f_2 \ldots f_n x \not = 0$。滤过
$$
0 \subset R f_1 \ldots f_n x \subset R f_1 \ldots f_{n - 1} x
\subset \ldots \subset R x \subset M
$$
的前 $n$ 步彼此不同。例如，若
$R f_1 x = R f_1 f_2 x$，则对某个 $g$，
$f_1 x = g f_1 f_2 x$，从而 $(1 - gf_2) f_1 x = 0$；
又因为 $1 - gf_2$ 是单位，所以 $f_1 x = 0$，
这与 $x$ 和 $f_1, \ldots, f_n$ 的选取矛盾。因此长度为无穷。
\end{proof}

\begin{lemma}
\label{algebra-lemma-length-independent}
设 $R \to S$ 为环映射，$M$ 为 $S$-模。总有
$\text{length}_R(M) \geq \text{length}_S(M)$。
若 $R \to S$ 是满射，则等号成立。
\end{lemma}

\begin{proof}
$M$ 的 $S$-子模滤过给出 $M$ 的 $R$-子模滤过。这证明不等式。
若 $R \to S$ 是满射，则 $M$ 的任意 $R$-子模自动也是 $S$-子模。
所以在此情形下等号成立。
\end{proof}

\begin{lemma}
\label{algebra-lemma-dimension-is-length}
设 $R$ 为极大理想是 $\mathfrak m$ 的环。假设 $M$ 是满足
$\mathfrak m M  =  0$ 的 $R$-模。则 $M$ 作为 $R$-模的长度，
等于 $M$ 作为 $R/\mathfrak m$-向量空间的维数。
该长度有限，当且仅当 $M$ 是有限 $R$-模。
\end{lemma}

\begin{proof}
第一部分是引理 \ref{algebra-lemma-length-independent} 的特例。
所以长度有限，当且仅当 $M$ 作为 $R/\mathfrak m$-向量空间有有限基，
当且仅当 $M$ 作为 $R$-模有有限生成集。
\end{proof}

\begin{lemma}
\label{algebra-lemma-length-localize}
设 $R$ 为环，$M$ 为 $R$-模，$S \subset R$ 为乘法子集。则
$\text{length}_R(M) \geq \text{length}_{S^{-1}R}(S^{-1}M)$。
\end{lemma}

\begin{proof}
由引理 \ref{algebra-lemma-submodule-localization}，任意子模
$N' \subset S^{-1}M$ 都形如 $S^{-1}N$，其中 $N \subset M$ 为 $R$-子模。
本引理得证。
\end{proof}

\begin{lemma}
\label{algebra-lemma-length-finite}
设 $R$ 为极大理想 $\mathfrak m$ 有限生成的环（例如 $R$ 为诺特环）。
假设 $M$ 是有限 $R$-模，并且对某个 $n$，
$\mathfrak m^n M  =  0$。则 $\text{length}_R(M) < \infty$。
\end{lemma}

\begin{proof}
考虑滤过
$0 = \mathfrak m^n M \subset
\mathfrak m^{n-1} M \subset
\ldots \subset \mathfrak m M \subset M$。
所有子商都是可应用引理 \ref{algebra-lemma-dimension-is-length} 的有限生成 $R$-模。
由可加性得出结论；参见引理 \ref{algebra-lemma-length-additive}。
\end{proof}

\begin{definition}
\label{algebra-definition-simple-module}
设 $R$ 为环，$M$ 为 $R$-模。若 $M \not = 0$，
且 $M$ 的每个子模等于 $M$ 或 $0$，则称 $M$ 是{it 单模}。
\end{definition}

\begin{lemma}
\label{algebra-lemma-characterize-length-1}
设 $R$ 为环，$M$ 为 $R$-模。下列条件等价：
\begin{enumerate}
\item $M$ 是单模；
\item $\text{length}_R(M) = 1$；
\item 对某个极大理想 $\mathfrak m \subset R$，有
$M \cong R/\mathfrak m$。
\end{enumerate}
\end{lemma}

\begin{proof}
设 $\mathfrak m$ 为 $R$ 的极大理想。由引理
\ref{algebra-lemma-dimension-is-length}，模 $R/\mathfrak m$ 的长度为 $1$。
前两个断言的等价性是重言式。假设 $M$ 是单模。
取 $x \in M$，$x \not = 0$。由于 $M$ 是单模，有 $M = R \cdot x$。
令 $I \subset R$ 为 $x$ 的零化理想，即
$I = \{f \in R \mid fx = 0\}$。映射 $R/I \to M$，
$f \bmod I \mapsto fx$ 是同构，因而 $R/I$ 是单 $R$-模。
由于 $R/I \not = 0$，可知 $I \not = R$。
令 $I \subset \mathfrak m$ 为包含 $I$ 的极大理想。
若 $I \not = \mathfrak m$，则 $\mathfrak m /I \subset R/I$
是非平凡子模，与 $R/I$ 的单性矛盾。因此如所需，$I = \mathfrak m$。
\end{proof}

\begin{lemma}
\label{algebra-lemma-simple-pieces}
设 $R$ 为环，$M$ 为有限长度 $R$-模。任取子模的极大链
$$
0 = M_0 \subset M_1 \subset M_2 \subset \ldots \subset M_n = M
$$
其中 $M_i \not = M_{i-1}$，$i = 1, \ldots, n$。则
\begin{enumerate}
\item $n = \text{length}_R(M)$；
\item 每个 $M_i/M_{i-1}$ 都是单模；
\item 每个 $M_i/M_{i-1}$ 都形如 $R/\mathfrak m_i$，
其中 $\mathfrak m_i$ 为某个极大理想；
\item 给定极大理想 $\mathfrak m \subset R$，有
$$
\# \{i \mid \mathfrak m_i = \mathfrak m\}
=
\text{length}_{R_{\mathfrak m}} (M_{\mathfrak m}).
$$
\end{enumerate}
\end{lemma}

\begin{proof}
若 $M_i/M_{i-1}$ 不是单模，则可以加细该滤过，因而该滤过不是极大的。
所以 $M_i/M_{i-1}$ 是单模。由引理
\ref{algebra-lemma-characterize-length-1}，模 $M_i/M_{i-1}$ 长度为 $1$，
并且形如 $R/\mathfrak m_i$，其中 $\mathfrak m_i$ 为某个极大理想。
由长度的可加性，即引理 \ref{algebra-lemma-length-additive}，可知
$n = \text{length}_R(M)$。由于局部化正合，
$$
0 = (M_0)_{\mathfrak m}
\subset (M_1)_{\mathfrak m}
\subset (M_2)_{\mathfrak m}
\subset \ldots
\subset (M_n)_{\mathfrak m} = M_{\mathfrak m}
$$
是 $M_{\mathfrak m}$ 的滤过，其相继商为
$(M_i/M_{i-1})_{\mathfrak m}$。所以，最后一个陈述直接由下述事实得出：
给定 $R$ 的极大理想 $\mathfrak m$、$\mathfrak m'$，有
$$
(R/\mathfrak m')_{\mathfrak m}
\cong
\left\{
\begin{matrix}
0 &
\text{若 } \mathfrak m \not = \mathfrak m', \\
R_{\mathfrak m}/\mathfrak m R_{\mathfrak m} &
\text{若 } \mathfrak m  = \mathfrak m'
\end{matrix}
\right.
$$
把这一点留给读者。
\end{proof}

\begin{lemma}
\label{algebra-lemma-pushdown-module}
设 $A$ 为极大理想是 $\mathfrak m$ 的局部环。
设 $B$ 为极大理想是 $\mathfrak m_i$，$i = 1, \ldots, n$ 的半局部环。
假设 $A \to B$ 是同态，每个 $\mathfrak m_i$ 都位于 $\mathfrak m$ 上，且
$$
[\kappa(\mathfrak m_i) : \kappa(\mathfrak m)] < \infty.
$$
设 $M$ 为有限长度 $B$-模。则
$$
\text{length}_A(M) = \sum\nolimits_{i = 1, \ldots, n}
[\kappa(\mathfrak m_i) : \kappa(\mathfrak m)]
\text{length}_{B_{\mathfrak m_i}}(M_{\mathfrak m_i}),
$$
特别地，$\text{length}_A(M) < \infty$。
\end{lemma}

\begin{proof}
选取引理 \ref{algebra-lemma-simple-pieces} 中那样的 $B$-子模极大链
$$
0 = M_0
\subset M_1
\subset M_2
\subset \ldots
\subset M_m = M
$$
则对某个 $i(j) \in \{1, \ldots, n\}$，每个商
$M_j/M_{j - 1}$ 都同构于 $\kappa(\mathfrak m_{i(j)})$。
此外，由引理 \ref{algebra-lemma-dimension-is-length}，还有：
$$
\text{length}_A(\kappa(\mathfrak m_i)) =
[\kappa(\mathfrak m_i) : \kappa(\mathfrak m)]
$$
本引理由长度的可加性得出（引理 \ref{algebra-lemma-length-additive}）。
\end{proof}

\begin{lemma}
\label{algebra-lemma-pullback-module}
设 $A \to B$ 为局部环的平坦局部同态。则对任意 $A$-模 $M$，有
$$
\text{length}_A(M) \text{length}_B(B/\mathfrak m_AB)
=
\text{length}_B(M \otimes_A B).
$$
特别地，若 $\text{length}_B(B/\mathfrak m_AB) < \infty$，
则 $M$ 长度有限，当且仅当 $M \otimes_A B$ 长度有限。
\end{lemma}

\begin{proof}
由引理 \ref{algebra-lemma-local-flat-ff}，环映射 $A \to B$ 忠实平坦。
所以，若 $0 = M_0 \subset M_1 \subset \ldots \subset M_n = M$
是 $M$ 中长度为 $n$ 的链，则相应的链
$0 = M_0 \otimes_A B \subset M_1 \otimes_A B \subset
\ldots \subset M_n \otimes_A B = M \otimes_A B$ 的长度也为 $n$。
这证明
$\text{length}_A(M) = \infty \Rightarrow
\text{length}_B(M \otimes_A B) = \infty$。
接着，假设 $\text{length}_A(M) < \infty$。
此时，$M$ 有一个长度为 $\ell = \text{length}_A(M)$ 的滤过，
其各商都是 $A/\mathfrak m_A$。如上论证，$M \otimes_A B$
有一个长度为 $\ell$ 的滤过，其各商同构于
$B \otimes_A A/\mathfrak m_A =  B/\mathfrak m_AB$。本引理得证。
\end{proof}

\begin{lemma}
\label{algebra-lemma-pullback-transitive}
设 $A \to B \to C$ 为局部环的平坦局部同态。则
$$
\text{length}_B(B/\mathfrak m_A B)
\text{length}_C(C/\mathfrak m_B C)
=
\text{length}_C(C/\mathfrak m_A C)
$$
\end{lemma}

\begin{proof}
将引理 \ref{algebra-lemma-pullback-module} 应用于环映射
$B \to C$ 和 $B$-模 $M = B/\mathfrak m_A B$ 即得。
\end{proof}












\section{阿廷环}
\label{algebra-section-artinian}

\noindent
阿廷环，特别是局部阿廷环，在代数几何中起重要作用，例如在形变理论中。

\begin{definition}
\label{algebra-definition-artinian}
若环 $R$ 的理想满足降链条件，则称它是{\it 阿廷的}。
\end{definition}

\begin{lemma}
\label{algebra-lemma-finite-dimensional-algebra}
假设 $R$ 是域上的有限维代数。则 $R$ 是阿廷环。
\end{lemma}

\begin{proof}
理想的降链条件显然成立。
\end{proof}

\begin{lemma}
\label{algebra-lemma-artinian-finite-nr-max}
若 $R$ 是阿廷环，则 $R$ 只有有限多个极大理想。
\end{lemma}

\begin{proof}
假设 $\mathfrak m_i$，$i = 1, 2, 3, \ldots$ 是两两不同的极大理想。
则
$\mathfrak m_1 \supset \mathfrak m_1\cap \mathfrak m_2
\supset \mathfrak m_1 \cap \mathfrak m_2 \cap \mathfrak m_3 \supset \ldots$
是无穷降链（因为由中国剩余定理，所有映射
$R \to \oplus_{i = 1}^n R/\mathfrak m_i$ 都是满射）。
\end{proof}

\begin{lemma}
\label{algebra-lemma-artinian-radical-nilpotent}
设 $R$ 为阿廷环。则 $R$ 的雅可布森根是幂零理想。
\end{lemma}

\begin{proof}
设 $I \subset R$ 为雅可布森根。注意，
$I \supset I^2 \supset I^3 \supset \ldots$ 是降链。
所以对某个 $n$，$I^n = I^{n + 1}$。置
$J = \{ x\in R \mid xI^n = 0\}$。须证明 $J = R$。
若非如此，选取满足 $J \subset J'$ 的理想 $J' \not = J$，
并使它极小（由阿廷性可行）。则 $J'/J$ 是单 $R$-模，
因而同构于某个极大理想 $\mathfrak m$ 所给出的 $R/\mathfrak m$；
参见引理 \ref{algebra-lemma-characterize-length-1}。
于是 $\mathfrak m I^n$ 零化 $J'$。由于 $I \subset \mathfrak m$，
可知 $I^{n + 1} = I^n$ 零化 $J'$。所以 $J' = J$，矛盾。
\end{proof}

\begin{lemma}
\label{algebra-lemma-product-local}
任意只有有限多个极大理想、且雅可布森根局部幂零的环，
都等于它在各极大理想处的局部化之积。此外，所有素理想都是极大理想。
\end{lemma}

\begin{proof}
设 $R$ 为极大理想只有 $\mathfrak m_1, \ldots, \mathfrak m_n$ 的环。
令 $I = \bigcap_{i = 1}^n \mathfrak m_i$ 为 $R$ 的雅可布森根。
假设 $I$ 局部幂零。设 $\mathfrak p$ 为 $R$ 的素理想。
由于每个素理想都包含 $R$ 的每个幂零元，可知
$ \mathfrak p \supset \mathfrak m_1 \cap \ldots \cap \mathfrak m_n$。
又由于 $\mathfrak m_1 \cap \ldots \cap \mathfrak m_n \supset
\mathfrak m_1 \ldots \mathfrak m_n$，可知
$\mathfrak p \supset \mathfrak m_1 \ldots \mathfrak m_n$。
故对某个 $i$，$\mathfrak p \supset \mathfrak m_i$，从而
$\mathfrak p = \mathfrak m_i$。因此 $R$ 的谱是离散拓扑空间
$\{\mathfrak m_1, \ldots, \mathfrak m_n\}$。
将引理 \ref{algebra-lemma-disjoint-implies-product} 应用 $n - 1$ 次，得到
$R = R_1 \times \ldots \times R_n$，其中对每个 $i$，$R_i$ 的谱都是单点集。
所以 $R_i$ 是局部环；又因为由该引理，它是 $R$ 的局部化，
故它如所需地是 $R$ 的这些局部环之一。
\end{proof}

\begin{lemma}
\label{algebra-lemma-artinian-finite-length}
环 $R$ 是阿廷环，当且仅当它作为自身上的模具有有限长度。
任意这样的环 $R$ 既是阿廷环又是诺特环；$R$ 的任意素理想都是极大理想；
并且 $R$ 等于它在各极大理想处的局部化的（有限）乘积。
\end{lemma}

\begin{proof}
若 $R$ 在自身上长度有限，则它的理想既满足升链条件又满足降链条件。
所以它既是诺特环又是阿廷环。由引理
\ref{algebra-lemma-artinian-finite-nr-max}、
\ref{algebra-lemma-artinian-radical-nilpotent} 和 \ref{algebra-lemma-product-local}，
任意阿廷环都等于它在极大理想处的局部化之积。

\medskip\noindent
假设 $R$ 是阿廷环。将证明 $R$ 在自身上长度有限。
只须给出一条子模链，使其相继商长度有限。
由上述结论，可以假设 $R$ 是极大理想为 $\mathfrak m$ 的局部环。
由引理 \ref{algebra-lemma-artinian-radical-nilpotent}，
对某个 $n$，有 $\mathfrak m^n =0$。考虑序列
$0 = \mathfrak m^n \subset \mathfrak m^{n-1} \subset
\ldots \subset \mathfrak m \subset R$。由引理
\ref{algebra-lemma-dimension-is-length}，每个子商
$\mathfrak m^j/\mathfrak m^{j + 1}$ 的长度，等于它作为
$\kappa(\mathfrak m)$-向量空间的维数。该维数必须有限；
否则会有向量子空间的无穷降链，它对应于 $R$ 中理想的无穷降链。
\end{proof}

\section{本质有限型同态}
\label{algebra-section-essentially-of-finite-type}

\noindent
下面给出关于有限型环同态之局部化的一些简单说明。

\begin{definition}
\label{algebra-definition-essentially-finite-p-t}
设 $R \to S$ 为环同态。
\begin{enumerate}
\item 若 $S$ 是某个有限型 $R$-代数的局部化，则称 $R \to S$
是{\it 本质有限型的}。
\item 若 $S$ 是某个有限表示 $R$-代数的局部化，则称 $R \to S$
是{\it 本质有限表示的}。
\end{enumerate}
\end{definition}

\begin{lemma}
\label{algebra-lemma-composition-essentially-of-finite-type}
本质有限型环同态的类在复合下保持不变。本质有限表示的情形亦然。
\end{lemma}

\begin{proof}
略。
\end{proof}

\begin{lemma}
\label{algebra-lemma-base-change-essentially-of-finite-type}
本质有限型环同态的类在基变换下保持不变。本质有限表示的情形亦然。
\end{lemma}

\begin{proof}
略。
\end{proof}

\begin{lemma}
\label{algebra-lemma-essentially-of-finite-type-into-artinian-local}
设 $R \to S$ 为环同态。假设 $S$ 是极大理想为 $\mathfrak m$
的阿廷局部环。则
\begin{enumerate}
\item $R \to S$ 有限，当且仅当 $R \to S/\mathfrak m$ 有限；
\item $R \to S$ 是有限型的，当且仅当 $R \to S/\mathfrak m$
是有限型的；
\item $R \to S$ 本质有限型，当且仅当复合
$R \to S/\mathfrak m$ 本质有限型。
\end{enumerate}
\end{lemma}

\begin{proof}
若 $R \to S$ 有限，则由引理 \ref{algebra-lemma-finite-transitive}，
$R \to S/\mathfrak m$ 有限。反过来，假设 $R \to S/\mathfrak m$ 有限。
由于 $S$ 在自身上具有有限长度（引理
\ref{algebra-lemma-artinian-finite-length}），可以选取由理想组成的滤过
$$
0 \subset I_1 \subset \ldots \subset I_n = S
$$
使得作为 $S$-模，有 $I_i/I_{i - 1} \cong S/\mathfrak m$。
因此，$S$ 有一条由 $R$-子模 $I_i$ 组成的滤过，其中每个相继商
都是有限 $R$-模。故由引理 \ref{algebra-lemma-extension}，$S$ 是有限
$R$-模。

\medskip\noindent
若 $R \to S$ 是有限型的，则由引理
\ref{algebra-lemma-compose-finite-type}，$R \to S/\mathfrak m$ 是有限型的。
反过来，假设 $R \to S/\mathfrak m$ 是有限型的。选取
$f_1, \ldots, f_n \in S$，使它们的像生成 $S/\mathfrak m$。
于是 $A = R[x_1, \ldots, x_n] \to S$、$x_i \mapsto f_i$
是环同态，且 $A \to S/\mathfrak m$ 满射（特别地，它是有限的）。
故由第 (1) 部分，$A \to S$ 有限；再由引理
\ref{algebra-lemma-compose-finite-type}，可知 $R \to S$ 是有限型的。

\medskip\noindent
若 $R \to S$ 本质有限型，则由引理
\ref{algebra-lemma-composition-essentially-of-finite-type}，
$R \to S/\mathfrak m$ 本质有限型。反过来，假设
$R \to S/\mathfrak m$ 本质有限型。设 $S/\mathfrak m$
是 $R[x_1, \ldots, x_n]/I$ 的局部化。选取
$f_1, \ldots, f_n \in S$，使它们模 $\mathfrak m$ 的剩余类
分别对应于 $x_1, \ldots, x_n$ 模 $I$ 的剩余类。考虑同态
$R[x_1, \ldots, x_n] \to S$、$x_i \mapsto f_i$，并记其核为 $J$。
置 $A = R[x_1, \ldots, x_n]/J \subset S$，并置
$\mathfrak p = A \cap \mathfrak m$。注意，
$A/\mathfrak p \subset S/\mathfrak m$ 等于
$R[x_1, \ldots, x_n]/I$ 在 $S/\mathfrak m$ 中的像。
所以 $\kappa(\mathfrak p) = S/\mathfrak m$。由第 (1) 部分，
$A_\mathfrak p \to S$ 有限。最后由引理
\ref{algebra-lemma-composition-essentially-of-finite-type}，
可知 $S$ 本质有限型。
\end{proof}

\noindent
下面的引理也可利用射影空间的固有性来证明，而不必使用这里给出的代数论证。

\begin{lemma}
\label{algebra-lemma-localization-at-closed-point-special-fibre}
设 $\varphi : R \to S$ 本质有限型，且 $R$ 与 $S$ 都是局部环
（但 $\varphi$ 不必是局部同态）。则存在某个 $n$ 和一个位于
$\mathfrak m_R$ 之上的极大理想
$\mathfrak m \subset R[x_1, \ldots, x_n]$，使得 $S$ 是
$R[x_1, \ldots, x_n]_\mathfrak m$ 的某个商环的局部化。
\end{lemma}

\begin{proof}
可以将 $S$ 写成 $R[x_1, \ldots, x_n]$ 的某个商环的局部化。
因此只须证明如下情形：对某个素理想
$\mathfrak q \subset R[x_1, \ldots, x_n]$，有
$S = R[x_1, \ldots, x_n]_\mathfrak q$。若
$\mathfrak q + \mathfrak m_R R[x_1, \ldots, x_n] \not = R[x_1, \ldots, x_n]$，
则可找到断言中的极大理想 $\mathfrak m$，使得
$\mathfrak q \subset \mathfrak m$；此时结论显然。

\medskip\noindent
选取支配 $R \to \kappa(\mathfrak q)$ 之像的赋值环
$A \subset \kappa(\mathfrak q)$（引理 \ref{algebra-lemma-dominate}）。
若 $x_i$ 的像 $\lambda_i \in \kappa(\mathfrak q)$ 属于 $A$，
则 $\mathfrak q$ 包含于 $\mathfrak m_A$ 在同态
$R[x_1, \ldots, x_n] \to A$ 下的逆像中，这就回到前一种情形。
因此存在某个 $i$，使得 $\lambda_i^{-1} \in A$，并且对所有
$j = 1, \ldots, n$，有 $\lambda_j/\lambda_i \in A$
（因为 $A$ 的值群是全序的，参见引理
\ref{algebra-lemma-valuation-group}）。于是考虑同态
$$
R[y_0, y_1, \ldots, \hat{y_i}, \ldots, y_n]
\to R[x_1, \ldots, x_n]_\mathfrak q,\quad
y_0 \mapsto 1/x_i,\quad y_j \mapsto x_j/x_i
$$
设 $\mathfrak q' \subset R[y_0, \ldots, \hat{y_i}, \ldots, y_n]$
为 $\mathfrak q$ 的逆像。由于 $y_0 \not \in \mathfrak q'$，
容易看出上述箭头在局部化后给出一个同构。另一方面，
第一段的结论适用于 $R[y_0, \ldots, \hat{y_i}, \ldots, y_n]$，
因为 $y_j$ 映到 $A$ 的一个元素。证明完毕。
\end{proof}






















\section{K-群}
\label{algebra-section-K-groups}

\noindent
设 $R$ 为环。下面将引入两个与 $R$ 相伴的阿贝尔群。第一个记作
$K'_0(R)$，它具有如下性质\footnote{该定义对任意环都有意义，
但除非 $R$ 是诺特环，否则很少使用。}：
\begin{enumerate}
\item 对每个有限 $R$-模 $M$，给定 $K'_0(R)$ 中的一个元素 $[M]$；
\item 对有限 $R$-模的每个短正合列
$0 \to M' \to M \to M'' \to 0$，都有关系
$[M] = [M'] + [M'']$；
\item 群 $K'_0(R)$ 由诸元素 $[M]$ 生成；
\item $K'_0(R)$ 中生成元 $[M]$ 之间的所有关系，都是由上述正合列
所产生关系的 $\mathbf{Z}$-线性组合。
\end{enumerate}
实际构造略为麻烦，因为必须顾及所有有限生成 $R$-模的汇集是一个真类。
不过，可以用下述方式克服这个问题：取诸元素 $[R^n/K]$ 作为群
$K'_0(R)$ 的生成元集合，其中 $n$ 遍历所有整数，而 $K$ 遍历所有
子模 $K \subset R^n$。施加于这些元素上的关系子群，可取为项均具有
$R^n/K$ 形式的短正合列所给出的关系。元素 $[M]$ 定义如下：
选取 $n$ 与 $K$，使得 $M \cong R^n/K$，并置
$[M] = [R^n/K]$。细节留给读者。

\begin{lemma}
\label{algebra-lemma-length-K0}
若 $R$ 是阿廷局部环，则长度函数定义一个自然的阿贝尔群同态
$\text{length}_R : K'_0(R) \to \mathbf{Z}$。
\end{lemma}

\begin{proof}
任意有限 $R$-模的长度都有限，因为它是 $R^n$ 的商，而由引理
\ref{algebra-lemma-artinian-finite-length}，后者长度有限。
此外，长度函数具有可加性；参见引理 \ref{algebra-lemma-length-additive}。
\end{proof}

\noindent
上述两个群中的第二个记作 $K_0(R)$，它具有如下性质：
\begin{enumerate}
\item 对每个有限投射 $R$-模 $M$，给定 $K_0(R)$ 中的一个元素 $[M]$；
\item 对有限投射 $R$-模的每个短正合列
$0 \to M' \to M \to M'' \to 0$，都有关系
$[M] = [M'] + [M'']$；
\item 群 $K_0(R)$ 由诸元素 $[M]$ 生成；
\item $K_0(R)$ 中的所有关系，都是由上述正合列所产生关系的
$\mathbf{Z}$-线性组合。
\end{enumerate}
这个群的构造与上面相同。

\medskip\noindent
注意，存在显然的映射 $K_0(R) \to K'_0(R)$，它一般不是同构。

\begin{example}
\label{algebra-example-K0-field}
注意，若 $R = k$ 是域，则显然有
$K_0(k) = K'_0(k) \cong \mathbf{Z}$；该同构由维数函数给出
（它也就是长度函数）。
\end{example}

\begin{example}
\label{algebra-example-K0-PID}
设 $R$ 为主理想整环。断言 $K_0(R) = K'_0(R) = \mathbf{Z}$。
事实上，任意有限投射 $R$-模都是有限自由模。由引理
\ref{algebra-lemma-rank}，有限自由模具有良定义的秩。给定有限自由模的短正合列
$$
0 \to M' \to M \to M'' \to 0
$$
有 $\text{rank}(M) = \text{rank}(M') + \text{rank}(M'')$，
因为在此情形下 $M \cong M' \oplus M'$（例如，由引理
\ref{algebra-lemma-lift-map}，该列分裂）。由此得到
$K_0(R) = \mathbf{Z}$。

\medskip\noindent
主理想整环上模的结构定理表明，任意有限生成 $R$-模都具有形式
$M = R^{\oplus r} \oplus R/(d_1) \oplus \ldots \oplus R/(d_k)$。
考虑短正合列
$$
0 \to (d_i) \to R \to R/(d_i) \to 0
$$
由于理想 $(d_i)$ 作为模同构于 $R$（它是由 $d_i$ 生成的自由模），
在 $K'_0(R)$ 中有 $[(d_i)] = [R]$。于是
$[R/(d_i)] = [(d_i)]-[R] = 0$。由此可知，挠模在
$K'_0(R)$ 中的类为零。利用自由部分的秩，可得到等同
$K'_0(R) = \mathbf{Z}$，且典范同态
$K_0(R) \to K'_0(R)$ 是同构。
\end{example}

\begin{example}
\label{algebra-example-K0-polynomial-ring}
设 $k$ 为域。则 $K_0(k[x]) = K'_0(k[x]) = \mathbf{Z}$。
这是例 \ref{algebra-example-K0-PID} 的直接推论，因为
$R = k[x]$ 是主理想整环。
\end{example}

\begin{example}
\label{algebra-example-K0-node}
设 $k$ 为域。令 $R = \{f \in k[x] \mid f(0) = f(1)\}$；
比较例 \ref{algebra-example-affine-open-not-standard}。在此情形下，
$K_0(R) \cong k^* \oplus \mathbf{Z}$，而
$K'_0(R) = \mathbf{Z}$。
\end{example}

\begin{lemma}
\label{algebra-lemma-K0-product}
设 $R = R_1 \times R_2$。则
$K_0(R) = K_0(R_1) \times K_0(R_2)$，且
$K'_0(R) = K'_0(R_1) \times K'_0(R_2)$。
\end{lemma}

\begin{proof}
略。
\end{proof}

\begin{lemma}
\label{algebra-lemma-K0prime-Artinian}
设 $R$ 为阿廷局部环。引理 \ref{algebra-lemma-length-K0} 中的映射
$\text{length}_R : K'_0(R) \to \mathbf{Z}$ 是同构。
\end{lemma}

\begin{proof}
略。
\end{proof}

\begin{lemma}
\label{algebra-lemma-K0-local}
设 $(R, \mathfrak m)$ 为局部环。每个有限投射 $R$-模都是有限自由模。
由 $[M] \to \text{rank}_R(M)$ 定义的映射
$\text{rank}_R : K_0(R) \to \mathbf{Z}$ 是良定义的同构。
\end{lemma}

\begin{proof}
设 $P$ 为有限投射 $R$-模。选取诸元素
$x_1, \ldots, x_n \in P$，使它们的像构成 $P/\mathfrak m P$ 的一组基。
由 Nakayama 引理 \ref{algebra-lemma-NAK}，这些元素生成 $P$。
相应的满射 $u : R^{\oplus n} \to P$ 有分裂，因为 $P$ 是投射模。
因此 $R^{\oplus n} = P \oplus Q$，其中 $Q = \Ker(u)$。
由此 $Q/\mathfrak m Q = 0$，故由 Nakayama 引理，$Q$ 为零。
这样就看出，每个有限投射 $R$-模都是有限自由模。
由引理 \ref{algebra-lemma-rank}，有限自由模具有良定义的秩。
给定有限自由 $R$-模的短正合列
$$
0 \to M' \to M \to M'' \to 0
$$
有 $\text{rank}(M) = \text{rank}(M') + \text{rank}(M'')$，
因为在此情形下 $M \cong M' \oplus M'$（例如，由引理
\ref{algebra-lemma-lift-map}，该列分裂）。由此得到
$K_0(R) = \mathbf{Z}$。
\end{proof}

\begin{lemma}
\label{algebra-lemma-K0-and-K0prime-Artinian-local}
设 $R$ 为阿廷局部环。存在交换图
$$
\xymatrix{
K_0(R) \ar[rr] \ar[d]_{\text{rank}_R} & &
K'_0(R) \ar[d]^{\text{length}_R} \\
\mathbf{Z} \ar[rr]^{\text{length}_R(R)} & &
\mathbf{Z}
}
$$
其中由引理 \ref{algebra-lemma-K0prime-Artinian} 和
\ref{algebra-lemma-K0-local}，两个竖直映射都是同构。
\end{lemma}

\begin{proof}
设 $P$ 为有限投射 $R$-模。须证明
$\text{length}_R(P) = \text{rank}_R(P) \text{length}_R(R)$。
由引理 \ref{algebra-lemma-K0-local}，模 $P$ 是有限自由模。所以对某个
$n \geq 0$，有 $P \cong R^{\oplus n}$。于是
$\text{rank}_R(P) = n$，并且由长度的可加性（引理
\ref{algebra-lemma-length-additive}），
$\text{length}_R(R^{\oplus n}) = n \text{length}_R(R)$。
因此结论成立。
\end{proof}

















\section{分次环}
\label{algebra-section-graded}

\noindent
这里所谓的{\it 分次环}，是指底层阿贝尔群具有直和分解
$S = \bigoplus_{d \geq 0} S_d$ 的环 $S$，并且
$S_d \cdot S_e \subset S_{d + e}$。注意，不允许负次数部分中有非零元。
{\it 无关理想}是理想 $S_{+} = \bigoplus_{d > 0} S_d$。
所谓{\it 分次模}，是指底层阿贝尔群具有直和分解
$M = \bigoplus_{n\in \mathbf{Z}} M_n$ 的 $S$-模 $M$，并且
$S_d \cdot M_e \subset M_{d + e}$。注意，对模允许负次数部分中有非零元。
置 $S_{-k} = (0)$（$k > 0$），便可将 $S$ 看成分次 $S$-模。
$M$（相应地，$S$）的元素 $x$（相应地，$f$）称为{\it 齐次的}，
如果对某个 $d$，有 $x \in M_d$（相应地，$f \in S_d$）。
分次 $S$-模之间的{\it 同态}，是满足
$\varphi(M_d) \subset M'_d$ 的 $S$-模同态 $\varphi : M \to M'$。
这里不允许同态移动次数。用 $\text{GrHom}_0(M, N)$ 表示从 $M$ 到 $N$
的分次模同态所成的 $S_0$-模。

\medskip\noindent
此时自然有分次理想、分次商环、分次子模、分次商模、分次张量积等概念。
相关定义和引理留给读者。例如，分次模的短正合列在每个次数上都是短正合的。

\medskip\noindent
给定分次环 $S$、分次 $S$-模 $M$ 和 $n \in \mathbf{Z}$，
用 $M(n)$ 表示满足 $M(n)_d = M_{n + d}$ 的分次 $S$-模。
它称为{\it $M$ 由 $n$ 的扭转}。特别地，得到诸模
$S(n)$、$n \in \mathbf{Z}$；它们在射影概形的研究中起重要作用。
存在一些显然的函子性同构，例如
$(M \oplus N)(n) = M(n) \oplus N(n)$、
$(M \otimes_S N)(n) = M \otimes_S N(n) = M(n) \otimes_S N$。
此外，还可以在下列 $S_0$-模上定义分次 $S$-模结构：
$$
\text{GrHom}(M, N) =
\bigoplus\nolimits_{n \in \mathbf{Z}} \text{GrHom}_n(M, N),
\quad
\text{GrHom}_n(M, N) = \text{GrHom}_0(M, N(n)).
$$
乘法的定义从略。

\begin{lemma}
\label{algebra-lemma-graded-NAK}
设 $S$ 为分次环，$M$ 为分次 $S$-模。
\begin{enumerate}
\item 若 $S_+M = M$ 且 $M$ 有限，则 $M = 0$。
\item 若 $N, N' \subset M$ 是分次子模，
$M = N + S_+N'$，且 $N'$ 有限，则 $M = N$。
\item 若 $N \to M$ 是分次模同态，$N/S_+N \to M/S_+M$
满射，且 $M$ 有限，则 $N \to M$ 满射。
\item 若齐次元 $x_1, \ldots, x_n \in M$ 生成 $M/S_+M$，
且 $M$ 有限，则 $x_1, \ldots, x_n$ 生成 $M$。
\end{enumerate}
\end{lemma}

\begin{proof}
证明 (1)。选取 $M$ 在 $S$ 上的一组生成元 $y_1, \ldots, y_r$。
可以假设 $y_i$ 是次数为 $d_i$ 的齐次元。重新编号后，可以假设
$d_r = \min(d_i)$。此时条件 $S_+M = M$ 蕴含 $y_r = 0$。
故对 $r$ 归纳即得 $M = 0$。将 (1) 应用于 $M/N$，即得 (2)。
将 (2) 应用于子模 $\Im(N \to M)$ 与 $M$，即得 (3)。
将 (3) 应用于模同态
$\bigoplus S(-d_i) \to M$、$(s_1, \ldots, s_n) \mapsto \sum s_i x_i$，
即得 (4)。
\end{proof}

\noindent
设 $S$ 为分次环，$d \geq 1$ 为整数。置
$S^{(d)} = \bigoplus_{n \geq 0} S_{nd}$。把 $S^{(d)}$ 看成分次环，
其次数为 $n$ 的分量是 $(S^{(d)})_n = S_{nd}$。给定分次
$S$-模 $M$，同样可以考虑
$M^{(d)} = \bigoplus_{n \in \mathbf{Z}} M_{nd}$；
这是一个分次 $S^{(d)}$-模。

\begin{lemma}
\label{algebra-lemma-uple-generated-degree-1}
设 $S$ 为在 $S_0$ 上有限生成的分次环。则对所有充分可除的 $d$，
代数 $S^{(d)}$ 在 $S_0$ 上由次数 $1$ 的部分生成。
\end{lemma}

\begin{proof}
设 $S$ 在 $S_0$ 上由 $f_1, \ldots, f_r \in S$ 生成。
用 $f_i$ 的齐次部分替换它们后，可以假设 $f_i$ 是次数
$d_i > 0$ 的齐次元。于是 $S_n$ 的任意元素都是单项式
$f_1^{e_1} \ldots f_r^{e_r}$ 的 $S_0$-系数线性组合，
其中 $\sum e_i d_i = n$。设 $m$ 是 $\text{lcm}(d_i)$ 的倍数。
对任意 $N \geq r$，若
$$
\sum e_i d_i = N m
$$
则由一个初等论证，对某个 $i$ 有 $e_i \geq m/d_i$。
所以次数为 $N m$ 的每个单项式，都是次数为 $m$ 的单项式
$f_i^{m/d_i}$ 与次数为 $(N - 1)m$ 的单项式之积。
由此，任意次数为 $nrm$、其中 $n \geq 2$ 的单项式，都是次数为
$rm$ 的诸单项式之积。因此 $S^{(rm)}$ 在 $S_0$ 上由次数 $1$ 的部分生成。
\end{proof}

\begin{lemma}
\label{algebra-lemma-integral-closure-graded}
\begin{reference}
\cite[Theorem 2.3.2]{Huneke-Swanson}
\end{reference}
设 $R \to S$ 为分次环同态，$S' \subset S$ 为 $R$ 在 $S$ 中的整闭包。
则
$$
S' = \bigoplus\nolimits_{d \geq 0} S' \cap S_d,
$$
即 $S'$ 是 $S$ 的分次 $R$-子代数。
\end{lemma}

\begin{proof}
须证明：若 $s = s_n + s_{n + 1} + \ldots + s_m \in S'$，
则每个齐次部分 $s_j \in S'$。对所有分次环同态
$R \to S$，关于 $m - n$ 归纳证明。先注意，$s_0$ 显然在 $R_0$
上整（因而在 $R$ 上整），因为存在与环同态 $R \to R_0$ 相容的
环同态 $S \to S_0$。故以 $s - s_0$ 代替 $s$ 后，可以假设
$n > 0$。考虑分次环扩张 $R[t, t^{-1}] \to S[t, t^{-1}]$，
其中 $t$ 的次数为 $0$。存在交换图
$$
\xymatrix{
S[t, t^{-1}] \ar[rr]_{s \mapsto t^{\deg(s)}s} & & S[t, t^{-1}] \\
R[t, t^{-1}] \ar[u] \ar[rr]^{r \mapsto t^{\deg(r)}r} & &  R[t, t^{-1}] \ar[u]
}
$$
其中横向映射都是环自同构。因此，$S[t, t^{-1}]$ 在
$R[t, t^{-1}]$ 上的整闭包 $C$ 被映入自身。于是
$$
t^m(s_n + s_{n + 1} + \ldots + s_m) -
(t^ns_n + t^{n + 1}s_{n + 1} + \ldots + t^ms_m) \in C
$$
由归纳假设可知，对 $i = n, \ldots, m - 1$，每个
$(t^m - t^i)s_i \in C$。注意，对任意环 $A$ 和
$m > i \geq n > 0$，有
$A[t, t^{-1}]/(t^m - t^i - 1) \cong A[t]/(t^m - t^i - 1) \supset A$，
因为在 $A[t]/(t^m - t^i - 1)$ 中，
$t(t^{m - 1} - t^{i - 1}) = 1$。由于 $t^m - t^i$ 映为 $1$，
可知对 $i = n, \ldots, m - 1$，$s_i$ 在环
$S[t]/(t^m - t^i - 1)$ 中的像，在 $R[t]/(t^m - t^i - 1)$ 上整。
由于 $R \to R[t]/(t^m - t^i - 1)$ 有限，由传递性可知
$s_i$ 在 $R$ 上整；参见引理 \ref{algebra-lemma-integral-transitive}。
最后还可得
$s_m = s - \sum_{i = n, \ldots, m - 1} s_i$ 在 $R$ 上整。
\end{proof}







\section{分次环的 Proj}
\label{algebra-section-proj}

\noindent
设 $S$ 为分次环。所谓{\it 齐次理想}，就是既是该环的理想、
又是 $S$ 的分次子模的 $I \subset S$。等价地，它是由齐次元生成的理想。
再等价地，若 $f \in I$ 且
$$
f = f_0 + f_1 + \ldots + f_n
$$
是 $f$ 在 $S$ 中的齐次部分分解，则对每个 $i$ 都有 $f_i \in I$。
要检验齐次理想 $\mathfrak p$ 是否为素理想，只须检验：
若 $ab \in \mathfrak p$，其中 $a, b$ 齐次，则
$a \in \mathfrak p$ 或 $b \in \mathfrak p$。

\begin{definition}
\label{algebra-definition-proj}
设 $S$ 为分次环。定义 $\text{Proj}(S)$ 为 $S$ 的所有满足
$S_{+} \not \subset \mathfrak p$ 的齐次素理想 $\mathfrak p$ 所成的集合。
集合 $\text{Proj}(S)$ 是 $\Spec(S)$ 的子集，赋予它诱导拓扑。
拓扑空间 $\text{Proj}(S)$ 称为分次环 $S$ 的{\it 齐次谱}。
\end{definition}

\noindent
注意，由构造有连续映射
$$
\text{Proj}(S) \longrightarrow \Spec(S_0).
$$

\medskip\noindent
设 $S = \oplus_{d \geq 0} S_d$ 为分次环。取 $f\in S_d$，
并假设 $d \geq 1$。定义 $S_{(f)}$ 为 $S_f$ 的子环，
其元素具有形式 $r/f^n$，其中 $r$ 齐次且 $\deg(r) = nd$。
若 $M$ 是分次 $S$-模，则定义 $S_{(f)}$-模 $M_{(f)}$
为 $M_f$ 的子模，其元素具有形式 $x/f^n$，其中 $x$ 是次数为
$nd$ 的齐次元。

\begin{lemma}
\label{algebra-lemma-Z-graded}
设 $S$ 是含有一个正次数齐次可逆元的 $\mathbf{Z}$-分次环。
则 $S$ 的 $\mathbf{Z}$-分次素理想所成集合
$G \subset \Spec(S)$（赋予诱导拓扑）同胚地映到 $\Spec(S_0)$。
\end{lemma}

\begin{proof}
先构造逆映射以证明该映射是双射。设 $f \in S_d$、$d > 0$
在 $S$ 中可逆。若 $\mathfrak p_0$ 是 $S_0$ 的素理想，则
$\mathfrak p_0S$ 是 $S$ 的 $\mathbf{Z}$-分次理想，且
$\mathfrak p_0S \cap S_0 = \mathfrak p_0$。若
$ab \in \mathfrak p_0S$，其中 $a$、$b$ 齐次，则
$a^db^d/f^{\deg(a) + \deg(b)} \in \mathfrak p_0$。
所以 $a^d/f^{\deg(a)} \in \mathfrak p_0$ 或
$b^d/f^{\deg(b)} \in \mathfrak p_0$；换言之，
$a^d \in \mathfrak p_0S$ 或 $b^d \in \mathfrak p_0S$。
因此 $\sqrt{\mathfrak p_0S}$ 是 $S$ 的 $\mathbf{Z}$-分次素理想，
且它与 $S_0$ 的交为 $\mathfrak p_0$。

\medskip\noindent
为证明该映射是同胚，只须证明 $G \cap D(g)$ 的像是开集。
若 $g = \sum g_i$，其中 $g_i \in S_i$，则由上述结论，
$G \cap D(g)$ 映满开集 $\bigcup D(g_i^d/f^i)$。
\end{proof}

\noindent
对齐次元 $f \in S$，若其次数为 $> 0$，则定义
$$
D_{+}(f) = \{ \mathfrak p \in \text{Proj}(S) \mid f \not\in \mathfrak p \}.
$$
最后，对齐次理想 $I \subset S$，定义
$$
V_{+}(I) = \{ \mathfrak p \in \text{Proj}(S) \mid I \subset \mathfrak p \}.
$$
更一般地，对齐次元的任意集合 $E$、$E \subset S$，也使用记号
$V_{+}(E)$。

\begin{lemma}[Proj 上的拓扑]
\label{algebra-lemma-topology-proj}
设 $S = \oplus_{d \geq 0} S_d$ 为分次环。
\begin{enumerate}
\item 诸集合 $D_{+}(f)$ 是 $\text{Proj}(S)$ 中的开集。
\item 有 $D_{+}(ff') = D_{+}(f) \cap D_{+}(f')$。
\item 设 $g = g_0 + \ldots + g_m$ 是 $S$ 的元素，且
$g_i \in S_i$。则
$$
D(g) \cap \text{Proj}(S) =
(D(g_0) \cap \text{Proj}(S))
\cup
\bigcup\nolimits_{i \geq 1} D_{+}(g_i).
$$
\item
设 $g_0\in S_0$ 是次数为 $0$ 的齐次元。则
$$
D(g_0) \cap \text{Proj}(S)
=
\bigcup\nolimits_{f \in S_d, \ d\geq 1} D_{+}(g_0 f).
$$
\item 开集 $D_{+}(f)$ 构成 $\text{Proj}(S)$ 的拓扑的一组基。
\item 设 $f \in S$ 是正次数齐次元。环 $S_f$ 有自然的
$\mathbf{Z}$-分次。环同态 $S \to S_f \leftarrow S_{(f)}$ 诱导同胚
$$
D_{+}(f)
\leftarrow
\{\mathbf{Z}\text{-分次素理想，位于 }S_f\}
\to
\Spec(S_{(f)}).
$$
\item 存在某个 $S$，使得 $\text{Proj}(S)$ 不是拟紧的。
\item 诸集合 $V_{+}(I)$ 是闭集。
\item 任意闭子集 $T \subset \text{Proj}(S)$ 都具有形式
$V_{+}(I)$，其中 $I \subset S$ 是某个齐次理想。
\item 对任意分次理想 $I \subset S$，有
$V_{+}(I) = \emptyset$，当且仅当 $S_{+} \subset \sqrt{I}$。
\end{enumerate}
\end{lemma}

\begin{proof}
由于 $D_{+}(f) = \text{Proj}(S) \cap D(f)$，这些集合是开集，
从而 (1) 得证。又因为 $D(ff') = D(f) \cap D(f')$，(2) 亦得证。
同理，诸集合 $V_{+}(I) = \text{Proj}(S) \cap V(I)$ 是闭集，
这证明了 (8)。

\medskip\noindent
假设 $T \subset \text{Proj}(S)$ 是闭集。则对某个理想
$J \subset S$，可写成 $T = \text{Proj}(S) \cap V(J)$。
由齐次理想的定义，若 $g \in J$、$g = g_0 + \ldots + g_m$，
其中 $g_d \in S_d$，则对所有 $\mathfrak p \in T$ 都有
$g_d \in \mathfrak p$。因此，令 $I \subset S$ 为由 $J$ 中诸元素的
齐次部分生成的理想，便有 $T = V_{+}(I)$。这证明了 (9)。

\medskip\noindent
由定义立即得到 $\text{Proj}(S) \cap D(g)$ 的公式，其中 $g \in S$。
这证明了 (3)。再考虑 $\text{Proj}(S) \cap D(g_0)$ 的公式。
右端包含于左端是显然的。反过来，假设
$g_0 \not \in \mathfrak p$，其中 $\mathfrak p \in \text{Proj}(S)$。
若对所有正次数齐次元 $f$ 都有 $g_0f \in \mathfrak p$，
则 $S_{+} \subset \mathfrak p$，矛盾。这给出反向包含，
从而证明 (4)。

\medskip\noindent
开集族 $D(g) \cap \text{Proj}(S)$ 构成一组拓扑基，因为标准开集
$D(g) \subset \Spec(S)$ 构成 $\Spec(S)$ 的一组拓扑基。
由上述公式，可将 $D(g) \cap \text{Proj}(S)$ 表成开集
$D_{+}(f)$ 之并。因此，开集族 $D_{+}(f)$ 也构成一组拓扑基。
这证明了 (5)。

\medskip\noindent
证明 (6)。首先，由引理 \ref{algebra-lemma-standard-open}，可将
$D_{+}(f)$ 等同为 $D(f) = \Spec(S_f)$ 的一个子集
（赋予诱导拓扑）。注意，环 $S_f$ 有一个 $\mathbf{Z}$-分次。
其齐次元具有形式 $r/f^n$，其中 $r \in S$ 齐次，且次数为
$\deg(r/f^n) = \deg(r) - n\deg(f)$。子集 $D_{+}(f)$ 恰好对应于
那些为 $\mathbf{Z}$-分次理想（即由齐次元生成）的素理想
$\mathfrak p \subset S_f$。因此只须证明，$S_f$ 的
$\mathbf{Z}$-分次素理想所成的集合同胚地映到
$\Spec(S_{(f)})$。这由引理 \ref{algebra-lemma-Z-graded} 得出。

\medskip\noindent
令 $S = \mathbf{Z}[X_1, X_2, X_3, \ldots]$，并赋予分次，
使每个 $X_i$ 的次数为 $1$。不难看出，覆盖
$$
\text{Proj}(S) = \bigcup\nolimits_{i = 1}^\infty D_{+}(X_i)
$$
不存在有限加细。这证明了 (7)。

\medskip\noindent
设 $I \subset S$ 为分次理想。若 $\sqrt{I} \supset S_{+}$，
则 $V_{+}(I) = \emptyset$，因为按照定义，每个素理想
$\mathfrak p \in \text{Proj}(S)$ 都不包含 $S_{+}$。
反过来，假设 $S_{+} \not \subset \sqrt{I}$。于是可以找到
$f \in S_{+}$，使得 $f$ 模 $I$ 不幂零。显然，这意味着 $f$
的某个齐次部分模 $I$ 不幂零；换言之，可以且确实假设 $f$ 齐次。
这蕴含 $I S_f \not = S_f$，也就是 $(S/I)_f$ 非零。
又因为 $(S/I)_{(f)} \not = 0$，它是映入 $(S/I)_f$ 的环。
选取素理想 $\mathfrak q \subset (S/I)_{(f)}$。它对应于
$S/I$ 的一个不包含无关理想 $(S/I)_{+}$ 的分次素理想，
后者又对应于 $S$ 的一个包含 $I$、但不包含 $S_{+}$ 的分次素理想
$\mathfrak p$，正合所需。这证明了 (10)，并完成证明。
\end{proof}

\begin{example}
\label{algebra-example-proj-polynomial-ring-1-variable}
设 $R$ 为环。若 $S = R[X]$ 且 $\deg(X) = 1$，则自然映射
$\text{Proj}(S) \to \Spec(R)$ 是双射，实际上是同胚。
事实上，假设 $\mathfrak p \in \text{Proj}(S)$。由于
$S_{+} \not \subset \mathfrak p$，可知 $X \not \in \mathfrak p$。
所以，若 $aX^n \in \mathfrak p$，其中 $a \in R$ 且 $n > 0$，
则 $a \in \mathfrak p$。由此
$\mathfrak p = \mathfrak p_0S$，其中
$\mathfrak p_0 = \mathfrak p \cap R$。
\end{example}

\noindent
若 $\mathfrak p \in \text{Proj}(S)$，则定义 $S_{(\mathfrak p)}$
为如下环：其元素是分式 $r/f$，其中 $r, f \in S$ 是同次数的齐次元，
且 $f \not\in \mathfrak p$。与通常一样，规定
$r/f = r'/f'$，当且仅当存在某个齐次元 $f'' \in S$，
满足 $f'' \not \in \mathfrak p$ 且 $f''(rf' - r'f) = 0$。
给定分次 $S$-模 $M$，令 $M_{(\mathfrak p)}$ 为如下
$S_{(\mathfrak p)}$-模：其元素是分式 $x/f$，其中
$x \in M$ 与 $f \in S$ 齐次且次数相同，并且
$f \not \in \mathfrak p$。规定 $x/f = x'/f'$，当且仅当存在某个
齐次元 $f'' \in S$，满足 $f'' \not \in \mathfrak p$ 且
$f''(xf' - x'f) = 0$。

\begin{lemma}
\label{algebra-lemma-proj-prime}
设 $S$ 为分次环，$M$ 为分次 $S$-模。设 $\mathfrak p$ 是
$\text{Proj}(S)$ 的一个元素。设 $f \in S$ 是正次数齐次元，
且 $f \not \in \mathfrak p$，即 $\mathfrak p \in D_{+}(f)$。
设 $\mathfrak p' \subset S_{(f)}$ 是 $\Spec(S_{(f)})$ 中按照
引理 \ref{algebra-lemma-topology-proj} 对应于 $\mathfrak p$ 的元素。则
$S_{(\mathfrak p)} = (S_{(f)})_{\mathfrak p'}$，并且相容地，
$M_{(\mathfrak p)} = (M_{(f)})_{\mathfrak p'}$。
\end{lemma}

\begin{proof}
定义映射 $\psi : M_{(\mathfrak p)} \to (M_{(f)})_{\mathfrak p'}$。
设 $x/g \in M_{(\mathfrak p)}$。置
$$
\psi(x/g) = (x g^{\deg(f) - 1}/f^{\deg(x)})/(g^{\deg(f)}/f^{\deg(g)}).
$$
这是有意义的，因为 $\deg(x) = \deg(g)$，且
$g^{\deg(f)}/f^{\deg(g)} \not \in \mathfrak p'$。
关于 $\psi$ 是良定义的模同态且为同构的验证从略。提示：
其逆映射把 $(x/f^n)/(g/f^m)$ 送到 $(xf^m)/(g f^n)$。
\end{proof}

\noindent
下面是引理 \ref{algebra-lemma-silly} 的一个分次版本。

\begin{lemma}
\label{algebra-lemma-graded-silly}
设 $S$ 为分次环，$\mathfrak p_i$（$i = 1, \ldots, r$）为齐次素理想，
$I \subset S_{+}$ 为分次理想。假设对所有 $i$ 都有
$I \not\subset \mathfrak p_i$。则存在正次数齐次元 $x\in I$，
使得对所有 $i$ 都有 $x\not\in \mathfrak p_i$。
\end{lemma}

\begin{proof}
可以假设诸 $\mathfrak p_i$ 之间不存在包含关系。当 $r = 1$ 时结论成立。
假设对 $r - 1$ 结论成立。选取正次数齐次元 $x \in I$，
使得对所有 $i = 1, \ldots, r - 1$ 都有
$x \not \in \mathfrak p_i$。若 $x \not\in \mathfrak p_r$，证明完毕。
故假设 $x \in \mathfrak p_r$。若
$I \mathfrak p_1 \ldots \mathfrak p_{r-1} \subset \mathfrak p_r$，
则 $I \subset \mathfrak p_r$，矛盾。选取齐次元
$y \in I\mathfrak p_1 \ldots \mathfrak p_{r-1}$，使得
$y \not \in \mathfrak p_r$。此时
$x^{\deg(y)} + y^{\deg(x)}$ 符合要求。
\end{proof}

\begin{lemma}
\label{algebra-lemma-smear-out}
设 $S$ 为分次环，$\mathfrak p \subset S$ 为素理想。
令 $\mathfrak q$ 为由 $\mathfrak p$ 的齐次元生成的 $S$ 的齐次理想。
则 $\mathfrak q$ 是 $S$ 的素理想。
\end{lemma}

\begin{proof}
为证明 $\mathfrak q$ 为素理想，只须检验：若 $f, g \in S$ 齐次且
$fg \in \mathfrak q$，则 $f$ 或 $g$ 属于 $\mathfrak q$。
因为 $\mathfrak p \subset \mathfrak q$，所以
$fg \in \mathfrak p$。由于 $\mathfrak p$ 为素理想，可知
$f \in \mathfrak p$ 或 $g \in \mathfrak p$。
又因为 $f$ 与 $g$ 齐次，显然 $f \in \mathfrak q$ 或
$g \in \mathfrak q$。
\end{proof}

\begin{lemma}
\label{algebra-lemma-graded-ring-minimal-prime}
设 $S$ 为分次环。
\begin{enumerate}
\item $S$ 的任意极小素理想都是 $S$ 的齐次理想。
\item 给定齐次理想 $I \subset S$，$I$ 上的任意极小素理想都是齐次的。
\end{enumerate}
\end{lemma}

\begin{proof}
第一条成立，因为引理 \ref{algebra-lemma-smear-out} 中构造的素理想
$\mathfrak q$ 满足 $\mathfrak q \subset \mathfrak p$。
第二条成立，因为可以考虑 $S/I$ 并应用第一条。
\end{proof}

\begin{lemma}
\label{algebra-lemma-dehomogenize-finite-type}
设 $R$ 为环，$S$ 为分次 $R$-代数，$f \in S_{+}$ 为齐次元。
假设 $S$ 在 $R$ 上是有限型的。则
\begin{enumerate}
\item 环 $S_{(f)}$ 在 $R$ 上是有限型的；
\item 对任意有限分次 $S$-模 $M$，模 $M_{(f)}$ 是有限
$S_{(f)}$-模。
\end{enumerate}
\end{lemma}

\begin{proof}
选取 $f_1, \ldots, f_n \in S$，使它们作为 $R$-代数生成 $S$。
可以假设每个 $f_i$ 都齐次（把每个 $f_i$ 分解成其齐次分量即可）。
$S_{(f)}$ 的一个元素是如下形式的和：
$$
\sum\nolimits_{e\deg(f) =
\sum e_i\deg(f_i)} \lambda_{e_1 \ldots e_n} f_1^{e_1} \ldots f_n^{e_n}/f^e
$$
其中 $\lambda_{e_1 \ldots e_n} \in R$。因此，作为 $R$-代数，
$S_{(f)}$ 由满足 $e\deg(f) = \sum e_i\deg(f_i)$ 的诸元素
$f_1^{e_1} \ldots f_n^{e_n} /f^e$ 生成。若
$e_i \geq \deg(f)$，则可将它写成
$$
f_1^{e_1} \ldots f_n^{e_n}/f^e =
f_i^{\deg(f)}/f^{\deg(f_i)} \cdot
f_1^{e_1} \ldots f_i^{e_i - \deg(f)} \ldots f_n^{e_n}/f^{e - \deg(f_i)}
$$
所以，只需要诸元素 $f_i^{\deg(f)}/f^{\deg(f_i)}$，以及满足
$e \deg(f) = \sum e_i \deg(f_i)$ 和 $e_i < \deg(f)$ 的诸元素
$f_1^{e_1} \ldots f_n^{e_n} /f^e$。这是一个有限表，故 (1) 成立。

\medskip\noindent
为证明 (2)，假设 $M$ 由齐次元 $x_1, \ldots, x_m$ 生成。
与上述论证相同，可知 $M_{(f)}$ 作为 $S_{(f)}$-模由如下有限表中的
元素生成：
$f_1^{e_1} \ldots f_n^{e_n} x_j /f^e$，其中
$e \deg(f) = \sum e_i \deg(f_i) + \deg(x_j)$ 且
$e_i < \deg(f)$。
\end{proof}

\begin{lemma}
\label{algebra-lemma-homogenize}
设 $R$ 为环，$R'$ 为有限型 $R$-代数，$M$ 为有限 $R'$-模。
则存在分次 $R$-代数 $S$、分次 $S$-模 $N$ 和一个次数为 $1$
的齐次元 $f \in S$，使得
\begin{enumerate}
\item $R' \cong S_{(f)}$ 且 $M \cong N_{(f)}$（作为模）；
\item $S_0 = R$，且 $S$ 在 $R$ 上由有限多个次数为 $1$ 的元素生成；
\item $N$ 是有限 $S$-模。
\end{enumerate}
\end{lemma}

\begin{proof}
对某个理想 $I$，可写 $R' = R[x_1, \ldots, x_n]/I$。
对元素 $g \in R[x_1, \ldots, x_n]$，用
$\tilde g \in R[X_0, \ldots, X_n]$ 表示满足
$g = \tilde g(1, x_1, \ldots, x_n)$ 的最小次数齐次元。
令 $\tilde I \subset R[X_0, \ldots, X_n]$ 为由所有元素
$\tilde g$（$g \in I$）生成的理想。置
$S = R[X_0, \ldots, X_n]/\tilde I$，并以 $f$ 表示 $X_0$
在 $S$ 中的像。由构造有同构
$$
S_{(f)} \longrightarrow R', \quad
X_i/X_0 \longmapsto x_i.
$$
为对模 $M$ 作同样的构造，选取表示
$$
M = (R')^{\oplus r}/\sum\nolimits_{j \in J} R'k_j
$$
其中 $k_j = (k_{1j}, \ldots, k_{rj})$。令
$d_{ij} = \deg(\tilde k_{ij})$，置 $d_j = \max\{d_{ij}\}$。
置 $K_{ij} = X_0^{d_j - d_{ij}}\tilde k_{ij}$；
它是次数为 $d_j$ 的齐次元。用这些记号，置
$$
N = \Coker\Big(
\bigoplus\nolimits_{j \in J} S(-d_j) \xrightarrow{(K_{ij})} S^{\oplus r}
\Big)
$$
即可。若干细节从略。
\end{proof}




\section{诺特分次环}
\label{algebra-section-noetherian-graded}

\noindent
本节给出一些诺特分次环的理论，包括关于希尔伯特多项式的若干材料。

\begin{lemma}
\label{algebra-lemma-S-plus-generated}
设 $S$ 为分次环。一族齐次元 $f_i \in S_{+}$ 作为 $S_0$-代数生成
$S$，当且仅当它们作为 $S$ 的理想生成 $S_{+}$。
\end{lemma}

\begin{proof}
若诸 $f_i$ 作为 $S_0$-代数生成 $S$，则 $S_{+}$ 中的每个元素
都是诸 $f_i$ 的无常数项多项式，因此 $S_{+}$ 作为理想由诸
$f_i$ 生成。反过来，假设 $S_{+} = \sum Sf_i$。将证明 $S$ 的任意元素
$f$ 都可以写成以诸 $f_i$ 为变量、系数在 $S_0$ 中的多项式。
只须对齐次元证明。设 $f$ 的次数为 $d$，对 $d$ 归纳。
$d = 0$ 的情形显然。若 $d > 0$，则 $f \in S_{+}$，故对某些
$g_i \in S$，可写成 $f = \sum g_i f_i$。由于 $S$ 分次，
可以用 $g_i$ 的次数为 $d - \deg(f_i)$ 的齐次分量代替它。
由归纳假设，每个 $g_i$ 都是诸 $f_i$ 的多项式，结论随即成立。
\end{proof}

\begin{lemma}
\label{algebra-lemma-graded-Noetherian}
分次环 $S$ 是诺特环，当且仅当 $S_0$ 是诺特环，且
$S_{+}$ 作为 $S$ 的理想有限生成。
\end{lemma}

\begin{proof}
显然，若 $S$ 是诺特环，则 $S_0 = S/S_{+}$ 是诺特环，且
$S_{+}$ 有限生成。反过来，假设 $S_0$ 是诺特环，且 $S_{+}$
作为 $S$ 的理想有限生成。选取生成元
$S_{+} = (f_1, \ldots, f_n)$。将诸 $f_i$ 分解为齐次部分后，
可以假设每个 $f_i$ 都齐次。由引理 \ref{algebra-lemma-S-plus-generated}，
把 $X_i$ 送到 $f_i$ 的映射
$S_0[X_1, \ldots X_n] \to S$ 是满射。因此由引理
\ref{algebra-lemma-Noetherian-permanence}，$S$ 是诺特环。
\end{proof}

\begin{definition}
\label{algebra-definition-numerical-polynomial}
设 $A$ 为阿贝尔群。设函数 $f : n \mapsto f(n) \in A$
对所有充分大的整数 $n$ 有定义。若存在 $r \geq 0$ 及元素
$a_0, \ldots, a_r\in A$，使得对所有 $n \gg 0$ 都有
$$
f(n) = \sum\nolimits_{i = 0}^r \binom{n}{i} a_i
$$
则称该函数为{\it 数值多项式}。
\end{definition}

\noindent
使用二项式系数的理由是如下初等事实：若多项式
$P \in \mathbf{Q}[T]$ 在所有整数点的值都是整数，则它等于某个和
$P(T) = \sum a_i \binom{T}{i}$，其中 $a_i \in \mathbf{Z}$。
特别地，注意表达式 $\binom{T + 1}{i + 1}$ 具有这种形式。

\begin{lemma}
\label{algebra-lemma-numerical-polynomial-functorial}
若 $A \to A'$ 是阿贝尔群同态，且
$f : n \mapsto f(n) \in A$ 是数值多项式，则其复合也是数值多项式。
\end{lemma}

\begin{proof}
这由定义立即得出。
\end{proof}

\begin{lemma}
\label{algebra-lemma-numerical-polynomial}
假设 $f: n \mapsto f(n) \in A$ 对所有充分大的 $n$ 有定义，
并假设 $n \mapsto f(n) - f(n-1)$ 是数值多项式。
则 $f$ 是数值多项式。
\end{lemma}

\begin{proof}
假设对所有 $n \gg 0$ 都有
$f(n) - f(n-1) = \sum\nolimits_{i = 0}^r \binom{n}{i} a_i$。
置 $g(n) = f(n) - \sum\nolimits_{i = 0}^r \binom{n + 1}{i + 1} a_i$。
于是对所有 $n \gg 0$，有 $g(n) - g(n-1) = 0$。
故 $g$ 最终为常值，设其等于 $a_{-1}$。留给读者证明
$a_{-1} + \sum\nolimits_{i = 0}^r \binom{n + 1}{i + 1} a_i$
具有所需形式（参见本引理上方的说明）。
\end{proof}

\begin{lemma}
\label{algebra-lemma-graded-module-fg}
若 $M$ 是有限生成分次 $S$-模，且 $S$ 在 $S_0$ 上有限生成，
则每个 $M_n$ 都是有限 $S_0$-模。
\end{lemma}

\begin{proof}
设 $M$ 的生成元为 $m_i$，$S$ 的生成元为 $f_i$。
取齐次分量后，可以假设诸 $m_i$ 与诸 $f_i$ 都齐次，
并可假设 $f_i \in S_{+}$。此时显然，每个 $M_n$ 作为
$S_0$-模由次数为 $n$ 的“单项式”$\prod f_i^{e_i} m_j$ 生成。
\end{proof}

\begin{proposition}
\label{algebra-proposition-graded-hilbert-polynomial}
假设 $S$ 是诺特分次环，$M$ 是有限分次 $S$-模。考虑函数
$$
\mathbf{Z} \longrightarrow K'_0(S_0), \quad
n \longmapsto [M_n]
$$
参见引理 \ref{algebra-lemma-graded-module-fg}。若 $S_{+}$ 由次数为 $1$
的元素生成，则该函数是数值多项式。
\end{proposition}

\begin{proof}
对 $S_1$ 的最小生成元数归纳证明。若该数为 $0$，则对所有
$n \gg 0$ 都有 $M_n = 0$，结论成立。为证明归纳步骤，
从一个极小生成元集中选取 $x\in S_1$，使归纳假设适用于分次环
$S/(x)$。

\medskip\noindent
先证明当 $x$ 在 $M$ 上幂零时结论成立。对此，再关于满足
$x^r M  = 0$ 的最小整数 $r$ 归纳。若 $r = 1$，则 $M$ 是
$S/xS$ 上的模，结论由另一个归纳假设成立。若 $r > 1$，
则可以找到短正合列 $0 \to M' \to M \to M'' \to 0$，
使整数 $r', r''$ 都严格小于 $r$。于是已知结论对 $M''$ 与
$M'$ 成立。由于 $K'_0(S_0)$ 中有关系
$
[M_d]  = [M'_d] + [M''_d]
$
故也得到 $M$ 的结论。

\medskip\noindent
若 $x$ 在 $M$ 上不幂零，设 $M' \subset M$ 是 $x$
在其上幂零的最大子模。考虑正合列
$0 \to M' \to M \to M/M' \to 0$，再次可见只须对 $M/M'$
证明结论。换言之，可以假设乘以 $x$ 是单射。

\medskip\noindent
令 $\overline{M} = M/xM$。注意，映射 $x : M \to M$
{\it 不是}分次 $S$-模同态，因为它不把 $M_d$ 映入 $M_d$。
具体而言，对每个 $d$ 都有短正合列
$$
0 \to M_d \xrightarrow{x} M_{d + 1} \to \overline{M}_{d + 1} \to 0
$$
这证明 $[M_{d + 1}] - [M_d] = [\overline{M}_{d + 1}]$。
故由引理 \ref{algebra-lemma-numerical-polynomial} 得到结论。
\end{proof}

\begin{remark}
\label{algebra-remark-period-polynomial}
若 $S$ 仍是诺特环，但 $S$ 不是由次数为 $1$ 的元素生成的，
则与分次 $S$-模相伴的函数是周期多项式
（即对某个 $n$，它在整数模 $n$ 的各同余类上都是数值多项式）。
\end{remark}

\begin{example}
\label{algebra-example-hilbert-function}
假设 $S = k[X_1, \ldots, X_d]$。由例 \ref{algebra-example-K0-field}，
可作等同 $K_0(k) = K'_0(k) = \mathbf{Z}$。因此，任意有限生成分次
$k[X_1, \ldots, X_d]$-模都给出数值多项式
$n \mapsto \dim_k(M_n)$。
\end{example}

\begin{lemma}
\label{algebra-lemma-quotient-smaller-d}
设 $k$ 为域。假设 $I \subset k[X_1, \ldots, X_d]$ 是非零分次理想。
令 $M = k[X_1, \ldots, X_d]/I$。则数值多项式
$n \mapsto \dim_k(M_n)$（参见例 \ref{algebra-example-hilbert-function}）
的次数为 $ < d - 1$（若 $d = 1$，则它为零）。
\end{lemma}

\begin{proof}
与分次模 $k[X_1, \ldots, X_d]$ 相伴的数值多项式是
$n \mapsto \binom{n - 1 + d}{d - 1}$。对任意次数为 $e$ 的非零齐次元
$f \in I$ 以及任意次数 $n >> e$，都有
$I_n \supset f \cdot k[X_1, \ldots, X_d]_{n-e}$，因而
$\dim_k(I_n) \geq \binom{n - e - 1 + d}{d - 1}$。所以
$\dim_k(M_n) \leq \binom{n - 1 + d}{d - 1} - \binom{n - e - 1 + d}{d - 1}$。
最后一个表达式的次数为 $ < d - 1$（若 $d = 1$，则它为零），
故结论成立。
\end{proof}








\section{诺特局部环}
\label{algebra-section-Noetherian-local}

\noindent
在本节中，$(R, \mathfrak m, \kappa)$ 始终表示诺特局部环。
本节将发展一些模的希尔伯特函数理论。设 $M$ 为有限 $R$-模。
定义 $M$ 的{\it 希尔伯特函数}为函数
$$
\varphi_M : n
\longmapsto
\text{length}_R(\mathfrak m^nM/{\mathfrak m}^{n + 1}M)
$$
它对所有整数 $n \geq 0$ 有定义。另一个重要不变量是函数
$$
\chi_M : n
\longmapsto
\text{length}_R(M/{\mathfrak m}^{n + 1}M)
$$
它也对所有整数 $n \geq 0$ 有定义。注意，由引理
\ref{algebra-lemma-length-additive} 有
$$
\chi_M(n) = \sum\nolimits_{i = 0}^n \varphi_M(i).
$$
这一构造还有一个使用定义理想的变体。

\begin{definition}
\label{algebra-definition-ideal-definition}
设 $(R, \mathfrak m)$ 为诺特局部环。满足
$\sqrt{I} = \mathfrak m$ 的理想 $I \subset R$ 称为
{\it $R$ 的定义理想}。
\end{definition}

\noindent
设 $I \subset R$ 为定义理想。因为 $R$ 是诺特环，这意味着对某个
$r$ 有 $\mathfrak m^r \subset I$；参见引理
\ref{algebra-lemma-Noetherian-power}。因此，被 $I$ 的某个幂零化的任意有限
$R$-模都有有限长度；参见引理 \ref{algebra-lemma-length-finite}。
所以可以定义
$$
\varphi_{I, M}(n) = \text{length}_R(I^nM/I^{n + 1}M)
\quad\text{且}\quad
\chi_{I, M}(n) = \text{length}_R(M/I^{n + 1}M)
$$
其中 $n \geq 0$。同样有
$$
\chi_{I, M}(n) = \sum\nolimits_{i = 0}^n \varphi_{I, M}(i).
$$

\begin{lemma}
\label{algebra-lemma-differ-finite}
假设 $M' \subset M$ 是有限 $R$-模，且商的长度有限。
则存在常数 $c_1, c_2$，使得对所有 $n \geq c_2$ 都有
$$
c_1 + \chi_{I, M'}(n - c_2) \leq \chi_{I, M}(n) \leq
c_1 + \chi_{I, M'}(n)
$$
\end{lemma}

\begin{proof}
由于 $M/M'$ 长度有限，存在 $c_2 \geq 0$，使得
$I^{c_2}M \subset M'$。令 $c_1 = \text{length}_R(M/M')$。
对 $n \geq c_2$，有
\begin{eqnarray*}
\chi_{I, M}(n)
& = &
\text{length}_R(M/I^{n + 1}M) \\
& = &
c_1 + \text{length}_R(M'/I^{n + 1}M) \\
& \leq &
c_1 + \text{length}_R(M'/I^{n + 1}M') \\
& = &
c_1 + \chi_{I, M'}(n)
\end{eqnarray*}
另一方面，由于 $I^{c_2}M \subset M'$，对 $n \geq c_2$ 有
$I^nM \subset I^{n - c_2}M'$。因此，对 $n \geq c_2$ 得
\begin{eqnarray*}
\chi_{I, M}(n)
& = &
\text{length}_R(M/I^{n + 1}M) \\
& = &
c_1 + \text{length}_R(M'/I^{n + 1}M) \\
& \geq &
c_1 + \text{length}_R(M'/I^{n + 1 - c_2}M') \\
& = &
c_1 + \chi_{I, M'}(n - c_2)
\end{eqnarray*}
证明完毕。
\end{proof}

\begin{lemma}
\label{algebra-lemma-hilbert-ses}
假设 $0 \to M' \to M \to M'' \to 0$ 是有限 $R$-模的短正合列。
则存在余长度为 $l$ 的子模 $N \subset M'$ 及 $c \geq 0$，使得
$$
\chi_{I, M}(n) = \chi_{I, M''}(n) + \chi_{I, N}(n - c) + l
$$
以及
$$
\varphi_{I, M}(n) = \varphi_{I, M''}(n) + \varphi_{I, N}(n - c)
$$
对所有 $n \geq c$ 成立。
\end{lemma}

\begin{proof}
注意，$M/I^nM \to M''/I^nM''$ 是满射，其核为
$M' / M' \cap I^nM$。由阿廷--里斯引理
\ref{algebra-lemma-Artin-Rees}，存在常数 $c$，使得
$M' \cap I^nM = I^{n - c}(M' \cap I^cM)$。记
$N = M' \cap I^cM$。注意，$I^c M' \subset N \subset M'$。
因此，对 $n \geq c$，
$\text{length}_R(M' / M' \cap I^nM)
= \text{length}_R(M'/N) + \text{length}_R(N/I^{n - c}N)$。
由短正合列
$$
0 \to M' / M' \cap I^nM \to M/I^nM \to M''/I^nM'' \to 0
$$
及长度的可加性（引理 \ref{algebra-lemma-length-additive}），得到等式
$$
\chi_{I, M}(n - 1)
=
\chi_{I, M''}(n - 1)
+
\chi_{I, N}(n - c - 1)
+
\text{length}_R(M'/N)
$$
它对 $n \geq c$ 成立。有
$\varphi_{I, M}(n) = \chi_{I, M}(n) - \chi_{I, M}(n - 1)$，
对模 $M''$ 与 $N$ 亦然。因此，对 $n \geq c$ 得
$\varphi_{I, M}(n) = \varphi_{I, M''}(n) + \varphi_{I, N}(n-c)$。
\end{proof}

\begin{lemma}
\label{algebra-lemma-hilbert-change-I}
假设 $I$、$I'$ 是诺特局部环 $R$ 的两个定义理想。
设 $M$ 为有限 $R$-模。则存在常数 $a$，使得对 $n \geq 1$ 有
$\chi_{I, M}(n) \leq \chi_{I', M}(an)$。
\end{lemma}

\begin{proof}
存在整数 $c \geq 1$，使得 $(I')^c \subset I$。
因此有满射 $M/(I')^{c(n + 1)}M \to M/I^{n + 1}M$。
所以取 $a = 2c - 1$ 即得结论。
\end{proof}

\begin{proposition}
\label{algebra-proposition-hilbert-function-polynomial}
设 $R$ 为诺特局部环，$M$ 为有限 $R$-模，$I \subset R$ 为定义理想。
希尔伯特函数 $\varphi_{I, M}$ 与函数 $\chi_{I, M}$ 都是数值多项式。
\end{proposition}

\begin{proof}
考虑分次环
$S = R/I \oplus I/I^2 \oplus I^2/I^3 \oplus
\ldots = \bigoplus_{d \geq 0} I^d/I^{d + 1}$。再考虑分次
$S$-模
$N = M/IM \oplus IM/I^2M \oplus \ldots =
\bigoplus_{d \geq 0} I^dM/I^{d + 1}M$。这一对 $(S, N)$ 满足命题
\ref{algebra-proposition-graded-hilbert-polynomial} 的假设。
因此，关于 $\varphi_{I, M}$ 的结论由该命题和引理
\ref{algebra-lemma-length-K0} 得出；关于 $\chi_{I, M}$ 的结论再由此及
引理 \ref{algebra-lemma-numerical-polynomial} 得出。
\end{proof}

\begin{definition}
\label{algebra-definition-hilbert-polynomial}
设 $R$ 为诺特局部环，$M$ 为有限 $R$-模。$M$ 在 $R$ 上的
{\it 希尔伯特多项式}，是满足对 $n \gg 0$ 有
$P(n) = \varphi_M(n)$ 的元素 $P(t) \in \mathbf{Q}[t]$。
\end{definition}

\noindent
由命题 \ref{algebra-proposition-hilbert-function-polynomial} 可知，
希尔伯特多项式存在。

\begin{lemma}
\label{algebra-lemma-d-independent}
设 $R$ 为诺特局部环，$M$ 为有限 $R$-模。
\begin{enumerate}
\item 数值多项式 $\varphi_{I, M}$ 的次数不依赖于定义理想 $I$。
\item 数值多项式 $\chi_{I, M}$ 的次数不依赖于定义理想 $I$。
\end{enumerate}
\end{lemma}

\begin{proof}
第 (2) 条由引理 \ref{algebra-lemma-hilbert-change-I} 立即得出。
第 (1) 条由第 (2) 条得出，因为对 $n \geq 1$ 有
$\varphi_{I, M}(n) = \chi_{I, M}(n) - \chi_{I, M}(n - 1)$。
\end{proof}

\begin{definition}
\label{algebra-definition-d}
设 $R$ 为诺特局部环，$M$ 为有限 $R$-模。用{\it $d(M)$}
表示 $\{-\infty, 0, 1, 2, \ldots \}$ 中按如下方式定义的元素：
\begin{enumerate}
\item 若 $M = 0$，置 $d(M) = -\infty$；
\item 若 $M \not = 0$，则 $d(M)$ 是数值多项式 $\chi_M$ 的次数。
\end{enumerate}
\end{definition}

\noindent
若对所有 $n$ 都有 $\mathfrak m^nM \not = 0$，则 $d(M)$ 是
$M$ 的希尔伯特多项式的次数 $+1$。

\begin{lemma}
\label{algebra-lemma-differ-finite-chi}
设 $R$ 为诺特局部环，$I \subset R$ 为定义理想，$M$ 为不具有
有限长度的有限 $R$-模。若 $M' \subset M$ 是余长度有限的子模，
则 $\chi_{I, M} - \chi_{I, M'}$ 的次数为
$<$ 这两个多项式中任一个的次数。
\end{lemma}

\begin{proof}
由引理 \ref{algebra-lemma-differ-finite} 及初等计算得出。
\end{proof}

\begin{lemma}
\label{algebra-lemma-hilbert-ses-chi}
设 $R$ 为诺特局部环，$I \subset R$ 为定义理想，并设
$0 \to M' \to M \to M'' \to 0$ 是有限 $R$-模的短正合列。则
\begin{enumerate}
\item 若 $M'$ 不具有有限长度，则
$\chi_{I, M} - \chi_{I, M''} - \chi_{I, M'}$ 是次数为
$<$ $\chi_{I, M'}$ 的次数的数值多项式；
\item $\max\{ \deg(\chi_{I, M'}), \deg(\chi_{I, M''}) \} = \deg(\chi_{I, M})$；
\item $\max\{d(M'), d(M'')\} = d(M)$。
\end{enumerate}
\end{lemma}

\begin{proof}
先证明 (1)。设 $N \subset M'$ 如引理 \ref{algebra-lemma-hilbert-ses} 中所述。
由引理 \ref{algebra-lemma-differ-finite-chi}，数值多项式
$\chi_{I, M'} - \chi_{I, N}$ 的次数为
$<$ $\chi_{I, M'}$ 与 $\chi_{I, N}$ 的共同次数。
由引理 \ref{algebra-lemma-hilbert-ses}，差
$$
\chi_{I, M}(n) - \chi_{I, M''}(n) - \chi_{I, N}(n - c)
$$
对 $n \gg 0$ 为常数。由初等计算，
$\chi_{I, N}(n) - \chi_{I, N}(n - c)$ 的次数为
$<$ $\chi_{I, N}$ 的次数，后者大于零（见上文）。
合并这些结论即得 (1)。

\medskip\noindent
注意，$\chi_{I, M'}$ 与 $\chi_{I, M''}$ 的首项系数非负。
所以 $\chi_{I, M'} + \chi_{I, M''}$ 的次数等于二者次数的最大值。
因此，若 $M'$ 不具有有限长度，则 (2) 由 (1) 得出。
若 $M'$ 长度有限，则由阿廷--里斯定理（引理
\ref{algebra-lemma-Artin-Rees}），对所有 $n \gg 0$，映射
$I^nM \to I^nM''$ 都是同构。所以对所有 $n \gg 0$，
$M/I^nM \to M''/I^nM''$ 是以 $M'$ 为核的满射，因而
$\chi_{I, M}(n) - \chi_{I, M''}(n) = \text{length}(M')$。
该等式对所有 $n \gg 0$ 成立，所以此时 (2) 也成立。

\medskip\noindent
证明 (3)。除非 $M$、$M'$ 或 $M''$ 之一为零，否则它由 (2) 得出。
这些特殊情形的证明从略。
\end{proof}



































\section{维数}
\label{algebra-section-dimension}

\noindent
请与拓扑学，第 \ref{topology-section-krull-dimension} 节比较。

\begin{definition}
\label{algebra-definition-chain-primes}
设 $R$ 为环。所谓一条{\it 素理想链}，是 $R$ 的素理想序列
$\mathfrak p_0 \subset \mathfrak p_1 \subset \ldots \subset \mathfrak p_n$，
满足对 $i = 0, \ldots, n - 1$ 有
$\mathfrak p_i \not = \mathfrak p_{i + 1}$。这条素理想链的
{\it 长度}为 $n$。
\end{definition}

\noindent
回忆，环 $R$ 的素理想与 $\Spec(R)$ 的不可约闭子集之间有一个
反转包含关系的双射；参见引理 \ref{algebra-lemma-irreducible}。

\begin{definition}
\label{algebra-definition-Krull}
环 $R$ 的{\it 克鲁尔维数}，是拓扑空间 $\Spec(R)$ 的克鲁尔维数；
参见拓扑学，定义 \ref{topology-definition-Krull}。
换言之，它是满足如下条件的整数 $n\geq 0$ 的上确界：
$R$ 有一条素理想链
$$
\mathfrak p_0
\subset
\mathfrak p_1
\subset
\ldots
\subset
\mathfrak p_n, \quad
\mathfrak p_i \not = \mathfrak p_{i + 1}.
$$
其长度为 $n$。
\end{definition}

\begin{definition}
\label{algebra-definition-height}
环 $R$ 的素理想 $\mathfrak p$ 的{\it 高度}，是局部环
$R_{\mathfrak p}$ 的维数。
\end{definition}

\begin{lemma}
\label{algebra-lemma-dimension-height}
$R$ 的克鲁尔维数，是它的（极大）素理想之高度的上确界。
\end{lemma}

\begin{proof}
这是因为总可以在一条素理想链的末端添加一个极大理想。
\end{proof}

\begin{lemma}
\label{algebra-lemma-Noetherian-dimension-0}
维数为 $0$ 的诺特环是阿廷环。反过来，任意阿廷环都是维数为零的诺特环。
\end{lemma}

\begin{proof}
假设 $R$ 是维数为 $0$ 的诺特环。由引理
\ref{algebra-lemma-Noetherian-topology}，空间 $\Spec(R)$ 是诺特空间。
由拓扑学，引理 \ref{topology-lemma-Noetherian} 可知，
$\Spec(R)$ 只有有限多个不可约分支，设
$\Spec(R) = Z_1 \cup \ldots \cup Z_r$。按照引理
\ref{algebra-lemma-irreducible}，每个 $Z_i = V(\mathfrak p_i)$，
其中 $\mathfrak p_i$ 是极小理想。由于维数为 $0$，
这些 $\mathfrak p_i$ 也是极大理想。因此 $\Spec(R)$ 是以
$\mathfrak p_i$ 为元素的离散拓扑空间。雅可布森根
$\bigcap \mathfrak p_i$ 的所有元素 $f$ 都是幂零的；
否则 $R_f$ 不是零环，从而会有另一个素理想。
由引理 \ref{algebra-lemma-product-local}，
$R$ 等于 $\prod R_{\mathfrak p_i}$。由于
$R_{\mathfrak p_i}$ 也是维数为 $0$ 的诺特环，上述论证表明，
它的根 $\mathfrak p_iR_{\mathfrak p_i}$ 局部幂零。
由引理 \ref{algebra-lemma-Noetherian-power}，对某个 $n \geq 1$ 有
$\mathfrak p_i^nR_{\mathfrak p_i} = 0$。
由引理 \ref{algebra-lemma-length-finite}，可知 $R_{\mathfrak p_i}$
在 $R$ 上具有有限长度。因此，由引理
\ref{algebra-lemma-artinian-finite-length} 可知 $R$ 是阿廷环。

\medskip\noindent
若 $R$ 是阿廷环，则由引理 \ref{algebra-lemma-artinian-finite-length}，
它是诺特环。由引理 \ref{algebra-lemma-artinian-finite-nr-max}、
\ref{algebra-lemma-artinian-radical-nilpotent} 和
\ref{algebra-lemma-product-local} 合并可知，它的所有素理想都是极大理想。
\end{proof}

\noindent
下文将使用定义 \ref{algebra-definition-d} 中定义的不变量 $d(-)$。
先给出一个预备引理。

\begin{lemma}
\label{algebra-lemma-dimension-0-d-0}
设 $R$ 为诺特局部环。则
$\dim(R) = 0 \Leftrightarrow d(R) = 0$。
\end{lemma}

\begin{proof}
这是因为 $d(R) = 0$，当且仅当 $R$ 作为 $R$-模具有有限长度。
参见引理 \ref{algebra-lemma-artinian-finite-length}。
\end{proof}

\begin{proposition}
\label{algebra-proposition-dimension-zero-ring}
设 $R$ 为环。下列条件等价：
\begin{enumerate}
\item $R$ 是阿廷环；
\item $R$ 是诺特环且 $\dim(R) = 0$；
\item $R$ 作为自身上的模具有有限长度；
\item $R$ 是有限多个阿廷局部环的乘积；
\item $R$ 是诺特环，且 $\Spec(R)$ 是有限离散拓扑空间；
\item $R$ 是有限多个维数为 $0$ 的诺特局部环的乘积；
\item $R$ 是有限多个满足 $d(R_i) = 0$ 的诺特局部环 $R_i$ 的乘积；
\item $R$ 是有限多个极大理想幂零的诺特局部环 $R_i$ 的乘积；
\item $R$ 是诺特环，只有有限多个极大理想，且它的雅可布森根理想幂零；
\item $R$ 是诺特环，且它的素理想之间不存在严格包含关系。
\end{enumerate}
\end{proposition}

\begin{proof}
这是引理 \ref{algebra-lemma-product-local}、
\ref{algebra-lemma-artinian-finite-length}、
\ref{algebra-lemma-Noetherian-dimension-0} 和
\ref{algebra-lemma-dimension-0-d-0} 的综合。
\end{proof}

\begin{lemma}
\label{algebra-lemma-height-1}
设 $R$ 为诺特局部环。下列条件等价：
\begin{enumerate}
\item
\label{algebra-item-dim-1}
$\dim(R) = 1$；
\item
\label{algebra-item-d-1}
$d(R) = 1$；
\item
\label{algebra-item-Vx}
存在 $x \in \mathfrak m$，其中 $x$ 非幂零，并使得
$V(x) = \{\mathfrak m\}$；
\item
\label{algebra-item-x}
存在 $x \in \mathfrak m$，其中 $x$ 非幂零，并使得
$\mathfrak m = \sqrt{(x)}$；
\item
\label{algebra-item-ideal-1}
存在由 $1$ 个元素生成的定义理想，且不存在由 $0$ 个元素生成的定义理想。
\end{enumerate}
\end{lemma}

\begin{proof}
首先，假设 $\dim(R) = 1$。设 $\mathfrak p_i$ 为 $R$ 的极小素理想。
由于维数为 $1$，$R$ 唯一的其他素理想是 $\mathfrak m$。
按照引理 \ref{algebra-lemma-Noetherian-irreducible-components}，
这样的极小素理想只有有限多个。因此，可以找到
$x \in \mathfrak m$、$x \not \in \mathfrak p_i$；
参见引理 \ref{algebra-lemma-silly}。所以唯一包含 $x$ 的素理想是
$\mathfrak m$，从而得到 (\ref{algebra-item-Vx})。

\medskip\noindent
若 (\ref{algebra-item-Vx}) 成立，则由引理
\ref{algebra-lemma-Zariski-topology}，有 $\mathfrak m = \sqrt{(x)}$，
从而得到 (\ref{algebra-item-x})。反命题也显然。
(\ref{algebra-item-x}) 与 (\ref{algebra-item-ideal-1}) 的等价性直接由定义得出。

\medskip\noindent
假设 (\ref{algebra-item-ideal-1}) 成立。设 $I = (x)$ 为定义理想。
注意，乘以 $x^n$ 给出从 $R/I$ 到 $I^n/I^{n + 1}$ 的满射，
所以 $\text{length}_R(I^n/I^{n + 1})$ 有界。
因此 $d(R) = 0$ 或 $d(R) = 1$；但由 $0$ 不是定义理想这一假设，
排除 $d(R) = 0$。

\medskip\noindent
假设 (\ref{algebra-item-d-1}) 成立。为导出矛盾，假设存在素理想
$\mathfrak p \subset \mathfrak q \subset \mathfrak m$，
且两个包含都严格。选取某个定义理想 $I \subset R$。
将反复使用引理 \ref{algebra-lemma-hilbert-ses-chi}。首先，将它应用于正合列
$0 \to \mathfrak p \to R \to R/\mathfrak p \to 0$，
可得 $d(R/\mathfrak p) \leq 1$。但它显然不可能为零。
选取 $x\in \mathfrak q$、$x\not \in \mathfrak p$。
考虑短正合列
$$
0 \to R/\mathfrak p \to R/\mathfrak p \to R/(xR + \mathfrak p) \to 0.
$$
这蕴含
$\chi_{I, R/\mathfrak p} - \chi_{I, R/\mathfrak p}
- \chi_{I, R/(xR + \mathfrak p)} = - \chi_{I, R/(xR + \mathfrak p)}$
的次数为 $ < 1$。换言之，
$d(R/(xR + \mathfrak p)) = 0$，因而由引理
\ref{algebra-lemma-dimension-0-d-0}，
$\dim(R/(xR + \mathfrak p)) = 0$。但
$R/(xR + \mathfrak p)$ 有两个不同的素理想
$\mathfrak q/(xR + \mathfrak p)$ 与
$\mathfrak m/(xR + \mathfrak p)$，这给出所需矛盾。
\end{proof}

\begin{proposition}
\label{algebra-proposition-dimension}
设 $R$ 为诺特局部环，$d \geq 0$ 为整数。下列条件等价：
\begin{enumerate}
\item
\label{algebra-item-dim-d}
$\dim(R) = d$；
\item
\label{algebra-item-d-d}
$d(R) = d$；
\item
\label{algebra-item-ideal-d}
存在由 $d$ 个元素生成的定义理想，且不存在由少于 $d$ 个元素生成的
定义理想。
\end{enumerate}
\end{proposition}

\begin{proof}
这个证明实际上与引理 \ref{algebra-lemma-height-1} 的证明相同。
对 $d$ 归纳证明该命题。由引理
\ref{algebra-lemma-dimension-0-d-0} 和 \ref{algebra-lemma-height-1}，
可以假设 $d > 1$。用 $d'(R)$ 表示 $R$ 的定义理想的最小生成元数。
将证明不等式
$\dim(R) \geq d'(R) \geq d(R) \geq \dim(R)$，
从而它们全都相等。

\medskip\noindent
首先，假设 $\dim(R) = d$。设 $\mathfrak p_i$ 为 $R$ 的极小素理想。
按照引理 \ref{algebra-lemma-Noetherian-irreducible-components}，
它们只有有限多个。因此，可以找到
$x \in \mathfrak m$、$x \not \in \mathfrak p_i$；
参见引理 \ref{algebra-lemma-silly}。注意，每条极大素理想链都从某个
$\mathfrak p_i$ 开始，所以 $R/xR$ 的维数至多为 $d-1$。
由归纳假设，存在 $x_2, \ldots, x_d$，使它们在 $R/xR$ 中
生成一个定义理想。因此，$R$ 有一个由（至多）$d$ 个元素生成的定义理想。

\medskip\noindent
假设 $d'(R) = d$。令 $I = (x_1, \ldots, x_d)$ 为定义理想。
注意，乘以 $x_1, \ldots, x_d$ 中所有次数为 $n$ 的单项式，
给出从 $\binom{d + n - 1}{d - 1}$ 个 $R/I$ 之直和到
$I^n/I^{n + 1}$ 的满射。所以
$\text{length}_R(I^n/I^{n + 1})$ 受一个次数为 $d-1$ 的多项式控制。
因此 $d(R) \leq d$。

\medskip\noindent
假设 $d(R) = d$。考虑素理想链
$\mathfrak p \subset \mathfrak q \subset
\mathfrak q_2 \subset \ldots \subset \mathfrak q_e = \mathfrak m$，
其中所有包含都严格，且 $e \geq 2$。选取某个定义理想
$I \subset R$。将反复使用引理 \ref{algebra-lemma-hilbert-ses-chi}。
首先，把它应用于正合列
$0 \to \mathfrak p \to R \to R/\mathfrak p \to 0$，
可得 $d(R/\mathfrak p) \leq d$；但它显然不可能为零。
选取 $x\in \mathfrak q$、$x\not \in \mathfrak p$。
考虑短正合列
$$
0 \to R/\mathfrak p \to R/\mathfrak p \to R/(xR + \mathfrak p) \to 0.
$$
这蕴含
$\chi_{I, R/\mathfrak p} - \chi_{I, R/\mathfrak p}
- \chi_{I, R/(xR + \mathfrak p)} = - \chi_{I, R/(xR + \mathfrak p)}$
的次数为 $ < d$。换言之，
$d(R/(xR + \mathfrak p)) \leq d - 1$，所以由归纳假设，
$\dim(R/(xR + \mathfrak p)) \leq d - 1$。现在，
$R/(xR + \mathfrak p)$ 有素理想链
$\mathfrak q/(xR + \mathfrak p) \subset \mathfrak q_2/(xR + \mathfrak p)
\subset \ldots \subset \mathfrak q_e/(xR + \mathfrak p)$，
由此 $e - 1 \leq d - 1$。由于开始时的素理想链是任意的，
这证明 $\dim(R) \leq d(R)$。

\medskip\noindent
回顾上述论证可见，我们已如所需地证明了这组循环不等式。
\end{proof}

\noindent
设 $(R, \mathfrak m)$ 为诺特局部环。由上述结果显然可知，
$\mathfrak m$ 不可能由少于 $\dim(R)$ 个变量生成。
由 Nakayama 引理 \ref{algebra-lemma-NAK}，$\mathfrak m$ 的最小生成元数等于
$\dim_{\kappa(\mathfrak m)} \mathfrak m/\mathfrak m^2$。
所以有如下基本不等式：
$$
\dim(R) \leq \dim_{\kappa(\mathfrak m)} \mathfrak m/\mathfrak m^2.
$$
事实证明，等号成立的环具有许多良好性质。它们称为正则局部环。

\begin{definition}
\label{algebra-definition-regular-local}
设 $(R, \mathfrak m)$ 为维数为 $d$ 的诺特局部环。
\begin{enumerate}
\item {\it $R$ 的一组参数系}，是生成 $R$ 的一个定义理想的元素序列
$x_1, \ldots, x_d \in \mathfrak m$；
\item 若存在 $x_1, \ldots, x_d \in \mathfrak m$，使得
$\mathfrak m = (x_1, \ldots, x_d)$，则称 $R$ 为
{\it 正则局部环}，并称 $x_1, \ldots, x_d$ 为一组
{\it 正则参数系}。
\end{enumerate}
\end{definition}

\noindent
下列引理由上面诸引理与命题的证明显然可得；这里明确写出，以便引用。

\begin{lemma}
\label{algebra-lemma-minimal-over-1}
设 $R$ 为诺特环，$x \in R$。
\begin{enumerate}
\item 若 $\mathfrak p$ 是 $(x)$ 上的极小素理想，则
$\mathfrak p$ 的高度为 $0$ 或 $1$。
\item 若 $\mathfrak p, \mathfrak q \in \Spec(R)$，且
$\mathfrak q$ 是 $(\mathfrak p, x)$ 上的极小素理想，则
$\mathfrak p$ 与 $\mathfrak q$ 之间不存在素理想。
\end{enumerate}
\end{lemma}

\begin{proof}
证明 (1)。若 $\mathfrak p$ 是 $x$ 上的极小素理想，则
$R_\mathfrak p$ 中唯一包含 $x$ 的素理想是极大理想
$\mathfrak p R_\mathfrak p$。这是因为 $R_\mathfrak p$ 的素理想
与 $R$ 中包含于 $\mathfrak p$ 的素理想作 $1$-对-$1$ 对应；参见引理
\ref{algebra-lemma-spec-localization}。因此，若 $x$ 在
$R_\mathfrak p$ 中不幂零，则引理 \ref{algebra-lemma-height-1} 表明
$\dim(R_\mathfrak p) = 1$。当然，若 $x$ 在
$R_\mathfrak p$ 中幂零，则论证表明
$\mathfrak pR_\mathfrak p$ 是唯一的素理想，因而高度为 $0$。

\medskip\noindent
证明 (2)。由第 (1) 部分可知，$\mathfrak q/\mathfrak p$
在 $R/\mathfrak p$ 中是高度为 $1$ 或 $0$ 的素理想。
这立即蕴含 $\mathfrak p$ 与 $\mathfrak q$ 之间不可能有素理想。
\end{proof}

\begin{lemma}
\label{algebra-lemma-minimal-over-r}
设 $R$ 为诺特环，$f_1, \ldots, f_r \in R$。
\begin{enumerate}
\item 若 $\mathfrak p$ 是 $(f_1, \ldots, f_r)$ 上的极小素理想，
则 $\mathfrak p$ 的高度为 $\leq r$。
\item 若 $\mathfrak p, \mathfrak q \in \Spec(R)$，且
$\mathfrak q$ 是 $(\mathfrak p, f_1, \ldots, f_r)$ 上的极小素理想，
则 $\mathfrak p$ 与 $\mathfrak q$ 之间的每条素理想链长度至多为 $r$。
\end{enumerate}
\end{lemma}

\begin{proof}
证明 (1)。若 $\mathfrak p$ 是 $f_1, \ldots, f_r$ 上的极小素理想，
则 $R_\mathfrak p$ 中唯一包含 $f_1, \ldots, f_r$ 的素理想是
极大理想 $\mathfrak p R_\mathfrak p$。这是因为
$R_\mathfrak p$ 的素理想与 $R$ 中包含于 $\mathfrak p$ 的素理想
作 $1$-对-$1$ 对应；参见引理 \ref{algebra-lemma-spec-localization}。
所以命题 \ref{algebra-proposition-dimension} 表明
$\dim(R_\mathfrak p) \leq r$。

\medskip\noindent
证明 (2)。由第 (1) 部分可知，$\mathfrak q/\mathfrak p$
是高度为 $\leq r$ 的素理想。这立即蕴含关于
$\mathfrak p$ 与 $\mathfrak q$ 之间素理想链的断言。
\end{proof}

\begin{lemma}
\label{algebra-lemma-one-equation}
假设 $R$ 为诺特局部环，$x\in \mathfrak m$ 是其极大理想中的元素。
则 $\dim R \leq \dim R/xR + 1$。若 $x$ 不属于 $R$ 的任何极小素理想，
则等号成立。（例如，$x$ 是非零因子时。）
\end{lemma}

\begin{proof}
若环中的下列元素：
$$
x_1, \ldots, x_{\dim R/xR} \in R
$$
在 $R/xR$ 中的像生成
$R/xR$ 的一个定义理想，则 $x, x_1, \ldots,
x_{\dim R/xR}$ 生成 $R$ 的一个定义理想。所以不等式由命题
\ref{algebra-proposition-dimension} 得出。另一方面，若 $x$ 不属于 $R$
的任何极小素理想，则 $R/xR$ 中的素理想链在 $R$ 中都产生一条
距离成为极大链至少还差一步的链。
\end{proof}

\begin{lemma}
\label{algebra-lemma-elements-generate-ideal-definition}
设 $(R, \mathfrak m)$ 为诺特局部环。假设
$x_1, \ldots, x_d \in \mathfrak m$ 生成一个定义理想，且
$d = \dim(R)$。则对所有 $i = 1, \ldots, d$，有
$\dim(R/(x_1, \ldots, x_i)) = d - i$。
\end{lemma}

\begin{proof}
或者直接由命题 \ref{algebra-proposition-dimension} 的证明得出，
或者对 $d$ 归纳并使用引理 \ref{algebra-lemma-one-equation}。
\end{proof}







\section{维数理论的应用}
\label{algebra-section-applications-dimension-theory}

\noindent
可以利用维数方面的结果证明某些环有无穷谱，并构造更多雅可布森环。

\begin{lemma}
\label{algebra-lemma-Noetherian-local-domain-dim-2-infinite-opens}
设 $R$ 为维数 $\geq 2$ 的诺特局部整环。
则任意非空开子集 $U \subset \Spec(R)$ 都是无限的。
\end{lemma}

\begin{proof}
为导出矛盾，假设 $U \subset \Spec(R)$ 有限。
此时 $(0) \in U$，且 $\{(0)\}$ 是 $U$ 的开子集
（因为 $\{(0)\}$ 的补集是其他各点之闭包的并）。
因此可以假设 $U = \{(0)\}$。设 $\mathfrak m \subset R$ 为极大理想。
可以找到 $x \in \mathfrak m$、$x \not = 0$，使得
$V(x) \cup U = \Spec(R)$。换言之，$D(x) = \{(0)\}$。
特别地，由引理 \ref{algebra-lemma-one-equation} 可知
$\dim(R/xR) = \dim(R) - 1 \geq 1$。
由命题 \ref{algebra-proposition-dimension}，选取
$\overline{y}_2, \ldots, \overline{y}_{\dim(R)} \in R/xR$，
使它们生成 $R/xR$ 的一个定义理想。选取提升
$y_2, \ldots, y_{\dim(R)} \in R$，从而
$x, y_2, \ldots, y_{\dim(R)}$ 在 $R$ 中生成一个定义理想。
由引理 \ref{algebra-lemma-elements-generate-ideal-definition}，这蕴含
$\dim(R/(y_2)) = \dim(R) - 1$ 且
$\dim(R/(y_2, x)) = \dim(R) - 2$。因此存在一个包含 $y_2$
但不包含 $x$ 的素理想 $\mathfrak p$。这与
$D(x) = \{(0)\}$ 矛盾。
\end{proof}

\noindent
若 $k$ 为域，则环 $k[[t]]$，或者 $p$-进数环，都是维数为 $1$、
且只有有限多个素理想的诺特环。这是可能出现这种现象的最大维数。

\begin{lemma}
\label{algebra-lemma-Noetherian-finite-nr-primes}
只有有限多个素理想的诺特环，其维数为 $\leq 1$。
\end{lemma}

\begin{proof}
设 $R$ 为只有有限多个素理想的诺特环。若 $R$ 是局部整环，
则结论由引理
\ref{algebra-lemma-Noetherian-local-domain-dim-2-infinite-opens} 得出。
若 $R$ 是整环，则由局部情形，对所有极大理想 $\mathfrak m$，
$R_\mathfrak m$ 的维数为 $\leq 1$。所以由引理
\ref{algebra-lemma-dimension-height}，$\dim(R) \leq 1$。
对一般的 $R$，则对 $R$ 的每个极小素理想 $\mathfrak q$，
都有 $\dim(R/\mathfrak q) \leq 1$。由于每个素理想都包含一个极小素理想
（引理 \ref{algebra-lemma-Zariski-topology}），这蕴含
$\dim(R) \leq 1$。
\end{proof}

\begin{lemma}
\label{algebra-lemma-finite-type-algebra-finite-nr-primes}
设 $S$ 是域 $k$ 上的非零有限型代数。下列条件等价：
\begin{enumerate}
\item $\dim(S) = 0$；
\item $S$ 只有有限多个素理想；
\item $S$ 只有有限多个极大理想；
\item $\Spec(S)$ 满足引理
\ref{algebra-lemma-ring-with-only-minimal-primes} 中任一等价条件；
\item $\dim_k(S) < \infty$；
\item $S$ 是阿廷环；
\item $\Spec(S)$ 是离散拓扑空间；
\item 在此添加更多条件。
\end{enumerate}
\end{lemma}

\begin{proof}
查看引理 \ref{algebra-lemma-ring-with-only-minimal-primes} 的第 (5) 部分，
由定义立即可知 (1) 等价于 (4)。回忆，$\Spec(S)$ 是清醒、诺特且
雅可布森的；参见引理 \ref{algebra-lemma-spec-spectral}、
\ref{algebra-lemma-Noetherian-topology}、
\ref{algebra-lemma-finite-type-field-Jacobson} 和 \ref{algebra-lemma-jacobson}。
若 $S$ 的维数为 $0$，则每个点都定义一个不可约分支，而不可约分支
只有有限多个（拓扑学，引理 \ref{topology-lemma-Noetherian}）。
所以 (1) 蕴含 (2)。显然，(2) 蕴含 (3)。
若 (3) 成立，则由拓扑学，引理
\ref{topology-lemma-finite-jacobson}，$\Spec(S)$ 离散，
因而 $S$ 的维数为 $0$。

\medskip\noindent
至此已知 (1)--(4) 等价。蕴含 (5) $\Rightarrow$ (6)
是引理 \ref{algebra-lemma-finite-dimensional-algebra}。
蕴含 (6) $\Rightarrow$ (7) 由命题
\ref{algebra-proposition-dimension-zero-ring} 得出。
蕴含 (7) $\Rightarrow$ (4) 是显然的。
反过来，若 $S$ 满足 (1)--(4)，则 $S$ 只有有限多个素理想
$\mathfrak m_1, \ldots, \mathfrak m_r$，且它们全都极大。
注意，由希尔伯特零点定理（定理
\ref{algebra-theorem-nullstellensatz}），$\kappa(\mathfrak m_i)$ 是
$k$ 的有限扩张。由命题 \ref{algebra-proposition-dimension-zero-ring}
还可知 $S$ 是阿廷环。接着，引理
\ref{algebra-lemma-artinian-finite-length} 告诉我们
$\text{length}_S(S) < \infty$。所以由引理
\ref{algebra-lemma-pushdown-module}，$\dim_k(S) < \infty$。
由此得出 (1)--(7) 等价。（注意，证明
$\dim(S) = 0 \Rightarrow \dim_k(S) < \infty$ 的另一种更标准的方法
是使用诺特正规化，但目前尚未建立它。）
\end{proof}

\begin{lemma}
\label{algebra-lemma-noetherian-dim-1-Jacobson}
诺特雅可布森环。
\begin{enumerate}
\item 任意维数为 $1$ 且有无穷多个素理想的诺特整环 $R$
都是雅可布森环。
\item 任意满足下述条件的诺特环都是雅可布森环：
每个素理想 $\mathfrak p$ 或者极大，或者包含于无穷多个素理想中。
\end{enumerate}
\end{lemma}

\begin{proof}
第 (1) 部分是引理 \ref{algebra-lemma-pid-jacobson} 的重新表述。

\medskip\noindent
设 $R$ 是如下诺特环：每个非极大素理想 $\mathfrak p$
都包含于无穷多个素理想中。为导出矛盾，假设
$\Spec(R)$ 不是雅可布森空间。由引理
\ref{algebra-lemma-irreducible} 和 \ref{algebra-lemma-Noetherian-topology}，
可知 $\Spec(R)$ 是清醒诺特拓扑空间。由拓扑学，引理
\ref{topology-lemma-non-jacobson-Noetherian-characterize}，
存在非极大理想 $\mathfrak p \subset R$，使得
$\{\mathfrak p\}$ 是 $\Spec(R)$ 的局部闭子集。
换言之，$\mathfrak p$ 非极大，且 $\{\mathfrak p\}$
是 $V(\mathfrak p)$ 的开子集。考虑满足
$\mathfrak p \subset \mathfrak q$ 的素理想
$\mathfrak q \subset R$。回忆，
$(R/\mathfrak p)_{\mathfrak q}
= R_{\mathfrak q}/\mathfrak pR_{\mathfrak q}$ 的谱上的拓扑
由 $\Spec(R)$ 的拓扑诱导；参见引理
\ref{algebra-lemma-spec-localization} 和 \ref{algebra-lemma-spec-closed}。
所以 $\{(0)\}$ 是
$\Spec((R/\mathfrak p)_{\mathfrak q})$ 的局部闭子集。
由引理 \ref{algebra-lemma-Noetherian-local-domain-dim-2-infinite-opens}，
可知 $\dim((R/\mathfrak p)_{\mathfrak q}) = 1$。
由于这对每个 $\mathfrak q \supset \mathfrak p$ 都成立，
可得 $\dim(R/\mathfrak p) = 1$。此时利用
$\mathfrak p$ 包含于无穷多个素理想中的假设，可知
$\Spec(R/\mathfrak p)$ 是无限的。因此，由本引理第 (1) 部分，
$V(\mathfrak p) \cong \Spec(R/\mathfrak p)$ 是其闭点的闭包。
这就是所需矛盾，因为它意味着
$\{\mathfrak p\} \subset V(\mathfrak p)$ 不可能是开集。
\end{proof}








\section{模的支集与维数}
\label{algebra-section-support}

\noindent
下面给出关于模的支集与维数的一些基本结果。

\begin{lemma}
\label{algebra-lemma-filter-Noetherian-module}
设 $R$ 为诺特环，$M$ 为有限 $R$-模。则存在一条由 $R$-子模组成的滤过
$$
0 = M_0 \subset M_1 \subset \ldots \subset M_n = M
$$
使得每个商 $M_i/M_{i-1}$ 都同构于某个 $R/\mathfrak p_i$，
其中 $\mathfrak p_i$ 是 $R$ 的素理想。
\end{lemma}

\begin{proof}[第一种证明]
由引理 \ref{algebra-lemma-trivial-filter-finite-module}，只须处理
$M = R/I$ 的情形，其中 $I$ 是理想。考虑如下理想 $J$ 所成的集合
$S$：本引理对模 $R/J$ 不成立；按包含关系排序。
为导出矛盾，假设 $S$ 非空。因为 $R$ 是诺特环，$S$ 有极大元 $J$。
按照 $S$ 的定义，理想 $J$ 不可能是素理想。选取
$a, b\in R$，使得 $ab \in J$，但 $a \in J$ 与 $b\in J$
都不成立。考虑滤过
$0 \subset aR/(J \cap aR) \subset R/J$。
注意，子模 $aR/(J \cap aR)$ 和商模
$(R/J)/(aR/(J \cap aR))$ 都是循环模；将它们写成
$R/J'$ 与 $R/J''$，便有短正合列
$0 \to R/J' \to R/J \to R/J'' \to 0$。
包含 $J \subset J'$ 是严格的，因为 $b \in J'$；
包含 $J \subset J''$ 也是严格的，因为 $a \in J''$。
因此，由 $J$ 的极大性，$R/J'$ 与 $R/J''$ 都有上述滤过，
从而 $R/J$ 也有。矛盾。
\end{proof}

\begin{proof}[第二种证明]
对 $R$-模 $M$，若存在本引理断言中的滤过，则称 $P(M)$ 成立。
注意，$P$ 在扩张下保持不变，且对 $0$ 成立。由引理
\ref{algebra-lemma-trivial-filter-finite-module}，只须证明对每个理想 $I$，
$P(R/I)$ 成立。否则，由于 $R$ 是诺特环，存在极大反例 $J$。
由例 \ref{algebra-example-oka-family-property-modules} 和命题
\ref{algebra-proposition-oka}，理想 $J$ 是素理想，矛盾。
\end{proof}

\begin{lemma}
\label{algebra-lemma-filter-primes-in-support}
设 $R$、$M$、$M_i$、$\mathfrak p_i$ 如引理
\ref{algebra-lemma-filter-Noetherian-module} 中所述。则
$\text{Supp}(M) = \bigcup V(\mathfrak p_i)$；特别地，
$\mathfrak p_i \in \text{Supp}(M)$。
\end{lemma}

\begin{proof}
这由引理 \ref{algebra-lemma-support-closed} 和
\ref{algebra-lemma-support-quotient} 得出。
\end{proof}

\begin{lemma}
\label{algebra-lemma-support-point}
假设 $R$ 是极大理想为 $\mathfrak m$ 的诺特局部环。
设 $M$ 是非零有限 $R$-模。则
$\text{Supp}(M) = \{ \mathfrak m\}$，当且仅当
$M$ 在 $R$ 上具有有限长度。
\end{lemma}

\begin{proof}
假设 $\text{Supp}(M) = \{ \mathfrak m\}$。
只须证明引理 \ref{algebra-lemma-filter-Noetherian-module} 的滤过中所有素理想
$\mathfrak p_i$ 都是极大理想。这由引理
\ref{algebra-lemma-filter-primes-in-support} 显然成立。

\medskip\noindent
假设 $M$ 在 $R$ 上具有有限长度。由引理
\ref{algebra-lemma-length-infinite}，$\mathfrak m^n M = 0$。
对 $R$ 中的素理想，若 $\mathfrak p \not = \mathfrak m$，
则 $\mathfrak m$ 的某个元素在 $R_{\mathfrak p}$
中映为单位，所以 $M_{\mathfrak p} = 0$。
\end{proof}

\begin{lemma}
\label{algebra-lemma-Noetherian-power-ideal-kills-module}
设 $R$ 为诺特环，$I \subset R$ 为理想，$M$ 为有限 $R$-模。
则对某个 $n \geq 0$ 有 $I^nM = 0$，当且仅当
$\text{Supp}(M) \subset V(I)$。
\end{lemma}

\begin{proof}
事实上，$I^nM = 0$ 等价于 $I^n \subset \text{Ann}(M)$。
由于 $R$ 是诺特环，由引理 \ref{algebra-lemma-Noetherian-power}，
这等价于 $I \subset \sqrt{\text{Ann}(M)}$。
再由引理 \ref{algebra-lemma-Zariski-topology}，这等价于
$V(I) \supset V(\text{Ann}(M))$。
由引理 \ref{algebra-lemma-support-closed}，后者等价于
$V(I) \supset \text{Supp}(M)$。
\end{proof}

\begin{lemma}
\label{algebra-lemma-filter-minimal-primes-in-support}
设 $R$、$M$、$M_i$、$\mathfrak p_i$ 如引理
\ref{algebra-lemma-filter-Noetherian-module} 中所述。集合
$\{\mathfrak p_i\}$ 的极小元正是 $\text{Supp}(M)$ 的极小元。
某个极小素理想 $\mathfrak p$ 出现的次数为
$$
\#\{i \mid \mathfrak p_i = \mathfrak p\}
=
\text{length}_{R_\mathfrak p} M_{\mathfrak p}.
$$
\end{lemma}

\begin{proof}
第一条由
$\text{Supp}(M) = \bigcup V(\mathfrak p_i)$ 得出；参见引理
\ref{algebra-lemma-filter-primes-in-support}。设
$\mathfrak p \in \text{Supp}(M)$ 极小。$M_{\mathfrak p}$ 的支集
是只含极大理想 $\mathfrak p R_{\mathfrak p}$ 的集合。
所以由引理 \ref{algebra-lemma-support-point}，$M_{\mathfrak p}$ 的长度有限且
$> 0$。接着注意，$M_{\mathfrak p}$ 有一条滤过，其子商为
$
(R/\mathfrak p_i)_{\mathfrak p}
=
R_{\mathfrak p}/{\mathfrak p_i}R_{\mathfrak p}
$。
当 $\mathfrak p_i \not \subset \mathfrak p$ 时，它们为零；
当 $\mathfrak p_i \subset \mathfrak p$ 时，它们等于
$\kappa(\mathfrak p)$，因为由 $\mathfrak p$ 的极小性，
此时 $\mathfrak p_i = \mathfrak p$。结论由
$\kappa(\mathfrak p)$ 的长度为 $1$ 得出。
\end{proof}

\begin{lemma}
\label{algebra-lemma-support-dimension-d}
设 $R$ 为诺特局部环，$M$ 为有限 $R$-模。则
$d(M) = \dim(\text{Supp}(M))$，其中 $d(M)$ 如定义
\ref{algebra-definition-d} 中所述。
\end{lemma}

\begin{proof}
设 $M_i, \mathfrak p_i$ 如引理
\ref{algebra-lemma-filter-Noetherian-module} 中所述。由引理
\ref{algebra-lemma-hilbert-ses-chi} 得等式
$d(M) = \max \{ d(R/\mathfrak p_i) \}$。
由命题 \ref{algebra-proposition-dimension}，有
$d(R/\mathfrak p_i) = \dim(R/\mathfrak p_i)$。
显然 $\dim(R/\mathfrak p_i) = \dim V(\mathfrak p_i)$。
由于 $\text{Supp}(M)$ 的所有极小素理想都出现在诸
$\mathfrak p_i$ 中（引理
\ref{algebra-lemma-filter-minimal-primes-in-support}），结论成立。
\end{proof}

\begin{lemma}
\label{algebra-lemma-ses-dimension}
设 $R$ 为诺特环，并设
$0 \to M' \to M \to M'' \to 0$ 是有限 $R$-模的短正合列。则
$\max\{\dim(\text{Supp}(M')), \dim(\text{Supp}(M''))\} =
\dim(\text{Supp}(M))$。
\end{lemma}

\begin{proof}
若 $R$ 是局部环，则结论立即由引理
\ref{algebra-lemma-support-dimension-d} 和 \ref{algebra-lemma-hilbert-ses-chi} 得出。
另一个更初等、且在 $R$ 非局部时也适用的论证，是利用
$\text{Supp}(M')$、$\text{Supp}(M'')$ 与 $\text{Supp}(M)$ 都是闭集
（引理 \ref{algebra-lemma-support-closed}），以及
$\text{Supp}(M) = \text{Supp}(M') \cup \text{Supp}(M'')$
（引理 \ref{algebra-lemma-support-quotient}）。
\end{proof}








\section{相伴素理想}
\label{algebra-section-ass}

\noindent
下面是标准定义。对于非诺特环和非有限模，采用
第 \ref{algebra-section-weakly-ass} 节中的定义可能更为合适。

\begin{definition}
\label{algebra-definition-associated}
设 $R$ 为环，$M$ 为 $R$-模。如果存在元素 $m \in M$，
其零化子为 $\mathfrak p$，则称 $R$ 的素理想 $\mathfrak p$
与 $M$ {\it 相伴}。所有这样的素理想组成的集合记作
$\text{Ass}_R(M)$ 或 $\text{Ass}(M)$。
\end{definition}

\begin{lemma}
\label{algebra-lemma-ass-support}
设 $R$ 为环，$M$ 为 $R$-模。则
$\text{Ass}(M) \subset \text{Supp}(M)$。
\end{lemma}

\begin{proof}
如果 $m \in M$ 的零化子为 $\mathfrak p$，那么特别地，
$R \setminus \mathfrak p$ 中没有元素零化 $m$。
因此 $m$ 是 $M_{\mathfrak p}$ 的非零元素，即
$\mathfrak p \in \text{Supp}(M)$。
\end{proof}

\begin{lemma}
\label{algebra-lemma-ass}
设 $R$ 为环，且 $0 \to M' \to M \to M'' \to 0$ 是
$R$-模的短正合列。则 $\text{Ass}(M') \subset \text{Ass}(M)$，且
$\text{Ass}(M) \subset \text{Ass}(M') \cup \text{Ass}(M'')$。
此外，$\text{Ass}(M' \oplus M'') = \text{Ass}(M') \cup \text{Ass}(M'')$。
\end{lemma}

\begin{proof}
若 $m' \in M'$，则 $m'$ 作为 $M'$ 的元素时的零化子，
与 $m'$ 作为 $M$ 的元素时的零化子相同。因此有包含
$\text{Ass}(M') \subset \text{Ass}(M)$。设 $m \in M$ 的零化子
是素理想 $\mathfrak p$。如果存在 $g \in R$，
$g \not \in \mathfrak p$，使得 $m' = gm \in M'$，
那么 $m'$ 的零化子就是 $\mathfrak p$。如果不存在
$g \in R$，$g \not \in \mathfrak p$，使得 $gm \in M'$，
那么像 $m'' \in M''$（即 $m$ 的像）的零化子就是
$\mathfrak p$。这就证明了包含
$\text{Ass}(M) \subset \text{Ass}(M') \cup \text{Ass}(M'')$。
最后一个断言的证明从略。
\end{proof}

\begin{lemma}
\label{algebra-lemma-ass-filter}
设 $R$ 为环，$M$ 为 $R$-模。假设存在一条由 $R$-子模组成的滤过
$$
0 = M_0 \subset M_1 \subset \ldots \subset M_n = M
$$
使得每个商 $M_i/M_{i-1}$ 都同构于 $R/\mathfrak p_i$，
其中 $\mathfrak p_i$ 是 $R$ 的某个素理想。则
$\text{Ass}(M) \subset \{\mathfrak p_1, \ldots, \mathfrak p_n\}$。
\end{lemma}

\begin{proof}
对滤过 $\{ M_i \}$ 的长度 $n$ 作归纳。选取 $m \in M$，
使其零化子为素理想 $\mathfrak p$。如果 $m \in M_{n-1}$，
则由归纳假设即得结论。否则，$m$ 在
$M/M_{n-1} \cong R/\mathfrak p_n$ 中映为非零元素。因此
$\mathfrak p \subset \mathfrak p_n$。如果二者不相等，则可找到
$f \in \mathfrak p_n$，$f \not\in \mathfrak p$。
此时 $fm$ 的零化子仍为 $\mathfrak p$，且 $fm \in M_{n-1}$。
于是归纳假设给出结论。
\end{proof}

\begin{lemma}
\label{algebra-lemma-finite-ass}
设 $R$ 为诺特环，$M$ 为有限 $R$-模。则
$\text{Ass}(M)$ 是有限集。
\end{lemma}

\begin{proof}
这立即由引理 \ref{algebra-lemma-ass-filter} 和
\ref{algebra-lemma-filter-Noetherian-module} 得出。
\end{proof}

\begin{proposition}
\label{algebra-proposition-minimal-primes-associated-primes}
设 $R$ 为诺特环，$M$ 为有限 $R$-模。下列素理想集合相同：
\begin{enumerate}
\item $M$ 的支集中的极小素理想；
\item $\text{Ass}(M)$ 中的极小素理想；
\item 对任意满足 $M_i/M_{i-1} \cong R/\mathfrak p_i$ 的滤过
$0 = M_0 \subset M_1 \subset \ldots
\subset M_{n-1} \subset M_n = M$，集合
$\{\mathfrak p_i\}$ 中的极小素理想。
\end{enumerate}
\end{proposition}

\begin{proof}
选取 (3) 中的滤过。在引理
\ref{algebra-lemma-filter-minimal-primes-in-support} 中，我们已经看到
(1) 与 (3) 中的集合相同。

\medskip\noindent
设 $\mathfrak p$ 是集合 $\{\mathfrak p_i\}$ 的极小元，
并取最小的 $i$ 使得 $\mathfrak p = \mathfrak p_i$。
选取 $m \in M_i$，$m \not \in M_{i-1}$。元素 $m$ 的零化子
包含于 $\mathfrak p_i = \mathfrak p$，并包含
$\mathfrak p_1 \mathfrak p_2 \ldots \mathfrak p_i$。
由 $i$ 与 $\mathfrak p$ 的选取，对 $j < i$ 有
$\mathfrak p_j \not \subset \mathfrak p$，因而
$\mathfrak p_1 \mathfrak p_2 \ldots \mathfrak p_{i - 1}
\not \subset \mathfrak p_i$。选取
$f \in \mathfrak p_1 \mathfrak p_2 \ldots \mathfrak p_{i - 1}$，
$f \not \in \mathfrak p$。则 $fm$ 的零化子为 $\mathfrak p$。
由此可见，$\mathfrak p$ 是 $M$ 的相伴素理想。
由引理 \ref{algebra-lemma-ass-support} 有
$\text{Ass}(M) \subset \text{Supp}(M)$，所以 $\mathfrak p$
在 $\text{Ass}(M)$ 中极小。因此，(1) 中的素理想集合
包含于 (2) 中的素理想集合。

\medskip\noindent
设 $\mathfrak p$ 是 $\text{Ass}(M)$ 的极小元。
由于 $\text{Ass}(M) \subset \text{Supp}(M)$，存在
$\text{Supp}(M)$ 的极小元 $\mathfrak q$，使得
$\mathfrak q \subset \mathfrak p$。我们刚刚已经证明
$\mathfrak q \in \text{Ass}(M)$。由 $\mathfrak p$ 的极小性，
$\mathfrak q = \mathfrak p$。因此，(2) 中的素理想集合
包含于 (1) 中的素理想集合。
\end{proof}

\begin{lemma}
\label{algebra-lemma-ass-zero}
\begin{slogan}
诺特环上的每个非零模都有相伴素理想。
\end{slogan}
设 $R$ 为诺特环，$M$ 为 $R$-模。则
$$
M = (0) \Leftrightarrow \text{Ass}(M) = \emptyset.
$$
\end{lemma}

\begin{proof}
如果 $M = (0)$，那么依定义
$\text{Ass}(M) = \emptyset$。如果 $M \not = 0$，选取任意非零有限生成
子模 $M' \subset M$，例如由单个非零元素生成的子模。由引理
\ref{algebra-lemma-support-zero} 可知 $\text{Supp}(M')$ 非空。
由命题 \ref{algebra-proposition-minimal-primes-associated-primes}，
这意味着 $\text{Ass}(M')$ 非空。再由引理
\ref{algebra-lemma-ass}，可知 $\text{Ass}(M) \not = \emptyset$。
\end{proof}

\begin{lemma}
\label{algebra-lemma-ass-minimal-prime-support}
设 $R$ 为诺特环，$M$ 为 $R$-模。在 $\text{Supp}(M)$
的元素中极小的任意 $\mathfrak p \in \text{Supp}(M)$ 都属于
$\text{Ass}(M)$。
\end{lemma}

\begin{proof}
如果 $M$ 是有限 $R$-模，这就是命题
\ref{algebra-proposition-minimal-primes-associated-primes} 的推论。
一般地，将 $M = \bigcup M_\lambda$ 写成其有限子模之并，并使用
$\text{Supp}(M) = \bigcup \text{Supp}(M_\lambda)$
以及
$\text{Ass}(M) = \bigcup \text{Ass}(M_\lambda)$。
\end{proof}

\begin{lemma}
\label{algebra-lemma-ass-zero-divisors}
设 $R$ 为诺特环，$M$ 为 $R$-模。并集
$\bigcup_{\mathfrak q \in \text{Ass}(M)} \mathfrak q$
正是 $R$ 中在 $M$ 上为零因子的元素所成的集合。
\end{lemma}

\begin{proof}
任意相伴素理想中的每个元素显然都是 $M$ 上的零因子。
反之，假设 $x \in R$ 是 $M$ 上的零因子。考虑子模
$N = \{m \in M \mid xm = 0\}$。由于 $N$ 非零，由引理
\ref{algebra-lemma-ass-zero}，它有相伴素理想 $\mathfrak q$。
于是 $x \in \mathfrak q$；又由引理
\ref{algebra-lemma-ass}，$\mathfrak q$ 是 $M$ 的相伴素理想。
\end{proof}

\begin{lemma}
\label{algebra-lemma-one-equation-module}
设 $R$ 为诺特局部环，$M$ 为有限 $R$-模，并设
$f \in \mathfrak m$ 是 $R$ 的极大理想中的元素。则
$$
\dim(\text{Supp}(M/fM)) \leq
\dim(\text{Supp}(M)) \leq
\dim(\text{Supp}(M/fM)) + 1
$$
如果 $f$ 不属于 $M$ 的支集中的任何极小素理想
（例如，如果 $f$ 是 $M$ 上的非零因子），
那么右侧不等式取等号。
\end{lemma}

\begin{proof}
（括号中的断言由引理 \ref{algebra-lemma-ass-zero-divisors} 得出。）
第一个不等式由
$\text{Supp}(M/fM) \subset \text{Supp}(M)$ 得出，
参见引理 \ref{algebra-lemma-support-quotient}。对于第二个不等式，注意
$\text{Supp}(M/fM) = \text{Supp}(M) \cap V(f)$，参见引理
\ref{algebra-lemma-support-quotient}。例如，由引理
\ref{algebra-lemma-filter-primes-in-support} 和维数的初等性质可知，只须证明
$\dim V(\mathfrak p) \leq \dim (V(\mathfrak p) \cap V(f)) + 1$
对 $R$ 的素理想 $\mathfrak p$ 成立。这是引理
\ref{algebra-lemma-one-equation} 的推论。最后，如果 $f$ 不包含于
$M$ 的支集中的任何极小素理想，那么
$\text{Supp}(M/fM)$ 中的每条素理想链都会产生
$\text{Supp}(M)$ 中的一条素理想链，而后者距极大链至少还差一步。
\end{proof}

\begin{lemma}
\label{algebra-lemma-ass-functorial}
设 $\varphi : R \to S$ 为环映射，$M$ 为 $S$-模。则
$\Spec(\varphi)(\text{Ass}_S(M)) \subset \text{Ass}_R(M)$。
\end{lemma}

\begin{proof}
如果 $\mathfrak q \in \text{Ass}_S(M)$，那么存在 $M$ 中的
$m$，使得 $m$ 在 $S$ 中的零化子为 $\mathfrak q$。
于是 $m$ 在 $R$ 中的零化子为 $\mathfrak q \cap R$。
\end{proof}

\begin{remark}
\label{algebra-remark-ass-reverse-functorial}
设 $\varphi : R \to S$ 为环映射，$M$ 为 $S$-模。
不一定有
$\Spec(\varphi)(\text{Ass}_S(M)) \supset \text{Ass}_R(M)$。
例如，考虑环映射
$R = k \to S = k[x_1, x_2, x_3, \ldots]/(x_i^2)$ 以及 $M = S$。
此时 $\text{Ass}_R(M)$ 非空，但 $\text{Ass}_S(S)$ 为空。
\end{remark}

\begin{lemma}
\label{algebra-lemma-ass-functorial-Noetherian}
设 $\varphi : R \to S$ 为环映射，$M$ 为 $S$-模。
如果 $S$ 是诺特环，则
$\Spec(\varphi)(\text{Ass}_S(M)) = \text{Ass}_R(M)$。
\end{lemma}

\begin{proof}
我们已经在引理
\ref{algebra-lemma-ass-functorial} 中看到
$\Spec(\varphi)(\text{Ass}_S(M)) \subset \text{Ass}_R(M)$。
反过来，选取素理想 $\mathfrak p \in \text{Ass}_R(M)$。
设 $m \in M$ 满足：$m$ 在 $R$ 中的零化子为
$\mathfrak p$。令
$I = \{g \in S \mid gm = 0\}$ 为 $m$ 在 $S$ 中的零化子。
于是 $R/\mathfrak p \subset S/I$ 是单射。
结合引理 \ref{algebra-lemma-injective-minimal-primes-in-image} 和
\ref{algebra-lemma-minimal-prime-image-minimal-prime}，可知存在包含
$I$ 的极小素理想 $\mathfrak q \subset S$，它映到
$\mathfrak p$。由命题
\ref{algebra-proposition-minimal-primes-associated-primes}，
$\mathfrak q$ 是 $S/I$ 的相伴素理想；再由引理
\ref{algebra-lemma-ass}，$\mathfrak q$ 是 $M$ 的相伴素理想，
结论得证。
\end{proof}

\begin{lemma}
\label{algebra-lemma-ass-quotient-ring}
设 $R$ 为环，$I$ 为理想，$M$ 为 $R/I$-模。
通过典范单射
$\Spec(R/I) \to \Spec(R)$，有
$\text{Ass}_{R/I}(M) = \text{Ass}_R(M)$。
\end{lemma}

\begin{proof}
证明从略。
\end{proof}

\begin{lemma}
\label{algebra-lemma-associated-primes-localize}
设 $R$ 为环，$M$ 为 $R$-模，并设
$\mathfrak p \subset R$ 为素理想。
\begin{enumerate}
\item 如果 $\mathfrak p \in \text{Ass}(M)$，则
$\mathfrak pR_{\mathfrak p} \in \text{Ass}(M_{\mathfrak p})$。
\item 如果 $\mathfrak p$ 是有限生成的，则逆命题也成立。
\end{enumerate}
\end{lemma}

\begin{proof}
如果 $\mathfrak p \in \text{Ass}(M)$，则存在元素 $m \in M$，
其零化子为 $\mathfrak p$。由于局部化正合
（命题 \ref{algebra-proposition-localization-exact}），
$m/1$ 在 $M_{\mathfrak p}$ 中的零化子为
$\mathfrak pR_{\mathfrak p}$，故 (1) 成立。
假设 $\mathfrak pR_{\mathfrak p} \in \text{Ass}(M_{\mathfrak p})$，
并且 $\mathfrak p = (f_1, \ldots, f_n)$。设 $m/g$ 是
$M_{\mathfrak p}$ 中零化子为
$\mathfrak pR_{\mathfrak p}$ 的元素。这意味着 $m$ 的零化子
包含于 $\mathfrak p$。由于 $f_i m/g = 0$ 在
$M_{\mathfrak p}$ 中成立，可知存在
$g_i \in R$，$g_i \not \in \mathfrak p$，使得
$g_i f_i m = 0$ 在 $M$ 中成立。合起来可知
$g_1\ldots g_nm$ 的零化子为 $\mathfrak p$。因此
$\mathfrak p \in \text{Ass}(M)$。
\end{proof}

\begin{lemma}
\label{algebra-lemma-localize-ass}
设 $R$ 为环，$M$ 为 $R$-模，并设 $S \subset R$ 为乘法子集。
通过典范单射 $\Spec(S^{-1}R) \to \Spec(R)$，有
\begin{enumerate}
\item $\text{Ass}_R(S^{-1}M) = \text{Ass}_{S^{-1}R}(S^{-1}M)$；
\item
$\text{Ass}_R(M) \cap \Spec(S^{-1}R) \subset \text{Ass}_R(S^{-1}M)$；且
\item 如果 $R$ 是诺特环，则此包含为等式。
\end{enumerate}
\end{lemma}

\begin{proof}
对 $m \in S^{-1}M$，设 $I \subset R$ 和 $J \subset S^{-1}R$
分别为 $m$ 的零化子。则 $I$ 是 $J$ 在映射
$R \to S^{-1}R$ 下的原像，且 $J = S^{-1}I$。
(1) 中的等式由引理 \ref{algebra-lemma-spec-localization} 对映射
$\Spec(S^{-1}R) \to \Spec(R)$ 的描述得出。
对 $m \in M$，设 $I \subset R$ 为 $m$ 在 $R$ 中的零化子，
并设 $J \subset S^{-1}R$ 为
$m/1 \in S^{-1}M$ 的零化子。我们有 $J = S^{-1}I$，这蕴含 (2)。
诺特情形中的等式由引理
\ref{algebra-lemma-associated-primes-localize} 得出，因为对
$\mathfrak p \in R$，若 $S \cap \mathfrak p = \emptyset$，则
$M_{\mathfrak p} = (S^{-1}M)_{S^{-1}\mathfrak p}$。
\end{proof}

\begin{lemma}
\label{algebra-lemma-localize-ass-nonzero-divisors}
设 $R$ 为环，$M$ 为 $R$-模，并设 $S \subset R$ 为乘法子集。
假设每个 $s \in S$ 都是 $M$ 上的非零因子。则
$$
\text{Ass}_R(M) = \text{Ass}_R(S^{-1}M).
$$
\end{lemma}

\begin{proof}
由假设有 $M \subset S^{-1}M$，故引理
\ref{algebra-lemma-ass} 给出包含
$\text{Ass}(M) \subset \text{Ass}(S^{-1}M)$。
反之，假设 $n/s \in S^{-1}M$ 的零化子为素理想
$\mathfrak p$。则 $n \in M$ 的零化子也是 $\mathfrak p$。
\end{proof}

\begin{lemma}
\label{algebra-lemma-ideal-nonzerodivisor}
设 $R$ 为极大理想是 $\mathfrak m$ 的诺特局部环，
$I \subset \mathfrak m$ 为理想，$M$ 为有限 $R$-模。
下列条件等价：
\begin{enumerate}
\item 存在 $x \in I$，它不是 $M$ 上的零因子；
\item 对所有 $\mathfrak q \in \text{Ass}(M)$，均有
$I \not \subset \mathfrak q$。
\end{enumerate}
\end{lemma}

\begin{proof}
如果 $I$ 中存在非零因子 $x$，那么 $x$ 显然不可能属于
$M$ 的任何相伴素理想。反之，假设对所有
$\mathfrak q \in \text{Ass}(M)$ 均有
$I \not \subset \mathfrak q$。在这种情形下，由引理
\ref{algebra-lemma-finite-ass} 和 \ref{algebra-lemma-silly}，可以选取
$x \in I$，使得对所有 $\mathfrak q \in \text{Ass}(M)$，
$x \not \in \mathfrak q$。由引理
\ref{algebra-lemma-ass-zero-divisors}，元素 $x$ 不是 $M$ 上的零因子。
\end{proof}

\begin{lemma}
\label{algebra-lemma-zero-at-ass-zero}
设 $R$ 为环，$M$ 为 $R$-模。如果 $R$ 是诺特环，则映射
$$
M
\longrightarrow
\prod\nolimits_{\mathfrak p \in \text{Ass}(M)} M_{\mathfrak p}
$$
是单射。
\end{lemma}

\begin{proof}
设 $x \in M$ 是该映射的核中的元素。如果 $\mathfrak p$ 是
$Rx \subset M$ 的相伴素理想，那么一方面
$\mathfrak p \in \text{Ass}(M)$（引理 \ref{algebra-lemma-ass}），
另一方面
$(Rx)_{\mathfrak p} \subset M_{\mathfrak p}$ 非零。
此矛盾说明 $\text{Ass}(Rx) = \emptyset$。因此由引理
\ref{algebra-lemma-ass-zero}，$Rx = 0$。
\end{proof}

\noindent
这个引理或许应该放在别处。

\begin{lemma}
\label{algebra-lemma-dim-not-zero-exists-nonzerodivisor-nonunit}
设 $k$ 为域，$S$ 为有限型 $k$-代数。如果 $\dim(S) > 0$，
则存在既是非零因子又是非单位的元素 $f \in S$。
\end{lemma}

\begin{proof}
由引理 \ref{algebra-lemma-finite-ass}，环 $S$ 只有有限多个相伴素理想。
由引理 \ref{algebra-lemma-finite-type-algebra-finite-nr-primes}，
环 $S$ 有无穷多个极大理想。因此，可以选取极大理想
$\mathfrak m \subset S$，它不是 $S$ 的相伴素理想。
由素理想回避（引理 \ref{algebra-lemma-silly}），可以选取非零元素
$f \in \mathfrak m$，它不属于 $S$ 的任何相伴素理想。
由引理 \ref{algebra-lemma-ass-zero-divisors}，元素 $f$ 是非零因子；
又因 $f \in \mathfrak m$，可知 $f$ 不是单位。
\end{proof}



\section{符号幂}
\label{algebra-section-symbolic-power}

\noindent
定义如下。

\begin{definition}
\label{algebra-definition-symbolic-power}
设 $R$ 为环，$\mathfrak p$ 为素理想。对 $n \geq 0$，
$\mathfrak p$ 的第 $n$ 个{\it 符号幂}是理想
$\mathfrak p^{(n)} = \Ker(R \to R_\mathfrak p/\mathfrak p^nR_\mathfrak p)$。
\end{definition}

\noindent
注意 $\mathfrak p^n \subset \mathfrak p^{(n)}$，但等式并非总成立。

\begin{lemma}
\label{algebra-lemma-symbolic-power-associated}
设 $R$ 为诺特环，$\mathfrak p$ 为素理想，且 $n > 0$。则
$\text{Ass}(R/\mathfrak p^{(n)}) = \{\mathfrak p\}$。
\end{lemma}

\begin{proof}
如果 $\mathfrak q$ 是 $R/\mathfrak p^{(n)}$ 的相伴素理想，
那么显然 $\mathfrak p \subset \mathfrak q$。另一方面，任意
$x \in R$，$x \not \in \mathfrak p$，都是
$R/\mathfrak p^{(n)}$ 上的非零因子。事实上，如果 $y \in R$ 且
$xy \in \mathfrak p^{(n)} = R \cap \mathfrak p^nR_{\mathfrak p}$，
那么 $y \in \mathfrak p^nR_{\mathfrak p}$，因而
$y \in \mathfrak p^{(n)}$。引理得证。
\end{proof}

\begin{lemma}
\label{algebra-lemma-symbolic-power-flat-extension}
设 $R \to S$ 为平坦环映射，$\mathfrak p \subset R$ 为素理想，
且 $\mathfrak q = \mathfrak p S$ 是 $S$ 的素理想。则
$\mathfrak p^{(n)} S = \mathfrak q^{(n)}$。
\end{lemma}

\begin{proof}
由于
$\mathfrak p^{(n)} = \Ker(R \to R_\mathfrak p/\mathfrak p^nR_\mathfrak p)$，
利用平坦性可知 $\mathfrak p^{(n)} S$ 是映射
$S \to S_\mathfrak p/\mathfrak p^nS_\mathfrak p$ 的核。
另一方面，$\mathfrak q^{(n)}$ 是映射
$S \to S_\mathfrak q/\mathfrak q^nS_\mathfrak q =
S_\mathfrak q/\mathfrak p^nS_\mathfrak q$ 的核。因此只须证明
$$
S_\mathfrak p/\mathfrak p^nS_\mathfrak p
\longrightarrow
S_\mathfrak q/\mathfrak p^nS_\mathfrak q
$$
是单射。注意，右侧的模是将左侧的模关于诸元素
$f \in S$，$f \not \in \mathfrak q$ 局部化所得。
因此只须证明这些元素是
$S_\mathfrak p/\mathfrak p^nS_\mathfrak p$ 上的非零因子。
由平坦性，模 $S_\mathfrak p/\mathfrak p^nS_\mathfrak p$
具有有限滤过，其子商为
$$
\mathfrak p^iS_\mathfrak p/\mathfrak p^{i + 1}S_\mathfrak p
\cong \mathfrak p^iR_\mathfrak p/\mathfrak p^{i + 1}R_\mathfrak p
\otimes_{R_\mathfrak p} S_\mathfrak p \cong
V \otimes_{\kappa(\mathfrak p)} (S/\mathfrak q)_\mathfrak p
$$
其中 $V$ 是 $\kappa(\mathfrak p)$-向量空间。因此，$f$
按所需方式可逆地作用。
\end{proof}

\section{相对相伴素理想集}
\label{algebra-section-relative-assassin}

\noindent
下面讨论相对相伴素理想集。设 $R \to S$ 为环映射，
$N$ 为 $S$-模。在这种情形下，可以引入如下 $S$ 的素理想
$\mathfrak q$ 的集合：
\begin{enumerate}
\item $A$：令 $\mathfrak p = R \cap \mathfrak q$，有
$\mathfrak q \in \text{Ass}_S(N \otimes_R \kappa(\mathfrak p))$；
\item $A'$：令 $\mathfrak p = R \cap \mathfrak q$，则在典范映射
$\Spec(S \otimes_R \kappa(\mathfrak p)) \to \Spec(S)$ 下，
$\mathfrak q$ 属于
$\text{Ass}_{S \otimes \kappa(\mathfrak p)}(N \otimes_R \kappa(\mathfrak p))$
的像；
\item $A_{fin}$：令 $\mathfrak p = R \cap \mathfrak q$，有
$\mathfrak q \in \text{Ass}_S(N/\mathfrak pN)$；
\item $A'_{fin}$：对某个素理想 $\mathfrak p' \subset R$，有
$\mathfrak q \in \text{Ass}_S(N/\mathfrak p'N)$；
\item $B$：对某个 $R$-模 $M$，有
$\mathfrak q \in \text{Ass}_S(N \otimes_R M)$；以及
\item $B_{fin}$：对某个有限 $R$-模 $M$，有
$\mathfrak q \in \text{Ass}_S(N \otimes_R M)$。
\end{enumerate}
下面确定这些集合之间的一些关系。

\begin{lemma}
\label{algebra-lemma-compare-relative-assassins}
设 $R \to S$ 为环映射，$N$ 为 $S$-模，并设
$A$、$A'$、$A_{fin}$、$B$ 与 $B_{fin}$ 为上面引入的
$\Spec(S)$ 的子集。
\begin{enumerate}
\item 总有 $A = A'$。
\item 总有 $A_{fin} \subset A$、
$B_{fin} \subset B$、$A_{fin} \subset A'_{fin} \subset B_{fin}$
以及 $A \subset B$。
\item 如果 $S$ 是诺特环，则 $A = A_{fin}$ 且 $B = B_{fin}$。
\item 如果 $N$ 在 $R$ 上平坦，则
$A = A_{fin} = A'_{fin}$ 且 $B = B_{fin}$。
\item 如果 $R$ 是诺特环且 $N$ 在 $R$ 上平坦，则所有这些集合相等，
即 $A = A' = A_{fin} = A'_{fin} = B = B_{fin}$。
\end{enumerate}
\end{lemma}

\begin{proof}
本证明中的一些论证会在后面更精确的引理的证明中重复出现
（因为那些引理处理给定的模 $M$ 或给定的素理想
$\mathfrak p$，而不是它们的全体）。

\medskip\noindent
(1) 的证明。设 $\mathfrak p$ 为 $R$ 的素理想。则有
$$
\text{Ass}_S(N \otimes_R \kappa(\mathfrak p)) =
\text{Ass}_{S/\mathfrak pS}(N \otimes_R \kappa(\mathfrak p)) =
\text{Ass}_{S \otimes_R \kappa(\mathfrak p)}(N \otimes_R \kappa(\mathfrak p))
$$
其中第一个等式由引理
\ref{algebra-lemma-ass-quotient-ring} 得出，第二个等式由引理
\ref{algebra-lemma-localize-ass} 的第 (1) 部分得出。这证明了
$A = A'$。包含 $A_{fin} \subset A'_{fin}$ 是清楚的。

\medskip\noindent
(2) 的证明。除 $A_{fin} \subset A$ 外，各个包含关系都立即由定义得出；
而该包含关系由引理 \ref{algebra-lemma-localize-ass} 以及在
$A_{fin}$ 的表述中要求 $\mathfrak p = R \cap \mathfrak q$
这一事实得出。

\medskip\noindent
(3) 的证明。如果 $S$ 是诺特环，则等式 $A = A_{fin}$ 由引理
\ref{algebra-lemma-localize-ass} 的第 (3) 部分得出。设
$\mathfrak q = (g_1, \ldots, g_m)$ 是 $S$ 的有限生成素理想。
设 $z \in N \otimes_R M$ 的零化子为 $\mathfrak q$。
可以选取有限子模 $M' \subset M$，使 $z$ 是某个
$z' \in N \otimes_R M'$ 的像。于是
$\text{Ann}_S(z') \subset \mathfrak q = \text{Ann}_S(z)$。
由于 $N \otimes_R -$ 与余极限交换，且 $M$ 是有限 $R$-模的
有向余极限，可以找到 $M' \subset M'' \subset M$，使得像
$z'' \in N \otimes_R M''$ 被
$g_1, \ldots, g_m$ 零化。因此
$\text{Ann}_S(z'') = \mathfrak q$。这证明了：如果 $S$ 是诺特环，
则 $B = B_{fin}$。

\medskip\noindent
(4) 的证明。如果 $N$ 平坦，则函子 $N \otimes_R -$ 正合。
特别地，如果 $M' \subset M$，则
$N \otimes_R M' \subset N \otimes_R M$。因此，如果
$z \in N \otimes_R M$ 的零化子
$\mathfrak q = \text{Ann}_S(z)$ 是素理想，那么可以选取任意有限
$R$-子模 $M' \subset M$，使得
$z \in N \otimes_R M'$；而将 $z$ 视为
$N \otimes_R M'$ 的元素时，其零化子仍等于
$\mathfrak q$。因此 $B = B_{fin}$。设 $\mathfrak p'$ 为
$R$ 的素理想，并设 $\mathfrak q$ 为 $S$ 的素理想，且它是
$N/\mathfrak p'N$ 的相伴素理想。
这意味着 $\mathfrak p'S \subset \mathfrak q$。由于 $N$
在 $R$ 上平坦，可知 $N/\mathfrak p'N$ 在整环
$R/\mathfrak p'$ 上平坦。因此，$R/\mathfrak p'$ 的每个非零元素
都是 $N/\mathfrak p'$ 上的非零因子。于是这些元素都不可能映到
$\mathfrak q$ 的元素，故 $\mathfrak p' = R \cap \mathfrak q$。
因此 $A_{fin} = A'_{fin}$。最后，由引理
\ref{algebra-lemma-localize-ass-nonzero-divisors} 可知
$\text{Ass}_S(N/\mathfrak p'N) =
\text{Ass}_S(N \otimes_R \kappa(\mathfrak p'))$，即
$A'_{fin} = A$。

\medskip\noindent
(5) 的证明。由于其他等式已在 (4) 中证明，只须证明
$A'_{fin} = B_{fin}$。为此，设 $M$ 为有限 $R$-模。
由引理 \ref{algebra-lemma-filter-Noetherian-module}，存在一条由
$R$-子模组成的滤过
$$
0 = M_0 \subset M_1 \subset \ldots \subset M_n = M
$$
使得每个商 $M_i/M_{i-1}$ 都同构于某个
$R/\mathfrak p_i$，其中 $\mathfrak p_i$ 是 $R$ 的素理想。
由于 $N$ 平坦，得到一条由 $S$-子模组成的滤过
$$
0 = N \otimes_R M_0 \subset N \otimes_R M_1 \subset \ldots \subset
N \otimes_R M_n = N \otimes_R M
$$
其每个子商都同构于 $N/\mathfrak p_iN$。由引理
\ref{algebra-lemma-ass}，可得
$\text{Ass}_S(N \otimes_R M) \subset \bigcup \text{Ass}_S(N/\mathfrak p_iN)$。
因此 $B_{fin} \subset A'_{fin}$。另一个包含关系是 (2) 的一部分，
故结论得证。
\end{proof}

\noindent
我们将 $N$ 在 $S/R$ 上的相对相伴素理想集定义为上面的集合
$A = A'$。作为动机，指出它只依赖于纤维环上的纤维模
$N \otimes_R \kappa(\mathfrak p)$。与模的相伴素理想集一样，
我们提醒读者：当纤维环
$S \otimes_R \kappa(\mathfrak p)$ 为诺特环时，这个概念最为适用；
例如，当 $R \to S$ 是有限型映射时。

\begin{definition}
\label{algebra-definition-relative-assassin}
设 $R \to S$ 为环映射，$N$ 为 $S$-模。
{\it $N$ 在 $S/R$ 上的相对相伴素理想集}是集合
$$
\text{Ass}_{S/R}(N)
=
\{ \mathfrak q \subset S \mid
\mathfrak q \in \text{Ass}_S(N \otimes_R \kappa(\mathfrak p))
\text{ 其中 }\mathfrak p = R \cap \mathfrak q\}.
$$
这就是引理 \ref{algebra-lemma-compare-relative-assassins} 中记作
$A$ 的集合。
\end{definition}

\noindent
下面几个结果的实质都与相对相伴素理想集有关，
尽管这一点未必显而易见。

\begin{lemma}
\label{algebra-lemma-bourbaki}
设 $R \to S$ 为环映射，$M$ 为 $R$-模，$N$ 为 $S$-模。
如果 $N$ 作为 $R$-模是平坦的，则
$$
\text{Ass}_S(M \otimes_R N)
\supset
\bigcup\nolimits_{\mathfrak p \in \text{Ass}_R(M)} \text{Ass}_S(N/\mathfrak pN)
$$
并且如果 $R$ 是诺特环，则等号成立。
\end{lemma}

\begin{proof}
如果 $\mathfrak p \in \text{Ass}_R(M)$，则存在单射
$R/\mathfrak p \to M$。由于 $N$ 在 $R$ 上平坦，得到单射
$R/\mathfrak p \otimes_R N \to M \otimes_R N$。又因为
$R/\mathfrak p \otimes_R N = N/\mathfrak pN$，由引理
\ref{algebra-lemma-ass} 可得
$\text{Ass}_S(N/\mathfrak pN) \subset \text{Ass}_S(M \otimes_R N)$。
因此右侧包含于左侧。

\medskip\noindent
将 $M = \bigcup M_\lambda$ 写成其有限生成
$R$-子模之并。于是也有
$N \otimes_R M = \bigcup N \otimes_R M_\lambda$
（因为 $N$ 在 $R$ 上平坦）。由相伴素理想的定义可知
$\text{Ass}_S(N \otimes_R M) = \bigcup \text{Ass}_S(N \otimes_R M_\lambda)$
且 $\text{Ass}_R(M) = \bigcup \text{Ass}(M_\lambda)$。
因此可以假设 $M$ 是有限生成的。

\medskip\noindent
设 $\mathfrak q \in \text{Ass}_S(M \otimes_R N)$，并假设
$R$ 是诺特环且 $M$ 是有限 $R$-模。为完成证明，须说明
$\mathfrak q$ 属于右侧。首先注意
$\mathfrak qS_{\mathfrak q} \in
\text{Ass}_{S_{\mathfrak q}}((M \otimes_R N)_{\mathfrak q})$，
参见引理 \ref{algebra-lemma-associated-primes-localize}。
设 $\mathfrak p$ 为 $R$ 中对应的素理想。注意
$$
(M \otimes_R N)_{\mathfrak q} = M \otimes_R N_{\mathfrak q}
= M_{\mathfrak p} \otimes_{R_{\mathfrak p}} N_{\mathfrak q}
$$
如果
$\mathfrak pR_{\mathfrak p} \not \in
\text{Ass}_{R_{\mathfrak p}}(M_{\mathfrak p})$，
那么存在 $x \in \mathfrak pR_{\mathfrak p}$，它是
$M_{\mathfrak p}$ 中的非零因子（参见引理
\ref{algebra-lemma-ideal-nonzerodivisor}）。由于
$N_{\mathfrak q}$ 在 $R_{\mathfrak p}$ 上平坦，可知 $x$ 在
$\mathfrak qS_{\mathfrak q}$ 中的像是
$(M \otimes_R N)_{\mathfrak q}$ 上的非零因子。
这与假设
$\mathfrak qS_{\mathfrak q} \in \text{Ass}_S((M \otimes_R N)_{\mathfrak q})$
矛盾。因此，$\mathfrak p$ 是 $M$ 的相伴素理想之一。

\medskip\noindent
继续这一论证，选取滤过
$$
0 = M_0 \subset M_1 \subset \ldots \subset M_n = M
$$
使得每个商 $M_i/M_{i-1}$ 都同构于某个
$R/\mathfrak p_i$，其中 $\mathfrak p_i$ 是 $R$ 的素理想；
参见引理 \ref{algebra-lemma-filter-Noetherian-module}。
（由引理 \ref{algebra-lemma-ass-filter}，至少对一个 $i$ 有
$\mathfrak p_i = \mathfrak p$。）这给出滤过
$$
0 = M_0 \otimes_R N \subset M_1 \otimes_R N \subset \ldots
\subset M_n \otimes_R N = M \otimes_R N
$$
其子商同构于 $N/\mathfrak p_iN$。如果
$\mathfrak p_i \not = \mathfrak p$，那么由前一段的结果，
$\mathfrak q$ 不可能与模 $N/\mathfrak p_iN$ 相伴
（因为 $\text{Ass}_R(R/\mathfrak p_i) = \{\mathfrak p_i\}$）。
因此，$\mathfrak q$ 按所需与 $N/\mathfrak pN$ 相伴。
\end{proof}

\begin{lemma}
\label{algebra-lemma-post-bourbaki}
设 $R \to S$ 为环映射，$N$ 为 $S$-模。假设 $N$
作为 $R$-模是平坦的，并且 $R$ 是分式域为 $K$ 的整环。则
$$
\text{Ass}_S(N) =
\text{Ass}_S(N \otimes_R K) =
\text{Ass}_{S \otimes_R K}(N \otimes_R K)
$$
其中使用典范包含
$\Spec(S \otimes_R K) \subset \Spec(S)$。
\end{lemma}

\begin{proof}
注意 $S \otimes_R K = (R \setminus \{0\})^{-1}S$ 且
$N \otimes_R K = (R \setminus \{0\})^{-1}N$。
对任意非零 $x \in R$，由于 $N$ 在 $R$ 上平坦，
乘以 $x$ 在 $N$ 上是单射。因此，本引理由引理
\ref{algebra-lemma-localize-ass-nonzero-divisors} 与引理
\ref{algebra-lemma-localize-ass} 的第 (1) 部分结合得出。
\end{proof}

\begin{lemma}
\label{algebra-lemma-bourbaki-fibres}
设 $R \to S$ 为环映射，$M$ 为 $R$-模，$N$ 为 $S$-模。
假设 $N$ 作为 $R$-模是平坦的。则
$$
\text{Ass}_S(M \otimes_R N)
\supset
\bigcup\nolimits_{\mathfrak p \in \text{Ass}_R(M)}
\text{Ass}_{S \otimes_R \kappa(\mathfrak p)}(N \otimes_R \kappa(\mathfrak p))
$$
这里使用注 \ref{algebra-remark-fundamental-diagram}，
将纤维环的谱视为 $\Spec(S)$ 的子集。
如果 $R$ 是诺特环，则此包含为等式。
\end{lemma}

\begin{proof}
这等价于引理
\ref{algebra-lemma-bourbaki}；这是由引理
\ref{algebra-lemma-ass-quotient-ring}、
\ref{algebra-lemma-flat-base-change} 及
\ref{algebra-lemma-post-bourbaki} 得出的。
\end{proof}

\begin{remark}
\label{algebra-remark-bourbaki}
设 $R \to S$ 为环映射，$N$ 为 $S$-模，并设
$\mathfrak p$ 为 $R$ 的素理想。则
$$
\text{Ass}_S(N \otimes_R \kappa(\mathfrak p)) =
\text{Ass}_{S/\mathfrak pS}(N \otimes_R \kappa(\mathfrak p)) =
\text{Ass}_{S \otimes_R \kappa(\mathfrak p)}(N \otimes_R \kappa(\mathfrak p)).
$$
第一个等式由引理
\ref{algebra-lemma-ass-quotient-ring} 得出，第二个等式由引理
\ref{algebra-lemma-localize-ass} 的第 (1) 部分得出。
\end{remark}

\section{弱相伴素理想}
\label{algebra-section-weakly-ass}

\noindent
这是相伴素理想概念的一个变体，对非诺特环和非有限模很有用。

\begin{definition}
\label{algebra-definition-weakly-associated}
设 $R$ 为环，$M$ 为 $R$-模。如果存在元素 $m \in M$，
使得 $\mathfrak p$ 在包含零化子
$\text{Ann}(m) = \{f \in R \mid fm = 0\}$ 的素理想中极小，
则称 $R$ 的素理想 $\mathfrak p$ 与 $M$ {\it 弱相伴}。
所有这样的素理想组成的集合记作
$\text{WeakAss}_R(M)$ 或 $\text{WeakAss}(M)$。
\end{definition}

\noindent
因此，相伴素理想必为弱相伴素理想。
下面用在该素理想处的局部化给出一个刻画。

\begin{lemma}
\label{algebra-lemma-weakly-ass-local}
设 $R$ 为环，$M$ 为 $R$-模，$\mathfrak p$ 为 $R$ 的素理想。
下列条件等价：
\begin{enumerate}
\item $\mathfrak p$ 与 $M$ 弱相伴；
\item $\mathfrak pR_{\mathfrak p}$ 与 $M_{\mathfrak p}$ 弱相伴；以及
\item $M_{\mathfrak p}$ 含有一个元素，其零化子的根等于
$\mathfrak pR_{\mathfrak p}$。
\end{enumerate}
\end{lemma}

\begin{proof}
假设 (1) 成立。则存在元素 $m \in M$，使
$\mathfrak p$ 在包含 $m$ 的零化子
$I = \{x \in R \mid xm = 0\}$ 的素理想中极小。
由于局部化正合，$m$ 在 $M_{\mathfrak p}$ 中的零化子为
$I_{\mathfrak p}$。因此，$\mathfrak pR_{\mathfrak p}$ 是
$R_{\mathfrak p}$ 中包含 $m$ 在 $M_{\mathfrak p}$ 中的零化子
$I_{\mathfrak p}$ 的极小素理想。这说明 (2) 成立；又因为它意味着
$\sqrt{I_{\mathfrak p}} = \mathfrak pR_{\mathfrak p}$，
(3) 也成立。

\medskip\noindent
将蕴含 (1) $\Rightarrow$ (3) 应用于 $R_{\mathfrak p}$ 上的
$M_{\mathfrak p}$，可知 (2) $\Rightarrow$ (3)。

\medskip\noindent
最后，假设 (3) 成立。这意味着存在元素
$m/f \in M_{\mathfrak p}$，其零化子的根等于
$\mathfrak pR_{\mathfrak p}$。于是，$m$ 在 $M$ 中的零化子
$I = \{x \in R \mid xm = 0\}$ 满足
$\sqrt{I_{\mathfrak p}} = \mathfrak pR_{\mathfrak p}$。
这显然意味着 $\mathfrak p$ 包含 $I$，且在包含 $I$ 的素理想中极小，
即 (1) 成立。
\end{proof}

\begin{lemma}
\label{algebra-lemma-reduced-weakly-ass-minimal}
对于约化环，该环的弱相伴素理想正是极小素理想。
\end{lemma}

\begin{proof}
设 $(R, \mathfrak m)$ 为约化局部环。假设 $x \in R$ 的零化子的根为
$\mathfrak m$。如果 $\mathfrak m \not = 0$，则 $x$ 不可能是单位，
所以 $x \in \mathfrak m$。特别地，对某个 $n \geq 0$ 有
$x^{1 + n} = 0$。因此 $x = 0$，这与 $x$ 的零化子包含于
$\mathfrak m$ 的假设矛盾。由此可见 $\mathfrak m = 0$，
即 $R$ 是域。由引理 \ref{algebra-lemma-weakly-ass-local}，
这就蕴含本引理的断言。
\end{proof}

\begin{lemma}
\label{algebra-lemma-weakly-ass}
设 $R$ 为环，并设
$0 \to M' \to M \to M'' \to 0$ 是 $R$-模的短正合列。
则 $\text{WeakAss}(M') \subset \text{WeakAss}(M)$ 且
$\text{WeakAss}(M) \subset \text{WeakAss}(M') \cup \text{WeakAss}(M'')$。
\end{lemma}

\begin{proof}
我们将使用引理 \ref{algebra-lemma-weakly-ass-local}
对弱相伴素理想的刻画。设 $\mathfrak p$ 为 $R$ 的素理想。
由于局部化正合，得到短正合列
$0 \to M'_{\mathfrak p} \to M_{\mathfrak p} \to M''_{\mathfrak p} \to 0$。
假设 $m \in M_{\mathfrak p}$ 的零化子的根为
$\mathfrak pR_{\mathfrak p}$。那么，要么 $m$ 在
$M''_{\mathfrak p}$ 中的像 $\overline{m}$ 为零且
$m \in M'_{\mathfrak p}$，要么 $\overline{m}$ 的零化子的根为
$\mathfrak pR_{\mathfrak p}$。这证明了
$\text{WeakAss}(M) \subset \text{WeakAss}(M') \cup \text{WeakAss}(M'')$。
包含 $\text{WeakAss}(M') \subset \text{WeakAss}(M)$
立即由定义得出。
\end{proof}

\begin{lemma}
\label{algebra-lemma-weakly-ass-zero}
\begin{slogan}
每个非零模都有弱相伴素理想。
\end{slogan}
设 $R$ 为环，$M$ 为 $R$-模。则
$$
M = (0) \Leftrightarrow \text{WeakAss}(M) = \emptyset
$$
\end{lemma}

\begin{proof}
如果 $M = (0)$，则依定义
$\text{WeakAss}(M) = \emptyset$。反之，假设 $M \not = 0$。
选取非零元素 $m \in M$。将 $m$ 的零化子写成
$I = \{x \in R \mid xm = 0\}$。于是 $R/I \subset M$。
因此，由引理 \ref{algebra-lemma-weakly-ass}，
$\text{WeakAss}(R/I) \subset \text{WeakAss}(M)$。
但由于 $I \not = R$，$V(I) = \Spec(R/I)$ 含有极小素理想；
参见引理
\ref{algebra-lemma-Zariski-topology} 和
\ref{algebra-lemma-spec-closed}。结论得证。
\end{proof}

\begin{lemma}
\label{algebra-lemma-weakly-ass-support}
设 $R$ 为环，$M$ 为 $R$-模。则
$$
\text{Ass}(M) \subset \text{WeakAss}(M) \subset \text{Supp}(M).
$$
\end{lemma}

\begin{proof}
第一个包含立即由定义得出。如果
$\mathfrak p \in \text{WeakAss}(M)$，则由引理
\ref{algebra-lemma-weakly-ass-local} 有
$M_{\mathfrak p} \not = 0$，故
$\mathfrak p \in \text{Supp}(M)$。
\end{proof}

\begin{lemma}
\label{algebra-lemma-weakly-ass-zero-divisors}
设 $R$ 为环，$M$ 为 $R$-模。并集
$\bigcup_{\mathfrak q \in \text{WeakAss}(M)} \mathfrak q$
正是 $R$ 中在 $M$ 上为零因子的元素所成的集合。
\end{lemma}

\begin{proof}
假设 $f \in \mathfrak q \in \text{WeakAss}(M)$。
则存在元素 $m \in M$，使得 $\mathfrak q$ 在包含
$I = \{x \in R \mid xm = 0\}$ 的素理想中极小。
因此，存在 $g \in R$，$g \not \in \mathfrak q$ 以及 $n > 0$，
使得 $f^ngm = 0$。注意，由于 $g \not \in I$，
$gm \not = 0$。如果取上述 $n$ 为最小，则
$f (f^{n - 1}gm) = 0$ 且 $f^{n - 1}gm \not = 0$，
所以 $f$ 是 $M$ 上的零因子。
反之，假设 $f \in R$ 是 $M$ 上的零因子。考虑子模
$N = \{m \in M \mid fm = 0\}$。由于 $N$ 非零，由引理
\ref{algebra-lemma-weakly-ass-zero}，它有弱相伴素理想 $\mathfrak q$。
显然 $f \in \mathfrak q$；又由引理
\ref{algebra-lemma-weakly-ass}，$\mathfrak q$ 是 $M$ 的弱相伴素理想。
\end{proof}

\begin{lemma}
\label{algebra-lemma-weakly-ass-minimal-prime-support}
设 $R$ 为环，$M$ 为 $R$-模。在 $\text{Supp}(M)$ 的元素中
极小的任意 $\mathfrak p \in \text{Supp}(M)$ 都属于
$\text{WeakAss}(M)$。
\end{lemma}

\begin{proof}
注意，在 $\Spec(R_{\mathfrak p})$ 中有
$\text{Supp}(M_{\mathfrak p}) = \{\mathfrak pR_{\mathfrak p}\}$。
特别地，$M_{\mathfrak p}$ 非零，所以由引理
\ref{algebra-lemma-weakly-ass-zero}，
$\text{WeakAss}(M_{\mathfrak p}) \not = \emptyset$。
又由引理
\ref{algebra-lemma-weakly-ass-support} 有
$\text{WeakAss}(M_{\mathfrak p}) \subset \text{Supp}(M_{\mathfrak p})$，
故
$\text{WeakAss}(M_{\mathfrak p}) = \{\mathfrak pR_{\mathfrak p}\}$。
于是由引理
\ref{algebra-lemma-weakly-ass-local}，
$\mathfrak p \in \text{WeakAss}(M)$。
\end{proof}

\begin{lemma}
\label{algebra-lemma-ass-weakly-ass}
设 $R$ 为环，$M$ 为 $R$-模，并设 $\mathfrak p$ 是 $R$
的有限生成素理想。则
$$
\mathfrak p \in \text{Ass}(M) \Leftrightarrow
\mathfrak p \in \text{WeakAss}(M).
$$
特别地，如果 $R$ 是诺特环，则
$\text{Ass}(M) = \text{WeakAss}(M)$。
\end{lemma}

\begin{proof}
写 $\mathfrak p = (g_1, \ldots, g_n)$，其中 $g_i \in R$。
由于另一个蕴含一般成立，只须证明蕴含“$\Leftarrow$”；
参见引理 \ref{algebra-lemma-weakly-ass-support}。
假设 $\mathfrak p \in \text{WeakAss}(M)$。由引理
\ref{algebra-lemma-weakly-ass-local}，存在元素 $m \in M_{\mathfrak p}$，
使得
$I = \{x \in R_{\mathfrak p} \mid xm = 0\}$ 的根为
$\mathfrak pR_{\mathfrak p}$。因此，对每个 $i$，存在最小的
$e_i > 0$，使得 $g_i^{e_i}m = 0$ 在 $M_{\mathfrak p}$ 中成立。
如果对某个 $i$ 有 $e_i > 1$，则可将 $m$ 替换为
$g_i^{e_i - 1} m \not = 0$，并减小 $\sum e_i$。
因此可以假设 $m \in M_{\mathfrak p}$ 的零化子为
$(g_1, \ldots, g_n)R_{\mathfrak p} = \mathfrak p R_{\mathfrak p}$。
由引理 \ref{algebra-lemma-associated-primes-localize} 可知
$\mathfrak p \in \text{Ass}(M)$。
\end{proof}

\begin{remark}
\label{algebra-remark-weakly-ass-not-functorial}
设 $\varphi : R \to S$ 为环映射，$M$ 为 $S$-模。
与相伴素理想的情形相反，并非总有
$\Spec(\varphi)(\text{WeakAss}_S(M)) \subset \text{WeakAss}_R(M)$
（参见引理 \ref{algebra-lemma-ass-functorial}）。
一个例子是考虑环映射
$$
R = k[x_1, x_2, x_3, \ldots] \to
S = k[x_1, x_2, x_3, \ldots, y_1, y_2, y_3, \ldots]/
(x_1y_1, x_2y_2, x_3y_3, \ldots)
$$
以及 $M = S$。在这种情形下，
$\mathfrak q = \sum x_iS$ 是 $S$ 的极小素理想，因而是
$M = S$ 的弱相伴素理想（参见引理
\ref{algebra-lemma-weakly-ass-minimal-prime-support}）。
但另一方面，$S$ 的任意非零元素在 $R$ 中的零化子都是有限生成的，
所以其根不可能等于
$R \cap \mathfrak q = (x_1, x_2, x_3, \ldots)$
（细节从略）。
\end{remark}

\begin{lemma}
\label{algebra-lemma-weakly-ass-reverse-functorial}
设 $\varphi : R \to S$ 为环映射，$M$ 为 $S$-模。则
$\Spec(\varphi)(\text{WeakAss}_S(M)) \supset \text{WeakAss}_R(M)$。
\end{lemma}

\begin{proof}
设 $\mathfrak p$ 是 $\text{WeakAss}_R(M)$ 的元素。
则存在 $m \in M_{\mathfrak p}$，其零化子
$I = \{x \in R_{\mathfrak p} \mid xm = 0\}$ 的根为
$\mathfrak pR_{\mathfrak p}$。考虑 $m$ 在 $S_{\mathfrak p}$ 中的零化子
$J = \{x \in S_{\mathfrak p} \mid xm = 0 \}$。
由于 $IS_{\mathfrak p} \subset J$，可知包含 $J$ 的任意极小素理想
$\mathfrak q \subset S_{\mathfrak p}$ 都位于 $\mathfrak p$ 之上。
此外，例如由引理
\ref{algebra-lemma-weakly-ass-local}，这样的 $\mathfrak q$
对应于 $M$ 的一个弱相伴素理想。
\end{proof}

\begin{remark}
\label{algebra-remark-ass-functorial}
设 $\varphi : R \to S$ 为环映射，$M$ 为 $S$-模。
将谱上相伴的映射记作
$f : \Spec(S) \to \Spec(R)$。则有
$$
f(\text{Ass}_S(M)) \subset
\text{Ass}_R(M) \subset
\text{WeakAss}_R(M) \subset
f(\text{WeakAss}_S(M))
$$
参见引理
\ref{algebra-lemma-ass-functorial}、
\ref{algebra-lemma-weakly-ass-reverse-functorial} 和
\ref{algebra-lemma-weakly-ass-support}。
一般而言，所有这些包含都可能是严格的；参见注
\ref{algebra-remark-ass-reverse-functorial} 和
\ref{algebra-remark-weakly-ass-not-functorial}。
如果 $S$ 是诺特环，那么所有这些包含都是等式，因为由引理
\ref{algebra-lemma-ass-weakly-ass}，最外侧的两个包含是等式。
\end{remark}

\begin{lemma}
\label{algebra-lemma-weakly-ass-finite-ring-map}
设 $\varphi : R \to S$ 为环映射，$M$ 为 $S$-模。
将谱上相伴的映射记作
$f : \Spec(S) \to \Spec(R)$。如果 $\varphi$ 是有限环映射，则
$$
\text{WeakAss}_R(M) = f(\text{WeakAss}_S(M)).
$$
\end{lemma}

\begin{proof}
其中一个包含已经证明；参见注
\ref{algebra-remark-ass-functorial}。为证明另一个包含，假设
$\mathfrak q \in \text{WeakAss}_S(M)$，并设
$\mathfrak p$ 为 $R$ 中对应的素理想。设 $m \in M$ 满足：
$\mathfrak q$ 是包含
$J = \{g \in S \mid gm = 0\}$ 的极小素理想。因此
$JS_{\mathfrak q}$ 的根为 $\mathfrak qS_{\mathfrak q}$。
由于 $R \to S$ 是有限的，$\mathfrak p$ 上方只有有限多个素理想
$\mathfrak q = \mathfrak q_1, \mathfrak q_2, \ldots, \mathfrak q_l$；
参见引理 \ref{algebra-lemma-finite-finite-fibres}。
选取 $x \in \mathfrak q$，使得对 $i > 1$ 有
$x \not \in \mathfrak q_i$；参见引理
\ref{algebra-lemma-silly}。由上述讨论，存在元素
$y \in S$，$y \not \in \mathfrak q$ 以及整数 $t > 0$，
使得 $y x^t m = 0$。因此，元素 $ym \in M$ 被
$x^t$ 零化，故对 $i = 2, \ldots, l$，$ym$ 在
$M_{\mathfrak q_i}$ 中映为零。须注意，$ym$ 在
$S_{\mathfrak q}$ 中并不映为零。

\medskip\noindent
由有限环映射的上升性质，环 $S_{\mathfrak p}$ 是半局部环，
其极大理想为 $\mathfrak q_i S_{\mathfrak p}$；参见引理
\ref{algebra-lemma-integral-going-up}。如果
$f \in \mathfrak pR_{\mathfrak p}$，则 $f$ 的某个幂落在
$JS_{\mathfrak q}$ 中，故对某个 $t > 0$，可知
$f^t ym$ 在 $M_{\mathfrak q}$ 中映为零。
由于 $ym$ 在 $S_{\mathfrak p}$ 的其他极大理想处为零，
由引理 \ref{algebra-lemma-characterize-zero-local} 可知
$f^t ym$ 在 $M_{\mathfrak p}$ 中为零。
由此可见，$\mathfrak p$ 是包含 $ym$ 在 $R$ 中的零化子的
极小素理想，结论得证。
\end{proof}

\begin{lemma}
\label{algebra-lemma-weakly-ass-quotient-ring}
设 $R$ 为环，$I$ 为理想，$M$ 为 $R/I$-模。
通过典范单射
$\Spec(R/I) \to \Spec(R)$，有
$\text{WeakAss}_{R/I}(M) = \text{WeakAss}_R(M)$。
\end{lemma}

\begin{proof}
这是引理 \ref{algebra-lemma-weakly-ass-finite-ring-map} 的特殊情形。
\end{proof}

\begin{lemma}
\label{algebra-lemma-localize-weakly-ass}
设 $R$ 为环，$M$ 为 $R$-模，并设 $S \subset R$ 为乘法子集。
通过典范单射 $\Spec(S^{-1}R) \to \Spec(R)$，有
$\text{WeakAss}_R(S^{-1}M) = \text{WeakAss}_{S^{-1}R}(S^{-1}M)$
以及
$$
\text{WeakAss}(M) \cap \Spec(S^{-1}R) = \text{WeakAss}(S^{-1}M).
$$
\end{lemma}

\begin{proof}
假设 $m \in S^{-1}M$。令
$I = \{x \in R \mid xm = 0\}$ 且
$I' = \{x' \in S^{-1}R \mid x'm = 0\}$。则
$I' = S^{-1}I$，并且除非 $I = R$，否则
$I \cap S = \emptyset$（验证从略）。
因此，$S^{-1}R$ 中包含 $I'$ 的极小素理想，与 $R$ 中
包含 $I$ 且避开 $S$ 的极小素理想一一对应。这证明了等式
$\text{WeakAss}_R(S^{-1}M) = \text{WeakAss}_{S^{-1}R}(S^{-1}M)$。
第二个等式由引理
\ref{algebra-lemma-weakly-ass-local} 得出，因为对
$\mathfrak p \in R$，若 $S \cap \mathfrak p = \emptyset$，则
$M_{\mathfrak p} = (S^{-1}M)_{S^{-1}\mathfrak p}$。
\end{proof}

\begin{lemma}
\label{algebra-lemma-localize-weakly-ass-nonzero-divisors}
设 $R$ 为环，$M$ 为 $R$-模，并设 $S \subset R$ 为乘法子集。
假设每个 $s \in S$ 都是 $M$ 上的非零因子。则
$$
\text{WeakAss}(M) = \text{WeakAss}(S^{-1}M).
$$
\end{lemma}

\begin{proof}
由假设有 $M \subset S^{-1}M$，故引理
\ref{algebra-lemma-weakly-ass} 给出
$\text{WeakAss}(M) \subset \text{WeakAss}(S^{-1}M)$。
反之，假设 $n/s \in S^{-1}M$ 的零化子为 $I$，且
$\mathfrak p$ 是包含 $I$ 的极小素理想。则
$n \in M$ 的零化子为 $I$，且 $\mathfrak p$ 是包含
$I$ 的极小素理想。
\end{proof}

\begin{lemma}
\label{algebra-lemma-zero-at-weakly-ass-zero}
设 $R$ 为环，$M$ 为 $R$-模。映射
$$
M
\longrightarrow
\prod\nolimits_{\mathfrak p \in \text{WeakAss}(M)} M_{\mathfrak p}
$$
是单射。
\end{lemma}

\begin{proof}
设 $x \in M$ 为该映射的核中的元素。令
$N = Rx \subset M$。如果 $\mathfrak p$ 是 $N$ 的弱相伴素理想，
那么一方面 $\mathfrak p \in \text{WeakAss}(M)$
（引理 \ref{algebra-lemma-weakly-ass}），另一方面
$N_{\mathfrak p} \subset M_{\mathfrak p}$ 非零。
此矛盾说明 $\text{WeakAss}(N) = \emptyset$。
因此，由引理
\ref{algebra-lemma-weakly-ass-zero}，$N = 0$，即 $x = 0$。
\end{proof}

\begin{lemma}
\label{algebra-lemma-weak-post-bourbaki}
设 $R \to S$ 为环映射，$N$ 为 $S$-模。
假设 $N$ 作为 $R$-模是平坦的，并且 $R$ 是分式域为 $K$ 的整环。
则
$$
\text{WeakAss}_S(N) = \text{WeakAss}_{S \otimes_R K}(N \otimes_R K)
$$
其中使用典范包含
$\Spec(S \otimes_R K) \subset \Spec(S)$。
\end{lemma}

\begin{proof}
注意 $S \otimes_R K = (R \setminus \{0\})^{-1}S$ 且
$N \otimes_R K = (R \setminus \{0\})^{-1}N$。
对任意非零 $x \in R$，由于 $N$ 在 $R$ 上平坦，
乘以 $x$ 在 $N$ 上是单射。因此，本引理由引理
\ref{algebra-lemma-localize-weakly-ass-nonzero-divisors} 得出。
\end{proof}

\begin{lemma}
\label{algebra-lemma-weakly-ass-change-fields}
设 $K/k$ 为域扩张，$R$ 为 $k$-代数，$M$ 为 $R$-模。
设 $\mathfrak q \subset R \otimes_k K$ 是位于
$\mathfrak p \subset R$ 上方的素理想。如果 $\mathfrak q$
与 $M \otimes_k K$ 弱相伴，那么 $\mathfrak p$ 与 $M$ 弱相伴。
\end{lemma}

\begin{proof}
设 $z \in M \otimes_k K$ 满足：$\mathfrak q$ 是包含 $z$ 的零化子
$J \subset R \otimes_k K$ 的极小素理想。
选择有限生成的子扩张 $K/L/k$，使
$z \in M \otimes_k L$。由于
$R \otimes_k L \to R \otimes_k K$ 平坦，由引理
\ref{algebra-lemma-annihilator-flat-base-change} 可知
$J = I(R \otimes_k K)$，其中 $I \subset R \otimes_k L$
是 $z$ 在较小环中的零化子。因此，由下降定理
（引理 \ref{algebra-lemma-flat-going-down}），
$\mathfrak q \cap (R \otimes_k L)$ 是包含 $I$ 的极小素理想。
由此归约到下一段所述的情形。

\medskip\noindent
假设 $K/k$ 是有限生成域扩张。设
$x_1, \ldots, x_r \in K$ 是 $K$ 在 $k$ 上的超越基；
参见《域》的第 \ref{fields-section-transcendence} 节。
令 $L = k(x_1, \ldots, x_r)$，并设 $[K : L] = n$。
于是 $R \otimes_k L \to R \otimes_k K$ 是有限环映射。
因此，由引理 \ref{algebra-lemma-weakly-ass-finite-ring-map}，
$\mathfrak q \cap (R \otimes_k L)$ 是 $M \otimes_k K$
作为 $R \otimes_k L$-模时的弱相伴素理想。
由于作为 $R \otimes_k L$-模有
$M \otimes_k K \cong (M \otimes_k L)^{\oplus n}$，
可知 $\mathfrak q \cap (R \otimes_k L)$ 是
$M \otimes_k L$ 的弱相伴素理想
（例如，使用引理 \ref{algebra-lemma-weakly-ass} 并作归纳）。
由此归约到下一段讨论的情形。

\medskip\noindent
假设 $K = k(x_1, \ldots, x_r)$ 是纯超越域扩张。
可以将 $R$ 替换为 $R_\mathfrak p$，将 $M$ 替换为
$M_\mathfrak p$，并将 $\mathfrak q$ 替换为
$\mathfrak q(R_\mathfrak p \otimes_k K)$。
参见引理 \ref{algebra-lemma-localize-weakly-ass}。
由此归约到下一段讨论的情形。

\medskip\noindent
假设 $K = k(x_1, \ldots, x_r)$ 是纯超越域扩张，且 $R$
是极大理想为 $\mathfrak p$ 的局部环。我们断言，任意
$f \in R \otimes_k K$，$f \not \in \mathfrak p(R \otimes_k K)$，
都是 $M \otimes_k K$ 上的非零因子。事实上，设
$z \in M \otimes_k K$。存在有限 $R$-子模
$M' \subset M$，使得 $z \in M' \otimes_k K$，而且 $M'$
在具有此性质的子模中极小：选取 $K$ 作为 $k$-向量空间的一组基
$\{t_\alpha\}$，写
$z = \sum m_\alpha \otimes t_\alpha$，并令 $M'$ 为诸
$m_\alpha$ 生成的 $R$-子模。如果
$z \in \mathfrak p(M' \otimes_k K) = \mathfrak p M' \otimes_k K$，
那么 $\mathfrak pM' = M'$，由引理 \ref{algebra-lemma-NAK}
可得 $M' = 0$，矛盾。
因此，$z$ 在 $M'/\mathfrak p M' \otimes_k K$ 中的像
$\overline{z}$ 非零。但
$R/\mathfrak p \otimes_k K$ 是
$\kappa(\mathfrak p)[x_1, \ldots, x_n]$ 的一个局部化，故为整环；
而 $M'/\mathfrak p M' \otimes_k K$ 是自由模，所以
$f\overline{z} \not = 0$。断言得证。

\medskip\noindent
最后，选取 $z \in M \otimes_k K$，使得 $\mathfrak q$
是包含 $z$ 的零化子 $J \subset R \otimes_k K$ 的极小素理想。
对 $f \in \mathfrak p$，存在 $n \geq 1$ 及
$g \in R \otimes_k K$，$g \not \in \mathfrak q$，使得
$g f^n z \in J$，即 $g f^n z = 0$。
（这是因为 $\mathfrak q$ 位于 $\mathfrak p$ 上方，且
$\mathfrak q$ 是包含 $J$ 的极小素理想。）
上面已经看到 $g$ 是非零因子，所以 $f^n z = 0$。
这意味着，将 $M \otimes_k K$ 视为 $R$-模时，
$\mathfrak p$ 是它的弱相伴素理想。由于
$M \otimes_k K$ 是若干个 $M$ 的直和，和上面一样可知
$\mathfrak p$ 是 $M$ 的弱相伴素理想。
\end{proof}

\section{嵌入素理想}
\label{algebra-section-embedded-primes}

\noindent
定义如下。

\begin{definition}
\label{algebra-definition-embedded-primes}
设 $R$ 为环，$M$ 为 $R$-模。
\begin{enumerate}
\item $M$ 的相伴素理想中，那些在 $M$ 的相伴素理想中不极小的，
称为 $M$ 的{\it 嵌入相伴素理想}。
\item {\it $R$ 的嵌入素理想}是将 $R$ 视为 $R$-模时的
嵌入相伴素理想。
\end{enumerate}
\end{definition}

\noindent
下面给出一种去除这些素理想的方法。

\begin{lemma}
\label{algebra-lemma-remove-embedded-primes}
设 $R$ 为诺特环，$M$ 为有限 $R$-模。考虑如下 $R$-子模集合：
$$
\{
K \subset M
\mid
\text{Supp}(K)
\text{ nowhere dense in }
\text{Supp}(M)
\}.
$$
这个集合有极大元 $K$，且商
$M' = M/K$ 具有下列性质：
\begin{enumerate}
\item $\text{Supp}(M) = \text{Supp}(M')$；
\item $M'$ 没有嵌入相伴素理想；
\item 对任意包含于 $M$ 的所有嵌入相伴素理想中的
$f \in R$，有 $M_f \cong M'_f$。
\end{enumerate}
\end{lemma}

\begin{proof}
下文将直接使用引理 \ref{algebra-lemma-finite-ass} 和命题
\ref{algebra-proposition-minimal-primes-associated-primes}，不再逐次说明。
令 $\mathfrak q_1, \ldots, \mathfrak q_t$ 表示 $M$ 的支集中的
极小素理想，令
$\mathfrak p_1, \ldots, \mathfrak p_s$ 表示 $M$ 的
嵌入相伴素理想。则
$\text{Ass}(M) = \{\mathfrak q_j, \mathfrak p_i\}$。
令
$$
K = \{m \in M \mid \text{Supp}(Rm) \subset \bigcup V(\mathfrak p_i)\}
$$
立即可见它是子模。由于 $M$ 是诺特环上的有限模，
$K$ 也是有限的。因此，$\text{Supp}(K)$ 在
$\text{Supp}(M)$ 中无处稠密。设 $K' \subset M$ 是另一个
支集在 $\text{Supp}(M)$ 中无处稠密的子模。这意味着
$K_{\mathfrak q_j} = 0$。因此，如果 $m \in K'$，
那么 $m$ 在 $M_{\mathfrak q_j}$ 中映为零，进而
$(Rm)_{\mathfrak q_j} = 0$。另一方面，有
$\text{Ass}(Rm) \subset \text{Ass}(M)$。
所以 $Rm$ 的支集包含于 $\bigcup V(\mathfrak p_i)$。
因此 $m \in K$；由于 $m$ 是 $K'$ 的任意元素，故
$K' \subset K$。

\medskip\noindent
令 $M' = M/K$。由于 $K_{\mathfrak q_j}=0$，可知对所有 $j$ 有
$M'_{\mathfrak q_j} = M_{\mathfrak q_j}$。
因此 $M$ 与 $M'$ 具有相同的支集。

\medskip\noindent
假设 $\mathfrak q = \text{Ann}(\overline{m}) \in \text{Ass}(M')$，
其中 $\overline{m} \in M'$ 是 $m \in M$ 的像。
于是 $m \not \in K$，所以 $Rm$ 的支集必含有某个
$\mathfrak q_j$。由于
$M_{\mathfrak q_j} = M'_{\mathfrak q_j}$，可知
$\overline{m}$ 在 $M'_{\mathfrak q_j}$ 中的像非零。
因此 $\mathfrak q \subset \mathfrak q_j$
（事实上二者相等），这意味着 $M'$ 的所有相伴素理想都不是嵌入的。

\medskip\noindent
设 $f$ 是包含于所有 $\mathfrak p_i$ 中的元素。
则 $D(f) \cap \text{supp}(K) = 0$。因此
$M_f = M'_f$，因为 $K_f = 0$。
\end{proof}

\begin{lemma}
\label{algebra-lemma-remove-embedded-primes-localize}
设 $R$ 为诺特环，$M$ 为有限 $R$-模。对任意 $f \in R$，有
$(M')_f = (M_f)'$，其中 $M \to M'$ 与
$M_f \to (M_f)'$ 是引理
\ref{algebra-lemma-remove-embedded-primes} 中构造的商。
\end{lemma}

\begin{proof}
证明从略。
\end{proof}

\begin{lemma}
\label{algebra-lemma-no-embedded-primes-endos}
设 $R$ 为诺特环，$M$ 为没有嵌入相伴素理想的有限 $R$-模。
令 $I = \{x \in R \mid xM = 0\}$。则环 $R/I$
没有嵌入素理想。
\end{lemma}

\begin{proof}
可以将 $R$ 替换为 $R/I$。因此，可以假设 $R$ 的每个非零元素
都在 $M$ 上非平凡地作用。由引理
\ref{algebra-lemma-support-closed}，这意味着 $\Spec(R)$ 等于 $M$ 的支集。
假设 $\mathfrak p$ 是 $R$ 的嵌入素理想。设 $x \in R$
的零化子为 $\mathfrak p$。考虑非零模
$N = xM \subset M$。它被 $\mathfrak p$ 零化。
因此，$N$ 的任意相伴素理想 $\mathfrak q$ 都包含
$\mathfrak p$，并且也是 $M$ 的相伴素理想。
于是 $\mathfrak q$ 将是 $M$ 的嵌入相伴素理想，
这与本引理的假设矛盾。
\end{proof}

\section{正则序列}
\label{algebra-section-regular-sequences}

\noindent
本节研究正则序列的一些基本性质。

\begin{definition}
\label{algebra-definition-regular-sequence}
设 $R$ 为环，$M$ 为 $R$-模。若 $R$ 中的元素序列
$f_1, \ldots, f_r$ 满足下列条件，则称它为
{\it $M$-正则序列}：
\begin{enumerate}
\item 对每个 $i = 1, \ldots, r$，$f_i$ 是
$M/(f_1, \ldots, f_{i - 1})M$ 上的非零因子；并且
\item 模 $M/(f_1, \ldots, f_r)M$ 非零。
\end{enumerate}
如果 $I$ 是 $R$ 的理想且 $f_1, \ldots, f_r \in I$，
则称 $f_1, \ldots, f_r$ 为{\it $I$ 中的 $M$-正则序列}。
如果 $M = R$，则简称 $f_1, \ldots, f_r$ 为
{\it 正则序列}（在 $I$ 中）。
\end{definition}

\noindent
请注意，这个定义依赖于元素 $f_1, \ldots, f_r$ 的次序
（参见下面的例子）。有些论文或书籍去掉了模
$M/(f_1, \ldots, f_r)M$ 非零这一要求。
这样做的优点是，正则序列这一性质在局部化下保持。
然而，我们主要用这个定义来定义 $R$ 为局部环时模的深度；
在这种情形下，要求 $f_i$ 属于极大理想——而从 $R$ 过渡到局部化
$R_\mathfrak p$ 时，这一条件并不保持。

\begin{example}
\label{algebra-example-global-regular}
设 $k$ 为域。在环 $k[x, y, z]$ 中，序列
$x, y(1-x), z(1-x)$ 是正则的，但序列
$y(1-x), z(1-x), x$ 不是。
\end{example}

\begin{example}
\label{algebra-example-local-regular}
设 $k$ 为域。考虑环
$k[x, y, w_0, w_1, w_2, \ldots]/I$，
其中 $I$ 由 $yw_i$（$i = 0, 1, 2, \ldots$）和
$w_i - xw_{i + 1}$（$i = 0, 1, 2, \ldots$）生成。
序列 $x, y$ 是正则的，但 $y$ 是零因子。
此外，还可以在极大理想 $(x, y, w_i)$ 处局部化，
仍得到一个例子。
\end{example}

\begin{lemma}
\label{algebra-lemma-permute-xi}
设 $R$ 为诺特局部环，$M$ 为有限 $R$-模，并设
$x_1, \ldots, x_c$ 为 $M$-正则序列。则 $x_i$ 的任意排列
也为正则序列。
\end{lemma}

\begin{proof}
先处理 $c = 2$ 的情形。考虑
$K \subset M$，即 $x_2 : M \to M$ 的核。
对任意 $z \in K$，由于 $x_2$ 是 $M/x_1M$ 上的非零因子，
存在某个 $z' \in M$，使得 $z = x_1 z'$。
由于 $x_1$ 是 $M$ 上的非零因子，可知 $x_2 z' = 0$。
因此 $x_1 : K \to K$ 是满射。由 Nakayama 引理
\ref{algebra-lemma-NAK}，$K = 0$。
接下来，考虑 $x_1$ 在 $M/x_2M$ 上的乘法。如果 $z \in M$
在此映射的核中映为元素 $\overline{z} \in M/x_2M$，
则对某个 $y \in M$ 有 $x_1 z = x_2 y$。
但由于 $x_1, x_2$ 是正则序列，存在某个 $y' \in M$，
使得 $y = x_1 y'$。所以
$x_1 ( z - x_2 y' ) =0$，进而 $z = x_2 y'$，
故按所需有 $\overline{z} = 0$。

\medskip\noindent
对一般情形，注意任意排列都是相邻指标换位的复合。
因此，只须证明
$$
x_1, \ldots, x_{i-2}, x_i, x_{i-1}, x_{i + 1}, \ldots, x_c
$$
是 $M$-正则序列。这由刚刚处理的情形得出：将它应用于模
$M/(x_1, \ldots, x_{i-2})$ 以及长度为 $2$ 的正则序列
$x_{i-1}, x_i$。
\end{proof}

\begin{lemma}
\label{algebra-lemma-flat-increases-depth}
\begin{slogan}
平坦局部环同态保持并反映正则序列。
\end{slogan}
设 $R, S$ 为局部环，$R \to S$ 为平坦局部环同态。
设 $x_1, \ldots, x_r$ 为 $R$ 中的序列，$M$ 为 $R$-模。
下列条件等价：
\begin{enumerate}
\item $x_1, \ldots, x_r$ 是 $R$ 中的 $M$-正则序列；并且
\item $x_1, \ldots, x_r$ 在 $S$ 中的像构成
$M \otimes_R S$-正则序列。
\end{enumerate}
\end{lemma}

\begin{proof}
这是因为由引理 \ref{algebra-lemma-local-flat-ff}，
$R \to S$ 忠实平坦。
\end{proof}

\begin{lemma}
\label{algebra-lemma-regular-sequence-in-neighbourhood}
设 $R$ 为诺特环，$M$ 为有限 $R$-模，$\mathfrak p$ 为素理想。
设 $x_1, \ldots, x_r$ 为 $R$ 中的序列，它在
$R_{\mathfrak p}$ 中的像构成 $M_{\mathfrak p}$-正则序列。
则存在 $g \in R$，$g \not \in \mathfrak p$，使得
$x_1, \ldots, x_r$ 在 $R_g$ 中的像构成 $M_g$-正则序列。
\end{lemma}

\begin{proof}
令
$$
K_i = \Ker\left(x_i : M/(x_1, \ldots, x_{i - 1})M \to
M/(x_1, \ldots, x_{i - 1})M\right).
$$
这是有限 $R$-模，依假设，它在 $\mathfrak p$ 处的局部化为零。
因此存在 $g \in R$，$g \not \in \mathfrak p$，使得对所有
$i = 1, \ldots, r$ 都有 $(K_i)_g = 0$。这个 $g$ 即满足要求。
\end{proof}

\begin{lemma}
\label{algebra-lemma-join-regular-sequences}
设 $A$ 为环，$I$ 为 $A$ 中由正则序列
$f_1, \ldots, f_n$ 生成的理想。设
$g_1, \ldots, g_m \in A$ 的像
$\overline{g}_1, \ldots, \overline{g}_m$ 在 $A/I$ 中构成正则序列。
则 $f_1, \ldots, f_n, g_1, \ldots, g_m$ 是 $A$ 中的正则序列。
\end{lemma}

\begin{proof}
这立即由定义得出。
\end{proof}

\begin{lemma}
\label{algebra-lemma-regular-sequence-short-exact-sequence}
设 $R$ 为环，$0 \to M_1 \to M_2 \to M_3 \to 0$
为 $R$-模的短正合列，并设 $f_1, \ldots, f_r \in R$。
如果 $f_1, \ldots, f_r$ 既为 $M_1$-正则序列又为
$M_3$-正则序列，则 $f_1, \ldots, f_r$ 也是 $M_2$-正则序列。
\end{lemma}

\begin{proof}
由引理 \ref{algebra-lemma-snake}，如果
$f_1 : M_1 \to M_1$ 和 $f_1 : M_3 \to M_3$ 都是单射，
那么 $f_1 : M_2 \to M_2$ 也是单射，并得到短正合列
$$
0 \to M_1/f_1M_1 \to M_2/f_1M_2 \to M_3/f_1M_3 \to 0
$$
本引理由此及对 $r$ 的归纳得出。一些细节从略。
\end{proof}

\begin{lemma}
\label{algebra-lemma-regular-sequence-powers}
设 $R$ 为环，$M$ 为 $R$-模。设
$f_1, \ldots, f_r \in R$，且 $e_1, \ldots, e_r > 0$ 为整数。
则 $f_1, \ldots, f_r$ 是 $M$-正则序列，当且仅当
$f_1^{e_1}, \ldots, f_r^{e_r}$ 是 $M$-正则序列。
\end{lemma}

\begin{proof}
对 $r$ 作归纳。如果 $r = 1$，结论由下面两个简单事实得出：
(a) $M$ 上非零因子的幂仍为 $M$ 上的非零因子；
(b) $M$ 上非零因子的因子仍为 $M$ 上的非零因子。
如果 $r > 1$，则将归纳假设应用于 $M/f_1M$，可知
$f_1, f_2, \ldots, f_r$ 是 $M$-正则序列，当且仅当
$f_1, f_2^{e_2}, \ldots, f_r^{e_r}$ 是 $M$-正则序列。
所以只须证明：给定 $e > 0$，
$f_1^e, f_2, \ldots, f_r$ 是 $M$-正则序列，当且仅当
$f_1, \ldots, f_r$ 是 $M$-正则序列。
我们对 $e$ 作归纳。$e = 1$ 的情形是平凡的。
在两种假设下，$f_1$ 都是非零因子
（由 $r = 1$ 的情形），所以有短正合列
$$
0 \to M/f_1M \xrightarrow{f_1^{e - 1}} M/f_1^eM \to M/f_1^{e - 1}M \to 0
$$
假设 $f_1, f_2, \ldots, f_r$ 是 $M$-正则序列。
则由归纳假设，元素 $f_2, \ldots, f_r$ 分别构成
$M/f_1M$-正则序列和 $M/f_1^{e - 1}M$-正则序列。
由引理 \ref{algebra-lemma-regular-sequence-short-exact-sequence}，
$f_2, \ldots, f_r$ 是 $M/f_1^eM$-正则序列。因此
$f_1^e, f_2, \ldots, f_r$ 是 $M$-正则序列。
反之，假设 $f_1^e, f_2, \ldots, f_r$ 是 $M$-正则序列。
那么 $f_2 : M/f_1^eM \to M/f_1^eM$ 是单射，故
$f_2 : M/f_1M \to M/f_1M$ 是单射，再由归纳（！），
$f_2 : M/f_1^{e - 1}M \to M/f_1^{e - 1}M$ 是单射，因而
$$
0 \to
M/(f_1, f_2)M \xrightarrow{f_1^{e - 1}}
M/(f_1^e, f_2)M \to
M/(f_1^{e - 1}, f_2)M \to 0
$$
由引理 \ref{algebra-lemma-snake} 是短正合列。这证明了 $r = 2$ 时的逆命题。
如果 $r > 2$，那么
$f_3 : M/(f_1^e, f_2)M \to M/(f_1^e, f_2)M$ 是单射，故
$f_3 : M/(f_1, f_2)M \to M/(f_1, f_2)M$ 是单射，依此类推。
一些细节从略。
\end{proof}

\begin{lemma}
\label{algebra-lemma-regular-sequence-in-polynomial-ring}
设 $R$ 为环，并设不生成单位理想的
$f_1, \ldots, f_r \in R$。下列条件等价：
\begin{enumerate}
\item $f_1, \ldots, f_r$ 的任意排列都是正则序列；
\item $f_1, \ldots, f_r$ 的任意子序列
（按给定次序）都是正则序列；以及
\item 在多项式环 $R[x_1, \ldots, x_r]$ 中，
$f_1x_1, \ldots, f_rx_r$ 是正则序列。
\end{enumerate}
\end{lemma}

\begin{proof}
显然 (1) 蕴含 (2)。下面对 $r$ 作归纳，证明 (2) 蕴含 (1)。
$r = 1$ 的情形是平凡的。$r = 2$ 的情形断言：如果
$a, b \in R$ 是正则序列且 $b$ 是非零因子，那么
$b, a$ 是正则序列。这是清楚的，因为如果 $a$ 与 $b$
都是非零因子，则
$a : R/(b) \to R/(b)$ 的核同构于
$b : R/(a) \to R/(a)$ 的核。
现在考虑 $r > 2$ 的情形。假设 (2) 成立，并设我们要证明
对某个排列 $\sigma$，$f_{\sigma(1)}, \ldots, f_{\sigma(r)}$
是正则序列。由归纳假设，已经知道
$f_{\sigma(1)}, \ldots, f_{\sigma(r - 1)}$ 是正则序列。
因此，只须说明当 $s = \sigma(r)$ 时，$f_s$ 模掉
$f_1, \ldots, \hat f_s, \ldots, f_r$ 后是非零因子。
如果 $s = r$，结论已经成立。如果 $s < r$，则注意
$f_s$ 和 $f_r$ 在环
$R/(f_1, \ldots, \hat f_s, \ldots, f_{r - 1})$
中都是非零因子（再次由归纳假设）。由于已知 $f_s, f_r$
在该环中是正则序列，由长度为 $2$ 的序列的情形可知
$f_r, f_s$ 也是正则序列。

\medskip\noindent
注意，作为 $R$-模，
$R[x_1, \ldots, x_r]/(f_1x_1, \ldots, f_ix_i)$
是如下诸模的直和：
$$
R/I_E \cdot x_1^{e_1} \ldots x_r^{e_r}
$$
其指标为多重指标 $E = (e_1, \ldots, e_r)$，其中
$I_E$ 是由满足 $1 \leq j \leq i$ 且 $e_j > 0$ 的 $f_j$
生成的理想。因此，$f_{i + 1}x_i$ 在此模上是非零因子，
当且仅当对所有 $E$，$f_{i + 1}$ 都是 $R/I_E$ 上的非零因子。
取所有分量均为正的 $E$，可知 $f_{i + 1}$ 是
$R/(f_1, \ldots, f_i)$ 上的非零因子。因此 (3) 蕴含 (2)。
反之，如果 (2) 成立，那么
$f_1, \ldots, f_i, f_{i + 1}$ 的任意子序列都是正则序列；
特别地，$f_{i + 1}$ 是所有 $R/I_E$ 上的非零因子。
由此可见 (2) 蕴含 (3)。
\end{proof}

\section{拟正则序列}
\label{algebra-section-quasi-regular}

\noindent
我们引入拟正则序列的概念；它比正则序列稍弱，但更容易使用。
设 $R$ 为环，且 $f_1, \ldots, f_c \in R$。
令 $J = (f_1, \ldots, f_c)$，并设 $M$ 为 $R$-模。
则存在分次 $R/J[X_1, \ldots, X_c]$-模的典范映射
\begin{equation}
\label{algebra-equation-quasi-regular}
M/JM \otimes_{R/J} R/J[X_1, \ldots, X_c]
\longrightarrow
\bigoplus\nolimits_{n \geq 0} J^nM/J^{n + 1}M
\end{equation}
它由如下规则定义：
$$
\overline{m} \otimes X_1^{e_1} \ldots X_c^{e_c} \longmapsto
f_1^{e_1} \ldots f_c^{e_c} m \bmod J^{e_1 + \ldots + e_c + 1}M.
$$
注意，(\ref{algebra-equation-quasi-regular}) 总是满射。

\begin{definition}
\label{algebra-definition-quasi-regular-sequence}
设 $R$ 为环，$M$ 为 $R$-模。如果
(\ref{algebra-equation-quasi-regular}) 是同构，则称 $R$ 的元素序列
$f_1, \ldots, f_c$ 为{\it $M$-拟正则序列}。
如果 $M = R$，则简称 $f_1, \ldots, f_c$ 为
{\it 拟正则序列}。
\end{definition}

\noindent
因此，如果 $f_1, \ldots, f_c$ 是拟正则序列，那么
$$
R/J[X_1, \ldots, X_c] = \bigoplus\nolimits_{n \geq 0} J^n/J^{n + 1}
$$
其中 $J = (f_1, \ldots, f_c)$。显然，拟正则序列这一性质
与 $f_1, \ldots, f_c$ 的次序无关。

\begin{lemma}
\label{algebra-lemma-regular-quasi-regular}
设 $R$ 为环。
\begin{enumerate}
\item $R$ 的正则序列 $f_1, \ldots, f_c$ 是拟正则序列。
\item 假设 $M$ 为 $R$-模，且 $f_1, \ldots, f_c$
是 $M$-正则序列。则 $f_1, \ldots, f_c$ 是
$M$-拟正则序列。
\end{enumerate}
\end{lemma}

\begin{proof}
令 $J = (f_1, \ldots, f_c)$。
我们对 $c$ 作归纳，证明第一个断言。须说明：给定任意关系
$\sum_{|I| = n} a_I f^I \in J^{n + 1}$，其中 $a_I \in R$，
实际上对所有多重指标 $I$ 都有 $a_I \in J$。
由于 $J^{n + 1}$ 的任意元素都可写成
$\sum_{|I| = n} b_I f^I$，其中 $b_I \in J$，
用 $a_I - b_I$ 替换 $a_I$ 后，可以假设该关系为
$\sum_{|I| = n} a_I f^I = 0$。可将它改写为
$$
\sum\nolimits_{e = 0}^n
\left(
\sum\nolimits_{|I'| = n - e}
a_{I', e} f^{I'}
\right)
f_c^e
=
0
$$
此处及下文中，“带撇”的多重指标 $I'$ 都要求具有形式
$I' = (i_1, \ldots, i_{c - 1}, 0)$。
我们将对 $l \in \{0, \ldots, n\}$ 作归纳，证明：如果有关系
$$
\sum\nolimits_{e = 0}^l
\left(
\sum\nolimits_{|I'| = n - e}
a_{I', e} f^{I'}
\right)
f_c^e
=
0
$$
那么对所有 $I', e$ 都有 $a_{I', e} \in J$。
事实上，令 $J' = (f_1, \ldots, f_{c-1})$。
注意
$\sum\nolimits_{|I'| = n - l} a_{I', l} f^{I'}$
在乘以 $f_c^{l}$ 后映入 $(J')^{n - l + 1}$。
由归纳假设（对 $c$ 的归纳），可知
$f_c^l a_{I', l} \in J'$。
由于 $f_1, \ldots, f_c$ 是正则序列，$f_c$
不是 $R/J'$ 上的零因子，所以
$a_{I', l} \in J'$。这使我们可以将项
$(\sum\nolimits_{|I'| = n - l} a_{I', l} f^{I'})f_c^l$
改写成
$(\sum\nolimits_{|I'| = n - l + 1} f_c b_{I', l - 1}
f^{I'})f_c^{l-1}$。于是得到如下形式的新关系：
$$
\left(\sum\nolimits_{|I'| = n - l + 1}
(a_{I', l-1} + f_c b_{I', l - 1}) f^{I'}\right)f_c^{l-1}
+
\sum\nolimits_{e = 0}^{l - 2}
\left(
\sum\nolimits_{|I'| = n - e}
a_{I', e} f^{I'}
\right)
f_c^e
=
0
$$
现在由归纳假设（这次是对 $l$ 的归纳），可知所有
$a_{I', l-1} + f_c b_{I', l - 1} \in J$，并且当
$e \leq l - 2$ 时所有 $a_{I', e} \in J$。
这与上面得到的 $a_{I', l} \in J' \subset J$ 合起来，
完成归纳步骤的证明。

\medskip\noindent
第二个断言意味着：给定任意形式表达式
$F = \sum_{|I| = n} m_I X^I$，其中 $m_I \in M$，且
$\sum m_I f^I \in J^{n + 1}M$，则所有系数 $m_I$ 都属于
$J$。其证明与上面对第一个断言的相应结果的证明完全相同。
\end{proof}

\begin{lemma}
\label{algebra-lemma-flat-base-change-quasi-regular}
设 $R \to R'$ 为平坦环映射，$M$ 为 $R$-模。
假设 $f_1, \ldots, f_r \in R$ 构成 $M$-拟正则序列。
则 $f_1, \ldots, f_r$ 在 $R'$ 中的像构成
$M \otimes_R R'$-拟正则序列。
\end{lemma}

\begin{proof}
令 $J = (f_1, \ldots, f_r)$、$J' = JR'$ 且
$M' = M \otimes_R R'$。须证明典范映射
$\mu : R'/J'[X_1, \ldots X_r] \otimes_{R'/J'} M'/J'M' \to
\bigoplus (J')^nM'/(J')^{n + 1}M'$ 是同构。
因为 $R \to R'$ 平坦，正合列
$0 \to J^nM \to M$ 和
$0 \to J^{n + 1}M \to J^nM \to J^nM/J^{n + 1}M \to 0$
在与 $R'$ 张量后仍正合。这首先蕴含
$J^nM \otimes_R R' = (J')^nM'$，继而蕴含
$(J')^nM'/(J')^{n + 1}M' = J^nM/J^{n + 1}M \otimes_R R'$。
因此，$\mu$ 是 (\ref{algebra-equation-quasi-regular}) 与
$\text{id}_{R'}$ 的张量积；依假设前者为同构，结论成立。
\end{proof}

\begin{lemma}
\label{algebra-lemma-quasi-regular-sequence-in-neighbourhood}
设 $R$ 为诺特环，$M$ 为有限 $R$-模，$\mathfrak p$ 为素理想。
设 $x_1, \ldots, x_c$ 为 $R$ 中的序列，它在
$R_{\mathfrak p}$ 中的像构成
$M_{\mathfrak p}$-拟正则序列。则存在
$g \in R$，$g \not \in \mathfrak p$，使得
$x_1, \ldots, x_c$ 在 $R_g$ 中的像构成
$M_g$-拟正则序列。
\end{lemma}

\begin{proof}
考虑映射 (\ref{algebra-equation-quasi-regular}) 的核 $K$。
由于 $M/JM \otimes_{R/J} R/J[X_1, \ldots, X_c]$ 是有限
$R/J[X_1, \ldots, X_c]$-模，且
$R/J[X_1, \ldots, X_c]$ 是诺特环，可知 $K$ 也是有限
$R/J[X_1, \ldots, X_c]$-模。选取齐次生成元
$k_1, \ldots, k_t \in K$。依假设，对每个
$i = 1, \ldots, t$，存在 $g_i \in R$，
$g_i \not \in \mathfrak p$，使得 $g_i k_i = 0$。
因此 $g = g_1 \ldots g_t$ 满足要求。
\end{proof}

\begin{lemma}
\label{algebra-lemma-truncate-quasi-regular}
设 $R$ 为环，$M$ 为 $R$-模，并设
$f_1, \ldots, f_c \in R$ 为 $M$-拟正则序列。
对任意 $i$，$\overline{R} = R/(f_1, \ldots, f_i)$ 中的序列
$\overline{f}_{i + 1}, \ldots, \overline{f}_c$
是 $\overline{M} = M/(f_1, \ldots, f_i)M$-拟正则序列。
\end{lemma}

\begin{proof}
只须证明 $i = 1$ 的情形。令
$\overline{J} = (\overline{f}_2, \ldots, \overline{f}_c) \subset \overline{R}$。
则
\begin{align*}
\overline{J}^n\overline{M}/\overline{J}^{n + 1}\overline{M}
& =
(J^nM + f_1M)/(J^{n + 1}M + f_1M) \\
& = J^nM / (J^{n + 1}M + J^nM \cap f_1M).
\end{align*}
因此，为证明本引理，只须说明
$J^{n + 1}M + J^nM \cap f_1M = J^{n + 1}M + f_1J^{n - 1}M$，
因为这将说明
$\bigoplus_{n \geq 0}
\overline{J}^n\overline{M}/\overline{J}^{n + 1}\overline{M}$
是
$\bigoplus_{n \geq 0} J^nM/J^{n + 1}M \cong M/JM[X_1, \ldots, X_c]$
关于 $X_1$ 的商。事实上，有
$J^nM \cap f_1M = f_1J^{n - 1}M$。
确实，如果 $m \not \in J^{n - 1}M$，那么
$f_1m \not \in J^nM$，因为依假设
$\bigoplus J^nM/J^{n + 1}M$ 是多项式代数
$M/J[X_1, \ldots, X_c]$。
\end{proof}

\begin{lemma}
\label{algebra-lemma-quasi-regular-regular}
设 $(R, \mathfrak m)$ 为诺特局部环，$M$ 为非零有限
$R$-模，并设 $f_1, \ldots, f_c \in \mathfrak m$
为 $M$-拟正则序列。则 $f_1, \ldots, f_c$ 是
$M$-正则序列。
\end{lemma}

\begin{proof}
令 $J = (f_1, \ldots, f_c)$。先说明 $f_1$ 是 $M$ 上的非零因子。
假设 $x \in M$ 非零。由 Krull 交定理，存在整数 $r$，
使得 $x \in J^rM$ 但 $x \not \in J^{r + 1}M$；参见引理
\ref{algebra-lemma-intersect-powers-ideal-module-zero}。
于是 $f_1 x \in J^{r + 1}M$，而由所假设的
$\bigoplus J^nM/J^{n + 1}M$ 的结构，它在
$J^{r + 1}M/J^{r + 2}M$ 中的类非零。因此
$f_1x \not = 0$。

\medskip\noindent
现在使用引理
\ref{algebra-lemma-truncate-quasi-regular} 并对 $c$ 作归纳，
即可完成证明。
\end{proof}

\begin{remark}[其他类型的正则序列]
\label{algebra-remark-koszul-regular}
在论文 \cite{Kabele} 中，作者讨论了环 $R$ 的元素序列
$x_1, \ldots, x_r$ 的另外两种正则性条件。具体而言，如果对
$i \geq 1$ 有 $H_i(K_{\bullet}(R, x_{\bullet})) = 0$，
其中 $K_{\bullet}(R, x_{\bullet})$ 是 Koszul 复形，
则称该序列为{\it Koszul-正则的}。如果
$H_1(K_{\bullet}(R, x_{\bullet})) = 0$，则称该序列为
{\it $H_1$-正则的}。有如下蕴含关系：
正则 $\Rightarrow$ Koszul-正则 $\Rightarrow$
$H_1$-正则 $\Rightarrow$ 拟正则。作者通过例子说明，
即使 $R$ 是（非诺特）局部环且该序列生成 $R$ 的极大理想，
这些蕴含一般也不能逆转。我们在《代数进阶》的第
\ref{more-algebra-section-koszul-regular} 节中更详细地引入这些概念。
\end{remark}

\begin{remark}
\label{algebra-remark-join-quasi-regular-sequences}
设 $k$ 为域。考虑环
$$
A = k[x, y, w, z_0, z_1, z_2, \ldots]/
(y^2z_0 - wx, z_0 - yz_1, z_1 - yz_2, \ldots)
$$
在此环中，$x$ 是非零因子，$y$ 在 $A/xA$ 中的像给出拟正则序列。
但 $x, y$ 在 $A$ 中并非拟正则序列，因为
$(x, y)/(x, y)^2$ 不是 $A/(x, y)$ 上秩为二的自由模：
在 $(x, y)/(x, y)^2$ 中有 $wx = 0$，但 $w$ 在
$A/(x, y)$ 中不为零。因此，引理
\ref{algebra-lemma-join-regular-sequences} 的类似断言对拟正则序列不成立。
\end{remark}

\begin{lemma}
\label{algebra-lemma-quasi-regular-on-quotient}
设 $R$ 为环，$J = (f_1, \ldots, f_r)$ 为 $R$ 的理想，
$M$ 为 $R$-模。令
$\overline{R} = R/\bigcap_{n \geq 0} J^n$、
$\overline{M} = M/\bigcap_{n \geq 0} J^nM$，并以
$\overline{f}_i$ 表示 $f_i$ 在 $\overline{R}$ 中的像。
则 $f_1, \ldots, f_r$ 是 $M$-拟正则序列，当且仅当
$\overline{f}_1, \ldots, \overline{f}_r$ 是
$\overline{M}$-拟正则序列。
\end{lemma}

\begin{proof}
这是因为
$J^nM/J^{n + 1}M \cong
\overline{J}^n\overline{M}/\overline{J}^{n + 1}\overline{M}$。
\end{proof}

\section{爆破代数}
\label{algebra-section-blow-up}

\noindent
本节对爆破作一些初等观察。

\begin{definition}
\label{algebra-definition-blow-up}
设 $R$ 为环。
设 $I \subset R$ 为理想。
\begin{enumerate}
\item 与二元组 $(R, I)$ 相关联的{\it 爆破代数}，或称
{\it Rees 代数}，是分次 $R$-代数
$$
\text{Bl}_I(R) =
\bigoplus\nolimits_{n \geq 0} I^n =
R \oplus I \oplus I^2 \oplus \ldots
$$
其中直和项 $I^n$ 置于次数 $n$。
\item 设 $a \in I$。以 $a^{(1)}$ 表示把 $a$ 看作 Rees 代数中
次数为 $1$ 的元素所得的元素。则{\it 仿射爆破代数}
$R[\frac{I}{a}]$ 是第 \ref{algebra-section-proj} 节所构造的代数
$(\text{Bl}_I(R))_{(a^{(1)})}$。
\end{enumerate}
\end{definition}

\noindent
换言之，$R[\frac{I}{a}]$ 的元素可由形如 $x/a^n$ 的表达式表示，
其中 $x \in I^n$。两个代表元 $x/a^n$ 与 $y/a^m$ 定义同一元素，
当且仅当存在某个 $k \geq 0$，使得
$a^k(a^mx - a^ny) = 0$。

\begin{lemma}
\label{algebra-lemma-affine-blowup}
设 $R$ 为环，$I \subset R$ 为理想，且 $a \in I$。
令 $R' = R[\frac{I}{a}]$ 为仿射爆破代数。则
\begin{enumerate}
\item $a$ 在 $R'$ 中的像是非零因子；
\item $IR' = aR'$；并且
\item $(R')_a = R_a$。
\end{enumerate}
\end{lemma}

\begin{proof}
这直接由上面对 $R[\frac{I}{a}]$ 的描述得到。
\end{proof}

\begin{lemma}
\label{algebra-lemma-blowup-base-change}
设 $R \to S$ 为环映射。设 $I \subset R$ 为理想且
$a \in I$。置 $J = IS$，并令 $b \in J$ 为 $a$ 的像。
则 $S[\frac{J}{b}]$ 是 $S \otimes_R R[\frac{I}{a}]$
关于由 $b$ 的某个幂零化的元素所成理想的商。
\end{lemma}

\begin{proof}
令 $S'$ 为 $S \otimes_R R[\frac{I}{a}]$ 关于其
$b$-幂挠元所成理想的商。环映射
$$
S \otimes_R R[\textstyle{\frac{I}{a}}]
\longrightarrow
S[\textstyle{\frac{J}{b}}]
$$
是满射；又因为 $b$ 在 $S[\frac{J}{b}]$ 中是非零因子，
它将 $a$-幂挠元映为零。因此得到满射 $S' \to S[\frac{J}{b}]$。
为说明其核为零，我们构造逆映射。具体地，设
$z = y/b^n$ 是 $S[\frac{J}{b}]$ 的元素，即 $y \in J^n$。
写成 $y = \sum x_is_i$，其中 $x_i \in I^n$ 且 $s_i \in S$。
把 $z$ 映到 $S'$ 中 $\sum s_i \otimes x_i/a^n$ 的类。
这是良定义的，因为映射 $S \otimes_R I^n \to J^n$ 的核中元素
由 $a^n$ 零化，因而在 $S'$ 中映为零。
\end{proof}

\begin{example}
\label{algebra-example-rees-algebra-polynomial}
设 $R$ 为环。设 $P = R[t_1, \ldots, t_n]$ 为多项式代数。
令 $I = (t_1, \ldots, t_n) \subset P$。沿用定义
\ref{algebra-definition-blow-up} 的记号，有同构
$$
P[T_1, \ldots, T_n]/(t_iT_j - t_jT_i) \longrightarrow \text{Bl}_I(P)
$$
把 $T_i$ 送到 $t_i^{(1)}$。由读者验证此映射是良定义的。
由于 $I$ 由 $t_1, \ldots, t_n$ 生成，可知该映射是满射。
为说明该映射是单射，需要证明：对每个 $e \geq 1$，
$P$-模 $I^e$ 由多重指标 $E = (e_1, \ldots, e_n)$ 中
次数 $|E| = e$ 的那些指标所对应的单项式
$t^E = t_1^{e_1} \ldots x_n^{e_n}$ 生成，而且它们只满足关系
$t_i t^E = t_j t^{E'}$，其中 $|E| = |E'| = e$ 且
$e_a + \delta_{a i} = e'_a + \delta_{a j},\ a = 1, \ldots, n$
（Kronecker delta）。细节从略。
\end{example}

\begin{example}
\label{algebra-example-affine-blowup-algebra-polynomial}
设 $R$ 为环。设 $P = R[t_1, \ldots, t_n]$ 为多项式代数。
令 $I = (t_1, \ldots, t_n) \subset P$，并令 $a = t_1$。
沿用定义 \ref{algebra-definition-blow-up} 的记号，有同构
$$
P[x_2, \ldots, x_n]/(t_1x_2 - t_2, \ldots, t_1x_n - t_n)
\longrightarrow
\textstyle{P[\frac{I}{a}] = P[\frac{I}{t_1}]}
$$
把 $x_i$ 送到 $t_i/t_1$。由读者验证此映射是良定义的。
由于 $I$ 由 $t_1, \ldots, t_n$ 生成，可知该映射是满射。
为说明该映射是单射，读者可以论证其源与靶都没有
$t_1$-挠元，并且在逆转 $t_1$ 后该映射成为同构；参见引理
\ref{algebra-lemma-affine-blowup}。或者，读者可使用例
\ref{algebra-example-rees-algebra-polynomial} 中对 Rees 代数的描述。
细节从略。
\end{example}

\begin{lemma}
\label{algebra-lemma-affine-blowup-quotient-description}
设 $R$ 为环，$I = (a_1, \ldots, a_n)$ 为 $R$ 的理想。
令 $a = a_1$。则存在满射
$$
R[x_2, \ldots, x_n]/(a x_2 - a_2, \ldots, a x_n - a_n)
\longrightarrow
\textstyle{R[\frac{I}{a}]}
$$
其核是源中的 $a$-幂挠元之集。
\end{lemma}

\begin{proof}
考虑环映射 $P = \mathbf{Z}[t_1, \ldots, t_n] \to R$，
它把 $t_i$ 送到 $a_i$。置 $J = (t_1, \ldots, t_n)$。
由例 \ref{algebra-example-affine-blowup-algebra-polynomial}，有
$P[\frac{J}{t_1}] =
P[x_2, \ldots, x_n]/(t_1 x_2 - t_2, \ldots, t_1 x_n - t_n)$。
将引理 \ref{algebra-lemma-blowup-base-change} 应用于映射 $P \to A$ 即得结论。
\end{proof}

\begin{lemma}
\label{algebra-lemma-blowup-in-principal}
设 $R$ 为环，$I \subset R$ 为理想，且 $a \in I$。
置 $R' = R[\frac{I}{a}]$。若 $f \in R$ 满足 $V(f) = V(I)$，
则 $f$ 在 $R'$ 中的像是非零因子，并且
$R'_f = R'_a = R_a$。
\end{lemma}

\begin{proof}
下文不再特别说明地使用引理 \ref{algebra-lemma-affine-blowup} 的结论。
假设 $V(f) = V(I)$ 蕴含 $V(fR') = V(IR') = V(aR')$。
因此存在 $b, c \in R'$ 和某些幂，使得
$a^n = fb$ 且 $f^m = ac$。引理随即成立。
\end{proof}

\begin{lemma}
\label{algebra-lemma-blowup-add-principal}
设 $R$ 为环，$I \subset R$ 为理想，$a \in I$，且 $f \in R$。
置 $R' = R[\frac{I}{a}]$ 与 $R'' = R[\frac{fI}{fa}]$。
则存在满的 $R$-代数映射 $R' \to R''$，其核为 $R'$ 中
$f$-幂挠元之集。
\end{lemma}

\begin{proof}
该映射把 $x \in I^n$ 所对应的 $x/a^n$ 送到
$f^nx/(fa)^n$。直接验证可知此映射良定义且为满射。
由于 $af$ 在 $R''$ 中是非零因子
（引理 \ref{algebra-lemma-affine-blowup}），可知 $f$-幂挠元之集映为零。
反之，若 $x \in R'$ 且对所有 $n > 0$ 都有
$f^n x \not = 0$，则由于 $a$ 在 $R'$ 中是非零因子，
对所有 $n$ 都有 $(af)^n x \not = 0$。根据定义
\ref{algebra-definition-blow-up} 之后对 $R''$ 的描述，$x$ 在 $R''$
中的像不为零。
\end{proof}

\begin{lemma}
\label{algebra-lemma-blowup-reduced}
\begin{slogan}
约化性在爆破下不变
\end{slogan}
若 $R$ 约化，则 $R$ 的每个（仿射）爆破代数都是约化的。
\end{lemma}

\begin{proof}
设 $I \subset R$ 为理想且 $a \in I$。假设 $x \in I^n$
所对应的 $x/a^n$ 是 $R[\frac{I}{a}]$ 的幂零元。
则 $(x/a^n)^m = 0$。故对某个 $N \geq 0$，在 $R$ 中有
$a^N x^m = 0$。必要时增大 $N$，可设 $N = me$，其中
$e \geq 0$。于是 $(a^e x)^m = 0$；由于 $R$ 约化，得到
$a^e x = 0$。这说明 $R[\frac{I}{a}]$ 中 $x/a^n = 0$。
\end{proof}

\begin{lemma}
\label{algebra-lemma-blowup-domain}
设 $R$ 为整环，$I \subset R$ 为理想，且 $a \in I$ 为非零元。
则仿射爆破代数 $R[\frac{I}{a}]$ 是整环。
\end{lemma}

\begin{proof}
假设 $x \in I^n$、$y \in I^m$ 所对应的
$x/a^n$、$y/a^m$ 是 $R[\frac{I}{a}]$ 中乘积为零的元素。
则在 $R$ 中有 $a^N x y = 0$。由于 $R$ 是整环，可知
$x = 0$ 或 $y = 0$。
\end{proof}

\begin{lemma}
\label{algebra-lemma-blowup-dominant}
设 $R$ 为环，$I \subset R$ 为理想，且 $a \in I$。
若 $a$ 不属于 $R$ 的任何极小素理想，则
$\Spec(R[\frac{I}{a}]) \to \Spec(R)$ 的像稠密。
\end{lemma}

\begin{proof}
若对 $x \in R$ 有 $a^k x = 0$，则 $x$ 属于 $R$ 的所有
极小素理想，因而是幂零元；参见引理
\ref{algebra-lemma-Zariski-topology}。所以 $R \to R[\frac{I}{a}]$
的核由幂零元组成。于是结论由引理
\ref{algebra-lemma-image-dense-generic-points} 得到。
\end{proof}

\begin{lemma}
\label{algebra-lemma-valuation-ring-colimit-affine-blowups}
设 $(R, \mathfrak m)$ 为分式域是 $K$ 的局部整环。
设 $R \subset A \subset K$ 为支配 $R$ 的赋值环。则
$$
A = \colim R[\textstyle{\frac{I}{a}}]
$$
是仿射爆破 $R \to R[\frac{I}{a}]$ 的有向余极限，
其中各仿射爆破满足下列性质：
\begin{enumerate}
\item $a \in I \subset \mathfrak m$；
\item $I$ 有限生成；并且
\item $R \to R[\frac{I}{a}]$ 在 $\mathfrak m$ 处的纤维环非零。
\end{enumerate}
\end{lemma}

\begin{proof}
由引理 \ref{algebra-lemma-blowup-domain}，任意爆破代数
$R[\frac{I}{a}]$ 都是包含于 $K$ 的整环。本引理就是说，
$A$ 是其中满足 $a \in I$ 及性质 (1)、(2)、(3) 的那些环的有向并。
若 $R[\frac{I}{a}] \subset A$ 且 $R[\frac{J}{b}] \subset A$，则
$$
R[\textstyle{\frac{I}{a}}] \cup R[\textstyle{\frac{J}{b}}]  \subset
R[\textstyle{\frac{IJ}{ab}}] \subset A
$$
第一个包含成立，是因为 $x/a^n = b^nx/(ab)^n$；第二个包含成立，
是因为若 $z \in (IJ)^n$，则可写成 $z = \sum x_iy_i$，其中
$x_i \in I^n$ 且 $y_i \in J^n$，从而
$z/(ab)^n = \sum (x_i/a^n)(y_i/b^n)$ 属于 $A$。

\medskip\noindent
考虑有限子集 $E \subset A$。设 $E = \{e_1, \ldots, e_n\}$。
选取非零的 $a \in R$，使得对所有 $i = 1, \ldots, n$，
均可写成 $e_i = f_i/a$。置 $I = (f_1, \ldots, f_n, a)$。
我们断言 $R[\frac{I}{a}] \subset A$。这是显然的，因为
$R[\frac{I}{a}]$ 的元素可表示为 $e_i$ 的多项式。
本引理立即由此观察得到。
\end{proof}

\section{Ext 群}
\label{algebra-section-ext}

\noindent
本节将使用一点同调代数，以便建立诺特局部环上深度的一些基本性质。

\begin{lemma}
\label{algebra-lemma-resolution-by-finite-free}
设 $R$ 为环，$M$ 为 $R$-模。
\begin{enumerate}
\item 存在正合复形
$$
\ldots \to F_2 \to F_1 \to F_0 \to M \to 0.
$$
其中 $F_i$ 都是自由 $R$-模。
\item 若 $R$ 是诺特环且 $M$ 是有限 $R$-模，则可以选取这个复形，
使得 $F_i$ 都是有限自由模。换言之，可以找到正合复形
$$
\ldots \to R^{\oplus n_2} \to R^{\oplus n_1} \to R^{\oplus n_0} \to M \to 0.
$$
\end{enumerate}
\end{lemma}

\begin{proof}
只说明诺特情形。第一步选取满射 $R^{n_0} \to M$。
若已构造出长度为 $e$ 的正合复形，只须选取满射
$R^{n_{e + 1}} \to
\Ker(R^{n_e} \to R^{n_{e-1}})$；由于 $R$ 是诺特环，这是可行的。
\end{proof}

\begin{definition}
\label{algebra-definition-finite-free-resolution}
设 $R$ 为环，$M$ 为 $R$-模。
\begin{enumerate}
\item $M$ 的一个（左）{\it 消解} $F_\bullet \to M$ 是 $R$-模的正合复形
$$
\ldots \to F_2 \to F_1 \to F_0 \to M \to 0
$$
\item $M$ 的一个{\it 自由 $R$-模消解}是消解
$F_\bullet \to M$，其中每个 $F_i$ 都是自由 $R$-模。
\item $M$ 的一个{\it 有限自由 $R$-模消解}是消解
$F_\bullet \to M$，其中每个 $F_i$ 都是有限自由 $R$-模。
\end{enumerate}
\end{definition}

\noindent
我们常用记号 $F_{\bullet}$ 表示形如
$$
\ldots \to F_i \to F_{i-1} \to \ldots
$$
的 $R$-模复形。此时常用 $d_i$ 或 $d_{F, i}$ 表示映射
$F_i \to F_{i-1}$。本节始终假设 $F_0$ 是复形中最后一个非零项。
复形 $F_{\bullet}$ 的{\it 第 $i$ 个同调群}是群
$H_i = \Ker(d_{F, i})/\Im(d_{F, i + 1})$。
一个{\it 复形映射 $\alpha : F_{\bullet} \to G_{\bullet}$}
由映射 $\alpha_i : F_i \to G_i$ 给出，并且满足
$\alpha_{i-1} \circ d_{F, i} = d_{G, i-1} \circ \alpha_i$。
这样的映射在同调上诱导映射 $H_i(\alpha) :
H_i(F_{\bullet}) \to H_i(G_{\bullet})$。若 $\alpha, \beta
:  F_{\bullet} \to G_{\bullet}$ 是复形映射，则 $\alpha$ 与
$\beta$ 之间的一个{\it 同伦}由一族映射
$h_i : F_i \to G_{i + 1}$ 给出，并且满足
$\alpha_i - \beta_i = d_{G, i + 1} \circ h_i +
h_{i-1} \circ d_{F, i}$。若 $\alpha$ 与 $\beta$ 之间存在同伦，
则称两个映射 $\alpha, \beta : F_{\bullet} \to G_{\bullet}$
{\it 同伦}。

\medskip\noindent
对于形如 $F^{\bullet}$ 的复形，我们将使用非常类似的记号；它形如
$$
\ldots \to F^i \xrightarrow{d^i} F^{i + 1} \to \ldots
$$
同样有复形映射、同伦等概念。在这种情形下，置
$H^i(F^{\bullet}) =
\Ker(d^i)/\Im(d^{i - 1})$，并称之为{\it 第 $i$ 个上同调群}。

\begin{lemma}
\label{algebra-lemma-homotopic-equal-homology}
任意两个同伦的复形映射在（上）同调群上诱导相同的映射。
\end{lemma}

\begin{proof}
从略。
\end{proof}

\begin{lemma}
\label{algebra-lemma-compare-resolutions}
设 $R$ 为环，$M \to N$ 为 $R$-模映射。
设 $N_\bullet \to N$ 为任意消解。设
$$
\ldots \to F_2 \to F_1 \to F_0 \to M
$$
为 $R$-模复形，其中每个 $F_i$ 都是自由 $R$-模。则
\begin{enumerate}
\item 存在复形映射 $F_\bullet \to N_\bullet$，使得
$$
\xymatrix{
F_0 \ar[r] \ar[d] & M \ar[d] \\
N_0 \ar[r] & N
}
$$
交换；并且
\item 任意两个如 (1) 所述的映射
$\alpha, \beta : F_\bullet \to N_\bullet$ 都同伦。
\end{enumerate}
\end{lemma}

\begin{proof}
证明 (1)。由于 $F_0$ 自由，可以找到提升映射
$F_0 \to M \to N$ 的映射 $F_0 \to N_0$。于是得到诱导映射
$F_1 \to F_0 \to N_0$，其像落在 $N_1 \to N_0$ 的像中。
由于 $F_1$ 自由，可以把它提升为映射 $F_1 \to N_1$。
这又诱导映射 $F_2 \to F_1 \to N_1$，它在 $N_0$ 中映为零。
由于 $N_\bullet$ 正合，该映射的像包含于 $N_2 \to N_1$ 的像中。
所以可以提升得到映射 $F_2 \to N_2$。如此重复。

\medskip\noindent
证明 (2)。为说明 $\alpha, \beta$ 同伦，只须说明差
$\gamma = \alpha - \beta$ 与零同伦。注意到
$\gamma_0 : F_0 \to N_0$ 的像包含于 $N_1 \to N_0$ 的像中。
因此可将 $\gamma_0$ 提升为映射 $h_0 : F_0 \to N_1$。
考虑映射 $\gamma_1' = \gamma_1 - h_0 \circ d_{F, 1}$。
由 $h_0$ 的选取，$\gamma_1'$ 的像包含于
$N_1 \to N_0$ 的核中。由于 $N_\bullet$ 正合，可以把
$\gamma_1'$ 提升为映射 $h_1 : F_1 \to N_2$。
至此有 $\gamma_1 = h_0 \circ d_{F, 1}
+ d_{N, 2} \circ h_1$。如此重复。
\end{proof}

\noindent
现在可以定义群 $\Ext^i_R(M, N)$。具体而言，选取 $M$ 的一个
自由 $R$-模消解 $F_{\bullet}$；参见引理
\ref{algebra-lemma-resolution-by-finite-free}。考虑（上同调）复形
$$
\Hom_R(F_\bullet, N) :
\Hom_R(F_0, N) \to
\Hom_R(F_1, N) \to
\Hom_R(F_2, N) \to \ldots
$$
对 $i \geq 0$，定义 $\Ext^i_R(M, N)$ 为该复形的第 $i$ 个
上同调群\footnote{此时也许说“一个”而不是“这个”Ext 群更为恰当。}。
对 $i < 0$，置 $\Ext^i_R(M, N) = 0$。在继续之前，先指出
$$
\Ext^0_R(M, N) = \Ker(\Hom_R(F_0, N) \to \Hom_R(F_1, N)) =
\Hom_R(M, N)
$$
因为可以把引理 \ref{algebra-lemma-hom-exact} 的 (1) 应用于正合序列
$F_1 \to F_0 \to M \to 0$。下一个引理说明此定义在何种意义下
是良定义的。

\begin{lemma}
\label{algebra-lemma-ext-welldefined}
设 $R$ 为环，$M_1, M_2, N$ 为 $R$-模。
假设 $F_{\bullet}$ 是模 $M_1$ 的自由消解，
$G_{\bullet}$ 是模 $M_2$ 的自由消解。
设 $\varphi : M_1 \to M_2$ 为模映射。
设 $\alpha : F_{\bullet} \to G_{\bullet}$ 为诱导
$M_1 = \Coker(d_{F, 1}) \to M_2 = \Coker(d_{G, 1})$ 上
$\varphi$ 的复形映射；参见引理 \ref{algebra-lemma-compare-resolutions}。
则诱导映射
$$
H^i(\alpha) :
H^i(\Hom_R(F_{\bullet}, N))
\longrightarrow
H^i(\Hom_R(G_{\bullet}, N))
$$
与 $\alpha$ 的选择无关。若 $\varphi$ 是同构，则所有映射
$H^i(\alpha)$ 也都是同构。若 $M_1 = M_2$、
$F_\bullet = G_\bullet$，且 $\varphi$ 是恒等映射，
则所有映射 $H_i(\alpha)$ 也都是恒等映射。
\end{lemma}

\begin{proof}
由引理 \ref{algebra-lemma-compare-resolutions}，另一个诱导 $\varphi$ 的映射
$\beta : F_{\bullet} \to G_{\bullet}$ 与 $\alpha$ 同伦。
因此映射 $\Hom_R(F_\bullet, N) \to
\Hom_R(G_\bullet, N)$ 同伦。于是独立性结论由引理
\ref{algebra-lemma-homotopic-equal-homology} 得到。

\medskip\noindent
假设 $\varphi$ 是同构。设 $\psi : M_2 \to M_1$ 为其逆。
选取诱导
$\psi : M_2 = \Coker(d_{G, 1}) \to M_1 = \Coker(d_{F, 1})$
的映射 $\beta : G_{\bullet} \to F_{\bullet}$；参见引理
\ref{algebra-lemma-compare-resolutions}。好，现在考虑映射
$H^i(\alpha) \circ H^i(\beta) =
H^i(\alpha \circ \beta)$。由上文，映射
$H^i(\alpha \circ \beta)$ 与映射
$H^i(\text{id}_{G_{\bullet}}) = \text{id}$ {\it 相同}。
复合 $H^i(\beta) \circ H^i(\alpha)$ 的情形也类似。
因此 $H^i(\alpha)$ 与 $H^i(\beta)$ 互为逆映射。
\end{proof}

\begin{lemma}
\label{algebra-lemma-long-exact-seq-ext}
设 $R$ 为环，$M$ 为 $R$-模。
设 $0 \to N' \to N \to N'' \to 0$ 为短正合序列。
则得到长正合序列
$$
\begin{matrix}
0
\to \Hom_R(M, N')
\to \Hom_R(M, N)
\to \Hom_R(M, N'')
\\
\phantom{0\ }
\to \Ext^1_R(M, N')
\to \Ext^1_R(M, N)
\to \Ext^1_R(M, N'')
\to \ldots
\end{matrix}
$$
\end{lemma}

\begin{proof}
选取自由消解 $F_{\bullet} \to M$。由于每个 $F_i$ 都自由，
得到复形的短正合序列
$$
0 \to
\Hom_R(F_{\bullet}, N') \to
\Hom_R(F_{\bullet}, N) \to
\Hom_R(F_{\bullet}, N'') \to
0
$$
将蛇引理应用于此序列，即得上述长正合序列。
\end{proof}

\begin{lemma}
\label{algebra-lemma-reverse-long-exact-seq-ext}
设 $R$ 为环，$N$ 为 $R$-模。
设 $0 \to M' \to M \to M'' \to 0$ 为短正合序列。
则得到长正合序列
$$
\begin{matrix}
0
\to \Hom_R(M'', N)
\to \Hom_R(M, N)
\to \Hom_R(M', N)
\\
\phantom{0\ }
\to \Ext^1_R(M'', N)
\to \Ext^1_R(M, N)
\to \Ext^1_R(M', N)
\to \ldots
\end{matrix}
$$
\end{lemma}

\begin{proof}
分别选取 $M'$ 与 $M''$ 的生成元集
$\{m'_{i'}\}_{i' \in I'}$ 和 $\{m''_{i''}\}_{i'' \in I''}$。
对每个 $i'' \in I''$，选取元素 $m''_{i''} \in M''$ 的一个提升
$\tilde m''_{i''} \in M$。置 $F' = \bigoplus_{i' \in I'} R$、
$F'' = \bigoplus_{i'' \in I''} R$ 以及 $F = F' \oplus F''$。
把这些自由模的生成元映到相应选定的生成元，得到满的
$R$-模映射 $F' \to M'$、$F'' \to M''$ 和 $F \to M$。
于是得到短正合序列之间的映射
$$
\begin{matrix}
0 & \to & M' & \to & M & \to & M'' & \to & 0 \\
& & \uparrow & & \uparrow & & \uparrow \\
0 & \to & F' & \to & F & \to & F'' & \to & 0 \\
\end{matrix}
$$
由蛇引理，核的序列 $0 \to K' \to K \to K'' \to 0$
是 $R$-模的短正合序列。因此可以无限地继续这一过程。
换言之，我们得到消解的短正合序列，它嵌入图表
$$
\begin{matrix}
0 & \to & M' & \to & M & \to & M'' & \to & 0 \\
& & \uparrow & & \uparrow & & \uparrow \\
0 & \to & F_\bullet' & \to & F_\bullet & \to & F_\bullet'' & \to & 0 \\
\end{matrix}
$$
由于每个序列 $0 \to F'_n \to F_n \to F''_n \to 0$
依构造都是分裂正合的，应用函子 $\Hom_R(-, N)$ 后，
得到复形的短正合序列
$$
0 \to
\Hom_R(F''_{\bullet}, N) \to
\Hom_R(F_{\bullet}, N) \to
\Hom_R(F'_{\bullet}, N) \to
0
$$
将蛇引理应用于此序列，即得上述长正合序列。
\end{proof}

\begin{lemma}
\label{algebra-lemma-annihilate-ext}
设 $R$ 为环，$M$, $N$ 为 $R$-模。
若 $x\in R$ 满足 $xN = 0$ 或 $xM = 0$，
则该元素零化每个模 $\Ext^i_R(M, N)$。
\end{lemma}

\begin{proof}
选取 $M$ 的自由消解 $F_{\bullet}$。由于 $\Ext^i_R(M, N)$
定义为复形 $\Hom_R(F_{\bullet}, N)$ 的上同调，
当 $xN = 0$ 时，引理显然成立。若 $xM = 0$，则
$F_{\bullet}$ 上的乘以 $x$ 提升了 $M$ 上的零映射。
故由引理 \ref{algebra-lemma-ext-welldefined}，它在 Ext 群上诱导的映射
与零映射诱导的映射相同。
\end{proof}

\begin{lemma}
\label{algebra-lemma-ext-noetherian}
设 $R$ 为诺特环，$M$, $N$ 为有限 $R$-模。
则对所有 $i$，$\Ext^i_R(M, N)$ 都是有限 $R$-模。
\end{lemma}

\begin{proof}
这是因为 $\Ext^i_R(M, N)$ 可计算为复形
$\Hom_R(F_\bullet, N)$ 的上同调群，其中每个 $F_n$
都是有限自由 $R$-模；参见引理
\ref{algebra-lemma-resolution-by-finite-free}。
\end{proof}

\section{深度}
\label{algebra-section-depth}

\noindent
定义如下。

\begin{definition}
\label{algebra-definition-depth}
设 $R$ 为环，$I \subset R$ 为理想，$M$ 为有限 $R$-模。
$M$ 的{\it $I$-深度}，记作 $\text{depth}_I(M)$，定义如下：
\begin{enumerate}
\item 若 $IM \not = M$，则 $\text{depth}_I(M)$ 是
$\{0, 1, 2, \ldots, \infty\}$ 中 $I$ 内 $M$-正则序列之长度的上确界；
\item 若 $IM = M$，则置 $\text{depth}_I(M) = \infty$。
\end{enumerate}
若 $(R, \mathfrak m)$ 是局部环，则简称
$\text{depth}_{\mathfrak m}(M)$ 为 $M$ 的{\it 深度}。
\end{definition}

\noindent
说明。依定义 \ref{algebra-definition-regular-sequence}，空序列不是零模上的
正则序列；但就实际应用而言，把 $0$ 模的深度置为 $+\infty$
是方便的。注意，若 $I = R$，则对所有有限 $R$-模 $M$，
都有 $\text{depth}_I(M) = \infty$。若 $I$ 包含于 $R$ 的
Jacobson 根中（例如，若 $R$ 是局部环且
$I \subset \mathfrak m_R$），则由 Nakayama 引理，
$M \not = 0 \Rightarrow IM \not = M$。
模 $M$ 的 $I$-深度为 $0$，当且仅当 $M$ 非零且
$I$ 不含 $M$ 上的非零因子。

\medskip\noindent
例 \ref{algebra-example-global-regular} 表明，即使环是诺特的，深度的表现
也不良好；例 \ref{algebra-example-local-regular} 表明，当环是局部但非诺特时，
其表现也不良好。我们将看到，对诺特局部环，深度的表现良好。

\begin{lemma}
\label{algebra-lemma-depth-weak-sequence}
设 $R$ 为环，$I \subset R$ 为理想，$M$ 为有限 $R$-模。
则 $\text{depth}_I(M)$ 等于满足下列条件的序列
$f_1, \ldots, f_r \in I$ 之长度的上确界：$f_i$ 是
$M/(f_1, \ldots, f_{i - 1})M$ 上的非零因子。
\end{lemma}

\begin{proof}
假设 $IM = M$。则引理 \ref{algebra-lemma-NAK} 表明，存在
$f \in I$，使得 $f : M \to M$ 是 $\text{id}_M$。
所以 $f, 0, 0, 0, \ldots$ 是由相继非零因子组成的无限序列，
由此可见此情形下二者相等。若 $IM \not =  M$，则本引理所述的
序列是 $M$-正则序列，故此情形下二者也相等。
\end{proof}

\begin{lemma}
\label{algebra-lemma-bound-depth}
设 $(R, \mathfrak m)$ 为诺特局部环，$M$ 为非零有限 $R$-模。
则 $\dim(\text{Supp}(M)) \geq \text{depth}(M)$。
\end{lemma}

\begin{proof}
对 $\dim(\text{Supp}(M))$ 作归纳。若
$\dim(\text{Supp}(M)) = 0$，则
$\text{Supp}(M) = \{\mathfrak m\}$，从而
$\text{Ass}(M) = \{\mathfrak m\}$（由引理
\ref{algebra-lemma-ass-support} 与 \ref{algebra-lemma-ass-zero}），
所以例如由引理 \ref{algebra-lemma-ideal-nonzerodivisor}，$M$ 的深度为零。
在归纳步骤中，假设 $\dim(\text{Supp}(M)) > 0$。
设 $f_1, \ldots, f_d$ 为 $\mathfrak m$ 中的元素序列，
且 $f_i$ 是 $M/(f_1, \ldots, f_{i - 1})M$ 上的非零因子。
依引理 \ref{algebra-lemma-depth-weak-sequence}，只须证明
$\dim(\text{Supp}(M)) \geq d$。不妨设 $d > 0$，否则引理显然成立。
由引理 \ref{algebra-lemma-one-equation-module}，有
$\dim(\text{Supp}(M/f_1M)) = \dim(\text{Supp}(M)) - 1$。
由归纳假设得到 $\dim(\text{Supp}(M/f_1M)) \geq d - 1$，正合所需。
\end{proof}

\begin{lemma}
\label{algebra-lemma-depth-finite-noetherian}
设 $R$ 为诺特环，$I \subset R$ 为理想，$M$ 为非零有限
$R$-模，且 $IM \not = M$。则
$\text{depth}_I(M) < \infty$。
\end{lemma}

\begin{proof}
由于 $M/IM$ 非零，由引理 \ref{algebra-lemma-support-zero}，
可以选取 $\mathfrak p \in \text{Supp}(M/IM)$。
于是 $(M/IM)_\mathfrak p \not = 0$，这蕴含
$I \subset \mathfrak p$；又因为局部化正合，还蕴含
$M_\mathfrak p \not = IM_\mathfrak p$。
设 $f_1, \ldots, f_r \in I$ 为 $M$-正则序列。
由于 $(f_1, \ldots, f_r) \subset I$，
$M_\mathfrak p/(f_1, \ldots, f_r)M_\mathfrak p$ 非零。
由于局部化平坦，$f_1, \ldots, f_r$ 的像在
$I_\mathfrak p$ 中构成 $M_\mathfrak p$-正则序列。
这对 $I$ 中每个 $M$-正则序列都成立，故
$\text{depth}_I(M) \leq \text{depth}_{I_\mathfrak p}(M_\mathfrak p)$。
后者的值 $\leq \text{depth}(M_\mathfrak p)$，而由引理
\ref{algebra-lemma-bound-depth}，该值 $< \infty$。
\end{proof}

\begin{lemma}
\label{algebra-lemma-depth-ext}
设 $R$ 为极大理想是 $\mathfrak m$ 的诺特局部环，
$M$ 为非零有限 $R$-模。则 $\text{depth}(M)$ 等于使得
$\Ext^i_R(R/\mathfrak m, M)$ 非零的最小整数 $i$。
\end{lemma}

\begin{proof}
以 $\delta(M)$ 表示 $M$ 的深度，并以 $i(M)$ 表示使得
$\Ext^i_R(R/\mathfrak m, M)$ 非零的最小整数 $i$。
马上会看到 $i(M) < \infty$。由引理
\ref{algebra-lemma-ideal-nonzerodivisor}，有 $\delta(M) = 0$
当且仅当 $i(M) = 0$，因为
$\mathfrak m \in \text{Ass}(M)$ 恰好意味着 $i(M) = 0$。
因此，若 $\delta(M)$ 或 $i(M)$ 的值 $> 0$，则可以选取
$x \in \mathfrak m$，使得 (a) $x$ 是 $M$ 上的非零因子，
且 (b) $\text{depth}(M/xM) = \delta(M) - 1$。
考虑由引理 \ref{algebra-lemma-long-exact-seq-ext} 从短正合序列
$0 \to M \to M \to M/xM \to 0$ 得到的 Ext 群长正合序列：
$$
\begin{matrix}
0
\to \Hom_R(\kappa, M)
\to \Hom_R(\kappa, M)
\to \Hom_R(\kappa, M/xM)
\\
\phantom{0\ }
\to \Ext^1_R(\kappa, M)
\to \Ext^1_R(\kappa, M)
\to \Ext^1_R(\kappa, M/xM)
\to \ldots
\end{matrix}
$$
由于 $x \in \mathfrak m$，所有映射 $\Ext^i_R(\kappa, M)
\to \Ext^i_R(\kappa, M)$ 都为零；参见引理
\ref{algebra-lemma-annihilate-ext}。因此显然有
$i(M/xM) = i(M) - 1$。对 $\delta(M)$ 归纳即完成证明。
\end{proof}

\begin{lemma}
\label{algebra-lemma-depth-in-ses}
设 $R$ 为诺特局部环，$0 \to N' \to N \to N'' \to 0$
为非零有限 $R$-模的短正合序列。
\begin{enumerate}
\item
$\text{depth}(N) \geq \min\{\text{depth}(N'), \text{depth}(N'')\}$
\item
$\text{depth}(N'') \geq \min\{\text{depth}(N), \text{depth}(N') - 1\}$
\item
$\text{depth}(N') \geq \min\{\text{depth}(N), \text{depth}(N'') + 1\}$
\end{enumerate}
\end{lemma}

\begin{proof}
使用引理 \ref{algebra-lemma-depth-ext} 中由 Ext 群
$\Ext^i(\kappa, N)$ 给出的深度刻画，并使用引理
\ref{algebra-lemma-long-exact-seq-ext} 给出的上同调长正合序列
$$
\begin{matrix}
0
\to \Hom_R(\kappa, N')
\to \Hom_R(\kappa, N)
\to \Hom_R(\kappa, N'')
\\
\phantom{0\ }
\to \Ext^1_R(\kappa, N')
\to \Ext^1_R(\kappa, N)
\to \Ext^1_R(\kappa, N'')
\to \ldots
\end{matrix}
$$
\end{proof}

\begin{lemma}
\label{algebra-lemma-depth-drops-by-one}
设 $R$ 为诺特局部环，$M$ 为非零有限 $R$-模。
\begin{enumerate}
\item 若 $x \in \mathfrak m$ 是 $M$ 上的非零因子，则
$\text{depth}(M/xM) = \text{depth}(M) - 1$。
\item 任意 $M$-正则序列 $x_1, \ldots, x_r$ 都可以扩充为
长度为 $\text{depth}(M)$ 的 $M$-正则序列。
\end{enumerate}
\end{lemma}

\begin{proof}
(2) 是 (1) 的形式推论。设 $x \in R$ 如 (1) 所述。
由短正合序列 $0 \to M \to M \to M/xM \to 0$
与引理 \ref{algebra-lemma-depth-in-ses}，可知深度至多下降 1。
另一方面，若 $x_1, \ldots, x_r \in \mathfrak m$
是 $M/xM$ 的正则序列，则 $x, x_1, \ldots, x_r$
是 $M$ 的正则序列。因此深度至少下降 1。
\end{proof}

\begin{lemma}
\label{algebra-lemma-inherit-minimal-primes}
设 $(R, \mathfrak m)$ 为诺特局部环，$M$ 为有限 $R$-模。
设 $x \in \mathfrak m$、$\mathfrak p \in \text{Ass}(M)$，并设
$\mathfrak q$ 是 $\mathfrak p + (x)$ 上的极小素理想。
则对某个 $n \geq 1$，有 $\mathfrak q \in \text{Ass}(M/x^nM)$。
\end{lemma}

\begin{proof}
选取子模 $N \subset M$，使得 $N \cong R/\mathfrak p$。
由 Artin--Rees 引理（引理 \ref{algebra-lemma-Artin-Rees}），
可以选取 $n > 0$，使得 $N \cap x^nM \subset xN$。
令 $\overline{N} \subset M/x^nM$ 为 $N \to M \to M/x^nM$ 的像。
由引理 \ref{algebra-lemma-ass}，只须证明
$\mathfrak q \in \text{Ass}(\overline{N})$。
根据 $n$ 的选择，存在满射
$\overline{N} \to N/xN = R/\mathfrak p + (x)$，
所以 $\mathfrak q$ 属于 $\overline{N}$ 的支撑。
由于 $\overline{N}$ 由 $x^n$ 和 $\mathfrak p$ 零化，
可知 $\mathfrak q$ 在 $\overline{N}$ 的支撑中的素理想里极小。
因此由引理 \ref{algebra-lemma-ass-minimal-prime-support}，
$\mathfrak q$ 是 $\overline{N}$ 的相伴素理想。
\end{proof}

\begin{lemma}
\label{algebra-lemma-depth-dim-associated-primes}
设 $(R, \mathfrak m)$ 为诺特局部环，$M$ 为有限 $R$-模。
对 $\mathfrak p \in \text{Ass}(M)$，有
$\dim(R/\mathfrak p) \geq \text{depth}(M)$。
\end{lemma}

\begin{proof}
若 $\mathfrak m \in \text{Ass}(M)$，则存在非零元
$x \in M$，它由 $\mathfrak m$ 的所有元素零化。
于是 $\text{depth}(M) = 0$。特别地，此时引理成立。

\medskip\noindent
若 $\text{depth}(M) = 1$，则由上一段可知
$\mathfrak m \not \in \text{Ass}(M)$。故对所有
$\mathfrak p \in \text{Ass}(M)$，有
$\dim(R/\mathfrak p) \geq 1$，此时引理也成立。

\medskip\noindent
一般情形下，对 $\text{depth}(M)$ 作归纳；可以且确实假设其值
$> 1$。选取 $M$ 上的非零因子 $x \in \mathfrak m$。
注意 $x \not \in \mathfrak p$（引理
\ref{algebra-lemma-ass-zero-divisors}）。由引理
\ref{algebra-lemma-one-equation}，有
$\dim(R/\mathfrak p + (x)) = \dim(R/\mathfrak p) - 1$。
因此存在 $\mathfrak p + (x)$ 上的极小素理想 $\mathfrak q$，使得
$\dim(R/\mathfrak q) = \dim(R/\mathfrak p) - 1$
（略去一个小论证；提示：诺特局部环 $A$ 的维数等于
$A/\mathfrak r$ 的维数在 $A$ 的极小素理想 $\mathfrak r$
上所取的最大值）。依引理 \ref{algebra-lemma-inherit-minimal-primes}
选取 $n$，使得 $\mathfrak q$ 是 $M/x^nM$ 的相伴素理想。
可以把归纳假设应用于 $M/x^nM$ 与 $\mathfrak q$，因为由引理
\ref{algebra-lemma-depth-drops-by-one}，
$\text{depth}(M/x^nM) = \text{depth}(M) - 1$。
于是得到 $\dim(R/\mathfrak q) \geq \text{depth}(M/x^nM)$，
结论得证。
\end{proof}

\begin{lemma}
\label{algebra-lemma-depth-localization}
设 $R$ 为诺特局部环，$M$ 为有限 $R$-模。
对素理想 $\mathfrak p \subset R$，有
$\text{depth}(M_\mathfrak p) + \dim(R/\mathfrak p) \geq \text{depth}(M)$。
\end{lemma}

\begin{proof}
若 $M_\mathfrak p = 0$，则
$\text{depth}(M_\mathfrak p) = \infty$，引理成立。
若 $\text{depth}(M) \leq \dim(R/\mathfrak p)$，则引理成立。
若 $\text{depth}(M) > \dim(R/\mathfrak p)$，则由引理
\ref{algebra-lemma-depth-dim-associated-primes}，$\mathfrak p$
不包含于 $M$ 的任何相伴素理想 $\mathfrak q$。
因此由引理 \ref{algebra-lemma-silly} 与
\ref{algebra-lemma-finite-ass}，可以找到 $x \in \mathfrak p$，
它不属于 $M$ 的任何相伴素理想。于是 $x$ 是 $M$ 上的非零因子；
参见引理 \ref{algebra-lemma-ass-zero-divisors}。
故 $\text{depth}(M/xM) = \text{depth}(M) - 1$，并且只要
$M_\mathfrak p$ 非零，就有
$\text{depth}(M_\mathfrak p / x M_\mathfrak p) =
\text{depth}(M_\mathfrak p) - 1$；参见引理
\ref{algebra-lemma-depth-drops-by-one}。因此，对
$\text{depth}(M)$ 作归纳即得结论。
\end{proof}

\begin{lemma}
\label{algebra-lemma-depth-goes-down-finite}
设 $(R, \mathfrak m)$ 为诺特局部环，$R \to S$ 为有限环映射。
设 $\mathfrak m_1, \ldots, \mathfrak m_n$ 为 $S$ 的极大理想，
$N$ 为有限 $S$-模。则
$$
\min\nolimits_{i = 1, \ldots, n} \text{depth}(N_{\mathfrak m_i}) =
\text{depth}_\mathfrak m(N)
$$
\end{lemma}

\begin{proof}
由引理 \ref{algebra-lemma-integral-no-inclusion}、
\ref{algebra-lemma-integral-going-up} 和
\ref{algebra-lemma-finite-finite-fibres}，$S$ 的极大理想恰为
$S$ 中位于 $\mathfrak m$ 上方的素理想，并且它们只有有限多个。
因此本引理的陈述有意义。对
$k = \min\nolimits_{i = 1, \ldots, n} \text{depth}(N_{\mathfrak m_i})$
作归纳。若 $k = 0$，则对某个 $i$ 有
$\text{depth}(N_{\mathfrak m_i}) = 0$。由引理
\ref{algebra-lemma-depth-ext}，这意味着
$\mathfrak m_i S_{\mathfrak m_i}$ 是
$N_{\mathfrak m_i}$ 的相伴素理想，因而
$\mathfrak m_i$ 是 $N$ 的相伴素理想（引理
\ref{algebra-lemma-localize-ass}）。由引理
\ref{algebra-lemma-ass-functorial-Noetherian}，把 $N$ 看作 $R$-模时，
$\mathfrak m$ 是它的相伴素理想。因此
$\text{depth}_\mathfrak m(N) = 0$。这就证明了初始情形。
若 $k > 0$，则 $\mathfrak m_i \not \in \text{Ass}_S(N)$。
再次由引理 \ref{algebra-lemma-ass-functorial-Noetherian}，
$\mathfrak m \not \in \text{Ass}_R(N)$。
因此可以找到 $f \in \mathfrak m$，它不是 $N$ 上的零因子；
参见引理 \ref{algebra-lemma-ideal-nonzerodivisor}。由引理
\ref{algebra-lemma-depth-drops-by-one}，从 $N$ 过渡到 $N/fN$ 时，
所有深度都恰好下降 $1$，其余由归纳假设得到。
\end{proof}

\section{Ext 的函子性}
\label{algebra-section-functoriality-ext}

\noindent
本节简要讨论 $\Ext$ 关于环变换等操作的函子性。
下面列出有待完成的事项。
\begin{enumerate}
\item 给定 $R \to R'$、一个 $R$-模 $M$ 和一个 $R'$-模 $N'$，
$R$-模 $\Ext^i_R(M, N')$ 具有自然的 $R'$-模结构。
此外，存在典范的 $R'$-线性映射
$\Ext^i_{R'}(M \otimes_R R', N') \to
\Ext^i_R(M, N')$。
\item 给定 $R \to R'$ 和 $R$-模 $M$, $N$，存在自然的
$R$-模映射
$\Ext^i_R(M, N) \to \text{Ext}^i_R(M, N \otimes_R R')$。
\end{enumerate}

\begin{lemma}
\label{algebra-lemma-flat-base-change-ext}
给定平坦环映射 $R \to R'$、一个 $R$-模 $M$ 和一个
$R'$-模 $N'$，自然映射
$$
\Ext^i_{R'}(M \otimes_R R', N') \to \text{Ext}^i_R(M, N')
$$
对 $i \geq 0$ 是同构。
\end{lemma}

\begin{proof}
选取 $M$ 的自由消解 $F_\bullet$。由于 $R \to R'$ 平坦，
$F_\bullet \otimes_R R'$ 是 $R'$ 上 $M \otimes_R R'$ 的自由消解。
本陈述即映射
$$
\Hom_{R'}(F_\bullet \otimes_R R', N') \to
\Hom_R(F_\bullet, N')
$$
在同调群上诱导同构；这是真的，因为由引理
\ref{algebra-lemma-adjoint-tensor-restrict}，它是复形的同构。
\end{proof}

\section{Ext 群的一个应用}
\label{algebra-section-ext-application}

\noindent
如下。

\begin{lemma}
\label{algebra-lemma-split-injection-after-completion}
设 $R$ 为诺特环，$I \subset R$ 为包含于 $R$ 的 Jacobson 根中的理想。
设 $N \to M$ 为有限 $R$-模之间的同态。
假设存在任意大的 $n$，使得
$N/I^nN \to M/I^nM$ 是分裂单射。
则 $N \to M$ 是分裂单射。
\end{lemma}

\begin{proof}
假设 $\varphi : N \to M$ 满足本引理的假设。
注意，这蕴含对任意大的 $n$ 都有 $\Ker(\varphi) \subset I^nN$。
因此由引理 \ref{algebra-lemma-intersection-powers-ideal-module}，
$\varphi$ 是单射。令 $Q = M/N$，于是有短正合序列
$$
0 \to N \to M \to Q \to 0.
$$
设
$$
F_2 \xrightarrow{d_2} F_1 \xrightarrow{d_1} F_0 \to Q \to 0
$$
为 $Q$ 的有限自由消解。可以选取提升映射 $F_0 \to Q$ 的映射
$\alpha : F_0 \to M$。这诱导映射
$\beta : F_1 \to N$，并且 $\beta \circ d_2 = 0$。
上面的扩张分裂，当且仅当存在映射 $\gamma : F_0 \to N$，
使得 $\beta = \gamma \circ d_1$。换言之，$\beta$ 在
$\Ext^1_R(Q, N)$ 中的类是上面的短正合序列分裂的障碍。

\medskip\noindent
假设 $n$ 是足够大的整数，使得
$N/I^nN \to M/I^nM$ 是分裂单射。这蕴含
$$
0 \to N/I^nN \to M/I^nM \to Q/I^nQ \to 0.
$$
仍是短正合的。此外，序列
$$
F_1/I^nF_1 \xrightarrow{d_1} F_0/I^nF_0 \to Q/I^nQ \to 0
$$
仍正合。与上面同样论证可知，由 $\beta$ 诱导的映射
$\overline{\beta} : F_1/I^nF_1 \to N/I^nN$
等于某个映射
$\overline{\gamma_n} : F_0/I^nF_0 \to N/I^nN$
所给出的复合 $\overline{\gamma_n} \circ d_1$。
由于 $F_0$ 自由，可以把 $\overline{\gamma_n}$ 提升为映射
$\gamma_n : F_0 \to N$；于是
$\beta - \gamma_n \circ d_1$ 是从 $F_1$ 到 $I^nN$ 的映射。
换言之，对这个 $n$，有
$$
\beta \in
\Im\Big(\Hom_R(F_0, N) \to \Hom_R(F_1, N)\Big) + I^n\Hom_R(F_1, N).
$$

\medskip\noindent
由于依假设，对任意大的 $n$ 都有这一性质，
由引理 \ref{algebra-lemma-intersection-powers-ideal-module}，
$\beta$ 在 $\Hom_R(F_0, N) \to \Hom_R(F_1, N)$ 的余核中的像为零。
因此按所需，$\beta$ 属于映射
$\Hom_R(F_0, N) \to \Hom_R(F_1, N)$ 的像。
\end{proof}

\section{Tor 群与平坦性}
\label{algebra-section-tor}

\noindent
本节使用上一节发展的一些同调代数来说明 Tor 群是什么。
具体地，设 $R$ 为环，$M$、$N$ 为两个 $R$-模。
选取 $M$ 的一个由自由 $R$-模组成的消解 $F_\bullet$，
参见引理 \ref{algebra-lemma-resolution-by-finite-free}。考虑同调复形
$$
F_\bullet \otimes_R N
:
\ldots
\to F_2 \otimes_R N
\to F_1 \otimes_R N
\to F_0 \otimes_R N
$$
我们把 $\text{Tor}^R_i(M, N)$ 定义为这个复形的第 $i$ 个同调群。
下面的引理说明这一对象在何种意义下是良定义的。

\begin{lemma}
\label{algebra-lemma-tor-welldefined}
设 $R$ 为环，$M_1, M_2, N$ 为 $R$-模。
假设 $F_\bullet$ 是模 $M_1$ 的自由消解，
$G_\bullet$ 是模 $M_2$ 的自由消解。
设 $\varphi : M_1 \to M_2$ 为模映射。
设 $\alpha : F_\bullet \to G_\bullet$ 为一个复形映射，
它在
$M_1 = \Coker(d_{F, 1}) \to M_2 = \Coker(d_{G, 1})$
上诱导 $\varphi$；参见引理 \ref{algebra-lemma-compare-resolutions}。
则诱导映射
$$
H_i(\alpha) :
H_i(F_\bullet \otimes_R N)
\longrightarrow
H_i(G_\bullet \otimes_R N)
$$
与 $\alpha$ 的选取无关。若 $\varphi$ 是同构，则所有映射
$H_i(\alpha)$ 都是同构。若 $M_1 = M_2$、
$F_\bullet = G_\bullet$ 且 $\varphi$ 是恒等映射，
则所有映射 $H_i(\alpha)$ 也都是恒等映射。
\end{lemma}

\begin{proof}
本引理的证明与引理 \ref{algebra-lemma-ext-welldefined} 的证明相同。
\end{proof}

\noindent
这个引理不仅蕴含 Tor 模是良定义的，还给出了构造
$(M, N) \mapsto \text{Tor}_i^R(M, N)$
关于第一个变量的函子性。关于第二个变量的函子性当然是显然的。
我们留给读者验证，每个 $\text{Tor}_i^R$ 实际上都是函子
$$
\text{Mod}_R \times \text{Mod}_R \to \text{Mod}_R.
$$
这里 $\text{Mod}_R$ 表示 $R$-模范畴；积范畴的定义见
《范畴论》，定义 \ref{categories-definition-product-category}。
具体地，给定 $R$-模态射 $M_1 \to M_2$ 和 $N_1 \to N_2$，
我们得到交换图
$$
\xymatrix{
\text{Tor}_i^R(M_1, N_1) \ar[r] \ar[d] &
\text{Tor}_i^R(M_1, N_2) \ar[d] \\
\text{Tor}_i^R(M_2, N_1) \ar[r] &
\text{Tor}_i^R(M_2, N_2) \\
}
$$

\begin{lemma}
\label{algebra-lemma-long-exact-sequence-tor}
设 $R$ 为环，$M$ 为 $R$-模。
假设 $0 \to N' \to N \to N'' \to 0$ 是 $R$-模的短正合序列。
存在长正合序列
$$
\text{Tor}_1^R(M, N')
\to \text{Tor}_1^R(M, N)
\to \text{Tor}_1^R(M, N'')
\to
M \otimes_R N'
\to M \otimes_R N
\to M \otimes_R N''
\to 0
$$
\end{lemma}

\begin{proof}
这个证明与引理 \ref{algebra-lemma-long-exact-seq-ext} 的证明相同。
\end{proof}

\noindent
考虑一个由 $R$-模组成的同调双复形
$$
\xymatrix{
\ldots \ar[r]^d &
A_{2, 0} \ar[r]^d &
A_{1, 0} \ar[r]^d &
A_{0, 0} \\
\ldots \ar[r]^d &
A_{2, 1} \ar[r]^d \ar[u]^\delta &
A_{1, 1} \ar[r]^d \ar[u]^\delta &
A_{0, 1} \ar[u]^\delta \\
\ldots \ar[r]^d &
A_{2, 2} \ar[r]^d \ar[u]^\delta &
A_{1, 2} \ar[r]^d \ar[u]^\delta &
A_{0, 2} \ar[u]^\delta \\
&
\ldots \ar[u]^\delta &
\ldots \ar[u]^\delta &
\ldots \ar[u]^\delta \\
}
$$
这意味着 $d_{i, j} : A_{i, j} \to A_{i-1, j}$ 和
$\delta_{i, j} : A_{i, j} \to A_{i, j-1}$ 具有下列性质：
\begin{enumerate}
\item 任意两个 $d_{i, j}$ 的复合为零。换言之，双复形的各行都是复形。
\item 任意两个 $\delta_{i, j}$ 的复合为零。换言之，双复形的各列都是复形。
\item 对任意指标对 $(i, j)$，有 $\delta_{i-1, j} \circ d_{i, j}
= d_{i, j-1} \circ \delta_{i, j}$。换言之，所有方块都交换。
\end{enumerate}
正确的做法是给每个这样的双复形配上一个谱序列。
不过目前我们用稍微容易一些的办法就足够了。

\medskip\noindent
具体地，仅为本节之目的，给定双复形
$(A_{\bullet, \bullet}, d, \delta)$，令
$R(A)_j = \Coker(A_{1, j} \to A_{0, j})$，并令
$U(A)_i = \Coker(A_{i, 1} \to A_{i, 0})$。（字母
$R$ 和 $U$ 分别提示 Right（右）与 Up（上）。）
我们用映射 $\delta$ 赋予 $R(A)_\bullet$ 复形结构；
类似地，用映射 $d$ 赋予 $U(A)_\bullet$ 复形结构。
换言之，我们得到如下巨大的交换图：
$$
\xymatrix{
\ldots \ar[r]^d &
U(A)_2 \ar[r]^d &
U(A)_1 \ar[r]^d &
U(A)_0 &
\\
\ldots \ar[r]^d &
A_{2, 0} \ar[r]^d \ar[u] &
A_{1, 0} \ar[r]^d \ar[u] &
A_{0, 0} \ar[r] \ar[u] &
R(A)_0 \\
\ldots \ar[r]^d &
A_{2, 1} \ar[r]^d \ar[u]^\delta &
A_{1, 1} \ar[r]^d \ar[u]^\delta &
A_{0, 1} \ar[r] \ar[u]^\delta &
R(A)_1 \ar[u]^\delta \\
\ldots \ar[r]^d &
A_{2, 2} \ar[r]^d \ar[u]^\delta &
A_{1, 2} \ar[r]^d \ar[u]^\delta &
A_{0, 2} \ar[r] \ar[u]^\delta &
R(A)_2 \ar[u]^\delta \\
&
\ldots \ar[u]^\delta &
\ldots \ar[u]^\delta &
\ldots \ar[u]^\delta &
\ldots \ar[u]^\delta \\
}
$$
（当然，这已不再是双复形。）双复形的态射
$\Phi : (A_{\bullet, \bullet}, d, \delta)
\to (B_{\bullet, \bullet}, d, \delta)$
是什么意思是清楚的；它会诱导复形态射
$R(\Phi) : R(A)_\bullet \to R(B)_\bullet$ 和
$U(\Phi) : U(A)_\bullet \to U(B)_\bullet$。

\begin{lemma}
\label{algebra-lemma-no-spectral-sequence}
设 $(A_{\bullet, \bullet}, d, \delta)$ 是满足下列条件的双复形：
\begin{enumerate}
\item 每一行 $A_{\bullet, j}$ 都是 $R(A)_j$ 的一个消解。
\item 每一列 $A_{i, \bullet}$ 都是 $U(A)_i$ 的一个消解。
\end{enumerate}
则存在典范同构
$$
H_i(R(A)_\bullet)
\cong
H_i(U(A)_\bullet).
$$
这些同构关于满足上述性质的双复形态射是函子性的。
\end{lemma}

\begin{proof}
我们将证明 $H_i(R(A)_\bullet))$ 与
$H_i(U(A)_\bullet)$ 都典范同构于第三个群，即
$$
\mathbf{H}_i(A) :=
\frac{
\{
(a_{i, 0}, a_{i-1, 1}, \ldots, a_{0, i})
\mid
d(a_{i, 0}) = \delta(a_{i-1, 1}), \ldots,
d(a_{1, i-1}) = \delta(a_{0, i})
\}}
{
\{
d(a_{i + 1, 0}) + \delta(a_{i, 1}),
d(a_{i, 1}) + \delta(a_{i-1, 2}),
\ldots,
d(a_{1, i}) + \delta(a_{0, i + 1})
\}
}
$$
这里采用记号约定：$a_{i, j}$ 表示 $A_{i, j}$ 的一个元素。
换言之，$\mathbf{H}_i$ 的一个元素由一条之字形表示；
当 $i = 2$ 时可画成
$$
\xymatrix{
a_{2, 0} \ar@{|->}[r] & d(a_{2, 0}) = \delta(a_{1, 1}) & \\
& a_{1, 1} \ar@{|->}[u] \ar@{|->}[r] & d(a_{1, 1}) = \delta(a_{0, 2}) \\
& & a_{0, 2} \ar@{|->}[u] \\
}
$$
自然地，我们要除去“平凡”之字形，也就是由形如
$(0, \ldots, 0, -\delta(a_{t + 1, t-i}),
d(a_{t + 1, t-i}), 0, \ldots, 0)$ 的元素生成的子模。
注意存在典范同态
$$
\mathbf{H}_i(A) \to H_i(R(A)_\bullet), \quad
(a_{i, 0}, a_{i-1, 1}, \ldots, a_{0, i}) \mapsto
\text{class of image of }a_{0, i}
$$
以及
$$
\mathbf{H}_i(A) \to H_i(U(A)_\bullet), \quad
(a_{i, 0}, a_{i-1, 1}, \ldots, a_{0, i}) \mapsto
\text{class of image of }a_{i, 0}
$$

\medskip\noindent
先证明这些映射是满射。设
$\overline{r} \in H_i(R(A)_\bullet)$。
令 $r \in R(A)_i$ 为表示 $\overline{r}$ 之类的一个余循环。
取映到 $r$ 的元素 $a_{0, i} \in A_{0, i}$。
因为 $\delta(r) = 0$，所以 $\delta(a_{0, i})$ 属于 $d$ 的像。
故存在 $a_{1, i-1} \in A_{1, i-1}$，使得
$d(a_{1, i-1}) = \delta(a_{0, i})$。
这又蕴含 $\delta(a_{1, i-1})$ 属于 $d$ 的核，因为
$d(\delta(a_{1, i-1})) = \delta(d(a_{1, i-1}))
= \delta(\delta(a_{0, i})) = 0$。由各行的正合性，
可找到元素 $a_{2, i-2}$，使得
$d(a_{2, i-2}) = \delta(a_{1, i-1})$。如此继续，
直到得到一条完整的之字形。当然，
$\mathbf{H}_i \to H_i(U(A))$ 的满射性也可同样证明。

\medskip\noindent
为证明单射性，我们采用完全相同的论证。
具体地，设给定一条映到 $H_i(R(A)_\bullet)$ 中零元的之字形
$(a_{i, 0}, a_{i-1, 1}, \ldots, a_{0, i})$。
这意味着 $a_{0, i}$ 映到
$\Coker(A_{i, 1} \to A_{i, 0})$ 的一个元素，
而此元素属于
$\delta : \Coker(A_{i + 1, 1} \to A_{i + 1, 0}) \to
\Coker(A_{i, 1} \to A_{i, 0})$ 的像。
换言之，$a_{0, i}$ 属于
$\delta \oplus d : A_{0, i + 1} \oplus A_{1, i} \to A_{0, i}$ 的像。
由平凡之字形的定义可知，我们可以用一条平凡之字形修改原之字形，
并假设 $a_{0, i} = 0$。这立刻蕴含
$d(a_{1, i-1}) = 0$。由于各行正合，
$a_{1, i-1}$ 属于 $d : A_{2, i-1} \to A_{1, i-1}$ 的像。
于是可以再次用一条平凡之字形修改，并假设原之字形形如
$(a_{i, 0}, a_{i-1, 1}, \ldots, a_{2, i-2}, 0, 0)$。
如此继续便得到所需的单射性。

\medskip\noindent
若 $\Phi : (A_{\bullet, \bullet}, d, \delta)
\to (B_{\bullet, \bullet}, d, \delta)$
是两个都满足本引理条件的双复形之间的态射，
则显然得到交换图
$$
\xymatrix{
H_i(U(A)_\bullet) \ar[d] &
\mathbf{H}_i(A) \ar[r] \ar[l] \ar[d] &
H_i(R(A)_\bullet) \ar[d] \\
H_i(U(B)_\bullet) &
\mathbf{H}_i(B) \ar[r] \ar[l] &
H_i(R(B)_\bullet) \\
}
$$
这证明了函子性。
\end{proof}

\begin{remark}
\label{algebra-remark-signs-double-complex}
上面构造的同构仅在相差符号的意义下才是“正确”的。
同调代数中很大一部分工作在于为各种映射选取符号，
并说明在插入适当符号后相应图表交换。
目前我们就使用上述证明给出的同构，稍后再处理符号问题。
\end{remark}

\begin{lemma}
\label{algebra-lemma-tor-left-right}
设 $R$ 为环。对任意 $i \geq 0$，函子
$\text{Mod}_R \times \text{Mod}_R \to \text{Mod}_R$，
$(M, N) \mapsto \text{Tor}_i^R(M, N)$ 与
$(M, N) \mapsto \text{Tor}_i^R(N, M)$ 典范同构。
\end{lemma}

\begin{proof}
设 $F_\bullet$ 是模 $M$ 的自由消解，
$G_\bullet$ 是模 $N$ 的自由消解。
考虑如下定义的双复形 $(A_{i, j}, d, \delta)$：
\begin{enumerate}
\item 令 $A_{i, j} = F_i \otimes_R G_j$；
\item 令 $d_{i, j} : F_i \otimes_R G_j \to F_{i-1} \otimes G_j$
等于 $d_{F, i} \otimes \text{id}$；
\item 令 $\delta_{i, j} : F_i \otimes_R G_j \to F_i \otimes G_{j-1}$
等于 $\text{id} \otimes d_{G, j}$。
\end{enumerate}
这个双复形通常简记为 $F_\bullet \otimes_R G_\bullet$。

\medskip\noindent
因为每个 $G_j$ 都自由，因而平坦，所以双复形的每一行除同调次数
$0$ 外都正合。因为每个 $F_i$ 都自由，因而平坦，所以双复形的每一列除同调次数
$0$ 外都正合。因此这个双复形满足引理
\ref{algebra-lemma-no-spectral-sequence} 的条件。

\medskip\noindent
为看清该引理给出什么，我们计算 $R(A)_\bullet$ 与 $U(A)_\bullet$。
具体地，
\begin{eqnarray*}
R(A)_i & = & \Coker(A_{1, i} \to A_{0, i}) \\
& = & \Coker(F_1 \otimes_R G_i \to F_0 \otimes_R G_i) \\
& = & \Coker(F_1 \to F_0) \otimes_R G_i \\
& = & M \otimes_R G_i
\end{eqnarray*}
事实上，这些同构与微分 $\delta$ 相容，因而作为同调复形有
$R(A)_\bullet = M \otimes_R G_\bullet$。
完全类似地，有 $U(A)_\bullet = F_\bullet \otimes_R N$。于是
\begin{eqnarray*}
\text{Tor}_i^R(M, N)
& = & H_i(F_\bullet \otimes_R N) \\
& = & H_i(U(A)_\bullet) \\
& = & H_i(R(A)_\bullet) \\
& = & H_i(M \otimes_R G_\bullet) \\
& = & H_i(G_\bullet \otimes_R M) \\
& = & \text{Tor}_i^R(N, M)
\end{eqnarray*}
这里第三个等号来自引理 \ref{algebra-lemma-no-spectral-sequence}，
第五个等号使用张量积的同构 $V \otimes W = W \otimes V$。

\medskip\noindent
函子性。假设有 $R$-模 $M_\nu$、$N_\nu$，其中
$\nu = 1, 2$。设 $\varphi : M_1 \to M_2$ 和
$\psi : N_1 \to N_2$ 为 $R$-模态射。
假设对 $M_\nu$ 有自由消解 $F_{\nu, \bullet}$，
对 $N_\nu$ 有自由消解 $G_{\nu, \bullet}$。
由引理 \ref{algebra-lemma-compare-resolutions}，可选取与
$\varphi$ 和 $\psi$ 相容的复形映射
$\alpha : F_{1, \bullet} \to F_{2, \bullet}$ 与
$\beta : G_{1, \bullet} \to G_{2, \bullet}$。
我们断言，映射对 $(\alpha, \beta)$ 诱导双复形态射
$$
\alpha \otimes \beta :
F_{1, \bullet} \otimes_R G_{1, \bullet}
\longrightarrow
F_{2, \bullet} \otimes_R G_{2, \bullet}
$$
这确实只需直接验证：映射
$F_{1, i} \otimes_R G_{1, j} \to F_{2, i} \otimes_R G_{2, j}$
由 $\alpha_i \otimes \beta_j$ 给出；其中 $\alpha_i$ 与
$\beta_j$ 分别是 $\alpha$ 的次数 $i$ 分量与 $\beta$ 的次数 $j$ 分量。
读者也容易验证，诱导映射
$R(F_{1, \bullet} \otimes_R G_{1, \bullet})_\bullet
\to R(F_{2, \bullet} \otimes_R G_{2, \bullet})_\bullet$
与由 $\varphi \otimes \beta$ 诱导的映射
$M_1 \otimes_R G_{1, \bullet}
\to M_2 \otimes_R G_{2, \bullet}$ 相同。
在 $U(-)_\bullet$ 复形上诱导的映射也类似。
因此，函子性断言来自引理
\ref{algebra-lemma-no-spectral-sequence} 中的函子性断言。
\end{proof}

\begin{remark}
\label{algebra-remark-curiosity-signs-swap}
上面的 $M = N$ 情形很有意思。此时得到典范映射
$\text{Tor}_i^R(M, M) \to \text{Tor}_i^R(M, M)$。
注意，这个映射不是恒等映射；即使 $i = 0$，它也不是恒等映射！
例如，若 $V$ 是域上的 $n$ 维向量空间，则交换映射
$V \otimes_k V \to V \otimes_k V$ 有
$(n^2 + n)/2$ 个特征值 $+1$ 和
$(n^2-n)/2$ 个特征值 $-1$。在特征 $2$ 时，它甚至不可对角化。
注意，即使改变映射的符号也无法消除这一现象。
\end{remark}

\begin{lemma}
\label{algebra-lemma-tor-noetherian}
设 $R$ 为诺特环，$M$、$N$ 为有限 $R$-模。
则对所有 $p$，$\text{Tor}_p^R(M, N)$ 都是有限 $R$-模。
\end{lemma}

\begin{proof}
这是因为 $\text{Tor}_p^R(M, N)$ 可计算为复形
$F_\bullet \otimes_R N$ 的上同调群，其中每个 $F_n$ 都是有限自由
$R$-模；参见引理 \ref{algebra-lemma-resolution-by-finite-free}。
\end{proof}

\begin{lemma}
\label{algebra-lemma-characterize-flat}
设 $R$ 为环，$M$ 为 $R$-模。下列条件等价：
\begin{enumerate}
\item 模 $M$ 在 $R$ 上平坦。
\item 对所有 $i > 0$，函子 $\text{Tor}_i^R(M, -)$ 为零。
\item 函子 $\text{Tor}_1^R(M, -)$ 为零。
\item 对所有理想 $I \subset R$，有 $\text{Tor}_1^R(M, R/I) = 0$。
\item 对所有有限生成理想 $I \subset R$，有
$\text{Tor}_1^R(M, R/I) = 0$。
\end{enumerate}
\end{lemma}

\begin{proof}
假设 $M$ 平坦。设 $N$ 为 $R$-模，并设 $F_\bullet$ 是
$N$ 的自由消解。由 $M$ 的平坦性，
$F_\bullet \otimes_R M$ 是 $N \otimes_R M$ 的一个消解。
故所有高阶 Tor 群都消失。

\medskip\noindent
现在只需证明最后一个条件蕴含 $M$ 平坦。
设 $I \subset R$ 为理想。考虑短正合序列
$0 \to I \to R \to R/I \to 0$。
应用引理 \ref{algebra-lemma-long-exact-sequence-tor}，得到正合序列
$$
\text{Tor}_1^R(M, R/I) \to
M \otimes_R I \to
M \otimes_R R \to
M \otimes_R R/I \to
0
$$
由于显然 $M \otimes_R R = M$，可知最后一个假设蕴含：
对每个有限生成理想 $I$，映射 $M \otimes_R I \to M$ 都是单射。
因此由引理 \ref{algebra-lemma-flat}，$M$ 平坦。
\end{proof}

\begin{remark}
\label{algebra-remark-Tor-ring-mod-ideal}
引理 \ref{algebra-lemma-characterize-flat} 的证明实际上表明
$$
\text{Tor}_1^R(M, R/I)
=
\Ker(I \otimes_R M \to M).
$$
\end{remark}

\section{Tor 的函子性}
\label{algebra-section-functoriality-tor}

\noindent
本节简要讨论 $\text{Tor}$ 关于环变换等操作的函子性。
下面列出若干需要完成的事项。
\begin{enumerate}
\item 给定环映射 $R \to R'$、一个 $R$-模 $M$ 和一个
$R'$-模 $N'$，$R$-模 $\text{Tor}_i^R(M, N')$ 具有自然的
$R'$-模结构。
\item 给定环映射 $R \to R'$ 以及 $R$-模 $M$、$N$，
存在自然的 $R$-模映射
$\text{Tor}_i^R(M, N) \to \text{Tor}_i^{R'}(M \otimes_R R', N \otimes_R R')$。
\item 给定环映射 $R \to R'$、一个 $R$-模 $M$ 和一个
$R'$-模 $N'$，存在自然的 $R'$-模映射
$\text{Tor}_i^R(M, N') \to \text{Tor}_i^{R'}(M \otimes_R R', N')$。
\end{enumerate}

\begin{lemma}
\label{algebra-lemma-flat-base-change-tor}
给定平坦环映射 $R \to R'$ 以及 $R$-模 $M$、$N$，
对所有 $i$，自然的 $R$-模映射
$\text{Tor}_i^R(M, N)\otimes_R R'
\to \text{Tor}_i^{R'}(M \otimes_R R', N \otimes_R R')$
都是同构。
\end{lemma}

\begin{proof}
略。这是因为 $M$ 在 $R$ 上的自由消解 $F_\bullet$ 在与
$R'$ 关于 $R$ 作张量积后仍正合，因而
$(F_\bullet \otimes_R N)\otimes_R R'$ 计算 $R'$ 上的 Tor 群。
\end{proof}

\noindent
下面的引理似乎放在其他地方都不合适。

\begin{lemma}
\label{algebra-lemma-tor-commutes-filtered-colimits}
设 $R$ 为环，$M = \colim M_i$ 是 $R$-模的滤过余极限，
$N$ 是 $R$-模。则对所有 $n$，
$\text{Tor}_n^R(M, N) = \colim \text{Tor}_n^R(M_i, N)$。
\end{lemma}

\begin{proof}
选取 $N$ 的自由消解 $F_\bullet$。由引理
\ref{algebra-lemma-tensor-products-commute-with-limits}，作为复形有
$F_\bullet \otimes_R M = \colim F_\bullet \otimes_R M_i$。
于是由引理 \ref{algebra-lemma-directed-colimit-exact} 得到结论。
\end{proof}

\section{投射模}
\label{algebra-section-projective}

\noindent
下面给出关于投射模的一些引理。

\begin{definition}
\label{algebra-definition-projective}
设 $R$ 为环。一个 $R$-模 $P$ 称为\textit{投射的}，
当且仅当函子
$\Hom_R(P, -) : \text{Mod}_R \to \text{Mod}_R$ 是正合函子。
\end{definition}

\noindent
对任意 $R$-模 $M$，函子 $\Hom_R(M, - )$ 都左正合，
参见引理 \ref{algebra-lemma-hom-exact}。
因此，$P$ 为投射模这一条件真正表达的是：
给定 $R$-模满射 $N \to N'$，映射
$\Hom_R(P, N) \to \Hom_R(P, N')$ 是满射。

\begin{lemma}
\label{algebra-lemma-characterize-projective}
设 $R$ 为环，$P$ 为 $R$-模。下列条件等价：
\begin{enumerate}
\item $P$ 是投射模；
\item $P$ 是某个自由 $R$-模的直和因子；
\item 对每个 $R$-模 $M$，都有 $\Ext^1_R(P, M) = 0$。
\end{enumerate}
\end{lemma}

\begin{proof}
假设 $P$ 是投射模。选取满射 $\pi : F \to P$，其中 $F$
是自由 $R$-模。因为 $P$ 投射，存在
$i \in \Hom_R(P, F)$，使得 $\pi \circ i = \text{id}_P$。
换言之，$F \cong \Ker(\pi) \oplus i(P)$，从而 $P$ 是 $F$ 的直和因子。

\medskip\noindent
反过来，假设 $P \oplus Q = F$ 是自由 $R$-模。
注意，自由模 $F = \bigoplus_{i \in I} R$ 是投射模，
因为有等式：
$$
\Hom_R(F, M) = \prod_{i \in I} M
$$
而函子 $M \mapsto \prod_{i \in I} M$ 正合。
又因为作为函子有下式：
$$
\Hom_R(F, -) = \Hom_R(P, -) \times \Hom_R(Q, -)
$$
所以 $P$ 与 $Q$ 都是投射模。

\medskip\noindent
假设 $P \oplus Q = F$ 是自由 $R$-模。则有形如
$$
\ldots F \xrightarrow{a} F \xrightarrow{b} F \to P \to 0
$$
的自由消解 $F_\bullet$，其中映射 $a, b$ 交替出现，
并且分别等于到 $P$ 和 $Q$ 上的投影算子。
因此复形 $\Hom_R(F_\bullet, M)$ 在次数 $\geq 1$ 上分裂正合，
从而得到条件 (3) 中的消失性。

\medskip\noindent
假设对每个 $R$-模 $M$ 都有 $\Ext^1_R(P, M) = 0$。
取一个自由消解 $F_\bullet \to P$，并令
$M = \Im(F_1 \to F_0) = \Ker(F_0 \to P)$。
考虑由商映射 $\pi : F_1 \to M$ 的类给出的元素
$\xi \in \Ext^1_R(P, M)$。由于 $\xi$ 为零，存在映射
$s : F_0 \to M$，使得 $\pi = s \circ (F_1 \to F_0)$。
显然，这意味着
$$
F_0 = \Ker(s) \oplus \Ker(F_0 \to P) =
P \oplus \Ker(F_0 \to P)
$$
结论得证。
\end{proof}

\begin{lemma}
\label{algebra-lemma-characterize-finite-projective-noetherian}
设 $R$ 为诺特环，$P$ 为有限 $R$-模。
如果对每个有限 $R$-模 $M$ 都有 $\Ext^1_R(P, M) = 0$，
则 $P$ 是投射模。
\end{lemma}

\noindent
这个引理可以加强：当不假设 $R$ 为诺特环时，
存在一个适用于有限表现 $R$-模的版本；在诺特情形，
也有一个让 $M$ 遍历所有有限长度模的版本。

\begin{proof}
选取满射 $R^{\oplus n} \to P$，其核为 $M$。
因为 $\Ext^1_R(P, M) = 0$，这个满射分裂；
于是由引理 \ref{algebra-lemma-characterize-projective} 得到结论。
\end{proof}

\begin{lemma}
\label{algebra-lemma-direct-sum-projective}
投射模的直和是投射模。
\end{lemma}

\begin{proof}
这由引理 \ref{algebra-lemma-characterize-projective} 中
“投射模是自由模的直和因子”这一刻画得到。
\end{proof}

\begin{lemma}
\label{algebra-lemma-lift-projective-module}
设 $R$ 为环，$I \subset R$ 为幂零理想，
$\overline{P}$ 为投射 $R/I$-模。
则存在投射 $R$-模 $P$，使得 $P/IP \cong \overline{P}$。
\end{lemma}

\begin{proof}
由引理 \ref{algebra-lemma-characterize-projective}，可以选取集合 $A$ 以及直和分解
$\bigoplus_{\alpha \in A} R/I = \overline{P} \oplus \overline{K}$，
其中 $\overline{K}$ 是某个 $R/I$-模。
记 $F = \bigoplus_{\alpha \in A} R$ 为以 $A$ 为基的自由 $R$-模。
设 $\overline{p}$ 是直和因子
$\overline{P}$ 在 $\bigoplus_{\alpha \in A} R/I$ 中对应的投影算子，
选取它的一个提升 $p : F \to F$。
注意，$p^2 - p \in \text{End}_R(F)$ 是 $F$ 的幂零自同态
（因为 $I$ 幂零且 $p^2 - p$ 的矩阵元素属于 $I$；更精确地，
若 $I^n = 0$，则 $(p^2 - p)^n = 0$）。
因此由引理 \ref{algebra-lemma-lift-idempotents-noncommutative}，
可以修改 $p$ 的选取并假设 $p$ 是投影算子。令 $P = \Im(p)$。
\end{proof}

\begin{lemma}
\label{algebra-lemma-lift-finite-projective-module}
设 $R$ 为环，$I \subset R$ 为局部幂零理想，
$\overline{P}$ 为有限投射 $R/I$-模。
则存在有限投射 $R$-模 $P$，使得 $P/IP \cong \overline{P}$。
\end{lemma}

\begin{proof}
回顾一下，由引理 \ref{algebra-lemma-characterize-projective}，
$\overline{P}$ 是自由 $R/I$-模
$\bigoplus_{\alpha \in A} R/I$ 的直和因子。
因为 $\overline{P}$ 有限，所以存在有限子集 $A' \subset A$，
使得 $\overline{P}$ 包含于 $\bigoplus_{\alpha \in A'} R/I$。
因此可以假设存在直和分解
$(R/I)^{\oplus n} = \overline{P} \oplus \overline{K}$，
其中 $n$ 是某个整数，$\overline{K}$ 是某个 $R/I$-模。
设 $\overline{p}$ 是直和因子 $\overline{P}$ 在
$(R/I)^{\oplus n}$ 中对应的投影算子，选取其提升
$p \in \text{Mat}(n \times n, R)$。
注意 $p^2 - p \in \text{Mat}(n \times n, R)$ 是幂零的：
因为 $I$ 局部幂零，且 $p^2 - p$ 的矩阵元素 $c_{ij}$ 属于 $I$，
所以对某个 $t > 0$ 有 $c_{ij}^t = 0$，进而
$(p^2 - p)^{tn^2} = 0$（考察矩阵系数即可）。
因此由引理 \ref{algebra-lemma-lift-idempotents-noncommutative}，
可以修改 $p$ 的选取并假设 $p$ 是投影算子。令 $P = \Im(p)$。
\end{proof}

\begin{lemma}
\label{algebra-lemma-lift-projective}
设 $R$ 为环，$I \subset R$ 为理想，$M$ 为 $R$-模。假设
\begin{enumerate}
\item $I$ 是幂零的；
\item $M/IM$ 是投射 $R/I$-模；
\item $M$ 是平坦 $R$-模。
\end{enumerate}
则 $M$ 是投射 $R$-模。
\end{lemma}

\begin{proof}
由引理 \ref{algebra-lemma-lift-projective-module}，
可以找到投射 $R$-模 $P$ 和同构 $P/IP \to M/IM$。
我们将证明 $M$ 同构于 $P$，从而完成证明。
因为 $P$ 投射，可以把映射 $P \to P/IP \to M/IM$
提升为 $R$-模映射 $P \to M$；它模 $I$ 后是同构。
对某个 $n$ 有 $I^n = 0$，所以可利用滤过
\begin{align*}
0 = I^nM \subset I^{n - 1}M \subset \ldots \subset IM \subset M \\
0 = I^nP \subset I^{n - 1}P \subset \ldots \subset IP \subset P
\end{align*}
看出只需证明诱导映射
$I^aP/I^{a + 1}P \to I^aM/I^{a + 1}M$ 都是双射。
因为 $P$ 与 $M$ 都是平坦 $R$-模，可以把这个映射等同于由
$P \to M$ 诱导的映射
$$
I^a/I^{a + 1} \otimes_{R/I} P/IP
\longrightarrow
I^a/I^{a + 1} \otimes_{R/I} M/IM
$$
我们选取 $P \to M$ 时已使诱导映射
$P/IP \to M/IM$ 成为同构，故结论成立。
\end{proof}

\begin{lemma}
\label{algebra-lemma-projective-modulo-two-ideals}
设 $R$ 为环，$I, J \subset R$ 是满足 $I \cap J = 0$ 的理想。
设 $P$ 为 $R$-模，并且 $P/IP$ 是投射 $R/I$-模，
$P/JP$ 是投射 $R/J$-模。则 $P$ 是投射 $R$-模。
\end{lemma}

\begin{proof}
选取满射 $p : F \to P$，其中 $F$ 是自由 $R$-模。
因为 $P/IP$ 是投射 $R/I$-模，可以选取映射
$f : P/IP \to F/IF$，它是 $p \bmod I$ 的右逆。
考虑由商映射 $F/JF \to F/(I + J)F$ 和 $p$ 诱导的映射
$q : F/JF \to F/(I + J)F \times_{P/(I + J)P} P/JP$。
注意，$q$ 是 $R/J$-模的满射（略去一个小细节）。
再考虑由 $f$ 与恒等映射诱导的映射
$f' : P/JP \to F/(I + J)F \times_{P/(I + J)P} P/JP$。
因为 $P/JP$ 是投射 $R/J$-模且 $q$ 为满射，
可以选取映射 $g : P/JP \to F/JF$，使得 $q \circ g = f'$。
于是 $f \bmod I + J = g \bmod I + J$。
又因为 $I \cap J = 0$，有
$F = F/JF \times_{F/(I + J)F} F/IF$；
所以得到良定义的 $R$-模同态 $h = (f, g) : P \to F$。
映射 $a = p \circ h : P \to P$ 模 $I$ 和模 $J$ 后都化为恒等映射。
于是 $a$ 在子模 $IP$ 上也诱导恒等映射：若
$x = \sum i_\alpha x_\alpha$，其中 $i_\alpha \in I$、
$x_\alpha \in P$，则
$a(x) = \sum i_\alpha a(x_\alpha) = \sum i_\alpha x_\alpha = x$，
因为 $a(x_\alpha) = x_\alpha + j_\alpha$，其中
$j_\alpha \in J$，且 $i_\alpha j_\alpha = 0$。
因此 $a : P \to P$ 在 $IP$ 以及 $P/IP$ 上都作为恒等映射，
从而 $a$ 是 $P$ 的自同构（略去一个小细节）。
故 $P$ 是 $F$ 的直和因子，因而是投射模。
\end{proof}

\section{有限投射模}
\label{algebra-section-finite-projective-modules}

\begin{definition}
\label{algebra-definition-locally-free}
设 $R$ 为环，$M$ 为 $R$-模。
\begin{enumerate}
\item 若可用标准开集 $D(f_i)$（$i \in I$）覆盖 $\Spec(R)$，
使得对所有 $i \in I$，$M_{f_i}$ 都是自由 $R_{f_i}$-模，
则称 $M$ \textit{局部自由}。
\item 若可选取上述覆盖，使每个 $M_{f_i}$ 都有限自由，
则称 $M$ \textit{有限局部自由}。
\item 若可选取上述覆盖，使每个 $M_{f_i}$ 都同构于
$R_{f_i}^{\oplus r}$，则称 $M$ 是\textit{秩为 $r$ 的有限局部自由模}。
\end{enumerate}
\end{definition}

\noindent
注意，由引理 \ref{algebra-lemma-cover}，有限局部自由 $R$-模自动是有限表现的。
此外，若 $M$ 是环 $R$ 上秩为 $r$ 的有限局部自由模且 $R$ 非零，
则由引理 \ref{algebra-lemma-rank}，$r$ 唯一确定
（因为局部化 $R_{f_i}$ 中至少有一个是非零环）。

\begin{lemma}
\label{algebra-lemma-finite-projective}
设 $R$ 为环，$M$ 为 $R$-模。下列条件等价：
\begin{enumerate}
\item $M$ 是有限表现且在 $R$ 上平坦；
\item $M$ 是有限投射模；
\item $M$ 是某个有限自由 $R$-模的直和因子；
\item $M$ 是有限表现的，并且对所有
$\mathfrak p \in \Spec(R)$，局部化 $M_{\mathfrak p}$ 都自由；
\item $M$ 是有限表现的，并且对所有极大理想
$\mathfrak m \subset R$，局部化 $M_{\mathfrak m}$ 都自由；
\item $M$ 有限且局部自由；
\item $M$ 有限局部自由；
\item $M$ 有限，对每个素理想 $\mathfrak p$，模
$M_{\mathfrak p}$ 都自由，并且函数
$$
\rho_M : \Spec(R) \to \mathbf{Z}, \quad
\mathfrak p
\longmapsto
\dim_{\kappa(\mathfrak p)} M \otimes_R \kappa(\mathfrak p)
$$
在 Zariski 拓扑下局部常值。
\end{enumerate}
\end{lemma}

\begin{proof}
先假设 $M$ 是有限投射模，即条件 (2) 成立。
取满射 $R^n \to M$，并令 $K$ 为其核。
因为 $M$ 投射，短正合序列
$0 \to K \to R^n \to M \to 0$ 分裂。
所以 (2) $\Rightarrow$ (3)。由投射模的直和因子仍投射这一事实，
(3) $\Rightarrow$ (2)；参见引理 \ref{algebra-lemma-characterize-projective}。

\medskip\noindent
假设 (3) 成立，因而可写 $K \oplus M \cong R^{\oplus n}$。
于是 $K$ 是 $R^n$ 的直和因子，从而有限生成。
这说明 $M = R^{\oplus n}/K$ 是有限表现的。
换言之，(3) $\Rightarrow$ (1)。

\medskip\noindent
假设 $M$ 是有限表现且平坦的，即 (1) 成立。
我们证明 (7)。任选素理想 $\mathfrak p$，并取
$x_1, \ldots, x_r \in M$，使它们的像构成
$M \otimes_R \kappa(\mathfrak p)$ 的一组基。
由 Nakayama 引理（采用引理 \ref{algebra-lemma-NAK-localization} 的形式），
存在 $g \in R$、$g \not \in \mathfrak p$，使这些元素生成 $M_g$。
相应满射 $\varphi : R_g^{\oplus r} \to M_g$ 有下列两个性质：
(a) $\Ker(\varphi)$ 是有限 $R_g$-模（参见引理 \ref{algebra-lemma-extension}）；
(b) 由 $M_g$ 在 $R_g$ 上的平坦性，
$\Ker(\varphi) \otimes \kappa(\mathfrak p) = 0$
（参见引理 \ref{algebra-lemma-flat-tor-zero}）。
再次使用 Nakayama 引理，存在 $g' \in R_g$，使得
$\Ker(\varphi)_{g'} = 0$。换言之，$M_{gg'}$ 自由。

\medskip\noindent
有限局部自由模是有限模，参见引理 \ref{algebra-lemma-cover}，
所以 (7) $\Rightarrow$ (6)。显然 (6) $\Rightarrow$ (7)，
并且 (7) $\Rightarrow$ (8)。

\medskip\noindent
有限局部自由模是有限表现模，参见引理 \ref{algebra-lemma-cover}，
所以 (7) $\Rightarrow$ (4)。当然 (4) 蕴含 (5)。
由于平坦性可以局部检验（参见引理 \ref{algebra-lemma-flat-localization}），
可知 (5) 蕴含 (1)。此时我们已有
$$
\xymatrix{
(2) \ar@{<=>}[r] & (3) \ar@{=>}[r] & (1) \ar@{=>}[r] &
(7)  \ar@{<=>}[r] \ar@{=>}[rd] \ar@{=>}[d] & (6) \\
& & (5) \ar@{=>}[u] & (4) \ar@{=>}[l] & (8)
}
$$

\medskip\noindent
假设 $M$ 满足 (1)、(4)、(5)、(6) 和 (7)。
我们证明 (3)。只需说明 $M$ 投射，即说明 $\Hom_R(M, -)$ 正合。
设 $0 \to N'' \to N \to N'\to 0$ 是 $R$-模的短正合序列。
我们必须证明
$0 \to \Hom_R(M, N'') \to \Hom_R(M, N) \to
\Hom_R(M, N') \to 0$ 正合。
因为 $M$ 有限局部自由，存在覆盖
$\Spec(R) = \bigcup D(f_i)$，使 $M_{f_i}$ 有限自由。
由引理 \ref{algebra-lemma-hom-from-finitely-presented}，序列
$$
0 \to \Hom_R(M, N'')_{f_i} \to \Hom_R(M, N)_{f_i} \to
\Hom_R(M, N')_{f_i} \to 0
$$
等于
$0 \to \Hom_{R_{f_i}}(M_{f_i}, N''_{f_i}) \to
\Hom_{R_{f_i}}(M_{f_i}, N_{f_i}) \to
\Hom_{R_{f_i}}(M_{f_i}, N'_{f_i}) \to 0$。
该序列正合，因为 $M_{f_i}$ 自由，并且局部化序列
$0 \to N''_{f_i} \to N_{f_i} \to N'_{f_i} \to 0$
正合（局部化正合）。于是由引理 \ref{algebra-lemma-cover}，
$0 \to \Hom_R(M, N'') \to \Hom_R(M, N) \to
\Hom_R(M, N') \to 0$ 正合。

\medskip\noindent
最后，假设 (8) 成立。选取极大理想 $\mathfrak m \subset R$，
并取 $x_1, \ldots, x_r \in M$，使它们的像构成
$M \otimes_R \kappa(\mathfrak m) = M/\mathfrak mM$ 的一组
$\kappa(\mathfrak m)$-基。特别地，$\rho_M(\mathfrak m) = r$。
由 Nakayama 引理 \ref{algebra-lemma-NAK}，存在
$f \in R$、$f \not \in \mathfrak m$，使
$x_1, \ldots, x_r$ 在 $R_f$ 上生成 $M_f$。
由 $\rho_M$ 局部常值的假设，存在
$g \in R$、$g \not \in \mathfrak m$，使
$\rho_M$ 在 $D(g)$ 上恒等于 $r$。我们断言
$$
\Psi : R_{fg}^{\oplus r} \longrightarrow M_{fg}, \quad
(a_1, \ldots, a_r) \longmapsto \sum a_i x_i
$$
是同构。这个断言将说明 $M$ 有限局部自由，即 (7) 成立。
为证明该断言，由引理 \ref{algebra-lemma-characterize-zero-local}，
只需证明对所有 $\mathfrak p \in D(fg)$，诱导的局部化映射
$\Psi_{\mathfrak p} : R_{\mathfrak p}^{\oplus r} \to M_{\mathfrak p}$
都是同构。由 $f$ 的选取，$\Psi_{\mathfrak p}$ 是满射。
由假设 (8)，有
$M_{\mathfrak p} \cong R_{\mathfrak p}^{\oplus \rho_M(\mathfrak p)}$，
而由 $g$ 的选取，有 $\rho_M(\mathfrak p) = r$。
因此 $\Psi_{\mathfrak p}$ 给出满射
$R_{\mathfrak p}^{\oplus r} \to
M_{\mathfrak p} \cong R_{\mathfrak p}^{\oplus r}$，
从而由引理 \ref{algebra-lemma-fun}，它是同构。
（当然，最后这个事实也可由简单的矩阵论证得到。）
\end{proof}

\begin{lemma}
\label{algebra-lemma-finite-projective-reduced}
设 $R$ 为约化环，$M$ 为 $R$-模。
则引理 \ref{algebra-lemma-finite-projective} 中的等价条件还等价于
\begin{enumerate}
\item[(9)] $M$ 有限，并且函数
$\rho_M : \Spec(R) \to \mathbf{Z}$，$\mathfrak p \mapsto
\dim_{\kappa(\mathfrak p)} M \otimes_R \kappa(\mathfrak p)$
在 Zariski 拓扑下局部常值。
\end{enumerate}
\end{lemma}

\begin{proof}
选取极大理想 $\mathfrak m \subset R$，并取
$x_1, \ldots, x_r \in M$，使它们的像构成
$M \otimes_R \kappa(\mathfrak m) = M/\mathfrak mM$ 的一组
$\kappa(\mathfrak m)$-基。特别地，$\rho_M(\mathfrak m) = r$。
由 Nakayama 引理 \ref{algebra-lemma-NAK}，存在
$f \in R$、$f \not \in \mathfrak m$，使
$x_1, \ldots, x_r$ 在 $R_f$ 上生成 $M_f$。
由 $\rho_M$ 局部常值的假设，存在
$g \in R$、$g \not \in \mathfrak m$，使
$\rho_M$ 在 $D(g)$ 上恒等于 $r$。我们断言
$$
\Psi : R_{fg}^{\oplus r} \longrightarrow M_{fg}, \quad
(a_1, \ldots, a_r) \longmapsto \sum a_i x_i
$$
是同构。这个断言将说明 $M$ 有限局部自由，即 (7) 成立。
因为 $\Psi$ 是满射，只需证明 $\Psi$ 是单射。
由于 $R_{fg}$ 约化，由引理
\ref{algebra-lemma-reduced-ring-sub-product-fields}，
只需证明在 $R_{fg}$ 的所有极小素理想 $\mathfrak p$ 处局部化后，
$\Psi$ 都是单射。然而，由引理
\ref{algebra-lemma-minimal-prime-reduced-ring}，有
$R_\mathfrak p = \kappa(\mathfrak p)$；又有
$\rho_M(\mathfrak p) = r$，所以
$\Psi_\mathfrak p :
R_\mathfrak p^{\oplus r} \to M \otimes_R \kappa(\mathfrak p)$
是相同维数的有限维向量空间之间的满射，因而是同构。
\end{proof}

\begin{remark}
\label{algebra-remark-warning}
有限且在 $R$ 上平坦的 $R$-模并不自动投射。
一个反例是：取 $R = \mathcal{C}^\infty(\mathbf{R})$ 为
$\mathbf{R}$ 上无穷次可微函数构成的环，并取
$M = R_{\mathfrak m} = R/I$，其中
$\mathfrak m = \{f \in R \mid f(0) = 0\}$，且
$I = \{f \in R \mid \exists \epsilon, \epsilon > 0 :
f(x) = 0\ \forall x, |x| < \epsilon\}$。
\end{remark}

\begin{lemma}
\label{algebra-lemma-finite-flat-local}
（警告：参见注 \ref{algebra-remark-warning}。）
假设 $R$ 是局部环，$M$ 是有限平坦 $R$-模。
则 $M$ 有限自由。
\end{lemma}

\begin{proof}
这由平坦性的方程判据得到；参见引理 \ref{algebra-lemma-flat-eq}。
具体地，假设 $x_1, \ldots, x_r \in M$ 的像构成
$M/\mathfrak mM$ 的一组基。由 Nakayama 引理
\ref{algebra-lemma-NAK}，这些元素生成 $M$。
我们要证明 $x_i$ 之间不存在关系。为此，改为对 $n$ 归纳证明：
若 $x_1, \ldots, x_n \in M$ 在向量空间
$M/\mathfrak mM$ 中线性无关，则它们在 $R$ 上线性无关。

\medskip\noindent
归纳起点是：有 $x \in M$、$x \not\in \mathfrak mM$，
以及关系 $fx = 0$。由方程判据，存在
$y_j \in M$ 和 $a_j \in R$，使得
$x = \sum a_j y_j$，并且对所有 $j$ 都有 $fa_j = 0$。
因为 $x \not\in \mathfrak mM$，至少有一个 $a_j$ 是单位，
故 $f = 0 $。

\medskip\noindent
假设 $\sum f_i x_i$ 是 $x_1, \ldots, x_n$ 之间的一个关系。
由 $x_i$ 的选取，有 $f_i \in \mathfrak m$。
根据平坦性的方程判据，存在 $a_{ij} \in R$ 与
$y_j \in M$，使得 $x_i = \sum a_{ij} y_j$，且
$\sum f_i a_{ij} = 0$。由于 $x_n \not \in \mathfrak mM$，
至少对一个 $j$ 有 $a_{nj}\not\in \mathfrak m$。
由 $\sum f_i a_{ij} = 0$，得到
$f_n = \sum_{i = 1}^{n-1} (-a_{ij}/a_{nj}) f_i$。
现在关系 $\sum f_i x_i = 0$ 可改写为
$\sum_{i = 1}^{n-1} f_i( x_i + (-a_{ij}/a_{nj}) x_n) = 0$。
注意，元素 $x_i + (-a_{ij}/a_{nj}) x_n$ 在
$M/\mathfrak mM$ 中的像是 $n-1$ 个线性无关元素。
由归纳假设，所有 $f_i$（$i \leq n-1$）都必须为零；
而且 $f_n = \sum_{i = 1}^{n-1} (-a_{ij}/a_{nj}) f_i$。
这完成归纳步骤。
\end{proof}

\begin{lemma}
\label{algebra-lemma-finite-projective-descends}
设 $R \to S$ 是局部环之间的平坦局部同态，
$M$ 是有限 $R$-模。则 $M$ 在 $R$ 上有限投射，
当且仅当 $M \otimes_R S$ 在 $S$ 上有限投射。
\end{lemma}

\begin{proof}
由引理 \ref{algebra-lemma-finite-projective}，
在局部环上有限投射等价于有限自由。
假设 $M \otimes_R S$ 是有限自由 $S$-模。
选取 $x_1, \ldots, x_r \in M$，使它们在
$M/\mathfrak m_RM$ 中的像构成一组
$\kappa(\mathfrak m)$-基。于是
$x_1 \otimes 1, \ldots, x_r \otimes 1$
构成 $M \otimes_R S$ 的一组基。
这蕴含映射 $R^{\oplus r} \to M, (a_i) \mapsto \sum a_i x_i$
与 $S$ 作张量积后成为同构。
由 $R \to S$ 的忠实平坦性（参见引理 \ref{algebra-lemma-local-flat-ff}），
该映射本身是同构。
\end{proof}

\begin{lemma}
\label{algebra-lemma-locally-free-semi-local-free}
设 $R$ 为半局部环，$M$ 为有限局部自由模。
若 $M$ 的秩为常数，则 $M$ 自由。特别地，若
$R$ 的谱连通，则 $M$ 自由。
\end{lemma}

\begin{proof}
略。提示：先利用秩为常数，证明对 $R$ 的所有极大理想
$\mathfrak m_1, \ldots, \mathfrak m_n$，
$M/\mathfrak m_iM$ 都有相同的维数 $d$。
然后利用中国剩余定理，证明存在元素
$x_1, \ldots, x_d \in M$，它们对每个
$M/\mathfrak m_iM$ 都构成一组基。最后证明
$x_1, \ldots, x_d$ 是 $M$ 的一组基。
\end{proof}

\noindent
下面这个技术引理将在群胚一章中使用。

\begin{lemma}
\label{algebra-lemma-semi-local-module-basis-in-submodule}
设 $R$ 为极大理想为 $\mathfrak m$、剩余域无限的局部环。
设 $R \to S$ 为环映射，$M$ 为 $S$-模，
$N \subset M$ 为 $R$-子模。假设
\begin{enumerate}
\item $S$ 是半局部的，并且 $\mathfrak mS$ 包含于 $S$ 的
Jacobson 根中；
\item $M$ 是有限自由 $S$-模；
\item $N$ 作为 $S$-模生成 $M$。
\end{enumerate}
则 $N$ 包含 $M$ 的一组 $S$-基。
\end{lemma}

\begin{proof}
假设 $M$ 是秩为 $n$ 的自由模，并设 $I \subset S$ 为 Jacobson 根。
由 Nakayama 引理 \ref{algebra-lemma-NAK}，元素序列
$m_1, \ldots, m_n$ 是 $M$ 的一组基，当且仅当
$\overline{m}_i \in M/IM$ 生成 $M/IM$。
因此可以把 $M$ 替换为 $M/IM$，把 $N$ 替换为
$N/(N \cap IM)$，把 $R$ 替换为 $R/\mathfrak m$，
并把 $S$ 替换为 $S/IS$。在此情形，$S$ 是域的有限积
$S = k_1 \times \ldots \times k_r$，且
$M = k_1^{\oplus n} \times \ldots \times k_r^{\oplus n}$。
$N \subset M$ 作为 $S$-模生成 $M$ 这一事实意味着：
存在 $x_j \in N$，使得某个线性组合
$\sum a_j x_j$（$a_j \in S$）在每个因子
$k_i^{\oplus n}$ 中都有非零分量。
因为 $R = k$ 是无限域，这又意味着存在某个线性组合
$y = \sum c_j x_j$（$c_j \in k$），
它在每个因子中都有非零分量。
因此 $y \in N$ 生成自由直和因子 $Sy \subset M$。
对 $n$ 作归纳，可知结论对 $M/Sy$ 及其子模
$\overline{N} = N/(N \cap Sy)$ 成立。
换言之，存在 $\overline{N}$ 中的
$\overline{y}_2, \ldots, \overline{y}_n$，
它们（自由地）生成 $M/Sy$。
于是 $y, y_2, \ldots, y_n$（自由地）生成 $M$，结论得证。
\end{proof}

\begin{lemma}
\label{algebra-lemma-evaluation-map-iso-finite-projective}
设 $R$ 为环，$L$、$M$、$N$ 为 $R$-模。若 $M$ 有限投射，
则典范映射
$$
\Hom_R(M, N) \otimes_R L \to \Hom_R(M, N \otimes_R L)
$$
是同构。
\end{lemma}

\begin{proof}
由引理 \ref{algebra-lemma-finite-projective}，$M$ 既有限表现又有限局部自由。
由引理 \ref{algebra-lemma-hom-from-finitely-presented} 与
\ref{algebra-lemma-tensor-product-localization}，箭头两边的构造都与局部化交换。
由引理 \ref{algebra-lemma-cover}，可在局部化后检验该映射是否为同构。
因此可以假设 $M$ 有限自由；此时引理立即成立。
\end{proof}

\section{由模映射定义的开集}
\label{algebra-section-loci-maps}

\noindent
给定模映射为满射或同构的素理想所成集合有时是开的。
在有限投射模的情形，可以考察映射的秩。

\begin{lemma}
\label{algebra-lemma-map-between-finite}
设 $R$ 为环，$\varphi : M \to N$ 为 $R$-模映射，
其中 $N$ 是有限 $R$-模。则有等式
\begin{align*}
U & = \{\mathfrak p \subset R \mid
\varphi_{\mathfrak p} : M_{\mathfrak p} \to N_{\mathfrak p}
\text{ 为满射}\} \\
& = \{\mathfrak p \subset R \mid
\varphi \otimes \kappa(\mathfrak p) :
M \otimes \kappa(\mathfrak p) \to N \otimes \kappa(\mathfrak p)
\text{ 为满射}\}
\end{align*}
并且 $U$ 是 $\Spec(R)$ 的开子集。此外，对任意满足
$D(f) \subset U$ 的 $f \in R$，映射 $M_f \to N_f$ 都是满射。
\end{lemma}

\begin{proof}
显示公式中的等式来自 Nakayama 引理。Nakayama 引理也蕴含
$U$ 是开的；参见引理 \ref{algebra-lemma-NAK}，特别是其中 (3)。
若 $D(f) \subset U$，则 $M_f \to N_f$ 在 $R_f$ 的所有素理想处
局部化后都是满射，因而由引理
\ref{algebra-lemma-characterize-zero-local}，它本身是满射。
\end{proof}

\begin{lemma}
\label{algebra-lemma-map-between-finitely-presented}
设 $R$ 为环，$\varphi : M \to N$ 为 $R$-模映射，
其中 $M$ 有限而 $N$ 有限表现。则
$$
U = \{\mathfrak p \subset R \mid
\varphi_{\mathfrak p} : M_{\mathfrak p} \to N_{\mathfrak p}
\text{ 为同构}\}
$$
是 $\Spec(R)$ 的开子集。
\end{lemma}

\begin{proof}
设 $\mathfrak p \in U$。选取表现
$N = R^{\oplus n}/\sum_{j = 1, \ldots, m} R k_j$。
以 $e_i$ 表示 $R^{\oplus n}$ 的第 $i$ 个基向量在 $N$ 中的像。
对每个 $i \in \{1, \ldots, n\}$，选取元素
$m_i \in M_{\mathfrak p}$，使得对某个
$f_i \in R$、$f_i \not \in \mathfrak p$ 有
$\varphi(m_i) = f_i e_i$。这可以做到，因为
$\varphi_{\mathfrak p}$ 是同构。令 $f = f_1 \ldots f_n$，
并令 $\psi : R_f^{\oplus n} \to M_f$ 为把第 $i$ 个基向量
映到 $m_i/f_i$ 的映射。注意，$\varphi_f \circ \psi$
是给定映射 $R^{\oplus n} \to N$ 在 $f$ 处的局部化。
因为 $\varphi_{\mathfrak p}$ 是同构，$\psi(k_j)$ 是 $M$ 中
映到 $M_{\mathfrak p}$ 中零元的元素。因此存在
$g_j \in R$、$g_j \not \in \mathfrak p$，使得
$g_j \psi(k_j) = 0$。令 $g = g_1 \ldots g_m$，
则 $\psi_g$ 经 $N_{fg}$ 分解，给出映射
$\chi : N_{fg} \to M_{fg}$。按构造，$\chi$ 是
$\varphi_{fg}$ 的右逆。于是 $\chi_\mathfrak p$ 是同构。
由引理 \ref{algebra-lemma-map-between-finite}，存在
$h \in R$、$h \not \in \mathfrak p$，使映射
$\chi_h : N_{fgh} \to M_{fgh}$ 为满射。
故 $\varphi_{fgh}$ 与 $\chi_h$ 互为逆映射，
这蕴含 $D(fgh) \subset U$，如所需。
\end{proof}

\begin{lemma}
\label{algebra-lemma-finitely-presented-localization-free}
设 $R$ 为环，$\mathfrak p \subset R$ 为素理想，
$M$ 为有限表现 $R$-模。若 $M_\mathfrak p$ 自由，
则存在 $f \in R$、$f \not \in \mathfrak p$，
使得 $M_f$ 是自由 $R_f$-模。
\end{lemma}

\begin{proof}
选取一组基 $x_1, \ldots, x_n \in M_\mathfrak p$。
可以选取 $f \in R$、$f \not \in \mathfrak p$，
使每个 $x_i$ 都是某个 $y_i \in M_f$ 的像。
用 $f^m y_i$（$m \gg 0$）替换 $y_i$ 后，可以假设
$y_i \in M$。事实上，这会把 $x_1, \ldots, x_n$ 替换为
$f^mx_1, \ldots, f^mx_n$；因为 $f$ 在 $R_\mathfrak p$
中的像是单位，它们仍是一组基。因此得到 $R$-模同态
$\varphi = (y_1, \ldots, y_n) : R^{\oplus n} \to M$，
其在 $\mathfrak p$ 处的局部化是同构。
由引理 \ref{algebra-lemma-map-between-finitely-presented}，
可以找到 $f \in R$、$f \not \in \mathfrak p$，使得对所有满足
$f \not \in \mathfrak q$ 的素理想 $\mathfrak q \subset R$，
$\varphi_\mathfrak q$ 都是同构。
于是由引理 \ref{algebra-lemma-characterize-zero-local}，
$\varphi_f$ 是同构，证明完成。
\end{proof}

\begin{lemma}
\label{algebra-lemma-cokernel-flat}
设 $R$ 为环，$\varphi : P_1 \to P_2$ 为有限投射模之间的映射。则
\begin{enumerate}
\item 使 $\varphi \otimes \kappa(\mathfrak p)$ 为单射的素理想
$\mathfrak p \in \Spec(R)$ 所成集合 $U$ 是开的；并且对任意满足
$D(f) \subset U$ 的 $f\in R$，有
\begin{enumerate}
\item $P_{1, f} \to P_{2, f}$ 是单射；
\item 模 $\Coker(\varphi)_f$ 在 $R_f$ 上有限投射。
\end{enumerate}
\item 使 $\varphi \otimes \kappa(\mathfrak p)$ 为满射的素理想
$\mathfrak p \in \Spec(R)$ 所成集合 $W$ 是开的；并且对任意满足
$D(f) \subset W$ 的 $f\in R$，有
\begin{enumerate}
\item $P_{1, f} \to P_{2, f}$ 是满射；
\item 模 $\Ker(\varphi)_f$ 在 $R_f$ 上有限投射。
\end{enumerate}
\item 使 $\varphi \otimes \kappa(\mathfrak p)$ 为同构的素理想
$\mathfrak p \in \Spec(R)$ 所成集合 $V$ 是开的；并且对任意满足
$D(f) \subset V$ 的 $f\in R$，映射
$\varphi : P_{1, f} \to P_{2, f}$ 是 $R_f$-模的同构。
\end{enumerate}
\end{lemma}

\begin{proof}
为证明集合 $U$ 是开的，可以在 $\Spec(R)$ 上局部工作。
因此可把 $R$ 替换为适当的局部化，并假设
$P_1 = R^{n_1}$ 且 $P_2 = R^{n_2}$；参见引理
\ref{algebra-lemma-finite-projective}。在此情形，
$\varphi \otimes \kappa(\mathfrak p)$ 为单射，当且仅当
$n_1 \leq n_2$，并且 $\varphi$ 的矩阵的某个
$n_1 \times n_1$ 子式 $f$ 在 $\kappa(\mathfrak p)$ 中可逆。
因此 $D(f) \subset U$。这个论证也说明，对
$\mathfrak p \in U$，映射
$P_{1, \mathfrak p} \to P_{2, \mathfrak p}$ 是单射。

\medskip\noindent
现在假设 $D(f) \subset U$。由上一段的说明以及引理
\ref{algebra-lemma-characterize-zero-local}，
$P_{1, f} \to P_{2, f}$ 是单射，即 (1)(a) 成立。
由引理 \ref{algebra-lemma-finite-projective}，为证明 (1)(b)，
只需证明 $\Coker(\varphi)$ 在 $D(f)$ 上局部有限投射。
因此如上所见，可以假设 $P_1 = R^{n_1}$、$P_2 = R^{n_2}$，
并且 $\varphi$ 的矩阵的某个子式在 $R$ 中可逆。
若该子式对应于 $R^{n_2}$ 的前 $n_1$ 个基向量，
则利用其余 $n_2 - n_1$ 个基向量得到映射
$R^{n_2 - n_1} \to R^{n_2} \to \Coker(\varphi)$，
容易看出它是同构。

\medskip\noindent
$W$ 的开性以及在 $D(f) \subset W$ 时的 (2)(a)
来自引理 \ref{algebra-lemma-map-between-finite}。
因为 $P_{2, f}$ 在 $R_f$ 上投射，映射
$\varphi_f : P_{1, f} \to P_{2, f}$ 有一个截面，
从而 $\Ker(\varphi)_f$ 是 $P_{2, f}$ 的直和因子。
因此 $\Ker(\varphi)_f$ 有限投射，故 (2)(b) 也成立。

\medskip\noindent
显然 $V = U \cap W$ 是开的，而 (3) 中的另一个断言
来自 (1)(a) 和 (2)(a)。
\end{proof}

\section{模投射性的忠实平坦下降}
\label{algebra-section-ffdescent-projectivity-introduction}

\medskip\noindent
接下来的几节依照 Raynaud 和 Gruson \cite{GruRay}，
证明模的投射性沿忠实平坦环映射下降。
证明思路是采用 Kaplansky 式的 d\'evissage（逐层剖析）\cite{Kaplansky}，
把问题约化到可数生成模的情形。给定模 $M$ 的一个性质良好的滤过，
d\'evissage 使我们可以把 $M$ 表为各滤过子模的相继商之直和
（参见第 \ref{algebra-section-transfinite-devissage} 节）。
利用这一技巧，我们证明投射模是可数生成模的直和
（定理 \ref{algebra-theorem-projective-direct-sum}）。
为证明可数生成模投射性的下降，我们引入模上的一个
“Mittag-Leffler”条件，证明可数生成模投射当且仅当它既平坦又是
Mittag-Leffler 模（定理 \ref{algebra-theorem-projectivity-characterization}），
然后证明成为 Mittag-Leffler 模这一性质可以下降
（引理 \ref{algebra-lemma-ffdescent-ML}）。最后，给定任意模 $M$，
若它沿某个忠实平坦环映射的基变换是投射模，
我们就用相继商为可数生成投射模的子模来滤过 $M$，
再由 d\'evissage 断定 $M$ 是投射模的直和，
因而自身也是投射模（定理 \ref{algebra-theorem-ffdescent-projectivity}）。

\medskip\noindent
我们指出，\cite{GruRay} 中投射性的忠实平坦下降之证明有一个错误。
该文由投射性沿一类更一般的环映射的下降，推出其沿忠实平坦环映射的下降
（\cite[Example 3.1.4(1) of Part II]{GruRay}）。
然而，沿这类更一般映射下降的证明并不正确。
在 \cite{G} 中，Gruson 解释了问题所在，但没有为这里关心的情形给出修补。
填补投射性忠实平坦下降证明中的这个缺口，归结为证明
Mittag-Leffler 模这一性质沿忠实平坦环映射下降。
我们将在引理 \ref{algebra-lemma-ffdescent-ML} 中完成这一点。

\section{平坦性的刻画}
\label{algebra-section-characterize-flatness}

\noindent
本节讨论平坦性的判据。本节的主要结果是下面的 Lazard 定理
（定理 \ref{algebra-theorem-lazard}）：平坦模是有限自由模所成有向系的余极限。
我们提醒读者回顾“平坦性的方程判据”，参见引理
\ref{algebra-lemma-flat-eq}。事实表明，可以把它整理成一个看似强得多的性质。

\begin{lemma}
\label{algebra-lemma-flat-factors-free}
设 $M$ 为 $R$-模。下列条件等价：
\begin{enumerate}
\item $M$ 平坦。
\item 若 $f: R^n \to M$ 是模映射且 $x \in \Ker(f)$，
则存在模映射 $h: R^n \to R^m$ 和 $g: R^m \to M$，
使得 $f = g \circ h$ 且 $x \in \Ker(h)$。
\item 假设 $f: R^n \to M$ 是模映射，
$N \subset \Ker(f)$ 是任意子模，而 $h: R^n \to R^{m}$
是满足 $N \subset \Ker(h)$ 且 $f$ 经 $h$ 分解的映射。
则给定任意 $x \in \Ker(f)$，都可找到映射
$h': R^n \to R^{m'}$，使得
$N + Rx \subset \Ker(h')$ 且 $f$ 经 $h'$ 分解。
\item 若 $f: R^n \to M$ 是模映射，而
$N \subset \Ker(f)$ 是有限生成子模，则存在模映射
$h: R^n \to R^m$ 和 $g: R^m \to M$，
使得 $f = g \circ h$ 且 $N \subset \Ker(h)$。
\end{enumerate}
\end{lemma}

\begin{proof}
(1) 与 (2) 的等价只是平坦性方程判据的另一种表述\footnote{
事实上，模映射 $f : R^n \to M$ 对应于选取 $M$ 中的元素
$x_1, x_2, \ldots, x_n$（即标准基向量
$e_1, e_2, \ldots, e_n$ 的像）；此外，元素
$x \in \Ker(f)$ 对应于这些 $x_1, x_2, \ldots, x_n$
之间的一个关系（即关系 $\sum_i f_i x_i = 0$，其中
$f_i$ 是 $x$ 的坐标）。模映射 $h$（表示成一个
$m \times n$ 矩阵）对应于引理 \ref{algebra-lemma-flat-eq} 中的矩阵
$(a_{ij})$，而引理 \ref{algebra-lemma-flat-eq} 中的 $y_j$
是 $R^m$ 的标准基向量在 $g$ 下的像。}。
为证明 (2) 蕴含 (3)，设 $g: R^m \to M$ 是使
$f$ 分解为 $f = g \circ h$ 的映射。
由 (2)，可找到 $h'': R^m \to R^{m'}$，使得
$h''$ 消去 $h(x)$，且 $g: R^m \to M$ 经 $h''$ 分解。
此时取 $h' = h'' \circ h$ 即可。
(3) 蕴含 (4) 可对 (4) 中 $N \subset \Ker(f)$ 的生成元个数作归纳得到。
显然 (4) 蕴含 (2)。
\end{proof}

\begin{lemma}
\label{algebra-lemma-flat-factors-fp}
设 $M$ 为 $R$-模。则 $M$ 平坦，当且仅当下列条件成立：
若 $P$ 是有限表现 $R$-模，$f: P \to M$ 是模映射，
则存在有限自由 $R$-模 $F$ 以及模映射
$h: P \to F$ 和 $g: F \to M$，使得 $f = g \circ h$。
\end{lemma}

\begin{proof}
这只是引理 \ref{algebra-lemma-flat-factors-free} 中条件 (4) 的另一种表述。
\end{proof}

\begin{lemma}
\label{algebra-lemma-flat-surjective-hom}
设 $M$ 为 $R$-模。则 $M$ 平坦，当且仅当下列条件成立：
对每个有限表现 $R$-模 $P$，若 $N \to M$ 是满的 $R$-模映射，
则诱导映射 $\Hom_R(P, N) \to \Hom_R(P, M)$ 是满射。
\end{lemma}

\begin{proof}
先假设 $M$ 平坦。我们必须证明：若 $P$ 有限表现，
则给定映射 $f: P \to M$，它经映射 $N \to M$ 分解。
由引理 \ref{algebra-lemma-flat-factors-fp}，$f$ 经某个映射
$F \to M$ 分解，其中 $F$ 有限自由。
因为 $F$ 自由，这个映射经 $N \to M$ 分解。
所以 $f$ 经 $N \to M$ 分解。

\medskip\noindent
反过来，假设本引理的条件成立。
设 $f: P \to M$ 是从有限表现模 $P$ 出发的映射。
选取自由模 $N$ 以及到 $M$ 上的满射 $N \to M$。
则 $f$ 经 $N \to M$ 分解；又因为 $P$ 有限生成，
$f$ 经 $N$ 的某个有限自由子模分解。
因此 $M$ 满足引理 \ref{algebra-lemma-flat-factors-fp} 的条件，故平坦。
\end{proof}

\begin{theorem}[Lazard 定理]
\label{algebra-theorem-lazard}
设 $M$ 为 $R$-模。则 $M$ 平坦，当且仅当它是有限自由
$R$-模所成有向系的余极限。
\end{theorem}

\begin{proof}
平坦模所成有向系的余极限平坦，因为取有向余极限正合且与张量积交换。
因此若 $M$ 是有限自由模所成有向系的余极限，则 $M$ 平坦。

\medskip\noindent
为证明反向蕴含，先回顾任意模 $M$ 都可以如下写成有限表现模所成
有向系的余极限。对某个集合 $I$ 选取满射
$f: R^I \to M$，并令 $K$ 为其核。
令 $E$ 为有序对 $(J, N)$ 所成集合，其中 $J$ 是 $I$ 的有限子集，
$N$ 是 $R^J \cap K$ 的有限生成子模。
规定 $(J, N) \leq (J', N')$ 当且仅当
$J \subset J'$ 且 $N \subset N'$，从而把 $E$ 变成有向偏序集。
对 $e = (J, N)$ 定义 $M_e = R^J/N$；当 $e \leq e'$ 时，
定义 $f_{ee'}: M_e \to M_{e'}$ 为自然映射。
于是 $(M_e, f_{ee'})$ 是有向系，并且自然映射
$f_e: M_e \to M$ 诱导同构
$\colim_{e
\in E} M_e \xrightarrow{\cong} M$。

\medskip\noindent
现在假设 $M$ 平坦。令 $I = M \times \mathbf{Z}$，
以 $(x_i)$ 表示 $R^{I}$ 的典范基；在上面的讨论中，
取 $f: R^I \to M$ 为把 $x_i$ 映到 $i$ 在 $M$ 上的投影的映射。
为证明定理，只需说明满足 $M_e$ 自由的 $e \in E$
构成 $E$ 的共尾子集。任取 $e = (J, N) \in E$。
由引理 \ref{algebra-lemma-flat-factors-fp}，存在有限自由模 $F$ 以及映射
$h: R^J/N \to F$ 和 $g: F \to M$，使自然映射
$f_e: R^J/N \to M$ 分解为
$R^J/N \xrightarrow{h} F \xrightarrow{g} M$。
我们将把 $F$ 实现为某个 $e' \geq e$ 的 $M_{e'}$。

\medskip\noindent
设 $\{ b_1, \ldots, b_n \}$ 是 $F$ 的一组有限基。
选取 $n$ 个互异元素 $i_1, \ldots, i_n \in I$，使得对所有
$\ell$ 都有 $i_{\ell} \notin J$，并使 $x_{i_{\ell}}$
在 $f: R^I \to M$ 下的像等于 $b_{\ell}$ 在
$g: F \to M$ 下的像。这可以做到，因为依我们对 $I$ 的选取，
$M$ 的每个元素都可对无穷多个互异的 $i \in I$ 写成 $f(x_i)$。
现在令 $J' = J \cup \{i_1, \ldots , i_n \}$，并如下定义
$R^{J'} \to F$：当 $i \in J$ 时令 $x_i \mapsto h(x_i)$，
而当 $\ell = 1, \ldots, n$ 时令
$x_{i_{\ell}} \mapsto b_{\ell}$。令
$N' = \Ker(R^{J'} \to F)$。注意：
\begin{enumerate}
\item 方块
$$
\xymatrix{
R^{J'} \ar[r] \ar@{^{(}->}[d] & F \ar[d]^{g} \\
R^{I} \ar[r]_{f} & M
}
$$
交换，
%$R^{J'} \to F$ factors $f: R^I \to M$,
因此 $N' \subset K = \Ker(f)$；
\item $R^{J'} \to F$ 是到有限自由模上的满射，因而分裂，
所以 $N'$ 有限生成；
\item $J \subset J'$ 且 $N \subset N'$。
\end{enumerate}
由 (1) 和 (2)，$e' = (J', N')$ 属于 $E$；
由 (3)，$e' \geq e$；而按构造
$M_{e'} = R^{J'}/N' \cong F$ 自由。
\end{proof}

\section{普遍单射模映射}
\label{algebra-section-universally-injective}

\noindent
下面讨论普遍单射模映射；在某种意义上，它们与平坦模相互补充
（见引理 \ref{algebra-lemma-flat-universally-injective}）。
我们遵循 Lazard 的学位论文 \cite{Autour}；另见 \cite{Lam}。

\begin{definition}
\label{algebra-definition-universally-injective}
设 $f: M \to N$ 是 $R$-模映射。若对每个 $R$-模 $Q$，映射
$f \otimes_R \text{id}_Q: M \otimes_R Q \to N \otimes_R Q$
都是单射，则称 $f$ 为{\it 普遍单射}。若 $R$-模序列
$0 \to M_1 \to M_2 \to M_3 \to 0$ 正合，且
$M_1 \to M_2$ 普遍单射，则称该序列{\it 普遍正合}。
\end{definition}

\begin{example}
\label{algebra-example-universally-exact}
普遍正合序列的例子如下。
\begin{enumerate}
\item 分裂短正合序列是普遍正合的，因为张量运算与取直和可交换。
\item 普遍正合序列的有向系之余极限仍然普遍正合。这是因为取有向余极限
是正合的，而张量运算与取余极限可交换。特别地，分裂正合序列的有向系
之余极限是普遍正合的。下面将看到，反过来，每个普遍正合序列都以这种
方式产生。
\end{enumerate}
\end{example}

\noindent
下面给出短正合序列普遍正合的若干判据。它们类似于上面给出的平坦性判据。
以下 (3)--(6) 分别对应引理 \ref{algebra-lemma-flat-eq}、
\ref{algebra-lemma-flat-factors-free}、\ref{algebra-lemma-flat-surjective-hom}
和定理 \ref{algebra-theorem-lazard} 中的平坦性判据。

\begin{theorem}
\label{algebra-theorem-universally-exact-criteria}
设
$$
0 \to M_1 \xrightarrow{f_1} M_2 \xrightarrow{f_2} M_3 \to 0
$$
是 $R$-模的正合序列。以下条件等价：
\begin{enumerate}
\item 序列 $0 \to M_1 \to M_2 \to M_3
\to 0$ 普遍正合。
\item 对每个有限表示的 $R$-模 $Q$，序列
$$
0 \to M_1 \otimes_R Q \to M_2 \otimes_R Q \to
M_3 \otimes_R Q \to 0
$$
正合。
\item 给定元素 $x_i \in M_1$ $(i = 1, \ldots, n)$、
$y_j \in M_2$ $(j = 1, \ldots, m)$ 和
$a_{ij} \in R$ $(i = 1, \ldots, n, j = 1, \ldots, m)$，并且对所有 $i$ 都有
$$
f_1(x_i) = \sum\nolimits_j a_{ij} y_j,
$$
则存在 $z_j \in M_1$ $(j =1, \ldots, m)$，使得对所有 $i$ 都有
$$
x_i = \sum\nolimits_j a_{ij} z_j .
$$
\item 给定 $R$-模映射的交换图
$$
\xymatrix{
R^n \ar[r] \ar[d] &  R^m \ar[d] \\
M_1 \ar[r]^{f_1}        &  M_2
}
$$
其中 $m$ 和 $n$ 是整数，则存在映射 $R^m \to M_1$，使上方三角形交换。
\item 对每个有限表示的 $R$-模 $P$，$R$-模映射
$\Hom_R(P, M_2) \to \Hom_R(P, M_3)$ 都是满射。
\item 序列 $0 \to M_1 \to M_2 \to M_3
\to 0$ 是如下形式的分裂正合序列之某个有向系的余极限：
$$
0 \to M_{1} \to M_{2, i} \to M_{3, i} \to 0
$$
其中各 $M_{3, i}$ 都是有限表示的。
\end{enumerate}
\end{theorem}

\begin{proof}
显然 (1) 蕴含 (2)。

\medskip\noindent
下面证明 (2) 蕴含 (3)。设 $f_1(x_i) = \sum_j a_{ij} y_j$ 是 (3)
中的关系。设 $(d_j)$ 是 $R^m$ 的一组基，$(e_i)$ 是 $R^n$ 的一组基，
并令映射 $R^m \to R^n$ 由
$d_j \mapsto \sum_i a_{ij} e_i$ 给出。令 $Q$ 为
$R^m \to R^n$ 的余核。将 $R^m \to R^n \to Q \to 0$
分别张量成上下两行，并用映射 $f_1: M_1 \to M_2$
连接所得序列，得到交换图
$$
\xymatrix{
M_1^{\oplus m} \ar[r] \ar[d] & M_1^{\oplus n} \ar[r] \ar[d] & M_1 \otimes_R Q
\ar[r] \ar[d] & 0 \\
M_2^{\oplus m} \ar[r] & M_2^{\oplus n} \ar[r] & M_2 \otimes_R Q \ar[r] & 0
}
$$
其中 $M_1^{\oplus m} \to M_1^{\oplus n}$ 由
$$
(z_1, \ldots, z_m) \mapsto
(\sum\nolimits_j a_{1j} z_j, \ldots, \sum\nolimits_j a_{nj} z_j),
$$
给出，而 $M_2^{\oplus m} \to M_2^{\oplus n}$ 的定义相同。
我们要证明 $x = (x_1, \ldots, x_n) \in M_1^{\oplus n}$
属于 $M_1^{\oplus m} \to M_1^{\oplus n}$ 的像。由 (2)，映射
$M_1 \otimes Q \to M_2 \otimes Q$ 是单射，故由上行的正合性，
只须证明 $x$ 在 $M_2 \otimes Q$ 中的像为 $0$；再由下行的正合性，
只须证明 $x$ 在 $M_2^{\oplus n}$ 中的像属于
$M_2^{\oplus m} \to M_2^{\oplus n}$ 的像。这由假设成立。

\medskip\noindent
条件 (4) 只是把 (3) 改写成图的形式。

\medskip\noindent
下面证明 (4) 蕴含 (5)。设 $\varphi : P \to M_3$ 是从有限表示的
$R$-模 $P$ 出发的映射。我们必须证明 $\varphi$ 可提升为映射
$P \to M_2$。选取 $P$ 的一个表示
$$
R^n \xrightarrow{g_1} R^m \xrightarrow{g_2} P \to 0.
$$
利用 $R^n$ 和 $R^m$ 的自由性，可以先构造 $h_2: R^m \to M_2$，
再构造 $h_1: R^n \to M_1$，使下图交换：
$$
\xymatrix{
         & R^n \ar[r]^{g_1} \ar[d]^{h_1} & R^m \ar[r]^{g_2} \ar[d]^{h_2} & P
\ar[r] \ar[d]^{\varphi} & 0 \\
0 \ar[r] & M_1 \ar[r]^{f_1} & M_2 \ar[r]^{f_2} & M_3 \ar[r] & 0 .
}
$$
由 (4)，存在映射 $k_1: R^m \to M_1$，使得
$k_1 \circ g_1 = h_1$。现在定义 $h'_2: R^m \to M_2$ 为
$h_2' = h_2 - f_1 \circ k_1$。于是
$$
h'_2 \circ g_1 = h_2 \circ g_1 - f_1 \circ k_1 \circ g_1 =
h_2 \circ g_1 - f_1 \circ h_1 = 0 .
$$
因此取商后，$h'_2$ 定义映射 $\varphi': P \to M_2$，满足
$\varphi' \circ g_2 = h_2'$。用图表示即为
$$
\xymatrix{
R^m \ar[r]^{g_2} \ar[d]_{h'_2} & P \ar[d]^{\varphi} \ar[dl]_{\varphi'} \\
M_2 \ar[r]^{f_2} & M_3.
}
$$
其中上方三角形交换。我们断言 $\varphi'$ 就是所求提升，即
$f_2 \circ \varphi' = \varphi$。由定义，
$$
f_2 \circ \varphi' \circ g_2 = f_2 \circ h'_2 =
f_2 \circ h_2 - f_2 \circ f_1 \circ k_1 = f_2 \circ h_2 =
\varphi \circ g_2.
$$
由于 $g_2$ 是满射，证明完成。

\medskip\noindent
现在证明 (5) 蕴含 (6)。把 $M_{3}$ 写成有限表示模
$M_{3, i}$ 的一个有向系的余极限，见引理
\ref{algebra-lemma-module-colimit-fp}。令 $M_{2, i}$ 为 $M_{3, i}$ 与
$M_{2}$ 在 $M_{3}$ 上的纤维积；按定义，它是 $M_2 \times M_{3, i}$
中由两个投影在 $M_3$ 中的像相等的元素组成的子模。令
$M_{1, i}$ 为投影 $M_{2, i} \to M_{3, i}$ 的核。
于是得到正合序列的有向系
$$
0 \to M_{1, i} \to M_{2, i} \to M_{3, i} \to 0,
$$
并且对每个 $i$ 都有正合序列的映射
$$
\xymatrix{
0  \ar[r] & M_{1, i} \ar[d] \ar[r]  &  M_{2, i}  \ar[r] \ar[d] & M_{3, i} \ar[d]
\ar[r] & 0 \\
0  \ar[r] & M_{1} \ar[r]  &  M_{2}  \ar[r] & M_{3} \ar[r] & 0
}
$$
且这些映射与有向系相容。由纤维积 $M_{2, i}$ 的定义可知，
映射 $M_{1, i} \to M_1$ 是同构。由 (5)，存在提升
$M_{3, i} \to M_{2}$ 的映射 $M_{3, i} \to M_3$；再由纤维积的
泛性质，这给出 $M_{2, i} \to M_{3, i}$ 的一个截面。因此序列
$$
0 \to M_{1, i} \to M_{2, i} \to M_{3, i} \to 0
$$
分裂。取余极限，得到交换图
$$
\xymatrix{
0  \ar[r] & \colim M_{1, i} \ar[d]^{\cong} \ar[r]  &  \colim M_{2, i}
 \ar[r]
\ar[d] & \colim M_{3, i} \ar[d]^{\cong} \ar[r] & 0 \\
0  \ar[r] & M_{1} \ar[r]  &  M_{2}  \ar[r] & M_{3} \ar[r] & 0
}
$$
其两行正合，且外侧竖直映射均为同构。因此
$\colim M_{2, i} \to M_2$ 也是同构，故 (6) 成立。

\medskip\noindent
由例 \ref{algebra-example-universally-exact} (2)，条件 (6) 蕴含 (1)。
\end{proof}

\noindent
上一条定理表明，普遍正合序列总是分裂短正合序列的余极限。
若一个普遍单射映射的余核是有限表示的，则该映射本身事实上分裂：

\begin{lemma}
\label{algebra-lemma-universally-exact-split}
设
$$
0 \to M_1 \to M_2 \to M_3 \to 0
$$
是 $R$-模的正合序列。假设 $M_3$ 是有限表示的。那么
$$
0 \to M_1 \to M_2 \to M_3 \to 0
$$
普遍正合，当且仅当它分裂。
\end{lemma}

\begin{proof}
分裂短正合序列总是普遍正合的，见例
\ref{algebra-example-universally-exact}。
反过来，若该序列普遍正合，则把定理
\ref{algebra-theorem-universally-exact-criteria} (5)
应用于 $P = M_3$，可知映射 $M_2 \to M_3$ 有截面。
\end{proof}

\noindent
下一个引理说明普遍单射映射如何与平坦模相互补充。

\begin{lemma}
\label{algebra-lemma-flat-universally-injective}
设 $M$ 是 $R$-模。那么 $M$ 平坦，当且仅当每个 $R$-模正合序列
$$
0 \to M_1 \to M_2 \to M \to 0
$$
都普遍正合。
\end{lemma}

\begin{proof}
这由引理 \ref{algebra-lemma-flat-surjective-hom} 和定理
\ref{algebra-theorem-universally-exact-criteria} (5) 得出。
\end{proof}

\begin{example}
\label{algebra-example-universally-exact-non-split-non-flat}
非分裂及非平坦的普遍正合序列。
\begin{enumerate}
\item 尽管有引理 \ref{algebra-lemma-universally-exact-split}，仍可能存在
$R$-模的短正合序列
$$
0 \to M_1 \to M_2 \to M_3 \to 0
$$
是普遍正合但非分裂的。例如，取 $R = \mathbf{Z}$，
令 $M_1 = \bigoplus_{n=1}^{\infty} \mathbf{Z}$，
令 $M_{2} = \prod_{n = 1}^{\infty} \mathbf{Z}$，并令
$M_{3}$ 为嵌入 $M_1 \to M_2$ 的余核。由于它们都无挠，
$M_1, M_2, M_3$ 都平坦（《代数进阶》，引理
\ref{more-algebra-lemma-dedekind-torsion-free-flat}），
故由引理 \ref{algebra-lemma-flat-universally-injective}，序列
$$
0 \to M_1 \to M_2 \to M_3 \to 0
$$
普遍正合。然而不存在截面 $s: M_3 \to M_2$。
事实上，若 $x$ 是 $(2, 2^2, 2^3, \ldots) \in M_2$
在 $M_3$ 中的像，则每个模映射 $s: M_3 \to M_2$ 都必须把
$x$ 映为零。这是因为对每个 $n \geq 1$ 都有
$x \in 2^n M_3$，所以对所有 $n \geq 1$，$s(x)$ 都可被
$2^n$ 整除，因而只能是 $0$。
\item 尽管有引理 \ref{algebra-lemma-flat-universally-injective}，仍可能存在
$R$-模的短正合序列
$$
0 \to M_1 \to M_2 \to M_3 \to 0
$$
是普遍正合的，但 $M_1, M_2, M_3$ 全都不平坦。事实上，若 $M$
是任意非平坦模，只须取分裂正合序列
$$
0 \to M \to M \oplus M \to M \to 0.
$$
例如在 $R = \mathbf{Z}$ 上，可取 $M$ 为任意挠模。
\item 将 (1) 中的一个正合序列与 (2) 中的一个正合序列取直和，
便得到 $R$-模的短正合序列
$$
0 \to M_1 \to M_2 \to M_3 \to 0
$$
它普遍正合、非分裂，并且 $M_1, M_2, M_3$ 全都不平坦。
\end{enumerate}
\end{example}

\begin{lemma}
\label{algebra-lemma-ui-flat-domain}
设 $0 \to M_1 \to M_2 \to M_3 \to 0$ 是 $R$-模的普遍正合序列，
并假设 $M_2$ 平坦。那么 $M_1$ 和 $M_3$ 都平坦。
\end{lemma}

\begin{proof}
设 $0 \to N \to N' \to N'' \to 0$ 是 $R$-模的短正合序列。
考虑交换图
$$
\xymatrix{
M_1 \otimes_R N \ar[r] \ar[d] &
M_2 \otimes_R N \ar[r] \ar[d] &
M_3 \otimes_R N \ar[d] \\
M_1 \otimes_R N' \ar[r] \ar[d] &
M_2 \otimes_R N' \ar[r] \ar[d] &
M_3 \otimes_R N' \ar[d] \\
M_1 \otimes_R N'' \ar[r] &
M_2 \otimes_R N'' \ar[r] &
M_3 \otimes_R N''
}
$$
（我们省略了边界上的 $0$）。由假设，各行都是短正合序列，并且箭头
$M_2 \otimes N \to M_2 \otimes N'$ 是单射。这显然蕴含
$M_1 \otimes N \to M_1 \otimes N'$ 是单射，故 $M_1$ 平坦。
特别地，左列和中列都给出短正合序列。作一次追图可知，箭头
$M_3 \otimes N \to M_3 \otimes N'$ 是单射。因此 $M_3$ 平坦。
\end{proof}

\begin{lemma}
\label{algebra-lemma-universally-injective-tensor}
设 $R$ 是环，$M \to M'$ 是普遍单射的 $R$-模映射。
那么对任意 $R$-模 $N$，映射
$M \otimes_R N \to M' \otimes_R N$ 都是普遍单射。
\end{lemma}

\begin{proof}
略。
\end{proof}

\begin{lemma}
\label{algebra-lemma-composition-universally-injective}
设 $R$ 是环。普遍单射的 $R$-模映射的复合仍然普遍单射。
\end{lemma}

\begin{proof}
略。
\end{proof}

\begin{lemma}
\label{algebra-lemma-universally-injective-permanence}
设 $R$ 是环，$M \to M'$ 和 $M' \to M''$ 是 $R$-模映射。
若其复合 $M \to M''$ 普遍单射，则 $M \to M'$ 普遍单射。
\end{lemma}

\begin{proof}
略。
\end{proof}

\begin{lemma}
\label{algebra-lemma-universally-injective-check-stalks}
设 $R \to S$ 是环映射，$M \to M'$ 是 $S$-模映射。
以下条件等价：
\begin{enumerate}
\item $M \to M'$ 作为 $R$-模映射普遍单射；
\item 对 $S$ 的每个素理想 $\mathfrak q$，映射
$M_{\mathfrak q} \to M'_{\mathfrak q}$
作为 $R$-模映射普遍单射；
\item 对 $S$ 的每个极大理想 $\mathfrak m$，映射
$M_{\mathfrak m} \to M'_{\mathfrak m}$
作为 $R$-模映射普遍单射；
\item 对 $S$ 的每个素理想 $\mathfrak q$，映射
$M_{\mathfrak q} \to M'_{\mathfrak q}$
作为 $R_{\mathfrak p}$-模映射普遍单射，其中 $\mathfrak p$ 是
$\mathfrak q$ 在 $R$ 中的逆像；
\item 对 $S$ 的每个极大理想 $\mathfrak m$，映射
$M_{\mathfrak m} \to M'_{\mathfrak m}$
作为 $R_{\mathfrak p}$-模映射普遍单射，其中 $\mathfrak p$ 是
$\mathfrak m$ 在 $R$ 中的逆像。
\end{enumerate}
\end{lemma}

\begin{proof}
设 $N$ 是 $R$-模。设 $\mathfrak q$ 是 $S$ 的素理想，位于
$R$ 的素理想 $\mathfrak p$ 之上。则
$$
(M \otimes_R N)_{\mathfrak q} =
M_{\mathfrak q} \otimes_R N =
M_{\mathfrak q} \otimes_{R_{\mathfrak p}} N_{\mathfrak p}.
$$
此外，对 $M'$ 也有同样的等式，而局部化是正合的。再者，若 $N$
是 $R_{\mathfrak p}$-模，则 $N_{\mathfrak p} = N$。利用这些事实，
可以直接证明上述等价性。

\medskip\noindent
例如，假设 (5) 成立。令
$K = \Ker(M \otimes_R N \to M' \otimes_R N)$。由上述说明，
对 $S$ 的每个极大理想 $\mathfrak m$ 都有
$K_{\mathfrak m} = 0$。因此由引理
\ref{algebra-lemma-characterize-zero-local}，$K = 0$，故 (1) 成立。
反过来，假设 (1) 成立。任取位于
$\mathfrak p \subset R$ 之上的 $\mathfrak q \subset S$，
并任取 $R_{\mathfrak p}$ 上的模 $N$。由假设，
$\Ker(M \otimes_R N \to M' \otimes_R N) = 0$。
于是由上述公式以及 $N = N_{\mathfrak p}$ 可知
$$
\Ker(M_{\mathfrak q} \otimes_{R_{\mathfrak p}} N \to
M'_{\mathfrak q} \otimes_{R_{\mathfrak p}} N) = 0
$$
换言之，(4) 成立。当然，(4) $\Rightarrow$ (5) 是立即的。
因此 (1)、(4) 和 (5) 彼此等价。其余等价性的证明从略。
\end{proof}

\begin{lemma}
\label{algebra-lemma-universally-injective-localize}
设 $\varphi : A \to B$ 是环映射。设 $S \subset A$ 和
$S' \subset B$ 是乘法子集，且 $\varphi(S) \subset S'$。
设 $M \to M'$ 是 $B$-模映射。
\begin{enumerate}
\item 若 $M \to M'$ 作为 $A$-模映射普遍单射，则
$(S')^{-1}M \to (S')^{-1}M'$ 作为 $A$-模映射和作为
$S^{-1}A$-模映射都是普遍单射。
\item 若 $M$ 和 $M'$ 是 $(S')^{-1}B$-模，则 $M \to M'$
作为 $A$-模映射普遍单射，当且仅当它作为 $S^{-1}A$-模映射
普遍单射。
\end{enumerate}
\end{lemma}

\begin{proof}
可以用引理 \ref{algebra-lemma-universally-injective-check-stalks} 证明这一点，
也可以如下直接证明。假设 $M \to M'$ 作为 $A$-模映射普遍单射。
设 $Q$ 是 $A$-模。则 $Q \otimes_A M \to Q \otimes_A M'$ 是单射。
因为局部化正合，所以
$(S')^{-1}(Q \otimes_A M) \to (S')^{-1}(Q \otimes_A M')$
是单射。又因为
$(S')^{-1}(Q \otimes_A M) = Q \otimes_A (S')^{-1}M$，
而对 $M'$ 也同样，所以
$Q \otimes_A (S')^{-1}M \to Q \otimes_A (S')^{-1}M'$ 是单射。
因此 $(S')^{-1}M \to (S')^{-1}M'$ 作为 $A$-模映射普遍单射。
这证明了 (1) 的第一部分。
为证明 (2)，可使用以下两个事实：(a) 若 $Q$ 是
$S^{-1}A$-模，则 $Q \otimes_A S^{-1}A = Q$，即在 $A$ 上与
$Q$ 张量和在 $S^{-1}A$ 上与 $Q$ 张量是一回事；(b) 若 $M$
是任意 $A$-模且 $S$ 中的元素在其上都可逆，则
$M \otimes_A Q = M \otimes_{S^{-1}A} S^{-1}Q$。
(2) 立即由此得出。
\end{proof}

\begin{lemma}
\label{algebra-lemma-check-universally-injective-into-flat}
设 $R$ 是环，$M \to M'$ 是 $R$-模映射。若 $M'$ 平坦，则
$M \to M'$ 普遍单射，当且仅当对 $R$ 的每个有限生成理想 $I$，
$M/IM \to M'/IM'$ 都是单射。
\end{lemma}

\begin{proof}
只须证明对每个有限 $R$-模 $Q$，映射
$M \otimes_R Q \to M' \otimes_R Q$ 都是单射，见定理
\ref{algebra-theorem-universally-exact-criteria}。
此时 $Q$ 有一个由子模组成的有限滤过
$0 = Q_0 \subset Q_1 \subset \ldots \subset Q_n = Q$，
其子商同构于循环模 $R/I_i$，见引理
\ref{algebra-lemma-trivial-filter-finite-module}。由于 $M'$ 平坦，
得到 $M' \otimes_R Q$ 的子模 $M' \otimes_R Q_i$ 所成的滤过
$$
\xymatrix{
M \otimes Q_1 \ar[r] \ar[d] &
M \otimes Q_2 \ar[r] \ar[d] &
\ldots \ar[r] &
M \otimes Q \ar[d] \\
M' \otimes Q_1 \ar@{^{(}->}[r] &
M' \otimes Q_2 \ar@{^{(}->}[r] &
\ldots \ar@{^{(}->}[r] &
M' \otimes Q
}
$$
其相继商为 $M' \otimes_R R/I_i = M'/I_iM'$。
简单的归纳表明，只须检验 $M/I_i M \to M'/I_i M'$ 是单射。
注意，有限生成理想 $I'_i \subset I_i$ 的全体是一个有向集。
因此 $M/I_iM = \colim M/I'_iM$ 是滤过余极限，$M'$ 的情形类似；
各映射 $M/I'_iM \to M'/I'_i M'$ 由假设都是单射，而滤过余极限
正合（引理 \ref{algebra-lemma-directed-colimit-exact}），故结论成立。
\end{proof}

\begin{lemma}
\label{algebra-lemma-base-change-along-universally-injective-ring-map}
设 $R \to S$ 是环映射，并且作为 $R$-模映射普遍单射。
那么 $R$-模上的函子 $M \mapsto M \otimes_R S$
反映单射、满射和同构。
\end{lemma}

\begin{proof}
设 $M \to N$ 是 $R$-模映射，其核为 $K$、余核为 $Q$。
若 $M \otimes_R S \to N \otimes_R S$ 是单射，则
$K \otimes_R S \to M \otimes_R S$ 的像为零。由于
$K \subset K \otimes_R S$ 且 $M \subset M \otimes_R S$，
可知 $K = 0$。另一方面，若
$M \otimes_R S \to N \otimes_R S$ 是满射，则由张量积的右正合性，
$Q \otimes_R S$ 为零。因此 $Q \subset Q \otimes_R S$ 也为零。
\end{proof}

\begin{lemma}
\label{algebra-lemma-faithfully-flat-universally-injective}
设 $R \to S$ 是忠实平坦环映射。那么作为 $R$-模映射，
$R \to S$ 普遍单射。特别地，对任意理想 $I \subset R$ 都有
$R \cap IS = I$。
\end{lemma}

\begin{proof}
设 $N$ 是 $R$-模。我们要证明 $N \to N \otimes_R S$ 是单射。
由于 $S$ 作为 $R$-模忠实平坦，只须在与 $S$ 张量后证明这一点。
因此只须证明映射
$N \otimes_R S \to N \otimes_R S \otimes_R S$，
$n \otimes s \mapsto n \otimes 1 \otimes s$ 是单射。
这是成立的，因为它有收缩映射
$n \otimes s \otimes s' \mapsto n \otimes ss'$。
\end{proof}

\section{有限投射模的下降}
\label{algebra-section-finite-projective}

\noindent
本节给出一个初等证明：作为{\it 有限}投射模的性质沿忠实平坦环映射
下降。去掉有限性条件后，这个证明便不再适用。不过，该方法预示了我们
将用来证明{\it 可数生成}投射模之下降的方法；见本节末尾的说明。

\begin{lemma}
\label{algebra-lemma-finite-projective-again}
设 $M$ 是 $R$-模。那么 $M$ 有限投射，当且仅当 $M$ 有限表示且平坦。
\end{lemma}

\begin{proof}
这是引理 \ref{algebra-lemma-finite-projective} 的一部分。不过在此处，我们可以
如下对该引理中的蕴含 (1) $\Rightarrow$ (2) 给出一个更优雅的证明。
若 $M$ 有限表示且平坦，取一个满射 $R^n \to M$。
把引理 \ref{algebra-lemma-flat-surjective-hom} 应用于 $P = M$，
可知映射 $R^n \to M$ 有截面。因此 $M$ 是自由模的直和因子，
从而是投射模。
\end{proof}

\noindent
以下是模的一些可下降性质。

\begin{lemma}
\label{algebra-lemma-descend-properties-modules}
设 $R \to S$ 是忠实平坦环映射，$M$ 是 $R$-模。那么
\begin{enumerate}
\item 若 $S$-模 $M \otimes_R S$ 是有限型的，则 $M$ 是有限型的；
\item 若 $S$-模 $M \otimes_R S$ 是有限表示的，则 $M$ 是有限表示的；
\item 若 $S$-模 $M \otimes_R S$ 平坦，则 $M$ 平坦；
\item 按需在此添加更多性质。
\end{enumerate}
\end{lemma}

\begin{proof}
假设 $M \otimes_R S$ 是有限型的。设 $y_1, \ldots, y_m$ 是
$M \otimes_R S$ 在 $S$ 上的一组生成元。对某些
$x_1, \ldots, x_n \in M$，写
$y_j = \sum x_i \otimes f_i$。于是映射
$\varphi : R^{\oplus n} \to M$ 具有如下性质：
$\varphi \otimes \text{id}_S : S^{\oplus n} \to M \otimes_R S$
是满射。由于 $R \to S$ 忠实平坦，可知 $\varphi$ 是满射，
故 $M$ 有限生成。

\medskip\noindent
假设 $M \otimes_R S$ 是有限表示的。由 (1)，$M$ 是有限型的。
选取满射 $R^{\oplus n} \to M$，并将其核记为 $K$。
由于 $R \to S$ 平坦，$K \otimes_R S$ 是基变换
$S^{\oplus n} \to M \otimes_R S$ 的核。
由于 $M \otimes_R S$ 是有限表示的，可知 $K \otimes_R S$
是有限型的。因此由 (1)，$K$ 是有限型的，从而 $M$ 是有限表示的。

\medskip\noindent
(3) 就是引理 \ref{algebra-lemma-flatness-descends}。
\end{proof}

\begin{proposition}
\label{algebra-proposition-ffdescent-finite-projectivity}
设 $R \to S$ 是忠实平坦环映射，$M$ 是 $R$-模。
若 $S$-模 $M \otimes_R S$ 有限投射，则 $M$ 有限投射。
\end{proposition}

\begin{proof}
这由引理 \ref{algebra-lemma-finite-projective-again} 和
\ref{algebra-lemma-descend-properties-modules} 得出。
\end{proof}

\noindent
下面数节将利用 d\'evissage（逐层剖析）把问题约化到可数生成情形，
以去掉有限性假设。在可数生成情形中，我们的策略是寻找一个类似于
引理 \ref{algebra-lemma-finite-projective-again} 的可数生成投射模刻画，
然后直接证明该刻画可下降。为此，我们引入 Mittag-Leffler 模的概念，
并证明：若模 $M$ 可数生成，则它是投射模，当且仅当它既平坦又是
Mittag-Leffler 模（定理 \ref{algebra-theorem-projectivity-characterization}）。
当 $M$ 有限生成时，该陈述约化为引理
\ref{algebra-lemma-finite-projective-again}，因为根据例
\ref{algebra-example-ML} (1)，有限生成模是 Mittag-Leffler 模，
当且仅当它有限表示。

\section{模的超限 d\'evissage}
\label{algebra-section-transfinite-devissage}

\noindent
本节引入一种把模分解为直和的 d\'evissage（逐层剖析）技术。
主要结果是投射模可以写成可数生成模的直和（见下文定理
\ref{algebra-theorem-projective-direct-sum}）。我们遵循 \cite{Kaplansky}。


\begin{definition}
\label{algebra-definition-devissage}
设 $M$ 是 $R$-模。$M$ 的一个{\it 直和 d\'evissage} 是一族子模
$(M_{\alpha})_{\alpha \in S}$，它由序数 $S$ 标号，并且按包含关系
递增，且满足：
\begin{enumerate}
\item[(0)] $M_0 = 0$；
\item[(1)] $M = \bigcup_{\alpha} M_{\alpha}$；
\item[(2)] 若 $\alpha \in S$ 是极限序数，则
$M_{\alpha} = \bigcup_{\beta < \alpha} M_{\beta}$；
\item[(3)] 若 $\alpha + 1 \in S$，则 $M_{\alpha}$ 是
$M_{\alpha + 1}$ 的直和因子。
\end{enumerate}
若还满足
\begin{enumerate}
\item[(4)] 对 $\alpha + 1 \in S$，$M_{\alpha + 1}/M_{\alpha}$
可数生成，
\end{enumerate}
则称 $(M_{\alpha})_{\alpha \in S}$ 为 $M$ 的一个
{\it Kaplansky d\'evissage}。
\end{definition}

\noindent
下一个引理说明了上述术语的理由。

\begin{lemma}
\label{algebra-lemma-direct-sum-devissage}
设 $M$ 是 $R$-模。若 $(M_{\alpha})_{\alpha \in S}$ 是 $M$ 的
一个直和 d\'evissage，则
$M \cong \bigoplus_{\alpha + 1 \in S} M_{\alpha + 1}/M_{\alpha}$。
\end{lemma}

\begin{proof}
由直和 d\'evissage 的性质 (3)，对每个 $\alpha \in S$，
都有嵌入 $M_{\alpha + 1}/M_{\alpha} \to M$。考虑由这些嵌入之和给出的映射
$$
f : \bigoplus\nolimits_{\alpha + 1\in S} M_{\alpha + 1}/M_{\alpha} \to M
$$
再对 $\beta\in S$ 考虑限制
$$
f_{\beta} :
\bigoplus\nolimits_{\alpha + 1 \leq \beta} M_{\alpha + 1}/M_{\alpha}
\longrightarrow
M
$$
对 $S$ 作超限归纳可知，$f_{\beta}$ 的像为 $M_{\beta}$。
当 $\beta=0$ 时，这由 $(0)$ 成立。若 $\beta+1$ 是后继序数，
且结论对 $\beta$ 成立，则由 (3)，结论对 $\beta + 1$ 成立。
若 $\beta$ 是极限序数，且结论对 $\alpha < \beta$ 成立，
则由 (2)，结论对 $\beta$ 成立。因此由 (1)，$f$ 是满射。

\medskip\noindent
对 $S$ 作超限归纳还可证明各限制 $f_{\beta}$ 都是单射。
当 $\beta = 0$ 时结论成立。若 $\beta+1$ 是后继序数，
且 $f_{\beta}$ 是单射，任取核中的 $x$，并按其分量写成
$x = (x_{\alpha + 1})_{\alpha + 1 \leq \beta + 1}$，其中
$x_{\alpha + 1} \in M_{\alpha + 1}/M_{\alpha}$。
由性质 (3) 以及 $f_{\beta}$ 的像为 $M_{\beta}$ 这一事实，
$(x_{\alpha + 1})_{\alpha + 1 \leq \beta}$ 和
$x_{\beta + 1}$ 都映为 $0$。因此 $x_{\beta+1} = 0$；
再由限制 $f_{\beta}$ 单射的假设，对每个
$\alpha + 1 \leq \beta$ 都有 $x_{\alpha + 1} = 0$。
故 $x = 0$，而 $f_{\beta+1}$ 是单射。
若 $\beta$ 是极限序数，考虑核中的元素 $x$。则 $x$ 已包含在某个
$\alpha < \beta$ 的 $f_{\alpha}$ 的定义域中，故 $x = 0$，
归纳完成。由于对每个 $\beta \in S$，$f_{\beta}$ 都单射，
可知 $f$ 单射。
\end{proof}

\begin{lemma}
\label{algebra-lemma-Kaplansky-devissage}
设 $M$ 是 $R$-模。那么 $M$ 是可数生成 $R$-模的直和，
当且仅当它有 Kaplansky d\'evissage。
\end{lemma}

\begin{proof}
上一引理给出“若”的方向。反过来，假设
$M = \bigoplus_{i \in I} N_i$，其中每个 $N_i$ 都是可数生成
$R$-模。给 $I$ 配备良序，从而可把它视为序数。令
$M_i = \bigoplus_{j < i} N_j$，便得到 $M$ 的一个
Kaplansky d\'evissage $(M_i)_{i \in I}$。
\end{proof}

\begin{theorem}
\label{algebra-theorem-kaplansky-direct-sum}
假设 $M$ 是可数生成 $R$-模的直和。若 $P$ 是 $M$ 的直和因子，
则 $P$ 也是可数生成 $R$-模的直和。
\end{theorem}

\begin{proof}
写成 $M = P \oplus Q$。我们将构造 $M$ 的一个 Kaplansky
d\'evissage $(M_{\alpha})_{\alpha \in S}$；除定义性质 (0)--(4)
外，它还满足：
\begin{enumerate}
\item[(5)] 每个 $M_{\alpha}$ 都是 $M$ 的直和因子；
\item[(6)] $M_{\alpha} = P_{\alpha} \oplus Q_{\alpha}$，其中
$P_{\alpha} =P \cap M_{\alpha}$ 且
$Q_\alpha = Q \cap M_{\alpha}$。
\end{enumerate}
（注意：若性质 (0)--(2) 成立，则性质 (3) 事实上等价于性质 (5)。）

\medskip\noindent
为说明这如何推出定理，只须证明 $(P_{\alpha})_{\alpha \in S}$
构成 $P$ 的一个 Kaplansky d\'evissage。性质 (0)、(1) 和 (2)
是显然的。由 $(M_{\alpha})$ 的 (5) 和 (6)，每个
$P_{\alpha}$ 都是 $M$ 的直和因子。由于
$P_{\alpha} \subset P_{\alpha + 1}$，这蕴含 $P_{\alpha}$ 是
$P_{\alpha + 1}$ 的直和因子；故 $(P_{\alpha})$ 满足 (3)。
关于 (4)，注意
$$
M_{\alpha + 1}/M_{\alpha} \cong P_{\alpha + 1}/P_{\alpha} \oplus
Q_{\alpha + 1}/Q_{\alpha},
$$
因此 $P_{\alpha + 1}/P_{\alpha}$ 可数生成，因为
$M_{\alpha + 1}/M_{\alpha}$ 可数生成。

\medskip\noindent
还需构造诸 $M_{\alpha}$。写成
$M = \bigoplus_{i \in I} N_i$，其中每个 $N_i$ 都是可数生成
$R$-模。选取 $I$ 的一个良序。我们将通过超限递归定义 $M$
的一族递增子模 $M_{\alpha}$，每个序数 $\alpha$ 对应一个，
并使 $M_{\alpha}$ 是某些 $N_i$ 的直和。

\medskip\noindent
当 $\alpha = 0$ 时，令 $M_{0} = 0$。若 $\alpha$ 是极限序数，
且对所有 $\beta < \alpha$ 已定义 $M_{\beta}$，则定义
$M_{\alpha} = \bigcup_{\beta < \alpha} M_{\beta}$。
由于每个 $\beta < \alpha$ 的 $M_{\beta}$ 都是某些 $N_i$ 的直和，
$M_{\alpha}$ 也有同样的性质。若 $\alpha + 1$ 是后继序数，
且已定义 $M_{\alpha}$，则如下定义 $M_{\alpha + 1}$。
若 $M_{\alpha} = M$，令 $M_{\alpha + 1} = M$。
否则，选取使 $N_j$ 不包含于 $M_{\alpha}$ 的最小 $j \in I$。
我们将构造无穷矩阵 $(x_{mn}), m, n = 1, 2, 3, \ldots$，使得：
\begin{enumerate}
\item $N_j$ 包含于由矩阵条目 $x_{mn}$ 生成的 $M$ 的子模；
\item 若把任意条目 $x_{k\ell}$ 按其 $P$-分量和 $Q$-分量写成
$x_{k\ell} = y_{k\ell} + z_{k\ell}$，则对每个 $N_i$，只要
$y_{k\ell}$ 或 $z_{k\ell}$ 在其中的分量非零，矩阵
$(x_{mn})$ 都含有它的一组生成元。
\end{enumerate}
然后定义 $M_{\alpha + 1}$ 为由 $M_{\alpha}$ 和所有 $x_{mn}$
生成的 $M$ 的子模；由矩阵 $(x_{mn})$ 的性质 (2)，
$M_{\alpha + 1}$ 将是某些 $N_i$ 的直和。
为构造矩阵 $(x_{mn})$，令 $x_{11}, x_{12}, x_{13}, \ldots$
为 $N_j$ 的一组可数生成元。然后，若
$x_{11} = y_{11} + z_{11}$ 是按 $P$-分量和 $Q$-分量的分解，
令 $x_{21}, x_{22}, x_{23}, \ldots$ 为如下诸 $N_i$ 之和的一组
可数生成元：$y_{11}$ 或 $z_{11}$ 在其中的分量非零。
对 $x_{12}$ 重复此过程，得到元素
$x_{31}, x_{32}, \ldots$，即矩阵的第三行。
再对 $x_{21}$ 重复，得到第四行；对 $x_{13}$ 重复，得到第五行；
依此类推，按下图所示沿相继的反对角线向下进行：
$$
\left(
\vcenter{
\xymatrix@R=2mm@C=2mm{
x_{11} & x_{12} \ar[dl] & x_{13} \ar[dl] & x_{14} \ar[dl] & \ldots  \\
x_{21} & x_{22} \ar[dl] & x_{23} \ar[dl] & \ldots  \\
x_{31} & x_{32} \ar[dl] & \ldots \\
x_{41} & \ldots \\
\ldots
}
}
\right).
$$

\medskip\noindent
对 $I$ 作超限归纳（使用如下事实：我们构造 $M_{\alpha + 1}$ 时，
令其包含使 $N_j$ 不包含于 $M_{\alpha}$ 的最小 $j$ 所对应的
$N_j$），可知对每个 $i \in I$，$N_i$ 都包含于某个
$M_{\alpha}$。因此存在充分大的序数 $S$，满足：对每个
$i \in I$，都存在 $\alpha \in S$，使 $N_i$ 包含于
$M_{\alpha}$。这意味着 $(M_{\alpha})_{\alpha \in S}$ 满足
$M$ 的 Kaplansky d\'evissage 的性质 (1)。
此外，族 $(M_{\alpha})_{\alpha \in S}$ 还满足其余定义性质以及
上述 (5) 和 (6)：性质 (0)、(2)、(4) 和 (6) 由构造显然成立；
性质 (5) 成立，因为按构造，每个 $M_{\alpha}$ 都是某些 $N_i$
的直和；而由 (5) 以及 $M_{\alpha} \subset M_{\alpha + 1}$，
可推出 (3)。
\end{proof}

\noindent
作为推论，得到本节开头所述的投射模结果。

\begin{theorem}
\label{algebra-theorem-projective-direct-sum}
\begin{slogan}
任意投射模都是可数生成投射模的直和。
\end{slogan}
若 $P$ 是投射 $R$-模，则 $P$ 是可数生成投射 $R$-模的直和。
\end{theorem}

\begin{proof}
一个模是投射模，当且仅当它是自由模的直和因子，
故结论由定理 \ref{algebra-theorem-kaplansky-direct-sum} 得出。
\end{proof}

\section{局部环上的投射模}
\label{algebra-section-projective-local-ring}

\noindent
本节证明一个非常漂亮的结果：局部环上的投射模 $M$ 是自由模
（见下文定理 \ref{algebra-theorem-projective-free-over-local-ring}）。
注意，在附加 $M$ 有限这一假设后，该结果就是引理
\ref{algebra-lemma-finite-flat-local}。一般地，有：

\begin{lemma}
\label{algebra-lemma-projective-free}
设 $R$ 是环。那么每个投射 $R$-模都是自由模，当且仅当每个
可数生成投射 $R$-模都是自由模。
\end{lemma}

\begin{proof}
立即由定理 \ref{algebra-theorem-projective-direct-sum} 得出。
\end{proof}

\noindent
以下是可数生成模为自由模的一个判据。

\begin{lemma}
\label{algebra-lemma-freeness-criteria}
设 $M$ 是可数生成 $R$-模，并具有如下性质：若
$M = N \oplus N'$，其中 $N'$ 是有限自由 $R$-模，则 $N$
的任意元素都包含在 $N$ 的某个自由直和因子中。那么 $M$ 是自由模。
\end{lemma}

\begin{proof}
设 $x_1, x_2, \ldots$ 是 $M$ 的一组可数生成元。
我们归纳构造 $M$ 的有限自由直和因子 $F_1, F_2, \ldots$，
使得对所有 $n$，$F_1 \oplus \ldots \oplus F_n$ 都是 $M$
的直和因子，并包含 $x_1, \ldots, x_n$。具体地，给定具有所需性质的
$F_1, \ldots, F_n$，写成
$$
M = F_1 \oplus \ldots \oplus F_n \oplus N
$$
并令 $x \in N$ 为 $x_{n + 1}$ 的像。由引理陈述中的假设，
可找到包含 $x$ 的自由直和因子 $F_{n + 1} \subset N$。
当然，可以把 $F_{n + 1}$ 换成 $F_{n + 1}$ 的一个有限自由直和因子，
归纳步骤遂告完成。于是
$M = \bigoplus_{i = 1}^{\infty} F_i$ 是自由模。
\end{proof}

\begin{lemma}
\label{algebra-lemma-projective-freeness-criteria}
设 $P$ 是局部环 $R$ 上的投射模。那么 $P$ 的任意元素都包含在
$P$ 的某个自由直和因子中。
\end{lemma}

\begin{proof}
因为 $P$ 投射，所以它是某个自由 $R$-模 $F$ 的直和因子；写成
$F = P \oplus Q$。设 $x \in P$ 是我们要证明包含于 $P$ 的某个
自由直和因子中的元素。选取 $F$ 的一组基 $B$，使表示 $x$ 所需的
基元素个数最少；设
$x = \sum_{i=1}^n a_i e_i$，其中 $e_i \in B$ 且 $a_i \in R$。
那么任何 $a_j$ 都不能写成其余 $a_i$ 的线性组合；否则，若对某些
$b_i \in R$ 有 $a_j = \sum_{i \neq  j} a_i b_i$，则对
$i \neq j$ 把 $e_i$ 换成 $e_i + b_ie_j$，并保持 $B$ 的其余
元素不变，便得到 $F$ 的一组新基，而 $x$ 关于这组基的表示更短。

\medskip\noindent
把 $e_i$ 按其 $P$-分量和 $Q$-分量分解为
$e_i = y_i + z_i, y_i \in P, z_i \in Q$。写成
$y_i = \sum_{j=1}^{n} b_{ij} e_j + t_i$，其中 $t_i$ 是
$B$ 中除 $e_1, \ldots, e_n$ 外的元素的线性组合。
为完成证明，只须说明矩阵 $(b_{ij})$ 可逆。因为在这种情况下，
对 $i=1, \ldots, n$ 按 $e_i \mapsto y_i$ 定义，并固定
$B \setminus \{e_1, \ldots, e_n\}$ 的映射 $F \to F$ 是同构，
所以 $y_1, \ldots, y_n$ 与
$B \setminus \{e_1, \ldots, e_n\}$ 一起构成 $F$ 的一组基。
于是由 $y_1, \ldots, y_n$ 张成的子模 $N$ 是 $P$ 的自由子模；
因为 $N \subset P$，并且 $N$ 和 $P$ 都是 $F$ 的直和因子，
$N$ 是 $P$ 的直和因子；又因为 $x \in P$ 蕴含
$x = \sum_{i=1}^n a_i e_i = \sum_{i=1}^n a_i y_i$，
所以 $x \in N$。

\medskip\noindent
现在证明 $(b_{ij})$ 可逆。把
$y_i = \sum_{j=1}^{n} b_{ij} e_j + t_i$ 代入
$\sum_{i=1}^n a_i e_i = \sum_{i=1}^n a_i y_i$，比较 $e_j$
的系数，得到 $a_j = \sum_{i=1}^n a_i b_{ij}$。
但如上所述，基 $B$ 的选取保证任何 $a_j$ 都不能写成其余
$a_i$ 的线性组合。因此，当 $i \neq j$ 时，$b_{ij}$ 是非单位；
而对所有 $i$，$1-b_{ii}$ 是非单位，因而特别地 $b_{ii}$ 是单位。
局部环上的矩阵若对角线元素均为单位而其余元素均为非单位，
则其行列式是单位，故矩阵可逆。
\end{proof}

\begin{theorem}
\label{algebra-theorem-projective-free-over-local-ring}
\begin{slogan}
局部环上的投射模是自由模。
\end{slogan}
若 $P$ 是局部环 $R$ 上的投射模，则 $P$ 是自由模。
\end{theorem}

\begin{proof}
结论由引理 \ref{algebra-lemma-projective-free}、
\ref{algebra-lemma-freeness-criteria} 和
\ref{algebra-lemma-projective-freeness-criteria} 得出。
\end{proof}

\section{Mittag-Leffler 系统}
\label{algebra-section-mittag-leffler}

\noindent
本节旨在定义 Mittag-Leffler 系统，并说明这一概念为何有用。

\medskip\noindent
下文中，$I$ 表示一个有向集，见《范畴》，定义
\ref{categories-definition-directed-set}。设
$(A_i, \varphi_{ji}: A_j \to A_i)$ 是以 $I$ 为指标的集合或模的逆系，
见《范畴》，定义 \ref{categories-definition-directed-system}。
由于假设 $I$ 有向，这是一个有向逆系（《范畴》，定义
\ref{categories-definition-directed-system}）。对每个 $i \in I$，当
$j \geq i$ 时，诸像 $\varphi_{ji}(A_j) \subset A_i$ 构成 $A_i$
的子集（或子模）的递减有向族。令
$A'_i = \bigcap_{j \geq i} \varphi_{ji}(A_j)$。
当 $j \geq i$ 时，$\varphi_{ji}(A'_j) \subset A'_i$，故限制映射给出
有向逆系 $(A'_i, \varphi_{ji}|_{A'_j})$。由集合范畴或模范畴中
逆系极限的构造，有 $\lim A_i = \lim A'_i$。逆系
$(A_i, \varphi_{ji})$ 的 Mittag-Leffler 条件是：$A'_i$ 等于某个
$j \geq i$ 所对应的 $\varphi_{ji}(A_j)$（因而对所有 $k \geq j$，
也等于 $\varphi_{ki}(A_k)$）：

\begin{definition}
\label{algebra-definition-ML-system}
设 $(A_i, \varphi_{ji})$ 是 $I$ 上集合的有向逆系。若对每个
$i \in I$，当 $j \geq i$ 时，族 $\varphi_{ji}(A_j) \subset A_i$
稳定，则称 $(A_i, \varphi_{ji})$ 是 {\it Mittag-Leffler} 的。
明确地说，这意味着对每个 $i \in I$，存在 $j \geq i$，使得对
$k \geq j$ 有 $\varphi_{ki}(A_k) = \varphi_{ji}( A_j)$。
若 $(A_i, \varphi_{ji})$ 是环 $R$ 上模的有向逆系，则当其底层集合逆系
是 Mittag-Leffler 的时候，称它是 Mittag-Leffler 的。
\end{definition}

\begin{example}
\label{algebra-example-ML-surjective-maps}
若 $(A_i, \varphi_{ji})$ 是集合或模的有向逆系，且映射
$\varphi_{ji}$ 都是满射，则该逆系显然是 Mittag-Leffler 的。
反之，设 $(A_i, \varphi_{ji})$ 是 Mittag-Leffler 的。令
$A'_i \subset A_i$ 为当 $j \geq i$ 时 $\varphi_{ji}(A_j)$ 的稳定像。
则当 $j \geq i$ 时，$\varphi_{ji}|_{A'_j}: A'_j \to A'_i$ 是满射，
且 $\lim A_i = \lim A'_i$。因此，Mittag-Leffler 逆系
$(A_i, \varphi_{ji})$ 的极限也可写成 $I$ 上一个转移映射均为满射的
有向逆系的极限。
\end{example}

\begin{lemma}
\label{algebra-lemma-ML-limit-nonempty}
设 $(A_i, \varphi_{ji})$ 是 $I$ 上的有向逆系。假设 $I$ 可数。
若 $(A_i, \varphi_{ji})$ 是 Mittag-Leffler 的，且各 $A_i$ 非空，
则 $\lim A_i$ 非空。
\end{lemma}

\begin{proof}
设 $i_1, i_2, i_3, \ldots$ 是 $I$ 的元素的一个枚举。按如下条件归纳定义
元素序列 $j_n \in I$，其中 $n = 1, 2, 3, \ldots$：$j_1 = i_1$，
并且 $j_n \geq i_n$，且对 $m < n$ 有 $j_n \geq j_m$。
于是序列 $j_n$ 递增，并构成 $I$ 的共尾子集。因此可假设
$I =\{1, 2, 3, \ldots \}$。由例
\ref{algebra-example-ML-surjective-maps}，问题归结为证明：由正整数指标、
转移映射均为满射的非空集合逆系，其极限非空。这由选择公理得出。
\end{proof}

\noindent
Mittag-Leffler 条件对我们很重要，因为它具有如下正合性性质。

\begin{lemma}
\label{algebra-lemma-ML-exact-sequence}
设
$$
0 \to A_i \xrightarrow{f_i} B_i \xrightarrow{g_i} C_i \to 0
$$
是 $I$ 上阿贝尔群的有向逆系的一个正合列。假设 $I$ 可数。
若 $(A_i)$ 是 Mittag-Leffler 的，则
$$
0 \to \lim A_i \to \lim B_i \to \lim C_i\to 0
$$
正合。
\end{lemma}

\begin{proof}
对有向逆系取极限是左正合的，故只须证明
$\lim B_i \to \lim C_i$ 满射。取 $(c_i) \in \lim C_i$。
对每个 $i \in I$，令 $E_i = g_i^{-1}(c_i)$；由于
$g_i: B_i \to C_i$ 满射，所得集合非空。逆系 $(B_i)$ 的映射系
$\varphi_{ji}: B_j \to B_i$ 限制为映射 $E_j \to E_i$，从而使
$(E_i)$ 成为非空集合的逆系。只须证明 $(E_i)$ 是 Mittag-Leffler 的。
因为这样一来，引理 \ref{algebra-lemma-ML-limit-nonempty} 就说明
$\lim E_i$ 非空，而取 $\lim E_i$ 的任意元素，便得到 $\lim B_i$
中映到 $(c_i)$ 的一个元素。

\medskip\noindent
借助单射 $f_i: A_i \to B_i$，我们把 $A_i$ 视为 $B_i$ 的子集。
由于 $(A_i)$ 是 Mittag-Leffler 的，若 $i \in I$，则存在
$j \geq i$，使得对 $k \geq j$ 有
$\varphi_{ki}(A_k) = \varphi_{ji}(A_j)$。我们断言，对 $k \geq j$
也有 $\varphi_{ki}(E_k) = \varphi_{ji}(E_j)$。对 $k \geq j$，总有
$\varphi_{ki}(E_k) \subset \varphi_{ji}(E_j)$。反过来，取
$e_j \in E_j$；需要找到 $x_k \in E_k$，使
$\varphi_{ki}(x_k) = \varphi_{ji}(e_j)$。任取 $e'_k \in E_k$，
并令 $e'_j = \varphi_{kj}(e'_k)$。于是
$g_j(e_j - e'_j) = c_j - c_j = 0$，故
$e_j - e'_j = a_j \in A_j$。由于
$\varphi_{ki}(A_k) = \varphi_{ji}(A_j)$，存在 $a_k \in A_k$，使
$\varphi_{ki}(a_k) = \varphi_{ji}(a_j)$。因此
$$
\varphi_{ki}(e'_k + a_k) = \varphi_{ji}(e'_j) + \varphi_{ji}(a_j) =
\varphi_{ji}(e_j),
$$
故可取 $x_k = e'_k + a_k$。
\end{proof}

\section{逆系}
\label{algebra-section-inverse-systems}

\noindent
在许多论文中（以及在本节中），术语{\it 逆系}专指偏序集
$(\mathbf{N}, \geq)$ 上的逆系。本节简要讨论这类逆系。
《同调代数》，第 \ref{homology-section-inverse-systems} 节将更广泛地讨论
这些材料。设给定一个环 $R$ 和一列 $R$-模
$$
M_1 \xleftarrow{\varphi_2} M_2 \xleftarrow{\varphi_3} M_3 \leftarrow \ldots
$$
以及图示的映射。复合相继映射后，只要 $i \geq i'$，就得到映射
$\varphi_{ii'} : M_i \to M_{i'}$；而且只要
$i \geq i' \geq i''$，就有
$\varphi_{ii''} = \varphi_{i'i''} \circ \varphi_{i i'}$。
反之，给定映射族 $\varphi_{ii'}$，令
$\varphi_i = \varphi_{i(i-1)}$，便恢复上面所示的映射。在这种情况下，
$$
\lim M_i
=
\{(x_i) \in \prod M_i \mid \varphi_i(x_i) = x_{i - 1}, \ i = 2, 3, \ldots\}
$$
参见《范畴》，第 \ref{categories-section-limit-sets} 节。
如《同调代数》，第 \ref{homology-section-inverse-systems} 节所说明的，
这实际上是 $R$-模范畴中的极限，其定义见《范畴》，第
\ref{categories-section-limits} 节。

\begin{lemma}
\label{algebra-lemma-Mittag-Leffler}
设 $R$ 是环。设 $0 \to K_i \to L_i \to M_i \to 0$ 是 $R$-模的短正合列，
$i \geq 1$，并且它们嵌入如下短正合列之间的映射：
$$
\xymatrix{
0 \ar[r] &
K_i \ar[r] &
L_i \ar[r] &
M_i \ar[r] &
0 \\
0 \ar[r] &
K_{i + 1} \ar[r] \ar[u] &
L_{i + 1} \ar[r] \ar[u] &
M_{i + 1} \ar[r] \ar[u] &
0}
$$
若对每个 $i$，都存在 $c = c(i) \geq i$，使得对所有 $j \geq c$ 有
$\Im(K_c \to K_i) = \Im(K_j \to K_i)$，则序列
$$
0 \to \lim K_i \to \lim L_i \to \lim M_i \to 0
$$
正合。
\end{lemma}

\begin{proof}
这是更一般的引理 \ref{algebra-lemma-ML-exact-sequence} 的特殊情形。
\end{proof}

\section{Mittag-Leffler 模}
\label{algebra-section-mittag-leffler-modules}

\noindent
粗略地说，Mittag-Leffler 模是一个可写成有向余极限、且其对偶是
Mittag-Leffler 系统的模。为了给出精确定义，还需要做一些准备。

\begin{definition}
\label{algebra-definition-ML-inductive-system}
设 $(M_i, f_{ij})$ 是 $R$-模的有向系。若每个 $M_i$ 都是有限表示
$R$-模，并且对每个 $R$-模 $N$，逆系
$$
(\Hom_R(M_i, N), \Hom_R(f_{ij}, N))
$$
都是 Mittag-Leffler 的，则称 $(M_i, f_{ij})$ 是一个
{\it 模的 Mittag-Leffler 有向系}。
\end{definition}

\noindent
下面将刻画哪些 $R$-模是模的 Mittag-Leffler 有向系的余极限。

\begin{definition}
\label{algebra-definition-domination}
设 $f: M \to N$ 和 $g: M \to M'$ 是 $R$-模映射。若对任意
$R$-模 $Q$ 都有 $\Ker(f
\otimes_R \text{id}_Q) \subset \Ker(g \otimes_R \text{id}_Q)$，
则称 $g$ {\it 支配} $f$。
\end{definition}

\noindent
只需对有限表示模检验这一条件。

\begin{lemma}
\label{algebra-lemma-domination-fp}
设 $f: M \to N$ 和 $g: M \to M'$ 是 $R$-模映射。那么 $g$ 支配 $f$，
当且仅当对任意有限表示 $R$-模 $Q$，都有
$\Ker(f \otimes_R \text{id}_Q) \subset \Ker(g \otimes_R
\text{id}_Q)$。
\end{lemma}

\begin{proof}
假设对所有有限表示模 $Q$，都有
$\Ker(f \otimes_R \text{id}_Q) \subset \Ker(g \otimes_R
\text{id}_Q)$。若 $Q$ 是任意模，把
$Q = \colim_{i \in I} Q_i$ 写成有限表示模 $Q_i$ 的一个有向系的余极限。
于是对所有 $i$，都有 $\Ker(f \otimes_R
\text{id}_{Q_i}) \subset \Ker(g \otimes_R \text{id}_{Q_i})$。
由于取有向余极限是正合的且与张量积交换，遂有
$\Ker(f \otimes_R \text{id}_Q) \subset \Ker(g
\otimes_R \text{id}_Q)$。
\end{proof}

\begin{lemma}
\label{algebra-lemma-domination-universally-injective}
设 $f : M \to N$ 和 $g : M \to M'$ 是 $R$-模映射。考虑 $f$ 与 $g$
的推出
$$
\xymatrix{
M  \ar[r]_f \ar[d]_g & N \ar[d]^{g'} \\
M' \ar[r]^{f'} & N'
}
$$
那么 $g$ 支配 $f$，当且仅当 $f'$ 普遍单射。
\end{lemma}

\begin{proof}
回顾 $N'$ 是 $M' \oplus N$ 对由所有 $x \in M$ 所对应的元素
$(g(x), -f(x))$ 组成的子模的商。由 $N'$ 的构造，有短正合列
$$
0 \to \Ker(f) \cap \Ker(g) \to \Ker(f) \to \Ker(f')
\to 0.
$$
因为张量运算与取推出交换，所以对每个 $R$-模 $Q$，都有这样的短正合列
$$
0 \to \Ker(f \otimes \text{id}_Q ) \cap \Ker(g \otimes
\text{id}_Q) \to \Ker(f \otimes \text{id}_Q)
\to \Ker(f' \otimes \text{id}_Q) \to 0
$$
因此 $f'$ 普遍单射，当且仅当对每个 $Q$ 都有
$\Ker(f \otimes \text{id}_Q ) \subset \Ker(g \otimes
\text{id}_Q)$，而这又当且仅当 $g$ 支配 $f$。
\end{proof}

\noindent
上述支配定义有时与通常的映射支配概念相关，如下引理所示。

\begin{lemma}
\label{algebra-lemma-domination}
设 $f: M \to N$ 和 $g: M \to M'$ 是 $R$-模映射。假设
$\Coker(f)$ 是有限表示的。那么 $g$ 支配 $f$，当且仅当 $g$ 经由
$f$ 分解，即存在模映射 $h: N
\to M'$，使得 $g = h \circ f$。
\end{lemma}

\begin{proof}
考虑引理 \ref{algebra-lemma-domination-universally-injective} 陈述中 $f$ 与 $g$
的推出。由推出的构造可知 $\Coker(f') = \Coker(f)$，故
$\Coker(f')$ 是有限表示的。于是由引理
\ref{algebra-lemma-universally-exact-split}，$f'$ 普遍单射，当且仅当
$$
0 \to M' \xrightarrow{f'} N' \to \Coker(f') \to 0
$$
分裂。这当且仅当存在映射 $h' : N' \to M'$，使
$h' \circ f' = \text{id}_{M'}$。由推出的泛性质，这样的 $h'$ 存在，
当且仅当 $g$ 经由 $f$ 分解。
\end{proof}


\begin{proposition}
\label{algebra-proposition-ML-characterization}
设 $M$ 是 $R$-模。设 $(M_i, f_{ij})$ 是以 $I$ 为指标的有限表示
$R$-模有向系，且 $M = \colim M_i$。令 $f_i:
M_i \to M$ 为典范映射。下列条件等价：
\begin{enumerate}
\item 对每个有限表示 $R$-模 $P$ 和模映射 $f: P
\to M$，都存在有限表示 $R$-模 $Q$ 和模映射 $g: P \to Q$，
使得 $g$ 与 $f$ 相互支配，即对每个 $R$-模 $N$ 都有
$\Ker(f \otimes_R \text{id}_N) = \Ker(g \otimes_R \text{id}_N)$。
\item 对每个 $i \in I$，都存在 $j \geq i$，使得
$f_{ij}: M_i \to M_j$ 支配 $f_i: M_i \to M$。
\item 对每个 $i \in I$，都存在 $j \geq i$，使得对所有
$k \geq i$，$f_{ij}: M_i \to M_j$ 都经由
$f_{ik}: M_i \to M_k$ 分解。
\item 对每个 $R$-模 $N$，逆系
$(\Hom_R(M_i, N), \Hom_R(f_{ij}, N))$ 都是 Mittag-Leffler 的。
\item 对 $N = \prod_{s \in I} M_s$，逆系
$(\Hom_R(M_i, N), \Hom_R(f_{ij}, N))$ 是 Mittag-Leffler 的。
\end{enumerate}
\end{proposition}

\begin{proof}
先证明 (1) 与 (2) 等价。假设 (1) 成立，并取 $i
\in I$。对应于映射 $f_i: M_i \to M$，可如 (1) 选取
$g: M_i \to Q$。由于 $M_i$ 和 $Q$ 都是有限表示的，
$\Coker(g)$ 也是有限表示的。于是由引理 \ref{algebra-lemma-domination}，
$f_i : M_i \to M$ 经由 $g: M_i \to Q$ 分解；设对某个
$h: Q \to M$ 有 $f_i = h \circ g$。又因为 $Q$ 有限表示，
$h$ 经由某个 $j \geq i$ 所对应的 $M_j \to M$ 分解；设对某个
$h': Q \to M_j$ 有 $h = f_j \circ h'$。综上得到交换图表
$$
\xymatrix{
  &  M  & \\
M_i \ar[dr]_g \ar[ur]^{f_i} \ar[rr]^{f_{ij}} &
& M_j \ar[ul]_{f_j} \\
  & Q  \ar[ur]_{h'} &
}
$$
因此 $f_{ij}$ 支配 $g$。而 $g$ 支配 $f_i$，所以 $f_{ij}$ 支配
$f_i$。

\medskip\noindent
反之，假设 (2) 成立。设 $P$ 是有限表示的，并设 $f: P
\to M$ 是模映射。那么对某个 $i \in I$，$f$ 经由
$f_i: M_i \to M$ 分解；设对某个 $g': P \to M_i$ 有
$f = f_i \circ g'$。由 (2) 选取 $j \geq i$，使 $f_{ij}$ 支配
$f_i$。我们有交换图表
$$
\xymatrix{
P \ar[d]_{g'} \ar[r]^{f}           & M  \\
M_i \ar[ur]^{f_i} \ar[r]_{f_{ij}} & M_j \ar[u]_{f_j}
}
$$
由该图表以及 $f_{ij}$ 支配 $f_i$ 这一事实可知，$f$ 与
$f_{ij} \circ g'$ 相互支配。因此取
$g = f_{ij} \circ g' : P \to M_j$ 即可。

\medskip\noindent
下面证明 (2) 与 (3) 等价。取 $i \in I$。对 $k \geq i$，
$f_i$ 总是支配 $f_{ik}$，因为 $f_i$ 经由 $f_{ik}$ 分解。
若 (2) 成立，选取 $j \geq i$，使得 $f_{ij}$ 支配 $f_i$。
由于支配关系具有传递性，对 $k \geq i$，$f_{ij}$ 支配 $f_{ik}$。
所有 $M_i$ 都是有限表示的，故对 $k \geq i$，
$\Coker(f_{ik})$ 是有限表示的。由引理 \ref{algebra-lemma-domination}，
对所有 $k \geq i$，$f_{ij}$ 经由 $f_{ik}$ 分解。因此 (2) 蕴含 (3)。
另一方面，若 (3) 成立，则对任意 $R$-模 $N$ 和 $k \geq i$，
$f_{ij} \otimes_R \text{id}_N$ 经由 $f_{ik}
\otimes_R \text{id}_N$ 分解。因此对 $k
\geq i$，有 $\Ker(f_{ik} \otimes_R
\text{id}_N) \subset \Ker(f_{ij} \otimes_R \text{id}_N)$。
而 $\Ker(f_i \otimes_R \text{id}_N: M_i \otimes_R N
\to M \otimes_R N)$ 是对 $k \geq i$ 的诸
$\Ker(f_{ik} \otimes_R
\text{id}_N)$ 之并。因此对任意 $R$-模 $N$，都有
$\Ker(f_i \otimes_R
\text{id}_N) \subset \Ker(f_{ij} \otimes_R \text{id}_N)$；
按定义，这意味着 $f_{ij}$ 支配 $f_i$。

\medskip\noindent
(3) 蕴含 (4)、(4) 蕴含 (5) 是显然的。下面证明 (5) 蕴含 (3)。令
$N = \prod_{s \in I} M_s$。若 (5) 成立，则给定 $i \in I$，选取
$j \geq i$，使得
$$
\Im( \Hom(M_j, N) \to  \Hom(M_i, N)) =
\Im( \Hom(M_k, N) \to  \Hom(M_i, N))
$$
对所有 $k \geq j$ 成立。把对 $s \in I$ 的乘积移到诸 $\Hom$ 外，
并考察乘积每个分量上的映射，这就是说
$$
\Im( \Hom(M_j, M_s) \to  \Hom(M_i, M_s)) =
\Im( \Hom(M_k, M_s) \to   \Hom(M_i, M_s))
$$
对所有 $k \geq j$ 和 $s \in I$ 成立。取 $s = j$，得到
$$
\Im( \Hom(M_j, M_j) \to  \Hom(M_i, M_j)) =
\Im( \Hom(M_k, M_j) \to   \Hom(M_i, M_j))
$$
对所有 $k \geq j$ 成立。由于 $f_{ij}$ 是
$\text{id} \in  \Hom(M_j, M_j)$ 在映射
$\Hom(M_j, M_j) \to  \Hom(M_i, M_j)$ 下的像，这说明对任意
$k \geq j$，都存在 $h \in \Hom(M_k, M_j)$，使
$f_{ij} = h \circ f_{ik}$。若 $j \geq k$，则可取
$h = f_{kj}$。因此 (3) 成立。
\end{proof}

\begin{definition}
\label{algebra-definition-mittag-leffler-module}
设 $M$ 是 $R$-模。若命题 \ref{algebra-proposition-ML-characterization}
中的诸等价条件成立，则称 $M$ 是 {\it Mittag-Leffler} 的。
\end{definition}

\noindent
特别地，有限表示模是 Mittag-Leffler 的。

\begin{remark}
\label{algebra-remark-flat-ML}
设 $M$ 是平坦 $R$-模。由 Lazard 定理（定理
\ref{algebra-theorem-lazard}），可把 $M = \colim M_i$ 写成有向系
$(M_i, f_{ij})$ 的余极限，其中各 $M_i$ 是有限自由 $R$-模。
要使 $M$ 成为 Mittag-Leffler 模，只需对偶逆系
$(\Hom_R(M_i, R), \Hom_R(f_{ij}, R))$ 是 Mittag-Leffler 的。
这由命题 \ref{algebra-proposition-ML-characterization} 的判据 (4)，以及如下事实
得出：对有限自由 $R$-模 $F$ 和任意 $R$-模 $N$，存在函子性同构
$\Hom_R(F, R) \otimes_R N \cong \Hom_R(F, N)$。
\end{remark}

\begin{lemma}
\label{algebra-lemma-tensor-ML-modules}
若 $R$ 是环，且 $M$、$N$ 是 $R$ 上的 Mittag-Leffler 模，
则 $M \otimes_R N$ 是 Mittag-Leffler 模。
\end{lemma}

\begin{proof}
把 $M = \colim_{i \in I} M_i$ 和 $N = \colim_{j \in J} N_j$
写成有限表示 $R$-模的有向余极限。分别以
$f_{ii'} : M_i \to M_{i'}$ 和 $g_{jj'} : N_j \to N_{j'}$ 表示
转移映射。于是 $M_i \otimes_R N_j$ 是有限表示 $R$-模（见引理
\ref{algebra-lemma-tensor-finiteness}），并且
% P01-E0452: frozen target preserves the source's M_j; see ERRATA_CANDIDATES.jsonl.
$M \otimes_R N = \colim_{(i, j) \in I \times J} M_i \otimes_R M_j$。
取 $(i, j) \in I \times J$。根据 Mittag-Leffler 模的定义，
两个系都满足命题 \ref{algebra-proposition-ML-characterization} (3)。
换言之，存在 $i' \geq i$ 和 $j' \geq j$，使得对每一组
$i'' \geq i$ 和 $j'' \geq j$，都存在映射
% P01-E0453: frozen target preserves the source's M-indexed b; see ERRATA_CANDIDATES.jsonl.
$a : M_{i''} \to M_{i'}$ 和 $b : M_{j''} \to M_{j'}$，满足
$f_{ii'} = a \circ f_{ii''}$ 和 $g_{jj'} = b \circ g_{jj''}$。
显然，
$a \otimes b : M_{i''} \otimes_R N_{j''} \to M_{i'} \otimes_R N_{j'}$
对系 $(M_i \otimes_R N_j, f_{ii'} \otimes g_{jj'})$ 起同样作用。
因此由刻画命题 \ref{algebra-proposition-ML-characterization} (3)，可知
$M \otimes_R N$ 是 Mittag-Leffler 的。
\end{proof}

\begin{lemma}
\label{algebra-lemma-ML-also}
设 $R$ 是环，$M$ 是 $R$-模。那么 $M$ 是 Mittag-Leffler 的，
当且仅当对每个有限自由 $R$-模 $F$ 和模映射 $f: F \to M$，
都存在有限表示 $R$-模 $Q$ 和模映射 $g : F \to Q$，使 $g$ 与
$f$ 相互支配，即对每个 $R$-模 $N$ 都有
$\Ker(f \otimes_R \text{id}_N) = \Ker(g \otimes_R \text{id}_N)$。
\end{lemma}

\begin{proof}
由于该条件显然弱于命题 \ref{algebra-proposition-ML-characterization} 的条件 (1)，
可见 Mittag-Leffler 模满足该条件。反之，假设 $M$ 满足该条件，且
$f : P \to M$ 是从有限表示 $R$-模 $P$ 到 $M$ 的 $R$-模映射。
选取满射 $F \to P$，其中 $F$ 是有限自由 $R$-模。由假设，
可找到映射 $F \to Q$，其中 $Q$ 是有限表示 $R$-模，并且
$F \to Q$ 与 $F \to M$ 相互支配。特别地，$F \to Q$ 的核包含
$F \to P$ 的核，因而得到 $R$-模映射 $g : P \to Q$，使
$F \to Q$ 等于复合 $F \to P \to Q$。设 $N$ 是任意 $R$-模，
并考虑交换图表
$$
\xymatrix{
F \otimes_R N \ar[d] \ar[r] & Q \otimes_R N \\
P \otimes_R N \ar[ru] \ar[r] & M \otimes_R N
}
$$
由假设，$F \otimes_R N \to Q \otimes_R N$ 与
$F \otimes_R N \to M \otimes_R N$ 的核相等。因此，因为
$F \otimes_R N \to P \otimes_R N$ 满射，
$P \otimes_R N \to Q \otimes_R N$ 与
$P \otimes_R N \to M \otimes_R N$ 的核也相等。
\end{proof}

\begin{lemma}
\label{algebra-lemma-restrict-ML-modules}
设 $R \to S$ 是有限且有限表示的环映射。设 $M$ 是 $S$-模。
若 $M$ 作为 $S$-模是 Mittag-Leffler 的，则 $M$ 作为 $R$-模也是
Mittag-Leffler 的。
\end{lemma}

\begin{proof}
假设 $M$ 作为 $S$-模是 Mittag-Leffler 的。把 $M = \colim M_i$
写成有限表示 $S$-模 $M_i$ 的一个有向余极限。由于 $M$ 在 $S$ 上
是 Mittag-Leffler 的，对每个 $i$，都存在指标 $j \geq i$，使得
对所有 $k \geq j$，都有分解 $f_{ij} = h \circ f_{ik}$（其中 $h$
依赖于 $i$、$j$ 的选取和 $k$）。注意，由引理
\ref{algebra-lemma-finite-finitely-presented-extension}，诸模 $M_i$ 作为
$R$-模也是有限表示的。此外，所有映射 $f_{ij}, f_{ik}, h$ 都是
$R$-模映射。因此，把系 $(M_i, f_{ij})$ 视为 $R$-模的系时，
它满足同一条件。所以 $M$ 作为 $R$-模是 Mittag-Leffler 的。
\end{proof}

\begin{lemma}
\label{algebra-lemma-mod-ideal-ML-modules}
设 $R$ 是环。设 $S = R/I$，其中 $I$ 是有限生成理想。
设 $M$ 是 $S$-模。那么 $M$ 作为 $R$-模是 Mittag-Leffler 的，
当且仅当 $M$ 作为 $S$-模是 Mittag-Leffler 的。
\end{lemma}

\begin{proof}
一个方向由引理 \ref{algebra-lemma-restrict-ML-modules} 得出。
为证明另一个方向，假设 $M$ 作为 $R$-模是 Mittag-Leffler 的。
把 $M = \colim M_i$ 写成有限表示 $S$-模的一个有向余极限。
由于 $I$ 有限生成，环 $S$ 作为 $R$-代数是有限且有限表示的，
故诸模 $M_i$ 作为 $R$-模是有限表示的，见引理
\ref{algebra-lemma-finite-finitely-presented-extension}。接着，设 $N$ 是任意
$S$-模。注意，对每个 $i$，都有
$\Hom_R(M_i, N) = \Hom_S(M_i, N)$，因为 $R \to S$ 满射。
因此，逆系 $(\Hom_R(M_i, N))_i$ 满足 Mittag-Leffler 条件，蕴含
逆系 $(\Hom_S(M_i, N))_i$ 满足 Mittag-Leffler 条件。故按定义，
$M$ 在 $S$ 上是 Mittag-Leffler 的。
\end{proof}

\begin{remark}
\label{algebra-remark-go-up-ML-modules}
设 $R \to S$ 是有限且有限表示的环映射。设 $M$ 是一个 $S$-模，
它作为 $R$-模是 Mittag-Leffler 的。一般而言，$M$ 作为 $S$-模未必
是 Mittag-Leffler 的。例如，假设 $S$ 是 $R$ 上的对偶数环，即
$S = R \oplus R\epsilon$，其中 $\epsilon^2 = 0$。那么一个 $S$-模
由一个 $R$-模 $M$ 连同一个平方为零的 $R$-线性自同态
$\epsilon : M \to M$ 构成。现在假设 $M_0$ 是一个非 Mittag-Leffler
的 $R$-模。选取表示 $F_1 \xrightarrow{u} F_0 \to M_0 \to 0$，
其中 $F_1$ 和 $F_0$ 是自由 $R$-模。令 $M = F_1 \oplus F_0$，并取
$$
\epsilon =
\left(
\begin{matrix}
0 & 0 \\
u & 0
\end{matrix}
\right) : M \longrightarrow M.
$$
那么 $M/\epsilon M \cong F_1 \oplus M_0$ 在
$R = S/\epsilon S$ 上不是 Mittag-Leffler 的，因而在 $S$ 上也不是
Mittag-Leffler 的（见引理 \ref{algebra-lemma-mod-ideal-ML-modules}）。
另一方面，$M/\epsilon M = M \otimes_S S/\epsilon S$；若 $M$ 是
Mittag-Leffler 的，则由引理 \ref{algebra-lemma-tensor-ML-modules}，它在
$S$ 上也会是 Mittag-Leffler 的。
\end{remark}

\section{直积与张量积的交换}
\label{algebra-section-products-tensor}

\noindent
设 $M$ 是 $R$-模，且 $(Q_{\alpha})_{\alpha \in A}$ 是一族 $R$-模。
那么有典范映射 $M \otimes_R \left( \prod_{\alpha
\in A} Q_{\alpha} \right) \to \prod_{\alpha \in A} ( M \otimes_R
Q_{\alpha})$，它在纯张量上由 $x \otimes (q_{\alpha}) \mapsto (x
\otimes q_{\alpha})$ 给出。如下例所示，该映射不一定单射，也不一定满射。

\begin{example}
\label{algebra-example-Q-not-ML}
取 $R = \mathbf{Z}$、$M = \mathbf{Q}$，并考虑当 $n \geq 1$ 时的族
$Q_n = \mathbf{Z}/n$。那么 $\prod_n (M \otimes Q_n) = 0$。
然而，将单射 $\mathbf{Z} \to \prod_n Q_n$ 与 $M$ 作张量，得到单射
$\mathbf{Q} \to M \otimes (\prod_n Q_n)$，所以
$M \otimes (\prod_n Q_n)$ 非零。因此
$M \otimes (\prod_n Q_n) \to \prod_n (M \otimes Q_n)$ 不是单射。

\medskip\noindent
另一方面，再取 $R = \mathbf{Z}$、$M = \mathbf{Q}$，并令当
$n \geq 1$ 时 $Q_n = \mathbf{Z}$。映射
$M \otimes (\prod_n Q_n)
\to \prod_n (M \otimes Q_n) = \prod_n M$ 的像恰由形如
$(a_n/m)_{n \geq 1}$ 的序列组成，其中 $a_n \in \mathbf{Z}$，而 $m$
是某个非零整数。因此该映射不是满射。
\end{example}

\noindent
下面确定使映射 $M
\otimes_R \left( \prod_{\alpha} Q_{\alpha} \right) \to \prod_{\alpha}
(M \otimes_R Q_{\alpha})$ 对 $(Q_{\alpha})_{\alpha \in A}$ 的所有选择
分别成为满射、双射或单射时，$M$ 所需满足的精确条件。这一点很重要，
因为使该映射单射的模恰好就是 Mittag-Leffler 模（命题
\ref{algebra-proposition-ML-tensor}）。下文中，若 $M$ 是 $R$-模且 $A$ 是集合，
我们用 $M^A$ 表示直积 $\prod_{\alpha \in A} M$。

\begin{proposition}
\label{algebra-proposition-fg-tensor}
设 $M$ 是 $R$-模。下列条件等价：
\begin{enumerate}
\item $M$ 有限生成。
\item 对每族 $R$-模 $(Q_{\alpha})_{\alpha \in A}$，典范映射
$M \otimes_R \left( \prod_{\alpha} Q_{\alpha} \right)
\to \prod_{\alpha} (M \otimes_R Q_{\alpha})$ 都是满射。
\item 对每个 $R$-模 $Q$ 和每个集合 $A$，典范映射
$M \otimes_R Q^{A} \to (M \otimes_R Q)^{A}$ 都是满射。
\item 对每个集合 $A$，典范映射 $M \otimes_R R^{A} \to
M^{A}$ 都是满射。
\end{enumerate}
\end{proposition}

\begin{proof}
先证明 (1) 蕴含 (2)。选取满射 $R^n \to M$，并考虑交换图表
$$
\xymatrix{
R^n \otimes_R (\prod_{\alpha} Q_{\alpha})  \ar[r]^{\cong} \ar[d] &
\prod_{\alpha} (R^n \otimes_R Q_{\alpha}) \ar[d] \\
M \otimes_R (\prod_{\alpha} Q_{\alpha})  \ar[r] & \prod_{\alpha} ( M
\otimes_R Q_{\alpha}).
}
$$
上方箭头是同构，两个竖直箭头是满射。因此下方箭头是满射。

\medskip\noindent
显然 (2) 蕴含 (3)，(3) 蕴含 (4)，所以只须证明 (4) 蕴含 (1)。
事实上，要使 (1) 成立，只须 $M^M$ 的元素
$d = (x)_{x \in M}$ 属于映射 $f: M \otimes_R R^{M} \to M^M$ 的像。
在这种情况下，对某些 $x_i \in M$ 和 $a_i \in R^M$，有
$d = \sum_{i = 1}^{n} f(x_i \otimes a_i)$。若对 $x \in M$，用
$p_x: M^M \to M$ 表示到第 $x$ 个因子的投影，则
% P01-E0455: p_x was typed only on M^M but is also applied to a_i in R^M; see ERRATA_CANDIDATES.jsonl.
$$
x = p_x(d) = \sum\nolimits_{i = 1}^{n} p_x(f(x_i \otimes a_i)) =
\sum\nolimits_{i=1}^{n} p_x(a_i) x_i.
$$
因此 $x_1, \ldots, x_n$ 生成 $M$。
\end{proof}

\begin{proposition}
\label{algebra-proposition-fp-tensor}
设 $M$ 是 $R$-模。下列条件等价：
\begin{enumerate}
\item $M$ 是有限表示的。
\item 对每族 $R$-模 $(Q_{\alpha})_{\alpha \in A}$，典范映射
$M \otimes_R \left( \prod_{\alpha} Q_{\alpha} \right)
\to \prod_{\alpha} (M \otimes_R Q_{\alpha})$ 都是双射。
\item 对每个 $R$-模 $Q$ 和每个集合 $A$，典范映射
$M \otimes_R Q^{A} \to (M \otimes_R Q)^{A}$ 都是双射。
\item 对每个集合 $A$，典范映射 $M \otimes_R R^{A} \to
M^{A}$ 都是双射。
\end{enumerate}
\end{proposition}

\begin{proof}
先证明 (1) 蕴含 (2)。选取表示 $R^m \to R^n
\to M$，并考虑交换图表
$$
\xymatrix{
R^m \otimes_R (\prod_{\alpha} Q_{\alpha}) \ar[r] \ar[d]^{\cong} & R^n
\otimes_R (\prod_{\alpha} Q_{\alpha}) \ar[r] \ar[d]^{\cong} & M \otimes_R
(\prod_{\alpha} Q_{\alpha}) \ar[r] \ar[d] & 0 \\
\prod_{\alpha} (R^m \otimes_R Q_{\alpha}) \ar[r] & \prod_{\alpha} (R^n
\otimes_R Q_{\alpha}) \ar[r] & \prod_{\alpha} (M \otimes_R Q_{\alpha})
\ar[r] & 0.
}
$$
前两个竖直箭头是同构，两行都正合。这说明映射
$M \otimes_R (\prod_{\alpha} Q_{\alpha}) \to
\prod_{\alpha} ( M \otimes_R Q_{\alpha})$
是满射；图追踪还说明它是单射。因此 (2) 成立。

\medskip\noindent
显然 (2) 蕴含 (3)，(3) 蕴含 (4)，所以只须证明 (4) 蕴含 (1)。
由命题 \ref{algebra-proposition-fg-tensor}，若 (4) 成立，我们已经知道 $M$
有限生成。因此可选取满射 $F \to M$，其中 $F$ 是有限自由模。
令 $K$ 为其核。必须证明 $K$ 有限生成。对任意集合 $A$，有交换图表
$$
\xymatrix{
& K \otimes_R R^A \ar[r] \ar[d]_{f_3} & F \otimes_R R^A \ar[r]
\ar[d]_{f_2}^{\cong} & M \otimes_R R^A \ar[r] \ar[d]_{f_1}^{\cong} & 0 \\
0 \ar[r] & K^A \ar[r] & F^A \ar[r] & M^A \ar[r] & 0 .
}
$$
由假设，映射 $f_1$ 是同构；因为 $F$ 是有限自由模，映射 $f_2$
是同构；两行都正合。图追踪说明 $f_3$ 满射，故由命题
\ref{algebra-proposition-fg-tensor}，$K$ 有限生成。
\end{proof}

\noindent
下一个命题需要如下引理。

\begin{lemma}
\label{algebra-lemma-kernel-tensored-fp}
设 $M$ 是 $R$-模，$P$ 是有限表示 $R$-模，且 $f: P
\to M$ 是映射。设 $Q$ 是 $R$-模，并假设 $x \in \Ker(P
\otimes Q \to M \otimes Q)$。那么存在有限表示 $R$-模 $P'$ 和映射
$f': P \to P'$，使得 $f$ 经由 $f'$ 分解，并且
$x \in \Ker(P \otimes Q \to P' \otimes Q)$。
\end{lemma}

\begin{proof}
把 $M$ 写成有限表示模 $M_i$ 的一个有向系的余极限
$M = \colim_{i \in I} M_i$。因为 $P$ 有限表示，映射 $f:
P \to M$ 经由某个 $j \in I$ 所对应的 $M_j \to M$ 分解。
与 $Q$ 作张量后，得到交换图表
$$
\xymatrix{
& M_j \otimes Q \ar[dr] & \\
P \otimes Q \ar[ur] \ar[rr] & & M \otimes Q .
}
$$
$x$ 在 $M_j \otimes Q$ 中的像 $y$ 属于
$M_j \otimes Q \to M \otimes Q$ 的核。由于
$M \otimes Q = \colim_{i \in I} (M_i
\otimes
Q)$，这意味着对某个 $j' \geq j$，$y$ 在 $M_{j'} \otimes Q$ 中
映为 $0$。因此可取 $P' = M_{j'}$，并取 $f'$ 为复合
$P \to M_j \to M_{j'}$。
\end{proof}

\begin{proposition}
\label{algebra-proposition-ML-tensor}
设 $M$ 是 $R$-模。下列条件等价：
\begin{enumerate}
\item $M$ 是 Mittag-Leffler 的。
\item 对每族 $R$-模 $(Q_{\alpha})_{\alpha \in A}$，典范映射
$M \otimes_R \left( \prod_{\alpha} Q_{\alpha} \right)
\to \prod_{\alpha} (M \otimes_R Q_{\alpha})$ 都是单射。
\end{enumerate}
\end{proposition}

\begin{proof}
先证明 (1) 蕴含 (2)。假设 $M$ 是 Mittag-Leffler 的，并设 $x$ 属于
$M \otimes_R (\prod_{\alpha} Q_{\alpha}) \to
\prod_{\alpha} (M \otimes_R Q_{\alpha})$ 的核。把 $M$ 写成有限表示模
$M_i$ 的一个有向系的余极限 $M =
\colim_{i \in I} M_i$。那么
$M \otimes_R (\prod_{\alpha} Q_{\alpha})$ 是诸
$M_i \otimes_R (\prod_{\alpha} Q_{\alpha})$ 的余极限。因此 $x$
是某个元素 $x_i \in M_i \otimes_R (\prod_{\alpha} Q_{\alpha})$ 的像。
必须说明对某个 $j \geq i$，$x_i$ 在
$M_j \otimes_R (\prod_{\alpha} Q_{\alpha})$ 中映为 $0$。
因为 $M$ 是 Mittag-Leffler 的，可选取 $j \geq i$，使
$M_i \to M_j$ 与 $M_i \to M$ 相互支配。然后考虑交换图表
$$
\xymatrix{
M \otimes_R (\prod_{\alpha} Q_{\alpha}) \ar[r] & \prod_{\alpha} (M
\otimes_R Q_{\alpha}) \\
M_i \otimes_R (\prod_{\alpha} Q_{\alpha}) \ar[r]^{\cong} \ar[d] \ar[u] &
\prod_{\alpha} (M_i \otimes_R Q_{\alpha}) \ar[d] \ar[u] \\
M_j \otimes_R (\prod_{\alpha} Q_{\alpha}) \ar[r]^{\cong} & \prod_{\alpha}
(M_j \otimes_R Q_{\alpha})
}
$$
根据命题 \ref{algebra-proposition-fp-tensor}，下方两个水平映射都是同构。
由于 $x_i$ 在 $\prod_{\alpha} (M
\otimes_R Q_{\alpha})$ 中映为 $0$，它在
$\prod_{\alpha} (M_i \otimes_R
Q_{\alpha})$ 中的像属于映射
$\prod_{\alpha} (M_i \otimes_R
Q_{\alpha}) \to \prod_{\alpha} (M \otimes_R Q_{\alpha})$ 的核。
由 $j$ 的选择，该核等于映射
$\prod_{\alpha} (M_i \otimes_R Q_{\alpha})
\to \prod_{\alpha} (M_j \otimes_R Q_{\alpha})$ 的核。
因此 $x_i$ 在 $\prod_{\alpha} (M_j \otimes_R
Q_{\alpha})$ 中映为 $0$，从而在
$M_j \otimes_R (\prod_{\alpha} Q_{\alpha})$ 中也映为 $0$。

\medskip\noindent
现在假设 (2) 成立。我们证明 $M$ 满足命题
\ref{algebra-proposition-ML-characterization} 中 Mittag-Leffler 性的表述 (1)。
设 $f: P \to M$ 是从有限表示模 $P$ 到 $M$ 的映射。
选取集合 $B$，其中含有有限表示 $R$-模的每个同构类的一个代表。
令 $A$ 为诸偶对 $(Q, x)$ 的集合，其中 $Q \in B$ 且
$x \in \Ker(P \otimes Q \to M \otimes Q)$。对
$\alpha = (Q, x) \in A$，记 $Q_{\alpha}$ 为 $Q$，记 $x_{\alpha}$ 为
$x$。考虑交换图表
$$
\xymatrix{
M \otimes_R (\prod_{\alpha} Q_{\alpha}) \ar[r] &
\prod_{\alpha} (M \otimes_R Q_{\alpha}) \\
P \otimes_R (\prod_{\alpha} Q_{\alpha}) \ar[r]^{\cong} \ar[u] &
\prod_{\alpha} (P \otimes_R Q_{\alpha}) \ar[u] .
}
$$
由假设，上方箭头是单射；由命题 \ref{algebra-proposition-fp-tensor}，
下方箭头是同构。令 $x \in P
\otimes_R (\prod_{\alpha} Q_{\alpha})$ 为该同构下对应于
$(x_{\alpha}) \in \prod_{\alpha} (P \otimes_R Q_{\alpha})$ 的元素。
由于图表中的上方箭头单射，有
$x \in \Ker( P \otimes_R (\prod_{\alpha} Q_{\alpha})
\to M \otimes_R (\prod_{\alpha} Q_{\alpha}))$。由引理
\ref{algebra-lemma-kernel-tensored-fp}，得到有限表示模 $P'$ 和映射
$f': P \to P'$，使 $f: P
\to M$ 经由 $f'$ 分解，并且
$x \in \Ker(P \otimes_R
(\prod_{\alpha} Q_{\alpha}) \to P' \otimes_R (\prod_{\alpha}
Q_{\alpha}))$。我们有交换图表
$$
\xymatrix{
P' \otimes_R (\prod_{\alpha} Q_{\alpha}) \ar[r]^{\cong} &
\prod_{\alpha} (P' \otimes_R Q_{\alpha}) \\
P \otimes_R (\prod_{\alpha} Q_{\alpha}) \ar[r]^{\cong} \ar[u] &
\prod_{\alpha} (P \otimes_R Q_{\alpha}) \ar[u] .
}
$$
其中上下两个箭头由命题 \ref{algebra-proposition-fp-tensor} 都是同构。
因此，因为 $x$ 属于左侧竖直映射的核，$(x_{\alpha})$ 属于右侧竖直
映射的核。这意味着对每个 $\alpha \in A$，都有
$x_{\alpha} \in \Ker(P \otimes_R Q_{\alpha} \to P' \otimes_R
Q_{\alpha})$。由 $A$ 的定义，这意味着对所有有限表示 $Q$，都有
$\Ker(P \otimes_R Q \to P' \otimes_R Q) \supset \Ker(P \otimes_R
Q \to M \otimes_R Q)$；而因为 $f: P
\to M$ 经由 $f': P \to P'$ 分解，实际上等号成立。
由引理 \ref{algebra-lemma-domination-fp}，$f$ 与 $f'$ 相互支配。
\end{proof}

\begin{lemma}
\label{algebra-lemma-minimal-contains}
设 $M$ 是 $R$ 上的平坦 Mittag-Leffler 模。设 $F$ 是 $R$-模，且
$x \in F \otimes_R M$。那么存在使 $x \in F' \otimes_R M$
成立的最小子模 $F' \subset F$。此外，$F'$ 是有限 $R$-模。
\end{lemma}

\begin{proof}
因为 $M$ 平坦，若 $F' \subset F$ 是子模，就有
$F' \otimes_R M \subset F \otimes_R M$，故该陈述有意义。
令 $I = \{F' \subset F \mid x \in F' \otimes_R M\}$；对
$i \in I$，用 $F_i \subset F$ 表示相应的子模。那么 $x$ 在映射
$$
F \otimes_R M \longrightarrow \prod (F/F_i \otimes_R M)
$$
下映为零，因而由命题 \ref{algebra-proposition-ML-tensor}，$x$ 在映射
$$
F \otimes_R M \longrightarrow \left(\prod F/F_i\right) \otimes_R M
$$
下映为零。由于 $M$ 平坦，该箭头的核是
$(\bigcap F_i) \otimes_R M$，这就证明 $F' = \bigcap F_i$。
为说明 $F'$ 是有限模，设
$x = \sum_{j = 1, \ldots, m} f_j \otimes m_j$，其中 $f_j \in F'$
且 $m_j \in M$。那么 $x \in F'' \otimes_R M$，其中
$F'' \subset F'$ 是由 $f_1, \ldots, f_m$ 生成的子模。
于是当然有 $F'' = F'$，最终陈述得证。
\end{proof}

\begin{lemma}
\label{algebra-lemma-pure-submodule-ML}
设 $0 \to M_1 \to M_2 \to M_3 \to 0$ 是 $R$-模的一个普遍正合列。
那么：
\begin{enumerate}
\item 若 $M_2$ 是 Mittag-Leffler 的，则 $M_1$ 是 Mittag-Leffler 的。
\item 若 $M_1$ 和 $M_3$ 是 Mittag-Leffler 的，则 $M_2$ 是
Mittag-Leffler 的。
\end{enumerate}
\end{lemma}

\begin{proof}
对任意一族 $R$-模 $(Q_{\alpha})_{\alpha \in A}$，都有交换图表
$$
\xymatrix{
0 \ar[r] & M_1 \otimes_R (\prod_{\alpha} Q_{\alpha}) \ar[r] \ar[d] & M_2
\otimes_R (\prod_{\alpha} Q_{\alpha}) \ar[r] \ar[d] & M_3 \otimes_R
(\prod_{\alpha} Q_{\alpha}) \ar[r] \ar[d] & 0 \\
0 \ar[r] & \prod_{\alpha}(M_1 \otimes Q_{\alpha}) \ar[r] & \prod_{\alpha}(M_2
\otimes Q_{\alpha}) \ar[r] & \prod_{\alpha}(M_3 \otimes Q_{\alpha})\ar[r] & 0
}
$$
两行都正合。因此 (1) 和 (2) 由命题 \ref{algebra-proposition-ML-tensor} 得出。
\end{proof}

\begin{lemma}
\label{algebra-lemma-quotient-module-ML}
设 $M_1 \to M_2 \to M_3 \to 0$ 是 $R$-模的正合列。
若 $M_1$ 有限生成且 $M_2$ 是 Mittag-Leffler 的，则 $M_3$ 是
Mittag-Leffler 的。
\end{lemma}

\begin{proof}
对任意一族 $R$-模 $(Q_{\alpha})_{\alpha \in A}$，由于张量积右正合，
有交换图表
$$
\xymatrix{
M_1 \otimes_R (\prod_{\alpha} Q_{\alpha}) \ar[r] \ar[d] & M_2
\otimes_R (\prod_{\alpha} Q_{\alpha}) \ar[r] \ar[d] & M_3 \otimes_R
(\prod_{\alpha} Q_{\alpha}) \ar[r] \ar[d] & 0 \\
\prod_{\alpha}(M_1 \otimes Q_{\alpha}) \ar[r] & \prod_{\alpha}(M_2
\otimes Q_{\alpha}) \ar[r] & \prod_{\alpha}(M_3 \otimes Q_{\alpha})\ar[r] & 0
}
$$
两行都正合。由命题 \ref{algebra-proposition-fg-tensor}，左侧竖直箭头满射。
由命题 \ref{algebra-proposition-ML-tensor}，中间竖直箭头单射。
图追踪说明右侧竖直箭头单射。故由命题
\ref{algebra-proposition-ML-tensor}，$M_3$ 是 Mittag-Leffler 的。
\end{proof}


\begin{lemma}
\label{algebra-lemma-colimit-universally-injective-ML}
若 $M = \colim M_i$ 是由 Mittag-Leffler $R$-模 $M_i$ 组成、且转移映射
普遍单射的有向系的余极限，则 $M$ 是 Mittag-Leffler 的。
\end{lemma}

\begin{proof}
设 $(Q_{\alpha})_{\alpha \in A}$ 是一族 $R$-模。必须说明
$M \otimes_R (\prod Q_\alpha) \to \prod M \otimes_R Q_\alpha$
单射；而对每个 $i$，已知
$M_i \otimes_R (\prod Q_\alpha) \to \prod M_i \otimes_R Q_\alpha$
单射，见命题 \ref{algebra-proposition-ML-tensor}。由于 $\otimes$ 与滤过余极限
交换，只须说明 $\prod M_i \otimes_R Q_\alpha \to \prod M \otimes_R Q_\alpha$
单射。这是显然的，因为根据转移映射普遍单射这一假设，每个映射
$M_i \otimes_R Q_\alpha \to M \otimes_R Q_\alpha$ 都是单射。
\end{proof}

\begin{lemma}
\label{algebra-lemma-direct-sum-ML}
若 $M = \bigoplus_{i \in I} M_i$ 是 $R$-模的直和，则 $M$ 是
Mittag-Leffler 的，当且仅当每个 $M_i$ 都是 Mittag-Leffler 的。
\end{lemma}

\begin{proof}
“仅当”方向由引理 \ref{algebra-lemma-pure-submodule-ML} (1) 以及分裂短正合列
普遍正合这一事实得出。反方向由引理
\ref{algebra-lemma-colimit-universally-injective-ML} 得出，但也可如下直接论证。
先注意，若 $I$ 有限，则结论由引理
\ref{algebra-lemma-pure-submodule-ML} (2) 得出。对一般的 $I$，若所有 $M_i$
都是 Mittag-Leffler 的，则通过验证命题
\ref{algebra-proposition-ML-characterization} 的条件 (1)，证明 $M$ 也如此。
设 $f: P \to M$ 是从有限表示模 $P$ 出发的映射。那么对 $I$ 的某个
有限子集 $I'$，$f$ 分解为
$P \xrightarrow{f'} \bigoplus_{i' \in I'} M_{i'} \hookrightarrow
\bigoplus_{i \in I} M_i$。由有限情形，
$\bigoplus_{i' \in I'} M_{i'}$ 是 Mittag-Leffler 的，故存在有限表示模
$Q$ 和映射 $g: P \to Q$，使 $g$ 与 $f'$ 相互支配。
于是 $g$ 与 $f$ 也相互支配。
\end{proof}

\begin{lemma}
\label{algebra-lemma-flat-ML-over-ML-ring}
设 $R \to S$ 是环映射。设 $M$ 是 $S$-模。若 $S$ 作为 $R$-模是
Mittag-Leffler 的，且 $M$ 作为 $S$-模是平坦并且 Mittag-Leffler 的，
则 $M$ 作为 $R$-模是 Mittag-Leffler 的。
\end{lemma}

\begin{proof}
我们由命题 \ref{algebra-proposition-ML-tensor} 的刻画推出这一点。
具体地，假设 $Q_\alpha$ 是一族 $R$-模。考虑复合
$$
\xymatrix{
M \otimes_R \prod_\alpha Q_\alpha =
M \otimes_S S \otimes_R \prod_\alpha Q_\alpha \ar[d] \\
M \otimes_S \prod_\alpha (S \otimes_R Q_\alpha) \ar[d] \\
\prod_\alpha (M \otimes_S S \otimes_R Q_\alpha) =
\prod_\alpha (M \otimes_R Q_\alpha)
}
$$
第一个箭头是单射，因为 $M$ 在 $S$ 上平坦，且 $S$ 在 $R$ 上是
Mittag-Leffler 的；第二个箭头是单射，因为 $M$ 在 $S$ 上是
Mittag-Leffler 的。因此 $M$ 在 $R$ 上是 Mittag-Leffler 的。
\end{proof}

\section{凝聚环}
\label{algebra-section-coherent}

\noindent
我们利用关于交换 $\prod$ 与 $\otimes$ 的讨论，确定哪些环具有如下性质：
平坦模的直积仍然平坦。结果表明，这些恰是所谓的凝聚环。
你或许更熟悉环化空间上的凝聚 $\mathcal{O}_X$-模这一概念，见
Modules，第 \ref{modules-section-coherent} 节。

\begin{definition}
\label{algebra-definition-coherent}
设 $R$ 是环。设 $M$ 是 $R$-模。
\begin{enumerate}
\item 若 $M$ 有限生成，且 $M$ 的每个有限生成子模在 $R$ 上都是
有限表示的，则称其为{\it 凝聚模}。
\item 若 $R$ 作为自身上的模是凝聚的，则称其为{\it 凝聚环}。
\end{enumerate}
\end{definition}

\noindent
因此，一个环是凝聚的，当且仅当每个有限生成理想作为模都是有限表示的。

\begin{example}
\label{algebra-example-valuation-ring-coherent}
赋值环是凝聚环。事实上，每个非零有限生成理想都是主理想（引理
\ref{algebra-lemma-characterize-valuation-ring}）；又因为赋值环是整环，
该理想是自由的，因而是有限表示的。
\end{example}

\noindent
凝聚模范畴是阿贝尔范畴。

\begin{lemma}
\label{algebra-lemma-coherent}
设 $R$ 是环。
\begin{enumerate}
\item 凝聚模的有限子模是凝聚的。
\item 设 $\varphi : N \to M$ 是从有限模到凝聚模的同态。那么
$\Ker(\varphi)$ 有限，$\Im(\varphi)$ 凝聚，并且
$\Coker(\varphi)$ 凝聚。
\item 设 $\varphi : N \to M$ 是凝聚模之间的同态。那么
$\Ker(\varphi)$ 和 $\Coker(\varphi)$ 都是凝聚模。
\item 给定 $R$-模的短正合列
$0 \to M_1 \to M_2 \to M_3 \to 0$，若其中两个模凝聚，
则第三个也凝聚。
\end{enumerate}
\end{lemma}

\begin{proof}
第一条陈述由定义立即得出。

\medskip\noindent
设 $\varphi : N \to M$ 满足 (2) 的假设。首先，
$\Im(\varphi)$ 有限，因而由 (1) 是凝聚的。特别地，
$\Im(\varphi)$ 有限表示；于是把引理 \ref{algebra-lemma-extension} 应用于
正合列 $0 \to \Ker(\varphi) \to N \to \Im(\varphi) \to 0$，
可知 $\Ker(\varphi)$ 有限。为证明 $\Coker(\varphi)$ 凝聚，
设 $E \subset \Coker(\varphi)$ 是有限子模，并令 $E'$ 为它在 $M$
中的逆像。由正合列
$0 \to \Im(\varphi) \to E' \to E \to 0$，并且由于
% P01-E0456: the source invokes finite Ker(phi), but this exact sequence requires finite Im(phi); see ERRATA_CANDIDATES.jsonl.
$\Ker(\varphi)$ 有限，由引理 \ref{algebra-lemma-extension} 可知
$E' \subset M$ 有限，因而因为 $M$ 凝聚，它是有限表示的。
同一正合列继而说明 $E$ 有限表示，断言得证。

\medskip\noindent
(3) 立即由 (1) 和 (2) 得出。

\medskip\noindent
设 $0 \to M_1 \xrightarrow{i} M_2 \xrightarrow{p} M_3 \to 0$
是 (4) 中的 $R$-模短正合列。只须证明：若 $M_1$ 和 $M_3$ 凝聚，
则 $M_2$ 也凝聚。由引理 \ref{algebra-lemma-extension}，$M_2$ 有限。
设 $N_2 \subset M_2$ 是有限子模。令
$N_3 = p(N_2) \subset M_3$，并令
$N_1 = i^{-1}(N_2) \subset M_1$。我们有正合列
$0 \to N_1 \to N_2 \to N_3 \to 0$。显然 $N_3$ 有限（因为它是
$N_2$ 的商），因而是有限表示的（因为它是 $M_3$ 的有限子模）。
由引理 \ref{algebra-lemma-extension} (5)，$N_1$ 有限，因而是有限表示的
（因为它是 $M_1$ 的有限子模）。再由引理
\ref{algebra-lemma-extension} (2)，$N_2$ 有限表示。
\end{proof}

\begin{lemma}
\label{algebra-lemma-coherent-ring}
设 $R$ 是环。若 $R$ 凝聚，则一个模凝聚，当且仅当它有限表示。
\end{lemma}

\begin{proof}
显然，凝聚模是有限表示的（在任意环上都是如此）。反之，若 $R$ 凝聚，
则 $R^{\oplus n}$ 凝聚，且任意映射
$R^{\oplus m} \to R^{\oplus n}$ 的余核也凝聚，见引理
\ref{algebra-lemma-coherent}。
\end{proof}

\begin{lemma}
\label{algebra-lemma-Noetherian-coherent}
Noether 环是凝聚环。
\end{lemma}

\begin{proof}
由引理 \ref{algebra-lemma-Noetherian-finite-type-is-finite-presentation}，
任意有限 $R$-模都是有限表示的。特别地，$R$ 的任意理想都是有限表示的。
\end{proof}

\begin{proposition}
\label{algebra-proposition-characterize-coherent}
\begin{reference}
这是 \cite[Theorem 2.1]{Chase}。
\end{reference}
设 $R$ 是环。下列条件等价：
\begin{enumerate}
\item $R$ 凝聚；
\item 平坦 $R$-模的任意直积都平坦；
\item 对每个集合 $A$，模 $R^A$ 都平坦。
\end{enumerate}
\end{proposition}

\begin{proof}
假设 $R$ 凝聚，并设 $Q_\alpha$（$\alpha \in A$）是一族平坦
$R$-模。必须证明，对 $R$ 的每个有限生成理想 $I$，映射
$I \otimes_R \prod_\alpha Q_\alpha \to \prod Q_\alpha$ 都单射，
见引理 \ref{algebra-lemma-flat}。由于 $R$ 凝聚，$I$ 是有限表示 $R$-模。
因此由命题 \ref{algebra-proposition-fp-tensor}，有
$I \otimes_R \prod_\alpha Q_\alpha = \prod I \otimes_R Q_\alpha$。
所需的单射性由此得出，因为根据 $Q_\alpha$ 的平坦性，
$I \otimes_R Q_\alpha \to Q_\alpha$ 单射。

\medskip\noindent
(2) $\Rightarrow$ (3) 是显然的。

\medskip\noindent
假设对每个集合 $A$，$R$-模 $R^A$ 都平坦。设 $I$ 是 $R$ 的有限生成
理想。由假设，$I \otimes_R R^A \to R^A$ 单射。由命题
\ref{algebra-proposition-fg-tensor} 以及 $I$ 的有限性，其像等于 $I^A$。
因此对每个集合 $A$，都有 $I \otimes_R R^A = I^A$；再由命题
\ref{algebra-proposition-fp-tensor} 可知 $I$ 有限表示。
\end{proof}

\section{Mittag-Leffler 模的例与反例}
\label{algebra-section-examples}

\noindent
本节最后给出一些 Mittag-Leffler 模的例与反例。

\begin{example}
\label{algebra-example-ML}
Mittag-Leffler 模。
\begin{enumerate}
\item 任意有限表示模都是 Mittag-Leffler 的。例如，这由命题
\ref{algebra-proposition-ML-characterization} (1) 得出。一般而言，有限生成模
是 Mittag-Leffler 的，当且仅当它是有限表示的。这由命题
\ref{algebra-proposition-fg-tensor}、\ref{algebra-proposition-fp-tensor} 和
\ref{algebra-proposition-ML-tensor} 得出。
\item 自由模是 Mittag-Leffler 的，因为它满足命题
\ref{algebra-proposition-ML-characterization} 的条件 (1)。
\item 由上一个例子和引理 \ref{algebra-lemma-direct-sum-ML}，投射模是
Mittag-Leffler 的。
\end{enumerate}
\end{example}

\noindent
我们还想把 Noether 环 $R$ 上的幂级数环加入例子列表。这将由如下引理推出。

\begin{lemma}
\label{algebra-lemma-flat-ML-criterion}
设 $M$ 是平坦 $R$-模。下列条件等价：
\begin{enumerate}
\item $M$ 是 Mittag-Leffler 的；
\item 若 $F$ 是有限自由 $R$-模，且 $x \in F \otimes_R M$，则存在
$F$ 的最小子模 $F'$，使得 $x \in F' \otimes_R M$。
\end{enumerate}
\end{lemma}

\begin{proof}
(1) $\Rightarrow$ (2) 是引理 \ref{algebra-lemma-minimal-contains} 的特殊情形。
假设 (2)。由定理 \ref{algebra-theorem-lazard}，可把 $M$ 写成有限自由
$R$-模有向系 $(M_i, f_{ij})$ 的余极限
$M = \colim_{i \in I} M_i$。由评注 \ref{algebra-remark-flat-ML}，只须说明
对偶逆系
$$
(\Hom_R(M_i, R), \Hom_R(f_{ij}, R))
$$
是 Mittag-Leffler 的。
换言之，固定 $i \in I$，并对 $j \geq i$，令 $Q_j$ 为
$\Hom_R(M_j, R)
\to \Hom_R(M_i, R)$ 的像；必须证明诸 $Q_j$ 稳定。

\medskip\noindent
因为 $M_i$ 有限自由，对所有 $j$，都有等同
$\Hom_R(M_i, M_j) = \Hom_R(M_i, R) \otimes_R  M_j$。
利用诸 $M_j$ 自由这一事实可知，对 $j \geq i$，$Q_j$ 是满足
$f_{ij} \in Q_j
\otimes_R
M_j$ 的 $\Hom_R(M_i, R)$ 的最小子模。在等同
$\Hom_R(M_i, M) = \Hom_R(M_i, R)
\otimes_R
M$ 下，典范映射 $f_i: M_i \to M$ 属于
$\Hom_R(M_i, R)
\otimes_R M$。由关于 $M$ 的假设，存在 $\Hom_R(M_i, R)$ 的最小子模
$Q$，使得 $f_i \in Q \otimes_R M$。下面证明诸 $Q_j$ 稳定为 $Q$。

\medskip\noindent
对 $j \geq i$，有交换图表
$$
\xymatrix{
Q_j \otimes_R M_j \ar[r] \ar[d] & \Hom_R(M_i, R) \otimes_R M_j
\ar[d] \\
Q_j \otimes_R M \ar[r] & \Hom_R(M_i, R) \otimes_R M.
}
$$
由于 $f_{ij} \in Q_j \otimes_R M_j$ 映到
$f_i \in \Hom_R(M_i, R)
\otimes_R M$，可知 $f_i \in Q_j \otimes_R M$。因此，由 $Q$ 的选择，
对所有 $j \geq i$，都有 $Q \subset Q_j$。

\medskip\noindent
因为诸 $Q_j$ 递减，且对所有 $j \geq i$ 都有 $Q \subset Q_j$，
要说明诸 $Q_j$ 稳定为 $Q$，只须找到 $j \geq i$，使
$Q_j \subset Q$。把这个典范映射看作下式的元素：
% P01-E0459: the source changes the established index set I to undefined J in both colimits; see ERRATA_CANDIDATES.jsonl.
$$
\Hom_R(M_i, R) \otimes_R M = \colim_{j \in J}
(\Hom_R(M_i, R) \otimes_R
M_j),
$$
则 $f_i$ 是当 $j \geq i$ 时诸 $f_{ij}$ 的余极限，并且 $f_i$ 也属于子模
$$
\colim_{j \in J} (Q \otimes_R M_j) \subset \colim_{j \in J}
(\Hom_R(M_i, R)
\otimes_R M_j).
$$
因此，对某个 $j \geq i$，$f_{ij}$ 属于 $Q \otimes_R M_j$。
由于 $Q_j$ 是满足 $f_{ij}
\in
Q_j \otimes_R M_j$ 的 $\Hom_R(M_i, R)$ 的最小子模，故
$Q_j\subset Q$。
\end{proof}

\begin{lemma}
\label{algebra-lemma-product-over-Noetherian-ring}
设 $R$ 是 Noether 环，$A$ 是集合。那么
$M = R^A$ 是平坦的 Mittag-Leffler $R$-模。
\end{lemma}

\begin{proof}
结合引理 \ref{algebra-lemma-Noetherian-coherent} 与命题
\ref{algebra-proposition-characterize-coherent}，可知 $M$ 在 $R$ 上平坦。
下面说明 $M$ 满足引理 \ref{algebra-lemma-flat-ML-criterion} 的条件。
设 $F$ 是有限自由 $R$-模。若 $F'$ 是 $F$ 的任意子模，则由于 $R$
是 Noether 环，它是有限表示的。因此由命题
\ref{algebra-proposition-fp-tensor}，有交换图表
$$
\xymatrix{
F' \otimes_R M \ar[r] \ar[d]^{\cong} & F \otimes_R M \ar[d]^{\cong} \\
(F')^A \ar[r] & F^A
}
$$
借此可把映射 $F' \otimes_R M \to F \otimes_R M$ 与
$(F')^A \to F^A$ 等同。因此，若 $x \in F \otimes_R M$ 对应于
$(x_\alpha) \in F^A$，则由诸 $x_\alpha$ 生成的 $F$ 的子模 $F'$
是 $F$ 的满足 $x \in F' \otimes_R M$ 的最小子模。
% P01-E0460: the translated sentence renders the clear intended meaning of the malformed source phrase; see ERRATA_CANDIDATES.jsonl.
\end{proof}

\begin{lemma}
\label{algebra-lemma-power-series-ML}
设 $R$ 是 Noether 环，$n$ 是正整数。那么 $R$-模
$M = R[[t_1, \ldots, t_n]]$ 平坦且是 Mittag-Leffler 的。
\end{lemma}

\begin{proof}
作为 $R$-模，对某个（可数）集合 $A$，有 $M = R^A$。
因此该引理是引理 \ref{algebra-lemma-product-over-Noetherian-ring} 的特殊情形。
\end{proof}

\begin{example}
\label{algebra-example-not-ML}
非 Mittag-Leffler 模。
\begin{enumerate}
\item 由例 \ref{algebra-example-Q-not-ML} 和命题
\ref{algebra-proposition-ML-tensor}，$\mathbf{Q}$ 不是 Mittag-Leffler
$\mathbf{Z}$-模。
\item 下文将证明（定理 \ref{algebra-theorem-projectivity-characterization}）：
对平坦可数生成模，投射性等价于 Mittag-Leffler 性。因此，任意平坦、
可数生成且非投射的模 $M$ 都是非 Mittag-Leffler 模的例子。
这样的例子见评注 \ref{algebra-remark-warning}。
\item 设 $k$ 是域。令 $R = k[[x]]$。$R$-模
$M = \prod_{n \in \mathbf{N}} R/(x^n)$ 不是 Mittag-Leffler 的。
具体地，考虑元素 $\xi = (\xi_1, \xi_2, \xi_3, \ldots)$，其定义为
$\xi_{2^m} = x^{2^{m - 1}}$，并在其余情形取 $\xi_n = 0$，所以
$$
\xi = (0, x, 0, x^2, 0, 0, 0, x^4, 0, 0, 0, 0, 0, 0, 0, x^8, \ldots)
$$
当 $m \gg 0$ 时，$\xi$ 在 $M/x^{2^m}M$ 中的零化子由
$x^{2^{m - 1}}$ 生成。但是，若 $M$ 是 Mittag-Leffler 的，
则存在有限 $R$-模 $Q$ 和元素 $\xi' \in Q$，使得对所有
$l \geq 1$，$\xi'$ 在 $Q/x^l Q$ 中的零化子与 $\xi$ 在
$M/x^l M$ 中的零化子一致，见命题
\ref{algebra-proposition-ML-characterization} (1)。现在可以证明，存在整数
$a \geq 0$，使得对所有 $l \gg 0$，$\xi'$ 在 $Q/x^l Q$ 中的零化子
由 $x^a$ 或 $x^{l - a}$ 生成（取决于 $\xi' \in Q$ 是否为挠元）。
结合以上结论，会得到对所有 $l = 2^m >> 0$，都有
% P01-E0462: frozen target preserves >>, which renders as two greater-than signs rather than \gg; see ERRATA_CANDIDATES.jsonl.
$a = l/2$ 或 $l - a = l/2$，这是荒谬的。
\item 同一论证说明，$\bigoplus_{n \in \mathbf{N}} R/(x^n)$ 的
$(x)$-进完备化在 $R = k[[x]]$ 上不是 Mittag-Leffler 的
（提示：$\xi$ 实际上是该完备化的一个元素）。
\item 令 $R = k[a, b]/(a^2, ab, b^2)$。令 $S$ 是具有表示
$S = R[t]/(at - b)$ 的有限表示 $R$-代数。那么作为 $R$-模，
$S$ 可数生成且不可分解（细节略）。另一方面，$R$ 是 Artin 局部环，
因而是完备局部环，进而是 Hensel 局部环，见引理
\ref{algebra-lemma-complete-henselian}。若 $S$ 作为 $R$-模是 Mittag-Leffler 的，
则由引理 \ref{algebra-lemma-split-ML-henselian}，它会是有限 $R$-模的直和。
因此，$S$ 作为 $R$-模不是 Mittag-Leffler 的。
\end{enumerate}
\end{example}

\section{可数生成的 Mittag-Leffler 模}
\label{algebra-section-ML-countable}

\noindent
事实证明，可数生成的 Mittag-Leffler 模具有特别简单的结构。

\begin{lemma}
\label{algebra-lemma-ML-countable-colimit}
设 $M$ 是 $R$-模。写成 $M = \colim_{i \in I} M_i$，其中
$(M_i, f_{ij})$ 是由有限表示 $R$-模组成的有向系。若 $M$ 是
Mittag-Leffler 且可数生成，则存在有向可数子集 $I' \subset I$，
使得 $M \cong \colim_{i \in I'} M_i$。
\end{lemma}

\begin{proof}
设 $x_1, x_2, \ldots$ 是 $M$ 的一个可数生成集。对每个 $x_n$，
选取 $i \in I$，使 $x_n$ 属于典范映射 $f_i:
M_i \to M$ 的像；令 $I'_{0} \subset I$ 为所有这些 $i$ 的集合。
现在，由于 $M$ 是 Mittag-Leffler 的，对每个 $i \in I'_{0}$，
可以选取 $j \in I$，使 $j \geq i$，并且对所有 $k \geq i$，
$f_{ij}: M_i \to M_j$ 都经由 $f_{ik}: M_i \to M_k$ 分解
（命题 \ref{algebra-proposition-ML-characterization} 的条件 (3)）；
令 $I'_1$ 为 $I'_0$ 与所有这些 $j$ 的并。由于 $I'_1$ 可数，
可以把它扩大为可数有向集 $I'_{2} \subset I$。现在对 $I'_{2}$
施行与对 $I'_{0}$ 相同的步骤，得到新的可数集
$I'_{3} \subset I$。再把 $I'_{3}$ 扩大为可数有向集 $I'_{4}$。
% P01-E0464: the frozen source reverses the even/odd stages relative to its explicit I'_0 -> I'_1 -> I'_2 construction; see ERRATA_CANDIDATES.jsonl.
以这种方式继续下去——若 $\ell$ 为奇数，则对每个
$ i \in I'_{\ell}$ 加入命题
\ref{algebra-proposition-ML-characterization} (3) 中的一个 $j$；若 $\ell$
为偶数，则把 $I'_{\ell}$ 扩大为有向集——便得到一列子集
$I'_{\ell} \subset I$，其中 $\ell \geq 0$。并集 $I' =
\bigcup I'_{\ell}$ 满足：
\begin{enumerate}
\item $I'$ 可数且有向；
\item 每个 $x_n$ 都属于某个 $i \in I'$ 所对应的
$f_i: M_i \to M$ 的像；
\item 若 $i \in I'$，则存在 $j \in I'$，使得 $j \geq i$，并且对所有
满足 $k \geq i$ 的 $k \in I$，$f_{ij}:
M_i \to M_j$ 都经由 $f_{ik}: M_i \to M_k$ 分解。特别地，对
$k \geq i$ 有 $\Ker(f_{ik}) \subset \Ker(f_{ij})$。
\end{enumerate}
我们断言典范映射 $\colim_{i \in I'} M_i \to
\colim_{i
\in I} M_i = M$ 是同构。由 (2)，它是满射。为证其为单射，
设 $x \in \colim_{i \in I'} M_i$ 在 $\colim_{i \in
I} M_i$ 中映为 $0$。
用某个 $i \in I'$ 的元素 $\tilde{x} \in M_i$ 表示 $x$，
这意味着对某个 $k \in I, k \geq i$，有
$f_{ik}(\tilde{x}) = 0$。但是，由 (3)，存在
$j \in I', j \geq i,$，使 $f_{ij}(\tilde{x}) = 0$。因此，
$x = 0$ 在 $\colim_{i \in I'} M_i$ 中成立。
\end{proof}

\noindent
引理 \ref{algebra-lemma-ML-countable-colimit}
表明，环 $R$ 上可数生成的 Mittag-Leffler 模 $M$ 是如下系统的余极限：
$$
M_1 \to M_2 \to M_3 \to M_4 \to \ldots
$$
其中每个 $M_n$ 都是有限表示 $R$-模。为说明这一点，仿照引理
\ref{algebra-lemma-ML-limit-nonempty}
% P01-E0465: the frozen source gives the direct-system cofinal order as (N, >=), conflicting with the displayed forward system and the referenced increasing construction; see ERRATA_CANDIDATES.jsonl.
的证明，可知可数有向集具有一个与 $(\mathbf{N}, \geq)$
同构的共尾子集。设
$R = k[x_1, x_2, x_3, \ldots]$ 且 $M = R/(x_i)$。则 $M$
有限生成但非有限表示，因而不是 Mittag-Leffler 的（见例
\ref{algebra-example-ML} 的 (1)）。
不过，当然可以取
$M_n = R/(x_1, \ldots, x_n)$，从而写成 $M = \colim_n M_n$；
因此，能把 $M$ 写成这样的极限并不蕴含 $M$ 是 Mittag-Leffler 的。

\begin{lemma}
\label{algebra-lemma-ML-countable}
设 $R$ 是环。
设 $M$ 是 $R$-模。
假设 $M$ 是 Mittag-Leffler 且可数生成。
对任意 $R$-模映射 $f : P \to M$，若 $P$ 有限生成，则存在
自同态 $\alpha : M \to M$，使得
\begin{enumerate}
\item $\alpha : M \to M$ 经由一个有限表示 $R$-模分解，并且
\item $\alpha \circ f = f$。
\end{enumerate}
\end{lemma}

\begin{proof}
把 $M = \colim_{i \in I} M_i$ 写成由有限表示 $R$-模组成的
有向余极限，其中 $I$ 可数，见引理
\ref{algebra-lemma-ML-countable-colimit}。
转移映射记为 $f_{ij}$，并用 $f_i : M_i \to M$
表示映入 $M$ 的典范映射。置 $N = \prod_{s \in I} M_s$。记
$$
M_i^* = \Hom_R(M_i, N) = \prod\nolimits_{s \in I} \Hom_R(M_i, M_s)
$$
于是 $(M_i^*)$ 是指标集为 $I$ 的 $R$-模逆系。
注意 $\Hom_R(M, N) = \lim M_i^*$。
由于 $M$ 是 Mittag-Leffler 的，对每个
$i \in I$，可以找到指标 $k(i) \geq i$，使得
$$
E_i := \bigcap\nolimits_{i' \geq i} \Im(M_{i'}^* \to M_i^*)
=
\Im(M_{k(i)}^* \to M_i^*)
$$
选定 $j \in I$，使 $\Im(P \to M) \subset \Im(M_j \to M)$。
这是可能的，因为 $P$ 有限生成。置 $k = k(j)$。令
$x = (0, \ldots, 0, \text{id}_{M_k}, 0, \ldots, 0) \in M_k^*$，并注意
它映到 $y = (0, \ldots, 0, f_{jk}, 0, \ldots, 0) \in M_j^*$。
由 $k$ 的选取可知 $y \in E_j$。由例
\ref{algebra-example-ML-surjective-maps}，
对每个 $i \geq j$，转移映射 $E_i \to E_j$ 都是满射，并且
$\lim E_i = \lim M_i^* = \Hom_R(M, N)$。因此，引理
\ref{algebra-lemma-ML-limit-nonempty}
保证存在元素 $z \in \Hom_R(M, N)$，
它在 $E_j \subset M_j^*$ 中映到 $y$。令 $z_k$ 为 $z$ 的第 $k$ 个分量。
那么 $z_k : M \to M_k$ 是同态，并且下图
$$
\xymatrix{
M \ar[r]_{z_k} & M_k \\
M_j \ar[ru]_{f_{jk}} \ar[u]^{f_j}
}
$$
交换。令 $\alpha : M \to M$ 为复合
$f_k \circ z_k : M \to M_k \to M$。
按构造，$\alpha$ 经由有限表示模分解，并且
$\alpha \circ f_j = f_j$。由于 $f$ 的像包含在 $f_j$ 的像中，
这也蕴含 $\alpha \circ f = f$。
\end{proof}

\noindent
稍后将看到（见引理
\ref{algebra-lemma-split-ML-henselian}），
引理 \ref{algebra-lemma-ML-countable}
意味着 Hensel 局部环上的可数生成 Mittag-Leffler 模是有限表示模的直和。

\section{投射模的刻画}
\label{algebra-section-characterize-projective}

\noindent
本节的目标是证明，一个模是投射的，当且仅当它平坦、
是 Mittag-Leffler 的，并且是可数生成模的直和
（见下文定理 \ref{algebra-theorem-projectivity-characterization}）。

\begin{lemma}
\label{algebra-lemma-countgen-projective}
设 $M$ 是 $R$-模。若 $M$ 平坦、是 Mittag-Leffler 的且可数生成，
则 $M$ 是投射的。
\end{lemma}

\begin{proof}
由 Lazard 定理（定理 \ref{algebra-theorem-lazard}），可以把 $M =
\colim_{i
\in I} M_i$ 写成由有限自由 $R$-模构成、指标集为集合 $I$ 的有向系
$(M_i, f_{ij})$。由引理 \ref{algebra-lemma-ML-countable-colimit}，
可以假设 $I$ 可数。现在令
$$
0 \to N_1 \to N_2 \to N_3 \to 0
$$
是 $R$-模的正合列。我们必须证明，施用
$\Hom_R(M, -)$
后仍保持正合性。由于 $M_i$ 是有限自由的，
$$
0 \to \Hom_R(M_i, N_1) \to \Hom_R(M_i, N_2) \to
\Hom_R(M_i, N_3) \to 0
$$
对每个 $i$ 都是正合的。由于 $M$ 是 Mittag-Leffler 的，
$(\Hom_R(M_i, N_{1}))$ 是 Mittag-Leffler 逆系。因此由引理
\ref{algebra-lemma-ML-exact-sequence}，
$$
0 \to \lim_{i \in I} \Hom_R(M_i, N_1) \to
\lim_{i \in I} \Hom_R(M_i, N_2) \to
\lim_{i \in I} \Hom_R(M_i, N_3) \to 0
$$
正合。但对任意 $R$-模 $N$，都有函子性的同构
$\Hom_R(M, N) \cong \lim_{i \in I} \Hom_R(M_i, N)$，所以
$$
0 \to \Hom_R(M, N_1) \to \Hom_R(M, N_2) \to
\Hom_R(M, N_3) \to 0
$$
正合。
\end{proof}

\begin{remark}
\label{algebra-remark-characterize-projective}
若没有可数生成这一假设，引理 \ref{algebra-lemma-countgen-projective}
不成立。例如，$\mathbf Z$-模 $M =
\mathbf{Z}[[x]]$ 平坦且是 Mittag-Leffler 的，但并不投射。
由引理 \ref{algebra-lemma-power-series-ML}，它是 Mittag-Leffler 的。
自由 Abel 群的子群是自由的，因而投射 $\mathbf Z$-模事实上是自由的，
其子模也自由。因此，为证明 $M$ 不投射，只须构造一个非自由子模。
固定一个素数 $p$，并考虑由如下幂级数组成的子模 $N$：
$f(x) = \sum a_i x^i$，使得对每个整数 $m
\geq 1$，除有限多个 $i$ 外，$p^m$ 都整除 $a_i$。于是对所有
$a_i \in \mathbf{Z}$，$\sum a_i p^i x^i$ 都属于 $N$，所以 $N$ 不可数。
因此，若 $N$ 自由，则它的秩不可数，并且 $N/pN$ 在
$\mathbf{Z}/p$ 上的维数不可数。这并不成立，因为诸元素
$x^i \in N/pN$（其中 $i \geq 0$）张成 $N/pN$。
\end{remark}

\begin{theorem}
\label{algebra-theorem-projectivity-characterization}
设 $M$ 是 $R$-模。那么 $M$ 投射，当且仅当
\begin{enumerate}
\item $M$ 平坦；
\item $M$ 是 Mittag-Leffler 的；
\item $M$ 是可数生成 $R$-模的直和。
\end{enumerate}
\end{theorem}

\begin{proof}
先设 $M$ 投射。则 $M$ 是自由模的直和因子，因此 $M$ 平坦且是
Mittag-Leffler 的，因为这些性质传递到直和因子。由 Kaplansky 定理
（定理 \ref{algebra-theorem-projective-direct-sum}），$M$ 满足 (3)。

\medskip\noindent
反过来，设 $M$ 满足 (1)--(3)。由于平坦性和 Mittag-Leffler 性传递到
直和因子，$M$ 是平坦、Mittag-Leffler、可数生成 $R$-模的直和。
引理 \ref{algebra-lemma-countgen-projective}
表明 $M$ 是投射模的直和。因此 $M$ 是投射的。
\end{proof}

\begin{lemma}
\label{algebra-lemma-ML-ui-descent}
设 $f: M \to N$ 是 $R$-模的普遍单射映射。假设
$M$ 是可数生成 $R$-模的直和，并假设 $N$ 平坦且是 Mittag-Leffler 的。
则 $M$ 是投射的。
\end{lemma}

\begin{proof}
由引理 \ref{algebra-lemma-ui-flat-domain} 和
\ref{algebra-lemma-pure-submodule-ML}，
$M$ 平坦且是 Mittag-Leffler 的，故结论由定理
\ref{algebra-theorem-projectivity-characterization} 得出。
\end{proof}

\begin{lemma}
\label{algebra-lemma-universally-injective-submodule-powerseries}
设 $R$ 是 Noether 环，$M$ 是 $R$-模。假设 $M$ 是可数生成
$R$-模的直和，并假设对某个 $n$ 存在普遍单射映射
$M \to R[[t_1, \ldots, t_n]]$。则 $M$ 是投射的。
\end{lemma}

\begin{proof}
这由引理 \ref{algebra-lemma-ML-ui-descent} 和
\ref{algebra-lemma-power-series-ML} 得出。
\end{proof}

\section{模的性质的上升}
\label{algebra-section-ascent}

\noindent
定理 \ref{algebra-theorem-projectivity-characterization}
中模的所有性质都沿任意环映射上升：

\begin{lemma}
\label{algebra-lemma-ascend-properties-modules}
设 $R \to S$ 是环映射。设 $M$ 是 $R$-模。则：
\begin{enumerate}
\item 若 $M$ 平坦，则 $S$-模 $M \otimes_R S$ 平坦。
\item 若 $M$ 是 Mittag-Leffler 的，则 $S$-模 $M \otimes_R S$
是 Mittag-Leffler 的。
\item 若 $M$ 是可数生成 $R$-模的直和，则 $S$-模
$M \otimes_R S$ 是可数生成 $S$-模的直和。
\item 若 $M$ 投射，则 $S$-模 $M \otimes_R S$ 投射。
\end{enumerate}
\end{lemma}

\begin{proof}
除 (2) 外都很显然。对 (2)，使用命题
\ref{algebra-proposition-ML-characterization} 中 Mittag-Leffler 性的表述 (3)，
以及张量运算与取余极限可交换这一事实。也可以使用命题
\ref{algebra-proposition-ML-tensor} 的刻画。
\end{proof}

\section{模的性质的下降}
\label{algebra-section-descent}

\noindent
我们讨论定理 \ref{algebra-theorem-projectivity-characterization}
中刻画投射性的诸性质的忠实平坦下降。
在平坦性成立时，Mittag-Leffler 模这一性质可以下降：

\begin{lemma}
\label{algebra-lemma-ffdescent-ML}
\begin{reference}
Juan Pablo Acosta Lopez 于 2014 年 12 月 20 日发送的电子邮件。
\end{reference}
设 $R \to S$ 是忠实平坦环映射。设 $M$ 是 $R$-模。若
$S$-模 $M \otimes_R S$ 是 Mittag-Leffler 的，则 $M$
是 Mittag-Leffler 的。
\end{lemma}

\begin{proof}
把 $M = \colim_{i\in I} M_i$ 写成有限表示
$R$-模 $M_i$ 的有向余极限。
由命题 \ref{algebra-proposition-ML-characterization} 可知，我们必须证明：
对每个 $i \in I$，存在 $i \leq j$、$j\in I$，使得
$M_i\rightarrow M_j$ 支配 $M_i\rightarrow M$。

\medskip\noindent
取 $N$ 为推出
$$
\xymatrix{
M_i  \ar[r] \ar[d] & M_j \ar[d] \\
M \ar[r] & N
}
$$
于是由引理 \ref{algebra-lemma-domination-universally-injective}，
该引理等价于存在 $j$，使 $M_j\rightarrow N$ 普遍单射。
注意，与 $S$ 张量所得的图表
$$
\xymatrix{
M_i\otimes_R S  \ar[r] \ar[d] & M_j\otimes_R S \ar[d] \\
M\otimes_R S \ar[r] & N\otimes_R S
}
$$
是推出图表。由于
$M \otimes_R S = \colim_{i\in I} M_i \otimes_R S$
把 $M\otimes_R S$ 表为有限表示 $S$-模的余极限，并且
$M\otimes_R S$ 是 Mittag-Leffler 的，所以存在 $j \geq i$，
使 $M_j\otimes_R S\rightarrow N\otimes_R S$ 普遍单射。
再利用 $R\rightarrow S$ 忠实平坦，即知
$M_j\rightarrow N$ 也普遍单射。
\end{proof}

\begin{lemma}
\label{algebra-lemma-ffdescent-countable}
设 $R \to S$ 是忠实平坦环映射。设 $M$ 是 $R$-模。若
$S$-模 $M \otimes_R S$ 可数生成，则 $M$ 可数生成。
\end{lemma}

\begin{proof}
设 $M \otimes_R S$ 由诸元素
$y_i$（$i = 1, 2, 3, \ldots$）生成。写成
$y_i = \sum_{j = 1, \ldots, n_i}
x_{ij} \otimes s_{ij}$，其中 $n_i \geq 0$、$x_{ij} \in M$
且 $s_{ij} \in S$。记 $M' \subset M$ 为由可数个元素 $x_{ij}$
生成的子模。于是 $M' \otimes_R S \to M \otimes_R S$ 是满射，
因为它的像包含诸生成元 $y_i$。由于 $S$ 在 $R$ 上忠实平坦，
可得 $M' = M$，如所欲。
\end{proof}

\noindent
至此，可数生成投射模的忠实平坦下降很容易得出。

\begin{lemma}
\label{algebra-lemma-ffdescent-countable-projectivity}
设 $R \to S$ 是忠实平坦环映射。设 $M$ 是 $R$-模。
若 $S$-模 $M \otimes_R S$ 可数生成且投射，则 $M$
可数生成且投射。
\end{lemma}

\begin{proof}
这由引理 \ref{algebra-lemma-descend-properties-modules}、
\ref{algebra-lemma-ffdescent-ML}、\ref{algebra-lemma-ffdescent-countable}
和定理 \ref{algebra-theorem-projectivity-characterization} 得出。
\end{proof}

\noindent
剩下的只是利用 d\'evissage，把一般情形下投射性的下降约化到
可数生成情形。先给出两个简单引理。

\begin{lemma}
\label{algebra-lemma-lift-countably-generated-submodule}
设 $R \to S$ 是环映射，$M$ 是 $R$-模，并设 $Q$ 是
$M \otimes_R S$ 的可数生成 $S$-子模。则存在 $M$ 的可数生成
$R$-子模 $P$，使得
$\Im(P \otimes_R S \to M \otimes_R S)$ 包含 $Q$。
\end{lemma}

\begin{proof}
设 $y_1, y_2, \ldots$ 是 $Q$ 的生成元，并写成
$y_j = \sum_k x_{jk} \otimes s_{jk}$，其中 $x_{jk} \in M$
且 $s_{jk} \in S$。取 $P$ 为 $M$ 中由诸 $x_{jk}$ 生成的子模。
\end{proof}

\begin{lemma}
\label{algebra-lemma-adapted-submodule}
设 $R \to S$ 是环映射，$M$ 是 $R$-模。假设
$M
\otimes_R S = \bigoplus_{i \in I} Q_i$ 是可数生成 $S$-模 $Q_i$
的直和。若 $N$ 是 $M$ 的可数生成子模，则存在 $M$ 的可数生成子模
$N'$，使得 $N' \supset N$，且对某个子集 $I' \subset I$，
有 $\Im(N' \otimes_R S \to M \otimes_R S) =
\bigoplus_{i \in I'} Q_i$。
\end{lemma}

\begin{proof}
令 $N'_0 = N$。归纳构造可数生成子模的递增列
$N'_{\ell} \subset M$（$\ell = 0, 1, 2, \ldots$），满足：
若 $I'_{\ell}$ 是这样的 $i \in I$ 的集合，使得
$\Im(N'_{\ell} \otimes_R S \to M \otimes_R S)$
在 $Q_i$ 上的投影非零，则
$\Im(N'_{\ell + 1} \otimes_R S \to M \otimes_R
S)$ 对所有 $i \in I'_{\ell}$ 都包含 $Q_i$。为从
$N'_\ell$ 构造 $N'_{\ell + 1}$，令 $Q$ 为
$i \in I'_{\ell}$ 时诸（可数多个）$Q_i$ 之和，按引理
\ref{algebra-lemma-lift-countably-generated-submodule} 选取 $P$，再令
$N'_{\ell + 1} = N'_{\ell} + P$。构造出诸 $N'_{\ell}$ 后，只需取
$N' = \bigcup_{\ell} N'_{\ell}$ 和 $I' = \bigcup_{\ell} I'_{\ell}$。
\end{proof}

\begin{theorem}
\label{algebra-theorem-ffdescent-projectivity}
设 $R \to S$ 是忠实平坦环映射。设 $M$ 是 $R$-模。
若 $S$-模 $M \otimes_R S$ 投射，则 $M$ 投射。
\end{theorem}

\begin{proof}
我们将构造 $M$ 的 Kaplansky d\'evissage，以证明它是投射模的直和，
因而投射。由定理 \ref{algebra-theorem-projective-direct-sum}，可以把
$M \otimes_R S =
\bigoplus_{i \in I} Q_i$ 写成可数生成 $S$-模 $Q_i$ 的直和。
在 $M$ 上选取一个良序。我们将利用超限递归定义 $M$ 的子模递增族
$M_{\alpha}$，每个序数 $\alpha$ 对应一个子模，使得
$M_{\alpha} \otimes_R S$ 是某些 $Q_i$ 的直和。

\medskip\noindent
当 $\alpha = 0$ 时，令 $M_0 = 0$。若 $\alpha$ 是极限序数，并且
对所有 $\beta < \alpha$ 已定义 $M_{\beta}$，则定义
$M_\alpha = \bigcup_{\beta < \alpha} M_{\beta}$。由于每个
$\beta < \alpha$ 所对应的 $M_{\beta} \otimes_R S$
都是某些 $Q_i$ 的直和，$M_{\alpha} \otimes_R S$ 也如此。
若 $\alpha + 1$ 是后继序数且 $M_{\alpha}$ 已定义，则如下定义
$M_{\alpha + 1}$。若 $M_{\alpha} = M$，则令
$M_{\alpha +1} = M$。否则选取满足 $x \notin
M_{\alpha}$ 的最小 $x \in M$（相对于固定的良序）。由于 $S$
在 $R$ 上平坦，$(M/M_{\alpha}) \otimes_R S = M
\otimes_R S/M_{\alpha} \otimes_R S$；又因为
$M_{\alpha} \otimes_R S$ 是某些 $Q_i$ 的直和，
$(M/M_{\alpha}) \otimes_R S$ 也如此。由引理
\ref{algebra-lemma-adapted-submodule}，可以在 $M/M_{\alpha}$ 中找到
可数生成 $R$-子模 $P$，它包含 $x$ 在 $M/M_{\alpha}$ 中的像，
并且 $P \otimes_R S$（因为 $S$ 在 $R$ 上平坦，它等于
% P01-E0471: the frozen source maps P tensor S into M tensor S although P is a submodule of M/M_alpha; see ERRATA_CANDIDATES.jsonl.
$\Im(P
\otimes_R S \to M \otimes_R S)$）是某些 $Q_i$ 的直和。由于
$M \otimes_R S = \bigoplus_{i \in I} Q_i$ 投射，且投射性传递到
直和因子，$P \otimes_R S$ 也投射。因此，由引理
\ref{algebra-lemma-ffdescent-countable-projectivity}，$P$ 投射。最后，把
$M_{\alpha + 1}$ 定义为 $P$ 在 $M$ 中的逆像，于是
$M_{\alpha + 1}/M_{\alpha} = P$ 可数生成且投射。特别地，
$M_{\alpha}$ 是 $M_{\alpha + 1}$ 的直和因子，因为
$M_{\alpha + 1}/M_{\alpha}$ 的投射性意味着正合列
$0 \to M_{\alpha} \to M_{\alpha + 1} \to
M_{\alpha + 1}/M_{\alpha} \to 0$ 分裂。

\medskip\noindent
对 $M$ 作超限归纳（利用如下事实：在构造 $M_{\alpha + 1}$ 时，
我们使其包含不属于 $M_{\alpha}$ 的最小元素 $x \in M$），
可知每个 $x \in M$ 都属于某个 $M_{\alpha}$。因此，存在足够大的序数
% P01-E0473: the frozen source reuses S, already the faithfully flat target ring, as an ordinal; see ERRATA_CANDIDATES.jsonl.
$S$，满足：对每个 $x \in M$，存在 $\alpha \in S$，使得
$x \in M_{\alpha}$。这意味着 $(M_{\alpha})_{\alpha \in S}$
满足 $M$ 的 Kaplansky d\'evissage 的性质 (1)。其余性质由构造即明。
我们得到
$M = \bigoplus_{\alpha + 1 \in S} M_{\alpha + 1}/M_{\alpha}$。
由于按构造每个 $M_{\alpha + 1}/M_{\alpha}$ 都投射，$M$ 投射。
\end{proof}

\section{完备化}
\label{algebra-section-completion}

\noindent
设 $R$ 是环，$I$ 是理想。定义{\it $R$ 关于 $I$ 的完备化}为极限
$$
R^\wedge = \lim_n R/I^n.
$$
$R^\wedge$ 的一个元素由一列元素 $f_n \in R/I^n$ 给出，
它们对所有 $n$ 满足 $f_n \equiv f_{n + 1} \bmod I^n$。
我们把 $R^\wedge$ 看作 $R$-代数。类似地，若 $M$ 是 $R$-模，
则定义{\it $M$ 关于 $I$ 的完备化}为极限
$$
M^\wedge = \lim_n M/I^nM.
$$
$M^\wedge$ 的一个元素由一列元素 $m_n \in M/I^nM$ 给出，
它们对所有 $n$ 满足 $m_n \equiv m_{n + 1} \bmod I^nM$。
我们把 $M^\wedge$ 看作 $R^\wedge$-模。由这一描述可立即看出，
总有典范映射
$$
M \longrightarrow M^\wedge
\quad\text{且}\quad
M \otimes_R R^\wedge \longrightarrow M^\wedge.
$$
此外，给定模映射 $\varphi : M \to N$，会在完备化上诱导映射
$\varphi^\wedge : M^\wedge \to N^\wedge$，使图表
$$
\xymatrix{
M \ar[r] \ar[d] & N \ar[d] \\
M^\wedge \ar[r] & N^\wedge
}
$$
交换。一般而言，完备化不是正合函子，见
《例》，第 \ref{examples-section-completion-not-exact} 节。
下面给出一些初步的肯定结果。

\begin{lemma}
\label{algebra-lemma-completion-generalities}
设 $R$ 是环。设 $I \subset R$ 是理想。
设 $\varphi : M \to N$ 是 $R$-模映射。
\begin{enumerate}
\item 若 $M/IM \to N/IN$ 满射，则 $M^\wedge \to N^\wedge$ 满射。
\item 若 $M \to N$ 满射，则 $M^\wedge \to N^\wedge$ 满射。
\item 若 $0 \to K \to M \to N \to 0$ 是 $R$-模的短正合列且
$N$ 平坦，则
$0 \to K^\wedge \to M^\wedge \to N^\wedge \to 0$ 是短正合列。
\item 对任意有限 $R$-模 $M$，映射
$M \otimes_R R^\wedge \to M^\wedge$ 都满射。
\end{enumerate}
\end{lemma}

\begin{proof}
假设 $M/IM \to N/IN$ 满射。由 Nakayama 引理，对每个
$n \geq 1$，映射 $M/I^nM \to N/I^nN$ 都满射。更确切地说，
在环 $R/I^n$ 及其幂零理想 $I/I^n$ 上，对映射
$M/I^nM \to N/I^nN$ 应用引理 \ref{algebra-lemma-NAK} 的 (11)
即可看出这一点。置
$K_n = \{x \in M \mid \varphi(x) \in I^nN\}$。于是得到短正合列
$$
0 \to K_n/I^nM \to M/I^nM \to N/I^nN \to 0
$$
我们断言典范映射 $K_{n + 1}/I^{n + 1}M \to K_n/I^nM$
是满射。事实上，若 $x \in K_n$，写成
$\varphi(x) = \sum z_j n_j$，其中 $z_j \in I^n$、$n_j \in N$。
由假设，可写成 $n_j = \varphi(m_j) + \sum z_{jk}n_{jk}$，
其中 $m_j \in M$、$z_{jk} \in I$ 且 $n_{jk} \in N$。因此
$$
\varphi(x - \sum z_j m_j) = \sum z_jz_{jk} n_{jk}.
$$
这意味着 $x' = x - \sum z_j m_j \in K_{n + 1}$ 映到
$x \bmod I^nM$，从而证明了断言。现在可对上述短正合列逆系应用
引理 \ref{algebra-lemma-Mittag-Leffler}，得出 (1)。
(2) 是 (1) 的特殊情形。
若 (3) 的假设成立，则对每个 $n$，序列
$$
0 \to K/I^nK \to M/I^nM \to N/I^nN \to 0
$$
由引理 \ref{algebra-lemma-flat-tor-zero} 是短正合的。
因此可以直接应用引理 \ref{algebra-lemma-Mittag-Leffler}，得出 (3)。
为证 (4)，选取生成元 $x_i \in M$，$i = 1, \ldots, n$。
于是映射 $R^{\oplus n} \to M$，
$(a_1, \ldots, a_n) \mapsto \sum a_ix_i$ 满射。因此由 (2)，
映射 $(R^\wedge)^{\oplus n} \to M^\wedge$，
$(a_1, \ldots, a_n) \mapsto \sum a_ix_i$ 满射。(4) 由此得出。
\end{proof}

\begin{definition}
\label{algebra-definition-complete}
设 $R$ 是环。设 $I \subset R$ 是理想。
设 $M$ 是 $R$-模。若映射
$$
M \longrightarrow M^\wedge = \lim_n M/I^nM
$$
是同构\footnote{这包括条件 $\bigcap I^nM = 0$。}，则称 $M$
是{\it $I$-进完备的}。若 $R$ 作为 $R$-模是 $I$-进完备的，
则称 $R$ 是{\it $I$-进完备的}。
\end{definition}

\noindent
$R$-模 $M$ 关于 $I$ 的完备化未必是 $I$-进完备的。例子见
《例》，第 \ref{examples-section-noncomplete-completion} 节。
若理想有限生成，则完备化是完备的。

\begin{lemma}
\label{algebra-lemma-hathat-finitely-generated}
\begin{reference}
\cite[定理 15]{Matlis}。这里给出的简洁证明来自 Bjorn Poonen
于 2016 年 11 月 5 日发送的一封电子邮件。
\end{reference}
设 $R$ 是环。设 $I$ 是 $R$ 的有限生成理想。
设 $M$ 是 $R$-模。则
\begin{enumerate}
\item 完备化 $M^\wedge$ 是 $I$-进完备的，并且
\item 对所有 $n \geq 1$，
$I^nM^\wedge = \Ker(M^\wedge \to M/I^nM) = (I^nM)^\wedge$。
\end{enumerate}
特别地，$R^\wedge$ 是 $I$-进完备的，
$I^nR^\wedge = (I^n)^\wedge$，并且
$R^\wedge/I^nR^\wedge = R/I^n$。
\end{lemma}

\begin{proof}
由于 $I$ 有限生成，$I^n$ 也有限生成；设其生成元为
$f_1, \ldots, f_r$。把引理 \ref{algebra-lemma-completion-generalities} 的
(2) 应用于满射 $(f_1, \ldots, f_r) : M^{\oplus r} \to I^n M$，
得到满射
$$
(M^\wedge)^{\oplus r} \xrightarrow{(f_1, \ldots, f_r)} (I^n M)^\wedge =
\lim_{m \geq n} I^n M/I^m M = \Ker(M^\wedge \to M/I^n M).
$$
另一方面，映射
$(f_1, \ldots, f_r) : (M^\wedge)^{\oplus r} \to M^\wedge$
的像是 $I^n M^\wedge$。
因此 $M^\wedge / I^n M^\wedge \simeq M/I^n M$。
取逆极限得到 $(M^\wedge)^\wedge \simeq M^\wedge$；
也就是说，$M^\wedge$ 是 $I$-进完备的。
\end{proof}

\begin{lemma}
\label{algebra-lemma-completion-differ-by-torsion}
设 $R$ 是环。设 $I \subset R$ 是理想。设
$0 \to M \to N \to Q \to 0$ 是 $R$-模的正合列，
并且 $Q$ 被 $I$ 的某个幂零化。则完备化产生正合列
$0 \to M^\wedge \to N^\wedge \to Q \to 0$。
\end{lemma}

\begin{proof}
设 $I^c Q = 0$。则当 $n \geq c$ 时，$Q/I^nQ = Q$。
另一方面，当 $n \geq c$ 时，显然
$I^nM \subset M \cap I^nN \subset I^{n - c}M$。
因此 $M^\wedge = \lim M/(M \cap I^n N)$。对 $n \geq c$ 时的
正合列逆系
$$
0 \to M/(M \cap I^n N) \to N/I^n N \to Q \to 0
$$
应用引理 \ref{algebra-lemma-Mittag-Leffler} 即得结论。
\end{proof}

\begin{lemma}
\label{algebra-lemma-hathat}
\begin{reference}
取自 Lenstra 和 de Smit 的一份未发表笔记。
\end{reference}
设 $R$ 是环。设 $I \subset R$ 是理想。设 $M$ 是 $R$-模。
记 $K_n = \Ker(M^\wedge \to M/I^nM)$。则 $M^\wedge$
是 $I$-进完备的，当且仅当对所有 $n \geq 1$，
$K_n$ 等于 $I^nM^\wedge$。
\end{lemma}

\begin{proof}
模 $I^n M^\wedge$ 包含在 $K_n$ 中。
因此，对每个 $n \geq 1$，都有典范正合列
$$
0 \to K_n/I^nM^\wedge \to M^\wedge/I^nM^\wedge \to M/I^nM \to 0.
$$
由于 $I^nM^\wedge$ 满射到 $I^nM/I^{n + 1}M$，可知
$K_{n + 1} + I^n M^\wedge = K_n$。因此逆系
$\{K_n/I^n M^\wedge\}_{n \geq 1}$ 的转移映射都是满射。
由引理 \ref{algebra-lemma-Mittag-Leffler}，得到短正合列
$$
0 \to
\lim_n K_n/I^n M^\wedge \to
(M^\wedge)^\wedge \to
M^\wedge \to 0
$$
因此，$M^\wedge$ 完备，当且仅当对所有 $n \geq 1$，
$K_n/I^n M^\wedge = 0$。
\end{proof}

\begin{lemma}
\label{algebra-lemma-radical-completion}
设 $R$ 是环，$I \subset R$ 是理想，并令
$R^\wedge = \lim R/I^n$。
\begin{enumerate}
\item $R^\wedge$ 中任何映到 $R/I$ 的单位的元素都是单位；
\item $1 + I$ 中任何元素都映到 $R^\wedge$ 的可逆元素；
\item $1 + IR^\wedge$ 中任何元素在 $R^\wedge$ 中都可逆；并且
\item 理想 $IR^\wedge$ 和 $\Ker(R^\wedge \to R/I)$
都包含在 $R^\wedge$ 的 Jacobson 根中。
\end{enumerate}
\end{lemma}

\begin{proof}
设 $x \in R^\wedge$ 映到 $R/I$ 中的单位 $x_1$。
由引理 \ref{algebra-lemma-locally-nilpotent-unit}，对每个 $n$，
$x$ 都映到 $R/I^n$ 中的单位 $x_n$。
因此 $y = (x_n^{-1}) \in \lim R/I^n = R^\wedge$ 是 $x$ 的逆元。
(2) 和 (3) 立即由 (1) 得出。
(4) 由 (1) 和引理 \ref{algebra-lemma-contained-in-radical} 得出。
\end{proof}

\begin{lemma}
\label{algebra-lemma-when-surjective-to-completion}
设 $A$ 是环。设 $I = (f_1, \ldots, f_r)$ 是有限生成理想。
% P01-E0474: the frozen lemma uses M throughout without stating that M is an A-module; see ERRATA_CANDIDATES.jsonl.
若对每个 $i$，$M \to \lim M/f_i^nM$ 都满射，则
$M \to \lim M/I^nM$ 满射。
\end{lemma}

\begin{proof}
注意有等式：
$$
\lim M/I^nM = \lim M/(f_1^n, \ldots, f_r^n)M
$$
这是因为：
$$
I^n \supset (f_1^n, \ldots, f_r^n) \supset I^{rn}
$$
于是 $\lim M/(f_1^n, \ldots, f_r^n)M$ 的元素 $\xi$
可以符号化地写成
$$
\xi = \sum\nolimits_{n \geq 0} \sum\nolimits_i f_i^n x_{n, i}
$$
其中 $x_{n, i} \in M$。若 $M \to \lim M/f_i^nM$ 满射，
则存在 $x_i \in M$，它在 $\lim M/f_i^nM$ 中映到
$\sum x_{n, i} f_i^n$。于是 $x = \sum x_i$ 在
$\lim M/I^nM$ 中映到 $\xi$。
\end{proof}

\begin{lemma}
\label{algebra-lemma-complete-by-sub}
设 $A$ 是环。设 $I \subset J \subset A$ 是理想。
% P01-E0475: the frozen lemma uses M without stating that M is an A-module; see ERRATA_CANDIDATES.jsonl.
若 $M$ 是 $J$-进完备的且 $I$ 有限生成，则 $M$ 是 $I$-进完备的。
\end{lemma}

\begin{proof}
假设 $M$ 是 $J$-进完备的且 $I$ 有限生成。
因为 $\bigcap J^nM = 0$，所以 $\bigcap I^nM = 0$。由引理
\ref{algebra-lemma-when-surjective-to-completion}，只须在 $I$
由单个元素生成时证明 $M \to \lim M/I^nM$ 满射。设 $I = (f)$。
令 $x_n \in M$ 满足 $x_{n + 1} - x_n \in f^nM$。我们必须证明，
存在 $x \in M$，使得对所有 $n$，都有 $x_n - x \in f^nM$。
由于 $x_{n + 1} - x_n \in J^nM$，且 $M$ 是 $J$-进完备的，
存在元素 $x \in M$，使 $x_n - x \in J^nM$。
用 $x_n - x$ 代替 $x_n$，可以假设 $x_n \in J^nM$。
为完成证明，我们将说明这蕴含 $x_n \in I^nM$。
事实上，写成 $x_n - x_{n + 1} = f^nz_n$。则
$$
x_n = f^n(z_n + fz_{n + 1} + f^2z_{n + 2} + \ldots)
$$
由于 $f^c \in J^c$，和
$z_n + fz_{n + 1} + f^2z_{n + 2} + \ldots$ 在 $M$ 中收敛。
和 $f^n(z_n + fz_{n + 1} + f^2z_{n + 2} + \ldots)$
在 $M$ 中收敛到 $x_n$，因为其部分和等于
$x_n - x_{n + c}$，而 $x_{n + c} \in J^{n + c}M$。
\end{proof}

\begin{lemma}
\label{algebra-lemma-change-ideal-completion}
设 $R$ 是环。
设 $I$、$J$ 是 $R$ 的理想。
假设存在整数 $c, d > 0$，使
$I^c \subset J$ 且 $J^d \subset I$。
则对任意 $R$-模，关于 $I$ 的完备化与关于 $J$ 的完备化一致。
特别地，$R$-模 $M$ 是 $I$-进完备的，当且仅当它是
$J$-进完备的。
\end{lemma}

\begin{proof}
考虑映射系
$M/I^nM \to M/J^{\lfloor n/c \rfloor}M$ 与映射系
$M/J^mM \to M/I^{\lfloor m/d \rfloor}M$，
得到两个完备化之间互为逆的映射。
\end{proof}

\begin{lemma}
\label{algebra-lemma-quotient-complete}
设 $R$ 是环。设 $I$ 是 $R$ 的理想。
设 $M$ 是 $I$-进完备 $R$-模，并设 $K \subset M$ 是 $R$-子模。
以下条件等价：
\begin{enumerate}
\item $K = \bigcap (K + I^nM)$；并且
\item $M/K$ 是 $I$-进完备的。
\end{enumerate}
\end{lemma}

\begin{proof}
置 $N = M/K$。由引理 \ref{algebra-lemma-completion-generalities}，
映射 $M = M^\wedge \to N^\wedge$ 满射。
因此 $N \to N^\wedge$ 满射。不难看出，映射
$N \to N^\wedge$ 的核是模 $\bigcap (K + I^nM) / K$。
\end{proof}

\begin{lemma}
\label{algebra-lemma-when-finite-module-complete-over-complete-ring}
设 $R$ 是环。设 $I$ 是 $R$ 的理想。
设 $M$ 是 $R$-模。
若 (a) $R$ 是 $I$-进完备的，(b) $M$ 是有限 $R$-模，
且 (c) $\bigcap I^nM = (0)$，则 $M$ 是 $I$-进完备的。
\end{lemma}

\begin{proof}
由引理 \ref{algebra-lemma-completion-generalities}，映射
$M = M \otimes_R R = M \otimes_R R^\wedge \to M^\wedge$
满射。该映射的核是 $\bigcap I^nM$，因而由假设为零。
所以 $M \cong M^\wedge$，且 $M$ 完备。
\end{proof}

\begin{lemma}
\label{algebra-lemma-finite-over-complete-ring}
设 $R$ 是环。设 $I \subset R$ 是理想。设 $M$ 是 $R$-模。
假设
\begin{enumerate}
\item $R$ 是 $I$-进完备的；
\item $\bigcap_{n \geq 1} I^nM = (0)$；并且
\item $M/IM$ 是有限 $R/I$-模。
\end{enumerate}
则 $M$ 是有限 $R$-模。
\end{lemma}

\begin{proof}
设 $x_1, \ldots, x_n \in M$ 是这样的元素：它们在 $M/IM$
中的像作为 $R/I$-模生成 $M/IM$。记 $M' \subset M$ 为由
$x_1, \ldots, x_n$ 生成的 $R$-子模。由引理
\ref{algebra-lemma-completion-generalities}，映射
$(M')^\wedge \to M^\wedge$ 满射。
由于 $\bigcap I^nM = 0$，特别有 $\bigcap I^nM' = (0)$。
因此由引理 \ref{algebra-lemma-when-finite-module-complete-over-complete-ring}，
$M'$ 完备，从而 $M' \to M^\wedge$ 满射。最后，
$M \to M^\wedge$ 的核为零，因为它等于 $\bigcap I^nM = (0)$。
所以 $M \cong M' \cong M^\wedge$ 有限生成。
\end{proof}

\section{Noether 环的完备化}
\label{algebra-section-completion-noetherian}

\noindent
本节讨论 Noether 环中关于理想的完备化。

\begin{lemma}
\label{algebra-lemma-completion-tensor}
设 $I$ 是 Noether 环 $R$ 的理想。
用 ${}^\wedge$ 表示关于 $I$ 的完备化。
\begin{enumerate}
\item 若 $K \to N$ 是有限 $R$-模之间的单射映射，
则完备化上的映射 $K^\wedge \to N^\wedge$ 是单射。
\item 若 $0 \to K \to N \to M \to 0$ 是有限 $R$-模的短正合列，
则 $0 \to K^\wedge \to N^\wedge \to M^\wedge \to 0$ 是短正合列。
\item 若 $M$ 是有限 $R$-模，则 $M^\wedge = M \otimes_R R^\wedge$。
\end{enumerate}
\end{lemma}

\begin{proof}
置 $M = N/K$，可知 (1) 由 (2) 得出。
设 $0 \to K \to N \to M \to 0$ 如 (2) 所述。
对每个 $n$，得到短正合列
$$
0 \to K/(I^nN \cap K) \to N/I^nN \to M/I^nM \to 0.
$$
由引理 \ref{algebra-lemma-Mittag-Leffler}，得到正合列
$$
0 \to \lim K/(I^nN \cap K) \to N^\wedge \to M^\wedge \to 0.
$$
由 Artin--Rees 引理 \ref{algebra-lemma-Artin-Rees}，可以选取 $c$，
使得当 $n \geq c$ 时，
$I^nK \subset I^n N \cap K \subset I^{n-c} K$。
因此 $K^\wedge = \lim K/I^nK = \lim K/(I^nN \cap K)$，
从而 (2) 成立。

\medskip\noindent
设 $M$ 如 (3) 所述，并设
$0 \to K \to R^{\oplus t} \to M \to 0$ 是 $M$ 的一个表示。
得到交换图表
$$
\xymatrix{
&
K \otimes_R R^\wedge \ar[r] \ar[d] &
R^{\oplus t} \otimes_R R^\wedge \ar[r] \ar[d] &
M \otimes_R R^\wedge \ar[r] \ar[d] &
0 \\
0 \ar[r] &
K^\wedge \ar[r] &
(R^{\oplus t})^\wedge \ar[r] &
M^\wedge \ar[r] & 0
}
$$
上行正合，见第 \ref{algebra-section-flat} 节。
下行由 (2) 正合。由引理
\ref{algebra-lemma-completion-generalities}，竖直箭头都满射。
中间的竖直箭头是同构。由蛇引理 \ref{algebra-lemma-snake} 得出 (3)。
\end{proof}

\begin{lemma}
\label{algebra-lemma-completion-flat}
设 $I$ 是 Noether 环 $R$ 的理想。
用 ${}^\wedge$ 表示关于 $I$ 的完备化。
\begin{enumerate}
\item 环映射 $R \to R^\wedge$ 平坦。
\item 函子 $M \mapsto M^\wedge$ 在有限生成 $R$-模范畴上正合。
\end{enumerate}
\end{lemma}

\begin{proof}
考虑 $J \otimes_R R^\wedge \to R \otimes_R R^\wedge = R^\wedge$，
其中 $J$ 是 $R$ 的任意理想。
依照引理 \ref{algebra-lemma-completion-tensor}，该映射等同于
$J^\wedge \to R^\wedge$，而 $J^\wedge \to R^\wedge$ 是单射。
(1) 由引理 \ref{algebra-lemma-flat} 得出。
(2) 是引理 \ref{algebra-lemma-completion-tensor} 的 (2) 的重新表述。
\end{proof}

\begin{lemma}
\label{algebra-lemma-completion-faithfully-flat}
设 $I$ 是 Noether 环 $R$ 的理想。用 $R^\wedge$ 表示
$R$ 关于 $I$ 的完备化。若 $I$ 包含在 $R$ 的 Jacobson 根中，
则环映射 $R \to R^\wedge$ 忠实平坦。特别地，若
$(R, \mathfrak m)$ 是 Noether 局部环，则完备化
$\lim_n R/\mathfrak m^n$ 忠实平坦。
\end{lemma}

\begin{proof}
由引理 \ref{algebra-lemma-completion-flat}，它平坦。
考虑复合 $R \to R^\wedge \to R/I$，其中后一个映射是投影映射
$R^\wedge \to R/I$；这表明 $R$ 的任意极大理想都属于
$\Spec(R^\wedge) \to \Spec(R)$ 的像。因此，由引理
\ref{algebra-lemma-ff}，该映射忠实平坦。
\end{proof}

\begin{lemma}
\label{algebra-lemma-completion-complete}
设 $R$ 是 Noether 环。
设 $I$ 是 $R$ 的理想。
设 $M$ 是 $R$-模。
则 $M$ 关于 $I$ 的完备化 $M^\wedge$ 是 $I$-进完备的，
$I^n M^\wedge = (I^nM)^\wedge$，并且
$M^\wedge/I^nM^\wedge = M/I^nM$。
\end{lemma}

\begin{proof}
这是引理 \ref{algebra-lemma-hathat-finitely-generated} 的特殊情形，
因为 $I$ 是有限生成理想。
\end{proof}

\begin{lemma}
\label{algebra-lemma-completion-Noetherian}
设 $I$ 是环 $R$ 的理想。假设
\begin{enumerate}
\item $R/I$ 是 Noether 环；
\item $I$ 有限生成。
\end{enumerate}
则 $R$ 关于 $I$ 的完备化 $R^\wedge$ 是关于
$IR^\wedge$ 完备的 Noether 环。
\end{lemma}

\begin{proof}
由引理 \ref{algebra-lemma-hathat-finitely-generated} 可知，$R^\wedge$
是 $I$-进完备的。因此它也是 $IR^\wedge$-进完备的。
由于 $R^\wedge/IR^\wedge = R/I$ 是 Noether 的，把 $R$
替换为 $R^\wedge$ 后，除假设 (1) 和 (2) 外，还可以假设
$R$ 是 $I$-进完备的。

\medskip\noindent
设 $f_1, \ldots, f_t$ 是 $I$ 的生成元。
则有环的满射
$R/I[T_1, \ldots, T_t] \to \bigoplus I^n/I^{n + 1}$，
它把 $T_i$ 映到元素 $\overline{f}_i \in I/I^2$。
因此 $\bigoplus I^n/I^{n + 1}$ 是 Noether 环。
设 $J \subset R$ 是理想。考虑理想
$$
\bigoplus J \cap I^n/J \cap I^{n + 1} \subset \bigoplus I^n/I^{n + 1}.
$$
设 $\overline{g}_1, \ldots, \overline{g}_m$ 是该理想的生成元。
可以选取 $\overline{g}_j$ 为次数 $d_j$ 的齐次元素，并可选取
$g_j \in J \cap I^{d_j}$，使其映到
$\overline{g}_j \in J \cap I^{d_j}/J \cap I^{d_j + 1}$。
我们断言 $g_1, \ldots, g_m$ 生成 $J$。

\medskip\noindent
设 $x \in J \cap I^n$。存在
$a_j \in I^{\max(0, n - d_j)}$，使
$x - \sum a_j g_j \in J \cap I^{n + 1}$。
其理由是，按 $g_1, \ldots, g_m$ 的选取，
$J \cap I^n/J \cap I^{n + 1}$ 等于
% P01-E0480: the frozen source writes I^(n-d_j) for every j without restricting to d_j <= n, leaving negative ideal powers undefined; see ERRATA_CANDIDATES.jsonl.
$\sum \overline{g}_j I^{n - d_j}/I^{n - d_j + 1}$。
因此，从 $x \in J$ 出发，可以找到一列向量
$(a_{1, n}, \ldots, a_{m, n})_{n \geq 0}$，其中
$a_{j, n} \in I^{\max(0, n - d_j)}$，使得
$$
x =
\sum\nolimits_{n = 0, \ldots, N}
\sum\nolimits_{j = 1, \ldots, m} a_{j, n} g_j \bmod I^{N + 1}
$$
置 $A_j = \sum_{n \geq 0} a_{j, n}$。由于 $R$ 完备，可知
$x = \sum A_j g_j$。所以 $J$ 有限生成，结论得证。
\end{proof}

\begin{lemma}
\label{algebra-lemma-completion-Noetherian-Noetherian}
设 $R$ 是 Noether 环。
设 $I$ 是 $R$ 的理想。
$R$ 关于 $I$ 的完备化 $R^\wedge$ 是 Noether 的。
\end{lemma}

\begin{proof}
这是引理 \ref{algebra-lemma-completion-Noetherian} 的推论。
也可以如下直接看出。选取 $I$ 的生成元
$f_1, \ldots, f_n$。考虑映射
$$
R[[x_1, \ldots, x_n]] \longrightarrow R^\wedge,
\quad
x_i \longmapsto f_i.
$$
这是定义良好的满射环映射（细节略）。
由于 $R[[x_1, \ldots, x_n]]$ 是 Noether 的（见引理
\ref{algebra-lemma-Noetherian-power-series}），结论得证。
\end{proof}

\noindent
设 $R \to S$ 是局部环 $(R, \mathfrak m)$ 与
$(S, \mathfrak n)$ 之间的局部同态。令 $S^\wedge$ 为
$S$ 关于 $\mathfrak n$ 的完备化。一般而言，$S^\wedge$
不是 $S$ 的 $\mathfrak m$-进完备化。若对某个 $t \geq 1$，
有 $\mathfrak n^t \subset \mathfrak mS$，则由引理
\ref{algebra-lemma-change-ideal-completion}，确有
$S^\wedge = \lim S/\mathfrak m^nS$。在某些情形下，这甚至蕴含
$S^\wedge$ 在 $R^\wedge$ 上有限。

\begin{lemma}
\label{algebra-lemma-finite-after-completion}
设 $R \to S$ 是局部环 $(R, \mathfrak m)$ 与
$(S, \mathfrak n)$ 之间的局部同态。分别令
$R^\wedge$、$S^\wedge$ 为 $R$、$S$ 关于
$\mathfrak m$、$\mathfrak n$ 的完备化。若
$\mathfrak m$ 和 $\mathfrak n$ 有限生成，且
$\dim_{\kappa(\mathfrak m)} S/\mathfrak mS < \infty$，则
\begin{enumerate}
\item $S^\wedge$ 等于 $S$ 的 $\mathfrak m$-进完备化；并且
\item $S^\wedge$ 是有限 $R^\wedge$-模。
\end{enumerate}
\end{lemma}

\begin{proof}
因为 $R \to S$ 是局部环映射，所以
$\mathfrak mS \subset \mathfrak n$。
假设 $\dim_{\kappa(\mathfrak m)} S/\mathfrak mS < \infty$
蕴含 $S/\mathfrak mS$ 是 Artin 环，见引理
\ref{algebra-lemma-finite-dimensional-algebra}。
因此它的维数为 $0$，见引理
\ref{algebra-lemma-Noetherian-dimension-0}，从而
$\mathfrak n = \sqrt{\mathfrak mS}$。
这与 $\mathfrak n$ 有限生成这一事实蕴含：对某个
$t \geq 1$，有 $\mathfrak n^t \subset \mathfrak mS$。
由引理 \ref{algebra-lemma-change-ideal-completion}，可把 $S^\wedge$
等同于 $S$ 的 $\mathfrak m$-进完备化。由于 $\mathfrak m$
有限生成，由引理 \ref{algebra-lemma-hathat-finitely-generated} 可知，
$S^\wedge$ 与 $R^\wedge$ 都是 $\mathfrak m$-进完备的。
此时可以把引理 \ref{algebra-lemma-finite-over-complete-ring}
应用于作为 $R^\wedge$-模的 $S^\wedge$，得出结论。
\end{proof}

\begin{lemma}
\label{algebra-lemma-completion-finite-extension}
设 $R$ 是 Noether 环，$R \to S$ 是有限环映射。
设 $\mathfrak p \subset R$ 是素理想，并设
$\mathfrak q_1, \ldots, \mathfrak q_m$ 是 $S$ 中位于
$\mathfrak p$ 上方的素理想
（引理 \ref{algebra-lemma-finite-finite-fibres}）。则
$$
R_\mathfrak p^\wedge \otimes_R S =
(S_\mathfrak p)^\wedge =
S_{\mathfrak q_1}^\wedge \times \ldots \times S_{\mathfrak q_m}^\wedge
$$
其中 $(S_\mathfrak p)^\wedge$ 是关于 $\mathfrak p$ 的完备化，
而局部环 $R_\mathfrak p$ 和 $S_{\mathfrak q_i}$
都关于各自的极大理想完备化。
\end{lemma}

\begin{proof}
可以把 $R$ 替换为局部化 $R_\mathfrak p$，并把 $S$ 替换为
$S_\mathfrak p = S \otimes_R R_\mathfrak p$。
因此可以假设 $R$ 是 Noether 局部环，并且
$\mathfrak p = \mathfrak m$ 是其极大理想。
由引理 \ref{algebra-lemma-finite-after-completion}，
$\mathfrak q_iS_{\mathfrak q_i}$-进完备化
$S_{\mathfrak q_i}^\wedge$ 等于 $\mathfrak m$-进完备化。
对每个 $n \geq 1$，$S/\mathfrak m^nS$ 的素理想与
$S$ 的极大理想 $\mathfrak q_1, \ldots, \mathfrak q_m$
一一对应（由 $S$ 在 $R$ 上的上升定理，见引理
\ref{algebra-lemma-integral-going-up}）。因此，由引理
\ref{algebra-lemma-artinian-finite-length}，
$$
S/\mathfrak m^nS = \prod S_{\mathfrak q_i}/\mathfrak m^nS_{\mathfrak q_i}
$$
（例如使用命题 \ref{algebra-proposition-dimension-zero-ring}
可知 $S/\mathfrak m^nS$ 是 Artin 的）。
所以 $S$ 的 $\mathfrak m$-进完备化 $S^\wedge$ 等于
$\prod S_{\mathfrak q_i}^\wedge$。最后，由引理
\ref{algebra-lemma-completion-tensor}，
$R^\wedge \otimes_R S = S^\wedge$。
\end{proof}

\begin{lemma}
\label{algebra-lemma-split-completed-sequence}
设 $R$ 是环，$I \subset R$ 是理想。
设 $0 \to K \to P \to M \to 0$ 是 $R$-模的短正合列。
若 $M$ 在 $R$ 上平坦，且 $M/IM$ 是投射 $R/I$-模，
则 $I$-进完备化序列
$$
0 \to K^\wedge \to P^\wedge \to M^\wedge \to 0
$$
是分裂正合列。
\end{lemma}

\begin{proof}
由于 $M$ 平坦，每个序列
$$
0 \to K/I^nK \to P/I^nP \to M/I^nM \to 0
$$
都是短正合的，见引理 \ref{algebra-lemma-flat-tor-zero}；
而序列 $0 \to K^\wedge \to P^\wedge \to M^\wedge \to 0$
是短正合的，见引理 \ref{algebra-lemma-completion-generalities}。
只须证明，可以找到分裂
$s_n : M/I^nM \to P/I^nP$，使
$s_{n + 1} \bmod I^n = s_n$。
我们对 $n$ 归纳构造这些 $s_n$。
任取一个分裂 $s_1$；它存在，因为 $M/IM$ 是投射 $R/I$-模。
假设对某个 $n > 0$ 已给定 $s_n$。置
$P_{n + 1} = \{x \in P \mid x \bmod I^nP \in \Im(s_n)\}$。
映射
$\pi : P_{n + 1}/I^{n + 1}P_{n + 1} \to M/I^{n + 1}M$
满射（细节略）。由引理 \ref{algebra-lemma-lift-projective}，
$M/I^{n + 1}M$ 作为 $R/I^{n + 1}$-模是投射的，因而可以选取
$\pi$ 的截面
$t : M/I^{n + 1}M \to P_{n + 1}/I^{n + 1}P_{n + 1}$。
令 $s_{n + 1}$ 等于 $t$ 与典范映射
$P_{n + 1}/I^{n + 1}P_{n + 1} \to P/I^{n + 1}P$
的复合即可。
\end{proof}

\begin{lemma}
\label{algebra-lemma-complete-modulo-nilpotent}
设 $A$ 是 Noether 环。设 $I, J \subset A$ 是理想。
若 $A$ 是 $I$-进完备的且 $A/I$ 是 $J$-进完备的，
则 $A$ 是 $J$-进完备的。
\end{lemma}

\begin{proof}
设 $B$ 是 $A$ 的 $(I + J)$-进完备化。由引理
\ref{algebra-lemma-completion-flat}，$B/IB$ 是 $A/I$ 的 $J$-进完备化，
故由假设同构于 $A/I$。此外，由引理
\ref{algebra-lemma-complete-by-sub}，$B$ 是 $I$-进完备的。因此，由引理
\ref{algebra-lemma-finite-over-complete-ring}，$B$ 是有限 $A$-模。
由 Nakayama 引理（引理 \ref{algebra-lemma-NAK}；这里使用引理
\ref{algebra-lemma-radical-completion} 所给出的 $I$ 包含在 $A$ 的
Jacobson 根中这一事实），可知 $A \to B$ 满射。
由引理 \ref{algebra-lemma-completion-flat}，映射 $A \to B$ 平坦。
$\Spec(B) \to \Spec(A)$ 的像包含 $V(I)$；又由于 $I$
包含在 $A$ 的 Jacobson 根中，由引理
\ref{algebra-lemma-ff-rings} 可知 $A \to B$ 忠实平坦。因此
$A \to B$ 单射。于是 $A$ 关于 $I + J$ 完备，因而更是关于
$J$ 完备的。
\end{proof}

\section{模的极限}
\label{algebra-section-limits}

\noindent
本节讨论对模取极限时会发生什么。

\begin{lemma}
\label{algebra-lemma-limit-complete-pre}
设 $I \subset A$ 是一个环的有限生成理想。
设 $(M_n)$ 是满足 $I^n M_n = 0$ 的 $A$-模逆系统。
则 $M = \lim M_n$ 是 $I$-进完备的。
\end{lemma}

\begin{proof}
有 $M \to M/I^nM \to M_n$。取极限得到
$M \to M^\wedge \to M$。因此 $M$ 是 $M^\wedge$ 的直和因子。
由引理 \ref{algebra-lemma-hathat-finitely-generated}，$M^\wedge$
是 $I$-进完备的，故 $M$ 也是如此。
\end{proof}

\begin{lemma}
\label{algebra-lemma-limit-complete}
设 $I \subset A$ 是一个环的有限生成理想。
设 $(M_n)$ 是满足
$M_n = M_{n + 1}/I^nM_{n + 1}$ 的 $A$-模逆系统。置 $M = \lim M_n$。
则 $M/I^nM = M_n$，且 $M$ 是 $I$-进完备的。
\end{lemma}

\begin{proof}
由引理 \ref{algebra-lemma-limit-complete-pre} 可知 $M$ 是 $I$-进完备的。
由于转移映射都是满射，映射 $M \to M_n$ 也是满射。
考虑定义 $N_n$ 的短正合列逆系统
$$
0 \to N_n \to M \to M_n \to 0
$$
由于 $M_n = M_{n + 1}/I^nM_{n + 1}$，映射
$N_{n + 1} + I^nM \to N_n$ 是满射。因此
$N_{n + 1}/(N_{n + 1} \cap I^{n + 1}M) \to N_n/(N_n \cap I^nM)$
是满射。对短正合列
$$
0 \to N_n/(N_n \cap I^nM) \to M/I^nM \to M_n \to 0
$$
取逆极限，得到正合列
$$
0 \to \lim N_n/(N_n \cap I^nM) \to M^\wedge \to M
$$
由于 $M$ 是 $I$-进完备的，可知
$\lim N_n/(N_n \cap I^nM) = 0$；再由转移映射的满射性，
对所有 $n$ 都有 $N_n/(N_n \cap I^nM) = 0$。
故 $M_n = M/I^nM$，如所求。
\end{proof}

\begin{lemma}
\label{algebra-lemma-finiteness-graded}
设 $A$ 是 Noether 分次环。设 $I \subset A_+$ 是齐次理想。
设 $(N_n)$ 是有限分次 $A$-模的逆系统，并满足
$N_n = N_{n + 1}/I^n N_{n + 1}$。则存在有限分次 $A$-模
$N$，使得对所有 $n$，作为分次模都有 $N_n = N/I^nN$。
\end{lemma}

\begin{proof}
选取 $r$ 以及次数分别为 $d_1, \ldots, d_r$、
生成 $N_1$ 的齐次元素 $x_{1, 1}, \ldots, x_{1, r} \in N_1$。
由于转移映射都是满射，可以选取提升 $x_{1, i}$ 的相容齐次元素系
$x_{n, i} \in N_n$。由分次 Nakayama 引理
（引理 \ref{algebra-lemma-graded-NAK}）可知，$N_n$ 由次数分别为
$d_1, \ldots, d_r$ 的元素 $x_{n, 1}, \ldots, x_{n, r}$ 生成。
因此，当 $m \leq n$ 时，映射 $N_n \to N_n/I^m N_n$
在次数 $< \min(d_i) + m$ 上是同构（因为 $I^mN_n$
在这些次数上为零）。从而次数为 $d$ 的各部分所成的逆系统
% P01-E0483: frozen source stabilization indices are retained; see part-local ledger.
$$
\ldots
= N_{2 + d - \min(d_i), d}
= N_{1 + d - \min(d_i), d}
= N_{d - \min(d_i), d} \to N_{-1 + d - \min(d_i), d} \to \ldots
$$
按所示方式稳定。令 $N$ 为如下分次 $A$-模：其 $d$ 次齐次部分
就是这个稳定值。特别地，在 $N$ 中有元素
$x_i = \lim x_{n, i}$。我们断言诸 $x_i$ 生成 $N$：
任意 $x \in N_d$ 都是 $x_1, \ldots, x_r$ 的线性组合，
因为可以在 $N_{d - \min(d_i), d}$ 中检验这一点；在那里结论成立，
这是由于诸 $x_{d - \min(d_i), i}$ 生成
$N_{d - \min(d_i)}$。最后，读者可逐次检查，验证满射
$N/I^nN \to N_n$ 是同构。细节从略。
\end{proof}

\begin{lemma}
\label{algebra-lemma-daniel-litt}
设 $A$ 是分次环。设 $I \subset A_+$ 是齐次理想。
记 $A' = \lim A/I^n$。设 $(G_n)$ 是分次 $A$-模的逆系统，
其中 $G_n$ 被 $I^n$ 零化。设 $M$ 是分次 $A$-模，并设
$\varphi_n : M \to G_n$ 是相容的分次 $A$-模映射系。若诱导映射
$$
\varphi : M \otimes_A A' \longrightarrow \lim G_n
$$
是同构，则对每个 $d \in \mathbf{Z}$，映射
$M_d \to \lim G_{n, d}$ 都是同构。
\end{lemma}

\begin{proof}
按约定，分次环集中在次数 $\geq 0$ 上，而分次模可以在任意次数上
具有非零部分，见第 \ref{algebra-section-graded} 节。映射 $\varphi$ 存在，
因为 $G_n$ 被 $I^n$ 零化，所以 $\lim G_n$ 是 $A'$-模。
另一个值得记住的事实是
$$
\bigoplus\nolimits_{d \in \mathbf{Z}} \lim G_{n, d} \subset
\lim G_n \subset
\prod\nolimits_{d \in \mathbf{Z}} \lim G_{n, d}
$$
其中下标 ${\ }_d$ 表示 $d$ 次齐次部分。

\medskip\noindent
单射性。设 $x \in M_d$。若在 $\lim G_{n, d}$ 中有 $x \mapsto 0$，
则在 $M \otimes_A A'$ 中有 $x \otimes 1 = 0$。于是可以找到
有限生成子模 $M' \subset M$，且 $x \in M'$，使得
$x \otimes 1$ 在 $M' \otimes_A A'$ 中为零。
设 $M'$ 由次数为 $d_1, \ldots, d_r$ 的齐次元素生成。
令 $n = d - \min(d_i) + 1$。由于 $A'$ 有到 $A/I^n$ 的映射，
且 $A \to A/I^n$ 在次数 $\leq n - 1$ 上是同构，故可知
% P01-E0484: frozen source degree bound is retained; see part-local ledger.
$M' \to M' \otimes_A A'$ 在次数 $\leq n - 1$ 上是单射。
因此 $x = 0$，如所求。

\medskip\noindent
满射性。设 $y \in \lim G_{n, d}$。在 $M \otimes_A A'$ 中选取
一个映到 $y$ 的有限和 $\sum x_i \otimes f'_i$。
可以假设 $x_i$ 是齐次的，设其次数为 $d_i$。
注意，虽然 $A'$ 不是分次环，但它是分次环 $A/I^nA$ 的极限；
而且在任一给定次数上，转移映射最终成为同构（见上文）。这给出
$$
A = \bigoplus\nolimits_{d \geq 0} A_d \subset A' \subset
\prod\nolimits_{d \geq 0} A_d
$$
因此可以写成
$$
f'_i = \sum\nolimits_{j = 0, \ldots, d - d_i - 1} f_{i, j} + f_i + g'_i
$$
其中 $f_{i, j} \in A_j$、$f_i \in A_{d - d_i}$，且
$g'_i \in A'$ 在 $\prod_{j \leq d - d_i} A_j$ 中的像为零。
现在计算 $\varphi_n(\sum_{i, j} f_{i, j}x_i) \in G_n$，
得到的是次数 $< d$ 的齐次元素之和。因此
$\varphi(\sum x_i \otimes f_{i, j})$ 在 $\lim G_{n, d}$ 中的像为零。
类似地，计算表明元素 $\varphi(\sum x_i \otimes g'_i)$
在 $\prod_{d' \leq d} \lim G_{n, d'}$ 中的像为零。
由于已知 $\varphi(\sum x_i \otimes f'_i)$ 等于 $y$，
可知 $\sum f_ix_i \in M_d$ 映到 $y$，如所求。
\end{proof}

\section{平坦性判据}
\label{algebra-section-criteria-flatness}

\noindent
本节证明 Noether 情形下的一些重要技术引理。我们将在第
\ref{algebra-section-more-flatness-criteria} 节把它们（部分地）推广到
非 Noether 情形。

\begin{lemma}
\label{algebra-lemma-mod-injective}
设 $R \to S$ 是局部环的局部同态，且 $S$ 是 Noether 的。
% P01-E0485: frozen source omits "by" in three notation declarations; see ledger.
记 $\mathfrak m$ 为 $R$ 的极大理想。设 $M$ 是平坦 $R$-模，
$N$ 是有限 $S$-模。设 $u : N \to M$ 是 $R$-模映射。
若 $\overline{u} : N/\mathfrak m N \to M/\mathfrak m M$ 是单射，
则 $u$ 是单射。在这种情形下，$M/u(N)$ 在 $R$ 上平坦。
\end{lemma}

\begin{proof}
首先断言，对所有 $n \geq 1$，
$u_n : N/{\mathfrak m}^nN \to M/{\mathfrak m}^nM$ 都是单射。
我们作归纳；基础情形即 $\overline{u} = u_1$ 是单射。
由 $M$ 在 $R$ 上平坦这一假设，有短正合列
$0 \to M \otimes_R {\mathfrak m}^n/{\mathfrak m}^{n + 1}
\to M/{\mathfrak m}^{n + 1}M \to M/{\mathfrak m}^n M \to 0$。
此外，
$M \otimes_R {\mathfrak m}^n/{\mathfrak m}^{n + 1}
= M/{\mathfrak m}M \otimes_{R/{\mathfrak m}}
{\mathfrak m}^n/{\mathfrak m}^{n + 1}$。对 $N$ 也有类似的正合列
$N \otimes_R {\mathfrak m}^n/{\mathfrak m}^{n + 1}
\to N/{\mathfrak m}^{n + 1}N \to N/{\mathfrak m}^n N \to 0$，
只是左端没有零。还有
$N \otimes_R {\mathfrak m}^n/{\mathfrak m}^{n + 1}
= N/{\mathfrak m}N \otimes_{R/{\mathfrak m}}
{\mathfrak m}^n/{\mathfrak m}^{n + 1}$。因此，
$u_n$ 与映射
$\overline{u} \otimes \text{id}_{{\mathfrak m}^n/{\mathfrak m}^{n + 1}}$
都是单射，从而 $u_{n + 1}$ 也是单射。

\medskip\noindent
把 Krull 交定理（引理 \ref{algebra-lemma-intersect-powers-ideal-module-zero}）
应用于环 $S$ 上的 $N$ 及理想 $\mathfrak mS$，得到
$\bigcap \mathfrak m^nN = 0$。因此，对所有 $n$，$u_n$ 的单射性
蕴含 $u$ 是单射。

\medskip\noindent
为证明 $M/u(N)$ 在 $R$ 上平坦，只须证明对每个理想
$I \subset R$ 都有 $\text{Tor}_1^R(M/u(N), R/I) = 0$，
见引理 \ref{algebra-lemma-characterize-flat}。由短正合列
$$
0 \to N \xrightarrow{u} M \to M/u(N) \to 0
$$
以及 $M$ 的平坦性，得到 Tor 的正合列
$$
0 \to \text{Tor}_1^R(M/u(N), R/I) \to N/IN \to M/IM
$$
见引理 \ref{algebra-lemma-long-exact-sequence-tor}。因此，只须证明
$N/IN$ 单射到 $M/IM$ 中。注意，$R/I \to S/IS$ 是局部环的
局部同态，且 $S/IS$ 是 Noether 的；$N/IN \to M/IM$ 是
$R/I$-模映射；$N/IN$ 在 $S/IS$ 上有限；$M/IM$ 在 $R/I$ 上
平坦；而 $u \bmod I : N/IN \to M/IM$ 模 $\mathfrak m$ 是单射。
于是可以把证明的第一部分应用于 $u \bmod I$，结论即得。
\end{proof}

\begin{lemma}
\label{algebra-lemma-grothendieck}
设 $R \to S$ 是 Noether 局部环之间的平坦局部环同态。
记 $\mathfrak m$ 为 $R$ 的极大理想。设 $f \in S$ 在
$S/{\mathfrak m}S$ 中是非零因子。则 $S/fS$ 在 $R$ 上平坦，
且 $f$ 在 $S$ 中是非零因子。
\end{lemma}

\begin{proof}
直接由引理 \ref{algebra-lemma-mod-injective} 得出。
\end{proof}

\begin{lemma}
\label{algebra-lemma-grothendieck-regular-sequence}
设 $R \to S$ 是 Noether 局部环之间的平坦局部环同态。
记 $\mathfrak m$ 为 $R$ 的极大理想。设 $f_1, \ldots, f_c$
是 $S$ 中的元素序列，使其像
$\overline{f}_1, \ldots, \overline{f}_c$ 在
$S/{\mathfrak m}S$ 中构成正则序列。则 $f_1, \ldots, f_c$
在 $S$ 中构成正则序列，且每个商 $S/(f_1, \ldots, f_i)$
都在 $R$ 上平坦。
\end{lemma}

\begin{proof}
由归纳法和引理 \ref{algebra-lemma-grothendieck} 得出。
\end{proof}

\begin{lemma}
\label{algebra-lemma-free-fibre-flat-free}
设 $R \to S$ 是 Noether 局部环之间的局部同态。
设 $\mathfrak m$ 是 $R$ 的极大理想。设 $M$ 是非零有限 $S$-模。
假设 (a) $M/\mathfrak mM$ 是自由 $S/\mathfrak mS$-模，且
(b) $M$ 在 $R$ 上平坦。则 $M$ 是自由的，且 $S$ 在 $R$ 上平坦。
\end{lemma}

\begin{proof}
设 $\overline{x}_1, \ldots, \overline{x}_n$ 是自由模
$M/\mathfrak mM$ 的一组基。选取 $x_1, \ldots, x_n \in M$，
使 $x_i$ 映到 $\overline{x}_i$。令
$u : S^{\oplus n} \to M$ 为把第 $i$ 个标准基向量映到 $x_i$ 的映射。
由引理 \ref{algebra-lemma-mod-injective}，$u$ 是单射。另一方面，
由 Nakayama 引理 \ref{algebra-lemma-NAK}，该映射是满射。结论随即得出。
\end{proof}

\begin{lemma}
\label{algebra-lemma-complex-exact-mod}
设 $R \to S$ 是局部 Noether 环之间的局部同态。设
$\mathfrak m$ 是 $R$ 的极大理想。设
$0 \to F_e \to F_{e-1} \to \ldots \to F_0$
是有限 $S$-模的有限复形。假设每个 $F_i$ 都在 $R$ 上平坦，
且复形
$0 \to F_e/\mathfrak m F_e \to F_{e-1}/\mathfrak m F_{e-1}
\to \ldots \to F_0 / \mathfrak m F_0$ 正合。则
$0 \to F_e \to F_{e-1} \to \ldots \to F_0$ 正合，
而且模 $\Coker(F_1 \to F_0)$ 在 $R$ 上平坦。
\end{lemma}

\begin{proof}
对 $e$ 作归纳。若 $e = 1$，这正是引理
\ref{algebra-lemma-mod-injective}。若 $e > 1$，则由引理
\ref{algebra-lemma-mod-injective} 可知 $F_e \to F_{e-1}$ 是单射，且
$C = \Coker(F_e \to F_{e-1})$ 是在 $R$ 上平坦的有限 $S$-模。
因此可以把归纳假设应用于复形
$0 \to C \to F_{e-2} \to \ldots \to F_0$。
由此得出 $C \to F_{e-2}$ 是单射，进而得到该复形的正合性，
以及 $F_1 \to F_0$ 的余核的平坦性。
\end{proof}

\noindent
本节余下部分证明所谓“{\it 平坦性的局部判据}”的两个版本。
还请注意下文颇有意思的引理 \ref{algebra-lemma-CM-over-regular-flat}。

\begin{lemma}
\label{algebra-lemma-prepare-local-criterion-flatness}
设 $R$ 是极大理想为 $\mathfrak m$、剩余域为
$\kappa = R/\mathfrak m$ 的局部环。设 $M$ 是 $R$-模。
若 $\text{Tor}_1^R(\kappa, M) = 0$，则对每个有限长度
$R$-模 $N$，都有 $\text{Tor}_1^R(N, M) = 0$。
\end{lemma}

\begin{proof}
% P01-E0486: frozen source says descending induction although the proof is ordinary strong induction.
对 $N$ 的长度作归纳。若 $N$ 的长度为 $1$，则
$N \cong \kappa$，结论成立。若 $N$ 的长度大于 $1$，
则可将 $N$ 置于短正合列
$0 \to N' \to N \to N'' \to 0$ 中，其中 $N'$、$N''$
都是长度更小的有限长度 $R$-模。由
$\text{Tor}_1^R(N', M)$ 和 $\text{Tor}_1^R(N'', M)$ 的消没
（归纳假设），再结合 Tor 群的长正合列，可得
$\text{Tor}_1^R(N, M)$ 消没，见引理
\ref{algebra-lemma-long-exact-sequence-tor}。
\end{proof}

\begin{lemma}[平坦性的局部判据]
\label{algebra-lemma-local-criterion-flatness}
设 $R \to S$ 是局部 Noether 环之间的局部同态。
设 $\mathfrak m$ 是 $R$ 的极大理想，并设
$\kappa = R/\mathfrak m$。设 $M$ 是有限 $S$-模。
若 $\text{Tor}_1^R(\kappa, M) = 0$，则 $M$ 在 $R$ 上平坦。
\end{lemma}

\begin{proof}
设 $I \subset R$ 是理想。由引理 \ref{algebra-lemma-flat}，只须证明
$I \otimes_R M \to M$ 是单射。由注记
\ref{algebra-remark-Tor-ring-mod-ideal}，这个核等于
$\text{Tor}_1^R(M, R/I)$。由引理
\ref{algebra-lemma-prepare-local-criterion-flatness} 可知，
% P01-E0487: frozen source uses J without quantifying it; see ledger.
对所有有限余长度的理想，$J \otimes_R M \to M$ 都是单射。

\medskip\noindent
选取 $n >> 0$，并考虑短正合列
$$
0
\to I \cap \mathfrak m^n
\to I \oplus \mathfrak m^n
\to I + \mathfrak m^n
\to 0
$$
这是短正合列 $0 \to R \to R^{\oplus 2} \to R \to 0$ 的子序列。
因此得到交换图
$$
\xymatrix{
(I\cap \mathfrak m^n) \otimes_R M \ar[r] \ar[d] &
I \otimes_R M \oplus \mathfrak m^n \otimes_R M \ar[r] \ar[d] &
(I + \mathfrak m^n) \otimes_R M \ar[d] \\
M \ar[r] &
M \oplus M \ar[r] &
M
}
$$
注意，$I + \mathfrak m^n$ 和 $\mathfrak m^n$ 都是有限余长度的理想。
因此，追图表明
$\Ker((I \cap \mathfrak m^n)\otimes_R M \to M)
\to \Ker(I \otimes_R M \to M)$ 是满射。特别地，由此可知
$K = \Ker(I \otimes_R M \to M)$ 包含在
$(I \cap \mathfrak m^n) \otimes_R M$ 于 $I \otimes_R M$ 中的像内。
由 Artin--Rees 引理 \ref{algebra-lemma-Artin-Rees}，存在某个 $c > 0$，
使得对所有 $n >> 0$，$K$ 包含在
$\mathfrak m^{n-c}(I \otimes_R M)$ 中。由于
$I \otimes_R M$ 是有限 $S$-模（！），而 $S$ 是 Noether 的，
可知这蕴含 $K = 0$。具体而言，上述结论蕴含 $K$ 在
$I \otimes_R M$ 的 $\mathfrak mS$-进完备化中的像为零。
但由引理 \ref{algebra-lemma-completion-faithfully-flat}，从 $S$ 到其
$\mathfrak mS$-进完备化的映射是忠实平坦的。因此 $K = 0$，如所求。
\end{proof}

\noindent
下文经常遇到条件“$M/IM$ 在 $R/I$ 上平坦且
$\text{Tor}_1^R(R/I, M) = 0$”。下面的引理给出这些条件的一些推论
（它是引理 \ref{algebra-lemma-prepare-local-criterion-flatness} 的推广）。

\begin{lemma}
\label{algebra-lemma-what-does-it-mean}
设 $R$ 是环。设 $I \subset R$ 是理想。设 $M$ 是 $R$-模。
若 $M/IM$ 在 $R/I$ 上平坦，且
$\text{Tor}_1^R(R/I, M) = 0$，则
\begin{enumerate}
\item 对所有 $n \geq 1$，$M/I^nM$ 在 $R/I^n$ 上平坦；并且
\item 对任意被某个 $m \geq 0$ 的 $I^m$ 零化的模 $N$，都有
$\text{Tor}_1^R(N, M) = 0$。
\end{enumerate}
特别地，若 $I$ 是幂零理想，则 $M$ 在 $R$ 上平坦。
\end{lemma}

\begin{proof}
假设 $M/IM$ 在 $R/I$ 上平坦且 $\text{Tor}_1^R(R/I, M) = 0$。
设 $N$ 是 $R/I$-模。选取短正合列
% P01-E0488: frozen source incorrectly uses the ideal I as an arbitrary free-module index set.
$$
0 \to K \to \bigoplus\nolimits_{i \in I} R/I \to N \to 0
$$
由 $\text{Tor}$ 的长正合列和 $\text{Tor}_1^R(R/I, M)$ 的消没，得到
$$
0 \to \text{Tor}_1^R(N, M) \to K \otimes_R M \to
(\bigoplus\nolimits_{i \in I} R/I) \otimes_R M \to N \otimes_R M \to 0
$$
但由于 $K$、$\bigoplus_{i \in I} R/I$ 和 $N$ 都被 $I$ 零化，
可知
\begin{align*}
K \otimes_R M & = K \otimes_{R/I} M/IM, \\
(\bigoplus\nolimits_{i \in I} R/I) \otimes_R M & =
(\bigoplus\nolimits_{i \in I} R/I) \otimes_{R/I} M/IM, \\
N \otimes_R M & = N \otimes_{R/I} M/IM.
\end{align*}
因为 $M/IM$ 在 $R/I$ 上平坦，所以
% P01-E0489: frozen source has malformed tensor base R/; retained exactly.
$$
0 \to K \otimes_{R/I} M/IM \to
(\bigoplus\nolimits_{i \in I} R/I) \otimes_{R/I} M/IM \to
N \otimes_{R/} M/IM \to 0
$$
正合。结合上文，可知对任意被 $I$ 零化的 $R$-模 $N$，都有
$\text{Tor}_1^R(N, M) = 0$。

\medskip\noindent
下面对 $m$ 归纳证明 (2)。$m = 1$ 的情形已在上一段证明。
若 $N$ 被 $I^m$ 零化且 $m > 1$，可以选取正合列
$0 \to N' \to N \to N'' \to 0$，其中 $N'$ 和 $N''$
都被 $I^{m - 1}$ 零化。例如可以取 $N' = IN$ 和 $N'' = N/IN$。
此时，正合列
$$
\text{Tor}_1^R(N', M) \to
\text{Tor}_1^R(N, M) \to
\text{Tor}_1^R(N'', M)
$$
与归纳假设证明了所需的消没。

\medskip\noindent
最后证明 (1)。给定 $n \geq 1$，须证明 $M/I^nM$ 在 $R/I^n$ 上平坦。
换言之，须证明函子
$N \mapsto N \otimes_{R/I^n} M/I^nM$ 在由被 $I^n$ 零化的
$R$-模 $N$ 构成的范畴上正合。但对这样的 $N$，有
$N \otimes_{R/I^n} M/I^nM = N \otimes_R M$。
由 (2) 中 $\text{Tor}_1$ 的消没可知，函子
$N \mapsto N \otimes_R M$ 在被 $I$ 的某个幂零化的 $N$
所构成的范畴上正合，结论即得。
\end{proof}

\begin{lemma}
\label{algebra-lemma-what-does-it-mean-again}
设 $R$ 是环。设 $I \subset R$ 是理想。设 $M$ 是 $R$-模。
\begin{enumerate}
\item 若 $M/IM$ 在 $R/I$ 上平坦，且
$M \otimes_R I/I^2 \to IM/I^2M$ 是单射，则
$M/I^2M$ 在 $R/I^2$ 上平坦。
\item 若 $M/IM$ 在 $R/I$ 上平坦，且对
$n = 1, \ldots, k$，映射
$M \otimes_R I^n/I^{n + 1} \to I^nM/I^{n + 1}M$ 都是单射，
则 $M/I^{k + 1}M$ 在 $R/I^{k + 1}$ 上平坦。
\end{enumerate}
\end{lemma}

\begin{proof}
把引理 \ref{algebra-lemma-what-does-it-mean} 应用于以 $R/I^2$ 代替 $R$、
以 $M/I^2M$ 代替 $M$ 的情形，并使用
$$
\text{Tor}_1^{R/I^2}(M/I^2M, R/I) =
\Ker(M \otimes_R I/I^2 \to IM/I^2M),
$$
即可得到第一条，见注记 \ref{algebra-remark-Tor-ring-mod-ideal}。
第二条以同样方式得出：对 $n$ 作归纳，证明当
$n = 1, \ldots, k$ 时，$M/I^{n + 1}M$ 在 $R/I^{n + 1}$ 上平坦。
这里使用对每个 $n$ 都成立的等式
$$
\text{Tor}_1^{R/I^{n + 1}}(M/I^{n + 1}M, R/I^n) =
\Ker(M \otimes_R I^n/I^{n + 1} \to I^nM/I^{n + 1}M)
$$
\end{proof}

\begin{lemma}[局部判据的变式]
\label{algebra-lemma-variant-local-criterion-flatness}
设 $R \to S$ 是 Noether 局部环之间的局部同态。
设 $I \not = R$ 是 $R$ 中的理想。设 $M$ 是有限 $S$-模。
若 $\text{Tor}_1^R(M, R/I) = 0$，且 $M/IM$ 在 $R/I$ 上平坦，
则 $M$ 在 $R$ 上平坦。
\end{lemma}

\begin{proof}
第一种证明：由引理 \ref{algebra-lemma-what-does-it-mean}，
$\text{Tor}_1^R(\kappa, M)$ 为零，其中 $\kappa$ 是 $R$ 的剩余域。
于是由引理 \ref{algebra-lemma-local-criterion-flatness}，$M$ 在 $R$ 上平坦。

\medskip\noindent
第二种证明：设 $\mathfrak m$ 是 $R$ 的极大理想。我们将证明
$\mathfrak m \otimes_R M \to M$ 是单射，然后应用引理
\ref{algebra-lemma-local-criterion-flatness}。设
$\sum f_i \otimes x_i \in \mathfrak m \otimes_R M$，且在 $M$ 中
$\sum f_i x_i = 0$。把平坦性的等式判据——引理
\ref{algebra-lemma-flat-eq}——应用于 $R/I$ 上的 $M/IM$，可知存在
$\overline{a}_{ij} \in R/I$ 和 $\overline{y}_j \in M/IM$，使得
$x_i \bmod IM = \sum_j \overline{a}_{ij} \overline{y}_j $，且
$0 = \sum_i (f_i \bmod I) \overline{a}_{ij}$。
设 $a_{ij} \in R$ 是 $\overline{a}_{ij}$ 的一个提升；同样地，
设 $y_j \in M$ 是 $\overline{y}_j$ 的一个提升。于是
\begin{eqnarray*}
\sum f_i \otimes x_i
& = &
\sum f_i \otimes x_i +
\sum f_ia_{ij} \otimes y_j -
\sum f_i \otimes a_{ij} y_j
\\
& = &
\sum f_i \otimes (x_i - \sum a_{ij} y_j) +
\sum (\sum f_i a_{ij}) \otimes y_j
\end{eqnarray*}
由于 $x_i - \sum a_{ij} y_j \in IM$，且
$\sum f_i a_{ij} \in I$，可知存在 $I \otimes_R M$ 中的元素，
其在 $\mathfrak m \otimes_R M$ 中的像等于给定元素
$\sum f_i \otimes x_i$。但由假设，$I \otimes_R M \to M$ 是单射
（见注记 \ref{algebra-remark-Tor-ring-mod-ideal}），故结论成立。
\end{proof}

\noindent
特别地，在引理 \ref{algebra-lemma-variant-local-criterion-flatness} 的情形下，
假设 $I = (x)$ 由单个元素 $x$ 生成，且该元素在 $R$ 中是非零因子。
那么，$\text{Tor}_1^R(M, R/(x)) = (0)$ 当且仅当 $x$ 在 $M$ 上是非零因子。

\begin{lemma}
\label{algebra-lemma-flat-module-powers}
设 $R \to S$ 是环映射。设 $I \subset R$ 是理想。
设 $M$ 是 $S$-模。假设
\begin{enumerate}
\item $R$ 是 Noether 环；
\item $S$ 是 Noether 环；
\item $M$ 是有限 $S$-模；并且
\item 对每个 $n \geq 1$，模 $M/I^n M$ 在 $R/I^n$ 上平坦。
\end{enumerate}
则对每个 $\mathfrak q \in V(IS)$，局部化 $M_{\mathfrak q}$
在 $R$ 上平坦。特别地，若 $S$ 是局部环且 $IS$ 包含在其极大理想中，
则 $M$ 在 $R$ 上平坦。
\end{lemma}

\begin{proof}
我们将使用引理 \ref{algebra-lemma-variant-local-criterion-flatness}。
由假设，$M/IM$ 在 $R/I$ 上平坦。因此，只须检验
$\text{Tor}_1^R(M, R/I)$ 在 $\mathfrak q$ 处局部化后为零。
由注记 \ref{algebra-remark-Tor-ring-mod-ideal}，这个 Tor 群等于
$K = \Ker(I \otimes_R M \to M)$。已知对所有 $n \geq 1$，映射
$I/I^n \otimes_{R/I^n} M/I^nM \to M/I^nM$ 的核为零。
因此，$K$ 中任意元素在 $I/I^n \otimes_{R/I^n} M/I^nM$ 中的像为零。
由于
$$
I/I^n \otimes_{R/I^n} M/I^nM = I/I^n \otimes_R M =
(I \otimes_R M)/I^{n - 1}(I \otimes_R M)
$$
可知对所有 $n \geq 1$，都有
$K \subset I^{n - 1}(I \otimes_R M)$。由 Artin--Rees 引理，
更确切地说，由引理 \ref{algebra-lemma-intersection-powers-ideal-module}，
可知 $K_{\mathfrak q} = 0$，如所求。
\end{proof}

\begin{lemma}
\label{algebra-lemma-surjective-on-tor-one}
设 $R \to R' \to R''$ 是环映射。设 $M$ 是 $R$-模。
假设 $M \otimes_R R'$ 在 $R'$ 上平坦。则自然映射
$\text{Tor}_1^R(M, R') \otimes_{R'} R'' \to
\text{Tor}_1^R(M, R'')$ 是满射。
\end{lemma}

\begin{proof}
设 $F_\bullet$ 是 $M$ 在 $R$ 上的自由分解。复形
$F_2 \otimes_R R' \to F_1\otimes_R R' \to F_0 \otimes_R R'$
计算 $\text{Tor}_1^R(M, R')$，而复形
$F_2 \otimes_R R'' \to F_1\otimes_R R'' \to F_0 \otimes_R R''$
计算 $\text{Tor}_1^R(M, R'')$。注意，
$F_i \otimes_R R' \otimes_{R'} R'' = F_i \otimes_R R''$。
令 $K' = \Ker(F_1\otimes_R R' \to F_0 \otimes_R R')$，并类似地令
$K'' = \Ker(F_1\otimes_R R'' \to F_0 \otimes_R R'')$。
于是有正合列
$$
0 \to K' \to F_1\otimes_R R' \to F_0 \otimes_R R' \to M \otimes_R R' \to 0.
$$
由 $M \otimes_R R'$ 在 $R'$ 上平坦这一假设，序列
$$
K' \otimes_{R'} R'' \to
F_1 \otimes_R R'' \to
F_0 \otimes_R R'' \to
M \otimes_R R'' \to 0
$$
仍然正合。这意味着 $K' \otimes_{R'} R'' \to K''$ 是满射。
由于 $\text{Tor}_1^R(M, R')$ 是 $K'$ 的商，而
$\text{Tor}_1^R(M, R'')$ 是 $K''$ 的商，结论即得。
\end{proof}

\begin{lemma}
\label{algebra-lemma-surjective-on-tor-one-trivial}
设 $R \to R'$ 是环映射。设 $I \subset R$ 是理想，并令
$I' = IR'$。设 $M$ 是 $R$-模，并置 $M' = M \otimes_R R'$。
自然映射
$\text{Tor}_1^R(R'/I', M) \to \text{Tor}_1^{R'}(R'/I', M')$
是满射。
\end{lemma}

\begin{proof}
设 $F_2 \to F_1 \to F_0 \to M \to 0$ 是 $M$ 在 $R$ 上的自由分解。
置 $F_i' = F_i \otimes_R R'$。序列
$F_2' \to F_1' \to F_0' \to M' \to 0$ 在 $F_1'$ 处可能不再正合。
因此，$M'$ 在 $R'$ 上的一个自由分解形如
$$
F_2' \oplus F_2'' \to F_1' \to F_0' \to M' \to 0
$$
其中 $F_2''$ 是适当的自由 $R'$-模。其次注意，
$F_i \otimes_R R'/I' = F_i' / IF_i' = F_i'/I'F_i'$。
所以复形 $F_2'/I'F_2' \to F_1'/I'F_1' \to F_0'/I'F_0'$
计算 $\text{Tor}_1^R(M, R'/I')$。另一方面，
$F_i' \otimes_{R'} R'/I' = F_i'/I'F_i'$，而 $F_2''$ 也类似。
因此，复形
$F_2'/I'F_2' \oplus F_2''/I'F_2'' \to F_1'/I'F_1' \to F_0'/I'F_0'$
计算 $\text{Tor}_1^{R'}(M', R'/I')$。由于复形间的竖直映射
$$
\xymatrix{
F_2'/I'F_2' \ar[r] \ar[d] &
F_1'/I'F_1' \ar[r] \ar[d] &
F_0'/I'F_0' \ar[d] \\
F_2'/I'F_2' \oplus F_2''/I'F_2'' \ar[r] &
F_1'/I'F_1' \ar[r] &
F_0'/I'F_0'
}
$$
% P01-E0490: frozen source says cohomology for these homological Tor complexes.
显然在同调上诱导满射，结论即得。
\end{proof}

\begin{lemma}
\label{algebra-lemma-another-variant-local-criterion-flatness}
设
$$
\xymatrix{
S \ar[r] & S' \\
R \ar[r] \ar[u] & R' \ar[u]
}
$$
是局部 Noether 环的局部同态所成的交换图。设 $I \subset R$
是正常理想。设 $M$ 是有限 $S$-模。
% P01-E0491: frozen source omits "by" in the notation declaration; see ledger.
记 $I' = IR'$ 且 $M' = M \otimes_S S'$。假设
\begin{enumerate}
\item $S'$ 是张量积 $S \otimes_R R'$ 的一个局部化；
\item $M/IM$ 在 $R/I$ 上平坦；
\item 映射
$\text{Tor}_1^R(M, R/I) \to \text{Tor}_1^{R'}(M', R'/I')$ 为零。
\end{enumerate}
则 $M'$ 在 $R'$ 上平坦。
\end{lemma}

\begin{proof}
由于 $S'$ 是 $S \otimes_R R'$ 的局部化，$M'$ 是
$M \otimes_R R'$ 的局部化。注意，由引理
\ref{algebra-lemma-flat-base-change}，模
$M/IM \otimes_{R/I} R'/I'
= M \otimes_R R' /I'(M \otimes_R R')$ 在 $R'/I'$ 上平坦。
因此，$M'/I'M'$ 作为平坦模的局部化，在 $R'/I'$ 上也是平坦的。
由引理 \ref{algebra-lemma-variant-local-criterion-flatness}，只须证明
$\text{Tor}_1^{R'}(M', R'/I')$ 为零。由于 $M'$ 是
$M \otimes_R R'$ 的局部化，最后一个假设蕴含只须证明
$\text{Tor}_1^R(M, R/I) \otimes_R R'
\to
\text{Tor}_1^{R'}(M \otimes_R R', R'/I')$
是满射。

\medskip\noindent
由引理 \ref{algebra-lemma-surjective-on-tor-one-trivial}，映射
$\text{Tor}_1^R(M, R'/I') \to \text{Tor}_1^{R'}(M \otimes_R R', R'/I')$
是满射。因此现在只须证明
$\text{Tor}_1^R(M, R/I) \otimes_R R'
\to
\text{Tor}_1^R(M, R'/I')$
是满射。这由引理 \ref{algebra-lemma-surjective-on-tor-one} 得出；
只需考察环映射 $R \to R/I \to R'/I'$ 和模 $M$。
\end{proof}

\noindent
请把下面的引理与引理
\ref{algebra-lemma-criterion-flatness-fibre-nilpotent}（幂零理想的情形）
及引理 \ref{algebra-lemma-criterion-flatness-fibre}（有限表示代数的情形）比较。

\begin{lemma}[按纤维的平坦性判据；Noether 情形]
\label{algebra-lemma-criterion-flatness-fibre-Noetherian}
设 $R$、$S$、$S'$ 是 Noether 局部环，且 $R \to S \to S'$
是局部环同态。设 $\mathfrak m \subset R$ 是极大理想。
设 $M$ 是 $S'$-模。假设
\begin{enumerate}
\item 模 $M$ 在 $S'$ 上有限。
\item 模 $M$ 非零。
\item 模 $M/\mathfrak m M$ 在 $S/\mathfrak m S$ 上平坦。
\item 模 $M$ 是平坦 $R$-模。
\end{enumerate}
则 $S$ 在 $R$ 上平坦，且 $M$ 是平坦 $S$-模。
\end{lemma}

\begin{proof}
置 $I = \mathfrak mS \subset S$。由 (3)，$M/IM$ 是平坦
$S/I$-模。由于 $\mathfrak m \otimes_R S' \to I \otimes_S S'$
是满射，$\mathfrak m \otimes_R M \to I \otimes_S M$ 也是满射。
考虑
$$
\mathfrak m \otimes_R M \to I \otimes_S M \to M.
$$
因为 $M$ 在 $R$ 上平坦，复合映射是单射，故两个箭头都是单射。
特别地，$\text{Tor}_1^S(S/I, M) = 0$，见注记
\ref{algebra-remark-Tor-ring-mod-ideal}。由引理
\ref{algebra-lemma-variant-local-criterion-flatness}，可知 $M$ 在 $S$ 上平坦。
注意，由 Nakayama 引理 \ref{algebra-lemma-NAK}，
$M/\mathfrak m_{S'}M$ 非零，故由引理 \ref{algebra-lemma-ff}，$M$
实际上在 $S$ 上忠实平坦（因为这迫使
$M/\mathfrak m_SM \not = 0$）。

\medskip\noindent
考虑正合列 $0 \to \mathfrak m \to R \to \kappa \to 0$。
它给出正合列
$0 \to \text{Tor}_1^R(\kappa, S) \to \mathfrak m \otimes_R S \to I \to 0$。
由于 $M$ 在 $S$ 上平坦，又得到正合列
$0 \to \text{Tor}_1^R(\kappa, S)\otimes_S M \to
\mathfrak m \otimes_R M \to I \otimes_S M \to 0$。
由上文，这蕴含 $\text{Tor}_1^R(\kappa, S)\otimes_S M = 0$。
由于 $M$ 在 $S$ 上忠实平坦，可知
$\text{Tor}_1^R(\kappa, S) = 0$；再由引理
\ref{algebra-lemma-local-criterion-flatness}，可知 $S$ 在 $R$ 上平坦。
\end{proof}

\begin{lemma}
\label{algebra-lemma-flatness-scallop}
设 $A$ 是环，$M$ 是 $A$-模，且 $f \in A$。若
\begin{enumerate}
\item $f$ 在 $A$ 和 $M$ 上都是非零因子；
\item $M_f$ 是平坦 $A_f$-模；并且
\item $M/fM$ 是平坦 $A/fA$-模，
\end{enumerate}
则 $M$ 是平坦 $A$-模。把“平坦”换成“忠实平坦”时结论仍成立。
\end{lemma}

\begin{proof}
由于 $0 \to A \to A \to A/fA \to 0$ 和
$0 \to M \to M \to M/fM \to 0$ 正合，可知对 $i = 1$ 和 $i = 2$，
都有 $\text{Tor}_i^A(M, A/fA) = 0$。由引理
\ref{algebra-lemma-what-does-it-mean}，对所有被 $f$ 零化的 $A$-模 $N$，
都有 $\text{Tor}_1^A(M, N) = 0$（这里使用 $M/fM$ 在 $A/fA$
上的平坦性）。给定被 $f$ 零化的 $A$-模 $N$，可以选取
$A$-模的短正合列 $0 \to N' \to F \to N \to 0$，其中 $F$
是若干份 $A/fA$ 的直和。由正合列
$$
\text{Tor}_2^A(M, F) \to
\text{Tor}_2^A(M, N) \to
\text{Tor}_1^A(M, N')
$$
可知，对所有被 $f$ 零化的 $A$-模 $N$，都有
$\text{Tor}_2^A(M, N) = 0$。其次，设 $K$ 是任意 $A$-模。
可以把映射 $f : K \to K$ 分解成两个短正合列
$$
0 \to K[f] \to K \to K' \to 0
\quad\text{且}\quad
0 \to K' \to K \to K/fK \to 0
$$
其中 $K' = K/K[f] \cong fK$。应用 Tor 的正合列，得到正合列
$$
\text{Tor}_1^A(K[f], M) \to
\text{Tor}_1^A(K, M) \to
\text{Tor}_1^A(K', M)
$$
以及
$$
\text{Tor}_2^A(K/fK, M) \to
\text{Tor}_1^A(K', M) \to
\text{Tor}_1^A(K, M)
$$
利用 $\text{Tor}_1^A(K[f], M)$ 和
$\text{Tor}_2^A(K/fK, M)$ 的消没，可知 $f : K \to K$
在 $\text{Tor}^A_1(M, K)$ 上诱导单射。换言之，
$\text{Tor}^A_1(M, K)$ 为零，当且仅当
$\text{Tor}^A_1(M, K) \otimes_A A_f$ 为零。
由引理 \ref{algebra-lemma-flat-base-change-tor}，有
$$
\text{Tor}^A_1(M, K) \otimes_A A_f = \text{Tor}^{A_f}_1(M_f, K_f)  = 0
$$
% P01-E0492: frozen source has this sentence fragment twice; Chinese supplies the predicate.
最后一个消没来自 $M_f$ 在 $A_f$ 上的平坦性
（引理 \ref{algebra-lemma-characterize-flat}）。可知对所有 $A$-模 $K$，
都有 $\text{Tor}_1^A(M, K) = 0$。因此，由引理
\ref{algebra-lemma-characterize-flat}，$M$ 在 $A$ 上平坦。

\medskip\noindent
忠实平坦模情形的论证从略。
\end{proof}

\begin{lemma}
\label{algebra-lemma-flatness-scallop-pre}
设 $A$ 是环，$M$ 是 $A$-模。设 $I = (f_1, \ldots, f_r)$
是 $A$ 中由 $r \geq 1$ 个元素生成的理想。若
\begin{enumerate}
\item 对 $i = 1, \ldots, r$，$M_{f_i}$ 是平坦 $A_{f_i}$-模；
\item $M/IM$ 是平坦 $A/I$-模；
\item 对 $i = 1, \ldots, r + 1$，有
$\text{Tor}_i^A(M, A/I) = 0$。
\end{enumerate}
则 $M$ 是平坦 $A$-模。把“平坦”换成“忠实平坦”时结论仍成立。
\end{lemma}

\begin{proof}
由引理 \ref{algebra-lemma-what-does-it-mean}，对所有被 $I$ 零化的
$A$-模 $K$，都有 $\text{Tor}_1^A(M, K) = 0$。
给定被 $I$ 零化的 $A$-模 $K$，可以选取 $A$-模的短正合列
$0 \to N \to F \to K \to 0$，其中 $F$ 是若干份 $A/I$ 的直和。
得到正合列
$$
\text{Tor}_i^A(M, F) \to
\text{Tor}_i^A(M, K) \to
\text{Tor}_{i - 1}^A(M, N)
$$
因此，利用假设 (3) 并对 $i$ 归纳，可知对所有被 $I$ 零化的
$A$-模 $K$ 以及 $i = 1, \ldots, r + 1$，都有
$\text{Tor}_i^A(M, K) = 0$。

\medskip\noindent
假设对于某个 $1 \leq j \leq r$，已经证明：对所有被
$f_1, \ldots, f_j$ 零化的 $A$-模 $K$ 以及
$i = 1, \ldots, j + 1$，都有 $\text{Tor}_i^A(M, K) = 0$。
设 $K$ 是被 $f_1, \ldots, f_{j - 1}$ 零化的 $A$-模。
可以把映射 $f_j : K \to K$ 分解成两个短正合列
$$
0 \to K[f_j] \to K \to K' \to 0
\quad\text{且}\quad
0 \to K' \to K \to K/f_jK \to 0
$$
其中 $K' = K/K[f_j] \cong f_jK$。设 $1 \leq i \leq j$。
应用 Tor 的正合列，得到正合列
$$
\text{Tor}_i^A(K[f_j], M) \to
\text{Tor}_i^A(K, M) \to
\text{Tor}_i^A(K', M)
$$
以及
$$
\text{Tor}_{i + 1}^A(K/f_jK, M) \to
\text{Tor}_i^A(K', M) \to
\text{Tor}_i^A(K, M)
$$
利用 $\text{Tor}_i^A(K[f_j], M)$ 和
$\text{Tor}_{i + 1}^A(K/f_jK, M)$ 的消没，可知
$f_j : K \to K$ 在 $\text{Tor}^A_i(M, K)$ 上诱导单射。
换言之，$\text{Tor}^A_i(M, K)$ 为零，当且仅当
$\text{Tor}^A_i(M, K) \otimes_A A_{f_j}$ 为零。
由引理 \ref{algebra-lemma-flat-base-change-tor}，有
$$
\text{Tor}^A_i(M, K) \otimes_A A_{f_j} =
\text{Tor}^{A_{f_j}}_i(M_{f_j}, K_{f_j})  = 0
$$
最后一个消没来自 $M_{f_j}$ 在 $A_{f_j}$ 上的平坦性
（引理 \ref{algebra-lemma-characterize-flat}）。可知对所有被
$f_1, \ldots, f_{j - 1}$ 零化的 $A$-模 $K$ 以及
$i = 1, \ldots, j$，都有 $\text{Tor}_i^A(M, K) = 0$。
对 $j$ 作降归纳，得到 $j = 0$ 时的结论；也就是说，
对所有 $A$-模 $K$，都有 $\text{Tor}_1^A(M, K) = 0$。
因此，由引理 \ref{algebra-lemma-characterize-flat}，$M$ 在 $A$ 上平坦。

\medskip\noindent
忠实平坦模情形的论证从略。
\end{proof}

\section{基变换与平坦性}
\label{algebra-section-base-change-flat}

\noindent
下面给出若干处理基变换下平坦性变化的引理。

\begin{lemma}
\label{algebra-lemma-base-change-flat-up-down}
设
$$
\xymatrix{
S \ar[r] & S' \\
R \ar[r] \ar[u] & R' \ar[u]
}
$$
是局部环的局部同态所成的交换图。假设 $S'$ 是张量积
$S \otimes_R R'$ 的一个局部化。设 $M$ 是 $S$-模，并置
$M' = S' \otimes_S M$。
\begin{enumerate}
\item 若 $M$ 在 $R$ 上平坦，则 $M'$ 在 $R'$ 上平坦。
\item 若 $M'$ 在 $R'$ 上平坦，且 $R \to R'$ 平坦，则
$M$ 在 $R$ 上平坦。
\end{enumerate}
特别地，有
\begin{enumerate}
\item[(3)] 若 $S$ 在 $R$ 上平坦，则 $S'$ 在 $R'$ 上平坦。
\item[(4)] 若 $R' \to S'$ 和 $R \to R'$ 平坦，则 $S$ 在 $R$ 上平坦。
\end{enumerate}
\end{lemma}

\begin{proof}
(1) 的证明。若 $M$ 在 $R$ 上平坦，则由引理
\ref{algebra-lemma-flat-base-change}，$M \otimes_R R'$ 在 $R'$ 上平坦。
若 $W \subset S \otimes_R R'$ 是满足
$W^{-1}(S \otimes_R R') = S'$ 的乘法子集，则
$M' = W^{-1}(M \otimes_R R')$。因此，$M'$ 作为平坦模的局部化，
在 $R'$ 上平坦，见引理 \ref{algebra-lemma-flat-localization} 的 (5)。
这证明了 (1)，并且特别证明了 (3)。

\medskip\noindent
(2) 的证明。假设 $M'$ 在 $R'$ 上平坦，且 $R \to R'$ 平坦。
把 (3) 应用于沿西北对角线反射后的图，可知 $S \to S'$ 平坦。
于是，由引理 \ref{algebra-lemma-local-flat-ff}，$S \to S'$ 忠实平坦。
我们将用引理 \ref{algebra-lemma-flat} 的判据（\ref{algebra-item-f-ideal}）
证明 $M$ 平坦。设 $I \subset R$ 是理想。
% P01-E0493: frozen source says "has a kernel" although every map has one; intended nonzero kernel.
若 $I \otimes_R M \to M$ 的核非零，则
$(I \otimes_R M) \otimes_S S' \to M \otimes_S S' = M'$
的核也非零。注意，由于 $R \to R'$ 平坦，
$I \otimes_R R' = IR'$，并且
$$
(I \otimes_R M) \otimes_S S' =
(I \otimes_R R') \otimes_{R'} (M \otimes_S S') =
IR' \otimes_{R'} M'.
$$
由 $M'$ 在 $R'$ 上的平坦性，可知它单射到 $M'$ 中。
这完成了 (2) 的证明，因而 (4) 也成立。
\end{proof}

\noindent
下面是平坦性局部判据的又一个应用。

\begin{lemma}
\label{algebra-lemma-yet-another-variant-local-criterion-flatness}
考虑局部环和局部同态所成的交换图
$$
\xymatrix{
S \ar[r] & S' \\
R \ar[r] \ar[u] & R' \ar[u]
}
$$
设 $M$ 是有限 $S$-模。假设
\begin{enumerate}
\item 水平箭头是平坦环映射；
\item $M$ 在 $R$ 上平坦；
\item $\mathfrak m_R R' = \mathfrak m_{R'}$；
\item $R'$ 和 $S'$ 都是 Noether 环。
\end{enumerate}
则 $M' = M \otimes_S S'$ 在 $R'$ 上平坦。
\end{lemma}

\begin{proof}
由于 $\mathfrak m_R \subset R$ 且 $R \to R'$ 平坦，由假设 (3) 得到
$\mathfrak m_R \otimes_R R' = \mathfrak m_R R' = \mathfrak m_{R'}$。
注意，由引理 \ref{algebra-lemma-flatness-descends-more-general}，$M'$ 是有限
$S'$-模，且在 $R$ 上平坦。因此
$\mathfrak m_R \otimes_R M' \to M'$ 是单射。又有
$$
\mathfrak m_R \otimes_R M' =
\mathfrak m_R \otimes_R R' \otimes_{R'} M' =
\mathfrak m_{R'} \otimes_{R'} M'
$$
故 $\mathfrak m_{R'} \otimes_{R'} M' \to M'$ 是单射。
这表明 $\text{Tor}_1^{R'}(\kappa_{R'}, M') = 0$
（注记 \ref{algebra-remark-Tor-ring-mod-ideal}）。因此，由引理
\ref{algebra-lemma-local-criterion-flatness}，$M'$ 在 $R'$ 上平坦。
\end{proof}

\section{Artin 环上的平坦性判据}
\label{algebra-section-flatness-artinian}

\noindent
我们讨论 Artin 环上模的一些平坦性判据。注意，Artin 局部环的
极大理想是幂零的，所以以下两个引理适用于 Artin 局部环。

\begin{lemma}
\label{algebra-lemma-local-artinian-basis-when-flat}
设 $(R, \mathfrak m)$ 是极大理想 $\mathfrak m$ 幂零的局部环。
设 $M$ 是平坦 $R$-模。若 $A$ 是集合，而
$x_\alpha \in M$（$\alpha \in A$）是 $M$ 中的一族元素，
则以下条件等价：
\begin{enumerate}
\item $\{\overline{x}_\alpha\}_{\alpha \in A}$ 构成
$R/\mathfrak m$ 上向量空间 $M/\mathfrak mM$ 的一组基；
\item $\{x_\alpha\}_{\alpha \in A}$ 构成 $M$ 在 $R$ 上的一组基。
\end{enumerate}
\end{lemma}

\begin{proof}
(2) $\Rightarrow$ (1) 是立即的。假设 (1)。由 Nakayama 引理
\ref{algebra-lemma-NAK}，诸 $x_\alpha$ 生成 $M$。于是得到短正合列
$$
0 \to K \to \bigoplus\nolimits_{\alpha \in A} R \to M \to 0
$$
与 $R/\mathfrak m$ 张量，并使用引理 \ref{algebra-lemma-flat-tor-zero}，
得到 $K/\mathfrak mK = 0$。由 Nakayama 引理 \ref{algebra-lemma-NAK}，
可知 $K = 0$。
\end{proof}

\begin{lemma}
\label{algebra-lemma-local-artinian-characterize-flat}
设 $R$ 是极大理想幂零的局部环。设 $M$ 是 $R$-模。
以下条件等价：
\begin{enumerate}
\item $M$ 在 $R$ 上平坦；
\item $M$ 是自由 $R$-模；
\item $M$ 是投射 $R$-模。
\end{enumerate}
\end{lemma}

\begin{proof}
由于任意投射模都是平坦的（它是自由模的直和因子），而每个自由模
都是投射的，只须证明平坦模是自由的。设 $M$ 是平坦模。
设 $A$ 是集合，并设 $x_\alpha \in M$（$\alpha \in A$）是使
$\overline{x_\alpha} \in M/\mathfrak m M$ 构成 $R$ 的剩余域上的
一组基的元素。由引理 \ref{algebra-lemma-local-artinian-basis-when-flat}，
诸 $x_\alpha$ 构成 $M$ 在 $R$ 上的一组基，结论即得。
\end{proof}

\begin{lemma}
\label{algebra-lemma-lift-basis}
设 $R$ 是环。设 $I \subset R$ 是理想。设 $M$ 是 $R$-模。
设 $A$ 是集合，并设 $x_\alpha \in M$（$\alpha \in A$）是
$M$ 中的一族元素。假设
\begin{enumerate}
\item $I$ 是幂零理想；
\item $\{\overline{x}_\alpha\}_{\alpha \in A}$ 构成
$M/IM$ 在 $R/I$ 上的一组基；
\item $\text{Tor}_1^R(R/I, M) = 0$。
\end{enumerate}
则 $M$ 是以 $\{x_\alpha\}_{\alpha \in A}$ 为基的自由 $R$-模。
\end{lemma}

\begin{proof}
设 $R$、$I$、$M$、$\{x_\alpha\}_{\alpha \in A}$ 如引理中所述，
并满足假设 (1)、(2)、(3)。由 Nakayama 引理 \ref{algebra-lemma-NAK}，
诸 $x_\alpha$ 在 $R$ 上生成 $M$。假设
$\text{Tor}_1^R(R/I, M) = 0$ 蕴含有短正合列
$$
0 \to I \otimes_R M \to M \to M/IM \to 0.
$$
设 $\sum f_\alpha x_\alpha = 0$ 是 $M$ 中的一个关系。
由 $x_\alpha$ 的选取可知 $f_\alpha \in I$。因此，在
$I \otimes_R M$ 中有 $\sum f_\alpha \otimes x_\alpha = 0$。
映射 $I \otimes_R M \to I/I^2 \otimes_{R/I} M/IM$ 与
$\{x_\alpha\}_{\alpha \in A}$ 构成 $M/IM$ 的一组基这一事实
蕴含 $f_\alpha \in I^2$！因此，诸 $x_\alpha$ 在
$M/I^2M$ 中的像之间没有关系。换言之，$M/I^2M$ 是自由的，
并以诸 $x_\alpha$ 的像为基。然后使用映射
$I \otimes_R M \to I/I^3 \otimes_{R/I^2} M/I^2M$，
可知 $f_\alpha \in I^3$！如此继续。由假设 (1)，对某个 $n$
有 $I^n = 0$，结论即得。
\end{proof}

\begin{lemma}
\label{algebra-lemma-prepare-lift-flatness}
设 $\varphi : R \to R'$ 是环映射。设 $I \subset R$ 是理想。
设 $M$ 是 $R$-模。假设
\begin{enumerate}
\item $M/IM$ 在 $R/I$ 上平坦；
\item $R' \otimes_R M$ 在 $R'$ 上平坦。
\end{enumerate}
置 $I_2 = \varphi^{-1}(\varphi(I^2)R')$。则 $M/I_2M$
在 $R/I_2$ 上平坦。
\end{lemma}

\begin{proof}
可以分别用 $R/I_2$、$M/I_2M$、$R'/\varphi(I)^2R'$ 代替
$R$、$M$、$R'$。此时 $I^2 = 0$，且 $\varphi$ 是单射。
由引理 \ref{algebra-lemma-what-does-it-mean} 及 $I^2 = 0$ 这一事实，
只须证明
$\text{Tor}^R_1(R/I, M) = K = \Ker(I \otimes_R M \to M)$ 为零。
置 $M' = M \otimes_R R'$ 且 $I' = IR'$。由假设，映射
$I' \otimes_{R'} M' \to M'$ 是单射。因此，$K$ 在
$$
I' \otimes_{R'} M' = I' \otimes_R M = I' \otimes_{R/I} M/IM.
$$
中的像为零。此时 $I \to I'$ 是 $R/I$-模的单射。
由于 $M/IM$ 在 $R/I$ 上平坦，映射
$$
I \otimes_{R/I} M/IM \longrightarrow I' \otimes_{R/I} M/IM
$$
是单射。这蕴含 $K$ 在 $I \otimes_R M = I \otimes_{R/I} M/IM$
中为零，如所求。
\end{proof}

\begin{lemma}
\label{algebra-lemma-lift-flatness}
设 $\varphi : R \to R'$ 是环映射。设 $I \subset R$ 是理想。
设 $M$ 是 $R$-模。假设
\begin{enumerate}
\item $I$ 是幂零理想；
\item $R \to R'$ 是单射；
\item $M/IM$ 在 $R/I$ 上平坦；
\item $R' \otimes_R M$ 在 $R'$ 上平坦。
\end{enumerate}
则 $M$ 在 $R$ 上平坦。
\end{lemma}

\begin{proof}
归纳定义 $I_1 = I$ 以及
$I_{n + 1} = \varphi^{-1}(\varphi(I_n)^2R')$（$n \geq 1$）。
注意，由引理 \ref{algebra-lemma-prepare-lift-flatness}，对每个 $n \geq 1$，
$M/I_nM$ 都在 $R/I_n$ 上平坦。显然，
$\varphi(I_{n + 1}) \subset \varphi(I)^{2^n}R'$。由于 $I$ 幂零，
对某个 $n$ 有 $\varphi(I_n) = 0$。又因 $\varphi$ 是单射，
可知对某个 $n$ 有 $I_n = 0$，结论即得。
\end{proof}

\noindent
下面是平坦性局部判据的局部 Artin 版本。

\begin{lemma}
\label{algebra-lemma-artinian-variant-local-criterion-flatness}
设 $R$ 是 Artin 局部环。设 $M$ 是 $R$-模。
设 $I \subset R$ 是正常理想。以下条件等价：
\begin{enumerate}
\item $M$ 在 $R$ 上平坦；
\item $M/IM$ 在 $R/I$ 上平坦，且 $\text{Tor}_1^R(R/I, M) = 0$。
\end{enumerate}
\end{lemma}

\begin{proof}
(1) $\Rightarrow$ (2) 直接由定义得出。假设 $M/IM$ 在 $R/I$
上平坦，且 $\text{Tor}_1^R(R/I, M) = 0$。由引理
\ref{algebra-lemma-local-artinian-characterize-flat}，这蕴含 $M/IM$
在 $R/I$ 上自由。选取集合 $A$ 和元素 $x_\alpha \in M$，
使它们在 $M/IM$ 中的像构成一组基。由引理
\ref{algebra-lemma-lift-basis}，可知 $M$ 是自由的，因而特别是平坦的。
\end{proof}

\noindent
事实上，平坦性可沿源为 Artin 环的单射同态下降。

\begin{lemma}
\label{algebra-lemma-descent-flatness-injective-map-artinian-rings}
设 $R \to S$ 是环映射。设 $M$ 是 $R$-模。假设
\begin{enumerate}
\item $R$ 是 Artin 环；
\item $R \to S$ 是单射；
\item $M \otimes_R S$ 是平坦 $S$-模。
\end{enumerate}
则 $M$ 是平坦 $R$-模。
\end{lemma}

\begin{proof}
第一种证明：设 $I \subset R$ 是 $R$ 的 Jacobson 根。
则 $I$ 幂零；而 $M/IM$ 在 $R/I$ 上平坦，因为 $R/I$
是域的乘积，见第 \ref{algebra-section-artinian} 节。因此，应用引理
\ref{algebra-lemma-lift-flatness}，可知 $M$ 平坦。

\medskip\noindent
第二种证明：由引理 \ref{algebra-lemma-artinian-finite-length}，可将
$R = \prod R_i$ 写成有限个局部 Artin 环的乘积。这对 $R$ 和 $S$
都诱导出类似的乘积分解。因此可约化到 $R$ 是局部 Artin 环的情形
（细节略）。

\medskip\noindent
假设 $R \to S$、$M$ 如引理中所述，满足 (1)、(2)、(3)，并且
$R$ 还是极大理想为 $\mathfrak m$ 的局部环。设 $A$ 是集合，并设
% P01-E0494: frozen source types x_alpha as an element of the index set A, not of M.
$x_\alpha \in A$ 是使 $\overline{x}_\alpha$ 构成
$M/\mathfrak mM$ 在 $R/\mathfrak m$ 上的一组基的元素。
由 Nakayama 引理 \ref{algebra-lemma-NAK}，诸 $x_\alpha$ 作为 $R$-模生成
$M$。置 $N = S \otimes_R M$ 且 $I = \mathfrak mS$。
于是 $\{1 \otimes x_\alpha\}_{\alpha \in A}$ 是 $N$ 中的一族元素，
并构成 $N/IN$ 的一组基。此外，由于 $N$ 在 $S$ 上平坦，
有 $\text{Tor}_1^S(S/I, N) = 0$。因此，由引理
\ref{algebra-lemma-lift-basis}，$N$ 是以
$\{1 \otimes x_\alpha\}_{\alpha \in A}$ 为基的自由模。
$R \to S$ 的单射性于是保证，诸 $x_\alpha$ 之间不可能有系数在
$R$ 中的非平凡关系。
\end{proof}

\noindent
请把下面的引理与引理
\ref{algebra-lemma-criterion-flatness-fibre-Noetherian}（Noether 局部环的情形）、
引理 \ref{algebra-lemma-criterion-flatness-fibre}（有限表示代数的情形），
以及引理 \ref{algebra-lemma-criterion-flatness-fibre-locally-nilpotent}
（局部幂零理想的情形）比较。

\begin{lemma}[按纤维的平坦性判据：幂零情形]
\label{algebra-lemma-criterion-flatness-fibre-nilpotent}
设
$$
\xymatrix{
S \ar[rr] & & S' \\
& R \ar[lu] \ar[ru]
}
$$
是环范畴中的交换图。设 $I \subset R$ 是幂零理想，
$M$ 是 $S'$-模。假设
\begin{enumerate}
\item 模 $M/IM$ 是平坦 $S/IS$-模。
\item 模 $M$ 是平坦 $R$-模。
\end{enumerate}
则 $M$ 是平坦 $S$-模；而且对每个满足
$M \otimes_S \kappa(\mathfrak q)$ 非零的 $\mathfrak q \subset S$，
$S_{\mathfrak q}$ 在 $R$ 上平坦。
\end{lemma}

\begin{proof}
由于 $M$ 在 $R$ 上平坦，把短正合列
$0 \to I \to R \to R/I \to 0$ 与之张量，得到短正合列
$$
0 \to I \otimes_R M \to M \to M/IM \to 0.
$$
注意，$I \otimes_R M \to IS \otimes_S M$ 是满射。结合上式，
这意味着复合映射中的两个映射
$$
I \otimes_R M \to IS \otimes_S M \to M
$$
都是单射。因此
% P01-E0495: frozen source has Tor_1^S(IS,M); the ideal-kernel criterion requires S/IS.
$\text{Tor}_1^S(IS, M) = 0$（见注记
\ref{algebra-remark-Tor-ring-mod-ideal}）；再由引理
\ref{algebra-lemma-what-does-it-mean}，可知 $M$ 是平坦 $S$-模。
最后，须证明对任何满足 $M \otimes_S \kappa(\mathfrak q)$ 非零的
素理想 $\mathfrak q \subset S$，$S_{\mathfrak q}$ 在 $R$ 上平坦。
这由引理 \ref{algebra-lemma-ff} 和 \ref{algebra-lemma-flat-permanence} 得出。
\end{proof}

\section{复形何时正合？}
\label{algebra-section-complex-exact}

\noindent
本节的部分材料可见 Buchsbaum 和 Eisenbud 的论文
\cite{WhatExact}。

\begin{situation}
\label{algebra-situation-complex}
这里 $R$ 是环，并且有复形
$$
0
\to
R^{n_e}
\xrightarrow{\varphi_e}
R^{n_{e-1}}
\xrightarrow{\varphi_{e-1}}
\ldots
\xrightarrow{\varphi_{i + 1}}
R^{n_i}
\xrightarrow{\varphi_i}
R^{n_{i-1}}
\xrightarrow{\varphi_{i-1}}
\ldots
\xrightarrow{\varphi_1}
R^{n_0}
$$
换言之，对 $i = 1, \ldots, e - 1$，要求
$\varphi_i \circ \varphi_{i + 1} = 0$。
\end{situation}

\begin{lemma}
\label{algebra-lemma-add-trivial-complex}
设 $R$ 是环。设
$$
\ldots
\xrightarrow{\varphi_{i + 1}}
R^{n_i}
\xrightarrow{\varphi_i}
R^{n_{i-1}}
\xrightarrow{\varphi_{i-1}}
\ldots
$$
是有限自由 $R$-模的复形。假设对某个 $i$，映射 $\varphi_i$
的某个矩阵系数可逆。则所示复形同构于一个复形
$$
\ldots \to
R^{n_{i + 2}} \xrightarrow{\varphi_{i + 2}}
R^{n_{i + 1}} \to
R^{n_i - 1} \to
R^{n_{i - 1} - 1} \to
R^{n_{i - 2}} \xrightarrow{\varphi_{i - 2}}
R^{n_{i - 3}} \to
\ldots
$$
与复形 $\ldots \to 0 \to R \to R \to 0 \to \ldots$ 的直和，
其中映射 $R \to R$ 是恒等映射。
\end{lemma}

\begin{proof}
该假设意味着，在对 $R^{n_i}$ 和 $R^{n_{i-1}}$ 换基之后，
$R^{n_i}$ 的第一个基向量经 $\varphi_i$ 映到
$R^{n_{i-1}}$ 的第一个基向量。用 $e_j$ 表示 $R^{n_i}$
的第 $j$ 个基向量，用 $f_k$ 表示 $R^{n_{i-1}}$ 的第 $k$
个基向量。写成 $\varphi_i(e_j) = \sum a_{jk} f_k$。
于是除 $k = 1$ 外均有 $a_{1k} = 0$，且 $a_{11} = 1$。
再次更换 $R^{n_i}$ 的基：对 $j > 1$ 置
$e'_j = e_j - a_{j1} e_1$。换坐标后，对 $j > 1$ 有
$a_{j1} = 0$。注意，映射 $R^{n_{i + 1}} \to R^{n_i}$
的像包含于由 $e_j$（$j > 1$）张成的
% P01-E0497: frozen source says "subspace" although R is an arbitrary ring.
子模。还要注意，因为 $f_1$ 在像中，映射
$R^{n_{i-1}} \to R^{n_{i-2}}$ 必须将它零化。
这些条件以及 $\varphi_i$ 的矩阵 $(a_{jk})$ 的形状即蕴含本引理。
\end{proof}

\noindent
在情形 \ref{algebra-situation-complex} 中，我们称形如
$$
0 \to \ldots \to 0 \to R \xrightarrow{1} R \to 0 \to \ldots \to 0
$$
或形如
$$
0 \to \ldots \to 0 \to R
$$
的复形为{\it 平凡的}。更准确地说，若存在
$e \geq i \geq 1$，使得 $n_i = n_{i - 1} = 1$、
$\varphi_i = \text{id}_R$，且对
$j \not \in \{i, i - 1\}$ 有 $n_j = 0$；或者若
$n_0 = 1$ 且对 $i > 0$ 有 $n_i = 0$，则称
$0 \to R^{n_e} \to R^{n_{e-1}} \to \ldots \to R^{n_0}$
平凡。上面的引理显然说明，局部环上的任意有限自由模有限复形，
在与平凡复形作直和的意义下，都等同于一个所有映射的全部矩阵系数
均属于极大理想的复形。

\begin{lemma}
\label{algebra-lemma-exact-depth-zero-local}
在情形 \ref{algebra-situation-complex} 中，假设 $R$ 是极大理想为
$\mathfrak m$ 的 Noether 局部环。假设
$\mathfrak m \in \text{Ass}(R)$，换言之，$R$ 的深度为 $0$。
假设
$0 \to R^{n_e} \to R^{n_{e-1}} \to \ldots \to R^{n_0}$
在 $R^{n_e}, \ldots, R^{n_1}$ 处正合。则该复形同构于平凡复形
的一个直和。
\end{lemma}

\begin{proof}
选取满足 $x \not = 0$ 且 $\mathfrak m x = 0$ 的 $x \in R$。
设 $i$ 是满足 $n_i > 0$ 的最大指标。若 $i = 0$，则命题成立。
若 $i > 0$，用 $f_1$ 表示 $R^{n_i}$ 的第一个基向量。
由复形的正合性，$xf_1$ 不会映到零。因此，映射
$R^{n_i} \to R^{n_{i - 1}}$ 的某个矩阵系数不属于
$\mathfrak m$。于是，引理 \ref{algebra-lemma-add-trivial-complex}
使我们可以减小 $n_e + \ldots + n_1$。归纳即完成证明。
\end{proof}

\begin{lemma}
\label{algebra-lemma-exact-artinian-local}
在情形 \ref{algebra-situation-complex} 中，设 $R$ 是 Artin 局部环。
假设 $0 \to R^{n_e} \to R^{n_{e-1}} \to \ldots \to R^{n_0}$
在 $R^{n_e}, \ldots, R^{n_1}$ 处正合。则该复形同构于平凡复形
的一个直和。
\end{lemma}

\begin{proof}
这是引理 \ref{algebra-lemma-exact-depth-zero-local} 的特殊情形，因为
Artin 局部环的深度为 $0$。
\end{proof}

\noindent
下面定义有限自由模之间映射的秩。这只是秩的一种可能定义，
也只是适合本节使用的定义；在其他场合，别的定义或许更方便。

\begin{definition}
\label{algebra-definition-rank}
设 $R$ 是环。设 $\varphi : R^m \to R^n$ 是有限自由模之间的映射。
\begin{enumerate}
\item $\varphi$ 的{\it 秩}是使
$\wedge^r \varphi : \wedge^r R^m \to \wedge^r R^n$
非零的最大 $r$。
\item 令 $I(\varphi) \subset R$ 为 $\varphi$ 的矩阵的所有
$r \times r$ 子式生成的理想，其中 $r$ 是上面定义的秩。
\end{enumerate}
\end{definition}

\noindent
映射 $\varphi : R^m \to R^n$ 的秩为 $0$ 当且仅当
$\varphi = 0$；在这种情形下，$I(\varphi) = R$。

\begin{lemma}
\label{algebra-lemma-trivial-case-exact}
在情形 \ref{algebra-situation-complex} 中，假设该复形同构于平凡复形
的一个直和。则
\begin{enumerate}
\item 映射 $\varphi_i$ 的秩为
$r_i = n_i - n_{i + 1} + \ldots + (-1)^{e-i-1} n_{e-1} + (-1)^{e-i} n_e$；
\item 对所有 $i$，$1 \leq i \leq e - 1$，有
$\text{rank}(\varphi_{i + 1}) + \text{rank}(\varphi_i) = n_i$；
\item 每个 $I(\varphi_i) = R$。
\end{enumerate}
\end{lemma}

\begin{proof}
可以假设该复形就是平凡复形的直和。于是，对每个 $i$，
可以把 $R^{n_i}$ 的标准基元素分成两类：一类映到
$R^{n_{i-1}}$ 的基元素，另一类映到零（并且当 $i > 0$ 时，
后一类元素是 $R^{n_{i + 1}}$ 中一些基元素的像）。
从 $i = e$ 开始作降归纳，很容易证明：$R^{n_i}$ 中有
$r_{i + 1}$ 个基元素映到零，并有 $r_i$ 个基元素映到
$R^{n_{i-1}}$ 的基元素。结论由此得出。
\end{proof}

\begin{lemma}
\label{algebra-lemma-div-x-exact-one-less}
在情形 \ref{algebra-situation-complex} 中，假设 $R$ 是极大理想为
$\mathfrak m$ 的局部环。假设
$0 \to R^{n_e} \to R^{n_{e-1}} \to \ldots \to R^{n_0}$
在 $R^{n_e}, \ldots, R^{n_1}$ 处正合。设
$x \in \mathfrak m$ 是非零因子。则复形
$0 \to (R/xR)^{n_e} \to \ldots \to (R/xR)^{n_1}$
在 $(R/xR)^{n_e}, \ldots, (R/xR)^{n_2}$ 处正合。
\end{lemma}

\begin{proof}
用 $F_\bullet$ 表示项为 $F_i = R^{n_i}$、微分由
$\varphi_i$ 给出的复形。于是有复形的短正合列
$$
0 \to F_\bullet \xrightarrow{x} F_\bullet \to F_\bullet/xF_\bullet \to 0
$$
应用蛇引理，得到长正合列
$$
H_i(F_\bullet) \xrightarrow{x} H_i(F_\bullet) \to
H_i(F_\bullet/xF_\bullet) \to H_{i - 1}(F_\bullet)
\xrightarrow{x} H_{i - 1}(F_\bullet)
$$
引理由此得出。
\end{proof}

\begin{lemma}[无环性引理]
\label{algebra-lemma-acyclic}
\begin{reference}
\cite[Lemma 1.8]{Peskine-Szpiro}
\end{reference}
设 $R$ 是 Noether 局部环。设
$0 \to M_e \to M_{e-1} \to \ldots \to M_0$
是有限 $R$-模的复形。假设 $\text{depth}(M_i) \geq i$。
设 $i$ 是使复形在 $M_i$ 处不正合的最大指标。若 $i > 0$，则
$$
\Ker(M_i \to M_{i-1})/\Im(M_{i + 1} \to M_i)
$$
的深度 $\geq 1$。
\end{lemma}

\begin{proof}
令 $H = \Ker(M_i \to M_{i-1})/\Im(M_{i + 1} \to M_i)$ 为所讨论的
% P01-E0501: frozen source calls H a cohomology group although this is a homological complex.
同调群。可以把该复形拆成正合列
% P01-E0502: frozen source calls every listed sequence short exact, but one ends at M_{i-1}.
$0 \to M_e \to M_{e-1} \to K_{e-2} \to 0$、
$0 \to K_j \to M_j \to K_{j-1} \to 0$（$i + 2 \leq j \leq e-2 $）、
$0 \to K_{i + 1} \to M_{i + 1} \to B_i \to 0$、
$0 \to K_i \to M_i \to M_{i-1}$，以及
$0 \to B_i \to K_i \to H \to 0$。
我们沿这些序列逐步向上，反复使用引理
\ref{algebra-lemma-depth-in-ses} 来证明关于深度的断言。
首先，由于 $\text{depth}(M_e) \geq e$ 且
$\text{depth}(M_{e-1}) \geq e-1$，可知
$\text{depth}(K_{e-2}) \geq e - 1$。此时，对
$i + 2 \leq j \leq e-2$，序列
$0 \to K_j \to M_j \to K_{j-1} \to 0$ 同样蕴含
$\text{depth}(K_{j-1}) \geq j$（$i + 2 \leq j \leq e-2$）。
序列
$0 \to K_{i + 1} \to M_{i + 1} \to B_i \to 0$
于是表明 $\text{depth}(B_i) \geq i + 1$。序列
$0 \to K_i \to M_i \to M_{i-1}$ 表明
$\text{depth}(K_i) \geq 1$，因为由假设
$M_i$ 的深度 $\geq i \geq 1$。最后，序列
$0 \to B_i \to K_i \to H \to 0$ 蕴含结论。
\end{proof}

\begin{proposition}
\label{algebra-proposition-what-exact}
\begin{reference}
\cite[Corollary 1]{WhatExact}
\end{reference}
在情形 \ref{algebra-situation-complex} 中，假设 $R$ 是 Noether 局部环。
以下条件等价：
\begin{enumerate}
\item $0 \to R^{n_e} \to R^{n_{e-1}} \to \ldots \to R^{n_0}$
在 $R^{n_e}, \ldots, R^{n_1}$ 处正合；
\item 对所有 $i$，$1 \leq i \leq e$，以下两个条件成立：
\begin{enumerate}
\item $\text{rank}(\varphi_i) = r_i$，其中
$r_i = n_i - n_{i + 1} + \ldots + (-1)^{e-i-1} n_{e-1} + (-1)^{e-i} n_e$；
\item $I(\varphi_i) = R$，或 $I(\varphi_i)$ 包含一个长度为 $i$
的正则序列。
\end{enumerate}
\end{enumerate}
\end{proposition}

\begin{proof}
若对某个 $i$，$\varphi_i$ 的某个矩阵系数不属于
$\mathfrak m$，则应用引理 \ref{algebra-lemma-add-trivial-complex}。
不难看出，一个复形的命题与给该复形加上一个平凡复形后所得复形
的命题等价。因此，可以假设每个 $\varphi_i$ 的所有矩阵元素
都属于极大理想。还可以假设 $e \geq 1$。

\medskip\noindent
假设该复形在 $R^{n_e}, \ldots, R^{n_1}$ 处正合。
设 $\mathfrak q \in \text{Ass}(R)$。注意，环
$R_{\mathfrak q}$ 的深度为 $0$，而该复形在
$\mathfrak q$ 处局部化后仍然正合。对 $R_{\mathfrak q}$
上的局部化复形应用引理 \ref{algebra-lemma-exact-depth-zero-local}
和引理 \ref{algebra-lemma-trivial-case-exact}，可知对所有 $i$，
$\varphi_{i, \mathfrak q}$ 的秩为 $r_i$。由于
$R \to \bigoplus_{\mathfrak q \in \text{Ass}(R)} R_\mathfrak q$
是单射（引理 \ref{algebra-lemma-zero-at-ass-zero}），由定义
\ref{algebra-definition-rank} 中给出的秩的定义可知，$\varphi_i$
在 $R$ 上的秩为 $r_i$。因此，由于秩不变，可见
$I(\varphi_i)_\mathfrak q = I(\varphi_{i, \mathfrak q})$。
上面引用的引理说明，对所有 $e \geq i \geq 1$，理想
$I(\varphi_i)_{\mathfrak q}$ 都等于 $R_{\mathfrak q}$；
故没有任何 $I(\varphi_i)$ 包含于 $\mathfrak q$。这蕴含
$I(\varphi_e)I(\varphi_{e-1})\ldots I(\varphi_1)$ 不包含于
$R$ 的任何伴随素理想。由引理 \ref{algebra-lemma-silly}，可以选取
$$
x \in I(\varphi_e)I(\varphi_{e - 1})\ldots I(\varphi_1)
$$
使得对所有 $\mathfrak q \in \text{Ass}(R)$，
$x \not \in \mathfrak q$。注意，$x$ 是非零因子
（引理 \ref{algebra-lemma-ass-zero-divisors}）。由引理
\ref{algebra-lemma-div-x-exact-one-less}，复形
$0 \to (R/xR)^{n_e} \to \ldots \to (R/xR)^{n_1}$
在 $(R/xR)^{n_e}, \ldots, (R/xR)^{n_2}$ 处正合。
对 $e$ 归纳可知，所有理想 $I(\varphi_i)/xR$ 都含有一个
长度为 $i - 1$ 的正则序列。这证明了 $I(\varphi_i)$
含有一个长度为 $i$ 的正则序列。

\medskip\noindent
假设 (2)(a) 和 (2)(b) 成立。我们对 $\dim(R)$ 归纳证明 (1)。
若 $\dim(R) = 0$，则由 (2)(b)，对 $1 \leq i \leq e$
必有 $I(\varphi_i) = R$。由于 $\varphi_i$ 的系数属于极大理想，
这只有在对所有 $i$ 都有 $r_i = 0$ 时才可能。由 (2)(a)，
可知 $e = 0$，故 (1) 成立。假设 $\dim(R) > 0$。
我们断言，对任何素理想 $\mathfrak p \subset R$，条件
(2)(a) 和 (2)(b) 对 $R_\mathfrak p$ 上的复形
$$
0 \to R_\mathfrak p^{n_e} \to R_\mathfrak p^{n_{e - 1}} \to
\ldots \to R_\mathfrak p^{n_0}
$$
以及映射 $\varphi_{i, \mathfrak p}$ 仍然成立。事实上，由于
$I(\varphi_i)$ 含有非零因子，$I(\varphi_i)$ 在
$R_\mathfrak p$ 中的像非零。
这蕴含 $\varphi_{i, \mathfrak p}$ 的秩与 $\varphi_i$ 的秩相同：
上面定义的环 $R$ 上矩阵 $\varphi$ 的秩，在传递到 $R$-代数
$R'$ 时只能下降，而且下降当且仅当 $I(\varphi)$ 在 $R'$
中的像为零。因此 (2)(a) 成立。由此可知
$I(\varphi_{i, \mathfrak p}) = I(\varphi_i)_\mathfrak p$，
并且 (2)(b) 在局部化下也保持成立。
由对 $R$ 的维数的归纳，可以假设该复形在 $R$ 的任何非极大
素理想 $\mathfrak p$ 处局部化后都正合。因此
$\Ker(\varphi_i)/\Im(\varphi_{i + 1})$ 的支撑包含于
$\{\mathfrak m\}$，故它若非零，深度便为 $0$。
由于第一段所述的理由，对所有 $i$ 都有
% P01-E0503: frozen source overlooks rank-zero maps, for which I(varphi_i)=R.
$I(\varphi_i) \subset \mathfrak m$，所以 (2)(b) 蕴含
$\text{depth}(R) \geq e$。由引理 \ref{algebra-lemma-acyclic}，
该复形在 $R^{n_e}, \ldots, R^{n_1}$ 处正合，证明完成。
\end{proof}

\begin{remark}
\label{algebra-remark-what-exact}
若命题 \ref{algebra-proposition-what-exact} 中的等价条件 (1) 和 (2)
成立，则存在 $j$，使得 $I(\varphi_i) = R$ 当且仅当
$i \geq j$。与命题的证明一样，只须在所有矩阵的系数均属于
$R$ 的极大理想 $\mathfrak m$ 时验证这一点。在这种情形下，
$I(\varphi_j) = R$ 当且仅当 $\varphi_j = 0$。但若
$\varphi_j = 0$，则得到任意长的正合复形
$0 \to R^{n_e} \to R^{n_{e-1}} \to \ldots
\to R^{n_j} \to 0 \to 0 \to \ldots \to 0$，
因而由该命题可知，对 $i > j$，$I(\varphi_i)$ 必须为 $R$
（否则它是 Noether 局部环的一个真理想，却包含任意长的正则序列，
这是不可能的）。
\end{remark}
\section{Cohen--Macaulay 模}
\label{algebra-section-CM}

\noindent
这里将说明 Cohen--Macaulay 模具有良好的性质。我们暂不使用
Ext 群来建立它们与对偶性等概念的联系。

\begin{definition}
\label{algebra-definition-CM}
设 $R$ 是 Noether 局部环。设 $M$ 是有限 $R$-模。若
$\dim(\text{Supp}(M)) = \text{depth}(M)$，则称 $M$ 是
{\it Cohen--Macaulay 模}。
\end{definition}

\noindent
我们的第一个目标是证明命题 \ref{algebra-proposition-CM-module}。
为此，将给出一系列（或许并不标准的）初等引理，其中几乎不使用
前面关于深度的结果。先引入一些记号。

\medskip\noindent
设 $R$ 是 Noether 局部环。设 $M$ 是 Cohen--Macaulay 模，
$f_1, \ldots, f_d$ 是一个 $M$-正则序列，其中
$d = \dim(\text{Supp}(M))$。若对所有
$i = 0, 1, \ldots, d-1$ 都有
$\dim (\text{Supp}(M) \cap V(g, f_1, \ldots, f_i))
= d - i - 1$，则称 $g \in \mathfrak m$ {\it 关于
$(M, f_1, \ldots, f_d)$ 良好}。这等价于：对
$i = 0, 1, \ldots, d - 1$，有
$\dim(\text{Supp}(M/(f_1, \ldots, f_i)M) \cap V(g)) =
d - i - 1$。

\begin{lemma}
\label{algebra-lemma-good-element}
记号和假设同上。若 $g$ 关于 $(M, f_1, \ldots, f_d)$ 良好，
则 (a) $g$ 是 $M$ 上的非零因子，并且 (b) $M/gM$ 是
Cohen--Macaulay 模，其极大正则序列为
$f_1, \ldots, f_{d - 1}$。
\end{lemma}

\begin{proof}
对 $d$ 归纳证明该引理。
% P01-E0504: for d=0 the definition of good is vacuous, contrary to the source's claim that no case applies.
若 $d = 0$，则 $M$ 有限，并且不存在本引理适用的情形。
若 $d = 1$，则须证明 $g : M \to M$ 是单射。由于假设
$\dim \text{Supp}(M) \cap V(g) = 0$，其核 $K$ 的支撑为
$\{\mathfrak m\}$，故 $K$ 具有有限长度。因此，
$f_1 : K \to K$ 的单射性蕴含其像与 $K$ 长度相同，从而
$f_1 K = K$；再由 Nakayama 引理 \ref{algebra-lemma-NAK} 可知
$K = 0$。此外，$\dim \text{Supp}(M/gM) = 0$，所以
$M/gM$ 是深度为 $0$ 的 Cohen--Macaulay 模。

\medskip\noindent
假设 $d > 1$。由定义容易看出，$g$ 关于
$(M/f_1M, f_2, \ldots, f_d)$ 良好。由归纳假设，
(a) $g$ 是 $M/f_1M$ 上的非零因子，并且
(b) $M/(g, f_1)M$ 是 Cohen--Macaulay 模，其极大正则序列为
$f_2, \ldots, f_{d - 1}$。由引理 \ref{algebra-lemma-permute-xi}，
可知 $g, f_1$ 是一个 $M$-正则序列。因此，$g$ 是 $M$ 上的
非零因子，而 $f_1, \ldots, f_{d - 1}$ 是一个
$M/gM$-正则序列。
\end{proof}

\begin{lemma}
\label{algebra-lemma-CM-one-g}
设 $R$ 是 Noether 局部环。设 $M$ 是 $R$ 上的
Cohen--Macaulay 模。假设 $g \in \mathfrak m$ 满足
$\dim(\text{Supp}(M) \cap V(g))
= \dim(\text{Supp}(M)) - 1$。则 (a) $g$ 是 $M$ 上的非零因子，
并且 (b) $M/gM$ 是深度少一的 Cohen--Macaulay 模。
\end{lemma}

\begin{proof}
选取一个 $M$-正则序列 $f_1, \ldots, f_d$，其中
$d = \dim(\text{Supp}(M))$。若 $g$ 关于
$(M, f_1, \ldots, f_d)$ 良好，则由引理
\ref{algebra-lemma-good-element} 即得结论。特别地，当 $d = 1$ 时
引理成立。（$d = 0$ 的情形不会出现。）假设 $d > 1$。
选取元素 $h \in R$，使得
(\romannumeral1) $h$ 关于 $(M, f_1, \ldots, f_d)$ 良好，并且
(\romannumeral2) $\dim(\text{Supp}(M) \cap V(h, g)) = d - 2$。
为说明这样的 $h$ 存在，令 $\{\mathfrak q_j\}$ 为以下闭集的
极小素理想所成的（有限）集合：
$\text{Supp}(M)$，
$\text{Supp}(M)\cap V(f_1, \ldots, f_i)$
（$i = 1, \ldots, d - 1$），以及
$\text{Supp}(M) \cap V(g)$。这些 $\mathfrak q_j$ 都不等于
$\mathfrak m$，故由引理 \ref{algebra-lemma-silly}，可以找到
$h \in \mathfrak m$ 使 $h \not \in \mathfrak q_j$。
显然，$h$ 满足 (\romannumeral1) 和 (\romannumeral2)。
由引理 \ref{algebra-lemma-good-element}，$M/hM$ 是 Cohen--Macaulay 模。
由 (\romannumeral2)，二元组 $(M/hM, g)$ 满足归纳假设。
因此，$M/(h, g)M$ 是 Cohen--Macaulay 模，且
$g : M/hM \to M/hM$ 是单射。由引理
\ref{algebra-lemma-permute-xi} 可知，$g : M \to M$ 和
$h : M/gM \to M/gM$ 都是单射。结合 $M/(g, h)M$
是 Cohen--Macaulay 模这一事实，证明完成。
\end{proof}

\begin{proposition}
\label{algebra-proposition-CM-module}
设 $R$ 是极大理想为 $\mathfrak m$ 的 Noether 局部环。
设 $M$ 是 $R$ 上的 Cohen--Macaulay 模，其支撑的维数为 $d$。
假设 $g_1, \ldots, g_c$ 是 $\mathfrak m$ 中的元素，并且
$\dim(\text{Supp}(M/(g_1, \ldots, g_c)M))
= d - c$。则 $g_1, \ldots, g_c$ 是一个 $M$-正则序列，
并且可以延拓为一个极大 $M$-正则序列。
\end{proposition}

\begin{proof}
令 $Z = \text{Supp}(M) \subset \Spec(R)$。由引理
\ref{algebra-lemma-one-equation}，在链
$Z \supset Z \cap V(g_1) \supset \ldots \supset
Z \cap V(g_1, \ldots, g_c)$ 中，每一步至多使维数下降 $1$。
因此，由假设，每一步都恰好使维数下降 $1$。于是，可以依次把
引理 \ref{algebra-lemma-CM-one-g} 应用于模
% P01-E0505: frozen source omits the terminal M in this module quotient.
$M/(g_1, \ldots, g_i)$ 和元素 $g_{i + 1}$。

\medskip\noindent
若 $c < d$，为了用一个元素延拓 $g_1, \ldots, g_c$，只须使用
引理 \ref{algebra-lemma-silly}，选取一个不属于
$Z \cap V(g_1, \ldots, g_c)$ 的有限多个极小素理想中任何一个的
元素 $g_{c + 1} \in \mathfrak m$。
\end{proof}

\noindent
命题 \ref{algebra-proposition-CM-module} 已得证；下面继续发展标准理论。

\begin{lemma}
\label{algebra-lemma-nonzerodivisor-on-CM}
设 $R$ 是极大理想为 $\mathfrak m$ 的 Noether 局部环。
设 $M$ 是有限 $R$-模。设 $x \in \mathfrak m$ 是 $M$ 上的
非零因子。则 $M$ 是 Cohen--Macaulay 模，当且仅当
$M/xM$ 是 Cohen--Macaulay 模。
\end{lemma}

\begin{proof}
由引理 \ref{algebra-lemma-depth-drops-by-one}，有
$\text{depth}(M/xM) = \text{depth}(M)-1$。
由引理 \ref{algebra-lemma-one-equation-module}，有
$\dim(\text{Supp}(M/xM)) = \dim(\text{Supp}(M)) - 1$。
\end{proof}

\begin{lemma}
\label{algebra-lemma-CM-over-quotient}
设 $R \to S$ 是 Noether 局部环的满同态。设 $N$ 是有限
$S$-模。则 $N$ 作为 $S$-模是 Cohen--Macaulay 模，当且仅当
$N$ 作为 $R$-模是 Cohen--Macaulay 模。
\end{lemma}

\begin{proof}
略。
\end{proof}

\begin{lemma}
\label{algebra-lemma-CM-ass-minimal-support}
\begin{reference}
\cite[Chapter 0, Proposition 16.5.4]{EGA}
\end{reference}
设 $R$ 是 Noether 局部环。设 $M$ 是有限 Cohen--Macaulay
$R$-模。若 $\mathfrak p \in \text{Ass}(M)$，则
$\dim(R/\mathfrak p) = \dim(\text{Supp}(M))$，并且
$\mathfrak p$ 是 $M$ 的支撑中的极小素理想。特别地，
$M$ 没有嵌入伴随素理想。
\end{lemma}

\begin{proof}
由引理 \ref{algebra-lemma-depth-dim-associated-primes}，有
$\text{depth}(M) \leq \dim(R/\mathfrak p)$。
当然，由于 $\mathfrak p \in \text{Supp}(M)$
（引理 \ref{algebra-lemma-ass-support}），有
$\dim(R/\mathfrak p) \leq \dim(\text{Supp}(M))$。
因为 $M$ 是 Cohen--Macaulay 模，这两个不等式都取等号。
于是 $\mathfrak p$ 必须在 $\text{Supp}(M)$ 中极小，否则会有
$\dim(R/\mathfrak p) < \dim(\text{Supp}(M))$。
最后，$M$ 的支撑中的极小素理想等于 $\text{Ass}(M)$ 的极小元素
（命题 \ref{algebra-proposition-minimal-primes-associated-primes}），
所以 $M$ 没有嵌入伴随素理想
（定义 \ref{algebra-definition-embedded-primes}）。
\end{proof}

\begin{definition}
\label{algebra-definition-maximal-CM}
设 $R$ 是 Noether 局部环。若有限 $R$-模 $M$ 满足
$\text{depth}(M) = \dim(R)$，则称它为{\it 极大
Cohen--Macaulay 模}。
\end{definition}

\noindent
换言之，Noether 局部环上的极大 Cohen--Macaulay 模就是该环上
深度尽可能大的有限模。等价地，Noether 局部环 $R$ 上的极大
Cohen--Macaulay 模就是维数等于环的维数的 Cohen--Macaulay 模。
特别地，若 $M$ 是满足 $\Spec(R) = \text{Supp}(M)$ 的
Cohen--Macaulay $R$-模，则 $M$ 是极大 Cohen--Macaulay 模。
所以下面的两个引理讨论极大 Cohen--Macaulay 模。

\begin{lemma}
\label{algebra-lemma-maximal-chain-maximal-CM}
\begin{slogan}
在 Cohen--Macaulay 局部环中，任意极大素理想链的长度都等于
环的维数。
\end{slogan}
设 $R$ 是 Noether 局部环。假设存在满足
$\Spec(R) = \text{Supp}(M)$ 的 Cohen--Macaulay 模 $M$。
则任意极大素理想链
$\mathfrak p_0 \subset \mathfrak p_1 \subset \ldots
\subset \mathfrak p_n$ 的长度均为 $n = \dim(R)$。
\end{lemma}

\begin{proof}
对 $\dim(R)$ 归纳证明。若 $\dim(R) = 0$，结论显然。
假设 $\dim(R) > 0$。于是 $n > 0$。选取
$x \in \mathfrak p_1$，使 $x$ 不属于 $R$ 的任何极小素理想，
特别地，$x \not \in \mathfrak p_0$。（见引理
\ref{algebra-lemma-silly}。）由引理 \ref{algebra-lemma-one-equation}，
$\dim(R/xR) = \dim(R) - 1$。由命题
\ref{algebra-proposition-CM-module} 和引理
\ref{algebra-lemma-CM-over-quotient}，$M/xM$ 是 $R/xR$ 上的
Cohen--Macaulay 模。由引理 \ref{algebra-lemma-support-quotient}，
$M/xM$ 的支撑为 $\Spec(R/xR)$。把 $x$ 替换为
$x^n$（对某个 $n$）后，可以假设 $\mathfrak p_1$ 是
$M/xM$ 的伴随素理想，
见引理 \ref{algebra-lemma-inherit-minimal-primes}。由引理
\ref{algebra-lemma-CM-ass-minimal-support}，可知
$\mathfrak p_1/(x)$ 是 $R/xR$ 的极小素理想。因此，链
$\mathfrak p_1/(x) \subset \ldots \subset \mathfrak p_n/(x)$
是 $R/xR$ 中的极大素理想链。由归纳假设，该链的长度为
$\dim(R/xR) = \dim(R) - 1$，正如所需。
\end{proof}

\begin{lemma}
\label{algebra-lemma-dim-formula-maximal-CM}
设 $R$ 是 Noether 局部环。假设存在满足
$\Spec(R) = \text{Supp}(M)$ 的 Cohen--Macaulay 模 $M$。
则对素理想 $\mathfrak p \subset R$，有
$$
\dim(R) = \dim(R_{\mathfrak p}) + \dim(R/\mathfrak p).
$$
\end{lemma}

\begin{proof}
这立即由引理 \ref{algebra-lemma-maximal-chain-maximal-CM} 得出。
\end{proof}

\begin{lemma}
\label{algebra-lemma-localize-CM-module}
设 $R$ 是 Noether 局部环。设 $M$ 是 $R$ 上的
Cohen--Macaulay 模。对任意素理想 $\mathfrak p \subset R$，
模 $M_{\mathfrak p}$ 都是 $R_\mathfrak p$ 上的
Cohen--Macaulay 模。
\end{lemma}

\begin{proof}
可以且确实假设 $\mathfrak p \not = \mathfrak m$ 且 $M$ 非零。
选取极大素理想链
$\mathfrak p = \mathfrak p_c \subset
\mathfrak p_{c - 1} \subset \ldots \subset \mathfrak p_1
\subset \mathfrak m$。若能证明 $M_{\mathfrak p_1}$ 在
$R_{\mathfrak p_1}$ 上的结论，则对 $c$ 归纳即可得到本引理。
因此，可以假设 $\mathfrak p$ 与 $\mathfrak m$ 之间不存在
严格位于二者之间的素理想。注意，
$\dim(\text{Supp}(M_\mathfrak p))
\leq \dim(\text{Supp}(M)) - 1$，因为 $M_\mathfrak p$
的支撑中的任意素理想链都可以在 $M$ 的支撑中再延长一个素理想
（即 $\mathfrak m$）。另一方面，由引理
\ref{algebra-lemma-depth-localization} 以及对 $\mathfrak p$ 的选择，有
$\text{depth}(M_\mathfrak p) \geq
\text{depth}(M) - \dim(R/\mathfrak p) =
\text{depth}(M) - 1$。因此，正如所需，
$\text{depth}(M_\mathfrak p) \geq
\dim(\text{Supp}(M_\mathfrak p))$
（另一个不等式是引理 \ref{algebra-lemma-bound-depth}）。
\end{proof}

\begin{definition}
\label{algebra-definition-module-CM}
设 $R$ 是 Noether 环。设 $M$ 是有限 $R$-模。若对 $R$ 的所有
素理想 $\mathfrak p$，$M_\mathfrak p$ 都是
$R_\mathfrak p$ 上的 Cohen--Macaulay 模，则称 $M$ 是
{\it Cohen--Macaulay 模}。
\end{definition}

\noindent
由引理 \ref{algebra-lemma-localize-CM-module}，只须在 $R$ 的极大理想
处检验这一条件。

\begin{lemma}
\label{algebra-lemma-maximal-CM-polynomial-algebra}
设 $R$ 是 Noether 环。设 $M$ 是 $R$ 上的 Cohen--Macaulay 模。
则 $M \otimes_R R[x_1, \ldots, x_n]$ 是
$R[x_1, \ldots, x_n]$ 上的 Cohen--Macaulay 模。
\end{lemma}

\begin{proof}
对变量个数归纳，只须证明 $R[x]$ 上的
$M[x] = M \otimes_R R[x]$ 的情形。设
$\mathfrak m \subset R[x]$ 是极大理想，并令
$\mathfrak p = R \cap \mathfrak m$。设
$f_1, \ldots, f_d$ 是 $R_{\mathfrak p}$ 的极大理想中的一个
$M_\mathfrak p$-正则序列，其长度为
$d = \dim(\text{Supp}(M_{\mathfrak p}))$。注意，因为 $R[x]$
在 $R$ 上平坦，局部化 $R[x]_{\mathfrak m}$ 在
$R_{\mathfrak p}$ 上平坦。因此，由引理
\ref{algebra-lemma-flat-increases-depth}，序列 $f_1, \ldots, f_d$
是 $R[x]_{\mathfrak m}$ 中一个长度为 $d$ 的
$M[x]_{\mathfrak m}$-正则序列。商
$$
Q = M[x]_{\mathfrak m}/(f_1, \ldots, f_d)M[x]_{\mathfrak m} =
M_{\mathfrak p}/(f_1, \ldots, f_d)M_{\mathfrak p}
\otimes_{R_\mathfrak p} R[x]_{\mathfrak m}
$$
的支撑等于位于 $\mathfrak p$ 上方的素理想，因为
$R_\mathfrak p \to R[x]_\mathfrak m$ 平坦，而
$M_{\mathfrak p}/(f_1, \ldots, f_d)M_{\mathfrak p}$
的支撑等于 $\{\mathfrak p\}$（细节略；提示：由引理
\ref{algebra-lemma-annihilator-flat-base-change} 和
\ref{algebra-lemma-support-closed} 得出）。所以其维数为 $1$。
为完成证明，只须找到 $f \in \mathfrak m$，使它是 $Q$ 上的
非零因子。由于 $\mathfrak m$ 是极大理想，域扩张
$\kappa(\mathfrak m)/\kappa(\mathfrak p)$ 是有限的
（定理 \ref{algebra-theorem-nullstellensatz}）。因此，可以找到
$f \in \mathfrak m$，使它被视为 $x$ 的多项式时，首项系数
不属于 $\mathfrak p$。这样的 $f$ 作用在
$$
M_{\mathfrak p}/(f_1, \ldots, f_d)M_{\mathfrak p} \otimes_R R[x] =
\bigoplus\nolimits_{n \geq 0}
M_{\mathfrak p}/(f_1, \ldots, f_d)M_{\mathfrak p} \cdot x^n
$$
上是非零因子，因而作用在 $Q$ 上也是非零因子。
\end{proof}

\section{Cohen--Macaulay 环}
\label{algebra-section-CM-ring}

\noindent
本节的大多数结果都是第 \ref{algebra-section-CM} 节中结果的特殊情形。

\begin{definition}
\label{algebra-definition-local-ring-CM}
若 Noether 局部环 $R$ 作为自身上的模是 Cohen--Macaulay 模，
则称它为{\it Cohen--Macaulay 环}。
\end{definition}

\noindent
注意，这等价于要求极大理想中存在一个 $R$-正则序列
$x_1, \ldots, x_d$，使得 $R/(x_1, \ldots, x_d)$ 的维数为
$0$。通常只说“正则序列”，而不说“$R$-正则序列”。

\begin{lemma}
\label{algebra-lemma-reformulate-CM}
\begin{slogan}
Cohen--Macaulay 局部环中的正则序列，可以由它截出的对象是否具有
正确的维数来刻画。
\end{slogan}
设 $R$ 是极大理想为 $\mathfrak m $ 的 Noether
Cohen--Macaulay 局部环。设 $x_1, \ldots, x_c \in \mathfrak m$
为若干元素。则
$$
x_1, \ldots, x_c
\text{ is a regular sequence }
\Leftrightarrow
\dim(R/(x_1, \ldots, x_c)) = \dim(R) - c
$$
若这些条件成立，则 $x_1, \ldots, x_c$ 可以延拓为一个长度为
$\dim(R)$ 的正则序列，并且每个商
$R/(x_1, \ldots, x_i)$ 都是维数为 $\dim(R) - i$ 的
Cohen--Macaulay 环。
\end{lemma}

\begin{proof}
这是命题 \ref{algebra-proposition-CM-module} 的特殊情形。
\end{proof}

\begin{lemma}
\label{algebra-lemma-maximal-chain-CM}
设 $R$ 是 Noether 局部环。假设 $R$ 是维数为 $d$ 的
Cohen--Macaulay 环。任何极大
% P01-E0510: frozen source says a maximal chain of ideals; the dimension theorem requires prime ideals.
理想链 $\mathfrak p_0 \subset
\mathfrak p_1 \subset \ldots \subset \mathfrak p_n$
的长度都为 $n = d$。
\end{lemma}

\begin{proof}
这是引理 \ref{algebra-lemma-maximal-chain-maximal-CM} 的特殊情形。
\end{proof}

\begin{lemma}
\label{algebra-lemma-CM-dim-formula}
假设 $R$ 是维数为 $d$ 的 Noether Cohen--Macaulay 局部环。
对任何素理想 $\mathfrak p \subset R$，有
$$
\dim(R) = \dim(R_{\mathfrak p}) + \dim(R/\mathfrak p).
$$
\end{lemma}

\begin{proof}
这立即由引理 \ref{algebra-lemma-maximal-chain-CM} 得出。
（这也是引理 \ref{algebra-lemma-dim-formula-maximal-CM} 的特殊情形。）
\end{proof}

\begin{lemma}
\label{algebra-lemma-localize-CM}
假设 $R$ 是 Cohen--Macaulay 局部环。对任何素理想
$\mathfrak p \subset R$，环 $R_{\mathfrak p}$ 也是
Cohen--Macaulay 环。
\end{lemma}

\begin{proof}
这是引理 \ref{algebra-lemma-localize-CM-module} 的特殊情形。
\end{proof}

\begin{definition}
\label{algebra-definition-ring-CM}
若 Noether 环 $R$ 的所有局部环都是 Cohen--Macaulay 环，
则称它为{\it Cohen--Macaulay 环}。
\end{definition}

\begin{lemma}
\label{algebra-lemma-CM-polynomial-algebra}
假设 $R$ 是 Noether Cohen--Macaulay 环。则 $R$ 上的任意
多项式代数都是 Cohen--Macaulay 环。
\end{lemma}

\begin{proof}
这是引理 \ref{algebra-lemma-maximal-CM-polynomial-algebra} 的特殊情形。
\end{proof}

\begin{lemma}
\label{algebra-lemma-dimension-shift}
设 $R$ 是维数为 $d$ 的 Noether Cohen--Macaulay 局部环。
设 $0 \to K \to R^{\oplus n} \to M \to 0$ 是 $R$-模的正合列。
则或者 $M = 0$，或者 $\text{depth}(K) > \text{depth}(M)$，
或者 $\text{depth}(K) = \text{depth}(M) = d$。
\end{lemma}

\begin{proof}
这是引理 \ref{algebra-lemma-depth-in-ses} 的特殊情形。
\end{proof}

\begin{lemma}
\label{algebra-lemma-mcm-resolution}
设 $R$ 是维数为 $d$ 的 Noether Cohen--Macaulay 局部环。
设 $M$ 是深度为 $e$ 的有限 $R$-模。存在正合复形
% P01-E0511: for e=d the displayed index d-e-1 is -1; the source needs a separate maximal-CM case.
$$
0 \to K \to F_{d-e-1} \to \ldots \to F_0 \to M \to 0
$$
其中每个 $F_i$ 都是有限自由模，而 $K$ 是极大
Cohen--Macaulay 模。
\end{lemma}

\begin{proof}
这立即由定义和引理 \ref{algebra-lemma-dimension-shift} 得出。
\end{proof}

\begin{lemma}
\label{algebra-lemma-find-sequence-image-regular}
设 $\varphi : A \to B$ 是局部环映射。假设 $B$ 是 Noether
Cohen--Macaulay 环，并且
$\mathfrak m_B = \sqrt{\varphi(\mathfrak m_A) B}$。
则 $A$ 中存在元素序列 $f_1, \ldots, f_{\dim(B)}$，使得
$\varphi(f_1), \ldots, \varphi(f_{\dim(B)})$ 是 $B$ 中的
正则序列。
\end{lemma}

\begin{proof}
对 $\dim(B)$ 归纳，只须证明：若 $\dim(B) \geq 1$，则可以找到
$A$ 的一个元素 $f$，使它在 $B$ 中的像为非零因子。由引理
\ref{algebra-lemma-reformulate-CM}，只须找到 $f \in A$，使其在 $B$
中的像不属于 $B$ 的有限多个极小素理想
$\mathfrak q_1, \ldots, \mathfrak q_r$ 中的任何一个。
由假设 $\mathfrak m_B = \sqrt{\varphi(\mathfrak m_A) B}$，
可知 $\mathfrak m_A \not \subset \varphi^{-1}(\mathfrak q_i)$。
因此，由引理 \ref{algebra-lemma-silly} 可以找到 $f$。
\end{proof}

\section{链状环}
\label{algebra-section-catenary}

\noindent
与拓扑学第 \ref{topology-section-catenary-spaces} 节比较。

\begin{definition}
\label{algebra-definition-catenary}
若对任意一对素理想
$\mathfrak p \subset \mathfrak q$，都存在一个整数作为所有有限素理想链
$\mathfrak p = \mathfrak p_0 \subset \mathfrak p_1 \subset \ldots \subset
\mathfrak p_e = \mathfrak q$ 的长度上界，并且所有这样的极大链都具有相同
长度，则称环 $R$ 是{\it 链状的}。
\end{definition}

\begin{lemma}
\label{algebra-lemma-catenary}
环 $R$ 是链状的，当且仅当拓扑空间 $\Spec(R)$ 是链状的（参见
拓扑学定义 \ref{topology-definition-catenary}）。
\end{lemma}

\begin{proof}
这立即由定义以及引理 \ref{algebra-lemma-irreducible} 中对不可约闭子集的刻画得出。
\end{proof}

\noindent
一般而言，即使 $R$ 是链状的，有限生成 $R$-代数也未必是链状的。
因此我们作如下定义。

\begin{definition}
\label{algebra-definition-universally-catenary}
若 Noether 环 $R$ 的每个有限型 $R$-代数都是链状的，则称该环
是{\it 普遍链状的}。
\end{definition}

\noindent
我们将范围限制在 Noether 环上，因为尚不清楚对非 Noether 环而言，
这个定义是否恰当。由引理 \ref{algebra-lemma-quotient-catenary}，为了检验
Noether 环 $R$ 是否普遍链状，只须检验每个多项式代数
$R[x_1, \ldots, x_n]$ 是否链状。

\begin{lemma}
\label{algebra-lemma-localization-catenary}
链状环的任意局部化都是链状环。Noether 普遍链状环的任意局部化
都是普遍链状环。
\end{lemma}

\begin{proof}
设 $A$ 是环，$S \subset A$ 是乘法子集。引理
\ref{algebra-lemma-spec-localization} 中对 $\Spec(S^{-1}A)$ 的描述表明，
若 $A$ 是链状的，则 $S^{-1}A$ 也是链状的。若
$S^{-1}A \to C$ 是有限型的，则存在有限型环映射 $A \to B$，使得
$C = S^{-1}B$。因此，若 $A$ 是 Noether 且普遍链状的，则 $B$
是链状的，从而 $C$ 也是链状的。这就证明了该引理。
\end{proof}

\begin{lemma}
\label{algebra-lemma-universally-catenary}
设 $A$ 是 Noether 普遍链状环。任意在 $A$ 上本质有限型的
$A$-代数都是普遍链状环。
\end{lemma}

\begin{proof}
若 $B$ 是有限型 $A$-代数，则由引理 \ref{algebra-lemma-Noetherian-permanence}，
$B$ 是 Noether 环。任意有限型 $B$-代数都是有限型 $A$-代数，因而由
$A$ 普遍链状这一假设，它是链状的。因此 $B$ 是普遍链状的。由引理
\ref{algebra-lemma-localization-catenary}，$B$ 的任意局部化都是普遍链状的，
这就完成了证明。
\end{proof}

\begin{lemma}
\label{algebra-lemma-catenary-check-local}
设 $R$ 是环。以下条件等价：
\begin{enumerate}
\item $R$ 是链状的；
\item 对每个素理想 $\mathfrak p$，$R_\mathfrak p$ 都是链状的；
\item 对每个极大理想 $\mathfrak m$，$R_\mathfrak m$ 都是链状的。
\end{enumerate}
假设 $R$ 是 Noether 环。以下条件等价：
\begin{enumerate}
\item $R$ 是普遍链状的；
\item 对每个素理想 $\mathfrak p$，$R_\mathfrak p$ 都是普遍链状的；
\item 对每个极大理想 $\mathfrak m$，$R_\mathfrak m$ 都是普遍链状的。
\end{enumerate}
\end{lemma}

\begin{proof}
在两种情形中，(1) $\Rightarrow$ (2) 都由引理
\ref{algebra-lemma-localization-catenary} 得出，而 (2) $\Rightarrow$ (3)
都是立即的。假设对 $R$ 的每个极大理想 $\mathfrak m$，
$R_\mathfrak m$ 都是链状的。若 $\mathfrak p \subset \mathfrak q$
是 $R$ 的素理想，则选择一个满足 $\mathfrak q \subset \mathfrak m$
的极大理想。$\mathfrak p$ 与 $\mathfrak q$ 之间的素理想链，和
% P01-E0512: the source says “Chains of primes ideals”; Chinese uses the unambiguous intended “prime ideals”.
$\mathfrak pR_\mathfrak m$ 与 $\mathfrak qR_\mathfrak m$ 之间的素理想链
一一对应，故 $R$ 是链状的。现在假设 $R$ 是 Noether 环，并且对
$R$ 的每个极大理想 $\mathfrak m$，$R_\mathfrak m$ 都是普遍链状的。
设 $R \to S$ 是有限型环映射。设 $\mathfrak q$ 是 $S$ 的素理想，
它位于素理想 $\mathfrak p \subset R$ 之上。在 $R$ 中选择一个满足
$\mathfrak p \subset \mathfrak m$ 的极大理想。于是
$R_\mathfrak p$ 是 $R_\mathfrak m$ 的局部化，因而由引理
\ref{algebra-lemma-localization-catenary}，它是普遍链状的。于是
$S_\mathfrak p$ 作为 $R_\mathfrak p$ 上的有限型环是链状的，
而 $S_\mathfrak q$ 作为它的局部化也是链状的。因此，由上面处理的
第一种情形，$S$ 是链状的。
\end{proof}

\begin{lemma}
\label{algebra-lemma-quotient-catenary}
链状环的任意商环都是链状环。Noether 普遍链状环的任意商环都是
普遍链状环。
\end{lemma}

\begin{proof}
设 $A$ 是环，$I \subset A$ 是理想。引理 \ref{algebra-lemma-spec-closed}
中对 $\Spec(A/I)$ 的描述表明，若 $A$ 是链状的，则 $A/I$
也是链状的。第二个断言是引理 \ref{algebra-lemma-universally-catenary}
的特殊情形。
\end{proof}

\begin{lemma}
\label{algebra-lemma-catenary-check-irreducible}
设 $R$ 是 Noether 环。
\begin{enumerate}
\item $R$ 是链状的，当且仅当对每个极小素理想 $\mathfrak p$，
$R/\mathfrak p$ 都是链状的。
\item $R$ 是普遍链状的，当且仅当对每个极小素理想 $\mathfrak p$，
$R/\mathfrak p$ 都是普遍链状的。
\end{enumerate}
\end{lemma}

\begin{proof}
若 $\mathfrak a \subset \mathfrak b$ 是 $R$ 的素理想包含关系，
则可以找到一个满足 $\mathfrak p \subset \mathfrak a$ 的极小素理想，
第一个断言于是显然。第二个断言的证明从略。
\end{proof}

\begin{lemma}
\label{algebra-lemma-CM-ring-catenary}
Noether Cohen--Macaulay 环是普遍链状的。更一般地，若 $R$ 是
Noether 环，而 $M$ 是满足 $\text{Supp}(M) = \Spec(R)$ 的
Cohen--Macaulay $R$-模，则 $R$ 是普遍链状的。
\end{lemma}

\begin{proof}
由引理 \ref{algebra-lemma-CM-polynomial-algebra}，$R$ 上的多项式代数是
Cohen--Macaulay 环，所以只须证明 Cohen--Macaulay 环是链状的。
设 $R$ 是 Cohen--Macaulay 环，$\mathfrak p \subset \mathfrak q$
是 $R$ 的素理想。由定义，$R_{\mathfrak q}$ 和 $R_{\mathfrak p}$
都是 Cohen--Macaulay 环。取一条极大素理想链
$\mathfrak p = \mathfrak p_0 \subset \mathfrak p_1 \subset
\ldots \subset \mathfrak p_n = \mathfrak q$。再选择一条极大素理想链
$\mathfrak q_0 \subset \mathfrak q_1 \subset \ldots \subset
\mathfrak q_m = \mathfrak p$。由引理 \ref{algebra-lemma-maximal-chain-CM}，
有 $n + m = \dim(R_{\mathfrak q})$。由同一引理，还有
$m = \dim(R_{\mathfrak p})$。因此
$n = \dim(R_{\mathfrak q}) - \dim(R_{\mathfrak p})$ 与选择无关。

\medskip\noindent
为了证明更一般的断言，完全照上面论证，但改用引理
\ref{algebra-lemma-maximal-CM-polynomial-algebra} 和
\ref{algebra-lemma-maximal-chain-maximal-CM}。
\end{proof}

\begin{lemma}
\label{algebra-lemma-catenary-Noetherian-local}
设 $(A, \mathfrak m)$ 是 Noether 局部环。以下条件等价：
\begin{enumerate}
\item $A$ 是链状的；
\item $\mathfrak p \mapsto \dim(A/\mathfrak p)$ 是 $\Spec(A)$
上的维数函数。
\end{enumerate}
\end{lemma}

\begin{proof}
若 $A$ 是链状的，则由拓扑学引理
\ref{topology-lemma-locally-dimension-function}（以及引理
\ref{algebra-lemma-catenary}），$\Spec(A)$ 有一个维数函数 $\delta$。
可以假设 $\delta(\mathfrak m) = 0$。于是由拓扑学引理
\ref{topology-lemma-dimension-function-catenary}，可知
$$
\delta(\mathfrak p) = \text{codim}(V(\mathfrak m), V(\mathfrak p)) =
\dim(A/\mathfrak p)
$$
由此可见 (1) 推出 (2)。反向蕴含也由拓扑学引理
\ref{topology-lemma-dimension-function-catenary} 得出。
\end{proof}

\section{正则局部环}
\label{algebra-section-regular}

\noindent
正则局部环定义于定义 \ref{algebra-definition-regular-local}。要证明正则局部环的
所有素局部化都是正则的并不容易。事实上，迄今发展的大量材料都是为
证明这一事实而准备的。参见命题
\ref{algebra-proposition-finite-gl-dim-regular}，并沿其中的引用向前追溯。

\begin{lemma}
\label{algebra-lemma-regular-graded}
设 $(R, \mathfrak m, \kappa)$ 是维数为 $d$ 的正则局部环。
分次环 $\bigoplus \mathfrak m^n / \mathfrak m^{n + 1}$ 同构于
分次多项式代数 $\kappa[X_1, \ldots, X_d]$。
\end{lemma}

\begin{proof}
设 $x_1, \ldots, x_d$ 是极大理想 $\mathfrak m$ 的极小生成集，
参见定义 \ref{algebra-definition-regular-local}。存在满射
$\kappa[X_1, \ldots, X_d]
\to \bigoplus \mathfrak m^n/\mathfrak m^{n + 1}$，它将 $X_i$
映到 $x_i$ 在 $\mathfrak m/\mathfrak m^2$ 中的类。由命题
\ref{algebra-proposition-dimension} 有 $d(R) = d$，故数值多项式
$n \mapsto \dim_\kappa \mathfrak m^n/\mathfrak m^{n + 1}$
的次数为 $d - 1$。由引理 \ref{algebra-lemma-quotient-smaller-d} 可知，
该满射 $\kappa[X_1, \ldots, X_d]
\to \bigoplus \mathfrak m^n/\mathfrak m^{n + 1}$ 是同构。
\end{proof}

\begin{lemma}
\label{algebra-lemma-regular-domain}
任意正则局部环都是整环。
\end{lemma}

\begin{proof}
我们将使用引理 \ref{algebra-lemma-intersect-powers-ideal-module-zero} 给出的
$\bigcap \mathfrak m^n = 0$。设 $f, g \in R$ 且 $fg = 0$。
% P01-E0513: the source implicitly needs f and g nonzero before choosing maximal m-adic orders.
只须考虑二者都非零的情形。假设 $f \in \mathfrak m^a$ 且
$g \in \mathfrak m^b$，其中 $a, b$ 均为最大的。由于
$fg = 0 \in \mathfrak m^{a + b + 1}$，由引理
\ref{algebra-lemma-regular-graded} 的结论可知，或者
$f \in \mathfrak m^{a + 1}$，或者 $g \in \mathfrak m^{b + 1}$。
矛盾。
\end{proof}

\begin{lemma}
\label{algebra-lemma-regular-ring-CM}
设 $R$ 是正则局部环，$x_1, \ldots, x_d$ 是极大理想
$\mathfrak m$ 的极小生成集。则 $x_1, \ldots, x_d$ 是正则序列，
并且每个 $R/(x_1, \ldots, x_c)$ 都是维数为 $d - c$ 的正则局部环。
特别地，$R$ 是 Cohen--Macaulay 环。
\end{lemma}

\begin{proof}
注意，由引理 \ref{algebra-lemma-one-equation}，$R/x_1R$ 是维数
$\geq d - 1$ 的 Noether 局部环，而 $x_2, \ldots, x_d$ 生成其
极大理想。因此，按定义它是正则局部环。由引理
\ref{algebra-lemma-regular-domain}，$R$ 是整环，故 $x_1$ 是非零因子。
\end{proof}

\begin{lemma}
\label{algebra-lemma-regular-quotient-regular}
设 $R$ 是正则局部环。设 $I \subset R$ 是理想，并且 $R/I$
也是正则局部环。则存在 $R$ 的极大理想 $\mathfrak m$ 的极小生成集
$x_1, \ldots, x_d$，使得对某个 $0 \leq c \leq d$，有
$I = (x_1, \ldots, x_c)$。
\end{lemma}

\begin{proof}
记 $\dim(R) = d$ 且 $\dim(R/I) = d - c$。用
$\overline{\mathfrak m} = \mathfrak m/I$ 表示 $R/I$ 的极大理想。
设 $\kappa = R/\mathfrak m$。由正则局部环的定义，有
% P01-E0514: the source omits the residue-field subscript on the third vector-space dimension.
$$
\dim_\kappa((I + \mathfrak m^2)/\mathfrak m^2) =
\dim_\kappa(\mathfrak m/\mathfrak m^2)
- \dim(\overline{\mathfrak m}/\overline{\mathfrak m}^2) = d - (d - c) = c
$$
因此可以选择 $x_1, \ldots, x_c \in I$，使它们在
$\mathfrak m/\mathfrak m^2$ 中的像线性无关，再补上
$x_{c + 1}, \ldots, x_d$，得到 $\mathfrak m$ 的极小生成系。
诱导映射 $R/(x_1, \ldots, x_c) \to R/I$ 是相同维数的正则局部环
之间的满射（引理 \ref{algebra-lemma-regular-ring-CM}）。因此其核为零，
也就是说，$I = (x_1, \ldots, x_c)$。具体地，若非如此，则由引理
\ref{algebra-lemma-regular-domain} 和 \ref{algebra-lemma-one-equation}，将有
$\dim(R/I) < \dim(R/(x_1, \ldots, x_c))$。
\end{proof}

\begin{lemma}
\label{algebra-lemma-free-mod-x}
设 $R$ 是 Noether 局部环，$x \in \mathfrak m$。设 $M$ 是有限
$R$-模，使得 $x$ 在 $M$ 上是非零因子，并且 $M/xM$ 作为
$R/xR$-模是自由的。则 $M$ 作为 $R$-模是自由的。
\end{lemma}

\begin{proof}
设 $m_1, \ldots, m_r$ 是 $M$ 的元素，它们的像构成 $M/xM$ 的一组
$R/xR$-基。由 Nakayama 引理 \ref{algebra-lemma-NAK}，
$m_1, \ldots, m_r$ 生成 $M$。若 $\sum a_i m_i = 0$ 是一个关系，
则对所有 $i$ 都有 $a_i \in xR$。因此，对某些 $b_i \in R$，
有 $a_i = b_i x$。于是 $R^r \to M$ 的核 $K$ 满足 $xK = K$，
所以由 Nakayama 引理，它为零。
\end{proof}

\begin{lemma}
\label{algebra-lemma-regular-mcm-free}
设 $R$ 是正则局部环。任意极大 Cohen--Macaulay $R$-模都是自由的。
\end{lemma}

\begin{proof}
设 $M$ 是极大 Cohen--Macaulay $R$-模。设 $x \in \mathfrak m$
是生成 $\mathfrak m$ 的一个正则序列的一部分。由命题
\ref{algebra-proposition-CM-module}，$x$ 在 $M$ 上是非零因子，而
$M/xM$ 是 $R/xR$ 上的极大 Cohen--Macaulay 模。对
$\dim(R)$ 归纳可知 $M/xM$ 是自由的。由引理
\ref{algebra-lemma-free-mod-x} 即得结论。
\end{proof}

\begin{lemma}
\label{algebra-lemma-regular-mod-x}
假设 $R$ 是 Noether 局部环。设 $x \in \mathfrak m$ 是非零因子，
并且 $R/xR$ 是正则局部环。则 $R$ 是正则局部环。更一般地，若
$x_1, \ldots, x_r$ 是 $R$ 中的正则序列，并且
$R/(x_1, \ldots, x_r)$ 是正则局部环，则 $R$ 是正则局部环。
\end{lemma}

\begin{proof}
这是因为把 $x$ 与 $R/xR$ 的极大理想的一组极小生成元的提升放在一起，
就会得到 $\mathfrak m$ 的 $\dim(R)$ 个生成元。使用引理
\ref{algebra-lemma-one-equation}。最后一个断言由第一个断言和归纳法得出。
\end{proof}

\begin{lemma}
\label{algebra-lemma-colimit-regular}
设 $(R_i, \varphi_{ii'})$ 是局部环的有向系，其过渡映射都是局部环映射。
若每个 $R_i$ 都是正则局部环，并且 $R = \colim R_i$ 是 Noether 环，
则 $R$ 是正则局部环。
\end{lemma}

\begin{proof}
设 $\mathfrak m \subset R$ 是极大理想；它是各极大理想
$\mathfrak m_i \subset R_i$ 的余极限。我们对
% P01-E0516: the source omits the residue-field subscript and grouping in this tangent-space dimension.
$d = \dim \mathfrak m/\mathfrak m^2$ 归纳证明该引理。若 $d = 0$，
则 $R = R/\mathfrak m$ 是域，且 $R$ 是正则局部环。若 $d > 0$，
选择 $x \in \mathfrak m$ 且 $x \not \in \mathfrak m^2$。对某个 $i$，
可以找到映到 $x$ 的 $x_i \in \mathfrak m_i$。注意
$R/xR = \colim_{i' \geq i} R_{i'}/x_iR_{i'}$ 是 Noether 局部环。
由引理 \ref{algebra-lemma-regular-ring-CM}，可知
$R_{i'}/x_iR_{i'}$ 是正则局部环。因此由归纳可知 $R/xR$
是正则局部环。由于每个 $R_i$ 都是整环
% P01-E0515: the cited lemma proves the graded-ring statement; regular-domain is the direct domain lemma.
（引理 \ref{algebra-lemma-regular-graded}），可知 $R$ 是整环。因此 $x$
是非零因子，由引理 \ref{algebra-lemma-regular-mod-x} 可知 $R$ 是正则局部环。
\end{proof}

\section{环的满态射}
\label{algebra-section-epimorphism}

\noindent
在任意范畴中都有{\it 满态射}的概念。本节部分材料取自
\cite{Autour} 和 \cite{Mazet}。

\begin{lemma}
\label{algebra-lemma-epimorphism}
设 $R \to S$ 是环映射。以下条件等价：
\begin{enumerate}
\item $R \to S$ 是满态射；
\item 两个环映射 $S \to S \otimes_R S$ 相等；
\item 任一个环映射 $S \to S \otimes_R S$ 都是同构；
\item 环映射 $S \otimes_R S \to S$ 是同构。
\end{enumerate}
\end{lemma}

\begin{proof}
从略。
\end{proof}

\begin{lemma}
\label{algebra-lemma-composition-epimorphism}
两个环满态射的复合仍是满态射。
\end{lemma}

\begin{proof}
从略。提示：这在任意范畴中都成立。
\end{proof}

\begin{lemma}
\label{algebra-lemma-base-change-epimorphism}
若 $R \to S$ 是环满态射，而 $R \to R'$ 是任意环映射，则
$R' \to R' \otimes_R S$ 是满态射。
\end{lemma}

\begin{proof}
从略。提示：这在任意具有推出的范畴中都成立。
\end{proof}

\begin{lemma}
\label{algebra-lemma-permanence-epimorphism}
若 $A \to B \to C$ 是环映射，而 $A \to C$ 是满态射，则
$B \to C$ 也是满态射。
\end{lemma}

\begin{proof}
从略。提示：这在任意范畴中都成立。
\end{proof}

\noindent
特别地，这意味着：若 $R \to S$ 是满态射，其像为
$\overline{R} \subset S$，则 $\overline{R} \to S$ 是满态射。
% P01-E0517: the source has a spurious comma in “in particular, that”.
因此在证明关于满态射的结果时，常可假设映射是单射。下面的引理特别
蕴含每个局部化都是满态射。

\begin{lemma}
\label{algebra-lemma-epimorphism-local}
设 $R \to S$ 是环映射。以下条件等价：
\begin{enumerate}
\item $R \to S$ 是满态射；
\item 对 $R$ 的每个素理想 $\mathfrak p$，
$R_{\mathfrak p} \to S_{\mathfrak p}$ 都是满态射。
\end{enumerate}
\end{lemma}

\begin{proof}
由于 $S_{\mathfrak p} = R_{\mathfrak p} \otimes_R S$（参见引理
\ref{algebra-lemma-tensor-localization}），由引理
\ref{algebra-lemma-base-change-epimorphism} 可知 (1) 推出 (2)。反过来，
假设 (2) 成立。设 $a, b : S \to A$ 是从 $S$ 到环 $A$ 的两个环映射，
并且它们与 $R \to S$ 复合后相等。由假设可知，对 $R$ 的每个素理想
% P01-E0518: the source phrase “two ring maps ... equalizing the map” is category-theoretically malformed.
$\mathfrak p$，诱导映射
$a_{\mathfrak p}, b_{\mathfrak p} : S_{\mathfrak p} \to A_{\mathfrak p}$
相同。由于 $A \subset \prod_{\mathfrak p} A_{\mathfrak p}$，故
$a = b$，参见引理 \ref{algebra-lemma-characterize-zero-local}。
\end{proof}

\begin{lemma}
\label{algebra-lemma-finite-epimorphism-surjective}
\begin{slogan}
一个环映射是满射，当且仅当它是有限满态射。
\end{slogan}
设 $R \to S$ 是环映射。以下条件等价：
\begin{enumerate}
\item $R \to S$ 是满态射且有限；
\item $R \to S$ 是满射。
\end{enumerate}
\end{lemma}

\begin{proof}
（这个引理似乎已在文献中被多次重新证明，并且有许多不同的证明。）
满射环映射显然是满态射。假设 $R \to S$ 是有限环映射，并且
$S \otimes_R S \to S$ 是同构。我们的目标是证明 $R \to S$ 是满射。
假设 $S/R$ 非零。正合列 $R \to S \to S/R \to 0$ 导出正合列
$$
R \otimes_R S \to S \otimes_R S \to S/R \otimes_R S \to 0.
$$
我们的假设蕴含第一个箭头是同构，故 $S/R \otimes_R S = 0$。
从而还有 $S/R \otimes_R S/R = 0$。由引理
\ref{algebra-lemma-trivial-filter-finite-module}，对某个真理想 $I \subset R$，
存在 $R$-模满射 $S/R \to R/I$。于是存在满射
$S/R \otimes_R S/R \to R/I \otimes_R R/I = R/I \not = 0$，
矛盾。
\end{proof}

\begin{lemma}
\label{algebra-lemma-faithfully-flat-epimorphism}
忠实平坦满态射是同构。
\end{lemma}

\begin{proof}
这由引理 \ref{algebra-lemma-epimorphism} 的第 (3) 部分立即得出，因为映射
$S \to S \otimes_R S$ 是映射 $R \to S$ 与 $S$ 作张量积所得的映射。
\end{proof}

\begin{lemma}
\label{algebra-lemma-epimorphism-over-field}
若 $k \to S$ 是满态射，且 $k$ 是域，则 $S = k$ 或 $S = 0$。
\end{lemma}

\begin{proof}
这由引理 \ref{algebra-lemma-faithfully-flat-epimorphism} 的结论立即得出
（因为 $k$ 上任意非零代数都是忠实平坦的）；也可以直接论证：
% P01-E0519: the source switches to undefined R; the intended map is S -> S tensor_k S.
$R \to R \otimes_k R$ 不可能是满射，除非 $\dim_k(R) \leq 1$。
\end{proof}

\begin{lemma}
\label{algebra-lemma-epimorphism-injective-spec}
设 $R \to S$ 是环满态射。则
\begin{enumerate}
\item $\Spec(S) \to \Spec(R)$ 是单射；
\item 对位于 $\mathfrak p \subset R$ 之上的 $\mathfrak q \subset S$，
有 $\kappa(\mathfrak p) = \kappa(\mathfrak q)$。
\end{enumerate}
\end{lemma}

\begin{proof}
设 $\mathfrak p$ 是 $R$ 的素理想。该映射的纤维是纤维环
$S \otimes_R \kappa(\mathfrak p)$ 的谱。由引理
\ref{algebra-lemma-base-change-epimorphism}，映射
$\kappa(\mathfrak p) \to S \otimes_R \kappa(\mathfrak p)$ 是满态射，
所以由引理 \ref{algebra-lemma-epimorphism-over-field}，或者
$S \otimes_R \kappa(\mathfrak p) = 0$，或者
$S \otimes_R \kappa(\mathfrak p) = \kappa(\mathfrak p)$。
这证明了 (1) 和 (2)。
\end{proof}

\begin{lemma}
\label{algebra-lemma-relations}
设 $R$ 是环。设 $M$、$N$ 是 $R$-模。设
$\{x_i\}_{i \in I}$ 是 $M$ 的一组生成元，
$\{y_j\}_{j \in J}$ 是 $N$ 的一组生成元。设
$\{m_j\}_{j \in J}$ 是 $M$ 中的一族元素，其中除有限多个 $j$ 外，
均有 $m_j = 0$。则
$$
\sum\nolimits_{j \in J} m_j \otimes y_j = 0 \text{ 于 } M \otimes_R N
$$
等价于下述条件：存在 $a_{i, j} \in R$，并且除有限多对
$(i, j)$ 外均有 $a_{i, j} = 0$，使得
\begin{align*}
m_j & = \sum\nolimits_{i \in I} a_{i, j} x_i \quad\text{对所有 } j \in J, \\
0 & = \sum\nolimits_{j \in J} a_{i, j} y_j \quad\text{对所有 } i \in I.
\end{align*}
\end{lemma}

\begin{proof}
充分性是立即的。假设 $\sum_{j \in J} m_j \otimes y_j = 0$。
考虑短正合列
$$
0 \to K \to \bigoplus\nolimits_{j \in J} R \to N \to 0
$$
其中 $\bigoplus\nolimits_{j \in J} R$ 的第 $j$ 个基向量映到 $y_j$。
将它与 $M$ 作张量积，得到正合列
$$
K \otimes_R M \to \bigoplus\nolimits_{j \in J} M \to N \otimes_R M \to 0.
$$
假设蕴含存在元素 $k_i \in K$，使得 $\sum k_i \otimes x_i$ 映到
中间项的元素 $(m_j)_{j \in J}$。写成
% P01-E0520: the source says “Writing ... and we obtain”, leaving the sentence grammatically broken.
$k_i = (a_{i, j})_{j \in J}$，即得所需结论。
\end{proof}

\begin{lemma}
\label{algebra-lemma-kernel-difference-projections}
设 $\varphi : R \to S$ 是环映射，$g \in S$。以下条件等价：
\begin{enumerate}
\item 在 $S \otimes_R S$ 中有 $g \otimes 1 = 1 \otimes g$；
\item 存在 $n \geq 0$、元素 $y_i, z_j \in S$ 以及
$x_{i, j} \in R$（其中 $1 \leq i, j \leq n$），使得
\begin{enumerate}
\item $g = \sum_{i, j \leq n} x_{i, j} y_i z_j$；
\item 对每个 $j$，有 $\sum x_{i, j}y_i \in \varphi(R)$；
\item 对每个 $i$，有 $\sum x_{i, j}z_j \in \varphi(R)$。
\end{enumerate}
\end{enumerate}
\end{lemma}

\begin{proof}
(2) 推出 (1) 是显然的。反过来，假设
$g \otimes 1 = 1 \otimes g$。选择 $S$ 作为 $R$-模的一组生成元
$\{s_i\}_{i \in I}$，使得 $0, 1 \in I$，且 $s_0 = 1$、$s_1 = g$。
将引理 \ref{algebra-lemma-relations} 应用于关系
$g \otimes s_0 + (-1) \otimes s_1 = 0$。可知存在
$a_{i, j} \in R$，使得 $g = \sum_i a_{i, 0} s_i$、
$-1 = \sum_i a_{i, 1} s_i$；当 $j \not = 0, 1$ 时，有
$0 = \sum_i a_{i, j} s_i$；并且对所有 $i$，有
$\sum_j a_{i, j}s_j = 0$。于是
$$
\sum\nolimits_{i, j \not = 0} a_{i, j} s_i s_j = -g + a_{0, 0}
$$
且对每个 $j \not = 0$，有
% P01-E0521: the source writes R here although the lemma consistently uses the image varphi(R).
$\sum_{i \not = 0} a_{i, j}s_i \in R$。这证明了
$-g + a_{0, 0}$ 可写成 (2) 中的形式，从而 $g$ 也可写成 (2)
中的形式。细节从略。提示：证明 $S$ 中具有 (2) 所述表达式的元素
构成 $S$ 的一个 $R$-子代数。
\end{proof}

\begin{remark}
\label{algebra-remark-matrices-associated-to-elements-epicenter}
设 $R \to S$ 是环映射。有时把满足
$g \in S$ 且 $g \otimes 1 = 1 \otimes g$ 的元素集合称为 $S$ 的
{\it 满态中心（epicenter）}。它是一个 $R$-代数。由引理
\ref{algebra-lemma-kernel-difference-projections} 的构造，对满态中心中的每个
$g$，得到矩阵分解
$$
(g) = Y X Z
$$
其中 $X \in \text{Mat}(n \times n, R)$、
$Y \in \text{Mat}(1 \times n, S)$，且
$Z \in \text{Mat}(n \times 1, S)$。具体地，设 $x_{i, j}, y_i, z_j$
如该引理第 (2) 部分所述。令 $X = (x_{i, j})$，令 $y$ 为以 $y_i$
为分量的行向量，令 $z$ 为以 $z_j$ 为分量的列向量。
% P01-E0522: the source defines lowercase y,z but immediately uses uppercase Y,Z in the factorization.
在此记号下，引理 \ref{algebra-lemma-kernel-difference-projections} 的条件
(b) 和 (c) 恰好表示
$Y X \in \text{Mat}(1 \times n, R)$、
$X Z \in \text{Mat}(n \times 1, R)$。考虑矩阵三元组
$(X, YX, XZ)$ 十分方便。给定 $n \in \mathbf{N}$ 和三元组
$(P, U, V)$，若存在如上的矩阵分解，使得
$P = X$、$U = YX$ 且 $V = XZ$，则称 $(P, U, V)$ 是
{\it 与 $g$ 关联的 $n$-三元组}。
\end{remark}

\begin{lemma}
\label{algebra-lemma-epimorphism-cardinality}
设 $R \to S$ 是环满态射。则 $S$ 的基数至多为 $R$ 的基数。
用公式表示：$|S| \leq |R|$。
\end{lemma}

\begin{proof}
$R \to S$ 是满态射这一条件意味着每个 $g \in S$ 都满足
$g \otimes 1 = 1 \otimes g$，参见引理 \ref{algebra-lemma-epimorphism}。
我们将使用评注 \ref{algebra-remark-matrices-associated-to-elements-epicenter}
中引入的记号。假设 $g, g' \in S$，并且 $(P, U, V)$ 是同时与
$g$ 和 $g'$ 关联的 $n$-三元组。我们断言 $g = g'$。事实上，
把与 $g$ 的矩阵分解 $(g) = YXZ$ 写成
$(P, U, V) = (X, YX, XZ)$，把与 $g'$ 的矩阵分解
$(g') = Y'X'Z'$ 写成 $(P, U, V) = (X', Y'X', X'Z')$。
于是
$$
(g) = YXZ = UZ = Y'X'Z = Y'PZ = Y'XZ = Y'V = Y'X'Z' = (g')
$$
故 $g = g'$。这说明 $S$ 的基数被所有可能三元组的数目所控制，
后者的基数至多为 $\sup_{n \in \mathbf{N}} |R|^n$。若 $R$ 是无限的，
则它至多为 $|R|$，参见 \cite[Ch. I, 10.13]{Kunen}。

\medskip\noindent
若 $R$ 是有限环，则上述论证只证明 $S$ 至多可数。事实上，在这种情形
$R$ 是 Artin 环，并且映射 $R \to S$ 是满射。我们略去这一情形的证明。
\end{proof}

\begin{lemma}
\label{algebra-lemma-epimorphism-modules}
对环同态 $R \to S$，以下条件等价：
\begin{enumerate}
\item $R \to S$ 是环满态射；
\item 对任意 $S$-模 $N_1, N_2$，都有
$\Hom_S(N_1, N_2) = \Hom_R(N_1, N_2)$；
\item 限制函子 $\text{Mod}_S \to \text{Mod}_R$ 是全忠实的。
\end{enumerate}
\end{lemma}

\begin{proof}
按定义，(2) 与 (3) 等价。

\medskip\noindent
假设 (1) 成立，设 $N_1, N_2$ 是 $S$-模，并设
$\varphi : N_1 \to N_2$ 是 $R$-线性映射。对任意 $x \in N_1$，
考虑由规则 $g \otimes g' \mapsto g\varphi(g'x)$ 定义的映射
$S \otimes_R S \to N_2$。由于两个映射 $S \to S \otimes_R S$
都是同构（引理 \ref{algebra-lemma-epimorphism}），可得
$g \varphi(g'x) = gg'\varphi(x) = \varphi(gg' x)$。因此 $\varphi$
是 $S$-线性的。

\medskip\noindent
假设 (2) 成立。令 $N_1 = S \otimes_R S$，通过左 $S$-模作用把它
看作 $S$-模；令 $N_2 = S \otimes_R S$，赋予右 $S$-模作用。
由于作为 $R$-模有 $N_1 = N_2$，由 (2) 可知
$S \otimes_R S$ 上的两个 $S$-模结构相同。由引理
\ref{algebra-lemma-epimorphism}，这蕴含 (1)。
\end{proof}

\section{纯理想}
\label{algebra-section-pure-ideals}

\noindent
许多论文讨论了本节材料，例如参见 \cite{Lazard}、\cite{Bkouche}
和 \cite{DeMarco}。

\begin{definition}
\label{algebra-definition-pure-ideal}
设 $R$ 是环。若商环 $R/I$ 作为 $R$-模是平坦的，则称
$I \subset R$ 是{\it 纯的}。
\end{definition}

\begin{lemma}
\label{algebra-lemma-pure}
设 $R$ 是环，$I \subset R$ 是理想。以下条件等价：
\begin{enumerate}
\item $I$ 是纯的；
\item 对每个理想 $J \subset R$，都有 $J \cap I = IJ$；
\item 对每个有限生成理想 $J \subset R$，都有 $J \cap I = JI$；
\item 对每个 $x \in R$，都有 $(x) \cap I = xI$；
\item 对每个 $x \in I$，存在某个 $y \in I$，使得 $x = yx$；
\item 对每个 $x_1, \ldots, x_n \in I$，存在 $y \in I$，使得对所有
$i = 1, \ldots, n$ 都有 $x_i = yx_i$；
\item 对 $R$ 的每个素理想 $\mathfrak p$，有
$IR_{\mathfrak p} = 0$ 或 $IR_{\mathfrak p} = R_{\mathfrak p}$；
\item $\text{Supp}(I) = \Spec(R) \setminus V(I)$；
\item $I$ 是映射 $R \to (1 + I)^{-1}R$ 的核；
\item 对 $R$ 的某个乘法子集 $S$，作为 $R$-代数有
$R/I \cong S^{-1}R$；
\item 作为 $R$-代数有 $R/I \cong (1 + I)^{-1}R$。
\end{enumerate}
\end{lemma}

\begin{proof}
对 $R$ 的任意理想 $J$，都有短正合列
$0 \to J \to R \to R/J \to 0$。与 $R/I$ 作张量积，得到正合列
% P01-E0523: the source's R/I + J term should be the quotient R/(I + J).
$J \otimes_R R/I \to R/I \to R/I + J \to 0$，并且
$J \otimes_R R/I = J/JI$。因此，由引理 \ref{algebra-lemma-flat}，
(1)、(2) 和 (3) 等价。此外，它们蕴含 (4)。

\medskip\noindent
(4) $\Rightarrow$ (5) 是平凡的。假设 (5) 成立，并设
$x_1, \ldots, x_n \in I$。选择 $y_i \in I$，使得
$x_i = y_ix_i$。设 $y \in I$ 为满足
$1 - y = \prod_{i = 1, \ldots, n} (1 - y_i)$ 的元素。则对所有
$i = 1, \ldots, n$ 都有 $x_i = yx_i$。故 (6) 成立，从而
(5) $\Leftrightarrow$ (6)。

\medskip\noindent
假设 (5) 成立。设 $x \in I$。则对某个 $y \in I$ 有 $x = yx$。
所以 $x(1 - y) = 0$，这表明 $x$ 在 $(1 + I)^{-1}R$ 中的像为零。
映射 $R \to (1 + I)^{-1}R$ 的核当然总包含于 $I$。因此 (5)
蕴含 (9)。假设 (9) 成立。则对任意 $x \in I$，存在 $y \in I$
使得 $x(1 - y) = 0$。换言之，$x = yx$。故 (5) 与 (9) 等价。

\medskip\noindent
假设 (5) 成立。设 $\mathfrak p$ 是 $R$ 的素理想。若
$\mathfrak p \not \in V(I)$，则 $IR_{\mathfrak p} = R_{\mathfrak p}$。
若 $\mathfrak p \in V(I)$，换言之，若 $I \subset \mathfrak p$，
则对 $x \in I$，存在 $y \in I$ 使得 $x(1 - y) = 0$，因而
% P01-E0524: the source repeats “implies” and leaves this conditional sentence malformed.
$x$ 在 $R_{\mathfrak p}$ 中的像为零，即 $IR_{\mathfrak p} = 0$。
由此可见 (7) 成立。

\medskip\noindent
假设 (7) 成立。则对 $R$ 的任意素理想 $\mathfrak p$，
$(R/I)_{\mathfrak p}$ 或者是 $0$，或者是 $R_{\mathfrak p}$。
因此，由引理 \ref{algebra-lemma-flat-localization} 可知 (1) 成立。
至此可见 (1) -- (7) 和 (9) 全部等价。

\medskip\noindent
由于 $IR_{\mathfrak p} = I_{\mathfrak p}$，可知 (7) 蕴含 (8)。
最后，若 (8) 成立，则它恰好意味着：$I_{\mathfrak p}$ 是零模，
当且仅当 $\mathfrak p \in V(I)$；这显然就是 (7)。现在
(1) -- (9) 等价。

\medskip\noindent
假设 (1) -- (9) 成立。则由 (9) 有
$R/I \subset (1 + I)^{-1}R$；由 (7) 给出的素理想处局部化描述，映射
$R/I \to (1 + I)^{-1}R$ 也是满射。因此 (11) 成立。

\medskip\noindent
(11) $\Rightarrow$ (10) 是平凡的。而 (10) 蕴含 (1)，因为 $R$
的局部化作为 $R$-模是平坦的，参见引理
\ref{algebra-lemma-flat-localization}。
\end{proof}

\begin{lemma}
\label{algebra-lemma-pure-ideal-determined-by-zero-set}
\begin{slogan}
纯理想由其零点集决定。
\end{slogan}
设 $R$ 是环。若 $I, J \subset R$ 是纯理想，则
$V(I) = V(J)$ 蕴含 $I = J$。
\end{lemma}

\begin{proof}
例如，由引理 \ref{algebra-lemma-pure} 的性质 (7)，可知
$I = \Ker(R \to \prod_{\mathfrak p \in V(I)} R_{\mathfrak p})$，
所以可以从与它关联的闭子集恢复该理想。
\end{proof}

\begin{lemma}
\label{algebra-lemma-pure-open-closed-specializations}
设 $R$ 是环。规则 $I \mapsto V(I)$ 确定双射
$$
\{I \subset R \text{ pure}\}
\leftrightarrow
\{Z \subset \Spec(R)\text{ closed and closed under generalizations}\}
$$
\end{lemma}

\begin{proof}
设 $I$ 是纯理想。由于 $R \to R/I$ 平坦，由降链定理，一般化可以沿
映射 $\Spec(R/I) \to \Spec(R)$ 提升。因此 $V(I)$ 对一般化封闭，
这说明该映射良定义。由引理
\ref{algebra-lemma-pure-ideal-determined-by-zero-set}，该映射是单射。假设
$Z \subset \Spec(R)$ 是闭集且对一般化封闭。设 $J \subset R$
是满足 $Z = V(J)$ 的根理想。令
$I = \{x \in R : x \in xJ\}$。注意 $I$ 是理想：若
$x, y \in I$，则存在 $f, g \in J$，使得 $x = xf$ 且 $y = yg$。
于是
$$
x + y = (x + y)(f + g - fg)
$$
验证留给读者。我们断言 $I$ 是纯的，且 $V(I) = V(J)$。若断言成立，
则引理中的映射是满射，从而引理成立。

\medskip\noindent
注意 $I \subset J$，所以 $V(J) \subset V(I)$。设
$I \subset \mathfrak p$ 是素理想。考虑乘法子集
$S = (R \setminus \mathfrak p)(1 + J)$。由 $I$ 的定义以及
$I \subset \mathfrak p$，可知 $0 \not \in S$。因此可以找到一个与
$S$ 不交的 $R$ 的素理想 $\mathfrak q$，参见引理
\ref{algebra-lemma-localization-zero} 和 \ref{algebra-lemma-spec-localization}。
于是 $\mathfrak q \subset \mathfrak p$ 且
$\mathfrak q \cap (1 + J) = \emptyset$。这蕴含
$\mathfrak q + J$ 是 $R$ 的真理想。设 $\mathfrak m$ 是包含
$\mathfrak q + J$ 的极大理想。则 $\mathfrak m \in V(J)$；由于假设
$Z$ 对一般化封闭，故 $\mathfrak q \in V(J) = Z$。进而，由
$\mathfrak q \subset \mathfrak p$ 可知 $\mathfrak p \in V(J)$。
因此 $V(I) = V(J)$。

\medskip\noindent
最后，由于 $V(I) = V(J)$（且 $J$ 是根理想），可知
$J = \sqrt{I}$。取 $x \in I$；按定义，对某个 $y \in J$ 有
$x = xy$。于是 $x = xy = xy^2 = \ldots = xy^n$。因为对某个
$n > 0$ 有 $y^n \in I$，故引理 \ref{algebra-lemma-pure} 的性质 (5)
成立，从而 $I$ 的确是纯的。
\end{proof}

\begin{lemma}
\label{algebra-lemma-finitely-generated-pure-ideal}
设 $R$ 是环，$I \subset R$ 是理想。以下条件等价：
\begin{enumerate}
\item $I$ 是纯的且有限生成；
\item $I$ 由一个幂等元生成；
\item $I$ 是纯的且 $V(I)$ 是开集；
\item $R/I$ 是投射 $R$-模。
\end{enumerate}
\end{lemma}

\begin{proof}
若 (1) 成立，则由引理 \ref{algebra-lemma-pure} 有
$I = I \cap I = I^2$。于是由引理
\ref{algebra-lemma-ideal-is-squared-union-connected}，$I$ 由一个幂等元生成。
故 (1) $\Rightarrow$ (2)。若 (2) 成立，则 $I = (e)$，并且作为
$R$-模有 $R = (1 - e) \oplus (e)$；因此 $R/I$ 是平坦的，$I$
是纯的，且 $V(I) = D(1 - e)$ 是开集。故 (2) $\Rightarrow$
(1) $+$ (3)。最后，假设 (3) 成立。则 $V(I)$ 既开又闭，所以对
$R$ 的某个幂等元 $e$，有 $V(I) = D(1 - e)$，参见引理
\ref{algebra-lemma-disjoint-decomposition}。理想 $J = (e)$ 是满足
$V(J) = V(I)$ 的纯理想，故由引理
\ref{algebra-lemma-pure-ideal-determined-by-zero-set} 有 $I = J$。由此可见
(3) $\Rightarrow$ (2)。由引理 \ref{algebra-lemma-finite-projective}，
(4) 等价于 $I$ 是纯的且 $R/I$ 有限表现。又由引理
\ref{algebra-lemma-extension}，$R/I$ 有限表现当且仅当 $I$ 有限生成。
因此 (4) 等价于 (1)。
\end{proof}

\noindent
可以用上面的结果刻画使每个有限平坦模都有限表现的那些环。

\begin{lemma}
\label{algebra-lemma-finite-flat-module-finitely-presented}
设 $R$ 是环。以下条件等价：
\begin{enumerate}
\item 每个闭且对一般化封闭的 $Z \subset \Spec(R)$ 也是开集；
\item 任意有限平坦 $R$-模都是有限局部自由的。
\end{enumerate}
\end{lemma}

\begin{proof}
若任意有限平坦 $R$-模都有限局部自由，则对纯理想 $I$，$R/I$
的支撑是开集。因此，(2) $\Rightarrow$ (1) 由引理
% P01-E0525: the cited uniqueness lemma does not supply the pure-ideal/closed-set correspondence used here.
\ref{algebra-lemma-pure-ideal-determined-by-zero-set} 得出。

\medskip\noindent
反过来，假设 $R$ 满足 (1)。设 $M$ 是有限平坦 $R$-模。
$M$ 的支撑 $Z = \text{Supp}(M)$ 是闭集，参见引理
\ref{algebra-lemma-support-closed}。另一方面，若
$\mathfrak p \subset \mathfrak p'$，则由引理
\ref{algebra-lemma-finite-flat-local}，$M_{\mathfrak p'}$ 是自由的，并且
% P01-E0526: the source sentence lacks terminal punctuation after this displayed inline equality.
$M_{\mathfrak p} = M_{\mathfrak p'} \otimes_{R_{\mathfrak p'}} R_{\mathfrak p}$。
因此
$\mathfrak p' \in \text{Supp}(M) \Rightarrow \mathfrak p \in \text{Supp}(M)$；
换言之，该支撑对一般化封闭。由于 $R$ 满足 (1)，可知 $M$ 的支撑
既开又闭。假设 $M$ 由 $r$ 个元素 $m_1, \ldots, m_r$ 生成。
各模 $\wedge^i(M)$（$i = 1, \ldots, r$）也是有限平坦 $R$-模，
因为 $\wedge^i(M)_{\mathfrak p} = \wedge^i(M_{\mathfrak p})$
在 $R_{\mathfrak p}$ 上是自由的。注意
$\text{Supp}(\wedge^{i + 1}(M)) \subset \text{Supp}(\wedge^i(M))$。
因此存在由开闭子集给出的分解
$$
\Spec(R) = U_0 \amalg U_1 \amalg \ldots \amalg U_r
$$
使得对所有 $i = 0, \ldots, r$，$\wedge^i(M)$ 的支撑为
$U_r \cup \ldots \cup U_i$。设 $\mathfrak p$ 是 $R$ 的素理想，
并设 $\mathfrak p \in U_i$。注意
$\wedge^i(M) \otimes_R \kappa(\mathfrak p) =
\wedge^i(M \otimes_R \kappa(\mathfrak p))$。必要时重排
$m_1, \ldots, m_r$，可以假设 $m_1, \ldots, m_i$ 生成
$M \otimes_R \kappa(\mathfrak p)$。由 Nakayama 引理
\ref{algebra-lemma-NAK}，对某个 $f \in R$、$f \not \in \mathfrak p$，得到满射
$$
R_f^{\oplus i} \longrightarrow M_f, \quad
(a_1, \ldots, a_i) \longmapsto \sum a_im_i
$$
还可以假设 $D(f) \subset U_i$。这意味着
$\wedge^i(M_f) = \wedge^i(M)_f$ 是平坦 $R_f$-模，其支撑是整个
$\Spec(R_f)$。由上可知，它由单个元素
$m_1 \wedge \ldots \wedge m_i$ 生成。因此，对某个纯理想
$J \subset R_f$，有 $\wedge^i(M)_f \cong R_f/J$，并且
$V(J) = \Spec(R_f)$。显然这意味着 $J = (0)$，参见引理
\ref{algebra-lemma-pure-ideal-determined-by-zero-set}。所以
$m_1 \wedge \ldots \wedge m_i$ 是 $\wedge^i(M_f)$ 的一组基，
从而所显示的映射既是单射也是满射。这按要求证明了 $M$ 有限局部自由。
\end{proof}

\section{有限整体维数的环}
\label{algebra-section-ring-finite-gl-dim}

\noindent
下面的引理常用于比较给定模的不同投射消解。

\begin{lemma}[Schanuel 引理]
\label{algebra-lemma-Schanuel}
设 $R$ 是环，$M$ 是 $R$-模。假设
$$
0 \to K \xrightarrow{c_1} P_1 \xrightarrow{p_1} M \to 0
\quad\text{且}\quad
0 \to L \xrightarrow{c_2} P_2 \xrightarrow{p_2} M \to 0
$$
是两个短正合列，其中 $P_i$ 是投射模。则
$K \oplus P_2 \cong L \oplus P_1$。更精确地，存在交换图
$$
\xymatrix{
0 \ar[r] &
K \oplus P_2 \ar[r]_{(c_1, \text{id})} \ar[d] &
P_1 \oplus P_2 \ar[r]_{(p_1, 0)} \ar[d] &
M \ar[r] \ar@{=}[d] &
0 \\
0 \ar[r] &
P_1 \oplus L \ar[r]^{(\text{id}, c_2)} &
P_1 \oplus P_2 \ar[r]^{(0, p_2)} &
M \ar[r] &
0
}
$$
使得各竖直箭头都是同构。
\end{lemma}

\begin{proof}
考虑由短正合列
$0 \to N \to P_1 \oplus P_2 \to M \to 0$
定义的模 $N$，其中最后一个映射是两个映射 $P_i \to M$ 之和。
容易看出，投影 $N \to P_1$ 是满射，其核为 $L$；而
$N \to P_2$ 是满射，其核为 $K$。由于 $P_i$ 是投射模，故
$N \cong K \oplus P_2 \cong L \oplus P_1$。这证明了第一个断言。

\medskip\noindent
为证明第二个断言（同时再次证明第一个断言），选择
$a : P_1 \to P_2$ 和 $b : P_2 \to P_1$，使得
$p_1 = p_2 \circ a$ 且 $p_2 = p_1 \circ b$。这是可行的，因为
$P_1$ 和 $P_2$ 都是投射模。于是得到交换图
$$
\xymatrix{
0 \ar[r] &
K \oplus P_2 \ar[r]_{(c_1, \text{id})} &
P_1 \oplus P_2 \ar[r]_{(p_1, 0)} &
M \ar[r] &
0 \\
0 \ar[r] &
N \ar[r] \ar[d] \ar[u] &
P_1 \oplus P_2 \ar[r]_{(p_1, p_2)}
\ar[d]_S \ar[u]^T &
M \ar[r] \ar@{=}[d] \ar@{=}[u] &
0 \\
0 \ar[r] &
P_1 \oplus L \ar[r]^{(\text{id}, c_2)} &
P_1 \oplus P_2 \ar[r]^{(0, p_2)} &
M \ar[r] &
0
}
$$
其中 $T$ 和 $S$ 由矩阵
$$
S = \left(
\begin{matrix}
\text{id} & 0 \\
a & \text{id}
\end{matrix}
\right)
\quad\text{且}\quad
T = \left(
\begin{matrix}
\text{id} & b \\
0 & \text{id}
\end{matrix}
\right)
$$
给出。于是 $S$、$T$ 以及映射 $N \to P_1 \oplus L$ 和
$N \to K \oplus P_2$ 都是所需的同构。
\end{proof}

\begin{definition}
\label{algebra-definition-finite-proj-dim}
设 $R$ 是环，$M$ 是 $R$-模。若 $M$ 有一个由投射 $R$-模组成的
有限长度消解，则称其具有{\it 有限投射维数}。这种消解的最小长度
称为 $M$ 的{\it 投射维数}。
\end{definition}

\noindent
显然，$M$ 的投射维数为 $0$ 当且仅当 $M$ 是投射模。下面的引理说明，
投射维数在何种程度上不依赖于投射消解的选择。

\begin{lemma}
\label{algebra-lemma-independent-resolution}
设 $R$ 是环。假设 $M$ 是投射维数为 $d$ 的 $R$-模。假设
$F_e \to F_{e-1} \to \ldots \to F_0 \to M \to 0$
正合，其中 $F_i$ 是投射模，并且 $e \geq d - 1$。则
$F_e \to F_{e-1}$ 的核是投射模（当 $e = 0$ 时，则指
$F_0 \to M$ 的核是投射模）。
\end{lemma}

\begin{proof}
对 $d$ 归纳。若 $d = 0$，则 $M$ 是投射模。在这种情形下有分裂
$F_0 = \Ker(F_0 \to M) \oplus M$，所以
$\Ker(F_0 \to M)$ 是投射模。若 $e = 0$，证明至此完成；若
$e > 0$，则以 $\Ker(F_0 \to M)$ 代替 $M$，便使 $e$ 减少。

\medskip\noindent
下面假设 $d > 0$。设
$0 \to P_d \to P_{d-1} \to \ldots \to P_0 \to M \to 0$
是由投射模 $P_i$ 组成的最短有限消解。由 Schanuel 引理
\ref{algebra-lemma-Schanuel}，
$P_0 \oplus \Ker(F_0 \to M) \cong F_0 \oplus \Ker(P_0 \to M)$。
当 $d = 1$、$e = 0$ 时，这证明了结论，因为此时右端是投射模
$F_0 \oplus P_1$。因此现在可以假设 $e > 0$。模
$F_0 \oplus \Ker(P_0 \to M)$ 有长度为 $d - 1$ 的有限投射消解
$$
0 \to P_d \to P_{d-1} \to \ldots \to
P_2 \to P_1 \oplus F_0 \to \Ker(P_0 \to M) \oplus F_0 \to 0
$$
把归纳假设应用于长度为 $e - 1$ 的正合列
$$
F_e \to F_{e-1} \to \ldots \to F_2 \to P_0 \oplus F_1 \to
P_0 \oplus \Ker(F_0 \to M) \to 0
$$
可知，当 $e \geq 2$ 时 $\Ker(F_e \to F_{e - 1})$ 是投射模；
或者，$\Ker(F_1 \oplus P_0 \to F_0 \oplus P_0)$ 是投射模。
这蕴含引理。
\end{proof}

\begin{lemma}
\label{algebra-lemma-what-kind-of-resolutions}
设 $R$ 是环，$M$ 是 $R$-模，并设 $d \geq 0$。以下条件等价：
\begin{enumerate}
\item $M$ 的投射维数 $\leq d$；
\item 存在由投射模 $P_i$ 组成的消解
$0 \to P_d \to P_{d - 1} \to \ldots \to P_0 \to M \to 0$；
\item 对某个由投射模 $P_i$ 组成的消解
$\ldots \to P_2 \to P_1 \to P_0 \to M \to 0$，当
$d \geq 2$ 时 $\Ker(P_{d - 1} \to P_{d - 2})$ 是投射模；当
$d = 1$ 时 $\Ker(P_0 \to M)$ 是投射模；当 $d = 0$ 时
$M$ 是投射模；
\item 对任意由投射模 $P_i$ 组成的消解
$\ldots \to P_2 \to P_1 \to P_0 \to M \to 0$，当
$d \geq 2$ 时 $\Ker(P_{d - 1} \to P_{d - 2})$ 是投射模；当
$d = 1$ 时 $\Ker(P_0 \to M)$ 是投射模；当 $d = 0$ 时
$M$ 是投射模。
\end{enumerate}
\end{lemma}

\begin{proof}
(1) 与 (2) 的等价性就是投射维数的定义，参见定义
\ref{algebra-definition-finite-proj-dim}。由引理
\ref{algebra-lemma-independent-resolution} 有 (2) $\Rightarrow$ (4)。
(4) $\Rightarrow$ (3) 和 (3) $\Rightarrow$ (2) 都是直接的。
\end{proof}

\begin{lemma}
\label{algebra-lemma-what-kind-of-resolutions-local}
设 $R$ 是局部环，$M$ 是 $R$-模，并设 $d \geq 0$。引理
\ref{algebra-lemma-what-kind-of-resolutions} 中的等价条件 (1) -- (4)
还等价于
\begin{enumerate}
\item[(5)] 存在由自由模 $P_i$ 组成的消解
$0 \to P_d \to P_{d - 1} \to \ldots \to P_0 \to M \to 0$。
\end{enumerate}
\end{lemma}

\begin{proof}
由引理 \ref{algebra-lemma-what-kind-of-resolutions} 和定理
\ref{algebra-theorem-projective-free-over-local-ring} 得出。
\end{proof}

\begin{lemma}
\label{algebra-lemma-what-kind-of-resolutions-Noetherian}
设 $R$ 是诺特环，$M$ 是有限 $R$-模，并设 $d \geq 0$。引理
\ref{algebra-lemma-what-kind-of-resolutions} 中的等价条件 (1) -- (4)
还等价于
\begin{enumerate}
\item[(6)] 存在由有限投射模 $P_i$ 组成的消解
$0 \to P_d \to P_{d - 1} \to \ldots \to P_0 \to M \to 0$。
\end{enumerate}
\end{lemma}

\begin{proof}
选择由有限自由模 $F_i$ 组成的消解
$\ldots \to F_2 \to F_1 \to F_0 \to M \to 0$
（引理 \ref{algebra-lemma-resolution-by-finite-free}）。由引理
\ref{algebra-lemma-what-kind-of-resolutions} 可知，至少当 $d \geq 2$ 时，
% P01-E0527: the source proof does not address the d = 0 and d = 1 cases included in the statement.
$P_d = \Ker(F_{d - 1} \to F_{d - 2})$ 是投射模。由于 $R$ 是诺特环，
并且 $P_d \subset F_{d - 1}$、后者有限自由，故 $P_d$
是有限 $R$-模。因此
$0 \to P_d \to F_{d - 1} \to \ldots \to F_1 \to F_0 \to M \to 0$
就是所需的消解。
\end{proof}

\begin{lemma}
\label{algebra-lemma-what-kind-of-resolutions-Noetherian-local}
设 $R$ 是诺特局部环，$M$ 是有限 $R$-模，并设 $d \geq 0$。引理
\ref{algebra-lemma-what-kind-of-resolutions} 中的等价条件 (1) -- (4)、
引理 \ref{algebra-lemma-what-kind-of-resolutions-local} 的条件 (5)，以及
引理 \ref{algebra-lemma-what-kind-of-resolutions-Noetherian} 的条件 (6)
还等价于
\begin{enumerate}
\item[(7)] 存在由有限自由模 $F_i$ 组成的消解
$0 \to F_d \to F_{d - 1} \to \ldots \to F_0 \to M \to 0$。
\end{enumerate}
\end{lemma}

\begin{proof}
这由引理 \ref{algebra-lemma-what-kind-of-resolutions}、
\ref{algebra-lemma-what-kind-of-resolutions-local} 和
\ref{algebra-lemma-what-kind-of-resolutions-Noetherian} 得出；还要用到局部环上的
有限投射模是有限自由模，参见引理 \ref{algebra-lemma-finite-projective}。
\end{proof}

\begin{lemma}
\label{algebra-lemma-projective-dimension-ext}
设 $R$ 是环，$M$ 是 $R$-模，并设 $n \geq 0$。以下条件等价：
\begin{enumerate}
\item $M$ 的投射维数 $\leq n$；
\item 对所有 $R$-模 $N$ 和所有 $i \geq n + 1$，都有
$\Ext^i_R(M, N) = 0$；
\item 对所有 $R$-模 $N$，都有 $\Ext^{n + 1}_R(M, N) = 0$。
\end{enumerate}
\end{lemma}

\begin{proof}
假设 (1) 成立。选择 $M$ 的自由消解 $F_\bullet \to M$。记
$d_e : F_e \to F_{e - 1}$。由引理
\ref{algebra-lemma-independent-resolution} 可知，对 $e \geq n - 1$，
$P_e = \Ker(d_e)$ 是投射模。于是对 $e \geq n$ 有
% P01-E0528: for n = 0 the source uses P_{e - 1} at e = 0 without defining P_{-1}.
$F_e \cong P_e \oplus P_{e - 1}$，其中 $d_e$ 把直和项
$P_{e - 1}$ 同构地映到 $F_{e - 1}$ 中的 $P_{e - 1}$。
因此，对任意 $R$-模 $N$，复形 $\Hom_R(F_\bullet, N)$
在次数 $\geq n + 1$ 上分裂正合。故 (2) 成立。
(2) $\Rightarrow$ (3) 是平凡的。

\medskip\noindent
假设 (3) 成立。若 $n = 0$，则由引理
\ref{algebra-lemma-characterize-projective}，$M$ 是投射模，所以 (1) 成立。
若 $n > 0$，选择自由 $R$-模 $F$ 和核为 $K$ 的满射 $F \to M$。
由引理 \ref{algebra-lemma-reverse-long-exact-seq-ext}，以及根据 (1)
对所有 $i > 0$ 都有 $\Ext_R^i(F, N)$ 消失，可知对所有 $R$-模
$N$ 都有 $\Ext_R^n(K, N) = 0$。故由归纳假设，$K$ 的投射维数
$\leq n - 1$。在 $K$ 的任意有限投射消解前添加 $F$，便得到长度
多一的 $M$ 的投射消解，所以 $M$ 的投射维数 $\leq n$。
\end{proof}

\begin{lemma}
\label{algebra-lemma-exact-sequence-projective-dimension}
设 $R$ 是环，并设 $0 \to M' \to M \to M'' \to 0$ 是 $R$-模的短正合列。
\begin{enumerate}
\item 若 $M$ 的投射维数 $\leq n$，且 $M''$ 的投射维数
$\leq n + 1$，则 $M'$ 的投射维数 $\leq n$。
\item 若 $M'$ 和 $M''$ 的投射维数都 $\leq n$，则 $M$
的投射维数 $\leq n$。
\item 若 $M'$ 的投射维数 $\leq n$，且 $M$ 的投射维数
$\leq n + 1$，则 $M''$ 的投射维数 $\leq n + 1$。
\end{enumerate}
\end{lemma}

\begin{proof}
把引理 \ref{algebra-lemma-projective-dimension-ext} 中对投射维数的刻画与引理
\ref{algebra-lemma-reverse-long-exact-seq-ext} 中的 Ext 群长正合列结合起来。
\end{proof}

\begin{definition}
\label{algebra-definition-finite-gl-dim}
设 $R$ 是环。若存在整数 $n$，使得每个 $R$-模都有一个长度至多为
$n$ 的投射 $R$-模消解，则称环 $R$ 具有{\it 有限整体维数}。
满足此条件的最小 $n$ 称为 $R$ 的{\it 整体维数}。
\end{definition}

\noindent
下面引理的证明论证可见 Auslander 的论文 \cite{Auslander}。

\begin{lemma}
\label{algebra-lemma-colimit-projective-dimension}
设 $R$ 是环。假设有模 $M = \bigcup_{e \in E} M_e$，其中各子模
$M_e$ 按包含关系良序排列。假设各商模
$M_e/\bigcup\nolimits_{e' < e} M_{e'}$ 的投射维数 $\leq n$。
则 $M$ 的投射维数 $\leq n$。
\end{lemma}

\begin{proof}
对 $n$ 归纳。

\medskip\noindent
基本情形：$n = 0$。此时
$P_e = M_e/\bigcup_{e' < e} M_{e'}$ 是投射模。因此可以选择投影
$M_e \to P_e$ 的截面 $P_e \to M_e$。我们断言，诱导映射
$\psi : \bigoplus_{e \in E} P_e \to M$ 是同构。事实上，若
$x = \sum x_e \in \bigoplus P_e$ 非零，则令 $e_{max}$ 为使
$x_{e_{max}}$ 非零的最大指标，并令
$y = \psi(x) = \psi(\sum x_e)$。由于
$y \in M_{e_{max}}$ 在 $P_{e_{max}}$ 中的像 $x_{e_{max}}$ 非零，
故该元素非零。
另一方面，设 $y \in M$。则对某个 $e$ 有 $y \in M_e$。
我们对 $e$ 作超限归纳，证明 $y \in \Im(\psi)$。设
$x_e \in P_e$ 是 $y$ 的像。则
$y - \psi(x_e) \in \bigcup_{e' < e} M_{e'}$。由归纳假设，
$y - \psi(x_e) \in \Im(\psi)$，所以 $y \in \Im(\psi)$。
因此断言成立，$\psi$ 是同构。由引理
\ref{algebra-lemma-direct-sum-projective}，$M$ 作为投射模的直和是投射模。

\medskip\noindent
若 $n > 0$，则对 $e \in E$，令 $F_e$ 为以 $M_e$ 的元素集为基的
自由 $R$-模。于是，在良序集 $E$ 上有一族短正合列
$$
0 \to K_e \to F_e \to M_e \to 0
$$
注意，转移映射 $F_{e'} \to F_e$ 和 $K_{e'} \to K_e$ 也都是单射。
令 $F = \bigcup F_e$、$K = \bigcup K_e$。则
$$
0 \to
K_e/\bigcup\nolimits_{e' < e} K_{e'} \to
F_e/\bigcup\nolimits_{e' < e} F_{e'} \to
M_e/\bigcup\nolimits_{e' < e} M_{e'} \to 0
$$
也是 $R$-模的短正合列，并且
$F_e/\bigcup_{e' < e} F_{e'}$ 是以 $M_e$ 中不属于
$\bigcup_{e' < e} M_{e'}$ 的元素集为基的自由 $R$-模。因此，由引理
\ref{algebra-lemma-exact-sequence-projective-dimension} 可知
$K_e/\bigcup_{e' < e} K_{e'}$ 的投射维数至多为 $n - 1$。
由归纳假设，$K$ 的投射维数至多为 $n - 1$。因而 $M$
的投射维数至多为 $n$，得证。
\end{proof}

\begin{lemma}
\label{algebra-lemma-finite-gl-dim}
设 $R$ 是环。以下条件等价：
\begin{enumerate}
\item $R$ 的整体维数有限且 $\leq n$；
\item 每个有限 $R$-模的投射维数 $\leq n$；
\item 每个循环 $R$-模 $R/I$ 的投射维数 $\leq n$。
\end{enumerate}
\end{lemma}

\begin{proof}
显然 (1) $\Rightarrow$ (2)，且 (2) $\Rightarrow$ (3)。
假设 (3) 成立。选择 $M$ 的生成元集 $E \subset M$，并在 $E$
上选择一个良序。对 $e \in E$，记 $M_e$ 为由满足 $e' \leq e$
的元素 $e' \in E$ 生成的 $M$ 的子模。则
$M = \bigcup_{e \in E} M_e$。注意，对每个 $e \in E$，商模
$$
M_e/\bigcup\nolimits_{e' < e} M_{e'}
$$
或者为零，或者由一个元素生成，故由 (3)，其投射维数 $\leq n$。
由引理 \ref{algebra-lemma-colimit-projective-dimension}，这意味着 $M$
的投射维数 $\leq n$。
\end{proof}

\begin{lemma}
\label{algebra-lemma-localize-finite-gl-dim}
设 $R$ 是环，$M$ 是 $R$-模，并设 $S \subset R$ 是乘法子集。
\begin{enumerate}
\item 若 $M$ 的投射维数 $\leq n$，则 $S^{-1}M$ 作为
$S^{-1}R$-模的投射维数 $\leq n$。
\item 若 $R$ 的整体维数有限且 $\leq n$，则 $S^{-1}R$
的整体维数有限且 $\leq n$。
\end{enumerate}
\end{lemma}

\begin{proof}
设
$0 \to P_n \to P_{n - 1} \to \ldots \to P_0 \to M \to 0$
是投射消解。局部化是正合的，参见命题
\ref{algebra-proposition-localization-exact}；且每个 $S^{-1}P_i$ 都是投射
$S^{-1}R$-模，参见引理 \ref{algebra-lemma-ascend-properties-modules}。
因此
$0 \to S^{-1}P_n \to \ldots \to S^{-1}P_0 \to S^{-1}M \to 0$
是 $S^{-1}M$ 的投射消解。这证明了 (1)。设 $M'$ 是
$S^{-1}R$-模。注意 $M' = S^{-1}M'$。因此 (2) 由 (1) 得出。
\end{proof}

\section{正则环与整体维数}
\label{algebra-section-regular-finite-gl-dim}

\noindent
可以用有限整体维数环的有关材料给出正则局部环的另一个刻画。

\begin{proposition}
\label{algebra-proposition-regular-finite-gl-dim}
设 $R$ 是维数为 $d$ 的正则局部环。每个深度为 $e$ 的有限
$R$-模 $M$ 都有有限自由消解
$$
0 \to F_{d-e} \to \ldots \to F_0 \to M \to 0.
$$
特别地，正则局部环的整体维数 $\leq d$。
\end{proposition}

\begin{proof}
第一部分由引理 \ref{algebra-lemma-regular-mcm-free} 和引理
\ref{algebra-lemma-mcm-resolution} 得出。最后一部分由此及引理
\ref{algebra-lemma-finite-gl-dim} 得出。
\end{proof}

\begin{lemma}
\label{algebra-lemma-finite-gl-dim-primes}
设 $R$ 是诺特环，$n \geq 0$ 是整数。则 $R$ 的整体维数有限且
$\leq n$，当且仅当对 $R$ 的所有极大理想 $\mathfrak m$，环
$R_{\mathfrak m}$ 的整体维数 $\leq n$。
\end{lemma}

\begin{proof}
由引理 \ref{algebra-lemma-localize-finite-gl-dim} 已知，若 $R$
的有限整体维数为 $n$，则所有局部化 $R_{\mathfrak m}$
的有限整体维数至多为 $n$。
% P01-E0529: the source reads “We saw, Lemma ... that”, with a malformed citation clause.
反过来，假设所有 $R_{\mathfrak m}$ 的整体维数都 $\leq n$。
设 $M$ 是有限 $R$-模，并设
% P01-E0530: for n = 0 the source resolution uses F_{n-1}=F_{-1} without a convention.
$0 \to K_n \to F_{n-1} \to \ldots \to F_0 \to M \to 0$
是由有限自由模 $F_i$ 组成的消解。则 $K_n$ 是有限 $R$-模。
由引理 \ref{algebra-lemma-independent-resolution} 和假设，所有模
$K_n \otimes_R R_{\mathfrak m}$ 都是投射模。因此，由引理
\ref{algebra-lemma-finite-projective}，模 $K_n$ 是有限投射模。
\end{proof}

\begin{lemma}
\label{algebra-lemma-length-resolution-residue-field}
假设 $R$ 是极大理想为 $\mathfrak m$、剩余域为 $\kappa$
的诺特局部环。在这种情形下，$\kappa$ 的投射维数
$\geq \dim_\kappa \mathfrak m / \mathfrak m^2$。
\end{lemma}

\begin{proof}
设 $x_1 , \ldots, x_n$ 是 $\mathfrak m$ 中的元素，其在
$\mathfrak m / \mathfrak m^2$ 中的像构成一组基。考虑
$x_1, \ldots, x_n$ 上的{\it Koszul 复形}。这是复形
$$
0 \to \wedge^n R^n \to \wedge^{n-1} R^n \to \wedge^{n-2} R^n \to
\ldots \to \wedge^i R^n \to \ldots \to R^n \to R
$$
其映射由
$$
e_{j_1} \wedge \ldots \wedge e_{j_i}
\longmapsto
\sum_{a = 1}^i (-1)^{a + 1} x_{j_a} e_{j_1} \wedge \ldots
\wedge \hat e_{j_a} \wedge \ldots \wedge e_{j_i}
$$
给出。容易看出这是一个复形 $K_{\bullet}(R, x_{\bullet})$。
注意，由引理 \ref{algebra-lemma-NAK} 第 (8) 部分，
$K_{\bullet}(R, x_{\bullet})$ 的最后一个映射的余核是 $\kappa$。

\medskip\noindent
若 $\kappa$ 的投射维数有限且为 $d$，则可以找到由有限自由
$R$-模组成的长度为 $d$ 的消解 $F_{\bullet} \to \kappa$
（引理 \ref{algebra-lemma-what-kind-of-resolutions-Noetherian-local}）。
由引理 \ref{algebra-lemma-add-trivial-complex}，可以假设复形
$F_{\bullet}$ 中的所有映射都满足
$\Im(F_i \to F_{i-1}) \subset \mathfrak m F_{i-1}$，因为从消解中
移除平凡直和项至多只会缩短消解。由引理
\ref{algebra-lemma-compare-resolutions}，可以找到诱导 $\kappa$ 上恒等映射的
复形映射 $\alpha : K_{\bullet}(R, x_{\bullet}) \to F_{\bullet}$。
我们将归纳证明各映射
$\alpha_i : \wedge^i R^n = K_i(R, x_{\bullet}) \to F_i$ 都满足
$\alpha_i \otimes \kappa : \wedge^i \kappa^n \to F_i \otimes \kappa$
是单射。这表明 $F_n \not = 0$，因而按要求有 $d \geq n$。

\medskip\noindent
当 $i = 0$ 时，结论是清楚的，因为复合
$R \xrightarrow{\alpha_0} F_0 \to \kappa$ 非零。注意，$F_0$
的秩必为 $1$，否则余核为单个 $\kappa$ 的映射
$F_1 \to F_0$ 不可能具有包含在 $\mathfrak m F_0$ 中的像。

\medskip\noindent
接着检验 $i = 1$ 的情形，因为我们认为这颇具启发性；读者可以跳过，
因为归纳步骤会由 $i = 0$ 的情形推出 $i = 1$ 的情形。上面已经看到
$F_0 = R$，并且 $F_1 \to F_0 = R$ 的像是 $\mathfrak m$。
有交换图
$$
\begin{matrix}
R^n & = & K_1(R, x_{\bullet}) & \to & K_0(R, x_{\bullet}) & = & R \\
& & \downarrow & & \downarrow & & \downarrow \\
& & F_1 & \to & F_0 & = & R
\end{matrix}
$$
其中最右边的竖直箭头由乘以一个单位给出。因此，复合
$R^n \to F_1 \to F_0 = R$ 的像也等于 $\mathfrak m$。
由于 $\dim_\kappa (\mathfrak m / \mathfrak m^2) = n$，映射
$R^n \otimes \kappa \to F_1 \otimes \kappa$ 必为单射。

\medskip\noindent
设 $i \geq 1$，并假设已经对所有 $j \leq i - 1$ 证明了
$\alpha_j \otimes \kappa$ 的单射性。考虑交换图
$$
\begin{matrix}
\wedge^i R^n & = & K_i(R, x_{\bullet}) & \to & K_{i-1}(R, x_{\bullet})
& = & \wedge^{i-1} R^n \\
& & \downarrow & & \downarrow & & \\
& & F_i & \to & F_{i-1} & &
\end{matrix}
$$
已知 $\wedge^{i-1} \kappa^n \to F_{i-1} \otimes \kappa$
是单射。这证明了
$\wedge^{i-1} \kappa^n \otimes_{\kappa} \mathfrak m/\mathfrak m^2
\to F_{i-1} \otimes \mathfrak m/\mathfrak m^2$ 是单射。
此外，由复形的选择，$F_i$ 映入 $\mathfrak mF_{i-1}$；
Koszul 复形也同样如此。因此得到交换图
$$
\begin{matrix}
\wedge^i \kappa^n & \to &
\wedge^{i-1} \kappa^n \otimes \mathfrak m/\mathfrak m^2 \\
\downarrow & & \downarrow \\
F_i \otimes \kappa & \to & F_{i-1} \otimes \mathfrak m/\mathfrak m^2
\end{matrix}
$$
此时，只须验证映射
$\wedge^i \kappa^n \to
\wedge^{i-1} \kappa^n \otimes \mathfrak m/\mathfrak m^2$
是单射，而这可以直接检验。
\end{proof}

\begin{lemma}
\label{algebra-lemma-dim-gl-dim}
设 $R$ 是诺特局部环。假设剩余域 $\kappa$ 在 $R$
上的投射维数有限且为 $n$。在这种情形下，$\dim(R) \geq n$。
\end{lemma}

\begin{proof}
设 $F_{\bullet}$ 是由有限自由 $R$-模组成的 $\kappa$ 的有限消解
（引理 \ref{algebra-lemma-what-kind-of-resolutions-Noetherian-local}）。
由引理 \ref{algebra-lemma-add-trivial-complex}，可以假设复形
$F_{\bullet}$ 中的所有映射都具有性质
% P01-E0531: the source says “have to property” instead of “have the property”.
$\Im(F_i \to F_{i-1}) \subset \mathfrak m F_{i-1}$，因为从消解中移除
平凡直和项至多只会缩短消解。设 $F_n \not = 0$，且对 $i > n$
有 $F_i = 0$，从而 $\kappa$ 的投射维数为 $n$。由命题
\ref{algebra-proposition-what-exact}，可知
$\text{depth}_{I(\varphi_n)}(R) \geq n$，因为按复形的选择，
$I(\varphi_n)$ 不可能等于 $R$。故由引理
\ref{algebra-lemma-bound-depth} 还有 $\dim(R) \geq n$。
\end{proof}

\begin{proposition}
\label{algebra-proposition-finite-gl-dim-regular}
设 $(R, \mathfrak m, \kappa)$ 是诺特局部环。以下条件等价：
\begin{enumerate}
\item $\kappa$ 作为 $R$-模具有有限投射维数；
\item $R$ 具有有限整体维数；
\item $R$ 是正则局部环。
\end{enumerate}
此外，在这种情形下，$R$ 的整体维数等于
$\dim(R) = \dim_\kappa(\mathfrak m/\mathfrak m^2)$。
\end{proposition}

\begin{proof}
由命题 \ref{algebra-proposition-regular-finite-gl-dim} 有
(3) $\Rightarrow$ (2)。(2) $\Rightarrow$ (1) 是平凡的。
假设 (1) 成立。由引理 \ref{algebra-lemma-length-resolution-residue-field}
和 \ref{algebra-lemma-dim-gl-dim}，可知
$\dim(R) \geq \dim_\kappa(\mathfrak m /\mathfrak m^2)$。
因此 $R$ 是正则的，参见定义 \ref{algebra-definition-regular-local}
及其前面的讨论。假设等价条件 (1) -- (3) 成立。由命题
\ref{algebra-proposition-regular-finite-gl-dim}，$R$ 的整体维数至多为
$\dim(R)$；由引理 \ref{algebra-lemma-length-resolution-residue-field}，
它至少为 $\dim_\kappa(\mathfrak m/\mathfrak m^2)$。故所述等式成立。
\end{proof}

\begin{lemma}
\label{algebra-lemma-localization-of-regular-local-is-regular}
诺特局部环 $R$ 是正则局部环，当且仅当它具有有限整体维数。
在这种情形下，对所有素理想 $\mathfrak p$，$R_{\mathfrak p}$
都是正则局部环。
\end{lemma}

\begin{proof}
由命题 \ref{algebra-proposition-finite-gl-dim-regular} 和
\ref{algebra-proposition-regular-finite-gl-dim} 可知，诺特局部环是正则局部环，
当且仅当它具有有限整体维数。此外，任意局部化
$R_{\mathfrak p}$ 都具有有限整体维数，参见引理
\ref{algebra-lemma-localize-finite-gl-dim}，因而是正则局部环。
\end{proof}

\noindent
由引理 \ref{algebra-lemma-localization-of-regular-local-is-regular}，
作如下定义是合理的，因为它不与此前的正则局部环定义冲突。

\begin{definition}
\label{algebra-definition-regular}
若诺特环 $R$ 在所有素理想处的局部化 $R_{\mathfrak p}$
都是正则局部环，则称其是{\it 正则的}。
\end{definition}

\noindent
只要求极大理想处的局部环正则就已足够。注意，即使假设 $R$
是诺特环，这也不同于要求 $R$ 具有有限整体维数。这是因为存在
不具有有限整体维数的正则诺特环；其原因正是它不具有有限维数。

\begin{lemma}
\label{algebra-lemma-finite-gl-dim-finite-dim-regular}
设 $R$ 是诺特环。以下条件等价：
\begin{enumerate}
\item $R$ 的有限整体维数为 $n$；
\item $R$ 是维数为 $n$ 的正则环；
\item 存在整数 $n$，使得极大理想处的所有局部化
$R_{\mathfrak m}$ 都是维数 $\leq n$ 的正则环，并且至少对一个
$\mathfrak m$ 取到等号；
\item 存在整数 $n$，使得素理想处的所有局部化
$R_{\mathfrak p}$ 都是维数 $\leq n$ 的正则环，并且至少对一个
$\mathfrak p$ 取到等号。
\end{enumerate}
\end{lemma}

\begin{proof}
这由上面的讨论得出。更精确地说，这是把定义
\ref{algebra-definition-regular}、引理 \ref{algebra-lemma-finite-gl-dim-primes}
和命题 \ref{algebra-proposition-finite-gl-dim-regular} 结合起来的结果。
\end{proof}

\begin{lemma}
\label{algebra-lemma-flat-under-regular}
设 $R \to S$ 是诺特局部环的局部同态。假设 $R \to S$
平坦，且 $S$ 正则。则 $R$ 正则。
\end{lemma}

\begin{proof}
设 $\mathfrak m \subset R$ 是极大理想，令
$\kappa = R/\mathfrak m$ 为剩余域，并令 $d = \dim S$。
选择任意消解 $F_\bullet \to \kappa$，其中每个 $F_i$
都是有限自由 $R$-模。令
% P01-E0532: when d = 0 the source defines K_d using F_{-1} and F_{-2}.
$K_d = \Ker(F_{d - 1} \to F_{d - 2})$。由 $R \to S$ 的平坦性，复形
$0 \to K_d \otimes_R S \to F_{d - 1} \otimes_R S \to \ldots
\to F_0 \otimes_R S \to \kappa \otimes_R S \to 0$
仍然正合。因为 $S$ 的整体维数为 $d$，
% P01-E0533: the cited proposition gives only the upper bound; the equality is established by the preceding proposition-finite-gl-dim-regular.
参见命题 \ref{algebra-proposition-regular-finite-gl-dim}，可知
$K_d \otimes_R S$ 是有限自由 $S$-模（另见引理
\ref{algebra-lemma-independent-resolution}）。由引理
\ref{algebra-lemma-finite-projective-descends}，可知 $K_d$
是有限自由 $R$-模。因此 $\kappa$ 具有有限投射维数；由命题
\ref{algebra-proposition-finite-gl-dim-regular}，$R$ 正则。
\end{proof}

\section{Auslander--Buchsbaum 公式}
\label{algebra-section-Auslander-Buchsbaum}

\noindent
下述结果见 \cite{Auslander-Buchsbaum}。

\begin{proposition}
\label{algebra-proposition-Auslander-Buchsbaum}
设 $R$ 是诺特局部环。设 $M$ 是非零有限 $R$-模，且具有有限投射维数
$\text{pd}_R(M)$。则
$$
\text{depth}(R) = \text{pd}_R(M) + \text{depth}(M)
$$
成立。
\end{proposition}

\begin{proof}
我们对 $\text{depth}(M)$ 作归纳。最有意思的情形是
$\text{depth}(M) = 0$。在此情形下，设
$$
0 \to R^{n_e} \to R^{n_{e-1}} \to \ldots \to R^{n_0} \to M \to 0
$$
是一个极小有限自由消解，于是 $e = \text{pd}_R(M)$。
由引理 \ref{algebra-lemma-add-trivial-complex}，可以假设该复形中各映射的
所有矩阵系数都包含在 $R$ 的极大理想中。一方面，由命题
\ref{algebra-proposition-what-exact} 可知 $\text{depth}(R) \geq e$。
另一方面，把这个长正合列分解成短正合列
% P01-E0534: for e = 0 or e = 1, the displayed decomposition uses negative-indexed K terms and needs separate low-dimensional cases.
\begin{align*}
0 \to R^{n_e} \to R^{n_{e - 1}} \to K_{e - 2} \to 0,\\
0 \to K_{e - 2} \to R^{n_{e - 2}} \to K_{e - 3} \to 0,\\
\ldots,\\
0 \to K_0 \to R^{n_0} \to M \to 0
\end{align*}
并使用引理 \ref{algebra-lemma-depth-in-ses}，可得
\begin{align*}
\text{depth}(K_{e - 2}) \geq \text{depth}(R) - 1,\\
\text{depth}(K_{e - 3}) \geq \text{depth}(R) - 2,\\
\ldots,\\
\text{depth}(K_0) \geq \text{depth}(R) - (e - 1),\\
\text{depth}(M) \geq \text{depth}(R) - e
\end{align*}
又因为 $\text{depth}(M) = 0$，所以 $\text{depth}(R) \leq e$。
这就完成了 $\text{depth}(M) = 0$ 情形的证明。

\medskip\noindent
归纳步骤。若 $\text{depth}(M) > 0$，则取 $x \in \mathfrak m$，
使其同时是 $M$ 和 $R$ 上的非零因子。这是可能的，因为要么
$\text{pd}_R(M) > 0$，并且由前述命题
\ref{algebra-proposition-what-exact} 有 $\text{depth}(R) > 0$；要么
$\text{pd}_R(M) = 0$，此时 $M$ 是有限自由模，因而同样有
$\text{depth}(R) = \text{depth}(M) > 0$。
% P01-E0535: the source immediately chooses the same x a second time after already choosing it above.
于是由引理 \ref{algebra-lemma-depth-in-ses}（例如）可知
$\text{depth}(R \oplus M) > 0$，从而可以找到 $x \in \mathfrak m$，
使其同时是 $R$ 和 $M$ 上的非零因子。设
$$
0 \to R^{n_e} \to R^{n_{e-1}} \to \ldots \to R^{n_0} \to M \to 0
$$
是如上的极小消解。应用蛇引理可知
$$
0 \to (R/xR)^{n_e} \to (R/xR)^{n_{e-1}} \to \ldots \to (R/xR)^{n_0} \to
M/xM \to 0
$$
也是极小消解。因此 $\text{pd}_R(M) = \text{pd}_{R/xR}(M/xM)$。
由引理 \ref{algebra-lemma-depth-drops-by-one}，有
$\text{depth}(R/xR) = \text{depth}(R) - 1$ 和
$\text{depth}(M/xM) = \text{depth}(M) - 1$。
到此为止，所有深度都指作为 $R$-模的深度；但我们注意到
$\text{depth}_R(M/xM) = \text{depth}_{R/xR}(M/xM)$，对 $R/xR$ 也同样如此。
由归纳假设，$M/xM$ 作为 $R/xR$-模满足 Auslander--Buchsbaum 公式。
由于 $R/xR$ 和 $M/xM$ 的深度都减一，而投射维数没有改变，
结论随即得出。
\end{proof}

\section{同态与维数}
\label{algebra-section-homomorphism-dimension}

\noindent
本节汇集若干简单结果，用以联系存在环映射时各环的维数。

\begin{lemma}
\label{algebra-lemma-dimension-going-up}
设环映射 $R \to S$ 满足上升（见定义
\ref{algebra-definition-going-up-down}）或满足下降（见定义
\ref{algebra-definition-going-up-down}）。再假设
$\Spec(S) \to \Spec(R)$ 是满射。则 $\dim(R) \leq \dim(S)$。
\end{lemma}

\begin{proof}
先假设满足上升。任取 $R$ 中的素理想链
$\mathfrak p_0 \subset \mathfrak p_1 \subset \ldots
\subset \mathfrak p_e$。由满射性，可以选取映到
$\mathfrak p_0$ 的素理想 $\mathfrak q_0$。由上升，可以把它扩充为
长度为 $e$ 的素理想链，其中 $\mathfrak q_i$ 位于
$\mathfrak p_i$ 上方。因此 $\dim(S) \geq \dim(R)$。
下降情形完全相同。纯拓扑版本另见拓扑一章的引理
\ref{topology-lemma-dimension-specializations-lift}。
\end{proof}

\begin{lemma}
\label{algebra-lemma-going-up-maximal-on-top}
设 $R \to S$ 是满足上升性质的环映射，见定义
\ref{algebra-definition-going-up-down}。若
% P01-E0536: the source ends the conditional clause as a sentence before “Then”.
$\mathfrak q \subset S$ 是极大理想，则 $\mathfrak q$ 在 $R$
中的逆像也是极大理想。
\end{lemma}

\begin{proof}
显然。
\end{proof}

\begin{lemma}
\label{algebra-lemma-integral-dim-up}
设 $R \to S$ 是环映射，且 $S$ 在 $R$ 上整。则
$\dim (R) \geq \dim(S)$，并且 $\Spec(S)$ 的每个闭点都映到
$\Spec(R)$ 的闭点。
\end{lemma}

\begin{proof}
这由引理 \ref{algebra-lemma-integral-no-inclusion}、
\ref{algebra-lemma-going-up-maximal-on-top} 和各定义立即得出。
\end{proof}

\begin{lemma}
\label{algebra-lemma-integral-sub-dim-equal}
设 $R \subset S$，且 $S$ 在 $R$ 上整。则 $\dim(R) = \dim(S)$。
\end{lemma}

\begin{proof}
这是引理 \ref{algebra-lemma-integral-going-up}、
\ref{algebra-lemma-integral-overring-surjective}、
\ref{algebra-lemma-dimension-going-up} 和
\ref{algebra-lemma-integral-dim-up} 的组合。
\end{proof}

\begin{definition}
\label{algebra-definition-fibre}
设 $R \to S$ 是环映射。设 $\mathfrak q \subset S$ 是位于
$R$ 的素理想 $\mathfrak p$ 上方的素理想。所谓
{\it $\mathfrak q$ 处纤维的局部环}，是指局部环
$$
S_{\mathfrak q}/\mathfrak pS_{\mathfrak q}
=
(S/\mathfrak pS)_{\mathfrak q}
=
(S \otimes_R \kappa(\mathfrak p))_{\mathfrak q}
$$
\end{definition}

\begin{lemma}
\label{algebra-lemma-dimension-base-fibre-total}
设 $R \to S$ 是诺特环的同态。设 $\mathfrak q \subset S$
是位于素理想 $\mathfrak p$ 上方的素理想。则
$$
\dim(S_{\mathfrak q})
\leq
\dim(R_{\mathfrak p})
+
\dim(S_{\mathfrak q}/\mathfrak pS_{\mathfrak q}).
$$
\end{lemma}

\begin{proof}
我们使用命题 \ref{algebra-proposition-dimension} 对维数的刻画。
取 $\mathfrak p$ 中的元素 $x_1, \ldots, x_d$，使它们生成
$R_{\mathfrak p}$ 的一个定义理想，其中
$d = \dim(R_{\mathfrak p})$。取 $\mathfrak q$ 中的元素
$y_1, \ldots, y_e$，使它们生成
$S_{\mathfrak q}/\mathfrak pS_{\mathfrak q}$ 的一个定义理想，其中
$e = \dim(S_{\mathfrak q}/\mathfrak pS_{\mathfrak q})$。
显然，$S_{\mathfrak q}/(x_1, \ldots, x_d, y_1, \ldots, y_e)$
具有幂零极大理想。因此
$x_1, \ldots, x_d, y_1, \ldots, y_e$ 生成
$S_{\mathfrak q}$ 的一个定义理想。
\end{proof}

\begin{lemma}
\label{algebra-lemma-dimension-base-fibre-equals-total}
设 $R \to S$ 是诺特环的同态。设 $\mathfrak q \subset S$
是位于素理想 $\mathfrak p$ 上方的素理想。假设 $R \to S$
满足下降性质（例如，当 $R \to S$ 平坦时，引理
\ref{algebra-lemma-flat-going-down} 表明这一点）。则
$$
\dim(S_{\mathfrak q})
=
\dim(R_{\mathfrak p})
+
\dim(S_{\mathfrak q}/\mathfrak pS_{\mathfrak q}).
$$
\end{lemma}

\begin{proof}
由引理 \ref{algebra-lemma-dimension-base-fibre-total}，有不等式
$\dim(S_{\mathfrak q}) \leq
\dim(R_{\mathfrak p}) + \dim(S_{\mathfrak q}/\mathfrak pS_{\mathfrak q})$。
为了得到等式，选取素理想链
$\mathfrak pS \subset \mathfrak q_0 \subset \mathfrak q_1 \subset \ldots
\subset \mathfrak q_d = \mathfrak q$，其中
$d = \dim(S_{\mathfrak q}/\mathfrak pS_{\mathfrak q})$。
另一方面，选取素理想链
$\mathfrak p_0 \subset \mathfrak p_1 \subset \ldots \subset \mathfrak p_e
= \mathfrak p$，其中 $e = \dim(R_{\mathfrak p})$。
由下降定理，可以选取位于 $\mathfrak p_{e-1}$ 上方的
$\mathfrak q_{-1} \subset \mathfrak q_0$。继而可以选取位于
$\mathfrak p_{e-2}$ 上方的
$\mathfrak q_{-2} \subset \mathfrak q_{-1}$。依此归纳，最终得到长度为
$e + d$ 的链
$\mathfrak q_{-e} \subset \ldots \subset \mathfrak q_d$。
\end{proof}

\begin{lemma}
\label{algebra-lemma-flat-over-regular-with-regular-fibre}
设 $R \to S$ 是诺特局部环的局部同态。假设
\begin{enumerate}
\item $R$ 正则；
\item $S/\mathfrak m_RS$ 正则；
\item $R \to S$ 平坦。
\end{enumerate}
则 $S$ 正则。
\end{lemma}

\begin{proof}
由引理 \ref{algebra-lemma-dimension-base-fibre-equals-total}，有
$\dim(S) = \dim(R) + \dim(S/\mathfrak m_RS)$。
选取生成元 $x_1, \ldots, x_d \in \mathfrak m_R$，其中
$d = \dim(R)$；再选取 $y_1, \ldots, y_e \in \mathfrak m_S$，
使其在 $S/\mathfrak m_RS$ 中生成该环的极大理想，其中
$e = \dim(S/\mathfrak m_RS)$。于是
$x_1, \ldots, x_d, y_1, \ldots, y_e$ 生成 $S$ 的极大理想，
并且 $e + d = \dim(S)$。
\end{proof}

\noindent
下面的引理稍后将用于证明：域上的有限型环是 Cohen--Macaulay 环，
当且仅当它在某个多项式环上拟有限且平坦。它是引理
\ref{algebra-lemma-CM-over-regular-flat} 的一个部分逆命题。

\begin{lemma}
\label{algebra-lemma-finite-flat-over-regular-CM}
设 $R \to S$ 是诺特局部环的局部同态。假设 $R$ 是
Cohen--Macaulay 环。若 $S$ 在 $R$ 上有限平坦，或者 $S$ 在
$R$ 上平坦且 $\dim(S) \leq \dim(R)$，则 $S$ 是 Cohen--Macaulay 环，
并且 $\dim(R) = \dim(S)$。
\end{lemma}

\begin{proof}
设 $x_1, \ldots, x_d \in \mathfrak m_R$ 是长度为
$d = \dim(R)$ 的正则序列。由引理 \ref{algebra-lemma-flat-increases-depth}，
它在 $S$ 中的像是正则序列。因此，只要 $\dim(S) \leq d$，
$S$ 就是 Cohen--Macaulay 环。若 $S$ 在 $R$ 上有限平坦，
则由引理 \ref{algebra-lemma-integral-sub-dim-equal} 可知这一不等式成立；
第二种情形中，该不等式就是假设。
\end{proof}

\section{维数公式}
\label{algebra-section-dimension-formula}

\noindent
回顾链状（定义 \ref{algebra-definition-catenary}）和普遍链状
（定义 \ref{algebra-definition-universally-catenary}）的定义。

\begin{lemma}
\label{algebra-lemma-dimension-formula}
设 $R \to S$ 是环映射。设 $\mathfrak q$ 是 $S$ 的素理想，
位于 $R$ 的素理想 $\mathfrak p$ 上方。假设
\begin{enumerate}
\item $R$ 是诺特环；
\item $R \to S$ 是有限型的；
\item $R$、$S$ 都是整环；
\item $R \subset S$。
\end{enumerate}
则
$$
\text{高度}(\mathfrak q)
\leq
\text{高度}(\mathfrak p) + \text{trdeg}_R(S)
- \text{trdeg}_{\kappa(\mathfrak p)} \kappa(\mathfrak q)
$$
成立；若 $R$ 普遍链状，则等号成立。
\end{lemma}

\begin{proof}
设 $R \subset S' \subset S$，其中 $S'$ 是 $S$ 的有限生成
$R$-子代数。此时令 $\mathfrak q' = S' \cap \mathfrak q$。
对环映射 $R \to S'$ 和 $S' \to S$ 应用本引理，再使用域塔中
超越次数的可加性（域一章的引理
\ref{fields-lemma-transcendence-degree-tower}），即可推出
$R \to S$ 的情形。因此，可以对 $S$ 在 $R$ 上的生成元数目作归纳，
把问题约化到 $S$ 由 $R$ 上一个元素生成的情形。

\medskip\noindent
情形 I：$S = R[x]$ 是 $R$ 上的多项式代数。
此时 $\text{trdeg}_R(S) = 1$。此外，$R \to S$ 平坦，因而
$$
\dim(S_{\mathfrak q}) =
\dim(R_{\mathfrak p}) + \dim(S_{\mathfrak q}/\mathfrak pS_{\mathfrak q})
$$
，见引理 \ref{algebra-lemma-dimension-base-fibre-equals-total}。
令 $\mathfrak r = \mathfrak pS$。那么
$\text{trdeg}_{\kappa(\mathfrak p)} \kappa(\mathfrak q) = 1$
等价于 $\mathfrak q = \mathfrak r$，并且蕴含
$\dim(S_{\mathfrak q}/\mathfrak pS_{\mathfrak q}) = 0$。
同理，$\text{trdeg}_{\kappa(\mathfrak p)} \kappa(\mathfrak q) = 0$
等价于存在严格包含
$\mathfrak r \subset \mathfrak q$，而这蕴含
$\dim(S_{\mathfrak q}/\mathfrak pS_{\mathfrak q}) = 1$。
至此情形 I 得证，并且每种情况下等号都成立。

\medskip\noindent
情形 II：$S = R[x]/\mathfrak n$，其中 $\mathfrak n \not = 0$。
此时 $\text{trdeg}_R(S) = 0$。
% P01-E0537: the source omits “by” in “Denote q' ... the prime corresponding to q”.
以 $\mathfrak q' \subset R[x]$ 表示与 $\mathfrak q$ 对应的素理想。于是
$$
S_{\mathfrak q} = (R[x])_{\mathfrak q'}/\mathfrak n(R[x])_{\mathfrak q'}
$$
由上一情形，有
$\dim((R[x])_{\mathfrak q'}) =
\dim(R_{\mathfrak p}) + 1
- \text{trdeg}_{\kappa(\mathfrak p)} \kappa(\mathfrak q)$。
由于 $\mathfrak n \not = 0$，可知 $S_{\mathfrak q}$ 的维数至少减一，
见引理 \ref{algebra-lemma-one-equation}；这证明了引理中的不等式。
% P01-E0538: the source phrase “in case R is universally catenary” is missing the article “the”.
为了在 $R$ 普遍链状时证明等号，注意到
$\mathfrak n \subset R[x]$ 是高度为一的素理想，因为它对应于
$F[x]$ 中的非零素理想，其中 $F$ 是 $R$ 的分式域。因此，
$S_\mathfrak q = R[x]_{\mathfrak q'}/\mathfrak nR[x]_{\mathfrak q'}$
中的任意极大素理想链，都对应于 $R[x]$ 中 $\mathfrak q'$ 与 $(0)$
之间长度大一的极大素理想链。
% P01-E0539: the source's “with length 1 greater between q' and (0)” is grammatically misordered.
若 $R$ 普遍链状，则这些链的长度全都相同，并等于
$\mathfrak q'$ 的高度。这证明
$\dim(S_\mathfrak q) = \dim(R[x]_{\mathfrak q'}) - 1$，
进而按所需得到等号。
\end{proof}

\noindent
下面的引理说明，一般有限映射往往在余维 $1$ 处拟有限。

\begin{lemma}
\label{algebra-lemma-finite-in-codim-1}
设 $A \to B$ 是环映射。假设
\begin{enumerate}
\item $A \subset B$ 是整环的扩张；
\item 所诱导的分式域扩张是有限的；
\item $A$ 是诺特环；
\item $A \to B$ 是有限型的。
\end{enumerate}
设 $\mathfrak p \subset A$ 是高度为 $1$ 的素理想。则 $B$ 中位于
$\mathfrak p$ 上方的素理想至多只有有限多个，并且它们的高度全都为 $1$。
\end{lemma}

\begin{proof}
由维数公式（引理 \ref{algebra-lemma-dimension-formula}），对位于
$\mathfrak p$ 上方的任意素理想 $\mathfrak q$，有
$$
\dim(B_{\mathfrak q}) \leq
\dim(A_{\mathfrak p}) - \text{trdeg}_{\kappa(\mathfrak p)} \kappa(\mathfrak q).
$$
由于整环 $B_\mathfrak q$ 至少有 $2$ 个素理想，可知
$\dim(B_{\mathfrak q}) \geq 1$。因此
$\dim(B_{\mathfrak q}) = 1$，并且扩张
$\kappa(\mathfrak p) \subset \kappa(\mathfrak q)$ 是代数的。
所以，$\mathfrak q$ 给出其纤维
$\Spec(B \otimes_A \kappa(\mathfrak p))$ 的一个闭点，见引理
\ref{algebra-lemma-finite-residue-extension-closed}。
由于 $B \otimes_A \kappa(\mathfrak p)$ 是诺特环，该纤维
$\Spec(B \otimes_A \kappa(\mathfrak p))$ 是诺特拓扑空间，见引理
\ref{algebra-lemma-Noetherian-topology}。只由闭点组成的诺特拓扑空间是有限的，
例如见拓扑一章的引理 \ref{topology-lemma-Noetherian}。
\end{proof}

\section{域上有限型代数的维数}
\label{algebra-section-dimension-finite-type-algebras}

\noindent
本节计算域上多项式环的维数。我们还将证明，域上的有限型整环的维数
等于它在极大理想处各局部环的维数。它与基域上超越次数的联系将在
第 \ref{algebra-section-dimension-finite-type-algebras-reprise} 节建立。

\begin{lemma}
\label{algebra-lemma-dim-affine-space}
设 $\mathfrak m$ 是 $k[x_1, \ldots, x_n]$ 中的极大理想。
理想 $\mathfrak m$ 可由 $n$ 个元素生成。
$k[x_1, \ldots, x_n]_{\mathfrak m}$ 的维数为 $n$。
因此，$k[x_1, \ldots, x_n]_{\mathfrak m}$ 是维数为 $n$ 的正则局部环。
\end{lemma}

\begin{proof}
由 Hilbert 零点定理（定理 \ref{algebra-theorem-nullstellensatz}），
剩余域 $\kappa = \kappa(\mathfrak m)$ 是 $k$ 的有限扩张。
% P01-E0540: the source omits “by” in “Denote alpha_i ... the image of x_i”.
以 $\alpha_i \in \kappa$ 表示 $x_i$ 的像。
% P01-E0541: the source uses the ungrammatical “Denote kappa_i = ...” construction.
记 $\kappa_i = k(\alpha_1, \ldots, \alpha_i)
\subset \kappa$，其中 $i = 1, \ldots, n$，并令 $\kappa_0 = k$。
由域论可知 $\kappa_i = k[\alpha_1, \ldots, \alpha_i]$。
按如下方式归纳定义元素
$f_i \in \mathfrak m \cap k[x_1, \ldots, x_i]$：
设 $P_i(T) \in \kappa_{i-1}[T]$ 是 $\alpha_i $ 在
$\kappa_{i-1}$ 上的首一极小多项式。设
$Q_i(T) \in k[x_1, \ldots, x_{i-1}][T]$ 是 $P_i(T)$
的一个首一提升（次数相同）。令 $f_i = Q_i(x_i)$。
注意，若
$d_i = \deg_T(P_i) = \deg_T(Q_i) = \deg_{x_i}(f_i)$，则由域一章的引理
\ref{fields-lemma-multiplicativity-degrees} 和
\ref{fields-lemma-degree-minimal-polynomial}，
有 $d_1d_2\ldots d_i = [\kappa_i : k]$。

\medskip\noindent
我们断言，对所有 $i = 0, 1, \ldots, n$，都存在同构
$$
\psi_i : k[x_1, \ldots, x_i] /(f_1, \ldots, f_i) \cong \kappa_i.
$$
按构造，复合映射
$k[x_1, \ldots, x_i] \to k[x_1, \ldots, x_n] \to \kappa$
满射到 $\kappa_i$，并且 $f_1, \ldots, f_i$ 都在核中。
这给出一个满同态。我们归纳证明 $\psi_i$ 是单射。
$i = 0$ 时显然成立。假设结论对 $i$ 成立，我们来证明
$i + 1$ 的情形。环扩张
$k[x_1, \ldots, x_i]/(f_1, \ldots, f_i) \to
k[x_1, \ldots, x_{i + 1}]/(f_1, \ldots, f_{i + 1})$
是在一个域上添加 $1$ 个元素并施加一个不可约方程而生成的。
由初等域论，
$k[x_1, \ldots, x_{i + 1}]/(f_1, \ldots, f_{i + 1})$
是域，因而
% P01-E0542: the induction step proves injectivity of psi_{i+1}, but the source says psi_i.
$\psi_i$ 是单射。

\medskip\noindent
这蕴含 $\mathfrak m = (f_1, \ldots, f_n)$。此外，还可推出
$$
k[x_1, \ldots, x_n]/(f_1, \ldots, f_i)
\cong
\kappa_i[x_{i + 1}, \ldots, x_n].
$$
因此 $(f_1, \ldots, f_i)$ 是素理想。于是
$$
(0) \subset (f_1) \subset (f_1, f_2) \subset \ldots \subset
(f_1, \ldots, f_n) = \mathfrak m
$$
是长度为 $n$ 的素理想链。引理得证。
\end{proof}

\begin{proposition}
\label{algebra-proposition-finite-gl-dim-polynomial-ring}
域上的 $n$ 元多项式代数是正则环。它的整体维数为 $n$。
所有在极大理想处的局部化都是维数为 $n$ 的正则局部环。
\end{proposition}

\begin{proof}
由引理 \ref{algebra-lemma-dim-affine-space}，在极大理想处的所有局部化
$k[x_1, \ldots, x_n]_{\mathfrak m}$ 都是维数为 $n$ 的正则局部环。
因此，由引理 \ref{algebra-lemma-finite-gl-dim-finite-dim-regular} 得出结论。
\end{proof}

\begin{lemma}
\label{algebra-lemma-dimension-height-polynomial-ring}
设 $k$ 是域。设
$\mathfrak p \subset \mathfrak q \subset k[x_1, \ldots, x_n]$
是一对素理想。$\mathfrak p$ 与 $\mathfrak q$ 之间的任意极大素理想链
长度均为 $\text{高度}(\mathfrak q) - \text{高度}(\mathfrak p)$。
\end{lemma}

\begin{proof}
由命题 \ref{algebra-proposition-finite-gl-dim-polynomial-ring}，
$k[x_1, \ldots, x_n]$ 的任意局部环都是正则的。
因此，所有局部环都是 Cohen--Macaulay 环，见引理
\ref{algebra-lemma-regular-ring-CM}。极大理想处的局部环维数为 $n$，
所以 $k[x_1, \ldots, x_n]$ 中的每条极大素理想链长度均为 $n$，
见引理 \ref{algebra-lemma-maximal-chain-CM}。于是，$(0)$ 与 $\mathfrak p$
之间的每条极大素理想链长度均为 $\text{高度}(\mathfrak p)$，
例如见引理 \ref{algebra-lemma-CM-dim-formula}。把这些结论合在一起，
即可得到本引理的断言。
\end{proof}

\begin{lemma}
\label{algebra-lemma-dimension-spell-it-out}
设 $k$ 是域。设 $S$ 是有限型 $k$-代数，并且是整环。
则对 $S$ 的任意极大理想 $\mathfrak m$，都有
$\dim(S) = \dim(S_{\mathfrak m})$。换言之，每条极大素理想链的长度
都等于 $S$ 的维数。
\end{lemma}

\begin{proof}
写成 $S = k[x_1, \ldots, x_n]/\mathfrak p$。
由命题 \ref{algebra-proposition-finite-gl-dim-polynomial-ring} 和引理
\ref{algebra-lemma-dimension-height-polynomial-ring}，$S$ 中的所有极大素理想链
（它们必然终止于极大理想）的长度都是
$n - \text{高度}(\mathfrak p)$。因此，这个数既是 $S$ 的维数，
也是 $S$ 的任意极大理想 $\mathfrak m$ 所对应的
$S_{\mathfrak m}$ 的维数。
\end{proof}

\noindent
回顾我们在拓扑一章的定义 \ref{topology-definition-Krull} 中定义了
拓扑空间 $X$ 在点 $x$ 处的维数 $\dim_x(X)$。

\begin{lemma}
\label{algebra-lemma-dimension-at-a-point-finite-type-over-field}
设 $k$ 是域。设 $S$ 是有限型 $k$-代数。令 $X = \Spec(S)$。
设 $\mathfrak p \subset S$ 是素理想，并令 $x \in X$ 为对应的点。
以下各数相等：
\begin{enumerate}
\item $\dim_x(X)$；
\item $\max \dim(Z)$，其中最大值取遍所有经过 $x$ 的
$X$ 的不可约分支 $Z$；
\item $\min \dim(S_{\mathfrak m})$，其中最小值取遍所有满足
$\mathfrak p \subset \mathfrak m$ 的极大理想 $\mathfrak m$。
\end{enumerate}
\end{lemma}

\begin{proof}
设 $X = \bigcup_{i \in I} Z_i$ 是 $X$ 的不可约分支分解。
这些分支只有有限多个（见引理 \ref{algebra-lemma-obvious-Noetherian} 和
\ref{algebra-lemma-Noetherian-topology}）。令
$I' = \{i \mid x \in Z_i\}$，并令
$T = \bigcup_{i \not \in I'} Z_i$。于是 $U = X \setminus T$
是 $X$ 中包含点 $x$ 的开子集。数 (2) 是
$\max_{i \in I'} \dim(Z_i)$。对任意满足 $x \in W$ 的开集
$W \subset U$，$W$ 的不可约分支是不可约集
$W_i = Z_i \cap W$，其中 $i \in I'$，并且 $x$ 属于其中每一个。
注意，每个 $W_i$（$i \in I'$）都包含闭点，因为 $X$ 是 Jacobson 的，
见第 \ref{algebra-section-ring-jacobson} 节。由于 $W_i \subset Z_i$，
有 $\dim(W_i) \leq \dim(Z_i)$。闭点的存在结合引理
\ref{algebra-lemma-dimension-spell-it-out} 表明，开集 $W_i$ 中存在长度为
$\dim(Z_i)$ 的不可约闭子集链。因此，对任意 $i \in I'$，
都有 $\dim(W_i) = \dim(Z_i)$。故 $\dim(W)$ 等于数 (2)。
这证明 (1) $ = $ (2)。

\medskip\noindent
设 $\mathfrak m \supset \mathfrak p$ 是包含 $\mathfrak p$
的任意极大理想。令 $x_0 \in X$ 为对应的点。首先，$x_0$
包含在所有不可约分支 $Z_i$（$i \in I'$）中。
令 $\mathfrak q_i$ 表示 $S$ 中对应于不可约分支 $Z_i$ 的极小素理想。
对每个满足 $x_0 \in Z_i$ 的 $i$（这等价于
$\mathfrak m \supset \mathfrak q_i$），都有满射
$$
S_{\mathfrak m} \longrightarrow
S_\mathfrak m/\mathfrak q_i S_\mathfrak m =(S/\mathfrak q_i)_{\mathfrak m}
$$
此外，这样得到的素理想 $\mathfrak q_i S_\mathfrak m$
穷尽诺特局部环 $S_{\mathfrak m}$ 的极小素理想，见引理
\ref{algebra-lemma-irreducible-components-containing-x}。
由引理 \ref{algebra-lemma-dimension-spell-it-out}，可知
$S_{\mathfrak m}$ 的维数是经过 $x_0$ 的各 $Z_i$ 的维数中的最大值。
为完成引理的证明，只需说明可以选取 $x_0$，使
$x_0 \in Z_i \Rightarrow i \in I'$。由于 $S$ 是 Jacobson 环
（如上所见），只需说明 $V(\mathfrak p) \setminus T$
（其中 $T$ 同上）非空。这是显然的，因为它包含点 $x$
（即 $\mathfrak p$）。
\end{proof}

\begin{lemma}
\label{algebra-lemma-dimension-closed-point-finite-type-field}
设 $k$ 是域。设 $S$ 是有限型 $k$-代数。令 $X = \Spec(S)$。
设 $\mathfrak m \subset S$ 是极大理想，并令 $x \in X$
为相应的闭点。则 $\dim_x(X) = \dim(S_{\mathfrak m})$。
\end{lemma}

\begin{proof}
这是引理 \ref{algebra-lemma-dimension-at-a-point-finite-type-over-field}
的一个特例。
\end{proof}

\begin{lemma}
\label{algebra-lemma-disjoint-decomposition-CM-algebra}
设 $k$ 是域。设 $S$ 是有限型
% P01-E0543: the source omits the hyphen in “finite type k algebra”.
$k$-代数。假设 $S$ 是 Cohen--Macaulay 环。则
$\Spec(S) = \coprod T_d$ 是开闭子集 $T_d$ 的有限不交并，
其中 $T_d$ 是维数为 $d$ 的等维空间（见拓扑一章的定义
\ref{topology-definition-equidimensional}）。等价地，$S$ 是环
$S_d$（$d = 0, \ldots, \dim(S)$）的乘积，并且 $S_d$
的每个极大理想 $\mathfrak m$ 都具有高度 $d$。
\end{lemma}

\begin{proof}
两个陈述的等价性由引理 \ref{algebra-lemma-disjoint-implies-product} 得出。
设 $\mathfrak m \subset S$ 是极大理想。$S_{\mathfrak m}$
中的每条极大素理想链长度都相同，并等于
$\dim(S_{\mathfrak m})$，见引理 \ref{algebra-lemma-maximal-chain-CM}。
因此，经过与 $\mathfrak m$ 对应之点的不可约分支
% P01-E0544: the source has singular “the dimension ... all have”.
维数全都等于 $\dim(S_{\mathfrak m})$，见引理
\ref{algebra-lemma-dimension-spell-it-out}。由于 $\Spec(S)$ 是
Jacobson 拓扑空间，它的任意两个不可约分支若相交，其交集便包含闭点，
见引理 \ref{algebra-lemma-finite-type-field-Jacobson} 和
\ref{algebra-lemma-jacobson}。所以，我们已经证明任意两个相交的不可约分支
具有相同维数。再结合 $\Spec(S)$ 只有有限多个不可约分支这一事实，
即可容易地推出本引理；后者见引理 \ref{algebra-lemma-obvious-Noetherian}
和 \ref{algebra-lemma-Noetherian-topology}。
\end{proof}

\section{Noether 正规化}
\label{algebra-section-Noether-normalization}

\noindent
本节证明 Noether 正规化引理的若干变体。我们将使用的关键成分包含在
下面两个引理中。

\begin{lemma}
\label{algebra-lemma-helper}
设 $n \in \mathbf{N}$。设 $N$ 是由多重指标
$\nu = (\nu_1, \ldots, \nu_n)$ 组成的有限非空集合。
给定 $e = (e_1, \ldots, e_n)$，令
$e \cdot \nu = \sum e_i\nu_i$。则当
$e_1 \gg e_2 \gg \ldots \gg e_{n-1} \gg e_n$ 时，有：
若 $\nu, \nu' \in N$，则
$$
(e \cdot \nu = e \cdot \nu')
\Leftrightarrow
(\nu = \nu')
$$
\end{lemma}

\begin{proof}
设 $N = \{\nu_j\}$，其中
$\nu_j = (\nu_{j1}, \ldots, \nu_{jn})$。令
$A_i = \max_j \nu_{ji} - \min_j \nu_{ji}$。
% P01-E0545: “for each i” makes the displayed sufficient condition use the undefined e_0 at i = 1; the intended range starts at i = 2.
若对每个 $i$ 都有
$e_{i - 1} > A_ie_i + A_{i + 1}e_{i + 1} + \ldots + A_ne_n$，
则引理成立。事实上，假设 $e \cdot (\nu - \nu') = 0$。当
$n \ge 2$ 时，
$$
e_1(\nu_1 - \nu'_1) = \sum\nolimits_{i = 2}^n e_i(\nu'_i - \nu_i).
$$
可以假设 $(\nu_1 - \nu'_1) \ge 0$。若
$(\nu_1 - \nu'_1) > 0$，则
$$
e_1(\nu_1 - \nu'_1) \ge e_1 >
A_2e_2 + \ldots + A_ne_n \ge
\sum\nolimits_{i = 2}^n e_i|\nu'_i - \nu_i| \ge
\sum\nolimits_{i = 2}^n e_i(\nu'_i - \nu_i).
$$
这个矛盾说明 $\nu'_1 = \nu_1$。归纳可得，对
$2 \le i \le n$，都有 $\nu'_i = \nu_i$。
\end{proof}

\begin{lemma}
\label{algebra-lemma-helper-polynomial}
设 $R$ 是环。设 $g \in R[x_1, \ldots, x_n]$ 是非常数元素，
即 $g \not \in R$。当
$e_1 \gg e_2 \gg \ldots \gg e_{n-1} \gg e_n = 1$ 时，多项式
$$
g(x_1 + x_n^{e_1}, x_2 + x_n^{e_2}, \ldots, x_{n - 1} + x_n^{e_{n - 1}}, x_n)
=
ax_n^d + \text{lower order terms in }x_n
$$
，其中 $d > 0$，而 $a \in R$ 是 $g$ 的一个非零系数。
\end{lemma}

\begin{proof}
写成 $g = \sum_{\nu \in N} a_\nu x^\nu$，其中 $a_\nu \in R$
% P01-E0546: the source uses the nonstandard phrase “not zero” for a coefficient.
非零。这里 $N$ 是引理 \ref{algebra-lemma-helper} 中那样的有限多重指标集，
并且 $x^\nu = x_1^{\nu_1} \ldots x_n^{\nu_n}$。注意，
$$
(x_1 + x_n^{e_1})^{\nu_1} \ldots (x_{n-1} + x_n^{e_{n-1}})^{\nu_{n-1}}
x_n^{\nu_n}
\quad\text{is}\quad
x_n^{e_1\nu_1 + \ldots + e_{n-1}\nu_{n-1} + \nu_n}.
$$
中的首项如上所示。因此，引理由引理 \ref{algebra-lemma-helper} 得出：
它保证 $g$ 中恰有一个非零项 $a_\nu x^\nu$ 产生
$g(x_1 + x_n^{e_1}, x_2 + x_n^{e_2}, \ldots, x_{n - 1} + x_n^{e_{n - 1}},
x_n)$ 的首项；也就是说，对唯一满足 $e \cdot \nu$ 取最大值的
$\nu \in N$，有 $a = a_\nu$。
\end{proof}

\begin{lemma}
\label{algebra-lemma-one-relation}
设 $k$ 是域。设 $S = k[x_1, \ldots, x_n]/I$，其中 $I$ 是真理想。
若 $I \not = 0$，则存在
$y_1, \ldots, y_{n-1} \in k[x_1, \ldots, x_n]$，使得 $S$ 在
$k[y_1, \ldots, y_{n-1}]$ 上有限。此外，可以选取 $y_i$ 属于由
$x_1, \ldots, x_n$ 生成的 $k[x_1, \ldots, x_n]$ 的
$\mathbf{Z}$-子代数。
\end{lemma}

\begin{proof}
选取 $f \in I$，$f\not = 0$。因为 $S$ 是
$k[x_1, \ldots, x_n]/(f)$ 的商环，只需对后一个环证明本引理。
对适当的整数 $e_i$，我们将取
$y_i = x_i - x_n^{e_i}$，$i = 1, \ldots, n-1$。
何时这会奏效？只需证明
$\overline{x_n} \in k[x_1, \ldots, x_n]/(f)$ 在环
$k[y_1, \ldots, y_{n-1}]$ 上整。$\overline{x_n}$ 在此环上满足的方程是
$$
f(y_1 + x_n^{e_1}, \ldots, y_{n-1} + x_n^{e_{n-1}}, x_n) = 0.
$$
所以，只要能证明
% P01-E0547: the source has the subject–verb mismatch “there exists integers”.
存在整数 $e_i$，使上述方程关于 $x_n$ 的首项系数是 $k$ 的非零元素，
就完成了证明。例如，取
$e_1 \gg e_2 \gg \ldots \gg e_{n-1}$ 即可，见引理
\ref{algebra-lemma-helper-polynomial}。
\end{proof}

\begin{lemma}
\label{algebra-lemma-Noether-normalization}
\begin{slogan}
Noether 正规化
\end{slogan}
设 $k$ 是域。设 $S = k[x_1, \ldots, x_n]/I$，其中 $I$ 是理想。
若 $I \neq (1)$，则存在 $r\geq 0$ 和
$y_1, \ldots, y_r \in k[x_1, \ldots, x_n]$，使得：
(a) 映射 $k[y_1, \ldots, y_r] \to S$ 是单射，其中源是以
$y_1, \ldots, y_r$ 为变量的多项式环；(b) 映射
$k[y_1, \ldots, y_r] \to S$ 是有限的。在此情形下，整数 $r$
就是 $S$ 的维数。此外，可以选取 $y_i$ 属于由
$x_1, \ldots, x_n$ 生成的 $k[x_1, \ldots, x_n]$ 的
$\mathbf{Z}$-子代数。
\end{lemma}

\begin{proof}
对 $n$ 作归纳，$n = 0$ 时显然。若 $I = 0$，则取
$r = n$ 和 $y_i = x_i$。若 $I \not = 0$，则按引理
\ref{algebra-lemma-one-relation} 选取 $y_1, \ldots, y_{n-1}$。令
$S' \subset S$ 为由各 $y_i$ 的像生成的子环。由归纳假设，
可以选取 $r$ 和
$z_1, \ldots, z_r \in k[y_1, \ldots, y_{n-1}]$，使 (a)、(b)
对 $k[z_1, \ldots, z_r] \to S'$ 成立。由于 $S' \to S$
是单射且有限，可知 (a)、(b) 对
$k[z_1, \ldots, z_r] \to S$ 也成立。
% P01-E0548: “The last assertion” literally points to the Z-subalgebra choice, but the cited dimension lemma proves the preceding assertion that r = dim(S).
关于维数的断言由引理 \ref{algebra-lemma-integral-sub-dim-equal} 得出。
\end{proof}

\begin{lemma}
\label{algebra-lemma-Noether-normalization-at-point}
设 $k$ 是域。设 $S$ 是有限型
% P01-E0549: the source omits the hyphen in “finite type k algebra”.
$k$-代数，并记 $X = \Spec(S)$。设 $\mathfrak q$ 是 $S$ 的素理想，
并令 $x \in X$ 为对应的点。存在 $g \in S$、$g \not \in \mathfrak q$，
使得 $\dim(S_g) = \dim_x(X) =: d$，并且存在有限单射环映射
$k[y_1, \ldots, y_d] \to S_g$。
\end{lemma}

\begin{proof}
注意，按定义，$\dim_x(X)$ 是所有 $S_g$ 的维数中的最小值，
其中 $g \in S$、$g \not \in \mathfrak q$；也就是说，该最小值能够取到。
因此，引理由引理 \ref{algebra-lemma-Noether-normalization} 得出。
\end{proof}

\begin{lemma}
\label{algebra-lemma-refined-Noether-normalization}
设 $k$ 是域。设 $\mathfrak q \subset k[x_1, \ldots, x_n]$
是素理想。令 $r = \text{trdeg}_k\ \kappa(\mathfrak q)$。
则存在有限环映射
$\varphi : k[y_1, \ldots, y_n] \to k[x_1, \ldots, x_n]$，使得
$\varphi^{-1}(\mathfrak q) = (y_{r + 1}, \ldots, y_n)$。
\end{lemma}

\begin{proof}
对 $n$ 作归纳。$n = 0$ 时显然。假设 $n > 0$。
若 $r = n$，则 $\mathfrak q = (0)$，结论显然。
选取非零的 $f \in \mathfrak q$。当然，$f$ 不是常数。
应用如下形式的自同构后
$$
k[x_1, \ldots, x_n] \longrightarrow k[x_1, \ldots, x_n],
\quad
x_n \mapsto x_n,
\quad
x_i \mapsto x_i + x_n^{e_i}\ (i < n)
$$
，可以假设
% P01-E0550: the source says f is monic in x_n over a coefficient ring that itself contains x_n; the intended coefficient ring is k[x_1, ..., x_{n-1}].
$f$ 在 $k[x_1, \ldots, x_n]$ 上关于 $x_n$ 首一，见引理
\ref{algebra-lemma-helper-polynomial}。因此，环映射
$$
k[y_1, \ldots, y_n] \longrightarrow k[x_1, \ldots, x_n],
\quad
y_n \mapsto f,
\quad
y_i \mapsto x_i\ (i < n)
$$
是有限的。此外，按构造，
$y_n \in \mathfrak q \cap k[y_1, \ldots, y_n]$。因此
$\mathfrak q \cap k[y_1, \ldots, y_n] = \mathfrak pk[y_1, \ldots, y_n] + (y_n)$，
其中 $\mathfrak p \subset k[y_1, \ldots, y_{n - 1}]$ 是素理想。
注意，$\kappa(\mathfrak p) \subset \kappa(\mathfrak q)$ 是有限扩张，
因而 $r = \text{trdeg}_k\ \kappa(\mathfrak p)$。
对二元组 $(k[y_1, \ldots, y_{n - 1}], \mathfrak p)$ 应用归纳假设，
得到有限环映射
$k[z_1, \ldots, z_{n - 1}] \to k[y_1, \ldots, y_{n - 1}]$，使得
$\mathfrak p \cap k[z_1, \ldots, z_{n - 1}] = (z_{r + 1}, \ldots, z_{n - 1})$。
把环映射
$k[z_1, \ldots, z_{n - 1}] \to k[y_1, \ldots, y_{n - 1}]$
扩张为环映射
$k[z_1, \ldots, z_n] \to k[y_1, \ldots, y_n]$，方法是把
$z_n$ 映到 $y_n$。环映射的复合
$$
k[z_1, \ldots, z_n] \to k[y_1, \ldots, y_n] \to k[x_1, \ldots, x_n]
$$
解决了问题。
\end{proof}

\begin{lemma}
\label{algebra-lemma-Noether-normalization-over-a-domain}
设 $R \to S$ 是单射有限型环映射，并假设 $R$ 是整环。
则存在整数 $d$ 以及由单射组成的分解
$$
R \to R[y_1, \ldots, y_d] \to S' \to S
$$
，使得 $S'$ 在 $R[y_1, \ldots, y_d]$ 上有限，并且对某个
非零的 $f \in R$，有 $S'_f \cong S_f$。
\end{lemma}

\begin{proof}
选取 $x_1, \ldots, x_n \in S$，使其在 $R$ 上生成 $S$。
令 $K$ 为 $R$ 的分式域，并令 $S_K = S \otimes_R K$。
由引理 \ref{algebra-lemma-Noether-normalization}，可以找到
$y_1, \ldots, y_d \in S$，使
$K[y_1, \ldots, y_d] \to S_K$ 是有限单射环映射。
注意，$y_i \in S$，因为可以选取
% P01-E0551: the source switches from y_i to y_j in the explanation without a change of quantified index.
$y_j$ 属于由 $x_1, \ldots, x_n$ 生成的 $\mathbf{Z}$-代数。
由于有限环映射是整的（见引理 \ref{algebra-lemma-finite-is-integral}），
可以找到首一多项式 $P_i \in K[y_1, \ldots, y_d][T]$，使得
$P_i(x_i) = 0$ 在 $S_K$ 中成立。取非零元素 $f \in R$，使得对所有
$i$ 都有 $fP_i \in R[y_1, \ldots, y_d][T]$。于是
$fP_i(x_i)$ 在 $S_K$ 中映到零。因此，把 $f$ 换成 $R$
的另一个非零元素后，还可以假设 $fP_i(x_i)$ 在 $S$ 中为零。
令 $x_i' = fx_i$，并令 $S' \subset S$ 为由
$y_1, \ldots, y_d$ 和 $x'_1, \ldots, x'_n$ 生成的
$R$-子代数。注意，$x'_i$ 在 $R[y_1, \ldots, y_d]$ 上整，
因为 $Q_i(x_i') = 0$，其中
$Q_i = f^{\deg_T(P_i)}P_i(T/f)$；按 $f$ 的选取，这是系数属于
$R[y_1, \ldots, y_d]$ 的关于 $T$ 的首一多项式。因此，由引理
\ref{algebra-lemma-characterize-finite-in-terms-of-integral}，
$R[y_1, \ldots, y_d] \subset S'$ 是有限的。
因为 $S' \subset S$，所以 $S'_f \subset S_f$（局部化是正合的）。
另一方面，$S'_f$ 中的元素 $x_i = x'_i/f$ 在 $R_f$ 上生成
$S_f$，故 $S'_f \to S_f$ 是满射。因此
$S'_f \cong S_f$，证明完成。
\end{proof}

\section{域上有限型代数的维数（再论）}
\label{algebra-section-dimension-finite-type-algebras-reprise}

\noindent
本节是第 \ref{algebra-section-dimension-finite-type-algebras} 节的延续。
我们将在本节建立域上的有限型整环之维数与其在基域上的超越次数之间的联系。

\begin{lemma}
\label{algebra-lemma-dimension-prime-polynomial-ring}
设 $k$ 是域。设 $S$ 是有限型
% P01-E0552: the source omits the hyphen in “finite type k algebra”.
$k$-代数，并且是整环。设 $K$ 是 $S$ 的分式域。
设 $r = \text{trdeg}(K/k)$ 是 $K$ 在 $k$ 上的超越次数。
则 $\dim(S) = r$。此外，$S$ 在每个极大理想处的局部环维数均为 $r$。
\end{lemma}

\begin{proof}
可以写成 $S = k[x_1, \ldots, x_n]/\mathfrak p$。
由引理 \ref{algebra-lemma-dimension-height-polynomial-ring}，$S$
在各极大理想处的所有局部环维数相同。应用引理
\ref{algebra-lemma-Noether-normalization}，得到有限单射环映射
$$
k[y_1, \ldots, y_d] \to S
$$
，其中 $d = \dim(S)$。显然，
$k(y_1, \ldots, y_d) \subset K$ 是有限扩张，结论得证。
\end{proof}

\begin{lemma}
\label{algebra-lemma-tr-deg-specialization}
设 $k$ 是域。设 $S$ 是有限型 $k$-代数。设
$\mathfrak q \subset \mathfrak q' \subset S$ 是不同的素理想。则
$\text{trdeg}_k\ \kappa(\mathfrak q') < \text{trdeg}_k\ \kappa(\mathfrak q)$。
\end{lemma}

\begin{proof}
由引理 \ref{algebra-lemma-dimension-prime-polynomial-ring}，有
$\dim V(\mathfrak q) = \text{trdeg}_k\ \kappa(\mathfrak q)$，
对 $\mathfrak q'$ 也同样成立。因此，严格包含
$V(\mathfrak q') \subset V(\mathfrak q)$ 蕴含维数的严格不等式，
从而得到结论。
\end{proof}

\noindent
下面的引理推广了引理
\ref{algebra-lemma-dimension-closed-point-finite-type-field}。

\begin{lemma}
\label{algebra-lemma-dimension-at-a-point-finite-type-field}
设 $k$ 是域。设 $S$ 是有限型
% P01-E0553: the source omits the hyphen in “finite type k algebra”.
$k$-代数。令 $X = \Spec(S)$。设 $\mathfrak p \subset S$ 是素理想，
并令 $x \in X$ 为对应的点。则
$$
\dim_x(X) = \dim(S_{\mathfrak p}) + \text{trdeg}_k\ \kappa(\mathfrak p).
$$
\end{lemma}

\begin{proof}
由引理 \ref{algebra-lemma-dimension-prime-polynomial-ring}，可知
$r = \text{trdeg}_k\ \kappa(\mathfrak p)$ 等于
$V(\mathfrak p)$ 的维数。在 $S$ 中任取一条以 $\mathfrak p$ 开始的
极大素理想链
$\mathfrak p \subset \mathfrak p_1 \subset \ldots \subset \mathfrak p_r$。
由引理 \ref{algebra-lemma-dimension-spell-it-out}，它的长度为 $r$。
设 $\mathfrak q_j$（$j \in J$）为 $S$ 中包含在 $\mathfrak p$
内的极小素理想。按规则
$\mathfrak q_j \mapsto \mathfrak q_jS_{\mathfrak p}$，它们与
$S_{\mathfrak p}$ 中的极小素理想成 $1-1$ 对应。
由引理 \ref{algebra-lemma-dimension-at-a-point-finite-type-over-field}，
可知 $\dim_x(X)$ 等于各环 $S/\mathfrak q_j$ 的维数中的最大值。
对每个 $j$，选取极大素理想链
$\mathfrak q_j \subset \mathfrak p'_1 \subset \ldots \subset \mathfrak p'_{s(j)}
= \mathfrak p$。于是
$\dim(S_{\mathfrak p}) = \max_{j \in J} s(j)$。
现在，每条链
$$
\mathfrak q_j \subset \mathfrak p'_1 \subset \ldots \subset
\mathfrak p'_{s(j)} = \mathfrak p \subset
\mathfrak p_1 \subset \ldots \subset \mathfrak p_r
$$
都是 $S/\mathfrak q_j$ 中的极大链，并且由前述讨论有
$\dim_x(X) = \max_{j \in J} r + s(j)$。引理得证。
\end{proof}

\noindent
下面的引理说明，一个有限型谱在另一个有限型谱中的余维等于高度之差。

\begin{lemma}
\label{algebra-lemma-codimension}
设 $k$ 是域。设 $S' \to S$ 是有限型
% P01-E0554: the source omits the hyphen in “finite type k algebras”.
$k$-代数之间的满射。设 $\mathfrak p \subset S$ 是素理想，
并设 $\mathfrak p'$ 是 $S'$ 中对应的素理想。令
$X = \Spec(S)$、相应地 $X' = \Spec(S')$，并令
$x \in X$、相应地 $x'\in X'$ 为分别对应于
$\mathfrak p$、相应地 $\mathfrak p'$ 的点。则
$$
\dim_{x'} X' - \dim_x X =
\text{高度}(\mathfrak p') - \text{高度}(\mathfrak p).
$$
\end{lemma}

\begin{proof}
这由引理 \ref{algebra-lemma-dimension-at-a-point-finite-type-field} 立即得出。
\end{proof}

\begin{lemma}
\label{algebra-lemma-dimension-preserved-field-extension}
设 $k$ 是域。设 $S$ 是有限型 $k$-代数。设 $K/k$ 是域扩张。
则 $\dim(S) = \dim(K \otimes_k S)$。
\end{lemma}

\begin{proof}
由引理 \ref{algebra-lemma-Noether-normalization}，存在有限单射映射
$k[y_1, \ldots, y_d] \to S$，其中 $d = \dim(S)$。
由于 $K$ 在 $k$ 上平坦，还得到有限单射映射
$K[y_1, \ldots, y_d] \to K \otimes_k S$。
结论由引理 \ref{algebra-lemma-integral-sub-dim-equal} 得出。
\end{proof}

\begin{lemma}
\label{algebra-lemma-dimension-at-a-point-preserved-field-extension}
设 $k$ 是域。设 $S$ 是有限型 $k$-代数。令 $X = \Spec(S)$。
设 $K/k$ 是域扩张。令 $S_K = K \otimes_k S$，并令
$X_K = \Spec(S_K)$。设 $\mathfrak q \subset S$ 是对应于
$x \in X$ 的素理想，并设 $\mathfrak q_K \subset S_K$ 是位于
$\mathfrak q$ 上方且对应于 $x_K \in X_K$ 的素理想。则
$\dim_x X = \dim_{x_K} X_K$。
\end{lemma}

\begin{proof}
选取表示 $S = k[x_1, \ldots, x_n]/I$。这给出表示
$K \otimes_k S = K[x_1, \ldots, x_n]/(K \otimes_k I)$。
设 $\mathfrak q_K' \subset K[x_1, \ldots, x_n]$、相应地
$\mathfrak q' \subset k[x_1, \ldots, x_n]$ 为对应的素理想。
考虑下列诺特局部环的交换图
$$
\xymatrix{
K[x_1, \ldots, x_n]_{\mathfrak q_K'} \ar[r] &
(K \otimes_k S)_{\mathfrak q_K} \\
k[x_1, \ldots, x_n]_{\mathfrak q'} \ar[r] \ar[u] &
S_{\mathfrak q} \ar[u]
}
$$
两个竖直箭头都是平坦的，因为它们分别是平坦环映射
$S \to S_K$ 和
$k[x_1, \ldots, x_n] \to K[x_1, \ldots, x_n]$ 的局部化。
此外，两个竖直箭头具有相同的纤维环。因此，由引理
\ref{algebra-lemma-dimension-base-fibre-equals-total} 可知
$\text{高度}(\mathfrak q') - \text{高度}(\mathfrak q)
= \text{高度}(\mathfrak q_K') - \text{高度}(\mathfrak q_K)$。
% P01-E0555: the source omits “by” in the two-object “Denote ... the points” construction.
以 $x' \in X' = \Spec(k[x_1, \ldots, x_n])$
和 $x'_K \in X'_K = \Spec(K[x_1, \ldots, x_n])$
表示分别对应于 $\mathfrak q'$ 和 $\mathfrak q_K'$ 的点。
由引理 \ref{algebra-lemma-codimension} 和上面证明的结论，有
\begin{eqnarray*}
n - \dim_x X & = & \dim_{x'} X' - \dim_x X \\
& = & \text{高度}(\mathfrak q') - \text{高度}(\mathfrak q) \\
& = & \text{高度}(\mathfrak q_K') - \text{高度}(\mathfrak q_K) \\
& = & \dim_{x'_K} X'_K - \dim_{x_K} X_K \\
& = & n - \dim_{x_K} X_K
\end{eqnarray*}
，引理随即得证。
\end{proof}

\begin{lemma}
\label{algebra-lemma-inequalities-under-field-extension}
设 $k$ 是域。设 $S$ 是有限型 $k$-代数。设 $K/k$ 是域扩张。
令 $S_K = K \otimes_k S$。设 $\mathfrak q \subset S$ 是素理想，
并设 $\mathfrak q_K \subset S_K$ 是位于 $\mathfrak q$ 上方的素理想。
则
$$
\dim (S_K \otimes_S \kappa(\mathfrak q))_{\mathfrak q_K} =
\dim (S_K)_{\mathfrak q_K} - \dim S_\mathfrak q =
\text{trdeg}_k \kappa(\mathfrak q) - \text{trdeg}_K \kappa(\mathfrak q_K)
$$
此外，给定 $\mathfrak q$，总可以选取 $\mathfrak q_K$，使得上面的数为零。
\end{lemma}

\begin{proof}
注意，$S_\mathfrak q \to (S_K)_{\mathfrak q_K}$ 是诺特局部环的平坦局部同态，
其特殊纤维为
$(S_K \otimes_S \kappa(\mathfrak q))_{\mathfrak q_K}$。因此，
% P01-E0556: the source sentence “Hence the first equality by Lemma ...” is missing its finite verb.
由引理 \ref{algebra-lemma-dimension-base-fibre-equals-total} 得到第一个等式。
第二个等式来自以下事实：采用引理
\ref{algebra-lemma-dimension-at-a-point-preserved-field-extension} 中的记号，有
$\dim_x X = \dim_{x_K} X_K$；并且由引理
\ref{algebra-lemma-dimension-at-a-point-finite-type-field}，有
$\dim_x X = \dim S_\mathfrak q + \text{trdeg}_k \kappa(\mathfrak q)$，
对 $\dim_{x_K} X_K$ 也同样成立。
若选取 $\mathfrak q_K$ 为 $\mathfrak q S_K$ 上的极小素理想，
则纤维环的维数为零。
\end{proof}

\section{域上分次代数的维数}
\label{algebra-section-dimension-graded}

\noindent
下面给出一个基本结果。

\begin{lemma}
\label{algebra-lemma-dimension-graded}
设 $k$ 是域。设 $S$ 是分次 $k$-代数，并且在 $k$ 上由有限多个
次数为 $1$ 的元素生成。假设 $S_0 = k$。设
$P(T) \in \mathbf{Q}[T]$ 是满足如下性质的多项式：
对所有 $d \gg 0$，都有 $\dim(S_d) = P(d)$。见命题
\ref{algebra-proposition-graded-hilbert-polynomial}。则
\begin{enumerate}
\item 无关理想 $S_{+}$ 是一个极大理想 $\mathfrak m$。
\item $S$ 的任意极小素理想都是齐次理想，并包含于
$S_{+} = \mathfrak m$。
\item 有 $\dim(S) = \deg(P) + 1 = \dim_x\Spec(S)$
（约定 $\deg(0) = -1$），其中 $x$ 是对应于极大理想
$S_{+} = \mathfrak m$ 的点。
\item 局部环 $R = S_{\mathfrak m}$ 的 Hilbert 函数等于 $S$
的 Hilbert 函数。
\end{enumerate}
\end{lemma}

\begin{proof}
第一个陈述显然。第二个陈述由引理
\ref{algebra-lemma-graded-ring-minimal-prime} 得出。由 (2)，每个不可约分支
都经过 $x$。因此，由引理
\ref{algebra-lemma-dimension-at-a-point-finite-type-over-field}，有
$\dim(S) = \dim_x\Spec(S) = \dim(S_\mathfrak m)$。由于
$\mathfrak m^d/\mathfrak m^{d + 1} \cong
\mathfrak m^dS_\mathfrak m/\mathfrak m^{d + 1}S_\mathfrak m$，
可知局部环 $S_\mathfrak m$ 的 Hilbert 函数等于 $S$ 的 Hilbert 函数，
这就是 (4)。最后，由命题 \ref{algebra-proposition-dimension}
得到 (3) 中的最后一个等式。
\end{proof}

\section{一般平坦性}
\label{algebra-section-generic-flatness}

\noindent
这个结果大致是说，整环上的有限型代数在把该整环的一个元素可逆化后会变得平坦。
这个结果有若干版本，强度依次递增。

\begin{lemma}
\label{algebra-lemma-generic-flatness-Noetherian}
设 $R \to S$ 是环映射。
设 $M$ 是 $S$-模。假设
\begin{enumerate}
\item $R$ 是诺特环，
\item $R$ 是整环，
\item $R \to S$ 是有限型的，并且
\item $M$ 是有限型 $S$-模。
\end{enumerate}
则存在非零元 $f \in R$，使得
$M_f$ 是自由 $R_f$-模。
\end{lemma}

\begin{proof}
设 $K$ 为 $R$ 的分式域。置 $S_K = K \otimes_R S$。这是 $K$ 上的
有限型代数。我们对 $d = \dim(S_K)$ 作归纳（这是有限的，例如可由
Noether 正规化得到，见第 \ref{algebra-section-Noether-normalization} 节）。
固定 $d \geq 0$。假设在所有满足 $\dim(S_K) < d$ 的情形中引理均已成立。

\medskip\noindent
设给定引理中的 $R \to S$ 和 $M$，且 $\dim(S_K) = d$。由引理
\ref{algebra-lemma-filter-Noetherian-module}，存在滤过
$0 \subset M_1 \subset M_2 \subset \ldots \subset M_n = M$，
使得 $M_i/M_{i - 1}$ 同构于 $S/\mathfrak q$，其中 $\mathfrak q$
是 $S$ 的某个素理想。注意
$\dim((S/\mathfrak q)_K) \leq \dim(S_K)$。还要注意，自由模的扩张仍是自由模
（见基本概念 \ref{algebra-item-extension-free}）。因此可以假设 $M = S$，并且
$S$ 是 $R$ 上的有限型整环。

\medskip\noindent
若 $R \to S$ 有非平凡的核，则在该核中取非零元 $f \in R$。
此时 $S_f = 0$，引理成立。（这实际上是 $d = -1$ 的情形，也是归纳的起点。）
所以可以假设 $R \to S$ 是诺特整环之间的有限型扩张。

\medskip\noindent
应用引理 \ref{algebra-lemma-Noether-normalization-over-a-domain}，并把 $R$
替换为 $R_f$（其中 $f$ 如该引理所述），得到分解
$$
R \subset R[y_1, \ldots, y_d] \subset S
$$
其中第二个扩张是有限的。选取 $z_1, \ldots, z_r \in S$，使其构成
$S$ 的分式域在 $R[y_1, \ldots, y_d]$ 的分式域上的一组基。这给出短正合列
$$
0 \to
R[y_1, \ldots, y_d]^{\oplus r} \xrightarrow{(z_1, \ldots, z_r)}
S \to N \to 0
$$
由构造，$N$ 是有限 $R[y_1, \ldots, y_d]$-模，其支集不包含
$\Spec(R[y_1, \ldots, y_d])$ 的泛点 $(0)$。由引理
\ref{algebra-lemma-support-closed}，存在非零元 $g \in R[y_1, \ldots, y_d]$
消去 $N$，所以可以把 $N$ 视为
$S' = R[y_1, \ldots, y_d]/(g)$ 上的有限模。因为 $\dim(S'_K) < d$，
归纳假设给出非零元 $f \in R$，使得 $N_f$ 是自由 $R_f$-模。又因为
% P01-E0557: the source clause “is free also we conclude” is grammatically malformed.
$(R[y_1, \ldots, y_d])_f \cong R_f[y_1, \ldots, y_d]$ 是自由的，
由前述自由模的扩张仍自由这一事实可知结论成立。
\end{proof}

\begin{lemma}
\label{algebra-lemma-generic-flatness-finitely-presented}
\begin{slogan}
一般自由性。
\end{slogan}
设 $R \to S$ 是环映射。
设 $M$ 是 $S$-模。假设
\begin{enumerate}
\item $R$ 是整环，
\item $R \to S$ 是有限表示的，并且
\item $M$ 是有限表示 $S$-模。
\end{enumerate}
则存在非零元 $f \in R$，使得
$M_f$ 是自由 $R_f$-模。
\end{lemma}

\begin{proof}
写成 $S = R[x_1, \ldots, x_n]/(g_1, \ldots, g_m)$。
% P01-E0558: the source designation “denote ... its image” omits “by”.
对 $g \in R[x_1, \ldots, x_n]$，以 $\overline{g}$ 表示它在 $S$ 中的像。
可以写成 $M = S^{\oplus t}/\sum Sn_i$，其中
$n_i \in S^{\oplus t}$。把它们写成
$n_i = (\overline{g}_{i1}, \ldots, \overline{g}_{it})$，其中
$g_{ij} \in R[x_1, \ldots, x_n]$。设 $R_0 \subset R$ 是由所有元素
$g_i, g_{ij} \in R[x_1, \ldots, x_n]$ 的全部系数生成的子环。定义
$S_0 = R_0[x_1, \ldots, x_n]/(g_1, \ldots, g_m)$，并定义
$M_0 = S_0^{\oplus t}/\sum S_0n_i$。于是 $R_0$ 是
$\mathbf{Z}$ 上的有限型整环，因而是诺特环（见引理
\ref{algebra-lemma-Noetherian-permanence}）。此外，通过单射 $R_0 \to R$，有
$S \cong R \otimes_{R_0} S_0$ 和 $M \cong R \otimes_{R_0} M_0$。
应用引理 \ref{algebra-lemma-generic-flatness-Noetherian}，得到非零元
$f \in R_0$，使得 $(M_0)_f$ 是自由 $(R_0)_f$-模。因此
$M_f = R_f \otimes_{(R_0)_f} (M_0)_f$ 是自由 $R_f$-模。
\end{proof}

\begin{lemma}
\label{algebra-lemma-generic-flatness}
设 $R \to S$ 是环映射。
设 $M$ 是 $S$-模。假设
\begin{enumerate}
\item $R$ 是整环，
\item $R \to S$ 是有限型的，并且
\item $M$ 是有限型 $S$-模。
\end{enumerate}
则存在非零元 $f \in R$，使得
\begin{enumerate}
\item[(a)] $M_f$ 和 $S_f$ 作为 $R_f$-模都是自由的，并且
\item[(b)] $S_f$ 是有限表示 $R_f$-代数，而 $M_f$ 是有限表示
$S_f$-模。
\end{enumerate}
\end{lemma}

\begin{proof}
我们先对 $S = R[x_1, \ldots, x_n]$ 证明引理，再由此推出一般情形。

\medskip\noindent
假设 $S = R[x_1, \ldots, x_n]$。
选取生成 $M$ 的元素 $m_1, \ldots, m_t$。这给出短正合列
$$
0 \to N \to S^{\oplus t} \xrightarrow{(m_1, \ldots, m_t)} M \to 0.
$$
% P01-E0559: the source twice uses “Denote” without the required preposition “by”.
以 $K$ 表示 $R$ 的分式域。记
$S_K = K \otimes_R S = K[x_1, \ldots, x_n]$，并类似地记
$N_K = K \otimes_R N$、$M_K = K \otimes_R M$。
由于 $R \to K$ 平坦，该序列在与 $K$ 张量后仍正合。由于
$S_K = K[x_1, \ldots, x_n]$ 是诺特环（见引理
\ref{algebra-lemma-Noetherian-permanence}），可以找到有限多个生成它的元素
$n'_1, \ldots, n'_s \in N_K$。选取 $n_1, \ldots, n_r \in N$，
使得对某些 $a_{ij} \in K$ 有 $n'_i = \sum a_{ij}n_j$。置
$$
M' = S^{\oplus t}/\sum\nolimits_{i = 1, \ldots, r} Sn_i
$$
由构造，$M'$ 是有限表示 $S$-模，并且存在满射 $M' \to M$，
它诱导同构 $M'_K \cong M_K$。可以对 $R \to S$ 和 $M'$
应用引理 \ref{algebra-lemma-generic-flatness-finitely-presented}，得到
$f \in R$，使得 $M'_f$ 是自由 $R_f$-模。因此 $M'_f \to M_f$
是整环 $R_f$ 上的模满射，其源是自由模，并且与 $K$ 张量后成为同构。
由于 $M'_f$ 自由，并且 $R_f$ 是分式域为 $K$ 的整环，故
$M'_f \subset M'_K$，从而该映射是单射。因此 $M'_f \to M_f$
是同构，结论得证。

\medskip\noindent
在一般情形中，选取满射 $R[x_1, \ldots, x_n] \to S$。
把 $S$ 和 $M$ 都看作有限 $R[x_1, \ldots, x_n]$-模。
由上面证明的特殊情形，存在非零元 $f \in R$，使得 $S_f$ 和 $M_f$
作为 $R_f$-模都是自由的，并且作为 $R_f[x_1, \ldots, x_n]$-模
都是有限表示的。这显然蕴含 $S_f$ 是有限表示 $R_f$-代数，且
$M_f$ 是有限表示 $S_f$-模。
\end{proof}

\noindent
设 $R \to S$ 是环映射。设 $M$ 是 $S$-模。考虑元素 $f \in R$
所满足的下列条件：
\begin{equation}
\label{algebra-equation-flat-and-finitely-presented}
\left\{
\begin{matrix}
S_f & \text{is of finite presentation over }R_f\\
M_f & \text{is of finite presentation as }S_f\text{-模}\\
S_f, M_f & \text{are free as }R_f\text{-模}
\end{matrix}
\right.
\end{equation}
我们定义
\begin{equation}
\label{algebra-equation-good-locus}
U(R \to S, M)
=
\bigcup\nolimits_{f \in
R\text{ 其中 }(\ref{algebra-equation-flat-and-finitely-presented})}
D(f)
\end{equation}
这是 $\Spec(R)$ 的一个开子集。

\begin{lemma}
\label{algebra-lemma-generic-flatness-locus-extension}
设 $R \to S$ 是环映射。
设 $0 \to M_1 \to M_2 \to M_3 \to 0$ 是 $S$-模的短正合列。
则
$$
U(R \to S, M_1) \cap U(R \to S, M_3) \subset U(R \to S, M_2).
$$
\end{lemma}

\begin{proof}
设 $u \in U(R \to S, M_1) \cap U(R \to S, M_3)$。选取
$f_1, f_3 \in R$，使得 $u \in D(f_1)$、$u \in D(f_3)$，并且
(\ref{algebra-equation-flat-and-finitely-presented}) 对 $f_1$ 和 $M_1$
以及对 $f_3$ 和 $M_3$ 都成立。令 $f = f_1f_3$。于是
$u \in D(f)$，并且 (\ref{algebra-equation-flat-and-finitely-presented})
对 $f$ 以及 $M_1$、$M_3$ 都成立。自由模的扩张仍是自由模，
而有限表示模的扩张仍是有限表示模（引理 \ref{algebra-lemma-extension}）。
因此 (\ref{algebra-equation-flat-and-finitely-presented}) 对 $f$ 和 $M_2$
成立。故 $u \in U(R \to S, M_2)$，结论得证。
\end{proof}

\begin{lemma}
\label{algebra-lemma-generic-flatness-locus-localize}
设 $R \to S$ 是环映射。
设 $M$ 是 $S$-模。设 $f \in R$。
在等同 $\Spec(R_f) = D(f)$ 之下，有
$U(R_f \to S_f, M_f) = D(f) \cap U(R \to S, M)$。
\end{lemma}

\begin{proof}
假设 $u \in U(R_f \to S_f, M_f)$。则存在元素 $g \in R_f$，
使得 $u \in D(g)$，并且 (\ref{algebra-equation-flat-and-finitely-presented})
对二元组 $((R_f)_g \to (S_f)_g, (M_f)_g)$ 成立。把
$g$ 写成 $g = a/f^n$，其中 $a \in R$。令 $h = af$。则
$R_h = (R_f)_g$、$S_h = (S_f)_g$ 且 $M_h = (M_f)_g$。
此外 $u \in D(h)$。故 $u \in U(R \to S, M)$。
反过来，假设 $u \in D(f) \cap U(R \to S, M)$。则存在元素
$g \in R$，使得 $u \in D(g)$，并且
(\ref{algebra-equation-flat-and-finitely-presented}) 对二元组
$(R_g \to S_g, M_g)$ 成立。于是显然，
(\ref{algebra-equation-flat-and-finitely-presented}) 对二元组
$(R_{fg} \to S_{fg}, M_{fg}) = ((R_f)_g \to (S_f)_g, (M_f)_g)$
也成立。故 $u \in U(R_f \to S_f, M_f)$，结论得证。
\end{proof}

\begin{lemma}
\label{algebra-lemma-generic-flatness-locus-reduce}
设 $R \to S$ 是环映射。
设 $M$ 是 $S$-模。设 $U \subset \Spec(R)$ 是稠密开集。
假设存在由开集组成的覆盖 $U = \bigcup_{i \in I} D(f_i)$，
使得对每个 $i \in I$，$U(R_{f_i} \to S_{f_i}, M_{f_i})$
在 $D(f_i)$ 中稠密。则 $U(R \to S, M)$ 在 $\Spec(R)$ 中稠密。
\end{lemma}

\begin{proof}
由引理 \ref{algebra-lemma-generic-flatness-locus-localize}，这纯粹是一个拓扑陈述。
事实上，由该引理可知，对每个 $i \in I$，
$U(R \to S, M) \cap D(f_i)$ 在 $D(f_i)$ 中稠密。
由拓扑学，引理 \ref{topology-lemma-nowhere-dense-local}，可知
$U(R \to S, M) \cap U$ 在 $U$ 中稠密。又因为 $U$ 在
$\Spec(R)$ 中稠密，所以 $U(R \to S, M)$ 在 $\Spec(R)$ 中稠密。
\end{proof}

\begin{lemma}
\label{algebra-lemma-generic-flatness-reduced}
设 $R \to S$ 是环映射。设 $M$ 是 $S$-模。假设
\begin{enumerate}
\item $R \to S$ 是有限型的，
\item $M$ 是有限 $S$-模，并且
\item $R$ 是约化环。
\end{enumerate}
则存在子集 $U \subset \Spec(R)$，使得
\begin{enumerate}
\item $U$ 是 $\Spec(R)$ 中的稠密开集；
\item 对每个 $u \in U$，存在 $f \in R$，使得
$u \in D(f) \subset U$，并且
\begin{enumerate}
\item $M_f$ 和 $S_f$ 在 $R_f$ 上自由，
\item $S_f$ 是有限表示 $R_f$-代数，并且
\item $M_f$ 是有限表示 $S_f$-模。
\end{enumerate}
\end{enumerate}
\end{lemma}

\begin{proof}
注意，该引理等价于开集 $U(R \to S, M)$ 在 $\Spec(R)$ 中稠密；
见方程 (\ref{algebra-equation-good-locus})。我们先对
$S = R[x_1, \ldots, x_n]$ 证明引理，再由此推出一般情形。

\medskip\noindent
先证明 $S = R[x_1, \ldots, x_n]$ 且 $M$ 是 $S$ 上任意有限模的情形。
注意，此时 $S_f = R_f[x_1, \ldots, x_n]$ 在 $R_f$ 上自由且有限表示，
所以无须考虑关于 $S$ 的条件，只须考虑关于 $M$ 的条件。
我们对 $n$ 作归纳。

\medskip\noindent
存在有限滤过
$$
0 \subset M_1 \subset M_2 \subset \ldots \subset M_t = M
$$
使得 $M_i/M_{i - 1} \cong S/J_i$，其中 $J_i \subset S$ 是某个理想，见引理
\ref{algebra-lemma-trivial-filter-finite-module}。由于有限多个稠密开集的交仍是稠密开集，
由引理 \ref{algebra-lemma-generic-flatness-locus-extension} 可知，只需对每个模
% P01-E0560: the frozen source says R/J_i although J_i is an ideal of S and the filtration factors are S/J_i.
$R/J_i$ 证明引理。因此可以假设 $M = S/J$，其中 $J$ 是
$S = R[x_1, \ldots, x_n]$ 的一个理想。

\medskip\noindent
设 $I \subset R$ 是由 $J$ 中元素的系数生成的理想。令
$U_1 = \Spec(R) \setminus V(I)$，并令
$$
U_2 = \Spec(R) \setminus \overline{U_1}.
$$
显然 $U = U_1 \cup U_2$ 在 $\Spec(R)$ 中稠密。设 $f \in R$
满足 (a) $D(f) \subset U_1$ 或 (b) $D(f) \subset U_2$。
若对每个这样的 $f$，引理都对二元组
$(R_f \to R_f[x_1, \ldots, x_n], M_f)$ 成立，则由引理
\ref{algebra-lemma-generic-flatness-locus-reduce} 可知
$U(R \to S, M)$ 在 $\Spec(R)$ 中稠密。因此可以假设
(a) $I = R$，或 (b) $V(I) = \Spec(R)$。

\medskip\noindent
在情形 (b) 中，由于 $R$ 约化，实际上有 $I = 0$！因此
$J = 0$ 且 $M = S$，引理在此情形中成立。

\medskip\noindent
在情形 (a) 中还需多做一点工作。注意，$I$ 的每个元素实际上都是
$J$ 中某个元素的某个单项式的系数，因为 $J$ 中元素的系数之集构成一个理想
（细节略）。所以可以找到元素
$$
g = \sum\nolimits_{K \in E} a_K x^K \in J
$$
其中 $E$ 是多重指标 $K = (k_1, \ldots, k_n)$ 的有限集合，且至少有一个系数
$a_{K_0}$ 是 $R$ 中的单位。事实上，由于在情形 (a) 中 $1 \in I$，
可以找到一个系数等于 $1$ 的这种元素。令
$m = \#\{K \in E \mid a_K \text{ is not a unit}\}$。
注意 $0 \leq m \leq \# E - 1$。我们对 $m$ 作归纳。

\medskip\noindent
情形 $m = 0$。此时 $g$ 的所有系数 $a_K$（$K \in E$）都是单位，且
$E \not = \emptyset$。若 $E = \{K_0\}$ 是单元素集，并且
$K_0 = (0, \ldots, 0)$，则 $g$ 是单位且 $J = S$，所以结论当然成立。
（特别地，当 $n = 0$ 时会出现这种情况，它给出了关于 $n$ 的归纳起点。）
若并非 $E = \{(0, \ldots, 0)\}$，则至少有一个 $K$ 不等于
$(0, \ldots, 0)$，亦即 $g \not \in R$。此时采用 Noether 正规化的惯用技巧。
具体地，考虑
$$
G(y_1, \ldots, y_n) =
g(y_1 + y_n^{e_1}, y_2 + y_n^{e_2}, \ldots, y_{n - 1} + y_n^{e_{n - 1}}, y_n)
$$
其中 $0 \ll e_{n -1} \ll e_{n - 2} \ll \ldots \ll e_1$。由引理
\ref{algebra-lemma-helper-polynomial}，把 $G(y_1, \ldots, y_n)$ 视为
$y_n$ 的多项式时，它形如
$$
a_K y_n^{k_n + \sum_{i = 1, \ldots, n - 1} e_i k_i} +
\text{lower order terms in }y_n
$$
由于 $a_K$ 是单位，可知 $M = R[x_1, \ldots, x_n]/J$ 是
$R[y_1, \ldots, y_{n - 1}]$ 上的有限模。因此
$U(R \to R[x_1, \ldots, x_n], M) = U(R \to R[y_1, \ldots, y_{n - 1}], M)$，
由对 $n$ 的归纳即得结论。

\medskip\noindent
情形 $m > 0$。选取多重指标 $K \in E$，使得 $a_K$ 不是单位。
同前，令
$U_1 = \Spec(R_{a_K}) = \Spec(R) \setminus V(a_K)$，
并令
$$
U_2 = \Spec(R) \setminus \overline{U_1}.
$$
显然 $U = U_1 \cup U_2$ 在 $\Spec(R)$ 中稠密。设 $f \in R$
满足 (a) $D(f) \subset U_1$ 或 (b) $D(f) \subset U_2$。
若对每个这样的 $f$，引理都对二元组
$(R_f \to R_f[x_1, \ldots, x_n], M_f)$ 成立，则由引理
\ref{algebra-lemma-generic-flatness-locus-reduce} 可知
$U(R \to S, M)$ 在 $\Spec(R)$ 中稠密。因此可以假设
(a) $a_KR = R$，或 (b) $V(a_K) = \Spec(R)$。
在情形 (a) 中，$a_K$ 已成为单位，所以数 $m$ 减小。
在情形 (b) 中，由于 $R$ 约化，可知 $a_K = 0$。
于是集合 $E$ 缩小，数 $m$ 也随之减小。在两种情形中都由对 $m$ 的归纳得到结论。

\medskip\noindent
至此已经证明了 $S = R[x_1, \ldots, x_n]$ 的情形。假设
$(R \to S, M)$ 是满足引理条件的任意二元组。选取满射
$R[x_1, \ldots, x_n] \to S$。采用 (\ref{algebra-equation-good-locus})
中引入的记号，有
$$
U(R \to S, M) =
U(R \to R[x_1, \ldots, x_n], S)
\cap
U(R \to R[x_1, \ldots, x_n], M)
$$
% P01-E0561: the source final sentence is grammatically malformed.
右侧两个开集由刚证明的结果都是稠密的，故左侧的开集也是稠密的。
\end{proof}

\section{Krull--Akizuki 定理及相关结果}
\label{algebra-section-krull-akizuki}

\noindent
Krull--Akizuki 定理的一个应用是说明离散赋值环非常丰富。
本节还将更一般地说明如何构造支配诺特局部环的离散赋值环。

\medskip\noindent
首先说明如何通过对极大理想作爆破，用一个 $1$ 维诺特局部整环
支配给定的诺特局部整环。

\begin{lemma}
\label{algebra-lemma-dominate-by-dimension-1}
设 $R$ 是分式域为 $K$ 的诺特局部整环。假设 $R$ 不是域。
则存在 $R \subset R' \subset K$，满足
\begin{enumerate}
\item $R'$ 是维数为 $1$ 的诺特局部环；
\item $R \to R'$ 是局部环映射，亦即 $R'$ 支配 $R$；并且
\item $R \to R'$ 本质有限型。
\end{enumerate}
\end{lemma}

\begin{proof}
选取支配 $R$ 的任意赋值环 $A \subset K$（其存在性由引理
\ref{algebra-lemma-dominate} 保证）。
% P01-E0562: the source designation “Denote v ...” omits the required preposition “by”.
以 $v$ 表示相应的赋值。设 $x_1, \ldots, x_r$ 是 $R$ 的极大理想
$\mathfrak m$ 的一组极小生成元。可以并且确实假设
$v(x_r) = \min\{v(x_1), \ldots, v(x_r)\}$。考虑环
$$
S = R[x_1/x_r, x_2/x_r, \ldots, x_{r - 1}/x_r] \subset K.
$$
注意 $\mathfrak mS = x_rS$ 是主理想。又有 $S \subset A$ 且
$v(x_r) > 0$，因此 $x_rS \not = S$。选取 $x_rS$ 上方的一个极小素理想
$\mathfrak q$。由引理 \ref{algebra-lemma-minimal-over-1}，
$\text{高度}(\mathfrak q) = 1$，并且 $\mathfrak q$ 位于
$\mathfrak m$ 上方。因此 $R' = S_{\mathfrak q}$ 即为所求。
\end{proof}

\begin{lemma}[Koll\'ar]
\label{algebra-lemma-hart-serre-loc-thm}
\begin{reference}
这取自 J\'anos Koll\'ar 即将发表的论文
“Variants of normality for Noetherian schemes”。
\end{reference}
设 $(R, \mathfrak m)$ 是诺特局部环。则下列情形恰有一个成立：
\begin{enumerate}
\item $(R, \mathfrak m)$ 是阿廷环；
\item $(R, \mathfrak m)$ 是维数为 $1$ 的正则环；
\item $\text{depth}(R) \geq 2$；或
\item 存在有限环映射 $R \to R'$，它不是同构，其核与余核都被
$\mathfrak m$ 的某个幂消去，并且 $\mathfrak m$ 不是 $R'$ 的相关素理想，
且 $R' \not = 0$。
\end{enumerate}
\end{lemma}

\begin{proof}
注意，$(R, \mathfrak m)$ 不是阿廷环，当且仅当
$V(\mathfrak m) \subset \Spec(R)$ 无处稠密。见命题
\ref{algebra-proposition-dimension-zero-ring}。以下始终作此假设。

\medskip\noindent
设 $J \subset R$ 是被 $\mathfrak m$ 的某个幂消去的最大理想。
若 $J \not = 0$，则 $R \to R/J$ 表明 $(R, \mathfrak m)$ 属于情形 (4)。

\medskip\noindent
否则 $J = 0$。特别地，$\mathfrak m$ 不是 $R$ 的相关素理想，
由引理 \ref{algebra-lemma-ideal-nonzerodivisor}，存在非零因子
$x \in \mathfrak m$。若 $\mathfrak m$ 不是 $R/xR$ 的相关素理想，
则由同一引理有 $\text{depth}(R) \geq 2$。所以只剩下如下情形：
存在 $y \in R$、$y \not \in xR$，使得 $y \mathfrak m \subset xR$。

\medskip\noindent
若 $y \mathfrak m \subset x \mathfrak m$，则可以考虑映射
$\varphi : \mathfrak m \to \mathfrak m$、$f \mapsto yf/x$
（由于 $x$ 是非零因子，该映射定义良好）。由引理
\ref{algebra-lemma-charpoly-module} 的行列式技巧，存在系数在 $R$ 中的首一多项式
$P$，使得 $P(\varphi) = 0$。于是 $P(y/x) = 0$ 在 $R_x$ 中成立。
设 $R' \subset R_x$ 是由 $R$ 和 $y/x$ 生成的环。则
$R \subset R'$，并且 $R'/R$ 是被 $\mathfrak m$ 的某个幂消去的有限
$R$-模。因此 $R$ 属于情形 (4)。

\medskip\noindent
否则，存在 $t \in \mathfrak m$，使得对 $R$ 的某个单位 $u$ 有
$y t = u x$。把 $t$ 替换为 $u^{-1}t$ 后，得到 $yt = x$。
特别地，$y$ 是非零因子。对任意 $t' \in \mathfrak m$，
都有 $y t' = x s$，其中 $s \in R$。所以
$y (t' - s t ) = x s - x s = 0$。由于 $y$ 不是零因子，
这蕴含 $t' = ts$，因而 $\mathfrak m = (t)$。所以
$(R, \mathfrak m)$ 是维数为 1 的正则环。

\medskip\noindent
到目前为止，论证表明每个 $R$ 都属于上述四种情形之一。
为完成证明，必须说明任意两个性质不能同时成立，也就是检验六种性质组合。
这里只说明 (2) 和 (3) 都排除 (4)；其余情形留给读者。

\medskip\noindent
假设 $R$ 是维数为 $1$ 的正则环，并且 $R \to R'$ 是有限环映射，
其核与余核被 $\mathfrak m$ 的某个幂消去，而 $\mathfrak m$ 不是
$R'$ 的相关素理想。于是 $R'$ 是深度至少为 $1$ 的有限 $R$-模，
因而由引理 \ref{algebra-lemma-regular-mcm-free} 可知它是自由的。
% P01-E0563: the source says “in the generic point” instead of “at the generic point”.
由于 $R \to R'$ 在泛点处是同构，$R'$ 作为 $R$-模必为秩 $1$ 的自由模。
于是 $R' \to \text{End}_R(R') \cong R$ 是映射 $R \to R'$ 的逆。
故 (4) 不可能成立。

\medskip\noindent
假设 $R$ 的深度 $\geq 2$，并且 $R \to R'$ 是有限环映射，
其核与余核被 $\mathfrak m$ 的某个幂消去，而 $\mathfrak m$ 不是
$R'$ 的相关素理想。则 $R \to R'$ 必为单射（否则其核的深度会同时为
$0$ 和 $\geq 1$），并且其余核的深度至少为 $1$（由引理
\ref{algebra-lemma-depth-in-ses} 以及假设所给的 $R'$ 深度 $\geq 1$）。
因此，除非该余核也是 $0$，否则它不可能支撑在
$\{\mathfrak m\}$ 上。故 (4) 不可能成立。
\end{proof}

\begin{lemma}
\label{algebra-lemma-nonregular-dimension-one}
设 $R$ 是极大理想为 $\mathfrak m$ 的局部环。假设 $R$ 是诺特环、
维数为 $1$，并且 $\dim(\mathfrak m/\mathfrak m^2) > 1$。
则存在环映射 $R \to R'$，使得
\begin{enumerate}
\item $R \to R'$ 是有限的；
\item $R \to R'$ 不是同构；
\item $R \to R'$ 的核与余核被 $\mathfrak m$ 的某个幂消去；并且
\item $\mathfrak m$ 不是 $R'$ 的相关素理想。
\end{enumerate}
\end{lemma}

\begin{proof}
这由引理 \ref{algebra-lemma-hart-serre-loc-thm} 以及下列事实得出：
$R$ 不是阿廷环，不正则，并且深度不 $\geq 2$（最后一点是因为，
由引理 \ref{algebra-lemma-bound-depth}，深度不超过维数）。
\end{proof}

\begin{example}
\label{algebra-example-nonreduced}
考虑诺特局部环
$$
R = k[[x, y]]/(y^2)
$$
它的维数为 1，并且是 Cohen--Macaulay 环。
引理 \ref{algebra-lemma-nonregular-dimension-one} 中扩张的一个例子是
$$
k[[x, y]]/(y^2) \subset k[[x, z]]/(z^2), \ \ y \mapsto xz
$$
换言之，这是通过向 $R$ 添加 $y/x$ 得到的扩张。把该构造重复
$n > 1$ 次，其效果是添加元素 $y/x^n$。
\end{example}

\begin{example}
\label{algebra-example-bad-dvr-char-p}
设 $k$ 是特征为 $p > 0$ 的域，并且 $k$ 在其 $p$ 次幂子域
$k^p$ 上的次数为无穷。例如，$k = \mathbf{F}_p(t_1, t_2, t_3, \ldots)$。
考虑环
$$
A =
\left\{
\sum a_i x^i \in k[[x]] \text{ 使得 }
[k^p(a_0, a_1, a_2, \ldots) : k^p] < \infty
\right\}
$$
则 $A$ 是离散赋值环，其完备化为 $A^\wedge = k[[x]]$。
注意，$A \subset k[[x]]$ 所诱导的分式域扩张是无限纯不可分的。
任取 $f \in k[[x]]$、$f \not \in A$。令 $R = A[f] \subset k[[x]]$。
则 $R$ 是维数为 $1$ 的诺特局部整环，而其完备化 $R^\wedge$
非约化（请思考！）。
\end{example}

\begin{remark}
\label{algebra-remark-resolution-dim-1}
假设 $R$ 是 $1$ 维半局部诺特整环。若存在极大理想
$\mathfrak m \subset R$，使得 $R_{\mathfrak m}$ 不正则，
则可以对 $(R, \mathfrak m)$ 应用引理
\ref{algebra-lemma-nonregular-dimension-one}，得到有限环扩张
$R \subset R_1$。（例如，可以这样进行构造，使得
$\Spec(R_1) \to \Spec(R)$ 是 $\Spec(R)$ 沿理想 $\mathfrak m$
的爆破。）当然，$R_1$ 是 $1$ 维半局部诺特整环，并且与 $R$
有相同的分式域。若 $R_1$ 不是正则半局部环，则可以重复该构造，
得到 $R_1 \subset R_2$。这样得到有限环扩张序列
$$
R \subset R_1 \subset R_2 \subset R_3 \subset \ldots
$$
若对某个 $n$，$R_n$ 正则，则该序列可以停止。奇点消解所断言的是，
$R_n$ 最终确实正则。事实上，情况并非如此。具体而言，存在特征为
$0$、维数为 $1$ 的诺特局部整环 $A$，其完备化非约化；见
\cite[Proposition 3.1]{Ferrand-Raynaud}，或参见我们的例子，第
\ref{examples-section-local-completion-nonreduced} 节。
特征 $p > 0$ 的例子见例 \ref{algebra-example-bad-dvr-char-p}。
由于爆破构造与完备化交换，很容易看出该序列永不稳定。
相关讨论见 \cite{Bennett}（主要针对正特征情形）。另一方面，
如果 $R$ 在其所有极大理想处的完备化都是约化的，则该过程会停止
% P01-E0564: the source contains the unresolved editorial placeholder “insert future reference here”.
（在此插入未来参考文献）。
\end{remark}

\begin{lemma}
\label{algebra-lemma-characterize-dvr}
设 $A$ 是环。下列条件等价。
\begin{enumerate}
\item 环 $A$ 是离散赋值环。
\item 环 $A$ 是诺特赋值环，但不是域。
\item 环 $A$ 是维数为 $1$ 的正则局部环。
\item 环 $A$ 是诺特局部整环，其极大理想
$\mathfrak m$ 由一个非零元生成。
\item 环 $A$ 是维数为 $1$ 的诺特正规局部整环。
\end{enumerate}
在这种情形中，若 $\pi$ 是 $A$ 的极大理想的生成元，则 $A$
的每个非零元都能唯一地写成 $u\pi^n$，其中 $u \in A$ 是单位。
\end{lemma}

\begin{proof}
(1) 与 (2) 的等价性就是引理
\ref{algebra-lemma-valuation-ring-Noetherian-discrete}。此外，在引理
\ref{algebra-lemma-valuation-ring-Noetherian-discrete} 的证明中已经看到：
若 $A$ 是离散赋值环，则 $A$ 是 PID，因而满足 (3)。
注意正则局部环是整环（见引理 \ref{algebra-lemma-regular-domain}）。
利用这一点，(3) 与 (4) 的等价性由维数理论得到，见第
\ref{algebra-section-dimension} 节。

\medskip\noindent
假设 (3) 成立，并设 $\pi$ 是极大理想 $\mathfrak m$ 的生成元。
对所有 $n \geq 0$，都有
$\dim_{A/\mathfrak m} \mathfrak m^n/\mathfrak m^{n + 1} = 1$，
因为该空间由 $\pi^n$ 生成（且不可能为零）。特别地，
$\mathfrak m^n = (\pi^n)$，并且分次环
$\bigoplus \mathfrak m^n/\mathfrak m^{n + 1}$ 同构于多项式环
% P01-E0565: the source omits parentheses around the residue-field quotient in A/m[T].
$A/\mathfrak m[T]$。对 $x \in A \setminus \{0\}$，定义
$v(x) = \max\{n \mid x \in \mathfrak m^n\}$。
换言之，$x = u \pi^{v(x)}$，其中 $u \in A^*$。
由上述说明，对所有 $x, y \in A \setminus \{0\}$，都有
$v(xy) = v(x) + v(y)$。通过规定
$v(a/b) = v(a) - v(b)$，把它扩张到 $A$ 的分式域 $K$
（由上面证明的乘法性可知其定义良好）。于是显然，$A$ 恰好由
$K$ 中赋值 $\geq 0$ 的元素组成。因此，由引理
\ref{algebra-lemma-valuation-valuation-ring} 可知 $A$ 是赋值环。

\medskip\noindent
由引理 \ref{algebra-lemma-valuation-ring-normal}，赋值环是正规整环。
所以等价条件 (1) -- (3) 蕴含 (5)。假设 (5) 成立。
为导出矛盾，假设 $\mathfrak m$ 不能由 $1$ 个元素生成。
由引理 \ref{algebra-lemma-nonregular-dimension-one}，存在有限环映射
$A \to A'$，它在把 $\mathfrak m$ 的任意非零元可逆化后成为同构，
但本身不是同构。特别地，可以把 $A'$ 视为 $A$ 的分式域的子集。
由于 $A \to A'$ 有限，它是整的（见引理
\ref{algebra-lemma-finite-is-integral}）。因为 $A$ 正规，得到
$A = A'$，矛盾。
\end{proof}

\begin{definition}
\label{algebra-definition-uniformizer}
设 $A$ 是离散赋值环。若元素 $\pi \in A$ 生成 $A$ 的极大理想，
则称它为{\it 一致化元}。
\end{definition}

\noindent
由引理 \ref{algebra-lemma-characterize-dvr}，离散赋值环的任意两个一致化元都相伴。

\begin{lemma}
\label{algebra-lemma-finite-length}
设 $R$ 是分式域为 $K$ 的整环。设 $M$ 是 $K^{\oplus r}$
的 $R$-子模。假设 $R$ 是维数为 $1$ 的诺特局部环。
对任意非零元 $x \in R$，有 $\text{length}_R(R/xR) < \infty$，并且
$$
\text{length}_R(M/xM) \leq r \cdot \text{length}_R(R/xR).
$$
\end{lemma}

\begin{proof}
若 $x$ 是单位，则结论成立。所以可以假设 $x \in \mathfrak m$，
其中后者是 $R$ 的极大理想。由于 $x$ 非零且 $R$ 是整环，
有 $\dim(R/xR) = 0$，因而 $R/xR$ 的长度有限。
考虑引理中的 $M \subset K^{\oplus r}$。必要时把
$K^{\oplus r}$ 替换为较小的子空间后，可以假设 $M$ 中的元素
生成 $K^{\oplus r}$ 这个 $K$-向量空间。

\medskip\noindent
先假设 $M$ 是有限 $R$-模。在此情形中，可以清除分母并假设
$M \subset R^{\oplus r}$。因为 $M$ 生成 $K^{\oplus r}$ 这个
% P01-E0566: the source has the number-agreement typo “a vectors space”.
向量空间，所以 $R^{\oplus r}/M$ 的长度有限。特别地，存在整数
$c \geq 0$，使得 $x^cR^{\oplus r} \subset M$。注意
$M \supset xM \supset x^2M \supset \ldots$ 是一个模序列，
其逐次商均同构于 $M/xM$。因此
$$
n \text{length}_R(M/xM) = \text{length}_R(M/x^nM).
$$
对 $M = R^{\oplus r}$ 作同样的论证，得到
$$
n \text{length}_R(R^{\oplus r}/xR^{\oplus r}) =
\text{length}_R(R^{\oplus r}/x^nR^{\oplus r}).
$$
由上面对 $c$ 的选取，$x^nM$ 夹在 $x^n R^{\oplus r}$ 与
$x^{n + c}R^{\oplus r}$ 之间。由此容易得到
$$
r(n + c) \text{length}_R(R/xR)
\geq
n \text{length}_R(M/xM)
\geq
r (n - c) \text{length}_R(R/xR)
$$
因此，在有限情形中，引理的结论事实上取等号。

\medskip\noindent
现在假设 $M$ 不是有限模。假设对某个自然数 $k$，
$M/xM$ 的长度 $\geq k$。则可以找到
% P01-E0567: the source display omits the final subset relation before N_k.
$$
0 \subset N_0 \subset N_1 \subset N_2 \subset \ldots N_k \subset M/xM
$$
使得对 $i = 0, \ldots k - 1$，有 $N_i \not = N_{i + 1}$。
对 $i = 1, \ldots, k$，选取元素 $m_i \in M$，使其模 $xM$
的同余类落在 $N_i$ 中但不落在 $N_{i - 1}$ 中。
考虑有限 $R$-模 $M' = Rm_1 + \ldots + Rm_k \subset M$。
设 $N'_i \subset M'/xM'$ 是 $N_i$ 的逆像。
由 $m_i$ 的选取，显然 $N'_i \not =N'_{i + 1}$。
所以 $\text{length}_R(M'/xM') \geq k$。由有限情形可知
$k \leq r\text{length}_R(R/xR)$，如欲得。
\end{proof}

\noindent
下面是第一个应用。

\begin{lemma}
\label{algebra-lemma-finite-extension-residue-fields-dimension-1}
设 $R \to S$ 是整环同态，并诱导分式域的嵌入
$K \subset L$。若 $R$ 是维数为 $1$ 的诺特局部环，并且
$[L : K] < \infty$，则
\begin{enumerate}
\item $S$ 中位于 $R$ 的极大理想 $\mathfrak m$ 上方的每个素理想
$\mathfrak n_i$ 都是极大理想；
\item 这样的素理想只有有限多个；并且
\item 对每个 $i$，都有
$[\kappa(\mathfrak n_i) : \kappa(\mathfrak m)] < \infty$。
\end{enumerate}
\end{lemma}

\begin{proof}
选取非零元 $x \in \mathfrak m$。对
$S \subset L \cong K^{\oplus n}$ 应用引理
\ref{algebra-lemma-finite-length}，其中 $n = [L : K]$。
于是环 $S/xS$ 作为 $R$-模的长度有限。因此
$S/\mathfrak m S$ 作为 $\kappa(\mathfrak m)$-向量空间的长度有限。
换言之，$\dim_{\kappa(\mathfrak m)} S/\mathfrak m S$ 有限
（引理 \ref{algebra-lemma-dimension-is-length}）。所以
$S/\mathfrak mS$ 是阿廷环（引理
\ref{algebra-lemma-finite-dimensional-algebra}）。
% P01-E0568: the source has subject–verb disagreement in “results ... implies”.
阿廷环的结构结果蕴含 (1) 与 (2)，例如可见引理
\ref{algebra-lemma-artinian-finite-length}。上面已经证明的有限性蕴含 (3)。
\end{proof}

\begin{lemma}
\label{algebra-lemma-finite-length-global}
设 $R$ 是分式域为 $K$ 的整环。设 $M$ 是 $K^{\oplus r}$
的 $R$-子模。假设 $R$ 是维数为 $1$ 的诺特环。
对任意非零元 $x \in R$，都有
$\text{length}_R(M/xM) < \infty$。
\end{lemma}

\begin{proof}
由于 $R$ 的维数为 $1$，$x$ 包含于有限多个素理想
$\mathfrak m_i$ 中，其中 $i = 1, \ldots, n$，并且它们都是极大理想。
由于 $R$ 是诺特环，$R/xR$ 是阿廷环，并且由命题
\ref{algebra-proposition-dimension-zero-ring} 和引理
\ref{algebra-lemma-artinian-finite-length}，有
$R/xR = \prod_{i = 1, \ldots, n} (R/xR)_{\mathfrak m_i}$。
因此 $M/xM$ 类似地分解为它在各 $\mathfrak m_i$ 处的局部化之积
$M/xM = \prod (M/xM)_{\mathfrak m_i}$。
对 $R_{\mathfrak m_i}$ 上的 $M_{\mathfrak m_i}$ 应用引理
\ref{algebra-lemma-finite-length}，可知每个
$M_{\mathfrak m_i}/xM_{\mathfrak m_i} = (M/xM)_{\mathfrak m_i}$
作为 $R_{\mathfrak m_i}$-模的长度有限。因此 $M/xM$ 作为 $R$-模
的长度有限，因为上述结论蕴含 $M/xM$ 有一个由 $R$-子模组成的有限滤过，
其逐次商同构于剩余域 $\kappa(\mathfrak m_i)$。
\end{proof}

\begin{lemma}[Krull--Akizuki]
\label{algebra-lemma-krull-akizuki}
设 $R$ 是分式域为 $K$ 的整环。设 $L/K$ 是有限域扩张。
假设 $R$ 是诺特环且 $\dim(R) = 1$。则任意满足
$R \subset A \subset L$ 的环 $A$ 都是诺特环。
\end{lemma}

\begin{proof}
设 $I \subset A$ 是非零理想。由引理
\ref{algebra-lemma-intersection-not-zero}，可以找到非零元
$x \in I \cap R$。于是有 $I/xA \subset A/xA$。
由引理 \ref{algebra-lemma-finite-length-global}，$R$-模 $A/xA$
作为 $R$-模的长度有限。因此 $I/xA$ 作为 $R$-模的长度有限。
所以 $I$ 作为 $A$ 的理想是有限生成的。
\end{proof}

\begin{lemma}
\label{algebra-lemma-exists-dvr}
设 $R$ 是分式域为 $K$ 的诺特局部整环。假设 $R$ 不是域。
设 $L/K$ 是有限生成域扩张。则存在分式域为 $L$ 的离散赋值环
% P01-E0569: the source statement omits the indefinite article before “discrete valuation ring”.
$A$，它支配 $R$。
\end{lemma}

\begin{proof}
若 $L$ 在 $K$ 上不是有限的，则选取 $L$ 在 $K$ 上的一组超越基
$x_1, \ldots, x_r$，并把 $R$ 替换为 $R[x_1, \ldots, x_r]$
在由 $\mathfrak m_R$ 和 $x_1, \ldots, x_r$ 生成的极大理想处的局部化。
于是可以假设 $K \subset L$ 有限。

\medskip\noindent
由引理 \ref{algebra-lemma-dominate-by-dimension-1}，可以假设 $\dim(R) = 1$。

\medskip\noindent
设 $A \subset L$ 是 $R$ 在 $L$ 中的整闭包。由引理
\ref{algebra-lemma-krull-akizuki}，它是诺特环。由引理
\ref{algebra-lemma-integral-overring-surjective}，存在素理想
$\mathfrak q \subset A$ 位于 $R$ 的极大理想上方。
由引理 \ref{algebra-lemma-characterize-dvr}，环 $A_{\mathfrak q}$
是支配 $R$ 的离散赋值环，如欲得。
\end{proof}

\section{因子分解}
\label{algebra-section-factoring}

\noindent
下面给出一些通常在本科一年级代数课程中讲授的概念及其相互关系。

\begin{definition}
\label{algebra-definition-irreducible-prime-element}
设 $R$ 是整环。
\begin{enumerate}
\item 若存在单位 $u \in R^*$，使得 $x = uy$，则称元素
$x, y \in R$ {\it 相伴}。
\item 若元素 $x \in R$ 非零、不是单位，并且每当
$x = yz$、$y, z \in R$ 时，$y$ 是单位或与 $x$ 相伴，
则称 $x$ 为{\it 不可约元}。
\item 若由元素 $x \in R$ 生成的理想是素理想，则称它为{\it 素元}。
\end{enumerate}
\end{definition}

\begin{lemma}
\label{algebra-lemma-easy-divisibility}
设 $R$ 是整环。设 $x, y \in R$。则 $x$、$y$ 相伴，当且仅当
$(x) = (y)$。
\end{lemma}

\begin{proof}
若对某个单位 $u \in R$ 有 $x = uy$，则 $(x) \subset (y)$；
又有 $y = u^{-1}x$，所以也有 $(y) \subset (x)$。
反过来，假设 $(x) = (y)$。则对某些
% P01-E0570: the frozen source uses the undefined ring A instead of R.
$f, g \in A$，有 $x = fy$ 且 $y = gx$。于是 $x = fg x$；
由于 $R$ 是整环，$fg = 1$。所以 $x$ 与 $y$ 相伴。
\end{proof}

\begin{lemma}
\label{algebra-lemma-factorization-exists}
设 $R$ 是整环。考虑下列条件：
\begin{enumerate}
\item 环 $R$ 满足主理想升链条件。
\item 每个非零非单位元 $a \in R$ 都有分解
$a = b_1 \ldots b_k$，其中每个 $b_i$ 都是 $R$ 的不可约元。
\end{enumerate}
则 (1) 蕴含 (2)。
\end{lemma}

\begin{proof}
设 $x$ 是一个不能分解为不可约元之积的非零非单位元。置
$x_1 = x$。可以写成 $x = yz$，其中
% P01-E0571: the source incorrectly requires both factors to be nonirreducible, rather than only nonunits.
$y$ 和 $z$ 都既不是不可约元，也不是单位。
于是，或者 $y$ 没有不可约分解，此时置 $x_2 = y$；
或者 $z$ 没有不可约分解，此时置 $x_2 = z$。
如此继续，得到 $R$ 中的元素序列
$$
\ldots | x_3 | x_2 | x_1
$$
其中 $x_n/x_{n + 1}$ 不是单位。这给出严格递增的主理想序列
$(x_1) \subset (x_2) \subset (x_3) \subset \ldots$，证明完成。
\end{proof}

\begin{definition}
\label{algebra-definition-UFD}
若整环 $R$ 满足如下性质，则称它为{\it 唯一分解整环}，简称 {\it UFD}：
若 $x \in R$ 是非零非单位元，则 $x$ 有不可约分解；并且，若
$$
x = a_1 \ldots a_m = b_1 \ldots b_n
$$
是两个不可约分解，则 $n = m$，并且存在置换
$\sigma : \{1, \ldots, n\} \to \{1, \ldots, n\}$，
使得 $a_i$ 与 $b_{\sigma(i)}$ 相伴。
\end{definition}

\begin{lemma}
\label{algebra-lemma-characterize-UFD}
设 $R$ 是整环。假设每个非零非单位元都能分解为不可约元之积。
则 $R$ 是 UFD，当且仅当每个不可约元都是素元。
\end{lemma}

\begin{proof}
假设 $R$ 是 UFD，并设 $x \in R$ 是不可约元。
设 $ab \in (x)$，亦即 $ab = cx$。选取分解
$a = a_1 \ldots a_n$、$b = b_1 \ldots b_m$ 和
$c = c_1 \ldots c_r$。由分解
$$
a_1 \ldots a_n b_1 \ldots b_m = c_1 \ldots c_r x
$$
的唯一性，可知 $x$ 与元素
% P01-E0572: the source range “a_1, ..., b_m” omits the intermediate endpoint and start.
$a_1, \ldots, b_m$ 中的一个相伴。换言之，
$a \in (x)$ 或 $b \in (x)$，所以 $x$ 是素元。

\medskip\noindent
假设每个不可约元都是素元。需要证明不可约分解在相差置换及取相伴元的意义下唯一。
设 $a_1 \ldots a_m = b_1 \ldots b_n$，其中 $a_i$ 和 $b_j$
均不可约。由于 $a_1$ 是素元，对某个 $j$ 有 $b_j \in (a_1)$。
重新编号后，可以假设 $b_1 \in (a_1)$。于是 $b_1 = a_1 u$；
由于 $b_1$ 不可约，$u$ 是单位。所以 $a_1$ 与 $b_1$ 相伴，并且
% P01-E0573: the frozen source swaps the a- and b-factor endpoints in the cancelled equality.
$a_2 \ldots a_n = ub_2\ldots b_m$。对 $n + m$ 作归纳，
可知 $n = m$，并且对 $i = 2, \ldots, n$，$a_i$ 与
$b_{\sigma(i)}$ 相伴，如欲得。
\end{proof}

\begin{lemma}
\label{algebra-lemma-characterize-UFD-height-1}
设 $R$ 是诺特整环。则 $R$ 是 UFD，当且仅当每个高度为 $1$
的素理想都是主理想。
\end{lemma}

\begin{proof}
假设 $R$ 是 UFD，并设 $\mathfrak p$ 是高度为 1 的素理想。
取非零元 $x \in \mathfrak p$，并设 $x = a_1 \ldots a_n$
是不可约分解。由于 $\mathfrak p$ 是素理想，对某个 $i$ 有
$a_i \in \mathfrak p$。由引理 \ref{algebra-lemma-characterize-UFD}，
理想 $(a_i)$ 是素理想。因为 $\mathfrak p$ 的高度为 $1$，
可知 $(a_i) = \mathfrak p$。

\medskip\noindent
假设每个高度为 $1$ 的素理想都是主理想。由于 $R$ 是诺特环，
每个非零非单位元 $x$ 都有不可约分解，见引理
\ref{algebra-lemma-factorization-exists}。由引理
\ref{algebra-lemma-characterize-UFD}，只须证明不可约元 $x$ 是素元。
设 $(x) \subset \mathfrak p$，其中后者是
$(x)$ 上方的极小素理想。由引理 \ref{algebra-lemma-minimal-over-1}，
$\mathfrak p$ 的高度为 $1$。根据假设，$\mathfrak p = (y)$。
于是 $x = yz$；由于 $x$ 不可约，$z$ 是单位。
所以 $(x) = (y)$，从而 $x$ 是素元。
\end{proof}

\begin{lemma}[Nagata 唯一分解性判据]
\label{algebra-lemma-invert-prime-elements}
\begin{reference}
\cite[Lemma 2]{Nagata-UFD}
\end{reference}
设 $A$ 是整环。设 $S \subset A$ 是由素元生成的乘法子集。
设 $x \in A$ 不可约。则
\begin{enumerate}
\item $x$ 在 $S^{-1}A$ 中的像是不可约元或单位；并且
\item $x$ 是素元，当且仅当 $x$ 在 $S^{-1}A$ 中的像是
$S^{-1}A$ 中的素元或单位。
\end{enumerate}
此外，$A$ 是 UFD，当且仅当 $A$ 的每个非零非单位元都有不可约分解，
并且 $S^{-1}A$ 是 UFD。
\end{lemma}

\begin{proof}
设对 $\alpha, \beta \in S^{-1}A$ 有 $x = \alpha \beta$。
则对某些 $a, b \in A$、$s, s' \in S$，有
$\alpha = a/s$ 且 $\beta = b/s'$。所以 $ss'x = ab$。
根据假设，可以对某些素元 $p_i$ 写成
$ss' = p_1 \ldots p_r$。对每个 $i$，元素 $p_i$ 整除
$a$ 或 $b$。作除法后，得到分解 $x = a' b'$，并且对某些
$s'', s''' \in S$，有 $a = s'' a'$、$b = s''' b'$。
由于 $x$ 不可约，$a'$ 或 $b'$ 是单位。回溯上述步骤可知
$\alpha$ 或 $\beta$ 是单位。这证明了 (1)。

\medskip\noindent
假设 $x$ 是素元。则 $A/(x)$ 是整环。因此
$S^{-1}A/xS^{-1}A = S^{-1}(A/(x))$ 是整环或零环。
所以 $x$ 映到素元或单位。

\medskip\noindent
假设 $x$ 在 $S^{-1}A$ 中的像是单位。则对某些
$s \in S$ 和 $y \in A$，有 $y x = s$。根据假设，
$s = p_1 \ldots p_r$，其中 $p_i$ 是素元。对每个 $i$，
$p_i$ 整除 $y$ 或 $p_i$ 整除 $x$。在后一种情形中，
$p_i$ 与 $x$ 相伴（因为 $x$ 不可约），结论得证。
但若对所有 $i = 1, \ldots, r$ 都出现前一种情形，则 $x$
是单位，矛盾。

\medskip\noindent
假设 $x$ 在 $S^{-1}A$ 中的像是素元。设 $a, b \in A$ 且
$ab \in (x)$。则对某些 $s \in S$ 和 $y \in A$，有
$sa = xy$ 或 $sb = xy$。设出现第一种情形。根据假设，
$s = p_1 \ldots p_r$，其中 $p_i$ 是素元。对每个 $i$，
$p_i$ 整除 $y$ 或 $p_i$ 整除 $x$。在后一种情形中，
$p_i$ 与 $x$ 相伴（因为 $x$ 不可约），结论得证。
若对所有 $i = 1, \ldots, r$ 都出现前一种情形，则
$a \in (x)$，如欲得。这完成了 (2) 的证明。

\medskip\noindent
引理的最后一个陈述由 (1)、(2) 以及引理
\ref{algebra-lemma-characterize-UFD} 得出。
\end{proof}

\begin{lemma}
\label{algebra-lemma-UFD-ascending-chain-condition-principal-ideals}
UFD 满足主理想升链条件。
\end{lemma}

\begin{proof}
考虑 $R$ 中主理想的升链
$(a_1) \subset (a_2) \subset (a_3) \subset \ldots$。
写成 $a_1 = p_1^{e_1} \ldots p_r^{e_r}$，其中 $p_i$ 是素元。
于是对某些 $0 \leq c_i \leq e_i$，$a_n$ 与
$p_1^{c_1} \ldots p_r^{c_r}$ 相伴。因为只有有限种可能，结论成立。
\end{proof}

\begin{lemma}
\label{algebra-lemma-factoring-in-polynomial}
设 $R$ 是整环。假设 $R$ 满足主理想升链条件。
则 $R$ 上的多项式环也满足这一性质。
\end{lemma}

\begin{proof}
考虑 $R[x]$ 中主理想的升链
$(f_1) \subset (f_2) \subset (f_3) \subset \ldots$。
由于 $f_{n + 1}$ 整除 $f_n$，该序列中的次数递减。
所以对所有 $n \gg 0$，$f_n$ 都有固定次数 $d \geq 0$。
设 $a_n$ 是 $f_n$ 的首项系数。条件 $f_n \in (f_{n + 1})$
蕴含对所有 $n$，$a_{n + 1}$ 整除 $a_n$。由对 $R$ 的假设，
对所有充分大的 $n$，$a_{n + 1}$ 与 $a_n$ 相伴（引理
\ref{algebra-lemma-easy-divisibility}）。所以对充分大的 $n$，
有 $f_n = u f_{n + 1}$，其中 $u \in R$（由次数原因）是单位
（因为 $a_n$ 与 $a_{n + 1}$ 相伴）。
\end{proof}

\begin{lemma}
\label{algebra-lemma-polynomial-ring-UFD}
UFD 上的多项式环是 UFD。特别地，若 $k$ 是域，则
$k[x_1, \ldots, x_n]$ 是 UFD。
\end{lemma}

\begin{proof}
设 $R$ 是 UFD。则 $R$ 满足主理想升链条件
（引理 \ref{algebra-lemma-UFD-ascending-chain-condition-principal-ideals}），
所以 $R[x]$ 满足主理想升链条件（引理
\ref{algebra-lemma-factoring-in-polynomial}），从而 $R[x]$ 的每个元素
都有不可约分解（引理 \ref{algebra-lemma-factorization-exists}）。
设 $S \subset R$ 是由素元生成的乘法子集。因为 $R$ 的每个非单位元
都是素元之积，可知 $K = S^{-1}R$ 是 $R$ 的分式域。
注意，$R$ 的每个素元都映到 $R[x]$ 的素元，并且
$S^{-1}(R[x]) = S^{-1}R[x] = K[x]$ 是 UFD（甚至是 PID）。
所以可以应用引理 \ref{algebra-lemma-invert-prime-elements} 得出结论。
\end{proof}

\begin{lemma}
\label{algebra-lemma-UFD-normal}
唯一分解整环是正规的。
\end{lemma}

\begin{proof}
设 $R$ 是 UFD。设 $x$ 是 $R$ 的分式域中整于 $R$ 的元素。
设 $x^d - a_1 x^{d - 1} - \ldots - a_d = 0$，其中
$a_i \in R$。可以写成
$x = u p_1^{e_1} \ldots p_r^{e_r}$，其中 $u$ 是单位，
$e_i \in \mathbf{Z}$，而 $p_1, \ldots, p_r$ 是两两不相伴的不可约元。
为证明引理，需要说明 $e_i \geq 0$。否则，不妨设 $e_1 < 0$；
则对 $N \gg 0$，有
$$
u^d p_2^{de_2 + N} \ldots p_r^{de_r + N} =
p_1^{-de_1}p_2^N \ldots p_r^N(
\sum\nolimits_{i = 1, \ldots, d} a_i x^{d - i} ) \in (p_1)
$$
这与 $R$ 中分解的唯一性矛盾。
\end{proof}

\begin{definition}
\label{algebra-definition-PID}
若整环 $R$ 的每个理想都是主理想，则称它为{\it 主理想整环}，
简称 {\it PID}。
\end{definition}

\begin{lemma}
\label{algebra-lemma-PID-UFD}
主理想整环是唯一分解整环。
\end{lemma}

\begin{proof}
因为 PID 是诺特环，这由引理
\ref{algebra-lemma-characterize-UFD-height-1} 得出。
\end{proof}

\begin{definition}
\label{algebra-definition-dedekind-domain}
若整环 $R$ 的每个非零理想 $I \subset R$ 都能写成非零素理想的乘积
$$
I = \mathfrak p_1 \ldots \mathfrak p_r
$$
并且在相差 $\mathfrak p_i$ 的置换的意义下唯一，
则称它为 {\it Dedekind 整环}。
\end{definition}

\begin{lemma}
\label{algebra-lemma-PID-dedekind}
PID 是 Dedekind 整环。
\end{lemma}

\begin{proof}
设 $R$ 是 PID。由于 $R$ 的每个非零理想都是主理想，并且 $R$
是 UFD（引理 \ref{algebra-lemma-PID-UFD}），结论由如下事实得出：
$R$ 中的每个不可约元都是素元（引理
\ref{algebra-lemma-characterize-UFD}），所以元素的分解转化为素理想分解。
\end{proof}

\begin{lemma}
\label{algebra-lemma-product-ideals-principal}
\begin{slogan}
理想的乘积作为模可逆，当且仅当两个因子都可逆。
\end{slogan}
设 $A$ 是环。设 $I$ 和 $J$ 是 $A$ 的非零理想，并且对某个
非零因子 $f \in A$ 有 $IJ = (f)$。则 $I$ 和 $J$ 都是有限生成理想，
并且作为 $A$-模都是秩为 $1$ 的有限局部自由模。
\end{lemma}

\begin{proof}
由引理 \ref{algebra-lemma-finite-projective}，只须证明 $I$ 和 $J$ 是秩为 $1$
的有限局部自由 $A$-模。为此，写成
$f = \sum_{i = 1, \ldots, n} x_i y_i$，其中 $x_i \in I$ 且
$y_i \in J$。还可以对某些 $a_i \in A$ 写成
$x_i y_i = a_i f$。因为 $f$ 是非零因子，可知 $\sum a_i = 1$。
所以只须说明每个 $I_{a_i}$ 和 $J_{a_i}$ 都是 $A_{a_i}$ 上秩为 $1$
的自由模。把 $A$ 替换为 $A_{a_i}$ 后，可知对某些
$x \in I$ 和 $y \in J$ 有 $f = xy$。注意 $x$ 和 $y$ 都是非零因子。
我们断言 $I = (x)$ 且 $J = (y)$，这就完成了证明。事实上，
若 $x' \in I$，则对某个 $a \in A$ 有 $x'y = af = axy$。
所以 $x' = ax$，结论得证。
\end{proof}

\begin{lemma}
\label{algebra-lemma-characterize-Dedekind}
设 $R$ 是环。下列条件等价：
% P01-E0574: the frozen source omits punctuation after “equivalent”.
\begin{enumerate}
\item $R$ 是 Dedekind 整环；
\item $R$ 是诺特整环，并且对每个非零极大理想 $\mathfrak m$，
局部环 $R_{\mathfrak m}$ 都是离散赋值环；并且
\item $R$ 是诺特正规整环，并且 $\dim(R) \leq 1$。
\end{enumerate}
\end{lemma}

\begin{proof}
假设 (1) 成立。由于在 Dedekind 整环的定义中并未假设 $R$ 是诺特环，
所以论证并非平凡。设 $\mathfrak p \subset R$ 是素理想。
由定义中分解的唯一性，注意到
$\mathfrak p \not = \mathfrak p^2$。选取
$x \in \mathfrak p$，使得 $x \not \in \mathfrak p^2$。
设 $y \in \mathfrak p$ 是第二个元素（例如 $y = 0$）。
写成 $(x, y) = \mathfrak p_1 \ldots \mathfrak p_r$。
由于 $(x, y) \subset \mathfrak p$，至少有一个素理想
$\mathfrak p_i$ 包含于 $\mathfrak p$。但因为
$x \not \in \mathfrak p^2$，至多只有一个。因此
$\mathfrak p_1, \ldots, \mathfrak p_r$ 中恰有一个包含于
$\mathfrak p$，不妨设 $\mathfrak p_1 \subset \mathfrak p$。
可知对 $y$ 的每种选取，
$(x, y)R_\mathfrak p = \mathfrak p_1R_\mathfrak p$ 都是素理想。
我们断言 $(x)R_\mathfrak p = \mathfrak pR_\mathfrak p$。
事实上，选取 $y \in \mathfrak p$。把上述论证应用于 $y^2$，
可知 $(x, y^2)R_\mathfrak p$ 是素理想。因此
$y \in (x, y^2)R_\mathfrak p$，亦即在 $R_\mathfrak p$ 中
$y = ax + by^2$。所以
$(1 - by)y = ax \in (x)R_\mathfrak p$，亦即
$y \in (x)R_\mathfrak p$，如欲得。

\medskip\noindent
重新写成 $(x) = \mathfrak p_1 \ldots \mathfrak p_r$，其中
$\mathfrak p_1 \subset \mathfrak p$。可知
$\mathfrak p_1 R_\mathfrak p = \mathfrak p R_\mathfrak p$，
亦即 $\mathfrak p_1 = \mathfrak p$。此外，由引理
\ref{algebra-lemma-product-ideals-principal}，
$\mathfrak p_1 = \mathfrak p$ 是 $R$ 的有限生成理想。
由引理 \ref{algebra-lemma-cohen}，可知 $R$ 是诺特环。
% P01-E0575: the source mismatches m/p and overstates “every prime ideal”, including the zero prime.
此外，对每个素理想 $\mathfrak p$，$R_\mathfrak m$ 都是离散赋值环；
见引理 \ref{algebra-lemma-characterize-dvr}。

\medskip\noindent
(2) 与 (3) 的等价性由引理
\ref{algebra-lemma-normality-is-local} 和
\ref{algebra-lemma-characterize-dvr} 得出。假设 (2) 和 (3) 成立。
设 $I \subset R$ 是理想。我们构造 $I$ 的分解。
若 $I$ 是素理想，则无须证明。否则，选取
$I \subset \mathfrak p$，其中 $\mathfrak p \subset R$ 是极大理想。
令 $J = \{x \in R \mid x \mathfrak p \subset I\}$。
我们断言 $J \mathfrak p = I$。只须在 $R$ 的各极大理想
$\mathfrak m$ 处局部化后检验这一点（$J$ 的形成与局部化交换，
并且使用引理 \ref{algebra-lemma-characterize-zero-local}）。
于是，或者 $\mathfrak p R_\mathfrak m = R_\mathfrak m$，
此时结论显然；或者
$\mathfrak p R_\mathfrak m = \mathfrak m R_\mathfrak m$。
在后一种情形中，$\mathfrak p R_\mathfrak m = (\pi)$，
而 $\mathfrak p$ 为主理想时结论立即成立。
由诺特归纳，理想 $J$ 有分解，从而得到所需的 $I$ 的分解。
我们略去分解唯一性的证明。
\end{proof}

\noindent
下面是 Krull--Akizuki 引理的一个变体。

\begin{lemma}
\label{algebra-lemma-integral-closure-Dedekind}
设 $A$ 是维数为 $1$、分式域为 $K$ 的诺特整环。
设 $L/K$ 是有限扩张。设 $B$ 是 $A$ 在 $L$ 中的整闭包。
则 $B$ 是 Dedekind 整环，并且
$\Spec(B) \to \Spec(A)$ 是满射、具有有限纤维，并诱导有限剩余域扩张。
\end{lemma}

\begin{proof}
由 Krull--Akizuki 定理（引理 \ref{algebra-lemma-krull-akizuki}），
环 $B$ 是诺特环。由引理
\ref{algebra-lemma-integral-sub-dim-equal}，$\dim(B) = 1$。
所以由引理 \ref{algebra-lemma-characterize-Dedekind}，$B$ 是 Dedekind 整环。
谱上映射的满性由引理
\ref{algebra-lemma-integral-overring-surjective} 得出。
最后两个陈述由引理
\ref{algebra-lemma-finite-extension-residue-fields-dimension-1} 得出。
\end{proof}

\section{消失阶}
\label{algebra-section-orders-of-vanishing}

\begin{lemma}
\label{algebra-lemma-ord-additive}
设 $R$ 是维数为 $1$ 的半局部诺特环。
若 $a, b \in R$ 是非零因子，则
$$
\text{length}_R(R/(ab)) =
\text{length}_R(R/(a)) +
\text{length}_R(R/(b))
$$
且这些长度都是有限的。
\end{lemma}

\begin{proof}
我们已在引理 \ref{algebra-lemma-finite-length-global} 中证明了有限性。
% P01-E0576: the frozen source says “Use length is additive” instead of “Use that length is additive”.
可加性来自短正合列
$0 \to R/(a) \to R/(ab) \to R/(b) \to 0$，其中第一个映射
由乘以 $b$ 给出。（使用长度的可加性；
见引理 \ref{algebra-lemma-length-additive}。）
\end{proof}

\begin{definition}
\label{algebra-definition-ord}
设 $K$ 是域，且 $R \subset K$ 是一个
局部的\footnote{当 $R$ 仅为半局部环时，我们也可以作此定义，
但这大概绝不会真正是你想要的！}
诺特子环，其维数为 $1$，分式域为 $K$。
在这种情形下，我们定义沿 $R$ 的\textit{消失阶}
$$
\text{ord}_R : K^* \longrightarrow \mathbf{Z}
$$
如下：若 $x \in R$，则
$$
\text{ord}_R(x) = \text{length}_R(R/(x))
$$
并且对两个非零元 $x, y \in R$，置
$\text{ord}_R(x/y) = \text{ord}_R(x) - \text{ord}_R(y)$。
\end{definition}

\noindent
我们可以用消失阶比较向量空间中的格。定义如下。

\begin{definition}
\label{algebra-definition-lattice}
设 $R$ 是维数为 $1$、分式域为 $K$ 的诺特局部整环。
设 $V$ 是有限维 $K$-向量空间。
$V$ 中的一个\textit{格}是有限 $R$-子模 $M \subset V$，
满足 $V = K \otimes_R M$。
\end{definition}

\noindent
条件 $V = K \otimes_R M$ 表示 $M$ 包含向量空间 $V$ 的一组基。
我们指出，文献中许多地方仅在环 $R$ 是离散赋值环时定义格。
若 $R$ 是离散赋值环，则任一格都是自由 $R$-模；一般情形下
则未必如此。

\begin{lemma}
\label{algebra-lemma-compare-lattices}
设 $R$ 是维数为 $1$、分式域为 $K$ 的诺特局部整环。
设 $V$ 是有限维 $K$-向量空间。
\begin{enumerate}
\item 若 $M$ 是 $V$ 中的格，且 $M \subset M' \subset V$ 是
包含 $M$ 的 $V$ 的一个 $R$-子模，则下列条件等价：
\begin{enumerate}
\item $M'$ 是格；
\item $\text{length}_R(M'/M)$ 有限；
\item $M'$ 是有限生成的。
\end{enumerate}
\item 若 $M$ 是 $V$ 中的格，而 $M' \subset M$ 是 $M$ 的一个
$R$-子模，则 $M'$ 是格当且仅当
$\text{length}_R(M/M')$ 有限。
\item 若 $M$、$M'$ 是 $V$ 中的格，则
$M \cap M'$ 和 $M + M'$ 也是格。
\item 若 $M \subset M' \subset M'' \subset V$ 是 $V$ 中的格，则
$$
\text{length}_R(M''/M) =
\text{length}_R(M'/M) +
\text{length}_R(M''/M').
$$
\item 若 $M$、$M'$、$N$、$N'$ 是 $V$ 中的格，并且
$N \subset M \cap M'$、$M + M' \subset N'$，则
\begin{eqnarray*}
& & \text{length}_R(M/M \cap M') - \text{length}_R(M'/M \cap M')\\
& = &
\text{length}_R(M/N) - \text{length}_R(M'/N) \\
& = &
\text{length}_R(M + M' / M') - \text{length}_R(M + M'/M) \\
& = &
\text{length}_R(N' / M') - \text{length}_R(N'/M)
\end{eqnarray*}
\end{enumerate}
\end{lemma}

\begin{proof}
证明 (1)。假设 (1)(a) 成立。设 $y_1, \ldots, y_m$ 生成 $M'$。
则每个 $y_i = x_i/f_i$，其中 $x_i \in M$，
而 $f_i \in R$ 非零。
因此 $f_1 \ldots f_m M' \subset M$。
由于 $R$ 是维数为 $1$ 的诺特局部环，
对某个 $n$ 有 $\mathfrak m^n \subset (f_1 \ldots f_m)$
（例如可合用引理 \ref{algebra-lemma-one-equation} 与命题
\ref{algebra-proposition-dimension-zero-ring}，或合用引理
\ref{algebra-lemma-finite-length} 与 \ref{algebra-lemma-length-infinite}）。
% P01-E0577: the frozen source omits sentence punctuation after “for some n”.
换言之，对某个 $n$ 有 $\mathfrak m^nM' \subset M$。
所以由引理 \ref{algebra-lemma-length-finite}，
$\text{length}(M'/M) < \infty$；也就是说 (1)(b) 成立。
假设 (1)(b) 成立。则 $M'/M$ 是有限 $R$-模
（见引理 \ref{algebra-lemma-finite-length-finite}）。
所以 $M'$ 作为两个有限 $R$-模的扩张也是有限 $R$-模，
从而 (1)(c) 成立。蕴含
(1)(c) $\Rightarrow$ (1)(a) 来自定义 \ref{algebra-definition-lattice}
之后的说明。

\medskip\noindent
证明 (2)。设
$M$ 是 $V$ 中的格，而 $M' \subset M$ 是一个 $R$-子模。
我们已在 (1) 中看到，若 $M'$ 是格，则
$\text{length}_R(M/M') < \infty$。反之，假设
$\text{length}_R(M/M') < \infty$。由于 $R$ 是诺特环，
$M'$ 是有限生成的；并且由引理 \ref{algebra-lemma-length-infinite}，
对某个 $n$ 有 $\mathfrak m^n M \subset M'$。
因此 $M'$ 包含 $V$ 的一组基，故 $M'$ 是格。

\medskip\noindent
证明 (3)。设 $M$、$M'$ 是 $V$ 中的格。
由于 $R$ 是诺特环，$M$ 的子模 $M \cap M'$ 是有限的。
因为 $M$ 是格，我们可以取
$x_1, \ldots, x_n \in M$，使其构成 $V$ 的一组 $K$-基。
因为 $M'$ 是格，可以写成 $x_i = y_i/f_i$，其中
$y_i \in M'$ 且 $f_i \in R$。于是 $f_ix_i \in M \cap M'$。
所以 $M \cap M'$ 也是格。
$M + M'$ 是格这一事实由 (1) 得出。

\medskip\noindent
(4) 来自长度的可加性（引理 \ref{algebra-lemma-length-additive}）以及如下正合列。
% P01-E0578: the frozen source omits punctuation after this displayed exact sequence.
$$
0 \to M'/M \to M''/M \to M''/M' \to 0
$$
(5) 由反复应用 (4) 得出。
\end{proof}

\begin{definition}
\label{algebra-definition-distance}
设 $R$ 是维数为 $1$、分式域为 $K$ 的诺特局部整环。
设 $V$ 是有限维 $K$-向量空间。
设 $M$、$M'$ 是 $V$ 中的两个格。$M$ 与 $M'$ 之间的\textit{距离}
是引理 \ref{algebra-lemma-compare-lattices} 第 (5) 部分中的整数
$$
d(M, M') = \text{length}_R(M/M \cap M') - \text{length}_R(M'/M \cap M')
$$
\end{definition}

\noindent
特别地，若 $M' \subset M$，则
$d(M, M') = \text{length}_R(M/M')$。

\begin{lemma}
\label{algebra-lemma-properties-distance-function}
设 $R$ 是维数为 $1$、分式域为 $K$ 的诺特局部整环。
设 $V$ 是有限维 $K$-向量空间。
这个距离函数具有如下性质：给定 $V$ 中任意三个格
$M$、$M'$、$M''$，都有
$$
d(M, M'') = d(M, M') + d(M', M'')
$$
特别地，$d(M, M') = - d(M', M)$。
\end{lemma}

\begin{proof}
略。
\end{proof}

\begin{lemma}
\label{algebra-lemma-order-vanishing-determinant}
设 $R$ 是维数为 $1$、分式域为 $K$ 的诺特局部整环。
设 $V$ 是有限维 $K$-向量空间。
设 $\varphi : V \to V$ 是 $K$-线性同构。
对任意格 $M \subset V$，有
$$
d(M, \varphi(M)) = \text{ord}_R(\det(\varphi))
$$
\end{lemma}

\begin{proof}
如下可见，整数 $d(M, \varphi(M))$ 与格 $M$ 的选择无关。
设 $M'$ 是另一个这样的格。则
\begin{eqnarray*}
d(M, \varphi(M)) & = & d(M, M') + d(M', \varphi(M)) \\
& = & d(M, M') + d(\varphi(M'), \varphi(M)) + d(M', \varphi(M'))
\end{eqnarray*}
由于 $\varphi$ 是同构，有
$d(\varphi(M'), \varphi(M)) = d(M', M) = -d(M, M')$，因而
$d(M, \varphi(M)) = d(M', \varphi(M'))$。此外，等式（即本引理的等式）
两边关于 $\varphi$ 都是可加的，也就是说
$$
\text{ord}_R(\det(\varphi \circ \psi))
=
\text{ord}_R(\det(\varphi))
+
\text{ord}_R(\det(\psi))
$$
以及
% P01-E0579: the frozen display is missing the final parenthesis closing d(...).
\begin{eqnarray*}
d(M, \varphi(\psi((M))) & = &
d(M, \psi(M)) + d(\psi(M), \varphi(\psi(M))) \\
& = & d(M, \psi(M)) + d(M, \varphi(M))
\end{eqnarray*}
其中第二个等式使用了上面证明的独立性。因此，只须对
$\text{GL}(V)$ 的生成元证明本引理。选取同构
$K^{\oplus n} \cong V$。则 $\text{GL}(V) = \text{GL}_n(K)$
由初等矩阵 $E$ 生成。
当 $E$ 为单位矩阵时，结论显然。
若 $E = E_{ij}(\lambda)$，其中 $i \not = j$、$\lambda \in K$、
$\lambda \not = 0$，例如
$$
E_{12}(\lambda) =
\left(
\begin{matrix}
1 & \lambda & \ldots \\
0 & 1 & \ldots \\
\ldots & \ldots & \ldots
\end{matrix}
\right)
$$
则相对于另一组基，它变为 $E_{12}(1)$。
取格 $R^{\oplus n} \subset K^{\oplus n}$，结论对
$E = E_{12}(1)$ 是显然的。最后，若 $E = E_i(a)$，
其中 $a \in K^*$，例如
$$
E_1(a) =
\left(
\begin{matrix}
a & 0 & \ldots \\
0 & 1 & \ldots \\
\ldots & \ldots & \ldots
\end{matrix}
\right)
$$
% P01-E0580: the frozen source uses undefined b in R^{\oplus b}; the ambient rank is n.
则 $E_1(a)(R^{\oplus b}) = aR \oplus R^{\oplus n - 1}$，并且显然
$d(R^{\oplus n}, aR \oplus R^{\oplus n - 1})
= \text{ord}_R(a)$，如所需。
\end{proof}

\begin{lemma}
\label{algebra-lemma-finite-extension-dim-1}
设 $A \to B$ 是环映射。假设
\begin{enumerate}
\item $A$ 是维数为 $1$ 的诺特局部整环；
\item $A \subset B$ 是整环的有限扩张。
\end{enumerate}
设 $L/K$ 是相应的分式域有限扩张。
设 $y \in L^*$，并令 $x = \text{Nm}_{L/K}(y)$。
在这种情形下，$B$ 是半局部环。
设 $\mathfrak m_i$（$i = 1, \ldots, n$）是 $B$ 的极大理想。
则
$$
\text{ord}_A(x) =
\sum\nolimits_i
[\kappa(\mathfrak m_i) : \kappa(\mathfrak m_A)]
\text{ord}_{B_{\mathfrak m_i}}(y)
$$
其中 $\text{ord}$ 如定义 \ref{algebra-definition-ord} 所定义。
\end{lemma}

\begin{proof}
由引理 \ref{algebra-lemma-finite-in-codim-1}，环 $B$ 是半局部的。
将其写成 $y = b/b'$，其中 $b, b' \in B$。
由 $\text{ord}$ 的可加性和 $\text{Nm}$ 的乘法性，只须对
$y = b$ 或 $y = b'$ 证明本引理。换言之，可以假设 $y \in B$。
此时公式右边为
$$
\sum [\kappa(\mathfrak m_i) : \kappa(\mathfrak m_A)]
\text{length}_{B_{\mathfrak m_i}}((B/yB)_{\mathfrak m_i})
$$
由引理 \ref{algebra-lemma-pushdown-module}，它等于
$\text{length}_A(B/yB)$。由引理 \ref{algebra-lemma-order-vanishing-determinant}，
有
$$
\text{length}_A(B/yB) = d(B, yB) =
\text{ord}_A(\det\nolimits_K(L \xrightarrow{y} L)).
$$
根据定义，$x = \text{Nm}_{L/K}(y) = \det\nolimits_K(L \xrightarrow{y} L)$，
故引理得证。
\end{proof}

\section{拟有限映射}
\label{algebra-section-quasi-finite}

\noindent
考虑一个有限型环映射 $R \to S$。
若某点在其纤维中是孤立点，就称映射
$\Spec(S) \to \Spec(R)$ 在该点拟有限。
这表示该纤维在该点处为零维。
本节研究这一重要但技术性的概念的基本性质。
更深入的内容可见下一节。

\begin{lemma}
\label{algebra-lemma-isolated-point}
设 $k$ 是域。
设 $S$ 是有限型 $k$-代数。
设 $\mathfrak q$ 是 $S$ 的素理想。
下列条件等价：
\begin{enumerate}
\item $\mathfrak q$ 是 $\Spec(S)$ 的孤立点；
\item $S_{\mathfrak q}$ 在 $k$ 上有限；
\item 存在 $g \in S$、$g \not\in \mathfrak q$，使得
$D(g) = \{ \mathfrak q \}$；
\item $\dim_{\mathfrak q} \Spec(S) = 0$；
\item $\mathfrak q$ 是 $\Spec(S)$ 的闭点，并且
$\dim(S_{\mathfrak q}) = 0$；
\item 域扩张 $\kappa(\mathfrak q)/k$ 是有限的，并且
$\dim(S_{\mathfrak q}) = 0$。
\end{enumerate}
在这种情形下，对某个有限型 $k$-代数 $S'$，有
$S = S_{\mathfrak q} \times S'$。此外，(3) 中的元素 $g$
还满足 $S_{\mathfrak q} = S_g$。
\end{lemma}

\begin{proof}
假设 $\mathfrak q$ 是 $\Spec(S)$ 的孤立点，即
$\{\mathfrak q\}$ 在 $\Spec(S)$ 中是开集。
因为 $\Spec(S)$ 是 Jacobson 空间（见引理
\ref{algebra-lemma-finite-type-field-Jacobson} 和
\ref{algebra-lemma-jacobson}），所以 $\mathfrak q$ 是闭点。
因此 $\{\mathfrak q\}$ 在 $\Spec(S)$ 中既开又闭。
由引理 \ref{algebra-lemma-disjoint-decomposition} 和
\ref{algebra-lemma-disjoint-implies-product}，可以写成
$S = S_1 \times S_2$，其中 $\mathfrak q$ 对应于
$\Spec(S_1)$ 的唯一一点。
所以 $S_1 = S_{\mathfrak q}$ 是有限型 $k$-代数，且维数为零。
因此它在 $k$ 上有限，例如可由引理
\ref{algebra-lemma-Noether-normalization} 得出。
我们已经证明 (1) 蕴含 (2)。

\medskip\noindent
假设 $S_{\mathfrak q}$ 在 $k$ 上有限。
则由引理 \ref{algebra-lemma-finite-dimensional-algebra}，
$S_{\mathfrak q}$ 是阿廷局部环。因此由引理
\ref{algebra-lemma-artinian-finite-length}，
$\Spec(S_{\mathfrak q}) = \{\mathfrak qS_{\mathfrak q}\}$。
考虑正合列 $0 \to K \to S \to S_{\mathfrak q}
\to Q \to 0$。显然 $K_{\mathfrak q} = Q_{\mathfrak q} = 0$。
又因为 $S$ 是诺特环，$K$ 是有限 $S$-模；而由于
$S_{\mathfrak q}$ 在 $k$ 上有限，$Q$ 是有限 $S$-模。
所以存在 $g \in S$、$g \not \in \mathfrak q$，使得
$K_g = Q_g = 0$。于是 $S_{\mathfrak q} = S_g$，且
$D(g) = \{ \mathfrak q \}$。我们已经证明 (2) 蕴含 (3)。

\medskip\noindent
假设 $D(g) =  \{ \mathfrak q \}$。按照 $\Spec(S)$ 上拓扑的构造，
$D(g)$ 是开集，因而 $\mathfrak q$ 是 $\Spec(S)$ 的孤立点。
我们已经证明 (3) 蕴含 (1)。
换言之，(1)、(2)、(3) 等价。

\medskip\noindent
假设 $\dim_{\mathfrak q} \Spec(S) = 0$。这表示
$\mathfrak q$ 在 $\Spec(S)$ 中有某个零维开邻域。
于是它有一个形如 $D(g)$ 的零维开邻域。因为 $S_g$ 是诺特环，
可知 $S_g$ 是阿廷环，且
$D(g) = \Spec(S_g)$ 是有限离散集；见命题
\ref{algebra-proposition-dimension-zero-ring}。
所以 $\mathfrak q$ 是 $D(g)$ 的孤立点；把上面 (1) 与 (2) 的等价性
用于 $\mathfrak qS_g \subset S_g$，便得到
$S_{\mathfrak q} = (S_g)_{\mathfrak qS_g}$ 在 $k$ 上有限。
所以 (4) 蕴含 (2)。显然 (1) 蕴含 (4)。
故 (1)--(4) 全部等价。

\medskip\noindent
引理 \ref{algebra-lemma-dimension-closed-point-finite-type-field}
给出蕴含 (5) $\Rightarrow$ (4)。
蕴含 (4) $\Rightarrow$ (6) 来自引理
\ref{algebra-lemma-dimension-at-a-point-finite-type-field}。
蕴含 (6) $\Rightarrow$ (5) 来自引理
\ref{algebra-lemma-finite-residue-extension-closed}。
至此，我们知道 (1)--(6) 等价。

\medskip\noindent
% P01-E0581: the frozen source has the object before the finite verb (“The two statements ... we saw”).
在上面证明 (1)、(2)、(3) 等价的过程中，
我们已经看到引理末尾的两个陈述。
\end{proof}

\begin{lemma}
\label{algebra-lemma-isolated-point-fibre}
\begin{slogan}
纤维中孤立点的等价条件
\end{slogan}
设 $R \to S$ 是有限型环映射。
设 $\mathfrak q \subset S$ 是位于
$\mathfrak p \subset R$ 之上的素理想。令
$F = \Spec(S \otimes_R \kappa(\mathfrak p))$
为 $\Spec(S) \to \Spec(R)$ 的纤维；见注
\ref{algebra-remark-fundamental-diagram}。
% P01-E0582: the frozen source says “Denote qbar the point” rather than “Denote by qbar ...”.
以 $\overline{\mathfrak q} \in F$ 表示对应于
$\mathfrak q$ 的点。下列条件等价：
% P01-E0583: the frozen source omits punctuation after “equivalent”.
\begin{enumerate}
\item $\overline{\mathfrak q}$ 是 $F$ 的孤立点；
\item $S_{\mathfrak q}/\mathfrak pS_{\mathfrak q}$ 在
$\kappa(\mathfrak p)$ 上有限；
\item 存在 $g \in S$、$g \not \in \mathfrak q$，使得
$D(g)$ 中映到 $\mathfrak p$ 的素理想只有 $\mathfrak q$；
\item $\dim_{\overline{\mathfrak q}}(F) = 0$；
\item $\overline{\mathfrak q}$ 是 $F$ 的闭点，并且
$\dim(S_{\mathfrak q}/\mathfrak pS_{\mathfrak q}) = 0$；
\item 域扩张 $\kappa(\mathfrak q)/\kappa(\mathfrak p)$
是有限的，并且 $\dim(S_{\mathfrak q}/\mathfrak pS_{\mathfrak q}) = 0$。
\end{enumerate}
\end{lemma}

\begin{proof}
注意，$S_{\mathfrak q}/\mathfrak pS_{\mathfrak q} =
(S \otimes_R \kappa(\mathfrak p))_{\overline{\mathfrak q}}$。
此外，$S \otimes_R \kappa(\mathfrak p)$ 是
$\kappa(\mathfrak p)$ 上的有限型代数。
这些条件恰好对应于引理 \ref{algebra-lemma-isolated-point} 中
$\kappa(\mathfrak p)$-代数 $S \otimes_R \kappa(\mathfrak p)$
及其素理想 $\overline{\mathfrak q}$ 的各个条件，所以它们等价。
\end{proof}

\begin{definition}
\label{algebra-definition-quasi-finite}
设 $R \to S$ 是有限型环映射。
设 $\mathfrak q \subset S$ 是素理想。
\begin{enumerate}
\item 若引理 \ref{algebra-lemma-isolated-point-fibre} 中的等价条件成立，
就称 $R \to S$ \textit{在 $\mathfrak q$ 处拟有限}。
\item 若环映射 $A \to B$ 是有限型的，并且在 $B$ 的每个素理想处
都拟有限，就称它是\textit{拟有限}的。
\end{enumerate}
\end{definition}

\begin{lemma}
\label{algebra-lemma-quasi-finite-above-p}
设 $R \to S$ 是有限型环映射，
$\mathfrak p$ 是 $R$ 的素理想。
则下列条件等价：
\begin{enumerate}
\item $R \to S$ 在 $S$ 中位于 $\mathfrak p$ 之上的所有素理想处拟有限；
\item $S \otimes_R \kappa(\mathfrak p)$ 是有限
$\kappa(\mathfrak p)$-代数；
\item $\Spec(S \otimes_R \kappa(\mathfrak p))$ 是有限集。
\end{enumerate}
\end{lemma}

\begin{proof}
条件 (1) 表示 $\Spec(S \otimes_R \kappa(\mathfrak p))$
上的拓扑是离散的（因为每一点都是开点）。
所以 (1)、(2)、(3) 的等价性来自引理
\ref{algebra-lemma-finite-type-algebra-finite-nr-primes}。
\end{proof}

\begin{lemma}
\label{algebra-lemma-quasi-finite}
设 $R \to S$ 是有限型环映射。则 $R \to S$ 拟有限，当且仅当
对所有素理想 $\mathfrak p \subset R$，环
$S \otimes_R \kappa(\mathfrak p)$ 在 $\kappa(\mathfrak p)$ 上有限。
\end{lemma}

\begin{proof}
这立即由更一般的引理 \ref{algebra-lemma-quasi-finite-above-p} 得出。
\end{proof}

\begin{lemma}
\label{algebra-lemma-quasi-finite-local}
设 $R \to S$ 是有限型环映射。设 $\mathfrak q \subset S$
是位于 $\mathfrak p \subset R$ 之上的素理想。设
$f \in R$、$f \not \in \mathfrak p$，且 $g \in S$、$g \not \in \mathfrak q$。
则 $R \to S$ 在 $\mathfrak q$ 处拟有限，当且仅当
$R_f \to S_{fg}$ 在 $\mathfrak qS_{fg}$ 处拟有限。
\end{lemma}

\begin{proof}
$\Spec(S_{fg}) \to \Spec(R_f)$ 的纤维同胚于
$\Spec(S) \to \Spec(R)$ 的纤维的一个开子集。
因此，本引理由引理 \ref{algebra-lemma-isolated-point-fibre}
等价条件中的 (1) 得出。
\end{proof}

\begin{lemma}
\label{algebra-lemma-four-rings}
设
$$
\xymatrix{
S \ar[r] & S' & &
\mathfrak q \ar@{-}[r] & \mathfrak q' \\
R \ar[u] \ar[r] &  R' \ar[u] & &
\mathfrak p \ar@{-}[r] \ar@{-}[u] & \mathfrak p' \ar@{-}[u]
}
$$
是带有所示素理想的环的交换图。
% P01-E0584: the frozen source omits “is” in “Assume R to S of finite type”.
假设 $R \to S$ 是有限型的，并且
$S \otimes_R R' \to S'$ 是满射。
若 $R \to S$ 在 $\mathfrak q$ 处拟有限，则
$R' \to S'$ 在 $\mathfrak q'$ 处拟有限。
\end{lemma}

\begin{proof}
写成 $S \otimes_R \kappa(\mathfrak p) = S_1 \times S_2$，
其中 $S_1$ 在 $\kappa(\mathfrak p)$ 上有限，且按照引理
\ref{algebra-lemma-isolated-point}，$\mathfrak q$ 对应于 $S_1$ 的一点。
这个乘积分解在任意 $S \otimes_R \kappa(\mathfrak p)$-代数上
诱导相应的乘积分解。特别地，我们得到
$S' \otimes_{R'} \kappa(\mathfrak p') = S'_1 \times S'_2$。
因为 $S \otimes_R R' \to S'$ 是满射，典范映射
$(S \otimes_R \kappa(\mathfrak p)) \otimes_{\kappa(\mathfrak p)}
\kappa(\mathfrak p') \to S' \otimes_{R'} \kappa(\mathfrak p')$
是满射，因而
$S_i \otimes_{\kappa(\mathfrak p)} \kappa(\mathfrak p') \to S'_i$
是满射。所以 $S'_1$ 在 $\kappa(\mathfrak p')$ 上有限。
映射 $S' \otimes_{R'} \kappa(\mathfrak p') \to
\kappa(\mathfrak q')$ 通过 $S_1'$ 分解
（即它在因子 $S_2'$ 上为零），
因为映射 $S \otimes_R \kappa(\mathfrak p) \to
\kappa(\mathfrak q)$ 通过 $S_1$ 分解
（即它在因子 $S_2$ 上为零）。所以在纤维的无交并分解
$\Spec(S' \otimes_{R'} \kappa(\mathfrak p'))
= \Spec(S_1') \amalg \Spec(S_2')$ 中，$\mathfrak q'$ 对应于
$\Spec(S_1')$ 的一点；见引理 \ref{algebra-lemma-spec-product}。
% P01-E0585: the frozen source omits the article in “it is Artinian ring”.
由于 $S_1'$ 在域上有限，它是阿廷环，
从而 $\Spec(S_1')$ 是有限离散集。
（见命题 \ref{algebra-proposition-dimension-zero-ring}。）
因此 $\mathfrak q'$ 在其纤维中是孤立的，如所需。
\end{proof}

\begin{lemma}
\label{algebra-lemma-quasi-finite-composition}
拟有限环映射的复合是拟有限的。
\end{lemma}

\begin{proof}
假设 $A \to B$ 与 $B \to C$ 是拟有限环映射。由引理
\ref{algebra-lemma-compose-finite-type}，$A \to C$ 是有限型的。
设 $\mathfrak r \subset C$ 是 $C$ 的素理想，位于
$\mathfrak q \subset B$ 和 $\mathfrak p \subset A$ 之上。
% P01-E0586: the frozen “Since ..., then ...” construction has a redundant then.
由于 $A \to B$ 与 $B \to C$ 分别在 $\mathfrak q$ 和
$\mathfrak r$ 处拟有限，存在 $b \in B$ 与 $c \in C$，使得
$\mathfrak q$ 是 $D(b)$ 中映到 $\mathfrak p$ 的唯一素理想，且类似地，
$\mathfrak r$ 是 $D(c)$ 中映到 $\mathfrak q$ 的唯一素理想。
若 $c' \in C$ 是 $b \in B$ 的像，则 $\mathfrak r$ 是
$D(cc')$ 中映到 $\mathfrak p$ 的唯一素理想。
所以 $A \to C$ 在 $\mathfrak r$ 处拟有限。
\end{proof}

\begin{lemma}
\label{algebra-lemma-quasi-finite-base-change}
设 $R \to S$ 是有限型环映射。
设 $R \to R'$ 是任意环映射。置 $S' = R' \otimes_R S$。
\begin{enumerate}
\item 集合
$\{\mathfrak q' \mid R' \to S' \text{ quasi-finite at }\mathfrak q'\}$
是 $\Spec(S)$ 中相应集合在典范映射
$\Spec(S') \to \Spec(S)$ 下的逆像。
\item 若 $\Spec(R') \to \Spec(R)$ 是满射，则
$R \to S$ 拟有限当且仅当 $R' \to S'$ 拟有限。
\item 拟有限环映射的任意基变换都是拟有限的。
\end{enumerate}
\end{lemma}

\begin{proof}
设 $\mathfrak p' \subset R'$ 是位于 $\mathfrak p \subset R$
之上的素理想。则纤维环 $S' \otimes_{R'} \kappa(\mathfrak p')$
是纤维环 $S \otimes_R \kappa(\mathfrak p)$
沿域扩张 $\kappa(\mathfrak p) \to \kappa(\mathfrak p')$ 的基变换。
因此，第一个断言来自维数在域扩张下的不变性
（引理 \ref{algebra-lemma-dimension-at-a-point-preserved-field-extension}）
以及引理 \ref{algebra-lemma-isolated-point}。
拟有限映射在基变换下的稳定性，由此和有限型性质在基变换下的稳定性得出。
第二个断言成立，是因为该假设蕴含：给定素理想
$\mathfrak q \subset S$，总能找到位于它之上的素理想
$\mathfrak q' \subset S'$。
\end{proof}

\begin{lemma}
\label{algebra-lemma-quasi-finite-permanence}
设 $A \to B$ 与 $B \to C$ 是环同态，且 $A \to C$
是有限型的。设 $\mathfrak r$ 是 $C$ 的素理想，位于
$\mathfrak q \subset B$ 和 $\mathfrak p \subset A$ 之上。
若 $A \to C$ 在 $\mathfrak r$ 处拟有限，则
$B \to C$ 在 $\mathfrak r$ 处拟有限。
\end{lemma}

\begin{proof}
注意，$B \to C$ 是有限型的
（引理 \ref{algebra-lemma-compose-finite-type}），所以此陈述有意义。
我们使用引理 \ref{algebra-lemma-isolated-point-fibre} 的刻画 (3)。
若 $A \to C$ 在 $\mathfrak r$ 处拟有限，则存在某个
$c \in C$，使得
$$
\{\mathfrak r' \subset C \text{ 位于下列对象之上： }\mathfrak p\} \cap D(c) =
\{\mathfrak{r}\}.
$$
由于位于 $\mathfrak q$ 之上的素理想
$\mathfrak r' \subset C$ 构成位于 $\mathfrak p$ 之上的素理想
$\mathfrak r' \subset C$ 的一个子集，所以 $B \to C$ 在
$\mathfrak r$ 处拟有限。
\end{proof}

\begin{lemma}
\label{algebra-lemma-generically-finite}
设 $R \to S$ 是有限型环映射。
设 $\mathfrak p \subset R$ 是极小素理想。
假设 $S$ 中位于 $\mathfrak p$ 之上的素理想至多有限多个。
则存在 $g \in R$、$g \not \in \mathfrak p$，使得环映射
$R_g \to S_g$ 是有限的。
\end{lemma}

\begin{proof}
设 $x_1, \ldots, x_n$ 是 $S$ 在 $R$ 上的一组生成元。
由引理 \ref{algebra-lemma-quasi-finite-above-p}，该假设表示
$S \otimes_R \kappa(\mathfrak p) = S_{\mathfrak p}/\mathfrak pS_{\mathfrak p}$
是有限 $\kappa(\mathfrak p)$-代数。
因此可以找到首一多项式
$P_i \in R_{\mathfrak p}[X]$，使得 $P_i(x_i)$ 在
$S_{\mathfrak p}/\mathfrak pS_{\mathfrak p}$ 中的像为零。
由于 $\mathfrak p$ 是极小素理想，
$\mathfrak pR_{\mathfrak p}$ 是局部幂零理想；见引理
\ref{algebra-lemma-minimal-prime-reduced-ring}。
所以由引理 \ref{algebra-lemma-locally-nilpotent}，
$\mathfrak pS_{\mathfrak p}$ 是局部幂零理想。
% P01-E0587: the frozen source twice uses undefined P where the indexed polynomial is P_i.
因此存在 $e_i \geq 1$，使得 $P(x_i)^{e_i} = 0$
在 $S_{\mathfrak p}$ 中成立。取 $g_1 \in R$、
$g_1 \not \in \mathfrak p$，使 $P_i$ 对所有 $i$ 都有系数属于
$R[1/g_1]$。再取 $g_2 \in R$、$g_2 \not \in \mathfrak p$，
使得 $P(x_i)^{e_i} = 0$ 在 $S_{g_1g_2}$ 中成立。
置 $g = g_1g_2$ 即得结论。
\end{proof}

\section{Zariski 主定理}
\label{algebra-section-Zariski}

\noindent
本节的目标是证明 Zariski 主定理的代数版本。
在 Stacks project 后续关于概形及概形态射的理论发展中，
这个定理将成为许多结果的基础。

\medskip\noindent
设 $R \to S$ 是有限型环映射。
本节的目标是证明，$\Spec(S)$ 中映射拟有限的点所组成的集合
是\textit{开集}（定理 \ref{algebra-theorem-main-theorem}）。
事实上，将会看到存在有限环映射 $R \to S'$，使得在某种意义下，
$S/R$ 的拟有限点集在 $\Spec(S')$ 中是开的
（但我们不会在代数章中证明这一点，因为这里不发展概形语言；
当 $R \to S$ 拟有限时，见引理
\ref{algebra-lemma-quasi-finite-open-integral-closure}）。
这些陈述的证明颇为棘手，我们通过一长串关于环的整扩张与有限扩张的引理
来完成证明。这些材料可见 \cite{Henselian} 和 \cite{Peskine}。
Thierry Coquand 的笔记对我们也很有帮助。

\begin{lemma}
\label{algebra-lemma-make-integral-trivial}
设 $\varphi : R \to S$ 是环映射。
假设 $t \in S$ 满足关系
$\varphi(a_0) + \varphi(a_1)t + \ldots + \varphi(a_n) t^n = 0$。
则 $\varphi(a_n)t$ 在 $R$ 上是整的。
\end{lemma}

\begin{proof}
事实上，将等式
$\varphi(a_0) + \varphi(a_1)t + \ldots + \varphi(a_n) t^n = 0$
乘以 $\varphi(a_n)^{n-1}$，并将其写成
$\varphi(a_0 a_n^{n-1}) +
\varphi(a_1 a_n^{n-2}) (\varphi(a_n)t) +
\ldots +
(\varphi(a_n) t)^n = 0$。
\end{proof}

\noindent
下面的引理在某种意义下是本节的关键引理。

\begin{lemma}
\label{algebra-lemma-make-integral-trick}
设 $R$ 是环。设 $\varphi : R[x] \to S$ 是环映射。
设 $t \in S$。
假设 (a) $t$ 在 $R[x]$ 上是整的，且
(b) 存在首一多项式 $p \in R[x]$，使得
$t \varphi(p) \in \Im(\varphi)$。则存在 $q \in R[x]$，
使得 $t - \varphi(q)$ 在 $R$ 上是整的。
\end{lemma}

\begin{proof}
对某个 $r \in R[x]$，写成 $t \varphi(p) = \varphi(r)$。
利用欧几里得除法，写成 $r = qp + r'$，其中
$q, r' \in R[x]$ 且 $\deg(r') < \deg(p)$。
可以用 $t - \varphi(q)$ 代替 $t$；它在 $R[x]$ 上仍是整的，
并且这样便有 $t \varphi(p) = \varphi(r')$。
在环 $S_t$ 中，可将它写成
$\varphi(p) - (1/t) \varphi(r') = 0$。
这说明 $\varphi(x)$ 给出局部化 $S_t$ 的一个元素，
它在 $\varphi(R)[1/t] \subset S_t$ 上是整的。
另一方面，$t$ 在子环 $\varphi(R)[\varphi(x)] \subset S$ 上是整的。
合起来，由引理 \ref{algebra-lemma-integral-transitive}，可知 $t$ 在子环
$\varphi(R)[1/t] \subset S_t$ 上是整的。换言之，在 $S_t$ 中
存在形如
$$
t^d + \sum\nolimits_{i < d}
\left(\sum\nolimits_{j = 0, \ldots, n_i} \varphi(r_{i, j})/t^j\right) t^i = 0
$$
的等式，其中 $r_{i, j} \in R$。这表示，对某个充分大的 $N$，
$S$ 中有
$t^{d + N} +
\sum_{i < d} \sum_{j = 0, \ldots, n_i} \varphi(r_{i, j}) t^{i + N - j} = 0$。
换言之，$t$ 在 $R$ 上是整的。
\end{proof}

\begin{lemma}
\label{algebra-lemma-combine-lemmas}
设 $R$ 是环。设 $\varphi : R[x] \to S$ 是环映射。
设 $t \in S$。假设 $t$ 在 $R[x]$ 上是整的。
设 $p \in R[x]$、$p = a_0 + a_1x + \ldots +
a_k x^k$，并且 $t \varphi(p) \in \Im(\varphi)$。
则存在 $q \in R[x]$ 和 $n \geq 0$，使得
$\varphi(a_k)^n t - \varphi(q)$ 在 $R$ 上是整的。
\end{lemma}

\begin{proof}
设 $R'$ 和 $S'$ 分别是 $R$ 和 $S$ 在元素 $a_k$ 处的局部化。
设 $\varphi' : R'[x] \to S'$ 是 $\varphi$ 的局部化。
设 $t' \in S'$ 是 $t$ 的像。置 $p' = p/a_k \in R'[x]$。
由于 $t \varphi(p) \in \Im(\varphi)$，有
$t' \varphi'(p') \in \Im(\varphi')$。
因为 $p'$ 是首一的，由引理 \ref{algebra-lemma-make-integral-trick}，
存在 $q' \in R'[x]$，使得 $t' - \varphi'(q')$ 在 $R'$ 上是整的。
可以选取 $n \geq 0$ 和元素 $q \in R[x]$，使得
$a_k^n q'$ 是 $q$ 的像。
于是 $\varphi(a_k)^n t - \varphi(q)$ 是 $S$ 的一个元素，
它在 $S'$ 中的像在 $R'$ 上是整的。
由引理 \ref{algebra-lemma-integral-closure-localize}，
存在 $m \geq 0$，使得
$\varphi(a_k)^m(\varphi(a_k)^n t - \varphi(q))$ 在 $R$ 上是整的。
因此 $\varphi(a_k)^{m + n}t - \varphi(a_k^m q)$
在 $R$ 上是整的，如所需。
\end{proof}

\begin{situation}
\label{algebra-situation-one-transcendental-element}
设 $R$ 是环。
设 $\varphi : R[x] \to S$ 是有限的。令
$$
J = \{ g \in S \mid gS \subset \Im(\varphi)\}
$$
为 $\varphi$ 的“导子理想”。
% P01-E0588: the frozen source omits “is” in the integrally-closed assumption.
假设 $\varphi(R) \subset S$ 在 $S$ 中是整闭的。
\end{situation}

\begin{lemma}
\label{algebra-lemma-leading-coefficient-in-J}
% P01-E0589: the frozen source makes “In Situation ...” a sentence fragment.
在情形 \ref{algebra-situation-one-transcendental-element} 中，
假设 $u \in S$、$a_0, \ldots, a_k \in R$，且
$u \varphi(a_0 + a_1x + \ldots + a_k x^k) \in J$。
则存在 $m \geq 0$，使得 $u \varphi(a_k)^m \in J$。
\end{lemma}

\begin{proof}
假设 $S$ 作为 $R[x]$-模由 $t_1, \ldots, t_n$ 生成。
在这种情形下，
$J = \{ g \in S \mid gt_i \in \Im(\varphi)\text{ 对所有 }i\}$。
注意，每个元素 $u t_i$ 都在 $R[x]$ 上整；见引理
\ref{algebra-lemma-finite-is-integral}。
有 $\varphi(a_0 + a_1x + \ldots + a_k x^k) u t_i \in
\Im(\varphi)$。由引理 \ref{algebra-lemma-combine-lemmas}，
对每个 $i$，存在整数 $n_i$ 和元素 $q_i \in R[x]$，使得
$\varphi(a_k^{n_i}) u t_i - \varphi(q_i)$ 在 $R$ 上是整的。
由假设，此元素属于 $\varphi(R)$，所以
$\varphi(a_k^{n_i}) u t_i \in \Im(\varphi)$。
因此 $m = \max\{n_1, \ldots, n_n\}$ 可用。
\end{proof}

\begin{lemma}
\label{algebra-lemma-all-coefficients-in-J}
在情形 \ref{algebra-situation-one-transcendental-element} 中，
假设 $u \in S$、$a_0, \ldots, a_k \in R$，且
$u \varphi(a_0 + a_1x + \ldots + a_k x^k) \in \sqrt{J}$。
则对所有 $i$，有 $u \varphi(a_i) \in \sqrt{J}$。
\end{lemma}

\begin{proof}
在本引理的假设下，对某个 $n \geq 1$，有
$u^n \varphi(a_0 + a_1x + \ldots + a_k x^k)^n \in J$。
由引理 \ref{algebra-lemma-leading-coefficient-in-J}，可知对某个
$m \geq 1$，有 $u^n \varphi(a_k^{nm}) \in J$。
所以 $u \varphi(a_k) \in \sqrt{J}$，从而
$u \varphi(a_0 + a_1x + \ldots + a_k x^k) - u \varphi(a_k x^k) =
u \varphi(a_0 + a_1x + \ldots + a_{k-1} x^{k-1}) \in \sqrt{J}$。
对 $k$ 作归纳即得结论。
\end{proof}

\noindent
这个引理启发了如下定义。

\begin{definition}
\label{algebra-definition-strongly-transcendental}
给定环的包含 $R \subset S$ 和元素 $x \in S$，
如果每当
$u(a_0 + a_1 x + \ldots + a_k x^k) = 0$，其中
$u \in S$ 且 $a_i \in R$，就有对所有 $i$ 都成立的
$ua_i = 0$，那么称 $x$ \textit{在 $R$ 上强超越}。
\end{definition}

\noindent
注意，若 $S$ 是整环，则这等价于说，$x$ 作为 $S$ 的分式域的元素，
在 $R$ 的分式域上超越。

\begin{lemma}
\label{algebra-lemma-reduced-strongly-transcendental-minimal-prime}
假设 $R \subset S$ 是约化环的包含，并且
$x \in S$ 在 $R$ 上强超越。
设 $\mathfrak q \subset S$ 是极小素理想，并令
$\mathfrak p = R \cap \mathfrak q$。
则 $x$ 在 $S/\mathfrak q$ 中的像，在子环
$R/\mathfrak p$ 上强超越。
\end{lemma}

\begin{proof}
假设 $u(a_0 + a_1x + \ldots + a_k x^k) \in \mathfrak q$。
由引理 \ref{algebra-lemma-minimal-prime-reduced-ring}，
局部环 $S_{\mathfrak q}$ 是域，所以
$u(a_0 + a_1x + \ldots + a_k x^k) $ 在 $S_{\mathfrak q}$ 中为零。
因此对某个 $u' \in S$、$u' \not\in \mathfrak q$，有
$uu'(a_0 + a_1x + \ldots + a_k x^k) = 0$。
因为 $x$ 在 $R$ 上强超越，所以对所有 $i$ 都有
$uu'a_i = 0$。这又蕴含 $ua_i \in \mathfrak q$。
\end{proof}

\begin{lemma}
\label{algebra-lemma-domains-transcendental-not-quasi-finite}
假设 $R\subset S$ 是整环的包含，并设 $x \in S$。
假设 $x$ 在 $R$ 上（强）超越，并且 $S$ 在 $R[x]$ 上有限。
则 $R \to S$ 在 $S$ 的任何素理想处都不拟有限。
\end{lemma}

\begin{proof}
先假设 $R$ 正规；见定义 \ref{algebra-definition-ring-normal}。
由引理 \ref{algebra-lemma-polynomial-ring-normal}，$R[x]$ 正规。
取素理想 $\mathfrak q \subset S$，并置
$\mathfrak p = R \cap \mathfrak q$。
假设扩张 $\kappa(\mathfrak p)
\subset \kappa(\mathfrak q)$ 是有限的。
若 $R \to S$ 在 $\mathfrak q$ 处拟有限，就会是这种情形。
令 $\mathfrak r = R[x] \cap \mathfrak q$。
由于 $\kappa(\mathfrak p)
\subset \kappa(\mathfrak r) \subset \kappa(\mathfrak q)$，
扩张 $\kappa(\mathfrak p)
\subset \kappa(\mathfrak r)$ 也是有限的。
所以包含 $\mathfrak r \supset \mathfrak p R[x]$ 是严格的。
由 $R[x] \subset S$ 的下降性质（见命题
\ref{algebra-proposition-going-down-normal-integral}），
可找到位于素理想 $\mathfrak pR[x]$ 之上的素理想
$\mathfrak q' \subset \mathfrak q$。
因此，纤维 $\Spec(S \otimes_R \kappa(\mathfrak p))$
包含一个不等于 $\mathfrak q$ 的点，即 $\mathfrak q'$；
它的闭包包含 $\mathfrak q$，故 $\mathfrak q$ 在其纤维中不是孤立点。

\medskip\noindent
若 $R$ 不正规，设 $R \subset R' \subset K$，其中 $R'$ 是
$R$ 在其分式域 $K$ 中的整闭包。设 $S \subset S' \subset L$，
其中 $S'$ 是在 $S$ 的分式域 $L$ 中由 $R'$ 与 $S$ 生成的子环。
注意，根据构造，映射 $S \otimes_R R' \to S'$ 是满射。
这蕴含 $R'[x] \subset S'$ 是有限的。此外，映射
$S \subset S'$ 在 $\Spec$ 上诱导满射；见引理
\ref{algebra-lemma-integral-overring-surjective}。
由引理 \ref{algebra-lemma-four-rings} 和刚讨论的正规情形即得结论。
\end{proof}

\begin{lemma}
\label{algebra-lemma-reduced-strongly-transcendental-not-quasi-finite}
假设 $R \subset S$ 是约化环的包含。
% P01-E0590: the frozen source says “Assume x in S be strongly transcendental”.
假设 $x \in S$ 在 $R$ 上强超越，并且 $S$ 在 $R[x]$ 上有限。
则 $R \to S$ 在 $S$ 的任何素理想处都不拟有限。
\end{lemma}

\begin{proof}
任取素理想 $\mathfrak q \subset S$。
选取极小素理想 $\mathfrak q' \subset \mathfrak q$。
根据引理
\ref{algebra-lemma-reduced-strongly-transcendental-minimal-prime} 和
\ref{algebra-lemma-domains-transcendental-not-quasi-finite}，扩张
$R/(R \cap \mathfrak q') \subset
S/\mathfrak q'$ 在对应于 $\mathfrak q$ 的素理想处不拟有限。
由引理 \ref{algebra-lemma-four-rings}，扩张 $R \to S$
在 $\mathfrak q$ 处不拟有限。
\end{proof}

\begin{lemma}
\label{algebra-lemma-quasi-finite-monogenic}
设 $R$ 是环。设 $S = R[x]/I$。
设 $\mathfrak q \subset S$ 是素理想。
假设 $R \to S$ 在 $\mathfrak q$ 处拟有限。
设 $S' \subset S$ 是 $R$ 在 $S$ 中的整闭包。
则存在元素 $g \in S'$、$g \not\in \mathfrak q$，使得
$S'_g \cong S_g$。
\end{lemma}

\begin{proof}
设 $\mathfrak p$ 是 $\mathfrak q$ 在 $\Spec(R)$ 中的像。
存在 $f \in I$、$f = a_nx^n + \ldots + a_0$，使得
对某个 $i$ 有 $a_i \not \in \mathfrak p$。事实上，否则纤维环
$S \otimes_R \kappa(\mathfrak p)$ 将是 $\kappa(\mathfrak p)[x]$，
而映射在位于 $\mathfrak p$ 之上的任何素理想处都不会拟有限。
由此可知，存在关系
$b_m x^m + \ldots + b_0 = 0$，其中 $b_j \in S'$、
$j = 0, \ldots, m$，并且对某个 $j$ 有
$b_j \not \in \mathfrak q \cap S'$。
我们对 $m$ 作归纳来证明本引理。
% P01-E0591: the frozen source has the duplicated copula “The base case is m = 0 is vacuous”.
基始情形 $m = 0$ 是空的（因为陈述 $b_0 = 0$ 与
$b_0 \not \in \mathfrak q$ 相矛盾）。

\medskip\noindent
考虑 $b_m \not \in \mathfrak q$ 的情形。此时 $x$ 在
$S'_{b_m}$ 上是整的；事实上，由引理
\ref{algebra-lemma-make-integral-trivial}，有 $b_mx \in S'$。
因此，单射 $S'_{b_m} \to S_{b_m}$ 也是满射，
即它是所需的同构。

\medskip\noindent
考虑 $b_m \in \mathfrak q$ 的情形。此时由引理
\ref{algebra-lemma-make-integral-trivial}，有 $b_mx \in S'$。
置 $b'_{m - 1} = b_mx + b_{m - 1}$。则
% P01-E0592: the frozen source omits punctuation after this displayed relation.
$$
b'_{m - 1}x^{m - 1} + b_{m - 2}x^{m - 2} + \ldots + b_0 = 0
$$
由于 $b'_{m - 1}$ 模 $S' \cap \mathfrak q$ 同余于 $b_{m - 1}$，
仍有 $b'_{m - 1}, b_{m - 2}, \ldots, b_0$ 中的某个元素
不属于 $S' \cap \mathfrak q$。
所以对 $m$ 应用归纳即可。
\end{proof}

\begin{theorem}[Zariski 主定理]
\label{algebra-theorem-main-theorem}
设 $R$ 是环。
% P01-E0593: the frozen source calls the ring map itself “a finite type R-algebra”.
设 $R \to S$ 是有限型 $R$-代数。
设 $S' \subset S$ 是 $R$ 在 $S$ 中的整闭包。
设 $\mathfrak q \subset S$ 是 $S$ 的素理想。
若 $R \to S$ 在 $\mathfrak q$ 处拟有限，则存在
$g \in S'$、$g \not \in \mathfrak q$，使得 $S'_g \cong S_g$。
\end{theorem}

\begin{proof}
存在有限多个元素 $x_1, \ldots, x_n \in S$，使得 $S$
在由 $x_1, \ldots, x_n$ 生成的 $R$-子代数上有限
% P01-E0594: the frozen source writes the standard compound “subalgebra” as “sub algebra”.
（例如可取 $S$ 在 $R$ 上的生成元）。
% P01-E0595: the frozen proof calls this theorem “the proposition”.
我们对这样的最小数目 $n$ 作归纳来证明该命题。

\medskip\noindent
情形 $n = 0$ 是平凡的，因为此时 $S' = S$；
见引理 \ref{algebra-lemma-finite-is-integral}。

\medskip\noindent
考虑 $n = 1$ 的情形。可以把 $R$ 替换为其在 $S$ 中的整闭包
（引理 \ref{algebra-lemma-quasi-finite-permanence} 保证 $R \to S$
在 $\mathfrak q$ 处仍拟有限）。所以可以假设
$R \subset S$ 在 $S$ 中整闭，换言之，$R = S'$。
考虑映射 $\varphi : R[x] \to S$、$x \mapsto x_1$。
（下面将看到 $\varphi$ 不是单射。）
由假设，$\varphi$ 是有限的。因此我们处于情形
\ref{algebra-situation-one-transcendental-element}。
设 $J \subset S$ 为情形
\ref{algebra-situation-one-transcendental-element} 中定义的“导子理想”。
考虑图
$$
\xymatrix{
R[x] \ar[r] & S \ar[r] & S/\sqrt{J} & R/(R \cap \sqrt{J})[x] \ar[l]
\\
& R \ar[lu] \ar[r] \ar[u] & R/(R \cap \sqrt{J}) \ar[u] \ar[ru] &
}
$$
根据引理 \ref{algebra-lemma-all-coefficients-in-J}，
$x$ 在商 $S/\sqrt{J}$ 中的像在
$R/ (R \cap \sqrt{J})$ 上强超越。
因此，由引理 \ref{algebra-lemma-reduced-strongly-transcendental-not-quasi-finite}，
环映射 $R/ (R \cap \sqrt{J}) \to S/\sqrt{J}$
在 $S/\sqrt{J}$ 的任何素理想处都不拟有限。
由引理 \ref{algebra-lemma-four-rings}，可知 $\mathfrak q$
不属于 $V(J) \subset \Spec(S)$。
所以存在元素 $s \in J$、$s \not\in \mathfrak q$。
由 $J$ 的定义，可以对某个多项式 $f \in R[x]$ 写成
$s = \varphi(f)$。令 $I = \Ker(\varphi : R[x] \to S)$。
由于 $\varphi(f) \in J$，有
$(R[x]/I)_f \cong S_{\varphi(f)}$。又因为
$s \not \in \mathfrak q$，有 $f \not \in \varphi^{-1}(\mathfrak q)$。
因此 $\varphi^{-1}(\mathfrak q)/I$ 是 $R[x]/I$ 的一个素理想，
$R \to R[x]/I$ 在该素理想处拟有限；见引理
\ref{algebra-lemma-quasi-finite-local}。注意，$R$ 在 $R[x]/I$ 中整闭，
因为 $R$ 在 $S$ 中整闭。由引理
\ref{algebra-lemma-quasi-finite-monogenic}，存在元素
$h \in R$、$h \not \in R \cap \mathfrak q$，使得
$R_h \cong (R[x]/I)_h$。所以
$(R[x]/I)_{fh} = S_{\varphi(fh)}$ 同构于 $R$ 的某个主局部化
$R_{h'}$，其中 $h' \in R$、$h' \not \in \mathfrak q$。

\medskip\noindent
考虑 $n > 1$ 的情形。令子环 $R' \subset S$ 为
$R[x_1, \ldots, x_{n-1}]$ 在 $S$ 中的整闭包。
由引理 \ref{algebra-lemma-quasi-finite-permanence}，扩张
$S/R'$ 在 $\mathfrak q$ 处拟有限。还要注意，$S$ 在
$R'[x_n]$ 上有限。由上面的 $n = 1$ 情形，存在
$g' \in R'$、$g' \not \in \mathfrak q$，使得
$(R')_{g'} \cong S_{g'}$。此时不能对 $R \to R'$ 应用归纳，
因为 $R'$ 在 $R$ 上可能不是有限型的。
由于 $S$ 在 $R$ 上有限生成，特别可知
$(R')_{g'}$ 在 $R$ 上有限生成。设元素
$g'$ 以及 $y_1/(g')^{n_1}, \ldots, y_N/(g')^{n_N}$，
其中 $y_i \in R'$，生成 $(R')_{g'}$ 作为 $R$-代数。
令 $R''$ 为 $R'$ 的 $R$-子代数，由
$x_1, \ldots, x_{n-1}, y_1, \ldots, y_N, g'$ 生成。
它具有性质 $(R'')_{g'} \cong S_{g'}$。
% P01-E0596: the frozen source uses two sentence fragments for surjectivity and injectivity.
满性来自 $y_i$ 的选取方式；单性来自
$R'' \subset R'$ 以及局部化的正合性。
注意，由 $R'$ 的选取及引理 \ref{algebra-lemma-characterize-integral}，
$R''$ 在 $R[x_1, \ldots, x_{n-1}]$ 上有限。
令 $\mathfrak q'' = R'' \cap \mathfrak q$。由于
$(R'')_{\mathfrak q''} = S_{\mathfrak q}$，由引理
\ref{algebra-lemma-isolated-point-fibre}，$R \to R''$
在 $\mathfrak q''$ 处拟有限。
对 $R \to R''$、$\mathfrak q''$ 以及
$x_1, \ldots, x_{n-1} \in R''$ 应用归纳假设，便得到
在 $R$ 上整的子环 $R''' \subset R''$ 和元素
$g'' \in R'''$、$g'' \not \in \mathfrak q''$，使得
$(R''')_{g''} \cong (R'')_{g''}$。
把 $g'$ 在 $(R'')_{g''}$ 中的像写成
$g'''/(g'')^n$，其中 $g''' \in  R'''$。
置 $g = g''g''' \in R'''$。则显然
$g \not\in \mathfrak q$，且 $(R''')_g \cong S_g$。
根据构造，$R''' \subset S'$，所以也有 $S'_g \cong S_g$，如所需。
\end{proof}

\begin{lemma}
\label{algebra-lemma-quasi-finite-open}
设 $R \to S$ 是有限型环映射。
$\Spec(S)$ 中 $S/R$ 拟有限的点 $\mathfrak q$ 所组成的集合
在 $\Spec(S)$ 中是开的。
\end{lemma}

\begin{proof}
设 $\mathfrak q \subset S$ 是环映射拟有限的一个点。
由定理 \ref{algebra-theorem-main-theorem}，存在整环映射
$R \to S'$、$S' \subset S$ 以及元素
$g \in S'$、$g\not \in \mathfrak q$，使得
$S'_g \cong S_g$。由于 $S$、从而 $S_g$ 在 $R$ 上是有限型的，
可以找到 $S'$ 的有限多个元素 $y_1, \ldots, y_N$，使得
$S''_g \cong S_g$，其中 $S'' \subset S'$ 是由
$g, y_1, \ldots, y_N$ 生成的 $R$-子代数。
由于 $S''$ 在 $R$ 上有限
（见引理 \ref{algebra-lemma-characterize-integral}），由引理
\ref{algebra-lemma-quasi-finite}，$S''$ 在 $R$ 上拟有限。
容易看出，这蕴含 $S''_g$ 在 $R$ 上拟有限；
例如，因为在一个素理想处拟有限这一性质仅依赖于该素理想处的局部环。
所以 $S_g$ 在 $R$ 上拟有限。同理，这蕴含
$R \to S$ 在 $S$ 中属于 $D(g)$ 的每个素理想处拟有限。
\end{proof}

\begin{lemma}
\label{algebra-lemma-quasi-finite-open-integral-closure}
设 $R \to S$ 是有限型环映射。
假设 $S$ 在 $R$ 上拟有限。
设 $S' \subset S$ 是 $R$ 在 $S$ 中的整闭包。则
\begin{enumerate}
\item $\Spec(S) \to \Spec(S')$ 是到一个开子集上的同胚；
\item 若 $g \in S'$ 且 $D(g)$ 包含于该映射的像中，
则 $S'_g \cong S_g$；
\item 存在有限 $R$-代数 $S'' \subset S'$，使得环映射
$S'' \to S$ 满足 (1) 和 (2)。
\end{enumerate}
\end{lemma}

\begin{proof}
因为 $S/R$ 拟有限，可以对 $\Spec(S)$ 的每个点
$\mathfrak q$ 应用定理 \ref{algebra-theorem-main-theorem}。
由于 $\Spec(S)$ 拟紧（见引理 \ref{algebra-lemma-quasi-compact}），
可以选取有限多个 $g_i \in S'$、$i = 1, \ldots, n$，使得
$S'_{g_i} = S_{g_i}$，并且 $g_1, \ldots, g_n$ 在 $S$
中生成单位理想（换言之，$\Spec(S)$ 中与
$g_1, \ldots, g_n$ 相应的标准开集覆盖整个 $\Spec(S)$）。

\medskip\noindent
假设 $D(g) \subset \Spec(S')$ 包含于该映射的像中。
则 $D(g) \subset \bigcup D(g_i)$。换言之，
$g_1, \ldots, g_n$ 在 $S'_g$ 中生成单位理想。
注意，根据 $g_i$ 的选取，有 $S'_{gg_i} \cong S_{gg_i}$。
所以由引理 \ref{algebra-lemma-cover}，$S'_g \cong S_g$。

\medskip\noindent
我们构造 (3) 中的有限代数 $S'' \subset S'$。
为此，注意每个 $S'_{g_i} \cong S_{g_i}$ 都是有限型
$R$-代数。对每个 $i$，选取一些元素 $y_{ij} \in S'$，
使得每个 $S'_{g_i}$ 作为 $R$-代数由 $1/g_i$ 与各个
$y_{ij}$ 生成。然后令 $S''$ 等于 $S'$ 的 $R$-子代数，
由所有 $g_i$ 和所有 $y_{ij}$ 生成。细节略去。
\end{proof}

\section{Zariski 主定理的应用}
\label{algebra-section-apply-main-theorem}

\noindent
下面是一个直接应用：在 $1$ 维半局部环的拟有限映射中，
有限映射恰好是引理 \ref{algebra-lemma-finite-extension-dim-1}
的公式总取等号的那些映射。

\begin{lemma}
\label{algebra-lemma-quasi-finite-extension-dim-1}
设 $A \subset B$ 是整环的扩张。假设
\begin{enumerate}
\item $A$ 是维数为 $1$ 的诺特局部环；
\item $A \to B$ 是有限型的；
\item 所诱导的分式域扩张 $L/K$ 是有限的。
\end{enumerate}
则 $B$ 是半局部环。
设 $x \in \mathfrak m_A$、$x \not = 0$。
设 $\mathfrak m_i$（$i = 1, \ldots, n$）是 $B$ 的极大理想。
则
$$
[L : K]\text{ord}_A(x)
\geq
\sum\nolimits_i
[\kappa(\mathfrak m_i) : \kappa(\mathfrak m_A)]
\text{ord}_{B_{\mathfrak m_i}}(x)
$$
其中 $\text{ord}$ 如定义 \ref{algebra-definition-ord} 所定义。
等号成立当且仅当 $A \to B$ 是有限的。
\end{lemma}

\begin{proof}
由引理 \ref{algebra-lemma-finite-in-codim-1}，环 $B$ 是半局部的。
设 $B'$ 为 $A$ 在 $B$ 中的整闭包。由引理
\ref{algebra-lemma-quasi-finite-open-integral-closure}，
可以找到有限 $A$-子代数 $C \subset B'$，使得置
$\mathfrak n_i = C \cap \mathfrak m_i$ 后，有
$C_{\mathfrak n_i} \cong B_{\mathfrak m_i}$，并且各素理想
$\mathfrak n_1, \ldots, \mathfrak n_n$ 两两不同。
由引理 \ref{algebra-lemma-finite-in-codim-1}，环 $C$ 是半局部的。
设 $\mathfrak p_j$（$j = 1, \ldots, m$）是 $C$ 的其他极大理想
（即“缺失点”）。由引理 \ref{algebra-lemma-finite-extension-dim-1}，有
$$
\text{ord}_A(x^{[L : K]}) =
\sum\nolimits_i
[\kappa(\mathfrak n_i) : \kappa(\mathfrak m_A)]
\text{ord}_{C_{\mathfrak n_i}}(x)
+
\sum\nolimits_j
[\kappa(\mathfrak p_j) : \kappa(\mathfrak m_A)]
\text{ord}_{C_{\mathfrak p_j}}(x)
$$
故不等式成立。在等号成立的情形下，可知 $m = 0$
（没有“缺失点”）。因此 $C \subset B$ 是半局部环的包含，
它在极大理想上诱导双射，并在所有极大理想处的局部化上诱导同构。
所以若 $b \in B$，则
$I = \{x \in C \mid xb \in C\}$ 是 $C$ 的一个理想，
不包含于 $C$ 的任何极大理想，因而 $I = C$，从而 $b \in C$。
所以 $B = C$，且 $B$ 在 $A$ 上有限。
\end{proof}

\noindent
下面是 Zariski 主定理对局部环之间局部同态结构的一个更标准的应用。

\begin{lemma}
\label{algebra-lemma-essentially-finite-type-fibre-dim-zero}
设 $(R, \mathfrak m_R) \to (S, \mathfrak m_S)$ 是局部环的局部同态。
假设
\begin{enumerate}
\item $R \to S$ 本质上是有限型的；
\item $\kappa(\mathfrak m_R) \subset \kappa(\mathfrak m_S)$ 是有限扩张；
\item $\dim(S/\mathfrak m_RS) = 0$。
\end{enumerate}
则 $S$ 是某个有限 $R$-代数的局部化。
\end{lemma}

\begin{proof}
设 $S'$ 是有限型 $R$-代数，使得对 $S'$ 的某个素理想
$\mathfrak q'$ 有 $S = S'_{\mathfrak q'}$。
由定义 \ref{algebra-definition-quasi-finite}，可知
$R \to S'$ 在 $\mathfrak q'$ 处拟有限。
用某个 $g' \in S'$、$g' \not \in \mathfrak q'$ 的
$S'_{g'}$ 代替 $S'$ 后，可以假设 $R \to S'$ 拟有限；
见引理 \ref{algebra-lemma-quasi-finite-open}。
于是由引理 \ref{algebra-lemma-quasi-finite-open-integral-closure}，
存在有限 $R$-代数 $S''$ 以及元素
$g' \in S'$、$g' \not \in \mathfrak q'$ 和 $g'' \in S''$，
使得 $S'_{g'} \cong S''_{g''}$ 为 $R$-代数。
这就证明了本引理。
\end{proof}

\begin{lemma}
\label{algebra-lemma-completion-at-quasi-finite-prime}
设 $R \to S$ 是环映射，$\mathfrak q$ 是 $S$ 的素理想，
位于 $R$ 的 $\mathfrak p$ 之上。若
\begin{enumerate}
\item $R$ 是诺特环；
\item $R \to S$ 是有限型的；
\item $R \to S$ 在 $\mathfrak q$ 处拟有限，
\end{enumerate}
则对某个 $R_\mathfrak p^\wedge$-代数 $B$，有
$R_\mathfrak p^\wedge \otimes_R S = S_\mathfrak q^\wedge \times B$。
\end{lemma}

\begin{proof}
存在有限 $R$-代数 $S' \subset S$ 和元素
$g \in S'$、$g \not \in \mathfrak q' = S' \cap \mathfrak q$，
使得 $S'_g = S_g$，特别地，
$S'_{\mathfrak q'} = S_\mathfrak q$；见引理
\ref{algebra-lemma-quasi-finite-open-integral-closure}。
由引理 \ref{algebra-lemma-completion-finite-extension}，有
$$
R_\mathfrak p^\wedge \otimes_R S' = (S'_{\mathfrak q'})^\wedge \times B'
$$
注意，在这个乘积分解下，$g$ 映到一对 $(u, b')$，其中
$u \in (S'_{\mathfrak q'})^\wedge$ 是单位，
因为 $g \not \in \mathfrak q'$。
$R_\mathfrak p^\wedge \otimes_R S'$ 的乘积分解诱导乘积分解
% P01-E0597: the frozen source omits punctuation after this displayed product decomposition.
$$
R_\mathfrak p^\wedge \otimes_R S = A \times B
$$
由于 $S'_g = S_g$，还有
$(R_\mathfrak p^\wedge \otimes_R S')_g = (R_\mathfrak p^\wedge \otimes_R S)_g$；
又因为 $g \mapsto (u, b')$ 且 $u$ 是单位，所以
$(S'_{\mathfrak q'})^\wedge = A$。
同构 $S'_{\mathfrak q'} = S_\mathfrak q$ 在完备化上确定一个同构，
故这也说明 $A = S_\mathfrak q^\wedge$。
证明至此完成，不过还应作一个合理性检验：诱导映射
$\phi : R_\mathfrak p^\wedge \otimes_R S \to A = S_\mathfrak q^\wedge$
确实是自然映射。对于形如 $x \otimes 1$、
$x \in R_\mathfrak p^\wedge$ 的元素，这是显然的，
因为自然映射 $R_\mathfrak p^\wedge \to S_\mathfrak q^\wedge$
通过 $(S'_{\mathfrak q'})^\wedge$ 分解。
对于形如 $1 \otimes y$、$y \in S$ 的元素，可以这样论证：
对某个 $n \geq 1$，元素 $g^ny$ 是某个 $y' \in S'$ 的像。
% P01-E0598: the frozen source uses “image ... by the composition” rather than “under”.
因此，$\phi(1 \otimes g^ny)$ 是 $y'$ 在复合映射
$S' \to (S'_{\mathfrak q'})^\wedge \to S_\mathfrak q^\wedge$ 下的像，
它等于 $g^ny$ 在映射 $S \to S_\mathfrak q^\wedge$ 下的像。
由于 $g$ 映到单位，这也蕴含 $\phi(1 \otimes y)$ 具有正确的值，
即 $y$ 在 $S \to S_\mathfrak q^\wedge$ 下的像。
\end{proof}

\section{纤维的维数}
\label{algebra-section-dimension-fibres}

\noindent
我们利用 Zariski 主定理研究纤维维数的行为。
回顾一下，我们已在拓扑学定义 \ref{topology-definition-Krull} 中
定义了拓扑空间 $X$ 在点 $x$ 处的维数 $\dim_x(X)$。

\begin{definition}
\label{algebra-definition-relative-dimension}
假设 $R \to S$ 是有限型的，并设
$\mathfrak q \subset S$ 是位于 $R$ 的素理想
$\mathfrak p$ 之上的素理想。
我们把 $\Spec(S \otimes_R \kappa(\mathfrak p))$
在对应于 $\mathfrak q$ 的点处的维数，定义为
$S/R$ \textit{在 $\mathfrak q$ 处的相对维数}，记作
$\dim_{\mathfrak q}(S/R)$。令 $\dim(S/R)$ 为所有
$\dim_{\mathfrak q}(S/R)$ 的上确界，其中 $\mathfrak q$ 取遍所有素理想。
它称为 $S/R$ 的\textit{相对维数}。
\end{definition}

\noindent
特别地，$R \to S$ 在 $\mathfrak q$ 处拟有限，当且仅当
$\dim_{\mathfrak q}(S/R) = 0$。下面的引理大致是
Zariski 主定理的改述。

\begin{lemma}
\label{algebra-lemma-quasi-finite-over-polynomial-algebra}
设 $R \to S$ 是有限型环映射。
设 $\mathfrak q \subset S$ 是素理想。
假设 $\dim_{\mathfrak q}(S/R) = n$。
则存在 $g \in S$、$g \not\in \mathfrak q$，使得
$S_g$ 在多项式代数 $R[t_1, \ldots, t_n]$ 上拟有限。
\end{lemma}

\begin{proof}
环 $\overline{S} = S \otimes_R \kappa(\mathfrak p)$ 是
$\kappa(\mathfrak p)$ 上的有限型代数。
设 $\overline{\mathfrak q}$ 为 $\overline{S}$ 中对应于
$\mathfrak q$ 的素理想。
根据拓扑空间在一点处的维数的定义，存在开集
$U \subset \Spec(\overline{S})$，满足
% P01-E0599: the frozen source drops mathfrak from the barred prime symbol.
$\overline{q} \in U$ 且 $\dim(U) = n$。
由于 $\Spec(\overline{S})$ 上的拓扑由 $\Spec(S)$ 上的拓扑诱导
（见注 \ref{algebra-remark-fundamental-diagram}），可以找到
$g \in S$、$g \not \in \mathfrak q$，其像为
$\overline{g} \in \overline{S}$，并且
$D(\overline{g}) \subset U$。
因此，用 $S_g$ 代替 $S$ 后，有 $\dim(\overline{S}) = n$。

\medskip\noindent
接着，选取 $S$ 作为 $R$-代数的一组生成元
$x_1, \ldots, x_N$。由引理 \ref{algebra-lemma-Noether-normalization}，
在由 $x_1, \ldots, x_N$ 生成的 $S$ 的 $\mathbf{Z}$-子代数中，
存在元素 $y_1, \ldots, y_n$，使得映射
$R[t_1, \ldots, t_n] \to S$、$t_i \mapsto y_i$ 具有如下性质：
% P01-E0600: the frozen polynomial algebra omits the comma after t_1.
$\kappa(\mathfrak p)[t_1\ldots, t_n] \to \overline{S}$ 是有限的。
特别地，$S$ 在 $R[t_1, \ldots, t_n]$ 上于
$\mathfrak q$ 处拟有限。因此，由引理 \ref{algebra-lemma-quasi-finite-open}，
可以用某个 $g\in S$、$g \not \in \mathfrak q$ 的 $S_g$ 代替 $S$，
使得 $R[t_1, \ldots, t_n] \to S$ 拟有限。
\end{proof}

\begin{lemma}
\label{algebra-lemma-refined-quasi-finite-over-polynomial-algebra}
设 $R \to S$ 是环映射。设 $\mathfrak q \subset S$
是位于 $R$ 的素理想 $\mathfrak p$ 之上的素理想。
假设
\begin{enumerate}
\item $R \to S$ 是有限型的；
\item $\dim_{\mathfrak q}(S/R) = n$；
\item $\text{trdeg}_{\kappa(\mathfrak p)}\kappa(\mathfrak q) = r$。
\end{enumerate}
则存在 $f \in R$、$f \not \in \mathfrak p$，
$g \in S$、$g \not\in \mathfrak q$，以及拟有限环映射
$$
\varphi : R_f[x_1, \ldots, x_n] \longrightarrow S_g
$$
使得
% P01-E0601: the frozen lemma statement omits punctuation after this displayed equality.
$\varphi^{-1}(\mathfrak qS_g) =
(\mathfrak p, x_{r + 1}, \ldots, x_n)R_f[x_1, \ldots, x_n]$。
\end{lemma}

\begin{proof}
用一个主局部化代替 $S$ 后，可以假设存在拟有限环映射
$\varphi : R[t_1, \ldots, t_n] \to S$；见引理
\ref{algebra-lemma-quasi-finite-over-polynomial-algebra}。
置 $\mathfrak q' = \varphi^{-1}(\mathfrak q)$。
设 $\overline{\mathfrak q}' \subset \kappa(\mathfrak p)[t_1, \ldots, t_n]$
为对应于 $\mathfrak q'$ 的素理想。
由引理 \ref{algebra-lemma-refined-Noether-normalization}，存在有限环映射
$\kappa(\mathfrak p)[x_1, \ldots, x_n] \to
\kappa(\mathfrak p)[t_1, \ldots, t_n]$，使得
$\overline{\mathfrak q}'$ 的逆像为
$(x_{r + 1}, \ldots, x_n)$。
设 $\overline{h}_i \in \kappa(\mathfrak p)[t_1, \ldots, t_n]$
为 $x_i$ 的像。可以找到元素
$f \in R$、$f \not \in \mathfrak p$ 以及
$h_i \in R_f[t_1, \ldots, t_n]$，它们在
$\kappa(\mathfrak p)[t_1, \ldots, t_n]$ 中映到 $\overline{h}_i$。
于是环映射
$$
R_f[x_1, \ldots, x_n] \longrightarrow R_f[t_1, \ldots, t_n]
$$
与 $\kappa(\mathfrak p)$ 张量后成为有限映射。
特别地，$R_f[t_1, \ldots, t_n]$ 在
$R_f[x_1, \ldots, x_n]$ 上于素理想
$\mathfrak q'R_f[t_1, \ldots, t_n]$ 处拟有限。
因此，由引理 \ref{algebra-lemma-quasi-finite-open}，存在
$g \in R_f[t_1, \ldots, t_n]$、
$g \not \in \mathfrak q'R_f[t_1, \ldots, t_n]$，使得
$R_f[x_1, \ldots, x_n] \to R_f[t_1, \ldots, t_n, 1/g]$
拟有限。所以复合
$$
R_f[x_1, \ldots, x_n] \longrightarrow
R_f[t_1, \ldots, t_n, 1/g] \longrightarrow S_{\varphi(g)}
$$
拟有限，结论得证。
\end{proof}

\begin{lemma}
\label{algebra-lemma-dimension-inequality-quasi-finite}
设 $R \to S$ 是有限型环映射。
设 $\mathfrak q \subset S$ 是位于 $\mathfrak p \subset R$
之上的素理想。若 $R \to S$ 在 $\mathfrak q$ 处拟有限，则
$\dim(S_{\mathfrak q}) \leq \dim(R_{\mathfrak p})$。
\end{lemma}

\begin{proof}
若 $R_{\mathfrak p}$ 是诺特环
（从而 $S_{\mathfrak q}$ 也是诺特环，因为它在
$R_{\mathfrak p}$ 上本质上为有限型），
则这立即由引理 \ref{algebra-lemma-dimension-base-fibre-total} 和定义得出。
在一般情形下，设 $S'$ 是 $R_\mathfrak p$ 在 $S_\mathfrak p$
中的整闭包。由 Zariski 主定理 \ref{algebra-theorem-main-theorem}，
% P01-E0602: the frozen source reverses the direction by saying q' in S' lies over q in S.
对某个位于 $\mathfrak q$ 之上的 $\mathfrak q' \subset S'$，有
$S_{\mathfrak q} = S'_{\mathfrak q'}$。
由引理 \ref{algebra-lemma-integral-dim-up}，有
$\dim(S') \leq \dim(R_\mathfrak p)$，从而更有
$\dim(S_\mathfrak q) = \dim(S'_{\mathfrak q'}) \leq \dim(R_\mathfrak p)$。
\end{proof}

\begin{lemma}
\label{algebra-lemma-dimension-quasi-finite-over-polynomial-algebra}
\begin{slogan}
仿射 n 维空间的拟有限覆盖的维数至多为 n。
\end{slogan}
设 $k$ 是域。设 $S$ 是有限型 $k$-代数。
假设存在拟有限 $k$-代数映射
$k[t_1, \ldots, t_n] \subset S$。则 $\dim(S) \leq n$。
\end{lemma}

\begin{proof}
由引理 \ref{algebra-lemma-dim-affine-space}，
$k[t_1, \ldots, t_n]$ 的任意局部环的维数至多为 $n$。
所以结论由引理 \ref{algebra-lemma-dimension-inequality-quasi-finite} 得出。
\end{proof}

\begin{lemma}
\label{algebra-lemma-dimension-fibres-bounded-open-upstairs}
设 $R \to S$ 是有限型环映射。
设 $\mathfrak q \subset S$ 是素理想。
假设 $\dim_{\mathfrak q}(S/R) = n$。
则 $\mathfrak q$ 在 $\Spec(S)$ 中有一个开邻域 $V$，
使得对所有 $\mathfrak q' \in V$，都有
$\dim_{\mathfrak q'}(S/R) \leq n$。
\end{lemma}

\begin{proof}
由引理 \ref{algebra-lemma-quasi-finite-over-polynomial-algebra}，
可知可以假设 $S$ 在多项式代数
$R[t_1, \ldots, t_n]$ 上拟有限。
考察纤维，便归结为引理
\ref{algebra-lemma-dimension-quasi-finite-over-polynomial-algebra}。
\end{proof}

\noindent
换言之，本引理说明，纤维维数 $\leq n$ 的点所组成的集合
在 $\Spec(S)$ 中是开的。下一个引理说明，这个开集的形成与基变换交换。
若环映射是有限表示的，则此集合是拟紧开集（见下文）。

\begin{lemma}
\label{algebra-lemma-dimension-fibres-bounded-open-upstairs-base-change}
设 $R \to S$ 是有限型环映射。
设 $R \to R'$ 是任意环映射。
置 $S' = R' \otimes_R S$，并以
$f : \Spec(S') \to \Spec(S)$ 表示相应的谱上映射。
设 $n \geq 0$。
逆像
$f^{-1}(\{\mathfrak q \in \Spec(S) \mid
\dim_{\mathfrak q}(S/R) \leq n\})$
等于
$\{\mathfrak q' \in \Spec(S') \mid
\dim_{\mathfrak q'}(S'/R') \leq n\}$。
\end{lemma}

\begin{proof}
该条件用纤维环的维数表述，而这些纤维环是域上的有限型代数。
合用引理 \ref{algebra-lemma-dimension-at-a-point-preserved-field-extension}
即可得到本引理。
\end{proof}

\begin{lemma}
\label{algebra-lemma-dimension-fibres-bounded-quasi-compact-open-upstairs}
设 $R \to S$ 是有限表示环同态。
设 $n \geq 0$。集合
$$
V_n = \{\mathfrak q \in \Spec(S) \mid \dim_{\mathfrak q}(S/R) \leq n\}
$$
是 $\Spec(S)$ 的拟紧开子集。
\end{lemma}

\begin{proof}
由引理 \ref{algebra-lemma-dimension-fibres-bounded-open-upstairs}，它是开的。
% P01-E0603: the frozen proof reuses n both for the fibre-dimension bound and for the number of presentation variables.
设 $S = R[x_1, \ldots, x_n]/(f_1, \ldots, f_m)$ 是 $S$ 的一个表示。
设 $R_0$ 是 $R$ 的 $\mathbf{Z}$-子代数，由各多项式
$f_i$ 的系数生成。令
$S_0 = R_0[x_1, \ldots, x_n]/(f_1, \ldots, f_m)$。
则 $S = R \otimes_{R_0} S_0$。由引理
\ref{algebra-lemma-dimension-fibres-bounded-open-upstairs-base-change}，
$V_n$ 是开集 $V_{0, n}$ 在拟紧连续映射
$\Spec(S) \to \Spec(S_0)$ 下的逆像。
由于 $S_0$ 是诺特环，$V_{0, n}$ 是拟紧的。
\end{proof}

\begin{lemma}
\label{algebra-lemma-finite-type-domain-over-valuation-ring-dim-fibres}
设 $R$ 是剩余域为 $k$、分式域为 $K$ 的赋值环。
设 $S$ 是包含 $R$ 的整环，并且 $S$ 在 $R$ 上为有限型。
若 $S \otimes_R k$ 不是零环，则
$$
\dim(S \otimes_R k) = \dim(S \otimes_R K)
$$
% P01-E0604: the frozen source omits punctuation after the displayed equality.
事实上，$\Spec(S \otimes_R k)$ 是等维的。
\end{lemma}

\begin{proof}
只须证明，对 $S$ 中包含 $\mathfrak m_RS$ 的每个素理想
$\mathfrak q$，$\dim_{\mathfrak q}(S/k)$ 都等于
$\dim(S \otimes_R K)$。取这样的素理想。
由引理 \ref{algebra-lemma-dimension-fibres-bounded-open-upstairs}，
不等式 $\dim_{\mathfrak q}(S/k) \geq \dim(S \otimes_R K)$ 成立。
置 $n = \dim_{\mathfrak q}(S/k)$。
由引理 \ref{algebra-lemma-quasi-finite-over-polynomial-algebra}，
用某个 $g \in S$、$g \not \in \mathfrak q$ 的 $S_g$ 代替 $S$ 后，
存在拟有限环映射 $R[t_1, \ldots, t_n] \to S$。
若 $\dim(S \otimes_R K) < n$，则
$K[t_1, \ldots, t_n] \to S \otimes_R K$ 有非零核。
设 $f = \sum a_I t_1^{i_1}\ldots t_n^{i_n}$。
用 $f$ 的一个具有最小赋值的非零系数除 $f$，可以假设
$f\in R[t_1, \ldots, t_n]$，并且某个 $a_I$ 在 $k$ 中的像非零。
所以环映射 $k[t_1, \ldots, t_n] \to S \otimes_R k$
有非零核，这蕴含 $\dim(S \otimes_R k) < n$，矛盾。
\end{proof}

\section{有限表示的代数与模}
\label{algebra-section-finite-presentation}

\noindent
本节讨论一些标准结果，其中的关键特征是：所用假设包含有限型或有限表示条件。

\begin{lemma}
\label{algebra-lemma-finite-type-descends}
设 $R \to S$ 是环映射。
设 $R \to R'$ 是忠实平坦环映射。
置 $S' = R'\otimes_R S$。
则 $R \to S$ 为有限型的，当且仅当 $R' \to S'$ 为有限型的。
\end{lemma}

\begin{proof}
若 $R \to S$ 为有限型的，则 $R' \to S'$ 显然为有限型的。
假设 $R' \to S'$ 为有限型的。
设 $y_1, \ldots, y_m$ 在 $R'$ 上生成 $S'$。
对某些 $a_{ij} \in R'$ 和 $x_{ji} \in S$，写成
$y_j = \sum_i a_{ij} \otimes x_{ji}$。
% P01-E0605: the frozen proof switches the two indices from x_{ji} to x_{ij}.
设 $A \subset S$ 是由这些 $x_{ij}$ 生成的 $R$-子代数。
由平坦性，我们有 $A' := R' \otimes_R A \subset S'$，而且
按构造 $y_j \in A'$。因此 $A' = S'$。
由忠实平坦性，$A = S$。
\end{proof}

\begin{lemma}
\label{algebra-lemma-finite-presentation-descends}
设 $R \to S$ 是环映射。
设 $R \to R'$ 是忠实平坦环映射。
置 $S' = R'\otimes_R S$。
则 $R \to S$ 为有限表示的，当且仅当 $R' \to S'$ 为有限表示的。
\end{lemma}

\begin{proof}
若 $R \to S$ 为有限表示的，则 $R' \to S'$ 显然为有限表示的。
假设 $R' \to S'$ 为有限表示的。
由引理 \ref{algebra-lemma-finite-type-descends}，$R \to S$ 为有限型的。
写成 $S = R[x_1, \ldots, x_n]/I$。
由平坦性，$S' = R'[x_1, \ldots, x_n]/R'\otimes I$。
设 $g_1, \ldots, g_m$ 在 $R'[x_1, \ldots, x_n]$ 上生成
$R'\otimes I$。对某些 $a_{ij} \in R'$ 和 $f_{ji} \in I$，写成
$g_j = \sum_i a_{ij} \otimes f_{ji}$。
% P01-E0606: the frozen proof switches the two indices from f_{ji} to f_{ij}.
设 $J \subset I$ 是由这些 $f_{ij}$ 生成的理想。
由平坦性，$R' \otimes_R J \subset R'\otimes_R I$，而且两者都是
$R'[x_1, \ldots, x_n]$ 上的理想。
按构造 $g_j \in R' \otimes_R J$。因此
$R' \otimes_R J = R'\otimes_R I$。
由忠实平坦性，$J = I$。
\end{proof}

\begin{lemma}
\label{algebra-lemma-construct-fp-module}
设 $R$ 是环。
设 $I \subset R$ 是理想。
设 $S \subset R$ 是乘法子集。
置 $R' = S^{-1}(R/I) = S^{-1}R/S^{-1}I$。
\begin{enumerate}
\item 对任意有限 $R'$-模 $M'$，存在有限 $R$-模 $M$，使得
$S^{-1}(M/IM) \cong M'$。
\item 对任意有限表示 $R'$-模 $M'$，存在有限表示 $R$-模 $M$，
使得 $S^{-1}(M/IM) \cong M'$。
\end{enumerate}
\end{lemma}

\begin{proof}
证明 (1)。选取短正合列
$0 \to K' \to (R')^{\oplus n} \to M' \to 0$。
设 $K \subset R^{\oplus n}$ 是 $K'$ 在映射
$R^{\oplus n} \to (R')^{\oplus n}$ 下的逆像。
则 $M = R^{\oplus n}/K$ 符合要求。

\medskip\noindent
证明 (2)。
选取表示 $(R')^{\oplus m} \to (R')^{\oplus n} \to M' \to 0$。
假设第一个映射由矩阵 $A' = (a'_{ij})$ 给出，第二个映射由生成元
$x'_i \in M'$、$i = 1, \ldots, n$ 确定。由于 $R' = S^{-1}(R/I)$，
可以选取 $s \in S$ 以及一个系数在 $R$ 中的矩阵 $A = (a_{ij})$，使得
$a'_{ij} = a_{ij} / s \bmod S^{-1}I$。设 $M$ 是具有表示
$R^{\oplus m} \to R^{\oplus n} \to M \to 0$
的有限表示 $R$-模，其中第一个映射由矩阵 $A$ 给出，
第二个映射由生成元 $x_i \in M$、$i = 1, \ldots, n$ 确定。
则映射 $M \to M'$、$x_i \mapsto x'_i$ 诱导同构
$S^{-1}(M/IM) \cong M'$。
\end{proof}

\begin{lemma}
\label{algebra-lemma-construct-fp-module-from-localization}
设 $R$ 是环。
设 $S \subset R$ 是乘法子集。
设 $M$ 是 $R$-模。
\begin{enumerate}
\item 若 $S^{-1}M$ 是有限 $S^{-1}R$-模，则存在有限 $R$-模 $M'$ 以及映射
$M' \to M$，它诱导同构 $S^{-1}M' \to S^{-1}M$。
\item 若 $S^{-1}M$ 是有限表示 $S^{-1}R$-模，则存在有限表示
$R$-模 $M'$ 以及映射 $M' \to M$，它诱导同构
$S^{-1}M' \to S^{-1}M$。
\end{enumerate}
\end{lemma}

\begin{proof}
证明 (1)。设 $x_1, \ldots, x_n \in M$ 是作为 $S^{-1}R$-模生成
$S^{-1}M$ 的元。设 $M'$ 是由 $x_1, \ldots, x_n$ 生成的 $M$ 的
$R$-子模。

\medskip\noindent
证明 (2)。设 $x_1, \ldots, x_n \in M$ 是作为 $S^{-1}R$-模生成
$S^{-1}M$ 的元。设
$K = \Ker(R^{\oplus n} \to M)$，其中该映射由规则
$(a_1, \ldots, a_n) \mapsto \sum a_i x_i$ 给出。由引理 \ref{algebra-lemma-extension}，
$S^{-1}K$ 是有限 $S^{-1}R$-模。
由 (1)，可以找到有限子模 $K' \subset K$，使得
$S^{-1}K' = S^{-1}K$。取
$M' = \Coker(K' \to R^{\oplus n})$。
\end{proof}

\begin{lemma}
\label{algebra-lemma-construct-fp-module-from-stalk}
设 $R$ 是环。
设 $\mathfrak p \subset R$ 是素理想。
设 $M$ 是 $R$-模。
\begin{enumerate}
\item 若 $M_{\mathfrak p}$ 是有限 $R_{\mathfrak p}$-模，则存在有限 $R$-模 $M'$ 以及映射
$M' \to M$，它诱导同构 $M'_{\mathfrak p} \to M_{\mathfrak p}$。
\item 若 $M_{\mathfrak p}$ 是有限表示 $R_{\mathfrak p}$-模，则存在有限表示
$R$-模 $M'$ 以及映射 $M' \to M$，它诱导同构
$M'_{\mathfrak p} \to M_{\mathfrak p}$。
\end{enumerate}
\end{lemma}

\begin{proof}
% P01-E0607: the frozen proof omits terminal punctuation.
这是引理 \ref{algebra-lemma-construct-fp-module-from-localization} 的特殊情形。
\end{proof}

\begin{lemma}
\label{algebra-lemma-local-isomorphism}
设 $\varphi : R \to S$ 是环映射。设 $\mathfrak q \subset S$ 是位于
$\mathfrak p \subset R$ 之上的素理想。假设
\begin{enumerate}
\item $S$ 在 $R$ 上为有限表示的，
\item $\varphi$ 诱导同构 $R_\mathfrak p \cong S_\mathfrak q$。
\end{enumerate}
则存在 $f \in R$、$f \not \in \mathfrak p$ 以及一个 $R_f$-代数 $C$，使得
$S_f \cong R_f \times C$ 为 $R_f$-代数同构。
\end{lemma}

\begin{proof}
写成 $S = R[x_1, \ldots, x_n]/(g_1, \ldots, g_m)$。设 $a_i \in R_\mathfrak p$
是映到 $x_i$ 在 $S_\mathfrak q$ 中之像的元。
对某个 $f \in R$、$f \not \in \mathfrak p$，写成 $a_i = b_i/f$。
用 $R_f$ 代替 $R$，并用 $x_i - a_i$ 代替 $x_i$ 后，可以假设
$S = R[x_1, \ldots, x_n]/(g_1, \ldots, g_m)$，且 $x_i$ 在 $S_\mathfrak q$ 中映到零。
若 $c_j$ 表示 $g_j$ 的常数项，则 $c_j$ 在 $R_\mathfrak p$ 中映到零。
再一次代替 $R$ 后，可以假设这些 $g_j$ 的常数系数 $c_j$ 为零。
因而得到 $R$-代数映射 $S \to R$、$x_i \mapsto 0$，它的核为理想
$(x_1, \ldots, x_n)$。

\medskip\noindent
我们有同构 $R_\mathfrak p \to S_\mathfrak q \to R_\mathfrak p$，而映射
$S \to R$ 把 $x_i$ 送到零。因此有：
$$
S_\mathfrak q = R_\mathfrak p[x_1, \ldots, x_n]/(x_1, \ldots, x_n)
$$
从而有：
$$
S_\mathfrak q = S_\mathfrak p/(x_1, \ldots, x_n)S_\mathfrak p
$$
这意味着有限生成理想 $(x_1, \ldots, x_n)S_\mathfrak p$ 在
$S_\mathfrak p$ 中是纯的，见定义 \ref{algebra-definition-pure-ideal}。
因此，由引理 \ref{algebra-lemma-finitely-generated-pure-ideal}，
$(x_1, \ldots, x_n)S_\mathfrak p$ 由 $S_\mathfrak p$ 中的一个幂等元 $e$ 生成。
对某个 $f \in R$、$f \not \in \mathfrak p$，用 $R_f \to S_f$ 代替
$R \to S$ 后，可以找到一个映到 $e$ 的幂等元 $e' \in S$。
此时 $e'S$ 和 $(x_1, \ldots, x_n)S$ 是有限生成理想，它们在
$S_\mathfrak p$ 中变得相等。因此，对某个 $f \in R$、$f \not \in \mathfrak p$，
用 $R_f \to S_f$ 代替 $R \to S$ 后，可以假设
$e'S = (x_1, \ldots, x_n)S$。置 $C = e'S$ 即完成证明。
\end{proof}

\begin{lemma}
\label{algebra-lemma-isomorphic-local-rings}
设 $R$ 是环。
设 $S$、$S'$ 在 $R$ 上为有限表示的。
设 $\mathfrak q \subset S$ 和 $\mathfrak q' \subset S'$ 是素理想。
若 $S_{\mathfrak q} \cong S'_{\mathfrak q'}$ 为 $R$-代数同构，则存在
$g \in S$、$g \not \in \mathfrak q$ 和 $g' \in S'$、$g' \not \in \mathfrak q'$，
使得 $S_g \cong S'_{g'}$ 为 $R$-代数同构。
\end{lemma}

\begin{proof}
设 $\psi : S_{\mathfrak q} \to S'_{\mathfrak q'}$ 是引理假设中的同构。
写成 $S = R[x_1, \ldots, x_n]/(f_1, \ldots, f_r)$ 和
$S' = R[y_1, \ldots, y_m]/J$。
对每个 $i = 1, \ldots, n$，选取分式 $h_i/g_i$，其中
$h_i, g_i \in R[y_1, \ldots, y_m]$，且 $g_i \bmod J$ 不属于
$\mathfrak q'$，使它表示 $x_i$ 在 $\psi$ 下的像。
% P01-E0608: the frozen replacement sentence omits the second “replacing ... by ...” construction.
用 $S'_{g_1 \ldots g_n}$ 代替 $S'$，并换用
$R[y_1, \ldots, y_m, y_{m + 1}]$（把 $y_{m + 1}$ 映到 $1/(g_1\ldots g_n)$）后，
可以假设 $\psi(x_i)$ 是某个 $h_i \in R[y_1, \ldots, y_m]$ 的像。
考虑这些元
$f_j(h_1, \ldots, h_n) \in R[y_1, \ldots, y_m]$。
由于 $\psi$ 消去每个 $f_j$，存在
$g \in R[y_1, \ldots, y_m]$、$g \bmod J \not \in \mathfrak q'$，使得对每个
$j = 1, \ldots, r$ 都有 $g f_j(h_1, \ldots, h_n) \in J$。
像前面一样用 $S'_g$ 代替 $S'$，并换用
$R[y_1, \ldots, y_m, y_{m + 1}]$ 后，可以假设
$f_j(h_1, \ldots, h_n) \in J$。因而得到环映射
$S \to S'$、$x_i \mapsto h_i$，它在局部环上诱导 $\psi$。
由引理 \ref{algebra-lemma-compose-finite-type}，映射 $S \to S'$ 为有限表示的。
由引理 \ref{algebra-lemma-local-isomorphism}，可以假设 $S' = S \times C$。
% P01-E0609: localizing at the idempotent of the C factor yields C, not the required S factor.
因此，在对应于因子 $C$ 的幂等元处局部化 $S'$，即得结论。
\end{proof}

\begin{lemma}
\label{algebra-lemma-finite-type-mod-nilpotent}
设 $R$ 是环。设 $I \subset R$ 是幂零理想。
设 $S$ 是 $R$-代数，且 $R/I \to S/IS$ 为有限型的。
则 $R \to S$ 为有限型的。
\end{lemma}

\begin{proof}
选取 $s_1, \ldots, s_n \in S$，使它们在 $S/IS$ 中的像作为
$R/I$-代数生成 $S/IS$。由引理 \ref{algebra-lemma-NAK} 的 (11)，
$R$-代数映射 $R[x_1, \ldots, x_n] \to S$、$x_i \mapsto s_i$ 为满射，
结论随之得证。
\end{proof}

\begin{lemma}
\label{algebra-lemma-surjective-mod-locally-nilpotent}
设 $R$ 是环。设 $I \subset R$ 是局部幂零理想。
设 $S \to S'$ 是 $R$-代数映射，且 $S \to S'/IS'$ 为满射，
并且 $S'$ 在 $R$ 上为有限型的。则 $S \to S'$ 为满射。
\end{lemma}

\begin{proof}
对某个理想 $K$，写成 $S' = R[x_1, \ldots, x_m]/K$。
% P01-E0610: the frozen proof uses m variables, switches to an undefined n in three polynomial rings, and then returns to m.
由假设，存在 $g_j = x_j + \sum \delta_{j, J} x^J \in R[x_1, \ldots, x_n]$，
其中 $\delta_{j, J} \in I$，且 $g_j \bmod K \in \Im(S \to S')$。
因此，只须证明 $g_1, \ldots, g_m$ 生成
$R[x_1, \ldots, x_n]$。设 $R_0 \subset R$ 是 $R$ 的一个有限生成
$\mathbf{Z}$-子代数，它至少包含这些 $\delta_{j, J}$。
则 $R_0 \cap I$ 是幂零理想（由引理 \ref{algebra-lemma-Noetherian-power}）。
从而 $R_0[x_1, \ldots, x_n]$ 由 $g_1, \ldots, g_m$ 生成
（因为 $x_j \mapsto g_j$ 定义 $R_0[x_1, \ldots, x_m]$ 的一个自同构；
细节从略）。由于 $R$ 是这些子环 $R_0$ 的并，结论得证。
\end{proof}

\begin{lemma}
\label{algebra-lemma-isomorphism-modulo-ideal}
设 $R$ 是环。设 $I \subset R$ 是理想。设 $S \to S'$ 是
$R$-代数映射。设 $IS \subset \mathfrak q \subset S$ 是素理想。
假设
\begin{enumerate}
\item $S \to S'$ 为满射，
\item $S_\mathfrak q/IS_\mathfrak q \to S'_\mathfrak q/IS'_\mathfrak q$ 是同构，
\item $S$ 在 $R$ 上为有限型的，
% P01-E0611: the frozen fourth item omits the finite verb “is”.
\item $S'$ 在 $R$ 上为有限表示的，以及
\item $S'_\mathfrak q$ 在 $R$ 上平坦。
\end{enumerate}
则对某个 $g \in S$、$g \not \in \mathfrak q$，映射
$S_g \to S'_g$ 是同构。
\end{lemma}

\begin{proof}
设 $J = \Ker(S \to S')$。由引理 \ref{algebra-lemma-compose-finite-type}，
$J$ 是有限生成理想。由于 $S'_\mathfrak q$ 在 $R$ 上平坦，我们有
$J_\mathfrak q/IJ_\mathfrak q \subset S_\mathfrak q/IS_{\mathfrak q}$
（对 $0 \to J \to S \to S' \to 0$ 应用引理 \ref{algebra-lemma-flat-tor-zero}）。
由假设 (2)，$J_\mathfrak q/IJ_\mathfrak q$ 为零。
由 Nakayama 引理（引理 \ref{algebra-lemma-NAK}），存在
$g \in S$、$g \not \in \mathfrak q$，使得 $J_g = 0$。
因此 $S_g \cong S'_g$，正如所求。
\end{proof}

\begin{lemma}
\label{algebra-lemma-isomorphism-modulo-locally-nilpotent}
设 $R$ 是环。设 $I \subset R$ 是理想。设 $S \to S'$ 是
$R$-代数映射。假设
\begin{enumerate}
\item $I$ 是局部幂零的，
\item $S/IS \to S'/IS'$ 是同构，
\item $S$ 在 $R$ 上为有限型的，
\item $S'$ 在 $R$ 上为有限表示的，以及
\item $S'$ 在 $R$ 上平坦。
\end{enumerate}
则 $S \to S'$ 是同构。
\end{lemma}

\begin{proof}
由引理 \ref{algebra-lemma-surjective-mod-locally-nilpotent}，映射
$S \to S'$ 是满射。由于 $I$ 是局部幂零的，理想
$IS$ 和 $IS'$ 也是局部幂零的（引理 \ref{algebra-lemma-locally-nilpotent}）。
因此，$S$ 的每个素理想 $\mathfrak q$ 都包含 $IS$，并且（显然地）
$S_\mathfrak q/IS_\mathfrak q \cong S'_\mathfrak q/IS'_\mathfrak q$。
于是引理 \ref{algebra-lemma-isomorphism-modulo-ideal} 可用，并且对每个素理想
$\mathfrak q \subset S$，映射 $S_\mathfrak q \to S'_\mathfrak q$ 都是同构。
% P01-E0612: the frozen conclusion has the awkward word order “injective for example by”.
例如由引理 \ref{algebra-lemma-characterize-zero-local}，可知 $S \to S'$ 为单射。
\end{proof}

\section{余极限与有限表示映射}
\label{algebra-section-colimits-flat}

\noindent
% P01-E0613: the frozen introduction says “prove result” without an article or plural ending.
本节证明一些预备引理，它们最终将帮助我们使用绝对 Noether 化约
来证明结果。在范畴篇第 \ref{categories-section-directed-colimits} 节中，
我们一般地讨论了滤过余极限。下面是这一十分一般的概念的一个例子。

\begin{lemma}
\label{algebra-lemma-ring-colimit-fp-category}
设 $R \to A$ 是环映射。考虑由所有 $R$-代数映射图式
$A' \to A$ 组成的范畴 $\mathcal{I}$，其中 $A'$ 在 $R$ 上为有限表示的。
则 $\mathcal{I}$ 是滤过的，而且这些 $A'$ 在 $\mathcal{I}$ 上的余极限同构于 $A$。
\end{lemma}

\begin{proof}
范畴\footnote{为避免集合论上的困难，我们只考虑这样的 $A' \to A$：
$A'$ 是 $R[x_1, x_2, x_3, \ldots]$ 的商。}
$\mathcal{I}$ 非空，因为 $R \to R$ 是它的对象。
考虑 $\mathcal{I}$ 的一对对象 $A' \to A$、$A'' \to A$。
则 $A' \otimes_R A'' \to A$ 属于 $\mathcal{I}$（使用引理
\ref{algebra-lemma-compose-finite-type} 和 \ref{algebra-lemma-base-change-finiteness}）。
环映射 $A' \to A' \otimes_R A''$ 和 $A'' \to A' \otimes_R A''$
定义了 $\mathcal{I}$ 中的态射，从而证明了滤过范畴的第二个定义性质；
见范畴篇定义 \ref{categories-definition-directed}。
最后，假设有 $\mathcal{I}$ 中的两个态射
$\sigma, \tau : A' \to A''$。若 $x_1, \ldots, x_r \in A'$ 作为
$R$-代数生成 $A'$，则可以考虑
$A''' = A''/(\sigma(x_i) - \tau(x_i))$。
这是有限表示 $R$-代数，而给定的 $R$-代数映射
$A'' \to A$ 通过满射 $\nu : A'' \to A'''$ 分解。
因此 $\nu$ 是 $\mathcal{I}$ 中使 $\sigma$ 和 $\tau$ 相等的态射，正如所求。

\medskip\noindent
% P01-E0614: the frozen proof calls the filtered index category cofiltered.
我们的指标范畴是余滤过的，这一事实意味着，可以在集合范畴中计算
$B = \colim_{A' \to A} A'$ 的值（略去一些细节；可与范畴篇第
\ref{categories-section-directed-colimits} 节的讨论比较）。
要看出 $B \to A$ 为满射，对每个 $a \in A$，可以使用
$R[x] \to A$、$x \mapsto a$ 来看出 $a$ 在 $B \to A$ 的像中。
反之，若 $b \in B$ 在 $A$ 中映到零，则可以找到
$\mathcal{I}$ 中的 $A' \to A$ 以及 $a' \in A'$，使得该元映到 $b$。
那么 $A'/(a') \to A$ 也属于 $\mathcal{I}$，而映射 $A' \to B$
分解为 $A' \to A'/(a') \to B$，这说明 $b = 0$，正如所求。
\end{proof}

\noindent
通常，在预序集上思考余极限更容易。
% P01-E0615: the frozen sentence “Let ... a preordered set” omits “be”.
设 $(\Lambda, \geq)$ 是预序集。
$\Lambda$ 上的一个环系统由以下数据给出：对每个
$\lambda \in \Lambda$ 给出一个环 $R_\lambda$，并且当
$\lambda \leq \mu$ 时给出一个态射 $R_\lambda \to R_\mu$。
这些态射必须满足如下规则：对所有 $\lambda \leq \mu \leq \nu$，
$R_\lambda \to R_\mu \to R_\nu$ 等于映射 $R_\lambda \to R_\nu$。
见范畴篇第 \ref{categories-section-posets-limits} 节。
% P01-E0616: the frozen text switches from Lambda to an undefined I in the directedness assumption.
我们将经常假设 $(I, \leq)$ 是 {\it 有向的}，这意味着 $\Lambda$ 非空，
并且对给定的 $\lambda, \mu \in \Lambda$，存在 $\nu \in \Lambda$，使得
$\lambda \leq \nu$ 和 $\mu \leq \nu$。
回忆，在这种情形下，余极限 $\colim_\lambda R_\lambda$
有时称为“直极限”（但我们不使用这一术语）。

\medskip\noindent
注意，范畴篇引理 \ref{categories-lemma-directed-category-system} 告诉我们，
滤过指标范畴上的余极限与有向集上的余极限是同一回事。

\begin{lemma}
\label{algebra-lemma-ring-colimit-fp}
设 $R \to A$ 是环映射。存在一个由有限表示 $R$-代数 $A_\lambda$
组成的有向系统，使得 $A = \colim_\lambda A_\lambda$。
若 $A$ 在 $R$ 上为有限型的，则可以把它安排成：
$A_\lambda$ 的该系统中所有转移映射都是满射。
\end{lemma}

\begin{proof}
第一种证明是：这由引理 \ref{algebra-lemma-ring-colimit-fp-category} 和
范畴篇引理 \ref{categories-lemma-directed-category-system} 得出。

\medskip\noindent
第二种证明。
与引理 \ref{algebra-lemma-module-colimit-fp} 的证明比较。
考虑任意有限子集 $S \subset A$，以及 $S$ 中元之间的任意有限个
多项式关系所组成的集合 $E$。
于是每个 $s \in S$ 对应于 $x_s \in A$，而每个 $e \in E$
由一个多项式 $f_e \in R[X_s; s\in S]$ 组成，它满足 $f_e(x_s) = 0$。
设 $A_{S, E} = R[X_s; s\in S]/(f_e; e\in E)$，
它是有限表示 $R$-代数。
存在典范映射 $A_{S, E} \to A$。
若 $S \subset S'$，并且 $E$ 的各元经映射
$R[X_s; s \in S] \to R[X_s; s\in S']$ 对应于 $E'$ 的一个子集，
则有一个明显的映射 $A_{S, E} \to A_{S', E'}$，它与到 $A$ 的映射可换。
% P01-E0617: the frozen phrase “setting Lambda equal the set” omits “to”.
因此，令 $\Lambda$ 等于按上述包含关系排序的对 $(S, E)$ 所组成的集合，
就得到一个有向偏序集。
显然，这一有向系统的余极限是 $A$。

\medskip\noindent
对最后一个断言，假设 $A = R[x_1, \ldots, x_n]/I$。
在这种情形下，考虑子集 $\Lambda' \subset \Lambda$，它由上述那些满足
$S = \{x_1, \ldots, x_n\}$ 的系统 $(S, E)$ 组成。
容易看出，仍然有 $A = \colim_{\lambda' \in \Lambda'} A_{\lambda'}$。
此外，转移映射显然都是满射。
\end{proof}

\noindent
事实证明，可以如下刻画有限表示环映射。
在某种意义下，这说有限表示代数是 $R$-代数范畴中的“紧”对象。

\begin{lemma}
\label{algebra-lemma-characterize-finite-presentation}
设 $\varphi : R \to S$ 是环映射。下列条件等价：
\begin{enumerate}
\item $\varphi$ 为有限表示的；
\item 对每个 $R$-代数的有向系统 $A_\lambda$，映射
$$
\colim_\lambda \Hom_R(S, A_\lambda) \longrightarrow
\Hom_R(S, \colim_\lambda A_\lambda)
$$
是双射；以及
\item 对每个 $R$-代数的有向系统 $A_\lambda$，映射
$$
\colim_\lambda \Hom_R(S, A_\lambda) \longrightarrow
\Hom_R(S, \colim_\lambda A_\lambda)
$$
是满射。
\end{enumerate}
\end{lemma}

\begin{proof}
假设 (1)，并写成 $S = R[x_1, \ldots, x_n] / (f_1, \ldots, f_m)$。
设 $A = \colim A_\lambda$。注意，$R$-代数同态 $S \to A$ 或
$S \to A_\lambda$ 由 $x_1, \ldots, x_n$ 的像确定。
因此，$\colim_\lambda \Hom_R(S, A_\lambda) \to \Hom_R(S, A)$ 显然为单射。
要看出它为满射，设 $\chi : S \to A$ 是 $R$-代数同态。
则每个 $x_i$ 都映到某个 $A_{\lambda_i}$ 的像中的一个元。
可以选取 $\mu \geq \lambda_i$、$i = 1, \ldots, n$，并假设对
$i = 1, \ldots, n$，$\chi(x_i)$ 是 $y_i \in A_\mu$ 的像。
考虑 $z_j = f_j(y_1, \ldots, y_n) \in A_\mu$。
因为 $\chi$ 是同态，$z_j$ 在 $A = \colim_\lambda A_\lambda$ 中的像为零。
因此存在 $\mu_j \geq \mu$，使得 $z_j$ 在 $A_{\mu_j}$ 中映到零。
选取 $\nu \geq \mu_j$、$j = 1, \ldots, m$。则 $z_1, \ldots, z_m$
在 $A_\nu$ 中的像都是零。这恰好意味着，这些 $y_i$ 映到满足关系
$f_j(y'_1, \ldots, y'_n) = 0$ 的元 $y'_i \in A_\nu$。
从而得到环映射 $S \to A_\nu$。这证明 (1) 蕴含 (2)。

\medskip\noindent
(2) 显然蕴含 (3)。假设 (3)。
由引理 \ref{algebra-lemma-ring-colimit-fp}，可以写成
$S = \colim_\lambda S_\lambda$，其中 $S_\lambda$ 在 $R$ 上为有限表示的。
那么对某个 $\lambda$，恒等映射分解为
$$
S \to S_\lambda \to S
$$
对 $S \to S_\lambda \to S$ 应用引理 \ref{algebra-lemma-compose-finite-type} 的 (4)，
这蕴含 $S$ 在 $S_\lambda$ 上为有限表示的。
对 $R \to S_\lambda \to S$ 应用同一引理的 (2)，得出
$S$ 在 $R$ 上为有限表示的。
\end{proof}

\noindent
使用上面的基本材料，可以如下给出判别一个代数 $A$
何时是某一给定类型代数的滤过余极限的准则。

\begin{lemma}
\label{algebra-lemma-when-colimit}
设 $R \to \Lambda$ 是环映射。设 $\mathcal{E}$ 是一个 $R$-代数集合，且每个
$A \in \mathcal{E}$ 在 $R$ 上为有限表示的。则下列两个断言等价：
\begin{enumerate}
\item $\Lambda$ 是 $\mathcal{E}$ 中元的滤过余极限；以及
% P01-E0619: the frozen phrase “R algebra map” omits the compound hyphen.
\item 对任意 $R$-代数映射 $A \to \Lambda$，若 $A$ 在 $R$ 上为有限表示的，
则可以找到分解 $A \to B \to \Lambda$，其中 $B \in \mathcal{E}$。
\end{enumerate}
\end{lemma}

\begin{proof}
假设 $\mathcal{I} \to \mathcal{E}$、$i \mapsto A_i$ 是滤过图式，且
$\Lambda = \colim_i A_i$。设 $A \to \Lambda$ 是 $R$-代数映射，
且 $A$ 在 $R$ 上为有限表示的。
应用引理 \ref{algebra-lemma-characterize-finite-presentation}，得到分解
$A \to A_i \to \Lambda$。因此 (1) 蕴含 (2)。

\medskip\noindent
考虑引理 \ref{algebra-lemma-ring-colimit-fp-category} 中的范畴 $\mathcal{I}$。
由范畴篇引理 \ref{categories-lemma-cofinal-in-filtered}，
由这些满足 $A \in \mathcal{E}$ 的 $A \to \Lambda$ 组成的充满子范畴
$\mathcal{J}$ 在 $\mathcal{I}$ 中是共终的，并且它是滤过范畴。
再由范畴篇引理 \ref{categories-lemma-cofinal}，$\Lambda$ 也是
$\mathcal{J}$ 上的余极限。
\end{proof}

\noindent
但是还有更强的结果。也就是说，给定 $R = \colim_\lambda R_\lambda$，
% P01-E0620: the frozen text says “the limit of the category” although the following lemma states a colimit of categories.
有限表示 $R$-模的范畴等价于有限表示 $R_\lambda$-模范畴的极限。
对有限表示 $R$-代数的范畴也同样。

\begin{lemma}
\label{algebra-lemma-module-map-property-in-colimit}
设 $A$ 是环，设 $M, N$ 是 $A$-模。
假设 $R = \colim_{i \in I} R_i$ 是 $A$-代数的有向余极限。
\begin{enumerate}
\item 若 $M$ 是有限 $A$-模，且 $u, u' : M \to N$ 是 $A$-模映射，并且
$u \otimes 1 = u' \otimes 1 : M \otimes_A R \to N \otimes_A R$，
则对某个 $i$，有
$u \otimes 1 = u' \otimes 1 : M \otimes_A R_i \to N \otimes_A R_i$。
\item 若 $N$ 是有限 $A$-模，且 $u : M \to N$ 是 $A$-模映射，并且
$u \otimes 1 : M \otimes_A R \to N \otimes_A R$ 为满射，
则对某个 $i$，映射
$u \otimes 1 : M \otimes_A R_i \to N \otimes_A R_i$ 为满射。
\item 若 $N$ 是有限表示 $A$-模，且
$v : N \otimes_A R \to M \otimes_A R$ 是 $R$-模映射，则存在 $i$
和 $R_i$-模映射 $v_i : N \otimes_A R_i \to M \otimes_A R_i$，
使得 $v = v_i \otimes 1$。
\item 若 $M$ 是有限 $A$-模，$N$ 是有限表示 $A$-模，而且
$u : M \to N$ 是 $A$-模映射，使得
$u \otimes 1 : M \otimes_A R \to N \otimes_A R$ 是同构，则对某个 $i$，映射
$u \otimes 1 : M \otimes_A R_i \to N \otimes_A R_i$ 是同构。
\end{enumerate}
\end{lemma}

\begin{proof}
为证明 (1)，假设 $u$ 如 (1) 所述，并设 $x_1, \ldots, x_m \in M$
是生成元。由于 $N \otimes_A R = \colim_i N \otimes_A R_i$，
% P01-E0621: u(x_j) and u'(x_j) belong to N, but the frozen proof puts their equality in M tensor R_i.
可以选取 $i \in I$，使得对 $j = 1, \ldots, m$，
$u(x_j) \otimes 1 = u'(x_j) \otimes 1$ 在 $M \otimes_A R_i$ 中成立。
对这样的 $i$，我们有
$u \otimes 1 = u' \otimes 1 : M \otimes_A R_i \to N \otimes_A R_i$。

\medskip\noindent
为证明 (2)，假设 $u \otimes 1$ 为满射，并设
$y_1, \ldots, y_m \in N$ 是生成元。由于
$N \otimes_A R = \colim_i N \otimes_A R_i$，
% P01-E0622: the frozen proof says elements z_j of M tensor R_i have images in N tensor R, omitting the map u tensor 1.
可以选取 $i \in I$ 和 $z_j \in M \otimes_A R_i$、$j = 1, \ldots, m$，
使得它们在 $N \otimes_A R$ 中的像等于 $y_j \otimes 1$。
对这样的 $i$，映射 $u \otimes 1 : M \otimes_A R_i \to N \otimes_A R_i$
为满射。

\medskip\noindent
为证明 (3)，设 $y_1, \ldots, y_m \in N$ 是生成元。设
$K = \Ker(A^{\oplus m} \to N)$，其中该映射由规则
% P01-E0623: the frozen rule uses undefined x_j although the generators were named y_j.
$(a_1, \ldots, a_m) \mapsto \sum a_j x_j$ 给出。设 $k_1, \ldots, k_t$
是 $K$ 的生成元。写成 $k_s = (k_{s1}, \ldots, k_{sm})$。
由于 $M \otimes_A R = \colim_i M \otimes_A R_i$，可以选取
$i \in I$ 和 $z_j \in M \otimes_A R_i$、$j = 1, \ldots, m$，使得它们在
$M \otimes_A R$ 中的像等于 $v(y_j \otimes 1)$。
我们想用这些 $z_j$ 定义映射
$v_i : N \otimes_A R_i \to M \otimes_A R_i$。
由于 $K \otimes_A R_i \to R_i^{\oplus m} \to N \otimes_A R_i \to 0$
是一个表示，只须验证对每个 $s = 1, \ldots, t$，
$\xi_s = \sum_j k_{sj}z_j$ 在 $M \otimes_A R_i$ 中为零。
这不一定已经成立，但由于 $\xi_s$ 在 $M \otimes_A R$ 中的像为零，
把 $i$ 适当增大后它将成立。

\medskip\noindent
为证明 (4)，假设 $u \otimes 1$ 是同构、$M$ 是有限的，
且 $N$ 是有限表示的。设
$v : N \otimes_A R \to M \otimes_A R$ 是 $u \otimes 1$ 的逆。
应用 (3)，对某个 $i$ 得到映射
$v_i : N \otimes_A R_i \to M \otimes_A R_i$。
% P01-E0624: the frozen identities change the tensor base from A to R, making M tensor_R R_i and N tensor_R R_i generally undefined.
应用 (1)可知，增大 $i$ 后，有
$v_i \circ (u \otimes 1) = \text{id}_{M \otimes_R R_i}$ 和
$(u \otimes 1) \circ v_i = \text{id}_{N \otimes_R R_i}$。
\end{proof}

\begin{lemma}
\label{algebra-lemma-colimit-category-fp-modules}
假设 $R = \colim_{\lambda \in \Lambda} R_\lambda$ 是环的有向余极限。
则有限表示 $R$-模的范畴是有限表示 $R_\lambda$-模范畴的余极限。
更精确地说：
\begin{enumerate}
\item 给定有限表示 $R$-模 $M$，存在 $\lambda \in \Lambda$ 和有限表示
$R_\lambda$-模 $M_\lambda$，使得 $M \cong M_\lambda \otimes_{R_\lambda} R$。
\item 给定 $\lambda \in \Lambda$、有限表示 $R_\lambda$-模
$M_\lambda, N_\lambda$ 以及 $R$-模映射
$\varphi : M_\lambda \otimes_{R_\lambda} R \to N_\lambda \otimes_{R_\lambda} R$，
存在 $\mu \geq \lambda$ 和 $R_\mu$-模映射
$\varphi_\mu : M_\lambda \otimes_{R_\lambda} R_\mu \to
N_\lambda \otimes_{R_\lambda} R_\mu$，使得 $\varphi = \varphi_\mu \otimes 1_R$。
\item 给定 $\lambda \in \Lambda$、有限表示 $R_\lambda$-模
$M_\lambda, N_\lambda$ 以及 $R_\lambda$-模映射
$\varphi, \psi : M_\lambda \to N_\lambda$，若
$\varphi \otimes 1_R = \psi \otimes 1_R$，则对某个 $\mu \geq \lambda$，
$\varphi \otimes 1_{R_\mu} = \psi \otimes 1_{R_\mu}$。
\end{enumerate}
\end{lemma}

\begin{proof}
为证明 (1)，选取表示
$R^{\oplus m} \to R^{\oplus n} \to M \to 0$。
假设第一个映射由矩阵 $A = (a_{ij})$ 给出。
可以选取 $\lambda \in \Lambda$ 以及系数在 $R_\lambda$ 中的矩阵
$A_\lambda = (a_{\lambda, ij})$，使它在 $R$ 中映到 $A$。
然后，只需令 $M_\lambda$ 是具有表示
$R_\lambda^{\oplus m} \to R_\lambda^{\oplus n} \to M_\lambda \to 0$
的 $R_\lambda$-模，其中第一个箭头由 $A_\lambda$ 给出。

\medskip\noindent
(2) 和 (3) 由引理 \ref{algebra-lemma-module-map-property-in-colimit} 得出。
\end{proof}

\begin{lemma}
\label{algebra-lemma-algebra-map-property-in-colimit}
设 $A$ 是环，设 $B, C$ 是 $A$-代数。
假设 $R = \colim_{i \in I} R_i$ 是 $A$-代数的有向余极限。
\begin{enumerate}
\item 若 $B$ 是有限型 $A$-代数，且 $u, u' : B \to C$ 是 $A$-代数映射，并且
$u \otimes 1 = u' \otimes 1 : B \otimes_A R \to C \otimes_A R$，
则对某个 $i$，有
$u \otimes 1 = u' \otimes 1 : B \otimes_A R_i \to C \otimes_A R_i$。
\item 若 $C$ 是有限型 $A$-代数，且 $u : B \to C$ 是 $A$-代数映射，并且
$u \otimes 1 : B \otimes_A R \to C \otimes_A R$ 为满射，
则对某个 $i$，映射
$u \otimes 1 : B \otimes_A R_i \to C \otimes_A R_i$ 为满射。
\item 若 $C$ 在 $A$ 上为有限表示的，且
$v : C \otimes_A R \to B \otimes_A R$ 是 $R$-代数映射，则存在 $i$
和 $R_i$-代数映射 $v_i : C \otimes_A R_i \to B \otimes_A R_i$，
使得 $v = v_i \otimes 1$。
\item 若 $B$ 是有限型 $A$-代数，$C$ 是有限表示 $A$-代数，而且
$u \otimes 1 : B \otimes_A R \to C \otimes_A R$ 是同构，则对某个 $i$，映射
$u \otimes 1 : B \otimes_A R_i \to C \otimes_A R_i$ 是同构。
\end{enumerate}
\end{lemma}

\begin{proof}
为证明 (1)，假设 $u$ 如 (1) 所述，并设
$x_1, \ldots, x_m \in B$ 是生成元。
% P01-E0625: u(x_j) and u'(x_j) belong to C, but the frozen proof invokes the colimit for B and places their equality in B tensor R_i.
由于 $B \otimes_A R = \colim_i B \otimes_A R_i$，可以选取 $i \in I$，
使得对 $j = 1, \ldots, m$，$u(x_j) \otimes 1 = u'(x_j) \otimes 1$
在 $B \otimes_A R_i$ 中成立。对这样的 $i$，我们有
$u \otimes 1 = u' \otimes 1 : B \otimes_A R_i \to C \otimes_A R_i$。

\medskip\noindent
为证明 (2)，假设 $u \otimes 1$ 为满射，并设
$y_1, \ldots, y_m \in C$ 是生成元。由于
$B \otimes_A R = \colim_i B \otimes_A R_i$，
% P01-E0626: the frozen proof omits u tensor 1 when describing the images of z_j in C tensor R.
可以选取 $i \in I$ 和 $z_j \in B \otimes_A R_i$、$j = 1, \ldots, m$，
使得它们在 $C \otimes_A R$ 中的像等于 $y_j \otimes 1$。
对这样的 $i$，映射 $u \otimes 1 : B \otimes_A R_i \to C \otimes_A R_i$
为满射。

\medskip\noindent
为证明 (3)，设 $c_1, \ldots, c_m \in C$ 是生成元。设
% P01-E0627: the frozen presentation maps the polynomial ring to an undefined N instead of C.
$K = \Ker(A[x_1, \ldots, x_m] \to N)$，其中该映射由规则
% P01-E0628: the frozen rule sends every x_j to the index-free sum of c_j rather than to c_j.
$x_j \mapsto \sum c_j$ 给出。设 $f_1, \ldots, f_t$
是 $K$ 作为 $A[x_1, \ldots, x_m]$ 中理想的生成元。
我们把 $f_j = f_j(x_1, \ldots, x_m)$ 看作多项式。
由于 $B \otimes_A R = \colim_i B \otimes_A R_i$，可以选取
$i \in I$ 和 $z_j \in B \otimes_A R_i$、$j = 1, \ldots, m$，使得它们在
$B \otimes_A R$ 中的像等于 $v(c_j \otimes 1)$。
我们想用这些 $z_j$ 定义映射
$v_i : C \otimes_A R_i \to B \otimes_A R_i$。
由于 $K \otimes_A R_i \to R_i[x_1, \ldots, x_m] \to C \otimes_A R_i \to 0$
是一个表示，只须验证：
% P01-E0629: the frozen definition uses f_j while the element and quantifier are indexed by s.
$\xi_s = f_j(z_1, \ldots, z_m)$ 对每个 $s = 1, \ldots, t$
都在 $B \otimes_A R_i$ 中为零。这不一定已经成立，但由于
$\xi_s$ 在 $B \otimes_A R$ 中的像为零，把 $i$ 适当增大后它将成立。

\medskip\noindent
为证明 (4)，假设 $u \otimes 1$ 是同构、$B$ 是有限型
$A$-代数，并且 $C$ 是有限表示 $A$-代数。
% P01-E0630: the frozen inverse v is typed B tensor R to C tensor R, the same direction as u, while the next sentence correctly types v_i in the opposite direction.
设 $v : B \otimes_A R \to C \otimes_A R$ 是 $u \otimes 1$ 的逆。
设 $v_i : C \otimes_A R_i \to B \otimes_A R_i$ 如 (3) 中所述。
% P01-E0631: the frozen identities change the tensor base from A to R.
应用 (1) 可知，增大 $i$ 后，有
$v_i \circ (u \otimes 1) = \text{id}_{B \otimes_R R_i}$ 和
$(u \otimes 1) \circ v_i = \text{id}_{C \otimes_R R_i}$。
\end{proof}

\begin{lemma}
\label{algebra-lemma-colimit-category-fp-algebras}
假设 $R = \colim_{\lambda \in \Lambda} R_\lambda$ 是环的有向余极限。
则有限表示 $R$-代数的范畴是有限表示 $R_\lambda$-代数范畴的余极限。
更精确地说：
\begin{enumerate}
\item 给定有限表示 $R$-代数 $A$，存在 $\lambda \in \Lambda$ 和有限表示
$R_\lambda$-代数 $A_\lambda$，使得 $A \cong A_\lambda \otimes_{R_\lambda} R$。
\item 给定 $\lambda \in \Lambda$、有限表示 $R_\lambda$-代数
$A_\lambda, B_\lambda$ 以及 $R$-代数映射
$\varphi : A_\lambda \otimes_{R_\lambda} R \to B_\lambda \otimes_{R_\lambda} R$，
存在 $\mu \geq \lambda$ 和 $R_\mu$-代数映射
$\varphi_\mu : A_\lambda \otimes_{R_\lambda} R_\mu \to
B_\lambda \otimes_{R_\lambda} R_\mu$，使得 $\varphi = \varphi_\mu \otimes 1_R$。
\item 给定 $\lambda \in \Lambda$、有限表示 $R_\lambda$-代数
$A_\lambda, B_\lambda$ 以及 $R_\lambda$-代数映射
$\varphi_\lambda, \psi_\lambda : A_\lambda \to B_\lambda$，
% P01-E0632: the frozen equalities drop the lambda subscripts from the maps just introduced.
若 $\varphi \otimes 1_R = \psi \otimes 1_R$，则对某个 $\mu \geq \lambda$，
$\varphi \otimes 1_{R_\mu} = \psi \otimes 1_{R_\mu}$。
\end{enumerate}
\end{lemma}

\begin{proof}
为证明 (1)，选取表示
$A = R[x_1, \ldots, x_n]/(f_1, \ldots, f_m)$。
可以选取 $\lambda \in \Lambda$ 和映到
$f_j \in R[x_1, \ldots, x_n]$ 的元
$f_{\lambda, j} \in R_\lambda[x_1, \ldots, x_n]$。
然后只需令
$A_\lambda =
R_\lambda[x_1, \ldots, x_n]/(f_{\lambda, 1}, \ldots, f_{\lambda, m})$。

\medskip\noindent
(2) 和 (3) 由引理 \ref{algebra-lemma-algebra-map-property-in-colimit} 得出。
\end{proof}

\begin{lemma}
\label{algebra-lemma-limit-no-condition-local}
假设 $R \to S$ 是局部环的局部同态。
存在有向集 $(\Lambda, \leq)$ 以及局部环的局部同态
$R_\lambda \to S_\lambda$ 所组成的系统，使得
\begin{enumerate}
\item 系统 $R_\lambda \to S_\lambda$ 的余极限等于 $R \to S$。
\item 每个 $R_\lambda$ 在 $\mathbf{Z}$ 上本质有限型。
\item 每个 $S_\lambda$ 在 $R_\lambda$ 上本质有限型。
\end{enumerate}
\end{lemma}

\begin{proof}
% P01-E0633: the frozen phrase “Denote phi ... the ring map” omits “by”.
以 $\varphi : R  \to S$ 表示该环映射。
设 $\mathfrak m \subset R$ 是 $R$ 的极大理想，设
$\mathfrak n \subset S$ 是 $S$ 的极大理想。令
$$
\Lambda = \{
(A, B)
\mid
A \subset R, B \subset S, \# A < \infty, \# B < \infty, \varphi(A) \subset B
\}.
$$
我们取包含关系作为偏序。对每个
$\lambda = (A, B) \in \Lambda$，设 $R'_\lambda$ 是由 $a \in A$ 生成的
$\mathbf{Z}$-子代数，并设 $S'_\lambda$ 是由 $b$、$b \in B$ 生成的
$\mathbf{Z}$-子代数。设 $R_\lambda$ 是 $R'_\lambda$ 在素理想
$R'_\lambda \cap \mathfrak m$ 处的局部化，并设 $S_\lambda$ 是
$S'_\lambda$ 在素理想 $S'_\lambda \cap \mathfrak n$ 处的局部化。
用图表示为
$$
\xymatrix{
B \ar[r] &
S'_\lambda \ar[r] &
S_\lambda \ar[r] &
S \\
A \ar[r] \ar[u] &
R'_\lambda \ar[r] \ar[u] &
R_\lambda \ar[r] \ar[u] &
R \ar[u]
}.
$$
转移映射是清楚的。其他断言的证明留给读者。
\end{proof}

\begin{lemma}
\label{algebra-lemma-limit-essentially-finite-type}
假设 $R \to S$ 是局部环的局部同态。
假设 $S$ 在 $R$ 上本质有限型。
则存在有向集 $(\Lambda, \leq)$ 以及局部环的局部同态
$R_\lambda \to S_\lambda$ 所组成的系统，使得
\begin{enumerate}
\item 系统 $R_\lambda \to S_\lambda$ 的余极限等于 $R \to S$。
\item 每个 $R_\lambda$ 在 $\mathbf{Z}$ 上本质有限型。
\item 每个 $S_\lambda$ 在 $R_\lambda$ 上本质有限型。
\item 对每个 $\lambda \leq \mu$，映射
$S_\lambda \otimes_{R_\lambda} R_\mu \to S_\mu$
把 $S_\mu$ 表示为 $S_\lambda \otimes_{R_\lambda} R_\mu$ 的一个商的局部化。
\end{enumerate}
\end{lemma}

\begin{proof}
以 $\varphi : R  \to S$ 表示该环映射。
设 $\mathfrak m \subset R$ 是 $R$ 的极大理想，设
$\mathfrak n \subset S$ 是 $S$ 的极大理想。
设 $x_1, \ldots, x_n \in S$ 是这样的元：$S$ 是 $S$ 的由
$x_1, \ldots, x_n$ 生成的 $R$-子代数的一个局部化。
换言之，$S$ 是多项式环 $R[x_1, \ldots, x_n]$ 的一个局部化的商。

\medskip\noindent
设 $\Lambda = \{ A \subset R \mid \# A < \infty\}$ 是 $R$ 的有限子集所组成的集合。
我们取包含关系作为偏序。对每个
$\lambda = A \in \Lambda$，设 $R'_\lambda$ 是由 $a \in A$ 生成的
$\mathbf{Z}$-子代数，并设 $S'_\lambda$ 是由 $\varphi(a)$、$a \in A$
和元 $x_1, \ldots, x_n$ 生成的 $\mathbf{Z}$-子代数。
设 $R_\lambda$ 是 $R'_\lambda$ 在素理想 $R'_\lambda \cap \mathfrak m$ 处的局部化，
并设 $S_\lambda$ 是 $S'_\lambda$ 在素理想
$S'_\lambda \cap \mathfrak n$ 处的局部化。用图表示为
$$
\xymatrix{
\varphi(A) \amalg \{x_i\} \ar[r] &
S'_\lambda \ar[r] &
S_\lambda \ar[r] &
S \\
A \ar[r] \ar[u] &
R'_\lambda \ar[r] \ar[u] &
R_\lambda \ar[r] \ar[u] &
R \ar[u]
}
$$
显然，若 $A \subset B$ 对应于 $\Lambda$ 中的 $\lambda \leq \mu$，
则存在典范映射 $R_\lambda \to R_\mu$ 和 $S_\lambda \to S_\mu$，
从而在有向集 $\Lambda$ 上得到一个系统。

\medskip\noindent
断言 $R = \colim R_\lambda$ 是清楚的，因为所有映射
$R_\lambda \to R$ 都是单射，而 $R$ 的任意元最终都在像中。
同样的论证适用于 $S = \colim S_\lambda$。
断言 (2)、(3) 按构造成立。最后一个断言成立，因为映射
$S'_\lambda \otimes_{R'_\lambda} R'_\mu \to S'_\mu$ 显然是满射。
\end{proof}

\begin{lemma}
\label{algebra-lemma-limit-essentially-finite-presentation}
假设 $R \to S$ 是局部环的局部同态。
假设 $S$ 在 $R$ 上本质有限表示。
% P01-E0634: the frozen statement uses singular “a system of local homomorphism”.
则存在有向集 $(\Lambda, \leq)$ 以及局部环的局部同态
$R_\lambda \to S_\lambda$ 所组成的系统，使得
\begin{enumerate}
\item 系统 $R_\lambda \to S_\lambda$ 的余极限等于 $R \to S$。
\item 每个 $R_\lambda$ 在 $\mathbf{Z}$ 上本质有限型。
\item 每个 $S_\lambda$ 在 $R_\lambda$ 上本质有限型。
\item 对每个 $\lambda \leq \mu$，映射
$S_\lambda \otimes_{R_\lambda} R_\mu \to S_\mu$
把 $S_\mu$ 表示为 $S_\lambda \otimes_{R_\lambda} R_\mu$
在一个素理想处的局部化。
\end{enumerate}
\end{lemma}

\begin{proof}
由假设，可以选取同构
$\Phi : (R[x_1, \ldots, x_n]/I)_{\mathfrak q} \to S$，
其中 $I \subset R[x_1, \ldots, x_n]$ 是有限生成理想，
$\mathfrak q \subset R[x_1, \ldots, x_n]/I$ 是素理想。
（注意，$R \cap \mathfrak q$ 等于 $R$ 的极大理想 $\mathfrak m$。）
我们还选取理想 $I$ 的生成元 $f_1, \ldots, f_m \in I$。
以任意方式把 $R$ 写成有向集 $(\Lambda, \leq )$ 上的余极限
$R = \colim R_\lambda$，其中每个 $R_\lambda$ 都是局部环，
并在 $\mathbf{Z}$ 上本质有限型。
存在某个 $\lambda_0 \in \Lambda$，使得 $f_j$ 是某个
$f_{j, \lambda_0} \in R_{\lambda_0}[x_1, \ldots, x_n]$ 的像。
对所有 $\lambda \geq \lambda_0$，以
$f_{j, \lambda} \in R_{\lambda}[x_1, \ldots, x_n]$ 表示 $f_{j, \lambda_0}$ 的像。
从而得到环映射系统
$$
R_\lambda[x_1, \ldots, x_n]/(f_{1, \lambda}, \ldots, f_{m, \lambda})
\to
R[x_1, \ldots, x_n]/(f_1, \ldots, f_m) \to S
$$
% P01-E0635: the frozen phrase “Set q_lambda the inverse image” omits “to be”.
设 $\mathfrak q_\lambda$ 为 $\mathfrak q$ 的逆像。置
$S_\lambda = (R_\lambda[x_1, \ldots, x_n]/
(f_{1, \lambda}, \ldots, f_{m, \lambda}))_{\mathfrak q_\lambda}$。
这一构造确实符合要求，其验证留给读者。
\end{proof}

\begin{remark}
\label{algebra-remark-suitable-systems-limits}
假设 $R \to S$ 是局部环的局部同态，并且它本质上是有限表示的。
取任意系统 $(\Lambda, \leq)$、$R_\lambda \to S_\lambda$，使它具有引理
\ref{algebra-lemma-limit-essentially-finite-type} 中列出的性质。
可能出现的情况是，这是“错误的”系统；也就是说，
引理 \ref{algebra-lemma-limit-essentially-finite-presentation} 的性质 (4) 可能不成立。
下面是一个例子。设 $k = \mathbf{F}_2$。考虑环
$$
R = \text{localization of }
k[z, y_1, y_2, \ldots]/(y_i^2 - zy_{i + 1})
\text{ 在 }(z, y_1, y_2, \ldots)
$$
置 $S = R/zR$。取 $\Lambda = \mathbf{N}$ 作为系统，并令
$$
R_n = \text{localization of }
k[z, y_1, \ldots, y_n]/(\{y_i^2 - zy_{i + 1}\}_{i \leq n-1})
\text{ 在 }(z, y_1, \ldots, y_n)
$$
以及 $S_n = R_n/(z, y_n^2)$。
% P01-E0636: “All the maps ... are not” has ambiguous quantifier scope; the proof means none of them is a localization.
所有映射 $S_n \otimes_{R_n} R_{n + 1} \to S_{n + 1}$
都不是局部化（在本情形中也就是同构），因为
$1 \otimes y_{n + 1}^2$ 映到零。
若改取 $S_n' = R_n/zR_n$，则映射
$S'_n \otimes_{R_n} R_{n + 1} \to S'_{n + 1}$ 是同构。
本备注的要点是：选择这些系统时确实必须稍加小心。
\end{remark}

\begin{lemma}
\label{algebra-lemma-limit-module-essentially-finite-presentation}
假设 $R \to S$ 是局部环的局部同态。
假设 $S$ 在 $R$ 上本质有限表示。
设 $M$ 是有限表示 $S$-模。
则存在有向集 $(\Lambda, \leq)$，以及局部环的局部同态
$R_\lambda \to S_\lambda$ 所组成的系统和 $S_\lambda$-模 $M_\lambda$，使得
\begin{enumerate}
\item 系统 $R_\lambda \to S_\lambda$ 的余极限等于 $R \to S$。
系统 $M_\lambda$ 的余极限是 $M$。
\item 每个 $R_\lambda$ 在 $\mathbf{Z}$ 上本质有限型。
\item 每个 $S_\lambda$ 在 $R_\lambda$ 上本质有限型。
\item 每个 $M_\lambda$ 在 $S_\lambda$ 上都是有限的。
\item 对每个 $\lambda \leq \mu$，映射
$S_\lambda \otimes_{R_\lambda} R_\mu \to S_\mu$
把 $S_\mu$ 表示为 $S_\lambda \otimes_{R_\lambda} R_\mu$
在一个素理想处的局部化。
\item 对每个 $\lambda \leq \mu$，映射
$M_\lambda \otimes_{S_\lambda} S_\mu \to M_\mu$ 是同构。
\end{enumerate}
\end{lemma}

\begin{proof}
与引理 \ref{algebra-lemma-limit-essentially-finite-presentation} 的证明一样，
首先可以把 $R = \colim R_\lambda$ 写成局部 $\mathbf{Z}$-代数的有向余极限，
其中各代数本质上是有限型的。
接着，可以假设对某个 $\lambda_1 \in \Lambda$，存在
$f_{j, \lambda_1} \in R_{\lambda_1}[x_1, \ldots, x_n]$，使得
$$
S =
\colim_{\lambda \geq \lambda_1} S_\lambda, \text{ 其中 }
S_\lambda =
(R_\lambda[x_1, \ldots, x_n]/
(f_{1, \lambda}, \ldots, f_{m, \lambda}))_{\mathfrak q_\lambda}
$$
选取 $M$ 在 $S$ 上的一个表示
$$
S^{\oplus s} \to S^{\oplus t} \to M \to 0
$$
设 $A \in \text{Mat}(t \times s, S)$ 是该表示的矩阵。
对某个 $\lambda_2 \in \Lambda$、$\lambda_2 \geq \lambda_1$，
可以找到映到 $A$ 的矩阵
$A_{\lambda_2} \in \text{Mat}(t \times s, S_{\lambda_2})$。
对所有 $\lambda \geq \lambda_2$，令
$M_\lambda = \Coker(S_\lambda^{\oplus s} \xrightarrow{A_\lambda}
S_\lambda^{\oplus t})$。这一构造确实符合要求，其验证留给读者。
\end{proof}

\begin{lemma}
\label{algebra-lemma-limit-no-condition}
假设 $R \to S$ 是环映射。
则存在有向集 $(\Lambda, \leq)$ 以及环映射
$R_\lambda \to S_\lambda$ 所组成的系统，使得
\begin{enumerate}
\item 系统 $R_\lambda \to S_\lambda$ 的余极限等于 $R \to S$。
\item 每个 $R_\lambda$ 在 $\mathbf{Z}$ 上为有限型的。
\item 每个 $S_\lambda$ 在 $R_\lambda$ 上为有限型的。
\end{enumerate}
\end{lemma}

\begin{proof}
这是引理 \ref{algebra-lemma-limit-no-condition-local} 的非局部版本。
证明类似，留给读者。
\end{proof}

\begin{lemma}
\label{algebra-lemma-limit-integral}
假设 $R \to S$ 是环映射。
假设 $S$ 在 $R$ 上是整的。
则存在有向集 $(\Lambda, \leq)$ 以及环映射
$R_\lambda \to S_\lambda$ 所组成的系统，使得
\begin{enumerate}
\item 系统 $R_\lambda \to S_\lambda$ 的余极限等于 $R \to S$。
\item 每个 $R_\lambda$ 在 $\mathbf{Z}$ 上为有限型的。
\item 每个 $S_\lambda$ 在 $R_\lambda$ 上都是有限的。
\end{enumerate}
\end{lemma}

\begin{proof}
考虑由对 $(E, F)$ 组成的集合 $\Lambda$，其中 $E \subset R$ 是有限子集，
$F \subset S$ 是有限子集，并且每个元 $f \in F$ 都是某个首一多项式
$P(X) \in R[X]$ 的根，而该多项式的系数属于 $E$。
定义 $(E, F) \leq (E', F')$ 当且仅当 $E \subset E'$ 且 $F \subset F'$。
给定 $\lambda = (E, F) \in \Lambda$，设 $R_\lambda \subset R$ 等于由 $E$ 生成的
$R$ 的 $\mathbf{Z}$-子代数，并设 $S_\lambda \subset S$ 等于由 $F$ 和 $E$
在 $S$ 中的像生成的 $\mathbf{Z}$-子代数。
显然，$R = \colim R_\lambda$。
% P01-E0637: the frozen proof says every element of S is integral over S, which is tautological and does not justify the construction; the hypothesis is integral over R.
我们有 $S = \colim S_\lambda$，因为 $S$ 的每个元在 $S$ 上都是整的。
由引理 \ref{algebra-lemma-characterize-finite-in-terms-of-integral} 以及下述事实，
环映射 $R_\lambda \to S_\lambda$ 是有限的：$S_\lambda$ 在 $R_\lambda$
上由 $F$ 中的元生成，而按我们对对 $(E, F)$ 施加的条件，
这些元在 $R_\lambda$ 上是整的。引理得证。
\end{proof}

\begin{lemma}
\label{algebra-lemma-limit-finite-type}
假设 $R \to S$ 是环映射。
假设 $S$ 在 $R$ 上为有限型的。
则存在有向集 $(\Lambda, \leq)$ 以及环映射
$R_\lambda \to S_\lambda$ 所组成的系统，使得
\begin{enumerate}
\item 系统 $R_\lambda \to S_\lambda$ 的余极限等于 $R \to S$。
\item 每个 $R_\lambda$ 在 $\mathbf{Z}$ 上为有限型的。
\item 每个 $S_\lambda$ 在 $R_\lambda$ 上为有限型的。
\item 对每个 $\lambda \leq \mu$，映射
$S_\lambda \otimes_{R_\lambda} R_\mu \to S_\mu$
把 $S_\mu$ 表示为 $S_\lambda \otimes_{R_\lambda} R_\mu$ 的商。
\end{enumerate}
\end{lemma}

\begin{proof}
这是引理 \ref{algebra-lemma-limit-essentially-finite-type} 的非局部版本。
证明类似，留给读者。
\end{proof}

\begin{lemma}
\label{algebra-lemma-limit-finite-presentation}
假设 $R \to S$ 是环映射。
假设 $S$ 在 $R$ 上为有限表示的。
则存在有向集 $(\Lambda, \leq)$ 以及环映射
$R_\lambda \to S_\lambda$ 所组成的系统，使得
\begin{enumerate}
\item 系统 $R_\lambda \to S_\lambda$ 的余极限等于 $R \to S$。
\item 每个 $R_\lambda$ 在 $\mathbf{Z}$ 上为有限型的。
\item 每个 $S_\lambda$ 在 $R_\lambda$ 上为有限型的。
\item 对每个 $\lambda \leq \mu$，映射
$S_\lambda \otimes_{R_\lambda} R_\mu \to S_\mu$ 是同构。
\end{enumerate}
\end{lemma}

\begin{proof}
这是引理 \ref{algebra-lemma-limit-essentially-finite-presentation} 的非局部版本。
证明类似，留给读者。
\end{proof}

\begin{lemma}
\label{algebra-lemma-limit-module-finite-presentation}
假设 $R \to S$ 是环映射。
假设 $S$ 在 $R$ 上为有限表示的。
设 $M$ 是有限表示 $S$-模。
则存在有向集 $(\Lambda, \leq)$，以及环映射
$R_\lambda \to S_\lambda$ 所组成的系统和 $S_\lambda$-模 $M_\lambda$，使得
\begin{enumerate}
\item 系统 $R_\lambda \to S_\lambda$ 的余极限等于 $R \to S$。
系统 $M_\lambda$ 的余极限是 $M$。
\item 每个 $R_\lambda$ 在 $\mathbf{Z}$ 上为有限型的。
\item 每个 $S_\lambda$ 在 $R_\lambda$ 上为有限型的。
\item 每个 $M_\lambda$ 在 $S_\lambda$ 上都是有限的。
\item 对每个 $\lambda \leq \mu$，映射
$S_\lambda \otimes_{R_\lambda} R_\mu \to S_\mu$ 是同构。
\item 对每个 $\lambda \leq \mu$，映射
$M_\lambda \otimes_{S_\lambda} S_\mu \to M_\mu$ 是同构。
\end{enumerate}
特别地，对每个 $\lambda \in \Lambda$，我们有
$$
M = M_\lambda \otimes_{S_\lambda} S
= M_\lambda \otimes_{R_\lambda} R.
$$
\end{lemma}

\begin{proof}
这是引理 \ref{algebra-lemma-limit-module-essentially-finite-presentation} 的非局部版本。
证明类似，留给读者。
\end{proof}

\section{更多平坦性准则}
\label{algebra-section-more-flatness-criteria}

\noindent
下面的引理在代数几何中常用于证明：从正规曲面到光滑曲面的有限态射是平坦的。
它是引理 \ref{algebra-lemma-finite-flat-over-regular-CM} 的一个部分逆命题，
因为单射有限局部环映射当然满足条件 (3)。

\begin{lemma}
\label{algebra-lemma-CM-over-regular-flat}
\begin{slogan}
奇迹平坦性
\end{slogan}
设 $R \to S$ 是 Noether 局部环的局部同态。假设
\begin{enumerate}
\item $R$ 是正则的，
% P01-E0638: the frozen second assumption omits “is”.
\item $S$ Cohen--Macaulay，
\item $\dim(S) = \dim(R) + \dim(S/\mathfrak m_R S)$。
\end{enumerate}
则 $R \to S$ 是平坦的。
\end{lemma}

\begin{proof}
对 $\dim(R)$ 作归纳。情形 $\dim(R) = 0$ 是显然的，因为此时
$R$ 是域。假设 $\dim(R) > 0$。由 (3)，这蕴含 $\dim(S) > 0$。
设 $\mathfrak q_1, \ldots, \mathfrak q_r$ 是 $S$ 的极小素理想。
注意 $\mathfrak q_i \not \supset \mathfrak m_R S$，因为
$$
\dim(S/\mathfrak q_i) = \dim(S) > \dim(S/\mathfrak m_R S)
$$
% P01-E0639: the frozen sentence lacks punctuation after the display before its explanation continues.
其中第一个等式由引理 \ref{algebra-lemma-maximal-chain-CM} 得出，
不等式由 (3) 得出。因此
$\mathfrak p_i = R \cap \mathfrak q_i$ 不等于 $\mathfrak m_R$。
选取 $x \in \mathfrak m_R$，使得 $x \not \in \mathfrak m_R^2$ 且
$x \not \in \mathfrak p_i$；见引理 \ref{algebra-lemma-silly}。
从而可知 $x$ 不包含于 $S$ 的任一极小素理想中。
因此由 (2)，$x$ 是 $S$ 上的非零因子；见引理
\ref{algebra-lemma-reformulate-CM}，并且 $S/xS$ 是 Cohen--Macaulay 的，而且
$\dim(S/xS) = \dim(S) - 1$。
由 (1) 和引理 \ref{algebra-lemma-regular-ring-CM}，环 $R/xR$ 是正则的，且
$\dim(R/xR) = \dim(R) - 1$。
由归纳可知 $R/xR \to S/xS$ 是平坦的。
因此，由引理 \ref{algebra-lemma-variant-local-criterion-flatness} 及其后的备注，得出结论。
\end{proof}

\begin{lemma}
\label{algebra-lemma-flat-over-regular}
设 $R \to S$ 是 Noether 局部环的同态。
假设 $R$ 是正则局部环，并且一个正则参数系映到 $S$ 中的一个正则列。
则 $R \to S$ 是平坦的。
\end{lemma}

\begin{proof}
假设 $x_1, \ldots, x_d$ 是 $R$ 的一个参数系，它在 $S$ 中映到一个正则列。
注意，$S/(x_1, \ldots, x_d)S$ 在 $R/(x_1, \ldots, x_d)$ 上平坦，
因为后者是域。此时 $x_d$ 是 $S/(x_1, \ldots, x_{d - 1})S$ 中的非零因子，
因此由平坦性局部准则，$S/(x_1, \ldots, x_{d - 1})S$ 在
$R/(x_1, \ldots, x_{d - 1})$ 上平坦（见引理
\ref{algebra-lemma-variant-local-criterion-flatness} 及其后的备注）。
接着，$x_{d - 1}$ 是 $S/(x_1, \ldots, x_{d - 2})S$ 中的非零因子，
因此由平坦性局部准则，$S/(x_1, \ldots, x_{d - 2})S$ 在
$R/(x_1, \ldots, x_{d - 2})$ 上平坦（见引理
\ref{algebra-lemma-variant-local-criterion-flatness} 及其后的备注）。
继续如此，直到得出 $S$ 在 $R$ 上平坦的结论。
\end{proof}

\noindent
下面的引理是证明如下事实的关键：基环上有限表示环上的有限表示模的结果，
可以从 Noether 情形中有限模的相应结果推出。

\begin{lemma}
\label{algebra-lemma-colimit-eventually-flat}
设 $R \to S$、$M$、$\Lambda$、$R_\lambda \to S_\lambda$、$M_\lambda$
如引理 \ref{algebra-lemma-limit-module-essentially-finite-presentation} 中所述。
假设 $M$ 在 $R$ 上平坦。则对某个 $\lambda \in \Lambda$，模
$M_\lambda$ 在 $R_\lambda$ 上平坦。
\end{lemma}

\begin{proof}
选取某个 $\lambda \in \Lambda$，并考虑
$$
\text{Tor}_1^{R_\lambda}(M_\lambda, R_\lambda/\mathfrak m_\lambda)
=
\Ker(\mathfrak m_\lambda \otimes_{R_\lambda} M_\lambda
\to M_\lambda).
$$
见备注 \ref{algebra-remark-Tor-ring-mod-ideal}。右端表明这是有限生成
$S_\lambda$-模（因为 $S_\lambda$ 是 Noether 的，而所涉模都是有限的）。
设 $\xi_1, \ldots, \xi_n$ 是生成元。
由于 $M$ 在 $R$ 上平坦，我们有
$0 = \Ker(\mathfrak m_\lambda R \otimes_R M \to M)$。
由于 $\otimes$ 与余极限可换，存在 $\lambda' \geq \lambda$，
使得每个 $\xi_i$ 都在
$\mathfrak m_{\lambda}R_{\lambda'} \otimes_{R_{\lambda'}} M_{\lambda'}$ 中映到零。
因此映射
$$
\text{Tor}_1^{R_\lambda}(M_\lambda, R_\lambda/\mathfrak m_\lambda)
\longrightarrow
\text{Tor}_1^{R_{\lambda'}}(M_{\lambda'},
R_{\lambda'}/\mathfrak m_{\lambda}R_{\lambda'})
$$
为零。注意，$M_\lambda \otimes_{R_\lambda} R_\lambda/\mathfrak m_\lambda$
在 $R_\lambda/\mathfrak m_\lambda$ 上平坦，因为后一个环是域。
因此可以应用引理 \ref{algebra-lemma-another-variant-local-criterion-flatness}，得到
$M_{\lambda'}$ 在 $R_{\lambda'}$ 上平坦。
\end{proof}

\noindent
使用上面的引理，我们可以开始在非 Noether 情形中重新证明第
\ref{algebra-section-criteria-flatness} 节的各个结果。

\begin{lemma}
\label{algebra-lemma-mod-injective-general}
假设 $R \to S$ 是局部环的局部同态。
% P01-E0640: three nearby statements use “Denote m the maximal ideal” without “by/as”.
以 $\mathfrak m$ 表示 $R$ 的极大理想。
设 $u : M \to N$ 是 $S$-模的映射。假设
\begin{enumerate}
\item $S$ 在 $R$ 上本质有限表示，
\item $M$、$N$ 在 $S$ 上为有限表示的，
\item $N$ 在 $R$ 上平坦，以及
\item $\overline{u} : M/\mathfrak mM \to N/\mathfrak mN$ 是单射。
\end{enumerate}
则 $u$ 是单射，且 $N/u(M)$ 在 $R$ 上平坦。
\end{lemma}

\begin{proof}
由引理 \ref{algebra-lemma-limit-module-essentially-finite-presentation} 及其证明，
可以找到局部环映射系统 $R_\lambda \to S_\lambda$ 以及
$S_\lambda$-模映射 $u_\lambda : M_\lambda \to N_\lambda$，它们对该引理中的
% P01-E0641: the frozen phrase “conclusions (1)--(6) for both N and M of that lemma” has malformed attachment.
$N$ 和 $M$ 都满足结论 (1) -- (6)，并且这些映射 $u_\lambda$ 的余极限是 $u$。
由引理 \ref{algebra-lemma-colimit-eventually-flat}，可以假设对所有充分大的
$\lambda$，$N_\lambda$ 在 $R_\lambda$ 上平坦。
以 $\mathfrak m_\lambda \subset R_\lambda$ 表示极大理想，并以
$\kappa_\lambda = R_\lambda / \mathfrak m_\lambda$，分别以
$\kappa = R/\mathfrak m$ 表示剩余域。

\medskip\noindent
考虑映射
$$
\Psi_\lambda :
M_\lambda/\mathfrak m_\lambda M_\lambda \otimes_{\kappa_\lambda} \kappa
\longrightarrow
M/\mathfrak m M.
$$
由于 $S_\lambda/\mathfrak m_\lambda S_\lambda$ 在域 $\kappa_\lambda$ 上本质有限型，
张量积 $S_\lambda/\mathfrak m_\lambda S_\lambda \otimes_{\kappa_\lambda} \kappa$
在 $\kappa$ 上本质有限型。因此它是 Noether 环，从而得出
$\Psi_\lambda$ 的核是有限生成的。
由于 $M/\mathfrak m M$ 是系统 $M_\lambda/\mathfrak m_\lambda M_\lambda$ 的余极限，
而 $\kappa$ 是这些域 $\kappa_\lambda$ 的余极限，
% P01-E0642: a directed preorder guarantees a later index >= lambda, not necessarily a strict index > lambda.
存在 $\lambda' > \lambda$，使得 $\Psi_\lambda$ 的核由以下映射的核生成：
$$
\Psi_{\lambda, \lambda'} :
M_\lambda/\mathfrak m_\lambda M_\lambda
\otimes_{\kappa_\lambda}
\kappa_{\lambda'}
\longrightarrow
M_{\lambda'}/\mathfrak m_{\lambda'} M_{\lambda'}.
$$
按构造，存在乘法子集
$W \subset S_\lambda \otimes_{R_\lambda} R_{\lambda'}$，使得
$S_{\lambda'} = W^{-1}(S_\lambda \otimes_{R_\lambda} R_{\lambda'})$，且
$$
W^{-1}(M_\lambda/\mathfrak m_\lambda M_\lambda
\otimes_{\kappa_\lambda}
\kappa_{\lambda'})
=
M_{\lambda'}/\mathfrak m_{\lambda'} M_{\lambda'}.
$$
现在，假设 $x$ 是以下映射的核中的元：
$$
\Psi_{\lambda'} :
M_{\lambda'}/\mathfrak m_{\lambda'} M_{\lambda'}
\otimes_{\kappa_{\lambda'}} \kappa
\longrightarrow
M/\mathfrak m M.
$$
% P01-E0643: the frozen membership for wx omits the tensor base and places a localized element back in the earlier module without specifying a representative.
则对某个 $w \in W$，有
$wx \in M_\lambda/\mathfrak m_\lambda M_\lambda \otimes \kappa$。
因此 $wx \in \Ker(\Psi_\lambda)$。所以 $wx$ 是
$\Psi_{\lambda, \lambda'}$ 的核中各元的线性组合。
从而 $wx = 0$ 在 $M_{\lambda'}/\mathfrak m_{\lambda'} M_{\lambda'}
\otimes_{\kappa_{\lambda'}} \kappa$ 中成立；又因为 $w$ 在 $S_{\lambda'}$ 中可逆，有 $x = 0$。
我们得出，对所有充分大的 $\lambda'$，$\Psi_{\lambda'}$ 的核都是零！

\medskip\noindent
由前一段的结果，可以假设对所有充分大的 $\lambda$，
$\Psi_\lambda$ 的核都是零，这蕴含映射
$M_\lambda/\mathfrak m_\lambda M_\lambda \to M/\mathfrak m M$ 是单射。
结合 $\overline{u}$ 为单射这一事实，形式地推出
$\overline{u_\lambda} : M_\lambda/\mathfrak m_\lambda M_\lambda
\to N_\lambda/\mathfrak m_\lambda N_\lambda$ 也是单射。
由引理 \ref{algebra-lemma-mod-injective}，得出（对所有充分大的 $\lambda$）
映射 $u_\lambda$ 是单射，且 $N_\lambda/u_\lambda(M_\lambda)$ 在
$R_\lambda$ 上平坦。引理得证。
\end{proof}

\begin{lemma}
\label{algebra-lemma-grothendieck-general}
假设 $R \to S$ 是局部环的局部环同态。
以 $\mathfrak m$ 表示 $R$ 的极大理想。假设
\begin{enumerate}
\item $S$ 在 $R$ 上本质有限表示，
\item $S$ 在 $R$ 上平坦，以及
\item $f \in S$ 是 $S/{\mathfrak m}S$ 中的非零因子。
\end{enumerate}
则 $S/fS$ 在 $R$ 上平坦，而且 $f$ 是 $S$ 中的非零因子。
\end{lemma}

\begin{proof}
% P01-E0645: the frozen one-sentence proof omits its subject.
这直接由引理 \ref{algebra-lemma-mod-injective-general} 得出。
\end{proof}

\begin{lemma}
\label{algebra-lemma-grothendieck-regular-sequence-general}
假设 $R \to S$ 是局部环的局部环同态。
以 $\mathfrak m$ 表示 $R$ 的极大理想。假设
\begin{enumerate}
\item $R \to S$ 本质上是有限表示的，
\item $R \to S$ 是平坦的，以及
\item $f_1, \ldots, f_c$ 是 $S$ 中的元列，且这些像
$\overline{f}_1, \ldots, \overline{f}_c$ 在 $S/{\mathfrak m}S$ 中构成正则列。
\end{enumerate}
则 $f_1, \ldots, f_c$ 是 $S$ 中的正则列，而且每个商
$S/(f_1, \ldots, f_i)$ 在 $R$ 上平坦。
\end{lemma}

\begin{proof}
% P01-E0646: the frozen proof is the noun phrase “Induction and Lemma ...”.
对列长作归纳，并应用引理 \ref{algebra-lemma-grothendieck-general}。
\end{proof}

\noindent
下面给出平坦性局部准则在局部有限表示的局部环映射情形的版本。

\begin{lemma}
\label{algebra-lemma-variant-local-criterion-flatness-general}
设 $R \to S$ 是局部环的局部同态。
设 $I \not = R$ 是 $R$ 中的理想，并设 $M$ 是 $S$-模。假设
\begin{enumerate}
\item $S$ 在 $R$ 上本质有限表示，
\item $M$ 在 $S$ 上有限表示，
\item $\text{Tor}_1^R(M, R/I) = 0$，以及
\item $M/IM$ 在 $R/I$ 上平坦。
\end{enumerate}
则 $M$ 在 $R$ 上平坦。
\end{lemma}

\begin{proof}
设 $\Lambda$、$R_\lambda \to S_\lambda$、$M_\lambda$ 如引理
\ref{algebra-lemma-limit-module-essentially-finite-presentation} 中所述。
% P01-E0647: the frozen declaration “Denote I_lambda ... the inverse image” omits “by/as”.
以 $I_\lambda \subset R_\lambda$ 表示 $I$ 的逆像。
% P01-E0648: the frozen approximation data uses R_lambda rather than R_lambda/I_lambda as the source of the quotient map.
此时，下列目标与系统数据
$R/I \to S/IS$、$M/IM$、$R_\lambda \to S_\lambda/I_\lambda S_\lambda$
以及 $M_\lambda/I_\lambda M_\lambda$ 也满足引理
\ref{algebra-lemma-limit-module-essentially-finite-presentation} 的各项结论。
因此由引理 \ref{algebra-lemma-colimit-eventually-flat}，可以假设（缩小索引集
$\Lambda$ 之后）对所有 $\lambda$，$M_\lambda/I_\lambda M_\lambda$ 都是平坦的。
% P01-E0649: the frozen statement does not name the base over which this quotient module is flat.
选取某个 $\lambda$，并考虑
$$
\text{Tor}_1^{R_\lambda}(M_\lambda, R_\lambda/I_\lambda)
=
\Ker(I_\lambda \otimes_{R_\lambda} M_\lambda
\to M_\lambda).
$$
见备注 \ref{algebra-remark-Tor-ring-mod-ideal}。右端表明这是有限生成
$S_\lambda$-模（因为 $S_\lambda$ 是 Noether 的，而所涉模都是有限的）。
设 $\xi_1, \ldots, \xi_n$ 是生成元。
由于 $\text{Tor}_1^R(M, R/I) = 0$，并且 $\otimes$ 与余极限可换，
可知存在 $\lambda' \geq \lambda$，使得每个 $\xi_i$ 都在
$\text{Tor}_1^{R_{\lambda'}}(M_{\lambda'}, R_{\lambda'}/I_{\lambda'})$
中映到零。以下各映射的复合
$$
\xymatrix{
R_{\lambda'} \otimes_{R_\lambda}
\text{Tor}_1^{R_\lambda}(M_\lambda, R_\lambda/I_\lambda)
\ar[d]^{\text{surjective by Lemma \ref{algebra-lemma-surjective-on-tor-one}}} \\
\text{Tor}_1^{R_\lambda}(M_\lambda, R_{\lambda'}/I_\lambda R_{\lambda'})
\ar[d]^{\text{surjective up to localization by
Lemma \ref{algebra-lemma-surjective-on-tor-one-trivial}}} \\
\text{Tor}_1^{R_{\lambda'}}(M_{\lambda'}, R_{\lambda'}/I_\lambda R_{\lambda'})
\ar[d]^{\text{surjective by Lemma \ref{algebra-lemma-surjective-on-tor-one}}} \\
\text{Tor}_1^{R_{\lambda'}}(M_{\lambda'}, R_{\lambda'}/I_{\lambda'}).
}
$$
由所示理由，在局部化之后是满射。
这里需要局部化，是因为 $M_{\lambda'}$ 不等于
$M_\lambda \otimes_{R_\lambda} R_{\lambda'}$；事实上，它等于
$M_\lambda \otimes_{S_\lambda} S_{\lambda'}$，且 $S_{\lambda'}$
是 $S_{\lambda} \otimes_{R_\lambda} R_{\lambda'}$ 的局部化，
% P01-E0650: the frozen clause “whence the statement up to a localization” is grammatically malformed.
因而得到局部化之后的断言（或等价地，与 $S_{\lambda'}$ 作张量积之后的断言）。
注意，引理 \ref{algebra-lemma-surjective-on-tor-one} 可用于第一支和第三支箭头，
因为 $M_\lambda/I_\lambda M_\lambda$ 在 $R_\lambda/I_\lambda$ 上平坦，
而 $M_{\lambda'}/I_\lambda M_{\lambda'}$ 在
$R_{\lambda'}/I_\lambda R_{\lambda'}$ 上平坦，后者是平坦模
$M_\lambda/I_\lambda M_\lambda$ 的基变换。
如上所述，该复合把生成元 $\xi_i$ 映到零。
最后得出
$\text{Tor}_1^{R_{\lambda'}}(M_{\lambda'}, R_{\lambda'}/I_{\lambda'})$
为零。由引理 \ref{algebra-lemma-variant-local-criterion-flatness}，这蕴含
$M_{\lambda'}$ 在 $R_{\lambda'}$ 上平坦。
\end{proof}

\noindent
请将下面的引理与引理 \ref{algebra-lemma-criterion-flatness-fibre-Noetherian}
（Noether 局部环的情形）及引理
\ref{algebra-lemma-criterion-flatness-fibre-nilpotent}
（基环中幂零理想的情形）比较。

\begin{lemma}[纤维平坦性准则]
\label{algebra-lemma-criterion-flatness-fibre}
设 $R$、$S$、$S'$ 是局部环，并设 $R \to S \to S'$ 是局部环同态。
设 $M$ 是 $S'$-模。设 $\mathfrak m \subset R$ 是极大理想。假设
\begin{enumerate}
\item 环映射 $R \to S$ 和 $R \to S'$ 本质上都是有限表示的。
\item 模 $M$ 在 $S'$ 上有限表示。
\item 模 $M$ 非零。
\item 模 $M/\mathfrak mM$ 是平坦的 $S/\mathfrak mS$-模。
\item 模 $M$ 是平坦的 $R$-模。
\end{enumerate}
则 $S$ 在 $R$ 上平坦，且 $M$ 是平坦的 $S$-模。
\end{lemma}

\begin{proof}
如引理 \ref{algebra-lemma-limit-essentially-finite-presentation} 的证明中那样，
首先可将 $R = \colim R_\lambda$ 写成局部有限型
$\mathbf{Z}$-代数的有向余极限。
% P01-E0651: the frozen declaration “Denote p_lambda the maximal ideal” omits “by/as”.
以 $\mathfrak p_\lambda$ 表示 $R_\lambda$ 的极大理想。
其次，可以假设对某个 $\lambda_1 \in \Lambda$，存在
$f_{j, \lambda_1} \in R_{\lambda_1}[x_1, \ldots, x_n]$，使得
$$
S =
\colim_{\lambda \geq \lambda_1} S_\lambda, \text{ 其中 }
S_\lambda =
(R_\lambda[x_1, \ldots, x_n]/
(f_{1, \lambda}, \ldots, f_{u, \lambda}))_{\mathfrak q_\lambda}
$$
对某个 $\lambda_2 \in \Lambda$，$\lambda_2 \geq \lambda_1$，存在
$g_{j, \lambda_2} \in R_{\lambda_2}[x_1, \ldots, x_n, y_1, \ldots, y_m]$，
其像为
$\overline{g}_{j, \lambda_2} \in S_{\lambda_2}[y_1, \ldots, y_m]$，
并且
$$
S' =
\colim_{\lambda \geq \lambda_2} S'_\lambda, \text{ 其中 }
S'_\lambda =
(S_\lambda[y_1, \ldots, y_m]/
(\overline{g}_{1, \lambda}, \ldots,
\overline{g}_{v, \lambda}))_{\overline{\mathfrak q}'_\lambda}
$$
注意，这也蕴含
% P01-E0652: the frozen R_lambda-presentation omits the defining f-relations of S_lambda.
$$
S'_\lambda =
(R_\lambda[x_1, \ldots, x_n, y_1, \ldots, y_m]/
(g_{1, \lambda}, \ldots, g_{v, \lambda}))_{\mathfrak q'_\lambda}
$$
选取 $M$ 在 $S'$ 上的一个表示
$$
(S')^{\oplus s} \to (S')^{\oplus t} \to M \to 0
$$
设 $A \in \text{Mat}(t \times s, S')$ 是该表示的矩阵。
对某个 $\lambda_3 \in \Lambda$，$\lambda_3 \geq \lambda_2$，可以找到映到
$A$ 的矩阵
% P01-E0653: the frozen approximating presentation matrix is placed over S_lambda rather than S'_lambda.
$A_{\lambda_3} \in \text{Mat}(t \times s, S_{\lambda_3})$。
对所有 $\lambda \geq \lambda_3$，令
$M_\lambda = \Coker((S'_\lambda)^{\oplus s} \xrightarrow{A_\lambda}
(S'_\lambda)^{\oplus t})$。

\medskip\noindent
按这些选取，对每个 $\lambda_3 \leq \lambda \leq \mu$，
$S_\lambda \otimes_{R_{\lambda}} R_\mu \to S_\mu$ 是局部化，
$S'_\lambda \otimes_{S_{\lambda}} S_\mu \to S'_\mu$ 是局部化，且映射
$M_\lambda \otimes_{S'_\lambda} S'_\mu \to M_\mu$ 是同构。
这还蕴含 $S'_\lambda \otimes_{R_{\lambda}} R_\mu \to S'_\mu$ 是局部化。
因此，由于 $M$ 在 $R$ 上平坦，由引理
\ref{algebra-lemma-colimit-eventually-flat} 可知，对所有充分大的 $\lambda$，
模 $M_\lambda$ 在 $R_\lambda$ 上平坦。
此外，注意
$
\mathfrak m = \colim \mathfrak p_\lambda
$、
$
S/\mathfrak mS = \colim S_\lambda/\mathfrak p_\lambda S_\lambda
$、
$
S'/\mathfrak mS' = \colim S'_\lambda/\mathfrak p_\lambda S'_\lambda
$，并且
$
M/\mathfrak mM = \colim M_\lambda/\mathfrak p_\lambda M_\lambda
$。另外，对每个 $\lambda_3 \leq \lambda \leq \mu$，由上列性质可知
$$
S'_\lambda/\mathfrak p_\lambda S'_\lambda
\otimes_{S_{\lambda}/\mathfrak p_\lambda S_\lambda}
S_\mu/\mathfrak p_\mu S_\mu
\longrightarrow
S'_\mu/\mathfrak p_\mu S'_\mu
$$
是局部化，而映射
$$
M_\lambda / \mathfrak p_\lambda M_\lambda
\otimes_{S'_\lambda/\mathfrak p_\lambda S'_\lambda}
S'_\mu /\mathfrak p_\mu S'_\mu
\longrightarrow
M_\mu/\mathfrak p_\mu M_\mu
$$
是同构。因此，系统
$(S_\lambda/\mathfrak p_\lambda S_\lambda \to
S'_\lambda/\mathfrak p_\lambda S'_\lambda,
M_\lambda/\mathfrak p_\lambda M_\lambda)$
也如引理 \ref{algebra-lemma-limit-module-essentially-finite-presentation} 中所述。
由于假设 $M/\mathfrak m M$ 在 $S/\mathfrak mS$ 上平坦，
可以再次应用引理 \ref{algebra-lemma-colimit-eventually-flat}，得到对所有充分大的
$\lambda$，$M_\lambda/\mathfrak p_\lambda M_\lambda$ 在
$S_\lambda/\mathfrak p_\lambda S_\lambda$ 上平坦。
因此，对充分大的 $\lambda$，数据
$R_\lambda \to S_\lambda \to S'_\lambda, M_\lambda$ 满足引理
\ref{algebra-lemma-criterion-flatness-fibre-Noetherian} 的假设。
选取这样的 $\lambda$。于是
% P01-E0654: the frozen proof asserts equality with the unlocalized base change although the constructed transition map is only a localization.
$S = S_\lambda \otimes_{R_\lambda} R$ 在 $R$ 上平坦，并且
$M = M_\lambda \otimes_{S_\lambda} S$ 在 $S$ 上平坦
（因为平坦模的基变换仍平坦）。
\end{proof}

\noindent
下面的结论容易从“纤维平坦性准则”引理
\ref{algebra-lemma-criterion-flatness-fibre} 推出。关于此类结论的更多内容，
见“更多平坦性”第 \ref{flat-section-introduction} 节。

\begin{lemma}
\label{algebra-lemma-criterion-flatness-fibre-fp-over-ft}
设 $R$、$S$、$S'$ 是局部环，并设 $R \to S \to S'$ 是局部环同态。
设 $M$ 是 $S'$-模。设 $\mathfrak m \subset R$ 是极大理想。假设
\begin{enumerate}
\item $R \to S'$ 本质上是有限表示的，
\item $R \to S$ 本质上是有限型的，
\item $M$ 在 $S'$ 上有限表示，
\item $M$ 非零，
\item $M/\mathfrak mM$ 是平坦的 $S/\mathfrak mS$-模，以及
\item $M$ 是平坦的 $R$-模。
\end{enumerate}
则 $S$ 在 $R$ 上本质有限表示且平坦，并且 $M$ 是平坦的 $S$-模。
\end{lemma}

\begin{proof}
% P01-E0655: the frozen proof assumes the finite-presentation conclusion instead of the finite-type hypothesis.
由于 $S$ 在 $R$ 上本质有限表示，可对某个有限型 $R$-代数 $C$ 写成
$S = C_{\overline{\mathfrak q}}$。
写 $C = R[x_1, \ldots, x_n]/I$。
% P01-E0656: the frozen declaration “Denote q ... the prime ideal” omits “by/as”.
以 $\mathfrak q \subset R[x_1, \ldots, x_n]$ 表示对应于
$\overline{\mathfrak q}$ 的素理想。于是 $S = B/J$，其中
$B = R[x_1, \ldots, x_n]_{\mathfrak q}$ 在 $R$ 上本质有限表示，且
$J = IB$。可以找到 $f_1, \ldots, f_k \in J$，使得像
$\overline{f}_i \in B/\mathfrak mB$ 生成 $J$ 在 Noether 环
$B/\mathfrak mB$ 中的像 $\overline{J}$。因此存在有限生成理想
% P01-E0657: the frozen plural existential is paired with the single symbol J'.
$J' \subset J$，使得 $B/J' \to B/J$ 诱导同构
$$
(B/J') \otimes_R R/\mathfrak m \longrightarrow
B/J \otimes_R R/\mathfrak m = S/\mathfrak mS.
$$
对任意上述 $J'$，引理 \ref{algebra-lemma-criterion-flatness-fibre} 可用于环映射
$$
R \longrightarrow B/J' \longrightarrow S'
$$
和模 $M$。因此，对上述任意选取的 $J'$，$B/J'$ 在 $R$ 上平坦。
现在，若
% P01-E0658: the frozen chain repeats J' where the middle ideal must be J''.
$J' \subset J' \subset J$ 是两个如上的有限生成理想，则可知
$B/J' \to B/J''$ 是本质有限表示的平坦 $R$-代数之间的满射，
% P01-E0659: the frozen relative-clause chain has two “which” clauses with conflicting antecedents.
而且模掉 $\mathfrak m$ 后它是同构。因此，引理
\ref{algebra-lemma-mod-injective-general} 蕴含 $B/J' = B/J''$，即
$J' = J''$。这显然意味着 $J$ 是有限生成的，也就是说，
$S$ 在 $R$ 上本质有限表示。因此可以将引理
\ref{algebra-lemma-criterion-flatness-fibre} 用于 $R \to S \to S'$，结论得证。
\end{proof}

\begin{lemma}[纤维平坦性准则：局部幂零情形]
\label{algebra-lemma-criterion-flatness-fibre-locally-nilpotent}
设
$$
\xymatrix{
S \ar[rr] & & S' \\
& R \ar[lu] \ar[ru]
}
$$
是环范畴中的交换图。设 $I \subset R$ 是局部幂零理想，
$M$ 是 $S'$-模。假设
\begin{enumerate}
\item $R \to S$ 是有限型的，
\item $R \to S'$ 是有限表示的，
\item $M$ 是有限表示的 $S'$-模，
\item $M/IM$ 作为 $S/IS$-模是平坦的，以及
\item $M$ 作为 $R$-模是平坦的。
\end{enumerate}
则 $M$ 是平坦的 $S$-模，并且对每个满足
$M \otimes_S \kappa(\mathfrak q)$ 非零的 $\mathfrak q \subset S$，
$S_\mathfrak q$ 在 $R$ 上平坦且本质有限表示。
\end{lemma}

\begin{proof}
如果 $M \otimes_S \kappa(\mathfrak q)$ 非零，那么
$S' \otimes_S \kappa(\mathfrak q)$ 非零，因而存在位于
$\mathfrak q$ 上方的素理想 $\mathfrak q' \subset S'$（引理
\ref{algebra-lemma-in-image}）。设 $\mathfrak p \subset R$ 是
$\mathfrak q$ 在 $\Spec(R)$ 中的像。由于 $I$ 局部幂零，
$I \subset \mathfrak p$，故 $M/\mathfrak p M$ 在
$S/\mathfrak pS$ 上平坦。因此，可以把引理
\ref{algebra-lemma-criterion-flatness-fibre-fp-over-ft} 用于
$R_\mathfrak p \to S_\mathfrak q \to S'_{\mathfrak q'}$
和 $M_{\mathfrak q'}$。得出 $M_{\mathfrak q'}$ 在 $S$ 上平坦，且
$S_\mathfrak q$ 在 $R$ 上平坦并本质有限表示。
由于 $\mathfrak q'$ 是 $S'$ 的任意素理想，还可知 $M$ 在 $S$ 上平坦
（引理 \ref{algebra-lemma-flat-localization}）。
\end{proof}

\section{平坦轨迹的开性}
\label{algebra-section-open-flat}

\noindent
我们使用引理 \ref{algebra-lemma-colimit-eventually-flat} 化约到 Noether 情形。
Noether 情形则用第 \ref{algebra-section-complex-exact} 节给出的正合复形刻画来处理。

\begin{lemma}
\label{algebra-lemma-CM-dim-finite-type}
设 $k$ 是域，$S$ 是有限型 $k$-代数。设 $f_1, \ldots, f_i$ 是
$S$ 中的元。假设 $S$ 是 Cohen--Macaulay 的、等维且维数为 $d$，并且
$\dim V(f_1, \ldots, f_i) \leq d - i$。则等号成立，而且
% P01-E0660: the frozen plural subject “f_1, ..., f_i” takes singular “forms”.
$f_1, \ldots, f_i$ 对 $V(f_1, \ldots, f_i)$ 的每个素理想
$\mathfrak q$，都在 $S_{\mathfrak q}$ 中构成正则列。
% P01-E0661: the frozen phrase “prime of V(...)” uses the wrong point/set preposition.
\end{lemma}

\begin{proof}
如果 $S$ 是维数为 $d$ 的 Cohen--Macaulay 等维环，则对 $S$ 的所有极大理想
$\mathfrak m$ 都有 $\dim(S_{\mathfrak m}) = d$；见引理
\ref{algebra-lemma-disjoint-decomposition-CM-algebra}。
由命题 \ref{algebra-proposition-CM-module} 可知，对所有极大理想
$\mathfrak m \in V(f_1, \ldots, f_i)$，该序列在 $S_{\mathfrak m}$ 中是正则列，
而局部环 $S_{\mathfrak m}/(f_1, \ldots, f_i)$ 是维数为 $d - i$ 的
Cohen--Macaulay 环。这实际上意味着 $S/(f_1, \ldots, f_i)$ 是维数为
$d - i$ 的 Cohen--Macaulay 等维环。
\end{proof}

\begin{lemma}
\label{algebra-lemma-open-regular-sequence}
设 $R \to S$ 是有限型环映射。设 $d$ 是整数，使得所有纤维
$S \otimes_R \kappa(\mathfrak p)$ 都是维数为 $d$ 的 Cohen--Macaulay 等维环。
设 $f_1, \ldots, f_i$ 是 $S$ 中的元。集合
$$
\{ \mathfrak q \in V(f_1, \ldots, f_i)
\mid f_1, \ldots, f_i
\text{ are a regular sequence in }
S_{\mathfrak q}/\mathfrak p S_{\mathfrak q}
\text{ 其中 }\mathfrak p = R \cap \mathfrak q
\}
$$
在 $V(f_1, \ldots, f_i)$ 中是开的。
\end{lemma}

\begin{proof}
写 $\overline{S} = S/(f_1, \ldots, f_i)$。
假设 $\mathfrak q$ 是引理中所定义集合的元，$\mathfrak p$ 是 $R$ 中相应的素理想。
我们将使用定义 \ref{algebra-definition-relative-dimension} 中定义的相对维数。
首先，由引理 \ref{algebra-lemma-dimension-at-a-point-finite-type-field}，注意到
$d = \dim_{\mathfrak q}(S/R) =
\dim(S_{\mathfrak q}/\mathfrak pS_{\mathfrak q}) +
\text{trdeg}_{\kappa(\mathfrak p)}\ \kappa(\mathfrak q)$。
由于 $f_1, \ldots, f_i$ 在 Noether 局部环
$S_{\mathfrak q}/\mathfrak pS_{\mathfrak q}$ 中构成正则列，
引理 \ref{algebra-lemma-one-equation} 告诉我们
$\dim(\overline{S}_{\mathfrak q}/\mathfrak p\overline{S}_{\mathfrak q})
= \dim(S_{\mathfrak q}/\mathfrak pS_{\mathfrak q}) - i$。
由引理 \ref{algebra-lemma-dimension-at-a-point-finite-type-field} 得出
$\dim_{\mathfrak q}(\overline{S}/R)
= \dim(\overline{S}_{\mathfrak q}/\mathfrak p\overline{S}_{\mathfrak q}) +
\text{trdeg}_{\kappa(\mathfrak p)}\ \kappa(\mathfrak q)
= d - i$。由引理 \ref{algebra-lemma-dimension-fibres-bounded-open-upstairs}，
在 $\mathfrak q$ 的一个邻域中，对所有
$\mathfrak q' \in V(f_1, \ldots, f_i) = \Spec(\overline{S})$，都有
$\dim_{\mathfrak q'}(\overline{S}/R) \leq d - i$。
因此，用某个 $g \in S$、$g \not \in \mathfrak q$ 的 $S_g$ 代替 $S$ 后，
可以假设该不等式对所有 $\mathfrak q'$ 都成立。
结论由引理 \ref{algebra-lemma-CM-dim-finite-type} 得出。
\end{proof}

\begin{lemma}
\label{algebra-lemma-exact-on-fibres-open}
设 $R \to S$ 是环映射。考虑由有限自由 $S$-模组成的有限同调复形：
$$
F_{\bullet} :
0
\to
S^{n_e}
\xrightarrow{\varphi_e}
S^{n_{e-1}}
\xrightarrow{\varphi_{e-1}}
\ldots
\xrightarrow{\varphi_{i + 1}}
S^{n_i}
\xrightarrow{\varphi_i}
S^{n_{i-1}}
\xrightarrow{\varphi_{i-1}}
\ldots
\xrightarrow{\varphi_1}
S^{n_0}
$$
对 $S$ 的每个素理想 $\mathfrak q$，考虑复形
$\overline{F}_{\bullet, \mathfrak q} =
F_{\bullet, \mathfrak q} \otimes_R \kappa(\mathfrak p)$，其中
$\mathfrak p$ 是 $\mathfrak q$ 在 $R$ 中的逆像。
% P01-E0662: the frozen hypothesis omits coordinated finite verbs around “finite type, flat with fibres ... Cohen-Macaulay”.
假设 $R$ 是 Noether 的，并且存在整数 $d$，使得 $R \to S$ 是有限型、平坦的，
其纤维 $S \otimes_R \kappa(\mathfrak p)$ 是维数为 $d$ 的 Cohen--Macaulay 环。
集合
$$
\{\mathfrak q \in \Spec(S) \mid
\overline{F}_{\bullet, \mathfrak q}\text{ is exact}\}
$$
在 $\Spec(S)$ 中是开的。
\end{lemma}

\begin{proof}
设 $\mathfrak q$ 是引理所定义集合的元。
我们将使用命题 \ref{algebra-proposition-what-exact} 来证明：存在
$g \in S$、$g \not \in \mathfrak q$，使得 $D(g)$ 包含于引理所定义的集合中。
换言之，我们将证明：用 $S_g$ 代替 $S$ 后，引理中的集合是整个
$\Spec(S)$。因此，在证明过程中，我们会有限多次地用这样的局部化代替 $S$。
回顾命题 \ref{algebra-proposition-what-exact}：在环为局部环时，它用映射
$\varphi_i$ 的秩和理想 $I(\varphi_i)$ 刻画复形的正合性。
我们先处理秩条件。令
$r_i = n_i - n_{i + 1} + \ldots + (-1)^{e - i} n_e$。
注意 $r_i + r_{i + 1} = n_i$，并且在正合情形下，$r_i$ 是
$\varphi_i$ 的预期秩。

\medskip\noindent
由引理 \ref{algebra-lemma-complex-exact-mod} 可知，如果
$\overline{F}_{\bullet, \mathfrak q}$ 正合，则局部化
$F_{\bullet, \mathfrak q}$ 正合。特别地，由某个
$g \in S$、$g \not \in \mathfrak q$ 局部化后，复形 $F_\bullet$ 变为正合。
在这种情况下，把命题 \ref{algebra-proposition-what-exact} 用于 $S$ 在所有素理想处的局部化，
可知 $\varphi_i$ 的所有 $(r_i + 1) \times (r_i + 1)$-子式均为零。
因此 $\varphi_i$ 的秩至多为 $r_i$。

\medskip\noindent
设 $I_i \subset S$ 表示由 $\varphi_i$ 的矩阵的
$r_i \times r_i$-子式生成的理想。由命题
\ref{algebra-proposition-what-exact}，复形
$\overline{F}_{\bullet, \mathfrak q}$ 正合，当且仅当对每个
$1 \leq i \leq e$，或者 $(I_i)_{\mathfrak q} = S_{\mathfrak q}$，
或者 $(I_i)_{\mathfrak q}$ 包含一个长度为 $i$ 的
$S_{\mathfrak q}/\mathfrak p S_{\mathfrak q}$-正则列。
事实上，由上面对 $r_i$ 的选取以及对 $\varphi_i$ 秩的界，
这是命题 \ref{algebra-proposition-what-exact} 的条件得到满足的唯一方式。

\medskip\noindent
如果 $(I_i)_{\mathfrak q} = S_{\mathfrak q}$，则在某个
$g \not\in \mathfrak q$ 处局部化 $S$ 后，可以假设 $I_i = S$。
显然，这是一个开条件。

\medskip\noindent
如果 $(I_i)_{\mathfrak q} \not = S_{\mathfrak q}$，则有序列
% P01-E0663: the frozen proof chooses f_j in the localized ideal but then treats them as global elements at every prime without explicitly clearing denominators.
$f_1, \ldots, f_i \in (I_i)_{\mathfrak q}$，它们在
$S_{\mathfrak q}/\mathfrak pS_{\mathfrak q}$ 中构成正则列。
注意，对任意满足 $(f_1, \ldots, f_i) \not \subset \mathfrak q'$ 的素理想
$\mathfrak q' \subset S$，都有 $(I_i)_{\mathfrak q'} = S_{\mathfrak q'}$。
因此结论由引理 \ref{algebra-lemma-open-regular-sequence} 得出。
\end{proof}


\begin{theorem}
\label{algebra-theorem-openness-flatness}
设 $R$ 是环，$R \to S$ 是有限表示的环映射，$M$ 是有限表示的 $S$-模。
集合
$$
\{ \mathfrak q \in \Spec(S) \mid
M_{\mathfrak q}\text{ is flat over }R\}
$$
在 $\Spec(S)$ 中是开的。
\end{theorem}

\begin{proof}
设 $\mathfrak q \in \Spec(S)$ 是素理想，$\mathfrak p \subset R$ 是
$\mathfrak q$ 在 $R$ 中的逆像。注意，$M_{\mathfrak q}$ 在 $R$ 上平坦，
当且仅当它在 $R_{\mathfrak p}$ 上平坦。
假设 $M_{\mathfrak q}$ 在 $R$ 上平坦。我们断言：存在
$g \in S$、$g \not \in \mathfrak q$，使得 $M_g$ 在 $R$ 上平坦。

\medskip\noindent
首先化约到 $R$ 和 $S$ 都是在 $\mathbf{Z}$ 上有限型的情形。
选取有向集 $\Lambda$ 以及引理
\ref{algebra-lemma-limit-module-finite-presentation} 中的系统
$(R_\lambda \to S_\lambda, M_\lambda)$。
令 $\mathfrak p_\lambda$ 等于 $\mathfrak p$ 在 $R_\lambda$ 中的逆像，
并令 $\mathfrak q_\lambda$ 等于 $\mathfrak q$ 在 $S_\lambda$ 中的逆像。
则系统
$$
((R_\lambda)_{\mathfrak p_\lambda},
(S_\lambda)_{\mathfrak q_\lambda},
(M_\lambda)_{\mathfrak q_{\lambda}})
$$
如引理 \ref{algebra-lemma-limit-module-essentially-finite-presentation} 中所述。
因此由引理 \ref{algebra-lemma-colimit-eventually-flat} 可知，对某个 $\lambda$，
模 $M_\lambda$ 在素理想 $\mathfrak q_{\lambda}$ 处在 $R_\lambda$ 上平坦。
假设可以对系统
$(R_\lambda \to S_\lambda, M_\lambda, \mathfrak q_{\lambda})$
证明我们的断言。换言之，假设可以找到
$g \in S_\lambda$、$g \not\in \mathfrak q_\lambda$，使得
$(M_\lambda)_g$ 在 $R_\lambda$ 上平坦。
由引理 \ref{algebra-lemma-limit-module-finite-presentation}，有
$M = M_\lambda \otimes_{R_\lambda} R$，因而还有
$M_g = (M_\lambda)_g \otimes_{R_\lambda} R$。
于是由引理 \ref{algebra-lemma-flat-base-change}，得出系统
$(R \to S, M, \mathfrak q)$ 的断言。

\medskip\noindent
至此，可以假设 $R$ 和 $S$ 都在 $\mathbf{Z}$ 上有限型。
可以把 $S$ 写成多项式环 $R[x_1, \ldots, x_n]$ 的商。
当然，可以用 $R[x_1, \ldots, x_n]$ 代替 $S$，并假设 $S$ 是
$R$ 上的多项式环。特别地，$R \to S$ 是平坦的，而且所有纤维环
% P01-E0664: the frozen phrase “all fibres rings” has the wrong plural modifier.
$S \otimes_R \kappa(\mathfrak p)$ 的整体维数都是 $n$。

\medskip\noindent
选取 $M$ 在 $S$ 上的分解 $F_\bullet$，其中每个 $F_i$ 都有限自由；
见引理 \ref{algebra-lemma-resolution-by-finite-free}。
令 $K_n = \Ker(F_{n-1} \to F_{n-2})$。
注意，$(K_n)_{\mathfrak q}$ 在 $R$ 上平坦，因为每个 $F_i$ 在 $R$ 上平坦，
并且对 $M$ 的假设成立；见引理 \ref{algebra-lemma-flat-ses}。
此外，序列
$$
0 \to
K_n/\mathfrak p K_n \to
F_{n-1}/ \mathfrak p F_{n-1} \to
\ldots \to
F_0 / \mathfrak p F_0 \to
M/\mathfrak p M \to
0
$$
在 $\mathfrak q$ 处局部化后是正合的，这是因为
$\text{Tor}_i^{R_\mathfrak p}(\kappa(\mathfrak p), M_{\mathfrak q})$ 消失。
% P01-E0665: the frozen proof says the localized polynomial fibre has global dimension exactly n; it is only bounded above by n in general.
由于 $S_\mathfrak q/\mathfrak p S_{\mathfrak q}$ 的整体维数是 $n$，
可知 $K_n / \mathfrak p K_n$ 在 $\mathfrak q$ 处的局部化是
$S_\mathfrak q/\mathfrak p S_{\mathfrak q}$ 上的有限自由模。
由引理 \ref{algebra-lemma-free-fibre-flat-free}，$(K_n)_{\mathfrak q}$ 在
$S_{\mathfrak q}$ 上自由。特别地，存在
$g \in S$、$g \not \in \mathfrak q$，使得 $(K_n)_g$ 在 $S_g$ 上有限自由。

\medskip\noindent
由引理 \ref{algebra-lemma-exact-on-fibres-open}，存在进一步的局部化 $S_g$，
使得复形
$$
0 \to K_n \to F_{n-1} \to \ldots \to F_0
$$
在 $R \to S$ 的{\it 所有纤维}上正合。
由引理 \ref{algebra-lemma-complex-exact-mod}，这蕴含 $F_1 \to F_0$ 的余核平坦。
这就在 Noether 情形中证明了该定理。
\end{proof}

\section{Cohen--Macaulay 轨迹的开性}
\label{algebra-section-CM-open}

\noindent
本节用平坦性刻画有限型代数的 Cohen--Macaulay 性质。
% P01-E0666: the frozen introduction omits “that” before the object clause “the set ... is open”.
然后利用这一刻画证明：这样的代数为 Cohen--Macaulay 的点集是开的。

\begin{lemma}
\label{algebra-lemma-where-CM}
设 $S$ 是域 $k$ 上的有限型代数。设
$\varphi : k[y_1, \ldots, y_d] \to S$ 是拟有限环映射。
作为 $\Spec(S)$ 的子集，有
% P01-E0667: the frozen left-hand predicate omits “is” before “flat”.
% P01-E0668: the frozen right-hand predicate omits “is” before “CM”.
$$
\{ \mathfrak q \mid
S_{\mathfrak q} \text{ flat over }k[y_1, \ldots, y_d]\}
=
\{ \mathfrak q \mid
S_{\mathfrak q} \text{ CM and }\dim_{\mathfrak q}(S/k) = d\}
$$
记号见定义 \ref{algebra-definition-relative-dimension}。
\end{lemma}

\begin{proof}
设 $\mathfrak q \subset S$ 是素理想。
% P01-E0669: the frozen “Denote p = ...” construction lacks a grammatical complement.
记 $\mathfrak p = k[y_1, \ldots, y_d] \cap \mathfrak q$。
注意，例如由引理 \ref{algebra-lemma-dimension-inequality-quasi-finite}，总有
$\dim(S_{\mathfrak q}) \leq \dim(k[y_1, \ldots, y_d]_{\mathfrak p})$。
此外，域扩张 $\kappa(\mathfrak q)/\kappa(\mathfrak p)$ 是有限的，因而
$\text{trdeg}_k(\kappa(\mathfrak p)) = \text{trdeg}_k(\kappa(\mathfrak q))$。

\medskip\noindent
设 $\mathfrak q$ 是左端集合的元。此时引理
\ref{algebra-lemma-finite-flat-over-regular-CM} 适用，得到
$S_{\mathfrak q}$ 是 Cohen--Macaulay 的，并且
$\dim(S_{\mathfrak q}) = \dim(k[y_1, \ldots, y_d]_{\mathfrak p})$。
结合上面超越次数的等式及引理
\ref{algebra-lemma-dimension-at-a-point-finite-type-field}，这蕴含
$\dim_{\mathfrak q}(S/k) = d$。因此 $\mathfrak q$ 是右端集合的元。

\medskip\noindent
设 $\mathfrak q$ 是右端集合的元。由上面超越次数的等式、假设
$\dim_{\mathfrak q}(S/k) = d$ 及引理
\ref{algebra-lemma-dimension-at-a-point-finite-type-field}，得出
$\dim(S_{\mathfrak q}) = \dim(k[y_1, \ldots, y_d]_{\mathfrak p})$。
因此引理 \ref{algebra-lemma-CM-over-regular-flat} 适用，从而
$\mathfrak q$ 是左端集合的元。
\end{proof}

\begin{lemma}
\label{algebra-lemma-finite-type-over-field-CM-open}
设 $S$ 是域 $k$ 上的有限型代数。使得 $S_{\mathfrak q}$ 为
Cohen--Macaulay 的素理想 $\mathfrak q$ 的集合在 $S$ 中是开的。
% P01-E0670: the frozen statement says “open in S” although the locus is a subset of Spec(S).
\end{lemma}

\noindent
该引理是下方引理
\ref{algebra-lemma-finite-presentation-flat-CM-locus-open} 的特殊情形，
所以若愿意，可以直接跳到该引理的证明。

\begin{proof}
设 $\mathfrak q \subset S$ 是使得 $S_{\mathfrak q}$ 为
Cohen--Macaulay 的素理想。必须证明存在
$g \in S$、$g \not \in \mathfrak q$，使得环 $S_g$ 是 Cohen--Macaulay 的。
对任意 $g \in S$、$g \not \in \mathfrak q$，可以用 $S_g$ 代替 $S$，
并用 $\mathfrak qS_g$ 代替 $\mathfrak q$。结合引理
\ref{algebra-lemma-Noether-normalization-at-point} 和
\ref{algebra-lemma-dimension-at-a-point-finite-type-field}，可以假设存在有限单射环映射
$k[y_1, \ldots, y_d] \to S$，其中
$d = \dim(S_{\mathfrak q}) + \text{trdeg}_k(\kappa(\mathfrak q))$。
令 $\mathfrak p = k[y_1, \ldots, y_d] \cap \mathfrak q$。
按构造，$\mathfrak q$ 是引理 \ref{algebra-lemma-where-CM} 所示等式右端集合的元，
因而也是左端集合的元。

\medskip\noindent
由定理 \ref{algebra-theorem-openness-flatness}，可知对某个
$g \in S$、$g \not \in \mathfrak q$，环 $S_g$ 在
$k[y_1, \ldots, y_d]$ 上平坦。因此再次由引理
\ref{algebra-lemma-where-CM} 的等式，得出 $S_g$ 的所有局部环都如所需那样是
Cohen--Macaulay 的。
\end{proof}

\begin{lemma}
\label{algebra-lemma-generic-CM}
设 $k$ 是域，$S$ 是有限型 $k$ 代数。
% P01-E0671: the frozen compound “finite type k algebra” lacks standard hyphenation.
Cohen--Macaulay 素理想的集合构成稠密开集 $U \subset \Spec(S)$。
\end{lemma}

\begin{proof}
由引理 \ref{algebra-lemma-finite-type-over-field-CM-open}，该集合是开的。
它包含所有极小素理想 $\mathfrak q \subset S$，因为极小素理想处的局部环
$S_{\mathfrak q}$ 维数为零，因而是 Cohen--Macaulay 的。
\end{proof}

\begin{lemma}
\label{algebra-lemma-finite-presentation-flat-CM-locus-open}
设 $R$ 是环。设 $R \to S$ 是有限表示且平坦的。对任意 $d \geq 0$，集合
$$
\left\{
\begin{matrix}
\mathfrak q \in \Spec(S)
\text{ such that setting }\mathfrak p = R \cap \mathfrak q
\text{ the fibre ring}\\
S_{\mathfrak q}/\mathfrak pS_{\mathfrak q}
\text{ 为 Cohen--Macaulay}
\text{ 且 } \dim_{\mathfrak q}(S/R) = d
\end{matrix}
\right\}
$$
在 $\Spec(S)$ 中是开的。
\end{lemma}

\begin{proof}
设 $\mathfrak q$ 是所示集合的元，$\mathfrak p$ 是 $R$ 中相应的素理想。
% P01-E0672: the frozen proof goal mentions Cohen-Macaulay fibres but omits the fixed relative-dimension condition defining the displayed locus.
必须找到 $g \in S$、$g \not \in \mathfrak q$，使得
$R \to S_g$ 的所有纤维环都是 Cohen--Macaulay 的。
在证明过程中，可以（有限多次地）用某个
$g \in S$、$g \not \in \mathfrak q$ 的 $S_g$ 代替 $S$。
因此，由引理 \ref{algebra-lemma-quasi-finite-over-polynomial-algebra}，
可以假设存在拟有限环映射 $R[t_1, \ldots, t_d] \to S$，其中
$d = \dim_{\mathfrak q}(S/R)$。令
$\mathfrak q' = R[t_1, \ldots, t_d] \cap \mathfrak q$。
由引理 \ref{algebra-lemma-where-CM}，可知环映射
$$
R[t_1, \ldots, t_d]_{\mathfrak q'} /
\mathfrak p R[t_1, \ldots, t_d]_{\mathfrak q'}
\longrightarrow
S_{\mathfrak q}/\mathfrak p S_{\mathfrak q}
$$
是平坦的。因此，由纤维平坦性准则引理
\ref{algebra-lemma-criterion-flatness-fibre}，可知
$R[t_1, \ldots, t_d]_{\mathfrak q'} \to S_{\mathfrak q}$ 是平坦的。
由定理 \ref{algebra-theorem-openness-flatness}，可知对某个
% P01-E0673: the frozen sentence lacks a comma after the localized-point condition before “the ring map”.
$g \in S$、$g \not \in \mathfrak q$，环映射
$R[t_1, \ldots, t_d] \to S_g$ 是平坦的。用 $S_g$ 代替 $S$ 后，
可知对每个素理想 $\mathfrak r \subset S$，令
$\mathfrak r' = R[t_1, \ldots, t_d] \cap \mathfrak r$ 及
$\mathfrak p' = R \cap \mathfrak r$，局部环映射
$R[t_1, \ldots, t_d]_{\mathfrak r'} \to S_{\mathfrak r}$ 是平坦的。
因此其基变换
$$
R[t_1, \ldots, t_d]_{\mathfrak r'} /
\mathfrak p' R[t_1, \ldots, t_d]_{\mathfrak r'}
\longrightarrow
S_{\mathfrak r}/\mathfrak p' S_{\mathfrak r}
$$
也是平坦的。因此，把引理 \ref{algebra-lemma-where-CM} 用于
$k = \kappa(\mathfrak p')$，可知 $\mathfrak r$ 如所需那样属于引理中的集合。
\end{proof}

\begin{lemma}
\label{algebra-lemma-generic-CM-flat-finite-presentation}
设 $R$ 是环，$R \to S$ 是有限表示的平坦映射。
使得纤维环 $S_{\mathfrak q} \otimes_R \kappa(\mathfrak p)$
为 Cohen--Macaulay 的素理想 $\mathfrak q$ 的集合，其中
$\mathfrak p = R \cap \mathfrak q$，在 $\Spec(S) \to \Spec(R)$ 的每条纤维中
都是开且稠密的。
\end{lemma}

\begin{proof}
由引理 \ref{algebra-lemma-finite-presentation-flat-CM-locus-open}，该集合（记作 $W$）是开的。
它在各纤维中稠密，因为 $W$ 与一条纤维的交正是该纤维中可应用引理
\ref{algebra-lemma-generic-CM} 的相应集合。
\end{proof}

\begin{lemma}
\label{algebra-lemma-extend-field-CM-locus}
设 $k$ 是域，$S$ 是有限型 $k$-代数。设 $K/k$ 是域扩张，并令
$S_K = K \otimes_k S$。设 $\mathfrak q \subset S$ 是 $S$ 的素理想，
并设 $\mathfrak q_K \subset S_K$ 是位于 $\mathfrak q$ 上方的 $S_K$ 的素理想。
则 $S_{\mathfrak q}$ 是 Cohen--Macaulay 的，当且仅当
$(S_K)_{\mathfrak q_K}$ 是 Cohen--Macaulay 的。
\end{lemma}

\begin{proof}
在证明过程中，可以（有限多次地）对任意
$g \in S$、$g \not \in \mathfrak q$ 用 $S_g$ 代替 $S$。
因此，利用引理 \ref{algebra-lemma-Noether-normalization-at-point}，
可以假设 $\dim(S) = \dim_{\mathfrak q}(S/k) =: d$，并找到有限单射
$k[x_1, \ldots, x_d] \to S$。注意，基变换还诱导有限单射
$K[x_1, \ldots, x_d] \to S_K$。
由引理 \ref{algebra-lemma-dimension-at-a-point-preserved-field-extension}，有
$\dim_{\mathfrak q_K}(S_K/K) = d$。令
$\mathfrak p = k[x_1, \ldots, x_d] \cap \mathfrak q$ 且
$\mathfrak p_K = K[x_1, \ldots, x_d] \cap \mathfrak q_K$。
考虑下列 Noether 局部环交换图：
$$
\xymatrix{
S_{\mathfrak q} \ar[r] &
(S_K)_{\mathfrak q_K} \\
k[x_1, \ldots, x_d]_{\mathfrak p} \ar[r] \ar[u] &
K[x_1, \ldots, x_d]_{\mathfrak p_K} \ar[u]
}
$$
由引理 \ref{algebra-lemma-where-CM}，必须证明左侧竖直箭头平坦，当且仅当右侧竖直箭头平坦。
由于底部箭头平坦，该等价性由引理 \ref{algebra-lemma-base-change-flat-up-down} 成立。
\end{proof}

\begin{lemma}
\label{algebra-lemma-CM-locus-commutes-base-change}
设 $R$ 是环，$R \to S$ 是有限型的。设 $R \to R'$ 是任意环映射，令
$S' = R' \otimes_R S$。
% P01-E0674: the frozen declaration “Denote f ... the map” omits “by/as”.
以 $f : \Spec(S') \to \Spec(S)$ 表示与环映射 $S \to S'$ 相关联的映射。
令 $W$ 等于满足如下条件的素理想 $\mathfrak q$ 的集合：纤维环
$S_{\mathfrak q} \otimes_R \kappa(\mathfrak p)$ 是 Cohen--Macaulay 的，
其中 $\mathfrak p = R \cap \mathfrak q$；并令 $W'$ 表示 $S'/R'$ 的类似集合。
则 $W' = f^{-1}(W)$。
\end{lemma}

\begin{proof}
% P01-E0675: the frozen one-sentence proof omits its subject.
这直接由引理 \ref{algebra-lemma-extend-field-CM-locus} 和定义得出。
\end{proof}

\begin{lemma}
\label{algebra-lemma-relative-dimension-CM}
设 $R$ 是环。设 $R \to S$ 是满足如下条件的环映射：(a) 平坦；
(b) 有限表示；(c) 具有 Cohen--Macaulay 纤维。
% P01-E0676: the frozen list is grammatically nonparallel: “is ... (c) has ...”.
则可以把 $S = S_0 \times \ldots \times S_n$ 写成 $R$-代数 $S_d$ 的乘积，
使得每个 $S_d$ 都满足 (a)、(b)、(c)，并且所有纤维都是维数为 $d$ 的等维环。
\end{lemma}

\begin{proof}
% P01-E0677: the frozen declaration “denote W_d ... the set” omits “by/as”.
对每个整数 $d$，以 $W_d \subset \Spec(S)$ 表示引理
\ref{algebra-lemma-finite-presentation-flat-CM-locus-open} 中定义的集合。
显然有 $\Spec(S) = \coprod W_d$，且由刚引用的引理，每个 $W_d$ 都是开的。
因此结论由引理 \ref{algebra-lemma-disjoint-implies-product} 得出。
\end{proof}

\section{微分}
\label{algebra-section-differentials}

\noindent
本节定义环映射的微分模。

\begin{definition}
\label{algebra-definition-derivation}
设 $\varphi : R \to S$ 是环映射，$M$ 是 $S$-模。
一个{\it 导子}，更准确地说，一个取值于 $M$ 的
{\it $R$-导子}，是映射 $D : S \to M$，它具有可加性，
在 $\varphi(R)$ 的元上为零，并且满足{\it Leibniz 法则}：
$D(ab) = aD(b) + bD(a)$。
\end{definition}

\noindent
注意，如果 $r \in R$ 且 $a \in S$，则 $D(ra) = rD(a)$。
等价地，$R$-导子是满足 Leibniz 法则的 $R$-线性映射 $D : S \to M$。
所有 $R$-导子的集合构成一个 $S$-模：给定两个 $R$-导子 $D, D'$，
其和 $D + D' : S \to M$，$a \mapsto D(a)+D'(a)$ 是 $R$-导子；
给定 $R$-导子 $D$ 和元 $c\in S$，其标量倍
$cD : S \to M$，$a \mapsto cD(a)$ 也是 $R$-导子。这个 $S$-模记作
$$
\text{Der}_R(S, M).
$$
另外，如果 $\alpha : M \to N$ 是 $S$-模映射，则复合
$\alpha \circ D$ 是到 $N$ 的 $R$-导子。这样，赋值
$M \mapsto \text{Der}_R(S, M)$ 是协变函子。

\medskip\noindent
考虑下列自由 $S$-模映射
$$
\bigoplus\nolimits_{(a, b)\in S^2} S[(a, b)] \oplus
\bigoplus\nolimits_{(f, g)\in S^2} S[(f, g)] \oplus
\bigoplus\nolimits_{r\in R} S[r]
\longrightarrow
\bigoplus\nolimits_{a\in S} S[a]
$$
它由如下规则定义：
$$
[(a, b)] \longmapsto [a + b] - [a] - [b],\quad
[(f, g)] \longmapsto [fg] -f[g] - g[f],\quad
[r]      \longmapsto [\varphi(r)]
$$
记号的含义是显然的。令 $\Omega_{S/R}$ 为该映射的余核。
有映射 $\text{d} : S \to \Omega_{S/R}$，它把 $a$ 映到
$[a]$ 在余核中的类 $\text{d}a$。由 $\Omega_{S/R}$ 上所施加的关系，
这是一个 $R$-导子；换言之，
$$
\text{d}(a + b) = \text{d}a + \text{d}b, \quad
\text{d}(fg) = f\text{d}g + g\text{d}f, \quad
\text{d}\varphi(r) = 0
$$
其中 $a,b,f,g \in S$ 且 $r \in R$。

\begin{definition}
\label{algebra-definition-differentials}
二元组 $(\Omega_{S/R}, \text{d})$ 称为 $S$ 在 $R$ 上的
{\it K"ahler 微分模}，或简称{\it 微分模}。
\end{definition}

\begin{lemma}
\label{algebra-lemma-universal-omega}
\begin{slogan}
从微分模出发的映射与导子相同。
\end{slogan}
$S$ 在 $R$ 上的微分模具有如下泛性质。映射
$$
\Hom_S(\Omega_{S/R}, M)
\longrightarrow
\text{Der}_R(S, M), \quad
\alpha
\longmapsto
\alpha \circ \text{d}
$$
是函子的同构。
\end{lemma}

\begin{proof}
按定义，$R$-导子是一条规则，它为每个 $a \in S$ 配置一个元
$D(a) \in M$。因此 $D$ 给出映射 $[D] : \bigoplus S[a] \to M$。
而成为 $R$-导子的各项条件恰好意味着，$[D]$ 在上面所示
$\Omega_{S/R}$ 的表示映射的像上为零。
\end{proof}

\begin{lemma}
\label{algebra-lemma-trivial-differential-surjective}
假设 $R \to S$ 是满射。则 $\Omega_{S/R} = 0$。
\end{lemma}

\begin{proof}
这或者可由所有 $R$-导子显然必须为零看出，或者可由
$\Omega_{S/R}$ 的表示中的映射是满射看出。
\end{proof}

\noindent
假设
\begin{equation}
\label{algebra-equation-functorial-omega}
\vcenter{
\xymatrix{
S \ar[r]_\varphi
&
S'
\\
R \ar[r]^\psi \ar[u]^\alpha
&
R' \ar[u]_\beta
}
}
\end{equation}
是环的交换图。此时有微分模之间的自然映射，它嵌入交换图
$$
\xymatrix{
\Omega_{S/R} \ar[r] &
\Omega_{S'/R'}
\\
S \ar[u]^{\text{d}} \ar[r]^{\varphi}
&
S' \ar[u]_{\text{d}}
}
$$
要构造该映射，只须使用 $\Omega_{S/R}$ 与 $\Omega_{S'/R'}$
的表示之间显然的映射。具体地说，
\begin{equation}
\label{algebra-equation-map-presentations}
\vcenter{
\xymatrix{
\bigoplus S'[(a', b')]
\oplus
\bigoplus S'[(f', g')]
\oplus
\bigoplus S'[r'] \ar[r]
&
\bigoplus S' [a'] \\
 \\
\bigoplus S[(a, b)]
\oplus
\bigoplus S[(f, g)]
\oplus
\bigoplus S[r] \ar[r]
\ar[uu]^{
\begin{matrix}
[(a, b)] \mapsto [(\varphi(a), \varphi(b))] \\
[(f, g)] \mapsto [(\varphi(f), \varphi(g))] \\
[r]\mapsto [\psi(r)]
\end{matrix}
} &
\bigoplus S[a] \ar[uu]_{[a] \mapsto [\varphi(a)]}
}
}
\end{equation}
所得结果就是：$f\text{d}g \in \Omega_{S/R}$ 映到
$\varphi(f)\text{d}\varphi(g)$。


\begin{lemma}
\label{algebra-lemma-colimit-differentials}
设 $I$ 是有向集。设 $(R_i \to S_i, \varphi_{ii'})$ 是 $I$ 上的环映射系统；
见“范畴”第 \ref{categories-section-posets-limits} 节。则
% P01-E0678: the frozen display ends with a period before the following “where” clause.
$$
\Omega_{S/R} =
\colim_i \Omega_{S_i/R_i}.
$$
其中 $R \to S = \colim (R_i \to S_i)$。
\end{lemma}

\begin{proof}
这由 $\Omega_{S/R}$ 的定义表示以及上面所述的这一构造的函子性立即得出。
% P01-E0679: the frozen phrase “the functoriality of this described above” lacks a noun modified by “of this”.
\end{proof}

\begin{lemma}
\label{algebra-lemma-differential-surjective}
在图 (\ref{algebra-equation-functorial-omega}) 中，假设 $S \to S'$ 是满射，
其核为 $I \subset S$。则 $\Omega_{S/R} \to \Omega_{S'/R'}$ 是满射，
其核作为 $S$-模由如下元生成：$\text{d}a$，其中 $a \in S$ 满足
$\varphi(a) \in \beta(R')$。（特别地，这包括元 $\text{d}(i)$，$i \in I$。）
\end{lemma}

\begin{proof}[第一种证明]
考虑表示映射 (\ref{algebra-equation-map-presentations})。
显然，自由模之间右侧的竖直映射是满射。因此该映射是满射。
假设 $\Omega_{S/R}$ 的某个元 $\eta$ 在 $\Omega_{S'/R'}$ 中映到零。
把 $\eta$ 写成某个 $s_i, a_i \in S$ 的 $\sum s_i[a_i]$ 的像。
那么 $\sum \varphi(s_i)[\varphi(a_i)]$ 是上方左角中元
$$
\theta = \sum s_j'[a_j', b_j'] + \sum s_k'[f_k', g_k'] + \sum s_l'[r_l']
$$
的像。由于 $\varphi$ 是满射，各项 $s_j'[a_j', b_j']$ 和
$s_k'[f_k', g_k']$ 都是下方右角中各元的像。
% P01-E0680: the frozen proof names the lower right corner, but relation generators lift from the lower left corner.
因此，从 $\eta$ 和 $\theta$ 中减去这些元的像来修改它们后，可以假设
$\theta = \sum s_l'[r_l']$。换言之，
$\sum \varphi(s_i)[\varphi(a_i)]$ 具有形式
$\sum s'_l [\beta(r'_l)]$。
接着，可以假设存在某个 $a' \in S'$，使得对所有 $i$ 都有
$a' = \varphi(a_i)$，并且对所有 $l$ 都有 $a' = \beta(r_l')$。
% P01-E0681: the frozen proof suppresses the needed componentwise decomposition before assuming one common basis element a'.
这由图中右上角的直和分解立即得出。选取满足
$\varphi(a) = a'$ 的 $a \in S$。于是可以对某些 $x_i \in I$ 写成
$a_i = a + x_i$。因此利用容许的关系，可以假设所有 $a_i$ 都等于 $a$。
但此时可假设我们的元具有形式 $s[a]$。仍然知道
% P01-E0682: the frozen formula applies phi to coefficients s'_l that already lie in S'.
$\varphi(s)[a'] = \sum \varphi(s_l')[\beta(r_l')]$。
所以，或者 $\varphi(s) = 0$，此时已得结论；或者
$a' = \varphi(a)$ 在 $\beta$ 的像中，此时同样已得结论。
\end{proof}

\begin{proof}[第二种证明]
我们将使用引理 \ref{algebra-lemma-universal-omega} 给出的微分模泛性质，
不再另行说明。

\medskip\noindent
在 (\ref{algebra-equation-functorial-omega}) 中令 $R'' = S \times_{S'} R'$。
则有下图：
% P01-E0683: the frozen phrase “we have following diagram” omits “the”.
$$
\xymatrix{
S \ar[r] & S \ar[r] & S' \\
R \ar[r] \ar[u] & R'' \ar[r] \ar[u] & R' \ar[u]
}
$$
设 $M$ 是 $S$-模。由定义立即可知，$R$-导子 $D : S \to M$ 是
$R''$-导子，当且仅当它在 $R'' \to S$ 的像中的元上为零。
泛性质把这转化为如下断言：自然映射
$\Omega_{S/R} \to \Omega_{S/R''}$ 是满射，其核作为 $S$-模由
$R''$ 的像生成。
% P01-E0684: the frozen kernel description names elements of R'' rather than their differentials in Omega.

\medskip\noindent
由前一段可知，只须证明：当 $S \to S'$ 是满射且
$R = S \times_{S'} R'$ 时，$\Omega_{S/R} \to \Omega_{S'/R'}$ 是同构。
设 $M'$ 是 $S'$-模。注意，任意 $R'$-导子 $D' : S' \to M'$ 与
$S \to S'$ 前复合后，都给出一个 $R$-导子。
反过来，假设 $M$ 是 $S$-模，$D : S \to M$ 是 $R$-导子。
如果 $i \in I$，则存在 $a \in R$ 使得 $\alpha(a) = i$
（因为 $R = S \times_{S'} R'$）。于是 $D(i) = 0$，从而对所有
$s \in S$，有 $0 = D(is) = iD(s)$。因此，$D$ 的像包含于
由 $I$ 零化的子模 $M' \subset M$ 中，而且诱导映射
$S \to M'$ 分解为一个 $R'$-导子 $S' \to M'$。
利用泛性质看出这意味着 $\Omega_{S/R} \to \Omega_{S'/R'}$ 是同构，
这是一个练习；细节从略。
\end{proof}

\begin{lemma}
\label{algebra-lemma-exact-sequence-differentials}
设 $A \to B \to C$ 是环映射。则有典范的 $C$-模正合列
$$
C \otimes_B \Omega_{B/A} \to
\Omega_{C/A} \to
\Omega_{C/B} \to 0
$$
\end{lemma}

\begin{proof}
在图 (\ref{algebra-equation-functorial-omega}) 中置
$R = A$、$S = C$、$R' = B$ 及 $S' = C$，得到一个图。
由引理 \ref{algebra-lemma-differential-surjective}，映射
$\Omega_{C/A} \to \Omega_{C/B}$ 是满射，其核由元
$\text{d}(c)$ 生成，其中 $c \in C$ 属于 $B \to C$ 的像。引理得证。
\end{proof}

\begin{lemma}
\label{algebra-lemma-differentials-localize}
设 $\varphi : A \to B$ 是环映射。
\begin{enumerate}
\item 若 $S \subset A$ 是一个映到 $B$ 的可逆元中的乘法子集，则
$\Omega_{B/A} = \Omega_{B/S^{-1}A}$。
\item 若 $S \subset B$ 是乘法子集，则
$S^{-1}\Omega_{B/A} = \Omega_{S^{-1}B/A}$。
\end{enumerate}
\end{lemma}

\begin{proof}
要证明 (1) 中的等式，只须证明任意 $A$-导子
$D : B \to M$ 都在元 $\varphi(s)^{-1}$ 上为零。
这由把 Leibniz 法则应用于 $1 = \varphi(s) \varphi(s)^{-1}$ 立即可得。
% P01-E0697: the frozen introductory phrase “To show (2) note” lacks a comma after “(2)”.
要证明 (2)，注意有显然的映射
$S^{-1}\Omega_{B/A} \to \Omega_{S^{-1}B/A}$。
要证明它是同构，只须证明存在一个取值于 $S^{-1}\Omega_{B/A}$ 的
$S^{-1}B$ 的 $A$-导子 $\text{d}'$。
% P01-E0685: the frozen source says “a A-derivation” instead of “an A-derivation”.
为定义它，只须置
$\text{d}'(b/s) = (1/s)\text{d}b - (1/s^2)b\text{d}s$。
细节从略。
\end{proof}

\begin{lemma}
\label{algebra-lemma-differential-seq}
在图 (\ref{algebra-equation-functorial-omega}) 中，假设 $S \to S'$ 是满射，
其核为 $I \subset S$，并且假设 $R' = R$。
则有典范的 $S'$-模正合列
$$
I/I^2
\longrightarrow
\Omega_{S/R} \otimes_S S'
\longrightarrow
\Omega_{S'/R}
\longrightarrow
0
$$
% P01-E0692: the frozen source omits terminal punctuation after this displayed exact sequence.
最左侧映射由如下规则刻画：$f \in I$ 映到
$\text{d}f \otimes 1$。
\end{lemma}

\begin{proof}
中间项是 $\Omega_{S/R} \otimes_S S/I$。
对 $f \in I$，用 $\overline{f}$ 表示 $f$ 在 $I/I^2$ 中的像。
% P01-E0686: the frozen phrase “denote \bar f the image” omits “by/as”.
要证明映射 $\overline{f} \mapsto \text{d}f \otimes 1$ 是良定义的，
只须验证当 $f_1, f_2 \in I$ 时
$\text{d} f_1f_2 \otimes 1 = 0$。
这由 Leibniz 法则立即可得：
$\text{d} f_1f_2 \otimes 1 = (f_1 \text{d}f_2 + f_2 \text{d} f_1 )\otimes 1 =
\text{d}f_2 \otimes f_1 + \text{d}f_1 \otimes f_2 = 0$。
类似的计算表明此映射是 $S' = S/I$-线性的。
% P01-E0687: the frozen source has subject-verb disagreement in “A similar computation show”.

\medskip\noindent
映射 $\Omega_{S/R} \otimes_S S' \to \Omega_{S'/R}$ 是与
$S$-线性映射 $\Omega_{S/R} \to \Omega_{S'/R}$ 相伴的典范
$S'$-线性映射。它是满射，因为由引理
\ref{algebra-lemma-differential-surjective}，映射
$\Omega_{S/R} \to \Omega_{S'/R}$ 是满射。

\medskip\noindent
两个映射的复合为零，因为当 $f \in I$ 时，
$\text{d}f$ 在 $\Omega_{S'/R}$ 中映到零。
注意，正合性恰好是说 $\Omega_{S/R} \to \Omega_{S'/R}$ 的核
作为 $S$-子模由子模 $I\Omega_{S/R}$ 与元 $\text{d}f$（$f \in I$）
共同生成。由引理 \ref{algebra-lemma-differential-surjective}，我们知道
此核由元 $\text{d}(a)$ 生成，其中对某个 $r \in R$ 有
$\varphi(a) = \beta(r)$。但这时
$a = \alpha(r) + a - \alpha(r)$，所以
$\text{d}(a) = \text{d}(a - \alpha(r))$。并且
$a - \alpha(r) \in I$，因为 $\varphi(a - \alpha(r)) =
\varphi(a) - \varphi(\alpha(r)) = \beta(r) - \beta(r) = 0$。
由此可知，元 $\text{d}f$（$f \in I$）已经作为 $S$-模生成此核，
正如所需。
% P01-E0696: the frozen sentence “We conclude the elements...” omits “that”.
\end{proof}

\begin{lemma}
\label{algebra-lemma-differential-seq-split}
在图 (\ref{algebra-equation-functorial-omega}) 中，假设 $S \to S'$ 是满射，
其核为 $I \subset S$，并假设 $R' = R$。进一步，假设存在一个
$R$-代数映射 $S' \to S$，它是 $S \to S'$ 的右逆。
则引理 \ref{algebra-lemma-differential-seq} 的 $S'$-模正合列变为短正合列
$$
0 \longrightarrow
I/I^2
\longrightarrow
\Omega_{S/R} \otimes_S S'
\longrightarrow
\Omega_{S'/R}
\longrightarrow
0
$$
而且它甚至是分裂短正合列。
\end{lemma}

\begin{proof}
令 $\beta : S' \to S$ 是满射 $\alpha : S \to S'$ 的右逆。
考虑映射
$$
D : S \longrightarrow I/I^2,
\quad
x \longmapsto x - \beta(\alpha(x))
$$
% P01-E0693: the frozen source omits terminal punctuation after the displayed map.
容易证明 $D$ 是 $R$-导子（从略）。
此外，若 $x \in I, s \in S$，则 $x D(s) = 0$。
因此由泛性质，$D$ 诱导映射
$\tau : \Omega_{S/R} \otimes_S S' \to I/I^2$。
它是 $\text{d} : I/I^2  \to \Omega_{S/R} \otimes_S S'$ 的左逆；
此验证从略。因此结论成立。
\end{proof}

\begin{lemma}
\label{algebra-lemma-differential-mod-power-ideal}
设 $R \to S$ 是环映射，$I \subset S$ 是理想，$n \geq 1$ 是整数。
置 $S' = S/I^{n + 1}$。映射
$\Omega_{S/R} \to \Omega_{S'/R}$ 诱导同构
$$
\Omega_{S/R} \otimes_S S/I^n
\longrightarrow
\Omega_{S'/R} \otimes_{S'} S/I^n.
$$
\end{lemma}

\begin{proof}
这由引理 \ref{algebra-lemma-differential-seq} 以及
$\text{d}$ 的 Leibniz 法则所给出的
$\text{d}(I^{n + 1}) \subset I^n\Omega_{S/R}$ 得出。
\end{proof}

\begin{lemma}
\label{algebra-lemma-differentials-base-change}
假设有环映射 $R \to R'$ 和 $R \to S$。
置 $S' = S \otimes_R R'$，于是得到图
(\ref{algebra-equation-functorial-omega})。则上面定义的典范映射诱导同构
$\Omega_{S/R} \otimes_R R' = \Omega_{S'/R'}$。
\end{lemma}

\begin{proof}
令 $\text{d}' : S' = S \otimes_R R' \to \Omega_{S/R} \otimes_R R'$
表示映射
$\text{d}'( \sum a_i \otimes x_i ) = \sum \text{d}(a_i) \otimes x_i$。
它存在，因为映射
$S \times R' \to \Omega_{S/R} \otimes_R R'$，
$(a, x)\mapsto \text{d}a \otimes_R x$ 是 $R$-双线性的。
这是一个 $R'$-导子；直接计算即可验证。
我们将证明 $(\Omega_{S/R} \otimes_R R', \text{d}')$ 满足泛性质。
设 $D : S' \to M'$ 是一个取值于某个 $S'$-模的 $R'$-导子。
复合 $S \to S' \to M'$ 是 $R$-导子，所以得到一个
$S$-线性映射 $\varphi_D : \Omega_{S/R} \to M'$。
将它与 $R'$ 作张量积，得到映射
$\varphi'_D : \Omega_{S/R} \otimes_R R' \to M'$，
$\varphi'_D(\eta \otimes x) = x\varphi_D(\eta)$。
显然 $D = \varphi'_D \circ \text{d}'$。
\end{proof}

\noindent
乘法映射 $S \otimes_R S \to S$ 是把 $a \otimes b$ 映到
$S$ 中 $ab$ 的 $R$-代数映射。若通过两个映射
$S \to S \otimes_R S$ 中的任一个，把 $S \otimes_R S$
视为 $S$-代数，则它也是 $S$-代数映射。

\begin{lemma}
\label{algebra-lemma-differentials-diagonal}
设 $R \to S$ 是环映射，令 $J = \Ker(S \otimes_R S \to S)$
为乘法映射的核。存在典范的 $S$-模同构
$\Omega_{S/R} \to J/J^2$，
$a \text{d} b \mapsto a \otimes b - ab \otimes 1$。
\end{lemma}

\begin{proof}[第一种证明]
把引理 \ref{algebra-lemma-differential-seq-split} 应用于交换图
$$
\xymatrix{
S \otimes_R S \ar[r] & S \\
S \ar[r] \ar[u] & S \ar[u]
}
$$
其中左侧竖直箭头是 $a \mapsto a \otimes 1$。得到正合列
$0 \to J/J^2 \to
\Omega_{S \otimes_R S/S} \otimes_{S \otimes_R S} S \to \Omega_{S/S} \to 0$。
由引理 \ref{algebra-lemma-trivial-differential-surjective}，
项 $\Omega_{S/S}$ 是 $0$，因而得到另两项之间的同构。
由引理 \ref{algebra-lemma-differentials-base-change}，
由于 $S \to S \otimes_R S$ 是 $R \to S$ 的基变换，有
$\Omega_{S \otimes_R S/S} = \Omega_{S/R} \otimes_S (S \otimes_R S)$，
所以
$$
\Omega_{S \otimes_R S/S} \otimes_{S \otimes_R S} S =
\Omega_{S/R} \otimes_S (S \otimes_R S) \otimes_{S \otimes_R S} S =
\Omega_{S/R}
$$
% P01-E0694: the frozen source omits terminal punctuation after this displayed chain of equalities.
映射由引理中的规则给出；此验证从略。
\end{proof}

\begin{proof}[第二种证明]
先证明规则 $a \text{d} b \mapsto a \otimes b - ab \otimes 1$
是良定义的。为此必须证明 $\text{d}r$ 和
$a\text{d}b + b \text{d}a - d(ab)$ 映到零。
% P01-E0688: the frozen relation uses d(ab) instead of the established \text{d}(ab) notation.
第一项为零，因为由张量积的定义，
$r \otimes 1 - 1 \otimes r = 0$；第二项为零，因为
% P01-E0689: the frozen “The first because...” clause is a sentence fragment.
$$
(a \otimes b - ab \otimes 1) +
(b \otimes a - ba \otimes 1) -
(1 \otimes ab - ab \otimes 1)
=
(a \otimes 1 - 1\otimes a)(1\otimes b - b \otimes 1)
$$
属于 $J^2$。

\medskip\noindent
我们构造反方向的映射。可以把
$S \to S \otimes_R S$，$a \mapsto a \otimes 1$
视为 $R \to S$ 的基变换。因此由引理
\ref{algebra-lemma-differentials-base-change} 有
$\Omega_{S \otimes_R S/S} = \Omega_{S/R} \otimes_S (S \otimes_R S)$。
此时，引理 \ref{algebra-lemma-differential-seq} 的序列给出映射
$$
J/J^2  \to \Omega_{S \otimes_R S/ S} \otimes_{S \otimes_R S} S
= (\Omega_{S/R} \otimes_S (S \otimes_R S))\otimes_{S \otimes_R S} S
= \Omega_{S/R}.
$$
其为上述映射的逆；留给读者验证。
% P01-E0690: the frozen phrase “see it is the inverse” omits “that”.
\end{proof}

\begin{lemma}
\label{algebra-lemma-differentials-polynomial-ring}
若 $S = R[x_1, \ldots, x_n]$，则 $\Omega_{S/R}$ 是一个有限自由
$S$-模，其基为 $\text{d}x_1, \ldots, \text{d}x_n$。
\end{lemma}

\begin{proof}
先证明 $\text{d}x_1, \ldots, \text{d}x_n$ 作为 $S$-模生成
$\Omega_{S/R}$。为此，对任意 $g \in R[x_1, \ldots, x_n]$，
我们证明 $\text{d}g$ 可表示为和
$\sum g_i \text{d}x_i$。对 $g$ 的（全）次数作归纳。
若 $g$ 的次数为 $0$，则结论显然，因为此时 $\text{d}g = 0$。
若 $g$ 的次数为 $> 0$，则可把 $g$ 写成
$c + \sum g_i x_i$，其中 $c\in R$ 且
$\deg(g_i) < \deg(g)$。由 Leibniz 法则，
$\text{d}g = \sum g_i \text{d} x_i + \sum x_i \text{d}g_i$，
因而由归纳法得到结论。

\medskip\noindent
考虑 $R$-导子
$\partial / \partial x_i :
R[x_1, \ldots, x_n] \to R[x_1, \ldots, x_n]$。
（它的定义留给读者；其刻画性质是
$\partial / \partial x_i (x_j) = \delta_{ij}$。）
由泛性质，它对应于一个 $S$-模映射
$l_i : \Omega_{S/R} \to R[x_1, \ldots, x_n]$，
后者把 $\text{d}x_i$ 映到 $1$，并在 $j \not = i$ 时把
$\text{d}x_j$ 映到 $0$。因此，元
$\text{d}x_1, \ldots, \text{d}x_n$ 之间显然不存在
$S$-线性关系。
\end{proof}

\begin{lemma}
\label{algebra-lemma-differentials-finitely-presented}
假设 $R \to S$ 是有限表现的。则 $\Omega_{S/R}$ 是有限表现
$S$-模。
\end{lemma}

\begin{proof}
写成 $S = R[x_1, \ldots, x_n]/(f_1, \ldots, f_m)$，
并置 $I = (f_1, \ldots, f_m)$。由引理
\ref{algebra-lemma-differential-seq}，有 $S$-模正合列
$$
I/I^2
\to
\Omega_{R[x_1, \ldots, x_n]/R} \otimes_{R[x_1, \ldots, x_n]} S
\to
\Omega_{S/R}
\to
0
$$
% P01-E0695: the frozen source omits terminal punctuation after this displayed exact sequence.
结论由以下事实得出：$I/I^2$ 是有限 $S$-模（由 $f_i$ 的像生成），
而由引理 \ref{algebra-lemma-differentials-polynomial-ring}，中间项是有限自由的。
\end{proof}

\begin{lemma}
\label{algebra-lemma-differentials-finitely-generated}
假设 $R \to S$ 是有限型的。则 $\Omega_{S/R}$ 是有限生成
$S$-模。
% P01-E0691: the frozen source omits the article in “is finitely generated S-module”.
\end{lemma}

\begin{proof}
这与引理 \ref{algebra-lemma-differentials-finitely-presented} 的证明非常相似，
但更为容易。
\end{proof}

\section{de Rham 复形}
\label{algebra-section-de-rham-complex}

\noindent
设 $A \to B$ 是环映射。用 $\text{d} : B \to \Omega_{B/A}$
表示第 \ref{algebra-section-differentials} 节构造的微分模及其泛 $A$-导子。
% P01-E0698: the frozen designation grammatically identifies the map d as “the module of differentials” rather than as its universal derivation.
对 $i \geq 0$，令
$\Omega_{B/A}^i = \wedge^i_B(\Omega_{B/A})$ 为第
\ref{algebra-section-tensor-algebra} 节所述的第 $i$ 个外幂。
{\it $B$ 在 $A$ 上的 de Rham 复形}是复形
$$
\Omega_{B/A}^0 \to \Omega_{B/A}^1 \to \Omega_{B/A}^2 \to \ldots
$$
其 $A$-线性微分将在下面构造并描述。

\medskip\noindent
映射 $\text{d} : \Omega^0_{B/A} \to \Omega^1_{B/A}$
就是泛导子 $\text{d} : B \to \Omega_{B/A}$。
注意它确实是 $A$-线性的。

\medskip\noindent
对 $p \geq 1$，我们断言存在唯一的 $A$-线性映射
$\text{d} : \Omega_{B/A}^p \to \Omega_{B/A}^{p + 1}$，使得
\begin{equation}
\label{algebra-equation-rule}
\text{d}\left(b_0\text{d}b_1 \wedge \ldots \wedge \text{d}b_p\right) =
\text{d}b_0 \wedge \text{d}b_1 \wedge \ldots \wedge \text{d}b_p
\end{equation}
% P01-E0699: the frozen numbered equation lacks terminal punctuation before the next sentence.
回忆 $\Omega_{B/A}$ 作为 $B$-模由元 $\text{d}b$ 生成。
因此，$\Omega^p_{B/A}$ 作为 $A$-模由元
$b_0\text{d}b_1 \wedge \ldots \wedge \text{d}b_p$ 生成，
从而映射 $\text{d} : \Omega^p_{B/A} \to \Omega^{p + 1}_{B/A}$
若存在便是唯一的。

\medskip\noindent
构造 $\text{d} : \Omega_{B/A}^1 \to \Omega_{B/A}^2$。
由定义 \ref{algebra-definition-differentials}，元 $\text{d}b$
作为 $B$-模自由生成 $\Omega_{B/A}$，但须满足关系：
当 $a \in A$ 时 $\text{d}a = 0$，以及当 $b', b'' \in B$ 时
$\text{d}(b' + b'') = \text{d}b' + \text{d}b''$ 和
$\text{d}(b'b'') = b'\text{d}b'' + b''\text{d}b'$。
因此，要证明规则
$$
\sum b'_i \text{d}b_i \longmapsto \sum \text{d}b'_i \wedge \text{d}b_i
$$
是良定义的，必须证明诸元
$$
b\text{d}a,
\quad\text{且}\quad
b\text{d}(b' + b'') - b\text{d}b' - b\text{d}b''
\quad\text{且}\quad
b\text{d}(b'b'') - bb'\text{d}b'' - bb''\text{d}b'
$$
在 $a \in A$ 且 $b, b', b'' \in B$ 时映到零。
使用 $\text{d}$ 的 Leibniz 法则直接计算即可看出这一点。

\medskip\noindent
注意，复合
$\Omega^0_{B/A} \to \Omega^1_{B/A} \to \Omega^2_{B/A}$
为零，因为 $\text{d}(\text{d}(b)) = \text{d}(1 \text{d}b) =
\text{d}(1) \wedge \text{d}(b) = 0 \wedge \text{d}b = 0$。
这里 $\text{d}(1) = 0$，因为 $1 \in B$ 属于 $A \to B$ 的像。
下面将使用这一点。

\medskip\noindent
对 $p \geq 2$ 构造
$\text{d} : \Omega_{B/A}^p \to \Omega_{B/A}^{p + 1}$。
我们将证明由公式
$$
\gamma : \Omega^1_{B/A} \otimes_A \ldots \otimes_A \Omega^1_{B/A}
\longrightarrow
\Omega_{B/A}^{p + 1}
$$
定义的 $A$-线性映射
$$
\omega_1 \otimes \ldots \otimes \omega_p
\longmapsto
\sum (-1)^{i + 1}
\omega_1 \wedge \ldots \wedge \text{d}(\omega_i) \wedge \ldots \wedge \omega_p
$$
% P01-E0700: the frozen alternating-sum formula omits the summation bounds i = 1,...,p.
通过自然满射
$\Omega^1_{B/A} \otimes_A \ldots \otimes_A \Omega^1_{B/A} \to \Omega^p_{B/A}$
分解，从而给出所需映射
$\text{d} : \Omega^p_{B/A} \to \Omega^{p + 1}_{B/A}$。
由引理 \ref{algebra-lemma-present-wedge}，
$\Omega^1_{B/A} \otimes_A \ldots \otimes_A \Omega^1_{B/A} \to \Omega^p_{B/A}$
的核作为 $A$-模由以下元生成：某些 $i \not = j$ 满足
$\omega_i = \omega_j$ 的元
$\omega_1 \otimes \ldots \otimes \omega_p$，以及某个 $f \in B$ 对应的元
$\omega_1 \otimes \ldots \otimes f\omega_i \otimes \ldots \otimes \omega_p -
\omega_1 \otimes \ldots \otimes f\omega_j \otimes \ldots \otimes \omega_p$。
直接计算表明，第一类元被 $\gamma$ 映到 $0$；换言之，
$\gamma$ 是交替的。为完成证明，还须证明
$$
\gamma(
\omega_1 \otimes \ldots \otimes f\omega_i \otimes \ldots \otimes \omega_p)
=
\gamma(
\omega_1 \otimes \ldots \otimes f\omega_j \otimes \ldots \otimes \omega_p)
$$
对 $f \in B$ 成立。由 $A$-线性性和交替性，只须对
$p = 2$、$i = 1$、$j = 2$、
$\omega_1 = b \text{d}b'$ 及 $\omega_2 = c \text{d} c'$
证明，其中 $b, b', c, c' \in B$。因此须证明
\begin{align*}
& \text{d}(fb) \wedge \text{d}b' \wedge c \text{d}c'
- fb \text{d}b' \wedge \text{d}c \wedge \text{d}c' \\
& =
\text{d}b \wedge \text{d}b' \wedge fc\text{d}c'
- b \text{d}b' \wedge \text{d}(fc) \wedge \text{d}c'
\end{align*}
% P01-E0701: the frozen align display lacks punctuation before the continuing “in other words” clause.
换言之，须证明
$$
(c \text{d}(fb) + fb \text{d}c - fc \text{d}b - b \text{d}(fc))
\wedge \text{d}b' \wedge \text{d}c' = 0.
$$
这由 Leibniz 法则得出。注意，$\gamma$ 在元
$b_0\text{d}b_1 \otimes \text{d}b_2 \otimes \ldots \otimes \text{d}b_p$
上的值是
$\text{d}b_0 \wedge \text{d}b_1 \wedge \ldots \wedge \text{d}b_p$，
所以如此得到的映射
$\text{d} : \Omega^p_{B/A} \to \Omega^{p + 1}_{B/A}$
满足 (\ref{algebra-equation-rule})。

\medskip\noindent
最后，由于 $\Omega^p_{B/A}$ 作为加法群由元
$b_0\text{d}b_1 \wedge \ldots \wedge \text{d}b_p$ 生成，
且
$\text{d}(b_0\text{d}b_1 \wedge \ldots \wedge \text{d}b_p) =
\text{d}b_0 \wedge \ldots \wedge \text{d}b_p$，
完全相同的论证表明，对 $p \geq 1$，复合
$
\Omega^p_{B/A} \to
\Omega^{p + 1}_{B/A} \to
\Omega^{p + 2}_{B/A}
$
为零。因此，de Rham 复形确实是一个复形。

\medskip\noindent
给定一个环 $R$，置 $\Omega_R = \Omega_{R/\mathbf{Z}}$。
它有时称为 $R$ 的绝对微分模；此称谓是合理的：若
$\Omega_R$ 是只要求 Leibniz 法则而不要求 $\text{d}1$ 为零的微分模，
则 Leibniz 法则给出
$\text{d}1 = \text{d}(1 \cdot 1) =
1 \text{d}1 + 1 \text{d}1 = 2 \text{d}1$，从而
$\Omega_R$ 中 $\text{d}1 = 0$。在这种情形下，
{\it $R$ 的绝对 de Rham 复形}是相应的复形
$$
\Omega_R^0 \to \Omega_R^1 \to \Omega_R^2 \to \ldots
$$
% P01-E0702: the frozen display lacks a comma before the following “where” clause.
其中置 $\Omega^i_R = \Omega^i_{R/\mathbf{Z}}$，其余类推。

\medskip\noindent
假设有环的交换图
$$
\xymatrix{
B \ar[r] & B' \\
A \ar[r] \ar[u] & A' \ar[u]
}
$$
% P01-E0703: the frozen diagram lacks terminal punctuation before the next sentence.
存在自然的 de Rham 复形映射
$$
\Omega^\bullet_{B/A} \longrightarrow \Omega^\bullet_{B'/A'}
$$
% P01-E0704: the frozen displayed map lacks terminal punctuation before “Namely”.
具体而言，在次数 $0$ 它是映射 $B \to B'$；在次数 $1$，
它是第 \ref{algebra-section-differentials} 节构造的映射
$\Omega_{B/A} \to \Omega_{B'/A'}$；对 $p \geq 2$，它是诱导映射
$\Omega^p_{B/A} = \wedge^p_B(\Omega_{B/A}) \to
\wedge^p_{B'}(\Omega_{B'/A'}) = \Omega^p_{B'/A'}$。
与微分的相容性由公式 (\ref{algebra-equation-rule}) 对微分的刻画得出。

\begin{lemma}
\label{algebra-lemma-de-rham-complex-base-change}
假设有环映射 $A \to A'$ 和 $A \to B$。
置 $B' = B \otimes_A A'$，于是得到如上的图。
则上面定义的典范映射诱导复形同构
$\Omega^\bullet_{B/A} \otimes_A A' = \Omega^\bullet_{B'/A'}$。
\end{lemma}

\begin{proof}
这由引理 \ref{algebra-lemma-differentials-base-change} 以及取外幂与基变换可交换
这一事实得出。
\end{proof}

\begin{lemma}
\label{algebra-lemma-de-rham-complex}
设 $A \to B$ 是环映射，$\pi : \Omega_{B/A} \to \Omega$
是满的 $B$-模映射。用 $\text{d} : B \to \Omega$ 表示 $\pi$
与泛导子 $\text{d}_{B/A} : B \to \Omega_{B/A}$ 的复合。
% P01-E0705: the frozen designation “Denote d ... the composition” omits “by/as”.
置 $\Omega^i = \wedge_B^i(\Omega)$。
假设 $\pi$ 的核作为 $B$-模由如下元
$\omega \in \Omega_{B/A}$ 生成：$\text{d}_{B/A}(\omega) \in
\Omega_{B/A}^2$ 在 $\Omega^2$ 中映到零。
则存在 de Rham 复形
$$
\Omega^0 \to \Omega^1 \to \Omega^2 \to \ldots
$$
其微分由规则
$$
\text{d} : \Omega^p \to \Omega^{p + 1},\quad
\text{d}\left(f_0\text{d}f_1 \wedge \ldots \wedge \text{d}f_p\right) =
\text{d}f_0 \wedge \text{d}f_1 \wedge \ldots \wedge \text{d}f_p
$$
% P01-E0706: the frozen defining display lacks terminal punctuation at the end of the lemma.
\end{lemma}

\begin{proof}
我们将证明存在交换图
$$
\xymatrix{
\Omega_{B/A}^0 \ar[d] \ar[r]_{\text{d}_{B/A}} &
\Omega_{B/A}^1 \ar[d]_\pi \ar[r]_{\text{d}_{B/A}} &
\Omega_{B/A}^2 \ar[d]_{\wedge^2\pi} \ar[r]_{\text{d}_{B/A}} &
\ldots \\
\Omega^0 \ar[r]^{\text{d}} &
\Omega^1 \ar[r]^{\text{d}} &
\Omega^2 \ar[r]^{\text{d}} &
\ldots
}
$$
% P01-E0707: the frozen source runs the diagram directly into a lowercase independent clause without punctuation or a conjunction.
映射 $\text{d}$ 的描述将由上述
$B$ 在 $A$ 上的 de Rham 复形的微分
$\text{d}_{B/A} : \Omega^p_{B/A} \to \Omega^{p + 1}_{B/A}$
的构造得出。由于最左侧竖直箭头是同构，第一个方块成立。
% P01-E0708: the frozen source writes the closed compound adjective “left most” instead of “leftmost”.
由于 $\pi$ 是满射，要得到第二个方块，只须证明
$\text{d}_{B/A}$ 把 $\pi$ 的核映到 $\wedge^2\pi$ 的核中。
已知 $\pi$ 的核的任意元都具有形式
$\sum b_i\omega_i$，其中 $\pi(\omega_i) = 0$ 且
$\wedge^2\pi(\text{d}_{B/A}(\omega_i)) = 0$。
由 $\text{d}_{B/A}$ 的 Leibniz 法则，
$\text{d}_{B/A}(\sum b_i\omega_i) = \sum b_i \text{d}_{B/A}(\omega_i) +
\sum \text{d}_{B/A}(b_i) \wedge \omega_i$。因此它在
$\wedge^2\pi$ 下映到零。

\medskip\noindent
对 $i > 1$，注意 $\wedge^i \pi$ 是满射，其核为映射
$\Ker(\pi) \wedge \Omega^{i - 1}_{B/A}
\to \Omega_{B/A}^i$ 的像。
对 $\omega_1 \in \Ker(\pi)$ 和
$\omega_2 \in \Omega^{i - 1}_{B/A}$，有
$$
\text{d}_{B/A}(\omega_1 \wedge \omega_2) =
\text{d}_{B/A}(\omega_1) \wedge \omega_2 -
\omega_1 \wedge \text{d}_{B/A}(\omega_2)
$$
% P01-E0709: the frozen display lacks a comma before the following “which” clause.
由上面刚刚证明的结论，它属于 $\wedge^{i + 1}\pi$ 的核。
因此得到上图中的第 $(i + 1)$ 个方块。证明完毕。
\end{proof}

\section{有限阶微分算子}
\label{algebra-section-differential-operators}

\noindent
本节引入有限阶微分算子。

\begin{definition}
\label{algebra-definition-differential-operators}
设 $R \to S$ 是环映射，$M$、$N$ 是 $S$-模，$k \geq 0$ 是整数。
归纳地定义一个{\it $k$ 阶微分算子 $D : M \to N$}：
它是一个 $R$-线性映射，使得对所有 $g \in S$，映射
$m \mapsto D(gm) - gD(m)$ 是 $k - 1$ 阶微分算子。
对于基始情形 $k = 0$，定义 $0$ 阶微分算子为 $S$-线性映射。
\end{definition}

\noindent
若 $D : M \to N$ 是 $k$ 阶微分算子，则对所有 $g \in S$，
映射 $gD$ 是 $k$ 阶微分算子。两个 $k$ 阶微分算子的和仍是
这类算子。因此，全体这样的算子所成的集合
$$
\text{Diff}^k(M, N) = \text{Diff}^k_{S/R}(M, N)
$$
是 $S$-模。我们有
$$
\text{Diff}^0(M, N) \subset
\text{Diff}^1(M, N) \subset
\text{Diff}^2(M, N) \subset \ldots
$$
% P01-E0710: the frozen inclusion display lacks terminal punctuation at the end of the paragraph.

\begin{lemma}
\label{algebra-lemma-composition-differential-operators}
设 $R \to S$ 是环映射，$L, M, N$ 是 $S$-模。
若 $D : L \to M$ 和 $D' : M \to N$ 分别是 $k$ 阶和 $k'$ 阶
微分算子，则 $D' \circ D$ 是 $k + k'$ 阶微分算子。
\end{lemma}

\begin{proof}
设 $g \in S$。则把 $x \in L$ 映到
$$
D'(D(gx)) - gD'(D(x)) = D'(D(gx)) - D'(gD(x)) + D'(gD(x)) - gD'(D(x))
$$
的映射是两个较低阶微分算子复合的和。因此对 $k + k'$ 作归纳即可。
\end{proof}

\begin{lemma}
\label{algebra-lemma-module-principal-parts}
设 $R \to S$ 是环映射，$M$ 是 $S$-模，并且 $k \geq 0$。
存在一个 $S$-模 $P^k_{S/R}(M)$ 以及典范同构
$$
\text{Diff}^k_{S/R}(M, N) = \Hom_S(P^k_{S/R}(M), N)
$$
它关于 $S$-模 $N$ 是函子的。
\end{lemma}

\noindent
通过 Yoneda 引理，这说明存在一个 $k$ 阶的“泛”微分算子
$D_{univ} : M \to P^k_{S/R}(M)$，使得对每个 $k$ 阶微分算子
$D : M \to N$，存在唯一的 $S$-线性映射
$\alpha : P^k_{S/R}(M) \to N$ 满足
$D = D_{univ} \circ \alpha$。
% P01-E0711: the frozen universal factorization reverses the composition order; the typed composite is alpha after D_univ.
偶对 $(P^k_{S/R}(M), D_{univ})$ 在唯一同构的意义下唯一。

\begin{proof}
$P^k_{S/R}(M)$ 的存在性可由一般范畴论论证推出
（在此插入未来的参考文献），但我们也给出一个构造。
% P01-E0712: the frozen source contains the editorial placeholder “insert future reference here”.
置 $F = \bigoplus_{m \in M} S[m]$，其中 $[m]$ 是表示与 $m$
对应的直和项中基元的符号。给定任意微分算子
$D : M \to N$，得到 $S$-线性映射 $L_D : F \to N$，
它把 $[m]$ 映到 $D(m)$。若 $D$ 的阶为 $0$，则 $L_D$ 零化诸元
$$
[m + m'] - [m] - [m'],\quad
g_0[m] - [g_0m]
$$
其中 $g_0 \in S$ 且 $m, m' \in M$。
若 $D$ 的阶为 $1$，则 $L_D$ 零化诸元
$$
[m + m'] - [m] - [m'],\quad
f[m] - [fm], \quad
g_0g_1[m] - g_0[g_1m] - g_1[g_0m] + [g_1g_0m]
$$
其中 $f \in R$、$g_0, g_1 \in S$ 且 $m \in M$。
若 $D$ 的阶为 $k$，则 $L_D$ 零化诸元
$[m + m'] - [m] - [m']$、$f[m] - [fm]$ 以及诸元
$$
g_0g_1\ldots g_k[m] - \sum g_0 \ldots \hat g_i \ldots g_k[g_im] + \ldots
+(-1)^{k + 1}[g_0\ldots g_km]
$$
% P01-E0713: the frozen general relation display lacks terminal punctuation before “Conversely”.
反之，若 $S$-线性映射 $L : F \to N$ 零化前一句所列的所有元，
则 $m \mapsto L([m])$ 是 $k$ 阶微分算子。
因此，$P^k_{S/R}(M)$ 是 $F$ 关于这些元所生成子模的商。
\end{proof}

\begin{definition}
\label{algebra-definition-module-principal-parts}
设 $R \to S$ 是环映射，$M$ 是 $S$-模。
引理 \ref{algebra-lemma-module-principal-parts} 中构造的模
$P^k_{S/R}(M)$ 称为 $M$ 的{\it $k$ 阶主部模}。
\end{definition}

\noindent
注意，包含关系
$$
\text{Diff}^0(M, N) \subset
\text{Diff}^1(M, N) \subset
\text{Diff}^2(M, N) \subset \ldots
$$
通过 Yoneda 引理（范畴，引理 \ref{categories-lemma-yoneda}）
对应于满射
$$
\ldots \to P^2_{S/R}(M) \to P^1_{S/R}(M) \to P^0_{S/R}(M) = M
$$
% P01-E0714: the frozen surjection display lacks terminal punctuation at the end of the paragraph.

\begin{example}
\label{algebra-example-derivations-and-differential-operators}
设 $R \to S$ 是环映射，$N$ 是 $S$-模。
注意 $\text{Diff}^1(S, N) = \text{Der}_R(S, N) \oplus N$。
事实上，若 $D : S \to N$ 是 $1$ 阶微分算子，则由
$\sigma_D(g) := D(g) - gD(1)$ 定义的
$\sigma_D : S \to N$ 是 $R$-导子，并且
$D = \sigma_D + \lambda_{D(1)}$，其中
$\lambda_x : S \to N$ 是把 $g$ 映到 $gx$ 的线性映射。
由 $\Omega_{S/R}$ 的泛性质可得
$P^1_{S/R} = \Omega_{S/R} \oplus S$。
\end{example}

\begin{lemma}
\label{algebra-lemma-sequence-of-principal-parts}
设 $R \to S$ 是环映射，$M$ 是 $S$-模。存在典范短正合列
$$
0 \to \Omega_{S/R} \otimes_S M \to P^1_{S/R}(M) \to M \to 0
$$
它关于 $M$ 是函子的，称为{\it 主部列}。
\end{lemma}

\begin{proof}
映射 $P^1_{S/R}(M) \to M$ 已在上面给出。
设 $N$ 是 $S$-模，$D : M \to N$ 是 $1$ 阶微分算子。
对 $m \in M$，映射
$$
g \longmapsto D(gm) - gD(m)
$$
依照 $1$ 阶微分算子的公理，是一个 $R$-导子 $S \to N$。
因此，它对应于线性映射 $D_m : \Omega_{S/R} \to N$，
后者由规则 $a\text{d}b \mapsto aD(bm) - abD(m)$ 决定
（见引理 \ref{algebra-lemma-universal-omega}）。映射
$$
\Omega_{S/R} \times M \longrightarrow N,\quad
(\eta, m) \longmapsto D_m(\eta)
$$
是 $S$-双线性的（细节从略），因而决定一个 $S$-线性映射
$$
\sigma_D : \Omega_{S/R} \otimes_S M \to N
$$
% P01-E0715: the frozen sigma_D display lacks terminal punctuation before the next sentence.
这样便得到关于 $N$ 函子的映射
$\text{Diff}^1(M, N) \to \Hom_S(\Omega_{S/R} \otimes_S M, N)$，
$D \mapsto \sigma_D$。由 Yoneda 引理，这对应于一个映射
$\Omega_{S/R} \otimes_S M \to P^1_{S/R}(M)$。
% P01-E0716: the frozen phrase “this corresponds a map” omits “to”.
由构造立即可见，此映射关于 $M$ 是函子的。序列
$$
\Omega_{S/R} \otimes_S M \to P^1_{S/R}(M) \to M \to 0
$$
是正合的，因为对每个模 $N$，序列
$$
0 \to \Hom_S(M, N) \to
\text{Diff}^1(M, N) \to
\Hom_S(\Omega_{S/R} \otimes_S M, N)
$$
由直接检验是正合的。

\medskip\noindent
要看出 $\Omega_{S/R} \otimes_S M \to P^1_{S/R}(M)$ 是单射，
论证如下。选取正合列
$$
0 \to M' \to F \to M \to 0
$$
其中 $F$ 是自由 $S$-模。这对所有 $N$ 诱导正合列
$$
0 \to \text{Diff}^1(M, N) \to \text{Diff}^1(F, N) \to \text{Diff}^1(M', N)
$$
这证明在交换图
$$
\xymatrix{
0 \ar[r] &
\Omega_{S/R} \otimes_S M' \ar[r] \ar[d] &
P^1_{S/R}(M') \ar[r] \ar[d] &
M' \ar[r] \ar[d] &
0 \\
0 \ar[r] &
\Omega_{S/R} \otimes_S F \ar[r] \ar[d] &
P^1_{S/R}(F) \ar[r] \ar[d] &
F \ar[r] \ar[d] &
0 \\
0 \ar[r] &
\Omega_{S/R} \otimes_S M \ar[r] \ar[d] &
P^1_{S/R}(M) \ar[r] \ar[d] &
M \ar[r] \ar[d] &
0 \\
& 0 & 0 & 0
}
$$
中，中间列是正合的。由
$\Omega_{S/R} \otimes_S -$ 的右正合性，左列是正合的。
由蛇形引理（见第 \ref{algebra-section-snake} 节），只须对自由模 $F$
证明左侧的正合性。利用 $P^1_{S/R}(-)$ 与直和可交换，
归结到 $M = S$ 的情形。此情形由例
\ref{algebra-example-derivations-and-differential-operators} 中的讨论得出。
\end{proof}

\begin{remark}
\label{algebra-remark-functoriality-principal-parts}
假设给定环的交换图
$$
\xymatrix{
B \ar[r] & B' \\
A \ar[u] \ar[r] & A' \ar[u]
}
$$
% P01-E0717: the frozen multi-object “Suppose given ...” sentence lacks a comma after the displayed diagram.
一个 $B$-模 $M$、一个 $B'$-模 $M'$，以及一个 $B$-线性映射
$M \to M'$。则得到相容的模映射系统
$$
\xymatrix{
\ldots \ar[r] &
P^2_{B'/A'}(M') \ar[r] &
P^1_{B'/A'}(M') \ar[r] &
P^0_{B'/A'}(M') \\
\ldots \ar[r] &
P^2_{B/A}(M) \ar[r] \ar[u] &
P^1_{B/A}(M) \ar[r] \ar[u] &
P^0_{B/A}(M) \ar[u]
}
$$
% P01-E0718: the frozen system diagram lacks terminal punctuation before the next sentence.
这些映射与这类映射的进一步复合相容。
看出这一点最容易的方法是使用引理
\ref{algebra-lemma-module-principal-parts} 的证明中通过生成元和关系给出的
$P^k_{B/A}(M)$ 的描述；也可以直接从这些模的泛性质看出。
% P01-E0719: the frozen compound sentence lacks a comma before “but”.
此外，这些映射与引理 \ref{algebra-lemma-sequence-of-principal-parts}
的短正合列相容。
\end{remark}

\begin{lemma}
\label{algebra-lemma-principal-parts-base-change}
假设有环映射 $A \to A'$ 和 $A \to B$，并有一个 $B$-模 $M$。
置 $B' = B \otimes_A A'$，并把 $M' = M \otimes_A A'$
视为 $B'$-模。评注 \ref{algebra-remark-functoriality-principal-parts} 的映射
诱导同构 $P^k_{B/A}(M) \otimes_A A' = P^k_{B'/A'}(M')$。
\end{lemma}

\begin{proof}
令 $D_{univ} : M \to P^k_{B/A}(M)$ 是泛微分算子。
诱导的 $A'$-线性映射
$D_{univ} \otimes 1 : M' = M \otimes_A A' \to P^k_{B/A}(M) \otimes_A A'$
容易验证为关于 $M'/B'/A'$ 的 $k$ 阶微分算子。
我们将证明 $D_{univ} \otimes 1$ 满足泛性质。
设 $D' : M' \to N'$ 是一个取值于 $B'$-模 $N'$ 的 $k$ 阶微分算子。
则复合 $D : M \to M' \to N'$ 是微分算子，
其中把其靶模 $N'$ 视为 $B$-模。因此存在 $B$-线性映射
$\gamma : P^k_{B/A}(M) \to N'$ 使得
$D = \gamma \circ D_{univ}$。读者容易验证，诱导的 $B'$-线性映射
$\gamma' : P^k_{B/A}(M) \otimes_A A' \to N'$ 满足
$\gamma' \circ (D_{univ} \otimes 1) = D'$。
$\gamma'$ 的唯一性证明从略。
\end{proof}

\begin{lemma}
\label{algebra-lemma-principal-parts-diagonal}
设 $R \to S$ 是环映射，$M$ 是 $S$-模。令
$J = \Ker(S \otimes_R S \to S)$ 为乘法映射的核。
存在典范的 $S$-模同构
$$
P^k_{S/R}(M) \longrightarrow (S \otimes_R M)/J^{k + 1}(S \otimes_R M)
$$
其中 $s \in S$ 通过乘以 $s \otimes 1$ 作用在靶上，并且
$k$ 阶泛微分算子对应于映射
$m \mapsto \text{下列对象的类： }1 \otimes m$。
% P01-E0720: the frozen statement omits the verb “corresponds” and mismatches singular/plural in its universal-operator clause.
\end{lemma}

\begin{proof}
考虑映射 $T : M \to S \otimes M$，$m \mapsto 1 \otimes m$。
% P01-E0721: the frozen map writes S tensor M without the base-ring subscript R used by every surrounding tensor product.
% P01-E0722: the frozen sentence also omits punctuation after the displayed map before “Since”.
由于 $T$ 是 $R$-线性的，而且
\begin{align*}
g_0g_1\ldots g_k T(m)
- \sum g_0 \ldots \hat g_i \ldots g_k T(g_im) + \ldots
+(-1)^{k + 1}T(g_0\ldots g_km) \\
=
\prod_{i = 0,\ldots,k} (g_i \otimes 1 - 1 \otimes g_i) 1 \otimes m
\end{align*}
对 $m \in M$ 和 $g_0, \ldots, g_k \in S$ 成立，所以可知引理陈述中的
规则 $m \mapsto \text{下列对象的类： }1 \otimes m$ 是 $k$ 阶微分算子
（见引理 \ref{algebra-lemma-module-principal-parts} 的证明）。
由 $P^k_{S/R}(M)$ 的泛性质，得到引理陈述中的箭头。
另一方面，若 $D : M \to N$ 是 $k$ 阶微分算子，则可考虑
$S$-线性映射
$L : S \otimes_R M \to N$，$g \otimes m \mapsto gD(m)$。
读者用与上述类似的计算以及 $J$ 由形如
$g \otimes 1 - 1 \otimes g$ 的元生成这一事实，可以验证
$L$ 零化 $J^{k + 1}(S \otimes_R M)$。
特别地，把它应用于泛微分算子
$D_{univ} : M \to P^k_{S/R}(M)$，便得到 $S$-线性映射
$(S \otimes_R M)/J^{k + 1}(S \otimes_R M) \to P^k_{S/R}(M)$。
此映射是引理陈述中映射的逆；留给读者验证。
\end{proof}

\begin{lemma}
\label{algebra-lemma-differentials-de-rham-complex-order-1}
设 $A \to B$ 是环映射。微分
$\text{d} : \Omega^i_{B/A}  \to \Omega^{i + 1}_{B/A}$
是 $1$ 阶微分算子。
\end{lemma}

\begin{proof}
给定 $b \in B$，必须证明
$\text{d} \circ b - b \circ \text{d}$ 是线性算子。
因此必须证明
$$
\text{d} \circ b \circ b' - b \circ \text{d} \circ b' -
b' \circ \text{d} \circ b + b' \circ b \circ \text{d} = 0
$$
% P01-E0723: the frozen commutator display lacks terminal punctuation before the next sentence.
为此，只须在 $\Omega^i_{B/A}$ 的加法生成元上检验。
因而只须证明
$$
\text{d}(bb'b_0\text{d}b_1 \wedge \ldots \wedge \text{d}b_i) -
b\text{d}(b'b_0\text{d}b_1 \wedge \ldots \wedge \text{d}b_i) -
b'\text{d}(bb_0\text{d}b_1 \wedge \ldots \wedge \text{d}b_i) +
bb'\text{d}(b_0\text{d}b_1 \wedge \ldots \wedge \text{d}b_i)
$$
为零。这是使用 Leibniz 法则的一项愉快计算，留给读者完成。
\end{proof}

\begin{lemma}
\label{algebra-lemma-check-differential-operators}
设 $A \to B$ 是环映射。设 $g_i \in B$（$i \in I$）是
$B$ 作为 $A$-代数的一组生成元。设 $M, N$ 是 $B$-模，
$D : M \to N$ 是 $A$-线性映射。要证明 $D$ 是 $k$ 阶微分算子，
只须证明对 $i \in I$，$D \circ g_i - g_i \circ D$
是 $k - 1$ 阶微分算子。
% P01-E0724: the frozen criterion omits k >= 1, although order k - 1 is undefined by the preceding definition when k = 0.
\end{lemma}

\begin{proof}
也就是说，我们断言满足
$D \circ g - g \circ D$ 是 $k - 1$ 阶微分算子的元
$g \in B$ 构成 $B$ 的一个 $A$-子代数。这由关系
$$
D \circ (g + g') - (g + g') \circ D =
(D \circ g - g \circ D) + (D \circ g' - g' \circ D)
$$
以及
$$
D \circ gg' - gg' \circ D =
(D \circ g - g \circ D) \circ g' + g \circ (D \circ g' - g' \circ D)
$$
% P01-E0725: the frozen product display lacks terminal punctuation before “Strictly speaking”.
得出。严格来说，为了对乘积得出结论，还要使用引理
\ref{algebra-lemma-composition-differential-operators}。
\end{proof}

\begin{lemma}
\label{algebra-lemma-invert-system-differential-operators}
设 $A \to B$ 是环映射，$M, N$ 是 $B$-模，
$S \subset B$ 是乘法子集。任意 $k$ 阶微分算子
$D : M \to N$ 唯一地延拓为 $k$ 阶微分算子
$E : S^{-1}M \to S^{-1}N$。
\end{lemma}

\begin{proof}
对 $k$ 作归纳。若 $k = 0$，则 $D$ 是 $B$-线性的，
所以由局部化的函子性得到延拓。给定 $b \in B$，
算子 $L_b : m \mapsto D(bm) - bD(m)$ 的阶为 $k - 1$。
因此，由归纳假设，它唯一地延拓为 $k - 1$ 阶微分算子
$E_b : S^{-1}M \to S^{-1}N$。
此外，计算表明
$L_{b'b} = L_{b'} \circ b + b' \circ L_b$，所以由唯一性得到
$E_{b'b} = E_{b'} \circ b + b' \circ E_b$。
% P01-E0728: the frozen sentence lacks punctuation before “hence”.
类似地，得到
$E_{b'} \circ b - b \circ E_{b'} =
E_b \circ b' - b' \circ E_b$。
现在，对 $m \in M$ 和 $g \in S$，置
$$
E(m/g) = (1/g)(D(m) - E_g(m/g))
$$
% P01-E0729: the frozen defining display lacks terminal punctuation before the next sentence.
要证明这是良定义的，只须证明：对 $g' \in S$，
若使用代表元 $g'm/g'g$，便得到同一结果。计算如下：
\begin{align*}
(1/g'g)(D(g'm) - E_{g'g}(g'm/gg'))
& =
(1/gg')(g'D(m) + E_{g'}(m) - E_{g'g}(g'm/gg')) \\
& =
(1/g'g)(g'D(m) - g' E_g(m/g))
\end{align*}
% P01-E0730: the frozen align display lacks a comma before the following “which” clause.
这与先前结果相同。显然 $E$ 是 $R$-线性的，因为
$D$ 和 $E_g$ 都是 $R$-线性的。
% P01-E0726: the frozen proof uses an undefined ring R; the lemma's base ring is A.
取 $g = 1$ 并使用 $E_1 = 0$，看出 $E$ 延拓 $D$。
由引理 \ref{algebra-lemma-check-differential-operators}，要证明 $E$ 是
$k$ 阶微分算子，现在只须证明以下两个算子都是 $k - 1$ 阶
微分算子：对 $b \in B$ 的 $E \circ b - b \circ E$，以及对
$g' \in S$ 的 $E \circ 1/g' - 1/g' \circ E$。
对于前者，选取 $S^{-1}M$ 中的元 $m/g$，并注意
\begin{align*}
E(b m/g) - bE(m/g)
& =
(1/g)(D(bm) - bD(m) - E_g(bm/g) + bE_g(m/g)) \\
& =
(1/g)(L_b(m) - E_b(m) + gE_b(m/g)) \\
& =
E_b(m/g)
\end{align*}
% P01-E0731: the frozen align display lacks a comma before the following “which” clause.
这是 $k - 1$ 阶微分算子。最后，
\begin{align*}
E(m/g'g) - (1/g')E(m/g)
& =
(1/g'g)(D(m) - E_{g'g}(m/g'g)) -
(1/g'g)(D(m) - E_g(m/g)) \\
& =
-(1/g')E_{g'}(m/g'g)
\end{align*}
% P01-E0732: the frozen align display lacks a comma before the following “which” clause.
这也是 $k - 1$ 阶微分算子，因为它是线性映射
（乘以 $1/g'$ 和符号）与 $E_{g'}$ 的复合。
唯一性的证明从略。
\end{proof}

\begin{lemma}
\label{algebra-lemma-base-change-differential-operators}
设 $R \to A$ 和 $R \to B$ 是环映射。设 $M$ 和 $M'$ 是 $A$-模，
$D : M \to M'$ 是关于 $R \to A$ 的 $k$ 阶微分算子。
设 $N$ 是任意 $B$-模。则映射
$$
D \otimes \text{id}_N : M \otimes_R N \to M' \otimes_R N
$$
是关于 $B \to A \otimes_R B$ 的 $k$ 阶微分算子。
\end{lemma}

\begin{proof}
显然 $D' = D \otimes \text{id}_N$ 是 $B$-线性的。
由引理 \ref{algebra-lemma-check-differential-operators}，只须证明
$$
D' \circ a \otimes 1 -  a \otimes 1 \circ D' =
(D \circ a - a \circ D) \otimes \text{id}_N
$$
是 $k - 1$ 阶微分算子；这由对 $k$ 作归纳得出。
% P01-E0727: the frozen proof invokes the k - 1 criterion without treating k = 0, where negative order was never defined.
\end{proof}

\section{朴素余切复形}
\label{algebra-section-netherlander}

\noindent
设 $R \to S$ 是环映射。用 $R[S]$ 表示以元 $s \in S$ 为变量的多项式环。
% P01-E0733: the frozen designation “Denote R[S] the polynomial ring” omits “by/as”.
用 $[s] \in R[S]$ 表示与 $s \in S$ 对应的变量。
于是，$R[S]$ 是自由 $R$-模，其基元为
$[s_1] \ldots [s_n]$，其中 $s_1, \ldots, s_n$
遍历 $S$ 的所有无序元列。
% P01-E0734: the frozen source has subject-verb disagreement in “s_1,...,s_n ranges”.
存在典范满射
\begin{equation}
\label{algebra-equation-canonical-presentation}
R[S] \longrightarrow S,\quad [s] \longmapsto s
\end{equation}
其核记为 $I \subset R[S]$。容易观察到，$I$ 由诸元
$[s + s'] - [s] - [s']$、$[s][s'] - [ss']$ 和 $[r] - r$ 生成。
由引理 \ref{algebra-lemma-differential-seq}，存在典范映射
\begin{equation}
\label{algebra-equation-naive-cotangent-complex}
I/I^2 \longrightarrow \Omega_{R[S]/R} \otimes_{R[S]} S
\end{equation}
其余核典范同构于 $\Omega_{S/R}$。
注意，$S$-模 $\Omega_{R[S]/R} \otimes_{R[S]} S$
以诸元 $\text{d}[s]$ 为自由生成元。

\begin{definition}
\label{algebra-definition-naive-cotangent-complex}
设 $R \to S$ 是环映射。{\it 朴素余切复形} $\NL_{S/R}$
是链复形 (\ref{algebra-equation-naive-cotangent-complex})
$$
\NL_{S/R} = \left(I/I^2 \longrightarrow \Omega_{R[S]/R} \otimes_{R[S]} S\right)
$$
其中把 $I/I^2$ 放在（同调）次数 $1$，把
$\Omega_{R[S]/R} \otimes_{R[S]} S$ 放在次数 $0$。
用 $H_1(L_{S/R}) = H_1(\NL_{S/R})$\footnote{文献中有时把此模记为
$\Gamma_{S/R}$。}表示次数 $1$ 的同调。
% P01-E0735: the frozen notation sentence “denote ... the homology” omits “by/as”.
\end{definition}

\noindent
在继续之前，先略述真正的余切复形（余切，第
\ref{cotangent-section-cotangent-ring-map} 节）。
给定环映射 $R \to S$，存在典范的单纯 $R$-代数 $P_\bullet$，
其各项都是多项式代数，并带有典范同伦等价
$$
P_\bullet \longrightarrow S
$$
% P01-E0736: the frozen display lacks terminal punctuation before the next sentence.
$S$ 在 $R$ 上的余切复形 $L_{S/R}$ 定义为与单纯模
$$
\Omega_{P_\bullet/R} \otimes_{P_\bullet} S
$$
% P01-E0737: the frozen display lacks terminal punctuation before the next sentence.
相伴的链复形。上面定义的朴素余切复形典范同构于截断
$\tau_{\leq 1}L_{S/R}$（见同调，第
\ref{homology-section-truncations} 节，以及余切，第
\ref{cotangent-section-surjections} 节）。
特别地，确有 $H_1(\NL_{S/R}) = H_1(L_{S/R})$，所以我们的定义
与使用余切复形的定义相容。此外，如上所见，
$H_0(L_{S/R}) = H_0(\NL_{S/R}) = \Omega_{S/R}$。

\medskip\noindent
设 $R \to S$ 是环映射。$S$ 在 $R$ 上的一个{\it 表现}
是 $R$-代数满射 $\alpha : P \to S$，其中 $P$ 是一个
（以某变量集为变量的）多项式代数。通常，当 $S$ 在 $R$ 上有限型时，
我们会说：“设 $R[x_1, \ldots, x_n] \to S$ 是 $S/R$ 的一个表现”；
若还想指明 $I$ 是表现的核，则说：“设
$0 \to I \to R[x_1, \ldots, x_n] \to S \to 0$ 是 $S/R$ 的一个表现”。
注意，用来定义朴素余切复形的映射 $R[S] \to S$ 是一个表现的例子。

\medskip\noindent
注意，对每个表现 $\alpha$，都得到一个两项 $S$-模链复形
% P01-E0738: the frozen compound modifier “two term” should be hyphenated.
$$
\NL(\alpha) : I/I^2 \longrightarrow \Omega_{P/R} \otimes_P S.
$$
这里把项 $I/I^2$ 放在次数 $1$，把项
$\Omega_{P/R} \otimes S$ 放在次数 $0$。
$I/I^2$ 中 $f \in I$ 的类映到
$\Omega_{P/R} \otimes S$ 中的 $\text{d}f \otimes 1$。
% P01-E0739: the frozen two subsequent tensor products omit the base P shown in the defining display.
此复形的余核典范地为 $\Omega_{S/R}$，见引理
\ref{algebra-lemma-differential-seq}。把复形 $\NL(\alpha)$ 称为
{\it 与 $S/R$ 的表现 $\alpha : P \to S$ 相伴的朴素余切复形}。
注意，若 $P = R[S]$ 且取其到 $S$ 的典范满射，则恢复
$\NL_{S/R}$。若 $P = R[x_1, \ldots, x_n]$，则有时用记号
$I/I^2 \to \bigoplus_{i = 1, \ldots, n} S\text{d}x_i$
表示此复形。
% P01-E0740: the frozen phrase “then will sometimes use” omits the subject “we”.

\medskip\noindent
假设给定环的交换图
\begin{equation}
\label{algebra-equation-functoriality-NL}
\vcenter{
\xymatrix{
S \ar[r]_{\phi} & S' \\
R \ar[r] \ar[u] & R' \ar[u]
}
}
\end{equation}
设 $\alpha : P \to S$ 是 $S$ 在 $R$ 上的表现，
$\alpha' : P' \to S'$ 是 $S'$ 在 $R'$ 上的表现。
从 $\alpha : P \to S$ 到 $\alpha' : P' \to S'$ 的一个
{\it 表现态射}定义为 $R$-代数映射
$$
\varphi : P \to P'
$$
使得 $\phi \circ \alpha = \alpha' \circ \varphi$。
注意，这时 $\varphi(I) \subset I'$，其中
$I = \Ker(\alpha)$ 且 $I' = \Ker(\alpha')$。
因此，$\varphi$ 诱导 $S$-模映射
$I/I^2 \to I'/(I')^2$，并由微分的函子性还诱导 $S$-模映射
$\Omega_{P/R} \otimes S \to \Omega_{P'/R'} \otimes S'$。
% P01-E0741: the frozen functoriality formula omits the tensor bases P and P'.
这些映射与 $\NL(\alpha)$ 和 $\NL(\alpha')$ 的微分相容，
从而得到朴素余切复形的映射
$$
\NL(\alpha) \longrightarrow \NL(\alpha').
$$
考虑诱导映射 $\NL(\alpha) \otimes_S S' \to \NL(\alpha')$
往往很方便。

\medskip\noindent
在特殊情形 $P = R[S]$、$P' = R'[S']$ 下，映射
$\phi : S \to S'$ 通过规则 $[s] \mapsto [\phi(s)]$
诱导典范环映射 $\varphi : P \to P'$。
因此上述构造确定了与图 (\ref{algebra-equation-functoriality-NL}) 相伴的
典范（！）链复形映射
$$
\NL_{S/R} \longrightarrow \NL_{S'/R'},\quad\text{且}\quad
\NL_{S/R} \otimes_S S' \longrightarrow \NL_{S'/R'}
$$
注意，此构造与复合相容：给定交换图
$$
\xymatrix{
S \ar[r]_{\phi} & S' \ar[r]_{\phi'} & S'' \\
R \ar[r] \ar[u] & R' \ar[u] \ar[r] & R'' \ar[u]
}
$$
% P01-E0742: the frozen diagram lacks a comma before the continuing “we see” clause.
可见，复合
$$
\NL_{S/R} \longrightarrow \NL_{S'/R'} \longrightarrow \NL_{S''/R''}
$$
就是外方块给出的映射 $\NL_{S/R} \to \NL_{S''/R''}$。

\medskip\noindent
事实证明，$\NL(\alpha)$ 与 $\NL_{S/R}$ 同伦等价，
且上面构造的映射在同伦意义下良定义（复形映射的同伦见同调，第
\ref{homology-section-complexes} 节；但在下面引理的证明中，
我们也明确说明这些陈述的确切含义）。

\begin{lemma}
\label{algebra-lemma-NL-homotopy}
假设给定图 (\ref{algebra-equation-functoriality-NL})。
设 $\alpha : P \to S$ 和 $\alpha' : P' \to S'$ 是表现。
\begin{enumerate}
\item 存在一个从 $\alpha$ 到 $\alpha'$ 的表现态射。
\item 任意两个表现态射诱导同伦的复形态射
$\NL(\alpha) \to \NL(\alpha')$。
\item 该构造与表现态射的复合相容（确切陈述见证明）。
\item 若 $R \to R'$ 和 $S \to S'$ 是同构，则对任意从
$\alpha$ 到 $\alpha'$ 的表现映射 $\varphi$，诱导映射
$\NL(\alpha) \to \NL(\alpha')$ 是同伦等价和拟同构。
\end{enumerate}
特别地，把 $\alpha$ 与典范表现
(\ref{algebra-equation-canonical-presentation}) 比较，得到一个在同伦意义下
良定义并与所有函子性相容（不计同伦）的拟同构
$\NL(\alpha) \to \NL_{S/R}$。
\end{lemma}

\begin{proof}
由于 $P$ 是 $R$ 上的多项式代数，可对某个集合 $A$ 写成
$P = R[x_a, a \in A]$。由于 $\alpha'$ 是满射，可对每个
$a \in A$ 选取 $f_a \in P'$，使得
$\alpha'(f_a) = \phi(\alpha(x_a))$。令
$\varphi : P = R[x_a, a \in A] \to P'$ 为满足
$\varphi(x_a) = f_a$ 的唯一 $R$-代数映射。
这给出 (1) 中的态射。

\medskip\noindent
设 $\varphi$ 和 $\varphi'$ 是从 $\alpha$ 到 $\alpha'$ 的表现态射。
% P01-E0743: the frozen source omits “be” in “Let phi and phi' morphisms”.
令 $I = \Ker(\alpha)$ 且 $I' = \Ker(\alpha')$。
必须构造图中的对角映射 $h$：
$$
\xymatrix{
I/I^2 \ar[r]^-{\text{d}}
\ar@<1ex>[d]^{\varphi'_1} \ar@<-1ex>[d]_{\varphi_1}
&
\Omega_{P/R} \otimes_P S
\ar@<1ex>[d]^{\varphi'_0} \ar@<-1ex>[d]_{\varphi_0}
\ar[ld]_h
\\
I'/(I')^2 \ar[r]^-{\text{d}}
&
\Omega_{P'/R'} \otimes_{P'} S'
}
$$
其中竖直映射由 $\varphi$、$\varphi'$ 诱导，并且该对角映射满足
% P01-E0744: the frozen “where ... induced by phi, phi' such that” clause omits a conjunction and the subject h.
$$
\varphi_1 - \varphi'_1 = h \circ \text{d}
\quad\text{且}\quad
\varphi_0 - \varphi'_0 = \text{d} \circ h
$$
考虑映射 $\varphi - \varphi' : P \to P'$。
由于 $\varphi$ 和 $\varphi'$ 都与 $\alpha$、$\alpha'$ 相容，
得到 $\varphi - \varphi' : P \to I'$。
这说明 $\varphi, \varphi' : P \to P'$ 在 $I'/(I')^2$ 上诱导同一
$P$-模结构，因为
$\varphi(p)i' - \varphi'(p)i' = (\varphi - \varphi')(p)i' \in (I')^2$。
此外，$\varphi - \varphi'$ 是 $R$-线性的，并且
$$
(\varphi - \varphi')(fg) =
\varphi(f)(\varphi - \varphi')(g) + (\varphi - \varphi')(f)\varphi'(g)
$$
所以诱导映射 $D : P \to I'/(I')^2$ 是一个 $R$-导子。
% P01-E0745: the frozen source says “a R-derivation” instead of “an R-derivation”.
从而得到典范映射
$h : \Omega_{P/R} \otimes_P S \to I'/(I')^2$，
使得 $D = h \circ \text{d}$。计算（从略）表明 $h$ 正是所需同伦。

\medskip\noindent
假设有交换图
$$
\xymatrix{
S \ar[r]_{\phi} & S' \ar[r]_{\phi'} & S'' \\
R \ar[r] \ar[u] & R' \ar[u] \ar[r] & R'' \ar[u]
}
$$
并且
\begin{enumerate}
\item $\alpha : P \to S$，
\item $\alpha' : P' \to S'$，以及
\item $\alpha'' : P'' \to S''$
\end{enumerate}
是表现。假设
\begin{enumerate}
\item $\varphi : P \to P'$ 是从 $\alpha$ 到 $\alpha'$ 的表现态射，并且
\item $\varphi' : P' \to P''$ 是从 $\alpha'$ 到 $\alpha''$ 的表现态射。
\end{enumerate}
则立即可见，$\varphi' \circ \varphi : P \to P''$
是从 $\alpha$ 到 $\alpha''$ 的表现态射，并且由
$\varphi$ 和 $\varphi'$ 诱导的映射
$\NL(\alpha) \to \NL(\alpha')$ 与
$\NL(\alpha') \to \NL(\alpha'')$ 的复合，正是诱导的朴素余切复形映射
$\NL(\alpha) \to \NL(\alpha'')$。

\medskip\noindent
在两项复形的简单情形下，拟同构就是在两项间映射的核和余核上都诱导
同构的映射。
% P01-E0746: the frozen compound modifier “2 term” should be hyphenated.
注意，同伦的两项复形映射（如上所述）在核和余核上定义相同的映射。
因此，若 $\varphi$ 是从 $S$ 在 $R$ 上的表现 $\alpha$
到自身的映射，则诱导映射
$\NL(\alpha) \to \NL(\alpha)$ 是拟同构，因为由 (2) 它同伦于恒等映射。
% P01-E0747: the frozen phrase “is a quasi-isomorphism being homotopic” uses an awkward/dangling participial construction.
为了在一般情形下证明 (4)，考虑由 (1) 存在的从
$\alpha'$ 到 $\alpha$ 的态射 $\varphi'$。复合
$\NL(\alpha) \to \NL(\alpha') \to \NL(\alpha)$ 和
$\NL(\alpha') \to \NL(\alpha) \to \NL(\alpha')$
由 (3) 同伦于恒等映射，因此按照定义，这些映射是同伦等价。
% P01-E0748: compatibility in (3) alone does not give homotopy to the identity; the argument also needs part (2).
形式上立即推出，两个映射
$\NL(\alpha) \to \NL(\alpha')$ 和
$\NL(\alpha') \to \NL(\alpha)$ 都是拟同构。略去一些细节。
\end{proof}

\begin{lemma}
\label{algebra-lemma-NL-polynomial-algebra}
设 $A \to B$ 是多项式代数。则 $\NL_{B/A}$ 与链复形
$(0 \to \Omega_{B/A})$ 同伦等价，其中 $\Omega_{B/A}$ 位于次数 $0$。
\end{lemma}

\begin{proof}
这由引理 \ref{algebra-lemma-NL-homotopy} 以及如下事实得出：
$\text{id}_B : B \to B$ 是核为零的 $B$ 在 $A$ 上的表现。
% P01-E0749: the frozen proof is the sentence fragment “Follows from ...”.
\end{proof}

\noindent
下面的引理是引入朴素余切复形的部分动机。
余切复形把它扩张为真正的长正合上同调列。
若 $B \to C$ 是局部完全交，则可在序列左侧补上一个零；
见代数进阶，引理 \ref{more-algebra-lemma-transitive-lci-at-end}。

\begin{lemma}[Jacobi--Zariski 序列]
\label{algebra-lemma-exact-sequence-NL}
设 $A \to B \to C$ 是环映射。选取核为 $I$ 的表现
$\alpha : A[x_s, s \in S] \to B$，并选取核为 $J$ 的表现
$\beta : B[y_t, t \in T] \to C$。令
$\gamma : A[x_s, y_t] \to C$ 为所诱导的 $C$ 的表现，其核为 $K$。
则得到典范交换图
$$
\xymatrix{
0 \ar[r] &
\Omega_{A[x_s]/A} \otimes C \ar[r] &
\Omega_{A[x_s, y_t]/A} \otimes C \ar[r] &
\Omega_{B[y_t]/B} \otimes C \ar[r] &
0 \\
&
I/I^2 \otimes C \ar[r] \ar[u] &
K/K^2 \ar[r] \ar[u] &
J/J^2 \ar[r] \ar[u] &
0
}
$$
其两行正合。得到如下同调群正合列
$$
H_1(\NL_{B/A} \otimes_B C) \to
H_1(L_{C/A}) \to
H_1(L_{C/B}) \to
C \otimes_B \Omega_{B/A} \to
\Omega_{C/A} \to
\Omega_{C/B} \to 0
$$
它是扩张引理 \ref{algebra-lemma-exact-sequence-differentials} 中序列的
$C$-模序列。若
$\text{Tor}_1^B(\Omega_{B/A}, C) = 0$ 且
$\text{Tor}_2^B(\Omega_{B/A}, C) = 0$，则
$H_1(\NL_{B/A} \otimes_B C) = H_1(L_{B/A}) \otimes_B C$。
\end{lemma}

\begin{proof}
映射的精确定义从略。上行的正合性来自如下事实：
$\text{d}x_s$、$\text{d}y_t$ 构成中间模的一组基。
映射 $\gamma$ 分解为
$$
A[x_s, y_t] \to B[y_t] \to C
$$
其中第一支箭头是满射，第二支箭头等于 $\beta$。
因此看出 $K \to J$ 是满射。此外，第一个显示箭头的核是
$IA[x_s, y_t]$。故 $I/I^2 \otimes C$ 满射到
$K/K^2 \to J/J^2$ 的核上。最后，可用引理
\ref{algebra-lemma-NL-homotopy} 把各项认作朴素余切复形的同调群。

\medskip\noindent
最后一个断言是同调代数中的陈述。回忆
$\NL_{B/A} = (N^{-1} \to N^0)$ 是一个两项 $B$-模复形，
其中 $N^0$ 自由，而上同调模为
$H^0 = \Omega_{B/A}$ 和 $H^{-1} = H_1(L_{B/A})$。
记 $M \subset N^0$ 为微分的像。若
$\text{Tor}_1^B(H^0, C) = 0$，则有正合列
$$
0 \to M \otimes_B C \to N^0 \otimes_B C \to H^0 \otimes_B C \to 0
$$
% P01-E0750: the frozen display lacks terminal punctuation before the next sentence.
由于 $N^0$ 自由，还看出
$\text{Tor}_2^B(H^0, C) = \text{Tor}_1^B(M, C)$。
因此，若 $\text{Tor}_2^B(H^0, C) = 0$，则还有正合列
$$
0 \to H^{-1} \otimes_B C \to N^{-1} \otimes_B C \to M \otimes_B C \to 0
$$
% P01-E0751: the frozen display lacks terminal punctuation before the next sentence.
综合起来可见，若 $\text{Tor}_1^B(H^0, C) = 0$ 且
$\text{Tor}_2^B(H^0, C) = 0$，则
$H^{-1} \otimes_B C$ 正是
$N^{-1} \otimes_B C \to N^0 \otimes_B C$ 的核，正如所需。
\end{proof}

\begin{remark}
\label{algebra-remark-composition-homotopy-equivalent-to-zero}
设 $A \to B$ 和 $\phi : B \to C$ 是环映射。
则复合 $\NL_{B/A} \to \NL_{C/A} \to \NL_{C/B}$ 同伦于零。
具体而言，此复合是朴素余切复形关于方块
$$
\xymatrix{
B \ar[r]_\phi & C \\
A \ar[r] \ar[u] & B \ar[u]
}
$$
% P01-E0752: the frozen diagram lacks terminal punctuation before the next sentence.
的函子性。写 $J = \Ker(B[C] \to C)$。
一个显式同伦由映射
$\Omega_{A[B]/A} \otimes_{A[B]} B \to J/J^2$ 给出，
它把基元 $\text{d}[b]$ 映到 $J/J^2$ 中
$[\phi(b)] - b$ 的类。
\end{remark}

\begin{lemma}
\label{algebra-lemma-NL-surjection}
设 $A \to B$ 是核为 $I$ 的满环映射。
则 $\NL_{B/A}$ 与链复形 $(I/I^2 \to 0)$ 同伦等价，
其中 $I/I^2$ 位于次数 $1$。特别地，
$H_1(L_{B/A}) = I/I^2$。
\end{lemma}

\begin{proof}
这由引理 \ref{algebra-lemma-NL-homotopy} 以及如下事实得出：
$A \to B$ 是 $B$ 在 $A$ 上的一个表现。
% P01-E0753: the frozen proof is the sentence fragment “Follows from ...”.
\end{proof}

\begin{lemma}
\label{algebra-lemma-application-NL}
设 $A \to B \to C$ 是环映射。假设 $A \to C$ 是满射
（所以 $B \to C$ 也是满射）。记
$I = \Ker(A \to C)$ 且 $J = \Ker(B \to C)$。
则序列
$$
I/I^2 \to J/J^2 \to \Omega_{B/A} \otimes_B B/J \to 0
$$
是正合的。
\end{lemma}

\begin{proof}
这由引理 \ref{algebra-lemma-exact-sequence-NL} 以及引理
\ref{algebra-lemma-NL-surjection} 中对朴素余切复形
$\NL_{C/B}$ 和 $\NL_{C/A}$ 的描述得出。
% P01-E0754: the frozen proof is the sentence fragment “Follows from ...”.
\end{proof}

\begin{lemma}[平坦基变换]
\label{algebra-lemma-change-base-NL}
设 $R \to S$ 是环映射，$\alpha : P \to S$ 是一个表现，
$R \to R'$ 是平坦环映射。令
$\alpha' : P \otimes_R R' \to S' = S \otimes_R R'$
为诱导表现。则
$\NL(\alpha) \otimes_R R' = \NL(\alpha) \otimes_S S' = \NL(\alpha')$。
特别地，若 $R \to R'$ 平坦，则典范映射
$$
\NL_{S/R} \otimes_S S' \longrightarrow \NL_{S \otimes_R R'/R'}
$$
是同伦等价。
\end{lemma}

\begin{proof}
这是因为 $R \to R'$ 平坦，故
$\Ker(\alpha') = R' \otimes_R \Ker(\alpha)$。
\end{proof}

\begin{lemma}
\label{algebra-lemma-colimits-NL}
设 $R_\lambda \to S_\lambda$ 是有向集 $\Lambda$ 上的环映射系统。
置 $R = \colim R_\lambda$ 且 $S = \colim S_\lambda$。
则 $\NL_{S/R} = \colim \NL_{S_\lambda/R_\lambda}$。
\end{lemma}

\begin{proof}
回忆 $\NL_{S/R}$ 是复形
$I/I^2 \to \bigoplus_{s \in S} S\text{d}[s]$，其中
$I \subset R[S]$ 是典范表现 $R[S] \to S$ 的核。
现在显然 $R[S] = \colim R_\lambda[S_\lambda]$；类似地，
$I = \colim I_\lambda$，其中
$I_\lambda = \Ker(R_\lambda[S_\lambda] \to S_\lambda)$。
引理即得。
\end{proof}

\begin{lemma}
\label{algebra-lemma-NL-of-localization}
若 $S \subset A$ 是 $A$ 的一个乘法子集，则
$\NL_{S^{-1}A/A}$ 与零复形同伦等价。
\end{lemma}

\begin{proof}
由于 $A \to S^{-1}A$ 平坦，由平坦基变换
（引理 \ref{algebra-lemma-change-base-NL}）可知
$\NL_{S^{-1}A/A} \otimes_A S^{-1}A \to \NL_{S^{-1}A/S^{-1}A}$
是同伦等价。箭头的源同构于 $\NL_{S^{-1}A/A}$，
而箭头的目标为零（由引理 \ref{algebra-lemma-NL-surjection}），故得证。
\end{proof}

\begin{lemma}
\label{algebra-lemma-NL-localize-bottom}
设 $S \subset A$ 是 $A$ 的一个乘法子集。
设 $S^{-1}A \to B$ 是环映射。
则 $\NL_{B/A} \to \NL_{B/S^{-1}A}$ 是同伦等价。
% P01-E0755: the frozen statement has “Let S ... is” instead of “be”.
\end{lemma}

\begin{proof}
选取 $B$ 在 $A$ 上的一个表现 $\alpha : P \to B$。
则 $\beta : S^{-1}P \to B$ 是 $B$ 在 $S^{-1}A$ 上的一个表现。
直接计算表明 $\NL(\alpha) = \NL(\beta)$；又由引理
\ref{algebra-lemma-NL-homotopy}，朴素余切复形在同伦意义下良定义，
故引理成立。
% P01-E0756: the frozen concluding clause is weakly joined by an unpunctuated “as”.
\end{proof}

\begin{lemma}
\label{algebra-lemma-principal-localization-NL}
\begin{slogan}
朴素余切复形的构造与在一个元素处的局部化交换。
\end{slogan}
设 $A \to B$ 是环映射。设 $g \in B$。假设 $\alpha : P \to B$
是核为 $I$ 的一个表现。则 $B_g$ 在 $A$ 上的一个表现是映射
$$
\beta : P[x] \longrightarrow B_g
$$
它延拓 $\alpha$ 并把 $x$ 映到 $1/g$。
$\beta$ 的核 $J$ 由 $I$ 和元素 $f x - 1$ 生成，
其中 $f \in P$ 是由 $\alpha$ 映到 $g \in B$ 的一个元素。在此情形下，
\begin{enumerate}
\item $J/J^2 = (I/I^2)_g \oplus B_g (f x - 1)$,
\item $\Omega_{P[x]/A} \otimes_{P[x]} B_g =
\Omega_{P/A} \otimes_P B_g \oplus B_g \text{d}x$,
\item $\NL(\beta) \cong
\NL(\alpha) \otimes_B B_g \oplus (B_g \xrightarrow{g} B_g)$
\end{enumerate}
% P01-E0757: the frozen enumeration has no terminal punctuation.
因此典范映射 $\NL_{B/A} \otimes_B B_g \to \NL_{B_g/A}$
是同伦等价。
\end{lemma}

\begin{proof}
由于 $P[x]/(I, fx - 1) = B[x]/(gx - 1) = B_g$，便得到关于
$I$ 和 $fx - 1$ 生成 $J$ 的断言。考虑交换图
$$
\xymatrix{
0 \ar[r] &
\Omega_{P/A} \otimes B_g \ar[r] &
\Omega_{P[x]/A} \otimes B_g \ar[r] &
\Omega_{B[x]/B} \otimes B_g \ar[r] &
0 \\
&
(I/I^2)_g \ar[r] \ar[u] &
J/J^2 \ar[r] \ar[u] &
(gx - 1)/(gx - 1)^2 \ar[r] \ar[u] &
0
}
$$
其两行由引理 \ref{algebra-lemma-exact-sequence-NL} 知为正合。
% P01-E0758: several tensor products in the frozen diagram and prose omit their inferable base rings.
$B_g$-模 $\Omega_{B[x]/B} \otimes B_g$ 是以
$\text{d}x$ 为基的秩 $1$ 自由模。
$B_g$-模 $\Omega_{P[x]/A} \otimes B_g$ 中的元素 $\text{d}x$
给出上面一行的一个分裂。元素 $gx - 1 \in (gx - 1)/(gx - 1)^2$
被映到 $\Omega_{B[x]/B} \otimes B_g$ 中的 $g\text{d}x$，
因而 $(gx - 1)/(gx - 1)^2$ 是秩为 $1$ 的自由 $B_g$-模。
（也可如下说明：$gx - 1$ 在 $B[x]$ 中是非零因子，
因为它是常数项可逆的多项式；而由引理
\ref{algebra-lemma-regular-quasi-regular}，任何非零因子均给出长度为 $1$
的拟正则序列。）

\medskip\noindent
我们来证明 $(I/I^2)_g \to J/J^2$ 是单射。考虑 $P$-代数映射
$$
\pi : P[x] \to (P/I^2)_f = P_f/I_f^2
$$
它把 $x$ 映到 $1/f$。由于 $J$ 由 $I$ 和 $fx - 1$ 生成，
可知 $\pi(J) \subset (I/I^2)_f = (I/I^2)_g$。由于后者是平方为零的理想，
可知 $\pi(J^2) = 0$。
% P01-E0759: the frozen sentence “Let us prove ... injective” omits “is”.
若 $a \in I$ 在 $J$ 中映到 $J^2$ 的一个元素，则
$\pi(a) = 0$，这蕴含 $a$ 在 $I_f/I_f^2$ 中的像为零。
% P01-E0760: the frozen phrase “maps to an element of J^2 in J” is grammatically malformed.
这证明了所需的单射性。

\medskip\noindent
因此有二项复形的短正合序列
% P01-E0761: the frozen compound adjective “two term” is unhyphenated.
$$
0 \to \NL(\alpha) \otimes_B B_g \to \NL(\beta)
\to (B_g \xrightarrow{g} B_g) \to 0
$$
这样的短正合序列在复形范畴中总可分裂。
% P01-E0762: the frozen blanket splitting claim is false for arbitrary short exact sequences of complexes; the explicit termwise splittings below justify only this case.
在本例中，可取如下分裂：
$$
J/J^2 = (I/I^2)_g \oplus B_g (fx - 1)\quad\text{且}\quad
\Omega_{P[x]/A} \otimes B_g = \Omega_{P/A} \otimes B_g \oplus
B_g (g^{-2}\text{d}f + \text{d}x)
$$
% P01-E0763: the frozen display lacks terminal punctuation before “This works”.
这是因为在 $\Omega_{P[x]/A} \otimes B_g$ 中，
$\text{d}(fx - 1) = x\text{d}f + f \text{d}x =
g(g^{-2}\text{d}f + \text{d}x)$。
\end{proof}

\begin{lemma}
\label{algebra-lemma-localize-NL}
设 $A \to B$ 是环映射。设 $S \subset B$ 是乘法子集。
典范映射 $\NL_{B/A} \otimes_B S^{-1}B \to \NL_{S^{-1}B/A}$
是拟同构。
\end{lemma}

\begin{proof}
我们有 $S^{-1}B = \colim_{g \in S} B_g$，这里把 $S$
视为有向集（以整除关系为序）；见引理
\ref{algebra-lemma-localization-colimit}。
由引理 \ref{algebra-lemma-principal-localization-NL}，每个映射
$\NL_{B/A} \otimes_B B_g \to \NL_{B_g/A}$
都是拟同构。
% P01-E0764: the frozen subject “each” incorrectly takes the plural verb “are”.
本引理由引理 \ref{algebra-lemma-colimits-NL} 得出。
\end{proof}

\begin{lemma}
\label{algebra-lemma-sum-two-terms}
设 $R$ 是环。
设 $A_1 \to A_0$ 和 $B_1 \to B_0$ 是二项复形。
% P01-E0765: the frozen sentence has both a spurious comma after the first complex and an unhyphenated “two term”.
假设存在复形态射 $\varphi : A_\bullet \to B_\bullet$
和 $\psi : B_\bullet \to A_\bullet$，使得
$\varphi \circ \psi$ 和 $\psi \circ \varphi$
均与恒等映射同伦。
则作为 $R$-模，$A_1 \oplus B_0 \cong B_1 \oplus A_0$。
\end{lemma}

\begin{proof}
选取映射 $h : A_0 \to A_1$，使得
$$
\text{id}_{A_1} - \psi_1 \circ \varphi_1 = h \circ d_A
\text{ 且 }
\text{id}_{A_0} - \psi_0 \circ \varphi_0 = d_A \circ h.
$$
类似地，选取映射 $h' : B_0 \to B_1$，使得
$$
\text{id}_{B_1} - \varphi_1 \circ \psi_1 = h' \circ d_B
\text{ 且 }
\text{id}_{B_0} - \varphi_0 \circ \psi_0 = d_B \circ h'.
$$
直接计算表明
$$
\left(
\begin{matrix}
\text{id}_{A_1} & -h' \circ \psi_1 + h \circ \psi_0 \\
0 & \text{id}_{B_0}
\end{matrix}
\right)
=
\left(
\begin{matrix}
\psi_1 & h \\
-d_B & \varphi_0
\end{matrix}
\right)
\left(
\begin{matrix}
\varphi_1 & - h' \\
d_A & \psi_0
\end{matrix}
\right)
$$
% P01-E0766: the frozen matrix display lacks terminal punctuation.
这表明右边的两个矩阵都可逆，从而证明了引理。
% P01-E0767: invertibility of the displayed product alone does not imply both factors are invertible over arbitrary modules; the reverse product or explicit inverses are not supplied.
\end{proof}

\begin{lemma}
\label{algebra-lemma-conormal-module}
设 $R \to S$ 是有限型环映射。
对任意下列两个表现：
$$
\alpha : R[x_1, \ldots, x_n] \to S
$$
以及
$$
\beta : R[y_1, \ldots, y_m] \to S
$$
均有
% P01-E0768: the frozen statement has a spurious comma before “and” between the two presentations.
$$
I/I^2 \oplus S^{\oplus m} \cong J/J^2 \oplus S^{\oplus n}
$$
作为 $S$-模成立，其中 $I = \Ker(\alpha)$ 且 $J = \Ker(\beta)$。
% P01-E0769: the frozen phrase “as S-modules where” lacks a comma before the defining clause.
\end{lemma}

\begin{proof}
见引理 \ref{algebra-lemma-NL-homotopy} 和 \ref{algebra-lemma-sum-two-terms}。
\end{proof}

\begin{lemma}
\label{algebra-lemma-conormal-module-localize}
设 $R \to S$ 是有限型环映射。
设 $g \in S$。对任意表现
$\alpha : R[x_1, \ldots, x_n] \to S$ 和
$\beta : R[y_1, \ldots, y_m] \to S_g$，都有
% P01-E0770: the frozen statement has a spurious comma before “and” between the two presentations.
$$
(I/I^2)_g \oplus S^{\oplus m}_g \cong J/J^2 \oplus S_g^{\oplus n}
$$
作为 $S_g$-模成立，其中
$I = \Ker(\alpha)$ 且 $J = \Ker(\beta)$。
% P01-E0771: the frozen phrase “as S_g-modules where” lacks a comma before the defining clause.
\end{lemma}

\begin{proof}
令 $\beta' : R[x_1, \ldots, x_n, x] \to S_g$ 为从 $\alpha$
出发、按引理 \ref{algebra-lemma-principal-localization-NL} 构造的表现。
于是已知 $\NL(\alpha) \otimes_S S_g$ 与 $\NL(\beta')$ 同伦等价。
由引理 \ref{algebra-lemma-NL-homotopy}，$\NL(\beta)$ 与
$\NL(\beta')$ 同伦等价。因此 $\NL(\alpha) \otimes_S S_g$
与 $\NL(\beta)$ 同伦等价。最后应用引理
\ref{algebra-lemma-conormal-module}。
\end{proof}

\section{局部完全交}
\label{algebra-section-lci}

\noindent
作为局部完全交这一性质是 Noether 局部环的一个内蕴性质。
这将在《除幂代数》的第 \ref{dpa-section-lci} 节讨论。
不过眼下我们只对域上的有限型代数定义这一性质。

\begin{definition}
\label{algebra-definition-lci-field}
设 $k$ 是域。
设 $S$ 是有限型 $k$-代数。
\begin{enumerate}
\item 如果存在一个表现
$S = k[x_1, \ldots, x_n]/(f_1, \ldots, f_c)$，使得
$\dim(S) = n - c$，则称 $S$ 是一个\textit{$k$ 上的全局完全交}。
\item 如果存在覆盖 $\Spec(S) = \bigcup D(g_i)$，使得每个环
$S_{g_i}$ 都是 $k$ 上的全局完全交，则称 $S$ 是一个
\textit{$k$ 上的局部完全交}。
\end{enumerate}
我们还约定零环是 $k$ 上的全局完全交。
\end{definition}

\noindent
假设 $S$ 是如定义 \ref{algebra-definition-lci-field} 所述的全局完全交
$S = k[x_1, \ldots, x_n]/(f_1, \ldots, f_c)$。
对极大理想 $\mathfrak m \subset k[x_1, \ldots, x_n]$，
有 $\dim(k[x_1, \ldots, x_n]_\mathfrak m) = n$
（引理 \ref{algebra-lemma-dim-affine-space}）。
若 $(f_1, \ldots, f_c) \subset \mathfrak m$，则由引理
\ref{algebra-lemma-one-equation} 可知 $\dim(S_\mathfrak m) \geq n - c$。
% P01-E0772: the frozen paragraph indexes S by a maximal ideal of the polynomial ring without first replacing it by the corresponding maximal ideal of S.
由于定义 \ref{algebra-definition-lci-field} 给出 $\dim(S) = n - c$，
便得对 $S$ 的所有极大理想都有 $\dim(S_\mathfrak m) = n - c$，
而且 $\Spec(S)$ 是维数为 $n - c$ 的等维空间
（《拓扑》，定义 \ref{topology-definition-equidimensional}）；见引理
\ref{algebra-lemma-dimension-at-a-point-finite-type-over-field}。
下文会经常使用这一点而不再说明。

\begin{lemma}
\label{algebra-lemma-localize-lci}
设 $k$ 是域。
设 $S$ 是有限型 $k$-代数。
设 $g \in S$。
\begin{enumerate}
\item 若 $S$ 是全局完全交，则 $S_g$ 也是。
\item 若 $S$ 是局部完全交，则 $S_g$ 也是。
\end{enumerate}
\end{lemma}

\begin{proof}
第二个断言立即由第一个断言得出。现证明第一个断言。
% P01-E0773: the frozen “Proof of the first statement.” is a sentence fragment.
若 $S_g$ 是零环，断言成立。假设 $S_g$ 非零。
按定义 \ref{algebra-definition-lci-field}，写成
$S = k[x_1, \ldots, x_n]/(f_1, \ldots, f_c)$，其中
$n - c = \dim(S)$。由定义后的说明，$\dim(S_g) = n - c$。
令 $g' \in k[x_1, \ldots, x_n]$ 是剩余类对应于 $g$ 的一个元素。
则
$S_g =  k[x_1, \ldots, x_n, x_{n + 1}]/(f_1, \ldots, f_c, x_{n + 1}g' - 1)$，
正合所需。
\end{proof}

\begin{lemma}
\label{algebra-lemma-lci-CM}
设 $k$ 是域。设 $S$ 是有限型 $k$-代数。
若 $S$ 是局部完全交，则 $S$ 是 Cohen--Macaulay 环。
\end{lemma}

\begin{proof}
选取 $S$ 的一个极大素理想 $\mathfrak m$。
% P01-E0774: the frozen proof says “maximal prime” where the local criterion uses a maximal ideal.
我们须证明 $S_\mathfrak m$ 是 Cohen--Macaulay 的。
由假设，不妨设 $S = k[x_1, \ldots, x_n]/(f_1, \ldots, f_c)$，
其中 $\dim(S) = n - c$。令
$\mathfrak m' \subset k[x_1, \ldots, x_n]$ 为对应于
$\mathfrak m$ 的极大理想。由命题
\ref{algebra-proposition-finite-gl-dim-polynomial-ring}，局部环
$k[x_1, \ldots, x_n]_{\mathfrak m'}$ 是维数为 $n$ 的正则局部环。
特别地，由引理 \ref{algebra-lemma-regular-ring-CM}，它是 Cohen--Macaulay 的。
对引理 \ref{algebra-lemma-one-equation} 应用 $c$ 次可知，局部环
$S_{\mathfrak m} = k[x_1, \ldots, x_n]_{\mathfrak m'}/(f_1, \ldots, f_c)$
的维数 $\geq n - c$。由假设，$\dim(S_{\mathfrak m}) \leq n - c$。
故等号成立。这蕴含 $f_1, \ldots, f_c$ 在
$k[x_1, \ldots, x_n]_{\mathfrak m'}$ 中是正则序列，且
$S_{\mathfrak m}$ 是 Cohen--Macaulay 的；见命题
\ref{algebra-proposition-CM-module}。
\end{proof}

\noindent
下面的引理是本节其余内容的技术关键。它的一个重要特征是：
我们可以为环 $S$ 任取一个表现，而条件 (1) 不依赖于这一选择。

\begin{lemma}
\label{algebra-lemma-lci}
设 $k$ 是域。
设 $S$ 是有限型 $k$-代数。
设 $\mathfrak q$ 是 $S$ 的一个素理想。
任取一个表现 $S = k[x_1, \ldots, x_n]/I$。
令 $\mathfrak q'$ 为 $k[x_1, \ldots, x_n]$ 中对应于
$\mathfrak q$ 的素理想。置
$c = \text{高度}(\mathfrak q') - \text{高度}(\mathfrak q)$，
换言之，$\dim_{\mathfrak q}(S) = n - c$
（见引理 \ref{algebra-lemma-codimension}）。下列条件等价：
\begin{enumerate}
\item 存在 $g \in S$，$g \not \in \mathfrak q$，使得
$S_g$ 是 $k$ 上的全局完全交。
\item 理想 $I_{\mathfrak q'} \subset k[x_1, \ldots, x_n]_{\mathfrak q'}$
可由 $c$ 个元素生成。
\item 余法模 $(I/I^2)_{\mathfrak q}$ 作为
$S_{\mathfrak q}$-模可由 $c$ 个元素生成。
\item 余法模 $(I/I^2)_{\mathfrak q}$ 是秩为 $c$ 的自由
$S_{\mathfrak q}$-模。
\item 理想 $I_{\mathfrak q'}$ 在正则局部环
$k[x_1, \ldots, x_n]_{\mathfrak q'}$ 中可由一个正则序列生成。
\end{enumerate}
在此情形下，$I_{\mathfrak q'}$ 中任意 $c$ 个生成
$I_{\mathfrak q'}/\mathfrak q'I_{\mathfrak q'}$ 的元素，
都在局部环 $k[x_1, \ldots, x_n]_{\mathfrak q'}$ 中构成正则序列。
\end{lemma}

\begin{proof}
置 $R = k[x_1, \ldots, x_n]_{\mathfrak q'}$。这是维数为
$\text{高度}(\mathfrak q')$ 的 Cohen--Macaulay 局部环；
例如见引理 \ref{algebra-lemma-lci-CM}。此外，
$\overline{R} = R/IR = R/I_{\mathfrak q'} = S_{\mathfrak q}$
是维数为 $\text{高度}(\mathfrak q)$ 的商环。
取 $f_1, \ldots, f_c \in I_{\mathfrak q'}$ 为生成
$(I/I^2)_{\mathfrak q}$ 的元素。由引理 \ref{algebra-lemma-NAK} 可知，
$f_1, \ldots, f_c$ 生成 $I_{\mathfrak q'}$。
由维数恰好吻合，再由命题 \ref{algebra-proposition-CM-module} 可知
$f_1, \ldots, f_c$ 在 $R$ 中是正则序列。
由引理 \ref{algebra-lemma-regular-quasi-regular} 可知
$(I/I^2)_{\mathfrak q}$ 是自由的。
这些论证表明 (2)、(3)、(4) 等价，并且它们蕴含引理的最后一个断言，
从而也蕴含 (5)。

\medskip\noindent
若 (5) 成立，设 $I_{\mathfrak q'}$ 由长度为 $e$ 的一个正则序列生成，
则由维数理论
$$
\text{高度}(\mathfrak q) = \dim(S_{\mathfrak q}) =
\dim(k[x_1, \ldots, x_n]_{\mathfrak q'}) - e =
\text{高度}(\mathfrak q') - e
$$
（见第 \ref{algebra-section-dimension} 节）。因此 $e = c$，
故 (5) 蕴含 (2)。

\medskip\noindent
继续沿用第一段引入的记号。对每个 $f_i$，可找到
$d_i \in k[x_1, \ldots, x_n]$，$d_i \not \in \mathfrak q'$，
使得 $f_i' = d_i f_i \in k[x_1, \ldots, x_n]$。
仍有 $I_{\mathfrak q'} = (f_1', \ldots, f_c')R$。
因此存在 $g' \in k[x_1, \ldots, x_n]$，$g' \not \in \mathfrak q'$，
使得 $I_{g'} = (f_1', \ldots, f_c')$。
此外，选取 $g'' \in k[x_1, \ldots, x_n]$，
$g'' \not \in \mathfrak q'$，使得
$\dim(S_{g''}) = \dim_{\mathfrak q} \Spec(S)$。
由引理 \ref{algebra-lemma-codimension}，这一维数等于 $n - c$。
最后，令 $g$ 为 $g'g''$ 在 $S$ 中的像。于是
$$
S_g \cong k[x_1, \ldots, x_n, x_{n + 1}]
/
(f_1', \ldots, f_c', x_{n + 1}g'g'' - 1)
$$
并且由 $g''$ 的选取，此环的维数为 $n - c$。
故它是全局完全交。
因此 (2)、(3)、(4) 中的每一个都蕴含 (1)。

\medskip\noindent
假设 (1) 成立。令
$S_g \cong k[y_1, \ldots, y_m]/(f_1, \ldots, f_t)$
为 $S_g$ 作为全局完全交的一个表现。写
$J = (f_1, \ldots, f_t)$。令
$\mathfrak q'' \subset k[y_1, \ldots, y_m]$
为对应于 $\mathfrak qS_g$ 的素理想。注意
$t = m - \dim(S_g) =
\text{高度}(\mathfrak q'') - \text{高度}(\mathfrak q)$；
最后一个等号见引理 \ref{algebra-lemma-codimension}。
如引理 \ref{algebra-lemma-lci-CM} 的证明（以及上文）所示，
$f_1, \ldots, f_t$ 在局部环
$k[y_1, \ldots, y_m]_{\mathfrak q''}$ 中构成正则序列。
由引理 \ref{algebra-lemma-regular-quasi-regular} 可知
$(J/J^2)_{\mathfrak q}$ 是秩为 $t$ 的自由模。
由引理 \ref{algebra-lemma-conormal-module-localize}，有
$$
J/J^2 \oplus S_g^n \cong (I/I^2)_g \oplus S_g^m
$$
% P01-E0775: the frozen free-module summands use ordinary powers S_g^n and S_g^m instead of direct-sum notation.
因此 $(I/I^2)_{\mathfrak q}$ 是秩为
$t + n - m = m - \dim(S_g) + n - m = n - \dim(S_g) =
\text{高度}(\mathfrak q') - \text{高度}(\mathfrak q) = c$
的自由模。故得到 (4)。
\end{proof}

\noindent
引理 \ref{algebra-lemma-lci} 的结果提示我们作如下定义。

\begin{definition}
\label{algebra-definition-lci-local-ring}
设 $k$ 是域。设 $S$ 是在 $k$ 上本质有限型的局部 $k$-代数。
若存在一个局部 $k$-代数 $R$ 以及元素
$f_1, \ldots, f_c \in \mathfrak m_R$，使得
\begin{enumerate}
\item $R$ 在 $k$ 上本质有限型；
\item $R$ 是正则局部环；
\item $f_1, \ldots, f_c$ 在 $R$ 中构成正则序列；并且
\item 作为 $k$-代数，$S \cong R/(f_1, \ldots, f_c)$，
\end{enumerate}
则称 $S$ 是一个\textit{（$k$ 上的）完全交}。
\end{definition}

\noindent
由 Cohen 结构定理（见定理 \ref{algebra-theorem-cohen-structure-theorem}），
任意完备 Noether 局部环都可以写成某个正则完备局部环的商。
因此可以利用上述定义，为任意完备 Noether 局部环定义完全交环的概念。
我们将在《除幂代数》的第 \ref{dpa-section-lci} 节讨论这一点。
与此同时，下面的引理表明这样的定义是有意义的。

\begin{lemma}
\label{algebra-lemma-ci-well-defined}
设 $A \to B \to C$ 是满的局部环同态。
假设 $A$ 和 $B$ 是正则局部环。下列条件等价：
% P01-E0776: the frozen sentence introducing the equivalence list lacks a colon.
\begin{enumerate}
\item $\Ker(A \to C)$ 由一个正则序列生成；
\item $\Ker(A \to C)$ 可由 $\dim(A) - \dim(C)$ 个元素生成；
\item $\Ker(B \to C)$ 由一个正则序列生成；并且
\item $\Ker(B \to C)$ 可由 $\dim(B) - \dim(C)$ 个元素生成。
\end{enumerate}
\end{lemma}

\begin{proof}
正则局部环是 Cohen--Macaulay 的；见引理
\ref{algebra-lemma-regular-ring-CM}。因此 (1) $\Leftrightarrow$ (2) 以及
(3) $\Leftrightarrow$ (4)；见命题 \ref{algebra-proposition-CM-module}。
% P01-E0777: the frozen “Hence the equivalences ...” clause has no finite verb.
由引理 \ref{algebra-lemma-regular-quotient-regular}，理想
$\Ker(A \to B)$ 可由 $\dim(A) - \dim(B)$ 个元素生成。
故 (4) 蕴含 (2)。

\medskip\noindent
还须证明 (1) 蕴含 (4)。我们对 $\dim(A) - \dim(B)$ 作归纳。
$\dim(A) - \dim(B) = 0$ 的情形是平凡的。
假设 $\dim(A) > \dim (B)$。
写 $I = \Ker(A \to C)$ 且 $J = \Ker(A \to B)$。
注意 $J \subset I$。我们的假设是 $I$ 的最小生成元数为
$\dim(A) - \dim(C)$。令 $\mathfrak m \subset A$ 为极大理想。
考虑映射
$$
J/ \mathfrak m J \to  I / \mathfrak m I \to \mathfrak m /\mathfrak m^2
$$
% P01-E0778: the frozen display lacks terminal punctuation before a new sentence.
由引理 \ref{algebra-lemma-regular-quotient-regular} 及其证明，上述复合是单射。
任取 $x \in J$，使其在 $J /\mathfrak mJ$ 中非零。
由上述结论和 Nakayama 引理，$x$ 是 $I$ 的某个最小生成集中的元素。
因此可以把 $A$ 替换为 $A/xA$，把 $I$ 替换为 $I/xA$；
这使 $\dim(A)$ 以及 $I$ 的最小生成元数都减少 $1$。故得证。
\end{proof}

\begin{lemma}
\label{algebra-lemma-lci-local}
设 $k$ 是域。设 $S$ 是在 $k$ 上本质有限型的局部 $k$-代数。
下列条件等价：
\begin{enumerate}
\item $S$ 是 $k$ 上的完全交；
\item 对任意满射 $R \to S$，若 $R$ 是在 $k$ 上本质有限表示的
正则局部环，则理想 $\Ker(R \to S)$ 可由一个正则序列生成；
\item 对某个满射 $R \to S$，若 $R$ 是在 $k$ 上本质有限表示的
正则局部环，则理想 $\Ker(R \to S)$ 可由
$\dim(R) - \dim(S)$ 个元素生成；
\item 存在 $k$ 上的全局完全交 $A$ 以及 $A$ 的一个素理想
$\mathfrak a$，使得 $S \cong A_{\mathfrak a}$；并且
\item 存在 $k$ 上的局部完全交 $A$ 以及 $A$ 的一个素理想
$\mathfrak a$，使得 $S \cong A_{\mathfrak a}$。
\end{enumerate}
\end{lemma}

\begin{proof}
显然 (2) 蕴含 (1)，而 (1) 蕴含 (3)。同样显然 (4) 蕴含 (5)。
我们来证明 (3) 蕴含 (4)。因此假设存在满射
$R \to S$，其中 $R$ 是在 $k$ 上本质有限表示的正则局部环，
并且理想 $\Ker(R \to S)$ 可由
$\dim(R) - \dim(S)$ 个元素生成。
可写成 $R = (k[x_1, \ldots, x_n]/J)_{\mathfrak q}$，
其中 $J \subset k[x_1, \ldots, x_n]$，而
$\mathfrak q \subset k[x_1, \ldots, x_n]$ 是满足
$J \subset \mathfrak q$ 的某个素理想。令
$I \subset k[x_1, \ldots, x_n]$ 为映射
$k[x_1, \ldots, x_n] \to S$ 的核，从而
$S \cong (k[x_1, \ldots, x_n]/I)_{\mathfrak q}$。
由假设，$(I/J)_{\mathfrak q}$ 可由
$\dim(R) - \dim(S)$ 个元素生成。引理
\ref{algebra-lemma-ci-well-defined} 蕴含 $I_{\mathfrak q}$ 可由
$\dim(k[x_1, \ldots, x_n]_{\mathfrak q}) - \dim(S)$
个元素生成。由引理 \ref{algebra-lemma-lci} 可知，存在
$g \in k[x_1, \ldots, x_n]$，$g \not \in \mathfrak q$，使得代数
$(k[x_1, \ldots, x_n]/I)_g$ 是全局完全交，而 $S$
同构于它的一个局部环。

\medskip\noindent
为了完成引理的证明，还须证明 (5) 蕴含 (2)。
假设 (5)，并令 $\pi : R \to S$ 为满射，其中 $R$
是在 $k$ 上本质有限型的正则局部 $k$-代数。
由假设，对 $k$ 上的某个局部完全交 $A$，有
$S = A_{\mathfrak a}$。选取表现
$R = (k[y_1, \ldots, y_m]/J)_{\mathfrak q}$，其中
$J \subset \mathfrak q \subset k[y_1, \ldots, y_m]$。
可以且确实假设 $J$ 是映射 $k[y_1, \ldots, y_m] \to R$ 的核。
令 $I \subset k[y_1, \ldots, y_m]$ 为映射
$k[y_1, \ldots, y_m] \to S = A_{\mathfrak a}$ 的核。
则 $J \subset I$，而 $(I/J)_{\mathfrak q}$ 是满射
$\pi : R \to S$ 的核。因此
$S = (k[y_1, \ldots, y_m]/I)_{\mathfrak q}$。

\medskip\noindent
由引理 \ref{algebra-lemma-isomorphic-local-rings} 可知，存在
$g \in A$，$g \not \in \mathfrak a$，以及
$g' \in k[y_1, \ldots, y_m]$，$g' \not \in \mathfrak q$，
使得 $A_g \cong (k[y_1, \ldots, y_m]/I)_{g'}$。
把 $A$ 替换为 $A_g$，把 $k[y_1, \ldots, y_m]$ 替换为
$k[y_1, \ldots, y_{m + 1}]$ 后，可以假设
$A \cong k[y_1, \ldots, y_m]/I$。考虑局部环的满映射
$$
k[y_1, \ldots, y_m]_{\mathfrak q} \to R \to S.
$$
我们须证明 $R \to S$ 的核由一个正则序列生成。
由引理 \ref{algebra-lemma-lci} 可知，
$k[y_1, \ldots, y_m]_{\mathfrak q} \to A_{\mathfrak a} = S$
具有这一性质（因为 $A$ 是 $k$ 上的局部完全交）。
再由引理 \ref{algebra-lemma-ci-well-defined} 即得证。
\end{proof}

\begin{lemma}
\label{algebra-lemma-lci-at-prime}
设 $k$ 是域。设 $S$ 是有限型 $k$-代数。
设 $\mathfrak q$ 是 $S$ 的一个素理想。下列条件等价：
\begin{enumerate}
\item 局部环 $S_{\mathfrak q}$ 是完全交环
（定义 \ref{algebra-definition-lci-local-ring}）。
\item 存在 $g \in S$，$g \not \in \mathfrak q$，使得
$S_g$ 是 $k$ 上的局部完全交。
\item 存在 $g \in S$，$g \not \in \mathfrak q$，使得
$S_g$ 是 $k$ 上的全局完全交。
\item 对任意表现 $S = k[x_1, \ldots, x_n]/I$，若
$\mathfrak q' \subset k[x_1, \ldots, x_n]$ 对应于 $\mathfrak q$，
则引理 \ref{algebra-lemma-lci} 的等价条件 (1) -- (5) 中的任一个成立。
% P01-E0779: the frozen plural phrase “any of the equivalent conditions ... hold” has number disagreement and is logically clearer as “any one ... holds”.
\end{enumerate}
\end{lemma}

\begin{proof}
这是引理 \ref{algebra-lemma-lci}、\ref{algebra-lemma-lci-local} 与各定义的结合。
\end{proof}

\begin{lemma}
\label{algebra-lemma-lci-global}
设 $k$ 是域。设 $S$ 是有限型 $k$-代数。
下列条件等价：
\begin{enumerate}
\item 环 $S$ 是 $k$ 上的局部完全交。
\item $S$ 的所有局部环都是 $k$ 上的完全交环。
\item $S$ 在极大理想处的所有局部化都是 $k$ 上的完全交环。
\end{enumerate}
\end{lemma}

\begin{proof}
这由引理 \ref{algebra-lemma-lci-at-prime}、$\Spec(S)$ 的拟紧性以及各定义得出。
\end{proof}

\noindent
下面的引理说明，完全交这一性质在基域变换下保持（而且是在强意义下）。

\begin{lemma}
\label{algebra-lemma-lci-field-change-local}
设 $K/k$ 是域扩张。
设 $S$ 是 $k$ 上的有限型代数。
令 $\mathfrak q_K$ 为 $S_K = K \otimes_k S$ 的一个素理想，
并令 $\mathfrak q$ 为 $S$ 中对应的素理想。
则 $S_{\mathfrak q}$ 是 $k$ 上的完全交
（定义 \ref{algebra-definition-lci-local-ring}），当且仅当
$(S_K)_{\mathfrak q_K}$ 是 $K$ 上的完全交。
\end{lemma}

\begin{proof}
选取表现 $S = k[x_1, \ldots, x_n]/I$。
这给出表现 $S_K = K[x_1, \ldots, x_n]/I_K$，其中
$I_K = K \otimes_k I$。令
$\mathfrak q_K' \subset K[x_1, \ldots, x_n]$ 与
$\mathfrak q' \subset k[x_1, \ldots, x_n]$ 分别为对应的素理想。
我们将证明：引理 \ref{algebra-lemma-lci} 的等价条件对二元组
$(S = k[x_1, \ldots, x_n]/I, \mathfrak q)$ 成立，当且仅当它们对二元组
$(S_K = K[x_1, \ldots, x_n]/I_K, \mathfrak q_K)$ 成立。
由此即可得出引理（见引理 \ref{algebra-lemma-lci-at-prime}）。

\medskip\noindent
由引理 \ref{algebra-lemma-dimension-at-a-point-preserved-field-extension}，有
$\dim_{\mathfrak q} S = \dim_{\mathfrak q_K} S_K$。
因此，引理 \ref{algebra-lemma-lci} 中出现的整数 $c$ 对二元组
$(S = k[x_1, \ldots, x_n]/I, \mathfrak q)$ 与二元组
$(S_K = K[x_1, \ldots, x_n]/I_K, \mathfrak q_K)$ 相同。
另一方面，
\begin{eqnarray*}
I \otimes_{k[x_1, \ldots, x_n]} \kappa(\mathfrak q')
\otimes_{\kappa(\mathfrak q')} \kappa(\mathfrak q_K')
& = &
I \otimes_{k[x_1, \ldots, x_n]} \kappa(\mathfrak q_K') \\
& = &
I \otimes_{k[x_1, \ldots, x_n]} K[x_1, \ldots, x_n]
\otimes_{K[x_1, \ldots, x_n]} \kappa(\mathfrak q_K') \\
& = &
(K \otimes_k I) \otimes_{K[x_1, \ldots, x_n]} \kappa(\mathfrak q_K') \\
& = &
I_K \otimes_{K[x_1, \ldots, x_n]} \kappa(\mathfrak q'_K).
\end{eqnarray*}
因此，
$\dim_{\kappa(\mathfrak q')}
I \otimes_{k[x_1, \ldots, x_n]} \kappa(\mathfrak q')
=
\dim_{\kappa(\mathfrak q'_K)}
I_K \otimes_{K[x_1, \ldots, x_n]} \kappa(\mathfrak q_K')$。
% P01-E0780: the frozen display alternates between q_K' and q'_K for the same extended prime.
于是由 Nakayama 引理 \ref{algebra-lemma-NAK}，
$I_{\mathfrak q'}$ 的最小生成元数等于
$(I_K)_{\mathfrak q'_K}$ 的最小生成元数。
因此，引理由引理 \ref{algebra-lemma-lci} 的刻画 (2) 得出。
\end{proof}

\begin{lemma}
\label{algebra-lemma-lci-field-change}
设 $k \to K$ 是域扩张。
设 $S$ 是有限型 $k$-代数。
则 $S$ 是 $k$ 上的局部完全交，当且仅当
$S \otimes_k K$ 是 $K$ 上的局部完全交。
\end{lemma}

\begin{proof}
这由引理 \ref{algebra-lemma-lci-global} 与
\ref{algebra-lemma-lci-field-change-local} 结合得出。
不过，我们也给出另一个证明（它基于同样的原理）。

\medskip\noindent
置 $S' = S \otimes_k K$。令 $\alpha : k[x_1, \ldots, x_n] \to S$
是核为 $I$ 的一个表现。令
$\alpha' : K[x_1, \ldots, x_n] \to S'$ 是所诱导的表现，其核为 $I'$。

\medskip\noindent
假设 $S$ 是局部完全交。取素理想 $\mathfrak q \subset S'$。
% P01-E0781: the frozen designation sentence omits “by/as” after “Denote”.
分别用 $\mathfrak q'$、$\mathfrak p$ 和 $\mathfrak p'$ 表示
$K[x_1, \ldots, x_n]$ 中对应的素理想、$S$ 中对应的素理想以及
$k[x_1, \ldots, x_n]$ 中对应的素理想。考虑如下 Noether 局部环的图
$$
\xymatrix{
S'_{\mathfrak q} &  K[x_1, \ldots, x_n]_{\mathfrak q'} \ar[l] \\
S_{\mathfrak p}\ar[u] &  k[x_1, \ldots, x_n]_{\mathfrak p'} \ar[u] \ar[l]
}
$$
% P01-E0782: the frozen diagram lacks terminal punctuation before the next sentence.
由引理 \ref{algebra-lemma-lci}，我们知道 $S_{\mathfrak p}$ 在
$k[x_1, \ldots, x_n]_{\mathfrak p'}$ 中由某个正则序列
$f_1, \ldots, f_c$ 截出。由于右边的竖直箭头是平坦的，
$f_1, \ldots, f_c$ 的像在
$K[x_1, \ldots, x_n]_{\mathfrak q'}$ 中构成正则序列。
由于在 $k$ 上与 $K$ 作张量是正合函子，我们有
$S'_{\mathfrak q} = K[x_1, \ldots, x_n]_{\mathfrak q'}/(f_1, \ldots, f_c)$。
故再次由引理 \ref{algebra-lemma-lci} 可知，$S'$ 在 $\mathfrak q$ 的一个邻域内
是局部完全交。由于 $\mathfrak q$ 是任意的，$S'$ 是 $K$ 上的局部完全交。

\medskip\noindent
反过来，假设 $S'$ 是局部完全交。取 $S$ 的一个极大理想
$\mathfrak m$。令 $\mathfrak m'$ 表示 $k[x_1, \ldots, x_n]$
中对应的极大理想。
% P01-E0783: the frozen residue-field designation omits “by/as”.
记 $\kappa = \kappa(\mathfrak m)$ 为剩余域。
由备注 \ref{algebra-remark-fundamental-diagram}，$S'$ 中位于
$\mathfrak m$ 上方的素理想对应于 $K \otimes_k \kappa$ 中的素理想。
由 Hilbert 零点定理 \ref{algebra-theorem-nullstellensatz}，有
$[\kappa : k] < \infty$。因此 $K \otimes_k \kappa$
是 $K$ 上非零的有限代数。于是 $K \otimes_k \kappa$ 有有限个
（其个数 $> 0$）素理想；它们全是极大理想，并且各自的剩余域都是 $K$ 的有限扩张
（见第 \ref{algebra-section-artinian} 节）。
% P01-E0784: the frozen prose twice says “a finite number > 0” of primes.
因此，位于 $\mathfrak m$ 上方的素理想
$\mathfrak n \subset S'$ 有有限个（其个数 $> 0$）；它们全是极大理想，并且各自的剩余域
都是 $K$ 的有限扩张。取其中一个，记为 $\mathfrak n \subset S'$，
并令 $\mathfrak n' \subset K[x_1, \ldots, x_n]$ 表示
$K[x_1, \ldots, x_n]$ 中对应的素理想。
注意，由于 $V(\mathfrak mS')$ 是有限的，$\mathfrak n$
是其中的一个孤立闭点，从而 $\mathfrak mS'_{\mathfrak n}$
是 $S'_{\mathfrak n}$ 的一个定义理想。例如，由引理
\ref{algebra-lemma-dimension-base-fibre-equals-total}，这蕴含
$\dim(S_{\mathfrak m}) = \dim(S'_{\mathfrak n})$。
（也可以利用引理
\ref{algebra-lemma-dimension-at-a-point-preserved-field-extension} 看出这一点。）
考虑相应的 Noether 局部环之图
$$
\xymatrix{
S'_{\mathfrak n} &  K[x_1, \ldots, x_n]_{\mathfrak n'} \ar[l] \\
S_{\mathfrak m}\ar[u] &  k[x_1, \ldots, x_n]_{\mathfrak m'} \ar[u] \ar[l]
}
$$
% P01-E0785: the frozen diagram lacks terminal punctuation before the next sentence.
由引理 \ref{algebra-lemma-change-base-NL}，有
$\NL(\alpha) \otimes_S S' = \NL(\alpha')$，特别地，
$I'/(I')^2 = I/I^2 \otimes_S S'$。因此
$(I/I^2)_{\mathfrak m} \otimes_{S_{\mathfrak m}} \kappa$
与
$(I'/(I')^2)_{\mathfrak n} \otimes_{S'_{\mathfrak n}} \kappa(\mathfrak n)$
具有相同的维数。由于 $(I'/(I')^2)_{\mathfrak n}$
是秩为 $n - \dim S'_{\mathfrak n}$ 的自由模，我们推出
$(I/I^2)_{\mathfrak m}$ 可由
$n - \dim S'_{\mathfrak n} = n - \dim S_{\mathfrak m}$ 个元素生成。
由引理 \ref{algebra-lemma-lci} 可知，$S$ 在 $\mathfrak m$ 的一个邻域内
是局部完全交。由于 $\mathfrak m$ 是任意极大理想，我们得出
$S$ 是局部完全交。
\end{proof}

\noindent
最后给出一个引理。稍后我们将用它证明：给定环映射
$T \to A \to B$，若 $B$ 在 $T$ 上是 syntomic 的，且
$B$ 在 $A$ 上是 syntomic 的，则 $A$ 在 $T$ 上是 syntomic 的。

\begin{lemma}
\label{algebra-lemma-lci-permanence-initial}
设
$$
\xymatrix{
B & S \ar[l] \\
A \ar[u] & R \ar[l] \ar[u]
}
$$
% P01-E0786: the frozen diagram lacks terminal punctuation before “be a commutative square”.
是局部环的交换方块。假设
\begin{enumerate}
\item $R$ 和 $\overline{S} = S/\mathfrak m_R S$ 是正则局部环，
\item 对某些理想 $I$、$J$，有 $A = R/I$ 和 $B = S/J$，
\item $J \subset S$ 以及
% P01-E0787: the frozen quotient omits parentheses and likely the extension of m_R to S.
$\overline{J} = J/\mathfrak m_R \cap J \subset \overline{S}$
由正则序列生成，并且
\item $A \to B$ 和 $R \to S$ 是平坦的。
\end{enumerate}
则 $I$ 由一个正则序列生成。
\end{lemma}

\begin{proof}
置 $\overline{B} = B/\mathfrak m_RB = B/\mathfrak m_AB$，于是
$\overline{B} = \overline{S}/\overline{J}$。
取元素 $f_1, \ldots, f_{\overline{c}} \in J$，使得
$\overline{f}_1, \ldots, \overline{f}_{\overline{c}} \in \overline{J}$
构成生成 $\overline{J}$ 的正则序列。注意
$\overline{c} = \dim(\overline{S}) - \dim(\overline{B})$；
见引理 \ref{algebra-lemma-ci-well-defined}。
由引理 \ref{algebra-lemma-grothendieck-regular-sequence}，环
$S/(f_1, \ldots, f_{\overline{c}})$ 在 $R$ 上平坦。
% P01-E0788: the frozen quotient by the sum lacks parentheses in both occurrences.
因此 $S/(f_1, \ldots, f_{\overline{c}}) + IS$ 在 $A$ 上平坦。
映射 $S/(f_1, \ldots, f_{\overline{c}}) + IS \to B$
于是是有限 $S/IS$-模之间的满射，这些模在 $A$ 上平坦；它模去
$\mathfrak m_A$ 后是同构，故由引理 \ref{algebra-lemma-mod-injective}
可知它是同构。换言之，
$J = (f_1, \ldots, f_{\overline{c}}) + IS$。

\medskip\noindent
再次由引理 \ref{algebra-lemma-ci-well-defined}，理想 $J$ 由一个长度为
$c = \dim(S) - \dim(B)$ 的正则序列生成。因此
$J/\mathfrak m_SJ$ 是维数为 $c$ 的向量空间。
由上述对 $J$ 的描述，存在
$g_1, \ldots, g_{c - \overline{c}} \in I$，使得 $J$ 由
$f_1, \ldots, f_{\overline{c}}, g_1, \ldots, g_{c - \overline{c}}$
生成（使用 Nakayama 引理 \ref{algebra-lemma-NAK}）。考虑环
$A' = R/(g_1, \ldots, g_{c - \overline{c}})$ 以及满射
$A' \to A$。由上述结论可知，
$B = S/(f_1, \ldots, f_{\overline{c}}, g_1, \ldots, g_{c - \overline{c}})$
在 $A'$ 上平坦（因为 $S/(f_1, \ldots, f_{\overline{c}})$
在 $R$ 上平坦）。因此 $A' \to B$ 是单射（因为它忠实平坦；
见引理 \ref{algebra-lemma-local-flat-ff}）。
由于该映射经由 $A$ 分解，我们得到 $A' = A$。
注意 $\dim(B) = \dim(A) + \dim(\overline{B})$，且
$\dim(S) = \dim(R) + \dim(\overline{S})$；见引理
\ref{algebra-lemma-dimension-base-fibre-equals-total}。
% P01-E0789: the frozen dimension difference omits spacing before the second dim command.
由初等代数可知 $c - \overline{c} = \dim(R) -\dim(A)$。
故由引理 \ref{algebra-lemma-ci-well-defined}，
$I = (g_1, \ldots, g_{c - \overline{c}})$ 由一个正则序列生成。
\end{proof}

\section{syntomic 态射}
\label{algebra-section-syntomic}

\noindent
syntomic 环映射是平坦、有限表示，并且所有纤维都是局部完全交的环映射。
一般的局部完全交环映射将在《代数进阶》的第
\ref{more-algebra-section-lci} 节中讨论。

\begin{definition}
\label{algebra-definition-lci}
若环映射 $R \to S$ 是平坦且有限表示的，并且它的所有纤维环
$S \otimes_R \kappa(\mathfrak p)$ 都是局部完全交（见定义
\ref{algebra-definition-lci-field}），则称它是 {\it syntomic} 的，
也称 $S$ 是 $R$ 上的一个{\it 平坦局部完全交}。
\end{definition}

\noindent
显然，域上的代数在该域上是 syntomic 的，当且仅当它是局部完全交。
这个定义具有如下令人满意的性质。

\begin{lemma}
\label{algebra-lemma-syntomic-descends}
\begin{slogan}
syntomic 性质在基底上是 fpqc 局部的。
\end{slogan}
设 $R \to S$ 是环映射。
设 $R \to R'$ 是忠实平坦环映射。
置 $S' = R'\otimes_R S$。
则 $R \to S$ 是 syntomic 的，当且仅当 $R' \to S'$ 是 syntomic 的。
\end{lemma}

\begin{proof}
由引理 \ref{algebra-lemma-finite-presentation-descends} 和
\ref{algebra-lemma-flatness-descends}，结论对平坦性和有限表示性成立。
映射 $\Spec(R') \to \Spec(R)$ 是满射；见引理
\ref{algebra-lemma-ff-rings}。因此，只须证明：给定位于
$\mathfrak p \subset R$ 上方的素理想 $\mathfrak p' \subset R'$，
% P01-E0790: the frozen “it suffices to show given primes ... that” construction is grammatically malformed.
$S \otimes_R \kappa(\mathfrak p)$ 是局部完全交，当且仅当
$S' \otimes_{R'} \kappa(\mathfrak p')$ 是局部完全交。注意
$S' \otimes_{R'} \kappa(\mathfrak p') =
S \otimes_R \kappa(\mathfrak p)
\otimes_{\kappa(\mathfrak p)} \kappa(\mathfrak p')$。
于是可应用引理 \ref{algebra-lemma-lci-field-change}。
\end{proof}

\begin{lemma}
\label{algebra-lemma-base-change-syntomic}
syntomic 映射的任意基变换仍是 syntomic 映射。
\end{lemma}

\begin{proof}
平坦性、有限表示性以及纤维为局部完全交这三个性质分别由引理
\ref{algebra-lemma-flat-base-change}、\ref{algebra-lemma-compose-finite-type} 和
\ref{algebra-lemma-lci-field-change} 得到。
\end{proof}

\begin{lemma}
\label{algebra-lemma-local-syntomic}
设 $R \to S$ 是环映射。
假设有生成单位理想的 $g_1, \ldots g_m \in S$，并且每个
% P01-E0791: the frozen mathematical list omits the comma after \ldots.
$R \to S_{g_i}$ 都是 syntomic 的。
则 $R \to S$ 是 syntomic 的。
\end{lemma}

\begin{proof}
由引理 \ref{algebra-lemma-flat-localization} 和 \ref{algebra-lemma-cover-upstairs}，
这对平坦性和有限表示性成立。纤维环为局部完全交这一性质按照其定义
在 $S$ 上是局部的；见定义 \ref{algebra-definition-lci-field}。
\end{proof}

\begin{definition}
\label{algebra-definition-relative-global-complete-intersection}
设 $R \to S$ 是环映射。若存在表现
$S = R[x_1, \ldots, x_n]/(f_1, \ldots, f_c)$，并且
$\Spec(S) \to \Spec(R)$ 的每个非空纤维都具有维数 $n - c$，
则称 $R \to S$ 是一个{\it 相对全局完全交}。
我们将用“令
$S = R[x_1, \ldots, x_n]/(f_1, \ldots, f_c)$
为一个相对全局完全交”来表示这种情形。
\end{definition}

\noindent
下面的引理有时可用于寻找全局表现。

\begin{lemma}
\label{algebra-lemma-huber}
设 $S$ 是有限表示 $R$-代数，并且它有一个表现
$S = R[x_1, \ldots, x_n]/I$，使得 $I/I^2$ 是自由 $S$-模。
则 $S$ 有一个表现
$S = R[y_1, \ldots, y_m]/(f_1, \ldots, f_c)$，使得
$(f_1, \ldots, f_c)/(f_1, \ldots, f_c)^2$ 是自由模，且
$f_1, \ldots, f_c$ 的类给出其一组基。
\end{lemma}

\begin{proof}
由引理 \ref{algebra-lemma-finite-presentation-independent}，$I$ 是有限生成理想。
取映到 $I/I^2$ 的一组基的元素 $f_1, \ldots, f_c \in I$。
由 Nakayama 引理（引理 \ref{algebra-lemma-NAK}），存在
$g \in 1 + I$，使得
$$
g \cdot I \subset (f_1, \ldots, f_c)
$$
并且 $I_g \cong (f_1, \ldots, f_c)_g$。因此，如所要求的，
$$
S \cong R[x_1, \ldots, x_n]/(f_1, \ldots, f_c)[1/g]
\cong R[x_1, \ldots, x_n, x_{n + 1}]/(f_1, \ldots, f_c, gx_{n + 1} - 1)
$$
进而可知 $f_1, \ldots, f_c,gx_{n + 1} - 1$ 构成
% P01-E0792: the frozen source omits a space after the comma before gx_{n+1}-1.
$(f_1, \ldots, f_c, gx_{n + 1} - 1)/(f_1, \ldots, f_c, gx_{n + 1} - 1)^2$
的一组基。例如，这可由应用引理
\ref{algebra-lemma-principal-localization-NL} 得出。
\end{proof}



\begin{example}
\label{algebra-example-factor-polynomials}
设 $n , m \geq 1$ 是整数。考虑环映射
\begin{eqnarray*}
R = \mathbf{Z}[a_1, \ldots, a_{n + m}]
& \longrightarrow &
S = \mathbf{Z}[b_1, \ldots, b_n, c_1, \ldots, c_m] \\
a_1 & \longmapsto & b_1 + c_1 \\
a_2 & \longmapsto & b_2 + b_1 c_1 + c_2 \\
\ldots & \ldots & \ldots \\
a_{n + m} & \longmapsto & b_n c_m
\end{eqnarray*}
换言之，这是如上所示的多项式环之间唯一的环映射，使得多项式分解
$$
x^{n + m} + a_1 x^{n + m - 1} + \ldots + a_{n + m}
=
(x^n + b_1 x^{n - 1} + \ldots + b_n)
(x^m + c_1 x^{m - 1} + \ldots + c_m)
$$
成立。注意 $S$ 在 $R$ 上由 $n + m$ 个元素（即 $b_i, c_j$）
生成，并且有 $n + m$ 个方程（即 $a_k = a_k(b_i, c_j)$）。
为了证明 $S$ 是 $R$ 上的相对全局完全交，只须证明所有纤维的维数为 $0$。

\medskip\noindent
为此，设 $R \to k$ 是到域 $k$ 的环映射。
设 $a_i$ 映到 $\alpha_i \in k$。
考虑纤维环 $S_k = k \otimes_R S$。设 $k \to K$ 是域扩张。
一个 $k$-代数映射 $S_k \to K$ 等价于找到
% P01-E0793: the frozen prose says “a map of S_k to K” instead of “from S_k to K”.
$\beta_1, \ldots, \beta_n, \gamma_1, \ldots, \gamma_m \in K$，使得
$$
x^{n + m} + \alpha_1 x^{n + m - 1} + \ldots + \alpha_{n + m}
=
(x^n + \beta_1 x^{n - 1} + \ldots + \beta_n)
(x^m + \gamma_1 x^{m - 1} + \ldots + \gamma_m).
$$
因此，在 $K$ 中这样的 $n + m$-元组只有有限种选择。
% P01-E0794: the frozen “at most finitely many” is redundant.
这证明所有纤维都只有有限个闭点（例如，使用 Hilbert 零点定理可知，
它们都对应于 $\overline{k}$ 中的解），因而 $R \to S$
是相对全局完全交。

\medskip\noindent
另一种论证方式是证明
$\mathbf{Z}[a_1, \ldots, a_{n + m}] \to
\mathbf{Z}[b_1, \ldots, b_n, c_1, \ldots, c_m]$
实际上也是一个{\it 有限}环映射。事实上，由引理
\ref{algebra-lemma-polynomials-divide}，每个 $b_i, c_j$ 在 $R$ 上都是整的，
故由引理 \ref{algebra-lemma-characterize-integral}，$R \to S$ 是有限的。
\end{example}

\begin{example}
\label{algebra-example-roots-universal-polynomial}
考虑环映射
\begin{eqnarray*}
R = \mathbf{Z}[a_1, \ldots, a_n]
& \longrightarrow &
S = \mathbf{Z}[\alpha_1, \ldots, \alpha_n] \\
a_1 & \longmapsto &
\alpha_1 + \ldots + \alpha_n \\
\ldots & \ldots & \ldots \\
a_n & \longmapsto & \alpha_1 \ldots \alpha_n
\end{eqnarray*}
换言之，这是如上所示的多项式环之间唯一的环映射，使得
$$
x^n + a_1 x^{n - 1} + \ldots + a_n
=
\prod\nolimits_{i = 1}^n (x + \alpha_i)
$$
在 $\mathbf{Z}[\alpha_i, x]$ 中成立。也就是说，$a_i$ 映到
$\alpha_1, \ldots, \alpha_n$ 的第 $i$ 个初等对称函数。
% P01-E0795: the frozen English typesets “i” and “th” without a hyphen.
由初等对称多项式的通常理论（略去细节），$S$ 是有限自由 $R$-模，
它的一组基由满足 $0 \leq e_i \leq n - i$ 的元素
$\alpha_1^{e_1} \alpha_2^{e_2} \ldots \alpha_n^{e_n}$ 给出。
尤其，$R \to S$ 的纤维环是有限的，因而具有维数 $0$。
另一方面，$S$ 在 $R$ 上由 $n$ 个元素和 $n$ 个方程生成。
因此 $S$ 是 $R$ 上的相对全局完全交。
由于 $S$ 在 $R$ 上的秩为正，还可看出 $S$ 在 $R$ 上忠实平坦。
\end{example}

\begin{lemma}
\label{algebra-lemma-base-change-relative-global-complete-intersection}
设 $S = R[x_1, \ldots, x_n]/(f_1, \ldots, f_c)$
是一个相对全局完全交（定义
\ref{algebra-definition-relative-global-complete-intersection}）。
% P01-E0796: the frozen statement lacks punctuation before the enumeration.
\begin{enumerate}
\item 对任意 $R \to R'$，基变换
$R' \otimes_R S = R'[x_1, \ldots, x_n]/(f_1, \ldots, f_c)$
是相对全局完全交。
\item 对任意 $g \in S$，若它是 $h \in R[x_1, \ldots, x_n]$ 的像，
则环
$$
S_g = R[x_1, \ldots, x_n, x_{n + 1}]/(f_1, \ldots, f_c, hx_{n + 1} - 1)
$$
是相对全局完全交。
\item 若对某个 $f \in R$，$R \to S$ 分解为 $R \to R_f \to S$，
% P01-E0797: the frozen source incorrectly separates the if-clause and “Then” with a period.
则环 $S = R_f[x_1, \ldots, x_n]/(f_1, \ldots, f_c)$
是 $R_f$ 上的相对全局完全交。
\end{enumerate}
\end{lemma}

\begin{proof}
由引理 \ref{algebra-lemma-dimension-preserved-field-extension}，基变换的纤维
与原映射的纤维具有相同的维数。此外，下式成立：
$$
R' \otimes_R R[x_1, \ldots, x_n]/(f_1, \ldots, f_c)
= R'[x_1, \ldots, x_n]/(f_1, \ldots, f_c)
$$
故 (1) 成立。
(2) 的证明在于，关于一个元素的局部化可以写成
$S_g \cong S[x_{n + 1}]/(gx_{n + 1} - 1)$。
断言 (3) 由 (1) 得出，因为在 (3) 的假设下有
$R_f \otimes_R S \cong S$。
\end{proof}

\begin{lemma}
\label{algebra-lemma-localize-relative-complete-intersection}
设 $R$ 是环。设 $S = R[x_1, \ldots, x_n]/(f_1, \ldots, f_c)$。
在以下每种情形中，我们都将找到
$h \in R[x_1, \ldots, x_n]$，其像为 $g \in S$，使得
$$
S_g = R[x_1, \ldots, x_n, x_{n + 1}]/(f_1, \ldots, f_c, hx_{n + 1} - 1)
$$
是相对全局完全交，并且具有定义
\ref{algebra-definition-relative-global-complete-intersection} 中的表现：
\begin{enumerate}
\item 设 $I \subset R$ 是理想。若
$\Spec(S/IS) \to \Spec(R/I)$ 的纤维具有维数 $n - c$，
则可以找到如上的 $(h, g)$，使得 $g$ 的像为 $1 \in S/IS$。
\item 设 $\mathfrak p \subset R$ 是素理想。若
$\dim(S \otimes_R \kappa(\mathfrak p)) = n - c$，
则可以找到如上的 $(h, g)$，使得 $g$ 在
$S \otimes_R \kappa(\mathfrak p)$ 中的像为单位。
\item 设 $\mathfrak q \subset S$ 是位于
$\mathfrak p \subset R$ 上方的素理想。若
$\dim_{\mathfrak q}(S/R) = n - c$，则可以找到如上的
$(h, g)$，使得 $g \not \in \mathfrak q$。
\end{enumerate}
\end{lemma}

\begin{proof}
对于 (1)，由引理 \ref{algebra-lemma-dimension-fibres-bounded-open-upstairs}，
存在包含 $V(IS)$ 的开子集 $W \subset \Spec(S)$，使得
$W \to \Spec(R)$ 的所有纤维都具有维数 $\leq n - c$。
写成 $W = \Spec(S) \setminus V(J)$。于是
$V(J) \cap V(IS) = \emptyset$，故可以找到 $g \in J$，
其像为 $1 \in S/IS$。
任取 $h \in R[x_1, \ldots, x_n]$ 为 $g$ 的原像。

\medskip\noindent
对于 (2)，由引理 \ref{algebra-lemma-dimension-fibres-bounded-open-upstairs}，
存在包含 $\Spec(S \otimes_R \kappa(\mathfrak p))$ 的开子集
$W \subset \Spec(S)$，使得 $W \to \Spec(R)$ 的所有纤维
都具有维数 $\leq n - c$。
写成 $W = \Spec(S) \setminus V(J)$。于是
$V(J \cdot S \otimes_R \kappa(\mathfrak p)) = \emptyset$。
故可以找到 $g \in J$，其在
$S \otimes_R \kappa(\mathfrak p)$ 中的像为单位（略去细节）。
任取 $h \in R[x_1, \ldots, x_n]$ 为 $g$ 的原像。

\medskip\noindent
对于 (3)，由引理 \ref{algebra-lemma-dimension-fibres-bounded-open-upstairs}，
存在 $g \in S$，$g \not \in \mathfrak q$，使得
$R \to S_g$ 的所有非空纤维都具有维数 $\leq n - c$。
任取 $h \in R[x_1, \ldots, x_n]$ 为映到 $g$ 的元素。
\end{proof}

\noindent
下面的引理说明，可以对相对全局完全交作绝对 Noether 逼近。

\begin{lemma}
\label{algebra-lemma-relative-global-complete-intersection-Noetherian}
设 $R$ 是环。设
$S = R[x_1, \ldots, x_n]/(f_1, \ldots, f_c)$
是相对全局完全交（定义
\ref{algebra-definition-relative-global-complete-intersection}）。
存在一个有限型 $\mathbf{Z}$-子代数 $R_0 \subset R$，使得
% P01-E0798: the frozen “There exist a ... subalgebra” has subject-verb disagreement.
$f_i \in R_0[x_1, \ldots, x_n]$，且
$$
S_0 = R_0[x_1, \ldots, x_n]/(f_1, \ldots, f_c)
$$
是相对全局完全交。
\end{lemma}

\begin{proof}
令 $R_0 \subset R$ 为 $R$ 中由多项式
$f_1, \ldots, f_c$ 的所有系数生成的 $\mathbf{Z}$-代数。定义：
$$
S_0 = R_0[x_1, \ldots, x_n]/(f_1, \ldots, f_c)
$$
显然 $S = R \otimes_{R_0} S_0$。
取素理想 $\mathfrak q \subset S$，并用
$\mathfrak p \subset R$、$\mathfrak q_0 \subset S_0$ 和
$\mathfrak p_0 \subset R_0$ 表示它所位于其上的素理想。
% P01-E0799: the frozen “denote ... the primes” construction omits by/as and compresses a chain of contractions into “it lies over”.
由于 $\dim (S \otimes_R \kappa(\mathfrak p) ) = n - c$，由引理
\ref{algebra-lemma-dimension-preserved-field-extension}，还有
$\dim (S_0 \otimes_{R_0} \kappa(\mathfrak p_0)) = n - c$。
由引理 \ref{algebra-lemma-dimension-fibres-bounded-open-upstairs}，
存在 $g \in S_0$，$g \not \in \mathfrak q_0$，使得
$R_0 \to (S_0)_g$ 的所有非空纤维都具有维数 $\leq n - c$。
由于 $\mathfrak q$ 是任意的且 $\Spec(S)$ 拟紧，可以找到有限多个
$g_1, \ldots, g_m \in S_0$，使得：
(a) 对 $j = 1, \ldots, m$，映射
$R_0 \to (S_0)_{g_j}$ 的非空纤维都具有维数 $\leq n - c$；
(b) $\Spec(S) \to \Spec(S_0)$ 的像包含于
$D(g_1) \cup \ldots \cup D(g_m)$。
换言之，$g_1, \ldots, g_m$ 在
$S = R \otimes_{R_0} S_0$ 中的像生成单位理想。
扩大 $R_0$ 后，可以假设 $g_1, \ldots, g_m$ 在 $S_0$
中生成单位理想。由 (a)，$R_0 \to S_0$ 的非空纤维
都具有维数 $\leq n - c$，结论得证。
\end{proof}

\begin{lemma}
\label{algebra-lemma-relative-global-complete-intersection-conormal}
设 $R$ 是环。设
$S = R[x_1, \ldots, x_n]/(f_1, \ldots, f_c)$
是相对全局完全交（定义
\ref{algebra-definition-relative-global-complete-intersection}）。
对 $S$ 的每个素理想 $\mathfrak q$，令 $\mathfrak q'$ 表示
$R[x_1, \ldots, x_n]$ 中对应的素理想。则
\begin{enumerate}
\item $f_1, \ldots, f_c$ 在局部环
$R[x_1, \ldots, x_n]_{\mathfrak q'}$ 中构成正则序列；
% P01-E0800: the frozen English uses singular “is” with the list f_1,...,f_c.
\item 每个环
$R[x_1, \ldots, x_n]_{\mathfrak q'}/(f_1, \ldots, f_i)$
在 $R$ 上平坦；并且
\item $S$-模 $(f_1, \ldots, f_c)/(f_1, \ldots, f_c)^2$
是自由的，且元素 $f_i \bmod (f_1, \ldots, f_c)^2$ 给出其一组基。
\end{enumerate}
\end{lemma}

\begin{proof}
先假设 $R$ 是 Noether 环。令 $\mathfrak p = R \cap \mathfrak q'$。
例如，由引理 \ref{algebra-lemma-lci} 可知，$f_1, \ldots, f_c$
在局部环
$R[x_1, \ldots, x_n]_{\mathfrak q'} \otimes_R \kappa(\mathfrak p)$
中构成正则序列。此外，局部环
$R[x_1, \ldots, x_n]_{\mathfrak q'}$ 在 $R_{\mathfrak p}$ 上平坦。
由于 $R$ 是 Noether 环，因而
$R[x_1, \ldots, x_n]_{\mathfrak q'}$ 也是 Noether 环，
由引理 \ref{algebra-lemma-grothendieck-regular-sequence} 可知 (1) 和 (2) 成立。

\medskip\noindent
现在设 $R$ 是一般的环。将
$R = \colim_{\lambda \in \Lambda} R_\lambda$ 写成有限型
$\mathbf{Z}$-子代数的滤过余极限（与第
\ref{algebra-section-colimits-flat} 节比较）。可以假设对所有 $\lambda$ 都有
$f_1, \ldots, f_c \in R_\lambda[x_1, \ldots, x_n]$。
令 $R_0 \subset R$ 如引理
\ref{algebra-lemma-relative-global-complete-intersection-Noetherian} 中所述。
于是可以假设对所有 $\lambda$ 都有 $R_0 \subset R_\lambda$。
由此，$S_\lambda = R_\lambda[x_1, \ldots, x_n]/(f_1, \ldots, f_c)$
是相对全局完全交（它是 $S_0$ 经由 $R_0 \to R_\lambda$
作基变换所得；见引理
\ref{algebra-lemma-base-change-relative-global-complete-intersection}）。
分别用 $\mathfrak p_\lambda$、$\mathfrak q_\lambda$、
$\mathfrak q'_\lambda$ 表示由 $\mathfrak p$、$\mathfrak q$、
$\mathfrak q'$ 诱导的 $R_\lambda$、$S_\lambda$、
$R_\lambda[x_1, \ldots, x_n]$ 中的素理想。
% P01-E0801: the frozen designation omits by/as and uses singular “the prime” for three primes.
在此记号下，对每个 $\lambda$，(1) 和 (2) 都成立。由于
$$
R[x_1, \ldots, x_n]_{\mathfrak q'}/(f_1, \ldots, f_i)
=
\colim R_\lambda[x_1, \ldots, x_n]_{\mathfrak q_\lambda'}/(f_1, \ldots, f_i)
$$
由引理 \ref{algebra-lemma-colimit-rings-flat}，可推出 (2) 中在 $R$ 上的平坦性。
又因为
\begin{align*}
R[x_1, \ldots, x_n]_{\mathfrak q'}/(f_1, \ldots, f_i)
\xrightarrow{f_{i + 1}}
R[x_1, \ldots, x_n]_{\mathfrak q'}/(f_1, \ldots, f_i) \\
=
\colim
\left(
R_\lambda[x_1, \ldots, x_n]_{\mathfrak q_\lambda'}/(f_1, \ldots, f_i)
\xrightarrow{f_{i + 1}}
R_\lambda[x_1, \ldots, x_n]_{\mathfrak q_\lambda'}/(f_1, \ldots, f_i)
\right)
\end{align*}
并且滤过余极限是正合的（引理
\ref{algebra-lemma-directed-colimit-exact}），所以 (1) 成立。

\medskip\noindent
最后证明 (3)。记 $N$ 为 $S$-模
$(f_1, \ldots, f_c)/(f_1, \ldots, f_c)^2$，并记
$e_i \in N$ 为 $f_i$ 的像。
% P01-E0802: the frozen “Denote N the S-module” construction omits by/as.
由引理 \ref{algebra-lemma-regular-quasi-regular} 和 (1)，对 $S$
的所有素理想 $\mathfrak q$，$e_1, \ldots, e_c$ 都是
$N_\mathfrak q$ 的一组基。
% P01-E0803: the frozen subject e_1,...,e_c takes singular “is a basis”; “form a basis” is required.
由引理 \ref{algebra-lemma-characterize-zero-local}，得出 (3)。
\end{proof}

\begin{lemma}
\label{algebra-lemma-relative-global-complete-intersection}
相对全局完全交是 syntomic 的，即它是平坦的。
% P01-E0804: the frozen “syntomic, i.e., flat” can misleadingly identify syntomicity with flatness outside the already-constrained relative-global-CI context.
\end{lemma}

\begin{proof}
设 $R \to S$ 是相对全局完全交。
它的纤维都是全局完全交，并且 $S$ 是有限表示 $R$-代数。
因此只须证明 $R \to S$ 平坦。这由引理
\ref{algebra-lemma-relative-global-complete-intersection-conormal} 的 (2) 得出。
\end{proof}

\begin{lemma}
\label{algebra-lemma-adjoin-roots}
设 $A$ 是环，并且
$P(x) = x^n + b_1 x^{n-1} + \ldots + b_n \in A[x]$
是 $A$ 上的首一多项式。则存在一个 syntomic、有限自由且忠实平坦的
环扩张 $A \subset A'$，使得对某些 $\beta_i \in A'$ 有
$P(x) = \prod_{i = 1, \ldots, n} (x - \beta_i)$。
\end{lemma}

\begin{proof}
取 $A' = A \otimes_R S$，其中 $R$ 和 $S$ 如例
\ref{algebra-example-roots-universal-polynomial} 中所述，$R \to A$
把 $a_i$ 映到 $b_i$；并令 $\beta_i = -1 \otimes \alpha_i$。
注意，由引理 \ref{algebra-lemma-relative-global-complete-intersection}，
$R \to S$ 忠实平坦、有限自由且 syntomic。
基变换 $A \to A'$ 继承这些性质；略去一些细节。
\end{proof}

\begin{lemma}
\label{algebra-lemma-syntomic}
设 $R \to S$ 是环映射。
设 $\mathfrak q \subset S$ 是位于 $R$ 的素理想 $\mathfrak p$
上方的素理想。下列条件等价：
\begin{enumerate}
\item 存在元素 $g \in S$，$g \not \in \mathfrak q$，使得
$R \to S_g$ 是 syntomic 的。
\item 存在元素 $g \in S$，$g \not \in \mathfrak q$，使得
$S_g$ 是 $R$ 上的相对全局完全交。
\item 存在元素 $g \in S$，$g \not \in \mathfrak q$，使得
$R \to S_g$ 是有限表示的，局部环映射
$R_{\mathfrak p} \to S_{\mathfrak q}$ 是平坦的，并且局部环
$S_{\mathfrak q}/\mathfrak pS_{\mathfrak q}$ 是
$\kappa(\mathfrak p)$ 上的完全交环（见定义
\ref{algebra-definition-lci-local-ring}）。
\end{enumerate}
\end{lemma}

\begin{proof}
蕴含 (1) $\Rightarrow$ (3) 是引理 \ref{algebra-lemma-lci-at-prime}。
蕴含 (2) $\Rightarrow$ (1) 是引理
\ref{algebra-lemma-relative-global-complete-intersection}。
还须证明 (3) 蕴含 (2)。

\medskip\noindent
假设 (3)。对某个 $g \in S$，$g\not\in \mathfrak q$，
% P01-E0805: the frozen sentence lacks a comma after the localization condition.
把 $S$ 替换为 $S_g$ 后，可以假设 $S$ 是有限表示 $R$-代数。
选取表现 $S = R[x_1, \ldots, x_n]/I$。令
$\mathfrak q' \subset R[x_1, \ldots, x_n]$ 为对应于
$\mathfrak q$ 的素理想。写 $\kappa(\mathfrak p) = k$。
注意 $S \otimes_R k = k[x_1, \ldots, x_n]/\overline{I}$，其中
$\overline{I} \subset k[x_1, \ldots, x_n]$ 是由 $I$ 的像生成的理想。
令 $\overline{\mathfrak q}' \subset k[x_1, \ldots, x_n]$
为由 $\mathfrak q'$ 的像生成的素理想。
由引理 \ref{algebra-lemma-lci-at-prime}，引理 \ref{algebra-lemma-lci}
的等价条件对 $\overline{I}$ 和 $\overline{\mathfrak q}'$ 成立。
设
$\overline{I}_{\overline{\mathfrak q}'}/
\overline{\mathfrak q}'\overline{I}_{\overline{\mathfrak q}'}$
在 $\kappa(\overline{\mathfrak q}')$ 上的维数为 $c$。
取 $f_1, \ldots, f_c \in I$，其像构成该向量空间的一组基。
像 $\overline{f}_j \in \overline{I}$ 生成
$\overline{I}_{\overline{\mathfrak q}'}$（由引理
\ref{algebra-lemma-lci}）。置
$S' = R[x_1, \ldots, x_n]/(f_1, \ldots, f_c)$。
令 $J$ 为满射 $S' \to S$ 的核。
由于 $S$ 是有限表示的，$J$ 是有限生成理想
（引理 \ref{algebra-lemma-compose-finite-type}）。考虑短正合列
$$
0 \to J \to S' \to S \to 0
$$
% P01-E0806: the frozen short exact display lacks terminal punctuation before the next sentence.
由于 $S_\mathfrak q$ 在 $R$ 上平坦，
$J_{\mathfrak q'} \otimes_R k \to S'_{\mathfrak q'} \otimes_R k$
是单射（引理 \ref{algebra-lemma-flat-tor-zero}）。
然而，按构造，$S'_{\mathfrak q'} \otimes_R k$ 同构地映到
$S_\mathfrak q \otimes_R k$。因此
$J_{\mathfrak q'} \otimes_R k =
J_{\mathfrak q'}/\mathfrak pJ_{\mathfrak q'} = 0$。
由 Nakayama 引理（引理 \ref{algebra-lemma-NAK}），存在
$g \in R[x_1, \ldots, x_n]$，$g \not \in \mathfrak q'$，使得
$J_g = 0$。换言之，$S'_g \cong S_g$。
进一步局部化后，由引理
\ref{algebra-lemma-localize-relative-complete-intersection} 可知，
$S'$（因而 $S$）成为相对全局完全交，如所要求。
\end{proof}

\begin{lemma}
\label{algebra-lemma-syntomic-presentation-ideal-mod-squares}
设 $R$ 是环。设 $S = R[x_1, \ldots, x_n]/I$，其中
$I$ 是有限生成理想。若 $g \in S$ 使得 $S_g$ 在 $R$ 上
是 syntomic 的，则 $(I/I^2)_g$ 是有限投射 $S_g$-模。
\end{lemma}

\begin{proof}
由引理 \ref{algebra-lemma-syntomic}，存在有限多个元素
$g_1, \ldots, g_m \in S$，它们在 $S_g$ 中生成单位理想，
并且每个 $S_{gg_j}$ 都是 $R$ 上的相对全局完全交。
由引理 \ref{algebra-lemma-finite-projective}，只须证明
$(I/I^2)_{gg_j}$ 有限投射，故可以假设 $S_g$
是相对全局完全交。在这种情况下，结论由引理
\ref{algebra-lemma-conormal-module-localize} 和
\ref{algebra-lemma-relative-global-complete-intersection-conormal} 得出。
\end{proof}

\begin{lemma}
\label{algebra-lemma-composition-syntomic}
设 $R \to S$、$S \to S'$ 是环映射。
\begin{enumerate}
\item 若 $R \to S$ 和 $S \to S'$ 是 syntomic 的，则
$R \to S'$ 是 syntomic 的。
\item 若 $R \to S$ 和 $S \to S'$ 是相对全局完全交，则
$R \to S'$ 是相对全局完全交。
\end{enumerate}
\end{lemma}

\begin{proof}
先证明 (2)。设 $R \to S$ 和 $S \to S'$ 是相对全局完全交，
% P01-E0807: the frozen “Proof of (2).” is a sentence fragment.
并且有定义 \ref{algebra-definition-relative-global-complete-intersection}
中的表现
$S =  R[x_1, \ldots, x_n]/(f_1, \ldots, f_c)$ 和
$S' = S[y_1, \ldots, y_m]/(h_1, \ldots, h_d)$。
则
$$
S' \cong
R[x_1, \ldots, x_n, y_1, \ldots, y_m]/(f_1, \ldots, f_c, h'_1, \ldots, h'_d)
$$
其中 $h_j' \in R[x_1, \ldots, x_n, y_1, \ldots, y_m]$
是 $h_j$ 的某些提升。
% P01-E0808: the frozen display and prose place the prime on h in inconsistent source order (h'_j versus h_j').
因此，只须对纤维环的维数给出界。于是可以假设 $R = k$ 是域。
在这种情况下，环 $S$ 是有限型 $k$-代数，且等维、维数为
$n - c$；有限型环映射 $S \to S'$ 的每个非空纤维环
都是等维的，且维数为 $m - d$。
例如，把引理 \ref{algebra-lemma-dimension-base-fibre-total}
应用于 $S'$ 在极大理想处的局部化，便可看出
$\dim(S') \leq n - c + m - d$，如所要求。

\medskip\noindent
我们把 (1) 归约到 (2)。假设 $R \to S$ 和 $S \to S'$
是 syntomic 的。设 $\mathfrak q' \subset S$ 是位于
$\mathfrak q \subset S$ 上方的素理想。
% P01-E0809: the frozen source places both q' and q in S; q' must be a prime of S' for the subsequent localization and containment statements.
由引理 \ref{algebra-lemma-syntomic}，存在
$g' \in S'$，$g' \not \in \mathfrak q'$，使得
$S \to S'_{g'}$ 是相对全局完全交。
类似地，可以找到 $g \in S$，$g \not \in \mathfrak q$，使得
$R \to S_g$ 是相对全局完全交。
由引理 \ref{algebra-lemma-base-change-relative-global-complete-intersection}，
环映射 $S_g \to S_{gg'}$ 是相对全局完全交。
% P01-E0810: the frozen base-change target is S_{gg'}, but g' belongs to S'; the target should be S'_{gg'}.
由 (2)，$R \to S_{gg'}$ 是相对全局完全交，并且
$gg' \not \in \mathfrak q'$。
% P01-E0811: the frozen composite repeats S_{gg'} where the target must be the localization of S'.
由于 $\mathfrak q'$ 是任意的，结合引理 \ref{algebra-lemma-syntomic}
和 \ref{algebra-lemma-local-syntomic}，可知 $R \to S'$ 是 syntomic 的
（这里也使用了 $S'$ 的谱拟紧；见引理
\ref{algebra-lemma-quasi-compact}）。
\end{proof}

\noindent
下面的引理稍后会得到加强；见《环同态的平滑化》的命题
\ref{smoothing-proposition-lift-smooth}。

\begin{lemma}
\label{algebra-lemma-lift-syntomic}
设 $R$ 是环，$I \subset R$ 是理想。
设 $R/I \to \overline{S}$ 是 syntomic 映射。
则存在生成 $\overline{S}$ 的单位理想的元素
$\overline{g}_i \in \overline{S}$，使得每个
% P01-E0812: the frozen “there exists elements” has subject-verb disagreement.
$\overline{S}_{g_i} \cong S_i/IS_i$，其中 $S_i$
是 $R$ 上的某个相对全局完全交。
% P01-E0813: the frozen localization uses g_i although only \overline{g}_i is introduced; the same mismatch recurs in the proof.
\end{lemma}

\begin{proof}
由引理 \ref{algebra-lemma-syntomic}，可以找到生成 $\overline{S}$
的单位理想的一族元素 $\overline{g}_i \in \overline{S}$，
使得每个 $\overline{S}_{g_i}$ 都是 $R/I$ 上的相对全局完全交。
因此可以假设 $\overline{S}$ 是相对全局完全交。写成
$\overline{S} =
(R/I)[x_1, \ldots, x_n]/(\overline{f}_1, \ldots, \overline{f}_c)$，
如定义 \ref{algebra-definition-relative-global-complete-intersection} 中所述。
选取 $f_1, \ldots, f_c \in R[x_1, \ldots, x_n]$，
它们分别提升 $\overline{f}_1, \ldots, \overline{f}_c$。
置 $S = R[x_1, \ldots, x_n]/(f_1, \ldots, f_c)$。
注意 $S/IS \cong \overline{S}$。
由引理 \ref{algebra-lemma-localize-relative-complete-intersection}，
可以找到 $g \in S$，它在 $\overline{S}$ 中的像为 $1$，
并且 $S_g$ 是 $R$ 上的相对全局完全交。
由于 $\overline{S} \cong S_g/IS_g$，证明完成。
\end{proof}


\section{光滑环映射}
\label{algebra-section-smooth}

\noindent
先用一个例子说明光滑环映射定义的动机。
设 $R$ 是环，且对某个非零 $f \in R[x, y]$ 有
$S = R[x, y]/(f)$。此时有正合列
$$
S \to
S\text{d}x \oplus S\text{d}y \to
\Omega_{S/R} \to 0
$$
其中第一支箭头把 $1$ 映到
$\frac{\partial f}{\partial x} \text{d}x +
\frac{\partial f}{\partial y} \text{d}y$；见第
% P01-E0814: the frozen sentence lacks punctuation before “see Section”.
\ref{algebra-section-netherlander} 节。
由此可知，若 $f$ 的偏导数在 $S$ 中生成单位理想，则
$\Omega_{S/R}$ 是秩为 $1$ 的局部自由模。
在这种情况下，$S$ 在 $R$ 上是相对维数为 $1$ 的光滑代数。
但 $\Omega_{S/R}$ 也可能是秩为 $2$ 的局部自由模，即 $f$
的两个偏导数都为零时如此。例如，若对某个素数 $p$，
$R$ 中有 $p = 0$，并且 $f = x^p + y^p$，就会出现这种情况。
这里 $R \to S$ 是相对维数为 $1$ 的相对全局完全交，但并不光滑。
因此，要检验环映射是否光滑，仅检验微分模是否自由是不够的。
正确的条件如下。

\begin{definition}
\label{algebra-definition-smooth}
若环映射 $R \to S$ 是有限表示的，并且朴素余切复形
$\NL_{S/R}$ 拟同构于置于次数 $0$ 的某个有限投射 $S$-模，
则称它是{\it 光滑的}；这意味着
$H_1(\NL_{S/R}) = 0$，并且 $\Omega_{S/R}$ 是有限投射 $S$-模。
\end{definition}

\noindent
设 $R \to S$ 是环映射。由引理 \ref{algebra-lemma-NL-homotopy}，
对任意表现 $\alpha : P \to S$，都有
$H_1(\NL(\alpha)) = H_1(\NL_{S/R})$ 以及
$H_0(\NL(\alpha)) = \Omega_{S/R}$。
因此，若 $R \to S$ 光滑，则对任意核为 $I$ 的满射
$\alpha : R[x_1, \ldots, x_n] \to S$，映射
$$
I/I^2
\longrightarrow
\Omega_{R[x_1, \ldots, x_n]/R} \otimes_{R[x_1, \ldots, x_n]} S
$$
都是单射，并且其余核是有限投射 $S$-模。
故所显示的箭头是分裂单射。换言之，作为 $S$-模，有
$\bigoplus_{i = 1}^n S \text{d}x_i \cong I/I^2 \oplus \Omega_{S/R}$，
而 $I/I^2$ 也是有限投射 $S$-模。
反过来，若 $R \to S$ 是有限表示的，并且存在核为 $I$ 的满射
$\alpha : R[x_1, \ldots, x_n] \to S$，使得所显示的箭头
是分裂单射，则 $R \to S$ 光滑。

\begin{lemma}
\label{algebra-lemma-localize-smooth}
设 $R \to S$ 是光滑环映射。
任意局部化 $S_g$ 在 $R$ 上都是光滑的。
若 $f \in R$ 在 $S$ 中的像可逆，则 $R_f \to S$ 光滑。
\end{lemma}

\begin{proof}
由引理 \ref{algebra-lemma-localize-NL}，$S_g$ 在 $R$ 上的朴素余切复形
是 $S$ 在 $R$ 上的朴素余切复形的基变换。
假设是 $S/R$ 的朴素余切复形拟同构于 $\Omega_{S/R}$，
% P01-E0816: the frozen proof says the naive cotangent complex “is” Omega although the definition only gives a quasi-isomorphism.
并且后者是有限投射 $S$-模。因此，它的基变换也具有同样的性质。
故 $S_g$ 在 $R$ 上光滑。

\medskip\noindent
第二个断言以同样方式由引理 \ref{algebra-lemma-NL-localize-bottom} 得出。
\end{proof}

\begin{lemma}
\label{algebra-lemma-base-change-smooth}
\begin{slogan}
光滑性在基变换下保持。
% P01-E0815: the frozen slogan lacks terminal punctuation.
\end{slogan}
设 $R \to S$ 是光滑环映射。
设 $R \to R'$ 是任意环映射。
则基变换 $R' \to S' = R' \otimes_R S$ 光滑。
\end{lemma}

\begin{proof}
设 $\alpha : R[x_1, \ldots, x_n] \to S$ 是核为 $I$ 的一个表现。
设 $\alpha' : R'[x_1, \ldots, x_n] \to R' \otimes_R S$
为所诱导的表现。令 $I' = \Ker(\alpha')$。
由于 $0 \to I \to R[x_1, \ldots, x_n] \to S \to 0$ 正合，序列
$R' \otimes_R I \to R'[x_1, \ldots, x_n] \to R' \otimes_R S \to 0$
正合。因此 $R' \otimes_R I \to I'$ 是满射。
由定义 \ref{algebra-definition-smooth}，有短正合列
$$
0 \to I/I^2 \to
\Omega_{R[x_1, \ldots, x_n]/R} \otimes_{R[x_1, \ldots, x_n]} S \to
\Omega_{S/R} \to
0
$$
并且 $S$-模 $\Omega_{S/R}$ 是有限投射的。特别地，$I/I^2$
是 $\Omega_{R[x_1, \ldots, x_n]/R}
\otimes_{R[x_1, \ldots, x_n]} S$ 的直和项。
考虑交换图
$$
\xymatrix{
R' \otimes_R (I/I^2) \ar[r] \ar[d] &
R' \otimes_R (\Omega_{R[x_1, \ldots, x_n]/R} \otimes_{R[x_1, \ldots, x_n]} S)
\ar[d] \\
I'/(I')^2 \ar[r] &
\Omega_{R'[x_1, \ldots, x_n]/R'}
\otimes_{R'[x_1, \ldots, x_n]} (R' \otimes_R S)
}
$$
由于右边的竖直映射是同构，结合上面所述，可知左边的竖直映射
既是单射也是满射。因此，$\NL(\alpha')$ 拟同构于置于次数 $0$ 的
$\Omega_{S'/R'} \cong S' \otimes_S \Omega_{S/R}$。
该模是有限投射的，因为它是有限投射模的基变换。
\end{proof}

\begin{lemma}
\label{algebra-lemma-smooth-over-field}
设 $k$ 是域。
设 $S$ 是光滑 $k$-代数。
则 $S$ 是局部完全交。
\end{lemma}

\begin{proof}
由引理 \ref{algebra-lemma-base-change-smooth} 和
\ref{algebra-lemma-lci-field-change}，只须证明 $k$ 为代数闭域的情形。
选取核为 $I$ 的表现
$\alpha : k[x_1, \ldots, x_n] \to S$。
设 $\mathfrak m$ 是 $S$ 的极大理想，并令
$\mathfrak m' \supset I$ 为 $k[x_1, \ldots, x_n]$
中对应的极大理想。我们将证明引理 \ref{algebra-lemma-lci}
的条件 (5) 成立（以 $\mathfrak m$ 代替 $\mathfrak q$）。
因为 $k$ 是代数闭域，可以对某些 $a_i \in k$ 写成
$\mathfrak m' = (x_1 - a_1, \ldots, x_n - a_n)$；见定理
\ref{algebra-theorem-nullstellensatz}。
由 $k \to S$ 光滑这一假设，$S$-模映射
$\text{d} : I/I^2 \to \bigoplus_{i = 1}^n S \text{d}x_i$
是分裂单射。因此，相应的映射
$I/\mathfrak m' I \to \bigoplus \kappa(\mathfrak m') \text{d}x_i$
是单射。设
$\dim_{\kappa(\mathfrak m')}(I/\mathfrak m' I) = c$，
并取 $f_1, \ldots, f_c \in I$，其像构成
$I/\mathfrak m' I$ 的一组 $\kappa(\mathfrak m')$-基。
由 Nakayama 引理 \ref{algebra-lemma-NAK}，$f_1, \ldots, f_c$
在 $k[x_1, \ldots, x_n]_{\mathfrak m'}$ 上生成
$I_{\mathfrak m'}$。考虑交换图
$$
\xymatrix{
I \ar[r] \ar[d] & I/I^2 \ar[rr] \ar[d] & &
I/\mathfrak m'I \ar[d] \\
\Omega_{k[x_1, \ldots, x_n]/k} \ar[r] &
\bigoplus S\text{d}x_i \ar[rr]^{\text{d}x_i \mapsto x_i - a_i} & &
\mathfrak m'/(\mathfrak m')^2
}
$$
（略去交换性的证明）。中间的竖直映射就是定义
$\alpha$ 的朴素余切复形的映射。注意，右下方的水平箭头诱导同构
$\bigoplus \kappa(\mathfrak m') \text{d}x_i
\to \mathfrak m'/(\mathfrak m')^2$。
因此，$I_{\mathfrak m'}$ 的生成元 $f_1, \ldots, f_c$
映到 $k[x_1, \ldots, x_n]_{\mathfrak m'}$ 中的一族元素，
它们在 $\mathfrak m'/(\mathfrak m')^2$ 中的类在
$\kappa(\mathfrak m')$ 上线性无关。
故由引理 \ref{algebra-lemma-regular-ring-CM}，它们在环
$k[x_1, \ldots, x_n]_{\mathfrak m'}$ 中构成正则序列。
这验证了引理 \ref{algebra-lemma-lci} 的条件 (5)，
% P01-E0817: the frozen source joins the lemma citation directly to “hence” without punctuation.
因而对某个 $g \in S$，$g \not \in \mathfrak m$，$S_g$
是 $k$ 上的全局完全交。由于这对 $S$ 的任意极大理想都成立，
可知 $S$ 是 $k$ 上的局部完全交。
\end{proof}

\begin{definition}
\label{algebra-definition-standard-smooth}
设 $R$ 是环。给定整数 $n \geq c \geq 0$ 以及
$f_1, \ldots, f_c \in R[x_1, \ldots, x_n]$，称
% P01-E0818: the frozen fronted data lack a comma before “we say”.
$$
R[x_1, \ldots, x_n]/(f_1, \ldots, f_c)
$$
为 $R$ 上的一个{\it 标准光滑代数}，如果多项式
$$
g =
\det
\left(
\begin{matrix}
\partial f_1/\partial x_1 &
\partial f_2/\partial x_1 &
\ldots &
\partial f_c/\partial x_1 \\
\partial f_1/\partial x_2 &
\partial f_2/\partial x_2 &
\ldots &
\partial f_c/\partial x_2 \\
\ldots & \ldots & \ldots & \ldots \\
\partial f_1/\partial x_c &
\partial f_2/\partial x_c &
\ldots &
\partial f_c/\partial x_c
\end{matrix}
\right)
$$
在 $R[x_1, \ldots, x_n]/(f_1, \ldots, f_c)$ 中的像可逆。
若存在 $n \geq c \geq 0$ 和
$f_1, \ldots, f_c \in R[x_1, \ldots, x_n]$，使得下列代数：
$$
R[x_1, \ldots, x_n]/(f_1, \ldots, f_c)
$$
是 $R$ 上的标准光滑代数，并且 $S$ 作为 $R$-代数与下列代数同构：
$$
R[x_1, \ldots, x_n]/(f_1, \ldots, f_c)
$$
则称 $R$-代数 $S$ 是{\it 标准光滑的}，也称环映射
$R \to S$ 是{\it 标准光滑的}。
\end{definition}

\begin{lemma}
\label{algebra-lemma-standard-smooth}
设
$S = R[x_1, \ldots, x_n]/(f_1, \ldots, f_c)
= R[x_1, \ldots, x_n]/I$
是标准光滑代数。则
\begin{enumerate}
\item 环映射 $R \to S$ 光滑；
\item $S$-模 $\Omega_{S/R}$ 以
$\text{d}x_{c + 1}, \ldots, \text{d}x_n$ 为一组基；
\item $S$-模 $I/I^2$ 以 $f_1, \ldots, f_c$ 的类为一组基；
\item 对任意 $g \in S$，环映射 $R \to S_g$ 标准光滑；
\item 对任意环映射 $R \to R'$，基变换
$R' \to R'\otimes_R S$ 标准光滑；
\item 若 $f \in R$ 在 $S$ 中的像可逆，则 $R_f \to S$ 标准光滑；并且
\item 环 $S$ 是 $R$ 上的相对全局完全交。
\end{enumerate}
\end{lemma}

\begin{proof}
考虑给定表现的朴素余切复形
$$
(f_1, \ldots, f_c)/(f_1, \ldots, f_c)^2
\longrightarrow
\bigoplus\nolimits_{i = 1}^n S \text{d}x_i
$$
% P01-E0819: the frozen display lacks terminal punctuation before the next sentence.
将此映射与到直和的前 $c$ 个直和项上的投影复合。
按照标准光滑代数的定义，类
$f_i \bmod (f_1, \ldots, f_c)^2$ 映到
$\bigoplus_{i = 1}^c S\text{d}x_i$ 的一组基。
由此得出 $(f_1, \ldots, f_c)/(f_1, \ldots, f_c)^2$
是秩为 $c$ 的自由模，元素
$f_i \bmod (f_1, \ldots, f_c)^2$ 给出其一组基；并且朴素余切复形
在次数 $0$ 的同调，即 $\Omega_{S/R}$，是一个自由 $S$-模，
其一组基为 $\text{d}x_{c + j}$（$j = 1, \ldots, n - c$）的像。
特别地，这证明 $R \to S$ 光滑。

\medskip\noindent
略去 (4) 和 (6) 的证明。不过，可参见下面的例子以及引理
\ref{algebra-lemma-base-change-relative-global-complete-intersection} 的证明。

\medskip\noindent
设 $\varphi : R \to R'$ 是任意环映射。
% P01-E0820: the frozen “Denote S' = ... where ...” construction omits by/as.
置 $S' = R'[x_1, \ldots, x_n]/(f_1^\varphi, \ldots, f_c^\varphi)$，
其中 $f^\varphi$ 是对 $f \in R[x_1, \ldots, x_n]$ 的所有系数
应用 $\varphi$ 后得到的多项式。则 $S' \cong R' \otimes_R S$。
此外，定义 \ref{algebra-definition-standard-smooth} 中 $S'/R'$
的行列式等于 $g^\varphi$。因此，它在 $S'$ 中的像就是 $g$
经由 $R[x_1, \ldots, x_n] \to S \to S'$ 的像，故可逆。
这证明了 (5)。

\medskip\noindent
要证明 (7)，只须证明对每个素理想 $\mathfrak p \subset R$，
$S \otimes_R \kappa(\mathfrak p)$ 的维数都是 $n - c$。
由 (5)，只须证明域 $k$ 上任意标准光滑代数
$k[x_1, \ldots, x_n]/(f_1, \ldots, f_c)$ 的维数为 $n - c$。
由引理 \ref{algebra-lemma-smooth-over-field}，我们已经知道
$k[x_1, \ldots, x_n]/(f_1, \ldots, f_c)$ 是局部完全交。
因此，由于 $I/I^2$ 是秩为 $c$ 的自由模，例如由引理
\ref{algebra-lemma-lci} 可知
$k[x_1, \ldots, x_n]/(f_1, \ldots, f_c)$ 的维数为 $n - c$。
\end{proof}

\begin{example}
\label{algebra-example-make-standard-smooth}
设 $R$ 是环。
设 $f_1, \ldots, f_c \in R[x_1, \ldots, x_n]$。
令
$$
h =
\det
\left(
\begin{matrix}
\partial f_1/\partial x_1 &
\partial f_2/\partial x_1 &
\ldots &
\partial f_c/\partial x_1 \\
\partial f_1/\partial x_2 &
\partial f_2/\partial x_2 &
\ldots &
\partial f_c/\partial x_2 \\
\ldots & \ldots & \ldots & \ldots \\
\partial f_1/\partial x_c &
\partial f_2/\partial x_c &
\ldots &
\partial f_c/\partial x_c
\end{matrix}
\right).
$$
置 $S = R[x_1, \ldots, x_{n + 1}]/(f_1, \ldots, f_c, x_{n + 1}h - 1)$。
这是标准光滑代数的一个例子，只不过该表现中的变量次序不对，
应当依次为
$x_1, \ldots, x_c, x_{n + 1}, x_{c + 1}, \ldots, x_n$。
\end{example}

\begin{lemma}
\label{algebra-lemma-compose-standard-smooth}
标准光滑环映射的合成仍是标准光滑环映射。
\end{lemma}

\begin{proof}
假设 $R \to S$ 和 $S \to S'$ 标准光滑。选取表现
$S =  R[x_1, \ldots, x_n]/(f_1, \ldots, f_c)$
以及
$S' = S[y_1, \ldots, y_m]/(g_1, \ldots, g_d)$。
选取映到 $g_j$ 的元素
$g_j' \in R[x_1, \ldots, x_n, y_1, \ldots, y_m]$。
由此可见
$S' = R[x_1, \ldots, x_n, y_1, \ldots, y_m]/
(f_1, \ldots, f_c, g'_1, \ldots, g'_d)$。
% P01-E0821: the frozen source alternates g_j' and g'_j for the same lifts.
要证明 $S'$ 标准光滑，只须验证行列式
$$
\det
\left(
\begin{matrix}
\partial f_1/\partial x_1 &
\ldots &
\partial f_c/\partial x_1 &
\partial g_1/\partial x_1 &
\ldots &
\partial g_d/\partial x_1 \\
\ldots &
\ldots &
\ldots &
\ldots &
\ldots &
\ldots  \\
\partial f_1/\partial x_c &
\ldots &
\partial f_c/\partial x_c &
\partial g_1/\partial x_c &
\ldots &
\partial g_d/\partial x_c \\
0 &
\ldots &
0 &
\partial g_1/\partial y_1 &
\ldots &
\partial g_d/\partial y_1 \\
\ldots &
\ldots &
\ldots &
\ldots &
\ldots &
\ldots  \\
0 &
\ldots &
0 &
\partial g_1/\partial y_d &
\ldots &
\partial g_d/\partial y_d
\end{matrix}
\right)
$$
在 $S'$ 中可逆。这是显然的，因为它是由假设可逆的两个行列式
的乘积。
\end{proof}

\begin{lemma}
\label{algebra-lemma-smooth-syntomic}
设 $R \to S$ 是光滑环映射。
存在由标准开集 $D(g)$ 构成的 $\Spec(S)$ 的开覆盖，使得每个
$S_g$ 在 $R$ 上标准光滑。特别地，$R \to S$ 是 syntomic 的。
\end{lemma}

\begin{proof}
选取核为 $I = (f_1, \ldots, f_m)$ 的表现
$\alpha : R[x_1, \ldots, x_n] \to S$。
对每个子集 $E \subset \{1, \ldots, m\}$，考虑开子集
$U_E$：在其上，类 $f_e, e\in E$ 自由生成有限投射
$S$-模 $I/I^2$；见引理 \ref{algebra-lemma-cokernel-flat}。
可以用标准开集 $D(g)$ 覆盖 $\Spec(S)$，使每个这样的开集
都完全包含于某个 $U_E$。对这样的 $g$，考虑表现
$$
\beta : R[x_1, \ldots, x_n, x_{n + 1}] \longrightarrow S_g
$$
它把 $x_{n + 1}$ 映到 $1/g$。置 $J = \Ker(\beta)$。
由引理 \ref{algebra-lemma-principal-localization-NL}，
$J/J^2 \cong (I/I^2)_g \oplus S_g$ 是自由模。
我们可以且确实把 $S$ 替换为 $S_g$。于是，利用引理
\ref{algebra-lemma-huber}，可以假设有核为
$I = (f_1, \ldots, f_c)$ 的表现
$\alpha : R[x_1, \ldots, x_n] \to S$，并且 $I/I^2$
以 $f_1, \ldots, f_c$ 的类为一组基。

\medskip\noindent
利用上一段末尾得到的表现 $\alpha$，对
$\Omega_{R[x_1, \ldots, x_n]/R}$ 的基元素
$\text{d}x_1, \ldots, \text{d}x_n$ 作大致相同的论证。
具体地，对任意基数为 $c$ 的子集
$E \subset \{1, \ldots, n\}$，考虑 $\Spec(S)$ 的开子集
$U_E$：在其上，$\NL(\alpha)$ 的微分与投影的复合
$$
S^{\oplus c} \cong I/I^2
\longrightarrow
\Omega_{R[x_1, \ldots, x_n]/R} \otimes_{R[x_1, \ldots, x_n]} S
\longrightarrow
\bigoplus\nolimits_{i \in E} S\text{d}x_i
$$
是同构。再次可以找到由（有限多个）标准开集 $D(g)$ 构成的
$\Spec(S)$ 的覆盖，使每个 $D(g)$ 完全包含于某个 $U_E$。
重编号后，可以假设 $E = \{1, \ldots, c\}$。
对满足 $D(g) \subset U_E$ 的 $g$，考虑表现
$$
\beta : R[x_1, \ldots, x_n, x_{n + 1}] \to S_g
$$
它把 $x_{n + 1}$ 映到 $1/g$。置 $J = \Ker(\beta)$。
由引理 \ref{algebra-lemma-principal-localization-NL} 可知
$J = (f_1, \ldots, f_c, fx_{n + 1} - 1)$，其中
$\alpha(f) = g$，并且复合
$$
J/J^2 \longrightarrow
\Omega_{R[x_1, \ldots, x_{n + 1}]/R}
\otimes_{R[x_1, \ldots, x_{n + 1}]} S_g
\longrightarrow
\bigoplus\nolimits_{i = 1}^c S_g\text{d}x_i \oplus S_g \text{d}x_{n + 1}
$$
是同构。将坐标重排为
$x_1, \ldots, x_c, x_{n + 1}, x_{c + 1}, \ldots, x_n$，
便得 $S_g$ 在 $R$ 上标准光滑，如所要求。

\medskip\noindent
这就完成了证明，因为标准光滑代数是 syntomic 的
（引理 \ref{algebra-lemma-standard-smooth} 和
\ref{algebra-lemma-relative-global-complete-intersection}），
而在 $R$ 上为 syntomic 的性质在 $S$ 上是局部的
（引理 \ref{algebra-lemma-local-syntomic}）。
\end{proof}

\begin{definition}
\label{algebra-definition-smooth-at-prime}
设 $R \to S$ 是环映射。
设 $\mathfrak q$ 是 $S$ 的素理想。
若存在 $g \in S$，$g \not \in \mathfrak q$，使得
$R \to S_g$ 光滑，则称 $R \to S$ {\it 在 $\mathfrak q$ 处光滑}。
\end{definition}

\noindent
对有限表示环映射，可以如下刻画这一性质。

\begin{lemma}
\label{algebra-lemma-smooth-at-point}
设 $R \to S$ 是有限表示的。设 $\mathfrak q$ 是 $S$ 的素理想。
下列条件等价：
% P01-E0822: the frozen statement lacks a colon before the enumeration.
\begin{enumerate}
\item $R \to S$ 在 $\mathfrak q$ 处光滑；
\item $H_1(L_{S/R})_\mathfrak q = 0$，并且
$\Omega_{S/R, \mathfrak q}$ 是有限自由 $S_\mathfrak q$-模；
\item $H_1(L_{S/R})_\mathfrak q = 0$，并且
$\Omega_{S/R, \mathfrak q}$ 是投射 $S_\mathfrak q$-模；并且
\item $H_1(L_{S/R})_\mathfrak q = 0$，并且
$\Omega_{S/R, \mathfrak q}$ 是平坦 $S_\mathfrak q$-模。
% P01-E0823: the frozen criterion switches from the section's \NL_{S/R} notation to an otherwise unexplained L_{S/R}.
\end{enumerate}
\end{lemma}

\begin{proof}
我们将不再另行说明地使用朴素余切复形的形成与局部化可交换这一事实；
见第 \ref{algebra-section-netherlander} 节，特别是引理
\ref{algebra-lemma-localize-NL}。
注意 $\Omega_{S/R}$ 是有限表示 $S$-模；见引理
\ref{algebra-lemma-differentials-finitely-presented}。
因此，由引理 \ref{algebra-lemma-finite-projective}，(2)、(3) 和 (4) 等价。
显然 (1) 蕴含等价条件 (2)、(3) 和 (4)。
假设 (2) 成立。把 $S_\mathfrak q$ 写成主局部化的余极限，
由引理 \ref{algebra-lemma-colimit-category-fp-modules} 可知，存在
$g \in S$，$g \not \in \mathfrak q$，使得
$(\Omega_{S/R})_g$ 有限自由。
选取核为 $I$ 的表现 $\alpha : R[x_1, \ldots, x_n] \to S$。
由引理 \ref{algebra-lemma-NL-homotopy}，可以用 $\NL(\alpha)$
代替 $\NL_{S/R}$。满射
$$
\Omega_{R[x_1, \ldots, x_n]/R} \otimes_{R[x_1, \ldots, x_n]} S
\to \Omega_{S/R} \to 0
$$
在把 $g$ 变为可逆元后有右逆，因为 $(\Omega_{S/R})_g$ 是投射模。
因此
$\text{d} : (I/I^2)_g \to
\Omega_{R[x_1, \ldots, x_n]/R}
\otimes_{R[x_1, \ldots, x_n]} S_g$
的像是直和项，而且该映射也有右逆。
% P01-E0824: a split cokernel makes Im(d) a direct summand, but d itself need not have a right inverse onto the whole free target; only the surjection onto its image splits.
由此得出 $H_1(L_{S/R})_g$ 是 $(I/I^2)_g$ 的商。
特别地，$H_1(L_{S/R})_g$ 是有限 $S_g$-模。
因此，$H_1(L_{S/R})_{\mathfrak q}$ 的消失蕴含：对某个
$g' \in S$，$g' \not \in \mathfrak q$，
$H_1(L_{S/R})_{gg'}$ 消失。于是按照定义，
$R \to S_{gg'}$ 光滑。
\end{proof}

\begin{lemma}
\label{algebra-lemma-locally-smooth}
\begin{slogan}
环映射光滑，当且仅当它在目标的所有素理想处光滑。
% P01-E0825: the frozen slogan lacks terminal punctuation.
\end{slogan}
设 $R \to S$ 是环映射。
则 $R \to S$ 光滑，当且仅当 $R \to S$ 在 $S$
的每个素理想 $\mathfrak q$ 处光滑。
\end{lemma}

\begin{proof}
正向蕴含是平凡的。假设 $R \to S$ 在 $S$ 的每个素理想
$\mathfrak q$ 处都光滑。由于 $\Spec(S)$ 拟紧（见引理
\ref{algebra-lemma-quasi-compact}），存在有限覆盖
$\Spec(S) = \bigcup D(g_i)$，使得每个 $S_{g_i}$ 都光滑。
由引理 \ref{algebra-lemma-cover-upstairs}，这蕴含 $S$ 是有限表示
$R$-代数。由引理 \ref{algebra-lemma-localize-NL}，
$\NL_{S/R} \otimes_S S_{g_i}$ 拟同构于置于次数 $0$
的某个有限投射 $S_{g_i}$-模。
由引理 \ref{algebra-lemma-finite-projective}，这蕴含 $\NL_{S/R}$
拟同构于置于次数 $0$ 的某个有限投射 $S$-模。
\end{proof}

\begin{lemma}
\label{algebra-lemma-compose-smooth}
光滑环映射的合成仍是光滑环映射。
\end{lemma}

\begin{proof}
这可以用许多不同的方法证明。一种方法是使用蛇形引理（引理
\ref{algebra-lemma-snake}）、Jacobi--Zariski 序列（引理
\ref{algebra-lemma-exact-sequence-NL}），并结合投射模是自由模直和项的刻画
（引理 \ref{algebra-lemma-characterize-projective}）。
另一个证明可以通过结合引理 \ref{algebra-lemma-smooth-syntomic}、
\ref{algebra-lemma-compose-standard-smooth} 和
\ref{algebra-lemma-locally-smooth} 得到。
\end{proof}

\begin{lemma}
\label{algebra-lemma-product-smooth}
设 $R$ 是环。设 $S = S' \times S''$ 是 $R$-代数的乘积。
则 $S$ 在 $R$ 上光滑，当且仅当 $S'$ 和 $S''$
都在 $R$ 上光滑。
\end{lemma}

\begin{proof}
略。提示：由引理 \ref{algebra-lemma-locally-smooth}，可以逐个素理想检验
光滑性。由引理 \ref{algebra-lemma-spec-product}，
$\Spec(S)$ 是 $\Spec(S')$ 与 $\Spec(S'')$ 的不交并；
因此，$R \to S$ 在 $\mathfrak q$ 处的光滑性对应于
$R \to S'$ 在相应素理想处的光滑性，或者
$R \to S''$ 在相应素理想处的光滑性。
\end{proof}

\begin{lemma}
\label{algebra-lemma-relative-global-complete-intersection-smooth}
设 $R$ 是环。设
$S = R[x_1, \ldots, x_n]/(f_1, \ldots, f_c)$
是相对全局完全交。
设 $\mathfrak q \subset S$ 是素理想。则 $R \to S$
在 $\mathfrak q$ 处光滑，当且仅当存在基数为 $c$ 的子集
$I \subset \{1, \ldots, n\}$，使得多项式
$$
g_I = \det (\partial f_j/\partial x_i)_{j = 1, \ldots, c, \ i \in I}.
$$
不映到 $\mathfrak q$ 的元素。
% P01-E0826: the frozen display includes a period and is followed by the predicate fragment “does not map”.
\end{lemma}

\begin{proof}
由引理 \ref{algebra-lemma-relative-global-complete-intersection-conormal}，
与 $S$ 的给定表现相伴的朴素余切复形为
$$
\bigoplus\nolimits_{j = 1}^c S \cdot f_j
\longrightarrow
\bigoplus\nolimits_{i = 1}^n S \cdot \text{d}x_i, \quad
f_j \longmapsto \sum \frac{\partial f_j}{\partial x_i} \text{d}x_i.
$$
% P01-E0827: the frozen summation has no index range even though the target direct sum runs i=1,...,n.
给出该映射的矩阵之极大子式正是多项式 $g_I$。

\medskip\noindent
假设 $g_I$ 映到 $g \in S$，其中 $g \not \in \mathfrak q$。
则代数 $S_g$ 在 $R$ 上光滑。事实上，由引理
\ref{algebra-lemma-localize-NL}，它的朴素余切复形拟同构于上面的复形
在 $g$ 处的局部化。而按构造，它拟同构于置于次数 $0$
的秩为 $n - c$ 的自由模。
% P01-E0828: the frozen phrase “a free rank n-c module” is nonidiomatic.

\medskip\noindent
反过来，假设所有 $g_I$ 最终都属于 $\mathfrak q$。
在这种情况下，上面的复形与 $\kappa(\mathfrak q)$ 作张量后
不具有极大秩；因而不存在元素 $g \in S$、$g \not \in \mathfrak q$，
使得在它处局部化后该映射成为分裂单射。
再次由引理 \ref{algebra-lemma-localize-NL}，不存在这样的在 $R$ 上光滑的
局部化。
\end{proof}

\begin{lemma}
\label{algebra-lemma-flat-fibre-smooth}
设 $R \to S$ 是环映射。
设 $\mathfrak q \subset S$ 是位于 $R$ 的素理想
$\mathfrak p$ 上方的素理想。假设
\begin{enumerate}
\item 存在 $g \in S$，$g \not\in \mathfrak q$，使得
$R \to S_g$ 是有限表示的；
\item 局部环同态 $R_{\mathfrak p} \to S_{\mathfrak q}$ 是平坦的；
\item 纤维 $S \otimes_R \kappa(\mathfrak p)$ 在对应于
$\mathfrak q$ 的素理想处是 $\kappa(\mathfrak p)$ 上光滑的。
\end{enumerate}
则 $R \to S$ 在 $\mathfrak q$ 处光滑。
\end{lemma}

\begin{proof}
由引理 \ref{algebra-lemma-syntomic} 和 \ref{algebra-lemma-smooth-over-field}，
可知存在 $g \in S$，使得 $S_g$ 是相对全局完全交。
% P01-E0829: the frozen application omits the required condition g notin q before replacing S by S_g.
把 $S$ 替换为 $S_g$ 后，可以假设
$S = R[x_1, \ldots, x_n]/(f_1, \ldots, f_c)$
是相对全局完全交。
对任意基数为 $c$ 的子集 $I \subset \{1, \ldots, n\}$，
考虑引理 \ref{algebra-lemma-relative-global-complete-intersection-smooth}
中的多项式
$g_I = \det (\partial f_j/\partial x_i)_{j = 1, \ldots, c, i \in I}$。
注意，$g_I$ 在多项式环
$\kappa(\mathfrak p)[x_1, \ldots, x_n]$ 中的像
$\overline{g}_I$，是 $f_j$ 的像
$\overline{f}_j$ 在环
$\kappa(\mathfrak p)[x_1, \ldots, x_n]$ 中的偏导数之行列式。
因此，把引理
\ref{algebra-lemma-relative-global-complete-intersection-smooth}
分别应用于 $R \to S$ 以及
$\kappa(\mathfrak p) \to S \otimes_R \kappa(\mathfrak p)$，
即可得出本引理。
\end{proof}

\noindent
注意，下面引理中的集合 $U, V$ 按定义是开集。

\begin{lemma}
\label{algebra-lemma-flat-base-change-locus-smooth}
设 $R \to S$ 是有限表示环映射。
设 $R \to R'$ 是平坦环映射。
% P01-E0830: the frozen designation “Denote S'=... the base change” omits by/as.
记 $S' = R' \otimes_R S$ 为基变换。
令 $U \subset \Spec(S)$ 为 $R \to S$ 光滑处的素理想集合。
令 $V \subset \Spec(S')$ 为 $R' \to S'$ 光滑处的素理想集合。
% P01-E0831: the frozen “Let V ... the set” omits the verb “be”.
则在映射 $f : \Spec(S') \to \Spec(S)$ 下，$V$ 是 $U$ 的逆像。
\end{lemma}

\begin{proof}
由引理 \ref{algebra-lemma-change-base-NL}，$\NL_{S/R} \otimes_S S'$
同伦等价于 $\NL_{S'/R'}$。这已经蕴含
$f^{-1}(U) \subset V$。

\medskip\noindent
设 $\mathfrak q' \subset S'$ 是位于
$\mathfrak q \subset S$ 上方的素理想。假设 $\mathfrak q' \in V$。
须证明 $\mathfrak q \in U$。
由于 $S \to S'$ 平坦，$S_{\mathfrak q} \to S'_{\mathfrak q'}$
忠实平坦（引理 \ref{algebra-lemma-local-flat-ff}）。
因此，$H_1(L_{S'/R'})_{\mathfrak q'}$ 的消失蕴含
$H_1(L_{S/R})_{\mathfrak q}$ 的消失。
% P01-E0832: the frozen proof again switches from \NL to the unexplained L notation.
把引理 \ref{algebra-lemma-finite-projective-descends} 应用于
$S_{\mathfrak q}$-模 $(\Omega_{S/R})_{\mathfrak q}$
以及映射 $S_{\mathfrak q} \to S'_{\mathfrak q'}$，可知
$(\Omega_{S/R})_{\mathfrak q}$ 是投射模。
故由引理 \ref{algebra-lemma-smooth-at-point}，$R \to S$
在 $\mathfrak q$ 处光滑。
\end{proof}

\begin{lemma}
\label{algebra-lemma-smooth-field-change-local}
设 $K/k$ 是域扩张。
设 $S$ 是有限型 $k$-代数。
设 $\mathfrak q_K$ 是 $S_K = K \otimes_k S$ 的素理想，
并令 $\mathfrak q$ 为 $S$ 中对应的素理想。
则 $S$ 在 $\mathfrak q$ 处是 $k$ 上光滑的，当且仅当
$S_K$ 在 $\mathfrak q_K$ 处是 $K$ 上光滑的。
\end{lemma}

\begin{proof}
这是引理 \ref{algebra-lemma-flat-base-change-locus-smooth} 的特殊情形。
\end{proof}

\begin{lemma}
\label{algebra-lemma-lift-smooth}
设 $R$ 是环，$I \subset R$ 是理想。
设 $R/I \to \overline{S}$ 是光滑环映射。
则存在生成 $\overline{S}$ 的单位理想的元素
$\overline{g}_i \in \overline{S}$，使得每个
% P01-E0833: the frozen “there exists elements” has subject-verb disagreement.
$\overline{S}_{g_i} \cong S_i/IS_i$，其中 $S_i$ 是
$R$ 上的某个（标准）光滑环。
% P01-E0834: the frozen localization drops the bar from the only introduced localizer, and the proof repeats the mismatch.
\end{lemma}

\begin{proof}
由引理 \ref{algebra-lemma-smooth-syntomic}，可以找到生成
$\overline{S}$ 的单位理想的一族元素
$\overline{g}_i \in \overline{S}$，使得每个
$\overline{S}_{g_i}$ 都是 $R/I$ 上的标准光滑代数。
因此可以假设 $\overline{S}$ 是 $R/I$ 上的标准光滑代数。
按照定义 \ref{algebra-definition-standard-smooth}，写成
$\overline{S} =
(R/I)[x_1, \ldots, x_n]/(\overline{f}_1, \ldots, \overline{f}_c)$。
选取 $f_1, \ldots, f_c \in R[x_1, \ldots, x_n]$，
它们分别提升 $\overline{f}_1, \ldots, \overline{f}_c$。置
$S = R[x_1, \ldots, x_n, x_{n + 1}]/(f_1, \ldots, f_c, x_{n + 1}\Delta - 1)$，
其中
$\Delta = \det(\frac{\partial f_j}{\partial x_i})_{i, j = 1, \ldots, c}$，
如例 \ref{algebra-example-make-standard-smooth} 中所述。这证明了引理。
\end{proof}

\section{形式光滑映射}
\label{algebra-section-formally-smooth}

\noindent
本节定义形式光滑环映射。我们将会看到，有限表现的环映射形式光滑，
当且仅当它光滑，见命题 \ref{algebra-proposition-smooth-formally-smooth}。

\begin{definition}
\label{algebra-definition-formally-smooth}
设 $R \to S$ 是环映射。若对每个由实线箭头组成的交换图
$$
\xymatrix{
S \ar[r] \ar@{-->}[rd] & A/I \\
R \ar[r] \ar[u] & A \ar[u]
}
$$
其中 $I \subset A$ 是平方为零的理想，都存在使图交换的虚线箭头，
则称 $S$ 在 $R$ 上\textit{形式光滑}。
\end{definition}

\begin{lemma}
\label{algebra-lemma-base-change-fs}
设 $R \to S$ 是形式光滑环映射，$R \to R'$ 是任意环映射。
则基变换 $S' = R' \otimes_R S$ 在 $R'$ 上形式光滑。
\end{lemma}

\begin{proof}
给定定义 \ref{algebra-definition-formally-smooth} 中的实线图
$$
\xymatrix{
S \ar[r] \ar@{-->}[rrd] & R' \otimes_R S \ar[r] \ar@{-->}[rd] & A/I \\
R  \ar[u] \ar[r] & R' \ar[r] \ar[u] & A \ar[u]
}
$$
依假设，较长的虚线箭头存在。由张量积的泛性质，得到较短的虚线箭头。
\end{proof}

\begin{lemma}
\label{algebra-lemma-compose-formally-smooth}
形式光滑环映射的合成形式光滑。
\end{lemma}

\begin{proof}
略。（提示：这是完全形式的，由考察一个适当的图即可得到。）
\end{proof}

\begin{lemma}
\label{algebra-lemma-polynomial-ring-formally-smooth}
$R$ 上的多项式环在 $R$ 上形式光滑。
\end{lemma}

\begin{proof}
假设有定义 \ref{algebra-definition-formally-smooth} 中的图，且
$S = R[x_j; j \in J]$。只需对上方水平箭头下每个 $x_j$ 在 $A/I$
中的像选取一个提升 $a_j \in A$，便得到一条虚线箭头。
% P01-E0835: the frozen English has the redundant construction “to which ... map to”.
\end{proof}

\begin{lemma}
\label{algebra-lemma-characterize-formally-smooth}
设 $R \to S$ 是环映射。设 $P \to S$ 是从多项式环 $P$ 到 $S$ 的满
$R$-代数映射，并以 $J \subset P$ 表示其核。则 $R \to S$ 形式光滑，
当且仅当存在 $R$-代数映射 $
\sigma : S \to P/J^2$，它是满射 $P/J^2 \to S$ 的右逆。
% P01-E0836: the frozen English “Denote J ... the kernel” omits “by/as”.
\end{lemma}

\begin{proof}
假设 $R \to S$ 形式光滑。考察交换图
$$
\xymatrix{
S \ar[r] \ar@{-->}[rd] & P/J \\
R \ar[r] \ar[u] &  P/J^2\ar[u]
}
$$
依假设，虚线箭头存在。这证明了 $
\sigma$ 存在。

\medskip\noindent
反之，假设有引理中的 $
\sigma$。给定定义 \ref{algebra-definition-formally-smooth} 中的实线图
$$
\xymatrix{
S \ar[r] \ar@{-->}[rd] & A/I \\
R \ar[r] \ar[u] & A \ar[u]
}
$$
由引理 \ref{algebra-lemma-polynomial-ring-formally-smooth}，$P$ 形式光滑，
故存在提升映射 $P \to S \to A/I$ 的 $R$-代数同态
$\psi : P \to A$。显然 $
\psi(J) \subset I$，而 $I^2 = 0$，所以 $
\psi(J^2) = 0$。因而 $
\psi$ 分解为 $
\overline{\psi} : P/J^2 \to A$。所求虚线箭头即合成
$\overline{\psi} \circ \sigma : S \to A$。
\end{proof}

\begin{remark}
\label{algebra-remark-lemma-characterize-formally-smooth}
只要 $P$ 在 $R$ 上形式光滑，引理
\ref{algebra-lemma-characterize-formally-smooth} 便仍然成立。
\end{remark}

\begin{lemma}
\label{algebra-lemma-characterize-formally-smooth-again}
设 $R \to S$ 是环映射。设 $P \to S$ 是从多项式环 $P$ 到 $S$ 的满
$R$-代数映射，并以 $J \subset P$ 表示其核。则 $R \to S$ 形式光滑，
当且仅当引理 \ref{algebra-lemma-differential-seq} 中的序列
$$
0 \to J/J^2 \to \Omega_{P/R} \otimes_P S \to \Omega_{S/R} \to 0
$$
是分裂正合列。
% P01-E0837: the second frozen “Denote J ... the kernel” again omits “by/as”.
\end{lemma}

\begin{proof}
假设 $S$ 在 $R$ 上形式光滑。由引理
\ref{algebra-lemma-characterize-formally-smooth}，这意味着存在 $R$-代数映射
$S \to P/J^2$，它是典范映射 $P/J^2 \to S$ 的右逆。由引理
\ref{algebra-lemma-differential-mod-power-ideal}，有
$\Omega_{P/R} \otimes_P S = \Omega_{(P/J^2)/R} \otimes_{P/J^2} S$。
由引理 \ref{algebra-lemma-differential-seq-split}，该序列分裂。

\medskip\noindent
假设引理中的正合列分裂。选取分裂映射
$\sigma : \Omega_{S/R} \to \Omega_{P/R} \otimes_P S$。
对每个 $
\lambda \in S$，选取映到 $
\lambda$ 的 $x_\lambda \in P$。
然后对每个 $
\lambda \in S$，选取 $f_\lambda \in J$，使得
$$
\text{d}f_\lambda = \text{d}x_\lambda - \sigma(\text{d}\lambda)
$$
在正合列的中项中成立。我们断言，映射
$s : \lambda \mapsto x_\lambda - f_\lambda \mod J^2$ 是 $R$-代数同态
$s : S \to P/J^2$。为证明这一点，我们将反复使用如下事实：若
$h \in J$，且 $
\text{d}h = 0$ 在 $
\Omega_{P/R} \otimes_R S$ 中成立，则 $h \in J^2$。
% P01-E0838: the frozen tensor base is R here, while the exact sequence's middle term is tensored over P.
设 $
\lambda, \mu \in S$。则
$\sigma(\text{d}\lambda + \text{d}\mu - \text{d}(\lambda + \mu)) = 0$。
这推出
$$
\text{d}(x_\lambda + x_\mu - x_{\lambda + \mu}
- f_\lambda - f_\mu + f_{\lambda + \mu}) = 0
$$
因而 $x_\lambda + x_\mu - x_{\lambda + \mu}
- f_\lambda - f_\mu + f_{\lambda + \mu} \in J^2$，进而
$s(\lambda) + s(\mu) = s(\lambda + \mu)$。同理，
$\sigma(\lambda \text{d}\mu + \mu \text{d}\lambda - \text{d}\lambda \mu) = 0$，
所以
% P01-E0839: the frozen final differential reads d lambda mu rather than d(lambda mu).
$$
\mu(\text{d}x_\lambda - \text{d}f_\lambda) +
\lambda(\text{d}x_\mu - \text{d}f_\mu) -
\text{d}x_{\lambda\mu} + \text{d}f_{\lambda\mu} = 0
$$
在正合列的中项中成立。此外，在该中项中还有
$$
\text{d}(x_\lambda x_\mu) =
x_\lambda \text{d}x_\mu + x_\mu \text{d}x_\lambda =
\lambda \text{d}x_\mu + \mu \text{d} x_\lambda
$$
合并这些等式可知
$x_\lambda x_\mu - x_{\lambda\mu}
- \mu f_\lambda - \lambda f_\mu + f_{\lambda\mu} \in J^2$，
从而由于 $f_\lambda f_\mu \in J^2$，有
$(x_\lambda - f_\lambda)(x_\mu - f_\mu) -
(x_{\lambda\mu} - f_{\lambda\mu}) \in J^2$；这就是说
$s(\lambda)s(\mu) = s(\lambda\mu)$。若 $
\lambda \in R$，则 $
\text{d}\lambda = 0$，于是 $
\text{d}f_\lambda = \text{d}x_\lambda$，故
$\lambda - x_\lambda + f_\lambda \in J^2$，进而如所需有
$s(\lambda) = \lambda$。现在应用引理
\ref{algebra-lemma-characterize-formally-smooth}，便得 $S/R$ 形式光滑。
\end{proof}

\begin{proposition}
\label{algebra-proposition-characterize-formally-smooth}
设 $R \to S$ 是环映射。考察形式光滑的 $R$-代数 $P$ 以及核为 $J$ 的满射
$P \to S$。下列条件等价：
\begin{enumerate}
\item $S$ 在 $R$ 上形式光滑；
\item 对某个上述 $P \to S$，映射 $P/J^2 \to S$ 有截面；
\item 对每个上述 $P \to S$，映射 $P/J^2 \to S$ 有截面；
\item 对某个上述 $P \to S$，序列
$0 \to J/J^2 \to \Omega_{P/R} \otimes S \to \Omega_{S/R} \to 0$ 分裂正合；
\item 对每个上述 $P \to S$，序列
$0 \to J/J^2 \to \Omega_{P/R} \otimes S \to \Omega_{S/R} \to 0$ 分裂正合；以及
% P01-E0840: the two frozen tensor products omit the base P used everywhere else in the section.
\item 朴素余切复形 $
\NL_{S/R}$ 拟同构于置于 $0$ 次的投射 $S$-模：
这意味着 $H_1(\NL_{S/R}) = 0$，且 $
\Omega_{S/R}$ 是投射 $S$-模。
\end{enumerate}
\end{proposition}

\begin{proof}
显然 (1) 推出 (3)，(3) 推出 (2)，见引理
\ref{algebra-lemma-characterize-formally-smooth} 证明的第一部分。
又有 (3) 推出 (5)，(5) 推出 (4)，且 (2) 推出 (4)，见引理
\ref{algebra-lemma-characterize-formally-smooth-again} 证明的第一部分。
最后，把引理 \ref{algebra-lemma-characterize-formally-smooth-again}
应用于典范满射 $R[S] \to S$（\ref{algebra-equation-canonical-presentation}），
便知 (1) 推出 (6)。

\medskip\noindent
假设 (4)，来证明 (6)。考察与环映射 $R \to P \to S$ 相关的引理
\ref{algebra-lemma-exact-sequence-NL} 中的序列。由上面已证明的蕴含
(1) $\Rightarrow$ (6)，可知 $
\NL_{P/R} \otimes_R S$ 拟同构于置于 $0$ 次的
$\Omega_{P/R} \otimes_P S$。因此
$H_1(\NL_{P/R} \otimes_P S) = 0$。由于 $P \to S$ 满，
由引理 \ref{algebra-lemma-NL-surjection}，$
\NL_{S/P}$ 同伦等价于置于 $1$ 次的 $J/J^2$。于是得到正合列
$0 \to H_1(L_{S/R}) \to J/J^2 \to \Omega_{P/R} \otimes_P S \to
\Omega_{S/R} \to 0$。由假设可知 $H_1(L_{S/R}) = 0$，且
$\Omega_{S/R}$ 是投射 $S$-模。因此 (6) 成立。
% P01-E0841: the first frozen tensor base is R, while the following occurrence is P.
% P01-E0842: the frozen proof switches from NL to unexplained L in both H_1 terms.

\medskip\noindent
最后证明 (6) 推出 (1)。假设意味着复形
$J/J^2 \to \Omega_{P/R} \otimes S$ 拟同构于置于 $0$ 次的投射
$S$-模，其中 $P = R[S]$，且 $P \to S$ 是典范满射
（\ref{algebra-equation-canonical-presentation}）。因此引理
\ref{algebra-lemma-characterize-formally-smooth-again} 表明 $S$ 在 $R$ 上形式光滑。
% P01-E0843: this frozen tensor product again omits the base P.
\end{proof}

\begin{lemma}
\label{algebra-lemma-ses-formally-smooth}
设 $A \to B \to C$ 是环映射，并假设 $B \to C$ 形式光滑。
则引理 \ref{algebra-lemma-exact-sequence-differentials} 中的序列
$$
0 \to \Omega_{B/A} \otimes_B C \to \Omega_{C/A} \to \Omega_{C/B} \to 0
$$
是分裂短正合列。
\end{lemma}

\begin{proof}
这由命题 \ref{algebra-proposition-characterize-formally-smooth}
和引理 \ref{algebra-lemma-exact-sequence-NL} 得到。
\end{proof}

\begin{lemma}
\label{algebra-lemma-differential-seq-formally-smooth}
设 $A \to B \to C$ 是环映射，其中 $A \to C$ 形式光滑，
$B \to C$ 满且核为 $J \subset B$。则引理
\ref{algebra-lemma-differential-seq} 中的正合列
$$
0 \to J/J^2 \to \Omega_{B/A} \otimes_B C \to \Omega_{C/A} \to 0
$$
分裂正合。
\end{lemma}

\begin{proof}
这由命题 \ref{algebra-proposition-characterize-formally-smooth}、
引理 \ref{algebra-lemma-exact-sequence-NL} 和引理
\ref{algebra-lemma-differential-seq} 得到。
\end{proof}

\begin{lemma}
\label{algebra-lemma-application-NL-formally-smooth}
设 $A \to B \to C$ 是环映射。假设 $A \to C$ 满（因而
$B \to C$ 也满），且 $A \to B$ 形式光滑。记
$I = \Ker(A \to C)$ 和 $J = \Ker(B \to C)$。则引理
\ref{algebra-lemma-application-NL} 中的序列
$$
0 \to I/I^2 \to J/J^2 \to \Omega_{B/A} \otimes_B B/J \to 0
$$
分裂正合。
% P01-E0844: the frozen “Denote I = ... and J = ...” omits “by/as”.
\end{lemma}

\begin{proof}
由于 $A \to B$ 形式光滑，存在环映射
$\sigma : B \to A/I^2$，使它与 $A \to B$ 的合成等于商映射
$A \to A/I^2$。于是 $\sigma$ 诱导映射 $J/J^2 \to I/I^2$，
它是映射 $I/I^2 \to J/J^2$ 的逆。
\end{proof}

\begin{lemma}
\label{algebra-lemma-lift-formal-smoothness}
设 $R \to S$ 是环映射，$I \subset R$ 是理想。假设
\begin{enumerate}
\item $I^2 = 0$；
\item $R \to S$ 平坦；并且
\item $R/I \to S/IS$ 形式光滑。
\end{enumerate}
则 $R \to S$ 形式光滑。
\end{lemma}

\begin{proof}
假设 (1)、(2) 和 (3)。取 $R$-代数满射
$P = R[\{x_t\}_{t \in T}] \to S$，核为 $J$。于是
$0 \to J \to P \to S \to 0$ 是平坦 $R$-模的短正合列。
这推出 $I \otimes_R S = IS$、$I \otimes_R P = IP$ 和
$I \otimes_R J = IJ$，以及 $J \cap IP = IJ$。在整个证明中我们将使用
$$
\Omega_{(S/IS)/(R/I)} = \Omega_{S/R} \otimes_S (S/IS)
= \Omega_{S/R} \otimes_R R/I = \Omega_{S/R} / I\Omega_{S/R}
$$
对 $P$ 也有类似结论（见引理 \ref{algebra-lemma-differentials-base-change}）。
由引理 \ref{algebra-lemma-characterize-formally-smooth-again}，序列
\begin{equation}
\label{algebra-equation-split}
0 \to J/(IJ + J^2) \to
\Omega_{P/R} \otimes_P S/IS \to
\Omega_{S/R} \otimes_S S/IS \to 0
\end{equation}
分裂正合。当然，中项是
$\bigoplus_{t \in T} S/IS \text{d}x_t$。选取分裂映射
$\sigma : \Omega_{P/R} \otimes_P S/IS \to J/(IJ + J^2)$。
对每个 $t \in T$，选取 $f_t \in J$，使它在 $J/(IJ + J^2)$
中的像为 $\sigma(\text{d}x_t)$。这唯一确定了 $S$-模映射
$$
\tilde \sigma : \Omega_{P/R} \otimes_R S
= \bigoplus S\text{d}x_t \longrightarrow J/J^2
$$
并满足 $\tilde\sigma(\text{d}x_t) = f_t$。
% P01-E0845: the frozen tensor base is R although the displayed direct-sum identification requires tensoring over P.
由于 $\sigma$ 是 $\text{d}$ 的截面，差
$$
\Delta = \text{id}_{J/J^2} - \tilde \sigma \circ \text{d}
$$
是自映射 $J/J^2 \to J/J^2$，其像包含在
$(IJ + J^2)/J^2$ 中。特别地，因为 $I^2 = 0$，有
$\Delta((IJ + J^2)/J^2) = 0$。这意味着 $\Delta$ 分解为
$$
J/J^2 \to J/(IJ + J^2) \xrightarrow{\overline{\Delta}}
(IJ + J^2)/J^2 \to J/J^2
$$
其中 $\overline{\Delta}$ 是 $S/IS$-模映射。再次利用序列
（\ref{algebra-equation-split}）分裂，可以找到 $S/IS$-模映射
$\overline{\delta} : \Omega_{P/R} \otimes_P S/IS \to (IJ + J^2)/J^2$，
使得 $\overline{\delta} \circ d$ 等于 $\overline{\Delta}$。
与上面相同，映射 $\overline{\delta}$ 确定一个 $S$-模映射
$\delta : \Omega_{P/R} \otimes_P S \to J/J^2$。
% P01-E0846: the frozen source uses bare d here, unlike the section's \text{d} notation.
把 $\tilde \sigma$ 替换为 $\tilde \sigma + \delta$ 后，简单计算表明
$\Delta = 0$。换言之，$\tilde \sigma$ 是
$J/J^2 \to \Omega_{P/R} \otimes_P S$ 的截面。
由引理 \ref{algebra-lemma-characterize-formally-smooth-again}，
可知 $R \to S$ 形式光滑。
\end{proof}

\begin{proposition}
\label{algebra-proposition-smooth-formally-smooth}
设 $R \to S$ 是环映射。下列条件等价：
% P01-E0847: the frozen proposition omits the colon before its enumeration.
\begin{enumerate}
\item $R \to S$ 是有限表现且形式光滑的；
\item $R \to S$ 光滑。
\end{enumerate}
\end{proposition}

\begin{proof}
这由命题 \ref{algebra-proposition-characterize-formally-smooth}
和定义 \ref{algebra-definition-smooth} 得到。
（注意，若 $R \to S$ 是有限表现的，则 $\Omega_{S/R}$ 是有限表现的
$S$-模，见引理 \ref{algebra-lemma-differentials-finitely-presented}。）
\end{proof}

\begin{lemma}
\label{algebra-lemma-finite-presentation-fs-Noetherian}
设 $R \to S$ 是光滑环映射。则存在在 $\mathbf{Z}$ 上有限型的子环
$R_0 \subset R$ 和光滑环映射 $R_0 \to S_0$，使得
$S \cong R \otimes_{R_0} S_0$。
\end{lemma}

\begin{proof}
我们将使用光滑等价于有限表现且形式光滑这一事实，见命题
\ref{algebra-proposition-smooth-formally-smooth}。写成
$S = R[x_1, \ldots, x_n]/(f_1, \ldots, f_m)$，并记
$I = (f_1, \ldots, f_m)$。按引理
\ref{algebra-lemma-characterize-formally-smooth}，选取投影到 $S$ 的右逆
$\sigma : S \to R[x_1, \ldots, x_n]/I^2$。选取
$h_i \in R[x_1, \ldots, x_n]$，使
$\sigma(x_i \bmod I) = h_i \bmod I^2$。由于 $x_i - h_i \in I$，
存在 $b_{ij} \in R[x_1, \ldots, x_n]$，使得
$$
x_i - h_i = \sum\nolimits_j b_{ij} f_j
$$
$\sigma$ 是 $R$-代数同态
$R[x_1, \ldots, x_n]/I \to R[x_1, \ldots, x_n]/I^2$
这一事实等价于如下条件：对某些
$a_{kl} \in R[x_1, \ldots, x_n]$，有
$$
f_j(h_1, \ldots, h_n) = \sum\nolimits_{j_1 j_2} a_{j_1 j_2} f_{j_1} f_{j_2}
$$
令 $R_0 \subset R$ 为所有多项式 $f_j, h_i, a_{kl}, b_{ij}$ 的系数
在 $\mathbf{Z}$ 上生成的子环。定义：
$$
S_0 = R_0[x_1, \ldots, x_n]/(f_1, \ldots, f_m)
$$
并令：
$$
I_0 = (f_1, \ldots, f_m)
$$
由于第二个陈列等式在
$R_0[x_1, \ldots, x_n]$ 中成立，可以令
$$
\sigma_0 : S_0 \to R_0[x_1, \ldots, x_n]/I_0^2
$$
为由规则 $x_i \mapsto h_i \bmod I_0^2$ 定义的 $R_0$-代数映射。
由于第一个陈列等式在 $R_0[x_1, \ldots, x_n]$ 中成立，
可知 $\sigma_0$ 是投影
$$
R_0[x_1, \ldots, x_n] / I_0^2 \to R_0[x_1, \ldots, x_n] / I_0 = S_0
$$
的右逆。因此由引理 \ref{algebra-lemma-characterize-formally-smooth}，
环 $S_0$ 在 $R_0$ 上形式光滑。
\end{proof}

\begin{lemma}
\label{algebra-lemma-smooth-descends-through-colimit}
设 $A = \colim A_i$ 是环的滤过余极限，$A \to B$ 是光滑环映射。
则存在某个 $i$ 和光滑环映射 $A_i \to B_i$，使得
$B = B_i \otimes_{A_i} A$。
\end{lemma}

\begin{proof}
这由引理 \ref{algebra-lemma-finite-presentation-fs-Noetherian} 得到，
因为由引理 \ref{algebra-lemma-characterize-finite-presentation}，
对某个 $i$，$R_0 \to A$ 将分解经过 $A_i$。
\end{proof}

\begin{lemma}
\label{algebra-lemma-descent-formally-smooth}
设 $R \to S$ 是环映射，$R \to R'$ 是忠实平坦环映射。置
$S' = S \otimes_R R'$。则 $R \to S$ 形式光滑，当且仅当
$R' \to S'$ 形式光滑。
\end{lemma}

\begin{proof}
若 $R \to S$ 形式光滑，则由引理 \ref{algebra-lemma-base-change-fs}，
$R' \to S'$ 形式光滑。为证逆命题，假设 $R' \to S'$ 形式光滑。
注意，对任意 $S$-模 $N$，有
$N \otimes_R R' = N \otimes_S S'$。特别地，$S \to S'$ 也忠实平坦。
选取多项式环 $P = R[\{x_i\}_{i \in I}]$ 以及核为 $J$ 的
$R$-代数满射 $P \to S$。注意，$P' = P \otimes_R R'$ 是 $R'$ 上的
多项式代数。由于 $R \to R'$ 平坦，满射 $P' \to S'$ 的核
$J'$ 是 $J \otimes_R R'$。因此，分裂正合列（见引理
\ref{algebra-lemma-characterize-formally-smooth-again}）
$$
0 \to J'/(J')^2 \to \Omega_{P'/R'} \otimes_{P'} S' \to \Omega_{S'/R'} \to 0
$$
是相应序列
$$
J/J^2 \to \Omega_{P/R} \otimes_P S \to \Omega_{S/R} \to 0
$$
经 $S \to S'$ 的基变换，见引理
\ref{algebra-lemma-differential-seq}。
由于 $S \to S'$ 忠实平坦，得到两件事：
(1) 这个不带 ${}'$ 的序列也正合；(2) $\Omega_{S/R}$ 是投射 $S$-模。
事实上，$\Omega_{S'/R'}$ 作为自由模
$\Omega_{P'/R'} \otimes_{P'} S'$ 的直和因子是投射的，并且由上文有
$\Omega_{S/R} \otimes_S {S'} = \Omega_{S'/R'}$。因此 (2) 由投射性沿忠实平坦环
映射的下降得到，见定理 \ref{algebra-theorem-ffdescent-projectivity}。
所以序列
$0 \to J/J^2 \to \Omega_{P/R} \otimes_P S \to \Omega_{S/R} \to 0$
也正合，再次应用引理
\ref{algebra-lemma-characterize-formally-smooth-again} 即得结论。
\end{proof}

\noindent
事实证明，光滑环映射满足下面这个很强的提升性质。

\begin{lemma}
\label{algebra-lemma-smooth-strong-lift}
设 $R \to S$ 是光滑环映射。给定由实线箭头组成的交换图
$$
\xymatrix{
S \ar[r] \ar@{-->}[rd] & A/I \\
R \ar[r] \ar[u] & A \ar[u]
}
$$
其中 $I \subset A$ 是局部幂零理想，则存在使图交换的虚线箭头。
\end{lemma}

\begin{proof}
由引理 \ref{algebra-lemma-finite-presentation-fs-Noetherian}，可把该图扩张成交换图
$$
\xymatrix{
S_0 \ar[r] & S \ar[r] \ar@{-->}[rd] & A/I \\
R_0 \ar[r] \ar[u] & R \ar[r] \ar[u] & A \ar[u]
}
$$
其中 $R_0 \to S_0$ 光滑，$R_0$ 在 $\mathbf{Z}$ 上有限型，且
$S = S_0 \otimes_{R_0} R$。设 $x_1, \ldots, x_n \in S_0$
是在 $R_0$ 上生成 $S_0$ 的元。设 $a_1, \ldots, a_n$ 是 $A$ 中的元，
它们在 $A/I$ 中的像与 $x_1, \ldots, x_n$ 的像分别相同。
记 $A_0 \subset A$ 为由 $R_0$ 的像和元 $a_1, \ldots, a_n$ 生成的子环。
置 $I_0 = A_0 \cap I$。则 $A_0/I_0 \subset A/I$，且
$S_0 \to A/I$ 映入 $A_0/I_0$。因此，只需在图
% P01-E0848: the frozen “Denote A_0 ... the subring” omits “by/as”.
$$
\xymatrix{
S_0 \ar[r] \ar@{-->}[rd] & A_0/I_0 \\
R_0 \ar[r] \ar[u] & A_0 \ar[u]
}
$$
中找到虚线箭头。按构造，环 $A_0$ 在 $\mathbf{Z}$ 上有限型。
故 $A_0$ 是 Noether 环，从而 $I_0$ 幂零，见引理
\ref{algebra-lemma-Noetherian-power}。设 $I_0^n = 0$。由命题
\ref{algebra-proposition-smooth-formally-smooth}，可以依次把 $R_0$-代数映射
$S_0 \to A_0/I_0$ 提升为
$S_0 \to A_0/I_0^2$、$S_0 \to A_0/I_0^3$、$\ldots$，
最终提升为 $S_0 \to A_0/I_0^n = A_0$。
\end{proof}

\section{光滑性与微分}
\label{algebra-section-smooth-differential}

\noindent
下面给出一些关于微分与光滑环映射的结果。

\begin{lemma}
\label{algebra-lemma-triangle-differentials-smooth}
设有环映射 $A \to B \to C$，且 $B \to C$ 光滑，则引理
\ref{algebra-lemma-exact-sequence-differentials} 中的序列
$$
0 \to C \otimes_B \Omega_{B/A} \to \Omega_{C/A} \to \Omega_{C/B} \to 0
$$
正合。
% P01-E0849: the frozen English combines “Given ...” with an unnecessary “then”.
\end{lemma}

\begin{proof}
这由更一般的引理 \ref{algebra-lemma-ses-formally-smooth} 得到，
因为光滑环映射形式光滑，见命题
\ref{algebra-proposition-smooth-formally-smooth}。它也可以直接由引理
\ref{algebra-lemma-exact-sequence-NL} 得到，因为
$H_1(L_{C/B}) = 0$ 是 $B \to C$ 光滑定义的一部分。
% P01-E0850: the frozen proof switches from the section's NL notation to unexplained L.
\end{proof}

\begin{lemma}
\label{algebra-lemma-differential-seq-smooth}
设 $A \to B \to C$ 是环映射，其中 $A \to C$ 光滑，
$B \to C$ 满且核为 $J \subset B$。则引理
\ref{algebra-lemma-differential-seq} 中的正合列
$$
0 \to J/J^2 \to \Omega_{B/A} \otimes_B C \to \Omega_{C/A} \to 0
$$
分裂正合。
\end{lemma}

\begin{proof}
这由更一般的引理
\ref{algebra-lemma-differential-seq-formally-smooth} 得到，
因为光滑环映射形式光滑，见命题
\ref{algebra-proposition-smooth-formally-smooth}。
\end{proof}

\begin{lemma}
\label{algebra-lemma-application-NL-smooth}
设 $A \to B \to C$ 是环映射。假设 $A \to C$ 满（因而
$B \to C$ 也满），且 $A \to B$ 光滑。记
$I = \Ker(A \to C)$ 和 $J = \Ker(B \to C)$。则引理
\ref{algebra-lemma-application-NL} 中的序列
$$
0 \to I/I^2 \to J/J^2 \to \Omega_{B/A} \otimes_B B/J \to 0
$$
正合。
% P01-E0851: the frozen “Denote I = ... and J = ...” omits “by/as”.
\end{lemma}

\begin{proof}
这由更一般的引理
\ref{algebra-lemma-application-NL-formally-smooth} 得到，
因为光滑环映射形式光滑，见命题
\ref{algebra-proposition-smooth-formally-smooth}。
\end{proof}

\begin{lemma}
\label{algebra-lemma-section-smooth}
\begin{slogan}
若 $R$ 是 $S$ 的直和因子，且 $S$ 在 $R$ 上光滑，则 $S$ 的
$I$-进完备化通常是 $R$ 上的幂级数环，其中 $I$ 是从 $S$ 到 $R$
的投影映射的核。
% P01-E0852: the frozen slogan says “a power series over R” rather than a power series ring.
% P01-E0853: the frozen slogan lacks terminal punctuation.
\end{slogan}
设 $\varphi : R \to S$ 是光滑环映射，$
\sigma : S \to R$ 是 $
\varphi$ 的左逆。置 $I = \Ker(\sigma)$。则
\begin{enumerate}
\item $I/I^2$ 是有限局部自由 $R$-模；并且
\item 若 $I/I^2$ 自由，则作为 $R$-代数有
$S^\wedge \cong R[[t_1, \ldots, t_d]]$，其中 $S^\wedge$ 是 $S$ 的
$I$-进完备化。
\end{enumerate}
\end{lemma}

\begin{proof}
把引理 \ref{algebra-lemma-differential-seq-split} 应用于 $R \to S \to R$，
可知 $I/I^2 = \Omega_{S/R} \otimes_{S, \sigma} R$。
按光滑态射的定义，模 $
\Omega_{S/R}$ 在 $S$ 上有限局部自由，故 (1) 成立。
若 $I/I^2$ 自由，则选取 $f_1, \ldots, f_d \in I$，使它们在
$I/I^2$ 中的像构成一个 $R$-基。考察由下式定义的 $R$-代数映射：
$$
\Psi : R[[x_1, \ldots, x_d]] \longrightarrow S^\wedge, \quad
x_i \longmapsto f_i.
$$
记 $P = R[[x_1, \ldots, x_d]]$ 和 $J = (x_1, \ldots, x_d) \subset P$。
以 $
\Psi_n : P/J^n \to S/I^n$ 表示诱导的商环映射。注意
$S/I^2 = \varphi(R) \oplus I/I^2$，故 $
\Psi_2$ 是同构。记
$\sigma_2 : S/I^2 \to P/J^2$ 为 $
\Psi_2$ 的逆。
% P01-E0854: the frozen “Denote P = ... and J = ...” omits “by/as”.
% P01-E0855: the frozen “Denote sigma_2 ... the inverse” omits “by/as”.
我们将对 $n$ 归纳证明：对每个 $n > 2$，$
\Psi_n$ 都有逆映射
$\sigma_n : S/I^n \to P/J^n$。事实上，由命题
\ref{algebra-proposition-smooth-formally-smooth}，$S$ 在 $R$ 上形式光滑，
所以在 $R$-代数的实线图
$$
\xymatrix{
S \ar@{..>}[r] \ar[rd]_{\sigma_{n - 1}} & P/J^n \ar[d] \\
 & P/J^{n - 1}
}
$$
中，可以填入某个 $R$-代数映射
$\tau : S \to P/J^n$ 作为虚线箭头，使该图交换。它诱导一个
$R$-代数映射 $
\overline{\tau} : S/I^n \to P/J^n$，模 $J^{n - 1}$ 后等于
$\sigma_{n - 1}$。按构造，映射 $
\Psi_n$ 是满射，而现在
$\overline{\tau} \circ \Psi_n$ 是 $P/J^n$ 的 $R$-代数自同态，
它把 $x_i$ 映到 $x_i + \delta_{i, n}$，其中
$\delta_{i, n} \in J^{n - 1}/J^n$。由此可知 $
\Psi_n$ 是同构，
因而它有逆 $
\sigma_n$。这证明了引理。
\end{proof}

\section{域上的光滑代数}
\label{algebra-section-smooth-over-field}

\noindent
警告：一般而言，下面两个引理在非完美域上并不成立。

\begin{lemma}
\label{algebra-lemma-rank-omega}
设 $k$ 是代数闭域，$S$ 是有限型 $k$-代数，
$\mathfrak m \subset S$ 是极大理想。则
$$
\dim_{\kappa(\mathfrak m)} \Omega_{S/k} \otimes_S \kappa(\mathfrak m)
=
\dim_{\kappa(\mathfrak m)} \mathfrak m/\mathfrak m^2.
$$
\end{lemma}

\begin{proof}
考察引理 \ref{algebra-lemma-differential-seq} 中的正合列
$$
\mathfrak m/\mathfrak m^2 \to
\Omega_{S/k} \otimes_S \kappa(\mathfrak m) \to
\Omega_{\kappa(\mathfrak m)/k} \to 0
$$
我们要证明第一个映射是同构。由于 $k$ 代数闭，由定理
\ref{algebra-theorem-nullstellensatz}，合成 $k \to \kappa(\mathfrak m)$
是同构。因此满射 $S \to \kappa(\mathfrak m)$ 作为 $k$-代数映射分裂，
而引理 \ref{algebra-lemma-differential-seq-split} 表明上述序列在左端也正合。
由于 $\Omega_{\kappa(\mathfrak m)/k} = 0$，结论成立。
\end{proof}

\begin{lemma}
\label{algebra-lemma-characterize-smooth-kbar}
设 $k$ 是代数闭域，$S$ 是有限型 $k$-代数，
$\mathfrak m \subset S$ 是极大理想。下列条件等价：
\begin{enumerate}
\item 环 $S_{\mathfrak m}$ 是正则局部环。
\item 有
$\dim_{\kappa(\mathfrak m)} \Omega_{S/k} \otimes_S \kappa(\mathfrak m)
\leq \dim(S_{\mathfrak m})$。
\item 有
$\dim_{\kappa(\mathfrak m)} \Omega_{S/k} \otimes_S \kappa(\mathfrak m)
= \dim(S_{\mathfrak m})$。
\item 存在 $g \in S$、$g \not \in \mathfrak m$，使得 $S_g$ 在 $k$ 上光滑。
换言之，$S/k$ 在 $\mathfrak m$ 处光滑。
\end{enumerate}
\end{lemma}

\begin{proof}
由引理 \ref{algebra-lemma-rank-omega} 和定义 \ref{algebra-definition-regular}，
(1)、(2) 和 (3) 等价。

\medskip\noindent
假设 $S$ 在 $\mathfrak m$ 处光滑。由引理
\ref{algebra-lemma-smooth-syntomic}，对某个适当的
$g \in S$、$g \not \in \mathfrak m$，$S_g$ 在 $k$ 上标准光滑。
因此由引理 \ref{algebra-lemma-standard-smooth}，
$\Omega_{S_g/k}$ 自由且秩为 $\dim(S_g)$。再由引理
\ref{algebra-lemma-rank-omega}，有
$\dim(S_{\mathfrak m}) = \dim (\mathfrak m/\mathfrak m^2)$，
换言之，$S_\mathfrak m$ 正则。

\medskip\noindent
反之，假设 $S_{\mathfrak m}$ 正则。令
$d = \dim(S_{\mathfrak m}) = \dim \mathfrak m/\mathfrak m^2$。
选取表现 $S = k[x_1, \ldots, x_n]/I$，使对所有 $i$，$x_i$ 都映到
$\mathfrak m$ 的一个元。换言之，
$\mathfrak m'' = (x_1, \ldots, x_n)$ 是
$k[x_1, \ldots, x_n]$ 中相应的极大理想。注意，有序列
$$
I/\mathfrak m''I \to \mathfrak m''/(\mathfrak m'')^2
\to \mathfrak m/(\mathfrak m)^2 \to 0
$$
% P01-E0856: the frozen prose calls this a short exact sequence although the display has no initial 0 and the first map need not be injective.
选取 $c = n - d$ 个元 $f_1, \ldots, f_c \in I$，使它们在
$\mathfrak m''/(\mathfrak m'')^2$ 中的像张成到
$\mathfrak m/\mathfrak m^2$ 的映射的核。这显然可以做到。记
$J = (f_1, \ldots, f_c)$，于是 $J \subset I$。置
$S' = k[x_1, \ldots, x_n]/J$，从而有满射 $S' \to S$。记
$\mathfrak m' = \mathfrak m''S'$ 为 $S'$ 中相应的极大理想。于是有
% P01-E0857: three consecutive frozen “Denote ...” clauses omit “by/as”.
$$
\xymatrix{
k[x_1, \ldots, x_n] \ar[r] & S' \ar[r] & S \\
\mathfrak m'' \ar[u] \ar[r] & \mathfrak m' \ar[r] \ar[u] &
\mathfrak m \ar[u]
}
$$
由 $J$ 的选取，正合列
$$
J/\mathfrak m''J \to \mathfrak m''/(\mathfrak m'')^2
\to \mathfrak m'/(\mathfrak m')^2 \to 0
$$
表明 $\dim( \mathfrak m'/(\mathfrak m')^2 ) = d$。
由于 $S'_{\mathfrak m'}$ 满射到 $S_{\mathfrak m}$，可知
$\dim(S_{\mathfrak m'}) \geq d$。因此由定义
\ref{algebra-definition-regular-local} 前的讨论，$S'_{\mathfrak m'}$
也正则且维数为 $d$。
% P01-E0858: the frozen dimension term localizes S at m', although m' is an ideal of S'; the surrounding argument requires S'_{m'}.
由于 $S'$ 由 $c = n - d$ 个方程切出，由引理
\ref{algebra-lemma-lci}，存在 $g' \in S'$、$g' \not \in \mathfrak m'$，
使 $S'_{g'}$ 是 $k$ 上的全局完全交。映射
$S'_{\mathfrak m'} \to S_{\mathfrak m}$ 又是相同维数的 Noether 局部整环
之间的满射，故为同构。因此 $S' \to S$ 是有限生成核的满射，
且在 $\mathfrak m'$ 处局部化后成为同构。于是可以找到
$g' \in S'$、$g \not \in \mathfrak m'$，使
$S'_{g'} \to S_{g'}$ 是同构。综上，把 $S$ 替换为某个主局部化后，
可以假设 $S$ 是全局完全交。
% P01-E0859: the frozen localization sentence introduces g' but states g notin m', and then localizes S at g' although g' belongs to S'.

\medskip\noindent
此时可写
$S = k[x_1, \ldots, x_n]/(f_1, \ldots, f_c)$，且
$\dim S = n - c$。回顾一下，这个代数的朴素余切复形为
$$
\bigoplus S \cdot f_j
\to
\bigoplus S \cdot \text{d}x_i
$$
见引理 \ref{algebra-lemma-relative-global-complete-intersection-conormal}。
由引理 \ref{algebra-lemma-relative-global-complete-intersection-smooth}，
为证明 $S$ 在 $\mathfrak m$ 处光滑，只需证明给出上述映射的矩阵
``$A$'' 的某个 $c \times c$ 子式 $g_I$ 在 $\mathfrak m$ 处不消失。
由引理 \ref{algebra-lemma-rank-omega}，矩阵 $A \bmod \mathfrak m$ 的秩为 $c$。
因此结论成立。
\end{proof}

\begin{lemma}
\label{algebra-lemma-characterize-smooth-over-field}
设 $k$ 是任意域，$S$ 是有限型 $k$-代数，$X = \Spec(S)$。
设素理想 $\mathfrak q \subset S$ 对应于 $x \in X$。下列条件等价：
\begin{enumerate}
\item $k$-代数 $S$ 在 $\mathfrak q$ 处相对于 $k$ 光滑。
\item 有
$\dim_{\kappa(\mathfrak q)} \Omega_{S/k} \otimes_S \kappa(\mathfrak q)
\leq \dim_x X$。
\item 有
$\dim_{\kappa(\mathfrak q)} \Omega_{S/k} \otimes_S \kappa(\mathfrak q)
= \dim_x X$。
\end{enumerate}
此外，在这些条件成立时，局部环 $S_{\mathfrak q}$ 正则。
\end{lemma}

\begin{proof}
若 $S$ 在 $\mathfrak q$ 处相对于 $k$ 光滑，则由引理
\ref{algebra-lemma-smooth-syntomic}，存在
$g \in S$、$g \not \in \mathfrak q$，使 $S_g$ 在 $k$ 上标准光滑。
$k$ 上的标准光滑代数，其微分模自由且秩等于维数，见引理
\ref{algebra-lemma-standard-smooth}（这里使用了域上的相对全局完全交的维数
等于变量数减去方程数这一事实）。因此 (1) 推出 (3)。
为完成引理的证明，只需证明 (2) 推出 (1)，并推出
$S_{\mathfrak q}$ 正则。

\medskip\noindent
假设 (2)。由 Nakayama 引理 \ref{algebra-lemma-NAK}，
$\Omega_{S/k, \mathfrak q}$ 可由 $\leq \dim_x X$ 个元生成。
可以把 $S$ 替换为某个 $S_g$，其中
$g \in S$、$g \not \in \mathfrak q$，使 $\Omega_{S/k}$
由至多 $\dim_x X$ 个元生成。取代数闭域扩张 $K/k$，使得存在
$k$-代数映射 $\psi : \kappa(\mathfrak q) \to K$。
考察 $S_K = K \otimes_k S$。令 $\mathfrak m \subset S_K$
为下列满射对应的极大理想：
$$
\xymatrix{
S_K = K \otimes_k S \ar[r] &
K \otimes_k \kappa(\mathfrak q)
\ar[r]^-{\text{id}_K \otimes \psi} &
K.
}
$$
注意，$\mathfrak m \cap S = \mathfrak q$；换言之，
$\mathfrak m$ 位于 $\mathfrak q$ 之上。由引理
\ref{algebra-lemma-dimension-at-a-point-preserved-field-extension}，
$X_K = \Spec(S_K)$ 在 $\mathfrak m$ 对应点处的维数为 $\dim_x X$。
由引理 \ref{algebra-lemma-dimension-closed-point-finite-type-field}，
它等于 $\dim((S_K)_{\mathfrak m})$。由引理
\ref{algebra-lemma-differentials-base-change}，$S_K$ 在 $K$ 上的微分模
是 $\Omega_{S/k}$ 的基变换，因此也可由至多
$\dim_x X = \dim((S_K)_{\mathfrak m})$ 个元生成。由引理
\ref{algebra-lemma-characterize-smooth-kbar}，$S_K$ 在 $\mathfrak m$ 处相对于
$K$ 光滑。由引理 \ref{algebra-lemma-flat-base-change-locus-smooth}，
这推出 $S$ 在 $\mathfrak q$ 处相对于 $k$ 光滑，从而证明了 (1)。
此外，由引理 \ref{algebra-lemma-characterize-smooth-kbar}，局部环
$(S_K)_{\mathfrak m}$ 正则。由于
$S_{\mathfrak q} \to (S_K)_{\mathfrak m}$ 平坦，由引理
\ref{algebra-lemma-flat-under-regular} 可知 $S_{\mathfrak q}$ 正则。
\end{proof}

\noindent
下面的引理可以显著推广（而且有几种不同的推广方式）。

\begin{lemma}
\label{algebra-lemma-computation-differential}
设 $k$ 是域，$R$ 是包含 $k$ 的 Noether 局部环。假设剩余域
$\kappa = R/\mathfrak m$ 是 $k$ 的有限生成可分扩张。则映射
$$
\text{d} :
\mathfrak m/\mathfrak m^2
\longrightarrow
\Omega_{R/k} \otimes_R \kappa(\mathfrak m)
$$
是单射。
\end{lemma}

\begin{proof}
可以把 $R$ 替换为 $R/\mathfrak m^2$，故可假设
$\mathfrak m^2 = 0$。由假设可写
$\kappa = k(\overline{x}_1, \ldots, \overline{x}_r, \overline{y})$，
其中 $\overline{x}_1, \ldots, \overline{x}_r$ 是 $\kappa$ 在 $k$ 上的
超越基，而 $\overline{y}$ 在
$k(\overline{x}_1, \ldots, \overline{x}_r)$ 上可分代数。令
$P(T)$ in $k(\overline{x}_1, \ldots, \overline{x}_r)[T]$ 为
$\overline{y}$ 在 $k(\overline{x}_1, \ldots, \overline{x}_r)$ 上的
最小多项式。于是 $P(\overline{y}) = 0$，但
$P'(\overline{y}) \not = 0$。任取 $x_i \in R$，提升
$\overline{x}_i \in \kappa$。这给出 $k$-代数的交换图
$$
\xymatrix{
R \ar[r] & \kappa \\
& k(\overline{x}_1, \ldots, \overline{x}_r) \ar[lu]^\varphi \ar[u]
}
$$
我们要把向左上的箭头 $\varphi$ 扩张为从 $\kappa$ 到 $R$ 的
$k$-代数映射。为此，任取提升 $\overline{y}$ 的 $y \in R$。
为说明它定义了以
$\kappa \cong k(\overline{x}_1, \ldots, \overline{x}_r)[T]/(P)$
为定义域的 $k$-代数映射，只需证明可以选取 $y$，使
$P^\varphi(y) = 0$。否则，对 $\delta \in \mathfrak m$ 计算得到
$$
P^\varphi(y + \delta) = P^\varphi(y) + (P')^\varphi(y)\delta
$$
因为 $\mathfrak m^2 = 0$。由于 $(P')^\varphi(y)$ 是 $R$ 的单位，
可以按所需调整我们的选择。这表明作为 $k$-代数有
$R \cong \kappa \oplus \mathfrak m$！现在，直接计算
$\Omega_{\kappa \oplus \mathfrak m/k}$，或者应用引理
\ref{algebra-lemma-differential-seq-split}，都可完成证明。
\end{proof}

\begin{lemma}
\label{algebra-lemma-separable-smooth}
设 $k$ 是域，$S$ 是有限型 $k$-代数，$\mathfrak q \subset S$ 是素理想。
假设 $\kappa(\mathfrak q)$ 在 $k$ 上可分。下列条件等价：
\begin{enumerate}
\item 代数 $S$ 在 $\mathfrak q$ 处相对于 $k$ 光滑。
\item 环 $S_{\mathfrak q}$ 正则。
\end{enumerate}
\end{lemma}

\begin{proof}
记 $R = S_{\mathfrak q}$，以 $\mathfrak m$ 表示其极大理想，
以 $\kappa$ 表示其剩余域。由引理
\ref{algebra-lemma-computation-differential} 和
\ref{algebra-lemma-differential-seq}，有短正合列
% P01-E0860: the frozen designation omits “ideal” after “its maximal” and uses an incomplete “denote ... by” construction.
% P01-E0861: the frozen citation says “By Lemma X and [bare reference Y]”, omitting “Lemma” before the second reference.
$$
0 \to \mathfrak m/\mathfrak m^2 \to
\Omega_{R/k} \otimes_R \kappa \to
\Omega_{\kappa/k} \to 0
$$
注意，$\Omega_{R/k} = \Omega_{S/k, \mathfrak q}$，见引理
\ref{algebra-lemma-differentials-localize}。此外，由于 $\kappa$ 在 $k$ 上可分，
有 $\dim_{\kappa} \Omega_{\kappa/k} = \text{trdeg}_k(\kappa)$。
因而得到
$$
\dim_{\kappa} \Omega_{R/k} \otimes_R \kappa
=
\dim_\kappa \mathfrak m/\mathfrak m^2 + \text{trdeg}_k (\kappa)
\geq
\dim R + \text{trdeg}_k (\kappa)
=
\dim_{\mathfrak q} S
$$
（最后一个等式见引理
\ref{algebra-lemma-dimension-at-a-point-finite-type-field}），
且等号成立当且仅当 $R$ 正则。因此应用引理
\ref{algebra-lemma-characterize-smooth-over-field} 即得结论。
% P01-E0862: the frozen parenthetical citation lacks punctuation before the following “with equality” clause.
\end{proof}

\begin{lemma}
\label{algebra-lemma-characteristic-zero}
设 $R \to S$ 是 $\mathbf{Q}$-代数映射。设 $f \in S$，并且对某个
$S$-子模 $C$ 有 $\Omega_{S/R} = S \text{d}f \oplus C$。则
\begin{enumerate}
\item $f$ 不是幂零元；并且
\item 若 $S$ 是 Noether 局部环，则 $f$ 是 $S$ 中的非零因子。
\end{enumerate}
\end{lemma}

\begin{proof}
对 $a \in S$，写成
$\text{d}(a) = \theta(a)\text{d}f + c(a)$，其中
$\theta(a) \in S$ 且 $c(a) \in C$。考察 $R$-导子
$S \to S$、$a \mapsto \theta(a)$。注意，$\theta(f) = 1$。

\medskip\noindent
若 $f^n = 0$，且 $n > 1$ 最小，则
$0 = \theta(f^n) = n f^{n - 1}$，这与 $n$ 的最小性矛盾。
所以 $f$ 不是幂零元。

\medskip\noindent
假设 $fa = 0$。若 $f$ 是单位，则 $a = 0$，结论成立。
假设 $f$ 不是单位。由 Leibniz 法则，
$0 = \theta(fa) = f\theta(a) + a$，故 $a \in (f)$。
归纳地，假设已经证明 $fa = 0 \Rightarrow a \in (f^n)$。
写 $a = f^nb$，便有
$0 = \theta(f^{n + 1}b) = (n + 1)f^nb + f^{n + 1}\theta(b)$。
因此
$a = f^n b = -f^{n + 1}\theta(b)/(n + 1) \in (f^{n + 1})$。
在 Noether 局部环 $S$ 中有 $\bigcap (f^n) = 0$，见引理
\ref{algebra-lemma-intersect-powers-ideal-module-zero}，所以结论成立。
\end{proof}

\noindent
下面这个结论在应用中大概没有什么用。

\begin{lemma}
\label{algebra-lemma-characteristic-zero-local-smooth}
设 $k$ 是特征 $0$ 的域，$S$ 是有限型 $k$-代数，
$\mathfrak q \subset S$ 是素理想。下列条件等价：
\begin{enumerate}
\item 代数 $S$ 在 $\mathfrak q$ 处相对于 $k$ 光滑。
\item $S_{\mathfrak q}$-模 $\Omega_{S/k, \mathfrak q}$（有限）自由。
\item 环 $S_{\mathfrak q}$ 正则。
\end{enumerate}
\end{lemma}

\begin{proof}
在特征零时，任意域扩张都可分，因此 (1) 与 (3) 的等价性由引理
\ref{algebra-lemma-separable-smooth} 得到。又由光滑代数的定义，(1) 推出 (2)。
假设 $\Omega_{S/k, \mathfrak q}$ 在 $S_{\mathfrak q}$ 上自由。
我们将使用引理 \ref{algebra-lemma-separable-smooth} 的证明中所用的记号和观察。
因此 $R = S_{\mathfrak q}$，其极大理想为 $\mathfrak m$，
剩余域为 $\kappa$。目标是证明 $R$ 正则。

\medskip\noindent
若 $\mathfrak m/\mathfrak m^2 = 0$，则 $\mathfrak m = 0$，且
$R \cong \kappa$。所以 $R$ 正则，结论成立。

\medskip\noindent
若 $\mathfrak m/ \mathfrak m^2 \not = 0$，则任取
$f \in \mathfrak m$，使其在 $\mathfrak m/ \mathfrak m^2$ 中的像非零。
由引理 \ref{algebra-lemma-computation-differential}，可知
$\text{d}f$ 在 $\Omega_{R/k}/\mathfrak m\Omega_{R/k}$ 中的像非零。
按假设，$\Omega_{R/k} = \Omega_{S/k, \mathfrak q}$ 有限自由，
因而由 Nakayama 引理 \ref{algebra-lemma-NAK}，可知
$\text{d}f$ 生成一个直和因子。应用引理
\ref{algebra-lemma-characteristic-zero}，得到 $f$ 是 $R$ 中的非零因子。
再由引理 \ref{algebra-lemma-differential-seq}，得到正合列
$$
(f)/(f^2) \to \Omega_{R/k} \otimes_R R/fR \to \Omega_{(R/fR)/k} \to 0
$$
这推出 $\Omega_{(R/fR)/k}$ 也有限自由。由归纳可知
$R/fR$ 是正则局部环。由于 $f \in \mathfrak m$ 是非零因子，
由引理 \ref{algebra-lemma-regular-mod-x} 可知 $R$ 正则。
\end{proof}

\begin{example}
\label{algebra-example-characteristic-p}
引理 \ref{algebra-lemma-characteristic-zero-local-smooth} 在特征
$p > 0$ 时不成立。标准例子是环映射
$$
\mathbf{F}_p \longrightarrow \mathbf{F}_p[x]/(x^p)
$$
其微分模自由，但显然不光滑；另一个例子是环映射（$p > 2$）
$$
\mathbf{F}_p(t) \to \mathbf{F}_p(t)[x, y]/(x^p + y^2 + \alpha)
$$
它在素理想 $\mathfrak q = (y, x^p + \alpha)$ 处不光滑，但正则。
% P01-E0863: alpha is not defined or constrained in the frozen characteristic-p example.
\end{example}

\noindent
利用上面的材料，可以用域扩张刻画一般点处的光滑性。

\begin{lemma}
\label{algebra-lemma-smooth-at-generic-point}
设 $R \to S$ 是单射的有限型环映射，且 $R$、$S$ 都是整环。
则 $R \to S$ 在 $\mathfrak q = (0)$ 处光滑，当且仅当其分式域的诱导扩张
$L/K$ 可分。
\end{lemma}

\begin{proof}
假设 $R \to S$ 在 $(0)$ 处光滑。可以把 $S$ 替换为某个
$S_g$，其中 $g \in S$ 非零，并假设 $R \to S$ 光滑。
于是 $K \to S \otimes_R K$ 光滑（引理
\ref{algebra-lemma-base-change-smooth}）。此外，对任意域扩张
$K'/K$，环映射 $K' \to S \otimes_R K'$ 也光滑。因此由引理
\ref{algebra-lemma-characterize-smooth-over-field}，
$S \otimes_R K'$ 是正则环，特别地是约化环。这表明
$S \otimes_R K$ 在 $K$ 上几何约化。因此由引理
\ref{algebra-lemma-geometrically-reduced-permanence}，$L$ 在 $K$ 上几何约化。
所以由引理 \ref{algebra-lemma-characterize-separable-field-extensions}，
$L/K$ 可分。
% P01-E0864: the frozen sentence has the extra article “a” in “is a geometrically reduced over K”.

\medskip\noindent
反之，假设 $L/K$ 可分。由引理
\ref{algebra-lemma-generic-finite-presentation}，可以假设 $R \to S$ 是有限表现的。
由引理 \ref{algebra-lemma-flat-base-change-locus-smooth}，只需证明
$K \to S \otimes_R K$ 在 $(0)$ 处光滑。这由引理
\ref{algebra-lemma-separable-smooth}、域是正则环这一事实以及 $L/K$ 可分的假设得到。
\end{proof}

\section{Noether 情形的光滑环映射}
\label{algebra-section-smooth-Noetherian}

\begin{definition}
\label{algebra-definition-small-extension}
设 $\varphi : B' \to B$ 是环映射。若 $B'$ 和 $B$ 是局部 Artin 环，
$\varphi$ 满，且 $I = \Ker(\varphi)$ 作为 $B'$-模长度为 $1$，
则称 $\varphi$ 是\textit{小扩张}。
\end{definition}

\noindent
显然，这意味着 $I^2 = 0$，并且对某个
$x \in B'$ 有 $I = (x)$ 及 $\mathfrak m' x = 0$，其中
$\mathfrak m' \subset B'$ 是极大理想。

\begin{lemma}
\label{algebra-lemma-smooth-test-artinian}
设 $R \to S$ 是环映射。设 $\mathfrak q$ 是 $S$ 的素理想，
位于 $\mathfrak p \subset R$ 之上。假设 $R$ 是 Noether 环，
且 $R \to S$ 是有限型的。下列条件等价：
% P01-E0865: the frozen premise omits “is” before “of finite type”.
\begin{enumerate}
\item $R \to S$ 在 $\mathfrak q$ 处光滑；
\item 对每个局部 $R$-代数满射
$(B', \mathfrak m') \to (B, \mathfrak m)$，若
$\Ker(B' \to B)$ 平方为零，并且对每个实线部分交换图
$$
\xymatrix{
S \ar[r] \ar@{-->}[rd] & B \\
R \ar[r] \ar[u] & B' \ar[u]
}
$$
满足 $\mathfrak q = S \cap \mathfrak m$，则存在使图交换的虚线箭头；
% P01-E0866: the frozen equality uses an intersection although S -> B need not be injective; it should state inverse image of m equals q.
\item 与 (2) 相同，但 $B' \to B$ 只遍历小扩张；并且
\item 与 (2) 相同，但 $B' \to B$ 只遍历还满足如下条件的小扩张：
$S \to B$ 诱导同构
$\kappa(\mathfrak q) \cong \kappa(\mathfrak m)$。
\end{enumerate}
\end{lemma}

\begin{proof}
假设 (1)。这意味着存在
$g \in S$、$g \not \in \mathfrak q$，使 $R \to S_g$ 光滑。
由命题 \ref{algebra-proposition-smooth-formally-smooth}，
$R \to S_g$ 形式光滑。注意，给定 (2) 中的任意图，
映射 $S \to B$ 自动分解经过 $S_{\mathfrak q}$，因而更分解经过 $S_g$。
$S_g$ 在 $R$ 上形式光滑，给出一个态射 $S_g \to B'$，
使其嵌入以 $S_g$ 为左上角的类似图中。与 $S \to S_g$ 合成，
便得到所需箭头。换言之，我们已经证明 (1) 推出 (2)。

\medskip\noindent
显然 (2) 推出 (3)，(3) 推出 (4)。

\medskip\noindent
假设 (4)。我们来证明 (1)，从而完成引理的证明。选取表现
$S = R[x_1, \ldots, x_n]/(f_1, \ldots, f_m)$。
这是可行的，因为 $S$ 在 Noether 环 $R$ 上有限型，因而有限表现
（见引理
\ref{algebra-lemma-Noetherian-finite-type-is-finite-presentation}）。
置 $I = (f_1, \ldots, f_m)$。考察这个表现的朴素余切复形
$$
\text{d} : I/I^2
\longrightarrow
\bigoplus\nolimits_{j = 1}^n S\text{d}x_j
$$
（见第 \ref{algebra-section-netherlander} 节）。只需证明把这个复形在
$\mathfrak q$ 处局部化后，该映射成为分裂单射，见引理
\ref{algebra-lemma-smooth-at-point}。记
$S' = R[x_1, \ldots, x_n]/I^2$。由引理
\ref{algebra-lemma-differential-mod-power-ideal}，有
% P01-E0867: the frozen “Denote S' = ...” omits “by/as”.
$$
S \otimes_{S'} \Omega_{S'/R} =
S \otimes_{R[x_1, \ldots, x_n]} \Omega_{R[x_1, \ldots, x_n]/R} =
\bigoplus\nolimits_{j = 1}^m S\text{d}x_j.
$$
% P01-E0868: the frozen direct sum runs to m (number of equations), but the differential module of the n-variable polynomial ring has basis dx_1,...,dx_n.
因此映射
$$
\text{d} :
I/I^2
\longrightarrow
S \otimes_{S'} \Omega_{S'/R}
$$
与上述朴素余切复形中的映射相同。特别地，我们要证明的断言是否成立，
只依赖三个环 $R \to S' \to S$。设
$\mathfrak q' \subset R[x_1, \ldots, x_n]$ 是对应于 $\mathfrak q$ 的素理想。
由于取微分模与局部化交换（引理
\ref{algebra-lemma-differentials-localize}），只需证明来自
$R \to S'_{\mathfrak q'} \to S_{\mathfrak q}$ 的映射
\begin{equation}
\label{algebra-equation-target-map}
\text{d} :
I_{\mathfrak q'}/I_{\mathfrak q'}^2
\longrightarrow
S_{\mathfrak q} \otimes_{S'_{\mathfrak q'}} \Omega_{S'_{\mathfrak q'}/R}
\end{equation}
是分裂单射。

\medskip\noindent
设 $N \in \mathbf{N}$ 是整数。考察环
$$
B'_N = S'_{\mathfrak q'} / (\mathfrak q')^N S'_{\mathfrak q'}
= (S'/(\mathfrak q')^N S')_{\mathfrak q'}
$$
及其商 $B_N = B'_N/IB'_N$。注意，
$B_N \cong S_{\mathfrak q}/\mathfrak q^NS_{\mathfrak q}$。
观察到 $B'_N$ 是局部 Artin 环，因为它是一个局部 Noether 环
对其极大理想的幂取商。取 $B'_N \to B_N$ 的核 $I_N$ 的一个
$B'_N$-子模滤过
$$
0 \subset J_{N, 1} \subset J_{N, 2} \subset \ldots \subset J_{N, n(N)} = I_N
$$
使每个相继商 $J_{N, i}/J_{N, i - 1}$ 的长度为 $1$。
（由于 $B'_N$ 是 Artin 环，这样的滤过存在。）这给出一列小扩张
$$
B'_N  \to B'_N/J_{N, 1} \to B'_N/J_{N, 2} \to \ldots \to
B'_N/J_{N, n(N)} = B'_N/I_N
= B_N = S_{\mathfrak q}/\mathfrak q^NS_{\mathfrak q}
$$
从映射 $S \to B_N$ 开始，对这些小扩张依次应用条件 (4)，
可知存在交换图
% P01-E0869: the frozen phrase “we see there exists” omits “that”.
$$
\xymatrix{
S \ar[r] \ar[rd] & B_N  \\
R \ar[r] \ar[u] & B'_N \ar[u]
}
$$
显然，环映射 $S \to B'_N$ 分解为
$S \to S_{\mathfrak q} \to B'_N$，其中
$S_{\mathfrak q} \to B'_N$ 是局部环的局部同态。此外，由于
$B'_N$ 的极大理想的 $N$ 次幂为零，
$S_{\mathfrak q} \to B'_N$ 分解经过
$S_{\mathfrak q}/(\mathfrak q)^NS_{\mathfrak q} = B_N$。
换言之，我们已经证明，对每个 $N \in \mathbf{N}$，
$R$-代数满射 $B'_N \to B_N$ 都有分裂映射。

\medskip\noindent
考察由核为 $I_N$ 的满射 $B'_N \to B_N$ 得到的表现
$$
I_N \to B_N \otimes_{B'_N} \Omega_{B'_N/R} \to \Omega_{B_N/R} \to 0
$$
（见引理 \ref{algebra-lemma-differential-seq}）。由上可知，$R$-代数映射
$B'_N \to B_N$ 有右逆。因此由引理
\ref{algebra-lemma-differential-seq-split}，上述序列分裂正合！于是对每个 $N$，
映射
$$
I_N \longrightarrow B_N \otimes_{B'_N} \Omega_{B'_N/R}
$$
是分裂单射。证明的其余部分只需展开这句话的确切含义。注意
$$
I_N = I_{\mathfrak q'}/
(I_{\mathfrak q'}^2 + (\mathfrak q')^N \cap I_{\mathfrak q'})
$$
由 Artin--Rees 引理（引理 \ref{algebra-lemma-Artin-Rees}），
可找到 $c \geq 0$，使得对所有 $N \geq c$ 都有
$$
S_{\mathfrak q}/\mathfrak q^{N - c}S_{\mathfrak q}
\otimes_{S_{\mathfrak q}} I_N =
S_{\mathfrak q}/\mathfrak q^{N - c}S_{\mathfrak q}
\otimes_{S_{\mathfrak q}}
I_{\mathfrak q'}/I_{\mathfrak q'}^2
$$
（这些张量积不过是除以
$\mathfrak q^{N - c}$ 的一种华丽写法）。当然可以假设
$c \geq 1$。
% P01-E0870: the frozen parenthetical says “these tensor product are”, with singular “product” after plural “these”.
由引理 \ref{algebra-lemma-differential-mod-power-ideal}，有
$$
S'_{\mathfrak q'}/(\mathfrak q')^{N - c}S'_{\mathfrak q'}
\otimes_{S'_{\mathfrak q'}} \Omega_{B'_N/R} =
S'_{\mathfrak q'}/(\mathfrak q')^{N - c}S'_{\mathfrak q'}
\otimes_{S'_{\mathfrak q'}} \Omega_{S'_{\mathfrak q'}/R}
$$
进一步与 $B_N = S_{\mathfrak q}/\mathfrak q^N$ 作张量积，得到
% P01-E0871: the frozen source continues with lowercase “we” after the display and has no punctuation ending the preceding sentence.
$$
S_{\mathfrak q}/\mathfrak q^{N - c}S_{\mathfrak q}
\otimes_{S'_{\mathfrak q'}} \Omega_{B'_N/R} =
S_{\mathfrak q}/\mathfrak q^{N - c}S_{\mathfrak q}
\otimes_{S'_{\mathfrak q'}} \Omega_{S'_{\mathfrak q'}/R}.
$$
由于分裂单射与任意对象作张量积后仍是分裂单射，可知
$$
S_{\mathfrak q}/\mathfrak q^{N - c}S_{\mathfrak q}
\otimes_{S_{\mathfrak q}}
(\ref{algebra-equation-target-map}) =
S_{\mathfrak q}/\mathfrak q^{N - c}S_{\mathfrak q}
\otimes_{S_{\mathfrak q}/\mathfrak q^N S_{\mathfrak q}}
(I_N \longrightarrow B_N \otimes_{B'_N} \Omega_{B'_N/R})
$$
对所有 $N \geq c$ 都是分裂单射。由引理
\ref{algebra-lemma-split-injection-after-completion}，可知
（\ref{algebra-equation-target-map}）是分裂单射。这完成了证明。
\end{proof}

\section{光滑环映射结果概览}
\label{algebra-section-smooth-overview}

\noindent
下面列出我们在前几节证明的关于光滑环映射的结果。
更精确的陈述和定义请查阅所给参考文献。
\begin{enumerate}
\item 环映射 $R \to S$ 光滑，是指它有限表现，且 $S/R$ 的朴素余切复形
拟同构于置于 $0$ 次的有限投射 $S$-模，见定义
\ref{algebra-definition-smooth}。
\item 若 $S$ 在 $R$ 上光滑，则 $\Omega_{S/R}$ 是有限投射
$S$-模，见定义 \ref{algebra-definition-smooth} 后的讨论。
\item 光滑性在 $S$ 上是局部的，见引理
\ref{algebra-lemma-locally-smooth}。
\item 光滑性在基变换下稳定，见引理
\ref{algebra-lemma-base-change-smooth}。
\item 光滑性在合成下稳定，见引理
\ref{algebra-lemma-compose-smooth}。
\item 光滑环映射是 syntomic 的，特别地平坦，见引理
\ref{algebra-lemma-smooth-syntomic}。
\item 光滑纤维环的有限表现平坦环映射光滑，见引理
\ref{algebra-lemma-flat-fibre-smooth}。
\item 有限表现环映射 $R \to S$ 光滑，当且仅当它形式光滑，见命题
\ref{algebra-proposition-smooth-formally-smooth}。
\item 若 $R \to S$ 是有限型环映射，且 $R$ 是 Noether 环，
则为检验 $R \to S$ 光滑，只需检验沿局部 Artin 环的小扩张的形式光滑
提升性质，见引理 \ref{algebra-lemma-smooth-test-artinian}。
\item 光滑环映射 $R \to S$ 是某个光滑环映射
$R_0 \to S_0$ 的基变换，其中 $R_0$ 在 $\mathbf{Z}$ 上有限型，见引理
\ref{algebra-lemma-finite-presentation-fs-Noetherian}。
\item 环映射的光滑点集的形成与平坦基变换交换，见引理
\ref{algebra-lemma-flat-base-change-locus-smooth}。
\item 若 $S$ 在代数闭域 $k$ 上有限型，且
$\mathfrak m \subset S$ 是极大理想，则下列条件等价：
% P01-E0872: the frozen premise omits “is” before “a maximal ideal”, and the list introduction omits a colon.
\begin{enumerate}
\item $S$ 在 $\mathfrak m$ 的某个邻域中相对于 $k$ 光滑；
\item $S_{\mathfrak m}$ 是正则局部环；
\item $\dim(S_{\mathfrak m}) =
\dim_{\kappa(m)} \Omega_{S/k} \otimes_S \kappa(\mathfrak m)$。
\end{enumerate}
见引理 \ref{algebra-lemma-characterize-smooth-kbar}。
% P01-E0873: the frozen dimension formula uses kappa(m) once although only the ideal mathfrak m is defined and the same formula ends with kappa(mathfrak m).
\item 若 $S$ 在域 $k$ 上有限型，且
$\mathfrak q \subset S$ 是素理想，则下列条件等价：
\begin{enumerate}
\item $S$ 在 $\mathfrak q$ 的某个邻域中相对于 $k$ 光滑；
\item $\dim_{\mathfrak q}(S/k) =
\dim_{\kappa(\mathfrak q)} \Omega_{S/k} \otimes_S \kappa(\mathfrak q)$。
\end{enumerate}
见引理 \ref{algebra-lemma-characterize-smooth-over-field}。
% P01-E0874: the frozen prime-ideal premise omits “is” and the list introduction omits a colon; the same omitted predicate recurs in the next two items.
\item 若 $S$ 在域上光滑，则它的所有局部环都正则，见引理
\ref{algebra-lemma-characterize-smooth-over-field}。
\item 若 $S$ 在域 $k$ 上有限型，$\mathfrak q \subset S$ 是素理想，
域扩张 $\kappa(\mathfrak q)/k$ 可分，且 $S_{\mathfrak q}$ 正则，
则 $S$ 在 $\mathfrak q$ 处相对于 $k$ 光滑，见引理
\ref{algebra-lemma-separable-smooth}。
\item 若 $S$ 在域 $k$ 上有限型，$k$ 的特征为 $0$，
$\mathfrak q \subset S$ 是素理想，且
$\Omega_{S/k, \mathfrak q}$ 自由，则 $S$ 在 $\mathfrak q$ 处相对于
$k$ 光滑，见引理 \ref{algebra-lemma-characteristic-zero-local-smooth}。
\end{enumerate}
其中一些结果是用标准光滑环映射的概念证明的，见定义
\ref{algebra-definition-standard-smooth}。它对应于 syntomic 态射情形中的
相对全局完全交映射，也是构造例子的最简便方式。

\section{平展环映射}
\label{algebra-section-etale}

\noindent
平展环映射是相对维数等于零的光滑环映射。这与下面稍微更直接的定义
等价。

\begin{definition}
\label{algebra-definition-etale}
设 $R \to S$ 是环映射。若它是有限表现的，并且朴素余切复形
$\NL_{S/R}$ 拟同构于零，即 $H_1(\NL_{S/R}) = 0$
且 $\Omega_{S/R} = 0$，则称 $R \to S$ 是{\it 平展的}。给定 $S$ 的素理想
$\mathfrak q$，若存在 $g \in S$，$g \not \in \mathfrak q$，使得
$R \to S_g$ 平展，则称 $R \to S$ {\it 在 $\mathfrak q$ 处平展}。
\end{definition}

\noindent
特别地，若 $S$ 在 $R$ 上平展，则 $\Omega_{S/R} = 0$。
若 $R \to S$ 光滑，则 $R \to S$ 平展当且仅当 $\Omega_{S/R} = 0$。
由关于光滑环映射的结果，我们自动得到大量关于平展映射的结果。
下面在引理 \ref{algebra-lemma-etale} 中汇总这些结果。但在此之前，我们先证明
{\it 任意}平展环映射都是标准光滑的。

\begin{lemma}
\label{algebra-lemma-etale-standard-smooth}
任意平展环映射都是标准光滑的。更确切地说，若 $R \to S$ 平展，
则存在一个表现
$S = R[x_1, \ldots, x_n]/(f_1, \ldots, f_n)$，使得
$\det(\partial f_j/\partial x_i)$ 在 $S$ 中的像可逆。
\end{lemma}

\begin{proof}
设 $R \to S$ 平展。选取一个表现 $S = R[x_1, \ldots, x_n]/I$。
由 $R \to S$ 平展可知
$$
\text{d} :
I/I^2
\longrightarrow
\bigoplus\nolimits_{i = 1, \ldots, n} S\text{d}x_i
$$
是同构，特别地，$I/I^2$ 是自由 $S$-模。因此由引理
\ref{algebra-lemma-huber}，可以假设（必要时更换表现）
$I = (f_1, \ldots, f_c)$，并且各类 $f_i \bmod I^2$
构成 $I/I^2$ 的一组基。由上面显示的映射是同构，立刻得到
$c = n$，且 $\det(\partial f_j/\partial x_i)$ 在 $S$ 中可逆。
\end{proof}

\begin{lemma}
\label{algebra-lemma-etale}
关于平展环映射的结果如下。
\begin{enumerate}
\item 对任意环 $R$ 和任意 $f \in R$，环映射 $R \to R_f$ 平展。
\item 平展环映射的合成平展。
\item 平展环映射的基变换平展。
\item 平展性是局部性质：给定环映射 $R \to S$ 和生成单位理想的元素
$g_1, \ldots, g_m \in S$，若对 $j = 1, \ldots, m$，
$R \to S_{g_j}$ 平展，则 $R \to S$ 平展。
\item 设 $R \to S$ 是有限表现环映射，$R \to R'$ 是平坦环映射，
并令 $S' = R' \otimes_R S$。$R' \to S'$ 的平展点集，是
$R \to S$ 的平展点集在 $\Spec(S') \to \Spec(S)$ 下的逆像。
\item 平展环映射是 syntomic 的，特别地平坦。
\item 若 $S$ 在域 $k$ 上有限型，则 $S$ 在 $k$ 上平展，当且仅当
$\Omega_{S/k} = 0$。
\item 任意平展环映射 $R \to S$ 都是某个平展环映射
$R_0 \to S_0$ 的基变换，其中 $R_0$ 在 $\mathbf{Z}$ 上有限型。
\item 设 $A = \colim A_i$ 是环的滤过余极限，并设 $A \to B$
是平展环映射。则对某个 $i$，存在平展环映射 $A_i \to B_i$，
使得 $B \cong A \otimes_{A_i} B_i$。
\item 设 $A$ 是环，$S$ 是 $A$ 的乘法子集，并设
$S^{-1}A \to B'$ 平展。则存在平展环映射 $A \to B$，使得
$B' \cong S^{-1}B$。
\item 设 $A$ 是环，并设 $B = B' \times B''$ 是 $A$-代数的乘积。
则 $B$ 在 $A$ 上平展，当且仅当 $B'$ 和 $B''$ 都在 $A$ 上平展。
\end{enumerate}
\end{lemma}

\begin{proof}
在每一种情形中，我们都使用光滑环映射的相应结果，再加上一个说明
$\Omega_{S/R}$ 为零的小论证。

\medskip\noindent
(1) 的证明。环映射 $R \to R_f$ 光滑，且 $\Omega_{R_f/R} = 0$。

\medskip\noindent
(2) 的证明。光滑映射 $A \to B$ 与 $B \to C$ 的合成
$A \to C$ 光滑，见引理 \ref{algebra-lemma-compose-smooth}。由引理
\ref{algebra-lemma-exact-sequence-differentials}，由于 $\Omega_{C/B}$ 和
$\Omega_{B/A}$ 都为零，可知 $\Omega_{C/A}$ 为零。

\medskip\noindent
(3) 的证明。设 $R \to S$ 平展，而 $R \to R'$ 任意。
则 $R' \to S' = R' \otimes_R S$ 光滑，见引理
\ref{algebra-lemma-base-change-smooth}。由引理
\ref{algebra-lemma-differentials-base-change}，
$\Omega_{S'/R'} = S' \otimes_S \Omega_{S/R}$，故
$\Omega_{S'/R'} = 0$。因此 $R' \to S'$ 平展。

\medskip\noindent
(4) 的证明。假设 (4) 的条件成立。由引理
\ref{algebra-lemma-locally-smooth}，$R \to S$ 光滑。又已知对所有 $i$，
$\Omega_{S_{g_i}/R} = (\Omega_{S/R})_{g_i} = 0$。
于是 $\Omega_{S/R} = 0$，见引理 \ref{algebra-lemma-cover}。

\medskip\noindent
(5) 的证明。光滑映射的相应结果是引理
\ref{algebra-lemma-flat-base-change-locus-smooth}。在该引理的证明中，
我们用到了 $\NL_{S/R} \otimes_S S'$ 与 $\NL_{S'/R'}$
同伦等价。这样只需证明：若 $M$ 是有限表现 $S$-模，则 $S'$
的满足 $(M \otimes_S S')_{\mathfrak q'} = 0$ 的素理想
$\mathfrak q'$ 所成集合，是 $S$ 的满足 $M_{\mathfrak q} = 0$
的素理想 $\mathfrak q$ 所成集合的逆像。这由引理
\ref{algebra-lemma-support-base-change} 得到。

\medskip\noindent
(6) 的证明。直接由光滑环映射的相应结果（引理
\ref{algebra-lemma-smooth-syntomic}）得到。

\medskip\noindent
(7) 的证明。由引理 \ref{algebra-lemma-characterize-smooth-over-field}
和定义得到。

\medskip\noindent
(8) 的证明。引理 \ref{algebra-lemma-finite-presentation-fs-Noetherian}
给出光滑环映射的相应结果。所得光滑环映射 $R_0 \to S_0$
满足引理 \ref{algebra-lemma-relative-dimension-CM} 的条件，因而可以把
$S_0$ 换成它在 $R_0$ 上相对维数为 $0$ 的因子。

\medskip\noindent
(9) 的证明。由 (8) 得到，因为由引理
\ref{algebra-lemma-characterize-finite-presentation}，对某个 $i$，
$R_0 \to A$ 会通过 $A_i$ 分解。

\medskip\noindent
(10) 的证明。由 (9)、(1) 和 (2) 得到，因为 $S^{-1}A$
是 $A$ 的主局部化的滤过余极限。

\medskip\noindent
(11) 的证明。先用引理 \ref{algebra-lemma-product-smooth} 得到光滑性的
相应结果，再用 $\Omega_{B/A}$ 为零当且仅当
$\Omega_{B'/A}$ 和 $\Omega_{B''/A}$ 都为零。
\end{proof}

\noindent
下面更详细地说明在域上平展意味着什么。

\begin{lemma}
\label{algebra-lemma-etale-over-field}
设 $k$ 是域。环映射 $k \to S$ 平展，当且仅当 $S$ 作为
$k$-代数同构于 $k$ 的若干有限可分扩张的有限乘积。
\end{lemma}

\begin{proof}
以下事实将不再另行说明：若 $S = S_1 \times \ldots \times S_n$
是 $k$-代数的有限乘积，则 $S$ 在 $k$ 上平展，当且仅当每个
$S_i$ 都在 $k$ 上平展。见引理 \ref{algebra-lemma-etale} 的 (11)。

\medskip\noindent
若 $k'/k$ 是有限可分域扩张，则可写成
$k' = k(\alpha) \cong k[x]/(f)$，其中 $f$ 是元素 $\alpha$
的极小多项式。由于 $k'$ 在 $k$ 上可分，有
$\gcd(f, f') = 1$。这说明
$\text{d} : k'\cdot f \to k' \cdot \text{d}x$ 是同构，故
$k \to k'$ 平展。因此，若 $S$ 是 $k$ 的若干有限可分扩张的
有限乘积，则 $S$ 在 $k$ 上平展。

\medskip\noindent
反过来，假设 $k \to S$ 平展。则 $S$ 在 $k$ 上光滑，且
$\Omega_{S/k} = 0$。由引理
\ref{algebra-lemma-characterize-smooth-over-field}，对 $S$ 的每个极大理想
$\mathfrak m$，有 $\dim_\mathfrak m \Spec(S) = 0$。因此
$\dim(S) = 0$。由命题 \ref{algebra-proposition-dimension-zero-ring}，
$S$ 是有限多个 Artin 局部环的乘积。由已经使用过的引理
\ref{algebra-lemma-characterize-smooth-over-field}，这些局部环都是域。
所以可以假设 $S = k'$ 是域。由 Hilbert 零点定理（定理
\ref{algebra-theorem-nullstellensatz}），扩张 $k'/k$ 有限。环映射
$k \to k'$ 的光滑性由引理 \ref{algebra-lemma-smooth-at-generic-point}
推出 $k'/k$ 是可分扩张，证明完成。
\end{proof}

\begin{lemma}
\label{algebra-lemma-etale-at-prime}
设 $R \to S$ 是环映射，$\mathfrak q \subset S$ 是位于 $R$
中的 $\mathfrak p$ 之上的素理想。若 $S/R$ 在 $\mathfrak q$
处平展，则
\begin{enumerate}
\item $\mathfrak p S_{\mathfrak q} = \mathfrak qS_{\mathfrak q}$，
它是局部环 $S_{\mathfrak q}$ 的极大理想；
\item 域扩张 $\kappa(\mathfrak q)/\kappa(\mathfrak p)$ 有限可分。
\end{enumerate}
\end{lemma}

\begin{proof}
首先可以把 $S$ 换成某个 $S_g$，其中 $g \in S$ 且
$g \not \in \mathfrak q$，从而假设 $R \to S$ 平展。展开
$S \otimes_R \kappa(\mathfrak p)$ 在 $\kappa(\mathfrak p)$
上平展这一事实，再应用引理 \ref{algebra-lemma-etale-over-field}，
即可得到本引理。
\end{proof}

\begin{lemma}
\label{algebra-lemma-etale-quasi-finite}
平展环映射是拟有限的。
\end{lemma}

\begin{proof}
设 $R \to S$ 是平展环映射。由定义，$R \to S$ 是有限型的。
对任意素理想 $\mathfrak p \subset R$，纤维环
$S \otimes_R \kappa(\mathfrak p)$ 在 $\kappa(\mathfrak p)$ 上平展，
因而是 $\kappa(\mathfrak p)$ 的若干有限可分域扩张的有限乘积；
特别地，它在 $\kappa(\mathfrak p)$ 上有限。
% P01-E0875: the frozen source says “a finite products of fields”.
故由引理 \ref{algebra-lemma-quasi-finite}，$R \to S$ 拟有限。
\end{proof}

\begin{lemma}
\label{algebra-lemma-characterize-etale}
设 $R \to S$ 是环映射。设 $\mathfrak q$ 是 $S$ 的素理想，
它位于 $R$ 的素理想 $\mathfrak p$ 之上。若
\begin{enumerate}
\item $R \to S$ 是有限表现的；
\item $R_{\mathfrak p} \to S_{\mathfrak q}$ 平坦；
\item $\mathfrak p S_{\mathfrak q}$ 是局部环
$S_{\mathfrak q}$ 的极大理想；
\item 域扩张 $\kappa(\mathfrak q)/\kappa(\mathfrak p)$ 有限可分，
\end{enumerate}
则 $R \to S$ 在 $\mathfrak q$ 处平展。
\end{lemma}

\begin{proof}
应用引理 \ref{algebra-lemma-isolated-point-fibre}，可找到
$g \in S$，$g \not \in \mathfrak q$，使得 $\mathfrak q$ 是
$S_g$ 中唯一位于 $\mathfrak p$ 之上的素理想。我们可以并且确实
把 $S$ 换成 $S_g$。于是 $S \otimes_R \kappa(\mathfrak p)$
只有一个素理想，因而是局部环，并且等于
$S_{\mathfrak q}/\mathfrak pS_{\mathfrak q}
\cong \kappa(\mathfrak q)$。
由引理 \ref{algebra-lemma-flat-fibre-smooth}，存在
$g \in S$，$g \not \in \mathfrak q$，使得 $R \to S_g$ 光滑。
再次把 $S$ 换成 $S_g$，便可假设 $R \to S$ 光滑。
% P01-E0876: the frozen source has the malformed transition “Replace S by S_g again we may assume”.
由引理 \ref{algebra-lemma-smooth-syntomic}，甚至可以假设
$R \to S$ 标准光滑，比如
$S = R[x_1, \ldots, x_n]/(f_1, \ldots, f_c)$。
由于 $S \otimes_R \kappa(\mathfrak p) = \kappa(\mathfrak q)$
的维数为 $0$，可知 $n = c$，即 $R \to S$ 平展。
\end{proof}

\noindent
下面是一个全新的现象。

\begin{lemma}
\label{algebra-lemma-map-between-etale}
设 $R \to S$ 和 $R \to S'$ 平展。则任意 $R$-代数映射
$S' \to S$ 都平展。
\end{lemma}

\begin{proof}
首先，由引理 \ref{algebra-lemma-compose-finite-type}，$S' \to S$
是有限表现的。设素理想 $\mathfrak q \subset S$ 位于
$\mathfrak q' \subset S'$ 和 $\mathfrak p \subset R$ 之上。
由引理 \ref{algebra-lemma-etale-at-prime}，环映射
$S'_{\mathfrak q'}/\mathfrak p S'_{\mathfrak q'} \to
S_{\mathfrak q}/\mathfrak p S_{\mathfrak q}$
是 $\kappa(\mathfrak p)$ 的有限可分扩张之间的映射。
特别地，它平坦。因此由引理 \ref{algebra-lemma-criterion-flatness-fibre}，
$S'_{\mathfrak q'} \to S_{\mathfrak q}$ 平坦。所以 $S' \to S$
平坦。此外，上述论证还说明 $\mathfrak q'S_{\mathfrak q}$ 是
$S_{\mathfrak q}$ 的极大理想，并且
$S'_{\mathfrak q'} \to S_{\mathfrak q}$ 的剩余域扩张有限可分。
故由引理 \ref{algebra-lemma-characterize-etale}，$S' \to S$ 在
$\mathfrak q$ 处平展。由于平展性是局部的（见引理
\ref{algebra-lemma-etale}），结论成立。
\end{proof}

\begin{lemma}
\label{algebra-lemma-surjective-flat-finitely-presented}
设 $\varphi : R \to S$ 是环映射。若 $R \to S$ 满射、平坦且
有限表现，则存在幂等元 $e \in R$，使得 $S = R_e$。
% P01-E0877: the frozen source says “there exist an idempotent”.
\end{lemma}

\begin{proof}[第一种证明]
设 $I$ 是 $\varphi$ 的核。由于 $\varphi$ 是有限表现的，
由引理 \ref{algebra-lemma-finite-presentation-independent}，$I$ 有限生成。
此外，由于 $S$ 在 $R$ 上平坦，把正合列
$0 \to I \to R \to S \to 0$ 在 $R$ 上与 $S$ 作张量积，得到
$I/I^2 = 0$。现在由引理
\ref{algebra-lemma-ideal-is-squared-union-connected} 即得结论。
\end{proof}

\begin{proof}[第二种证明]
由于 $\Spec(S) \to \Spec(R)$ 是到某个闭子集上的同胚（见引理
\ref{algebra-lemma-spec-closed}），并且是开映射（见命题
\ref{algebra-proposition-fppf-open}），可知其像等于某个幂等元
$e \in R$ 所确定的 $D(e)$（见引理
\ref{algebra-lemma-disjoint-decomposition}）。因此 $R_e \to S$ 在谱上诱导
双射。比如由引理 \ref{algebra-lemma-finite-flat-local} 和
\ref{algebra-lemma-NAK}，该映射在所有局部环上都诱导同构。于是
$R_e \to S$ 也为单射，例如可见引理
\ref{algebra-lemma-characterize-zero-local}。
\end{proof}

\begin{lemma}
\label{algebra-lemma-lift-etale}
\begin{slogan}
平展环映射可以沿环的满射提升
\end{slogan}
设 $R$ 是环，$I \subset R$ 是理想。设
$R/I \to \overline{S}$ 是平展环映射。则存在平展环映射
$R \to S$，使得作为 $R/I$-代数，
$\overline{S} \cong S/IS$。
\end{lemma}

\begin{proof}
由引理 \ref{algebra-lemma-etale-standard-smooth}，可以如定义
\ref{algebra-definition-standard-smooth} 中那样写成
$\overline{S} =
(R/I)[x_1, \ldots, x_n]/(\overline{f}_1, \ldots, \overline{f}_n)$，
其中
$\overline{\Delta} =
\det(\frac{\partial \overline{f}_i}{\partial x_j})_{i, j = 1, \ldots, n}$
在 $\overline{S}$ 中可逆。任取一些提升 $f_i$，
并令
$S = R[x_1, \ldots, x_n, x_{n+1}]/(f_1, \ldots, f_n, x_{n + 1}\Delta - 1)$，
其中
$\Delta = \det(\frac{\partial f_i}{\partial x_j})_{i, j = 1, \ldots, n}$，
如例 \ref{algebra-example-make-standard-smooth} 所示。这就证明了引理。
\end{proof}

\begin{lemma}
\label{algebra-lemma-lift-etale-infinitesimal}
考虑交换图
$$
\xymatrix{
0 \ar[r] &
J \ar[r] &
B' \ar[r] &
B \ar[r] & 0 \\
0 \ar[r] &
I \ar[r] \ar[u] &
A' \ar[r] \ar[u] &
A \ar[r] \ar[u] & 0
}
$$
其各行正合，其中 $B' \to B$ 和 $A' \to A$ 是满环映射，
且其核都是平方为零的理想。若 $A \to B$ 平展，且
$J = I \otimes_A B$，则 $A' \to B'$ 平展。
\end{lemma}

\begin{proof}
由引理 \ref{algebra-lemma-lift-etale}，存在平展环映射 $A' \to C$，
使得 $C/IC = B$。于是 $A' \to C$ 形式光滑（由命题
\ref{algebra-proposition-smooth-formally-smooth}），故得到 $A'$-代数映射
$\varphi : C \to B'$。由于 $A' \to C$ 平坦，有
$I \otimes_A B = I \otimes_A C/IC = IC$。因此，假设
$J = I \otimes_A B$ 蕴含 $\varphi$ 诱导同构 $IC \to J$
和同构 $C/IC \to B'/IB'$，从而 $\varphi$ 是同构。
\end{proof}

\begin{example}
\label{algebra-example-factor-polynomials-etale}
设 $n , m \geq 1$ 是整数。考虑例
\ref{algebra-example-factor-polynomials} 中的环映射
% P01-E0878: the frozen source has an anomalous space before the comma in “n , m”.
\begin{eqnarray*}
R = \mathbf{Z}[a_1, \ldots, a_{n + m}]
& \longrightarrow &
S = \mathbf{Z}[b_1, \ldots, b_n, c_1, \ldots, c_m] \\
a_1 & \longmapsto & b_1 + c_1 \\
a_2 & \longmapsto & b_2 + b_1 c_1 + c_2 \\
\ldots & \ldots & \ldots \\
a_{n + m} & \longmapsto & b_n c_m
\end{eqnarray*}
符号化地写成
$$
S = R[b_1, \ldots, c_m]/(\{a_k(b_i, c_j) - a_k\}_{k = 1, \ldots, n + m})
$$
其中，例如 $a_1(b_i, c_j) = b_1 + c_1$。偏导数矩阵为
$$
\left(
\begin{matrix}
1 & c_1 & \ldots & c_m & 0 & \ldots & \ldots & 0 \\
0 & 1 & c_1 & \ldots & c_m & 0 & \ldots & 0 \\
\ldots & \ldots & \ldots & \ldots & \ldots & \ldots & \ldots & \ldots \\
0 & \ldots & 0 & 1 & c_1 & c_2 & \ldots & c_m \\
1 & b_1 & \ldots & b_{n - 1} & b_n & 0 & \ldots & 0 \\
0 & 1 & b_1 & \ldots & b_{n - 1} & b_n & \ldots & 0 \\
\ldots & \ldots & \ldots & \ldots & \ldots & \ldots & \ldots & \ldots \\
0 & \ldots & \ldots & 0 & 1 & b_1 & \ldots & b_n
\end{matrix}
\right)
$$
这个矩阵的行列式 $\Delta$ 更常称为多项式
$g = x^n + b_1 x^{n - 1} + \ldots + b_n$ 与
$h = x^m + c_1 x^{m - 1} + \ldots + c_m$ 的{\it 结式}，
上面的矩阵称为与 $g, h$ 相伴的 {\it Sylvester 矩阵}。
用公式表示就是 $\Delta = \text{Res}_x(g, h)$。Sylvester 矩阵是
线性映射矩阵的转置：
\begin{eqnarray*}
S[x]_{< m} \oplus S[x]_{< n} & \longrightarrow & S[x]_{< n + m} \\
a \oplus b & \longmapsto & ag + bh
\end{eqnarray*}
设 $\mathfrak q \subset S$ 是任意素理想。由上述讨论，下列条件等价：
\begin{enumerate}
\item $R \to S$ 在 $\mathfrak q$ 处平展；
\item $\Delta = \text{Res}_x(g, h) \not \in \mathfrak q$；
\item 多项式 $g, h$ 在 $\kappa(\mathfrak q)[x]$ 中的像
$\overline{g}, \overline{h} \in \kappa(\mathfrak q)[x]$ 互素。
\end{enumerate}
(2) 与 (3) 等价，是因为 Sylvester 矩阵在
$\text{Mat}(n + m, \kappa(\mathfrak q))$ 中的像有核，
当且仅当多项式 $\overline{g}, \overline{h}$ 有公共因子。
由此可知环映射
$$
R \longrightarrow S[\frac{1}{\Delta}] = S[\frac{1}{\text{Res}_x(g, h)}]
$$
平展。
\end{example}


\begin{lemma}
\label{algebra-lemma-factor-mod-lift-etale}
设 $R$ 是环，$f \in R[x]$ 是首一多项式，$\mathfrak p$ 是
$R$ 的素理想。设
$f \bmod \mathfrak p = \overline{g} \overline{h}$ 是 $f$ 在
$\kappa(\mathfrak p)[x]$ 中的像的一个分解式。若
$\gcd(\overline{g}, \overline{h}) = 1$，则存在
\begin{enumerate}
\item 平展环映射 $R \to R'$；
\item 位于 $\mathfrak p$ 之上的素理想 $\mathfrak p' \subset R'$；
\item $R'[x]$ 中的分解式 $f = g h$，
\end{enumerate}
使得
\begin{enumerate}
\item $\kappa(\mathfrak p) = \kappa(\mathfrak p')$；
\item $\overline{g} = g \bmod \mathfrak p'$，
$\overline{h} = h \bmod \mathfrak p'$；
\item 多项式 $g, h$ 生成 $R'[x]$ 中的单位理想。
\end{enumerate}
\end{lemma}

\begin{proof}
假设
$\overline{g} = \overline{b}_0 x^n + \overline{b}_1 x^{n - 1} + \ldots
+ \overline{b}_n$，且
$\overline{h} = \overline{c}_0 x^m + \overline{c}_1 x^{m - 1} + \ldots
+ \overline{c}_m$，其中
$\overline{b}_0, \overline{c}_0 \in \kappa(\mathfrak p)$ 非零。
在 $R$ 的某个不属于 $\mathfrak p$ 的元素处对 $R$ 局部化后，可以假设
$\overline{b}_0$ 是某个可逆元 $b_0 \in R$ 的像。用
$\overline{g}/b_0$ 替换 $\overline{g}$，用
$b_0\overline{h}$ 替换 $\overline{h}$，便约化到
$\overline{g}$、$\overline{h}$ 都首一的情形（验证从略）。设
$\overline{g} = x^n + \overline{b}_1 x^{n - 1} + \ldots + \overline{b}_n$，
且
$\overline{h} = x^m + \overline{c}_1 x^{m - 1} + \ldots + \overline{c}_m$。
写成
$f = x^{n + m} + a_1 x^{n + m - 1} + \ldots + a_{n + m}$。
考虑纤维积
$$
R' = R \otimes_{\mathbf{Z}[a_1, \ldots, a_{n + m}]}
\mathbf{Z}[b_1, \ldots, b_n, c_1, \ldots, c_m]
$$
其中映射 $\mathbf{Z}[a_k] \to \mathbf{Z}[b_i, c_j]$
如例 \ref{algebra-example-factor-polynomials} 和
\ref{algebra-example-factor-polynomials-etale} 所示。由构造，存在
$R$-代数映射
$$
R' = R \otimes_{\mathbf{Z}[a_1, \ldots, a_{n + m}]}
\mathbf{Z}[b_1, \ldots, b_n, c_1, \ldots, c_m]
\longrightarrow
\kappa(\mathfrak p)
$$
它把 $b_i$ 映到 $\overline{b}_i$，把 $c_j$ 映到
$\overline{c}_j$。记这个映射的核为 $\mathfrak p' \subset R'$。
% P01-E0879: the frozen source says “Denote p' the kernel” and omits “by/as”.
由于按假设多项式 $\overline{g}, \overline{h}$ 互素，可知元素
$\Delta = \text{Res}_x(g, h) \in \mathbf{Z}[b_i, c_j]$
（见例 \ref{algebra-example-factor-polynomials-etale}）在上述映射下
在 $\kappa(\mathfrak p)$ 中的像不为零。因此 $R \to R'$
在 $\mathfrak p'$ 处平展。事实上，本引理所提问题的一个解是
环映射 $R \to R'[1/\Delta]$ 和素理想
$\mathfrak p' R'[1/\Delta]$。由于在此环中
$\text{Res}_x(f, g)$ 可逆，Sylvester 矩阵在 $R'[1/\Delta]$ 上可逆，因而对某些
$a, b \in R'[1/\Delta][x]$ 有 $1 = a g +  b h$，见例
\ref{algebra-example-factor-polynomials-etale}。
% P01-E0880: the frozen source invokes Res_x(f,g), but Delta and the Sylvester matrix are for Res_x(g,h); since f=gh, Res_x(f,g) cannot support the stated invertibility.
\end{proof}

\section{平展环映射的局部结构}
\label{algebra-section-etale-local-structure}

\noindent
引理 \ref{algebra-lemma-etale-standard-smooth} 告诉我们，把标准平展态射定义为
相对维数为 $0$ 的标准光滑态射其实没有多大意义。作为平展态射的模型，
我们取域的有限可分扩张 $k'/k$ 给出的例子。也就是说，总能找到元素
$\alpha \in k'$，使得 $k' = k(\alpha)$，并且 $\alpha$ 的极小多项式
$f(x) \in k[x]$ 的导数 $f'$ 与 $f$ 互素。

\begin{definition}
\label{algebra-definition-standard-etale}
设 $R$ 是环。设 $g , f  \in R[x]$。假设 $f$ 首一，并且导数
$f'$ 在局部化 $R[x]_g/(f)$ 中可逆。此时称环映射
$R \to R[x]_g/(f)$ 是{\it 标准平展的}。
% P01-E0881: the frozen source has anomalous spaces in the mathematical list “g , f  in R[x]”.
\end{definition}

\noindent
在命题 \ref{algebra-proposition-etale-locally-standard} 中，我们将证明每个平展
环映射都局部标准平展。

\begin{lemma}
\label{algebra-lemma-standard-etale}
设 $R \to R[x]_g/(f)$ 标准平展。
\begin{enumerate}
\item 环映射 $R \to R[x]_g/(f)$ 平展。
\item 对任意环映射 $R \to R'$，标准平展环映射
$R \to R[x]_g/(f)$ 的基变换 $R' \to R'[x]_g/(f)$ 标准平展。
\item $R[x]_g/(f)$ 的任意主局部化在 $R$ 上标准平展。
\item 一般而言，标准平展映射的合成{\bf 不}标准平展。
\end{enumerate}
\end{lemma}

\begin{proof}
从略。下面给出 (4) 的一个例子。环映射
$\mathbf{F}_2 \to \mathbf{F}_{2^2}$ 标准平展。环映射
$\mathbf{F}_{2^2} \to \mathbf{F}_{2^2} \times \mathbf{F}_{2^2}
\times \mathbf{F}_{2^2} \times \mathbf{F}_{2^2}$ 标准平展。
但是环映射
$\mathbf{F}_2 \to \mathbf{F}_{2^2} \times \mathbf{F}_{2^2}
\times \mathbf{F}_{2^2} \times \mathbf{F}_{2^2}$ 不标准平展。
\end{proof}

\noindent
标准平展态射是构造平展映射的一种方便方法。下面给出一个例子。

\begin{lemma}
\label{algebra-lemma-make-etale-map-prescribed-residue-field}
设 $R$ 是环，$\mathfrak p$ 是 $R$ 的素理想，并设
$L/\kappa(\mathfrak p)$ 是有限可分域扩张。则存在平展环映射
$R \to R'$ 以及位于 $\mathfrak p$ 之上的素理想 $\mathfrak p'$，
使得域扩张 $\kappa(\mathfrak p')/\kappa(\mathfrak p)$ 同构于
$\kappa(\mathfrak p) \subset L$。
\end{lemma}

\begin{proof}
由本原元定理，可以写成 $L = \kappa(\mathfrak p)[\alpha]$。
令 $\overline{f} \in \kappa(\mathfrak p)[x]$ 表示 $\alpha$
的极小多项式（特别地，它首一）。用某个 $c \in R$，
$c\not \in \mathfrak p$ 的 $c\alpha$ 替换 $\alpha$ 后，可以假设
$\overline{f}$ 的所有系数都在 $R \to \kappa(\mathfrak p)$ 的像中
（验证从略）。所以可以找到首一多项式 $f \in R[x]$，它在
$\kappa(\mathfrak p)[x]$ 中映到 $\overline{f}$。由于
$\kappa(\mathfrak p) \subset L$ 可分，可知
$\gcd(\overline{f}, \overline{f}') = 1$。因此存在
$\gamma \in L$，使得 $\overline{f}'(\alpha) \gamma = 1$。
于是得到 $R$-代数映射
\begin{eqnarray*}
R[x, 1/f']/(f) & \longrightarrow & L \\
x & \longmapsto & \alpha \\
1/f' & \longmapsto & \gamma
\end{eqnarray*}
左端是 $R$ 上的标准平展代数 $R'$，而该环映射的核给出所需的素理想。
\end{proof}

\begin{proposition}
\label{algebra-proposition-etale-locally-standard}
设 $R \to S$ 是环映射，$\mathfrak q \subset S$ 是素理想。
若 $R \to S$ 在 $\mathfrak q$ 处平展，则存在
$g \in S$，$g \not \in \mathfrak q$，使得 $R \to S_g$
标准平展。
\end{proposition}

\begin{proof}
下面的证明稍显迂回，或许有缩短它的方法。

\medskip\noindent
第 1 步。由定义 \ref{algebra-definition-etale}，存在
$g \in S$，$g \not \in \mathfrak q$，使得 $R \to S_g$ 平展。
所以可以假设 $S$ 在 $R$ 上平展。

\medskip\noindent
第 2 步。由引理 \ref{algebra-lemma-etale}，存在平展环映射
$R_0 \to S_0$，其中 $R_0$ 在 $\mathbf{Z}$ 上有限型；还存在环映射
$R_0 \to R$，使得 $R = R \otimes_{R_0} S_0$。记
$\mathfrak q_0$ 为 $S_0$ 中对应于 $\mathfrak q$ 的素理想。
% P01-E0882: the frozen base-change identity has R on the left, although the construction and the next sentence require S = R tensor_{R_0} S_0.
% P01-E0883: the frozen source says “Denote q_0 the prime” and omits “by/as”.
若对 $(R_0 \to S_0, \mathfrak q_0)$ 证明了结论，则由基变换，
结论也对 $(R \to S, \mathfrak q)$ 成立。因此可以假设 $R$
是 Noether 环。

\medskip\noindent
第 3 步。注意，由引理 \ref{algebra-lemma-etale-quasi-finite}，
$R \to S$ 拟有限。由引理
\ref{algebra-lemma-quasi-finite-open-integral-closure}，存在有限环映射
$R \to S'$、$R$-代数映射 $S' \to S$，以及元素
$g' \in S'$，满足 $g' \not \in \mathfrak q$，使得
$S' \to S$ 诱导同构 $S'_{g'} \cong S_{g'}$。
% P01-E0884: the frozen sentence repeats “such that”, and writes g' notin q although g' is introduced in S' while q is an ideal of S without explicitly invoking its image.
（当然，一般而言，$S'$ 在 $R$ 上并不平展。）所以可以假设：
(a) $R$ 是 Noether 环；(b) $R \to S$ 有限；(c) $R \to S$
在 $\mathfrak q$ 处平展（但不再必然在所有素理想处平展）。

\medskip\noindent
第 4 步。设 $\mathfrak p \subset R$ 是对应于 $\mathfrak q$
的素理想。考虑纤维环 $S \otimes_R \kappa(\mathfrak p)$。
这是 $\kappa(\mathfrak p)$ 上的有限代数，因此是 Artin 环
（见引理 \ref{algebra-lemma-finite-dimensional-algebra}），从而是局部环的
有限乘积
$$
S \otimes_R \kappa(\mathfrak p) = \prod\nolimits_{i = 1}^n A_i
$$
（见命题 \ref{algebra-proposition-dimension-zero-ring}）。其中一个因子，比如
$A_1$，是局部环 $S_{\mathfrak q}/\mathfrak pS_{\mathfrak q}$，
它同构于 $\kappa(\mathfrak q)$，见引理
\ref{algebra-lemma-etale-at-prime}。其他因子对应于 $S$ 中位于
$\mathfrak p$ 之上的其他素理想，比如
$\mathfrak q_2, \ldots, \mathfrak q_n$。
% P01-E0885: the frozen display is followed directly by “see Proposition” without sentence punctuation.

\medskip\noindent
第 5 步。可以选择非零元素 $\alpha \in \kappa(\mathfrak q)$，
它生成有限可分域扩张
$\kappa(\mathfrak q)/\kappa(\mathfrak p)$（因此即使域扩张平凡，
也不允许 $\alpha = 0$）。注意，对任意
$\lambda \in \kappa(\mathfrak p)^*$，元素 $\lambda \alpha$ 也在
$\kappa(\mathfrak p)$ 上生成 $\kappa(\mathfrak q)$。考虑元素
$$
\overline{t} =
(\alpha, 0, \ldots, 0) \in
\prod\nolimits_{i = 1}^n A_i =
S \otimes_R \kappa(\mathfrak p).
$$
必要时如上用 $\lambda \alpha$ 替换 $\alpha$，可以假设
$\overline{t}$ 是某个 $t \in S$ 的像。设 $I \subset R[x]$
是 $R$-代数映射 $R[x] \to S$ 的核，其中该映射把 $x$ 映到
$t$。令 $S' = R[x]/I$，于是 $S' \subset S$。有交换图
$$
\xymatrix{
R[x] \ar[r] & S' \ar[r] & S \\
R \ar[u] \ar[ru] \ar[rru] & &
}
$$
由构造，$S$ 的素理想 $\mathfrak q_j$（$j \geq 2$）都位于
$R[x]$ 的素理想 $(\mathfrak p, x)$ 之上，而素理想
$\mathfrak q$ 位于 $R[x]$ 的另一个素理想之上，因为
$\alpha \not = 0$。

\medskip\noindent
第 6 步。记 $\mathfrak q' \subset S'$ 为 $S'$ 中对应于
$\mathfrak q$ 的素理想。由上述讨论，$\mathfrak q$ 是 $S$
中唯一位于 $\mathfrak q'$ 之上的素理想。因此
$S_{\mathfrak q} = S_{\mathfrak q'}$，见引理
\ref{algebra-lemma-unique-prime-over-localize-below}（对 $S' \to S$
有上升性质，见引理 \ref{algebra-lemma-integral-going-up}，因为
$S' \to S$ 有限，这是由于 $R \to S$ 有限）。
% P01-E0886: the frozen source says “Denote q' ... the prime” and omits “by/as”.
% P01-E0887: the frozen parenthesis says “we have going up” rather than “going up holds”.
由此，$S'_{\mathfrak q'} \to S_{\mathfrak q}$ 有限且为单射，
因为它是有限单射环映射 $S' \to S$ 的局部化。考虑局部环映射
$$
R_{\mathfrak p} \to S'_{\mathfrak q'} \to S_{\mathfrak q}
$$
第二个映射有限且为单射。由引理 \ref{algebra-lemma-etale-at-prime}，
$S_{\mathfrak q}/\mathfrak pS_{\mathfrak q} = \kappa(\mathfrak q)$。
从而更有
$S_{\mathfrak q}/\mathfrak q'S_{\mathfrak q} = \kappa(\mathfrak q)$。
由于
$$
\kappa(\mathfrak p) \subset \kappa(\mathfrak q') \subset \kappa(\mathfrak q)
$$
并且 $\alpha$ 属于 $\kappa(\mathfrak q')$ 在
$\kappa(\mathfrak q)$ 中的像，可知
$\kappa(\mathfrak q') = \kappa(\mathfrak q)$。于是，对
$S'_{\mathfrak q'}$-模映射
$S'_{\mathfrak q'} \to S_{\mathfrak q}$ 应用 Nakayama 引理
\ref{algebra-lemma-NAK}，可知映射
$S'_{\mathfrak q'} \to S_{\mathfrak q}$ 满。换言之，
$S'_{\mathfrak q'} \cong S_{\mathfrak q}$。

\medskip\noindent
第 7 步。由引理 \ref{algebra-lemma-isomorphic-local-rings}，存在
$g \in S$，$g \not \in \mathfrak q$ 和
$g' \in S'$，$g' \not \in \mathfrak q'$，使得
$S'_{g'} \cong S_g$。由于 $R$ 是 Noether 环，$S'$ 在 $R$
上有限，因为它是有限 $R$-模 $S$ 的 $R$-子模。因此用 $S'$
替换 $S$ 后，可以假设：(a) $R$ 是 Noether 环；
(b) $S$ 在 $R$ 上有限；(c) $S$ 在 $\mathfrak q$ 处相对于
$R$ 平展；(d) $S = R[x]/I$。
% P01-E0888: item (b) in the frozen source omits the verb “is” before “finite over R”.

\medskip\noindent
第 8 步。考虑环
$S \otimes_R \kappa(\mathfrak p) = \kappa(\mathfrak p)[x]/\overline{I}$，
其中 $\overline{I} = I \cdot \kappa(\mathfrak p)[x]$ 是 $I$ 在
$\kappa(\mathfrak p)[x]$ 中生成的理想。由于
$\kappa(\mathfrak p)[x]$ 是主理想整环，可知对某个首一多项式
$\overline{h} \in \kappa(\mathfrak p)[x]$，有
$\overline{I} = (\overline{h})$。用某个
$\lambda \in \kappa(\mathfrak p)$ 的
$\lambda \cdot \overline{h}$ 替换 $\overline{h}$ 后，可以假设
$\overline{h}$ 是某个 $h \in I \subset R[x]$ 的像。
（问题在于我们不知道能否选择首一的 $h$。）此外，与第 4 步中一样，
有 $S \otimes_R \kappa(\mathfrak p) = A_1 \times \ldots \times A_n$，
其中 $A_1 = \kappa(\mathfrak q)$ 是
$\kappa(\mathfrak p)$ 的有限可分扩张，而
$A_2, \ldots, A_n$ 是局部环。这蕴含
$$
\overline{h} = \overline{h}_1 \overline{h}_2^{e_2} \ldots \overline{h}_n^{e_n}
$$
其中 $\overline{h}_i \in \kappa(\mathfrak p)[x]$ 是两两互素的
不可约首一多项式，$e_2, \ldots, e_n \geq 1$。编号的选取使得
作为 $\kappa(\mathfrak p)[x]$-代数，
$A_i = \kappa(\mathfrak p)[x]/(\overline{h}_i^{e_i})$。
% P01-E0889: the frozen formula uses e_i without defining e_1, while only e_2,...,e_n are introduced and the A_1 factor is separately a field.
注意，$\overline{h}_1$ 是 $\alpha \in \kappa(\mathfrak q)$
的极小多项式，因而是可分多项式（其导数与它本身互素）。

\medskip\noindent
第 9 步。设 $m \in I$ 是首一元素；这样的元素存在，因为环扩张
$R \to R[x]/I$ 有限，因而整。记 $\overline{m}$ 为它在
$\kappa(\mathfrak p)[x]$ 中的像。可以分解为
% P01-E0890: the frozen source says “Denote mbar the image” and omits “by/as”.
$$
\overline{m} = \overline{k}
\overline{h}_1^{d_1} \overline{h}_2^{d_2} \ldots \overline{h}_n^{d_n}
$$
其中 $d_1 \geq 1$，$d_j \geq e_j$，$j = 2, \ldots, n$，
并且 $\overline{k} \in \kappa(\mathfrak p)[x]$ 与所有
$\overline{h}_i$ 互素。令 $f = m^l + h$，其中
$l \deg(m) > \deg(h)$，且 $l \geq 2$。于是 $f$ 是 $R$
上的首一多项式。此外，$f$ 在 $\kappa(\mathfrak p)[x]$
中的像 $\overline{f}$ 分解为
$$
\overline{f} =
\overline{h}_1 \overline{h}_2^{e_2} \ldots \overline{h}_n^{e_n}
+
\overline{k}^l \overline{h}_1^{ld_1} \overline{h}_2^{ld_2}
\ldots \overline{h}_n^{ld_n}
=
\overline{h}_1(\overline{h}_2^{e_2} \ldots \overline{h}_n^{e_n}
+
\overline{k}^l
\overline{h}_1^{ld_1 - 1} \overline{h}_2^{ld_2} \ldots \overline{h}_n^{ld_n})
= \overline{h}_1 \overline{w}
$$
其中 $\overline{w}$ 是与 $\overline{h}_1$ 互素的多项式。
令 $g = f'$（关于 $x$ 的导数）。

\medskip\noindent
第 10 步。环映射 $R[x] \to S = R[x]/I$ 具有以下性质：
(1) 它把 $f$ 映到零；(2) 它把 $g$ 映到
$S \setminus \mathfrak q$ 中的元素。第一条断言是显然的，因为
$f$ 是 $I$ 的元素。要证明第二条断言，只需说明 $g$ 在
$\kappa(\mathfrak q) = \kappa(\mathfrak p)[x]/(\overline{h}_1)$
中的像不为零。$g$ 在 $\kappa(\mathfrak p)[x]$ 中的像是
$\overline{f}$ 的导数。因此 (2) 是显然的，因为
$$
\overline{g} =
\frac{\text{d}\overline{f}}{\text{d}x} =
\overline{w}\frac{\text{d}\overline{h}_1}{\text{d}x} +
\overline{h}_1\frac{\text{d}\overline{w}}{\text{d}x},
$$
$\overline{w}$ 与 $\overline{h}_1$ 互素，而
$\overline{h}_1$ 可分。

\medskip\noindent
第 11 步。由此可知，$\varphi : R[x]/(f) \to S$ 是满环映射，
$R[x]_g/(f)$ 在 $R$ 上平展（因为它标准平展，见引理
\ref{algebra-lemma-standard-etale}），且 $\varphi(g) \not \in \mathfrak q$。
选择元素 $g' \in R[x]/(f)$，使得
$\varphi(g') \not \in \mathfrak q$，且 $S_{\varphi(g')}$
在 $R$ 上也平展（这样的元素存在，因为 $S$ 在 $\mathfrak q$
处相对于 $R$ 平展）。于是环映射
$R[x]_{gg'}/(f) \to S_{\varphi(gg')}$ 是 $R$ 上平展代数之间的
满射。因此由引理 \ref{algebra-lemma-map-between-etale}，它平展。
再由引理 \ref{algebra-lemma-surjective-flat-finitely-presented}，
它是局部化。这样，$S$ 在某个不属于 $\mathfrak q$ 的元素处的
局部化，同构于 $R$ 上某个标准平展代数的局部化，这正是所要证明的。
\end{proof}

\noindent
下面两个引理说明，平展拓扑比由 Zariski 覆盖和有限平坦态射生成的
拓扑更粗。初次阅读时应当跳过它们。

\begin{lemma}
\label{algebra-lemma-standard-etale-finite-flat-Zariski}
设 $R \to S$ 是标准平展态射。则存在具有下列性质的环映射
$R \to S'$：
% P01-E0891: the frozen source omits a colon after “the following properties”.
\begin{enumerate}
\item $R \to S'$ 有限、有限表现且平坦
（换言之，$S'$ 作为 $R$-模有限投射）；
\item $\Spec(S') \to \Spec(R)$ 满射；
\item 对每个位于 $\mathfrak p \subset R$ 之上的素理想
$\mathfrak q \subset S$，以及每个位于 $\mathfrak p$ 之上的素理想
$\mathfrak q' \subset S'$，存在 $g' \in S'$，
$g' \not \in \mathfrak q'$，使得环映射 $R \to S'_{g'}$
通过映射 $\varphi : S \to S'_{g'}$ 分解，并满足
$\varphi^{-1}(\mathfrak q'S'_{g'}) = \mathfrak q$。
\end{enumerate}
\end{lemma}

\begin{proof}
设 $S = R[x]_g/(f)$ 是定义 \ref{algebra-definition-standard-etale}
中那样的 $S$ 的表现。写成
$f = x^n + a_1 x^{n - 1} + \ldots + a_n$，其中 $a_i \in R$。
由引理 \ref{algebra-lemma-adjoin-roots}，存在有限自由且忠实平坦的环映射
$R \to S'$，使得对某些 $\alpha_i \in S'$，有
$f = \prod (x - \alpha_i)$。因此 $R \to S'$ 满足条件 (1)、(2)。
设 $\mathfrak q \subset R[x]/(f)$ 是满足
$g \not \in \mathfrak q$ 的素理想（即它对应于 $S$ 的一个素理想）。
令 $\mathfrak p = R \cap \mathfrak q$，并设
$\mathfrak q' \subset S'$ 是位于 $\mathfrak p$ 之上的素理想。
注意，存在 $n$ 个 $R$-代数映射
\begin{eqnarray*}
\varphi_i : R[x]/(f) & \longrightarrow & S' \\
x & \longmapsto & \alpha_i
\end{eqnarray*}
为完成证明，需要说明对某个 $i$，有
(a) $\varphi_i(g)$ 在 $\kappa(\mathfrak q')$ 中的像不为零；
(b) $\varphi_i^{-1}(\mathfrak q') = \mathfrak q$。
因为这时只需取 $g' = \varphi_i(g)$，并对该 $i$ 取
$\varphi = \varphi_i$。

\medskip\noindent
设 $\overline{f}$ 表示 $f$ 在 $\kappa(\mathfrak p)[x]$ 中的像。
注意，作为 $\Spec(\kappa(\mathfrak p)[x]/(\overline{f}))$
的一个点，素理想 $\mathfrak q$ 对应于 $\overline{f}$ 的某个
不可约因子 $f_1$。此外，$g \not \in \mathfrak q$ 意味着
$f_1$ 不整除 $g$ 在 $\kappa(\mathfrak p)[x]$ 中的像
$\overline{g}$。记
$\overline{\alpha}_1, \ldots, \overline{\alpha}_n$ 为
$\alpha_1, \ldots, \alpha_n$ 在 $\kappa(\mathfrak q')$ 中的像。
% P01-E0892: the frozen source says “Denote alpha-bar_1,... the images” and omits “by/as”.
注意，多项式 $\overline{f}$ 在 $\kappa(\mathfrak q')[x]$
中完全分裂，即
$$
\overline{f} = \prod\nolimits_i (x - \overline{\alpha}_i)
$$
此外，$\varphi_i(g)$ 约化为
$\overline{g}(\overline{\alpha}_i)$。因此可以选择 $i$，使得
$f_1(\overline{\alpha}_i) = 0$ 且
$\overline{g}(\overline{\alpha}_i) \not = 0$。
对这个 $i$，性质 (a)、(b) 成立。略去一些细节。
\end{proof}

\begin{lemma}
\label{algebra-lemma-etale-finite-flat-zariski}
设 $R \to S$ 是环映射。假设
\begin{enumerate}
\item $R \to S$ 平展；
\item $\Spec(S) \to \Spec(R)$ 满射。
\end{enumerate}
则存在环映射 $R \to S'$，使得
\begin{enumerate}
\item $R \to S'$ 有限、有限表现且平坦
（换言之，它作为 $R$-模有限投射）；
\item $\Spec(S') \to \Spec(R)$ 满射；
\item 对每个素理想 $\mathfrak q' \subset S'$，存在
$g' \in S'$，$g' \not \in \mathfrak q'$，使得环映射
$R \to S'_{g'}$ 分解为 $R \to S \to S'_{g'}$。
\end{enumerate}
\end{lemma}

\begin{proof}
由命题 \ref{algebra-proposition-etale-locally-standard} 和
$\Spec(S)$ 的拟紧性（见引理 \ref{algebra-lemma-quasi-compact}），
可以找到生成 $S$ 的单位理想的 $g_1, \ldots, g_n \in S$，
使得每个 $R \to S_{g_i}$ 都标准平展。若对环映射
$R \to \prod_{i = 1, \ldots, n} S_{g_i}$ 证明了引理，
则对环映射 $R \to S$ 也得到引理。因此可以假设
$S = \prod_{i = 1, \ldots, n} S_i$ 是标准平展态射的有限乘积。

\medskip\noindent
对每个 $i$，选取引理
\ref{algebra-lemma-standard-etale-finite-flat-Zariski} 中适用于标准平展态射
$R \to S_i$ 的环映射 $R \to S_i'$。令
$S' = S_1' \otimes_R \ldots \otimes_R S_n'$；下文不再另行说明地
使用 $R$-代数映射 $S_i' \to S'$。我们断言这给出所需构造。
性质 (1)、(2) 是直接的。对性质 (3)，假设
$\mathfrak q' \subset S'$ 是素理想。记 $\mathfrak p$ 为它在
$\Spec(R)$ 中的像。选择 $i \in \{1, \ldots, n\}$，使得
$\mathfrak p$ 属于 $\Spec(S_i) \to \Spec(R)$ 的像；由假设，
这是可能的。令 $\mathfrak q_i' \subset S_i'$ 为
$\mathfrak q'$ 在 $S_i'$ 的谱中的像。
% P01-E0893: the frozen source says “Denote p its image” and omits “by/as”.
% P01-E0894: the frozen source says “Set q_i' ... the image” and omits “to be/as”.
由 $S'_i$ 的构造，存在 $g'_i \in S_i'$，使得
$R \to (S_i')_{g_i'}$ 分解为
$R \to S_i \to (S_i')_{g_i'}$。因此
$R \to S'_{g_i'}$ 也分解为
% P01-E0895: the frozen proof invokes the preceding lemma but omits the inherited condition g'_i notin q'_i, which is needed to match property (3)'s neighborhood statement.
$$
R \to S_i \to (S_i')_{g_i'} \to S'_{g_i'}
$$
如所要求。
\end{proof}

\section{拟有限环映射的平展局部结构}
\label{algebra-section-etale-local-quasi-finite}

\noindent
下面这些引理粗略地说，经过平展扩张后，
拟有限环映射会成为有限环映射。
为帮助理解这些结果，回顾有限型环映射为拟有限的轨迹是开集
（见引理 \ref{algebra-lemma-quasi-finite-open}），而且这一轨迹的形成
与任意基变换交换
（见引理 \ref{algebra-lemma-quasi-finite-base-change}）。

\begin{lemma}
\label{algebra-lemma-produce-finite}
设 $R \to S' \to S$ 是环映射。
设 $\mathfrak p \subset R$ 是素理想。
设 $g \in S'$ 是元素。
假设
\begin{enumerate}
\item $R \to S'$ 是整的；
\item $R \to S$ 是有限型的；
\item $S'_g \cong S_g$；并且
\item $g$ 在 $S' \otimes_R \kappa(\mathfrak p)$ 中可逆。
% P01-E0896: the frozen source says “g invertible” and omits “is”.
\end{enumerate}
则存在 $f \in R$，$f \not \in \mathfrak p$，使得
$R_f \to S_f$ 有限。
% P01-E0897: the frozen source uses the ungrammatical article phrase “a f”.
\end{lemma}

\begin{proof}
由假设，$V(g) \subset \Spec(S')$ 在态射
$\Spec(S') \to \Spec(R)$ 下的像 $T$ 不含
$\mathfrak p$。由第 \ref{algebra-section-going-up} 节，
特别是引理 \ref{algebra-lemma-going-up-closed}，可见 $T$ 是闭集。
% P01-E0898: the frozen source has the malformed citation phrase “By Section ..., especially, Lemma ...”.
取 $f \in R$，$f \not \in \mathfrak p$，使得
$T \cap D(f) = \emptyset$。于是 $g$ 在 $S'_f$ 中变得可逆，
因而 $S'_f \cong S_f$。所以 $S_f$ 在 $R_f$ 上既是有限型的
又是整的，从而是有限的。
\end{proof}

\begin{lemma}
\label{algebra-lemma-etale-makes-quasi-finite-finite-one-prime}
设 $R \to S$ 是环映射。
设 $\mathfrak q \subset S$ 是位于
素理想 $\mathfrak p \subset R$ 之上的素理想。
假设 $R \to S$ 是有限型的，并且在 $\mathfrak q$ 处拟有限。
% P01-E0899: the frozen source says “Assume R -> S finite type” and omits “is”.
则存在：
\begin{enumerate}
\item 平展环映射 $R \to R'$；
\item 位于 $\mathfrak p$ 之上的素理想 $\mathfrak p' \subset R'$；
\item 乘积分解
$$
R' \otimes_R S = A \times B
$$
\end{enumerate}
% P01-E0900: the frozen source omits punctuation after “with the following properties”.
具有下列性质：
\begin{enumerate}
\item $\kappa(\mathfrak p) = \kappa(\mathfrak p')$；
\item $R' \to A$ 有限；
\item $A$ 恰有一个位于 $\mathfrak p'$ 之上的素理想
$\mathfrak r$；
\item $\mathfrak r$ 位于 $\mathfrak q$ 之上；并且
\item $B$ 不含同时位于 $\mathfrak q$ 和 $\mathfrak p'$ 之上的素理想。
\end{enumerate}
\end{lemma}

\begin{proof}
设 $S' \subset S$ 是 $R$ 在 $S$ 中的整闭包。
令 $\mathfrak q' = S' \cap \mathfrak q$。
由 Zariski 主定理 \ref{algebra-theorem-main-theorem}，
存在 $g \in S'$，$g \not \in \mathfrak q'$，使得
$S'_g \cong S_g$。考虑纤维环
$F = S \otimes_R \kappa(\mathfrak p)$ 和
$F' = S' \otimes_R \kappa(\mathfrak p)$。记
$\overline{\mathfrak q}'$ 为 $F'$ 中对应于 $\mathfrak q'$ 的素理想。
% P01-E0901: the frozen source says “Denote q-bar' the prime” and omits “by/as”.
因为 $F'$ 在 $\kappa(\mathfrak p)$ 上整，所以
$\overline{\mathfrak q}'$ 是 $\Spec(F')$ 的闭点，
见引理 \ref{algebra-lemma-integral-over-field}。
注意，$\mathfrak q$ 定义了 $\Spec(F)$ 的一个孤立闭点
$\overline{\mathfrak q}$
（见定义 \ref{algebra-definition-quasi-finite}）。
因为 $S'_g \cong S_g$，所以 $F'_g \cong F_g$，
从而 $\overline{\mathfrak q}$ 和 $\overline{\mathfrak q}'$
在 $\Spec(F)$ 与 $\Spec(F')$ 中有同构的开邻域。
我们断定集合
$\{\overline{\mathfrak q}'\} \subset \Spec(F')$
是开集。再结合 $\mathfrak q'$ 是闭点（已在上面证明），
便断定 $\overline{\mathfrak q}'$ 也定义了
$\Spec(F')$ 的一个孤立闭点。
% P01-E0902: the frozen source says q' is closed although the preceding argument proves q-bar' is closed.

\medskip\noindent
另有一个小观察：在映射
$\Spec(F) \to \Spec(F')$ 下，$\overline{\mathfrak q}$
是映到 $\overline{\mathfrak q}'$ 的唯一一点。
这由上面的讨论推出。

\medskip\noindent
由引理 \ref{algebra-lemma-disjoint-implies-product}，可以写成
$F' = F'_1 \times F'_2$，其中
$\Spec(F'_1) = \{\overline{\mathfrak q}'\}$。
因为 $F' = S' \otimes_R \kappa(\mathfrak p)$，
存在 $s' \in S'$，它映到元素
$(r, 0) \in F'_1 \times F'_2 = F'$，
其中某个 $r \in R$ 满足 $r \not \in \mathfrak p$。
事实上，我们将用到的关于 $s'$ 的事实是：它是 $S'$ 的元素，
不属于 $\mathfrak q'$，而属于任何其他位于
$\mathfrak p$ 之上的素理想。

\medskip\noindent
取首一多项式 $f(x) \in R[x]$，使得 $f(s') = 0$。
记 $\overline{f} \in \kappa(\mathfrak p)[x]$ 为其像。
% P01-E0903: the frozen source says “Denote f-bar the image” and omits “by/as”.
可以将它分解为 $\overline{f} = x^e \overline{h}$，其中
$\overline{h}(0) \not = 0$。必要时将 $f$ 替换为
$x f$，便可假设 $e \geq 1$。
由引理 \ref{algebra-lemma-factor-mod-lift-etale}，
可以找到平展环扩张 $R \to R'$、
位于 $\mathfrak p$ 之上的素理想 $\mathfrak p'$，以及
$R'[x]$ 中的分解 $f = h i$，使得
$\kappa(\mathfrak p) = \kappa(\mathfrak p')$、
$\overline{h} = h \bmod \mathfrak p'$、
$x^e = i \bmod \mathfrak p'$，并且
可以在 $R'[x]$ 中写成 $a h + b i = 1$
（其中 $a, b$ 适当）。

\medskip\noindent
考虑元素 $h(s'), i(s') \in R' \otimes_R S'$。
由构造，$h(s')i(s') = f(s') = 0$。另一方面，
它们生成单位理想，因为 $a(s')h(s') + b(s')i(s') = 1$。
所以 $R' \otimes_R S'$ 是在这些元素处的局部化之积：
$$
R' \otimes_R S'
=
(R' \otimes_R S')_{i(s')}
\times
(R' \otimes_R S')_{h(s')}
=
S'_1 \times S'_2
$$
此外，这一乘积分解与我们为纤维环 $F'$ 得到的乘积分解相容；
这是由我们对 $s', i, h$ 的选择得到的，因为这些选择保证
$\overline{\mathfrak q}'$ 是 $F'$ 中唯一不含
$i(s')$ 在 $F'$ 中之像的素理想。这里使用了如下事实：
$R'\otimes_R S'$ 在 $R'$ 的 $\mathfrak p'$ 处的纤维环
与 $F'$ 相同，因为
$\kappa(\mathfrak p) = \kappa(\mathfrak p')$。
由此，$S'_1$ 恰有一个位于 $\mathfrak p'$ 之上的素理想，
记为 $\mathfrak r'$，并且这个素理想位于
$\mathfrak q'$ 之上。
% P01-E0904: the frozen source contains doubled spacing in “S'_1  has”.
因此元素 $g \in S'$ 在 $S'_1$ 中的像不属于
$\mathfrak r'$。

\medskip\noindent
基变换 $R'\otimes_R S$ 继承类似的乘积分解
$$
R' \otimes_R S
=
(R' \otimes_R S)_{i(s')}
\times
(R' \otimes_R S)_{h(s')}
=
S_1 \times S_2
$$
由上可知，$S_1$ 恰有一个位于 $\mathfrak p'$ 之上的素理想，
记为 $\mathfrak r$（如上考虑纤维环即可），
并且这个素理想位于 $\mathfrak q$ 之上。

\medskip\noindent
现在对环映射 $R' \to S'_1 \to S_1$、
素理想 $\mathfrak p'$ 和元素 $g$ 应用引理
\ref{algebra-lemma-produce-finite}，可知在将 $R'$ 替换为某个主局部化后，
可以假设 $S_1$ 在 $R'$ 上有限，正如所需。
\end{proof}

\begin{lemma}
\label{algebra-lemma-etale-makes-quasi-finite-finite}
设 $R \to S$ 是环映射。
设 $\mathfrak p \subset R$ 是素理想。
假设 $R \to S$ 是有限型的。
则存在：
% P01-E0905: the frozen source introduces the following enumeration without punctuation after “there exists”.
\begin{enumerate}
\item 平展环映射 $R \to R'$；
\item 位于 $\mathfrak p$ 之上的素理想 $\mathfrak p' \subset R'$；
\item 乘积分解
$$
R' \otimes_R S = A_1 \times \ldots \times A_n \times B
$$
\end{enumerate}
具有下列性质：
% P01-E0906: the frozen source omits punctuation after “with the following properties”.
\begin{enumerate}
\item 有 $\kappa(\mathfrak p) = \kappa(\mathfrak p')$；
\item 每个 $A_i$ 在 $R'$ 上有限；
\item 每个 $A_i$ 恰有一个位于 $\mathfrak p'$ 之上的素理想
$\mathfrak r_i$；并且
\item $R' \to B$ 在任何位于 $\mathfrak p'$ 之上的素理想处
都不是拟有限的。
% P01-E0907: the frozen source says “R' -> B not quasi-finite” and omits “is”.
\end{enumerate}
\end{lemma}

\begin{proof}
记 $F = S \otimes_R \kappa(\mathfrak p)$ 为 $S/R$
在素理想 $\mathfrak p$ 处的纤维环。
% P01-E0908: the frozen source says “Denote F ... the fibre ring” and omits “by/as”.
由于 $F$ 在 $\kappa(\mathfrak p)$ 上有限型，所以它是 Noether 环，
因而 $\Spec(F)$ 只有有限多个孤立闭点。若没有孤立闭点，
即不存在 $S$ 中位于 $\mathfrak p$ 之上的素理想
$\mathfrak q$，使得 $S/R$ 在 $\mathfrak q$ 处拟有限，
那么引理成立。
若至少存在一个这样的素理想 $\mathfrak q$，则可以应用引理
\ref{algebra-lemma-etale-makes-quasi-finite-finite-one-prime}。
这给出该引理中的图表
$$
\xymatrix{
S \ar[r] & R'\otimes_R S \ar@{=}[r] & A_1 \times B' \\
R \ar[r] \ar[u] & R' \ar[u] \ar[ru]
}
$$
因为 $\mathfrak p$ 与 $\mathfrak p'$ 处的剩余域相同，
$S/R$ 与 $(A_1 \times B')/R'$ 的纤维环相同。
因此，对纤维中孤立闭点的数目作归纳，可以假设
引理对 $R' \to B'$ 和 $\mathfrak p'$ 成立。
于是得到平展环映射 $R' \to R''$、
素理想 $\mathfrak p'' \subset R''$，以及分解
$$
R'' \otimes_{R'} B' = A_2 \times \ldots \times A_n \times B
$$
% P01-E0909: the frozen source leaves the displayed decomposition without terminal punctuation.
略去如下验证：环映射 $R \to R''$、素理想
$\mathfrak p''$ 以及所得分解
$$
R'' \otimes_R S = (R'' \otimes_{R'} A_1) \times
A_2 \times \ldots \times A_n \times B
$$
给出引理中问题的一个解。
\end{proof}

\begin{lemma}
\label{algebra-lemma-etale-makes-quasi-finite-finite-variant}
设 $R \to S$ 是环映射。
设 $\mathfrak p \subset R$ 是素理想。
假设 $R \to S$ 是有限型的。
% P01-E0910: the frozen source says “Assume R -> S finite type” and omits “is”.
则存在：
% P01-E0911: the frozen source introduces the following enumeration without punctuation after “there exists”.
\begin{enumerate}
\item 平展环映射 $R \to R'$；
\item 位于 $\mathfrak p$ 之上的素理想 $\mathfrak p' \subset R'$；
\item 乘积分解
$$
R' \otimes_R S = A_1 \times \ldots \times A_n \times B
$$
\end{enumerate}
具有下列性质：
% P01-E0912: the frozen source omits punctuation after “with the following properties”.
\begin{enumerate}
\item 每个 $A_i$ 在 $R'$ 上有限；
\item 每个 $A_i$ 恰有一个位于 $\mathfrak p'$ 之上的素理想
$\mathfrak r_i$；
\item 有限域扩张
$\kappa(\mathfrak r_i)/\kappa(\mathfrak p')$
是纯不可分的；并且
\item $R' \to B$ 在任何位于 $\mathfrak p'$ 之上的素理想处
都不是拟有限的。
% P01-E0913: the frozen source says “R' -> B not quasi-finite” and omits “is”.
\end{enumerate}
\end{lemma}

\begin{proof}
证明的策略是作两次平展环扩张：先控制剩余域，
再应用引理 \ref{algebra-lemma-etale-makes-quasi-finite-finite}。

\medskip\noindent
记 $F = S \otimes_R \kappa(\mathfrak p)$ 为 $S/R$
在素理想 $\mathfrak p$ 处的纤维环。
% P01-E0914: the frozen source says “Denote F ... the fibre ring” and omits “by/as”.
与引理 \ref{algebra-lemma-etale-makes-quasi-finite-finite} 的证明中一样，
存在有限多个素理想，记为
$\mathfrak q_1, \ldots, \mathfrak q_n$，它们是 $S$ 中位于
$R$ 之上的素理想，并且环映射 $R \to S$ 在这些素理想处拟有限。
% P01-E0915: the frozen source has the typo “finitely may primes”.
% P01-E0916: the frozen source says these primes of S lie over R, but the fibre argument requires them to lie over p.
设 $\kappa(\mathfrak p) \subset L_i \subset \kappa(\mathfrak q_i)$
是这样的子域：$\kappa(\mathfrak p) \subset L_i$
是可分扩张，而域扩张 $\kappa(\mathfrak q_i)/L_i$
是纯不可分的。设 $L/\kappa(\mathfrak p)$
是一个有限 Galois 扩张，使得对 $i = 1, \ldots, n$，
$L_i$ 可嵌入其中。
由引理 \ref{algebra-lemma-make-etale-map-prescribed-residue-field}，
可以找到平展环扩张
$R \to R'$ 以及位于 $\mathfrak p$ 之上的素理想
$\mathfrak p'$，使得域扩张
$\kappa(\mathfrak p')/\kappa(\mathfrak p)$ 同构于
$\kappa(\mathfrak p) \subset L$。
因此 $R' \otimes_R S$ 在 $\mathfrak p'$ 处的纤维环同构于
$F \otimes_{\kappa(\mathfrak p)} L$。
位于 $\mathfrak q_i$ 之上的素理想对应于
$\kappa(\mathfrak q_i) \otimes_{\kappa(\mathfrak p)} L$
的素理想；由 $L$ 的选择和初等域论，后者是
在 $L$ 上纯不可分的域之积。
由引理 \ref{algebra-lemma-quasi-finite-base-change}，
这些也是位于 $\mathfrak p'$ 之上、使
$R' \to R' \otimes_R S$ 在其处拟有限的全部素理想。
因此，将 $R$ 替换为 $R'$、将 $\mathfrak p$ 替换为
$\mathfrak p'$、并将 $S$ 替换为 $R' \otimes_R S$ 后，
可以假设：对所有位于 $\mathfrak p$ 之上且使
$S/R$ 在其处拟有限的素理想 $\mathfrak q$，
域扩张
$\kappa(\mathfrak q)/\kappa(\mathfrak p)$
都是纯不可分的。

\medskip\noindent
接着应用引理 \ref{algebra-lemma-etale-makes-quasi-finite-finite}。
这就得到所需结果，因为在这次平展环扩张下域扩张不变。
\end{proof}

\section{局部同态}
\label{algebra-section-local-homomorphisms}

\noindent
这里收录若干没有自然所属章节的引理。
第一个引理粗略地说，局部环之间的平展映射，
模掉一个由非单位元生成的主理想的任意次幂后都是同构。

\begin{lemma}
\label{algebra-lemma-lindel}
\begin{reference}
\cite[Lemma on page 321]{Lindel}, \cite[Lemma 4.1.5]{KC}
\end{reference}
设 $(R, \mathfrak m_R) \to (S, \mathfrak m_S)$ 是局部环的
局部同态。假设 $S$ 是 $R$ 的某个平展环扩张的局部化，
并且 $\kappa(\mathfrak m_R) \to \kappa(\mathfrak m_S)$
是同构。则存在 $t \in \mathfrak m_R$，使得对所有
$n \geq 1$，$R/t^nR \to S/t^nS$ 都是同构。
% P01-E0917: the frozen source uses the ungrammatical article phrase “an t”.
% P01-E0918: the frozen source misspells “isomorphism” as “isomorphsm”.
\end{lemma}

\begin{proof}
写成 $S = T_{\mathfrak q}$，其中 $T$ 是某个平展
$R$-代数，$\mathfrak q \subset T$ 是位于
$\mathfrak m_R$ 之上的素理想。
由命题 \ref{algebra-proposition-etale-locally-standard}，
可以假设 $R \to T$ 标准平展。
按照定义 \ref{algebra-definition-standard-etale}，写成
$T = R[x]_g/(f)$。由关于剩余域的假设，可以选择
$a \in R$，使得 $x$ 和 $a$ 在
$\kappa(\mathfrak q) = \kappa(\mathfrak m_S) = \kappa(\mathfrak m_R)$
中的像相同。于是用 $x - a$ 替换 $x$ 后，可以假设
$\mathfrak q$ 在 $T$ 中由 $x$ 和 $\mathfrak m_R$ 生成。
特别地，$t = f(0) \in \mathfrak m_R$。
下面证明 $t = f(0)$ 满足要求。

\medskip\noindent
写成 $f = x^d + \sum_{i = 1, \ldots, d - 1} a_i x^i + t$。
因为 $R \to T$ 标准平展，可知 $a_1$ 是 $R$ 中的单位元：
$f$ 的导数在 $T$ 中可逆，特别地，它不属于
$\mathfrak q$。
% P01-E0919: the frozen source joins the derivative clauses without the needed conjunction or punctuation.
令 $h = a_1 + a_2 x + \ldots + a_{d - 1} x^{d - 2} + x^{d - 1} \in R[x]$，
从而在 $R[x]$ 中有 $f = t + xh$。可见
$h \not \in \mathfrak q$，因此可以用
$R[x]_{hg}/(f)$ 替换 $T$。替换后可见
$$
T/tT = (R/tR)[x]_{hg}/(f) = (R/tR)[x]_{hg}/(xh) = (R/tR)[x]_{hg}/(x)
$$
是 $R/tR$ 的商。由引理
\ref{algebra-lemma-surjective-mod-locally-nilpotent}，可知对所有
$n \geq 1$，$R/t^nR \to T/t^nT$ 都满射。
另一方面，已知平坦局部环映射
$R/t^nR \to S/t^nS$ 对每个 $n$ 都通过
$R/t^nR \to T/t^nT$ 分解，因此这些映射也单射
（局部环之间的平坦局部同态忠实平坦，因而单射，见
引理 \ref{algebra-lemma-local-flat-ff} 和
\ref{algebra-lemma-faithfully-flat-universally-injective}）。
因为 $S$ 是 $T$ 的局部化，所以 $S/t^nS$ 是
$T/t^nT = R/t^nR$ 在位于极大理想之上的某个素理想处的局部化；
但这个环已经是局部环，证明完成。
\end{proof}

\begin{lemma}
\label{algebra-lemma-etale-under-finite-flat}
设 $(R, \mathfrak m_R) \to (S, \mathfrak m_S)$ 是局部环的
局部同态。假设 $S$ 是 $R$ 的某个平展环扩张的局部化。
则存在有限、有限表现且忠实平坦的环映射 $R \to S'$，
使得对 $S'$ 的每个极大理想 $\mathfrak m'$，环映射
$R \to S'_{\mathfrak m'}$ 都有分解
$$
R \to S \to S'_{\mathfrak m'}.
$$
\end{lemma}

\begin{proof}
写成 $S = T_{\mathfrak q}$，其中 $T$ 是某个平展
$R$-代数。由命题 \ref{algebra-proposition-etale-locally-standard}，
可以假设 $T$ 标准平展。
对环映射 $R \to T$ 应用引理
\ref{algebra-lemma-standard-etale-finite-flat-Zariski}，得到
$R \to S'$。特别地，对 $S'$ 的每个极大理想
$\mathfrak m'$，都得到某个 $g' \not \in \mathfrak m'$ 和分解
$\varphi : T \to S'_{g'}$，使得
$\mathfrak q = \varphi^{-1}(\mathfrak m'S'_{g'})$。
因此 $\varphi$ 诱导所需的局部环映射
$S \to S'_{\mathfrak m'}$。
\end{proof}

\section{整闭包与光滑基变换}
\label{algebra-section-integral-closure-smooth-base-change}

\begin{lemma}
\label{algebra-lemma-trick}
设 $R$ 是环。
设 $f \in R[x]$ 是首一多项式。
设 $R \to B$ 是环映射。
若 $h \in B[x]/(f)$ 在 $R$ 上整，则元素
$f' h$ 可以写成 $f'h = \sum_i b_i x^i$，其中
$b_i \in B$ 在 $R$ 上整。
\end{lemma}

\begin{proof}
设在环 $B[x]/(f)$ 中有
$h^e + r_1 h^{e - 1} + \ldots + r_e = 0$，其中
$r_i \in R$。
存在有限自由环扩张 $B \subset B'$，使得对某些
$\alpha_i \in B'$ 有
$f = (x - \alpha_1) \ldots (x - \alpha_d)$，
见引理 \ref{algebra-lemma-adjoin-roots}。
注意，每个 $\alpha_i$ 都在 $R$ 上整。
可以用 $h_i \in B$ 将 $h$ 表成
$h = h_0 + h_1 x + \ldots + h_{d - 1} x^{d - 1}$。
于是有普遍恒等式
$$
f' h =
\sum\nolimits_{i = 1, \ldots, d}
h(\alpha_i)
(x - \alpha_1) \ldots \widehat{(x - \alpha_i)} \ldots (x - \alpha_d)
$$
作为 $B'[x]/(f)$ 的元素。要证明这一点，只需在环
$B_{univ} = \mathbf{Z}[\alpha_i, h_j]$ 上的各点
$\alpha_i$ 处计算等式两边（略去一些细节）。
由假设，h 在环 $B[x]/(f)$ 中满足
$h^e + r_1 h^{e - 1} + \ldots + r_e = 0$，可见
$$
h(\alpha_i)^e + r_1 h(\alpha_i)^{e - 1} + \ldots + r_e = 0
$$
在 $B'$ 中成立。因此 $h(\alpha_i)$ 在 $R$ 上整。
使用上面的公式可知，在 $B'[x]/(f)$ 中有
$f'h \equiv \sum_{j = 0, \ldots, d - 1} b'_j x^j$，
其中 $b'_j \in B'$ 在 $R$ 上整。然而，因为
$f' h \in B[x]/(f)$，而且 $1, x, \ldots, x^{d - 1}$
是 $B'[x]/(f)$ 的一个 $B'$-基，所以如所需，
$b'_j \in B$。
\end{proof}

\begin{lemma}
\label{algebra-lemma-integral-closure-commutes-etale}
设 $R \to S$ 是平展环映射。
设 $R \to B$ 是任意环映射。
设 $A \subset B$ 是 $R$ 在 $B$ 中的整闭包。
设 $A' \subset S \otimes_R B$ 是 $S$ 在
$S \otimes_R B$ 中的整闭包。则典范映射
$S \otimes_R A \to A'$ 是同构。
\end{lemma}

\begin{proof}
映射 $S \otimes_R A \to A'$ 单射，因为 $A \subset B$，
且 $R \to S$ 平坦。我们将反复使用如下事实：
取整闭包与局部化交换，见引理
\ref{algebra-lemma-integral-closure-localize}。
因此，由引理 \ref{algebra-lemma-cover}（判断 $S$-模映射是否为同构
的准则），可以在 $S$ 上局部化。
于是可以假设
$S = R[x]_g/(f) = (R[x]/(f))_g$
在 $R$ 上标准平展，见命题
\ref{algebra-proposition-etale-locally-standard}。
再作一次局部化，可见 $A'$ 是 $(A'')_g$，其中
$A''$ 是 $R[x]/(f)$ 在 $B[x]/(f)$ 中的整闭包。
假设 $a \in A''$。只需证明 $a$ 属于
$S \otimes_R A$。由引理 \ref{algebra-lemma-trick}，
$f' a = \sum a_i x^i$，其中 $a_i \in A$。
因为 $f'$ 在 $S$ 中可逆（由标准平展环映射的定义），
所以如所需，$a \in S \otimes_R A$。
\end{proof}

\begin{example}
\label{algebra-example-fourier}
设 $p$ 是素数。环扩张
$$
R = \mathbf{Z}[1/p] \subset
R' = \mathbf{Z}[1/p][x]/(x^{p - 1} + \ldots + x + 1)
$$
具有如下性质：对 $d < p$，存在元素
$\alpha_0, \ldots, \alpha_{d - 1} \in R'$，使得
$$
\prod\nolimits_{0 \leq i < j < d} (\alpha_i - \alpha_j)
$$
是 $R'$ 中的单位元。具体地，对
$i = 0, \ldots, p - 1$，取 $\alpha_i$ 为
$R'$ 中 $x^i$ 的类。于是，在 $R'[T]$ 中有
$$
T^p - 1 = \prod\nolimits_{i = 0, \ldots, p - 1} (T - \alpha_i)
$$
事实上，环 $\mathbf{Q}[x]/(x^{p - 1} + \ldots + x + 1)$
是域，因为分圆多项式
$x^{p - 1} + \ldots + x + 1$ 在 $\mathbf{Q}$ 上不可约，
而 $\alpha_i$ 是 $T^p - 1$ 的两两不同的根，故有该等式。
对等式两边求导，再代入 $T = \alpha_i$，得到
$$
p \alpha_i^{p - 1}
=
(\alpha_i - \alpha_1) \ldots
\widehat{(\alpha_i - \alpha_i)} \ldots
(\alpha_i - \alpha_1)
$$
% P01-E0920: the frozen derivative-product formula repeats alpha_1 at both endpoints; the full product should run from alpha_0 through alpha_{p-1}, omitting alpha_i.
由此可见，它在 $R'$ 中可逆。
\end{example}

\begin{lemma}
\label{algebra-lemma-integral-closure-commutes-smooth}
设 $R \to S$ 是光滑环映射。
设 $R \to B$ 是任意环映射。
设 $A \subset B$ 是 $R$ 在 $B$ 中的整闭包。
设 $A' \subset S \otimes_R B$ 是 $S$ 在
$S \otimes_R B$ 中的整闭包。则典范映射
$S \otimes_R A \to A'$ 是同构。
\end{lemma}

\begin{proof}
像引理 \ref{algebra-lemma-integral-closure-commutes-etale}
的证明那样论证，可以在 $S$ 上局部化。因此可以假设
$R \to S$ 是标准光滑环映射，见引理
\ref{algebra-lemma-smooth-syntomic}。由标准光滑环映射的定义，
$S$ 在多项式环 $R[x_1, \ldots, x_n]$ 上平展。
由于已经对平展环扩张证明了该结果（引理
\ref{algebra-lemma-integral-closure-commutes-etale}），
问题归结为 $S = R[x]$ 的情形。因此需要证明
$$
f = \sum b_i x^i
\text{ integral over }R[x]
\Leftrightarrow
\text{每个 }b_i\text{ integral over }R.
$$
从右向左的蕴含成立，因为 $B[x]$ 中在 $R[x]$ 上整的元素之集
是一个环（引理 \ref{algebra-lemma-integral-closure-is-ring}），
并且包含 $x$。

\medskip\noindent
假设 $f \in B[x]$ 在 $R[x]$ 上整，并且
$f = \sum_{i < d} b_i x^i$ 的次数 $< d$。
因为整闭包与局部化交换，只需证明存在不同的素数
$p, q$，使得每个 $b_i$ 既在 $R[1/p]$ 上整，
又在 $R[1/q]$ 上整。因此，可以找到有限自由环扩张
$R \subset R'$，使得 $R'$ 含有
$\alpha_1, \ldots, \alpha_d$，并且
$\prod_{i < j} (\alpha_i - \alpha_j)$ 是 $R'$ 中的单位元，
见例 \ref{algebra-example-fourier}。
% P01-E0921: the frozen proof invokes Example Fourier over R without explicitly replacing R by one of the localizations R[1/p] or R[1/q], although invertibility of the chosen integer prime is required.
在这种情形下，有普遍等式
$$
f
=
\sum_i
f(\alpha_i)
\frac{(x - \alpha_1) \ldots \widehat{(x - \alpha_i)} \ldots (x - \alpha_d)}
{(\alpha_i - \alpha_1) \ldots \widehat{(\alpha_i - \alpha_i)} \ldots
(\alpha_i - \alpha_d)}.
$$
此外，元素 $f(\alpha_i)$ 在 $R'$ 上整，因为
$(R' \otimes_R B)[x] \to R' \otimes_R B$，
$h \mapsto h(\alpha_i)$ 是环映射。
% P01-E0922: the frozen source begins this sentence with the editorial artifact “OK,”.
因此可见，$f$ 在 $(R' \otimes_R B)[x]$ 中的系数
在 $R'$ 上整。由于 $R'$ 在 $R$ 上有限（因而在 $R$ 上整），
可知这些系数也在 $R$ 上整，正如所需。
\end{proof}

\begin{lemma}
\label{algebra-lemma-integral-closure-commutes-colim-smooth}
设 $R \to S$ 和 $R \to B$ 是环映射。
设 $A \subset B$ 是 $R$ 在 $B$ 中的整闭包。
设 $A' \subset S \otimes_R B$ 是 $S$ 在
$S \otimes_R B$ 中的整闭包。若 $S$ 是光滑
$R$-代数的滤过余极限，则典范映射
$S \otimes_R A \to A'$ 是同构。
\end{lemma}

\begin{proof}
这由如下直接事实以及引理
\ref{algebra-lemma-integral-closure-commutes-smooth} 推出：
取张量积和取整闭包都与滤过余极限交换。
% P01-E0923: the frozen source has compound-subject agreement “taking ... and taking ... commutes”.
\end{proof}

\section{形式非分歧映射}
\label{algebra-section-formally-unramified}

\noindent
事实表明，在引入形式平展映射的概念之前先定义形式非分歧映射的概念，
在逻辑上更为高效。

\begin{definition}
\label{algebra-definition-formally-unramified}
设 $R \to S$ 是环映射。若对于每个由实线箭头组成的交换图
$$
\xymatrix{
S \ar[r] \ar@{-->}[rd] & A/I \\
R \ar[r] \ar[u] & A \ar[u]
}
$$
其中 $I \subset A$ 是平方为零的理想，使图交换的虚线箭头至多存在一条，
则称 $S$ 在 $R$ 上\textit{形式非分歧}。
\end{definition}

\begin{lemma}
\label{algebra-lemma-base-change-formally-unramified}
设 $R \to S$ 是形式非分歧映射，$R \to R'$ 是任意环映射。
则基变换 $S' = R' \otimes_R S$ 在 $R'$ 上形式非分歧。
\end{lemma}

\begin{proof}
给定定义 \ref{algebra-definition-formally-unramified} 中如下实线图
$$
\xymatrix{
S \ar[r] \ar@{-->}[rrd] & R' \otimes_R S \ar[r] \ar@{-->}[rd] & A/I \\
R  \ar[u] \ar[r] & R' \ar[r] \ar[u] & A \ar[u]
}
$$
依假设，较长的虚线箭头至多存在一条。
由张量积的泛性质，较短的虚线箭头也至多存在一条。
\end{proof}

\begin{lemma}
\label{algebra-lemma-characterize-formally-unramified}
设 $R \to S$ 是环映射。下列条件等价：
\begin{enumerate}
\item $R \to S$ 形式非分歧；
\item 微分模 $\Omega_{S/R}$ 为零。
\end{enumerate}
\end{lemma}

\begin{proof}
设 $J = \Ker(S \otimes_R S \to S)$ 是乘法映射的核，
并设 $A_{univ} = S \otimes_R S/J^2$。回忆
$I_{univ} = J/J^2$ 同构于 $\Omega_{S/R}$，见引理
\ref{algebra-lemma-differentials-diagonal}。此外，两个 $R$-代数映射
$\sigma_1, \sigma_2 : S \to A_{univ}$，其中
$\sigma_1(s) = s \otimes 1 \bmod J^2$ 且
$\sigma_2(s) = 1 \otimes s \bmod J^2$，相差泛导子
$\text{d} : S \to \Omega_{S/R} = I_{univ}$。

\medskip\noindent
假设 $R \to S$ 形式非分歧。于是 $\sigma_1 = \sigma_2$，
从而对所有 $s \in S$ 都有 $\text{d}(s) = 0$。
故 $\Omega_{S/R} = 0$。

\medskip\noindent
假设 $\Omega_{S/R} = 0$。设 $A, I, R \to A, S \to A/I$
构成定义 \ref{algebra-definition-formally-unramified} 中的实线图。
设 $\tau_1, \tau_2 : S \to A$ 是两个使图交换的虚线箭头。
考察由规则 $s_1 \otimes s_2 \mapsto \tau_1(s_1)\tau_2(s_2)$
定义的 $R$-代数映射 $A_{univ} \to A$。
这里省略验证该映射良定义。由于
$I_{univ} = \Omega_{S/R} = 0$，故 $A_{univ} \cong S$，
从而 $\tau_1 = \tau_2$。
\end{proof}

\begin{lemma}
\label{algebra-lemma-formally-unramified-local}
设 $R \to S$ 是环映射。下列条件等价：
\begin{enumerate}
\item $R \to S$ 形式非分歧；
\item 对 $S$ 的所有素理想 $\mathfrak q$，
$R \to S_{\mathfrak q}$ 形式非分歧；
\item 对 $S$ 的所有素理想 $\mathfrak q$，令
$\mathfrak p = R \cap \mathfrak q$，则
$R_{\mathfrak p} \to S_{\mathfrak q}$ 形式非分歧。
\end{enumerate}
\end{lemma}

\begin{proof}
引理 \ref{algebra-lemma-characterize-formally-unramified} 已表明，
(1) 等价于 $\Omega_{S/R} = 0$。类似地，由引理
\ref{algebra-lemma-differentials-localize} 可知，(2) 和 (3) 都等价于：
对所有 $\mathfrak q$ 都有
$(\Omega_{S/R})_{\mathfrak q} = 0$。
因此，该等价性由引理 \ref{algebra-lemma-characterize-zero-local} 推出。
\end{proof}

\begin{lemma}
\label{algebra-lemma-formally-unramified-localize}
设 $A \to B$ 是形式非分歧环映射。
\begin{enumerate}
\item 若 $S \subset A$ 是乘法子集，则
$S^{-1}A \to S^{-1}B$ 形式非分歧。
\item 若 $S \subset B$ 是乘法子集，则
$A \to S^{-1}B$ 形式非分歧。
\end{enumerate}
\end{lemma}

\begin{proof}
这由引理 \ref{algebra-lemma-formally-unramified-local} 推出。
% P01-E0924: the frozen proof begins with the sentence fragment “Follows from”.
（也可以由引理 \ref{algebra-lemma-characterize-formally-unramified}
结合引理 \ref{algebra-lemma-differentials-localize} 推出。）
\end{proof}

\begin{lemma}
\label{algebra-lemma-colimit-formally-unramified}
设 $R$ 是环，$I$ 是有向集，并设
$(S_i, \varphi_{ii'})$ 是 $I$ 上的 $R$-代数系。
若每个 $R \to S_i$ 都形式非分歧，则
$S = \colim_{i \in I} S_i$ 在 $R$ 上形式非分歧。
% P01-E0925: the frozen lemma conclusion lacks a terminal period.
\end{lemma}

\begin{proof}
考察定义 \ref{algebra-definition-formally-unramified} 中的图。
依假设，提升复合映射 $S_i \to S \to A/I$ 的
$R$-代数映射 $S_i \to A$ 至多存在一个。
由于 $S$ 的每个元素都属于某个映射 $S_i \to S$ 的像，
可知嵌入该图的映射 $S \to A$ 至多存在一个。
\end{proof}

\section{余法模与泛加厚}
\label{algebra-section-conormal}

\noindent
事实表明，第一无穷小邻域不仅可对概形的闭浸入定义，
而且已经可对任意形式非分歧态射定义。其基础是如下代数事实。

\begin{lemma}
\label{algebra-lemma-universal-thickening}
设 $R \to S$ 是形式非分歧环映射。存在 $R$-代数满射
$S' \to S$，其核为平方为零的理想，并具有如下泛性质：
给定任意交换图
$$
\xymatrix{
S \ar[r]_a & A/I \\
R \ar[r]^b \ar[u] & A \ar[u]
}
$$
其中 $I \subset A$ 是平方为零的理想，则存在唯一的
$R$-代数映射 $a' : S' \to A$，使得
$S' \to A \to A/I$ 等于 $S' \to S \to A/I$。
\end{lemma}

\begin{proof}
选取 $S$ 作为 $R$-代数的一族生成元 $z_i \in S$，$i \in I$。
% P01-E0926: the frozen proof reuses I first as this generator index set and later as the square-zero ideal in an arbitrary diagram.
以 $P = R[\{x_i\}_{i \in I}]$ 表示以 $x_i$（$i \in I$）
为生成元的多项式环。考察把 $x_i$ 映到 $z_i$ 的
$R$-代数映射 $P \to S$，并令 $J = \Ker(P \to S)$。
考察映射
$$
\text{d} : J/J^2 \longrightarrow \Omega_{P/R} \otimes_P S
$$
见引理 \ref{algebra-lemma-differential-seq}。
由于依假设 $\Omega_{S/R} = 0$，该映射是满射，见引理
\ref{algebra-lemma-characterize-formally-unramified}。
注意 $\Omega_{P/R}$ 是以 $\text{d}x_i$ 为基的自由模，
故 $\Omega_{P/R} \otimes_P S$ 是自由 $S$-模。
因此可选取上述满射的一个分裂，并写成
$$
J/J^2 = K \oplus \Omega_{P/R} \otimes_P S
$$
% P01-E0927: the frozen decomposition display lacks terminal punctuation before the next sentence.
令 $J^2 \subset J' \subset J$ 为 $P$ 的理想，使得
$J'/J^2$ 是上述分解中的第二个直和项。置 $S' = P/J'$。
得到短正合列
$$
0 \to J/J' \to S' \to S \to 0
$$
并且 $J/J' \cong K$ 是 $S'$ 中平方为零的理想。因此
$$
\xymatrix{
S \ar[r]_1 & S \\
R \ar[r] \ar[u] & S' \ar[u]
}
$$
是上述类型的图。事实上，我们断言它是这些图所成范畴中的初始对象。
确切地，设 $(I \subset A, a, b)$ 是任意一个这样的图。
可以选取 $R$-代数映射 $\beta : P \to A$，使图
$$
\xymatrix{
S \ar[r]_1 & S \ar[r]_a & A/I \\
R \ar[r] \ar@/_/[rr]_b \ar[u] & P \ar[u] \ar[r]^\beta & A \ar[u]
}
$$
交换。此时未必有 $\beta(J') = 0$；换言之，$\beta$ 未必通过
$S' = P/J'$ 分解。但显然 $\beta(J') \subset I$，而且由于
$\beta(J) \subset I$ 且 $I^2 = 0$，有 $\beta(J^2) = 0$。
因此，从 $(J/J' \subset S', 1, R \to S')$ 到
$(I \subset A, a, b)$ 的态射存在的“阻碍”，就是相应的
$S$-线性映射 $\overline{\beta} : J'/J^2 \to I$。
选取 $\beta$ 的自由度在于选取 $\beta(x_i)$。
若对 $\beta$ 作另一选择，记作 $\beta'$，则它由
$\beta'(x_i) = \beta(x_i) + \delta_i$ 给出，其中 $\delta_i \in I$。
在这种情形下，对 $g \in J'$ 有
$$
\beta'(g) =
\beta(g) + \sum\nolimits_i \delta_i \beta(\frac{\partial g}{\partial x_i})
$$
% P01-E0928: the frozen display lacks terminal punctuation before the following sentence.
由于按构造，映射
$\text{d}|_{J'/J^2} : J'/J^2 \to \Omega_{P/R} \otimes_P S$
由 $g \mapsto \frac{\partial g}{\partial x_i}\text{d}x_i \otimes 1$
给出且是同构，可知存在唯一的一族 $\delta_i \in I$，使得对所有
$g \in J'$ 都有 $\beta'(g) = 0$。
% P01-E0929: the frozen differential formula omits the summation over i.
（确切地，$\delta_i$ 是 $-\overline{\beta}(g)$，其中
$g \in J'/J^2$ 是在 $\text{d}|_{J'/J^2}$ 下映到
$\text{d}x_i \otimes 1$ 的唯一元素。）
% P01-E0930: the frozen parenthetical uses an unindexed g for the family defining each delta_i; g_i would make the dependence explicit.
解的唯一性蕴含引理所要求的唯一性。
\end{proof}

\noindent
在引理 \ref{algebra-lemma-universal-thickening} 的情形下，
$R$-代数映射 $S' \to S$ 在唯一同构意义下唯一。

\begin{definition}
\label{algebra-definition-universal-thickening}
设 $R \to S$ 是形式非分歧环映射。
\begin{enumerate}
\item $S$ 在 $R$ 上的\textit{泛第一阶加厚}，是引理
\ref{algebra-lemma-universal-thickening} 中的 $R$-代数满射 $S' \to S$。
\item $R \to S$ 的\textit{余法模}，是泛第一阶加厚
$S' \to S$ 的核 $I$，视为 $S$-模。
\end{enumerate}
在这种情形下，我们常以 \textit{$C_{S/R}$} 表示余法模。
\end{definition}

\begin{lemma}
\label{algebra-lemma-universal-thickening-quotient}
设 $I \subset R$ 是环的理想。$R/I$ 在 $R$ 上的泛第一阶加厚是满射
$R/I^2 \to R/I$。$R/I$ 在 $R$ 上的余法模为
$C_{(R/I)/R} = I/I^2$。
\end{lemma}

\begin{proof}
略。
\end{proof}

\begin{lemma}
\label{algebra-lemma-universal-thickening-localize}
设 $A \to B$ 是形式非分歧环映射，并设
$\varphi : B' \to B$ 是 $B$ 在 $A$ 上的泛第一阶加厚。
\begin{enumerate}
\item 设 $S \subset A$ 是乘法子集。则
$S^{-1}B' \to S^{-1}B$ 是 $S^{-1}B$ 在 $S^{-1}A$ 上的
泛第一阶加厚。特别地，
$S^{-1}C_{B/A} = C_{S^{-1}B/S^{-1}A}$。
\item 设 $S \subset B$ 是乘法子集。则
$S' = \varphi^{-1}(S)$ 是 $B'$ 中的乘法子集，并且
$(S')^{-1}B' \to S^{-1}B$ 是 $S^{-1}B$ 在 $A$ 上的
泛第一阶加厚。特别地，$S^{-1}C_{B/A} = C_{S^{-1}B/A}$。
\end{enumerate}
由引理 \ref{algebra-lemma-formally-unramified-localize}，该引理中的说法确有意义。
\end{lemma}

\begin{proof}
采用 (1) 中的记号和假设。设 $(S^{-1}B)' \to S^{-1}B$ 是
$S^{-1}B$ 在 $S^{-1}A$ 上的泛第一阶加厚。
% P01-E0931: both proof parts in the frozen source begin with the sentence fragment “With notation and assumptions as in ...”.
注意 $S^{-1}B' \to S^{-1}B$ 是 $S^{-1}A$-代数的满射，
且其核平方为零。因此按定义得到映射
$(S^{-1}B)' \to S^{-1}B'$，它与到 $S^{-1}B$ 的映射相容。
考察任意交换图
$$
\xymatrix{
B \ar[r] & S^{-1}B \ar[r] & D/I \\
A \ar[r] \ar[u] & S^{-1}A \ar[r] \ar[u] & D \ar[u]
}
$$
其中 $I \subset D$ 是平方为零的理想。由于 $B'$ 是 $B$ 在 $A$ 上的
泛第一阶加厚，得到 $A$-代数映射 $B' \to D$。
但显然，$S$ 在 $D$ 中的像由此映到 $D$ 的可逆元素，
因而得到相容的映射 $S^{-1}B' \to D$。
取 $D = (S^{-1}B)'$，可知得到映射
$S^{-1}B' \to (S^{-1}B)'$。这里省略验证该映射是上述映射的逆。

\medskip\noindent
采用 (2) 中的记号和假设。设 $(S^{-1}B)' \to S^{-1}B$
是 $S^{-1}B$ 在 $A$ 上的泛第一阶加厚。
注意 $(S')^{-1}B' \to S^{-1}B$ 是 $A$-代数的满射，
且其核平方为零。因此按定义得到映射
$(S^{-1}B)' \to (S')^{-1}B'$，它与到 $S^{-1}B$ 的映射相容。
考察任意交换图
$$
\xymatrix{
B \ar[r] & S^{-1}B \ar[r] & D/I \\
A \ar[r] \ar[u] & A \ar[r] \ar[u] & D \ar[u]
}
$$
其中 $I \subset D$ 是平方为零的理想。由于 $B'$ 是 $B$ 在 $A$ 上的
泛第一阶加厚，得到 $A$-代数映射 $B' \to D$。
但显然，$S'$ 在 $D$ 中的像由此映到 $D$ 的可逆元素，
因而得到相容的映射 $(S')^{-1}B' \to D$。
取 $D = (S^{-1}B)'$，可知得到映射
$(S')^{-1}B' \to (S^{-1}B)'$。这里省略验证该映射是上述映射的逆。
\end{proof}

\begin{lemma}
\label{algebra-lemma-differentials-universal-thickening}
设 $R \to A \to B$ 是环映射，并假设 $A \to B$ 形式非分歧。
设 $B' \to B$ 是 $B$ 在 $A$ 上的泛第一阶加厚。
则 $B'$ 在 $A$ 上形式非分歧，并且典范映射
$\Omega_{A/R} \otimes_A B \to \Omega_{B'/R} \otimes_{B'} B$
是同构。
\end{lemma}

\begin{proof}
我们将使用引理 \ref{algebra-lemma-universal-thickening} 的证明中对 $B'$ 的构造，
尽管原则上应该能从定义形式地推出这些结果。确切地，选取表现
$B = P/J$，其中 $P = A[x_i]$ 是 $A$ 上的多项式环。
然后选取元素 $f_i \in J$，使得在
$\Omega_{P/A} \otimes_P B$ 中有
$\text{d}f_i = \text{d}x_i \otimes 1$。作出这些选择后，
$B' = P/J'$，其中 $J' = (f_i) + J^2$，见引理
\ref{algebra-lemma-universal-thickening} 的证明。

\medskip\noindent
考察典范正合列
$$
J'/(J')^2 \to \Omega_{P/A} \otimes_P B' \to \Omega_{B'/A} \to 0
$$
见引理 \ref{algebra-lemma-differential-seq}。
按构造，$f_i \in J'$ 的类在模
$\Omega_{P/A} \otimes_P B'$ 中的像模掉 $J'/J^2$ 后生成该模。
由于 $J'/J^2$ 是幂零理想，可知这些元素生成整个模（由
Nakayama 引理 \ref{algebra-lemma-NAK}）。
% P01-E0932: the frozen source uses J'/J^2 as the modulus and calls it a nilpotent ideal of B'; the relevant square-zero kernel of B' -> B is J/J'.
这证明 $\Omega_{B'/A} = 0$，因而 $B'$ 在 $A$ 上形式非分歧，
见引理 \ref{algebra-lemma-characterize-formally-unramified}。

\medskip\noindent
由于 $P$ 是 $A$ 上的多项式环，有
$\Omega_{P/R} = \Omega_{A/R} \otimes_A P \oplus \bigoplus P\text{d}x_i$。
我们将使用这一分解。考察如下正合列
$$
J'/(J')^2 \to
\Omega_{P/R} \otimes_P B' \to
\Omega_{B'/R} \to 0
$$
见引理 \ref{algebra-lemma-differential-seq}。
可以将其与 $B$ 作张量积，得到正合列
$$
J'/(J')^2 \otimes_{B'} B \to
\Omega_{P/R} \otimes_P B \to
\Omega_{B'/R} \otimes_{B'} B \to 0
$$
% P01-E0933: the frozen exact-sequence display lacks terminal punctuation before the following sentence.
回忆 $J' = (f_i) + J^2$，可知第一条箭头把子模
$J^2/(J')^2$ 化为零。就给出的直和分解
$\Omega_{P/R} \otimes_P B =
\Omega_{A/R} \otimes_A B \oplus \bigoplus B\text{d}x_i $
而言，可知子模 $(f_i)/(J')^2 \otimes_{B'} B$ 同构地映到直和项
$\bigoplus B\text{d}x_i$ 上。
% P01-E0934: the frozen phrase “the direct sum decomposition ... given we see” is missing “above” and punctuation after “given”.
因此，该正合列剩下的部分就是同构
$\Omega_{A/R} \otimes_A B \to \Omega_{B'/R} \otimes_{B'} B$，
正如所需。
\end{proof}

\section{形式平展映射}
\label{algebra-section-formally-etale}

\begin{definition}
\label{algebra-definition-formally-etale}
设 $R \to S$ 是环映射。若对每个由实线箭头组成的交换图
$$
\xymatrix{
S \ar[r] \ar@{-->}[rd] & A/I \\
R \ar[r] \ar[u] & A \ar[u]
}
$$
其中 $I \subset A$ 是平方为零的理想，都存在唯一一条使图交换的虚线箭头，
则称 $S$ \textit{在 $R$ 上形式平展}。
\end{definition}

\noindent
显然，环映射形式平展，当且仅当它既形式光滑又形式非分歧。

\begin{lemma}
\label{algebra-lemma-base-change-formally-etale}
设 $R \to S$ 是形式平展映射，$R \to R'$ 是任意环映射。
则基变换 $S' = R' \otimes_R S$ 在 $R'$ 上形式平展。
\end{lemma}

\begin{proof}
结合引理 \ref{algebra-lemma-base-change-fs} 和
\ref{algebra-lemma-base-change-formally-unramified}。
\end{proof}

\begin{lemma}
\label{algebra-lemma-formally-etale-etale}
设 $R \to S$ 是有限表现环映射。下列条件等价：
\begin{enumerate}
\item $R \to S$ 形式平展；
\item $R \to S$ 平展。
\end{enumerate}
\end{lemma}

\begin{proof}
假设 $R \to S$ 形式平展。由命题
\ref{algebra-proposition-smooth-formally-smooth}，$R \to S$ 光滑。
由引理 \ref{algebra-lemma-characterize-formally-unramified}，有
$\Omega_{S/R} = 0$。故按定义，$R \to S$ 平展。

\medskip\noindent
假设 $R \to S$ 平展。由命题
\ref{algebra-proposition-smooth-formally-smooth}，$R \to S$ 形式光滑。
由引理 \ref{algebra-lemma-characterize-formally-unramified}，它形式非分歧。
因此 $R \to S$ 形式平展。
\end{proof}

\begin{lemma}
\label{algebra-lemma-colimit-formally-etale}
设 $R$ 是环，$I$ 是有向集，并设
$(S_i, \varphi_{ii'})$ 是 $I$ 上的 $R$-代数系。
若每个 $R \to S_i$ 都形式平展，则
$S = \colim_{i \in I} S_i$ 在 $R$ 上形式平展。
% P01-E0935: the frozen lemma conclusion lacks terminal punctuation.
\end{lemma}

\begin{proof}
考察定义 \ref{algebra-definition-formally-etale} 中的图。
依假设，存在唯一的一族 $R$-代数映射 $S_i \to A$
提升复合映射 $S_i \to S \to A/I$。
% P01-E0936: the frozen source says “we get unique ... maps” without a determiner or explicit per-index quantifier.
因此这些映射与转移映射 $\varphi_{ii'}$ 相容，并定义提升
$S \to A$。这证明了存在性。把映射限制到每个 $S_i$ 上，
唯一性便显然成立。
\end{proof}

\begin{lemma}
\label{algebra-lemma-localization-formally-etale}
设 $R$ 是环，$S \subset R$ 是任意乘法子集。
则环映射 $R \to S^{-1}R$ 形式平展。
\end{lemma}

\begin{proof}
设 $I \subset A$ 是平方为零的理想。这里所说的是：
给定环映射 $\varphi : R \to A$，若对所有 $f \in S$，
$\varphi(f) \mod I$ 都可逆，则对所有 $f \in S$，
$\varphi(f)$ 在 $A$ 中也可逆。这是因为，在典范映射
$A \to A/I$ 下，$A^*$ 是 $(A/I)^*$ 的逆像。
\end{proof}

\begin{lemma}
\label{algebra-lemma-formally-etale-lift-infinitesimal}
设 $R \to S$ 是环映射。设 $J \subset S$ 是理想，使得
$R \to S/J$ 为满射；令 $I \subset R$ 为其核。
若 $R \to S$ 形式平展，则对所有 $n$，
$R/I^n \to S/J^n$ 都是同构，并且
$\bigoplus I^n/I^{n + 1} \to \bigoplus J^n/J^{n + 1}$
是分次环的同构。
\end{lemma}

\begin{proof}
归纳使用提升性质，得到虚线箭头
$$
\xymatrix{
S \ar[r] \ar@{-->}[rd] & S/J = R/I \\
R \ar[r] \ar[u] & R/I^2 \ar[u]
}
\quad
\xymatrix{
S \ar[r] \ar@{-->}[rd] & R/I^2 \\
R \ar[r] \ar[u] & R/I^3 \ar[u]
}
\quad
\xymatrix{
S \ar[r] \ar@{-->}[rd] & R/I^3 \\
R \ar[r] \ar[u] & R/I^4 \ar[u]
}
$$
% P01-E0937: the frozen three-diagram display lacks terminal punctuation.
相应的映射 $S/J^n \to R/I^n$ 是同构，因为复合映射
$S/J^n \to R/I^n \to S/J^n$ 由形式平展环映射提升性质中的
唯一性可（归纳地）知为恒等映射。
\end{proof}

\begin{lemma}
\label{algebra-lemma-formally-etale-omega}
设 $R \to S \to S'$ 是环映射。令 $J$ 和 $J'$ 分别为乘法映射
$S \otimes_R S \to S$ 和 $S' \otimes_{R'} S' \to S'$ 的核。
% P01-E0938: R' is undefined in the frozen source; the second tensor product should be over R.
若 $S \to S'$ 形式平展，则对所有 $k \geq 0$，映射
$$
S' \otimes_S \left((S \otimes_R S)/J^{k + 1}\right)
\longrightarrow
(S' \otimes_R S')/(J')^{k + 1}
$$
都是同构。特别地，映射
$S' \otimes_S \Omega_{S/R} \to \Omega_{S'/R}$ 是同构。
\end{lemma}

\begin{proof}
注意 $S' \otimes_S (S \otimes_R S) = S' \otimes_R S$，
而 $J$ 在 $S' \otimes_R S$ 中生成满射
$S' \otimes_R S \to S'$ 的核 $I$。因此左边等于
$S' \otimes_R S/I^{k + 1}$。由引理
\ref{algebra-lemma-base-change-formally-etale}，映射
$S' \otimes_R S \to S' \otimes_R S'$ 形式平展。
于是由引理 \ref{algebra-lemma-formally-etale-lift-infinitesimal}，
引理陈述中的显示箭头是同构。最后一个断言由此以及如下事实推出：
由引理 \ref{algebra-lemma-differentials-diagonal}，
$\Omega_{S/R} = J/J^2$ 且 $\Omega_{S'/R} = J'/(J')^2$。
\end{proof}

\begin{lemma}
\label{algebra-lemma-formally-etale-principal-parts}
设 $R \to S \to S'$ 是环映射，且 $S \to S'$ 形式平展
（例如平展）。设 $M$ 是 $S$-模。置 $M' = S' \otimes_R M$。则
% P01-E0939: the frozen statement tensors M over R, while the lemma and its proof require M' = S' tensor_S M.
$$
S' \otimes_S P^k_{S/R}(M) = P^k_{S'/R}(M')
$$
% P01-E0940: the frozen display lacks terminal punctuation.
由此，对任意 $S$-模 $N$ 以及任意有限阶微分算子
$D : M \to N$，都存在唯一的 $D$ 到同阶或更低阶微分算子的延拓
$D' : M' \to S' \otimes_S N$。
\end{lemma}

\begin{proof}
令 $J$ 和 $J'$ 如引理 \ref{algebra-lemma-formally-etale-omega} 的陈述中所设。
于是
\begin{align*}
S' \otimes_S P^k_{S/R}(M)
& =
S' \otimes_S \left((S \otimes_R M)/J^{k + 1}(S \otimes_R M)\right) \\
& =
S' \otimes_S \left((S \otimes_R S)/J^{k + 1}\right) \otimes_S M \\
& =
\left(S' \otimes_R S'/(J')^{k + 1}\right) \otimes_S M \\
& =
\left(S' \otimes_R S'/(J')^{k + 1}\right) \otimes_{S'} M' \\
& =
S' \otimes_R M'/(J')^{k + 1}(S' \otimes_R M') \\
& =
P^k_{S'/R'}(M')
\end{align*}
% P01-E0941: the frozen final term uses undefined R' instead of R.
第一个和最后一个等式来自引理
\ref{algebra-lemma-principal-parts-diagonal}，第三个等式是引理
\ref{algebra-lemma-formally-etale-omega}。最后一个断言成立，是因为若
$D$ 对应于线性映射 $\gamma : P^k_{S/R}(M) \to N$，
则可令 $D' : M' \to S' \otimes_R N$ 对应于线性映射
$1 \otimes \gamma : S' \otimes_R P^k_{S/R}(M) \to S' \otimes_R N$。
% P01-E0942: the frozen proof tensors both N and the principal-parts module over R, contradicting the stated S-base-change targets.
\end{proof}

\begin{remark}
\label{algebra-remark-formally-etale-differential-operators}
设 $R \to S \to S'$ 是环映射，且 $S \to S'$ 形式平展
（例如平展）。设 $M_i$（$i = 1, 2, 3$）是 $S$-模，并设
$D_i : M_i \to M_{i + 1}$（$i = 1, 2$）是有限阶微分算子。
若 $D'_i : M'_i \to M'_{i + 1}$（$i = 1, 2$）是按引理
\ref{algebra-lemma-formally-etale-principal-parts} 将 $D_i$ 延拓到
$M'_i = S' \otimes_S M_i$ 所得的算子，则
$D'_2 \circ D'_1$ 是 $D_2 \circ D_1$ 的延拓。
特别地，若 $M$ 是 $S$-模，则
$M' = S' \otimes_S M$ 是 $S$-代数
$\text{Diff}_{S/R}(M, M)$ 上的模。
\end{remark}

\section{非分歧环映射}
\label{algebra-section-unramified}

\noindent
我们对非分歧环映射所采用的定义来自 \cite{Henselian}。
我们所谓的 G-非分歧环映射，就是 EGA 所称的非分歧映射。
% P01-E0943: the frozen source has the incorrect article “an G-unramified ring map”.

\begin{definition}
\label{algebra-definition-unramified}
设 $R \to S$ 是环映射。
\begin{enumerate}
\item 若 $R \to S$ 为有限型且 $\Omega_{S/R} = 0$，
则称 $R \to S$ \textit{非分歧}。
\item 若 $R \to S$ 为有限表现且 $\Omega_{S/R} = 0$，
则称 $R \to S$ \textit{G-非分歧}。
\item 给定 $S$ 的素理想 $\mathfrak q$，若存在
$g \in S$、$g \not \in \mathfrak q$，使得 $R \to S_g$ 非分歧，
则称 $S$ \textit{在 $\mathfrak q$ 处非分歧}。
\item 给定 $S$ 的素理想 $\mathfrak q$，若存在
$g \in S$、$g \not \in \mathfrak q$，使得 $R \to S_g$ G-非分歧，
则称 $S$ \textit{在 $\mathfrak q$ 处 G-非分歧}。
\end{enumerate}
\end{definition}

\noindent
当然，G-非分歧映射必非分歧。

\begin{lemma}
\label{algebra-lemma-formally-unramified-unramified}
设 $R \to S$ 是环映射。下列条件等价：
\begin{enumerate}
\item $R \to S$ 形式非分歧且为有限型；
\item $R \to S$ 非分歧。
\end{enumerate}
此外，下列条件也等价：
\begin{enumerate}
\item $R \to S$ 形式非分歧且为有限表现；
\item $R \to S$ G-非分歧。
\end{enumerate}
\end{lemma}

\begin{proof}
由引理 \ref{algebra-lemma-characterize-formally-unramified} 和定义立即得到。
\end{proof}

\begin{lemma}
\label{algebra-lemma-unramified}
非分歧与 G-非分歧环映射具有下列性质。
\begin{enumerate}
\item 非分歧环映射的基变换非分歧。
G-非分歧环映射的基变换 G-非分歧。
\item 非分歧环映射的复合非分歧。
G-非分歧环映射的复合 G-非分歧。
\item 任意主局部化 $R \to R_f$ 都既 G-非分歧又非分歧。
\item 若 $I \subset R$ 是理想，则 $R \to R/I$ 非分歧。
若 $I \subset R$ 是有限生成理想，则 $R \to R/I$ G-非分歧。
\item 平展环映射既 G-非分歧又非分歧。
\item 若 $R \to S$ 为有限型（相应地有限表现），
$\mathfrak q \subset S$ 为素理想且 $(\Omega_{S/R})_{\mathfrak q} = 0$，
则 $R \to S$ 在 $\mathfrak q$ 处非分歧（相应地 G-非分歧）。
\item 若 $R \to S$ 为有限型（相应地有限表现），
$\mathfrak q \subset S$ 为素理想且
$\Omega_{S/R} \otimes_S \kappa(\mathfrak q) = 0$，则
$R \to S$ 在 $\mathfrak q$ 处非分歧（相应地 G-非分歧）。
\item 若 $R \to S$ 为有限型（相应地有限表现），
$\mathfrak q \subset S$ 是位于 $\mathfrak p \subset R$ 之上的素理想，且
$(\Omega_{S \otimes_R \kappa(\mathfrak p)/\kappa(\mathfrak p)})_{\mathfrak q}
= 0$，则 $R \to S$ 在 $\mathfrak q$ 处非分歧（相应地 G-非分歧）。
\item 若 $R \to S$ 为有限型（相应地有限表现），
% P01-E0944: the frozen source says “resp. presentation”, omitting “finite”.
$\mathfrak q \subset S$ 是位于 $\mathfrak p \subset R$ 之上的素理想，且
$(\Omega_{S \otimes_R \kappa(\mathfrak p)/\kappa(\mathfrak p)})
\otimes_{S \otimes_R \kappa(\mathfrak p)} \kappa(\mathfrak q) = 0$，
则 $R \to S$ 在 $\mathfrak q$ 处非分歧（相应地 G-非分歧）。
\item 若 $R \to S$ 是环映射，$g_1, \ldots, g_m \in S$ 生成单位理想，
并且对 $j = 1, \ldots, m$，$R \to S_{g_j}$ 都非分歧
（相应地 G-非分歧），则 $R \to S$ 非分歧（相应地 G-非分歧）。
\item 若环映射 $R \to S$ 在 $S$ 的每个素理想处都非分歧
（相应地 G-非分歧），则 $R \to S$ 非分歧（相应地 G-非分歧）。
\item 若 $R \to S$ G-非分歧，则存在有限型
$\mathbf{Z}$-代数 $R_0$、G-非分歧环映射 $R_0 \to S_0$，
以及环映射 $R_0 \to R$，使得 $S = R \otimes_{R_0} S_0$。
\item 若 $R \to S$ 非分歧，则存在有限型
$\mathbf{Z}$-代数 $R_0$、非分歧环映射 $R_0 \to S_0$，
以及环映射 $R_0 \to R$，使得 $S$ 是
$R \otimes_{R_0} S_0$ 的商。
\end{enumerate}
\end{lemma}

\begin{proof}
依次证明各项。

\medskip\noindent
对 (1)。由引理 \ref{algebra-lemma-differentials-base-change}
和 \ref{algebra-lemma-base-change-finiteness} 得到。

\medskip\noindent
对 (2)。由引理 \ref{algebra-lemma-exact-sequence-differentials}
和 \ref{algebra-lemma-base-change-finiteness} 得到。
% P01-E0945: the frozen proof of composition cites the base-change finiteness lemma; finite-type/finite-presentation composition is lemma-compose-finite-type.

\medskip\noindent
对 (3)。直接计算 $\Omega_{R_f/R}$ 即得；计算从略。

\medskip\noindent
对 (4)。由引理 \ref{algebra-lemma-trivial-differential-surjective}，有
$\Omega_{(R/I)/R} = 0$，且环映射 $R \to R/I$ 为有限型。
若 $I$ 是有限生成理想，则 $R \to R/I$ 为有限表现。

\medskip\noindent
对 (5)。见定义 \ref{algebra-definition-etale} 后的讨论。

\medskip\noindent
对 (6)。此时 $\Omega_{S/R}$ 是有限 $S$-模（见引理
\ref{algebra-lemma-differentials-finitely-generated}），故存在
$g \in S$、$g \not \in \mathfrak q$，使得
$(\Omega_{S/R})_g = 0$。由引理 \ref{algebra-lemma-differentials-localize}，
这意味着 $\Omega_{S_g/R} = 0$，因而 $R \to S_g$ 如所需地非分歧。

\medskip\noindent
对 (7)。使用 Nakayama 引理（引理 \ref{algebra-lemma-NAK}），
可知该条件等价于 (6) 的条件。

\medskip\noindent
对 (8) 和 (9)。二者等价的方式与 (6) 和 (7) 相同。此外，由引理
\ref{algebra-lemma-differentials-base-change}，有
$\Omega_{S \otimes_R \kappa(\mathfrak p)/\kappa(\mathfrak p)} =
\Omega_{S/R} \otimes_S (S \otimes_R \kappa(\mathfrak p))$。
因此 (9) 等价于 (7)，因为两项中的 $\kappa(\mathfrak q)$-向量空间
典范同构。

\medskip\noindent
对 (10)。由引理 \ref{algebra-lemma-cover}
和 \ref{algebra-lemma-differentials-localize} 得到。

\medskip\noindent
对 (11)。由 (6)、(7) 以及 $S$ 的谱拟紧这一事实得到。
% P01-E0946: the frozen proof cites (6) and (7), which assume global finite type/presentation; the local-at-every-prime hypothesis plus quasi-compactness instead yields a finite principal cover and invokes (10).

\medskip\noindent
对 (12)。写成 $S = R[x_1, \ldots, x_n]/(g_1, \ldots, g_m)$。
由于 $\Omega_{S/R} = 0$，可在
$\Omega_{R[x_1, \ldots, x_n]/R}$ 中写成
$$
\text{d}x_i = \sum h_{ij}\text{d}g_j + \sum a_{ijk}g_j\text{d}x_k
$$
其中 $h_{ij}, a_{ijk} \in R[x_1, \ldots, x_n]$。
选取有限生成的 $\mathbf{Z}$-子代数 $R_0 \subset R$，
使其包含多项式 $g_i, h_{ij}, a_{ijk}$ 的所有系数。置
$S_0 = R_0[x_1, \ldots, x_n]/(g_1, \ldots, g_m)$。这就满足要求。

\medskip\noindent
对 (13)。写成 $S = R[x_1, \ldots, x_n]/I$。
由于 $\Omega_{S/R} = 0$，可在
$\Omega_{R[x_1, \ldots, x_n]/R}$ 中写成
$$
\text{d}x_i = \sum h_{ij}\text{d}g_{ij} + \sum g'_{ik}\text{d}x_k
$$
其中 $h_{ij} \in R[x_1, \ldots, x_n]$，且 $g_{ij}, g'_{ik} \in I$。
选取有限生成的 $\mathbf{Z}$-子代数 $R_0 \subset R$，
使其包含多项式 $g_{ij}, h_{ij}, g'_{ik}$ 的所有系数。置
$S_0 = R_0[x_1, \ldots, x_n]/(g_{ij}, g'_{ik})$。这就满足要求。
\end{proof}

\begin{lemma}
\label{algebra-lemma-diagonal-unramified}
设 $R \to S$ 是环映射。若 $R \to S$ 非分歧，则存在幂等元
$e \in S \otimes_R S$，使得 $S \otimes_R S \to S$ 同构于
$S \otimes_R S \to (S \otimes_R S)_e$。
\end{lemma}

\begin{proof}
令 $J = \Ker(S \otimes_R S \to S)$。由假设及引理
\ref{algebra-lemma-differentials-diagonal}，有 $J/J^2 = 0$。
由于 $S$ 在 $R$ 上为有限型，$J$ 有限生成；具体而言，
若 $x_i$ 在 $R$ 上生成 $S$，则生成元可以取为
$x_i \otimes 1 - 1 \otimes x_i$ 生成。
应用引理 \ref{algebra-lemma-ideal-is-squared-union-connected} 即得结论。
\end{proof}

\begin{lemma}
\label{algebra-lemma-unramified-at-prime}
设 $R \to S$ 是环映射。设 $\mathfrak q \subset S$ 是位于
$R$ 中 $\mathfrak p$ 之上的素理想。若 $S/R$ 在 $\mathfrak q$ 处非分歧，
则
\begin{enumerate}
\item $\mathfrak p S_{\mathfrak q} = \mathfrak qS_{\mathfrak q}$，
且它是局部环 $S_{\mathfrak q}$ 的极大理想；
% P01-E0947: the frozen source reads “we have pS_q = qS_q is the maximal ideal”, joining two predicates ungrammatically.
\item 域扩张 $\kappa(\mathfrak q)/\kappa(\mathfrak p)$ 有限可分。
\end{enumerate}
\end{lemma}

\begin{proof}
先用 $S_g$ 替换 $S$，其中 $g \in S$、$g \not \in \mathfrak q$，
便可假设 $R \to S$ 非分歧。由引理 \ref{algebra-lemma-unramified}，
基变换 $S \otimes_R \kappa(\mathfrak p)$ 在
$\kappa(\mathfrak p)$ 上非分歧。由引理
\ref{algebra-lemma-characterize-smooth-over-field}，它在
$\kappa(\mathfrak p)$ 上光滑，因而平展。因此
$F = S \otimes_R \kappa(\mathfrak p)$ 是
$\kappa(\mathfrak p)$ 的有限可分域扩张的有限乘积，见引理
\ref{algebra-lemma-etale-over-field}。采用注记 \ref{algebra-remark-local-ring-fibre}
的记号和结论，可得
$F_{\overline{\mathfrak q}} = S_\mathfrak q/\mathfrak pS_\mathfrak q$
等于 $\kappa(\mathfrak q)$。这就推出引理。
\end{proof}

\begin{lemma}
\label{algebra-lemma-unramified-quasi-finite}
设 $R \to S$ 是有限型环映射，$\mathfrak q$ 是 $S$ 的素理想。
若 $R \to S$ 在 $\mathfrak q$ 处非分歧，则
$R \to S$ 在 $\mathfrak q$ 处拟有限。
特别地，非分歧环映射拟有限。
\end{lemma}

\begin{proof}
非分歧环映射是有限型的，因此第二个断言显然由第一个断言推出。
为证明第一个断言，应用引理 \ref{algebra-lemma-isolated-point-fibre}
的 (2) 中的刻画，并使用引理 \ref{algebra-lemma-unramified-at-prime}。
\end{proof}

\begin{lemma}
\label{algebra-lemma-characterize-unramified}
设 $R \to S$ 是环映射，$\mathfrak q$ 是 $S$ 的素理想，
位于 $R$ 的素理想 $\mathfrak p$ 之上。若
\begin{enumerate}
\item $R \to S$ 为有限型；
\item $\mathfrak p S_{\mathfrak q}$ 是局部环
$S_{\mathfrak q}$ 的极大理想；
\item 域扩张 $\kappa(\mathfrak q)/\kappa(\mathfrak p)$ 有限可分，
\end{enumerate}
则 $R \to S$ 在 $\mathfrak q$ 处非分歧。
\end{lemma}

\begin{proof}
由引理 \ref{algebra-lemma-unramified} (8)，只须证明
$\Omega_{S \otimes_R \kappa(\mathfrak p) / \kappa(\mathfrak p)}$
在 $\mathfrak q$ 处局部化后为零。因此可以用
$S \otimes_R \kappa(\mathfrak p)$ 替换 $S$，并用
$\kappa(\mathfrak p)$ 替换 $R$。换言之，可以假设 $R = k$ 是域，
而 $S$ 是有限型 $k$-代数。此时，假设推出
$S_{\mathfrak q} \cong \kappa(\mathfrak q)$。因此
$(\Omega_{S/k})_{\mathfrak q} = \Omega_{S_\mathfrak q/k} =
\Omega_{\kappa(\mathfrak q)/k}$ 如所需地为零（第一个等号是引理
\ref{algebra-lemma-differentials-localize}）。
\end{proof}

\begin{lemma}
\label{algebra-lemma-etale-flat-unramified-finite-presentation}
设 $R \to S$ 是环映射。下列条件等价：
\begin{enumerate}
\item $R \to S$ 平展；
\item $R \to S$ 平坦且 G-非分歧；
\item $R \to S$ 平坦、非分歧且为有限表现。
\end{enumerate}
\end{lemma}

\begin{proof}
按定义，(2) 和 (3) 等价。由平展环映射为有限表现、引理
\ref{algebra-lemma-etale}（平展映射的平坦性）以及引理
\ref{algebra-lemma-unramified}（平展映射非分歧），可得
(1) $\Rightarrow$ (3)。反过来，引理 \ref{algebra-lemma-characterize-etale}
中对平展环映射的刻画，以及引理 \ref{algebra-lemma-unramified-at-prime}
中非分歧环映射的结构，说明 (3) 推出 (1)。
% P01-E0948: the frozen compound subject “the characterization ... and the structure ... shows” has singular agreement.
（这里使用了如下事实：若 $R \to S$ 在每个素理想
$\mathfrak q \subset S$ 处都平展，则 $R \to S$ 平展；见引理
\ref{algebra-lemma-etale}。）
\end{proof}

\begin{lemma}
\label{algebra-lemma-characterize-etale-over-polynomial-ring}
设 $k$ 是域，并设
$$
\varphi : k[x_1, \ldots, x_n] \to A, \quad x_i \longmapsto a_i
$$
是有限型环映射。则 $\varphi$ 平展，当且仅当下列两个条件成立：
(a) $A$ 在各极大理想处的局部环维数为 $n$；
(b) 元素 $\text{d}(a_1), \ldots, \text{d}(a_n)$ 作为 $A$-模生成
$\Omega_{A/k}$。
\end{lemma}

\begin{proof}
假设 (a) 和 (b) 成立。条件 (b) 推出
$\Omega_{A/k[x_1, \ldots, x_n]} = 0$，因而 $\varphi$ 非分歧。
于是由引理 \ref{algebra-lemma-etale-flat-unramified-finite-presentation}，
只须证明 $\varphi$ 平坦。设 $\mathfrak m \subset A$ 是极大理想。
置 $X = \Spec(A)$，并以 $x \in X$ 表示与 $\mathfrak m$ 对应的闭点。
由引理 \ref{algebra-lemma-dimension-closed-point-finite-type-field}，
$\dim(A_\mathfrak m)$ 就是 $\dim_x X$。因此由引理
\ref{algebra-lemma-characterize-smooth-over-field} 可知，若 (a) 和 (b) 成立，
则对每个极大理想 $\mathfrak m$，$A_\mathfrak m$ 都是正则局部环。
于是由引理 \ref{algebra-lemma-CM-over-regular-flat}，
$k[x_1, \ldots, x_n]_{\varphi^{-1}(\mathfrak m)} \to A_\mathfrak m$
平坦（并使用了正则局部环是 CM 环这一事实，见引理
\ref{algebra-lemma-regular-ring-CM}）。故由引理
\ref{algebra-lemma-flat-localization}，$\varphi$ 平坦。

\medskip\noindent
假设 $\varphi$ 平展。则 $\Omega_{A/k[x_1, \ldots, x_n]} = 0$，
故 (b) 成立。另一方面，平展环映射平坦（引理
\ref{algebra-lemma-etale}）且拟有限（引理 \ref{algebra-lemma-etale-quasi-finite}）。
因此，对 $A$ 的每个极大理想 $\mathfrak m$，可以将引理
\ref{algebra-lemma-dimension-base-fibre-equals-total} 应用于
$k[x_1, \ldots, x_n]_{\varphi^{-1}(\mathfrak m)} \to A_\mathfrak m$，
得到 $\dim(A_\mathfrak m) = n$，故 (a) 成立。
% P01-E0949: the frozen source has the typo “we my apply” instead of “we may apply”.
\end{proof}

\section{非分歧环映射的局部结构}
\label{algebra-section-local-structure-unramified}

\noindent
非分歧态射局部地（在适当意义下）是闭浸入与平展态射的复合。
本节讨论这一事实的代数基础。

\begin{proposition}
\label{algebra-proposition-unramified-locally-standard}
设 $R \to S$ 是环映射，$\mathfrak q \subset S$ 是素理想。
若 $R \to S$ 在 $\mathfrak q$ 处非分歧，则存在
\begin{enumerate}
\item $g \in S$、$g \not \in \mathfrak q$；
\item 标准平展环映射 $R \to S'$；
\item 满 $R$-代数映射 $S' \to S_g$。
\end{enumerate}
\end{proposition}

\begin{proof}
这个证明与命题 \ref{algebra-proposition-etale-locally-standard} 的证明“相同”。
证明稍显迂回，也许可以缩短。

\medskip\noindent
步骤 1。由定义 \ref{algebra-definition-unramified}，存在
$g \in S$、$g \not \in \mathfrak q$，使得 $R \to S_g$ 非分歧。
因此可以假设 $S$ 在 $R$ 上非分歧。

\medskip\noindent
步骤 2。由引理 \ref{algebra-lemma-unramified}，存在非分歧环映射
$R_0 \to S_0$，其中 $R_0$ 是有限型 $\mathbf{Z}$-代数；还存在环映射
$R_0 \to R$，使得 $S$ 是 $R \otimes_{R_0} S_0$ 的商。
记 $\mathfrak q_0$ 为对应于 $\mathfrak q$ 的 $S_0$ 的素理想。
若对 $(R_0 \to S_0, \mathfrak q_0)$ 证明结论，则经基变换即可对
$(R \to S, \mathfrak q)$ 得到结论。因此可以假设 $R$ 是 Noether 环。

\medskip\noindent
步骤 3。由引理 \ref{algebra-lemma-unramified-quasi-finite}，
$R \to S$ 拟有限。由引理 \ref{algebra-lemma-quasi-finite-open-integral-closure}，
存在有限环映射 $R \to S'$、$R$-代数映射 $S' \to S$，以及元素
$g' \in S'$，其中 $g' \not \in \mathfrak q$，使得 $S' \to S$ 诱导同构
$S'_{g'} \cong S_{g'}$。
% P01-E0950: the frozen source compares g' in S' directly with q in S and repeats “such that”; it should refer to the image of g' (or the contracted prime) once.
（注意，$S'$ 未必在 $R$ 上非分歧。）
因此可以假设：(a) $R$ 是 Noether 环；(b) $R \to S$ 有限；
(c) $R \to S$ 在 $\mathfrak q$ 处非分歧
（但不再必然在所有素理想处非分歧）。

\medskip\noindent
步骤 4。设 $\mathfrak p \subset R$ 是与 $\mathfrak q$ 对应的素理想。
考察纤维环 $S \otimes_R \kappa(\mathfrak p)$。它是
$\kappa(\mathfrak p)$ 上的有限代数，因而是 Artin 环
（见引理 \ref{algebra-lemma-finite-dimensional-algebra}），所以是局部环的有限乘积
$$
S \otimes_R \kappa(\mathfrak p) = \prod\nolimits_{i = 1}^n A_i
$$
见命题 \ref{algebra-proposition-dimension-zero-ring}。其中一个因子，例如 $A_1$，
是局部环 $S_{\mathfrak q}/\mathfrak pS_{\mathfrak q}$，
并且它同构于 $\kappa(\mathfrak q)$，见引理
\ref{algebra-lemma-unramified-at-prime}。其余因子对应于 $S$ 中位于
$\mathfrak p$ 之上的其他素理想，记为
$\mathfrak q_2, \ldots, \mathfrak q_n$。

\medskip\noindent
步骤 5。可以选取非零元素 $\alpha \in \kappa(\mathfrak q)$，
使其生成有限可分域扩张 $\kappa(\mathfrak q)/\kappa(\mathfrak p)$
（所以即使该域扩张平凡，也不允许 $\alpha = 0$）。注意，对任意
$\lambda \in \kappa(\mathfrak p)^*$，元素 $\lambda \alpha$ 也在
$\kappa(\mathfrak p)$ 上生成 $\kappa(\mathfrak q)$。考察元素
$$
\overline{t} =
(\alpha, 0, \ldots, 0) \in
\prod\nolimits_{i = 1}^n A_i =
S \otimes_R \kappa(\mathfrak p).
$$
必要时如上用 $\lambda \alpha$ 替换 $\alpha$，可以假设
$\overline{t}$ 是某个 $t \in S$ 的像。令 $I \subset R[x]$ 为
把 $x$ 映到 $t$ 的 $R$-代数映射 $R[x] \to S$ 的核。
置 $S' = R[x]/I$，于是 $S' \subset S$。有交换图
$$
\xymatrix{
R[x] \ar[r] & S' \ar[r] & S \\
R \ar[u] \ar[ru] \ar[rru] & &
}
$$
按构造，$S$ 的素理想 $\mathfrak q_j$（$j \geq 2$）都位于
$R[x]$ 的素理想 $(\mathfrak p, x)$ 之上，而 $\mathfrak q$ 位于
$R[x]$ 的另一个素理想之上，因为 $\alpha \not = 0$。

\medskip\noindent
步骤 6。以 $\mathfrak q' \subset S'$ 表示对应于 $\mathfrak q$ 的
$S'$ 的素理想。由上可知，$\mathfrak q$ 是 $S$ 中位于
$\mathfrak q'$ 之上的唯一素理想。因此由引理
\ref{algebra-lemma-unique-prime-over-localize-below}，有
$S_{\mathfrak q} = S_{\mathfrak q'}$（由于 $R \to S$ 有限，
$S' \to S$ 也有限，故由引理 \ref{algebra-lemma-integral-going-up}，
$S' \to S$ 满足上升性质）。由此，
$S'_{\mathfrak q'} \to S_{\mathfrak q}$ 是有限且单射的，
因为它是有限单射环映射 $S' \to S$ 的局部化。考察局部环映射
$$
R_{\mathfrak p} \to S'_{\mathfrak q'} \to S_{\mathfrak q}
$$
第二个映射有限且单射。由引理 \ref{algebra-lemma-unramified-at-prime}，有
$S_{\mathfrak q}/\mathfrak pS_{\mathfrak q} = \kappa(\mathfrak q)$。
因此更有
$S_{\mathfrak q}/\mathfrak q'S_{\mathfrak q} = \kappa(\mathfrak q)$。
由于
$$
\kappa(\mathfrak p) \subset \kappa(\mathfrak q') \subset \kappa(\mathfrak q)
$$
且 $\alpha$ 属于 $\kappa(\mathfrak q')$ 在 $\kappa(\mathfrak q)$
中的像，可得 $\kappa(\mathfrak q') = \kappa(\mathfrak q)$。
于是把 Nakayama 引理 \ref{algebra-lemma-NAK} 应用于
$S'_{\mathfrak q'}$-模映射 $S'_{\mathfrak q'} \to S_{\mathfrak q}$，
可知映射 $S'_{\mathfrak q'} \to S_{\mathfrak q}$ 是满射。换言之，
$S'_{\mathfrak q'} \cong S_{\mathfrak q}$。

\medskip\noindent
步骤 7。由引理 \ref{algebra-lemma-isomorphic-local-rings}，存在
$g \in S$、$g \not \in \mathfrak q$ 和
$g' \in S'$、$g' \not \in \mathfrak q'$，使得 $S'_{g'} \cong S_g$。
由于 $R$ 是 Noether 环，$S'$ 是有限 $R$-模 $S$ 的 $R$-子模，
故它在 $R$ 上有限。因此，用 $S'$ 替换 $S$ 后，可以假设：
(a) $R$ 是 Noether 环；(b) $S$ 在 $R$ 上有限；
% P01-E0951: the frozen source omits “is” in “(b) S finite over R”.
(c) $S$ 在 $R$ 上于 $\mathfrak q$ 处非分歧；
(d) $S = R[x]/I$。

\medskip\noindent
步骤 8。考察环
$S \otimes_R \kappa(\mathfrak p) = \kappa(\mathfrak p)[x]/\overline{I}$，
其中 $\overline{I} = I \cdot \kappa(\mathfrak p)[x]$ 是 $I$ 在
$\kappa(\mathfrak p)[x]$ 中生成的理想。由于
$\kappa(\mathfrak p)[x]$ 是主理想整环，存在首一的
$\overline{h} \in \kappa(\mathfrak p)$，使得
$\overline{I} = (\overline{h})$。
% P01-E0952: the frozen source places the monic polynomial h-bar in kappa(p) instead of kappa(p)[x].
对某个 $\lambda \in \kappa(\mathfrak p)$，以
$\lambda \cdot \overline{h}$ 替换 $\overline{h}$ 后，可以假设
$\overline{h}$ 是某个 $h \in R[x]$ 的像。
% P01-E0953: the scaling factor must be nonzero to preserve the principal generator, but the frozen source omits the star on kappa(p).
（问题在于我们不知道能否选取首一的 $h$。）此外，与步骤 4 相同，
已知 $S \otimes_R \kappa(\mathfrak p) = A_1 \times \ldots \times A_n$，
其中 $A_1 = \kappa(\mathfrak q)$ 是 $\kappa(\mathfrak p)$ 的有限可分扩张，
而 $A_2, \ldots, A_n$ 是局部环。这推出
$$
\overline{h} = \overline{h}_1 \overline{h}_2^{e_2} \ldots \overline{h}_n^{e_n}
$$
% P01-E0954: after the frozen source replaces the monic generator by an arbitrary nonzero scalar multiple, monicity is no longer assured; this factorization into monic factors needs the resulting leading unit as a coefficient.
其中 $\overline{h}_i \in \kappa(\mathfrak p)[x]$ 是两两互素的不可约
首一多项式，且 $e_2, \ldots, e_n \geq 1$。编号使得
$A_i = \kappa(\mathfrak p)[x]/(\overline{h}_i^{e_i})$ 是
$\kappa(\mathfrak p)[x]$-代数。
% P01-E0955: the frozen formula uses e_i also for i = 1, although only e_2,...,e_n are defined; it should set e_1 = 1 or state the i = 1 case separately.
注意，$\overline{h}_1$ 是 $\alpha \in \kappa(\mathfrak q)$ 的极小多项式，
因而是可分多项式（其导数与自身互素）。

\medskip\noindent
步骤 9。取首一元素 $m \in I$；这种元素存在，因为环扩张
$R \to R[x]/I$ 有限，因而整。以 $\overline{m}$ 表示它在
$\kappa(\mathfrak p)[x]$ 中的像。可以分解为
$$
\overline{m} = \overline{k}
\overline{h}_1^{d_1} \overline{h}_2^{d_2} \ldots \overline{h}_n^{d_n}
$$
其中 $d_1 \geq 1$、$d_j \geq e_j$（$j = 2, \ldots, n$），
且 $\overline{k} \in \kappa(\mathfrak p)[x]$ 与所有
$\overline{h}_i$ 互素。置 $f = m^l + h$，其中
$l \deg(m) > \deg(h)$ 且 $l \geq 2$。于是 $f$ 是 $R$ 上的首一多项式。
此外，$f$ 在 $\kappa(\mathfrak p)[x]$ 中的像 $\overline{f}$ 分解为
$$
\overline{f} =
\overline{h}_1 \overline{h}_2^{e_2} \ldots \overline{h}_n^{e_n}
+
\overline{k}^l \overline{h}_1^{ld_1} \overline{h}_2^{ld_2}
\ldots \overline{h}_n^{ld_n}
=
\overline{h}_1(\overline{h}_2^{e_2} \ldots \overline{h}_n^{e_n}
+
\overline{k}^l
\overline{h}_1^{ld_1 - 1} \overline{h}_2^{ld_2} \ldots \overline{h}_n^{ld_n})
= \overline{h}_1 \overline{w}
$$
其中 $\overline{w}$ 是与 $\overline{h}_1$ 互素的多项式。
置 $g = f'$（对 $x$ 的导数）。

\medskip\noindent
步骤 10。环映射 $R[x] \to S = R[x]/I$ 具有下列性质：
(1) 它把 $f$ 映到零；
(2) 它把 $g$ 映到 $S \setminus \mathfrak q$ 中的元素。
第一个断言显然，因为 $f$ 是 $I$ 的元素。
% P01-E0956: the frozen source only says h lies in R[x], but Step 10 needs h in I to conclude f = m^l + h lies in I.
为证明第二个断言，只须证明 $g$ 在
$\kappa(\mathfrak q) = \kappa(\mathfrak p)[x]/(\overline{h}_1)$
中的像不为零。$g$ 在 $\kappa(\mathfrak p)[x]$ 中的像是
$\overline{f}$ 的导数。因此 (2) 显然成立，因为
$$
\overline{g} =
\frac{\text{d}\overline{f}}{\text{d}x} =
\overline{w}\frac{\text{d}\overline{h}_1}{\text{d}x} +
\overline{h}_1\frac{\text{d}\overline{w}}{\text{d}x},
$$
$\overline{w}$ 与 $\overline{h}_1$ 互素，且
$\overline{h}_1$ 可分。

\medskip\noindent
步骤 11。由此可知，$\varphi : R[x]/(f) \to S$ 是满环映射，
$R[x]_g/(f)$ 在 $R$ 上平展（因为它标准平展；见引理
\ref{algebra-lemma-standard-etale}），并且 $\varphi(g) \not \in \mathfrak q$。
因此映射 $(R[x]/(f))_g \to S_{\varphi(g)}$ 就是所需的满射。
\end{proof}



\begin{lemma}
\label{algebra-lemma-etale-makes-unramified-closed-at-prime}
设 $R \to S$ 是环映射。设 $\mathfrak q$ 是 $S$ 的素理想，
位于 $\mathfrak p \subset R$ 之上。假设 $R \to S$ 为有限型，
且在 $\mathfrak q$ 处非分歧。则存在
\begin{enumerate}
\item 平展环映射 $R \to R'$；
\item 位于 $\mathfrak p$ 之上的素理想 $\mathfrak p' \subset R'$；
\item 乘积分解
$$
R' \otimes_R S = A \times B
$$
\end{enumerate}
满足下列性质：
\begin{enumerate}
\item $R' \to A$ 是满射；
\item $\mathfrak p'A$ 是 $A$ 的素理想，位于 $\mathfrak p'$ 之上，
并且位于 $\mathfrak q$ 之上。
\end{enumerate}
\end{lemma}

\begin{proof}
可以用沿平展环映射 $R \to R'$ 的任意基变换
$(R' \to R'\otimes_R S, \mathfrak p', \mathfrak q')$ 替换
$(R \to S, \mathfrak p, \mathfrak q)$；其中 $\mathfrak p'$ 是位于
$\mathfrak p$ 之上的素理想，并选取同时位于 $\mathfrak q$ 和
$\mathfrak p'$ 之上的 $\mathfrak q'$。还要注意，给定
$R \to R'$ 和 $\mathfrak p'$ 后，总能找到合适的 $\mathfrak q'$。

\medskip\noindent
$R \to S$ 为有限型这一假设意味着可以应用引理
\ref{algebra-lemma-etale-makes-quasi-finite-finite-variant}。因此可以假设
$S = A_1 \times \ldots \times A_n \times B$，每个 $R \to A_i$ 都有限，
并恰有一个位于 $\mathfrak p$ 之上的素理想 $\mathfrak r_i$，使得
$\kappa(\mathfrak p) \subset \kappa(\mathfrak r_i)$ 纯不可分；而
$R \to B$ 在位于 $\mathfrak p$ 之上的任一素理想处都不拟有限。
于是 $\mathfrak q = \mathfrak r_i$ 对某个 $i$ 成立，因为非分歧态射拟有限
（见引理 \ref{algebra-lemma-unramified-quasi-finite}）。不妨设
$\mathfrak q = \mathfrak r_1$。由引理 \ref{algebra-lemma-unramified-at-prime}，
$\kappa(\mathfrak r_1)/\kappa(\mathfrak p)$ 可分，因而是平凡域扩张，
且 $\mathfrak p(A_1)_{\mathfrak r_1}$ 是极大理想。此外，由引理
\ref{algebra-lemma-unique-prime-over-localize-below}（它适用于
$R \to A_1$，因为由引理 \ref{algebra-lemma-integral-going-up}，
有限环映射满足上升性质），有
$(A_1)_{\mathfrak r_1} = (A_1)_{\mathfrak p}$。
由 Nakayama 引理 \ref{algebra-lemma-NAK}，局部环映射
$R_{\mathfrak p} \to (A_1)_{\mathfrak p} = (A_1)_{\mathfrak r_1}$
是满射。由于 $A_1$ 在 $R$ 上有限，存在
$f \in R$、$f \not \in \mathfrak p$，使得
$R_f \to (A_1)_f$ 是满射。用 $R_f$ 替换 $R$ 后即得结论。
\end{proof}

\begin{lemma}
\label{algebra-lemma-etale-makes-unramified-closed}
\begin{slogan}
对于非分歧环映射，可以通过取平展邻域分离一个纤维中的点。
\end{slogan}
设 $R \to S$ 是环映射，$\mathfrak p$ 是 $R$ 的素理想。
若 $R \to S$ 非分歧，则存在
\begin{enumerate}
\item 平展环映射 $R \to R'$；
\item 位于 $\mathfrak p$ 之上的素理想 $\mathfrak p' \subset R'$；
\item 乘积分解
$$
R' \otimes_R S = A_1 \times \ldots \times A_n \times B
$$
\end{enumerate}
满足下列性质：
\begin{enumerate}
\item $R' \to A_i$ 是满射；
\item $\mathfrak p'A_i$ 是 $A_i$ 的素理想，位于 $\mathfrak p'$ 之上；
\item $B$ 中没有位于 $\mathfrak p'$ 之上的素理想。
\end{enumerate}
\end{lemma}

\begin{proof}
可以应用引理 \ref{algebra-lemma-etale-makes-quasi-finite-finite-variant}。
因此，经一次平展基变换后，可以假设
$S = A_1 \times \ldots \times A_n \times B$，每个 $R \to A_i$ 都有限，
并恰有一个位于 $\mathfrak p$ 之上的素理想 $\mathfrak r_i$，使得
$\kappa(\mathfrak p) \subset \kappa(\mathfrak r_i)$ 纯不可分；而
$R \to B$ 在位于 $\mathfrak p$ 之上的任一素理想处都不拟有限。
由于 $R \to S$ 拟有限（见引理 \ref{algebra-lemma-unramified-quasi-finite}），
可知 $B$ 中没有位于 $\mathfrak p$ 之上的素理想。
由引理 \ref{algebra-lemma-unramified-at-prime}，
$\kappa(\mathfrak r_i)/\kappa(\mathfrak p)$ 可分，因而是平凡域扩张，
且 $\mathfrak p(A_i)_{\mathfrak r_i}$ 是极大理想。此外，由引理
\ref{algebra-lemma-unique-prime-over-localize-below}（它适用于
$R \to A_i$，因为由引理 \ref{algebra-lemma-integral-going-up}，
有限环映射满足上升性质），有
$(A_i)_{\mathfrak r_i} = (A_i)_{\mathfrak p}$。
由 Nakayama 引理 \ref{algebra-lemma-NAK}，局部环映射
$R_{\mathfrak p} \to (A_i)_{\mathfrak p} = (A_i)_{\mathfrak r_i}$
是满射。由于 $A_i$ 在 $R$ 上有限，存在
$f \in R$、$f \not \in \mathfrak p$，使得
$R_f \to (A_i)_f$ 是满射。用 $R_f$ 替换 $R$ 后即得结论。
\end{proof}

\section{Hensel 局部环}
\label{algebra-section-henselian}

\noindent
本节简要讨论 Hensel 局部环的概念。设
$(R, \mathfrak m, \kappa)$ 是局部环。对 $a \in R$，以
$\overline{a}$ 表示 $a$ 在 $\kappa$ 中的像。对多项式
$f \in R[T]$，常以 $\overline{f}$ 表示 $f$ 在 $\kappa[T]$ 中的像。
给定多项式 $f \in R[T]$，以 $f'$ 表示 $f$ 关于 $T$ 的导数。
注意 $\overline{f}' = \overline{f'}$。

\begin{definition}
\label{algebra-definition-henselian}
设 $(R, \mathfrak m, \kappa)$ 是局部环。
\begin{enumerate}
\item 若对每个首一多项式 $f \in R[T]$ 以及
$\overline{f}$ 的每个根 $a_0 \in \kappa$，只要
$\overline{f'}(a_0) \not = 0$，就存在 $a \in R$，使得
$f(a) = 0$ 且 $a_0 = \overline{a}$，则称 $R$ \textit{Hensel}。
\item 若 $R$ Hensel，且其剩余域可分代数闭，则称 $R$
\textit{严格 Hensel}。
\end{enumerate}
\end{definition}

\noindent
注意，条件 $\overline{f'}(a_0) \not = 0$ 等价于 $a_0$ 是多项式
$\overline{f}$ 的单根。事实上，它还意味着：若提升
$a \in R$ 存在，则该提升唯一。

\begin{lemma}
\label{algebra-lemma-uniqueness}
设 $(R, \mathfrak m, \kappa)$ 是局部环。设 $f \in R[T]$。
设 $a, b \in R$ 满足 $f(a) = f(b) = 0$、
$a = b \bmod \mathfrak m$，并且 $f'(a) \not \in \mathfrak m$。
则 $a = b$。
\end{lemma}

\begin{proof}
在 $R[x, y]$ 中写成
$f(x + y) - f(x) = f'(x)y + g(x, y) y^2$（展开 $f(x + y)$ 即可看出；
细节从略）。于是对某个 $c \in R$，有
$0 = f(b) - f(a) = f(a + (b - a)) - f(a) =
f'(a)(b - a) + c (b - a)^2$。由假设，$f'(a)$ 是 $R$ 中的单位。
故 $(b - a)(1 + f'(a)^{-1}c(b - a)) = 0$。
又由假设，$b - a \in \mathfrak m$，所以
$1 + f'(a)^{-1}c(b - a)$ 是 $R$ 中的单位。因此 $b - a = 0$ 于 $R$。
\end{proof}

\noindent
下面给出 Hensel 局部环的刻画。

\begin{lemma}
\label{algebra-lemma-characterize-henselian}
\begin{slogan}
Hensel 局部环的刻画
\end{slogan}
设 $(R, \mathfrak m, \kappa)$ 是局部环。下列条件等价：
\begin{enumerate}
\item $R$ Hensel；
\item 对每个 $f \in R[T]$ 以及 $\overline{f}$ 的每个根
$a_0 \in \kappa$，只要 $\overline{f'}(a_0) \not = 0$，
就存在 $a \in R$，使得 $f(a) = 0$ 且 $a_0 = \overline{a}$；
\item 对任意首一 $f \in R[T]$ 以及任意分解
$\overline{f} = g_0 h_0$，若 $\gcd(g_0, h_0) = 1$，则存在
$R[T]$ 中的分解 $f = gh$，使得
$g_0 = \overline{g}$ 且 $h_0 = \overline{h}$；
\item 对任意首一 $f \in R[T]$ 以及任意分解
$\overline{f} = g_0 h_0$，若 $\gcd(g_0, h_0) = 1$，则存在
$R[T]$ 中的分解 $f = gh$，使得
$g_0 = \overline{g}$ 且 $h_0 = \overline{h}$，并且
$\deg_T(g) = \deg_T(g_0)$；
\item 对任意 $f \in R[T]$ 以及任意分解
$\overline{f} = g_0 h_0$，若 $\gcd(g_0, h_0) = 1$，则存在
$R[T]$ 中的分解 $f = gh$，使得
$g_0 = \overline{g}$ 且 $h_0 = \overline{h}$；
\item 对任意 $f \in R[T]$ 以及任意分解
$\overline{f} = g_0 h_0$，若 $\gcd(g_0, h_0) = 1$，则存在
$R[T]$ 中的分解 $f = gh$，使得
$g_0 = \overline{g}$ 且 $h_0 = \overline{h}$，并且
$\deg_T(g) = \deg_T(g_0)$；
\item 对任意平展环映射 $R \to S$ 以及 $S$ 中位于
$\mathfrak m$ 之上的素理想 $\mathfrak q$，若
$\kappa = \kappa(\mathfrak q)$，则存在 $R \to S$ 的收缩
$\tau : S \to R$；
\item 对任意平展环映射 $R \to S$ 以及 $S$ 中位于
$\mathfrak m$ 之上的素理想 $\mathfrak q$，若
$\kappa = \kappa(\mathfrak q)$，则存在唯一的 $R \to S$ 的收缩
$\tau : S \to R$，满足 $\mathfrak q = \tau^{-1}(\mathfrak m)$；
\item 任意有限 $R$-代数都是局部环的乘积；
\item 任意有限 $R$-代数都是局部环的有限乘积；
\item 任意有限型 $R$-代数 $S$ 都可以写成 $A \times B$，
其中 $R \to A$ 有限，而 $R \to B$ 在位于
$\mathfrak m$ 之上的任一素理想处都不拟有限；
\item 任意有限型 $R$-代数 $S$ 都可以写成 $A \times B$，
其中 $R \to A$ 有限，并且
$\Spec(B \otimes_R \kappa)$ 的每个不可约分支维数都 $\geq 1$；
\item 任意拟有限 $R$-代数 $S$ 都可以写成 $S = A \times B$，
其中 $R \to A$ 有限，且 $B \otimes_R \kappa = 0$。
\end{enumerate}
\end{lemma}

\begin{proof}
先列出较容易的蕴含：
\begin{enumerate}
\item 2$\Rightarrow$1，因为 (2) 考察所有多项式，
而 (1) 只考察首一多项式；
\item 5$\Rightarrow$3，因为 (5) 考察所有多项式，
而 (3) 只考察首一多项式；
\item 6$\Rightarrow$4，因为 (6) 考察所有多项式，
而 (4) 只考察首一多项式；
\item 4$\Rightarrow$3 显然；
\item 6$\Rightarrow$5 显然；
\item 8$\Rightarrow$7 显然；
\item 10$\Rightarrow$9 显然；
\item 11$\Leftrightarrow$12，由在素理想处拟有限的定义；
\item 11$\Rightarrow$13，由拟有限的定义。
\end{enumerate}

\noindent
证明 1$\Rightarrow$8。假设 (1) 成立。设 $R \to S$ 平展，
并设 $\mathfrak q \subset S$ 是满足
$\kappa(\mathfrak q) \cong \kappa$ 的素理想。由命题
\ref{algebra-proposition-etale-locally-standard}，可以找到
$g \in S$、$g \not \in \mathfrak q$，使得 $R \to S_g$ 标准平展。
用 $S_g$ 替换 $S$ 后，可以假设
$S = R[t]_g/(f)$ 标准平展（细节从略）。由于素理想
$\mathfrak q$ 的剩余域是 $\kappa$，它对应于 $\overline{f}$ 的根
$a_0$，而该根不是 $\overline{g}$ 的根。由标准平展代数的定义，
这还意味着 $\overline{f'}(a_0) \not = 0$。此外，再次由标准平展代数
的定义，$f$ 首一；所以可以利用 $R$ Hensel，得到
$a \in R$，满足 $a_0 = \overline{a}$ 且 $f(a) = 0$。
这意味着 $g(a)$ 是 $R$ 的单位，并可按规则 $t \mapsto a$ 得到所需映射
$\tau : S = R[t]_g/(f) \to R$。按构造，
$\tau^{-1}(\mathfrak m) = \mathfrak q$。由引理
\ref{algebra-lemma-uniqueness}，映射 $\tau$ 唯一。这证明了 (8)。

\medskip\noindent
证明 7$\Rightarrow$8。（这一部分其实不重要，可以跳过。）
假设 (7) 成立，并假设 $R \to S$ 平展。设
$\mathfrak q_1, \ldots, \mathfrak q_r$ 是 $S$ 中位于
$\mathfrak m$ 之上的其他素理想。可以找到
$g \in S$、$g \not \in \mathfrak q$，且对
$i = 1, \ldots, r$ 有 $g \in \mathfrak q_i$。事实上，可以论证
$\bigcap_{i=1}^{r} \mathfrak{q}_{i} \not\subset \mathfrak{q}$；
否则对某个 $i$ 有 $\mathfrak{q}_{i} \subset \mathfrak{q}$，
但这是不可能的，因为平展态射的纤维是离散的（例如使用引理
\ref{algebra-lemma-etale-over-field}）。将 (7) 应用于平展环映射
$R \to S_g$ 和素理想 $\mathfrak qS_g$，得到收缩
$\tau_g : S_g \to R$，使得复合映射 $\tau : S \to S_g \to R$
具有性质 $\tau^{-1}(\mathfrak m) = \mathfrak q$。细节从略。

\medskip\noindent
证明 8$\Rightarrow$11。假设 (8) 成立，并设 $R \to S$ 是有限型环映射。
应用引理 \ref{algebra-lemma-etale-makes-quasi-finite-finite}。得到平展环映射
$R \to R'$ 以及位于 $\mathfrak m$ 之上的素理想
$\mathfrak m' \subset R'$，满足 $\kappa = \kappa(\mathfrak m')$，
并且 $R' \otimes_R S = A' \times B'$，其中 $A'$ 在 $R'$ 上有限，
而 $B'$ 在位于 $\mathfrak m'$ 之上的任一素理想处都不在 $R'$ 上拟有限。
应用 (8)，得到收缩 $\tau : R' \to R$，满足
$\mathfrak m = \tau^{-1}(\mathfrak m')$。
% P01-E0957: the frozen inverse-image equality is ill-typed and reversed; for tau : R' -> R it should be m' = tau^{-1}(m).
然后使用
$$
S = (S \otimes_R R') \otimes_{R', \tau} R
= (A' \times B') \otimes_{R', \tau} R
= (A' \otimes_{R', \tau} R)  \times  (B' \otimes_{R', \tau} R)
$$
这就给出 (11) 中的分解。

\medskip\noindent
证明 8$\Rightarrow$10。假设 (8) 成立，并设 $R \to S$ 是有限环映射。
应用引理 \ref{algebra-lemma-etale-makes-quasi-finite-finite}。得到平展环映射
$R \to R'$ 以及位于 $\mathfrak m$ 之上的素理想
$\mathfrak m' \subset R'$，满足 $\kappa = \kappa(\mathfrak m')$，
并且 $R' \otimes_R S = A'_1 \times \ldots \times A'_n \times B'$，
其中 $A'_i$ 在 $R'$ 上有限且恰有一个位于 $\mathfrak m'$ 之上的素理想，
而 $B'$ 在位于 $\mathfrak m'$ 之上的任一素理想处都不在 $R'$ 上拟有限。
应用 (8)，得到收缩 $\tau : R' \to R$，满足
$\mathfrak m' = \tau^{-1}(\mathfrak m)$。于是
\begin{align*}
S & = (S \otimes_R R') \otimes_{R', \tau} R \\
& = (A'_1 \times \ldots \times A'_n \times B') \otimes_{R', \tau} R \\
& = (A'_1 \otimes_{R', \tau} R)  \times
\ldots \times (A'_1 \otimes_{R', \tau} R) \times
(B' \otimes_{R', \tau} R) \\
& = A_1 \times \ldots \times A_n \times B
\end{align*}
% P01-E0958: the frozen third line repeats A'_1 in the final A-factor; it should end with A'_n.
因子 $B$ 在 $R$ 上有限，但 $R \to B$ 在位于 $\mathfrak m$ 之上的
任一素理想处都不拟有限。因此 $B = 0$。各因子 $A_i$ 都是有限
$R$-代数，且恰有一个位于 $\mathfrak m$ 之上的素理想，因而都是局部环。
这证明 $S$ 是局部环的有限乘积。

\medskip\noindent
证明 9$\Rightarrow$10。这是因为，若 $S$ 在局部环 $R$ 上有限，
则它至多有有限多个极大理想。事实上，由 $R \to S$ 的上升性质，
$S$ 的所有极大理想都位于 $\mathfrak m$ 之上，而
$S/\mathfrak mS$ 是 Artin 环，因而只有有限多个素理想。

\medskip\noindent
证明 10$\Rightarrow$1。假设 (10) 成立。设 $f \in R[T]$ 是首一多项式，
$a_0 \in \kappa$ 是 $\overline{f}$ 的单根。则
$S = R[T]/(f)$ 是有限 $R$-代数。应用 (10)，得到
$S = A_1 \times \ldots \times A_r$，它是局部 $R$-代数的有限乘积。
特别地，$S/\mathfrak mS = \prod A_i/\mathfrak mA_i$ 是
$\kappa[T]/(\overline{f})$ 分解为局部环乘积的分解。这意味着其中一个因子，
例如 $A_1/\mathfrak mA_1$，就是商
$\kappa[T]/(\overline{f}) \to \kappa[T]/(T - a_0)$。
由于 $A_1$ 是有限自由 $R$-模 $S$ 的直和项，它自身也是有限自由
$R$-模。又因 $A_1/\mathfrak mA_1$ 是维数为 1 的
$\kappa$-向量空间，可知作为 $R$-模有 $A_1 \cong R$。
这显然意味着 $R \to A_1$ 是同构。令 $a \in R$ 为 $T$ 在映射
$R[T] \to S \to A_1 \to R$ 下的像。则如所需，
$f(a) = 0$ 且 $\overline{a} = a_0$。

\medskip\noindent
证明 13$\Rightarrow$1。假设 (13) 成立。设 $f \in R[T]$ 是首一多项式，
$a_0 \in \kappa$ 是 $\overline{f}$ 的单根。则
$S_1 = R[T]/(f)$ 是有限 $R$-代数。取任意 $g \in R[T]$，满足
$\overline{g} = \overline{f}/(T - a_0)$。则 $S = (S_1)_g$ 是拟有限
$R$-代数，并且
$S \otimes_R \kappa \cong \kappa[T]_{\overline{g}}/(\overline{f})
\cong \kappa[T]/(T - a_0) \cong \kappa$。
将 (13) 应用于 $S$，得到 $S = A \times B$，其中 $A$ 在 $R$ 上有限，
且 $B \otimes_R \kappa = 0$。特别地，
$\kappa \cong S/\mathfrak mS = A/\mathfrak mA$。由于 $A$ 是平坦
$R$-代数 $S$ 的直和项，可知它有限平坦，因而在 $R$ 上自由。
又因 $A/\mathfrak mA$ 是维数为 1 的 $\kappa$-向量空间，
可知作为 $R$-模有 $A \cong R$。这显然意味着 $R \to A$ 是同构。
令 $a \in R$ 为 $T$ 在映射 $R[T] \to S \to A \to R$ 下的像。
则如所需，$f(a) = 0$ 且 $\overline{a} = a_0$。

\medskip\noindent
证明 8$\Rightarrow$2。假设 (8) 成立。设 $f \in R[T]$ 是任意多项式，
并设 $a_0 \in \kappa$ 是单根。则代数
$S = R[T]_{f'}/(f)$ 在 $R$ 上平展。
% P01-E0959: the frozen proof calls a_0 “a simple root” without specifying that it is a root of f-bar.
设 $\mathfrak q \subset S$ 是由 $\mathfrak m$ 和 $T - b$ 生成的素理想，
其中 $b \in R$ 是满足 $\overline{b} = a_0$ 的任意元素。
将 (8) 应用于 $S$ 和 $\mathfrak q$，得到 $\tau : S \to R$。
于是像 $\tau(T) = a \in R$ 满足 (2)。

\medskip\noindent
至此可知，(1)、(2)、(7)、(8)、(9)、(10)、(11)、(12)、(13)
全都等价。(3)、(4)、(5)、(6) 中最弱的断言是 (3)，最强的是 (6)。
所以还须证明 (3) 推出 (1)，以及 (1) 推出 (6)。

\medskip\noindent
证明 3$\Rightarrow$1。假设 (3) 成立。设 $f \in R[T]$ 首一，
并设 $a_0 \in \kappa$ 是 $\overline{f}$ 的单根。这给出分解
$\overline{f} = (T - a_0)h_0$，其中 $h_0(a_0) \not = 0$，
故 $\gcd(T - a_0, h_0) = 1$。应用 (3)，得到分解
$f = gh$，其中 $\overline{g} = T - a_0$ 且
$\overline{h} = h_0$。置 $S = R[T]/(f)$，它是有限自由 $R$-代数。
我们也用 $g$、$h$ 表示 $g$、$h$ 在 $S$ 中的像。由于模掉
$\mathfrak m$ 后等式成立，由 Nakayama 引理 \ref{algebra-lemma-NAK}，有
$gS + hS = S$。又因为在 $S$ 中 $gh = f = 0$，还可得
$gS \cap hS = 0$。故由中国剩余定理，得到
$S = S/(g) \times S/(h)$。这意味着 $A = S/(g)$ 是有限自由
$R$-模的直和项，因而有限自由。此外，由
$A/\mathfrak mA = \kappa[T]/(T - a_0)$ 可知 $A$ 的秩为 $1$。
因此映射 $R \to A$ 是同构。令 $a \in R$ 等于 $T$ 在映射
$R[T] \to S \to A \to R$ 下的像，即得到满足
$f(a) = 0$ 且 $\overline{a} = a_0$ 的 $R$ 中元素。

\medskip\noindent
证明 1$\Rightarrow$6。假设 (1) 成立，或者等价地，假设
(1)、(2)、(7)、(8)、(9)、(10)、(11)、(12)、(13) 全部成立。
设 $f \in R[T]$ 是多项式。假设
$\overline{f} = g_0h_0$ 是满足 $\gcd(g_0, h_0) = 1$ 的分解。
可以且将要假设 $g_0$ 首一。由于有这一分解，可知 $f$ 的系数
在 $R$ 中生成单位理想。
% P01-E0960: if g_0 = 1 and h_0 = 0, the coprime factorization is valid but f-bar is zero, so the coefficients of f need not generate the unit ideal; the trivial factorization must be handled separately.
考察 $S = R[T]/(f)$。这意味着 $S$ 在 $R$ 上具有有限纤维，
因而在 $R$ 上拟有限。由引理 \ref{algebra-lemma-grothendieck-general}，
还可知 $S$ 在 $R$ 上平坦。结合 (13) 和 (10)，可写成
$S = A_1 \times \ldots \times A_n \times B$，其中每个 $A_i$
都是在 $R$ 上有限的局部环，且 $B \otimes_R \kappa = 0$。
重新排列因子 $A_1, \ldots, A_n$ 后，可以假设
$$
\kappa[T]/(g_0) =
A_1/\mathfrak m A_1 \times \ldots \times A_r/\mathfrak mA_r,
\ \kappa[T]/(h_0) =
A_{r + 1}/\mathfrak mA_{r + 1} \times \ldots \times A_n/\mathfrak mA_n
$$
作为 $\kappa[T]$ 的商。有限平坦 $R$-代数
$A = A_1 \times \ldots \times A_r$ 作为 $R$-模自由，见引理
\ref{algebra-lemma-finite-flat-local}。它的秩是 $\deg_T(g_0)$。
设 $g \in R[T]$ 是 $R$-线性算子 $T : A \to A$ 的特征多项式。
则 $g$ 是次数 $\deg_T(g) = \deg_T(g_0)$ 的首一多项式，
且 $\overline{g} = g_0$。由 Cayley--Hamilton 定理（引理
\ref{algebra-lemma-charpoly}），有 $g(T_A) = 0$，其中 $T_A$ 表示 $T$ 在
$A$ 中的像。因此得到良定义的满射 $R[T]/(g) \to A$，
它由 Nakayama 引理 \ref{algebra-lemma-NAK} 是同构。按构造，
映射 $R[T] \to A$ 经由 $R[T]/(f)$ 分解，因而对某个 $h$，
可以写成 $f = gh$。证明完成。
\end{proof}

\begin{lemma}
\label{algebra-lemma-finite-over-henselian}
设 $(R, \mathfrak m, \kappa)$ 是 Hensel 局部环。
\begin{enumerate}
\item 若 $R \to S$ 是有限环映射，则 $S$ 是 Hensel 局部环的有限乘积，
其中每个因子都在 $R$ 上有限。
\item 若 $R \to S$ 是有限环映射且 $S$ 是局部环，则 $S$ 是
Hensel 局部环，并且 $R \to S$ 是（有限）局部环映射。
\item 若 $R \to S$ 是有限型环映射，$\mathfrak q$ 是 $S$ 中位于
$\mathfrak m$ 之上的素理想，且 $R \to S$ 在该处拟有限，则
$S_{\mathfrak q}$ Hensel 且在 $R$ 上有限。
\item 若 $R \to S$ 拟有限，则对每个位于 $\mathfrak m$ 之上的素理想
$\mathfrak q$，$S_{\mathfrak q}$ 都 Hensel 且在 $R$ 上有限。
\end{enumerate}
\end{lemma}

\begin{proof}
(2) 推出 (1)，因为由引理 \ref{algebra-lemma-characterize-henselian} (10)，
(1) 中的 $S$ 是它在位于 $\mathfrak m$ 之上的素理想处的局部化的有限乘积。
(2) 也由引理 \ref{algebra-lemma-characterize-henselian} (10) 推出，
因为任意有限 $S$-代数也都是有限 $R$-代数（当然，局部环之间的任意
有限环映射都是局部映射）。

\medskip\noindent
设 $R \to S$ 和 $\mathfrak q$ 如 (3) 中所设。写成
$S = A \times B$，其中 $A$ 在 $R$ 上有限，而 $B$ 在位于
$\mathfrak m$ 之上的任一素理想处都不在 $R$ 上拟有限；见引理
\ref{algebra-lemma-characterize-henselian} (11)。因此 $S_\mathfrak q$
是 $A$ 在某个极大理想处的局部化，由 (1) 可得 (3)。
(4) 由 (3) 推出。
\end{proof}

\begin{lemma}
\label{algebra-lemma-mop-up}
设 $(R, \mathfrak m, \kappa)$ 是 Hensel 局部环。
任意有限型 $R$-代数 $S$ 都可写成
$S = A_1 \times \ldots \times A_n \times B$，其中 $A_i$ 是在
$R$ 上有限的局部环，而 $R \to B$ 在 $B$ 中位于
$\mathfrak m$ 之上的任一素理想处都不拟有限。
\end{lemma}

\begin{proof}
这是引理 \ref{algebra-lemma-characterize-henselian} 的 (11) 与 (10) 的结合。
\end{proof}

\begin{lemma}
\label{algebra-lemma-mop-up-strictly-henselian}
设 $(R, \mathfrak m, \kappa)$ 是严格 Hensel 局部环。
任意有限型 $R$-代数 $S$ 都可写成
$S = A_1 \times \ldots \times A_n \times B$，其中 $A_i$ 是在
$R$ 上有限的局部环，$\kappa \subset \kappa(\mathfrak m_{A_i})$
有限纯不可分，而 $R \to B$ 在 $B$ 中位于 $\mathfrak m$ 之上的
任一素理想处都不拟有限。
\end{lemma}

\begin{proof}
先按引理 \ref{algebra-lemma-mop-up} 写成
$S = A_1 \times \ldots \times A_n \times B$。
域扩张 $\kappa(\mathfrak m_{A_i})/\kappa$ 有限，而 $\kappa$
可分代数闭，故该扩张有限纯不可分。
\end{proof}

\begin{lemma}
\label{algebra-lemma-henselian-cat-finite-etale}
设 $(R, \mathfrak m, \kappa)$ 是 Hensel 局部环。有限平展环扩张
$R \to S$ 的范畴，借助函子 $S \mapsto S/\mathfrak mS$，
等价于有限平展代数 $\kappa \to \overline{S}$ 的范畴。
\end{lemma}

\begin{proof}
以 $\mathcal{C} \to \mathcal{D}$ 表示陈述中的范畴函子。
假设 $R \to S$ 有限平展。利用引理 \ref{algebra-lemma-mop-up} 或引理
\ref{algebra-lemma-characterize-henselian} (10)，可以写成
$$
S = A_1 \times \ldots \times A_n
$$
其中 $A_i$ 是在 $S$ 上有限平展的局部环。特别地，
$A_i/\mathfrak mA_i$ 是 $\kappa$ 的有限可分域扩张，见引理
\ref{algebra-lemma-etale-at-prime}。由此可见，$\mathcal{C}$ 和
$\mathcal{D}$ 中的每个对象都典范分解为不可约部分，
这些部分经给定函子相互对应。
其次，假设 $S_1$、$S_2$ 在 $R$ 上有限平展，并且
$\kappa_1 = S_1/\mathfrak mS_1$ 与
$\kappa_2 = S_2/\mathfrak mS_2$ 都是域（即 $\kappa$ 的有限可分扩张）。
则 $S_1 \otimes_R S_2$ 在 $R$ 上有限平展，并且可以像上面一样写成
$$
S_1 \otimes_R S_2 = A_1 \times \ldots \times A_n
$$
于是 $\Hom_R(S_1, S_2)$ 与满足 $S_2 \to A_i$ 为同构的指标
$i \in \{1, \ldots, n\}$ 的集合相等同。为说明这一点，注意给定任意
$R$-代数映射 $\varphi : S_1 \to S_2$，映射
$\varphi \times 1 : S_1 \otimes_R S_2 \to S_2$ 是满射，
因而等于向某个因子 $A_i$ 的投影。完全相同地，
$\Hom_\kappa(\kappa_1, \kappa_2)$ 与满足
$\kappa_2 \to A_i/\mathfrak mA_i$ 为同构的指标
$i \in \{1, \ldots, n\}$ 的集合相等同。由上面的讨论，
这两个指标集合相同，故该函子完全忠实。
最后，设 $\kappa'/\kappa$ 是有限可分域扩张。由引理
\ref{algebra-lemma-make-etale-map-prescribed-residue-field}，存在平展环映射
$R \to S$ 以及 $S$ 中位于 $\mathfrak m$ 之上的素理想
$\mathfrak q$，使得 $\kappa \subset \kappa(\mathfrak q)$
同构于给定扩张。由 (1)，可以写成
$S = A_1 \times \ldots \times A_n \times B$。
% P01-E0961: the frozen proof cites an undefined “part (1)” here; the required global finite-plus-non-quasi-finite decomposition is lemma-mop-up.
由于 $R \to S$ 拟有限，可知 $B$ 中不存在位于 $\mathfrak m$ 之上的素理想。
因此 $S_{\mathfrak q}$ 等于某个指标 $i$ 所对应的 $A_i$。故 $R \to A_i$ 有限平展，
并给出所指定的剩余域扩张。于是该函子本质满，结论得证。
\end{proof}

\begin{lemma}
\label{algebra-lemma-unramified-over-strictly-henselian}
设 $(R, \mathfrak m, \kappa)$ 是严格 Hensel 局部环。
设 $R \to S$ 是非分歧环映射。则
$$
S = A_1 \times \ldots \times A_n \times B
$$
其中每个 $R \to A_i$ 都是满射，且 $B$ 中没有位于
$\mathfrak m$ 之上的素理想。
\end{lemma}

\begin{proof}
先按引理 \ref{algebra-lemma-mop-up} 写成
$S = A_1 \times \ldots \times A_n \times B$。
此时 $R \to A_i$ 有限非分歧，而 $A_i$ 是局部环。故 $A_i$
的极大理想是 $\mathfrak mA_i$，其剩余域
$A_i / \mathfrak m A_i$ 是 $\kappa$ 的有限可分扩张，见引理
\ref{algebra-lemma-unramified-at-prime}。然而，$R$ 严格 Hensel 意味着
$\kappa$ 可分代数闭，所以 $\kappa = A_i / \mathfrak m A_i$。
由 Nakayama 引理 \ref{algebra-lemma-NAK}，可知 $R \to A_i$ 如所需地为满射。
\end{proof}

\begin{lemma}
\label{algebra-lemma-complete-henselian}
\begin{slogan}
完备局部环由 Newton 法可知为 Hensel 环
\end{slogan}
设 $(R, \mathfrak m, \kappa)$ 是完备局部环，见定义
\ref{algebra-definition-complete-local-ring}。则 $R$ Hensel。
\end{lemma}

\begin{proof}
设 $f \in R[T]$ 首一。以
$f_n \in R/\mathfrak m^{n + 1}[T]$ 表示其像，以 $f'_n$ 表示
$f_n$ 关于 $T$ 的导数。设 $a_0 \in \kappa$ 是 $f_0$ 的单根。
如下归纳地把它提升为 $f$ 在 $R$ 上的解：假设已经给定
$a_n \in R/\mathfrak m^{n + 1}$，满足
$a_n \bmod \mathfrak m = a_0$ 且 $f_n(a_n) = 0$。
任取元素 $b \in R/\mathfrak m^{n + 2}$，满足
$a_n = b \bmod \mathfrak m^{n + 1}$。则
$f_{n + 1}(b) \in \mathfrak m^{n + 1}/\mathfrak m^{n + 2}$。置
$$
a_{n + 1} = b - f_{n + 1}(b)/f'_{n + 1}(b)
$$
（Newton 法）。这是有意义的，因为由对 $a_0$ 的条件，
$f'_{n + 1}(b) \in R/\mathfrak m^{n + 2}$ 可逆。于是计算得
$f_{n + 1}(a_{n + 1}) = f_{n + 1}(b) - f_{n + 1}(b) = 0$
于 $R/\mathfrak m^{n + 2}$。如此构造的元素系
$a_n \in R/\mathfrak m^{n + 1}$ 相容，故得到元素
$a \in \lim R/\mathfrak m^n = R$（这里使用了 $R$ 完备）。
此外，$f(a) = 0$，因为它在每个 $R/\mathfrak m^n$ 中的像都为零。
最后，$\overline{a} = a_0$，结论得证。
\end{proof}

\begin{lemma}
\label{algebra-lemma-local-dimension-zero-henselian}
\begin{slogan}
零维局部环是 Hensel 环。
\end{slogan}
设 $(R, \mathfrak m)$ 是维数为 $0$ 的局部环。则 $R$ Hensel。
\end{lemma}

\begin{proof}
设 $R \to S$ 是有限环映射。由引理
\ref{algebra-lemma-characterize-henselian}，只须证明 $S$ 是局部环的乘积。
由引理 \ref{algebra-lemma-finite-finite-fibres}，$S$ 有有限多个素理想
$\mathfrak m_1, \ldots, \mathfrak m_r$，它们全都位于
$\mathfrak m$ 之上。由引理 \ref{algebra-lemma-integral-no-inclusion}，
这些素理想之间不存在包含关系，故它们全都极大。由引理
\ref{algebra-lemma-Zariski-topology}，
$\mathfrak m_1 \cap \ldots \cap \mathfrak m_r$ 的每个元素都幂零。
于是由引理 \ref{algebra-lemma-product-local}，$S$ 是 $S$ 在各素理想
$\mathfrak m_i$ 处局部化的乘积。
\end{proof}

\noindent
下一个引理将是 Hensel 化与严格 Hensel 化的唯一性和函子性质的关键。

\begin{lemma}
\label{algebra-lemma-map-into-henselian}
设 $R \to S$ 是环映射，且 $S$ Hensel 局部。给定
\begin{enumerate}
\item 平展环映射 $R \to A$；
\item $A$ 中位于 $\mathfrak p = R \cap \mathfrak m_S$ 之上的素理想
$\mathfrak q$；
% P01-E0962: for a possibly noninjective map R -> S, the frozen expression R intersect m_S should be the inverse image of m_S in R.
\item $\kappa(\mathfrak p)$-代数映射
$\tau : \kappa(\mathfrak q) \to S/\mathfrak m_S$，
\end{enumerate}
则存在唯一的 $R$-代数同态 $f : A \to S$，满足
$\mathfrak q = f^{-1}(\mathfrak m_S)$，且 $f$ 在剩余域上诱导映射
$\tau$。
\end{lemma}

\begin{proof}
考察 $A \otimes_R S$。由引理 \ref{algebra-lemma-etale}，
它是在 $S$ 上的平展代数。此外，核
$$
\mathfrak q' = \Ker(A \otimes_R S \to
\kappa(\mathfrak q) \otimes_{\kappa(\mathfrak p)} \kappa(\mathfrak m_S)
\xrightarrow{\tau \otimes 1}
\kappa(\mathfrak m_S))
$$
是位于 $\mathfrak m_S$ 之上的素理想，且其剩余域等于 $S$ 的剩余域。
因此，由引理 \ref{algebra-lemma-characterize-henselian}，存在唯一收缩
$\sigma : A \otimes_R S \to S$，满足
$\sigma^{-1}(\mathfrak m_S) = \mathfrak q'$。令 $f$ 等于复合映射
$A \to A \otimes_R S \to S$。略去对 $f$ 的各项性质的验证；
$f$ 的唯一性来自 $\sigma$ 的唯一性（细节从略）。
\end{proof}

\begin{lemma}
\label{algebra-lemma-strictly-henselian-solutions}
设 $\varphi : R \to S$ 是严格 Hensel 局部环的局部同态。
设 $P_1, \ldots, P_n \in R[x_1, \ldots, x_n]$ 是多项式，
使得 $R[x_1, \ldots, x_n]/(P_1, \ldots, P_n)$ 在 $R$ 上平展。
则映射
$$
R^n \longrightarrow S^n, \quad
(h_1, \ldots, h_n) \longmapsto (\varphi(h_1), \ldots, \varphi(h_n))
$$
在下列两个集合之间诱导双射：
$$
\{
(r_1, \ldots, r_n) \in R^n
\mid
P_i(r_1, \ldots, r_n) = 0, \ i = 1, \ldots, n
\}
$$
以及
$$
\{
(s_1, \ldots, s_n) \in S^n
\mid
P^\varphi_i(s_1, \ldots, s_n) = 0, \ i = 1, \ldots, n
\}
$$
其中 $P^\varphi_i \in S[x_1, \ldots, x_n]$ 是 $P_i$ 在
$\varphi$ 下的像。
\end{lemma}

\begin{proof}
第一个解集典范同构于集合
$$
\Hom_R(R[x_1, \ldots, x_n]/(P_1, \ldots, P_n), R).
$$
由于 $R$ Hensel，映射 $R \to R/\mathfrak m_R$ 在此集合与剩余域
$R/\mathfrak m_R$ 中的解集之间诱导双射，见引理
\ref{algebra-lemma-characterize-henselian}。对 $S$ 同样如此。
现在，由于 $R[x_1, \ldots, x_n]/(P_1, \ldots, P_n)$ 在 $R$ 上平展，
且 $R/\mathfrak m_R$ 可分代数闭，可知
$R/\mathfrak m_R[x_1, \ldots, x_n]/(\overline{P}_1, \ldots, \overline{P}_n)$
是 $R/\mathfrak m_R$ 的若干副本的有限乘积；其中
$\overline{P}_i$ 是 $P_i$ 在
$R/\mathfrak m_R[x_1, \ldots, x_n]$ 中的像。
% P01-E0963: the frozen residue-field polynomial-ring expressions omit parentheses around R/m_R (and later S/m_S), making the displayed quotients formally ill-typed or ambiguous.
因此张量积
$$
R/\mathfrak m_R[x_1, \ldots, x_n]/(\overline{P}_1, \ldots, \overline{P}_n)
\otimes_{R/\mathfrak m_R} S/\mathfrak m_S
=
S/\mathfrak m_S[x_1, \ldots, x_n]/
(\overline{P}^\varphi_1, \ldots, \overline{P}^\varphi_n)
$$
也是 $S/\mathfrak m_S$ 的若干副本在同一个指标集上的有限乘积。
这证明了引理。
\end{proof}

\begin{lemma}
\label{algebra-lemma-split-ML-henselian}
设 $R$ 是 Hensel 局部环。任意可数生成的 Mittag--Leffler
$R$-模都是有限表现 $R$-模的直和。
\end{lemma}

\begin{proof}
设 $M$ 是可数生成的 Mittag--Leffler $R$-模。我们断言：
对任意元素 $x \in M$，都存在直和分解 $M = N \oplus K$，
使得 $x \in N$，模 $N$ 有限表现，而 $K$ 是 Mittag--Leffler 模。

\medskip\noindent
假设该断言成立。选取 $M$ 的生成元 $x_1, x_2, x_3, \ldots$。
由该断言，可以归纳地找到直和分解
$$
M = N_1 \oplus N_2 \oplus \ldots \oplus N_n \oplus K_n
$$
其中 $N_i$ 有限表现，
$x_1, \ldots, x_n \in N_1 \oplus \ldots \oplus N_n$，
且 $K_n$ 是 Mittag--Leffler 模。无限重复即可得到
$M = \bigoplus N_i$。

\medskip\noindent
还须证明该断言。设 $x \in M$。由引理 \ref{algebra-lemma-ML-countable}，
存在自同态 $\alpha : M \to M$，使得 $\alpha$ 经由有限表现模分解，
且 $\alpha (x) = x$。设 $\alpha$ 分解为
$$
\xymatrix{
M \ar[r]^\pi & P \ar[r]^i & M
}
$$
置 $a = \pi \circ \alpha \circ i : P \to P$，则
$i \circ a \circ \pi = \alpha^3$。由引理
\ref{algebra-lemma-charpoly-module}，存在首一多项式 $P \in R[T]$，
使得 $P(a) = 0$。
% P01-E0964: the frozen proof reuses P for both the finitely presented module in the diagram and the annihilating monic polynomial, making later references such as “N is a quotient of P” ambiguous.
注意，这形式地推出 $\alpha^2 P(\alpha) = 0$。因此可以把 $M$
看作 $R[T]/(T^2P)$-模。假设 $x \not = 0$。则
$\alpha(x) = x$ 推出 $0 = \alpha^2P(\alpha)x = P(1)x$，
故 $P(1) = 0$ 于 $R/I$，其中
$I = \{r \in R \mid rx = 0\}$ 是 $x$ 的零化子。
由于 $x \not = 0$，有 $I \subset \mathfrak m_R$，所以 $1$ 是
$\overline{P} = P \bmod \mathfrak m_R \in R/\mathfrak m_R[T]$ 的根。
由于 $R$ Hensel，可以找到分解
$$
T^2P = (T^2 Q_1) Q_2
$$
其中 $Q_1, Q_2 \in R[T]$，并且
$Q_2 = (T - 1)^e \bmod \mathfrak m_R R[T]$、
$Q_1(1) \not = 0 \bmod \mathfrak m_R$；见引理
\ref{algebra-lemma-characterize-henselian}。令
$N = \Im(\alpha^2Q_1(\alpha) : M \to M)$，并令
$K = \Im(Q_2(\alpha) : M \to M)$。由于 $T^2Q_1$ 与 $Q_2$
生成 $R[T]$ 的单位理想，得到直和分解 $M = N \oplus K$。
此外，$Q_2$ 在 $N$ 上作用为零，$T^2Q_1$ 在 $K$ 上作用为零。
注意，$N$ 是 $P$ 的商，因而有限生成。又有 $x \in N$，因为
$\alpha^2Q_1(\alpha)x = Q_1(1)x$，而 $Q_1(1)$ 是 $R$ 中的单位。
由引理 \ref{algebra-lemma-direct-sum-ML}，模 $N$ 和 $K$ 都是
Mittag--Leffler 模。最后，有限生成模 $N$ 有限表现，因为有限生成的
Mittag--Leffler 模有限表现，见例 \ref{algebra-example-ML} (1)。
\end{proof}


\section{平展环映射的滤过余极限}
\label{algebra-section-ind-etale}

\noindent
本节是关于 ind-平展环映射一节（《pro-平展上同调》，第
\ref{proetale-section-ind-etale} 节）的先导。
这些材料也将用于证明局部环的 Hensel 化与严格 Hensel 化的唯一性。

\begin{lemma}
\label{algebra-lemma-base-change-colimit-etale}
设 $R \to A$ 与 $R \to R'$ 是环映射。若 $A$ 是平展
$R$-代数的滤过余极限，则 $R' \otimes_R A$ 也是平展
$R'$-代数的滤过余极限。
% P01-E0965: the frozen English statement duplicates its predicate (“then so is ... is a filtered colimit”).
\end{lemma}

\begin{proof}
这是因为余极限与张量积交换，而平展环映射在基变换下保持
（引理 \ref{algebra-lemma-etale}）。
\end{proof}

\begin{lemma}
\label{algebra-lemma-composition-colimit-etale}
设 $A \to B \to C$ 是环映射。若 $B$ 是平展
$A$-代数的滤过余极限，且 $C$ 是平展 $B$-代数的滤过余极限，
则 $C$ 是平展 $A$-代数的滤过余极限。
\end{lemma}

\begin{proof}
我们将使用引理 \ref{algebra-lemma-when-colimit} 的判别准则。
设 $A \to P \to C$ 是 $A \to C$ 的一个分解，其中 $P$
在 $A$ 上有限表现。写
$B = \colim_{i \in I} B_i$，其中 $I$ 是有向集，而 $B_i$
是平展 $A$-代数。写 $C = \colim_{j \in J} C_j$，其中
$J$ 是有向集，而 $C_j$ 是平展 $B$-代数。由引理
\ref{algebra-lemma-characterize-finite-presentation}，对某个 $j$，
可以把 $P \to C$ 分解为 $P \to C_j \to C$。
由引理 \ref{algebra-lemma-etale}，可以找到 $i \in I$ 以及平展环映射
$B_i \to C'_j$，使得 $C_j = B \otimes_{B_i} C'_j$。
于是
$C_j = \colim_{i' \geq i} B_{i'} \otimes_{B_i} C'_j$，
再次可见 $P \to C_j$ 分解为
$P \to B_{i'} \otimes_{B_i} C'_j \to C$。
由于平展环映射的复合与张量积仍为平展，
$A \to C' = B_{i'} \otimes_{B_i} C'_j$ 是平展的。
% P01-E0966: the frozen English source begins this sentence with “As” and then starts a second causal “as”, leaving a sentence fragment.
因此我们已经把 $P \to C$ 分解为 $P \to C' \to C$，
其中 $C'$ 在 $A$ 上平展，故引理 \ref{algebra-lemma-when-colimit}
的判别准则适用。
\end{proof}

\begin{lemma}
\label{algebra-lemma-colimit-colimit-etale}
设 $R$ 是环。设 $A = \colim A_i$ 是 $R$-代数的滤过余极限，
并且每个 $A_i$ 都是平展 $R$-代数的滤过余极限。
则 $A$ 是平展 $R$-代数的滤过余极限。
\end{lemma}

\begin{proof}
写 $A_i = \colim_{j \in J_i} A_j$，其中 $J_i$ 是有向集，
而 $A_j$ 是平展 $R$-代数。对每个 $i \leq i'$ 和
$j \in J_i$，存在 $j' \in J_{i'}$ 以及 $R$-代数映射
$\varphi_{jj'} : A_j \to A_{j'}$，使图表
$$
\xymatrix{
A_i \ar[r] & A_{i'} \\
A_j \ar[u] \ar[r]^{\varphi_{jj'}} & A_{j'} \ar[u]
}
$$
交换。这是因为 $R \to A_j$ 有限表现，所以引理
\ref{algebra-lemma-characterize-finite-presentation} 适用。
令 $\mathcal{J}$ 为如下范畴：其对象为
$\coprod_{i \in I} J_i$，其态为如上的三元组
$(j, j', \varphi_{jj'})$（复合律是显然的）。
则 $\mathcal{J}$ 是滤过范畴，并且
$A = \colim_\mathcal{J} A_j$。细节从略。
\end{proof}

\begin{lemma}
\label{algebra-lemma-colimit-colimit-etale-better}
设 $I$ 是有向集。设 $i \mapsto (R_i \to A_i)$ 是 $I$
上的环箭头系统。置 $R = \colim R_i$、$A = \colim A_i$。
若每个 $A_i$ 都是平展 $R_i$-代数的滤过余极限，
则 $A$ 是平展 $R$-代数的滤过余极限。
\end{lemma}

\begin{proof}
这是因为
$A = A \otimes_R R = \colim A_i \otimes_{R_i} R$，
因而可以应用引理 \ref{algebra-lemma-colimit-colimit-etale}：
由引理 \ref{algebra-lemma-base-change-colimit-etale}，
$R \to A_i \otimes_{R_i} R$ 是平展环映射的滤过余极限。
\end{proof}

\begin{lemma}
\label{algebra-lemma-colimits-of-etale}
设 $R$ 是环。设 $A \to B$ 是 $R$-代数同态。
若 $A$ 与 $B$ 都是平展 $R$-代数的滤过余极限，
则 $B$ 是平展 $A$-代数的滤过余极限。
\end{lemma}

\begin{proof}
把 $A = \colim A_i$ 与 $B = \colim B_j$ 写成滤过余极限，
其中 $A_i$ 和 $B_j$ 在 $R$ 上平展。对每个 $i$，
可以找到 $j$，使 $A_i \to B$ 经 $B_j$ 分解，见引理
\ref{algebra-lemma-characterize-finite-presentation}。
由引理 \ref{algebra-lemma-map-between-etale}，分解映射
$A_i \to B_j$ 是平展的。由于
$A \to A \otimes_{A_i} B_j$ 平展（引理 \ref{algebra-lemma-etale}），
只须证明
$B = \colim A \otimes_{A_i} B_j$，其中余极限取遍
二元组 $(i, j)$ 以及 $A_i \to B$ 的分解
$A_i \to B_j \to B$（这是一个有向系统；细节从略）。
这是清楚的，因为余极限与张量积交换，故
$\colim A \otimes_{A_i} B_j = A \otimes_A B = B$。
\end{proof}

\begin{lemma}
\label{algebra-lemma-map-into-henselian-colimit}
设 $R \to S$ 是环映射，且 $S$ 是 Hensel 局部环。给定
\begin{enumerate}
\item 一个 $R$-代数 $A$，它是平展 $R$-代数的滤过余极限；
\item $A$ 的一个素理想 $\mathfrak q$，它位于
$\mathfrak p = R \cap \mathfrak m_S$ 之上；
% P01-E0967: as in P01-E0962, the frozen source uses an intersection for the contraction along a possibly noninjective ring map.
\item 一个 $\kappa(\mathfrak p)$-代数映射
$\tau : \kappa(\mathfrak q) \to S/\mathfrak m_S$。
\end{enumerate}
则存在唯一的 $R$-代数同态 $f : A \to S$，使得
$\mathfrak q = f^{-1}(\mathfrak m_S)$ 且
$f \bmod \mathfrak q = \tau$。
\end{lemma}

\begin{proof}
把 $A = \colim A_i$ 写成平展 $R$-代数的滤过余极限。
置 $\mathfrak q_i = A_i \cap \mathfrak q$。
% P01-E0968: the frozen source again writes an intersection although the transition map A_i -> A need not be injective; the contraction is the inverse image.
应用引理 \ref{algebra-lemma-map-into-henselian}，得到
$f_i : A_i \to S$。置 $f = \colim f_i$。
\end{proof}

\begin{lemma}
\label{algebra-lemma-uniqueness-henselian}
设 $R$ 是环。给定环映射的交换图表
$$
\xymatrix{
S \ar[r] & K \\
R \ar[u] \ar[r] & S' \ar[u]
}
$$
其中 $S$、$S'$ 是 Hensel 局部环，$S$、$S'$ 是平展
$R$-代数的滤过余极限，$K$ 是域，并且箭头 $S \to K$ 与
$S' \to K$ 都把 $K$ 等同于 $S$、$S'$ 的剩余域。
则存在唯一的 $R$-代数同构 $S \to S'$，它与到 $K$ 的映射相容。
\end{lemma}

\begin{proof}
这立即由引理 \ref{algebra-lemma-map-into-henselian-colimit} 得出。
\end{proof}

\noindent
下面的引理严格说来并不涉及平展环映射的余极限。

\begin{lemma}
\label{algebra-lemma-colimit-henselian}
沿局部同态所取的（严格）Hensel 局部环的滤过余极限仍是
（严格）Hensel 的。
\end{lemma}

\begin{proof}
《范畴》，引理 \ref{categories-lemma-directed-category-system}
说明，实际上只须考虑有向集上的（严格）Hensel 局部环的余极限。
设 $(R_i, \varphi_{ii'})$ 是这样的系统，每个
$\varphi_{ii'}$ 都是局部的。则 $R = \colim_i R_i$ 是局部环，
其剩余域 $\kappa$ 为 $\colim \kappa_i$（论证从略）。
容易看出，若每个 $\kappa_i$ 都可分代数闭，则
$\colim \kappa_i$ 也可分代数闭；因此只须证明：
若每个 $R_i$ 都是 Hensel 的，则 $R$ 是 Hensel 的。
假设 $f \in R[T]$ 首一，并且 $a_0 \in \kappa$ 是
$\overline{f}$ 的单根。则对某个充分大的 $i$，
存在映到 $f$ 的 $f_i \in R_i[T]$，以及映到 $a_0$ 的
$a_{0, i} \in \kappa_i$。由于
$\overline{f_i}(a_{0, i}) \in \kappa_i$ 与
$\overline{f_i'}(a_{0, i}) \in \kappa_i$ 分别映到
$0 = \overline{f}(a_0) \in \kappa$ 与
$0 \not = \overline{f'}(a_0) \in \kappa$，
可知 $a_{0, i}$ 是 $\overline{f_i}$ 的单根。
由于 $R_i$ 是 Hensel 的，可以找到 $a_i \in R_i$，使得
$f_i(a_i) = 0$ 且 $a_{0, i} = \overline{a_i}$。
于是 $a_i$ 的像 $a \in R$ 就是所求解。
因此 $R$ 是 Hensel 的。
\end{proof}


\section{Hensel 化与严格 Hensel 化}
\label{algebra-section-henselization}

\noindent
本节构造 Hensel 化。阅读这些材料时，我们建议读者始终记住引理
\ref{algebra-lemma-uniqueness-henselian} 已经证明的唯一性，以及引理
\ref{algebra-lemma-map-into-henselian-colimit} 指出的函子性。

\begin{lemma}
\label{algebra-lemma-henselization}
设 $(R, \mathfrak m, \kappa)$ 是局部环。存在局部环映射
$R \to R^h$，它具有下列性质：
\begin{enumerate}
\item $R^h$ 是 Hensel 的；
\item $R^h$ 是平展 $R$-代数的滤过余极限；
\item $\mathfrak m R^h$ 是 $R^h$ 的极大理想；且
\item $\kappa = R^h/\mathfrak m R^h$。
\end{enumerate}
\end{lemma}

\begin{proof}
考虑由偶对 $(S, \mathfrak q)$ 构成的范畴，其中 $R \to S$
是平展环映射，而 $\mathfrak q$ 是 $S$ 中位于
$\mathfrak m$ 之上的素理想，并且
$\kappa = \kappa(\mathfrak q)$。偶对之间的态
$(S, \mathfrak q) \to (S', \mathfrak q')$ 由
$R$-代数映射 $\varphi : S \to S'$ 给出，并且满足
$\varphi^{-1}(\mathfrak q') = \mathfrak q$。置
$$
R^h = \colim_{(S, \mathfrak q)} S.
$$
我们来证明这个偶对范畴是滤过的；见《范畴》，定义
\ref{categories-definition-directed}。该范畴含有偶对
$(R, \mathfrak m)$，因而非空；这证明了《范畴》，定义
\ref{categories-definition-directed} 的 (1)。
对任意偶对 $(S, \mathfrak q)$，素理想 $\mathfrak q$
都是极大理想，且其剩余域为 $\kappa$，因为复合
$\kappa \to S/\mathfrak q \to \kappa(\mathfrak q)$ 是同构。
假设 $(S, \mathfrak q)$ 与 $(S', \mathfrak q')$ 是两个对象。
置
$S'' = S \otimes_R S'$ 且
$\mathfrak q'' = \mathfrak qS'' + \mathfrak q'S''$。
由上面的说明，
$S''/\mathfrak q'' = S/\mathfrak q \otimes_R S'/\mathfrak q' = \kappa$。
此外，由引理 \ref{algebra-lemma-etale}，$R \to S''$ 平展。
这证明了《范畴》，定义 \ref{categories-definition-directed} 的 (2)。

其次，假设
$\varphi, \psi : (S, \mathfrak q) \to (S', \mathfrak q')$
是偶对的两个态。由引理 \ref{algebra-lemma-map-between-etale}，
$\varphi$、$\psi$ 以及 $S' \otimes_R S' \to S'$
都是平展环映射。考虑
$$
S'' = (S' \otimes_{\varphi, S, \psi} S')
\otimes_{S' \otimes_R S'} S'
$$
及其中的素理想
$$
\mathfrak q'' =
(\mathfrak q' \otimes S' + S' \otimes \mathfrak q') \otimes S'
+
(S' \otimes_{\varphi, S, \psi} S') \otimes \mathfrak q'
$$
与上面同样论证（平展映射的基变换是平展的，平展映射的复合也是
平展的），可知 $S''$ 在 $R$ 上平展。此外，典范映射
$S' \to S''$（例如使用最右边的因子）使 $\varphi$ 与 $\psi$
相等。这证明了《范畴》，定义
\ref{categories-definition-directed} 的 (3)。
因此可知，$R^h$ 的元由三元组 $(S, \mathfrak q, f)$ 表示，
其中 $f \in S$；两个这样的三元组
$(S, \mathfrak q, f)$、$(S', \mathfrak q', f')$
定义 $R^h$ 的同一个元，当且仅当存在偶对
$(S'', \mathfrak q'')$ 以及偶对的态
$\varphi : (S, \mathfrak q) \to (S'', \mathfrak q'')$ 和
$\varphi' : (S', \mathfrak q') \to (S'', \mathfrak q'')$，
使得 $\varphi(f) = \varphi'(f')$。

\medskip\noindent
假设 $x \in R^h$。用三元组 $(S, \mathfrak q, f)$ 表示 $x$。
设 $\mathfrak q_1, \ldots, \mathfrak q_r$ 是 $S$ 中位于
$\mathfrak m$ 之上的其他素理想。由于上面已经看到
$\mathfrak q$ 极大，故 $\mathfrak q \not \subset \mathfrak q_i$。
于是，因为 $\mathfrak q$ 是素理想，可以找到
$g \in S$、$g \not \in \mathfrak q$，并且对
$i = 1, \ldots, r$ 有 $g \in \mathfrak q_i$。
考虑偶对的态
$(S, \mathfrak q) \to (S_g, \mathfrak qS_g)$。
这样就看出，总可假设 $x$ 由三元组
$(S, \mathfrak q, f)$ 给出，其中 $\mathfrak q$ 是 $S$ 中
位于 $\mathfrak m$ 之上的唯一素理想，即
$\sqrt{\mathfrak mS} = \mathfrak q$。但由于 $R \to S$
平展，由引理 \ref{algebra-lemma-etale-at-prime}，
$\mathfrak mS_{\mathfrak q} = \mathfrak qS_{\mathfrak q}$。
因此实际上有 $\mathfrak mS = \mathfrak q$。

\medskip\noindent
假设 $x \not \in \mathfrak mR^h$。用满足
$\mathfrak mS = \mathfrak q$ 的三元组
$(S, \mathfrak q, f)$ 表示 $x$。
则 $f \not \in \mathfrak mS$，即 $f \not \in \mathfrak q$。
因此 $(S, \mathfrak q) \to (S_f, \mathfrak qS_f)$
是偶对的态，且 $f$ 的像在其中可逆。
所以 $x$ 可逆，其逆由三元组
$(S_f, \mathfrak qS_f, 1/f)$ 表示。
由此得出 $R^h$ 是局部环，其极大理想为 $\mathfrak mR^h$。
剩余域为 $\kappa$，因为可以定义
$R^h/\mathfrak mR^h \to \kappa$：把三元组
$(S, \mathfrak q, f)$ 映到 $f$ 模 $\mathfrak q$ 的剩余类。

\medskip\noindent
还须证明 $R^h$ 是 Hensel 的。事实上，假设
$P \in R^h[T]$ 是首一多项式，而 $a_0 \in \kappa$
是其约化 $\overline{P} \in \kappa[T]$ 的单根。
则可以找到偶对 $(S, \mathfrak q)$，使得 $P$ 是某个首一多项式
$Q \in S[T]$ 的像。置 $S' = S[T]/(Q)$，并令
$\mathfrak q' \subset S'$ 为极大理想
$\mathfrak q' = \mathfrak qS' + (T - a')S'$，
其中 $a' \in S$ 是 $a_0$ 的任意提升。
由构造，$S \to S'$ 在 $\mathfrak q'$ 处平展，并且
$\kappa = \kappa(\mathfrak q')$。选取
$g \in S'$、$g \not \in \mathfrak q'$，使得
$S'' = S'_g$ 在 $S$ 上平展。于是
$(S, \mathfrak q) \to (S'', \mathfrak q'S'')$ 是偶对的态。
三元组 $(S'', \mathfrak q'S'', \text{下列对象的类： }T)$ 现在确定
$a \in R^h$，满足 $P(a) = 0$ 且
$\overline{a} = a_0$，正如所需。
% P01-E0969: the frozen English source says “Now that triple”, a grammatical substitution of “that” for the article “the”.
\end{proof}

\begin{lemma}
\label{algebra-lemma-strict-henselization}
设 $(R, \mathfrak m, \kappa)$ 是局部环。
设 $\kappa \subset \kappa^{sep}$ 是可分代数闭包。
存在交换图表
$$
\xymatrix{
\kappa \ar[r] & \kappa \ar[r] & \kappa^{sep} \\
R \ar[r] \ar[u] & R^h \ar[r] \ar[u] & R^{sh} \ar[u]
}
$$
它具有下列性质：
\begin{enumerate}
\item 映射 $R^h \to R^{sh}$ 是局部的；
\item $R^{sh}$ 是严格 Hensel 的；
\item $R^{sh}$ 是平展 $R$-代数的滤过余极限；
\item $\mathfrak m R^{sh}$ 是 $R^{sh}$ 的极大理想；且
\item $\kappa^{sep} = R^{sh}/\mathfrak m R^{sh}$。
\end{enumerate}
\end{lemma}

\begin{proof}
证明与引理 \ref{algebra-lemma-henselization} 所用的证明完全相同。
唯一差别在于，不使用偶对，而使用三元组
$(S, \mathfrak q, \alpha)$，其中 $R \to S$ 平展，
$\mathfrak q$ 是 $S$ 中位于 $\mathfrak m$ 之上的素理想，
而 $\alpha : \kappa(\mathfrak q) \to \kappa^{sep}$
是 $\kappa$ 的扩张之间的嵌入。
\end{proof}

\begin{definition}
\label{algebra-definition-henselization}
设 $(R, \mathfrak m, \kappa)$ 是局部环。
\begin{enumerate}
\item 引理 \ref{algebra-lemma-henselization} 中构造的局部环映射
$R \to R^h$ 称为 $R$ 的{\it Hensel 化}。
\item 给定可分代数闭包 $\kappa \subset \kappa^{sep}$，
引理 \ref{algebra-lemma-strict-henselization} 中构造的局部环映射
$R \to R^{sh}$ 称为{\it $R$ 关于
$\kappa \subset \kappa^{sep}$ 的严格 Hensel 化}。
\item 若局部环映射 $R \to R^{sh}$ 同构于引理
\ref{algebra-lemma-strict-henselization} 中构造的某个局部环映射，
则称它为 $R$ 的一个{\it 严格 Hensel 化}
% P01-E0970: the frozen third definition item lacks terminal punctuation.
\end{enumerate}
\end{definition}

\noindent
映射 $R \to R^h \to R^{sh}$ 是平坦局部环同态。
由引理 \ref{algebra-lemma-uniqueness-henselian}，条件
“它们是 Hensel 局部环、是平展 $R$-代数的
滤过余极限，且剩余域分别为 $\kappa$ 与 $\kappa^{sep}$”
把 $R$-代数 $R^h$ 与 $R^{sh}$ 在唯一同构意义下确定。
本节余下部分主要讨论（严格）Hensel 化的函子性。
关于 $R$ 与其 Hensel 化之间关系的更细致结果，
将在《代数进阶》，第
\ref{more-algebra-section-permanence-henselization} 节讨论。

\begin{remark}
\label{algebra-remark-construct-sh-from-h}
也可以由 $R^h$ 构造 $R^{sh}$。事实上，对任意有限可分子扩张
$\kappa^{sep}/\kappa'/\kappa$，存在唯一的（在唯一同构意义下）
有限平展局部环扩张 $R^h \subset R^h(\kappa')$，
其剩余域扩张重现给定扩张；见引理
\ref{algebra-lemma-henselian-cat-finite-etale}。因此可以置
$$
R^{sh} =
\bigcup\nolimits_{\kappa \subset \kappa' \subset \kappa^{sep}}
R^h(\kappa')
$$
该系统中的箭头与剩余域层面上的箭头相容，并由引理
\ref{algebra-lemma-henselian-cat-finite-etale} 保证存在。
由引理 \ref{algebra-lemma-colimit-henselian}，这将产生 Hensel 局部环，
因为每个环 $R^h(\kappa')$ 由引理
\ref{algebra-lemma-finite-over-henselian} 都是 Hensel 的。
由构造，$R^h \to R^{sh}$ 诱导的剩余域扩张
就是域扩张 $\kappa^{sep}/\kappa$。所以这样构造的
$R^{sh}$ 是严格 Hensel 的。由引理
\ref{algebra-lemma-composition-colimit-etale}，$R$-代数 $R^{sh}$
是平展 $R$-代数的余极限。因此，引理
\ref{algebra-lemma-uniqueness-henselian} 的唯一性说明
$R^{sh}$ 就是严格 Hensel 化。
\end{remark}

\begin{lemma}
\label{algebra-lemma-henselian-functorial-prepare}
设 $R \to S$ 是局部环之间的局部映射。
设 $S \to S^h$ 是 Hensel 化。设 $R \to A$ 是平展环映射，
并设 $\mathfrak q$ 是 $A$ 中位于 $\mathfrak m_R$ 之上的素理想，
使得 $R/\mathfrak m_R \cong \kappa(\mathfrak q)$。
则存在唯一的环态 $f : A \to S^h$，使交换图表
$$
\xymatrix{
A \ar[r]_f & S^h \\
R \ar[u] \ar[r] & S \ar[u]
}
$$
成立，并且 $f^{-1}(\mathfrak m_{S^h}) = \mathfrak q$。
\end{lemma}

\begin{proof}
这是引理 \ref{algebra-lemma-map-into-henselian} 的特殊情形。
\end{proof}

\begin{lemma}
\label{algebra-lemma-henselian-functorial}
设 $R \to S$ 是局部环之间的局部映射。
设 $R \to R^h$ 与 $S \to S^h$ 是 Hensel 化。
存在唯一的局部环映射 $R^h \to S^h$，使交换图表
$$
\xymatrix{
R^h \ar[r]_f & S^h \\
R \ar[u] \ar[r] & S \ar[u]
}
$$
成立。
\end{lemma}

\begin{proof}
这立即由引理 \ref{algebra-lemma-map-into-henselian-colimit} 得出。
\end{proof}

\noindent
下面给出 Hensel 化的一个略有不同的构造。

\begin{lemma}
\label{algebra-lemma-henselization-different}
设 $R$ 是环，$\mathfrak p \subset R$ 是素理想。
考虑由偶对 $(S, \mathfrak q)$ 构成的范畴，其中
$R \to S$ 平展，而 $\mathfrak q$ 是位于 $\mathfrak p$
之上的素理想，使得
$\kappa(\mathfrak p) = \kappa(\mathfrak q)$。
这个范畴是滤过的，并且典范地有
$$
(R_{\mathfrak p})^h = \colim_{(S, \mathfrak q)} S
= \colim_{(S, \mathfrak q)} S_{\mathfrak q}
$$
\end{lemma}

\begin{proof}
偶对之间的态 $(S, \mathfrak q) \to (S', \mathfrak q')$
由 $R$-代数映射 $\varphi : S \to S'$ 给出，并且满足
$\varphi^{-1}(\mathfrak q') = \mathfrak q$。
我们来证明偶对范畴是滤过的；见《范畴》，定义
\ref{categories-definition-directed}。该范畴含有偶对
$(R, \mathfrak p)$，因而非空；这证明了《范畴》，定义
\ref{categories-definition-directed} 的 (1)。
假设 $(S, \mathfrak q)$ 与 $(S', \mathfrak q')$ 是两个偶对。
注意，$\mathfrak q$ 与 $\mathfrak q'$ 分别对应于纤维环
$S \otimes \kappa(\mathfrak p)$ 与
$S' \otimes \kappa(\mathfrak p)$ 的素理想，其剩余域为
$\kappa(\mathfrak p)$；因此它们分别对应于
$S \otimes \kappa(\mathfrak p)$ 与
$S' \otimes \kappa(\mathfrak p)$ 的极大理想。
置 $S'' = S \otimes_R S'$。由上可知，存在唯一的素理想
$\mathfrak q'' \subset S''$，它位于 $\mathfrak q$ 与
$\mathfrak q'$ 之上，并且剩余域为 $\kappa(\mathfrak p)$。
由引理 \ref{algebra-lemma-etale}，环映射 $R \to S''$ 平展。
这证明了《范畴》，定义 \ref{categories-definition-directed} 的 (2)。

其次，假设
$\varphi, \psi : (S, \mathfrak q) \to (S', \mathfrak q')$
是偶对的两个态。由引理 \ref{algebra-lemma-map-between-etale}，
$\varphi$、$\psi$ 以及 $S' \otimes_R S' \to S'$
都是平展环映射。考虑
$$
S'' = (S' \otimes_{\varphi, S, \psi} S')
\otimes_{S' \otimes_R S'} S'
$$
与 $S''$ 在 $\mathfrak p$ 上的纤维环
$$
F'' = (F' \otimes_{\varphi, F, \psi} F')
\otimes_{F' \otimes_{\kappa(\mathfrak p)} F'} F'
$$
其中 $F', F$ 分别是 $S'$、$S$ 的纤维环。
与上面同样论证（平展映射的基变换是平展的，平展映射的复合也是
平展的），可知 $S''$ 在 $R$ 上平展。
由于 $\varphi$ 与 $\psi$ 是偶对的态，对应于
$\mathfrak p'$ 的映射 $F' \to \kappa(\mathfrak p)$
延拓为映射 $F'' \to \kappa(\mathfrak p)$，进而对应于素理想
$\mathfrak q'' \subset S''$，其剩余域为
$\kappa(\mathfrak p)$。
% P01-E0971: the frozen source refers to an undefined prime p' here; the pair carries q', which is the prime corresponding to the residue map.
典范映射 $S' \to S''$（例如使用最右边的因子）是偶对的态
$(S', \mathfrak q') \to (S'', \mathfrak q'')$，
并使 $\varphi$ 与 $\psi$ 相等。这证明了《范畴》，定义
\ref{categories-definition-directed} 的 (3)。
因此该范畴是滤过的。

\medskip\noindent
回忆在引理 \ref{algebra-lemma-henselization} 的证明中，我们把
$(R_{\mathfrak p})^h$ 构造成相应的余极限，不过是从
$R_{\mathfrak p}$ 及其极大理想
$\mathfrak pR_{\mathfrak p}$ 出发。现在，对
$(R, \mathfrak p)$ 的任意偶对 $(S, \mathfrak q)$，
得到 $(R_{\mathfrak p}, \mathfrak pR_{\mathfrak p})$ 的偶对
$(S_{\mathfrak p}, \mathfrak qS_{\mathfrak p})$。
此外，在这种情况下
$$
S_{\mathfrak p} = \colim_{f \in R, f \not \in \mathfrak p} S_f.
$$
所以，为了证明引理中的等式，只须证明
$(R_{\mathfrak p}, \mathfrak pR_{\mathfrak p})$ 的任意偶对
$(S_{loc}, \mathfrak q_{loc})$ 都具有形式
$(S_{\mathfrak p}, \mathfrak qS_{\mathfrak p})$，
其中 $(S, \mathfrak q)$ 是
$(R, \mathfrak p)$ 上的某个偶对（略去一些细节）。
这由引理 \ref{algebra-lemma-etale} 得出。
\end{proof}

\begin{lemma}
\label{algebra-lemma-henselian-functorial-improve}
设 $R \to S$ 是环映射。设 $\mathfrak q \subset S$
是位于 $\mathfrak p \subset R$ 之上的素理想。
设 $R \to R^h$ 与 $S \to S^h$ 分别是
$R_\mathfrak p$ 与 $S_\mathfrak q$ 的 Hensel 化。
引理 \ref{algebra-lemma-henselian-functorial} 的局部环映射
$R^h \to S^h$ 把 $S^h$ 等同于
$R^h \otimes_R S$ 在位于 $\mathfrak m^h$ 与
$\mathfrak q$ 之上的唯一素理想处的 Hensel 化。
\end{lemma}

\begin{proof}
由引理 \ref{algebra-lemma-henselization-different}，可知
$R^h$ 与 $S^h$ 分别是平展 $R$-代数与平展 $S$-代数的滤过余极限。
所以 $R^h \otimes_R S$ 是平展 $S$-代数 $A_i$ 的滤过余极限
（引理 \ref{algebra-lemma-etale}）。由引理
\ref{algebra-lemma-colimits-of-etale}，可知 $S^h$ 是平展
$R^h \otimes_R S$-代数的滤过余极限。
又因为 $S^h$ 是 Hensel 局部环，且剩余域等于
$\kappa(\mathfrak q)$，所述结论由引理
\ref{algebra-lemma-uniqueness-henselian} 的唯一性结果得出。
\end{proof}

\begin{lemma}
\label{algebra-lemma-strictly-henselian-functorial-prepare}
设 $\varphi : R \to S$ 是局部环之间的局部映射。
设 $S/\mathfrak m_S \subset \kappa^{sep}$ 是可分代数闭包。
设 $S \to S^{sh}$ 是 $S$ 关于
$S/\mathfrak m_S \subset \kappa^{sep}$ 的严格 Hensel 化。
设 $R \to A$ 是平展环映射，并设 $\mathfrak q$
是 $A$ 中位于 $\mathfrak m_R$ 之上的素理想。
给定任意交换图表
$$
\xymatrix{
\kappa(\mathfrak q) \ar[r]_{\phi} & \kappa^{sep} \\
R/\mathfrak m_R \ar[r]^{\varphi} \ar[u] & S/\mathfrak m_S \ar[u]
}
$$
存在唯一的环态 $f : A \to S^{sh}$，使交换图表
$$
\xymatrix{
A \ar[r]_f & S^{sh} \\
R \ar[u] \ar[r]^{\varphi} & S \ar[u]
}
$$
成立，并且 $f^{-1}(\mathfrak m_{S^h}) = \mathfrak q$，
而所诱导的映射 $\kappa(\mathfrak q) \to \kappa^{sep}$
就是给定映射。
% P01-E0972: the frozen strict-Henselization statement takes the inverse image of m_{S^h}; the codomain is S^{sh}, so the ideal should be m_{S^{sh}}.
\end{lemma}

\begin{proof}
这是引理 \ref{algebra-lemma-map-into-henselian} 的特殊情形。
\end{proof}

\begin{lemma}
\label{algebra-lemma-strictly-henselian-functorial}
设 $R \to S$ 是局部环之间的局部映射。选取可分代数闭包
$R/\mathfrak m_R \subset \kappa_1^{sep}$
以及
$S/\mathfrak m_S \subset \kappa_2^{sep}$。
设 $R \to R^{sh}$ 与 $S \to S^{sh}$ 是相应的严格
Hensel 化。给定任意交换图表
$$
\xymatrix{
\kappa_1^{sep} \ar[r]_{\phi} & \kappa_2^{sep} \\
R/\mathfrak m_R \ar[r]^{\varphi} \ar[u] & S/\mathfrak m_S \ar[u]
}
$$
存在唯一的局部环映射 $R^{sh} \to S^{sh}$，使交换图表
$$
\xymatrix{
R^{sh} \ar[r]_f & S^{sh} \\
R \ar[u] \ar[r] & S \ar[u]
}
$$
成立，并在 $R^{sh}$ 与 $S^{sh}$ 的剩余域上诱导 $\phi$。
% P01-E0973: the frozen English sentence beginning “Given any commutative diagram” is followed by capitalized “There exists” without completing or punctuating the introductory clause.
\end{lemma}

\begin{proof}
这立即由引理 \ref{algebra-lemma-map-into-henselian-colimit} 得出。
\end{proof}

\begin{lemma}
\label{algebra-lemma-strict-henselization-different}
设 $R$ 是环，$\mathfrak p \subset R$ 是素理想。
设 $\kappa(\mathfrak p) \subset \kappa^{sep}$ 是可分代数闭包。
考虑由三元组 $(S, \mathfrak q, \phi)$ 构成的范畴，其中
$R \to S$ 平展，$\mathfrak q$ 是位于 $\mathfrak p$
之上的素理想，而
$\phi : \kappa(\mathfrak q) \to \kappa^{sep}$
是 $\kappa(\mathfrak p)$-代数映射。
这个范畴是滤过的，并且典范地有
$$
(R_{\mathfrak p})^{sh} =
\colim_{(S, \mathfrak q, \phi)} S =
\colim_{(S, \mathfrak q, \phi)} S_{\mathfrak q}
$$
\end{lemma}

\begin{proof}
三元组之间的态
$(S, \mathfrak q, \phi) \to (S', \mathfrak q', \phi')$
由 $R$-代数映射 $\varphi : S \to S'$ 给出，并满足
$\varphi^{-1}(\mathfrak q') = \mathfrak q$ 以及
$\phi' \circ \varphi = \phi$。
我们来证明偶对范畴是滤过的；见《范畴》，定义
\ref{categories-definition-directed}。
% P01-E0974: the frozen proof says “category of pairs” although the objects and morphisms in this lemma are triples.
该范畴含有三元组
$(R, \mathfrak p, \kappa(\mathfrak p) \subset \kappa^{sep})$，
因而非空；这证明了《范畴》，定义
\ref{categories-definition-directed} 的 (1)。
假设 $(S, \mathfrak q, \phi)$ 与
$(S', \mathfrak q', \phi')$ 是两个三元组。
注意，$\mathfrak q$ 与 $\mathfrak q'$ 分别对应于纤维环
$S \otimes \kappa(\mathfrak p)$ 与
$S' \otimes \kappa(\mathfrak p)$ 的素理想，其剩余域是
$\kappa(\mathfrak p)$ 的有限可分扩张；而 $\phi$ 与 $\phi'$
分别对应于到 $\kappa^{sep}$ 的映射。因此这些数据对应于
$\kappa(\mathfrak p)$-代数映射
$$
\phi : S \otimes_R \kappa(\mathfrak p) \longrightarrow \kappa^{sep},
\quad
\phi' : S' \otimes_R \kappa(\mathfrak p) \longrightarrow \kappa^{sep}.
$$
置 $S'' = S \otimes_R S'$。组合上述映射，得到唯一的
$\kappa(\mathfrak p)$-代数映射
% P01-E0975: the frozen source says “Combining the maps the above”, with the article and modifier in the wrong order.
$$
\phi'' = \phi \otimes \phi' :
S'' \otimes_R \kappa(\mathfrak p)
\longrightarrow
\kappa^{sep}
$$
其核对应于素理想 $\mathfrak q'' \subset S''$，
它位于 $\mathfrak q$ 与 $\mathfrak q'$ 之上；其剩余域经
$\phi''$ 映到 $\kappa^{sep}$ 中
$\phi(\kappa(\mathfrak q))$ 与
$\phi'(\kappa(\mathfrak q'))$ 的合成域。
由引理 \ref{algebra-lemma-etale}，环映射 $R \to S''$ 平展。
因此 $(S'', \mathfrak q'', \phi'')$ 是同时支配
$(S, \mathfrak q, \phi)$ 与 $(S', \mathfrak q', \phi')$
的三元组。这证明了《范畴》，定义
\ref{categories-definition-directed} 的 (2)。

其次，假设
$\varphi, \psi : (S, \mathfrak q, \phi) \to
(S', \mathfrak q', \phi')$ 是偶对的两个态。
% P01-E0976: the frozen source again says “morphisms of pairs” although these are morphisms of triples.
由引理 \ref{algebra-lemma-map-between-etale}，$\varphi$、$\psi$ 以及
$S' \otimes_R S' \to S'$ 都是平展环映射。考虑
$$
S'' = (S' \otimes_{\varphi, S, \psi} S')
\otimes_{S' \otimes_R S'} S'
$$
与 $S''$ 在 $\mathfrak p$ 上的纤维环
$$
F'' = (F' \otimes_{\varphi, F, \psi} F')
\otimes_{F' \otimes_{\kappa(\mathfrak p)} F'} F'
$$
其中 $F', F$ 是 $S'$、$S$ 的纤维环。
与上面同样论证（平展映射的基变换是平展的，平展映射的复合也是
平展的），可知 $S''$ 在 $R$ 上平展。
由于 $\varphi$ 与 $\psi$ 是三元组的态，映射
$\phi' : F' \to \kappa^{sep}$ 延拓为映射
$\phi'' : F'' \to \kappa^{sep}$，后者进而对应于素理想
$\mathfrak q'' \subset S''$。典范映射 $S' \to S''$
（例如使用最右边的因子）是三元组之间的态
$(S', \mathfrak q', \phi') \to
(S'', \mathfrak q'', \phi'')$，并使 $\varphi$ 与 $\psi$
相等。这证明了《范畴》，定义
\ref{categories-definition-directed} 的 (3)。
因此该范畴是滤过的。

\medskip\noindent
还须证明，该系统的余极限 $R_{colim}$ 等于
$R_{\mathfrak p}$ 关于 $\kappa^{sep}$ 的严格 Hensel 化。
为此，注意三元组 $(S, \mathfrak q, \phi)$ 的系统含有引理
\ref{algebra-lemma-henselization-different} 的偶对
$(S, \mathfrak q)$ 作为子系统。因此，由该引理的结果，
$R_{colim}$ 含有 $R_{\mathfrak p}^h$。
此外，很清楚，$R_{\mathfrak p}^h \subset R_{colim}$
是平展环扩张的有向余极限。由引理
\ref{algebra-lemma-finite-over-henselian} 与
\ref{algebra-lemma-colimit-henselian}，可知 $R_{colim}$ 是 Hensel 的。
最后，由引理 \ref{algebra-lemma-make-etale-map-prescribed-residue-field}，
可知 $R_{colim}$ 的剩余域等于 $\kappa^{sep}$。
因此 $R_{colim}$ 严格 Hensel，从而如所需，等于
$R_{\mathfrak p}$ 的严格 Hensel 化。略去一些细节。
\end{proof}

\begin{lemma}
\label{algebra-lemma-strictly-henselian-functorial-improve}
设 $R \to S$ 是环映射。设 $\mathfrak q \subset S$
是位于 $\mathfrak p \subset R$ 之上的素理想。
选取可分代数闭包
$\kappa(\mathfrak p) \subset \kappa_1^{sep}$
以及
$\kappa(\mathfrak q) \subset \kappa_2^{sep}$。
设 $R^{sh}$ 与 $S^{sh}$ 分别为
$R_\mathfrak p$ 与 $S_\mathfrak q$ 的相应严格 Hensel 化。
给定任意交换图表
$$
\xymatrix{
\kappa_1^{sep} \ar[r]_{\phi} & \kappa_2^{sep} \\
\kappa(\mathfrak p) \ar[r]^{\varphi} \ar[u] & \kappa(\mathfrak q) \ar[u]
}
$$
引理 \ref{algebra-lemma-strictly-henselian-functorial} 的局部环映射
$R^{sh} \to S^{sh}$ 把 $S^{sh}$ 等同于
$R^{sh} \otimes_R S$ 在一个位于 $\mathfrak q$ 以及极大理想
$\mathfrak m^{sh} \subset R^{sh}$ 之上的素理想处的严格
Hensel 化。
\end{lemma}

\begin{proof}
证明与引理 \ref{algebra-lemma-henselian-functorial-improve} 的证明相同，
只是使用引理 \ref{algebra-lemma-strict-henselization-different}，
而不是引理 \ref{algebra-lemma-henselization-different}。
\end{proof}

\begin{lemma}
\label{algebra-lemma-sh-from-h-map}
设 $R \to S$ 是环映射。设 $\mathfrak q \subset S$
是位于 $\mathfrak p \subset R$ 之上的素理想，使得
$\kappa(\mathfrak p) \to \kappa(\mathfrak q)$ 是同构。
选取 $\kappa(\mathfrak p) = \kappa(\mathfrak q)$ 的一个
可分代数闭包 $\kappa^{sep}$。则
$$
(S_\mathfrak q)^{sh} =
(S_\mathfrak q)^h \otimes_{(R_\mathfrak p)^h} (R_\mathfrak p)^{sh}
$$
\end{lemma}

\begin{proof}
这由备注 \ref{algebra-remark-construct-sh-from-h} 中局部环严格
Hensel 化的另一种构造以及剩余域相等这一事实得出。
略去一些细节。
\end{proof}

\section{Hensel 化与拟有限环映射}
\label{algebra-section-henselization-quasi-finite}

\noindent
本节证明关于拟有限环映射的（严格）Hensel 化之间函子性映射的一些结果。

\begin{lemma}
\label{algebra-lemma-quasi-finite-henselization}
设 $R \to S$ 是环映射。
设 $S$ 的素理想 $\mathfrak q$ 位于 $R$ 的素理想 $\mathfrak p$ 之上。
假设 $R \to S$ 在 $\mathfrak q$ 处是拟有限的。
引理 \ref{algebra-lemma-henselian-functorial} 的交换图表
$$
\xymatrix{
R_{\mathfrak p}^h \ar[r] & S_{\mathfrak q}^h \\
R_{\mathfrak p} \ar[u] \ar[r] & S_{\mathfrak q} \ar[u]
}
$$
把 $S_{\mathfrak q}^h$ 等同于
$R_{\mathfrak p}^h \otimes_{R_{\mathfrak p}} S_{\mathfrak q}$
在由 $\mathfrak q$ 生成的素理想处的局部化。此外，环映射
$R_{\mathfrak p}^h \to S_{\mathfrak q}^h$ 是有限的。
\end{lemma}

\begin{proof}
注意，$R_{\mathfrak p}^h \otimes_R S$ 在与 $\mathfrak q$
对应的素理想处是 $R_{\mathfrak p}^h$ 上拟有限的，见引理
\ref{algebra-lemma-four-rings}。因此，
$R_{\mathfrak p}^h \otimes_{R_{\mathfrak p}} S_{\mathfrak q}$
的局部化 $S'$ 是 Hensel 的，并且在 $R_{\mathfrak p}^h$
上有限，见引理 \ref{algebra-lemma-finite-over-henselian}。
作为一个局部化，$S'$ 是一族 'etale
$R_{\mathfrak p}^h \otimes_{R_{\mathfrak p}} S_{\mathfrak q}$-代数的
滤过余极限。由引理 \ref{algebra-lemma-henselian-functorial-improve}，
可知 $S_\mathfrak q^h$ 是
$R_{\mathfrak p}^h \otimes_{R_{\mathfrak p}} S_{\mathfrak q}$
的 Hensel 化。因此，由引理 \ref{algebra-lemma-uniqueness-henselian}
的唯一性结果可得 $S' = S_\mathfrak q^h$。
\end{proof}

\begin{lemma}
\label{algebra-lemma-quotient-henselization}
\begin{slogan}
Hensel 化与商相容。
\end{slogan}
设 $R$ 是局部环，其 Hensel 化为 $R^h$。
设 $I \subset \mathfrak m_R$。则 $R^h/IR^h$ 是 $R/I$ 的
Hensel 化。
\end{lemma}

\begin{proof}
这是引理 \ref{algebra-lemma-quasi-finite-henselization} 的一个特例。
\end{proof}

\begin{lemma}
\label{algebra-lemma-quasi-finite-strict-henselization}
设 $R \to S$ 是环映射。
设 $S$ 的素理想 $\mathfrak q$ 位于 $R$ 的素理想 $\mathfrak p$ 之上。
假设 $R \to S$ 在 $\mathfrak q$ 处是拟有限的。
设 $\kappa_2^{sep}/\kappa(\mathfrak q)$ 是可分代数闭包，并以
$\kappa_1^{sep} \subset \kappa_2^{sep}$ 表示其中在
$\kappa(\mathfrak p)$ 上可分代数的元素所成的子域
（域，引理 \ref{fields-lemma-separable-first}）。
引理 \ref{algebra-lemma-strictly-henselian-functorial} 的交换图表
$$
\xymatrix{
R_{\mathfrak p}^{sh} \ar[r] & S_{\mathfrak q}^{sh} \\
R_{\mathfrak p} \ar[u] \ar[r] & S_{\mathfrak q} \ar[u]
}
$$
把 $S_{\mathfrak q}^{sh}$ 等同于
$R_{\mathfrak p}^{sh} \otimes_{R_{\mathfrak p}} S_{\mathfrak q}$
在下列映射的核所给出的素理想处的局部化：
$$
R_{\mathfrak p}^{sh} \otimes_{R_{\mathfrak p}} S_{\mathfrak q}
\longrightarrow
\kappa_1^{sep} \otimes_{\kappa(\mathfrak p)} \kappa(\mathfrak q)
\longrightarrow
\kappa_2^{sep}
$$
此外，环映射 $R_{\mathfrak p}^{sh} \to S_{\mathfrak q}^{sh}$
是局部环之间的有限局部同态，其剩余域扩张为
$\kappa_2^{sep}/\kappa_1^{sep}$；该扩张既有限又纯不可分。
\end{lemma}

\begin{proof}
由于 $R \to S$ 在 $\mathfrak q$ 处拟有限，由定义
\ref{algebra-definition-quasi-finite} 和引理
\ref{algebra-lemma-isolated-point-fibre} 可知扩张
$\kappa(\mathfrak q)/\kappa(\mathfrak p)$ 是有限的。
因此，$\kappa_1^{sep}$ 是 $\kappa(\mathfrak p)$ 的一个可分代数闭包
（略去一个小细节）。特别地，引理
\ref{algebra-lemma-strictly-henselian-functorial} 确实适用。
其次，$\kappa(\mathfrak q)$ 与 $\kappa_1^{sep}$ 在
$\kappa_2^{sep}$ 中的合成域是可分代数闭的，因而等于
$\kappa_2^{sep}$。由此可知
$\kappa_2^{sep}/\kappa_1^{sep}$ 是有限的。按构造，扩张
$\kappa_2^{sep}/\kappa_1^{sep}$ 是纯不可分的。
环映射 $R_{\mathfrak p}^{sh} \to S_{\mathfrak q}^{sh}$
确为局部的，并且诱导确为有限纯不可分的剩余域扩张
$\kappa_2^{sep}/\kappa_1^{sep}$。

\medskip\noindent
注意，$R_{\mathfrak p}^{sh} \otimes_R S$ 在引理陈述中给出的素理想
$\mathfrak q'$ 处是 $R_{\mathfrak p}^{sh}$ 上拟有限的，见引理
\ref{algebra-lemma-four-rings}。因此，
$R_{\mathfrak p}^{sh} \otimes_{R_{\mathfrak p}} S_{\mathfrak q}$
在 $\mathfrak q'$ 处的局部化 $S'$ 是 Hensel 的，并且在
$R_{\mathfrak p}^{sh}$ 上有限，见引理
\ref{algebra-lemma-finite-over-henselian}。
注意，$S'$ 的剩余域是 $\kappa_2^{sep}$，因为由上一段的讨论，映射
$\kappa_1^{sep} \otimes_{\kappa(\mathfrak p)} \kappa(\mathfrak q)
\to \kappa_2^{sep}$ 是满射。
此外，作为一个局部化，$S'$ 是一族 'etale
$R_{\mathfrak p}^{sh} \otimes_{R_{\mathfrak p}} S_{\mathfrak q}$-代数的
滤过余极限。由引理
\ref{algebra-lemma-strictly-henselian-functorial-improve}，可知
$S_{\mathfrak q}^{sh}$ 是
$R_{\mathfrak p}^{sh} \otimes_{R_{\mathfrak p}} S_{\mathfrak q}$
在 $\mathfrak q'$ 处的严格 Hensel 化。因此，由引理
\ref{algebra-lemma-uniqueness-henselian} 的唯一性结果可得
$S' = S_\mathfrak q^{sh}$。
\end{proof}

\begin{lemma}
\label{algebra-lemma-quotient-strict-henselization}
设 $R$ 是局部环，其严格 Hensel 化为 $R^{sh}$。
设 $I \subset \mathfrak m_R$。则 $R^{sh}/IR^{sh}$ 是 $R/I$
的一个严格 Hensel 化。
\end{lemma}

\begin{proof}
这是引理 \ref{algebra-lemma-quasi-finite-strict-henselization} 的一个特例。
\end{proof}

\begin{lemma}
\label{algebra-lemma-local-tensor-with-integral}
设 $A \to B$ 和 $A \to C$ 是局部环之间的局部同态。
如果 $A \to C$ 是整的，并且
$\kappa(\mathfrak m_C)/\kappa(\mathfrak m_A)$ 或
$\kappa(\mathfrak m_B)/\kappa(\mathfrak m_A)$ 是纯不可分的，
则 $D = B \otimes_A C$ 是局部环，且 $B \to D$ 和
$C \to D$ 都是局部的。
\end{lemma}

\begin{proof}
由整环映射 $B \to D$ 的上升定理（引理
\ref{algebra-lemma-integral-going-up}），$D$ 的任意极大理想均位于
$B$ 的极大理想之上。现在
$D/\mathfrak m_B D = \kappa(\mathfrak m_B) \otimes_A C =
\kappa(\mathfrak m_B) \otimes_{\kappa(\mathfrak m_A)} C/\mathfrak m_A C$。
$C/\mathfrak m_A C$ 的谱只含一个点，即 $\mathfrak m_C$。
因此，$D/\mathfrak m_B D$ 的谱与
$\kappa(\mathfrak m_B) \otimes_{\kappa(\mathfrak m_A)}
\kappa(\mathfrak m_C)$ 的谱相同。根据我们的假设，
$\kappa(\mathfrak m_C)/\kappa(\mathfrak m_A)$ 或
$\kappa(\mathfrak m_B)/\kappa(\mathfrak m_A)$ 是纯不可分的，
故后一个谱只含一个点。这就证明了 $D$ 是局部的，并且环映射
$B \to D$ 和 $C \to D$ 都是局部的。
\end{proof}

\begin{lemma}
\label{algebra-lemma-base-change-strict-henselization-quasi-finite}
设 $A \to B$ 和 $A \to C$ 是环映射。设 $\kappa$ 是可分代数闭域，
并设 $B \otimes_A C \to \kappa$ 是环同态。以下列图表表示严格
Hensel 化之间相应的映射（见证明）：
$$
\xymatrix{
B^{sh} \ar[r] & (B \otimes_A C)^{sh} \\
A^{sh} \ar[u] \ar[r] & C^{sh} \ar[u]
}
$$
如果
\begin{enumerate}
\item $A \to B$ 在素理想
$\mathfrak p_B = \Ker(B \to \kappa)$ 处是拟有限的；或
\item $B$ 是拟有限 $A$-代数的滤过余极限；或
\item $B_{\mathfrak p_B}$ 是 $A_{\mathfrak p_A}$ 上拟有限代数的
滤过余极限；或
\item $B$ 在 $A$ 上整，
\end{enumerate}
则 $B^{sh} \otimes_{A^{sh}} C^{sh} \to (B \otimes_A C)^{sh}$
是同构。
\end{lemma}

\begin{proof}
记 $D = B \otimes_A C$。记
$\mathfrak p_A = \Ker(A \to \kappa)$，并类似地定义
$\mathfrak p_B$、$\mathfrak p_C$ 和 $\mathfrak p_D$。
以 $\kappa_A \subset \kappa$ 表示 $\kappa(\mathfrak p_A)$
在 $\kappa$ 中的可分代数闭包，并类似地定义
$\kappa_B$、$\kappa_C$ 和 $\kappa_D$。
以 $A^{sh}$ 表示利用可分代数闭包
$\kappa_A/\kappa(\mathfrak p_A)$ 构造的
$A_{\mathfrak p_A}$ 的严格 Hensel 化。类似地定义
$B^{sh}$、$C^{sh}$ 和 $D^{sh}$。由引理
\ref{algebra-lemma-strictly-henselian-functorial} 的函子性，我们得到引理中的
交换图表。

\medskip\noindent
考虑由该交换图表得到的映射
$$
c : B^{sh} \otimes_{A^{sh}} C^{sh} \to D^{sh} = (B \otimes_A C)^{sh}
$$
如果 $A \to B$ 在
$\mathfrak p_B = \Ker(B \to \kappa)$ 处拟有限，则由引理
\ref{algebra-lemma-four-rings}，环映射 $C \to D$ 在 $\mathfrak p_D$
处拟有限。因此，由引理
\ref{algebra-lemma-quasi-finite-strict-henselization}
（以及引理 \ref{algebra-lemma-base-change-integral}），环映射 $c$
是有限 $C^{sh}$-代数之间的同态，并且
$$
B^{sh} = (B \otimes_A A^{sh})_{\mathfrak q}
\quad\text{且}\quad
D^{sh} = (D \otimes_C C^{sh})_{\mathfrak r} =
(B \otimes_A C^{sh})_{\mathfrak r}
$$
对某些素理想 $\mathfrak q$ 和 $\mathfrak r$ 成立。由于
$$
B^{sh} \otimes_{A^{sh}} C^{sh} =
(B \otimes_A A^{sh})_{\mathfrak q} \otimes_{A^{sh}} C^{sh} =
\text{a localization of }
B \otimes_A C^{sh}
$$
可知 $c$ 的源和靶都是 $B \otimes_A C^{sh}$ 的局部化
（且与该映射相容）。因此，只需证明
$B^{sh} \otimes_{A^{sh}} C^{sh}$ 是局部的（略去一个小细节）。
这由引理 \ref{algebra-lemma-local-tensor-with-integral} 以及下述事实得出：
根据上面已经使用过的引理
\ref{algebra-lemma-quasi-finite-strict-henselization}，
$A^{sh} \to B^{sh}$ 是有限的，其剩余域扩张纯不可分。
这证明了引理的情形 (1)。

\medskip\noindent
在情形 (2) 中，把 $B = \colim B_i$ 写成拟有限 $A$-代数的
滤过余极限。相应地得到 $D = \colim D_i$，其中
$D_i = B_i \otimes_A C$。注意 $B^{sh} = \colim B_i^{sh}$。
事实上，由引理 \ref{algebra-lemma-colimit-henselian}，环
$\colim B_i^{sh}$ 是严格 Hensel 局部环。
由引理 \ref{algebra-lemma-colimit-colimit-etale-better}，
$\colim B_i^{sh}$ 也是一族 'etale $B$-代数的滤过余极限。
最后，$\colim B_i^{sh}$ 的剩余域是
$\kappa(\mathfrak p_B)$ 的一个可分代数闭包（略去细节）。
因此，由定义 \ref{algebra-definition-henselization} 后的讨论可得
$B^{sh} = \colim B_i^{sh}$。类似地，
$D^{sh} = \colim D_i^{sh}$。于是由情形 (1) 可得
$$
D^{sh} = \colim D_i^{sh} = \colim B_i^{sh} \otimes_{A^{sh}} C^{sh} =
B^{sh} \otimes_{A^{sh}} C^{sh}
$$
% P01-E0977: frozen English source has subject--verb disagreement in “filtered colimit commute”.
因为滤过余极限与张量积可交换。

\medskip\noindent
情形 (3)。可以分别在 $\mathfrak p_A$、$\mathfrak p_B$ 和
$\mathfrak p_C$ 处用 $A$、$B$、$C$ 的局部化替换它们。
因此，(3) 由 (2) 得出。

\medskip\noindent
由于整环映射是有限环映射的滤过余极限，(4) 也由 (2) 得出。
\end{proof}

\section{Serre 正规性判据}
\label{algebra-section-serre-criterion}

\noindent
我们引入诺特环的下列性质。

\begin{definition}
\label{algebra-definition-conditions}
设 $R$ 是诺特环，$k \geq 0$ 是整数。
\begin{enumerate}
\item 如果对每个高度 $\leq k$ 的素理想 $\mathfrak p$，局部环
$R_{\mathfrak p}$ 都是正则的，则称 $R$ 具有性质
{\it $(R_k)$}。我们也称 $R$ 在余维 $\leq k$ 处{\it 正则}。
\item 如果对每个素理想 $\mathfrak p$，局部环
$R_{\mathfrak p}$ 的深度至少为
$\min\{k, \dim(R_{\mathfrak p})\}$，则称 $R$ 具有性质
{\it $(S_k)$}。
\item 设 $M$ 是有限 $R$-模。如果对每个素理想
$\mathfrak p$，模 $M_{\mathfrak p}$ 的深度至少为
$\min\{k, \dim(\text{Supp}(M_{\mathfrak p}))\}$，
则称 $M$ 具有性质 $(S_k)$。
\end{enumerate}
\end{definition}

\noindent
任意诺特环都具有性质 $(S_0)$，其上的任意有限模也如此。
我们的约定是零模的深度为 $\infty$（见第
\ref{algebra-section-depth} 节），而空集的维数为 $-\infty$
（见拓扑，第 \ref{topology-section-krull-dimension} 节）；
这些约定保证零模对所有 $k$ 都具有性质 $(S_k)$。

\begin{lemma}
\label{algebra-lemma-criterion-no-embedded-primes}
设 $R$ 是诺特环，$M$ 是有限 $R$-模。下列条件等价：
\begin{enumerate}
\item $M$ 没有嵌入相伴素理想；
\item $M$ 具有性质 $(S_1)$。
\end{enumerate}
\end{lemma}

\begin{proof}
设 $\mathfrak p$ 是 $M$ 的一个嵌入相伴素理想。
则 $M$ 还有另一个相伴素理想 $\mathfrak q$，使得
$\mathfrak p \supset \mathfrak q$。特别地，这意味着
$\dim(\text{Supp}(M_{\mathfrak p})) \geq 1$
（因为 $\mathfrak q$ 也在支集中）。另一方面，
$\mathfrak pR_{\mathfrak p}$ 是 $M_{\mathfrak p}$ 的相伴素理想
（引理 \ref{algebra-lemma-associated-primes-localize}），因而
$\text{depth}(M_{\mathfrak p}) = 0$
（见引理 \ref{algebra-lemma-ideal-nonzerodivisor}）。
换言之，$(S_1)$ 不成立。
反过来，如果 $(S_1)$ 不成立，则存在素理想 $\mathfrak p$，使得
$\dim(\text{Supp}(M_{\mathfrak p})) \geq 1$ 且
$\text{depth}(M_{\mathfrak p}) = 0$。由于
$\text{depth}(M_{\mathfrak p}) = 0$，由引理
\ref{algebra-lemma-associated-primes-localize} 和
\ref{algebra-lemma-ideal-nonzerodivisor} 可知
$\mathfrak p \in \text{Ass}(M)$。
由于 $\dim(\text{Supp}(M_{\mathfrak p})) \geq 1$，
存在 $\mathfrak q \in \text{Supp}(M)$，使得
$\mathfrak q \subset \mathfrak p$ 且
$\mathfrak q \not = \mathfrak p$。可以选取这样的
$\mathfrak q$，使其在 $\text{Supp}(M)$ 中极小。于是由命题
\ref{algebra-proposition-minimal-primes-associated-primes}，有
$\mathfrak q \in \text{Ass}(M)$，因此 $\mathfrak p$
是嵌入相伴素理想。
\end{proof}

\begin{lemma}
\label{algebra-lemma-criterion-reduced}
\begin{slogan}
约化等价于 R0 加 S1。
\end{slogan}
设 $R$ 是诺特环。下列条件等价：
\begin{enumerate}
\item $R$ 是约化的；
\item $R$ 具有性质 $(R_0)$ 和 $(S_1)$。
\end{enumerate}
\end{lemma}

\begin{proof}
假设 $R$ 是约化的。根据引理
\ref{algebra-lemma-minimal-prime-reduced-ring}，对 $R$ 的每个极小素理想
$\mathfrak p$，$R_{\mathfrak p}$ 都是域。因此，$(R_0)$ 成立。
设 $\mathfrak p$ 是高度 $\geq 1$ 的素理想。则
$A = R_{\mathfrak p}$ 是维数 $\geq 1$ 的约化局部环。
因此，它的极大理想 $\mathfrak m$ 不是相伴素理想，因为否则将存在
零化子为 $\mathfrak m$ 的 $x \in \mathfrak m$，从而
$x^2 = 0$。于是由引理 \ref{algebra-lemma-ass-zero-divisors}，
$A = R_{\mathfrak p}$ 的深度至少为一。这表明 $(S_1)$ 成立。

\medskip\noindent
反过来，假设 $R$ 满足 $(R_0)$ 和 $(S_1)$。
如果 $\mathfrak p$ 是 $R$ 的极小素理想，则由 $(R_0)$，
$R_{\mathfrak p}$ 是域，因而是约化的。
如果 $\mathfrak p$ 不是极小的，则由 $(S_1)$ 可知
$R_{\mathfrak p}$ 的深度 $\geq 1$，因而存在元素
$t \in \mathfrak pR_{\mathfrak p}$，使得
$R_{\mathfrak p} \to R_{\mathfrak p}[1/t]$ 是单射。
由引理 \ref{algebra-lemma-characterize-zero-local}，
$R_\mathfrak p[1/t]$ 包含于它在各素理想处局部化所得环的乘积中。
这意味着 $R_{\mathfrak p}$ 是下列各局部化之乘积的子环：
在满足 $\mathfrak p \supset \mathfrak q$ 且
$t \not \in \mathfrak q$ 的素理想处对 $R$ 作局部化。
这些素理想的高度更小，故对高度作归纳可知 $R$ 是约化的。
\end{proof}

\begin{lemma}[Serre 正规性判据]
\label{algebra-lemma-criterion-normal}
\begin{reference}
\cite[IV, Theorem 5.8.6]{EGA}
\end{reference}
\begin{slogan}
正规等价于 R1 加 S2。
\end{slogan}
设 $R$ 是诺特环。下列条件等价：
\begin{enumerate}
\item $R$ 是正规环；
\item $R$ 具有性质 $(R_1)$ 和 $(S_2)$。
\end{enumerate}
\end{lemma}

\begin{proof}
证明 (1) $\Rightarrow$ (2)。假设 $R$ 正规，即在各素理想处的所有
局部化 $R_{\mathfrak p}$ 都是正规整环。特别地，由引理
\ref{algebra-lemma-criterion-reduced} 可知 $R$ 具有 $(R_0)$ 和
$(S_1)$。因此，只需证明维数为 $d$ 的局部诺特正规整环 $R$
的深度 $\geq \min(2, d)$，并且当 $d = 1$ 时它是正则的。
当 $d = 1$ 时，该断言由引理 \ref{algebra-lemma-characterize-dvr} 得出。

\medskip\noindent
设 $R$ 是极大理想为 $\mathfrak m$、维数为 $d \geq 2$ 的局部
诺特正规整环。对 $R$ 应用引理 \ref{algebra-lemma-hart-serre-loc-thm}。
显然，$R$ 不属于该引理的情形 (1) 或 (2)。
设 $R \to R'$ 如该引理的 (4) 所述。
由于 $R$ 是整环，有 $R \subset R'$。由于 $\mathfrak m$
不是 $R'$ 的相伴素理想，存在 $x \in \mathfrak m$，它在
$R'$ 上是非零因子。于是 $R_x = R'_x$，所以 $R$ 和 $R'$
是具有相同分式域的整环。然而，$R \subset R'$ 的有限性意味着
$R'$ 的每个元素都整于 $R$（引理
\ref{algebra-lemma-finite-is-integral}）；由于 $R$ 正规，故
$R = R'$。这意味着 (4) 不会发生。因而只剩情形 (3)，即
$\text{depth}(R) \geq 2$，如所需。

\medskip\noindent
证明 (2) $\Rightarrow$ (1)。假设 $R$ 满足 $(R_1)$ 和
$(S_2)$。由引理 \ref{algebra-lemma-criterion-reduced} 可知 $R$ 是约化的。
因此，只需证明：如果 $R$ 是维数为 $d$、满足 $(S_2)$ 和
$(R_1)$ 的约化局部诺特环，则 $R$ 是正规整环。
如果 $d = 0$，结论显然。如果 $d = 1$，则结论由引理
\ref{algebra-lemma-characterize-dvr} 得出。

\medskip\noindent
设 $R$ 是极大理想为 $\mathfrak m$、维数为 $d \geq 2$
并且满足 $(R_1)$ 和 $(S_2)$ 的约化局部诺特环。
由引理 \ref{algebra-lemma-characterize-reduced-ring-normal}，只需证明
$R$ 在其全分式环 $Q(R)$ 中整闭。选取整于 $R$ 的
$x \in Q(R)$。则 $R' = R[x]$ 是 $R$ 的有限环扩张
（引理 \ref{algebra-lemma-characterize-finite-in-terms-of-integral}）。
由于对每个非极大素理想 $\mathfrak p \subset R$ 都有
$\dim(R_\mathfrak p) < d$，归纳可得
$R_\mathfrak p = R'_\mathfrak p$。因此，$R'/R$ 的支集为
$\{\mathfrak m\}$。由此可知 $R'/R$ 被 $\mathfrak m$
的某个幂零化（引理
\ref{algebra-lemma-Noetherian-power-ideal-kills-module}）。
由引理 \ref{algebra-lemma-hart-serre-loc-thm}，这与 $R$ 的深度
$\geq 2 = \min(2, d)$ 的假设矛盾，证明完毕。
\end{proof}

\begin{lemma}
\label{algebra-lemma-regular-normal}
正则环是正规的。
\end{lemma}

\begin{proof}
设 $R$ 是正则环。由引理 \ref{algebra-lemma-criterion-normal}，
只需证明 $R$ 满足 $(R_1)$ 和 $(S_2)$。
正则局部环是 Cohen-Macaulay 的，见引理
\ref{algebra-lemma-regular-ring-CM}，故 $R$ 显然满足 $(S_2)$。
性质 $(R_1)$ 立即成立。
\end{proof}

\begin{lemma}
\label{algebra-lemma-normal-domain-intersection-localizations-height-1}
设 $R$ 是分式域为 $K$ 的诺特正规整环。则
\begin{enumerate}
\item 对任意非零 $a \in R$，商 $R/aR$ 没有嵌入素理想，
且它的所有相伴素理想都具有高度 $1$；
% P01-E0978: frozen English item (1) has no terminal punctuation before the next item.
\item
$$
R = \bigcap\nolimits_{\text{高度}(\mathfrak p) = 1} R_{\mathfrak p}
$$
\item 对任意非零 $x \in K$，商 $R/(R \cap xR)$
没有嵌入素理想，且它的所有相伴素理想都具有高度 $1$。
% P01-E0979: frozen English has the malformed phrase “associates primes”.
\end{enumerate}
\end{lemma}

\begin{proof}
由引理 \ref{algebra-lemma-criterion-normal} 可知 $R$ 具有 $(S_2)$。
因此，对任意非零元素 $a \in R$，$R/aR$ 都具有 $(S_1)$
（例如使用引理 \ref{algebra-lemma-depth-in-ses}）。
% P01-E0980: frozen English omits sentence punctuation after the parenthetical citation.
所以 $R/aR$ 没有嵌入素理想
（引理 \ref{algebra-lemma-criterion-no-embedded-primes}）。
由此可知，$R/aR$ 的相伴素理想恰为 $(a)$ 上方的极小素理想
$\mathfrak p$；由于 $a$ 非零，它们的高度为 $1$
（引理 \ref{algebra-lemma-minimal-over-1}）。这证明了 (1)。

\medskip\noindent
因此，给定 $b \in R$，$b \in aR$ 当且仅当对 $(a)$ 上方的
每个极小素理想 $\mathfrak p$ 都有
$b \in aR_{\mathfrak p}$（见引理
\ref{algebra-lemma-zero-at-ass-zero}）。如上所见，这些素理想的高度均为
$1$，故 $b/a \in R$ 当且仅当对所有高度为 1 的素理想都有
$b/a \in R_{\mathfrak p}$。因此 (2) 成立。

\medskip\noindent
为证明 (3)，写 $x = a/b$。设
$\mathfrak p_1, \ldots, \mathfrak p_r$ 是 $(ab)$ 上方的极小素理想。
由上可知，它们都具有高度 1。于是由本引理的 (2) 可知
$R \cap xR = \bigcap_{i = 1, \ldots, r} (R \cap xR_{\mathfrak p_i})$。
因此，$R/(R \cap xR)$ 是
$\bigoplus R/(R \cap xR_{\mathfrak p_i})$ 的子模。
由于 $R_{\mathfrak p_i}$ 是离散赋值环（这是诺特正规整环
$R$ 的性质 $(R_1)$ 所致，见引理
\ref{algebra-lemma-criterion-normal}），对某个 $e_i \in \mathbf{Z}$ 有
$xR_{\mathfrak p_i} = \mathfrak p_i^{e_i}R_{\mathfrak p_i}$。
因此，该直和等于
$\bigoplus_{e_i > 0} R/\mathfrak p_i^{(e_i)}$，见定义
\ref{algebra-definition-symbolic-power}。
由引理 \ref{algebra-lemma-symbolic-power-associated}，模
$R/\mathfrak p^{(n)}$ 的唯一相伴素理想是 $\mathfrak p$。
因此，$R/(R \cap xR)$ 的相伴素理想集合是
$\{\mathfrak p_i\}$ 的子集，且它们之间不存在包含关系。
% P01-E0981: frozen English has “the set of associate primes” instead of “associated primes”.
这证明了 (3)。
\end{proof}

\section{域的形式光滑性}
\label{algebra-section-p-bases}

\noindent
本节将证明，域扩张形式光滑当且仅当它可分。不过，我们先证明：
有限生成域扩张是可分代数扩张，当且仅当它形式非分歧。

\begin{lemma}
\label{algebra-lemma-characterize-separable-algebraic-field-extensions}
设 $K/k$ 是有限生成域扩张。下列条件等价：
\begin{enumerate}
\item $K$ 是 $k$ 的有限可分域扩张；
\item $\Omega_{K/k} = 0$；
\item $K$ 在 $k$ 上形式非分歧；
\item $K$ 在 $k$ 上非分歧；
\item $K$ 在 $k$ 上形式平展；
\item $K$ 在 $k$ 上平展。
\end{enumerate}
\end{lemma}

\begin{proof}
(2) 与 (3) 的等价性是引理
\ref{algebra-lemma-characterize-formally-unramified}。
由引理 \ref{algebra-lemma-etale-over-field} 可知 (1) 等价于 (6)。
性质 (6) 蕴含 (5) 和 (4)，而后二者又都蕴含 (3)
（引理 \ref{algebra-lemma-formally-etale-etale}、
\ref{algebra-lemma-unramified} 和
\ref{algebra-lemma-formally-unramified-unramified}）。
因此，只需证明 (2) 蕴含 (1)。
选取有限生成 $k$-子代数 $A \subset K$，使得 $K$ 是整环
$A$ 的分式域。令 $S = A \setminus \{0\}$。
由于 $0 = \Omega_{K/k} = S^{-1}\Omega_{A/k}$
（引理 \ref{algebra-lemma-differentials-localize}），并且
$\Omega_{A/k}$ 是有限生成的（引理
\ref{algebra-lemma-differentials-finitely-generated}），可以用某个局部化
$A_f$ 替换 $A$，从而归约到 $\Omega_{A/k} = 0$ 的情形
（略去细节）。于是 $A$ 在 $k$ 上非分歧，故 $K/k$
有限可分；例如可对 $\mathfrak q = (0)$ 应用引理
\ref{algebra-lemma-unramified-at-prime}。
\end{proof}

\begin{lemma}
\label{algebra-lemma-derivative-zero-pth-power}
设 $k$ 是特征 $p > 0$ 的完美域，$K/k$ 是扩张，且
$a \in K$。则 $\text{d}a = 0$ 在 $\Omega_{K/k}$ 中成立，
当且仅当 $a$ 是一个 $p$ 次幂。
\end{lemma}

\begin{proof}
由引理 \ref{algebra-lemma-colimit-differentials} 可知，存在子域
$k \subset L \subset K$，使得 $L/k$ 是有限生成域扩张，
且 $\text{d}a$ 在 $\Omega_{L/k}$ 中为零。
因此，可以假设 $K$ 是 $k$ 的有限生成域扩张。

\medskip\noindent
选取超越基 $x_1, \ldots, x_r \in K$，使得 $K$ 是
$k(x_1, \ldots, x_r)$ 的有限可分扩张。根据定义，这是可能的，
见定义 \ref{algebra-definition-perfect} 和
\ref{algebra-definition-separable-field-extension}。
我们指出，该结果对纯超越子域
$k(x_1, \ldots, x_r) \subset K$ 成立。事实上，
$$
\Omega_{k(x_1, \ldots, x_r)/k} =
\bigoplus\nolimits_{i = 1}^r k(x_1, \ldots, x_r) \text{d}x_i
$$
并且偏导数全为零的任意有理函数都是 $p$ 次幂。此外，还有
$$
\Omega_{K/k} =
\bigoplus\nolimits_{i = 1}^r K\text{d}x_i
$$
因为 $k(x_1, \ldots, x_r) \subset K$ 有限可分
（略去计算）。假设 $a \in K$ 且 $\text{d}a = 0$
在微分模中成立。根据 $x_i$ 的选取，$a$ 的最小多项式
$P(T) \in k(x_1, \ldots, x_r)[T]$ 是可分的。写成
$$
P(T) = T^d + \sum\nolimits_{i = 1}^d a_i T^{d - i}
$$
从而
$$
0 = \text{d}P(a) = \sum\nolimits_{i = 1}^d a^{d - i}\text{d}a_i
$$
在 $\Omega_{K/k}$ 中成立。由上述对 $\Omega_{K/k}$ 的描述以及
$P$ 是 $a$ 的最小多项式这一事实，可知这蕴含
$\text{d}a_i = 0$。因此，对每个 $i$ 都有
$a_i = b_i^p$。于是由域，引理
\ref{fields-lemma-pth-root} 可知 $a$ 是 $p$ 次幂。
\end{proof}

\begin{lemma}
\label{algebra-lemma-size-extension-pth-roots}
设 $k$ 是特征 $p > 0$ 的域。设
$a_1, \ldots, a_n \in k$，且
$\text{d}a_1, \ldots, \text{d}a_n$ 在
$\Omega_{k/\mathbf{F}_p}$ 中线性无关。则域扩张
$k(a_1^{1/p}, \ldots, a_n^{1/p})$ 在 $k$ 上的次数为
$p^n$。
\end{lemma}

\begin{proof}
对 $n$ 归纳。当 $n = 1$ 时，结论是引理
\ref{algebra-lemma-derivative-zero-pth-power}。
在归纳步骤中，假设
$k(a_1^{1/p}, \ldots, a_{n - 1}^{1/p})$ 在 $k$ 上的次数为
$p^{n - 1}$。需要证明 $a_n$ 在
$k(a_1^{1/p}, \ldots, a_{n - 1}^{1/p})$ 中的像不是 $p$ 次幂。
如果它是，则可写成
\begin{align*}
a_n & =
\left(\sum\nolimits_{I = (i_1, \ldots, i_{n - 1}),\ 0 \leq i_j \leq p - 1}
\lambda_I a_1^{i_1/p} \ldots a_{n - 1}^{i_{n - 1}/p}\right)^p \\
& = \sum\nolimits_{I = (i_1, \ldots, i_{n - 1}),\ 0 \leq i_j \leq p - 1}
\lambda_I^p a_1^{i_1} \ldots a_{n - 1}^{i_{n - 1}}
\end{align*}
应用 $\text{d}$，可知 $\text{d}a_n$ 线性依赖于
$\text{d}a_i$，$i < n$。这产生矛盾。
\end{proof}

\begin{lemma}
\label{algebra-lemma-separable-differentials}
设 $k$ 是特征 $p > 0$ 的域。下列条件等价：
\begin{enumerate}
\item 域扩张 $K/k$ 可分
（见定义 \ref{algebra-definition-separable-field-extension}）；
\item 映射
$K \otimes_k \Omega_{k/\mathbf{F}_p} \to \Omega_{K/\mathbf{F}_p}$
是单射。
\end{enumerate}
\end{lemma}

\begin{proof}
把 $K$ 写成有限生成域扩张 $K_i/k$ 的有向余极限
$K = \colim_i K_i$。按定义，$K$ 可分当且仅当每个
$K_i$ 都在 $k$ 上可分；而由引理
\ref{algebra-lemma-colimit-differentials} 可知，
$K \otimes_k \Omega_{k/\mathbf{F}_p} \to \Omega_{K/\mathbf{F}_p}$
是单射，当且仅当每个
$K_i \otimes_k \Omega_{k/\mathbf{F}_p} \to
\Omega_{K_i/\mathbf{F}_p}$ 都是单射。因此，可以假设
$K/k$ 是有限生成域扩张。

\medskip\noindent
假设 $K/k$ 是有限生成可分域扩张。按照引理
\ref{algebra-lemma-generating-finitely-generated-separable-field-extensions}
选取 $x_1, \ldots, x_{r + 1} \in K$。
在此情形下，存在不可约多项式
$G(X_1, \ldots, X_{r + 1}) \in k[X_1, \ldots, X_{r + 1}]$，
使得 $G(x_1, \ldots, x_{r + 1}) = 0$，且
$\partial G/\partial X_{r + 1}$ 不恒为零。
此外，$K$ 是该整环的分式域。
% P01-E0982: the frozen English splits an incomplete domain sentence from a presentation with the wrong base field K instead of k.
$S = K[X_1, \ldots, X_{r + 1}]/(G)$。
写成
$$
G = \sum a_I X^I, \quad X^I = X_1^{i_1}\ldots X_{r + 1}^{i_{r + 1}}.
$$
利用上面对 $S$ 的表现，可知
$$
\Omega_{S/\mathbf{F}_p}
=
\frac{
S \otimes_k \Omega_k \oplus
\bigoplus\nolimits_{i = 1, \ldots, r + 1} S\text{d}X_i
}{
\langle
\sum X^I \text{d}a_I + \sum \partial G/\partial X_i \text{d}X_i
\rangle
}
$$
% P01-E0983: the two frozen presentations use the undefined abbreviation Omega_k instead of Omega_{k/F_p}.
由于 $\Omega_{K/\mathbf{F}_p}$ 是 $S$-模
$\Omega_{S/\mathbf{F}_p}$ 的局部化（见引理
\ref{algebra-lemma-differentials-localize}），可得
$$
\Omega_{K/\mathbf{F}_p}
=
\frac{
K \otimes_k \Omega_k \oplus
\bigoplus\nolimits_{i = 1, \ldots, r + 1} K\text{d}X_i
}{
\langle
\sum X^I \text{d}a_I + \sum \partial G/\partial X_i \text{d}X_i
\rangle
}
$$
现在，由于多项式 $\partial G/\partial X_{r + 1}$ 不恒为零，
可知映射
$K \otimes_k \Omega_{k/\mathbf{F}_p} \to \Omega_{S/\mathbf{F}_p}$
是单射，如所需。
% P01-E0984: the frozen conclusion names Omega_{S/F_p}, although the desired and just-presented codomain is Omega_{K/F_p}.

\medskip\noindent
假设 $K/k$ 是有限生成域扩张，并且
$K \otimes_k \Omega_{k/\mathbf{F}_p} \to \Omega_{K/\mathbf{F}_p}$
是单射。（证明的这一部分与证明引理
\ref{algebra-lemma-characterize-separable-field-extensions} 的论证相同。）
设 $x_1, \ldots, x_r$ 是 $K$ 在 $k$ 上的超越基，并且有限扩张
$k(x_1, \ldots, x_r) \subset K$ 的不可分次数最小。
如果 $K$ 在 $k(x_1, \ldots, x_r)$ 上可分，则已经得到结论。
假设不是如此，以导出矛盾。于是存在元素
$\alpha \in K$，它在 $k(x_1, \ldots, x_r)$ 上不可分。
设 $P(T) \in k(x_1, \ldots, x_r)[T]$ 是它的最小多项式。
因为 $\alpha$ 不可分，实际上 $P$ 是 $T^p$ 的多项式。
清除分母，得到不可约多项式
$$
G(X_1, \ldots, X_r, T) = \sum a_{I, i} X^I T^i \in k[X_1, \ldots, X_r, T]
$$
使得 $G(x_1, \ldots, x_r, \alpha) = 0$ 在 $K$ 中成立。
注意，这意味着
$k[X_1, \ldots, X_r, T]/(G) \subset K$。
可以假设对某个二元组 $(I_0, i_0)$ 有
$a_{I_0, i_0} = 1$。
我们断言，至少对一个 $i$，$\text{d}G/\text{d}X_i$
不恒为零。事实上，如果不是这样，则 $G$ 实际上是
$X_1^p, \ldots, X_r^p, T^p$ 的多项式。于是
$$
\sum\nolimits_{(I, i) \not = (I_0, i_0)} x^I\alpha^i \text{d}a_{I, i}
$$
在 $\Omega_{K/\mathbf{F}_p}$ 中为零。注意，$K$ 的元素集合
$$
\{x^I\alpha^i \mid a_{I, i} \not = 0 \text{ 且 } (I, i) \not = (I_0, i_0)\}
$$
中不存在 $k$-线性关系。因此，映射
$K \otimes_k \Omega_{k/\mathbf{F}_p} \to \Omega_{K/\mathbf{F}_p}$
是单射这一假设蕴含：对所有 $(I, i)$，
$\text{d}a_{I, i} = 0$ 在 $\Omega_{k/\mathbf{F}_p}$ 中成立。
由引理 \ref{algebra-lemma-derivative-zero-pth-power} 可知，每个
$a_{I, i}$ 都是 $p$ 次幂，这意味着 $G$ 是 $p$ 次幂，
与 $G$ 的不可约性矛盾。因此，重新编号后可以假设
$\text{d}G/\text{d}X_1$ 非零。
于是可知 $x_1$ 在 $k(x_2, \ldots, x_r, \alpha)$ 上可分代数，
并且 $x_2, \ldots, x_r, \alpha$ 是 $K$ 在 $k$ 上的超越基。
这意味着有限扩张
$k(x_2, \ldots, x_r, \alpha) \subset K$ 的不可分次数小于有限扩张
$k(x_1, \ldots, x_r) \subset K$ 的不可分次数，产生矛盾。
\end{proof}

\begin{lemma}
\label{algebra-lemma-formally-smooth-implies-separable}
设 $K/k$ 是域扩张。如果 $K$ 在 $k$ 上形式光滑，
则 $K$ 是 $k$ 的可分扩张。
\end{lemma}

\begin{proof}
假设 $K$ 在 $k$ 上形式光滑。由引理
\ref{algebra-lemma-ses-formally-smooth} 可知
$K \otimes_k \Omega_{k/\mathbf{Z}} \to \Omega_{K/\mathbf{Z}}$
是单射。于是由引理 \ref{algebra-lemma-separable-differentials}，
$K$ 在 $k$ 上可分。
% P01-E0985: the frozen proof invokes a characteristic-p lemma without separately handling characteristic zero.
\end{proof}

\begin{lemma}
\label{algebra-lemma-characterize-formally-smooth-field-extension}
设 $K/k$ 是域扩张。则 $K$ 在 $k$ 上形式光滑，当且仅当
$H_1(L_{K/k}) = 0$。
\end{lemma}

\begin{proof}
这由命题 \ref{algebra-proposition-characterize-formally-smooth} 以及
向量空间是自由模（因而是投射模）这一事实得出。
% P01-E0986: frozen English has the number disagreement “a vector spaces is free”.
\end{proof}

\begin{lemma}
\label{algebra-lemma-formally-smooth-extensions-easy}
设 $K/k$ 是域扩张。
\begin{enumerate}
\item 如果 $K$ 在 $k$ 上纯超越，则 $K$ 在 $k$ 上形式光滑。
\item 如果 $K$ 在 $k$ 上可分代数，则 $K$ 在 $k$ 上形式光滑。
\item 如果 $K$ 在 $k$ 上可分，则 $K$ 在 $k$ 上形式光滑。
\end{enumerate}
\end{lemma}

\begin{proof}
对 (1)，写 $K = k(x_j; j \in J)$。假设 $A$ 是 $k$-代数，
$I \subset A$ 是平方为零的理想。设
$\varphi : K \to A/I$ 是 $k$-代数映射。
设 $a_j \in A$ 满足 $a_j \mod I = \varphi(x_j)$。
那么容易看出，存在唯一的 $k$-代数映射 $K \to A$，
它把 $x_j$ 映到 $a_j$，且模 $I$ 约化为 $\varphi$。
因此，$k \subset K$ 形式光滑。

\medskip\noindent
在情形 (2) 中，$k \subset K$ 是平展环扩张的余极限。
平展环映射是形式平展的（引理
\ref{algebra-lemma-formally-etale-etale}）。因此，该情形由引理
\ref{algebra-lemma-colimit-formally-etale} 以及形式平展环映射形式光滑这一
显然事实得出。

\medskip\noindent
在情形 (3) 中，把 $K = \colim K_i$ 写成其有限生成
$k$-子扩张的滤过余极限。
% P01-E0987: frozen English says “finitely generated sub k-extensions” instead of “finitely generated k-subextensions”.
由定义 \ref{algebra-definition-separable-field-extension}，每个 $K_i$
都是 $k$ 的某个纯超越扩张上的可分代数扩张。
因此，由情形 (1)、(2) 以及引理
\ref{algebra-lemma-compose-formally-smooth}，$K_i/k$ 形式光滑。
于是由引理
\ref{algebra-lemma-characterize-formally-smooth-field-extension}，
$H_1(L_{K_i/k}) = 0$。由引理 \ref{algebra-lemma-colimits-NL}，
$H_1(L_{K/k}) = 0$。再次应用引理
\ref{algebra-lemma-characterize-formally-smooth-field-extension}，可知
$K/k$ 形式光滑。
\end{proof}

\begin{lemma}
\label{algebra-lemma-fields-are-formally-smooth}
\begin{slogan}
对域扩张，形式光滑等价于可分。
\end{slogan}
设 $k$ 是域。
\begin{enumerate}
\item 如果 $k$ 的特征为零，则 $k$ 的任意扩域都在 $k$
上形式光滑。
\item 如果 $k$ 的特征为 $p > 0$，则 $K/k$ 形式光滑，
当且仅当它是可分域扩张。
\end{enumerate}
\end{lemma}

\begin{proof}
合并引理 \ref{algebra-lemma-formally-smooth-implies-separable} 和
\ref{algebra-lemma-formally-smooth-extensions-easy}。
\end{proof}

\noindent
下面汇集可分域扩张的所有不同刻画。

\begin{proposition}
\label{algebra-proposition-characterize-separable-field-extensions}
设 $K/k$ 是域扩张。如果 $k$ 的特征为零，则
% P01-E0988: the frozen proposition omits the phrase “the following are equivalent” before this list.
\begin{enumerate}
\item $K$ 在 $k$ 上可分；
\item $K$ 在 $k$ 上几何约化；
\item $K$ 在 $k$ 上形式光滑；
\item $H_1(L_{K/k}) = 0$；
\item 映射 $K \otimes_k \Omega_{k/\mathbf{Z}} \to \Omega_{K/\mathbf{Z}}$
是单射。
\end{enumerate}
如果 $k$ 的特征为 $p > 0$，则下列条件等价：
\begin{enumerate}
\item $K$ 在 $k$ 上可分；
\item 环 $K \otimes_k k^{1/p}$ 约化；
\item $K$ 在 $k$ 上几何约化；
\item 映射 $K \otimes_k \Omega_{k/\mathbf{F}_p} \to \Omega_{K/\mathbf{F}_p}$
是单射；
\item $H_1(L_{K/k}) = 0$；
\item $K$ 在 $k$ 上形式光滑。
\end{enumerate}
\end{proposition}

\begin{proof}
这是下列引理的合并：
\ref{algebra-lemma-characterize-separable-field-extensions}、
\ref{algebra-lemma-fields-are-formally-smooth}
\ref{algebra-lemma-formally-smooth-implies-separable} 以及
\ref{algebra-lemma-separable-differentials}。
% P01-E0989: frozen English omits punctuation or a conjunction between the middle two lemma references.
\end{proof}

\noindent
下面给出有限生成可分域扩张的又一个刻画。

\begin{lemma}
\label{algebra-lemma-localization-smooth-separable}
设 $K/k$ 是有限生成域扩张。则 $K$ 在 $k$ 上可分，
当且仅当 $K$ 是某个光滑 $k$-代数的局部化。
\end{lemma}

\begin{proof}
选取有限型 $k$-代数 $R$，使其为分式域等于 $K$ 的整环。
引理 \ref{algebra-lemma-smooth-at-generic-point} 说明，
$k \to R$ 在 $(0)$ 处光滑，当且仅当 $K/k$ 可分。
这证明了该引理。
\end{proof}

\begin{lemma}
\label{algebra-lemma-colimit-syntomic}
设 $K/k$ 是域扩张。则 $K$ 是 $k$ 上全局完全交代数的
滤过余极限。如果 $K/k$ 可分，则 $K$ 是 $k$ 上光滑代数的
滤过余极限。
\end{lemma}

\begin{proof}
假设 $E \subset K$ 是有限子集。只需证明，存在包含 $E$
的 $k$-子代数 $A \subset K$，并且它是 $k$ 上的全局完全交
（相应地，光滑）。可分/光滑情形由引理
\ref{algebra-lemma-localization-smooth-separable} 得出。
一般地，设 $L \subset K$ 是由 $E$ 生成的子域。
选取它在 $k$ 上的超越基
$x_1, \ldots, x_d \in L$。扩张
$L/k(x_1, \ldots, x_d)$ 是有限的。设
$L = k(x_1, \ldots, x_d)[y_1, \ldots, y_r]$。
归纳选取多项式
$P_i \in k(x_1, \ldots, x_d)[Y_1, \ldots, Y_r]$，
使得 $P_i = P_i(Y_1, \ldots, Y_i)$ 是
$k(x_1, \ldots, x_d)[Y_1, \ldots, Y_{i - 1}]$
上关于 $Y_i$ 的首一多项式，并且映到 $y_i$ 在
$k(x_1, \ldots, x_d)[y_1, \ldots, y_{i - 1}][Y_i]$
中的最小多项式。
% P01-E0990: frozen English says “minimum polynomial” instead of “minimal polynomial”.
于是显然，$P_1, \ldots, P_r$ 是
$k(x_1, \ldots, x_r)[Y_1, \ldots, Y_r]$ 中的正则序列，且
$L = k(x_1, \ldots, x_r)[Y_1, \ldots, Y_r]/(P_1, \ldots, P_r)$。
% P01-E0991: both frozen formulas use x_1,...,x_r where the established transcendence basis is x_1,...,x_d.
如果 $h \in k[x_1, \ldots, x_d]$ 是使得
$P_i \in k[x_1, \ldots, x_d, 1/h, Y_1, \ldots, Y_r]$
的多项式，则 $P_1, \ldots, P_r$ 是
$k[x_1, \ldots, x_d, 1/h, Y_1, \ldots, Y_r]$ 中的正则序列，且
$A = k[x_1, \ldots, x_d, 1/h, Y_1, \ldots, Y_r]/(P_1, \ldots, P_r)$
是全局完全交。适当调整 $h$ 的选取后，可以假设
$E \subset A$，从而得到结论。
\end{proof}

\section{构造平坦环映射}
\label{algebra-section-constructing-flat}

\noindent
下面的引理有时很有用。

\begin{lemma}
\label{algebra-lemma-flat-local-given-residue-field}
设 $(R, \mathfrak m, k)$ 是局部环，$K/k$ 是域扩张。
则存在局部环 $(R', \mathfrak m', k')$ 以及平坦局部环映射
$R \to R'$，使得 $\mathfrak m' = \mathfrak mR'$，并且
$k'$ 作为 $k$ 的扩张同构于 $K$。
\end{lemma}

\begin{proof}
假设 $k' = k(\alpha)$ 是 $k$ 的单生成扩张。
则 $k'$ 是引理所述某平坦局部扩张 $R \subset R'$ 的剩余域。
具体而言，如果 $\alpha$ 在 $k$ 上超越，则令 $R'$ 为
$R[x]$ 在素理想 $\mathfrak mR[x]$ 处的局部化。
如果 $\alpha$ 是代数的，其最小多项式为
$T^d + \sum \overline{\lambda}_iT^{d - i}$，则令
$R' = R[T]/(T^d + \sum \lambda_i T^{d - i})$。

\medskip\noindent
考虑所有三元组 $(k', R \to R', \phi)$ 所成的集合，其中
$k \subset k' \subset K$ 是子域，$R \to R'$ 是引理所述的
局部环映射，并且
$\phi : R' \to k'$ 诱导 $k$-扩张的同构
$R'/\mathfrak mR' \cong k'$。这些对象构成一个``大''范畴
$\mathcal{C}$；态射
$(k_1, R_1, \phi_1) \to (k_2, R_2, \phi_2)$ 由环映射
$\psi : R_1 \to R_2$ 给出，并要求交换图表
$$
\xymatrix{
R_1 \ar[d]_\psi \ar[r]_{\phi_1} & k_1 \ar[r] & K \ar@{=}[d] \\
R_2 \ar[r]^{\phi_2} & k_2 \ar[r] & K
}
$$
交换。这蕴含 $k_1 \subset k_2$。

\medskip\noindent
假设 $I$ 是有向集，而
$((R_i, k_i, \phi_i), \psi_{ii'})$ 是 $I$ 上的系统，见范畴，
第 \ref{categories-section-posets-limits} 节。
% P01-E0992: frozen object order changes from (k', R -> R', phi) to (R_i, k_i, phi_i), and the stated upper bound repeats the reversed order.
在此情形下，可以考虑
$$
R' = \colim_{i \in I} R_i
$$
这是极大理想为 $\mathfrak mR'$ 的局部环，剩余域为
$k' = \bigcup_{i \in I} k_i$。此外，环映射 $R \to R'$
是平坦的，因为它是平坦映射的余极限（且张量积与有向余极限交换）。
因此，$(R', k', \phi')$ 是该系统的一个``上界''。

\medskip\noindent
如果 $\mathcal{C}$ 是集合，则几乎平凡地应用 Zorn 引理即可完成证明，
但它并不是集合。
% P01-E0993: frozen English uses “if C was a set” rather than the irrealis “if C were a set”.
（实际上，可以通过找出所出现局部环之基数的合理界来使这一论证成立。）
为绕开这个问题，我们在 $K$ 上选取一个良序。
对 $x \in K$，令 $K(x)$ 为 $K$ 的如下子域：由 $K$ 中所有
$\leq x$ 的元素所生成。对 $x \in K$ 作超限递归，我们将构造引理所述的
环映射 $R \subset R(x)$，其剩余域扩张为 $K(x)/k$。
此外，按构造，对所有 $y \leq x$，$R(x)$ 都将包含 $R(y)$。
具体而言，如果 $x$ 有前驱 $x'$，则 $K(x) = K(x')[x]$，
因而可以令 $R(x') \subset R(x)$ 为证明第一段所构造的局部环扩张。
如果 $x$ 没有前驱，则先按照证明第三段令
$R'(x) = \colim_{x' < x} R(x')$。
$R'(x)$ 的剩余域为
$K'(x) = \bigcup_{x' < x} K(x')$。由于
$K(x) = K'(x)[x]$，可用证明第一段的构造得到
$R'(x) \subset R(x)$。
% P01-E0994: both frozen field-generation formulas use square brackets, which fail when x is transcendental.
这就完成了引理的证明。
\end{proof}

\begin{lemma}
\label{algebra-lemma-colimit-finite-etale-given-residue-field}
设 $(R, \mathfrak m, k)$ 是局部环。如果 $k \subset K$
是可分代数扩张，则存在有向集 $I$ 以及局部环的有限平展扩张系统
$R \subset R_i$，$i \in I$，使得 $R' = \colim R_i$
的剩余域为 $K$（作为 $k$ 的扩张）。
\end{lemma}

\begin{proof}
设 $R \subset R'$ 是引理
\ref{algebra-lemma-flat-local-given-residue-field} 的证明中构造的扩张。
按构造，$R' = \colim_{\alpha \in A} R_\alpha$，其中 $A$
是良序集，过渡映射 $R_\alpha \to R_{\alpha + 1}$ 是有限平展的，
并且当 $\alpha$ 不是后继元时，
$R_\alpha = \colim_{\beta < \alpha} R_\beta$。
我们将用超限归纳证明该结果。

\medskip\noindent
假设结论对 $R_\alpha$ 成立，即
$R_\alpha = \colim R_i$，其中 $R_i$ 在 $R$ 上有限平展。
由于 $R_\alpha \to R_{\alpha + 1}$ 有限平展，存在一个 $i$
以及有限平展扩张 $R_i \to R_{i, 1}$，使得
$R_{\alpha + 1} = R_\alpha \otimes_{R_i} R_{i, 1}$。
因此，
$R_{\alpha + 1} = \colim_{i' \geq i} R_{i'} \otimes_{R_i} R_{i, 1}$，
结论对 $\alpha + 1$ 成立。
假设 $\alpha$ 不是后继元，并且结论对所有
$\beta < \alpha$ 的 $R_\beta$ 成立。
由于每个有限子集 $E \subset R_\alpha$ 都包含于某个
$\beta < \alpha$ 对应的 $R_\beta$ 中，由假设可知 $E$
包含于某个有限平展子扩张中。因此，结论对 $R_\alpha$ 成立。
% P01-E0995: frozen English inserts an ungrammatical “and” between the finite-subset premise and its conclusion.
\end{proof}

\begin{lemma}
\label{algebra-lemma-finite-free-given-residue-field-extension}
设 $R$ 是环，$\mathfrak p \subset R$ 是素理想，并设
$L/\kappa(\mathfrak p)$ 是有限域扩张。
则存在有限自由环映射 $R \to S$，使得
$\mathfrak q = \mathfrak pS$ 是素理想，并且
$\kappa(\mathfrak q)/\kappa(\mathfrak p)$ 同构于给定扩张
$L/\kappa(\mathfrak p)$。
\end{lemma}

\begin{proof}
对 $\kappa(\mathfrak p) \subset L$ 的次数作归纳。
% P01-E0996: frozen English says “induction of the degree” instead of “induction on the degree”.
如果次数为 $1$，则取 $R = S$。
一般地，如果存在子扩张
$\kappa(\mathfrak p) \subset L' \subset L$，则对次数归纳即可得到
结论（先构造与 $L'/\kappa(\mathfrak p)$ 对应的
$R \subset S'$，再构造与 $L/L'$ 对应的 $S' \subset S$）。
% P01-E0997: the frozen conditional is vacuous without requiring a proper intermediate field, and “then construction” is also malformed.
因此，可以假设 $L \supset \kappa(\mathfrak p)$ 由单个元素
$\alpha \in L$ 生成。设
$X^d + \sum_{i < d} a_iX^i$ 是 $\alpha$ 在
$\kappa(\mathfrak p)$ 上的最小多项式，所以
$a_i \in \kappa(\mathfrak p)$。
可以把 $a_i$ 写成某个 $f_i, g \in R$ 的 $f_i/g$ 的像，
其中 $g \not \in \mathfrak p$。
用 $g\alpha$ 替换 $\alpha$（并相应地用
$g^{d - i}a_i$ 替换 $a_i$）后，可以假设 $a_i$ 是某个
$f_i \in R$ 的像。此时只需取
$S = R[x]/(x^d + \sum f_ix^i)$。
\end{proof}

\begin{lemma}
\label{algebra-lemma-cofinal-system-flat}
设 $A$ 是环，令 $\kappa = \max(|A|, \aleph_0)$。
则每个平坦 $A$-代数 $B$ 都是其满足
$B' \subset B$、$|B'| \leq \kappa$ 的平坦 $A$-子代数的
滤过余极限。（注意，如果 $B$ 在 $A$ 上忠实平坦，
则 $B'$ 也在 $A$ 上忠实平坦。）
\end{lemma}

\begin{proof}
如果 $B$ 的基数 $\leq \kappa$，则结论成立。
设 $E \subset B$ 是满足 $|E| \leq \kappa$ 的 $A$-子代数。
我们将证明 $E$ 包含于某个平坦 $A$-子代数 $B'$ 中，且
$|B'| \leq \kappa$。于是引理成立，因为
(a) $B$ 的每个有限子集都包含于某个基数至多为 $\kappa$ 的
$A$-子代数中，并且 (b) $B$ 的任意两个基数至多为
$\kappa$ 的 $A$-子代数都包含于某个基数至多为 $\kappa$
的 $A$-子代数中。略去细节。

\medskip\noindent
我们将归纳构造一列 $A$-子代数
$$
E = E_0 \subset E_1 \subset E_2 \subset \ldots
$$
其中每个子代数的基数都 $\leq \kappa$；最后证明
$B' = \bigcup E_k$ 在 $A$ 上平坦。

\medskip\noindent
构造如下。令 $E_0 = E$。给定 $k \geq 0$ 时的 $E_k$，
考虑 $E_k$ 中元素之间以 $A$ 中元素为系数的关系所成的集合
$S_k$。因此，元素 $s \in S_k$ 由整数 $n \geq 1$、
$a_1, \ldots, a_n \in A$ 以及
$e_1, \ldots, e_n \in E_k$ 给出，并满足
$\sum a_i e_i = 0$ 在 $E_k$ 中成立。
$A \to B$ 的平坦性以及引理 \ref{algebra-lemma-flat-eq} 表明，
对每个
$s = (n, a_1, \ldots, a_n, e_1, \ldots, e_n) \in S_k$，
可以选取
$$
(m_s, b_{s, 1}, \ldots, b_{s, m_s}, a_{s, 11}, \ldots, a_{s, nm_s})
$$
其中 $m_s \geq 0$ 是整数，$b_{s, j} \in B$，
$a_{s, ij} \in A$，并且
$$
e_i = \sum\nolimits_j a_{s, ij} b_{s, j}, \forall i,
\quad\text{且}\quad
0 = \sum\nolimits_i a_i a_{s, ij}, \forall j.
$$
给定这些选取后，令 $E_{k + 1} \subset B$ 为下列元素生成的
$A$-子代数：
% P01-E0998: frozen English misspells “choices” as “choicse”.
\begin{enumerate}
\item $E_k$；
\item 对每个 $s \in S_k$，元素
$b_{s, 1}, \ldots, b_{s, m_s}$。
\end{enumerate}
一些集合论（略去）表明，$E_{k + 1}$ 的基数至多为
$\kappa$（这里用到归纳已知 $|E_k| \leq \kappa$，因而
$S_k$ 的基数也至多为 $\kappa$）。
% P01-E0999: frozen English says “has at most cardinality kappa” instead of “has cardinality at most kappa”.

\medskip\noindent
为了证明 $B' = \bigcup E_k$ 在 $A$ 上平坦，考虑 $B'$ 中
系数属于 $A$ 的关系
$\sum_{i = 1, \ldots, n} a_i b'_i = 0$。
选取充分大的 $k$，使得对 $i = 1, \ldots, n$ 有
$b'_i \in E_k$。则
$(n, a_1, \ldots, a_n, b'_1, \ldots, b'_n) \in S_k$，
所以该关系在 $E_{k + 1}$ 中是平凡的，从而在 $B'$ 中更是如此。
因此，由引理 \ref{algebra-lemma-flat-eq}，$A \to B'$ 是平坦的。
\end{proof}

\section{Cohen 结构定理}
\label{algebra-section-cohen-structure-theorem}

\noindent
下面是交换代数中的一个基本概念。

\begin{definition}
\label{algebra-definition-complete-local-ring}
设 $(R, \mathfrak m)$ 是局部环。如果到 $R$ 关于
$\mathfrak m$ 的完备化的典范映射
$$
R \longrightarrow \lim_n R/\mathfrak m^n
$$
是同构，则称 $R$ 是{\it 完备局部环}\footnote{这包括条件
$\bigcap \mathfrak m^n = (0)$；在一些文献中，这可能表述为
$R$ 完备且分离。警告：局部环的完备化
$\lim_n R/\mathfrak m^n$ 可能不是完备的，见例，引理
\ref{examples-lemma-noncomplete-completion}。
当 $\mathfrak m$ 有限生成时不会发生这种情况，见引理
\ref{algebra-lemma-hathat-finitely-generated}；在这种情况下，完备化是
Noether 环，见引理 \ref{algebra-lemma-completion-Noetherian}。}。
\end{definition}

\noindent
注意，Artin 局部环 $R$ 是完备局部环，因为对某个
$n > 0$ 有 $\mathfrak m_R^n = 0$。本节主要讨论 Noether
完备局部环。

\begin{lemma}
\label{algebra-lemma-quotient-complete-local}
设 $R$ 是 Noether 完备局部环。$R$ 的任意商也是 Noether
完备局部环。给定有限环映射 $R \to S$，则 $S$ 是 Noether
完备局部环的乘积。
% P01-E1000: frozen English uses the doubled construction “Given ..., then ...”.
\end{lemma}

\begin{proof}
由引理 \ref{algebra-lemma-Noetherian-permanence}，环 $S$ 是 Noether 的。
由引理 \ref{algebra-lemma-completion-tensor}，$S$ 作为 $R$-模是完备的。
因此，由引理 \ref{algebra-lemma-completion-finite-extension}，$S$
是它在各极大理想处的完备化的乘积。
\end{proof}

\begin{lemma}
\label{algebra-lemma-complete-local-ring-Noetherian}
设 $(R, \mathfrak m)$ 是完备局部环。如果 $\mathfrak m$
是有限生成理想，则 $R$ 是 Noether 环。
\end{lemma}

\begin{proof}
见引理 \ref{algebra-lemma-completion-Noetherian}。
\end{proof}

\begin{definition}
\label{algebra-definition-coefficient-ring}
设 $(R, \mathfrak m)$ 是完备局部环。如果子环
$\Lambda \subset R$ 满足下列条件，则称它为{\it 系数环}：
\begin{enumerate}
\item $\Lambda$ 是完备局部环，其极大理想为
$\Lambda \cap \mathfrak m$；
\item $\Lambda$ 的剩余域同构地映到 $R$ 的剩余域；
\item $\Lambda \cap \mathfrak m = p\Lambda$，其中 $p$ 是
$R$ 的剩余域的特征。
\end{enumerate}
\end{definition}

\noindent
下面对该定义作几点说明。分成下列情形讨论：
\begin{enumerate}
\item 局部环 $R$ 包含一个域。如果
$\mathbf{Q} \subset R$，或者 $pR = 0$，其中 $p$ 是
$R/\mathfrak m$ 的特征，就会出现这种情况。此时，系数环
$\Lambda$ 是包含于 $R$ 中的域，并同构地映到
$R/\mathfrak m$。
\item $R/\mathfrak m$ 的特征为 $p > 0$，但 $p$ 的任何幂
在 $R$ 中都不为零。此时，$\Lambda$ 是以 $p$ 为一致化元、
剩余域为 $R/\mathfrak m$ 的完备离散赋值环。
\item $R/\mathfrak m$ 的特征为 $p > 0$，并且对某个
$n > 1$，在 $R$ 中有 $p^{n - 1} \not = 0$、$p^n = 0$。
此时，$\Lambda$ 是 Artin 局部环，其极大理想由 $p$ 生成，
且剩余域为 $R/\mathfrak m$。
\end{enumerate}
上述以 $p$ 为一致化元的完备离散赋值环具有特殊作用，
我们作如下命名。

\begin{definition}
\label{algebra-definition-cohen-ring}
{it Cohen 环}是以某个素数 $p$ 为一致化元的完备离散赋值环。
\end{definition}

\begin{lemma}
\label{algebra-lemma-cohen-rings-exist}
设 $p$ 是素数，$k$ 是特征为 $p$ 的域。
则存在 Cohen 环 $\Lambda$，使得 $\Lambda/p\Lambda \cong k$。
\end{lemma}

\begin{proof}
首先注意，$p$ 进整数 $\mathbf{Z}_p$ 构成
$\mathbf{F}_p$ 的一个 Cohen 环。设 $k$ 是特征 $p$ 的任意域。
取平坦局部环映射 $\mathbf{Z}_p \to R$，使得
$\mathfrak m_R = pR$ 且 $R/pR = k$，见引理
\ref{algebra-lemma-flat-local-given-residue-field}。
由引理 \ref{algebra-lemma-completion-Noetherian}，完备化
$\Lambda = R^\wedge$ 是 Noether 的。由于
$\Lambda/p\Lambda = R/pR$ 是域，可知 $\Lambda$ 是极大理想为
$(p)$ 的完备 Noether 局部环（该等式及完备性由引理
\ref{algebra-lemma-hathat-finitely-generated} 得出）。
对每个 $n$，映射
$\mathbf{Z}/p^n\mathbf{Z} \to \Lambda/p^n\Lambda = R/p^nR$
都是平坦的（该等式再次由引理
\ref{algebra-lemma-hathat-finitely-generated} 得出）。
因此，例如由引理 \ref{algebra-lemma-flat-module-powers} 可知
$\mathbf{Z}_p \to \Lambda$ 平坦。所以 $p$ 是 $\Lambda$
中的非零因子。于是 $\Lambda$ 的维数 $\geq 1$
（引理 \ref{algebra-lemma-one-equation}），从而 $\Lambda$ 是维数为
$1$ 的正则环；即由引理 \ref{algebra-lemma-characterize-dvr}，
它是离散赋值环。因此，$\Lambda$ 是 $k$ 的 Cohen 环。
\end{proof}

\begin{lemma}
\label{algebra-lemma-cohen-ring-formally-smooth}
设 $p > 0$ 是素数。设 $\Lambda$ 是剩余域特征为 $p$
的 Cohen 环。对每个 $n \geq 1$，环映射
$$
\mathbf{Z}/p^n\mathbf{Z} \to \Lambda/p^n\Lambda
$$
都是形式光滑的。
\end{lemma}

\begin{proof}
当 $n = 1$ 时，这由命题
\ref{algebra-proposition-characterize-separable-field-extensions} 得出。
对一般的 $n$，对 $n$ 作归纳。具体而言，如果
$\mathbf{Z}/p^n\mathbf{Z} \to \Lambda/p^n\Lambda$ 形式光滑，
则可将引理 \ref{algebra-lemma-lift-formal-smoothness} 应用于环映射
$\mathbf{Z}/p^{n + 1}\mathbf{Z} \to \Lambda/p^{n + 1}\Lambda$
以及理想
$I = (p^n) \subset \mathbf{Z}/p^{n + 1}\mathbf{Z}$。
\end{proof}

\begin{theorem}[Cohen 结构定理]
\label{algebra-theorem-cohen-structure-theorem}
设 $(R, \mathfrak m)$ 是完备局部环。
\begin{enumerate}
\item $R$ 有一个系数环（见定义
\ref{algebra-definition-coefficient-ring}）；
\item 如果 $\mathfrak m$ 是有限生成理想，则 $R$ 同构于商
$$
\Lambda[[x_1, \ldots, x_n]]/I
$$
其中 $\Lambda$ 是域或 Cohen 环。
\end{enumerate}
\end{theorem}

\begin{proof}
先证明系数环存在。首先考虑剩余域 $\kappa$ 的特征为零的情形。
具体而言，在这种情况下，我们对 $n > 0$ 归纳证明存在映射
$$
\varphi_n : \kappa \longrightarrow R/\mathfrak m^n
$$
它是典范映射 $R/\mathfrak m^n \to \kappa = R/\mathfrak m$
的截面。当 $n = 1$ 时，这显然成立。如果 $n > 1$，
设已给定 $\varphi_{n - 1}$。
由命题 \ref{algebra-proposition-characterize-separable-field-extensions}，
域扩张 $\kappa/\mathbf{Q}$ 形式光滑。因此，可以在下图中找到
虚线箭头：
$$
\xymatrix{
R/\mathfrak m^{n - 1} &
R/\mathfrak m^n \ar[l] \\
\kappa \ar[u]^{\varphi_{n - 1}} \ar@{..>}[ru] & \mathbf{Q} \ar[l] \ar[u]
}
$$
这证明了归纳步骤。把这些映射组合起来，得到
$$
\lim_n\ \varphi_n : \kappa \longrightarrow
R = \lim_n\ R/\mathfrak m^n
$$
其像就是所需的系数环。

\medskip\noindent
其次，证明剩余域 $\kappa$ 的特征为 $p > 0$ 时系数环存在。
具体而言，选取满足 $\kappa = \Lambda/p\Lambda$ 的 Cohen 环
$\Lambda$，见引理 \ref{algebra-lemma-cohen-rings-exist}。
在这种情况下，我们对 $n > 0$ 归纳证明存在映射
$$
\varphi_n :
\Lambda/p^n\Lambda
\longrightarrow
R/\mathfrak m^n
$$
它与约化映射 $R/\mathfrak m^n \to \kappa$ 的合成给出
已知同构 $\Lambda/p\Lambda = \kappa$。
当 $n = 1$ 时，这显然成立。如果 $n > 1$，设已给定
$\varphi_{n - 1}$。由引理
\ref{algebra-lemma-cohen-ring-formally-smooth}，环映射
$\mathbf{Z}/p^n\mathbf{Z} \to \Lambda/p^n\Lambda$ 形式光滑。
因此，可以在下图中找到虚线箭头：
$$
\xymatrix{
R/\mathfrak m^{n - 1} &
R/\mathfrak m^n \ar[l] \\
\Lambda/p^n\Lambda \ar[u]^{\varphi_{n - 1}} \ar@{..>}[ru] &
\mathbf{Z}/p^n\mathbf{Z} \ar[l] \ar[u]
}
$$
这证明了归纳步骤。把这些映射组合起来，得到
$$
\lim_n\ \varphi_n :
\Lambda = \lim_n\ \Lambda/p^n\Lambda
\longrightarrow
R = \lim_n\ R/\mathfrak m^n
$$
其像就是所需的系数环。

\medskip\noindent
定理的最后一个断言容易得出。具体而言，如果
$y_1, \ldots, y_n$ 是理想 $\mathfrak m$ 的生成元，则可利用
刚构造的映射 $\Lambda \to R$ 得到映射
$$
\Lambda[[x_1, \ldots, x_n]] \longrightarrow R,
\quad x_i \longmapsto y_i.
$$
由于两边关于 $(x_1, \ldots, x_n)$ 都是完备的，并且按构造该映射
模 $(x_1, \ldots, x_n)$ 是满射，由引理
\ref{algebra-lemma-completion-generalities} 可知该映射是满射。
\end{proof}

\begin{remark}
\label{algebra-remark-Noetherian-complete-local-ring-universally-catenary}
如果 $k$ 是域，则幂级数环 $k[[X_1, \ldots, X_d]]$ 是维数为
$d$ 的 Noether 完备正则局部环。如果 $\Lambda$ 是 Cohen 环，
则 $\Lambda[[X_1, \ldots, X_d]]$ 是维数为 $d + 1$ 的
完备 Noether 正则局部环。因此，Cohen 结构定理蕴含：
任意 Noether 完备局部环都是正则局部环的商。
特别地，Noether 完备局部环是普遍链状的，见引理
\ref{algebra-lemma-CM-ring-catenary} 和引理
\ref{algebra-lemma-regular-ring-CM}。
\end{remark}

\begin{lemma}
\label{algebra-lemma-regular-complete-containing-coefficient-field}
设 $(R, \mathfrak m)$ 是 Noether 完备局部环。假设 $R$ 正则。
\begin{enumerate}
\item 如果 $R$ 包含 $\mathbf{F}_p$ 或 $\mathbf{Q}$，则 $R$
同构于其剩余域上的幂级数环。
\item 如果 $k$ 是域，并且环映射 $k \to R$ 诱导同构
$k \to R/\mathfrak m$，则 $R$ 作为 $k$-代数同构于 $k$
上的幂级数环。
\end{enumerate}
\end{lemma}

\begin{proof}
在情形 (1) 中，由 Cohen 结构定理（定理
\ref{algebra-theorem-cohen-structure-theorem}），存在一个系数环；
它必须是同构地映到剩余域的域。因此，只需证明 (2)。
在情形 (2) 中，选取
$f_1, \ldots, f_d \in \mathfrak m$，使它们映到
$\mathfrak m/\mathfrak m^2$ 的一组基，并考虑把 $x_i$ 映到
$f_i$ 的连续 $k$-代数映射
$k[[x_1, \ldots, x_d]] \to R$。
由于源和靶都关于 $(x_1, \ldots, x_d)$ 完备，由引理
\ref{algebra-lemma-completion-generalities} 可知该映射是满射。
另一方面，它必为单射，因为否则由引理
\ref{algebra-lemma-one-equation}，$R$ 的维数将 $< d$。
\end{proof}

\begin{lemma}
\label{algebra-lemma-complete-local-Noetherian-domain-finite-over-regular}
设 $(R, \mathfrak m)$ 是 Noether 完备局部整环。
则存在 $R_0 \subset R$，具有下列性质：
% P01-E1001: frozen English says “there exists a R_0 subset R”.
\begin{enumerate}
\item $R_0$ 是正则完备局部环；
\item $R_0 \subset R$ 是有限的，并诱导剩余域上的同构；
\item $R_0$ 同构于 $k[[X_1, \ldots, X_d]]$ 或
$\Lambda[[X_1, \ldots, X_d]]$，其中 $k$ 是域，
$\Lambda$ 是 Cohen 环。
\end{enumerate}
\end{lemma}

\begin{proof}
设 $\Lambda$ 是 $R$ 的系数环。由于 $R$ 是整环，
$\Lambda$ 是域，或者 $\Lambda$ 是 Cohen 环。

\medskip\noindent
情形 I：$\Lambda = k$ 是域。令 $d = \dim(R)$。
选取生成定义理想 $I \subset R$ 的元素
$x_1, \ldots, x_d \in \mathfrak m$。
（见第 \ref{algebra-section-dimension} 节。）
由引理 \ref{algebra-lemma-change-ideal-completion} 可知，$R$
关于 $I$ 也是完备的。考虑把 $X_i$ 映到 $x_i$ 的映射
$R_0 = k[[X_1, \ldots, X_d]] \to R$。
注意，$R_0$ 关于理想 $I_0 = (X_1, \ldots, X_d)$ 完备，
并且 $R/I_0R \cong R/IR$ 在 $k = R_0/I_0$ 上有限
（因为 $\dim(R/I) = 0$，见第 \ref{algebra-section-dimension} 节）。
因此，由引理 \ref{algebra-lemma-finite-over-complete-ring} 可知
$R_0 \to R$ 是有限的。由于 $\dim(R) = \dim(R_0)$，
这蕴含 $R_0 \to R$ 是单射（见引理
\ref{algebra-lemma-integral-dim-up}），引理得证。

\medskip\noindent
情形 II：$\Lambda$ 是 Cohen 环。令 $d + 1 = \dim(R)$。
设 $p > 0$ 是剩余域 $k$ 的特征。由于 $R$ 是整环，
$p$ 是 $R$ 中的非零因子。因此，
$\dim(R/pR) = d$，见引理 \ref{algebra-lemma-one-equation}。
选取 $x_1, \ldots, x_d \in R$，使它们在 $R/pR$
中生成一个定义理想。则
$I = (p, x_1, \ldots, x_d)$ 是 $R$ 的定义理想。
由引理 \ref{algebra-lemma-change-ideal-completion} 可知，$R$
关于 $I$ 也是完备的。考虑把 $X_i$ 映到 $x_i$ 的映射
$R_0 = \Lambda[[X_1, \ldots, X_d]] \to R$。
注意，$R_0$ 关于理想
$I_0 = (p, X_1, \ldots, X_d)$ 完备，并且
$R/I_0R \cong R/IR$ 在 $k = R_0/I_0$ 上有限
（因为 $\dim(R/I) = 0$，见第 \ref{algebra-section-dimension} 节）。
因此，由引理 \ref{algebra-lemma-finite-over-complete-ring} 可知
$R_0 \to R$ 是有限的。由于 $\dim(R) = \dim(R_0)$，
这蕴含 $R_0 \to R$ 是单射（见引理
\ref{algebra-lemma-integral-dim-up}），引理得证。
\end{proof}

\section{日本型环}
\label{algebra-section-japanese}

\noindent
本节开始讨论整闭包的有限性。

\begin{definition}
\label{algebra-definition-N}
\begin{reference}
\cite[Chapter 0, Definition 23.1.1]{EGA}
\end{reference}
设 $R$ 是分式域为 $K$ 的整环。
\begin{enumerate}
\item 若 $R$ 在 $K$ 中的整闭包是有限 $R$-模，则称 $R$ 为
{\it N-1}；
\item 若对任意有限域扩张 $L/K$，$R$ 在 $L$ 中的整闭包
在 $R$ 上有限，则称 $R$ 为 {\it N-2} 或{\it 日本型的}。
\end{enumerate}
\end{definition}

\noindent
这些概念主要用于 Noether 环，不过下面给出一个非 Noether 的例子。

\begin{example}
\label{algebra-example-Japanese-not-Noetherian}
设 $k$ 是域。整环 $R = k[x_1, x_2, x_3, \ldots]$ 是 N-2 的，
但不是 Noether 的。理由如下。假设 $R \subset L$，并且域 $L$
是 $R$ 的分式域的有限扩张。则存在整数 $n$，使得 $L$ 可由有限扩张
$L_0/k(x_1, \ldots, x_n)$ 添加（超越）元素
$x_{n + 1}, x_{n + 2}$ 等而得到。令 $S_0$ 为
$k[x_1, \ldots, x_n]$ 在 $L_0$ 中的整闭包。由下面的命题
\ref{algebra-proposition-ubiquity-nagata}，$S_0$ 在
$k[x_1, \ldots, x_n]$ 上有限。此外，$R$ 在 $L$ 中的整闭包是
$S = S_0[x_{n + 1}, x_{n + 2}, \ldots]$（使用引理
\ref{algebra-lemma-polynomial-domain-normal}），因而在 $R$ 上有限。
同一论证也适用于 $R = \mathbf{Z}[x_1, x_2, x_3, \ldots]$。
\end{example}

\begin{lemma}
\label{algebra-lemma-localize-N}
设 $R$ 是整环。若 $R$ 是 N-1 的，则 $R$ 的任何局部化也是 N-1 的。
N-2 的情形亦然。
\end{lemma}

\begin{proof}
这是因为取整闭包与局部化交换，见引理
\ref{algebra-lemma-integral-closure-localize}。
\end{proof}

\begin{lemma}
\label{algebra-lemma-Japanese-local}
设 $R$ 是整环。设 $f_1, \ldots, f_n \in R$ 生成单位理想。
若每个整环 $R_{f_i}$ 都是 N-1 的，则 $R$ 也是 N-1 的。
N-2 的情形亦然。
\end{lemma}

\begin{proof}
假设 $R_{f_i}$ 是 N-2 的（或 N-1 的）。设 $L$ 是 $R$ 的分式域的
有限扩张（在 N-1 情形中就取该分式域本身）。令 $S$ 为 $R$ 在 $L$
中的整闭包。由引理 \ref{algebra-lemma-integral-closure-localize}，
$S_{f_i}$ 是 $R_{f_i}$ 在 $L$ 中的整闭包。因此按假设，
$S_{f_i}$ 在 $R_{f_i}$ 上有限。于是由引理 \ref{algebra-lemma-cover}，
$S$ 在 $R$ 上有限。
\end{proof}

\begin{lemma}
\label{algebra-lemma-quasi-finite-over-Noetherian-japanese}
设 $R$ 是整环，$R \subset S$ 是整环的拟有限扩张（例如有限扩张）。
假设 $R$ 是 N-2 且 Noether 的。则 $S$ 是 N-2 的。
\end{lemma}

\begin{proof}
设 $L/K$ 为所诱导的分式域扩张。注意，这是有限域扩张（例如，对纤维
$S \otimes_R K$ 应用引理 \ref{algebra-lemma-isolated-point-fibre} (2)，
并使用拟有限环映射的定义）。
令 $S'$ 为 $R$ 在 $S$ 中的整闭包。则 $S'$ 包含于 $R$ 在 $L$
中的整闭包；后者按假设在 $R$ 上有限。由于 $R$ 是 Noether 的，
这蕴含 $S'$ 在 $R$ 上有限。由引理
\ref{algebra-lemma-quasi-finite-open-integral-closure}，存在元素
$g_1, \ldots, g_n \in S'$，使得 $S'_{g_i} \cong S_{g_i}$，
并且 $g_1, \ldots, g_n$ 在 $S$ 中生成单位理想。因此，由引理
\ref{algebra-lemma-localize-N} 和引理 \ref{algebra-lemma-Japanese-local}，
只须证明 $S'$ 是 N-2 的。于是已经约化到 $S$ 在 $R$ 上有限的情形。

\medskip\noindent
假设 $R \subset S$ 满足引理的各项假设，并进一步假设 $S$ 在 $R$
上有限。设 $M$ 是 $S$ 的分式域的有限域扩张。则 $M$ 也是 $K$
的有限域扩张，从而 $R$ 在 $M$ 中的整闭包 $T$ 在 $R$ 上有限。
由引理 \ref{algebra-lemma-integral-closure-transitive}，$T$ 也是 $S$ 在 $M$
中的整闭包；再由引理 \ref{algebra-lemma-integral-permanence} 即得结论。
\end{proof}

\begin{lemma}
\label{algebra-lemma-Laurent-ring-N-1}
设 $R$ 是 Noether 整环。若 $R[z, z^{-1}]$ 是 N-1 的，
则 $R$ 也是 N-1 的。
\end{lemma}

\begin{proof}
设 $R'$ 为 $R$ 在其分式域 $K$ 中的整闭包。设 $S'$ 为
$R[z, z^{-1}]$ 在其分式域中的整闭包。显然 $R' \subset S'$。
由于 $K[z, z^{-1}]$ 是正规整环，可知 $S' \subset K[z, z^{-1}]$。
假设 $f_1, \ldots, f_n \in S'$ 作为 $R[z, z^{-1}]$-模生成 $S'$。
写成 $f_i = \sum a_{ij}z^j$（有限和），其中 $a_{ij} \in K$。
对任意 $x \in R'$，可以写成
$$
x = \sum h_i f_i
$$
其中 $h_i \in R[z, z^{-1}]$。因此，$R'$ 包含于有限 $R$-子模
$\sum Ra_{ij} \subset K$。由于 $R$ 是 Noether 的，可知 $R'$
是有限 $R$-模。
\end{proof}

\begin{lemma}
\label{algebra-lemma-finite-extension-N-2}
设 $R$ 是 Noether 整环，$R \subset S$ 是整环的有限扩张。
若 $S$ 是 N-1 的，则 $R$ 也是 N-1 的。
若 $S$ 是 N-2 的，则 $R$ 也是 N-2 的。
\end{lemma}

\begin{proof}
从略。（提示：$R$ 在扩域中的整闭包包含于 $S$ 在扩域中的整闭包。）
\end{proof}

\begin{lemma}
\label{algebra-lemma-Noetherian-normal-domain-finite-separable-extension}
设 $R$ 是分式域为 $K$ 的 Noether 正规整环。
设 $L/K$ 是有限可分域扩张。则 $R$ 在 $L$ 中的整闭包
在 $R$ 上有限。
\end{lemma}

\begin{proof}
考虑迹配对（《域》，定义 \ref{fields-definition-trace-pairing}）
$$
L \times L \longrightarrow K,
\quad (x, y) \longmapsto \langle x, y\rangle := \text{Trace}_{L/K}(xy).
$$
由于 $L/K$ 可分，该配对非退化（《域》，引理
\ref{fields-lemma-separable-trace-pairing}）。此外，若 $x \in L$
整于 $R$，则 $\text{Trace}_{L/K}(x)$ 属于 $R$。这是因为 $x$ 在 $K$
上的最小多项式的系数属于 $R$（引理
\ref{algebra-lemma-minimal-polynomial-normal-domain}），并且
$\text{Trace}_{L/K}(x)$ 是其中某个系数的整数倍（《域》，引理
\ref{fields-lemma-trace-and-norm-from-minimal-polynomial}）。
选取整于 $R$ 且构成 $L$ 的一组 $K$-基的元素
$x_1, \ldots, x_n \in L$。则整闭包 $S \subset L$ 包含于 $R$-模
$$
M = \{y \in L \mid \langle x_i, y\rangle \in R, \ i = 1, \ldots, n\}
$$
中。由线性代数，$M \cong R^{\oplus n}$ 是 $R$-模。因此，由于
$R$ 是 Noether 的，$S \subset R^{\oplus n}$ 是有限生成 $R$-模。
\end{proof}

\begin{example}
\label{algebra-example-bad-invariants}
若环不是 Noether 的，引理
\ref{algebra-lemma-Noetherian-normal-domain-finite-separable-extension}
不再成立。例如，考虑 $G = \{+1, -1\}$ 在
$A = \mathbf{C}[x_1, x_2, x_3, \ldots]$ 上的作用，其中 $-1$
把 $x_i$ 映到 $-x_i$。不变量环 $R = A^G$ 是由所有 $x_ix_j$
生成的 $\mathbf{C}$-代数。因此，$R \subset A$ 不是有限的。
但是 $R$ 是正规整环，其分式域 $K = L^G$ 是 $A$ 的分式域 $L$
中的 $G$-不变量；而且显然，$A$ 是 $R$ 在 $L$ 中的整闭包。
% P01-E1002: frozen English lacks a separator after K = L^G.
\end{example}

\noindent
对于纯不可分扩张，下面的引理有时可以替代引理
\ref{algebra-lemma-Noetherian-normal-domain-finite-separable-extension}。

\begin{lemma}
\label{algebra-lemma-Noetherian-normal-domain-insep-extension}
设 $R$ 是分式域为 $K$、特征为 $p > 0$ 的 Noether 正规整环。
设 $a \in K$，并且存在导子 $D : R \to R$ 使得
$D(a) \not = 0$。则 $R$ 在 $L = K[x]/(x^p - a)$ 中的整闭包
在 $R$ 上有限。
\end{lemma}

\begin{proof}
对某个 $f \in R$，把 $x$ 替换为 $fx$、把 $a$ 替换为 $f^pa$ 后，
可以假设 $a \in R$。于是也有 $D(a) \in R$。我们对
$i \leq p - 1$ 归纳证明：如果
$$
y = a_0 + a_1x + \ldots + a_i x^i,\quad a_j \in K
$$
整于 $R$，则 $D(a)^i a_j \in R$。因此，整闭包包含于以
$D(a)^{-p + 1}x^j$（$j = 0, \ldots, p - 1$）为基的有限
$R$-模中。由于 $R$ 是 Noether 的，引理得证。

\medskip\noindent
若 $i = 0$，则 $y = a_0$ 整于 $R$ 当且仅当 $a_0 \in R$，
所以断言成立。假设断言对某个 $i < p - 1$ 成立，并假设
$$
y = a_0 + a_1x + \ldots + a_{i + 1} x^{i + 1},\quad a_j \in K
$$
整于 $R$。则
$$
y^p = a_0^p + a_1^p a + \ldots + a_{i + 1}^pa^{i + 1}
$$
属于 $R$（因为它属于 $K$ 且整于 $R$）。应用 $D$，得到
$$
(a_1^p + 2a_2^p a + \ldots + (i + 1)a_{i + 1}^p a^i)D(a)
$$
属于 $R$。因此
$$
D(a)a_1 + 2D(a) a_2 x + \ldots + (i + 1)D(a) a_{i + 1} x^i
$$
整于 $R$。由归纳假设，对 $j = 1, \ldots, i + 1$，有
$D(a)^{i + 1}a_j \in R$。（这里用到 $1, \ldots, i + 1$
均可逆。）进而 $D(a)^{i + 1}a_0$ 也属于 $R$，因为它是
$D(a)^{i + 1}y$ 与 $\sum_{j > 0} D(a)^{i + 1}a_jx^j$ 之差，
而二者都整于 $R$（这是因为 $a \in R$，所以 $x$ 整于 $R$）。
\end{proof}

\begin{lemma}
\label{algebra-lemma-domain-char-zero-N-1-2}
分式域特征为零的 Noether 整环是 N-1 的，当且仅当它是 N-2 的
（即日本型的）。
\end{lemma}

\begin{proof}
这由引理 \ref{algebra-lemma-Noetherian-normal-domain-finite-separable-extension}
立即得出，因为特征为零的每个域扩张都是可分的。
\end{proof}

\begin{lemma}
\label{algebra-lemma-domain-char-p-N-1-2}
设 $R$ 是分式域为 $K$、特征为 $p > 0$ 的 Noether 整环。
则 $R$ 是 N-2 的，当且仅当对每个有限纯不可分扩张 $L/K$，
$R$ 在 $L$ 中的整闭包在 $R$ 上有限。
\end{lemma}

\begin{proof}
假设在 $K$ 的每个有限纯不可分域扩张中，$R$ 的整闭包都是有限的。
设 $L/K$ 是任意有限扩张。我们必须证明 $R$ 在 $L$ 中的整闭包
在 $R$ 上有限。选取包含 $L$ 的有限正规域扩张 $M/K$。
由于 $R$ 是 Noether 的，只须证明 $R$ 在 $M$ 中的整闭包在 $R$
上有限。由《域》，引理 \ref{fields-lemma-normal-case}，存在子扩张
$M/M_{insep}/K$，使得 $M_{insep}/K$ 纯不可分，而
$M/M_{insep}$ 可分。按假设，$R$ 在 $M_{insep}$ 中的整闭包
$R'$ 在 $R$ 上有限。由引理
\ref{algebra-lemma-Noetherian-normal-domain-finite-separable-extension}，
$R'$ 在 $M$ 中的整闭包 $R''$ 在 $R'$ 上有限。再由引理
\ref{algebra-lemma-finite-transitive}，$R''$ 在 $R$ 上有限。
由于 $R''$ 也就是 $R$ 在 $M$ 中的整闭包（见引理
\ref{algebra-lemma-integral-closure-transitive}），结论得证。
\end{proof}

\begin{lemma}
\label{algebra-lemma-polynomial-ring-N-2}
设 $R$ 是 Noether 整环。
若 $R$ 是 N-1 的，则 $R[x]$ 是 N-1 的。
若 $R$ 是 N-2 的，则 $R[x]$ 是 N-2 的。
\end{lemma}

\begin{proof}
假设 $R$ 是 N-1 的。设 $R'$ 为 $R$ 的整闭包；它在 $R$
上有限。因此，$R'[x]$ 也在 $R[x]$ 上有限。环 $R'[x]$
是正规的（见引理 \ref{algebra-lemma-polynomial-domain-normal}），
因而是 N-1 的。这证明了第一个断言。

\medskip\noindent
对于第二个断言，由引理 \ref{algebra-lemma-finite-extension-N-2}，
只须证明 $R'[x]$ 是 N-2 的。换言之，可以并且确实假设 $R$
是正规 N-2 整环。在特征零情形中，由引理
\ref{algebra-lemma-domain-char-zero-N-1-2} 即得结论。
在特征 $p > 0$ 情形中，我们必须证明：对任意有限纯不可分域扩张
$L/K(x)$，$R[x]$ 在该扩域中的整闭包都是有限的，其中 $K$ 是 $R$
的分式域。
% P01-E1003: frozen English says “extension of L/K(x)”.
存在有限纯不可分域扩张 $L'/K$ 和 $q = p^e$，使得
$L \subset L'(x^{1/q})$；一些细节从略。由于 $R[x]$ 是 Noether 的，
只须证明 $R[x]$ 在 $L'(x^{1/q})$ 中的整闭包在 $R[x]$ 上有限。
该整闭包等于 $R'[x^{1/q}]$，其中 $R \subset R' \subset L'$
是 $R$ 在 $L'$ 中的整闭包。
% P01-E1004: frozen English omits “where” before the R' apposition.
由于 $R$ 是 N-2 的，可知 $R'$ 在 $R$ 上有限，故
$R'[x^{1/q}]$ 在 $R[x]$ 上有限。
\end{proof}

\begin{lemma}
\label{algebra-lemma-openness-normal-locus}
设 $R$ 是 Noether 整环。若存在非零 $f \in R$ 使得 $R_f$
正规，则
$$
U = \{\mathfrak p \in \Spec(R) \mid R_{\mathfrak p} \text{ is normal}\}
$$
在 $\Spec(R)$ 中是开集。
\end{lemma}

\begin{proof}
显然，标准开集 $D(f)$ 包含于 $U$。由 Serre 判据，引理
\ref{algebra-lemma-criterion-normal}，$\mathfrak p \not \in U$ 蕴含：
对某个 $\mathfrak q \subset \mathfrak p$，有下列情形之一：
\begin{enumerate}
\item 情形 I：$\text{depth}(R_{\mathfrak q}) < 2$ 且
$\dim(R_{\mathfrak q}) \geq 2$；
\item 情形 II：$R_{\mathfrak q}$ 不正则且
$\dim(R_{\mathfrak q}) = 1$。
\end{enumerate}
特别地，这也意味着 $R_{\mathfrak q}$ 不正规，因而
$f \in \mathfrak q$。在情形 I 中，有
$\text{depth}(R_{\mathfrak q}) =
\text{depth}(R_{\mathfrak q}/fR_{\mathfrak q}) + 1$。
因此，这样的素理想 $\mathfrak q$ 恰是 $R/fR$ 的嵌入相伴素理想。
在情形 II 中，$\mathfrak q$ 是 $R/fR$ 的高度为 1 的相伴素理想。
所以这样的素理想 $\mathfrak q$ 构成有限集合 $E$（见引理
\ref{algebra-lemma-finite-ass}），并且
$$
\Spec(R) \setminus U
=
\bigcup\nolimits_{\mathfrak q \in E} V(\mathfrak q)
$$
如所需。
\end{proof}

\begin{lemma}
\label{algebra-lemma-characterize-N-1}
设 $R$ 是 Noether 整环。则 $R$ 是 N-1 的，当且仅当下列两个条件成立：
\begin{enumerate}
\item 存在非零 $f \in R$，使得 $R_f$ 正规；
\item 对每个极大理想 $\mathfrak m \subset R$，局部环
$R_{\mathfrak m}$ 都是 N-1 的。
\end{enumerate}
\end{lemma}

\begin{proof}
先假设 $R$ 是 N-1 的。设 $R'$ 为 $R$ 在其分式域 $K$ 中的整闭包。
按假设，可以找到 $R'$ 中的 $x_1, \ldots, x_n$，使它们作为
$R$-模生成 $R'$。由于 $R' \subset K$，可以找到非零
$f_i \in R$，使得 $f_i x_i \in R$。于是 $R_f \cong R'_f$，其中
$f = f_1 \ldots f_n$。所以 $R_f$ 正规，从而得到 (1)。
第 (2) 项由引理 \ref{algebra-lemma-localize-N} 得出。

\medskip\noindent
假设 (1) 和 (2)。设 $K$ 是 $R$ 的分式域。假设
$R \subset R' \subset K$ 是包含于 $K$ 的有限 $R$-扩张。
注意，$R_f = R'_f$，因为 $R_f$ 已经正规。因此，由引理
\ref{algebra-lemma-openness-normal-locus}，使 $R'_{\mathfrak p'}$ 不正规的
素理想 $\mathfrak p' \in \Spec(R')$ 所构成的集合在
$\Spec(R')$ 中闭。由于 $\Spec(R') \to \Spec(R)$ 是闭映射，
该集合在 $\Spec(R)$ 中的像是闭集。对这样的环 $R'$，把这个像记作
$Z_{R'} \subset \Spec(R)$。
% P01-E1005: frozen English has malformed word order in the Z_{R'} definition.

\medskip\noindent
选取极大理想 $\mathfrak m \subset R$。设
$R_{\mathfrak m} \subset R_{\mathfrak m}'$ 为该局部环在 $K$
中的整闭包。按假设，这是有限环扩张。由引理
\ref{algebra-lemma-integral-closure-localize}，可以找到有限多个整于 $R$
的元素 $x_1, \ldots, x_n \in K$，使得 $x_1, \ldots, x_n$
在 $R_{\mathfrak m}$ 上生成 $R_{\mathfrak m}'$。
令 $R' = R[x_1, \ldots, x_n] \subset K$。按此选取，显然有
$\mathfrak m \not \in Z_{R'}$。

\medskip\noindent
由于 $\Spec(R)$ 拟紧，上述论证表明，可以找到有限多个
$R \subset R'_i \subset K$，使得 $\bigcap Z_{R'_i} = \emptyset$。
令 $R'$ 为 $K$ 中由所有这些环生成的子环。它在 $R$ 上有限。
而且 $Z_{R'} = \emptyset$。事实上，每个素理想 $\mathfrak p'$
都位于某个素理想 $\mathfrak p'_i$ 之上，而
$(R'_i)_{\mathfrak p'_i}$ 正规。这蕴含
$R'_{\mathfrak p'} = (R'_i)_{\mathfrak p'_i}$ 也正规。
因此 $R'$ 正规；换言之，$R'$ 是 $R$ 在 $K$ 中的整闭包。
\end{proof}

\begin{lemma}[Tate]
\label{algebra-lemma-tate-japanese}
\begin{reference}
\cite[Theorem 23.1.3]{EGA}
\end{reference}
设 $R$ 是环，$x \in R$。假设
\begin{enumerate}
\item $R$ 是 Noether 正规整环；
\item $R/xR$ 是整环且为 N-2；
\item $R \cong \lim_n R/x^nR$ 关于 $x$ 完备。
\end{enumerate}
则 $R$ 是 N-2 的。
\end{lemma}

\begin{proof}
可以假设 $x \not = 0$，否则引理显然。设 $K$ 为 $R$ 的分式域。
若 $K$ 的特征为零，则由 (1) 和引理
\ref{algebra-lemma-domain-char-zero-N-1-2} 即得结论。因此，可以假设
$K$ 的特征为 $p > 0$，并应用引理
\ref{algebra-lemma-domain-char-p-N-1-2}。于是，给定有限纯不可分域扩张
$L/K$，必须证明 $R$ 在 $L$ 中的整闭包 $S$ 在 $R$ 上有限。

\medskip\noindent
设 $q$ 是 $p$ 的幂，并使得 $L^q \subset K$。必要时扩大 $L$，
可以假设存在 $y \in L$ 使得 $y^q = x$。由于 $R \to S$
诱导谱上的同胚（见引理 \ref{algebra-lemma-p-ring-map}），存在唯一素理想
$\mathfrak q \subset S$ 位于素理想 $\mathfrak p = xR$ 之上。
显然
$$
\mathfrak q = \{f \in S \mid f^q \in \mathfrak p\} = yS
$$
因为 $y^q = x$。由引理 \ref{algebra-lemma-characterize-dvr}，
$R_{\mathfrak p}$ 是离散赋值环。于是由 Krull--Akizuki
（引理 \ref{algebra-lemma-krull-akizuki}），$S_{\mathfrak q}$ 是
Noether 的。再次应用引理 \ref{algebra-lemma-characterize-dvr}，
可知 $S_{\mathfrak q}$ 是离散赋值环。由引理
\ref{algebra-lemma-finite-extension-residue-fields-dimension-1}，
$\kappa(\mathfrak q)/\kappa(\mathfrak p)$ 是有限域扩张。
因此，按假设 (2)，$R/xR$ 的整闭包 $S' \subset \kappa(\mathfrak q)$
在 $R/xR$ 上有限。
由于 $S/yS \subset S'$，这蕴含 $S/yS$ 在 $R$ 上有限。
注意，$S/y^nS$ 有有限滤过，其子商为模
$y^iS/y^{i + 1}S \cong S/yS$。因此，每个 $S/y^nS$ 都在 $R$
上有限。特别地，$S/xS$ 在 $R$ 上有限。此外，显然
$\bigcap x^nS = (0)$，因为交中元素的 $q$ 次幂属于
$\bigcap x^nR = (0)$（引理
\ref{algebra-lemma-intersect-powers-ideal-module-zero}）。
所以可以应用引理 \ref{algebra-lemma-finite-over-complete-ring}，得出
$S$ 在 $R$ 上有限，结论得证。
\end{proof}

\begin{lemma}
\label{algebra-lemma-power-series-over-N-2}
设 $R$ 是环。若 $R$ 是 Noether 的整环且为 N-2，
则 $R[[x]]$ 也是如此。
\end{lemma}

\begin{proof}
由引理 \ref{algebra-lemma-Noetherian-power-series}，$R[[x]]$ 是 Noether 的。
设 $R' \supset R$ 为 $R$ 在其分式域中的整闭包。由于 $R$ 是 N-2 的，
这个整闭包在 $R$ 上有限。因此，$R'[[x]]$ 在 $R[[x]]$ 上有限。
由引理 \ref{algebra-lemma-power-series-over-Noetherian-normal-domain}，
$R'[[x]]$ 是正规整环。对元素 $x \in R'[[x]]$ 应用引理
\ref{algebra-lemma-tate-japanese}，可知 $R'[[x]]$ 是 N-2 的。
再由引理 \ref{algebra-lemma-finite-extension-N-2}，$R[[x]]$ 是 N-2 的。
\end{proof}

\section{Nagata 环}
\label{algebra-section-nagata}

\noindent
定义如下。

\begin{definition}
\label{algebra-definition-nagata}
设 $R$ 是环。
\begin{enumerate}
\item 如果对任意有限型环映射 $R \to S$，只要 $S$ 是整环，
就有 $S$ 为 N-2（即日本型的），则称 $R$ 是{\it 广泛日本型的}。
\item 如果 $R$ 是 Noether 的，并且对每个素理想 $\mathfrak p$，
环 $R/\mathfrak p$ 都是 N-2 的，则称 $R$ 是{\it Nagata 环}。
\end{enumerate}
\end{definition}

\noindent
显然，Noether 的广泛日本型环是 Nagata 环。我们的目标是证明
Nagata 环是广泛日本型的。这一点丝毫不显然，需要做一些工作。
不过先给出一个有用的引理。

\begin{lemma}
\label{algebra-lemma-nagata-in-reduced-finite-type-finite-integral-closure}
设 $R$ 是 Nagata 环。设 $R \to S$ 本质上为有限型，且 $S$ 约化。
则 $R$ 在 $S$ 中的整闭包 $A$ 在 $R$ 上有限。
\end{lemma}

\begin{proof}
由于 $S$ 本质上在 $R$ 上有限型，它是 Noether 的，并且只有有限多个
极小素理想 $\mathfrak q_1, \ldots, \mathfrak q_m$；见引理
\ref{algebra-lemma-Noetherian-irreducible-components}。由于 $S$ 约化，
有 $S \subset \prod S_{\mathfrak q_i}$，且每个
$S_{\mathfrak q_i} = K_i$ 都是域；见引理
\ref{algebra-lemma-total-ring-fractions-no-embedded-points}
和 \ref{algebra-lemma-minimal-prime-reduced-ring}。只需证明 $R$ 在每个 $K_i$
中的整闭包 $A_i'$ 在 $R$ 上有限。之所以如此，是因为 $R$ 是 Noether 的，
且 $A \subset \prod A_i'$。设 $\mathfrak p_i \subset R$ 是与
$\mathfrak q_i$ 对应的 $R$ 的素理想。由于 $S$ 本质上在 $R$ 上有限型，
可知 $K_i = S_{\mathfrak q_i} = \kappa(\mathfrak q_i)$ 是
$\kappa(\mathfrak p_i)$ 的有限生成域扩张。因此，
$\kappa(\mathfrak p_i)$ 在 $K_i$ 中的代数闭包 $L_i$ 是
$\kappa(\mathfrak p_i)$ 的有限扩张；见《域》，引理
\ref{fields-lemma-algebraic-closure-in-finitely-generated}。
显然，$A_i'$ 是 $R/\mathfrak p_i$ 在 $L_i$ 中的整闭包，
所以由 Nagata 环的定义即得结论。
\end{proof}

\begin{lemma}
\label{algebra-lemma-check-universally-japanese}
设 $R$ 是环。为检验 $R$ 是广泛日本型的，只需证明如下命题：
若 $R \to S$ 为有限型且 $S$ 是整环，则 $S$ 是 N-1 的。
\end{lemma}

\begin{proof}
假设引理中的条件成立。设 $R \to S$ 是有限型环映射，且 $S$ 是整环。
设 $L$ 是 $S$ 的分式域的有限扩张。则存在有限环扩张
$S \subset S' \subset L$，使 $L$ 为 $S'$ 的分式域。
由假设，$S'$ 是 N-1 的，因此 $S'$ 在 $L$ 中的整闭包 $S''$
在 $S'$ 上有限。于是 $S''$ 在 $S$ 上有限（引理
\ref{algebra-lemma-finite-transitive}），并且 $S''$ 是 $S$ 在 $L$ 中的整闭包
（引理 \ref{algebra-lemma-integral-closure-transitive}）。所以 $R$ 是广泛日本型的。
\end{proof}

\begin{lemma}
\label{algebra-lemma-universally-japanese}
若 $R$ 是广泛日本型的，则任何本质上在 $R$ 上有限型的代数也都是
广泛日本型的。
\end{lemma}

\begin{proof}
在代数在 $R$ 上有限型的情形，结论直接来自定义。一般情形由应用引理
\ref{algebra-lemma-localize-N} 得到。
\end{proof}

\begin{lemma}
\label{algebra-lemma-quasi-finite-over-nagata}
设 $R$ 是 Nagata 环。若 $R \to S$ 是拟有限环映射（例如有限环映射），
则 $S$ 也是 Nagata 环。
\end{lemma}

\begin{proof}
先注意到 $S$ 是 Noether 的，因为 $R$ 是 Noether 的，而拟有限环映射
是有限型的。设 $\mathfrak q \subset S$ 是素理想，并令
$\mathfrak p = R \cap \mathfrak q$。则
$R/\mathfrak p \subset S/\mathfrak q$ 是拟有限的，因而由引理
\ref{algebra-lemma-quasi-finite-over-Noetherian-japanese} 可知
$S/\mathfrak q$ 是 N-2 的，如所求。
\end{proof}

\begin{lemma}
\label{algebra-lemma-nagata-localize}
Nagata 环的局部化仍是 Nagata 环。
\end{lemma}

\begin{proof}
直接由引理 \ref{algebra-lemma-localize-N} 得到。
\end{proof}

\begin{lemma}
\label{algebra-lemma-nagata-local}
设 $R$ 是环。设 $f_1, \ldots, f_n \in R$ 生成单位理想。
\begin{enumerate}
\item  如果每个 $R_{f_i}$ 都是广泛日本型的，则 $R$ 也是如此。
\item  如果每个 $R_{f_i}$ 都是 Nagata 环，则 $R$ 也是 Nagata 环。
\end{enumerate}
\end{lemma}

\begin{proof}
设 $\varphi : R \to S$ 是有限型环映射，且 $S$ 是整环。
则 $\varphi(f_1), \ldots, \varphi(f_n)$ 生成 $S$ 的单位理想。
因此，若每个 $S_{f_i} = S_{\varphi(f_i)}$ 都是 N-1 的，
则 $S$ 也是 N-1 的；见引理 \ref{algebra-lemma-Japanese-local}。这证明了 (1)。

\medskip\noindent
若每个 $R_{f_i}$ 都是 Nagata 环，则每个 $R_{f_i}$ 都是 Noether 的，
因而 $R$ 是 Noether 的；见引理 \ref{algebra-lemma-cover}。再设
$\mathfrak p \subset R$ 是素理想，则每个
$R_{f_i}/\mathfrak pR_{f_i} = (R/\mathfrak p)_{f_i}$ 都是 N-2 的，
所以由引理 \ref{algebra-lemma-Japanese-local} 可知 $R/\mathfrak p$ 是 N-2 的。
这证明了 (2)。
\end{proof}

\begin{lemma}
\label{algebra-lemma-Noetherian-complete-local-Nagata}
Noether 完备局部环是 Nagata 环。
\end{lemma}

\begin{proof}
设 $R$ 是完备的 Noether 局部环。设 $\mathfrak p \subset R$ 是素理想。
则 $R/\mathfrak p$ 也是完备的 Noether 局部环；见引理
\ref{algebra-lemma-quotient-complete-local}。因此，只需证明完备的 Noether
局部整环 $R$ 是 N-2 的。由引理
\ref{algebra-lemma-quasi-finite-over-Noetherian-japanese}
和 \ref{algebra-lemma-complete-local-Noetherian-domain-finite-over-regular}，
可化归到 $R = k[[X_1, \ldots, X_d]]$，其中 $k$ 是域；或者化归到
$R = \Lambda[[X_1, \ldots, X_d]]$，其中 $\Lambda$ 是 Cohen 环。

\medskip\noindent
在 $k[[X_1, \ldots, X_d]]$ 的情形，由引理
\ref{algebra-lemma-power-series-over-N-2} 可化归到域为 N-2 的命题，
这是显然的。在 $\Lambda[[X_1, \ldots, X_d]]$ 的情形，
可化归到 Cohen 环 $\Lambda$ 是 N-2 的命题。再次对
$x = p \in \Lambda$ 应用引理 \ref{algebra-lemma-tate-japanese}，
又一次化归到域的情形。因此结论成立。
\end{proof}

\begin{definition}
\label{algebra-definition-analytically-unramified}
设 $(R, \mathfrak m)$ 是 Noether 局部环。如果它的完备化
$R^\wedge = \lim_n R/\mathfrak m^n$ 是约化的，则称 $R$
是{\it 解析非分歧的}。若素理想 $\mathfrak p \subset R$ 使
$R/\mathfrak p$ 解析非分歧，则称该素理想是
{\it 解析非分歧的}。
\end{definition}

\noindent
至此，对任意 Noether 局部环 $R$，我们已经知道以下事实成立：
映射 $R \to R^\wedge$ 是忠实平坦局部环同态（引理
\ref{algebra-lemma-completion-faithfully-flat}）。完备化 $R^\wedge$
是 Noether 的（引理 \ref{algebra-lemma-completion-Noetherian}），
并且是完备的（引理 \ref{algebra-lemma-completion-complete}）。因此，
完备化 $R^\wedge$ 是 Nagata 环（引理
\ref{algebra-lemma-Noetherian-complete-local-Nagata}）。此外，我们已经在节
\ref{algebra-section-cohen-structure-theorem} 中看到，$R^\wedge$ 是正则局部环的商
（定理 \ref{algebra-theorem-cohen-structure-theorem}），从而是普遍链状的
（评注 \ref{algebra-remark-Noetherian-complete-local-ring-universally-catenary}）。

\begin{lemma}
\label{algebra-lemma-analytically-unramified-easy}
设 $(R, \mathfrak m)$ 是 Noether 局部环。
\begin{enumerate}
\item 若 $R$ 解析非分歧，则 $R$ 约化。
\item 若 $R$ 解析非分歧，则 $R$ 的每个极小素理想都解析非分歧。
\item 若 $R$ 约化，其极小素理想为
$\mathfrak q_1, \ldots, \mathfrak q_t$，并且每个 $\mathfrak q_i$
都解析非分歧，则 $R$ 解析非分歧。
\item 若 $R$ 解析非分歧，则 $R$ 在其全分式环 $Q(R)$ 中的整闭包
在 $R$ 上有限。
\item 若 $R$ 是整环且解析非分歧，则 $R$ 是 N-1 的。
\end{enumerate}
\end{lemma}

\begin{proof}
在本证明中，我们将使用紧接在定义
\ref{algebra-definition-analytically-unramified} 后的评注。由于
$R \to R^\wedge$ 是忠实平坦局部环同态，它是单射，因此 (1) 成立。

\medskip\noindent
设 $\mathfrak q$ 是 $R$ 的极小素理想，并假设 $R$ 解析非分歧。
则 $\mathfrak q$ 是 $R$ 的相伴素理想（见命题
\ref{algebra-proposition-minimal-primes-associated-primes}）。因此存在
$f \in R$，使得 $\{x \in R \mid fx = 0\} = \mathfrak q$。
注意 $(R/\mathfrak q)^\wedge = R^\wedge/\mathfrak q^\wedge$，
并且 $\{x \in R^\wedge \mid fx = 0\} = \mathfrak q^\wedge$，
因为完备化是正合的（引理 \ref{algebra-lemma-completion-flat}）。若
$x \in R^\wedge$ 满足 $x^2 \in \mathfrak q^\wedge$，则
$fx^2 = 0$，从而 $(fx)^2 = 0$，从而 $fx = 0$，从而
$x \in \mathfrak q^\wedge$。所以 $\mathfrak q$ 解析非分歧，
(2) 成立。

\medskip\noindent
假设 $R$ 约化，其极小素理想为
$\mathfrak q_1, \ldots, \mathfrak q_t$，且每个 $\mathfrak q_i$
都解析非分歧。则映射
$R \to R/\mathfrak q_1 \times \ldots \times R/\mathfrak q_t$ 是单射。
由于完备化正合（见引理 \ref{algebra-lemma-completion-flat}），可知
$R^\wedge \subset (R/\mathfrak q_1)^\wedge \times \ldots \times
(R/\mathfrak q_t)^\wedge$。因此 (3) 显然。

\medskip\noindent
假设 $R$ 解析非分歧。设 $\mathfrak p_1, \ldots, \mathfrak p_s$
为 $R^\wedge$ 的极小素理想。则
$$
Q(R^\wedge) =
R^\wedge_{\mathfrak p_1} \times \ldots \times R^\wedge_{\mathfrak p_s}
$$
且每个 $R^\wedge_{\mathfrak p_i}$ 都是域，因为 $R^\wedge$ 约化
（见引理 \ref{algebra-lemma-total-ring-fractions-no-embedded-points}）。
因此，$R^\wedge$ 在 $Q(R^\wedge)$ 中的整闭包 $S$ 等于
$S = S_1 \times \ldots \times S_s$，其中 $S_i$ 是
$R^\wedge/\mathfrak p_i$ 在其分式域中的整闭包。由引理
\ref{algebra-lemma-Noetherian-complete-local-Nagata}，环 $S_i$ 在
$R^\wedge/\mathfrak p_i$ 上有限。因此，$S$ 在 $R^\wedge$ 上有限。
% P01-E1006
用 $R'$ 表示 $R$ 在 $Q(R)$ 中的整闭包。由于
$R \to R^\wedge$ 平坦，可知
$R' \otimes_R R^\wedge \subset Q(R) \otimes_R R^\wedge \subset Q(R^\wedge)$。
此外，$R' \otimes_R R^\wedge$ 整于 $R^\wedge$
（引理 \ref{algebra-lemma-base-change-integral}）。
% P01-E1007
所以 $R' \otimes_R R^\wedge \subset S$ 是 $R^\wedge$-子模。
由于 $R^\wedge$ 是 Noether 的，它是有限 $R^\wedge$-模。
因此可取 $f_1, \ldots, f_n \in R'$，使
$R' \otimes_R R^\wedge$ 作为 $R^\wedge$-模由元素
$f_i \otimes 1$ 生成。由忠实平坦性可知，$R'$ 作为 $R$-模由
$f_1, \ldots, f_n$ 生成。这证明了 (4)。

\medskip\noindent
(5) 是 (4) 的特例。
\end{proof}

\begin{lemma}
\label{algebra-lemma-codimension-1-analytically-unramified}
设 $R$ 是 Noether 局部环。设 $\mathfrak p \subset R$ 是素理想。假设
\begin{enumerate}
\item $R_{\mathfrak p}$ 是离散赋值环；并且
\item $\mathfrak p$ 解析非分歧。
\end{enumerate}
则对 $R^\wedge/\mathfrak pR^\wedge$ 的任意相伴素理想
$\mathfrak q$，局部环 $(R^\wedge)_{\mathfrak q}$ 是离散赋值环。
\end{lemma}

\begin{proof}
假设 (2) 表明 $R^\wedge/\mathfrak pR^\wedge$ 是约化环。
因此，$R^\wedge/\mathfrak pR^\wedge$ 的相伴素理想
$\mathfrak q \subset R^\wedge$，等价于 $\mathfrak pR^\wedge$
之上的极小素理想。特别地，$(R^\wedge)_{\mathfrak q}$ 的极大理想为
$\mathfrak p(R^\wedge)_{\mathfrak q}$。选取 $x \in R$，使得
$xR_{\mathfrak p} = \mathfrak pR_{\mathfrak p}$。由上可知，
$x \in (R^\wedge)_{\mathfrak q}$ 生成极大理想。由于
$R \to R^\wedge$ 忠实平坦，$x$ 是
$(R^\wedge)_{\mathfrak q}$ 中的非零因子。因此结论成立。
\end{proof}

\begin{lemma}
\label{algebra-lemma-criterion-analytically-unramified}
设 $(R, \mathfrak m)$ 是 Noether 局部整环。设 $x \in \mathfrak m$。假设
\begin{enumerate}
\item $x \not = 0$；
\item $R/xR$ 没有嵌入素理想；并且
\item 对 $R/xR$ 的每个相伴素理想 $\mathfrak p \subset R$，都有
\begin{enumerate}
\item 局部环 $R_{\mathfrak p}$ 是正则的；并且
\item $\mathfrak p$ 解析非分歧。
\end{enumerate}
\end{enumerate}
则 $R$ 解析非分歧。
\end{lemma}

\begin{proof}
设 $\mathfrak p_1, \ldots, \mathfrak p_t$ 为 $R$-模 $R/xR$
的相伴素理想。由于 $R/xR$ 没有嵌入素理想，可知每个
$\mathfrak p_i$ 都具有高度 $1$，并且是 $(x)$ 之上的极小素理想。
对每个 $i$，设
$\mathfrak q_{i1}, \ldots, \mathfrak q_{is_i}$ 为 $R^\wedge$-模
$R^\wedge/\mathfrak p_iR^\wedge$ 的相伴素理想。由引理
\ref{algebra-lemma-codimension-1-analytically-unramified}，
$(R^\wedge)_{\mathfrak q_{ij}}$ 是正则的。由引理
\ref{algebra-lemma-bourbaki}，有
$$
\text{Ass}_{R^\wedge}(R^\wedge/xR^\wedge)
=
\bigcup\nolimits_{\mathfrak p \in \text{Ass}_R(R/xR)}
\text{Ass}_{R^\wedge}(R^\wedge/\mathfrak pR^\wedge)
=
\{\mathfrak q_{ij}\}.
$$
设 $y \in R^\wedge$ 且 $y^2 = 0$。由于
$(R^\wedge)_{\mathfrak q_{ij}}$ 正则，从而是整环
（引理 \ref{algebra-lemma-regular-domain}），可知 $y$ 在
$(R^\wedge)_{\mathfrak q_{ij}}$ 中的像为零。因此，由引理
\ref{algebra-lemma-zero-at-ass-zero}，$y$ 在 $R^\wedge/xR^\wedge$
中的像为零。所以 $y = xy'$。由于 $x$ 是非零因子
（因为 $R \to R^\wedge$ 平坦），可知 $(y')^2 = 0$。
因此
$y \in \bigcap x^nR^\wedge = (0)$
（引理 \ref{algebra-lemma-intersect-powers-ideal-module-zero}）。
\end{proof}

\begin{lemma}
\label{algebra-lemma-local-nagata-domain-analytically-unramified}
设 $(R, \mathfrak m)$ 是局部环。若 $R$ 是 Noether 的 Nagata 整环，
则 $R$ 解析非分歧。
\end{lemma}

\begin{proof}
对 $\dim(R)$ 作归纳。当 $\dim(R) = 0$ 时结论平凡。
因此，假设 $\dim(R) = d$，并假设引理对维数 $< d$ 的所有
Noether Nagata 整环都成立。

\medskip\noindent
设 $R \subset S$ 是 $R$ 在 $R$ 的分式域中的整闭包。由假设，$S$ 是有限
$R$-模。由引理 \ref{algebra-lemma-quasi-finite-over-nagata}，$S$ 是 Nagata 环。
由引理 \ref{algebra-lemma-integral-sub-dim-equal}，有
$\dim(R) = \dim(S)$。设 $\mathfrak m_1, \ldots, \mathfrak m_t$
是 $S$ 的极大理想。它们都位于 $R$ 的极大理想 $\mathfrak m$ 之上。
此外，由于 $S/\mathfrak mS$ 是 Artin 的，对充分大的 $n$ 有
$$
(\mathfrak m_1 \cap \ldots \cap \mathfrak m_t)^n \subset \mathfrak mS
$$
由引理 \ref{algebra-lemma-completion-flat}，$R^\wedge \to S^\wedge$
是单射；再把中国剩余引理 \ref{algebra-lemma-chinese-remainder} 与引理
\ref{algebra-lemma-change-ideal-completion} 结合起来，可得
$S^\wedge = \prod S^\wedge_i$，其中 $S^\wedge_i$ 是 $S$
关于极大理想 $\mathfrak m_i$ 的完备化。因此，只需证明
$S_{\mathfrak m_i}$ 解析非分歧。换言之，我们已经化归到
$R$ 是 Noether 正规 Nagata 整环的情形。

\medskip\noindent
假设 $R$ 是 Noether 正规局部 Nagata 整环。在极大理想中选取
非零元素 $x \in \mathfrak m$。我们将应用引理
\ref{algebra-lemma-criterion-analytically-unramified}。必须检验性质
(1)、(2)、(3)(a) 和 (3)(b)。性质 (1) 显然。由引理
\ref{algebra-lemma-normal-domain-intersection-localizations-height-1}，
$R/xR$ 没有嵌入素理想，所以性质 (2) 成立。同一个引理还说明，
$R/xR$ 的每个相伴素理想 $\mathfrak p$ 都具有高度 $1$。
因此，$R_{\mathfrak p}$ 是 $1$-维正规整环，从而是正则的
（引理 \ref{algebra-lemma-characterize-dvr}）。于是 (3)(a) 成立。
最后，(3)(b) 由归纳假设成立，因为 $R/\mathfrak p$ 是 Nagata 环
（由引理 \ref{algebra-lemma-quasi-finite-over-nagata}，或直接由定义）。
所以 $R$ 解析非分歧。
\end{proof}

\begin{lemma}
\label{algebra-lemma-local-nagata-and-analytically-unramified}
% P01-E1008
设 $(R, \mathfrak m)$ 是 Noether 局部环。以下条件等价
\begin{enumerate}
\item $R$ 是 Nagata 环；
\item 对有限映射 $R \to S$，若 $S$ 是整环且
$\mathfrak m' \subset S$ 是极大理想，则局部环
$S_{\mathfrak m'}$ 解析非分歧；
% P01-E1009
\item 对有限局部同态 $(R, \mathfrak m) \to (S, \mathfrak m')$，
若 $S$ 是整环，则 $S$ 解析非分歧。
\end{enumerate}
\end{lemma}

\begin{proof}
假设 $R$ 是 Nagata 环，并设 $R \to S$ 和 $\mathfrak m' \subset S$
如 (2) 所述。由引理 \ref{algebra-lemma-quasi-finite-over-nagata}，$S$ 是
Nagata 环。因此，局部环 $S_{\mathfrak m'}$ 是 Nagata 环
（引理 \ref{algebra-lemma-nagata-localize}）。再由引理
\ref{algebra-lemma-local-nagata-domain-analytically-unramified}，
它解析非分歧。显然，(2) 蕴含 (3)。

\medskip\noindent
假设 (3) 成立。设 $\mathfrak p \subset R$ 是素理想，并设
$L/\kappa(\mathfrak p)$ 是有限域扩张。
% P01-E1010
为证明 (1)，必须证明 $R/\mathfrak p$ 的整闭包在
$R/\mathfrak p$ 上有限。选取 $x_1, \ldots, x_n \in L$，
使它们生成 $L$ 作为 $\kappa(\mathfrak p)$ 的扩张。对每个 $i$，设
$P_i(T) = T^{d_i} + a_{i, 1} T^{d_i - 1} + \ldots + a_{i, d_i}$
是 $x_i$ 在 $\kappa(\mathfrak p)$ 上的极小多项式。
对适当的 $f_i \in R$，$f_i \not \in \mathfrak p$，
用 $f_i x_i$ 代替 $x_i$ 后，可以假设
$a_{i, j} \in R/\mathfrak p$。事实上，再乘以 $\mathfrak m$
中的元素后，可以假设对所有 $i, j$ 都有
$a_{i, j} \in \mathfrak m/\mathfrak p \subset R/\mathfrak p$。
% P01-E1011
完成这些操作后，令 $S = R/\mathfrak p[x_1, \ldots, x_n] \subset L$。
则 $S$ 在 $R$ 上有限，是整环，并且
% P01-E1012
$S/\mathfrak m S$ 是
$R/\mathfrak m[T_1, \ldots, T_n]/(T_1^{d_1}, \ldots, T_n^{d_n})$
的商。因此 $S$ 是局部环。由 (3)，$S$ 解析非分歧；再由引理
\ref{algebra-lemma-analytically-unramified-easy}，它在 $L$ 中的整闭包
$S'$ 在 $S$ 上有限。由于 $S'$ 也是 $R/\mathfrak p$ 在 $L$
中的整闭包，结论得证。
\end{proof}

\noindent
下面的命题特别说明，Nagata 环上的有限型代数仍是 Nagata 环。

\begin{proposition}[Nagata]
\label{algebra-proposition-nagata-universally-japanese}
设 $R$ 是环。以下条件等价：
\begin{enumerate}
\item $R$ 是 Nagata 环；
\item 任意有限型 $R$-代数都是 Nagata 环；并且
\item $R$ 是广泛日本型的且为 Noether 的。
\end{enumerate}
\end{proposition}

\begin{proof}
显然，Noether 的广泛日本型环是广泛 Nagata 的（即条件 (2) 成立）。
设 $R$ 是 Nagata 环。我们将证明任意有限生成 $R$-代数 $S$
都是 Nagata 环。这将证明命题。

\medskip\noindent
第 1 步。存在环映射列
$R = R_0 \to R_1 \to R_2 \to \ldots \to R_n = S$，
使得每个 $R_i \to R_{i + 1}$ 都由单个元素生成。因此，通过归纳，
只需在 $S \cong R[x]/I$ 时证明 $S$ 是 Nagata 环。

\medskip\noindent
第 2 步。设 $\mathfrak q \subset S$ 是 $S$ 的素理想，并设
$\mathfrak p \subset R$ 是相应的 $R$ 的素理想。必须证明
$S/\mathfrak q$ 是 N-2 的。因此，我们已经化归到证明如下命题：
(*) 给定 Nagata 整环 $R$ 和整环的单生成扩张 $R \subset S$，
则 $S$ 是 N-2 的。

\medskip\noindent
第 3 步：设 $R$ 是 Nagata 整环，并设 $R \subset S$ 是整环的
单生成扩张。假设 $R$ 和 $S$ 的分式域之间所诱导的扩张是纯超越的。
在这种情况下，$S = R[x]$。由引理 \ref{algebra-lemma-polynomial-ring-N-2}，
$S$ 是 N-2 的。因此，我们已经化归到证明如下命题：
(**) 给定 Nagata 整环 $R$ 和整环的单生成扩张 $R \subset S$，
若所诱导的分式域扩张有限，则 $S$ 是 N-2 的。

\medskip\noindent
第 4 步。设 $R$ 是正规 Nagata 整环，并设 $R \subset S$ 是整环的
单生成扩张，所诱导的分式域扩张为有限扩张 $L/K$。
选取生成元 $x \in S$，使 $S$ 作为 $R$-代数由它生成。
设 $M/L$ 是有限域扩张。设 $R'$ 是 $R$ 在 $M$ 中的整闭包。
则 $S$ 在 $M$ 中的整闭包 $S'$ 等于 $R'[x]$ 在 $M$ 中的整闭包。
此外，$R'$ 的分式域是 $M$，且 $R \subset R'$ 有限
（由 $R$ 的 Nagata 性质）。这说明 $R'$ 是 Nagata 环
（引理 \ref{algebra-lemma-quasi-finite-over-nagata}）。证明 $S'$ 在 $S$
上有限，等价于证明 $S'$ 在 $R'[x]$ 上有限。用 $R'$ 代替 $R$，
并用 $R'[x]$ 代替 $S$，即可化归到如下命题：
(***) 给定分式域为 $K$ 的正规 Nagata 整环 $R$，以及 $x \in K$，
则由 $R$ 和 $x$ 生成的子环 $S \subset K$ 是 N-1 的。

\medskip\noindent
第 5 步。设 $R$ 是分式域为 $K$ 的正规 Nagata 整环。
设 $x = b/a \in K$。必须证明由 $R$ 和 $x$ 生成的子环
$S \subset K$ 是 N-1 的。注意，$S_a \cong R_a$ 是正规的。
因此，由引理 \ref{algebra-lemma-characterize-N-1}，只需对 $S$ 的每个极大理想
$\mathfrak m$ 证明 $S_{\mathfrak m}$ 是 N-1 的。

\medskip\noindent
在前一段的假设下，选定这样的极大理想，并令
$\mathfrak n = R \cap \mathfrak m$。剩余域扩张
$\kappa(\mathfrak m)/\kappa(\mathfrak n)$ 是有限的
（定理 \ref{algebra-theorem-nullstellensatz}），并由 $x$ 的像生成。
% P01-E1013
对某个 $a \in R$，$a \not \in \mathfrak n$，用 $R_a$ 代替 $R$，
并用 $S_a$ 代替 $S$ 后，可以假设存在首一多项式
% P01-E1014 P01-E1015
$f(X) = X^d + \sum_{i = 1, \ldots, d} a_iX^{d - i}$ 属于 $R[x]$，
使得 $f(x) \in \mathfrak m$（细节略）。设 $K''/K$ 是有限域扩张，
使多项式 $f(X)$ 在 $K''[X]$ 中完全分裂。设 $R'$ 是 $R$ 在
$K''$ 中的整闭包。设 $S' \subset K''$ 是由 $R'$ 和 $x$ 生成的子环。
由于 $R$ 是 Nagata 环，可知 $R'$ 在 $R$ 上有限且为 Nagata 环
（引理 \ref{algebra-lemma-quasi-finite-over-nagata}）。此外，$S'$ 在 $S$
上有限。如果对 $S'$ 的每个极大理想 $\mathfrak m'$，局部环
$S'_{\mathfrak m'}$ 都是 N-1 的，则由引理
\ref{algebra-lemma-characterize-N-1}，$S'_{\mathfrak m}$ 是 N-1 的；
进而由引理 \ref{algebra-lemma-finite-extension-N-2}，
$S_{\mathfrak m}$ 是 N-1 的。用 $R'$ 代替 $R$，用 $S'$ 代替 $S$，
再用位于 $\mathfrak m$ 之上的任意极大理想 $\mathfrak m'$ 代替
$\mathfrak m$ 后，我们达到上述多项式 $f$ 完全分裂的情形：
$f(X) = \prod_{i = 1, \ldots, d} (X - a_i)$，其中 $a_i \in R$。
由于 $f(x) \in \mathfrak m$，对某个 $i$ 有
$x - a_i \in \mathfrak m$。最后，用 $x - a_i$ 代替 $x$ 后，
可以假设 $x \in \mathfrak m$。

\medskip\noindent
概括起来：$R$ 是分式域为 $K$ 的正规 Nagata 整环，$x \in K$，
$S$ 是 $K$ 中由 $x$ 和 $R$ 生成的子环；最后，
$\mathfrak m \subset S$ 是极大理想，且 $x \in \mathfrak m$。必须证明
$S_{\mathfrak m}$ 是 N-1 的。

\medskip\noindent
我们将证明引理 \ref{algebra-lemma-criterion-analytically-unramified}
可应用于局部环 $S_{\mathfrak m}$ 和元素 $x$。这将说明
$S_{\mathfrak m}$ 解析非分歧；随后由引理
\ref{algebra-lemma-analytically-unramified-easy} 可知它是 N-1 的。

\medskip\noindent
必须检验性质 (1)、(2)、(3)(a) 和 (3)(b)。性质 (1) 平凡。
设 $I = \Ker(R[X] \to S)$，其中 $X \mapsto x$。我们断言，
$I$ 由所有满足 $ax = b$ 于 $K$ 中的线性式 $aX - b$ 生成。
显然，所有这些线性式都属于 $I$。若
$g = a_d X^d + \ldots a_1 X + a_0 \in I$，则 $a_dx$ 整于 $R$
（引理 \ref{algebra-lemma-make-integral-trivial}），从而
$b := a_dx \in R$，因为 $R$ 正规。于是
$g - (a_dX - b)X^{d - 1} \in I$，再对次数作归纳即得断言。
作为推论，
$$
S/xS = R[X]/(X, I) = R/J
$$
其中
$$
J = \{b \in R \mid ax = b \text{ 对某个 }a \in R\} = xR \cap R
$$
由引理 \ref{algebra-lemma-normal-domain-intersection-localizations-height-1}，
$S/xS = R/J$ 作为 $R$-模没有嵌入素理想，因而作为 $R/J$-模没有，
因而作为 $S/xS$-模没有，进而作为 $S$-模也没有。这证明了性质 (2)。
取这样的相伴素理想 $\mathfrak q \subset S$，并要求
$\mathfrak q \subset \mathfrak m$（从而它是
$S_{\mathfrak m}/xS_{\mathfrak m}$ 的相伴素理想——这对论证并不重要）。
则 $\mathfrak q$ 是 $xS$ 之上的极小素理想，因而具有高度 $1$。
由上述一系列等式可知，$\mathfrak p = R \cap \mathfrak q$ 是
$R/J$ 的相伴素理想，因此具有高度 $1$（见引理
\ref{algebra-lemma-normal-domain-intersection-localizations-height-1}）。
所以 $R_{\mathfrak p}$ 是离散赋值环，从而包含
$R_{\mathfrak p} \subset S_{\mathfrak q}$ 实为等式。
这说明 $S_{\mathfrak q}$ 正则，证明了性质 (3)(a)。最后，
$(S/\mathfrak q)_{\mathfrak m}$ 是 $S/\mathfrak q$ 的局部化，
而后者是 $S/xS = R/J$ 的商。因此，
$(S/\mathfrak q)_{\mathfrak m}$ 是 Nagata 环 $R$ 的商的局部化，
所以是 Nagata 环（引理 \ref{algebra-lemma-quasi-finite-over-nagata}
和 \ref{algebra-lemma-nagata-localize}），从而解析非分歧
（引理 \ref{algebra-lemma-local-nagata-domain-analytically-unramified}）。
这说明 (3)(b) 成立，证明完毕。
\end{proof}

\begin{proposition}
\label{algebra-proposition-ubiquity-nagata}
以下各类环都是 Nagata 环，因而尤其是广泛日本型环：
\begin{enumerate}
\item 域；
\item Noether 完备局部环；
\item $\mathbf{Z}$；
\item 分式域具有特征零的 Dedekind 整环；
\item 上述任一种环的有限型环扩张。
\end{enumerate}
\end{proposition}

\begin{proof}
Noether 完备局部环的情形是引理
\ref{algebra-lemma-Noetherian-complete-local-Nagata}。在其他情形，
只需对环的每个素理想 $\mathfrak p$ 检验 $R/\mathfrak p$ 是否为 N-2。
只要 $R/\mathfrak p$ 是域，即 $\mathfrak p$ 是极大理想，这就是显然的。
因此，在 Dedekind 环的情形，只需检验 $\mathfrak p = (0)$。
但由于假设分式域具有特征零，引理
\ref{algebra-lemma-domain-char-zero-N-1-2} 给出结论。
\end{proof}

\begin{example}
\label{algebra-example-nonjapanese-dvr}
离散赋值环是 Nagata 环，当且仅当它是 N-2 的（因为模极大理想的商
是域，从而是 N-2 的）。例 \ref{algebra-example-bad-dvr-char-p} 中的离散赋值环
$A$ 不是 Nagata 环，即不是 N-2 的。事实上，有限扩张
$A \subset R = A[f]$ 不是 N-1 的。为此，写
$f = \sum a_i x^i$。对每个 $n \geq 1$，令
$g_n = \sum_{i < n} a_i x^i \in A$。则
$h_n = (f - g_n)/x^n$ 是 $R$ 的分式域中的元素，并且
$h_n^p \in k^p[[x]] \subset A$。所以 $R$ 的整闭包 $R'$ 包含
$h_1, h_2, h_3, \ldots$。现在，若 $R'$ 在 $R$ 上有限，从而在
$A$ 上有限，则对所有 $n$，$f  = x^n h_n + g_n$ 都属于子模
$A + x^nR'$。由 Artin--Rees，这将蕴含 $f \in A$
（引理 \ref{algebra-lemma-intersect-powers-ideal-module-zero}），矛盾。
\end{example}

\begin{lemma}
\label{algebra-lemma-nagata-pth-roots}
设 $(A, \mathfrak m)$ 是 Noether 的 Nagata 局部整环，其分式域
具有特征 $p$。若 $a \in A$ 在 $A^\wedge$ 中有 $p$ 次根，
则 $a$ 在 $A$ 中也有 $p$ 次根。
\end{lemma}

\begin{proof}
设 $\alpha \in A^\wedge$ 是 $a$ 的 $p$ 次根。反设 $a$ 在 $A$
中没有 $p$ 次根。则 $a$ 在 $A$ 的分式域 $K$ 中也没有 $p$ 次根。
确实，若 $\beta = b/c$，其中 $b, c \in A$ 且 $c \not = 0$，
满足 $\beta^p = a$，则在 $A^\wedge$ 中有 $\alpha = b/c$；
由 $A \to A^\wedge$ 的忠实平坦性，这蕴含 $c$ 在 $A$ 中整除 $b$，
因而 $\beta \in A$，矛盾。所以
$B = A[x]/(x^p - a)$ 是整环，因为 $K[x]/(x^p - a)$ 是域。
然而，由所假设的 $\alpha$ 的存在，$B$ 的完备化不是约化的。
这与前面的结果矛盾：$B$ 是 Nagata 环
（命题 \ref{algebra-proposition-nagata-universally-japanese}），
因而由引理 \ref{algebra-lemma-local-nagata-domain-analytically-unramified}
解析非分歧。
\end{proof}

\section{性质的上升}
\label{algebra-section-ascending-properties}

\noindent
本节开始证明一些关于环的性质“上升”的代数事实。为处理环的深度，
我们要用到下述关于模的深度上升的结果，见
\cite[IV, Proposition 6.3.1]{EGA}。

\begin{lemma}
\label{algebra-lemma-apply-grothendieck-module}
\begin{reference}
\cite[IV, Proposition 6.3.1]{EGA}
\end{reference}
有
$$
\text{depth}(M \otimes_R N)
=
\text{depth}(M) + \text{depth}(N/\mathfrak m_RN)
$$
其中 $R \to S$ 是 Noether 局部环的局部同态，$M$ 是有限
$R$-模，而 $N$ 是在 $R$ 上平坦的有限 $S$-模。
\end{lemma}

\begin{proof}
在命题及下面的证明中，我们把 $M$ 的深度取为其作为 $R$-模的深度，
把 $M \otimes_R N$ 的深度取为其作为 $S$-模的深度，并把
$N/\mathfrak m_RN$ 的深度取为其作为 $S/\mathfrak m_RS$-模的深度。
% P01-E1016
以 $n$ 表示右端。先假设 $n$ 为零。于是
$\text{depth}(M) = 0$ 且
$\text{depth}(N/\mathfrak m_RN) = 0$。
这意味着存在 $z \in M$，其零化理想是 $\mathfrak m_R$；还存在
$\overline{y} \in N/\mathfrak m_RN$，其零化理想是
$\mathfrak m_S/\mathfrak m_RS$。取 $\overline{y}$ 的提升 $y \in N$。
由于 $N$ 在 $R$ 上平坦，由该元素给出的映射
$z : R/\mathfrak m_R \to M$ 诱导单射
$N/\mathfrak m_RN \to M \otimes_R N$。因此，
$z \otimes y$ 的零化理想是 $\mathfrak m_S$。于是
$\text{depth}(M \otimes_R N) = 0$。

\medskip\noindent
假设 $n > 0$。若 $\text{depth}(N/\mathfrak m_RN) > 0$，则可选取
$f \in \mathfrak m_S$，其像
$\overline{f} \in S/\mathfrak m_RS$ 是 $N/\mathfrak m_RN$
上的非零因子。由引理 \ref{algebra-lemma-depth-drops-by-one}，
$$
\text{depth}(N/\mathfrak m_RN) =
\text{depth}(N/(f, \mathfrak m_R)N) + 1
$$
根据引理 \ref{algebra-lemma-mod-injective}，元素 $f \in S$ 是 $N$
上的非零因子，且 $N/fN$ 在 $R$ 上平坦。因此，对 $n$ 作归纳可得
$$
\text{depth}(M \otimes_R N/fN) =
\text{depth}(M) + \text{depth}(N/(f, \mathfrak m_R)N).
$$
由于 $N/fN$ 在 $R$ 上平坦，序列
$$
0 \to M \otimes_R N \to M \otimes_R N \to M \otimes_R N/fN \to 0
$$
正合，其中第一个映射是乘以 $f$（引理
\ref{algebra-lemma-flat-tor-zero}）。因此，由引理
\ref{algebra-lemma-depth-drops-by-one} 可得
$\text{depth}(M \otimes_R N) = \text{depth}(M \otimes_R N/fN) + 1$，
从而引理中的公式成立。

\medskip\noindent
若 $n > 0$，但 $\text{depth}(N/\mathfrak m_RN) = 0$，则可选取
$f \in \mathfrak m_R$，使其为 $M$ 上的非零因子。由于 $N$ 在 $R$
上平坦，$f$ 也是 $M \otimes_R N$ 上的非零因子。再次对 $n$ 作归纳，得
$$
\text{depth}(M/fM \otimes_R N) =
\text{depth}(M/fM) + \text{depth}(N/\mathfrak m_RN).
$$
在此情形下，由引理 \ref{algebra-lemma-depth-drops-by-one}，
$\text{depth}(M \otimes_R N) = \text{depth}(M/fM \otimes_R N) + 1$，
且 $\text{depth}(M) = \text{depth}(M/fM) + 1$，
从而引理中的公式成立。
\end{proof}

\begin{lemma}
\label{algebra-lemma-apply-grothendieck}
设 $R \to S$ 是 Noether 局部环之间的平坦局部环同态。则
$$
\text{depth}(S) = \text{depth}(R) + \text{depth}(S/\mathfrak m_RS).
$$
\end{lemma}

\begin{proof}
这是引理 \ref{algebra-lemma-apply-grothendieck-module} 的特例。
\end{proof}

\begin{lemma}
\label{algebra-lemma-CM-goes-up}
% P01-E1017
设 $R \to S$ 是 Noether 局部环之间的平坦局部同态。以下条件等价：
\begin{enumerate}
\item $S$ 是 Cohen--Macaulay 的；并且
\item $R$ 和 $S/\mathfrak m_RS$ 都是 Cohen--Macaulay 的。
\end{enumerate}
\end{lemma}

\begin{proof}
这由定义以及引理 \ref{algebra-lemma-apply-grothendieck} 和
\ref{algebra-lemma-dimension-base-fibre-equals-total} 得出。
\end{proof}

\begin{lemma}
\label{algebra-lemma-Sk-goes-up}
设 $\varphi : R \to S$ 是环映射。假设
\begin{enumerate}
\item $R$ 是 Noether 的；
\item $S$ 是 Noether 的；
\item $\varphi$ 平坦；
\item 纤维环 $S \otimes_R \kappa(\mathfrak p)$ 满足 $(S_k)$；并且
\item $R$ 具有性质 $(S_k)$。
\end{enumerate}
则 $S$ 具有性质 $(S_k)$。
\end{lemma}

\begin{proof}
设 $\mathfrak q$ 是 $S$ 的素理想，位于 $R$ 的素理想
$\mathfrak p$ 之上。由引理 \ref{algebra-lemma-apply-grothendieck}，有
$$
\text{depth}(S_{\mathfrak q}) =
\text{depth}(S_{\mathfrak q}/\mathfrak pS_{\mathfrak q}) +
\text{depth}(R_{\mathfrak p}).
$$
另一方面，由引理 \ref{algebra-lemma-dimension-base-fibre-total}，有
$$
\dim(R_{\mathfrak p})
+
\dim(S_{\mathfrak q}/\mathfrak pS_{\mathfrak q})
\geq
\dim(S_{\mathfrak q})
$$
（实际上，由引理 \ref{algebra-lemma-dimension-base-fibre-equals-total}，
这里等号成立，但严格地说我们并不需要这一点）。最后，由于假设映射的
纤维环满足 $(S_k)$，可知
$\text{depth}(S_{\mathfrak q}/\mathfrak pS_{\mathfrak q})
\geq \min(k, \dim(S_{\mathfrak q}/\mathfrak pS_{\mathfrak q}))$。
因此，引理由下列不等式链得出：
\begin{eqnarray*}
\text{depth}(S_{\mathfrak q}) & = &
\text{depth}(S_{\mathfrak q}/\mathfrak pS_{\mathfrak q}) +
\text{depth}(R_{\mathfrak p}) \\
& \geq &
\min(k, \dim(S_{\mathfrak q}/\mathfrak pS_{\mathfrak q})) +
\min(k, \dim(R_{\mathfrak p})) \\
& = &
\min(2k, \dim(S_{\mathfrak q}/\mathfrak pS_{\mathfrak q}) + k,
k + \dim(R_\mathfrak p),
\dim(S_{\mathfrak q}/\mathfrak pS_{\mathfrak q}) +
\dim(R_{\mathfrak p})) \\
& \geq &
\min(k, \dim(S_{\mathfrak q}))
\end{eqnarray*}
如所需。
\end{proof}

\begin{lemma}
\label{algebra-lemma-Rk-goes-up}
设 $\varphi : R \to S$ 是环映射。假设
\begin{enumerate}
\item $R$ 是 Noether 的；
% P01-E1018
\item $S$ 是 Noether 的；
\item $\varphi$ 平坦；
\item 纤维环 $S \otimes_R \kappa(\mathfrak p)$ 具有性质 $(R_k)$；并且
\item $R$ 具有性质 $(R_k)$。
\end{enumerate}
则 $S$ 具有性质 $(R_k)$。
\end{lemma}

\begin{proof}
设 $\mathfrak q$ 是 $S$ 的素理想，位于 $R$ 的素理想
$\mathfrak p$ 之上。假设 $\dim(S_{\mathfrak q}) \leq k$。
由于由引理 \ref{algebra-lemma-dimension-base-fibre-equals-total} 有
$\dim(S_{\mathfrak q}) = \dim(R_{\mathfrak p})
+ \dim(S_{\mathfrak q}/\mathfrak pS_{\mathfrak q})$，可知
$\dim(R_{\mathfrak p}) \leq k$ 且
$\dim(S_{\mathfrak q}/\mathfrak pS_{\mathfrak q}) \leq k$。
因此，由假设，$R_{\mathfrak p}$ 和
$S_{\mathfrak q}/\mathfrak pS_{\mathfrak q}$ 都正则。再由引理
\ref{algebra-lemma-flat-over-regular-with-regular-fibre}，
$S_{\mathfrak q}$ 正则。
\end{proof}

\begin{lemma}
\label{algebra-lemma-reduced-goes-up-noetherian}
设 $\varphi : R \to S$ 是环映射。假设
\begin{enumerate}
\item $R$ 是 Noether 的；
% P01-E1019
\item $S$ 是 Noether 的；
\item $\varphi$ 平坦；
\item 纤维环 $S \otimes_R \kappa(\mathfrak p)$ 约化；
\item $R$ 约化。
\end{enumerate}
则 $S$ 约化。
\end{lemma}

\begin{proof}
对 Noether 环而言，约化等价于具有性质 $(S_1)$ 和 $(R_0)$，
见引理 \ref{algebra-lemma-criterion-reduced}。因此，已知 $R$ 和纤维环具有这些性质。
于是可应用引理 \ref{algebra-lemma-Sk-goes-up} 和 \ref{algebra-lemma-Rk-goes-up}，
得到 $S$ 具有 $(S_1)$ 和 $(R_0)$；换言之，再由引理
\ref{algebra-lemma-criterion-reduced}，它约化。
\end{proof}

\begin{lemma}
\label{algebra-lemma-reduced-goes-up}
设 $\varphi : R \to S$ 是环映射。假设
\begin{enumerate}
\item $\varphi$ 光滑；
\item $R$ 约化。
\end{enumerate}
则 $S$ 约化。
\end{lemma}

\begin{proof}
注意，$R \to S$ 平坦且纤维正则（见节
\ref{algebra-section-smooth-overview} 中关于光滑环映射的结果列表）。
特别地，纤维约化。因此，若 $R$ 是 Noether 的，则 $S$ 是 Noether 的，
结论由引理 \ref{algebra-lemma-reduced-goes-up-noetherian} 得出。

\medskip\noindent
一般情形下，可找到有限生成的 $\mathbf{Z}$-子代数 $R_0 \subset R$
和光滑环映射 $R_0 \to S_0$，使得
$S \cong R \otimes_{R_0} S_0$，见节
\ref{algebra-section-smooth-overview} 中的评注 (10)。现在，若 $x \in S$
满足 $x^2 = 0$，则可扩大 $R_0$，并假设 $x$ 来自元素
$x_0 \in S_0$。再次扩大 $R_0$ 后，可假设 $x_0^2 = 0$ 于 $S_0$ 中。
% P01-E1020
然而，由于 $R_0 \subset R$ 且后者约化，前者也约化，故可知
$S_0$ 约化，因而如所需 $x_0 = 0$。
\end{proof}

\begin{lemma}
\label{algebra-lemma-normal-goes-up-noetherian}
设 $\varphi : R \to S$ 是环映射。假设
\begin{enumerate}
\item $R$ 是 Noether 的；
\item $S$ 是 Noether 的；
\item $\varphi$ 平坦；
\item 纤维环 $S \otimes_R \kappa(\mathfrak p)$ 正规；并且
\item $R$ 正规。
\end{enumerate}
则 $S$ 正规。
\end{lemma}

\begin{proof}
对 Noether 环而言，正规等价于具有性质 $(S_2)$ 和 $(R_1)$，
见引理 \ref{algebra-lemma-criterion-normal}。因此，已知 $R$ 和纤维环具有这些性质。
于是可应用引理 \ref{algebra-lemma-Sk-goes-up} 和 \ref{algebra-lemma-Rk-goes-up}，
得到 $S$ 具有 $(S_2)$ 和 $(R_1)$；换言之，再由引理
\ref{algebra-lemma-criterion-normal}，它正规。
\end{proof}

\begin{lemma}
\label{algebra-lemma-normal-goes-up}
设 $\varphi : R \to S$ 是环映射。假设
\begin{enumerate}
\item $\varphi$ 光滑；
\item $R$ 正规。
\end{enumerate}
则 $S$ 正规。
\end{lemma}

\begin{proof}
注意，$R \to S$ 平坦且纤维正则（见节
\ref{algebra-section-smooth-overview} 中关于光滑环映射的结果列表）。
特别地，纤维正规。因此，若 $R$ 是 Noether 的，则 $S$ 是 Noether 的，
结论由引理 \ref{algebra-lemma-normal-goes-up-noetherian} 得出。

\medskip\noindent
现考虑一般情形。先注意 $R$ 约化，因而由引理
\ref{algebra-lemma-reduced-goes-up}，$S$ 约化。设 $\mathfrak q$ 是 $S$
的素理想，并设 $\mathfrak p$ 是 $R$ 的相应素理想。注意，
$R_{\mathfrak p}$ 是正规整环。必须证明 $S_{\mathfrak q}$ 是正规整环。
为此，可以用 $R_{\mathfrak p}$ 代替 $R$，并用 $S_{\mathfrak p}$
代替 $S$。因此，可以假设 $R$ 是正规整环。

\medskip\noindent
假设 $R \to S$ 光滑，且 $R$ 是正规整环。可找到有限生成的
$\mathbf{Z}$-子代数 $R_0 \subset R$ 和光滑环映射
$R_0 \to S_0$，使得 $S \cong R \otimes_{R_0} S_0$，见节
\ref{algebra-section-smooth-overview} 中的评注 (10)。由于 $R_0$ 是 Nagata 整环
（见命题 \ref{algebra-proposition-ubiquity-nagata}），可知其整闭包 $R_0'$
在 $R_0$ 上有限。此外，由于 $R$ 是正规整环，显然有
$R_0' \subset R$。因此，可以用 $R_0'$ 代替 $R_0$，并用
$R_0' \otimes_{R_0} S_0$ 代替 $S_0$，从而假设 $R_0$ 是 Noether
正规整环。由证明的第一段可得，$S_0$ 是正规环（当然，它不一定是整环）。
如此可见，$R = \bigcup R_\lambda$ 是 Noether 正规整环的并，
相应地，$S = \colim R_\lambda \otimes_{R_0} S_0$ 是正规环的余极限。
这蕴含 $S$ 是正规环。略去一些细节。
\end{proof}

\begin{lemma}
\label{algebra-lemma-regular-goes-up}
\begin{slogan}
正则性沿光滑环映射上升。
\end{slogan}
设 $\varphi : R \to S$ 是环映射。假设
\begin{enumerate}
\item $\varphi$ 光滑；
\item $R$ 是正则环。
\end{enumerate}
则 $S$ 正则。
\end{lemma}

\begin{proof}
% P01-E1021
对所有性质 $(R_k)$ 应用引理 \ref{algebra-lemma-Rk-goes-up}，并用引理
\ref{algebra-lemma-characterize-smooth-over-field} 验证假设，即得结论。
\end{proof}

\section{性质的下降}
\label{algebra-section-descending-properties}

\noindent
本节开始证明一些关于环的性质“下降”的代数事实。事实表明，性质的下降
往往比其上升“更容易”。换言之，
% P01-E1022
对环映射 $R \to S$ 的假设常常弱于上一节相应引理中的假设。
不过，我们提醒读者：除非也能证明相应的上升结果，否则下降结果往往毫无用处！
下面是说明这种现象的典型结果。

\begin{lemma}
\label{algebra-lemma-descent-Noetherian}
设 $R \to S$ 是环映射。假设
\begin{enumerate}
\item $R \to S$ 忠实平坦；并且
\item $S$ 是 Noether 的。
\end{enumerate}
则 $R$ 是 Noether 的。
\end{lemma}

\begin{proof}
设 $I_0 \subset I_1 \subset I_2 \subset \ldots$ 是 $R$ 的递增理想列。
由假设，对某个 $n$ 有
$I_nS = I_{n+1}S = I_{n+2}S = \ldots$。由忠实平坦性，扩张后再收缩
得到原理想，即对 $R$ 的每个理想 $I$ 都有
$I = R \cap IS$（引理
\ref{algebra-lemma-faithfully-flat-universally-injective}）。因此，如所需，
$I_n = I_{n+1} = I_{n+2} = \ldots$。
\end{proof}

\begin{lemma}
\label{algebra-lemma-descent-reduced}
设 $R \to S$ 是环映射。假设
\begin{enumerate}
\item $R \to S$ 忠实平坦；并且
\item $S$ 约化。
\end{enumerate}
则 $R$ 约化。
\end{lemma}

\begin{proof}
由引理 \ref{algebra-lemma-faithfully-flat-universally-injective}，
$R \to S$ 是单射，故结论显然。
\end{proof}

\begin{lemma}
\label{algebra-lemma-descent-normal}
设 $R \to S$ 是环映射。假设
\begin{enumerate}
\item $R \to S$ 忠实平坦；并且
\item $S$ 是正规环。
\end{enumerate}
则 $R$ 是正规环。
\end{lemma}

\begin{proof}
由于 $S$ 约化，可知 $R$ 约化。设 $\mathfrak p$ 是 $R$ 的素理想。
必须证明 $R_{\mathfrak p}$ 是正规整环。
% P01-E1023
由于 $S_{\mathfrak p}$ 在 $R_{\mathfrak p}$ 上也忠实平坦，
可以假设 $R$ 是极大理想为 $\mathfrak m$ 的局部环。设
$\mathfrak q$ 是 $S$ 的位于 $\mathfrak m$ 之上的素理想。于是
$R \to S_{\mathfrak q}$ 忠实平坦（引理
\ref{algebra-lemma-local-flat-ff}）。因此，可以进一步假设 $S$ 也是局部环。
特别地，$S$ 是正规整环。由于 $R \to S$ 忠实平坦且 $S$
是正规整环，可知 $R$ 是整环。接着，设 $a/b$ 整于 $R$，其中
$a, b \in R$。由于 $S$ 正规，有 $a/b \in S$，因而 $a \in bS$。
这意味着 $a : R \to R/bR$ 经向 $S$ 基变换后成为零映射。
由忠实平坦性，$a \in bR$，所以 $a/b \in R$。因此 $R$ 正规。
\end{proof}

\begin{lemma}
\label{algebra-lemma-descent-regular}
设 $R \to S$ 是环映射。假设
\begin{enumerate}
\item $R \to S$ 忠实平坦；并且
\item $S$ 是正则环。
\end{enumerate}
则 $R$ 是正则环。
\end{lemma}

\begin{proof}
由引理 \ref{algebra-lemma-descent-Noetherian}，$R$ 是 Noether 的。
设 $\mathfrak p \subset R$ 是素理想。选取位于 $\mathfrak p$ 之上的素理想
$\mathfrak q \subset S$。于是，引理 \ref{algebra-lemma-flat-under-regular}
可应用于 $R_\mathfrak p \to S_\mathfrak q$，从而
$R_\mathfrak p$ 正则。由于 $\mathfrak p$ 任意，可知 $R$ 正则。
\end{proof}

\begin{lemma}
\label{algebra-lemma-descent-Sk}
设 $R \to S$ 是环映射。假设
\begin{enumerate}
\item $R \to S$ 忠实平坦；并且
\item $S$ 是 Noether 的且具有性质 $(S_k)$。
\end{enumerate}
则 $R$ 是 Noether 的且具有性质 $(S_k)$。
\end{lemma}

\begin{proof}
我们已经看到，(1) 和 (2) 蕴含 $R$ 是 Noether 的，见引理
\ref{algebra-lemma-descent-Noetherian}。设 $\mathfrak p \subset R$ 是素理想。
选取位于 $\mathfrak p$ 之上的素理想 $\mathfrak q \subset S$，
使其对应于纤维环 $S \otimes_R \kappa(\mathfrak p)$ 的一个极小素理想。
于是，$A = R_{\mathfrak p} \to S_{\mathfrak q} = B$ 是 Noether
局部环之间的平坦局部环同态，且 $\mathfrak m_AB$ 是 $B$ 的定义理想。
因此，$\dim(A) = \dim(B)$（引理
\ref{algebra-lemma-dimension-base-fibre-equals-total}），且
$\text{depth}(A) = \text{depth}(B)$（引理
\ref{algebra-lemma-apply-grothendieck}）。所以，由 $B$ 具有 $(S_k)$，
可知 $A$ 具有 $(S_k)$。
\end{proof}

\begin{lemma}
\label{algebra-lemma-descent-Rk}
设 $R \to S$ 是环映射。假设
\begin{enumerate}
\item $R \to S$ 忠实平坦；并且
\item $S$ 是 Noether 的且具有性质 $(R_k)$。
\end{enumerate}
则 $R$ 是 Noether 的且具有性质 $(R_k)$。
\end{lemma}

\begin{proof}
我们已经看到，(1) 和 (2) 蕴含 $R$ 是 Noether 的，见引理
\ref{algebra-lemma-descent-Noetherian}。设 $\mathfrak p \subset R$ 是素理想，
并假设 $\dim(R_{\mathfrak p}) \leq k$。选取位于 $\mathfrak p$ 之上的
素理想 $\mathfrak q \subset S$，使其对应于纤维环
$S \otimes_R \kappa(\mathfrak p)$ 的一个极小素理想。于是，
$A = R_{\mathfrak p} \to S_{\mathfrak q} = B$ 是 Noether 局部环之间的
平坦局部环同态，且 $\mathfrak m_AB$ 是 $B$ 的定义理想。因此，
$\dim(A) = \dim(B)$（引理
\ref{algebra-lemma-dimension-base-fibre-equals-total}）。由于 $S$ 具有 $(R_k)$，
可知 $B$ 是正则局部环。由引理 \ref{algebra-lemma-flat-under-regular}，
$A$ 正则。
\end{proof}

\begin{lemma}
\label{algebra-lemma-descent-nagata}
设 $R \to S$ 是环映射。假设
\begin{enumerate}
\item $R \to S$ 光滑且在谱上满；并且
\item $S$ 是 Nagata 环。
\end{enumerate}
则 $R$ 是 Nagata 环。
\end{lemma}

\begin{proof}
回忆，Nagata 环等价于 Noether 的广泛日本型环
（命题 \ref{algebra-proposition-nagata-universally-japanese}）。我们已经在引理
\ref{algebra-lemma-descent-Noetherian} 中看到 $R$ 是 Noether 的。设
$R \to A$ 是到整环的有限型环映射。根据引理
\ref{algebra-lemma-check-universally-japanese}，只需检验 $A$ 是 N-1 的。
显然，$B = A \otimes_R S$ 是有限型 $S$-代数，因而是 Nagata 环
（命题 \ref{algebra-proposition-nagata-universally-japanese}）。由于
$A \to B$ 光滑（引理 \ref{algebra-lemma-base-change-smooth}），可知
$B$ 约化（引理 \ref{algebra-lemma-reduced-goes-up}）。由于 $B$ 是 Noether 的，
它只有有限多个极小素理想 $\mathfrak q_1, \ldots, \mathfrak q_t$
（见引理 \ref{algebra-lemma-Noetherian-irreducible-components}）。由于
$A \to B$ 平坦，它们都位于 $(0) \subset A$ 之上
（由下降定理，见引理 \ref{algebra-lemma-flat-going-down}）。
% P01-E1024
全分式环 $Q(B)$ 是 $L_i = \kappa(\mathfrak q_i)$ 的乘积
（引理 \ref{algebra-lemma-total-ring-fractions-no-embedded-points} 和
\ref{algebra-lemma-minimal-prime-reduced-ring}）。此外，$B$ 在 $Q(B)$ 中的整闭包
$B'$ 是各个 $B/\mathfrak q_i$ 在相应因子 $L_i$ 中的整闭包
$B_i'$ 的乘积（与引理
\ref{algebra-lemma-characterize-reduced-ring-normal} 比较）。由于 $B$
是广泛日本型环，环扩张 $B/\mathfrak q_i \subset B_i'$ 有限，
从而 $B' = \prod B_i'$ 在 $B$ 上有限。由于 $A \to B$ 平坦，
$A$ 上的任意非零因子都映为 $B$ 上的非零因子。相应的映射
$$
Q(A) \otimes_A B = (A \setminus \{0\})^{-1}A \otimes_A B
= (A \setminus \{0\})^{-1}B \to Q(B)
$$
是单射（这里使用了引理 \ref{algebra-lemma-tensor-localization}）。
% P01-E1025
经此映射，$A'$ 映入 $B'$。这诱导映射
$$
A' \otimes_A B \longrightarrow B'
$$
它是单射（由上文以及 $A \to B$ 的平坦性）。由于 $B'$ 是有限
$B$-模且 $B$ 是 Noether 的，可知 $A' \otimes_A B$ 是有限 $B$-模。
因此，存在有限多个元素 $x_i \in A'$，使得 $x_i \otimes 1$
作为 $B$-模生成 $A' \otimes_A B$。最后，由 $A \to B$ 的忠实平坦性，
% P01-E1026
这些 $x_i$ 也作为 $A$-模生成 $A'$，结论得证。
\end{proof}

\begin{remark}
\label{algebra-remark-universally-catenary-does-not-descend}
“普遍链状”这一性质不下降，即使沿平展环映射也不下降。在
《例》的节 \ref{examples-section-non-catenary-Noetherian-local}
中，构造了有限环映射 $A \to B$，其中 $A$ 是非普遍链状的
Noether 局部环，而 $B$ 是具有两个极大理想
$\mathfrak m$, $\mathfrak n$ 的半局部环；$B_{\mathfrak m}$ 和
$B_{\mathfrak n}$ 分别是维数 $2$ 和 $1$ 的正则环，
% P01-E1027
并且其剩余域与 $A$ 的剩余域相同。此外，$\mathfrak m_A$ 在
$B_{\mathfrak m}$ 和 $B_{\mathfrak n}$ 中都生成极大理想
（所以 $A \to B$ 既非分歧又有限）。由引理
\ref{algebra-lemma-etale-makes-unramified-closed}，存在局部平展环映射
$A \to A'$，使得 $B \otimes_A A' = B_1 \times B_2$ 分解，
且 $A' \to B_i$ 满。因此，$A'$ 具有两个极小素理想
$\mathfrak q_i$，且 $A'/\mathfrak q_i \cong B_i$。由于 $B_i$
是正则局部环（因为它平展于 $B_{\mathfrak m}$ 或
$B_{\mathfrak n}$ 之一），可得 $A'$ 普遍链状。
\end{remark}

\section{几何正规代数}
\label{algebra-section-geometrically-normal}

\noindent
本节给出环的性质上升和下降的一些应用。

\begin{lemma}
\label{algebra-lemma-geometrically-normal}
设 $k$ 是域。设 $A$ 是 $k$-代数。$A$ 的以下性质等价：
\begin{enumerate}
\item 对每个域扩张 $k'/k$，$k' \otimes_k A$ 都是正规环；
\item 对每个有限生成域扩张 $k'/k$，$k' \otimes_k A$ 都是正规环；
\item 对每个有限纯不可分扩张 $k'/k$，$k' \otimes_k A$ 都是正规环；
\item $k^{perf} \otimes_k A$ 是正规环。
\end{enumerate}
这里的正规环定义见定义 \ref{algebra-definition-ring-normal}。
\end{lemma}

\begin{proof}
显然，(1) $\Rightarrow$ (2) $\Rightarrow$ (3)，且
(1) $\Rightarrow$ (4)。

\medskip\noindent
若 $k'/k$ 是有限纯不可分扩张，则存在 $k$-扩张的嵌入
$k' \to k^{perf}$。环映射
$k' \otimes_k A \to k^{perf} \otimes_k A$ 忠实平坦，
所以若 $k^{perf} \otimes_k A$ 正规，则由引理
\ref{algebra-lemma-descent-normal}，$k' \otimes_k A$ 正规。
由此可见，(4) $\Rightarrow$ (3)。

\medskip\noindent
假设 (2)，并设 $k'/k$ 是任意域扩张。可将
$k' = \colim_i k_i$ 写成有限生成域扩张的有向余极限。因此，
$k' \otimes_k A = \colim_i k_i \otimes_k A$ 是正规环的有向余极限。
由引理 \ref{algebra-lemma-colimit-normal-ring}，$k' \otimes_k A$ 是正规环。
所以 (1) 成立。

\medskip\noindent
假设 (3)，并设 $K/k$ 是有限生成域扩张。由引理
\ref{algebra-lemma-make-separable}，可找到图
$$
\xymatrix{
K \ar[r] & K' \\
k \ar[u] \ar[r] & k' \ar[u]
}
$$
其中 $k'/k$、$K'/K$ 是有限纯不可分域扩张，并且 $K'/k'$ 可分。
由引理 \ref{algebra-lemma-localization-smooth-separable}，存在光滑
$k'$-代数 $B$，使 $K'$ 是 $B$ 的分式域。现在可如下论证：
第 1 步：$k' \otimes_k A$ 是正规环，因为我们假设了 (3)。
第 2 步：$B \otimes_{k'} k' \otimes_k A$ 是正规环，因为
$k' \otimes_k A \to B \otimes_{k'} k' \otimes_k A$ 光滑
（引理 \ref{algebra-lemma-base-change-smooth}），并且正规性沿光滑映射上升
（引理 \ref{algebra-lemma-normal-goes-up}）。
第 3 步。$K' \otimes_{k'} k' \otimes_k A = K' \otimes_k A$
是正规环，因为它是正规环的局部化
（引理 \ref{algebra-lemma-localization-normal-ring}）。
第 4 步。最后，正规性沿忠实平坦环映射
$K \otimes_k A \to K' \otimes_k A$ 下降，所以由引理
\ref{algebra-lemma-descent-normal}，$K \otimes_k A$ 是正规环。
这就证明了引理。
\end{proof}

\begin{definition}
\label{algebra-definition-geometrically-normal}
设 $k$ 是域。若引理 \ref{algebra-lemma-geometrically-normal} 的等价条件成立，
则称 $k$-代数 $R$ 在 $k$ 上{it 几何正规}。
\end{definition}

\begin{lemma}
\label{algebra-lemma-localization-geometrically-normal-algebra}
\begin{slogan}
局部化保持几何正规性。
\end{slogan}
设 $k$ 是域。几何正规 $k$-代数的局部化仍几何正规。
\end{lemma}

\begin{proof}
这是显然的，因为一个环正规与否可在各个素理想处的局部化上检验。
\end{proof}

\begin{lemma}
\label{algebra-lemma-separable-field-extension-geometrically-normal}
设 $k$ 是域。设 $K/k$ 是可分域扩张。则 $K$ 在 $k$ 上几何正规。
\end{lemma}

\begin{proof}
这是因为 $k^{perf} \otimes_k K$ 是域。事实上，由引理
\ref{algebra-lemma-separable-extension-preserves-reducedness}，它约化。
由引理 \ref{algebra-lemma-perfection}（或定义 \ref{algebra-definition-perfection}），
域扩张 $k^{perf}/k$ 纯不可分。因此，由引理
\ref{algebra-lemma-radicial-integral-bijective}，环
$k^{perf} \otimes_k K$ 具有唯一素理想。具有唯一素理想的约化环是域。
\end{proof}

\begin{lemma}
\label{algebra-lemma-geometrically-normal-tensor-normal}
设 $k$ 是域。设 $A, B$ 是 $k$-代数。假设 $A$ 在 $k$ 上几何正规，
且 $B$ 是正规环。则 $A \otimes_k B$ 是正规环。
\end{lemma}

\begin{proof}
% P01-E1028
设 $\mathfrak r$ 是 $A \otimes_k B$ 的素理想。记
$\mathfrak p$、$\mathfrak q$ 分别为 $A$、$B$ 中的相应素理想。则
$(A \otimes_k B)_{\mathfrak r}$ 是
$A_{\mathfrak p} \otimes_k B_{\mathfrak q}$ 的局部化。因此，由引理
\ref{algebra-lemma-localization-normal-ring} 和
\ref{algebra-lemma-localization-geometrically-normal-algebra}，
只需对环 $A_{\mathfrak p} \otimes_k B_{\mathfrak q}$ 证明结论。
所以可以假设 $A$ 和 $B$ 都是整环。

\medskip\noindent
% P01-E1029
假设 $A$ 和 $B$ 是分式域分别为 $K$ 和 $L$ 的整环。注意，
% P01-E1030
$B$ 是其有限型正规 $k$-子代数的滤过余极限（因为 $k$ 是 Nagata 环，
见命题 \ref{algebra-proposition-ubiquity-nagata}；从而，由命题
\ref{algebra-proposition-nagata-universally-japanese}，有限型
$k$-子代数的整闭包仍是有限型 $k$-子代数）。由引理
\ref{algebra-lemma-colimit-normal-ring}，可化归到 $B$ 在 $k$ 上有限型的情形。

\medskip\noindent
% P01-E1029 P01-E1031
假设 $A$ 和 $B$ 是分式域分别为 $K$ 和 $L$ 的整环，且 $B$
在 $k$ 上有限型。在此情形下，环 $K \otimes_k B$ 在 $K$ 上有限型，
因而是 Noether 的（引理 \ref{algebra-lemma-Noetherian-permanence}）。特别地，
$K \otimes_k B$ 只有有限多个极小素理想
（引理 \ref{algebra-lemma-Noetherian-irreducible-components}）。由于
$A \to A \otimes_k B$ 平坦，这蕴含 $A \otimes_k B$ 只有有限多个
极小素理想（由平坦环映射的下降定理——引理
\ref{algebra-lemma-flat-going-down}——这些素理想都位于 $(0) \subset A$ 之上）。
因此，只需证明 $A \otimes_k B$ 在其全分式环中整闭
（引理 \ref{algebra-lemma-characterize-reduced-ring-normal}）。

\medskip\noindent
我们断言，$K \otimes_k B$ 和 $A \otimes_k L$ 都是正规环。
若此断言成立，则 $Q(A \otimes_k B)$ 中任一整于
$A \otimes_k B$ 的元素 $x$ 都属于
$K \otimes_k B \cap A \otimes_k L = A \otimes_k B$
（由引理 \ref{algebra-lemma-normal-ring-integrally-closed}），结论得证。
% P01-E1032
由于由假设 $A \otimes_K L$ 是正规环，只需证明
$K \otimes_k B$ 正规。

\medskip\noindent
由于 $A$ 在 $k$ 上几何正规，可知 $K$ 在 $k$ 上几何正规
（引理 \ref{algebra-lemma-localization-geometrically-normal-algebra}），
因而 $K$ 在 $k$ 上几何约化。所以，$K = \bigcup K_i$ 是
$k$ 的几何约化有限生成域扩张的并
（引理 \ref{algebra-lemma-subalgebra-separable}）。每个 $K_i$ 都是某个光滑
$k$-代数的局部化（引理
\ref{algebra-lemma-localization-smooth-separable}）。因此，
$K_i \otimes_k B$ 是某个光滑 $B$-代数的局部化，故正规
（引理 \ref{algebra-lemma-normal-goes-up}）。于是，
$K \otimes_k B$ 是正规环（引理
\ref{algebra-lemma-colimit-normal-ring}），结论得证。
\end{proof}

\begin{lemma}
\label{algebra-lemma-geometrically-normal-over-separable-algebraic}
设 $k'/k$ 是可分代数域扩张。设 $A$ 是 $k'$ 上的代数。
则 $A$ 在 $k$ 上几何正规，当且仅当它在 $k'$ 上几何正规。
\end{lemma}

\begin{proof}
设 $L/k$ 是有限纯不可分域扩张。则
$L' = k' \otimes_k L$ 是域（见《域》的节
\ref{fields-section-algebraic} 中的材料），且
$A \otimes_k L = A \otimes_{k'} L'$。因此，若 $A$ 在 $k'$
上几何正规，则 $A$ 在 $k$ 上几何正规。

\medskip\noindent
假设 $A$ 在 $k$ 上几何正规。设 $K/k'$ 是域扩张。则
$$
K \otimes_{k'} A = (K \otimes_k A) \otimes_{(k' \otimes_k k')} k'
$$
由引理 \ref{algebra-lemma-separable-algebraic-diagonal}，
$k' \otimes_k k' \to k'$ 是局部化，所以
$K \otimes_{k'} A$ 是正规环的局部化，因而正规。
\end{proof}

\section{几何正则代数}
\label{algebra-section-geometrically-regular}

\noindent
设 $k$ 是域。
设 $A$ 是 Noether 的 $k$-代数。
设 $K/k$ 是有限生成域扩张。
则环 $K \otimes_k A$ 也是 Noether 的，见
引理 \ref{algebra-lemma-Noetherian-field-extension}。
因此，以下引理是有意义的。

\begin{lemma}
\label{algebra-lemma-geometrically-regular}
设 $k$ 是域。设 $A$ 是 $k$-代数。
假设 $A$ 是 Noether 的。
$A$ 的以下性质等价：
\begin{enumerate}
\item 对每个有限生成域扩张 $k'/k$，$k' \otimes_k A$ 都是正则环；
\item 对每个有限纯不可分域扩张 $k'/k$，$k' \otimes_k A$ 都是正则环。
\end{enumerate}
此处正则环的含义如定义 \ref{algebra-definition-regular} 所述。
\end{lemma}

\begin{proof}
由引理前面的说明，此引理是有意义的。
显然 (1) $\Rightarrow$ (2)。

\medskip\noindent
假设 (2)，并设 $K/k$ 是有限生成域扩张。
由引理 \ref{algebra-lemma-make-separable}，可以找到交换图
$$
\xymatrix{
K \ar[r] & K' \\
k \ar[u] \ar[r] & k' \ar[u]
}
$$
其中 $k'/k$、$K'/K$ 是有限纯不可分域扩张，且 $K'/k'$ 可分。由
引理 \ref{algebra-lemma-localization-smooth-separable}，
存在光滑 $k'$-代数 $B$，使 $K'$ 是 $B$ 的分式域。现在论证如下：
步骤 1：$k' \otimes_k A$ 是正则环，因为我们假设了 (2)。
步骤 2：$B \otimes_{k'} k' \otimes_k A$ 是正则环，这是因为
$k' \otimes_k A \to B \otimes_{k'} k' \otimes_k A$ 光滑
（引理 \ref{algebra-lemma-base-change-smooth}），并且正则性沿光滑映射上升
（引理 \ref{algebra-lemma-regular-goes-up}）。
步骤 3：$K' \otimes_{k'} k' \otimes_k A = K' \otimes_k A$ 是正则环，
因为它是正则环的局部化（直接由定义可得）。
步骤 4：最后，由正则性沿忠实平坦环映射
$K \otimes_k A \to K' \otimes_k A$ 的下降
（引理 \ref{algebra-lemma-descent-regular}），$K \otimes_k A$ 是正则环。
这就证明了引理。
\end{proof}

\begin{definition}
\label{algebra-definition-geometrically-regular}
设 $k$ 是域。设 $R$ 是 Noether 的 $k$-代数。
若引理 \ref{algebra-lemma-geometrically-regular} 的等价条件成立，则称
$k$-代数 $R$ 在 $k$ 上\textit{几何正则}。
\end{definition}

\noindent
由定义显然可知，对于 $k$ 的任意有限生成域扩张 $K$，
$K \otimes_k R$ 是在 $K$ 上几何正则的代数。
我们以后将会看到（《代数进阶》，命题
\ref{more-algebra-proposition-characterization-geometrically-regular}），
只需在 $k \subset k' \subset k^{1/p}$（有限）时检验
$R \otimes_k k'$ 正则即可。

\begin{lemma}
\label{algebra-lemma-geometrically-regular-descent}
\begin{slogan}
几何正则性沿代数的忠实平坦映射下降
\end{slogan}
设 $k$ 是域。设 $A \to B$ 是忠实平坦的 $k$-代数映射。
若 $B$ 在 $k$ 上几何正则，则 $A$ 也如此。
\end{lemma}

\begin{proof}
假设 $B$ 在 $k$ 上几何正则。
设 $k'/k$ 是有限纯不可分扩张。
由于 $A \to B$ 的基变换（由
引理 \ref{algebra-lemma-surjective-spec-radical-ideal} 和
\ref{algebra-lemma-flat-base-change}），
$A \otimes_k k' \to B \otimes_k k'$ 忠实平坦；而由我们对 $B$
在 $k$ 上的假设，$B \otimes_k k'$ 正则。于是由引理
\ref{algebra-lemma-descent-regular}，$A \otimes_k k'$ 正则。
\end{proof}

\begin{lemma}
\label{algebra-lemma-geometrically-regular-goes-up}
设 $k$ 是域。设 $A \to B$ 是 $k$-代数间的光滑环映射。
若 $A$ 在 $k$ 上几何正则，则 $B$ 在 $k$ 上几何正则。
\end{lemma}

\begin{proof}
设 $k'/k$ 是有限生成域扩张。
则 $A \otimes_k k' \to B \otimes_k k'$ 是光滑环映射
（引理 \ref{algebra-lemma-base-change-smooth}），且 $A \otimes_k k'$ 正则。
故由引理 \ref{algebra-lemma-regular-goes-up}，$B \otimes_k k'$ 正则。
\end{proof}

\begin{lemma}
\label{algebra-lemma-geometrically-regular-over-subfields}
设 $k$ 是域。设 $A$ 是 $k$ 上的代数。
设 $k = \colim k_i$ 是由子域组成的有向余极限。
若 $A$ 在每个 $k_i$ 上都几何正则，则 $A$ 在 $k$ 上几何正则。
\end{lemma}

\begin{proof}
设 $k'/k$ 是有限纯不可分域扩张。
可以通过向 $k$ 添入有限多个变量并施加有限多个多项式关系得到 $k'$。
因而可见，存在 $i$ 和有限纯不可分域扩张 $k_i'/k_i$，使得
% P01-E1033
$k_i = k \otimes_{k_i} k_i'$。
于是 $A \otimes_k k' = A \otimes_{k_i} k_i'$，引理即得证。
\end{proof}

\begin{lemma}
\label{algebra-lemma-geometrically-regular-over-separable-algebraic}
设 $k'/k$ 是可分代数域扩张。
设 $A$ 是 $k'$ 上的代数。则 $A$ 在 $k$ 上几何正则，当且仅当
它在 $k'$ 上几何正则。
\end{lemma}

\begin{proof}
设 $L/k$ 是有限纯不可分域扩张。
则 $L' = k' \otimes_k L$ 是域（见《域》的节
\ref{fields-section-algebraic} 中的材料），且
$A \otimes_k L = A \otimes_{k'} L'$。因此，若 $A$ 在 $k'$
上几何正则，则 $A$ 在 $k$ 上几何正则。

\medskip\noindent
假设 $A$ 在 $k$ 上几何正则。由于 $k'$ 是 $k$ 的有限扩张的滤过余极限，
由引理 \ref{algebra-lemma-geometrically-regular-over-subfields}，可以假设
$k'/k$ 是有限可分的。考虑环映射
$$
k' \to A \otimes_k k' \to A.
$$
注意，作为 $A$ 到 $k'$ 的基变换，$A \otimes_k k'$ 在 $k'$ 上几何正则。
注意，$A \otimes_k k' \to A$ 是 $k' \otimes_k k' \to k'$
沿映射 $k' \to A$ 的基变换。由于 $k'/k$ 是环的平展扩张，可知
$k' \otimes_k k' \to k'$ 是平展的
（引理 \ref{algebra-lemma-etale}）。故由引理
\ref{algebra-lemma-geometrically-regular-goes-up}，$A$ 在 $k'$ 上几何正则。
\end{proof}

\section{几何 Cohen--Macaulay 代数}
\label{algebra-section-geometrically-CM}

\noindent
本节的名称略有不当，因为 Cohen--Macaulay 代数自动是几何
Cohen--Macaulay 的。具体见下面的
引理 \ref{algebra-lemma-extend-field-CM-locus}
和
引理 \ref{algebra-lemma-CM-geometrically-CM}。

\begin{lemma}
\label{algebra-lemma-tensor-fields-CM}
设 $k$ 是域，$K/k$ 和 $L/k$ 是两个域扩张，并且其中一个是有限型域扩张。
则 $K \otimes_k L$ 是 Noether 的 Cohen--Macaulay 环。
\end{lemma}

\begin{proof}
由引理 \ref{algebra-lemma-Noetherian-field-extension}，
环 $K \otimes_k L$ 是 Noether 的。
不妨设 $K$ 是纯超越扩张 $k(t_1, \ldots, t_r)$ 的有限扩张。则
$k(t_1, \ldots, t_r) \otimes_k L \to K \otimes_k L$
是有限自由环映射。由
引理 \ref{algebra-lemma-finite-flat-over-regular-CM}，
只需证明 $k(t_1, \ldots, t_r) \otimes_k L$ 是 Cohen--Macaulay 的。
这是显然的，因为它是多项式环 $L[t_1, \ldots, t_r]$ 的局部化。
（关于多项式环是 Cohen--Macaulay 的事实，例如见
引理 \ref{algebra-lemma-CM-polynomial-algebra}。）
\end{proof}

\begin{lemma}
\label{algebra-lemma-CM-geometrically-CM}
设 $k$ 是域。设 $S$ 是 Noether 的 $k$-代数。
设 $K/k$ 是有限生成域扩张，并置 $S_K = K \otimes_k S$。
设 $\mathfrak q \subset S$ 是 $S$ 的素理想。设
$\mathfrak q_K \subset S_K$ 是 $S_K$ 中位于 $\mathfrak q$ 之上的素理想。
则 $S_{\mathfrak q}$ 是 Cohen--Macaulay 的，当且仅当
$(S_K)_{\mathfrak q_K}$ 是 Cohen--Macaulay 的。
\end{lemma}

\begin{proof}
由
引理 \ref{algebra-lemma-Noetherian-field-extension}，
环 $S_K$ 是 Noether 的。因此
$S_{\mathfrak q} \to (S_K)_{\mathfrak q_K}$ 是 Noether 局部环之间的
平坦局部同态。注意，其纤维
$$
(S_K)_{\mathfrak q_K} / \mathfrak q (S_K)_{\mathfrak q_K}
\cong (\kappa(\mathfrak q) \otimes_k K)_{\mathfrak q'}
$$
是 Cohen--Macaulay 环（引理 \ref{algebra-lemma-tensor-fields-CM}）
$\kappa(\mathfrak q) \otimes_k K$ 在某个适当素理想
$\mathfrak q'$ 处的局部化。因此，引理由
引理 \ref{algebra-lemma-CM-goes-up} 得出。
\end{proof}

\section{余极限与有限表示映射（二）}
\label{algebra-section-colimits-finite-presentation}

\noindent
本节是节 \ref{algebra-section-colimits-flat} 的续篇。

\medskip\noindent
我们先给出平坦性开性的一个应用。
它说明可以用平坦模逼近平坦模；这一点很有用。

\begin{lemma}
\label{algebra-lemma-flat-finite-presentation-limit-flat}
设 $R \to S$ 是环映射。
设 $M$ 是 $S$-模。
假设
\begin{enumerate}
\item $R \to S$ 是有限表示的；
\item $M$ 是有限表示 $S$-模；
\item $M$ 在 $R$ 上平坦。
\end{enumerate}
在此情形下有：
\begin{enumerate}
\item 存在有限型 $\mathbf{Z}$-代数 $R_0$、有限型环映射
$R_0 \to S_0$ 和有限 $S_0$-模 $M_0$，使得 $M_0$ 在 $R_0$ 上平坦；
% P01-E1034
还存在环映射 $R_0 \to R$、$S_0 \to S$ 和 $S_0$-模映射 $M_0 \to M$，
使得 $S \cong R \otimes_{R_0} S_0$ 且
$M = S \otimes_{S_0} M_0$。
\item 若将 $R = \colim_{\lambda \in \Lambda} R_\lambda$ 写成有向余极限，
则存在 $\lambda$、有限表示环映射 $R_\lambda \to S_\lambda$ 和
有限表示 $S_\lambda$-模 $M_\lambda$，使得 $M_\lambda$ 在
$R_\lambda$ 上平坦，并且
$S = R \otimes_{R_\lambda} S_\lambda$ 且
$M = S \otimes_{S_{\lambda}} M_\lambda$。
\item 若
$$
(R \to S, M) =
\colim_{\lambda \in \Lambda}
(R_\lambda \to S_\lambda, M_\lambda)
$$
写成满足以下条件的有向余极限：
\begin{enumerate}
\item 对 $\mu \geq \lambda$，
$R_\mu \otimes_{R_\lambda} S_\lambda \to S_\mu$ 和
$S_\mu \otimes_{S_\lambda} M_\lambda \to M_\mu$ 是同构；
\item $R_\lambda \to S_\lambda$ 是有限表示的；
\item $M_\lambda$ 是有限表示 $S_\lambda$-模；
\end{enumerate}
则对于所有充分大的 $\lambda$，模 $M_\lambda$ 在 $R_\lambda$ 上平坦。
\end{enumerate}
\end{lemma}

\begin{proof}
先将 $(R \to S, M)$ 写成引理
\ref{algebra-lemma-limit-module-finite-presentation} 所述系统
$(R_\lambda \to S_\lambda, M_\lambda)$ 的有向余极限。
设 $\mathfrak q \subset S$ 是素理想。
设 $\mathfrak p \subset R$、$\mathfrak q_\lambda \subset S_\lambda$
和 $\mathfrak p_\lambda \subset R_\lambda$ 是相应的素理想。
如定理 \ref{algebra-theorem-openness-flatness} 的证明中所见，
$$
((R_\lambda)_{\mathfrak p_\lambda},
(S_\lambda)_{\mathfrak q_\lambda},
(M_\lambda)_{\mathfrak q_{\lambda}})
$$
是引理 \ref{algebra-lemma-limit-module-essentially-finite-presentation}
所述的系统，因而由引理 \ref{algebra-lemma-colimit-eventually-flat}，
可知存在某个 $\lambda_{\mathfrak q} \in \Lambda$，使得对所有
$\lambda \geq \lambda_{\mathfrak q}$，模 $M_\lambda$ 在素理想
$\mathfrak q_{\lambda}$ 处于 $R_\lambda$ 上平坦。

\medskip\noindent
由定理 \ref{algebra-theorem-openness-flatness}，得到开子集
$U_\lambda \subset \Spec(S_\lambda)$，使得 $M_\lambda$ 在
$U_\lambda$ 的所有素理想处于 $R_\lambda$ 上平坦。
记 $V_\lambda \subset \Spec(S)$ 为 $U_\lambda$ 在映射
$\Spec(S) \to \Spec(S_\lambda)$ 下的逆像。
上述论证说明，对每个 $\mathfrak q \in \Spec(S)$，存在
$\lambda_{\mathfrak q}$，使得对所有
$\lambda \geq \lambda_{\mathfrak q}$，都有
$\mathfrak q \in V_\lambda$。由于 $\Spec(S)$ 拟紧，可知这蕴含
存在单个 $\lambda_0 \in \Lambda$，使得
$V_{\lambda_0} = \Spec(S)$。

\medskip\noindent
补集 $\Spec(S_{\lambda_0}) \setminus U_{\lambda_0}$ 等于某个理想
$I \subset S_{\lambda_0}$ 的 $V(I)$。由于
$V_{\lambda_0} = \Spec(S)$，可知 $IS = S$。
% P01-E1035
选取 $f_1, \ldots, f_r \in I$ 和 $s_1, \ldots, s_n \in S$，
使得 $\sum f_i s_i = 1$。由于 $\colim S_\lambda = S$，增大
$\lambda_0$ 后，可以假设存在 $s_{i, \lambda_0} \in S_{\lambda_0}$，
使得 $\sum f_i s_{i, \lambda_0} = 1$。
因此，对这个 $\lambda_0$ 有
$U_{\lambda_0} = \Spec(S_{\lambda_0})$。
这证明了 (1)。

\medskip\noindent
证明 (2)。设 $(R_0 \to S_0, M_0)$ 如 (1) 所述，并假设
$R = \colim R_\lambda$。由于 $R_0$ 是有限型 $\mathbf{Z}$-代数，
存在 $\lambda$ 和映射 $R_0 \to R_\lambda$，使得
$R_0 \to R_\lambda \to R$ 是给定映射 $R_0 \to R$
（见引理 \ref{algebra-lemma-characterize-finite-presentation}）。
于是，取 $S_\lambda = R_\lambda \otimes_{R_0} S_0$ 和
$M_\lambda = S_\lambda \otimes_{S_0} M_0$，即得 (2)。

\medskip\noindent
最后证明 (3)。设 $(R_\lambda \to S_\lambda, M_\lambda)$ 如 (3) 所述。
按 (1) 选取 $(R_0 \to S_0, M_0)$ 和 $R_0 \to R$。
如 (2) 的证明所述，存在 $\lambda_0$ 和环映射
$R_0 \to R_{\lambda_0}$，使得
$R_0 \to R_{\lambda_0} \to R$ 是给定映射 $R_0 \to R$。
由于 $S_0$ 在 $R_0$ 上有限表示且 $S = \colim S_\lambda$，
可知对某个 $\lambda_1 \geq \lambda_0$，存在 $R_0$-代数映射
$S_0 \to S_{\lambda_1}$，使得复合
$S_0 \to S_{\lambda_1} \to S$ 是给定映射 $S_0 \to S$
（见引理 \ref{algebra-lemma-characterize-finite-presentation}）。
对所有 $\lambda \geq \lambda_1$，这给出映射
$$
\Psi_{\lambda} :
R_\lambda \otimes_{R_0} S_0
\longrightarrow
R_\lambda \otimes_{R_{\lambda_1}} S_{\lambda_1}
\cong
S_\lambda
$$
其中最后一个同构来自假设。由构造，
$\colim_\lambda \Psi_\lambda$ 是同构。因此，由引理
\ref{algebra-lemma-colimit-category-fp-algebras}，对所有充分大的 $\lambda$，
$\Psi_\lambda$ 是同构。
同理，存在 $\lambda_2 \geq \lambda_1$ 和 $S_0$-模映射
$M_0 \to M_{\lambda_2}$，使得
$M_0 \to M_{\lambda_2} \to M$ 是给定映射
$M_0 \to M$（见引理 \ref{algebra-lemma-module-map-property-in-colimit}）。
对 $\lambda \geq \lambda_2$，有诱导映射
$$
S_\lambda \otimes_{S_0} M_0
\longrightarrow
S_\lambda \otimes_{S_{\lambda_2}} M_{\lambda_2}
\cong
M_\lambda
$$
而由引理 \ref{algebra-lemma-colimit-category-fp-modules}，对充分大的
$\lambda$，此映射是同构。由于 $M_0$ 在 $R_0$ 上平坦，这就蕴含 (3)。
\end{proof}

\begin{lemma}
\label{algebra-lemma-descend-faithfully-flat-finite-presentation}
设 $R \to A \to B$ 是环映射。
假设 $A \to B$ 是有限表示的忠实平坦映射。
则存在交换图
$$
\xymatrix{
R \ar[r] \ar@{=}[d] &
A_0 \ar[d] \ar[r] &
B_0 \ar[d] \\
R \ar[r] & A \ar[r] & B
}
$$
其中 $R \to A_0$ 是有限表示的，
$A_0 \to B_0$ 是有限表示的忠实平坦映射，并且
$B = A \otimes_{A_0} B_0$。
\end{lemma}

\begin{proof}
% P01-E1036
先证明以 $\mathbf{Z}$ 替代 $R$ 时的引理。
由引理 \ref{algebra-lemma-flat-finite-presentation-limit-flat}，存在交换图
$$
\xymatrix{
A_0 \ar[r] \ar[d] & A \ar[d] \\
B_0 \ar[r] & B
}
$$
其中 $A_0$ 在 $\mathbf{Z}$ 上有限型，$B_0$ 在 $A_0$ 上平坦且有限表示，
并且 $B = A \otimes_{A_0} B_0$。
由于 $A_0 \to B_0$ 平坦且有限表示，可知映射
$\Spec(B_0) \to \Spec(A_0)$ 的像是开的，见
命题 \ref{algebra-proposition-fppf-open}。因此，像的补集等于某个理想
$I_0 \subset A_0$ 的 $V(I_0)$。
由于 $A \to B$ 忠实平坦，映射
$\Spec(B) \to \Spec(A)$ 满射，见引理 \ref{algebra-lemma-ff-rings}。
现在利用基变换的像等于像的基变换这一事实。
于是 $I_0A = A$。选取关系
$\sum f_i r_i = 1$，其中 $r_i \in A$、$f_i \in I_0$。
于是，在扩大 $A_0$ 使其包含元素 $r_i$（并相应扩大 $B_0$）后，
可知 $A_0 \to B_0$ 在谱上也满射，即忠实平坦。

\medskip\noindent
因此，当 $R = \mathbf{Z}$ 时引理成立。
在一般情形下，取刚得到的解 $A_0' \to B_0'$，并置
$A_0 = A_0' \otimes_{\mathbf{Z}} R$、
$B_0 = B_0' \otimes_{\mathbf{Z}} R$。
\end{proof}

\begin{lemma}
\label{algebra-lemma-colimit-finite}
设 $A = \colim_{i \in I} A_i$ 是环的有向余极限。
% P01-E1037
设 $0 \in I$，且 $\varphi_0 : B_0 \to C_0$ 是 $A_0$-代数的映射。
假设
\begin{enumerate}
\item $A \otimes_{A_0} B_0 \to A \otimes_{A_0} C_0$ 是有限的；
\item $C_0$ 在 $B_0$ 上有限型。
\end{enumerate}
则存在 $i \geq 0$，使得映射
$A_i \otimes_{A_0} B_0 \to A_i \otimes_{A_0} C_0$
是有限的。
\end{lemma}

\begin{proof}
设 $x_1, \ldots, x_m$ 是 $C_0$ 在 $B_0$ 上的一组生成元。
选取首一多项式 $P_j \in A \otimes_{A_0} B_0[T]$，使得
$P_j(1 \otimes x_j) = 0$ 在 $A \otimes_{A_0} C_0$ 中成立。
对某个 $i \geq 0$，可以找到映到 $P_j$ 的
$P_{j, i} \in A_i \otimes_{A_0} B_0[T]$。
由于 $\otimes$ 与余极限交换，必要时增大 $i$ 后，可知
$P_{j, i}(1 \otimes x_j)$ 在 $A_i \otimes_{A_0} C_0$ 中为零。
于是这个 $i$ 符合要求。
\end{proof}

\begin{lemma}
\label{algebra-lemma-colimit-surjective}
设 $A = \colim_{i \in I} A_i$ 是环的有向余极限。
设 $0 \in I$，且 $\varphi_0 : B_0 \to C_0$ 是 $A_0$-代数的映射。
假设
\begin{enumerate}
\item $A \otimes_{A_0} B_0 \to A \otimes_{A_0} C_0$ 是满射；
\item $C_0$ 在 $B_0$ 上有限型。
\end{enumerate}
则对某个 $i \geq 0$，映射
$A_i \otimes_{A_0} B_0 \to A_i \otimes_{A_0} C_0$
是满射。
\end{lemma}

\begin{proof}
设 $x_1, \ldots, x_m$ 是 $C_0$ 在 $B_0$ 上的一组生成元。
选取 $b_j \in A \otimes_{A_0} B_0$，使其在
$A \otimes_{A_0} C_0$ 中映到 $1 \otimes x_j$。
对某个 $i \geq 0$，可以找到映到 $b_j$ 的
$b_{j, i} \in A_i \otimes_{A_0} B_0$。
增大 $i$ 后，可以假设对所有 $j = 1, \ldots, m$，
$b_{j, i}$ 在 $A_i \otimes_{A_0} C_0$ 中映到 $1 \otimes x_j$。
于是这个 $i$ 符合要求。
\end{proof}

\begin{lemma}
\label{algebra-lemma-colimit-unramified}
设 $A = \colim_{i \in I} A_i$ 是环的有向余极限。
设 $0 \in I$，且 $\varphi_0 : B_0 \to C_0$ 是 $A_0$-代数的映射。
假设
\begin{enumerate}
\item $A \otimes_{A_0} B_0 \to A \otimes_{A_0} C_0$ 是非分歧的；
\item $C_0$ 在 $B_0$ 上有限型。
\end{enumerate}
则对某个 $i \geq 0$，映射
$A_i \otimes_{A_0} B_0 \to A_i \otimes_{A_0} C_0$
是非分歧的。
\end{lemma}

\begin{proof}
置 $B_i = A_i \otimes_{A_0} B_0$、$C_i = A_i \otimes_{A_0} C_0$、
$B = A \otimes_{A_0} B_0$ 和 $C = A \otimes_{A_0} C_0$。
设 $x_1, \ldots, x_m$ 是 $C_0$ 在 $B_0$ 上的一组生成元。
则 $\text{d}x_1, \ldots, \text{d}x_m$ 在 $C_0$ 上生成
$\Omega_{C_0/B_0}$，并且它们的像在 $C_i$ 上生成
$\Omega_{C_i/B_i}$（引理 \ref{algebra-lemma-differentials-polynomial-ring} 和
\ref{algebra-lemma-differential-seq}）。
注意
$0 = \Omega_{C/B} = \colim \Omega_{C_i/B_i}$
（引理 \ref{algebra-lemma-colimit-differentials}）。
因而存在 $i$，使得 $\text{d}x_1, \ldots, \text{d}x_m$ 映到零，
从而如所需有 $\Omega_{C_i/B_i} = 0$。
\end{proof}

\begin{lemma}
\label{algebra-lemma-colimit-isomorphism}
设 $A = \colim_{i \in I} A_i$ 是环的有向余极限。
设 $0 \in I$，且 $\varphi_0 : B_0 \to C_0$ 是 $A_0$-代数的映射。
假设
\begin{enumerate}
\item $A \otimes_{A_0} B_0 \to A \otimes_{A_0} C_0$
是同构；
\item $B_0 \to C_0$ 是有限表示的。
\end{enumerate}
则对某个 $i \geq 0$，映射
$A_i \otimes_{A_0} B_0 \to A_i \otimes_{A_0} C_0$
是同构。
\end{lemma}

\begin{proof}
由引理 \ref{algebra-lemma-colimit-surjective}，存在 $i$，使得
$A_i \otimes_{A_0} B_0 \to A_i \otimes_{A_0} C_0$
是满射。由于此映射是有限表示的，其核是有限生成理想。设
$g_1, \ldots, g_r \in A_i \otimes_{A_0} B_0$ 生成该核。
于是可以选取 $i' \geq i$，使得 $g_j$ 在
$A_{i'} \otimes_{A_0} B_0$ 中映到零。这样，
$A_{i'} \otimes_{A_0} B_0 \to A_{i'} \otimes_{A_0} C_0$
就是同构。
\end{proof}

\begin{lemma}
\label{algebra-lemma-colimit-etale}
设 $A = \colim_{i \in I} A_i$ 是环的有向余极限。
设 $0 \in I$，且 $\varphi_0 : B_0 \to C_0$ 是 $A_0$-代数的映射。
假设
\begin{enumerate}
\item $A \otimes_{A_0} B_0 \to A \otimes_{A_0} C_0$ 是平展的；
\item $B_0 \to C_0$ 是有限表示的。
\end{enumerate}
则对某个 $i \geq 0$，映射
$A_i \otimes_{A_0} B_0 \to A_i \otimes_{A_0} C_0$
是平展的。
\end{lemma}

\begin{proof}
写成 $C_0 = B_0[x_1, \ldots, x_n]/(f_{1, 0}, \ldots, f_{m, 0})$。
记 $B_i = A_i \otimes_{A_0} B_0$、$C_i = A_i \otimes_{A_0} C_0$。
注意，
$C_i = B_i[x_1, \ldots, x_n]/(f_{1, i}, \ldots, f_{m, i})$，
其中 $f_{j, i}$ 是 $f_{j, 0}$ 在 $B_i$ 上多项式环中的像。
记 $B = A \otimes_{A_0} B_0$、$C = A \otimes_{A_0} C_0$。
注意，
$C = B[x_1, \ldots, x_n]/(f_1, \ldots, f_m)$，
其中 $f_j$ 是 $f_{j, 0}$ 在 $B$ 上多项式环中的像。假设是映射
$$
\text{d} :
(f_1, \ldots, f_m)/(f_1, \ldots, f_m)^2
\longrightarrow
\bigoplus C \text{d}x_k
$$
为同构。因此，对充分大的 $i$，可以找到元素
$$
\xi_{k, i} \in (f_{1, i}, \ldots, f_{m, i})/(f_{1, i}, \ldots, f_{m, i})^2
$$
使得 $\text{d}\xi_{k, i} = \text{d}x_k$ 在
$\bigoplus C_i \text{d}x_k$ 中成立。
此外，必要时增大 $i$ 后，可知
$\sum (\partial f_{j, i}/\partial x_k) \xi_{k, i} =
f_{j, i} \bmod (f_{1, i}, \ldots, f_{m, i})^2$，
因为这在极限中成立。于是这个 $i$ 符合要求。
\end{proof}


\begin{lemma}
\label{algebra-lemma-colimit-smooth}
设 $A = \colim_{i \in I} A_i$ 是环的有向余极限。
设 $0 \in I$，且 $\varphi_0 : B_0 \to C_0$ 是 $A_0$-代数的映射。
假设
\begin{enumerate}
\item $A \otimes_{A_0} B_0 \to A \otimes_{A_0} C_0$ 是光滑的；
\item $B_0 \to C_0$ 是有限表示的。
\end{enumerate}
则对某个 $i \geq 0$，映射
$A_i \otimes_{A_0} B_0 \to A_i \otimes_{A_0} C_0$
是光滑的。
\end{lemma}

\begin{proof}
写成 $C_0 = B_0[x_1, \ldots, x_n]/(f_{1, 0}, \ldots, f_{m, 0})$。
记 $B_i = A_i \otimes_{A_0} B_0$、$C_i = A_i \otimes_{A_0} C_0$。
注意，
$C_i = B_i[x_1, \ldots, x_n]/(f_{1, i}, \ldots, f_{m, i})$，
其中 $f_{j, i}$ 是 $f_{j, 0}$ 在 $B_i$ 上多项式环中的像。
记 $B = A \otimes_{A_0} B_0$、$C = A \otimes_{A_0} C_0$。
注意，
$C = B[x_1, \ldots, x_n]/(f_1, \ldots, f_m)$，
其中 $f_j$ 是 $f_{j, 0}$ 在 $B$ 上多项式环中的像。假设是映射
$$
\text{d} :
(f_1, \ldots, f_m)/(f_1, \ldots, f_m)^2
\longrightarrow
\bigoplus C \text{d}x_k
$$
为分裂单射。设
$\xi_k \in (f_1, \ldots, f_m)/(f_1, \ldots, f_m)^2$ 是满足
$\sum (\partial f_j/\partial x_k) \xi_k =
f_j \bmod (f_1, \ldots, f_m)^2$ 的元素。则对充分大的 $i$，
可以找到元素
$$
\xi_{k, i} \in (f_{1, i}, \ldots, f_{m, i})/(f_{1, i}, \ldots, f_{m, i})^2
$$
使得
$\sum (\partial f_{j, i}/\partial x_k) \xi_{k, i} =
f_{j, i} \bmod (f_{1, i}, \ldots, f_{m, i})^2$，
因为这在极限中成立。于是这个 $i$ 符合要求。
\end{proof}


\begin{lemma}
\label{algebra-lemma-colimit-lci}
设 $A = \colim_{i \in I} A_i$ 是环的有向余极限。
设 $0 \in I$，且 $\varphi_0 : B_0 \to C_0$ 是 $A_0$-代数的映射。
假设
\begin{enumerate}
\item $A \otimes_{A_0} B_0 \to A \otimes_{A_0} C_0$
是 syntomic 的（相应地，是相对全局完全交）；
\item $C_0$ 在 $B_0$ 上有限表示。
\end{enumerate}
则存在 $i \geq 0$，使得映射
$A_i \otimes_{A_0} B_0 \to A_i \otimes_{A_0} C_0$
是 syntomic 的（相应地，是相对全局完全交）。
\end{lemma}

\begin{proof}
假设 $A \otimes_{A_0} B_0 \to A \otimes_{A_0} C_0$
是相对全局完全交。
由引理 \ref{algebra-lemma-relative-global-complete-intersection-Noetherian}，
存在有限型 $\mathbf{Z}$-代数 $R$、环映射
$R \to A \otimes_{A_0} B_0$、相对全局完全交 $R \to S$ 和同构
$$
(A \otimes_{A_0} B_0) \otimes_R S
\longrightarrow
A \otimes_{A_0} C_0
$$
由于 $R$ 在 $\mathbf{Z}$ 上有限型（从而有限表示），存在
$i$ 和映射 $R \to A_i \otimes_{A_0} B_0$，提升映射
$R \to A \otimes_{A_0} B_0$，见引理
\ref{algebra-lemma-characterize-finite-presentation}。
再次应用同一引理，存在 $i' \geq i$，使得
$(A_i \otimes_{A_0} B_0) \otimes_R S \to A \otimes_{A_0} C_0$
来自映射
$(A_i \otimes_{A_0} B_0) \otimes_R S \to A_{i'} \otimes_{A_0} C_0$。
因此，以 $i'$ 替换 $i$ 后，可以假设上述映射来自
$A_i \otimes_{A_0} B_0$-代数映射
$$
(A_i \otimes_{A_0} B_0) \otimes_R S
\longrightarrow
A_i \otimes_{A_0} C_0
$$
由引理 \ref{algebra-lemma-colimit-isomorphism}，增大 $i$ 后，此映射是同构。
这完成了此种情形的证明，因为由引理
\ref{algebra-lemma-base-change-relative-global-complete-intersection}，
相对全局完全交的基变换仍是相对全局完全交。

\medskip\noindent
假设 $A \otimes_{A_0} B_0 \to A \otimes_{A_0} C_0$
是 syntomic 的。则存在元素 $g_1, \ldots, g_m$，它们在
$A \otimes_{A_0} C_0$ 中生成单位理想，并且
$A \otimes_{A_0} B_0 \to (A \otimes_{A_0} C_0)_{g_j}$
是相对全局完全交，见引理 \ref{algebra-lemma-syntomic}。
可以找到 $i$ 和映到 $g_j$ 的元素
$g_{i, j} \in A_i \otimes_{A_0} C_0$。
增大 $i$ 后，可以假设
$g_{i, 1}, \ldots, g_{i, m}$ 在
$A_i \otimes_{A_0} C_0$ 中生成单位理想。
由前一段的结果，再次增大 $i$ 后，可以假设映射
$A_i \otimes_{A_0} B_0 \to (A_i \otimes_{A_0} C_0)_{g_{i, j}}$
都是相对全局完全交。
于是，由引理 \ref{algebra-lemma-local-syntomic}
（以及已经用过的引理 \ref{algebra-lemma-syntomic}），
$A_i \otimes_{A_0} B_0 \to A_i \otimes_{A_0} C_0$
是 syntomic 的。
\end{proof}













\noindent
下面的引理是上述结果的一个应用，似乎不太适合放在其他地方。

\begin{lemma}
\label{algebra-lemma-fppf-fpqf}
设 $R \to S$ 是有限表示的忠实平坦环映射。
则存在交换图
$$
\xymatrix{
S \ar[rr] & & S' \\
& R \ar[lu] \ar[ru]
}
$$
其中 $R \to S'$ 拟有限、忠实平坦且有限表示。
\end{lemma}

\begin{proof}
第一步，将此引理化归到 $R$ 在 $\mathbf{Z}$ 上有限型的情形。
由引理 \ref{algebra-lemma-descend-faithfully-flat-finite-presentation}，
存在交换图
$$
\xymatrix{
S_0 \ar[r] & S \\
R_0 \ar[u] \ar[r] & R \ar[u]
}
$$
其中 $R_0$ 在 $\mathbf{Z}$ 上有限型，而 $S_0$ 在 $R_0$ 上
忠实平坦且有限表示，并且 $S = R \otimes_{R_0} S_0$。
若对环映射 $R_0 \to S_0$ 证明了引理，则由基变换，对
$R \to S$ 也得到引理，因为拟有限环映射的基变换仍拟有限，见
引理 \ref{algebra-lemma-quasi-finite-base-change}。（当然，还使用了平坦映射的
基变换仍平坦，以及有限表示映射的基变换仍有限表示。）

\medskip\noindent
假设 $R \to S$ 是有限表示的忠实平坦环映射，且 $R$ 是 Noether 的
（由前一段可以如此假设）。设 $W \subset \Spec(S)$ 是引理
\ref{algebra-lemma-finite-presentation-flat-CM-locus-open} 的开集。
由于 $R \to S$ 忠实平坦，映射
$\Spec(S) \to \Spec(R)$ 满射，见引理 \ref{algebra-lemma-ff-rings}。
由引理 \ref{algebra-lemma-generic-CM-flat-finite-presentation}，
映射 $W \to \Spec(R)$ 也满射。
因此，将 $S$ 替换为乘积
$S_{g_1} \times \ldots \times S_{g_m}$ 后，可以假设
$W = \Spec(S)$；这里使用了 $\Spec(R)$ 拟紧
（引理 \ref{algebra-lemma-quasi-compact}），以及映射
$\Spec(S) \to \Spec(R)$ 是开的
（命题 \ref{algebra-proposition-fppf-open}）。
假设 $\mathfrak p \subset R$ 是素理想。选取位于
$\mathfrak p$ 之上的素理想 $\mathfrak q \subset S$，它对应于
纤维环 $S \otimes_R \kappa(\mathfrak p)$ 的一个极大理想。
Noether 局部环
$\overline{S}_{\mathfrak q} = S_{\mathfrak q}/\mathfrak pS_{\mathfrak q}$
是 Cohen--Macaulay 的，设其维数为 $d$。可以在
$S_{\mathfrak q}$ 的极大理想中选取 $f_1, \ldots, f_d$，使它们在
$\overline{S}_{\mathfrak q}$ 中的像构成正则序列。
选取 $f_1, \ldots, f_d$ 的一个公分母
$g \in S$，其中 $g \not \in \mathfrak q$，并考虑 $R$-代数
$$
S' = S_g/(f_1, \ldots, f_d).
$$
由构造，存在位于 $\mathfrak p$ 之上的素理想
$\mathfrak q' \subset S'$，它对应于 $\mathfrak q$（通过
$S_g \to S'_g$）。同样由构造，环映射 $R \to S'$ 在
% P01-E1038
$\mathfrak q$ 处拟有限，因为局部环
% P01-E1039
$$
S'_{\mathfrak q'}/\mathfrak pS'_{\mathfrak q'} =
S_{\mathfrak q}/(f_1, \ldots, f_d) + \mathfrak pS_{\mathfrak q} =
\overline{S}_{\mathfrak q}/(\overline{f}_1, \ldots, \overline{f}_d)
$$
维数为零，见引理 \ref{algebra-lemma-isolated-point-fibre}。
同样由构造，$R \to S'$ 是有限表示的。
最后，由引理 \ref{algebra-lemma-grothendieck-regular-sequence}，局部环映射
$R_{\mathfrak p} \to S'_{\mathfrak q'}$ 平坦（这里使用了 $R$
是 Noether 的）。因此，由平坦性的开性
（定理 \ref{algebra-theorem-openness-flatness}）和拟有限性的开性
（引理 \ref{algebra-lemma-quasi-finite-open}），可以用
$gg'$ 替换 $g$，其中某个合适的 $g' \in S$ 满足
$g' \not \in \mathfrak q$，从而假设 $R \to S'$ 平坦且拟有限。
像 $\Spec(S') \to \Spec(R)$ 是开集并包含 $\mathfrak p$。
换言之，我们已经证明，存在引理陈述所述的环 $S'$
（可能忠实性除外），其像包含任意给定的素理想。
再次使用 $\Spec(R)$ 的拟紧性，可知有限多个这样的环的乘积符合要求。
\end{proof}
% Generated mechanically from the frozen brauer.tex target.
% Do not edit this generated file; edit the bound locale source instead.

% OK, start here.
%

\setcounter{chapter}{10}
\chapter{Brauer 群}



\phantomsection
\label{brauer-section-phantom}


\section{引言}
\label{brauer-section-introduction}

\noindent
一个参考资料是 Serre 在 Séminaire Cartan 上的讲义，见
\cite{Serre-Cartan}。Serre 又引用了
\cite{Deuring} 和 \cite{ANT}。我们改写了其中的一些证明，特别是使用
Rieffel 的一个巧妙论证来证明 Wedderburn 定理。
这一改写很可能并非改进；我们强烈建议读者阅读 Serre 的原始论述。


\section{非交换代数}
\label{brauer-section-algebras}

\noindent
设 $k$ 为域。本章所谓域 $k$ 上的一个{\it 代数} $A$，是指一个可能
非交换的环 $A$，连同一个环同态 $k \to A$，使得 $k$ 的像包含于
$A$ 的中心，且 $1$ 映到 $A$ 的单位元。所谓一个 {\it $A$-模}，
是指 $A$ 的单位元按恒等映射作用的右 $A$-模。

\begin{definition}
\label{brauer-definition-finite}
设 $A$ 为 $k$-代数。若 $\dim_k(A) < \infty$，则称 $A$ 为
{\it 有限维的}。此时记 $[A : k] = \dim_k(A)$。
\end{definition}

\begin{definition}
\label{brauer-definition-skew-field}
{\it 除环}是一个带单位元 $1$ 的、可能非交换的环，其中
$1 \not = 0$，且每个非零元素都有乘法逆元。
\end{definition}

\noindent
一个除环总是某个 $k$ 上的代数（例如可取它所含的素域作为 $k$）。
下文将使用如下事实：除环上的任意模都是自由模，因为任一极大的
线性无关向量集构成一组基，而这样的集合由 Zorn 引理存在。

\begin{definition}
\label{brauer-definition-simple}
设 $A$ 为 $k$-代数。
若一个 $A$-模 $M$ 非零，且其仅有的 $A$-子模为 $0$ 和 $M$，
则称 $M$ 是{\it 单模}。
若 $A$ 的双边理想只有 $0$ 和 $A$，则称 $A$ 是{\it 单的}。
\end{definition}

\begin{definition}
\label{brauer-definition-central}
若 $A$ 的中心恰为 $k \to A$ 的像，则称 $k$-代数 $A$ 是
{\it 中心的}。
\end{definition}

\begin{definition}
\label{brauer-definition-opposite}
给定一个 $k$-代数 $A$，以 $A^{op}$ 表示将 $A$ 中乘法次序反转
所得的 $k$-代数。它称为{\it 反代数}。
\end{definition}




\section{Wedderburn 定理}
\label{brauer-section-wedderburn}

\noindent
以下巧妙论证可见于 Rieffel 的一篇论文，见
\cite{Rieffel}。这个证明不能再简单了（引自 Carl Faith 的评论）。

\begin{lemma}
\label{brauer-lemma-rieffel}
设 $A$ 是一个带 $1$ 的、可能非交换的环，且没有非平凡双边理想。
设 $M$ 是 $A$ 的一个非零右理想，并将 $M$ 视为右 $A$-模。
则 $A$ 与 $M$ 的双交换子相同。
\end{lemma}

\begin{proof}
令 $A' = \text{End}_A(M)$，于是 $M$ 是左 $A'$-模。
置 $A'' = \text{End}_{A'}(M)$（即 $M$ 的双交换子）。
我们将 $M$ 视为右 $A''$-模\footnote{这意味着，给定
$a'' \in A''$ 与 $m \in M$，有乘积 $m a'' \in M$。特别地，
$A''$ 中的乘法次序与把 $A''$ 的元素写成从左作用的自同态时
得到的次序相反。}。
令 $R : A \to A''$ 为满足 $mR(a) = ma$ 的自然同态。
由于 $R(1) = \text{id}_M$ 且 $A$ 没有非平凡双边理想，$R$ 是单射。
我们断言 $R(M)$ 是 $A''$ 的一个右理想。事实上，对
$a'' \in A''$ 和 $m$ 属于 $M$ 时有 $R(m)a'' = R(ma'')$，因为由
$M$ 的任一元素 $n$ 给出的左乘表示 $A'$ 的一个元素，从而对
所有属于 $M$ 的 $n$ 均有 $(nm)a'' = n(ma'')$。
最后，乘积理想 $AM$ 是双边理想，故 $A = AM$。于是
$R(A) = R(A)R(M)$，所以 $R(A)$ 是 $A''$ 的右理想。
但 $R(A)$ 含有 $A''$ 的单位元，故 $R(A) = A''$。
\end{proof}

\begin{lemma}
\label{brauer-lemma-simple-module}
设 $A$ 为 $k$-代数。若 $A$ 是有限维的，则
\begin{enumerate}
\item $A$ 有单模；
\item 任一非零模都含有单子模；
\item $A$ 上的单模在 $k$ 上是有限维的；并且
\item 若 $M$ 是单 $A$-模，则 $\text{End}_A(M)$ 是除环。
\end{enumerate}
\end{lemma}

\begin{proof}
当然，由于 $A$ 自身是一个非零 $A$-模，(1) 可由 (2) 推出。
对于 (2)，任一作为 $k$-向量空间具有极小（有限）维数的子模都是
单模。这样的有限维子模存在，因为循环子模就是有限维的。
若 $M$ 是单模，则 $mA \subset M$ 是一个子模，由此得到 (3)。
$\text{End}_A(M)$ 的任一非零元素都是同构，故 (4) 成立。
\end{proof}

\begin{theorem}
\label{brauer-theorem-wedderburn}
\begin{slogan}
域上的有限维单代数是除环上的矩阵代数。
\end{slogan}
设 $A$ 是一个有限维单 $k$-代数。则 $A$ 是某个有限维
$k$-代数 $K$ 上的矩阵代数，其中后者是除环。
\end{theorem}

\begin{proof}
可以选择一个单子模 $M \subset A$，此时
$k$-代数 $K = \text{End}_A(M)$ 是除环，见
引理 \ref{brauer-lemma-simple-module}。
由
引理 \ref{brauer-lemma-rieffel}
可知 $A = \text{End}_K(M)$。由于 $K$ 是除环且
$M$ 有限生成（因为 $\dim_k(M) < \infty$），可知
$M$ 作为左 $K$-模是有限秩自由模。立即得到
$A \cong \text{Mat}(n \times n, K^{op})$。
\end{proof}






\section{代数的若干引理}
\label{brauer-section-lemmas}

\noindent
设 $A$ 为 $k$-代数，$B \subset A$ 为子代数。
{\it $B$ 在 $A$ 中的中心化子}是子代数
$$
C  = \{y \in A \mid xy = yx \text{ 对所有 }x \in B\}.
$$
它是一个 $k$-代数。

\begin{lemma}
\label{brauer-lemma-centralizer}
设 $A$、$A'$ 为 $k$-代数。设 $B \subset A$、$B' \subset A'$ 为
子代数，其中心化子分别为 $C$、$C'$。则
$B \otimes_k B'$ 在 $A \otimes_k A'$ 中的中心化子为
$C \otimes_k C'$。
\end{lemma}

\begin{proof}
以 $C'' \subset A \otimes_k A'$ 表示 $B \otimes_k B'$ 的
中心化子。显然 $C \otimes_k C' \subset C''$。反过来，
$C''$ 的每个元素都与 $B \otimes 1$ 交换，因而属于
$C \otimes_k A'$。类似地，$C'' \subset A \otimes_k C'$。故
$C'' \subset C \otimes_k A' \cap A \otimes_k C' = C \otimes_k C'$。
\end{proof}

\begin{lemma}
\label{brauer-lemma-center-csa}
设 $A$ 是有限维单 $k$-代数。则 $A$ 的中心 $k'$ 是 $k$ 的
有限域扩张。
\end{lemma}

\begin{proof}
由
定理 \ref{brauer-theorem-wedderburn}，
可写成 $A = \text{Mat}(n \times n, K)$，其中 $K$ 是
$k$ 上的有限维除环。
由
引理 \ref{brauer-lemma-centralizer}，
$A$ 的中心为 $k \otimes_k k'$，其中 $k' \subset K$ 是 $K$ 的
中心。除环的中心是域，结论得证。
\end{proof}

\begin{lemma}
\label{brauer-lemma-generate-two-sided-sub}
设 $V$ 为 $k$-向量空间。设 $K$ 是中心 $k$-代数且为除环。
设 $W \subset V \otimes_k K$ 是一个双边 $K$-子向量空间。
则作为左 $K$-向量空间，$W$ 由 $W \cap (V \otimes 1)$ 生成。
\end{lemma}

\begin{proof}
令 $V' \subset V$ 为所有满足 $v \otimes 1 \in W$ 的
$v \in V$ 所生成的 $k$-子向量空间。于是
$V' \otimes_k K \subset W$，并且
$$
W/(V' \otimes_k K)  \subset  (V/V') \otimes_k K.
$$
若 $\overline{v} \in V/V'$ 是一个非零向量，且
$\overline{v} \otimes 1$ 属于 $W/(V' \otimes_k K)$，
则对 $\overline{v}$ 的一个提升 $v \in V$，有
$v \otimes 1 \in W$。这与 $V'$ 的构造矛盾。因此可将
$V$ 替换为 $V/V'$，将 $W$ 替换为 $W/(V' \otimes_k K)$；
只需证明当 $W$ 非零时，$W \cap (V \otimes 1)$ 非零。

\medskip\noindent
为此，取非零元素 $w \in W$，并将其写成
$w = \sum_{i = 1, \ldots, n} v_i \otimes k_i$，其中 $n$ 最小。
右乘 $k_1^{-1}$ 后，可设 $k_1 = 1$。
若 $n = 1$，则因 $v_1 \otimes 1 \in W$，结论成立。
若 $n > 1$，则对任意 $c \in K$，
$$
c w - w c = \sum\nolimits_{i = 2, \ldots, n} v_i \otimes (c k_i - k_i c) \in W
$$
而由 $n$ 的极小性，必有 $c k_i - k_i c = 0$。
这说明 $k_i$ 位于 $K$ 的中心；按假设该中心就是 $k$。
于是 $w = (v_1 + \sum k_i v_i) \otimes 1$，与 $n$ 的极小性矛盾。
\end{proof}

\begin{lemma}
\label{brauer-lemma-generate-two-sided-ideal}
设 $A$ 为 $k$-代数。设 $K$ 是中心 $k$-代数且为除环。
则任一双边理想 $I \subset A \otimes_k K$ 都可写成
$J \otimes_k K$，其中 $J \subset A$ 是双边理想。
特别地，若 $A$ 是单的，则 $A \otimes_k K$ 也是单的。
\end{lemma}

\begin{proof}
令 $J = \{a \in A \mid a \otimes 1 \in I\}$。这是 $A$ 的一个
双边理想。由
引理 \ref{brauer-lemma-generate-two-sided-sub}
有 $I = J \otimes_k K$。
\end{proof}

\begin{lemma}
\label{brauer-lemma-matrix-algebras}
设 $R$ 是一个可能非交换的环，$n \geq 1$ 是整数。
令 $R_n = \text{Mat}(n \times n, R)$。
\begin{enumerate}
\item 函子 $M \mapsto M^{\oplus n}$ 与
$N \mapsto Ne_{11}$ 给出范畴之间的一对拟逆等价
$\text{Mod}_R \leftrightarrow \text{Mod}_{R_n}$。
\item $R_n$ 的任一双边理想都形如 $IR_n$，其中
$I$ 是 $R$ 的双边理想。
\item $R_n$ 的中心等于 $R$ 的中心。
\end{enumerate}
\end{lemma}

\begin{proof}
第 (1) 项是显然的。若 $J \subset R_n$ 是双边理想，则
$J = \bigoplus e_{ii}Je_{jj}$，且所有分量 $e_{ii}Je_{jj}$
彼此相等，并构成 $R$ 的一个双边理想 $I$。这证明了 (2)。
第 (3) 项显然。
\end{proof}

\begin{lemma}
\label{brauer-lemma-simple-module-unique}
设 $A$ 是有限维单 $k$-代数。
\begin{enumerate}
\item 在同构意义下，恰有一个单 $A$-模 $M$。
\item 任一有限维 $A$-模都是若干个单模副本的直和。
\item 两个有限维 $A$-模同构，当且仅当它们在 $k$ 上的维数相同。
\item 若 $A = \text{Mat}(n \times n, K)$，其中 $K$ 是 $k$ 的一个
有限维除环扩张，则 $M = K^{\oplus n}$ 是单 $A$-模，且
$\text{End}_A(M) = K^{op}$。
\item 若 $M$ 是单 $A$-模，则 $L = \text{End}_A(M)$ 是 $k$ 上的
有限维除环，并从左作用于 $M$；此外
$A = \text{End}_L(M)$，且 $A$ 与 $L$ 的中心相同。
并且 $[A : k] [L : k] = \dim_k(M)^2$。
\item 对于有限维 $A$-模 $N$，代数 $B = \text{End}_A(N)$ 是
(5) 中除环 $L$ 上的矩阵代数。此外 $\text{End}_B(N) = A$。
\end{enumerate}
\end{lemma}

\begin{proof}
由
定理 \ref{brauer-theorem-wedderburn}，
可写成 $A = \text{Mat}(n \times n, K)$，其中 $K$ 是
$k$ 的一个有限维除环扩张。由
引理 \ref{brauer-lemma-matrix-algebras}，
$A$-模范畴等价于 $K$-模范畴。由于 $K$ 上的每个模都是自由模，
(1)、(2)、(3) 成立。第 (4) 项成立，因为该等价将 $K$-模 $K$
变换为 $M = K^{\oplus n}$。在 (5) 中取 $M = K^{\oplus n}$，
可见 $L = K^{op}$。关于 $L = K^{op}$ 的中心的断言来自
引理 \ref{brauer-lemma-matrix-algebras}。
关于 $\text{End}_L(M)$ 的断言由 $M$ 的显式形式得到。
维数公式显然。第 (6) 项成立，因为 $N$ 同构于若干个单模副本的直和。
\end{proof}

\begin{lemma}
\label{brauer-lemma-tensor-simple}
设 $A$、$A'$ 是两个单 $k$-代数，其中一个在 $k$ 上有限维且中心。
则 $A \otimes_k A'$ 是单的。
\end{lemma}

\begin{proof}
设 $A'$ 在 $k$ 上有限维且中心。
由
定理 \ref{brauer-theorem-wedderburn}，
可写 $A' = \text{Mat}(n \times n, K')$。
此时 $K'$ 的中心为 $k$，由
引理 \ref{brauer-lemma-generate-two-sided-ideal}
可知 $A \otimes_k K'$ 是单的。
因此，由引理 \ref{brauer-lemma-matrix-algebras}，
$A \otimes_k A' = \text{Mat}(n \times n, A \otimes_k K')$ 是单的。
\end{proof}

\begin{lemma}
\label{brauer-lemma-tensor-central-simple}
域 $k$ 上两个有限维中心单代数的张量积仍是有限维、中心且单的。
\end{lemma}

\begin{proof}
合并引理 \ref{brauer-lemma-centralizer} 与 \ref{brauer-lemma-tensor-simple}。
\end{proof}

\begin{lemma}
\label{brauer-lemma-base-change}
设 $A$ 是 $k$ 上的有限维中心单代数，$k'/k$ 是域扩张。
则 $A' = A \otimes_k k'$ 是 $k'$ 上的有限维中心单代数。
\end{lemma}

\begin{proof}
合并引理 \ref{brauer-lemma-centralizer} 与 \ref{brauer-lemma-tensor-simple}。
\end{proof}

\begin{lemma}
\label{brauer-lemma-inverse}
设 $A$ 是 $k$ 上的有限维中心单代数。则
$A \otimes_k A^{op} \cong \text{Mat}(n \times n, k)$，
其中 $n = [A : k]$。
\end{lemma}

\begin{proof}
由引理 \ref{brauer-lemma-tensor-central-simple}，
代数 $A \otimes_k A^{op}$ 是单的。因此映射
$$
A \otimes_k A^{op} \longrightarrow \text{End}_k(A),\quad
a \otimes a' \longmapsto (x \mapsto axa')
$$
是单射。箭头两端维数相同，故结论成立。
\end{proof}





\section{域的 Brauer 群}
\label{brauer-section-brauer}

\noindent
设 $k$ 为域。考虑 $k$ 上两个有限维中心单代数 $A$ 与 $B$。
若存在 $n, m > 0$，使得作为 $k$-代数有
$\text{Mat}(n \times n, A) \cong \text{Mat}(m \times m, B)$，
则称 $A$ 与 $B$ {\it 相似}。

\begin{lemma}
\label{brauer-lemma-similar}
关于相似性，有：
\begin{enumerate}
\item 相似性在 $k$ 上有限维中心单代数的同构类集合上定义了一个
等价关系。
\item 每个相似类都含有唯一的（在同构意义下）$k$ 上有限维中心除环。
\item 若 $A = \text{Mat}(n \times n, K)$ 且
$B = \text{Mat}(m \times m, K')$，其中 $K$、$K'$ 是 $k$ 上的
有限维中心除环，则 $A$ 与 $B$ 相似，当且仅当
$K \cong K'$（作为 $k$-代数）。
\end{enumerate}
\end{lemma}

\begin{proof}
注意，由 Wedderburn 定理（定理 \ref{brauer-theorem-wedderburn}），
任一有限维中心单代数都可写成某个有限维中心除环上的矩阵代数。
因此只需证明第三个断言。为此，只需证明若
$A = \text{Mat}(n \times n, K) \cong \text{Mat}(m \times m, K') = B$，
则 $K \cong K'$。事实上，对 $A$ 的一个单模 $M$，由
引理 \ref{brauer-lemma-simple-module-unique}
有 $\text{End}_A(M) = K^{op}$。
所以 $A \cong B$ 蕴含 $K^{op} \cong (K')^{op}$，结论得证。
\end{proof}

\noindent
给定两个有限维中心单 $k$-代数 $A$、$B$，其张量积
$A \otimes_k B$ 仍是这样的代数，见
引理 \ref{brauer-lemma-tensor-central-simple}。
此外，若 $A$ 与 $A'$ 相似，则 $A \otimes_k B$ 与
$A' \otimes_k B$ 相似，因为张量积与取矩阵代数可交换。
因此张量积在有限维中心单代数的等价类上定义了一个显然满足
结合律和交换律的运算。最后，
引理 \ref{brauer-lemma-inverse}
表明 $A \otimes_k A^{op}$ 同构于矩阵代数，也就是说，
$A \otimes_k A^{op}$ 属于 $k$ 的相似类。
由此得到一个阿贝尔群。

\begin{definition}
\label{brauer-definition-brauer-group}
设 $k$ 为域。$k$ 的 {\it Brauer 群}是上述有限维中心单
$k$-代数的相似类所成的阿贝尔群，记作 $\text{Br}(k)$。
\end{definition}

\noindent
对任意域同态 $k \to k'$，由引理 \ref{brauer-lemma-base-change}
得到群同态
$$
\text{Br}(k) \longrightarrow \text{Br}(k'),\quad
A \longmapsto A \otimes_k k'
$$
。换言之，$\text{Br}(-)$ 是从域范畴到阿贝尔群范畴的函子。
注意，一个域的 Brauer 群为零，当且仅当每个有限维中心除环扩张
$k \subset K$ 都是平凡的。

\begin{lemma}
\label{brauer-lemma-brauer-algebraically-closed}
代数闭域的 Brauer 群为零。
\end{lemma}

\begin{proof}
设 $k \subset K$ 为有限维中心除环扩张。
对任意元素 $x \in K$，子环 $k[x] \subset K$ 是一个交换的、
有限维且整的 $k$-子代数，因而是域，见
《交换代数》，引理 \ref{algebra-lemma-integral-over-field}。
由于 $k$ 代数闭，得到 $k[x] = k$。因 $x$ 任意，遂有 $k = K$。
\end{proof}

\begin{lemma}
\label{brauer-lemma-dimension-square}
设 $A$ 是域 $k$ 上的有限维中心单代数，则 $[A : k]$ 是完全平方数。
\end{lemma}

\begin{proof}
这是因为由
引理 \ref{brauer-lemma-brauer-algebraically-closed}，
$A \otimes_k \overline{k}$ 是 $\overline{k}$ 上的矩阵代数。
\end{proof}




\section{Skolem--Noether 定理}
\label{brauer-section-skolem-noether}



\begin{theorem}
\label{brauer-theorem-skolem-noether}
设 $A$ 是有限维中心单 $k$-代数，$B$ 是单 $k$-代数，并设
$f, g : B \to A$ 是两个 $k$-代数同态。则存在可逆元素
$x \in A$，使得对所有 $b \in B$ 都有
$f(b) = xg(b)x^{-1}$。
\end{theorem}

\begin{proof}
选择一个单 $A$-模 $M$，并令 $L = \text{End}_A(M)$。
由引理 \ref{brauer-lemma-simple-module} 和
\ref{brauer-lemma-simple-module-unique}，
$L$ 是中心为 $k$ 的除环，并从左作用于 $M$。
此时 $M$ 上有两个 $B \otimes_k L^{op}$-模结构，分别定义为
$m \cdot_1 (b \otimes l) = lmf(b)$ 和
$m \cdot_2 (b \otimes l) = lmg(b)$。
由引理 \ref{brauer-lemma-tensor-simple}，
$k$-代数 $B \otimes_k L^{op}$ 是单的。由于 $B$ 是单的，
$k$-代数同态 $B \to A$ 的存在性蕴含 $B$ 有限维。
因此 $B \otimes_k L^{op}$ 是有限维单代数；由
引理 \ref{brauer-lemma-simple-module-unique}，
$M$ 上这两个 $B \otimes_k L^{op}$-模结构同构。
所以存在交织这两个作用的 $\varphi : M \to M$。
特别地，$\varphi$ 属于 $L$ 的交换子，这说明 $\varphi$
是由某个 $x \in A$ 给出的乘法，见
引理 \ref{brauer-lemma-simple-module-unique}。展开定义可见，$x$
正是所求解。
\end{proof}

\begin{lemma}
\label{brauer-lemma-automorphism-inner}
设 $A$ 是有限维中心单 $k$-代数。$A$ 的任一自同构都是内自同构。
特别地，$\text{Mat}(n \times n, k)$ 的任一自同构都是内自同构。
\end{lemma}

\begin{proof}
应用 Skolem--Noether 定理（定理 \ref{brauer-theorem-skolem-noether}）。
\end{proof}



\section{中心化子定理}
\label{brauer-section-centralizer}



\begin{theorem}
\label{brauer-theorem-centralizer}
设 $A$ 是 $k$ 上的有限维中心单代数，$B$ 是 $A$ 的单子代数。
则
\begin{enumerate}
\item $B$ 在 $A$ 中的中心化子 $C$ 是单的；
\item $[A : k] = [B : k][C : k]$；并且
\item $C$ 在 $A$ 中的中心化子是 $B$。
\end{enumerate}
\end{theorem}

\begin{proof}
在本证明中，我们将自由使用
引理 \ref{brauer-lemma-simple-module-unique} 的结论。
选择一个单 $A$-模 $M$，并令 $L = \text{End}_A(M)$。
则 $L$ 是中心为 $k$ 的除环，从左作用于 $M$，并且
$A = \text{End}_L(M)$。
于是 $M$ 是右 $B \otimes_k L^{op}$-模，且
$C = \text{End}_{B \otimes_k L^{op}}(M)$。
由引理 \ref{brauer-lemma-tensor-simple}，代数
$B \otimes_k L^{op}$ 是单的，故 $C$ 是单的（再次应用
引理 \ref{brauer-lemma-simple-module-unique}）。

\medskip\noindent
写成 $B \otimes_k L^{op} = \text{Mat}(m \times m, K)$，
其中 $K$ 是 $k$ 上的有限维除环。若 $M$ 同构于单
$B \otimes_k L^{op}$-模 $K^{\oplus m}$ 的 $n$ 个副本的直和，
则 $C = \text{Mat}(n \times n, K^{op})$（再次应用该引理）。
因此 $\dim_k(M) = nm [K : k]$，
$[B : k] [L : k] = m^2 [K : k]$，
$[C : k] = n^2 [K : k]$，且
$[A : k] [L : k] = \dim_k(M)^2$（仍由该引理）。
由此得到 (2)。

\medskip\noindent
对 $C \subset A$ 应用 (2)，可知 (3) 成立：若 $C'$ 是
$C$ 在 $A$ 中的中心化子，则 $[B : k] = [C' : k]$，
而显然 $B \subset C'$。
\end{proof}

\begin{lemma}
\label{brauer-lemma-when-tensor-is-equal}
设 $A$ 是 $k$ 上的有限维中心单代数，$B$ 是 $A$ 的单子代数。
若 $B$ 是中心 $k$-代数，则 $A = B \otimes_k C$，其中
$C$ 是 $B$ 在 $A$ 中的（中心单）中心化子。
\end{lemma}

\begin{proof}
由定理 \ref{brauer-theorem-centralizer}，
$\dim_k(A) = \dim_k(B \otimes_k C)$。
由引理 \ref{brauer-lemma-tensor-simple}，该张量积是单的。
因此自然映射 $B \otimes_k C \to A$ 是单射，从而是同构。
\end{proof}

\begin{lemma}
\label{brauer-lemma-self-centralizing-subfield}
设 $A$ 是 $k$ 上的有限维中心单代数。若 $K \subset A$ 是子域，
则以下条件等价：
\begin{enumerate}
\item $[A : k] = [K : k]^2$；
\item $K$ 等于它自身的中心化子；并且
\item $K$ 是极大交换子环。
\end{enumerate}
\end{lemma}

\begin{proof}
定理 \ref{brauer-theorem-centralizer} 表明 (1) 与 (2) 等价。
显然，(3) 与 (2) 等价。
\end{proof}

\begin{lemma}
\label{brauer-lemma-maximal-subfield}
\begin{slogan}
有限维中心除环的维数等于其任一极大子域维数的平方。
\end{slogan}
设 $A$ 是 $k$ 上的有限维中心除环。
则每个极大子域 $K \subset A$ 都满足
$[A : k] = [K : k]^2$。
\end{lemma}

\begin{proof}
这是引理 \ref{brauer-lemma-self-centralizing-subfield} 的特殊情形。
\end{proof}





\section{分裂域}
\label{brauer-section-splitting}



\begin{definition}
\label{brauer-definition-splitting}
设 $A$ 是有限维中心单 $k$-代数。若 $A \otimes_k k'$ 是
$k'$ 上的矩阵代数，则称域扩张 $k'/k$ {\it 分裂} $A$，
或称 $k'$ 是 $A$ 的一个{\it 分裂域}。
\end{definition}

\noindent
等价地说，$A$ 的类在映射
$\text{Br}(k) \to \text{Br}(k')$ 下映为零。

\begin{theorem}
\label{brauer-theorem-splitting}
设 $A$ 是有限维中心单 $k$-代数，$k'/k$ 是有限域扩张。
以下条件等价：
\begin{enumerate}
\item $k'$ 分裂 $A$；并且
\item 存在一个与 $A$ 相似的有限维中心单代数 $B$，使得
$k' \subset B$ 且 $[B : k] = [k' : k]^2$。
\end{enumerate}
\end{theorem}

\begin{proof}
假设 (2) 成立。只需证明 $B \otimes_k k'$ 是矩阵代数。
已知 $B \otimes_k B^{op} \cong \text{End}_k(B)$。
由引理 \ref{brauer-lemma-self-centralizing-subfield}，$k'$ 是
$B^{op}$ 中 $k'$ 的中心化子，故 $B \otimes_k k'$ 是
$k \otimes k'$ 在
$B \otimes_k B^{op} = \text{End}_k(B)$ 中的中心化子。
当然，这个中心化子就是 $\text{End}_{k'}(B)$，其中借助嵌入
$k' \to B$ 将 $B$ 视为 $k'$-向量空间。结论随即成立。

\medskip\noindent
假设 (1) 成立。这意味着对某个 $k'$-向量空间 $V$，存在同构
$A \otimes_k k' \cong \text{End}_{k'}(V)$。
令 $B$ 为 $A$ 在 $\text{End}_k(V)$ 中的交换子。注意
$k'$ 包含于 $B$。由
引理 \ref{brauer-lemma-when-tensor-is-equal}，
$A$ 与 $B$ 的类在 $\text{Br}(k)$ 中之和为零。
由
定理 \ref{brauer-theorem-centralizer}
中的维数公式，
$$
[B : k] [A : k] =
\dim_k(V)^2 =
[k' : k]^2 \dim_{k'}(V)^2 =
[k' : k]^2 [A : k].
$$
因此 $[B : k] = [k' : k]^2$。这样，我们已对 $A$ 的 Brauer
类的相反类证明了结论。但 $k'$ 分裂 $A$ 的 Brauer 类，当且仅当
它分裂反代数的 Brauer 类，故结论仍然成立。
\end{proof}

\begin{lemma}
\label{brauer-lemma-maximal-subfield-splits}
$k$ 上有限维中心除环 $K$ 的任一极大子域都是 $K$ 的分裂域。
\end{lemma}

\begin{proof}
合并引理 \ref{brauer-lemma-maximal-subfield} 与
定理 \ref{brauer-theorem-splitting}。
\end{proof}

\begin{lemma}
\label{brauer-lemma-splitting-field-degree}
设 $K$ 是 $k$ 上的有限维中心除环，并令 $d^2 = [K : k]$。
对 $K$ 的任一有限分裂域 $k'$，次数 $[k' : k]$ 可被 $d$ 整除。
\end{lemma}

\begin{proof}
由定理 \ref{brauer-theorem-splitting}，存在 Brauer 类与 $K$ 相同的
有限维中心单代数 $B$，使得 $[B : k] = [k' : k]^2$。
由
引理 \ref{brauer-lemma-similar}，
对某个 $n$ 有 $B = \text{Mat}(n \times n, K)$。
于是 $[k' : k]^2 = n^2d^2$，结论得证。
\end{proof}

\begin{proposition}
\label{brauer-proposition-separable-splitting-field}
设 $K$ 是 $k$ 上的有限维中心除环。存在极大子域
$k \subset k' \subset K$，使得它在 $k$ 上可分。
% P02-E0001: frozen English source says “spitting field” at brauer.tex:711.
特别地，每个 Brauer 类都有一个有限可分分裂域。
\end{proposition}

\begin{proof}
由于每个 Brauer 类都可由 $k$ 上的有限维中心除环表示，
第二个断言由第一个断言和
引理 \ref{brauer-lemma-maximal-subfield-splits}
立即得到。

\medskip\noindent
为证明第一个断言，设已给定可分子域 $k' \subset K$。
则 $k'$ 在 $K$ 中的中心化子 $K'$ 以 $k'$ 为中心，问题归结为
寻找 $K'$ 的一个在 $k'$ 上可分的极大子域。因此只需证明：
若 $k \not = K$，则存在 $x \in K$、$x \not \in k$，且该元素
在 $k$ 上可分。在特征零时，这是显然的。因此可设 $k$ 的特征为
$p > 0$。若基域 $k$ 有限，结论同样显然（因为有限域的扩张总是
可分的）。故可设 $k$ 是正特征无限域。

\medskip\noindent
% P02-E0002: source statement quantifies over all of K although elements of k are separable over k.
反设 $K$ 中没有元素在 $k$ 上可分，以求矛盾。由
《域》，第 \ref{fields-section-algebraic} 节
中的讨论，这意味着任意 $x \in K$ 的最小多项式都形如
$T^q - a$，其中 $q$ 是 $p$ 的幂且 $a \in k$。
显然，$K$ 中每个元素的最小多项式，其次数都满足
$\leq \dim_k(K)$，故存在一个固定的 $p$ 的幂 $q$，使得对所有
$x \in K$ 都有 $x^q \in k$。

\medskip\noindent
考虑如下映射：
$$
(-)^q : K \longrightarrow K
$$
取 $K$ 的一组 $k$-基 $\{a_1, \ldots, a_n\}$，其中
$a_1 = 1$，用这组基将该映射写成
$$
(\sum x_i a_i)^q = \sum f_i(x_1, \ldots, x_n)a_i.
$$
由于 $K$ 上的乘法是 $k$-双线性的，每个 $f_i$ 都是
$x_1, \ldots, x_n$ 的多项式（细节从略）。
上述 $q$ 的选择以及 $k$ 为无限域这一事实表明，对 $i \geq 2$，
$f_i$ 恒为零。因此，将 $k$ 扩张到其代数闭包
$\overline{k}$ 后，它仍恒为零。但代数
$K \otimes_k \overline{k}$ 是矩阵代数（例如由
引理 \ref{brauer-lemma-base-change} 与
\ref{brauer-lemma-brauer-algebraically-closed}），
这意味着存在某些元素，其 $q$ 次幂不是中心元素
（例如 $e_{11}$）。这就是所需的矛盾。
\end{proof}

\noindent
以上结果使我们能够如下刻画有限维中心单代数。

\begin{lemma}
\label{brauer-lemma-finite-central-simple-algebra}
设 $k$ 为域。对 $k$-代数 $A$，以下条件等价：
\begin{enumerate}
\item $A$ 是有限维中心单 $k$-代数；
\item $A$ 是有限维 $k$-向量空间，$k$ 是 $A$ 的中心，且
$A$ 没有非平凡双边理想；
\item 存在 $d \geq 1$，使得
$A \otimes_k \bar k \cong \text{Mat}(d \times d, \bar k)$；
\item 存在 $d \geq 1$，使得
$A \otimes_k k^{sep} \cong \text{Mat}(d \times d, k^{sep})$；
\item 存在 $d \geq 1$ 及有限 Galois 扩张 $k'/k$，使得
$A \otimes_k k' \cong \text{Mat}(d \times d, k')$；
\item 存在 $n \geq 1$ 及 $k$ 上的有限维中心除环 $K$，使得
$A \cong \text{Mat}(n \times n, K)$。
\end{enumerate}
整数 $d$ 称为 $A$ 的{\it 次数}。
\end{lemma}

\begin{proof}
(1) 与 (2) 的等价性直接来自定义，见第 \ref{brauer-section-algebras} 节。
设 (1) 成立。由
命题 \ref{brauer-proposition-separable-splitting-field}，
$A$ 有一个可分分裂域 $k \subset k'$。
当然，$k'/k$ 的一个 Galois 闭包也是分裂域。
所以 (1) 蕴含 (5)。显然，(5) $\Rightarrow$ (4)
$\Rightarrow$ (3)。设 (3) 成立。则
$A \otimes_k \overline{k}$ 是有限维中心单
$\overline{k}$-代数，例如可由
引理 \ref{brauer-lemma-matrix-algebras} 得到。
这显然蕴含 $A$ 是有限维中心单 $k$-代数。
最后，(1) 与 (6) 的等价性就是 Wedderburn 定理，见
定理 \ref{brauer-theorem-wedderburn}。
\end{proof}
% Generated mechanically from the frozen homology.tex target.
% Do not edit this generated file; edit the bound locale source instead.

% OK, start here.
%

\setcounter{chapter}{11}
\chapter{同调代数}



\phantomsection
\label{homology-section-phantom}


\section{引言}
\label{homology-section-introduction}

\noindent
本文将讲解基础同调代数。我们会依其他各篇所需而继续补充，因为这方面的材料显然浩如烟海。可参见 \cite{Maclane}。

\section{基本概念}
\label{homology-section-topology-basic}

\noindent
下列概念视为基础概念，本文不再定义或证明。这并不意味着它们必然都很容易或广为人知。

\begin{enumerate}
\item 暂无。
\end{enumerate}




\section{预加性范畴与加性范畴}
\label{homology-section-additive-categories}

\noindent
先给出预加性范畴的定义。

\begin{definition}
\label{homology-definition-preadditive}
若范畴 $\mathcal{A}$ 的每个态射集 $\Mor_\mathcal{A}(x, y)$ 都带有阿贝尔群结构，并且复合映射
$$
\Mor(x, y) \times \Mor(y, z)
\longrightarrow
\Mor(x, z)
$$
是双线性的，则称该范畴为{\it 预加性范畴}。预加性范畴之间的函子 $F : \mathcal{A} \to \mathcal{B}$ 称为{\it 加性函子}，当且仅当对所有 $x, y \in \Ob(\mathcal{A})$，映射 $F : \Mor(x, y) \to \Mor(F(x), F(y))$ 都是阿贝尔群同态。
\end{definition}

\noindent
特别地，对每一对 $x, y$，至少存在一个态射 $x \to y$，即零映射。

\begin{lemma}
\label{homology-lemma-preadditive-zero}
设 $\mathcal{A}$ 为预加性范畴，$x$ 为 $\mathcal{A}$ 的对象。下列条件等价：
\begin{enumerate}
\item $x$ 是始对象；
\item $x$ 是终对象；
\item 在 $\Mor_\mathcal{A}(x, x)$ 中有 $\text{id}_x = 0$。
\end{enumerate}
此外，若这样的对象 $0$ 存在，则态射 $\alpha : y \to z$ 可通过 $0$ 分解，当且仅当 $\alpha = 0$。
\end{lemma}

\begin{proof}
先假设 $x$ 是 (1) 始对象或 (2) 终对象。两种情形下，$\Mor(x,x)$ 都是包含 $\text{id}_x$ 的平凡阿贝尔群，故在 $\Mor(x, x)$ 中 $\text{id}_x = 0$。因此 (1) 与 (2) 均蕴含 (3)。

\medskip\noindent
现假设在 $\Mor(x,x)$ 中 $\text{id}_x = 0$。任取 $\mathcal{A}$ 的对象 $y$ 及 $f \in \Mor(x ,y)$。以 $C : \Mor(x,x) \times \Mor(x,y) \to \Mor(x,y)$ 表示复合映射。于是 $f = C(0, f)$；又因 $C$ 双线性，$C(0, f) = 0$。故 $f = 0$，从而 $x$ 是 $\mathcal{A}$ 的始对象。对 $f \in \Mor(y, x)$ 作同样论证，可知 $x$ 也是终对象。因此 (3) 蕴含 (1) 与 (2)。
\end{proof}

\begin{definition}
\label{homology-definition-zero-object}
在预加性范畴 $\mathcal{A}$ 中，同时为终对象与始对象（如引理 \ref{homology-lemma-preadditive-zero}）的对象称为{\it 零对象}，记作 $0$。
\end{definition}

\begin{lemma}
\label{homology-lemma-preadditive-direct-sum}
设 $\mathcal{A}$ 为预加性范畴，$x, y \in \Ob(\mathcal{A})$。若积 $x \times y$ 存在，则余积 $x \amalg y$ 也存在；若余积 $x \amalg y$ 存在，则积 $x \times y$ 也存在。在这种情形下还有 $x \amalg y \cong x \times y$。
\end{lemma}

\begin{proof}
设 $z = x \times y$，投影为 $p : z \to x$ 与 $q : z \to y$。以 $i : x \to z$ 表示对应于 $(1, 0)$ 的态射，以 $j : y \to z$ 表示对应于 $(0, 1)$ 的态射。于是有交换图
$$
\xymatrix{
x \ar[rr]^1 \ar[rd]^i & & x \\
& z \ar[ru]^p \ar[rd]^q & \\
y \ar[rr]^1 \ar[ru]^j & & y
}
$$
其中两条对角线上的复合均为零。由此可知 $i \circ p + j \circ q : z \to z$ 是恒等态射，因为它与 $p$ 复合得到 $p$，与 $q$ 复合得到 $q$。
给定态射 $a : x \to w$ 与 $b : y \to w$，可构造映射 $a \circ p + b \circ q : z \to w$。这样便得到双射 $\Mor(z, w) = \Mor(x, w) \times \Mor(y, w)$，从而说明 $z = x \amalg y$。

\medskip\noindent
若从余积 $x \amalg y$ 而非积出发，我们把态射 $p, q$ 的构造留给读者。
\end{proof}

\begin{definition}
\label{homology-definition-direct-sum}
给定预加性范畴 $\mathcal{A}$ 中的一对对象 $x, y$，$x$ 与 $y$ 的
{\it 直和} $x \oplus y$ 是直积 $x \times y$，并配备引理
\ref{homology-lemma-preadditive-direct-sum} 中的态射 $i, j, p, q$。
\end{definition}

\begin{remark}
\label{homology-remark-direct-sum}
注意，引理 \ref{homology-lemma-preadditive-direct-sum} 的证明表明：给定 $p$ 与 $q$ 后，态射 $i$, $j$ 由关系 $p \circ i = \text{id}_x$、$q \circ j = \text{id}_y$、$p \circ j = 0$、$q \circ i = 0$ 唯一确定。此外自动有 $i \circ p + j \circ q = \text{id}_{x \oplus y}$。类似地，给定 $i$, $j$ 后，态射 $p$ 与 $q$ 也唯一确定。

最后，若给定对象 $x, y, z$ 以及态射 $i : x \to z$、$j : y \to z$、$p : z \to x$、$q : z \to y$，并满足 $p \circ i = \text{id}_x$、$q \circ j = \text{id}_y$、$p \circ j = 0$、$q \circ i = 0$ 及 $i \circ p + j \circ q = \text{id}_z$，则 $z$ 是 $x$ 与 $y$ 的直和，其四个结构态射分别为 $i, j, p, q$。
\end{remark}

\begin{lemma}
\label{homology-lemma-additive-additive}
设 $\mathcal{A}$、$\mathcal{B}$ 为预加性范畴，$F : \mathcal{A} \to \mathcal{B}$ 为加性函子。则 $F$ 将直和变为直和，并将零对象变为零对象。
\end{lemma}

\begin{proof}
设 $F$ 为加性函子。依注记 \ref{homology-remark-direct-sum}，$x$ 与 $y$ 的直和 $z$ 可由态射 $i : x \to z$、$j : y \to z$、$p : z \to x$、$q : z \to y$ 及关系 $p \circ i = \text{id}_x$、$q \circ j = \text{id}_y$、$p \circ j = 0$、$q \circ i = 0$、$i \circ p + j \circ q = \text{id}_z$ 刻画。显然，$F(x), F(y), F(z)$ 与态射 $F(i), F(j), F(p), F(q)$ 满足完全相同的关系（由加性），故 $F(z)$ 是 $F(x)$ 与 $F(y)$ 的直和。因此，$F$ 将直和变为直和。

\medskip\noindent
要看出 $F$ 将零对象变为零对象，使用引理 \ref{homology-lemma-preadditive-zero} 中零对象的刻画 (3) 即可。
\end{proof}

\begin{definition}
\label{homology-definition-additive-category}
若范畴 $\mathcal{A}$ 是预加性的且有限积存在，换言之，它有零对象和直和，则称它为{\it 加性范畴}。
\end{definition}

\noindent
事实上，空积是有限积；若其存在，它就是终对象。

\begin{definition}
\label{homology-definition-kernel}
设 $\mathcal{A}$ 为预加性范畴，$f : x \to y$ 为态射。
\begin{enumerate}
\item $f$ 的一个{\it 核}是态射 $i : z \to x$，满足：(a) $f \circ i = 0$；(b) 对任何满足 $f \circ i' = 0$ 的 $i' : z' \to x$，存在唯一态射 $g : z' \to z$ 使得 $i' = i \circ g$。
\item 若 $f$ 的核存在，则将其记作 $\Ker(f) \to x$。
\item $f$ 的一个{\it 余核}是态射 $p : y \to z$，满足：(a) $p \circ f = 0$；(b) 对任何满足 $p' \circ f = 0$ 的 $p' : y \to z'$，存在唯一态射 $g : z \to z'$ 使得 $p' = g \circ p$。
\item 若 $f$ 的余核存在，则将其记作 $y \to \Coker(f)$。
\item 若 $f$ 的核存在，则 $f$ 的一个{\it 余像}是态射 $\Ker(f) \to x$ 的余核。
\item 若核和余像均存在，则将其记作 $x \to \Coim(f)$。
\item 若 $f$ 的余核存在，则 $f$ 的{\it 像}是态射 $y \to \Coker(f)$ 的核。
\item 若 $f$ 的余核和像均存在，则将其记作 $\Im(f) \to y$。
\end{enumerate}
\end{definition}

\noindent
在上述定义中，我们谈到“该核”和“该余核”时，默认使用了它们在唯一同构意义下的唯一性。这由 Yoneda 引理推出（《范畴》，第 \ref{categories-section-opposite} 节），因为 $f : x \to y$ 的核表示这样一个函子：它把对象 $z$ 送到集合 $\Ker(\Mor_\mathcal{A}(z, x) \to \Mor_\mathcal{A}(z, y))$。余核的情形与之对偶。

\medskip\noindent
首先如下说明直和与核之间的关系。

\begin{lemma}
\label{homology-lemma-additive-cat-biproduct-kernel}
设 $\mathcal{C}$ 为预加性范畴。设 $x \oplus y$ 连同态射 $i, j, p, q$ 如引理 \ref{homology-lemma-preadditive-direct-sum} 所述，是 $\mathcal{C}$ 中的直和。则 $i : x \to x \oplus y$ 是 $q : x \oplus y \rightarrow y$ 的核。对偶地，$p$ 是 $j$ 的余核。
\end{lemma}

\begin{proof}
设 $f : z' \to x \oplus y$ 满足 $q \circ f = 0$。须证明存在唯一态射 $g : z' \to x$ 使 $f = i \circ g$。由于 $i \circ p + j \circ q$ 是 $x \oplus y$ 上的恒等态射，故
$$
f = (i \circ p + j \circ q) \circ f = i \circ p \circ f
$$
所以取 $g = p \circ f$ 即可。唯一性来自 $p \circ i$ 是 $x$ 上的恒等态射。第二个断言的证明与此对偶。
\end{proof}

\begin{lemma}
\label{homology-lemma-kernel-mono}
设 $\mathcal{C}$ 为预加性范畴，$f : x \to y$ 为 $\mathcal{C}$ 中的态射。
\begin{enumerate}
\item 若 $f$ 的核存在，则该核是单态射。
\item 若 $f$ 的余核存在，则该余核是满态射。
\item 若 $f$ 的核和余像存在，则该余像是满态射。
\item 若 $f$ 的余核和像存在，则该像是单态射。
\end{enumerate}
\end{lemma}

\begin{proof}
(1) 直接来自核定义中的唯一性。(2) 的证明与之对偶。(3) 由 (2) 推出，因为余像是一个余核；类似地，(4) 由 (1) 推出。
\end{proof}

\begin{lemma}
\label{homology-lemma-coim-im-map}
设 $f : x \to y$ 是预加性范畴中的态射，并且其核、余核、像和余像均存在。则 $f$ 唯一分解为
$x \to \Coim(f) \to \Im(f) \to y$。
\end{lemma}

\begin{proof}
因为 $\Ker(f) \to x \to y$ 为零，所以有典范态射 $\Coim(f) \to y$。复合 $\Coim(f) \to y \to \Coker(f)$ 为零，因为它是唯一诱导出零态射 $x \to y \to \Coker(f)$ 的态射（唯一性由引理 \ref{homology-lemma-kernel-mono} (3) 得到）。因此，$\Coim(f) \to y$ 唯一地通过 $\Im(f) \to y$ 分解，从而得到所需映射。
\end{proof}

\begin{example}
\label{homology-example-not-abelian}
设 $k$ 为域。考虑 $k$ 上滤过向量空间所成的范畴（见定义 \ref{homology-definition-filtered}）。考虑滤过向量空间 $(V, F)$ 与 $(W, F)$，其中 $V = W = k$ 且
$$
F^iV
=
\left\{
\begin{matrix}
V & \text{若} & i < 0 \\
0 & \text{若} & i \geq 0
\end{matrix}
\right.
\text{ 且 }
F^iW
=
\left\{
\begin{matrix}
W & \text{若} & i \leq 0 \\
0 & \text{若} & i > 0
\end{matrix}
\right.
$$
底层向量空间上的 $\text{id}_k$ 所对应的映射 $f : V \to W$ 具有平凡的核与余核，却不是同构。还要注意，$\Coim(f) = V$ 且 $\Im(f) = W$。这说明 $k$ 上滤过向量空间所成的范畴不是阿贝尔范畴。
\end{example}

\noindent
最后以若干关于分裂幂等自同态之像的引理结束本节。回顾：若自同态 $f$ 满足 $f \circ f = f$，则称其为幂等的。
% P02-E0003: frozen source line 354 has “endormorphisms”.

\begin{remark}
\label{homology-remark-idempotent-symmetry}
设 $\mathcal{A}$ 为预加性范畴，$x$ 为 $\mathcal{A}$ 的对象，$f : x \to x$ 为幂等态射。令 $g = \text{id} _ x - f$，则 $f \circ g = 0$、$g \circ f = 0$ 且 $g \circ g = g$。
\end{remark}

\begin{lemma}
\label{homology-lemma-idempotent-kernel-cokernel}
设 $\mathcal{A}$ 为预加性范畴，$x$ 为 $\mathcal{A}$ 的对象，$f : x \to x$ 为幂等态射。令 $g = \text{id} _ x - f$；由注记 \ref{homology-remark-idempotent-symmetry}，有 $f \circ g = g \circ f = 0$。
\begin{enumerate}
\item 若 $f$ 有核 $i : \Ker (f) \to x$，且 $p$ 是满足 $g = i \circ p$ 的唯一态射，则 $p \circ i = \text{id} _ {\Ker (f)}$，并且 $p$ 是 $f$ 的余核。
\item 若 $f$ 有余核 $p : x \to \Coker (f)$，且 $i$ 是满足 $g = i \circ p$ 的唯一态射，则 $p \circ i = \text{id} _ {\Coker (f)}$，并且 $i$ 是 $f$ 的核。
\item 若 $g$ 有核 $j : \Ker (g) \to x$，且 $q$ 是满足 $f = j \circ q$ 的唯一态射，则 $q \circ j = \text{id} _ {\Ker (g)}$，并且 $q$ 是 $g$ 的余核。
\item 若 $g$ 有余核 $q : x \to \Coker (g)$，且 $j$ 是满足 $f = j \circ q$ 的唯一态射，则 $q \circ j = \text{id} _ {\Coker (g)}$，并且 $j$ 是 $g$ 的核。
\end{enumerate}
\end{lemma}

\begin{proof}
先证 (1)。首先计算 $i = (f + g) \circ i = g \circ i$。由 $p$ 的构造，这蕴含
$i \circ p \circ i = g \circ i = i \circ \text{id} _ {\Ker (f)}$。
由引理 \ref{homology-lemma-kernel-mono} 得 $p \circ i = \text{id} _ {\Ker (f)}$。

\medskip\noindent
尚须证明 $p : x \to \Ker (f)$ 满足 $f$ 的余核之泛性质：即给定满足 $a \circ f = 0$ 的态射 $a : x \to y$，存在唯一态射 $b : \Ker (f) \to y$ 使 $a = b \circ p$。先计算 $a = a \circ (f + g) = a \circ g$；再由 $p$ 的构造得到 $a = a \circ g = a \circ i \circ p$，故存在性可取 $b = a \circ i$。唯一性来自 $b = b \circ p \circ i = a \circ i$，即 $a$ 决定 $b$。

\medskip\noindent
分别借助对偶性、对称性（即 $g$ 本身幂等；参见注记 \ref{homology-remark-idempotent-symmetry}）以及二者兼用，(2)、(3)、(4) 可按 (1) 同样推出。
\end{proof}

\begin{lemma}
\label{homology-lemma-idempotent-splitting}
设 $\mathcal{A}$ 为预加性范畴，$x$ 为 $\mathcal{A}$ 的对象，$f : x \to x$ 为幂等态射。令 $g = \text{id} _ x - f$。若 $f$ 与 $g$ 各自至少有核或余核之一，并依引理 \ref{homology-lemma-idempotent-kernel-cokernel} 的陈述构造态射 $i , j , p , q$，则四个（余）核全部存在，且
\begin{align*}
x
& \simeq \Ker (f) \oplus \Ker (g) \\
& \simeq \Coker (f) \oplus \Ker (g) \\
& \simeq \Ker (f) \oplus \Coker (g) \\
& \simeq \Coker (f) \oplus \Coker (g),
\end{align*}
其中每个直积的结构均由 $i , j , p , q$ 给出。
\end{lemma}

\begin{proof}
由引理 \ref{homology-lemma-idempotent-kernel-cokernel}，
$p \circ i = \text{id} _ {\Ker (f)} = \text{id} _ {\Coker (f)}$，且 $q \circ j = \text{id} _ {\Ker (g)} = \text{id} _ {\Coker (g)}$。由构造，$i \circ p + j \circ q = f + g = \text{id} _ x$。结论由注记 \ref{homology-remark-direct-sum} 得出。
\end{proof}

\begin{lemma}
\label{homology-lemma-split-morphism-kernel-cokernel}
设 $\mathcal{A}$ 为预加性范畴，$x$ 与 $y$ 为 $\mathcal{A}$ 的对象，态射 $j : y \to x$ 与 $q : x \to y$ 满足 $q \circ j = \text{id} _ y$。则
\begin{enumerate}
\item $\text{id} _ x - j \circ q$ 以 $j$ 为核；
\item $\text{id} _ x - j \circ q$ 以 $q$ 为余核。
\end{enumerate}
\end{lemma}

\begin{proof}
先证 (1)。首先计算 $(\text{id} _ x - j \circ q) \circ j = 0$。须证明 $j$ 满足 $\text{id} _ {x} - j \circ q$ 的核之泛性质：即给定满足 $(\text{id} _ x - j \circ q) \circ a = 0$ 的态射 $a : z \to x$，存在唯一态射 $b : z \to y$ 使 $a = j \circ b$。由 $a = j \circ q \circ a$ 可知存在性成立，可取 $b = q \circ a$。由 $b = q \circ j \circ b = q \circ a$ 可知唯一性成立，即 $a$ 决定 $b$。

\medskip\noindent
(2) 由对偶性按 (1) 推出。
\end{proof}

\begin{lemma}
\label{homology-lemma-split-morphism-splitting}
设 $\mathcal{A}$ 为预加性范畴，$x$ 与 $y$ 为 $\mathcal{A}$ 的对象，态射 $j : y \to x$ 与 $q : x \to y$ 满足 $q \circ j = \text{id} _ y$。若 $j \circ q$ 至少有核或余核之一，则两者均存在，且
\begin{align*}
x
& \simeq \Ker (j \circ q) \oplus y \\
& \simeq \Coker (j \circ q) \oplus y,
\end{align*}
其中任一直积的结构都包含 $j , q$ 与典范（余）核态射，而剩余映射则唯一确定。
\end{lemma}

\begin{proof}
令 $f = j \circ q$，计算得 $f \circ f = f$，再应用引理 \ref{homology-lemma-idempotent-splitting} 与引理 \ref{homology-lemma-split-morphism-kernel-cokernel}。
\end{proof}

\section{Karoubian 范畴}
\label{homology-section-karoubian}

\noindent
初次阅读时可略过本节。

\begin{definition}
\label{homology-definition-karoubian}
设 $\mathcal{C}$ 为预加性范畴。若 $\mathcal{C}$ 中每个对象的每个幂等自同态都有核，则称 $\mathcal{C}$ 为 {\it Karoubian 范畴}。
\end{definition}

\noindent
对偶概念应是每个对象的每个幂等自同态都有余核。然而，由下述引理（的对偶）可知，这给出等价的概念。

\begin{lemma}
\label{homology-lemma-karoubian}
设 $\mathcal{C}$ 为预加性范畴。下列条件等价：
\begin{enumerate}
\item $\mathcal{C}$ 是 Karoubian 范畴；
\item $\mathcal{C}$ 中每个对象的每个幂等自同态都有余核；
\item 给定 $\mathcal{C}$ 中的幂等自同态 $p : z \to z$，存在直和分解 $z = x \oplus y$，使 $p$ 对应于到 $y$ 上的投影。
\end{enumerate}
\end{lemma}

\begin{proof}
由引理 \ref{homology-lemma-idempotent-kernel-cokernel}、\ref{homology-lemma-idempotent-splitting} 和 \ref{homology-lemma-additive-cat-biproduct-kernel} 立即得到。
\end{proof}

\begin{lemma}
\label{homology-lemma-projectors-have-images}
设 $\mathcal{D}$ 为预加性范畴。
\begin{enumerate}
\item 若 $\mathcal{D}$ 有可数积，且有右逆的映射均有核，则 $\mathcal{D}$ 是 Karoubian 范畴。
\item 若 $\mathcal{D}$ 有可数余积，且有左逆的映射均有余核，则 $\mathcal{D}$ 是 Karoubian 范畴。
\end{enumerate}
\end{lemma}

\begin{proof}
设 $X$ 为 $\mathcal{D}$ 的对象，$e : X \to X$ 为幂等态射。函子
$$
W \longmapsto \Ker(
\Mor_\mathcal{D}(W, X)
\xrightarrow{e}
\Mor_\mathcal{D}(W, X)
)
$$
可表示，当且仅当 $e$ 有核。注意，对任意阿贝尔群 $A$ 与幂等自同态 $e : A \to A$，都有
$$
\Ker(e : A \to A)
= \Ker(\Phi :
\prod\nolimits_{n \in \mathbf{N}} A
\to
\prod\nolimits_{n \in \mathbf{N}} A
)
$$
其中
$$
\Phi(a_1, a_2, a_3, \ldots) = (ea_1 + (1 - e)a_2, ea_2 + (1 - e)a_3, \ldots)
$$
此外，$\Phi$ 有右逆
$$
\Psi(a_1, a_2, a_3, \ldots) =
(a_1, (1 - e)a_1 + ea_2, (1 - e)a_2 + ea_3, \ldots).
$$
故 (1) 成立。(2) 的证明与之对偶（使用 Karoubian 范畴的对偶定义，即引理 \ref{homology-lemma-karoubian} 的条件 (2)）。
\end{proof}








\section{阿贝尔范畴}
\label{homology-section-abelian-categories}

\noindent
阿贝尔范畴是满足恰好足以使蛇引理成立的一组公理的范畴。有一条公理有时会被忘记：引理 \ref{homology-lemma-coim-im-map} 中的典范映射 $\Coim(f) \to \Im(f)$ 总是同构。例 \ref{homology-example-not-abelian} 表明这条公理确有必要。

\begin{definition}
\label{homology-definition-abelian-category}
若范畴 $\mathcal{A}$ 是加性范畴，所有核和余核均存在，并且对其中的每个态射 $f$，自然映射 $\Coim(f) \to \Im(f)$ 都是同构，则称 $\mathcal{A}$ 为{\it 阿贝尔范畴}。
\end{definition}

\begin{lemma}
\label{homology-lemma-abelian-opposite}
设 $\mathcal{A}$ 为预加性范畴。态射集上的加法使 $\mathcal{A}^{opp}$ 成为预加性范畴。此外，$\mathcal{A}$ 是加性范畴当且仅当 $\mathcal{A}^{opp}$ 是加性范畴；$\mathcal{A}$ 是阿贝尔范畴当且仅当 $\mathcal{A}^{opp}$ 是阿贝尔范畴。
\end{lemma}

\begin{proof}
第一个断言直接可得。为说明 $\mathcal{A}$ 是加性范畴当且仅当 $\mathcal{A}^{opp}$ 是加性范畴，回顾加性可由零对象和直和的存在性刻画，而二者在取对偶范畴时均得以保留。最后，为说明 $\mathcal{A}$ 是阿贝尔范畴当且仅当 $\mathcal{A}^{opp}$ 是阿贝尔范畴，只需观察 $\mathcal{A}^{opp}$ 中的核、余核、像和余像分别对应于 $\mathcal{A}$ 中的余核、核、余像和像。
% P02-E0004: frozen source line 634 has “observes that”.
\end{proof}

\begin{definition}
\label{homology-definition-injective-surjective}
设 $f : x \to y$ 为阿贝尔范畴中的态射。
\begin{enumerate}
\item 若 $\Ker(f) = 0$，则称 $f$ 为{\it 单射态射}。
\item 若 $\Coker(f) = 0$，则称 $f$ 为{\it 满射态射}。
\end{enumerate}
若 $x \to y$ 为单射态射，则称 $x$ 为 $y$ 的{\it 子对象}，并记作 $x \subset y$。若 $x \to y$ 为满射态射，则称 $y$ 为 $x$ 的一个{\it 商}。
\end{definition}

\begin{lemma}
\label{homology-lemma-characterize-injective}
设 $f : x \to y$ 为阿贝尔范畴 $\mathcal{A}$ 中的态射。则
\begin{enumerate}
\item $f$ 为单射态射，当且仅当 $f$ 为单态射；
\item $f$ 为满射态射，当且仅当 $f$ 为满态射；
\item $f$ 为同构，当且仅当 $f$ 既为单射态射又为满射态射。
\end{enumerate}
\end{lemma}

\begin{proof}
证明 (1)。回顾 $\Ker(f)$ 表示把 $z$ 送到 $\Ker(\Mor_\mathcal{A}(z, x) \to \Mor_\mathcal{A}(z, y))$ 的函子，见定义 \ref{homology-definition-kernel}。因此，$\Ker(f)$ 为 $0$，当且仅当对所有 $z$，$\Mor_\mathcal{A}(z, x) \to \Mor_\mathcal{A}(z, y)$ 都是单射，也当且仅当 $f$ 是单态射。(2) 的证明类似。对 (3)，同构既是单态射又是满态射，由 (1)、(2) 可知 $f$ 既为单射态射又为满射态射。反之，若 $f$ 既为单射态射又为满射态射，则 $x = \Coim(f)$ 且 $y = \Im(f)$，所以 $f$ 是同构。
\end{proof}

\noindent
在阿贝尔范畴中，若 $x \subset y$ 为子对象，则记
$$
y/x = \Coker(x \to y).
$$

\begin{lemma}
\label{homology-lemma-colimit-abelian-category}
设 $\mathcal{A}$ 为阿贝尔范畴。$\mathcal{A}$ 中所有有限极限与有限余极限都存在。
\end{lemma}

\begin{proof}
要证明有限极限存在，只需证明有限积与等化子存在，见《范畴》，引理 \ref{categories-lemma-finite-limits-exist}。有限积由定义存在，而 $a, b : x \to y$ 的等化子是 $a - b$ 的核。有限余极限的论证与此类似且对偶。
\end{proof}

\begin{example}
\label{homology-example-fibre-product-pushouts}
设 $\mathcal{A}$ 为阿贝尔范畴。$\mathcal{A}$ 中推出与纤维积有如下简单描述：
\begin{enumerate}
\item 若 $a : x \to y$、$b : z \to y$ 是 $\mathcal{A}$ 中的态射，则纤维积为 $x \times_y z = \Ker((a, -b) : x \oplus z \to y)$。
\item 若 $a : y \to x$、$b : y \to z$ 是 $\mathcal{A}$ 中的态射，则推出为 $x \amalg_y z = \Coker((a, -b) : y \to x \oplus z)$。
\end{enumerate}
\end{example}

\begin{definition}
\label{homology-definition-exact}
设 $\mathcal{A}$ 为加性范畴。考虑 $\mathcal{A}$ 中的态射序列
$$
\ldots \to x \to y \to z \to \ldots
\quad\text{或}\quad
x_1 \to x_2 \to \ldots \to x_n
$$
若任意两个相邻（画出的）箭头之复合为零，则称这样的序列为{\it 复形}。若 $\mathcal{A}$ 为阿贝尔范畴，则称上面第一类复形在 $y$ 处{\it 正合}，若 $\Im(x \to y) = \Ker(y \to z)$；称第二类复形在 $x_i$ 处（其中 $1 < i < n$）{\it 正合}，若 $\Im(x_{i - 1} \to x_i) = \Ker(x_i \to x_{i + 1})$。若上述序列是复形并且在每个对象处（第一种情形）或在所有满足 $1 < i < n$ 的 $x_i$ 处（第二种情形）正合，则称其{\it 正合}，或称为{\it 正合序列}、{\it 正合复形}。
% P02-E0005: frozen source lines 729-730 read “We a sequence as above is”.
对如下形状的序列也有相应的概念：
$$
\ldots \to x_{-3} \to x_{-2} \to x_{-1}
\quad\text{且}\quad
x_1 \to x_2 \to x_3 \to \ldots
$$
形如
$$
0 \to A  \to B \to C \to 0.
$$
的正合复形称为{\it 短正合序列}。
\end{definition}

\noindent
在下述引理中，我们假定读者知道阿贝尔群序列正合的含义。

\begin{lemma}
\label{homology-lemma-check-exactness}
设 $\mathcal{A}$ 为阿贝尔范畴。设 $0 \to M_1 \to M_2 \to M_3 \to 0$ 为 $\mathcal{A}$ 中的复形。
\begin{enumerate}
\item $M_1 \to M_2 \to M_3 \to 0$ 正合，当且仅当对 $\mathcal{A}$ 的所有对象 $N$，阿贝尔群序列
$$
0 \to \Hom_\mathcal{A}(M_3, N) \to
\Hom_\mathcal{A}(M_2, N) \to \Hom_\mathcal{A}(M_1, N)
$$
正合；
\item $0 \to M_1 \to M_2 \to M_3$ 正合，当且仅当对 $\mathcal{A}$ 的所有对象 $N$，阿贝尔群序列
$$
0 \to \Hom_\mathcal{A}(N, M_1) \to \Hom_\mathcal{A}(N, M_2) \to
\Hom_\mathcal{A}(N, M_3)
$$
正合。
\end{enumerate}
\end{lemma}

\begin{proof}
略。提示：见《代数》，引理 \ref{algebra-lemma-hom-exact}。
\end{proof}

\begin{definition}
\label{homology-definition-ses-split}
设 $\mathcal{A}$ 为阿贝尔范畴，$i : A \to B$ 与 $q : B \to C$ 为 $\mathcal{A}$ 中的态射，使 $0 \to A \to B \to C \to 0$ 成为短正合序列。若存在态射 $j : C \to B$ 与 $p : B \to A$，使 $(B, i, j, p, q)$ 是 $A$ 与 $C$ 的直和，则称该短正合序列{\it 分裂}。
\end{definition}

\begin{lemma}
\label{homology-lemma-ses-split}
设 $\mathcal{A}$ 为阿贝尔范畴，$0 \to A \to B \to C \to 0$ 为短正合序列。
\begin{enumerate}
\item 给定 $B \to C$ 的右逆态射 $s : C \to B$，存在唯一 $\pi : B \to A$，使 $(s, \pi)$ 按定义 \ref{homology-definition-ses-split} 分裂该短正合序列。
\item 给定 $A \to B$ 的左逆态射 $\pi : B \to A$，存在唯一 $s : C \to B$，使 $(s, \pi)$ 按定义 \ref{homology-definition-ses-split} 分裂该短正合序列。
\end{enumerate}
\end{lemma}

\begin{proof}
这是引理 \ref{homology-lemma-split-morphism-splitting} 的直接推论。
\end{proof}

\begin{lemma}
\label{homology-lemma-characterize-cartesian}
设 $\mathcal{A}$ 为阿贝尔范畴。设
$$
\xymatrix{
w\ar[r]^f\ar[d]_g
& y\ar[d]^h\\
x\ar[r]^k
& z
}
$$
为交换图。
\begin{enumerate}
\item 该图是笛卡尔图，当且仅当
$$
0 \to w \xrightarrow{(g, f)} x \oplus y \xrightarrow{(k, -h)} z
$$
正合。
\item 该图是余笛卡尔图，当且仅当
$$
w \xrightarrow{(g, -f)} x \oplus y \xrightarrow{(k, h)} z \to 0
$$
正合。
\end{enumerate}
\end{lemma}

\begin{proof}
令 $u = (g, f) : w \to x \oplus y$，$v = (k, -h) : x \oplus y \to z$。令 $p : x \oplus y \to x$ 与 $q : x \oplus y \to y$ 为典范投影，$i : \Ker(v) \to x \oplus y$ 为典范单射。由例 \ref{homology-example-fibre-product-pushouts}，该图为笛卡尔图，当且仅当存在同构 $r : \Ker(v) \to w$，满足 $f \circ r = q \circ i$ 且 $g \circ r = p \circ i$。序列 $0 \to w \overset{u} \to x \oplus y \overset{v} \to z$ 正合，当且仅当存在同构 $r : \Ker(v) \to w$ 使 $u \circ r = i$。但给定 $r : \Ker(v) \to w$，有 $f \circ r = q \circ i$ 且 $g \circ r = p \circ i$，当且仅当 $q \circ u \circ r= f \circ r = q \circ i$ 且 $p \circ u \circ r = g \circ r = p \circ i$，亦即当且仅当 $u \circ r = i$。这证明了 (1)，而 (2) 由对偶性得出。
\end{proof}

\begin{lemma}
\label{homology-lemma-cartesian-kernel}
设 $\mathcal{A}$ 为阿贝尔范畴。设
$$
\xymatrix{
w\ar[r]^f\ar[d]_g
& y\ar[d]^h\\
x\ar[r]^k
& z
}
$$
为交换图。
\begin{enumerate}
\item 若该图为笛卡尔图，则由 $g$ 诱导的态射 $\Ker(f)\to\Ker(k)$ 是同构。
\item 若该图为余笛卡尔图，则由 $h$ 诱导的态射 $\Coker(f)\to\Coker(k)$ 是同构。
\end{enumerate}
\end{lemma}

\begin{proof}
设该图为笛卡尔图。令 $e:\Ker(f)\to\Ker(k)$ 为 $g$ 诱导的态射，$i:\Ker(f)\to w$ 与 $j:\Ker(k)\to x$ 为典范单射。存在 $t:\Ker(k)\to w$，满足 $f\circ t=0$ 且 $g\circ t=j$。因此存在 $u:\Ker(k)\to\Ker(f)$，使 $i\circ u=t$。于是
$g\circ i\circ u\circ e=g\circ t\circ e=j\circ e=g\circ i$ 且 $f\circ i\circ u\circ e=0=f\circ i$，故 $i\circ u\circ e=i$。因 $i$ 为单态射，$u\circ e=\text{id}_{\Ker(f)}$。又有 $j\circ e\circ u=g\circ i\circ u=g\circ t=j$；因 $j$ 为单态射，$e\circ u=\text{id}_{\Ker(k)}$。这证明 (1)，(2) 由对偶性得出。
\end{proof}

\begin{lemma}
\label{homology-lemma-cartesian-cocartesian}
设 $\mathcal{A}$ 为阿贝尔范畴。设
$$
\xymatrix{
w\ar[r]^f\ar[d]_g
& y\ar[d]^h\\
x\ar[r]^k
& z
}
$$
为交换图。
\begin{enumerate}
\item 若该图为笛卡尔图且 $k$ 为满态射，则该图为余笛卡尔图，且 $f$ 为满态射。
\item 若该图为余笛卡尔图且 $g$ 为单态射，则该图为笛卡尔图，且 $h$ 为单态射。
\end{enumerate}
\end{lemma}

\begin{proof}
设该图为笛卡尔图且 $k$ 为满态射。令 $u = (g, f) : w \to x \oplus y$，$v = (k, -h) : x \oplus y \to z$。因为 $k$ 为满态射，$v$ 也为满态射。因此结合引理 \ref{homology-lemma-characterize-cartesian}，序列 $0\to w\overset{u}\to x\oplus y\overset{v}\to z\to 0$ 正合。故由引理 \ref{homology-lemma-characterize-cartesian}，该图为余笛卡尔图。最后，由引理 \ref{homology-lemma-cartesian-kernel} 与引理 \ref{homology-lemma-characterize-injective}，$f$ 为满态射。这证明 (1)，(2) 由对偶性得出。
\end{proof}

\begin{lemma}
\label{homology-lemma-epimorphism-universal-abelian-category}
设 $\mathcal{A}$ 为阿贝尔范畴。
\begin{enumerate}
\item 若 $x \to y$ 为满射态射，则对每个 $z \to y$，投影 $x \times_y z \to z$ 为满射态射。
\item 若 $x \to y$ 为单射态射，则对每个 $x \to z$，态射 $z \to z \amalg_x y$ 为单射态射。
\end{enumerate}
\end{lemma}

\begin{proof}
立即由引理 \ref{homology-lemma-characterize-injective} 与引理 \ref{homology-lemma-cartesian-cocartesian} 得到。
\end{proof}

\begin{lemma}
\label{homology-lemma-check-exactness-fibre-product}
设 $\mathcal{A}$ 为阿贝尔范畴。设 $f:x\to y$ 与 $g:y\to z$ 满足 $g\circ f=0$。则下列陈述等价：
\begin{enumerate}
\item 序列 $x\overset{f}\to y\overset{g}\to z$ 正合。
\item 对每个满足 $g\circ h=0$ 的 $h:w\to y$，存在对象 $v$、满态射 $k:v\to w$ 与态射 $l:v\to x$，使 $h\circ k=f\circ l$。
\end{enumerate}
\end{lemma}

\begin{proof}
令 $i:\Ker(g)\to y$ 为典范单射，$p:x\to\Coim(f)$ 为典范投影，$j:\Im(f)\to\Ker(g)$ 为典范单射。

\medskip\noindent
假设 (1) 成立。设 $h:w\to y$ 且 $g\circ h=0$。存在 $c:w\to\Ker(g)$ 使 $i\circ c=h$。令 $v=x\times_{\Ker(g)}w$，其典范投影为 $k:v\to w$ 与 $l:v\to x$，于是 $c\circ k=j\circ p\circ l$。故 $h\circ k=i\circ c\circ k=i\circ j\circ p\circ l=f\circ l$。由假设，$j\circ p$ 为满态射，故由引理 \ref{homology-lemma-cartesian-cocartesian}，$k$ 为满态射。这蕴含 (2)。

\medskip\noindent
假设 (2) 成立。则 $g\circ i=0$。故存在对象 $w$、满态射 $k:w\to\Ker(g)$ 与态射 $l:w\to x$，使 $f\circ l=i\circ k$。于是 $i\circ j\circ p\circ l=f\circ l=i\circ k$。由于 $i$ 为单态射，$j\circ p\circ l=k$ 为满态射。因此 $j$ 为满态射，从而是同构。这蕴含 (1)。
% P02-E0006: frozen source line 965 has “an epimorphisms”.
\end{proof}

\begin{lemma}
\label{homology-lemma-exact-kernel-sequence}
设 $\mathcal{A}$ 为阿贝尔范畴。设
$$
\xymatrix{
x \ar[r]^f \ar[d]^\alpha &
y \ar[r]^g \ar[d]^\beta &
z \ar[d]^\gamma\\
u \ar[r]^k & v \ar[r]^l & w
}
$$
为交换图。
\begin{enumerate}
\item 若第一行正合且 $k$ 为单态射，则诱导序列 $\Ker(\alpha) \to \Ker(\beta) \to \Ker(\gamma)$ 正合。
\item 若第二行正合且 $g$ 为满态射，则诱导序列 $\Coker(\alpha) \to \Coker(\beta) \to \Coker(\gamma)$ 正合。
\end{enumerate}
\end{lemma}

\begin{proof}
设第一行正合且 $k$ 为单态射。令 $a:\Ker(\alpha)\to\Ker(\beta)$ 与 $b:\Ker(\beta)\to\Ker(\gamma)$ 为诱导态射；令 $h:\Ker(\alpha)\to x$、$i:\Ker(\beta)\to y$、$j:\Ker(\gamma)\to z$ 为典范单射。因 $j$ 为单态射，有 $b\circ a=0$。设 $c:s\to\Ker(\beta)$ 且 $b\circ c=0$。则 $g\circ i\circ c=j\circ b\circ c=0$。由引理 \ref{homology-lemma-check-exactness-fibre-product}，存在对象 $t$、满态射 $d:t\to s$ 与态射 $e:t\to x$，使 $i\circ c\circ d=f\circ e$。于是
$k\circ \alpha\circ e=\beta\circ f\circ e=\beta\circ i\circ c\circ d=0$。因 $k$ 为单态射，得 $\alpha\circ e=0$。故存在 $m:t\to\Ker(\alpha)$ 使 $h\circ m=e$。于是
$i\circ a\circ m=f\circ h\circ m=f\circ e=i\circ c\circ d$。因 $i$ 为单态射，$a\circ m=c\circ d$。因此引理 \ref{homology-lemma-check-exactness-fibre-product} 蕴含 (1)，而 (2) 由对偶性得出。
\end{proof}

\begin{lemma}
\label{homology-lemma-snake}
设 $\mathcal{A}$ 为阿贝尔范畴。设
$$
\xymatrix{
& x \ar[r]^f \ar[d]^\alpha &
y \ar[r]^g \ar[d]^\beta &
z \ar[r] \ar[d]^\gamma &
0 \\
0 \ar[r] & u \ar[r]^k & v \ar[r]^l & w
}
$$
为两行均正合的交换图。
\begin{enumerate}
\item 存在唯一态射 $\delta : \Ker(\gamma) \to \Coker(\alpha)$，使图
$$
\xymatrix{
y \ar[d]_\beta &
y \times_z \Ker(\gamma) \ar[l]_{\pi'} \ar[r]^{\pi} &
\Ker(\gamma) \ar[d]^\delta \\
v \ar[r]^{\iota'} & \Coker(\alpha) \amalg_u v &
\Coker(\alpha) \ar[l]_\iota
}
$$
交换，其中 $\pi$、$\pi'$ 为典范投影，$\iota$、$\iota'$ 为典范余注入。
\item 诱导序列
$$
\Ker(\alpha) \xrightarrow{f'} \Ker(\beta) \xrightarrow{g'}
\Ker(\gamma) \xrightarrow{\delta} \Coker(\alpha) \xrightarrow{k'}
\Coker(\beta) \xrightarrow{l'} \Coker(\gamma)
$$
正合。若 $f$ 为单射态射，则 $f'$ 亦然；若 $l$ 为满射态射，则 $l'$ 亦然。
\end{enumerate}
\end{lemma}

\begin{proof}
由引理 \ref{homology-lemma-cartesian-cocartesian}，$\pi$ 为满态射而 $\iota$ 为单态射，故 $\delta$ 的唯一性显然。令 $p = y \times_z \Ker(\gamma)$，$q = \Coker(\alpha) \amalg_u v$。令 $h : \Ker(\beta) \to y$、$i : \Ker(\gamma) \to z$ 与 $j : \Ker(\pi) \to p$ 为典范单射，并令 $\pi'' : u \to \Coker(\alpha)$ 为典范投影。结合引理 \ref{homology-lemma-cartesian-cocartesian}，得到如下各行正合的交换图：
$$
\xymatrix{
0 \ar[r] &
\Ker(\pi) \ar[r]^j &
p \ar[r]^{\pi} \ar[d]_{\pi'} &
\Ker(\gamma) \ar[d]_i \ar[r] & 0 \\
& x \ar[r]^f \ar[d]_\alpha & y \ar[r]^g \ar[d]_\beta &
z \ar[d]_\gamma \ar[r] & 0 \\
0 \ar[r] & u \ar[r]^k \ar[d]_{\pi''} &
v \ar[r]^l \ar[d]_{\iota'} & w & \\
0 \ar[r] & \Coker(\alpha) \ar[r]^\iota & q & &
}
$$
由于 $l \circ \beta \circ \pi' = \gamma \circ i \circ \pi = 0$ 且上图第三行正合，存在 $a : p \to u$ 使 $k \circ a = \beta \circ \pi'$。上图右上方框为笛卡尔图，故引理 \ref{homology-lemma-cartesian-kernel} 给出满态射 $b : x \to \Ker(\pi)$，满足 $\pi' \circ j \circ b = f$。于是
$k \circ a \circ j \circ b = \beta \circ \pi' \circ j \circ b = \beta \circ f = k \circ \alpha$。因 $k$ 为单态射，得 $a \circ j \circ b = \alpha$，进而 $\pi'' \circ a \circ j \circ b = \pi'' \circ \alpha = 0$。因 $b$ 为满态射，这蕴含 $\pi'' \circ a \circ j = 0$。所以，由上图第一行的正合性，存在 $\delta : \Ker(\gamma) \to \Coker(\alpha)$，使 $\delta \circ \pi = \pi'' \circ a$。于是
$\iota \circ \delta \circ \pi = \iota \circ \pi'' \circ a = \iota' \circ k \circ a = \iota' \circ \beta \circ \pi'$，正如所需。

\medskip\noindent
% Frozen source repeats \ref{homology-lemma-check-exactness-fibre-product} within this proof.
由于上图右上方框为笛卡尔图，存在 $c : \Ker(\beta) \to p$，使 $\pi' \circ c = h$ 且 $\pi \circ c = g'$。于是
$\iota \circ \delta \circ g' = \iota \circ \delta \circ \pi \circ c = \iota' \circ \beta \circ \pi' \circ c = \iota' \circ \beta \circ h = 0$。因 $\iota$ 为单态射，得 $\delta \circ g' = 0$。

\medskip\noindent
再设 $d : r \to \Ker(\gamma)$ 且 $\delta \circ d = 0$。把引理 \ref{homology-lemma-check-exactness-fibre-product} 应用于正合序列 $p \xrightarrow{\pi} \Ker(\gamma) \to 0$ 与 $d$，得到对象 $s$、满态射 $m : s \to r$ 和态射 $n : s \to p$，满足 $\pi \circ n = d \circ m$。由于 $\pi'' \circ a \circ n = \delta \circ d \circ m = 0$，把同一引理应用于正合序列 $x \xrightarrow{\alpha} u \xrightarrow{p} \Coker(\alpha)$ 与 $a \circ n$，得到对象 $t$、满态射 $\varepsilon : t \to s$ 与态射 $\zeta : t \to x$，满足 $a \circ n \circ \varepsilon = \alpha \circ \zeta$。有
$\beta \circ \pi' \circ n \circ \varepsilon = k \circ \alpha \circ \zeta = \beta \circ f \circ \zeta$。令 $\eta = \pi' \circ n \circ \varepsilon - f \circ \zeta : t \to y$，则 $\beta \circ \eta = 0$。故存在 $\vartheta : t \to \Ker(\beta)$ 使 $\eta = h \circ \vartheta$。于是
$i \circ g' \circ \vartheta = g \circ h \circ \vartheta = g \circ \pi' \circ n \circ \varepsilon - g \circ f \circ \zeta = i \circ \pi \circ n \circ \varepsilon = i \circ d \circ m \circ \varepsilon$。因 $i$ 为单态射，得 $g' \circ \vartheta = d \circ m \circ \varepsilon$。因此，由于 $m \circ \varepsilon$ 为满态射，引理 \ref{homology-lemma-check-exactness-fibre-product} 蕴含 $\Ker(\beta) \xrightarrow{g'} \Ker(\gamma) \xrightarrow{\delta} \Coker(\alpha)$ 正合。其余结论由引理 \ref{homology-lemma-exact-kernel-sequence} 与对偶性得到。
\end{proof}

\begin{lemma}
\label{homology-lemma-snake-natural}
设 $\mathcal{A}$ 为阿贝尔范畴。设
$$
\xymatrix{
& & & x\ar[ld]\ar[rr]\ar[dd]^(.4)\alpha
& & y\ar[ld]\ar[rr]\ar[dd]^(.4)\beta
& & z\ar[ld]\ar[rr]\ar[dd]^(.4)\gamma
& & 0\\
& & x'\ar[rr]\ar[dd]^(.4){\alpha'}
& & y'\ar[rr]\ar[dd]^(.4){\beta'}
& & z'\ar[rr]\ar[dd]^(.4){\gamma'}
& & 0
& \\
& 0\ar[rr]
& & u\ar[ld]\ar[rr]
& & v\ar[ld]\ar[rr]
& & w\ar[ld]
& & \\
0\ar[rr]
& & u'\ar[rr]
& & v'\ar[rr]
& & w'
& & &
}
$$
为各行正合的交换图。则诱导图
$$
\xymatrix@C=15pt{
\Ker(\alpha) \ar[r] \ar[d] &
\Ker(\beta) \ar[r] \ar[d] &
\Ker(\gamma) \ar[r]^(.45){\delta} \ar[d] &
\Coker(\alpha) \ar[r] \ar[d] &
\Coker(\beta) \ar[r] \ar[d] &
\Coker(\gamma) \ar[d] \\
\Ker(\alpha') \ar[r] &
\Ker(\beta') \ar[r] &
\Ker(\gamma') \ar[r]^(.45){\delta'} &
\Coker(\alpha') \ar[r] &
\Coker(\beta') \ar[r] &
\Coker(\gamma')
}
$$
交换。
\end{lemma}

\begin{proof}
略。
\end{proof}

\begin{lemma}
\label{homology-lemma-four-lemma}
设 $\mathcal{A}$ 为阿贝尔范畴。设
$$
\xymatrix{
w \ar[r] \ar[d]^\alpha & x \ar[r] \ar[d]^\beta & y \ar[r] \ar[d]^\gamma &
z \ar[d]^\delta \\
w' \ar[r] & x' \ar[r] & y' \ar[r] & z'
}
$$
为两行正合的交换图。
\begin{enumerate}
\item 若 $\alpha, \gamma$ 为满射态射且 $\delta$ 为单射态射，则 $\beta$ 为满射态射。
\item 若 $\beta, \delta$ 为单射态射且 $\alpha$ 为满射态射，则 $\gamma$ 为单射态射。
\end{enumerate}
\end{lemma}

\begin{proof}
假设 $\alpha, \gamma$ 为满射态射且 $\delta$ 为单射态射。可将 $w'$ 替换为 $\Im(w' \to x')$，即可以假设 $w' \to x'$ 为单射态射。可将 $z$ 替换为 $\Im(y \to z)$，即可以假设 $y \to z$ 为满射态射。于是可将引理 \ref{homology-lemma-snake} 应用于
$$
\xymatrix{
& \Ker(y \to z) \ar[r] \ar[d] & y \ar[r] \ar[d] & z \ar[r] \ar[d] & 0 \\
0 \ar[r] & \Ker(y' \to z') \ar[r] & y' \ar[r] & z'
}
$$
从而得出 $\Ker(y \to z) \to \Ker(y' \to z')$ 为满射态射。最后，把引理 \ref{homology-lemma-snake} 应用于
$$
\xymatrix{
& w \ar[r] \ar[d] & x \ar[r] \ar[d] & \Ker(y \to z) \ar[r] \ar[d] & 0 \\
0 \ar[r] & w' \ar[r] & x' \ar[r] & \Ker(y' \to z')
}
$$
即可得出 $x \to x'$ 为满射态射。这证明 (1)；(2) 的证明与之对偶。
\end{proof}

\begin{lemma}
\label{homology-lemma-five-lemma}
\begin{reference}
\cite[Lemma 4.5 page 16]{Eilenberg-Steenrod}
\end{reference}
设 $\mathcal{A}$ 为阿贝尔范畴。设
$$
\xymatrix{
v \ar[r] \ar[d]^\alpha &
w \ar[r] \ar[d]^\beta &
x \ar[r] \ar[d]^\gamma &
y \ar[r] \ar[d]^\delta &
z \ar[d]^\epsilon \\
v' \ar[r] & w' \ar[r] & x' \ar[r] & y' \ar[r] & z'
}
$$
为两行正合的交换图。若 $\beta, \delta$ 是同构，$\epsilon$ 为单射态射且 $\alpha$ 为满射态射，则 $\gamma$ 是同构。
\end{lemma}

\begin{proof}
这是引理 \ref{homology-lemma-four-lemma} 的直接推论。
\end{proof}

\section{扩张}
\label{homology-section-extensions}

\begin{definition}
\label{homology-definition-extension}
设 $\mathcal{A}$ 为阿贝尔范畴，$A, B \in \Ob(\mathcal{A})$。$B$ 被 $A$ 的一个{\it 扩张 $E$} 是短正合序列
$$
0 \to A \to E \to B \to 0.
$$
两个扩张 $0 \to A \to E \to B \to 0$ 与 $0 \to A \to F \to B \to 0$ 之间的一个{\it 扩张态射}，是 $\mathcal{A}$ 中使下图交换的态射 $f : E \to F$：
$$
\xymatrix{
0 \ar[r] &
A \ar[r] \ar[d]^{\text{id}} &
E \ar[r] \ar[d]^f &
B \ar[r] \ar[d]^{\text{id}} &
0 \\
0 \ar[r] &
A \ar[r] &
F \ar[r] &
B \ar[r] &
0
}
$$
因此，$B$ 被 $A$ 的扩张构成一个范畴。
\end{definition}

\noindent
作为一种语言上的简略，我们常省略对态射 $A \to E$ 与 $E \to B$ 的提及，尽管它们确实是扩张结构的一部分。

\begin{definition}
\label{homology-definition-ext-group}
设 $\mathcal{A}$ 为阿贝尔范畴，$A, B \in \Ob(\mathcal{A})$。$B$ 被 $A$ 的扩张之同构类所成的集合记作
$$
\Ext_\mathcal{A}(B, A).
$$
它称为 {\it $\Ext$ 群}。
\end{definition}

\noindent
这个定义是可行的，因为按我们的约定，$\Ob(\mathcal{A})$ 是集合，从而 $\Ext_\mathcal{A}(B, A)$ 也是集合。对《范畴》，注记 \ref{categories-remark-big-categories} 中列出的各种“大”阿贝尔范畴，也可直接验证 $\Ext_\mathcal{A}(B, A)$ 是集合。稍后还将看到，当 $\mathcal{A}$ 有足够多投射对象或足够多内射对象时，总是如此。此处待补将来的引用。

\medskip\noindent
事实上，可如下把 $\Ext_\mathcal{A}(-, -)$ 变为函子
$$
\mathcal{A} \times \mathcal{A}^{opp} \longrightarrow \textit{Sets}, \quad
(A, B) \longmapsto \Ext_\mathcal{A}(B, A)
$$
\begin{enumerate}
\item 给定态射 $B' \to B$ 与 $B$ 被 $A$ 的扩张 $E$，定义 $E' = E \times_B B'$。由引理 \ref{homology-lemma-cartesian-kernel}，有如下短正合序列的交换图：
$$
\xymatrix{
0 \ar[r] & A \ar[r] \ar[d] & E' \ar[r] \ar[d] & B' \ar[r] \ar[d] & 0 \\
0 \ar[r] & A \ar[r] & E \ar[r] & B \ar[r] & 0
}
$$
扩张 $E'$ 称为 $E$ 沿 $B' \to B$ 的{\it 拉回}。
\item 给定态射 $A \to A'$ 与 $B$ 被 $A$ 的扩张 $E$，定义 $E' = A' \amalg_A E$。由引理 \ref{homology-lemma-cartesian-cocartesian}，有如下短正合序列的交换图：
$$
\xymatrix{
0 \ar[r] & A \ar[r] \ar[d] & E \ar[r] \ar[d] & B \ar[r] \ar[d] & 0 \\
0 \ar[r] & A' \ar[r] & E' \ar[r] & B \ar[r] & 0
}
$$
扩张 $E'$ 称为 $E$ 沿 $A \to A'$ 的{\it 推出}。
\end{enumerate}
要说明这确实定义了上述函子，需要验证若干事项。首先，变量 $B$ 上的函子性要求
$(E \times_B B') \times_{B'} B'' = E \times_B B''$，这是纤维积的一般性质。变量 $A$ 上的函子性由对偶方式处理。最后，给定 $A \to A'$ 与 $B' \to B$，须证明
$$
A' \amalg_A (E \times_B B')
\cong
(A' \amalg_A E)\times_B B'
$$
作为 $B'$ 被 $A'$ 的扩张成立。回顾 $A' \amalg_A E$ 是 $A' \oplus E$ 的一个商。因此右边是 $A' \oplus E \times_B B'$ 的一个商；容易看出其核恰好给出左边。

\medskip\noindent
注意，若 $E_1$ 与 $E_2$ 都是 $B$ 被 $A$ 的扩张，则 $E_1\oplus E_2$ 是 $B \oplus B$ 被 $A\oplus A$ 的扩张。沿求和映射 $A \oplus A \to A$ 推出，再沿对角映射 $B \to B \oplus B$ 拉回，便得到 $B$ 被 $A$ 的扩张 $E_1 + E_2$：
$$
\xymatrix{
0 \ar[r] &
A \oplus A \ar[r] \ar[d]_\Sigma &
E_1 \oplus E_2 \ar[r] \ar[d] &
B \oplus B \ar[r] \ar[d] &
0 \\
0 \ar[r] &
A \ar[r] &
E' \ar[r] &
B \oplus B \ar[r] &
0\\
0 \ar[r] &
A \ar[r] \ar[u] &
E_1 + E_2 \ar[r] \ar[u] &
B \ar[r] \ar[u]^\Delta &
0
}
$$
扩张 $E_1 + E_2$ 称为给定扩张的 {\it Baer 和}。

\begin{lemma}
\label{homology-lemma-baer-sum}
上述构造 $(E_1, E_2) \mapsto E_1 + E_2$ 在 $\Ext_\mathcal{A}(B, A)$ 上定义交换群运算，并且对两个变量均具函子性。
\end{lemma}

\begin{proof}
略。
\end{proof}

\begin{lemma}
\label{homology-lemma-six-term-sequence-ext}
设 $\mathcal{A}$ 为阿贝尔范畴，$0 \to M_1 \to M_2 \to M_3 \to 0$ 为 $\mathcal{A}$ 中的短正合序列。
\begin{enumerate}
\item 对 $\mathcal{A}$ 的每个对象 $N$，有典范的六项阿贝尔群正合序列
$$
\xymatrix{
0 \ar[r] &
\Hom_\mathcal{A}(M_3, N) \ar[r] &
\Hom_\mathcal{A}(M_2, N) \ar[r] &
\Hom_\mathcal{A}(M_1, N) \ar[lld] \\
& \Ext_\mathcal{A}(M_3, N) \ar[r] &
\Ext_\mathcal{A}(M_2, N) \ar[r] &
\Ext_\mathcal{A}(M_1, N)
}
$$
；
\item 对 $\mathcal{A}$ 的每个对象 $N$，有典范的六项阿贝尔群正合序列
$$
\xymatrix{
0 \ar[r] &
\Hom_\mathcal{A}(N, M_1) \ar[r] &
\Hom_\mathcal{A}(N, M_2) \ar[r] &
\Hom_\mathcal{A}(N, M_3) \ar[lld] \\
& \Ext_\mathcal{A}(N, M_1) \ar[r] &
\Ext_\mathcal{A}(N, M_2) \ar[r] &
\Ext_\mathcal{A}(N, M_3)
}
$$
。
\end{enumerate}
\end{lemma}

\begin{proof}
略。提示：边界映射由给定短正合序列的推出或拉回来定义。
\end{proof}

\section{加性函子}
\label{homology-section-functors}

\noindent
先给出一个刻画加性范畴之间加性函子的完全平凡的引理。

\begin{lemma}
\label{homology-lemma-additive-functor}
设 $\mathcal{A}$ 与 $\mathcal{B}$ 为加性范畴，$F : \mathcal{A} \to \mathcal{B}$ 为函子。下列条件等价：
\begin{enumerate}
\item $F$ 是加性函子；
\item 对所有 $A, B \in \mathcal{A}$，$F(A) \oplus F(B) \to F(A \oplus B)$ 是同构；
\item 对所有 $A, B \in \mathcal{A}$，$F(A \oplus B) \to F(A) \oplus F(B)$ 是同构。
\end{enumerate}
\end{lemma}

\begin{proof}
由引理 \ref{homology-lemma-additive-additive}，加性函子与直和可交换，故 (1) 蕴含 (2) 与 (3)。若 (2) 成立且 $0$ 表示 $\mathcal{A}$ 的零对象，则映射 $F(0) \oplus F(0) \to F(0 \oplus 0) \cong F(0)$ 为同构，这蕴含对 $\mathcal{B}$ 中所有 $C$，对角映射 $\Hom(F(0), C) \to \Hom(F(0), C) \oplus \Hom(F(0), C)$ 均为双射。因此 $F(0)$ 是 $\mathcal{B}$ 的零对象，且 $F$ 把任意两个对象间的零态射映为零态射。若 (3) 成立，结论同样成立。于是 (2) 与 (3) 等价，因为复合
$F(A) \oplus F(B) \to F(A \oplus B) \to F(A) \oplus F(B)$
由矩阵
$$
\left(
\begin{matrix}
F(1) & F(0) \\
F(0) & F(1)
\end{matrix}
\right)
$$
给出，因而是恒等映射。现假设 (2) 与 (3) 成立。设 $f, g : A \to B$ 为映射。则 $f + g$ 等于复合
$$
A \to A \oplus A \xrightarrow{\text{diag}(f, g)} B \oplus B \to B
$$
应用函子 $F$ 并考虑图
$$
\xymatrix{
F(A) \ar[r] \ar[rd] &
F(A \oplus A) \ar[rr]_{F(\text{diag}(f, g))} & &
F(B \oplus B) \ar[r] \ar[d] &
F(B) \\
&
F(A) \oplus F(A) \ar[u] \ar[rr]^{\text{diag}(F(f), F(g))} & &
F(B) \oplus F(B) \ar[ru]
}
$$
我们断言此图交换。中间方框可分别在四个诱导映射 $F(A) \to F(B)$ 上验证；这由 $F$ 是函子且把零态射映为零态射得出。对左侧三角形，利用 $F(A \oplus A) \to F(A) \oplus F(A)$ 是同构，可知只需与该映射复合后验证，而此时验证是平凡的。另一三角形由对偶方式处理。因此沿底部复合等于 $F(f + g)$，结论随之得出。
\end{proof}

\noindent
回顾我们在《范畴》，定义 \ref{categories-definition-exact} 中，对具有有限（余）极限的范畴之间的函子定义了“右正合”“左正合”和“正合”。因此这些概念特别适用于阿贝尔范畴之间的函子。

\begin{lemma}
\label{homology-lemma-exact-functor}
设 $\mathcal{A}$ 与 $\mathcal{B}$ 为阿贝尔范畴，$F : \mathcal{A} \to \mathcal{B}$ 为函子。
\begin{enumerate}
\item 若 $F$ 左正合或右正合，则它是加性函子。
\item $F$ 左正合，当且仅当对每个短正合序列 $0 \to A \to B \to C \to 0$，序列 $0 \to F(A) \to F(B) \to F(C)$ 正合。
\item $F$ 右正合，当且仅当对每个短正合序列 $0 \to A \to B \to C \to 0$，序列 $F(A) \to F(B) \to F(C) \to 0$ 正合。
\item $F$ 正合，当且仅当对每个短正合序列 $0 \to A \to B \to C \to 0$，序列 $0 \to F(A) \to F(B) \to F(C) \to 0$ 正合。
\end{enumerate}
\end{lemma}

\begin{proof}
若 $F$ 左正合，即 $F$ 与有限极限可交换，则 $F$ 把积送到积，故保直和，因而由引理 \ref{homology-lemma-additive-functor}，$F$ 是加性函子。另一方面，假设对每个短正合序列 $0 \to A \to B \to C \to 0$，序列 $0 \to F(A) \to F(B) \to F(C)$ 都正合。对任意两个对象 $A, B$，有短正合序列
$$
0 \to A \to A \oplus B \to B \to 0
$$
例如见引理 \ref{homology-lemma-additive-cat-biproduct-kernel}。由假设，下列交换图的下行正合：
$$
\xymatrix{
0 \ar[r] &
F(A) \ar[d] \ar[r] &
F(A) \oplus F(B) \ar[r] \ar[d] &
F(B) \ar[d] \ar[r] &
0 \\
0 \ar[r] &
F(A) \ar[r] &
F(A \oplus B) \ar[r] &
F(B)
}
$$
因此由蛇引理（引理 \ref{homology-lemma-snake}），$F(A) \oplus F(B) \to F(A \oplus B)$ 是同构。故在这种情形下 $F$ 也是加性函子。于是以下可假设 $F$ 为加性函子。

\medskip\noindent
记 $f : B \to C$ 为从 $B$ 到 $C$ 的映射。$0 \to A \to B \to C$ 正合恰好意味着 $A = \Ker(f)$。显然，$f$ 的核是从 $B$ 到 $C$ 的两个映射 $f$ 与 $0$ 的等化子。因此，若 $F$ 与极限可交换，则 $F(\Ker(f)) = \Ker(F(f))$，这恰好意味着 $0 \to F(A) \to F(B) \to F(C)$ 正合。

\medskip\noindent
反过来，假设 $F$ 是加性函子，并将任一短正合序列 $0 \to A \to B \to C \to 0$ 变为正合序列 $0 \to F(A) \to F(B) \to F(C)$。由于 $F$ 加性，它与直和、因而与 $\mathcal{A}$ 中的有限积可交换。要证明它与有限极限可交换，只需证明它与等化子可交换。但阿贝尔范畴中的等化子就是差映射的核，故只需证明 $F$ 与取核可交换。设 $f : A \to B$ 为态射。把 $f$ 分解为 $A \to I \to B$，其中 $f' : A \to I$ 为满射态射，$i : I \to B$ 为单射态射（由阿贝尔范畴的定义可作此分解）。显然 $\Ker(f) = \Ker(f')$。此外，
$0 \to \Ker(f') \to A \to I \to 0$ 与 $0 \to I \to B \to B/I \to 0$ 均为短正合序列。由对 $F$ 的条件可知
$0 \to F(\Ker(f')) \to F(A) \to F(I)$ 与 $0 \to F(I) \to F(B) \to F(B/I)$ 正合。因此，$F(\Ker(f'))$ 也是映射 $F(A) \to F(B)$ 的核，结论得证。

\medskip\noindent
(3) 的证明与 (2) 类似。(4) 是 (2) 与 (3) 的合并。
\end{proof}

\begin{lemma}
\label{homology-lemma-exact-functor-ext}
设 $\mathcal{A}$ 与 $\mathcal{B}$ 为阿贝尔范畴，$F : \mathcal{A} \to \mathcal{B}$ 为正合函子。对 $\mathcal{A}$ 的每一对对象 $A, B$，函子 $F$ 诱导阿贝尔群同态
$$
\Ext_\mathcal{A}(B, A)
\longrightarrow
\Ext_\mathcal{B}(F(B), F(A))
$$
它把扩张 $E$ 映为 $F(E)$。
\end{lemma}

\begin{proof}
略。
\end{proof}

\noindent
下述引理用于证明一个位点上的阿贝尔层所成的范畴是阿贝尔范畴，其中函子 $b$ 是层化函子。

\begin{lemma}
\label{homology-lemma-adjoint-get-abelian}
设 $a : \mathcal{A} \to \mathcal{B}$ 与 $b : \mathcal{B} \to \mathcal{A}$ 为函子。假设：
\begin{enumerate}
\item $\mathcal{A}$、$\mathcal{B}$ 为加性范畴，$a$、$b$ 为加性函子，且 $a$ 是 $b$ 的右伴随；
\item $\mathcal{B}$ 为阿贝尔范畴且 $b$ 左正合；
\item $ba \cong \text{id}_\mathcal{A}$。
\end{enumerate}
则 $\mathcal{A}$ 是阿贝尔范畴。
\end{lemma}

\begin{proof}
由 $\mathcal{B}$ 为阿贝尔范畴及引理 \ref{homology-lemma-colimit-abelian-category}，$\mathcal{B}$ 中所有有限极限和余极限都存在。由于 $b$ 是左伴随，$b$ 也右正合，因而正合，见《范畴》，引理 \ref{categories-lemma-exact-adjoint}。设 $\varphi : B_1 \to B_2$ 为 $\mathcal{B}$ 的态射。特别地，若 $K = \Ker(B_1 \to B_2)$，则 $K$ 是 $0$ 与 $\varphi$ 的等化子，故 $bK$ 是 $0$ 与 $b\varphi$ 的等化子，即 $bK$ 是 $b\varphi$ 的核。类似地，若 $Q = \Coker(B_1 \to B_2)$，则 $Q$ 是 $0$ 与 $\varphi$ 的余等化子，故 $bQ$ 是 $0$ 与 $b\varphi$ 的余等化子，即 $bQ$ 是 $b\varphi$ 的余核。因此，$\mathcal{A}$ 中每个形如 $b\varphi$ 的态射都有核和余核。又因 $ba \cong \text{id}$，$\mathcal{A}$ 的每个态射都具有这种形式，于是 $\mathcal{A}$ 中的核和余核均存在。事实上，上述论证表明，若 $\psi : A_1 \to A_2$ 为态射，则
$$
\Ker(\psi) = b\Ker(a\psi),
\quad\text{且}\quad
\Coker(\psi) = b\Coker(a\psi).
$$
还须证明 $\Coim(\psi)= \Im(\psi)$，如下进行。

首先注意，由于 $\mathcal{A}$ 有核和余核，它有所有有限极限和余极限（见引理 \ref{homology-lemma-colimit-abelian-category} 的证明）。于是由《范畴》，引理 \ref{categories-lemma-exact-adjoint}，$a$ 左正合，因而把核（即等化子）变为核。
\begin{align*}
\Coim(\psi)
& =
\Coker(\Ker(\psi) \to A_1)
& \text{由定义} \\
& =
b\Coker(a(\Ker(\psi) \to A_1))
& \text{由上式} \\
& =
b\Coker(\Ker(a\psi) \to aA_1))
& a\text{ 保持核} \\
& =
b\Coim(a\psi)
& \text{由定义} \\
& =
b\Im(a\psi)
& \mathcal{B}\text{ 为 Abel 范畴} \\
& =
b\Ker(aA_2 \to \Coker(a\psi))
& \text{由定义} \\
& =
\Ker(baA_2 \to b\Coker(a\psi))
& b\text{ 保持核} \\
& =
\Ker(A_2 \to b\Coker(a\psi))
& ba = \text{id}_\mathcal{A} \\
& =
\Ker(A_2 \to \Coker(\psi))
& \text{由上式} \\
& =
\Im(\psi)
& \text{由定义}.
\end{align*}
故引理成立。
\end{proof}

\section{局部化}
\label{homology-section-localization}

\noindent
本节说明 Gabriel--Zisman 局部化如何与范畴上的加性结构相互作用。

\begin{lemma}
\label{homology-lemma-localization-preadditive}
设 $\mathcal{C}$ 为预加性范畴，$S$ 为左乘法系或右乘法系。$S^{-1}\mathcal{C}$ 上存在唯一的预加性结构，使局部化函子 $Q : \mathcal{C} \to S^{-1}\mathcal{C}$ 为加性函子。
\end{lemma}

\begin{proof}
证明 $S$ 为左乘法系的情形；$S$ 为右乘法系的情形与之对偶。设 $X, Y$ 为 $\mathcal{C}$ 的对象，$\alpha, \beta : X \to Y$ 为 $S^{-1}\mathcal{C}$ 中的态射。由《范畴》，引理 \ref{categories-lemma-morphisms-left-localization}，可用具有公分母 $s$ 的两对 $s^{-1}f, s^{-1}g$ 表示它们。假设 $S^{-1}\mathcal{C}$ 上已有预加性结构且 $Q$ 为加性函子，则
\begin{align*}
s^{-1}f+s^{-1}g
&=Q(s)^{-1}\circ Q(f)+Q(s)^{-1}\circ Q(g)\\
&=Q(s)^{-1}\circ(Q(f)+Q(g))\\
&=Q(s)^{-1}\circ Q(f+g)\\
&=s^{-1}(f+g).
\end{align*}
因此，若 $S^{-1}\mathcal{C}$ 上的预加性结构存在，它必唯一。相应地，把 $\alpha + \beta$ 定义为 $s^{-1}(f + g)$ 的等价类；以下证明此定义良定且复合为双线性。一旦完成这些验证，$Q$ 是加性函子便显然成立。

\medskip\noindent
先证明上述构造良定。抽象地说，这使用了阿贝尔群的滤过余极限与集合的滤过余极限相同，并应用《范畴》，公式 (\ref{categories-equation-left-localization-morphisms-colimit})。更具体地，设 $s : Y \to Y_1$ 且 $f, g : X \to Y_1$。再设 $\alpha, \beta$ 有第二组表示 $(s')^{-1}f', (s')^{-1}g'$，其中 $s' : Y \to Y_2$，$f', g' : X \to Y_2$。由《范畴》，注记 \ref{categories-remark-left-localization-morphisms-colimit}，可找到态射 $s_3 : Y \to Y_3$ 以及 $a_1 : Y_1 \to Y_3$、$a_2 : Y_2 \to Y_3$，满足 $a_1 \circ s = s_3 = a_2 \circ s'$，并且 $a_1 \circ f = a_2 \circ f'$、$a_1 \circ g = a_2 \circ g'$。于是 $s^{-1}(f + g)$ 等价于
\begin{align*}
s_3^{-1}(a_1 \circ (f + g)) & =
s_3^{-1}(a_1 \circ f + a_1 \circ g) \\
& = s_3^{-1}(a_2 \circ f' + a_2 \circ g') \\
& = s_3^{-1}(a_2 \circ (f' + g')),
\end{align*}
而后者等价于 $(s')^{-1}(f' + g')$。

\medskip\noindent
固定 $s : Y \to Y'$ 与 $f, g : X \to Y'$，使得在 $S^{-1}\mathcal{C}$ 中作为 $X \to Y$ 的态射有 $\alpha = s^{-1}f$、$\beta = s^{-1}g$。为证明复合双线性，先考虑 $S^{-1}\mathcal{C}$ 中态射 $\gamma : Y \to Z$ 的情形。设某个 $h : Y \to Z'$ 与 $S$ 中的 $t : Z \to Z'$ 使 $\gamma = t^{-1}h$。利用 LMS2，选取态射 $a : Y' \to Z''$ 与 $S$ 中的 $t' : Z' \to Z''$，使 $a \circ s = t' \circ h$。图示为
$$
\xymatrix{
& & Z \ar[d]^t \\
& Y \ar[r]^h \ar[d]^s & Z' \ar[d]^{t'} \\
X \ar[r]^{f, g} & Y' \ar[r]^a & Z''
}
$$
于是 $\gamma \circ \alpha = (t' \circ t)^{-1}(a \circ f)$，且 $\gamma \circ \beta = (t' \circ t)^{-1}(a \circ g)$。所以 $\gamma \circ (\alpha + \beta)$ 由
$(t' \circ t)^{-1}(a \circ (f + g)) = (t' \circ t)^{-1}(a \circ f + a \circ g)$ 表示，而它表示 $\gamma \circ \alpha + \gamma \circ \beta$。

\medskip\noindent
最后，设 $\delta : W \to X$ 是 $S^{-1}\mathcal{C}$ 的另一态射。设某个 $i : W \to X'$ 与 $S$ 中的 $r : X \to X'$ 使 $\delta = r^{-1}i$。我们断言可找到 $S$ 中的态射 $s' : Y' \to Y''$ 与态射 $a'', b'' : X' \to Y''$，使下图交换：
$$
\xymatrix{
& & & Y \ar[d]^s \\
& X \ar[rr]^{f, g, f + g} \ar[d]^r & & Y' \ar[d]^{s'} \\
W \ar[r]^i & X' \ar[rr]^{a'', b'', a'' + b''} & & Y''
}
$$
确实，先由 LMS2 选取 $S$ 中的 $s_1 : Y' \to Y_1$、$s_2 : Y' \to Y_2$ 以及 $a : X' \to Y_1$、$b : X' \to Y_2$，使 $a \circ r = s_1 \circ f$ 且 $b \circ r = s_2 \circ f$。
% P02-E0007: frozen source line 1848 repeats f; the second occurrence appears to require g.
然后利用范畴 $Y'/S$ 为滤过范畴（见《范畴》，注记 \ref{categories-remark-left-localization-morphisms-colimit}），可找到 $s' : Y' \to Y''$ 及态射 $a' : Y_1 \to Y''$、$b' : Y_2 \to Y''$，满足 $s' = a' \circ s_1$ 且 $s' = b' \circ s_2$。令 $a'' = a' \circ a$、$b'' = b' \circ b$ 即可。此时，复合 $\alpha \circ \delta$ 与 $\beta \circ \delta$ 分别由 $(s' \circ s)^{-1}(a'' \circ i)$ 与 $(s' \circ s)^{-1}(b'' \circ i)$ 表示。因此 $\alpha \circ \delta + \beta \circ \delta$ 由
$(s' \circ s)^{-1}(a'' \circ i + b'' \circ i) = (s' \circ s)^{-1}((a'' + b'') \circ i)$ 表示；由上图，它又是 $(\alpha + \beta) \circ \delta$ 的一个代表。
\end{proof}

\begin{lemma}
\label{homology-lemma-localization-additive}
设 $\mathcal{C}$ 为加性范畴，$S$ 为左乘法系或右乘法系。则 $S^{-1}\mathcal{C}$ 为加性范畴，且局部化函子 $Q : \mathcal{C} \to S^{-1}\mathcal{C}$ 为加性函子。
\end{lemma}

\begin{proof}
由引理 \ref{homology-lemma-localization-preadditive}，$S^{-1}\mathcal{C}$ 为预加性范畴且 $Q$ 为加性函子。由引理 \ref{homology-lemma-additive-additive}，$S^{-1}\mathcal{C}$ 有零对象与直和。
\end{proof}

\noindent
下述引理描述在反演一个乘法系时局部化函子的“核”。

\begin{lemma}
\label{homology-lemma-kernel-localization}
设 $\mathcal{C}$ 为加性范畴，$S$ 为乘法系，$X$ 为 $\mathcal{C}$ 的对象。下列条件等价：
\begin{enumerate}
\item 在 $S^{-1}\mathcal{C}$ 中 $Q(X) = 0$；
\item 存在 $Y \in \Ob(\mathcal{C})$，使 $0 : X \to Y$ 属于 $S$；
\item 存在 $Z \in \Ob(\mathcal{C})$，使 $0 : Z \to X$ 属于 $S$。
\end{enumerate}
\end{lemma}

\begin{proof}
若 (2) 成立，则 $0 = Q(0) : Q(X) \to Q(Y)$ 是同构。在加性范畴 $S^{-1}\mathcal{C}$ 中，这蕴含 $Q(X) = 0$。故 (2) $\Rightarrow$ (1)。类似地，(3) $\Rightarrow$ (1)。设 $Q(X) = 0$，则态射 $f : 0 \to X$ 在 $S^{-1}\mathcal{C}$ 中变为同构。因此由《范畴》，引理 \ref{categories-lemma-what-gets-inverted}，存在态射 $g : Z \to 0$，使 $fg \in S$。这证明 (1) $\Rightarrow$ (3)。类似地，(1) $\Rightarrow$ (2)。
\end{proof}

\begin{lemma}
\label{homology-lemma-localization-abelian}
设 $\mathcal{A}$ 为阿贝尔范畴。
\begin{enumerate}
\item 若 $S$ 为左乘法系，则范畴 $S^{-1}\mathcal{A}$ 有余核，且函子 $Q : \mathcal{A} \to S^{-1}\mathcal{A}$ 与取余核可交换。
\item 若 $S$ 为右乘法系，则范畴 $S^{-1}\mathcal{A}$ 有核，且函子 $Q : \mathcal{A} \to S^{-1}\mathcal{A}$ 与取核可交换。
\item 若 $S$ 为乘法系，则范畴 $S^{-1}\mathcal{A}$ 为阿贝尔范畴，且函子 $Q : \mathcal{A} \to S^{-1}\mathcal{A}$ 正合。
\end{enumerate}
\end{lemma}

\begin{proof}
假设 $S$ 为左乘法系。设 $a : X \to Y$ 为 $S^{-1}\mathcal{A}$ 中的态射，则某个 $S$ 中的 $s : Y \to Y'$ 与 $f : X \to Y'$ 使 $a = s^{-1}f$。由于 $Q(s)$ 是同构，$\Coker(a : X \to Y)$ 的存在性等价于 $\Coker(Q(f) : X \to Y')$ 的存在性。因为 $\Coker(Q(f))$ 是 $0$ 与 $Q(f)$ 的余等化子，由《范畴》，引理 \ref{categories-lemma-left-localization-limits}，$\Coker(Q(f))$ 由 $Q(\Coker(f))$ 表示。这证明 (1)。

\medskip\noindent
(2) 与 (1) 对偶。

\medskip\noindent
若 $S$ 为乘法系，则 $S$ 同时是左乘法系和右乘法系。因此 $S^{-1}\mathcal{A}$ 有核和余核，且 $Q$ 与取核、取余核可交换。要完成 (3) 的证明，还须说明 $S^{-1}\mathcal{A}$ 中 $\Coim = \Im$。再次利用 $S^{-1}\mathcal{A}$ 的任意箭头都同构于某个箭头 $Q(f)$，结论由 $\mathcal{A}$ 中的相应结果推出。
\end{proof}

\section{Jordan--H\"older 定理}
\label{homology-section-jordan-holder}

\noindent
Jordan--H\"older 引理即引理 \ref{homology-lemma-jordan-holder}。先给出一些定义。

\begin{definition}
\label{homology-definition-simple}
设 $\mathcal{A}$ 为阿贝尔范畴。若 $\mathcal{A}$ 的对象 $A$ 非零，且 $A$ 的子对象只有 $0$ 与 $A$，则称该对象为{\it 单对象}。
\end{definition}

\begin{definition}
\label{homology-definition-Artinian}
设 $\mathcal{A}$ 为阿贝尔范畴。
\begin{enumerate}
\item 若 $\mathcal{A}$ 的对象 $A$ 对子对象满足降链条件，则称该对象为 {\it Artin 对象}。
\item 若 $\mathcal{A}$ 的每个对象都是 Artin 对象，则称 $\mathcal{A}$ 为 {\it Artin 范畴}。
\end{enumerate}
\end{definition}

\begin{definition}
\label{homology-definition-Noetherian}
设 $\mathcal{A}$ 为阿贝尔范畴。
\begin{enumerate}
\item 若 $\mathcal{A}$ 的对象 $A$ 对子对象满足升链条件，则称该对象为 {\it Noether 对象}。
\item 若 $\mathcal{A}$ 的每个对象都是 Noether 对象，则称 $\mathcal{A}$ 为 {\it Noether 范畴}。
\end{enumerate}
\end{definition}

\begin{lemma}
\label{homology-lemma-ses-artinian}
设 $\mathcal{A}$ 为阿贝尔范畴，$0 \to A_1 \to A_2 \to A_3 \to 0$ 为 $\mathcal{A}$ 中的短正合序列。则 $A_2$ 是 Artin 对象，当且仅当 $A_1$ 与 $A_3$ 都是 Artin 对象。
\end{lemma}

\begin{proof}
略。
\end{proof}

\begin{lemma}
\label{homology-lemma-ses-noetherian}
设 $\mathcal{A}$ 为阿贝尔范畴，$0 \to A_1 \to A_2 \to A_3 \to 0$ 为 $\mathcal{A}$ 中的短正合序列。则 $A_2$ 是 Noether 对象，当且仅当 $A_1$ 与 $A_3$ 都是 Noether 对象。
\end{lemma}

\begin{proof}
略。
\end{proof}

\begin{lemma}
\label{homology-lemma-finite-length}
设 $\mathcal{A}$ 为阿贝尔范畴，$A$ 为 $\mathcal{A}$ 的对象。下列条件等价：
\begin{enumerate}
\item $A$ 既为 Artin 对象又为 Noether 对象；
\item 存在由子对象组成的滤过 $0 \subset A_1 \subset A_2 \subset \ldots \subset A_n = A$，使 $i = 1, \ldots, n$ 时 $A_i/A_{i - 1}$ 为单对象。
\end{enumerate}
\end{lemma}

\begin{proof}
假设 (1)。若 $A$ 为零，则 (2) 成立。若 $A$ 非零，则由 Artin 性质，存在最小非零对象 $A_1 \subset A$；显然 $A_1$ 为单对象。若 $A_1 = A$，则已经完成。否则，可找到在满足 $A_2 \not = A_1$ 的条件下极小的 $A_1 \subset A_2 \subset A$；于是 $A_2/A_1$ 为单对象。如此继续，可得到 $A$ 的子对象序列 $0 \subset A_1 \subset A_2 \subset \ldots$，且 $A_i/A_{i - 1}$ 为单对象。由于 $A$ 为 Noether 对象，该过程必终止，故 (2) 成立。

\medskip\noindent
假设 (2)。对 $n$ 归纳证明 (1)。若 $n = 1$，则 $A$ 为单对象，显然既是 Noether 对象又是 Artin 对象。若结论对 $n - 1$ 成立，则利用短正合序列 $0 \to A_{n - 1} \to A_n \to A_n/A_{n - 1} \to 0$ 以及引理 \ref{homology-lemma-ses-artinian}、\ref{homology-lemma-ses-noetherian}，即可对 $n$ 得出结论。
\end{proof}

\begin{lemma}[Jordan--H\"older]
\label{homology-lemma-jordan-holder}
设 $\mathcal{A}$ 为阿贝尔范畴，$A$ 为 $\mathcal{A}$ 的对象并满足引理 \ref{homology-lemma-finite-length} 中的等价条件。给定两个滤过
$$
0 \subset A_1 \subset A_2 \subset \ldots \subset A_n = A
\quad\text{且}\quad
0 \subset B_1 \subset B_2 \subset \ldots \subset B_m = A
$$
并设 $S_i = A_i/A_{i - 1}$、$T_j = B_j/B_{j - 1}$ 都是单对象，则 $n = m$，且存在 $\{1, \ldots, n\}$ 的置换 $\sigma$，使对所有 $i \in \{1, \ldots, n\}$ 都有 $S_i \cong T_{\sigma(i)}$。
\end{lemma}

\begin{proof}
令 $j$ 为满足 $A_1 \subset B_j$ 的最小指标。则映射 $S_1 = A_1 \to B_j/B_{j - 1} = T_j$ 是同构。此外，对象 $A/A_1 = A_n/A_1 = B_m/A_1$ 有两个滤过
$$
0 \subset A_2/A_1 \subset A_3/A_1 \subset \ldots \subset A_n/A_1
$$
与
$$
0 \subset (B_1 + A_1)/A_1 \subset \ldots \subset
(B_{j - 1} + A_1)/A_1 = B_j/A_1 \subset B_{j + 1}/A_1
\subset \ldots \subset B_m/A_1
$$
由归纳法得出结论。
\end{proof}








\section{Serre 子范畴}
\label{homology-section-serre-subcategories}

\noindent
在 \cite[Chapter I, Section 1]{Serre_homotopie_classes} 中定义了阿贝尔群的一个“类”。许多作者把这一概念推广到阿贝尔范畴（方式略有不同）。我们采用下面这个与 Serre 原定义几乎完全相同的版本。

\begin{definition}
\label{homology-definition-serre-subcategory}
\begin{reference}
\cite[Condition (I) on page 259]{Serre_homotopie_classes}
\end{reference}
设 $\mathcal{A}$ 为阿贝尔范畴。
\begin{enumerate}
\item $\mathcal{A}$ 的一个 {\it Serre 子范畴}是 $\mathcal{A}$ 的非空满子范畴 $\mathcal{C}$，满足：给定正合序列\footnote{依定义 \ref{homology-definition-exact}，这意味着 $\Im(A \to B) = \Ker(B \to C)$。}
$$
A \to B \to C
$$
若 $A, C \in \Ob(\mathcal{C})$，则 $B \in \Ob(\mathcal{C})$。
\item $\mathcal{A}$ 的一个{\it 弱 Serre 子范畴}是 $\mathcal{A}$ 的非空满子范畴 $\mathcal{C}$，满足：给定正合序列
$$
A_0 \to A_1 \to A_2 \to A_3 \to A_4
$$
若 $A_0, A_1, A_3, A_4$ 属于 $\mathcal{C}$，则 $A_2$ 也属于 $\mathcal{C}$。
\end{enumerate}
\end{definition}

\noindent
有些文献把第二个概念称为“厚”子范畴，另一些文献则把第一个概念称为“厚”子范畴。不过，Serre 子范畴这一名称似乎已普遍指上述第一个概念。注意在两种情形下，$\mathcal{C}$ 都是阿贝尔范畴，且包含函子 $\mathcal{C} \to \mathcal{A}$ 是全忠实正合函子。下面进一步刻画这两类子范畴。

\begin{lemma}
\label{homology-lemma-characterize-serre-subcategory}
设 $\mathcal{A}$ 为阿贝尔范畴，$\mathcal{C}$ 为 $\mathcal{A}$ 的子范畴。则 $\mathcal{C}$ 是 Serre 子范畴，当且仅当满足：
\begin{enumerate}
\item $0 \in \Ob(\mathcal{C})$；
\item $\mathcal{C}$ 是 $\mathcal{A}$ 的严格满子范畴；
\item $\mathcal{C}$ 的对象的任意子对象或商对象仍为 $\mathcal{C}$ 的对象；
\item 若 $A \in \Ob(\mathcal{A})$ 是 $\mathcal{C}$ 中对象的扩张，则 $A \in \Ob(\mathcal{C})$。
\end{enumerate}
此外，Serre 子范畴是阿贝尔范畴，且包含函子正合。
\end{lemma}

\begin{proof}
略。
\end{proof}

\begin{lemma}
\label{homology-lemma-characterize-weak-serre-subcategory}
设 $\mathcal{A}$ 为阿贝尔范畴，$\mathcal{C}$ 为 $\mathcal{A}$ 的子范畴。则 $\mathcal{C}$ 是弱 Serre 子范畴，当且仅当满足：
\begin{enumerate}
\item $0 \in \Ob(\mathcal{C})$；
\item $\mathcal{C}$ 是 $\mathcal{A}$ 的严格满子范畴；
\item $\mathcal{C}$ 中对象间态射在 $\mathcal{A}$ 中的核与余核仍在 $\mathcal{C}$ 中；
\item 若 $A \in \Ob(\mathcal{A})$ 是 $\mathcal{C}$ 中对象的扩张，则 $A \in \Ob(\mathcal{C})$。
\end{enumerate}
此外，弱 Serre 子范畴是阿贝尔范畴，且包含函子正合。
\end{lemma}

\begin{proof}
略。
\end{proof}

\begin{lemma}
\label{homology-lemma-kernel-exact-functor}
设 $\mathcal{A}$、$\mathcal{B}$ 为阿贝尔范畴，$F : \mathcal{A} \to \mathcal{B}$ 为正合函子。则 $\mathcal{A}$ 中满足 $F(C) = 0$ 的对象 $C$ 所成的满子范畴构成 $\mathcal{A}$ 的 Serre 子范畴。
\end{lemma}

\begin{proof}
略。
\end{proof}

\begin{definition}
\label{homology-definition-kernel-category}
设 $\mathcal{A}$、$\mathcal{B}$ 为阿贝尔范畴，$F : \mathcal{A} \to \mathcal{B}$ 为正合函子。$\mathcal{A}$ 中满足 $F(C) = 0$ 的对象 $C$ 所成的满子范畴称为{\it 函子 $F$ 的核}，有时记作 $\Ker(F)$。
\end{definition}

\noindent
阿贝尔范畴的任一 Serre 子范畴都是某个正合函子的核。在《例》，第 \ref{examples-section-serre-quotient-modulo-torsion-modules} 节中，我们对 Serre 最初的挠群例子讨论这一点。

\begin{lemma}
\label{homology-lemma-serre-subcategory-is-kernel}
设 $\mathcal{A}$ 为阿贝尔范畴，$\mathcal{C} \subset \mathcal{A}$ 为 Serre 子范畴。存在阿贝尔范畴 $\mathcal{A}/\mathcal{C}$ 与正合函子
$$
F : \mathcal{A} \longrightarrow \mathcal{A}/\mathcal{C}
$$
其中该函子本质满且核为 $\mathcal{C}$。范畴 $\mathcal{A}/\mathcal{C}$ 与函子 $F$ 由如下泛性质刻画：对任何满足 $\mathcal{C} \subset \Ker(G)$ 的正合函子 $G : \mathcal{A} \to \mathcal{B}$，存在唯一正合函子 $H : \mathcal{A}/\mathcal{C} \to \mathcal{B}$，使 $G = H \circ F$。
\end{lemma}

\begin{proof}
考虑 $\mathcal{A}$ 的如下箭头集合：
$$
S = \{f \in \text{Arrows}(\mathcal{A}) \mid
\Ker(f), \Coker(f) \in \Ob(\mathcal{C}) \}.
$$
我们断言 $S$ 是乘法系。为此须检验 MS1、MS2、MS3，见《范畴》，定义 \ref{categories-definition-multiplicative-system}。

\medskip\noindent
恒等态射显然属于 $S$。设 $f : A \to B$ 与 $g : B \to C$ 属于 $S$。有正合序列
$$
\begin{matrix}
0 \to \Ker(f) \to \Ker(gf) \to \Ker(g) \\
\Coker(f) \to \Coker(gf) \to \Coker(g) \to 0
\end{matrix}
$$
故 $gf \in S$，这证明 MS1。（事实上，类似论证表明 $S$ 是饱和乘法系，见《范畴》，定义 \ref{categories-definition-saturated-multiplicative-system}。）

\medskip\noindent
考虑实线图
$$
\xymatrix{
A \ar[d]_t \ar[r]_g & B \ar@{..>}[d]^s \\
C \ar@{..>}[r]^f & C \amalg_A B
}
$$
其中 $t \in S$。令 $W = C \amalg_A B =  \Coker((t, -g) : A \to C \oplus B)$。则 $\Ker(t) \to \Ker(s)$ 为满射态射，而 $\Coker(t) \to \Coker(s)$ 为同构。因此 $s$ 属于 $S$。这证明 LMS2；RMS2 的证明与之对偶。

\medskip\noindent
最后，设 $f, g : B \to C$ 为态射，$s : A \to B$ 属于 $S$，且 $f \circ s = g \circ s$。这意味着 $(f - g) \circ s = 0$，继而意味着 $I = \Im(f - g) \subset C$ 是 $\Coker(s)$ 的商，因而是 $\mathcal{C}$ 的对象。于是 $t : C \to C' = C/I$ 属于 $S$，且 $t \circ (f - g) = 0$，即 $t \circ f = t \circ g$。这证明 LMS3；RMS3 的证明与之对偶。

\medskip\noindent
已证 $S$ 为乘法系，令 $\mathcal{A}/\mathcal{C} = S^{-1}\mathcal{A}$，并令 $F$ 为局部化函子 $Q$。由引理 \ref{homology-lemma-localization-abelian}，范畴 $\mathcal{A}/\mathcal{C}$ 为阿贝尔范畴且 $F$ 正合。若 $X$ 属于 $F = Q$ 的核，则由引理 \ref{homology-lemma-kernel-localization}，$0 : X \to Z$ 属于 $S$，故 $X$ 是 $\mathcal{C}$ 的对象，即 $F$ 的核为 $\mathcal{C}$。

最后，若 $G$ 如引理所述，则 $G$ 把 $S$ 的每个元素变为同构。因此由局部化的泛性质（见《范畴》，引理 \ref{categories-lemma-properties-left-localization}），得到函子 $H : \mathcal{A}/\mathcal{C} \to \mathcal{B}$。还须证明 $H$ 正合。由引理 \ref{homology-lemma-exact-functor}，只需证明 $H$ 与取核、取余核可交换。设 $A \to B$ 为 $\mathcal{A}/\mathcal{C}$ 中的态射。可将 $A \to B$ 表示为 $fs^{-1}$，其中 $s : A' \to A$ 属于 $S$，$f : A' \to B$ 为 $\mathcal{A}$ 的任意态射。因为 $F = Q$ 把 $s$ 映为商范畴 $\mathcal{A}/\mathcal{C}$ 中的同构，只需证明 $H$ 与 $\mathcal{A}$ 的态射 $f : A \to B$ 的核和余核可交换。但此时 $H(f) = G(f)$，结论由 $G$ 正合得到。
\end{proof}

\begin{lemma}
\label{homology-lemma-quotient-by-kernel-exact-functor}
设 $\mathcal{A}$、$\mathcal{B}$ 为阿贝尔范畴，$F : \mathcal{A} \to \mathcal{B}$ 为正合函子。设 $\mathcal{C} \subset \mathcal{A}$ 为包含于 $F$ 的核中的 Serre 子范畴。则 $\mathcal{C} = \Ker(F)$，当且仅当诱导函子 $\overline{F} : \mathcal{A}/\mathcal{C} \to \mathcal{B}$（引理 \ref{homology-lemma-serre-subcategory-is-kernel}）是忠实函子。
\end{lemma}

\begin{proof}
以下不再另行说明地使用引理 \ref{homology-lemma-serre-subcategory-is-kernel} 的结果。“仅当”方向成立，因为按构造 $\overline{F}$ 的核为零。确实，若 $f : X \to Y$ 为 $\mathcal{A}/\mathcal{C}$ 中满足 $\overline{F}(f) = 0$ 的态射，则 $\overline{F}(\Im(f)) = \Im(\overline{F}(f)) = 0$；由对 $F$ 的核的假设，$\Im(f) = 0$，故 $f = 0$。

\medskip\noindent
对“若”方向，设 $X$ 为 $\mathcal{A}$ 中满足 $F(X) = 0$ 的对象。则 $\overline{F}(\text{id}_X) = \text{id}_{\overline{F}(X)} = 0$；由 $\overline{F}$ 的忠实性，$\mathcal{A}/\mathcal{C}$ 中 $\text{id}_X = 0$。所以 $\mathcal{A}/\mathcal{C}$ 中 $X = 0$，即 $X \in \Ob(\mathcal{C})$。
\end{proof}

\section{K 群}
\label{homology-section-K-groups}

\noindent
简单介绍阿贝尔范畴的 $K_0$。

\begin{definition}
\label{homology-definition-K-zero}
设 $\mathcal{A}$ 为阿贝尔范畴。以 $K_0(\mathcal{A})$ 表示{\it $\mathcal{A}$ 的第零 $K$ 群}。它是如下构造的阿贝尔群：取以 $\mathcal{A}$ 的对象为生成元的自由阿贝尔群，并对每个短正合序列 $0 \to A \to B \to C \to 0$ 加入关系 $[B] - [A] - [C] = 0$。
\end{definition}

\noindent
换言之，有表示
$$
\bigoplus_{A \to B \to C\text{ 短正合列}}
\mathbf{Z}[A \to B \to C]
\longrightarrow
\bigoplus_{A \in \Ob(\mathcal{A})}
\mathbf{Z}[A]
\longrightarrow
K_0(\mathcal{A})
\longrightarrow
0
$$
其中 $[A \to B \to C] \mapsto [B] - [A] - [C]$ 给出 $K_0(\mathcal{A})$ 的这一表示。短正合序列 $0 \to 0 \to 0 \to 0 \to 0$ 在 $K_0(\mathcal{A})$ 中给出关系 $[0] = 0$。除非另有说明，我们的范畴都是“小”范畴，故不存在集合论问题。《代数》，第 \ref{algebra-section-K-groups} 节计算了若干环上模范畴的 $K$ 群例子。

\begin{lemma}
\label{homology-lemma-exact-functor-K-groups}
设 $F : \mathcal{A} \to \mathcal{B}$ 为阿贝尔范畴之间的正合函子。则 $F$ 通过 $K_0(F)([A]) = [F(A)]$ 诱导 $K$ 群同态 $K_0(F) : K_0(\mathcal{A}) \to K_0(\mathcal{B})$。
\end{lemma}

\begin{proof}
不证自明。
\end{proof}

\noindent
设给定阿贝尔范畴 $\mathcal{A}$ 的对象 $M$ 与形如
\begin{equation}
\label{homology-equation-cyclic-complex}
\xymatrix{
\ldots \ar[r] &
M \ar[r]^\varphi &
M \ar[r]^\psi &
M \ar[r]^\varphi &
M \ar[r] & \ldots
}
\end{equation}
的复形。在这种情形下定义
$$
H^0(M, \varphi, \psi) = \Ker(\psi)/\Im(\varphi)
, \quad\text{且}\quad
H^1(M, \varphi, \psi) = \Ker(\varphi)/\Im(\psi).
$$

\begin{lemma}
\label{homology-lemma-serre-subcategory-K-groups}
设 $\mathcal{A}$ 为阿贝尔范畴，$\mathcal{C} \subset \mathcal{A}$ 为 Serre 子范畴，并令 $\mathcal{B} = \mathcal{A}/\mathcal{C}$。
\begin{enumerate}
\item 正合函子 $\mathcal{C} \to \mathcal{A}$ 与 $\mathcal{A} \to \mathcal{B}$ 诱导 $K$ 群正合序列
$$
K_0(\mathcal{C}) \to
K_0(\mathcal{A}) \to
K_0(\mathcal{B}) \to
0
$$
\item $K_0(\mathcal{C}) \to K_0(\mathcal{A})$ 的核等于所有形如
$$
[H^0(M, \varphi, \psi)] - [H^1(M, \varphi, \psi)]
$$
的元素之集合，其中 $(M, \varphi, \psi)$ 是 (\ref{homology-equation-cyclic-complex}) 中那样的复形，且它在 $\mathcal{B}$ 中变为正合；换言之，$H^0(M, \varphi, \psi)$ 与 $H^1(M, \varphi, \psi)$ 都是 $\mathcal{C}$ 的对象。
\end{enumerate}
\end{lemma}

\begin{proof}
证明 (1)。$K_0(\mathcal{A}) \to K_0(\mathcal{B})$ 显然为满射，且复合 $K_0(\mathcal{C}) \to K_0(\mathcal{A}) \to K_0(\mathcal{B})$ 为零。设 $x \in K_0(\mathcal{A})$ 映为 $K_0(\mathcal{B})$ 中的零。可写 $x = [A] - [A']$，其中 $A, A'$ 属于 $\mathcal{A}$（这是个有趣的练习）。以 $B, B'$ 表示 $\mathcal{B}$ 中的相应对象。$x$ 映为 $K_0(\mathcal{B})$ 中的零，意味着存在有限集 $I = I^+ \amalg I^{-}$，并且对每个 $i \in I$ 有 $\mathcal{B}$ 中的短正合序列
$$
0 \to B_i \to B'_i \to B''_i \to 0
$$
使在以 $\mathcal{B}$ 对象的同构类为生成元的自由阿贝尔群中
$$
[B] - [B'] = \sum\nolimits_{i \in I^{+}} ([B'_i] - [B_i] - [B''_i])
-
\sum\nolimits_{i \in I^{-}} ([B'_i] - [B_i] - [B''_i])
$$
可改写为
$$
[B]
+ \sum\nolimits_{i \in I^{+}} ([B_i] + [B''_i])
+ \sum\nolimits_{i \in I^{-}} [B'_i]
=
[B']
+ \sum\nolimits_{i \in I^{-}} ([B_i] + [B''_i])
+ \sum\nolimits_{i \in I^{+}} [B'_i].
$$
由于等式两边必须含有按重数计算相同的 $\mathcal{B}$ 对象同构类，这意味着存在双射
$$
\tau :
\{B\} \amalg \{B_i, B''_i; i \in I^+\} \amalg \{B'_i; i \in I^-\}
\longrightarrow
\{B'\} \amalg \{B_i, B''_i; i \in I^-\} \amalg \{B'_i; i \in I^+\}
$$
使 $N$ 与 $\tau(N)$ 在 $\mathcal{B}$ 中同构。引理 \ref{homology-lemma-serre-subcategory-is-kernel} 与 \ref{homology-lemma-localization-abelian} 的证明表明，可对 $i \in I$ 选取 $\mathcal{A}$ 中的短正合序列
$$
0 \to A_i \to A'_i \to A''_i \to 0
$$
使 $B_i, B'_i, B''_i$ 分别同构于 $A_i, A'_i, A''_i$ 在 $\mathcal{B}$ 中的像。这蕴含相应双射
$$
\tau :
\{A\} \amalg \{A_i, A''_i; i \in I^+\} \amalg \{A'_i; i \in I^-\}
\longrightarrow
\{A'\} \amalg \{A_i, A''_i; i \in I^-\} \amalg \{A'_i; i \in I^+\}
$$
满足：$M$ 与 $\tau(M)$ 是 $\mathcal{A}$ 的对象，且在 $\mathcal{B}$ 中变为同构。这意味着 $[M] - [\tau(M)]$ 属于 $K_0(\mathcal{C}) \to K_0(\mathcal{A})$ 的像。确实，$\mathcal{B}$ 中该同构由 $\mathcal{A}$ 中的图 $M \leftarrow M' \rightarrow \tau(M)$ 给出，其中 $M' \to M$ 与 $M' \to \tau(M)$ 的核和余核都在 $\mathcal{C}$ 中。逆推以上过程可知 $x = [A] - [A']$ 属于 $K_0(\mathcal{C}) \to K_0(\mathcal{A})$ 的像，(1) 得证。

\medskip\noindent
证明 (2)。证明与 (1) 类似，但需要稍多的记账。先注意，任何形如 $[H^0(M, \varphi, \psi)] - [H^1(M, \varphi, \psi)]$ 的类在 $K_0(\mathcal{A})$ 中为零，因为
\begin{align*}
0 & = [M] - [M] \\
& =  [\Ker(\varphi)] + [\Im(\varphi)]
- [\Ker(\psi)] - [\Im(\psi)] \\
& =
[\Ker(\varphi)/\Im(\psi)] -
[\Ker(\psi)/\Im(\varphi)] \\
& = [H^1(M, \varphi, \psi)] - [H^0(M, \varphi, \psi)],
\end{align*}
正如所需。因此只须证明 $K_0(\mathcal{C}) \to K_0(\mathcal{A})$ 的任一核中元素都有这种形式。

\medskip\noindent
$K_0(\mathcal{C})$ 中任一元素 $x$ 都可表示为 $\mathcal{C}$ 的两个对象之差 $x = [P] - [Q]$（有趣的练习）。假设该元素在 $K_0(\mathcal{A})$ 中映为零。这意味着存在：
\begin{enumerate}
\item 有限集 $I = I^{+} \amalg I^{-}$；
\item 对每个 $i \in I$，$\mathcal{A}$ 中的短正合序列 $0 \to A_i \to B_i \to C_i \to 0$；
\end{enumerate}
使在以 $\mathcal{A}$ 的对象为生成元的自由阿贝尔群中
$$
[P] - [Q] =
\sum\nolimits_{i \in I^{+}} ([B_i] - [A_i] - [C_i])
-
\sum\nolimits_{i \in I^{-}} ([B_i] - [A_i] - [C_i])
$$
可改写为
$$
[P]
+ \sum\nolimits_{i \in I^{+}} ([A_i] + [C_i])
+ \sum\nolimits_{i \in I^{-}} [B_i]
=
[Q]
+ \sum\nolimits_{i \in I^{-}} ([A_i] + [C_i])
+ \sum\nolimits_{i \in I^{+}} [B_i].
$$
由于等式两边必须含有按重数计算相同的 $\mathcal{A}$ 对象，故其中出现的项之间存在双射 $\tau$。令
$$
T^{+} =
\{p\}\ \amalg\ \{a, c\} \times I^{+}\ \amalg\ \{b\} \times I^{-}
$$
及
$$
T^{-} =
\{q\}\ \amalg\ \{a, c\} \times I^{-}\ \amalg\ \{b\} \times I^{+}.
$$
令 $T = T^{+} \amalg T^{-} = \{p, q\} \amalg \{a, b, c\} \times I$。对 $t \in T$ 定义
$$
O(t)
=
\left\{
\begin{matrix}
P & \text{若} & t = p \\
Q & \text{若} & t = q \\
A_i & \text{若} & t = (a, i) \\
B_i & \text{若} & t = (b, i) \\
C_i & \text{若} & t = (c, i)
\end{matrix}
\right.
$$
于是可把 $\tau : T^{+} \to T^{-}$ 看作双射，使对所有 $t \in T^{+}$ 都有 $O(t) = O(\tau(t))$。令 $t^{-}_0 = \tau(p)$，并令 $t^{+}_0 \in T^{+}$ 为满足 $\tau(t^{+}_0) = q$ 的唯一元素。考虑对象
$$
M^{+} = \bigoplus\nolimits_{t \in T^{+}} O(t)
$$
借助 $\tau$ 可知它等于对象
$$
M^{-} = \bigoplus\nolimits_{t \in T^{-}} O(t)
$$
考虑映射
$$
\varphi : M^{+} \longrightarrow M^{-}
$$
它在对应于 $t = (a, i)$、$i \in I^{+}$ 的直和项 $O(t) = A_i$ 上，使用映射 $A_i \to B_i$ 进入 $M^{-}$ 的直和项 $O((b, i)) = B_i$；在对应于 $(b, i)$、$i \in I^{-}$ 的直和项 $O(t) = B_i$ 上，使用映射 $B_i \to C_i$ 进入 $M^{-}$ 的直和项 $O((c, i)) = C_i$。该映射在对应于 $p$ 与 $(c, i)$、$i \in I^{+}$ 的直和项上为零。

类似地，考虑映射
$$
\psi : M^{-} \longrightarrow M^{+}
$$
它在对应于 $t = (a, i)$、$i \in I^{-}$ 的直和项 $O(t) = A_i$ 上，使用映射 $A_i \to B_i$ 进入 $M^{+}$ 的直和项 $O((b, i)) = B_i$；在对应于 $(b, i)$、$i \in I^{+}$ 的直和项 $O(t) = B_i$ 上，使用映射 $B_i \to C_i$ 进入 $M^{+}$ 的直和项 $O((c, i)) = C_i$。该映射在对应于 $q$ 与 $(c, i)$、$i \in I^{-}$ 的直和项上为零。

\medskip\noindent
注意，$\varphi$ 的核等于直和项 $P$、各直和项 $O((c, i)) = C_i$（$i \in I^{+}$），以及各直和项 $O((b, i)) = B_i$ 内的子对象 $A_i$（$i \in I^{-}$）之直和。$\psi$ 的像等于各直和项 $O((c, i)) = C_i$（$i \in I^{+}$）及各直和项 $O((b, i)) = B_i$ 内的子对象 $A_i$（$i \in I^{-}$）之直和。换言之，
$$
P \cong \Ker(\varphi)/\Im(\psi).
$$
完全同样地，
$$
Q \cong \Ker(\psi)/\Im(\varphi).
$$
如上所述，双射 $\tau$ 的存在表明 $M^{+} = M^{-}$，引理遂得证。
\end{proof}

\section{上同调 \texorpdfstring{\zhdelta}{δ}-函子}
\label{homology-section-cohomological-delta-functor}

\begin{definition}
\label{homology-definition-cohomological-delta-functor}
设 $\mathcal{A}, \mathcal{B}$ 为阿贝尔范畴。从 $\mathcal{A}$ 到
$\mathcal{B}$ 的一个{\it 上同调 $\delta$-函子}，简称一个
{\it $\delta$-函子}，由下列数据给出：
\begin{enumerate}
\item 一族加性函子 $F^n : \mathcal{A} \to \mathcal{B}$，$n \geq 0$；以及
\item 对 $\mathcal{A}$ 中每个短正合列 $0 \to A \to B \to C \to 0$，
一族 $\mathcal{B}$ 中的态射
$\delta_{A \to B \to C} : F^n(C) \to F^{n + 1}(A)$，$n \geq 0$。
\end{enumerate}
要求这些数据满足下列公理：
\begin{enumerate}
\item 对每个如上的短正合列，序列
$$
\xymatrix{
0 \ar[r] &
F^0(A) \ar[r] &
F^0(B) \ar[r] &
F^0(C) \ar[lld]^{\delta_{A \to B \to C}} \\
 &
F^1(A) \ar[r] &
F^1(B) \ar[r] &
F^1(C) \ar[lld]^{\delta_{A \to B \to C}} \\
 &
F^2(A) \ar[r] &
F^2(B) \ar[r] &
\ldots
}
$$
正合；以及
\item 对 $\mathcal{A}$ 中短正合列的每个态射
$(A \to B \to C) \to (A' \to B' \to C')$，图表
$$
\xymatrix{
F^n(C) \ar[d] \ar[rr]_{\delta_{A \to B \to C}} & & F^{n + 1}(A) \ar[d] \\
F^n(C') \ar[rr]^{\delta_{A' \to B' \to C'}} & & F^{n + 1}(A')
}
$$
交换。
\end{enumerate}
\end{definition}

\noindent
特别地，注意这蕴含 $F^0$ 左正合。

\begin{definition}
\label{homology-definition-morphism-delta-functors}
设 $\mathcal{A}, \mathcal{B}$ 为阿贝尔范畴，并设
$(F^n, \delta_F)$ 与 $(G^n, \delta_G)$ 是从 $\mathcal{A}$ 到
$\mathcal{B}$ 的 $\delta$-函子。从 $F$ 到 $G$ 的一个
{\it $\delta$-函子态射}，是函子变换
$t^n : F^n \to G^n$（$n \geq 0$）的一族，使得对 $\mathcal{A}$ 中
每个短正合列 $0 \to A \to B \to C \to 0$，图表
$$
\xymatrix{
F^n(C) \ar[d]_{t^n} \ar[rr]_{\delta_{F, A \to B \to C}} &
& F^{n + 1}(A) \ar[d]^{t^{n + 1}} \\
G^n(C) \ar[rr]^{\delta_{G, A \to B \to C}} & & G^{n + 1}(A)
}
$$
交换。
\end{definition}

\begin{definition}
\label{homology-definition-universal-delta-functor}
设 $\mathcal{A}, \mathcal{B}$ 为阿贝尔范畴，并设
$F = (F^n, \delta_F)$ 是从 $\mathcal{A}$ 到 $\mathcal{B}$ 的
$\delta$-函子。若对每个 $\delta$-函子 $G = (G^n, \delta_G)$ 以及
任意函子态射 $t : F^0 \to G^0$，都存在唯一的 $\delta$-函子态射
$\{t^n\}_{n \geq 0} : F \to G$ 使得 $t = t^0$，则称 $F$ 是一个
{\it 泛 $\delta$-函子}。
\end{definition}

\begin{lemma}
\label{homology-lemma-efface-implies-universal}
设 $\mathcal{A}, \mathcal{B}$ 为阿贝尔范畴，并设
$F = (F^n, \delta_F)$ 是从 $\mathcal{A}$ 到 $\mathcal{B}$ 的
$\delta$-函子。假设对每个 $n > 0$ 以及任意
$A \in \Ob(\mathcal{A})$，都存在一个单射态射 $u : A \to B$
（依赖于 $A$ 与 $n$），使得
$F^n(u) : F^n(A) \to F^n(B)$ 为零。那么 $F$ 是泛 $\delta$-函子。
\end{lemma}

\begin{proof}
设 $G = (G^n, \delta_G)$ 是从 $\mathcal{A}$ 到 $\mathcal{B}$ 的
$\delta$-函子，并设 $t : F^0 \to G^0$ 为函子态射。我们须证明存在
唯一的 $\delta$-函子态射 $\{t^n\}_{n \geq 0} : F \to G$，使得
$t = t^0$。我们对 $n$ 归纳构造 $t^n$。当 $n = 0$ 时令
$t^0 = t$。假设已经唯一地构造了 $i \leq n$ 时的函子变换序列
$t^i$，并且它与次数 $\leq n$ 中的映射 $\delta$ 相容。

\medskip\noindent
令 $A \in \Ob(\mathcal{A})$。由假设可选取一个嵌入
$u : A \to B$，使得 $F^{n + 1}(u) = 0$。令 $C = B/u(A)$。
短正合列 $0 \to A \to B \to C \to 0$ 对 $\delta$-函子 $F$
给出的长正合上同调列，连同我们对 $u$ 的选择，表明
$F^{n + 1}(A) = \Coker(F^n(B) \to F^n(C))$。由于已经定义了
$t^n$，我们可令
$$
t^{n + 1}_A : F^{n + 1}(A) \to G^{n + 1}(A)
$$
为使下图交换的唯一映射：
$$
\xymatrix{
\Coker(F^n(B) \to F^n(C)) \ar[r]_{t^n}
\ar[d]_{\delta_{F, A \to B \to C}} &
\Coker(G^n(B) \to G^n(C))
\ar[d]^{\delta_{G, A \to B \to C}} \\
F^{n + 1}(A) \ar[r]^{t^{n + 1}_A} &
G^{n + 1}(A)
}
$$
显然，所加要求唯一地决定了这一映射。我们略去它定义函子变换的验证。
\end{proof}

\begin{lemma}
\label{homology-lemma-uniqueness-universal-delta-functor}
设 $\mathcal{A}, \mathcal{B}$ 为阿贝尔范畴，并设
$F : \mathcal{A} \to \mathcal{B}$ 为函子。若存在一个从
$\mathcal{A}$ 到 $\mathcal{B}$ 的泛 $\delta$-函子
$(F^n, \delta_F)$，且 $F^0 = F$，则它在 $\delta$-函子的唯一同构意义下确定。
\end{lemma}

\begin{proof}
由定义立即可得。
\end{proof}

\section{复形}
\label{homology-section-complexes}

\noindent
当然，链复形与上链复形的概念互为对偶，本节两部分只需选读其一即可；
因此任选你喜欢的一部分。（严格说来，这并不完全奏效，因为各项的编号约定
并不适合在链复形与上链复形之间直接转换。）

\medskip\noindent
加性范畴 $\mathcal{A}$ 中的一个{\it 链复形 $A_\bullet$}，是
$\mathcal{A}$ 中的复形
$$
\ldots \to
A_{n + 1} \xrightarrow{d_{n + 1}}
A_n \xrightarrow{d_n}
A_{n - 1} \to
\ldots
$$
换言之，对每个 $i \in \mathbf{Z}$ 给定 $\mathcal{A}$ 的对象 $A_i$，
并对每个 $i \in \mathbf{Z}$ 给定态射 $d_i : A_i \to A_{i - 1}$，
且对所有 $i$ 均有 $d_{i - 1} \circ d_i = 0$。一个
{\it 链复形态射 $f : A_\bullet \to B_\bullet$}，由一族态射
$f_i : A_i \to B_i$ 给出，并要求所有图表
$$
\xymatrix{
A_i \ar[r]_{d_i} \ar[d]_{f_i} & A_{i - 1} \ar[d]^{f_{i - 1}} \\
B_i \ar[r]^{d_i} & B_{i - 1}
}
$$
交换。$\mathcal{A}$ 的{\it 链复形范畴}记作
$\text{Ch}(\mathcal{A})$。其中由形如
$$
\ldots \to A_2 \to A_1 \to A_0 \to 0 \to 0 \to \ldots
$$
的对象组成的全子范畴记作 $\text{Ch}_{\geq 0}(\mathcal{A})$。
换言之，链复形 $A_\bullet$ 属于
$\text{Ch}_{\geq 0}(\mathcal{A})$，当且仅当对所有 $i < 0$ 均有
$A_i = 0$。

\medskip\noindent
给定加性范畴 $\mathcal{A}$，通过函子
$$
\mathcal{A} \longrightarrow \text{Ch}(\mathcal{A}),\quad
A \longmapsto (\ldots \to 0 \to A \to 0 \to \ldots)
$$
把 $\mathcal{A}$ 等同于 $\text{Ch}(\mathcal{A})$ 中由仅在次数 $0$
处非零的链复形组成的全子范畴。略用记号时，我们常把右端对象简记为
$A$。若想强调把 $A$ 看作链复形，有时可使用记号 $A[0]$；见
第 \ref{homology-section-homotopy-shift} 节。

\medskip\noindent
链复形态射 $f, g : A_\bullet \to B_\bullet$ 之间的一个
{\it 同伦 $h$}，是一族态射 $h_i : A_i \to B_{i + 1}$，使得对所有
$i$ 均有
$$
f_i - g_i = d_{i + 1} \circ h_i + h_{i - 1} \circ d_i
$$
若 $f$ 与 $g$ 之间存在同伦，则称两个态射
$f, g : A_\bullet \to B_\bullet$ {\it 同伦}。显然，链复形、链复形态射
以及链复形态射之间的同伦这些概念，即使在预加性范畴中也有意义。

\begin{lemma}
\label{homology-lemma-compose-homotopy}
\begin{slogan}
$\text{Ch}(\mathcal{A})$ 的 Hom 函子保持同伦关系。
\end{slogan}
设 $\mathcal{A}$ 为加性范畴。设 $f, g : B_\bullet \to C_\bullet$
为链复形态射，并给定链复形态射
$a : A_\bullet \to B_\bullet$ 与 $c : C_\bullet \to D_\bullet$。
若 $\{h_i : B_i \to C_{i + 1}\}$ 定义了 $f$ 与 $g$ 之间的同伦，
则 $\{c_{i + 1} \circ h_i \circ a_i\}$ 定义了
$c \circ f \circ a$ 与 $c \circ g \circ a$ 之间的同伦。
\end{lemma}

\begin{proof}
略。
\end{proof}

\noindent
特别地，这意味着可以定义以同伦类为态射的链复形范畴。稍后我们将回到这一点。

\begin{definition}
\label{homology-definition-homotopy-equivalent}
设 $\mathcal{A}$ 为加性范畴。若存在态射
$b : B_\bullet \to A_\bullet$，使得 $b \circ a$ 与 $\text{id}_A$
之间存在同伦，并且 $a \circ b$ 与 $\text{id}_B$ 之间存在同伦，
则称态射 $a : A_\bullet \to B_\bullet$ 为{\it 同伦等价}。
若 $A_\bullet$ 与 $B_\bullet$ 之间存在这样的态射，则称
$A_\bullet$ 与 $B_\bullet$ {\it 同伦等价}。
\end{definition}

\noindent
换言之，两个复形同伦等价，当且仅当它们在以同伦类为态射的复形范畴中同构。

\begin{lemma}
\label{homology-lemma-cat-chain-abelian}
设 $\mathcal{A}$ 为阿贝尔范畴。
\begin{enumerate}
\item $\mathcal{A}$ 中链复形的范畴是阿贝尔范畴。
\item 复形态射 $f : A_\bullet \to B_\bullet$ 是单射态射，当且仅当每个
$f_n : A_n \to B_n$ 都是单射态射。
\item 复形态射 $f : A_\bullet \to B_\bullet$ 是满射态射，当且仅当每个
$f_n : A_n \to B_n$ 都是满射态射。
\item 链复形序列
$$
A_\bullet \xrightarrow{f} B_\bullet \xrightarrow{g} C_\bullet
$$
在 $B_\bullet$ 处正合，当且仅当每个序列
$$
A_i \xrightarrow{f_i} B_i \xrightarrow{g_i} C_i
$$
都在 $B_i$ 处正合。
\end{enumerate}
\end{lemma}

\begin{proof}
略。
\end{proof}

\noindent
对任意 $i \in \mathbf{Z}$，阿贝尔范畴中链复形 $A_\bullet$ 的第
$i$ 个{\it 同调群}定义为
$$
H_i(A_\bullet) = \Ker(d_i)/\Im(d_{i + 1}).
$$
若 $f : A_\bullet \to B_\bullet$ 是 $\mathcal{A}$ 中的链复形态射，
则得到诱导态射 $H_i(f) : H_i(A_\bullet) \to H_i(B_\bullet)$，因为显然
$f_i(\Ker(d_i : A_i \to A_{i - 1})) \subset
\Ker(d_i : B_i \to B_{i - 1})$，而 $\Im(d_{i + 1})$ 也同样如此。
因此得到函子
$$
H_i : \text{Ch}(\mathcal{A}) \longrightarrow \mathcal{A}.
$$

\begin{definition}
\label{homology-definition-quasi-isomorphism}
设 $\mathcal{A}$ 为阿贝尔范畴。
\begin{enumerate}
\item 若诱导映射 $H_i(f) : H_i(A_\bullet) \to H_i(B_\bullet)$ 对所有
$i \in \mathbf{Z}$ 都是同构，则称链复形态射
$f : A_\bullet \to B_\bullet$ 为{\it 拟同构}。
\item 若链复形 $A_\bullet$ 的所有同调对象 $H_i(A_\bullet)$ 都为零，
则称它{\it 无环}。
\end{enumerate}
\end{definition}

\begin{lemma}
\label{homology-lemma-map-homology-homotopy}
设 $\mathcal{A}$ 为阿贝尔范畴。
\begin{enumerate}
\item 若映射 $f, g : A_\bullet \to B_\bullet$ 同伦，则诱导映射
$H_i(f)$ 与 $H_i(g)$ 相等。
\item 若映射 $f : A_\bullet \to B_\bullet$ 是同伦等价，则 $f$ 是拟同构。
\end{enumerate}
\end{lemma}

\begin{proof}
略。
\end{proof}

\begin{lemma}
\label{homology-lemma-long-exact-sequence-chain}
设 $\mathcal{A}$ 为阿贝尔范畴。假设
$$
0 \to
A_\bullet \to
B_\bullet \to
C_\bullet \to
0
$$
是 $\mathcal{A}$ 中链复形的短正合列。那么存在典范的长正合同调列
$$
\xymatrix{
\ldots & \ldots & \ldots \ar[lld] \\
H_i(A_\bullet) \ar[r] & H_i(B_\bullet) \ar[r] & H_i(C_\bullet) \ar[lld] \\
H_{i - 1}(A_\bullet) \ar[r] &
H_{i - 1}(B_\bullet) \ar[r] &
H_{i - 1}(C_\bullet) \ar[lld] \\
\ldots & \ldots & \ldots \\
}
$$
\end{lemma}

\begin{proof}
略。各映射来自把蛇引理 \ref{homology-lemma-snake} 应用于图表
$$
\xymatrix{
&
A_i/\Im(d_{A, i + 1}) \ar[r] \ar[d]^{d_{A, i}} &
B_i/\Im(d_{B, i + 1}) \ar[r] \ar[d]^{d_{B, i}} &
C_i/\Im(d_{C, i + 1}) \ar[r] \ar[d]^{d_{C, i}} &
0 \\
0 \ar[r] &
\Ker(d_{A, i - 1}) \ar[r] &
\Ker(d_{B, i - 1}) \ar[r] &
\Ker(d_{C, i - 1}) &
}
$$
\end{proof}

\noindent
加性范畴 $\mathcal{A}$ 中的一个{\it 上链复形 $A^\bullet$}，是
$\mathcal{A}$ 中的复形
$$
\ldots \to
A^{n - 1} \xrightarrow{d^{n - 1}}
A^n \xrightarrow{d^n}
A^{n + 1} \to
\ldots
$$
换言之，对每个 $i \in \mathbf{Z}$ 给定 $\mathcal{A}$ 的对象 $A^i$，
并对每个 $i \in \mathbf{Z}$ 给定态射 $d^i : A^i \to A^{i + 1}$，
且对所有 $i$ 均有 $d^{i + 1} \circ d^i = 0$。一个
{\it 上链复形态射 $f : A^\bullet \to B^\bullet$}，由一族态射
$f^i : A^i \to B^i$ 给出，并要求所有图表
$$
\xymatrix{
A^i \ar[r]_{d^i} \ar[d]_{f^i} & A^{i + 1} \ar[d]^{f^{i + 1}} \\
B^i \ar[r]^{d^i} & B^{i + 1}
}
$$
交换。$\mathcal{A}$ 的{\it 上链复形范畴}记作
$\text{CoCh}(\mathcal{A})$。其中由形如
$$
\ldots \to 0 \to 0 \to A^0 \to A^1 \to A^2 \to \ldots
$$
的对象组成的全子范畴记作 $\text{CoCh}_{\geq 0}(\mathcal{A})$。
换言之，上链复形 $A^\bullet$ 属于子范畴
$\text{CoCh}_{\geq 0}(\mathcal{A})$，当且仅当对所有 $i < 0$ 均有
$A^i = 0$。

\medskip\noindent
给定加性范畴 $\mathcal{A}$，通过函子
$$
\mathcal{A} \longrightarrow \text{CoCh}(\mathcal{A}),\quad
A \longmapsto (\ldots \to 0 \to A \to 0 \to \ldots)
$$
把 $\mathcal{A}$ 等同于 $\text{CoCh}(\mathcal{A})$ 中由仅在次数 $0$
处非零的上链复形组成的全子范畴。略用记号时，我们常把右端对象简记为
$A$。若想强调把 $A$ 看作上链复形，有时可使用记号 $A[0]$；见
第 \ref{homology-section-homotopy-shift} 节。

\medskip\noindent
上链复形态射 $f, g : A^\bullet \to B^\bullet$ 之间的一个
{\it 同伦 $h$}，是一族态射 $h^i : A^i \to B^{i - 1}$，使得对所有
$i$ 均有
$$
f^i - g^i = d^{i - 1} \circ h^i + h^{i + 1} \circ d^i
$$
若 $f$ 与 $g$ 之间存在同伦，则称两个态射
$f, g : A^\bullet \to B^\bullet$ {\it 同伦}。显然，上链复形、上链复形态射
以及上链复形态射之间的同伦这些概念，即使在预加性范畴中也有意义。

\begin{lemma}
\label{homology-lemma-compose-homotopy-cochain}
设 $\mathcal{A}$ 为加性范畴。设 $f, g : B^\bullet \to C^\bullet$
为上链复形态射，并给定上链复形态射
$a : A^\bullet \to B^\bullet$ 与 $c : C^\bullet \to D^\bullet$。
若 $\{h^i : B^i \to C^{i - 1}\}$ 定义了 $f$ 与 $g$ 之间的同伦，
则 $\{c^{i - 1} \circ h^i \circ a^i\}$ 定义了
$c \circ f \circ a$ 与 $c \circ g \circ a$ 之间的同伦。
\end{lemma}

\begin{proof}
略。
\end{proof}

\noindent
特别地，这意味着可以定义以同伦类为态射的上链复形范畴。稍后我们将回到这一点。

\begin{definition}
\label{homology-definition-homotopy-equivalent-cochain}
设 $\mathcal{A}$ 为加性范畴。若存在态射
$b : B^\bullet \to A^\bullet$，使得 $b \circ a$ 与 $\text{id}_A$
之间存在同伦，并且 $a \circ b$ 与 $\text{id}_B$ 之间存在同伦，
则称态射 $a : A^\bullet \to B^\bullet$ 为{\it 同伦等价}。
若 $A^\bullet$ 与 $B^\bullet$ 之间存在这样的态射，则称
$A^\bullet$ 与 $B^\bullet$ {\it 同伦等价}。
\end{definition}

\noindent
换言之，两个复形同伦等价，当且仅当它们在以同伦类为态射的复形范畴中同构。

\begin{lemma}
\label{homology-lemma-cat-cochain-abelian}
设 $\mathcal{A}$ 为阿贝尔范畴。
\begin{enumerate}
\item $\mathcal{A}$ 中上链复形的范畴是阿贝尔范畴。
\item 上链复形态射 $f : A^\bullet \to B^\bullet$ 是单射态射，当且仅当每个
$f^n : A^n \to B^n$ 都是单射态射。
\item 上链复形态射 $f : A^\bullet \to B^\bullet$ 是满射态射，当且仅当每个
$f^n : A^n \to B^n$ 都是满射态射。
\item 上链复形序列
$$
A^\bullet \xrightarrow{f} B^\bullet \xrightarrow{g} C^\bullet
$$
在 $B^\bullet$ 处正合，当且仅当每个序列
$$
A^i \xrightarrow{f^i} B^i \xrightarrow{g^i} C^i
$$
都在 $B^i$ 处正合。
\end{enumerate}
\end{lemma}

\begin{proof}
略。
\end{proof}

\noindent
对任意 $i \in \mathbf{Z}$，上链复形 $A^\bullet$ 的第
$i$ 个{\it 上同调群}定义为
$$
H^i(A^\bullet) = \Ker(d^i)/\Im(d^{i - 1}).
$$
若 $f : A^\bullet \to B^\bullet$ 是 $\mathcal{A}$ 中的上链复形态射，
则得到诱导态射 $H^i(f) : H^i(A^\bullet) \to H^i(B^\bullet)$，因为显然
$f^i(\Ker(d^i : A^i \to A^{i + 1})) \subset
\Ker(d^i : B^i \to B^{i + 1})$，而 $\Im(d^{i - 1})$ 也同样如此。
因此得到函子
$$
H^i : \text{CoCh}(\mathcal{A}) \longrightarrow \mathcal{A}.
$$

\begin{definition}
\label{homology-definition-quasi-isomorphism-cochain}
设 $\mathcal{A}$ 为阿贝尔范畴。
\begin{enumerate}
\item 若诱导映射 $H^i(f) : H^i(A^\bullet) \to H^i(B^\bullet)$ 对所有
$i \in \mathbf{Z}$ 都是同构，则称 $\mathcal{A}$ 的上链复形态射
$f : A^\bullet \to B^\bullet$ 为{\it 拟同构}。
\item 若上链复形 $A^\bullet$ 的所有上同调对象 $H^i(A^\bullet)$ 都为零，
则称它{\it 无环}。
\end{enumerate}
\end{definition}

\begin{lemma}
\label{homology-lemma-map-cohomology-homotopy-cochain}
设 $\mathcal{A}$ 为阿贝尔范畴。
\begin{enumerate}
\item 若映射 $f, g : A^\bullet \to B^\bullet$ 同伦，则诱导映射
$H^i(f)$ 与 $H^i(g)$ 相等。
\item 若 $f : A^\bullet \to B^\bullet$ 是同伦等价，则 $f$ 是拟同构。
\end{enumerate}
\end{lemma}

\begin{proof}
略。
\end{proof}

\begin{lemma}
\label{homology-lemma-long-exact-sequence-cochain}
\begin{slogan}
复形的短正合列导出（上）同调的长正合列。
\end{slogan}
设 $\mathcal{A}$ 为阿贝尔范畴。假设
$$
0 \to
A^\bullet \to
B^\bullet \to
C^\bullet \to
0
$$
是 $\mathcal{A}$ 中上链复形的短正合列。那么存在长正合上同调列
$$
\xymatrix{
\ldots & \ldots & \ldots \ar[lld] \\
H^i(A^\bullet) \ar[r] &
H^i(B^\bullet) \ar[r] &
H^i(C^\bullet) \ar[lld] \\
H^{i + 1}(A^\bullet) \ar[r] &
H^{i + 1}(B^\bullet) \ar[r] &
H^{i + 1}(C^\bullet) \ar[lld] \\
\ldots & \ldots & \ldots \\
}
$$
此构造产生关于短正合列函子性且与
定义 \ref{homology-definition-cohomology-shift} 中的平移相容的长正合上同调列。
\end{lemma}

\begin{proof}
对水平映射 $H^i(A^\bullet) \to H^i(B^\bullet)$ 与
$H^i(B^\bullet) \to H^i(C^\bullet)$，使用上文所述 $H^i$ 的函子性。
对“边界映射” $H^i(C^\bullet) \to H^{i + 1}(A^\bullet)$，使用把蛇引理
\ref{homology-lemma-snake} 应用于图表
$$
\xymatrix{
&
A^i/\Im(d_A^{i - 1}) \ar[r] \ar[d]^{d_A^i} &
B^i/\Im(d_B^{i - 1}) \ar[r] \ar[d]^{d_B^i} &
C^i/\Im(d_C^{i - 1}) \ar[r] \ar[d]^{d_C^i} &
0 \\
0 \ar[r] &
\Ker(d_A^{i + 1}) \ar[r] &
\Ker(d_B^{i + 1}) \ar[r] &
\Ker(d_C^{i + 1}) &
}
$$
所得的映射 $\delta$。这是可行的，因为右侧竖直映射的核等于
$H^i(C^\bullet)$，而左侧竖直映射的余核等于
$H^{i + 1}(A^\bullet)$。长序列的正合性就是引理 \ref{homology-lemma-snake}
第 (2) 部分中的正合性。% P02-E0008: source line 3246 has “exactnesss”.
函子性由引理 \ref{homology-lemma-snake-natural} 给出；与平移的相容性由定义立即可得。
\end{proof}

\section{同伦与平移函子}
\label{homology-section-homotopy-shift}

\noindent
令人烦恼的是，任何同调代数的讨论都不得不涉及符号与指标\footnote{如果你发现
符号错误，或能改进我们的约定，请告知我们。}。

\begin{definition}
\label{homology-definition-shift}
设 $\mathcal{A}$ 为加性范畴，并设 $A_\bullet$ 为边界映射为
$d_{A, n} : A_n \to A_{n - 1}$ 的链复形。对任意
$k \in \mathbf{Z}$，如下定义{\it 平移 $k$ 的链复形 $A[k]_\bullet$}：
\begin{enumerate}
\item 令 $A[k]_n = A_{n + k}$；以及
\item 令 $d_{A[k], n} : A[k]_n \to A[k]_{n - 1}$ 等于
$d_{A[k], n} = (-1)^k d_{A, n + k}$。
\end{enumerate}
若 $f : A_\bullet \to B_\bullet$ 是链复形态射，则令
$f[k] : A[k]_\bullet \to B[k]_\bullet$ 为满足
$f[k]_n = f_{k + n}$ 的链复形态射。
\end{definition}

\noindent
当然，这意味着有函子
$[k] : \text{Ch}(\mathcal{A}) \to \text{Ch}(\mathcal{A})$，它们彼此
严格交换（无需插入任何函子同构），并且
$A[k][l]_\bullet = A[k + l]_\bullet$、
$[0] = \text{id}_{\text{Ch}(\mathcal{A})}$。

\medskip\noindent
回忆我们把 $\mathcal{A}$ 看作 $\text{Ch}(\mathcal{A})$ 的全子范畴；
见第 \ref{homology-section-complexes} 节。因此，对 $\mathcal{A}$ 的任意对象
$A$，记号 $A[k]$ 表示唯一的链复形：除次数 $-k$ 处为 $A$ 外，
其他各次均为零。

\begin{definition}
\label{homology-definition-homology-shift}
设 $\mathcal{A}$ 为阿贝尔范畴，并设 $A_\bullet$ 为边界映射为
$d_{A, n} : A_n \to A_{n - 1}$ 的链复形。对任意
$k \in \mathbf{Z}$，借助等同 $A_{i + k} = A[k]_i$，我们作等同
{\it $H_{i + k}(A_\bullet) \rightarrow H_i(A[k]_\bullet)$}。
\end{definition}

\noindent
这一等同关于 $A_\bullet$ 具有函子性。注意，因为该定义不涉及符号，
实际上我们得到所有同调对象 $H_{i - k}(A[k]_\bullet)$ 的相容等同系；
它们还与等同 $A[k][l]_\bullet = A[k + l]_\bullet$ 以及
$[0] = \text{id}_{\text{Ch}(\mathcal{A})}$ 相容。

\medskip\noindent
设 $\mathcal{A}$ 为加性范畴。假设 $A_\bullet$ 与 $B_\bullet$ 是链复形，
$a, b : A_\bullet \to B_\bullet$ 是链复形态射，并且
$\{h_i : A_i \to B_{i + 1}\}$ 是 $a$ 与 $b$ 之间的同伦。回忆这意味着
$a_i - b_i = d_{i + 1} \circ h_i + h_{i - 1} \circ d_i$。
如果 $a = b$，则得到公式
$0 = d_{i + 1} \circ h_i + h_{i - 1} \circ d_i$；换言之，
$ - d_{i + 1} \circ h_i = h_{i - 1} \circ d_i$。由上面的定义，这意味着
$\{h_i\}$ 定义了链复形态射
$$
A_\bullet \longrightarrow B[1]_\bullet.
$$
由上面的说明，这与态射 $A[-1]_\bullet \to B_\bullet$ 是同一回事。
这就证明了下面的引理。

\begin{lemma}
\label{homology-lemma-homotopy-shift}
设 $\mathcal{A}$ 为加性范畴，并假设 $A_\bullet$ 与 $B_\bullet$ 是链复形。
给定任意链复形态射 $a : A_\bullet \to B_\bullet$，从 $a$ 到 $a$ 的
同伦之集合与
$\Mor_{\text{Ch}(\mathcal{A})}(A_\bullet, B[1]_\bullet)$ 之间存在双射。
更一般地，$a$ 与 $b$ 之间的同伦之集合或者为空，或者是群
$\Mor_{\text{Ch}(\mathcal{A})}(A_\bullet, B[1]_\bullet)$ 的主齐性空间。
\end{lemma}

\begin{proof}
见上文。
\end{proof}

\begin{lemma}
\label{homology-lemma-ses-termwise-split}
设 $\mathcal{A}$ 为阿贝尔范畴。设
$$
0 \to A_\bullet \to B_\bullet \to C_\bullet \to 0
$$
为复形的短正合列。假设 $\{s_n : C_n \to B_n\}$ 是一族态射，
它们分裂短正合列 $0 \to A_n \to B_n \to C_n \to 0$。
令 $\pi_n : B_n \to A_n$ 为相应的投影；见引理
\ref{homology-lemma-ses-split}。那么态射族
$$
\pi_{n - 1} \circ d_{B, n} \circ s_n
:
C_n \to A_{n - 1}
$$
定义了复形态射 $\delta(s) : C_\bullet \to A[-1]_\bullet$。
\end{lemma}

\begin{proof}
用 $i : A_\bullet \to B_\bullet$ 与 $q : B_\bullet \to C_\bullet$
表示短正合列中的复形态射。于是
$i_{n - 1} \circ \pi_{n - 1} \circ d_{B, n} \circ s_n =
d_{B, n} \circ s_n - s_{n - 1} \circ d_{C, n}$。因此
$i_{n - 2} \circ d_{A, n - 1} \circ \pi_{n - 1} \circ d_{B, n} \circ s_n =
d_{B, n - 1} \circ (d_{B, n} \circ s_n - s_{n - 1} \circ d_{C, n}) =
- d_{B, n - 1} \circ s_{n - 1} \circ d_{C, n}$，正如所需。
\end{proof}

\begin{lemma}
\label{homology-lemma-ses-termwise-split-long}
记号与假设同上面的引理 \ref{homology-lemma-ses-termwise-split}。复形态射
$\delta(s) : C_\bullet \to A[-1]_\bullet$ 诱导的映射
$$
H_i(\delta(s)) :
H_i(C_\bullet) \longrightarrow H_i(A[-1]_\bullet) = H_{i - 1}(A_\bullet)
$$
就是由引理 \ref{homology-lemma-long-exact-sequence-chain} 关联于该链复形短正合列的
长正合同调列中出现的映射。
\end{lemma}

\begin{proof}
略。
\end{proof}

\begin{lemma}
\label{homology-lemma-ses-termwise-split-homotopy}
记号与假设同上面的引理 \ref{homology-lemma-ses-termwise-split}。假设
$\{s'_n : C_n \to B_n\}$ 是对分裂的另一种选择。对某些唯一的态射
$h_n : C_n \to A_n$，写成 $s'_n = s_n + i_n \circ h_n$。
映射族 $\{h_n : C_n \to A[-1]_{n + 1}\}$ 是相应态射
$\delta(s), \delta(s') : C_\bullet \to A[-1]_\bullet$ 之间的同伦。
\end{lemma}

\begin{proof}
略。
\end{proof}

\begin{definition}
\label{homology-definition-shift-cochain}
设 $\mathcal{A}$ 为加性范畴，并设 $A^\bullet$ 为边界映射为
$d_A^n : A^n \to A^{n + 1}$ 的上链复形。对任意
$k \in \mathbf{Z}$，如下定义{\it 平移 $k$ 的上链复形 $A[k]^\bullet$}：
\begin{enumerate}
\item 令 $A[k]^n = A^{n + k}$；以及
\item 令 $d_{A[k]}^n : A[k]^n \to A[k]^{n + 1}$ 等于
$d_{A[k]}^n = (-1)^k d_A^{n + k}$。
\end{enumerate}
若 $f : A^\bullet \to B^\bullet$ 是上链复形态射，则令
$f[k] : A[k]^\bullet \to B[k]^\bullet$ 为满足
$f[k]^n = f^{k + n}$ 的上链复形态射。
\end{definition}

\noindent
当然，这意味着有函子
$[k] : \text{CoCh}(\mathcal{A}) \to \text{CoCh}(\mathcal{A})$，
它们彼此严格交换（无需插入任何函子同构），并且
$A[k][l]^\bullet = A[k + l]^\bullet$、
$[0] = \text{id}_{\text{CoCh}(\mathcal{A})}$。

\medskip\noindent
回忆我们把 $\mathcal{A}$ 看作 $\text{CoCh}(\mathcal{A})$ 的全子范畴；
见第 \ref{homology-section-complexes} 节。因此，对 $\mathcal{A}$ 的任意对象
$A$，记号 $A[k]$ 表示唯一的上链复形：除次数 $-k$ 处为 $A$ 外，
其他各次均为零。

\begin{definition}
\label{homology-definition-cohomology-shift}
设 $\mathcal{A}$ 为阿贝尔范畴，并设 $A^\bullet$ 为边界映射为
$d_A^n : A^n \to A^{n + 1}$ 的上链复形。对任意
$k \in \mathbf{Z}$，借助等同 $A^{i + k} = A[k]^i$，我们作等同
{\it $H^{i + k}(A^\bullet) \longrightarrow H^i(A[k]^\bullet)$}。
\end{definition}

\noindent
这一等同关于 $A^\bullet$ 具有函子性。注意，因为该定义不涉及符号，
实际上我们得到所有同调对象 $H^{i - k}(A[k]^\bullet)$ 的相容等同系；
它们还与等同 $A[k][l]^\bullet = A[k + l]^\bullet$ 以及
$[0] = \text{id}_{\text{CoCh}(\mathcal{A})}$ 相容。

\medskip\noindent
设 $\mathcal{A}$ 为加性范畴。假设 $A^\bullet$ 与 $B^\bullet$ 是上链复形，
$a, b : A^\bullet \to B^\bullet$ 是上链复形态射，并且
$\{h^i : A^i \to B^{i - 1}\}$ 是 $a$ 与 $b$ 之间的同伦。回忆这意味着
$a^i - b^i = d^{i - 1} \circ h^i + h^{i + 1} \circ d^i$。
如果 $a = b$，则得到公式
$0 = d^{i - 1} \circ h^i + h^{i + 1} \circ d^i$；换言之，
$ - d^{i - 1} \circ h^i = h^{i + 1} \circ d^i$。由上面的定义，这意味着
$\{h^i\}$ 定义了上链复形态射
$$
A^\bullet \longrightarrow B[-1]^\bullet.
$$
由上面的说明，这与态射 $A[1]^\bullet \to B^\bullet$ 是同一回事。
这就证明了下面的引理。

\begin{lemma}
\label{homology-lemma-homotopy-shift-cochain}
设 $\mathcal{A}$ 为加性范畴，并假设 $A^\bullet$ 与 $B^\bullet$ 是上链复形。
给定任意上链复形态射 $a : A^\bullet \to B^\bullet$，从 $a$ 到 $a$ 的
同伦之集合与
$\Mor_{\text{CoCh}(\mathcal{A})}(A^\bullet, B[-1]^\bullet)$ 之间存在双射。
更一般地，$a$ 与 $b$ 之间的同伦之集合或者为空，或者是群
$\Mor_{\text{CoCh}(\mathcal{A})}(A^\bullet, B[-1]^\bullet)$ 的主齐性空间。
\end{lemma}

\begin{proof}
见上文。
\end{proof}

\begin{lemma}
\label{homology-lemma-ses-termwise-split-cochain}
设 $\mathcal{A}$ 为加性范畴。设
$$
0 \to A^\bullet \to B^\bullet \to C^\bullet \to 0
$$
为复形组成的一个复形（！）。假设给定与所示序列中映射相容的分裂
$B^n = A^n \oplus C^n$。令 $s^n : C^n \to B^n$ 与
$\pi^n : B^n \to A^n$ 为相应映射。那么态射族
$$
\pi^{n + 1} \circ d_B^n \circ s^n
:
C^n \to A^{n + 1}
$$
定义了复形态射 $\delta : C^\bullet \to A[1]^\bullet$。
\end{lemma}

\begin{proof}
用 $i : A^\bullet \to B^\bullet$ 与 $q : B^\bullet \to C^\bullet$
表示短正合列中的复形态射。于是
$i^{n + 1} \circ \pi^{n + 1} \circ d_B^n \circ s^n =
d_B^n \circ s^n - s^{n + 1} \circ d_C^n$。因此
$i^{n + 2} \circ d_A^{n + 1} \circ \pi^{n + 1} \circ d_B^n \circ s^n =
d_B^{n + 1} \circ (d_B^n \circ s^n - s^{n + 1} \circ d_C^n) =
- d_B^{n + 1} \circ s^{n + 1} \circ d_C^n$，正如所需。
\end{proof}

\begin{lemma}
\label{homology-lemma-ses-termwise-split-long-cochain}
记号与假设同上面的引理 \ref{homology-lemma-ses-termwise-split-cochain}，并进一步
假设 $\mathcal{A}$ 是阿贝尔范畴。复形态射
$\delta : C^\bullet \to A[1]^\bullet$ 诱导的映射
$$
H^i(\delta) :
H^i(C^\bullet) \longrightarrow H^i(A[1]^\bullet) = H^{i + 1}(A^\bullet)
$$
就是由引理 \ref{homology-lemma-long-exact-sequence-cochain} 关联于该上链复形
短正合列的长正合同调列中出现的映射。
\end{lemma}

\begin{proof}
略。
\end{proof}

\begin{lemma}
\label{homology-lemma-ses-termwise-split-homotopy-cochain}
记号与假设同引理 \ref{homology-lemma-ses-termwise-split-cochain}。令
$\alpha : A^\bullet \to B^\bullet$、
$\beta : B^\bullet \to C^\bullet$ 为给定的复形态射。假设
$(s')^n : C^n \to B^n$ 与 $(\pi')^n : B^n \to A^n$ 是对分裂的
另一种选择。对某些唯一的态射 $h^n : C^n \to A^n$ 与
$g^n : C^n \to A^n$，写成
$(s')^n = s^n + \alpha^n \circ h^n$ 及
$(\pi')^n = \pi^n + g^n \circ \beta^n$。那么
\begin{enumerate}
\item $g^n = - h^n$；并且
\item 映射族 $\{g^n : C^n \to A[1]^{n - 1}\}$ 是
$\delta, \delta' : C^\bullet \to A[1]^\bullet$ 之间的同伦；更精确地，
$(\delta')^n = \delta^n + g^{n + 1} \circ d_C^n + d_{A[1]}^{n - 1} \circ g^n$。
\end{enumerate}
\end{lemma}

\begin{proof}
由于 $(s')^n$ 与 $(\pi')^n$ 是分裂，我们有
$(\pi')^n \circ (s')^n = 0$。因此
$$
0 = ( \pi^n + g^n \circ \beta^n ) \circ ( s^n + \alpha^n \circ h^n ) =
g^n \circ \beta^n \circ s^n + \pi^n \circ \alpha^n \circ h^n =
g^n + h^n
$$
这证明了 (1)。如下计算 $(\delta')^n$：
$$
( \pi^{n + 1} + g^{n + 1} \circ \beta^{n + 1} )
\circ d_B^n \circ
( s^n + \alpha^n \circ h^n )
= \delta^n + g^{n + 1} \circ d_C^n + d_A^n \circ h^n
$$
由于 $h^n = -g^n$ 且 $d_{A[1]}^{n - 1} = -d_A^n$，遂得 (2)。
\end{proof}

\section{复形的截断}
\label{homology-section-truncations}

\noindent
设 $\mathcal{A}$ 为阿贝尔范畴，并设 $A_\bullet$ 为链复形。
复形 $A_\bullet$ 有若干种{\it 截断}方式。
\begin{enumerate}
\item {\it “朴素”截断 $\sigma_{\leq n}$} 是子复形
$\sigma_{\leq n} A_\bullet$，其定义规则为：若 $i > n$，则
$(\sigma_{\leq n} A_\bullet)_i = 0$；若 $i \leq n$，则
$(\sigma_{\leq n} A_\bullet)_i = A_i$。图示为
$$
\xymatrix{
\sigma_{\leq n}A_\bullet \ar[d]  &
\ldots \ar[r] &
0 \ar[r] \ar[d] &
A_n \ar[r] \ar[d] &
A_{n - 1} \ar[r] \ar[d] &
\ldots \\
A_\bullet  &
\ldots \ar[r] &
A_{n + 1} \ar[r] &
A_n \ar[r] &
A_{n - 1} \ar[r] &
\ldots
}
$$
注意性质
$\sigma_{\leq n}A_\bullet / \sigma_{\leq n - 1}A_\bullet = A_n[-n]$。
\item {\it “朴素”截断 $\sigma_{\geq n}$} 是商复形
$\sigma_{\geq n} A_\bullet$，其定义规则为：若 $i \geq n$，则
$(\sigma_{\geq n} A_\bullet)_i = A_i$；若 $i < n$，则
$(\sigma_{\geq n} A_\bullet)_i = 0$。图示为
$$
\xymatrix{
A_\bullet \ar[d]  &
\ldots \ar[r] &
A_{n + 1} \ar[r] \ar[d] &
A_n \ar[r] \ar[d] &
A_{n - 1} \ar[r] \ar[d] &
\ldots \\
\sigma_{\geq n}A_\bullet  &
\ldots \ar[r] &
A_{n + 1} \ar[r] &
A_n \ar[r] &
0 \ar[r] &
\ldots
}
$$
复形态射 $\sigma_{\geq n}A_\bullet \to \sigma_{\geq n + 1}A_\bullet$
是满射态射，其核为 $A_n[-n]$。
\item {\it 典范截断} $\tau_{\geq n}A_\bullet$ 由下图定义：
$$
\xymatrix{
\tau_{\geq n}A_\bullet \ar[d]  &
\ldots \ar[r] &
A_{n + 1} \ar[r] \ar[d] &
\Ker(d_n) \ar[r] \ar[d] &
0 \ar[r] \ar[d] &
\ldots \\
A_\bullet  &
\ldots \ar[r] &
A_{n + 1} \ar[r] &
A_n \ar[r] &
A_{n - 1} \ar[r] &
\ldots
}
$$
注意这些复形具有性质
$$
H_i(\tau_{\geq n}A_\bullet) =
\left\{
\begin{matrix}
H_i(A_\bullet) & \text{若} & i \geq n \\
0 & \text{若} & i < n
\end{matrix}
\right.
$$
\item {\it 典范截断} $\tau_{\leq n}A_\bullet$ 由下图定义：
$$
\xymatrix{
A_\bullet \ar[d]  &
\ldots \ar[r] &
A_{n + 1} \ar[r] \ar[d] &
A_n \ar[r] \ar[d] &
A_{n - 1} \ar[r] \ar[d] &
\ldots \\
\tau_{\leq n}A_\bullet  &
\ldots \ar[r] &
0 \ar[r] &
\Coker(d_{n + 1}) \ar[r] &
A_{n - 1} \ar[r] &
\ldots
}
$$
注意这些复形具有性质
$$
H_i(\tau_{\leq n}A_\bullet) =
\left\{
\begin{matrix}
H_i(A_\bullet) & \text{若} & i \leq n \\
0 & \text{若} & i > n
\end{matrix}
\right.
$$
\end{enumerate}

\noindent
设 $\mathcal{A}$ 为阿贝尔范畴，并设 $A^\bullet$ 为上链复形。
复形 $A^\bullet$ 有四种截断方式。
\begin{enumerate}
\item {\it “朴素”截断 $\sigma_{\geq n}$} 是子复形
$\sigma_{\geq n} A^\bullet$，其定义规则为：若 $i < n$，则
$(\sigma_{\geq n} A^\bullet)^i = 0$；若 $i \geq n$，则
$(\sigma_{\geq n} A^\bullet)^i = A^i$。图示为
$$
\xymatrix{
\sigma_{\geq n}A^\bullet \ar[d]  &
\ldots \ar[r] &
0 \ar[r] \ar[d] &
A^n \ar[r] \ar[d] &
A^{n + 1} \ar[r] \ar[d] &
\ldots \\
A^\bullet  &
\ldots \ar[r] &
A^{n - 1} \ar[r] &
A^n \ar[r] &
A^{n + 1} \ar[r] &
\ldots
}
$$
注意性质
$\sigma_{\geq n}A^\bullet / \sigma_{\geq n + 1}A^\bullet
= A^n[-n]$。
\item {\it “朴素”截断 $\sigma_{\leq n}$} 是商复形
$\sigma_{\leq n} A^\bullet$，其定义规则为：若 $i > n$，则
$(\sigma_{\leq n} A^\bullet)^i = 0$；若 $i \leq n$，则
$(\sigma_{\leq n} A^\bullet)^i = A^i$。图示为
$$
\xymatrix{
A^\bullet \ar[d]  &
\ldots \ar[r] &
A^{n - 1} \ar[r] \ar[d] &
A^n \ar[r] \ar[d] &
A^{n + 1} \ar[r] \ar[d] &
\ldots \\
\sigma_{\leq n}A^\bullet &
\ldots \ar[r] &
A^{n - 1} \ar[r] &
A^n \ar[r] &
0 \ar[r] &
\ldots \\
}
$$
复形态射 $\sigma_{\leq n}A^\bullet \to \sigma_{\leq n - 1}A^\bullet$
是满射态射，其核为 $A^n[-n]$。
\item {\it 典范截断} $\tau_{\leq n}A^\bullet$ 由下图定义：
$$
\xymatrix{
\tau_{\leq n}A^\bullet \ar[d]  &
\ldots \ar[r] &
A^{n - 1} \ar[r] \ar[d] &
\Ker(d^n) \ar[r] \ar[d] &
0 \ar[r] \ar[d] &
\ldots \\
A^\bullet  &
\ldots \ar[r] &
A^{n - 1} \ar[r] &
A^n \ar[r] &
A^{n + 1} \ar[r] &
\ldots
}
$$
注意这些复形具有性质
$$
H^i(\tau_{\leq n}A^\bullet) =
\left\{
\begin{matrix}
H^i(A^\bullet) & \text{若} & i \leq n \\
0 & \text{若} & i > n
\end{matrix}
\right.
$$
\item {\it 典范截断} $\tau_{\geq n}A^\bullet$ 由下图定义：
$$
\xymatrix{
A^\bullet \ar[d] &
\ldots \ar[r] &
A^{n - 1} \ar[r] \ar[d] &
A^n \ar[r] \ar[d] &
A^{n + 1} \ar[r] \ar[d] &
\ldots \\
\tau_{\geq n}A^\bullet &
\ldots \ar[r] &
0 \ar[r] &
\Coker(d^{n - 1}) \ar[r] &
A^{n + 1} \ar[r] &
\ldots
}
$$
注意这些复形具有性质
$$
H^i(\tau_{\geq n}A^\bullet) =
\left\{
\begin{matrix}
0 & \text{若} & i < n \\
H^i(A^\bullet) & \text{若} & i \geq n
\end{matrix}
\right.
$$
\end{enumerate}

\section{分次对象}
\label{homology-section-graded}

\noindent
作如下定义。

\begin{definition}
\label{homology-definition-graded}
设 $\mathcal{A}$ 为加性范畴。$\mathcal{A}$ 的{\it 分次对象范畴}记作
$\text{Gr}(\mathcal{A})$，它是具有下列对象与态射的范畴：
\begin{enumerate}
\item 对象 $A = (A^i)$ 是 $\mathcal{A}$ 的对象 $A^i$（$i \in \mathbf{Z}$）
组成的族；以及
\item 态射 $f : A = (A^i) \to B = (B^i)$ 是 $\mathcal{A}$ 的态射
$f^i : A^i \to B^i$ 组成的族。
\end{enumerate}
\end{definition}

\noindent
若 $\mathcal{A}$ 有可数直和，则可把对象
$$
A = \bigoplus\nolimits_{i \in \mathbf{Z}} A^i
$$
关联于 $\text{Gr}(\mathcal{A})$ 的对象 $A = (A^i)$，并置
$k^iA = A^i$。在这种情况下，$\text{Gr}(\mathcal{A})$ 等价于下述
二元组 $(A, k)$ 的范畴：$A$ 是 $\mathcal{A}$ 的对象，并带有由
$\mathbf{Z}$ 标号的直和项所成的直和分解
$$
A = \bigoplus\nolimits_{i \in \mathbf{Z}} k^iA
$$
此类对象之间的态射 $(A, k) \to (B, k)$，由 $\mathcal{A}$ 中满足
对所有 $i \in \mathbf{Z}$ 均有 $\varphi(k^iA) \subset k^iB$ 的态射
$\varphi : A \to B$ 给出。只要加性范畴 $\mathcal{A}$ 有可数直和，
我们都会使用这一等价而不再说明。

\medskip\noindent
然而，按我们的定义，加性范畴或阿贝尔范畴不一定具有所有（可数）直和。
在这种情况下，上述定义仍然有意义。例如，若
$\mathcal{A} = \text{Vect}_k$ 是域 $k$ 上有限维向量空间的范畴，则
$\text{Gr}(\text{Vect}_k)$ 是带有给定分次且每个分次部分均有限维的
向量空间之范畴，而不是带有给定分次的有限维向量空间之范畴。

\begin{lemma}
\label{homology-lemma-graded}
设 $\mathcal{A}$ 为阿贝尔范畴。分次对象范畴
$\text{Gr}(\mathcal{A})$ 是阿贝尔范畴。
\end{lemma}

\begin{proof}
设 $f : A = (A^i) \to B = (B^i)$ 是 $\mathcal{A}$ 的分次对象之间
由 $\mathcal{A}$ 的态射 $f^i : A^i \to B^i$ 给出的态射。那么在范畴
$\text{Gr}(\mathcal{A})$ 中有
$\Ker(f) = (\Ker(f^i))$ 及 $\Coker(f) = (\Coker(f^i))$。
由于 $\mathcal{A}$ 中 $\Im = \Coim$，在
$\text{Gr}(\mathcal{A})$ 中也同样成立。
\end{proof}

\begin{remark}[警告]
\label{homology-remark-direct-sums-not-exact}
存在具有可数直和，但可数直和不正合的阿贝尔范畴 $\mathcal{A}$。
一个例子是 $\mathbf{R}$ 上阿贝尔群层的范畴之反范畴。事实上，
$\mathbf{R}$ 上阿贝尔群层的范畴具有可数积，但可数积不正合。
对这样的范畴，上述函子
$\text{Gr}(\mathcal{A}) \to \mathcal{A}$、
$(A^i) \mapsto \bigoplus A^i$ 不正合。仍然成立的是，
$\text{Gr}(\mathcal{A})$ 等价于 $\mathcal{A}$ 的分次对象
$(A, k)$ 所成的范畴；但在分次对象范畴中，态射
$\varphi : (A, k) \to (B, k)$ 的核并不等于配备直和分解的
$\Ker(\varphi)$，而是各映射 $k^iA \to k^iB$ 的核之直和。
\end{remark}

\begin{definition}
\label{homology-definition-graded-shift}
设 $\mathcal{A}$ 为加性范畴。若 $A = (A^i)$ 是分次对象，则它的第
$k$ 次{\it 平移} $A[k]$ 是满足 $A[k]^i = A^{k + i}$ 的分次对象。
\end{definition}

\noindent
若 $A$ 与 $B$ 是 $\mathcal{A}$ 的分次对象，则有
\begin{equation}
\label{homology-equation-hom-into-shift}
\Hom_{\text{Gr}(\mathcal{A})}(A, B[k]) =
\Hom_{\text{Gr}(\mathcal{A})}(A[-k], B)
\end{equation}
而这个群的元素有时称为分次对象之间的一个{\it $k$ 次齐次映射}。

\medskip\noindent
给定任意集合 $G$，可把 $\mathcal{A}$ 的 $G$-分次对象定义为如下
范畴：其对象是由 $G$ 的元素参数化的对象族
$A = (A^g)_{g \in G}$；态射 $f : A \to B$ 定义为映射
$f^g : A^g \to B^g$ 的族，其中 $g$ 遍历 $G$ 的元素。若 $G$ 是阿贝尔群，
则可无歧义地在 $G$-分次对象范畴上按规则
$(A[g])^{g_0} = A^{g + g_0}$ 定义平移函子 $[g]$。
这种构造的一个特例是 $G = \mathbf{Z} \times \mathbf{Z}$。
此时范畴的对象称为 $\mathcal{A}$ 的{\it 双分次}对象。双分次对象
$A$ 的 $(p, q)$ 分量通常记作 $A^{p, q}$。对
$(a, b) \in \mathbf{Z} \times \mathbf{Z}$，我们写 $A[a, b]$ 而不是
$A[(a, b)]$。态射 $A \to A[a, b]$ 有时称为一个
{\it 双次数为 $(a, b)$ 的映射}。

\section{加性幺半范畴}
\label{homology-section-monoidal}

\noindent
本节给出范畴上的幺半结构与加性结构相互作用的一些材料。

\begin{definition}
\label{homology-definition-additive-monoidal}
一个{\it 加性幺半范畴}，是配备幺半结构 $\otimes, \phi$
（范畴论，定义 \ref{categories-definition-monoidal-category}）的加性范畴
$\mathcal{A}$，并要求 $\otimes$ 关于每个变量都是加性函子。
\end{definition}

\begin{lemma}
\label{homology-lemma-additive-dual}
设 $\mathcal{A}$ 为加性幺半范畴。若 $Y_i$（$i = 1, 2$）是
$X_i$（$i = 1, 2$）的左对偶，则 $Y_1 \oplus Y_2$ 是
$X_1 \oplus X_2$ 的左对偶。
\end{lemma}

\begin{proof}
由伴随的唯一性及范畴论，注记
\ref{categories-remark-left-dual-adjoint} 可得。
\end{proof}

\begin{lemma}
\label{homology-lemma-Karoubian-dual}
在 Karoubian 加性幺半范畴中，具有左对偶的对象之每个直和项都有左对偶。
\end{lemma}

\begin{proof}
我们将使用范畴论，引理 \ref{categories-lemma-left-dual} 而不再说明。
设 $X$ 是具有左对偶 $Y$ 的对象。我们有
$$
\Hom(X, X) = \Hom(\mathbf{1}, X \otimes Y) = \Hom(Y, Y)
$$
若 $a : X \to X$ 对应于 $b : Y \to Y$，则 $b$ 是 $Y$ 的唯一自同态，
使得在
$$
\Hom(Z' \otimes X, Z) = \Hom(Z', Z \otimes Y)
$$
上预复合 $a$，等同于后复合 $1 \otimes b$。因此双射
$\Hom(X, X) \to \Hom(Y, Y)$、$a \mapsto b$ 是代数
$\Hom(X, X)$ 的反代数与代数 $\Hom(Y, Y)$ 之间的同构。
特别地，若 $X = X_1 \oplus X_2$，则相应投影元 $e_1, e_2$ 被映到
$\Hom(Y, Y)$ 中的幂等元。若 $Y = Y_1 \oplus Y_2$ 是 $Y$ 的相应
直和分解（第 \ref{homology-section-karoubian} 节），则在双射
$\Hom(Z' \otimes X, Z) = \Hom(Z', Z \otimes Y)$ 下，对
$i = 1, 2$，我们函子性地得到子群之间的等式
$\Hom(Z' \otimes X_i, Z) = \Hom(Z', Z \otimes Y_i)$。
由范畴论，注记 \ref{categories-remark-left-dual-adjoint} 中的讨论可知，
$Y_i$ 是 $X_i$ 的左对偶。
\end{proof}

\begin{example}
\label{homology-example-graded-vector-spaces}
设 $F$ 为域，并设 $\mathcal{C}$ 为分次 $F$-向量空间的范畴。
给定分次向量空间 $V$ 与 $W$，令 $V \otimes W$ 表示其次数 $n$ 部分为
$$
(V \otimes W)^n = \bigoplus\nolimits_{n = p + q} V^p \otimes_F W^q
$$
的分次 $F$-向量空间。给定第三个分次向量空间 $U$，作为结合约束
$\phi : U \otimes (V \otimes W) \to (U \otimes V) \otimes W$，
我们使用向量空间的“通常”同构
$$
U^p \otimes_F (V^q \otimes_F W^r) \to (U^p \otimes_F V^q) \otimes_F W^r
$$
作为单位对象，使用在次数 $0$ 处为 $F$、其他次数处为零的分次
$F$-向量空间 $\mathbf{1}$。$\mathcal{C}$ 上有两个交换约束，使
$\mathcal{C}$ 成为对称幺半范畴：一个涉及符号，另一个不涉及。我们通常
使用涉及符号的那个。
明确地说，若 $V$ 与 $W$ 是分次 $F$-向量空间，我们使用同构
$\psi : V \otimes W \to W \otimes V$，它在次数 $n$ 处使用
$$
V^p \otimes_F W^q \to W^q \otimes_F V^p,\quad
v \otimes w \mapsto (-1)^{pq} w \otimes v
$$
我们略去其有效性的验证。
\end{example}

\begin{lemma}
\label{homology-lemma-left-dual-graded-vector-spaces}
设 $F$ 为域，并按例 \ref{homology-example-graded-vector-spaces} 把分次
$F$-向量空间范畴 $\mathcal{C}$ 看作幺半范畴。若 $\mathcal{C}$ 中的
$V$ 具有左对偶 $W$，则 $\sum_n \dim_F V^n < \infty$，且映射
$\epsilon$ 定义非退化配对 $W^{-n} \times V^n \to F$。
\end{lemma}

\begin{proof}
作为单位对象，我们取。% P02-E0009: source line 4074 is the incomplete fragment “As unit we take”.
由范畴论，定义 \ref{categories-definition-dual}，我们有映射
$$
\eta : \mathbf{1} \to V \otimes W\quad
\epsilon : W \otimes V \to \mathbf{1}
$$
由于 $\mathbf{1} = F$ 置于次数 $0$，可以把 $\epsilon$ 看作引理陈述中的
一列配对 $W^{-n} \times V^n \to F$。对所有 $n$，为 $V^n$ 选取基
$\{e_{n, i}\}_{i \in I_n}$。写成
$$
\eta(1) = \sum e_{n, i} \otimes w_{-n, i}
$$
其中 $w_{-n, i} \in W^{-n}$，且除有限多个外均为零！条件
$(\epsilon \otimes 1) \circ (1 \otimes \eta)$ 是 $W$ 上的恒等映射，
意味着
$$
\sum\nolimits_{n, i} \epsilon(w, e_{n, i})w_{-n, i} = w
$$
因此，作为分次向量空间，$W$ 由有限多个非零向量 $w_{-n, i}$ 生成。
条件 $(1 \otimes \epsilon) \circ (\eta \otimes 1)$ 是 $V$ 上的恒等映射，
意味着
$$
\sum\nolimits_{n, i} e_{n, i}\ \epsilon(w_{-n, i}, v) = v
$$
特别地，置 $v = e_{n, i}$，可得
$\epsilon(w_{-n, i}, e_{n, i'}) = \delta_{ii'}$。因此引理的陈述成立，
并且 $\{w_{-n, i}\}_{i \in I_n}$ 是 $W^{-n}$ 中与所选 $V^n$ 的基
对偶的基。
\end{proof}

\section{双复形及其相伴总复形}
\label{homology-section-double-complexes}

\noindent
我们讨论双复形及其相伴总复形。

\begin{definition}
\label{homology-definition-double-complex}
设 $\mathcal{A}$ 为加性范畴。$\mathcal{A}$ 中的一个{\it 双复形}，
由系统
$(\{A^{p, q}, d_1^{p, q}, d_2^{p, q}\}_{p, q\in \mathbf{Z}})$ 给出；
其中每个 $A^{p, q}$ 都是 $\mathcal{A}$ 的对象，而
$d_1^{p, q} : A^{p, q} \to A^{p + 1, q}$ 及
$d_2^{p, q} : A^{p, q} \to A^{p, q + 1}$ 是 $\mathcal{A}$ 的态射，
并且对所有 $p, q \in \mathbf{Z}$ 均满足下列规则：
\begin{enumerate}
\item $d_1^{p + 1, q} \circ d_1^{p, q} = 0$
\item $d_2^{p, q + 1} \circ d_2^{p, q} = 0$
\item $d_1^{p, q + 1} \circ d_2^{p, q} = d_2^{p + 1, q} \circ d_1^{p, q}$
\end{enumerate}
\end{definition}

\noindent
这正是该定义的上链版本。它说明每个 $A^{p, \bullet}$ 都是上链复形，
每个 $d_1^{p, \bullet}$ 都是复形态射
$A^{p, \bullet} \to A^{p + 1, \bullet}$，并且作为复形态射有
$d_1^{p + 1, \bullet} \circ d_1^{p, \bullet} = 0$。
换言之，可把双复形看作复形组成的复形。因此，在图表
$$
\xymatrix{
\ldots &
\ldots &
\ldots &
\ldots \\
\ldots \ar[r] &
A^{p, q + 1} \ar[r]^{d_1^{p, q + 1}} \ar[u] &
A^{p + 1, q + 1} \ar[r] \ar[u] &
\ldots \\
\ldots \ar[r] &
A^{p, q} \ar[r]^{d_1^{p, q}} \ar[u]^{d_2^{p, q}} &
A^{p + 1, q} \ar[r] \ar[u]_{d_2^{p + 1, q}} &
\ldots \\
\ldots &
\ldots \ar[u] &
\ldots \ar[u] &
\ldots
}
$$
中，每个方块都交换。警告：文献中还可见另一种定义，其中“双复形”
（bicomplex 或 double complex）的上述方块反交换。

\begin{example}
\label{homology-example-double-complex-as-tensor-product-of}
设 $\mathcal{A}$、$\mathcal{B}$、$\mathcal{C}$ 为加性范畴。假设
$$
\otimes : \mathcal{A} \times \mathcal{B} \longrightarrow \mathcal{C},
\quad
(X, Y) \longmapsto X \otimes Y
$$
是一个在态射上双线性的函子；$\mathcal{A} \times \mathcal{B}$ 的定义
见范畴论，定义 \ref{categories-definition-product-category}。
给定 $\mathcal{A}$ 的复形 $X^\bullet$ 及 $\mathcal{B}$ 的复形
$Y^\bullet$，可在 $\mathcal{C}$ 中得到双复形
$$
K^{\bullet, \bullet} = X^\bullet \otimes Y^\bullet
$$
这里第一微分 $K^{p, q} \to K^{p + 1, q}$ 是由态射
$X^p \to X^{p + 1}$ 以及 $Y^q$ 上的恒等映射诱导的态射
$X^p \otimes Y^q \to X^{p + 1} \otimes Y^q$。第二微分类似。
\end{example}

\begin{definition}
\label{homology-definition-associated-simple-complex}
设 $\mathcal{A}$ 为加性范畴，并设 $A^{\bullet, \bullet}$ 为双复形。
它的{\it 相伴单复形}记作 $sA^\bullet$，也常称为{\it 相伴总复形}，
记作 $\text{Tot}(A^{\bullet, \bullet})$；其定义为
$$
sA^n = \text{Tot}^n(A^{\bullet, \bullet}) =
\bigoplus\nolimits_{n = p + q} A^{p, q}
$$
（如果存在），并带有微分
$$
d_{sA^\bullet}^n = d_{\text{Tot}(A^{\bullet, \bullet})}^n =
\sum\nolimits_{n = p + q} (d_1^{p, q} + (-1)^p d_2^{p, q})
$$
\end{definition}

\noindent
若 $\mathcal{A}$ 中存在可数直和，或对每个 $n$，至多有限多个
$A^{p, n - p}$ 非零，则 $\text{Tot}(A^{\bullet, \bullet})$ 存在。
注意该定义关于指标 $(p, q)$ {\it 不}对称。

\begin{remark}
\label{homology-remark-triple-complex}
设 $\mathcal{A}$ 为加性范畴，并设 $A^{\bullet, \bullet, \bullet}$
为三重复形。相伴总复形的项为
$$
\text{Tot}^n(A^{\bullet, \bullet, \bullet}) =
\bigoplus\nolimits_{p + q + r = n} A^{p, q, r}
$$
而其微分为
$$
d^n_{\text{Tot}(A^{\bullet, \bullet, \bullet})} =
\sum\nolimits_{p + q + r = n}
d_1^{p, q, r} + (-1)^pd_2^{p, q, r} + (-1)^{p + q}d_3^{p, q, r}
$$
按此定义，一个简单计算表明相伴总复形等于
$$
\text{Tot}(A^{\bullet, \bullet, \bullet}) =
\text{Tot}(\text{Tot}_{12}(A^{\bullet, \bullet, \bullet})) =
\text{Tot}(\text{Tot}_{23}(A^{\bullet, \bullet, \bullet}))
$$
换言之，可以先合并前两个变量，再把它们的和与最后一个变量合并；
也可以先合并后两个变量，再把第一个变量与后两者之和合并。
\end{remark}

\begin{remark}
\label{homology-remark-shift-double-complex}
设 $\mathcal{A}$ 为加性范畴，并设 $A^{\bullet, \bullet}$ 是微分为
$d_1^{p, q}$ 与 $d_2^{p, q}$ 的双复形。用
$A^{\bullet, \bullet}[a, b]$ 表示满足
$$
(A^{\bullet, \bullet}[a, b])^{p, q} = A^{p + a, q + b}
$$
且微分为
$$
d_{A^{\bullet, \bullet}[a, b], 1}^{p, q} = (-1)^a d_1^{p + a, q + b}
\quad\text{且}\quad
d_{A^{\bullet, \bullet}[a, b], 2}^{p, q} = (-1)^b d_2^{p + a, q + b}
$$
的双复形。在这种情况下，存在良定义的同构
$$
\gamma :
\text{Tot}(A^{\bullet, \bullet})[a + b]
\longrightarrow
\text{Tot}(A^{\bullet, \bullet}[a, b])
$$
它在次数 $n$ 处由映射
$$
\xymatrix{
(\text{Tot}(A^{\bullet, \bullet})[a + b])^n =
\bigoplus_{p + q = n + a + b} A^{p, q}
\ar[d]^{\epsilon(p, q, a, b)\text{id}_{A^{p, q}}} \\
\text{Tot}(A^{\bullet, \bullet}[a, b])^n =
\bigoplus_{p' + q' = n} A^{p' + a, q' + b}
}
$$
给出，其中 $\epsilon(p, q, a, b)$ 是某个符号。当然，当
$p = p' + a$ 且 $q = q' + b$ 时，直和项 $A^{p, q}$ 映到直和项
$A^{p' + a, q' + b}$。为求出这些符号应满足的条件，注意在源上有
$$
d|_{A^{p, q}} = (-1)^{a + b}\left(d_1^{p, q} + (-1)^pd_2^{p, q}\right)
$$
而在目标上有
$$
d|_{A^{p' + a, q' + b}} =
(-1)^ad_1^{p' + a, q' + b} + (-1)^{p'}(-1)^bd_2^{p' + a, q' + b}
$$
因此约束为
$$
(-1)^a \epsilon(p, q, a, b) = \epsilon(p + 1, q, a, b)(-1)^{a + b}
\Leftrightarrow
\epsilon(p + 1, q, a, b) = (-1)^b \epsilon(p, q, a, b)
$$
以及
$$
(-1)^{p' + b}\epsilon(p, q, a, b) =
\epsilon(p, q + 1, a, b) (-1)^{a + b + p}
\Leftrightarrow
\epsilon(p, q, a, b) = \epsilon(p, q + 1, a, b)
$$
因此选择 $\epsilon(p, q, a, b) = (-1)^{pb}$。
\end{remark}

\begin{remark}
\label{homology-remark-double-complex-complex-of-complexes-first}
设 $\mathcal{A}$ 为具有可数直和的加性范畴。用
$\text{DoubleComp}(\mathcal{A})$ 表示双复形范畴。可把
$\text{DoubleComp}(\mathcal{A})$ 的对象 $A^{\bullet, \bullet}$ 如下看作
复形组成的复形：
$$
\ldots \to A^{\bullet, -1} \to A^{\bullet, 0} \to A^{\bullet, 1} \to \ldots
$$
交换两个指标之角色的变体见注记
\ref{homology-remark-double-complex-complex-of-complexes-second}。本注记说明，
取相伴总复形与本章此前研究过的复形上的所有结构相容。

\medskip\noindent
首先，注意按上述方式把双复形看作复形组成的复形时，其上的平移函子就是
注记 \ref{homology-remark-shift-double-complex} 中定义的函子 $[0, 1]$。
由注记 \ref{homology-remark-shift-double-complex}，函子
$$
\text{Tot} : \text{DoubleComp}(\mathcal{A}) \to \text{Comp}(\mathcal{A})
$$
与平移函子相容；也就是说，存在函子性同构
$\gamma : \text{Tot}(A^{\bullet, \bullet})[1] \to
\text{Tot}(A^{\bullet, \bullet}[0, 1])$。

\medskip\noindent
其次，若
$$
f, g : A^{\bullet, \bullet} \to B^{\bullet, \bullet}
$$
在按上述方式把 $f$ 与 $g$ 看作复形组成的复形之间的态射时同伦，则
$$
\text{Tot}(f), \text{Tot}(g) :
\text{Tot}(A^{\bullet, \bullet}) \to \text{Tot}(B^{\bullet, \bullet})
$$
是同伦的复形映射。事实上，设 $h = (h^q)$ 为 $f$ 与 $g$ 之间的同伦。
若用 $h^{p, q} : A^{p, q} \to B^{p, q - 1}$ 表示 $h^q$ 的次数 $p$
分量，则这意味着
$$
f^{p, q} - g^{p, q} = d_2^{p, q - 1} \circ h^{p, q} +
h^{p, q + 1} \circ d_2^{p, q}
$$
而 $h^q : A^{\bullet, q} \to B^{\bullet, q - 1}$ 是复形映射这一事实意味着
$$
d_1^{p, q - 1} \circ h^{p, q} = h^{p + 1, q} \circ d_1^{p, q}
$$
定义同伦 $h' = ((h')^n)$ 如下：映射
$(h')^n : \text{Tot}^n(A^{\bullet, \bullet}) \to
\text{Tot}^{n - 1}(B^{\bullet, \bullet})$ 在 $p + q = n$ 的直和项
$A^{p, q}$ 上使用 $(-1)^ph^{p, q}$。于是可见
$$
d_{\text{Tot}(B^{\bullet, \bullet})} \circ h' +
h' \circ d_{\text{Tot}(A^{\bullet, \bullet})}
$$
限制在直和项 $A^{p, q}$ 上等于
$$
d_1^{p, q - 1} \circ (-1)^p h^{p, q} +
(-1)^p d_2^{p, q - 1} \circ (-1)^p h^{p, q} +
(-1)^{p + 1} h^{p + 1, q} \circ d_1^{p, q} +
(-1)^p h^{p, q + 1} \circ (-1)^p d_2^{p, q}
$$
由上面的等式，其值为 $f^{p, q} - g^{p, q}$。这证明了第二个相容性。

\medskip\noindent
第三，假设上一段中 $f = g$。那么上述对应 $h \leadsto h'$ 与引理
\ref{homology-lemma-homotopy-shift-cochain} 的等同相容。更精确地，若通过该引理
把 $h$ 看作复形组成的复形之间的态射
$A^{\bullet, \bullet} \to B^{\bullet, \bullet}[0, -1]$，则
$$
\text{Tot}(A^{\bullet, \bullet})
\xrightarrow{\text{Tot}(h)}
\text{Tot}(B^{\bullet, \bullet}[0, -1])
\xrightarrow{\gamma^{-1}}
\text{Tot}(B^{\bullet, \bullet})[-1]
$$
等于把 $h'$ 通过该引理看作的复形态射。这里 $\gamma$ 是注记
\ref{homology-remark-shift-double-complex} 的等同。第三点的验证是立即的。

\medskip\noindent
第四，设
$$
0 \to A^{\bullet, \bullet} \to B^{\bullet, \bullet} \to
C^{\bullet, \bullet} \to 0
$$
是双复形组成的一个复形；并假设按上述方式把双复形看作复形组成的复形时，
给定了如引理 \ref{homology-lemma-ses-termwise-split-cochain} 中的分裂
$s^q : C^{\bullet, q} \to B^{\bullet, q}$ 与
$\pi^q : B^{\bullet, q} \to A^{\bullet, q}$。一方面，引理
\ref{homology-lemma-ses-termwise-split-cochain} 的过程产生映射
$$
\delta : C^{\bullet, \bullet} \longrightarrow A^{\bullet, \bullet}[0, 1]
$$
另一方面，取 $\text{Tot}$ 得到逐项分裂（见下文）的复形
$$
0 \to \text{Tot}(A^{\bullet, \bullet}) \to
\text{Tot}(B^{\bullet, \bullet}) \to
\text{Tot}(C^{\bullet, \bullet}) \to 0
$$
因而由引理 \ref{homology-lemma-ses-termwise-split-cochain} 与
\ref{homology-lemma-ses-termwise-split-homotopy-cochain}，它带有一个在同伦意义下
良定义的态射
$$
\delta' :
\text{Tot}(C^{\bullet, \bullet})
\longrightarrow
\text{Tot}(A^{\bullet, \bullet})[1]
$$
断言：这些映射在下述意义下相同，即
$$
\text{Tot}(C^{\bullet, \bullet})
\xrightarrow{\text{Tot}(\delta)}
\text{Tot}(A^{\bullet, \bullet}[0, 1]) \xrightarrow{\gamma^{-1}}
\text{Tot}(A^{\bullet, \bullet})[1]
$$
等于 $\delta'$，其中 $\gamma$ 如注记 \ref{homology-remark-shift-double-complex}
所述。为说明这一点，用 $s^{p, q} : C^{p, q} \to B^{\bullet, q}$ 与
$\pi^{p, q} : B^{p, q} \to A^{p, q}$ 表示 $s^q$ 与 $\pi^q$ 的分量。
% P02-E0010: source line 4435 types the component codomain as B^{\bullet,q}, not B^{p,q}.
作为分裂
$(s')^n : \text{Tot}^n(C^{\bullet, \bullet}) \to
\text{Tot}^n(B^{\bullet, \bullet})$ 与
$(\pi')^n : \text{Tot}^n(B^{\bullet, \bullet}) \to
\text{Tot}^n(A^{\bullet, \bullet})$，我们使用在 $p + q = n$ 时分量为
$s^{p, q}$ 与 $\pi^{p, q}$ 的映射。回忆
$$
(\delta')^n =
(\pi')^{n + 1} \circ d_{\text{Tot}(B^{\bullet, \bullet})}^n \circ (s')^n :
\text{Tot}^n(C^{\bullet, \bullet}) \to
\text{Tot}^{n + 1}(A^{\bullet, \bullet})
$$
它在直和项 $C^{p, q}$ 上的限制等于
$$
\pi^{p + 1, q} \circ
d_1^{p, q} \circ
s^{p, q} +
\pi^{p, q + 1} \circ
(-1)^p d_2^{p, q} \circ
s^{p, q} =
\pi^{p, q + 1} \circ
(-1)^p d_2^{p, q} \circ
s^{p, q}
$$
等式成立，是因为 $s^q$ 是以 $d_1$ 为微分的复形态射，并且在每个双次数中
$s$ 与 $\pi$ 对应于 $B$ 的直和分解，故
$\pi^{p + 1, q} \circ s^{p + 1, q} = 0$。另一方面，对 $\delta$ 有
$$
\delta^q =  \pi^q \circ d_2 \circ s^q :
C^{\bullet, q} \to A^{\bullet, q + 1}
$$
其在直和项 $C^{p, q}$ 上的限制等于
$\pi^{p, q + 1} \circ d_2^{p, q} \circ s^{p, q}$。符号之差恰由同构
$\gamma$ 中的符号 $(-1)^p$ 抵消，第四个断言得证。
\end{remark}

\begin{remark}
\label{homology-remark-double-complex-complex-of-complexes-second}
设 $\mathcal{A}$ 为具有可数直和的加性范畴。用
$\text{DoubleComp}(\mathcal{A})$ 表示双复形范畴。可把
$\text{DoubleComp}(\mathcal{A})$ 的对象 $A^{\bullet, \bullet}$ 如下看作
复形组成的复形：
$$
\ldots \to A^{-1, \bullet} \to A^{0, \bullet} \to A^{1, \bullet} \to \ldots
$$
交换两个指标之角色的变体见注记
\ref{homology-remark-double-complex-complex-of-complexes-first}。本注记说明，
取相伴总复形与本章此前研究过的复形上的所有结构相容。

\medskip\noindent
首先，注意按上述方式把双复形看作复形组成的复形时，其上的平移函子就是
注记 \ref{homology-remark-shift-double-complex} 中定义的函子 $[1, 0]$。
由注记 \ref{homology-remark-shift-double-complex}，函子
$$
\text{Tot} : \text{DoubleComp}(\mathcal{A}) \to \text{Comp}(\mathcal{A})
$$
与平移函子相容；也就是说，存在函子性同构
$\gamma : \text{Tot}(A^{\bullet, \bullet})[1] \to
\text{Tot}(A^{\bullet, \bullet}[1, 0])$。

\medskip\noindent
其次，若
$$
f, g : A^{\bullet, \bullet} \to B^{\bullet, \bullet}
$$
在按上述方式把 $f$ 与 $g$ 看作复形组成的复形之间的态射时同伦，则
$$
\text{Tot}(f), \text{Tot}(g) :
\text{Tot}(A^{\bullet, \bullet}) \to \text{Tot}(B^{\bullet, \bullet})
$$
是同伦的复形映射。事实上，设 $h = (h^p)$ 为 $f$ 与 $g$ 之间的同伦。
若用 $h^{p, q} : A^{p, q} \to B^{p - 1, q}$ 表示 $h^q$ 的次数 $p$
分量，则这意味着
% P02-E0011: source line 4518 says “degree p of h^q” although the family was h=(h^p).
$$
f^{p, q} - g^{p, q} = d_1^{p - 1, q} \circ h^{p, q} +
h^{p + 1, q} \circ d_1^{p, q}
$$
而 $h^p : A^{p, \bullet} \to B^{p - 1, \bullet}$ 是复形映射这一事实意味着
$$
d_2^{p - 1, q} \circ h^{p, q} = h^{p, q + 1} \circ d_2^{p, q}
$$
定义同伦 $h' = ((h')^n)$ 如下：映射
$(h')^n : \text{Tot}^n(A^{\bullet, \bullet}) \to
\text{Tot}^{n - 1}(B^{\bullet, \bullet})$ 在 $p + q = n$ 的直和项
$A^{p, q}$ 上使用 $h^{p, q}$。于是可见
$$
d_{\text{Tot}(B^{\bullet, \bullet})} \circ h' +
h' \circ d_{\text{Tot}(A^{\bullet, \bullet})}
$$
限制在直和项 $A^{p, q}$ 上等于
$$
d_1^{p - 1, q} \circ h^{p, q} +
(-1)^{p - 1} d_2^{p - 1, q} \circ h^{p, q} +
h^{p + 1, q} \circ d_1^{p, q} +
h^{p, q + 1} \circ (-1)^p d_2^{p, q}
$$
由上面的等式，其值为 $f^{p, q} - g^{p, q}$。这证明了第二个相容性。

\medskip\noindent
第三，假设上一段中 $f = g$。那么上述对应 $h \leadsto h'$ 与引理
\ref{homology-lemma-homotopy-shift-cochain} 的等同相容。更精确地，若通过该引理
把 $h$ 看作复形组成的复形之间的态射
$A^{\bullet, \bullet} \to B^{\bullet, \bullet}[-1, 0]$，则
$$
\text{Tot}(A^{\bullet, \bullet})
\xrightarrow{\text{Tot}(h)}
\text{Tot}(B^{\bullet, \bullet}[-1, 0])
\xrightarrow{\gamma^{-1}}
\text{Tot}(B^{\bullet, \bullet})[-1]
$$
等于把 $h'$ 通过该引理看作的复形态射。这里 $\gamma$ 是注记
\ref{homology-remark-shift-double-complex} 的等同。第三点的验证是立即的。

\medskip\noindent
第四，设
$$
0 \to A^{\bullet, \bullet} \to B^{\bullet, \bullet} \to
C^{\bullet, \bullet} \to 0
$$
是双复形组成的一个复形；并假设按上述方式把双复形看作复形组成的复形时，
给定了如引理 \ref{homology-lemma-ses-termwise-split-cochain} 中的分裂
$s^p : C^{p, \bullet} \to B^{p, \bullet}$ 与
$\pi^p : B^{p, \bullet} \to A^{p, \bullet}$。一方面，引理
\ref{homology-lemma-ses-termwise-split-cochain} 的过程产生映射
$$
\delta : C^{\bullet, \bullet} \longrightarrow A^{\bullet, \bullet}[0, 1]
$$
% P02-E0012: source line 4581 uses [0,1] although this second-index-switched construction and lines 4603/4648 require [1,0].
另一方面，取 $\text{Tot}$ 得到逐项分裂（见下文）的复形
$$
0 \to \text{Tot}(A^{\bullet, \bullet}) \to
\text{Tot}(B^{\bullet, \bullet}) \to
\text{Tot}(C^{\bullet, \bullet}) \to 0
$$
因而由引理 \ref{homology-lemma-ses-termwise-split-cochain} 与
\ref{homology-lemma-ses-termwise-split-homotopy-cochain}，它带有一个在同伦意义下
良定义的态射
$$
\delta' :
\text{Tot}(C^{\bullet, \bullet})
\longrightarrow
\text{Tot}(A^{\bullet, \bullet})[1]
$$
断言：这些映射在下述意义下相同，即
$$
\text{Tot}(C^{\bullet, \bullet})
\xrightarrow{\text{Tot}(\delta)}
\text{Tot}(A^{\bullet, \bullet}[1, 0]) \xrightarrow{\gamma^{-1}}
\text{Tot}(A^{\bullet, \bullet})[1]
$$
等于 $\delta'$，其中 $\gamma$ 如注记 \ref{homology-remark-shift-double-complex}
所述。为说明这一点，用 $s^{p, q} : C^{p, q} \to B^{\bullet, q}$ 与
$\pi^{p, q} : B^{p, q} \to A^{p, q}$ 表示 $s^q$ 与 $\pi^q$ 的分量。
% P02-E0013: source lines 4608-4609 retain s^q/pi^q and B^{\bullet,q} despite the second construction using s^p/pi^p and bidegree components.
作为分裂
$(s')^n : \text{Tot}^n(C^{\bullet, \bullet}) \to
\text{Tot}^n(B^{\bullet, \bullet})$ 与
$(\pi')^n : \text{Tot}^n(B^{\bullet, \bullet}) \to
\text{Tot}^n(A^{\bullet, \bullet})$，我们使用在 $p + q = n$ 时分量为
$s^{p, q}$ 与 $\pi^{p, q}$ 的映射。回忆
$$
(\delta')^n =
(\pi')^{n + 1} \circ d_{\text{Tot}(B^{\bullet, \bullet})}^n \circ (s')^n :
\text{Tot}^n(C^{\bullet, \bullet}) \to
\text{Tot}^{n + 1}(A^{\bullet, \bullet})
$$
它在直和项 $C^{p, q}$ 上的限制等于
$$
\pi^{p + 1, q} \circ
d_1^{p, q} \circ
s^{p, q} +
\pi^{p, q + 1} \circ
(-1)^p d_2^{p, q} \circ
s^{p, q} =
\pi^{p + 1, q} \circ
d_1^{p, q} \circ
s^{p, q}
$$
等式成立，是因为 $s^p$ 是以 $d_2$ 为微分的复形态射，并且在每个双次数中
$s$ 与 $\pi$ 对应于 $B$ 的直和分解，故
$\pi^{p, q + 1} \circ s^{p, q + 1} = 0$。另一方面，对 $\delta$ 有
$$
\delta^p =  \pi^p \circ d_1 \circ s^p :
C^{p, \bullet} \to A^{p + 1, \bullet}
$$
其在直和项 $C^{p, q}$ 上的限制等于
$\pi^{p + 1, q} \circ d_1^{p, q} \circ s^{p, q}$。
因此所得与前面相同，这也符合下述事实：同构
$\gamma : \text{Tot}(A^{\bullet, \bullet})[1] \to
\text{Tot}(A^{\bullet, \bullet}[1, 0])$ 的定义不涉及符号。
\end{remark}

\section{滤过}
\label{homology-section-filtrations}

\noindent
这部分材料可参见很好的参考文献 \cite[Section 1]{HodgeII}。
（注意，我们关于阿贝尔范畴的约定不同。）

\begin{definition}
\label{homology-definition-filtered}
设 $\mathcal{A}$ 为阿贝尔范畴。
\begin{enumerate}
\item 对象 $A$ 上的一个{\it 递减滤过} $F$，是一族 $A$ 的子对象
$(F^nA)_{n \in \mathbf{Z}}$，使得
$$
A \supset \ldots \supset F^nA \supset F^{n + 1}A \supset \ldots \supset 0
$$
\item $\mathcal{A}$ 的一个{\it 滤过对象}，是由 $\mathcal{A}$ 的对象
$A$ 及 $A$ 上的递减滤过 $F$ 组成的二元组 $(A, F)$。
\item 一个{\it 滤过对象态射 $(A, F) \to (B, F)$}，由
$\mathcal{A}$ 中满足对所有 $i \in \mathbf{Z}$ 均有
$\varphi(F^iA) \subset F^iB$ 的态射 $\varphi : A \to B$ 给出。
\item 滤过对象的范畴记作 $\text{Fil}(\mathcal{A})$。
\item 给定滤过对象 $(A, F)$ 与子对象 $X \subset A$，$X$ 上的
{\it 诱导滤过}是满足 $F^nX = X \cap F^nA$ 的滤过。
\item 给定滤过对象 $(A, F)$ 与满射态射 $\pi : A \to Y$，
{\it 商滤过}是满足 $F^nY = \pi(F^nA)$ 的滤过。
\item 若存在 $n, m$ 使得 $F^nA = A$ 且 $F^mA = 0$，则称对象
$A$ 上的滤过 $F$ {\it 有限}。
\item 给定滤过对象 $(A, F)$，若存在 $A$ 的、包含于所有 $F^iA$ 中的最大子对象，
则称 $\bigcap F^iA$ 存在。若存在 $A$ 的、包含所有 $F^iA$ 的最小子对象，
则称 $\bigcup F^iA$ 存在。
\item 若 $\bigcap F^iA = 0$，则称滤过对象 $(A, F)$ 上的滤过
{\it 分离}；若 $\bigcup F^iA = A$，则称其{\it 穷尽}。
\end{enumerate}
\end{definition}

\noindent
略用记号时，若滤过对象态射 $f : (A, F) \to (B, F)$ 所对应的
$f : A \to B$ 在阿贝尔范畴 $\mathcal{A}$ 中是单射态射，我们就称
该态射{\it 单射}。类似地，若 $f : A \to B$ 在范畴 $\mathcal{A}$ 中是
满射态射，就称 $f$ {\it 满射}。单射（相应地，满射）等价于在
$\text{Fil}(\mathcal{A})$ 中是单态射（相应地，满态射）。由引理
\ref{homology-lemma-filtered}，这也等价于具有零核（相应地，零余核）。

\begin{lemma}
\label{homology-lemma-filtered}
设 $\mathcal{A}$ 为阿贝尔范畴。滤过对象范畴
$\text{Fil}(\mathcal{A})$ 具有下列性质：
\begin{enumerate}
\item 它是加性范畴。
\item 它有零对象。
\item 它有核与余核、像与余像。
\item 一般而言，它不是阿贝尔范畴。
\end{enumerate}
\end{lemma}

\begin{proof}
显然 $\text{Fil}(\mathcal{A})$ 是加性的；其直和由
$(A, F) \oplus (B, F) = (A \oplus B, F)$ 给出，其中
$F^p(A \oplus B) = F^pA \oplus F^pB$。滤过对象态射
$f : (A, F) \to (B, F)$ 的核，是嵌入 $\Ker(f) \subset A$，
其中 $\Ker(f)$ 配备诱导滤过。滤过对象态射 $f : A \to B$ 的余核，
是满射态射 $B \to \Coker(f)$，其中 $\Coker(f)$ 配备商滤过。
既然所有核与余核都存在，所有余像与像也存在。最后一个陈述见例
\ref{homology-example-not-abelian}。
\end{proof}

\begin{definition}
\label{homology-definition-strict}
设 $\mathcal{A}$ 为阿贝尔范畴。若对所有 $i \in \mathbf{Z}$ 均有
$f(F^iA) = f(A) \cap F^iB$，则称 $\mathcal{A}$ 的滤过对象之间的态射
$f : A \to B$ 是{\it 严格的}。
\end{definition}

\noindent
这也等价于要求对所有 $i \in \mathbf{Z}$ 均有
$f^{-1}(F^iB) = F^iA + \Ker(f)$。如下刻画严格态射。

\begin{lemma}
\label{homology-lemma-characterize-strict-general}
设 $\mathcal{A}$ 为阿贝尔范畴，并设 $f : A \to B$ 是
$\mathcal{A}$ 的滤过对象之间的态射。下列条件等价：
\begin{enumerate}
\item $f$ 是严格的；
\item 引理 \ref{homology-lemma-coim-im-map} 的态射
$\Coim(f) \to \Im(f)$ 是同构。
\end{enumerate}
\end{lemma}

\begin{proof}
注意 $\Coim(f) \to \Im(f)$ 是 $\mathcal{A}$ 的对象之间的同构，而
(2) 表示它是滤过对象之间的同构。由引理 \ref{homology-lemma-filtered} 的证明中
对核与余核的描述，$\Coim(f)$ 上的滤过是来自
$A \to \Coim(f)$ 的商滤过。类似地，$\Im(f)$ 上的滤过是来自嵌入
$\Im(f) \to B$ 的诱导滤过。严格性的定义恰好说明该商滤过就是该诱导滤过。
\end{proof}

\begin{lemma}
\label{homology-lemma-add-summand-strict-monomorphism}
设 $\mathcal{A}$ 为阿贝尔范畴。设 $f : A \to B$ 是滤过对象之间的
严格单态射，并设 $g : A \to C$ 是滤过对象态射。那么
$f \oplus g : A \to B \oplus C$ 是严格单态射。
\end{lemma}

\begin{proof}
由定义显然。
\end{proof}

\begin{lemma}
\label{homology-lemma-add-summand-strict-epimorphism}
设 $\mathcal{A}$ 为阿贝尔范畴。设 $f : B \to A$ 是滤过对象之间的
严格满态射，并设 $g : C \to A$ 是滤过对象态射。那么
$f \oplus g : B \oplus C \to A$ 是严格满态射。
\end{lemma}

\begin{proof}
由定义显然。
\end{proof}

\begin{lemma}
\label{homology-lemma-induced-and-quotient-strict}
设 $\mathcal{A}$ 为阿贝尔范畴，并设 $(A, F)$、$(B, F)$ 为滤过对象，
$u : A \to B$ 为滤过对象态射。若 $u$ 是单射态射，则 $u$ 严格当且仅当
$A$ 上的滤过是诱导滤过。若 $u$ 是满射态射，则 $u$ 严格当且仅当
$B$ 上的滤过是商滤过。
\end{lemma}

\begin{proof}
由定义立即可得。
\end{proof}

\begin{lemma}
\label{homology-lemma-composition-strict}
设 $\mathcal{A}$ 为阿贝尔范畴。设 $f : A \to B$、$g : B \to C$
为滤过对象之间的严格态射。
\begin{enumerate}
\item 一般而言，复合 $g \circ f$ 不严格。
\item 若 $g$ 是单射态射，则 $g \circ f$ 严格。
\item 若 $f$ 是满射态射，则 $g \circ f$ 严格。
\end{enumerate}
\end{lemma}

\begin{proof}
设 $B$ 为域 $k$ 上以 $e_1, e_2$ 为基的向量空间，并带有如下滤过：
当 $n < 0$ 时 $F^nB = B$，$F^0B = ke_1$，而当 $n > 0$ 时
$F^nB = 0$。现取 $A = k(e_1 + e_2)$ 与 $C = B/ke_2$，并赋予由
$B$ 诱导的滤过，即使得 $A \to B$ 与 $B \to C$ 严格的滤过
（引理 \ref{homology-lemma-induced-and-quotient-strict}）。于是，当 $n < 0$ 时
$F^n(A) = A$，而当 $n \geq 0$ 时 $F^n(A) = 0$。另外，当
$n \leq 0$ 时 $F^n(C) = C$，而当 $n > 0$ 时 $F^n(C) = 0$。
故非零复合 $A \to C$ 不严格。

\medskip\noindent
假设 $g$ 是单射态射。则
\begin{align*}
g(f(F^pA)) & = g(f(A) \cap F^pB) \\
& = g(f(A)) \cap g(F^p(B)) \\
& = (g \circ f)(A) \cap (g(B) \cap F^pC) \\
& = (g \circ f)(A) \cap F^pC.
\end{align*}
第一个等式来自 $f$ 严格；第二个来自 $g$ 是单射态射；第三个来自
$g$ 严格；第四个来自 $(g \circ f)(A) \subset g(B)$。

\medskip\noindent
假设 $f$ 是满射态射。则
\begin{align*}
(g \circ f)^{-1}(F^iC) & = f^{-1}(F^iB + \Ker(g)) \\
& = f^{-1}(F^iB) + f^{-1}(\Ker(g)) \\
& = F^iA + \Ker(f) + \Ker(g \circ f) \\
& = F^iA + \Ker(g \circ f)
\end{align*}
第一个等式来自 $g$ 严格；第二个来自 $f$ 是满射态射；第三个来自
$f$ 严格；最后一个来自 $\Ker(f) \subset \Ker(g \circ f)$。
\end{proof}

\noindent
下面的引理说明，滤过对象的子对象具有良定义的滤过，与把该对象写成余核的
方式无关。

\begin{lemma}
\label{homology-lemma-filtration-subobject}
设 $\mathcal{A}$ 为阿贝尔范畴，$(A, F)$ 为 $\mathcal{A}$ 的滤过对象，
并设 $X \subset Y \subset A$ 为 $A$ 的子对象。在对象
$$
Y/X = \Ker(A/X \to A/Y)
$$
上，来自 $Y$ 的诱导滤过的商滤过，与来自 $A/X$ 的商滤过的诱导滤过
相同。态射 $X \to Y$、$X \to A$、$Y \to A$、$Y \to A/X$、
$Y \to Y/X$、$Y/X \to A/X$ 中任一个，在配备诱导滤过或商滤过后都是
严格的。
\end{lemma}

\begin{proof}
$Y/X$ 上的商滤过由
$F^p(Y/X) = F^pY/(X \cap F^pY) = F^pY/F^pX$ 给出，因为
$F^pY = Y \cap F^pA$ 且 $F^pX = X \cap F^pA$。由嵌入
$Y/X \to A/X$ 得到的诱导滤过为
\begin{align*}
F^p(Y/X) & = Y/X \cap F^p(A/X) \\
& = Y/X \cap (F^pA + X)/X \\
& = (Y \cap F^pA)/(X \cap F^pA) \\
& = F^pY/F^pX.
\end{align*}
这证明了引理的第一个陈述。其余各情形的证明类似。
\end{proof}

\begin{lemma}
\label{homology-lemma-pushout-filtered}
设 $\mathcal{A}$ 为阿贝尔范畴，并设
$A, B, C \in \text{Fil}(\mathcal{A})$。设 $f : A \to B$ 与
$g : A \to C$ 为态射。那么在 $\text{Fil}(\mathcal{A})$ 中存在推出
$$
\xymatrix{
A \ar[r]_f \ar[d]_g & B \ar[d]^{g'} \\
C \ar[r]^{f'} & C \amalg_A B
}
$$
若 $f$ 严格，则 $f'$ 也严格。
\end{lemma}

\begin{proof}
令 $C \amalg_A B$ 等于 $\text{Fil}(\mathcal{A})$ 中的
$\Coker((g, -f) : A \to C \oplus B)$。由引理 \ref{homology-lemma-filtered}，
该余核存在。它是推出；见例 \ref{homology-example-fibre-product-pushouts}。
注意 $F^p(C \amalg_A B)$ 是 $F^pC \oplus F^pB$ 的像。因此
$$
(f')^{-1}(F^p(C \amalg_A B)) = g(f^{-1}(F^pB)) + F^pC
$$
最后一个陈述由此可得。
\end{proof}

\begin{lemma}
\label{homology-lemma-fibre-product-filtered}
设 $\mathcal{A}$ 为阿贝尔范畴，并设
$A, B, C \in \text{Fil}(\mathcal{A})$。设 $f : B \to A$ 与
$g : C \to A$ 为态射。那么在 $\text{Fil}(\mathcal{A})$ 中存在纤维积
$$
\xymatrix{
B \times_A C \ar[r]_{g'} \ar[d]_{f'} & B \ar[d]^f \\
C \ar[r]^g & A
}
$$
若 $f$ 严格，则 $f'$ 也严格。
\end{lemma}

\begin{proof}
本引理与引理 \ref{homology-lemma-pushout-filtered} 对偶。
\end{proof}

\noindent
设 $\mathcal{A}$ 为阿贝尔范畴，并设 $(A, F)$ 为 $\mathcal{A}$ 的滤过对象。
用 $\text{gr}^p_F(A) = \text{gr}^p(A)$ 表示 $\mathcal{A}$ 的对象
$F^pA/F^{p + 1}A$。这定义了加性函子
$$
\text{gr}^p :
\text{Fil}(\mathcal{A})
\longrightarrow
\mathcal{A}, \quad
(A, F)
\longmapsto
\text{gr}^p(A).
$$
回忆我们已在第 \ref{homology-section-graded} 节定义 $\mathcal{A}$ 的分次对象范畴
$\text{Gr}(\mathcal{A})$。对 $\text{Fil}(\mathcal{A})$ 中的
$(A, F)$，可置
$$
\text{gr}(A) = \text{下列对象的分次对象： }\mathcal{A}\text{ 其 }
p\text{次分次部分为 }\text{gr}^p(A)
$$
若 $\mathcal{A}$ 有可数直和，则简单地有
$$
\text{gr}(A) = \bigoplus \text{gr}^p(A)
$$
这定义了加性函子
$$
\text{gr} :
\text{Fil}(\mathcal{A})
\longrightarrow
\text{Gr}(\mathcal{A}), \quad
(A, F) \longmapsto \text{gr}(A).
$$

\begin{lemma}
\label{homology-lemma-ses-gr}
设 $\mathcal{A}$ 为阿贝尔范畴。
\begin{enumerate}
\item 设 $A$ 为滤过对象且 $X \subset A$。那么对每个 $p$，序列
$$
0 \to \text{gr}^p(X) \to \text{gr}^p(A) \to \text{gr}^p(A/X) \to 0
$$
正合（其中 $X$ 配备诱导滤过，$A/X$ 配备商滤过）。
\item 设 $f : A \to B$ 为 $\mathcal{A}$ 的滤过对象之间的态射。
那么对每个 $p$，序列
$$
0 \to \text{gr}^p(\Ker(f)) \to \text{gr}^p(A) \to
\text{gr}^p(\Coim(f)) \to 0
$$
及
$$
0 \to \text{gr}^p(\Im(f)) \to \text{gr}^p(B) \to
\text{gr}^p(\Coker(f)) \to 0
$$
正合。
\end{enumerate}
\end{lemma}

\begin{proof}
我们有 $F^{p + 1}X = X \cap F^{p + 1}A$，故映射
$\text{gr}^p(X) \to \text{gr}^p(A)$ 是单射态射。对偶地，映射
$\text{gr}^p(A) \to \text{gr}^p(A/X)$ 是满射态射。映射
$F^pA/F^{p + 1}A \to A/X + F^{p + 1}A$ 的核显然为
$F^{p + 1}A + X \cap F^pA/F^{p + 1}A = F^pX/F^{p + 1}X$，
故在中间正合。(2) 的两个短正合列是 (1) 的短正合列之特例。
\end{proof}

\begin{lemma}
\label{homology-lemma-characterize-strict}
设 $\mathcal{A}$ 为阿贝尔范畴，并设 $f : A \to B$ 为
$\mathcal{A}$ 的有限滤过对象之间的态射。下列条件等价：
\begin{enumerate}
\item $f$ 严格；
\item 态射 $\Coim(f) \to \Im(f)$ 是同构；
\item $\text{gr}(\Coim(f)) \to \text{gr}(\Im(f))$ 是同构；
\item 序列 $\text{gr}(\Ker(f)) \to \text{gr}(A) \to \text{gr}(B)$ 正合；
\item 序列 $\text{gr}(A) \to \text{gr}(B) \to
\text{gr}(\Coker(f))$ 正合；以及
\item 序列
$$
0 \to
\text{gr}(\Ker(f)) \to
\text{gr}(A) \to
\text{gr}(B) \to
\text{gr}(\Coker(f)) \to 0
$$
正合。
\end{enumerate}
\end{lemma}

\begin{proof}
(1) 与 (2) 的等价性是引理 \ref{homology-lemma-characterize-strict-general}。
由引理 \ref{homology-lemma-ses-gr}，可见 (4)、(5)、(6) 蕴含 (3)，且 (3)
蕴含 (4)、(5)、(6)。因此只须证明 (3) 蕴含 (2)。于是我们要证明：
若 $f : A \to B$ 是有限滤过对象之间的单射且满射的映射，并诱导同构
$\text{gr}(A) \to \text{gr}(B)$，则 $f$ 诱导滤过对象之间的同构。
换言之，要证明对所有 $p$ 均有 $f(F^pA) = F^pB$。由于滤过有限，
可对 $p$ 作降归纳。假设 $f(F^{p + 1}A) = F^{p + 1}B$。那么交换图
$$
\xymatrix{
0 \ar[r] &
F^{p + 1}A \ar[r] \ar[d]^f &
F^pA \ar[r] \ar[d]^f &
\text{gr}^p(A) \ar[r] \ar[d]^{\text{gr}^p(f)} &
0 \\
0 \ar[r] &
F^{p + 1}B \ar[r] &
F^pB \ar[r] &
\text{gr}^p(B) \ar[r] &
0
}
$$
与五引理蕴含 $f(F^pA) = F^pB$。
\end{proof}

\begin{lemma}
\label{homology-lemma-filtered-complex}
设 $\mathcal{A}$ 为阿贝尔范畴。设 $A \to B \to C$ 是
$\mathcal{A}$ 的滤过对象所成的复形。假设
$\alpha : A \to B$ 与 $\beta : B \to C$ 是滤过对象之间的严格态射。
那么
$\text{gr}(\Ker(\beta)/\Im(\alpha)) =
\Ker(\text{gr}(\beta))/\Im(\text{gr}(\alpha)))$。
% P02-E0014: source line 5082 contains an unmatched closing parenthesis.
\end{lemma}

\begin{proof}
这形式地来自引理 \ref{homology-lemma-ses-gr}，以及由引理
\ref{homology-lemma-characterize-strict-general} 得到的
$\Coim(\alpha) \cong \Im(\alpha)$ 与
$\Coim(\beta) \cong \Im(\beta)$。
\end{proof}

\begin{lemma}
\label{homology-lemma-filtered-acyclic}
设 $\mathcal{A}$ 为阿贝尔范畴。设 $A \to B \to C$ 是
$\mathcal{A}$ 的滤过对象所成的复形。假设 $A, B, C$ 具有有限滤过，
且 $\text{gr}(A) \to \text{gr}(B) \to \text{gr}(C)$ 正合。那么
\begin{enumerate}
\item 对每个 $p \in \mathbf{Z}$，序列
$\text{gr}^p(A) \to \text{gr}^p(B) \to \text{gr}^p(C)$ 正合；
\item 对每个 $p \in \mathbf{Z}$，序列
$F^p(A) \to F^p(B) \to F^p(C)$ 正合；
\item 对每个 $p \in \mathbf{Z}$，序列
$A/F^p(A) \to B/F^p(B) \to C/F^p(C)$ 正合；
\item 映射 $A \to B$ 与 $B \to C$ 严格；以及
\item $A \to B \to C$ 正合（作为 $\mathcal{A}$ 中的序列）。
\end{enumerate}
\end{lemma}

\begin{proof}
(1) 由定义立即可得。我们对滤过的长度归纳证明 (3)。若 $A$、$B$、$C$
都只有一个非零的分次部分，则 (3) 成立，因为 $\text{gr}(A) = A$ 等。
令 $n$ 为使 $F^nA, F^nB, F^nC$ 中至少一个非零的最大整数。令
$A' = A/F^nA$、$B' = B/F^nB$、$C' = C/F^nC$，并配备诱导滤过。
注意 $\text{gr}(A) = F^nA \oplus \text{gr}(A')$，对 $B$ 与 $C$ 亦然。
归纳假设可应用于 $A' \to B' \to C'$，从而蕴含当 $p \geq n$ 时
$A/F^p(A) \to B/F^p(B) \to C/F^p(C)$ 正合。为对 $p = n + 1$
得到同一结论，即证明 $A \to B \to C$ 正合，使用交换图
$$
\xymatrix{
0 \ar[r] & F^nA \ar[r] \ar[d] & A \ar[r] \ar[d] & A' \ar[r] \ar[d] & 0 \\
0 \ar[r] & F^nB \ar[r] \ar[d] & B \ar[r] \ar[d] & B' \ar[r] \ar[d] & 0 \\
0 \ar[r] & F^nC \ar[r] & C \ar[r] & C' \ar[r] & 0
}
$$
其各行是 $\mathcal{A}$ 中对象的短正合列。(2) 的证明与此对偶。
当然，(5) 由 (2) 得到。

\medskip\noindent
为证明 (4)，用 $f : A \to B$ 与 $g : B \to C$ 表示给定态射。
由 (2) 知 $f(F^p(A)) = \Ker(F^p(B) \to F^p(C))$，由 (5) 知
$f(A) = \Ker(g)$。因此
$f(F^p(A)) =  \Ker(F^p(B) \to F^p(C)) =
\Ker(g) \cap F^p(B) = f(A) \cap F^p(B)$，这证明 $f$ 严格。
$g$ 严格的证明与此对偶。
\end{proof}

\section{谱序列}
\label{homology-section-spectral-sequence}

\noindent
关于谱序列的精彩讨论可见 \cite{Eisenbud}。亦见 \cite{McCleary}、
\cite{Lang} 等。

\begin{definition}
\label{homology-definition-spectral-sequence}
设 $\mathcal{A}$ 为阿贝尔范畴。
\begin{enumerate}
\item $\mathcal{A}$ 中的一个{\it 谱序列}由系统
$(E_r, d_r)_{r \geq 1}$ 给出，其中每个 $E_r$ 是 $\mathcal{A}$ 的对象，
每个 $d_r : E_r \to E_r$ 是满足 $d_r \circ d_r = 0$ 的态射，且对
$r \geq 1$ 有 $E_{r + 1} = \Ker(d_r)/\Im(d_r)$。
\item 一个{\it 谱序列态射}
$f : (E_r, d_r)_{r \geq 1} \to (E'_r, d'_r)_{r \geq 1}$，
由一族态射 $f_r : E_r \to E'_r$ 给出，并满足
$f_r \circ d_r = d'_r \circ f_r$；此外，在等同
$E_{r + 1} = \Ker(d_r)/\Im(d_r)$ 与
$E'_{r + 1} = \Ker(d'_r)/\Im(d'_r)$ 下，$f_{r + 1}$ 是由 $f_r$
诱导的态射。
\end{enumerate}
\end{definition}

\noindent
有时我们会稍微放宽这一定义，允许 $E_{r + 1}$ 是配备给定同构
$E_{r + 1} \to \Ker(d_r)/\Im(d_r)$ 的对象。此外，有时对某个
$r_0 \in \mathbf{Z}$，我们有系统 $(E_r, d_r)_{r \geq r_0}$，
它对指标 $\geq r_0$ 满足上面定义的性质。我们也称其为谱序列，因为简单地
重新编号后，它无论如何都落在该定义之内。事实上，文献中可见
$r_0 = 0$ 与 $r_0 = -1$ 的情形。

\medskip\noindent
给定谱序列 $(E_r, d_r)_{r \geq 1}$，按如下简单过程定义
$$
0 = B_1 \subset B_2 \subset \ldots \subset B_r \subset \ldots
\subset Z_r \subset \ldots \subset Z_2 \subset Z_1 = E_1
$$
令 $B_2 = \Im(d_1)$、$Z_2 = \Ker(d_1)$。于是显然有
$d_2 : Z_2/B_2 \to Z_2/B_2$。因此可把 $B_3$ 定义为包含 $B_2$
且使 $B_3/B_2$ 为 $d_2$ 之像的 $E_1$ 的唯一子对象。类似地，可把
$Z_3$ 定义为包含 $B_2$ 且使 $Z_3/B_2$ 为 $d_2$ 之核的 $E_1$
的唯一子对象。依此类推。特别地，对所有 $r \geq 1$ 有
$$
E_r = Z_r/B_r
$$
若谱序列始于 $r = r_0$，则类似地可把 $B_i$、$Z_i$ 构造为
$E_{r_0}$ 的子对象。事实上，文献中有时使用记号
$$
0 = B_r(E_r) \subset B_{r + 1}(E_r) \subset B_{r + 2}(E_r) \subset \ldots
\subset Z_{r + 2}(E_r) \subset Z_{r + 1}(E_r) \subset Z_r(E_r) = E_r
$$
表示上面所述但以 $E_r$ 开始的滤过。

\begin{definition}
\label{homology-definition-limit-spectral-sequence}
设 $\mathcal{A}$ 为阿贝尔范畴，并设 $(E_r, d_r)_{r \geq 1}$ 为谱序列。
\begin{enumerate}
\item 若 $E_1$ 的子对象 $Z_{\infty} = \bigcap Z_r$ 与
$B_{\infty} = \bigcup B_r$ 存在，则把谱序列的
{\it 极限}\footnote{这一记号并未得到普遍接受。在一些参考文献中，
$E_1$ 的一对子对象 $Z_\infty$ 与 $B_\infty$，满足
$0 = B_1 \subset B_2 \subset \ldots \subset B_\infty \subset Z_\infty
\subset \ldots \subset Z_2 \subset Z_1 = E_1$，是构成谱序列的数据的一部分！}
定义为对象 $E_{\infty} = Z_{\infty}/B_{\infty}$。
\item 若微分 $d_r, d_{r + 1}, \ldots$ 全都为零，则称该谱序列
{\it 在 $E_r$ 处退化}。
\end{enumerate}
\end{definition}

\noindent
注意，若谱序列在 $E_r$ 处退化，则
$E_r = E_{r + 1} = \ldots = E_{\infty}$（而且极限当然存在）。
此外，我们将遇到的几乎每个阿贝尔范畴都有可数和与交。

\begin{remark}[变体]
\label{homology-remark-allow-translation-functors}
谱序列的各项通常带有附加结构，例如分次或双分次。为容纳这种情形
（并绕开某些技术问题），引入如下概念。设 $\mathcal{A}$ 为阿贝尔范畴，
并设 $(T_r)_{r \geq 1}$ 为一列{\it 平移}函子或{\it 移位}函子，
也就是说 $T_r : \mathcal{A} \to \mathcal{A}$ 是范畴同构。
在这一设定下，一个{\it 谱序列}由系统 $(E_r, d_r)_{r \geq 1}$ 给出，
其中每个 $E_r$ 是 $\mathcal{A}$ 的对象，每个
$d_r : E_r \to T_rE_r$ 是满足 $T_rd_r \circ d_r = 0$ 的态射，从而
$$
\xymatrix{
\ldots \ar[r] &
T_r^{-1}E_r \ar[r]^-{T_r^{-1}d_r} &
E_r \ar[r]^-{d_r} &
T_rE_r \ar[r]^{T_r d_r} &
T_r^2E_r \ar[r] & \ldots
}
$$
是复形，并且对 $r \geq 1$ 有
$E_{r + 1} = \Ker(d_r)/\Im(T_r^{-1}d_r)$。在这一设定下，
{\it 谱序列态射}的含义是清楚的。在这一设定下仍可定义
$$
0 = B_1 \subset B_2 \subset \ldots \subset B_r \subset \ldots
\subset Z_r \subset \ldots \subset Z_2 \subset Z_1 = E_1
$$
以及如上的 $Z_\infty$ 与 $B_\infty$（如果它们存在）。
\end{remark}

\section{谱序列：正合偶}
\label{homology-section-exact-couple}

\begin{definition}
\label{homology-definition-exact-couple}
设 $\mathcal{A}$ 为阿贝尔范畴。
\begin{enumerate}
\item 一个{\it 正合偶}是数据 $(A, E, \alpha, f, g)$，其中
$A$、$E$ 是 $\mathcal{A}$ 的对象，而 $\alpha$、$f$、$g$ 是下图中的态射：
$$
\xymatrix{
A \ar[rr]_{\alpha} & & A \ar[ld]^g \\
& E \ar[lu]^f &
}
$$
并要求每个箭头的核都是其前一箭头的像。因此
$\Ker(\alpha) = \Im(f)$、$\Ker(f) = \Im(g)$，且
$\Ker(g) = \Im(\alpha)$。
\item 一个{\it 正合偶态射}
$t : (A, E, \alpha, f, g) \to (A', E', \alpha', f', g')$，
由态射 $t_A : A \to A'$ 与 $t_E : E \to E'$ 给出，并满足
$\alpha' \circ t_A = t_A \circ \alpha$、
$f' \circ t_E = t_A \circ f$，且 $g' \circ t_A = t_E \circ g$。
\end{enumerate}
\end{definition}

\begin{lemma}
\label{homology-lemma-derived-exact-couple}
设 $(A, E, \alpha, f, g)$ 是阿贝尔范畴 $\mathcal{A}$ 中的正合偶。
置
\begin{enumerate}
\item $d = g \circ f : E \to E$，从而 $d \circ d = 0$；
\item $E' = \Ker(d)/\Im(d)$；
\item $A' = \Im(\alpha)$；
\item $\alpha' : A' \to A'$ 为 $\alpha$ 诱导的态射；
\item $f' : E' \to A'$ 为 $f$ 诱导的态射；
\item $g' : A' \to E'$ 为“$g \circ \alpha^{-1}$”诱导的态射。
\end{enumerate}
那么
\begin{enumerate}
\item $\Ker(d) = f^{-1}(\Ker(g)) = f^{-1}(\Im(\alpha))$；
\item $\Im(d) = g(\Im(f)) = g(\Ker(\alpha))$；
\item $(A', E', \alpha', f', g')$ 是正合偶。
\end{enumerate}
\end{lemma}

\begin{proof}
略。
\end{proof}

\noindent
因此显然，给定正合偶 $(A, E, \alpha, f, g)$，令
$E_1 = E$、$d_1 = d$、$E_2 = E'$、$d_2 = d' = g' \circ f'$、
$E_3 = E''$、$d_3 = d'' = g'' \circ f''$，依此类推，就得到谱序列。

\begin{definition}
\label{homology-definition-spectral-sequence-associated-exact-couple}
设 $\mathcal{A}$ 为阿贝尔范畴，并设 $(A, E, \alpha, f, g)$ 为正合偶。
{\it 关联于该正合偶的谱序列}，是谱序列
$(E_r, d_r)_{r \geq 1}$，其中 $E_1 = E$、$d_1 = d$、
$E_2 = E'$、$d_2 = d' = g' \circ f'$、$E_3 = E''$、
$d_3 = d'' = g'' \circ f''$，依此类推。
\end{definition}

\begin{lemma}
\label{homology-lemma-spectral-sequence-associated-exact-couple}
设 $\mathcal{A}$ 为阿贝尔范畴，$(A, E, \alpha, f, g)$ 为正合偶，
并设 $(E_r, d_r)_{r \geq 1}$ 为关联于该正合偶的谱序列。在这种情况下有
$$
0 = B_1 \subset \ldots \subset
B_{r + 1} = g(\Ker(\alpha^r))
\subset \ldots \subset
Z_{r + 1} = f^{-1}(\Im(\alpha^r))
\subset \ldots \subset Z_1 = E
$$
且映射 $d_{r + 1} : E_{r + 1} \to E_{r + 1}$ 由下列规则描述：
对 $\mathcal{A}$ 的任意（测试）对象 $T$ 以及任意元素
$x : T \to Z_{r + 1}$、$y : T \to A$，若
$f \circ x = \alpha^r \circ y$，则
$$
d_{r + 1} \circ \overline{x} = \overline{g \circ y}
$$
其中 $\overline{x} : T \to E_{r + 1}$ 是诱导态射。
\end{lemma}

\begin{proof}
略。
\end{proof}

\noindent
注意，在引理的情形中，只要 $\bigcup \Ker(\alpha^r)$ 与
$\bigcap \Im(\alpha^r)$ 存在，显然有
$$
B_\infty = g\left(\bigcup\nolimits_r \Ker(\alpha^r)\right)
\subset
Z_\infty = f^{-1}\left(\bigcap\nolimits_r \Im(\alpha^r)\right)
$$
这产生极限 $E_\infty = Z_\infty / B_\infty$；见定义
\ref{homology-definition-limit-spectral-sequence}。

\begin{remark}[变体]
\label{homology-remark-shifted-exact-couple}
设 $\mathcal{A}$ 为阿贝尔范畴，并设
$S, T : \mathcal{A} \to \mathcal{A}$ 为移位函子，即范畴同构。
对 $A \in \Ob(\mathcal{A})$ 与 $n \in \mathbf{Z}$，我们用
$S^nA$ 与 $T^nA$ 表示 $n$ 重复合。在这种情况下，一个{\it 正合偶}
是数据 $(A, E, \alpha, f, g)$，其中 $A$、$E$ 是 $\mathcal{A}$ 的对象，
而 $\alpha : A \to T^{-1}A$、$f : E \to A$、$g : A \to SE$
为态射，使得
$$
\xymatrix{
TE \ar[r]^-{Tf} &
TA \ar[r]^-{T\alpha} &
A \ar[r]^-{g} &
SE \ar[r]^{Sf} & SA
}
$$
是正合复形。可将其图示如下：
$$
\xymatrix{
TA \ar[rrr]_{T\alpha} & & &
A \ar[ld]^g \ar[rrr]_\alpha & & &
T^{-1}A \ar[ld]^{T^{-1}g} \\
& TE \ar[lu]^{Tf} \ar@{..}[r] & SE & &
E \ar[lu]^f \ar@{..}[r] & T^{-1}SE
}
$$
置 $d = g \circ f : E \to SE$。那么
$d \circ S^{-1}d = g \circ f \circ S^{-1}g \circ S^{-1}f = 0$，
因为 $f \circ S^{-1}g = 0$。置
$E' = \Ker(d)/\Im(S^{-1}d)$、$A' = \Im(T\alpha)$。
令 $\alpha' : A' \to T^{-1}A'$ 为 $\alpha$ 诱导的态射。令
$f' : E' \to A'$ 为 $f$ 诱导的态射；这是可行的，因为
$f(\Ker(d)) \subset \Ker(g) = \Im(T\alpha)$。最后，令
$g' : A' \to TSE'$ 为“$Tg \circ (T\alpha)^{-1}$”诱导的态射\footnote{这是
可行的，因为 $TSE' = \Ker(TSd)/\Im(Td)$，且
$Tg(\Ker(T\alpha)) = Tg(\Im(Tf)) = \Im(T(d))$，以及
$TS(d)(\Im(Tg)) = \Im(TSg \circ TSf \circ Tg) = 0$。}。

\medskip\noindent
与上文完全相同地可得
\begin{enumerate}
\item $\Ker(d) = f^{-1}(\Ker(g)) = f^{-1}(\Im(T\alpha))$；
\item $\Im(d) = g(\Im(f)) = g(\Ker(\alpha))$；
\item $(A', E', \alpha', f', g')$ 是关于移位函子 $TS$ 与 $T$ 的正合偶。
\end{enumerate}
我们得到一个谱序列（如注记 \ref{homology-remark-allow-translation-functors}），
其 $E_1 = E$、$E_2 = E'$ 等，并且对所有 $r \geq 1$ 有
$d_r : E_r \to T^{r - 1}SE_r$。引理
\ref{homology-lemma-spectral-sequence-associated-exact-couple} 告诉我们，在这种情况下
$$
SB_{r + 1} =
g(\Ker(T^{-r + 1}\alpha \circ \ldots \circ T^{-1}\alpha \circ \alpha))
$$
以及
$$
Z_{r + 1} = f^{-1}(\Im(T\alpha \circ T^2\alpha \circ \ldots \circ T^r\alpha))
$$
映射 $d_{r + 1}$ 的描述与引理中给出的类似。（也许使用这些显式描述来证明
由这样的正合偶得到谱序列会更容易。）
\end{remark}

\section{谱序列：微分对象}
\label{homology-section-differential-object}

\begin{definition}
\label{homology-definition-differential-object}
设 $\mathcal{A}$ 为阿贝尔范畴。$\mathcal{A}$ 的一个{\it 微分对象}，
是二元组 $(A, d)$，其中 $A$ 是 $\mathcal{A}$ 的对象，并配备满足
$d \circ d = 0$ 的自映射 $d$。一个{\it 微分对象态射}
$(A, d) \to (B, d)$，由满足 $d \circ \alpha = \alpha \circ d$ 的
态射 $\alpha : A \to B$ 给出。
\end{definition}

\begin{lemma}
\label{homology-lemma-differential-objects-abelian}
\begin{slogan}
阿贝尔范畴的微分对象之范畴本身是阿贝尔范畴。
\end{slogan}
设 $\mathcal{A}$ 为阿贝尔范畴。$\mathcal{A}$ 的微分对象之范畴是阿贝尔范畴。
\end{lemma}

\begin{proof}
略。
\end{proof}

\begin{definition}
\label{homology-definition-differential-object-homology}
对微分对象 $(A, d)$，用
$$
H(A, d) = \Ker(d)/\Im(d)
$$
表示它的{\it 同调}。
\end{definition}

\begin{lemma}
\label{homology-lemma-differential-objects-ses}
设 $\mathcal{A}$ 为阿贝尔范畴，并设
$0 \to (A, d) \to (B, d) \to (C, d) \to 0$ 为微分对象的短正合列。
那么得到正合同调列
$$
\ldots \to H(C, d) \to H(A, d) \to H(B, d) \to H(C, d) \to \ldots
$$
\end{lemma}

\begin{proof}
把引理 \ref{homology-lemma-long-exact-sequence-cochain} 应用于复形短正合列
$$
\begin{matrix}
0 & \to & A & \to & B & \to & C & \to & 0 \\
& & \downarrow & & \downarrow & & \downarrow \\
0 & \to & A & \to & B & \to & C & \to & 0 \\
& & \downarrow & & \downarrow & & \downarrow \\
0 & \to & A & \to & B & \to & C & \to & 0
\end{matrix}
$$
其中竖直箭头为 $d$。
\end{proof}

\noindent
下面给出谱序列的一个重要例子。设 $\mathcal{A}$ 为阿贝尔范畴，
$(A, d)$ 为 $\mathcal{A}$ 的微分对象，并设
$\alpha : (A, d) \to (A, d)$ 为该微分对象的自同态。若假设
$\alpha$ 是单射态射，则得到微分对象的短正合列
$$
0 \to (A, d) \to (A, d) \to (A/\alpha A, d) \to 0
$$
由引理 \ref{homology-lemma-differential-objects-ses}，得到正合偶
$$
\xymatrix{
H(A, d) \ar[rr]_{\overline{\alpha}} & & H(A, d) \ar[ld]^g \\
& H(A/\alpha A, d) \ar[lu]^f &
}
$$
其中 $g$ 是典范映射，$f$ 是蛇引理中定义的映射。由第
\ref{homology-section-exact-couple} 节，得到一个相伴谱序列。由于在此情形中
$E_1 = H(A/\alpha A, d)$，定义 $E_0 = A/\alpha A$ 与 $d_0 = d$
是有意义的。换言之，我们从 $r = 0$ 开始谱序列。按照第
\ref{homology-section-spectral-sequence} 节的约定，定义一列子对象
$$
0 = B_0 \subset \ldots \subset B_r \subset \ldots
\subset Z_r \subset \ldots \subset Z_0 = E_0
$$
并满足 $E_r = Z_r/B_r$。由引理
\ref{homology-lemma-spectral-sequence-associated-exact-couple}，对 $r \geq 1$ 有
\begin{enumerate}
\item $B_r$ 是 $(\alpha^{r - 1})^{-1}(d A)$ 在自然映射
$A \to A/\alpha A$ 下的像；
\item $Z_r$ 是 $d^{-1}(\alpha^r A)$ 在自然映射
$A \to A/\alpha A$ 下的像；以及
\item $d_r : E_r \to E_r$ 如下给出：给定元素 $z \in Z_r$，选取
$y \in A$ 使得 $d(z) = \alpha^r(y)$。那么
$d_r(z + B_r + \alpha A) = y + B_r + \alpha A$。
\end{enumerate}
警告：不一定有 $\alpha A \subset (\alpha^{r - 1})^{-1}(dA)$，也不一定有
$\alpha A \subset d^{-1}(\alpha^r A)$。成立的是
$(\alpha^{r - 1})^{-1}(dA) \subset d^{-1}(\alpha^r A)$。我们有
$$
E_r
=
\frac{d^{-1}(\alpha^r A) + \alpha A}{(\alpha^{r - 1})^{-1}(dA) + \alpha A}.
$$
直接验证 (1)--(3) 给出谱序列并不困难。

\begin{definition}
\label{homology-definition-differential-object-selfmap}
设 $\mathcal{A}$ 为阿贝尔范畴，$(A, d)$ 为 $\mathcal{A}$ 的微分对象，
并设 $\alpha : A \to A$ 为 $A$ 的、与 $d$ 交换的单射自映射。{\it 关联于
$(A, d, \alpha)$ 的谱序列}，是上面描述的谱序列
$(E_r, d_r)_{r \geq 0}$。
\end{definition}

\begin{remark}[变体]
\label{homology-remark-differential-object-selfmap}
设 $\mathcal{A}$ 为阿贝尔范畴，并设
$S, T : \mathcal{A} \to \mathcal{A}$ 为移位函子，即范畴同构。
假设作为函子有 $TS = ST$。考虑二元组 $(A, d)$，其中 $A$ 是
$\mathcal{A}$ 的对象，而 $d : A \to SA$ 是满足
$d \circ S^{-1}d = 0$ 的态射。这些对象所成的范畴是阿贝尔范畴。
定义 $H(A, d) = \Ker(d)/\Im(S^{-1}d)$，并注意
$H(SA, Sd) = SH(A, d)$（典范同构）。给定短正合列
$$
0 \to (A, d) \to (B, d) \to (C, d) \to 0
$$
得到长正合同调列
$$
\ldots \to S^{-1}H(C, d) \to
H(A, d) \to H(B, d) \to H(C, d) \to SH(A, d) \to \ldots
$$
（注意边界映射中的移位）。由于 $ST = TS$，函子 $T$ 通过置
$T(A, d) = (TA, Td)$ 定义二元组范畴上的移位函子。其次，设
$\alpha : (A, d) \to T^{-1}(A, d)$ 是单射态射，其余核为 $(Q, d)$。
那么得到如注记 \ref{homology-remark-shifted-exact-couple} 中关于移位函子
$TS$ 与 $T$ 的正合偶
$$
(H(A, d), S^{-1}H(Q, d), \overline{\alpha}, f, g)
$$
其中 $\overline{\alpha} : H(A, d) \to T^{-1}H(A, d)$ 由
$\alpha$ 诱导，映射 $f : S^{-1}H(Q, d) \to H(A, d)$ 是边界映射，
而 $g : H(A, d) \to TH(Q, d) = TS(S^{-1}H(Q, d))$ 由商映射
$A \to TQ$ 诱导。因此得到如上的谱序列，其
$E_1 = S^{-1}H(Q, d)$，微分为 $d_r : E_r \to T^rSE_r$。
如上置 $E_0 = S^{-1}Q$，并令 $d_0 : E_0 \to SE_0$ 由
$S^{-1}d : S^{-1}Q \to Q$ 给出。若按照约定定义
$B_r \subset Z_r \subset E_0$，则对 $r \geq 1$ 有
\begin{enumerate}
\item $SB_r$ 是
$$
(T^{-r + 1}\alpha \circ \ldots \circ T^{-1}\alpha)^{-1}
\Im(T^{-r}S^{-1}d)
$$
在自然映射 $T^{-1}A \to Q$ 下的像；
\item $Z_r$ 是
$$
(S^{-1}T^{-1}d)^{-1}
\Im(\alpha \circ \ldots \circ T^{r - 1}\alpha)
$$
在自然映射 $S^{-1}T^{-1}A \to S^{-1}Q$ 下的像。
\end{enumerate}
微分可描述如下：若 $x \in Z_r$，选取映到 $x$ 的
$x' \in S^{-1}T^{-1}A$。那么对某个 $y \in T^{r - 1}A$，
$S^{-1}T^{-1}d(x')$ 是
$(\alpha \circ \ldots \circ T^{r - 1}\alpha)(y)$。于是
$d_r(x) \in T^rSE_r$ 由 $y$ 在 $T^rSE_0 = T^rQ$ 中的像之模
$T^rSB_r$ 的类表示。
\end{remark}

\section{谱序列：滤过微分对象}
\label{homology-section-filtered-differential}

\noindent
可以从一个滤过微分对象出发构造谱序列。

\begin{definition}
\label{homology-definition-filtered-differential}
设 $\mathcal{A}$ 为阿贝尔范畴。一个{\it 滤过微分对象}
$(K, F, d)$，是 $\mathcal{A}$ 的滤过对象 $(K, F)$，并配备平方为零的
自同态 $d : (K, F) \to (K, F)$，即 $d \circ d = 0$。
\end{definition}

\noindent
为描述关联于这种对象的谱序列，暂时假设 $\mathcal{A}$ 是具有可数直和且
可数直和正合的阿贝尔范畴（这不是自动成立的；见注记
\ref{homology-remark-direct-sums-not-exact}）。设 $(K, F, d)$ 为
$\mathcal{A}$ 的滤过微分对象。注意每个 $F^nK$ 本身都是微分对象。
考虑对象 $A = \bigoplus F^nK$，并在每个直和项上使用 $d$，从而为其
配备微分 $d$。那么 $(A, d)$ 是 $\mathcal{A}$ 的带分次微分对象。
考虑映射
$$
\alpha : A \to A
$$
它由各嵌入 $F^nK \to F^{n - 1}K$ 给出。这显然是微分对象之间的
单射态射 $\alpha : (A, d) \to (A, d)$。因此，由定义
\ref{homology-definition-differential-object-selfmap} 得到一个谱序列。
我们称它为{\it 关联于滤过微分对象 $(K, F, d)$ 的谱序列}。

\medskip\noindent
来求出这个谱序列的各项。首先，注意
$A/\alpha A = \text{gr}(K)$，并带有微分 $d = \text{gr}(d)$。
因此可见
$$
E_0 = \text{gr}(K), \quad d_0 = \text{gr}(d).
$$
于是分次微分对象 $\text{gr}(K)$ 的同调是下一项：
$$
E_1 = H(\text{gr}(K), \text{gr}(d)).
$$
此外，$E_0$ 是 $\mathcal{A}$ 的分次对象，而 $d_0$ 与该分次相容，
故 $E_1$ 显然也是分次对象。然而，微分 $d_1$ 并不保持这一分次，
而是把次数平移 $1$。

\medskip\noindent
为精确说明这一点，定义
$$
Z_r^p =
\frac{F^pK \cap d^{-1}(F^{p + r}K) + F^{p + 1}K}{F^{p + 1}K}
$$
以及
$$
B_r^p =
\frac{F^pK \cap d(F^{p - r + 1}K) + F^{p + 1}K}{F^{p + 1}K}.
$$
这一记号虽然相当自然，但似乎不同于文献中大多数地方的记号。
也许这并不重要，因为文献本身似乎也没有一致的记号选择。
按这些选择，第 \ref{homology-section-differential-object} 节中定义的
$B_r \subset E_0$（相应地，$Z_r \subset E_0$）等于
$\bigoplus_p B_r^p$（相应地，$\bigoplus_p Z_r^p$）。因此，若对
$r \geq 0$ 与 $p \in \mathbf{Z}$ 定义
$$
E_r^p = Z_r^p/B_r^p
$$
则 $E_r = \bigoplus_p E_r^p$。可按规则
$$
z + F^{p + 1}K
\longmapsto
dz + F^{p + r + 1}K
$$
定义微分 $d_r^p : E_r^p \to E_r^{p + r}$，其中
$z \in F^pK \cap d^{-1}(F^{p + r}K)$。

\begin{lemma}
\label{homology-lemma-spectral-sequence-filtered-differential}
设 $\mathcal{A}$ 为阿贝尔范畴，并设 $(K, F, d)$ 为 $\mathcal{A}$ 的
滤过微分对象。存在 $\text{Gr}(\mathcal{A})$ 中关联于 $(K, F, d)$
的谱序列 $(E_r, d_r)_{r \geq 0}$，使得对所有 $r$ 均有
$d_r : E_r \to E_r[r]$，并且分次部分 $E_r^p$ 及映射
$d_r^p : E_r^p \to E_r^{p + r}$ 如上给出。此外，
$E_0^p = \text{gr}^p K$、$d_0^p = \text{gr}^p(d)$，且
$E_1^p = H(\text{gr}^pK, \text{gr}^p(d))$。
\end{lemma}

\begin{proof}
若 $\mathcal{A}$ 有可数直和且可数直和正合，则结论来自上面的讨论。
一般情形如下处理；我们强烈建议读者跳过这个证明。考虑
$\text{Gr}(\mathcal{A})$ 的对象 $A = (F^{p + 1}K)$，即把
$F^{p + 1}K$ 置于次数 $p$（编号中的奇特平移是为了使后面的编号正确）。
通过在每个分量上使用 $d$，为其配备微分 $d$。那么 $(A, d)$ 是
$\text{Gr}(\mathcal{A})$ 的微分对象。考虑映射
$$
\alpha : A \to A[-1]
$$
它在次数 $p$ 处由嵌入 $F^{p + 1}A \to F^pA$ 给出。
% P02-E0015: source line 5773 applies F to A; the construction defined A^p=F^{p+1}K and appears to require the inclusions F^{p+1}K -> F^pK.
这显然是微分对象之间的单射态射
$\alpha : (A, d) \to (A, d)[-1]$。因此可在
$S = \text{id}$、$T = [1]$ 时应用注记
\ref{homology-remark-differential-object-selfmap}。所得
$\text{Gr}(\mathcal{A})$ 中的谱序列
$(E_r, d_r)_{r \geq 0}$ 正是我们所求的谱序列。稍微展开这些定义。
首先，$E_r = (E_r^p)$ 是 $\text{Gr}(\mathcal{A})$ 的对象。
其次，由于 $T^rS = [r]$，有 $d_r : E_r \to E_r[r]$，即
$d_r^p : E_r^p \to E_r^{p + r}$。

\medskip\noindent
为看出分次部分的描述成立，按上面的方法论证。首先有
$E_0 = \Coker(\alpha : A \to A[-1])$，由上面的编号选择，这给出
$E_0^p = \text{gr}^pK$。第一个微分为
$d_0^p = \text{gr}^pd : E_0^p \to E_0^p$。其次，注记
\ref{homology-remark-differential-object-selfmap} 中对边界 $B_r$ 与余闭元
$Z_r$ 的描述，直接转化为上面给出的 $Z_r^p$ 与 $B_r^p$ 的公式。
\end{proof}

\begin{lemma}
\label{homology-lemma-spectral-sequence-filtered-differential-d1}
设 $\mathcal{A}$ 为阿贝尔范畴，并设 $(K, F, d)$ 为 $\mathcal{A}$ 的
滤过微分对象。关联于 $(K, F, d)$ 的谱序列
$(E_r, d_r)_{r \geq 0}$ 中的
$$
d_1^p :
E_1^p = H(\text{gr}^pK)
\longrightarrow
H(\text{gr}^{p + 1}K) = E_1^{p + 1}
$$
等于关联于微分对象短正合列
$$
0 \to \text{gr}^{p + 1}K \to F^pK/F^{p + 2}K \to \text{gr}^pK \to 0.
$$
的同调边界映射。
\end{lemma}

\begin{proof}
这由引理 \ref{homology-lemma-spectral-sequence-filtered-differential} 前给出的
微分 $d_1^p$ 的公式显然可得。
\end{proof}

\begin{definition}
\label{homology-definition-filtration-cohomology-filtered-differential}
设 $\mathcal{A}$ 为阿贝尔范畴，并设 $(K, F, d)$ 为 $\mathcal{A}$ 的
滤过微分对象。$H(K, d)$ 上的{\it 诱导滤过}定义为
$F^pH(K, d) = \Im(H(F^pK, d) \to H(K, d))$。
\end{definition}

\noindent
展开其含义可见
$$
F^pH(K, d) =
\frac{\Ker(d) \cap F^pK + \Im(d)}{\Im(d)}
$$
因而
$$
\text{gr}^p H(K) =
\frac{\Ker(d) \cap F^pK + \Im(d)}{\Ker(d) \cap F^{p + 1}K + \Im(d)} =
\frac{\Ker(d) \cap F^pK}{\Ker(d) \cap F^{p + 1}K + \Im(d) \cap F^pK}
$$

\begin{lemma}
\label{homology-lemma-compute-filtered-cohomology}
设 $\mathcal{A}$ 为阿贝尔范畴，并设 $(K, F, d)$ 为 $\mathcal{A}$ 的
滤过微分对象。若 $Z_\infty^p$ 与 $B_\infty^p$ 存在（见证明），则
\begin{enumerate}
\item 极限 $E_\infty$ 存在并且是分次的，其次数 $p$ 部分为
$E_\infty^p = Z_\infty^p/B_\infty^p$；以及
\item $K$ 的同调的相伴分次对象 $\text{gr}(H(K))$ 是分次极限对象
$E_\infty$ 的分次子商。
\end{enumerate}
\end{lemma}

\begin{proof}
定义 \ref{homology-definition-limit-spectral-sequence} 的对象 $Z_\infty$、
$B_\infty$ 及极限 $E_\infty = Z_\infty/B_\infty$，按照引理
\ref{homology-lemma-spectral-sequence-filtered-differential} 的证明中对谱序列的
构造，是 $\text{Gr}(\mathcal{A})$ 的对象。
由于 $Z_r^p$（相应地，$B_r^p$）是 $Z_r$（相应地，$B_r$）的第 $p$
个分次部分，若假设
$$
Z_\infty^p = \bigcap\nolimits_r Z_r^p =
\frac{\bigcap_r (F^pK \cap d^{-1}(F^{p + r}K) + F^{p + 1}K)}{F^{p + 1}K}
$$
及
$$
B_\infty^p = \bigcup\nolimits_r B_r^p =
\frac{\bigcup_r (F^pK \cap d(F^{p - r + 1}K) + F^{p + 1}K)}{F^{p + 1}K}.
$$
存在，则 $Z_\infty$ 与 $B_\infty$ 存在，且其次数 $p$ 部分分别为
$Z_\infty^p$ 与 $B_\infty^p$（这来自关于分次子对象之并与交的初等论证）。
因此
$$
E_\infty^p =
\frac{\bigcap_r (F^pK \cap d^{-1}(F^{p + r}K) + F^{p + 1}K)}
{\bigcup_r (F^pK \cap d(F^{p - r + 1}K) + F^{p + 1}K)}.
$$
其中分子与分母都存在。我们有
\begin{equation}
\label{homology-equation-on-top}
\Ker(d) \cap F^pK + F^{p + 1}K
\subset
\bigcap\nolimits_r \left(F^pK \cap d^{-1}(F^{p + r}K) + F^{p + 1}K\right)
\end{equation}
以及
\begin{equation}
\label{homology-equation-at-bottom}
\bigcup\nolimits_r \left(F^pK \cap d(F^{p - r + 1}K) + F^{p + 1}K\right)
\subset
\Im(d) \cap F^pK + F^{p + 1}K.
\end{equation}
因此，$E_\infty^p$ 的一个子商为
$$
\frac{\Ker(d) \cap F^pK + F^{p + 1}K}{\Im(d) \cap F^pK + F^{p + 1}K} =
\frac{\Ker(d) \cap F^pK}{\Im(d) \cap F^pK + \Ker(d) \cap F^{p + 1}K}
$$
将其与定义 \ref{homology-definition-filtration-cohomology-filtered-differential}
后的讨论中给出的 $\text{gr}^pH(K)$ 公式比较，即得结论。
\end{proof}

\begin{definition}
\label{homology-definition-filtered-differential-ss-converges}
设 $\mathcal{A}$ 为阿贝尔范畴，并设 $(K, F, d)$ 为 $\mathcal{A}$ 的
滤过微分对象。称关联于 $(K, F, d)$ 的谱序列
\begin{enumerate}
\item {\it 弱收敛到 $H(K)$}，若通过引理
\ref{homology-lemma-compute-filtered-cohomology} 有
$\text{gr}H(K) = E_{\infty}$；
\item {\it 汇聚到 $H(K)$}，若它弱收敛到 $H(K)$，并且
$\bigcap F^pH(K) = 0$、$\bigcup F^pH(K) = H(K)$。
\end{enumerate}
\end{definition}

\noindent
遗憾的是，文献中似乎很难找到关于这些概念的一套一致术语。

\begin{lemma}
\label{homology-lemma-filtered-differential-ss-converges}
设 $\mathcal{A}$ 为阿贝尔范畴，并设 $(K, F, d)$ 为 $\mathcal{A}$ 的
滤过微分对象。关联谱序列
\begin{enumerate}
\item 弱收敛到 $H(K)$，当且仅当对每个 $p \in \mathbf{Z}$，
等式 (\ref{homology-equation-at-bottom}) 与 (\ref{homology-equation-on-top}) 中均取等号；
\item 汇聚到 $H(K)$，当且仅当它弱收敛到 $H(K)$，并且
$\bigcap_p (\Ker(d) \cap F^pK + \Im(d)) = \Im(d)$ 及
$\bigcup_p (\Ker(d) \cap F^pK + \Im(d)) = \Ker(d)$。
\end{enumerate}
\end{lemma}

\begin{proof}
由上面的讨论立即可得。
\end{proof}

\section{谱序列：滤过复形}
\label{homology-section-filtered-complex}

\begin{definition}
\label{homology-definition-filtered-complex}
设 $\mathcal{A}$ 为阿贝尔范畴。$\mathcal{A}$ 的一个
{\it 滤过复形 $K^\bullet$}，是 $\text{Fil}(\mathcal{A})$ 中的复形
（见定义 \ref{homology-definition-filtered}）。
\end{definition}

\noindent
我们用 $F$ 表示各对象上的滤过。因此，$F^pK^n$ 表示复形第 $n$ 项之
滤过的第 $p$ 步。注意，每个 $F^pK^\bullet$ 都是 $\mathcal{A}$ 中的
复形。因此，也可把滤过复形定义为 $\mathcal{A}$ 的复形所成的
（阿贝尔）范畴中的滤过对象。特别地，$\text{gr} K^\bullet$ 是
$\mathcal{A}$ 的复形范畴中的分次对象。

\medskip\noindent
为描述关联于这种对象的谱序列，暂时假设 $\mathcal{A}$ 是具有可数直和且
可数直和正合的阿贝尔范畴（这不是自动成立的；见注记
\ref{homology-remark-direct-sums-not-exact}）。用 $d$ 表示 $K$ 的微分。
忘掉分次后，可把 $\bigoplus K^n$ 看作 $\mathcal{A}$ 的滤过微分对象。
因此，按照第 \ref{homology-section-filtered-differential} 节，得到谱序列
$(E_r, d_r)_{r \geq 0}$。本节将求出该谱序列的各项，并用 $K$ 的分次
所给出的附加结构装备这些项。

\medskip\noindent
首先指出，$E_0^p = \text{gr}^p K^\bullet$ 是复形，因而是分次的。
所以 $E_0$ 自然地是双分次的。通常使用双分次
$$
E_0 = \bigoplus\nolimits_{p, q} E_0^{p, q},
\quad
E_0^{p, q} = \text{gr}^p K^{p + q}
$$
其思想是把 $p + q$ 看作（上）同调类的{\it 总次数}。此外，$p$ 称为
{\it 滤过次数}，$q$ 称为{\it 余次数}。微分 $d_0$ 按如下方式与此
双分次相容：
$$
d_0  = \bigoplus d_0^{p, q},
\quad
d_0^{p, q} : E_0^{p, q} \to E_0^{p, q + 1}.
$$
事实上，$d_0^p$ 正是复形 $\text{gr}^p K^\bullet$ 上的微分
（该复形以略作平移的方式作为 $\text{gr}^pE_0$ 出现）。

\medskip\noindent
为进一步讨论，把第 \ref{homology-section-filtered-differential} 节引入的对象
$B_r^p$ 与 $Z_r^p$ 视为分次对象，并求出微分的相应分解。
我们完全直接地进行，但再次警告读者，我们的记号与其他地方的记号不同。
定义
$$
Z_r^{p, q} =
\frac{F^pK^{p + q} \cap d^{-1}(F^{p + r}K^{p + q + 1}) + F^{p + 1}K^{p + q}}
{F^{p + 1}K^{p + q}}
$$
及
$$
B_r^{p, q} =
\frac{F^pK^{p + q} \cap d(F^{p - r + 1}K^{p + q - 1}) + F^{p + 1}K^{p + q}}
{F^{p + 1}K^{p + q}}
$$
当然，$E_r^{p, q} = Z_r^{p, q}/B_r^{p, q}$。按这些定义，显然有
$Z_r^p = \bigoplus_q Z_r^{p, q}$、
$B_r^p = \bigoplus_q B_r^{p, q}$ 及
$E_r^p = \bigoplus_q E_r^{p, q}$。而且
$$
0 \subset \ldots \subset B_r^{p, q} \subset
\ldots
\subset Z_r^{p, q} \subset \ldots \subset E_0^{p, q}
$$
映射 $d_r^p$ 也分解为映射
$$
d_r^{p, q} : E_r^{p, q} \longrightarrow E_r^{p + r, q - r + 1},
\quad
z + F^{p + 1}K^{p + q}
\mapsto
dz + F^{p + r + 1}K^{p + q + 1}
$$
的直和，其中
$z \in F^pK^{p + q} \cap d^{-1}(F^{p + r}K^{p + q + 1})$。

\begin{lemma}
\label{homology-lemma-spectral-sequence-filtered-complex}
设 $\mathcal{A}$ 为阿贝尔范畴，并设 $(K^\bullet, F)$ 为
$\mathcal{A}$ 的滤过复形。存在 $\mathcal{A}$ 的双分次对象范畴中关联于
$(K^\bullet, F)$ 的谱序列 $(E_r, d_r)_{r \geq 0}$，使 $d_r$ 的
双次数为 $(r, - r + 1)$，并且 $E_r$ 的双分次部分 $E_r^{p, q}$ 及
映射 $d_r^{p, q} : E_r^{p, q} \to E_r^{p + r, q - r + 1}$ 如上给出。
此外，有 $E_0^{p, q} = \text{gr}^p(K^{p + q})$、
$d_0^{p, q} = \text{gr}^p(d^{p + q})$，且
$E_1^{p, q} = H^{p + q}(\text{gr}^p(K^\bullet))$。
\end{lemma}

\begin{proof}
若 $\mathcal{A}$ 有可数直和且可数直和正合，则结论来自上面的讨论。
一般情形如下处理；我们强烈建议读者跳过这个证明。考虑
$\mathcal{A}$ 的双分次对象
$A = (F^{p + 1}K^{p + 1 + q})$，即把 $F^{p + 1}K^{p + 1 + q}$
置于双次数 $(p, q)$（编号中的奇特平移是为了使后面的编号正确）。
通过在每个分量上使用 $d$，为其配备微分 $d : A \to A[0, 1]$。
那么 $(A, d)$ 是微分双分次对象。考虑映射
$$
\alpha : A \to A[-1, 1]
$$
它在双次数 $(p, q)$ 处由嵌入
$F^{p + 1}K^{p + 1 + q} \to F^pK^{p + 1 + q}$ 给出。
这是微分对象之间的单射态射
$\alpha : (A, d) \to (A, d)[-1, 1]$。因此可在
$S = [0, 1]$、$T = [1, -1]$ 时应用注记
\ref{homology-remark-differential-object-selfmap}。所得双分次对象的谱序列
$(E_r, d_r)_{r \geq 0}$ 正是我们所求的谱序列。稍微展开这些定义。
首先有 $E_r = (E_r^{p, q})$。其次，由于
$T^rS = [r, -r + 1]$，有 $d_r : E_r \to E_r[r, -r + 1]$，
即 $d_r^{p, q} : E_r^{p, q} \to E_r^{p + r, q - r + 1}$。

\medskip\noindent
为看出分次部分的描述成立，按上面的方法论证。首先有
$$
E_0 = \Coker(\alpha : A \to A[-1, 1])[0, -1] =
\Coker(\alpha[0, -1] : A[0, -1] \to A[-1, 0])
$$
由上面的编号选择，这给出
$$
E_0^{p, q} = \Coker(F^{p + 1}K^{p + q} \to F^pK^{p + q}) = \text{gr}^pK^{p + q}
$$
第一个微分为
$d_0^{p, q} = \text{gr}^pd^{p + q} : E_0^{p, q} \to E_0^{p, q + 1}$。
其次，注记 \ref{homology-remark-differential-object-selfmap} 中对边界 $B_r$
与余闭元 $Z_r$ 的描述，直接转化为上面给出的
$Z_r^{p, q}$ 与 $B_r^{p, q}$ 的公式。
\end{proof}

\begin{lemma}
\label{homology-lemma-spectral-sequence-filtered-complex-d1}
设 $\mathcal{A}$ 为阿贝尔范畴，并设 $(K^\bullet, F)$ 为
$\mathcal{A}$ 的滤过复形。假设 $\mathcal{A}$ 有可数直和，并设
$(E_r, d_r)_{r \geq 0}$ 为关联于 $(K^\bullet, F)$ 的谱序列。
\begin{enumerate}
\item 映射
$$
d_1^{p, q} :
E_1^{p, q} = H^{p + q}(\text{gr}^p(K^\bullet))
\longrightarrow
E_1^{p + 1, q} = H^{p + q + 1}(\text{gr}^{p + 1}(K^\bullet))
$$
等于关联于复形短正合列
$$
0 \to \text{gr}^{p + 1}(K^\bullet) \to
F^pK^\bullet/F^{p + 2}K^\bullet \to \text{gr}^p(K^\bullet) \to 0.
$$
的上同调边界映射。
\item 假设对所有 $p \in \mathbf{Z}$ 均有
$d(F^pK) \subset F^{p + 1}K$。那么 $d$ 在
$\text{gr}^p(K^\bullet)$ 上诱导零微分，因而
$E_1^{p, q} = \text{gr}^p(K^\bullet)^{p + q}$。此外，此时
$$
d_1^{p, q} :
E_1^{p, q} = \text{gr}^p(K^\bullet)^{p + q}
\longrightarrow
E_1^{p + 1, q} = \text{gr}^{p + 1}(K^\bullet)^{p + q + 1}
$$
是由 $d$ 诱导的态射。
\end{enumerate}
\end{lemma}

\begin{proof}
这由引理 \ref{homology-lemma-spectral-sequence-filtered-complex} 前给出的
微分 $d_1^{p, q}$ 的公式显然可得。
\end{proof}

\begin{lemma}
\label{homology-lemma-spectral-sequence-filtered-complex-functorial}
设 $\mathcal{A}$ 为阿贝尔范畴，并设
$\alpha : (K^\bullet, F) \to (L^\bullet, F)$ 为 $\mathcal{A}$ 的
滤过复形之间的态射。设 $(E_r(K), d_r)_{r \geq 0}$（相应地，
$(E_r(L), d_r)_{r \geq 0}$）为关联于 $(K^\bullet, F)$（相应地，
$(L^\bullet, F)$）的谱序列。态射 $\alpha$ 诱导与双分次相容的典范
谱序列态射 $\{\alpha_r : E_r(K) \to E_r(L)\}_{r \geq 0}$。
\end{lemma}

\begin{proof}
由谱序列各项的显式表示显然可得。
\end{proof}

\begin{definition}
\label{homology-definition-filtration-cohomology-filtered-complex}
设 $\mathcal{A}$ 为阿贝尔范畴，并设 $(K^\bullet, F)$ 为
$\mathcal{A}$ 的滤过复形。$H^n(K^\bullet)$ 上的{\it 诱导滤过}定义为
$F^pH^n(K^\bullet) = \Im(H^n(F^pK^\bullet) \to H^n(K^\bullet))$。
\end{definition}

\noindent
展开其含义可见
\begin{equation}
\label{homology-equation-filtration-cohomology}
F^pH^n(K^\bullet, d) =
\frac{\Ker(d) \cap F^pK^n + \Im(d) \cap K^n}{\Im(d) \cap K^n}
\end{equation}
因而
\begin{equation}
\label{homology-equation-graded-cohomology}
\text{gr}^p H^n(K^\bullet) =
\frac{\Ker(d) \cap F^pK^n}{\Ker(d) \cap F^{p + 1}K^n + \Im(d) \cap F^pK^n}
\end{equation}
（略去一个中间步骤）。

\begin{lemma}
\label{homology-lemma-compute-cohomology-filtered-complex}
设 $\mathcal{A}$ 为阿贝尔范畴，并设 $(K^\bullet, F)$ 为
$\mathcal{A}$ 的滤过复形。若 $Z_\infty^{p, q}$ 与
$B_\infty^{p, q}$ 存在（见证明），则
\begin{enumerate}
\item 极限 $E_\infty$ 存在并且是双分次对象，其双次数 $(p, q)$ 部分为
$E_\infty^{p, q} = Z_\infty^{p, q}/B_\infty^{p, q}$；
\item $K^\bullet$ 的第 $n$ 个上同调对象之第 $p$ 个分次部分
$\text{gr}^pH^n(K^\bullet)$ 是 $E_\infty^{p, n - p}$ 的子商。
\end{enumerate}
\end{lemma}

\begin{proof}
定义 \ref{homology-definition-limit-spectral-sequence} 的对象
$Z_\infty$、$B_\infty$ 及极限 $E_\infty = Z_\infty/B_\infty$，
按照引理 \ref{homology-lemma-spectral-sequence-filtered-complex} 中对谱序列的
构造，是 $\mathcal{A}$ 的双分次对象。由于 $Z_r^{p, q}$（相应地，
$B_r^{p, q}$）是 $Z_r$（相应地，$B_r$）的双次数 $(p, q)$ 部分，
若假设
$$
Z_\infty^{p, q} = \bigcap\nolimits_r Z_r^{p, q} =
\bigcap\nolimits_r
\frac{F^pK^{p + q} \cap d^{-1}(F^{p + r}K^{p + q + 1}) + F^{p + 1}K^{p + q}}
{F^{p + 1}K^{p + q}}
$$
以及
$$
B_\infty^{p, q} = \bigcup\nolimits_r B_r^{p, q} =
\bigcup\nolimits_r
\frac{F^pK^{p + q} \cap d(F^{p - r + 1}K^{p + q - 1}) + F^{p + 1}K^{p + q}}
{F^{p + 1}K^{p + q}}
$$
存在，则 $Z_\infty$ 与 $B_\infty$ 存在，且双次数 $(p, q)$ 部分分别为
$Z_\infty^{p, q}$ 与 $B_\infty^{p, q}$（这来自关于双分次对象之并与交的
初等论证）。因此
$$
E_\infty^{p, q} =
\frac{\bigcap_r (F^pK^{p + q} \cap d^{-1}(F^{p + r}K^{p + q + 1})
+ F^{p + 1}K^{p + q})}
{\bigcup_r (F^pK^{p + q} \cap d(F^{p - r + 1}K^{p + q - 1})
+ F^{p + 1}K^{p + q})}.
$$
其中分子与分母都存在。置 $n = p + q$，有
\begin{equation}
\label{homology-equation-on-top-bigraded}
\Ker(d) \cap F^pK^{n} + F^{p + 1}K^{n}
\subset
\bigcap\nolimits_r
\left(
F^pK^{n} \cap d^{-1}(F^{p + r}K^{n + 1}) + F^{p + 1}K^{n}
\right)
\end{equation}
以及
\begin{equation}
\label{homology-equation-at-bottom-bigraded}
\bigcup\nolimits_r
\left(
F^pK^{n} \cap d(F^{p - r + 1}K^{n - 1}) + F^{p + 1}K^{n}
\right)
\subset
\Im(d) \cap F^pK^{n} + F^{p + 1}K^{n}.
\end{equation}
因此，$E_\infty^{p, q}$ 的一个子商为
$$
\frac{\Ker(d) \cap F^pK^{n} + F^{p + 1}K^n}
{\Im(d) \cap F^pK^{n} + F^{p + 1}K^{n}} =
\frac{\Ker(d) \cap F^pK^n}{\Im(d) \cap F^pK^n + \Ker(d) \cap F^{p + 1}K^n}
$$
与 (\ref{homology-equation-graded-cohomology}) 比较即得结论。
\end{proof}

\begin{definition}
\label{homology-definition-bounded-ss}
设 $\mathcal{A}$ 为阿贝尔范畴，并设 $(E_r, d_r)_{r \geq r_0}$ 为
$\mathcal{A}$ 的双分次对象的谱序列，其中 $d_r$ 的双次数为
$(r, -r + 1)$。称这样的谱序列
\begin{enumerate}
\item {\it 正则}，若对所有 $p, q \in \mathbf{Z}$，存在
$b = b(p, q)$，使得当 $r \geq b$ 时映射
$d_r^{p, q} : E_r^{p, q} \to E_r^{p + r, q - r + 1}$ 均为零；
\item {\it 余正则}，若对所有 $p, q \in \mathbf{Z}$，存在
$b = b(p, q)$，使得当 $r \geq b$ 时映射
$d_r^{p - r, q + r - 1} : E_r^{p - r, q + r - 1} \to E_r^{p, q}$
均为零；
\item {\it 有界}，若对所有 $n$，只有有限多个
$E_{r_0}^{p, n - p}$ 非零；
\item {\it 下有界}，若对所有 $n$，存在 $b = b(n)$，使得当
$p \geq b$ 时 $E_{r_0}^{p, n - p} = 0$；
\item {\it 上有界}，若对所有 $n$，存在 $b = b(n)$，使得当
$p \leq b$ 时 $E_{r_0}^{p, n - p} = 0$。
\end{enumerate}
\end{definition}

\noindent
下有界意味着：观察直线 $p + q = n$（其斜率为 $-1$）上的
$E_r^{p, q}$ 时，当 $(p, q)$ 向右下移动便得到零。如上所述，
文献中关于这些概念没有一致术语。

\begin{lemma}
\label{homology-lemma-relate-boundedness}
在定义 \ref{homology-definition-bounded-ss} 的情形中，令
$Z_r^{p, q}, B_r^{p, q} \subset E_{r_0}^{p, q}$ 为按第
\ref{homology-section-spectral-sequence} 节定义的 $Z_r, B_r$ 的
$(p, q)$-分次部分。
\begin{enumerate}
\item 谱序列正则，当且仅当对所有 $p, q$，存在 $r = r(p, q)$ 使得
$Z_r^{p, q} = Z_{r + 1}^{p, q} = \ldots$
\item 谱序列余正则，当且仅当对所有 $p, q$，存在 $r = r(p, q)$ 使得
$B_r^{p, q} = B_{r + 1}^{p, q} = \ldots$
\item 谱序列有界，当且仅当它既下有界又上有界。
\item 若谱序列下有界，则它正则。
\item 若谱序列上有界，则它余正则。
\end{enumerate}
\end{lemma}

\begin{proof}
略。提示：若 $E_r^{p, q} = 0$，则对所有 $r' \geq r$ 均有
$E_{r'}^{p, q} = 0$。
\end{proof}

\begin{definition}
\label{homology-definition-filtered-complex-ss-converges}
设 $\mathcal{A}$ 为阿贝尔范畴，并设 $(K^\bullet, F)$ 为
$\mathcal{A}$ 的滤过复形。称关联于 $(K^\bullet, F)$ 的谱序列
\begin{enumerate}
\item {\it 弱收敛到 $H^*(K^\bullet)$}，若通过引理
\ref{homology-lemma-compute-cohomology-filtered-complex}，对所有
$p, n \in \mathbf{Z}$ 均有
$\text{gr}^pH^n(K^\bullet) = E_{\infty}^{p, n - p}$；
\item {\it 汇聚到 $H^*(K^\bullet)$}，若它弱收敛到
$H^*(K^\bullet)$，并且对所有 $n$ 均有
$\bigcap_p F^pH^n(K^\bullet) = 0$ 及
$\bigcup_p F^p H^n(K^\bullet) = H^n(K^\bullet)$；
\item {\it 收敛到 $H^*(K^\bullet)$}，若它正则、汇聚到
$H^*(K^\bullet)$，并且
$H^n(K^\bullet) = \lim_p H^n(K^\bullet)/F^pH^n(K^\bullet)$。
\end{enumerate}
\end{definition}

\noindent
弱收敛、汇聚或收敛都用记号
$E_r^{p, q} \Rightarrow H^{p + q}(K^\bullet)$ 表示。如上所述，
文献中关于这些概念没有一致术语。

\begin{lemma}
\label{homology-lemma-filtered-complex-ss-converges}
设 $\mathcal{A}$ 为阿贝尔范畴，并设 $(K^\bullet, F)$ 为
$\mathcal{A}$ 的滤过复形。关联谱序列
\begin{enumerate}
\item 弱收敛到 $H^*(K^\bullet)$，当且仅当对每个
$p, q \in \mathbf{Z}$，等式 (\ref{homology-equation-at-bottom-bigraded}) 与
(\ref{homology-equation-on-top-bigraded}) 中均取等号；
\item 汇聚到 $H^*(K)$，当且仅当它弱收敛到 $H^*(K^\bullet)$，并且有
$\bigcap_p (\Ker(d) \cap F^pK^n + \Im(d) \cap K^n) = \Im(d) \cap K^n$
及
$\bigcup_p (\Ker(d) \cap F^pK^n + \Im(d) \cap K^n) = \Ker(d) \cap K^n$。
\end{enumerate}
\end{lemma}

\begin{proof}
由上面的讨论立即可得。
\end{proof}

\begin{lemma}
\label{homology-lemma-biregular-ss-converges}
设 $\mathcal{A}$ 为阿贝尔范畴，并设 $(K^\bullet, F)$ 为
$\mathcal{A}$ 的滤过复形。假设每个 $K^n$ 上的滤过有限
（见定义 \ref{homology-definition-filtered}）。那么
\begin{enumerate}
\item 关联于 $(K^\bullet, F)$ 的谱序列有界；
\item 每个 $H^n(K^\bullet)$ 上的滤过有限；
\item 关联于 $(K^\bullet, F)$ 的谱序列收敛到 $H^*(K^\bullet)$；
\item 若 $\mathcal{C} \subset \mathcal{A}$ 是弱 Serre 子范畴，并且对
某个 $r$，所有 $p, q \in \mathbf{Z}$ 都有
$E_r^{p, q} \in \mathcal{C}$，则 $H^n(K^\bullet)$ 属于
$\mathcal{C}$。
\end{enumerate}
\end{lemma}

\begin{proof}
(1) 来自 $E_0^{p, n - p} = \text{gr}^p K^n$。(2) 由等式
(\ref{homology-equation-filtration-cohomology}) 显然可得。我们用引理
\ref{homology-lemma-filtered-complex-ss-converges} 证明谱序列弱收敛。
固定 $p, n \in \mathbf{Z}$。等式 (\ref{homology-equation-on-top-bigraded})
的右端等于 $F^pK^n \cap \Ker(d) + F^{p + 1}K^n$，因为当
$r \gg 0$ 时 $F^{p + r}K^n = 0$。因此
(\ref{homology-equation-on-top-bigraded}) 是等式。等式
(\ref{homology-equation-at-bottom-bigraded}) 的左端等于
$F^pK^n \cap \Im(d) + F^{p + 1}K^n$，因为当 $r \gg 0$ 时
$F^{p - r + 1}K^{n - 1} = K^{n - 1}$。因此
(\ref{homology-equation-at-bottom-bigraded}) 是等式。由 (2)，
$H^n(K^\bullet)$ 上的滤过有限，所以得到汇聚。为证明收敛，还须证明谱序列
正则；这来自其有界性（引理 \ref{homology-lemma-relate-boundedness}）。还须证明
$H^n(K^\bullet) = \lim_p H^n(K^\bullet)/F^pH^n(K^\bullet)$，
这来自 (2) 中已证的 $H^*(K^\bullet)$ 上的滤过有限。

\medskip\noindent
证明 (4)。假设对某个 $r \geq 0$ 及 $\mathcal{A}$ 的某个弱 Serre
子范畴 $\mathcal{C}$，有 $E_r^{p, q} \in \mathcal{C}$。那么
$E_{r + 1}^{p, q}$ 也属于 $\mathcal{C}$；见引理
\ref{homology-lemma-characterize-weak-serre-subcategory}。由上面证明的有界性
（由引理 \ref{homology-lemma-relate-boundedness}，它蕴含谱序列既正则又余正则），
可找到 $r' \geq r$，使得对所有满足 $p + q = n$ 的 $p, q$，均有
$E_\infty^{p, q} = E_{r'}^{p, q}$。于是 $H^n(K^\bullet)$ 是
$\mathcal{A}$ 中具有有限滤过且各分次部分都属于 $\mathcal{C}$ 的对象。
由引理 \ref{homology-lemma-characterize-weak-serre-subcategory}，这蕴含
$H^n(K^\bullet)$ 属于 $\mathcal{C}$。
\end{proof}

\begin{lemma}
\label{homology-lemma-biregular-ss-relation-in-K0}
设 $\mathcal{A}$ 为阿贝尔范畴，并设 $(K^\bullet, F)$ 为
$\mathcal{A}$ 的滤过复形。假设每个 $K^n$ 上的滤过有限
（见定义 \ref{homology-definition-filtered}），并且对某个 $r$，只有有限多个
$E_r^{p, q}$ 非零。那么只有有限多个 $H^n(K^\bullet)$ 非零，并且
$$
\sum (-1)^n[H^n(K^\bullet)] = \sum (-1)^{p + q} [E_r^{p, q}]
$$
在 $K_0(\mathcal{A}')$ 中成立，其中 $\mathcal{A}'$ 是包含各对象
$E_r^{p, q}$ 的 $\mathcal{A}$ 的最小弱 Serre 子范畴。
\end{lemma}

\begin{proof}
用 $E_r^{even}$ 与 $E_r^{odd}$ 表示 $E_r$ 的偶部分与奇部分，它们定义为
所有 $p + q$ 分别为偶数与奇数的 $(p, q)$ 分量之直和。微分 $d_r$
定义映射 $\varphi : E_r^{even} \to E_r^{odd}$ 与
$\psi : E_r^{odd} \to E_r^{even}$，其两个方向的复合都为零。
于是
\begin{align*}
[E_r^{even}] - [E_r^{odd}] & =
[\Ker(\varphi)] + [\Im(\varphi)] - [\Ker(\psi)] - [\Im(\psi)] \\
& =
[\Ker(\varphi)/\Im(\psi)] - [\Ker(\psi)/\Im(\varphi)] \\
& =
[E_{r + 1}^{even}] - [E_{r + 1}^{odd}]
\end{align*}
注意所有中间对象都属于包含各对象 $E_r^{p, q}$ 的最小 Serre 子范畴。
继续这一过程，可任意增大 $r$。由于只有有限多对 $(p, q)$ 使
$E_r^{p, q}$ 非零，而这一性质由 $E_{r + 1}, E_{r + 2}, \ldots$
继承，故可假设 $d_r = 0$。此时，$H^n(K^\bullet)$ 具有有限滤过
（引理 \ref{homology-lemma-biregular-ss-converges}），其分次部分恰为
$E_r^{p, n - p}$，结论显然。
\end{proof}

\noindent
下面的引理更多是对我们定义的一种合理性检查。显然，若有一个滤过复形，
使得对每个 $n$ 均有
$$
H^n(F^pK^\bullet) = 0\text{ 对 }p \gg 0
\quad\text{且}\quad
H^n(F^pK^\bullet) = H^n(K^\bullet)\text{ 对 }p \ll 0,
$$
那么相应谱序列应当收敛，对吗？

\begin{lemma}
\label{homology-lemma-ss-converges-trivial}
设 $\mathcal{A}$ 为阿贝尔范畴，并设 $(K^\bullet, F)$ 为
$\mathcal{A}$ 的滤过复形。假设
\begin{enumerate}
\item 对每个 $n$，存在 $p_0(n)$，使得当 $p \geq p_0(n)$ 时
$H^n(F^pK^\bullet) = 0$；
\item 对每个 $n$，存在 $p_1(n)$，使得当 $p \leq p_1(n)$ 时
$H^n(F^pK^\bullet) \to H^n(K^\bullet)$ 是同构。
\end{enumerate}
那么
\begin{enumerate}
\item 关联于 $(K^\bullet, F)$ 的谱序列有界；
\item 每个 $H^n(K^\bullet)$ 上的滤过有限；
\item 关联于 $(K^\bullet, F)$ 的谱序列收敛到 $H^*(K^\bullet)$。
\end{enumerate}
\end{lemma}

\begin{proof}
固定 $n$。使用关联于复形短正合列
$$
0 \to F^{p + 1}K^\bullet \to F^pK^\bullet \to \text{gr}^pK^\bullet \to 0
$$
的长正合上同调列，可得当
$p \geq \max(p_0(n), p_0(n + 1))$ 以及
$p < \min(p_1(n), p_1(n + 1))$ 时，$E_1^{p, n - p} = 0$。
因此谱序列有界（定义 \ref{homology-definition-bounded-ss}）。这证明 (1)。

\medskip\noindent
由假设与定义 \ref{homology-definition-filtration-cohomology-filtered-complex}，
$H^n(K^\bullet)$ 上的滤过显然有限。这证明 (2)。

\medskip\noindent
接着用引理 \ref{homology-lemma-filtered-complex-ss-converges} 证明谱序列弱收敛到
$H^*(K^\bullet)$。先证明 (\ref{homology-equation-on-top-bigraded}) 中取等号。
事实上，当 $p + r > p_0(n + 1)$ 时，映射
$$
d : F^pK^{n} \cap d^{-1}(F^{p + r}K^{n + 1}) \to F^{p + r}K^{n + 1}
$$
的像落在 $d : F^{p + r}K^n \to F^{p + r}K^{n + 1}$ 的像中，
因为复形 $F^{p + r}K^\bullet$ 在次数 $n + 1$ 处正合。
由此
$F^pK^{n} \cap d^{-1}(F^{p + r}K^{n + 1}) =
d(F^{p + r}K^n) + \Ker(d) \cap F^pK^n$。因此对这样的 $r$ 有
$$
\Ker(d) \cap F^pK^{n} + F^{p + 1}K^{n} =
F^pK^{n} \cap d^{-1}(F^{p + r}K^{n + 1}) + F^{p + 1}K^{n}
$$
这证明了所需等式。为证明 (\ref{homology-equation-at-bottom-bigraded}) 中取等号，
使用如下事实：当 $p - r + 1 < p_1(n - 1)$ 时有
$$
d(F^{p - r + 1}K^{n - 1}) = \Im(d) \cap F^{p - r + 1}K^n
$$
因为映射 $F^{p - r + 1}K^\bullet \to K^\bullet$ 在次数 $n - 1$
的上同调上诱导同构。这说明对这样的 $r$ 有
$$
F^pK^{n} \cap d(F^{p - r + 1}K^{n - 1}) + F^{p + 1}K^{n} =
\Im(d) \cap F^pK^{n} + F^{p + 1}K^{n}
$$
从而证明所需等式。

\medskip\noindent
为由引理 \ref{homology-lemma-filtered-complex-ss-converges} 看出谱序列汇聚到
$H^*(K^\bullet)$，须证明
$\bigcap_p (\Ker(d) \cap F^pK^n + \Im(d) \cap K^n) = \Im(d) \cap K^n$
及
$\bigcup_p (\Ker(d) \cap F^pK^n + \Im(d) \cap K^n) = \Ker(d) \cap K^n$。
当 $p \geq p_0(n)$ 时有
$\Ker(d) \cap F^pK^n + \Im(d) \cap K^n = \Im(d) \cap K^n$；
当 $p \leq p_1(n)$ 时有
$\Ker(d) \cap F^pK^n + \Im(d) \cap K^n = \Ker(d) \cap K^n$。
结合弱收敛、汇聚与有界性，可见 (2) 与 (3) 成立。
\end{proof}

\section{谱序列：双复形}
\label{homology-section-double-complex}

\noindent
设 $K^{\bullet, \bullet}$ 为双复形；见第
\ref{homology-section-double-complexes} 节。通常用
$H^p_I(K^{\bullet, \bullet})$ 表示如下复形：其各项为
$\Ker(d_1^{p, q})/\Im(d_1^{p - 1, q})$（$q$ 变化），微分由 $d_2$
诱导。于是 $H^q_{II}(H^p_I(K^{\bullet, \bullet}))$ 表示它在次数
$q$ 处的上同调。通常还用 $H^q_{II}(K^{\bullet, \bullet})$ 表示如下
复形：其各项为 $\Ker(d_2^{p, q})/\Im(d_2^{p, q - 1})$（$p$ 变化），
微分由 $d_1$ 诱导。于是
$H^p_I(H^q_{II}(K^{\bullet, \bullet}))$ 表示它在次数 $p$ 处的上同调。
后面将会看到，这些上同调群作为关联总复形或单复形上的某个滤过之谱序列
的各项出现；见定义 \ref{homology-definition-associated-simple-complex}。

\medskip\noindent
关联于双复形 $K^{\bullet, \bullet}$ 的总复形
$\text{Tot}(K^{\bullet, \bullet})$ 上有两个自然滤过。即定义
$$
F_I^p(\text{Tot}^n(K^{\bullet, \bullet})) =
\bigoplus\nolimits_{i + j = n, \ i \geq p} K^{i, j}
\quad
\text{且}
\quad
F_{II}^p(\text{Tot}^n(K^{\bullet, \bullet})) =
\bigoplus\nolimits_{i + j = n, \ j \geq p} K^{i, j}.
$$
立即可验证 $(\text{Tot}(K^{\bullet, \bullet}), F_I)$ 与
$(\text{Tot}(K^{\bullet, \bullet}), F_{II})$ 是滤过复形。
由第 \ref{homology-section-filtered-complex} 节，得到两个谱序列。通常用
$({}'E_r, {}'d_r)_{r \geq 0}$ 表示关联于滤过 $F_I$ 的谱序列，
用 $({}''E_r, {}''d_r)_{r \geq 0}$ 表示关联于滤过 $F_{II}$ 的谱序列。
下面描述这些谱序列。

\begin{lemma}
\label{homology-lemma-ss-double-complex}
设 $\mathcal{A}$ 为阿贝尔范畴，并设 $K^{\bullet, \bullet}$ 为双复形。
关联于 $K^{\bullet, \bullet}$ 的谱序列具有下列各项：
\begin{enumerate}
\item ${}'E_0^{p, q} = K^{p, q}$，其中
${}'d_0^{p, q} = (-1)^p d_2^{p, q} : K^{p, q} \to K^{p, q + 1}$；
\item ${}''E_0^{p, q} = K^{q, p}$，其中
${}''d_0^{p, q} = d_1^{q, p} : K^{q, p} \to K^{q + 1, p}$；
\item ${}'E_1^{p, q} = H^q(K^{p, \bullet})$，其中
${}'d_1^{p, q} = H^q(d_1^{p, \bullet})$；
\item ${}''E_1^{p, q} = H^q(K^{\bullet, p})$，其中
${}''d_1^{p, q} = (-1)^q H^q(d_2^{\bullet, p})$；
\item ${}'E_2^{p, q} = H^p_I(H^q_{II}(K^{\bullet, \bullet}))$；
\item ${}''E_2^{p, q} = H^p_{II}(H^q_I(K^{\bullet, \bullet}))$。
\end{enumerate}
\end{lemma}

\begin{proof}
略。
\end{proof}

\noindent
这些谱序列在 $H^n(\text{Tot}(K^{\bullet, \bullet}))$ 上定义两个滤过，
我们将其记作 $F_I$ 与 $F_{II}$。

\begin{definition}
\label{homology-definition-ss-double-complex-converge}
设 $\mathcal{A}$ 为阿贝尔范畴，并设 $K^{\bullet, \bullet}$ 为双复形。
若定义 \ref{homology-definition-filtered-complex-ss-converges} 适用，则称谱序列
$({}'E_r, {}'d_r)_{r \geq 0}$ {\it 弱收敛到
$H^n(\text{Tot}(K^{\bullet, \bullet}))$}、{\it 汇聚到
$H^n(\text{Tot}(K^{\bullet, \bullet}))$} 或{\it 收敛到
$H^n(\text{Tot}(K^{\bullet, \bullet}))$}。类似地，若定义
\ref{homology-definition-filtered-complex-ss-converges} 适用，则称谱序列
$({}''E_r, {}''d_r)_{r \geq 0}$ {\it 弱收敛到
$H^n(\text{Tot}(K^{\bullet, \bullet}))$}、{\it 汇聚到
$H^n(\text{Tot}(K^{\bullet, \bullet}))$} 或{\it 收敛到
$H^n(\text{Tot}(K^{\bullet, \bullet}))$}。
\end{definition}

\noindent
如上所述，文献中关于这些概念没有一致术语。在定义的情形中，若对所有
$n$，通过引理 \ref{homology-lemma-compute-cohomology-filtered-complex} 的典范比较有
$$
\text{gr}_{F_I}(H^n(\text{Tot}(K^{\bullet, \bullet}))) =
\oplus_{p + q = n} {}'E_\infty^{p, q}
$$
则第一个谱序列弱收敛。类似地，若对所有 $n$，通过引理
\ref{homology-lemma-compute-cohomology-filtered-complex} 的典范比较有
$$
\text{gr}_{F_{II}}(H^n(\text{Tot}(K^{\bullet, \bullet}))) =
\oplus_{p + q = n} {}''E_\infty^{p, q}
$$
则第二个谱序列 $({}''E_r, {}''d_r)_{r \geq 0}$ 弱收敛。

\begin{lemma}
\label{homology-lemma-first-quadrant-ss}
设 $\mathcal{A}$ 为阿贝尔范畴，并设 $K^{\bullet, \bullet}$ 为双复形。
假设对每个 $n \in \mathbf{Z}$，满足 $p + q = n$ 的非零
$K^{p, q}$ 只有有限多个。那么
\begin{enumerate}
\item 关联于 $K^{\bullet, \bullet}$ 的两个谱序列有界；
\item 每个 $H^n(\text{Tot}(K^{\bullet, \bullet}))$ 上的滤过
$F_I$、$F_{II}$ 有限；
\item 谱序列 $({}'E_r, {}'d_r)_{r \geq 0}$ 与
$({}''E_r, {}''d_r)_{r \geq 0}$ 收敛到
$H^*(\text{Tot}(K^{\bullet, \bullet}))$；
\item 若 $\mathcal{C} \subset \mathcal{A}$ 是弱 Serre 子范畴，并且对
某个 $r$，所有 $p, q \in \mathbf{Z}$ 都有
${}'E_r^{p, q} \in \mathcal{C}$，则
$H^n(\text{Tot}(K^{\bullet, \bullet}))$ 属于 $\mathcal{C}$。
对 $({}''E_r, {}''d_r)_{r \geq 0}$ 亦然。
\end{enumerate}
\end{lemma}

\begin{proof}
立即来自引理 \ref{homology-lemma-biregular-ss-converges}。
\end{proof}

\noindent
下面给出谱序列的第一个应用。

\begin{lemma}
\label{homology-lemma-double-complex-gives-resolution}
设 $\mathcal{A}$ 为阿贝尔范畴，$K^\bullet$ 为复形，
$A^{\bullet, \bullet}$ 为双复形，并设
$\alpha^p : K^p \to A^{p, 0}$ 为态射。假设
\begin{enumerate}
\item 对每个 $n \in \mathbf{Z}$，满足 $p + q = n$ 的非零
$A^{p, q}$ 只有有限多个。
\item 若 $q < 0$，则 $A^{p, q} = 0$。
\item 态射 $\alpha^p$ 给出复形态射
$\alpha : K^\bullet \to A^{\bullet, 0}$。
\item 复形 $A^{p, \bullet}$ 在所有次数 $q \not = 0$ 处正合，且态射
$K^p \to A^{p, 0}$ 诱导同构 $K^p \to \Ker(d_2^{p, 0})$。
\end{enumerate}
那么 $\alpha$ 诱导复形的拟同构
$$
K^\bullet \longrightarrow \text{Tot}(A^{\bullet, \bullet})
$$
此外，还有一个以第二个变量 $q$ 代替 $p$ 的引理变体。
\end{lemma}

\begin{proof}
该映射就是由态射
$K^n \to A^{n, 0} \to \text{Tot}^n(A^{\bullet, \bullet})$ 给出的映射；
容易看出它们定义复形态射。考虑关联于双复形
$A^{\bullet, \bullet}$ 的谱序列
$({}'E_r, {}'d_r)_{r \geq 0}$。由引理 \ref{homology-lemma-first-quadrant-ss}，
该谱序列收敛，并且每个 $n$ 上
$H^n(\text{Tot}(A^{\bullet, \bullet}))$ 的诱导滤过有限。
由引理 \ref{homology-lemma-ss-double-complex} 及假设 (4)，除 $q = 0$ 外有
${}'E_1^{p, q} = 0$，并且 $'E_1^{p, 0} = K^p$，其微分
${}'d_1^{p, 0}$ 等同于 $d_K^p$。因此
${}'E_2^{p, 0} = H^p(K^\bullet)$，其余均为零。由次数理由，这显然蕴含
${}'d_2^{p, q} = {}'d_3^{p, q} = \ldots = 0$。故
$H^n(\text{Tot}(A^{\bullet, \bullet})) = H^n(K^\bullet)$。
我们略去如下验证：这一等同由上面引入的复形态射
$K^\bullet \to \text{Tot}(A^{\bullet, \bullet})$ 给出。
\end{proof}

\begin{lemma}
\label{homology-lemma-homotopy-complex-complexes}
设 $\mathcal{A}$ 为阿贝尔范畴，$M^\bullet$ 为 $\mathcal{A}$ 的复形，
并设
$$
a :
M^\bullet[0]
\longrightarrow
\left(A^{0, \bullet} \to A^{1, \bullet} \to A^{2, \bullet} \to \ldots \right)
$$
是 $\mathcal{A}$ 的复形组成的复形之范畴中的同伦等价。那么由
$M^\bullet \to A^{0, \bullet}$ 诱导的映射
$\alpha : M^\bullet \to \text{Tot}(A^{\bullet, \bullet})$ 是同伦等价。
\end{lemma}

\begin{proof}
该陈述是有意义的，因为复形组成的复形与双复形是同一回事。假设意味着存在映射
$$
b :
\left(A^{0, \bullet} \to A^{1, \bullet} \to A^{2, \bullet} \to \ldots \right)
\longrightarrow
M^\bullet[0]
$$
使得在复形组成的复形之范畴中，$a \circ b$ 与 $b \circ a$ 都同伦于
恒等映射。这意味着 $b \circ a$ 是 $M^\bullet[0]$ 上的恒等映射
（因为次数 $0$ 中只有一项）。另注意，$b$ 由映射
$b^0 : A^{0, \bullet} \to M^\bullet$ 给出，而在所有其他次数中为零。
因此 $b$ 诱导映射
$\beta : \text{Tot}(A^{\bullet, \bullet}) \to M^\bullet$，且
$\beta \circ \alpha$ 是 $M^\bullet$ 上的恒等映射。最后须证明
$\alpha \circ \beta$ 同伦于恒等映射。为此，选取复形映射
$h^n : A^{n, \bullet} \to A^{n - 1, \bullet}$，使得
$a \circ b - \text{id} = d_1 \circ h + h \circ d_1$；这些映射由假设存在。
这里 $d_1 : A^{n, \bullet} \to A^{n + 1, \bullet}$ 是复形组成的复形之
微分。对所有 $n$，还用 $d_2$ 表示复形 $A^{n, \bullet}$ 的微分。
令 $h^{n, m} : A^{n, m} \to A^{n - 1, m}$ 为 $h^n$ 的分量。
于是可考虑
$$
h' : \text{Tot}(A^{\bullet, \bullet})^k =
\bigoplus\nolimits_{n + m = k} A^{n, m}
\to
\bigoplus\nolimits_{n + m = k - 1} A^{n, m} =
\text{Tot}(A^{\bullet, \bullet})^{k - 1}
$$
它在直和项 $A^{n, m}$ 上由 $h^{n, m}$ 给出。计算可得，映射
$$
d_{\text{Tot}(A^{\bullet, \bullet})} \circ h' +
h' \circ d_{\text{Tot}(A^{\bullet, \bullet})}
$$
限制在直和项 $A^{n, m}$ 上等于
$$
d_1^{n - 1, m} \circ h^{n, m} +
(-1)^{n - 1} d_2^{n - 1, m} \circ h^{n, m} +
h^{n + 1, m} \circ d_1^{n, m} + h^{n, m + 1} \circ (-1)^nd_2^{n, m}
$$
由于 $h^n$ 是复形映射，项
$(-1)^{n - 1} d_2^{n - 1, m} \circ h^{n, m}$ 与
$h^{n, m + 1} \circ (-1)^nd_2^{n, m}$ 相消。另两项给出
$(\alpha \circ \beta)|_{A^{n, m}} - \text{id}_{A^{n, m}}$，因为
$a \circ b - \text{id} = d_1 \circ h + h \circ d_1$。证明完成。
\end{proof}

\section{阿贝尔群的双复形}
\label{homology-section-double-complexes-abelian-groups}

\noindent
本节给出关于阿贝尔群双复形的一些结果；对于一般阿贝尔范畴，我们（尚）没有
相应结果。请注意，只有当底层阿贝尔范畴是阿贝尔群范畴或类似范畴
（例如某个环上的模范畴）时，才能使用这些引理。若不在餐巾上画出“之字形”，
某些论证将难以理解——可与代数，引理
\ref{algebra-lemma-no-spectral-sequence} 的证明比较。

\begin{lemma}
\label{homology-lemma-right-resolution-gives-qis}
设 $M^\bullet$ 为阿贝尔群的复形，并设
$$
0 \to M^\bullet \to A_0^\bullet \to A_1^\bullet \to A_2^\bullet \to \ldots
$$
为阿贝尔群复形组成的正合复形。置 $A^{p, q} = A_p^q$，从而得到双复形。
那么，由 $M^\bullet \to A_0^\bullet$ 诱导的映射
$M^\bullet \to \text{Tot}(A^{\bullet, \bullet})$ 是拟同构。
\end{lemma}

\begin{proof}
若存在 $t \in \mathbf{Z}$，使得当 $q < t$ 时 $A_0^q = 0$，则结论立即
来自引理 \ref{homology-lemma-double-complex-gives-resolution}（按该引理最后的
陈述交换 $p$ 与 $q$）。好，不过对每个 $t \in \mathbf{Z}$，都有朴素截断
组成的复形
$$
0 \to
\sigma_{\geq t}M^\bullet \to
\sigma_{\geq t}A_0^\bullet \to
\sigma_{\geq t}A_1^\bullet \to
\sigma_{\geq t}A_2^\bullet \to \ldots
$$
用 $A(t)^{\bullet, \bullet}$ 表示相应双复形。
$H^n(\text{Tot}(A^{\bullet, \bullet}))$ 的每个元素 $\xi$，都是某个 $t$
下 $H^n(\text{Tot}(A(t)^{\bullet, \bullet}))$ 中元素的像
（考察上同调类的显式代表元）。因此 $\xi$ 属于
$H^n(\sigma_{\geq t}M^\bullet)$ 的像。故映射
$H^n(M^\bullet) \to H^n(\text{Tot}(A^{\bullet, \bullet}))$ 满射。
它也是单射，因为对所有 $t$，映射
$H^n(\sigma_{\geq t}M^\bullet) \to
H^n(\text{Tot}(A(t)^{\bullet, \bullet}))$ 都是单射，且可作类似论证。
\end{proof}

\begin{lemma}
\label{homology-lemma-good-resolution-gives-qis}
设 $M^\bullet$ 为阿贝尔群的复形，并设
$$
\ldots \to A_2^\bullet \to A_1^\bullet \to A_0^\bullet \to M^\bullet \to 0
$$
为阿贝尔群复形组成的正合复形，并且对所有 $p \in \mathbf{Z}$，复形
$$
\ldots \to \Ker(d_{A_2^\bullet}^p) \to \Ker(d_{A_1^\bullet}^p)
\to \Ker(d_{A_0^\bullet}^p) \to \Ker(d_{M^\bullet}^p) \to 0
$$
也正合。置 $A^{p, q} = A_{-p}^q$，从而得到双复形。那么，由
$A_0^\bullet \to M^\bullet$ 诱导的映射
$\text{Tot}(A^{\bullet, \bullet}) \to M^\bullet$ 是拟同构。
\end{lemma}

\begin{proof}
使用短正合列
$0 \to \Ker(d^p_{A_n^\bullet}) \to A_n^p \to \Im(d^p_{A_n^\bullet}) \to 0$
及假设，可见对所有 $p \in \mathbf{Z}$，
$$
\ldots \to \Im(d_{A_2^\bullet}^p) \to \Im(d_{A_1^\bullet}^p)
\to \Im(d_{A_0^\bullet}^p) \to \Im(d_{M^\bullet}^p) \to 0
$$
正合。再对正合列
$0 \to \Im(d^{p - 1}_{A_n^\bullet}) \to \Ker(d^p_{A_n^\bullet})
\to H^p(A_n^\bullet) \to 0$ 重复这一论证，得到对所有
$p \in \mathbf{Z}$，
$$
\ldots \to H^p(A_2^\bullet) \to H^p(A_1^\bullet)
\to H^p(A_0^\bullet) \to H^p(M^\bullet) \to 0
$$
正合。

\medskip\noindent
写 $T^\bullet = \text{Tot}(A^{\bullet, \bullet})$。我们将证明
$H^0(T^\bullet) \to H^0(M^\bullet)$ 是同构。相同论证适用于其他次数。
令 $x \in \Ker(\text{d}_{T^\bullet}^0)$ 代表元素
$\xi \in H^0(T^\bullet)$。写 $x = \sum_{i = n, \ldots, 0} x_i$，
其中 $x_i \in A_i^i$。假设 $n > 0$。那么 $x_n$ 属于
$d_{A_n^\bullet}^n$ 的核，并且映到 $H^n(A_{n - 1}^\bullet)$ 中的零，
因为它映到的元素在相差一个符号后是 $x_{n - 1}$ 的边界。
由证明的第一段，
$x_n \bmod \Im(d^{n - 1}_{A_n^\bullet})$ 属于
$H^n(A_{n + 1}^\bullet) \to H^n(A_n^\bullet)$ 的像。
因此可用一个边界修改 $x$，达到 $x_n$ 是边界的情形。再修改一次 $x$，
可假设 $x_n = 0$。归纳可见，每个上同调类 $\xi$ 都由余闭元
$x = x_0$ 表示。最后，关于核的正合性条件告诉我们，两个这样的余闭元
$x_0$ 与 $x_0'$ 同调，当且仅当它们在 $H^0(M^\bullet)$ 中的像相同。
\end{proof}

\begin{lemma}
\label{homology-lemma-good-right-resolution-gives-qis}
设 $M^\bullet$ 为阿贝尔群的复形，并设
$$
0 \to M^\bullet \to A_0^\bullet \to A_1^\bullet \to A_2^\bullet \to \ldots
$$
为阿贝尔群复形组成的正合复形，并且对所有 $p \in \mathbf{Z}$，复形
$$
0 \to
\Coker(d_{M^\bullet}^p) \to
\Coker(d_{A_0^\bullet}^p) \to
\Coker(d_{A_1^\bullet}^p) \to
\Coker(d_{A_2^\bullet}^p) \to \ldots
$$
也正合。置 $A^{p, q} = A_p^q$，从而得到双复形。令
$\text{Tot}_\pi(A^{\bullet, \bullet})$ 为关联于该双复形的积总复形
（见证明）。那么，由 $M^\bullet \to A_0^\bullet$ 诱导的映射
$M^\bullet \to \text{Tot}_\pi(A^{\bullet, \bullet})$ 是拟同构。
\end{lemma}

\begin{proof}
简写 $T^\bullet = \text{Tot}_\pi(A^{\bullet, \bullet})$，定义
$$
T^n = \prod\nolimits_{p + q = n} A^{p, q} =
\prod\nolimits_{p + q = n} A_p^q
\quad\text{其中}\quad
\text{d}_{T^\bullet}^n =
\prod\nolimits_{n = p + q} (f_p^q + (-1)^pd_{A_p^\bullet}^q)
$$
其中 $f_p^\bullet : A_p^\bullet \to A_{p + 1}^\bullet$ 是引理中的
复形态射。

\medskip\noindent
我们将证明 $H^0(M^\bullet) \to H^0(T^\bullet)$ 是同构。
相同论证适用于其他次数。令 $x \in \Ker(\text{d}_{T^\bullet}^0)$
代表 $\xi \in H^0(T^\bullet)$。写 $x = (x_i)$，其中
$x_i \in A_i^{-i}$。注意，$x_0$ 在
$\Coker(A_1^{-1} \to A_1^0)$ 中映为零。因此，对某个
$m_0 \in M^0$ 与 $y \in A_0^{-1}$，有
$x_0 = m_0 + d_{A_0^\bullet}^{-1}(y)$。因为
$\text{d}_{T^\bullet}(x) = 0$ 蕴含
$\text{d}_{A_0^\bullet}(x_0) = 0$，所以
$d_{M^\bullet}(m_0) = 0$。因此，从 $\xi$ 中减去
$H^0(M^\bullet) \to H^0(T^\bullet)$ 像内的某个元素后，可假设
$x_0$ 属于 $\Im(d^{-1}_{A_0^\bullet})$。

\medskip\noindent
假设 $x_0 \in \Im(d^{-1}_{A_0^\bullet})$。断言此时 $\xi = 0$。
为证明这一点，对 $n$ 归纳寻找元素 $y_0, y_1, \ldots, y_n$，其中
$y_i \in A_i^{-i - 1}$，使得 $x_0 = \text{d}_{A_0}^{-1}(y_0)$，且当
$j = 1, \ldots, n$ 时
$x_j = f_{j - 1}^{-j}(y_{j - 1}) + (-1)^j d^{-j - 1}_{A_j^\bullet}(y_j)$。
当 $n = 0$ 时显然。证明归纳步骤：假设已经找到
$y_0, \ldots, y_{n - 1}$。那么
$w_n = x_n - f_{n - 1}^{-n}(y_{n - 1})$ 属于
$d^{-n}_{A_n^\bullet}$ 的核，并映到 $H^{-n}(A_{n + 1}^\bullet)$ 中的零
（因为它映到的元素在相差一个符号后是 $x_{n + 1}$ 的边界）。
完全如引理 \ref{homology-lemma-good-resolution-gives-qis} 的证明，假设蕴含对所有
$p \in \mathbf{Z}$，
$$
0 \to
H^p(M^\bullet) \to
H^p(A_0^\bullet) \to
H^p(A_1^\bullet) \to
H^p(A_2^\bullet) \to \ldots
$$
正合。因此，在用 $\Ker(d^{n - 1}_{A_{n - 1}^\bullet})$ 中的元素修改
$y_{n - 1}$ 后，可假设 $w_n$ 在 $H^{-n}(A_n^\bullet)$ 中映为零。
这意味着可以找到所需的 $y_n$。注意该过程不改变
$y_0, \ldots, y_{n - 2}$。因此无限继续下去，可在 $T^{n - 1}$ 中找到
元素 $y = (y_i)$，使得 $d_{T^\bullet}(y) = \xi$。
% P02-E0016: source line 7003 says y lies in T^{n-1}; the constructed components y_i have total degree -1, so T^{-1} appears required.
这说明 $H^0(M^\bullet) \to H^0(T^\bullet)$ 满射。

\medskip\noindent
假设 $m_0 \in \Ker(d^0_{M^\bullet})$ 在 $H^0(T^\bullet)$ 中映为零。
设它映到对 $y = (y_i) \in T^{-1}$ 施加微分所得的元素。
那么 $y_0 \in A_0^{-1}$ 在 $\Coker(d^{-2}_{A_1^\bullet})$ 中映为零。
由假设，这意味着 $y_0 \bmod \Im(d^{-2}_{A_0^\bullet})$ 是某个
$z \in M^{-1}$ 的像。于是 $m_0 = d^{-1}_{M^\bullet}(z)$。
这证明单射性，证明完成。
\end{proof}

\begin{lemma}
\label{homology-lemma-resolution-gives-qis}
设 $M^\bullet$ 为阿贝尔群的复形，并设
$$
\ldots \to A_2^\bullet \to A_1^\bullet \to A_0^\bullet \to M^\bullet \to 0
$$
为阿贝尔群复形组成的正合复形。置 $A^{p, q} = A_{-p}^q$，从而得到
双复形。令 $\text{Tot}_\pi(A^{\bullet, \bullet})$ 为关联于该双复形的
积总复形（见证明）。那么，由 $A_0^\bullet \to M^\bullet$ 诱导的映射
$\text{Tot}_\pi(A^{\bullet, \bullet}) \to M^\bullet$ 是拟同构。
\end{lemma}

\begin{proof}
简写 $T^\bullet = \text{Tot}_\pi(A^{\bullet, \bullet})$，定义
$$
T^n = \prod\nolimits_{p + q = n} A^{p, q} =
\prod\nolimits_{p + q = n} A_{-p}^q
\quad\text{其中}\quad
\text{d}_{T^\bullet}^n =
\prod\nolimits_{n = p + q} (f_{-p}^q + (-1)^pd_{A_{-p}^\bullet}^q)
$$
其中 $f_p^\bullet : A_p^\bullet \to A_{p - 1}^\bullet$ 是引理中的
复形态射。先证明当 $M^\bullet$ 为零复形时，$T^\bullet$ 无环。
下面的技巧说明这已足够。置 $B_n^\bullet = A_{n + 1}^\bullet$ 及
$B_0^\bullet = M^\bullet$。那么有如引理中的正合列
$$
\ldots \to B_2^\bullet \to B_1^\bullet \to B_0^\bullet \to 0 \to 0
$$
令 $S^\bullet = \text{Tot}_\pi(B^{\bullet, \bullet})$。于是有显然的
复形短正合列
$$
0 \to M^\bullet \to S^\bullet \to T^\bullet[1] \to 0
$$
再由长正合上同调列得到结论。略去一些细节。

\medskip\noindent
假设 $M^\bullet = 0$。我们将证明 $H^0(T^\bullet) = 0$。
相同论证适用于其他次数。令
$x =(x_n) \in \Ker(d_{T^\bullet})$ 映到 $\xi \in H^0(T^\bullet)$，
其中 $x_n \in A^{-n, n} = A_n^n$。由于 $M^0 = 0$，可得对某个
$y_0 \in A_1^0$ 有 $x_0 = f_1^0(y_0)$。于是
$x_1 - d^0_{A_1^\bullet}(y_0) = f_2^1(y_1)$，因为 $x$ 是余闭元，
故它被 $f_1^1$ 映为零；这里某个 $y_1 \in A_2^1$。
继续归纳，找到 $y = (y_i) \in T^{-1}$，使得
$d_{T^\bullet}(y) = x$，正如所需。
\end{proof}

\section{内射对象}
\label{homology-section-injectives}

\begin{definition}
\label{homology-definition-injective}
设 $\mathcal{A}$ 为阿贝尔范畴。若对每个单射态射
$A \hookrightarrow B$ 及每个态射 $A \to J$，都存在使下图交换的态射
$B \to J$，则称对象 $J \in \Ob(\mathcal{A})$ {\it 内射}：
$$
\xymatrix{
A \ar[r] \ar[d] & B \ar@{-->}[ld] \\
J &
}
$$
\end{definition}

\noindent
下面给出不可缺少的内射对象刻画。

\begin{lemma}
\label{homology-lemma-characterize-injectives}
设 $\mathcal{A}$ 为阿贝尔范畴，并设 $I$ 为 $\mathcal{A}$ 的对象。
下列条件等价：
\begin{enumerate}
\item 对象 $I$ 内射。
\item 函子 $B \mapsto \Hom_\mathcal{A}(B, I)$ 正合。
\item $\mathcal{A}$ 中任意短正合列
$$
0 \to I \to A \to B \to 0
$$
都分裂。
\item 对所有 $B \in \Ob(\mathcal{A})$，均有
$\Ext_\mathcal{A}(B, I) = 0$。
\end{enumerate}
\end{lemma}

\begin{proof}
略。
\end{proof}

\begin{lemma}
\label{homology-lemma-product-injectives}
设 $\mathcal{A}$ 为阿贝尔范畴。假设
$I_\omega$（$\omega \in \Omega$）是 $\mathcal{A}$ 的一族内射对象。
若 $\prod_{\omega \in \Omega} I_\omega$ 存在，则它内射。
\end{lemma}

\begin{proof}
略。
\end{proof}

\begin{definition}
\label{homology-definition-enough-injectives}
设 $\mathcal{A}$ 为阿贝尔范畴。若每个对象 $A$ 都有一个进入某个内射对象
$J$ 的单射态射 $A \to J$，则称 $\mathcal{A}$ {\it 有足够多的内射对象}。
\end{definition}

\begin{definition}
\label{homology-definition-functorial-injective-embedding}
设 $\mathcal{A}$ 为阿贝尔范畴。若存在函子
$$
J : \mathcal{A} \longrightarrow \text{Arrows}(\mathcal{A})
$$
满足
\begin{enumerate}
\item $s \circ J = \text{id}_\mathcal{A}$；
\item 对任意对象 $A \in \Ob(\mathcal{A})$，态射 $J(A)$ 是单射态射；
\item 对任意对象 $A \in \Ob(\mathcal{A})$，对象 $t(J(A))$ 是
$\mathcal{A}$ 的内射对象，
\end{enumerate}
则称 $\mathcal{A}$ {\it 具有函子性内射嵌入}。我们将这样的函子记作
$A \mapsto (A \to J(A))$。
\end{definition}

\section{投射对象}
\label{homology-section-projectives}

\begin{definition}
\label{homology-definition-projective}
设 $\mathcal{A}$ 为阿贝尔范畴。若对每个满射态射
$A \rightarrow B$ 及每个态射 $P \to B$，都存在使下图交换的态射
$P \to A$，则称对象 $P \in \Ob(\mathcal{A})$ {\it 投射}：
$$
\xymatrix{
A \ar[r] & B \\
P \ar@{-->}[u] \ar[ru] &
}
$$
\end{definition}

\noindent
下面给出不可缺少的投射对象刻画。

\begin{lemma}
\label{homology-lemma-characterize-projectives}
设 $\mathcal{A}$ 为阿贝尔范畴，并设 $P$ 为 $\mathcal{A}$ 的对象。
下列条件等价：
\begin{enumerate}
\item 对象 $P$ 投射。
\item 函子 $B \mapsto \Hom_\mathcal{A}(P, B)$ 正合。
\item $\mathcal{A}$ 中任意短正合列
$$
0 \to A \to B \to P \to 0
$$
都分裂。
\item 对所有 $A \in \Ob(\mathcal{A})$，均有
$\Ext_\mathcal{A}(P, A) = 0$。
\end{enumerate}
\end{lemma}

\begin{proof}
略。
\end{proof}

\begin{lemma}
\label{homology-lemma-coproduct-projectives}
设 $\mathcal{A}$ 为阿贝尔范畴。假设
$P_\omega$（$\omega \in \Omega$）是 $\mathcal{A}$ 的一族投射对象。
若 $\coprod_{\omega \in \Omega} P_\omega$ 存在，则它投射。
\end{lemma}

\begin{proof}
略。
\end{proof}

\begin{definition}
\label{homology-definition-enough-projectives}
设 $\mathcal{A}$ 为阿贝尔范畴。若每个对象 $A$ 都有一个从某个投射对象
$P$ 出发到它的满射态射 $P \to A$，则称 $\mathcal{A}$
{\it 有足够多的投射对象}。
\end{definition}

\begin{definition}
\label{homology-definition-functorial-projective-surjections}
设 $\mathcal{A}$ 为阿贝尔范畴。若存在函子
$$
P : \mathcal{A} \longrightarrow \text{Arrows}(\mathcal{A})
$$
满足
\begin{enumerate}
\item $t \circ P = \text{id}_\mathcal{A}$；
\item 对任意对象 $A \in \Ob(\mathcal{A})$，态射 $P(A)$ 是满射态射；
\item 对任意对象 $A \in \Ob(\mathcal{A})$，对象 $s(P(A))$ 是
$\mathcal{A}$ 的投射对象，
\end{enumerate}
则称 $\mathcal{A}$ {\it 具有函子性投射满射}。我们将这样的函子记作
$A \mapsto (P(A) \to A)$。
\end{definition}

\section{内射对象与伴随函子}
\label{homology-section-adjoint}

\noindent
下面给出关于伴随函子及其与内射对象之关系的若干引理。亦见引理
\ref{homology-lemma-adjoint-get-abelian}。

\begin{lemma}
\label{homology-lemma-adjoint-preserve-injectives}
\begin{slogan}
具有正合左伴随的函子保持内射对象。
\end{slogan}
设 $\mathcal{A}$ 与 $\mathcal{B}$ 为阿贝尔范畴，并设
$u : \mathcal{A} \to \mathcal{B}$ 与
$v : \mathcal{B} \to \mathcal{A}$ 为加性函子，其中 $u$ 是 $v$ 的右伴随。
考虑下列条件：
\begin{enumerate}
\item[(a)] $v$ 把单射态射变为单射态射；
\item[(b)] $v$ 正合；
\item[(c)] $u$ 把内射对象变为内射对象。
\end{enumerate}
那么 (a) $\Leftrightarrow$ (b) $\Rightarrow$ (c)。若 $\mathcal{A}$
有足够多的内射对象，则三个条件全都等价。
\end{lemma}

\begin{proof}
注意，作为左伴随，$v$ 右正合（范畴论，引理
\ref{categories-lemma-exact-adjoint}）。结合引理 \ref{homology-lemma-exact-functor}，
这解释了为何 (a) $\Leftrightarrow$ (b)。

\medskip\noindent
假设 (a)。设 $I$ 为 $\mathcal{A}$ 的内射对象，
$\varphi : N \to M$ 为 $\mathcal{B}$ 中的单射态射，并设
$\alpha : N \to uI$ 为态射。由伴随性得到态射
$\alpha : vN \to I$，而由假设，$v\varphi : vN \to vM$ 是单射态射。
由于 $I$ 内射，得到延拓 $\alpha$ 的态射 $\beta : vM \to I$。
再次由伴随性，它对应于延拓 $\alpha$ 的态射 $\beta : M \to uI$。
故 (c) 成立。

\medskip\noindent
假设 $\mathcal{A}$ 有足够多的内射对象且 (c) 成立。设
$f : B \to B'$ 为 $\mathcal{B}$ 中的单态射，并令
$A = \Ker(v(f))$。选取到内射对象 $I$ 的单态射 $g : A \to I$
（由假设它存在）。那么 $g$ 延拓为 $g' : v(B) \to I$，由伴随得到
态射 $B \to u(I)$。由于 $u(I)$ 内射，该态射延拓为
$h : B' \to u(I)$；由伴随得到延拓 $g'$ 的态射
$k : v(B') \to I$。但此时 $k$ “延拓” $g$；由于 $g$ 是单态射，
这迫使 $A = 0$。故 (a) 成立。
\end{proof}

\begin{remark}
\label{homology-remark-need-left-exactness}
设 $R \to S$ 为环同态。令 $u : \text{Mod}_S \to \text{Mod}_R$
由 $u(N) = N_R$ 给出，令 $v : \text{Mod}_R \to \text{Mod}_S$
由 $v(M) = M \otimes_R S$ 给出。那么 $u$ 是 $v$ 的右伴随，且
$u$ 正合、$v$ 右正合。但一般而言，引理
\ref{homology-lemma-adjoint-preserve-injectives} 的条件 (a)、(b)、(c) 并不成立。
例如，若 $R = \mathbf{Z}$、$S = \mathbf{Z}/p\mathbf{Z}$，则内射
$S$-模 $\mathbf{Z}/p\mathbf{Z}$ 不是内射 $\mathbf{Z}$-模。
事实上，该引理表明，所有内射 $S$-模作为 $R$-模都内射，当且仅当
$R \to S$ 是平坦环同态。
\end{remark}

\begin{lemma}
\label{homology-lemma-adjoint-enough-injectives}
设 $\mathcal{A}$ 与 $\mathcal{B}$ 为阿贝尔范畴，并设
$u : \mathcal{A} \to \mathcal{B}$ 与
$v : \mathcal{B} \to \mathcal{A}$ 为加性函子。假设
\begin{enumerate}
\item $u$ 是 $v$ 的右伴随；
\item $v$ 把单射态射变为单射态射；
\item $\mathcal{A}$ 有足够多的内射对象；
\item 对任意 $B \in \Ob(\mathcal{B})$，$vB = 0$ 蕴含 $B = 0$。
\end{enumerate}
那么 $\mathcal{B}$ 有足够多的内射对象。
\end{lemma}

\begin{proof}
取 $B \in \Ob(\mathcal{B})$。取到 $\mathcal{A}$ 的某个内射对象
$I$ 的单射态射 $vB \to I$。由引理
\ref{homology-lemma-adjoint-preserve-injectives} 与各假设，相应映射
$B \to uI$ 是把 $B$ 嵌入内射对象的单射态射。
\end{proof}

\begin{remark}
\label{homology-remark-faithfulness-needed}
设 $\mathcal{A}$、$\mathcal{B}$、$u : \mathcal{A} \to \mathcal{B}$ 与
$v : \mathcal{B} \to \mathcal{A}$ 如引理
\ref{homology-lemma-adjoint-enough-injectives}。在条件 (1) 与 (2) 下，条件 (4)
等价于 $v$ 忠实。而且，条件 (4) 是必要的。一个例子是函子 $u$ 与 $v$
都为零函子的情形。
\end{remark}

\begin{lemma}
\label{homology-lemma-adjoint-functorial-injectives}
设 $\mathcal{A}$ 与 $\mathcal{B}$ 为阿贝尔范畴，并设
$u : \mathcal{A} \to \mathcal{B}$ 与
$v : \mathcal{B} \to \mathcal{A}$ 为加性函子。假设
\begin{enumerate}
\item $u$ 是 $v$ 的右伴随；
\item $v$ 把单射态射变为单射态射；
\item $\mathcal{A}$ 有足够多的内射对象；
\item 对任意 $B \in \Ob(\mathcal{B})$，$vB = 0$ 蕴含 $B = 0$；
\item $\mathcal{A}$ 具有函子性内射嵌入。
\end{enumerate}
那么 $\mathcal{B}$ 具有函子性内射嵌入。
\end{lemma}

\begin{proof}
设 $A \mapsto (A \to J(A))$ 是 $\mathcal{A}$ 上的函子性内射嵌入。
那么 $B \mapsto (B \to uJ(vB))$ 是 $\mathcal{B}$ 上的函子性内射嵌入。
与引理 \ref{homology-lemma-adjoint-enough-injectives} 的证明比较。
\end{proof}

\begin{lemma}
\label{homology-lemma-partially-defined-adjoint}
设 $\mathcal{A}$ 与 $\mathcal{B}$ 为阿贝尔范畴，并设
$u : \mathcal{A} \to \mathcal{B}$ 为函子。若存在子集
$\mathcal{P} \subset \Ob(\mathcal{B})$，使得
\begin{enumerate}
\item $\mathcal{B}$ 的每个对象都是 $\mathcal{P}$ 中某个元素的商；
\item 对每个 $P \in \mathcal{P}$，存在 $\mathcal{A}$ 的对象 $Q$，
使得关于 $A$ 函子性地有
$\Hom_\mathcal{A}(Q, A) = \Hom_\mathcal{B}(P, u(A))$，
\end{enumerate}
则 $u$ 存在左伴随 $v$。
\end{lemma}

\begin{proof}
由 Yoneda 引理（范畴论，引理 \ref{categories-lemma-yoneda}），
对应于 $P$ 的 $\mathcal{A}$ 的对象 $Q$ 由公式
$\Hom_\mathcal{A}(Q, A) = \Hom_\mathcal{B}(P, u(A))$ 在唯一同构意义下
确定。写 $Q = v(P)$。用 $i_P : P \to u(v(P))$ 表示
$\Hom_\mathcal{A}(v(P), v(P))$ 中 $\text{id}_{v(P)}$ 对应的映射。
(2) 中的函子性蕴含该双射由
$$
\Hom_\mathcal{A}(v(P), A) \to \Hom_\mathcal{B}(P, u(A)),\quad
\varphi \mapsto u(\varphi) \circ i_P
$$
给出。对任意一对元素 $P_1, P_2 \in \mathcal{P}$，存在典范映射
$$
\Hom_\mathcal{B}(P_2, P_1)
\to
\Hom_\mathcal{A}(v(P_2), v(P_1)),\quad
\varphi \mapsto v(\varphi)
$$
其特征规则是在 $\Hom_\mathcal{B}(P_2, u(v(P_1)))$ 中有
$u(v(\varphi)) \circ i_{P_2} = i_{P_1} \circ \varphi$。
注意 $\varphi \mapsto v(\varphi)$ 与复合相容；这可直接由该特征看出。
因此 $P \mapsto v(P)$ 是从 $\mathcal{B}$ 中以 $\mathcal{P}$ 的元素为
对象的全子范畴出发的函子。

\medskip\noindent
给定 $\mathcal{B}$ 的任意对象 $B$，选择正合列
$$
P_2 \to P_1 \to B \to 0
$$
这由假设 (1) 可行。把 $v(B)$ 定义为适合正合列
$$
v(P_2) \to v(P_1) \to v(B) \to 0
$$
的 $\mathcal{A}$ 的对象。那么
\begin{align*}
\Hom_\mathcal{A}(v(B), A)
& =
\Ker(\Hom_\mathcal{A}(v(P_1), A) \to \Hom_\mathcal{A}(v(P_2), A)) \\
& =
\Ker(\Hom_\mathcal{B}(P_1, u(A)) \to \Hom_\mathcal{B}(P_2, u(A))) \\
& =
\Hom_\mathcal{B}(B, u(A))
\end{align*}
因此可取 $\mathcal{P} = \Ob(\mathcal{B})$；也就是说，$v$ 处处有定义。
\end{proof}

\section{本质常值系统}
\label{homology-section-essentially-constant}

\noindent
本节讨论取值于加性范畴的本质常值系统。

\begin{lemma}
\label{homology-lemma-essentially-constant-into-karoubian}
设 $\mathcal{I}$ 为范畴，$\mathcal{A}$ 为预加性 Karoubian 范畴，并设
$M : \mathcal{I} \to \mathcal{A}$ 为图式。
\begin{enumerate}
\item 假设 $\mathcal{I}$ 是滤过的。下列条件等价：
\begin{enumerate}
\item $M$ 本质常值；
\item $X = \colim M$ 存在，且存在余终滤过子范畴
$\mathcal{I}' \subset \mathcal{I}$，使得对
$i' \in \Ob(\mathcal{I}')$ 有直和分解
$M_{i'} = X_{i'} \oplus Z_{i'}$，其中 $X_{i'}$ 同构地映到 $X$，
而对 $\mathcal{I}'$ 中的某个 $i' \to i''$，$Z_{i'}$ 在
$M_{i''}$ 中的像为零。
\end{enumerate}
\item 假设 $\mathcal{I}$ 是余滤过的。下列条件等价：
\begin{enumerate}
\item $M$ 本质常值；
\item $X = \lim M$ 存在，且存在共始余滤过子范畴
$\mathcal{I}' \subset \mathcal{I}$，使得对
$i' \in \Ob(\mathcal{I}')$ 有直和分解
$M_{i'} = X_{i'} \oplus Z_{i'}$，其中 $X$ 同构地映到
$X_{i'}$，而对 $\mathcal{I}'$ 中的某个 $i'' \to i'$，
$M_{i''} \to Z_{i'}$ 为零。
\end{enumerate}
\end{enumerate}
\end{lemma}

\begin{proof}
假设 (1)(a) 成立，即 $\mathcal{I}$ 是滤过的且 $M$ 本质常值。
令 $X = \colim M_i$。按《范畴》，定义
\ref{categories-definition-essentially-constant-diagram} 选取
$i$ 与 $X \to M_i$。令 $\mathcal{I}'$ 为如下对象所组成的全子范畴：
它们都是以 $i$ 为源的某个态射的靶。设
$i' \in \Ob(\mathcal{I}')$，并选取态射 $i \to i'$。
那么 $X \to M_i \to M_{i'}$ 与 $M_{i'} \to X$ 的复合是
$X$ 上的恒等态射。由于 $\mathcal{A}$ 是 Karoubian 范畴，可得直和分解
$M_{i'} = X_{i'} \oplus Z_{i'}$，其中
$Z_{i'} = \Ker(M_{i'} \to X)$，且 $X_{i'}$ 同构地映到 $X$。
按《范畴》，定义
\ref{categories-definition-essentially-constant-diagram}，选取
$i \to k$ 与 $i' \to k$，使
$M_{i'} \to X \to M_i \to M_k$ 等于 $M_{i'} \to M_k$。
于是 $M_{i'} \to M_k$ 在 $Z_{i'}$ 上为零。因此 (1)(b) 成立。

\medskip\noindent
假设 (1)(b) 成立，即 $\mathcal{I}$ 是滤过的，并且有
$\mathcal{I}' \subset \mathcal{I}$ 以及对
$i' \in \Ob(\mathcal{I}')$ 的直和分解
$M_{i'} = X_{i'} \oplus Z_{i'}$，满足引理所述条件。为证明
$M$ 本质常值，可把 $\mathcal{I}$ 换成 $\mathcal{I}'$；见
《范畴》，引理 \ref{categories-lemma-cofinal-essentially-constant}。
任取 $i \in \Ob(\mathcal{I})$，以 $X \to M_i$ 表示同构
$X_i \to X$ 的逆与包含映射 $X_i \to M_i$ 的复合。若 $j$ 是另一对象，
则选取 $j \to k$，使 $Z_j \to M_k$ 为零。由于
$\mathcal{I}$ 是滤过的，在必要时增大 $k$ 后，还可假设存在态射
$i \to k$。于是 $M_j \to X \to M_i \to M_k$ 与
$M_j \to M_k$ 都在 $Z_j$ 上为零。再后复合一个使直和项
$Z_k$ 为零的态射 $M_k \to M_l$，便得到
$M_j \to X \to M_i \to M_l$ 与 $M_j \to M_l$ 相等；也就是说，
$M$ 本质常值。

\medskip\noindent
(2) 的证明与之对偶。
\end{proof}

\begin{lemma}
\label{homology-lemma-direct-sum-from-product-colimit}
设 $\mathcal{I}$ 为范畴，$\mathcal{A}$ 为加性 Karoubian 范畴，并设
$F : \mathcal{I} \to \mathcal{A}$ 与
$G : \mathcal{I} \to \mathcal{A}$ 为函子。下列条件等价：
\begin{enumerate}
\item $\colim_\mathcal{I} F \oplus G$ 存在；
\item $\colim_\mathcal{I} F$ 与 $\colim_\mathcal{I} G$ 均存在。
\end{enumerate}
在这种情形下，
$\colim_\mathcal{I} F \oplus G =
\colim_\mathcal{I} F \oplus \colim_\mathcal{I} G$。
\end{lemma}

\begin{proof}
假设 (1) 成立。令 $W = \colim_\mathcal{I} F \oplus G$。
注意，到 $F$ 上的投影定义了幂等自然变换
$F \oplus G \to F \oplus G$。因此由《范畴》，引理
\ref{categories-lemma-functorial-colimit}，得到幂等自同态
$W \to W$。由于 $\mathcal{A}$ 是 Karoubian 范畴，得到相应的直和分解
$W = X \oplus Y$；见引理 \ref{homology-lemma-karoubian}。一个直接的论证
（略去）表明 $X = \colim_\mathcal{I} F$ 且
$Y = \colim_\mathcal{I} G$。故 (2) 成立。我们略去 (2) 蕴含 (1)
的证明。
\end{proof}

\begin{lemma}
\label{homology-lemma-direct-sum-from-product-essentially-constant}
设 $\mathcal{I}$ 为滤过范畴，$\mathcal{A}$ 为加性 Karoubian 范畴，
并设 $F : \mathcal{I} \to \mathcal{A}$ 与
$G : \mathcal{I} \to \mathcal{A}$ 为函子。下列条件等价：
\begin{enumerate}
\item $F \oplus G : \mathcal{I} \to \mathcal{A}$ 本质常值；
\item $F$ 与 $G$ 都本质常值。
\end{enumerate}
\end{lemma}

\begin{proof}
假设 (1) 成立。特别地，
$W = \colim_\mathcal{I} F \oplus G$ 存在，因而由引理
\ref{homology-lemma-direct-sum-from-product-colimit}，有
$W = X \oplus Y$，其中 $X = \colim_\mathcal{I} F$ 且
$Y = \colim_\mathcal{I} G$。例如利用《范畴》，引理
\ref{categories-lemma-characterize-essentially-constant-ind}
中的刻画，一个直接的论证（略去）表明 $F$ 以 $X$ 为值本质常值，
而 $G$ 以 $Y$ 为值本质常值。故 (2) 成立。略去 (2) 蕴含 (1)
的证明。
\end{proof}





\section{逆系统}
\label{homology-section-inverse-systems}

\noindent
设 $\mathcal{C}$ 为范畴。在《范畴》，第
\ref{categories-section-posets-limits} 节中，我们定义了预序集上的
（取值于范畴 $\mathcal{C}$ 的）逆系统。若该预序集是带通常次序的
$\mathbf{N} = \{1, 2, 3, \ldots\}$，则 $\mathbf{N}$ 上的这种逆系统
通常简称为{\it 逆系统}。它无非由一对 $(M_i, f_{ii'})$ 构成，其中对每个
$i \in \mathbf{N}$，$M_i$ 是 $\mathcal{C}$ 的对象；并且对每个
$i > i'$，$i, i' \in \mathbf{N}$，都有态射
$f_{ii'} : M_i \to M_{i'}$，还要求在复合有意义时
$f_{i'i''} \circ f_{ii'} = f_{ii''}$。显然，实际上只需给出态射
$M_2 \to M_1$、$M_3 \to M_2$，依此类推。因此逆系统经常画成
$$
M_1 \xleftarrow{\varphi_2} M_2 \xleftarrow{\varphi_3} M_3 \leftarrow \ldots
$$
此外，我们常从记号中省去转移映射 $\varphi_i$，而只说
“设 $(M_i)$ 为逆系统”。

\medskip\noindent
所有取值于 $\mathcal{C}$ 的逆系统在显然的态射概念下构成范畴。

\begin{lemma}
\label{homology-lemma-inverse-systems-abelian}
设 $\mathcal{C}$ 为范畴。
\begin{enumerate}
\item 若 $\mathcal{C}$ 是加性范畴，则取值于 $\mathcal{C}$ 的逆系统范畴
是加性范畴。
\item 若 $\mathcal{C}$ 是阿贝尔范畴，则取值于 $\mathcal{C}$ 的逆系统
范畴是阿贝尔范畴。逆系统的序列
$(K_i) \to (L_i) \to (M_i)$ 正合，当且仅当每个
% P02-E0017: frozen source names N_i here although the displayed system names M_i.
$K_i \to L_i \to N_i$ 都正合。
\end{enumerate}
\end{lemma}

\begin{proof}
略。
\end{proof}

\noindent
这种逆系统的极限（见《范畴》，第
\ref{categories-section-posets-limits} 节）记为 $\lim M_i$ 或
$\lim_i M_i$。若 $\mathcal{C}$ 是阿贝尔群（或集合）的范畴，则极限
总是存在，并且实际上可描述为
$$
\lim_i M_i
=
\{(x_i) \in \prod M_i \mid \varphi_i(x_i) = x_{i - 1}, \ i = 2, 3, \ldots\}
$$
参见《范畴》，第 \ref{categories-section-limit-sets} 节。然而，给定
阿贝尔群逆系统的短正合列
$$
0 \to (A_i) \to (B_i) \to (C_i) \to 0
$$
相应的极限列并不总是正合。为进一步讨论这一问题，引入如下概念。

\begin{definition}
\label{homology-definition-Mittag-Leffler}
设 $\mathcal{C}$ 为阿贝尔范畴。若对每个 $i$ 都存在
$c = c(i) \geq i$，使得
$$
\Im(A_k \to A_i) = \Im(A_c \to A_i)
$$
对所有 $k \geq c$ 成立，则称逆系统 $(A_i)$ 满足
{\it Mittag-Leffler 条件}，或简称为 {\it ML}。
\end{definition}

\noindent
事实表明，只要在阿贝尔群、某环上的模等所构成的阿贝尔范畴中工作，
Mittag-Leffler 条件就足以保证 $\lim$-函子正合。A.\ Neeman 的一篇论文
（见 \cite{Neeman-Counterexample}）表明，在具有 AB4*（即积正合）的
阿贝尔范畴中，该条件并不足够强。

\begin{lemma}
\label{homology-lemma-Mittag-Leffler}
设
$$
0 \to (A_i) \to (B_i) \to (C_i) \to 0
$$
为阿贝尔群逆系统的短正合列。
\begin{enumerate}
\item 无论如何，序列
$$
0 \to \lim_i A_i \to \lim_i B_i \to \lim_i C_i
$$
是正合的。
\item 若 $(B_i)$ 是 ML，则 $(C_i)$ 也是 ML。
\item 若 $(A_i)$ 是 ML，则
$$
0 \to \lim_i A_i \to \lim_i B_i \to \lim_i C_i \to 0
$$
是正合的。
\end{enumerate}
\end{lemma}

\begin{proof}
这是一个很好的练习。关于 (3)，见《代数》，引理
\ref{algebra-lemma-Mittag-Leffler}。
\end{proof}

\begin{lemma}
\label{homology-lemma-apply-Mittag-Leffler}
设
$$
(A_i) \to (B_i) \to (C_i) \to (D_i)
$$
为阿贝尔群逆系统的正合列。若系统 $(A_i)$ 是 ML，则序列
$$
\lim_i B_i \to \lim_i C_i \to \lim_i D_i
$$
正合。
\end{lemma}

\begin{proof}
令 $Z_i = \Ker(C_i \to D_i)$ 且 $I_i = \Im(A_i \to B_i)$。
那么 $\lim Z_i = \Ker(\lim C_i \to \lim D_i)$，并且得到系统的短正合列
$$
0 \to (I_i) \to (B_i) \to (Z_i) \to 0
$$
此外，由引理 \ref{homology-lemma-Mittag-Leffler} 可知 $(I_i)$ 具有 (ML)；
再次应用引理 \ref{homology-lemma-Mittag-Leffler}，便知
$\lim B_i \to \lim Z_i$ 满射，从而证明引理。
\end{proof}

\noindent
下述对本质常值逆系统的刻画特别表明它们具有 ML。

\begin{lemma}
\label{homology-lemma-essentially-constant}
设 $\mathcal{A}$ 为阿贝尔范畴，$(A_i)$ 为 $\mathcal{A}$ 中的逆系统，
其极限为 $A = \lim A_i$。那么 $(A_i)$ 本质常值（见《范畴》，定义
\ref{categories-definition-essentially-constant-diagram}），当且仅当存在
某个 $i$，使得对所有 $j \geq i$ 都有直和分解
$A_j = A \oplus Z_j$，并且：(a) 映射 $A_{j'} \to A_j$ 与这些直和分解
相容；(b) 对所有 $j$，存在某个 $j' \geq j$，使
$Z_{j'} \to Z_j$ 为零。
\end{lemma}

\begin{proof}
假设 $(A_i)$ 本质常值。那么存在某个 $i$ 以及态射
$A_i \to A$，使 $A \to A_i \to A$ 为恒等态射；并且对所有
$j \geq i$，存在 $j' \geq j$，使 $A_{j'} \to A_j$ 分解为
$A_{j'} \to A_i \to A \to A_j$（最后一个映射来自
$A = \lim A_i$）。于是，对所有 $j \geq i$ 令
$Z_j = \Ker(A_j \to A)$ 即可。略去逆命题的证明。
\end{proof}

\noindent
我们将在《更多代数》，引理
\ref{more-algebra-lemma-Mittag-Leffler} 中改进下述引理。

\begin{lemma}
\label{homology-lemma-exact-sequence-ML}
设
$$
0 \to (A_i) \to (B_i) \to (C_i) \to 0
$$
为阿贝尔群逆系统的正合列。若 $(C_i)$ 本质常值，则 $(A_i)$ 具有 ML
当且仅当 $(B_i)$ 具有 ML。
\end{lemma}

\begin{proof}
重新编号后，可假设 $C_i = C \oplus Z_i$，该分解与转移映射相容，
并且对所有 $i$ 都存在 $i' \geq i$，使 $Z_{i'} \to Z_i$ 为零；
见引理 \ref{homology-lemma-essentially-constant}。

\medskip\noindent
先假设 $C = 0$，即 $C_i = Z_i$。在这种情形下，选取
$1 = n_1 < n_2 < n_3 < \ldots$，使
$Z_{n_{i + 1}} \to Z_{n_i}$ 为零。于是
$B_{n_{i + 1}} \to B_{n_i}$ 通过 $A_{n_i} \subset B_{n_i}$ 分解。
因此，当 $j \geq i + 1$ 时，在 $A_{n_i}$ 中有
$$
\Im(A_{n_j} \to A_{n_i}) \subset \Im(B_{n_j} \to B_{n_i}) \subset
\Im(A_{n_{j - 1}} \to A_{n_i})
$$
所以当 $j \geq i + 1$ 时，像 $\Im(A_{n_j} \to A_{n_i})$ 稳定，
当且仅当像 $\Im(B_{n_j} \to B_{n_i})$ 稳定。所需等价性由此推出
（略去一个小细节）。

\medskip\noindent
若 $C \not = 0$，以 $B'_i \subset B_i$ 表示映射
$B_i \to C \oplus Z_i$ 下 $C$ 的逆像。由前一段可知，$(B'_i)$
具有 ML 当且仅当 $(B_i)$ 具有 ML。因此可用 $(B'_i)$ 替换
$(B_i)$。此时，对所有 $i$ 都有正合列
$0 \to A_i \to B_i \to C \to 0$。由此，对所有 $j \geq i$，
$$
0 \to \Im(A_j \to A_i) \to \Im(B_j \to B_i) \to C \to 0
$$
均为短正合列。故当 $j \geq i$ 时，像
$\Im(A_j \to A_i)$ 稳定，当且仅当像
$\Im(B_j \to B_i)$ 稳定，正如所需。
\end{proof}

\noindent
下述引理的“正确”版本见《更多代数》，引理
\ref{more-algebra-lemma-apply-Mittag-Leffler-again}。

\begin{lemma}
\label{homology-lemma-apply-Mittag-Leffler-again}
设
$$
(A^{-2}_i \to A^{-1}_i \to A^0_i \to A^1_i)
$$
为阿贝尔群复形的逆系统，并以
$A^{-2} \to A^{-1} \to A^0 \to A^1$ 表示其极限。以
$(H_i^{-1})$、$(H_i^0)$ 表示各上同调所成的逆系统，并以
$H^{-1}$、$H^0$ 表示复形
$A^{-2} \to A^{-1} \to A^0 \to A^1$ 的上同调。若
$(A^{-2}_i)$ 与 $(A^{-1}_i)$ 是 ML，且
$(H^{-1}_i)$ 本质常值，则 $H^0 = \lim H_i^0$。
\end{lemma}

\begin{proof}
令 $Z^j_i = \Ker(A^j_i \to A^{j + 1}_i)$ 且
$I^j_i = \Im(A^{j - 1}_i \to A^j_i)$。注意
$\lim Z^0_i = \Ker(\lim A^0_i \to \lim A^1_i)$，因为取核与极限可交换。
系统 $(I^{-1}_i)$ 与 $(I^0_i)$ 作为系统 $(A^{-2}_i)$ 与
$(A^{-1}_i)$ 的商而具有 ML；见引理 \ref{homology-lemma-Mittag-Leffler}。
因而有逆系统的正合列
$$
0 \to (I^{-1}_i) \to (Z^{-1}_i) \to (H^{-1}_i) \to 0
$$
其中 $(I^{-1}_i)$ 具有 ML，而 $(H^{-1}_i)$ 依假设本质常值。
故由引理 \ref{homology-lemma-exact-sequence-ML}，$(Z^{-1}_i)$ 具有 ML。
正合列
$$
0 \to (Z^{-1}_i) \to (A^{-1}_i) \to (I^0_i) \to 0
$$
以及引理 \ref{homology-lemma-Mittag-Leffler} 的一次应用表明
$\lim A^{-1}_i \to \lim I^0_i$ 满射。最后，正合列
$$
0 \to (I^0_i) \to (Z^0_i) \to (H^0_i) \to 0
$$
和引理 \ref{homology-lemma-Mittag-Leffler} 表明
$\lim I^0_i \to \lim Z^0_i \to \lim H^0_i \to 0$ 正合。
综合起来即得结论。
\end{proof}

\noindent
有时我们需要上述引理在大序数上取极限的版本。

\begin{lemma}
\label{homology-lemma-ML-over-ordinals}
设 $\alpha$ 为序数，并设 $K_\beta^\bullet$，$\beta < \alpha$，
为 $\alpha$ 上的阿贝尔群复形逆系统。若对所有
$\beta < \alpha$，复形 $K_\beta^\bullet$ 都无环，且映射
$$
K^n_\beta \longrightarrow \lim_{\gamma < \beta} K^n_\gamma
$$
满射，则复形 $\lim_{\beta < \alpha} K_\beta^\bullet$ 无环。
\end{lemma}

\begin{proof}
用超限归纳证明，对每个序数 $\alpha$ 以及引理所述的每个系统，结论均
成立。特别地，在证明关于 $\alpha$ 的结论时，可以假设复形
% P02-E0018: frozen source writes K^n here although it calls the limit a complex.
$\lim_{\gamma < \beta} K^n_\gamma$ 无环。

\medskip\noindent
设 $x \in \lim_{\beta < \alpha} K^0_\beta$ 且
$\text{d}(x) = 0$。我们要找
$y \in \lim_{\beta < \alpha} K^{-1}_\beta$，使
$\text{d}(y) = x$。写 $x = (x_\beta)$，其中对
$\beta < \alpha$，$x_\beta \in K_\beta^0$ 是 $x$ 的像。
我们将用超限递归构造 $y = (y_\beta)$。

\medskip\noindent
当 $\beta = 0$ 时，任取 $y_0 \in K_0^{-1}$，使
$\text{d}(y_0) = x_0$。

\medskip\noindent
当 $\beta = \gamma + 1$ 为后继序数时，要找元素 $y_\beta$，使它在转移
映射 $f : K^\bullet_\beta \to K^\bullet_\gamma$ 下映到
$y_\gamma$，并在微分下映到 $x_\beta$。作为第一次近似，选取
$y'_\beta$，使 $\text{d}(y'_\beta) = x_\beta$。于是差
$y_\gamma - f(y'_\beta)$ 位于微分的核中，故等于
$\text{d}(z_\gamma)$，其中某个 $z_\gamma \in K^{-2}_\gamma$。
由假设，映射 $f^{-2} : K^{-2}_\beta \to K^{-2}_\gamma$ 满射。
因而可写 $z_\gamma = f(z_\beta)$，并把 $y'_\beta$ 改为
$y_\beta = y'_\beta + \text{d}(z_\beta)$，这便满足要求。

\medskip\noindent
若 $\beta$ 是极限序数，则有
$\lim_{\gamma < \beta} K^{-1}_\gamma$ 中的元素
$(y_\gamma)_{\gamma < \beta}$，其微分是 $x_\beta$ 的像。因此可与上面
完全同样地论证：使用逐项满射的复形映射
$f : K_\beta^\bullet \to \lim_{\gamma < \beta} K_\gamma^\bullet$，
以及由归纳可假设
$\lim_{\gamma < \beta} K_\gamma^\bullet$ 无环这一事实（见证明第一段）。
\end{proof}


\section{积的正合性}
\label{homology-section-product-exact}


\begin{lemma}
\label{homology-lemma-product-abelian-groups-exact}
设 $I$ 为集合。对 $i \in I$，设 $L_i \to M_i \to N_i$ 为阿贝尔群
复形，并设 $H_i = \Ker(M_i \to N_i)/\Im(L_i \to M_i)$ 为其同调。
那么
$$
\prod L_i \to \prod M_i \to \prod N_i
$$
是阿贝尔群复形，其同调为 $\prod H_i$。
\end{lemma}

\begin{proof}
略。
\end{proof}
% Generated mechanically from the frozen derived.tex target.
% Do not edit this generated file; edit the bound locale source instead.

% OK, start here.
%

\setcounter{chapter}{12}
\chapter{导出范畴}



\phantomsection
\label{derived-section-phantom}


\section{引言}
\label{derived-section-introduction}

\noindent
我们首先讨论三角范畴以及三角范畴中的局部化。随后证明：加性范畴中的复形所成的
同伦范畴是三角范畴。在此基础上，我们把阿贝尔范畴的导出范畴定义为其同伦范畴
关于拟同构的局部化。Verdier 的学位论文 \cite{Verdier} 是一份很好的参考文献。



\section{三角范畴}
\label{derived-section-triangulated-categories}

\noindent
三角范畴是描述阿贝尔范畴的导出范畴所固有结构的一种便利工具。可参见
\cite{Verdier}、\cite{KS} 和 \cite{Neeman}。




\section{三角范畴的定义}
\label{derived-section-triangulated-definitions}

\noindent
本节汇集有关三角范畴和预三角范畴的大部分定义。

\begin{definition}
\label{derived-definition-triangle}
设 $\mathcal{D}$ 为加性范畴。设
$[1] : \mathcal{D} \to \mathcal{D}$，$E \mapsto E[1]$
为 $\mathcal{D}$ 的一个加性自等价函子。
\begin{enumerate}
\item 一个{\it 三角}是六元组
$(X, Y, Z, f, g, h)$，其中 $X, Y, Z \in \Ob(\mathcal{D})$，且
$f : X \to Y$、$g : Y \to Z$ 和 $h : Z \to X[1]$ 是
$\mathcal{D}$ 中的态射。
\item 三角之间的一个{\it 态射}
$(X, Y, Z, f, g, h) \to (X', Y', Z', f', g', h')$
由 $\mathcal{D}$ 中的态射 $a : X \to X'$、$b : Y \to Y'$ 和
$c : Z \to Z'$ 给出，并满足
$b \circ f = f' \circ a$、$c  \circ g = g' \circ b$ 以及
$a[1] \circ h = h' \circ c$。
\end{enumerate}
\end{definition}

\noindent
三角的态射可由下列交换图表示：
$$
\xymatrix{
X \ar[r] \ar[d]^a &
Y \ar[r] \ar[d]^b &
Z \ar[r] \ar[d]^c &
X[1] \ar[d]^{a[1]} \\
X' \ar[r] &
Y' \ar[r] &
Z' \ar[r] &
X'[1]
}
$$
在定义 \ref{derived-definition-triangle} 的情形中，记
$[0] = \text{id}$；当 $n > 0$ 时，以 $[n]$ 表示 $[1]$ 的 $n$ 次复合；
选择 $[1]$ 的一个拟逆 $[-1]$；并令 $[-n]$ 等于 $[-1]$ 的 $n$ 次复合。
于是 $\{[n]\}_{n \in \mathbf{Z}}$ 是一族以 $n \in \mathbf{Z}$ 为指标的
$\mathcal{D}$ 的加性自等价，并且给定函子同构
$[n] \circ [m] \cong [n + m]$。

\medskip\noindent
下面是 Verdier 学位论文中给出的三角范畴定义。

\begin{definition}
\label{derived-definition-triangulated-category}
一个{\it 三角范畴}由三元组
$(\mathcal{D}, \{[n]\}_{n\in \mathbf{Z}}, \mathcal{T})$
构成，其中
\begin{enumerate}
\item $\mathcal{D}$ 是加性范畴；
\item $[1] : \mathcal{D} \to \mathcal{D}$，$E \mapsto E[1]$
是加性自等价，而 $n \in \mathbf{Z}$ 时的 $[n]$ 如上所述；
\item $\mathcal{T}$ 是一组三角（定义 \ref{derived-definition-triangle}），
称为{\it 特异三角}。
\end{enumerate}
它们须满足下列条件：
\begin{enumerate}
\item[TR1] 与特异三角同构的任何三角都是特异三角。形如
$(X, X, 0, \text{id}, 0, 0)$ 的任何三角都是特异三角。
对 $\mathcal{D}$ 的任一态射 $f : X \to Y$，存在形如
$(X, Y, Z, f, g, h)$ 的特异三角。
\item[TR2] 三角 $(X, Y, Z, f, g, h)$ 是特异三角，当且仅当三角
$(Y, Z, X[1], g, h, -f[1])$ 是特异三角。
\item[TR3] 给定实线图
$$
\xymatrix{
X \ar[r]^f \ar[d]^a &
Y \ar[r]^g \ar[d]^b &
Z \ar[r]^h \ar@{-->}[d] &
X[1] \ar[d]^{a[1]} \\
X' \ar[r]^{f'} &
Y' \ar[r]^{g'} &
Z' \ar[r]^{h'} &
X'[1]
}
$$
其两行均为特异三角，并且满足 $b \circ f = f' \circ a$，则存在态射
$c : Z \to Z'$，使得 $(a, b, c)$ 是三角的态射。
\item[TR4] 给定 $\mathcal{D}$ 的对象 $X$、$Y$、$Z$，态射
$f : X \to Y$、$g : Y \to Z$，以及特异三角
$(X, Y, Q_1, f, p_1, d_1)$、
$(X, Z, Q_2, g \circ f, p_2, d_2)$ 和
$(Y, Z, Q_3, g, p_3, d_3)$，则存在态射
$a : Q_1 \to Q_2$ 和 $b : Q_2 \to Q_3$，使得
\begin{enumerate}
\item $(Q_1, Q_2, Q_3, a, b, p_1[1] \circ d_3)$ 是特异三角；
\item 三元组 $(\text{id}_X, g, a)$ 是三角之间的态射
$(X, Y, Q_1, f, p_1, d_1) \to (X, Z, Q_2, g \circ f, p_2, d_2)$；并且
\item 三元组 $(f, \text{id}_Z, b)$ 是三角之间的态射
$(X, Z, Q_2, g \circ f, p_2, d_2) \to (Y, Z, Q_3, g, p_3, d_3)$。
\end{enumerate}
\end{enumerate}
若 TR1、TR2 和 TR3 成立，则称
$(\mathcal{D}, [\ ], \mathcal{T})$ 为{\it 预三角范畴}。\footnote{这里用
$[\ ]$ 作为族 $\{[n]\}_{n\in \mathbf{Z}}$ 的缩写。}
\end{definition}

\noindent
TR4 的含义如下：若把 $Q_1$ 看作 $Y/X$，把 $Q_2$ 看作 $Z/X$，
把 $Q_3$ 看作 $Z/Y$，则 TR4(a) 表达同构
$(Z/X)/(Y/X) \cong Z/Y$，而 TR4(b) 与 TR4(c) 表明我们可以用三角的态射
来比较三角 $X \to Y \to Q_1 \to X[1]$ 等。关于这一思想更精确的重述，
参见引理 \ref{derived-lemma-two-split-injections} 的证明。

\medskip\noindent
TR2 中的符号意味着：若 $(X, Y, Z, f, g, h)$ 是特异三角，则长序列
\begin{equation}
\label{derived-equation-rotate}
\ldots \to
Z[-1] \xrightarrow{-h[-1]}
X \xrightarrow{f}
Y \xrightarrow{g}
Z \xrightarrow{h}
X[1] \xrightarrow{-f[1]}
Y[1] \xrightarrow{-g[1]}
Z[1] \to \ldots
\end{equation}
中每个四项序列都给出一个特异三角。

\medskip\noindent
照例，我们略微滥用记号，只说一个（预）三角范畴 $\mathcal{D}$，而不显式引入
附加数据的记号。预三角范畴这一概念可用于寻找与 TR4 等价的命题。

\medskip\noindent
下面定义三角函子。

\begin{definition}
\label{derived-definition-exact-functor-triangulated-categories}
设 $\mathcal{D}$、$\mathcal{D}'$ 为预三角范畴。从 $\mathcal{D}$ 到
$\mathcal{D}'$ 的一个{\it 正合函子}，或称{\it 三角函子}，是函子
$F : \mathcal{D} \to \mathcal{D}'$ 连同给定的函子同构
$\xi_X : F(X[1]) \to F(X)[1]$，并且对 $\mathcal{D}$ 的每个特异三角
$(X, Y, Z, f, g, h)$，三角
$(F(X), F(Y), F(Z), F(f), F(g), \xi_X \circ F(h))$
是 $\mathcal{D}'$ 的特异三角。
\end{definition}

\noindent
正合函子是加性函子，见引理 \ref{derived-lemma-exact-functor-additive}。
所谓两个三角范畴等价，是指它们在三角范畴的 $2$-范畴中等价。
该 $2$-范畴中的一个 $2$-态射 $a : (F, \xi) \to (F', \xi')$，
就是与 $\xi$ 和 $\xi'$ 相容的函子变换 $a : F \to F'$，亦即下图交换：
$$
\xymatrix{
F \circ [1] \ar[r]_\xi \ar[d]_{a \star 1} & [1] \circ F \ar[d]^{1 \star a} \\
F' \circ [1] \ar[r]^{\xi'} & [1] \circ F'
}
$$

\begin{definition}
\label{derived-definition-triangulated-subcategory}
设 $(\mathcal{D}, [\ ], \mathcal{T})$ 为预三角范畴。一个
{\it 预三角子范畴}\footnote{这一定义可能并非标准定义。若
$\mathcal{D}'$ 是满子范畴，则 $\mathcal{T}'$ 是 $\mathcal{D}'$ 中三角之集
与 $\mathcal{T}$ 的交，见引理 \ref{derived-lemma-triangulated-subcategory}。
在这种情形下，我们在记号中省略 $\mathcal{T}'$。}
是一个偶 $(\mathcal{D}', \mathcal{T}')$，满足
\begin{enumerate}
\item $\mathcal{D}'$ 是 $\mathcal{D}$ 的加性子范畴，它在 $[1]$ 下保持，
且 $[1] : \mathcal{D}' \to \mathcal{D}'$ 是自等价；
\item $\mathcal{T}' \subset \mathcal{T}$，并且对每个
$(X, Y, Z, f, g, h) \in \mathcal{T}'$，有
$X, Y, Z \in \Ob(\mathcal{D}')$，且
$f, g, h \in \text{Arrows}(\mathcal{D}')$；
\item $(\mathcal{D}', [\ ], \mathcal{T}')$ 是预三角范畴。
\end{enumerate}
若 $\mathcal{D}$ 是三角范畴，则当
$(\mathcal{D}', \mathcal{T}')$ 是预三角子范畴且
$(\mathcal{D}', [\ ], \mathcal{T}')$ 是三角范畴时，称其为
{\it 三角子范畴}。
\end{definition}

\noindent
在这种情形下，嵌入函子 $\mathcal{D}' \to \mathcal{D}$ 是正合函子，
其中 $\xi_X : X[1] \to X[1]$ 取为 $X[1]$ 上的恒等态射。

\medskip\noindent
我们将在引理 \ref{derived-lemma-composition-zero} 中看到：在预三角范畴的特异三角
$(X, Y, Z, f, g, h)$ 中，复合 $g \circ f : X \to Z$ 为零。
因此序列 (\ref{derived-equation-rotate}) 是复形。同调函子就是把这一复形变为长正合序列的函子。

\begin{definition}
\label{derived-definition-homological}
设 $\mathcal{D}$ 为预三角范畴，$\mathcal{A}$ 为阿贝尔范畴。
加性函子 $H : \mathcal{D} \to \mathcal{A}$ 称为{\it 同调的}，
如果对每个特异三角 $(X, Y, Z, f, g, h)$，序列
$$
H(X) \to H(Y) \to H(Z)
$$
在阿贝尔范畴 $\mathcal{A}$ 中正合。加性函子
$H : \mathcal{D}^{opp} \to \mathcal{A}$ 称为{\it 上同调的}，
如果相应的函子 $\mathcal{D} \to \mathcal{A}^{opp}$ 是同调的。
\end{definition}

\noindent
若 $H : \mathcal{D} \to \mathcal{A}$ 是同调函子，我们常记
$H^n(X) = H(X[n])$，从而 $H(X) = H^0(X)$。上面对 TR2 的讨论表明，
特异三角 $(X, Y, Z, f, g, h)$ 决定长正合序列
\begin{equation}
\label{derived-equation-long-exact-cohomology-sequence}
\xymatrix@C=3pc{
H^{-1}(Z) \ar[r]^{H(h[-1])} &
H^0(X) \ar[r]^{H(f)} &
H^0(Y) \ar[r]^{H(g)} &
H^0(Z) \ar[r]^{H(h)} &
H^1(X)
}
\end{equation}
它称为与该特异三角及同调函子相伴的{\it 长正合序列}。
如记号所示，我们不在长正合序列的态射上使用任何符号。其副作用是：
与一个特异三角的旋转（TR2）相伴的长正合序列中的映射，与上述序列中的映射
会相差某些符号。

\begin{definition}
\label{derived-definition-delta-functor}
设 $\mathcal{A}$ 为阿贝尔范畴，$\mathcal{D}$ 为三角范畴。
一个{\it 从 $\mathcal{A}$ 到 $\mathcal{D}$ 的 $\delta$-函子}由函子
$G : \mathcal{A} \to \mathcal{D}$ 以及如下规则给出：对每个短正合序列
$$
0 \to A \xrightarrow{a} B \xrightarrow{b} C \to 0
$$
指定态射 $\delta = \delta_{A \to B \to C} : G(C) \to G(A)[1]$，
使得
\begin{enumerate}
\item 对上述每个短正合序列，三角
$(G(A), G(B), G(C), G(a), G(b), \delta_{A \to B \to C})$
都是 $\mathcal{D}$ 的特异三角；并且
\item 对短正合序列的每个态射
$(A \to B \to C) \to (A' \to B' \to C')$，下图交换：
$$
\xymatrix{
G(C) \ar[d] \ar[rr]_{\delta_{A \to B \to C}} & &
G(A)[1] \ar[d] \\
G(C') \ar[rr]^{\delta_{A' \to B' \to C'}} & &
G(A')[1]
}
$$
\end{enumerate}
在这种情形下，称
$(G(A), G(B), G(C), G(a), G(b), \delta_{A \to B \to C})$
为给定 $\delta$-函子下该短正合序列的{\it 像}。
\end{definition}

\noindent
注意，$\delta$-函子带有附加结构。严格说来，称某个给定函子
$\mathcal{A} \to \mathcal{D}$ 为 $\delta$-函子并无意义，不过我们仍会经常这样说。










\section{三角范畴的初步结果}
\label{derived-section-elementary-results}

\noindent
本节的大多数结果都是对预三角范畴证明的，因而自然也适用于任何三角范畴。

\begin{lemma}
\label{derived-lemma-composition-zero}
设 $\mathcal{D}$ 为预三角范畴，$(X, Y, Z, f, g, h)$ 为特异三角。则
$g \circ f = 0$、$h \circ g = 0$ 且 $f[1] \circ h = 0$。
\end{lemma}

\begin{proof}
由 TR1 可知 $(X, X, 0, 1, 0, 0)$ 是特异三角。对下图应用 TR3：
$$
\xymatrix{
X \ar[r] \ar[d]^1 &
X \ar[r] \ar[d]^f &
0 \ar[r] \ar@{-->}[d] &
X[1] \ar[d]^{1[1]} \\
X \ar[r]^f &
Y \ar[r]^g &
Z \ar[r]^h &
X[1]
}
$$
虚线箭头当然是零映射。因此图的交换性给出 $g \circ f = 0$。
其余情形只需旋转三角，即应用 TR2。
\end{proof}

\begin{lemma}
\label{derived-lemma-representable-homological}
设 $\mathcal{D}$ 为预三角范畴。对 $\mathcal{D}$ 的任一对象 $W$，
函子 $\Hom_\mathcal{D}(W, -)$ 是同调函子，而函子
$\Hom_\mathcal{D}(-, W)$ 是上同调函子。
\end{lemma}

\begin{proof}
考虑特异三角 $(X, Y, Z, f, g, h)$。我们已经看到
$g \circ f = 0$，见引理 \ref{derived-lemma-composition-zero}。
设 $a : W \to Y$ 为满足 $g \circ a = 0$ 的态射。于是得到交换图
$$
\xymatrix{
W \ar[r]_1 \ar@{..>}[d]^b &
W \ar[r] \ar[d]^a &
0 \ar[r] \ar[d]^0 &
W[1] \ar@{..>}[d]^{b[1]} \\
X \ar[r] & Y \ar[r] & Z \ar[r] & X[1]
}
$$
两行都是特异三角（上行使用 TR1）。因此可补出虚线箭头 $b$
（先用 TR2 旋转，再应用 TR3，最后旋转回来）。这就证明了引理。
\end{proof}

\begin{lemma}
\label{derived-lemma-third-isomorphism-triangle}
设 $\mathcal{D}$ 为预三角范畴。设
$$
(a, b, c) : (X, Y, Z, f, g, h) \to (X', Y', Z', f', g', h')
$$
为特异三角的态射。若 $a, b, c$ 中有两个是同构，则第三个也是同构。
\end{lemma}

\begin{proof}
假设 $a$ 和 $c$ 是同构。对 $\mathcal{D}$ 的任一对象 $W$，记
$H_W( - ) = \Hom_\mathcal{D}(W, -)$。于是得到阿贝尔群的交换图
$$
\xymatrix{
H_W(Z[-1]) \ar[r] \ar[d] &
H_W(X) \ar[r] \ar[d] &
H_W(Y) \ar[r] \ar[d] &
H_W(Z) \ar[r] \ar[d] &
H_W(X[1]) \ar[d] \\
H_W(Z'[-1]) \ar[r] &
H_W(X') \ar[r] &
H_W(Y') \ar[r] &
H_W(Z') \ar[r] &
H_W(X'[1])
}
$$
由假设，右边两个和左边两个竖直箭头都是双射。由引理
\ref{derived-lemma-representable-homological}，$H_W$ 是同调函子；再由五引理
（同调，引理 \ref{homology-lemma-five-lemma}），中间的竖直箭头是同构。
于是由 Yoneda 引理（见范畴，引理 \ref{categories-lemma-yoneda}），
$b$ 是同构。旋转三角（使用 TR2）即可得到其他情形。
\end{proof}

\begin{remark}
\label{derived-remark-special-triangles}
设 $\mathcal{D}$ 为带有定义 \ref{derived-definition-triangle} 中平移函子 $[n]$
的加性范畴。若对 $\mathcal{D}$ 的每个对象 $W$，阿贝尔群长序列
$$
\ldots \to
\Hom_\mathcal{D}(W, X) \to
\Hom_\mathcal{D}(W, Y) \to
\Hom_\mathcal{D}(W, Z) \to
\Hom_\mathcal{D}(W, X[1]) \to \ldots
$$
正合，则称三角 $(X, Y, Z, f, g, h)$ 为{\it 特殊的}\footnote{这不是标准术语。}。
引理 \ref{derived-lemma-third-isomorphism-triangle} 的证明表明：若
$$
(a, b, c) : (X, Y, Z, f, g, h) \to (X', Y', Z', f', g', h')
$$
是特殊三角的态射，且 $a, b, c$ 中有两个是同构，则第三个也是同构。
对{\it 余特殊}三角也有对偶命题；所谓余特殊三角，是指应用函子
$\Hom_\mathcal{D}(-, W)$ 后变成长正合序列的三角。因此特异三角既特殊又余特殊，
但一般而言，（余）特殊三角远多于特异三角。
\end{remark}

\begin{lemma}
\label{derived-lemma-third-map-square-zero}
设 $\mathcal{D}$ 为预三角范畴。设
$$
(0, b, 0), (0, b', 0) : (X, Y, Z, f, g, h) \to (X, Y, Z, f, g, h)
$$
为某特异三角的自同态。则 $bb' = 0$。
\end{lemma}

\begin{proof}
考虑图
$$
\xymatrix{
X \ar[r] \ar[d]^0 &
Y \ar[r] \ar[d]^{b, b'} \ar@{..>}[ld]^\alpha &
Z \ar[r] \ar[d]^0 \ar@{..>}[ld]^\beta &
X[1] \ar[d]^0 \\
X \ar[r] & Y \ar[r] & Z \ar[r] & X[1]
}
$$
应用引理 \ref{derived-lemma-representable-homological}，可找到虚线箭头
$\alpha$ 和 $\beta$，使得 $b' = f \circ \alpha$ 且
$b = \beta \circ g$。于是
$bb' = \beta \circ g \circ f \circ \alpha = 0$，因为由引理
\ref{derived-lemma-composition-zero} 有 $g \circ f = 0$。
\end{proof}

\begin{lemma}
\label{derived-lemma-third-map-idempotent}
设 $\mathcal{D}$ 为预三角范畴，$(X, Y, Z, f, g, h)$ 为特异三角。
若
$$
\xymatrix{
Z \ar[r]_h \ar[d]_c & X[1] \ar[d]^{a[1]} \\
Z \ar[r]^h & X[1]
}
$$
交换，且 $a^2 = a$、$c^2 = c$，则存在满足 $b^2 = b$ 的态射
$b : Y \to Y$，使得 $(a, b, c)$ 是三角 $(X, Y, Z, f, g, h)$ 的自同态。
\end{lemma}

\begin{proof}
由 TR3，存在态射 $b'$，使得 $(a, b', c)$ 是
$(X, Y, Z, f, g, h)$ 的自同态。于是 $(0, (b')^2 - b', 0)$ 也是自同态。
由引理 \ref{derived-lemma-third-map-square-zero}，$(b')^2 - b'$ 的平方为零。
令 $b = b' - (2b' - 1)((b')^2 - b') = 3(b')^2 - 2(b')^3$。
直接计算表明 $(a, b, c)$ 是自同态，且
$b^2 - b = (4(b')^2 - 4b' - 3)((b')^2 - b')^2 = 0$。
\end{proof}

\begin{lemma}
\label{derived-lemma-cone-triangle-unique-isomorphism}
设 $\mathcal{D}$ 为预三角范畴，$f : X \to Y$ 为 $\mathcal{D}$ 的态射。
存在特异三角 $(X, Y, Z, f, g, h)$，而且它在三角的（不唯一）同构意义下唯一。
更精确地，给定另一个这样的特异三角 $(X, Y, Z', f, g', h')$，存在同构
$$
(1, 1, c) : (X, Y, Z, f, g, h) \longrightarrow (X, Y, Z', f, g', h')
$$
\end{lemma}

\begin{proof}
存在性由 TR1 得到。在同构意义下的唯一性由 TR3 和引理
\ref{derived-lemma-third-isomorphism-triangle} 得到。
\end{proof}

\begin{lemma}
\label{derived-lemma-uniqueness-third-arrow}
设 $\mathcal{D}$ 为预三角范畴。设
$$
(a, b, c) : (X, Y, Z, f, g, h) \to (X', Y', Z', f', g', h')
$$
为特异三角的态射。若下列条件之一成立：
\begin{enumerate}
\item $\Hom(Y, X') = 0$；
\item $\Hom(Z, Y') = 0$；
\item $\Hom(X, X') = \Hom(Z, X') = 0$；
\item $\Hom(Z, X') = \Hom(Z, Z') = 0$；或
\item $\Hom(X[1], Z') = \Hom(Z, X') = 0$，
\end{enumerate}
则在使 $(a, b, c)$ 成为三角态射的所有 $Y \to Y'$ 态射中，$b$ 是唯一的。
\end{lemma}

\begin{proof}
若还有第二个三角态射 $(a, b', c)$，则 $(0, b - b', 0)$ 是三角态射。
因此只需证明：若态射 $b : Y \to Y'$ 使得
$X \to Y \to Y'$ 与 $Y \to Y' \to Z'$ 均为零，则其必为 $0$。
下文将不再说明地使用引理 \ref{derived-lemma-representable-homological}。
特别地，条件 (3) 蕴含条件 (1)。在条件 (1) 下，若复合
$g' \circ b : Y \to Y' \to Z'$ 为零，则 $b$ 可提升为态射
$Y \to X'$，而后者必为零。这证明了 (1)。

\medskip\noindent
(2) 与 (4) 的证明与此论证对偶。

\medskip\noindent
假设 (5)。考虑图
$$
\xymatrix{
X \ar[r]_f \ar[d]^0 &
Y \ar[r]_g \ar[d]^b &
Z \ar[r]_h \ar[d]^0 \ar@{..>}[ld]^\epsilon &
X[1] \ar[d]^0 \\
X' \ar[r]^{f'} &
Y' \ar[r]^{g'} &
Z' \ar[r]^{h'} &
X'[1]
}
$$
可选择 $\epsilon$ 使 $b = \epsilon \circ g$。于是
$g' \circ \epsilon \circ g = 0$，故对某个
$\delta \in \Hom(X[1], Z')$ 有 $g' \circ \epsilon = \delta \circ h$。
由于 $\Hom(X[1], Z') = 0$，得到 $g' \circ \epsilon = 0$。
因而对某个 $\gamma \in \Hom(Z, X')$ 有
$\epsilon = f' \circ \gamma$。又因 $\Hom(Z, X') = 0$，可知
$\epsilon = 0$，从而如愿得到 $b = 0$。
\end{proof}

\begin{lemma}
\label{derived-lemma-third-object-zero}
设 $\mathcal{D}$ 为预三角范畴，$f : X \to Y$ 为 $\mathcal{D}$ 的态射。
下列条件等价：
\begin{enumerate}
\item $f$ 是同构；
\item $(X, Y, 0, f, 0, 0)$ 是特异三角；
\item 对任一特异三角 $(X, Y, Z, f, g, h)$，都有 $Z = 0$。
\end{enumerate}
\end{lemma}

\begin{proof}
由 TR1，三角 $(X, X, 0, 1, 0, 0)$ 是特异三角。设
$(X, Y, Z, f, g, h)$ 为特异三角。由 TR3，存在特异三角的态射
$(1, f, 0) : (X, X, 0) \to (X, Y, Z)$。若 $f$ 是同构，则由引理
\ref{derived-lemma-third-isomorphism-triangle}，$(1, f, 0)$ 是三角的同构，
故 $Z = 0$。反过来，若 $Z = 0$，则 $(1, f, 0)$ 同样是三角的同构，
因而 $f$ 是同构。
\end{proof}

\begin{lemma}
\label{derived-lemma-direct-sum-triangles}
设 $\mathcal{D}$ 为预三角范畴。设 $(X, Y, Z, f, g, h)$ 和
$(X', Y', Z', f', g', h')$ 为三角。下列条件等价：
\begin{enumerate}
\item $(X \oplus X', Y \oplus Y', Z \oplus Z',
f \oplus f', g \oplus g', h \oplus h')$
是特异三角；
\item $(X, Y, Z, f, g, h)$ 和 $(X', Y', Z', f', g', h')$ 都是特异三角。
\end{enumerate}
\end{lemma}

\begin{proof}
假设 (2)。由 TR1，可选取特异三角
$(X \oplus X', Y \oplus Y', Q, f \oplus f', g'', h'')$。
由 TR3，可找到特异三角的态射
$(X, Y, Z, f, g, h) \to
(X \oplus X', Y \oplus Y', Q, f \oplus f', g'', h'')$
以及
$(X', Y', Z', f', g', h') \to
(X \oplus X', Y \oplus Y', Q, f \oplus f', g'', h'')$。
取这些态射的直和，得到三角的态射
$$
\xymatrix{
(X \oplus X', Y \oplus Y', Z \oplus Z',
f \oplus f', g \oplus g', h \oplus h')
\ar[d]^{(1, 1, c)} \\
(X \oplus X', Y \oplus Y', Q, f \oplus f', g'', h'').
}
$$
用注 \ref{derived-remark-special-triangles} 的术语说，这是特殊三角之间的映射
（因为特殊三角的直和仍然特殊），因此 $c$ 是同构。故 (1) 成立。

\medskip\noindent
假设 (1)。我们证明 $(X, Y, Z, f, g, h)$ 是特异三角。先注意到
$(X, Y, Z, f, g, h)$ 是特殊三角
（使用注 \ref{derived-remark-special-triangles} 的术语），因为它是特异、因而特殊的三角
$(X \oplus X', Y \oplus Y', Z \oplus Z',
f \oplus f', g \oplus g', h \oplus h')$ 的直和因子。
用 TR1 取特异三角 $(X, Y, Q, f, g'', h'')$。由 TR3，存在特异三角的态射
$(X \oplus X', Y \oplus Y', Z \oplus Z',
f \oplus f', g \oplus g', h \oplus h') \to (X, Y, Q, f, g'', h'')$。
将它与嵌入映射复合，得到三角的态射
$$
(1, 1, c) :
(X, Y, Z, f, g, h)
\longrightarrow
(X, Y, Q, f, g'', h'')
$$
由注 \ref{derived-remark-special-triangles}，$c$ 是同构，因此 (2) 成立。
\end{proof}

\begin{lemma}
\label{derived-lemma-split}
设 $\mathcal{D}$ 为预三角范畴，$(X, Y, Z, f, g, h)$ 为特异三角。
\begin{enumerate}
\item 若 $h = 0$，则 $g$ 存在右逆 $s : Z \to Y$。
\item 对 $g$ 的任一右逆 $s : Z \to Y$，映射
$f \oplus s : X \oplus Z \to Y$ 是同构。
\item 对 $\mathcal{D}$ 的任意对象 $X', Z'$，三角
$(X', X' \oplus Z', Z', (1, 0), (0, 1), 0)$ 是特异三角。
\end{enumerate}
\end{lemma}

\begin{proof}
为证明 (1)，使用引理 \ref{derived-lemma-representable-homological} 所给出的正合性
$\Hom_\mathcal{D}(Z, Y) \to \Hom_\mathcal{D}(Z, Z) \to
\Hom_\mathcal{D}(Z, X[1])$。
同理，若 $s$ 如 (2) 所述，则 $h = 0$，且序列
$$
0 \to \Hom_\mathcal{D}(W, X) \to \Hom_\mathcal{D}(W, Y)
\to \Hom_\mathcal{D}(W, Z) \to 0
$$
是分裂正合的（由 $s : Z \to Y$ 分裂）。因此由 Yoneda 引理，
$X \oplus Z \to Y$ 是同构。最后一个断言由 TR1 和引理
\ref{derived-lemma-direct-sum-triangles} 得出。
\end{proof}

\begin{lemma}
\label{derived-lemma-when-split}
设 $\mathcal{D}$ 为预三角范畴，$f : X \to Y$ 为 $\mathcal{D}$ 的态射。
下列条件等价：
\begin{enumerate}
\item $f$ 有核；
\item $f$ 有余核；
% P02-E0019: frozen source has the grammatical typo “is the isomorphic”.
\item 对 $\mathcal{D}$ 的某些对象 $K, Z, Q$，$f$ 同构于投影和余投影的复合
$K \oplus Z \to Z \to Z \oplus Q$。
\end{enumerate}
\end{lemma}

\begin{proof}
任何同构于形如 $X' \oplus Z \to Z \oplus Y'$ 之映射的态射都有核和余核。
故 (3) $\Rightarrow$ (1), (2)。下面证明 (1) $\Rightarrow$ (3)。
先设 $f : X \to Y$ 是单态射，即其核为零。由 TR1，存在特异三角
$(X, Y, Z, f, g, h)$。由引理 \ref{derived-lemma-composition-zero}，复合
$f \circ h[-1] = 0$。由于 $f$ 是单态射，得到 $h[-1] = 0$，从而
$h = 0$。于是引理 \ref{derived-lemma-split} 给出 $Y = X \oplus Z$，即 (3) 成立。
再假设 $f$ 有核 $K$。由于 $K \to X$ 是单态射，可得
$X = K \oplus X'$，并且 $f|_{X'} : X' \to Y$ 是单态射。
因此 $Y = X' \oplus Y'$，结论成立。
(2) $\Rightarrow$ (3) 与此对偶。
\end{proof}

\begin{lemma}
\label{derived-lemma-products-sums-shifts-triangles}
设 $\mathcal{D}$ 为预三角范畴，$I$ 为集合。
\begin{enumerate}
\item 设 $X_i$，$i \in I$，为 $\mathcal{D}$ 的一族对象。
\begin{enumerate}
\item 若 $\prod X_i$ 存在，则 $(\prod X_i)[1] = \prod X_i[1]$。
\item 若 $\bigoplus X_i$ 存在，则 $(\bigoplus X_i)[1] = \bigoplus X_i[1]$。
\end{enumerate}
\item 设 $X_i \to Y_i \to Z_i \to X_i[1]$ 为 $\mathcal{D}$ 的一族特异三角。
\begin{enumerate}
\item 若 $\prod X_i$、$\prod Y_i$、$\prod Z_i$ 存在，则
$\prod X_i \to \prod Y_i \to \prod Z_i \to \prod X_i[1]$
是特异三角。
\item 若 $\bigoplus X_i$、$\bigoplus Y_i$、$\bigoplus Z_i$ 存在，则
$\bigoplus X_i \to \bigoplus Y_i \to \bigoplus Z_i \to \bigoplus X_i[1]$
是特异三角。
\end{enumerate}
\end{enumerate}
\end{lemma}

\begin{proof}
(1) 成立，因为 $[1]$ 是 $\mathcal{D}$ 的自等价，而直和与积都由范畴结构定义。
我们来证明 (2)(a)。选取特异三角
$\prod X_i \to \prod Y_i \to Z \to \prod X_i[1]$。
对每个 $j$，可用 TR3 选取态射 $p_j : Z \to Z_j$，使其与投影映射
$\prod X_i \to X_j$ 和 $\prod Y_i \to Y_j$ 一起构成特异三角的态射。
由积的定义，得到映射 $\prod p_i : Z \to \prod Z_i$，它构成从上述
特异三角到由积组成之三角的三角态射。注意，用注
\ref{derived-remark-special-triangles} 的术语说，``积''三角
$\prod X_i \to \prod Y_i \to \prod Z_i \to \prod X_i[1]$
是特殊的，因为阿贝尔群正合序列的积仍正合。因此注
\ref{derived-remark-special-triangles} 表明该三角态射是同构，再由 TR1 得出结论。
(2)(b) 的证明与此对偶。
\end{proof}

\begin{lemma}
\label{derived-lemma-projectors-have-images-triangulated}
设 $\mathcal{D}$ 为预三角范畴。若 $\mathcal{D}$ 有可数积，则
$\mathcal{D}$ 是 Karoubi 范畴。若 $\mathcal{D}$ 有可数余积，则
$\mathcal{D}$ 是 Karoubi 范畴。
\end{lemma}

\begin{proof}
假设 $\mathcal{D}$ 有可数积。由同调，引理
\ref{homology-lemma-projectors-have-images}，只需验证有右逆的态射都有核。
任何有右逆的态射都是满态射，故由引理 \ref{derived-lemma-when-split} 有核。
第二个断言与第一个对偶。
\end{proof}

\noindent
下面的引理使证明一个预三角范畴是三角范畴略为容易。

\begin{lemma}
\label{derived-lemma-easier-axiom-four}
设 $\mathcal{D}$ 为预三角范畴。为证明 TR4，只需证明：给定任意一对可复合态射
$f : X \to Y$ 与 $g : Y \to Z$，存在
\begin{enumerate}
\item 同构 $i : X' \to X$、$j : Y' \to Y$ 和 $k : Z' \to Z$；进而令
$f' = j^{-1}fi : X' \to Y'$、$g' = k^{-1}gj : Y' \to Z'$ 后，存在
\item 特异三角
$(X', Y', Q_1, f', p_1, d_1)$、
$(X', Z', Q_2, g' \circ f', p_2, d_2)$ 和
$(Y', Z', Q_3, g', p_3, d_3)$，
使 TR4 的断言成立。
\end{enumerate}
\end{lemma}

\begin{proof}
按我们对特异三角及其同构的定义，用 $X', Y', Z'$ 替换 $X, Y, Z$ 无关紧要。
由引理 \ref{derived-lemma-cone-triangle-unique-isomorphism}，特异三角
$(X', Y', Q_1, f', p_1, d_1)$、
$(X', Z', Q_2, g' \circ f', p_2, d_2)$ 和
$(Y', Z', Q_3, g', p_3, d_3)$ 在同构意义下唯一，故引理成立。
\end{proof}

\begin{lemma}
\label{derived-lemma-triangulated-subcategory}
设 $\mathcal{D}$ 为预三角范畴。假设 $\mathcal{D}'$ 是
$\mathcal{D}$ 的加性满子范畴。下列条件等价：
\begin{enumerate}
\item 存在一组三角 $\mathcal{T}'$，使得
$(\mathcal{D}', \mathcal{T}')$ 是 $\mathcal{D}$ 的预三角子范畴；
\item $\mathcal{D}'$ 在 $[1]$ 下保持，
$[1] : \mathcal{D}' \to \mathcal{D}'$ 是自等价，并且给定
$\mathcal{D}'$ 中的任一态射 $f : X \to Y$，在 $\mathcal{D}$ 中存在特异三角
$(X, Y, Z, f, g, h)$，其中 $Z$ 同构于 $\mathcal{D}'$ 的某个对象。
\end{enumerate}
在这种情形下，(1) 中的 $\mathcal{T}'$ 正是 $\mathcal{D}$ 的所有满足
$X, Y, Z \in \Ob(\mathcal{D}')$ 的特异三角
$(X, Y, Z, f, g, h)$ 之集。最后，若 $\mathcal{D}$ 是三角范畴，
则 (1) 和 (2) 还等价于
\begin{enumerate}
\item[(3)] $\mathcal{D}'$ 是三角子范畴。
\end{enumerate}
\end{lemma}

\begin{proof}
略。
\end{proof}
\begin{lemma}
\label{derived-lemma-exact-functor-additive}
预三角范畴之间的正合函子是加性函子。
\end{lemma}

\begin{proof}
设 $F : \mathcal{D} \to \mathcal{D}'$ 为预三角范畴之间的正合函子。
由于 $(0, 0, 0, 1_0, 1_0, 0)$ 是 $\mathcal{D}$ 的特异三角，三角
$$
(F(0), F(0), F(0), 1_{F(0)}, 1_{F(0)}, F(0))
$$
在 $\mathcal{D}'$ 中是特异三角。由引理 \ref{derived-lemma-composition-zero}，
这意味着 $1_{F(0)} \circ 1_{F(0)}$ 为零。因此 $F(0)$ 是
$\mathcal{D}'$ 的零对象。这也意味着 $F$ 把任一零态射映为零
（因为加性范畴中的态射为零，当且仅当它经过零对象分解）。接着，由引理
\ref{derived-lemma-split} 第 (3) 部分，
$(X, X \oplus Y, Y, (1, 0), (0, 1), 0)$ 是特异三角；故
$(F(X), F(X \oplus Y), F(Y), F(1, 0), F(0, 1), 0)$
也是特异三角。由引理 \ref{derived-lemma-split} 第 (2) 部分，映射
$F(X) \oplus F(Y) \to F(X \oplus Y)$ 是同构。最后应用同调，引理
\ref{homology-lemma-additive-functor} 即可。
\end{proof}

\begin{lemma}
\label{derived-lemma-exact-equivalence}
设 $F : \mathcal{D} \to \mathcal{D}'$ 为预三角范畴之间的全忠实正合函子。
则 $\mathcal{D}$ 的三角 $(X, Y, Z, f, g, h)$ 是特异三角，当且仅当
$(F(X), F(Y), F(Z), F(f), F(g), F(h))$ 是 $\mathcal{D}'$ 中的特异三角。
\end{lemma}

\begin{proof}
``仅当''部分显然。假设 $(F(X), F(Y), F(Z))$ 是 $\mathcal{D}'$ 中的特异三角。
在 $\mathcal{D}$ 中选取特异三角 $(X, Y, Z', f, g', h')$。
由引理 \ref{derived-lemma-cone-triangle-unique-isomorphism}，存在三角同构
$$
(1, 1, c') : (F(X), F(Y), F(Z)) \longrightarrow (F(X), F(Y), F(Z')).
$$
由于 $F$ 全忠实，存在态射 $c : Z \to Z'$，使得 $F(c) = c'$。
于是 $(1, 1, c)$ 是 $(X, Y, Z)$ 与 $(X, Y, Z')$ 之间的同构。
因此由 TR1，$(X, Y, Z)$ 是特异三角。
\end{proof}

\begin{lemma}
\label{derived-lemma-composition-exact}
设 $\mathcal{D}, \mathcal{D}', \mathcal{D}''$ 为预三角范畴，
$F : \mathcal{D} \to \mathcal{D}'$ 和
$F' : \mathcal{D}' \to \mathcal{D}''$ 为正合函子。则
$F' \circ F$ 是正合函子。
\end{lemma}

\begin{proof}
略。
\end{proof}

\begin{lemma}
\label{derived-lemma-exact-compose-homological-functor}
设 $\mathcal{D}$ 为预三角范畴，$\mathcal{A}$ 为阿贝尔范畴，
$H : \mathcal{D} \to \mathcal{A}$ 为同调函子。
\begin{enumerate}
\item 设 $\mathcal{D}'$ 为预三角范畴，
$F : \mathcal{D}' \to \mathcal{D}$ 为正合函子。则复合
$H \circ F$ 也是同调函子。
\item 设 $\mathcal{A}'$ 为阿贝尔范畴，
$G : \mathcal{A} \to \mathcal{A}'$ 为正合函子。则
$G \circ H$ 也是同调函子。
\end{enumerate}
\end{lemma}

\begin{proof}
略。
\end{proof}

\begin{lemma}
\label{derived-lemma-exact-compose-delta-functor}
设 $\mathcal{D}$ 为三角范畴，$\mathcal{A}$ 为阿贝尔范畴，
$G : \mathcal{A} \to \mathcal{D}$ 为 $\delta$-函子。
\begin{enumerate}
\item 设 $\mathcal{D}'$ 为三角范畴，
$F : \mathcal{D} \to \mathcal{D}'$ 为正合函子。则复合
$F \circ G$ 也是 $\delta$-函子。
\item 设 $\mathcal{A}'$ 为阿贝尔范畴，
$H : \mathcal{A}' \to \mathcal{A}$ 为正合函子。则
$G \circ H$ 也是 $\delta$-函子。
\end{enumerate}
\end{lemma}

\begin{proof}
略。
\end{proof}

\begin{lemma}
\label{derived-lemma-compose-delta-functor-homological}
设 $\mathcal{D}$ 为三角范畴，$\mathcal{A}$ 和 $\mathcal{B}$ 为阿贝尔范畴。
设 $G : \mathcal{A} \to \mathcal{D}$ 为 $\delta$-函子，
$H : \mathcal{D} \to \mathcal{B}$ 为同调函子。假设对
$\mathcal{A}$ 中所有 $A$ 都有 $H^{-1}(G(A)) = 0$。则族
$$
\{H^n \circ G, H^n(\delta_{A \to B \to C})\}_{n \geq 0}
$$
是从 $\mathcal{A} \to \mathcal{B}$ 的 $\delta$-函子；见同调，定义
\ref{homology-definition-cohomological-delta-functor}。
\end{lemma}

\begin{proof}
该记号表示如下内容。若
$0 \to A \xrightarrow{a} B \xrightarrow{b} C \to 0$ 是
$\mathcal{A}$ 中的短正合序列，则
$$
\delta = \delta_{A \to B \to C} : G(C) \to G(A)[1]
$$
是 $\mathcal{D}$ 中的态射，使得
$(G(A), G(B), G(C), a, b, \delta)$ 是特异三角；见定义
\ref{derived-definition-delta-functor}。于是
$H^n(\delta) : H^n(G(C)) \to H^n(G(A)[1]) = H^{n + 1}(G(A))$
显然关于短正合序列具有函子性。最后，长正合上同调序列
(\ref{derived-equation-long-exact-cohomology-sequence}) 与
$H^{-1}(G(C))$ 的消失合在一起，给出 $\mathcal{B}$ 中的长正合序列
$$
0 \to H^0(G(A)) \to H^0(G(B)) \to H^0(G(C))
\xrightarrow{H^0(\delta)} H^1(G(A)) \to \ldots
$$
如所需。
\end{proof}

\noindent
下面结果的证明使用 TR4。

\begin{proposition}
\label{derived-proposition-9}
设 $\mathcal{D}$ 为三角范畴。任何交换图
$$
\xymatrix{
X \ar[r] \ar[d] & Y \ar[d] \\
X' \ar[r] & Y'
}
$$
都可扩张为图
$$
\xymatrix{
X \ar[r] \ar[d] & Y \ar[r] \ar[d] & Z \ar[r] \ar[d] & X[1] \ar[d] \\
X' \ar[r] \ar[d] & Y' \ar[r] \ar[d] & Z' \ar[r] \ar[d] & X'[1] \ar[d] \\
X'' \ar[r] \ar[d] & Y'' \ar[r] \ar[d] & Z'' \ar[r] \ar[d] & X''[1] \ar[d] \\
X[1] \ar[r] & Y[1] \ar[r] & Z[1] \ar[r] & X[2]
}
$$
其中除右下方形反交换外，所有方形都交换。此外，每一行和每一列都是特异三角。
最后，底行（相应地，右列）中的态射，是对顶行（相应地，左列）的态射应用
$[1]$ 得到的。
\end{proposition}

\begin{proof}
为使证明易读，本证明中不写出箭头。选取特异三角
$(X, Y, Z)$、$(X', Y', Z')$、$(X, X', X'')$、$(Y, Y', Y'')$ 和
$(X, Y', A)$。注意，态射 $X \to Y'$ 既等于复合
$X \to Y \to Y'$，也等于复合 $X \to X' \to Y'$。因此可按 TR4 找到态射
\begin{enumerate}
\item $a : Z \to A$ 和 $b : A \to Y''$；以及
\item $a' : X'' \to A$ 和 $b' : A \to Z'$。
\end{enumerate}
记 $c : Y'' \to Z[1]$ 为复合 $Y'' \to Y[1] \to Z[1]$，并记
$c' : Z' \to X''[1]$ 为复合 $Z' \to X'[1] \to X''[1]$。
此次应用 TR4 的结论是
\begin{enumerate}
\item $(Z, A, Y'', a, b, c)$、$(X'', A, Z', a', b', c')$
是特异三角；
\item $(X, Y, Z) \to (X, Y', A)$、
$(X, Y', A) \to (Y, Y', Y'')$、
$(X, X', X'') \to (X, Y', A)$、
$(X, Y', A) \to (X', Y', Z')$ 是三角的态射。
\end{enumerate}
首先使用 $(X, X', X'') \to (X, Y', A)$ 和
$(X, Y', A) \to (Y, Y', Y'')$ 是三角态射这一事实，可知下列两图中的第一个交换：
$$
\vcenter{
\xymatrix{
X' \ar[r] \ar[d] & Y' \ar[d] \\
X'' \ar[r]^{b \circ a'} \ar[d] & Y'' \ar[d] \\
X[1] \ar[r] & Y[1]
}
}
\quad\text{且}\quad
\vcenter{
\xymatrix{
Y \ar[r] \ar[d] & Z \ar[d]^{b' \circ a} \ar[r] & X[1] \ar[d] \\
Y' \ar[r] & Z' \ar[r] & X'[1]
}
}
$$
使用 $(X, Y, Z) \to (X, Y', A)$ 和
$(X, Y', A) \to (X', Y', Z')$ 是三角态射这一事实，第二个图也交换。
此时，选取以映射 $b \circ a' : X'' \to Y''$ 开始的特异三角
$(X'', Y'' , Z'')$。

\medskip\noindent
接着，对态射 $X'' \to A \to Y''$ 及三角
$(X'', A, Z', a', b', c')$、$(X'', Y'', Z'')$ 和
$(A, Y'', Z[1], b, c , -a[1])$ 再应用一次 TR4，得到态射
$a'' : Z' \to Z''$ 和 $b'' : Z'' \to Z[1]$。
于是 $(Z', Z'', Z[1], a'', b'', - b'[1] \circ a[1])$ 是特异三角，
因而 $(Z, Z', Z'', -b' \circ a, a'', -b'')$ 也是特异三角，进而
$(Z, Z', Z'', b' \circ a, a'', b'')$ 也是特异三角。
此外，$(X'', A, Z') \to (X'', Y'', Z'')$ 和
$(X'', Y'', Z'') \to (A, Y'', Z[1], b, c , -a[1])$
都是三角的态射。至此，所有特异三角和所有态射都已定义，只需验证若干交换关系。

\medskip\noindent
为证明图的中间方形交换，注意箭头 $Y' \to Z'$ 分解为
$Y' \to A \to Z'$，因为 $(X, Y', A) \to (X', Y', Z')$
是三角的态射。同理，态射 $Y' \to Y''$ 分解为
$Y' \to A \to Y''$，因为 $(X, Y', A) \to (Y, Y', Y'')$
是三角的态射。因此，边为 $(A, Z', Z'', Y'')$ 的方形交换，因为
$(X'', A, Z') \to (X'', Y'', Z'')$ 是三角的态射（由 TR4）。
边为 $(Y'', Z'', Y[1], Z[1])$ 的方形交换，因为
$(X'', Y'', Z'') \to (A, Y'', Z[1], b, c , -a[1])$
是三角的态射，且 $c : Y'' \to Z[1]$ 是复合
$Y'' \to Y[1] \to Z[1]$。边为 $(Z', X'[1], X''[1], Z'')$ 的方形交换，
因为 $(X'', A, Z') \to (X'', Y'', Z'')$ 是三角的态射，且
$c' : Z' \to X''[1]$ 是复合 $Z' \to X'[1] \to X''[1]$。
最后必须证明边为 $(Z'', X''[1], Z[1], X[2])$ 的方形反交换。
这是因为 $(X'', Y'', Z'') \to (A, Y'', Z[1], b, c , -a[1])$
是三角的态射，证明完毕。
\end{proof}







\section{三角范畴的局部化}
\label{derived-section-localization}

\noindent
为从复形的同伦范畴构造导出范畴，我们将使用局部化过程。

\begin{definition}
\label{derived-definition-localization}
设 $\mathcal{D}$ 为预三角范畴。若乘法系 $S$ 满足下列两个条件，
则称它{\it 与三角结构相容}：
\begin{enumerate}
\item[MS5] 对 $\mathcal{D}$ 的态射 $f$，有
$f \in S \Leftrightarrow f[1] \in S$\footnote{见注 \ref{derived-remark-MS5}。}。
\item[MS6] 给定实线交换图
$$
\xymatrix{
X \ar[r] \ar[d]^s &
Y \ar[r] \ar[d]^{s'} &
Z \ar[r] \ar@{-->}[d] &
X[1] \ar[d]^{s[1]} \\
X' \ar[r] &
Y' \ar[r] &
Z' \ar[r] &
X'[1]
}
$$
其两行均为特异三角且 $s, s' \in S$，则存在 $S$ 中的态射
$s'' : Z \to Z'$，使得 $(s, s', s'')$ 是三角的态射。
\end{enumerate}
\end{definition}

\noindent
事实证明，这些公理并不独立于定义乘法系的公理。

\begin{lemma}
\label{derived-lemma-localization-conditions}
设 $\mathcal{D}$ 为预三角范畴，$S \subset \text{Arrows}(\mathcal{D})$。
\begin{enumerate}
\item 若 $S$ 包含所有恒等态射并且 MS6 成立
（定义 \ref{derived-definition-localization}），则 $\mathcal{D}$ 的每个同构都属于 $S$。
\item 若 MS1、MS5（范畴，定义
\ref{categories-definition-multiplicative-system}）和 MS6 成立，则 MS2 成立。
\end{enumerate}
\end{lemma}

\begin{proof}
假设 $S$ 包含所有恒等态射且 MS6 成立。设 $f : X \to Y$ 为
$\mathcal{D}$ 的同构。考虑图
$$
\xymatrix{
0 \ar[r] \ar[d]^1 &
X \ar[r]_1 \ar[d]^1 &
X \ar[r] \ar@{-->}[d] &
0[1] \ar[d]^{1[1]} \\
0 \ar[r] &
X \ar[r]^f &
Y \ar[r] &
0[1]
}
$$
由引理 \ref{derived-lemma-third-object-zero}，两行都是特异三角。
由 MS6，虚线箭头存在且属于 $S$，故 $f$ 属于 $S$。

\medskip\noindent
假设 MS1、MS5、MS6 成立。设 $f : X \to Y$ 为 $\mathcal{D}$ 的态射，
$t : X \to X'$ 为 $S$ 的元素。选取特异三角
$(X, Y, Z, f, g, h)$。再选取特异三角
$(X', Y', Z, f', g', t[1] \circ h)$（这里使用 TR1 与 TR2）。
由 MS5、MS6（并用 TR2 旋转），可在交换图
$$
\xymatrix{
X \ar[r] \ar[d]^t &
Y \ar[r] \ar@{..>}[d]^{s'} &
Z \ar[r] \ar[d]^1 &
X[1] \ar[d]^{t[1]} \\
X' \ar[r] &
Y' \ar[r] &
Z \ar[r] &
X'[1]
}
$$
中找到虚线箭头，并且 $s' \in S$。这证明了 LMS2。RMS2 的证明与此对偶。
\end{proof}

\begin{remark}
\label{derived-remark-MS5}
在 MS1 与 MS6 成立时，条件 MS5 等价于要求：对所有
$s \in S$ 和 $n \in \mathbf{Z}$，都有 $s[n] \in S$。
例如，假设 MS5 成立，已知 $s \in S$，而我们想证明 $s[-1] \in S$。
这并非立即成立，因为 $s[-1][1]$ 并不等于 $s$，而只是在
$\mathcal{D}$ 的箭头意义下同构于 $s$。不过，这仍意味着对某些同构 $f$、$g$ 有
$s[-1][1] = f \circ s \circ g$。由引理
\ref{derived-lemma-localization-conditions} (1)，$f, g \in S$；故由 MS1，
$s[-1][1] \in S$；再由 MS5，$s[-1] \in S$。
完整证明留给读者作为练习。
\end{remark}

\begin{lemma}
\label{derived-lemma-triangle-functor-localize}
设 $F : \mathcal{D} \to \mathcal{D}'$ 为预三角范畴之间的正合函子。令
$$
S = \{f \in \text{Arrows}(\mathcal{D})
\mid F(f)\text{ 为同构}\}
$$
则 $S$ 是饱和的（见范畴，定义
\ref{categories-definition-saturated-multiplicative-system}）乘法系，
并且与 $\mathcal{D}$ 上的三角结构相容。
\end{lemma}

\begin{proof}
需要证明公理 MS1 -- MS6；见范畴，定义
\ref{categories-definition-multiplicative-system}、
\ref{categories-definition-saturated-multiplicative-system} 以及定义
\ref{derived-definition-localization}。MS1、MS4 和 MS5 直接来自定义。
MS6 由 TR3 和引理 \ref{derived-lemma-third-isomorphism-triangle} 得出。
由引理 \ref{derived-lemma-localization-conditions}，MS2 成立。为完成证明，
还需证明 MS3。设 $f, g : X \to Y$ 为 $\mathcal{D}$ 的态射，
$t : Z \to X$ 为 $S$ 的元素，并满足 $f \circ t = g \circ t$。
由于 $\mathcal{D}$ 是加性范畴，这只意味着 $a \circ t = 0$，其中
$a = f - g$。用 TR1 选取特异三角 $(Z, X, Q, t, d, h)$。
由于 $a \circ t = 0$，由引理 \ref{derived-lemma-representable-homological}，
存在态射 $i : Q \to Y$，使得 $i \circ d = a$。最后，再次使用 TR1，
可选取三角 $(Q, Y, W, i, j, k)$。图示如下：
$$
\xymatrix{
Z \ar[r]_t & X \ar[r]_d \ar[d]^1 & Q \ar[r] \ar[d]^i & Z[1] \\
& X \ar[r]_a & Y \ar[d]^j \\
& & W
}
$$
现在对该图应用函子 $F$。由于 $t \in S$，由引理
\ref{derived-lemma-third-object-zero} 可知 $F(Q) = 0$。
故由同一引理，$F(j)$ 是同构，即 $j \in S$。最后，
$j \circ a = j \circ i \circ d = 0$，因为 $j \circ i = 0$。
因此 $j \circ f = j \circ g$，从而 LMS3 成立。RMS3 的证明与此对偶。
\end{proof}

\begin{lemma}
\label{derived-lemma-homological-functor-localize}
设 $H : \mathcal{D} \to \mathcal{A}$ 为从预三角范畴到阿贝尔范畴的同调函子。令
$$
S = \{f \in \text{Arrows}(\mathcal{D})
\mid \text{对所有 }i \in \mathbf{Z},\ H^i(f)\text{ 均为同构}\}
$$
则 $S$ 是饱和的（见范畴，定义
\ref{categories-definition-saturated-multiplicative-system}）乘法系，
并且与 $\mathcal{D}$ 上的三角结构相容。
\end{lemma}

\begin{proof}
需要证明公理 MS1 -- MS6；见范畴，定义
\ref{categories-definition-multiplicative-system}、
\ref{categories-definition-saturated-multiplicative-system} 以及定义
\ref{derived-definition-localization}。MS1、MS4 和 MS5 直接来自定义。
MS6 由 TR3 和长正合上同调序列
(\ref{derived-equation-long-exact-cohomology-sequence}) 得出。
由引理 \ref{derived-lemma-localization-conditions}，MS2 成立。为完成证明，
还需证明 MS3。设 $f, g : X \to Y$ 为 $\mathcal{D}$ 的态射，
$t : Z \to X$ 为 $S$ 的元素，并满足 $f \circ t = g \circ t$。
由于 $\mathcal{D}$ 是加性范畴，这只意味着 $a \circ t = 0$，其中
$a = f - g$。用 TR1 和 TR2 选取特异三角 $(Z, X, Q, t, u, v)$。
% P02-E0020: frozen source types the factorization as i \circ g = a; later lines and the diagram require u.
由于 $a \circ t = 0$，由引理 \ref{derived-lemma-representable-homological}，
存在态射 $i : Q \to Y$，使得 $i \circ g = a$。最后，再次使用 TR1，
可选取三角 $(Q, Y, W, i, j, k)$。图示如下：
$$
\xymatrix{
Z \ar[r]_t & X \ar[r]_u \ar[d]^1 & Q \ar[r]_v \ar[d]^i & Z[1] \\
& X \ar[r]_a & Y \ar[d]^j \\
& & W
}
$$
现在对该图应用各函子 $H^i$。由于 $t \in S$，由长正合上同调序列
(\ref{derived-equation-long-exact-cohomology-sequence}) 可知 $H^i(Q) = 0$。
故由同一论证，对所有 $i$，$H^i(j)$ 都是同构，即 $j \in S$。
最后，$j \circ a = j \circ i \circ u = 0$，因为 $j \circ i = 0$。
因此 $j \circ f = j \circ g$，从而 LMS3 成立。RMS3 的证明与此对偶。
\end{proof}

\begin{proposition}
\label{derived-proposition-construct-localization}
设 $\mathcal{D}$ 为预三角范畴，$S$ 为与三角结构相容的乘法系。
则 $S^{-1}\mathcal{D}$ 上存在唯一的预三角范畴结构，使得
$[1] \circ Q = Q \circ [1]$，并且局部化函子
$Q : \mathcal{D} \to S^{-1}\mathcal{D}$ 是正合的。
此外，若 $\mathcal{D}$ 是三角范畴，则 $S^{-1}\mathcal{D}$ 也是三角范畴。
\end{proposition}

\begin{proof}
我们已经在同调，引理 \ref{homology-lemma-localization-additive} 中看到，
$S^{-1}\mathcal{D}$ 是加性范畴，且局部化函子 $Q$ 是加性函子。
由 MS5，存在唯一的加性自等价
$[1] : S^{-1}\mathcal{D} \to S^{-1}\mathcal{D}$，使得
$Q \circ [1] = [1] \circ Q$（函子的等式）；细节从略。
若 $S^{-1}\mathcal{D}$ 的一个三角同构于某个特异三角在局部化函子 $Q$
下的像，则称其为特异三角。

\medskip\noindent
证明 TR1。这里只需证明：若 $a : Q(X) \to Q(Y)$ 是
$S^{-1}\mathcal{D}$ 的态射，则 $a$ 可嵌入一个特异三角。
对某个 $S$ 中的 $s : Y \to Y'$ 和某个 $f : X \to Y'$，写成
$a = Q(s)^{-1} \circ Q(f)$。在 $\mathcal{D}$ 中选取特异三角
$(X, Y', Z, f, g, h)$。于是
$(Q(X), Q(Y), Q(Z), a, Q(g) \circ Q(s), Q(h))$
是 $S^{-1}\mathcal{D}$ 的特异三角。

\medskip\noindent
证明 TR2。由定义立即得到。

\medskip\noindent
证明 TR3。注意，在把两个三角替换为与之同构的三角之后，也可以证明所需虚线箭头的存在性。
因此可设给定 $\mathcal{D}$ 的特异三角
$(X, Y, Z, f, g, h)$ 和 $(X', Y', Z', f', g', h')$，以及交换图
$$
\xymatrix{
Q(X) \ar[r]_{Q(f)} \ar[d]_a & Q(Y) \ar[d]^b \\
Q(X') \ar[r]^{Q(f')} & Q(Y')
}
$$
它位于 $S^{-1}\mathcal{D}$ 中。现在应用范畴，引理
\ref{categories-lemma-left-localization-diagram}，在 $\mathcal{D}$ 中找到态射
$f'' : X'' \to Y''$ 以及交换图
$$
\xymatrix{
X \ar[d]_f \ar[r]_k & X'' \ar[d]^{f''} & X' \ar[d]^{f'} \ar[l]^s \\
Y \ar[r]^l & Y'' & Y' \ar[l]_t
}
$$
其中 $s, t \in S$，且 $a = s^{-1}k$、$b = t^{-1}l$。
此时可在 $\mathcal{D}$ 中使用 TR3 和 MS6，得到交换图
% P02-E0021: frozen source labels this vertical arrow g[1], although k[1] is type-correct.
$$
\xymatrix{
X \ar[r] \ar[d]^k &
Y \ar[r] \ar[d]^l &
Z \ar[r] \ar[d]^m &
X[1] \ar[d]^{g[1]} \\
X'' \ar[r] &
Y'' \ar[r] &
Z'' \ar[r] &
X''[1] \\
X' \ar[r] \ar[u]_s &
Y' \ar[r] \ar[u]_t &
Z' \ar[r] \ar[u]_r &
X'[1] \ar[u]_{s[1]}
}
$$
其中 $r \in S$。于是令 $c = Q(r)^{-1}Q(m)$，便得到所需的三角态射
$$
\xymatrix{
(Q(X), Q(Y), Q(Z), Q(f), Q(g), Q(h))
\ar[d]^{(a, b, c)} \\
(Q(X'), Q(Y'), Q(Z'), Q(f'), Q(g'), Q(h'))
}
$$

\medskip\noindent
这证明了命题的第一个断言。若 $\mathcal{D}$ 还是三角范畴，则为证明
$S^{-1}\mathcal{D}$ 也是三角范畴，还需证明 TR4。由引理
\ref{derived-lemma-easier-axiom-four}，可归结为下列断言：给定可复合态射
$a : Q(X) \to Q(Y)$ 与 $b : Q(Y) \to Q(Z)$，在可能用同构对象替换
$Q(X), Q(Y), Q(Z)$ 后，必须构造一个八面体。为此，可先把 $Y$ 替换为某个对象，
使得对 $\mathcal{D}$ 中某个态射 $f : X \to Y$ 有 $a = Q(f)$。
（更精确地，写成 $a = s^{-1}f$，其中 $s : Y \to Y'$ 属于 $S$，
$f : X \to Y'$，然后用 $Y'$ 替换 $Y$。）此后同样替换 $Z$，使得对
$\mathcal{D}$ 中某个态射 $g : Y \to Z$ 有 $b = Q(g)$。
现在可在 $\mathcal{D}$ 中找到特异三角
$(X, Y, Q_1, f, p_1, d_1)$、
$(X, Z, Q_2, g \circ f, p_2, d_2)$ 和
$(Y, Z, Q_3, g, p_3, d_3)$（由 TR1），以及 TR4 中的态射
$a : Q_1 \to Q_2$ 和 $b : Q_2 \to Q_3$。
立即可验证，对所有这些数据应用函子 $Q$，便在
$S^{-1}\mathcal{D}$ 中得到相应结构。
\end{proof}

\noindent
三角范畴局部化的泛性质如下（我们对预三角范畴表述，因此它自然也适用于三角范畴）。

\begin{lemma}
\label{derived-lemma-universal-property-localization}
设 $\mathcal{D}$ 为预三角范畴，$S$ 为与三角结构相容的乘法系。
设 $Q : \mathcal{D} \to S^{-1}\mathcal{D}$ 为局部化函子，见命题
\ref{derived-proposition-construct-localization}。
\begin{enumerate}
\item 若 $H : \mathcal{D} \to \mathcal{A}$ 是取值于阿贝尔范畴
$\mathcal{A}$ 的同调函子，并且对所有 $s \in S$，$H(s)$ 都是同构，
则满足 $H = H' \circ Q$ 的唯一分解
$H' : S^{-1}\mathcal{D} \to \mathcal{A}$（见范畴，引理
\ref{categories-lemma-properties-left-localization}）也是同调函子。
\item 若 $F : \mathcal{D} \to \mathcal{D}'$ 是取值于预三角范畴
$\mathcal{D}'$ 的正合函子，并且对所有 $s \in S$，$F(s)$ 都是同构，
则满足 $F = F' \circ Q$ 的唯一分解
$F' : S^{-1}\mathcal{D} \to \mathcal{D}'$（见范畴，引理
\ref{categories-lemma-properties-left-localization}）也是正合函子。
\end{enumerate}
\end{lemma}

\begin{proof}
本引理不证自明。细节从略。
\end{proof}

\begin{lemma}
\label{derived-lemma-localization-subcategory}
设 $\mathcal{D}$ 为预三角范畴，$\mathcal{D}' \subset \mathcal{D}$
为满的预三角子范畴。设 $S$ 为 $\mathcal{D}$ 的饱和乘法系，并与三角结构相容。
假设对 $\mathcal{D}$ 中每个 $X$，都存在 $S$ 中的 $s : X' \to X$，
使得 $X'$ 是 $\mathcal{D}'$ 的对象。则
$S' = S \cap \text{Arrows}(\mathcal{D}')$ 是与三角结构相容的饱和乘法系，且函子
$$
(S')^{-1}\mathcal{D}' \longrightarrow S^{-1}\mathcal{D}
$$
是预三角范畴的等价。
\end{lemma}

\begin{proof}
考虑命题 \ref{derived-proposition-construct-localization} 的商函子
$Q : \mathcal{D} \to S^{-1}\mathcal{D}$。由于 $S$ 饱和，态射
$f : X \to Y$ 属于 $S$，当且仅当 $Q(f)$ 可逆；见范畴，引理
\ref{categories-lemma-what-gets-inverted}。因此，$S'$ 正是被复合
$\mathcal{D}' \to \mathcal{D} \to S^{-1}\mathcal{D}$
变为同构的箭头之族。
% P02-E0022: frozen source says “S' is is a saturated”.
故由引理 \ref{derived-lemma-triangle-functor-localize}，$S'$ 是与三角结构相容的饱和乘法系。
由引理 \ref{derived-lemma-universal-property-localization}，得到预三角范畴的正合函子
$(S')^{-1}\mathcal{D}' \to S^{-1}\mathcal{D}$。由假设，此函子本质满。
设 $X', Y'$ 为 $\mathcal{D}'$ 的对象。由范畴，注
\ref{categories-remark-right-localization-morphisms-colimit}，有
$$
\Mor_{S^{-1}\mathcal{D}}(X', Y') =
\colim_{s : X \to X'\text{ 于 }S} \Mor_\mathcal{D}(X, Y')
$$
我们的假设蕴含：对 $S$ 中任意 $s : X \to X'$，都能找到 $S$ 中态射
$s' : X'' \to X$，其中 $X''$ 属于 $\mathcal{D}'$。于是
$s \circ s' : X'' \to X'$ 属于 $S'$，故上面的余极限等于
$$
\colim_{s'' : X'' \to X'\text{ 于 }S'} \Mor_{\mathcal{D}'}(X'', Y') =
\Mor_{(S')^{-1}\mathcal{D}'}(X', Y')
$$
这证明该函子也全忠实，从而证明完成。
\end{proof}

\noindent
下面的引理描述局部化函子的核（见定义 \ref{derived-definition-kernel-category}）。

\begin{lemma}
\label{derived-lemma-kernel-localization}
设 $\mathcal{D}$ 为预三角范畴，$S$ 为与三角结构相容的乘法系，
$Z$ 为 $\mathcal{D}$ 的对象。下列条件等价：
\begin{enumerate}
\item 在 $S^{-1}\mathcal{D}$ 中，$Q(Z) = 0$；
\item 存在 $Z' \in \Ob(\mathcal{D})$，使得 $0 : Z \to Z'$ 是 $S$ 的元素；
\item 存在 $Z' \in \Ob(\mathcal{D})$，使得 $0 : Z' \to Z$ 是 $S$ 的元素；并且
\item 存在对象 $Z'$ 和特异三角 $(X, Y, Z \oplus Z', f, g, h)$，
使得 $f \in S$。
\end{enumerate}
若 $S$ 饱和，则它们还等价于
\begin{enumerate}
\item[(5)] 态射 $0 \to Z$ 是 $S$ 的元素；
\item[(6)] 态射 $Z \to 0$ 是 $S$ 的元素；
\item[(7)] 存在特异三角 $(X, Y, Z, f, g, h)$，使得 $f \in S$。
\end{enumerate}
\end{lemma}

\begin{proof}
(1)、(2)、(3) 的等价性即同调，引理
\ref{homology-lemma-kernel-localization}。若 (2) 成立，则
$(Z'[-1], Z'[-1] \oplus Z, Z, (1, 0), (0, 1), 0)$ 是特异三角
（见引理 \ref{derived-lemma-split}），并带有“$0 \in S$”。旋转后得到 (4)。
若 $(X, Y, Z \oplus Z', f, g, h)$ 是满足 $f \in S$ 的特异三角，
则 $Q(f)$ 是同构，因而 $Q(Z \oplus Z') = 0$，进而 $Q(Z) = 0$。
所以 (1) -- (4) 全部等价。

\medskip\noindent
下面假设 $S$ 饱和。注意，(5)、(6)、(7) 各自都蕴含等价条件 (1) -- (4)
之一。设 $Q(Z) = 0$。于是 $0 \to Z$ 是 $\mathcal{D}$ 中的态射，
并在 $S^{-1}\mathcal{D}$ 中变为同构。根据范畴，引理
\ref{categories-lemma-what-gets-inverted}，$S$ 的饱和性蕴含
$0 \to Z$ 属于 $S$。故 (1) $\Rightarrow$ (5)。对偶地，
(1) $\Rightarrow$ (6)。最后，若 $0 \to Z$ 属于 $S$，则三角
$(0, Z, Z, 0, \text{id}_Z, 0)$ 由 TR1 与 TR2 是特异三角，且正是 (7) 中的三角。
\end{proof}

\begin{lemma}
\label{derived-lemma-limit-triangles}
设 $\mathcal{D}$ 为预三角范畴，$S$ 为 $\mathcal{D}$ 中与三角结构相容的
饱和乘法系。设 $(X, Y, Z, f, g, h)$ 为 $\mathcal{D}$ 中的特异三角。
考虑三角态射所成的范畴
$$
\mathcal{I} =
\{(s, s', s'') : (X, Y, Z, f, g, h) \to (X', Y', Z', f', g', h')
\mid s, s', s'' \in S\}
$$
则 $\mathcal{I}$ 是滤过范畴，并且函子
$\mathcal{I} \to X/S$、$\mathcal{I} \to Y/S$ 和 $\mathcal{I} \to Z/S$
都是余终的。
\end{lemma}

\begin{proof}
我们强烈建议读者跳过本引理的证明，改在餐巾纸上自行推导。

\medskip\noindent
第一个观察是：旋转特异三角（TR2）给出 $\mathcal{I}$ 与特异三角
$(Y, Z, X[1], g, h, -f[1])$ 所对应范畴之间的等价。例如，由此可见，
若我们证明函子 $\mathcal{I} \to X/S$ 余终，那么函子
$\mathcal{I} \to Y/S$ 和 $\mathcal{I} \to Z/S$ 也余终。

\medskip\noindent
注意，若 $s : X \to X'$ 是 $S$ 的态射，则由 MS2 可找到
$s' : Y \to Y'$ 和 $f' : X' \to Y'$，使得
$f' \circ s = s' \circ f$；随后可用 MS6 将其补成 $\mathcal{I}$ 的一个对象。
因此函子 $\mathcal{I} \to X/S$ 在对象上满。由上述旋转，对
$\mathcal{I} \to Y/S$ 和 $\mathcal{I} \to Z/S$ 也同样成立。

\medskip\noindent
设给定 $X/S$ 中的对象 $s_1 : X \to X_1 $、$s_2 : X \to X_2$，
以及 $X/S$ 中的态射 $a : X_1 \to X_2$。由于 $S$ 饱和，由范畴，引理
\ref{categories-lemma-what-gets-inverted} 可知 $a \in S$。
按照上一段的论证，可把 $s_1 : X \to X_1$ 补成 $\mathcal{I}$ 中的对象
$(s_1, s'_1, s''_1) : (X, Y, Z, f, g, h) \to (X_1, Y_1, Z_1, f_1, g_1, h_1)$。
然后重复此步骤，找到
$(a, b, c) : (X_1, Y_1, Z_1, f_1, g_1, h_1) \to (X_2, Y_2, Z_2, f_2, g_2, h_2)$，
其中 $a, b, c \in S$，并补全给定的 $a : X_1 \to X_2$。
于是 $(a, b, c)$ 是 $\mathcal{I}$ 中的态射。因此函子
$\mathcal{I} \to X/S$ 在箭头上也满。再由上述旋转，对函子
$\mathcal{I} \to Y/S$ 和 $\mathcal{I} \to Z/S$ 亦然。

\medskip\noindent
范畴 $\mathcal{I}$ 非空，因为恒等态射给出一个对象。这证明了滤过范畴定义的条件 (1)；
见范畴，定义 \ref{categories-definition-directed}。

\medskip\noindent
我们为范畴 $\mathcal{I}$ 验证范畴，定义
\ref{categories-definition-directed} 的条件 (2)。设给定 $\mathcal{I}$ 中的对象
$(s_1, s'_1, s''_1) : (X, Y, Z, f, g, h) \to (X_1, Y_1, Z_1, f_1, g_1, h_1)$
和
$(s_2, s'_2, s''_2) : (X, Y, Z, f, g, h) \to (X_2, Y_2, Z_2, f_2, g_2, h_2)$。
我们要找 $\mathcal{I}$ 中的一个对象，使它同时是从
$(X_1, Y_1, Z_1, f_1, g_1, h_1)$ 和 $(X_2, Y_2, Z_2, f_2, g_2, h_2)$
出发之箭头的靶。由范畴，注
\ref{categories-remark-left-localization-morphisms-colimit}，范畴
$X/S$、$Y/S$、$Z/S$ 都是滤过的。因此可在 $X/S$ 中找到
$X \to X_3$ 以及态射 $s : X_2 \to X_3$ 和 $a : X_1 \to X_3$。
由上文可找到三角态射
$(s, s', s'') : (X_2, Y_2, Z_2, f_2, g_2, h_2) \to
(X_3, Y_3, Z_3, f_3, g_3, h_3)$，其中 $s', s'' \in S$。
把 $(X_2, Y_2, Z_2)$ 替换为 $(X_3, Y_3, Z_3)$ 后，可假设
$X/S$ 中存在态射 $a : X_1 \to X_2$。对 $Y$ 和 $Z$ 重复此论证
（用上述旋转），可假设存在 $X/S$ 中的态射 $a : X_1 \to X_2$、
$Y/S$ 中的态射 $b : Y_1 \to Y_2$，以及 $Z/S$ 中的态射
$c : Z_1 \to Z_2$。然而，这些态射未必给出特异三角的态射。
另一方面，所需各图在 $S^{-1}\mathcal{D}$ 中确实交换。因此，例如可见，
存在 $S$ 中的态射 $s'_2 : Y_2 \to Y_3$，使得
$s'_2 \circ f_2 \circ a = s'_2 \circ b \circ f_1$。
再像上面那样替换一次 $(X_2, Y_2, Z_2)$，便可进入
$f_2 \circ a = b \circ f_1$ 的情形。旋转并再应用两次同一论证，
可假设 $(a, b, c)$ 是三角的态射。这证明了条件 (2)。

\medskip\noindent
接着验证范畴，定义 \ref{categories-definition-directed} 的条件 (3)。
设 $(s_1, s_1', s_1'') : (X, Y, Z) \to (X_1, Y_1, Z_1)$ 和
$(s_2, s_2', s_2'') : (X, Y, Z) \to (X_2, Y_2, Z_2)$
是 $\mathcal{I}$ 的对象，并设 $(a, b, c), (a', b', c')$ 是它们之间的两个态射。
由于 $a \circ s_1 = a' \circ s_1$，存在态射 $s_3 : X_2 \to X_3$，使得
$s_3 \circ a = s_3 \circ a'$。使用上述满性断言，可将它补成三角态射
$(s_3, s_3', s_3'') : (X_2, Y_2, Z_2) \to (X_3, Y_3, Z_3)$，
其中 $s_3, s_3', s_3'' \in S$。因此
$(s_3 \circ s_2, s_3' \circ s_2', s_3'' \circ s_2'') :
(X, Y, Z) \to (X_3, Y_3, Z_3)$ 也是 $\mathcal{I}$ 的对象；将映射
$(a, b, c), (a', b', c')$ 与 $(s_3, s_3', s_3'')$ 复合后，得到
$a = a'$。通过旋转可同样得到 $b = b'$ 和 $c = c'$。

\medskip\noindent
最后验证 $\mathcal{I} \to X/S$ 余终；见范畴，定义
\ref{categories-definition-cofinal}。第一个条件成立，因为该函子满。
假设在 $X/S$ 中有对象 $s : X \to X'$，并且在 $\mathcal{I}$ 中有两个对象
$(s_1, s'_1, s''_1) : (X, Y, Z, f, g, h) \to (X_1, Y_1, Z_1, f_1, g_1, h_1)$
和
$(s_2, s'_2, s''_2) : (X, Y, Z, f, g, h) \to (X_2, Y_2, Z_2, f_2, g_2, h_2)$，
并有 $X/S$ 中的态射 $t_1 : X' \to X_1$ 和 $t_2 : X' \to X_2$。
由上面已证的 $\mathcal{I}$ 的性质 (2)，可找到 $\mathcal{I}$ 中的态射
$(s_3, s'_3, s''_3) : (X_1, Y_1, Z_1, f_1, g_1, h_1) \to
(X_3, Y_3, Z_3, f_3, g_3, h_3)$
和
$(s_4, s'_4, s''_4) : (X_2, Y_2, Z_2, f_2, g_2, h_2) \to
(X_3, Y_3, Z_3, f_3, g_3, h_3)$。
若复合 $X' \to X_1 \to X_3$ 与 $X' \to X_2 \to X_3$ 相等，
则已经完成（见范畴，定义 \ref{categories-definition-cofinal} 中的显示等式）。
若不相等，则因 $X/S$ 滤过，可选取 $X/S$ 中的态射 $X_3 \to X_4$，
使复合 $X' \to X_1 \to X_3 \to X_4$ 与
$X' \to X_2 \to X_3 \to X_4$ 相等。最后，把 $X_3 \to X_4$
补成 $\mathcal{I}$ 中的态射
$(X_3, Y_3, Z_3) \to (X_4, Y_4, Z_4)$，再与它复合，便得到结论。
\end{proof}







\section{三角范畴的商}
\label{derived-section-quotients}

\noindent
给定一个三角范畴及其一个三角子范畴，可以取“商”来构造另一个三角范畴。
这一构造使用局部化。它类似于阿贝尔范畴对 Serre 子范畴的商；见同调，
第 \ref{homology-section-serre-subcategories} 节。在实际构造之前，我们先简要讨论
正合函子的核。

\begin{definition}
\label{derived-definition-saturated}
设 $\mathcal{D}$ 为预三角范畴。若每当 $X \oplus Y$ 同构于
$\mathcal{D}'$ 的一个对象时，$X$ 和 $Y$ 都分别同构于
$\mathcal{D}'$ 的对象，则称 $\mathcal{D}$ 的满预三角子范畴
$\mathcal{D}'$ 是{\it 饱和的}。
\end{definition}

\noindent
饱和三角子范畴有时称为{\it 厚三角子范畴}。在某些文献中，
这一术语只用于严格满三角子范畴（有时定义本身即蕴含严格满性）。
还有另一个概念，称为 {\it 'epaisse 三角子范畴}。其定义是：给定交换图
$$
\xymatrix{
& S \ar[rd] \\
X \ar[ru] \ar[rr] & & Y \ar[r] & T \ar[r] & X[1]
}
$$
若第二行是特异三角，且 $S$ 与 $T$ 均同构于 $\mathcal{D}'$ 的对象，
则 $X$ 与 $Y$ 也同构于 $\mathcal{D}'$ 的对象。事实证明，这等价于饱和性
（这是初等事实，可见 \cite{Rickard-derived}），而饱和范畴这一概念更便于使用。

\begin{lemma}
\label{derived-lemma-triangle-functor-kernel}
设 $F : \mathcal{D} \to \mathcal{D}'$ 为预三角范畴之间的正合函子。
设 $\mathcal{D}''$ 为 $\mathcal{D}$ 的满子范畴，其对象为
$$
\Ob(\mathcal{D}'') =
\{X \in \Ob(\mathcal{D}) \mid F(X) = 0\}
$$
则 $\mathcal{D}''$ 是 $\mathcal{D}$ 的严格满、饱和预三角子范畴。
若 $\mathcal{D}$ 是三角范畴，则 $\mathcal{D}''$ 是三角子范畴。
\end{lemma}

\begin{proof}
显然，$\mathcal{D}''$ 在 $[1]$ 与 $[-1]$ 下保持。
若 $(X, Y, Z, f, g, h)$ 是 $\mathcal{D}$ 的特异三角，并且
$F(X) = F(Y) = 0$，则由于
$(F(X), F(Y), F(Z), F(f), F(g), F(h))$ 是特异三角，也有
$F(Z) = 0$。因此可应用引理 \ref{derived-lemma-triangulated-subcategory}，
得到 $\mathcal{D}''$ 是预三角子范畴（相应地，若 $\mathcal{D}$ 是三角范畴，
则它是三角子范畴）。最后，饱和性由
$F(X) \oplus F(Y) = 0 \Rightarrow F(X) = F(Y) = 0$ 得出。
\end{proof}

\begin{lemma}
\label{derived-lemma-homological-functor-kernel}
设 $H : \mathcal{D} \to \mathcal{A}$ 为从预三角范畴到阿贝尔范畴的同调函子。
设 $\mathcal{D}'$ 为 $\mathcal{D}$ 的满子范畴，其对象为
$$
\Ob(\mathcal{D}') =
\{X \in \Ob(\mathcal{D}) \mid
H(X[n]) = 0\text{ 对所有 }n \in \mathbf{Z}\}
$$
则 $\mathcal{D}'$ 是 $\mathcal{D}$ 的严格满、饱和预三角子范畴。
若 $\mathcal{D}$ 是三角范畴，则 $\mathcal{D}'$ 是三角子范畴。
\end{lemma}

\begin{proof}
显然，$\mathcal{D}'$ 在 $[1]$ 与 $[-1]$ 下保持。若
$(X, Y, Z, f, g, h)$ 是 $\mathcal{D}$ 的特异三角，并且对所有 $n$ 都有
$H(X[n]) = H(Y[n]) = 0$，则由长正合序列
(\ref{derived-equation-long-exact-cohomology-sequence})，对所有 $n$ 也有
$H(Z[n]) = 0$。因此可应用引理 \ref{derived-lemma-triangulated-subcategory}，
得到 $\mathcal{D}'$ 是预三角子范畴（相应地，若 $\mathcal{D}$ 是三角范畴，
则它是三角子范畴）。饱和性由
\begin{align*}
H((X \oplus Y)[n]) = 0 & \Rightarrow H(X[n] \oplus Y[n]) = 0 \\
& \Rightarrow H(X[n]) \oplus H(Y[n]) = 0 \\
& \Rightarrow H(X[n]) = H(Y[n]) = 0
\end{align*}
对所有 $n \in \mathbf{Z}$ 成立而得。
\end{proof}

\begin{lemma}
\label{derived-lemma-homological-functor-bounded}
设 $H : \mathcal{D} \to \mathcal{A}$ 为从预三角范畴到阿贝尔范畴的同调函子。
设 $\mathcal{D}_H^{+}, \mathcal{D}_H^{-}, \mathcal{D}_H^b$
为 $\mathcal{D}$ 的满子范畴，其对象为
$$
\begin{matrix}
\Ob(\mathcal{D}_H^{+}) =
\{X \in \Ob(\mathcal{D}) \mid
H(X[n]) = 0\text{ 对所有 }n \ll 0\} \\
\Ob(\mathcal{D}_H^{-}) =
\{X \in \Ob(\mathcal{D}) \mid
H(X[n]) = 0\text{ 对所有 }n \gg 0\} \\
\Ob(\mathcal{D}_H^b) =
\{X \in \Ob(\mathcal{D}) \mid
H(X[n]) = 0\text{ 对所有 }|n| \gg 0\}
\end{matrix}
$$
它们每一个都是 $\mathcal{D}$ 的严格满、饱和预三角子范畴。
若 $\mathcal{D}$ 是三角范畴，则它们每一个都是三角子范畴。
\end{lemma}

\begin{proof}
我们对 $\mathcal{D}_H^{+}$ 证明此事。显然，它在 $[1]$ 与 $[-1]$ 下保持。
若 $(X, Y, Z, f, g, h)$ 是 $\mathcal{D}$ 的特异三角，并且对所有
$n \ll 0$ 有 $H(X[n]) = H(Y[n]) = 0$，则由长正合序列
(\ref{derived-equation-long-exact-cohomology-sequence})，对所有 $n \ll 0$
也有 $H(Z[n]) = 0$。因此可应用引理
\ref{derived-lemma-triangulated-subcategory}，得到 $\mathcal{D}_H^{+}$ 是预三角子范畴
（相应地，若 $\mathcal{D}$ 是三角范畴，则它是三角子范畴）。饱和性由
\begin{align*}
H((X \oplus Y)[n]) = 0 & \Rightarrow H(X[n] \oplus Y[n]) = 0 \\
& \Rightarrow H(X[n]) \oplus H(Y[n]) = 0 \\
& \Rightarrow H(X[n]) = H(Y[n]) = 0
\end{align*}
对所有 $n \in \mathbf{Z}$ 成立而得。
\end{proof}

\begin{definition}
\label{derived-definition-kernel-category}
设 $\mathcal{D}$ 为（预）三角范畴。
\begin{enumerate}
\item 设 $F : \mathcal{D} \to \mathcal{D}'$ 为正合函子。
{\it $F$ 的核}是引理 \ref{derived-lemma-triangle-functor-kernel} 所描述的严格满、
饱和（预）三角子范畴。
\item 设 $H : \mathcal{D} \to \mathcal{A}$ 为同调函子。
{\it $H$ 的核}是引理 \ref{derived-lemma-homological-functor-kernel} 所描述的严格满、
饱和（预）三角子范畴。
\end{enumerate}
它们有时记为 $\Ker(F)$ 或 $\Ker(H)$。
\end{definition}

\noindent
下面引理的证明使用 TR4。

\begin{lemma}
\label{derived-lemma-construct-multiplicative-system}
设 $\mathcal{D}$ 为三角范畴，$\mathcal{D}' \subset \mathcal{D}$
为满三角子范畴。令
\begin{equation}
\label{derived-equation-multiplicative-system}
S =
\left\{
\begin{matrix}
f \in \text{Arrows}(\mathcal{D})
\text{ 使得存在一个特异三角 }\\
(X, Y, Z, f, g, h) \text{ 属于 }\mathcal{D}\text{ 其中 }
Z\text{ 同构于下列范畴中的对象： }\mathcal{D}'
\end{matrix}
\right\}
\end{equation}
则 $S$ 是与 $\mathcal{D}$ 上三角结构相容的乘法系。在此情形下，下列条件等价：
\begin{enumerate}
\item $S$ 是饱和乘法系；
\item $\mathcal{D}'$ 是饱和三角子范畴。
\end{enumerate}
\end{lemma}

\begin{proof}
为证明第一个断言，需要证明 MS1、MS2、MS3、MS5、MS6。

\medskip\noindent
证明 MS1。显然，恒等态射属于 $S$，因为对 $\mathcal{D}$ 的每个对象 $X$，
$(X, X, 0, 1, 0, 0)$ 都是特异三角，而且 $0$ 是 $\mathcal{D}'$ 的对象。
设 $f : X \to Y$ 和 $g : Y \to Z$ 为 $S$ 中的可复合态射。
选取特异三角 $(X, Y, Q_1, f, p_1, d_1)$、
$(X, Z, Q_2, g \circ f, p_2, d_2)$ 和
$(Y, Z, Q_3, g, p_3, d_3)$。由假设，$Q_1$ 和 $Q_3$ 同构于
$\mathcal{D}'$ 的对象。由 TR4，存在特异三角
$(Q_1, Q_2, Q_3, a, b, c)$。由于 $\mathcal{D}'$ 是三角子范畴，
$Q_2$ 同构于 $\mathcal{D}'$ 的对象。因此 $g \circ f \in S$。

\medskip\noindent
证明 MS3。设 $a : X \to Y$ 为态射，$t : Z \to X$ 为 $S$ 中满足
$a \circ t = 0$ 的元素。为证明 LMS3，只需找到 $s \in S$ 使得
$s \circ a = 0$；可与引理 \ref{derived-lemma-triangle-functor-localize} 的证明比较。
用 TR1 与 TR2 选取特异三角 $(Z, X, Q, t, g, h)$。
由于 $a \circ t = 0$，由引理 \ref{derived-lemma-representable-homological}，
存在态射 $i : Q \to Y$ 使得 $i \circ g = a$。最后再次使用 TR1，
选取三角 $(Q, Y, W, i, s, k)$。图示如下：
$$
\xymatrix{
Z \ar[r]_t & X \ar[r]_g \ar[d]^1 & Q \ar[r] \ar[d]^i & Z[1] \\
& X \ar[r]_a & Y \ar[d]^s \\
& & W
}
$$
由于 $t \in S$，$Q$ 同构于 $\mathcal{D}'$ 的对象。因此 $s \in S$。
最后，由引理 \ref{derived-lemma-composition-zero}，$s \circ i = 0$，从而
$s \circ a = s \circ i \circ g = 0$。故 LMS3 成立。RMS3 的证明与此对偶。

\medskip\noindent
证明 MS5。这由特异三角及 $\mathcal{D}'$ 在平移下稳定得出。

\medskip\noindent
证明 MS6。设给定交换图
$$
\xymatrix{
X \ar[r] \ar[d]^s &
Y \ar[d]^{s'} \\
X' \ar[r] &
Y'
}
$$
其中 $s, s' \in S$。由命题 \ref{derived-proposition-9}，可将其扩张为九方图。
由于 $s, s'$ 是 $S$ 的元素，$X'', Y''$ 同构于 $\mathcal{D}'$ 的对象。
又因 $\mathcal{D}'$ 是满三角子范畴，$Z''$ 也同构于 $\mathcal{D}'$ 的对象。
因而态射 $Z \to Z'$ 是 $S$ 的元素。这证明了 MS6。

\medskip\noindent
MS2 是 MS1、MS5、MS6 的形式推论；见引理
\ref{derived-lemma-localization-conditions}。这完成了引理第一个断言的证明。

\medskip\noindent
假设 $S$ 饱和。（下文将不再说明地使用特异三角的旋转。）
设 $X \oplus Y$ 是同构于 $\mathcal{D}'$ 中某对象的对象。考虑态射
$f : 0 \to X$。复合 $0 \to X \to X \oplus Y$ 属于 $S$，因为
$(0, X \oplus Y, X \oplus Y, 0, 1, 0)$ 是特异三角。
复合 $Y[-1] \to 0 \to X$ 属于 $S$，因为
$(X, X \oplus Y, Y, (1, 0), (0, 1), 0)$ 是特异三角；见引理
\ref{derived-lemma-split}。由于 $S$ 饱和，$0 \to X$ 属于 $S$。
因此，如所需，$X$ 同构于 $\mathcal{D}'$ 的对象。

\medskip\noindent
最后，假设 $\mathcal{D}'$ 是饱和三角子范畴。设
$$
W \xrightarrow{h}
X \xrightarrow{g}
Y \xrightarrow{f} Z
$$
为 $\mathcal{D}$ 的可复合态射，且 $fg, gh \in S$。
我们将逐步构造下图中的对象：
$$
\xymatrix{
 & &
Q_{12} \ar[rd] & &
Q_{23} \ar[rd] \\
 &
Q_1 \ar[ld]_{\! + \! 1} \ar[ru] & &
Q_2 \ar[ld]_{\! + \! 1} \ar[ll]_{\! + \! 1} \ar[ru] & &
Q_3 \ar[ld]_{\! + \! 1} \ar[ll]_{\! + \! 1} \\
W \ar[rr] & &
X \ar[lu] \ar[rr] & &
Y \ar[lu] \ar[rr] & &
Z \ar[lu]
}
$$
先选取特异三角 $(W, X, Q_1)$、$(X, Y, Q_2)$、$(Y, Z, Q_3)$、
$(W, Y, Q_{12})$ 和 $(X, Z, Q_{23})$。记
$s : Q_2 \to Q_1[1]$ 为复合 $Q_2 \to X[1] \to Q_1[1]$；记
$t : Q_3 \to Q_2[1]$ 为复合 $Q_3 \to Y[1] \to Q_2[1]$。
把 TR4 应用于复合 $W \to X \to Y$ 与复合 $X \to Y \to Z$，
得到使用态射 $s$ 和 $t$ 的特异三角
$(Q_1, Q_{12}, Q_2)$ 和 $(Q_2, Q_{23}, Q_3)$。
由于已假设 $W \to Y$ 与 $X \to Z$ 属于 $S$，对象 $Q_{12}$ 和
$Q_{23}$ 同构于 $\mathcal{D}'$ 的对象。又因 $S$ 在复合和平移下封闭，
$s[1]t$ 也属于 $S$。注意 $s[1]t = 0$，因为
$Y[1] \to Q_2[1] \to X[2]$ 为零；见引理
\ref{derived-lemma-composition-zero}。因此由引理 \ref{derived-lemma-split}，
$Q_3[1] \oplus Q_1[2]$ 同构于 $\mathcal{D}'$ 的对象。
由对 $\mathcal{D}'$ 的假设，$Q_3$ 和 $Q_1$ 同构于
$\mathcal{D}'$ 的对象。考察特异三角 $(Q_1, Q_{12}, Q_2)$，可知
$Q_2$ 也同构于 $\mathcal{D}'$ 的对象。最后考察特异三角
$(X, Y, Q_2)$，得到 $g \in S$。
% P02-E0023: frozen source has “It is also follows”.
（还可推出 $h, f \in S$，但我们不需要这一点。）
\end{proof}

\begin{definition}
\label{derived-definition-quotient-category}
设 $\mathcal{D}$ 为三角范畴，$\mathcal{B}$ 为满三角子范畴。
定义{\it 商范畴 $\mathcal{D}/\mathcal{B}$} 为
$\mathcal{D}/\mathcal{B} = S^{-1}\mathcal{D}$，其中 $S$ 是通过引理
\ref{derived-lemma-construct-multiplicative-system} 与 $\mathcal{B}$ 相伴的
$\mathcal{D}$ 的乘法系。此时局部化函子
$Q : \mathcal{D} \to \mathcal{D}/\mathcal{B}$ 称为{\it 商函子}。
\end{definition}

\noindent
注意，商函子 $Q : \mathcal{D} \to \mathcal{D}/\mathcal{B}$
是三角范畴之间的正合函子；见命题 \ref{derived-proposition-construct-localization}。
这一构造的泛性质如下。

\begin{lemma}
\label{derived-lemma-universal-property-quotient}
\begin{slogan}
Verdier 商的泛性质。
\end{slogan}
设 $\mathcal{D}$ 为三角范畴，$\mathcal{B}$ 为 $\mathcal{D}$ 的满三角子范畴，
$Q : \mathcal{D} \to \mathcal{D}/\mathcal{B}$ 为商函子。
\begin{enumerate}
\item 若 $H : \mathcal{D} \to \mathcal{A}$ 是取值于阿贝尔范畴
$\mathcal{A}$ 的同调函子，并且 $\mathcal{B} \subset \Ker(H)$，
则存在唯一分解 $H' : \mathcal{D}/\mathcal{B} \to \mathcal{A}$，
使得 $H = H' \circ Q$，并且 $H'$ 也是同调函子。
\item 若 $F : \mathcal{D} \to \mathcal{D}'$ 是取值于预三角范畴
$\mathcal{D}'$ 的正合函子，并且 $\mathcal{B} \subset \Ker(F)$，
则存在唯一分解 $F' : \mathcal{D}/\mathcal{B} \to \mathcal{D}'$，
使得 $F = F' \circ Q$，并且 $F'$ 也是正合函子。
\end{enumerate}
\end{lemma}

\begin{proof}
本引理由引理 \ref{derived-lemma-universal-property-localization} 得出。
事实上，若 $f : X \to Y$ 是 $\mathcal{D}$ 的态射，并且对某个特异三角
$(X, Y, Z, f, g, h)$，对象 $Z$ 同构于 $\mathcal{B}$ 的对象，
则在 (1)（相应地，(2)）的假设下，$H(f)$（相应地，$F(f)$）是同构。
细节从略。
\end{proof}

\noindent
商函子的核可描述如下。

\begin{lemma}
\label{derived-lemma-kernel-quotient}
设 $\mathcal{D}$ 为三角范畴，$\mathcal{B}$ 为满三角子范畴。
商函子 $Q : \mathcal{D} \to \mathcal{D}/\mathcal{B}$ 的核是
$\mathcal{D}$ 的严格满子范畴，其对象为
$$
\Ob(\Ker(Q)) =
\left\{
\begin{matrix}
Z \in \Ob(\mathcal{D})
\text{ 使得存在一个 }Z' \in \Ob(\mathcal{D}) \\
\text{ 使得 }Z \oplus Z'\text{ 同构于下列范畴中的对象： }\mathcal{B}
\end{matrix}
\right\}
$$
换言之，它是包含 $\mathcal{B}$ 的 $\mathcal{D}$ 的最小严格满、饱和三角子范畴。
\end{lemma}

\begin{proof}
先注意，核自动是严格满三角子范畴，并且包含其任一对象的所有直和因子；
见引理 \ref{derived-lemma-triangle-functor-kernel}。其对象的上述描述由定义和引理
\ref{derived-lemma-kernel-localization} 第 (4) 部分得出。
\end{proof}

\noindent
设 $\mathcal{D}$ 为三角范畴。现在已有一些构造，它们在下列两个偏序集之间诱导保序映射：
\begin{enumerate}
\item $\mathcal{D}$ 中与三角结构相容的乘法系 $S$ 所成的偏序集；
\item 满三角子范畴 $\mathcal{B} \subset \mathcal{D}$ 所成的偏序集。
\end{enumerate}
具体地，这些构造由
$S \mapsto \mathcal{B}(S) = \Ker(Q : \mathcal{D} \to S^{-1}\mathcal{D})$
和 $\mathcal{B} \mapsto S(\mathcal{B})$ 给出，其中 $S(\mathcal{B})$
是 (\ref{derived-equation-multiplicative-system}) 中的乘法集，即
$$
S(\mathcal{B}) =
\left\{
\begin{matrix}
f \in \text{Arrows}(\mathcal{D})
\text{ 使得存在一个特异三角 }\\
(X, Y, Z, f, g, h) \text{ 属于 }\mathcal{D}\text{ 其中 }
Z\text{ 同构于下列范畴中的对象： }\mathcal{B}
\end{matrix}
\right\}
$$
注意，这两个操作并不互为逆操作。

\begin{lemma}
\label{derived-lemma-operations}
设 $\mathcal{D}$ 为三角范畴。上述操作具有下列性质：
\begin{enumerate}
\item $S(\mathcal{B}(S))$ 是 $S$ 的“饱和”，即包含 $S$ 的
$\mathcal{D}$ 中最小饱和乘法系；
\item $\mathcal{B}(S(\mathcal{B}))$ 是 $\mathcal{B}$ 的“饱和”，
即包含 $\mathcal{B}$ 的 $\mathcal{D}$ 中最小严格满、饱和三角子范畴。
\end{enumerate}
特别地，这些构造在以下两个（偏序）集合之间定义互逆映射：
$\mathcal{D}$ 中与 $\mathcal{D}$ 上三角结构相容的饱和乘法系所成的集合，
以及 $\mathcal{D}$ 的严格满、饱和三角子范畴所成的集合。
\end{lemma}

\begin{proof}
先从满三角子范畴 $\mathcal{B}$ 出发。此时
$\mathcal{B}(S(\mathcal{B})) =
\Ker(Q : \mathcal{D} \to \mathcal{D}/\mathcal{B})$，
所以 (2) 正是引理 \ref{derived-lemma-kernel-quotient} 的内容。

\medskip\noindent
再设 $S$ 是 $\mathcal{D}$ 中与 $\mathcal{D}$ 上三角结构相容的乘法系。
于是 $\mathcal{B}(S) = \Ker(Q : \mathcal{D} \to S^{-1}\mathcal{D})$。
因此（在局部化范畴中使用引理 \ref{derived-lemma-third-object-zero}）
\begin{align*}
S(\mathcal{B}(S))
& =
\left\{
\begin{matrix}
f \in \text{Arrows}(\mathcal{D})
\text{ 使得存在一个特异}\\
\text{特异三角 }(X, Y, Z, f, g, h) \text{ 属于 }\mathcal{D}\text{ 其中 }Q(Z) = 0
\end{matrix}
\right\}.
\\
& =
\{f \in \text{Arrows}(\mathcal{D}) \mid Q(f)\text{ 为同构}\} \\
& = \hat S = S'
\end{align*}
其中记号取自范畴，引理 \ref{categories-lemma-what-gets-inverted}。
该引理的最后一个断言完成了证明。
\end{proof}

\begin{lemma}
\label{derived-lemma-acyclic-general}
设 $H : \mathcal{D} \to \mathcal{A}$ 为从三角范畴 $\mathcal{D}$ 到
阿贝尔范畴 $\mathcal{A}$ 的同调函子；见定义 \ref{derived-definition-homological}。
$\mathcal{D}$ 的子范畴 $\Ker(H)$ 是 $\mathcal{D}$ 的严格满、饱和三角子范畴，
而与其相应的饱和乘法系（见引理 \ref{derived-lemma-operations}）是集合
$$
S = \{f \in \text{Arrows}(\mathcal{D}) \mid
\text{对所有 }i \in \mathbf{Z},\ H^i(f)\text{ 均为同构}\}.
$$
函子 $H$ 经过商函子 $Q : \mathcal{D} \to \mathcal{D}/\Ker(H)$ 分解。
\end{lemma}

\begin{proof}
由引理 \ref{derived-lemma-homological-functor-kernel}，范畴 $\Ker(H)$
是 $\mathcal{D}$ 的严格满、饱和三角子范畴。由引理
\ref{derived-lemma-homological-functor-localize}，集合 $S$ 是与三角结构相容的
饱和乘法系。回忆与 $\Ker(H)$ 相应的乘法系是集合
$$
\left\{
\begin{matrix}
f \in \text{Arrows}(\mathcal{D})
\text{ 使得存在一个特异三角 }\\
(X, Y, Z, f, g, h)\text{ 其中 } H^i(Z) = 0 \text{ 对所有 }i
\end{matrix}
\right\}.
$$
由长正合上同调序列（见 (\ref{derived-equation-long-exact-cohomology-sequence})），
显然，$f$ 属于该集合当且仅当 $f$ 属于 $S$。最后，$H$ 经过 $Q$ 的分解
是引理 \ref{derived-lemma-universal-property-quotient} 的推论。
\end{proof}









\section{正合函子的伴随}
\label{derived-section-adjoints}

\noindent
本节给出三角范畴之间伴随函子的有关结果。

\begin{lemma}
\label{derived-lemma-adjoint-is-exact}
设 $F : \mathcal{D} \to \mathcal{D}'$ 为三角范畴之间的正合函子。
若 $F$ 有右伴随 $G: \mathcal{D'} \to \mathcal{D}$，则 $G$ 也是正合函子。
\end{lemma}

\begin{proof}
设 $\xi_X : F(X[1]) \to F(X)[1]$ 如定义
\ref{derived-definition-exact-functor-triangulated-categories} 所述。
设 $\epsilon_A : F(G(A)) \to A$ 为伴随映射。考虑复合
$$
F(G(A)[1]) \xrightarrow{\xi_{G(A)}} F(G(A))[1] \xrightarrow{\epsilon_A[1]} A[1]
$$
此映射的伴随是一个映射 $G(A)[1] \to G(A[1])$，我们断言它是同构。
为此，由 Yoneda 引理，只需证明：对 $\mathcal{D}$ 的每个对象 $X$ 应用
$\Mor_\mathcal{D}(X[1], -)$ 后，得到双射。
现在，每个态射 $f : X[1] \to G(A)[1]$ 都唯一地形如 $g[1]$，其中
$g : X \to G(A)$；每个 $g$ 又唯一地伴随于某个 $h : F(X) \to A$；
进而 $h[1] \circ \xi_X$ 唯一地伴随于某个 $i : X[1] \to G(A[1])$。
把 $f$ 送到 $i$ 的规则给出
$\Mor_\mathcal{D}(X[1], G(A)[1])$ 与
$\Mor_\mathcal{D}(X[1], G(A[1]))$ 之间的双射。
该双射确由上述态射给出，这可由范畴，第
\ref{categories-section-adjoint} 节的讨论以及下列计算得出：
\begin{align*}
\epsilon_A[1] \circ \xi_{G(A)} \circ F(f)
& =
\epsilon_A[1] \circ \xi_{G(A)} \circ F(g[1]) \\
& =
\epsilon_A[1] \circ F(g)[1] \circ \xi_X \\
& = 
(\epsilon_A \circ F(g))[1] \circ \xi_X \\
& =
h[1] \circ \xi_X
\end{align*}
略去一些细节。我们将用同构 $G(A)[1] \to G(A[1])$ 的逆证明
$G$ 是正合函子。这些同构关于 $A$ 具有函子性（细节从略）。

\medskip\noindent
设 $A \to B \to C \to A[1]$ 为 $\mathcal{D}'$ 中的特异三角。
在 $\mathcal{D}$ 中选取特异三角
$$
G(A) \to G(B) \to X \to G(A)[1]
$$
于是 $F(G(A)) \to F(G(B)) \to F(X) \to F(G(A))[1]$
是 $\mathcal{D}'$ 中的特异三角。由 TR3，可选取特异三角的态射
$$
\xymatrix{
F(G(A)) \ar[r] \ar[d]^{\epsilon_A} &
F(G(B)) \ar[r] \ar[d]^{\epsilon_B} &
F(X) \ar[r] \ar[d] &
F(G(A))[1] \ar[d]^{\epsilon_A[1]} \\
A \ar[r] & B \ar[r] & C \ar[r] & A[1]
}
$$
回忆，箭头 $F(X) \to F(G(A))[1]$ 是
$F(X) \to F(G(A)[1])$ 与 $\xi_{G(A)}$ 的复合。
因此利用 $G$ 是 $F$ 的右伴随及上述讨论，得到态射 $X \to G(C)$，
使下图交换：
$$
\xymatrix{
G(A) \ar[r] \ar[d] & G(B) \ar[r] \ar[d] & X \ar[r] \ar[d] & G(A)[1] \ar[d] \\
G(A) \ar[r] & G(B) \ar[r] & G(C) \ar[r] & G(A[1])
}
$$
且右边竖直箭头正是上面构造的态射。
% P02-E0024: frozen source applies Hom_{D'} to X and G(C), although both lie in D.
对 $\mathcal{D}'$ 的对象 $W$ 应用同调函子
$\Hom_{\mathcal{D}'}(W, -)$，由 $5$ 引理推得
$$
\Hom_{\mathcal{D}'}(W, X) \to \Hom_{\mathcal{D}'}(W, G(C))
$$
是双射；再次使用 Yoneda 引理，得到 $X \to G(C)$ 是同构。
因此 $G(A) \to G(B) \to G(C) \to G(A)[1]$ 是特异三角，正如所需。
\end{proof}

\begin{lemma}
\label{derived-lemma-fully-faithful-adjoint-kernel-zero}
设 $\mathcal{D}$、$\mathcal{D}'$ 为三角范畴，
$F : \mathcal{D} \to \mathcal{D}'$ 和
$G : \mathcal{D}' \to \mathcal{D}$ 为函子。假设
\begin{enumerate}
\item $F$ 与 $G$ 都是正合函子；
\item $F$ 全忠实；
\item $G$ 是 $F$ 的右伴随；
\item $G$ 的核为零。
\end{enumerate}
则 $F$ 是范畴等价。
\end{lemma}

\begin{proof}
由于 $F$ 全忠实，伴随映射 $\text{id} \to G \circ F$ 是同构
（范畴，引理 \ref{categories-lemma-adjoint-fully-faithful}）。
设 $X$ 为 $\mathcal{D}'$ 的对象。选取 $\mathcal{D}'$ 中的特异三角
$$
F(G(X)) \to X \to Y \to F(G(X))[1]
$$
应用 $G$，并使用 $G(F(G(X))) = G(X)$，得到特异三角
$$
G(X) \to G(X) \to G(Y) \to G(X)[1]
$$
因此 $G(Y) = 0$，故 $Y = 0$。所以 $F(G(X)) \to X$ 是同构。
\end{proof}








\section{同伦范畴}
\label{derived-section-homotopy}

\noindent
设 $\mathcal{A}$ 为加性范畴。$\mathcal{A}$ 的同伦范畴
$K(\mathcal{A})$ 是由 $\mathcal{A}$ 的复形组成的范畴，其态射是复形态射模同伦。
下面给出形式定义。

\begin{definition}
\label{derived-definition-complexes-notation}
设 $\mathcal{A}$ 为加性范畴。
\begin{enumerate}
\item 令 $\text{Comp}(\mathcal{A}) = \text{CoCh}(\mathcal{A})$
为{\it （上链）复形范畴}。
\item 若对所有 $n \ll 0$ 都有 $K^n = 0$，则称复形 $K^\bullet$
{\it 下有界}。
\item 若对所有 $n \gg 0$ 都有 $K^n = 0$，则称复形 $K^\bullet$
{\it 上有界}。
\item 若对所有 $|n| \gg 0$ 都有 $K^n = 0$，则称复形 $K^\bullet$
{\it 有界}。
\item 令 $\text{Comp}^{+}(\mathcal{A})$、$\text{Comp}^{-}(\mathcal{A})$、
相应地 $\text{Comp}^b(\mathcal{A})$ 为 $\text{Comp}(\mathcal{A})$
的满子范畴，其对象分别为下有界、上有界、相应地有界的复形。
\item 令 $K(\mathcal{A})$ 为与 $\text{Comp}(\mathcal{A})$ 有相同对象的范畴，
而其态射为复形映射的同伦类（见同调，引理
\ref{homology-lemma-compose-homotopy-cochain}）。
\item 令 $K^{+}(\mathcal{A})$、$K^{-}(\mathcal{A})$、
相应地 $K^b(\mathcal{A})$ 为 $K(\mathcal{A})$ 的满子范畴，
其对象分别为 $\mathcal{A}$ 的下有界、上有界、相应地有界复形。
\end{enumerate}
\end{definition}

\noindent
我们将看到，范畴 $K(\mathcal{A})$、$K^{+}(\mathcal{A})$、
$K^{-}(\mathcal{A})$ 和 $K^b(\mathcal{A})$
都是三角范畴。为了证明这一点，我们先发展一些
与锥和分裂正合序列有关的工具。




\section{锥与逐项分裂序列}
\label{derived-section-cones}

\noindent
设 $\mathcal{A}$ 为加性范畴，并设 $K(\mathcal{A})$ 表示由
$\mathcal{A}$ 的复形组成的范畴，其态射是模同伦的复形态射。
注意，复形上的平移函子 $[n]$（见同调，定义
\ref{homology-definition-shift-cochain}）诱导出函子
$[n] : K(\mathcal{A}) \to K(\mathcal{A})$，满足
$[n] \circ [m] = [n + m]$ 且 $[0] = \text{id}$。

\begin{definition}
\label{derived-definition-cone}
设 $\mathcal{A}$ 为加性范畴。设
$f : K^\bullet \to L^\bullet$ 为 $\mathcal{A}$ 的复形的态射。
$f$ 的{\it 锥}是复形 $C(f)^\bullet$，其定义为
$C(f)^n = L^n \oplus K^{n + 1}$，微分为
$$
d_{C(f)}^n =
\left(
\begin{matrix}
d^n_L & f^{n + 1} \\
0 & -d_K^{n + 1}
\end{matrix}
\right)
$$
它带有典范复形态射
$i : L^\bullet \to C(f)^\bullet$ 和 $p : C(f)^\bullet \to K^\bullet[1]$，
它们由显然的映射 $L^n \to C(f)^n \to K^{n + 1}$ 诱导。
\end{definition}

\noindent
换言之，$(K, L, C(f), f, i, p)$ 构成一个三角：
$$
K^\bullet \to L^\bullet \to C(f)^\bullet \to K^\bullet[1]
$$
这个三角的构造在下述意义下具有函子性。

\begin{lemma}
\label{derived-lemma-functorial-cone}
假设
$$
\xymatrix{
K_1^\bullet \ar[r]_{f_1} \ar[d]_a & L_1^\bullet \ar[d]^b \\
K_2^\bullet \ar[r]^{f_2} & L_2^\bullet
}
$$
是一个复形态射图，并且它交换至同伦。那么存在态射
$c : C(f_1)^\bullet \to C(f_2)^\bullet$，它在 $K(\mathcal{A})$ 中给出三角的态射
$(a, b, c) : (K_1^\bullet, L_1^\bullet, C(f_1)^\bullet, f_1, i_1, p_1)
\to
(K_2^\bullet, L_2^\bullet, C(f_2)^\bullet, f_2, i_2, p_2)$。
\end{lemma}

\begin{proof}
设 $h^n : K_1^n \to L_2^{n - 1}$ 为一族态射，满足
$b \circ f_1 - f_2 \circ a= d \circ h + h \circ d$。
用下列矩阵定义 $c^n$：
$$
c^n =
\left(
\begin{matrix}
b^n & h^{n + 1} \\
0 & a^{n + 1}
\end{matrix}
\right) :
L_1^n \oplus K_1^{n + 1} \to L_2^n \oplus K_2^{n + 1}
$$
% P02-E0025: frozen source has the grammatical error "A matrix computation show".
矩阵计算表明 $c$ 是复形态射。显然有
$c \circ i_1 = i_2 \circ b$，同样不难验证
$p_2 \circ c = a \circ p_1$。
\end{proof}

\noindent
注意，引理 \ref{derived-lemma-functorial-cone} 的证明中构造的态射
$c : C(f_1)^\bullet \to C(f_2)^\bullet$ 一般依赖于
$f_2 \circ a$ 与 $b \circ f_1$ 之间所选的同伦 $h$。

\begin{lemma}
\label{derived-lemma-map-from-cone}
假设 $f: K^\bullet \to L^\bullet$ 和 $g : L^\bullet \to M^\bullet$
是复形态射，且 $g \circ f$ 零同伦。那么
\begin{enumerate}
\item $g$ 经由某个态射 $C(f)^\bullet \to M^\bullet$ 分解；
\item $f$ 经由某个态射 $K^\bullet \to C(g)^\bullet[-1]$ 分解。
\end{enumerate}
\end{lemma}

\begin{proof}
这些假设表明图
$$
\xymatrix{
K^\bullet \ar[r]_f \ar[d] & L^\bullet \ar[d]^g \\
0 \ar[r] & M^\bullet
}
$$
交换至同伦。由于 $0 \to M^\bullet$ 的锥是 $M^\bullet$，
引理 \ref{derived-lemma-functorial-cone} 中的映射
$C(f)^\bullet \to C(0 \to M^\bullet) = M^\bullet$
就是 (1) 中的映射。$K^\bullet \to 0$ 的锥是
$K^\bullet[1]$，而应用引理 \ref{derived-lemma-functorial-cone}
得到映射 $K^\bullet[1] \to C(g)^\bullet$。应用 $[-1]$
便得到 (2) 中的映射。
\end{proof}

\noindent
注意，引理 \ref{derived-lemma-map-from-cone} 的证明中构造的态射
$C(f)^\bullet \to M^\bullet$ 和 $K^\bullet \to C(g)^\bullet[-1]$
一般依赖于所选的同伦。

\begin{definition}
\label{derived-definition-termwise-split-map}
设 $\mathcal{A}$ 为加性范畴。
{\it 逐项分裂单射 $\alpha : A^\bullet \to B^\bullet$}
是一个复形态射，使得每个 $A^n \to B^n$
都同构于某个直和项的嵌入。
{\it 逐项分裂满射 $\beta : B^\bullet \to C^\bullet$}
是一个复形态射，使得每个 $B^n \to C^n$
都同构于到某个直和项的投影。
\end{definition}

\begin{lemma}
\label{derived-lemma-make-commute-map}
设 $\mathcal{A}$ 为加性范畴。设
$$
\xymatrix{
A^\bullet \ar[r]_f \ar[d]_a & B^\bullet \ar[d]^b \\
C^\bullet \ar[r]^g & D^\bullet
}
$$
是一个交换至同伦的复形态射图。
若 $f$ 是逐项分裂单射，则 $b$ 同伦于一个使该图交换的态射。
若 $g$ 是逐项分裂满射，则 $a$ 同伦于一个使该图交换的态射。
\end{lemma}

\begin{proof}
设 $h^n : A^n \to D^{n - 1}$ 为一族态射，满足
$bf - ga = dh + hd$。假设 $\pi^n : B^n \to A^n$
是分裂态射 $f^n$ 的态射。取
$b' = b - dh\pi - h\pi d$。假设
$s^n : D^n \to C^n$ 是分裂态射
$g^n : C^n \to D^n$ 的态射。取 $a' = a + dsh + shd$。
计算从略。
\end{proof}

\noindent
下一引理可用来将复形态射替换为一个复形态射，
使它在每一次数上都是直和项的嵌入。

\begin{lemma}
\label{derived-lemma-make-injective}
设 $\mathcal{A}$ 为加性范畴。设
$\alpha : K^\bullet \to L^\bullet$ 为 $\mathcal{A}$ 的复形的态射。
存在分解
$$
\xymatrix{
K^\bullet \ar[r]^{\tilde \alpha} \ar@/_1pc/[rr]_\alpha &
\tilde L^\bullet \ar[r]^\pi &
L^\bullet
}
$$
使得
\begin{enumerate}
\item $\tilde \alpha$ 是逐项分裂单射（见定义
\ref{derived-definition-termwise-split-map}）；
\item 存在复形映射 $s : L^\bullet \to \tilde L^\bullet$，使得
$\pi \circ s = \text{id}_{L^\bullet}$，并且
$s \circ \pi$ 与 $\text{id}_{\tilde L^\bullet}$ 同伦。
\end{enumerate}
此外，若 $K^\bullet$ 和 $L^\bullet$ 均属于
$K^{+}(\mathcal{A})$、$K^{-}(\mathcal{A})$ 或 $K^b(\mathcal{A})$，
则 $\tilde L^\bullet$ 也属于同一范畴。
\end{lemma}

\begin{proof}
我们令
$$
\tilde L^n = L^n \oplus K^n \oplus K^{n + 1}
$$
并定义
$$
d^n_{\tilde L} =
\left(
\begin{matrix}
d^n_L & 0 & 0 \\
0 & d^n_K & \text{id}_{K^{n + 1}} \\
0 & 0 & -d^{n + 1}_K
\end{matrix}
\right)
$$
换言之，$\tilde L^\bullet = L^\bullet \oplus C(1_{K^\bullet})$。
再令
$$
\tilde \alpha =
\left(
\begin{matrix}
\alpha \\
\text{id}_{K^n} \\
0
\end{matrix}
\right)
$$
这显然是分裂单射，而且显然定义了复形态射。定义
$$
\pi =
\left(
\begin{matrix}
\text{id}_{L^n} &
0 &
0
\end{matrix}
\right)
$$
于是显然有 $\pi \circ \tilde \alpha = \alpha$。我们令
$$
s =
\left(
\begin{matrix}
\text{id}_{L^n} \\
0 \\
0
\end{matrix}
\right)
$$
于是 $\pi \circ s = \text{id}_{L^\bullet}$。最后，设
$h^n : \tilde L^n \to \tilde L^{n - 1}$ 是如下映射：
它通过恒等态射将 $\tilde L^n$ 的直和项 $K^n$
映到 $\tilde L^{n - 1}$ 的直和项 $K^n$。那么很容易（见下面备注中的计算）
证明
$$
\text{id}_{\tilde L^\bullet} - s \circ \pi
=
d \circ h + h \circ d
$$
这就证完了引理。
\end{proof}

\begin{remark}
\label{derived-remark-compute-modules}
为了验证上述证明中最后的等式，可以如下逐元论证。我们有
$s\pi(l, k, k^{+}) = (l, 0, 0)$。
因此等式左边的态射将
$(l, k, k^{+})$ 映到 $(0, k, k^{+})$。
另一方面，$h(l, k, k^{+}) = (0, 0, k)$ 且
$d(l, k, k^{+}) = (dl, dk + k^{+}, -dk^{+})$。
因此 $(dh + hd)(l, k, k^{+}) =
d(0, 0, k) + h(dl, dk + k^{+}, -dk^{+}) =
(0, k, -dk) + (0, 0, dk + k^{+}) = (0, k, k^{+})$，
正如所需。
\end{remark}

\begin{lemma}
\label{derived-lemma-make-surjective}
设 $\mathcal{A}$ 为加性范畴。设
$\alpha : K^\bullet \to L^\bullet$ 为 $\mathcal{A}$ 的复形的态射。
存在分解
$$
\xymatrix{
K^\bullet \ar[r]^i \ar@/_1pc/[rr]_\alpha &
\tilde K^\bullet \ar[r]^{\tilde \alpha} &
L^\bullet
}
$$
使得
\begin{enumerate}
\item $\tilde \alpha$ 是逐项分裂满射（见定义
\ref{derived-definition-termwise-split-map}）；
\item 存在复形映射 $s : \tilde K^\bullet \to K^\bullet$，使得
$s \circ i = \text{id}_{K^\bullet}$，并且
$i \circ s$ 与 $\text{id}_{\tilde K^\bullet}$ 同伦。
\end{enumerate}
此外，若 $K^\bullet$ 和 $L^\bullet$ 均属于
$K^{+}(\mathcal{A})$、$K^{-}(\mathcal{A})$ 或 $K^b(\mathcal{A})$，
则 $\tilde K^\bullet$ 也属于同一范畴。
\end{lemma}

\begin{proof}
这是引理 \ref{derived-lemma-make-injective} 的对偶命题。取
$$
\tilde K^n = K^n \oplus L^{n - 1} \oplus L^n
$$
并定义
$$
d^n_{\tilde K} =
\left(
\begin{matrix}
d^n_K & 0 & 0 \\
0 & - d^{n - 1}_L & \text{id}_{L^n} \\
0 & 0 & d^n_L
\end{matrix}
\right)
$$
换言之，$\tilde K^\bullet = K^\bullet \oplus C(1_{L^\bullet[-1]})$。
再令
$$
\tilde \alpha =
\left(
\begin{matrix}
\alpha &
0 &
\text{id}_{L^n}
\end{matrix}
\right)
$$
这显然是分裂满射，而且显然定义了复形态射。定义
$$
i =
\left(
\begin{matrix}
\text{id}_{K^n} \\
0 \\
0
\end{matrix}
\right)
$$
于是显然有 $\tilde \alpha \circ i = \alpha$。我们令
$$
s =
\left(
\begin{matrix}
\text{id}_{K^n} &
0 &
0
\end{matrix}
\right)
$$
于是 $s \circ i = \text{id}_{K^\bullet}$。最后，设
$h^n : \tilde K^n \to \tilde K^{n - 1}$ 是如下映射：
它通过恒等态射将 $\tilde K^n$ 的直和项 $L^{n - 1}$
映到 $\tilde K^{n - 1}$ 的直和项 $L^{n - 1}$。那么很容易证明
$$
\text{id}_{\tilde K^\bullet} - i \circ s
=
d \circ h + h \circ d
$$
这就证完了引理。
\end{proof}

\begin{definition}
\label{derived-definition-split-ses}
设 $\mathcal{A}$ 为加性范畴。
$\mathcal{A}$ 的复形的{\it 逐项分裂正合序列}
是一个复形之复形
$$
0 \to
A^\bullet \xrightarrow{\alpha}
B^\bullet \xrightarrow{\beta}
C^\bullet \to 0
$$
连同给定的直和分解 $B^n = A^n \oplus C^n$，
并且这些分解与 $\alpha^n$ 和 $\beta^n$ 相容。
我们常将直和分解诱导的映射记为
$s^n : C^n \to B^n$ 和 $\pi^n : B^n \to A^n$。
根据同调，引理 \ref{homology-lemma-ses-termwise-split-cochain}，
我们得到相应的复形态射
$$
\delta : C^\bullet \longrightarrow A^\bullet[1]
$$
它在次数 $n$ 上的映射为
$\pi^{n + 1} \circ d_B^n \circ s^n$。换言之，
$(A^\bullet, B^\bullet, C^\bullet, \alpha, \beta, \delta)$
构成三角
$$
A^\bullet \to B^\bullet \to C^\bullet \to A^\bullet[1]
$$
它称为{\it 与该复形的逐项分裂序列相关联的三角}。
\end{definition}

\begin{lemma}
\label{derived-lemma-triangle-independent-splittings}
设 $\mathcal{A}$ 为加性范畴。设
$0 \to A^\bullet \to B^\bullet \to C^\bullet \to 0$
为定义 \ref{derived-definition-split-ses} 中的逐项分裂正合序列。
设 $(\pi')^n$、$(s')^n$ 为另一组分裂。以
$\delta' : C^\bullet \longrightarrow A^\bullet[1]$ 表示
与这第二组分裂相关联的态射。那么
$$
(1, 1, 1) :
(A^\bullet, B^\bullet, C^\bullet, \alpha, \beta, \delta)
\longrightarrow
(A^\bullet, B^\bullet, C^\bullet, \alpha, \beta, \delta')
$$
是 $K(\mathcal{A})$ 中的三角同构。
\end{lemma}

\begin{proof}
该命题无非是说 $\delta$ 和 $\delta'$ 是同伦的复形映射。
这就是同调，引理
\ref{homology-lemma-ses-termwise-split-homotopy-cochain}。
\end{proof}

\begin{remark}
\label{derived-remark-make-commute}
设 $\mathcal{A}$ 为加性范畴。设
$0 \to A_i^\bullet \to B_i^\bullet \to C_i^\bullet \to 0$、$i = 1, 2$
为逐项分裂正合序列。假设
$a : A_1^\bullet \to A_2^\bullet$、
$b : B_1^\bullet \to B_2^\bullet$ 和
$c : C_1^\bullet \to C_2^\bullet$ 是复形态射，使得
$$
\xymatrix{
A_1^\bullet \ar[d]_a \ar[r] &
B_1^\bullet \ar[r] \ar[d]_b &
C_1^\bullet \ar[d]_c \\
A_2^\bullet \ar[r] & B_2^\bullet \ar[r] & C_2^\bullet
}
$$
在 $K(\mathcal{A})$ 中交换。一般而言，{\bf 不}存在与 $b$
同伦的态射 $b' : B_1^\bullet \to B_2^\bullet$，
使上图在复形范畴中交换。事实上，考虑例，等式
(\ref{examples-equation-commutes-up-to-homotopy})。
若能将其中的中间映射替换为一个与它同伦且使该图交换的映射，
那么我们就会得到并不成立的迹的可加性。
\end{remark}

\begin{lemma}
\label{derived-lemma-nilpotent}
设 $\mathcal{A}$ 为加性范畴。设
$0 \to A_i^\bullet \to B_i^\bullet \to C_i^\bullet \to 0$、$i = 1, 2, 3$
为复形的逐项分裂正合序列。设
$b : B_1^\bullet \to B_2^\bullet$ 和 $b' : B_2^\bullet \to B_3^\bullet$
为复形态射，使得
$$
\vcenter{
\xymatrix{
A_1^\bullet \ar[d]_0 \ar[r] &
B_1^\bullet \ar[r] \ar[d]_b &
C_1^\bullet \ar[d]_0 \\
A_2^\bullet \ar[r] & B_2^\bullet \ar[r] & C_2^\bullet
}
}
\quad\text{且}\quad
\vcenter{
\xymatrix{
A_2^\bullet \ar[d]^0 \ar[r] &
B_2^\bullet \ar[r] \ar[d]^{b'} &
C_2^\bullet \ar[d]^0 \\
A_3^\bullet \ar[r] & B_3^\bullet \ar[r] & C_3^\bullet
}
}
$$
在 $K(\mathcal{A})$ 中交换。那么
$b' \circ b = 0$ 在 $K(\mathcal{A})$ 中成立。
\end{lemma}

\begin{proof}
根据引理 \ref{derived-lemma-make-commute-map}，我们可以用同伦的映射
替换 $b$ 和 $b'$，使左图的右方块交换、右图的左方块交换。
换言之，我们有
$\Im(b^n) \subset \Im(A_2^n \to B_2^n)$ 以及
$\Ker((b')^n) \supset \Im(A_2^n \to B_2^n)$。
因而 $b' \circ b = 0$ 作为复形映射成立。
\end{proof}

\begin{lemma}
\label{derived-lemma-third-isomorphism}
设 $\mathcal{A}$ 为加性范畴。设
$f_1 : K_1^\bullet \to L_1^\bullet$ 和
$f_2 : K_2^\bullet \to L_2^\bullet$ 为复形态射。设
$$
(a, b, c) :
(K_1^\bullet, L_1^\bullet, C(f_1)^\bullet, f_1, i_1, p_1)
\longrightarrow
(K_2^\bullet, L_2^\bullet, C(f_2)^\bullet, f_2, i_2, p_2)
$$
是 $K(\mathcal{A})$ 中任意的三角态射。
若 $a$ 和 $b$ 是同伦等价，则 $c$ 也是同伦等价。
\end{lemma}

\begin{proof}
设 $a^{-1} : K_2^\bullet \to K_1^\bullet$ 为复形态射，
它在 $K(\mathcal{A})$ 中是 $a$ 的逆。设
$b^{-1} : L_2^\bullet \to L_1^\bullet$ 为复形态射，
它在 $K(\mathcal{A})$ 中是 $b$ 的逆。设
$c' : C(f_2)^\bullet \to C(f_1)^\bullet$ 是把引理
\ref{derived-lemma-functorial-cone} 应用于
$f_1 \circ a^{-1} = b^{-1} \circ f_2$ 所得的态射。
若能证明 $c \circ c'$ 和 $c' \circ c$ 是 $K(\mathcal{A})$ 中的同构，
则即得结论。因而只需证明下述命题：对于 $K(\mathcal{A})$ 中任意的三角态射
$(1, 1, c) : (K^\bullet, L^\bullet, C(f)^\bullet, f, i, p)$，
态射 $c$ 是 $K(\mathcal{A})$ 中的同构。
由假设，图中的两个方块
$$
\xymatrix{
L^\bullet \ar[r] \ar[d]_1 &
C(f)^\bullet \ar[r] \ar[d]_c &
K^\bullet[1] \ar[d]_1 \\
L^\bullet \ar[r] &
C(f)^\bullet \ar[r] &
K^\bullet[1]
}
$$
交换至同伦。由 $C(f)^\bullet$ 的构造，两行均构成复形的逐项分裂序列。
因此，由引理 \ref{derived-lemma-nilpotent}，我们看到
$(c - 1)^2 = 0$ 在 $K(\mathcal{A})$ 中成立。
故 $c$ 是 $K(\mathcal{A})$ 中的同构，其逆为 $2 - c$。
\end{proof}

\noindent
因此，若 $a$ 和 $b$ 是同伦等价，则所得的三角态射
是 $K(\mathcal{A})$ 中的三角同构。
事实上，$K(\mathcal{A})$ 中由锥给出的三角之类，
与 $K(\mathcal{A})$ 中由复形的逐项分裂序列给出的三角之类，
在同构意义下是相同的，至少相差一个符号！

\begin{lemma}
\label{derived-lemma-the-same-up-to-isomorphisms}
设 $\mathcal{A}$ 为加性范畴。
\begin{enumerate}
\item 给定复形的逐项分裂序列
$(\alpha : A^\bullet \to B^\bullet,
\beta : B^\bullet \to C^\bullet, s^n, \pi^n)$，
存在同伦等价 $C(\alpha)^\bullet \to C^\bullet$，使图
$$
\xymatrix{
A^\bullet \ar[r] \ar[d] & B^\bullet \ar[d] \ar[r] &
C(\alpha)^\bullet \ar[r]_{-p} \ar[d] & A^\bullet[1] \ar[d] \\
A^\bullet \ar[r] & B^\bullet \ar[r] &
C^\bullet \ar[r]^\delta & A^\bullet[1]
}
$$
在 $K(\mathcal{A})$ 中定义一个三角同构。
\item 给定复形态射 $f : K^\bullet \to L^\bullet$，
存在三角同构
$$
\xymatrix{
K^\bullet \ar[r] \ar[d] & \tilde L^\bullet \ar[d] \ar[r] &
M^\bullet \ar[r]_{\delta} \ar[d] & K^\bullet[1] \ar[d] \\
K^\bullet \ar[r] & L^\bullet \ar[r] &
C(f)^\bullet \ar[r]^{-p} & K^\bullet[1]
}
$$
其中上方三角是与某个逐项分裂正合序列
$K^\bullet \to \tilde L^\bullet \to M^\bullet$ 相关联的三角。
\end{enumerate}
\end{lemma}

\begin{proof}
证明 (1)。我们有 $C(\alpha)^n = B^n \oplus A^{n + 1}$，
并将 $C(\alpha)^n \to C^n$ 简单定义为先投影到 $B^n$，
再应用 $\beta^n$。这定义了复形态射，因为复合
$A^{n + 1} \to B^{n + 1} \to C^{n + 1}$ 都为零。
为得到同伦逆，我们取 $C^\bullet \to C(\alpha)^\bullet$，
它在次数 $n$ 上由 $(s^n , -\delta^n)$ 给出。
这是复形态射，因为态射 $\delta^n$
可表征为使下式成立的唯一态射 $C^n \to A^{n + 1}$：
$d \circ s^n - s^{n + 1} \circ d = \alpha \circ \delta^n$，
见同调，引理 \ref{homology-lemma-ses-termwise-split-cochain} 的证明。
复合 $C^\bullet \to C(\alpha)^\bullet \to C^\bullet$ 是恒等态射。
复合 $C(\alpha)^\bullet \to C^\bullet \to C(\alpha)^\bullet$
等于态射
$$
\left(
\begin{matrix}
s^n \circ \beta^n & 0 \\
-\delta^n \circ \beta^n & 0
\end{matrix}
\right)
$$
为了说明它与恒等映射同伦，使用由矩阵
$$
\left(
\begin{matrix}
0 & 0 \\
\pi^n & 0
\end{matrix}
\right) : C(\alpha)^n = B^n \oplus A^{n + 1} \to
B^{n - 1} \oplus A^n = C(\alpha)^{n - 1}
$$
给出的同伦 $h^n : C(\alpha)^n \to C(\alpha)^{n - 1}$。
直接验证即得
$$
\left(
\begin{matrix}
1 & 0 \\
0 & 1
\end{matrix}
\right)
-
\left(
\begin{matrix}
s^n \\
-\delta^n
\end{matrix}
\right)
\left(
\begin{matrix}
\beta^n & 0
\end{matrix}
\right)
=
\left(
\begin{matrix}
d & \alpha^n \\
0 & -d
\end{matrix}
\right)
\left(
\begin{matrix}
0 & 0 \\
\pi^n & 0
\end{matrix}
\right)
+
\left(
\begin{matrix}
0 & 0 \\
\pi^{n + 1} & 0
\end{matrix}
\right)
\left(
\begin{matrix}
d & \alpha^{n + 1} \\
0 & -d
\end{matrix}
\right)
$$
为完成 (1) 的证明，还需说明态射
$-p : C(\alpha)^\bullet \to A^\bullet[1]$（见定义
\ref{derived-definition-cone}）与 $C(\alpha)^\bullet \to C^\bullet \to A^\bullet[1]$
在同伦意义下相等。这从上述讨论是清楚的。事实上，
可以使用同伦逆 $(s, -\delta) : C^\bullet \to C(\alpha)^\bullet$，
转而验证两个映射 $C^\bullet \to A^\bullet[1]$ 相等。又注意到
$p \circ (s, -\delta) = -\delta$，正如所需。

\medskip\noindent
证明 (2)。设 $\tilde f : K^\bullet \to \tilde L^\bullet$、
$s : L^\bullet \to \tilde L^\bullet$ 和
$\pi : \tilde L^\bullet \to L^\bullet$ 如引理
\ref{derived-lemma-make-injective} 中所定义。由引理
\ref{derived-lemma-functorial-cone} 和 \ref{derived-lemma-third-isomorphism}，三角
$(K^\bullet, L^\bullet, C(f), i, p)$ 和
$(K^\bullet, \tilde L^\bullet, C(\tilde f), \tilde i, \tilde p)$
同构。注意，三角同构可以复合。因此可以用
$\tilde L^\bullet$ 替换 $L^\bullet$，用 $\tilde f$ 替换 $f$。
换言之，可假设 $f$ 是逐项分裂单射。
在此情形下，结论由 (1) 得到。
\end{proof}


\begin{lemma}
\label{derived-lemma-sequence-maps-split}
设 $\mathcal{A}$ 为加性范畴。设
$A_1^\bullet \to A_2^\bullet \to \ldots \to A_n^\bullet$
为一列可复合的复形态射。存在交换图
$$
\xymatrix{
A_1^\bullet \ar[r] &
A_2^\bullet \ar[r] &
\ldots \ar[r] &
A_n^\bullet \\
B_1^\bullet \ar[r] \ar[u] &
B_2^\bullet \ar[r] \ar[u] &
\ldots \ar[r] &
B_n^\bullet \ar[u]
}
$$
使得每个态射 $B_i^\bullet \to B_{i + 1}^\bullet$
都是逐项分裂单射，而每个 $B_i^\bullet \to A_i^\bullet$
都是同伦等价。此外，若所有 $A_i^\bullet$ 都属于
$K^{+}(\mathcal{A})$、$K^{-}(\mathcal{A})$ 或 $K^b(\mathcal{A})$，
则所有 $B_i^\bullet$ 也属于同一范畴。
\end{lemma}

\begin{proof}
$n = 1$ 的情形没有实质内容。引理
\ref{derived-lemma-make-injective} 是 $n = 2$ 的情形。
假设除 $B_n^\bullet$ 之外的图已经构造出来。
将引理 \ref{derived-lemma-make-injective} 应用于复合
$B_{n - 1}^\bullet \to A_{n - 1}^\bullet \to A_n^\bullet$。
所得结果是所需的分解
$B_{n - 1}^\bullet \to B_n^\bullet \to A_n^\bullet$。
\end{proof}

\begin{lemma}
\label{derived-lemma-rotate-triangle}
设 $\mathcal{A}$ 为加性范畴。设
$(\alpha : A^\bullet \to B^\bullet, \beta : B^\bullet \to C^\bullet, s^n,
\pi^n)$ 为复形的逐项分裂序列。设
$(A^\bullet, B^\bullet, C^\bullet, \alpha, \beta, \delta)$
为相关联的三角。那么三角
$(C^\bullet[-1], A^\bullet, B^\bullet, \delta[-1], \alpha, \beta)$
同构于三角
$(C^\bullet[-1], A^\bullet, C(\delta[-1])^\bullet, \delta[-1], i, p)$。
\end{lemma}

\begin{proof}
我们写作 $B^n = A^n \oplus C^n$，并将 $\alpha^n$ 和 $\beta^n$
分别识别为自然嵌入与投影。由 $\delta$ 的构造，有
$$
d_B^n =
\left(
\begin{matrix}
d_A^n & \delta^n \\
0 & d_C^n
\end{matrix}
\right)
$$
另一方面，$\delta[-1] : C^\bullet[-1] \to A^\bullet$ 的锥满足
$C(\delta[-1])^n = A^n \oplus C^n$，其微分正是上述矩阵！
因此引理得证。
\end{proof}

\begin{lemma}
\label{derived-lemma-rotate-cone}
设 $\mathcal{A}$ 为加性范畴。设
$f : K^\bullet \to L^\bullet$ 为复形态射。三角
$(L^\bullet, C(f)^\bullet, K^\bullet[1], i, p, f[1])$ 是与
由 $f$ 的锥的定义得到的逐项分裂序列
$$
0 \to L^\bullet \to C(f)^\bullet \to K^\bullet[1] \to 0
$$
相关联的三角。
\end{lemma}

\begin{proof}
立即由定义得到。
\end{proof}






\section{同伦范畴中的特异三角}
\label{derived-section-homotopy-triangulated}

\noindent
由于我们希望上同调长正合序列中的边缘映射
由蛇引理中不带符号的映射给出，我们如下定义同伦范畴中的特异三角。

\begin{definition}
\label{derived-definition-distinguished-triangle}
设 $\mathcal{A}$ 为加性范畴。$K(\mathcal{A})$ 的三角
$(X, Y, Z, f, g, h)$ 称为 $K(\mathcal{A})$ 的{\it 特异三角}，
若它同构于与某个复形的逐项分裂正合序列相关联的三角；
见定义 \ref{derived-definition-split-ses}。对
$K^{+}(\mathcal{A})$、$K^{-}(\mathcal{A})$ 和 $K^b(\mathcal{A})$
作相同定义。
\end{definition}

\noindent
注意，根据引理 \ref{derived-lemma-the-same-up-to-isomorphisms}，
形如 $(K^\bullet, L^\bullet, C(f)^\bullet, f, i, -p)$ 的三角
是特异三角。这确实给出一个三角范畴，见命题
\ref{derived-proposition-homotopy-category-triangulated}。在证明该命题之前，
还需要一个引理才能证明 TR4。

\begin{lemma}
\label{derived-lemma-two-split-injections}
设 $\mathcal{A}$ 为加性范畴。假设
$\alpha : A^\bullet \to B^\bullet$ 和 $\beta : B^\bullet \to C^\bullet$
是复形的分裂单射。那么存在特异三角
$(A^\bullet, B^\bullet, Q_1^\bullet, \alpha, p_1, d_1)$、
$(A^\bullet, C^\bullet, Q_2^\bullet, \beta \circ \alpha, p_2, d_2)$
和
$(B^\bullet, C^\bullet, Q_3^\bullet, \beta, p_3, d_3)$，
对它们 TR4 成立。
\end{lemma}

\begin{proof}
设 $\pi_1^n : B^n \to A^n$ 和 $\pi_3^n : C^n \to B^n$ 是分裂。
那么 $A^\bullet \to C^\bullet$ 也是分裂单射，其分裂为
$\pi_2^n = \pi_1^n \circ \pi_3^n$。以 $Q_1^\bullet$、$Q_2^\bullet$ 和
$Q_3^\bullet$ 表示“商”复形。换言之，
$Q_1^n = \Ker(\pi_1^n)$、$Q_3^n = \Ker(\pi_3^n)$ 且
$Q_2^n = \Ker(\pi_2^n)$。注意这些核存在。于是
$B^n = A^n \oplus Q_1^n$ 且 $C^n = B^n \oplus Q_3^n$，
其中我们把 $A^n$ 看作 $B^n$ 的子对象，依此类推。这意味着
$C^n = A^n \oplus Q_1^n \oplus Q_3^n$。注意
$\pi_2^n = \pi_1^n \circ \pi_3^n$ 在 $Q_1^n$ 和 $Q_3^n$ 上均为零。因而
$Q_2^n = Q_1^n \oplus Q_3^n$。考虑交换图
$$
\begin{matrix}
0 & \to & A^\bullet & \to & B^\bullet & \to & Q_1^\bullet & \to & 0 \\
  &     & \downarrow &     & \downarrow &     & \downarrow  & \\
0 & \to & A^\bullet & \to & C^\bullet & \to & Q_2^\bullet & \to & 0 \\
  &     & \downarrow &     & \downarrow &     & \downarrow  & \\
0 & \to & B^\bullet & \to & C^\bullet & \to & Q_3^\bullet & \to & 0
\end{matrix}
$$
此图各行都是逐项分裂正合序列，因而按定义决定特异三角。
此外，上图中向下的箭头与所选分裂相容，故定义三角态射
$$
(A^\bullet \to B^\bullet \to Q_1^\bullet \to A^\bullet[1])
\longrightarrow
(A^\bullet \to C^\bullet \to Q_2^\bullet \to A^\bullet[1])
$$
和
$$
(A^\bullet \to C^\bullet \to Q_2^\bullet \to A^\bullet[1])
\longrightarrow
(B^\bullet \to C^\bullet \to Q_3^\bullet \to B^\bullet[1]).
$$
注意，上图底行分裂序列的分裂 $Q_3^n \to C^n$
与 $C^n \to Q_2^n$ 复合后，给出分裂序列
$0 \to Q_1^\bullet \to Q_2^\bullet \to Q_3^\bullet \to 0$ 的一个分裂。
由此容易看出，相应特异三角
$$
(Q_1^\bullet \to Q_2^\bullet \to Q_3^\bullet \to Q_1^\bullet[1])
$$
中的态射 $\delta : Q_3^\bullet \to Q_1^\bullet[1]$
等于复合 $Q_3^\bullet \to B^\bullet[1] \to Q_1^\bullet[1]$。
因而我们得到公理 TR4 的结论中所规定的结构。
\end{proof}

\begin{proposition}
\label{derived-proposition-homotopy-category-triangulated}
设 $\mathcal{A}$ 为加性范畴。复形的模同伦范畴
$K(\mathcal{A})$，连同它的自然平移函子和上面定义的特异三角，
是一个三角范畴。
\end{proposition}

\begin{proof}
证明 TR1。依定义，任意同构于特异三角的三角都是特异的。
另外，任意三角 $(A^\bullet, A^\bullet, 0, 1, 0, 0)$ 都是特异的，因为
$0 \to A^\bullet \to A^\bullet \to 0 \to 0$ 是复形的逐项分裂序列。
最后，给定任意复形态射 $f : K^\bullet \to L^\bullet$，三角
$(K, L, C(f), f, i, -p)$ 由引理
\ref{derived-lemma-the-same-up-to-isomorphisms} 是特异的。

\medskip\noindent
证明 TR2。设 $(X, Y, Z, f, g, h)$ 为三角。假设
$(Y, Z, X[1], g, h, -f[1])$ 是特异的。那么存在复形的逐项分裂序列
$A^\bullet \to B^\bullet \to C^\bullet$，使得相关联的三角
$(A^\bullet, B^\bullet, C^\bullet, \alpha, \beta, \delta)$
同构于 $(Y, Z, X[1], g, h, -f[1])$。反向旋转后看到
$(X, Y, Z, f, g, h)$ 同构于
$(C^\bullet[-1], A^\bullet, B^\bullet, -\delta[-1], \alpha, \beta)$。
由引理 \ref{derived-lemma-rotate-triangle}，三角
$(C^\bullet[-1], A^\bullet, B^\bullet, \delta[-1], \alpha, \beta)$
同构于
$(C^\bullet[-1], A^\bullet, C(\delta[-1])^\bullet, \delta[-1], i, p)$。
将前一个三角同构与 $Y$ 上的 $-1$ 预复合，得到
$(X, Y, Z, f, g, h)$ 同构于
$(C^\bullet[-1], A^\bullet, C(\delta[-1])^\bullet, \delta[-1], i, -p)$。
因此它由引理 \ref{derived-lemma-the-same-up-to-isomorphisms} 是特异的。
反过来，假设 $(X, Y, Z, f, g, h)$ 是特异的。
由引理 \ref{derived-lemma-the-same-up-to-isomorphisms}，这意味着它同构于形如
$(K^\bullet, L^\bullet, C(f), f, i, -p)$ 的三角，其中 $f$
是某个复形态射。那么旋转后的三角
$(Y, Z, X[1], g, h, -f[1])$ 同构于
$(L^\bullet, C(f), K^\bullet[1], i, -p, -f[1])$，后者又同构于三角
$(L^\bullet, C(f), K^\bullet[1], i, p, f[1])$。
由引理 \ref{derived-lemma-rotate-cone}，这个三角是特异的。
因此 $(Y, Z, X[1], g, h, -f[1])$ 如所需是特异的。

\medskip\noindent
证明 TR3。设
$(X, Y, Z, f, g, h)$ 和 $(X', Y', Z', f', g', h')$
是 $K(\mathcal{A})$ 的特异三角，并设
$a : X \to X'$ 和 $b : Y \to Y'$ 为满足
$f' \circ a = b \circ f$ 的态射。由引理
\ref{derived-lemma-the-same-up-to-isomorphisms}，可假设
$(X, Y, Z, f, g, h) = (X, Y, C(f), f, i, -p)$ 且
$(X', Y', Z', f', g', h') = (X', Y', C(f'), f', i', -p')$。
此时，只需将引理 \ref{derived-lemma-functorial-cone}
应用于由 $f, f', a, b$ 给出的交换图。

\medskip\noindent
证明 TR4。目前我们已知 $K(\mathcal{A})$ 是预三角范畴。
因此可使用引理 \ref{derived-lemma-easier-axiom-four}。设
$A^\bullet \to B^\bullet$ 和 $B^\bullet \to C^\bullet$
是 $K(\mathcal{A})$ 中可复合的态射。由引理
\ref{derived-lemma-sequence-maps-split}，可假设
$A^\bullet \to B^\bullet$ 和 $B^\bullet \to C^\bullet$
是分裂单射。在此情形下，结论由引理
\ref{derived-lemma-two-split-injections} 得到。
\end{proof}

\begin{remark}
\label{derived-remark-boundedness-conditions-triangulated}
设 $\mathcal{A}$ 为加性范畴。命题
\ref{derived-proposition-homotopy-category-triangulated} 的证明原封不动地表明，
范畴 $K^{+}(\mathcal{A})$、$K^{-}(\mathcal{A})$ 和 $K^b(\mathcal{A})$
都是三角范畴。实际上，有界（上有界、下有界）复形之间的态射的锥
仍然有界（上有界、下有界）。不过下面我们将证明它们是
$K(\mathcal{A})$ 的三角子范畴，这给出另一个证明。
\end{remark}

\begin{lemma}
\label{derived-lemma-bounded-triangulated-subcategories}
设 $\mathcal{A}$ 为加性范畴。范畴
$K^{+}(\mathcal{A})$、$K^{-}(\mathcal{A})$ 和 $K^b(\mathcal{A})$
都是 $K(\mathcal{A})$ 的满三角子范畴。
\end{lemma}

\begin{proof}
所提及的每个范畴都是满加性子范畴。我们使用引理
\ref{derived-lemma-triangulated-subcategory} 的判别准则说明它们是三角子范畴。
显然，每个范畴
$K^{+}(\mathcal{A})$、$K^{-}(\mathcal{A})$ 和 $K^b(\mathcal{A})$
都在平移函子 $[1], [-1]$ 下保持。最后，假设
$f : A^\bullet \to B^\bullet$ 是
$K^{+}(\mathcal{A})$、$K^{-}(\mathcal{A})$ 或 $K^b(\mathcal{A})$ 中的态射。
那么 $(A^\bullet, B^\bullet, C(f)^\bullet, f, i, -p)$
是 $K(\mathcal{A})$ 的特异三角，且
$C(f)^\bullet \in K^{+}(\mathcal{A})$、
$K^{-}(\mathcal{A})$ 或 $K^b(\mathcal{A})$，这从锥的构造是明显的。
因此引理得证。（或者，
% P02-E0026: frozen source says "an termwise split injection".
$K^\bullet \to L^\bullet$ 在
$K^{+}(\mathcal{A})$、$K^{-}(\mathcal{A})$ 或 $K^b(\mathcal{A})$ 中
同构于一个复形的逐项分裂单射，见引理
\ref{derived-lemma-make-injective}，然后可直接取相关联的特异三角。）
\end{proof}

\begin{lemma}
\label{derived-lemma-additive-exact-homotopy-category}
设 $\mathcal{A}$、$\mathcal{B}$ 为加性范畴，
$F : \mathcal{A} \to \mathcal{B}$ 为加性函子。所诱导的函子
$$
\begin{matrix}
F : K(\mathcal{A}) \longrightarrow K(\mathcal{B}) \\
F : K^{+}(\mathcal{A}) \longrightarrow K^{+}(\mathcal{B}) \\
F : K^{-}(\mathcal{A}) \longrightarrow K^{-}(\mathcal{B}) \\
F : K^b(\mathcal{A}) \longrightarrow K^b(\mathcal{B})
\end{matrix}
$$
都是三角范畴的正合函子。
\end{lemma}

\begin{proof}
假设 $A^\bullet \to B^\bullet \to C^\bullet$
是 $\mathcal{A}$ 的复形的逐项分裂序列，其分裂为
$(s^n, \pi^n)$，相关联的态射为
$\delta : C^\bullet \to A^\bullet[1]$；见定义
\ref{derived-definition-split-ses}。那么
$F(A^\bullet) \to F(B^\bullet) \to F(C^\bullet)$
是复形的逐项分裂序列，其分裂为
$(F(s^n), F(\pi^n))$，相关联的态射为
$F(\delta) : F(C^\bullet) \to F(A^\bullet)[1]$。
因而 $F$ 将特异三角变换为特异三角。
\end{proof}

\begin{lemma}
\label{derived-lemma-improve-distinguished-triangle-homotopy}
设 $\mathcal{A}$ 为加性范畴，并设
$(A^\bullet, B^\bullet, C^\bullet, a, b, c)$ 为 $K(\mathcal{A})$
中的特异三角。那么存在一个同构的特异三角
$(A^\bullet, (B')^\bullet, C^\bullet, a', b', c)$，使对所有 $n$，
$0 \to A^n \to (B')^n \to C^n \to 0$ 都是分裂短正合序列。
\end{lemma}

\begin{proof}
我们将使用命题 \ref{derived-proposition-homotopy-category-triangulated}
所证的 $K(\mathcal{A})$ 是三角范畴这一事实。
设 $W^\bullet$ 是 $c : C^\bullet \to A^\bullet[1]$ 的锥，
其映射为 $i : A^\bullet[1] \to W^\bullet$ 和
$p : W^\bullet \to C^\bullet[1]$。那么由引理
\ref{derived-lemma-the-same-up-to-isomorphisms}，
$(C^\bullet, A^\bullet[1], W^\bullet, c, i, -p)$ 是特异三角。
向后旋转两次，看到
$(A^\bullet, W^\bullet[-1], C^\bullet, -i[-1], p[-1], c)$
是特异三角。由 TR3，存在特异三角的态射
$(\text{id}, \beta, \text{id}) : (A^\bullet, B^\bullet, C^\bullet, a, b, c) \to
(A^\bullet, W^\bullet[-1], C^\bullet, -i[-1], p[-1], c)$，
由引理 \ref{derived-lemma-third-isomorphism-triangle}，它必为同构。
这就完成了证明，因为根据第 \ref{derived-section-cones} 节中锥的构造，
$0 \to A^\bullet \to W^\bullet[-1] \to C^\bullet \to 0$
是复形的逐项分裂短正合序列。
\end{proof}

\begin{remark}
\label{derived-remark-homotopy-double}
设 $\mathcal{A}$ 为具有可数直和的加性范畴。以
$\text{DoubleComp}(\mathcal{A})$ 表示 $\mathcal{A}$ 中的双复形范畴，
见同调，第 \ref{homology-section-double-complexes} 节。
我们可用这个范畴构造两个三角范畴。
\begin{enumerate}
\item 我们可将 $\text{DoubleComp}(\mathcal{A})$ 的对象
$A^{\bullet, \bullet}$ 如下看作复形之复形：
$$
\ldots \to A^{\bullet, -1} \to A^{\bullet, 0} \to A^{\bullet, 1} \to \ldots
$$
并取同伦范畴 $K_{first}(\text{DoubleComp}(\mathcal{A}))$，
其相应的三角结构由命题
\ref{derived-proposition-homotopy-category-triangulated} 给出。
由同调，备注
\ref{homology-remark-double-complex-complex-of-complexes-first}，函子
$$
\text{Tot} :
K_{first}(\text{DoubleComp}(\mathcal{A}))
\longrightarrow
K(\mathcal{A})
$$
是三角范畴的正合函子。
\item 我们可将 $\text{DoubleComp}(\mathcal{A})$ 的对象
$A^{\bullet, \bullet}$ 如下看作复形之复形：
$$
\ldots \to A^{-1, \bullet} \to A^{0, \bullet} \to A^{1, \bullet} \to \ldots
$$
并取同伦范畴 $K_{second}(\text{DoubleComp}(\mathcal{A}))$，
其相应的三角结构由命题
\ref{derived-proposition-homotopy-category-triangulated} 给出。
由同调，备注
\ref{homology-remark-double-complex-complex-of-complexes-second}，函子
$$
\text{Tot} :
K_{second}(\text{DoubleComp}(\mathcal{A}))
\longrightarrow
K(\mathcal{A})
$$
是三角范畴的正合函子。
\end{enumerate}
\end{remark}

\begin{remark}
\label{derived-remark-double-complex-as-tensor-product-of}
设 $\mathcal{A}$、$\mathcal{B}$、$\mathcal{C}$ 为加性范畴，
并假设 $\mathcal{C}$ 具有可数直和。假设
$$
\otimes : \mathcal{A} \times \mathcal{B} \longrightarrow \mathcal{C},
\quad
(X, Y) \longmapsto X \otimes Y
$$
是一个在态射上双线性的函子。这决定一个函子
$$
\text{Comp}(\mathcal{A}) \times \text{Comp}(\mathcal{B})
\longrightarrow
\text{DoubleComp}(\mathcal{C}), \quad
(X^\bullet, Y^\bullet)
\longmapsto
X^\bullet \otimes Y^\bullet
$$
见同调，例
\ref{homology-example-double-complex-as-tensor-product-of}。
\begin{enumerate}
\item 对于 $\text{Comp}(\mathcal{A})$ 的一个固定对象 $X^\bullet$，函子
$$
K(\mathcal{B}) \longrightarrow K(\mathcal{C}), \quad
Y^\bullet \longmapsto \text{Tot}(X^\bullet \otimes Y^\bullet)
$$
是三角范畴的正合函子。
\item 对于 $\text{Comp}(\mathcal{B})$ 的一个固定对象 $Y^\bullet$，函子
$$
K(\mathcal{A}) \longrightarrow K(\mathcal{C}), \quad
X^\bullet \longmapsto \text{Tot}(X^\bullet \otimes Y^\bullet)
$$
是三角范畴的正合函子。
\end{enumerate}
这由备注 \ref{derived-remark-homotopy-double} 得到，因为函子
$\text{Comp}(\mathcal{A}) \to \text{DoubleComp}(\mathcal{C})$、
$Y^\bullet \mapsto X^\bullet \otimes Y^\bullet$ 以及
$\text{Comp}(\mathcal{B}) \to \text{DoubleComp}(\mathcal{C})$、
$X^\bullet \mapsto X^\bullet \otimes Y^\bullet$
显然与同伦和逐项分裂短正合序列相容，因而诱导三角范畴的正合函子
$$
K(\mathcal{B}) \to K_{first}(\text{DoubleComp}(\mathcal{C}))
\quad\text{且}\quad
K(\mathcal{A}) \to K_{second}(\text{DoubleComp}(\mathcal{C}))
$$
注意，对于这两个函子中的第一个，同构
$$
\text{Tot}(X^\bullet \otimes Y^\bullet[1]) \cong
\text{Tot}(X^\bullet \otimes Y^\bullet)[1]
$$
涉及符号（这可追溯到同调，备注
\ref{homology-remark-shift-double-complex} 中所选的符号）。
\end{remark}






\section{导出范畴}
\label{derived-section-derived-categories}

\noindent
本节通过在 $K(\mathcal{A})$ 中取拟同构的逆，构造阿贝尔范畴
$\mathcal{A}$ 的导出范畴。为此先回忆，函子
$H^i : \text{Comp}(\mathcal{A}) \to \mathcal{A}$
经由 $K(\mathcal{A})$ 分解，见同调，引理
\ref{homology-lemma-map-cohomology-homotopy-cochain}。此外，在同调，定义
\ref{homology-definition-cohomology-shift} 中，我们已定义识别
$H^i(K^\bullet[n]) = H^{i + n}(K^\bullet)$。此时，为了避免混淆和可能的符号错误，
将其重新定义为
$$
H^i(K^\bullet) = H^0(K^\bullet[i])
$$
是合理的。

\begin{lemma}
\label{derived-lemma-cohomology-homological}
设 $\mathcal{A}$ 为阿贝尔范畴。函子
$$
H^0 : K(\mathcal{A}) \longrightarrow \mathcal{A}
$$
是同调函子。
\end{lemma}

\begin{proof}
由于 $H^0$ 是函子，按我们对特异三角的定义，只需证明：
对给定的复形的逐项分裂短正合序列
$0 \to A^\bullet \to B^\bullet \to C^\bullet \to 0$，序列
$H^0(A^\bullet) \to H^0(B^\bullet) \to H^0(C^\bullet)$
正合。这由同调，引理
\ref{homology-lemma-long-exact-sequence-cochain} 得到。
\end{proof}

\noindent
特别地，该引理意味着，$K(\mathcal{A})$ 中的特异三角
$(X, Y, Z, f, g, h)$ 诱导上同调长正合序列
\begin{equation}
\label{derived-equation-long-exact-cohomology-sequence-D}
\xymatrix{
\ldots \ar[r] &
H^i(X) \ar[r]^{H^i(f)} &
H^i(Y) \ar[r]^{H^i(g)} &
H^i(Z) \ar[r]^{H^i(h)} &
H^{i + 1}(X) \ar[r] & \ldots
}
\end{equation}
见 (\ref{derived-equation-long-exact-cohomology-sequence})。此外，它与复形短正合序列所相关联的
上同调长正合序列相容
% P02-E0027: frozen source contains the editorial placeholder "insert future reference here".
（此处插入未来的引用）。例如，若
$(A^\bullet, B^\bullet, C^\bullet, \alpha, \beta, \delta)$
是与复形的逐项分裂正合序列相关联的特异三角（见定义
\ref{derived-definition-split-ses}），则上述上同调序列与用蛇引理定义的序列一致；
见同调，引理 \ref{homology-lemma-long-exact-sequence-cochain}，
关于序列的一致性，见同调，引理
\ref{homology-lemma-ses-termwise-split-long-cochain}。

\medskip\noindent
回忆，若对所有 $i \in \mathbf{Z}$ 都有 $H^i(K^\bullet) = 0$，
则称复形 $K^\bullet$ {\it 无环}。又回忆，复形态射
$f : K^\bullet \to L^\bullet$ 是{\it 拟同构}，当且仅当
$H^i(f)$ 对所有 $i$ 都是同构。见同调，定义
\ref{homology-definition-quasi-isomorphism-cochain}。

\begin{lemma}
\label{derived-lemma-acyclic}
设 $\mathcal{A}$ 为阿贝尔范畴。$K(\mathcal{A})$ 中由无环复形组成的满子范畴
$\text{Ac}(\mathcal{A})$ 是 $K(\mathcal{A})$ 的严格满饱和三角子范畴。
$K(\mathcal{A})$ 中相应的饱和乘法系（见引理
\ref{derived-lemma-operations}）是拟同构集 $\text{Qis}(\mathcal{A})$。
特别地，局部化函子
$Q : K(\mathcal{A}) \to \text{Qis}(\mathcal{A})^{-1}K(\mathcal{A})$
的核是 $\text{Ac}(\mathcal{A})$，且函子 $H^0$ 经由 $Q$ 分解。
\end{lemma}

\begin{proof}
由引理 \ref{derived-lemma-cohomology-homological}，我们知道 $H^0$ 是同调函子。
因而本引理是引理 \ref{derived-lemma-acyclic-general} 的特殊情形。
\end{proof}

\begin{definition}
\label{derived-definition-unbounded-derived-category}
设 $\mathcal{A}$ 为阿贝尔范畴。设
$\text{Ac}(\mathcal{A})$ 和 $\text{Qis}(\mathcal{A})$ 如引理
\ref{derived-lemma-acyclic} 中所定义。$\mathcal{A}$ 的{\it 导出范畴}
是三角范畴
$$
D(\mathcal{A}) =
K(\mathcal{A})/\text{Ac}(\mathcal{A}) =
\text{Qis}(\mathcal{A})^{-1} K(\mathcal{A}).
$$
以 $H^0 : D(\mathcal{A}) \to \mathcal{A}$ 表示唯一的函子，
它与商函子的复合恢复为上面定义的函子 $H^0$。使用引理
\ref{derived-lemma-homological-functor-bounded}，我们引入严格满饱和三角子范畴
$D^{+}(\mathcal{A}), D^{-}(\mathcal{A}), D^b(\mathcal{A})$，它们的对象集为
$$
\begin{matrix}
\Ob(D^{+}(\mathcal{A})) =
\{X \in \Ob(D(\mathcal{A})) \mid
H^n(X) = 0\text{ 对所有 }n \ll 0\} \\
\Ob(D^{-}(\mathcal{A})) =
\{X \in \Ob(D(\mathcal{A})) \mid
H^n(X) = 0\text{ 对所有 }n \gg 0\} \\
\Ob(D^b(\mathcal{A})) =
\{X \in \Ob(D(\mathcal{A})) \mid
H^n(X) = 0\text{ 对所有 }|n| \gg 0\}
\end{matrix}
$$
范畴 $D^b(\mathcal{A})$ 称为 $\mathcal{A}$ 的{\it 有界导出范畴}。
\end{definition}

\noindent
若 $K^\bullet$ 和 $L^\bullet$ 是 $\mathcal{A}$ 的复形，那么我们有时说
“$K^\bullet$ 与 $L^\bullet$ {\it 拟同构}”，意指
$K^\bullet$ 和 $L^\bullet$ 是 $D(\mathcal{A})$ 中的同构对象。

\begin{remark}
\label{derived-remark-existence-derived}
本章一贯处理“小”阿贝尔范畴（这是 Stacks 项目中的约定）。
对于“大”阿贝尔范畴 $\mathcal{A}$，导出范畴
$D(\mathcal{A})$ 是否存在并不清楚，因为导出范畴中的态射是否构成集并不清楚。
实际上，它们一般不构成集，见例，引理
\ref{examples-lemma-big-abelian-category}。不过，若 $\mathcal{A}$ 是 Grothendieck 阿贝尔范畴，
且 $K^\bullet, L^\bullet$ 属于 $K(\mathcal{A})$，那么由内射，定理
\ref{injectives-theorem-K-injective-embedding-grothendieck}，存在到 K-内射复形
$I^\bullet$ 的拟同构 $L^\bullet \to I^\bullet$，而引理
\ref{derived-lemma-K-injective} 表明
$$
\Hom_{D(\mathcal{A})}(K^\bullet, L^\bullet) =
\Hom_{K(\mathcal{A})}(K^\bullet, I^\bullet)
$$
是一个集。Grothendieck 阿贝尔范畴的一些例子包括环上模的范畴，
或更一般地，带环址点上的模层范畴。
\end{remark}

\noindent
每个变体 $D^{+}(\mathcal{A}), D^{-}(\mathcal{A}), D^b(\mathcal{A})$
都可构造为相应同伦范畴的局部化。这依赖于下面的简单引理。

\begin{lemma}
\label{derived-lemma-complex-cohomology-bounded}
设 $\mathcal{A}$ 为阿贝尔范畴，$K^\bullet$ 为复形。
\begin{enumerate}
\item 若对所有 $n \ll 0$ 都有 $H^n(K^\bullet) = 0$，
则存在拟同构 $K^\bullet \to L^\bullet$，其中 $L^\bullet$ 下有界。
\item 若对所有 $n \gg 0$ 都有 $H^n(K^\bullet) = 0$，
则存在拟同构 $M^\bullet \to K^\bullet$，其中 $M^\bullet$ 上有界。
\item 若对所有 $|n| \gg 0$ 都有 $H^n(K^\bullet) = 0$，
则存在复形态射的交换图
$$
\xymatrix{
K^\bullet \ar[r] & L^\bullet \\
M^\bullet \ar[u] \ar[r] & N^\bullet \ar[u]
}
$$
其中所有箭头都是拟同构，$L^\bullet$ 下有界，
$M^\bullet$ 上有界，而 $N^\bullet$ 是有界复形。
\end{enumerate}
\end{lemma}

\begin{proof}
选取 $a \ll 0 \ll b$，并令
$M^\bullet = \tau_{\leq b}K^\bullet$、
$L^\bullet = \tau_{\geq a}K^\bullet$ 以及
$N^\bullet = \tau_{\leq b}L^\bullet = \tau_{\geq a}M^\bullet$。
关于截断函子，见同调，第
\ref{homology-section-truncations} 节。
\end{proof}

\noindent
为了陈述下一引理，以
$\text{Ac}^{+}(\mathcal{A})$、$\text{Ac}^{-}(\mathcal{A})$、
相应地 $\text{Ac}^b(\mathcal{A})$ 表示
$K^{+}(\mathcal{A})$、$K^{-}(\mathcal{A})$、相应地 $K^b(\mathcal{A})$
与 $\text{Ac}(\mathcal{A})$ 的交。以
$\text{Qis}^{+}(\mathcal{A})$、$\text{Qis}^{-}(\mathcal{A})$、
相应地 $\text{Qis}^b(\mathcal{A})$ 表示
$K^{+}(\mathcal{A})$、$K^{-}(\mathcal{A})$、相应地 $K^b(\mathcal{A})$
与 $\text{Qis}(\mathcal{A})$ 的交。

\begin{lemma}
\label{derived-lemma-bounded-derived}
设 $\mathcal{A}$ 为阿贝尔范畴。子范畴
$\text{Ac}^{+}(\mathcal{A})$、$\text{Ac}^{-}(\mathcal{A})$、
相应地 $\text{Ac}^b(\mathcal{A})$
分别是 $K^{+}(\mathcal{A})$、$K^{-}(\mathcal{A})$、
相应地 $K^b(\mathcal{A})$ 的严格满饱和三角子范畴。
相应的饱和乘法系（见引理 \ref{derived-lemma-operations}）
分别是集 $\text{Qis}^{+}(\mathcal{A})$、$\text{Qis}^{-}(\mathcal{A})$、
相应地 $\text{Qis}^b(\mathcal{A})$。
\begin{enumerate}
\item 函子 $K^{+}(\mathcal{A}) \to D^{+}(\mathcal{A})$ 的核是
$\text{Ac}^{+}(\mathcal{A})$，并诱导三角范畴等价
$$
K^{+}(\mathcal{A})/\text{Ac}^{+}(\mathcal{A}) =
\text{Qis}^{+}(\mathcal{A})^{-1}K^{+}(\mathcal{A})
\longrightarrow
D^{+}(\mathcal{A})
$$
\item 函子 $K^{-}(\mathcal{A}) \to D^{-}(\mathcal{A})$ 的核是
$\text{Ac}^{-}(\mathcal{A})$，并诱导三角范畴等价
$$
K^{-}(\mathcal{A})/\text{Ac}^{-}(\mathcal{A}) =
\text{Qis}^{-}(\mathcal{A})^{-1}K^{-}(\mathcal{A})
\longrightarrow
D^{-}(\mathcal{A})
$$
\item 函子 $K^b(\mathcal{A}) \to D^b(\mathcal{A})$ 的核是
$\text{Ac}^b(\mathcal{A})$，并诱导三角范畴等价
$$
K^b(\mathcal{A})/\text{Ac}^b(\mathcal{A}) =
\text{Qis}^b(\mathcal{A})^{-1}K^b(\mathcal{A})
\longrightarrow
D^b(\mathcal{A})
$$
\end{enumerate}
\end{lemma}

\begin{proof}
最初的各项陈述由引理 \ref{derived-lemma-acyclic-general}
应用于同调函子 $H^0$ 的限制得到。
(1)、(2)、(3) 中关于核的陈述是下述事实的结果：
每一种情形中的定义。由引理
\ref{derived-lemma-complex-cohomology-bounded}，每个函子都本质满。
为完成证明，还需说明这些函子全忠实。
我们先对下有界的版本证明这一点。

\medskip\noindent
假设 $K^\bullet, L^\bullet$ 是下有界复形。它们在
$D(\mathcal{A})$ 中的态射具有形式 $s^{-1}f$，其中的一对态射为
$f : K^\bullet \to (L')^\bullet$、$s : L^\bullet \to (L')^\bullet$，
且 $s$ 是拟同构。这意味着 $(L')^\bullet$ 的上同调下有界。
因而由引理 \ref{derived-lemma-complex-cohomology-bounded}，可选取拟同构
$s' : (L')^\bullet \to (L'')^\bullet$，使 $(L'')^\bullet$ 下有界。
那么这一对 $(s' \circ f, s' \circ s)$ 在
$\text{Qis}^{+}(\mathcal{A})^{-1}K^{+}(\mathcal{A})$ 中定义一个态射。
因此该函子是“满”的。最后，假设一对
$f : K^\bullet \to (L')^\bullet$、$s : L^\bullet \to (L')^\bullet$
在 $\text{Qis}^{+}(\mathcal{A})^{-1}K^{+}(\mathcal{A})$ 中定义一个态射，
而该态射在 $D(\mathcal{A})$ 中为零。这意味着存在拟同构
$s' : (L')^\bullet \to (L'')^\bullet$，使 $s' \circ f = 0$。
再次使用引理 \ref{derived-lemma-complex-cohomology-bounded}，得到拟同构
$s'' : (L'')^\bullet \to (L''')^\bullet$，其中 $(L''')^\bullet$ 下有界。
于是 $s'' \circ s' \circ f = 0$，这意味着 $s^{-1}f$ 在
$\text{Qis}^{+}(\mathcal{A})^{-1}K^{+}(\mathcal{A})$ 中为零。
这就证明了 (1) 中的函子是等价。

\medskip\noindent
(2) 的证明对偶于 (1) 的证明。为证明 (3)，可使用 (2) 的结果。
因此只需证明函子
$\text{Qis}^b(\mathcal{A})^{-1}K^b(\mathcal{A})
\to \text{Qis}^{-}(\mathcal{A})^{-1}K^{-}(\mathcal{A})$
全忠实。前一段的论证可直接用于证明此事，
只需全程使用已经上有界的复形。
\end{proof}







\section{典范 \zhdelta-函子}
\label{derived-section-canonical-delta-functor}

\noindent
导出范畴应是万有上同调函子的容纳对象。为陈述该结果，
我们使用从阿贝尔范畴到三角范畴的 $\delta$-函子概念，
见定义 \ref{derived-definition-delta-functor}。

\medskip\noindent
考虑函子 $\text{Comp}(\mathcal{A}) \to K(\mathcal{A})$。
一般而言，这个函子{\bf 不}是 $\delta$-函子。
看出这一点的最简单方法，是考虑 $\mathcal{A}$ 的对象的一个不分裂短正合序列
$0 \to A \to B \to C \to 0$。由于
$\Hom_{K(\mathcal{A})}(C[0], A[1]) = 0$，
我们看到，由这个短正合序列产生的任何特异三角必形如
$(A[0], B[0], C[0], a, b, 0)$。但 $K(\mathcal{A})$ 中存在这样的特异三角
意味着扩张分裂，矛盾。

\medskip\noindent
事实上，函子 $\text{Comp}(\mathcal{A}) \to D(\mathcal{A})$
是 $\delta$-函子。为看出这一点，必须定义与阿贝尔范畴
$\mathcal{A}$ 中复形的短正合序列
$$
0 \to A^\bullet \xrightarrow{a} B^\bullet \xrightarrow{b} C^\bullet \to 0
$$
相关联的态射 $\delta$。考虑态射 $a$ 的锥 $C(a)^\bullet$。
我们有 $C(a)^n = B^n \oplus A^{n + 1}$，并将
$q^n : C(a)^n \to C^n$ 定义为先投影到 $B^n$，再应用 $b^n$。
由此得到复形态射
$$
q : C(a)^\bullet \longrightarrow C^\bullet.
$$
显然 $q \circ i = b$，其中 $i$ 如定义 \ref{derived-definition-cone} 中所定义。
注意，由于 $a^\bullet$ 在每个次数上都是单射，$q$ 的核可识别为
$\text{id}_{A^\bullet}$ 的锥，而该锥无环。因而 $q$ 是拟同构。
根据引理 \ref{derived-lemma-the-same-up-to-isomorphisms}，三角
$$
(A, B, C(a), a, i, -p)
$$
是 $K(\mathcal{A})$ 中的特异三角。由于局部化函子
$K(\mathcal{A}) \to D(\mathcal{A})$ 正合，我们看到
$(A, B, C(a), a, i, -p)$ 是 $D(\mathcal{A})$ 中的特异三角。
由于 $q$ 是拟同构，$q$ 在 $D(\mathcal{A})$ 中是同构。
因而我们推得
$$
(A, B, C, a, b, -p \circ q^{-1})
$$
是 $D(\mathcal{A})$ 的特异三角。这引出下面的引理。

\begin{lemma}
\label{derived-lemma-derived-canonical-delta-functor}
设 $\mathcal{A}$ 为阿贝尔范畴。所定义的函子
$\text{Comp}(\mathcal{A}) \to D(\mathcal{A})$
具有一个自然的 $\delta$-函子结构，其中
$$
\delta_{A^\bullet \to B^\bullet \to C^\bullet} = - p \circ q^{-1}
$$
$p$ 和 $q$ 如上所述。同一构造将函子
$\text{Comp}^{+}(\mathcal{A}) \to D^{+}(\mathcal{A})$、
$\text{Comp}^{-}(\mathcal{A}) \to D^{-}(\mathcal{A})$ 和
$\text{Comp}^b(\mathcal{A}) \to D^b(\mathcal{A})$
都变为 $\delta$-函子。
\end{lemma}

\begin{proof}
我们已经看到，每当给定复形的短正合序列时，这个选择都产生特异三角。
还需证明，给定交换图
$$
\xymatrix{
0 \ar[r] &
A^\bullet \ar[r]_a \ar[d]_f &
B^\bullet \ar[r]_b \ar[d]_g &
C^\bullet \ar[r] \ar[d]_h &
0 \\
0 \ar[r] &
(A')^\bullet \ar[r]^{a'} &
(B')^\bullet \ar[r]^{b'} &
(C')^\bullet \ar[r] &
0
}
$$
后，我们得到定义 \ref{derived-definition-delta-functor} (2) 中所需的交换图。
由引理 \ref{derived-lemma-functorial-cone}，态射对 $(f, g)$ 诱导典范态射
$c : C(a)^\bullet \to C(a')^\bullet$。简单计算表明
$q' \circ c = h \circ q$ 且 $f[1] \circ p = p' \circ c$。
结论由此直接得到。
\end{proof}

\begin{lemma}
\label{derived-lemma-derived-compare-triangles-ses}
设 $\mathcal{A}$ 为阿贝尔范畴。设
$$
\xymatrix{
0 \ar[r] &
A^\bullet \ar[r] \ar[d] &
B^\bullet \ar[r] \ar[d] &
C^\bullet \ar[r] \ar[d] &
0 \\
0 \ar[r] &
D^\bullet \ar[r] &
E^\bullet \ar[r] &
F^\bullet \ar[r] &
0
}
$$
是一个复形态射的交换图，其各行是复形短正合序列，
而各竖直箭头都是拟同构。上述引理
\ref{derived-lemma-derived-canonical-delta-functor} 的 $\delta$-函子将短正合序列
$0 \to A^\bullet \to B^\bullet \to C^\bullet \to 0$ 和
$0 \to D^\bullet \to E^\bullet \to F^\bullet \to 0$
映到同构的特异三角。
\end{lemma}

\begin{proof}
这直接由下述事实得到：
$K(\mathcal{A}) \to D(\mathcal{A})$ 将拟同构变成同构，
且所相关联的特异三角具有函子性。
\end{proof}

\begin{lemma}
\label{derived-lemma-derived-compare-triangles-split-case}
设 $\mathcal{A}$ 为阿贝尔范畴。设
$$
\xymatrix{
0 \ar[r] &
A^\bullet \ar[r] &
B^\bullet \ar[r] &
C^\bullet \ar[r] &
0
}
$$
% P02-E0028: frozen source says "a short exact sequences".
为复形短正合序列。假设该短正合序列逐项分裂。设
$(A^\bullet, B^\bullet, C^\bullet, \alpha, \beta, \delta)$
为 $K(\mathcal{A})$ 中与该序列相关联的特异三角。上述引理
\ref{derived-lemma-derived-canonical-delta-functor} 的 $\delta$-函子
将短正合序列 $0 \to A^\bullet \to B^\bullet \to C^\bullet \to 0$
映到一个同构于特异三角
$$
(A^\bullet, B^\bullet, C^\bullet, \alpha, \beta, \delta).
$$
的三角。
\end{lemma}

\begin{proof}
由引理 \ref{derived-lemma-the-same-up-to-isomorphisms} 得到。
\end{proof}

\begin{remark}
\label{derived-remark-truncation-distinguished-triangle}
设 $\mathcal{A}$ 为阿贝尔范畴，$K^\bullet$ 为 $\mathcal{A}$ 的复形，
$a \in \mathbf{Z}$。我们声称，$D(\mathcal{A})$ 中存在典范特异三角
$$
\tau_{\leq a}K^\bullet \to K^\bullet \to \tau_{\geq a + 1}K^\bullet \to
(\tau_{\leq a}K^\bullet)[1]
$$
这里使用了同调，第 \ref{homology-section-truncations} 节中的典范截断函子 $\tau$。
事实上，我们先取与复形短正合序列
$$
0 \to \tau_{\leq a}K^\bullet \to K^\bullet \to
K^\bullet/\tau_{\leq a}K^\bullet \to 0
$$
由我们的 $\delta$-函子（引理
\ref{derived-lemma-derived-canonical-delta-functor}）相关联的特异三角。
然后，由同调，第 \ref{homology-section-truncations} 节中上同调群的描述，
我们使用映射 $K^\bullet \to \tau_{\geq a + 1}K^\bullet$
经由拟同构
$K^\bullet/\tau_{\leq a}K^\bullet \to \tau_{\geq a + 1}K^\bullet$ 分解这一事实。
以类似方式得到典范特异三角
$$
\tau_{\leq a}K^\bullet \to \tau_{\leq a + 1}K^\bullet \to
H^{a + 1}(K^\bullet)[-a-1] \to (\tau_{\leq a}K^\bullet)[1]
$$
和
$$
H^a(K^\bullet)[-a] \to \tau_{\geq a}K^\bullet \to \tau_{\geq a + 1}K^\bullet
\to H^a(K^\bullet)[-a + 1]
$$
\end{remark}

\begin{lemma}
\label{derived-lemma-trick-vanishing-composition}
设 $\mathcal{A}$ 为阿贝尔范畴。设
$$
K_0^\bullet \to K_1^\bullet \to \ldots \to K_n^\bullet
$$
为复形映射，且满足
\begin{enumerate}
\item 对 $i > 0$，$H^i(K_0^\bullet) = 0$；
\item $H^{-j}(K_j^\bullet) \to H^{-j}(K_{j + 1}^\bullet)$ 为零。
\end{enumerate}
那么在 $D(\mathcal{A})$ 中，复合 $K_0^\bullet \to K_n^\bullet$
经由 $\tau_{\leq -n}K_n^\bullet \to K_n^\bullet$ 分解。对偶地，
给定复形映射
$$
K_n^\bullet \to K_{n - 1}^\bullet \to \ldots \to K_0^\bullet
$$
且满足
\begin{enumerate}
\item 对 $i < 0$，$H^i(K_0^\bullet) = 0$；
\item $H^j(K_{j + 1}^\bullet) \to H^j(K_j^\bullet)$ 为零，
\end{enumerate}
则在 $D(\mathcal{A})$ 中，复合 $K_n^\bullet \to K_0^\bullet$
经由 $K_n^\bullet \to \tau_{\geq n}K_n^\bullet$ 分解。
\end{lemma}

\begin{proof}
考虑 $n = 1$ 的情形。由于在 $D(\mathcal{A})$ 中
$\tau_{\leq 0}K_0^\bullet = K_0^\bullet$，可以用
$\tau_{\leq 0}K_0^\bullet$ 替换 $K_0^\bullet$，用
$\tau_{\leq 0}K_1^\bullet$ 替换 $K_1^\bullet$。考虑特异三角
$$
\tau_{\leq -1}K_1^\bullet \to K_1^\bullet \to
H^0(K_1^\bullet)[0] \to (\tau_{\leq -1}K_1^\bullet)[1]
$$
（备注 \ref{derived-remark-truncation-distinguished-triangle}）。复合
$K_0^\bullet \to K_1^\bullet \to H^0(K_1^\bullet)[0]$ 为零，
因为它等于
$K_0^\bullet \to H^0(K_0^\bullet)[0] \to H^0(K_1^\bullet)[0]$，
而后者由假设为零。函子
$\Hom_{D(\mathcal{A})}(K_0^\bullet, -)$ 是同调函子（引理
\ref{derived-lemma-representable-homological}）这一事实允许我们找到所需的分解。
对 $n = 2$，由 $n = 1$ 的情形得到分解
$K_0^\bullet \to \tau_{\leq -1}K_1^\bullet$，然后将 $n = 1$ 的情形
应用于复形映射
$\tau_{\leq -1}K_1^\bullet \to \tau_{\leq -1}K_2^\bullet$，
便得到分解
$\tau_{\leq -1}K_1^\bullet \to \tau_{\leq -2}K_2^\bullet$。
一般情形完全以相同方式证明。
\end{proof}




\section{过滤导出范畴}
\label{derived-section-filtered-derived-category}

\noindent
本节的一个参考文献是 \cite[I, Chapter V]{cotangent}。设
$\mathcal{A}$ 为阿贝尔范畴。本节将定义 $\mathcal{A}$ 的过滤导出范畴
$DF(\mathcal{A})$。简言之，我们将它定义为 $\mathcal{A}$ 中带有有限过滤的对象所构成的
正合范畴的导出范畴。（因此，我们的构造是更一般的正合范畴之导出范畴
构造的特殊情形；例见 \cite{Buhler}、\cite{Keller}。）
Illusie 的过滤导出范畴是我们范畴中过滤有限的那些对象构成的满子范畴。
（在我们的范畴中，过滤在每个次数上仍有限，但可能不一致有界。）
我们这样选择的理由是：此构造并不更难，并且可将讨论应用于引理
\ref{derived-lemma-two-ss-complex-functor} 的谱序列，又见备注
\ref{derived-remark-functorial-ss}。

\medskip\noindent
我们将使用同调，第 \ref{homology-section-filtrations} 节中引入的过滤对象记号。
$\mathcal{A}$ 的过滤对象范畴记为 $\text{Fil}(\mathcal{A})$。
约定所有过滤都是递减的。

\begin{definition}
\label{derived-definition-finite-filtered}
设 $\mathcal{A}$ 为阿贝尔范畴。
$\mathcal{A}$ 的{\it 有限过滤对象范畴}
是 $\mathcal{A}$ 中过滤 $F$ 有限的过滤对象 $(A, F)$ 所构成的范畴。
我们将其记为 $\text{Fil}^f(\mathcal{A})$。
\end{definition}

\noindent
因而 $\text{Fil}^f(\mathcal{A})$ 是 $\text{Fil}(\mathcal{A})$ 的满子范畴。
对每个 $p \in \mathbf{Z}$，有函子
$\text{gr}^p : \text{Fil}^f(\mathcal{A}) \to \mathcal{A}$。
还有函子
$$
\text{gr} = \bigoplus\nolimits_{p \in \mathbf{Z}} \text{gr}^p :
\text{Fil}^f(\mathcal{A}) \to \text{Gr}(\mathcal{A})
$$
其中 $\text{Gr}(\mathcal{A})$ 是 $\mathcal{A}$ 的分次对象范畴，
见同调，定义 \ref{homology-definition-graded}。
最后，有函子
$$
(\text{遗忘 }F) : \text{Fil}^f(\mathcal{A}) \longrightarrow \mathcal{A}
$$
它将过滤对象 $(A, F)$ 映到其在 $\mathcal{A}$ 中的底对象。
范畴 $\text{Fil}^f(\mathcal{A})$ 是加性范畴，但一般不是阿贝尔范畴，
见同调，例 \ref{homology-example-not-abelian}。

\medskip\noindent
由于函子 $\text{gr}^p$、$\text{gr}$、$(\text{遗忘 }F)$ 都是加性函子，
由引理 \ref{derived-lemma-additive-exact-homotopy-category}，它们诱导三角范畴的正合函子
$$
\text{gr}^p, (\text{遗忘 }F) :
K(\text{Fil}^f(\mathcal{A}))
\to
K(\mathcal{A})
\quad\text{且}\quad
\text{gr} :
K(\text{Fil}^f(\mathcal{A}))
\to
K(\text{Gr}(\mathcal{A}))
$$
类比阿贝尔范畴的同伦范畴情形，我们作出下列定义。

\begin{definition}
\label{derived-definition-filtered-acyclic}
设 $\mathcal{A}$ 为阿贝尔范畴。
\begin{enumerate}
\item 设 $\alpha : K^\bullet \to L^\bullet$ 为
$K(\text{Fil}^f(\mathcal{A}))$ 的态射。若态射
$\text{gr}(\alpha)$ 是拟同构，则称 $\alpha$ 是{\it 过滤拟同构}。
\item 设 $K^\bullet$ 为 $K(\text{Fil}^f(\mathcal{A}))$ 的对象。
若复形 $\text{gr}(K^\bullet)$ 无环，则称 $K^\bullet$ {\it 过滤无环}。
\end{enumerate}
\end{definition}

\noindent
注意，$\alpha : K^\bullet \to L^\bullet$ 是过滤拟同构，当且仅当每个
$\text{gr}^p(\alpha)$ 都是拟同构。类似地，复形 $K^\bullet$ 过滤无环，
当且仅当每个 $\text{gr}^p(K^\bullet)$ 都无环。

\begin{lemma}
\label{derived-lemma-filtered-cohomology-homological}
设 $\mathcal{A}$ 为阿贝尔范畴。
\begin{enumerate}
\item 函子
$K(\text{Fil}^f(\mathcal{A})) \longrightarrow \text{Gr}(\mathcal{A})$、
$K^\bullet \longmapsto H^0(\text{gr}(K^\bullet))$
是同调函子。
\item 函子
$K(\text{Fil}^f(\mathcal{A})) \rightarrow \mathcal{A}$、
$K^\bullet \longmapsto H^0(\text{gr}^p(K^\bullet))$
是同调函子。
\item 函子
$K(\text{Fil}^f(\mathcal{A})) \longrightarrow \mathcal{A}$、
$K^\bullet \longmapsto H^0((\text{遗忘 }F)K^\bullet)$
是同调函子。
\end{enumerate}
\end{lemma}

\begin{proof}
这由下述事实得到：$H^0 : K(\mathcal{A}) \to \mathcal{A}$
是同调函子（见引理 \ref{derived-lemma-cohomology-homological}），而函子
$\text{gr}, \text{gr}^p, (\text{遗忘 }F)$ 是三角范畴的正合函子。
见引理 \ref{derived-lemma-exact-compose-homological-functor}。
\end{proof}

\begin{lemma}
\label{derived-lemma-filtered-acyclic}
设 $\mathcal{A}$ 为阿贝尔范畴。
$K(\text{Fil}^f(\mathcal{A}))$ 中由过滤无环复形组成的满子范畴
$\text{FAc}(\mathcal{A})$ 是 $K(\text{Fil}^f(\mathcal{A}))$
的严格满饱和三角子范畴。
$K(\text{Fil}^f(\mathcal{A}))$ 中相应的饱和乘法系（见引理
\ref{derived-lemma-operations}）是过滤拟同构集 $\text{FQis}(\mathcal{A})$。
特别地，局部化函子
$$
Q :
K(\text{Fil}^f(\mathcal{A}))
\longrightarrow
\text{FQis}(\mathcal{A})^{-1}K(\text{Fil}^f(\mathcal{A}))
$$
的核是 $\text{FAc}(\mathcal{A})$，且函子 $H^0 \circ \text{gr}$
经由 $Q$ 分解。
\end{lemma}

\begin{proof}
由引理 \ref{derived-lemma-filtered-cohomology-homological}，我们知道
$H^0 \circ \text{gr}$ 是同调函子。因而本引理是引理
\ref{derived-lemma-acyclic-general} 的特殊情形。
\end{proof}

\begin{definition}
\label{derived-definition-filtered-derived}
设 $\mathcal{A}$ 为阿贝尔范畴。设
$\text{FAc}(\mathcal{A})$ 和 $\text{FQis}(\mathcal{A})$ 如引理
\ref{derived-lemma-filtered-acyclic} 中所定义。
$\mathcal{A}$ 的{\it 过滤导出范畴}是三角范畴
$$
DF(\mathcal{A}) =
K(\text{Fil}^f(\mathcal{A}))/\text{FAc}(\mathcal{A}) =
\text{FQis}(\mathcal{A})^{-1} K(\text{Fil}^f(\mathcal{A})).
$$
\end{definition}

\begin{lemma}
\label{derived-lemma-filtered-derived-functors}
函子 $\text{gr}^p, \text{gr}, (\text{遗忘 }F)$ 诱导典范正合函子
$$
\text{gr}^p, (\text{遗忘 }F):
DF(\mathcal{A})
\longrightarrow
D(\mathcal{A})
$$
和
$$
\text{gr}:
DF(\mathcal{A})
\longrightarrow
D(\text{Gr}(\mathcal{A}))
$$
它们与局部化函子交换。
\end{lemma}

\begin{proof}
只要能证明每个函子
$\text{gr}^p, \text{gr}, (\text{遗忘 }F)$
都将过滤拟同构变为拟同构，这就由局部化的万有性质（见引理
\ref{derived-lemma-universal-property-localization}）得到。对前两个函子，这依定义成立。
对最后一个，需要做一点工作。设
$f : K^\bullet \to L^\bullet$ 是 $K(\text{Fil}^f(\mathcal{A}))$ 中的过滤拟同构。
选取包含 $f$ 的特异三角
$(K^\bullet, L^\bullet, M^\bullet, f, g, h)$。那么由引理
\ref{derived-lemma-filtered-acyclic}，$M^\bullet$ 过滤无环。
因此，由 $K(\mathcal{A})$ 的相应引理，只需证明过滤无环复形
在遗忘过滤后是无环复形。这由同调，引理
\ref{homology-lemma-filtered-acyclic} 得到。
\end{proof}

\begin{definition}
\label{derived-definition-filtered-derived-bounded}
设 $\mathcal{A}$ 为阿贝尔范畴。{\it 有界过滤导出范畴}
$DF^b(\mathcal{A})$ 是 $DF(\mathcal{A})$ 的满子范畴，其对象是满足
$\text{gr}(X) \in D^b(\mathcal{A})$ 的那些 $X$。
对下有界过滤导出范畴 $DF^{+}(\mathcal{A})$
和上有界过滤导出范畴 $DF^{-}(\mathcal{A})$ 作类似定义。
\end{definition}

\begin{lemma}
\label{derived-lemma-filtered-complex-cohomology-bounded}
设 $\mathcal{A}$ 为阿贝尔范畴，并设
$K^\bullet \in K(\text{Fil}^f(\mathcal{A}))$。
\begin{enumerate}
\item 若对所有 $n < a$ 都有 $H^n(\text{gr}(K^\bullet)) = 0$，
则存在过滤拟同构 $K^\bullet \to L^\bullet$，使得对所有 $n < a$
都有 $L^n = 0$。
\item 若对所有 $n > b$ 都有 $H^n(\text{gr}(K^\bullet)) = 0$，
则存在过滤拟同构 $M^\bullet \to K^\bullet$，使得对所有 $n > b$
都有 $M^n = 0$。
\item 若对所有 $|n| \gg 0$ 都有 $H^n(\text{gr}(K^\bullet)) = 0$，
则存在复形态射的交换图
$$
\xymatrix{
K^\bullet \ar[r] & L^\bullet \\
M^\bullet \ar[u] \ar[r] & N^\bullet \ar[u]
}
$$
其中所有箭头都是过滤拟同构，$L^\bullet$ 下有界，
$M^\bullet$ 上有界，而 $N^\bullet$ 是有界复形。
\end{enumerate}
\end{lemma}

\begin{proof}
假设对所有 $n < a$ 都有 $H^n(\text{gr}(K^\bullet)) = 0$。由同调，引理
\ref{homology-lemma-filtered-acyclic}，序列
$$
K^{a - 2} \xrightarrow{d^{a - 2}} K^{a - 1} \xrightarrow{d^{a - 1}} K^a
$$
是 $\mathcal{A}$ 的对象的正合序列，且态射
$d^{a - 2}$ 和 $d^{a - 1}$ 都是严格的。因而在
$\text{Fil}^f(\mathcal{A})$ 中有
$\Coim(d^{a - 1}) = \Im(d^{a - 1})$，而映射
$\text{gr}(\Im(d^{a - 1})) \to \text{gr}(K^a)$
是单射，其像等于映射
$\text{gr}(K^{a - 1}) \to \text{gr}(K^a)$ 的像，见同调，引理
\ref{homology-lemma-characterize-strict}。这意味着到截断
$$
\tau_{\geq a}K^\bullet =
(\ldots \to 0 \to K^a/\Im(d^{a - 1}) \to K^{a + 1} \to \ldots)
$$
的映射 $K^\bullet \to \tau_{\geq a}K^\bullet$ 是过滤拟同构。
这证明了 (1)。(2) 的证明对偶于 (1) 的证明。
(3) 形式地由 (1) 和 (2) 得到。
\end{proof}

\noindent
为陈述下一引理，以
$\text{FAc}^{+}(\mathcal{A})$、$\text{FAc}^{-}(\mathcal{A})$、
相应地 $\text{FAc}^b(\mathcal{A})$ 表示
$K^{+}(\text{Fil}^f\mathcal{A})$、$K^{-}(\text{Fil}^f\mathcal{A})$、
相应地 $K^b(\text{Fil}^f\mathcal{A})$ 与 $\text{FAc}(\mathcal{A})$ 的交。
以 $\text{FQis}^{+}(\mathcal{A})$、$\text{FQis}^{-}(\mathcal{A})$、
相应地 $\text{FQis}^b(\mathcal{A})$ 表示
$K^{+}(\text{Fil}^f\mathcal{A})$、$K^{-}(\text{Fil}^f\mathcal{A})$、
相应地 $K^b(\text{Fil}^f\mathcal{A})$ 与 $\text{FQis}(\mathcal{A})$ 的交。

\begin{lemma}
\label{derived-lemma-filtered-bounded-derived}
设 $\mathcal{A}$ 为阿贝尔范畴。子范畴
$\text{FAc}^{+}(\mathcal{A})$、$\text{FAc}^{-}(\mathcal{A})$、
相应地 $\text{FAc}^b(\mathcal{A})$
分别是 $K^{+}(\text{Fil}^f\mathcal{A})$、
$K^{-}(\text{Fil}^f\mathcal{A})$、相应地
$K^b(\text{Fil}^f\mathcal{A})$ 的严格满饱和三角子范畴。
相应的饱和乘法系（见引理 \ref{derived-lemma-operations}）
分别是集 $\text{FQis}^{+}(\mathcal{A})$、
$\text{FQis}^{-}(\mathcal{A})$、相应地 $\text{FQis}^b(\mathcal{A})$。
\begin{enumerate}
\item 函子
$K^{+}(\text{Fil}^f\mathcal{A}) \to DF^{+}(\mathcal{A})$
的核是 $\text{FAc}^{+}(\mathcal{A})$，并诱导三角范畴等价
$$
K^{+}(\text{Fil}^f\mathcal{A})/\text{FAc}^{+}(\mathcal{A}) =
\text{FQis}^{+}(\mathcal{A})^{-1}K^{+}(\text{Fil}^f\mathcal{A})
\longrightarrow
DF^{+}(\mathcal{A})
$$
\item 函子
$K^{-}(\text{Fil}^f\mathcal{A}) \to DF^{-}(\mathcal{A})$
的核是 $\text{FAc}^{-}(\mathcal{A})$，并诱导三角范畴等价
$$
K^{-}(\text{Fil}^f\mathcal{A})/\text{FAc}^{-}(\mathcal{A}) =
\text{FQis}^{-}(\mathcal{A})^{-1}K^{-}(\text{Fil}^f\mathcal{A})
\longrightarrow
DF^{-}(\mathcal{A})
$$
\item 函子
$K^b(\text{Fil}^f\mathcal{A}) \to DF^b(\mathcal{A})$
的核是 $\text{FAc}^b(\mathcal{A})$，并诱导三角范畴等价
$$
K^b(\text{Fil}^f\mathcal{A})/\text{FAc}^b(\mathcal{A}) =
\text{FQis}^b(\mathcal{A})^{-1}K^b(\text{Fil}^f\mathcal{A})
\longrightarrow
DF^b(\mathcal{A})
$$
\end{enumerate}
\end{lemma}

\begin{proof}
这由上述结果，特别是引理
\ref{derived-lemma-filtered-complex-cohomology-bounded}，通过与引理
\ref{derived-lemma-bounded-derived} 的证明中完全相同的论证得到。
\end{proof}











\section{一般导出函子}
\label{derived-section-derived-functors}

\noindent
本节的一个参考文献是 \cite{SGA4} 中 Deligne 的 expos\'e XVII。
存在一种非常一般的右导出函子和左导出函子概念：
当我们有三角范畴之间的正合函子、源范畴中的乘法系，
并希望找到该正合函子到局部化范畴的“正确”延拓时，便可使用它。

\begin{situation}
\label{derived-situation-derived-functor}
此处 $F : \mathcal{D} \to \mathcal{D}'$ 是三角范畴之间的正合函子，
$S$ 是 $\mathcal{D}$ 中的饱和乘法系，且与
$\mathcal{D}$ 上的三角范畴结构相容。
\end{situation}

\noindent
设 $X \in \Ob(\mathcal{D})$。回忆范畴，备注
\ref{categories-remark-left-localization-morphisms-colimit} 中源为 $X$ 的
$S$ 中箭头 $s : X \to X'$ 所组成的过滤范畴 $X/S$。
对偶地，在范畴，备注
\ref{categories-remark-right-localization-morphisms-colimit} 中，我们定义了目标为 $X$ 的
$S$ 中箭头 $s : X' \to X$ 所组成的余过滤范畴 $S/X$。

\begin{definition}
\label{derived-definition-right-derived-functor-defined}
假设和记号如情形 \ref{derived-situation-derived-functor} 中所述。设
$X \in \Ob(\mathcal{D})$。
\begin{enumerate}
\item 若 ind-对象
$$
(X/S) \longrightarrow \mathcal{D}', \quad
(s : X \to X') \longmapsto F(X')
$$
本质常值\footnote{关于范畴的 ind-对象或 pro-对象何时本质常值的讨论，
参见范畴，第 \ref{categories-section-essentially-constant} 节。}，
则称{\it 右导出函子 $RF$ 在} $X$ {\it 处有定义}；在此情形下，其值
$\mathcal{D}'$ 中的 $Y$ 称为{\it $RF$ 在 $X$ 处的值}。
\item 若 pro-对象
$$
(S/X) \longrightarrow \mathcal{D}', \quad
(s: X' \to X) \longmapsto F(X')
$$
本质常值，则称{\it 左导出函子 $LF$ 在} $X$
{\it 处有定义}；在此情形下，其值 $\mathcal{D}'$ 中的 $Y$
称为{\it $LF$ 在 $X$ 处的值}。
\end{enumerate}
我们常滥用记号，将这些值简记为 $RF(X)$ 或 $LF(X)$。
\end{definition}

\noindent
事实上，$\mathcal{D}$ 中由 $RF$ 有定义的对象组成的满子范畴
是三角子范畴，而 $RF$ 将在此子范畴上定义一个函子，
它将 $S$ 中的态射变为同构。

\begin{lemma}
\label{derived-lemma-derived-functor}
假设和记号如情形 \ref{derived-situation-derived-functor} 中所述。设
$f : X \to Y$ 为 $\mathcal{D}$ 的态射。
\begin{enumerate}
\item 若 $RF$ 在 $X$ 和 $Y$ 处都有定义，则值之间存在唯一态射
$RF(f) : RF(X) \to RF(Y)$，使对任意交换图
$$
\xymatrix{
X \ar[d]_f \ar[r]_s & X' \ar[d]^{f'} \\
Y \ar[r]^{s'} & Y'
}
$$
若 $s, s' \in S$，则图
$$
\xymatrix{
F(X) \ar[d] \ar[r] & F(X') \ar[d] \ar[r] & RF(X) \ar[d] \\
F(Y) \ar[r] & F(Y') \ar[r] & RF(Y)
}
$$
交换。
\item 若 $LF$ 在 $X$ 和 $Y$ 处都有定义，则值之间存在唯一态射
$LF(f) : LF(X) \to LF(Y)$，使对任意交换图
$$
\xymatrix{
X' \ar[d]_{f'} \ar[r]_s & X \ar[d]^f \\
Y' \ar[r]^{s'} & Y
}
$$
若 $s, s'$ 属于 $S$，则图
$$
\xymatrix{
LF(X) \ar[d] \ar[r] & F(X') \ar[d] \ar[r] & F(X) \ar[d] \\
LF(Y) \ar[r] & F(Y') \ar[r] & F(Y)
}
$$
交换。
\end{enumerate}
\end{lemma}

\begin{proof}
只要假设余极限
$$
RF(X) = \colim_{s : X \to X'} F(X')
\quad\text{且}\quad
RF(Y) = \colim_{s' : Y \to Y'} F(Y')
$$
存在，(1) 就成立。事实上，给出余极限之间的态射
$RF(X) \to RF(Y)$，等价于对 $\Ob(X/S)$ 中的每个
$s : X \to X'$ 给出一个态射 $F(X') \to RF(Y)$，
并要求它们与范畴 $X/S$ 中的态射相容。
为得到该态射，按 MS2 选取交换图
$$
\xymatrix{
X \ar[d]_f \ar[r]_s & X' \ar[d]^{f'} \\
Y \ar[r]^{s'} & Y'
}
$$
其中 $s, s'$ 属于 $S$，并将 $F(X') \to RF(Y)$ 定义为复合
$F(X') \to F(Y') \to RF(Y)$。使用 MS3 可看出，这与上述图的选择无关。
细节从略。(2) 的证明与之对偶。
\end{proof}

\begin{lemma}
\label{derived-lemma-derived-inverts}
假设和记号如情形 \ref{derived-situation-derived-functor} 中所述。设
$s : X \to Y$ 是 $S$ 的元素。
\begin{enumerate}
\item $RF$ 在 $X$ 处有定义，当且仅当它在 $Y$ 处有定义。
在此情形下，值之间的映射 $RF(s) : RF(X) \to RF(Y)$ 是同构。
\item $LF$ 在 $X$ 处有定义，当且仅当它在 $Y$ 处有定义。
在此情形下，值之间的映射 $LF(s) : LF(X) \to LF(Y)$ 是同构。
\end{enumerate}
\end{lemma}

\begin{proof}
从略。
\end{proof}

\begin{lemma}
\label{derived-lemma-derived-shift}
\begin{slogan}
导出函子与平移相容
\end{slogan}
假设和记号如情形 \ref{derived-situation-derived-functor} 中所述。
设 $X$ 为 $\mathcal{D}$ 的对象，$n \in \mathbf{Z}$。
\begin{enumerate}
\item $RF$ 在 $X$ 处有定义，当且仅当它在 $X[n]$ 处有定义。
在此情形下，值之间存在典范同构
$RF(X)[n]= RF(X[n])$。
\item $LF$ 在 $X$ 处有定义，当且仅当它在 $X[n]$ 处有定义。
在此情形下，值之间存在典范同构
$LF(X)[n] \to LF(X[n])$。
\end{enumerate}
\end{lemma}

\begin{proof}
从略。
\end{proof}

\begin{lemma}
\label{derived-lemma-2-out-of-3-defined}
假设和记号如情形 \ref{derived-situation-derived-functor} 中所述。
设 $(X, Y, Z, f, g, h)$ 为 $\mathcal{D}$ 的特异三角。
若 $RF$ 在 $X, Y, Z$ 中的任意两个处有定义，则它在第三个处也有定义。
此外，在这种情形下
$$
(RF(X), RF(Y), RF(Z), RF(f), RF(g), RF(h))
$$
是 $\mathcal{D}'$ 中的特异三角。对 $LF$ 也有类似结论。
\end{lemma}

\begin{proof}
假设 $RF$ 在 $X, Y$ 处有定义，值分别为 $A, B$。
设 $RF(f) : A \to B$ 为诱导态射，见引理
\ref{derived-lemma-derived-functor}。可在 $\mathcal{D}'$ 中选取特异三角
$(A, B, C, RF(f), b, c)$。我们声称 $C$ 是 $RF$ 在 $Z$ 处的一个值。

\medskip\noindent
为证明这一点，选取 $S$ 中的 $s : X \to X'$，使存在如范畴，定义
\ref{categories-definition-essentially-constant-diagram} 中的态射
$\alpha : A \to F(X')$。根据 MS2，可选取交换图
$$
\xymatrix{
X \ar[d]_f \ar[r]_s & X' \ar[d]^{f'} \\
Y \ar[r]^{s'} & Y'
}
$$
其中 $s' \in S$。利用 $Y/S$ 是过滤的，我们可以（把 $s'$ 替换为
$S$ 中某个 $s'' : Y \to Y''$ 之后）假设存在如范畴，定义
\ref{categories-definition-essentially-constant-diagram} 中的态射
$\beta : B \to F(Y')$。考虑图
$$
\xymatrix{
A \ar[d]_{RF(f)} \ar[r]_-\alpha &
F(X') \ar[r] \ar[d]^{F(f')} &
A \ar[d]^{RF(f)} \\
B \ar[r]^-\beta & F(Y') \ar[r] & B
}
$$
左方块未必交换，但外部方块和右方块交换。
% P02-E0029: frozen source indexes the ind-object by s':Y'->Y, opposite to its own definition of Y/S.
假设 ind-对象 $\{F(Y')\}_{s' : Y' \to Y}$ 本质常值，这意味着存在
$S$ 中的 $s'' : Y \to Y''$ 和态射 $h : Y' \to Y''$，
使 $s'' = h \circ s'$，并且 $F(h)$ 等于复合
$F(Y') \to B \to F(Y') \to F(Y'')$。因此，将 $Y'$ 替换为 $Y''$，
将 $\beta$ 替换为 $F(h) \circ \beta$ 之后，上图将交换
（这可对箭头作直接计算）。

\medskip\noindent
使用 MS6 选取三角态射
$$
(s, s', s'') : (X, Y, Z, f, g, h) \longrightarrow (X', Y', Z', f', g', h')
$$
其中 $s'' \in S$。由 TR3，选取三角态射
$$
(\alpha, \beta, \gamma) :
(A, B, C, RF(f), b, c)
\longrightarrow
(F(X'), F(Y'), F(Z'), F(f'), F(g'), F(h'))
$$

\medskip\noindent
由引理 \ref{derived-lemma-derived-inverts}，只需证明 $RF(Z')$ 有定义，
并且箭头 $\gamma : C \to F(Z')$ 诱导同构 $C \to RF(Z')$。
事实上，此时我们将得到三角同构
$$
(A, B, C, RF(f), b, c)
\longrightarrow
(RF(X'), RF(Y'), RF(Z'), RF(f'), RF(g'), RF(h'))
$$
并由 TR1 得出其目标是特异三角。考虑引理
\ref{derived-lemma-limit-triangles} 中由三角组成的范畴 $\mathcal{I}$：
$$
\mathcal{I} =
\{(t, t', t'') : (X', Y', Z', f', g', h') \to (X'', Y'', Z'', f'', g'', h'')
\mid (t, t', t'') \in S\}
$$
要说明系 $F(Z'')$ 在范畴 $Z'/S$ 上本质常值，
等价于说明 $F(Z'')$ 的系在 $\mathcal{I}$ 上本质常值，
因为 $\mathcal{I} \to Z'/S$ 是余终的；见范畴，引理
\ref{categories-lemma-cofinal-essentially-constant}（余终性在引理
\ref{derived-lemma-limit-triangles} 中证明）。对 $\mathcal{D}'$ 的任意对象 $W$，考虑图
$$
\xymatrix{
\colim_\mathcal{I} \Mor_{\mathcal{D}'}(W, F(X'')) \ar[d] &
\Mor_{\mathcal{D}'}(W, A) \ar[l] \ar[d] \\
\colim_\mathcal{I} \Mor_{\mathcal{D}'}(W, F(Y'')) \ar[d] &
\Mor_{\mathcal{D}'}(W, B) \ar[d] \ar[l] \\
\colim_\mathcal{I} \Mor_{\mathcal{D}'}(W, F(Z'')) \ar[d] &
\Mor_{\mathcal{D}'}(W, C) \ar[d] \ar[l] \\
\colim_\mathcal{I} \Mor_{\mathcal{D}'}(W, F(X''[1])) \ar[d] &
\Mor_{\mathcal{D}'}(W, A[1]) \ar[d] \ar[l] \\
\colim_\mathcal{I} \Mor_{\mathcal{D}'}(W, F(Y''[1])) &
\Mor_{\mathcal{D}'}(W, B[1]) \ar[l]
}
$$
其水平箭头由与 $(\alpha, \beta, \gamma)$ 复合给出。
由于过滤余极限正合（代数，引理
\ref{algebra-lemma-directed-colimit-exact}），左列是正合序列。
因而 $5$ 引理（同调，引理 \ref{homology-lemma-five-lemma}）表明映射
$$
\colim_\mathcal{I} \Mor_{\mathcal{D}'}(W, F(Z''))
\longrightarrow
\Mor_{\mathcal{D}'}(W, C)
$$
是双射。由范畴，引理
\ref{categories-lemma-characterize-essentially-constant-ind} 的第 (4) 部分，
我们得出 $F(Z'')$ 在 $\mathcal{I}$ 上以 $C$ 为值本质常值。
\end{proof}

\begin{lemma}
\label{derived-lemma-direct-sum-defined}
假设和记号如情形 \ref{derived-situation-derived-functor} 中所述。
设 $X, Y$ 为 $\mathcal{D}$ 的对象。
\begin{enumerate}
\item 若 $RF$ 在 $X$ 和 $Y$ 处有定义，则 $RF$ 在 $X \oplus Y$ 处有定义。
\item 若 $\mathcal{D}'$ 是 Karoubi 范畴，且 $RF$ 在 $X \oplus Y$ 处有定义，
则 $RF$ 在 $X$ 和 $Y$ 处都有定义。
\end{enumerate}
两种情形下都有 $RF(X \oplus Y) = RF(X) \oplus RF(Y)$。
对 $LF$ 也有类似结论。
\end{lemma}

\begin{proof}
若 $RF$ 在 $X$ 和 $Y$ 处有定义，则由特异三角
$X \to X \oplus Y \to Y \to X[1]$（引理 \ref{derived-lemma-split}）以及引理
\ref{derived-lemma-2-out-of-3-defined}，$RF$ 在 $X \oplus Y$ 处有定义，且有特异三角
$RF(X) \to RF(X \oplus Y) \to RF(Y) \to RF(X)[1]$。
再对其应用引理 \ref{derived-lemma-split}，得到
$RF(X \oplus Y) = RF(X) \oplus RF(Y)$。这证明了 (1) 和最后的断言。

\medskip\noindent
反过来，假设 $RF$ 在 $X \oplus Y$ 处有定义，且 $\mathcal{D}'$
是 Karoubi 范畴。由于 $S$ 是饱和系，$S$ 是在加性局部化函子
$Q : \mathcal{D} \to S^{-1}\mathcal{D}$ 下变为可逆的箭头之集，
见范畴，引理 \ref{categories-lemma-what-gets-inverted}。
因此，对 $S$ 中任意 $s : X \to X'$ 和 $s' : Y \to Y'$，态射
$s \oplus s' : X \oplus Y \to X' \oplus Y'$ 都是 $S$ 的元素。
由此得到函子
$$
X/S \times Y/S \longrightarrow (X \oplus Y)/S
$$
回忆，范畴 $X/S, Y/S, (X \oplus Y)/S$ 都是过滤范畴（范畴，备注
\ref{categories-remark-left-localization-morphisms-colimit}）。由范畴，引理
\ref{categories-lemma-essentially-constant-over-product}，$X/S \times Y/S$ 是过滤的，而
$F|_{X/S} : X/S \to \mathcal{D}'$
% P02-E0030: frozen source changes the only functor F to an undefined G on Y/S.
（相应地 $G|_{Y/S} : Y/S \to \mathcal{D}'$）本质常值，当且仅当
$F|_{X/S} \circ \text{pr}_1 : X/S \times Y/S \to \mathcal{D}'$
（相应地 $G|_{Y/S} \circ \text{pr}_2 : X/S \times Y/S \to \mathcal{D}'$）
本质常值。下面我们将证明上述函子余终，因此由范畴，引理
\ref{categories-lemma-cofinal-essentially-constant}，我们看到，
$F|_{(X \oplus Y)/S}$ 本质常值意味着
$F|_{X/S} \circ \text{pr}_1 \oplus F|_{Y/S} \circ \text{pr}_2 :
X/S \times Y/S \to \mathcal{D}'$
本质常值。由同调，引理
\ref{homology-lemma-direct-sum-from-product-essentially-constant}
（此处用到 $\mathcal{D}'$ 是 Karoubi 范畴），我们看到，
$F|_{X/S} \circ \text{pr}_1 \oplus F|_{Y/S} \circ \text{pr}_2$
本质常值意味着
$F|_{X/S} \circ \text{pr}_1$ 和 $F|_{Y/S} \circ \text{pr}_2$
都本质常值，这证明了 $RF$ 在 $X$ 和 $Y$ 处有定义。

\medskip\noindent
证明上述函子余终。为此，选取 $S$ 中任意的
$t : X \oplus Y \to Z$。使用 MS2，可找到态射
$Z \to X'$ 和 $Z \to Y'$，以及 $S$ 中的
$s : X \to X'$ 和 $s' : Y \to Y'$，使得
$$
\xymatrix{
X \ar[d]^s & X \oplus Y \ar[d] \ar[l] \ar[r] & Y \ar[d]_{s'} \\
X' & Z \ar[l] \ar[r] & Y'
}
$$
交换。这证明在 $(X \oplus Y)/S$ 中存在映射
$Z \to X' \oplus Y'$，也就得到范畴，定义
\ref{categories-definition-cofinal} 的第 (1) 部分。为证明第 (2) 部分，
只需证明：给定 $t : X \oplus Y \to Z$ 以及
$(X \oplus Y)/S$ 中的态射
$s_i \oplus s'_i : Z \to X'_i \oplus Y'_i$、$i = 1, 2$，
我们能找到 $S$ 中的态射
$a : X'_1 \to X'$、$b : X'_2 \to X'$、
$c : Y'_1 \to Y'$、$d : Y'_2 \to Y'$，使得
$a \circ s_1 = b \circ s_2$ 且 $c \circ s'_1 = d \circ s'_2$。
为此，先选取任意 $X'$ 和 $Y'$ 以及 $S$ 中的映射 $a, b, c, d$；
这是可能的，因为 $X/S$ 和 $Y/S$ 是过滤范畴。然后两个映射
$a \circ s_1, b \circ s_2 : Z \to X'$ 在 $S^{-1}\mathcal{D}$ 中变得相等。
因而可找到 $S$ 中的态射 $X' \to X''$ 使它们相等。
类似地，可找到 $S$ 中的 $Y' \to Y''$ 使
$c \circ s'_1$ 和 $d \circ s'_2$ 相等。用 $X''$ 替换 $X'$、
用 $Y''$ 替换 $Y'$，即得
$a \circ s_1 = b \circ s_2$ 且 $c \circ s'_1 = d \circ s'_2$。

\medskip\noindent
% P02-E0031: frozen source has subject-verb disagreement in "The proof ... are dual".
对 $LF$ 的相应陈述之证明是对偶的。
\end{proof}

\begin{proposition}
\label{derived-proposition-derived-functor}
假设和记号如情形 \ref{derived-situation-derived-functor} 中所述。
\begin{enumerate}
\item $\mathcal{D}$ 中由 $RF$ 有定义的对象组成的满子范畴
$\mathcal{E}$ 是 $\mathcal{D}$ 的严格满三角子范畴。
\item 我们得到三角范畴的正合函子
$RF : \mathcal{E} \longrightarrow \mathcal{D}'$。
\item $S$ 中源或目标属于 $\mathcal{E}$ 的元素都是 $\mathcal{E}$ 的态射。
\item $S_\mathcal{E} = \text{Arrows}(\mathcal{E}) \cap S$ 的任意元素都被
$RF$ 映到同构。
\item 集 $S_\mathcal{E}$ 是 $\mathcal{E}$ 中与三角结构相容的饱和乘法系。
\item 函子 $S_\mathcal{E}^{-1}\mathcal{E} \to S^{-1}\mathcal{D}$
是三角范畴的全忠实正合函子。
\item 我们得到正合函子
$$
RF : S_\mathcal{E}^{-1}\mathcal{E} \longrightarrow \mathcal{D}'.
$$
\item 若 $\mathcal{D}'$ 是 Karoubi 范畴，则 $\mathcal{E}$ 是 $\mathcal{D}$
的饱和三角子范畴。
\end{enumerate}
对 $LF$ 有类似结果。
\end{proposition}

\begin{proof}
由于 $S$ 饱和，它包含所有同构（见范畴，定义
\ref{categories-definition-saturated-multiplicative-system} 之后的备注）。
因而 (1) 由引理 \ref{derived-lemma-derived-inverts}、
\ref{derived-lemma-2-out-of-3-defined} 和 \ref{derived-lemma-derived-shift} 得到。
(2) 由引理 \ref{derived-lemma-derived-functor}、\ref{derived-lemma-derived-shift}
和 \ref{derived-lemma-2-out-of-3-defined} 得到。
(3) 由引理 \ref{derived-lemma-derived-inverts} 得到。
(4) 由引理 \ref{derived-lemma-derived-inverts} 得到。
(5) 由定义和 (3) 得到。(6) 中的全忠实性由 (3) 和定义得到。
$S_\mathcal{E}^{-1}\mathcal{E} \to S^{-1}\mathcal{D}$ 正合，是因为
$S_\mathcal{E}^{-1}\mathcal{E}$ 中的三角是特异的，当且仅当它同构于
$\mathcal{E}$ 中某个特异三角的像；见命题
\ref{derived-proposition-construct-localization} 的证明。
$RF : \mathcal{E} \to \mathcal{D}'$ 经由正合函子
$S_\mathcal{E}^{-1}\mathcal{E} \to \mathcal{D}'$ 分解，这由引理
\ref{derived-lemma-universal-property-localization} 得到。
最后，(8) 由引理 \ref{derived-lemma-direct-sum-defined} 得到。
\end{proof}

\noindent
命题 \ref{derived-proposition-derived-functor} 告诉我们，$RF$ 定义在
$S^{-1}\mathcal{D}$ 的一个极大严格满三角子范畴上，
并且在这个三角范畴上是正合函子。图示为：
$$
\xymatrix{
\mathcal{D} \ar[d]_Q \ar[rrr]_F & & & \mathcal{D}' \\
S^{-1}\mathcal{D} & &
S_\mathcal{E}^{-1}\mathcal{E}
\ar[ll]_{\text{全忠实}}^{\text{正合}} \ar[ur]_{RF}
}
$$

\begin{definition}
\label{derived-definition-everywhere-defined}
在情形 \ref{derived-situation-derived-functor} 中，
若 $RF$ 在 $\mathcal{D}$ 的每个对象处都有定义，
则称 $F$ {\it 右可导}，或称 $RF$ {\it 处处有定义}。
若 $LF$ 在 $\mathcal{D}$ 的每个对象处都有定义，
则称 $F$ {\it 左可导}，或称 $LF$ {\it 处处有定义}。
\end{definition}

\noindent
在这种情形下，我们得到右（相应地左）导出函子
\begin{equation}
\label{derived-equation-everywhere}
RF : S^{-1}\mathcal{D} \longrightarrow \mathcal{D}',
\quad\text{（分别 }
LF : S^{-1}\mathcal{D} \longrightarrow \mathcal{D}'),
\end{equation}
见命题 \ref{derived-proposition-derived-functor}。
在大多数有趣的情形中，$RF \circ Q$ 并不等于 $F$。
事实上，典范映射 $F(X) \to RF(X)$ 甚至可能从不是同构。
实际中这不会发生，因为我们实际上只会通过下述方法证明 $F$
右可导：证明可在三角范畴 $\mathcal{D}$ 中适当选择的对象上取 $F$ 的值，
以计算 $RF$。这促使我们作出一个定义。

\begin{definition}
\label{derived-definition-computes}
在情形 \ref{derived-situation-derived-functor} 中：
\begin{enumerate}
\item 若 $RF$ 在 $\mathcal{D}$ 的对象 $X$ 处有定义，且典范映射
$F(X) \to RF(X)$ 是同构，则称 $X$ {\it 计算} $RF$。
\item 若 $LF$ 在 $\mathcal{D}$ 的对象 $X$ 处有定义，且典范映射
$LF(X) \to F(X)$ 是同构，则称 $X$ {\it 计算} $LF$。
\end{enumerate}
\end{definition}

\begin{lemma}
\label{derived-lemma-computes-shift}
假设和记号如情形 \ref{derived-situation-derived-functor} 中所述。
设 $X$ 为 $\mathcal{D}$ 的对象，$n \in \mathbf{Z}$。
\begin{enumerate}
\item $X$ 计算 $RF$，当且仅当 $X[n]$ 计算 $RF$。
\item $X$ 计算 $LF$，当且仅当 $X[n]$ 计算 $LF$。
\end{enumerate}
\end{lemma}

\begin{proof}
从略。
\end{proof}

\begin{lemma}
\label{derived-lemma-2-out-of-3-computes}
假设和记号如情形 \ref{derived-situation-derived-functor} 中所述。
设 $(X, Y, Z, f, g, h)$ 为 $\mathcal{D}$ 的特异三角。
若 $X, Y$ 计算 $RF$，则 $Z$ 也如此。对 $LF$ 有类似结论。
\end{lemma}

\begin{proof}
由引理 \ref{derived-lemma-2-out-of-3-defined}，我们知道 $RF$ 在 $Z$ 处有定义，
且将 $RF$ 应用于该三角会产生特异三角。考虑特异三角的态射
$$
\xymatrix{
(F(X), F(Y), F(Z), F(f), F(g), F(h)) \ar[d] \\
(RF(X), RF(Y), RF(Z), RF(f), RF(g), RF(h))
}
$$
三个映射中的两个是同构，因而第三个也是同构。
\end{proof}

\begin{lemma}
\label{derived-lemma-direct-sum-computes}
假设和记号如情形 \ref{derived-situation-derived-functor} 中所述。
设 $X, Y$ 为 $\mathcal{D}$ 的对象。若 $X \oplus Y$ 计算 $RF$，则
$X$ 和 $Y$ 都计算 $RF$。对 $LF$ 也有类似结论。
\end{lemma}

\begin{proof}
若 $X \oplus Y$ 计算 $RF$，则 $RF(X \oplus Y) = F(X) \oplus F(Y)$。
在引理 \ref{derived-lemma-direct-sum-defined} 的证明中，我们已经看到函子
$X/S \times Y/S \to (X \oplus Y)/S$、$(s, s') \mapsto s \oplus s'$
是余终的。因而由范畴，引理
\ref{categories-lemma-cofinal-essentially-constant} 和范畴，引理
\ref{categories-lemma-characterize-essentially-constant-ind} 的表征 (4)，我们知道，
对 $\mathcal{D}'$ 的任意对象 $W$，映射
$$
\Hom_{\mathcal{D}'}(F(X \oplus Y), W)
\longrightarrow
\colim_{s : X \to X', s' : Y \to Y'} \Hom_{\mathcal{D}'}(F(X' \oplus Y'), W)
$$
是双射。由于这个箭头显然与两边的直和分解相容，我们得出映射
$$
\Hom_{\mathcal{D}'}(F(X), W)
\longrightarrow
\colim_{s : X \to X'} \Hom_{\mathcal{D}'}(F(X'), W)
$$
是双射（一个小细节从略）。因而由范畴，引理
\ref{categories-lemma-characterize-essentially-constant-ind}，我们得出
$RF$ 在 $X$ 处有定义，其值为 $F(X)$。对 $Y$ 也同样。
\end{proof}

\begin{lemma}
\label{derived-lemma-existence-computes}
假设和记号如情形 \ref{derived-situation-derived-functor} 中所述。
\begin{enumerate}
\item 若对每个对象 $X \in \Ob(\mathcal{D})$，都存在 $S$ 中的箭头
$s : X \to X'$，使 $X'$ 计算 $RF$，则 $RF$ 处处有定义。
\item 若对每个对象 $X \in \Ob(\mathcal{D})$，都存在 $S$ 中的箭头
$s : X' \to X$，使 $X'$ 计算 $LF$，则 $LF$ 处处有定义。
\end{enumerate}
\end{lemma}

\begin{proof}
这由定义是清楚的。
\end{proof}

\begin{lemma}
\label{derived-lemma-find-existence-computes}
假设和记号如情形 \ref{derived-situation-derived-functor} 中所述。
若存在子集 $\mathcal{I} \subset \Ob(\mathcal{D})$，使得
\begin{enumerate}
\item 对所有 $X \in \Ob(\mathcal{D})$，存在 $S$ 中的
$s : X \to X'$，且 $X' \in \mathcal{I}$；
\item 对 $S$ 中每个箭头 $s : X \to X'$，若 $X, X' \in \mathcal{I}$，
则映射 $F(s) : F(X) \to F(X')$ 是同构，
\end{enumerate}
那么 $RF$ 处处有定义，而每个 $X \in \mathcal{I}$ 都计算 $RF$。
对偶地，若存在子集 $\mathcal{P} \subset \Ob(\mathcal{D})$，使得
\begin{enumerate}
\item 对所有 $X \in \Ob(\mathcal{D})$，存在 $S$ 中的
$s : X' \to X$，且 $X' \in \mathcal{P}$；
\item 对 $S$ 中每个箭头 $s : X \to X'$，若 $X, X' \in \mathcal{P}$，
则映射 $F(s) : F(X) \to F(X')$ 是同构，
\end{enumerate}
那么 $LF$ 处处有定义，而每个 $X \in \mathcal{P}$ 都计算 $LF$。
\end{lemma}

\begin{proof}
设 $X$ 为 $\mathcal{D}$ 的对象。假设 (1) 意味着，$S$ 中满足
$X' \in \mathcal{I}$ 的箭头 $s : X \to X'$ 在范畴 $X/S$ 中余终。
假设 (2) 意味着 $F$ 在这个余终子范畴上常值。
这显然意味着 $F : (X/S) \to \mathcal{D}'$ 本质常值，
对 $S$ 中任意满足 $X' \in \mathcal{I}$ 的 $s : X \to X'$，其值为 $F(X')$。
\end{proof}

\begin{lemma}
\label{derived-lemma-compose-derived-functors-general}
设 $\mathcal{A}, \mathcal{B}, \mathcal{C}$ 为三角范畴。设 $S$、相应地 $S'$
分别是 $\mathcal{A}$、相应地 $\mathcal{B}$ 中与三角结构相容的饱和乘法系。
设 $F : \mathcal{A} \to \mathcal{B}$ 和 $G : \mathcal{B} \to \mathcal{C}$ 为正合函子。
以 $F' : \mathcal{A} \to (S')^{-1}\mathcal{B}$ 表示 $F$ 与局部化函子的复合。
\begin{enumerate}
\item 若 $RF'$、$RG$、$R(G \circ F)$ 都处处有定义，则存在函子的典范变换
$t : R(G \circ F) \longrightarrow RG \circ RF'$。
\item 若 $LF'$、$LG$、$L(G \circ F)$ 都处处有定义，则存在函子的典范变换
$t : LG \circ LF' \to L(G \circ F)$。
\end{enumerate}
\end{lemma}

\begin{proof}
在本证明中，我们力求审慎。因此把导出函子看作函子
$$
RF' : S^{-1}\mathcal{A} \to (S')^{-1}\mathcal{B}, \quad
R(G \circ F) : S^{-1}\mathcal{A} \to \mathcal{C}, \quad
RG : (S')^{-1}\mathcal{B} \to \mathcal{C}.
$$
以 $Q_A : \mathcal{A} \to S^{-1}\mathcal{A}$ 和
$Q_B : \mathcal{B} \to (S')^{-1}\mathcal{B}$ 表示局部化函子。
那么 $F' = Q_B \circ F$。注意，对 $\mathcal{B}$ 的每个对象 $Y$，
存在典范映射
$$
G(Y) \longrightarrow RG(Q_B(Y))
$$
换言之，存在函子变换 $t' : G \to RG \circ Q_B$。
设 $X$ 为 $\mathcal{A}$ 的对象。我们有
\begin{align*}
R(G \circ F)(Q_A(X))
& = \colim_{s : X \to X' \in S} G(F(X')) \\
& \xrightarrow{t'} \colim_{s : X \to X' \in S} RG(Q_B(F(X'))) \\
& = \colim_{s : X \to X' \in S} RG(F'(X')) \\
& = RG(\colim_{s : X \to X' \in S} F'(X')) \\
& = RG(RF'(X)).
\end{align*}
系 $F'(X')$ 在范畴 $(S')^{-1}\mathcal{B}$ 中本质常值。
因而在上图的第三个等号中，可将余极限拉入函子 $RG$ 内；
见范畴，引理 \ref{categories-lemma-image-essentially-constant} 及其证明。
我们略去这定义函子变换的证明。左导出函子的情形类似。
\end{proof}




\section{导出范畴上的导出函子}
\label{derived-section-derived-functors-classical}

\noindent
实际中，导出函子最常出现在给定阿贝尔范畴之间的加性函子时。

\begin{situation}
\label{derived-situation-classical}
此处 $F : \mathcal{A} \to \mathcal{B}$ 是阿贝尔范畴之间的加性函子。
它诱导正合函子
$$
F : K(\mathcal{A}) \to K(\mathcal{B}), \quad
K^{+}(\mathcal{A}) \to K^{+}(\mathcal{B}), \quad
K^{-}(\mathcal{A}) \to K^{-}(\mathcal{B}).
$$
见引理 \ref{derived-lemma-additive-exact-homotopy-category}。
我们还以 $F$ 表示 $F$ 与局部化函子
$K(\mathcal{B}) \to D(\mathcal{B})$ 等的复合
$K(\mathcal{A}) \to D(\mathcal{B})$、
$K^{+}(\mathcal{A}) \to D^{+}(\mathcal{B})$ 和
$K^{-}(\mathcal{A}) \to D^-(\mathcal{B})$。
这一情形导致下面要考虑的四个导出函子。
\begin{enumerate}
\item 相对于乘法系 $\text{Qis}(\mathcal{A})$ 的
$F : K(\mathcal{A}) \to D(\mathcal{B})$ 的右导出函子。
\item 相对于乘法系 $\text{Qis}^{+}(\mathcal{A})$ 的
$F : K^{+}(\mathcal{A}) \to D^{+}(\mathcal{B})$ 的右导出函子。
\item 相对于乘法系 $\text{Qis}(\mathcal{A})$ 的
$F : K(\mathcal{A}) \to D(\mathcal{B})$ 的左导出函子。
\item 相对于乘法系 $\text{Qis}^-(\mathcal{A})$ 的
$F : K^{-}(\mathcal{A}) \to D^{-}(\mathcal{B})$ 的左导出函子。
\end{enumerate}
每一种情形都是情形 \ref{derived-situation-derived-functor} 的一个例子。
\end{situation}

\noindent
下一引理消除了可能出现的一部分歧义。

\begin{lemma}
\label{derived-lemma-irrelevant}
在情形 \ref{derived-situation-classical} 中：
\begin{enumerate}
\item 设 $X$ 为 $K^{+}(\mathcal{A})$ 的对象。
$K(\mathcal{A}) \to D(\mathcal{B})$ 的右导出函子在 $X$ 处有定义，
当且仅当 $K^{+}(\mathcal{A}) \to D^{+}(\mathcal{B})$ 的右导出函子
在 $X$ 处有定义。此外，这些值典范同构。
\item 设 $X$ 为 $K^{+}(\mathcal{A})$ 的对象。
$X$ 计算 $K(\mathcal{A}) \to D(\mathcal{B})$ 的右导出函子，
当且仅当 $X$ 计算
$K^{+}(\mathcal{A}) \to D^{+}(\mathcal{B})$ 的右导出函子。
\item 设 $X$ 为 $K^{-}(\mathcal{A})$ 的对象。
$K(\mathcal{A}) \to D(\mathcal{B})$ 的左导出函子在 $X$ 处有定义，
当且仅当 $K^{-}(\mathcal{A}) \to D^{-}(\mathcal{B})$ 的左导出函子
在 $X$ 处有定义。此外，这些值典范同构。
\item 设 $X$ 为 $K^{-}(\mathcal{A})$ 的对象。
$X$ 计算 $K(\mathcal{A}) \to D(\mathcal{B})$ 的左导出函子，当且仅当 $X$ 计算
$K^{-}(\mathcal{A}) \to D^{-}(\mathcal{B})$ 的左导出函子。
\end{enumerate}
\end{lemma}

\begin{proof}
设 $X$ 为 $K^{+}(\mathcal{A})$ 的对象。
考虑 $K(\mathcal{A})$ 中的拟同构 $s : X \to X'$。
由引理 \ref{derived-lemma-complex-cohomology-bounded}，
存在拟同构 $X' \to X''$，其中 $X''$ 下有界。
因此 $X/\text{Qis}^+(\mathcal{A})$ 在
$X/\text{Qis}(\mathcal{A})$ 中余终。由此 (1) 显然成立。
(2) 直接由 (1) 推出。(3) 和 (4) 分别与 (1) 和 (2) 对偶。
\end{proof}

\noindent
给定阿贝尔范畴 $\mathcal{A}$ 的对象 $A$，可得到复形
$$
A[0] = ( \ldots \to 0 \to A \to 0 \to \ldots )
$$
其中 $A$ 位于零次。于是得到函子
$\mathcal{A} \to K(\mathcal{A})$，$A \mapsto A[0]$。
暂且把定义在某个子范畴上的函子称为部分函子。

\begin{definition}
\label{derived-definition-derived-functor}
在情形 \ref{derived-situation-classical} 中：
\begin{enumerate}
\item $F$ 的{\it 右导出函子}是与情形
\ref{derived-situation-classical} 的 (1) 和 (2) 相对应的部分函子 $RF$。
\item $F$ 的{\it 左导出函子}是与情形
\ref{derived-situation-classical} 的 (3) 和 (4) 相对应的部分函子 $LF$。
\item 若 $A[0]$ 计算 $RF$，则称 $\mathcal{A}$ 的对象 $A$
对 $F$ {\it 右无环}，或对 $RF$ {\it 无环}。
\item 若 $A[0]$ 计算 $LF$，则称 $\mathcal{A}$ 的对象 $A$
对 $F$ {\it 左无环}，或对 $LF$ {\it 无环}。
\end{enumerate}
\end{definition}

\noindent
下面几个引理给出存在足够多无环对象的一些判据。

\begin{lemma}
\label{derived-lemma-subcategory-left-resolution}
设 $\mathcal{A}$ 为阿贝尔范畴。设
$\mathcal{P} \subset \Ob(\mathcal{A})$ 为包含 $0$ 的子集，
且 $\mathcal{A}$ 的每个对象都是 $\mathcal{P}$ 中某个元素的商。
设 $a \in \mathbf{Z}$。
\begin{enumerate}
\item 给定满足 $n > a$ 时 $K^n = 0$ 的 $K^\bullet$，
存在拟同构 $P^\bullet \to K^\bullet$，使得对所有 $n$，
$P^n \in \mathcal{P}$ 且 $P^n \to K^n$ 为满射，并且
$n > a$ 时 $P^n = 0$。
\item 给定满足 $n > a$ 时 $H^n(K^\bullet) = 0$ 的 $K^\bullet$，
存在拟同构 $P^\bullet \to K^\bullet$，使得对所有 $n$，
$P^n \in \mathcal{P}$，并且 $n > a$ 时 $P^n = 0$。
\end{enumerate}
\end{lemma}

\begin{proof}
先证明 (1)。考虑如下归纳假设 $IH_n$：
存在 $P^j \in \mathcal{P}$（$j \geq n$），且 $j > a$ 时 $P^j = 0$；
存在映射 $d^j : P^j \to P^{j + 1}$（$j \geq n$）；
还存在满射 $\alpha^j : P^j \to K^j$（$j \geq n$），使得图
$$
\xymatrix{
& &
P^n \ar[d]^\alpha \ar[r] &
P^{n + 1} \ar[d]^\alpha \ar[r] &
P^{n + 2} \ar[d]^\alpha \ar[r] &
\ldots \\
\ldots \ar[r] &
K^{n - 1} \ar[r] &
K^n \ar[r] &
K^{n + 1} \ar[r] &
K^{n + 2} \ar[r] &
\ldots
}
$$
交换；对 $j \geq n$ 有 $d^{j + 1} \circ d^j = 0$；
对 $j > n$，$\alpha$ 诱导同构
$\Ker(d^j)/\Im(d^{j - 1}) \to H^j(K^\bullet)$；并且
$\alpha : \Ker(d^n) \to \Ker(d_K^n)$ 为满射。
于是选取满射
$$
P^{n - 1}
\longrightarrow
K^{n - 1} \times_{K^n} \Ker(d^n) =
K^{n - 1} \times_{\Ker(d_K^n)} \Ker(d^n)
$$
其中 $P^{n - 1}$ 属于 $\mathcal{P}$。这使我们能把上图扩张为
$$
\xymatrix{
& P^{n - 1} \ar[d]^\alpha \ar[r] &
P^n \ar[d]^\alpha \ar[r] &
P^{n + 1} \ar[d]^\alpha \ar[r] &
P^{n + 2} \ar[d]^\alpha \ar[r] &
\ldots \\
\ldots \ar[r] &
K^{n - 1} \ar[r] &
K^n \ar[r] &
K^{n + 1} \ar[r] &
K^{n + 2} \ar[r] &
\ldots
}
$$
读者不难验证，在这一选择下 $IH_{n - 1}$ 成立。

\medskip\noindent
如下完成 (1) 的证明。首先注意，$n = a + 1$ 时 $IH_n$ 成立，
因为可以直接令 $j > a$ 时 $P^j = 0$。
因此按下降归纳法可构造复形 $P^\bullet$，它满足
$n > a$ 时 $P^n = 0$，由 $\mathcal{P}$ 中的对象组成，
并带有逐项满射的拟同构
$\alpha : P^\bullet \to K^\bullet$，正合所需。

\medskip\noindent
再证明 (2)。假设意味着态射
$\tau_{\leq a}K^\bullet \to K^\bullet$
（同调，章节 \ref{homology-section-truncations}）是拟同构。
应用 (1) 得到 $P^\bullet \to \tau_{\leq a}K^\bullet$。
复合 $P^\bullet \to K^\bullet$ 即所需拟同构。
\end{proof}

\begin{lemma}
\label{derived-lemma-subcategory-right-resolution}
设 $\mathcal{A}$ 为阿贝尔范畴。设
$\mathcal{I} \subset \Ob(\mathcal{A})$ 为包含 $0$ 的子集，
且 $\mathcal{A}$ 的每个对象都是 $\mathcal{I}$ 中某个元素的子对象。
设 $a \in \mathbf{Z}$。
\begin{enumerate}
\item 给定满足 $n < a$ 时 $K^n = 0$ 的 $K^\bullet$，
存在拟同构 $K^\bullet \to I^\bullet$，使得对所有 $n$，
$K^n \to I^n$ 为单射且 $I^n \in \mathcal{I}$，并且
$n < a$ 时 $I^n = 0$；
\item 给定满足 $n < a$ 时 $H^n(K^\bullet) = 0$ 的 $K^\bullet$，
存在拟同构 $K^\bullet \to I^\bullet$，使得
$I^n \in \mathcal{I}$，并且 $n < a$ 时 $I^n = 0$。
\end{enumerate}
\end{lemma}

\begin{proof}
本引理与引理 \ref{derived-lemma-subcategory-left-resolution} 对偶。
\end{proof}

\begin{lemma}
\label{derived-lemma-subcategory-right-acyclics}
在情形 \ref{derived-situation-classical} 中，设
$\mathcal{I} \subset \Ob(\mathcal{A})$ 为具有下列性质的子集：
\begin{enumerate}
\item $\mathcal{A}$ 的每个对象都是 $\mathcal{I}$ 中某个元素的子对象；
\item 对 $\mathcal{A}$ 中任意短正合序列
$0 \to P \to Q \to R \to 0$，若
$P, Q \in \mathcal{I}$，则 $R \in \mathcal{I}$，且
$0 \to F(P) \to F(Q) \to F(R) \to 0$ 正合。
\end{enumerate}
则 $\mathcal{I}$ 的每个对象都对 $RF$ 无环。
\end{lemma}

\begin{proof}
取 $A \in \mathcal{I}$。设 $A[0] \to K^\bullet$ 为拟同构，
其中 $K^\bullet$ 下有界。于是可找到拟同构
$K^\bullet \to I^\bullet$，其中 $I^\bullet$ 下有界且每个
$I^n \in \mathcal{I}$；见引理
\ref{derived-lemma-subcategory-right-resolution}\footnote{由 (1) 可知
$\mathcal{I}$ 非空。取 $\mathcal{I}$ 中的 $P$。
短正合序列 $0 \to P \to P \to 0 \to 0$ 与假设 (2)
表明 $0$ 属于 $\mathcal{I}$，故可应用该引理。}。
因此这些消解在范畴 $A[0]/\text{Qis}^{+}(\mathcal{A})$ 中余终。
为完成证明，只须说明：对任意拟同构
$A[0] \to I^\bullet$，若 $I^\bullet$ 下有界且
$I^n \in \mathcal{I}$，则
$F(A)[0] \to F(I^\bullet)$ 是拟同构。
为此，设 $n < n_0$ 时 $I^n = 0$。当然可假设 $n_0 < 0$。
从 $n = n_0$ 开始，利用性质 (2) 和正合序列
$$
0 \to \Ker(d^n) \to I^n \to \Im(d^n) \to 0.
$$
归纳证明 $\Im(d^{n - 1}) = \Ker(d^n)$ 以及
$\Im(d^{-1})$ 都是 $\mathcal{I}$ 的元素。
此外，性质 (2) 还保证复形
$$
0 \to F(I^{n_0}) \to F(I^{n_0 + 1}) \to \ldots \to F(I^{-1}) \to
F(\Im(d^{-1})) \to 0
$$
正合。正合序列
$0 \to \Im(d^{-1}) \to I^0 \to I^0/\Im(d^{-1}) \to 0$
意味着 $I^0/\Im(d^{-1})$ 是 $\mathcal{I}$ 的元素。
% P02-E0034: frozen source has grammatical “is an elements of”.
正合序列 $0 \to A \to I^0/\Im(d^{-1}) \to \Im(d^0) \to 0$
继而意味着 $\Im(d^0) = \Ker(d^1)$ 是 $\mathcal{I}$ 的元素；
此后同前可继续证明，对所有 $n > 0$，
$\Im(d^{n - 1}) = \Ker(d^n)$ 都是 $\mathcal{I}$ 的元素。
对上述每个短正合序列应用 $F$ 并使用 (2)，可见
% P02-E0033: frozen source says “is an isomorphism”; preceding requirement and argument establish quasi-isomorphism.
$F(A)[0] \to F(I^\bullet)$ 是同构，正合所需。
\end{proof}

\begin{lemma}
\label{derived-lemma-subcategory-left-acyclics}
在情形 \ref{derived-situation-classical} 中，设
$\mathcal{P} \subset \Ob(\mathcal{A})$ 为具有下列性质的子集：
\begin{enumerate}
\item $\mathcal{A}$ 的每个对象都是 $\mathcal{P}$ 中某个元素的商；
\item 对 $\mathcal{A}$ 中任意短正合序列
$0 \to P \to Q \to R \to 0$，若
$Q, R \in \mathcal{P}$，则 $P \in \mathcal{P}$，且
$0 \to F(P) \to F(Q) \to F(R) \to 0$ 正合。
\end{enumerate}
则 $\mathcal{P}$ 的每个对象都对 $LF$ 无环。
\end{lemma}

\begin{proof}
与引理 \ref{derived-lemma-subcategory-right-acyclics} 的证明对偶。
\end{proof}







\section{高阶导出函子}
\label{derived-section-higher-derived}

\noindent
下面这个简单引理表明，右导出函子“向右移动”。

\begin{lemma}
\label{derived-lemma-negative-vanishing}
设 $F : \mathcal{A} \to \mathcal{B}$ 为阿贝尔范畴之间的加性函子。
设 $K^\bullet$ 为 $\mathcal{A}$ 的复形，且 $a \in \mathbf{Z}$。
\begin{enumerate}
\item 若对所有 $i < a$ 有 $H^i(K^\bullet) = 0$，且
$RF$ 在 $K^\bullet$ 处有定义，则对所有 $i < a$ 有
$H^i(RF(K^\bullet)) = 0$。
\item 若 $RF$ 在 $K^\bullet$ 和 $\tau_{\leq a}K^\bullet$ 处都有定义，
则对所有 $i \leq a$ 有
$H^i(RF(\tau_{\leq a}K^\bullet)) = H^i(RF(K^\bullet))$。
\end{enumerate}
\end{lemma}

\begin{proof}
假设 $K^\bullet$ 满足 (1) 的条件。设
$s : K^\bullet \to L^\bullet$ 为任意拟同构。
由对 $K^\bullet$ 的假设，$K^\bullet \to \tau_{\geq a}L^\bullet$
也是拟同构。因此，在范畴
$K^\bullet/\text{Qis}^{+}(\mathcal{A})$ 中，满足 $L^n = 0$
（$n < a$）的拟同构 $s : K^\bullet \to L^\bullet$ 是余终的。
由范畴，引理 \ref{categories-lemma-cofinal-essentially-constant}，
可知 $RF$ 是这些 $s$ 所对应的本质常值 ind-对象
$F(L^\bullet)$ 的值。这意味着
$\text{id} : RF(K^\bullet) \to RF(K^\bullet)$ 分解通过某个
$F(L^\bullet)$，其中复形 $L^\bullet$ 满足 $n < a$ 时 $L^n = 0$。
由此，对 $i < a$ 有 $H^i(RF(K^\bullet)) = 0$。

\medskip\noindent
为证明 (2)，使用注记 \ref{derived-remark-truncation-distinguished-triangle}
中的特异三角
$$
\tau_{\leq a}K^\bullet \to K^\bullet \to \tau_{\geq a + 1}K^\bullet
\to (\tau_{\leq a}K^\bullet)[1]
$$
并由引理 \ref{derived-lemma-2-out-of-3-defined} 得到：
$RF$ 在 $\tau_{\geq a + 1}K^\bullet$ 处也有定义，并且在
$D(\mathcal{B})$ 中有特异三角
$$
RF(\tau_{\leq a}K^\bullet) \to RF(K^\bullet) \to RF(\tau_{\geq a + 1}K^\bullet)
\to RF(\tau_{\leq a}K^\bullet)[1]
$$
由 (1)，$RF(\tau_{\geq a + 1}K^\bullet)$ 在次数 $< a + 1$ 上
上同调为零。该特异三角的长正合上同调序列于是给出所需结论。
\end{proof}

\begin{definition}
\label{derived-definition-higher-derived-functors}
设 $F : \mathcal{A} \to \mathcal{B}$ 为阿贝尔范畴之间的加性函子。
假设 $RF : D^{+}(\mathcal{A}) \to D^{+}(\mathcal{B})$ 处处有定义。
设 $i \in \mathbf{Z}$。$F$ 的{\it 第 $i$ 个右导出函子 $R^iF$}是函子
$$
R^iF = H^i \circ RF :
\mathcal{A}
\longrightarrow
\mathcal{B}
$$
\end{definition}

\noindent
下面的引理表明，除非函子左正合，否则取其右导出函子其实没有多大意义。

\begin{lemma}
\label{derived-lemma-left-exact-higher-derived}
设 $F : \mathcal{A} \to \mathcal{B}$ 为阿贝尔范畴之间的加性函子，
并假设 $RF : D^{+}(\mathcal{A}) \to D^{+}(\mathcal{B})$ 处处有定义。
\begin{enumerate}
\item 对 $i < 0$ 有 $R^iF = 0$；
\item $R^0F$ 左正合；
\item 映射 $F \to R^0F$ 为同构，当且仅当 $F$ 左正合。
\end{enumerate}
\end{lemma}

\begin{proof}
设 $A$ 为 $\mathcal{A}$ 的对象。
% P02-E0032: frozen source math has an unmatched opening parenthesis.
由引理 \ref{derived-lemma-negative-vanishing}，对 $i < 0$ 有
$H^i(RF(A[0]) = 0$。这证明了 (1)。

\medskip\noindent
设 $0 \to A \to B \to C \to 0$ 为 $\mathcal{A}$ 的短正合序列。
由引理 \ref{derived-lemma-derived-canonical-delta-functor}，可在
$D^{+}(\mathcal{A})$ 中得到特异三角
$(A[0], B[0], C[0], a, b, c)$。
由长正合上同调序列（以及上面证明的 $i < 0$ 时的消失性），
可知 $0 \to R^0F(A) \to R^0F(B) \to R^0F(C)$ 正合。
故 $R^0F$ 左正合。当然，这也证明了：若
$F \to R^0F$ 为同构，则 $F$ 左正合。

\medskip\noindent
假设 $F$ 左正合。回忆 $RF(A[0])$ 是本质常值系
$F(K^\bullet)$ 的值，其中 $s : A[0] \to K^\bullet$ 遍历拟同构。
由此 $R^0F(A)$ 是本质常值系 $H^0(F(K^\bullet))$ 的值，
其中 $s : A[0] \to K^\bullet$ 遍历拟同构；见范畴，引理
\ref{categories-lemma-image-essentially-constant}。
但若 $s : A[0] \to K^\bullet$ 是拟同构，则
$A[0] \to \tau_{\geq 0}K^\bullet$ 是拟同构。
因此在范畴 $A[0]/\text{Qis}^{+}(\mathcal{A})$ 中，满足
$n < 0$ 时 $K^n = 0$ 的拟同构 $s : A[0] \to K^\bullet$ 是余终的。
由范畴，引理 \ref{categories-lemma-cofinal-essentially-constant}，
可知只须考虑这样的 $s$。此外，对这样的 $s$，序列
$$
0 \to A \to K^0 \to K^1
$$
正合。由于 $F$ 左正合，
$0 \to F(A) \to F(K^0) \to F(K^1)$ 也正合。
于是 $F(A) \to H^0(F(K^\bullet))$ 为同构，
而且该系实际上是取值为 $F(A)$ 的常值系。
故 $R^0F = F$，正合所需。
\end{proof}

\begin{lemma}
\label{derived-lemma-F-acyclic}
设 $F : \mathcal{A} \to \mathcal{B}$ 为阿贝尔范畴之间的加性函子，
并假设 $RF : D^{+}(\mathcal{A}) \to D^{+}(\mathcal{B})$ 处处有定义。
设 $A$ 为 $\mathcal{A}$ 的对象。
\begin{enumerate}
\item $A$ 对 $F$ 右无环，当且仅当
$F(A) \to R^0F(A)$ 为同构，且对所有 $i > 0$ 有 $R^iF(A) = 0$；
\item 若 $F$ 左正合，则 $A$ 对 $F$ 右无环，当且仅当
对所有 $i > 0$ 有 $R^iF(A) = 0$。
\end{enumerate}
\end{lemma}

\begin{proof}
若 $A$ 对 $F$ 右无环，则 $RF(A[0]) = F(A)[0]$；特别地，
$F(A) \to R^0F(A)$ 为同构，且 $i \not = 0$ 时
$R^iF(A) = 0$。反之，若 $F(A) \to R^0F(A)$ 为同构，且对所有
$i > 0$ 有 $R^iF(A) = 0$，则由引理
\ref{derived-lemma-left-exact-higher-derived} 的 (1)，
$F(A[0]) \to RF(A[0])$ 为拟同构，故 $A$ 无环。
若 $F$ 左正合，则 $F = R^0F$；见引理
\ref{derived-lemma-left-exact-higher-derived}。
\end{proof}

\begin{lemma}
\label{derived-lemma-F-acyclic-ses}
设 $F : \mathcal{A} \to \mathcal{B}$ 为阿贝尔范畴之间的左正合函子，
并假设 $RF : D^{+}(\mathcal{A}) \to D^{+}(\mathcal{B})$ 处处有定义。
设 $0 \to A \to B \to C \to 0$ 为 $\mathcal{A}$ 的短正合序列。
\begin{enumerate}
\item 若 $A$ 和 $C$ 都对 $F$ 右无环，则 $B$ 亦然。
\item 若 $A$ 和 $B$ 都对 $F$ 右无环，则 $C$ 亦然。
\item 若 $B$ 和 $C$ 都对 $F$ 右无环，且
$F(B) \to F(C)$ 为满射，则 $A$ 对 $F$ 右无环。
\end{enumerate}
在上述三种情形中，
$$
0 \to F(A) \to F(B) \to F(C) \to 0
$$
都是 $\mathcal{B}$ 的短正合序列。
\end{lemma}

\begin{proof}
由引理 \ref{derived-lemma-derived-canonical-delta-functor}，
可在 $D^{+}(\mathcal{A})$ 中得到特异三角
$(A[0], B[0], C[0], a, b, c)$。
由于 $RF$ 是正合函子，并且 $i < 0$ 时 $R^iF = 0$、
$R^0F = F$（引理 \ref{derived-lemma-left-exact-higher-derived}），
可在阿贝尔范畴 $\mathcal{B}$ 中得到正合上同调序列
$$
0 \to F(A) \to F(B) \to F(C) \to R^1F(A) \to \ldots
$$
所以本引理由引理 \ref{derived-lemma-F-acyclic} 中对无环对象的刻画推出。
\end{proof}

\begin{lemma}
\label{derived-lemma-right-derived-delta-functor}
设 $F : \mathcal{A} \to \mathcal{B}$ 为阿贝尔范畴之间的加性函子，
并假设 $RF : D^{+}(\mathcal{A}) \to D^{+}(\mathcal{B})$ 处处有定义。
\begin{enumerate}
\item 函子 $R^iF$（$i \geq 0$）带有从
$\mathcal{A} \to \mathcal{B}$ 出发的典范 $\delta$-函子结构；见
同调，定义 \ref{homology-definition-cohomological-delta-functor}。
\item 若 $\mathcal{A}$ 的每个对象都是某个对 $F$ 右无环对象的子对象，
则 $\{R^iF, \delta\}_{i \geq 0}$ 是泛 $\delta$-函子；见
同调，定义 \ref{homology-definition-universal-delta-functor}。
\end{enumerate}
\end{lemma}

\begin{proof}
函子 $\mathcal{A} \to \text{Comp}^{+}(\mathcal{A})$，
$A \mapsto A[0]$ 是正合的。函子
$\text{Comp}^{+}(\mathcal{A}) \to D^{+}(\mathcal{A})$
是 $\delta$-函子；见引理 \ref{derived-lemma-derived-canonical-delta-functor}。
函子 $RF : D^{+}(\mathcal{A}) \to D^{+}(\mathcal{B})$ 是正合的。
最后，函子 $H^0 : D^{+}(\mathcal{B}) \to \mathcal{B}$
是同调函子；见定义 \ref{derived-definition-unbounded-derived-category}。
因此，由引理 \ref{derived-lemma-compose-delta-functor-homological} 和
引理 \ref{derived-lemma-exact-compose-delta-functor}，可得到一个
$\delta$-函子结构。(2) 由同调，引理
\ref{homology-lemma-efface-implies-universal} 以及引理
\ref{derived-lemma-F-acyclic} 中对无环对象的描述推出。
\end{proof}

\begin{lemma}[Leray 无环性引理]
\label{derived-lemma-leray-acyclicity}
设 $F : \mathcal{A} \to \mathcal{B}$ 为阿贝尔范畴之间的加性函子。
设 $A^\bullet$ 为由对 $F$ 右无环的对象组成的下有界复形，
且 $RF$ 在 $A^\bullet$ 处有定义\footnote{例如，当
$RF : D^{+}(\mathcal{A}) \to D^{+}(\mathcal{B})$ 处处有定义时成立。}。
典范映射
$$
F(A^\bullet) \longrightarrow RF(A^\bullet)
$$
是 $D^{+}(\mathcal{B})$ 中的同构；即 $A^\bullet$ 计算 $RF$。
\end{lemma}

\begin{proof}
设 $A^\bullet$ 为由对 $F$ 右无环的对象组成的有界复形。
我们断言 $RF$ 在 $A^\bullet$ 处有定义，且
$F(A^\bullet) \to RF(A^\bullet)$ 是 $D^+(\mathcal{B})$ 中的同构。
事实上，由定义 \ref{derived-definition-derived-functor}，
这对至多含有一个非零且对 $F$ 右无环对象的复形成立。
再设 $n \not \in [a, b]$ 时 $A^n = 0$。
使用“笨”截断可得到逐项分裂的复形短正合序列
$$
0 \to \sigma_{\geq a + 1} A^\bullet \to A^\bullet \to
\sigma_{\leq a} A^\bullet \to 0
$$
见同调，章节 \ref{homology-section-truncations}。
从而得到特异三角
$(\sigma_{\geq a + 1} A^\bullet, A^\bullet, \sigma_{\leq a} A^\bullet)$。
由归纳假设，$RF$ 在外侧两个复形处有定义，且这两个复形都计算 $RF$。
由引理 \ref{derived-lemma-2-out-of-3-computes}，中间的复形亦然。

\medskip\noindent
设 $A^\bullet$ 是由无环对象组成的下有界复形，
且 $RF$ 在 $A^\bullet$ 处有定义。为证明
$F(A^\bullet) \to RF(A^\bullet)$ 是 $D^{+}(\mathcal{B})$ 中的同构，
只须证明对所有 $i$，
$H^i(F(A^\bullet)) \to H^i(RF(A^\bullet))$ 都是同构。
取 $i$。考虑逐项分裂的复形短正合序列
$$
0 \to \sigma_{\geq i + 2} A^\bullet \to A^\bullet \to
\sigma_{\leq i + 1} A^\bullet \to 0.
$$
注意，它诱导逐项分裂的短正合序列
$$
0 \to \sigma_{\geq i + 2} F(A^\bullet) \to F(A^\bullet) \to
\sigma_{\leq i + 1} F(A^\bullet) \to 0.
$$
因此得到特异三角
$$
(\sigma_{\geq i + 2} A^\bullet, A^\bullet,
\sigma_{\leq i + 1} A^\bullet)
\quad\text{且}\quad
(\sigma_{\geq i + 2} F(A^\bullet), F(A^\bullet),
\sigma_{\leq i + 1} F(A^\bullet))
$$
由于 $RF$ 在 $A^\bullet$ 处有定义（依假设），并且在
$\sigma_{\leq i + 1}A^\bullet$ 处有定义（依第一段），
故 $RF$ 在 $\sigma_{\geq i + 2}A^\bullet$ 处也有定义，
而且得到特异三角
$$
(RF(\sigma_{\geq i + 2} A^\bullet), RF(A^\bullet),
RF(\sigma_{\leq i + 1} A^\bullet))
$$
见引理 \ref{derived-lemma-2-out-of-3-defined}。
利用这些特异三角，可得到正合序列之间的映射
$$
\xymatrix{
H^i(\sigma_{\geq i + 2} F(A^\bullet)) \ar[r] \ar[d] &
H^i(F(A^\bullet)) \ar[r] \ar[d]^\alpha &
H^i(\sigma_{\leq i + 1} F(A^\bullet)) \ar[r] \ar[d]^\beta &
H^{i + 1}(\sigma_{\geq i + 2} F(A^\bullet)) \ar[d] \\
H^i(RF(\sigma_{\geq i + 2} A^\bullet)) \ar[r] &
H^i(RF(A^\bullet)) \ar[r] &
H^i(RF(\sigma_{\leq i + 1} A^\bullet)) \ar[r] &
H^{i + 1}(RF(\sigma_{\geq i + 2} A^\bullet))
}
$$
由第一段的结果，映射 $\beta$ 是同构。
直接看出左上角和右上角的对象为零。因此，为完成证明，只须说明
$H^i(RF(\sigma_{\geq i + 2} A^\bullet)) = 0$ 且
$H^{i + 1}(RF(\sigma_{\geq i + 2} A^\bullet)) = 0$。
这立即由引理 \ref{derived-lemma-negative-vanishing} 推出。
\end{proof}

\begin{proposition}
\label{derived-proposition-enough-acyclics}
\begin{slogan}
阿贝尔范畴之间的函子可借助由该函子的无环对象组成的复形作消解，
扩张到（下有界或上有界）导出范畴。
\end{slogan}
设 $F : \mathcal{A} \to \mathcal{B}$ 为阿贝尔范畴之间的加性函子。
\begin{enumerate}
\item 若 $\mathcal{A}$ 的每个对象都嵌入某个对 $RF$ 无环的对象，
则 $RF$ 在整个 $K^{+}(\mathcal{A})$ 上有定义，并得到正合函子
$$
RF : D^{+}(\mathcal{A}) \longrightarrow D^{+}(\mathcal{B})
$$
见 (\ref{derived-equation-everywhere})。此外，任意由对 $RF$ 无环的项
组成的下有界复形 $A^\bullet$ 都计算 $RF$。
\item 若 $\mathcal{A}$ 的每个对象都是某个对 $LF$ 无环对象的商，
% P02-E0035: frozen source omits the article in “is quotient of”.
则 $LF$ 在整个 $K^{-}(\mathcal{A})$ 上有定义，并得到正合函子
$$
LF : D^{-}(\mathcal{A}) \longrightarrow D^{-}(\mathcal{B})
$$
见 (\ref{derived-equation-everywhere})。此外，任意由对 $LF$ 无环的项
组成的上有界复形 $A^\bullet$ 都计算 $LF$。
\end{enumerate}
\end{proposition}

\begin{proof}
假设 $\mathcal{A}$ 的每个对象都嵌入某个对 $RF$ 无环的对象。
设 $\mathcal{I}$ 为对 $RF$ 无环的对象所成的集合。
设 $K^\bullet$ 为 $\mathcal{A}$ 中的下有界复形。
由引理 \ref{derived-lemma-subcategory-right-resolution}，
存在拟同构 $\alpha : K^\bullet \to I^\bullet$，其中
$I^\bullet$ 下有界且 $I^n \in \mathcal{I}$。
所以，为证明 (1)，只须说明：若
$s : I^\bullet \to (I')^\bullet$ 是由 $\mathcal{I}$ 中对象组成的
下有界复形之间的拟同构，则
$F(I^\bullet) \to F((I')^\bullet)$ 是拟同构；见引理
\ref{derived-lemma-find-existence-computes}。
注意，锥 $C(s)^\bullet$ 是无环下有界复形，且每一项都属于
$\mathcal{I}$。故只须说明：给定每一项都属于 $\mathcal{I}$ 的
无环下有界复形 $I^\bullet$，复形 $F(I^\bullet)$ 是无环的。

\medskip\noindent
设 $n < n_0$ 时 $I^n = 0$。令 $J^n = \Im(d^n)$，
便把 $I^\bullet$ 分解成短正合序列
$0 \to J^n \to I^{n + 1} \to J^{n + 1} \to 0$（$n \geq n_0$）。
由引理 \ref{derived-lemma-derived-canonical-delta-functor}，
这些序列在 $D^+(\mathcal{A})$ 中诱导特异三角
$(J^n, I^{n + 1}, J^{n + 1})$。
对每个 $k \in \mathbf{Z}$，用 $H_k$ 表示如下断言：
对所有 $n \leq k$，对象 $J^n$ 属于 $\mathcal{I}$。
当 $k < n_0$ 时，$H_k$ 显然成立。若 $H_n$ 成立，
则引理 \ref{derived-lemma-2-out-of-3-computes} 表明
$J^{n + 1}$ 属于 $\mathcal{I}$，从而 $H_{n + 1}$ 成立。
由命题 \ref{derived-proposition-derived-functor}，有特异三角
$(RF(J^n), RF(I^{n + 1}), RF(J^{n + 1}))$。
由于 $J^n, I^{n + 1}, J^{n + 1}$ 都属于 $\mathcal{I}$，
与该特异三角相伴的长正合上同调序列
(\ref{derived-equation-long-exact-cohomology-sequence-D}) 缩并为正合序列
$$
0 \to F(J^n) \to F(I^{n + 1}) \to F(J^{n + 1}) \to 0
$$
这又证明了 $F(I^\bullet)$ 正合。

\medskip\noindent
$LF$ 的情形之证明与此对偶。
\end{proof}

\begin{lemma}
\label{derived-lemma-right-derived-exact-functor}
设 $F : \mathcal{A} \to \mathcal{B}$ 为阿贝尔范畴之间的正合函子。
则：
\begin{enumerate}
\item $\mathcal{A}$ 的每个对象都对 $F$ 右无环；
\item $RF : D^{+}(\mathcal{A}) \to D^{+}(\mathcal{B})$ 处处有定义；
\item $RF : D(\mathcal{A}) \to D(\mathcal{B})$ 处处有定义；
\item 每个复形都计算 $RF$；换言之，对所有复形，典范映射
$F(K^\bullet) \to RF(K^\bullet)$ 都是同构；
\item 对 $i \not = 0$ 有 $R^iF = 0$。
\end{enumerate}
\end{lemma}

\begin{proof}
这是因为 $F$ 把无环复形送到无环复形，并把拟同构送到拟同构。
细节从略。
\end{proof}









\section{导出范畴的三角子范畴}
\label{derived-section-triangulated-sub}

\noindent
设 $\mathcal{A}$ 为阿贝尔范畴。本节考察某些严格满饱和三角子范畴
$\mathcal{D}' \subset D(\mathcal{A})$。

\medskip\noindent
设 $\mathcal{B} \subset \mathcal{A}$ 为弱 Serre 子范畴；见
同调，定义 \ref{homology-definition-serre-subcategory} 和
引理 \ref{homology-lemma-characterize-weak-serre-subcategory}。
% P02-E0036: frozen source omits “be” in “We let ... the full subcategory”.
令 $D_\mathcal{B}(\mathcal{A})$ 为 $D(\mathcal{A})$ 的满子范畴，
其对象为
$$
\Ob(D_\mathcal{B}(\mathcal{A}))
=
\{X \in \Ob(D(\mathcal{A})) \mid
H^n(X) \text{ 是下列范畴的对象： }\mathcal{B}\text{ 对所有 }n\}
$$
还定义
$D^{+}_\mathcal{B}(\mathcal{A}) =
D^{+}(\mathcal{A}) \cap D_\mathcal{B}(\mathcal{A})$，
其他有界版本同理。

\begin{lemma}
\label{derived-lemma-cohomology-in-serre-subcategory}
设 $\mathcal{A}$ 为阿贝尔范畴，
$\mathcal{B} \subset \mathcal{A}$ 为弱 Serre 子范畴。
则 $D_\mathcal{B}(\mathcal{A})$ 是 $D(\mathcal{A})$ 的
严格满饱和三角子范畴。各有界版本亦然。
\end{lemma}

\begin{proof}
显然，$D_\mathcal{B}(\mathcal{A})$ 是在平移函子下保持的加性子范畴。
若 $X \oplus Y$ 属于 $D_\mathcal{B}(\mathcal{A})$，则
% P02-E0037: frozen source says “kernels of maps between maps of objects”; one “maps” is extraneous.
$H^n(X)$ 和 $H^n(Y)$ 都是 $\mathcal{B}$ 中对象之间映射的核，
因为 $H^n(X \oplus Y) = H^n(X) \oplus H^n(Y)$。
所以 $X$ 和 $Y$ 都属于 $D_\mathcal{B}(\mathcal{A})$。由
引理 \ref{derived-lemma-triangulated-subcategory}，只须证明：
给定特异三角 $(X, Y, Z, f, g, h)$，若 $X$ 和 $Y$ 属于
$D_\mathcal{B}(\mathcal{A})$，则 $Z$ 是
$D_\mathcal{B}(\mathcal{A})$ 的对象。
长正合上同调序列 (\ref{derived-equation-long-exact-cohomology-sequence-D})
和弱 Serre 子范畴的定义（见同调，定义
\ref{homology-definition-serre-subcategory}）表明，
对所有 $n$，$H^n(Z)$ 都是 $\mathcal{B}$ 的对象。
故 $Z$ 是 $D_\mathcal{B}(\mathcal{A})$ 的对象。
\end{proof}

\noindent
继续假设 $\mathcal{B}$ 是阿贝尔范畴 $\mathcal{A}$ 的弱 Serre 子范畴。
那么 $\mathcal{B}$ 是阿贝尔范畴，且包含函子
$\mathcal{B} \to \mathcal{A}$ 正合。
所以得到导出函子 $D(\mathcal{B}) \to D(\mathcal{A})$；见引理
\ref{derived-lemma-right-derived-exact-functor}。显然，函子
$D(\mathcal{B}) \to D(\mathcal{A})$ 分解通过典范正合函子
\begin{equation}
\label{derived-equation-compare}
D(\mathcal{B}) \longrightarrow D_\mathcal{B}(\mathcal{A})
\end{equation}
毕竟，由 $\mathcal{B}$ 中对象组成的复形当然给出
$D_\mathcal{B}(\mathcal{A})$ 的对象；而
$D_\mathcal{B}(\mathcal{A})$ 中的特异三角，恰好是顶点都属于
$D_\mathcal{B}(\mathcal{A})$ 的 $D(\mathcal{A})$ 中的特异三角。
由于 $D(\mathcal{B}) \to D(\mathcal{A})$ 正合，故上述函子正合。
类似地，对各有界版本得到函子
$D^+(\mathcal{B}) \to D^+_\mathcal{B}(\mathcal{A})$、
$D^-(\mathcal{B}) \to D^-_\mathcal{B}(\mathcal{A})$ 和
$D^b(\mathcal{B}) \to D^b_\mathcal{B}(\mathcal{A})$。
很多情形中的一个关键问题是所展示的函子是否为等价。

\medskip\noindent
现在假设 $\mathcal{B}$ 是 $\mathcal{A}$ 的 Serre 子范畴。
此时有商函子 $\mathcal{A} \to \mathcal{A}/\mathcal{B}$；见
同调，引理 \ref{homology-lemma-serre-subcategory-is-kernel}。
此时 $D_\mathcal{B}(\mathcal{A})$ 是函子
$D(\mathcal{A}) \to D(\mathcal{A}/\mathcal{B})$ 的核。
因此由引理 \ref{derived-lemma-universal-property-quotient} 得到典范函子
$$
D(\mathcal{A})/D_\mathcal{B}(\mathcal{A})
\longrightarrow
D(\mathcal{A}/\mathcal{B})
$$
各有界版本同理。

\begin{lemma}
\label{derived-lemma-derived-of-quotient}
设 $\mathcal{A}$ 为阿贝尔范畴，
$\mathcal{B} \subset \mathcal{A}$ 为 Serre 子范畴。
则 $D(\mathcal{A}) \to D(\mathcal{A}/\mathcal{B})$ 本质满。
\end{lemma}

\begin{proof}
我们使用同调，引理 \ref{homology-lemma-serre-subcategory-is-kernel}
的证明中对范畴 $\mathcal{A}/\mathcal{B}$ 的描述。
设 $(X^\bullet, d^\bullet)$ 为 $\mathcal{A}/\mathcal{B}$ 的复形。
这意味着 $X^i$ 是 $\mathcal{A}$ 的对象，
$d^i : X^i \to X^{i + 1}$ 是 $\mathcal{A}/\mathcal{B}$ 中的态射，
且在 $\mathcal{A}/\mathcal{B}$ 中有 $d^i \circ d^{i - 1} = 0$。

\medskip\noindent
对 $i \geq 0$，可写 $d^i = (s^i, f^i)$，其中
$s^i : Y^i \to X^i$ 是 $\mathcal{A}$ 中的态射，
它的核和余核都属于 $\mathcal{B}$（等价地，$s^i$ 在商范畴中成为同构），
而 $f^i : Y^i \to X^{i + 1}$ 是 $\mathcal{A}$ 中的态射。
我们将归纳构造交换图
$$
\xymatrix{
& (X')^1 \ar@{..>}[r] & (X')^2 \ar@{..>}[r] & \ldots \\
X^0 \ar@{..>}[ru] &
X^1 \ar@{..>}[u] &
X^2 \ar@{..>}[u] &
\ldots \\
Y^0 \ar[u]_{s^0} \ar[ru]_{f^0} &
Y^1 \ar[u]_{s^1} \ar[ru]_{f^1} &
Y^2 \ar[u]_{s^2} \ar[ru]_{f^2} &
\ldots
}
$$
其中竖直箭头 $X^i \to (X')^i$ 在商范畴中成为同构。
具体地，先令
$(X')^1 = \Coker(Y^0 \to X^0 \oplus X^1)$
（更准确地说，是箭头 $s^0$ 和 $f^0$ 所成图的推出），
从而得到第一个交换图。接着取
$(X')^2 = \Coker(Y^1 \to (X')^1 \oplus X^2)$，依此类推。
再令 $n \leq 0$ 时 $(X')^n = X^n$，便可看出映射
$(X^\bullet, d^\bullet) \to ((X')^\bullet, (d')^\bullet)$
是 $\mathcal{A}/\mathcal{B}$ 中复形的同构。
所以可以假设：对 $n \geq 0$，
$d^n : X^n \to X^{n + 1}$ 由 $\mathcal{A}$ 中的映射
$X^n \to X^{n + 1}$ 给出。

\medskip\noindent
对偶地，对 $i < 0$ 可写 $d^i = (g^i, t^{i + 1})$，其中
$t^{i + 1} : X^{i + 1} \to Z^{i + 1}$ 是商范畴中的同构，
而 $g^i : X^i \to Z^{i + 1}$ 是态射。
我们将归纳构造交换图
% P02-E0040: frozen diagram uses subscripted t and g indices where the declared maps require superscripts.
$$
\xymatrix{
\ldots &
Z^{-2} &
Z^{-1} &
Z^0 \\
\ldots &
X^{-2} \ar[u]_{t_{-2}} \ar[ru]_{g_{-2}} &
X^{-1} \ar[u]_{t_{-1}} \ar[ru]_{g_{-1}} &
X^0 \ar[u]_{t^0} \\
\ldots &
(X')^{-2} \ar@{..>}[u] \ar@{..>}[r] &
(X')^{-1} \ar@{..>}[u] \ar@{..>}[ru]
}
$$
其中竖直箭头 $(X')^i \to X^i$ 在商范畴中成为同构。
具体地，取 $(X')^{-1} = X^{-1} \times_{Z^0} X^0$，
再取 $(X')^{-2} = X^{-2} \times_{Z^{-1}} (X')^{-1}$，依此类推。
再令 $n \geq 0$ 时 $(X')^n = X^n$，便可看出映射
$((X')^\bullet, (d')^\bullet) \to (X^\bullet, d^\bullet)$
是 $\mathcal{A}/\mathcal{B}$ 中复形的同构。
所以可以假设：对所有 $n \in \mathbf{Z}$，
$d^n : X^n \to X^{n + 1}$ 由 $\mathcal{A}$ 中的映射
$d^n : X^n \to X^{n + 1}$ 给出。

\medskip\noindent
在这种情形下，复合 $d^n \circ d^{n - 1}$ 在
$\mathcal{A}/\mathcal{B}$ 中为零。若对 $n > 0$，把 $X^n$ 替换为
$$
(X')^n = X^n/\sum\nolimits_{0 < k \leq n} \Im(\Im(X^{k - 2} \to X^k) \to X^n)
$$
则对 $n \geq 0$，复合 $d^n \circ d^{n - 1}$ 为零。
（与上面第二段类似，可得到复形的同构
$(X^\bullet, d^\bullet) \to ((X')^\bullet, (d')^\bullet)$。）
最后，对 $n < 0$，把 $X^n$ 替换为
$$
(X')^n = \bigcap\nolimits_{n \leq k < 0}
(X^n \to X^k)^{-1}\Ker(X^k \to X^{k + 2})
$$
并用同样的方法论证，便得到 $\mathcal{A}$ 中的复形，
其在 $\mathcal{A}/\mathcal{B}$ 中的像与给定复形同构。
\end{proof}

\begin{lemma}
\label{derived-lemma-quotient-by-serre-easy}
设 $\mathcal{A}$ 为阿贝尔范畴，
$\mathcal{B} \subset \mathcal{A}$ 为 Serre 子范畴。
假设函子 $v : \mathcal{A} \to \mathcal{A}/\mathcal{B}$
有左伴随 $u : \mathcal{A}/\mathcal{B} \to \mathcal{A}$，
且 $vu \cong \text{id}$。则
$$
D(\mathcal{A})/D_\mathcal{B}(\mathcal{A}) = D(\mathcal{A}/\mathcal{B})
$$
各有界版本亦然。
\end{lemma}

\begin{proof}
由引理 \ref{derived-lemma-derived-of-quotient}，函子
$D(v) : D(\mathcal{A}) \to D(\mathcal{A}/\mathcal{B})$ 本质满。
对 $D(\mathcal{A})$ 的对象 $X$，伴随映射
$c_X : uvX \to X$ 在 $D(\mathcal{A}/\mathcal{B})$ 中映为同构，
因为假设 $vu \cong \text{id}$ 意味着 $vuv \cong v$。
所以，在特异三角 $(uvX, X, Z, c_X, g, h)$ 中，
考察长正合上同调序列可知，对象 $Z$ 属于
$D_\mathcal{B}(\mathcal{A})$。故 $c_X$ 是用于定义商范畴
$D(\mathcal{A})/D_\mathcal{B}(\mathcal{A})$ 的乘法系的元素。
从而在 $D(\mathcal{A})/D_\mathcal{B}(\mathcal{A})$ 中有
$uvX \cong X$。
% P02-E0038: frozen source has an extra closing parenthesis in Ob(A)).
对 $X, Y \in \Ob(\mathcal{A}))$，映射
$$
\Hom_{D(\mathcal{A})/D_\mathcal{B}(\mathcal{A})}(X, Y)
\longrightarrow
\Hom_{D(\mathcal{A}/\mathcal{B})}(vX, vY)
$$
是双射，因为 $u$ 给出逆映射（由上述讨论）。
\end{proof}

\noindent
对某些 Serre 子范畴 $\mathcal{B} \subset \mathcal{A}$，
可以证明函子 $D(\mathcal{B}) \to D_\mathcal{B}(\mathcal{A})$ 满忠实。

\begin{lemma}
\label{derived-lemma-fully-faithful-embedding}
设 $\mathcal{A}$ 为阿贝尔范畴，
$\mathcal{B} \subset \mathcal{A}$ 为 Serre 子范畴。
假设对每个满射 $X \to Y$，其中
$X \in \Ob(\mathcal{A})$ 且 $Y \in \Ob(\mathcal{B})$，
存在 $X' \subset X$，$X' \in \Ob(\mathcal{B})$，且后者满射到 $Y$。
则 (\ref{derived-equation-compare}) 中的函子
$D^-(\mathcal{B}) \to D^-_\mathcal{B}(\mathcal{A})$ 是等价。
\end{lemma}

\begin{proof}
设 $X^\bullet$ 为 $\mathcal{A}$ 的上有界复形，且对所有
$i \in \mathbf{Z}$ 有 $H^i(X^\bullet) \in \Ob(\mathcal{B})$。
再假设对所有 $i \in \mathbf{Z}$ 给定
$B^i \subset X^i$，$B^i \in \Ob(\mathcal{B})$。
断言：存在子复形 $Y^\bullet \subset X^\bullet$，使得
\begin{enumerate}
\item $Y^\bullet \to X^\bullet$ 是拟同构；
\item 对所有 $i \in \mathbf{Z}$ 有 $Y^i \in \Ob(\mathcal{B})$；
\item 对所有 $i \in \mathbf{Z}$ 有 $B^i \subset Y^i$。
\end{enumerate}
为证明该断言，先利用引理的假设选取
$C^i \subset \Ker(d^i : X^i \to X^{i + 1})$，
$C^i \in \Ob(\mathcal{B})$，使其满射到 $H^i(X^\bullet)$。
令 $D^i = C^i + d^{i - 1}(B^{i - 1}) + B^i$，
可得到满足 (2) 和 (3) 的子复形 $D^\bullet$，
并且对所有 $i \in \mathbf{Z}$，
$H^i(D^\bullet) \to H^i(X^\bullet)$ 都为满射。
任取满足 $E^i \subset X^i$、$E^i \in \Ob(\mathcal{B})$ 且
$d^i(E^i) \subset D^{i + 1} + E^{i + 1}$ 的一族对象。
令 $Y^i = D^i + E^i$，便得到各项属于 $\mathcal{B}$ 的子复形，
它的上同调满射到 $X^\bullet$ 的上同调。显然，若
$d^i(E^i) = (D^{i + 1} + E^{i + 1}) \cap \Im(d^i)$，
则上同调上的映射也是单射。当 $n \gg 0$ 时可令 $E^n$ 等于 $0$。
按下降归纳法，可以对所有 $i$ 选取具有所需性质的 $E^i$。
具体地，给定 $E^{i + 1}, E^{i + 2}, \ldots$ 后，选取
$E^i \subset X^i$，使得
$d^i(E^i) = (D^{i + 1} + E^{i + 1}) \cap \Im(d^i)$。
这是可行的，因为引理的假设与如下事实共同保证了它：
$(D^{i + 1} + E^{i + 1}) \cap \Im(d^i)$ 属于 $\mathcal{B}$，
这是由于 $\mathcal{B}$ 是 $\mathcal{A}$ 的 Serre 子范畴。

\medskip\noindent
上述断言推出本引理。本质满性立即由该断言得到。
下面证明忠实性。设 $f : U^\bullet \to V^\bullet$ 是
$\mathcal{B}$ 的上有界复形之间的态射，且其在
$D(\mathcal{A})$ 中的像为零。则存在到 $\mathcal{A}$ 的某个
上有界复形中的拟同构 $s : V^\bullet \to X^\bullet$，
使得 $s \circ f$ 零伦。选取 $0$ 与 $s \circ f$ 之间的同伦
$h^i : U^i \to X^{i - 1}$。
对 $B^i = h^{i + 1}(U^{i + 1}) + s^i(V^i)$ 应用该断言。
所得映射 $s' : V^\bullet \to Y^\bullet$ 也是拟同构，且
$s' \circ f$ 零伦；这由 $h^i$ 分解通过 $Y^{i - 1}$ 显然可见。
这证明了忠实性。
% P02-E0039: frozen source says “Fully faithfulness”; the noun is “Full faithfulness”.
满忠实性也用完全相同的方法证明。
\end{proof}





\section{内射消解}
\label{derived-section-injective-resolutions}

\noindent
本节证明若干关于具有足够多内射对象的阿贝尔范畴中
内射消解存在性的引理。

\begin{definition}
\label{derived-definition-injective-resolution}
设 $\mathcal{A}$ 为阿贝尔范畴，$A \in \Ob(\mathcal{A})$。
$A$ 的一个{\it 内射消解}是复形 $I^\bullet$，
连同映射 $A \to I^0$，满足：
\begin{enumerate}
\item 对 $n < 0$ 有 $I^n = 0$。
\item 每个 $I^n$ 都是 $\mathcal{A}$ 的内射对象。
\item 映射 $A \to I^0$ 是到 $\Ker(d^0)$ 上的同构。
\item 对 $i > 0$ 有 $H^i(I^\bullet) = 0$。
\end{enumerate}
因此 $A[0] \to I^\bullet$ 是拟同构。
换言之，复形
$$
\ldots \to 0 \to A \to I^0 \to I^1 \to \ldots
$$
无环。设 $K^\bullet$ 为 $\mathcal{A}$ 中的复形。
$K^\bullet$ 的一个{\it 内射消解}是复形 $I^\bullet$，
连同复形映射 $\alpha : K^\bullet \to I^\bullet$，满足：
\begin{enumerate}
\item 对 $n \ll 0$ 有 $I^n = 0$；即 $I^\bullet$ 下有界。
\item 每个 $I^n$ 都是 $\mathcal{A}$ 的内射对象。
\item 映射 $\alpha : K^\bullet \to I^\bullet$ 是拟同构。
\end{enumerate}
\end{definition}

\noindent
换言之，内射消解 $K^\bullet \to I^\bullet$ 给出图
$$
\xymatrix{
\ldots \ar[r] & K^{n - 1} \ar[d] \ar[r] & K^n \ar[d] \ar[r] &
K^{n + 1} \ar[d] \ar[r] & \ldots \\
\ldots \ar[r] & I^{n - 1} \ar[r] & I^n \ar[r] & I^{n + 1} \ar[r] & \ldots
}
$$
它在每个次数的上同调对象上诱导同构。
$\mathcal{A}$ 的对象 $A$ 的内射消解，
与复形 $A[0]$ 的内射消解几乎是同一回事。

\begin{lemma}
\label{derived-lemma-cohomology-bounded-below}
设 $\mathcal{A}$ 为阿贝尔范畴，$K^\bullet$ 为 $\mathcal{A}$ 的复形。
\begin{enumerate}
\item 若 $K^\bullet$ 有内射消解，则当 $n \ll 0$ 时
$H^n(K^\bullet) = 0$。
\item 若对所有 $n \ll 0$ 有 $H^n(K^\bullet) = 0$，
则存在拟同构 $K^\bullet \to L^\bullet$，其中 $L^\bullet$ 下有界。
\end{enumerate}
\end{lemma}

\begin{proof}
从略。对第二个断言，取某个 $n \ll 0$，并使用
$L^\bullet = \tau_{\geq n}K^\bullet$。
截断 $\tau_{\geq n}$ 的定义见同调，章节
\ref{homology-section-truncations}。
\end{proof}

\begin{lemma}
\label{derived-lemma-injective-resolutions-exist}
设 $\mathcal{A}$ 为阿贝尔范畴。假设 $\mathcal{A}$ 有足够多内射对象。
\begin{enumerate}
\item $\mathcal{A}$ 的任意对象都有内射消解。
\item 若对所有 $n \ll 0$ 有 $H^n(K^\bullet) = 0$，
则 $K^\bullet$ 有内射消解。
\item 若 $K^\bullet$ 是满足 $n < a$ 时 $K^n = 0$ 的复形，
则存在内射消解 $\alpha : K^\bullet \to I^\bullet$，
使得 $n < a$ 时 $I^n = 0$，且每个
$\alpha^n : K^n \to I^n$ 都是单射。
\end{enumerate}
\end{lemma}

\begin{proof}
先证明 (1)。首先选取从 $A$ 到 $\mathcal{A}$ 的某个内射对象中的单射
$A \to I^0$。接着，选取到 $\mathcal{A}$ 的某个内射对象中的单射
% P02-E0041: frozen source uses I_0/A, inconsistent with the previously defined I^0.
$I_0/A \to I^1$。以 $d^0$ 表示所诱导的映射 $I^0 \to I^1$。
再选取到 $\mathcal{A}$ 的某个内射对象中的单射
$I^1/\Im(d^0) \to I^2$。以 $d^1$ 表示所诱导的映射
$I^1 \to I^2$，依此类推。
由引理 \ref{derived-lemma-cohomology-bounded-below}，(2) 由 (3) 推出。
(3) 是引理 \ref{derived-lemma-subcategory-right-resolution} 的特例。
\end{proof}

\begin{lemma}
\label{derived-lemma-acyclic-is-zero}
设 $\mathcal{A}$ 为阿贝尔范畴，$K^\bullet$ 为无环复形。
设 $I^\bullet$ 下有界并由内射对象组成。
则任意态射 $K^\bullet \to I^\bullet$ 都零伦。
\end{lemma}

\begin{proof}
设 $\alpha : K^\bullet \to I^\bullet$ 为复形态射。
注意，由于 $I^\bullet$ 下有界，当 $j \ll 0$ 时 $\alpha^j = 0$。
特别地，可以找到 $n$，使得对 $j \leq n$ 存在
$h^j :  K^j \to I^{j - 1}$，并且对 $j < n$ 有
$\alpha^j = d^{j - 1} \circ h^j + h^{j + 1} \circ d^j$。
我们将说明存在态射 $h^{n + 1} : K^{n + 1} \to I^n$，使得
$\alpha^n = d^{n - 1} \circ h^n + h^{n + 1} \circ d^n$。
注意
% P02-E0042: the frozen calculation has an alpha-index error and several h-index errors; see the ledger.
\begin{align*}
(\alpha^n - d^{n - 1} \circ h^{n - 1}) \circ d^{n - 1}
& =
\alpha^{n - 1} \circ d^{n - 1} - d^{n - 1} \circ h^{n - 1} \circ d^{n - 1} \\
& =
d^{n - 1} \circ \alpha^{n - 1} - d^{n - 1} \circ h^{n - 1} \circ d^{n - 1} \\
& =
d^{n - 1} \circ (d^{n - 2} \circ h^{n - 1} + h^n \circ d^{n - 1})
- d^{n - 1} \circ h^{n - 1} \circ d^{n - 1} \\
& = 0
\end{align*}
由于 $K^\bullet$ 无环，有
$d^{n - 1}(K^{n - 1}) = \Ker(K^n \to K^{n + 1})$。
所以可把 $\alpha^n -  d^{n - 1} \circ h^{n - 1}$ 看作
定义在 $K^{n + 1}$ 的子对象
$\Im(K^n \to K^{n + 1})$ 上、取值于 $I^n$ 的映射。
利用对象 $I^n$ 的内射性，可把它扩张成映射
$h^{n + 1} : K^{n + 1} \to I^n$。
读者可验证，在这一选择下确有
$\alpha^n = d^{n - 1} \circ h^n + h^{n + 1} \circ d^n$。

\medskip\noindent
对 $n$ 归纳可知，能够找到
$h  = (h^j)_{j \in \mathbf{Z}}$，它给出 $\alpha$ 与 $0$ 之间的同伦，
正合所需。
\end{proof}

\begin{remark}
\label{derived-remark-easier-proofs}
设 $\mathcal{A}$ 为阿贝尔范畴。
利用 $K(\mathcal{A})$ 是三角范畴这一事实，可以用引理
\ref{derived-lemma-acyclic-is-zero} 证明下面一些通常通过追图证明的引理。
具体而言，假设 $\alpha : K^\bullet \to L^\bullet$ 是复形的拟同构。
则
$$
(K^\bullet, L^\bullet, C(\alpha)^\bullet, \alpha, i, -p)
$$
是 $K(\mathcal{A})$ 中的特异三角
（引理 \ref{derived-lemma-the-same-up-to-isomorphisms}），而
$C(\alpha)^\bullet$ 是无环复形（引理 \ref{derived-lemma-acyclic}）。
再设 $I^\bullet$ 为由内射对象组成的下有界复形。则
$$
\xymatrix{
\Hom_{K(\mathcal{A})}(C(\alpha)^\bullet, I^\bullet) \ar[r] &
\Hom_{K(\mathcal{A})}(L^\bullet, I^\bullet) \ar[r] &
\Hom_{K(\mathcal{A})}(K^\bullet, I^\bullet) \ar[lld] \\
\Hom_{K(\mathcal{A})}(C(\alpha)^\bullet[-1], I^\bullet)
}
$$
是阿贝尔群的正合序列；见引理 \ref{derived-lemma-representable-homological}。
此时，引理 \ref{derived-lemma-acyclic-is-zero} 保证外侧两个群为零，故
$\Hom_{K(\mathcal{A})}(L^\bullet, I^\bullet) =
\Hom_{K(\mathcal{A})}(K^\bullet, I^\bullet)$。
\end{remark}

\begin{lemma}
\label{derived-lemma-morphisms-lift}
设 $\mathcal{A}$ 为阿贝尔范畴。考虑实线图
$$
\xymatrix{
K^\bullet \ar[r]_\alpha \ar[d]_\gamma & L^\bullet \ar@{-->}[dl]^\beta \\
I^\bullet
}
$$
其中 $I^\bullet$ 下有界且由内射对象组成，而 $\alpha$ 是拟同构。
\begin{enumerate}
\item 存在复形映射 $\beta$，使图交换至同伦。
\item 若 $\alpha$ 在每个次数上都是单射，则可找到使图交换的 $\beta$。
\end{enumerate}
\end{lemma}

\begin{proof}
(1) 的“正确”证明见注记 \ref{derived-remark-easier-proofs}。
这里也给出一个直接证明。

\medskip\noindent
先说明 (2) 推出 (1)。设
% P02-E0043: frozen source writes K rather than K^\bullet in the complex map.
$\tilde \alpha : K \to \tilde L^\bullet$、$\pi$、$s$
如引理 \ref{derived-lemma-make-injective} 中所述。
由于 $\tilde \alpha$ 为单射，由 (2) 存在态射
$\tilde \beta : \tilde L^\bullet \to I^\bullet$，使得
$\gamma = \tilde \beta \circ \tilde \alpha$。令
$\beta = \tilde \beta \circ s$。则
$$
\beta \circ \alpha
=
\tilde \beta \circ s \circ \pi \circ \tilde \alpha
\sim
\tilde \beta \circ \tilde \alpha
=
\gamma
$$
正合所需。

\medskip\noindent
假设 $\alpha : K^\bullet \to L^\bullet$ 为单射。
设已在所有次数 $\leq n - 1$ 上定义了与微分相容的 $\beta$，
且对所有 $j \leq n - 1$ 有
$\gamma^j = \beta^j \circ \alpha^j$。
考虑交换的实线图
$$
\xymatrix{
K^{n - 1} \ar[r] \ar@/_2pc/[dd]_\gamma \ar[d]^\alpha &
K^n \ar@/^2pc/[dd]^\gamma \ar[d]^\alpha \\
L^{n - 1} \ar[r] \ar[d]^\beta &
L^n \ar@{-->}[d] \\
I^{n - 1} \ar[r] &
I^n
}
$$
因此虚线箭头在子对象 $\alpha(K^n)$ 和
$d^{n - 1}(L^{n - 1})$ 上已有规定。
此外，这两个箭头在 $\alpha(d^{n - 1}(K^{n - 1}))$ 上相同。
所以，若
\begin{equation}
\label{derived-equation-qis}
\alpha(d^{n - 1}(K^{n - 1}))
=
\alpha(K^n) \cap d^{n - 1}(L^{n - 1})
\end{equation}
则这些态射可粘合成态射
$\alpha(K^n) + d^{n - 1}(L^{n - 1}) \to I^n$；
再利用 $I^n$ 的内射性，可把它扩张为从整个 $L^n$ 到 $I^n$ 的态射。
此后通过归纳，可同时对所有 $n$ 得到态射 $\beta$
（注意，因为 $I^\bullet$ 下有界，当 $n \ll 0$ 时可令
$\beta^n = 0$，从而启动归纳）。

\medskip\noindent
还须证明等式 (\ref{derived-equation-qis})。
建议读者用适当的追图自行论证；不过这里仍给出我们的论证。
包含关系
$\alpha(d^{n - 1}(K^{n - 1})) \subset \alpha(K^n) \cap d^{n - 1}(L^{n - 1})$
显然成立。取 $\mathcal{A}$ 的对象 $T$ 以及态射
$x : T \to L^n$，其像包含在子对象
$\alpha(K^n) \cap d^{n - 1}(L^{n - 1})$ 中。
由于 $\alpha$ 为单射，有 $x = \alpha \circ x'$，
其中 $x' : T \to K^n$。又因 $x$ 位于
$d^{n - 1}(L^{n - 1})$ 中，有 $d^n \circ x = 0$。
再次利用 $\alpha$ 的单射性，得到 $d^n \circ x' = 0$。
因此 $x'$ 给出态射 $[x'] : T \to H^n(K^\bullet)$。
另一方面，由 $x$ 诱导的相应映射
$[x] : T \to H^n(L^\bullet)$ 依假设为零。
由于 $\alpha$ 是拟同构，必有 $[x'] = 0$。
这恰好意味着 $x'$ 的像包含在
$d^{n - 1}(K^{n - 1})$ 中，证明完成。
\end{proof}

\begin{lemma}
\label{derived-lemma-morphisms-equal-up-to-homotopy}
设 $\mathcal{A}$ 为阿贝尔范畴。考虑实线图
$$
\xymatrix{
K^\bullet \ar[r]_\alpha \ar[d]_\gamma & L^\bullet \ar@{-->}[dl]^{\beta_i} \\
I^\bullet
}
$$
其中 $I^\bullet$ 下有界且由内射对象组成，而 $\alpha$ 是拟同构。
任意两个使图交换至同伦的态射 $\beta_1, \beta_2$ 都彼此同伦。
\end{lemma}

\begin{proof}
这由注记 \ref{derived-remark-easier-proofs} 推出。
这里也给出一个直接论证。

\medskip\noindent
% P02-E0043: this second frozen locus also writes K rather than K^\bullet.
设 $\tilde \alpha : K \to \tilde L^\bullet$、$\pi$、$s$
如引理 \ref{derived-lemma-make-injective} 中所述。
若能证明 $\beta_1 \circ\pi$ 与 $\beta_2 \circ \pi$ 同伦，
则由 $\pi \circ s$ 是恒等映射可推出
$\beta_1 \sim \beta_2$。因此可以假设，对所有 $n$，
$\alpha^n : K^n \to L^n$ 都是直和项的包含映射。
于是得到逐项分裂的复形短正合序列
$$
0 \to K^\bullet \to L^\bullet \to M^\bullet \to 0
$$
且 $M^\bullet$ 无环。选取分裂 $L^n = K^n \oplus M^n$，
于是有 $\beta_i^n : K^n \oplus M^n \to I^n$ 和
$\gamma^n : K^n \to I^n$。此时对 $\beta_i$ 的条件是：
存在态射 $h_i^n : K^n \to I^{n - 1}$，使得
$$
\gamma^n - \beta_i^n|_{K^n} = d \circ h_i^n + h_i^{n + 1} \circ d
$$
因此
$$
\beta_1^n|_{K^n} - \beta_2^n|_{K^n}
=
d \circ (h_1^n - h_2^n) + (h_1^{n + 1} - h_2^{n + 1}) \circ d
$$
考虑映射 $h^n : K^n \oplus M^n \to I^{n - 1}$，
它在第一个直和项上等于 $h_1^n - h_2^n$，在第二个直和项上为零。
于是可见
% P02-E0044: frozen formula has (d h^n + h^{n+1}) d rather than d h^n + h^{n+1} d.
$$
\beta_1^n - \beta_2^n
-
(d \circ h^n + h^{n + 1}) \circ d
$$
是复形态射 $L^\bullet \to I^\bullet$，
它在子复形 $K^\bullet$ 上恒为零。因此它分解为
$L^\bullet \to M^\bullet \to I^\bullet$。
所以本引理的结论由引理 \ref{derived-lemma-acyclic-is-zero} 推出。
\end{proof}

\begin{lemma}
\label{derived-lemma-morphisms-into-injective-complex}
设 $\mathcal{A}$ 为阿贝尔范畴。
设 $I^\bullet$ 为由内射对象组成的下有界复形，
$L^\bullet \in K(\mathcal{A})$。则
$$
\Mor_{K(\mathcal{A})}(L^\bullet, I^\bullet)
=
\Mor_{D(\mathcal{A})}(L^\bullet, I^\bullet).
$$
\end{lemma}

\begin{proof}
设 $a$ 为右边的元素。可以把它表示为
$a = \gamma\alpha^{-1}$，其中
$\alpha : K^\bullet \to L^\bullet$ 是拟同构，
$\gamma : K^\bullet \to I^\bullet$ 是复形映射。
由引理 \ref{derived-lemma-morphisms-lift}，可找到态射
$\beta : L^\bullet \to I^\bullet$，使得
$\beta \circ \alpha$ 与 $\gamma$ 同伦。这证明该映射为满射。
设 $b$ 为左边的元素，且其在右边的像为零。
那么 $b$ 是某个态射 $\beta : L^\bullet \to I^\bullet$ 的同伦类，
并存在拟同构 $\alpha : K^\bullet \to L^\bullet$，使得
$\beta  \circ \alpha$ 零伦。由引理
\ref{derived-lemma-morphisms-equal-up-to-homotopy}，
$\beta$ 本身也零伦；即 $b = 0$。
\end{proof}

\begin{lemma}
\label{derived-lemma-injective-resolution-ses}
设 $\mathcal{A}$ 为阿贝尔范畴。假设 $\mathcal{A}$ 有足够多内射对象。
对 $\text{Comp}^{+}(\mathcal{A})$ 的任意短正合序列
$0 \to A^\bullet \to B^\bullet \to C^\bullet \to 0$，
存在 $\text{Comp}^{+}(\mathcal{A})$ 中的交换图
$$
\xymatrix{
0 \ar[r] &
A^\bullet \ar[r] \ar[d] &
B^\bullet \ar[r] \ar[d] &
C^\bullet \ar[r] \ar[d] &
0 \\
0 \ar[r] &
I_1^\bullet \ar[r] &
I_2^\bullet \ar[r] &
I_3^\bullet \ar[r] &
0
}
$$
其中竖直箭头是内射消解，各行是复形的短正合序列。此外：
\begin{enumerate}
\item 给定任意内射消解 $A^\bullet \to I^\bullet$，
可以假设 $I_1^\bullet = I^\bullet$；
\item 若对 $n < 0$ 有 $A^n = B^n = C^n = 0$，
则可以假设对 $n < 0$ 有 $I_j^n = 0$；
\item 若还对 $n < 0$ 有 $I^n = 0$，则可合并 (1) 和 (2)；
% P02-E0045: frozen source contains the unfinished placeholder “add more here.”
\item 在此添加更多内容。
\end{enumerate}
\end{lemma}

\begin{proof}
步骤 1。选取内射消解 $A^\bullet \to I^\bullet$
（见引理 \ref{derived-lemma-injective-resolutions-exist}），或使用给定的消解。
回忆 $\text{Comp}^{+}(\mathcal{A})$ 是阿贝尔范畴；见同调，引理
\ref{homology-lemma-cat-cochain-abelian}。
因此可沿映射 $A^\bullet \to I^\bullet$ 作推出，得到
$$
\xymatrix{
0 \ar[r] &
A^\bullet \ar[r] \ar[d] &
B^\bullet \ar[r] \ar[d] &
C^\bullet \ar[r] \ar[d] &
0 \\
0 \ar[r] &
I^\bullet \ar[r] &
E^\bullet \ar[r] &
C^\bullet \ar[r] &
0
}
$$
由 $5$-引理和同调，引理
\ref{homology-lemma-long-exact-sequence-cochain} 的最后一个断言，
映射 $B^\bullet \to E^\bullet$ 是拟同构。
注意，下面的短正合序列逐项分裂；见同调，引理
\ref{homology-lemma-characterize-injectives}。
所以只须在
$0 \to A^\bullet \to B^\bullet \to C^\bullet \to 0$
逐项分裂时证明本引理。

\medskip\noindent
步骤 2。选取分裂。换言之，写成 $B^n = A^n \oplus C^n$。
以 $\delta : C^\bullet \to A^\bullet[1]$ 表示同调，引理
\ref{homology-lemma-ses-termwise-split-cochain} 中的态射。
选取内射消解 $f_1 : A^\bullet \to I_1^\bullet$ 和
$f_3 : C^\bullet \to I_3^\bullet$。
（若 $A^\bullet$ 是内射对象的复形，则使用
$I_1^\bullet = A^\bullet$。若对 $n < 0$ 有
$A^n = C^n = 0$，则由引理
\ref{derived-lemma-injective-resolutions-exist}，可选取消解使得
对 $n < 0$ 有 $I_1^n = I_3^n = 0$。）
可以假设 $f_3$ 在每个次数上都是单射。
由引理 \ref{derived-lemma-morphisms-lift}，可找到态射
$\delta' : I_3^\bullet \to I_1^\bullet[1]$，使得
$\delta' \circ f_3 = f_1[1] \circ \delta$
（复形态射的等式）。令 $I_2^n = I_1^n \oplus I_3^n$，定义
$$
d_{I_2}^n =
\left(
\begin{matrix}
d_{I_1}^n & (\delta')^n \\
0 & d_{I_3}^n
\end{matrix}
\right)
$$
并把映射 $B^n \to I_2^n$ 定义为映射
$A^n \to I_1^n$ 与 $C^n \to I_3^n$ 之和。
一切均显然。
\end{proof}








\section{投射消解}
\label{derived-section-projective-resolutions}

\noindent
本节与章节 \ref{derived-section-injective-resolutions} 对偶。
我们给出定义并陈述结果，但不再证明这些引理。

\begin{definition}
\label{derived-definition-projective-resolution}
设 $\mathcal{A}$ 为阿贝尔范畴，$A \in \Ob(\mathcal{A})$。
$A$ 的一个{\it 投射消解}是复形 $P^\bullet$，
连同映射 $P^0 \to A$，满足：
\begin{enumerate}
\item 对 $n > 0$ 有 $P^n = 0$。
\item 每个 $P^n$ 都是 $\mathcal{A}$ 的投射对象。
\item 映射 $P^0 \to A$ 诱导同构 $\Coker(d^{-1}) \to A$。
\item 对 $i < 0$ 有 $H^i(P^\bullet) = 0$。
\end{enumerate}
因此 $P^\bullet \to A[0]$ 是拟同构。
换言之，复形
$$
\ldots \to P^{-1} \to P^0 \to A \to 0 \to \ldots
$$
无环。设 $K^\bullet$ 为 $\mathcal{A}$ 中的复形。
$K^\bullet$ 的一个{\it 投射消解}是复形 $P^\bullet$，
连同复形映射 $\alpha : P^\bullet \to K^\bullet$，满足：
\begin{enumerate}
\item 对 $n \gg 0$ 有 $P^n = 0$；即 $P^\bullet$ 上有界。
\item 每个 $P^n$ 都是 $\mathcal{A}$ 的投射对象。
\item 映射 $\alpha : P^\bullet \to K^\bullet$ 是拟同构。
\end{enumerate}
\end{definition}

\begin{lemma}
\label{derived-lemma-cohomology-bounded-above}
设 $\mathcal{A}$ 为阿贝尔范畴，$K^\bullet$ 为 $\mathcal{A}$ 的复形。
\begin{enumerate}
\item 若 $K^\bullet$ 有投射消解，则当 $n \gg 0$ 时
$H^n(K^\bullet) = 0$。
\item 若当 $n \gg 0$ 时 $H^n(K^\bullet) = 0$，
则存在拟同构 $L^\bullet \to K^\bullet$，其中 $L^\bullet$ 上有界。
\end{enumerate}
\end{lemma}

\begin{proof}
与引理 \ref{derived-lemma-cohomology-bounded-below} 对偶。
\end{proof}

\begin{lemma}
\label{derived-lemma-projective-resolutions-exist}
设 $\mathcal{A}$ 为阿贝尔范畴。假设 $\mathcal{A}$ 有足够多投射对象。
\begin{enumerate}
\item $\mathcal{A}$ 的任意对象都有投射消解。
\item 若对所有 $n \gg 0$ 有 $H^n(K^\bullet) = 0$，
则 $K^\bullet$ 有投射消解。
\item 若 $K^\bullet$ 是满足 $n > a$ 时 $K^n = 0$ 的复形，
则存在投射消解 $\alpha : P^\bullet \to K^\bullet$，
使得 $n > a$ 时 $P^n = 0$，且每个
$\alpha^n : P^n \to K^n$ 都为满射。
\end{enumerate}
\end{lemma}

\begin{proof}
与引理 \ref{derived-lemma-injective-resolutions-exist} 对偶。
\end{proof}

\begin{lemma}
\label{derived-lemma-projective-into-acyclic-is-zero}
设 $\mathcal{A}$ 为阿贝尔范畴，$K^\bullet$ 为无环复形。
设 $P^\bullet$ 上有界并由投射对象组成。
则任意态射 $P^\bullet \to K^\bullet$ 都零伦。
\end{lemma}

\begin{proof}
与引理 \ref{derived-lemma-acyclic-is-zero} 对偶。
\end{proof}

\begin{remark}
\label{derived-remark-easier-projective}
设 $\mathcal{A}$ 为阿贝尔范畴。
假设 $\alpha : K^\bullet \to L^\bullet$ 是复形的拟同构，
$P^\bullet$ 是投射对象的上有界复形。则
$$
\Hom_{K(\mathcal{A})}(P^\bullet, K^\bullet)
\longrightarrow
\Hom_{K(\mathcal{A})}(P^\bullet, L^\bullet)
$$
是同构。这与注记 \ref{derived-remark-easier-proofs} 对偶。
\end{remark}

\begin{lemma}
\label{derived-lemma-morphisms-lift-projective}
设 $\mathcal{A}$ 为阿贝尔范畴。考虑实线图
$$
\xymatrix{
K^\bullet & L^\bullet \ar[l]^\alpha \\
P^\bullet \ar[u] \ar@{-->}[ru]_\beta
}
$$
其中 $P^\bullet$ 上有界且由投射对象组成，而 $\alpha$ 是拟同构。
\begin{enumerate}
\item 存在复形映射 $\beta$，使图交换至同伦。
\item 若 $\alpha$ 在每个次数上都是满射，则可找到使图交换的 $\beta$。
\end{enumerate}
\end{lemma}

\begin{proof}
与引理 \ref{derived-lemma-morphisms-lift} 对偶。
\end{proof}

\begin{lemma}
\label{derived-lemma-morphisms-equal-up-to-homotopy-projective}
设 $\mathcal{A}$ 为阿贝尔范畴。考虑实线图
$$
\xymatrix{
K^\bullet & L^\bullet \ar[l]^\alpha \\
P^\bullet \ar[u] \ar@{-->}[ru]_{\beta_i}
}
$$
其中 $P^\bullet$ 上有界且由投射对象组成，而 $\alpha$ 是拟同构。
任意两个使图交换至同伦的态射 $\beta_1, \beta_2$ 都彼此同伦。
\end{lemma}

\begin{proof}
与引理 \ref{derived-lemma-morphisms-equal-up-to-homotopy} 对偶。
\end{proof}

\begin{lemma}
\label{derived-lemma-morphisms-from-projective-complex}
设 $\mathcal{A}$ 为阿贝尔范畴。
设 $P^\bullet$ 为由投射对象组成的上有界复形，
$L^\bullet \in K(\mathcal{A})$。则
$$
\Mor_{K(\mathcal{A})}(P^\bullet, L^\bullet)
=
\Mor_{D(\mathcal{A})}(P^\bullet, L^\bullet).
$$
\end{lemma}

\begin{proof}
与引理 \ref{derived-lemma-morphisms-into-injective-complex} 对偶。
\end{proof}

\begin{lemma}
\label{derived-lemma-projective-resolution-ses}
设 $\mathcal{A}$ 为阿贝尔范畴。假设 $\mathcal{A}$ 有足够多投射对象。
% P02-E0046: the projective dual is bounded above, but frozen source names Comp^+ rather than Comp^- here and below.
对 $\text{Comp}^{+}(\mathcal{A})$ 的任意短正合序列
$0 \to A^\bullet \to B^\bullet \to C^\bullet \to 0$，
存在 $\text{Comp}^{+}(\mathcal{A})$ 中的交换图
$$
\xymatrix{
0 \ar[r] &
P_1^\bullet \ar[r] \ar[d] &
P_2^\bullet \ar[r] \ar[d] &
P_3^\bullet \ar[r] \ar[d] &
0 \\
0 \ar[r] &
A^\bullet \ar[r] &
B^\bullet \ar[r] &
C^\bullet \ar[r] &
0
}
$$
其中竖直箭头为投射消解，各行是复形的短正合序列。
事实上，给定任意投射消解 $P^\bullet \to C^\bullet$，
可以假设 $P_3^\bullet = P^\bullet$。
\end{lemma}

\begin{proof}
与引理 \ref{derived-lemma-injective-resolution-ses} 对偶。
\end{proof}

\begin{lemma}
\label{derived-lemma-precise-vanishing}
设 $\mathcal{A}$ 为阿贝尔范畴，$P^\bullet$、$K^\bullet$ 为复形，
$n \in \mathbf{Z}$。假设：
\begin{enumerate}
\item $P^\bullet$ 是由投射对象组成的有界复形；
\item 对 $i < n$ 有 $P^i = 0$；
\item 对 $i \geq n$ 有 $H^i(K^\bullet) = 0$。
\end{enumerate}
则
$\Hom_{K(\mathcal{A})}(P^\bullet, K^\bullet) =
\Hom_{D(\mathcal{A})}(P^\bullet, K^\bullet) = 0$。
\end{lemma}

\begin{proof}
第一个等号由引理 \ref{derived-lemma-morphisms-from-projective-complex} 推出。
注意，由注记 \ref{derived-remark-truncation-distinguished-triangle}，有特异三角
$$
(\tau_{\leq n - 1}K^\bullet, K^\bullet, \tau_{\geq n}K^\bullet, f, g, h)
$$
因此，由引理 \ref{derived-lemma-representable-homological}，只须证明
$\Hom_{K(\mathcal{A})}(P^\bullet, \tau_{\leq n - 1}K^\bullet) = 0$ 和
$\Hom_{K(\mathcal{A})}(P^\bullet, \tau_{\geq n} K^\bullet) = 0$。
第一个消失结论显然，第二个就是引理
\ref{derived-lemma-projective-into-acyclic-is-zero}。
\end{proof}

\begin{lemma}
\label{derived-lemma-lift-map}
设 $\mathcal{A}$ 为阿贝尔范畴。设
$\beta : P^\bullet \to L^\bullet$ 和
$\alpha : E^\bullet \to L^\bullet$ 为复形映射，
$n \in \mathbf{Z}$。假设：
\begin{enumerate}
\item $P^\bullet$ 是投射对象的有界复形，且对 $i < n$ 有 $P^i = 0$；
\item 当 $i > n$ 时 $H^i(\alpha)$ 为同构，
当 $i = n$ 时为满射。
\end{enumerate}
则存在复形映射 $\gamma : P^\bullet \to E^\bullet$，
使得 $\alpha \circ \gamma$ 与 $\beta$ 同伦。
\end{lemma}

\begin{proof}
考虑带有映射 $i : L^\bullet \to C^\bullet$ 的锥
$C^\bullet = C(\alpha)^\bullet$。
由引理 \ref{derived-lemma-precise-vanishing}，$i \circ \beta$ 为零。
所以由引理 \ref{derived-lemma-representable-homological}，
可把 $\beta$ 提升到 $E^\bullet$。
\end{proof}













\section{右导出函子与内射消解}
\label{derived-section-right-derived-functor}

\noindent
现在可以利用上述材料，定义从具有足够多内射对象的阿贝尔范畴
到一般阿贝尔范畴的加性函子的右导出函子。

\begin{lemma}
\label{derived-lemma-injective-acyclic}
设 $\mathcal{A}$ 为阿贝尔范畴，
$I \in \Ob(\mathcal{A})$ 为内射对象，
$I^\bullet$ 为 $\mathcal{A}$ 中内射对象的下有界复形。
\begin{enumerate}
\item 对任意到任意三角范畴 $\mathcal{D}$ 的正合函子
$F : K^{+}(\mathcal{A}) \to \mathcal{D}$，
$I^\bullet$ 都相对于 $\text{Qis}^{+}(\mathcal{A})$ 计算 $RF$。
\item 对任意到任意阿贝尔范畴 $\mathcal{B}$ 的加性函子
$F : \mathcal{A} \to \mathcal{B}$，$I$ 都右无环。
\end{enumerate}
\end{lemma}

\begin{proof}
% P02-E0047: frozen source says “a direct consequences”.
(2) 是 (1) 和定义 \ref{derived-definition-derived-functor} 的直接推论。
为证明 (1)，设 $\alpha : I^\bullet \to K^\bullet$ 为到某复形中的拟同构。
由引理 \ref{derived-lemma-morphisms-lift}，$\alpha$ 有左逆。
所以范畴 $I^\bullet/\text{Qis}^{+}(\mathcal{A})$
本质常值，值为 $\text{id} : I^\bullet \to I^\bullet$。
因而 ind-对象
$$
I^\bullet/\text{Qis}^{+}(\mathcal{A}) \longrightarrow \mathcal{D}, \quad
(I^\bullet \to K^\bullet) \longmapsto F(K^\bullet)
$$
也本质常值，值为 $F(I^\bullet)$。这证明了 (1)；见定义
\ref{derived-definition-right-derived-functor-defined} 和
\ref{derived-definition-computes}。
\end{proof}

\begin{lemma}
\label{derived-lemma-enough-injectives-right-derived}
设 $\mathcal{A}$ 为具有足够多内射对象的阿贝尔范畴。
\begin{enumerate}
\item 对任意到三角范畴 $\mathcal{D}$ 的正合函子
$F : K^{+}(\mathcal{A}) \to \mathcal{D}$，右导出函子
$$
RF : D^{+}(\mathcal{A}) \longrightarrow \mathcal{D}
$$
处处有定义。
\item 对任意到阿贝尔范畴 $\mathcal{B}$ 的加性函子
$F : \mathcal{A} \to \mathcal{B}$，右导出函子
$$
RF : D^{+}(\mathcal{A}) \longrightarrow D^{+}(\mathcal{B})
$$
处处有定义。
\end{enumerate}
\end{lemma}

\begin{proof}
对第二个断言，结合引理 \ref{derived-lemma-injective-acyclic} 和
命题 \ref{derived-proposition-enough-acyclics}。
为得到第一个断言，结合引理 \ref{derived-lemma-injective-resolutions-exist}、
引理 \ref{derived-lemma-injective-acyclic}、
引理 \ref{derived-lemma-existence-computes} 和等式
(\ref{derived-equation-everywhere})。
\end{proof}

\begin{lemma}
\label{derived-lemma-right-derived-properties}
设 $\mathcal{A}$ 为具有足够多内射对象的阿贝尔范畴，
$F : \mathcal{A} \to \mathcal{B}$ 为加性函子。
\begin{enumerate}
\item 函子 $RF$ 是正合函子
$D^{+}(\mathcal{A}) \to D^{+}(\mathcal{B})$。
\item 函子 $RF$ 诱导正合函子
$K^{+}(\mathcal{A}) \to D^{+}(\mathcal{B})$。
\item 函子 $RF$ 诱导 $\delta$-函子
$\text{Comp}^{+}(\mathcal{A}) \to D^{+}(\mathcal{B})$。
\item 函子 $RF$ 诱导 $\delta$-函子
$\mathcal{A} \to D^{+}(\mathcal{B})$。
\end{enumerate}
\end{lemma}

\begin{proof}
本引理只是回顾目前得到的一些结果。注意，由引理
\ref{derived-lemma-enough-injectives-right-derived}，$RF$ 处处有定义。
下面给出一些参考：
\begin{enumerate}
\item 导出函子正合：这归结为引理 \ref{derived-lemma-2-out-of-3-defined}。
\item 这是因为 $K^{+}(\mathcal{A}) \to D^{+}(\mathcal{A})$
正合，而正合函子的复合正合。
\item 这是因为 $\text{Comp}^{+}(\mathcal{A}) \to D^{+}(\mathcal{A})$
是 $\delta$-函子；见引理 \ref{derived-lemma-derived-canonical-delta-functor}。
\item 这是因为 $\mathcal{A} \to \text{Comp}^{+}(\mathcal{A})$
正合，而以正合函子预复合 $\delta$-函子仍得到 $\delta$-函子。
\end{enumerate}
\end{proof}

\begin{lemma}
\label{derived-lemma-higher-derived-functors}
设 $\mathcal{A}$ 为具有足够多内射对象的阿贝尔范畴，
$F : \mathcal{A} \to \mathcal{B}$ 为左正合函子。
\begin{enumerate}
\item 对 $\text{Comp}^{+}(\mathcal{A})$ 中复形的任意短正合序列
$0 \to A^\bullet \to B^\bullet \to C^\bullet \to 0$，
有相伴的长正合序列
$$
\ldots \to
H^i(RF(A^\bullet)) \to
H^i(RF(B^\bullet)) \to
H^i(RF(C^\bullet)) \to
H^{i + 1}(RF(A^\bullet)) \to \ldots
$$
\item 函子 $R^iF : \mathcal{A} \to \mathcal{B}$ 在 $i < 0$ 时为零。
并且 $R^0F = F : \mathcal{A} \to \mathcal{B}$。
\item 若 $i > 0$ 且 $I$ 内射，则 $R^iF(I) = 0$。
\item 序列 $(R^iF, \delta)$ 构成从 $\mathcal{A}$ 到 $\mathcal{B}$ 的
泛 $\delta$-函子（见同调，定义
\ref{homology-definition-universal-delta-functor}）。
\end{enumerate}
\end{lemma}

\begin{proof}
本引理只是回顾目前得到的一些结果。注意，由引理
\ref{derived-lemma-enough-injectives-right-derived}，$RF$ 处处有定义。
下面给出一些参考：
\begin{enumerate}
\item 这由引理 \ref{derived-lemma-right-derived-properties} 的 (3) 与
$D^{+}(\mathcal{B})$ 的长正合上同调序列
(\ref{derived-equation-long-exact-cohomology-sequence-D}) 结合得到。
\item 这就是引理 \ref{derived-lemma-left-exact-higher-derived}。
\item 这就是内射对象无环这一事实。
\item 这就是引理 \ref{derived-lemma-right-derived-delta-functor}。
\end{enumerate}
\end{proof}





\section{Cartan--Eilenberg 消解}
\label{derived-section-cartan-eilenberg}

\noindent
本节还可以扩充。这里的材料可以推广并应用于更多情形。
消解不必使用内射对象，而且在某些情形下该方法也适用于无界情形。

\begin{definition}
\label{derived-definition-cartan-eilenberg}
设 $\mathcal{A}$ 为阿贝尔范畴，$K^\bullet$ 为下有界复形。
$K^\bullet$ 的一个{\it Cartan--Eilenberg 消解}由双复形
$I^{\bullet, \bullet}$ 和复形态射
$\epsilon : K^\bullet \to I^{\bullet, 0}$ 给出，满足以下性质：
\begin{enumerate}
% P02-E0048: frozen source has the article error “There exists a i”.
\item 存在 $i \ll 0$，使得对所有 $p < i$ 和所有 $q$，
$I^{p, q} = 0$。
\item 若 $q < 0$，则 $I^{p, q} = 0$。
\item 复形 $I^{p, \bullet}$ 是 $K^p$ 的内射消解。
\item 复形 $\Ker(d_1^{p, \bullet})$ 是 $\Ker(d_K^p)$ 的内射消解。
\item 复形 $\Im(d_1^{p, \bullet})$ 是 $\Im(d_K^p)$ 的内射消解。
\item 复形 $H^p_I(I^{\bullet, \bullet})$ 是
$H^p(K^\bullet)$ 的内射消解。
\end{enumerate}
\end{definition}

\begin{lemma}
\label{derived-lemma-cartan-eilenberg}
设 $\mathcal{A}$ 为具有足够多内射对象的阿贝尔范畴，
$K^\bullet$ 为下有界复形。则 $K^\bullet$ 有 Cartan--Eilenberg 消解。
\end{lemma}

\begin{proof}
假设 $p < n$ 时 $K^p = 0$。如下把 $K^\bullet$ 分解成短正合序列：
令 $Z^p = \Ker(d^p)$、$B^p = \Im(d^{p - 1})$、
$H^p = Z^p/B^p$，并考虑
$$
\begin{matrix}
0 \to Z^n \to K^n \to B^{n + 1} \to 0 \\
0 \to B^{n + 1} \to Z^{n + 1} \to H^{n + 1} \to 0 \\
0 \to Z^{n + 1} \to K^{n + 1} \to B^{n + 2} \to 0 \\
0 \to B^{n + 2} \to Z^{n + 2} \to H^{n + 2} \to 0 \\
\ldots
\end{matrix}
$$
对 $p < n$ 令 $I^{p, q} = 0$。按如下方式归纳选取内射消解：
\begin{enumerate}
\item 选取内射消解 $Z^n \to J_Z^{n, \bullet}$，
使得 $m < 0$ 时 $J_Z^{n, m} = 0$；见引理
\ref{derived-lemma-injective-resolutions-exist}。
\item 利用引理 \ref{derived-lemma-injective-resolution-ses}，
选取内射消解 $K^n \to I^{n, \bullet}$、
$B^{n + 1} \to J_B^{n + 1, \bullet}$，使得
$m < 0$ 时 $I^{n, m} = 0$，且 $m < 0$ 时 $J_B^{n + 1, m} = 0$；
并选取复形的正合序列
$0 \to J_Z^{n, \bullet} \to I^{n, \bullet} \to J_B^{n + 1, \bullet} \to 0$，
使之与短正合序列
$0 \to Z^n \to K^n \to B^{n + 1} \to 0$ 相容。
\item 利用引理 \ref{derived-lemma-injective-resolution-ses}，
选取内射消解 $Z^{n + 1} \to J_Z^{n + 1, \bullet}$、
$H^{n + 1} \to J_H^{n + 1, \bullet}$，使得
$m < 0$ 时 $J_Z^{n + 1, m} = 0$，且 $m < 0$ 时
$J_H^{n + 1, m} = 0$；并选取复形的正合序列
$0 \to J_B^{n + 1, \bullet} \to J_Z^{n + 1, \bullet}
\to J_H^{n + 1, \bullet} \to 0$，
使之与短正合序列
$0 \to B^{n + 1} \to Z^{n + 1} \to H^{n + 1} \to 0$ 相容。
\item 依此类推。
\end{enumerate}
把映射 $d_1^\bullet : I^{p, \bullet} \to I^{p + 1, \bullet}$
取为复合
$I^{p, \bullet} \to J_B^{p + 1, \bullet} \to
J_Z^{p + 1, \bullet} \to I^{p + 1, \bullet}$，则一切均显然。
\end{proof}

\begin{lemma}
\label{derived-lemma-two-ss-complex-functor}
设 $F : \mathcal{A} \to \mathcal{B}$ 为阿贝尔范畴之间的左正合函子。
设 $K^\bullet$ 为 $\mathcal{A}$ 的下有界复形，
$I^{\bullet, \bullet}$ 为 $K^\bullet$ 的 Cartan--Eilenberg 消解。
与双复形 $F(I^{\bullet, \bullet})$ 相伴的谱序列
$({}'E_r, {}'d_r)_{r \geq 0}$ 和 $({}''E_r, {}''d_r)_{r \geq 0}$
满足关系
$$
{}'E_1^{p, q} = R^qF(K^p)
\quad
\text{且}
\quad
{}''E_2^{p, q} = R^pF(H^q(K^\bullet))
$$
此外，这些谱序列有界，收敛到 $H^*(RF(K^\bullet))$，
且在 $H^n(RF(K^\bullet))$ 上相伴诱导的过滤是有限的。
\end{lemma}

\begin{proof}
下面这些注记将在不另加说明的情况下使用：
\begin{enumerate}
\item 由于 $I^{p, \bullet}$ 是 $K^p$ 的内射消解，
可知 $RF$ 在 $K^p[0]$ 处有定义，值为 $F(I^{p, \bullet})$。
\item 由于 $H^p_I(I^{\bullet, \bullet})$ 是
$H^p(K^\bullet)$ 的内射消解，导出函子 $RF$ 在
$H^p(K^\bullet)[0]$ 处有定义，值为
$F(H^p_I(I^{\bullet, \bullet}))$。
\item 由同调，引理 \ref{homology-lemma-double-complex-gives-resolution}，
全复形 $\text{Tot}(I^{\bullet, \bullet})$ 是 $K^\bullet$ 的内射消解。
所以 $RF$ 在 $K^\bullet$ 处有定义，值为
$F(\text{Tot}(I^{\bullet, \bullet}))$。
\end{enumerate}
考虑与双复形 $L^{\bullet, \bullet} = F(I^{\bullet, \bullet})$ 相伴的
两个谱序列；见同调，引理 \ref{homology-lemma-ss-double-complex}。
二者都有界，收敛到 $H^*(\text{Tot}(L^{\bullet, \bullet}))$，
并在 $H^n(\text{Tot}(L^{\bullet, \bullet}))$ 上诱导有限过滤；见
同调，引理 \ref{homology-lemma-first-quadrant-ss}。
由于
$\text{Tot}(L^{\bullet, \bullet}) =
\text{Tot}(F(I^{\bullet, \bullet})) =
F(\text{Tot}(I^{\bullet, \bullet}))$ 计算 $RF(K^\bullet)$，
便得到本引理的最后一个断言。

\medskip\noindent
计算第一个谱序列。由于
${}'E_1^{p, q} = H^q(L^{p, \bullet})$，换言之，
$$
{}'E_1^{p, q} = H^q(F(I^{p, \bullet})) = R^qF(K^p)
$$
正合所需。留意以后将用到：映射
${}'d_1^{p, q} : {}'E_1^{p, q} \to {}'E_1^{p + 1, q}$
就是由 $K^p \to K^{p + 1}$ 以及 $R^qF$ 是函子这一事实所诱导的映射
$R^qF(K^p) \to R^qF(K^{p + 1})$。

\medskip\noindent
计算第二个谱序列。有
${}''E_1^{p, q} = H^q(L^{\bullet, p}) = H^q(F(I^{\bullet, p}))$。
注意，复形 $I^{\bullet, p}$ 下有界，由内射对象组成，
而且各微分的每个核、像和上同调群都是 $\mathcal{A}$ 的内射对象。
所以可以分裂这些微分；即每个微分都是到某个直和项上的分裂满射。
应用 $F$ 后亦然。因此
${}''E_1^{p, q} = F(H^q(I^{\bullet, p})) = F(H^q_I(I^{\bullet, p}))$。
其上的微分是 $(-1)^q$ 乘以 $F$ 作用于复形
% P02-E0049: frozen source says H_I^p here; the E_2 target requires the resolution of H^q(K^\bullet).
$H^p_I(I^{\bullet, \bullet})$ 的微分，而该复形是
$H^p(K^\bullet)$ 的内射消解。由此得到对 $E_2$ 项的描述。
\end{proof}

\begin{remark}
\label{derived-remark-functorial-ss}
引理 \ref{derived-lemma-two-ss-complex-functor} 的谱序列关于复形
$K^\bullet$ 具有函子性。这由 Cartan--Eilenberg 消解的函子性性质推出。
另一方面，二者都是一种更一般的谱序列的例子；该谱序列可与
$\mathcal{A}$ 的过滤复形相伴。函子性将由其构造推出。
我们将在过滤导出范畴一节回到这一点；见注记
\ref{derived-remark-final-functorial}。
\end{remark}











\section{右导出函子的复合}
\label{derived-section-composition-right-derived-functors}

\noindent
有时可以计算复合的右导出函子。
% P02-E0050: frozen source says “Suppose ... be abelian categories”.
设 $\mathcal{A}, \mathcal{B}, \mathcal{C}$ 为阿贝尔范畴。
设 $F : \mathcal{A} \to \mathcal{B}$ 和
$G : \mathcal{B} \to \mathcal{C}$ 为左正合函子。
假设右导出函子
$RF : D^{+}(\mathcal{A}) \to D^{+}(\mathcal{B})$、
$RG : D^{+}(\mathcal{B}) \to D^{+}(\mathcal{C})$ 和
$R(G \circ F) : D^{+}(\mathcal{A}) \to D^{+}(\mathcal{C})$
处处有定义。那么存在从 $D^{+}(\mathcal{A})$ 到
$D^{+}(\mathcal{C})$ 的函子之间的典范变换
$$
t : R(G \circ F) \longrightarrow RG \circ RF
$$
见引理 \ref{derived-lemma-compose-derived-functors-general}。
此变换未必总是同构。

\begin{lemma}
\label{derived-lemma-compose-derived-functors}
设 $\mathcal{A}, \mathcal{B}, \mathcal{C}$ 为阿贝尔范畴。
设 $F : \mathcal{A} \to \mathcal{B}$ 和
$G : \mathcal{B} \to \mathcal{C}$ 为左正合函子。
假设 $\mathcal{A}$、$\mathcal{B}$ 有足够多内射对象。
下列条件等价：
\begin{enumerate}
\item 对 $\mathcal{A}$ 的每个内射对象 $I$，$F(I)$ 都对 $G$ 右无环；
\item 典范映射
$$
t : R(G \circ F) \longrightarrow RG \circ RF.
$$
% P02-E0051: frozen source has a premature display period and omits the article in “is isomorphism”.
是从 $D^{+}(\mathcal{A})$ 到 $D^{+}(\mathcal{C})$ 的函子同构。
\end{enumerate}
\end{lemma}

\begin{proof}
若 (2) 成立，则把同构 $t$ 在 $RF(I) = F(I)$ 处取值即可得到 (1)。
反之，假设 (1) 成立。设 $A^\bullet$ 为 $\mathcal{A}$ 的下有界复形，
选取内射消解 $A^\bullet \to I^\bullet$。
映射 $t$ 由如下映射给出（见引理
\ref{derived-lemma-compose-derived-functors-general} 的证明）：
% P02-E0052: frozen source has an extra closing parenthesis in G(F(I^\bullet))).
$$
R(G \circ F)(A^\bullet) =
(G \circ F)(I^\bullet) =
G(F(I^\bullet))) \to
RG(F(I^\bullet)) =
RG(RF(A^\bullet))
$$
其中箭头由引理 \ref{derived-lemma-leray-acyclicity} 是同构。
\end{proof}

\begin{lemma}[Grothendieck 谱序列]
\label{derived-lemma-grothendieck-spectral-sequence}
% P02-E0053: frozen opening is a sentence fragment (“With assumptions ... and assuming ...”).
在引理 \ref{derived-lemma-compose-derived-functors} 的假设下，
并假定其中等价条件 (1) 和 (2) 成立。
设 $X$ 为 $D^{+}(\mathcal{A})$ 的对象。
存在由 $\mathcal{C}$ 的双分次对象 $E_r$ 组成的谱序列
$(E_r, d_r)_{r \geq 0}$，其中 $d_r$ 的双次数为
$(r, - r + 1)$，并且
$$
E_2^{p, q} = R^pG(H^q(RF(X)))
$$
此外，该谱序列有界，收敛到 $H^*(R(G \circ F)(X))$，
并在每个 $H^n(R(G \circ F)(X))$ 上诱导有限过滤。
\end{lemma}

\noindent
对 $\mathcal{A}$ 的对象 $A$，得到
$E_2^{p, q} = R^pG(R^qF(A))$，它收敛到
$R^{p + q}(G \circ F)(A)$。

\begin{proof}
可以用下有界复形 $A^\bullet$ 表示 $X$。
选取内射消解 $A^\bullet \to I^\bullet$。
利用引理 \ref{derived-lemma-cartan-eilenberg}，选取 Cartan--Eilenberg 消解
$F(I^\bullet) \to I^{\bullet, \bullet}$。
应用引理 \ref{derived-lemma-two-ss-complex-functor} 的第二个谱序列。
\end{proof}






\section{消解函子}
\label{derived-section-derived-category}

\noindent
设 $\mathcal{A}$ 为具有足够多内射对象的阿贝尔范畴。
以 $\mathcal{I}$ 表示 $\mathcal{A}$ 的满加性子范畴，
其对象是 $\mathcal{A}$ 的内射对象。
在这种情形下，$K^{+}(\mathcal{I})$ 与 $D^{+}(\mathcal{A})$ 等价
（见命题 \ref{derived-proposition-derived-category}）。
所以，为了许多目的，可以把 $D^{+}(\mathcal{A})$ 看作
（较易理解的）范畴 $K^{+}(\mathcal{I})$。

\begin{proposition}
\label{derived-proposition-derived-category}
设 $\mathcal{A}$ 为阿贝尔范畴。假设 $\mathcal{A}$ 有足够多内射对象。
以 $\mathcal{I} \subset \mathcal{A}$ 表示严格满加性子范畴，
其对象是 $\mathcal{A}$ 的内射对象。函子
$$
K^{+}(\mathcal{I}) \longrightarrow D^{+}(\mathcal{A})
$$
正合、满忠实且本质满；即它是三角范畴的等价。
\end{proposition}

\begin{proof}
该函子显然正合。由引理 \ref{derived-lemma-injective-resolutions-exist}，
它本质满。满忠实性由引理
% P02-E0039: another frozen “Fully faithfulness” locus; the noun is “Full faithfulness”.
\ref{derived-lemma-morphisms-into-injective-complex} 推出。
\end{proof}

\noindent
命题 \ref{derived-proposition-derived-category} 表明可以找到消解函子。
事实上，即使在阿贝尔范畴 $\mathcal{A}$ 是“大”范畴
（即具有对象的类）的某些情形下，也能证明消解函子存在。

\begin{definition}
\label{derived-definition-localization-functor}
设 $\mathcal{A}$ 为具有足够多内射对象的阿贝尔范畴。
$\mathcal{A}$ 的一个{\it 消解函子}\footnote{这很可能不是标准术语。}
由以下数据给出：
\begin{enumerate}
\item 对所有 $K^\bullet \in \Ob(K^{+}(\mathcal{A}))$，
给定内射对象的下有界复形 $j(K^\bullet)$；
\item 对所有 $K^\bullet \in \Ob(K^{+}(\mathcal{A}))$，
给定拟同构 $i_{K^\bullet} : K^\bullet \to j(K^\bullet)$。
\end{enumerate}
\end{definition}

\begin{lemma}
\label{derived-lemma-resolution-functor}
设 $\mathcal{A}$ 为具有足够多内射对象的阿贝尔范畴。
给定消解函子 $(j, i)$，存在唯一方式把 $j$ 变成函子，
把 $i$ 变成 $2$-同构，从而产生 $2$-交换图
$$
\xymatrix{
K^{+}(\mathcal{A}) \ar[rd] \ar[rr]_j & & K^{+}(\mathcal{I}) \ar[ld] \\
& D^{+}(\mathcal{A})
}
$$
其中 $\mathcal{I}$ 是 $\mathcal{A}$ 的满加性子范畴，
由内射对象组成。
\end{lemma}

\begin{proof}
对 $K^{+}(\mathcal{A})$ 的每个态射
$\alpha : K^\bullet \to L^\bullet$，存在唯一的
$K^{+}(\mathcal{I})$ 中的态射
$j(\alpha) : j(K^\bullet) \to j(L^\bullet)$，使得图
$$
\xymatrix{
K^\bullet \ar[r]_\alpha \ar[d]_{i_{K^\bullet}} &
L^\bullet \ar[d]^{i_{L^\bullet}} \\
j(K^\bullet) \ar[r]^{j(\alpha)} & j(L^\bullet)
}
$$
在 $K^{+}(\mathcal{A})$ 中交换。为此，可以使用引理
\ref{derived-lemma-morphisms-lift} 和
\ref{derived-lemma-morphisms-equal-up-to-homotopy}，或等价的引理
\ref{derived-lemma-morphisms-into-injective-complex}。
唯一性意味着 $j$ 是函子，而图的交换性意味着 $i$ 给出
$2$-态射，它见证本引理所述范畴图的 $2$-交换性。
\end{proof}

\begin{lemma}
\label{derived-lemma-into-derived-category}
设 $\mathcal{A}$ 为阿贝尔范畴。假设 $\mathcal{A}$ 有足够多内射对象。
则消解函子 $j$ 存在，并且在函子的唯一同构意义下唯一。
\end{lemma}

\begin{proof}
考虑 $K^{+}(\mathcal{A})$ 的所有对象 $K^\bullet$ 所成的集合。
（回忆依我们的约定，除非另有说明，任意范畴都有对象的集合。）
由引理 \ref{derived-lemma-injective-resolutions-exist}，每个对象都有内射消解。
由选择公理，可以对每个 $K^\bullet$ 选取内射消解
$i_{K^\bullet} : K^\bullet \to j(K^\bullet)$。
\end{proof}

\begin{lemma}
\label{derived-lemma-j-is-exact}
设 $\mathcal{A}$ 为具有足够多内射对象的阿贝尔范畴。
任意消解函子
$j : K^{+}(\mathcal{A}) \to K^{+}(\mathcal{I})$
都是正合的。
\end{lemma}

\begin{proof}
以 $i_{K^\bullet} : K^\bullet \to j(K^\bullet)$ 表示定义
\ref{derived-definition-localization-functor} 的典范映射。
先讨论函子同构 $j(K^\bullet[1]) \to j(K^\bullet)[1]$ 的存在性。
考虑图
$$
\xymatrix{
K^\bullet[1] \ar[d]^{i_{K^\bullet[1]}} \ar@{=}[rr] & &
K^\bullet[1] \ar[d]^{i_{K^\bullet}[1]} \\
j(K^\bullet[1]) \ar@{..>}[rr]^{\xi_{K^\bullet}} & & j(K^\bullet)[1]
}
$$
由引理 \ref{derived-lemma-morphisms-lift} 和
\ref{derived-lemma-morphisms-equal-up-to-homotopy}，
在 $K^{+}(\mathcal{I})$ 中存在唯一虚线箭头
$\xi_{K^\bullet}$，使图在 $K^{+}(\mathcal{A})$ 中交换。
我们略去这给出函子同构的验证。
（提示：再次使用引理 \ref{derived-lemma-morphisms-equal-up-to-homotopy}。）

\medskip\noindent
设 $(K^\bullet, L^\bullet, M^\bullet, f, g, h)$ 为
$K^{+}(\mathcal{A})$ 的特异三角。须证明
$(j(K^\bullet), j(L^\bullet), j(M^\bullet), j(f), j(g),
\xi_{K^\bullet} \circ j(h))$ 是 $K^{+}(\mathcal{I})$ 的特异三角。
注意，在 $K^{+}(\mathcal{A})$ 中有交换图
$$
\xymatrix{
K^\bullet \ar[r]_f \ar[d] &
L^\bullet \ar[r]_g \ar[d] &
M^\bullet \ar[rr]_h \ar[d] & &
K^\bullet[1] \ar[d] \\
j(K^\bullet) \ar[r]^{j(f)} &
j(L^\bullet) \ar[r]^{j(g)} &
j(M^\bullet) \ar[rr]^{\xi_{K^\bullet} \circ j(h)} & &
j(K^\bullet)[1]
}
$$
其竖直箭头为拟同构 $i_K, i_L, i_M$。
所以 $(j(K^\bullet), j(L^\bullet), j(M^\bullet), j(f), j(g),
\xi_{K^\bullet} \circ j(h))$ 在 $D^{+}(\mathcal{A})$ 中的像
与某个特异三角同构，故由 TR1 它本身就是特异三角。
于是由引理 \ref{derived-lemma-exact-equivalence}，
$(j(K^\bullet), j(L^\bullet), j(M^\bullet), j(f), j(g),
\xi_{K^\bullet} \circ j(h))$ 是 $K^{+}(\mathcal{I})$ 中的特异三角。
\end{proof}

\begin{lemma}
\label{derived-lemma-resolution-functor-quasi-inverse}
设 $\mathcal{A}$ 为具有足够多内射对象的阿贝尔范畴，
$j$ 为消解函子。以
$Q : K^{+}(\mathcal{A}) \to D^{+}(\mathcal{A})$ 表示自然函子。
则 $j = j' \circ Q$，其中唯一的函子
$j' : D^{+}(\mathcal{A}) \to K^{+}(\mathcal{I})$
是典范函子 $K^{+}(\mathcal{I}) \to D^{+}(\mathcal{A})$ 的拟逆。
\end{lemma}

\begin{proof}
由引理 \ref{derived-lemma-bounded-derived}，函子 $Q$ 是局部化。
为证明 $j'$ 存在，只须说明 $\text{Qis}^{+}(\mathcal{A})$
的任意元素在函子 $j$ 下映为同构；见引理
\ref{derived-lemma-universal-property-localization}。
考虑引理 \ref{derived-lemma-j-is-exact} 的证明中的交换方块。
在该方块中，若 $\alpha$ 是拟同构，则
$i_{L^\bullet}\circ\alpha=j(\alpha)\circ i_{K^\bullet}$
也是拟同构，因而 $j(\alpha)$ 也是拟同构。
所以由命题 \ref{derived-proposition-derived-category}，态射
$j(\alpha)$ 在 $K^+(\mathcal{I})$ 中是同构。
我们略去 $j'$ 是 $K^{+}(\mathcal{I}) \to D^{+}(\mathcal{A})$
的拟逆这一验证。
\end{proof}

\begin{remark}
\label{derived-remark-big-localization}
假设 $\mathcal{A}$ 是具有足够多内射对象的“大”阿贝尔范畴，
例如阿贝尔群范畴。此时构造消解函子时必须稍加谨慎，
因为不能对量词遍历一个类的情形使用选择公理。
但注意，该引理的证明确实表明任意两个局部化函子典范同构。
具体地，给定从下有界复形 $K^\bullet$ 到内射对象的下有界复形中的
拟同构 $i : K^\bullet \to I^\bullet$ 和
$i' : K^\bullet \to J^\bullet$，在 $K^{+}(\mathcal{I})$ 中存在唯一（！）
态射 $a : I^\bullet \to J^\bullet$，使得作为
% P02-E0054: frozen source has the ill-typed composition i' = i \circ a; it should compose a after i.
$K^{+}(\mathcal{I})$ 中的态射，有 $i' = i \circ a$。
所以唯一的问题是存在性；下一节将说明如何处理它。
\end{remark}









\section{函子性内射嵌入与消解函子}
\label{derived-section-functorial-injective-resolutions}

\noindent
本节在范畴 $\mathcal{A}$ 有函子性内射嵌入的情形下，
重新构造消解函子
$K^{+}(\mathcal{A}) \to K^{+}(\mathcal{I})$。
这样做有两个理由：(1) 证明更容易；(2) 当 $\mathcal{A}$ 是“大”
阿贝尔范畴时，该构造也适用。见下面的注记
\ref{derived-remark-big-abelian-category}。

\medskip\noindent
设 $\mathcal{A}$ 为阿贝尔范畴。与前面一样，以 $\mathcal{I}$ 表示
$\mathcal{A}$ 的加性满子范畴，由内射对象组成。
考虑箭头 $\alpha : K^\bullet \to I^\bullet$ 所成的范畴
$\text{InjRes}(\mathcal{A})$，其中 $K^\bullet$ 是
$\mathcal{A}$ 的下有界复形，$I^\bullet$ 是
$\mathcal{A}$ 中内射对象的下有界复形，而 $\alpha$ 是拟同构。
换言之，$\alpha$ 是内射消解，且 $K^\bullet$ 下有界。
有显然的函子
$$
s : \text{InjRes}(\mathcal{A}) \longrightarrow \text{Comp}^{+}(\mathcal{A})
$$
由 $(\alpha : K^\bullet \to I^\bullet) \mapsto K^\bullet$ 定义。
还有函子
$$
t : \text{InjRes}(\mathcal{A}) \longrightarrow K^{+}(\mathcal{I})
$$
由 $(\alpha : K^\bullet \to I^\bullet) \mapsto I^\bullet$ 定义。

\begin{lemma}
\label{derived-lemma-functorial-injective-resolutions}
设 $\mathcal{A}$ 为阿贝尔范畴。假设 $\mathcal{A}$ 有函子性内射嵌入；
见同调，定义 \ref{homology-definition-functorial-injective-embedding}。
\begin{enumerate}
\item 存在函子
$inj : \text{Comp}^{+}(\mathcal{A}) \to \text{InjRes}(\mathcal{A})$，
使得 $s \circ inj = \text{id}$。
\item 对任意满足 $s \circ inj = \text{id}$ 的函子
$inj : \text{Comp}^{+}(\mathcal{A}) \to \text{InjRes}(\mathcal{A})$，
可得到一个消解函子；见定义 \ref{derived-definition-localization-functor}。
\end{enumerate}
\end{lemma}

\begin{proof}
设 $A \mapsto (A \to J(A))$ 为函子性内射嵌入；见同调，定义
\ref{homology-definition-functorial-injective-embedding}。
先注意，可以假设 $J(0) = 0$。事实上，若非如此，
则对任意对象 $A$，有 $0 \to A \to 0$，从而给出直和分解
$J(A) = J(0) \oplus \Ker(J(A) \to J(0))$。
注意，函子性态射 $A \to J(A)$ 必须映入第二个直和项。
因此必要时可把函子替换为
$J'(A) = \Ker(J(A) \to J(0))$。

\medskip\noindent
设 $K^\bullet$ 为 $\mathcal{A}$ 的下有界复形。
设 $p < B$ 时 $K^p = 0$。我们将构造内射对象的双复形
$I^{\bullet, \bullet}$，以及映射
$\alpha : K^\bullet \to I^{\bullet, 0}$，使得 $\alpha$ 诱导
$K^\bullet$ 到 $I^{\bullet, \bullet}$ 的相伴全复形的拟同构。
首先，只要 $q < 0$，就令 $I^{p, q} = 0$。
接着，令 $I^{p, 0} = J(K^p)$，并令
$\alpha^p : K^p \to I^{p, 0}$ 为函子性嵌入。
由于 $J$ 是函子，$I^{\bullet, 0}$ 是复形，而 $\alpha$ 是复形态射。
每个 $\alpha^p$ 都为单射。并且由于 $J(0) = 0$，
当 $p < B$ 时 $I^{p, 0} = 0$。
再令 $I^{p, 1} = J(\Coker(K^p \to I^{p, 0}))$。
再次由函子性，$I^{\bullet, 1}$ 是复形。
而且当 $p < B$ 时 $I^{p, 1} = 0$。
还显然有 $K^p$ 同构地映到
$\Ker(I^{p, 0} \to I^{p, 1})$ 上。
第三步，取 $I^{p, 2} = J(\Coker(I^{p, 0} \to I^{p, 1}))$，
依此类推。

\medskip\noindent
现在可以应用同调，引理
\ref{homology-lemma-double-complex-gives-resolution}，得到映射
$$
\alpha : K^\bullet \longrightarrow \text{Tot}(I^{\bullet, \bullet})
$$
是拟同构。
% P02-E0056: frozen source says “it rests to show”; “it remains to show” is intended.
为证明得到函子 $inj$，还须说明上述构造具有函子性。
略去这一验证。

\medskip\noindent
假设已有满足 $s \circ inj = \text{id}$ 的函子 $inj$。
对 $\text{Comp}^{+}(\mathcal{A})$ 的每个对象 $K^\bullet$，可写
$$
inj(K^\bullet) = (i_{K^\bullet} : K^\bullet \to j(K^\bullet))
$$
这就给出定义 \ref{derived-definition-localization-functor} 所述的消解函子。
\end{proof}

\begin{remark}
\label{derived-remark-match}
假设 $inj$ 是引理 \ref{derived-lemma-functorial-injective-resolutions}
的 (2) 中满足 $s \circ inj = \text{id}$ 的函子。
如该引理的证明中那样，写成
$inj(K^\bullet) = (i_{K^\bullet} : K^\bullet \to j(K^\bullet))$。
设 $\alpha : K^\bullet \to L^\bullet$ 是下有界复形的映射。
考虑范畴 $\text{InjRes}(\mathcal{A})$ 中的映射 $inj(\alpha)$。
它诱导复形态射的交换图
% P02-E0057: frozen diagram writes j(K)^\bullet and j(L)^\bullet, inconsistent with j(K^\bullet), j(L^\bullet).
$$
\xymatrix{
K^\bullet
\ar[rr]^-{\alpha}
\ar[d]_{i_K} & &
L^\bullet \ar[d]^{i_L} \\
j(K)^\bullet
\ar[rr]^-{inj(\alpha)}
& &
j(L)^\bullet
}
$$
因此，考察引理 \ref{derived-lemma-resolution-functor} 的证明，
可见函子 $j : K^{+}(\mathcal{A}) \to K^{+}(\mathcal{I})$ 由规则
$$
j(\alpha\text{ 在同伦意义下}) = inj(\alpha)\text{ 在同伦意义下}\in
\Hom_{K^{+}(\mathcal{I})}(j(K^\bullet), j(L^\bullet))
$$
给出。所以此时 $j$ 与 $t \circ inj$ 相符；即图
$$
\xymatrix{
\text{Comp}^{+}(\mathcal{A}) \ar[rr]_{t \circ inj} \ar[rd] & &
K^{+}(\mathcal{I}) \\
& K^{+}(\mathcal{A}) \ar[ru]_j
}
$$
交换。
\end{remark}

\begin{remark}
\label{derived-remark-big-abelian-category}
设 $\textit{Mod}(\mathcal{O}_X)$ 为环空间
$(X, \mathcal{O}_X)$（或更一般地，环化位点）上的
$\mathcal{O}_X$-模范畴。稍后将看到，
$\textit{Mod}(\mathcal{O}_X)$ 有足够多内射对象，
而且事实上有函子性内射嵌入；见内射对象，定理
\ref{injectives-theorem-sheaves-modules-injectives}。
注意，引理 \ref{derived-lemma-into-derived-category} 的证明不适用于
$\textit{Mod}(\mathcal{O}_X)$，但引理
\ref{derived-lemma-functorial-injective-resolutions} 的证明适用于
$\textit{Mod}(\mathcal{O}_X)$。
因此得到
$$
j : K^{+}(\textit{Mod}(\mathcal{O}_X))
\longrightarrow
K^{+}(\mathcal{I})
$$
这是消解函子，其中 $\mathcal{I}$ 是内射
$\mathcal{O}_X$-模所成的加性范畴。该论证还适用于以下情形：
\begin{enumerate}
\item 环 $R$ 上的 $R$-模范畴 $\text{Mod}_R$。
\item 带环预层的位点上 $\mathcal{O}$-模预层的范畴
$\textit{PMod}(\mathcal{O})$。
\item 环化位点上 $\mathcal{O}$-模层的范畴
$\textit{Mod}(\mathcal{O})$。
% P02-E0058: frozen source contains the unfinished placeholder “Add more here as needed.”
\item 需要时在此添加更多内容。
\end{enumerate}
\end{remark}










\section{通过消解函子得到右导出函子}
\label{derived-section-right-derived-functor-via-resolutions}

\noindent
下面引理的内容是：若给定消解函子 $j$，
便可以直接定义 $RF(K^\bullet) = F(j(K^\bullet))$。

\begin{lemma}
\label{derived-lemma-right-derived-functor}
% P02-E0059: frozen source lacks punctuation after “enough injectives”.
设 $\mathcal{A}$ 为具有足够多内射对象的阿贝尔范畴。
设 $F : \mathcal{A} \to \mathcal{B}$ 为到某阿贝尔范畴中的加性函子。
设 $(i, j)$ 为消解函子；见定义
\ref{derived-definition-localization-functor}。
$F$ 的右导出函子 $RF$ 嵌入如下 $2$-交换图：
$$
\xymatrix{
D^{+}(\mathcal{A}) \ar[rd]_{RF} \ar[rr]^{j'} & &
K^{+}(\mathcal{I}) \ar[ld]^F \\
& D^{+}(\mathcal{B})
}
$$
其中 $j'$ 是引理 \ref{derived-lemma-resolution-functor-quasi-inverse}
中的函子。
\end{lemma}

\begin{proof}
由引理 \ref{derived-lemma-injective-acyclic}，有
$RF(K^\bullet) = F(j(K^\bullet))$。
\end{proof}

\begin{remark}
\label{derived-remark-right-derived-functor}
在引理 \ref{derived-lemma-right-derived-functor} 的情形中，
实际上已把右导出函子提升为正合函子
$F \circ j' : D^{+}(\mathcal{A}) \to K^{+}(\mathcal{B})$。
有时使用这样的分解很有用。
\end{remark}








\section{过滤导出范畴与内射消解}
\label{derived-section-filtered-derived}

\noindent
设 $\mathcal{A}$ 为阿贝尔范畴。本节将说明：若 $\mathcal{A}$
有足够多内射对象，则在某种意义下范畴
$\text{Fil}^f(\mathcal{A})$ 亦然。
可以利用这一观察在 $\mathcal{A}$ 的过滤导出范畴中作计算。

\medskip\noindent
范畴 $\text{Fil}^f(\mathcal{A})$ 是正合范畴的一个例子；见
内射对象，注记 \ref{injectives-remark-embed-exact-category}。
严格态射扮演特殊角色；见同调，定义
\ref{homology-definition-strict}，即满足
$\Coim(f) = \Im(f)$ 的态射 $f$。
若 $\text{Fil}^f(\mathcal{A})$ 中的复形 $A \to B \to C$
满足序列 $\text{gr}(A) \to \text{gr}(B) \to \text{gr}(C)$
在 $\mathcal{A}$ 中正合，则称该复形{\it 正合}。
这意味着 $A \to B$ 和 $B \to C$ 是严格态射；见同调，引理
\ref{homology-lemma-filtered-acyclic}。

\begin{definition}
\label{derived-definition-filtered-complexes-notation}
设 $\mathcal{A}$ 为阿贝尔范畴。
若 $\text{Fil}^f(\mathcal{A})$ 的对象 $I$ 的每个
$\text{gr}^p(I)$ 都是 $\mathcal{A}$ 的内射对象，
则称该对象为{\it 过滤内射对象}。
\end{definition}

\begin{lemma}
\label{derived-lemma-filtered-injective}
设 $\mathcal{A}$ 为阿贝尔范畴。
$\text{Fil}^f(\mathcal{A})$ 的对象 $I$ 为过滤内射对象，
当且仅当存在 $a \leq b$、$\mathcal{A}$ 的内射对象
$I_n$（$a \leq n \leq b$）以及同构
$I \cong \bigoplus_{a \leq n \leq b} I_n$，使得
$F^pI = \bigoplus_{n \geq p} I_n$。
\end{lemma}

\begin{proof}
这由如下事实推出：若 $J$ 是内射对象，则 $\mathcal{A}$ 的任意单射
$J \to M$ 都分裂。细节从略。
\end{proof}

\begin{lemma}
\label{derived-lemma-split-strict-monomorphism}
设 $\mathcal{A}$ 为阿贝尔范畴。
$\text{Fil}^f(\mathcal{A})$ 的任意严格单态射
$u : I \to A$，若 $I$ 为过滤内射对象，则它是分裂单射。
\end{lemma}

\begin{proof}
设 $p$ 为满足 $F^pI \not = 0$ 的最大整数。
特别地，$\text{gr}^p(I) = F^pI$。
设 $I'$ 为 $\text{Fil}^f(\mathcal{A})$ 的对象，其
$\mathcal{A}$ 中的底层对象为 $F^pI$，其过滤由如下条件给出：
当 $n > p$ 时 $F^nI' = 0$，当 $n \leq p$ 时
$F^nI' = I' = F^pI$。注意，$I' \to I$ 也是严格单态射。
$u$ 是严格单态射这一事实意味着
$F^pI \to A/F^{p + 1}(A)$ 为单射；见同调，引理
\ref{homology-lemma-characterize-strict}。
在 $\mathcal{A}$ 中选取分裂 $s : A/F^{p + 1}A \to F^pI$。
所诱导的态射 $s' : A \to I'$ 是过滤对象的严格态射，
并分裂复合 $I' \to I \to A$。所以可写
$A = I' \oplus \Ker(s')$ 以及
$I = I' \oplus \Ker(s'|_I)$。注意，
$\Ker(s'|_I) \to \Ker(s')$ 是严格单态射，且
$\Ker(s'|_I)$ 是过滤内射对象。
对 $I$ 的过滤长度归纳，映射
$\Ker(s'|_I) \to \Ker(s')$ 是分裂单射。证毕。
\end{proof}

\begin{lemma}
\label{derived-lemma-injective-property-filtered-injective}
设 $\mathcal{A}$ 为阿贝尔范畴。设
$u : A \to B$ 为 $\text{Fil}^f(\mathcal{A})$ 的严格单态射，
$f : A \to I$ 为从 $A$ 到
$\text{Fil}^f(\mathcal{A})$ 的过滤内射对象中的态射。
则存在态射 $g : B \to I$，使得 $f = g \circ u$。
\end{lemma}

\begin{proof}
$f$ 沿 $u$ 的推出 $f' : I \to I \amalg_A B$ 是严格单态射；见
同调，引理 \ref{homology-lemma-pushout-filtered}。
所以结论形式地由引理 \ref{derived-lemma-split-strict-monomorphism} 推出。
\end{proof}

\begin{lemma}
\label{derived-lemma-strict-monomorphism-into-filtered-injective}
设 $\mathcal{A}$ 为具有足够多内射对象的阿贝尔范畴。
对 $\text{Fil}^f(\mathcal{A})$ 的任意对象 $A$，
存在严格单态射 $A \to I$，其中 $I$ 是过滤内射对象。
\end{lemma}

\begin{proof}
选取 $a \leq b$，使得除非
$p \in \{a, a + 1, \ldots, b\}$，否则
$\text{gr}^p(A) = 0$。对每个
$n \in \{a, a + 1, \ldots, b\}$，选取单射
$u_n : A/F^{n + 1}A \to I_n$，其中 $I_n$ 是内射对象。
令 $I = \bigoplus_{a \leq n \leq b} I_n$，其过滤为
$F^pI = \bigoplus_{n \geq p} I_n$；并令 $u : A \to I$
等于各映射 $u_n$ 的直和。
\end{proof}

\begin{lemma}
\label{derived-lemma-filtered-injective-right-resolution-single-object}
设 $\mathcal{A}$ 为具有足够多内射对象的阿贝尔范畴。
对 $\text{Fil}^f(\mathcal{A})$ 的任意对象 $A$，
存在过滤拟同构 $A[0] \to I^\bullet$，其中
$I^\bullet$ 是过滤内射对象的复形，且对 $n < 0$ 有 $I^n = 0$。
\end{lemma}

\begin{proof}
首先选取从 $A$ 到某个过滤内射对象中的严格单态射
$u_0 : A \to I^0$；见引理
\ref{derived-lemma-strict-monomorphism-into-filtered-injective}。
接着，选取到 $\mathcal{A}$ 的某个过滤内射对象中的严格单态射
% P02-E0060: filtered injectives are objects of Fil^f(A), not of A as frozen prose states here and below.
$u_1 : \Coker(u_0) \to I^1$。以 $d^0$ 表示所诱导的映射
$I^0 \to I^1$。再选取到 $\mathcal{A}$ 的某个过滤内射对象中的
严格单态射 $u_2 : \Coker(u_1) \to I^2$。
以 $d^1$ 表示所诱导的映射 $I^1 \to I^2$，依此类推。
这是可行的，因为每个序列
$$
0 \to \Coker(u_n) \to I^{n + 1} \to \Coker(u_{n + 1}) \to 0
$$
都是短正合的，即应用 $\text{gr}$ 后诱导短正合序列。
为此使用同调，引理 \ref{homology-lemma-characterize-strict}。
\end{proof}

\begin{lemma}
\label{derived-lemma-filtered-injective-right-resolution-map}
设 $\mathcal{A}$ 为具有足够多内射对象的阿贝尔范畴。
设 $f : A \to B$ 为 $\text{Fil}^f(\mathcal{A})$ 的态射。
给定过滤拟同构 $A[0] \to I^\bullet$ 和
$B[0] \to J^\bullet$，其中 $I^\bullet, J^\bullet$ 是过滤内射对象的
复形，且对 $n < 0$ 有 $I^n = J^n = 0$，则存在交换图
$$
\xymatrix{
A[0] \ar[r] \ar[d] &
B[0] \ar[d] \\
I^\bullet \ar[r] &
J^\bullet
}
$$
\end{lemma}

\begin{proof}
% P02-E0061: the second object is B[0] in the lemma, but the frozen proof says C[0].
由于 $A[0] \to I^\bullet$ 和 $C[0] \to J^\bullet$ 是过滤拟同构，
可知 $a : A \to I^0$、$b : B \to J^0$ 以及所有态射
$d_I^n$、$d_J^n$ 都是严格态射；见同调，引理
\ref{homology-lemma-filtered-acyclic}。
我们将在下列交换图中归纳构造映射 $f^n$：
$$
\xymatrix{
A \ar[r]_a \ar[d]_f &
I^0 \ar[r] \ar[d]^{f^0} &
I^1 \ar[r] \ar[d]^{f^1} &
I^2 \ar[r] \ar[d]^{f^2} &
\ldots \\
B \ar[r]^b &
J^0 \ar[r] &
J^1 \ar[r] &
J^2 \ar[r] &
\ldots
}
$$
由于 $A \to I^0$ 是严格单态射，且 $J^0$ 是过滤内射对象，
由引理 \ref{derived-lemma-injective-property-filtered-injective}，可找到态射
$f^0 : I^0 \to J^0$，使得 $f^0 \circ a = b \circ f$。
复合 $d_J^0 \circ b \circ f$ 为零，故
$d_J^0 \circ f^0 \circ a = 0$，从而 $d_J^0 \circ f^0$
唯一地经由如下态射分解：
$$
\Coker(a) = \Coim(d_I^0) = \Im(d_I^0) \longrightarrow J^1.
$$
由于 $\Im(d_I^0) \to I^1$ 是严格单态射，再次应用引理
\ref{derived-lemma-injective-property-filtered-injective}，可将所示箭头延拓为态射
$f^1 : I^1 \to J^1$。依此类推。
\end{proof}

\begin{lemma}
\label{derived-lemma-filtered-injective-right-resolution-ses}
设 $\mathcal{A}$ 为具有足够多内射对象的阿贝尔范畴。
设 $0 \to A \to B \to C \to 0$ 为
$\text{Fil}^f(\mathcal{A})$ 中的短正合序列。
给定过滤拟同构 $A[0] \to I^\bullet$ 和
$C[0] \to J^\bullet$，其中 $I^\bullet, J^\bullet$ 是过滤内射对象的
复形，且对 $n < 0$ 有 $I^n = J^n = 0$，则存在交换图
$$
\xymatrix{
0 \ar[r] &
A[0] \ar[r] \ar[d] &
B[0] \ar[r] \ar[d] &
C[0] \ar[r] \ar[d] &
0 \\
0 \ar[r] &
I^\bullet \ar[r] &
M^\bullet \ar[r] &
J^\bullet \ar[r] &
0
}
$$
其中下行是逐项分裂的复形序列。
\end{lemma}

\begin{proof}
由于 $A[0] \to I^\bullet$ 和 $C[0] \to J^\bullet$ 是过滤拟同构，
可知 $a : A \to I^0$、$c : C \to J^0$ 以及所有态射
$d_I^n$、$d_J^n$ 都是严格态射；见同调，引理
% P02-E0062: the frozen cross-reference omits the homology- namespace used by the actual label.
\ref{derived-lemma-filtered-acyclic}。
我们将逐步构造下列交换图中朝向东南和朝南的箭头：
$$
\xymatrix{
B \ar[r]_\beta \ar[rd]^b &
C \ar[r]_c \ar[rd]^{\overline{b}} &
J^0 \ar[d]^{\delta^0} \ar[r] &
J^1 \ar[d]^{\delta^1} \ar[r] & \ldots \\
A \ar[u]^\alpha \ar[r]^a &
I^0 \ar[r] &
I^1 \ar[r] &
I^2 \ar[r] & \ldots
}
$$
由于 $A \to B$ 是严格单态射，由引理
\ref{derived-lemma-injective-property-filtered-injective}，可找到态射
$b : B \to I^0$，使得 $b \circ \alpha = a$。
由于 $A$ 是严格态射 $I^0 \to I^1$ 的核，且
$\beta = \Coker(\alpha)$，得到唯一的态射
$\overline{b} : C \to I^1$，使其嵌入图中。
由于 $c$ 是严格单态射且 $I^1$ 是过滤内射对象，
由引理 \ref{derived-lemma-injective-property-filtered-injective}，可找到
$\delta^0 : J^0 \to I^1$。
因为 $B \to C$ 是严格满态射，且
$B \to I^0 \to I^1 \to I^2$ 为零，所以
$C \to I^1 \to I^2$ 为零。因此 $d_I^1 \circ \delta^0$
在 $C \cong \Im(c)$ 上为零。
故 $d_I^1 \circ \delta^0$ 唯一地经由如下态射分解：
$$
\Coker(c) = \Coim(d_J^0) = \Im(d_J^0) \longrightarrow I^2.
$$
由于 $I^2$ 是过滤内射对象，且 $\Im(d_J^0) \to J^1$ 是
严格单态射，由引理 \ref{derived-lemma-injective-property-filtered-injective}，
可将所示态射延拓为态射 $\delta^1 : J^1 \to I^2$。
依此类推。令 $M^\bullet = I^\bullet \oplus J^\bullet$，其微分为
$$
d_M^n =
\left(
\begin{matrix}
d_I^n & (-1)^{n + 1}\delta^n \\
0 & d_J^n
\end{matrix}
\right)
$$
最后，映射 $B[0] \to M^\bullet$ 由
% P02-E0063: the frozen formula gives domain M although the map just defined has domain B.
$b \oplus c \circ \beta : M \to I^0 \oplus J^0$ 给出。
\end{proof}

\begin{lemma}
\label{derived-lemma-right-resolution-by-filtered-injectives}
设 $\mathcal{A}$ 为具有足够多内射对象的阿贝尔范畴。
对每个 $K^\bullet \in K^{+}(\text{Fil}^f(\mathcal{A}))$，
存在过滤拟同构 $K^\bullet \to I^\bullet$，其中 $I^\bullet$ 有下界，
每个 $I^n$ 都是过滤内射对象，且每个
$K^n \to I^n$ 都是严格单态射。
\end{lemma}

\begin{proof}
把 $K^\bullet$ 替换为某个平移后（这对证明无碍），可假设对
$n < 0$ 有 $K^n = 0$。考虑正合范畴
$\text{Fil}^f(\mathcal{A})$ 中的短正合序列
$$
\begin{matrix}
0 \to \Ker(d_K^0) \to K^0 \to \Coim(d_K^0) \to 0 \\
0 \to \Ker(d_K^1) \to K^1 \to \Coim(d_K^1) \to 0 \\
0 \to \Ker(d_K^2) \to K^2 \to \Coim(d_K^2) \to 0 \\
\ldots
\end{matrix}
$$
以及映射 $u_i : \Coim(d_K^i) \to \Ker(d_K^{i + 1})$。
对每个 $i \geq 0$，可以选取过滤拟同构
$$
\begin{matrix}
\Ker(d_K^i)[0] \to I_{ker, i}^\bullet \\
\Coim(d_K^i)[0] \to I_{coim, i}^\bullet
\end{matrix}
$$
其中 $I_{ker, i}^n, I_{coim, i}^n$ 是过滤内射对象，并且对
$n < 0$ 为零；见引理
\ref{derived-lemma-filtered-injective-right-resolution-single-object}。
由引理 \ref{derived-lemma-filtered-injective-right-resolution-map}，
可将 $u_i$ 提升为复形态射
$u_i^\bullet : I_{coim, i}^\bullet \to I_{ker, i + 1}^\bullet$。
最后，对每个 $i \geq 0$，由引理
\ref{derived-lemma-filtered-injective-right-resolution-ses}，可补全各图
$$
\xymatrix{
0 \ar[r] &
\Ker(d_K^i)[0] \ar[r] \ar[d] &
K^i[0] \ar[r] \ar[d] &
\Coim(d_K^i)[0] \ar[r] \ar[d] &
0 \\
0 \ar[r] &
I_{ker, i}^\bullet \ar[r]^{\alpha_i} &
I_i^\bullet \ar[r]^{\beta_i} &
I_{coim, i}^\bullet \ar[r] &
0
}
$$
使得下行是逐项分裂的正合序列。
对 $i \geq 0$，令 $d_i : I_i^\bullet \to I_{i + 1}^\bullet$
等于 $d_i =  \alpha_{i + 1} \circ u_i^\bullet \circ \beta_i$。
注意，$d_i \circ d_{i - 1} = 0$，因为
$\beta_i \circ \alpha_i = 0$。因此我们构造了交换图
$$
\xymatrix{
I_0^\bullet \ar[r] &
I_1^\bullet \ar[r] &
I_2^\bullet \ar[r] & \ldots \\
K^0[0] \ar[r] \ar[u] &
K^1[0] \ar[r] \ar[u] &
K^2[0] \ar[r] \ar[u] &
\ldots
}
$$
这里竖直箭头是过滤拟同构。
上行是复形的复形，且每个复形由过滤内射对象组成，
在次数 $< 0$ 处没有非零对象。
因此令 $I^{a, b} = I_a^b$，并用
$$
d_1^{a, b} : I^{a, b} = I_a^b \to I_{a + 1}^b = I^{a + 1, b}
$$
表示映射 $d_a^b$，又用
$$
d_2^{a, b} : I^{a, b} = I_a^b \to I_a^{b + 1} = I^{a, b + 1}
$$
表示映射 $d_{I_a}^b$，便得到双复形。
以 $\text{Tot}(I^{\bullet, \bullet})$ 表示与此双复形关联的全复形；见
同调，定义 \ref{homology-definition-associated-simple-complex}。
注意，映射 $K^n[0] \to I_n^\bullet$ 来自映射
$K^n \to I^{n, 0}$，后者诱导复形态射
$$
K^\bullet \longrightarrow \text{Tot}(I^{\bullet, \bullet})
$$
我们断言这是过滤拟同构。
由于 $\text{gr}(-)$ 是加性函子，可见
$\text{gr}(\text{Tot}(I^{\bullet, \bullet})) =
\text{Tot}(\text{gr}(I^{\bullet, \bullet}))$。
因此可以应用同调，引理
\ref{homology-lemma-double-complex-gives-resolution}，得出
$\text{gr}(K^\bullet) \to \text{gr}(\text{Tot}(I^{\bullet, \bullet}))$
如所需是拟同构。
\end{proof}

\begin{lemma}
\label{derived-lemma-filtered-acyclic-is-zero}
设 $\mathcal{A}$ 为阿贝尔范畴。
设 $K^\bullet, I^\bullet \in K(\text{Fil}^f(\mathcal{A}))$。
假设 $K^\bullet$ 过滤无环，并且 $I^\bullet$ 有下界且由过滤内射对象组成。
任意态射 $K^\bullet \to I^\bullet$ 都零同伦：
$\Hom_{K(\text{Fil}^f(\mathcal{A}))}(K^\bullet, I^\bullet) = 0$。
\end{lemma}

\begin{proof}
设 $\alpha : K^\bullet \to I^\bullet$ 为复形态射。
假设对 $j < n$ 有 $\alpha^j = 0$。
我们将证明存在态射 $h : K^{n + 1} \to I^n$，使得
$\alpha^n = h \circ d$。于是 $\alpha$ 将同伦于由下式定义的复形态射
$\beta$：
$$
\beta^j =
\left\{
\begin{matrix}
0 & \text{若} & j \leq n \\
\alpha^{n + 1} - d \circ h & \text{若} & j = n + 1 \\
\alpha^j & \text{若} & j > n + 1
\end{matrix}
\right.
$$
这显然（通过归纳）证明了引理。
为证明 $h$ 的存在性，注意
$\alpha^n \circ d_K^{n - 1} = 0$，因为
$\alpha^{n - 1} = 0$。由于 $K^\bullet$ 过滤无环，
$d_K^{n - 1}$ 和 $d_K^n$ 都是严格态射，并且
$$
0 \to \Im(d_K^{n - 1}) \to K^n \to \Im(d_K^n) \to 0
$$
是正合范畴 $\text{Fil}^f(\mathcal{A})$ 中的正合序列；见
同调，引理 \ref{homology-lemma-filtered-acyclic}。
因此可把 $\alpha^n$ 看作定义在 $\Im(d_K^n)$ 上、取值于
$I^n$ 的映射。
利用 $\Im(d_K^n) \to K^{n + 1}$ 是严格单态射且 $I^n$
为过滤内射对象，由引理
\ref{derived-lemma-injective-property-filtered-injective}，可如所需把此映射提升为
$h : K^{n + 1} \to I^n$。
\end{proof}

\begin{lemma}
\label{derived-lemma-morphisms-into-filtered-injective-complex}
设 $\mathcal{A}$ 为阿贝尔范畴。
设 $I^\bullet \in K(\text{Fil}^f(\mathcal{A}))$ 是由过滤内射对象组成的
有下界复形。
\begin{enumerate}
\item 设 $\alpha : K^\bullet \to L^\bullet$ 是
$K(\text{Fil}^f(\mathcal{A}))$ 中的过滤拟同构。则映射
$$
\Hom_{K(\text{Fil}^f(\mathcal{A}))}(L^\bullet, I^\bullet)
\to
\Hom_{K(\text{Fil}^f(\mathcal{A}))}(K^\bullet, I^\bullet)
$$
是双射。
\item 设 $L^\bullet \in K(\text{Fil}^f(\mathcal{A}))$。则
$$
\Hom_{K(\text{Fil}^f(\mathcal{A}))}(L^\bullet, I^\bullet)
=
\Hom_{DF(\mathcal{A})}(L^\bullet, I^\bullet).
$$
\end{enumerate}
\end{lemma}

\begin{proof}
证明 (1)。注意
$$
(K^\bullet, L^\bullet, C(\alpha)^\bullet, \alpha, i, -p)
$$
是 $K(\text{Fil}^f(\mathcal{A}))$ 中的特异三角
（引理 \ref{derived-lemma-the-same-up-to-isomorphisms}），并且
$C(\alpha)^\bullet$ 是过滤无环复形
（引理 \ref{derived-lemma-filtered-acyclic}）。于是
$$
\xymatrix{
\Hom_{K(\text{Fil}^f(\mathcal{A}))}(C(\alpha)^\bullet, I^\bullet) \ar[r] &
\Hom_{K(\text{Fil}^f(\mathcal{A}))}(L^\bullet, I^\bullet) \ar[r] &
\Hom_{K(\text{Fil}^f(\mathcal{A}))}(K^\bullet, I^\bullet) \ar[lld] \\
\Hom_{K(\text{Fil}^f(\mathcal{A}))}(C(\alpha)^\bullet[-1], I^\bullet)
}
$$
是阿贝尔群的正合序列；见引理
\ref{derived-lemma-representable-homological}。
此时引理 \ref{derived-lemma-filtered-acyclic-is-zero} 保证外侧两个群为零，
因而
% P02-E0064: the frozen conclusion drops Fil^f(A) from both homotopy categories.
$\Hom_{K(\mathcal{A})}(L^\bullet, I^\bullet) =
\Hom_{K(\mathcal{A})}(K^\bullet, I^\bullet)$。

\medskip\noindent
证明 (2)。
设 $a$ 为右侧的一个元素。
可将其表示成 $a = \gamma\alpha^{-1}$，其中
$\alpha : K^\bullet \to L^\bullet$ 是过滤拟同构，而
$\gamma : K^\bullet \to I^\bullet$ 是复形映射。由 (1)，
可找到态射 $\beta : L^\bullet \to I^\bullet$，使得
$\beta \circ \alpha$ 同伦于 $\gamma$。这证明该映射是满射。
设 $b$ 为左侧的元素，且在右侧中的像为零。则 $b$ 是某个态射
$\beta : L^\bullet \to I^\bullet$ 的同伦类，并且存在过滤拟同构
$\alpha : K^\bullet \to L^\bullet$，使得
$\beta  \circ \alpha$ 零同伦。于是 (1) 表明
$\beta$ 也零同伦，即 $b = 0$。
\end{proof}

\begin{lemma}
\label{derived-lemma-filtered-localization-functor}
设 $\mathcal{A}$ 为具有足够多内射对象的阿贝尔范畴。
以 $\mathcal{I}^f \subset \text{Fil}^f(\mathcal{A})$ 表示对象为
过滤内射对象的严格满加性子范畴。典范函子
$$
K^{+}(\mathcal{I}^f)
\longrightarrow
DF^{+}(\mathcal{A})
$$
正合、全忠实且本质满，即它是三角范畴的等价。此外各图
$$
\xymatrix{
K^{+}(\mathcal{I}^f) \ar[d]_{\text{gr}^p} \ar[r] &
DF^{+}(\mathcal{A}) \ar[d]_{\text{gr}^p} \\
K^{+}(\mathcal{I}) \ar[r] &
D^{+}(\mathcal{A})
}
\quad
\xymatrix{
K^{+}(\mathcal{I}^f) \ar[d]^{\text{遗忘 }F} \ar[r] &
DF^{+}(\mathcal{A}) \ar[d]^{\text{遗忘 }F} \\
K^{+}(\mathcal{I}) \ar[r] &
D^{+}(\mathcal{A})
}
$$
交换，其中 $\mathcal{I} \subset \mathcal{A}$ 是对象为内射对象的
严格满加性子范畴。
\end{lemma}

\begin{proof}
由引理 \ref{derived-lemma-right-resolution-by-filtered-injectives}，函子
$K^{+}(\mathcal{I}^f) \to DF^{+}(\mathcal{A})$ 本质满。
由引理 \ref{derived-lemma-morphisms-into-filtered-injective-complex}，它全忠实。
由我们关于特异三角的定义，它是正合函子。
各方块的交换性是立即的。
\end{proof}

\begin{remark}
\label{derived-remark-filtered-localization-big}
只有当 $\mathcal{A}$ 是我们意义下的范畴，即它具有对象的集合时，
我们才能逆转引理中的箭头。不过，假设给定一个具有足够多内射对象的
大阿贝尔范畴 $\mathcal{A}$，例如 $\textit{Mod}(\mathcal{O}_X)$。
那么对任意给定的对象集合 $\{A_i\}_{i\in I}$，存在包含所有这些对象且
具有足够多内射对象的阿贝尔子范畴
$\mathcal{A}' \subset \mathcal{A}$；见集合，引理
\ref{sets-lemma-abelian-injectives}。
因此可以对 $\mathcal{A}'$ 使用上述引理。
这实质上意味着，若我们只使用一集合之多的图等等，
那么使用该引理绝不会遇到麻烦。
\end{remark}

\noindent
设 $\mathcal{A}, \mathcal{B}$ 为阿贝尔范畴。
设 $T : \mathcal{A} \to \mathcal{B}$ 为左正合函子。
（我们不能用字母 $F$ 表示该函子，因为这会与我们用字母 $F$
表示过滤严重冲突。）注意，$T$ 按规则
$T(A, F) = (T(A), F)$ 诱导加性函子
$$
T : \text{Fil}^f(\mathcal{A}) \to \text{Fil}^f(\mathcal{B})
$$
其中 $F^pT(A) = T(F^pA)$；这有意义，因为 $T$ 左正合。
（警告：未必有 $\text{gr}(T(A)) = T(\text{gr}(A))$。）
这诱导三角范畴的函子
\begin{equation}
\label{derived-equation-induced-T-filtered}
T :
K^{+}(\text{Fil}^f(\mathcal{A}))
\longrightarrow
K^{+}(\text{Fil}^f(\mathcal{B})).
\end{equation}
$T$ 的过滤右导出函子，是把定义
\ref{derived-definition-right-derived-functor-defined} 应用于这个正合函子和乘法系
$\text{FQis}^{+}(\mathcal{A})$ 所得的右导出函子，再与正合函子
$K^{+}(\text{Fil}^f(\mathcal{B})) \to DF^{+}(\mathcal{B})$ 复合。
假设 $\mathcal{A}$ 具有足够多内射对象。此时可以重复
第 \ref{derived-section-right-derived-functor} 节的讨论，定义函子 $T$ 的
{\it 过滤右导出函子}
\begin{equation}
\label{derived-equation-filtered-derived-functor}
RT : DF^{+}(\mathcal{A}) \longrightarrow DF^{+}(\mathcal{B})
\end{equation}

\medskip\noindent
不过，我们将改为像第
\ref{derived-section-right-derived-functor-via-resolutions} 节那样进行；结果表明，
即使 $T$ 只是加性函子，也能定义 $RT$。
具体地，首先选取引理 \ref{derived-lemma-filtered-localization-functor} 中等价的
拟逆
$j' : DF^{+}(\mathcal{A}) \to K^{+}(\mathcal{I}^f)$。
由引理 \ref{derived-lemma-exact-equivalence}，可知 $j'$ 是三角范畴的正合函子。
接着注意，对过滤内射对象 $I$，由引理
\ref{derived-lemma-filtered-injective}，我们有（非典范）分解
\begin{equation}
\label{derived-equation-decompose}
I \cong \bigoplus\nolimits_{p \in \mathbf{Z}} I_p,
\quad\text{其中}\quad
F^pI = \bigoplus\nolimits_{q \geq p} I_q
\end{equation}
因此，若 $T$ 是任意加性函子 $T : \mathcal{A} \to \mathcal{B}$，
则令 $T_{ext}(I) = \bigoplus T(I_p)$ 且
$F^pT_{ext}(I) = \bigoplus_{q \geq p} T(I_q)$，便得到加性函子
\begin{equation}
\label{derived-equation-extend-T}
T_{ext} : \mathcal{I}^f \to \text{Fil}^f(\mathcal{B})
\end{equation}
注意，由构造有性质 $\text{gr}(T_{ext}(I)) = T(\text{gr}(I))$。
于是得到与 $\text{gr}$ 交换的函子
\begin{equation}
\label{derived-equation-extend-T-complexes}
T_{ext} : K^{+}(\mathcal{I}^f) \to K^{+}(\text{Fil}^f(\mathcal{B}))
\end{equation}
然后把 (\ref{derived-equation-filtered-derived-functor}) 定义为复合
\begin{equation}
\label{derived-equation-definition-filtered-derived-functor}
RT = T_{ext} \circ j'.
\end{equation}
由于 $RT : D^{+}(\mathcal{A}) \to D^{+}(\mathcal{B})$
也由内射消解计算；见引理 \ref{derived-lemma-injective-acyclic}，
$T$ 与 $\text{gr}$ 的交换性以及引理
\ref{derived-lemma-filtered-localization-functor} 中的交换图蕴含
\begin{equation}
\label{derived-equation-commute-gr}
\text{gr}^p \circ RT \cong RT \circ \text{gr}^p
\end{equation}
以及
\begin{equation}
\label{derived-equation-commute-forget}
(\text{遗忘 }F) \circ RT \cong RT \circ (\text{遗忘 }F)
\end{equation}
是函子 $DF^{+}(\mathcal{A}) \to D^{+}(\mathcal{B})$ 的同构。

\medskip\noindent
过滤导出函子 $RT$（\ref{derived-equation-filtered-derived-functor}）诱导函子
$$
\begin{matrix}
RT : \text{Fil}^f(\mathcal{A}) \to DF^{+}(\mathcal{B}), \\
RT : \text{Comp}^{+}(\text{Fil}^f(\mathcal{A})) \to DF^{+}(\mathcal{B}), \\
RT : KF^{+}(\mathcal{A}) \to DF^{+}(\mathcal{B}).
\end{matrix}
$$
注意，由于 $\text{Fil}^f(\mathcal{A})$ 和
$\text{Comp}^{+}(\text{Fil}^f(\mathcal{A}))$ 不再是阿贝尔范畴，
说 $RT$ 在它们上限制为 $\delta$-函子没有意义。
（把这些范畴视为正合范畴，并表述从正合范畴到三角范畴的
$\delta$-函子概念，可以修复这一点。）
但是，说 $RT$ 是三角范畴 $KF^{+}(\mathcal{A})$ 上的正合函子是有意义的，
并且由构造确实如此。

\begin{lemma}
\label{derived-lemma-ss-filtered-derived}
设 $\mathcal{A}, \mathcal{B}$ 为阿贝尔范畴。设
$T : \mathcal{A} \to \mathcal{B}$ 为左正合函子。
假设 $\mathcal{A}$ 具有足够多内射对象。
设 $(K^\bullet, F)$ 为
$\text{Comp}^{+}(\text{Fil}^f(\mathcal{A}))$ 的对象。
存在谱序列 $(E_r, d_r)_{r\geq 0}$，它由
$\mathcal{B}$ 的双分次对象 $E_r$ 和双次数为 $(r, - r + 1)$ 的
$d_r$ 组成，并且
$$
E_1^{p, q} = R^{p + q}T(\text{gr}^p(K^\bullet))
$$
此外，该谱序列有界，收敛到 $R^*T(K^\bullet)$，并在每个
$R^nT(K^\bullet)$ 上诱导有限过滤。该谱序列的构造关于
$\text{Comp}^{+}(\text{Fil}^f(\mathcal{A}))$ 的对象
$K^\bullet$ 具有函子性，而且对 $r \geq 1$，各项
$(E_r, d_r)$ 不依赖任何选择。
\end{lemma}

\begin{proof}
选取过滤拟同构 $K^\bullet \to I^\bullet$，其中 $I^\bullet$
是过滤内射对象的有下界复形；见引理
\ref{derived-lemma-right-resolution-by-filtered-injectives}。
考虑复形 $RT(K^\bullet) = T_{ext}(I^\bullet)$；见
(\ref{derived-equation-definition-filtered-derived-functor})。
于是可考虑把它视为 $\mathcal{B}$ 中的过滤复形而关联的谱序列
$(E_r, d_r)_{r \geq 0}$；见同调，第
\ref{homology-section-filtered-complex} 节。
由同调，引理
\ref{homology-lemma-spectral-sequence-filtered-complex}，有
$E_1^{p, q} = H^{p + q}(\text{gr}^p(T(I^\bullet)))$。
由式 (\ref{derived-equation-decompose})，有
$E_1^{p, q} = H^{p + q}(T(\text{gr}^p(I^\bullet)))$；而由过滤内射消解的
定义，映射 $\text{gr}^p(K^\bullet) \to \text{gr}^p(I^\bullet)$
是内射消解。因此
$E_1^{p, q} = R^{p + q}T(\text{gr}^p(K^\bullet))$。

\medskip\noindent
另一方面，每个 $I^n$ 都有有限过滤，因而每个 $T(I^n)$ 也有有限过滤。
所以可以应用同调，引理
\ref{homology-lemma-biregular-ss-converges}，得出该谱序列有界且收敛到
$H^n(T(I^\bullet)) = R^nT(K^\bullet)$，并且还在每一项上诱导有限过滤。

\medskip\noindent
假设 $K^\bullet \to L^\bullet$ 是
$\text{Comp}^{+}(\text{Fil}^f(\mathcal{A}))$ 的态射。
选取过滤拟同构 $L^\bullet \to J^\bullet$，其中 $J^\bullet$
是过滤内射对象的有下界复形；见引理
\ref{derived-lemma-right-resolution-by-filtered-injectives}。
由上述结果，例如引理
\ref{derived-lemma-morphisms-into-filtered-injective-complex}，存在图
$$
\xymatrix{
K^\bullet \ar[r] \ar[d] & L^\bullet \ar[d] \\
I^\bullet \ar[r] & J^\bullet
}
$$
它在同伦意义下交换。因此得到过滤复形的态射
$T(I^\bullet) \to T(J^\bullet)$，它诱导谱序列的态射；见同调，引理
\ref{homology-lemma-spectral-sequence-filtered-complex-functorial}。
最后一个陈述由此得出。
\end{proof}

\begin{remark}
\label{derived-remark-final-functorial}
正如评注 \ref{derived-remark-functorial-ss} 中所承诺的，我们讨论上述引理与
使用 Cartan--Eilenberg 消解的构造之间的联系。
设 $T : \mathcal{A} \to \mathcal{B}$ 为阿贝尔范畴之间的左正合函子，
假设 $\mathcal{A}$ 具有足够多内射对象，并设 $K^\bullet$ 为
$\mathcal{A}$ 的有下界复形。我们给出引理
\ref{derived-lemma-two-ss-complex-functor} 中谱序列 ${}'E$ 和 ${}''E$
的另一种构造。

\medskip\noindent
第一个谱序列。考虑 $K^\bullet$ 上的“愚蠢”过滤，其定义为
$F^p(K^\bullet) = \sigma_{\geq p}(K^\bullet)$；见同调，第
\ref{homology-section-truncations} 节。
% P02-E0065: frozen prose reads “Note that this stupid” and omits “is”.
注意，这在 $d(F^p(K^\bullet)) \subset F^{p + 1}(K^\bullet)$ 的意义下
是愚蠢的；比较同调，引理
\ref{homology-lemma-spectral-sequence-filtered-complex-d1}。
注意，在此过滤下 $\text{gr}^p(K^\bullet) = K^p[-p]$。
由引理 \ref{derived-lemma-ss-filtered-derived}，存在谱序列，其 $E_1$ 项为
$$
E_1^{p, q} = R^{p + q}T(K^p[-p]) = R^qT(K^p)
$$
正如谱序列 ${}'E_r$ 中那样。此外，注意微分
$E_1^{p, q} \to E_1^{p + 1, q}$ 与 $'{}E_1$ 中的微分相同；见
同调，引理 \ref{homology-lemma-spectral-sequence-filtered-complex-d1}
的 (2) 以及引理 \ref{derived-lemma-two-ss-complex-functor} 证明中对
${}'d_1$ 的描述。

\medskip\noindent
第二个谱序列。考虑 $K^\bullet$ 上由
$F^p(K^\bullet) = \tau_{\leq -p}(K^\bullet)$ 定义的过滤；见同调，第
\ref{homology-section-truncations} 节。
负号是得到递减过滤所必需的。注意，在此过滤下
$\text{gr}^p(K^\bullet)$ 拟同构于 $H^{-p}(K^\bullet)[p]$。
由引理 \ref{derived-lemma-ss-filtered-derived}，存在谱序列，其 $E_1$ 项为
$$
E_1^{p, q} = R^{p + q}T(H^{-p}(K^\bullet)[p])
= R^{2p + q}T(H^{-p}(K^\bullet)) = {}''E_2^{i, j}
$$
其中 $i = 2p + q$ 且 $j = -p$。（这看起来不自然，但注意，我们也
完全可以用递增过滤来发展整个过滤复形理论；结果便是这里看起来自然，
而另一个却不自然。）留给读者验证各微分相符。

\medskip\noindent
事实上，给定 Cartan--Eilenberg 消解
$K^\bullet \to I^{\bullet, \bullet}$，诱导的态射
$K^\bullet \to \text{Tot}(I^{\bullet, \bullet})$ 取值于相应的全复形；
对 $\text{Tot}(I^{\bullet, \bullet})$ 使用适当过滤后，
它对任一过滤都是过滤内射消解。
这可用于精确匹配这些谱序列。
\end{remark}






\section{Ext 群}
\label{derived-section-ext}

\noindent
本节开始描述阿贝尔范畴中对象的 Ext 群。首先给出如下非常一般的定义。

\begin{definition}
\label{derived-definition-ext}
设 $\mathcal{A}$ 为阿贝尔范畴。设 $i \in \mathbf{Z}$。设
$X, Y$ 为 $D(\mathcal{A})$ 的对象。$X$ 被 $Y$ 扩张的
{\it 第 $i$ 个 Ext 群}是群
$$
\Ext^i_\mathcal{A}(X, Y) =
\Hom_{D(\mathcal{A})}(X, Y[i]) =
\Hom_{D(\mathcal{A})}(X[-i], Y).
$$
若 $A, B \in \Ob(\mathcal{A})$，则令
$\Ext^i_\mathcal{A}(A, B) = \text{Ext}^i_\mathcal{A}(A[0], B[0])$。
\end{definition}

\noindent
由于 $\Hom_{D(\mathcal{A})}(X, -)$ 是同调函子，而
$\Hom_{D(\mathcal{A})}(-, Y)$ 是上同调函子；见引理
\ref{derived-lemma-representable-homological}，所以特异三角
$(Y, Y', Y'')$ 分别诱导长正合序列
$$
\ldots \to
\Ext^i_\mathcal{A}(X, Y) \to
\Ext^i_\mathcal{A}(X, Y') \to
\Ext^i_\mathcal{A}(X, Y'') \to
\Ext^{i + 1}_\mathcal{A}(X, Y) \to \ldots
$$
而特异三角 $(X, X', X'')$ 分别诱导
$$
\ldots \to
\Ext^i_\mathcal{A}(X'', Y) \to
\Ext^i_\mathcal{A}(X', Y) \to
\Ext^i_\mathcal{A}(X, Y) \to
\Ext^{i + 1}_\mathcal{A}(X'', Y) \to \ldots
$$
注意，由于 $D^+(\mathcal{A})$、$D^-(\mathcal{A})$、$D^b(\mathcal{A})$
是满子范畴，只要 $X$、$Y$ 属于它们，就可以用这些范畴中的 Hom 群
计算 Ext 群。

\medskip\noindent
若范畴 $\mathcal{A}$ 具有足够多内射对象或足够多投射对象，
则可分别用内射消解或投射消解计算 Ext 群。为避免混淆，回忆：
具有内射（相应地，投射）消解蕴含所有低（相应地，高）次数的同调消失；
见引理 \ref{derived-lemma-cohomology-bounded-below} 和
\ref{derived-lemma-cohomology-bounded-above}。

\begin{lemma}
\label{derived-lemma-compute-ext-resolutions}
设 $\mathcal{A}$ 为阿贝尔范畴。
设 $X^\bullet, Y^\bullet \in \Ob(K(\mathcal{A}))$。
\begin{enumerate}
\item 设 $Y^\bullet \to I^\bullet$ 为内射消解
（定义 \ref{derived-definition-injective-resolution}）。则
$$
\Ext^i_\mathcal{A}(X^\bullet, Y^\bullet) =
\Hom_{K(\mathcal{A})}(X^\bullet, I^\bullet[i]).
$$
\item 设 $P^\bullet \to X^\bullet$ 为投射消解
（定义 \ref{derived-definition-projective-resolution}）。则
$$
\Ext^i_\mathcal{A}(X^\bullet, Y^\bullet) =
\Hom_{K(\mathcal{A})}(P^\bullet[-i], Y^\bullet).
$$
\end{enumerate}
\end{lemma}

\begin{proof}
立即由引理 \ref{derived-lemma-morphisms-into-injective-complex} 和
引理 \ref{derived-lemma-morphisms-from-projective-complex} 得出。
\end{proof}

\noindent
在本节其余部分，我们讨论阿贝尔范畴本身的对象的扩张。
首先注意如下事实。

\begin{lemma}
\label{derived-lemma-negative-exts}
设 $\mathcal{A}$ 为阿贝尔范畴。
\begin{enumerate}
\item 设 $X$、$Y$ 为 $D(\mathcal{A})$ 的对象。给定
$a, b \in \mathbf{Z}$，使得当 $i > a$ 时 $H^i(X) = 0$，
而当 $j < b$ 时 $H^j(Y) = 0$，则当 $n < b - a$ 时
$\Ext^n_\mathcal{A}(X, Y) = 0$，并且
$$
\Ext^{b - a}_\mathcal{A}(X, Y) = \Hom_\mathcal{A}(H^a(X), H^b(Y))
$$
\item 设 $A, B \in \Ob(\mathcal{A})$。
当 $i < 0$ 时有 $\Ext^i_\mathcal{A}(B, A) = 0$。
并有 $\Ext^0_\mathcal{A}(B, A) = \Hom_\mathcal{A}(B, A)$。
\end{enumerate}
\end{lemma}

\begin{proof}
选取表示 $X$ 和 $Y$ 的复形 $X^\bullet$ 和 $Y^\bullet$。
由于 $Y^\bullet \to \tau_{\geq b}Y^\bullet$ 是拟同构，
可假设当 $j < b$ 时 $Y^j = 0$。
设 $L^\bullet \to X^\bullet$ 为任意拟同构。
则 $\tau_{\leq a}L^\bullet \to X^\bullet$ 是拟同构。
因此，$D(\mathcal{A})$ 中的态射 $X \to Y[n]$ 可表示成
$fs^{-1}$，其中 $s : L^\bullet \to X^\bullet$ 是拟同构，
$f : L^\bullet \to Y^\bullet[n]$ 是态射，并且当 $i < a$ 时
% P02-E0067: after tau_{<=a}, the frozen proof should say L^i=0 for i>a, not i<a.
$L^i = 0$。注意，$f$ 把 $L^i$ 映到 $Y^{i + n}$。
故若 $n < b - a$，则 $f = 0$，因为 $L^i$ 与 $Y^{i + n}$
总有一个为零。若 $n = b - a$，则 $f$ 恰好对应于态射
$H^a(X) \to H^b(Y)$。(2) 是 (1) 的特例。
\end{proof}

\noindent
设 $\mathcal{A}$ 为阿贝尔范畴。
假设 $0 \to A \to A' \to A'' \to 0$ 是
$\mathcal{A}$ 中对象的短正合序列。则
$0 \to A[0] \to A'[0] \to A''[0] \to 0$ 在
$D(\mathcal{A})$ 中诱导特异三角（见引理
\ref{derived-lemma-derived-canonical-delta-functor}），因而诱导 Ext 群的长正合序列
$$
0 \to \Ext^0_\mathcal{A}(B, A) \to
\Ext^0_\mathcal{A}(B, A') \to
\Ext^0_\mathcal{A}(B, A'') \to
\Ext^1_\mathcal{A}(B, A) \to \ldots
$$
类似地，给定短正合序列 $0 \to B \to B' \to B'' \to 0$，
得到 Ext 群的长正合序列
$$
0 \to \Ext^0_\mathcal{A}(B'', A) \to
\Ext^0_\mathcal{A}(B', A) \to
\Ext^0_\mathcal{A}(B, A) \to
\Ext^1_\mathcal{A}(B'', A) \to \ldots
$$
可以把这些 Ext 群视为导出范畴构造的一种应用。它表明，无须存在足够多
内射对象或投射对象，也能定义 Ext 群并构造 Ext 群的长正合序列。
Yoneda 给出了避免使用导出范畴的另一种 Ext 群构造；见
\cite{Yoneda}。

\begin{definition}
\label{derived-definition-yoneda-extension}
设 $\mathcal{A}$ 为阿贝尔范畴。设 $A, B \in \Ob(\mathcal{A})$。
当 $i \geq 1$ 时，$B$ 被 $A$ 扩张的次数为 $i$ 的
{\it Yoneda 扩张}是 $\mathcal{A}$ 中的正合序列
$$
E : 0 \to A \to Z_{i - 1} \to Z_{i - 2} \to \ldots \to Z_0 \to B \to 0
$$
若存在交换图
$$
\xymatrix{
0 \ar[r] & A \ar[r] & Z_{i - 1} \ar[r] & \ldots \ar[r] &
Z_0 \ar[r] & B \ar[r] & 0 \\
0 \ar[r] &
A \ar[r] \ar[u]^{\text{id}} \ar[d]_{\text{id}} &
Z''_{i - 1} \ar[r] \ar[u] \ar[d] &
\ldots \ar[r] &
Z''_0 \ar[r] \ar[u] \ar[d] &
B \ar[r] \ar[u]_{\text{id}} \ar[d]^{\text{id}} & 0 \\
0 \ar[r] & A \ar[r] & Z'_{i - 1} \ar[r] & \ldots \ar[r] &
Z'_0 \ar[r] & B \ar[r] & 0
}
$$
且中间一行也是 Yoneda 扩张，则称同次数的两个 Yoneda 扩张
$E$ 和 $E'$ {\it 等价}。
\end{definition}

\noindent
定义中的等价性是否为等价关系并非立即清楚。直接证明这一点颇有启发性，
不过它也将由下面的引理 \ref{derived-lemma-yoneda-extension} 得出。

\medskip\noindent
设 $\mathcal{A}$ 为具有对象 $A$、$B$ 的阿贝尔范畴。
给定 Yoneda 扩张
$E : 0 \to A \to Z_{i - 1} \to Z_{i - 2} \to \ldots \to Z_0 \to B \to 0$，
我们把相应的元素 $\delta(E) \in \Ext^i_\mathcal{A}(B, A)$ 定义为态射
$\delta(E) = fs^{-1} : B[0] \to A[i]$，其中 $s$ 是拟同构
$$
(\ldots \to 0 \to A \to Z_{i - 1} \to \ldots \to Z_0 \to 0 \to \ldots)
\longrightarrow
B[0]
$$
而 $f$ 是复形态射
$$
(\ldots \to 0 \to A \to Z_{i - 1} \to \ldots \to Z_0 \to 0 \to \ldots)
\longrightarrow
A[i]
$$
称 $\delta(E) = fs^{-1}$ 为该 Yoneda 扩张的{\it 类}。
结果表明，这个类刻画 Yoneda 扩张的等价类。

\begin{lemma}
\label{derived-lemma-yoneda-extension}
设 $\mathcal{A}$ 为具有对象 $A$、$B$ 的阿贝尔范畴，并设
$i \geq 1$。$\Ext^i_\mathcal{A}(B, A)$ 中任意元素都等于某个
$B$ 被 $A$ 扩张的次数为 $i$ 的 Yoneda 扩张的 $\delta(E)$。
给定两个同次数的 Yoneda 扩张 $E$、$E'$，则 $E$ 与 $E'$ 等价，
当且仅当 $\delta(E) = \delta(E')$。
\end{lemma}

\begin{proof}
设 $\xi : B[0] \to A[i]$ 为 $\Ext^i_\mathcal{A}(B, A)$ 的元素。
可以把 $\xi = f s^{-1}$ 写成某个拟同构
$s : L^\bullet \to B[0]$ 和映射 $f : L^\bullet \to A[i]$。
把 $L^\bullet$ 替换为 $\tau_{\leq 0}L^\bullet$ 后，
可假设当 $j > 0$ 时 $L^j = 0$。图为
$$
\xymatrix{
L^{- i - 1} \ar[r] & L^{-i} \ar[r] \ar[d] & \ldots \ar[r] &
L^0 \ar[r] & B \ar[r] & 0 \\
& A
}
$$
令 $Z_{i - 1} = (L^{- i + 1} \oplus A)/L^{-i}$，并且对
$j = i - 2, \ldots, 0$ 令 $Z_j = L^{-j}$，便得到 $B$ 被 $A$
扩张的次数为 $i$ 的扩张 $E$，其类 $\delta(E)$ 等于 $\xi$。

\medskip\noindent
由定义立即可知，等价的 Yoneda 扩张具有相同的类。假设
$E : 0 \to A \to Z_{i - 1} \to Z_{i - 2} \to \ldots \to Z_0 \to B \to 0$
和
$E' : 0 \to A \to Z'_{i - 1} \to Z'_{i - 2} \to \ldots \to Z'_0 \to B \to 0$
是具有相同类的 Yoneda 扩张。
由把 $D(\mathcal{A})$ 构造成 $K(\mathcal{A})$ 关于拟同构集合的局部化，
这意味着存在复形 $L^\bullet$ 以及拟同构
$$
t : L^\bullet \to
(\ldots \to 0 \to A \to Z_{i - 1} \to \ldots \to Z_0 \to 0 \to \ldots)
$$
和
$$
t' : L^\bullet \to
(\ldots \to 0 \to A \to Z'_{i - 1} \to \ldots \to Z'_0 \to 0 \to \ldots)
$$
使得 $s \circ t = s' \circ t'$ 且 $f \circ t = f' \circ t'$；见
范畴，第 \ref{categories-section-localization} 节。
设 $E''$ 为从证明第一段中的一对映射
$L^\bullet \to B[0]$ 和 $L^\bullet \to A[i]$ 构造出的
$B$ 被 $A$ 扩张的次数为 $i$ 的扩张。读者容易看出，存在定义等价
Yoneda 扩张时所述的次数为 $i$ 的 Yoneda 扩张“态射”
$E'' \to E$ 和 $E'' \to E'$（细节略去）。这完成证明。
\end{proof}

\begin{lemma}
\label{derived-lemma-ext-1}
设 $\mathcal{A}$ 为阿贝尔范畴。设 $A$、$B$ 为 $\mathcal{A}$ 的对象。
则 $\Ext^1_\mathcal{A}(B, A)$ 是同调，定义
\ref{homology-definition-ext-group} 中构造的群
$\Ext_\mathcal{A}(B, A)$。
\end{lemma}

\begin{proof}
这是引理 \ref{derived-lemma-yoneda-extension} 中 $i = 1$ 的情形。
\end{proof}

\noindent
给定阿贝尔范畴 $\mathcal{A}$ 以及 $D(\mathcal{A})$ 的对象
$X, Y, Z$，存在双线性且结合的复合律
$$
\Ext^j_\mathcal{A}(Y, Z) \times \Ext^i_\mathcal{A}(X, Y)
\longrightarrow
\Ext^{i + j}_\mathcal{A}(X, Z),\quad
(\eta, \xi) \longmapsto \eta \circ \xi
$$
具体地，若 $\xi : X \to Y[i]$ 且 $\eta : Y \to Z[j]$，
则把 $\eta \circ \xi$ 定义为 $\xi$ 与
$\eta[i] : Y[i] \to Z[i + j]$ 的复合。
若 $A, B, C$ 是 $\mathcal{A}$ 的对象且 $i, j \geq 1$，
则这个复合律
$$
\Ext^j_\mathcal{A}(B, C) \times \Ext^i_\mathcal{A}(A, B)
\longrightarrow
\Ext^{i + j}_\mathcal{A}(A, C)
$$
可以用 Yoneda 扩张描述如下：扩张
$$
0 \to A \to Z_{i - 1} \to Z_{i - 2} \to \ldots \to Z_0 \to B \to 0
$$
与
$$
0 \to B \to Z'_{j - 1} \to Z'_{j - 2} \to \ldots \to Z'_0 \to C \to 0
$$
的复合是 Yoneda 扩张
$$
0 \to A \to Z_{i - 1} \to Z_{i - 2} \to \ldots \to Z_0 \to 
Z'_{j - 1} \to Z'_{j - 2} \to \ldots \to Z'_0 \to C \to 0
$$

\begin{lemma}
\label{derived-lemma-cup-ext-1-zero}
设 $\mathcal{A}$ 为阿贝尔范畴。设
$0 \to A \to Z \to B \to 0$ 和
$0 \to B \to Z' \to C \to 0$ 为 $\mathcal{A}$ 中的短正合序列。
以 $[Z] \in \Ext^1_\mathcal{A}(B, A)$ 和
$[Z'] \in \Ext^1_\mathcal{A}(C, B)$ 表示它们的类。
则 $[Z] \circ [Z'] \in \Ext^2_\mathcal{A}(C, A)$ 为 $0$，
当且仅当存在交换图
$$
\xymatrix{
&
&
0 \ar[d] &
0 \ar[d]
\\
0 \ar[r] &
A \ar[r] \ar[d]^1 &
Z \ar[r] \ar[d] &
B \ar[r] \ar[d] &
0 \\
0 \ar[r] &
A \ar[r]  &
W \ar[r] \ar[d] &
Z' \ar[r] \ar[d] &
0 \\
&
&
C \ar[r]^1  \ar[d]&
C \ar[d]\\
&
&
0 &
0 
}
$$
其各行各列在 $\mathcal{A}$ 中正合。
\end{lemma}

\begin{proof}
略。提示：可以应用引理 \ref{derived-lemma-yoneda-extension}，并具体分析
$2$-扩张的类为零意味着什么。
或者，可以利用如下事实：若
$[Z] \circ [Z'] \in \Ext^2_\mathcal{A}(C, A)$ 为零，则由
$\Ext$ 的上同调长正合序列，元素
$[Z] \in \Ext^1_\mathcal{A}(B, A)$ 是
$\Ext^1_\mathcal{A}(Z', A)$ 中某个元素的像。
\end{proof}

\begin{lemma}
\label{derived-lemma-higher-ext-zero}
设 $\mathcal{A}$ 为阿贝尔范畴，并设 $p \geq 0$。
若对 $\mathcal{A}$ 中任意一对对象 $A$、$B$，都有
$\Ext^p_\mathcal{A}(B, A) = 0$，则对 $i \geq p$ 以及
$\mathcal{A}$ 中任意一对对象 $A$、$B$，都有
$\Ext^i_\mathcal{A}(B, A) = 0$。
\end{lemma}

\begin{proof}
当 $i > p$ 时，把任意类 $\xi$ 写成 $\delta(E)$，其中 $E$ 是
Yoneda 扩张
$$
E : 0 \to A \to Z_{i - 1} \to Z_{i - 2} \to \ldots \to Z_0 \to B \to 0
$$
由引理 \ref{derived-lemma-yoneda-extension}，这是可行的。
令 $C = \Ker(Z_{p - 1} \to Z_{p - 2}) = \Im(Z_p \to Z_{p - 1})$。
则 $\delta(E)$ 是 $\delta(E')$ 与 $\delta(E'')$ 的复合，其中
$$
E' : 0 \to C \to Z_{p - 1} \to \ldots \to Z_0 \to B \to 0
$$
和
$$
E'' : 0 \to A \to Z_{i - 1} \to Z_{i - 2} \to \ldots \to Z_p \to C \to 0
$$
由于 $\delta(E') \in \Ext^p_\mathcal{A}(B, C) = 0$，结论得证。
\end{proof}

\begin{lemma}
\label{derived-lemma-ext-2-zero-pre}
设 $\mathcal{A}$ 为阿贝尔范畴，并设 $K$ 为 $D^b(\mathcal{A})$ 的对象，
使得对所有 $p \geq 2$ 和 $i > j$，都有
$\Ext^p_\mathcal{A}(H^i(K), H^j(K)) = 0$。
则 $K$ 同构于其上同调的直和：
$K \cong \bigoplus H^i(K)[-i]$。
\end{lemma}

\begin{proof}
选取 $a, b$，使得当 $i \not \in [a, b]$ 时 $H^i(K) = 0$。
我们对 $b - a$ 归纳证明该引理。若 $b - a \leq 0$，结论显然。
若 $b - a > 0$，则考察截断的特异三角
$$
\tau_{\leq b - 1}K \to K \to H^b(K)[-b] \to (\tau_{\leq b - 1}K)[1]
$$
见评注 \ref{derived-remark-truncation-distinguished-triangle}。
由引理 \ref{derived-lemma-split}，若最后一个箭头为零，则
$K \cong \tau_{\leq b - 1}K \oplus H^b(K)[-b]$，归纳即得结论。
再次使用归纳可见
$$
\Hom_{D(\mathcal{A})}(H^b(K)[-b], (\tau_{\leq b - 1}K)[1]) =
\bigoplus\nolimits_{i < b} \Ext_\mathcal{A}^{b - i + 1}(H^b(K), H^i(K))
$$
由假设，该直和为零，证明完成。
\end{proof}

\begin{lemma}
\label{derived-lemma-ext-2-zero}
设 $\mathcal{A}$ 为阿贝尔范畴。假设对 $\mathcal{A}$ 中任意一对对象
$A$、$B$，都有 $\Ext^2_\mathcal{A}(B, A) = 0$。
则 $D^b(\mathcal{A})$ 的任意对象 $K$ 都同构于其上同调的直和：
$K \cong \bigoplus H^i(K)[-i]$。
\end{lemma}

\begin{proof}
由引理 \ref{derived-lemma-higher-ext-zero}，该假设蕴含：对 $i \geq 2$ 以及
$\mathcal{A}$ 中任意一对对象 $A, B$，都有
$\Ext^i_\mathcal{A}(B, A) = 0$。因此本引理是引理
\ref{derived-lemma-ext-2-zero-pre} 的特例。
\end{proof}






\section{K-群}
\label{derived-section-K-groups}

\noindent
简要讨论三角范畴的 $K_0$。

\begin{definition}
\label{derived-definition-K-zero}
设 $\mathcal{D}$ 为三角范畴。以 $K_0(\mathcal{D})$ 表示
{\it $\mathcal{D}$ 的第零 $K$-群}。它是按如下方式构造的阿贝尔群：
取 $\mathcal{D}$ 的对象所生成的自由阿贝尔群，并对每个特异三角
$X \to Y \to Z$ 加上关系 $[Y] - [X] - [Z] = 0$。
% P02-E0068: frozen prose says “objects on D” rather than “objects of D”.
\end{definition}

\noindent
注意，这蕴含 $[X[n]] = (-1)^n[X]$，因为存在特异三角
$(X, 0, X[1], 0, 0, -\text{id}[1])$。

\begin{lemma}
\label{derived-lemma-K-bounded-derived}
设 $\mathcal{A}$ 为阿贝尔范畴。则第零 $K$-群之间有典范等同
$K_0(D^b(\mathcal{A})) = K_0(\mathcal{A})$。
\end{lemma}

\begin{proof}
给定 $\mathcal{A}$ 的对象 $A$，以 $A[0]$ 表示视为集中在次数 $0$
的复形的对象 $A$。
若 $0 \to A \to A' \to A'' \to 0$ 是短正合序列，
则得到特异三角 $A[0] \to A'[0] \to A''[0] \to A[1]$；见
第 \ref{derived-section-canonical-delta-functor} 节。
这表明，把 $[A]$ 送到 $[A[0]]$ 给出映射
$K_0(\mathcal{A}) \to K_0(D^b(\mathcal{A}))$；请原谅这个糟糕的记号。

\medskip\noindent
另一方面，给定 $D^b(\mathcal{A})$ 的对象 $X$，可考虑元素
$$
c(X) = \sum (-1)^i[H^i(X)] \in K_0(\mathcal{A})
$$
给定特异三角 $X \to Y \to Z$，上同调长正合序列
(\ref{derived-equation-long-exact-cohomology-sequence-D}) 以及
$K_0(\mathcal{A})$ 中的关系表明 $c(Y) = c(X) + c(Z)$。
所以 $c$ 经由映射
$c : K_0(D^b(\mathcal{A})) \to K_0(\mathcal{A})$ 分解。

\medskip\noindent
我们要证明上述两个映射互为逆。显然，复合
$K_0(\mathcal{A}) \to
K_0(D^b(\mathcal{A})) \to K_0(\mathcal{A})$ 是恒等映射。
假设 $X^\bullet$ 是 $\mathcal{A}$ 的有界复形。
“愚蠢截断”的特异三角（见同调，第
\ref{homology-section-truncations} 节）
$$
\sigma_{\geq n}X^\bullet \to \sigma_{\geq n - 1}X^\bullet \to
X^{n - 1}[-n + 1] \to (\sigma_{\geq n}X^\bullet)[1]
$$
的存在性和归纳表明，在 $K_0(D^b(\mathcal{A}))$ 中
$$
[X^\bullet] = \sum (-1)^i[X^i[0]]
$$
（再次请原谅这个记号）。于是复合
$K_0(\mathcal{A}) \to K_0(D^b(\mathcal{A}))$ 是满射，证明完成。
\end{proof}

\begin{lemma}
\label{derived-lemma-map-K}
设 $F : \mathcal{D} \to \mathcal{D}'$ 为三角范畴的正合函子。
则 $F$ 诱导群同态 $K_0(\mathcal{D}) \to K_0(\mathcal{D}')$。
\end{lemma}

\begin{proof}
略。
\end{proof}

\begin{lemma}
\label{derived-lemma-homological-map-K}
设 $H : \mathcal{D} \to \mathcal{A}$ 为从三角范畴到阿贝尔范畴的
同调函子。假设对 $\mathcal{D}$ 中任意 $X$，对象
$H(X[i])$ 中只有有限多个在 $\mathcal{A}$ 中非零。
则 $H$ 诱导群同态 $K_0(\mathcal{D}) \to K_0(\mathcal{A})$，
它把 $[X]$ 送到 $\sum (-1)^i[H(X[i])]$。
\end{lemma}

\begin{proof}
略。
\end{proof}

\begin{lemma}
\label{derived-lemma-DBA-map-K}
设 $\mathcal{B}$ 为阿贝尔范畴 $\mathcal{A}$ 的弱 Serre 子范畴。
存在典范同构
$$
K_0(\mathcal{B}) \longrightarrow
K_0(D^b_\mathcal{B}(\mathcal{A})),\quad
[B] \longmapsto [B[0]]
$$
其逆把 $X$ 的类 $[X]$ 送到元素 $\sum (-1)^i[H^i(X)]$。
\end{lemma}

\begin{proof}
我们略去验证逆映射的规则给出良定义映射
$K_0(D^b_\mathcal{B}(\mathcal{A})) \to K_0(\mathcal{B})$。
显然，复合
$K_0(\mathcal{B}) \to
K_0(D^b_\mathcal{B}(\mathcal{A})) \to
K_0(\mathcal{B})$ 是恒等映射。另一方面，利用评注
\ref{derived-remark-truncation-distinguished-triangle} 中的特异三角和归纳论证，
读者可证明引理陈述中所示箭头是满射（细节略去）。引理得证。
\end{proof}

\begin{lemma}
\label{derived-lemma-bilinear-map-K}
设 $\mathcal{D}$、$\mathcal{D}'$、$\mathcal{D}''$ 为三角范畴。
设
$$
\otimes : \mathcal{D} \times \mathcal{D}' \longrightarrow \mathcal{D}''
$$
为如下函子：固定 $\mathcal{D}$ 中的 $X$ 时，函子
$X \otimes - : \mathcal{D}' \to \mathcal{D}''$ 正合；固定
$\mathcal{D}'$ 中的 $X'$ 时，函子
$- \otimes X' : \mathcal{D} \to \mathcal{D}''$ 正合。
则 $\otimes$ 诱导双线性映射
$K_0(\mathcal{D}) \times K_0(\mathcal{D}') \to K_0(\mathcal{D}'')$，
它把 $([X], [X'])$ 送到 $[X \otimes X']$。
\end{lemma}

\begin{proof}
略。
\end{proof}










\section{无界复形}
\label{derived-section-unbounded}

\noindent
本节材料的一个参考文献是 \cite{Spaltenstein}。
下列引理可用于寻找无界复形的“良好”左消解。

\begin{lemma}
\label{derived-lemma-special-direct-system}
设 $\mathcal{A}$ 为阿贝尔范畴。设
$\mathcal{P} \subset \Ob(\mathcal{A})$ 为子集。
假设 $\mathcal{P}$ 包含 $0$、在（有限）直和下封闭，并且
$\mathcal{A}$ 的每个对象都是 $\mathcal{P}$ 中某个元素的商。
设 $K^\bullet$ 为复形。存在复形范畴中的交换图
$$
\xymatrix{
P_1^\bullet \ar[d] \ar[r] & P_2^\bullet \ar[d] \ar[r] & \ldots \\
\tau_{\leq 1}K^\bullet \ar[r] & \tau_{\leq 2}K^\bullet \ar[r] & \ldots
}
$$
使得
\begin{enumerate}
\item 竖直箭头是拟同构且逐项满；
\item $P_n^\bullet$ 是各项属于 $\mathcal{P}$ 的有上界复形；
\item 箭头 $P_n^\bullet \to P_{n + 1}^\bullet$ 是逐项分裂单射，
且每个余核 $P^i_{n + 1}/P^i_n$ 都是 $\mathcal{P}$ 的元素。
\end{enumerate}
\end{lemma}

\begin{proof}
我们将使用同伦范畴 $K(\mathcal{A})$ 是三角范畴这一事实；见命题
\ref{derived-proposition-homotopy-category-triangulated}。
由引理 \ref{derived-lemma-subcategory-left-resolution}，可找到逐项满的复形映射
$P_1^\bullet \to \tau_{\leq 1}K^\bullet$，它是拟同构，且
$P_1^\bullet$ 的各项属于 $\mathcal{P}$。
归纳地，只须在给定 $P_1^\bullet, \ldots, P_n^\bullet$ 后构造
$P_{n + 1}^\bullet$ 以及映射
$P_n^\bullet \to P_{n + 1}^\bullet$ 和
$P_{n + 1}^\bullet \to \tau_{\leq n + 1}K^\bullet$。

\medskip\noindent
在 $K(\mathcal{A})$ 中选取特异三角
$P_n^\bullet \to \tau_{\leq n + 1}K^\bullet \to C^\bullet \to P_n^\bullet[1]$。
应用引理 \ref{derived-lemma-subcategory-left-resolution}，选取拟同构的复形映射
$Q^\bullet \to C^\bullet$，使得 $Q^\bullet$ 的各项属于 $\mathcal{P}$。
由三角范畴的公理，可把复合
$Q^\bullet \to C^\bullet \to P_n^\bullet[1]$ 嵌入
$K(\mathcal{A})$ 中的特异三角
$P_n^\bullet \to P_{n + 1}^\bullet \to Q^\bullet \to P_n^\bullet[1]$。
由引理 \ref{derived-lemma-improve-distinguished-triangle-homotopy}，
可以并且确实假设
$0 \to P_n^\bullet \to P_{n + 1}^\bullet \to Q^\bullet \to 0$
是逐项分裂的短正合序列。这蕴含 $P_{n + 1}^\bullet$ 的各项属于
$\mathcal{P}$，并且 $P_n^\bullet \to P_{n + 1}^\bullet$ 是余核属于
$\mathcal{P}$ 的逐项分裂单射。
由三角范畴的公理，得到三角范畴 $K(\mathcal{A})$ 中的特异三角态射
$$
\xymatrix{
P_n^\bullet \ar[r] \ar[d] &
P_{n + 1}^\bullet \ar[r] \ar[d] &
Q^\bullet \ar[r] \ar[d] &
P_n^\bullet[1] \ar[d] \\
P_n^\bullet \ar[r] &
\tau_{\leq n + 1}K^\bullet \ar[r] &
C^\bullet \ar[r] &
P_n^\bullet[1]
}
$$
选取实际的复形态射
$f : P_{n + 1}^\bullet \to \tau_{\leq n + 1}K^\bullet$。
上图左方块交换至同伦；但由于
$P_n^\bullet \to P_{n + 1}^\bullet$ 是逐项分裂单射，
可以提升该同伦并修改对 $f$ 的选择，使该方块交换。
最后，$f$ 是拟同构，因为 $P_n^\bullet \to P_n^\bullet$ 和
$Q^\bullet \to C^\bullet$ 都是拟同构。

\medskip\noindent
此时已有所有所需性质，唯一尚不确定的是映射
$f : P_{n + 1}^\bullet \to \tau_{\leq n + 1}K^\bullet$
是否逐项满。由于存在复形的交换图
$$
\xymatrix{
P_n^\bullet \ar[d] \ar[r] & P_{n + 1}^\bullet \ar[d] \\
\tau_{\leq n}K^\bullet \ar[r] & \tau_{\leq n + 1}K^\bullet
}
$$
由归纳假设可知，除可能在次数 $n$ 和 $n + 1$ 之外，
$f$ 在所有次数的项上都是满射。选取对象
$P \in \mathcal{P}$ 以及满射 $q : P \to K^n$。
考虑映射
$$
g :
P^\bullet = (\ldots \to 0 \to P \xrightarrow{1} P \to 0 \to \ldots)
\longrightarrow
\tau_{\leq n + 1}K^\bullet
$$
其中第一个 $P$ 位于次数 $n$，次数 $n$ 的映射由 $q$ 给出，
次数 $n + 1$ 的映射由 $d_K \circ q$ 给出。
它在次数 $n$ 是满射，而次数 $n + 1$ 的余核是
$H^{n + 1}(\tau_{\leq n + 1}K^\bullet)$；要看到这一点，回忆
$\tau_{\leq n + 1}K^\bullet$ 在次数 $n + 1$ 的项为
$\Ker(d_K^{n + 1})$。
不过，由于 $f$ 是拟同构，我们知道 $H^{n + 1}(f)$ 是满射。
因此，把 $f : P_{n + 1}^\bullet \to \tau_{\leq n + 1}K^\bullet$
替换为
$f \oplus g : P_{n + 1}^\bullet \oplus P^\bullet \to \tau_{\leq n + 1}K^\bullet$
后即得结论。
\end{proof}

\noindent
在某些情形下，可以用上述引理证明左导出函子处处有定义。

\begin{proposition}
\label{derived-proposition-left-derived-exists}
设 $F : \mathcal{A} \to \mathcal{B}$ 为阿贝尔范畴之间的右正合函子。
设 $\mathcal{P} \subset \Ob(\mathcal{A})$ 为子集。假设
\begin{enumerate}
\item $\mathcal{P}$ 包含 $0$、在（有限）直和下封闭，并且
$\mathcal{A}$ 的每个对象都是 $\mathcal{P}$ 中某个元素的商；
\item 对 $\mathcal{A}$ 的任意有上界无环复形 $P^\bullet$，
若对所有 $n$ 有 $P^n \in \mathcal{P}$，则复形 $F(P^\bullet)$ 正合；
\item $\mathcal{A}$ 和 $\mathcal{B}$ 都有以 $\mathbf{N}$ 为指标的系统的余极限；
\item 以 $\mathbf{N}$ 为指标的余极限在 $\mathcal{A}$ 和
$\mathcal{B}$ 中都正合；并且
\item $F$ 与以 $\mathbf{N}$ 为指标的余极限交换。
\end{enumerate}
则 $LF$ 在整个 $D(\mathcal{A})$ 上有定义。
\end{proposition}

\begin{proof}
由 (1) 和引理 \ref{derived-lemma-subcategory-left-resolution}，对任意有上界复形
$K^\bullet$，存在拟同构 $P^\bullet \to K^\bullet$，其中
$P^\bullet$ 有上界，且对所有 $n$ 有 $P^n \in \mathcal{P}$。
假设 $s : P^\bullet \to (P')^\bullet$ 是由 $\mathcal{P}$ 中对象组成的
有上界复形之间的拟同构。则
$F(P^\bullet) \to F((P')^\bullet)$ 是拟同构，因为由假设 (2)，
$F(C(s)^\bullet)$ 无环。这已经表明 $LF$ 在 $D^{-}(\mathcal{A})$
上有定义，并且由 $\mathcal{P}$ 中对象组成的有上界复形计算 $LF$；见
引理 \ref{derived-lemma-find-existence-computes}。

\medskip\noindent
接着，设 $K^\bullet$ 为 $\mathcal{A}$ 的任意复形。
选取引理 \ref{derived-lemma-special-direct-system} 中的图
$$
\xymatrix{
P_1^\bullet \ar[d] \ar[r] & P_2^\bullet \ar[d] \ar[r] & \ldots \\
\tau_{\leq 1}K^\bullet \ar[r] & \tau_{\leq 2}K^\bullet \ar[r] & \ldots
}
$$
注意，映射 $\colim P_n^\bullet \to K^\bullet$ 是拟同构，
因为 $\mathcal{A}$ 中以 $\mathbf{N}$ 为指标的余极限正合，且当
$n > i$ 时 $H^i(P_n^\bullet) = H^i(K^\bullet)$。
我们断言
$$
F(\colim P_n^\bullet) = \colim F(P_n^\bullet)
$$
（逐项余极限）是 $LF(K^\bullet)$，即 $\colim P_n^\bullet$ 计算
$LF$。要看到这一点，由引理 \ref{derived-lemma-find-existence-computes}，
只须证明如下断言。假设
$$
\colim Q_n^\bullet = Q^\bullet
\xrightarrow{\ \alpha\ }
P^\bullet = \colim P_n^\bullet
$$
是复形的拟同构，使得每个 $P_n^\bullet$、$Q_n^\bullet$ 都是
各项属于 $\mathcal{P}$ 的有上界复形，并且映射
$P_n^\bullet \to \tau_{\leq n}P^\bullet$ 和
$Q_n^\bullet \to \tau_{\leq n}Q^\bullet$ 是拟同构。
断言：$F(\alpha)$ 是拟同构。

\medskip\noindent
困难在于，我们没有假设 $\alpha$ 由复形 $P_n^\bullet$ 和
$Q_n^\bullet$ 之间的映射取余极限而给出。不过，对每个 $n$，
我们知道图中的实线箭头
% P02-E0070: alpha is declared Q -> P, while the entire downstream construction consistently uses P -> Q.
$$
\xymatrix{
& R^\bullet \ar@{..>}[d] \\
P_n^\bullet \ar[d] &
L^\bullet \ar@{..>}[l] \ar@{..>}[r] &
Q_n^\bullet \ar[d] \\
\tau_{\leq n}P^\bullet \ar[rr]^{\tau_{\leq n}\alpha} & &
\tau_{\leq n}Q^\bullet
}
$$
都是拟同构。由于拟同构在 $K(\mathcal{A})$ 中构成乘法系
（见引理 \ref{derived-lemma-acyclic}），可以找到拟同构
$L^\bullet \to P_n^\bullet$ 和复形映射
$L^\bullet \to Q_n^\bullet$，使得上图交换至同伦。
于是 $\tau_{\leq n}L^\bullet \to L^\bullet$ 是拟同构。
因此（由证明的第一部分）可找到各项属于 $\mathcal{P}$ 的有上界复形
$R^\bullet$ 以及拟同构 $R^\bullet \to L^\bullet$（如图所示）。
使用证明第一段的结果，可知
$F(R^\bullet) \to F(P_n^\bullet)$ 和
$F(R^\bullet) \to F(Q_n^\bullet)$ 是拟同构。
因此得到同构
$H^i(F(P_n^\bullet)) \to H^i(F(Q_n^\bullet))$，它嵌入交换图
$$
\xymatrix{
H^i(F(P_n^\bullet)) \ar[r] \ar[d] &
H^i(F(Q_n^\bullet)) \ar[d] \\
H^i(F(P^\bullet)) \ar[r] &
H^i(F(Q^\bullet))
}
$$
完全相同的论证表明，随着 $n$ 变化，这些映射也彼此相容。
由于 (4) 和 (5) 给出
$$
H^i(F(P^\bullet)) =
H^i(F(\colim P_n^\bullet)) =
H^i(\colim F(P_n^\bullet)) = \colim H^i(F(P_n^\bullet))
$$
且对 $Q^\bullet$ 也类似，所以得出
% P02-E0071: the frozen displayed inline map has unmatched parentheses and omits F(alpha); direction follows P02-E0070.
$H^i(\alpha) : H^i(F(P^\bullet) \to H^i(F(Q^\bullet)$ 是同构，
断言随之得证。
\end{proof}

\begin{lemma}
\label{derived-lemma-special-inverse-system}
设 $\mathcal{A}$ 为阿贝尔范畴。设
$\mathcal{I} \subset \Ob(\mathcal{A})$ 为子集。
假设 $\mathcal{I}$ 包含 $0$、在（有限）积下封闭，并且
$\mathcal{A}$ 的每个对象都是 $\mathcal{I}$ 中某个元素的子对象。
设 $K^\bullet$ 为复形。存在复形范畴中的交换图
$$
\xymatrix{
\ldots \ar[r] &
\tau_{\geq -2}K^\bullet \ar[r] \ar[d] &
\tau_{\geq -1}K^\bullet \ar[d] \\
\ldots \ar[r] & I_2^\bullet \ar[r] & I_1^\bullet
}
$$
使得
\begin{enumerate}
\item 竖直箭头是拟同构且逐项单；
\item $I_n^\bullet$ 是各项属于 $\mathcal{I}$ 的有下界复形；
\item 箭头 $I_{n + 1}^\bullet \to I_n^\bullet$ 是逐项分裂满射，
且 $\Ker(I^i_{n + 1} \to I^i_n)$ 是 $\mathcal{I}$ 的元素。
\end{enumerate}
\end{lemma}

\begin{proof}
本引理对偶于引理 \ref{derived-lemma-special-direct-system}。
\end{proof}







\section{导出伴随函子}
\label{derived-section-deriving-adjoints}

\noindent
设 $F : \mathcal{D} \to \mathcal{D}'$ 和
$G : \mathcal{D}' \to \mathcal{D}$ 为三角范畴的正合函子。
设 $S$（相应地，$S'$）为 $\mathcal{D}$（相应地，$\mathcal{D}'$）
中与三角结构相容的乘法系。
以 $Q : \mathcal{D} \to S^{-1}\mathcal{D}$ 和
$Q' : \mathcal{D}' \to (S')^{-1}\mathcal{D}'$ 表示局部化函子。
在此情形下，滥用记号，常以 $RF$ 表示对应于
$Q' \circ F : \mathcal{D} \to (S')^{-1}\mathcal{D}'$ 和乘法系 $S$
的部分定义的右导出函子。类似地，以 $LG$ 表示对应于
$Q \circ G : \mathcal{D}' \to S^{-1}\mathcal{D}$ 和乘法系 $S'$
的部分定义的左导出函子。图为
$$
\vcenter{
\xymatrix{
\mathcal{D} \ar[r]_F \ar[d]_Q & \mathcal{D}' \ar[d]^{Q'} \\
S^{-1}\mathcal{D} \ar@{..>}[r]^{RF} &
(S')^{-1}\mathcal{D}'
}
}
\quad\text{且}\quad
\vcenter{
\xymatrix{
\mathcal{D}' \ar[r]_G \ar[d]_{Q'} & \mathcal{D} \ar[d]^Q \\
(S')^{-1}\mathcal{D}' \ar@{..>}[r]^{LG} &
S^{-1}\mathcal{D}
}
}
$$

\begin{lemma}
\label{derived-lemma-pre-derived-adjoint-functors-general}
在上述情形中，假设 $F$ 是 $G$ 的右伴随。设
$K \in \Ob(\mathcal{D})$ 且 $M \in \Ob(\mathcal{D}')$。
若 $RF$ 在 $K$ 处有定义，且 $LG$ 在 $M$ 处有定义，
则存在典范同构
$$
\Hom_{(S')^{-1}\mathcal{D}'}(M, RF(K)) =
\Hom_{S^{-1}\mathcal{D}}(LG(M), K)
$$
在 $RF$ 和 $LG$ 有定义的 $S^{-1}\mathcal{D}$ 与
$(S')^{-1}\mathcal{D}'$ 的三角子范畴上，该同构关于两个变量都具有函子性。
\end{lemma}

\begin{proof}
由于 $RF$ 在 $K$ 处有定义，把 $S$ 中的
$s : K \to I$ 送到对象 $F(I)$ 的规则，作为
$(S')^{-1}\mathcal{D}'$ 的 ind-对象，本质恒定且值为 $RF(K)$。
类似地，把 $S'$ 中的 $t : P \to M$ 送到对象 $G(P)$ 的规则，
作为 $S^{-1}\mathcal{D}$ 的 pro-对象，本质恒定且值为 $LG(M)$。
因此有
\begin{align*}
\Hom_{(S')^{-1}\mathcal{D}'}(M, RF(K))
& =
\colim_{s : K \to I} \Hom_{(S')^{-1}\mathcal{D}'}(M, F(I)) \\
& =
\colim_{s : K \to I} \colim_{t : P \to M} \Hom_{\mathcal{D}'}(P, F(I)) \\
& =
\colim_{t : P \to M} \colim_{s : K \to I} \Hom_{\mathcal{D}'}(P, F(I)) \\
& =
\colim_{t : P \to M} \colim_{s : K \to I} \Hom_{\mathcal{D}}(G(P), I) \\
& =
\colim_{t : P \to M} \Hom_{S^{-1}\mathcal{D}}(G(P), K) \\
& =
\Hom_{S^{-1}\mathcal{D}}(LG(M), K)
\end{align*}
第一个等号由范畴，引理
\ref{categories-lemma-characterize-essentially-constant-ind} 成立。
% P02-E0072: the generalized proof refers to D(B), although no abelian category B is present here.
第二个等号由 $D(\mathcal{B})$ 中态射的定义成立；见范畴，评注
\ref{categories-remark-right-localization-morphisms-colimit}。
第三个等号由范畴，引理
\ref{categories-lemma-colimits-commute} 成立。
第四个等号成立是因为 $F$ 与 $G$ 伴随。
% P02-E0073: similarly, the generalized proof refers to D(A), although only S^{-1}D is defined.
第五个等号由 $D(\mathcal{A})$ 中态射的定义成立；见范畴，评注
\ref{categories-remark-left-localization-morphisms-colimit}。
第六个等号由范畴，引理
\ref{categories-lemma-characterize-essentially-constant-pro} 成立。
我们略去函子性的证明。
\end{proof}

\begin{lemma}
\label{derived-lemma-pre-derived-adjoint-functors}
设 $F : \mathcal{A} \to \mathcal{B}$ 和
$G : \mathcal{B} \to \mathcal{A}$ 为阿贝尔范畴的函子，
并且 $F$ 是 $G$ 的右伴随。设 $K^\bullet$ 为 $\mathcal{A}$ 的复形，
$M^\bullet$ 为 $\mathcal{B}$ 的复形。若 $RF$ 在 $K^\bullet$ 处有定义，
且 $LG$ 在 $M^\bullet$ 处有定义，则存在典范同构
$$
\Hom_{D(\mathcal{B})}(M^\bullet, RF(K^\bullet)) =
\Hom_{D(\mathcal{A})}(LG(M^\bullet), K^\bullet)
$$
在 $RF$ 和 $LG$ 有定义的 $D(\mathcal{A})$ 与 $D(\mathcal{B})$
的三角子范畴上，该同构关于两个变量都具有函子性。
\end{lemma}

\begin{proof}
这是非常一般的引理
\ref{derived-lemma-pre-derived-adjoint-functors-general} 的特例。
\end{proof}

\noindent
下列引理说明了为何使用无界导出范畴更为容易：若没有无界导出函子，
甚至无法陈述该引理。

\begin{lemma}
\label{derived-lemma-derived-adjoint-functors}
设 $F : \mathcal{A} \to \mathcal{B}$ 和
$G : \mathcal{B} \to \mathcal{A}$ 为阿贝尔范畴的函子，
并且 $F$ 是 $G$ 的右伴随。若导出函子
$RF : D(\mathcal{A}) \to D(\mathcal{B})$ 和
$LG : D(\mathcal{B}) \to D(\mathcal{A})$ 存在，则
$RF$ 是 $LG$ 的右伴随。
\end{lemma}

\begin{proof}
立即由引理 \ref{derived-lemma-pre-derived-adjoint-functors} 得出。
\end{proof}





\section{K-内射复形}
\label{derived-section-K-injective}

\noindent
下列类型的复形可用于计算无界导出范畴上的右导出函子。

\begin{definition}
\label{derived-definition-K-injective}
设 $\mathcal{A}$ 为阿贝尔范畴。若对每个无环复形 $M^\bullet$ 都有
$\Hom_{K(\mathcal{A})}(M^\bullet, I^\bullet) = 0$，
则称复形 $I^\bullet$ 为{\it K-内射的}。
\end{definition}

\noindent
在该定义的情形中，事实上对所有 $i$ 都有
$\Hom_{K(\mathcal{A})}(M^\bullet[i], I^\bullet) = 0$，
因为无环复形的平移仍无环。

\begin{lemma}
\label{derived-lemma-K-injective}
设 $\mathcal{A}$ 为阿贝尔范畴。设 $I^\bullet$ 为复形。
下列条件等价：
\begin{enumerate}
\item $I^\bullet$ 是 K-内射的；
\item 对每个拟同构 $M^\bullet \to N^\bullet$，映射
$$
\Hom_{K(\mathcal{A})}(N^\bullet, I^\bullet)
\to \Hom_{K(\mathcal{A})}(M^\bullet, I^\bullet)
$$
是双射；并且
\item 对每个复形 $N^\bullet$，映射
$$
\Hom_{K(\mathcal{A})}(N^\bullet, I^\bullet)
\to \Hom_{D(\mathcal{A})}(N^\bullet, I^\bullet)
$$
是同构。
\end{enumerate}
\end{lemma}

\begin{proof}
假设 (1)。则 (2) 成立，因为函子
$\Hom_{K(\mathcal{A})}( - , I^\bullet)$ 是上同调函子，
而拟同构的锥无环。

\medskip\noindent
假设 (2)。$D(\mathcal{A})$ 中的态射
$N^\bullet \to I^\bullet$ 具有形式
$fs^{-1} : N^\bullet \to I^\bullet$，其中
$s : M^\bullet \to N^\bullet$ 是拟同构，而
$f : M^\bullet \to I^\bullet$ 是映射。由 (2)，它对应于
$K(\mathcal{A})$ 中唯一的态射 $N^\bullet \to I^\bullet$，
即 (3) 成立。

\medskip\noindent
假设 (3)。若 $M^\bullet$ 无环，则 $M^\bullet$ 在 $D(\mathcal{A})$
中同构于零复形，故
$\Hom_{D(\mathcal{A})}(M^\bullet, I^\bullet) = 0$，
从而由 (3)，
$\Hom_{K(\mathcal{A})}(M^\bullet, I^\bullet) = 0$，
即 (1) 成立。
\end{proof}

\begin{lemma}
\label{derived-lemma-triangle-K-injective}
设 $\mathcal{A}$ 为阿贝尔范畴。设 $(K, L, M, f, g, h)$ 为
$K(\mathcal{A})$ 的特异三角。若 $K$、$L$、$M$ 中有两个是
K-内射复形，则第三个也是。
\end{lemma}

\begin{proof}
由定义、引理 \ref{derived-lemma-representable-homological} 以及
$K(\mathcal{A})$ 是三角范畴这一事实得出（命题
\ref{derived-proposition-homotopy-category-triangulated}）。
\end{proof}

\begin{lemma}
\label{derived-lemma-bounded-below-injectives-K-injective}
设 $\mathcal{A}$ 为阿贝尔范畴。内射对象的有下界复形是 K-内射的。
\end{lemma}

\begin{proof}
由引理 \ref{derived-lemma-K-injective} 和
\ref{derived-lemma-morphisms-into-injective-complex} 得出。
\end{proof}

\begin{lemma}
\label{derived-lemma-product-K-injective}
设 $\mathcal{A}$ 为阿贝尔范畴。设 $T$ 为集合，并对每个
$t \in T$ 设 $I_t^\bullet$ 为 K-内射复形。
若对所有 $n$，$I^n = \prod_t I_t^n$ 都存在，则 $I^\bullet$
是 K-内射复形。此外，$I^\bullet$ 表示各对象 $I_t^\bullet$
在 $D(\mathcal{A})$ 中的积。
\end{lemma}

\begin{proof}
% P02-E0074: frozen prose reads “an complex” rather than “a complex”.
设 $K^\bullet$ 为一个复形。注意，复形
$$
C :
\prod\nolimits_b \Hom(K^{-b}, I^{b - 1}) \to
\prod\nolimits_b \Hom(K^{-b}, I^b) \to
\prod\nolimits_b \Hom(K^{-b}, I^{b + 1})
$$
在中间的上同调为
$\Hom_{K(\mathcal{A})}(K^\bullet, I^\bullet)$。
类似地，复形
$$
C_t :
\prod\nolimits_b \Hom(K^{-b}, I_t^{b - 1}) \to
\prod\nolimits_b \Hom(K^{-b}, I_t^b) \to
\prod\nolimits_b \Hom(K^{-b}, I_t^{b + 1})
$$
计算 $\Hom_{K(\mathcal{A})}(K^\bullet, I_t^\bullet)$。
接着注意，由对 $I$ 的选择，我们有阿贝尔群复形的等式
$$
C = \prod\nolimits_{t \in T} C_t
$$
取积是阿贝尔群范畴上的正合函子。因此，若 $K^\bullet$ 无环，
则 $\Hom_{K(\mathcal{A})}(K^\bullet, I_t^\bullet) = 0$，
故 $C_t$ 无环，故 $C$ 无环，进而
$\Hom_{K(\mathcal{A})}(K^\bullet, I^\bullet) = 0$。
于是 $I^\bullet$ 是 K-内射的。
由此，再用引理 \ref{derived-lemma-K-injective} 可得
$$
\Hom_{D(\mathcal{A})}(K^\bullet, I^\bullet)
=
\prod\nolimits_{t \in T} \Hom_{D(\mathcal{A})}(K^\bullet, I_t^\bullet)
$$
而且 $I^\bullet$ 的确表示导出范畴中的积。
\end{proof}

\begin{lemma}
\label{derived-lemma-K-injective-defined}
设 $\mathcal{A}$ 为阿贝尔范畴。
设 $F : K(\mathcal{A}) \to \mathcal{D}'$ 为三角范畴的正合函子。
则 $RF$ 在 $K(\mathcal{A})$ 中每个拟同构于 K-内射复形的复形处
都有定义。事实上，每个 K-内射复形都计算 $RF$。
\end{lemma}

\begin{proof}
由引理 \ref{derived-lemma-derived-inverts}，只须证明 $RF$ 在 K-内射复形处
有定义，即只须证明 K-内射复形 $I^\bullet$ 计算 $RF$。
考虑拟同构 $I^\bullet \to N^\bullet$。由引理
\ref{derived-lemma-K-injective}，它有左逆。
因此 $I^\bullet \to I^\bullet$ 是
$I^\bullet/\text{Qis}(\mathcal{A})$ 的终对象，结论得证。
\end{proof}

\begin{lemma}
\label{derived-lemma-enough-K-injectives-implies}
设 $\mathcal{A}$ 为阿贝尔范畴。
假设每个复形都有取值于某个 K-内射复形的拟同构。
则三角范畴的任意正合函子
$F : K(\mathcal{A}) \to \mathcal{D}'$ 都有右导出函子
$$
RF : D(\mathcal{A}) \longrightarrow \mathcal{D}'
$$
且对 K-内射复形 $I^\bullet$，有
$RF(I^\bullet) = F(I^\bullet)$。
\end{lemma}

\begin{proof}
为此，对 K-内射复形的类 $\mathcal{I}$ 应用引理
\ref{derived-lemma-find-existence-computes}。由于 (1) 由假设成立，只须证明：
若 $I^\bullet \to J^\bullet$ 是 K-内射复形之间的拟同构，
则 $F(I^\bullet) \to F(J^\bullet)$ 是同构。
这是显然的，因为由引理 \ref{derived-lemma-K-injective}，
$I^\bullet \to J^\bullet$ 是同伦等价，即 $K(\mathcal{A})$ 中的同构。
\end{proof}

\noindent
下列引理可推广到以更大序数为指标的极限。

\begin{lemma}
\label{derived-lemma-limit-K-injectives}
\begin{slogan}
K-内射复形的“分裂”塔的极限是 K-内射的。
\end{slogan}
设 $\mathcal{A}$ 为阿贝尔范畴。设
$$
\ldots \to I_3^\bullet \to I_2^\bullet \to I_1^\bullet
$$
为复形的逆系统。假设
\begin{enumerate}
\item 每个 $I_n^\bullet$ 都是 $K$-内射的；
\item 每个映射 $I_{n + 1}^m \to I_n^m$ 都是分裂满射；
\item 极限 $I^m = \lim I_n^m$ 存在。
\end{enumerate}
则复形 $I^\bullet$ 是 K-内射的。
\end{lemma}

\begin{proof}
我们强烈建议读者跳过本引理的证明。
设 $M^\bullet$ 为无环复形。简记
$H_n(a, b) = \Hom_\mathcal{A}(M^a, I_n^b)$。用此记号，
$\Hom_{K(\mathcal{A})}(M^\bullet, I^\bullet)$ 是复形
$$
\prod_m \lim\limits_n H_n(m, m - 2)
\to
\prod_m \lim\limits_n H_n(m, m - 1)
\to
\prod_m \lim\limits_n H_n(m, m)
\to
\prod_m \lim\limits_n H_n(m, m + 1)
$$
从左起第三处的上同调。
可以交换 $\prod$ 与 $\lim$ 的次序，而每个复形
$$
\prod_m H_n(m, m - 2)
\to
\prod_m H_n(m, m - 1)
\to
\prod_m H_n(m, m)
\to
\prod_m H_n(m, m + 1)
$$
都由假设 (1) 正合。由假设 (2)，系统中的映射
$$
\ldots \to
\prod_m H_3(m, m - 2) \to
\prod_m H_2(m, m - 2) \to
\prod_m H_1(m, m - 2)
$$
都是满射。因此本引理由同调，引理
\ref{homology-lemma-apply-Mittag-Leffler} 得出。
\end{proof}

\noindent
看来，引理 \ref{derived-lemma-special-inverse-system}、
\ref{derived-lemma-bounded-below-injectives-K-injective} 与
\ref{derived-lemma-limit-K-injectives} 的组合，会为任意具有足够多内射对象和
可数积的阿贝尔范畴产生“足够多 K-内射对象”。
事实上，这可能行不通！解释见引理
\ref{derived-lemma-difficulty-K-injectives}。

\begin{lemma}
\label{derived-lemma-adjoint-preserve-K-injectives}
设 $\mathcal{A}$ 和 $\mathcal{B}$ 为阿贝尔范畴。
设 $u : \mathcal{A} \to \mathcal{B}$ 和
$v : \mathcal{B} \to \mathcal{A}$ 为加性函子。假设
\begin{enumerate}
\item $u$ 是 $v$ 的右伴随；并且
\item $v$ 正合。
\end{enumerate}
则 $u$ 把 K-内射复形变为 K-内射复形。
\end{lemma}

\begin{proof}
设 $I^\bullet$ 为 $\mathcal{A}$ 的 K-内射复形。
% P02-E0075: frozen prose reads “a acyclic complex”.
设 $M^\bullet$ 为 $\mathcal{B}$ 的无环复形。
由于 $v$ 正合，$v(M^\bullet)$ 是无环复形。由伴随性得到
$$
0 = \Hom_{K(\mathcal{A})}(v(M^\bullet), I^\bullet) =
\Hom_{K(\mathcal{B})}(M^\bullet, u(I^\bullet))
$$
故引理得证。
\end{proof}





\section{有界上同调维数}
\label{derived-section-bounded}

\noindent
还有另一种无界导出函子存在的情形，即函子具有有界上同调维数时。

\begin{lemma}
\label{derived-lemma-replace-resolution}
设 $\mathcal{A}$ 为阿贝尔范畴。设
$d : \Ob(\mathcal{A}) \to \{0, 1, 2, \ldots, \infty\}$ 为函数。
假设
\begin{enumerate}
\item $\mathcal{A}$ 的每个对象都是某个满足 $d(A) = 0$ 的对象
$A$ 的子对象；
\item 对 $A, B \in \mathcal{A}$，有
$d(A \oplus B) \leq \max \{d(A), d(B)\}$；并且
\item 若 $0 \to A \to B \to C \to 0$ 短正合，则
$d(C) \leq \max\{d(A) - 1, d(B)\}$。
\end{enumerate}
设 $K^\bullet$ 为复形，使得当 $n \to -\infty$ 时，
$n + d(K^n)$ 趋于 $-\infty$。则存在拟同构
$K^\bullet \to L^\bullet$，使得对所有 $n \in \mathbf{Z}$，
都有 $d(L^n) = 0$。
\end{lemma}

\begin{proof}
由引理 \ref{derived-lemma-subcategory-right-resolution}，可找到拟同构
$\sigma_{\geq 0}K^\bullet \to M^\bullet$，其中当 $n < 0$ 时
$M^n = 0$，而当 $n \geq 0$ 时 $d(M^n) = 0$。于是 $K^\bullet$
拟同构于复形
$$
\ldots \to K^{-2} \to K^{-1} \to M^0 \to M^1 \to \ldots
$$
因此可假设当 $n \gg 0$ 时 $d(K^n) = 0$。
注意，这一替换不破坏当 $n \to -\infty$ 时
$n + d(K^n) \to -\infty$ 的条件。

\medskip\noindent
我们将通过（无限）一列基本替换来改进 $K^\bullet$。
所谓{\it 基本替换}如下。选取满足 $d(K^n) > 0$ 的指标 $n$。
选取单射 $K^n \to M$，其中 $d(M) = 0$。
令 $M' = \Coker(K^n \to M \oplus K^{n + 1})$。
考虑复形映射
$$
\xymatrix{
K^\bullet : \ar[d] &
K^{n - 1} \ar[d] \ar[r] &
K^n \ar[d] \ar[r] &
K^{n + 1} \ar[d] \ar[r] &
K^{n + 2} \ar[d] \\
(K')^\bullet : &
K^{n - 1} \ar[r] &
M \ar[r] &
M' \ar[r] &
K^{n + 2}
}
$$
显然 $K^\bullet \to (K')^\bullet$ 是拟同构。
此外，显然 $d((K')^n) = 0$，并且
$$
d((K')^{n + 1}) \leq \max\{d(K^n) - 1, d(M \oplus K^{n + 1})\} \leq
\max\{d(K^n) - 1, d(K^{n + 1})\}
$$
其他值不变。

\medskip\noindent
为完成证明，我们仔细选择进行基本替换的顺序，使得对每个整数 $m$，
复形 $\sigma_{\geq m}K^\bullet$ 只被改变有限多次。为此令
$$
\xi(K^\bullet) = \max \{n + d(K^n) \mid d(K^n) > 0\}
$$
并令
$$
I = \{n \in \mathbf{Z} \mid \xi(K^\bullet) = n + d(K^n)
\text{ 且 }
 d(K^n) > 0\}
$$
我们的假设：当 $n \to -\infty$ 时 $n + d(K^n)$ 趋于
$-\infty$，以及当 $n >> 0$ 时 $d(K^n) = 0$ 这一事实，
蕴含 $\xi(K^\bullet) < +\infty$ 且 $I$ 是有限集。
显然，对上述基本变换，有
$\xi((K')^\bullet) \leq \xi(K^\bullet)$。
一次基本变换在次数 $\leq \xi(K^\bullet) + 1$ 处改变复形。
因此，若能找到降低 $\xi(K^\bullet)$ 的有限列基本变换，结论即得。
% P02-E0076: frozen prose says “find finite sequence” and omits the article.
不过，注意，若从最小元素 $n \in I$ 开始进行基本变换，
则或者 $I$ 的大小减小，或者 $\min I$ 增大。
由于 $I$ 的每个元素都 $\leq \xi(K^\bullet)$，经过有限步便得结论。
\end{proof}

\begin{lemma}
\label{derived-lemma-unbounded-right-derived}
设 $F : \mathcal{A} \to \mathcal{B}$ 为阿贝尔范畴之间的左正合函子。
假设
\begin{enumerate}
\item $\mathcal{A}$ 的每个对象都是某个对 $F$ 右无环的对象的子对象；
\item 存在整数 $n \geq 0$，使得 $R^nF = 0$。
\end{enumerate}
则
\begin{enumerate}
\item $RF : D(\mathcal{A}) \to D(\mathcal{B})$ 存在；
\item 由对 $F$ 右无环的对象组成的任意复形都计算 $RF$；
\item 任意复形都是某个取值于由对 $F$ 右无环对象组成的复形中的
拟同构的源；
\item 对 $E \in D(\mathcal{A})$：
\begin{enumerate}
% P02-E0077: the frozen map omits a closing parenthesis after RF(tau_{<=a}E).
\item 当 $i \leq a$ 时，
$H^i(RF(\tau_{\leq a}E) \to H^i(RF(E))$ 是同构；
\item 当 $i \geq b$ 时，
$H^i(RF(E)) \to H^i(RF(\tau_{\geq b - n + 1}E))$ 是同构；
\item 若对某个 $-\infty \leq a \leq b \leq \infty$，
当 $i \not \in [a, b]$ 时 $H^i(E) = 0$，则当
$i \not \in [a, b + n - 1]$ 时 $H^i(RF(E)) = 0$。
\end{enumerate}
\end{enumerate}
\end{lemma}

\begin{proof}
注意，第一个假设蕴含
$RF : D^+(\mathcal{A}) \to D^+(\mathcal{B})$ 存在；见命题
\ref{derived-proposition-enough-acyclics}。
设 $A$ 为 $\mathcal{A}$ 的对象。选取单射 $A \to A'$，其中
$A'$ 无环。则由上同调长正合序列，
$R^{n + 1}F(A) = R^nF(A'/A) = 0$。
因此 $R^{n + 1}F = 0$。继续如此归纳，得到对所有
$m \geq n$ 都有 $R^mF = 0$。

\medskip\noindent
我们将对函数
$d : \Ob(\mathcal{A}) \to \{0, 1, 2, \ldots \}$ 应用引理
\ref{derived-lemma-replace-resolution}，其中
% P02-E0078: frozen d(A) is malformed as max {0} union {i | ...}, rather than max of the union.
$d(A) = \max \{0\} \cup \{i \mid R^iF(A) \not = 0\}$。
引理 \ref{derived-lemma-replace-resolution} 的第一个假设就是我们的假设 (1)。
引理 \ref{derived-lemma-replace-resolution} 的第二个假设由
$RF(A \oplus B) = RF(A) \oplus RF(B)$ 得出。
引理 \ref{derived-lemma-replace-resolution} 的第三个假设由上同调长正合序列得出。
因此，对每个复形 $K^\bullet$，都存在取值于由对 $F$ 右无环的对象组成的
复形中的拟同构 $K^\bullet \to L^\bullet$。这证明陈述 (3)。

\medskip\noindent
我们断言，若 $L^\bullet \to M^\bullet$ 是由对 $F$ 右无环的对象组成的
复形之间的拟同构，则
$F(L^\bullet) \to F(M^\bullet)$ 是拟同构。
若证明此断言，则由引理 \ref{derived-lemma-find-existence-computes}，
得到引理的陈述 (1) 和 (2)。
为证明断言，选取整数 $i \in \mathbf{Z}$。
考虑特异三角
$$
\sigma_{\geq i - n - 1}L^\bullet \to
\sigma_{\geq i - n - 1}M^\bullet \to Q^\bullet,
$$
即令 $Q^\bullet$ 为第一个映射的锥。
注意，$Q^\bullet$ 有下界，且除可能在
$j = i - n - 1$ 或 $j = i - n - 2$ 外，$H^j(Q^\bullet)$ 均为零。
可以对 $Q^\bullet$ 应用 $RF$。
利用引理 \ref{derived-lemma-two-ss-complex-functor} 的第二个谱序列和
所假设的上同调消失 (2)，得出除可能对
$j \in \{i - n - 2, \ldots, i - 1\}$ 外，
$H^j(RF(Q^\bullet))$ 均为零。因此
$RF(\sigma_{\geq i - n - 1}L^\bullet) \to RF(\sigma_{\geq i - n - 1}M^\bullet)$
在次数 $\geq i$ 的上同调对象上诱导同构。
由命题 \ref{derived-proposition-enough-acyclics}，我们知道
$RF(\sigma_{\geq i - n - 1}L^\bullet) = \sigma_{\geq i - n - 1}F(L^\bullet)$
以及
$RF(\sigma_{\geq i - n - 1}M^\bullet) = \sigma_{\geq i - n - 1}F(M^\bullet)$。
所以如所需，$F(L^\bullet) \to F(M^\bullet)$ 在次数 $i$ 是同构。

\medskip\noindent
(4)(a) 由引理 \ref{derived-lemma-negative-vanishing} 得出。

\medskip\noindent
为证明 (4)(b)，以由对 $F$ 右无环的对象组成的复形 $L^\bullet$
表示 $E$。由 (2)，$RF(E)$ 由复形 $F(L^\bullet)$ 表示，而
% P02-E0079: frozen proof introduces sigma_{>=c} although c is not defined in this argument.
$RF(\sigma_{\geq c}L^\bullet)$ 由
$\sigma_{\geq c}F(L^\bullet)$ 表示。考虑评注
\ref{derived-remark-truncation-distinguished-triangle} 的特异三角
$$
H^{b - n}(L^\bullet)[n - b] \to
\tau_{\geq b - n}L^\bullet \to
\tau_{\geq b - n + 1}L^\bullet
$$
上述消失性给出：当 $i \geq b$ 时，
$H^i(RF(\tau_{\geq b - n}L^\bullet))$ 与
$H^i(RF(\tau_{\geq b - n + 1}L^\bullet))$ 相同。
考虑复形的短正合序列
$$
0 \to
\Im(L^{b - n - 1} \to L^{b - n})[n - b] \to
\sigma_{\geq b - n}L^\bullet \to
\tau_{\geq b - n}L^\bullet \to 0
$$
利用与此关联的特异三角（见第
\ref{derived-section-canonical-delta-functor} 节）以及与此前相同的消失性，
得出当 $i \geq b$ 时，
$H^i(RF(\tau_{\geq b - n}L^\bullet))$ 与
$H^i(RF(\sigma_{\geq b - n}L^\bullet))$ 相同。
由于映射 $RF(\sigma_{\geq b - n}L^\bullet) \to RF(L^\bullet)$
由 $\sigma_{\geq b - n}F(L^\bullet) \to F(L^\bullet)$ 表示，
所以如所需，当 $i \geq b$ 时，它又与
$H^i(RF(L^\bullet))$ 相同。

\medskip\noindent
证明 (4)(c)。在关于 $E$ 的假设下，有
$\tau_{\leq a - 1}E = 0$，并由 (4)(a) 得到当
$i \leq a - 1$ 时 $H^i(RF(E))$ 消失。
类似地，有 $\tau_{\geq b + 1}E = 0$，因而由 (4)(b) 得到当
$i \geq b + n$ 时 $H^i(RF(E))$ 消失。
\end{proof}

\begin{lemma}
\label{derived-lemma-unbounded-left-derived}
设 $F : \mathcal{A} \to \mathcal{B}$ 为阿贝尔范畴之间的右正合函子。
若
\begin{enumerate}
\item $\mathcal{A}$ 的每个对象都是某个对 $F$ 左无环的对象的商；
\item 存在整数 $n \geq 0$，使得 $L^nF = 0$。
\end{enumerate}
则
\begin{enumerate}
\item $LF : D(\mathcal{A}) \to D(\mathcal{B})$ 存在；
\item 由对 $F$ 左无环的对象组成的任意复形都计算 $LF$；
\item 任意复形都是一个拟同构的靶，该拟同构的源是由对 $F$
左无环的对象组成的复形；
\item 对 $E \in D(\mathcal{A})$：
\begin{enumerate}
% P02-E0080: the frozen map omits a closing parenthesis after LF(tau_{<=a+n-1}E).
\item 当 $i \leq a$ 时，
$H^i(LF(\tau_{\leq a + n - 1}E) \to H^i(LF(E))$ 是同构；
\item 当 $i \geq b$ 时，
$H^i(LF(E)) \to H^i(LF(\tau_{\geq b}E))$ 是同构；
\item 若对某个 $-\infty \leq a \leq b \leq \infty$，
当 $i \not \in [a, b]$ 时 $H^i(E) = 0$，则当
$i \not \in [a - n + 1, b]$ 时 $H^i(LF(E)) = 0$。
\end{enumerate}
\end{enumerate}
\end{lemma}

\begin{proof}
这对偶于引理 \ref{derived-lemma-unbounded-right-derived}。
\end{proof}






\section{导出余极限}
\label{derived-section-derived-colimit}

\noindent
三角范畴中有导出余极限的概念。

\begin{definition}
\label{derived-definition-derived-colimit}
设 $\mathcal{D}$ 为三角范畴。
设 $(K_n, f_n)$ 为 $\mathcal{D}$ 中对象的系统。
若直和 $\bigoplus K_n$ 存在，并且存在特异三角
$$
\bigoplus K_n \to \bigoplus K_n \to K \to \bigoplus K_n[1]
$$
其中映射 $\bigoplus K_n \to \bigoplus K_n$ 在次数 $n$
由 $1 - f_n$ 给出，则称对象 $K$ 为系统 $(K_n)$ 的
{\it 导出余极限}，或{\it 同伦余极限}。
在此情形下，有时用记号 $K = \text{hocolim} K_n$ 表示。
\end{definition}

\noindent
由 TR3，若导出余极限存在，则它在（不唯一的）同构意义下唯一。
此外，由 TR1，只要 $\bigoplus K_n$ 存在，$K_n$ 的导出余极限就存在。
阿贝尔群范畴的导出范畴 $D(\textit{Ab})$ 是所有同伦余极限都存在的
三角范畴的一个例子。

\medskip\noindent
不唯一性使得难以确定导出余极限。在更多代数，引理
\ref{more-algebra-lemma-map-from-hocolim} 中，读者可找到正合序列
$$
0 \to R^1\lim \Hom(K_n, L[-1]) \to \Hom(\text{hocolim} K_n, L)
\to \lim \Hom(K_n, L) \to 0
$$
它用通常的 $\Hom$ 群描述从同伦余极限出发的 $\Hom$ 群。

\begin{remark}
\label{derived-remark-uniqueness-derived-colimit}
设 $\mathcal{D}$ 为三角范畴。
设 $(K_n, f_n)$ 为 $\mathcal{D}$ 中对象的系统。
可以把导出余极限看作 $\mathcal{D}$ 中的对象 $K$，连同态射
$i_n : K_n \to K$，使得 $i_{n + 1} \circ f_n = i_n$，
并且存在态射 $c : K \to \bigoplus K_n[1]$，满足
$$
\bigoplus K_n \xrightarrow{1 - f_n} \bigoplus K_n \xrightarrow{i_n}
K \xrightarrow{c} \bigoplus K_n[1]
$$
为特异三角。若 $(K', i'_n, c')$ 是第二个导出余极限，
则存在同构 $\varphi : K \to K'$，使得
$\varphi \circ i_n = i'_n$ 且 $c' \circ \varphi = c$。
$\varphi$ 的存在性是 TR3，而 $\varphi$ 是同构这一事实由引理
\ref{derived-lemma-third-isomorphism-triangle} 给出。
\end{remark}

\begin{remark}
\label{derived-remark-functoriality-derived-colimit}
设 $\mathcal{D}$ 为三角范畴。
设 $(a_n) : (K_n, f_n) \to (L_n, g_n)$ 为 $\mathcal{D}$ 中对象系统的态射。
设 $(K, i_n, c)$ 为第一个系统的导出余极限，
$(L, j_n, d)$ 为第二个系统的导出余极限，记号如评注
\ref{derived-remark-uniqueness-derived-colimit}。
则存在态射 $a : K \to L$，使得
$a \circ i_n = j_n$ 且 $d \circ a = (a_n[1]) \circ c$。
这由对定义这些对象的特异三角应用 TR3 得出。
\end{remark}

\begin{lemma}
\label{derived-lemma-hocolim-subsequence}
设 $\mathcal{D}$ 为三角范畴。
设 $(K_n, f_n)$ 为 $\mathcal{D}$ 中对象的系统。
设 $n_1 < n_2 < n_3 < \ldots$ 为整数序列。
假设 $\bigoplus K_n$ 和 $\bigoplus K_{n_i}$ 存在。
则存在同构
$\text{hocolim} K_{n_i} \to \text{hocolim} K_n$，使得对所有 $i$，图
$$
\xymatrix{
K_{n_i} \ar[r] \ar[d]_{\text{id}} & \text{hocolim} K_{n_i} \ar[d] \\
K_{n_i} \ar[r] & \text{hocolim} K_n
}
$$
交换。
\end{lemma}

\begin{proof}
设 $g_i : K_{n_i} \to K_{n_{i + 1}}$ 为复合
$f_{n_{i + 1} - 1} \circ \ldots \circ f_{n_i}$。
如下构造交换图
$$
\vcenter{
\xymatrix{
\bigoplus\nolimits_i K_{n_i} \ar[r]_{1 - g_i} \ar[d]_b &
\bigoplus\nolimits_i K_{n_i} \ar[d]^a \\
\bigoplus\nolimits_n K_n \ar[r]^{1 - f_n} &
\bigoplus\nolimits_n K_n
}
}
\quad\text{且}\quad
\vcenter{
\xymatrix{
\bigoplus\nolimits_n K_n \ar[r]_{1 - f_n} \ar[d]_d &
\bigoplus\nolimits_n K_n \ar[d]^c \\
\bigoplus\nolimits_i K_{n_i} \ar[r]^{1 - g_i} &
\bigoplus\nolimits_i K_{n_i}
}
}
$$
令 $a_i = a|_{K_{n_i}}$ 为 $K_{n_i}$ 到直和中的包含映射。
换言之，$a$ 是自然包含。令 $b_i = b|_{K_{n_i}}$ 为映射
$$
K_{n_i}
\xrightarrow{1,\ f_{n_i},\ f_{n_i + 1} \circ f_{n_i},
\ \ldots,\ f_{n_{i + 1} - 2} \circ \ldots \circ f_{n_i}}
K_{n_i} \oplus K_{n_i + 1} \oplus \ldots \oplus K_{n_{i + 1} - 1}
$$
若 $n_{i - 1} < j \leq n_i$，则令 $c_j = c|_{K_j}$ 为映射
$$
K_j \xrightarrow{f_{n_i - 1} \circ \ldots \circ f_j} K_{n_i}
$$
若对任意 $i$ 都有 $j \not = n_i$，则令 $d_j = d|_{K_j}$ 为零；
并令 $d_{n_i}$ 为 $K_{n_i}$ 到直和中的自然包含。
换言之，$d$ 是自然投影。
由 TR3，这些图定义态射
$$
\varphi : \text{hocolim} K_{n_i} \to \text{hocolim} K_n
\quad\text{且}\quad
\psi : \text{hocolim} K_n \to \text{hocolim} K_{n_i}
$$
由于 $c \circ a$ 和 $d \circ b$ 是恒等映射，由引理
\ref{derived-lemma-third-isomorphism-triangle}，可知
$\varphi \circ \psi$ 是同构。
反过来，得到态射 $a \circ c$ 和 $b \circ d$。
考虑由如下规则给出的态射
$h = (h_j) : \bigoplus K_n \to \bigoplus K_n$：
对 $n_{i - 1} < j < n_i$，令
$$
h_j : K_j
\xrightarrow{1,\ f_j,\ f_{j + 1} \circ f_j,
\ \ldots,\ f_{n_i - 1} \circ \ldots \circ f_j}
K_j \oplus \ldots \oplus K_{n_i}
$$
于是读者可验证
$(1 - f) \circ h = \text{id} - a \circ c$ 且
$h \circ (1 - f) = \text{id} - b \circ d$。
这意味着，由引理 \ref{derived-lemma-third-map-square-zero}，
$\text{id} - \psi \circ \varphi$ 的平方为零（略去一个小论证）。
换言之，$\psi \circ \varphi$ 与恒等映射之差为幂零自同态，
故它是同构。因此如所需，$\varphi$ 与 $\psi$ 都是同构。
\end{proof}

\begin{lemma}
\label{derived-lemma-direct-sums}
设 $\mathcal{A}$ 为阿贝尔范畴。
若 $\mathcal{A}$ 有正合的可数直和，则 $D(\mathcal{A})$ 有可数直和。
事实上，给定以可数指标集 $I$ 为指标的一族复形 $K_i^\bullet$，
逐项直和 $\bigoplus K_i^\bullet$ 是各 $K_i^\bullet$ 在
$D(\mathcal{A})$ 中的直和。
\end{lemma}

\begin{proof}
设 $L^\bullet$ 为复形。假设给定 $D(\mathcal{A})$ 中的映射
$\alpha_i : K_i^\bullet \to L^\bullet$。
这意味着存在复形的拟同构
$s_i : M_i^\bullet \to K_i^\bullet$ 和复形映射
$f_i : M_i^\bullet \to L^\bullet$，使得
$\alpha_i = f_is_i^{-1}$。由假设，复形映射
$$
s : \bigoplus M_i^\bullet \longrightarrow \bigoplus K_i^\bullet
$$
是拟同构。因此，令 $f = \bigoplus f_i$，可见
$\alpha = fs^{-1}$ 是 $D(\mathcal{A})$ 中的映射，
它与余投影 $K_i^\bullet \to \bigoplus K_i^\bullet$ 的复合为
$\alpha_i$。我们略去验证 $\alpha$ 的唯一性。
\end{proof}

\begin{lemma}
\label{derived-lemma-compute-colimit}
设 $\mathcal{A}$ 为阿贝尔范畴。假设以 $\mathbf{N}$ 为指标的余极限存在且
正合。则可数直和存在且正合。此外，若 $(A_n, f_n)$ 是以
$\mathbf{N}$ 为指标的系统，则存在短正合序列
$$
0 \to \bigoplus A_n \to \bigoplus A_n \to \colim A_n \to 0
$$
其中第一个映射在次数 $n$ 由 $1 - f_n$ 给出。
\end{lemma}

\begin{proof}
第一个陈述由
$\bigoplus A_n = \colim (A_1 \oplus \ldots \oplus A_n)$ 得出。
对于第二个陈述，注意对每个 $n$ 都有短正合序列
$$
0 \to
A_1 \oplus \ldots \oplus A_{n - 1} \to
A_1 \oplus \ldots \oplus A_n \to A_n \to 0
$$
其中第一个映射由映射 $1 - f_i$ 给出，第二个映射是转移映射之和。
取余极限便得到引理中的序列。
\end{proof}

\begin{lemma}
\label{derived-lemma-colim-hocolim}
设 $\mathcal{A}$ 为阿贝尔范畴。设 $L_n^\bullet$ 为
$\mathcal{A}$ 中复形的系统。假设 $\mathcal{A}$ 中以
$\mathbf{N}$ 为指标的余极限存在且正合。
则逐项余极限 $L^\bullet = \colim L_n^\bullet$ 是该系统在
$D(\mathcal{A})$ 中的同伦余极限。
\end{lemma}

\begin{proof}
由引理 \ref{derived-lemma-compute-colimit}，存在复形的正合序列
$$
0 \to \bigoplus L_n^\bullet \to \bigoplus L_n^\bullet \to L^\bullet \to 0
$$
由引理 \ref{derived-lemma-direct-sums}，这些直和是 $D(\mathcal{A})$ 中的直和。
因此，结论由定义 \ref{derived-definition-derived-colimit} 中导出余极限的定义，
以及复形的短正合序列给出特异三角这一事实得出
（引理 \ref{derived-lemma-derived-canonical-delta-functor}）。
\end{proof}

\begin{lemma}
\label{derived-lemma-cohomology-of-hocolim}
设 $\mathcal{D}$ 为具有可数直和的三角范畴。
设 $\mathcal{A}$ 为具有正合的、以 $\mathbf{N}$ 为指标的余极限的
阿贝尔范畴。设 $H : \mathcal{D} \to \mathcal{A}$ 为与可数直和交换的
同调函子。则对 $\mathcal{D}$ 中任意对象系统，都有
$H(\text{hocolim} K_n) = \colim H(K_n)$。
\end{lemma}

\begin{proof}
写 $K = \text{hocolim} K_n$。对定义它的特异三角应用 $H$，得到
$$
\bigoplus H(K_n) \to \bigoplus H(K_n)
\to H(K) \to
\bigoplus H(K_n[1]) \to \bigoplus H(K_n[1])
$$
其中第一个映射由 $1 - H(f_n)$ 给出，最后一个映射由
$1 - H(f_n[1])$ 给出。应用引理
\ref{derived-lemma-compute-colimit} 即证明本引理。
\end{proof}

\noindent
下列引理告诉我们，从紧对象（稍后定义）出发取映射与导出余极限交换。

\begin{lemma}
\label{derived-lemma-commutes-with-countable-sums}
设 $\mathcal{D}$ 为具有可数直和的三角范畴。
设 $K \in \mathcal{D}$ 为如下对象：对任意可数对象集
$E_n \in \mathcal{D}$，典范映射
$$
\bigoplus \Hom_\mathcal{D}(K, E_n)
\longrightarrow
\Hom_\mathcal{D}(K, \bigoplus E_n)
$$
都是双射。则给定 $\mathcal{D}$ 中任意以 $\mathbf{N}$ 为指标的系统
$L_n$，若其导出余极限 $L = \text{hocolim} L_n$ 存在，则映射
$$
\colim \Hom_\mathcal{D}(K, L_n) \longrightarrow \Hom_\mathcal{D}(K, L)
$$
是双射。
\end{lemma}

\begin{proof}
由引理 \ref{derived-lemma-representable-homological}，这是引理
\ref{derived-lemma-cohomology-of-hocolim} 的特例。
\end{proof}







\section{导出极限}
\label{derived-section-derived-limit}

\noindent
在三角范畴中有导出极限这一概念。

\begin{definition}
\label{derived-definition-derived-limit}
设 $\mathcal{D}$ 为三角范畴。
设 $(K_n, f_n)$ 为 $\mathcal{D}$ 中对象的逆系统。
若直积 $\prod K_n$ 存在，且有特异三角
$$
K \to \prod K_n \to \prod K_n \to K[1]
$$
其中映射 $\prod K_n \to \prod K_n$ 由
$(k_n) \mapsto (k_n - f_{n + 1}(k_{n + 1}))$ 给出，
则称对象 $K$ 为系统 $(K_n)$ 的一个{\it 导出极限}，
或一个{\it 同伦极限}。在这种情况下，有时记作 $K = R\lim K_n$。
\end{definition}

\noindent
由 TR3，导出极限若存在，则在（非唯一的）同构意义下唯一。
此外，由 TR1，只要 $\prod K_n$ 存在，导出极限 $R\lim K_n$ 就存在。
阿贝尔群范畴的导出范畴 $D(\textit{Ab})$ 是所有导出极限均存在的
三角范畴的一个例子。

\medskip\noindent
这种非唯一性使得难以精确刻画导出极限。在 More on Algebra，
引理 \ref{more-algebra-lemma-map-into-Rlim} 中，读者可找到正合列
$$
0 \to R^1\lim \Hom(L, K_n[-1]) \to \Hom(L, R\lim K_n)
\to \lim \Hom(L, K_n) \to 0
$$
它用通常的 $\Hom$ 描述了到导出极限的 $\Hom$。

\begin{lemma}
\label{derived-lemma-products}
设 $\mathcal{A}$ 为具有正合可数直积的阿贝尔范畴。则
\begin{enumerate}
\item $D(\mathcal{A})$ 有可数直积；
\item $D(\mathcal{A})$ 中的可数直积 $\prod K_i$，可通过对代表
$K_i$ 的任意复形逐项取直积而得到；
\item $H^p(\prod K_i) = \prod H^p(K_i)$。
\end{enumerate}
\end{lemma}

\begin{proof}
设 $K_i^\bullet$ 为 $D(\mathcal{A})$ 中代表 $K_i$ 的复形。
设 $L^\bullet$ 为复形。假设给定 $D(\mathcal{A})$ 中的映射
$\alpha_i : L^\bullet \to K_i^\bullet$。
这意味着存在复形的拟同构 $s_i : K_i^\bullet \to M_i^\bullet$
和复形映射 $f_i : L^\bullet \to M_i^\bullet$，使得
$\alpha_i = s_i^{-1}f_i$。由假设，复形映射
$$
s : \prod K_i^\bullet \longrightarrow \prod M_i^\bullet
$$
是拟同构。因此，令 $f = \prod f_i$，可见
$\alpha = s^{-1}f$ 是 $D(\mathcal{A})$ 中的映射，
它与投影 $\prod K_i^\bullet \to K_i^\bullet$ 的复合为 $\alpha_i$。
我们略去验证 $\alpha$ 的唯一性。
\end{proof}

\noindent
这里应当陈述并证明引理 \ref{derived-lemma-compute-colimit}、
\ref{derived-lemma-colim-hocolim} 和
\ref{derived-lemma-commutes-with-countable-sums} 的对偶形式。
不过目前我们不知道这些引理有什么应用。

\begin{lemma}
\label{derived-lemma-inverse-limit-bounded-below}
设 $\mathcal{A}$ 为具有可数直积和足够多内射对象的阿贝尔范畴。
设 $(K_n)$ 为 $D^+(\mathcal{A})$ 中的逆系统。
则 $R\lim K_n$ 存在。
\end{lemma}

\begin{proof}
只需证明 $\prod K_n$ 在 $D(\mathcal{A})$ 中存在。
对每个 $n$，可用内射对象组成的下有界复形 $I_n^\bullet$ 代表
$K_n$（引理 \ref{derived-lemma-injective-resolutions-exist}）。
于是 $\prod K_n$ 由 $\prod I_n^\bullet$ 代表，见引理
\ref{derived-lemma-product-K-injective}。
\end{proof}

\begin{remark}
\label{derived-remark-functoriality-derived-limit}
设 $\mathcal{D}$ 为三角范畴。设 $(K_n)$ 和 $(L_n)$ 为
$\mathcal{D}$ 中具有导出极限 $K$ 和 $L$ 的对象逆系统。
假设有 pro-对象的态射 $a : (K_n) \to (L_n)$，见 Categories，
例 \ref{categories-example-pro-morphism-inverse-systems}。
这意味着给定映射 $a_n : K_{m(n)} \to L_n$。不妨假设
$m(1) < m(2) < m(3) < \ldots$，且各映射 $a_n$ 与系统中的
转移映射相容。于是可以考虑映射
$$
a', a'' : \prod K_n \longrightarrow \prod L_n
$$
它们由规则 $a'((k_n)) = (a_n(k_{m(n)}))$ 和
$$
a''((k_n)) =
(a_n(k_{m(n)} + f(k_{m(n) + 1}) + \ldots + f(k_{m(n + 1) - 1})))
$$
定义，
% P02-E0081: frozen source misspells “occurrence” as “occurence”.
其中每次出现的 $f$ 都表示逆系统 $(K_n)$ 的一个适当转移映射。
于是实线图
$$
\xymatrix{
K \ar[r] \ar@{..>}[d] & \prod K_n \ar[r] \ar[d]^{a'} &
\prod K_n \ar[d]^{a''} \\
L \ar[r] & \prod L_n \ar[r] & \prod L_n
}
$$
交换，并且由 TR3 得到虚线箭头，从而产生特异三角的态射。
我们提醒读者，映射 $K \to L$ 不唯一。稍后会看到，若 $a$
是 pro-同构，则 $K \to L$ 是同构，见 More on Algebra，
引理 \ref{more-algebra-lemma-pro-isomorphism-Rlim}。
\end{remark}

\begin{remark}
\label{derived-remark-map-into-derived-limit-truncations}
设 $\mathcal{A}$ 为阿贝尔范畴。设 $K^\bullet$ 为
$\mathcal{A}$ 中的复形。则 $\tau_{\geq -n}K^\bullet$ 构成一个复形逆系统，
% P02-E0082: frozen source says “is an inverse system of complexes and which in particular determines”.
因而特别给出 $D(\mathcal{A})$ 中的逆系统。假设
$R\lim \tau_{\geq -n}K^\bullet$ 存在。则典范映射
$c_n : K^\bullet \to \tau_{\geq -n}K^\bullet$ 与逆系统的转移映射相容。
由定义 \ref{derived-definition-derived-limit} 中的定义特异三角以及引理
\ref{derived-lemma-representable-homological}，可知存在态射
$$
c : K^\bullet \longrightarrow R\lim \tau_{\geq -n}K^\bullet
$$
在 $D(\mathcal{A})$ 中，使得 $c$ 与投影
$R\lim \tau_{\geq -n}K^\bullet \to \tau_{\geq -m}K^\bullet$ 的复合等于 $c_m$。
态射 $c$ 可能不唯一，但我们断言，$c$ 是否为同构与 $c$ 的选择
（以及同伦极限的选择）无关。事实上，对 $i \in \mathbf{Z}$ 和
$m > -i$，复合
$$
H^i(K^\bullet) \xrightarrow{H^i(c)} H^i(R\lim \tau_{\geq -n}K^\bullet)
\to H^i(\tau_{\geq -m}K^\bullet) = H^i(K^\bullet)
$$
是恒等映射。因此，$H^i(c)$ 为同构当且仅当第二个映射为同构。
这与 $c$ 无关，也与同伦极限的选择无关（因为任意两个选择都同构）。
\end{remark}

\begin{lemma}
\label{derived-lemma-difficulty-K-injectives}
设 $\mathcal{A}$ 为具有可数直积和足够多内射对象的阿贝尔范畴。
设 $K^\bullet$ 为复形。设 $I_n^\bullet$ 为引理
\ref{derived-lemma-special-inverse-system} 所构造的由内射对象组成的下有界复形逆系统。
则 $I^\bullet = \lim I_n^\bullet$ 存在，是 K-内射的，并且在
$D(\mathcal{A})$ 中代表 $R\lim \tau_{\geq -n}K^\bullet$；此外，下列条件等价：
\begin{enumerate}
\item 映射 $K^\bullet \to I^\bullet$（见证明）是拟同构；
\item 评注 \ref{derived-remark-map-into-derived-limit-truncations} 的映射
$K^\bullet \to R\lim \tau_{\geq -n}K^\bullet$ 是 $D(\mathcal{A})$ 中的同构。
\end{enumerate}
\end{lemma}

\begin{proof}
由于 $R\lim \tau_{\geq -n}K^\bullet$ 由引理
\ref{derived-lemma-inverse-limit-bounded-below} 存在，本引理的陈述是有意义的。
由引理 \ref{derived-lemma-bounded-below-injectives-K-injective}，每个复形
$I_n^\bullet$ 都是 K-内射的。对所有 $n \geq 1$，选择直和分解
$I_{n + 1}^p = C_{n + 1}^p \oplus I_n^p$。令 $C_1^p = I_1^p$。
复形 $I^\bullet = \lim I_n^\bullet$ 存在，因为可以取
$I^p = \prod_{n \geq 1} C_n^p$。固定 $p \in \mathbf{Z}$。
我们断言存在 $\mathcal{A}$ 中对象的分裂短正合列
$$
0 \to I^p \to \prod I_n^p \to \prod I_n^p \to 0
$$
这里第一个映射由投影映射 $I^p \to I_n^p$ 给出，第二个映射由
$(x_n) \mapsto (x_n - f^p_{n + 1}(x_{n + 1}))$ 给出，其中
$f^p_n : I_n^p \to I_{n - 1}^p$ 是转移映射。分裂来自映射
$\prod I_n^p \to \prod C_n^p = I^p$。由此得到逐项分裂的复形短正合列
$$
0 \to I^\bullet \to \prod I_n^\bullet \to \prod I_n^\bullet \to 0
$$
从而得到 $K(\mathcal{A})$ 和 $D(\mathcal{A})$ 中相应的特异三角。
由引理 \ref{derived-lemma-product-K-injective}，这些直积是 K-内射的，
并代表 $D(\mathcal{A})$ 中相应的直积。
因此 $I^\bullet$ 代表 $R\lim I_n^\bullet$
（定义 \ref{derived-definition-derived-limit}）。由于导出极限是在导出范畴层面定义的，
$R\lim I_n^\bullet \cong R\lim \tau_{\geq -n}K^\bullet$，
故 $I^\bullet$ 代表 $R\lim \tau_{\geq -n}K^\bullet$。
此外，由引理 \ref{derived-lemma-triangle-K-injective}，复形 $I^\bullet$ 是 K-内射的。
由引理 \ref{derived-lemma-special-inverse-system} 的交换图，并且由于当
$n \gg 0$ 时 $K^i = (\tau_{\geq -n}K^\bullet)^i$，可知得到唯一映射
$\gamma : K^\bullet \to I^\bullet$，使得各图
$$
\xymatrix{
K^\bullet \ar[r] \ar[d]_\gamma & \tau_{\geq -n} K^\bullet \ar[d] \\
I^\bullet \ar[r] & I_n^\bullet
}
$$
交换。因此，$\gamma$ 是复形映射，并且在 $D(\mathcal{A})$ 中代表评注
\ref{derived-remark-map-into-derived-limit-truncations} 的映射
$c : K^\bullet \to R\lim \tau_{\geq -n}K^\bullet$。
换言之，图
$$
\xymatrix{
K^\bullet \ar[r]_-c \ar[d]_\gamma & R\lim \tau_{\geq -n} K^\bullet
\ar[d]^{\cong} \\
I^\bullet \ar[r]^-{\cong} & R\lim I_n^\bullet
}
$$
在 $D(\mathcal{A})$ 中交换。引理得证。
\end{proof}

\begin{lemma}
\label{derived-lemma-enough-K-injectives-Ab4-star}
设 $\mathcal{A}$ 为具有足够多内射对象和正合可数直积的阿贝尔范畴。
则每个复形都有一个到 K-内射复形的拟同构。
\end{lemma}

\begin{proof}
由引理 \ref{derived-lemma-difficulty-K-injectives}，只需证明对
$D(\mathcal{A})$ 中的每个 $K$，$K \to R\lim\tau_{\geq -n}K$ 都是同构。
考虑定义特异三角
$$
R\lim\tau_{\geq -n}K \to
\prod \tau_{\geq -n}K \to
\prod \tau_{\geq -n}K \to
(R\lim\tau_{\geq -n}K)[1]
$$
由引理 \ref{derived-lemma-products}，有
$$
H^p(\prod \tau_{\geq -n}K) = \prod\nolimits_{p \geq -n} H^p(K)
$$
由所示特异三角的上同调长正合列，直接可得
$H^p(R\lim \tau_{\geq -n}K) = H^p(K)$。
\end{proof}




\section{全子范畴上的运算}
\label{derived-section-operate-on-full}

\noindent
设 $\mathcal{T}$ 为三角范畴。我们把 $\mathcal{T}$ 的全子范畴
与 $\Ob(\mathcal{T})$ 的子集视同。给定全子范畴
$\mathcal{A}, \mathcal{B}, \ldots$，作如下约定：
\begin{enumerate}
\item 对 $-\infty \leq a \leq b \leq \infty$，令
$\mathcal{A}[a, b]$ 为 $\mathcal{T}$ 的全子范畴，其对象是所有形如
$A[-i]$ 的对象，其中 $i \in [a, b] \cap \mathbf{Z}$ 且
$A \in \Ob(\mathcal{A})$（注意负号！）；
\item 令 $smd(\mathcal{A})$ 为 $\mathcal{T}$ 的全子范畴，其对象是
与 $\mathcal{A}$ 中对象的直和因子同构的所有对象；
\item 令 $add(\mathcal{A})$ 为 $\mathcal{T}$ 的全子范畴，其对象是
与 $\mathcal{A}$ 中对象的有限直和同构的所有对象；
\item 令 $\mathcal{A} \star \mathcal{B}$ 为 $\mathcal{T}$ 的全子范畴，
其对象是 $\mathcal{T}$ 中所有能嵌入特异三角 $A \to X \to B$ 的对象
$X$，其中 $A \in \Ob(\mathcal{A})$ 且 $B \in \Ob(\mathcal{B})$；
\item 令 $\mathcal{A}^{\star n} = \mathcal{A} \star \ldots \star \mathcal{A}$，
其中有 $n \geq 1$ 个因子（下文将看到 $\star$ 满足结合律）；
\item 令 $smd(add(\mathcal{A})^{\star n}) =
smd(add(\mathcal{A}) \star \ldots \star add(\mathcal{A}))$，
其中有 $n \geq 1$ 个因子。
\end{enumerate}
若 $E$ 是 $\mathcal{T}$ 的对象，有时也把 $E$ 看作
$\mathcal{T}$ 的仅含对象 $E$ 的全子范畴。于是可以考虑
$add(E[-1, 2])$ 等对象。我们提醒读者，这套记号并未得到普遍采用。

\begin{lemma}
\label{derived-lemma-associativity-star}
设 $\mathcal{T}$ 为三角范畴。给定全子范畴
$\mathcal{A}$、$\mathcal{B}$、$\mathcal{C}$，有
$(\mathcal{A} \star \mathcal{B}) \star \mathcal{C} =
\mathcal{A} \star (\mathcal{B} \star \mathcal{C})$。
\end{lemma}

\begin{proof}
若有特异三角 $A \to X \to B$ 和 $X \to Y \to C$，
则由公理 TR4，有特异三角 $A \to Y \to Z$ 和 $B \to Z \to C$。
\end{proof}

\begin{lemma}
\label{derived-lemma-smd-star}
设 $\mathcal{T}$ 为三角范畴。给定全子范畴
$\mathcal{A}$、$\mathcal{B}$，有
$smd(\mathcal{A}) \star smd(\mathcal{B}) \subset
smd(\mathcal{A} \star \mathcal{B})$，以及
$smd(smd(\mathcal{A}) \star smd(\mathcal{B})) =
smd(\mathcal{A} \star \mathcal{B})$。
\end{lemma}

\begin{proof}
假设有特异三角 $A_1 \to X \to B_1$，其中
$A_1 \oplus A_2 \in \Ob(\mathcal{A})$ 且
$B_1 \oplus B_2 \in \Ob(\mathcal{B})$。于是得到特异三角
$A_1 \oplus A_2 \to A_2 \oplus X \oplus B_2 \to B_1 \oplus B_2$，
这证明 $X$ 属于 $smd(\mathcal{A} \star \mathcal{B})$。
由此得到包含关系，而等式显然由该包含关系推出。
\end{proof}

\begin{lemma}
\label{derived-lemma-add-star}
设 $\mathcal{T}$ 为三角范畴。给定全子范畴
$\mathcal{A}$、$\mathcal{B}$，全子范畴
$add(\mathcal{A}) \star add(\mathcal{B})$ 和
$smd(add(\mathcal{A}))$ 对直和封闭。
\end{lemma}

\begin{proof}
事实上，若 $A \to X \to B$ 和 $A' \to X' \to B'$ 是特异三角，
$A, A' \in add(\mathcal{A})$ 且 $B, B' \in add(\mathcal{B})$，则
$A \oplus A' \to X \oplus X' \to B \oplus B'$ 是特异三角，且
$A \oplus A' \in add(\mathcal{A})$、$B \oplus B' \in add(\mathcal{B})$。
对 $smd(add(\mathcal{A}))$ 的结论是显然的。
\end{proof}

\begin{lemma}
\label{derived-lemma-cone-n}
设 $\mathcal{T}$ 为三角范畴。给定全子范畴 $\mathcal{A}$，对
$n \geq 1$，上面定义的子范畴
$$
\mathcal{C}_n = smd(add(\mathcal{A})^{\star n}) =
smd(add(\mathcal{A}) \star \ldots \star add(\mathcal{A}))
$$
是 $\mathcal{T}$ 的严格全子范畴，并且对直和及直和因子封闭；而且对所有
$n, m \geq 1$，都有
$\mathcal{C}_{n + m} = smd(\mathcal{C}_n \star \mathcal{C}_m)$。
\end{lemma}

\begin{proof}
立即由引理 \ref{derived-lemma-associativity-star}、\ref{derived-lemma-smd-star} 和
\ref{derived-lemma-add-star} 得出。
\end{proof}

\begin{remark}
\label{derived-remark-operations-functor}
设 $F : \mathcal{T} \to \mathcal{T}'$ 为三角范畴的正合函子。
给定 $\mathcal{T}$ 的全子范畴 $\mathcal{A}$，记
$F(\mathcal{A})$ 为 $\mathcal{T}'$ 的全子范畴，其对象是所有满足
$A \in \Ob(\mathcal{A})$ 的对象 $F(A)$。
% P02-E0083: frozen source has subject--verb disagreement in “whose objects consists”.
我们有
$$
F(\mathcal{A}[a, b]) = F(\mathcal{A})[a, b]
$$
$$
F(smd(\mathcal{A})) \subset smd(F(\mathcal{A})),
$$
$$
F(add(\mathcal{A})) \subset add(F(\mathcal{A})),
$$
$$
F(\mathcal{A} \star \mathcal{B}) \subset F(\mathcal{A}) \star F(\mathcal{B}),
$$
$$
F(\mathcal{A}^{\star n}) \subset F(\mathcal{A})^{\star n}.
$$
我们略去这些简单验证。
\end{remark}

\begin{remark}
\label{derived-remark-operations-unions}
设 $\mathcal{T}$ 为三角范畴。给定 $\mathcal{T}$ 的全子范畴
$\mathcal{A}_1 \subset \mathcal{A}_2 \subset \mathcal{A}_3 \subset \ldots$
和 $\mathcal{B}$，有
$$
\left(\bigcup \mathcal{A}_i\right)[a, b] = \bigcup \mathcal{A}_i[a, b]
$$
$$
smd\left(\bigcup \mathcal{A}_i\right) = \bigcup smd(\mathcal{A}_i),
$$
$$
add\left(\bigcup \mathcal{A}_i\right) = \bigcup add(\mathcal{A}_i),
$$
$$
\left(\bigcup \mathcal{A}_i\right) \star \mathcal{B} =
\bigcup \mathcal{A}_i \star \mathcal{B},
$$
$$
\mathcal{B} \star \left(\bigcup \mathcal{A}_i\right) =
\bigcup \mathcal{B} \star \mathcal{A}_i,
$$
$$
\left(\bigcup \mathcal{A}_i\right)^{\star n} =
\bigcup \mathcal{A}_i^{\star n}.
$$
我们略去这些简单验证。
\end{remark}

\begin{lemma}
\label{derived-lemma-in-cone-n}
设 $\mathcal{A}$ 为阿贝尔范畴，令 $\mathcal{D} = D(\mathcal{A})$。
设 $\mathcal{E} \subset \Ob(\mathcal{A})$ 为一个子集，并也把它视作
$\Ob(\mathcal{D})$ 的子集。设 $K$ 为 $\mathcal{D}$ 的对象。
\begin{enumerate}
\item 设 $b \geq a$，并假设当 $i \not \in [a, b]$ 时 $H^i(K)$ 为零，
而当 $i \in [a, b]$ 时 $H^i(K) \in \mathcal{E}$。则 $K$ 属于
$smd(add(\mathcal{E}[a, b])^{\star (b - a + 1)})$。
\item 设 $b \geq a$，并假设当 $i \not \in [a, b]$ 时 $H^i(K)$ 为零，
而当 $i \in [a, b]$ 时 $H^i(K) \in smd(add(\mathcal{E}))$。则 $K$ 属于
$smd(add(\mathcal{E}[a, b])^{\star (b - a + 1)})$。
\item 设 $b \geq a$，并假设 $K$ 可由复形 $K^\bullet$ 代表，且当
$i \not \in [a, b]$ 时 $K^i = 0$，当 $i \in [a, b]$ 时
$K^i \in \mathcal{E}$。则 $K$ 属于
$smd(add(\mathcal{E}[a, b])^{\star (b - a + 1)})$。
\item 设 $b \geq a$，并假设 $K$ 可由复形 $K^\bullet$ 代表，且当
$i \not \in [a, b]$ 时 $K^i = 0$，当 $i \in [a, b]$ 时
$K^i \in smd(add(\mathcal{E}))$。则 $K$ 属于
$smd(add(\mathcal{E}[a, b])^{\star (b - a + 1)})$。
\end{enumerate}
\end{lemma}

\begin{proof}
以下不再特别说明地使用引理 \ref{derived-lemma-cone-n}。
我们证明 (2)，它显然推出 (1)。对 $b - a$ 作归纳。
若 $b - a = 0$，则 $K$ 在 $\mathcal{D}$ 中同构于
$H^i(K)[-a]$，结论立即成立。
% P02-E0086: frozen source uses H^i(K) although the one surviving degree is a.
若 $b - a > 0$，则考虑特异三角
$$
\tau_{\leq b - 1}K^\bullet \to K^\bullet \to K^b[-b]
$$
并对 $b - a$ 使用归纳即可。我们略去 (3) 和 (4) 的证明。
\end{proof}

\begin{lemma}
\label{derived-lemma-forward-cone-n}
设 $\mathcal{T}$ 为三角范畴。设
$H : \mathcal{T} \to \mathcal{A}$ 为到阿贝尔范畴 $\mathcal{A}$ 的同调函子。
设 $a \leq b$，且 $\mathcal{E} \subset \Ob(\mathcal{T})$ 为如下子集：
当 $E \in \mathcal{E}$ 且 $i \not \in [a, b]$ 时，$H^i(E) = 0$。
则对 $X \in smd(add(\mathcal{E}[-m, m])^{\star n})$，当
$i \not \in [-m + na, m + nb]$ 时有 $H^i(X) = 0$。
\end{lemma}

\begin{proof}
略。这是关于这些定义的一道有趣练习。
\end{proof}





\section{三角范畴的生成元}
\label{derived-section-generators}

\noindent
本节简要介绍三角范畴生成元的几种不同概念。我们的术语取自
\cite{BvdB}（但他们所称的 ``\'epaisse''，我们称为“饱和”，
见定义 \ref{derived-definition-saturated}；此外，我们对 $add(\mathcal{A})$ 的定义不同）。

\medskip\noindent
设 $\mathcal{D}$ 为三角范畴，$E$ 为 $\mathcal{D}$ 的对象。
记 $\langle E \rangle_1$ 为 $\mathcal{D}$ 的严格全子范畴，
其对象是 $\mathcal{D}$ 中与 $E$ 的平移的有限直和
$$
\bigoplus\nolimits_{i = 1, \ldots, r} E[n_i]
$$
的直和因子同构的对象。显然，采用节 \ref{derived-section-operate-on-full} 的记号，有
$$
\langle E \rangle_1 = smd(add(E[-\infty, \infty]))
$$
对 $n > 1$，记 $\langle E \rangle_n$ 为 $\mathcal{D}$ 的全子范畴，
其对象是与能嵌入特异三角
$$
A \to X \to B \to A[1]
$$
的对象 $X$ 的直和因子同构的 $\mathcal{D}$ 中对象，其中 $A$ 是
$\langle E \rangle_1$ 的对象，$B$ 是 $\langle E \rangle_{n - 1}$ 的对象。
采用节 \ref{derived-section-operate-on-full} 的记号，有
$$
\langle E \rangle_n = smd(\langle E \rangle_1 \star \langle E \rangle_{n - 1})
$$
每个范畴 $\langle E \rangle_n$ 都是 $\mathcal{D}$ 的严格全加性子范畴
（由引理 \ref{derived-lemma-add-star}），并且在平移和取直和因子下保持不变。
然而，$\langle E \rangle_n$ 不一定对“取锥”或“扩张”封闭，
因而不一定是三角子范畴。如下子范畴则具有这一性质：
$$
\langle E \rangle = \bigcup\nolimits_n \langle E \rangle_n
$$
下面的引理将证明这一点。

\begin{lemma}
\label{derived-lemma-generated-by-E-explicit}
设 $\mathcal{T}$ 为三角范畴，$E$ 为 $\mathcal{T}$ 的对象。
对 $n \geq 1$，有
$$
\langle E \rangle_n =
smd(\langle E \rangle_1 \star \ldots \star \langle E \rangle_1) =
smd({\langle E \rangle_1}^{\star n}) =
\bigcup\nolimits_{m \geq 1} smd(add(E[-m, m])^{\star n})
$$
对 $n, n' \geq 1$，有 $\langle E \rangle_{n + n'} =
smd(\langle E \rangle_n \star \langle E \rangle_{n'})$。
\end{lemma}

\begin{proof}
所示公式的左边等式由引理 \ref{derived-lemma-associativity-star}、
\ref{derived-lemma-smd-star} 和归纳法得出。中间的等式只是记号问题。
由于
$\langle E \rangle_1 = smd(add(E[-\infty, \infty])])$
% P02-E0084: frozen source has an extra closing bracket in this formula.
并且 $E[-\infty, \infty] = \bigcup_{m \geq 1} E[-m, m]$，
由评注 \ref{derived-remark-operations-unions} 和引理 \ref{derived-lemma-smd-star}，
得到右边的等式。最后一个陈述则由该评注及引理
\ref{derived-lemma-cone-n} 的相应陈述得出。
\end{proof}

\begin{lemma}
\label{derived-lemma-find-smallest-containing-E}
设 $\mathcal{D}$ 为三角范畴，$E$ 为 $\mathcal{D}$ 的对象。子范畴
$$
\langle E \rangle = \bigcup\nolimits_n \langle E \rangle_n
= \bigcup\nolimits_{n, m \geq 1} smd(add(E[-m, m])^{\star n})
$$
是 $\mathcal{D}$ 的严格全、饱和三角子范畴，并且是包含对象 $E$ 的
$\mathcal{D}$ 的最小此类子范畴。
\end{lemma}

\begin{proof}
右边的等式由引理 \ref{derived-lemma-generated-by-E-explicit} 得出。
显然，$\langle E \rangle = \bigcup \langle E \rangle_n$ 包含 $E$，
并且在平移、直和及直和因子下保持不变。若
$A \in \langle E \rangle_a$、$B \in \langle E \rangle_b$，且
$A \to X \to B \to A[1]$ 是特异三角，则由引理
\ref{derived-lemma-generated-by-E-explicit}，有
$X \in \langle E \rangle_{a + b}$。因此
$\bigcup \langle E \rangle_n$ 也对扩张封闭，从而是三角子范畴。

\medskip\noindent
最后，设 $\mathcal{D}' \subset \mathcal{D}$ 为包含 $E$ 的
$\mathcal{D}$ 的严格全、饱和三角子范畴。则
$\mathcal{D}'[-\infty, \infty] \subset \mathcal{D}'$、
$add(\mathcal{D}) \subset \mathcal{D}'$、
% P02-E0085: frozen source uses add(D), although preservation here requires add(D').
$smd(\mathcal{D}') \subset \mathcal{D}'$，且
$\mathcal{D}' \star \mathcal{D}' \subset \mathcal{D}'$。
换言之，从 $E$ 构造 $\langle E \rangle$ 所用的全部运算都保持
$\mathcal{D}'$。因此 $\langle E \rangle \subset \mathcal{D}'$，证明完毕。
\end{proof}

\begin{definition}
\label{derived-definition-generators}
设 $\mathcal{D}$ 为三角范畴，$E$ 为 $\mathcal{D}$ 的对象。
\begin{enumerate}
\item 若 $\mathcal{D}$ 中包含 $E$ 的最小严格全、饱和三角子范畴
等于 $\mathcal{D}$，换言之，若 $\langle E \rangle = \mathcal{D}$，
则称 $E$ 为 $\mathcal{D}$ 的{\it 经典生成元}。
\item 若对某个 $n \geq 1$ 有 $\langle E \rangle_n = \mathcal{D}$，
则称 $E$ 为 $\mathcal{D}$ 的{\it 强生成元}。
\item 若对 $\mathcal{D}$ 的任意非零对象 $K$，都存在整数 $n$
和非零映射 $E \to K[n]$，则称 $E$ 为 $\mathcal{D}$ 的
{\it 弱生成元}，或简称{\it 生成元}。
\end{enumerate}
\end{definition}

\noindent
这一定义可以推广到对象族的情形。

\begin{lemma}
\label{derived-lemma-right-orthogonal}
设 $\mathcal{D}$ 为三角范畴，$E, K$ 为 $\mathcal{D}$ 的对象。
下列条件等价：
\begin{enumerate}
\item 对所有 $i \in \mathbf{Z}$，有 $\Hom(E, K[i]) = 0$；
\item 对所有 $E' \in \langle E \rangle$，有 $\Hom(E', K) = 0$。
\end{enumerate}
\end{lemma}

\begin{proof}
(2) $\Rightarrow$ (1) 立即成立。反过来，假设 (1)。
则对 $\langle E \rangle_1$ 中所有 $X$，都有 $\Hom(X, K) = 0$。
对 $n$ 归纳并使用引理 \ref{derived-lemma-representable-homological}，
可见对 $\langle E \rangle_n$ 中所有 $X$，都有 $\Hom(X, K) = 0$。
\end{proof}

\begin{lemma}
\label{derived-lemma-classical-generator-generator}
设 $\mathcal{D}$ 为三角范畴，$E$ 为 $\mathcal{D}$ 的对象。
若 $E$ 是 $\mathcal{D}$ 的经典生成元，则 $E$ 是生成元。
\end{lemma}

\begin{proof}
假设 $E$ 是经典生成元。设 $K$ 为 $\mathcal{D}$ 的对象，且对所有
$i \in \mathbf{Z}$ 都有 $\Hom(E, K[i]) = 0$。由引理
\ref{derived-lemma-right-orthogonal}，对 $\langle E \rangle$ 中所有 $E'$，都有
$\Hom(E', K) = 0$。然而，由于 $\mathcal{D} = \langle E \rangle$，
可知 $\text{id}_K = 0$，即 $K = 0$。
\end{proof}

\begin{lemma}
\label{derived-lemma-classical-generator-strong-generator}
设 $\mathcal{D}$ 为具有强生成元的三角范畴，$E$ 为 $\mathcal{D}$ 的对象。
若 $E$ 是 $\mathcal{D}$ 的经典生成元，则 $E$ 是强生成元。
\end{lemma}

\begin{proof}
设 $E'$ 为 $\mathcal{D}$ 的对象，且
$\mathcal{D} = \langle E' \rangle_n$。由于
$\mathcal{D} = \langle E \rangle$，由引理
\ref{derived-lemma-find-smallest-containing-E}，对某个 $m \geq 1$ 有
$E' \in \langle E \rangle_m$。于是
$\langle E' \rangle_1 \subset \langle E \rangle_m$，故
$$
\mathcal{D} =
\langle E' \rangle_n = smd(
\langle E' \rangle_1 \star \ldots \star \langle E' \rangle_1)
\subset
smd(
\langle E \rangle_m \star \ldots \star \langle E \rangle_m)
=
\langle E \rangle_{nm}
$$
正如所需。这里使用了引理 \ref{derived-lemma-generated-by-E-explicit}。
\end{proof}

\begin{remark}
\label{derived-remark-check-on-generator}
设 $\mathcal{D}$ 为三角范畴，$E$ 为 $\mathcal{D}$ 的对象。
设 $T$ 为 $\mathcal{D}$ 中对象的一个性质。假设：
\begin{enumerate}
\item 若 $K_i \in \mathcal{D}$，$i = 1, \ldots, r$，且对
$i = 1, \ldots, r$ 都有 $T(K_i)$，则 $T(\bigoplus K_i)$ 成立；
\item 若 $K \to L \to M \to K[1]$ 是特异三角，且 $T$ 对其中两个对象成立，
则 $T$ 对第三个对象也成立；
\item 若 $T(K \oplus L)$ 成立，则 $T(K)$ 和 $T(L)$ 成立；
\item 对所有 $n$，$T(E[n])$ 成立。
\end{enumerate}
则 $T$ 对 $\langle E \rangle$ 的所有对象成立。
\end{remark}












\section{紧对象}
\label{derived-section-compact}

\noindent
定义如下。

\begin{definition}
\label{derived-definition-compact-object}
设 $\mathcal{D}$ 为具有任意直和的加性范畴。$\mathcal{D}$ 的
{\it 紧对象}是满足如下条件的对象 $K$：对任意集合 $I$ 以及由
$i \in I$ 参数化的对象 $E_i \in \Ob(\mathcal{D})$，映射
$$
\bigoplus\nolimits_{i \in I} \Hom_{\mathcal{D}}(K, E_i)
\longrightarrow
\Hom_{\mathcal{D}}(K, \bigoplus\nolimits_{i \in I} E_i)
$$
是双射。
\end{definition}

\noindent
这一概念在代数几何中非常有用。它是对象所满足的一个内在条件，
迫使这些对象确实具有某种“紧性”。

\begin{lemma}
\label{derived-lemma-compact-objects-subcategory}
设 $\mathcal{D}$ 为具有直和的（预）三角范畴。则 $\mathcal{D}$ 的紧对象
构成 $\mathcal{D}$ 的一个 Karoubi、饱和、严格全（预）三角子范畴
$\mathcal{D}_c$ 的全部对象。
\end{lemma}

\begin{proof}
设 $(X, Y, Z, f, g, h)$ 为 $\mathcal{D}$ 的特异三角，且
$X$ 和 $Y$ 都是紧的。则由引理 \ref{derived-lemma-representable-homological}
和五引理（Homology，引理 \ref{homology-lemma-five-lemma}），
$Z$ 也是紧对象。显然，若 $X \oplus Y$ 是紧的，则 $X$、$Y$
也都是紧对象。因此 $\mathcal{D}_c$ 是饱和三角子范畴。
由引理 \ref{derived-lemma-projectors-have-images-triangulated}，$\mathcal{D}$
是 Karoubi 范畴，故 $\mathcal{D}_c$ 亦然。
\end{proof}

\begin{lemma}
\label{derived-lemma-write-as-colimit}
设 $\mathcal{D}$ 为具有直和的三角范畴。设 $E_i$，$i \in I$，
为 $\mathcal{D}$ 的一族紧对象，且 $\bigoplus E_i$ 生成 $\mathcal{D}$。
则 $\mathcal{D}$ 的每个对象 $X$ 都可写成
$$
X = \text{hocolim} X_n
$$
其中 $X_1$ 是 $E_i$ 的平移的直和，并且每个转移态射都嵌入特异三角
$Y_n \to X_n \to X_{n + 1} \to Y_n[1]$，其中 $Y_n$ 是
$E_i$ 的平移的直和。
\end{lemma}

\begin{proof}
令 $X_1 = \bigoplus_{(i, m, \varphi)} E_i[m]$，其中直和遍历所有满足
$i \in I$、$m \in \mathbf{Z}$ 及 $\varphi : E_i[m] \to X$ 的三元组
$(i, m, \varphi)$。于是 $X_1$ 带有典范态射 $X_1 \to X$。
给定 $X_n \to X$，令 $Y_n = \bigoplus_{(i, m, \varphi)} E_i[m]$，
其中直和遍历所有满足 $i \in I$、$m \in \mathbf{Z}$，且
$\varphi : E_i[m] \to X_n$ 是使得 $E_i[m] \to X_n \to X$ 为零的态射的
三元组 $(i, m, \varphi)$。选择特异三角
$Y_n \to X_n \to X_{n + 1} \to Y_n[1]$，并任取态射
$X_{n + 1} \to X$，使 $X_n \to X_{n + 1} \to X$ 等于给定态射；
由 $Y_n$ 的选择，这样的态射存在。由各映射 $X_n \to X$ 的构造，
得到态射 $\text{hocolim} X_n \to X$。选择特异三角
$$
C \to \text{hocolim} X_n \to X \to C[1]
$$
设 $E_i[m] \to C$ 为态射。由于 $E_i$ 是紧的，引理
\ref{derived-lemma-commutes-with-countable-sums} 蕴含复合
$E_i[m] \to C \to \text{hocolim} X_n$ 对某个 $n$ 经 $X_n$ 分解，
设相应态射为 $E_i[m] \to X_n$。于是由 $Y_n$ 的构造，复合
$E_i[m] \to X_n \to X_{n + 1}$ 为零。换言之，复合
$E_i[m] \to C \to \text{hocolim} X_n$ 为零。这意味着态射
$E_i[m] \to C$ 来自某个态射 $E_i[m] \to X[-1]$。
$X_1$ 的构造又表明，该态射提升到 $\text{hocolim} X_n$，因而态射
$E_i[m] \to C$ 为零。假设 $\bigoplus E_i$ 生成 $\mathcal{D}$，
遂蕴含 $C$ 为零，证明完毕。
\end{proof}

\begin{lemma}
\label{derived-lemma-factor-through}
沿用引理 \ref{derived-lemma-write-as-colimit} 的假设和记号。
% P02-E0087: frozen source sentence “With assumptions and notation ...” lacks a finite verb.
若 $C$ 是紧对象，且 $C \to X_n$ 是态射，则存在分解
$C \to E \to X_n$，其中对某些 $i_1, \ldots, i_t \in I$，
$E$ 是 $\langle E_{i_1} \oplus \ldots \oplus E_{i_t} \rangle$ 的对象。
\end{lemma}

\begin{proof}
对 $n$ 归纳。基础情形 $n = 1$ 是显然的。若 $n > 1$，
考虑复合 $C \to X_n \to Y_{n - 1}[1]$。它可经某个
$E'[1] \to Y_{n - 1}[1]$ 分解，其中 $E'$ 是 $E_i$ 的平移的有限直和。
令 $I' \subset I$ 为该直和中出现的指标所成的有限集。于是得到
$$
\xymatrix{
E' \ar[r] \ar[d] &
C' \ar[r] \ar[d] &
C \ar[r] \ar[d] &
E'[1] \ar[d] \\
Y_{n - 1} \ar[r] &
X_{n - 1} \ar[r] &
X_n \ar[r] &
Y_{n - 1}[1]
}
$$
由归纳假设，态射 $C' \to X_{n - 1}$ 经
$E'' \to X_{n - 1}$ 分解，其中对某个有限子集 $I'' \subset I$，
$E''$ 是 $\langle \bigoplus_{i \in I''} E_i \rangle$ 的对象。
选择特异三角
$$
E' \to E'' \to E \to E'[1]
$$
则 $E$ 是 $\langle \bigoplus_{i \in I' \cup I''} E_i \rangle$ 的对象。
由构造及三角范畴公理，可选择态射 $C \to E$ 和 $E \to X_n$，
使之嵌入三角的态射 $(E', C', C) \to (E', E'', E)$ 及
$(E', E'', E) \to (Y_{n - 1}, X_{n - 1}, X_n)$。
复合 $C \to E \to X_n$ 可能不等于给定态射 $C \to X_n$，
但它们到 $Y_{n - 1}$ 的复合相等。设 $C \to X_{n - 1}$ 为提升其差的态射。
由归纳假设，它可经态射 $E''' \to X_{n - 1}$ 分解，其中
$E''$ 是 $\langle \bigoplus_{i \in I'''} E_i \rangle$ 的对象，
% P02-E0088: frozen source says E'' here, although the newly introduced object is E'''.
而 $I' \subset I$ 是某个有限子集。
% P02-E0089: frozen source repeats I' here, although the new index set is I'''.
因此，考虑 $E \oplus E''' \to X_n$ 即得到所求解，因为
$E \oplus E'''$ 是
$\langle \bigoplus_{i \in I' \cup I'' \cup I'''} E_i \rangle$ 的对象。
\end{proof}

\begin{definition}
\label{derived-definition-compactly-generated}
设 $\mathcal{D}$ 为具有任意直和的三角范畴。若存在以 $i \in I$
为指标的一族紧对象 $E_i$，使得 $\bigoplus E_i$ 生成 $\mathcal{D}$，
则称 $\mathcal{D}$ 是{\it 紧生成的}。
\end{definition}

\noindent
下面的命题阐明经典生成元与弱生成元之间的关系。

\begin{proposition}
\label{derived-proposition-generator-versus-classical-generator}
设 $\mathcal{D}$ 为具有直和的三角范畴，$E$ 为 $\mathcal{D}$ 的紧对象。
下列条件等价：
\begin{enumerate}
\item $E$ 是 $\mathcal{D}_c$ 的经典生成元，且 $\mathcal{D}$ 是紧生成的；
\item $E$ 是 $\mathcal{D}$ 的生成元。
\end{enumerate}
\end{proposition}

\begin{proof}
若 $E$ 是 $\mathcal{D}_c$ 的经典生成元，则
$\mathcal{D}_c = \langle E \rangle$。由 $\mathcal{D}$ 紧生成的假设和引理
\ref{derived-lemma-right-orthogonal}，形式地推出 $E$ 是 $\mathcal{D}$ 的生成元。

\medskip\noindent
反向更有意思。假设 $E$ 是 $\mathcal{D}$ 的生成元。
设 $X$ 为 $\mathcal{D}$ 的紧对象。应用引理
\ref{derived-lemma-write-as-colimit}，取 $I = \{1\}$ 和 $E_1 = E$，按该引理写成
$$
X = \text{hocolim} X_n
$$
由于 $X$ 是紧的，对某个 $n$，映射
$X \to \text{hocolim} X_n$ 经 $X_n$ 分解（引理
\ref{derived-lemma-commutes-with-countable-sums}）。因此 $X$ 是 $X_n$ 的直和因子。
由引理 \ref{derived-lemma-factor-through}，可见 $X$ 是
$\langle E \rangle$ 的对象，命题得证。
% P02-E0091: frozen source closes this proposition proof with “the lemma is proven”.
\end{proof}











\section{Brown 可表性}
\label{derived-section-brown}

\noindent
本节内容的参考文献是 \cite{Neeman-Grothendieck}。

\begin{lemma}
\label{derived-lemma-brown}
\begin{reference}
\cite[Theorem 3.1]{Neeman-Grothendieck}.
\end{reference}
设 $\mathcal{D}$ 为具有直和的紧生成三角范畴。设
$H : \mathcal{D} \to \textit{Ab}$ 为把直和变为直积的反变上同调函子。
则 $H$ 可表。
\end{lemma}

\begin{proof}
设 $E_i$，$i \in I$，为一族紧对象，且
$\bigoplus_{i \in I} E_i$ 生成 $\mathcal{D}$。不妨并且确实假设
对象集 $\{E_i\}$ 在平移下保持不变。考虑序对 $(i, a)$，其中
$i \in I$ 且 $a \in H(E_i)$，并令
$$
X_1 = \bigoplus\nolimits_{(i, a)} E_i
$$
由于 $H(X_1) = \prod_{(i, a)} H(E_i)$，可见 $(a)_{(i, a)}$
定义元素 $a_1 \in H(X_1)$。令 $H_1 = \Hom_\mathcal{D}(- , X_1)$。
由 Yoneda 引理（Categories，引理 \ref{categories-lemma-yoneda}），
元素 $a_1$ 定义自然变换 $H_1 \to H$。

\medskip\noindent
我们将归纳构造 $X_n$ 和变换 $a_n : H_n \to H$，其中
$H_n = \Hom_\mathcal{D}(-, X_n)$。具体地，对函子
$\Ker(H_n \to H)$ 应用上面的步骤，得到对象
$$
K_{n + 1} = \bigoplus\nolimits_{(i, k),\ k \in \Ker(H_n(E_i) \to H(E_i))} E_i
$$
以及变换 $\Hom_\mathcal{D}(-, K_{n + 1}) \to \Ker(H_n \to H)$。
由 Yoneda 引理，复合 $\Hom_\mathcal{D}(-, K_{n + 1}) \to H_n$
给出态射 $K_{n + 1} \to X_n$。在 $\mathcal{D}$ 中选择特异三角
$$
K_{n + 1} \to X_n \to X_{n + 1} \to K_{n + 1}[1]
$$
由构造，元素 $a_n \in H(X_n)$ 在 $H(K_{n + 1})$ 中的像为零。
由于 $H$ 是上同调函子，可将它提升为元素
$a_{n + 1} \in H(X_{n + 1})$。

\medskip\noindent
我们断言 $X = \text{hocolim} X_n$ 表示 $H$。对定义特异三角
$$
\bigoplus X_n \to
\bigoplus X_n \to X \to
\bigoplus X_n[1]
$$
应用 $H$，得到正合列
$$
\prod H(X_n) \leftarrow 
\prod H(X_n) \leftarrow 
H(X)
$$
因此存在元素 $a \in H(X)$，它在 $\prod H(X_n)$ 中的像为 $(a_n)$。
从而有自然变换 $\Hom_\mathcal{D}(- , X) \to H$，使得
% P02-E0092: frozen source has the sentence fragment “Hence a natural transformation ...”.
$$
\Hom_\mathcal{D}(-, X_1) \to
\Hom_\mathcal{D}(-, X_2) \to
\Hom_\mathcal{D}(-, X_3) \to \ldots \to
\Hom_\mathcal{D}(-, X) \to H
$$
交换。对每个 $i$，映射 $\Hom_\mathcal{D}(E_i, X) \to H(E_i)$
由 $X_1$ 的构造是满射。另一方面，由 $X_n \to X_{n + 1}$ 的构造，
$\Hom_\mathcal{D}(E_i, X_n) \to H(E_i)$ 的核被映射
$\Hom_\mathcal{D}(E_i, X_n) \to \Hom_\mathcal{D}(E_i, X_{n + 1})$ 消去。
由于
$$
\Hom_\mathcal{D}(E_i, X) = \colim \Hom_\mathcal{D}(E_i, X_n)
$$
由引理 \ref{derived-lemma-commutes-with-countable-sums}，可见
$\Hom_\mathcal{D}(E_i, X) \to H(E_i)$ 是单射。

\medskip\noindent
为完成证明，考虑子范畴
$$
\mathcal{D}' =
\{Y \in \Ob(\mathcal{D}) \mid \Hom_\mathcal{D}(Y[n], X) \to H(Y[n])
\text{ 对所有 }n\text{ 均为同构}\}
$$
由于 $\Hom_\mathcal{D}(-, X) \to H$ 是上同调函子之间的变换，
子范畴 $\mathcal{D}'$ 是 $\mathcal{D}$ 的严格全、饱和三角子范畴
（略去细节；见引理 \ref{derived-lemma-homological-functor-kernel} 的证明）。
此外，由于 $H$ 和 $\Hom_\mathcal{D}(-, X)$ 都把直和变为直积，
$\mathcal{D}'$ 中对象的直和仍在 $\mathcal{D}'$ 中。因此
$\mathcal{D}'$ 中对象的导出余极限仍在 $\mathcal{D}'$ 中。
由于 $\{E_i\}$ 在平移下保持不变，可知对所有 $i$，$E_i$
都是 $\mathcal{D}'$ 的对象。由引理 \ref{derived-lemma-write-as-colimit}，
$\mathcal{D}' = \mathcal{D}$，证明完毕。
\end{proof}

\begin{proposition}
\label{derived-proposition-brown}
\begin{reference}
\cite[Theorem 4.1]{Neeman-Grothendieck}.
\end{reference}
设 $\mathcal{D}$ 为具有直和的紧生成三角范畴。设
$F : \mathcal{D} \to \mathcal{D}'$ 为把直和变为直和的三角范畴正合函子。
则 $F$ 有正合右伴随。
\end{proposition}

\begin{proof}
对 $\mathcal{D}'$ 的对象 $Y$，考虑反变函子
$$
\mathcal{D} \to \textit{Ab},\quad W \mapsto \Hom_{\mathcal{D}'}(F(W), Y)
$$
由于 $F$ 正合，该函子是上同调函子；又因为 $F$ 把直和变为直和，
该函子把直和变为直积。因此，由引理 \ref{derived-lemma-brown}，存在
$\mathcal{D}$ 的对象 $X$，使得
$\Hom_\mathcal{D}(W, X) = \Hom_{\mathcal{D}'}(F(W), Y)$。
伴随的存在性由 Categories，引理 \ref{categories-lemma-adjoint-exists} 得出，
正合性由引理 \ref{derived-lemma-adjoint-is-exact} 得出。
\end{proof}









\section{Brown 可表性之二}
\label{derived-section-brown-bis}

\noindent
本节说明一类具有适当生成元集的三角范畴的 Brown 可表性版本；
其他版本见 \cite{Franke}、\cite{Neeman} 和 \cite{Krause}。

\begin{lemma}
\label{derived-lemma-brown-bis}
\begin{reference}
\cite[Theorem A]{Krause} {\stackszhgreek 的弱化版本。}
\end{reference}
设 $\mathcal{D}$ 为具有直和的三角范畴。假设给定 $\mathcal{D}$
的对象集 $\mathcal{E}$，满足：
\begin{enumerate}
\item 若 $X$ 是 $\mathcal{D}$ 的非零对象，则存在
$E \in \mathcal{E}$ 和非零映射 $E \to X$；
\item 给定 $\mathcal{D}$ 的对象 $X_n$，$n \in \mathbf{N}$，
$E \in \mathcal{E}$，以及 $\alpha : E \to \bigoplus X_n$，
存在 $E_n \in \mathcal{E}$、$\beta_n : E_n \to X_n$ 和态射
$\gamma : E \to \bigoplus E_n$，使得
$\alpha = (\bigoplus \beta_n) \circ \gamma$。
\end{enumerate}
设 $H : \mathcal{D} \to \textit{Ab}$ 为把直和变为直积的反变上同调函子。
则 $H$ 可表。
\end{lemma}

\begin{proof}
本证明与引理 \ref{derived-lemma-brown} 的证明非常相似。可将 $\mathcal{E}$
替换为 $\bigcup_{i \in \mathbf{Z}} \mathcal{E}[i]$，并假设
$\mathcal{E}$ 在平移下保持不变。考虑序对 $(E, a)$，其中
$E \in \mathcal{E}$ 且 $a \in H(E)$，并令
$$
X_1 = \bigoplus\nolimits_{(E, a)} E
$$
由于 $H(X_1) = \prod_{(E, a)} H(E)$，可见 $(a)_{(E, a)}$
定义元素 $a_1 \in H(X_1)$。令 $H_1 = \Hom_\mathcal{D}(- , X_1)$。
由 Yoneda 引理（Categories，引理 \ref{categories-lemma-yoneda}），
元素 $a_1$ 定义自然变换 $H_1 \to H$。

\medskip\noindent
我们将归纳构造 $X_n$ 和变换 $a_n : H_n \to H$，其中
$H_n = \Hom_\mathcal{D}(-, X_n)$。具体地，对函子
$\Ker(H_n \to H)$ 应用上面的步骤，得到对象
$$
K_{n + 1} = \bigoplus\nolimits_{(E, k),\ k \in \Ker(H_n(E) \to H(E))} E
$$
以及变换 $\Hom_\mathcal{D}(-, K_{n + 1}) \to \Ker(H_n \to H)$。
由 Yoneda 引理，复合 $\Hom_\mathcal{D}(-, K_{n + 1}) \to H_n$
给出态射 $K_{n + 1} \to X_n$。在 $\mathcal{D}$ 中选择特异三角
$$
K_{n + 1} \to X_n \to X_{n + 1} \to K_{n + 1}[1]
$$
由构造，元素 $a_n \in H(X_n)$ 在 $H(K_{n + 1})$ 中的像为零。
由于 $H$ 是上同调函子，可将它提升为元素
$a_{n + 1} \in H(X_{n + 1})$。

\medskip\noindent
令 $X = \text{hocolim} X_n$。对定义特异三角
$$
\bigoplus X_n \to
\bigoplus X_n \to X \to
\bigoplus X_n[1]
$$
应用 $H$，得到正合列
$$
\prod H(X_n) \leftarrow
\prod H(X_n) \leftarrow
H(X)
$$
因此存在元素 $a \in H(X)$，它在 $\prod H(X_n)$ 中的像为 $(a_n)$。
从而有自然变换 $\Hom_\mathcal{D}(- , X) \to H$，使得
% P02-E0093: frozen source repeats the sentence fragment “Hence a natural transformation ...”.
$$
\Hom_\mathcal{D}(-, X_1) \to
\Hom_\mathcal{D}(-, X_2) \to
\Hom_\mathcal{D}(-, X_3) \to \ldots \to
\Hom_\mathcal{D}(-, X) \to H
$$
交换。我们断言 $\Hom_\mathcal{D}(-, X) \to H(-)$ 是同构。

\medskip\noindent
设 $E \in \mathcal{E}$。先证明
$$
\Hom_\mathcal{D}(E, \bigoplus X_n) \to  \Hom_\mathcal{D}(E, \bigoplus X_n)
$$
是单射。事实上，设 $\alpha : E \to \bigoplus X_n$。
则由假设 (2)，得到分解 $\alpha = (\bigoplus \beta_n) \circ \gamma$。
由于由构造 $E_n \to X_n \to X_{n + 1}$ 为零，可见复合
$\bigoplus E_n \to \bigoplus X_n \to \bigoplus X_n$
等于 $\bigoplus \beta_n$。因而复合
$E \to \bigoplus X_n \to \bigoplus X_n$ 也等于 $\alpha$。
这证明所述单射性，并且也证明
$$
\Hom_\mathcal{D}(E, \bigoplus X_n[1]) \to \Hom_\mathcal{D}(E, \bigoplus X_n[1])
$$
是单射。因此，对所有 $E \in \mathcal{E}$，都有正合列
$$
\Hom_\mathcal{D}(E, \bigoplus X_n) \to
\Hom_\mathcal{D}(E, \bigoplus X_n) \to
\Hom_\mathcal{D}(E, X) \to 0
$$

\medskip\noindent
设 $E \in \mathcal{E}$，并设 $f : E \to X$ 为态射。
由上一段，可选择提升 $f$ 的 $\alpha : E \to \bigoplus X_n$。
再由假设 (2)，得到分解 $\alpha = (\bigoplus \beta_n) \circ \gamma$。
对每个 $n$，存在态射 $\delta_n : E_n \to X_1$，使得
$\delta_n$ 和 $\beta_n$ 在 $H(E_n)$ 中的像相同。于是由
$X_n \to X_{n + 1}$ 的构造，复合
$$
E_n \to X_n \to X_{n + 1}
\quad\text{且}\quad
E_n \to X_1 \to X_{n + 1}
$$
相等。因此复合
$$
\bigoplus E_n \to \bigoplus X_n \to X
\quad\text{且}\quad
\bigoplus E_n \to \bigoplus X_1 \to X
$$
也相同。注意到 $\bigoplus X_1 \to X$ 分解为
$\bigoplus X_1 \to X_1 \to X$，可知
$$
\Hom_\mathcal{D}(E, X_1) \to \Hom_\mathcal{D}(E, X)
$$
是满射。由构造，映射 $\Hom_\mathcal{D}(E, X_1) \to H(E)$ 是满射，
且该映射的核被 $\Hom_\mathcal{D}(E, X_1) \to \Hom_\mathcal{D}(E, X)$
消去，故对所有 $E \in \mathcal{E}$，
$\Hom_\mathcal{D}(E, X) \to H(E)$ 都是双射。

\medskip\noindent
为完成证明，考虑子范畴
$$
\mathcal{D}' =
\{Y \in \Ob(\mathcal{D}) \mid \Hom_\mathcal{D}(Y[n], X) \to H(Y[n])
\text{ 对所有 }n\text{ 均为同构}\}
$$
由于 $\Hom_\mathcal{D}(-, X) \to H$ 是上同调函子之间的变换，
子范畴 $\mathcal{D}'$ 是 $\mathcal{D}$ 的严格全、饱和三角子范畴
（略去细节；见引理 \ref{derived-lemma-homological-functor-kernel} 的证明）。
此外，由于 $H$ 和 $\Hom_\mathcal{D}(-, X)$ 都把直和变为直积，
$\mathcal{D}'$ 中对象的直和仍在 $\mathcal{D}'$ 中。因此
$\mathcal{D}'$ 中对象的导出余极限仍在 $\mathcal{D}'$ 中。
由于 $\mathcal{E}$ 在平移下保持不变，由上一段的结果可知
$\mathcal{E} \subset \Ob(\mathcal{D}')$。为完成证明，还需证明
$\mathcal{D}' = \mathcal{D}$。

\medskip\noindent
设 $Y$ 为 $\mathcal{D}$ 的对象，并令 $H(-) = \Hom_\mathcal{D}(-, Y)$。
则 $H$ 是把直和变为直积的上同调函子。由证明第一部分的构造，
得到态射 $\colim X_n = X \to Y$，使得对所有
% P02-E0094: frozen source writes colim X_n although X was defined as hocolim X_n.
$E \in \mathcal{E}$，$\Hom_\mathcal{D}(E, X) \to \Hom_\mathcal{D}(E, Y)$
都是双射。假设 (1) 遂表明 $X \to Y$ 是同构！另一方面，由构造，
$X_1, X_2, \ldots$ 都属于 $\mathcal{D}'$，因而 $X$ 也属于其中。
于是 $Y \in \mathcal{D}'$，证明完毕。
\end{proof}

\begin{proposition}
\label{derived-proposition-brown-bis}
设 $\mathcal{D}$ 为具有直和的三角范畴。假设存在 $\mathcal{D}$ 的对象集
$\mathcal{E}$，满足引理 \ref{derived-lemma-brown-bis} 的条件 (1) 和 (2)。
设 $F : \mathcal{D} \to \mathcal{D}'$ 为把直和变为直和的三角范畴正合函子。
则 $F$ 有正合右伴随。
\end{proposition}

\begin{proof}
对 $\mathcal{D}'$ 的对象 $Y$，考虑反变函子
$$
\mathcal{D} \to \textit{Ab},\quad W \mapsto \Hom_{\mathcal{D}'}(F(W), Y)
$$
由于 $F$ 正合，该函子是上同调函子；又因为 $F$ 把直和变为直和，
该函子把直和变为直积。因此，由引理 \ref{derived-lemma-brown-bis}，存在
$\mathcal{D}$ 的对象 $X$，使得
$\Hom_\mathcal{D}(W, X) = \Hom_{\mathcal{D}'}(F(W), Y)$。
伴随的存在性由 Categories，引理 \ref{categories-lemma-adjoint-exists} 得出，
正合性由引理 \ref{derived-lemma-adjoint-is-exact} 得出。
\end{proof}








\section{容许子范畴}
\label{derived-section-admissible}

\noindent
本节的参考文献是 \cite[Section 1]{Bondal-Kapranov}。

\begin{definition}
\label{derived-definition-orthogonal}
设 $\mathcal{D}$ 为加性范畴，$\mathcal{A} \subset \mathcal{D}$ 为全子范畴。
$\mathcal{A}$ 的{\it 右正交} $\mathcal{A}^\perp$ 是如下对象组成的全子范畴：
$\mathcal{D}$ 的对象 $X$ 满足，对所有 $A \in \Ob(\mathcal{A})$，
$\Hom(A, X) = 0$。$\mathcal{A}$ 的{\it 左正交}
${}^\perp\mathcal{A}$ 是如下对象组成的全子范畴：$\mathcal{D}$ 的对象
$X$ 满足，对所有 $A \in \Ob(\mathcal{A})$，$\Hom(X, A) = 0$。
\end{definition}

\begin{lemma}
\label{derived-lemma-pre-prepare-adjoint}
设 $\mathcal{D}$ 为三角范畴，$\mathcal{A} \subset \mathcal{D}$
为在所有平移下不变的全子范畴。考虑 $\mathcal{D}$ 的特异三角
$$
X \to Y \to Z \to X[1]
$$
下列条件等价：
\begin{enumerate}
\item $Z$ 属于 $\mathcal{A}^\perp$；
\item 对所有 $A \in \Ob(\mathcal{A})$，有
$\Hom(A, X) = \Hom(A, Y)$。
\end{enumerate}
\end{lemma}

\begin{proof}
由引理 \ref{derived-lemma-representable-homological}，函子
$\Hom(A, -)$ 是同调函子，因而得到如
(\ref{derived-equation-long-exact-cohomology-sequence}) 中的长正合列。
假设 (1)，并令 $A \in \Ob(\mathcal{A})$。考虑正合列
$$
\Hom(A[1], Z) \to \Hom(A, X) \to \Hom(A, Y) \to \Hom(A, Z)
$$
由于 $A[1] \in \Ob(\mathcal{A})$，第一个和最后一个群均为零，
故得到 (2)。假设 (2)，并令 $A \in \Ob(\mathcal{A})$。
考虑正合列
$$
\Hom(A, X) \to \Hom(A, Y) \to \Hom(A, Z) \to \Hom(A[-1], X) \to \Hom(A[-1], Y)
$$
遂如所需得到 $\Hom(A, Z) = 0$。
\end{proof}

\begin{lemma}
\label{derived-lemma-pre-prepare-adjoint-dual}
设 $\mathcal{D}$ 为三角范畴，$\mathcal{B} \subset \mathcal{D}$
为在所有平移下不变的全子范畴。考虑 $\mathcal{D}$ 的特异三角
$$
X \to Y \to Z \to X[1]
$$
下列条件等价：
\begin{enumerate}
\item $X$ 属于 ${}^\perp\mathcal{B}$；
\item 对所有 $B \in \Ob(\mathcal{B})$，有
$\Hom(Y, B) = \Hom(Z, B)$。
\end{enumerate}
\end{lemma}

\begin{proof}
这是引理 \ref{derived-lemma-pre-prepare-adjoint} 的对偶。
\end{proof}

\begin{lemma}
\label{derived-lemma-orthogonal-triangulated}
设 $\mathcal{D}$ 为三角范畴，$\mathcal{A} \subset \mathcal{D}$
为在所有平移下不变的全子范畴。则 $\mathcal{A}$ 的右正交
$\mathcal{A}^\perp$ 和左正交 ${}^\perp\mathcal{A}$ 都是
$\mathcal{D}$ 的严格全、饱和\footnote{定义 \ref{derived-definition-saturated}。}
三角子范畴。
% P02-E0095: frozen source misspells “subcategories” as “subcagories”.
\end{lemma}

\begin{proof}
由定义立即可知，正交子范畴在平移、直和和直和因子下保持不变。
考虑 $\mathcal{D}$ 的特异三角
$$
X \to Y \to Z \to X[1]
$$
由引理 \ref{derived-lemma-triangulated-subcategory}，只需证明若 $X$ 和 $Y$
属于 $\mathcal{A}^\perp$，则 $Z$ 也属于 $\mathcal{A}^\perp$。
这立即由引理 \ref{derived-lemma-pre-prepare-adjoint} 得出。
\end{proof}

\begin{lemma}
\label{derived-lemma-prepare-adjoint}
设 $\mathcal{D}$ 为三角范畴，$\mathcal{A}$ 为 $\mathcal{D}$ 的全三角子范畴。
对 $\mathcal{D}$ 的对象 $X$，考虑性质 $P(X)$：在 $\mathcal{D}$ 中存在
特异三角 $A \to X \to B \to A[1]$，其中 $A$ 属于 $\mathcal{A}$，
$B$ 属于 $\mathcal{A}^\perp$。
\begin{enumerate}
\item 若 $X_1 \to X_2 \to X_3 \to X_1[1]$ 是特异三角，且 $P$
对三个对象中的两个成立，则它对第三个也成立。
\item 若 $P$ 对 $X_1$ 和 $X_2$ 成立，则它对 $X_1 \oplus X_2$ 也成立。
\end{enumerate}
\end{lemma}

\begin{proof}
设 $X_1 \to X_2 \to X_3 \to X_1[1]$ 为特异三角，并假设 $P$
对 $X_1$ 和 $X_2$ 成立。选择性质 $P$ 中的特异三角
$$
A_1 \to X_1 \to B_1 \to A_1[1]
\quad\text{且}\quad
A_2 \to X_2 \to B_2 \to A_2[1]
$$
由于由引理 \ref{derived-lemma-pre-prepare-adjoint} 有
$\Hom(A_1, A_2) = \Hom(A_1, X_2)$，存在唯一态射
$A_1 \to A_2$，使图
$$
\xymatrix{
A_1 \ar[d] \ar[r] & X_1 \ar[d] \\
A_2 \ar[r] & X_2
}
$$
交换。选择它到下图的一个扩张
$$
\xymatrix{
A_1 \ar[r] \ar[d] & X_1 \ar[r] \ar[d] & Q_1 \ar[r] \ar[d] & A_1[1] \ar[d] \\
A_2 \ar[r] \ar[d] & X_2 \ar[r] \ar[d] & Q_2 \ar[r] \ar[d] & A_2[1] \ar[d] \\
A_3 \ar[r] \ar[d] & X_3 \ar[r] \ar[d] & Q_3 \ar[r] \ar[d] & A_3[1] \ar[d] \\
A_1[1] \ar[r] & X_1[1] \ar[r] & Q_1[1] \ar[r] & A_1[2]
}
$$
如命题 \ref{derived-proposition-9} 中所述。由 TR3，$Q_1 \cong B_1$ 且
$Q_2 \cong B_2$，故 $Q_1, Q_2 \in \Ob(\mathcal{A}^\perp)$。
由于 $Q_1 \to Q_2 \to Q_3 \to Q_1[1]$ 是特异三角，
由引理 \ref{derived-lemma-orthogonal-triangulated}，有
$Q_3 \in \Ob(\mathcal{A}^\perp)$。由于 $\mathcal{A}$ 是全三角子范畴，
$A_3$ 与 $\mathcal{A}$ 的某个对象同构。因此 $X_3$ 满足 $P$。
(1) 的其他情形由平移归结为此情形。(2) 经引理 \ref{derived-lemma-split}
是 (1) 的特例。
\end{proof}

\begin{lemma}
\label{derived-lemma-prepare-adjoint-dual}
设 $\mathcal{D}$ 为三角范畴，$\mathcal{B}$ 为 $\mathcal{D}$ 的全三角子范畴。
对 $\mathcal{D}$ 的对象 $X$，考虑性质 $P(X)$：在 $\mathcal{D}$ 中存在
特异三角 $A \to X \to B \to A[1]$，其中 $B$ 属于 $\mathcal{B}$，
$A$ 属于 ${}^\perp\mathcal{B}$。
\begin{enumerate}
\item 若 $X_1 \to X_2 \to X_3 \to X_1[1]$ 是特异三角，且 $P$
对三个对象中的两个成立，则它对第三个也成立。
\item 若 $P$ 对 $X_1$ 和 $X_2$ 成立，则它对 $X_1 \oplus X_2$ 也成立。
\end{enumerate}
\end{lemma}

\begin{proof}
这是引理 \ref{derived-lemma-prepare-adjoint} 的对偶。
\end{proof}

\begin{lemma}
\label{derived-lemma-right-adjoint}
设 $\mathcal{D}$ 为三角范畴，$\mathcal{A} \subset \mathcal{D}$
为全三角子范畴。下列条件等价：
\begin{enumerate}
\item 嵌入函子 $\mathcal{A} \to \mathcal{D}$ 有右伴随；
\item 对 $\mathcal{D}$ 中每个 $X$，存在 $\mathcal{D}$ 中的特异三角
$$
A \to X \to B \to A[1]
$$
其中 $A \in \Ob(\mathcal{A})$ 且 $B \in \Ob(\mathcal{A}^\perp)$。
\end{enumerate}
若这些条件成立，则 $\mathcal{A}$ 是饱和的
（定义 \ref{derived-definition-saturated}）；若 $\mathcal{A}$ 是
$\mathcal{D}$ 的严格全子范畴，则
$\mathcal{A} = {}^\perp(\mathcal{A}^\perp)$。
\end{lemma}

\begin{proof}
假设 (1)，并以 $v : \mathcal{D} \to \mathcal{A}$ 表示右伴随。
设 $X \in \Ob(\mathcal{D})$，令 $A = v(X)$。可将伴随映射
$A \to X$ 扩张为特异三角 $A \to X \to B \to A[1]$。
由于对 $A' \in \Ob(\mathcal{A})$，有
$$
\Hom_\mathcal{A}(A', A) =
\Hom_\mathcal{A}(A', v(X)) =
\Hom_\mathcal{D}(A', X)
$$
由引理 \ref{derived-lemma-pre-prepare-adjoint}，可知
$B \in \Ob(\mathcal{A}^\perp)$。

\medskip\noindent
假设 (2)。我们显式构造伴随 $v$。设 $X \in \Ob(\mathcal{D})$，
选择 (2) 中的 $A \to X \to B \to A[1]$，并令 $v(X) = A$。
设 $f : X \to Y$ 为 $\mathcal{D}$ 中的态射。选择 (2) 中的
$A' \to Y \to B' \to A'[1]$。由于由引理
\ref{derived-lemma-pre-prepare-adjoint} 有 $\Hom(A, A') = \Hom(A, Y)$，
存在唯一态射 $f' : A \to A'$，使图
$$
\xymatrix{
A \ar[d]_{f'} \ar[r] & X \ar[d]^f \\
A' \ar[r] & Y
}
$$
交换。因此可令 $v(f) = f'$，从而得到函子。
再次使用引理 \ref{derived-lemma-pre-prepare-adjoint}，可见 $v$ 是嵌入的伴随。

\medskip\noindent
证明最后一个陈述。为证明 $\mathcal{A}$ 饱和，可以把 $\mathcal{A}$
替换为与 $\mathcal{A}$ 具有相同同构类的严格全子范畴；略去细节。假设 $\mathcal{A}$
严格全。若能证明 $\mathcal{A} = {}^\perp(\mathcal{A}^\perp)$，
则由引理 \ref{derived-lemma-orthogonal-triangulated}，$\mathcal{A}$ 将是饱和的。
包含关系 $\mathcal{A} \subset {}^\perp(\mathcal{A}^\perp)$ 是显然的，
% P02-E0096: frozen source misspells “inclusion” as “incusion”.
故只需证明另一个包含关系。设 $X$ 为
${}^\perp(\mathcal{A}^\perp)$ 的对象。选择 (2) 中的特异三角
$A \to X \to B \to A[1]$。由于由假设有 $\Hom(X, B) = 0$，
由引理 \ref{derived-lemma-split}，可见 $A \cong X \oplus B[-1]$。
由于 $B \in \mathcal{A}^\perp$，有 $\Hom(A, B[-1]) = 0$，
这蕴含 $B[-1] = 0$ 且 $A \cong X$，正如所需。
\end{proof}

\begin{lemma}
\label{derived-lemma-left-adjoint}
设 $\mathcal{D}$ 为三角范畴，$\mathcal{B} \subset \mathcal{D}$
为全三角子范畴。下列条件等价：
\begin{enumerate}
\item 嵌入函子 $\mathcal{B} \to \mathcal{D}$ 有左伴随；
\item 对 $\mathcal{D}$ 中每个 $X$，存在 $\mathcal{D}$ 中的特异三角
$$
A \to X \to B \to A[1]
$$
其中 $B \in \Ob(\mathcal{B})$ 且 $A \in \Ob({}^\perp\mathcal{B})$。
\end{enumerate}
若这些条件成立，则 $\mathcal{B}$ 是饱和的
（定义 \ref{derived-definition-saturated}）；若 $\mathcal{B}$ 是
$\mathcal{D}$ 的严格全子范畴，则
$\mathcal{B} = ({}^\perp\mathcal{B})^\perp$。
\end{lemma}

\begin{proof}
这是引理 \ref{derived-lemma-right-adjoint} 的对偶。
\end{proof}

\begin{definition}
\label{derived-definition-admissible}
设 $\mathcal{D}$ 为三角范畴。$\mathcal{D}$ 的{\it 右容许}子范畴，
是满足引理 \ref{derived-lemma-right-adjoint} 中等价条件的严格全三角子范畴。
$\mathcal{D}$ 的{\it 左容许}子范畴，是满足引理
\ref{derived-lemma-left-adjoint} 中等价条件的严格全三角子范畴。
同时右容许和左容许的子范畴称为{\it 双侧容许}子范畴。
\end{definition}

\noindent
设 $\mathcal{A}$ 为三角范畴 $\mathcal{D}$ 的右容许子范畴。
注意到，对 $X \in \mathcal{D}$，特异三角
$$
A \to X \to B \to A[1]
$$
其中 $A \in \mathcal{A}$ 且 $B \in \mathcal{A}^\perp$，在如下意义下是
典范的：对任意另一个满足 $A' \in \mathcal{A}$、
$B' \in \mathcal{A}^\perp$ 的特异三角 $A' \to X \to B' \to A'[1]$，
都存在三角的同构
$(\alpha, \text{id}_X, \beta) : (A, X, B) \to (A', X, B')$。
下面的命题总结以上所述。

\begin{proposition}
\label{derived-proposition-summarize-admissible}
设 $\mathcal{D}$ 为三角范畴，$\mathcal{A} \subset \mathcal{D}$ 和
$\mathcal{B} \subset \mathcal{D}$ 为子范畴。下列条件等价：
\begin{enumerate}
\item $\mathcal{A}$ 右容许，且 $\mathcal{B} = \mathcal{A}^\perp$；
\item $\mathcal{B}$ 左容许，且 $\mathcal{A} = {}^\perp\mathcal{B}$；
\item 对所有 $A \in \mathcal{A}$ 和 $B \in \mathcal{B}$，有
$\Hom(A, B) = 0$；并且对 $\mathcal{D}$ 中每个 $X$，存在
$\mathcal{D}$ 中的特异三角 $A \to X \to B \to A[1]$，其中
$A \in \mathcal{A}$ 且 $B \in \mathcal{B}$。
\end{enumerate}
若这些条件成立，则 $\mathcal{A} \to \mathcal{D}/\mathcal{B}$ 和
$\mathcal{B} \to \mathcal{D}/\mathcal{A}$ 是三角范畴的等价；
嵌入函子 $\mathcal{A} \to \mathcal{D}$ 的右伴随是
$\mathcal{D} \to \mathcal{D}/\mathcal{B} \to \mathcal{A}$；
嵌入函子 $\mathcal{B} \to \mathcal{D}$ 的左伴随是
$\mathcal{D} \to \mathcal{D}/\mathcal{A} \to \mathcal{B}$。
\end{proposition}

\begin{proof}
(1)、(2)、(3) 的等价性直接由引理 \ref{derived-lemma-right-adjoint} 和
\ref{derived-lemma-left-adjoint} 得出（略去一个小细节）。以
$v : \mathcal{D} \to \mathcal{A}$ 表示嵌入函子
$i : \mathcal{A} \to \mathcal{D}$ 的右伴随。立即有
$\Ker(v) = \mathcal{A}^\perp = \mathcal{B}$。因此，由商的泛性质，
$v$ 经函子 $\overline{v} : \mathcal{D}/\mathcal{B} \to \mathcal{A}$ 分解。
由于由 Categories，引理 \ref{categories-lemma-adjoint-fully-faithful} 有
$v \circ i = \text{id}_\mathcal{A}$，可见 $\overline{v}$ 是
$\overline{i} : \mathcal{A} \to \mathcal{D}/\mathcal{B}$ 的左拟逆。
我们还断言复合 $\overline{i} \circ \overline{v}$ 同构于
$\text{id}_{\mathcal{D}/\mathcal{B}}$。事实上，假设 $X$ 嵌入 (3) 中的
特异三角 $A \to X \to B \to A[1]$。如引理
\ref{derived-lemma-right-adjoint} 的证明中所见，$v(X) = A$。
把 $X$ 看作 $\mathcal{D}/\mathcal{B}$ 的对象，则有
$\overline{i}(\overline{v}(X)) = A$，并且在 $\mathcal{D}/\mathcal{B}$ 中
存在函子性同构 $\overline{i}(\overline{v}(X)) = A \to X$。
因此，确实有 $\overline{v} : \mathcal{D}/\mathcal{B} \to \mathcal{A}$
为等价。证明 $\mathcal{B} \to \mathcal{D}/\mathcal{A}$ 是等价，
且嵌入函子 $\mathcal{B} \to \mathcal{D}$ 的左伴随是
$\mathcal{D} \to \mathcal{D}/\mathcal{A} \to \mathcal{B}$，
与上述论证对偶。
\end{proof}




\section{Postnikov 系统}
\label{derived-section-postnikov}

\noindent
本节的参考文献是 \cite{Orlov-K3}。设 $\mathcal{D}$ 为三角范畴，设
$$
X_n \to X_{n - 1} \to \ldots \to X_0
$$
为 $\mathcal{D}$ 中的复形。本节考虑构造该复形的“全复形化”的问题。

\begin{definition}
\label{derived-definition-postnikov-system}
设 $\mathcal{D}$ 为三角范畴，设
$$
X_n \to X_{n - 1} \to \ldots \to X_0
$$
为 $\mathcal{D}$ 中的复形。{\it Postnikov 系统}归纳定义如下。
\begin{enumerate}
\item 若 $n = 0$，它是同构 $Y_0 \to X_0$。
\item 若 $n = 1$，它由同构 $Y_0 \to X_0$ 的一个选择以及特异三角
$$
Y_1 \to X_1 \to Y_0 \to Y_1[1]
$$
的一个选择组成，其中 $X_1 \to Y_0$ 与 $Y_0 \to X_0$ 的复合是
给定态射 $X_1 \to X_0$。
\item 若 $n > 1$，它由复形 $X_{n - 1} \to \ldots \to X_0$
的一个 Postnikov 系统以及特异三角
$$
Y_n \to X_n \to Y_{n - 1} \to Y_n[1]
$$
的一个选择组成，其中态射 $X_n \to Y_{n - 1}$ 与
$Y_{n - 1} \to X_{n - 1}$ 的复合是给定态射 $X_n \to X_{n - 1}$。
\end{enumerate}
给定 $\mathcal{D}$ 中两个等长复形之间的态射
\begin{equation}
\label{derived-equation-map-complexes}
\vcenter{
\xymatrix{
X_n \ar[r] \ar[d] &
X_{n - 1} \ar[r] \ar[d] &
\ldots \ar[r] &
X_0 \ar[d] \\
X'_n \ar[r] &
X'_{n - 1} \ar[r] &
\ldots \ar[r] &
X'_0
}
}
\end{equation}
显然可定义{\it Postnikov 系统的态射}。
\end{definition}

\noindent
下面是一个关键例子。

\begin{example}
\label{derived-example-key-postnikov}
设 $\mathcal{A}$ 为阿贝尔范畴，设 $\ldots \to A_2 \to A_1 \to A_0$
为 $\mathcal{A}$ 中的链复形。于是可以考虑 $D(\mathcal{A})$ 中的对象
$$
X_n = A_n
\quad\text{且}\quad
Y_n = (A_n \to A_{n - 1} \to \ldots \to A_0)[-n]
$$
连同显然的典范映射 $Y_n \to X_n$ 和
$Y_0 \to Y_1[1] \to Y_2[2] \to \ldots$，特异三角
$Y_n \to X_n \to Y_{n - 1} \to Y_n[1]$ 按定义
\ref{derived-definition-postnikov-system} 给出
$\ldots \to X_2 \to X_1 \to X_0$ 的 Postnikov 系统。
这里使用了 Postnikov 系统对 $D(\mathcal{A})$ 中无限复形的显然推广。
最后，若 $\mathcal{A}$ 中以 $\mathbf{N}$ 为指标的余极限存在且正合，则
$$
\text{hocolim} Y_n[n] = (\ldots \to A_2 \to A_1 \to A_0 \to 0 \to \ldots)
$$
在 $D(\mathcal{A})$ 中成立。这立即由引理 \ref{derived-lemma-colim-hocolim} 得出。
\end{example}

\noindent
给定复形 $X_n \to X_{n - 1} \to \ldots \to X_0$ 以及定义
\ref{derived-definition-postnikov-system} 中的 Postnikov 系统，可以考虑映射
$$
Y_0 \to Y_1[1] \to \ldots \to Y_n[n]
$$
这些映射共同嵌入某些特异三角，并与各 $X_i$ 之间的给定映射相容。
当 $n = 3$ 时，图形如下：
$$
\xymatrix{
Y_0 \ar[rr] & &
Y_1[1] \ar[dl] \ar[rr] & &
Y_2[2] \ar[dl] \ar[rr] & &
Y_3[3] \ar[dl] \\
& X_1[1] \ar[lu]_{+1} & &
X_2[2] \ar[ll]_{+1} \ar[lu]_{+1} & &
X_3[3] \ar[ll]_{+1} \ar[lu]_{+1}
}
$$
我们建议读者把 $Y_n[n]$ 看作由 $X_0, X_1[1], \ldots, X_n[n]$ 得到；
例如，若各映射 $X_i \to X_{i - 1}$ 为零，则可以取
$Y_n[n] = \bigoplus_{i = 0, \ldots, n} X_i[i]$。
Postnikov 系统并非总是存在。下面是关于较小 $n$ 的简单引理。

\begin{lemma}
\label{derived-lemma-postnikov-system-small-cases}
设 $\mathcal{D}$ 为三角范畴。考虑长度为 $n$ 的复形的 Postnikov 系统。
\begin{enumerate}
\item 当 $n = 0$ 时，Postnikov 系统总是存在，并且复形的任意态射
(\ref{derived-equation-map-complexes}) 唯一地扩张为 Postnikov 系统的态射。
\item 当 $n = 1$ 时，Postnikov 系统总是存在，并且复形的任意态射
(\ref{derived-equation-map-complexes}) 扩张为 Postnikov 系统的一个（非唯一）态射。
\item 当 $n = 2$ 时，Postnikov 系统总是存在，但一般而言，复形的态射
(\ref{derived-equation-map-complexes}) 不会扩张为 Postnikov 系统的态射。
\item 当 $n > 2$ 时，Postnikov 系统并非总是存在。
\end{enumerate}
\end{lemma}

\begin{proof}
$n = 0$ 的情形立即成立，因为同构可逆。$n = 1$ 的情形立即由
TR1（三角的存在性）和 TR3（把态射扩张为三角的态射）得出。
对 $n = 2$，论证如下。令 $Y_0 = X_0$。由 $n = 1$ 的情形，
可选择 Postnikov 系统
$$
Y_1 \to X_1 \to Y_0 \to Y_1[1]
$$
由于复合 $X_2 \to X_1 \to X_0$ 为零，由引理
\ref{derived-lemma-representable-homological}，可将 $X_2 \to X_1$
（非唯一地）分解为 $X_2 \to Y_1 \to X_1$。然后只需把态射
$X_2 \to Y_1$ 嵌入特异三角
$$
Y_2 \to X_2 \to Y_1 \to Y_2[1]
$$
便得到 $n = 2$ 的 Postnikov 系统。对 $n > 2$，不能类似论证，
因为不知道复合 $X_n \to X_{n - 1} \to Y_{n - 1}$ 在
$\mathcal{D}$ 中是否为零。
% P02-E0097: frozen source's last composition is ill-typed for the preceding Postnikov triangle.
\end{proof}

\begin{lemma}
\label{derived-lemma-maps-postnikov-systems-vanishing}
设 $\mathcal{D}$ 为三角范畴。给定态射
(\ref{derived-equation-map-complexes})，考虑条件
\begin{equation}
\label{derived-equation-P}
\Hom(X_i[i - j - 1], X'_j) = 0 \text{ 对 }i > j + 1
\end{equation}
则：
\begin{enumerate}
\item 若已有复形 $X'_n \to X'_{n - 1} \to \ldots \to X'_0$
的 Postnikov 系统，则性质 (\ref{derived-equation-P}) 蕴含
$$
\Hom(X_i[i - j - 1], Y'_j) = 0 \text{ 对 }i > j + 1
$$
\item 若两个复形都配备了 Postnikov 系统并且
(\ref{derived-equation-P}) 成立，则该态射扩张为 Postnikov 系统的一个（非唯一）态射。
\end{enumerate}
\end{lemma}

\begin{proof}
先对 $j$ 归纳证明 (1)。基础情形 $j = 0$ 无需证明，因为
$Y'_0 \to X'_0$ 是同构。假设结论对 $j - 1$ 成立。考虑特异三角
$$
Y'_j \to X'_j \to Y'_{j - 1} \to Y'_j[1]
$$
引理 \ref{derived-lemma-representable-homological} 的长正合列给出正合列
$$
\Hom(X_i[i - j - 1], Y'_{j - 1}[-1]) \to
\Hom(X_i[i - j - 1], Y'_j) \to
\Hom(X_i[i - j - 1], X'_j)
$$
由归纳假设和 (\ref{derived-equation-P})，两端的群均为零，故结论成立。

\medskip\noindent
证明 (2)。当 $n = 1$ 时，态射的存在性已由引理
\ref{derived-lemma-postnikov-system-small-cases} 证明。若 $n > 1$，
由归纳可假设已经给定长度 $n - 1$ 的 Postnikov 系统态射。
问题在于不知道图
$$
\xymatrix{
X_n \ar[r] \ar[d] & Y_{n - 1} \ar[d] \\
X'_n \ar[r] & Y'_{n - 1}
}
$$
是否交换。以 $\alpha : X_n \to Y'_{n - 1}$ 表示其差。
我们确实知道 $\alpha$ 与 $Y'_{n - 1} \to X'_{n - 1}$ 的复合为零
（这是长度 $n - 1$ 的 Postnikov 系统态射的含义）。由特异三角
$Y'_{n - 1} \to X'_{n - 1} \to Y'_{n - 2} \to Y'_{n - 1}[1]$，
这意味着 $\alpha$ 是 $Y'_{n - 2}[-1] \to Y'_{n - 1}$ 与某个映射
$\alpha' : X_n \to Y'_{n - 2}[-1]$ 的复合。由本引理 (1)，条件
(\ref{derived-equation-P}) 保证 $\alpha'$ 为零，故 $\alpha$ 为零。
为完成存在性的证明，交换性保证可由 TR3 选择下图中的虚线箭头：
$$
\xymatrix{
Y_{n - 1}[-1] \ar[d] \ar[r] &
Y_n \ar[r] \ar@{..>}[d] &
X_n \ar[r] \ar[d] &
Y_{n - 1} \ar[d] \\
Y'_{n - 1}[-1] \ar[r] &
Y'_n \ar[r] &
X'_n \ar[r] &
Y'_{n - 1}
}
$$
\end{proof}

\begin{lemma}
\label{derived-lemma-uniqueness-maps-postnikov-systems}
设 $\mathcal{D}$ 为三角范畴。给定态射
(\ref{derived-equation-map-complexes})，并假设两个复形都配备了 Postnikov 系统。
若下列任一条件成立：
\begin{enumerate}
\item 对 $i = 1, \ldots, n$，$\Hom(X_i[i], Y'_n[n]) = 0$；或
\item 对 $i = 1, \ldots, n$，$\Hom(Y_n[n], X'_{n - i}[n - i]) = 0$；或
\item 对 $j \geq i > 0$，有 $\Hom(X_{j - i}[-i + 1], X'_j) = 0$ 且
$\Hom(X_j, X'_{j - i}[-i]) = 0$，
\end{enumerate}
则这些 Postnikov 系统之间至多存在一个态射。
\end{lemma}

\begin{proof}
证明 (1)。考察下图：
$$
\xymatrix{
Y_0 \ar[r] \ar[d] &
Y_1[1] \ar[r] \ar[ld] &
Y_2[2] \ar[r] \ar[lld] &
\ldots \ar[r] &
Y_n[n] \ar[lllld] \\
Y'_n[n]
}
$$
各箭头是态射 $Y_n[n] \to Y'_n[n]$ 与态射
$Y_i[i] \to Y_n[n]$ 的复合。箭头 $Y_0 \to Y'_n[n]$ 是确定的，
因为它是复合 $Y_0 = X_0 \to X'_0 = Y'_0 \to Y'_n[n]$。
由于有特异三角 $Y_0 \to Y_1[1] \to X_1[1]$，
$\Hom(X_1[1], Y'_n[n]) = 0$ 保证第二条竖直箭头唯一。
由于有特异三角 $Y_1[1] \to Y_2[2] \to X_2[2]$，
$\Hom(X_2[2], Y'_n[n]) = 0$ 保证第三条竖直箭头唯一。依此类推。

\medskip\noindent
证明 (2)。复合 $Y_n[n] \to Y'_n[n] \to X_n[n]$ 与复合
% P02-E0098: frozen source ends this composition at X_n[n], although the primed system requires X'_n[n].
$Y_n[n] \to X_n[n] \to X'_n[n]$ 相同，因而唯一。
再用特异三角 $Y'_{n - 1}[n - 1] \to Y'_n[n] \to X'_n[n]$，
可见只需证明 $\Hom(Y_n[n], Y'_{n - 1}[n - 1]) = 0$。
利用特异三角
$$
Y'_{n - i - 1}[n - i - 1] \to Y'_{n - i}[n - i] \to X'_{n - i}[n - i]
$$
由假设得到该消失。略去小细节。

\medskip\noindent
证明 (3)。考察引理 \ref{derived-lemma-maps-postnikov-systems-vanishing} 的证明并
对 $n$ 归纳，只需证明三角态射
$$
\xymatrix{
Y_{n - 1}[-1] \ar[d] \ar[r] &
Y_n \ar[r] \ar@{..>}[d] &
X_n \ar[r] \ar[d] &
Y_{n - 1} \ar[d] \\
Y'_{n - 1}[-1] \ar[r] &
Y'_n \ar[r] &
X'_n \ar[r] &
Y'_{n - 1}
}
$$
中的虚线箭头唯一。由引理 \ref{derived-lemma-uniqueness-third-arrow} (5)，
只需证明 $\Hom(Y_{n - 1}, X'_n) = 0$ 和
$\Hom(X_n, Y'_{n - 1}[-1]) = 0$。
为证明第一个消失，对 $i > 0$ 使用特异三角
$Y_{n - i - 1}[-i] \to Y_{n - i}[-(i - 1)] \to X_{n - i}[-(i - 1)]$，
并对 $i$ 归纳，可见所假设的
$\Hom(X_{n - i}[-i + 1], X'_n)$ 的消失已经足够。
对第二个消失，类似地使用特异三角
$Y'_{n - i - 1}[-i - 1] \to Y'_{n - i}[-i] \to X'_{n - i}[-i]$，
可见所假设的 $\Hom(X_n, X'_{n - i}[-i])$ 的消失也足够。
\end{proof}

\begin{lemma}
\label{derived-lemma-existence-postnikov-system}
设 $\mathcal{D}$ 为三角范畴。设
$X_n \to X_{n - 1} \to \ldots \to X_0$ 为 $\mathcal{D}$ 中的复形。
若
$$
\Hom(X_i[i - j - 2], X_j) = 0 \text{ 对 }i > j + 2
$$
则存在 Postnikov 系统。若
$$
\Hom(X_i[i - j - 1], X_j) = 0 \text{ 对 }i > j + 1
$$
则任意两个 Postnikov 系统同构。
\end{lemma}

\begin{proof}
对 $n$ 归纳。$n = 0, 1, 2$ 的情形由引理
\ref{derived-lemma-postnikov-system-small-cases} 得出。假设 $n > 2$。
假设已给定复形 $X_{n - 1} \to X_{n - 2} \to \ldots \to X_0$
的 Postnikov 系统。把它扩张为长度 $n$ 的 Postnikov 系统的唯一障碍，
是必须找到态射 $X_n \to Y_{n - 1}$，使复合
$X_n \to Y_{n - 1} \to X_{n - 1}$ 等于给定映射
$X_n \to X_{n - 1}$。考虑特异三角
$$
Y_{n - 1} \to X_{n - 1} \to Y_{n - 2} \to Y_{n - 1}[1]
$$
以及由它和函子 $\Hom(X_n, -)$ 给出的长正合列
（见引理 \ref{derived-lemma-representable-homological}），可知只需证明复合
$X_n \to X_{n - 1} \to Y_{n - 2}$ 为零。
由于已知 $X_n \to X_{n - 1} \to X_{n - 2}$ 为零，
应用特异三角
$$
Y_{n - 2} \to X_{n - 2} \to Y_{n - 3} \to Y_{n - 2}[1]
$$
可见只要 $\Hom(X_n, Y_{n - 3}[-1]) = 0$ 即可。
完全仿照引理 \ref{derived-lemma-maps-postnikov-systems-vanishing} (1) 的证明，
读者容易看出这由本引理中的条件推出。

\medskip\noindent
关于同构的陈述，由引理 \ref{derived-lemma-maps-postnikov-systems-vanishing} (2)
所证明的、扩张复形恒等映射的 Postnikov 系统之间态射的存在性，以及引理
\ref{derived-lemma-third-isomorphism-triangle}（它表明所有映射均为同构）得出。
\end{proof}









\section{本质常值系统}
\label{derived-section-essentially-constant}

\noindent
先给出三角范畴中本质常值系统的一些预备引理。

\begin{lemma}
\label{derived-lemma-essentially-constant}
设 $\mathcal{D}$ 为三角范畴，$(A_i)$ 为 $\mathcal{D}$ 中的逆系统。
则 $(A_i)$ 是本质常值的（见 Categories，定义
\ref{categories-definition-essentially-constant-diagram}），当且仅当存在
$i$，并且对所有 $j \geq i$，存在直和分解 $A_j = A \oplus Z_j$，使得
(a) 各映射 $A_{j'} \to A_j$ 与直和分解相容，并且在 $A$ 上为恒等映射；
(b) 对所有 $j \geq i$，存在某个 $j' \geq j$，使 $Z_{j'} \to Z_j$ 为零。
\end{lemma}

\begin{proof}
假设 $(A_i)$ 是值为 $A$ 的本质常值系统。则 $A = \lim A_i$，
并且存在 $i$ 和态射 $A_i \to A$，使得 (1) 复合
$A \to A_i \to A$ 在 $A$ 上为恒等映射；(2) 对所有 $j \geq i$，
存在 $j' \geq j$，使 $A_{j'} \to A_j$ 分解为
$A_{j'} \to A_i \to A \to A_j$。
由 (1)，可知对 $j \geq i$，映射 $A \to A_j$ 与
$A_j \to A_i \to A$ 的复合在 $A$ 上为恒等映射。因此
$A_j \to A$ 有核 $Z_j$，并且映射 $A \oplus Z_j \to A_j$ 是同构，
见引理 \ref{derived-lemma-when-split} 和 \ref{derived-lemma-split}。
这些直和分解显然满足 (a)。由 (2)，对所有 $j$，存在
$j' \geq j$，使 $Z_{j'} \to Z_j$ 为零，故 (b) 成立。
略去反向的证明。
\end{proof}

\begin{lemma}
\label{derived-lemma-essentially-constant-2-out-of-3}
设 $\mathcal{D}$ 为三角范畴。设
$$
A_n \to B_n \to C_n \to A_n[1]
$$
为 $\mathcal{D}$ 中特异三角的逆系统。若 $(A_n)$ 和 $(C_n)$
本质常值，则 $(B_n)$ 也本质常值，并且它们的值嵌入特异三角
$A \to B \to C \to A[1]$；此外，对某个 $n \geq 1$，存在特异三角的态射
$$
\xymatrix{
A_n \ar[d] \ar[r] &
B_n \ar[d] \ar[r] &
C_n \ar[d] \ar[r] &
A_n[1] \ar[d] \\
A \ar[r] &
B \ar[r] &
C \ar[r] &
A[1]
}
$$
它诱导同构 $\lim_{n' \geq n} A_{n'} \to A$，对 $B$ 和 $C$ 亦然。
\end{lemma}

\begin{proof}
重新编号后，可假设 $A_n = A \oplus A'_n$ 且
$C_n = C \oplus C'_n$，其中逆系统 $(A'_n)$ 和 $(C'_n)$
本质为零，见引理 \ref{derived-lemma-essentially-constant}。特别地，对所有 $n$，
态射
$$
C \oplus C'_n \to (A \oplus A'_n)[1]
$$
把直和因子 $C$ 映入直和因子 $A[1]$，其上的映射
$\delta : C \to A[1]$ 与 $n$ 无关。选择特异三角
$$
A \to B \to C \xrightarrow{\delta} A[1]
$$
接着选择特异三角的态射
$$
(A_1 \to B_1 \to C_1 \to A_1[1]) \to
(A \to B \to C \to A[1])
$$
这由 TR3 可以做到。对 $\mathcal{D}$ 的任意对象 $D$，它诱导交换图
$$
\xymatrix{
\ldots \ar[r] &
\Hom_\mathcal{D}(C, D) \ar[r] \ar[d] &
\Hom_\mathcal{D}(B, D) \ar[r] \ar[d] &
\Hom_\mathcal{D}(A, D) \ar[r] \ar[d] &
\ldots \\
\ldots \ar[r] &
\colim \Hom_\mathcal{D}(C_n, D) \ar[r] &
\colim \Hom_\mathcal{D}(B_n, D) \ar[r] &
\colim \Hom_\mathcal{D}(A_n, D) \ar[r] &
\ldots
}
$$
左、右竖直箭头是同构，它们左边和右边的箭头也都是同构。
因此，由五引理，中间的箭头是同构。由 Categories，评注
\ref{categories-remark-pro-category-copresheaves} 中的讨论，
可知 $(B_n)$ 同构于值为 $B$ 的常值逆系统。
由 Categories，评注 \ref{categories-remark-pro-category}，
这等价于 $(B_n)$ 是值为 $B$ 的本质常值系统，证明完毕。
\end{proof}

\begin{lemma}
\label{derived-lemma-essentially-constant-cohomology}
设 $\mathcal{A}$ 为阿贝尔范畴，$A_n$ 为 $D(\mathcal{A})$ 中对象的逆系统。
假设：
\begin{enumerate}
\item 存在整数 $a \leq b$，使得当 $i \not \in [a, b]$ 时
$H^i(A_n) = 0$；
\item 对所有 $i \in \mathbf{Z}$，$\mathcal{A}$ 中的逆系统
$H^i(A_n)$ 都本质常值。
\end{enumerate}
则 $A_n$ 是 $D(\mathcal{A})$ 中的本质常值系统；其值 $A$ 满足：
对每个 $i \in \mathbf{Z}$，$H^i(A)$ 是常值系统 $H^i(A_n)$ 的值。
\end{lemma}

\begin{proof}
由评注 \ref{derived-remark-truncation-distinguished-triangle}，得到特异三角的逆系统
$$
\tau_{\leq a}A_n \to A_n \to \tau_{\geq a + 1}A_n \to (\tau_{\leq a}A_n)[1]
$$
当然，在 $D(\mathcal{A})$ 中有
$\tau_{\leq a}A_n = H^a(A_n)[-a]$。因此，由假设，它们构成本质常值系统。
对 $b - a$ 归纳，可知逆系统 $\tau_{\geq a + 1}A_n$ 本质常值，
设其值为 $A'$。由引理 \ref{derived-lemma-essentially-constant-2-out-of-3}，
$A_n$ 是本质常值系统。略去关于上同调的陈述的证明
（提示：使用引理 \ref{derived-lemma-essentially-constant-2-out-of-3} 的最后一部分）。
\end{proof}

\begin{lemma}
\label{derived-lemma-pro-isomorphism}
设 $\mathcal{D}$ 为三角范畴。设
$$
A_n \to B_n \to C_n \to A_n[1]
$$
为特异三角的逆系统。若系统 $C_n$ 是 pro-零的（即值为 $0$ 的本质常值系统），
则各映射 $A_n \to B_n$ 确定 pro-对象 $(A_n)$ 与 pro-对象 $(B_n)$
之间的 pro-同构。
\end{lemma}

\begin{proof}
对 $\mathcal{D}$ 的任意对象 $X$，考虑正合列
$$
\colim \Hom(C_n, X) \to
\colim \Hom(B_n, X) \to
\colim \Hom(A_n, X) \to
\colim \Hom(C_n[-1], X) \to
$$
正合性由引理 \ref{derived-lemma-representable-homological} 与 Algebra，引理
\ref{algebra-lemma-directed-colimit-exact} 联合得出。由假设，第一项和最后一项
为零。因此，对所有 $X$，映射
$\colim \Hom(B_n, X) \to \colim \Hom(A_n, X)$ 是同构。
本引理由此及 Categories，评注
\ref{categories-remark-pro-category-copresheaves} 得出。
\end{proof}

\begin{lemma}
\label{derived-lemma-pro-isomorphism-bis}
设 $\mathcal{A}$ 为阿贝尔范畴。
$$
A_n \to B_n
$$
为 $D(\mathcal{A})$ 中映射的逆系统。
% P02-E0100: frozen source omits “Let” before this displayed inverse system.
假设：
\begin{enumerate}
\item 存在整数 $a \leq b$，使得当 $i \not \in [a, b]$ 时，
$H^i(A_n) = 0$ 且 $H^i(B_n) = 0$；
\item 对所有 $i \in \mathbf{Z}$，$\mathcal{A}$ 中映射的逆系统
$H^i(A_n) \to H^i(B_n)$ 定义 $\mathcal{A}$ 中 pro-对象的同构。
% P02-E0101: frozen source has subject--verb disagreement in “the inverse system ... define”.
\end{enumerate}
则各映射 $A_n \to B_n$ 确定 pro-对象 $(A_n)$ 与 pro-对象 $(B_n)$
之间的 pro-同构。
\end{lemma}

\begin{proof}
由公理 TR3，可以归纳地把映射 $A_n \to B_n$ 扩张为特异三角的逆系统
$A_n \to B_n \to C_n \to A_n[1]$。由引理 \ref{derived-lemma-pro-isomorphism}，
只需证明 $C_n$ 是 pro-零的。由引理
\ref{derived-lemma-essentially-constant-cohomology}，只需证明对每个 $p$，
$H^p(C_n)$ 都是 pro-零的。这由假设 (2) 及长正合列
$$
H^p(A_n) \xrightarrow{\alpha_n} H^p(B_n)
\xrightarrow{\beta_n}
H^p(C_n) \xrightarrow{\delta_n} H^{p + 1}(A_n)
\xrightarrow{\epsilon_n}
H^{p + 1}(B_n)
$$
得出。事实上，对每个 $n$，可找到 $m > n$，使得
$\Im(\beta_m)$ 在 $H^p(C_n)$ 中的像为零，因为可选择 $m$，使
$H^p(B_m) \to H^p(B_n)$ 经
$\alpha_n : H^p(A_n) \to H^p(B_n)$ 分解。出于类似理由，
接着可以选择 $k > m$，使 $\Im(\delta_k)$ 在
$H^{p + 1}(A_m)$ 中的像为零。于是 $H^p(C_k) \to H^p(C_n)$ 为零，
因为 $H^p(C_k) \to H^p(C_m)$ 的像包含于 $\Ker(\delta_m)$，而
$H^p(C_m) \to H^p(C_n)$ 消去
$\Ker(\delta_m) = \Im(\beta_m)$。
\end{proof}
% Generated mechanically from the frozen simplicial.tex target.
% Do not edit this generated file; edit the bound locale source instead.

% OK, start here.
%

\setcounter{chapter}{13}
\chapter{单纯方法}



\phantomsection
\label{simplicial-section-phantom}


\section{引言}
\label{simplicial-section-introduction}

\noindent
这是对单纯方法的一份最低限度的介绍。我们只在后文需要某项内容时，才把它加入
这里。关于这些材料，一份一般性的参考文献或许是 \cite{SimpHom}。这些方法能够
完成何种工作的例子包括 Quillen 关于同伦代数的论文，见 \cite{Quillen}，以及
Artin 与 Mazur 关于 'Etale 同伦的论文，见 \cite{ArtinMazur}。

\section{有限全序集的范畴}
\label{simplicial-section-Delta}

\noindent
范畴 $\Delta$ 定义如下：
\begin{enumerate}
\item 其对象为 $[0], [1], [2], \ldots$，其中
$[n] = \{0, 1, 2, \ldots, n\}$；
\item 态射 $[n] \to [m]$ 是相应集合之间的非减映射
$\{0, 1, 2, \ldots, n\} \to \{0, 1, 2, \ldots, m\}$。
\end{enumerate}
这里，映射 $\varphi : [n] \to [m]$ 为{\it 非减的}，按定义是指：若
$i \geq j$，则 $\varphi(i) \geq \varphi(j)$。换言之，$\Delta$ 等价于以非空有限
全序集为对象、以非减映射为态射的“大”范畴。恰有 $n + 1$ 个态射
$[0] \to [n]$，而恰有 $1$ 个态射 $[n] \to [0]$。恰有
$(n + 1)(n + 2)/2$ 个态射 $[1] \to [n]$，而恰有 $n + 2$ 个态射
$[n] \to [1]$。其余情形以此类推。

\begin{definition}
\label{simplicial-definition-face-degeneracy}
对任意整数 $n\geq 1$ 及任意 $0\leq j \leq n$，以
{\it $\delta^n_j : [n-1] \to [n]$}
表示跳过 $j$ 的保序单射。对任意整数 $n\geq 0$ 及任意
$0\leq j \leq n$，以
{\it $\sigma^n_j : [n + 1]  \to [n]$}
表示满足
$(\sigma^n_j)^{-1}(\{j\}) = \{j, j + 1\}$
的保序满射。
\end{definition}

\begin{lemma}
\label{simplicial-lemma-face-degeneracy}
$\Delta$ 中的任一态射都可写成若干态射 $\delta^n_j$ 与 $\sigma^n_j$ 的复合。
\end{lemma}

\begin{proof}
设 $\varphi : [n] \to [m]$ 为 $\Delta$ 的态射。若
$j \not \in \Im(\varphi)$，则对某个态射
$\psi : [n] \to [m - 1]$，可将 $\varphi$ 写成
$\delta^m_j \circ \psi$。若 $\varphi(j) = \varphi(j + 1)$，则对某个态射
$\psi : [n - 1] \to [m]$，可将 $\varphi$ 写成
$\psi \circ \sigma^{n - 1}_j$。上述每次替换都会减小 $n + m$，故在某一步
$\varphi$ 将同时为单射和满射，因而为恒等态射，结论随即成立。
\end{proof}

\begin{lemma}
\label{simplicial-lemma-relations-face-degeneracy}
态射 $\delta^n_j$ 与 $\sigma^n_j$ 满足下列关系。
\begin{enumerate}
\item 若 $0 \leq i < j \leq n + 1$，则
$\delta^{n + 1}_j \circ \delta^n_i =
\delta^{n + 1}_i \circ \delta^n_{j - 1}$。
换言之，下图交换：
$$
\xymatrix{
& [n] \ar[rd]^{\delta^{n + 1}_j} & \\
[n - 1] \ar[ru]^{\delta^n_i} \ar[rd]_{\delta^n_{j - 1}} & &
[n + 1] \\
& [n] \ar[ru]_{\delta^{n + 1}_i} &
}
$$
\item 若 $0 \leq i < j \leq n - 1$，则
$\sigma^{n - 1}_j \circ \delta^n_i =
\delta^{n - 1}_i \circ \sigma^{n - 2}_{j - 1}$。
换言之，下图交换：
$$
\xymatrix{
& [n] \ar[rd]^{\sigma^{n - 1}_j} & \\
[n - 1] \ar[ru]^{\delta^n_i} \ar[rd]_{\sigma^{n - 2}_{j - 1}} & &
[n - 1] \\
& [n - 2] \ar[ru]_{\delta^{n - 1}_i} &
}
$$
\item 若 $0 \leq j \leq n - 1$，则
$\sigma^{n - 1}_j \circ \delta^n_j = \text{id}_{[n - 1]}$
且
$\sigma^{n - 1}_j \circ \delta^n_{j + 1} = \text{id}_{[n - 1]}$。
换言之，下图交换：
$$
\xymatrix{
& [n] \ar[rd]^{\sigma^{n - 1}_j} & \\
[n - 1]
\ar[ru]^{\delta^n_j}
\ar[rd]_{\delta^n_{j + 1}}
\ar[rr]^{\text{id}_{[n - 1]}} & & [n - 1] \\
& [n] \ar[ru]_{\sigma^{n - 1}_j} &
}
$$
\item 若 $0 < j + 1 < i \leq n$，则
$\sigma^{n - 1}_j \circ \delta^n_i =
\delta^{n - 1}_{i - 1} \circ \sigma^{n - 2}_j$。
换言之，下图交换：
$$
\xymatrix{
& [n] \ar[rd]^{\sigma^{n - 1}_j} & \\
[n - 1] \ar[ru]^{\delta^n_i} \ar[rd]_{\sigma^{n - 2}_j} & &
[n - 1] \\
& [n - 2] \ar[ru]_{\delta^{n - 1}_{i - 1}} &
}
$$
\item 若 $0 \leq i \leq j \leq n - 1$，则
$\sigma^{n - 1}_j \circ \sigma^n_i =
\sigma^{n - 1}_i \circ \sigma^n_{j + 1}$。
换言之，下图交换：
$$
\xymatrix{
& [n] \ar[rd]^{\sigma^{n - 1}_j} & \\
[n + 1] \ar[ru]^{\sigma^n_i} \ar[rd]_{\sigma^n_{j + 1}} & &
[n - 1] \\
& [n] \ar[ru]_{\sigma^{n - 1}_i} &
}
$$
\end{enumerate}
\end{lemma}

\begin{proof}
略。
\end{proof}

\begin{lemma}
\label{simplicial-lemma-face-degeneracy-category}
范畴 $\Delta$ 是具有对象 $[n]$（$n \geq 0$）及态射
$\delta^n_j$、$\sigma^n_j$ 的泛范畴，并满足：(a) 每个态射都是这些态射的
复合；(b) 引理 \ref{simplicial-lemma-relations-face-degeneracy} 所列关系成立；
(c) 这些态射之间的任一关系都是上述关系的推论。
\end{lemma}

\begin{proof}
略。
\end{proof}




\section{单纯对象}
\label{simplicial-section-simplicial-object}

\begin{definition}
\label{simplicial-definition-simplicial-object}
设 $\mathcal{C}$ 为范畴。
\begin{enumerate}
\item $\mathcal{C}$ 的一个{\it 单纯对象 $U$} 是从 $\Delta$ 到
$\mathcal{C}$ 的反变函子 $U$，即
$$
U : \Delta^{opp} \longrightarrow \mathcal{C}
$$
\item 若 $\mathcal{C}$ 为集合范畴，则称 $U$ 为{\it 单纯集}。
\item 若 $\mathcal{C}$ 为阿贝尔群范畴，则称 $U$ 为{\it 单纯阿贝尔群}。
\item 单纯对象之间的一个{\it 态射 $U \to U'$} 是函子之间的自然变换。
\item $\mathcal{C}$ 的{\it 单纯对象范畴}记为 $\text{Simp}(\mathcal{C})$。
\end{enumerate}
\end{definition}

\noindent
这意味着存在对象 $U([0]), U([1]), U([2]), \ldots$，并且对任意非减映射
$\varphi$，即 $\varphi : [m] \to [n]$，存在态射
$U(\varphi) : U([n]) \to U([m])$，满足
$U(\varphi \circ \psi) = U(\psi) \circ U(\varphi)$。

\medskip\noindent
特别地，存在唯一态射 $U([0]) \to U([n])$，并且与 $n + 1$ 个映射
$[0] \to [n]$ 相对应，恰有 $n + 1$ 个态射 $U([n]) \to U([0])$。
显然，为了清楚地讨论这些单纯对象，还需要更多记号。为此我们使用上面
第 \ref{simplicial-section-Delta} 节中挑出的那些态射。

\begin{lemma}
\label{simplicial-lemma-characterize-simplicial-object}
设 $\mathcal{C}$ 为范畴。
\begin{enumerate}
\item 给定 $\mathcal{C}$ 中的单纯对象 $U$，得到对象列
$U_n = U([n])$，以及态射
$d^n_j = U(\delta^n_j) : U_n \to U_{n-1}$ 和
$s^n_j = U(\sigma^n_j) : U_n \to U_{n + 1}$。这些态射满足引理
\ref{simplicial-lemma-relations-face-degeneracy} 所列关系的反向关系，即
\begin{enumerate}
\item 若 $0 \leq i < j \leq n + 1$，则
$d^n_i \circ d^{n + 1}_j = d^n_{j - 1} \circ d^{n + 1}_i$。
\item 若 $0 \leq i < j \leq n - 1$，则
$d^n_i \circ s^{n - 1}_j = s^{n - 2}_{j - 1} \circ d^{n - 1}_i$。
\item 若 $0 \leq j \leq n - 1$，则
$\text{id} = d^n_j \circ s^{n - 1}_j = d^n_{j + 1} \circ s^{n - 1}_j$。
\item 若 $0 < j + 1 < i \leq n$，则
$d^n_i \circ s^{n - 1}_j = s^{n - 2}_j \circ d^{n - 1}_{i - 1}$。
\item 若 $0 \leq i \leq j \leq n - 1$，则
$s^n_i \circ s^{n - 1}_j = s^n_{j + 1} \circ s^{n - 1}_i$。
\end{enumerate}
\item 反过来，给定对象列 $U_n$ 及满足 (1)(a)--(e) 的态射
$d^n_j$、$s^n_j$，则存在唯一的 $\mathcal{C}$ 中的单纯对象 $U$，使得
$U_n = U([n])$、$d^n_j = U(\delta^n_j)$ 且
$s^n_j = U(\sigma^n_j)$。
\item 单纯对象 $U$ 与 $U'$ 之间的态射由一族与态射 $d^n_j$ 和
$s^n_j$ 交换的态射 $U_n \to U'_n$ 给出。
\end{enumerate}
\end{lemma}

\begin{proof}
这由引理 \ref{simplicial-lemma-face-degeneracy-category} 得出。
\end{proof}

\begin{remark}
\label{simplicial-remark-relations}
稍微滥用记号，我们有时以 $d_i : U_n \to U_{n - 1}$ 代替
$d^n_i$，并对 $s_i : U_n \to U_{n + 1}$ 使用类似记法。
态射 $d^n_i$ 与 $s^n_i$ 之间的关系可表为：
\begin{enumerate}
\item 若 $i < j$，则 $d_i \circ d_j = d_{j - 1} \circ d_i$。
\item 若 $i < j$，则 $d_i \circ s_j = s_{j - 1} \circ d_i$。
\item 有 $\text{id} = d_j \circ s_j = d_{j + 1} \circ s_j$。
\item 若 $i > j + 1$，则 $d_i \circ s_j = s_j \circ d_{i - 1}$。
\item 若 $i \leq j$，则 $s_i \circ s_j = s_{j + 1} \circ s_i$。
\end{enumerate}
这表示，只要等式左右两侧的复合都有定义，相应等式就成立。
\end{remark}

\noindent
存在唯一态射 $s^0_0 = U(\sigma^0_0) : U_0 \to U_1$，以及两个态射
$d^1_0 = U(\delta^1_0)$ 和 $d^1_1 = U(\delta^1_1)$，它们都是
$U_1 \to U_0$ 的态射。存在两个态射
$s^1_0 = U(\sigma^1_0)$、$s^1_1 = U(\sigma^1_1)$，它们都是
$U_1 \to U_2$ 的态射。还存在三个态射
$d^2_0 = U(\delta^2_0)$、$d^2_1 = U(\delta^2_1)$、$d^2_2 = U(\delta^2_2)$，
它们都是 $U_3 \to U_2$ 的态射。依此类推。
% SOURCE-ERRATUM-CANDIDATE P02-E0102: frozen source says d^2_i: U_3 -> U_2, but its own definition gives d^2_i: U_2 -> U_1.

\medskip\noindent
在图示上，我们把 $U$ 看作：
$$
\xymatrix{
U_2
\ar@<2ex>[r]
\ar@<0ex>[r]
\ar@<-2ex>[r]
&
U_1
\ar@<1ex>[r]
\ar@<-1ex>[r]
\ar@<1ex>[l]
\ar@<-1ex>[l]
&
U_0
\ar@<0ex>[l]
}
$$
这里，指向右侧的箭头是 $d$-态射，指向左侧的箭头是 $s$-态射。

\begin{example}
\label{simplicial-example-constant-simplicial-object}
最简单的例子是取值为 $X \in \Ob(\mathcal{C})$ 的{\it 常值}单纯对象。
换言之，$U_n = X$，且所有映射均为 $\text{id}_X$。
\end{example}

\begin{example}
\label{simplicial-example-fibre-products-simplicial-object}
设 $Y\to X$ 为 $\mathcal{C}$ 的态射，并且所有纤维积
$Y \times_X Y \times_X \ldots \times_X Y$ 都存在。令 $U_n$ 等于
$(n + 1)$ 重纤维积，并使 $\varphi : [n] \to [m]$ 对应于如下映射
（用“坐标”表示）：
$(y_0, \ldots, y_m) \mapsto (y_{\varphi(0)}, \ldots, y_{\varphi(n)})$。
换言之，映射 $U_0 = Y \to U_1 = Y \times_X Y$ 是对角映射，而两个映射
$U_1 = Y \times_X Y \to U_0 = Y$ 是投影映射。
\end{example}

\noindent
从几何上说，上面的例 \ref{simplicial-example-fibre-products-simplicial-object}
十分重要。它表明，最好把映射 $d^n_j : U_n \to U_{n - 1}$ 看成投影映射
（忘掉第 $j$ 个分量），并把映射 $s^n_j : U_n \to U_{n + 1}$ 看成对角映射
（重复第 $j$ 个坐标）。后面各节还会回到这一点。

\begin{lemma}
\label{simplicial-lemma-si-injective}
设 $\mathcal{C}$ 为范畴，$U$ 为 $\mathcal{C}$ 的单纯对象。每个态射
$s^n_i : U_n \to U_{n + 1}$ 都有左逆。特别地，$s^n_i$ 是单态射。
\end{lemma}

\begin{proof}
这是因为 $d_i^{n + 1} \circ s^n_i = \text{id}_{U_n}$。
\end{proof}

\section{作为预层的单纯对象}
\label{simplicial-section-simplicial-presheaves}

\noindent
另一个观察是，可以把 $\mathcal{C}$ 的单纯对象看作 $\Delta$ 上取值于
$\mathcal{C}$ 的预层。见 Sites，定义 \ref{sites-definition-presheaf}。
事实上，若 $U$、$U'$ 是 $\mathcal{C}$ 的单纯对象，则
\begin{equation}
\label{simplicial-equation-simplicial-set-presheaf}
\Mor(U, U') = \Mor_{\textit{PSh}(\Delta)}(U, U').
\end{equation}
下面的一些材料可以用位点一章中更一般的构造来替代。不过，采用单纯对象所特有的
论证似乎能给出更清楚的图景。















\section{余单纯对象}
\label{simplicial-section-cosimplicial-object}

\noindent
范畴 $\mathcal{C}$ 的余单纯对象当然可以直接定义为对偶范畴
$\mathcal{C}^{opp}$ 的单纯对象。不过，人的思维通常并不这样运作，
因此我们在这里单独引入余单纯对象，并指出它们的一些简单性质。

\begin{definition}
\label{simplicial-definition-cosimplicial-object}
设 $\mathcal{C}$ 为范畴。
\begin{enumerate}
\item $\mathcal{C}$ 的一个{\it 余单纯对象 $U$} 是从 $\Delta$ 到
$\mathcal{C}$ 的协变函子 $U$，即
$$
U : \Delta \longrightarrow \mathcal{C}
$$
\item 若 $\mathcal{C}$ 是集合范畴，则称 $U$ 为{\it 余单纯集}。
\item 若 $\mathcal{C}$ 是阿贝尔群范畴，则称 $U$ 为{\it 余单纯阿贝尔群}。
\item {\it 余单纯对象之间的态射 $U \to U'$} 是函子之间的变换。
\item $\mathcal{C}$ 的{\it 余单纯对象范畴}记作
$\text{CoSimp}(\mathcal{C})$。
\end{enumerate}
\end{definition}

\noindent
这就是说，我们有对象 $U([0]), U([1]), U([2]), \ldots$；对于任意非降映射
$\varphi$，即 $\varphi : [m] \to [n]$，都有一个态射
$U(\varphi) : U([m]) \to U([n])$，并且满足
$U(\varphi \circ \psi) = U(\varphi) \circ U(\psi)$。

\medskip\noindent
特别地，有唯一的态射 $U([n]) \to U([0])$，并且恰有 $n + 1$ 个态射
$U([0]) \to U([n])$，它们对应于 $n + 1$ 个映射 $[0] \to [n]$。
显然，要恰当地讨论这些单纯对象，还需要更多记号。
% P02-E0103: frozen source says "these simplicial objects" in the cosimplicial section.
为此，我们考察上面第 \ref{simplicial-section-Delta} 节中挑出的那些态射。

\begin{lemma}
\label{simplicial-lemma-characterize-cosimplicial-object}
设 $\mathcal{C}$ 为范畴。
\begin{enumerate}
\item 给定 $\mathcal{C}$ 中的余单纯对象 $U$，我们得到一列对象
$U_n = U([n])$，并带有态射
$\delta^n_j = U(\delta^n_j) : U_{n - 1} \to U_n$ 和
$\sigma^n_j = U(\sigma^n_j) : U_{n + 1} \to U_n$。这些态射满足
引理 \ref{simplicial-lemma-relations-face-degeneracy} 中列出的关系。
\item 反之，给定一列对象 $U_n$ 和满足这些关系的态射
$\delta^n_j$、$\sigma^n_j$，存在唯一的 $\mathcal{C}$ 中的余单纯对象 $U$，
使得 $U_n = U([n])$、$\delta^n_j = U(\delta^n_j)$ 且
$\sigma^n_j = U(\sigma^n_j)$。
\item 余单纯对象 $U$ 与 $U'$ 之间的态射由一族态射
$U_n \to U'_n$ 给出，并且它们与态射 $\delta^n_j$ 和 $\sigma^n_j$ 可交换。
\end{enumerate}
\end{lemma}

\begin{proof}
这由引理 \ref{simplicial-lemma-face-degeneracy-category} 推出。
\end{proof}

\begin{remark}
\label{simplicial-remark-relations-cosimplicial}
我们有时滥用记号，以 $\delta_i : U_{n - 1} \to U_n$ 代替
$\delta^n_i$，并对 $\sigma_i : U_{n + 1} \to U_n$ 使用类似记法。
态射 $\delta^n_i$ 与 $\sigma^n_i$ 之间的关系可以表述如下：
\begin{enumerate}
\item 若 $i < j$，则
$\delta_j \circ \delta_i = \delta_i \circ \delta_{j - 1}$。
\item 若 $i < j$，则
$\sigma_j \circ \delta_i = \delta_i \circ \sigma_{j - 1}$。
\item 有
$\text{id} = \sigma_j \circ \delta_j = \sigma_j \circ \delta_{j + 1}$。
\item 若 $i > j + 1$，则
$\sigma_j \circ \delta_i = \delta_{i - 1} \circ \sigma_j$。
\item 若 $i \leq j$，则
$\sigma_j \circ \sigma_i = \sigma_i \circ \sigma_{j + 1}$。
\end{enumerate}
这里的意思是，只要等式两边的复合都有定义，相应等式就应成立。
\end{remark}

\noindent
我们得到唯一的态射 $\sigma^0_0 = U(\sigma^0_0) : U_1 \to U_0$，以及
两个态射 $\delta^1_0 = U(\delta^1_0)$ 和
$\delta^1_1 = U(\delta^1_1)$，它们都是态射 $U_0 \to U_1$。
有两个态射 $\sigma^1_0 = U(\sigma^1_0)$、
$\sigma^1_1 = U(\sigma^1_1)$，它们都是态射 $U_2 \to U_1$。
还有三个态射 $\delta^2_0 = U(\delta^2_0)$、
$\delta^2_1 = U(\delta^2_1)$、$\delta^2_2 = U(\delta^2_2)$，
它们都是态射 $U_2 \to U_3$。依此类推。

\medskip\noindent
在图示上，我们把 $U$ 看作：
$$
\xymatrix{
U_0
\ar@<1ex>[r]
\ar@<-1ex>[r]
&
U_1
\ar@<0ex>[l]
\ar@<2ex>[r]
\ar@<0ex>[r]
\ar@<-2ex>[r]
&
U_2
\ar@<1ex>[l]
\ar@<-1ex>[l]
}
$$
这里，指向右侧的箭头是 $\delta$-态射，指向左侧的箭头是
$\sigma$-态射。

\begin{example}
\label{simplicial-example-constant-cosimplicial-object}
最简单的例子是取值为 $X \in \Ob(\mathcal{C})$ 的{\it 常值}余单纯对象。
换言之，$U_n = X$，且所有映射均为 $\text{id}_X$。
\end{example}

\begin{example}
\label{simplicial-example-push-outs-simplicial-object}
设 $X\to Y$ 为 $C$ 的态射，并且所有推出
$Y\amalg_X Y \amalg_X \ldots \amalg_X Y$ 都存在。
% P02-E0104: frozen source uses plain C although the ambient category is mathcal C.
令 $U_n$ 等于 $(n + 1)$ 重推出，并使 $\varphi : [n] \to [m]$ 对应于
“坐标”上的映射
$$
(y \text{ 于 }i\text{次分量})
\mapsto
(y \text{ 于 }\varphi(i)\text{次分量})
$$
换言之，映射 $U_1 = Y \amalg_X Y \to U_0 = Y$ 在每个分量上都是恒等映射。
两个映射 $U_0 = Y \to U_1 = Y \amalg_X Y$ 是两个余投影。
\end{example}

\begin{example}
\label{simplicial-example-simplex-cosimplicial-set}
对每个 $n \geq 0$，以 $C[n]$ 表示余单纯集
$$
\Delta \longrightarrow \textit{Sets},\quad
[k] \longmapsto \Mor_{\Delta}([n], [k])
$$
这个例子与例 \ref{simplicial-example-simplex-simplicial-set} 对偶。
\end{example}

\begin{lemma}
\label{simplicial-lemma-di-injective}
设 $\mathcal{C}$ 为范畴，$U$ 为 $\mathcal{C}$ 的余单纯对象。每个态射
$\delta^n_i : U_{n - 1} \to U_n$ 都有左逆。特别地，$\delta^n_i$ 是单态射。
\end{lemma}

\begin{proof}
这是因为当 $j < n$ 时，
$\sigma_i^{n - 1} \circ \delta^n_i = \text{id}_{U_n}$。
% P02-E0105: frozen source has a type/index mismatch: id_{U_n} and j<n.
\end{proof}













\section{单纯对象的积}
\label{simplicial-section-products}

\noindent
当然，我们应当把单纯对象的积定义为单纯对象范畴中的积。
这可能导致一种容易混淆的情形：积虽然存在，却不能按下面的方式描述。
为避免这种情形，我们直接作如下定义。

\begin{definition}
\label{simplicial-definition-product}
设 $\mathcal{C}$ 为范畴，$U$ 和 $V$ 为 $\mathcal{C}$ 的单纯对象。
假设积 $U_n \times V_n$ 在 $\mathcal{C}$ 中存在。
$U$ 与 $V$ 的{\it 积}是如下定义的单纯对象 $U \times V$：
\begin{enumerate}
\item $(U \times V)_n = U_n \times V_n$，
\item $d^n_i = (d^n_i, d^n_i)$，且
\item $s^n_i = (s^n_i, s^n_i)$。
\end{enumerate}
换言之，$U \times V$ 是 $\Delta$ 上预层 $U$ 与 $V$ 的积。
\end{definition}

\begin{lemma}
\label{simplicial-lemma-product}
若 $U$ 和 $V$ 是范畴 $\mathcal{C}$ 中的单纯对象，且 $U \times V$ 存在，
则对 $\mathcal{C}$ 的任意第三个单纯对象 $W$，都有
$$
\Mor(W, U \times V) =
\Mor(W, U) \times
\Mor(W, V)
$$
\end{lemma}

\begin{proof}
略。
\end{proof}


\section{单纯对象的纤维积}
\label{simplicial-section-fibre-products}

\noindent
当然，我们应当把单纯对象的纤维积定义为单纯对象范畴中的纤维积。
这可能导致一种容易混淆的情形：纤维积虽然存在，却不能按下面的方式描述。
为避免这种情形，我们直接作如下定义。

\begin{definition}
\label{simplicial-definition-fibre-product}
设 $\mathcal{C}$ 为范畴，$U, V, W$ 为 $\mathcal{C}$ 的单纯对象，并设
$a : V \to U$、$b : W \to U$ 为态射。假设纤维积
$V_n \times_{U_n} W_n$ 在 $\mathcal{C}$ 中存在。
$V$ 与 $W$ 在 $U$ 上的{\it 纤维积}是如下定义的单纯对象
$V \times_U W$：
\begin{enumerate}
\item $(V \times_U W)_n = V_n \times_{U_n} W_n$，
\item $d^n_i = (d^n_i, d^n_i)$，且
\item $s^n_i = (s^n_i, s^n_i)$。
\end{enumerate}
换言之，$V \times_U W$ 是 $\Delta$ 上预层 $V$ 与 $W$
在预层 $U$ 上的纤维积。
\end{definition}

\begin{lemma}
\label{simplicial-lemma-fibre-product}
若 $U, V, W$ 是范畴 $\mathcal{C}$ 中的单纯对象，
$a : V \to U$、$b : W \to U$ 是态射，并且 $V \times_U W$ 存在，
则对 $\mathcal{C}$ 的任意第四个单纯对象 $T$，都有
$$
\Mor(T, V \times_U W) =
\Mor(T, V) \times_{\Mor(T, U)}
\Mor(T, W)
$$
\end{lemma}

\begin{proof}
略。
\end{proof}

\section{单纯对象的推出}
\label{simplicial-section-push-outs}

\noindent
当然，我们应当把单纯对象的推出定义为单纯对象范畴中的推出。
这可能导致一种容易混淆的情形：推出虽然存在，却不能按下面的方式描述。
为避免这种情形，我们直接作如下定义。

\begin{definition}
\label{simplicial-definition-push-out}
设 $\mathcal{C}$ 为范畴，$U, V, W$ 为 $\mathcal{C}$ 的单纯对象，并设
$a : U \to V$、$b : U \to W$ 为态射。假设推出
$V_n \amalg_{U_n} W_n$ 在 $\mathcal{C}$ 中存在。
$V$ 与 $W$ 在 $U$ 上的{\it 推出}是如下定义的单纯对象
$V\amalg_U W$：
\begin{enumerate}
\item $(V \amalg_U W)_n = V_n \amalg_{U_n} W_n$，
\item $d^n_i = (d^n_i, d^n_i)$，且
\item $s^n_i = (s^n_i, s^n_i)$。
\end{enumerate}
换言之，$V\amalg_U W$ 是 $\Delta$ 上预层 $V$ 与 $W$
在预层 $U$ 上的推出。
\end{definition}

\begin{lemma}
\label{simplicial-lemma-push-out}
若 $U, V, W$ 是范畴 $\mathcal{C}$ 中的单纯对象，
$a : U \to V$、$b : U \to W$ 是态射，并且 $V\amalg_U W$ 存在，
则对 $\mathcal{C}$ 的任意第四个单纯对象 $T$，都有
$$
\Mor(V\amalg_U W, T) =
\Mor(V, T) \times_{\Mor(U, T)}
\Mor(W, T)
$$
\end{lemma}

\begin{proof}
略。
\end{proof}











\section{余单纯对象的积}
\label{simplicial-section-products-cosimplicial}

\noindent
当然，我们应当把余单纯对象的积定义为余单纯对象范畴中的积。
这可能导致一种容易混淆的情形：积虽然存在，却不能按下面的方式描述。
为避免这种情形，我们直接作如下定义。

\begin{definition}
\label{simplicial-definition-product-cosimplicial-objects}
设 $\mathcal{C}$ 为范畴，$U$ 和 $V$ 为 $\mathcal{C}$ 的余单纯对象。
假设积 $U_n \times V_n$ 在 $\mathcal{C}$ 中存在。
$U$ 与 $V$ 的{\it 积}是如下定义的余单纯对象 $U \times V$：
\begin{enumerate}
\item $(U \times V)_n = U_n \times V_n$，
\item 对任意 $\varphi : [n] \to [m]$，映射
$(U \times V)(\varphi) : U_n \times V_n \to U_m \times V_m$
是积 $U(\varphi) \times V(\varphi)$。
\end{enumerate}
\end{definition}

\begin{lemma}
\label{simplicial-lemma-product-cosimplicial-objects}
若 $U$ 和 $V$ 是范畴 $\mathcal{C}$ 中的余单纯对象，且 $U \times V$ 存在，
则对 $\mathcal{C}$ 的任意第三个余单纯对象 $W$，都有
$$
\Mor(W, U \times V) =
\Mor(W, U) \times
\Mor(W, V)
$$
\end{lemma}

\begin{proof}
略。
\end{proof}

\section{余单纯对象的纤维积}
\label{simplicial-section-fibre-products-cosimplicial}

\noindent
当然，我们应当把余单纯对象的纤维积定义为余单纯对象范畴中的纤维积。
这可能导致一种容易混淆的情形：积虽然存在，却不能按下面的方式描述。
为避免这种情形，我们直接作如下定义。

\begin{definition}
\label{simplicial-definition-fibre-product-cosimplicial-objects}
设 $\mathcal{C}$ 为范畴，$U, V, W$ 为 $\mathcal{C}$ 的余单纯对象，并设
$a : V \to U$ 和 $b : W \to U$ 为态射。假设纤维积
$V_n \times_{U_n} W_n$ 在 $\mathcal{C}$ 中存在。
$V$ 与 $W$ 在 $U$ 上的{\it 纤维积}是如下定义的余单纯对象
$V \times_U W$：
\begin{enumerate}
\item $(V \times_U W)_n = V_n \times_{U_n} W_n$，
\item 对任意 $\varphi : [n] \to [m]$，映射
$(V \times_U W)(\varphi) : V_n \times_{U_n} W_n \to V_m \times_{U_m} W_m$
是积 $V(\varphi) \times_{U(\varphi)} W(\varphi)$。
\end{enumerate}
\end{definition}

\begin{lemma}
\label{simplicial-lemma-fibre-product-cosimplicial-objects}
若 $U, V, W$ 是范畴 $\mathcal{C}$ 中的余单纯对象，
$a : V \to U$、$b : W \to U$ 是态射，并且 $V \times_U W$ 存在，
则对 $\mathcal{C}$ 的任意第四个余单纯对象 $T$，都有
$$
\Mor(T, V \times_U W) =
\Mor(T, V) \times_{\Mor(T, U)}
\Mor(T, W)
$$
\end{lemma}

\begin{proof}
略。
\end{proof}













\section{单纯集}
\label{simplicial-section-simplicial-set}

\noindent
设 $U$ 为单纯集。可以把 $U_0$ 看作{\it $0$-单形}的集合，
把 $U_1$ 看作{\it $1$-单形}的集合，把 $U_2$ 看作
{\it $2$-单形}的集合，依此类推。

\medskip\noindent
我们把映射 $s^n_j : U_n \to U_{n + 1}$ 看作这样的映射：它把
$n$-单形 $A$ 送到退化的 $(n + 1)$-单形 $B$，其中该单形的
$(j, j + 1)$-边被收缩到 $A$ 的第 $j$ 个顶点。我们把映射
$d^n_j : U_n \to U_{n - 1}$ 看作把 $n$-单形 $A$ 送到它的一个面，
即略去第 $j$ 个顶点所得的面。这样一来，就可以从几何上理解
映射 $s^n_j$ 与 $d^n_j$ 之间的关系。

\begin{definition}
\label{simplicial-definition-terminology-simplicial-sets}
设 $U$ 为单纯集。称 $x$ 是 $U$ 的一个{\it $n$-单形}，是指
$x$ 是 $U_n$ 的元素。称 $y$ 是 $x$ 的第 $j$ 个{\it 面}，是指
$d^n_jx = y$。若 $z = s^n_jx$，则称 $z$ 是 $x$ 的第 $j$ 个
{\it 退化}。若一个单形是另一单形的退化，则称它为{\it 退化的}。
\end{definition}

\noindent
下面给出几个基本例子。

\begin{example}
\label{simplicial-example-simplex-simplicial-set}
对每个 $n \geq 0$，以 $\Delta[n]$ 表示单纯集
$$
\Delta^{opp} \longrightarrow \textit{Sets},\quad
[k] \longmapsto \Mor_{\Delta}([k], [n])
$$
下列断言留给读者验证。若 $m > n$，则 $\Delta[n]$ 的每个
$m$-单形都是退化的。$\Delta[n]$ 有唯一的非退化 $n$-单形，
即 $\text{id}_{[n]}$。
\end{example}

\begin{lemma}
\label{simplicial-lemma-simplex-map}
设 $U$ 为单纯集，$n \geq 0$ 为整数。存在典范双射
$$
\Mor(\Delta[n], U)
\longrightarrow
U_n
$$
它把态射 $\varphi$ 送到 $\varphi$ 在 $\Delta[n]$ 的唯一非退化
$n$-单形上的取值。
\end{lemma}

\begin{proof}
略。
\end{proof}

\begin{example}
\label{simplicial-example-simplex-category}
考虑 $\Delta$ 中 $[n]$ 上的对象所成的范畴 $\Delta/[n]$，见
Categories，例 \ref{categories-example-category-over-X}。
有一个函子 $p : \Delta/[n] \to \Delta$。$p$ 在 $[k]$ 上的纤维范畴
（见 Categories，第 \ref{categories-section-fibred-groupoids} 节）
以 $\Delta[n]$ 中 $k$-单形的集合 $\Delta[n]_k$ 为对象，并且只有恒等态射。
对 $\Delta$ 的每个态射 $\varphi : [k] \to [l]$ 以及 $[l]$ 上纤维范畴中的
每个对象 $\psi : [l] \to [n]$，在 $[k]$ 上恰有一个对象带有覆盖
$\varphi$ 的态射，即 $\psi \circ \varphi : [k] \to [n]$。
因此，$\Delta/[n]$ 是 $\Delta$ 上以集合为纤维的范畴。换言之，
可把 $\Delta/[n]$ 看成 $\Delta$ 上的集合预层。另见 Categories，
例 \ref{categories-example-fibred-category-from-functor-of-points}。
这个集合预层与单纯集 $\Delta[n]$ 一致。特别地，由公式
(\ref{simplicial-equation-simplicial-set-presheaf}) 和上面的引理
\ref{simplicial-lemma-simplex-map}，对任意单纯集 $U$，得到公式
$$
\Mor_{\textit{PSh}(\Delta)}(\Delta/[n], U) = U_n
$$
\end{example}

\begin{lemma}
\label{simplicial-lemma-product-degenerate}
设 $U$、$V$ 为单纯集，$a, b \geq 0$ 为整数。假设当 $n > a$ 时，
$U$ 的每个 $n$-单形都是退化的；并假设当 $n > b$ 时，$V$ 的每个
$n$-单形都是退化的。则当 $n > a + b$ 时，$U \times V$ 的每个
$n$-单形都是退化的。
\end{lemma}

\begin{proof}
设 $n > a + b$，并取
$(u, v) \in (U \times V)_n = U_n \times V_n$。由假设，存在
$\alpha : [n] \to [a]$、$u' \in U_a$、$\beta : [n] \to [b]$
和 $v' \in V_b$，使得 $u = U(\alpha)(u')$ 且
$v = V(\beta)(v')$。由于 $n > a + b$，存在
$0 \leq i \leq a + b$，使得 $\alpha(i) = \alpha(i + 1)$ 且
$\beta(i) = \beta(i + 1)$。立即可知，$(u, v)$ 属于
$s^{n - 1}_i$ 的像。
\end{proof}



\section{截断单纯对象与骨架函子}
\label{simplicial-section-skeleton}

\noindent
以 $\Delta_{\leq n}$ 表示 $\Delta$ 的全子范畴，其对象为
$[0], [1], [2], \ldots, [n]$。设 $\mathcal{C}$ 为范畴。

\begin{definition}
\label{simplicial-definition-truncated-simplicial-object}
$\mathcal{C}$ 的一个{\it $n$-截断单纯对象}是从
$\Delta_{\leq n}$ 到 $\mathcal{C}$ 的反变函子。
{\it $n$-截断单纯对象之间的态射}是函子之间的变换。
$\mathcal{C}$ 的 $n$-截断单纯对象范畴记作
$\text{Simp}_n(\mathcal{C})$。
\end{definition}

\noindent
给定 $\mathcal{C}$ 的单纯对象 $U$，截断 $\text{sk}_n U$ 是
$U$ 在子范畴 $\Delta_{\leq n}$ 上的限制。由此定义一个{\it 骨架函子}
$$
\text{sk}_n :
\text{Simp}(\mathcal{C}) \longrightarrow \text{Simp}_n(\mathcal{C})
$$
它从 $\mathcal{C}$ 的单纯对象范畴到 $\mathcal{C}$ 的
$n$-截断单纯对象范畴。为避免与文献中的其他函子混淆，见备注
\ref{simplicial-remark-sk-literature}。




\section{与单纯集的积}
\label{simplicial-section-product-with-simplicial-sets}

\noindent
设 $\mathcal{C}$ 为范畴，$U$ 为单纯集，$V$ 为 $\mathcal{C}$ 的单纯对象。
可以考虑如下协变函子：它把 $\mathcal{C}$ 的单纯对象 $W$ 送到集合
\begin{equation}
\label{simplicial-equation-functor-product-with-simplicial-set}
\left\{
(f_{n, u} : V_n \to W_n)_{n \geq 0, u \in U_n}
\text{ 使得 }
\begin{matrix}
\forall \varphi : [m] \to [n] \\
f_{m, U(\varphi)(u)} \circ V(\varphi) = W(\varphi) \circ f_{n, u}
\end{matrix}
\right\}
\end{equation}
若此函子具有形式 $\Mor_{\text{Simp}(\mathcal{C})}(Q, -)$，则可以把
$Q$ 看作 $U$ 与 $V$ 的积。我们不作这种形式化，而是直接作如下定义。

\begin{definition}
\label{simplicial-definition-product-with-simplicial-set}
设 $\mathcal{C}$ 为范畴，且 $\mathcal{C}$ 的任意两个对象的余积都存在。
设 $U$ 为单纯集，
$V$ 为 $\mathcal{C}$ 的单纯对象，并假设每个 $U_n$ 都是有限非空的。
在这种情形下，定义 $U$ 与 $V$ 的{\it 积 $U \times V$}为
$\mathcal{C}$ 的如下单纯对象，其第 $n$ 项为
$$
(U \times V)_n = \coprod\nolimits_{u\in U_n} V_n
$$
对 $\varphi : [m] \to [n]$，相应映射为
$$
\coprod\nolimits_{u\in U_n} V_n
\longrightarrow
\coprod\nolimits_{u'\in U_m} V_m
$$
它通过态射 $V(\varphi)$，把对应于 $u$ 的分量 $V_n$ 映到对应于
$u' = U(\varphi)(u)$ 的分量 $V_m$。
更宽泛地，若上面列出的所有余积都存在（不对 $\mathcal{C}$ 作任何假设），
则称{\it 积 $U \times V$ 存在}。
\end{definition}

\begin{lemma}
\label{simplicial-lemma-check-product-with-simplicial-set}
设 $\mathcal{C}$ 为范畴，且 $\mathcal{C}$ 的任意两个对象的余积都存在。
设 $U$ 为单纯集，
$V$ 为 $\mathcal{C}$ 的单纯对象，并假设每个 $U_n$ 都是有限非空的。
函子 $W \mapsto \Mor_{\text{Simp}(\mathcal{C})}(U \times V, W)$
与把 $W$ 送到公式
(\ref{simplicial-equation-functor-product-with-simplicial-set}) 中集合的函子典范同构。
\end{lemma}

\begin{proof}
略。
\end{proof}

\begin{lemma}
\label{simplicial-lemma-back-to-U}
设 $\mathcal{C}$ 为范畴，且 $\mathcal{C}$ 的任意两个对象的余积都存在。
暂以
$\textit{FSSets}$ 表示各分量均有限非空的单纯集所成的范畴。
\begin{enumerate}
\item 规则 $(U, V) \mapsto U \times V$ 定义了函子
$\textit{FSSets} \times \text{Simp}(\mathcal{C})
\to \text{Simp}(\mathcal{C})$。
\item 对上面的每一对 $U$、$V$，都有单纯对象之间的典范映射
$$
U \times V \longrightarrow V
$$
它在 $(U \times V)_n = \coprod_u V_n$ 的每个分量上取恒等映射。
\end{enumerate}
\end{lemma}

\begin{proof}
略。
\end{proof}

\noindent
我们简要考察上述构造的一种特殊情形。设 $\mathcal{C}$ 为范畴，
$X$ 为 $\mathcal{C}$ 的对象，$k \geq 0$ 为整数。若所有余积
$X \amalg \ldots \amalg X$ 都存在，则按上述定义，积
$$
X \times \Delta[k]
$$
存在；这里把 $X$ 看作相应的常值单纯对象。

\begin{lemma}
\label{simplicial-lemma-morphism-from-coproduct}
沿用上面的 $X$ 和 $k$。对 $\mathcal{C}$ 的任意单纯对象 $V$，
有如下典范双射
$$
\Mor_{\text{Simp}(\mathcal{C})}(X \times \Delta[k], V)
\longrightarrow
\Mor_\mathcal{C}(X, V_k).
$$
它把 $\gamma$ 送到态射 $\gamma_k$ 在对应于
$\text{id}_{[k]}$ 的分量上的限制。类似地，对任意 $n \geq k$，
若 $W$ 是 $\mathcal{C}$ 的 $n$-截断单纯对象，则有
$$
\Mor_{\text{Simp}_n(\mathcal{C})}(\text{sk}_n(X \times \Delta[k]), W)
=
\Mor_\mathcal{C}(X, W_k).
$$
\end{lemma}

\begin{proof}
态射 $\gamma : X \times \Delta[k] \to V$ 由一族态射
$\gamma_\alpha : X \to V_n$ 给出，其中 $\alpha : [n] \to [k]$。
这些态射必须满足：对所有 $\varphi : [m] \to [n]$，图表
$$
\xymatrix{
X \ar[r]^{\gamma_\alpha} \ar[d]^{\text{id}_X} & V_n \ar[d]^{V(\varphi)} \\
X \ar[r]^{\gamma_{\alpha \circ \varphi}} & V_m
}
$$
可交换。取 $\alpha = \text{id}_{[k]}$，可见对任意
$\varphi : [m] \to [k]$，都有
$\gamma_\varphi = V(\varphi) \circ \gamma_{\text{id}_{[k]}}$。
因此，态射 $\gamma$ 由 $\gamma$ 在对应于 $\text{id}_{[k]}$ 的分量上的取值决定。
反之，给定这样的态射 $f : X \to V_k$，令
$\gamma_\alpha = V(\alpha) \circ f$，即可构造态射 $\gamma$。

\medskip\noindent
截断情形类似，留给读者验证。
\end{proof}

\noindent
一个特殊例子是 $k = 0$ 的情形。此时，引理中的公式只是在说
$$
\Mor_\mathcal{C}(X, V_0)
=
\Mor_{\text{Simp}(\mathcal{C})}(X, V)
$$
其中右边的 $X$ 表示取值为 $X$ 的常值单纯对象。下文将直接使用这个公式，
不再另行说明。




\section{从单纯集到余单纯对象的 Hom}
\label{simplicial-section-hom-from-simplicial-sets-into-cosimplicial}

\noindent
设 $\mathcal{C}$ 为范畴，$U$ 为 $\mathcal{C}$ 的单纯对象，
$V$ 为 $\mathcal{C}$ 的余单纯对象。于是得到如下余单纯集
$\Hom_\mathcal{C}(U, V)$：
\begin{enumerate}
\item 令 $\Hom_\mathcal{C}(U, V)_n = \Mor_\mathcal{C}(U_n, V_n)$；
\item 对 $\varphi : [m] \to [n]$，取映射
$\Hom_\mathcal{C}(U, V)_m \to \Hom_\mathcal{C}(U, V)_n$，
它由 $f \mapsto V(\varphi) \circ f \circ U(\varphi)$ 给出。
\end{enumerate}
这正是下述定义的动机。

\begin{definition}
\label{simplicial-definition-hom-deltak-cosimplicial}
设 $\mathcal{C}$ 为有有限积的范畴，$V$ 为 $\mathcal{C}$ 的余单纯对象，
$U$ 为单纯集，并且每个 $U_n$ 都是有限非空的。定义
{\it $\Hom(U, V)$} 为 $\mathcal{C}$ 的如下余单纯对象：
\begin{enumerate}
\item 令 $\Hom(U, V)_n = \prod_{u \in U_n} V_n$；换言之，
它是 $\mathcal{C}$ 中唯一满足下列 $X$-值点等式的对象：
$$
\Mor_\mathcal{C}(X, \Hom(U, V)_n)
=
\text{Map}(U_n, \Mor_\mathcal{C}(X, V_n))
$$
并且
\item 对 $\varphi : [m] \to [n]$，取映射
$\Hom(U, V)_m \to \Hom(U, V)_n$；在上述 $X$-值点上，
它由 $f \mapsto V(\varphi) \circ f \circ U(\varphi)$ 给出。
\end{enumerate}
\end{definition}

\noindent
把这个定义用积之间的映射完整写出，留给读者。还要指出，该构造关于 $U$
反变、关于 $V$ 协变，正如引理 \ref{simplicial-lemma-back-to-U} 中单纯集与单纯对象之积
的情形一样。




\section{从余单纯集到单纯对象的 Hom}
\label{simplicial-section-hom-from-cosimplicial-sets-into-simplicial}

\noindent
设 $\mathcal{C}$ 为范畴，$U$ 为 $\mathcal{C}$ 的余单纯对象，
$V$ 为 $\mathcal{C}$ 的单纯对象。于是得到如下单纯集
$\Hom_\mathcal{C}(U, V)$：
\begin{enumerate}
\item 令 $\Hom_\mathcal{C}(U, V)_n = \Mor_\mathcal{C}(U_n, V_n)$；
\item 对 $\varphi : [m] \to [n]$，取映射
$\Hom_\mathcal{C}(U, V)_n \to \Hom_\mathcal{C}(U, V)_m$，
它由 $f \mapsto V(\varphi) \circ f \circ U(\varphi)$ 给出。
\end{enumerate}
这正是下述定义的动机。

\begin{definition}
\label{simplicial-definition-hom-deltak-simplicial}
设 $\mathcal{C}$ 为有有限积的范畴，$V$ 为 $\mathcal{C}$ 的单纯对象，
$U$ 为余单纯集，并且每个 $U_n$ 都是有限非空的。定义
{\it $\Hom(U, V)$} 为 $\mathcal{C}$ 的如下单纯对象：
\begin{enumerate}
\item 令 $\Hom(U, V)_n = \prod_{u \in U_n} V_n$；换言之，
它是 $\mathcal{C}$ 中唯一满足下列 $X$-值点等式的对象：
$$
\Mor_\mathcal{C}(X, \Hom(U, V)_n)
=
\text{Map}(U_n, \Mor_\mathcal{C}(X, V_n))
$$
并且
\item 对 $\varphi : [m] \to [n]$，取映射
$\Hom(U, V)_n \to \Hom(U, V)_m$；在上述 $X$-值点上，
它由 $f \mapsto V(\varphi) \circ f \circ U(\varphi)$ 给出。
\end{enumerate}
\end{definition}

\noindent
把这个定义用积之间的映射完整写出，留给读者。还要指出，该构造关于 $U$
反变、关于 $V$ 协变，正如引理 \ref{simplicial-lemma-back-to-U} 中单纯集与单纯对象之积
的情形一样。

\medskip\noindent
我们在一种特殊情形下把上述构造明确写出。设 $X$ 为范畴
$\mathcal{C}$ 的对象，并假设自积 $X \times \ldots \times X$ 存在。
设 $k$ 为整数。考虑各项为
$$
U_n = \prod\nolimits_{\alpha \in \Mor([k], [n])} X
$$
的单纯对象 $U$；对 $\varphi : [m] \to [n]$，映射为
$$
U(\varphi) :
\prod\nolimits_{\alpha \in \Mor([k], [n])} X
\longrightarrow
\prod\nolimits_{\alpha' \in \Mor([k], [m])} X, \quad
(f_{\alpha})_{\alpha} \longmapsto
(f_{\varphi \circ \alpha'})_{\alpha'}
$$
用“坐标”表示，元素 $(x_\alpha)_\alpha$ 被映到元素
$(x_{\varphi \circ \alpha'})_{\alpha'}$。我们断言，这个对象等于
$\Hom(C[k], X)$；这里把 $X$ 看作取值为 $X$ 的常值单纯对象，而 $C[k]$ 是例
\ref{simplicial-example-simplex-cosimplicial-set} 中的余单纯集。

\begin{lemma}
\label{simplicial-lemma-morphism-into-product}
沿用上面的 $X$、$k$ 和 $U$。
\begin{enumerate}
\item 对 $\mathcal{C}$ 的任意单纯对象 $V$，有如下典范双射
$$
\Mor_{\text{Simp}(\mathcal{C})}(V, U)
\longrightarrow
\Mor_\mathcal{C}(V_k, X).
$$
它把 $\gamma$ 送到态射 $\gamma_k$ 与投影到对应于
$\text{id}_{[k]}$ 的因子的复合。
\item 类似地，若 $W$ 是 $\mathcal{C}$ 的 $k$-截断单纯对象，则有
$$
\Mor_{\text{Simp}_k(\mathcal{C})}(W, \text{sk}_k U)
=
\Mor_\mathcal{C}(W_k, X).
$$
\item 上面构造的对象 $U$ 是 $\Hom(C[k], X)$ 的一种实现；其中 $C[k]$
是例 \ref{simplicial-example-simplex-cosimplicial-set} 中的余单纯集。
\end{enumerate}
\end{lemma}

\begin{proof}
先证明 (1)。设 $\gamma : V \to U$ 为态射。它由一族态射
$\gamma_{\alpha} : V_n \to X$ 给出，其中 $\alpha : [k] \to [n]$。
这些态射必须满足：对所有 $\varphi : [m] \to [n]$，图表
$$
\xymatrix{
X \ar[d]^{\text{id}_X} &
V_n \ar[d]^{V(\varphi)}
\ar[l]^{\gamma_{\varphi \circ \alpha'}} \\
X &
V_m \ar[l]_{\gamma_{\alpha'}}
}
$$
对所有 $\alpha' : [k] \to [m]$ 都可交换。取
$\alpha' = \text{id}_{[k]}$，可见对任意 $\varphi : [k] \to [n]$，
都有 $\gamma_\varphi = \gamma_{\text{id}_{[k]}} \circ V(\varphi)$。
因此，态射 $\gamma$ 由 $\gamma_k$ 对应于 $\text{id}_{[k]}$ 的分量决定。
反之，给定这样的态射 $f : V_k \to X$，令
$\gamma_\alpha = f \circ V(\alpha)$，即可构造态射 $\gamma$。

\medskip\noindent
截断情形类似，留给读者验证。

\medskip\noindent
(3) 由 $U$ 的构造以及等式
$C[k]_n = \Mor([k], [n])$ 立即推出；后者正是构造 $U_n$ 时使用的指标集。
\end{proof}




\section{内部 Hom}
\label{simplicial-section-internal-hom}

\noindent
设 $\mathcal{C}$ 为有有限非空积的范畴，$U$、$V$ 为单纯对象
$\mathcal{C}$。
% P02-E0106: frozen source omits “of” before the ambient category.
在某些情形下，函子
$$
\text{Simp}(\mathcal{C})^{opp} \longrightarrow \textit{Sets}, \quad
W \longmapsto \Mor_{\text{Simp}(\mathcal{C})}(W \times V, U)
$$
是可表示的。在这种情形下，把所得的 $\mathcal{C}$ 的单纯对象记作
$\SheafHom(V, U)$，并称从 $V$ 到 $U$ 的{\it 内部 Hom 存在}。
此外，在这种情形下，给定 $\mathcal{C}$ 中的 $X$，有
\begin{align*}
\Mor_\mathcal{C}(X, \SheafHom(V, U)_n)
& =
\Mor_{\text{Simp}(\mathcal{C})}(X \times \Delta[n], \SheafHom(V, U)) \\
& =
\Mor_{\text{Simp}(\mathcal{C})}(X \times \Delta[n]\times V, U) \\
& =
\Mor_{\text{Simp}(\mathcal{C})}(X, \SheafHom(\Delta[n] \times V, U)) \\
& =
\Mor_\mathcal{C}(X, \SheafHom(\Delta[n] \times V, U)_0)
\end{align*}
只要 $\SheafHom(\Delta[n] \times V, U)$ 也存在即可。第一个和最后一个等式
由引理 \ref{simplicial-lemma-morphism-from-coproduct} 推出。

\medskip\noindent
由此得到的要点是：给定 $U$ 和 $V$，若要构造内部 Hom，就应尝试构造对象
$$
\SheafHom(\Delta[n] \times V, U)_0
$$
因为它们应当是 $\SheafHom(V, U)$ 的第 $n$ 项。下一节将研究单纯对象
“$\Hom(\Delta[n], U)$”的一种构造。


\section{从单纯集到单纯对象的 Hom}
\label{simplicial-section-hom-from-simplicial-sets}

\noindent
受上述内部 Hom 讨论的启发，我们定义一个单纯对象，它应当分类从单纯集到
范畴 $\mathcal{C}$ 的某个给定单纯对象的态射。

\begin{definition}
\label{simplicial-definition-hom-from-simplicial-set}
设 $\mathcal{C}$ 为任意两个对象的余积都存在的范畴。设 $U$ 为单纯集，
并且对所有 $n \geq 0$，$U_n$ 都有限非空。设 $V$ 为
$\mathcal{C}$ 的单纯对象。以 {\it $\Hom(U, V)$} 表示任意满足下式的
$\mathcal{C}$ 的单纯对象：
$$
\Mor_{\text{Simp}(\mathcal{C})}(W, \Hom(U, V))
=
\Mor_{\text{Simp}(\mathcal{C})}(W \times U, V)
$$
其中等式关于 $\mathcal{C}$ 的单纯对象 $W$ 函子性地成立。
\end{definition}

\noindent
当然，$\Hom(U, V)$ 未必存在。此外，由第 \ref{simplicial-section-internal-hom} 节
的讨论可预期，若它存在，则
$\Hom(U, V)_n = \Hom(U \times \Delta[n], V)_0$。
这里不对这些 Hom 对象使用斜体记号，因为 $\Hom(U, V)$ 不是内部 Hom。

\begin{lemma}
\label{simplicial-lemma-exists-hom-0-from-simplicial-set}
假设范畴 $\mathcal{C}$ 有任意两个对象的余积和可数极限。
设 $U$ 为单纯集，并且对所有 $n \geq 0$，$U_n$ 都有限非空。
设 $V$ 为 $\mathcal{C}$ 的单纯对象。则函子
\begin{eqnarray*}
\mathcal{C}^{opp} & \longrightarrow & \textit{Sets} \\
X
& \longmapsto &
\Mor_{\text{Simp}(\mathcal{C})}(X \times U, V)
\end{eqnarray*}
是可表示的。
\end{lemma}

\begin{proof}
从 $X \times U$ 到 $V$ 的态射由一族态射
$f_u : X \to V_n$ 给出，其中 $n \geq 0$ 且 $u \in U_n$。
而这样的一族态射恰好在下述条件下定义一个态射：对所有
$\varphi : [m] \to [n]$，图表
$$
\xymatrix{
X \ar[r]^{f_u} \ar[d]_{\text{id}_X} & V_n \ar[d]^{V(\varphi)} \\
X \ar[r]^{f_{U(\varphi)(u)}} & V_m
}
$$
都可交换。因此自然要引入范畴 $\mathcal{U}$ 和函子
$\mathcal{V} : \mathcal{U}^{opp} \to \mathcal{C}$，定义如下：
\begin{enumerate}
\item $\mathcal{U}$ 的对象集是 $\coprod_{n \geq 0} U_n$；
\item 从 $u' \in U_m$ 到 $u \in U_n$ 的态射，是满足
$U(\varphi)(u) = u'$ 的 $\varphi : [m] \to [n]$；
\item 对 $u \in U_n$，令 $\mathcal{V}(u) = V_n$；且
\item 对满足 $U(\varphi)(u) = u'$ 的 $\varphi : [m] \to [n]$，令
$\mathcal{V}(\varphi) = V(\varphi) : V_n \to V_m$。
\end{enumerate}
此时很清楚，我们的函子正是定义
$$
\lim_{\mathcal{U}^{opp}} \mathcal{V}
$$
的函子。因此，若 $\mathcal{C}$ 有可数极限，则此极限存在，从而表示引理中
函子的对象也存在。
\end{proof}

\begin{lemma}
\label{simplicial-lemma-exists-hom-0-from-simplicial-set-finite}
假设范畴 $\mathcal{C}$ 有任意两个对象的余积和有限极限。
设 $U$ 为单纯集，并且对所有 $n \geq 0$，$U_n$ 都有限非空。
假设对所有 $n \gg 0$，$U$ 的每个 $n$-单形都是退化的。
设 $V$ 为 $\mathcal{C}$ 的单纯对象。则函子
\begin{eqnarray*}
\mathcal{C}^{opp} & \longrightarrow & \textit{Sets} \\
X
& \longmapsto &
\Mor_{\text{Simp}(\mathcal{C})}(X \times U, V)
\end{eqnarray*}
是可表示的。
\end{lemma}

\begin{proof}
我们必须证明，引理 \ref{simplicial-lemma-exists-hom-0-from-simplicial-set} 的证明中
所述范畴 $\mathcal{U}$ 有一个有限子范畴 $\mathcal{U}'$，使得
$\mathcal{V}$ 在 $\mathcal{U}'$ 上的极限与 $\mathcal{V}$ 在
$\mathcal{U}$ 上的极限相同。我们将使用 Categories，引理
\ref{categories-lemma-initial}。对 $m > 0$，以
$\mathcal{U}_{\leq m}$ 表示对象为
$\coprod_{0 \leq n \leq m} U_m$ 的全子范畴。
% P02-E0107: frozen source uses U_m under an index n; likely U_n.
取整数 $m_0$，使得当 $n > m_0$ 时，单纯集 $U$ 的每个 $n$-单形都退化。
对任意充分大的 $m \geq m_0$，子范畴 $\mathcal{U}_{\leq m}$ 满足
Categories，定义 \ref{categories-definition-initial} 的性质 (1)。

\medskip\noindent
设 $u \in U_n$、$u' \in U_{n'}$，其中 $n, n' \leq m_0$；并设
$\varphi : [k] \to [n]$、$\varphi' : [k] \to [n']$ 为满足
$U(\varphi)(u) = U(\varphi')(u')$ 的态射。一个简单的组合论论证表明，
若 $k > 2m_0$，则存在指标 $0 \leq i \leq 2m_0$，使得
$\varphi(i) =\varphi(i + 1)$ 且
$\varphi'(i) = \varphi'(i + 1)$。（鸽巢原理会告诉你，当
$k > m_0^2$ 时这也成立；对下面的论证而言已经足够。）因此，若
$k > 2m_0$，则可写成
$\varphi = \psi \circ \sigma^{k - 1}_i$ 和
$\varphi' = \psi' \circ \sigma^{k - 1}_i$，其中
$\psi : [k - 1] \to [n]$、$\psi' : [k - 1] \to [n']$。
由于 $s^{k - 1}_i : U_{k - 1} \to U_k$ 是单射（见引理
\ref{simplicial-lemma-si-injective}），也可得
$U(\psi)(u) = U(\psi')(u')$。继续这一过程可知：给定
$\mathcal{U}$ 的态射 $u \to z$ 和 $u' \to z$，其中
$u, u' \in \mathcal{U}_{\leq m_0}$，存在可交换图表
$$
\xymatrix{
u \ar[rd] \ar[rrd] & & \\
& a \ar[r] & z \\
u' \ar[ru] \ar[rru]
}
$$
其中 $a \in \mathcal{U}_{\leq 2m_0}$。

\medskip\noindent
容易由此推出，有限子范畴 $\mathcal{U}_{\leq 2m_0}$ 符合要求。
事实上，设给定 $x' \in U_n$ 和 $x'' \in U_{n'}$，其中
$n, n' \leq 2m_0$；并给定 $\mathcal{U}$ 中具有同一目标的态射
$x' \to x$ 和 $x'' \to x$。由 $m_0$ 的选取，可以找到
$\mathcal{U}_{\leq m_0}$ 的对象 $u, u'$ 以及态射
$u \to x'$、$u' \to x''$。由上述结论，可以找到
$a \in \mathcal{U}_{\leq 2m_0}$ 以及态射 $u \to a$、$u' \to a$，
使得
$$
\xymatrix{
u \ar[rd] \ar[rrd] \ar[r] & x' \ar[rd] & \\
& a \ar[r] & x \\
u' \ar[ru] \ar[rru] \ar[r] & x'' \ar[ru] &
}
$$
可交换。把这个图顺时针旋转 90 度，就得到 Categories，定义
\ref{categories-definition-initial} 中 (2) 所要求的图表。
\end{proof}

\begin{lemma}
\label{simplicial-lemma-exists-hom-from-simplicial-set-finite}
假设范畴 $\mathcal{C}$ 有任意两个对象的余积和有限极限。
设 $U$ 为单纯集，并且对所有 $n \geq 0$，$U_n$ 都有限非空。
假设对所有 $n \gg 0$，$U$ 的每个 $n$-单形都是退化的。
设 $V$ 为 $\mathcal{C}$ 的单纯对象。则 $\Hom(U, V)$ 存在，而且有预期的等式
$$
\Hom(U, V)_n = \Hom(U \times \Delta[n], V)_0.
$$
\end{lemma}

\begin{proof}
如下构造这个单纯对象。对 $n \geq 0$，以 $\Hom(U, V)_n$ 表示
$\mathcal{C}$ 中表示函子
$$
X
\longmapsto
\Mor_{\text{Simp}(\mathcal{C})}(X \times U \times \Delta[n], V)
$$
的对象。由引理 \ref{simplicial-lemma-exists-hom-0-from-simplicial-set-finite}，
它存在，因为 $U \times \Delta[n]$ 是一个各单形集有限、且在足够高次数
没有非退化单形的单纯集；见引理 \ref{simplicial-lemma-product-degenerate}。
对 $\varphi : [m] \to [n]$，得到诱导的单纯集映射
$\varphi : \Delta[m] \to \Delta[n]$。因而得到关于 $X$ 函子性的态射
$X \times U \times \Delta[m] \to X \times U \times \Delta[n]$，
进而得到函子之间的变换，再进而给出
$$
\Hom(U, V)(\varphi) :
\Hom(U, V)_n
\longrightarrow
\Hom(U, V)_m.
$$
显然，这定义了从 $\Delta$ 到范畴 $\mathcal{C}$ 的反变函子
$\Hom(U, V)$。换言之，我们得到 $\mathcal{C}$ 的一个单纯对象。

\medskip\noindent
还需证明 $\Hom(U, V)$ 满足所需的泛性质
$$
\Mor_{\text{Simp}(\mathcal{C})}(W, \Hom(U, V))
=
\Mor_{\text{Simp}(\mathcal{C})}(W \times U, V)
$$
为此，给定 $f : W \to \Hom(U, V)$。我们要构造右端的元素
$f' : W \times U \to V$。由构造，每个
$f_n : W_n \to \Hom(U, V)_n$ 对应于态射
$f_n : W_n \times U \times \Delta[n] \to V$。此外，
对每个态射 $\varphi : [m] \to [n]$，图表
$$
\xymatrix{
W_n \times U \times \Delta[m]
\ar[rr]_{W(\varphi)\times \text{id} \times \text{id}}
\ar[d]_{\text{id} \times \text{id} \times \varphi} & &
W_m \times U \times \Delta[m] \ar[d]^{f_m} \\
W_n \times U \times \Delta[n] \ar[rr]^{f_n} & & V
}
$$
可交换。对 $(\Delta[n])_k$ 中的 $\psi : [n] \to [k]$，以
$(f_n)_{k, \psi} : W_n \times U_k \to V_k$ 表示 $(f_n)_k$
对应于元素 $\psi$ 的分量。定义
$f'_n : W_n \times U_n \to V_n$ 为
$f'_n = (f_n)_{n, \text{id}}$；换言之，它是
$(f_n)_n : W_n \times U_n \times (\Delta[n])_n \to V_n$
在 $W_n \times U_n \times \text{id}_{[n]}$ 上的限制。
为证明族 $(f'_n)$ 定义了单纯对象之间的态射，必须证明对任意
$\varphi : [m] \to [n]$，都有
$V(\varphi) \circ f'_n =
f'_m \circ W(\varphi) \times U(\varphi)$。
上面的可交换图表说明，
$(f_n)_{m, \varphi} : W_n \times U_m \to V_m$ 等于
$(f_m)_{m, \text{id}} \circ W(\varphi) :
W_n \times U_m \to V_m$。
另一方面，$f_n$ 是单纯对象之间的态射这一事实蕴含图表
$$
\xymatrix{
W_n \times U_n \times (\Delta[n])_n
\ar[r]_-{(f_n)_n}
\ar[d]_{\text{id} \times U(\varphi) \times \varphi}
& V_n \ar[d]^{V(\varphi)} \\
W_n \times U_m \times (\Delta[n])_m \ar[r]^-{(f_n)_m} & V_m
}
$$
可交换。这又蕴含 $(f_n)_{m, \varphi} \circ U(\varphi)$ 等于
$V(\varphi) \circ (f_n)_{n, \text{id}}$。综上，得到
$
V(\varphi) \circ (f_n)_{n, \text{id}}
=
(f_n)_{m, \varphi} \circ U(\varphi)
=
(f_m)_{m, \text{id}} \circ W(\varphi)\circ U(\varphi)
=
(f_m)_{m, \text{id}} \circ W(\varphi)\times U(\varphi)
$
如所欲证。

\medskip\noindent
反过来，给定态射 $f' : W \times U \to V$，我们如下定义态射
$f : W \to \Hom(U, V)$。由引理 \ref{simplicial-lemma-morphism-from-coproduct}，
态射 $\text{id} : W_n \to W_n$ 对应于唯一态射
$c_n : W_n \times \Delta[n] \to W$。因此可以考虑复合
$$
W_n \times \Delta[n] \times U
\xrightarrow{c_n}
W \times U
\xrightarrow{f'}
V.
$$
由构造，这对应于唯一态射 $f_n : W_n \to \Hom(U, V)_n$。
所需的这些态射定义了单纯集之间的态射这一点，留给读者验证。
% P02-E0108: frozen source uses plural “morphisms” with singular id.
% P02-E0109: frozen source says simplicial sets although these are C-valued simplicial objects.

\medskip\noindent
同样留给读者验证，$f \mapsto f'$ 和 $f' \mapsto f$ 是互逆操作。
\end{proof}

\begin{lemma}
\label{simplicial-lemma-hom-from-coprod}
假设范畴 $\mathcal{C}$ 有任意两个对象的余积和有限极限。
设 $a : U \to V$、$b : U \to W$ 为单纯集之间的态射。
假设对所有 $n \geq 0$，$U_n, V_n, W_n$ 都有限非空。
假设对所有 $n \gg 0$，$U, V, W$ 的每个 $n$-单形都是退化的。
设 $T$ 为 $\mathcal{C}$ 的单纯对象。则
$$
\Hom(V, T) \times_{\Hom(U, T)} \Hom(W, T)
=
\Hom(V \amalg_U W, T)
$$
换言之，左边的纤维积由右边的 Hom 对象表示。
\end{lemma}

\begin{proof}
由引理 \ref{simplicial-lemma-exists-hom-from-simplicial-set-finite}，所需的所有
$\Hom$ 对象都存在并满足正确的函子性。现在可以把左边的第 $n$ 项识别为
表示把 $X$ 送到如下函子性等式链中第一个集合的对象：
\begin{align*}
&
\Mor(X \times \Delta[n],
\Hom(V, T) \times_{\Hom(U, T)} \Hom(W, T)) \\
& =
\Mor(X \times \Delta[n], \Hom(V, T))
\times_{\Mor(X \times \Delta[n], \Hom(U, T))}
\Mor(X \times \Delta[n], \Hom(W, T)) \\
& =
\Mor(X \times \Delta[n] \times V, T)
\times_{\Mor(X \times \Delta[n] \times U, T)}
\Mor(X \times \Delta[n] \times W, T) \\
& =
\Mor(X \times \Delta[n] \times (V \amalg_U W), T))
\end{align*}
% P02-E0110: frozen source has an unmatched extra closing parenthesis in the final Hom term.
这里使用了等式
$$
(X \times \Delta[n] \times V)
\times_{X \times \Delta[n] \times U}
(X \times \Delta[n] \times W)
=
X \times \Delta[n] \times (V \amalg_U W)
$$
% P02-E0111: frozen source uses a fibre product on maps induced by U -> V and U -> W; the typed construction is a pushout.
逐项验证此等式很容易。引理的结论随即由上列等式的最后一项对应于
$\Hom(V \amalg_U W, T)_n$ 得出。
\end{proof}

\section{单纯对象的分裂}
\label{simplicial-section-splitting}

\noindent
范畴 $\mathcal{C}$ 的对象 $X$ 的子对象 $N$，是 $\mathcal{C}$ 的一个对象
$N$ 连同一个单态射 $N \to X$。当然，如果存在与到 $X$ 的态射相容的
同构 $N \to N'$，我们就（滥用记号地）称子对象 $N$、$N'$ 相等。
子对象全体构成偏序集。（这是由于我们对范畴的约定；例如，对差一个同伦
意义下的空间范畴并不成立。）

\begin{definition}
\label{simplicial-definition-split}
设 $\mathcal{C}$ 为有有限非空余积的范畴。称 $\mathcal{C}$ 的单纯对象
$U$ 是{\it 分裂的}，如果存在 $U_m$ 的子对象 $N(U_m)$（$m \geq 0$），
使得对所有 $n \geq 0$，映射
\begin{equation}
\label{simplicial-equation-splitting}
\coprod\nolimits_{\varphi : [n] \to [m]\text{ 满射}}
N(U_m)
\longrightarrow
U_n
\end{equation}
都是同构。若 $U$ 是 $\mathcal{C}$ 的 $r$-截断单纯对象，则称 $U$
是{\it 分裂的}，如果存在 $U_m$ 的子对象 $N(U_m)$（$r \geq m \geq 0$），
使得对 $r \geq n \geq 0$，映射 (\ref{simplicial-equation-splitting}) 是同构。
\end{definition}

\noindent
在这种情形下，$N(U_0) = U_0$。其次有
$U_1 = U_0 \amalg N(U_1)$。再者有
$$
U_2 = U_0 \amalg N(U_1) \amalg N(U_1) \amalg N(U_2).
$$
事实表明，在许多范畴 $\mathcal{C}$ 中，每个单纯对象都是分裂的。

\begin{lemma}
\label{simplicial-lemma-splitting-simplicial-sets}
设 $U$ 为单纯集。则 $U$ 有唯一的分裂，其中 $N(U_m)$ 等于非退化
$m$-单形的集合。
\end{lemma}

\begin{proof}
由定义立即可知，若存在分裂，则 $N(U_m)$ 必须是非退化单形的集合。
设 $x \in U_n$。假设存在满射 $\varphi : [n] \to [k]$ 和
$\psi : [n] \to [l]$，以及非退化单形 $y \in U_k$、$z \in U_l$，
使得 $x = U(\varphi)(y)$ 且 $x = U(\psi)(z)$。取 $\psi$ 的右逆
$\xi : [l] \to [n]$，即 $\psi \circ \xi = \text{id}_{[l]}$。
于是 $z = U(\xi)(x)$，从而
$z = U(\xi)(x) = U(\varphi \circ \xi)(y)$。由于 $z$ 非退化，
可知 $\varphi \circ \xi : [l] \to [k]$ 是满射，故 $l \geq k$。
类似地，$k \geq l$。因此，对 $\psi$ 的右逆 $\xi$ 的任意选取，
$\varphi \circ \xi : [l] \to [k]$ 都必须是恒等映射。
这很容易推出 $\psi = \varphi$。
\end{proof}

\noindent
当然，单纯集之间的映射可能把非退化 $n$-单形映到退化 $n$-单形。
因此，引理 \ref{simplicial-lemma-splitting-simplicial-sets} 的分裂不是函子性的。
下面给出一种它具有函子性的情形。

\begin{lemma}
\label{simplicial-lemma-injective-map-simplicial-sets}
设 $f : U \to V$ 为单纯集之间的态射。假设 (a) $U$ 的每个非退化单形
的像都是 $V$ 的非退化单形，并且 (b) $f$ 从 $U$ 的非退化单形集合
到 $V$ 的非退化单形集合的限制是单射。则对所有 $n$，$f_n$ 都是单射。
把“单射”替换成“满射”或“双射”时，结论同样成立。
\end{lemma}

\begin{proof}
在假设 (a) 下，映射 $f$ 保持引理
\ref{simplicial-lemma-splitting-simplicial-sets} 的分裂给出的不交并分解；
换言之，得到可交换图表
$$
\xymatrix{
\coprod\nolimits_{\varphi : [n] \to [m]\text{ 满射}}
N(U_m)
\ar[r] \ar[d] &
U_n \ar[d] \\
\coprod\nolimits_{\varphi : [n] \to [m]\text{ 满射}}
N(V_m)
\ar[r] &
V_n.
}
$$
于是 (b) 清楚地表明，左侧竖直箭头是单射（相应地，满射或双射）。
\end{proof}

\begin{lemma}
\label{simplicial-lemma-simplicial-set-n-skel-sub}
设 $U$ 为单纯集，$n \geq 0$ 为整数。规则
$$
U'_m = \bigcup\nolimits_{\varphi : [m] \to [i], \ i\leq n} \Im(U(\varphi))
$$
定义了子单纯集 $U' \subset U$，并且当 $i \leq n$ 时
$U'_i = U_i$。此外，对所有 $m > n$，$U'$ 的每个 $m$-单形都是退化的。
\end{lemma}

\begin{proof}
若 $x \in U_m$，且对某个 $y \in U_i$、$i \leq n$ 和某个
$\varphi : [m] \to [i]$ 有 $x = U(\varphi)(y)$，则对任意
$\psi : [m'] \to [m]$，像 $U(\psi)(x)$ 等于
$U(\varphi \circ \psi)(y)$，而且
$\varphi \circ \psi : [m'] \to [i]$。因此 $U'$ 是单纯集。
由构造，所有次数为 $n + 1$ 或更高的单形都是退化的。
\end{proof}

\begin{lemma}
\label{simplicial-lemma-splitting-simplicial-groups}
设 $U$ 为单纯阿贝尔群。则 $U$ 有如下分裂：取
$N(U_0) = U_0$，并对 $m \geq 1$ 取
$$
N(U_m) = \bigcap\nolimits_{i = 0}^{m - 1} \Ker(d^m_i).
$$
此外，这个分裂在单纯阿贝尔群范畴上是函子性的。
\end{lemma}

\begin{proof}
对 $n$ 归纳，我们将证明，引理中对 $N(U_m)$ 的选取保证
(\ref{simplicial-equation-splitting}) 对 $m \leq n$ 是同构。$n = 0$ 时显然。
在本证明的其余部分，为提高可读性，我们将省略映射 $d_i$ 和 $s_i$
的上标。还将反复使用备注 \ref{simplicial-remark-relations} 中的关系。

\medskip\noindent
先作一般性观察。对 $0 \leq i \leq m$ 和 $z \in U_m$，有
$d_i(s_i(z)) = z$。因此，通过取
$x' = (x - s_i(d_i(x)))$ 和 $x'' = s_i(d_i(x))$，可以把任意
$x \in U_{m + 1}$ 唯一写成 $x = x' + x''$，其中
$d_i(x') = 0$ 且 $x'' \in \Im(s_i)$。此外，满足
$x'' = s_i(z)$ 的元素 $z \in U_m$ 是唯一的，因为 $s_i$ 是单射。

\medskip\noindent
下面给出分解任意 $x \in U_{n + 1}$ 的步骤。先写成
$x = x_0 + s_0(z_0)$，其中 $d_0(x_0) = 0$。
再写成 $x_0 = x_1 + s_1(z_1)$，其中
$d_1(x_1) = 0$。如此继续，得到
\begin{eqnarray*}
x & = & x_0 + s_0(z_0), \\
x_0 & = & x_1 + s_1(z_1), \\
x_1 & = & x_2 + s_2(z_2), \\
\ldots & \ldots & \ldots \\
x_{n - 1} & = & x_n + s_n(z_n)
\end{eqnarray*}
其中对所有 $i = n, \ldots, 0$，都有 $d_i(x_i) = 0$。
由上面的一般性观察，所有 $x_i$ 和 $z_i$ 都由 $x$ 唯一确定。
我们断言
$x_i \in
\Ker(d_0) \cap
\Ker(d_1) \cap
\ldots \cap
\Ker(d_i)$
且
$z_i \in
\Ker(d_0) \cap
\ldots \cap
\Ker(d_{i - 1})$
对 $i = n, \ldots, 0$ 成立。在这里以及下文，核的空交表示整个空间；
也就是说，当 $i = 0$ 时，记号
$z_0 \in \Ker(d_0) \cap
\ldots \cap
\Ker(d_{i - 1})$
表示 $z_0 \in U_n$，没有任何限制。

\medskip\noindent
我们对 $i$ 作递增归纳来证明此断言。由 $x_0$ 和 $z_0$ 的构造，
$i = 0$ 时显然。假设对 $i - 1$ 结论成立，来证明
$0 < i \leq n$ 的情形。首先，由构造有 $d_i(x_i) = 0$。
取满足 $0 \leq j < i$ 的 $j$。由归纳假设，$d_j(x_{i - 1}) = 0$。
因此
$$
0 = d_j(x_{i - 1})
= d_j(x_i) + d_j(s_i(z_i))
= d_j(x_i) + s_{i - 1}(d_j(z_i)).
$$
最后一个等式由备注 \ref{simplicial-remark-relations} 的关系得到。
这些关系还蕴含
$d_{i - 1}(d_j(x_i)) = d_j(d_i(x_i)) = 0$，
因为由构造 $d_i(x_i)= 0$。于是，上面一般性观察中的唯一性表明，等式
$0 = x' + x'' = d_j(x_i) + s_{i - 1}(d_j(z_i))$
只能在两项都为零时成立。我们得到 $d_j(x_i) = 0$；又由
$s_{i - 1}$ 的单射性，得到 $d_j(z_i) = 0$。断言得证。

\medskip\noindent
该断言蕴含，可唯一写成
$$
x = s_0(z_0) + s_1(z_1) + \ldots + s_n(z_n) + x_n
$$
其中 $x_n \in N(U_{n + 1})$，且
$z_i \in \Ker(d_0) \cap \ldots \cap \Ker(d_{i - 1})$。
这也可以表述为我们找到了直和分解
$$
U_{n + 1}
=
N(U_{n + 1})
\oplus
\bigoplus\nolimits_{i = 0}^{i = n}
s_i\Big(\Ker(d_0) \cap \ldots \cap \Ker(d_{i - 1})\Big)
$$
并且它满足：对 $j = 0, \ldots, n$，有
$$
\Ker(d_0) \cap \ldots \cap \Ker(d_j)
=
N(U_{n + 1}) \oplus
\bigoplus\nolimits_{i = j + 1}^{i = n}
s_i\Big(\Ker(d_0) \cap \ldots \cap \Ker(d_{i - 1})\Big)
$$
结论可由此得到。$x$ 的表达式中的每个 $z_i$ 都可以唯一写成
$$
z_i = s_i(z'_{i, i}) + \ldots + s_{n - 1}(z'_{i, n - 1}) + z_{i, 0}
$$
其中 $z_{i, 0} \in N(U_n)$，且
$z'_{i, j} \in \Ker(d_0) \cap \ldots \cap \Ker(d_{j - 1})$。
因为 $z_i$ 已经属于 $d_0, \ldots, d_i$ 的核，分解 $z_i$ 时最初的几步
为零。这又唯一给出
$$
x = x_n + s_0(z_{0, 0}) + s_1(z_{1, 0}) + \ldots + s_n(z_{n, 0}) +
\sum\nolimits_{0 \leq i \leq j \leq n - 1} s_i(s_j(z'_{i, j})).
$$
继续这一过程，最终可见，$x$ 被分解成如下形式的项之和：
$$
s_{i_1} s_{i_2} \ldots s_{i_k} (z)
$$
其中 $0 \leq i_1 \leq i_2 \leq \ldots \leq i_k \leq n - k + 1$，
且 $z \in N(U_{n + 1 - k})$。这正是所需的分解，因为任意满射
$[n + 1] \to [n + 1 - k]$ 都可以唯一表示为
$$
\sigma^{n - k}_{i_k} \ldots \sigma^{n - 1}_{i_2} \sigma^n_{i_1}
$$
其中 $0 \leq i_1 \leq i_2 \leq \ldots \leq i_k \leq n - k + 1$。
\end{proof}

\begin{lemma}
\label{simplicial-lemma-splitting-abelian-category}
设 $\mathcal{A}$ 为阿贝尔范畴，$U$ 为 $\mathcal{A}$ 中的单纯对象。
则 $U$ 有如下分裂：取 $N(U_0) = U_0$，并对 $m \geq 1$ 取
$$
N(U_m) = \bigcap\nolimits_{i = 0}^{m - 1} \Ker(d^m_i).
$$
此外，这个分裂在 $\mathcal{A}$ 的单纯对象范畴上是函子性的。
\end{lemma}

\begin{proof}
对 $\mathcal{A}$ 的任意对象 $A$，得到单纯阿贝尔群
$\Mor_\mathcal{A}(A, U)$。其中每一个都由引理
\ref{simplicial-lemma-splitting-simplicial-groups} 典范地分裂。此外，
% P02-E0112: frozen source says “Each of these are” instead of “Each ... is”.
$$
N(\Mor_\mathcal{A}(A, U_m)) =
\bigcap\nolimits_{i = 0}^{m - 1} \Ker(d^m_i) =
\Mor_\mathcal{A}(A, N(U_m)).
$$
因此，对 $\mathcal{A}$ 的任意对象应用函子 $\Mor_\mathcal{A}(A, -)$ 后，
态射 (\ref{simplicial-equation-splitting}) 都成为同构。故由 Yoneda 引理，它本身是同构。
\end{proof}

\begin{lemma}
\label{simplicial-lemma-injective-map-simplicial-abelian}
\begin{slogan}
Dold--Kan 正规化函子反映单射性、满射性和同构性。
\end{slogan}
设 $\mathcal{A}$ 为阿贝尔范畴，$f : U \to V$ 为
$\mathcal{A}$ 的单纯对象之间的态射。若对所有 $i$，诱导态射
$N(f)_i : N(U)_i \to N(V)_i$ 都是单射，则对所有 $i$，$f_i$
都是单射。把“单射”替换成“满射”或“同构”时，结论同样成立。
\end{lemma}

\begin{proof}
这由引理 \ref{simplicial-lemma-splitting-abelian-category} 和分裂的定义立即推出。
\end{proof}


\begin{lemma}
\label{simplicial-lemma-N-d-in-N}
设 $\mathcal{A}$ 为阿贝尔范畴，$U$ 为 $\mathcal{A}$ 中的单纯对象。
令 $N(U_m)$ 如上面的引理 \ref{simplicial-lemma-splitting-abelian-category} 所述。
% P02-E0113: frozen source has “Let N(U_m) as in ...” without a finite verb.
则 $d^m_m(N(U_m)) \subset N(U_{m - 1})$。
\end{lemma}

\begin{proof}
对 $j = 0, \ldots, m - 2$，由备注 \ref{simplicial-remark-relations} 中的关系，有
$d^{m - 1}_j d^m_m = d^{m - 1}_{m - 1} d^m_j$。结论随即成立。
\end{proof}

\begin{lemma}
\label{simplicial-lemma-simplicial-abelian-n-skel-sub}
设 $\mathcal{A}$ 为阿贝尔范畴，$U$ 为 $\mathcal{A}$ 的单纯对象，
$n \geq 0$ 为整数。规则
$$
U'_m = \sum\nolimits_{\varphi : [m] \to [i], \ i\leq n} \Im(U(\varphi))
$$
定义了子单纯对象 $U' \subset U$，并且当 $i \leq n$ 时
$U'_i = U_i$。此外，对所有 $m > n$，都有 $N(U'_m) = 0$。
\end{lemma}

\begin{proof}
取 $m$、$i \leq n$ 和某个 $\varphi : [m] \to [i]$。
对任意 $\psi : [m'] \to [m]$，$\Im(U(\varphi))$ 在 $U(\psi)$
下的像等于 $U(\varphi \circ \psi)$ 的像，并且
$\varphi \circ \psi : [m'] \to [i]$。因此 $U'$ 是单纯对象。
取 $m > n$。必须证明 $N(U'_m) = 0$。由 $N(U_m)$ 和 $N(U'_m)$
的定义，有 $N(U'_m) = U'_m \cap N(U_m)$（子对象的交）。
由于引理 \ref{simplicial-lemma-splitting-abelian-category} 表明 $U$ 是分裂的，
只需证明 $U'_m$ 包含于和
$$
\sum\nolimits_{\varphi : [m] \to [m']\text{ 满射}, \ m' < m}
\Im(U(\varphi)|_{N(U_{m'})}).
$$
由分裂，每个 $U_{m'}$ 都是 $N(U_{m''})$ 在满射
$\psi : [m'] \to [m'']$ 所诱导的 $U(\psi)$ 下的像之和。
因此，上面显示的和等于
$$
\sum\nolimits_{\varphi : [m] \to [m']\text{ 满射}, \ m' < m}
\Im(U(\varphi)).
$$
显然，$U'_m$ 包含于其中：在 $U'_m$ 定义中出现的任意
$\varphi : [m] \to [i]$（$i \leq n$）都可分解为
$[m] \to [m'] \to [i]$，其中 $[m] \to [m']$ 是满射，且
$m' < m$，正如上面最后一个显示的和中那样。
\end{proof}

\section{余骨架函子}
\label{simplicial-section-coskeleton}

\noindent
设 $\mathcal{C}$ 为范畴。{\it 余骨架函子}（若存在）是函子
$$
\text{cosk}_n :
\text{Simp}_n(\mathcal{C}) \longrightarrow \text{Simp}(\mathcal{C})
$$
它是骨架函子的右伴随。用公式表示为
\begin{equation}
\label{simplicial-equation-cosk}
\Mor_{\text{Simp}(\mathcal{C})}(U, \text{cosk}_n V)
=
\Mor_{\text{Simp}_n(\mathcal{C})}(\text{sk}_n U, V)
\end{equation}
给定 $n$-截断单纯对象 $V$，如果存在
$\text{cosk}_nV \in \Ob(\text{Simp}(\mathcal{C}))$ 以及态射
$\text{sk}_n \text{cosk}_n V \to V$，使得上述公式成立，则称
{\it $\text{cosk}_nV$ 存在}；换言之，函子
$U \mapsto \Mor_{\text{Simp}_n(\mathcal{C})}(\text{sk}_n U, V)$
是可表示的。若它存在，则由 Yoneda 引理，它在唯一同构的意义下唯一。
见 Categories，第 \ref{categories-section-opposite} 节。

\begin{example}
\label{simplicial-example-cosk0}
假设范畴 $\mathcal{C}$ 有有限非空自积。$\mathcal{C}$ 的
$0$-截断单纯对象与 $\mathcal{C}$ 的对象 $X$ 是同一回事。
在这种情形下，我们断言 $\text{cosk}_0(X)$ 是单纯对象 $U$，其中
$U_n = X^{n + 1}$ 是 $X$ 的 $(n + 1)$ 重自积，而其单纯对象结构如例
\ref{simplicial-example-fibre-products-simplicial-object} 所述。事实上，态射
$V \to U$（其中 $V$ 是单纯对象）由态射 $V_n \to X^{n + 1}$ 给出，
并且所有图表
$$
\xymatrix{
V_n \ar[r] \ar[d]_{V([0] \to [n], 0 \mapsto i)} &
X^{n + 1} \ar[d]^{\text{pr}_i} \\
V_0 \ar[r] &
X
}
$$
都可交换。显然，这意味着该映射确定并且由唯一态射 $V_0 \to X$ 确定。
这就证明了公式 (\ref{simplicial-equation-cosk})。
\end{example}

\noindent
回忆范畴 $\Delta/[n]$，见例 \ref{simplicial-example-simplex-category}。
以 $(\Delta/[n])_{\leq m}$ 表示 $\Delta/[n]$ 的全子范畴，
其对象为 $\Delta/[n]$ 中满足 $k \leq m$ 的 $[k] \to [n]$。
换言之，有如下范畴与函子的可交换图表：
$$
\xymatrix{
(\Delta/[n])_{\leq m} \ar[r] \ar[d] &
\Delta/[n] \ar[d] \\
\Delta_{\leq m} \ar[r] &
\Delta
}
$$
给定 $\mathcal{C}$ 的 $m$-截断单纯对象 $U$，按下列规则定义函子
$$
U(n) : (\Delta/[n])_{\leq m}^{opp} \longrightarrow \mathcal{C}
$$
：
\begin{eqnarray*}
([k] \to [n]) & \longmapsto & U_k \\
\psi : ([k'] \to [n]) \to ([k] \to [n]) &
\longmapsto &
U(\psi) : U_k \to U_{k'}
\end{eqnarray*}
对 $\Delta$ 的给定态射 $\varphi : [n] \to [n']$，有相应函子
$$
\overline{\varphi} :
(\Delta/[n])_{\leq m}
\longrightarrow
(\Delta/[n'])_{\leq m}
$$
它把 $\alpha : [k] \to [n]$ 送到
$\varphi \circ \alpha : [k] \to [n']$。复合
$U(n') \circ \overline{\varphi}$ 等于函子 $U(n)$。

\begin{lemma}
\label{simplicial-lemma-existence-cosk}
若范畴 $\mathcal{C}$ 有有限极限，则对所有 $m$，函子
$\text{cosk}_m$ 都存在。此外，对任意 $m$-截断单纯对象 $U$，
单纯对象 $\text{cosk}_mU$ 由公式
$$
(\text{cosk}_mU)_n = \lim_{(\Delta/[n])_{\leq m}^{opp}} U(n)
$$
描述；而对 $\varphi : [n] \to [n']$，映射
$\text{cosk}_mU(\varphi)$ 由上面的等同
$U(n') \circ \overline{\varphi} = U(n)$，通过 Categories，引理
\ref{categories-lemma-functorial-limit} 得到。
\end{lemma}

\begin{proof}
在本引理的证明中，以 $\text{cosk}_mU$ 表示单纯对象，其中
$(\text{cosk}_mU)_n$ 等于
$\lim_{(\Delta/[n])_{\leq m}^{opp}} U(n)$。
证明最后将说明，它确实满足所需的映射性质。

\medskip\noindent
设 $V$ 为单纯对象。态射 $\gamma : V \to \text{cosk}_mU$
由一列态射 $\gamma_n : V_n \to (\text{cosk}_mU)_n$ 给出。
由极限的定义，这等价于给出一族态射
$\gamma(\alpha) : V_n \to U_k$，其中 $\alpha$ 取遍所有满足
$k \leq m$ 的 $\alpha : [k] \to [n]$。对任意
$0 \leq k, k' \leq m$、$0 \leq n, n'$ 以及 $\Delta$ 中满足
$\varphi \circ \alpha = \alpha' \circ \psi$ 的
$\psi : [k] \to [k']$、$\varphi : [n] \to [n']$、
$\alpha : [k] \to [n]$ 和 $\alpha' : [k'] \to [n']$，
这些态射还满足图表
$$
\xymatrix{
V_n \ar[r]_{\gamma(\alpha)} &  U_k \\
V_{n'} \ar[r]^{\gamma(\alpha')} \ar[u]^{V(\varphi)} & U_{k'} \ar[u]_{U(\psi)}
}
$$
可交换。取 $n = k = k'$、$\varphi = \alpha'$，以及
$\alpha = \psi = \text{id}_{[k]}$，得到
$\gamma(\alpha') = \gamma(\text{id}_{[k]}) \circ V(\alpha')$。
换言之，态射 $\gamma(\text{id}_{[k]})$（$k \leq m$）确定态射
$\gamma$。容易看出，这些态射组成一个态射
$\text{sk}_m V \to U$。

\medskip\noindent
反过来，给定态射 $\gamma : \text{sk}_m V \to U$，令
$\gamma(\alpha) = \gamma(\text{id}_{[k]}) \circ V(\alpha)$，
便得到一族态射 $\gamma(\alpha)$，其中 $\alpha$ 取遍所有满足
$k \leq m$ 的 $\alpha : [k] \to [n]$。这些态射满足上图所示的全部
可交换条件，因而给出态射 $V \to \text{cosk}_m U$。
\end{proof}

\begin{lemma}
\label{simplicial-lemma-trivial-cosk}
设 $\mathcal{C}$ 为范畴，$U$ 为 $\mathcal{C}$ 的 $m$-截断单纯对象。
对 $n \leq m$，极限
$\lim_{(\Delta/[n])_{\leq m}^{opp}} U(n)$ 存在，并典范同构于 $U_n$。
\end{lemma}

\begin{proof}
这是因为在这种情形下，范畴 $(\Delta/[n])_{\leq m}$ 有一个终对象，
即恒等映射 $[n] \to [n]$。
% P02-E0114: frozen source says “an final object”.
\end{proof}

\begin{lemma}
\label{simplicial-lemma-recover-cosk}
设 $\mathcal{C}$ 为有有限极限的范畴，$U$ 为 $\mathcal{C}$ 的
$n$-截断单纯对象。态射 $\text{sk}_n \text{cosk}_n U \to U$ 是同构。
\end{lemma}

\begin{proof}
合并引理 \ref{simplicial-lemma-existence-cosk} 与引理 \ref{simplicial-lemma-trivial-cosk}。
\end{proof}

\noindent
下面更详细地描述余骨架函子的一个特例。滥用记号，对任意
$n' \geq n$，也以 $\text{sk}_n$ 表示限制函子
$\text{Simp}_{n'}(\mathcal{C}) \to \text{Simp}_n(\mathcal{C})$。
我们将描述函子
$\text{sk}_n : \text{Simp}_{n + 1}(\mathcal{C})
\to \text{Simp}_n(\mathcal{C})$ 的一个右伴随。
对 $n \geq 1$、$0 \leq i < j \leq n + 1$，定义
$\delta^{n + 1}_{i, j} : [n - 1] \to [n + 1]$
为略去 $i$ 和 $j$ 的递增映射。注意
$\delta^{n + 1}_{i, j} =
\delta^{n + 1}_j \circ \delta^n_i =
\delta^{n + 1}_i \circ \delta^n_{j - 1}$；见引理
\ref{simplicial-lemma-relations-face-degeneracy}。这引出下述引理。

\begin{lemma}
\label{simplicial-lemma-formula-limit}
设整数 $n$，并要求其值 $\geq 1$。设 $U$ 为 $\mathcal{C}$ 的
$n$-截断单纯对象。考虑从 $\mathcal{C}$ 到 $\textit{Sets}$ 的反变函子，
它把对象 $T$ 送到集合
$$
\{ (f_0, \ldots, f_{n + 1}) \in \Mor_\mathcal{C}(T, U_n)
\mid
d^n_{j - 1} \circ f_i = d^n_i \circ f_j
\ \forall\ 0\leq i < j\leq n + 1\}
$$
若此函子由 $\mathcal{C}$ 的某个对象 $U_{n + 1}$ 表示，则
$$
U_{n + 1} = \lim_{(\Delta/[n + 1])_{\leq n}^{opp}} U(n)
$$
\end{lemma}

\begin{proof}
该极限若存在，就表示把对象 $T$ 送到集合
$$
\{
(f_\alpha)_{\alpha : [k] \to [n + 1], k \leq n}
\mid
f_{\alpha \circ \psi} = U(\psi) \circ f_\alpha\ \forall
\ \psi : [k'] \to [k], \alpha : [k] \to [n + 1]
\}.
$$
事实上，我们将证明此函子同构于引理中显示的函子。一个方向的映射由规则
$$
(f_\alpha)_{\alpha}
\longmapsto
(f_{\delta^{n + 1}_0}, \ldots, f_{\delta^{n + 1}_{n + 1}}).
$$
给出。由引理前回忆的关系，它满足引理中的条件，因为
$$
d^n_{j - 1} \circ f_{\delta^{n + 1}_i} =
f_{\delta^{n + 1}_i \circ \delta^n_{j - 1}} =
f_{\delta^{n + 1}_j \circ \delta^n_i} =
d^n_i \circ f_{\delta^{n + 1}_j}
$$
为构造反方向的映射，必须把引理中显示公式里的系统
$(f_0, \ldots, f_{n + 1})$ 送到映射系统 $f_\alpha$。
给定 $\alpha : [k] \to [n + 1]$。由于 $k \leq n$，映射 $\alpha$
不是满射。因此，可以写成
$\alpha = \delta^{n + 1}_i \circ \psi$，其中
$0 \leq i \leq n + 1$，且 $\psi : [k] \to [n]$。
我们只能定义
$$
f_\alpha = U(\psi) \circ f_i.
$$
当然，必须检查它与对偶 $(i, \psi)$ 的选择无关。首先注意，给定 $i$，
符合条件的 $\psi$ 唯一。其次，设 $(j, \phi)$ 为另一对。
则 $i \not = j$，并可假设 $i < j$。由于 $i, j$ 都不在
$\alpha$ 的像中，实际上可写
$\alpha = \delta^{n + 1}_{i, j} \circ \xi$；于是
$\psi = \delta^n_{j - 1} \circ \xi$，且
$\phi = \delta^n_i \circ \xi$。因此
\begin{eqnarray*}
U(\psi) \circ f_i & = & U(\delta^n_{j - 1} \circ \xi) \circ f_i \\
& = & U(\xi) \circ d^n_{j - 1} \circ f_i \\
& = & U(\xi) \circ d^n_i \circ f_j \\
& = & U(\delta^n_i \circ \xi) \circ f_j \\
& = & U(\phi) \circ f_j
\end{eqnarray*}
如所欲证。还需验证，如此定义的映射 $f_\alpha$ 满足映射系统
$(f_\alpha)_\alpha$ 的规则。为此，设
$\psi : [k'] \to [k]$、$\alpha : [k] \to [n + 1]$，其中
$k, k' \leq n$。令 $\alpha' = \alpha \circ \psi$。
取不在 $\alpha$ 的像中的 $i$，则 $i$ 显然也不在 $\alpha'$ 的像中。
写 $\alpha = \delta^{n + 1}_i \circ \phi$（这里不能使用字母 $\psi$，
因为它已经用过）。则显然
$\alpha' = \delta^{n + 1}_i \circ \phi \circ \psi$。
由上述构造，有
$$
U(\psi) \circ f_\alpha = U(\psi) \circ U(\phi) \circ f_i
= U(\phi \circ \psi) \circ f_i = f_{\alpha \circ \psi} = f_{\alpha'}
$$
如所欲证。留给读者一项愉快的工作：验证这些构造给出互逆双射，并且关于
$T$ 函子性地成立。
\end{proof}

\begin{lemma}
\label{simplicial-lemma-work-out}
设整数 $n$，并要求其值 $\geq 1$。设 $U$ 为 $\mathcal{C}$ 的
$n$-截断单纯对象。考虑从 $\mathcal{C}$ 到 $\textit{Sets}$ 的反变函子，
它把对象 $T$ 送到集合
$$
\{ (f_0, \ldots, f_{n + 1}) \in \Mor_\mathcal{C}(T, U_n)
\mid
d^n_{j - 1} \circ f_i = d^n_i \circ f_j
\ \forall\ 0\leq i < j\leq n + 1\}
$$
若此函子由 $\mathcal{C}$ 的某个对象 $U_{n + 1}$ 表示，则存在
$(n + 1)$-截断单纯对象 $\tilde U$，满足
$\text{sk}_n \tilde U = U$ 和 $\tilde U_{n + 1} = U_{n + 1}$，
并且有如下伴随关系：
$$
\Mor_{\text{Simp}_{n + 1}(\mathcal{C})}(V, \tilde U)
=
\Mor_{\text{Simp}_n(\mathcal{C})}(\text{sk}_nV, U)
$$
\end{lemma}

\begin{proof}
由引理 \ref{simplicial-lemma-trivial-cosk}，对 $0 \leq i \leq n$ 有等同
$$
U_i = \lim_{(\Delta/[i])_{\leq n}^{opp}} U(i)
$$
由引理 \ref{simplicial-lemma-formula-limit}，有
$$
U_{n + 1} = \lim_{(\Delta/[n + 1])_{\leq n}^{opp}} U(n).
$$
因此，对任意满足 $i, j \leq n + 1$ 的 $\varphi : [i] \to [j]$，
可以完全按照引理 \ref{simplicial-lemma-existence-cosk} 定义相应映射
$\tilde U(\varphi) : \tilde U_j \to \tilde U_i$。
这定义了满足 $\text{sk}_n \tilde U = U$ 的
$(n + 1)$-截断单纯对象 $\tilde U$。

\medskip\noindent
为证明伴随关系，论证如下。给定公式右端的任意元素
$\gamma : \text{sk}_n V \to U$，考虑态射
$f_i = \gamma_n \circ d^{n + 1}_i : V_{n + 1} \to V_n \to U_n$。
它们显然满足关系
$d^n_{j - 1} \circ f_i = d^n_i \circ f_j$，因而由 $U_{n + 1}$
的选取，定义唯一态射 $V_{n + 1} \to U_{n + 1}$。
反之，给定左端的态射 $\gamma' : V \to \tilde U$，只需限制到
$\Delta_{\leq n}$，便得到右端的一个元素。留给读者证明这两个构造互逆。
\end{proof}

\begin{remark}
\label{simplicial-remark-explicit-face-degeneracy}
设 $U$ 和 $U_{n + 1}$ 如引理 \ref{simplicial-lemma-work-out} 中所述。
在 $T$-值点上，可以容易地描述 $\tilde U$ 的面映射和退化映射。
具体地，映射 $d^{n + 1}_i : U_{n + 1} \to U_n$ 由
$$
(f_0, \ldots, f_{n + 1}) \longmapsto f_i.
$$
给出，而映射 $s^n_j : U_n \to U_{n + 1}$ 由
\begin{eqnarray*}
f & \longmapsto & (
s^{n - 1}_{j - 1} \circ d^{n - 1}_0 \circ f, \\
& &
s^{n - 1}_{j - 1} \circ d^{n - 1}_1 \circ f, \\
& &
\ldots\\
& &
s^{n - 1}_{j - 1} \circ d^{n - 1}_{j - 1} \circ f, \\
& &
f, \\
& &
f, \\
& &
s^{n - 1}_j \circ d^{n - 1}_{j + 1} \circ f, \\
& &
s^{n - 1}_j \circ d^{n - 1}_{j + 2} \circ f, \\
& &
\ldots\\
& &
s^{n - 1}_j \circ d^{n - 1}_n \circ f
)
\end{eqnarray*}
给出；留给读者验证，右端是引理 \ref{simplicial-lemma-work-out} 中显示集合的元素。
当 $n = 0$ 时，有一个映射，即 $f \mapsto (f, f)$。
当 $n = 1$ 时，有两个映射，即
$f \mapsto (f, f, s_0d_1f)$ 和
$f \mapsto (s_0d_0f, f, f)$。
当 $n = 2$ 时，有三个映射，即
$f \mapsto (f, f, s_0d_1f, s_0d_2f)$、
$f \mapsto (s_0d_0f, f, f, s_1d_2f)$ 和
$f \mapsto (s_1d_0f, s_1d_1f, f, f)$。依此类推。
\end{remark}


\begin{remark}
\label{simplicial-remark-cosk-simplicial-sets}
上面的引理 \ref{simplicial-lemma-work-out} 在单纯集情形下的构造如下。
给定 $n$-截断单纯集 $U$，如下作出典范的 $(n + 1)$-截断单纯集
$\tilde U$。利用引理中的公式添加一个 $(n + 1)$-单形集
$U_{n + 1}$。也就是说，$U_{n + 1}$ 的一个元素是一列编号的
$n$-单形 $(f_0, \ldots, f_{n + 1})$，它们具有与一个
$(n + 1)$-单形的各面相同的粘合性质。换言之，当 $i < j$ 时，
$f_j$ 的第 $i$ 个面是 $f_i$ 的第 $(j-1)$ 个面。
从几何上看，如何定义 $\tilde U$ 的面映射和退化映射是显然的。
若 $V$ 是 $(n + 1)$-截断单纯集，则它的 $(n + 1)$-单形给出相容的
$n$-单形族 $(f_0, \ldots, f_{n + 1})$，其中 $f_i \in V_n$。
因此有自然映射
$\Mor(\text{sk}_nV, U) \to \Mor(V, \tilde U)$，
它与反方向的典范限制映射互逆。

\medskip\noindent
此外，只需在截断单纯集的情形下完成该构造的组合论部分。
事实上，对范畴 $\mathcal{C}$ 的任意对象 $T$ 和 $\mathcal{C}$ 的任意
$n$-截断单纯对象 $U$，可以考虑 $n$-截断单纯集 $\Mor(T, U)$。
可以对此应用上述构造，取其 $(n + 1)$-单形集，并要求它可表示。
这是理解引理 \ref{simplicial-lemma-work-out} 结论的一种好方法。
\end{remark}

\begin{remark}
\label{simplicial-remark-inductive-coskeleton}
{\it 余骨架的归纳构造。}
假设 $\mathcal{C}$ 是有有限极限的范畴，$U$ 是 $\mathcal{C}$ 中的
$m$-截断单纯对象。可以如下归纳构造 $n$-截断对象 $U^n$：
\begin{enumerate}
\item 从 $U^m = U$ 开始。
\item 给定 $n \geq m$ 时的 $U^n$，令 $U^{n + 1} = \tilde U^n$，
其中 $\tilde U^n$ 按引理 \ref{simplicial-lemma-work-out} 由 $U^n$ 构造。
\end{enumerate}
由于引理 \ref{simplicial-lemma-work-out} 的构造保持 $U^n$ 的 $n$-骨架不变，
可以把 $\text{cosk}_m U$ 定义为满足
$(\text{cosk}_m U)_n = U^n_n = U^{n + 1}_n = \ldots$ 的单纯对象。
而且由引理 \ref{simplicial-lemma-work-out} 形式地推出，对所有 $n \geq m$，
$U^n$ 满足公式
$$
\Mor_{\text{Simp}_n(\mathcal{C})}(V, U^n)
=
\Mor_{\text{Simp}_m(\mathcal{C})}(\text{sk}_mV, U)
$$
进而又形式地推出
$$
\Mor_{\text{Simp}(\mathcal{C})}(V, \text{cosk}_mU)
=
\Mor_{\text{Simp}_m(\mathcal{C})}(\text{sk}_mV, U)
$$
其中 $\text{cosk}_mU$ 按上述方式选取。
\end{remark}

\begin{lemma}
\label{simplicial-lemma-cosk-up}
设 $\mathcal{C}$ 为有有限极限的范畴。
\begin{enumerate}
\item 对每个 $n$，函子 $\text{sk}_n : \text{Simp}(\mathcal{C})
\to \text{Simp}_n(\mathcal{C})$ 有右伴随 $\text{cosk}_n$。
\item 对每个 $n' \geq n$，函子
$\text{sk}_n : \text{Simp}_{n'}(\mathcal{C}) \to \text{Simp}_n(\mathcal{C})$
有右伴随，即 $\text{sk}_{n'}\text{cosk}_n$。
\item 对每个 $m \geq n \geq 0$ 以及 $\mathcal{C}$ 的每个
$n$-截断单纯对象 $U$，有
$\text{cosk}_m \text{sk}_m \text{cosk}_n U = \text{cosk}_n U$。
\item 若 $U$ 是 $\mathcal{C}$ 的单纯对象，并且对某个 $n \geq 0$，
典范映射 $U \to \text{cosk}_n \text{sk}_nU$ 是同构，则对所有
$m \geq n$，典范映射 $U \to \text{cosk}_m \text{sk}_mU$ 都是同构。
\end{enumerate}
\end{lemma}

\begin{proof}
(1) 中的存在性由上面的引理 \ref{simplicial-lemma-existence-cosk} 推出。
(2) 和 (3) 由备注 \ref{simplicial-remark-inductive-coskeleton} 中的讨论推出。
于是 (4) 显然。
\end{proof}

\begin{remark}
\label{simplicial-remark-existence-cosk}
定义余骨架函子并不需要所有有限极限。说明如下：
\begin{enumerate}
\item 在例 \ref{simplicial-example-cosk0} 中已经看到，若 $\mathcal{C}$ 有对象对的积，
则 $\text{cosk}_0$ 存在。
\item 对 $k > 0$，若 $\mathcal{C}$ 有有限连通极限，则函子
$\text{cosk}_k$ 存在。
\end{enumerate}
这由余骨架的归纳构造步骤（备注 \ref{simplicial-remark-cosk-simplicial-sets} 和
\ref{simplicial-remark-inductive-coskeleton}）清楚可见；也可以由下述事实推出：
引理 \ref{simplicial-lemma-existence-cosk} 中使用的范畴
$(\Delta/[n])_{\leq k}$ 在 $k \geq 1$ 且 $n \geq k + 1$ 时是连通的。
注意，由引理 \ref{simplicial-lemma-trivial-cosk} 或引理 \ref{simplicial-lemma-recover-cosk}，
我们不需要 $n \leq k$ 时的那些范畴。（随着 $k$ 增大，
在 $k \geq 1$ 且 $n \geq k + 1$ 时，范畴
$(\Delta/[n])_{\leq k}$ 在拓扑意义下越来越连通。）
\end{remark}

\begin{lemma}
\label{simplicial-lemma-cosk-product}
设 $U$、$V$ 为范畴 $\mathcal{C}$ 的 $n$-截断单纯对象。只要等式两端
都存在，就有
$$
\text{cosk}_n (U \times V) = \text{cosk}_nU \times \text{cosk}_nV
$$
\end{lemma}

\begin{proof}
设 $W$ 为单纯对象。有
\begin{eqnarray*}
\Mor(W, \text{cosk}_n (U \times V))
& = &
\Mor(\text{sk}_n W, U \times V) \\
& = &
\Mor(\text{sk}_n W, U)
\times
\Mor(\text{sk}_nW, V) \\
& = &
\Mor(W, \text{cosk}_n U)
\times
\Mor(W, \text{cosk}_n V) \\
& = &
\Mor(W, \text{cosk}_n U \times \text{cosk}_n V)
\end{eqnarray*}
引理得证。
\end{proof}

\begin{lemma}
\label{simplicial-lemma-cosk-fibre-product}
假设 $\mathcal{C}$ 有纤维积。设 $U \to V$ 和 $W \to V$ 为
范畴 $\mathcal{C}$ 的 $n$-截断单纯对象之间的态射。只要等式两端存在，
就有
$$
\text{cosk}_n (U \times_V W)
=
\text{cosk}_nU \times_{\text{cosk}_n V} \text{cosk}_nW
$$
\end{lemma}

\begin{proof}
略；证明与上面的引理 \ref{simplicial-lemma-cosk-product} 的证明非常相似。
\end{proof}

\begin{lemma}
\label{simplicial-lemma-cosk-above-object}
设 $\mathcal{C}$ 为有有限极限的范畴，并设
$X \in \Ob(\mathcal{C})$。对 $k \geq 1$，函子
$\mathcal{C}/X \to \mathcal{C}$ 与余骨架函子 $\text{cosk}_k$ 交换。
\end{lemma}

\begin{proof}
该断言的意思是：若 $U$ 是 $\mathcal{C}/X$ 的单纯对象，
可把它看成 $\mathcal{C}$ 的一个单纯对象，连同一个指向常值单纯对象
$X$ 的态射；那么在 $\mathcal{C}/X$ 中计算的 $\text{cosk}_k U$
与在 $\mathcal{C}$ 中计算的相同。例如，这由 Categories，引理
\ref{categories-lemma-connected-limit-over-X} 推出，因为引理
\ref{simplicial-lemma-existence-cosk} 中使用的范畴 $(\Delta/[n])_{\leq k}$
在 $k \geq 1$ 且 $n \geq k + 1$ 时是连通的。
注意，由引理 \ref{simplicial-lemma-trivial-cosk} 或引理 \ref{simplicial-lemma-recover-cosk}，
我们不需要 $n \leq k$ 时的那些范畴。
\end{proof}

\begin{lemma}
\label{simplicial-lemma-simplex-cosk}
典范映射 $\Delta[n] \to \text{cosk}_1 \text{sk}_1 \Delta[n]$ 是同构。
\end{lemma}

\begin{proof}
考虑单纯集 $U$ 和态射 $f : U \to \Delta[n]$。这是一条规则：
它把每个 $u \in U_i$ 关联到 $\Delta$ 中的映射
$f_u : [i] \to [n]$。此外，这些映射应满足：对任意
$\varphi : [j] \to [i]$，有
$f_u \circ \varphi = f_{U(\varphi)(u)}$。
以 $\epsilon^i_j : [0] \to [i]$ 表示把 $0$ 映到 $j$ 的映射。
以 $F : U_0 \to [n]$ 表示映射 $u \mapsto f_u(0)$。于是可见
$$
f_u(j) = F(\epsilon^i_j(u))
$$
对所有 $0 \leq j \leq i$ 和 $u \in U_i$ 成立。特别地，若知道函数
$F$，则对所有 $u\in U_i$ 和所有 $i$，都知道映射 $f_u$。
反之，给定映射 $F : U_0 \to [n]$，可以对任意 $i$、
任意 $u \in U_i$ 和任意 $0 \leq j \leq i$ 令
$$
f_u(j) = F(\epsilon^i_j(u))
$$
一般而言，这并不定义上面的单纯集态射 $f$。条件是所有映射
$f_u$ 都非降。这显然等价于：当 $0 \leq j \leq j' \leq i$
且 $u \in U_i$ 时，
$F(\epsilon^i_j(u)) \leq F(\epsilon^i_{j'}(u))$。
但在这种情形下，态射
$$
\epsilon^i_j, \epsilon^i_{j'} : [0] \to [i]
$$
都通过映射 $\epsilon^i_{j, j'} : [1] \to [i]$ 分解；后者由规则
$0 \mapsto j$、$1 \mapsto j'$ 定义。换言之，只需对 $i = 1$
和 $u \in X_1$ 检查这些不等式。
% P02-E0115: frozen source switches from the simplicial set U to undefined X_1.
也就是说，有
$$
\Mor(U, \Delta[n])
=
\Mor(\text{sk}_1 U, \text{sk}_1 \Delta[n])
$$
如所欲证。
\end{proof}



\section{增广}
\label{simplicial-section-augmentation}

\begin{definition}
\label{simplicial-definition-augmentation}
设 $\mathcal{C}$ 为范畴，$U$ 为 $\mathcal{C}$ 的单纯对象。
$U$ 向 $\mathcal{C}$ 的对象 $X$ 的一个{\it 增广
$\epsilon : U \to X$}，是从 $U$ 到常值单纯对象 $X$ 的态射。
\end{definition}

\begin{lemma}
\label{simplicial-lemma-augmentation-howto}
设 $\mathcal{C}$ 为范畴，$X \in \Ob(\mathcal{C})$，并设 $U$ 为
$\mathcal{C}$ 的单纯对象。给出 $U$ 向 $X$ 的一个增广，等价于给出
态射 $\epsilon_0 : U_0 \to X$，使得
$\epsilon_0 \circ d^1_0 = \epsilon_0 \circ d^1_1$。
\end{lemma}

\begin{proof}
给定态射 $\epsilon : U \to X$，当然得到引理中所述的
$\epsilon_0$。反之，给定引理中所述的 $\epsilon_0$，任选态射
$\alpha : [0] \to [n]$，并令
$\epsilon_n = \epsilon_0 \circ U(\alpha)$，从而定义
$\epsilon_n : U_n \to X$。事实上，若
$\beta : [0] \to [n]$ 是另一种选择，则存在态射
$\gamma : [1] \to [n]$，使得 $\alpha$ 和 $\beta$ 都分解为
$[0] \to [1] \to [n]$。因此，对 $\epsilon_0$ 的条件说明
$\epsilon_n$ 是良定义的。容易证明，
$(\epsilon_n) : U \to X$ 是单纯对象之间的态射。
\end{proof}

\begin{lemma}
\label{simplicial-lemma-cosk-minus-one}
设 $\mathcal{C}$ 为有纤维积的范畴。设 $f : Y\to X$ 为
$\mathcal{C}$ 的态射，并设 $U$ 为 $\mathcal{C}$ 的单纯对象，
其第 $n$ 项是 $(n + 1)$ 重纤维积
$Y \times_X Y \times_X \ldots \times_X Y$。
% P02-E0116: frozen source has “(n + 1)fold” without a hyphen.
见例 \ref{simplicial-example-fibre-products-simplicial-object}。
对 $\mathcal{C}$ 的任意单纯对象 $V$，有
\begin{align*}
\Mor_{\text{Simp}(\mathcal{C})}(V, U)
& =
\Mor_{\text{Simp}_1(\mathcal{C})}(\text{sk}_1 V, \text{sk}_1 U) \\
& =
\{g_0 : V_0 \to Y \mid f \circ g_0 \circ d^1_0 = f \circ g_0 \circ d^1_1\}
\end{align*}
特别地，有 $U = \text{cosk}_1 \text{sk}_1 U$。
\end{lemma}

\begin{proof}
设 $g : \text{sk}_1V \to \text{sk}_1U$ 为 $1$-截断单纯对象之间的态射。
则图表
$$
\xymatrix{
V_1 \ar@<1ex>[r]^{d^1_0} \ar@<-1ex>[r]_{d^1_1} \ar[d]_{g_1} &
V_0 \ar[d]^{g_0} \\
Y \times_X Y \ar@<1ex>[r]^{pr_1} \ar@<-1ex>[r]_{pr_0} &
Y \ar[r] & X
}
$$
可交换，这证明了引理中显示的关系。反过来，必须证明：给定满足关系
$f \circ g_0 \circ d^1_0 = f \circ g_0 \circ d^1_1$ 的态射
$g_0$，可得到唯一的单纯对象态射 $g : V \to U$。
作法如下。对任意 $n \geq 1$，令
$g_{n, i} = g_0 \circ V([0] \to [n], 0 \mapsto i) :
V_n \to Y$。由可交换图表
$$
\xymatrix{
[0] \ar[rd]_{0 \mapsto 0} \ar[rrrrrd]^{0 \mapsto i} \\
& [1] \ar[rrrr]^{0 \mapsto i, 1\mapsto i + 1} & & & & [n] \\
[0] \ar[ru]^{0 \mapsto 1} \ar[rrrrru]_{0 \mapsto i + 1}
}
$$
上述等式蕴含 $f \circ g_{n, i} = f \circ g_{n, i + 1}$。
因此得到
$(g_{n, 0}, \ldots, g_{n, n}) : V_n \to Y \times_X\ldots \times_X Y = U_n$。
留给读者验证，这是单纯对象之间的态射。引理最后一个断言等价于引理中
显示公式的第一个等式。
\end{proof}

\begin{remark}
\label{simplicial-remark-augmentation}
设 $\mathcal{C}$ 为有纤维积的范畴，$V$ 为单纯对象，
$\epsilon : V \to X$ 为增广。设 $U$ 为如下单纯对象：其第 $n$ 项
是 $V_0$ 在 $X$ 上的第 $(n + 1)$ 个纤维积。
% P02-E0117: frozen source says “(n + 1)st fibred product”; the construction is the (n + 1)-fold fibre product.
简单地结合引理 \ref{simplicial-lemma-augmentation-howto} 与引理
\ref{simplicial-lemma-cosk-minus-one}，得到典范态射 $V \to U$。
\end{remark}





\section{骨架函子的左伴随}
\label{simplicial-section-adjoint-left}

\noindent
本节在某些情形下构造骨架函子 $\text{sk}_m$ 的左伴随 $i_{m!}$。
伴随公式为
$$
\Mor_{\text{Simp}_m(\mathcal{C})}(U, \text{sk}_mV)
=
\Mor_{\text{Simp}(\mathcal{C})}(i_{m!}U, V).
$$
事实表明，当范畴 $\mathcal{C}$ 有有限余极限时，这个左伴随存在。

\medskip\noindent
我们使用与第 \ref{simplicial-section-skeleton} 节相似的构造。回忆 $[n]$ 下对象
所成的范畴 $[n]/\Delta$，见 Categories，例
\ref{categories-example-category-under-X}。其对象是态射
$\alpha : [n] \to [k]$，其态射是可交换三角形。
以 $([n]/\Delta)_{\leq m}$ 表示 $[n]/\Delta$ 的全子范畴，
由满足 $k \leq m$ 的对象 $[n] \to [k]$ 组成。
给定 $\mathcal{C}$ 的 $m$-截断单纯对象 $U$，按下列规则定义函子
$$
U(n) : ([n]/\Delta)_{\leq m}^{opp} \longrightarrow \mathcal{C}
$$
：
\begin{eqnarray*}
([n] \to [k]) & \longmapsto & U_k \\
\psi : ([n] \to [k']) \to ([n] \to [k])
& \longmapsto &
U(\psi) : U_k \to U_{k'}
\end{eqnarray*}
对 $\Delta$ 的给定态射 $\varphi : [n] \to [n']$，有相应函子
$$
\underline{\varphi} :
([n']/\Delta)_{\leq m}
\longrightarrow
([n]/\Delta)_{\leq m}
$$
它把 $\alpha : [n'] \to [k]$ 送到
$\varphi \circ \alpha : [n] \to [k]$。复合
$U(n) \circ \underline{\varphi}$ 等于函子 $U(n')$。

\begin{lemma}
\label{simplicial-lemma-left-adjoint-exists}
设 $\mathcal{C}$ 为有有限余极限的范畴。对所有 $m$，函子 $i_{m!}$
都存在。设 $U$ 为 $\mathcal{C}$ 的 $m$-截断单纯对象。
单纯对象 $i_{m!}U$ 由公式
$$
(i_{m!}U)_n = \colim_{([n]/\Delta)_{\leq m}^{opp}} U(n)
$$
描述；而对 $\varphi : [n] \to [n']$，映射 $i_{m!}U(\varphi)$
由上面的等同 $U(n) \circ \underline{\varphi} = U(n')$，通过
Categories，引理 \ref{categories-lemma-functorial-colimit} 得到。
\end{lemma}

\begin{proof}
在本证明中，以 $i_{m!}U$ 表示其第 $n$ 项由引理中显示公式给出的
单纯对象。我们将证明它满足伴随性质。

\medskip\noindent
设 $V$ 为 $\mathcal{C}$ 的单纯对象，并给定
$\gamma : U \to \text{sk}_mV$。态射
$$
\colim_{([n]/\Delta)_{\leq m}^{opp}} U(n) \to T
$$
由一个相容的态射系统 $f_\alpha : U_k \to T$ 给出，其中
$\alpha : [n] \to [k]$ 且 $k \leq m$。当然，取复合
$$
U_k \xrightarrow{\gamma_k} V_k \xrightarrow{V(\alpha)} V_n.
$$
便得到这样的态射系统。因此得到诱导态射
$(i_{m!}U)_n \to V_n$。留给读者验证，它们组成单纯对象之间的态射
$\gamma' : i_{m!}U \to V$。

\medskip\noindent
反过来，给定态射 $\gamma' : i_{m!}U \to V$，对
$0 \leq i \leq n$，令 $\gamma_i : U_i \to V_i$ 等于复合
$$
U_i
\xrightarrow{\text{id}_{[i]}}
\colim_{([i]/\Delta)_{\leq m}^{opp}} U(i)
\xrightarrow{\gamma'_i}
V_i
$$
% P02-E0118: frozen source bounds i by n here, but the reconstructed morphism has one component for every i <= m.
从而得到态射 $\gamma : U \to \text{sk}_m V$。留给读者验证，
这是上述构造的逆。
\end{proof}

\begin{lemma}
\label{simplicial-lemma-recovering-U}
设 $\mathcal{C}$ 为范畴，$U$ 为 $\mathcal{C}$ 的 $m$-截断单纯对象。
对任意 $n \leq m$，余极限
$$
\colim_{([n]/\Delta)_{\leq m}^{opp}} U(n)
$$
存在，并等于 $U_n$。
\end{lemma}

\begin{proof}
这是因为范畴 $([n]/\Delta)_{\leq m}$ 有初对象，即
$\text{id} : [n] \to [n]$。
\end{proof}

\begin{lemma}
\label{simplicial-lemma-recovering-U-for-real}
设 $\mathcal{C}$ 为有有限余极限的范畴，$U$ 为 $\mathcal{C}$ 的
$m$-截断单纯对象。映射 $U \to \text{sk}_m i_{m!}U$ 是同构。
\end{lemma}

\begin{proof}
合并引理 \ref{simplicial-lemma-left-adjoint-exists} 与引理 \ref{simplicial-lemma-recovering-U}。
\end{proof}

\begin{lemma}
\label{simplicial-lemma-imshriek-sets}
若 $U$ 是 $m$-截断单纯集且 $n > m$，则 $i_{m!}U$ 的所有
$n$-单形都是退化的。
\end{lemma}

\begin{proof}
这可以从引理 \ref{simplicial-lemma-left-adjoint-exists} 中 $i_{m!}U$ 的构造看出，
也可如下直接论证。令 $V = i_{m!}U$。设 $V' \subset V$ 为如下
单纯子集：当 $i \leq m$ 时 $V'_i = V_i$，而当 $i > m$ 时所有
$i$-单形都退化；见引理 \ref{simplicial-lemma-simplicial-set-n-skel-sub}。
由伴随公式，又因为 $\text{sk}_m V' = U$，单射 $V' \to V$ 有逆。
因此 $V' = V$。
\end{proof}

\begin{lemma}
\label{simplicial-lemma-n-skeleton-sets}
设 $U$ 为单纯集，$n \geq 0$ 为整数。态射
$i_{n!} \text{sk}_n U \to U$ 把 $i_{n!} \text{sk}_n U$ 与引理
\ref{simplicial-lemma-simplicial-set-n-skel-sub} 中定义的单纯集
$U' \subset U$ 等同起来。
\end{lemma}

\begin{proof}
由引理 \ref{simplicial-lemma-imshriek-sets}，$i_{n!} \text{sk}_n U$ 的非退化单形
只出现在次数 $\leq n$。映射 $i_{n!} \text{sk}_n U \to U$
在次数 $\leq n$ 上是同构。综合这两点，得到映射
$i_{n!} \text{sk}_n U \to U$ 把非退化单形映到非退化单形，
并且任意两个非退化单形的像都不相同。因此可应用引理
\ref{simplicial-lemma-injective-map-simplicial-sets}。于是
$i_{n!} \text{sk}_n U \to U$ 是单射。结论容易由此推出。
\end{proof}

\begin{remark}
\label{simplicial-remark-sk-literature}
在一些文献中，复合函子
$$
\text{Simp}(\mathcal{C})
\xrightarrow{\text{sk}_m}
\text{Simp}_m(\mathcal{C})
\xrightarrow{i_{m!}}
\text{Simp}(\mathcal{C})
$$
记作 $\text{sk}_m$。对单纯集而言这是合理的，因为此时引理
\ref{simplicial-lemma-n-skeleton-sets} 表明，$i_{m!} \text{sk}_m V$ 正是
$V$ 的如下子单纯集：它由 $V$ 的所有 $i$-单形（$i \leq m$）
及其退化组成。在这些文献中，通常还把复合
$$
\text{Simp}(\mathcal{C})
\xrightarrow{\text{sk}_m}
\text{Simp}_m(\mathcal{C})
\xrightarrow{\text{cosk}_m}
\text{Simp}(\mathcal{C})
$$
记作 $\text{cosk}_m$。
\end{remark}

\begin{lemma}
\label{simplicial-lemma-glue-simplex}
设 $U \subset V$ 为单纯集。假设 $n \geq 0$，且
$x \in V_n$、$x \not \in U_n$ 满足：
\begin{enumerate}
\item 当 $i < n$ 时，$V_i = U_i$；
\item $V_n = U_n \cup \{x\}$；
\item 当 $j > n$ 时，任意 $z \in V_j$、$z \not \in U_j$ 都退化。
\end{enumerate}
设 $\Delta[n] \to V$ 为把 $\Delta[n]$ 的非退化 $n$-单形映到
$x$ 的唯一态射。此时图表
$$
\xymatrix{
\Delta[n] \ar[r] & V \\
i_{(n - 1)!} \text{sk}_{n - 1} \Delta[n] \ar[r] \ar[u] & U \ar[u]
}
$$
是推出图表。
\end{lemma}

\begin{proof}
为方便起见，记
$\partial \Delta[n] = i_{(n - 1)!} \text{sk}_{n - 1} \Delta[n]$。
有自然映射 $U \amalg_{\partial \Delta[n]} \Delta[n] \to V$。
必须证明它在每个次数 $j$ 上都是双射，并且这对所有 $j$ 都成立。当 $j \leq n$ 时显然。
设 $j > n$。第三个条件意味着，任意 $z \in V_j$、
$z \not \in U_j$ 都是退化单形，例如
$z = s^{j - 1}_i(z')$。当然，$z' \not \in U_{j - 1}$。
由归纳可知，$z'$ 是 $x$ 的一个退化。因此，$V$ 的所有 $j$-单形
或者属于 $U$，或者是 $x$ 的退化。这蕴含映射
$U \amalg_{\partial \Delta[n]} \Delta[n] \to V$ 是满射。
注意，$U \amalg_{\partial \Delta[n]} \Delta[n]$ 的非退化单形，
或者是 $U$ 的某个非退化单形的像，或者是 $\Delta[n]$ 的（唯一）
非退化 $n$-单形的像。由于 $x$ 显然非退化，可推出
$U \amalg_{\partial \Delta[n]} \Delta[n] \to V$ 把非退化单形映到
非退化单形，并在非退化单形上是单射。因此，由引理
\ref{simplicial-lemma-injective-map-simplicial-sets}，它是单射。
\end{proof}

\begin{lemma}
\label{simplicial-lemma-add-simplices}
设 $U \subset V$ 为单纯集，并且对所有 $n$，$U_n, V_n$ 都有限非空。
假设 $U$ 和 $V$ 只有有限多个非退化单形。则存在子单纯集序列
$$
U = W^0 \subset W^1 \subset W^2 \subset \ldots W^r = V
$$
% P02-E0119: frozen source omits a subset relation between the ellipsis and W^r.
使得引理 \ref{simplicial-lemma-glue-simplex} 可应用于每个包含
$W^i \subset W^{i + 1}$。
\end{lemma}

\begin{proof}
设 $n$ 为使 $V$ 有一个不属于 $U$ 的非退化单形的最小整数。
取这样的非退化单形 $x \in V_n$、$x\not \in U_n$。
设 $W \subset V$ 为如下元素的集合：这些元素或者属于 $U$，
或者是 $x$ 的（反复）退化（换言之，具有形式 $V(\varphi)(x)$，
其中 $\varphi : [m] \to [n]$ 是满射）。容易看出，$W$ 是单纯集。
包含 $U \subset W$ 满足引理 \ref{simplicial-lemma-glue-simplex} 的条件。
此外，$V$ 中不属于 $W$ 的非退化单形数，恰比 $V$ 中不属于 $U$
的非退化单形数少一。因此，对这个数归纳即可。
\end{proof}

\begin{lemma}
\label{simplicial-lemma-imshriek-abelian}
设 $\mathcal{A}$ 为阿贝尔范畴
% P02-E0120: frozen source lacks sentence punctuation after “abelian category”.
设 $U$ 为 $\mathcal{A}$ 的 $m$-截断单纯对象。对 $n > m$，有
$N(i_{m!}U)_n = 0$。
\end{lemma}

\begin{proof}
令 $V = i_{m!}U$。设 $V' \subset V$ 为 $V$ 的如下单纯子对象：
当 $i \leq m$ 时 $V'_i = V_i$，而当 $i > m$ 时
$N(V'_i) = 0$；见引理 \ref{simplicial-lemma-simplicial-abelian-n-skel-sub}。
由伴随公式，又因为 $\text{sk}_m V' = U$，单射 $V' \to V$ 有逆。
因此 $V' = V$。
\end{proof}

\begin{lemma}
\label{simplicial-lemma-n-skeleton-abelian}
设 $\mathcal{A}$ 为阿贝尔范畴，$U$ 为 $\mathcal{A}$ 的单纯对象，
$n \geq 0$ 为整数。态射 $i_{n!} \text{sk}_n U \to U$
把 $i_{n!} \text{sk}_n U$ 与引理
\ref{simplicial-lemma-simplicial-abelian-n-skel-sub} 中定义的单纯子对象
$U' \subset U$ 等同起来。
\end{lemma}

\begin{proof}
由引理 \ref{simplicial-lemma-imshriek-abelian}，当 $i > n$ 时，有
$N(i_{n!} \text{sk}_n U)_i = 0$。映射
$i_{n!} \text{sk}_n U \to U$ 在次数 $\leq n$ 上是同构，
见引理 \ref{simplicial-lemma-recovering-U-for-real}。综合这两点，得到映射
$i_{n!} \text{sk}_n U \to U$ 对所有 $i$ 都诱导单射
$N(i_{n!} \text{sk}_n U)_i \to N(U)_i$。
因此可应用引理 \ref{simplicial-lemma-injective-map-simplicial-abelian}。
于是 $i_{n!} \text{sk}_n U \to U$ 是单射。结论容易由此推出。
\end{proof}

\noindent
下面利用上述材料给出理解余骨架函子的另一种方式。

\begin{lemma}
\label{simplicial-lemma-cosk-shriek}
设 $\mathcal{C}$ 为有有限余积和有限极限的范畴，$V$ 为
$\mathcal{C}$ 的单纯对象。此时
$$
(\text{cosk}_n \text{sk}_n V)_{n + 1}
=
\Hom(i_{n !}\text{sk}_n \Delta[n + 1], V)_0.
$$
\end{lemma}

\begin{proof}
由引理 \ref{simplicial-lemma-morphism-from-coproduct}，左边的对象表示把 $X$
送到下列等式中第一个集合的函子：
\begin{eqnarray*}
\Mor(X \times \Delta[n + 1], \text{cosk}_n \text{sk}_n V)
& = &
\Mor(X \times \text{sk}_n \Delta[n + 1], \text{sk}_n V) \\
& = &
\Mor(X \times i_{n !} \text{sk}_n \Delta[n + 1], V).
\end{eqnarray*}
引理公式右边的对象由把 $X$ 送到该等式链中最后一个集合的函子表示。
这证明了结论。

\medskip\noindent
在该等式链中，我们使用了
$\text{sk}_n (X \times \Delta[n + 1]) = X \times \text{sk}_n \Delta[n + 1]$
以及
$i_{n!}(X \times \text{sk}_n \Delta[n + 1]) =
X \times i_{n !} \text{sk}_n \Delta[n + 1]$。
第一个等式显然。对 $\mathcal{C}$ 的任意（可能截断的）单纯对象 $W$
和 $\mathcal{C}$ 的任意对象 $X$，暂以 $\Mor_\mathcal{C}(X, W)$ 表示（可能截断的）
单纯集 $[n] \mapsto \Mor_\mathcal{C}(X, W_n)$。
由定义，对任意（可能截断的）单纯集 $U$，有
$\Mor(U \times X, W) =
\Mor(U, \Mor_\mathcal{C}(X, W))$。因此
\begin{eqnarray*}
\Mor(X \times i_{n !} \text{sk}_n \Delta[n + 1], W)
& = &
\Mor(i_{n !} \text{sk}_n \Delta[n + 1], \Mor_\mathcal{C}(X, W)) \\
& = &
\Mor(\text{sk}_n \Delta[n + 1],
\text{sk}_n\Mor_\mathcal{C}(X, W)) \\
& = &
\Mor(X \times \text{sk}_n \Delta[n + 1], \text{sk}_nW) \\
& = &
\Mor(i_{n!}(X \times \text{sk}_n \Delta[n + 1]), W).
\end{eqnarray*}
这证明了所用的第二个等式，也结束了引理的证明。
\end{proof}

\section{阿贝尔范畴中的单纯对象}
\label{simplicial-section-abelian}

\noindent
回顾阿贝尔范畴的定义见《同调代数》第
\ref{homology-section-abelian-categories} 节。

\begin{lemma}
\label{simplicial-lemma-abelian}
设 $\mathcal{A}$ 为阿贝尔范畴。
\begin{enumerate}
\item 范畴 $\text{Simp}(\mathcal{A})$ 和
$\text{CoSimp}(\mathcal{A})$ 都是阿贝尔范畴。
\item （余）单纯对象的态射 $f : A \to B$ 为单射，当且仅当每个
$f_n : A_n \to B_n$ 都为单射。
\item （余）单纯对象的态射 $f : A \to B$ 为满射，当且仅当每个
$f_n : A_n \to B_n$ 都为满射。
\item （余）单纯对象的序列
$$
A \xrightarrow{f} B \xrightarrow{g} C
$$
在 $B$ 处正合，当且仅当每个序列
$$
A_i \xrightarrow{f_i} B_i \xrightarrow{g_i} C_i
$$
都在 $B_i$ 处正合。
\end{enumerate}
\end{lemma}

\begin{proof}
预加性很容易验证。终对象由所有次数中的 $U_n = 0$ 给出。
直积的存在性已见引理 \ref{simplicial-lemma-product} 和
\ref{simplicial-lemma-product-cosimplicial-objects}。
核与余核可逐项取核与余核得到。
\end{proof}

\noindent
对 $\mathcal{A}$ 的对象 $A$ 和整数 $k$，考虑满足下列条件的
$k$-截断单纯对象 $U$：
\begin{enumerate}
\item 当 $i < k$ 时，$U_i = 0$；
\item $U_k = A$；
\item 所有态射 $U(\varphi)$ 都等于零，唯有
$U(\text{id}_{[k]}) = \text{id}_A$。
\end{enumerate}
由于 $\mathcal{A}$ 既有有限极限也有有限余极限，
$\text{cosk}_k U$ 和 $i_{k!}U$ 都存在。
我们将描述二者以及典范映射 $i_{k!}U \to \text{cosk}_kU$。

\begin{lemma}
\label{simplicial-lemma-eilenberg-maclane-object}
设 $A$、$k$ 和 $U$ 如上，即当 $i < k$ 时 $U_i = 0$，且 $U_k = A$。
\begin{enumerate}
\item 给定 $k$-截断单纯对象 $V$，有
$$
\Mor(U, V)
=
\{ f : A \to V_k \mid d^k_i \circ f = 0, \ i = 0, \ldots, k \}
$$
以及
$$
\Mor(V, U)
=
\{ f : V_k \to A \mid f \circ s^{k - 1}_i = 0, \ i = 0, \ldots, k - 1 \}.
$$
\item 对象 $i_{k!} U$ 的第 $n$ 项等于
$\bigoplus_\alpha A$，其中 $\alpha$ 遍历所有满射
$\alpha : [n] \to [k]$。
\item 对任意 $\varphi : [m] \to [n]$，映射
$i_{k!} U(\varphi)$ 可描述为映射
$\bigoplus_\alpha A \to \bigoplus_{\alpha'} A$：
若 $\alpha \circ \varphi$ 不是满射，则它把与
$\alpha : [n] \to [k]$ 对应的分量映为零；若其为满射，
则以恒等映射把该分量映到与 $\alpha \circ \varphi$ 对应的分量。
\item 对象 $\text{cosk}_k U$ 的第 $n$ 项等于
$\bigoplus_\beta A$，其中 $\beta$ 遍历所有单射
$\beta : [k] \to [n]$。
\item 对任意 $\varphi : [m] \to [n]$，映射
$\text{cosk}_k U(\varphi)$ 可描述为映射
$\bigoplus_\beta A \to \bigoplus_{\beta'} A$：
若 $\beta$ 不经由 $\varphi$ 分解，则它把与
$\beta : [k] \to [n]$ 对应的分量映为零；若能够分解，
则以恒等映射把该分量映到每个满足
$\beta = \varphi \circ \beta'$ 的 $\beta'$ 所对应的分量。
\item 典范映射
$
c : i_{k !} U \to \text{cosk}_k U
$
在次数 $n$ 上的 $(\alpha, \beta)$ 系数 $A \to A$：
若 $\alpha \circ \beta$ 不是恒等映射，则等于零；
若其为恒等映射，则等于 $\text{id}_A$。
\item 典范映射
$
c : i_{k !} U \to \text{cosk}_k U
$
为单射。
\end{enumerate}
\end{lemma}

\begin{proof}
(1) 的证明留给读者。

\medskip\noindent
我们把 (2) 和 (3) 的规则作为一个单纯对象的定义，记为
$\tilde U$。下面证明它是 $i_{k!}U$ 的一种实现。
这将同时证明 (2) 和 (3)。必须证明：给定态射
$f : U \to \text{sk}_kV$，存在唯一态射
$\tilde f : \tilde U \to V$，使其取 $k$-骨架后恢复为 $f$。
由 (1)，$f$ 对应于态射 $f_k : A \to V_k$，且它对所有 $i$
都落在 $d^k_i$ 的核中。对任意满射
$\alpha : [n] \to [k]$，令
$\tilde f_\alpha : A \to V_n$ 等于复合
$\tilde f_\alpha = V(\alpha) \circ f_k : A \to V_n$。
定义 $\tilde f_n : \tilde U_n \to V_n$ 为所有满射
$\alpha : [n] \to [k]$ 所对应的 $\tilde f_\alpha$ 之和。
这样的一族 $\tilde f_\alpha$ 定义单纯对象的态射，当且仅当
对任意 $\varphi : [m] \to [n]$，图
$$
\xymatrix{
\bigoplus_{\alpha : [n] \to [k]\text{ 满射}} A
\ar[r]_-{\tilde f_n}
\ar[d]_{(3)} &
V_n \ar[d]^{V(\varphi)} \\
\bigoplus_{\alpha' : [m] \to [k]\text{ 满射}} A
\ar[r]^-{\tilde f_m} &
V_m
}
$$
交换。取 $\varphi = \alpha$ 可知，对 $\tilde f_\alpha$ 的选择由
$f_k$ 唯一决定。一般情形的交换性可以分别在左上角的每个直和分量上
检验。对于满足 $\alpha \circ \varphi$ 为满射的 $\alpha$ 所对应的
分量，这是显然的，因为这些分量通过 $\text{id}_A$ 映到
$\alpha' = \alpha \circ \varphi$ 所对应的分量，并且
$\tilde f_{\alpha'} = V(\alpha') \circ f_k =
V(\alpha \circ \varphi) \circ f_k = V(\varphi) \circ \tilde f_\alpha$。
对于 $\alpha \circ \varphi$ 不是满射的情形，必须证明
$V(\varphi) \circ \tilde f_\alpha = 0$。按定义，它等于
$V(\varphi) \circ V(\alpha) \circ f_k =
V(\alpha \circ \varphi) \circ f_k$。
由于 $\alpha \circ \varphi$ 不是满射，可以将其写成
$\delta^k_i \circ \psi$，于是由上可得
$V(\varphi) \circ V(\alpha) \circ f_k =
V(\psi) \circ d^k_i \circ f_k = 0$。

\medskip\noindent
我们把 (4) 和 (5) 的规则作为一个单纯对象的定义，也记为
$\tilde U$。下面证明它是 $\text{cosk}_k U$ 的一种实现。
这将同时证明 (4) 和 (5)。论证与上面 (2)、(3) 的证明完全对偶，
但我们仍将其写出。必须证明：给定态射
$f : \text{sk}_kV \to U$，存在唯一态射
$\tilde f : V \to \tilde U$，使其取 $k$-骨架后恢复为 $f$。
由 (1)，$f$ 对应于态射 $f_k : V_k \to A$，且对所有 $i$
都在 $s^{k - 1}_i$ 的像上为零。对任意单射 $\beta : [k] \to [n]$，
令 $\tilde f_\beta : V_n \to A$ 等于复合
$\tilde f_\beta = f_k \circ V(\beta) : V_n \to A$。
定义 $\tilde f_n : V_n \to \tilde U_n$ 为所有单射
$\beta : [k] \to [n]$ 所对应的 $\tilde f_\beta$ 之和。
这样的一族 $\tilde f_\beta$ 定义单纯对象的态射，当且仅当
对任意 $\varphi : [m] \to [n]$，图
$$
\xymatrix{
V_n
\ar[d]_{V(\varphi)}
\ar[r]_-{\tilde f_n}
&
\bigoplus_{\beta : [k] \to [n]\text{ 单射}} A
\ar[d]^{(5)}
\\
V_m
\ar[r]^-{\tilde f_m}
&
\bigoplus_{\beta' : [k] \to [m]\text{ 单射}} A
}
$$
交换。取 $\varphi = \beta$ 可知，对 $\tilde f_\beta$ 的选择由
$f_k$ 唯一决定。一般情形的交换性可以分别在右下角的每个直和分量上
检验。对于满足 $\varphi \circ \beta'$ 为单射的 $\beta'$ 所对应的
分量，这是显然的，因为恰有 $\beta = \varphi \circ \beta'$ 所对应的
分量映入其中，并且在这种情形下
$\tilde f_{\beta'} \circ V(\varphi) =
f_k \circ V(\beta') \circ V(\varphi) =
f_k \circ V(\beta) = \tilde f_\beta$。
对于 $\varphi \circ \beta'$ 不是单射的情形，必须证明
$\tilde f_{\beta'} \circ V(\varphi) = 0$。按定义，它等于
$f_k \circ V(\beta') \circ V(\varphi) =
f_k \circ V(\varphi \circ \beta')$。
由于 $\varphi \circ \beta'$ 不是单射，可以将其写成
$\psi \circ \sigma^{k - 1}_i$，于是由上可得
$f_k \circ V(\beta') \circ V(\varphi) =
f_k \circ s^{k - 1}_i \circ V(\psi) = 0$。

\medskip\noindent
复合 $i_{k!}U \to \text{cosk}_kU$ 是唯一的单纯对象态射，
并且在 $A = U_k = (i_{k!}U)_k = (\text{cosk}_kU)_k$ 上为恒等映射。
因此，只需检验所给规则定义了单纯对象的态射。为此，必须证明
对任意 $\varphi : [m] \to [n]$，图
$$
\xymatrix{
\bigoplus_{\alpha : [n] \to [k]\text{ 满射}} A
\ar[d]_{(3)}
\ar[r]_{(6)}
&
\bigoplus_{\beta : [k] \to [n]\text{ 单射}} A
\ar[d]^{(5)}
\\
\bigoplus_{\alpha' : [m] \to [k]\text{ 满射}} A
\ar[r]^{(6)}
&
\bigoplus_{\beta' : [k] \to [m]\text{ 单射}} A
}
$$
交换。现在可按如下方式把它理解为仅以 $0$ 和 $1$ 填充的矩阵：
(3) 的矩阵具有非零的 $(\alpha', \alpha)$ 项，当且仅当
$\alpha' = \alpha \circ \varphi$。同样，(5) 的矩阵具有非零的
$(\beta', \beta)$ 项，当且仅当
$\beta = \varphi \circ \beta'$。(6) 的上方矩阵具有非零的
$(\alpha, \beta)$ 项，当且仅当
$\alpha \circ \beta = \text{id}_{[k]}$；(6) 的下方矩阵也一样。
于是，该图的交换性归结为分别计算两个复合的
$(\alpha, \beta')$ 项并看出它们相等。这又归结为等式
$$
\# \left\{
\beta
\mid
\beta = \varphi \circ \beta'
\text{ 且 }
\alpha \circ \beta = \text{id}_{[k]}
\right\}
=
\# \left\{
\alpha'
\mid
\alpha' = \alpha \circ \varphi
\text{ 且 }
\alpha' \circ \beta' = \text{id}_{[k]}
\right\}
$$
其证明可以放心地留给读者。

\medskip\noindent
最后证明 (7)。它直接由引理
\ref{simplicial-lemma-injective-map-simplicial-abelian}、
\ref{simplicial-lemma-recover-cosk}、\ref{simplicial-lemma-recovering-U-for-real}
和 \ref{simplicial-lemma-imshriek-abelian} 得出。
\end{proof}

\begin{definition}
\label{simplicial-definition-eilenberg-maclane}
设 $\mathcal{A}$ 为阿贝尔范畴，$A$ 为 $\mathcal{A}$ 的对象，
且 $k$ 为整数并满足 $\geq 0$。{\it Eilenberg--Mac Lane 对象 $K(A, k)$}
是对象 $K(A, k) = i_{k!}U$，它已在上面的引理
\ref{simplicial-lemma-eilenberg-maclane-object} 中描述。
\end{definition}


\begin{lemma}
\label{simplicial-lemma-extension}
设 $\mathcal{A}$ 为阿贝尔范畴，$A$ 为 $\mathcal{A}$ 的对象，
且 $k$ 为整数并满足 $\geq 0$。考虑由下列规则定义的单纯对象 $E$：
\begin{enumerate}
\item $E_n = \bigoplus_\alpha A$，其中直和遍历所有
$\alpha : [n] \to [k + 1]$，其像为 $[k]$ 或 $[k + 1]$；
\item 给定 $\varphi : [m] \to [n]$，若
$\Im(\alpha \circ \varphi)$ 等于 $[k]$ 或 $[k + 1]$，
则映射 $E_n \to E_m$ 通过 $\text{id}_A$ 把
$\alpha$ 所对应的直和分量映到
$\alpha \circ \varphi$ 所对应的直和分量。
\end{enumerate}
则存在逐项分裂正合的短正合序列
$$
0 \to K(A, k) \to E \to K(A, k + 1) \to 0
$$
\end{lemma}

\begin{proof}
映射 $K(A, k)_n \to E_n$ 以及 $E_n \to K(A, k + 1)_n$
分别由直和的嵌入与投影给出，这从指标集的包含关系是显然的。
显然，它们是单纯对象的态射。
\end{proof}

\begin{lemma}
\label{simplicial-lemma-abelian-limit-skeleta}
设 $\mathcal{A}$ 为阿贝尔范畴。对 $\mathcal{A}$ 的任意单纯对象
$V$，有
$$
V = \colim_n i_{n!}\text{sk}_n V
$$
且所有转移映射都是单射。
\end{lemma}

\begin{proof}
这是因为一旦 $n \geq m$，每个 $V_m$ 就等于
$(i_{n!}\text{sk}_n V)_m$。关于转移映射，另见引理
\ref{simplicial-lemma-n-skeleton-abelian}。
\end{proof}

\section{单纯对象与链复形}
\label{simplicial-section-complexes}

\noindent
设 $\mathcal{A}$ 为阿贝尔范畴。关于链复形的约定和记号，
见《同调代数》第 \ref{homology-section-complexes} 节。
设 $U$ 为 $\mathcal{A}$ 的单纯对象。$U$ 的
{\it 伴随链复形 $s(U)$}，有时称为 {\it Moore 复形}，是链复形
$$
\ldots \to U_2 \to U_1 \to U_0 \to 0 \to 0 \to \ldots
$$
其边界映射 $d_n : U_n \to U_{n - 1}$ 由公式
$$
d_n = \sum\nolimits_{i = 0}^n (-1)^i d^n_i.
$$
给出。它确为复形，因为根据备注 \ref{simplicial-remark-relations} 中列出的关系，
有
\begin{eqnarray*}
d_n \circ d_{n + 1}
& = &
(\sum\nolimits_{i = 0}^n (-1)^i d^n_i) \circ
(\sum\nolimits_{j = 0}^{n + 1} (-1)^j d^{n + 1}_j) \\
& = &
\sum\nolimits_{0 \leq i < j \leq n + 1}
(-1)^{i + j} d^n_{j - 1} \circ d^{n + 1}_i
+
\sum\nolimits_{n \geq i \geq j \geq 0}
(-1)^{i + j} d^n_i \circ d^{n + 1}_j \\
& = & 0.
\end{eqnarray*}
各符号相互抵消！我们用 $s(U)$ 表示这个伴随链复形。
显然，该构造具有函子性，因此定义函子
$$
s : \text{Simp}(\mathcal{A}) \longrightarrow \text{Ch}_{\geq 0}(\mathcal{A}).
$$
于是得到这个看似令人困惑但正确的公式 $s(U)_n = U_n$。

\begin{lemma}
\label{simplicial-lemma-s-exact}
函子 $s$ 是正合的。
\end{lemma}

\begin{proof}
由引理 \ref{simplicial-lemma-abelian} 立即得出。
\end{proof}

\begin{lemma}
\label{simplicial-lemma-homology-extension}
设 $\mathcal{A}$ 为阿贝尔范畴，$A$ 为 $\mathcal{A}$ 的对象，
$k$ 为整数。设 $E$ 为引理 \ref{simplicial-lemma-extension} 中描述的对象。
则复形 $s(E)$ 是零调的。
\end{lemma}

\begin{proof}
对态射 $\alpha : [n] \to [k + 1]$，定义
$\alpha' : [n + 1] \to [k + 1]$ 为满足
$\alpha'|_{[n]} = \alpha$ 及 $\alpha'(n + 1) = k + 1$ 的映射。
注意，如果 $\alpha$ 的像为 $[k]$ 或 $[k + 1]$，
那么 $\alpha'$ 的像为 $[k + 1]$。考虑一族映射
$h_n : E_n \to E_{n + 1}$，它通过 $A$ 上的恒等映射，
把 $\alpha$ 所对应的分量映到 $\alpha'$ 所对应的分量。
我们来计算 $d_{n + 1} \circ h_n - h_{n - 1} \circ d_n$。
先在 $\mathcal{A}$ 为阿贝尔群范畴的情形下计算。
用记号 $x_\alpha$ 表示 $E_n$ 中由指标集内的映射 $\alpha$
所对应的直和分量里的元素 $x \in A$。我们还约定：
只要 $\Im(\alpha)$ 不是 $[k]$ 或 $[k + 1]$，
$x_\alpha$ 就表示 $E_n$ 的零元。由这些约定可见
$$
d_{n + 1}(h_n(x_\alpha)) =
\sum\nolimits_{i = 0}^{n + 1} (-1)^i x_{\alpha' \circ \delta^{n + 1}_i}
$$
以及
$$
h_{n - 1}(d_n(x_\alpha)) =
\sum\nolimits_{i = 0}^n (-1)^i x_{(\alpha \circ \delta_i^n)'}
$$
容易看出，当 $i = 0, \ldots, n$ 时，
$\alpha' \circ \delta^{n + 1}_i = (\alpha \circ \delta_i^n)'$；
也容易看出 $\alpha' \circ \delta^{n + 1}_{n + 1} = \alpha$。
因此
$$
(d_{n + 1} \circ h_n - h_{n - 1} \circ d_n)(x_\alpha)
=
(-1)^{n + 1} x_\alpha
$$
若 $\mathcal{A}$ 为任意阿贝尔范畴，这些恒等式仍然成立，
因为它们在单纯阿贝尔群 $[n] \mapsto \Hom(A, E_n)$ 中成立；
细节留给读者。由此，$E$ 上的恒等映射同伦于零，
同伦由一族映射
$h'_n = (-1)^{n + 1}h_n : E_n \to E_{n + 1}$ 给出。
因而 $E$ 是零调的；例如可由《同调代数》引理
\ref{homology-lemma-map-homology-homotopy} 得出。
\end{proof}

\begin{lemma}
\label{simplicial-lemma-homology-eilenberg-maclane}
设 $\mathcal{A}$ 为阿贝尔范畴，$A$ 为 $\mathcal{A}$ 的对象，
$k$ 为整数。当 $i = k$ 时 $H_i(s(K(A, k))) = A$，
其他情形为 $0$。
\end{lemma}

\begin{proof}
先证明 $k = 0$ 的情形。此时，对所有 $n$ 都有
$K(A, 0)_n = A$。此外，这个单纯阿贝尔群中的所有映射都是
$\text{id}_A$；换言之，$K(A, 0)$ 是取值为 $A$ 的常值单纯对象。
边界映射
$d_n = \sum_{i = 0}^n (-1)^i \text{id}_A = 0$
（当 $n$ 为奇数），而当 $n$ 为偶数时，则 $= \text{id}_A$。
因此 $s(K(A, 0))$ 形如
$$
\ldots \to A \xrightarrow{0} A \xrightarrow{1} A \xrightarrow{0} A \to 0
$$
结论显然。

\medskip\noindent
下面对 $k$ 归纳证明一般结论。假设 $k$ 的结论成立，考虑引理
\ref{simplicial-lemma-extension} 中的短正合序列
$$
0 \to K(A, k) \to E \to K(A, k + 1) \to 0
$$
由引理 \ref{simplicial-lemma-abelian}，相应的链复形序列正合。
由引理 \ref{simplicial-lemma-homology-extension}，$s(E)$ 是零调的。
于是 $k + 1$ 的结论由同调长正合序列得出；见《同调代数》引理
\ref{homology-lemma-long-exact-sequence-chain}。
\end{proof}

\noindent
还可以把第二个链复形关联给 $\mathcal{A}$ 的单纯对象。
回顾引理 \ref{simplicial-lemma-splitting-abelian-category}：
$\mathcal{A}$ 的任意单纯对象 $U$ 都典范分裂，并且
$N(U_m) = \bigcap_{i = 0}^{m - 1} \Ker(d^m_i)$。
定义 {\it 规范化链复形 $N(U)$} 为链复形
$$
\ldots \to N(U_2) \to N(U_1) \to N(U_0) \to 0 \to 0 \to \ldots
$$
其边界映射 $d_n : N(U_n) \to N(U_{n - 1})$
是 $(-1)^nd^n_n$ 在 $U_n$ 的直和分量 $N(U_n)$ 上的限制。
注意，引理 \ref{simplicial-lemma-N-d-in-N} 蕴含
$d^n_n(N(U_n)) \subset N(U_{n - 1})$。
它确为复形，因为
$d^n_n \circ d^{n + 1}_{n + 1} = d^n_n \circ d^{n + 1}_n$，
而按定义，$d^{n + 1}_n$ 在 $N(U_{n + 1})$ 上为零。
这样得到第二个函子
$$
N : \text{Simp}(\mathcal{A}) \longrightarrow \text{Ch}_{\geq 0}(\mathcal{A}).
$$
以下说明微分中符号的由来。

\begin{lemma}
\label{simplicial-lemma-map-associated-complexes}
设 $\mathcal{A}$ 为阿贝尔范畴，$U$ 为 $\mathcal{A}$ 的单纯对象。
典范映射 $N(U_n) \to U_n$ 给出复形态射 $N(U) \to s(U)$。
\end{lemma}

\begin{proof}
这是显然的，因为 $s(U)_n = U_n$ 上的微分是
$\sum (-1)^i d^n_i$，且当 $i < n$ 时，映射 $d^n_i$
在 $N(U_n)$ 上为零；而按定义，$(-1)^nd^n_n$ 的限制正是
$N(U)$ 的边界映射。
\end{proof}

\begin{lemma}
\label{simplicial-lemma-N-K}
设 $\mathcal{A}$ 为阿贝尔范畴，$A$ 为 $\mathcal{A}$ 的对象，
$k$ 为整数。当 $i = k$ 时，$N(K(A, k))_i = A$；
其他情形为 $0$。
\end{lemma}

\begin{proof}
当 $i < k$ 时，$N(K(A, k))_i = 0$ 是显然的，因为此时
$K(A, k)_i = 0$。又因 $K(A, k)_{k - 1} = 0$ 且
$K(A, k)_k = A$，显然有 $N(K(A, k))_k = A$。
当 $i > k$ 时，由引理 \ref{simplicial-lemma-imshriek-abelian} 及
$K(A, k)$ 的定义可得 $N(K(A, k))_i = 0$；见定义
\ref{simplicial-definition-eilenberg-maclane}。
\end{proof}

\begin{lemma}
\label{simplicial-lemma-decompose-associated-complexes}
设 $\mathcal{A}$ 为阿贝尔范畴，$U$ 为 $\mathcal{A}$ 的单纯对象。
链复形的典范态射 $N(U) \to s(U)$ 是分裂的。事实上，对某个复形
$D(U)$，有
$$
s(U) = N(U) \oplus D(U)
$$
构造 $U \mapsto D(U)$ 具有函子性。
\end{lemma}

\begin{proof}
定义 $D(U)_n$ 为映射
$$
\bigoplus\nolimits_{\varphi : [n] \to [m]\text{ 满射}, \ m < n} N(U_m)
\xrightarrow{\bigoplus U(\varphi)}
U_n
$$
的像。根据引理 \ref{simplicial-lemma-splitting-abelian-category} 和定义
\ref{simplicial-definition-split}，它是 $U_n$ 中与 $N(U_n)$ 互补的子对象。
下面证明 $D(U)$ 是子复形。取满射
$\varphi : [n] \to [m]$，其中 $m < n$，并考虑复合
$$
N(U_m) \xrightarrow{U(\varphi)} U_n \xrightarrow{d_n} U_{n - 1}.
$$
这个复合是映射
$$
N(U_m) \xrightarrow{U(\varphi \circ \delta^n_i)} U_{n - 1}
$$
之和，其符号为 $(-1)^i$，其中 $i = 0, \ldots, n$。

\medskip\noindent
先对 $m$ 作升序归纳，证明当 $0 \leq m < n - 1$ 时，
所有映射 $U(\varphi \circ \delta^n_i)$ 都把 $N(U_m)$ 映入
$D(U)_{n - 1}$。（情形 $m = n - 1$ 在下面处理。）
只要映射 $\varphi \circ \delta^n_i : [n - 1] \to [m]$
为满射，根据定义，$N(U_m)$ 在
$U(\varphi \circ \delta^n_i)$ 下的像就包含于 $D(U)_{n - 1}$。
若 $\varphi \circ \delta^n_i : [n - 1] \to [m]$ 不是满射，
令 $j = \varphi(i)$，并注意 $i$ 是在 $\varphi$ 下映为 $j$
的唯一指标。对某个 $\psi : [n - 1] \to [m - 1]$，可写成
$\varphi \circ \delta^n_i = \delta^m_j \circ \psi \circ \delta^n_i$。
% P02-E0121: the frozen factorization and its induced U-map contain an extra delta_i^n and are not type-correct.
因此
$U(\varphi \circ \delta^n_i) = U(\psi \circ \delta^n_i) \circ d^m_j$，
它在 $N(U_m)$ 上为零，除非 $j = m$。若 $j = m$，则
$d^m_m(N(U_m)) \subset N(U_{m - 1})$，从而
$U(\varphi \circ \delta^n_i)(N(U_m)) \subset
U(\psi \circ \delta^n_i)(N(U_{m - 1}))$，由归纳假设即得结论。

\medskip\noindent
为了完成 $D(U)$ 是子复形的证明，还须在 $m = n - 1$ 的情形下
处理复合
$$
N(U_m) \xrightarrow{U(\varphi)} U_n \xrightarrow{d_n} U_{n - 1}.
$$
。
% P02-E0122: frozen prose places a full stop between the displayed composition and the following “in case” clause.
此时，对某个 $0 \leq j \leq n - 1$，有
$\varphi = \sigma^{n - 1}_j$，且 $U(\varphi) = s^{n - 1}_j$。
因此，该复合由下列和给出：
$$
\sum (-1)^i d^n_i \circ s^{n - 1}_j
$$
回顾备注 \ref{simplicial-remark-relations}，
$d^n_j \circ s^{n - 1}_j = d^n_{j + 1} \circ s^{n - 1}_j = \text{id}$，
而相应两项符号相反，故相互抵消。若 $j < n - 1$，则映射
$d^n_n \circ s^{n - 1}_j$ 等于
$s^{n - 2}_j \circ d^{n - 1}_{n - 1}$。由于
$d^{n - 1}_{n - 1}$ 把 $N(U_{n - 1})$ 映入 $N(U_{n - 2})$，
可见像 $d^n_n ( s^{n - 1}_j (N(U_{n - 1}))$
% P02-E0123: frozen source leaves the outer parenthesis in this image expression unclosed.
包含于 $s^{n - 2}_j(N(U_{n - 2}))$，而根据定义，后者包含于
$D(U_{n - 1})$。
% P02-E0124: frozen source applies D to the component U_{n-1}; the established subcomplex term is D(U)_{n-1}.
对于所有其他的 $(i, j)$ 组合，或者
$d^n_i \circ s^{n - 1}_j = s^{n - 2}_{j - 1} \circ d^{n - 1}_i$
（若 $i < j$），或者
$d^n_i \circ s^{n - 1}_j = s^{n - 2}_j \circ d^{n - 1}_{i - 1}$
（若 $n > i > j + 1$）；在这些情形下，由 $N(U_{n - 1})$
的定义，该映射为零。
\end{proof}

\begin{remark}
\label{simplicial-remark-degenerate-subcomplex}
在引理 \ref{simplicial-lemma-decompose-associated-complexes} 的情形下，
子复形 $D(U) \subset s(U)$ 也可以定义为项
$$
D(U)_n = \Im\left(
\bigoplus\nolimits_{\varphi : [n] \to [m]\text{ 满射}, \ m < n} U_m
\xrightarrow{\bigoplus U(\varphi)}
U_n\right)
$$
所组成的子复形。事实上，$U_m$ 是子对象 $N(U_m)$
与所有满射 $[m] \to [k]$（其中 $k < m$）所给出的
$N(U_k)$ 之像的直和；因此这显然等同于引理
\ref{simplicial-lemma-decompose-associated-complexes} 的证明中给出的
$D(U)_n$ 的定义。由此可见，若 $U$ 是单纯阿贝尔群，
则 $D(U)_n$ 的元素恰为退化 $n$-单形之和。
\end{remark}

\begin{lemma}
\label{simplicial-lemma-N-exact}
函子 $N$ 是正合的。
\end{lemma}

\begin{proof}
由引理 \ref{simplicial-lemma-s-exact} 及引理
\ref{simplicial-lemma-decompose-associated-complexes} 的函子性分解即得。
\end{proof}

\begin{lemma}
\label{simplicial-lemma-quasi-isomorphism}
设 $\mathcal{A}$ 为阿贝尔范畴，$V$ 为 $\mathcal{A}$ 的单纯对象。
链复形的典范态射 $N(V) \to s(V)$ 是拟同构。换言之，引理
\ref{simplicial-lemma-decompose-associated-complexes} 中的复形 $D(V)$ 是零调的。
\end{lemma}

\begin{proof}
由引理 \ref{simplicial-lemma-homology-eilenberg-maclane} 和
\ref{simplicial-lemma-N-K}，对任意对象 $A$ 和任意 $k \geq 0$，
结论对 $K(A, k)$ 成立。考虑假设 $IH_{n, m}$：
对所有满足 $j \leq m$ 时 $V_j = 0$ 的 $V$，并对所有
$i \leq n$，映射 $N(V) \to s(V)$ 诱导同构
$H_i(N(V)) \to H_i(s(V))$。

\medskip\noindent
作为归纳的起点，$IH_{n, n}$ 显然成立，因为在这种情形下
$N(V)_n = 0$ 且 $s(V)_n = 0$。
% P02-E0125: frozen English says “To start of the induction”.

\medskip\noindent
假设 $IH_{n, m}$，其中 $m \leq n$。取满足当 $j < m$
时 $V_j = 0$ 的单纯对象 $V$。由引理
\ref{simplicial-lemma-eilenberg-maclane-object} 和定义
\ref{simplicial-definition-eilenberg-maclane}，有
$K(V_m, m) = i_{m!} \text{sk}_mV$。由引理
\ref{simplicial-lemma-n-skeleton-abelian}，自然态射
$$
K(V_m, m) = i_{m!} \text{sk}_mV \to V
$$
为单射。因此，对某个满足当 $i = 0, \ldots, m$ 时
$W_i = 0$ 的 $W$，得到短正合序列
$$
0 \to K(V_m, m) \to V \to W \to 0
$$
这个短正合序列诱导伴随复形的短正合序列之间的态射
% P02-E0126: frozen English has singular “a morphism of short exact sequence”.
$$
\xymatrix{
0 \ar[r] &
N(K(V_m, m)) \ar[r] \ar[d] &
N(V) \ar[r] \ar[d] &
N(W) \ar[r] \ar[d] &
0 \\
0 \ar[r] &
s(K(V_m, m)) \ar[r] &
s(V) \ar[r] &
s(W) \ar[r] &
0
}
$$
；见引理 \ref{simplicial-lemma-s-exact} 和 \ref{simplicial-lemma-N-exact}。
因此，由两端的结论可推出 $V$ 的结论。
\end{proof}






\section{Dold--Kan}
\label{simplicial-section-dold-kan}

\noindent
本节证明 Dold--Kan 定理，它把阿贝尔范畴中的单纯对象与链复形联系起来。

\begin{lemma}
\label{simplicial-lemma-N-faithful}
设 $\mathcal{A}$ 为阿贝尔范畴。函子 $N$ 是忠实的，并且反映同构、
单射和满射。
\end{lemma}

\begin{proof}
忠实性立即由引理 \ref{simplicial-lemma-splitting-abelian-category} 的典范分裂
得出。反映单射、满射和同构的断言由引理
\ref{simplicial-lemma-injective-map-simplicial-abelian} 得出。
\end{proof}

\begin{lemma}
\label{simplicial-lemma-S-N}
设 $\mathcal{A}$ 和 $\mathcal{B}$ 为阿贝尔范畴，并设
$N : \mathcal{A} \to \mathcal{B}$ 与
$S : \mathcal{B} \to \mathcal{A}$ 为函子。假设：
\begin{enumerate}
\item 函子 $S$ 和 $N$ 都正合；
\item 存在到 $\mathcal{B}$ 上恒等函子的同构
$g : N \circ S \to \text{id}_\mathcal{B}$；
\item $N$ 忠实；
\item $S$ 本质满。
\end{enumerate}
则 $S$ 和 $N$ 是互为拟逆的范畴等价。
\end{lemma}

\begin{proof}
只需构造函子性同构 $S(N(A)) \cong A$。为此，选取 $B$
和同构 $f : A  \to S(B)$。考虑映射
$$
f^{-1} \circ g_{S(B)} \circ S(N(f)) :
S(N(A)) \to S(N(S(B))) \to S(B) \to A.
$$
容易证明它不依赖于对 $f, B$ 的选择，并给出所需同构
$S \circ N \to \text{id}_\mathcal{A}$。
\end{proof}

\begin{theorem}
\label{simplicial-theorem-dold-kan}
设 $\mathcal{A}$ 为阿贝尔范畴。函子 $N$ 诱导范畴等价
$$
N :
\text{Simp}(\mathcal{A})
\longrightarrow
\text{Ch}_{\geq 0}(\mathcal{A})
$$
\end{theorem}

\begin{proof}
我们将描述一个反方向的函子，它受引理 \ref{simplicial-lemma-extension}
中的构造启发（但为了得到正确的边界映射，要加入一个符号）。
设 $A_\bullet$ 为边界映射
$d_{A, n} : A_n \to A_{n - 1}$ 的链复形。对每个 $n \geq 0$，记
$$
I_n =
\Big\{
\alpha : [n] \to \{0, 1, 2, \ldots\}
\mid
\Im(\alpha) = [k]\text{ 对某个 }k
\Big\}.
$$
对 $\alpha \in I_n$，以 $k(\alpha)$ 表示满足
$\Im(\alpha) = [k]$ 的唯一整数。如下定义单纯对象
$S(A_\bullet)$：
\begin{enumerate}
\item $S(A_\bullet)_n = \bigoplus_{\alpha \in I_n} A_{k(\alpha)}$，
并将其写成
$\bigoplus_{\alpha \in I_n} A_{k(\alpha)} \cdot \alpha$，
以提示把“$\alpha$”看作相应直和分量的基向量；
\item 给定 $\varphi : [m] \to [n]$，按如下规则，在
$S(A_\bullet)_n$ 的直和分量
$A_{k(\alpha)} \cdot \alpha$ 上定义 $S(A_\bullet)(\varphi)$：
\begin{enumerate}
\item 若 $\alpha \circ \varphi \not \in I_m$，则令它等于零；
% P02-E0127: frozen English omits “If” at the start of this case.
\item 若 $\alpha \circ \varphi \in I_m$，但
$k(\alpha \circ \varphi)$ 既不等于 $k(\alpha)$，也不等于
$k(\alpha) - 1$，则也令它等于零；
\item 若 $\alpha \circ \varphi \in I_m$ 且
$k(\alpha \circ \varphi) = k(\alpha)$，则通过恒等映射映到
$S(A_\bullet)_m$ 的分量
$A_{k(\alpha \circ \varphi)} \cdot (\alpha \circ \varphi)$；
\item 若 $\alpha \circ \varphi \in I_m$ 且
$k(\alpha \circ \varphi) = k(\alpha) - 1$，则使用
$(-1)^{k(\alpha)} d_{A, k(\alpha)}$ 映到
$S(A_\bullet)_m$ 的分量
$A_{k(\alpha \circ \varphi)}\cdot (\alpha \circ \varphi)$。
\end{enumerate}
\end{enumerate}
下面证明 $S(A_\bullet)$ 是 $\mathcal{A}$ 的单纯对象。
为此，设有映射 $\varphi : [m] \to [n]$ 和
$\psi : [n] \to [p]$。我们将证明
$S(A_\bullet)(\varphi) \circ S(A_\bullet)(\psi) =
S(A_\bullet)(\psi \circ \varphi)$。选取 $\beta \in I_p$，并令
$\alpha = \beta \circ \psi$ 和 $\gamma = \alpha \circ \varphi$，
将它们分别视为映射
$\alpha : [n] \to \{0, 1, 2, \ldots\}$ 和
$\gamma : [m] \to \{0, 1, 2, \ldots\}$。图示为
$$
\xymatrix{
[m] \ar[r]_\varphi \ar[d]_\gamma &
[n] \ar[r]_\psi \ar[d]_\alpha &
[p] \ar[d]^\beta \\
\Im(\gamma) \ar[r] &
\Im(\alpha) \ar[r] &
[k(\beta)]
}
$$
我们将证明映射
$S(A_\bullet)(\varphi) \circ S(A_\bullet)(\psi)$ 和
$S(A_\bullet)(\psi \circ \varphi)$
在分量 $A_{k(\beta)} \cdot \beta$ 上的限制相同。
% P02-E0128: frozen source breaks this sentence with a full stop before “to the summand”.
需要考虑若干情形：
\begin{enumerate}
\item 设 $\alpha \not \in I_n$，因而 $S(A_\bullet)(\psi)$
在 $A_{k(\beta)} \cdot \beta$ 上的限制为零。
此时，或者 $\gamma \not \in I_m$；或者
$[k(\gamma)] = \Im(\gamma) \subset \Im(\alpha) \subset [k(\beta)]$，
且 $[k(\beta)]$ 的子集 $\Im(\alpha)$ 有缺口，故
$k(\gamma) < k(\beta) - 1$。在两种情形下，
$S(A_\bullet)(\psi \circ \varphi)$ 在
$A_{k(\beta)} \cdot \beta$ 上的限制也为零。
\item 设 $\alpha \in I_n$ 且 $k(\alpha) < k(\beta) - 1$，
因而 $S(A_\bullet)(\psi)$ 在 $A_{k(\beta)} \cdot \beta$
上的限制为零。此时，或者 $\gamma \not \in I_m$；或者
$[k(\gamma)] \subset [k(\alpha)] \subset [k(\beta)]$，
从而 $k(\gamma) < k(\beta) - 1$。在两种情形下，
$S(A_\bullet)(\psi \circ \varphi)$ 在
$A_{k(\beta)} \cdot \beta$ 上的限制也为零。
\item 设 $\alpha \in I_n$ 且 $k(\alpha) = k(\beta)$，
因而 $S(A_\bullet)(\psi)$ 在 $A_{k(\beta)} \cdot \beta$
上的限制是从 $A_{k(\beta)} \cdot \beta$ 到
$A_{k(\alpha)} \cdot \alpha$ 的恒等映射。在这种情形下，
由于 $\Im(\alpha) =[k(\beta)]$，
$S(A_\bullet)(\psi \circ \varphi)$ 在分量
$A_{k(\beta)} \cdot \beta$ 上的限制规则，恰与
$S(A_\bullet)(\varphi)$ 在分量
$A_{k(\alpha)} \cdot \alpha$ 上的限制规则相同，故二者相合。
\item 设 $\alpha \in I_n$ 且 $k(\alpha) = k(\beta) - 1$，
因而 $S(A_\bullet)(\psi)$ 在 $A_{k(\beta)} \cdot \beta$
上的限制由到 $A_{k(\alpha)} \cdot \alpha$ 的
$(-1)^{k(\beta)}d_{A, k(\beta)}$ 给出。分以下子情形：
\begin{enumerate}
\item 若 $\gamma \not \in I_m$，则
$S(A_\bullet)(\psi \circ \varphi)$ 在分量
$A_{k(\beta)} \cdot \beta$ 上的限制，以及
$S(A_\bullet)(\varphi)$ 在分量
$A_{k(\alpha)} \cdot \alpha$ 上的限制都为零，故二者相合。
\item 若 $\gamma \in I_m$ 但 $k(\gamma) < k(\alpha) - 1$，
则两个限制同样都为零，故二者相合。
\item 若 $\gamma \in I_m$ 且 $k(\gamma) = k(\alpha)$，
则 $\Im(\gamma) = \Im(\alpha)$。此时，
$S(A_\bullet)(\psi \circ \varphi)$ 在分量
$A_{k(\beta)} \cdot \beta$ 上的限制由到
$A_{k(\gamma)} \cdot \gamma$ 的
$(-1)^{k(\beta)}d_{A, k(\beta)}$ 给出，而
$S(A_\bullet)(\varphi)$ 在分量
$A_{k(\alpha)} \cdot \alpha$ 上的限制是恒等映射
$A_{k(\alpha)} \cdot \alpha \to A_{k(\gamma)} \cdot \gamma$。
故二者相合。
\item 最后，若 $\gamma \in I_m$ 且
$k(\gamma) = k(\alpha) - 1$，则
$S(A_\bullet)(\varphi)$ 在分量
$A_{k(\alpha)} \cdot \alpha$ 上的限制由
$(-1)^{k(\alpha)} d_{A, k(\alpha)}$ 给出，作为映射
$A_{k(\alpha)} \cdot \alpha \to A_{k(\beta)} \cdot \beta$。
% P02-E0129: the frozen codomain here is the beta summand, but S(phi) should land in the gamma summand.
由于 $A_\bullet$ 是复形，可见复合
$A_{k(\beta)} \cdot \beta \to
A_{k(\alpha)} \cdot \alpha \to A_{k(\gamma)} \cdot \gamma$
为零；这与 $S(A_\bullet)(\psi \circ \varphi)$ 在分量
$A_{k(\beta)} \cdot \beta$ 上的限制所得结果相合，因为
$k(\gamma) = k(\beta) - 2 < k(\beta) - 1$。
\end{enumerate}
\end{enumerate}
所以 $S(A_\bullet)$ 是 $\mathcal{A}$ 的单纯对象。

\medskip\noindent
下面构造关于 $A_\bullet$ 具有函子性的同构
$A_\bullet \to N(S(A_\bullet))$。回顾引理
\ref{simplicial-lemma-decompose-associated-complexes} 给出的链复形分解
$$
S(A_\bullet) =  N(S(A_\bullet)) \oplus D(S(A_\bullet))
$$
另一方面，由备注 \ref{simplicial-remark-degenerate-subcomplex} 和
$S(A_\bullet)$ 的构造可得
$$
D(S(A_\bullet))_n =
\bigoplus\nolimits_{\alpha \in I_n,\ k(\alpha) < n} A_{k(\alpha)} \cdot \alpha
\subset
\bigoplus\nolimits_{\alpha \in I_n} A_{k(\alpha)} \cdot \alpha
$$
然而，若 $\alpha \in I_n$，则
$k(\alpha) \geq n \Leftrightarrow \alpha = \text{id}_{[n]} : [n] \to [n]$。
因此，$S(A_\bullet)_n$ 的分量
$A_n \cdot \text{id}_{[n]}$ 是分量 $D(S(A_\bullet))_n$ 的补。
所有映射 $d^n_i : S(A_\bullet)_n \to S(A_\bullet)_n$
% P02-E0130: frozen source gives the face maps codomain degree n instead of n-1.
在分量 $A_n \cdot \text{id}_{[n]}$ 上的限制都为零，
唯有 $d^n_n$ 从 $A_n \cdot \text{id}_{[n]}$ 到
$A_{n - 1} \cdot \text{id}_{[n - 1]}$ 给出
$(-1)^n d_{A, n}$。由此断定，$A_n \cdot \text{id}_{[n]}$
必定等于分量 $N(S(A_\bullet))_n$；而且微分
$d_n = \sum (-1)^id^n_i : S(A_\bullet)_n \to S(A_\bullet)_{n - 1}$
在分量 $A_n \cdot \text{id}_{[n]}$ 上的限制正给出所需映射！

\medskip\noindent
最后，必须证明 $S \circ N$ 同构于恒等函子。设 $U$ 为
$\mathcal{A}$ 的单纯对象。则可定义显然的映射
$$
S(N(U))_n = \bigoplus\nolimits_{\alpha \in I_n} N(U)_{k(\alpha)} \cdot \alpha
\longrightarrow
U_n
$$
：在 $\alpha$ 所对应的分量上使用
$U(\alpha) : N(U)_{k(\alpha)} \to U_n$。由定义
\ref{simplicial-definition-split}，这是同构。为了完成证明，还须证明它与单纯对象
中的映射相容。为此，设 $\varphi : [m] \to [n]$ 且
$\alpha \in I_n$，并令 $\beta = \alpha \circ \varphi$。图示为
$$
\xymatrix{
[m] \ar[r]_\varphi \ar[d]_\beta &
[n] \ar[d]_\alpha \\
\Im(\beta) \ar[r] &
[k(\alpha)]
}
$$
需要考虑若干情形：
\begin{enumerate}
\item 设 $\beta \not \in I_m$。则存在指标
$0 \leq j < k(\alpha)$，满足
$j \not \in \Im(\alpha \circ \varphi)$，从而可以选取分解
$\alpha \circ \varphi = \delta^{k(\alpha)}_j \circ \psi$，
其中 $\psi : [m] \to [k(\alpha) - 1]$。由此可知，
$U(\varphi)$ 在分量 $N(U)_{k(\alpha)} \cdot \alpha$ 的像上为零，
因为 $U(\varphi) \circ U(\alpha) =
U(\alpha \circ \varphi) = U(\psi) \circ d^{k(\alpha)}_j$
根据 $N$ 的构造在 $N(U)_{k(\alpha)}$ 上为零。这与上面对
$S(N(U))$ 所定的规则一致。
\item 设 $\beta \in I_m$ 且 $k(\beta) < k(\alpha) - 1$。
此时完全按情形 (1) 论证，取 $j = k(\alpha) - 1$。
\item 设 $\beta \in I_m$ 且 $k(\beta) = k(\alpha)$。
此时，分量 $N(U)_{k(\alpha)} \cdot \alpha$ 通过恒等映射
映到分量 $N(U)_{k(\beta)} \cdot \beta$。这与
$U(\varphi)$ 的作用相同，因为在这种情形下
$U(\varphi) \circ U(\alpha) = U(\beta)$。
\item 设 $\beta \in I_m$ 且 $k(\beta) = k(\alpha) - 1$。
此时，使用微分 $(-1)^{k(\alpha)} d_{N(U), k(\alpha)}$
把分量 $N(U)_{k(\alpha)} \cdot \alpha$ 映到分量
$N(U)_{k(\beta)} \cdot \beta$。另一方面，因为此时
$\Im(\beta) = [k(\beta)]$，有
$\alpha \circ \varphi = \delta^{k(\alpha)}_{k(\alpha)} \circ \beta$。
因此，$U(\varphi)$ 与 $U(\alpha)$ 在 $N(U)_{k(\alpha)}$
上的限制之复合，等于 $U(\beta)$ 前复合
$d^{k(\alpha)}_{k(\alpha)}$ 在 $N(U)_{k(\alpha)}$ 上的限制。
由于 $d_{N(U), k(\alpha)} = \sum (-1)^i d^{k(\alpha)}_i$，
并且当 $i < k(\alpha)$ 时，$d^{k(\alpha)}_i$ 在
$N(U)_{k(\alpha)}$ 上的限制为零，可见二者相等。
\end{enumerate}
定理得证。
\end{proof}







\section{余单纯对象的 Dold--Kan 对应}
\label{simplicial-section-dold-kan-cosimplicial}

\noindent
设 $\mathcal{A}$ 为阿贝尔范畴。根据《同调代数》引理
\ref{homology-lemma-abelian-opposite}，$\mathcal{A}^{opp}$
也是阿贝尔范畴。由定义形式地推出
$$
\text{CoSimp}(\mathcal{A}) = \text{Simp}(\mathcal{A}^{opp})^{opp}.
$$
于是 Dold--Kan 对应（定理 \ref{simplicial-theorem-dold-kan}）蕴含
$\text{CoSimp}(\mathcal{A})$ 等价于范畴
$\text{Ch}_{\geq 0}(\mathcal{A}^{opp})^{opp}$。
由定义还形式地推出
$$
\text{CoCh}_{\geq 0}(\mathcal{A}) =
\text{Ch}_{\geq 0}(\mathcal{A}^{opp})^{opp}.
$$
将这些箭头合在一起，得到等价
$$
Q :
\text{CoSimp}(\mathcal{A})
\longrightarrow
\text{CoCh}_{\geq 0}(\mathcal{A}).
$$
本节描述 $Q$。

\medskip\noindent
先定义余单纯对象 $U$ 的
{\it 伴随上链复形 $s(U)$}。它是在负次数中为零，并且当
$n \geq 0$ 时满足 $s(U)^n = U_n$ 的上链复形。
作为微分，使用映射
$d^n : s(U)^n \to s(U)^{n + 1}$，其定义为
$d^n = \sum_{i = 0}^{n + 1} (-1)^i \delta^{n + 1}_i$。
换言之，复形 $s(U)$ 形如
$$
\xymatrix{
0 \ar[r] &
U_0 \ar[rr]^{\delta^1_0 - \delta^1_1} & &
U_1 \ar[rr]^{\delta^2_0 - \delta^2_1 + \delta^2_2} & &
U_2 \ar[r] &
\ldots
}
$$
它有时也称为 $U$ 的 {\it Moore 复形}。

\medskip\noindent
另一方面，给定 $\mathcal{A}$ 的余单纯对象 $U$，令
$Q(U)^0 = U_0$，并令
$$
Q(U)^n = \Coker(
\xymatrix{
\bigoplus_{i = 0}^{n - 1} U_{n - 1} \ar[r]^-{\delta^n_i} &
U_n
}).
$$
微分 $d^n : Q(U)^n \to Q(U)^{n + 1}$ 由
$(-1)^{n + 1}\delta^{n + 1}_{n + 1}$ 诱导；也就是说，把态射
$(-1)^{n + 1}\delta^{n + 1}_{n + 1}$ 放入交换图
$$
\xymatrix{
U_n \ar[rr]_{(-1)^{n + 1}\delta^{n + 1}_{n + 1}} \ar[d] & &
U_{n + 1} \ar[d] \\
Q(U)^n \ar[rr]^{d_n} & &
Q(U)^{n + 1}.
}
$$
% P02-E0131: the frozen diagram labels the induced cochain differential d_n although the surrounding definition calls it d^n.
留给读者证明这个图是有意义的；即当
$i = 0, \ldots, n - 1$ 时，$\delta^n_i$ 的像映入右侧竖直箭头的核。
（这与引理 \ref{simplicial-lemma-N-d-in-N} 对偶。）
因此，上链复形 $Q(U)$ 形如
$$
0 \to Q(U)^0 \to Q(U)^1 \to Q(U)^2 \to \ldots
$$
它称为 $U$ 的 {\it 伴随规范化上链复形}。
Dold--Kan 定理 \ref{simplicial-theorem-dold-kan} 的对偶如下。

\begin{lemma}
\label{simplicial-lemma-dual-dold-kan}
设 $\mathcal{A}$ 为阿贝尔范畴。
\begin{enumerate}
\item 函子
$s : \text{CoSimp}(\mathcal{A}) \to \text{CoCh}_{\geq 0}(\mathcal{A})$
是正合的。
\item 映射 $s(U)^n \to Q(U)^n$ 定义上链复形的态射。
\item 在 $\text{CoCh}_{\geq 0}(\mathcal{A})$ 中存在函子性直和分解
$s(U) = D(U) \oplus Q(U)$。
\item 函子 $Q$ 是正合的。
\item 复形态射 $s(U) \to Q(U)$ 是拟同构。
\item 函子 $U \mapsto Q(U)^\bullet$ 定义范畴等价
$\text{CoSimp}(\mathcal{A}) \to \text{CoCh}_{\geq 0}(\mathcal{A})$。
\end{enumerate}
\end{lemma}

\begin{proof}
略。但这些结果正好分别是引理
\ref{simplicial-lemma-s-exact}、\ref{simplicial-lemma-map-associated-complexes}、
\ref{simplicial-lemma-decompose-associated-complexes}、
\ref{simplicial-lemma-N-exact}、\ref{simplicial-lemma-quasi-isomorphism}
以及定理 \ref{simplicial-theorem-dold-kan} 的对偶陈述。
\end{proof}

\section{同伦}
\label{simplicial-section-homotopy}

\noindent
考虑单纯集 $\Delta[0]$ 和 $\Delta[1]$。回顾有两个态射
$$
e_0, e_1 : \Delta[0] \longrightarrow \Delta[1],
$$
它们来自把 $0$ 映到 $[1] = \{0, 1\}$ 的一个元素的态射
$[0] \to [1]$。还回顾每个集合 $\Delta[1]_k$ 都是有限的。
因此，若范畴 $\mathcal{C}$ 有有限余积，则对 $\mathcal{C}$
的任意单纯对象 $U$，都可以构成积
$$
U \times \Delta[1]
$$
；见定义 \ref{simplicial-definition-product-with-simplicial-set}。
注意，$\Delta[0]$ 的性质是：对所有 $k \geq 0$，
$\Delta[0]_k = \{*\}$ 都是单点集。因此
$U \times \Delta[0] = U$。于是，上面的 $e_0, e_1$ 给出态射
% P02-E0132: frozen English uses singular “gives rise” with the compound subject e_0, e_1.
$$
e_0, e_1 : U \to U \times \Delta[1].
$$

\begin{definition}
\label{simplicial-definition-homotopy}
设 $\mathcal{C}$ 为有有限余积的范畴，$U$ 和 $V$ 为
$\mathcal{C}$ 的两个单纯对象，并设 $a, b : U \to V$ 为两个态射。
\begin{enumerate}
\item 如果态射
$$
h : U \times \Delta[1] \longrightarrow V
$$
满足 $a = h \circ e_0$ 且 $b = h \circ e_1$，则称它为
{\it 从 $a$ 到 $b$ 的同伦}。
\item 如果存在一列从 $U$ 到 $V$ 的态射
$a = a_0, a_1, \ldots, a_n = b$，使得对每个
$i = 1, \ldots, n$，或者存在从 $a_{i - 1}$ 到 $a_i$ 的同伦，
或者存在从 $a_i$ 到 $a_{i - 1}$ 的同伦，则称态射 $a$ 和 $b$
{\it 同伦}，或者称它们 {\it 属于同一个同伦类}。
\end{enumerate}
\end{definition}

\noindent
一般而言，“存在从 $a$ 到 $b$ 的同伦”这一关系既不传递也不对称；
例 \ref{simplicial-example-trivial-homotopy} 将说明它是自反的。当然，“同伦”
是集合 $\Mor(U, V)$ 上的等价关系，并且是由“存在从 $a$ 到 $b$
的同伦”这一关系生成的等价关系。事实上，可以在任意范畴中定义
单纯对象态射对之间的同伦。在详细分析存在态射
$h : U \times \Delta[1] \to V$ 的含义之后，我们将在备注
\ref{simplicial-remark-homotopy-better} 中完成这一点。

\medskip\noindent
设 $\mathcal{C}$ 为有有限余积的范畴，$U$、$V$ 为
$\mathcal{C}$ 的单纯对象，$a, b : U \to V$ 为态射。
进一步假设 $h : U \times \Delta[1] \to V$ 是从 $a$ 到 $b$
的同伦。对每个 $n \geq 0$，记
$$
\Delta[1]_n = \{\alpha^n_0, \ldots, \alpha^n_{n + 1}\}
$$
其中 $\alpha^n_i : [n] \to [1]$ 是满足
$$
\alpha^n_i(j)
=
\left\{
\begin{matrix}
0 & \text{若} & j < i\\
1 & \text{若} & j \geq i
\end{matrix}
\right.
$$
的映射。因此
$$
h_n : (U \times \Delta[1])_n = \coprod U_n \cdot \alpha^n_i
\longrightarrow
V_n
$$
具有分量 $h_{n, i} : U_n \to V_n$；对每个
$i = 0, \ldots, n + 1$，它是在 $\alpha^n_i$ 所对应的分量上的限制。

\begin{lemma}
\label{simplicial-lemma-relations-homotopy}
在上述情形下，有如下关系：
\begin{enumerate}
\item $h_{n, 0} = b_n$ 且 $h_{n, n + 1} = a_n$。
\item 当 $i > j$ 时，
$d^n_j \circ h_{n, i} = h_{n - 1, i - 1} \circ d^n_j$。
\item 当 $i \leq j$ 时，
$d^n_j \circ h_{n, i} = h_{n - 1, i} \circ d^n_j$。
\item 当 $i > j$ 时，
$s^n_j \circ h_{n, i} = h_{n + 1, i + 1} \circ s^n_j$。
\item 当 $i \leq j$ 时，
$s^n_j \circ h_{n, i} = h_{n + 1, i} \circ s^n_j$。
\end{enumerate}
反之，给定满足上述性质的一族映射 $h_{n, i}$，它们定义态射
$h$，并且这是从 $a$ 到 $b$ 的同伦。
\end{lemma}

\begin{proof}
略。最后一个断言可利用如下事实证明：根据引理
\ref{simplicial-lemma-face-degeneracy-category}，给出单纯对象的态射，
等价于给出与所有 $d^n_j$ 和 $s^n_j$ 交换的一列态射 $h_n$。
\end{proof}

\begin{example}
\label{simplicial-example-trivial-homotopy}
在上述情形下，假设 $a = b$。则存在从 $a$ 到 $b$ 的
{\it 平凡}同伦，即满足 $h_{n, i} = a_n = b_n$ 的同伦。
\end{example}

\begin{remark}
\label{simplicial-remark-homotopy-better}
设 $\mathcal{C}$ 为任意范畴（不作任何假设），$U$ 和 $V$
为 $\mathcal{C}$ 的单纯对象，$a, b : U \to V$ 为
$\mathcal{C}$ 中单纯对象的态射。一个 {\it 从 $a$ 到 $b$ 的同伦}
由态射\footnote{文献中经常使用映射
$h_{n + 1, i} \circ s_i : U_n \to V_{n + 1}$，而不是映射
$h_{n, i}$。当然，这些映射所满足的关系不同于引理
\ref{simplicial-lemma-relations-homotopy} 中的关系。}
$h_{n, i} : U_n \to V_n$ 给出，其中 $n \geq 0$ 且
$i = 0, \ldots, n + 1$，并要求它们满足引理
\ref{simplicial-lemma-relations-homotopy} 中的关系。与定义
\ref{simplicial-definition-homotopy} 一样，如果存在一列从 $U$ 到 $V$ 的态射
$a = a_0, a_1, \ldots, a_n = b$，使得对每个
$i = 1, \ldots, n$，或者存在从 $a_{i - 1}$ 到 $a_i$ 的同伦，
或者存在从 $a_i$ 到 $a_{i - 1}$ 的同伦，则称态射 $a$ 和 $b$
{\it 同伦}。显然，若 $F : \mathcal{C} \to \mathcal{C}'$
为任意函子，且 $\{h_{n, i}\}$ 是从 $a$ 到 $b$ 的同伦，
则 $\{F(h_{n, i})\}$ 是从 $F(a)$ 到 $F(b)$ 的同伦。
同样，若 $a$ 和 $b$ 同伦，则 $F(a)$ 和 $F(b)$ 同伦。
由于该引理表明，当有限余积存在时，新概念与旧概念相同，
% P02-E0133: frozen English says singular “finite coproduct exist”.
特别可知，只要两个范畴都有有限余积，函子就保持原来的同伦概念。
\end{remark}

\begin{remark}
\label{simplicial-remark-homotopy-pre-post-compose}
设 $\mathcal{C}$ 为任意范畴。假设单纯对象的两个态射
$a, a' : U \to V$ 同伦。则对任意态射 $b : V \to W$，
两个映射 $b \circ a, b \circ a' : U \to W$ 同伦。
同样，对任意态射 $c : X \to U$，两个映射
$a \circ c, a' \circ c : X \to V$ 同伦。事实上，映射
$b \circ a \circ c, b \circ a' \circ c : X \to W$ 同伦。
确切地说，如果映射 $h_{n, i} : U_n \to V_n$ 定义从
$a$ 到 $a'$ 的同伦，那么映射 $b \circ h_{n, i} \circ c$
定义从 $b \circ a \circ c$ 到 $b \circ a' \circ c$ 的同伦。
由此得到一个新范畴 $\text{hSimp}(\mathcal{C})$：它与
$\text{Simp}(\mathcal{C})$ 有相同的对象，而其态射是
$\text{Simp}(\mathcal{C})$ 中态射的同伦类。因此有典范函子
$$
\text{Simp}(\mathcal{C})
\longrightarrow
\text{hSimp}(\mathcal{C})
$$
它本质满，并且在态射集上满射。
\end{remark}

\begin{definition}
\label{simplicial-definition-homotopy-equivalent}
设 $U$ 和 $V$ 为范畴 $\mathcal{C}$ 的两个单纯对象。
如果存在态射 $b : V \to U$，使得 $a \circ b$ 同伦于
$\text{id}_V$，且 $b \circ a$ 同伦于 $\text{id}_U$，
则称态射 $a : U \to V$ 为 {\it 同伦等价}。
如果存在同伦等价 $a : U \to V$，则称 $U$ 和 $V$
{\it 同伦等价}。
\end{definition}

\begin{example}
\label{simplicial-example-simplex-contractible}
单纯集 $\Delta[m]$ 同伦等价于 $\Delta[0]$。事实上，考虑唯一态射
$f : \Delta[m] \to \Delta[0]$，以及由 $\Delta[m]$ 的最后一个
$0$-单形的包含给出的态射 $g : \Delta[0] \to \Delta[m]$。
有 $f \circ g = \text{id}$。下面给出从
$\text{id}_{\Delta[m]}$ 到 $g \circ f$ 的同伦
$h : \Delta[m] \times \Delta[1] \to \Delta[m]$。确切地说，
$h$ 由映射
$$
\Mor_\Delta([n], [m]) \times \Mor_\Delta([n], [1]) \to \Mor_\Delta([n], [m])
$$
给出，它们把 $(\varphi, \alpha)$ 映到
$$
k \mapsto
\left\{
\begin{matrix}
\varphi(k) & \text{若} & \alpha(k) = 0 \\
m & \text{若} & \alpha(k) = 1
\end{matrix}
\right.
$$
。注意，这只因我们取 $g$ 为最后一个 $0$-单形的包含而成立。
若取 $g$ 为第一个 $0$-单形的包含，则可找到从
$g \circ f$ 到 $\text{id}_{\Delta[m]}$ 的同伦。
这说明了单纯集范畴中同伦所固有的不对称性。
\end{example}

\noindent
下面的引理说明 $U \times \Delta[1]$ 同伦等价于 $U$。

\begin{lemma}
\label{simplicial-lemma-contractible}
设 $\mathcal{C}$ 为有有限余积的范畴，$U$ 为 $\mathcal{C}$
的单纯对象。考虑映射
$e_1, e_0 : U \to U \times \Delta[1]$ 和
$\pi : U \times \Delta[1] \to U$；见引理
\ref{simplicial-lemma-back-to-U}。
\begin{enumerate}
\item $\pi \circ e_1 = \pi \circ e_0 = \text{id}_U$；
\item 态射 $\text{id}_{U \times \Delta[1]}$ 与
$e_0 \circ \pi$ 同伦；
\item 态射 $\text{id}_{U \times \Delta[1]}$ 与
$e_1 \circ \pi$ 同伦。
\end{enumerate}
\end{lemma}

\begin{proof}
第一个断言平凡。考虑单纯集映射
$\Delta[1] \times \Delta[1] \longrightarrow \Delta[1]$，
它在次数 $n$ 上把一对 $(\beta_1, \beta_2)$，
其中 $\beta_i : [n] \to [1]$，映到由规则
$$
\beta(i) = \max\{\beta_1(i), \beta_2(i)\}.
$$
定义的态射 $\beta : [n] \to [1]$。
这是单纯集的态射，因为作用
$\Delta[1](\varphi) : \Delta[1]_n \to \Delta[1]_m$
（其中 $\varphi : [m] \to [n]$）是前复合。使用上面引入的记号
$\Delta[1]_n = \{\alpha^n_0, \ldots, \alpha^n_{n + 1}\}$，
若 $\beta_2 = \alpha^n_0$，则 $\beta = \alpha^n_0$；
若 $\beta_2 = \alpha^n_{n + 1}$，则 $\beta = \beta_1$。
由于 $\alpha^n_0$ 与 $\alpha^n_{n + 1}$ 分别是取值为
$1$ 与 $0$ 的常值映射，这意味着引理
\ref{simplicial-lemma-back-to-U} 的诱导态射
$$
U \times \Delta[1] \times \Delta[1]
\longrightarrow
U \times \Delta[1]
$$
是从 $\text{id}_{U \times \Delta[1]}$ 到
$e_1 \circ \pi$ 的同伦。$e_0 \circ \pi$ 的情形类似
（使用最小值而不是最大值）。
\end{proof}

\begin{lemma}
\label{simplicial-lemma-fibre-products-simplicial-object-w-section}
设 $f : Y \to X$ 为有纤维积的范畴 $\mathcal{C}$ 中的态射，
并假设 $f$ 有截面 $s$。考虑例
\ref{simplicial-example-fibre-products-simplicial-object} 中从 $f$
构造的单纯对象 $U$。态射 $U \to U$ 在每个次数上都是
$Y \times_X \ldots \times_X Y$ 的自映射
$(s \circ f)^{n + 1}$，即在每个因子上由 $s \circ f$ 给出；
这个态射同伦于 $U$ 上的恒等映射。特别地，$U$ 同伦等价于
常值单纯对象 $X$。
\end{lemma}

\begin{proof}
令 $g^0 = \text{id}_Y$ 且 $g^1 = s \circ f$。使用态射
\begin{eqnarray*}
Y \times_X \ldots \times_X Y \times \Mor([n], [1])
& \to &
Y \times_X \ldots \times_X Y \\
(y_0, \ldots, y_n) \times \alpha
& \mapsto &
(g^{\alpha(0)}(y_0), \ldots, g^{\alpha(n)}(y_n))
\end{eqnarray*}
这里采用函子点观点来定义这些映射。
% P02-E0134: frozen English repeats “point” in “functor of points point of view”.
换言之，
$h_{n, 0} = \text{id}$、
$h_{n, n + 1} = (s \circ f)^{n + 1}$，且
$h_{n, i} = \text{id}_Y^{i + 1} \times (s \circ f)^{n + 1 - i}$。
留给读者证明它们满足引理 \ref{simplicial-lemma-relations-homotopy} 的关系，
因而定义所需同伦。另见备注 \ref{simplicial-remark-homotopy-better}；
它说明不必对范畴 $\mathcal{C}$ 作任何其他假设。
\end{proof}

\begin{lemma}
\label{simplicial-lemma-products-homotopy}
设 $\mathcal{C}$ 为范畴，$T$ 为集合。对 $t \in T$，设
$X_t$、$Y_t$ 为 $\mathcal{C}$ 的单纯对象，并假设
$X = \prod_{t \in T} X_t$ 和 $Y = \prod_{t \in T} Y_t$ 存在。
\begin{enumerate}
\item 若对所有 $t \in T$，$X_t$ 与 $Y_t$ 同伦等价，且 $T$
有限，则 $X$ 与 $Y$ 同伦等价。
\end{enumerate}
对 $t \in T$，设 $a_t, b_t : X_t \to Y_t$ 为态射，并令
$a = \prod a_t : X \to Y$ 及 $b = \prod b_t : X \to Y$。
\begin{enumerate}
\item[(2)] 若对所有 $t \in T$ 都存在从 $a_t$ 到 $b_t$ 的同伦，
则存在从 $a$ 到 $b$ 的同伦。
\item[(3)] 若 $T$ 有限，且对 $t \in T$，
$a_t, b_t : X_t \to Y_t$ 同伦，则 $a$ 与 $b$ 同伦。
\end{enumerate}
\end{lemma}

\begin{proof}
若 $h_t = (h_{t, n , i})$ 是从 $a_t$ 到 $b_t$ 的同伦
（见备注 \ref{simplicial-remark-homotopy-better}），则
$h = (\prod_t h_{t, n, i})$ 是从 $\prod a_t$ 到
$\prod b_t$ 的同伦。这证明 (2)。

\medskip\noindent
证明 (3)。选取 $t \in T$。存在整数 $n \geq 0$ 和链
$a_t = a_{t, 0}, a_{t, 1}, \ldots, a_{t, n} = b_t$，
使得对每个 $1 \leq i \leq n$，或者存在从
$a_{t, i - 1}$ 到 $a_{t, i}$ 的同伦，或者存在从
$a_{t, i}$ 到 $a_{t, i - 1}$ 的同伦。若 $n = 0$，
就另取一个 $t$。（若对所有 $t \in T$ 都有 $a_t = b_t$，
则已经完成。）所以假设 $n > 0$。由例
\ref{simplicial-example-trivial-homotopy}，对所有
$t' \in T \setminus \{t\}$，存在从 $b_{t'}$ 到 $b_{t'}$
的同伦。
% P02-E0135: frozen English says “there are is a homotopy”.
因此由 (2)，或者存在从
$a_{t, n - 1} \times \prod_{t'} b_{t'}$ 到 $b$ 的同伦，
或者存在从 $b$ 到
$a_{t, n - 1} \times \prod_{t'} b_{t'}$ 的同伦。
如此可将 $n$ 减少 $1$。这证明 (3)。

\medskip\noindent
(1) 由 (3) 和定义得出。
\end{proof}




\section{阿贝尔范畴中的同伦}
\label{simplicial-section-homotopy-abelian}

\noindent
设 $\mathcal{A}$ 为加性范畴，$U$、$V$ 为 $\mathcal{A}$
的单纯对象，$a, b : U \to V$ 为态射。进一步假设
$h : U \times \Delta[1] \to V$ 是从 $a$ 到 $b$ 的同伦。
下面证明链复形的两个态射
$s(a), s(b) : s(U) \longrightarrow s(V)$
在《同调代数》第 \ref{homology-section-complexes} 节的意义下同伦。
使用第 \ref{simplicial-section-homotopy} 节引入的记号，按公式
\begin{equation}
\label{simplicial-equation-homotopy-to-homotopy}
s(h)_n = \sum\nolimits_{i = 0}^n (-1)^{i + 1} h_{n + 1, i + 1} \circ s^n_i.
\end{equation}
定义
$$
s(h)_n : U_n \longrightarrow V_{n + 1}
$$
。计算 $d_{n + 1} \circ s(h)_n + s(h)_{n - 1} \circ d_n$。
首先，
\begin{eqnarray*}
d_{n + 1} \circ s(h)_n & = &
\sum\nolimits_{j = 0}^{n + 1}
\sum\nolimits_{i = 0}^n
(-1)^{j + i + 1}
d^{n + 1}_j \circ h_{n + 1, i + 1} \circ s^n_i \\
& = &
\sum\nolimits_{1 \leq i + 1 \leq j \leq n + 1}
(-1)^{j + i + 1}
h_{n, i + 1} \circ d^{n + 1}_j \circ s^n_i \\
& &
+
\sum\nolimits_{n \geq i \geq j \geq 0}
(-1)^{i + j + 1}
h_{n, i} \circ d^{n + 1}_j \circ s^n_i \\
& = &
\sum\nolimits_{1 \leq i + 1 < j \leq n + 1}
(-1)^{j + i + 1}
h_{n, i + 1} \circ s^{n - 1}_i \circ d^n_{j - 1} \\
& &
+
\sum\nolimits_{1 \leq i + 1 = j \leq n + 1}
(-1)^{j + i + 1}
h_{n, i + 1} \\
& &
+
\sum\nolimits_{n \geq i = j \geq 0}
(-1)^{i + j + 1}
h_{n, i} \\
& &
+
\sum\nolimits_{n \geq i > j \geq 0}
(-1)^{i + j + 1}
h_{n, i} \circ s^{n - 1}_{i - 1} \circ d^n_j
\end{eqnarray*}
留给读者验证：四个和式中的第一个与最后一个，恰好同
$$
s(h)_{n - 1} \circ d_n
=
\sum_{i = 0}^{n - 1} \sum_{j = 0}^n
(-1)^{i + 1 + j} h_{n, i + 1} \circ s^{n - 1}_i \circ d^n_j.
$$
的所有项相消。因此得到
\begin{eqnarray*}
d_{n + 1} \circ s(h)_n + s(h)_{n - 1} \circ d_n
& = &
\sum_{j = 1}^{n + 1} (-1)^{2j} h_{n, j}
+
\sum_{i = 0}^n (-1)^{2i + 1} h_{n, i} \\
& = &
h_{n, n + 1} - h_{n , 0} \\
& = &
a_n - b_n
\end{eqnarray*}
，正是所需结论。

\begin{lemma}
\label{simplicial-lemma-homotopy-s-N}
设 $\mathcal{A}$ 为加性范畴，$a, b : U \to V$ 为
$\mathcal{A}$ 中单纯对象的态射。若 $a$、$b$ 同伦，则
$s(a), s(b) : s(U) \to s(V)$ 是链复形的同伦态射。
若 $\mathcal{A}$ 为阿贝尔范畴，则
$N(a), N(b) : N(U) \to N(V)$ 也是链复形的同伦态射。
\end{lemma}

\begin{proof}
可以选取一列从 $U$ 到 $V$ 的态射
$a = a_0, a_1, \ldots, a_n = b$，使得对每个
$i = 1, \ldots, n$，或者存在从 $a_i$ 到 $a_{i - 1}$
的同伦，或者存在从 $a_{i - 1}$ 到 $a_i$ 的同伦。
上面的计算说明，在此情形下，或者 $s(a_i)$ 作为链复形映射
同伦于 $s(a_{i - 1})$，或者 $s(a_{i - 1})$ 作为链复形映射
同伦于 $s(a_i)$。当然，这两个断言等价；而且同伦是链复形映射集
上的等价关系，见《同调代数》第 \ref{homology-section-complexes} 节。
这证明 $s(a)$ 和 $s(b)$ 作为链复形映射同伦。

\medskip\noindent
下面处理 $N(a)$ 和 $N(b)$。由引理
\ref{simplicial-lemma-decompose-associated-complexes}，
$N(a)$、$N(b)$ 是复合
$$
N(U) \to s(U) \to s(V) \to N(V)
$$
，其中中间使用 $s(a)$、$s(b)$。因此，断言由《同调代数》引理
\ref{homology-lemma-compose-homotopy} 得出。
\end{proof}

\begin{lemma}
\label{simplicial-lemma-homotopy-equivalence-s-N}
设 $\mathcal{A}$ 为加性范畴，$a : U \to V$ 为
$\mathcal{A}$ 中单纯对象的态射。若 $a$ 是同伦等价，则
$s(a) : s(U) \to s(V)$ 是链复形的同伦等价。
若进一步 $\mathcal{A}$ 为阿贝尔范畴，则
$N(a) : N(U) \to N(V)$ 也是链复形的同伦等价。
\end{lemma}

\begin{proof}
略。见上面的引理 \ref{simplicial-lemma-homotopy-s-N}。
\end{proof}




\section{同伦与余单纯对象}
\label{simplicial-section-homotopy-cosimplicial}

\noindent
设 $\mathcal{C}$ 为有有限积的范畴，$V$ 为余单纯对象，并考虑
$\Hom(\Delta[1], V)$；见第
\ref{simplicial-section-hom-from-simplicial-sets-into-cosimplicial} 节。
态射 $e_0, e_1 : \Delta[0] \to \Delta[1]$ 给出两个态射
$e_0, e_1 : \Hom(\Delta[1], V) \to V$。

\begin{definition}
\label{simplicial-definition-homotopy-cosimplicial}
设 $\mathcal{C}$ 为有有限积的范畴，$U$ 和 $V$ 为
$\mathcal{C}$ 的两个余单纯对象，$a, b : U \to V$ 为
$\mathcal{C}$ 中余单纯对象的两个态射。
\begin{enumerate}
\item 若态射
$$
h : U \longrightarrow \Hom(\Delta[1], V)
$$
满足 $a = e_0 \circ h$ 且 $b = e_1 \circ h$，则称它为
{\it 从 $a$ 到 $b$ 的同伦}。
\item 如果存在一列从 $U$ 到 $V$ 的态射
$a = a_0, a_1, \ldots, a_n = b$，使得对每个
$i = 1, \ldots, n$，或者存在从 $a_i$ 到 $a_{i - 1}$ 的同伦，
或者存在从 $a_{i - 1}$ 到 $a_i$ 的同伦，则称 $a$ 和 $b$
{\it 同伦}，或者称它们 {\it 属于同一个同伦类}。
\end{enumerate}
\end{definition}

\noindent
这与第 \ref{simplicial-section-homotopy} 节为单纯对象引入的概念对偶。
为作说明，考虑定义中的从 $a$ 到 $b$ 的同伦
$h : U \to \Hom(\Delta[1], V)$。回顾 $\Delta[1]_n$
是有限集。$h$ 的次数 $n$ 分量是态射
$$
h_n = (h_{n, \alpha}) :
U
\longrightarrow
\Hom(\Delta[1], V)_n =
\prod\nolimits_{\alpha \in \Delta[1]_n} V_n
$$
% P02-E0136: frozen degree-n component has source U rather than U_n, contradicted immediately by h_{n,alpha}: U_n -> V_n.
$\mathcal{C}$ 的态射 $h_{n, \alpha} : U_n \to V_n$ 具有如下性质：
对 $\Delta$ 的每个态射 $f : [n] \to [m]$，有
\begin{equation}
\label{simplicial-equation-property-homotopy-cosimplicial}
h_{m, \alpha} \circ U(f) = V(f) \circ h_{n, \alpha \circ f}
\end{equation}
此外，条件 $a = e_0 \circ h$ 意味着
$a_n = h_{n, 0 : [n] \to [1]}$，其中
$0 : [n] \to [1]$ 是取值为 $0$ 的常值映射。
同样，条件 $b = e_1 \circ h$ 意味着
$b_n = h_{n, 1 : [n] \to [1]}$，其中
$1 : [n] \to [1]$ 是取值为 $1$ 的常值映射。
反之，给定一族态射 $\{h_{n, \alpha}\}$，使得
(\ref{simplicial-equation-property-homotopy-cosimplicial})
对 $\Delta$ 的所有态射 $f$ 都成立，并且对所有 $n \geq 0$ 都有
$a_n = h_{n, 0 : [n] \to [1]}$ 和
$b_n = h_{n, 1 : [n] \to [1]}$，那么令
$h = \prod_{\alpha \in \Delta[1]_n} h_{n, \alpha}$，
便得到从 $a$ 到 $b$ 的同伦 $h$。

\begin{remark}
\label{simplicial-remark-homotopy-cosimplicial-better}
设 $\mathcal{C}$ 为任意范畴（不作任何假设），$U$ 和 $V$
为 $\mathcal{C}$ 的余单纯对象，$a, b : U \to V$ 为
$\mathcal{C}$ 中余单纯对象的态射。一个 {\it 从 $a$ 到 $b$ 的同伦}
由态射 $h_{n, \alpha} : U_n \to V_n$ 给出，其中
$n \geq 0$ 且 $\alpha \in \Delta[1]_n$；要求它们对
$\Delta$ 的所有态射 $f$ 都满足
(\ref{simplicial-equation-property-homotopy-cosimplicial})，
并且对所有 $n \geq 0$ 都有
$a_n = h_{n, 0 : [n] \to [1]}$ 与
$b_n = h_{n, 1 : [n] \to [1]}$。
与定义 \ref{simplicial-definition-homotopy-cosimplicial} 一样，
如果存在一列从 $U$ 到 $V$ 的态射
$a = a_0, a_1, \ldots, a_n = b$，使得对每个
$i = 1, \ldots, n$，或者存在从 $a_{i - 1}$ 到 $a_i$ 的同伦，
或者存在从 $a_i$ 到 $a_{i - 1}$ 的同伦，则称态射 $a$ 和 $b$
{\it 同伦}。显然，若 $F : \mathcal{C} \to \mathcal{C}'$
为任意函子，且 $\{h_{n, i}\}$ 是从 $a$ 到 $b$ 的同伦，
% P02-E0137: frozen cosimplicial remark switches the family index from alpha to i.
则 $\{F(h_{n, i})\}$ 是从 $F(a)$ 到 $F(b)$ 的同伦。
同样，若 $a$ 和 $b$ 同伦，则 $F(a)$ 和 $F(b)$ 同伦。
当有限积存在时，这个新概念与旧概念相同。因此特别可知，
只要两个范畴都有有限积，函子就保持原来的概念。
\end{remark}

\begin{lemma}
\label{simplicial-lemma-compare-homotopies}
设 $\mathcal{C}$ 为范畴，$U$ 和 $V$ 为 $\mathcal{C}$
的两个余单纯对象，$a, b : U \to V$ 为余单纯对象的态射。
回顾 $U$、$V$ 分别对应于 $\mathcal{C}^{opp}$ 的单纯对象
$U'$、$V'$，而 $a, b$ 分别对应于态射
$a', b' : V' \to U'$。下列条件等价：
\begin{enumerate}
\item 存在备注 \ref{simplicial-remark-homotopy-cosimplicial-better} 意义下，
从 $a$ 到 $b$ 的同伦 $h = \{h_{n, \alpha}\}$。
\item 存在备注 \ref{simplicial-remark-homotopy-better} 意义下，
从 $a'$ 到 $b'$ 的同伦 $h = \{h_{n, i}\}$。
\end{enumerate}
因此，$a$ 与 $b$ 在备注
\ref{simplicial-remark-homotopy-cosimplicial-better} 的意义下同伦，
当且仅当 $a'$ 与 $b'$ 在备注 \ref{simplicial-remark-homotopy-better}
的意义下同伦。
\end{lemma}

\begin{proof}
若 $\mathcal{C}$ 有有限积，则 $\mathcal{C}^{opp}$ 有有限余积，
并且可以用定义 \ref{simplicial-definition-homotopy-cosimplicial}
和 \ref{simplicial-definition-homotopy} 代替备注
\ref{simplicial-remark-homotopy-cosimplicial-better}
和 \ref{simplicial-remark-homotopy-better}。
此时，$h : U \to \Hom(\Delta[1], V)$ 等同于态射
$h' : \Hom(\Delta[1], V)' \to U'$。由于积与余积也相互交换，
立即有 $(\Hom(\Delta[1], V))' = V' \times \Delta[1]$。
此外，态射 $e_0, e_1 : \Hom(\Delta[1], V) \to V$
的“加撇”版本是态射
$e_0, e_1 : V' \to \Delta[1] \times V$。
因此，$e_0 \circ h = a$ 转化为 $h' \circ e_0 = a'$，
同样地，$e_1 \circ h = b$ 转化为 $h' \circ e_1 = b'$。
这证明了这种情形的引理。

\medskip\noindent
一般情形下，需要把
(\ref{simplicial-equation-property-homotopy-cosimplicial})
给出的关系转化为引理 \ref{simplicial-lemma-relations-homotopy} 中的关系。
细节从略。

\medskip\noindent
最后一个断言形式地由 (1) 与 (2) 的等价性得出。
\end{proof}

\begin{lemma}
\label{simplicial-lemma-functorial-homotopy}
\begin{slogan}
函子保持（余）单纯对象的同伦态射。
\end{slogan}
设 $\mathcal{C}, \mathcal{C}', \mathcal{D}, \mathcal{D}'$ 为范畴。
术语如备注 \ref{simplicial-remark-homotopy-cosimplicial-better} 和
\ref{simplicial-remark-homotopy-better} 所述。
% P02-E0138: frozen source presents “With terminology as in ...” as a sentence fragment.
\begin{enumerate}
\item 设 $a, b : U \to V$ 为 $\mathcal{D}$ 中单纯对象的态射，
$F : \mathcal{D} \to \mathcal{D}'$ 为协变函子。若 $a$ 与 $b$
同伦，则 $F(a)$、$F(b)$ 是单纯对象之间同伦的态射
$F(U) \to F(V)$。
\item 设 $a, b : U \to V$ 为 $\mathcal{C}$ 中余单纯对象的态射，
$F : \mathcal{C} \to \mathcal{C}'$ 为协变函子。若 $a$ 与 $b$
同伦，则 $F(a)$、$F(b)$ 是余单纯对象之间同伦的态射
$F(U) \to F(V)$。
\item 设 $a, b : U \to V$ 为 $\mathcal{D}$ 中单纯对象的态射，
$F : \mathcal{D} \to \mathcal{C}$ 为反变函子。若 $a$ 与 $b$
同伦，则 $F(a)$、$F(b)$ 是余单纯对象之间同伦的态射
$F(V) \to F(U)$。
\item 设 $a, b : U \to V$ 为 $\mathcal{C}$ 中余单纯对象的态射，
$F : \mathcal{C} \to \mathcal{D}$ 为反变函子。若 $a$ 与 $b$
同伦，则 $F(a)$、$F(b)$ 是单纯对象之间同伦的态射
$F(V) \to F(U)$。
\end{enumerate}
\end{lemma}

\begin{proof}
由上面的引理 \ref{simplicial-lemma-compare-homotopies}，可以把 $F$
变为一对范畴之间的协变函子，而只须证明函子保持同伦的映射对。
这已在备注 \ref{simplicial-remark-homotopy-better} 中说明。
\end{proof}

\begin{lemma}
\label{simplicial-lemma-push-outs-simplicial-object-w-section}
设 $f : X \to Y$ 为有推出的范畴 $\mathcal{C}$ 中的态射，
并假设存在态射 $s : Y \to X$，满足
$s \circ f = \text{id}_X$。考虑例
\ref{simplicial-example-push-outs-simplicial-object} 中从 $f$ 构造的
余单纯对象 $U$。态射 $U \to U$ 在每个次数上都是
$Y \amalg_X \ldots \amalg_X Y$ 的自映射，并且在每个因子上
由 $f \circ s$ 给出；它同伦于 $U$ 上的恒等映射。
特别地，$U$ 同伦等价于常值余单纯对象 $X$。
\end{lemma}

\begin{proof}
本引理与引理
\ref{simplicial-lemma-fibre-products-simplicial-object-w-section} 对偶。
因此，它由引理 \ref{simplicial-lemma-compare-homotopies} 得出。
\end{proof}

\begin{lemma}
\label{simplicial-lemma-homotopy-s-Q}
\begin{slogan}
（余单纯）Dold--Kan 函子把同伦映射映为同伦映射。
\end{slogan}
设 $\mathcal{A}$ 为加性范畴，$a, b : U \to V$ 为
$\mathcal{A}$ 中余单纯对象的态射。若 $a$、$b$ 同伦，则
$s(a), s(b) : s(U) \to s(V)$ 是上链复形的同伦映射。
若进一步 $\mathcal{A}$ 为阿贝尔范畴，则
$Q(a), Q(b) : Q(U) \to Q(V)$ 是上链复形的同伦映射。
\end{lemma}

\begin{proof}
设 $(-)' : \mathcal{A} \to \mathcal{A}^{opp}$ 为反变函子
$A \mapsto A$。根据引理
\ref{simplicial-lemma-push-outs-simplicial-object-w-section}，
映射 $a'$ 和 $b'$ 同伦。
% P02-E0139: frozen proof cites the pushout-with-section lemma here, but homotopy transfer to opposite categories is established by lemma-compare-homotopies.
由引理 \ref{simplicial-lemma-homotopy-s-N}，$s(a')$ 和 $s(b')$
是链复形的同伦映射。由于 $s(a') = (s(a))'$ 且
$s(b') = (s(b))'$，应用加性反变函子
$(-)'' : \mathcal{A}^{opp} \to \mathcal{A}$ 可知
$s(a)$ 和 $s(b)$ 也同伦。对 $Q$-复形的结论以相同方式得出，
这里使用 $Q(U)' = N(U')$。
\end{proof}

\begin{lemma}
\label{simplicial-lemma-homotopy-equivalence-s-Q}
设 $\mathcal{A}$ 为加性范畴，$a : U \to V$ 为
$\mathcal{A}$ 中余单纯对象的态射。若 $a$ 是同伦等价，则
$s(a) : s(U) \to s(V)$ 是链复形的同伦等价。
若进一步 $\mathcal{A}$ 为阿贝尔范畴，则
$Q(a) : Q(U) \to Q(V)$ 也是链复形的同伦等价。
% P02-E0140: frozen source calls both targets chain complexes although s(U) and Q(U) were defined here as cochain complexes.
\end{lemma}

\begin{proof}
略。见上面的引理 \ref{simplicial-lemma-homotopy-s-Q}。
\end{proof}




\section{阿贝尔范畴中的更多同伦}
\label{simplicial-section-backwards-homotopy}

\noindent
设 $\mathcal{A}$ 为阿贝尔范畴。本节将证明：
$\text{Ch}_{\geq 0}(\mathcal{A})$ 中态射之间的同伦，
总是来自单纯对象范畴中的态射 $U \times \Delta[1] \to V$。
在某种意义下，这将给出引理 \ref{simplicial-lemma-homotopy-s-N} 的逆命题。
先发展一些关于链复形态射之间同伦的材料。

\begin{lemma}
\label{simplicial-lemma-represent-homotopy}
设 $\mathcal{A}$ 为阿贝尔范畴，$A$ 为链复形。考虑协变函子
$$
B \longmapsto
\{
(a, b, h)
\mid
a, b : A \to B\text{ 且 }h\text{ 下列二者之间的同伦： }a, b
\}
$$
存在链复形 $\diamond A$，使得
$\Mor_{\text{Ch}(\mathcal{A})}(\diamond A, -)$
同构于所显示的函子。构造 $A \mapsto \diamond A$ 具有函子性。
\end{lemma}

\begin{proof}
令 $\diamond A_n = A_n \oplus A_n \oplus A_{n - 1}$，
并以矩阵
$$
d_{\diamond A, n}
=
\left(
\begin{matrix}
d_{A, n} & 0 & \text{id}_{A_{n - 1}} \\
0 & d_{A, n} & -\text{id}_{A_{n - 1}} \\
0 & 0 & -d_{A, n - 1}
\end{matrix}
\right) :
A_n \oplus A_n \oplus A_{n - 1} \to
A_{n - 1} \oplus A_{n - 1} \oplus A_{n - 2}
$$
定义 $d_{\diamond A, n}$。若 $\mathcal{A}$ 是阿贝尔群范畴，
且 $(x, y, z) \in A_n \oplus A_n \oplus A_{n - 1}$，则
$d_{\diamond A, n}(x, y, z) =
(d_n(x) + z, d_n(y) - z, -d_{n - 1}(z))$。
容易验证 $d^2 = 0$。显然有两个映射
$\diamond a, \diamond b : A \to \diamond A$
（第一与第二直和分量），以及映射
$\diamond A \to A[-1]$，它们给出逐项分裂的短正合序列
$$
0 \to
A \oplus A \to
\diamond A \to
A[-1] \to
0
$$
。此外，有一列映射
$\diamond h_n : A_n \to \diamond A_{n + 1}$，
即从 $A_n$ 到 $\diamond A_{n + 1}$ 的分量 $A_n$ 的恒等映射，
使得 $\diamond h$ 是 $\diamond a$ 与 $\diamond b$ 之间的同伦。

\medskip\noindent
由此，任意态射 $f : \diamond A \to B$ 都通过令
$a = f \circ \diamond a$、$b = f \circ \diamond b$ 以及
$h_n = f_{n + 1} \circ \diamond h_n$ 给出三元组 $(a, b, h)$。
反之，给定三元组 $(a, b, h)$，取
$$
f_n = (a_n, b_n, h_{n - 1}).
$$
即可得到态射 $f : \diamond A \to B$。
要看出这是链复形的态射，需要作一次计算。这里只在
$\mathcal{A}$ 为阿贝尔群范畴的情形下计算：
设 $(x, y, z) \in \diamond A_n = A_n \oplus A_n \oplus A_{n - 1}$。
则
\begin{eqnarray*}
f_{n - 1}(d_n(x, y, z)) & = &
f_{n - 1}(d_n(x) + z, d_n(y) - z, -d_{n - 1}(z)) \\
& = &
a_n(d_n(x)) + a_n(z) + b_n(d_n(y)) - b_n(z) - h_{n - 2}(d_{n - 1}(z))
\end{eqnarray*}
% P02-E0141: frozen calculation uses degree-n components a_n and b_n on degree-(n-1) inputs.
以及
\begin{eqnarray*}
d_n(f_n(x, y, z) & = &
% P02-E0142: frozen display leaves the parenthesis in d_n(f_n(x,y,z) unclosed.
d_n(a_n(x) + b_n(y) + h_{n - 1}(z)) \\
& = &
d_n(a_n(x)) + d_n(b_n(y)) + d_n(h_{n - 1}(z))
\end{eqnarray*}
根据同伦的定义，二者相同。
\end{proof}

\noindent
注意，扩张
$$
0 \to
A \oplus A \to
\diamond A \to
A[-1] \to
0
$$
配有态射 $\diamond A_n \to A[-1]_n$ 的截面，并且相应的态射
$\delta : A[-1] \to (A \oplus A)[-1]$（见《同调代数》引理
\ref{homology-lemma-ses-termwise-split}）等于态射
$(1, -1) : A[-1] \to A[-1] \oplus A[-1]$。

\begin{lemma}
\label{simplicial-lemma-map-into-diamond}
设 $\mathcal{A}$ 为阿贝尔范畴，并设
$$
0 \to A \oplus A \to B \to C \to 0
$$
为 $\mathcal{A}$ 中链复形的短正合序列。进一步假设给定态射
$s_n : C_n \to B_n$，它们在次数 $n$ 上分裂相应的短正合序列。
设 $\delta(s) : C \to (A \oplus A)[-1] = A[-1] \oplus A[-1]$
为相应的复形态射；见《同调代数》引理
\ref{homology-lemma-ses-termwise-split}。若 $\delta(s)$ 经态射
$(1, -1) : A[-1] \to A[-1] \oplus A[-1]$ 分解，则存在唯一态射
$B \to \diamond A$，使其嵌入交换图
$$
\xymatrix{
0 \ar[r] &
A \oplus A \ar[d] \ar[r] &
B \ar[r] \ar[d] &
C \ar[d] \ar[r] &
0 \\
0 \ar[r] &
A \oplus A \ar[r] &
\diamond A \ar[r] &
A[-1] \ar[r] &
0
}
$$
，其中竖直映射也分别与截面 $s_n$ 及
$\diamond A_n \to A[-1]_n$ 的截面相容。
\end{lemma}

\begin{proof}
以 $(p_n, q_n) : B_n \to A_n \oplus A_n$ 表示《同调代数》引理
\ref{homology-lemma-ses-termwise-split} 中的态射 $\pi_n$。
又将短正合序列中的映射写成
$(a, b) : A \oplus A \to B$ 和 $r : B \to C$。
把 $\delta(s)$ 的分解写成
$\delta(s) = (1, -1) \circ f$。这意味着
$p_{n - 1} \circ d_{B, n} \circ s_n = f_n$，并且
$q_{n - 1} \circ d_{B, n} \circ s_n = - f_n$，并且
% P02-E0143: frozen source leaves a dangling “and” before the imperative “Set”.
令 $B_n \to \diamond A_n = A_n \oplus A_n \oplus A_{n - 1}$
等于 $(p_n, q_n, f_n \circ r_n)$。

\medskip\noindent
现在必须检验这确实定义了复形态射。这里只在阿贝尔群情形下验证。
取 $x \in B_n$。则
$x = a_n(x_1) + b_n(x_2) + s_n(x_3)$，并且只须分别对每一项证明
我们的定义与微分交换。对于项 $a_n(x_1)$，有
$(p_n, q_n, f_n \circ r_n)(a_n(x_1)) = (x_1, 0, 0)$，
结论显然。项 $b_n(x_2)$ 的情形类似。对于项 $s_n(x_3)$，有
\begin{eqnarray*}
(p_n, q_n, f_n \circ r_n)(d_n(s_n(x_3))) & = &
(p_n, q_n, f_n \circ r_n)( \\
& & \ \ \ \ \ a_n(f_n(x_3)) - b_n(f_n(x_3)) + s_n(d_n(x_3))) \\
& = &
(f_n(x_3), -f_n(x_3), f_n(d_n(x_3)))
\end{eqnarray*}
% P02-E0144: frozen first calculation applies degree-n projections to d_n(s_n(x_3)), which lies in degree n-1.
，这是由 $f_n$ 的定义得到的。另一方面，
\begin{eqnarray*}
d_n(p_n, q_n, f_n \circ r_n)(s_n(x_3)) & = & d_n(0, 0, f_n(x_3)) \\
& = &
(f_n(x_3), - f_n(x_3), d_{A[-1], n}(f_n(x_3)))
\end{eqnarray*}
由于 $f$ 是复形态射，结论成立。
\end{proof}

\begin{lemma}
\label{simplicial-lemma-backwards-homotopy}
设 $\mathcal{A}$ 为阿贝尔范畴，$U$、$V$ 为 $\mathcal{A}$
的单纯对象，$a, b : U \to V$ 为一对态射。假设相应的链复形映射
$N(a), N(b) : N(U) \to N(V)$ 通过同伦
$\{N_n : N(U)_n \to N(V)_{n + 1}\}$ 同伦。
则存在定义 \ref{simplicial-definition-homotopy} 意义下从 $a$ 到 $b$
的同伦。此外，可以选取同伦
$h : U \times \Delta[1] \to V$，使得
$N_n = N(h)_n$，其中 $N(h)$ 是按第
\ref{simplicial-section-homotopy-abelian} 节从 $h$ 得到的同伦。
\end{lemma}

\begin{proof}
设 $(\diamond N(U), \diamond a, \diamond b, \diamond h)$
如引理 \ref{simplicial-lemma-represent-homotopy} 及其证明中所述。
由该引理，存在表示三元组 $(N(a), N(b), \{N_n\})$ 的态射
$\diamond N(U) \to N(V)$。下面证明存在态射
$\psi : N(U \times \Delta[1]) \to \diamond{N(U)}$，
使得 $\diamond a = \psi \circ N(e_0)$ 且
$\diamond b = \psi \circ N(e_1)$。还将证明：
由 (\ref{simplicial-equation-homotopy-to-homotopy}) 及引理
\ref{simplicial-lemma-homotopy-s-N}，取
$h = \text{id}_{U \times \Delta[1]}$ 得到的
$N(e_0), N(e_1) :
N(U) \to N(U \times \Delta[1])$ 之间的同伦，经 $\psi$
映到映射
$\diamond a, \diamond b : N(U) \to \diamond{N(U)}$
之间的典范同伦 $\diamond h$。这当然蕴含引理。

\medskip\noindent
注意，作为函子，
$N : \text{Simp}(\mathcal{A}) \to \text{Ch}_{\geq 0}(\mathcal{A})$
是函子
$s : \text{Simp}(\mathcal{A}) \to \text{Ch}_{\geq 0}(\mathcal{A})$
的直和分量。并且函子 $\diamond$ 与直和相容。
因此，只须改为构造具有相应性质的态射
$\Psi : s(U \times \Delta[1]) \to \diamond{s(U)}$。
下面正是这一构造。

\medskip\noindent
由定义 \ref{simplicial-definition-homotopy}，态射
$e_0 : U \to U \times \Delta[1]$ 与
$e_1 : U \to U \times \Delta[1]$
通过同伦 $\text{id}_{U \times \Delta[1]}$ 同伦。
由引理 \ref{simplicial-lemma-homotopy-s-N}，得到链复形态射
$s(e_0) : s(U) \to s(U \times \Delta[1])$ 与
$s(e_1) : s(U) \to s(U \times \Delta[1])$ 之间的显式同伦
$\{h_n : s(U)_n \to s(U \times \Delta[1])_{n + 1}\}$。
由上面的引理 \ref{simplicial-lemma-map-into-diamond}，得到相应态射
$$
\Phi : \diamond{s(U)} \to s(U \times \Delta[1])
$$
根据构造，$\Phi_n$ 在 $\diamond{s(U)}_n$ 的分量
$s(U)[-1]_n = s(U)_{n - 1}$ 上的限制等于 $h_{n - 1}$。并且
$$
h_{n - 1} = \sum\nolimits_{i = 0}^{n - 1}
(-1)^{i + 1} s^n_i \cdot \alpha^n_{i + 1} :
U_{n - 1} \to \bigoplus\nolimits_j U_n \cdot \alpha^n_j.
$$
% P02-E0145: frozen formula uses s_i^n on U_{n-1}; the degeneracy with this source degree is s_i^{n-1}.
这里采用显然的记号。

\medskip\noindent
另一方面，态射 $e_i : U \to U \times \Delta[1]$ 诱导态射
$(e_0, e_1) : U \oplus U \to U \times \Delta[1]$。
以 $W$ 表示余核。注意，若写成
$(U \times \Delta[1])_n =
\bigoplus_{\alpha : [n] \to [1]} U_n \cdot \alpha$，
则可以把
$W_n = \bigoplus_{i = 1}^n U_n \cdot \alpha^n_i$
作如第 \ref{simplicial-section-homotopy} 节中的 $\alpha^n_i$ 那样的识别。
有交换图
$$
\xymatrix{
0 \ar[r] &
U \oplus U \ar[rd]_{(1, 1)} \ar[r] &
U \times \Delta[1] \ar[d]^\pi \ar[r] &
W \ar[r] &
0 \\
& & U & &
}
$$
这意味着应用函子 $s$ 后有类似的交换图。然后选取分裂
$\sigma_n : s(W)_n \to s(U \times \Delta[1])_n$：
通过 $(-1, 1)$，把分量 $U_n \cdot \alpha^n_i \subset W_n$
映到分量
$U_n \cdot \alpha^n_0 \oplus U_n \cdot \alpha^n_i
\subset (U \times \Delta[1])_n$。
注意 $s(\pi)_n \circ \sigma_n = 0$，从而
$(1, 1) \circ \delta(\sigma)_n = 0$。因此
$\delta(\sigma)$ 如引理 \ref{simplicial-lemma-map-into-diamond} 中那样分解。
由该引理，得到典范态射
$\Psi : s(U \times \Delta[1]) \to \diamond{s(U)}$。

\medskip\noindent
为了计算 $\Psi$，先计算态射
$\delta(\sigma) : s(W) \to s(U)[-1] \oplus s(U)[-1]$。
根据《同调代数》引理
\ref{homology-lemma-ses-termwise-split} 及其证明，
为此必须计算
% P02-E0146: frozen source has “we have compute”, omitting “to”.
$$
d_{s(U \times \delta[1]), n} \circ \sigma_n
-
\sigma_{n - 1} \circ d_{s(W), n}
$$
% P02-E0147: frozen source writes lowercase \delta[1] in three places where the simplicial interval is \Delta[1].
并把它写成到
$U_{n - 1} \cdot \alpha^{n - 1}_0
\oplus U_{n - 1} \cdot \alpha^{n - 1}_n$ 的态射。
这里只在 $\mathcal{A}$ 是阿贝尔群范畴的情形下计算。
采用简写 $x_{\alpha}$：对 $x \in U_n$，以它表示
$(U \times \Delta[1])_n$ 的分量
$U_n \cdot \alpha$ 中的元素 $x$。回顾
$$
d_{s(U \times \delta[1]), n}
=
\sum\nolimits_{i = 0}^n (-1)^i d^n_i
$$
其中 $d^n_i$ 通过单纯对象 $U$ 的态射 $d^n_i$，
把分量 $U_n \cdot \alpha$ 映到分量
$U_{n - 1} \cdot (\alpha \circ \delta^n_i)$。
用上述记号，这意味着
$$
d_{s(U \times \delta[1]), n}(x_\alpha) =
\sum\nolimits_{i = 0}^n (-1)^i (d^n_i(x))_{\alpha \circ \delta^n_i}
$$
从 $x_\alpha \in W_n$ 出发，换言之，对某个
$j \in \{1, \ldots, n\}$ 有 $\alpha = \alpha^n_j$，
可见 $\sigma_n(x_\alpha) = x_\alpha - x_{\alpha^n_0}$，因而
$$
(d_{s(U \times \delta[1]), n} \circ \sigma_n)(x_\alpha)
=
\sum\nolimits_{i = 0}^n (-1)^i (d^n_i(x))_{\alpha \circ \delta^n_i}
-
\sum\nolimits_{i = 0}^n (-1)^i (d^n_i(x))_{\alpha^n_0 \circ \delta^n_i}
$$
为了计算 $d_{s(W), n}(x_\alpha)$，必须删去所有满足
$\alpha \circ \delta^n_i = \alpha^{n - 1}_0, \alpha^{n - 1}_n$
的项。因此得到
$$
\begin{matrix}
(\sigma_{n - 1} \circ d_{s(W), n})(x_\alpha) = \\
\sum\nolimits_{i = 0, \ldots, n\text{ 且 }
\alpha \circ \delta^n_i \not = \alpha^{n - 1}_0\text{ 或 }\alpha^{n - 1}_n}
\Big((-1)^i (d^n_i(x))_{\alpha \circ \delta^n_i}
-
(-1)^i (d^n_i(x))_{\alpha^{n - 1}_0}
\Big)
\end{matrix}
$$
% P02-E0148: frozen exclusion predicate uses “not equal A or B”, which is logically nonrestrictive; excluding both boundary values requires conjunction.
显然，两项之差是和
$$
\sum\nolimits_{i = 0, \ldots, n\text{ 且 }
\alpha \circ \delta^n_i = \alpha^{n - 1}_0\text{ 或 }\alpha^{n - 1}_n}
\Big((-1)^i (d^n_i(x))_{\alpha \circ \delta^n_i}
-
(-1)^i (d^n_i(x))_{\alpha^{n - 1}_0}
\Big)
$$
当然，若 $\alpha \circ \delta^n_i = \alpha^{n - 1}_0$，
则该项消失。回顾对某个 $j \in \{1, \ldots, n\}$，
有 $\alpha = \alpha^n_j$。等式
$\alpha^n_j \circ \delta^n_i = \alpha^{n - 1}_n$
成立的唯一可能是 $j = n$ 且 $i = n$。因此，除非
$j = n$，实际所得为 $0$；在该情形下得到
$(-1)^n(d^n_n(x))_{\alpha^{n - 1}_n}
- (-1)^n(d^n_n(x))_{\alpha^{n - 1}_0}$。
换言之，态射
$$
\delta(\sigma)_n :
W_n
\to
(s(U)[-1] \oplus s(U)[-1])_n = U_{n - 1} \oplus U_{n - 1}
$$
在除 $U_n \cdot \alpha^n_n$ 以外的所有分量上为零，
而在该分量上等于
$((-1)^nd^n_n, -(-1)^nd^n_n)$。
（也就是说，两个分量中的第一个对应于
$\alpha^{n - 1}_n$ 所在的因子，因为这是把所有元素映到 $0$
的映射 $[n - 1] \to [1]$，因而对应于 $e_0$。）

\medskip\noindent
得到典范图
$$
\xymatrix{
0 \ar[r] &
s(U) \oplus s(U) \ar[r] \ar[d] &
\diamond{s(U)} \ar[r] \ar[d]^{\Phi}&
s(U)[-1] \ar[r] \ar[d] &
0 \\
0 \ar[r] &
s(U) \oplus s(U) \ar[r] \ar[d] &
s(U \times \Delta[1]) \ar[r] \ar[d]^\Psi &
s(W) \ar[r] \ar[d] &
0 \\
0 \ar[r] &
s(U) \oplus s(U) \ar[r] &
\diamond{s(U)} \ar[r] &
s(U)[-1] \ar[r] &
0
}
$$
断言 $\Phi \circ \Psi$ 是恒等映射。为此，只须证明
$\Phi$ 与 $\delta(\sigma)$ 的复合，作为映射
$s(U)[-1] \to s(W) \to s(U)[-1] \oplus s(U)[-1]$，
在第一个因子上是恒等映射，在第二个因子上是负恒等映射。
根据上面的计算，它是
$((-1)^nd^n_0, -(-1)^nd^n_0) \circ (-1)^n s^n_n = (1, -1)$，
% P02-E0149: frozen final composition uses d_0^n after the preceding computation produced d_n^n; its displayed degrees are also not composable with s_n^n.
正如所需。
\end{proof}









\section{平凡 Kan 纤维化}
\label{simplicial-section-trivial-kan}

\noindent
回顾当 $n \geq 0$ 时，单纯集 $\Delta[n]$ 由规则
$[k] \mapsto \Mor_\Delta([k], [n])$ 给出；见例
\ref{simplicial-example-simplex-simplicial-set}。回顾 $\Delta[n]$ 有唯一的
非退化 $n$-单形，并且所有非退化单形都是这个 $n$-单形的面。
事实上，$\Delta[n]$ 的非退化单形恰好对应于单射
$[k] \to [n]$，可以把它们与 $[n]$ 的子集等同起来。
此外，回顾对任意单纯集 $X$，有
$\Mor(\Delta[n], X) = X_n$（引理 \ref{simplicial-lemma-simplex-map}）。
令
$$
\partial \Delta[n] = i_{(n - 1)!}\text{sk}_{n - 1}\Delta[n]
$$
并称它为 $\Delta[n]$ 的 {\it 边界}。
由引理 \ref{simplicial-lemma-n-skeleton-sets} 可知，
$\partial \Delta[n] \subset \Delta[n]$ 是这样的单纯子集：
它在次数 $\leq n - 1$ 中具有相同的非退化单形，
但不含那个非退化 $n$-单形。

\begin{definition}
\label{simplicial-definition-trivial-kan}
若单纯集映射 $X \to Y$ 满足 $X_0 \to Y_0$ 为满射，并且对所有
$n \geq 1$ 及任意交换实线图
$$
\xymatrix{
\partial \Delta[n] \ar[r] \ar[d] & X \ar[d] \\
\Delta[n] \ar[r] \ar@{-->}[ru] & Y
}
$$
都存在使图交换的虚线箭头，则称它为 {\it 平凡 Kan 纤维化}。
\end{definition}

\noindent
平凡 Kan 纤维化满足一种非常一般的提升性质。

\begin{lemma}
\label{simplicial-lemma-trivial-kan}
设 $f : X \to Y$ 为单纯集的平凡 Kan 纤维化。
对任意单纯集交换实线图
$$
\xymatrix{
Z \ar[r]_b \ar[d] & X \ar[d] \\
W \ar[r]^a \ar@{-->}[ru] & Y
}
$$
，若 $Z \to W$（逐项）为单射，则存在使图交换的虚线箭头。
\end{lemma}

\begin{proof}
假设 $Z \not = W$。设 $n$ 为满足 $Z_n \not = W_n$
的最小整数。取 $x \in W_n$，且 $x \not \in Z_n$。
以 $Z' \subset W$ 表示包含 $Z$、$x$ 以及 $x$ 的所有退化的
单纯子集。设 $\varphi : \Delta[n] \to Z'$ 为与 $x$
对应的态射（引理 \ref{simplicial-lemma-simplex-map}）。
由于 $\partial \Delta[n]$ 的所有非退化单形都映入 $Z$，
$\varphi|_{\partial \Delta[n]}$ 映入 $Z$。
由假设，可将 $b \circ \varphi|_{\partial \Delta[n]}$
延拓为 $\beta : \Delta[n] \to X$。
由引理 \ref{simplicial-lemma-glue-simplex}，单纯集 $Z'$ 是
$\Delta[n]$ 与 $Z$ 沿 $\partial \Delta[n]$ 的推出。
因此 $b$ 和 $\beta$ 定义态射 $b' : Z' \to X$。
换言之，我们已把态射 $b$ 延拓到 $Z$ 的一个更大单纯子集。
% P02-E0150: frozen source says the larger simplicial subset is “of Z”, although Z' strictly contains Z and is a subset of W.

\medskip\noindent
应用 Zorn 引理即可完成证明（从略）。
\end{proof}

\begin{lemma}
\label{simplicial-lemma-trivial-kan-base-change}
设 $f : X \to Y$ 为单纯集的平凡 Kan 纤维化，并设
$Y' \to Y$ 为单纯集态射。则
$X \times_Y Y' \to Y'$ 是平凡 Kan 纤维化。
\end{lemma}

\begin{proof}
这立即由纤维积的函子性（引理 \ref{simplicial-lemma-fibre-product}）和定义得出。
\end{proof}

\begin{lemma}
\label{simplicial-lemma-trivial-kan-composition}
两个平凡 Kan 纤维化的复合仍是平凡 Kan 纤维化。
\end{lemma}

\begin{proof}
略。
\end{proof}

\begin{lemma}
\label{simplicial-lemma-limit-trivial-kan}
设 $\ldots \to U^2 \to U^1 \to U^0$ 为一列平凡 Kan 纤维化。
令 $U = \lim U^t$，其定义为逐项取 $U_n = \lim U_n^t$。
则 $U \to U^0$ 是平凡 Kan 纤维化。
\end{lemma}

\begin{proof}
略。提示：使用集合的可数满射序列之逆极限非空这一事实。
\end{proof}

\begin{lemma}
\label{simplicial-lemma-product-trivial-kan}
\begin{slogan}
平凡 Kan 纤维化的积仍是平凡 Kan 纤维化。
\end{slogan}
设 $X_i \to Y_i$ 为一族平凡 Kan 纤维化。则
$\prod X_i \to \prod Y_i$ 是平凡 Kan 纤维化。
\end{lemma}

\begin{proof}
略。
\end{proof}

\begin{lemma}
\label{simplicial-lemma-filtered-colimit-trivial-kan}
平凡 Kan 纤维化的滤过余极限仍是平凡 Kan 纤维化。
\end{lemma}

\begin{proof}
略。提示：见《范畴论》第
\ref{categories-section-directed-colimits} 节中集合的滤过余极限之描述。
\end{proof}

\begin{lemma}
\label{simplicial-lemma-trivial-kan-homotopy}
设 $f : X \to Y$ 为单纯集的平凡 Kan 纤维化。
则 $f$ 是同伦等价。
\end{lemma}

\begin{proof}
由引理 \ref{simplicial-lemma-trivial-kan}，可以选取 $f$ 的右逆
$g : Y \to X$。
% P02-E0151: frozen source says “an right inverse”.
考虑图
$$
\xymatrix{
\partial \Delta[1] \times X \ar[d] \ar[r] & X \ar[d] \\
\Delta[1] \times X \ar[r] \ar@{-->}[ru] & Y
}
$$
这里，上方水平箭头由 $\text{id}_X$ 和 $g \circ f$ 给出；
使用了对所有 $n \geq 0$ 都有
$(\partial \Delta[1] \times X)_n = X_n \amalg X_n$。
下方水平箭头由映射 $\Delta[1] \to \Delta[0]$ 和
$f : X \to Y$ 给出。由于 $f \circ g \circ f = f$，该图交换。
由引理 \ref{simplicial-lemma-trivial-kan} 可以填入虚线箭头，结论即得。
\end{proof}





\section{Kan 纤维化}
\label{simplicial-section-kan}

\noindent
设 $n$、$k$ 为满足 $0 \leq k \leq n$ 且 $1 \leq n$ 的整数。
令 $\sigma_0, \ldots, \sigma_n$ 为 $\Delta[n]$ 的唯一非退化
$n$-单形 $\sigma$ 的 $n + 1$ 个面，即
$\sigma_i = d_i\sigma$。令
$$
\Lambda_k[n] \subset \Delta[n]
$$
为 $n$-单形 $\Delta[n]$ 的 {\it 第 $k$ 个角}。
它是由 $\sigma_0, \ldots, \hat \sigma_k, \ldots, \sigma_n$
生成的 $\Delta[n]$ 的单纯子集。换言之，所显示包含的像含有
$\Delta[n]$ 的所有非退化单形，唯独不含 $\sigma$ 和 $\sigma_k$。

\begin{definition}
\label{simplicial-definition-kan}
若单纯集映射 $X \to Y$ 满足：对所有整数 $k, n$，其中
$1 \leq n$、$0 \leq k \leq n$，以及任意交换实线图
$$
\xymatrix{
\Lambda_k[n] \ar[r] \ar[d] & X \ar[d] \\
\Delta[n] \ar[r] \ar@{-->}[ru] & Y
}
$$
都存在使图交换的虚线箭头，则称它为 {\it Kan 纤维化}。
若单纯集 $X$ 到单点集上的常值单纯集 $*$ 的映射
$X \to *$ 为 Kan 纤维化，则称它为 {\it Kan 复形}。
\end{definition}

\noindent
注意，$\Lambda_k[n]$ 总是非空的。因此，从空单纯集到任意单纯集
的态射总是 Kan 纤维化。由引理 \ref{simplicial-lemma-trivial-kan} 可知，
平凡 Kan 纤维化是 Kan 纤维化。

\begin{lemma}
\label{simplicial-lemma-kan-base-change}
设 $f : X \to Y$ 为单纯集的 Kan 纤维化，并设
$Y' \to Y$ 为单纯集态射。则
$X \times_Y Y' \to Y'$ 是 Kan 纤维化。
\end{lemma}

\begin{proof}
这立即由纤维积的函子性（引理 \ref{simplicial-lemma-fibre-product}）和定义得出。
\end{proof}

\begin{lemma}
\label{simplicial-lemma-kan-composition}
两个 Kan 纤维化的复合仍是 Kan 纤维化。
\end{lemma}

\begin{proof}
略。
\end{proof}

\begin{lemma}
\label{simplicial-lemma-limit-kan}
设 $\ldots \to U^2 \to U^1 \to U^0$ 为一列 Kan 纤维化。
令 $U = \lim U^t$，其定义为逐项取 $U_n = \lim U_n^t$。
则 $U \to U^0$ 是 Kan 纤维化。
\end{lemma}

\begin{proof}
略。提示：使用集合的可数满射序列之逆极限非空这一事实。
\end{proof}

\begin{lemma}
\label{simplicial-lemma-product-kan}
设 $X_i \to Y_i$ 为一族 Kan 纤维化。则
$\prod X_i \to \prod Y_i$ 是 Kan 纤维化。
\end{lemma}

\begin{proof}
略。
\end{proof}

\noindent
下面的引理由 J.C.\ Moore 给出，见 \cite{Moore-Cartan}。

\begin{lemma}
\label{simplicial-lemma-simplicial-group-kan}
设 $X$ 为单纯群。则 $X$ 是 Kan 复形。
\end{lemma}

\begin{proof}
下面的证明基本上只是把上述参考文献中的证明译成英语。
% The frozen English source explicitly characterizes this proof as an English translation.
使用本节引言中说明的术语，设
$f : \Lambda_k[n] \to X$ 为从一个角出发的态射。
对 $i = 0, \ldots, \hat k, \ldots, n$，令
$x_i = f(\sigma_i) \in X_{n - 1}$。这意味着，当 $i < j$
且 $i, j \not = k$ 时，有 $d_i x_j = d_{j - 1} x_i$。
必须找到 $x \in X_n$，使得对
$i = 0, \ldots, \hat k, \ldots, n$，有 $x_i = d_ix$。

\medskip\noindent
先证明存在 $u \in X_n$，使得对 $i < k$ 有 $d_i u = x_i$。
当 $k = 0$ 时这是平凡的。若 $k > 0$，则归纳定义
$u^r \in X_n$，使得对 $0 \leq i \leq r$ 有
$d_i u^r = x_i$。从 $u^0 = s_0x_0$ 开始。若 $r < k - 1$，令
$$
y^r = s_{r + 1}((d_{r + 1}u^r)^{-1}x_{r + 1}),\quad
u^{r + 1} = u^r y^r.
$$
简单计算表明，当 $i \leq r$ 时，$d_iy^r = 1$
（群 $X_{n - 1}$ 的单位元），且
$d_{r + 1}y^r = (d_{r + 1}u^r)^{-1}x_{r + 1}$。
因此，当 $i \leq r + 1$ 时，$d_iu^{r + 1} = x_i$。
最后取 $u = u^{k - 1}$，得到所需的 $u$。

\medskip\noindent
接着对整数 $r$ 归纳证明：当 $0 \leq r \leq n - k$ 时，
存在 $x^r \in X_n$，使得
% P02-E0152: frozen source says “there exists a x^r”.
$$
d_i x^r = x_i\quad\text{对 }i < k\text{ 且 }i > n - r.
$$
当 $r = 0$ 时，从 $x^0 = u$ 开始。若已对
$r \leq n - k - 1$ 定义了 $x^r$，令
$$
z^r = s_{n - r - 1}((d_{n - r}x^r)^{-1}x_{n - r}),\quad x^{r + 1} = x^rz^r
$$
利用给定关系作简单计算，可知当 $i < k$ 及 $i > n - r$
时 $d_iz^r = 1$，并且
$d_{n - r}(z^r) = (d_{n - r}x^r)^{-1}x_{n - r}$。
从而当 $i < k$ 及 $i > n - r - 1$ 时，
$d_ix^{r + 1} = x_i$。最后取 $x = x^{n - k}$，证明完成。
\end{proof}

\begin{lemma}
\label{simplicial-lemma-surjection-simplicial-abelian-groups-kan}
设 $f : X \to Y$ 为逐项满射的单纯阿贝尔群同态。
则 $f$ 是单纯集的 Kan 纤维化。
\end{lemma}

\begin{proof}
考虑定义 \ref{simplicial-definition-kan} 中的交换实线图
$$
\xymatrix{
\Lambda_k[n] \ar[r]_a \ar[d] & X \ar[d] \\
\Delta[n] \ar[r]^b \ar@{-->}[ru] & Y
}
$$
映射 $a$ 对应于满足
$d_i x_j = d_{j - 1} x_i$（当 $i < j$ 且 $i, j \not = k$）
的元素
$x_0, \ldots, \hat x_k, \ldots, x_n \in X_{n - 1}$。
映射 $b$ 对应于元素 $y \in Y_n$，使得当 $i \not = k$
时 $d_iy = f(x_i)$。任务是构造 $x \in X_n$，
使得当 $i \not = k$ 时 $d_ix = x_i$，且 $f(x) = y$。

\medskip\noindent
由于 $f$ 逐项满射，可取 $x \in X_n$ 使得 $f(x) = y$。
以 $0 = y - f(x)$ 代替 $y$，并对 $i \not = k$，
以 $x_i - d_ix$ 代替 $x_i$。于是可以假设 $y = 0$。
特别地，$f(x_i) = 0$。换言之，可以把 $X$ 替换为
$\Ker(f) \subset X$，并把 $Y$ 替换为 $0$。
在这种情形下，该陈述化为引理
\ref{simplicial-lemma-simplicial-group-kan}。
% P02-E0153: frozen source says “the statement become”.
\end{proof}

\begin{lemma}
\label{simplicial-lemma-qis-simplicial-abelian-groups}
设 $f : X \to Y$ 为逐项满射的单纯阿贝尔群同态，并且在伴随链复形上
诱导拟同构。则 $f$ 是单纯集的平凡 Kan 纤维化。
\end{lemma}

\begin{proof}
考虑定义 \ref{simplicial-definition-trivial-kan} 中的交换实线图
$$
\xymatrix{
\partial \Delta[n] \ar[r]_a \ar[d] & X \ar[d] \\
\Delta[n] \ar[r]^b \ar@{-->}[ru] & Y
}
$$
映射 $a$ 对应于满足当 $i < j$ 时
$d_i x_j = d_{j - 1} x_i$ 的元素
$x_0, \ldots, x_n \in X_{n - 1}$。
映射 $b$ 对应于元素 $y \in Y_n$，使得 $d_iy = f(x_i)$。
任务是构造 $x \in X_n$，使得 $d_ix = x_i$ 且 $f(x) = y$。

\medskip\noindent
由于 $f$ 逐项满射，可取 $x \in X_n$ 使得 $f(x) = y$。
以 $0 = y - f(x)$ 代替 $y$，并以 $x_i - d_ix$ 代替 $x_i$。
于是可以假设 $y = 0$。特别地，$f(x_i) = 0$。
换言之，可以把 $X$ 替换为 $\Ker(f) \subset X$，
并把 $Y$ 替换为 $0$。这是可行的，因为由《同调代数》引理
\ref{homology-lemma-long-exact-sequence-chain}，
$\Ker(f)$ 的伴随链复形之同调为零，因此
$\Ker(f) \to 0$ 在伴随链复形上诱导拟同构。

\medskip\noindent
由于 $X$ 是 Kan 复形（引理 \ref{simplicial-lemma-simplicial-group-kan}），
可取 $x \in X_n$，使得对 $i = 0, \ldots, n - 1$ 有
$d_i x = x_i$。对 $i = 0, \ldots, n$，以
$x_i - d_ix$ 代替 $x_i$ 后，可以假设
$x_0 = x_1 = \ldots = x_{n - 1} = 0$。
此时可见，当 $i = 0, \ldots, n - 1$ 时，$d_i x_n = 0$。
因此 $x_n \in N(X)_{n - 1}$，并属于微分
$N(X)_{n - 1} \to N(X)_{n - 2}$ 的核。
这里 $N(X)$ 是与 $X$ 伴随的规范化链复形，见第
\ref{simplicial-section-complexes} 节。由于 $N(X)$ 拟同构于 $s(X)$
（引理 \ref{simplicial-lemma-quasi-isomorphism}），因而是零调的，
可取 $x \in N(X_n)$，其微分为 $x_n$。
% P02-E0154: frozen source writes N(X_n) although the normalized complex term is N(X)_n (equivalently N(X_n) is not defined by this notation).
这个 $x$ 回答了引理提出的问题，证明完成。
\end{proof}

\begin{lemma}
\label{simplicial-lemma-homotopy-equivalence}
设 $f : X \to Y$ 为单纯阿贝尔群的映射。若 $f$ 是单纯集的同伦等价，
则 $f$ 在伴随链复形上诱导拟同构。
\end{lemma}

\begin{proof}
在本证明中，当 $Z$ 是单纯阿贝尔群时，记
$H_n(Z) = H_n(s(Z)) = H_n(N(Z))$，其中 $s$ 与 $N$
如第 \ref{simplicial-section-complexes} 节所述。
以 $\mathbf{Z}[X]$ 表示把 $X$ 视为单纯集后生成的自由阿贝尔群，
$\mathbf{Z}[Y]$ 亦然。考虑单纯阿贝尔群的交换图
$$
\xymatrix{
\mathbf{Z}[X] \ar[r]_g \ar[d] &
\mathbf{Z}[Y] \ar[d] \\
X \ar[r]^f & Y
}
$$
由于在集合上取自由阿贝尔群是函子，水平箭头是单纯阿贝尔群的
同伦等价；见引理 \ref{simplicial-lemma-functorial-homotopy}。
由引理 \ref{simplicial-lemma-homotopy-equivalence-s-N}，对所有
$n \geq 0$，映射
$H_n(g) : H_n(\mathbf{Z}[X]) \to H_n(\mathbf{Z}[Y])$ 都是双射。

\medskip\noindent
设 $\xi \in H_n(Y)$。根据 $N(Y)$ 的定义，可以用边界为零的元素
$y \in N(Y_n)$ 表示 $\xi$。这意味着 $y \in Y_n$，并且
$d^n_0(y) = \ldots = d^n_{n - 1}(y) = 0$，因为
$y \in N(Y_n)$；并且 $d^n_n(y) = 0$，因为 $y$ 的边界为零。
以 $0_n \in Y_n$ 表示零元。则
$$
\tilde y = [y] - [0_n] \in (\mathbf{Z}[Y])_n
$$
满足 $d^n_0(\tilde y) = \ldots = d^n_{n - 1}(\tilde y) = 0$
以及 $d^n_n(\tilde y) = 0$。因此，$\tilde y$ 属于
$N(\mathbf{Z}[Y])_n$，其边界为 $0$；也就是说，
$\tilde y$ 决定一个映到 $\xi$ 的类
$\tilde \xi \in H_n(\mathbf{Z}[Y])$。
% P02-E0155: frozen source joins “is in ... has boundary 0” without “and”.
因为 $H_n(\mathbf{Z}[X]) \to H_n(\mathbf{Z}[Y])$ 是双射，
可以把 $\tilde \xi$ 提升为 $H_n(\mathbf{Z}[X])$ 中的一个类。
由上面的交换图可知，$\xi$ 属于 $H_n(X) \to H_n(Y)$ 的像。

\medskip\noindent
设 $\xi \in H_n(X)$ 映为 $H_n(Y)$ 中的零。
与上一段完全一样，可以用边界为零的元素
$x \in N(X_n)$ 表示 $\xi$，即
$d^n_0(x) = \ldots = d^n_{n - 1}(x) = d^n_n(x) = 0$。
% P02-E0156: frozen source misspells “paragraph” as “parapgraph”.
特别地，$[x] - [0_n]$ 是 $N(\mathbf{Z}[X])_n$ 中边界为零的元素，
从而定义 $\xi$ 的提升
$\tilde \xi \in H_n(\mathbf{Z}[x])$。
% P02-E0157: frozen source uses lowercase x in \mathbf{Z}[x] instead of the simplicial set X.
$\xi$ 映为 $H_n(Y)$ 中的零意味着存在
$y \in N(Y_{n + 1})$，其边界为 $f_n(x)$。
这意味着 $d^{n + 1}_0(y) = \ldots = d^{n + 1}_n(y) = 0$ 且
$d^{n + 1}_{n + 1}(y) = f(x)$。
然而，这恰好意味着
$z = [y] - [0_{n + 1}]$ 属于
$N(\mathbf{Z}[y])_{n + 1}$，并且
% P02-E0158: frozen source uses \mathbf{Z}[y] where the ambient free simplicial abelian group is \mathbf{Z}[Y].
$$
g([x] - [0_n]) = [f(x)] - [0_n] = \text{下列对象的边界： }z
$$
这证明 $\tilde \xi$ 映为 $H_n(\mathbf{Z}[y])$ 中的零。
% P02-E0159: frozen source again uses \mathbf{Z}[y] instead of \mathbf{Z}[Y].
因为 $H_n(\mathbf{Z}[X]) \to H_n(\mathbf{Z}[Y])$ 是双射，
可得 $\tilde \xi = 0$，从而 $\xi = 0$。
\end{proof}






\section{一个同伦等价}
\label{simplicial-section-homotopy-equivalence}

\noindent
设 $A$、$B$ 为集合，$f : A \to B$ 为映射。考虑相应的单纯集映射
$$
\xymatrix{
\text{cosk}_0(A) \ar@{=}[r] &
\Big(\ldots A \times A \times A
\ar[d] \ar@<2ex>[r] \ar@<0ex>[r] \ar@<-2ex>[r] &
A \times A \ar[d] \ar@<1ex>[r] \ar@<-1ex>[r] \ar@<1ex>[l] \ar@<-1ex>[l] &
A \Big) \ar[d] \ar@<0ex>[l] \\
\text{cosk}_0(B) \ar@{=}[r] &
\Big( \ldots B \times B \times B \ar@<2ex>[r] \ar@<0ex>[r] \ar@<-2ex>[r] &
B \times B \ar@<1ex>[r] \ar@<-1ex>[r] \ar@<1ex>[l] \ar@<-1ex>[l] &
B \Big) \ar@<0ex>[l]
}
$$
见例 \ref{simplicial-example-cosk0}。下面引理在 $n = 0$ 的情形说明：
若 $f$ 为满射，则这个单纯集映射是平凡 Kan 纤维化。

\begin{lemma}
\label{simplicial-lemma-section}
设 $f : V \to U$ 为单纯集态射，$n \geq 0$ 为整数。假设：
\begin{enumerate}
\item 当 $i < n$ 时，映射 $f_i : V_i \to U_i$ 是双射；
\item 映射 $f_n : V_n \to U_n$ 是满射；
\item 典范态射 $U \to \text{cosk}_n \text{sk}_n U$ 是同构；
\item 典范态射 $V \to \text{cosk}_n \text{sk}_n V$ 是同构。
\end{enumerate}
则 $f$ 是平凡 Kan 纤维化。
\end{lemma}

\begin{proof}
考虑定义 \ref{simplicial-definition-trivial-kan} 中的实线图
$$
\xymatrix{
\partial \Delta[k] \ar[r] \ar[d] & V \ar[d] \\
\Delta[k] \ar[r] \ar@{-->}[ru] & U
}
$$
。设 $x \in U_k$ 为与下方水平箭头对应的 $k$-单形。
若 $k \leq n$，则虚线箭头与 $x$ 的提升
$y \in V_k$ 对应；该图交换，因为 $\Delta[k]$ 的其他非退化单形
都在次数 $< k$ 中，而 $f$ 在这些次数中为同构。
若 $k > n$，则由条件 (3) 和 (4)，利用骨架函子与余骨架函子的伴随性，
有
$$
\Mor(\Delta[k], U) =
\Mor(\text{sk}_n\Delta[k], \text{sk}_nU) =
\Mor(\text{sk}_n\partial\Delta[k], \text{sk}_nU) =
\Mor(\partial \Delta[k], U)
$$
，$V$ 也同样，因为当 $k > n$ 时，
$\text{sk}_n\Delta[k] = \text{sk}_n\partial\Delta[k]$。
因此，在这种情形下也得到唯一的虚线箭头嵌入该图。
\end{proof}

\noindent
设 $A, B$ 为集合，$f^0, f^1 : A \to B$ 为集合映射。
考虑同样记号所表示的诱导映射
$f^0, f^1 : \text{cosk}_0(A) \to \text{cosk}_0(B)$。
下面引理在 $n = 0$ 的情形说明 $f^0$ 同伦于 $f^1$。
事实上，存在从 $f^0$ 到 $f^1$ 的同伦
$h : \text{cosk}_0(A) \times \Delta[1] \to \text{cosk}_0(A)$，
% P02-E0160: frozen source gives the homotopy codomain cosk_0(A), but f^0 and f^1 land in cosk_0(B).
其分量为
\begin{eqnarray*}
h_m : A \times \ldots \times A \times \Mor_{\Delta}([m], [1])
& \longrightarrow &
B \times \ldots \times B, \\
(a_0, \ldots, a_m, \alpha) & \longmapsto &
(f^{\alpha(0)}(a_0), \ldots, f^{\alpha(m)}(a_m))
\end{eqnarray*}
为了检验其成立，注意对映射 $\varphi : [k] \to [m]$，
诱导映射为
$(a_0, \ldots, a_m) \mapsto
(a_{\varphi(0)}, \ldots, a_{\varphi(k)})$
以及 $\alpha \mapsto \alpha \circ \varphi$。
因此 $h = (h_m)_{m \geq 0}$ 显然是所需的单纯集映射。

\begin{lemma}
\label{simplicial-lemma-homotopy}
设 $f^0, f^1 : V \to U$ 为单纯集映射，$n \geq 0$ 为整数。假设：
\begin{enumerate}
\item 当 $i < n$ 时，映射
$f^j_i : V_i \to U_i$（$j = 0, 1$）相等；
\item 典范态射 $U \to \text{cosk}_n \text{sk}_n U$ 是同构；
\item 典范态射 $V \to \text{cosk}_n \text{sk}_n V$ 是同构。
\end{enumerate}
则 $f^0$ 同伦于 $f^1$。
\end{lemma}

\begin{proof}[第一个证明]
设 $W$ 为 $n$-截断单纯集，其中当 $i < n$ 时
$W_i = U_i$，而 $W_n = U_n / \sim$；这里 $\sim$
是由对 $y \in V_n$ 的关系 $f^0(y) \sim f^1(y)$ 生成的等价关系。
这是有意义的，因为与 $\varphi : [i] \to [n]$（$i < n$）
对应的态射 $U(\varphi) : U_n \to U_i$ 经商映射
$U_n \to W_n$ 分解：这是由于 $f^0$ 和 $f^1$
为单纯集态射，并且在次数 $< n$ 中相等。
接着，取 $\text{cosk}_n W$，把 $W$ 提升为单纯集。
由引理 \ref{simplicial-lemma-section}，态射 $g : U \to W$
是平凡 Kan 纤维化。根据构造，
$g \circ f^0 = g \circ f^1$；以 $f : V \to W$ 表示该态射。
考虑图
$$
\xymatrix{
\partial \Delta[1] \times V \ar[rr]_{f^0, f^1} \ar[d] & & U \ar[d] \\
\Delta[1] \times V \ar[rr]^f \ar@{-->}[rru] & & W
}
$$
由引理 \ref{simplicial-lemma-trivial-kan}，虚线箭头存在，证明完成。
\end{proof}

\begin{proof}[第二个证明]
必须构造单纯集态射
$h : V \times \Delta[1] \to U$，它与 $e_i$ 复合后恢复 $f^i$。
$n = 0$ 的情形已在引理之前处理。因此可假设 $n \geq 1$。
映射
$\Delta[1] \to \text{cosk}_1 \text{sk}_1 \Delta[1]$
是同构；见引理 \ref{simplicial-lemma-simplex-cosk}。
由引理 \ref{simplicial-lemma-cosk-up}，因 $n \geq 1$，
$\Delta[1] \to \text{cosk}_n \text{sk}_n \Delta[1]$
也是同构。因此，由引理 \ref{simplicial-lemma-cosk-product}，
$V \times \Delta[1] \to
\text{cosk}_n \text{sk}_n (V \times \Delta[1])$
也是同构。换言之，为了构造同伦，只须构造合适的
$n$-截断单纯集态射
$h : \text{sk}_n V \times \text{sk}_n \Delta[1] \to \text{sk}_n U$。

\medskip\noindent
对 $k = 0, \ldots, n - 1$，以公式
$h_k(v, \alpha) = f^0(v) = f^1(v)$ 定义 $h_k$。
映射 $h_n : V_n \times \Mor_{\Delta}([n], [1]) \to U_n$
如下定义。取 $v \in V_n$ 与 $\alpha : [n] \to [1]$：
\begin{enumerate}
\item 若 $\Im(\alpha) = \{0\}$，令 $h_n(v, \alpha) = f^0(v)$；
\item 若 $\Im(\alpha) = \{0, 1\}$，令 $h_n(v, \alpha) = f^0(v)$；
\item 若 $\Im(\alpha) = \{1\}$，令 $h_n(v, \alpha) = f^1(v)$。
\end{enumerate}
设 $\varphi : [k] \to [l]$ 为 $\Delta_{\leq n}$ 的态射。
我们将证明图
$$
\xymatrix{
V_{l} \times \Mor([l], [1]) \ar[r] \ar[d] & U_{l} \ar[d] \\
V_{k} \times \Mor([k], [1]) \ar[r] & U_{k}
}
$$
交换。取 $v \in V_{l}$ 和 $\alpha : [l] \to [1]$。
交换性意味着
$$
h_k(V(\varphi)(v), \alpha \circ \varphi)
=
U(\varphi)(h_l(v, \alpha)).
$$
几乎在所有情形下，这都由
$h_k(V(\varphi)(v), \alpha \circ \varphi) =
f^0(V(\varphi)(v))$ 和
$U(\varphi)(h_l(v, \alpha)) = U(\varphi)(f^0(v))$，
结合 $f^0$ 为单纯集态射这一事实得出。
唯一不满足这一描述的情形是：
或者 (A) $\Im(\alpha) = \{1\}$ 且 $l = n$，
或者 (B) $\Im(\alpha \circ \varphi) = \{1\}$ 且 $k = n$。
还注意，对 $V$ 的任意退化 $n$-单形，必有
$f^0(v) = f^1(v)$。因此，可把上述情形进一步缩小为：
(A) $\Im(\alpha) = \{1\}$、$l = n$ 且 $v$ 非退化；
以及 (B) $\Im(\alpha \circ \varphi) = \{1\}$、
$k = n$ 且 $V(\varphi)(v)$ 非退化。

\medskip\noindent
在情形 (A) 中，也有
$\Im(\alpha \circ \varphi) = \{1\}$。
因此，不仅 $h_l(v, \alpha) = f^1(v)$，而且
$h_k(V(\varphi)(v), \alpha \circ \varphi) =
f^1(V(\varphi)(v))$。由于 $f^1$ 是单纯集态射，该关系成立。

\medskip\noindent
在情形 (B) 中，断定 $l = k = n$ 且 $\varphi$ 为双射，
否则 $V(\varphi)(v)$ 是退化的。因此
$\varphi = \text{id}_{[n]}$，这是平凡情形。
\end{proof}

\begin{lemma}
\label{simplicial-lemma-cosk-minus-one-equivalence}
设 $A$、$B$ 为集合，并设 $f : A \to B$ 为映射。
% P02-E0161: frozen source says “Let A, B be sets, and that f ...”.
考虑第 $n$ 个单形为
$$
A \times_B A \times_B \ldots \times_B A\ (n + 1 \text{ 可分解）}.
$$
的单纯集 $U$；见例 \ref{simplicial-example-fibre-products-simplicial-object}。
若 $f$ 满射，则态射 $U \to B$ 是平凡 Kan 纤维化；
这里 $B$ 表示取值为 $B$ 的常值单纯集。
\end{lemma}

\begin{proof}
注意，$U$ 嵌入笛卡尔方块
$$
\xymatrix{
U \ar[d] \ar[r] & \text{cosk}_0(A) \ar[d] \\
B \ar[r] & \text{cosk}_0(B)
}
$$
由引理 \ref{simplicial-lemma-section}，右侧竖直箭头是平凡 Kan 纤维化；
再由引理 \ref{simplicial-lemma-trivial-kan-base-change}，左侧也如此。
\end{proof}










\section{标准预解式的准备}
\label{simplicial-section-standard-pre}

\noindent
本节材料可见 \cite[Appendix 1]{Godement}
% P02-E0162: frozen introductory sentence lacks terminal punctuation.

\begin{example}
\label{simplicial-example-godement}
设 $Y : \mathcal{C} \to \mathcal{C}$ 为一个范畴到自身的函子，并设给定
函子变换
$$
d : Y \longrightarrow \text{id}_\mathcal{C}
\quad\text{且}\quad
s : Y \longrightarrow Y \circ Y
$$
利用这些变换，我们可以构造一个形似单纯对象的东西。具体地，对
$n \geq 0$，定义
$$
X_n = Y \circ \ldots \circ Y \quad (n + 1\text{ 复合})
$$
注意，对 $n, m \geq 0$，有 $X_{n + m + 1} = X_n \circ X_m$。
接着，对 $n \geq 0$ 和 $0 \leq j \leq n$，采用《范畴论》第
\ref{categories-section-formal-cat-cat} 节中的记号，定义
$$
d^n_j = 1_{X_{j - 1}} \star d \star 1_{X_{n - j - 1}} :
X_n
\to
X_{n - 1}
\quad\text{且}\quad
s^n_j = 1_{X_{j - 1}} \star s \star 1_{X_{n - j - 1}} :
X_n
\to
X_{n + 1}
$$
因此，$d^n_j$（相应地，$s^n_j$）是在定义 $X_n$ 的复合中，
对从左数第 $j$ 个 $Y$ 使用 $d$（相应地，$s$）所得的自然变换。
\end{example}

\begin{lemma}
\label{simplicial-lemma-godement}
在例 \ref{simplicial-example-godement} 中，若
$$
1_Y = (d \star 1_Y) \circ s = (1_Y \star d) \circ s
\quad\text{且}\quad
(s \star 1) \circ s = (1 \star s) \circ s
$$
则 $X = (X_n, d^n_j, s^n_j)$ 是 $\mathcal{C}$ 的自函子范畴中的
单纯对象，且 $d : X_0 = Y \to \text{id}_\mathcal{C}$ 定义一个增广。
\end{lemma}

\begin{proof}
为证明所得对象是单纯对象，须验证引理
\ref{simplicial-lemma-characterize-simplicial-object} 的关系 (1)(a)--(e)。
我们将使用简写
$$
1_a = 1_{X_{a - 1}} = 1_Y \star \ldots \star 1_Y \quad (a\text{ 可分解})
$$
，其中 $a \geq 0$。采用此记号，有
$$
d^n_j = 1_j \star d \star 1_{n - j}
\quad\text{且}\quad
s^n_j = 1_j \star s \star 1_{n - j}
$$
% P02-E0163: frozen source misspells “functors” as “funtors”.
下文反复使用如下规则：对函子变换 $a, a', b, b'$，只要公式中的
$\star$ 与 $\circ$ 复合有定义，就有
$(a' \circ a) \star (b' \circ b) = (a' \star b') \circ (a \star b)$；
见《范畴论》引理 \ref{categories-lemma-properties-2-cat-cats}。

\medskip\noindent
条件 (1)(a) 总是成立（不需要对 $d$ 和 $s$ 作任何假设）。
事实上，设 $0 \leq i < j \leq n + 1$。我们须证明
$d^n_i \circ d^{n + 1}_j = d^n_{j - 1} \circ d^{n + 1}_i$，即
$$
(1_i \star d \star 1_{n - i}) \circ
(1_j \star d \star 1_{n + 1 - j}) =
(1_{j - 1} \star d \star 1_{n + 1 - j}) \circ
(1_i \star d \star 1_{n + 1 - i})
$$
左端可改写为
\begin{align*}
& (1_i \star d \star 1_{j - i - 1} \star 1_{n + 1 - j}) \circ
(1_i \star 1_1 \star 1_{j - i - 1} \star d \star 1_{n + 1 - j}) \\
& =
1_i \star
\left((d \star 1_{j - i - 1}) \circ
(1_1 \star 1_{j - i - 1} \star d)\right) \star 1_{n + 1 - j} \\
& =
1_i \star d \star 1_{j - i - 1} \star d \star 1_{n + 1 - j}
\end{align*}
第二个等号成立，是因为 $d \circ 1_1 = d$，且
$1_{j - i} \circ (1_{j - i - 1} \star d) = 1_{j - i - 1} \star d$。
对右端作相同计算也得到同一结果。

\medskip\noindent
验证条件 (1)(b)。设 $0 \leq i < j \leq n - 1$。须证明
$d^n_i \circ s^{n - 1}_j = s^{n - 2}_{j - 1} \circ d^{n - 1}_i$，即
$$
(1_i \star d \star 1_{n - i}) \circ
(1_j \star s \star 1_{n - 1 - j}) =
(1_{j - 1} \star s \star 1_{n - 1 - j}) \circ
(1_i \star d \star 1_{n - 1 - i})
$$
通过与情形 (1)(a) 相同的计算，两端都化为
$1_i \star d \star 1_{j - i - 1} \star s \star 1_{n - j - 1}$。

\medskip\noindent
验证条件 (1)(c)。设 $0 \leq j \leq n - 1$。须证明
$\text{id} = d^n_j \circ s^{n - 1}_j = d^n_{j + 1} \circ s^{n - 1}_j$，即
$$
1_n =
(1_j \star d \star 1_{n - j}) \circ
(1_j \star s \star 1_{n - 1 - j}) =
(1_{j + 1} \star d \star 1_{n - j - 1}) \circ
(1_j \star s \star 1_{n - 1 - j})
$$
容易看出，这由引理的第一个假设推出。

\medskip\noindent
验证条件 (1)(d)。设 $0 < j + 1 < i \leq n$。须证明
$d^n_i \circ s^{n - 1}_j = s^{n - 2}_j \circ d^{n - 1}_{i - 1}$，即
$$
(1_i \star d \star 1_{n - i}) \circ
(1_j \star s \star 1_{n - 1 - j}) =
(1_j \star s \star 1_{n - 2 - j}) \circ
(1_{i - 1} \star d \star 1_{n - i})
$$
通过与情形 (1)(a) 相同的计算，两端都化为
$1_j \star s \star 1_{i - j - 2} \star d \star 1_{n - i}$。

\medskip\noindent
验证条件 (1)(e)。设 $0 \leq i \leq j \leq n - 1$。须证明
$s^n_i \circ s^{n - 1}_j = s^n_{j + 1} \circ s^{n - 1}_i$，即
$$
(1_i \star s \star 1_{n - i}) \circ
(1_j \star s \star 1_{n - 1 - j}) =
(1_{j + 1} \star s \star 1_{n - 1 - j}) \circ
(1_i \star s \star 1_{n - 1 - i})
$$
通过与情形 (1)(a) 相同的计算，这化为
$$
(s \star 1_{j - i + 1}) \circ (1_{j - i} \star s) =
(1_{j - i + 1} \star s) \circ (s \star 1_{j - i})
$$
若 $j = i$，这恰为引理的两个假设之一。若 $j > i$，通过与
情形 (1)(a) 相似的计算，左端和右端都化为等式
$s \star 1_{j - i - 1} \star s$。

\medskip\noindent
最后，为证明 $d$ 定义一个增广，须证明
$d \circ (1_1 \star d) = d \circ (d \star 1_1)$；这是成立的，
因为两端都等于 $d \star d$。
\end{proof}

\begin{example}
\label{simplicial-example-godement-functorial}
设 $\mathcal{C}$、$Y$、$d$、$s$ 如例 \ref{simplicial-example-godement}，
并满足引理 \ref{simplicial-lemma-godement} 的等式。给定函子
$F : \mathcal{A} \to \mathcal{C}$ 和
$G : \mathcal{C} \to \mathcal{B}$，我们在从 $\mathcal{A}$
到 $\mathcal{B}$ 的函子范畴中得到单纯对象 $G \circ X \circ F$，
它带有到 $G \circ F$ 的增广。
\end{example}

\begin{lemma}
\label{simplicial-lemma-godement-section}
设 $\mathcal{A}$、$\mathcal{B}$、$\mathcal{C}$、$Y$、$d$、$s$、$F$、
$G$ 如例 \ref{simplicial-example-godement-functorial}。给定函子变换
$h_0 : G \circ F \to G \circ Y \circ F$，使得
$$
1_{G \circ F} = (1_G \star d \star 1_F) \circ h_0
$$
则存在单纯对象态射 $h : G \circ F \to G \circ X \circ F$，
满足 $\epsilon \circ h = \text{id}$；这里
$\epsilon : G \circ X \circ F \to G \circ F$ 是增广。
\end{lemma}

\begin{proof}
以 $u_n : Y = X_0 \to X_n$ 表示单纯对象 $X$ 中与
$\Delta$ 的唯一态射 $[n] \to [0]$ 对应的映射。令
$h_n : G \circ F \to G \circ X_n \circ F$ 等于
$(1_G \star u_n \star 1_F) \circ h_0$。

\medskip\noindent
对任意范畴中的任意单纯对象 $X = (X_n)$，
$u =(u_n) : X_0 \to X$ 是从取值为 $X_0$ 的常值单纯对象到 $X$
的态射。因此，$h$ 是单纯对象态射，因为它是
$1_G \star u \star 1_F$ 与 $h_0$ 的复合。

\medskip\noindent
验证 $\epsilon \circ h = \text{id}$。计算得
$$
\epsilon_n \circ (1_G \star u_n \star 1_F) \circ h_0 =
\epsilon_0 \circ h_0 = \text{id}
$$
第一个等号成立是因为 $\epsilon$ 是单纯对象态射；第二个等号成立是因为
$\epsilon_0 = (1_G \star d \star 1_F)$，故可应用引理陈述中的假设。
\end{proof}

\begin{lemma}
\label{simplicial-lemma-godement-two-maps}
设 $\mathcal{A}$、$\mathcal{B}$、$\mathcal{C}$、$Y$、$d$、$s$、$F$、
$G$ 如例 \ref{simplicial-example-godement-functorial}。设
$F' : \mathcal{A} \to \mathcal{C}$ 和
$G' : \mathcal{C} \to \mathcal{B}$ 为两个函子。
设 $(a_n) : G \circ X \to G' \circ X$ 为单纯对象态射，
且经增广与 $a : G \to G'$ 相容。设
$(b_n) : X \circ F \to X \circ F'$ 为单纯对象态射，
且经增广与 $b : F \to F'$ 相容。则两个映射
$$
a \star (b_n), (a_n) \star b : G \circ X \circ F \to G' \circ X \circ F'
$$
同伦。
\end{lemma}

\begin{proof}
为证明这些态射同伦，我们构造态射
$$
h_{n, i} : G \circ X_n \circ F \to G' \circ X_n \circ F'
$$
，其中 $n \geq 0$ 且 $0 \leq i \leq n + 1$，并使其满足引理
\ref{simplicial-lemma-relations-homotopy} 中所述关系；另见注
\ref{simplicial-remark-homotopy-better}。为满足引理
\ref{simplicial-lemma-relations-homotopy} 的条件 (1)，必须令
$h_{n, 0} = a \star b_n$ 和 $h_{n , n + 1} = a_n \star b$。
因此合理的选择是
$$
h_{n , i} = a_{i - 1} \star b_{n - i}
$$
，其中 $1 \leq i \leq n$。令 $a = a_{-1}$、$b = b_{-1}$，可见
% P02-E0164: frozen source misspells “formula” as “formular”.
上述公式对 $0 \leq i \leq n + 1$ 都成立。

\medskip\noindent
回顾
$$
d^n_j = 1_G \star 1_j \star d \star 1_{n - j} \star 1_F
$$
在 $G \circ X \circ F$ 上成立；这里采用引理
\ref{simplicial-lemma-godement} 证明中引入的记号
$1_a = 1_{Y \circ \ldots \circ Y}$。下文将使用其改写
\begin{align*}
d^n_j
& =
d^j_j \star 1_{n - j} =
d^{j + 1}_j \star 1_{n - j} = \ldots = d^{n - 1}_j \star 1_1 \\
& =
1_j \star d^{n - j}_0 = 1_{j - 1} \star d^{n - j + 1}_1 = \ldots =
1_1 \star d^{n - 1}_{j - 1}
\end{align*}
当然，在 $G' \circ X \circ F'$ 上的 $d^n_j$ 也有类似公式。

\medskip\noindent
验证引理 \ref{simplicial-lemma-relations-homotopy} 的条件 (2)。
设 $i > j$。须证明
$$
d^n_j \circ (a_{i - 1} \star b_{n - i}) =
(a_{i - 2} \star b_{n - i}) \circ d^n_j
$$
因 $i - 1 \geq j$，可用 $d^n_j$ 的一种描述把左端改写为
$$
(d^{i - 1}_j \star 1_{n - i + 1}) \circ (a_{i - 1} \star b_{n - i}) =
(d^{i - 1}_j \circ a_{i - 1}) \star b_{n - i} =
(a_{i - 2} \circ d^{i - 1}_j) \star b_{n - i}
$$
同样，右端变为
$$
(a_{i - 2} \star b_{n - i}) \circ (d^{i - 1}_j \star 1_{n - i + 1}) =
(a_{i - 2} \circ d^{i - 1}_j) \star b_{n - i}
$$
故两端结果相同，条件 (2) 得证。

\medskip\noindent
验证引理 \ref{simplicial-lemma-relations-homotopy} 的条件 (3)。
设 $i \leq j$。须证明
$$
d^n_j \circ (a_{i - 1} \star b_{n - i}) =
(a_{i - 1} \star b_{n - 1 - i}) \circ d^n_j
$$
因 $j \geq i$，可把左端改写为
$$
(1_i \star d^{n - i}_{j - i}) \circ (a_{i - 1} \star b_{n - i}) =
a_{i - 1} \star (b_{n - 1 - i} \circ d^{n - i}_{j - i})
$$
相同的变形表明它与右端相等。

\medskip\noindent
回顾
$$
s^n_j = 1_G \star 1_j \star s \star 1_{n - j} \star 1_F
$$
在 $G \circ X \circ F$ 上成立。下文将使用其改写
\begin{align*}
s^n_j
& =
s^j_j \star 1_{n - j} = s^{j + 1}_j \star 1_{n - j - 1} = \ldots =
s^{n - 1}_j \star 1_1 \\
& =
1_j \star s^{n - j}_0 = 1_{j - 1} \star s^{n - j + 1}_1 = \ldots =
1_1 \star s^{n - 1}_{j - 1}
\end{align*}
当然，在 $G' \circ X \circ F'$ 上的 $s^n_j$ 也有类似公式。

\medskip\noindent
验证引理 \ref{simplicial-lemma-relations-homotopy} 的条件 (4)。
设 $i > j$。须证明
$$
s^n_j \circ (a_{i - 1} \star b_{n - i}) =
(a_i \star b_{n - i}) \circ s^n_j
$$
因 $i - 1 \geq j$，可把左端改写为
$$
(s^{i - 1}_j \star 1_{n - i + 1}) \circ (a_{i - 1} \star b_{n - i}) =
(s^{i - 1}_j \circ a_{i - 1}) \star b_{n - i} =
(a_i \circ s^{i - 1}_j) \star b_{n - i}
$$
同样，右端变为
$$
(a_i \star b_{n - i}) \circ (s^{i - 1}_j \star 1_{n - i + 1}) =
(a_i \circ s^{i - 1}_j) \star b_{n - i}
$$
，正是所需等式。

\medskip\noindent
验证引理 \ref{simplicial-lemma-relations-homotopy} 的条件 (5)。
设 $i \leq j$。须证明
$$
s^n_j \circ (a_{i - 1} \star b_{n - i}) =
(a_{i - 1} \star b_{n + 1 - i}) \circ s^n_j
$$
与上面的论证完全相同，两端都化为
$a_{i - 1} \star (s^{n - i}_{j - i} \circ b_{n - i}) =
a_{i - 1} \star (b_{n + 1 - i} \circ s^{n - i}_{j - i})$，
故等式成立。
\end{proof}

\begin{lemma}
\label{simplicial-lemma-godement-before-after}
设 $\mathcal{C}$、$Y$、$d$、$s$ 如例 \ref{simplicial-example-godement}，
并满足引理 \ref{simplicial-lemma-godement} 的等式。设
$f : \text{id}_\mathcal{C} \to \text{id}_\mathcal{C}$ 为恒等函子的
自同态。则 $f \star 1_X, 1_X \star f : X \to X$ 是单纯对象映射，
且经增广 $\epsilon : X \to \text{id}_\mathcal{C}$ 与 $f$ 相容。
此外，$f \star 1_X$ 与 $1_X \star f$ 同伦。
\end{lemma}

\begin{proof}
映射 $f \star 1_X$ 的分量为
$$
X_n = \text{id}_\mathcal{C} \circ X_n
\xrightarrow{f \star 1_{X_n}}
\text{id}_\mathcal{C} \circ X_n = X_n
$$
。对 $\mathcal{C}$ 的自函子之间的变换 $a : F \to G$，有
$a \circ (f \star 1_F) = f \star a = (f \star 1_G) \circ a$。
因此 $f \star 1_X$ 确为单纯对象态射；$1_X \star f$ 的情形相同。

\medskip\noindent
为证明这些态射同伦，我们对 $n \geq 0$ 和
$0 \leq i \leq n + 1$ 构造态射 $h_{n, i} : X_n \to X_n$，
使其满足引理 \ref{simplicial-lemma-relations-homotopy} 中所述关系；另见注
\ref{simplicial-remark-homotopy-better}。可以取
$$
h_{n, i} = 1_i \star f \star 1_{n + 1 - i}
$$
，其中 $1_i$ 是引理 \ref{simplicial-lemma-godement} 证明中
$Y \circ \ldots \circ Y$ 上的恒等变换。我们有
$h_{n, 0} = f \star 1_{X_n}$ 和 $h_{n, n + 1} = 1_{X_n} \star f$，
这验证了第一个条件。验证其余条件时，我们使用引理
\ref{simplicial-lemma-godement-two-maps} 证明中关于映射 $d^n_j$ 与 $s^n_j$ 的说明。

\medskip\noindent
验证引理 \ref{simplicial-lemma-relations-homotopy} 的条件 (2)。
设 $i > j$。须证明
$$
d^n_j \circ (1_i \star f \star 1_{n + 1 - i}) =
(1_{i - 1} \star f \star 1_{n + 1 - i}) \circ d^n_j
$$
因 $i - 1 \geq j$，可用 $d^n_j$ 的一种描述把左端改写为
$$
(d^{i - 1}_j \star 1_{n - i + 1}) \circ (1_i \star f \star 1_{n + 1 - i}) =
d^{i - 1}_j \star f \star 1_{n + 1 - i}
$$
同样，右端变为
$$
(1_{i - 1} \star f \star 1_{n + 1 - i}) \circ
(d^{i - 1}_j \star 1_{n - i + 1}) =
d^{i - 1}_j \star f \star 1_{n + 1 - i}
$$
故两端结果相同，条件 (2) 得证。

\medskip\noindent
引理 \ref{simplicial-lemma-relations-homotopy} 的条件 (3)、(4)、(5)
可完全同样地验证，采用引理 \ref{simplicial-lemma-godement-two-maps}
证明中的策略即可。细节从略\footnote{当
$f$ 可逆时，只须证明 $(a_n) = 1_X$ 与
$(b_n) = f^{-1} \star 1_X \star f$ 同伦。
但这由引理 \ref{simplicial-lemma-godement-two-maps} 得出，
因为此时 $a = b = 1_{\text{id}_\mathcal{C}}$。}。
\end{proof}






\section{标准预解式}
\label{simplicial-section-standard}

\noindent
本节部分材料可见 \cite[Appendix 1]{Godement} 和
\cite[I 1.5]{cotangent}。

\begin{situation}
\label{simplicial-situation-adjoint-functors}
设 $\mathcal{A}$、$\mathcal{S}$ 为范畴，且
$V : \mathcal{A} \to \mathcal{S}$ 是具有左伴随
$U : \mathcal{S} \to \mathcal{A}$ 的函子。
\end{situation}

\noindent
在这一极一般的情形下，我们将在从 $\mathcal{A}$ 到
$\mathcal{A}$ 的函子范畴中构造一个单纯对象 $X$。
建议在阅读本节正文之前，先看后面给出的例子。

\medskip\noindent
构造中将使用《范畴论》第 \ref{categories-section-formal-cat-cat}
节所定义的水平复合。《范畴论》第
\ref{categories-section-adjoint} 节中伴随态射\footnote{不能以
$\epsilon$ 表示伴随的余单位，因为我们要用
$\epsilon$ 表示单纯对象的增广。}的定义为
$$
d : U \circ V \to \text{id}_\mathcal{A} \quad (\text{余单位})
\quad\text{且}\quad
\eta : \text{id}_\mathcal{S} \to V \circ U \quad (\text{单位})
$$
，并说明复合
\begin{equation}
\label{simplicial-equation-composition}
V \xrightarrow{\eta \star 1_V} V \circ U \circ V \xrightarrow{1_V \star d} V
\quad\text{且}\quad
U \xrightarrow{1_U \star \eta} U \circ V \circ U \xrightarrow{d \star 1_U} U
\end{equation}
都是恒等态射。这里，为定义态射 $\eta \star 1_V$，我们默认把
$V$ 与 $\text{id}_\mathcal{S} \circ V$ 识别，而
$1_V$ 表示 $\text{id}_V : V \to V$。下文将反复使用这些记号和关系。
对 $n \geq 0$，令
$$
X_n = (U \circ V)^{\circ (n + 1)} =
U \circ V \circ U \circ \ldots \circ U \circ V
$$
。换言之，$X_n$ 是 $U \circ V$ 与自身的 $(n + 1)$ 重复合。
还令 $X_{-1} = \text{id}_\mathcal{A}$。对所有 $n, m \geq -1$，
有 $X_{n + m + 1} = X_n \circ X_m$。我们将通过构造函子态射
$$
d^n_j : X_n \to X_{n - 1},\quad s^n_j : X_n \to X_{n + 1}
$$
，使其满足引理 \ref{simplicial-lemma-relations-face-degeneracy} 中所列关系，
从而赋予这个函子序列以 $\text{Fun}(\mathcal{A}, \mathcal{A})$
中单纯对象的结构。具体地，令
$$
d^n_j = 1_{X_{j - 1}} \star d \star 1_{X_{n - j - 1}}
\quad\text{且}\quad
s^n_j = 1_{X_{j - 1} \circ U} \star \eta \star 1_{V \circ X_{n - j - 1}}
$$
最后，记 $\epsilon_0 = d : X_0 \to X_{-1}$。

\begin{lemma}
\label{simplicial-lemma-standard-simplicial}
在情形 \ref{simplicial-situation-adjoint-functors} 中，
系统 $X = (X_n, d^n_j, s^n_j)$ 是
$\text{Fun}(\mathcal{A}, \mathcal{A})$ 中的单纯对象，
且 $\epsilon_0$ 定义从 $X$ 到取值为
$X_{-1} = \text{id}_\mathcal{A}$ 的常值单纯对象的增广 $\epsilon$。
\end{lemma}

\begin{proof}
考虑 $Y = U \circ V : \mathcal{A} \to \mathcal{A}$。
我们已有变换 $d : Y = U \circ V \to \text{id}_\mathcal{A}$。记
$$
s = 1_U \star \eta \star 1_V :
Y = U \circ \text{id}_\mathcal{S} \circ V
\longrightarrow
U \circ V \circ U \circ V = Y \circ Y
$$
这使我们处于例 \ref{simplicial-example-godement} 的情形。由公式立即可见，
上面构造的 $X, d^n_i, s^n_i$ 与从 $Y, d, s$ 按例
\ref{simplicial-example-godement} 构造的
% P02-E0165: frozen source repeats “X, s^n_i, s^n_i” where the first s-family should be d.
$X, s^n_i, s^n_i$ 一致。因此，由引理 \ref{simplicial-lemma-godement}，
只须证明
$$
1_Y = (d \star 1_Y) \circ s = (1_Y \star d) \circ s
\quad\text{且}\quad
(s \star 1) \circ s = (1 \star s) \circ s
$$
第一个等号转化为
$$
1_U \star 1_V = (d \star 1_U \star 1_V) \circ (1_U \star \eta \star 1_V)
$$
；只要有 $1_U = (d \star 1_U) \circ (1_U \star \eta)$，它便成立，
而后者由 (\ref{simplicial-equation-composition}) 成立。第二个等号同理。
对最后一个等式，须证明
$$
(1_U \star \eta \star 1_V \star 1_U \star 1_V) \circ
(1_U \star \eta \star 1_V) =
(1_U \star 1_V \star 1_U \star \eta \star 1_V) \circ
(1_U \star \eta \star 1_V)
$$
为此，只须证明
$(\eta \star 1_V \star 1_U) \circ \eta = (1_V \star 1_U \star \eta) \circ \eta$；
这是成立的，因为两端都与 $\eta \star \eta$ 相同。
\end{proof}

\noindent
阅读下一个引理的证明之前，建议先看例
\ref{simplicial-example-polynomial-algebra-homotopy} 中讨论的例子，以理解该引理的目的。

\begin{lemma}
\label{simplicial-lemma-standard-simplicial-homotopy}
在情形 \ref{simplicial-situation-adjoint-functors} 中，映射
$$
1_V \star \epsilon : V \circ X \to V,
\quad\text{且}\quad
\epsilon \star 1_U : X \circ U \to U
$$
都是同伦等价。
\end{lemma}

\begin{proof}
如引理 \ref{simplicial-lemma-standard-simplicial} 的证明，令
$Y = U \circ V$，于是处于例 \ref{simplicial-example-godement} 的情形。

\medskip\noindent
证明第一个同伦等价。由引理 \ref{simplicial-lemma-godement-section}，
为构造映射 $h : V \to V \circ X$，使它成为
$1_V \star \epsilon$ 的右逆，只须构造映射
$h_0 : V \to V \circ Y = V \circ U \circ V$，使
$1_V = (1_V \star d) \circ h_0$。当然，取
$h_0 = \eta \star 1_V$，等式由 (\ref{simplicial-equation-composition}) 成立。
为完成证明，须说明两个映射
$$
(1_V \star \epsilon) \circ h, 1_V \star \text{id}_X :
V \circ X \longrightarrow V \circ X
$$
同伦。这立即由引理 \ref{simplicial-lemma-godement-two-maps} 得出（取
$G = G' = V$ 且 $F = F' = \text{id}_\mathcal{S}$）。

\medskip\noindent
证明第二个同伦等价。由引理 \ref{simplicial-lemma-godement-section}，
为构造映射 $h : U \to X \circ U$，使它成为
$\epsilon \star 1_U$ 的右逆，只须构造映射
$h_0 : U \to Y \circ U = U \circ V \circ U$，使
$1_U = (d \star 1_U) \circ h_0$。当然，取
$h_0 = 1_U \star \eta$，等式由 (\ref{simplicial-equation-composition}) 成立。
为完成证明，须说明两个映射
$$
(\epsilon \star 1_U) \circ h, \text{id}_X \star 1_U :
X \circ U \longrightarrow X \circ U
$$
同伦。这立即由引理 \ref{simplicial-lemma-godement-two-maps} 得出（取
$G = G' = \text{id}_\mathcal{A}$ 且 $F = F' = U$）。
\end{proof}


\begin{example}
\label{simplicial-example-module}
设 $R$ 为环。作为上述构造的一个例子，可取遗忘函子
$i : \text{Mod}_R \to \textit{Sets}$，以及把集合 $E$ 送到
由 $E$ 生成的自由 $R$-模 $R[E]$ 的函子
$F : \textit{Sets} \to \text{Mod}_R$。对 $R$-模 $M$，
单纯 $R$-模 $X(M)$ 具有如下形状：
$$
X(M) =
\left( \ldots
\xymatrix{
R[R[R[M]]]
\ar@<2ex>[r]
\ar@<0ex>[r]
\ar@<-2ex>[r]
&
R[R[M]]
\ar@<1ex>[r]
\ar@<-1ex>[r]
\ar@<1ex>[l]
\ar@<-1ex>[l]
&
R[M]
\ar@<0ex>[l]
}
\right)
$$
，并带有到 $M$ 的增广。我们还将证明该增广在集合层面为同伦等价。
由引理 \ref{simplicial-lemma-trivial-kan-homotopy}、
\ref{simplicial-lemma-homotopy-equivalence} 和
\ref{simplicial-lemma-qis-simplicial-abelian-groups}，这等价于要求
% P02-E0166: frozen source says “cohomology group” for the homologically indexed chain complex here and again in Example polynomial-algebra.
$M$ 是与单纯模 $X(M)$ 伴随的链复形唯一的非零上同调群。
\end{example}

\begin{example}
\label{simplicial-example-polynomial-algebra}
设 $A$ 为环，$\textit{Alg}_A$ 为交换 $A$-代数的范畴。
作为上述构造的一个例子，可取遗忘函子
$i : \textit{Alg}_A \to \textit{Sets}$，以及把集合 $E$ 送到
$A$ 上由 $E$ 生成的多项式代数 $A[E]$ 的函子
$F : \textit{Sets} \to \textit{Alg}_A$。
（对于本例与例 \ref{simplicial-example-module} 之间的记号重叠，我们深表歉意。）
对 $A$-代数 $B$，单纯 $A$-代数 $X(B)$ 具有如下形状：
$$
X(B) = \left( \ldots
\xymatrix{
A[A[A[B]]]
\ar@<2ex>[r]
\ar@<0ex>[r]
\ar@<-2ex>[r]
&
A[A[B]]
\ar@<1ex>[r]
\ar@<-1ex>[r]
\ar@<1ex>[l]
\ar@<-1ex>[l]
&
A[B]
\ar@<0ex>[l]
}
\right)
$$
，并带有到 $B$ 的增广。我们还将证明该增广在集合层面为同伦等价。
由引理 \ref{simplicial-lemma-trivial-kan-homotopy}、
\ref{simplicial-lemma-homotopy-equivalence} 和
\ref{simplicial-lemma-qis-simplicial-abelian-groups}，这等价于要求
$B$ 是把 $X(B)$ 视为单纯 $A$-模后与之伴随的
$A$-模链复形唯一的非零上同调群。
\end{example}

\begin{example}
\label{simplicial-example-module-maps}
在例 \ref{simplicial-example-module} 中，
$X_n(M) = R[R[\ldots [M]\ldots]]$ 有 $n + 1$ 层方括号。
我们用 $X_1(M)$ 的一个典型元素
$$
\xi = \sum\nolimits_i r_i\left[\sum\nolimits_j r_{ij}[m_{ij}]\right]
$$
描述上面构造的映射。映射
$d_0, d_1 : R[R[M]] \to R[M]$ 由
$$
d_0(\xi) = \sum\nolimits_{i, j} r_ir_{ij}[m_{ij}]
\quad\text{且}\quad
d_1(\xi) = \sum\nolimits_i r_i\left[\sum\nolimits_j r_{ij}m_{ij}\right].
$$
给出。映射 $s_0, s_1 : R[R[M]] \to R[R[R[M]]]$ 由
$$
s_0(\xi) =
\sum\nolimits_i r_i\left[\left[\sum\nolimits_j r_{ij}[m_{ij}]\right]\right]
\quad\text{且}\quad
s_1(\xi) = \sum\nolimits_i r_i\left[\sum\nolimits_j r_{ij}[[m_{ij}]]\right].
$$
给出。
\end{example}

\begin{example}
\label{simplicial-example-polynomial-algebra-maps}
在例 \ref{simplicial-example-polynomial-algebra} 中，
$X_n(B) = A[A[\ldots [B]\ldots]]$ 有 $n + 1$ 层方括号。
我们用一个典型元素
$$
\xi = \sum\nolimits_i a_i [x_{i, 1}] \ldots [x_{i, m_i}] \in A[A[B]] = X_1(B)
$$
描述上面构造的映射，其中对每个 $i, j$，可写
$$
x_{i, j} = \sum a_{i, j, k} [b_{i, j, k, 1}] \ldots [b_{i, j, k, n_{i, j, k}}]
\in A[B]
$$
显然，这个记号极其可怕！为简化记号，也为了看清 $A$-代数映射
$d_0, d_1 : A[A[B]] \to A[B]$ 的作用，只须考察变量 $[x]$；
这里
$$
x = \sum a_k [b_{k, 1}] \ldots [b_{k, n_k}] \in A[B]
$$
是一般元素。对这些变量，有
$$
d_0([x]) = x = \sum a_k [b_{k, 1}] \ldots [b_{k, n_k}]
\quad\text{且}\quad
d_1([x]) = \left[\sum a_k b_{k, 1} \ldots b_{k, n_k}\right]
$$
映射 $s_0, s_1 : A[A[B]] \to A[A[A[B]]]$ 由
$$
s_0([x]) = \left[\left[\sum a_k [b_{k, 1}] \ldots [b_{k, n_k}]\right]\right]
\quad\text{且}\quad
s_1([x]) = \left[\sum a_k [[b_{k, 1}]] \ldots [[b_{k, n_k}]]\right]
$$
给出。
\end{example}

\begin{example}
\label{simplicial-example-polynomial-algebra-homotopy}
回到例 \ref{simplicial-example-polynomial-algebra} 中讨论的例子。
引理 \ref{simplicial-lemma-standard-simplicial-homotopy} 表明，对任意环映射
$A \to B$，单纯环映射
$$
\xymatrix{
A[A[A[B]]] \ar[d]
\ar@<2ex>[r]
\ar@<0ex>[r]
\ar@<-2ex>[r]
&
A[A[B]] \ar[d]
\ar@<1ex>[r]
\ar@<-1ex>[r]
\ar@<1ex>[l]
\ar@<-1ex>[l]
&
A[B] \ar[d]
\ar@<0ex>[l] \\
B
\ar@<2ex>[r]
\ar@<0ex>[r]
\ar@<-2ex>[r]
&
B
\ar@<1ex>[r]
\ar@<-1ex>[r]
\ar@<1ex>[l]
\ar@<-1ex>[l]
&
B
\ar@<0ex>[l]
}
$$
在底层单纯集上为同伦等价。此外，引理
\ref{simplicial-lemma-standard-simplicial-homotopy} 中构造的逆映射在次数
$n$ 上由
$$
b \longmapsto [\ldots[b]\ldots]
$$
给出，记号含义显然。反过来，该引理还说明，对每个集合 $E$，
存在环的同伦等价
$$
\xymatrix{
A[A[A[A[E]]]] \ar[d]
\ar@<2ex>[r]
\ar@<0ex>[r]
\ar@<-2ex>[r]
&
A[A[A[E]]] \ar[d]
\ar@<1ex>[r]
\ar@<-1ex>[r]
\ar@<1ex>[l]
\ar@<-1ex>[l]
&
A[A[E]] \ar[d]
\ar@<0ex>[l] \\
A[E]
\ar@<2ex>[r]
\ar@<0ex>[r]
\ar@<-2ex>[r]
&
A[E]
\ar@<1ex>[r]
\ar@<-1ex>[r]
\ar@<1ex>[l]
\ar@<-1ex>[l]
&
A[E]
\ar@<0ex>[l]
}
$$
。引理中构造的逆映射在次数 $n$ 上由环映射
$$
\sum a_{e_1,\ldots,e_p}[e_1][e_2] \ldots [e_p]
\longmapsto
\sum a_{e_1,\ldots,e_p}[\ldots[e_1]\ldots][\ldots[e_2]\ldots]
\ldots [\ldots[e_p]\ldots]
$$
给出（记号含义显然）。
\end{example}
% Generated mechanically from the frozen more-algebra.tex target.
% Do not edit this generated file; edit the bound locale source instead.

% OK, start here.
%

\setcounter{chapter}{14}
\chapter{代数进阶}



\phantomsection
\label{more-algebra-section-phantom}




\section{引言}
\label{more-algebra-section-introduction}

\noindent
本章证明交换代数中的若干结果；它们比第一章交换代数中的结果更为深入，见
Algebra, Section \ref{algebra-section-introduction}。
可参见 \cite{MatCA}。












\section{给读者的建议}
\label{more-algebra-section-advice}

\noindent
与交换代数一章相比，本章各节的独立性更强。从
Section \ref{more-algebra-section-derived-modules} 起，我们将自由使用环上模的（无界）
导出范畴及其所带的全部工具。





\section{稳定自由模}
\label{more-algebra-section-stably-free}

\noindent
下面给出看来已获普遍接受的定义。

\begin{definition}
\label{more-algebra-definition-stably-free}
设 $R$ 为环。
\begin{enumerate}
\item 称两个 $R$-模 $M$、$N$ {\it 稳定同构}，如果存在
$n, m \geq 0$，使得作为 $R$-模有
$M \oplus R^{\oplus m} \cong N \oplus R^{\oplus n}$。
\item 若模 $M$ 稳定同构于某个自由模，则称其 {\it 稳定自由}。
\end{enumerate}
\end{definition}
\noindent
注意，稳定自由模是投射模。

\begin{lemma}
\label{more-algebra-lemma-exact-category-stably-free}
设 $R$ 为环。设
$0 \to P' \to P \to P'' \to 0$ 是有限投射 $R$-模的短正合序列。
若这 $3$ 个模中的 $2$ 个稳定自由，则第三个也稳定自由。
\end{lemma}

\begin{proof}
由于这些模都是投射模，该序列分裂。因此可选取同构
$P = P' \oplus P''$。若 $P' \oplus R^{\oplus n}$ 与
$P'' \oplus R^{\oplus m}$ 均为自由模，则
$P \oplus R^{\oplus n + m}$ 为自由模。再假设 $P'$ 与 $P$
稳定自由，例如 $P \oplus R^{\oplus n}$ 和
$P' \oplus R^{\oplus m}$ 均为自由模。于是
$$
P'' \oplus (P' \oplus R^{\oplus m}) \oplus R^{\oplus n} =
(P'' \oplus P') \oplus R^{\oplus m} \oplus R^{\oplus n} =
(P \oplus R^{\oplus n}) \oplus R^{\oplus m}
$$
是自由模，故 $P''$ 稳定自由。由对称性也得到剩余一种情形。
\end{proof}
 
\begin{lemma}
\label{more-algebra-lemma-lift-stably-free}
设 $R$ 为环，$I \subset R$ 为理想。假设 $1 + I$ 中的每个元素都是
单位元（换言之，$I$ 包含于 $R$ 的 Jacobson 根）。对每个有限稳定自由
$R/I$-模 $E$，存在有限稳定自由 $R$-模 $M$，使得 $M/IM \cong E$。
\end{lemma}

\begin{proof}
选取 $n$、$m$ 以及一个同构
$E \oplus (R/I)^{\oplus n} \cong (R/I)^{\oplus m}$。
选取 $R$-线性映射 $\varphi : R^{\oplus m} \to R^{\oplus n}$ 和
$\psi : R^{\oplus n} \to R^{\oplus m}$，分别提升投影
$(R/I)^{\oplus m} \to (R/I)^{\oplus n}$ 与嵌入
$(R/I)^{\oplus n} \to (R/I)^{\oplus m}$。
于是 $\varphi \circ \psi : R^{\oplus n} \to R^{\oplus n}$
模 $I$ 后化为恒等映射。由关于 $I$ 的假设，该映射的行列式可逆。因此
$P = \Ker(\varphi)$ 稳定自由，并且提升 $E$。
\end{proof}

\begin{lemma}
\label{more-algebra-lemma-lift-projective}
设 $R$ 为环，$I \subset R$ 为理想。
假设 $1 + I$ 中的每个元素都是单位元
（换言之，$I$ 包含于 $R$ 的 Jacobson 根）。
设 $M$ 为有限平坦 $R$-模，并且 $M/IM$ 是投射 $R/I$-模。
则 $M$ 是有限投射 $R$-模。
\end{lemma}

\begin{proof}
由 Algebra, Lemma \ref{algebra-lemma-finite-flat-local}，
对所有素理想 $\mathfrak p \subset R$，$M_\mathfrak p$ 都是有限自由模。
由 Algebra, Lemma \ref{algebra-lemma-finite-projective}，
只需证明函数
$\rho_M : \mathfrak p \mapsto
\dim_{\kappa(\mathfrak p)} M \otimes_R \kappa(\mathfrak p)$
在 $\Spec(R)$ 上局部常值。由于 $M/IM$ 是有限投射模，这在
$V(I) \subset \Spec(R)$ 上成立。又因 $\Spec(R)$ 的每个闭点都属于
$V(I)$，且每当 $\mathfrak p \subset \mathfrak q \subset R$
为素理想时都有 $\rho_M(\mathfrak p) = \rho_M(\mathfrak q)$，
由一个此处略去的初等拓扑空间论证即可得出结论。
\end{proof}

\noindent
引理 \ref{more-algebra-lemma-lift-stably-free} 中的提升由下一个引理可知在同构意义下唯一。

\begin{lemma}
\label{more-algebra-lemma-isomorphic-finite-projective-lifts}
设 $R$ 为环，$I \subset R$ 为理想。假设 $1 + I$ 中的每个元素都是
单位元（换言之，$I$ 包含于 $R$ 的 Jacobson 根）。若 $P$ 和 $P'$
是有限投射 $R$-模，则
\begin{enumerate}
\item 若 $\varphi : P \to P'$ 是一个 $R$-模映射，并诱导同构
$\overline{\varphi} : P/IP \to P'/IP'$，则 $\varphi$ 是同构；
\item 若 $P/IP \cong P'/IP'$，则 $P \cong P'$。
\end{enumerate}
\end{lemma}

\begin{proof}
证明 (1)。由于 $P'$ 是投射 $R$-模，可选取映射
$P' \to P'/IP' \xrightarrow{\overline{\varphi}^{-1}} P/IP$
的一个提升 $\psi : P' \to P$。由 Nakayama 引理
（Algebra, Lemma \ref{algebra-lemma-NAK}），
$\psi \circ \varphi$ 和 $\varphi \circ \psi$ 都是满射。
因此这些映射都是同构
（Algebra, Lemma \ref{algebra-lemma-fun}），从而 $\varphi$ 是同构。

\medskip\noindent
证明 (2)。选取同构 $P/IP \cong P'/IP'$。由于 $P$ 是投射模，
可选取映射 $P \to P/IP \to P'/IP'$ 的一个提升
$\varphi : P \to P'$。由 (1)，$\varphi$ 是同构。
\end{proof}





\section{关于 Artin--Rees 性质的一点说明}
\label{more-algebra-section-artin-rees}

\noindent
本节部分材料取自 \cite{conrad-dejong}。在
\cite[Section 1]{Eis} 中可以找到含有更多参考文献的一般性讨论。

\medskip\noindent
设 $A$ 为 Noether 环，$I \subset A$ 为理想。给定有限 $A$-模的同态
$f : M \to N$，存在 $c \geq 0$，使得对所有 $n \geq c$ 都有
$$
f(M) \cap I^nN \subset f(I^{n - c}M)
$$
见 Algebra, Lemma \ref{algebra-lemma-map-AR}。在此情形下，我们称
{\it $c$ 对 $f$ 的 Artin--Rees 引理有效}。

\begin{lemma}
\label{more-algebra-lemma-approximate-complex}
设 $A$ 为 Noether 环，$I \subset A$ 为包含于 $A$ 的 Jacobson 根的理想。设
$$
S : L \xrightarrow{f} M \xrightarrow{g} N
\quad\text{且}\quad
S' : L \xrightarrow{f'} M \xrightarrow{g'} N
$$
为如上所示的两个有限 $A$-模复形。假设
\begin{enumerate}
\item $c$ 对 $f$ 和 $g$ 的 Artin--Rees 引理有效；
\item 复形 $S$ 正合；
\item $f' = f \bmod I^{c + 1}M$ 且 $g' = g \bmod I^{c + 1}N$。
\end{enumerate}
则 $c$ 对 $g'$ 的 Artin--Rees 引理有效，并且复形 $S'$ 正合。
\end{lemma}

\begin{proof}
先证明当 $n \geq c$ 时
$g'(M) \cap I^nN \subset g'(I^{n - c}M)$。设 $a$ 是 $M$ 的元素，
满足 $g'(a) \in I^nN$。我们希望用 $f'(L)$ 中的元素调整 $a$，
即不改变 $g'(a)$，使得 $a \in I^{n-c}M$。假设
$a \in I^rM$，其中 $r < n - c$。则
$$
g(a) = g'(a) + (g - g')(a) \in
I^n N + I^{r + c + 1}N = I^{r + c + 1}N.
$$
由 $g$ 的 Artin--Rees 性质，$g(a) \in g(I^{r + 1}M)$。设
$g(a) = g(a_1)$，其中 $a_1 \in I^{r + 1}M$。由于序列 $S$ 正合，
$a - a_1 \in f(L)$。相应地，对某个 $b \in L$，可写成
$a = f(b) + a_1$。于是 $f(b) = a - a_1 \in I^rM$。
$f$ 的 Artin--Rees 性质表明，若 $r \geq c$，则可用
$I^{r - c}L$ 中的元素替换 $b$。于是无论哪种情形，都有
$a = f'(b) + a_2$，其中
$a_2 = (f - f')(b) + a_1 \in I^{r + 1}M$。（确切地说，或是
$c \geq r$，由假设有 $(f - f')(b) \in I^{r + 1}M$；或是
$c < r$ 且 $b \in I^{r - c}$，从而仍有
$(f - f')(b) \in I^{c + 1} I^{r - c} M = I^{r + 1}M$。）
因此可用元素 $f'(b) \in f'(L)$ 调整 $a$，使 $r$ 增加 $1$。

\medskip\noindent
事实上，上述论证表明，对所有 $n \geq c$，都有
$(g')^{-1}(I^nN) \subset f'(L) + I^{n - c}M$。因此 $S'$ 正合，因为
$$
(g')^{-1}(0) = (g')^{-1}(\bigcap I^nN) \subset
\bigcap f'(L) + I^{n - c}M = f'(L)
$$
这里使用了 $I$ 包含于 $A$ 的 Jacobson 根这一事实，见
Algebra, Lemma \ref{algebra-lemma-intersection-powers-ideal-module}。
\end{proof}

\noindent
给定环 $A$ 的理想 $I \subset A$ 及 $A$-模 $M$，令
$$
\text{Gr}_I(M) = \bigoplus I^nM/I^{n + 1}M.
$$
我们把它视为一个分次 $\text{Gr}_I(A)$-模。

\begin{lemma}
\label{more-algebra-lemma-approximate-complex-graded}
假设同引理 \ref{more-algebra-lemma-approximate-complex}。令
$Q = \Coker(g)$ 和 $Q' = \Coker(g')$。则作为分次
$\text{Gr}_I(A)$-模，有
$\text{Gr}_I(Q) \cong \text{Gr}_I(Q')$。
\end{lemma}

\begin{proof}
在次数 $n$，有
$\text{Gr}_I(Q)_n = I^nN/(I^{n + 1}N + g(M) \cap I^nN)$，
$Q'$ 的情形也类似。我们断言
$$
g(M) \cap I^nN \subset I^{n + 1}N + g'(M) \cap I^nN.
$$
由对称性（断言的证明只会用到 $c$ 对 $g$ 有效，而由引理它对 $g'$
也有效），这将推出
$$
I^{n + 1}N + g(M) \cap I^nN = I^{n + 1}N + g'(M) \cap I^nN
$$
从而 $\text{Gr}_I(Q)_n$ 与 $\text{Gr}_I(Q')_n$ 作为 $N$ 的子商相同，
这就证明了引理。注意，当 $n \leq c$ 时，由
$g = g' \bmod I^{c + 1}N$，断言显然成立。若 $n > c$，设
$b \in g(M) \cap I^nN$。将其写为 $b = g(a)$，其中
$a \in I^{n - c}M$。令 $b' = g'(a)$。于是
$b - b' = (g - g')(a) \in I^{n + 1}N$，正合所需。
\end{proof}

\begin{lemma}
\label{more-algebra-lemma-works-flat-extension}
设 $A \to B$ 是 Noether 环的平坦映射，$I \subset A$ 为理想，
$f : M \to N$ 为有限 $A$-模的同态。假设 $c$ 对 $f$ 的
Artin--Rees 引理有效。则对于理想 $IB$，$c$ 对
$f \otimes 1 : M \otimes_A B \to N \otimes_A B$
的 Artin--Rees 引理也有效。
\end{lemma}

\begin{proof}
注意
$$
(f \otimes 1)(M) \cap I^n N \otimes_A B
= (f \otimes 1)\left((f \otimes 1)^{-1}(I^n N \otimes_A B)\right)
$$
另一方面，
\begin{align*}
(f \otimes 1)^{-1}(I^n N \otimes_A B) &
= \Ker(M \otimes_A B \to N \otimes_A B/(I^n N \otimes_A B)) \\
& =
\Ker(M \otimes_A B \to (N/I^nN) \otimes_A B).
\end{align*}
由于 $A \to B$ 平坦，取核与余核同张量 $B$ 可交换，故上式等于
$f^{-1}(I^nN) \otimes_A B$。由假设，$f^{-1}(I^nN)$ 包含于
$\Ker(f) + I^{n - c}M$，引理得证。
\end{proof}








\section{环的纤维积 I}
\label{more-algebra-section-fibre-products-rings}

\noindent
环的纤维积与概形的推出有关。More on Morphisms,
Section \ref{more-morphisms-section-pushouts} 讨论了概形推出的若干情形。

\begin{lemma}
\label{more-algebra-lemma-fibre-product-finite-type}
设 $R$ 为环，$A \to B$ 和 $C \to B$ 为 $R$-代数映射。假设
\begin{enumerate}
\item $R$ 是 Noether 环；
\item $A$、$B$、$C$ 在 $R$ 上为有限型；
\item $A \to B$ 是满射；
\item $B$ 在 $C$ 上有限。
\end{enumerate}
则 $A \times_B C$ 在 $R$ 上为有限型。
\end{lemma}

\begin{proof}
令 $D = A \times_B C$。有行正合的交换图
$$
\xymatrix{
0 &
B \ar[l] &
A \ar[l] &
I \ar[l] &
0 \ar[l] \\
0 &
C \ar[l] \ar[u] &
D \ar[l] \ar[u] &
I \ar[l] \ar[u] &
0 \ar[l]
}
$$
。选取 $y_1, \ldots, y_n \in B$，使其作为 $C$-模生成 $B$；
再选取映到 $y_i$ 的 $x_i \in A$。于是
$1, x_1, \ldots, x_n$ 作为 $D$-模生成 $A$。映射
$D \to A \times C$ 是单射，而环 $A \times C$ 作为 $D$-模有限
（因为它是两个有限 $D$-模 $A$ 和 $C$ 的直和）。故引理由
Artin--Tate 引理
（Algebra, Lemma \ref{algebra-lemma-Artin-Tate}）推出。
\end{proof}

\begin{lemma}
\label{more-algebra-lemma-formal-consequence}
设 $R$ 为 Noether 环，$I$ 为有限集。假设给定笛卡尔图
$$
\xymatrix{
\prod B_i &
\prod A_i \ar[l]^{\prod \varphi_i} \\
Q \ar[u]^{\prod \psi_i} &
P \ar[u] \ar[l]
}
$$
，其中 $\psi_i$ 和 $\varphi_i$ 都是满射，且
$Q$、$A_i$、$B_i$ 在 $R$ 上为有限型。则 $P$ 在 $R$ 上为有限型。
\end{lemma}

\begin{proof}
由引理 \ref{more-algebra-lemma-fibre-product-finite-type} 及对 $I$ 的大小作归纳即得。
确切地说，令 $I = I' \amalg \{i_0\}$，并令 $P'$ 为以 $I'$
代入引理中的图所定义的环。由归纳假设，$P'$ 为有限型。
最后，$P$ 位于纤维积图
$$
\xymatrix{
B_{i_0} &
A_{i_0} \ar[l] \\
P' \ar[u] &
P \ar[u] \ar[l] &
}
$$
中，可对该图应用上述引理。
\end{proof}

\begin{lemma}
\label{more-algebra-lemma-diagram-localize}
假设给定环的笛卡尔图
$$
\xymatrix{
R &
R' \ar[l]^t \\
B \ar[u]_s &
B'\ar[u] \ar[l]
}
$$
，即 $B' = B \times_R R'$。若 $h \in B'$ 对应于
$g \in B$ 与 $f \in R'$，并且 $s(g) = t(f)$，则图
$$
\xymatrix{
R_{s(g)} = R_{t(f)} &
(R')_f \ar[l]^-t \\
B_g \ar[u]_s &
(B')_h \ar[u] \ar[l]
}
$$
也是笛卡尔图。
\end{lemma}

\begin{proof}
等式 $B' = B \times_R R'$ 表明
$$
0 \to B' \to B \oplus R' \xrightarrow{s, -t} R
$$
是 $B'$-模的正合序列。作为 $B'$-模，有
$B_g = B_h$、$R'_f = R'_h$ 和 $R_{s(g)} = R_{t(f)} = R_h$。
由局部化的正合性
（Algebra, Proposition \ref{algebra-proposition-localization-exact}），得到正合序列
$$
0 \to B'_h \to B_g \oplus R'_f \xrightarrow{s, -t} R_{s(g)} = R_{t(f)}
$$
这就证明了引理。
\end{proof}

\noindent
考虑环的交换图
$$
\xymatrix{
R &
R' \ar[l] \\
B \ar[u] &
B' \ar[u] \ar[l]
}
$$
。考虑函子（其中范畴的纤维积按
Categories, Example \ref{categories-example-2-fibre-product-categories}
中的方式构造）
\begin{equation}
\label{more-algebra-equation-modules}
\text{Mod}_{B'} \longrightarrow
\text{Mod}_B \times_{\text{Mod}_R} \text{Mod}_{R'},\quad
L' \longmapsto (L' \otimes_{B'} B, L' \otimes_{B'} R', can)
\end{equation}
，其中 $can$ 是规范等同
$L' \otimes_{B'} B \otimes_B R = L' \otimes_{B'} R' \otimes_{R'} R$。
下文以 $(N, M', \varphi)$ 表示右端的对象，即 $N$ 是 $B$-模，
$M'$ 是 $R'$-模，而
$\varphi : N \otimes_B R \to M' \otimes_{R'} R$ 是同构。

\begin{lemma}
\label{more-algebra-lemma-modules}
给定环的交换图
$$
\xymatrix{
R &
R' \ar[l] \\
B \ar[u] &
B' \ar[u] \ar[l]
}
$$
，函子 (\ref{more-algebra-equation-modules}) 有右伴随，即函子
$$
F : (N, M', \varphi) \longmapsto N \times_\varphi M'
$$
（具体说明见证明）。
\end{lemma}

\begin{proof}
给定范畴
$\text{Mod}_B \times_{\text{Mod}_R} \text{Mod}_{R'}$
的对象 $(N, M', \varphi)$，令
$$
N \times_\varphi M' = \{(n, m') \in N \times M' \mid
\varphi(n \otimes 1) = m' \otimes 1\text{ 于 }M' \otimes_{R'} R\}
$$
并将其视为 $B'$-模。伴随性断言是：对 $B'$-模 $L'$ 和三元组
$(N, M', \varphi)$，有
$$
\Hom_{B'}(L', N \times_\varphi M') =
\Hom_B(L' \otimes_{B'} B, N)
\times_{\Hom_R(L' \otimes_{B'} R, M' \otimes_{R'} R)}
\Hom_{R'}(L' \otimes_{B'} R', M')
$$
由 Algebra, Lemma \ref{algebra-lemma-adjoint-tensor-restrict}，
右端等于
$$
\Hom_{B'}(L', N) \times_{\Hom_{B'}(L', M' \otimes_{R'} R)} \Hom_{B'}(L', M')
$$
因此很清楚，对此纤维积中的一对元素 $(g, f')$，可得到 $B'$-线性映射
$L' \to N \times_\varphi M'$，$l' \mapsto (g(l'), f'(l'))$。
反之，给定 $B'$-线性映射 $g' : L' \to N \times_\varphi M'$，
可令 $g$ 为复合 $L' \to N \times_\varphi M' \to N$，令 $f'$
为复合 $L' \to N \times_\varphi M' \to M'$。
这些构造互为逆，并定义所需同构。
\end{proof}






\section{环的纤维积 II}
\label{more-algebra-section-fibre-products-rings-II}

\noindent
本节讨论如下情形中的纤维积。

\begin{situation}
\label{more-algebra-situation-module-over-fibre-product}
下文考虑环映射
$$
\xymatrix{
B \ar[r] & A & A' \ar[l]
}
$$
，其中假设 $A' \to A$ 是核为 $I$ 的满射。在此情形中，令
$B' = B \times_A A'$，从而得到笛卡尔方块
$$
\xymatrix{
A & A' \ar[l] \\
B \ar[u] & B' \ar[l] \ar[u]
}
$$
。
\end{situation}

\begin{lemma}
\label{more-algebra-lemma-points-of-fibre-product}
在情形 \ref{more-algebra-situation-module-over-fibre-product} 中，作为拓扑空间有
$$
\Spec(B') = \Spec(B) \amalg_{\Spec(A)} \Spec(A')
$$
\end{lemma}

\begin{proof}
由 $B' = B \times_A A'$，得到谱的一个交换方块。由于源是拓扑空间范畴中的
推出（它由 Topology, Section \ref{topology-section-colimits} 存在），
该方块诱导连续映射
$$
can : \Spec(B) \amalg_{\Spec(A)} \Spec(A') \longrightarrow \Spec(B')
$$

\medskip\noindent
为证明映射 $can$ 是满射，设 $\mathfrak q' \subset B'$ 为素理想。
若 $I \subset \mathfrak q'$（此处及下文，我们也把 $I$ 视为 $B'$ 的理想，
而不仅是 $A'$ 的理想），则 $\mathfrak q'$ 对应于 $B$ 的一个素理想，因而
属于像。否则，选取 $h \in I$ 且 $h \not \in \mathfrak q'$。此时
$B_h = A_h = 0$，并且环映射 $B'_h \to A'_h$ 是同构，见
引理 \ref{more-algebra-lemma-diagram-localize}。因此 $\mathfrak q'$ 对应于
不包含 $I$ 的唯一素理想 $\mathfrak p' \subset A'$。

\medskip\noindent
由于 $B' \to B$ 是满射，$can$ 在分支 $\Spec(B)$ 上是单射。
上面已经看到，$\Spec(A') \to \Spec(B')$ 在
$V(I) \subset \Spec(A')$ 的补集上是单射。又因为
$V(I) \subset \Spec(A')$ 恰为 $\Spec(A) \to \Spec(A')$ 的像，
一个平凡的集合论论证表明 $can$ 是单射。

\medskip\noindent
最后还需证明 $can$ 是开映射。为此，注意推出中的开集形如
$V \amalg U'$，其中 $V \subset \Spec(B)$ 与
$U' \subset \Spec(A')$ 为开集，且它们在 $\Spec(A)$ 中的逆像相同。
设 $v \in V$。可选取 $g \in B$，使得 $v \in D(g) \subset V$。
令 $f \in A$ 为其像，并选取映到 $f$ 的 $f' \in A'$。则
$D(f') \cap U' \cap V(I) = D(f') \cap V(I)$。
因此 $V(I) \cap D(f')$ 与 $D(f') \cap (U')^c$ 是
$D(f') = \Spec(A'_{f'})$ 中互不相交的闭子集。对某个理想
$J \subset A'$，写 $(U')^c = V(J)$。由刚证明的不相交性，
映射
$A'_{f'} \to A'_{f'}/IA'_{f'} \times A'_{f'}/JA'_{f'}$
是满射，故可选取 $a'' \in A'_{f'}$，使其在
$A'_{f'}/IA'_{f'}$ 中映到 $1$，在 $A'_{f'}/JA'_{f'}$
中映到零。清除分母后，得到 $a' \in J$，它在 $A$ 中映到 $f^n$。
于是 $D(a'f') \subset U'$。令
$h' = (g^{n + 1}, a'f') \in B'$。由先前引用的引理，
$B'_{h'} = B_{g^{n + 1}} \times_{A_{f^{n + 1}}} A'_{a'f'}$，
故 $D(h')$ 在推出中的逆像是 $v$ 的一个开邻域；换言之，
$V \amalg U'$ 的像包含 $v$ 的像的一个开邻域。对于
$u' \in U'$ 且 $u' \not \in V(I)$ 的情形，结论也成立，其较为容易的证明略去。
\end{proof}

\begin{lemma}
\label{more-algebra-lemma-fibre-product-integral}
在情形 \ref{more-algebra-situation-module-over-fibre-product} 中，若
$B \to A$ 是整映射，则 $B' \to A'$ 是整映射。
\end{lemma}

\begin{proof}
设 $a' \in A'$，其像为 $a \in A$。设
$x^d + b_1 x^{d - 1} + \ldots + b_d$ 是系数在 $B$ 中且以 $a$
为根的首一多项式。选取映到 $b_i \in B$ 的 $b'_i \in B'$
（这是可能的）。于是
$(a')^d + b'_1 (a')^{d - 1} + \ldots + b'_d$
属于 $A' \to A$ 的核。由于
$\Ker(B' \to B) = \Ker(A' \to A)$，可修改 $b'_d$ 的选择，使得
$(a')^d + b'_1 (a')^{d - 1} + \ldots + b'_d = 0$，正合所需。
\end{proof}

\noindent
在情形 \ref{more-algebra-situation-module-over-fibre-product} 中，我们希望借助
$A'$、$A$ 与 $B$ 上的模来理解 $B'$-模。为此考虑函子
（其中范畴的纤维积按
Categories, Example \ref{categories-example-2-fibre-product-categories}
中的方式构造）
\begin{equation}
\label{more-algebra-equation-functor}
\text{Mod}_{B'} \longrightarrow
\text{Mod}_B \times_{\text{Mod}_A} \text{Mod}_{A'},\quad
L' \longmapsto (L' \otimes_{B'} B, L' \otimes_{B'} A', can)
\end{equation}
，其中 $can$ 是规范等同
$L' \otimes_{B'} B \otimes_B A = L' \otimes_{B'} A' \otimes_{A'} A$。
下文以 $(N, M', \varphi)$ 表示右端的对象，即 $N$ 是 $B$-模，
$M'$ 是 $A'$-模，而
$\varphi : N \otimes_B A \to M' \otimes_{A'} A$ 是同构。
不过，把 $\varphi$ 看成 $B$-线性映射
$\varphi : N \to M'/IM'$ 往往更为方便；它诱导同构
$N \otimes_B A \to M' \otimes_{A'} A = M'/IM'$。

\begin{lemma}
\label{more-algebra-lemma-module-over-fibre-product}
在情形 \ref{more-algebra-situation-module-over-fibre-product} 中，
函子 (\ref{more-algebra-equation-functor}) 有右伴随，即函子
$$
F : (N, M', \varphi) \longmapsto N \times_{\varphi, M} M'
$$
其中 $M = M'/IM'$。此外，$F$ 与 (\ref{more-algebra-equation-functor}) 的复合
是 $\text{Mod}_B \times_{\text{Mod}_A} \text{Mod}_{A'}$
上的恒等函子。换言之，令
$N' = N \times_{\varphi, M} M'$，则
$N' \otimes_{B'} B = N$ 且 $N' \otimes_{B'} A' = M'$。
\end{lemma}

\begin{proof}
伴随性断言由更一般的引理 \ref{more-algebra-lemma-modules} 推出。
为证明最后一个断言，回忆
$B' = B \times_A A'$ 和 $N' = N \times_{\varphi, M} M'$，
并把这些等式扩充为
$$
\vcenter{
\xymatrix{
A & A' \ar[l] & I \ar[l] \\
B \ar[u] & B' \ar[l] \ar[u] & J \ar[l] \ar[u]
}
}
\quad\text{且}\quad
\vcenter{
\xymatrix{
M & M' \ar[l] & K \ar[l] \\
N \ar[u]_\varphi & N' \ar[l] \ar[u] & L \ar[l] \ar[u]
}
}
$$
，其中 $I, J, K, L$ 分别是原来两个图中各水平映射的核。
把证明列为一系列观察：
\begin{enumerate}
\item $K = IM'$（见引理的陈述）；
\item $B' \to B$ 是核为 $J$ 的满射，且 $J \to I$ 是双射；
\item $N' \to N$ 是核为 $L$ 的满射，且 $L \to K$ 是双射；
\item $JN' \subset L$；
\item $\Im(N \to M)$ 作为 $A$-模生成 $M$
（因为 $N \otimes_B A = M$）；
\item $\Im(N' \to M')$ 作为 $A'$-模生成 $M'$
（因为模 $K$ 后这一点成立，且 $L$ 同构地映到 $K$）；
\item $JN' = L$（因为由上一项，$L \cong K = I M'$
由形如 $x n'$ 的元素的像生成，其中 $x \in I$ 且 $n' \in N'$）；
\item $N' \otimes_{B'} B = N$（因为 $N = N'/L$、$B = B'/J$
及上一项）；
\item 存在映射 $\gamma : N' \otimes_{B'} A' \to M'$；
\item $\gamma$ 是满射（见上文）；
\item 复合映射 $N' \otimes_{B'} A' \to M' \to M$ 的核
由形如 $l \otimes 1$ 和 $n' \otimes x$ 的元素生成，其中
$l \in K$、$n' \in N'$、$x \in I$（因为由假设
$M = N \otimes_B A$，并且 $N' \to N$ 与 $A' \to A$
分别是核为 $L$ 和 $I$ 的满射）；
\item $N' \otimes_{B'} A'$ 中由形如
$l \otimes 1$ 和 $n' \otimes x$ 的元素生成的子模的任意元素，
其中 $l \in L$、$n' \in N'$、$x \in I$，都可写成某个
$l \in L$ 的 $l \otimes 1$（因为 $J$ 同构地映到 $I$，
在 $N' \otimes_{B'} A'$ 中有
$n' \otimes x = n'x \otimes 1$；类似地，当
$n' \in N'$、$x \in J$、$a' \in A'$ 时，有
$x n' \otimes a' = n' \otimes xa' = n'(xa') \otimes 1$；
该等式成立于 $N' \otimes_{B'} A'$ 中；
再由已经证明的 $JN' = L$，便得断言）；
\item $\gamma$ 的核为零（因为由 (10) 和 (11)，核的每个元素都形如
$l \otimes 1$，其中 $l \in L$，而 $\gamma$ 将它映到
$l \in K \subset M'$）。
\end{enumerate}
证明完成。
\end{proof}

\begin{lemma}
\label{more-algebra-lemma-module-over-fibre-product-bis}
在引理 \ref{more-algebra-lemma-module-over-fibre-product} 的情形中，对 $B'$-模
$L'$，伴随映射
$$
L' \longrightarrow
(L' \otimes_{B'} B) \times_{(L' \otimes_{B'} A)} (L' \otimes_{B'} A')
$$
是满射，但一般不是单射。
\end{lemma}

\begin{proof}
与引理 \ref{more-algebra-lemma-module-over-fibre-product} 的证明相同，令
$J \subset B'$ 为映射 $B' \to B$ 的核。则
$L' \otimes_{B'} B = L'/JL'$。因此，为证明满性，只需说明纤维积中
形如 $(0, z)$ 的元素属于该引理映射的像。
映射 $L' \otimes_{B'} A' \to L' \otimes_{B'} A$ 的核是
$L' \otimes_{B'} I \to L' \otimes_{B'} A'$ 的像。
由于 $B' \to A'$ 诱导的映射 $J \to I$ 是同构，复合
$$
L' \otimes_{B'} J \to L' \to
(L' \otimes_{B'} B) \times_{(L' \otimes_{B'} A)} (L' \otimes_{B'} A')
$$
将 $L' \otimes_{B'} J$ 满射到形如 $(0, z)$ 的元素集合。
为看出此映射一般不是单射，给出一个简单例子。取域 $k$，令
$B' = k[x, y]/(xy)$、$A' = B'/(x)$、$B = B'/(y)$、
$A = B'/(x, y)$ 以及 $L' = B'/(x - y)$。此时 $L'$ 中 $x$
的类非零，但在上面展示的箭头下映到零。
\end{proof}

\begin{lemma}
\label{more-algebra-lemma-surjection-module-over-fibre-product}
在情形 \ref{more-algebra-situation-module-over-fibre-product} 中，设
$(N_1, M'_1, \varphi_1) \to (N_2, M'_2, \varphi_2)$
是 $\text{Mod}_B \times_{\text{Mod}_A} \text{Mod}_{A'}$
中的态射，并且 $N_1 \to N_2$ 和 $M'_1 \to M'_2$ 都是满射。则
$$
N_1 \times_{\varphi_1, M_1} M'_1 \to N_2 \times_{\varphi_2, M_2} M'_2
$$
是满射，其中 $M_1 = M'_1/IM'_1$ 且 $M_2 = M'_2/IM'_2$。
\end{lemma}

\begin{proof}
选取 $(x_2, y_2) \in N_2 \times_{\varphi_2, M_2} M'_2$。
选取映到 $x_2$ 的 $x_1 \in N_1$。由于 $M'_1 \to M_1$ 是满射，
可选取映到 $\varphi_1(x_1)$ 的 $y_1 \in M'_1$。于是
$(x_1, y_1)$ 在 $N_2 \times_{\varphi_2, M_2} M'_2$ 中映到
$(x_2, y'_2)$。因此只需证明形如 $(0, y_2)$ 的元素属于该映射的像。
这里 $y_2 \in IM'_2$。写
$y_2 = \sum t_i y_{2, i}$，其中 $t_i \in I$。选取映到
$y_{2, i}$ 的 $y_{1, i} \in M'_1$。于是
$y_1 = \sum t_iy_{1, i} \in IM'_1$，而元素 $(0, y_1)$ 满足要求。
\end{proof}

\begin{lemma}
\label{more-algebra-lemma-finite-module-over-fibre-product}
设 $A, A', B, B', I, M, M', N, \varphi$ 如引理
\ref{more-algebra-lemma-module-over-fibre-product}。若 $N$ 在 $B$ 上有限，且
$M'$ 在 $A'$ 上有限，则
$N' = N \times_{\varphi, M} M'$ 在 $B'$ 上有限。
\end{lemma}

\begin{proof}
下文将不再说明地使用引理 \ref{more-algebra-lemma-module-over-fibre-product}
的结论。选取 $N$ 在 $B$ 上的一组生成元 $y_1, \ldots, y_r$，
以及 $M'$ 在 $A'$ 上的一组生成元 $x_1, \ldots, x_s$。
利用 $N = N' \otimes_{B'} B$ 以及 $B' \to B$ 为满射，
可选取 $u_1, \ldots, u_r \in N'$，它们在 $N$ 中分别映到
$y_1, \ldots, y_r$。利用 $M' = N' \otimes_{B'} A'$，
可选取 $v_1, \ldots, v_t \in N'$，使得对某些
$a'_{ij} \in A'$ 有 $x_i = \sum v_j \otimes a'_{ij}$。
特别地，$v_j$ 的像 $\overline{v}_j \in M'$ 作为 $A'$-模生成 $M'$。
我们断言 $u_1, \ldots, u_r, v_1, \ldots, v_t$
作为 $B'$-模生成 $N'$。事实上，取 $\xi \in N'$。
先选取 $b'_1, \ldots, b'_r \in B'$，使得 $\xi$ 与
$\sum b'_i u_i$ 在 $N$ 中映到同一个元素。这是可能的，因为
$B' \to B$ 是满射，且 $y_1, \ldots, y_r$ 在 $B$ 上生成 $N$。
差 $\xi - \sum b'_i u_i$ 形如 $(0, \theta)$，其中某个
$\theta$ 属于 $IM'$。设 $\theta$ 为 $\sum t_j\overline{v}_j$，其中
$t_j \in I$。由于 $J = \Ker(B' \to B)$ 同构地映到 $I$，
可选取映到 $t_j$ 的 $s_j \in J \subset B'$。
由 $N' = N \times_{\varphi, M} M'$，得到
$\xi = \sum b'_i u_i + \sum s_j v_j$，正合所需。
\end{proof}

\begin{lemma}
\label{more-algebra-lemma-flat-module-over-fibre-product}
设 $A, A', B, B', I$ 如情形
\ref{more-algebra-situation-module-over-fibre-product}。
\begin{enumerate}
\item 设 $(N, M', \varphi)$ 是
$\text{Mod}_B \times_{\text{Mod}_A} \text{Mod}_{A'}$ 的对象。
若 $M'$ 在 $A'$ 上平坦，且 $N$ 在 $B$ 上平坦，则
$N' = N \times_{\varphi, M} M'$ 在 $B'$ 上平坦。
\item 若 $L'$ 是平坦 $B'$-模，则
$L' = (L \otimes_{B'} B) \times_{(L \otimes_{B'} A)} (L \otimes_{B'} A')$。
\item 平坦 $B'$-模的范畴等价于
$\text{Mod}_B \times_{\text{Mod}_A} \text{Mod}_{A'}$
的满子范畴；该满子范畴由满足 $N$ 在 $B$ 上平坦且 $M'$
在 $A'$ 上平坦的三元组 $(N, M', \varphi)$ 组成。
\end{enumerate}
\end{lemma}

\begin{proof}
证明中将不再说明地使用引理
\ref{more-algebra-lemma-module-over-fibre-product}。

\medskip\noindent
证明 (1)。令 $J = \Ker(B' \to B)$。这是 $B'$ 的理想，并同构地映到
$I = \Ker(A' \to A)$。设 $\mathfrak b' \subset B'$ 为理想。
必须证明 $\mathfrak b' \otimes_{B'} N' \to N'$ 是单射，见
Algebra, Lemma \ref{algebra-lemma-flat}。我们知道
$$
\mathfrak b'/(\mathfrak b' \cap J) \otimes_{B'} N' =
\mathfrak b'/(\mathfrak b' \cap J) \otimes_B N \to N
$$
是单射，因为 $N$ 在 $B$ 上平坦。由于
$\mathfrak b' \cap J \to \mathfrak b' \to
\mathfrak b'/(\mathfrak b' \cap J) \to 0$ 正合，可知只需证明
$(\mathfrak b' \cap J) \otimes_{B'} N' \to N'$ 是单射。
因此可假设 $\mathfrak b' \subset J$。其次，由于
$J \to I$ 是同构，有
$$
J \otimes_{B'} N' =
I \otimes_{A'} A' \otimes_{B'} N' =
I \otimes_{A'} M'
$$
而它单射地映入 $M'$，因为 $M'$ 是平坦 $A'$-模。因此
$J \otimes_{B'} N' \to N'$ 是单射，从而
$\text{Tor}_1^{B'}(B'/J, N') = 0$，见
Algebra, Remark \ref{algebra-remark-Tor-ring-mod-ideal}。
于是可将 Algebra, Lemma
\ref{algebra-lemma-what-does-it-mean} 应用于 $B'$ 上的 $N'$
及理想 $J$。
回到理想 $\mathfrak b' \subset J$，令
$\mathfrak b' \subset \mathfrak b'' \subset J$ 为其像在 $I$ 中
是 $I$ 的一个 $A'$-子模的最小理想。换言之，若把 $J = I$ 视为
$A'$-模，则 $\mathfrak b'' = A' \mathfrak b'$。于是
$\mathfrak b''/\mathfrak b'$ 被 $J$ 零化，并且由应用上述引理所得的
$\text{Tor}_1^{B'}(\mathfrak b''/\mathfrak b', N')$ 的消失，得到短正合序列
$$
0 \to \mathfrak b' \otimes_{B'} N' \to
\mathfrak b'' \otimes_{B'} N' \to
\mathfrak b''/\mathfrak b' \otimes_{B'} N' \to 0
$$
因此可用 $\mathfrak b''$ 代替 $\mathfrak b'$。特别地，可假设
$\mathfrak b'$ 是 $A'$-模，并映到 $A'$ 的一个理想。此时
$$
\mathfrak b' \otimes_{B'} N' =
\mathfrak b' \otimes_{A'} A' \otimes_{B'} N' =
\mathfrak b' \otimes_{A'} M'
$$
由 $M'$ 在 $A'$ 上平坦这一假设，该张量积单射地映入 $M'$。
由此得到 $\mathfrak b' \otimes_{B'} N' \to N' \to M'$ 是单射，
故第一个映射如所需也是单射。

\medskip\noindent
证明 (2)。将短正合序列
$0 \to B' \to B \oplus A' \to A \to 0$
在 $B'$ 上与 $L'$ 张量即可。

\medskip\noindent
证明 (3)。这是 (1) 与 (2) 的直接推论。
\end{proof}

\begin{lemma}
\label{more-algebra-lemma-finitely-presented-module-over-fibre-product}
设 $A, A', B, B', I$ 如情形
\ref{more-algebra-situation-module-over-fibre-product}。有限投射 $B'$-模的范畴
等价于 $\text{Mod}_B \times_{\text{Mod}_A} \text{Mod}_{A'}$
的满子范畴；该满子范畴由满足 $N$ 在 $B$ 上有限投射且 $M'$
在 $A'$ 上有限投射的三元组 $(N, M', \varphi)$ 组成。
\end{lemma}

\begin{proof}
回忆模有限投射当且仅当它有限表示且平坦，见
Algebra, Lemma \ref{algebra-lemma-finite-projective}。
利用引理 \ref{more-algebra-lemma-flat-module-over-fibre-product} 和
\ref{more-algebra-lemma-finite-module-over-fibre-product}，问题归结为证明：
若 $N$ 在 $B$ 上有限投射，且 $M'$ 在 $A'$ 上有限投射，则
$N' = N \times_{\varphi, M} M'$ 是有限表示的 $B'$-模。

\medskip\noindent
由引理 \ref{more-algebra-lemma-finite-module-over-fibre-product}，模 $N'$
在 $B'$ 上有限。选取核为 $K'$ 的满射
$(B')^{\oplus n} \to N'$。基变换给出映射
$B^{\oplus n} \to N$、$(A')^{\oplus n} \to M'$ 和
$A^{\oplus n} \to M$，其核分别为 $K_B$、$K_{A'}$ 和 $K_A$。
存在规范映射
$$
K' \longrightarrow K_B \times_{K_A} K_{A'}
$$
另一方面，由 $N' = N \times_{\varphi, M} M'$ 及
$B' = B \times_A A'$，也存在规范映射
$K_B \times_{K_A} K_{A'} \to K'$，它是上述箭头的逆。
故所展示的映射是同构。由
Algebra, Lemma \ref{algebra-lemma-extension}，模 $K_B$ 与
$K_{A'}$ 都有限。由引理
\ref{more-algebra-lemma-finite-module-over-fibre-product} 可知，只要
$K_B \to K_A$ 与 $K_{A'} \to K_A$ 诱导同构
$K_B \otimes_B A = K_A = K_{A'} \otimes_{A'} A$，
$K'$ 就是有限 $B'$-模。这一点成立，因为平坦性假设蕴含序列
$$
0 \to K_B \to B^{\oplus n} \to N \to 0
\quad\text{且}\quad
0 \to K_{A'} \to (A')^{\oplus n} \to M' \to 0
$$
在张量后仍然正合，见
Algebra, Lemma \ref{algebra-lemma-flat-tor-zero}。
\end{proof}

\section{环的纤维积 III}
\label{more-algebra-section-fibre-products-rings-III}

\noindent
本节讨论如下情形中的纤维积。

\begin{situation}
\label{more-algebra-situation-relative-module-over-fibre-product}
设 $A, A', B, B', I$ 如情形
\ref{more-algebra-situation-module-over-fibre-product}。设 $B' \to D'$ 为环映射。
令 $D = D' \otimes_{B'} B$、
$C' = D' \otimes_{B'} A'$ 和 $C = D' \otimes_{B'} A$。
由此得到环的一个大交换图
$$
\xymatrix{
C & & & C' \ar[lll] \\
& A \ar[ul] & A' \ar[l] \ar[ru] \\
& B \ar[u] \ar[ld] & B' \ar[l] \ar[u] \ar[rd] \\
D \ar[uuu] & & & D' \ar[lll] \ar[uuu]
}
$$
。注意，我们{\bf 不}假设映射 $D' \to D \times_C C'$ 是同构
\footnote{但由引理 \ref{more-algebra-lemma-module-over-fibre-product-bis}，
$D' \to D \times_C C'$ 是满射。}。
在此情形中，有与 (\ref{more-algebra-equation-functor}) 类似的函子
\begin{equation}
\label{more-algebra-equation-relative-functor}
\text{Mod}_{D'} \longrightarrow
\text{Mod}_D \times_{\text{Mod}_C} \text{Mod}_{C'},\quad
L' \longmapsto (L' \otimes_{D'} D, L' \otimes_{D'} C', can).
\end{equation}
注意
$L' \otimes_{D'} D = L \otimes_{D'} (D' \otimes_{B'} B) = L \otimes_{B'} B$
以及类似地
$L' \otimes_{D'} C' = L \otimes_{D'} (D' \otimes_{B'} A') = L \otimes_{B'} A'$，
故图
$$
\xymatrix{
\text{Mod}_{D'} \ar[r] \ar[d] &
\text{Mod}_D \times_{\text{Mod}_C} \text{Mod}_{C'} \ar[d] \\
\text{Mod}_{B'} \ar[r] &
\text{Mod}_B \times_{\text{Mod}_A} \text{Mod}_{A'}
}
$$
交换。下文以 $(N, M', \varphi)$ 表示
$\text{Mod}_D \times_{\text{Mod}_C} \text{Mod}_{C'}$ 的一个对象，
即 $N$ 是 $D$-模，$M'$ 是 $C'$-模，而
$\varphi : N \otimes_B A \to M' \otimes_{A'} A$
是 $C$-模同构。
不过，把 $\varphi$ 看成 $D$-线性映射
$\varphi : N \to M'/IM'$ 往往更为方便；它诱导同构
$N \otimes_B A \to M' \otimes_{A'} A = M'/IM'$。
\end{situation}

\begin{lemma}
\label{more-algebra-lemma-relative-module-over-fibre-product}
在情形 \ref{more-algebra-situation-relative-module-over-fibre-product} 中，
函子 (\ref{more-algebra-equation-relative-functor}) 有右伴随，即函子
$$
F : (N, M', \varphi) \longmapsto N \times_{\varphi, M} M'
$$
，其中 $M = M'/IM'$。此外，$F$ 与
(\ref{more-algebra-equation-relative-functor}) 的复合是
$\text{Mod}_D \times_{\text{Mod}_C} \text{Mod}_{C'}$
上的恒等函子。换言之，令
$N' = N \times_{\varphi, M} M'$，则
$N' \otimes_{D'} D = N$ 且 $N' \otimes_{D'} C' = M'$。
\end{lemma}

\begin{proof}
伴随性断言由更一般的引理 \ref{more-algebra-lemma-modules} 推出。
最后一个断言由引理 \ref{more-algebra-lemma-module-over-fibre-product}
中的相应断言推出，因为
$N' \otimes_{D'} D = N' \otimes_{D'} D' \otimes_{B'} B = N' \otimes_{B'} B$
且
$N' \otimes_{D'} C' = N' \otimes_{D'} D' \otimes_{B'} A' = N' \otimes_{B'} A'$。
\end{proof}

\begin{lemma}
\label{more-algebra-lemma-relative-surjection-ideals}
在情形 \ref{more-algebra-situation-relative-module-over-fibre-product} 中，映射
$JD' \to IC'$ 是满射，其中 $J = \Ker(B' \to B)$。
\end{lemma}

\begin{proof}
由于 $C' = D' \otimes_{B'} A'$，$IC'$ 是映射
$D' \otimes_{B'} I = C' \otimes_{A'} I \to C'$ 的像。
环映射 $B' \to A'$ 诱导同构 $J \to I$，故结论成立。
\end{proof}

\begin{lemma}
\label{more-algebra-lemma-relative-finite-module-over-fibre-product}
设 $A, A', B, B', C, C', D, D', I, M', M, N, \varphi$
如引理 \ref{more-algebra-lemma-relative-module-over-fibre-product}。
若 $N$ 在 $D$ 上有限，且 $M'$ 在 $C'$ 上有限，则
$N' = N \times_{\varphi, M} M'$ 在 $D'$ 上有限。
\end{lemma}

\begin{proof}
回忆由引理 \ref{more-algebra-lemma-module-over-fibre-product-bis}，
$D' \to D \times_C C'$ 是满射。注意
$N' = N \times_{\varphi, M} M'$ 是 $D \times_C C'$ 上的模。
把引理 \ref{more-algebra-lemma-finite-module-over-fibre-product} 应用于数据
$C, C', D, D', IC', M', M, N, \varphi$，可见
$N' = N \times_{\varphi, M} M'$ 在 $D \times_C C'$ 上有限，
故在 $D'$ 上也有限。
\end{proof}

\begin{lemma}
\label{more-algebra-lemma-relative-flat-module-over-fibre-product}
设 $A, A', B, B', C, C', D, D', I$ 如情形
\ref{more-algebra-situation-relative-module-over-fibre-product}。
\begin{enumerate}
\item 设 $(N, M', \varphi)$ 是
$\text{Mod}_D \times_{\text{Mod}_C} \text{Mod}_{C'}$ 的对象。
若 $M'$ 在 $A'$ 上平坦，且 $N$ 在 $B$ 上平坦，则
$N' = N \times_{\varphi, M} M'$ 在 $B'$ 上平坦。
\item 若 $L'$ 是在 $B'$ 上平坦的 $D'$-模，则
$L' = (L \otimes_{D'} D) \times_{(L \otimes_{D'} C)} (L \otimes_{D'} C')$。
\item 在 $B'$ 上平坦的 $D'$-模的范畴等价于
$\text{Mod}_D \times_{\text{Mod}_C} \text{Mod}_{C'}$
中满足 $N$ 在 $B$ 上平坦且 $M'$ 在 $A'$ 上平坦的对象
$(N, M', \varphi)$ 所成的范畴。
\end{enumerate}
\end{lemma}

\begin{proof}
第 (1) 项由引理 \ref{more-algebra-lemma-flat-module-over-fibre-product}
的第 (1) 项推出。

\medskip\noindent
第 (2) 项由引理 \ref{more-algebra-lemma-flat-module-over-fibre-product}
的第 (2) 项推出；这里使用
$L' \otimes_{D'} D = L' \otimes_{B'} B$、
$L' \otimes_{D'} C' = L' \otimes_{B'} A'$ 以及
$L' \otimes_{D'} C = L' \otimes_{B'} A$，见情形
\ref{more-algebra-situation-relative-module-over-fibre-product} 中的讨论。

\medskip\noindent
第 (3) 项是 (1) 与 (2) 的直接推论。
\end{proof}

\noindent
下面的引理比其绝对情形中的对应结论
（引理 \ref{more-algebra-lemma-finitely-presented-module-over-fibre-product}）
有趣得多，尽管证明本质上相同。

\begin{lemma}
\label{more-algebra-lemma-relative-finitely-presented-module-over-fibre-product}
设 $A, A', B, B', C, C', D, D', I, M', M, N, \varphi$
如引理 \ref{more-algebra-lemma-relative-module-over-fibre-product}。若
\begin{enumerate}
\item $N$ 在 $D$ 上有限表示，并在 $B$ 上平坦；
\item $M'$ 在 $C'$ 上有限表示，并在 $A'$ 上平坦；
\item 环映射 $B' \to D'$ 可分解为 $B' \to D'' \to D'$，
其中 $B' \to D''$ 平坦，而 $D'' \to D'$ 为有限表示；
\end{enumerate}
则 $N' = N \times_M M'$ 在 $D'$ 上有限表示。
\end{lemma}

\begin{proof}
选取核为有限生成理想 $J$ 的满射
$D''' = D''[x_1, \ldots, x_n] \to D'$。
由 Algebra, Lemma
\ref{algebra-lemma-finite-finitely-presented-extension}，
只需证明 $N'$ 作为 $D'''$-模有限表示。此外，
$D''' \otimes_{B'} B \to D' \otimes_{B'} B = D$ 以及
$D''' \otimes_{B'} A' \to D' \otimes_{B'} A' = C'$ 都是满射，
其核由 $J$ 的像生成。因此，再次应用
Algebra, Lemma
\ref{algebra-lemma-finite-finitely-presented-extension}，
$N$ 是有限表示的 $D''' \otimes_{B'} B$-模，而
$M'$ 是有限表示的 $D''' \otimes_{B'} A'$-模。
故可用 $D'''$ 代替 $D'$，并用 $D''' \otimes_{B'} B$
代替 $D$，其余类似。由于 $D'''$ 在 $B'$ 上平坦，
可归结到 $B' \to D'$ 平坦的情形。

\medskip\noindent
假设 $B' \to D'$ 平坦。由引理
\ref{more-algebra-lemma-relative-finite-module-over-fibre-product}，
模 $N'$ 在 $D'$ 上有限。选取核为 $K'$ 的满射
$(D')^{\oplus n} \to N'$。基变换给出映射
$D^{\oplus n} \to N$、$(C')^{\oplus n} \to M'$ 和
$C^{\oplus n} \to M$，其核分别为 $K_D$、$K_{C'}$ 和 $K_C$。
存在规范映射
$$
K' \longrightarrow K_D \times_{K_C} K_{C'}
$$
。另一方面，由于 $N' = N \times_M M'$，且
$D' = D \times_C C'$（把引理
\ref{more-algebra-lemma-flat-module-over-fibre-product} 应用于平坦 $B'$-模 $D'$），
还存在规范映射 $K_D \times_{K_C} K_{C'} \to K'$，
它是所展示箭头的逆。故该映射是同构。
由 Algebra, Lemma \ref{algebra-lemma-extension}，
模 $K_D$ 与 $K_{C'}$ 都有限。由引理
\ref{more-algebra-lemma-relative-finite-module-over-fibre-product}，只要
$K_D \to K_C$ 与 $K_{C'} \to K_C$ 诱导同构
$K_D \otimes_B A = K_C = K_{C'} \otimes_{A'} A$，
$K'$ 就是有限 $D'$-模。这一点成立，因为平坦性假设蕴含序列
$$
0 \to K_D \to D^{\oplus n} \to N \to 0
\quad\text{且}\quad
0 \to K_{C'} \to (C')^{\oplus n} \to M' \to 0
$$
在张量后仍然正合，见
Algebra, Lemma \ref{algebra-lemma-flat-tor-zero}。
\end{proof}

\begin{lemma}
\label{more-algebra-lemma-properties-algebras-over-fibre-product}
设 $A, A', B, B', I$ 如情形
\ref{more-algebra-situation-module-over-fibre-product}。设 $(D, C', \varphi)$
是由一个 $B$-代数 $D$、一个 $A'$-代数 $C'$ 以及一个同构
$D \otimes_B A \to C'/IC' = C$ 所组成的系统。令
$D' = D \times_C C'$（如引理
\ref{more-algebra-lemma-module-over-fibre-product}）。则
\begin{enumerate}
\item $B' \to D'$ 为有限型，当且仅当 $B \to D$ 与
$A' \to C'$ 为有限型；
\item $B' \to D'$ 平坦，当且仅当 $B \to D$ 与 $A' \to C'$ 平坦；
\item $B' \to D'$ 平坦且为有限表示，当且仅当
$B \to D$ 与 $A' \to C'$ 平坦且为有限表示；
\item $B' \to D'$ 光滑，当且仅当 $B \to D$ 与 $A' \to C'$ 光滑；
\item $B' \to D'$ 是 \'etale 的，当且仅当 $B \to D$ 与
$A' \to C'$ 是 \'etale 的。
\end{enumerate}
此外，若 $D'$ 是平坦 $B'$-代数，则
$D' \to (D' \otimes_{B'} B) \times_{(D' \otimes_{B'} A)} (D' \otimes_{B'} A')$
是同构。如此，平坦 $B'$-代数的范畴等价于上述系统
$(D, C', \varphi)$ 所成的范畴，其中 $D$ 在 $B$ 上平坦，而
$C'$ 在 $A'$ 上平坦。
\end{lemma}

\begin{proof}
蕴含关系“$\Rightarrow$”由下列引理推出：
Algebra, Lemmas \ref{algebra-lemma-base-change-finiteness},
\ref{algebra-lemma-flat-base-change},
\ref{algebra-lemma-base-change-smooth}, and
\ref{algebra-lemma-etale}；这是因为由引理
\ref{more-algebra-lemma-module-over-fibre-product} 有
$D' \otimes_{B'} B = D$ 及 $D' \otimes_{B'} A' = C'$。
因此只需证明反方向的蕴含关系。

\medskip\noindent
关于 (1)。假设 $D$ 在 $B$ 上为有限型，且 $C'$ 在 $A'$ 上为有限型。
下文将不再说明地使用引理
\ref{more-algebra-lemma-module-over-fibre-product} 的结论。
选取 $D$ 在 $B$ 上的一组生成元 $x_1, \ldots, x_r$，
以及 $C'$ 在 $A'$ 上的一组生成元 $y_1, \ldots, y_s$。
利用 $D = D' \otimes_{B'} B$ 以及 $B' \to B$ 为满射，
可选取 $u_1, \ldots, u_r \in D'$，它们在 $D$ 中分别映到
$x_1, \ldots, x_r$。利用 $C' = D' \otimes_{B'} A'$，
可选取 $v_1, \ldots, v_t \in D'$，使得对某些
$a'_{ij} \in A'$ 有 $y_i = \sum v_j \otimes a'_{ij}$。
特别地，$v_j$ 在 $C'$ 中的像作为 $A'$-代数生成 $C'$。
令 $N = r + t$，并考虑环的立方图
$$
\xymatrix{
A[x_1, \ldots, x_N] & & A'[x_1, \ldots, x_N] \ar[ll] \\
& A \ar[lu] & & A' \ar[ll] \ar[lu] \\
B[x_1, \ldots, x_N] \ar[uu] & & B'[x_1, \ldots, x_N] \ar[uu] \ar[ll] \\
& B \ar[uu] \ar[lu] & & B' \ar[ll] \ar[uu] \ar[lu]
}
$$
。注意后面的方块也是笛卡尔方块。考虑环映射
$$
B'[x_1, \ldots, x_N] \to D',\quad
x_i \mapsto u_i \quad\text{且}\quad x_{r + j} \mapsto v_j.
$$
于是所诱导的映射 $B[x_1, \ldots, x_N] \to D$ 与
$A'[x_1, \ldots, x_N] \to C'$ 都是满射，特别地都是有限映射。
由引理 \ref{more-algebra-lemma-relative-finite-module-over-fibre-product}，
$B'[x_1, \ldots, x_N] \to D'$ 是有限映射。因而 $D'$
在 $B'$ 上为有限型，例如可由
Algebra, Lemma \ref{algebra-lemma-compose-finite-type} 得到。

\medskip\noindent
关于 (2)。蕴含关系“$\Leftarrow$”由引理
\ref{more-algebra-lemma-relative-flat-module-over-fibre-product} 推出。
此外，最后一个断言由引理
\ref{more-algebra-lemma-relative-flat-module-over-fibre-product} 的最后一个断言推出。

\medskip\noindent
关于 (3)。假设 $B \to D$ 和 $A' \to C'$ 平坦且为有限表示。
$B' \to D'$ 的平坦性已经在 (2) 中得到。由 (1)，
$B' \to D'$ 为有限型。选取满射
$B'[x_1, \ldots, x_N] \to D'$。
由 Algebra, Lemma
\ref{algebra-lemma-finite-presentation-independent}，环 $D$
作为 $B[x_1, \ldots, x_N]$-模有限表示，而环 $C'$
作为 $A'[x_1, \ldots, x_N]$-模有限表示。
由引理
\ref{more-algebra-lemma-relative-finitely-presented-module-over-fibre-product}，
$D'$ 作为 $B'[x_1, \ldots, x_N]$-模有限表示，即
$B' \to D'$ 为有限表示。

\medskip\noindent
关于 (4)。假设 $B \to D$ 与 $A' \to C'$ 光滑。由 (3)，
$B' \to D'$ 平坦且为有限表示。由 Algebra, Lemma
\ref{algebra-lemma-flat-fibre-smooth}，只需对 $B'$ 上任意域 $k$
验证 $D' \otimes_{B'} k$ 光滑。若复合 $J \to B' \to k$ 为零，
则 $B' \to k$ 分解为 $B' \to B \to k$，且
$$
D' \otimes_{B'} k = D' \otimes_{B'} B \otimes_B k
= D \otimes_B k
$$
光滑，因为 $B \to D$ 光滑。若复合 $J \to B' \to k$ 非零，
则存在 $h \in J$，其在 $k$ 中的像非零。于是 $B' \to k$
分解为 $B' \to B'_h \to k$。注意 $h$ 在 $B$ 中映到零，故
$B_h = 0$。因此由引理 \ref{more-algebra-lemma-diagram-localize}，
$B'_h = A'_h$，并且
$$
D' \otimes_{B'} k = D' \otimes_{B'} B'_h \otimes_{B'_h} k
= C'_h \otimes_{A'_h} k
$$
光滑，因为 $A' \to C'$ 光滑。

\medskip\noindent
关于 (5)。假设 $B \to D$ 与 $A' \to C'$ 是 \'etale 的。
由 (4)，$B' \to D'$ 光滑。光滑映射是否为 \'etale 可由纤维的维数读出，
故 (5) 成立（仿照 (4) 的证明识别纤维；略去若干细节）。
\end{proof}

\begin{remark}
\label{more-algebra-remark-relative-modules-over-fibre-product}
在情形 \ref{more-algebra-situation-relative-module-over-fibre-product} 中，
假设 $B' \to D'$ 为有限表示，并设给定 $D'$-模 $L'$。
我们断言下列两类对象之间存在双射：
\begin{enumerate}
\item $D'$-模的满射 $L' \to Q'$，其中 $Q'$ 在 $D'$ 上有限表示，
且在 $B'$ 上平坦；
\item 模的成对满射
$(L' \otimes_{D'} D \to Q_1, L' \otimes_{D'} C' \to Q_2)$，
满足
\begin{enumerate}
\item $Q_1$ 在 $D$ 上有限表示，并在 $B$ 上平坦；
\item $Q_2$ 在 $C'$ 上有限表示，并在 $A'$ 上平坦；
\item 作为 $L' \otimes_{D'} C$ 的商，有
$Q_1 \otimes_D C = Q_2 \otimes_{C'} C$。
\end{enumerate}
\end{enumerate}
从前者到后者的对应由 $Q \mapsto (Q_1, Q_2)$ 给出，其中
$Q_1 = Q \otimes_{D'} D$ 且 $Q_2 = Q \otimes_{D'} C'$。
反过来，使用 $Q = Q_1 \times_{Q_{12}} Q_2$，其中 $Q_{12}$ 是
$L' \otimes_{D'} C$ 的公共商
$Q_1 \otimes_D C = Q_2 \otimes_{C'} C$。作为商映射，使用
$$
L' \longrightarrow
(L' \otimes_{D'} D) \times_{(L' \otimes_{D'} C)} (L' \otimes_{D'} C')
\longrightarrow Q_1 \times_{Q_{12}} Q_2 = Q
$$
，其中第一个箭头由引理
\ref{more-algebra-lemma-module-over-fibre-product-bis} 为满射，第二个箭头由引理
\ref{more-algebra-lemma-surjection-module-over-fibre-product} 为满射。
断言由引理 \ref{more-algebra-lemma-relative-flat-module-over-fibre-product} 与
\ref{more-algebra-lemma-relative-finitely-presented-module-over-fibre-product} 推出。
\end{remark}








\section{Fitting 理想}
\label{more-algebra-section-fitting-ideals}

\noindent
有限模的 Fitting 理想是由引理 \ref{more-algebra-lemma-fitting-ideal}
中的构造所确定的理想。

\begin{lemma}
\label{more-algebra-lemma-ideals-generated-by-minors}
设 $R$ 为环，$A$ 为系数在 $R$ 中的 $n \times m$ 矩阵。
令 $I_r(A)$ 为由 $A$ 的 $r \times r$ 子式生成的理想，并约定
$I_0(A) = R$；若 $r > \min(n, m)$，则约定 $I_r(A) = 0$。则
\begin{enumerate}
\item $I_0(A) \supset I_1(A) \supset I_2(A) \supset \ldots$；
\item 若 $B$ 是 $(n + n') \times m$ 矩阵，而 $A$ 是 $B$ 的前
$n$ 行，则 $I_{r + n'}(B) \subset I_r(A)$；
\item 若 $C$ 是 $n \times n$ 矩阵，则 $I_r(CA) \subset I_r(A)$。
\item 若 $A$ 是分块矩阵
$$
\left(
\begin{matrix}
A_1 & 0 \\
0 & A_2 
\end{matrix}
\right)
$$
，则 $I_r(A) = \sum_{r_1 + r_2 = r} I_{r_1}(A_1) I_{r_2}(A_2)$。
\item 此处待补充。
\end{enumerate}
\end{lemma}

\begin{proof}
略。（提示：行列式可沿一列或一行展开计算。）
\end{proof}

\begin{lemma}
\label{more-algebra-lemma-fitting-ideal}
设 $R$ 为环，$M$ 为有限 $R$-模。选取 $M$ 的一个表示
$$
\bigoplus\nolimits_{j \in J} R \longrightarrow R^{\oplus n}
\longrightarrow M \longrightarrow 0.
$$
令 $A = (a_{ij})_{i = 1, \ldots, n, j \in J}$ 为映射
$\bigoplus_{j \in J} R \to R^{\oplus n}$ 的矩阵。由 $A$ 的
$(n - k) \times (n - k)$ 子式生成的理想
$\text{Fit}_k(M)$ 不依赖于表示的选择。
\end{lemma}

\begin{proof}
设 $K \subset R^{\oplus n}$ 为满射 $R^{\oplus n} \to M$ 的核。
选取 $z_1, \ldots, z_{n - k} \in K$，并写
$z_j = (z_{1j}, \ldots, z_{nj})$。理想 $\text{Fit}_k(M)$
的另一种描述是：它由按这种方式得到的所有矩阵 $(z_{ij})$ 的
$(n - k) \times (n - k)$ 子式生成。

\medskip\noindent
假设把该满射换成核为 $K'$ 的满射
$R^{\oplus n + n'} \to M$，其中在 $R^{\oplus n + n'}$ 的前
$n$ 个标准基向量上使用原映射，而在最后 $n'$ 个基向量上取 $0$。
则相应的理想相同。事实上，若
$z_1, \ldots, z_{n - k} \in K$ 如上，令
$z'_j = (z_{1j}, \ldots, z_{nj}, 0, \ldots, 0) \in K'$，
其中 $j = 1, \ldots, n - k$，并令
$z'_{n + j'} = (0, \ldots, 0, 1, 0, \ldots, 0) \in K'$。
于是矩阵 $(z_{ij})$ 的 $(n - k) \times (n - k)$ 子式所生成的理想，
与 $(z'_{ij})$ 的 $(n + n' - k) \times (n + n' - k)$
子式所生成的理想相同。这给出一个包含关系。
反过来，给定 $K'$ 中的
$z'_1, \ldots, z'_{n + n' - k}$，可把它们投影到
$R^{\oplus n}$，得到 $K$ 中的
$z_1, \ldots, z_{n + n' - k}$。由引理
\ref{more-algebra-lemma-ideals-generated-by-minors}，$(z'_{ij})$ 的
$(n + n' - k) \times (n + n' - k)$ 子式所生成的理想，
包含于 $(z_{ij})$ 的 $(n - k) \times (n - k)$ 子式所生成的理想。
这给出另一个包含关系。

\medskip\noindent
设 $R^{\oplus m} \to M$ 是另一个核为 $L$ 的满射。
由 Schanuel 引理（Algebra, Lemma \ref{algebra-lemma-Schanuel}）
及上一段的结论，可假设 $m = n$，并且存在同构
$R^{\oplus n} \to R^{\oplus m}$，它与到 $M$ 的满射交换。
令 $C = (c_{li})$ 为该映射的（可逆）矩阵
（由于 $n = m$，这是方阵）。于是给定如上的
$z'_1, \ldots, z'_{n - k} \in L$，可选取
$z_1, \ldots, z_{n - k} \in K$，使得
$z_1' = Cz_1, \ldots, z'_{n - k} = Cz_{n - k}$。
由引理 \ref{more-algebra-lemma-ideals-generated-by-minors} 得到一个包含关系，
再由对称性得到另一个。
\end{proof}

\begin{definition}
\label{more-algebra-definition-fitting-ideal}
设 $R$ 为环，$M$ 为有限 $R$-模，$k \geq 0$。
$M$ 的{\it 第 $k$ 个 Fitting 理想}是引理
\ref{more-algebra-lemma-fitting-ideal} 中构造的理想 $\text{Fit}_k(M)$。
令 $\text{Fit}_{-1}(M) = 0$。
\end{definition}

\noindent
由于 Fitting 理想是一个大矩阵的子式理想
（其编号顺序与引理 \ref{more-algebra-lemma-ideals-generated-by-minors}
中的顺序相反），可见对某个 $t \gg 0$ 有
$$
0 = \text{Fit}_{-1}(M) \subset \text{Fit}_0(M) \subset \text{Fit}_1(M)
\subset \ldots \subset \text{Fit}_t(M) = R
$$
。下面给出 Fitting 理想的一些基本性质。

\begin{lemma}
\label{more-algebra-lemma-fitting-ideal-basics}
设 $R$ 为环，$M$ 为有限 $R$-模。
\begin{enumerate}
\item 若 $M$ 可由 $n$ 个元素生成，则 $\text{Fit}_n(M) = R$。
\item 给定另一个有限 $R$-模 $M'$，有
$\text{Fit}_0(M \oplus M') = \text{Fit}_0(M)\text{Fit}_0(M')$，
且更一般地，
$$
\text{Fit}_l(M \oplus M') =
\sum\nolimits_{k + k' = l} \text{Fit}_k(M)\text{Fit}_{k'}(M')
$$
\item 若 $R \to R'$ 为环映射，则 $\text{Fit}_k(M \otimes_R R')$
是 $R'$ 中由 $\text{Fit}_k(M)$ 的像生成的理想。
\item 若 $M$ 有限表示，则 $\text{Fit}_k(M)$ 是有限生成理想。
\item 若 $M \to M'$ 是满射，则
$\text{Fit}_k(M) \subset \text{Fit}_k(M')$。
\item 有 $\text{Fit}_0(M) \subset \text{Ann}_R(M)$。
\item 有 $V(\text{Fit}_0(M)) = \text{Supp}(M)$。
\item 若 $I$ 是 $R$ 的理想，则 $\text{Fit}_0(R/I) = I$。
\item 此处待补充。
\end{enumerate}
\end{lemma}

\begin{proof}
第 (1) 项由引理 \ref{more-algebra-lemma-ideals-generated-by-minors} 中
$I_0(A) = R$ 这一事实推出。

\medskip\noindent
第 (2) 项由引理 \ref{more-algebra-lemma-ideals-generated-by-minors} 中的相应断言推出。

\medskip\noindent
第 (3) 项由 $\otimes_R R'$ 右正合这一事实推出，因此 $M$ 的一个表示
作基变换后是 $M \otimes_R R'$ 的一个表示。

\medskip\noindent
证明 (4)。设
$R^{\oplus m} \xrightarrow{A} R^{\oplus n} \to M \to 0$
是一个表示。则 $\text{Fit}_k(M)$ 是由矩阵 $A$ 的
$(n - k) \times (n - k)$ 子式生成的理想。

\medskip\noindent
第 (5) 项由定义立即得到。

\medskip\noindent
证明 (6)。选取由矩阵 $A$ 给出的 $M$ 的一个表示，如引理
\ref{more-algebra-lemma-fitting-ideal}。设 $J' \subset J$ 是基数为 $n$ 的子集。
只需证明
$f = \det(a_{ij})_{i = 1, \ldots, n, j \in J'}$
零化 $M$。这是清楚的，因为
$$
R^{\oplus n} \xrightarrow{A' = (a_{ij})_{i = 1, \ldots, n, j \in J'}}
R^{\oplus n} \to M \to 0
$$
的余核被 $f$ 零化：存在矩阵 $B$，满足
$A' B = f1_{n \times n}$。

\medskip\noindent
证明 (7)。选取由矩阵 $A$ 给出的 $M$ 的一个表示，如引理
\ref{more-algebra-lemma-fitting-ideal}。由 Nakayama 引理
（Algebra, Lemma \ref{algebra-lemma-NAK}），有
$$
M_\mathfrak p \not = 0
\Leftrightarrow
M \otimes_R \kappa(\mathfrak p) \not = 0
\Leftrightarrow
\text{rank}(\text{像 }A\text{ 于 }\kappa(\mathfrak p)) < n
$$
。显然，$\text{Fit}_0(M)$ 恰好切出具有该性质的素理想集合。
\end{proof}

\begin{example}
\label{more-algebra-example-fitting-free}
设 $R$ 为环。有限自由模 $M = R^{\oplus n}$ 的 Fitting 理想满足：
当 $k < n$ 时 $\text{Fit}_k(M) = 0$，当 $k \geq n$ 时
$\text{Fit}_k(M) = R$。
\end{example}






\begin{example}
\label{more-algebra-example-laurent}
设 $R$ 为环，$I, J\subset R$ 为理想。则
$\text{Fit}_0(R/I) = I$、
$\text{Fit}_0(R/I \oplus R/J) = IJ$、
$\text{Fit}_1(R/I\oplus R/J) = I+J$。
\end{example}

\begin{lemma}
\label{more-algebra-lemma-fitting-ideal-generate-locally}
设 $R$ 为环，$M$ 为有限 $R$-模，$k \geq 0$。
设 $\mathfrak p \subset R$ 为素理想。下列条件等价：
\begin{enumerate}
\item $\text{Fit}_k(M) \not \subset \mathfrak p$；
\item $\dim_{\kappa(\mathfrak p)} M \otimes_R \kappa(\mathfrak p) \leq k$；
\item $M_\mathfrak p$ 作为 $R_\mathfrak p$-模可由 $k$ 个元素生成；
\item 对某个 $f \in R$、$f \not \in \mathfrak p$，
$M_f$ 作为 $R_f$-模可由 $k$ 个元素生成。
\end{enumerate}
\end{lemma}

\begin{proof}
由 Nakayama 引理（Algebra, Lemma \ref{algebra-lemma-NAK}），
若 $M \otimes_R \kappa(\mathfrak p)$ 可由 $k$ 个元素生成，
则对某个 $f \in R$、$f \not \in \mathfrak p$，
$M_f$ 作为 $R_f$-模可由 $k$ 个元素生成。因此 (2)、(3) 与 (4) 等价。
利用引理 \ref{more-algebra-lemma-fitting-ideal-basics} 第 (3) 项，
问题归结为 $R$ 是域且 $\mathfrak p = (0)$ 的情形。
此时结论由例 \ref{more-algebra-example-fitting-free} 推出。
\end{proof}

\begin{lemma}
\label{more-algebra-lemma-fitting-ideal-finite-locally-free}
设 $R$ 为环，$M$ 为有限 $R$-模，$r \geq 0$。下列条件等价：
\begin{enumerate}
\item $M$ 是秩为 $r$ 的有限局部自由模
（Algebra, Definition \ref{algebra-definition-locally-free}）；
\item $\text{Fit}_{r - 1}(M) = 0$ 且 $\text{Fit}_r(M) = R$；
\item 当 $k < r$ 时 $\text{Fit}_k(M) = 0$，而当 $k \geq r$ 时
$\text{Fit}_k(M) = R$。
\end{enumerate}
\end{lemma}

\begin{proof}
由于 Fitting 理想组成递增理想序列，(2) 与 (3) 显然等价。
由于 $\text{Fit}_k(M)$ 的形成与基变换可交换
（引理 \ref{more-algebra-lemma-fitting-ideal-basics}），例
\ref{more-algebra-example-fitting-free} 及粘合结果
（Algebra, Section \ref{algebra-section-more-glueing}）
表明 (1) 蕴含 (2)。反过来，假设 (2)。由引理
\ref{more-algebra-lemma-fitting-ideal-generate-locally}，可假设 $M$ 由 $r$
个元素生成。因此有一个表示
$\bigoplus_{j \in J} R \to R^{\oplus r} \to M \to 0$。
现在，假设 $\text{Fit}_{r - 1}(M) = 0$ 蕴含映射
$\bigoplus_{j \in J} R \to R^{\oplus r}$ 的矩阵的所有元素均为零。
因此 $M$ 是自由模。
\end{proof}

\begin{lemma}
\label{more-algebra-lemma-principal-fitting-ideal}
设 $R$ 为局部环，$M$ 为有限 $R$-模，$k \geq 0$。
假设对某个 $f \in R$ 有 $\text{Fit}_k(M) = (f)$。
令 $M'$ 为 $M$ 关于 $\{x \in M \mid fx = 0\}$ 的商。
则 $M'$ 可由 $k$ 个元素生成。
\end{lemma}

\begin{proof}
选取对应于满射 $R^{\oplus n} \to M$ 的生成元
$x_1, \ldots, x_n \in M$。由于 $R$ 是局部环，若元素集合
$E \subset (f)$ 生成 $(f)$，则某个 $e \in E$ 生成 $(f)$，见
Algebra, Lemma \ref{algebra-lemma-NAK}。因此可在
$R^{\oplus n} \to M$ 的核中选取 $z_1, \ldots, z_{n - k}$，
使得 $n \times (n - k)$ 矩阵 $A = (z_{ij})$ 的某个
$(n - k) \times (n - k)$ 子式生成 $(f)$。重新编号 $x_i$ 后，
可假设第一个子式
$\det(z_{ij})_{1 \leq i, j \leq n - k}$ 生成 $(f)$，即对某个
单位元 $u \in R$ 有
$\det(z_{ij})_{1 \leq i, j \leq n - k} = uf$。
其他每个子式都是 $f$ 的倍数。由 Algebra, Lemma
\ref{algebra-lemma-matrix-right-inverse}，存在一个
$n - k \times n - k$ 矩阵 $B$，使得
$$
AB = f
\left(
\begin{matrix}
u 1_{n - k \times n - k} \\
C
\end{matrix}
\right)
$$
，其中 $C$ 是系数在 $R$ 中的某个矩阵。这表明对每个
$i \leq n - k$，元素 $y_i = ux_i + \sum_j c_{ji}x_j$ 被 $f$ 零化。
由于 $M/\sum Ry_i$ 由 $x_{n - k + 1}, \ldots, x_n$ 的像生成，
结论成立。
\end{proof}

\begin{lemma}
\label{more-algebra-lemma-fitting-ideals-and-pd1}
设 $R$ 为环，$M$ 为有限表示 $R$-模，$k \geq 0$。
假设对某个非零因子 $f \in R$ 有
$\text{Fit}_k(M) = (f)$，并且 $\text{Fit}_{k - 1}(M) = 0$。则
\begin{enumerate}
\item $M$ 的投射维数 $\leq 1$；
\item $M' = \Ker(f : M \to M)$ 是 $M$ 的 $f$-幂挠子模；
\item $M'$ 的投射维数 $\leq 1$；
\item $M/M'$ 是秩为 $k$ 的有限局部自由模；
\item $M \cong M/M' \oplus M'$。
\end{enumerate}
\end{lemma}

\begin{proof}
选取由系数在 $R$ 中的某个矩阵 $A$ 给出的表示
$$
R^{\oplus m} \xrightarrow{A} R^{\oplus n} \to M \to 0
$$
。

\medskip\noindent
先在 $R$ 为局部环时证明该引理。如陈述中一样，令
$M' = \{x \in M \mid fx = 0\}$。由引理
\ref{more-algebra-lemma-principal-fitting-ideal}，可选取
$x_1, \ldots, x_k \in M$ 生成 $M/M'$。于是
$x_1, \ldots, x_k$ 生成 $M_f = (M/M')_f$。因此，若 $M$ 中存在关系
$\sum a_ix_i = 0$，则 $a_1, \ldots, a_k$ 在 $R_f$ 中映到零；
否则 $\text{Fit}_{k - 1}(M) R_f = \text{Fit}_{k - 1}(M_f)$
将非零。由于 $f$ 是非零因子，得到
$a_1 = \ldots = a_k = 0$。因此
$M \cong R^{\oplus k} \oplus M'$。对上述表示作基变换后，
可假设 $R^{\oplus n}$ 的前 $n - k$ 个基向量映入 $M$ 的直和分支
$M'$，而 $R^{\oplus n}$ 的后 $k$ 个基向量映到 $M$ 的直和分支 $R^{\oplus k}$
的基元素。如此，矩阵 $A$ 的后 $k$ 行为零。于是用 $M'$ 替换 $M$、
用 $0$ 替换 $k$、用 $n - k$ 替换 $n$，并用删除后 $k$ 行所得的
子矩阵替换 $A$，便归结为下一段讨论的情形。

\medskip\noindent
假设 $R$ 为局部环、$k = 0$，且 $M$ 被 $f$ 零化。
此时 $M$ 的第 $0$ 个 Fitting 理想是 $(f)$，并由大小为
$n \times m$ 的矩阵 $A$ 的 $n \times n$ 子式生成
（这特别蕴含 $m \geq n$）。由于 $R$ 是局部环，$A$ 的某个
$n \times n$ 子式等于 $uf$，其中 $u \in R$ 为单位元。
重新编号后，可假设这是第一个子式。此外，已知 $A$ 的所有其他
$n \times n$ 子式都被 $f$ 整除。以分块形式写
$A = (A_1 A_2)$，其中 $A_1$ 是 $n \times n$ 矩阵，
$A_2$ 是 $n \times (m - n)$ 矩阵。把
Algebra, Lemma \ref{algebra-lemma-matrix-right-inverse}
应用于 $A$ 的转置（！），得到一个 $n \times n$ 矩阵 $B$，使得
$$
BA = B(A_1 A_2) = f
\left(
\begin{matrix}
u 1_{n \times n} & C
\end{matrix}
\right)
$$
，其中 $C$ 是系数在 $R$ 中的某个 $n \times (m - n)$ 矩阵。
先得到 $BA_1 = fu 1_{n \times n}$。于是
$$
BA_2 = fC = u^{-1}fuC = u^{-1}BA_1C
$$
。由于 $B$ 的行列式是非零因子，得到
$A_2 = u^{-1}A_1C$。因此 $A$ 的像等于 $A_1$ 的像；
由于 $A_1$ 的行列式是非零因子，该像同构于 $R^{\oplus n}$。
故 $M$ 的投射维数 $\leq 1$。

\medskip\noindent
回到一般环 $R$ 的情形。由局部情形可知 $M/M'$ 是秩为 $k$
的有限局部自由模，见
Algebra, Lemma \ref{algebra-lemma-finite-projective}。
因此扩张 $0 \to M' \to M \to M/M' \to 0$ 分裂。
从而 $M'$ 是有限表示模。选取短正合序列
$0 \to K \to R^{\oplus a} \to M' \to 0$。则 $K$ 是有限 $R$-模，
见 Algebra, Lemma \ref{algebra-lemma-extension}。
由局部情形，对所有素理想都有
$K_\mathfrak p \cong R_\mathfrak p^{\oplus a}$。因此再次由
Algebra, Lemma \ref{algebra-lemma-finite-projective}，
$K$ 是秩为 $a$ 的有限局部自由模。由此 $M'$ 的投射维数
$\leq 1$，引理得证。
\end{proof}










\section{提升}
\label{more-algebra-section-lifting}

\noindent
本节汇集若干关于如下类型的提升陈述的引理：若 $A$ 是环、
$I \subset A$ 是理想，而 $\overline{\xi}$ 是 $A/I$ 上的某种结构，
则可把 $\overline{\xi}$ 提升为 $A$ 上或 $A$ 的某个 \'etale
扩张上的同类结构 $\xi$。对于下列结构，我们已经证明了一些结果：
\begin{enumerate}
\item 幂等元，见
Algebra, Lemmas \ref{algebra-lemma-lift-idempotents} and
\ref{algebra-lemma-lift-idempotents-noncommutative}；
\item 投射模，见
Algebra, Lemmas \ref{algebra-lemma-lift-projective-module} and
\ref{algebra-lemma-lift-finite-projective-module}；
\item 有限稳定自由模，见引理 \ref{more-algebra-lemma-lift-stably-free}；
\item 基元素，见
Algebra, Lemmas \ref{algebra-lemma-local-artinian-basis-when-flat} and
\ref{algebra-lemma-lift-basis}；
\item 环映射，即证明某些代数形式光滑，见
Algebra, Lemma \ref{algebra-lemma-polynomial-ring-formally-smooth},
Proposition \ref{algebra-proposition-smooth-formally-smooth}, and
Lemma \ref{algebra-lemma-smooth-strong-lift}；
\item syntomic 环映射，见
Algebra, Lemma \ref{algebra-lemma-lift-syntomic}；
\item 光滑环映射，见
Algebra, Lemma \ref{algebra-lemma-lift-smooth}；
\item \'etale 环映射，见
Algebra, Lemma \ref{algebra-lemma-lift-etale}；
\item 多项式的分解，见
Algebra, Lemma \ref{algebra-lemma-factor-mod-lift-etale}；
\item Algebra, Section \ref{algebra-section-henselian}
讨论 Hensel 局部环。
\end{enumerate}
感兴趣的读者可在 Smoothing Ring Maps,
Section \ref{smoothing-section-presentations} 中找到更多同类结果，
特别是 Smoothing Ring Maps, Proposition
\ref{smoothing-proposition-lift-smooth}。

\medskip\noindent
设 $A$ 为环，$I \subset A$ 为理想，$\overline{\xi}$
为 $A/I$ 上的某种结构。在下面的引理中，我们寻找诱导同构
$A/I \to A'/IA'$ 的 \'etale 环映射 $A \to A'$，以及 $A'$
上的对象 $\xi'$，使其提升 $\overline{\xi}$。一个一般性说明是：
给定 \'etale 环映射 $A \to A' \to A''$，若
$A/I \cong A'/IA'$ 且 $A'/IA' \cong A''/IA''$，则复合
$A \to A''$ 仍为 \'etale
（Algebra, Lemma \ref{algebra-lemma-etale}），并且仍满足
$A/I \cong A''/IA''$。下面将经常使用这一点而不再说明。
先给出一个我们所寻求类型的平凡例子。

\begin{lemma}
\label{more-algebra-lemma-lift-invertible-element}
设 $A$ 为环，$I \subset A$ 为理想，
$\overline{u} \in A/I$ 为可逆元素。存在一个诱导同构
$A/I \to A'/IA'$ 的 \'etale 环映射 $A \to A'$，以及提升
$\overline{u}$ 的可逆元素 $u' \in A'$。
\end{lemma}

\begin{proof}
任取 $\overline{u}$ 的提升 $f \in A$，并令 $A' = A_f$，
令 $u$ 为 $f$ 在 $A'$ 中的像。
\end{proof}

\begin{lemma}
\label{more-algebra-lemma-lift-idempotent}
设 $A$ 为环，$I \subset A$ 为理想，
$\overline{e} \in A/I$ 为幂等元。存在一个诱导同构
$A/I \to A'/IA'$ 的 \'etale 环映射 $A \to A'$，以及提升
$\overline{e}$ 的幂等元 $e' \in A'$。
\end{lemma}

\begin{proof}
任取 $\overline{e}$ 的提升 $x \in A$。令
$$
A' = A[t]/(t^2 - t)\left[\frac{1}{t - 1 + x}\right].
$$
环映射 $A \to A'$ 是 \'etale 的，因为
$(2t - 1)\text{d}t = 0$，且
$(2t - 1)(2t - 1) = 1$ 可逆。又有
$A'/IA' = A/I[t]/(t^2 - t)[\frac{1}{t - 1 + \overline{e}}] \cong A/I$；
最后一个映射把 $t$ 送到 $\overline{e}$，这是可行的，因为
$\overline{e}$ 是 $t^2 - t$ 的根。这也表明，令 $e'$ 为 $t$
在 $A'$ 中的类即可。
\end{proof}

\begin{lemma}
\label{more-algebra-lemma-lift-open-covering}
设 $A$ 为环，$I \subset A$ 为理想。设
$\Spec(A/I) = \coprod_{j \in J} \overline{U}_j$
为有限不交开覆盖。则存在一个诱导同构
$A/I \to A'/IA'$ 的 \'etale 环映射 $A \to A'$，以及提升给定覆盖的
有限不交开覆盖 $\Spec(A') = \coprod_{j \in J} U'_j$。
\end{lemma}

\begin{proof}
这由引理 \ref{more-algebra-lemma-lift-idempotent} 以及谱的开闭子集与幂等元对应这一事实
推出，见 Algebra, Lemma
\ref{algebra-lemma-disjoint-decomposition}。
\end{proof}

\begin{lemma}
\label{more-algebra-lemma-localize-upstairs}
设 $A \to B$ 为环映射，$J \subset B$ 为理想。
若 $A \to B$ 在 $V(J)$ 的每个素理想处都是 \'etale 的，
则存在 $g \in B$，它在 $B/J$ 中映到可逆元素，并且
$A' = B_g$ 在 $A$ 上为 \'etale。
\end{lemma}

\begin{proof}
$\Spec(B)$ 中 $A \to B$ 不是 \'etale 的点所成集合是
$\Spec(B)$ 的闭子集，见
Algebra, Definition \ref{algebra-definition-etale}。
把它写成 $V(J')$，其中 $J' \subset B$ 为某个理想。
由假设 $V(J') \cap V(J) = \emptyset$，故由
Algebra, Lemma \ref{algebra-lemma-Zariski-topology} 有
$J + J' = B$。写 $1 = f + g$，其中 $f \in J$ 且 $g \in J'$。
则 $g$ 满足要求。
\end{proof}

\noindent
下面三个引理说明可以提升多项式的分解。

\begin{lemma}
\label{more-algebra-lemma-lift-factorization-monic}
设 $A$ 为环，$I \subset A$ 为理想，$f \in A[x]$ 为首一多项式。
设 $\overline{f} = \overline{g} \overline{h}$ 是 $f$ 在
$A/I[x]$ 中的一个分解，其中 $\overline{g}$ 和 $\overline{h}$
首一且在 $A/I[x]$ 中生成单位理想。则存在一个诱导同构
$A/I \to A'/IA'$ 的 \'etale 环映射 $A \to A'$，以及 $A'[x]$
中的分解 $f = g' h'$，其中 $g'$、$h'$ 首一，并提升 $A/I$
上的给定分解。
\end{lemma}

\begin{proof}
我们将从先前证明的关于泛分解的结果推出这一点；不过，仍鼓励读者寻找一个
不使用该技巧的证明。设
$\deg(\overline{g}) = n$ 且 $\deg(\overline{h}) = m$，
从而 $\deg(f) = n + m$。对某些
$\alpha_1, \ldots, \alpha_{n + m} \in A$，写
$f = x^{n + m} + \sum \alpha_i x^{n + m - i}$。考虑
Algebra, Example \ref{algebra-example-factor-polynomials-etale} 中的环映射
$$
R = \mathbf{Z}[a_1, \ldots, a_{n + m}]
\longrightarrow
S = \mathbf{Z}[b_1, \ldots, b_n, c_1, \ldots, c_m]
$$
。令 $R \to A$ 为把 $a_i$ 送到 $\alpha_i$ 的环映射，并令
$$
B = A \otimes_R S
$$
。由构造，$f$ 在 $B[x]$ 中的像 $f_B$ 有一个分解，记为
$f_B = g_B h_B$，其中
$g_B = x^n + \sum (1 \otimes b_i) x^{n - i}$，而 $h_B$ 类似。
写
$\overline{g} = x^n + \sum \overline{\beta}_i x^{n - i}$ 以及
$\overline{h} = x^m + \sum \overline{\gamma}_i x^{m - i}$。
$A$-代数映射
$$
B \longrightarrow A/I, \quad
1 \otimes b_i \mapsto \overline{\beta}_i, \quad
1 \otimes c_i \mapsto \overline{\gamma}_i
$$
把 $g_B$ 和 $h_B$ 分别映到 $A/I[x]$ 中的
$\overline{g}$ 和 $\overline{h}$。所展示的映射是满射；
以 $J \subset B$ 表示其核。由
Algebra, Example \ref{algebra-example-factor-polynomials-etale}
中的讨论，显然 $A \to B$ 在 $V(J) \subset \Spec(B)$
的所有点处都是 etale 的。按引理 \ref{more-algebra-lemma-localize-upstairs}
选取 $g \in B$，并考虑 $A$-代数 $B_g$。由于 $g$ 在
$B/J = A/I$ 中映到单位元，还得到一个 $A/I$-代数映射
$B_g/I B_g \to A/I$。由于 $A/I \to B_g/I B_g$ 是 \'etale 的，
$B_g/IB_g \to A/I$ 也是 \'etale 的
（Algebra, Lemma \ref{algebra-lemma-map-between-etale}）。
因此存在幂等元 $e \in B_g/I B_g$，使得
$A/I = (B_g/I B_g)_e$
（Algebra, Lemma
\ref{algebra-lemma-surjective-flat-finitely-presented}）。
选取 $e$ 的提升 $h \in B_g$。于是 $A \to A' = (B_g)_h$，
连同 $f_B = g_B h_B$ 的分解在 $A'$ 中的像，给出了该引理所提问题的解。
\end{proof}

\begin{lemma}
\label{more-algebra-lemma-lift-factorization-easy}
设 $A$ 为环，$I \subset A$ 为理想，$f \in A[x]$ 为首一多项式。
设 $\overline{f} = \overline{g} \overline{h}$ 是 $f$ 在
$A/I[x]$ 中的一个分解，并假设
\begin{enumerate}
\item $\overline{g}$ 的首项系数是 $A/I$ 中的可逆元素；
\item $\overline{g}$、$\overline{h}$ 在 $A/I[x]$ 中生成单位理想。
\end{enumerate}
则存在一个诱导同构 $A/I \to A'/IA'$ 的 \'etale 环映射
$A \to A'$，以及 $A'[x]$ 中的分解 $f = g' h'$，
它提升 $A/I$ 上的给定分解。
\end{lemma}

\begin{proof}
应用引理 \ref{more-algebra-lemma-lift-invertible-element}，可假设
$\overline{g}$ 的首项系数是某个可逆元素 $u \in A$ 的约化。
于是可以用 $\overline{u}^{-1}\overline{g}$ 替换
$\overline{g}$，并用 $\overline{u}\overline{h}$ 替换
$\overline{h}$。故可假设 $\overline{g}$ 首一。由于 $f$ 首一，
得到 $\overline{h}$ 也首一。此时结论由引理
\ref{more-algebra-lemma-lift-factorization-monic} 推出。
\end{proof}

\begin{example}
\label{more-algebra-example-cannot-lift-factorization}
本例说明，引理 \ref{more-algebra-lemma-lift-factorization-easy} 中关于首项系数的
条件不可去掉。确切地说，令 $A = \mathbf{Z}$、
$I = 4\mathbf{Z}$、$f = 1$，并在 $A/I[x]$ 中令
$\bar{g} = \bar{h} = 2x^2 + 2x + 1$。若 $A \to A'$ 是 \'etale 的，
且 $A'/I A' = A/I$，则 $A'$ 的 $2$-进完备化同构于
$\mathbf{Z}_2$。然而，$\mathbf{Z}_2[x]$ 中任意模 $4$ 同余于
$2x^2 + 2x + 1$ 的多项式 $g'$ 都没有 $\mathbf{Z}_2[x]$ 中的
多项式乘法逆元，故该引理的结论不成立。
\end{example}

\begin{lemma}
\label{more-algebra-lemma-separate-image-closed-from-closed}
设 $R \to S$ 为环映射，$I \subset R$ 为 $R$ 的理想，
$J \subset S$ 为 $S$ 的理想。若 $V(J)$ 在 $\Spec(R)$ 中的像的闭包
与 $V(I)$ 不交，则存在元素 $f \in R$，它在 $R/I$ 中映到 $1$，
而在 $S$ 中映到 $J$ 的一个元素。
\end{lemma}

\begin{proof}
设 $I' \subset R$ 为理想，使 $V(I')$ 是 $V(J)$ 的像的闭包。
由假设 $V(I) \cap V(I') = \emptyset$，故由
Algebra, Lemma \ref{algebra-lemma-Zariski-topology} 有
$I + I' = R$。写 $1 = g + f$，其中 $g \in I$ 且 $f \in I'$。
设 $f'$ 是 $f$ 在 $S$ 中的像，则
$V(f') \supset V(J)$。因此对某个 $n$ 有 $(f')^n \in J$，见
Algebra, Lemma \ref{algebra-lemma-Zariski-topology}。
用 $f^n$ 替换 $f$ 即得结论。
\end{proof}

\begin{lemma}
\label{more-algebra-lemma-helper-integral}
设 $I$ 为环 $A$ 的理想，$A \to B$ 为整环映射。
设 $b \in B$ 在 $B/IB$ 中映到幂等元。则存在首一
$f \in A[x]$，使得 $f(b) = 0$，且对某个 $d \geq 1$ 有
$f \bmod I = x^d(x - 1)^d$。
\end{lemma}

\begin{proof}
注意 $z = b^2 - b$ 是 $IB$ 的元素。由 Algebra, Lemma
\ref{algebra-lemma-integral-integral-over-ideal}，存在次数为 $d$
的首一多项式 $g(x) = x^d + \sum a_j x^j$，其中 $a_j \in I$，
使得在 $B$ 中 $g(z) = 0$。因此
$f(x) = g(x^2 - x) \in A[x]$ 是首一多项式，满足
$f(x) \equiv x^d(x - 1)^d \bmod I$，并且在 $B$ 中
$f(b) = 0$。
\end{proof}

\begin{lemma}
\label{more-algebra-lemma-lift-idempotent-upstairs}
设 $A$ 为环，$I \subset A$ 为理想，$A \to B$ 为整环映射，
$\overline{e} \in B/IB$ 为幂等元。则存在一个诱导同构
$A/I \to A'/IA'$ 的 \'etale 环映射 $A \to A'$，以及提升
$\overline{e}$ 的幂等元 $e' \in B \otimes_A A'$。
\end{lemma}

\begin{proof}
选取 $\overline{e}$ 的提升 $y \in B$，并按引理
\ref{more-algebra-lemma-helper-integral} 为 $y$ 选取 $f \in A[x]$。
由引理 \ref{more-algebra-lemma-lift-factorization-easy}，可找到一个诱导同构
$A/I \to A'/IA'$ 的 \'etale 环映射 $A \to A'$，并使
$f = gh$ 在 $A[x]$ 中成立，其中
$g(x) = x^d \bmod IA'$ 且
$h(x) = (x - 1)^d \bmod IA'$。用 $A'$ 替换 $A$ 后，
可假设该分解定义在 $A$ 上。此时
$b_1 = g(y) \in B$ 是 $\overline{e}^d = \overline{e}$ 的提升，而
$b_2 = h(y) \in B$ 是
$(\overline{e} - 1)^d = (-1)^d (1 - \overline{e})^d = (-1)^d(1 - \overline{e})$
的提升，并且 $b_1b_2 = 0$。于是
$(b_1, b_2)B/IB = B/IB$，而
$V(b_1, b_2) \subset \Spec(B)$ 与 $V(IB)$ 不交。
由于 $\Spec(B) \to \Spec(A)$ 是闭映射（见
Algebra, Lemmas \ref{algebra-lemma-integral-going-up} and
\ref{algebra-lemma-going-up-closed}），由引理
\ref{more-algebra-lemma-separate-image-closed-from-closed} 可选取 $a \in A$，
使其在 $A/I$ 中映到可逆元素，而其在 $B$ 中的像属于
$(b_1, b_2)$。用局部化 $A_a$ 替换 $A$ 后，有
$(b_1, b_2) = B$。于是
$\Spec(B) = D(b_1) \amalg D(b_2)$；这是不交并，因为
$b_1b_2 = 0$；又覆盖 $\Spec(B)$，因为 $(b_1, b_2) = B$。
令 $e \in B$ 为对应于开闭子集 $D(b_1)$ 的幂等元，见
Algebra, Lemma \ref{algebra-lemma-disjoint-decomposition}。
由于 $b_1$ 是 $\overline{e}$ 的提升，而 $b_2$ 是
$\pm (1 - \overline{e})$ 的提升，由
Algebra, Lemma \ref{algebra-lemma-disjoint-decomposition}
中的唯一性断言，$e$ 是 $\overline{e}$ 的提升。
\end{proof}

\begin{lemma}
\label{more-algebra-lemma-lift-projective-module}
设 $A$ 为环，$I \subset A$ 为理想，
$\overline{P}$ 为有限投射 $A/I$-模。则存在一个诱导同构
$A/I \to A'/IA'$ 的 \'etale 环映射 $A \to A'$，以及提升
$\overline{P}$ 的有限投射 $A'$-模 $P'$。
\end{lemma}

\begin{proof}
可选取整数 $n$ 以及直和分解
$(A/I)^{\oplus n} = \overline{P} \oplus \overline{K}$，
其中 $\overline{K}$ 是某个 $R/I$-模。选取与直和分支
$\overline{P}$ 对应的投影算子 $\overline{p}$ 的一个提升
$\varphi : A^{\oplus n} \to A^{\oplus n}$。
令 $f \in A[x]$ 为 $\varphi$ 的特征多项式，并令
$B = A[x]/(f)$。由 Cayley--Hamilton
（Algebra, Lemma \ref{algebra-lemma-charpoly}），存在映射
$B \to \text{End}_A(A^{\oplus n})$，它把 $x$ 送到 $\varphi$。
对每个 $\mathfrak p \supset I$，$f$ 在 $\kappa(\mathfrak p)$
中的像是 $(x - 1)^rx^{n - r}$，其中 $r$ 是
$\overline{P} \otimes_{A/I} \kappa(\mathfrak p)$ 的维数。
因此对所有 $\mathfrak p \supset I$，
$(x - 1)^nx^n$ 在 $B \otimes_A \kappa(\mathfrak p)$ 中映到零。
故 $x(1 - x)$ 包含于 $B/IB$ 的每个素理想，从而对某个
$N \geq 1$，$x^N(1 - x)^N$ 包含于 $IB$。
由此 $x^N + (1 - x)^N$ 是 $B/IB$ 中的单位元，并且
$$
\overline{e} = \text{image of }\frac{x^N}{x^N + (1 - x)^N}\text{ 于 }B/IB
$$
是幂等元，因为这两个断言在
$\mathbf{Z}[x]/(x^N(x - 1)^N)$ 中都成立。
$\overline{e}$ 在 $\text{End}_{A/I}((A/I)^{\oplus n})$ 中的像是
$$
\frac{\overline{p}^N}{\overline{p}^N + (1 - \overline{p})^N} = \overline{p}
$$
，因为 $\overline{p}$ 是幂等元。按引理中的方式用一个 \'etale
扩张 $A'$ 替换 $A$ 后，可假设存在幂等元 $e \in B$，其在
$B/IB$ 中映到 $\overline{e}$，见引理
\ref{more-algebra-lemma-lift-idempotent-upstairs}。于是 $e$ 在映射
$$
B = A[x]/(f) \longrightarrow \text{End}_A(A^{\oplus n}).
$$
下的像是提升 $\overline{p}$ 的幂等元 $p$。令
$P = \Im(p)$ 即得结论。
\end{proof}

\begin{lemma}
\label{more-algebra-lemma-cotangent-complex-symmetric-algebra}
设 $A$ 为环，并设
$0 \to K \to A^{\oplus m} \to M \to 0$
为 $A$-模序列。考虑 $A$-代数
$C = \text{Sym}^*_A(M)$，以及由满射
$A^{\oplus m} \to M$ 给出的表示
$\alpha : A[y_1, \ldots, y_m] \to C$。则
$$
\NL(\alpha) =
(K \otimes_A C \to \bigoplus\nolimits_{j = 1, \ldots, m} C \text{d}y_j)
$$
（见 Algebra, Section \ref{algebra-section-netherlander}），
特别地 $\Omega_{C/A} = M \otimes_A C$。
\end{lemma}

\begin{proof}
令 $J = \Ker(\alpha)$。该引理断言
$J/J^2 \cong K \otimes_A C$。注意 $\alpha$ 是分次代数同态。
我们将证明在次数 $d$ 有
$(J/J^2)_d = K \otimes_A C_{d - 1}$。注意
$$
J_d = \Ker(\text{Sym}^d_A(A^{\oplus m}) \to \text{Sym}^d_A(M))
= \Im(K \otimes_A \text{Sym}^{d - 1}_A(A^{\oplus m})
\to \text{Sym}^d_A(A^{\oplus m})),
$$
见 Algebra, Lemma \ref{algebra-lemma-presentation-sym-exterior}。
由此 $(J^2)_d = \sum_{a + b = d} J_a \cdot J_b$ 是映射
$$
K \otimes_A K \otimes_A \text{Sym}^{d - 2}_A(A^{\otimes m})
\to \text{Sym}^d_A(A^{\oplus m}).
$$
的像。由上面引用的引理，映射
$K \otimes_A \text{Sym}^{d - 2}_A(A^{\otimes m}) \to
\text{Sym}^{d - 1}_A(A^{\oplus m})$
的余核是 $\text{Sym}^{d - 1}_A(M)$。
因此显然 $(J/J^2)_d = J_d/(J^2)_d$ 等于
\begin{align*}
\Coker(
K \otimes_A K \otimes_A \text{Sym}^{d - 2}_A(A^{\otimes m})
\to K \otimes_A \text{Sym}^{d - 1}_A(A^{\otimes m}))
& = K \otimes_A \text{Sym}^{d - 1}_A(M) \\
& = K \otimes_A C_{d -1}
\end{align*}
，正合所需。
\end{proof}

\begin{lemma}
\label{more-algebra-lemma-symmetric-algebra-smooth}
设 $A$ 为环，$M$ 为 $A$-模。则
$C = \text{Sym}_A^*(M)$ 在 $A$ 上光滑，当且仅当 $M$
是有限投射 $A$-模。
\end{lemma}

\begin{proof}
令 $\sigma : C \to A$ 为到 $C$ 的次数 $0$ 部分的投影。
于是 $J = \Ker(\sigma)$ 是次数 $> 0$ 的部分，并且作为 $A$-模有
$J/J^2 = M$。因此，若 $A \to C$ 光滑，则由
Algebra, Lemma \ref{algebra-lemma-section-smooth}，
$M$ 是有限投射 $A$-模。

\medskip\noindent
反过来，假设 $M$ 有限投射，并选取核为 $K$ 的满射
$A^{\oplus n} \to M$。当然，序列
$0 \to K \to A^{\oplus n} \to M \to 0$ 分裂，因为 $M$ 是投射模。
特别地，$K$ 是有限 $A$-模。因此 $C$ 在 $A$ 上为有限表示，
因为 $C$ 是 $A[x_1, \ldots, x_n]$ 关于由
$K \subset \bigoplus Ax_i$ 生成的理想的商。
引理 \ref{more-algebra-lemma-cotangent-complex-symmetric-algebra} 的计算表明，
$\NL_{C/A}$ 同伦等价于 $(K \to M) \otimes_A C$。
因此 $\NL_{C/A}$ 拟同构于置于次数 $0$ 的
$C \otimes_A M$；由 Algebra, Definition
\ref{algebra-definition-smooth}，这意味着 $C$ 在 $A$ 上光滑。
\end{proof}

\begin{lemma}
\label{more-algebra-lemma-lift-section-smooth-morphism}
设 $A$ 为环，$I \subset A$ 为理想。考虑交换图
$$
\xymatrix{
B \ar[rd] \\
A \ar[u] \ar[r] & A/I
}
$$
，其中 $B$ 是光滑 $A$-代数。则存在一个诱导同构
$A/I \to A'/IA'$ 的 \'etale 环映射 $A \to A'$，以及提升环映射
$B \to A/I$ 的 $A$-代数映射 $B \to A'$。
\end{lemma}

\begin{proof}
令 $J \subset B$ 为 $B \to A/I$ 的核，从而 $B/J = A/I$。
由 Algebra, Lemma \ref{algebra-lemma-application-NL-smooth}，序列
$$
0 \to I/I^2 \to J/J^2 \to \Omega_{B/A} \otimes_B B/J \to 0
$$
分裂正合。因此
$\overline{P} = J/(J^2 + IB) = \Omega_{B/A} \otimes_B B/J$
是有限投射 $A/I$-模。选取整数 $n$ 及直和分解
$A/I^{\oplus n} = \overline{P} \oplus \overline{K}$。
由引理 \ref{more-algebra-lemma-lift-projective-module}，可找到一个诱导同构
$A/I \to A'/IA'$ 的 \'etale 环映射 $A \to A'$，以及提升
$\overline{K}$ 的有限投射 $A$-模 $K$。我们可以且确实用 $A'$
替换 $A$。令
$B' = B \otimes_A \text{Sym}_A^*(K)$。由引理
\ref{more-algebra-lemma-symmetric-algebra-smooth}，
$A \to \text{Sym}_A^*(K)$ 光滑，故 $B \to B'$ 光滑，
进而 $A \to B'$ 光滑（见
Algebra, Lemmas \ref{algebra-lemma-base-change-smooth} and
\ref{algebra-lemma-locally-smooth}）。
此外，截面 $\text{Sym}^*_A(K) \to A$ 确定截面 $B' \to B$，
并令 $B' \to A/I$ 为复合 $B' \to B \to A/I$。
令 $J' \subset B'$ 为 $B' \to A/I$ 的核。有
$JB' \subset J'$ 且 $B \otimes_A K \subset J'$。
这些映射合并成同构
$$
(A/I)^{\oplus n} \cong J/(J^2 + IB) \oplus \overline{K}
\longrightarrow
J'/((J')^2 + IB')
$$
。因此，用 $B'$ 替换 $B$ 后，可假设
$J/(J^2 + IB) = \Omega_{B/A} \otimes_B B/J$
是秩为 $n$ 的自由 $A/I$-模。

\medskip\noindent
在此情形中，选取 $f_1, \ldots, f_n \in J$，使其像构成
$J/(J^2 + IB)$ 的一组基。考虑有限表示 $A$-代数
$C = B/(f_1, \ldots, f_n)$。注意有正合序列
$$
0 \to H_1(L_{C/A}) \to (f_1, \ldots, f_n)/(f_1, \ldots, f_n)^2
\to \Omega_{B/A} \otimes_B C \to \Omega_{C/A} \to 0
$$
，见 Algebra, Lemma \ref{algebra-lemma-exact-sequence-NL}
（注意 $H_1(L_{B/A}) = 0$，而 $\Omega_{B/A}$ 有限投射，
特别地平坦，故 Tor 群消失）。对 $B$ 的任意素理想
$\mathfrak q \supset J$，模 $\Omega_{B/A, \mathfrak q}$ 是秩为
$n$ 的自由模，因为 $\Omega_{B/A}$ 有限投射，且
$\Omega_{B/A} \otimes_B B/J$ 是秩为 $n$ 的自由模
（见 Algebra, Lemma \ref{algebra-lemma-finite-projective}）。
由 $f_1, \ldots, f_n$ 的选择，映射
$$
\left((f_1, \ldots, f_n)/(f_1, \ldots, f_n)^2\right)_{\mathfrak q}
\to
\Omega_{B/A, \mathfrak q}
$$
模 $J$ 后为满射。因此，这个局部环 $C_{\mathfrak q}$ 上的模映射
必为同构（由 Algebra, Lemma \ref{algebra-lemma-NAK}，该映射为满射；
而例如由 Algebra, Lemma \ref{algebra-lemma-fun}，由于
$$
((f_1, \ldots, f_n)/(f_1, \ldots, f_n)^2)_{\mathfrak q}
$$
可由 $n$ 个元素生成，该映射为单射）。故
$H_1(L_{C/A})_{\mathfrak q} = 0$ 且
$\Omega_{C/A, \mathfrak q} = 0$。由 Algebra, Lemma
\ref{algebra-lemma-smooth-at-point}，$A \to C$ 在 $C$ 中对应于
$\mathfrak q$ 的素理想 $\overline{\mathfrak q}$ 处光滑。
由于 $\Omega_{C/A, \mathfrak q} = 0$，它实际上在
$\overline{\mathfrak q}$ 处为 \'etale。因此 $A \to C$
在 $C$ 的所有包含 $JC$ 的素理想处都是 \'etale 的。
由引理 \ref{more-algebra-lemma-localize-upstairs}，可选取 $f \in C$，
它在 $C/JC$ 中映到单位元，并使 $A \to C_f$ 为 \'etale。
由 $f$ 的选择，仍有 $C_f/JC_f = A/I$。映射
$C_f/IC_f \to A/I$ 是满且 \'etale 的，见
Algebra, Lemma \ref{algebra-lemma-map-between-etale}。
因此 $A/I$ 同构于 $C_f/IC_f$ 在某个元素 $g \in C$ 处的局部化，
见 Algebra, Lemma
\ref{algebra-lemma-surjective-flat-finitely-presented}。
令 $A' = C_{fg}$，证明完成。
\end{proof}








\section{Zariski 对}
\label{more-algebra-section-zariski-pairs}

\noindent
在本节及下一节中，一个{\it 对}是指一个对 $(A, I)$，其中 $A$
是环，$I \subset A$ 是理想。{\it 对的态射}
$(A, I) \to (B, J)$ 是满足 $\varphi(I) \subset J$ 的环映射
$\varphi : A \to B$。

\begin{definition}
\label{more-algebra-definition-zariski-pair}
一个{\it Zariski 对}是满足 $I$ 包含于 $A$ 的 Jacobson 根中的对
$(A, I)$。
\end{definition}

\begin{lemma}
\label{more-algebra-lemma-idempotents-determined-modulo-radical}
设 $(A, I)$ 为 Zariski 对。则从 $A$ 的幂等元到 $A/I$ 的幂等元的映射
是单射。
\end{lemma}

\begin{proof}
局部环的幂等元只能是 $0$ 或 $1$。因此，由
Algebra, Lemma \ref{algebra-lemma-characterize-zero-local}，
一个幂等元由它在哪些极大理想处为零所确定。
\end{proof}

\begin{lemma}
\label{more-algebra-lemma-check-isomorphism-zariski}
设 $(A, I)$ 为 Zariski 对。设 $A \to B$ 为平坦、整且有限表示的
环映射，并且 $A/I \to B/IB$ 是同构。则 $A \to B$ 是同构。
\end{lemma}

\begin{proof}
由 Algebra, Lemma
\ref{algebra-lemma-characterize-finite-in-terms-of-integral}，环映射
$A \to B$ 是有限的。因此，由 Algebra, Lemma
\ref{algebra-lemma-finite-finitely-presented-extension}，$B$ 作为
$A$-模是有限表示的。进而，由 Algebra, Lemma
\ref{algebra-lemma-finite-projective}，$B$ 是有限局部自由 $A$-模。
由于模 $B$ 沿 $V(I)$ 的秩为 $1$（见 Algebra, Lemma
\ref{algebra-lemma-finite-projective} 中描述的秩函数），并且
$(A, I)$ 是 Zariski 对，我们推出它处处秩为 $1$。
由此 $A \to B$ 是同构：证明如下事实是一个有趣的练习：若环映射
$R \to S$ 使得 $S$ 是秩为 $1$ 的局部自由 $R$-模，则该映射是同构
（提示：考察局部环）。
\end{proof}

\begin{lemma}
\label{more-algebra-lemma-helper-finite}
设 $(A, I)$ 为 Zariski 对。设 $A \to B$ 为有限环映射。假设
\begin{enumerate}
\item $B/IB = B_1 \times B_2$ 是 $A/I$-代数的乘积；
\item $A/I \to B_1/IB_1$ 是满射；
\item $b \in B$ 在该乘积中的像为 $(1, 0)$。
\end{enumerate}
则存在首一多项式 $f \in A[x]$，使得 $f(b) = 0$，且对某个
$d \geq 1$ 有 $f \bmod I = (x - 1)x^d$。
\end{lemma}

\begin{proof}
由引理 \ref{more-algebra-lemma-lift-idempotent-upstairs}，可找到诱导同构
$A/I \to A'/IA'$ 的 \'etale 环映射 $A \to A'$，使得
$B' = B \otimes_A A'$ 包含一个幂等元 $e'$，它提升 $b$ 在
$B'/IB'$ 中的像。考虑相应的 $A'$-代数分解
$$
B' = B'_1 \times B'_2
$$
它在模 $I$ 约化后与引理中给定的分解相容。
映射 $A' \to B'_1$ 模 $IA'$ 后为满射。由 Nakayama 引理
（Algebra, Lemma \ref{algebra-lemma-NAK}），可找到 $i \in IA'$，
使得以 $A'_{1 + i}$ 替换 $A'$ 后，映射 $A' \to B'_1$ 为满射。
注意，$b$ 在第一分量中的像 $b'_1 \in B'_1$ 满足
$b'_1 - 1 \in IB'_1$。因此可选取映到 $b'_1 - 1$ 的
$a' \in IA'$。另一方面，$b$ 在第二分量中的像
$b'_2 \in B'_2$ 属于 $IB'_2$。由 Algebra, Lemma
\ref{algebra-lemma-integral-integral-over-ideal}，存在一个次数为
$d$ 的首一多项式 $g(x) = x^d + \sum a'_j x^j$，其中
$a'_j \in IA'$，并且 $g(b'_2) = 0$ 在 $B'_2$ 中成立。
因此，$b$ 的像 $b' = (b'_1, b'_2) \in B'$ 是多项式
$(x - 1 - a')g(x)$ 的根。我们得到
$$
(b' - 1)(b')^d \in \sum\nolimits_{j = 0, \ldots, d} IA' \cdot (b')^j
$$
我们断言，这蕴含 $B$ 中有
$$
(b - 1)b^d \in \sum\nolimits_{j = 0, \ldots, d} I \cdot b^j
$$
为此，只需证明环映射 $A \to A'$ 忠实平坦，因为该条件等价于
$(b - 1)b^d$ 的像在
$B/\sum_{j = 0, \ldots, d} Ib^j$ 中为零（使用 Algebra, Lemma
\ref{algebra-lemma-faithfully-flat-universally-injective}）。
映射 $A \to A'$ 是平坦的，因为它是 \'etale 的
（Algebra, Lemma \ref{algebra-lemma-etale}）。另一方面，诱导的谱映射
是开映射（见 Algebra, Proposition \ref{algebra-proposition-fppf-open}，
并使用前面引用的引理），而且其像包含 $V(I)$。由于 $I$ 包含于
$A$ 的 Jacobson 根中，结论随即成立。
\end{proof}

\begin{lemma}
\label{more-algebra-lemma-noetherian-zariski-jacobson-complement}
设 $(A, I)$ 为 Zariski 对，且 $A$ 为 Noether 环。设 $f \in I$。
则 $A_f$ 是 Jacobson 环。
\end{lemma}

\begin{proof}
我们使用 Algebra, Lemma
\ref{algebra-lemma-noetherian-dim-1-Jacobson} 的判别准则。
设 $\mathfrak p \subset A$ 为素理想，使得
$\mathfrak p_f = \mathfrak p A_f$ 是非极大的素理想。我们必须证明
$A_f/\mathfrak p_f = (A/\mathfrak p)_f$ 有无穷多个素理想。
以 $A/\mathfrak p$ 替换 $A$ 后，可以假设 $A$ 是整环，
$\dim A_f > 0$，而我们的目标是证明 $\Spec(A_f)$ 是无限集。
由于 $\dim A_f > 0$，可找到一个不包含 $f$ 的非零素理想
$\mathfrak q \subset A$。选取包含 $\mathfrak q$ 的极大理想
$\mathfrak m \subset A$。由于 $(A, I)$ 是 Zariski 对，有
$I \subset \mathfrak m$。故 $\mathfrak m \not = \mathfrak q$，且
$\dim(A_\mathfrak m) > 1$。于是由 Algebra, Lemma
\ref{algebra-lemma-Noetherian-local-domain-dim-2-infinite-opens}，
$\Spec((A_\mathfrak m)_f) \subset \Spec(A_f)$ 是无限集，证毕。
\end{proof}











\section{Hensel 对}
\label{more-algebra-section-henselian-pairs}

\noindent
第 \ref{more-algebra-section-lifting} 节的一些结果可以看作关于 Hensel 对的结果。
本节中，一个{\it 对}是一个对 $(A, I)$，其中 $A$ 是环且
$I \subset A$ 是理想。{\it 对的态射} $(A, I) \to (B, J)$ 是满足
$\varphi(I) \subset J$ 的环映射 $\varphi : A \to B$。与第
\ref{more-algebra-section-lifting} 节一样，给定 $A$ 上的对象 $\xi$，我们以
$\overline{\xi}$ 表示 $\xi$ 到 $A/I$ 上对象的``基变换''
（如果这有意义）。

\begin{definition}
\label{more-algebra-definition-henselian-pair}
一个{\it Hensel 对}是满足下列条件的对 $(A, I)$：
\begin{enumerate}
\item $I$ 包含于 $A$ 的 Jacobson 根中；
\item 对任意首一多项式 $f \in A[T]$ 以及分解
$\overline{f} = g_0h_0$，其中 $g_0, h_0 \in A/I[T]$ 首一且在
$A/I[T]$ 中生成单位理想，都存在 $A[T]$ 中的分解 $f = gh$，
其中 $g, h$ 首一，并且 $g_0 = \overline{g}$、
$h_0 = \overline{h}$。
\end{enumerate}
\end{definition}

\noindent
注意，若 $A$ 是局部环且 $I = \mathfrak m$ 是其极大理想，则
$(A, I)$ 是 Hensel 对当且仅当 $A$ 是 Hensel 局部环，见
Algebra, Lemma \ref{algebra-lemma-characterize-henselian}。
在引理 \ref{more-algebra-lemma-characterize-henselian-pair} 中，我们给出 Hensel 对的
若干等价刻画（并且以后还会继续补充）。

\begin{lemma}
\label{more-algebra-lemma-locally-nilpotent-henselian}
设 $(A, I)$ 为一个对，且 $I$ 局部幂零。则函子
$B \mapsto B/IB$ 诱导 $A$ 上 \'etale 代数范畴与 $A/I$ 上
\'etale 代数范畴之间的等价。此外，该对是 Hensel 对。
\end{lemma}

\begin{proof}
本质满性由 Algebra, Lemma \ref{algebra-lemma-lift-etale} 给出。
若 $B$、$B'$ 在 $A$ 上 \'etale，且 $B/IB \to B'/IB'$ 是
$A/I$-代数态射，则由 Algebra, Lemma
\ref{algebra-lemma-smooth-strong-lift} 可以提升该态射。
最后，设 $f, g : B \to B'$ 是两个满足
$f \bmod I = g \bmod I$ 的 $A$-代数映射。选取生成乘法映射
$B \otimes_A B \to B$ 之核的幂等元 $e \in B \otimes_A B$，见
Algebra, Lemmas \ref{algebra-lemma-diagonal-unramified} and
\ref{algebra-lemma-unramified}（后者用于说明 \'etale 蕴含非分歧）。
于是 $(f \otimes g)(e) \in IB'$。由于 $IB'$ 局部幂零
（Algebra, Lemma \ref{algebra-lemma-locally-nilpotent}），
由 Algebra, Lemma \ref{algebra-lemma-lift-idempotents} 可知
$(f \otimes g)(e) = 0$。故 $f = g$。

\medskip\noindent
显然，$I$ 包含于 $A$ 的 Jacobson 根中。设 $f \in A[T]$ 为首一
多项式，并设 $\overline{f} = g_0h_0$ 是
$\overline{f} = f \bmod I$ 的一个分解，其中
$g_0, h_0 \in A/I[T]$ 首一且在 $A/I[T]$ 中生成单位理想。
由引理 \ref{more-algebra-lemma-lift-factorization-monic}，存在诱导同构
$A/I \to A'/IA'$ 的 \'etale 环映射 $A \to A'$，使得该分解提升为
$A'$ 上首一多项式的分解。由上文可知 $A = A'$，故该分解定义在
$A$ 上。
\end{proof}

\begin{lemma}
\label{more-algebra-lemma-limit-henselian}
设 $A = \lim A_n$，其中 $(A_n)$ 是由满的转移映射组成的环的逆向系统，
并且这些转移映射的核局部幂零。令 $I_n = \Ker(A \to A_n)$。
则 $(A, I_n)$ 是 Hensel 对。
\end{lemma}

\begin{proof}
固定 $n$。设 $a \in A$ 是在 $A_n$ 中映到 $1$ 的元素。
由 Algebra, Lemma \ref{algebra-lemma-locally-nilpotent-unit}，
$a$ 在所有 $m \geq n$ 的 $A_m$ 中都映到单位元。因此 $a$ 是
$A$ 中的单位元。故由 Algebra, Lemma
\ref{algebra-lemma-contained-in-radical}，理想 $I_n$ 包含于
$A$ 的 Jacobson 根中。设 $f \in A[T]$ 为首一多项式，并设
$\overline{f} = g_nh_n$ 是 $\overline{f} = f \bmod I_n$ 的一个分解，
其中 $g_n, h_n \in A_n[T]$ 首一且在 $A_n[T]$ 中生成单位理想。
由引理 \ref{more-algebra-lemma-locally-nilpotent-henselian}，可以逐次把该分解提升为
$f \bmod I_m = g_m h_m$，其中对所有 $m \geq n$，
$g_m, h_m$ 都是 $A_m[T]$ 中的首一多项式。在每一步中，我们必须验证
所取提升 $g_m, h_m$ 在 $A_n[T]$ 中生成单位理想；这由
$g_n, h_n$ 的相应事实以及
$\Spec(A_n[T]) = \Spec(A_m[T])$ 得出，因为 $A_m \to A_n$ 的核
局部幂零。由于 $A = \lim A_m$，证明完成。
\end{proof}

\begin{lemma}
\label{more-algebra-lemma-complete-henselian}
设 $(A, I)$ 为一个对。若 $A$ 是 $I$-进完备的，则该对是 Hensel 对。
\end{lemma}

\begin{proof}
由 Algebra, Lemma \ref{algebra-lemma-radical-completion}，理想 $I$
包含于 $A$ 的 Jacobson 根中。设 $f \in A[T]$ 为首一多项式，且
$\overline{f} = g_0h_0$ 是 $\overline{f} = f \bmod I$ 的分解，
其中 $g_0, h_0 \in A/I[T]$ 首一且在 $A/I[T]$ 中生成单位理想。
由引理 \ref{more-algebra-lemma-locally-nilpotent-henselian}，可逐次将其提升为
$f \bmod I^n = g_n h_n$，其中对所有 $n \geq 1$，
$g_n, h_n$ 都是 $A/I^n[T]$ 中的首一多项式。
由于 $A = \lim A/I^n$，证明完成。
\end{proof}

\begin{lemma}
\label{more-algebra-lemma-helper-finite-type}
设 $(A, I)$ 为一个对。设 $A \to B$ 为有限型环映射，并且
$B/IB = C_1 \times C_2$，其中 $A/I \to C_1$ 是有限的。
令 $B'$ 为 $A$ 在 $B$ 中的整闭包。则可写成
$B'/IB' = C_1 \times C'_2$，使得映射 $B'/IB' \to B/IB$ 保持乘积分解，
并且存在映到 $C_1 \times C'_2$ 中 $(1, 0)$ 的 $g \in B'$，使
$B'_g \to B_g$ 为同构。
\end{lemma}

\begin{proof}
注意，$A \to B$ 在闭子集
$T = \Spec(C_1) \subset \Spec(B)$ 的每个素理想处都是拟有限的
（这可通过考察纤维环看出，见 Algebra, Definition
\ref{algebra-definition-quasi-finite}）。考虑拓扑空间的交换图
$$
\xymatrix{
\Spec(B) \ar[rr]_\phi \ar[rd]_\psi & & \Spec(B') \ar[ld]^{\psi'} \\
& \Spec(A)
}
$$
由 Algebra, Theorem \ref{algebra-theorem-main-theorem}，对每个
$\mathfrak p \in T$，存在 $h_\mathfrak p \in B'$，满足
$h_\mathfrak p \not \in \mathfrak p$，且 $B'_h \to B_h$ 是同构。
并集 $U = \bigcup D(h_\mathfrak p)$ 给出开集
$U \subset \Spec(B')$，使得 $\phi^{-1}(U) \to U$ 为同胚且
$T \subset \phi^{-1}(U)$。由于 $T$ 在 $\psi^{-1}(V(I))$ 中为开集，
可知 $\phi(T)$ 在 $U \cap (\psi')^{-1}(V(I))$ 中为开集。
因此 $\phi(T)$ 在 $(\psi')^{-1}(V(I))$ 中为开集。
另一方面，由于 $C_1$ 在 $A/I$ 上有限，它在 $B'$ 上也是有限的。
故由 Algebra, Lemmas \ref{algebra-lemma-going-up-closed} and
\ref{algebra-lemma-integral-going-up}，$\phi(T)$ 是 $\Spec(B')$ 的闭子集。
我们推出 $\Spec(B'/IB') \supset \phi(T)$ 既开又闭。由
Algebra, Lemma \ref{algebra-lemma-disjoint-implies-product}，得到相应的
乘积分解 $B'/IB' = C'_1 \times C'_2$。通过考察谱上的情形可见，映射
$B'/IB' \to B/IB$ 把 $C'_1$ 映入 $C_1$，把 $C'_2$ 映入 $C_2$
（提示：$\phi(T)$ 的逆像恰为 $T$；略去一些细节）。
选取映到 $C'_1 \times C'_2$ 中 $(1, 0)$ 的 $g \in B'$，使得
$D(g) \subset U$；这是可能的，因为 $\Spec(C'_1)$ 与
$\Spec(C'_2)$ 在 $\Spec(B')$ 中是不交闭集，且 $\Spec(C'_1)$ 包含于
$U$。于是 $B'_g \to B_g$ 在谱上定义同胚，并在局部环上定义同构
（由上面对 $U$ 的选取）。因此它是同构，例如这可由
Algebra, Lemma \ref{algebra-lemma-characterize-zero-local} 推出。
最后，由此有 $C'_1 = C_1$，证明完成。
\end{proof}

\begin{lemma}
\label{more-algebra-lemma-characterize-henselian-pair}
\begin{reference}
\cite[Chapter XI]{Henselian} and \cite[Proposition 1]{Gabber-henselian}
\end{reference}
设 $(A, I)$ 为一个对。下列条件等价：
\begin{enumerate}
\item $(A, I)$ 是 Hensel 对；
\item 给定 \'etale 环映射 $A \to A'$ 以及 $A$-代数映射
$\sigma : A' \to A/I$，存在提升 $\sigma$ 的 $A$-代数映射
$A' \to A$；
\item 对任意有限 $A$-代数 $B$，映射 $B \to B/IB$ 在幂等元上
诱导双射；
\item 对任意整 $A$-代数 $B$，映射 $B \to B/IB$ 在幂等元上
诱导双射；
\item （Gabber）$I$ 包含于 $A$ 的 Jacobson 根中，并且每个形如
$$
f(T) = T^n(T - 1) + a_n T^n + \ldots + a_1 T + a_0
$$
的首一多项式 $f(T) \in A[T]$，其中
$a_n, \ldots, a_0 \in I$ 且 $n \ge 1$，都有一个根
$\alpha \in 1 + I$。
\end{enumerate}
此外，第 (5) 项中的根是唯一的。
\end{lemma}

\begin{proof}
假设 (2) 成立。则 $I$ 包含于 $A$ 的 Jacobson 根中；否则，存在一个
模 $I$ 同余于 $1$ 的非单位元 $f \in A$，而映射 $A \to A_f$
将与 (2) 矛盾。若 $B$ 在 $A$ 上整，则 $IB \subset B$ 包含于
$B$ 的 Jacobson 根中，因为由 Algebra, Lemmas
\ref{algebra-lemma-going-up-closed} and
\ref{algebra-lemma-integral-going-up}，$\Spec(B) \to \Spec(A)$ 是闭映射。
因此，由引理 \ref{more-algebra-lemma-idempotents-determined-modulo-radical}，
从 $B$ 的幂等元到 $B/IB$ 的幂等元的映射是单射。另一方面，
由于 (2) 成立，由引理 \ref{more-algebra-lemma-lift-idempotent-upstairs}，
$B/IB$ 的每个幂等元都提升为 $B$ 的幂等元。因此 (2) 蕴含 (4)。

\medskip\noindent
(4) $\Rightarrow$ (3) 是平凡的。

\medskip\noindent
假设 (3)。设 $\mathfrak m$ 为极大理想，并考虑有限映射
$A \to B = A/(I \cap \mathfrak m)$。$B \to B/IB$ 在幂等元上诱导
双射这一条件蕴含 $I \subset \mathfrak m$（否则
$B = A/I \times A/\mathfrak m$ 且 $B/IB = A/I$）。因此 $I$ 包含于
$A$ 的 Jacobson 根中。设 $f \in A[T]$ 首一，并假设给定分解
$\overline{f} = g_0h_0$，其中 $g_0, h_0 \in A/I[T]$ 首一且在
$A/I[T]$ 中生成单位理想。令 $B = A[T]/(f)$。令
$\overline{e}$ 为 $B/IB$ 中对应于 $A$-代数分解
$$
B/IB = A/I[T]/(g_0) \times A/I[T]/(h_0)
$$
的幂等元。由假设 (3)，存在提升 $\overline{e}$ 的幂等元
$e \in B$。这给出乘积分解
$$
B = eB \times (1 - e)B
$$
注意，$B$ 作为 $A$-模是秩为 $\deg(f)$ 的自由模。因此 $eB$ 和
$(1 - e)B$ 是有限局部自由 $A$-模。然而，由于 $eB$ 和
$(1 - e)B$ 在 $A/I$ 上分别有常秩 $\deg(g_0)$ 和 $\deg(h_0)$，
我们发现它们在 $\Spec(A)$ 上也如此。由此
\begin{align*}
f & = \text{CharPol}_A(T : B \to B) \\
& =
\text{CharPol}_A(T : eB \to eB)
\text{CharPol}_A(T : (1 - e)B \to (1 - e)B)
\end{align*}
是首一多项式的分解，模 $I$ 后约化为给定的分解。这里
$\text{CharPol}_A$ 表示有限局部自由 $A$-模之自同态的特征多项式。
若该模自由，则 $\text{CharPol}_A$ 定义为相应矩阵的特征多项式；
一般情形使用 Algebra, Lemma \ref{algebra-lemma-standard-covering}
进行粘合。略去细节。故 (3) 蕴含 (1)。

\medskip\noindent
假设 (1)。令 $f$ 如 (5) 所述。$f \bmod I$ 分解为 $T^n$ 与
$T - 1$ 的乘积；由定义 \ref{more-algebra-definition-henselian-pair}，它提升为
首一多项式的分解 $f = gh$，其中 $g$ 与 $h$ 首一。于是 $h$ 必须次数为 $1$，从而
$f$ 有一个模 $I$ 约化为 $1$ 的根。最后，由 Hensel 对的定义，
$I$ 包含于 Jacobson 根中。因此 (1) 蕴含 (5)。

\medskip\noindent
在证明最后一步前，先说明 (5) 中的根 $\alpha$ 若存在则唯一。
事实上，由 $f(T)$ 的明确形状，有 $f'(\alpha) \in 1 + I$，其中
$f'$ 是 $f$ 关于 $T$ 的导数。初等论证给出
$$
f(T) = f(\alpha + T - \alpha) =
f(\alpha) + f'(\alpha) \cdot (T - \alpha) \bmod
(T - \alpha)^2 A[T]
$$
这表明，$f(T)$ 的任意另一个根 $\alpha' \in 1 + I$ 都满足
$0 = f(\alpha') - f(\alpha) = (\alpha' - \alpha)(1 + i)$，其中某个
$i \in I$。由于 $1 + i$ 是 $A$ 中的单位元，故
$\alpha = \alpha'$。

\medskip\noindent
假设 (5)。我们证明 (2) 成立；换言之，对每个 \'etale 映射
$A \to A'$，模 $I$ 的每个截面 $\sigma : A' \to A/I$ 都提升为
截面 $A' \to A$。由于 $A \to A'$ 是 \'etale 的，截面 $\sigma$
确定 $A/I$-代数的分解
\begin{equation}
\label{more-algebra-equation-GCHP}
A'/IA' \cong A/I \times C
\end{equation}
事实上，由 Algebra, Lemma \ref{algebra-lemma-map-between-etale}，满环映射
$A'/IA' \to A/I$ 是 \'etale 的；继而由 Algebra, Lemma
\ref{algebra-lemma-surjective-flat-finitely-presented} 得到所需幂等元。
我们将证明，该分解提升为 $A$-代数的分解
\begin{equation}
\label{more-algebra-equation-GCHP-want}
A' \cong A'_1 \times A'_2
\end{equation}
并且 $A'_1$ 在 $A$ 上整。于是 $A \to A'_1$ 既整又 \'etale，且
$A/I \to A'_1/IA'_1$ 是同构；故由引理
\ref{more-algebra-lemma-check-isomorphism-zariski}，$A \to A'_1$ 是同构
（这里还用到 \'etale 环映射平坦且有限表示，见 Algebra, Lemma
\ref{algebra-lemma-etale}）。

\medskip\noindent
令 $B'$ 为 $A$ 在 $A'$ 中的整闭包。由引理
\ref{more-algebra-lemma-helper-finite-type}，可把它相容于
(\ref{more-algebra-equation-GCHP}) 地分解为 $A/I$-代数
\begin{equation}
\label{more-algebra-equation-dec-mod-I}
B'/IB' \cong A/I \times C'
\end{equation}
并可找到提升 $(1, 0)$ 的 $b \in B'$，使得
$B'_b \to A'_b$ 是同构。若分解
(\ref{more-algebra-equation-dec-mod-I}) 提升为 $A$-代数的分解
\begin{equation}
\label{more-algebra-equation-want-2}
B' \cong B'_1 \times B'_2
\end{equation}
则诱导的分解 $A' = A'_1 \times A'_2$ 给出所需的
(\ref{more-algebra-equation-GCHP-want})：事实上，由于 $b$ 是 $B'_1$ 中的单位元
（略去细节），有 $B'_1 \cong A'_1$，因而 $A'_1$ 在 $A$ 上整。

\medskip\noindent
选取包含 $b$ 的有限 $A$-子代数 $B'' \subset B'$（注意，$B'$ 的
任意有限生成 $A$-子代数都在 $A$ 上有限）。扩大 $B''$ 后，可假设
$b$ 在 $B''/IB''$ 中映到给出如下分解的幂等元：
\begin{equation}
\label{more-algebra-equation-again-dec-mod-I}
B''/IB'' \cong C''_1 \times C''_2
\end{equation}
由于 $B'_b \cong A'_b$，可知 $B'_b$ 在 $A$ 上是有限型的。
设 $B'_b$ 由 $b_1/b^n, \ldots, b_t/b^n$ 在 $A$ 上生成；扩大
$B''$，使得 $b_1, \ldots, b_t \in B''$。于是
$B''_b \to B'_b$ 既满又单，故是同构。特别地，
$C''_1 = A/I$！因此 $A/I \to C''_1$ 是同构，尤其是满射。
由引理 \ref{more-algebra-lemma-helper-finite}，可找到形如
$$
f(T) = T^n(T - 1) + a_n T^n + \ldots + a_1 T + a_0
$$
的 $f(T) \in A[T]$，其中 $a_n, \ldots, a_0 \in I$、
$n \ge 1$，并且 $f(b) = 0$。特别地，$B'$ 是一个
$A[T]/(f)$-代数。由 (5)，存在 $f$ 的根 $a \in 1 + I$。
这产生乘积分解 $A[T]/(f) = A[T]/(T - a) \times D$，它与
$B'/IB'$ 的分裂 (\ref{more-algebra-equation-dec-mod-I}) 相容。由此诱导的
$B'$ 的分裂即为所需的 (\ref{more-algebra-equation-want-2})。
\end{proof}

\begin{lemma}
\label{more-algebra-lemma-change-ideal-henselian-pair}
设 $A$ 为环。设 $I, J \subset A$ 为满足 $V(I) = V(J)$ 的理想。
则 $(A, I)$ 是 Hensel 对当且仅当 $(A, J)$ 是 Hensel 对。
\end{lemma}

\begin{proof}
对任意整环映射 $A \to B$，有 $V(IB) = V(JB)$。故 $B/IB$ 与
$B/JB$ 的幂等元一一对应
（Algebra, Lemma \ref{algebra-lemma-disjoint-decomposition}）。
因此，$B \to B/IB$ 在幂等元集合上诱导双射当且仅当
$B \to B/JB$ 在幂等元集合上诱导双射。结论由引理
\ref{more-algebra-lemma-characterize-henselian-pair} 得出。
\end{proof}

\begin{lemma}
\label{more-algebra-lemma-integral-over-henselian-pair}
设 $(A, I)$ 为 Hensel 对，并设 $A \to B$ 为整环映射。
则 $(B, IB)$ 是 Hensel 对。
\end{lemma}

\begin{proof}
这立即由引理 \ref{more-algebra-lemma-characterize-henselian-pair} 中 Hensel 对的
第四个刻画以及整环映射的复合仍为整映射这一事实得出。
\end{proof}

\begin{lemma}
\label{more-algebra-lemma-henselian-henselian-pair}
设 $I \subset J \subset A$ 为环 $A$ 的理想。下列条件等价：
\begin{enumerate}
\item $(A, I)$ 与 $(A/I, J/I)$ 都是 Hensel 对；
\item $(A, J)$ 是 Hensel 对。
\end{enumerate}
\end{lemma}

\begin{proof}
假设 (1)。令 $B$ 为整 $A$-代数。考虑环映射
$$
B \to B/IB \to B/JB
$$
由引理 \ref{more-algebra-lemma-characterize-henselian-pair}，两个箭头在幂等元上都
诱导双射，故其复合也如此。由引理
\ref{more-algebra-lemma-characterize-henselian-pair}，$(A, J)$ 是 Hensel 对。

\medskip\noindent
反之，假设 (2) 成立。由引理 \ref{more-algebra-lemma-integral-over-henselian-pair}，
$(A/I, J/I)$ 是 Hensel 对。令 $B$ 为整 $A$-代数。考虑环映射
$$
B \to B/IB \to B/JB
$$
由引理 \ref{more-algebra-lemma-characterize-henselian-pair}，复合与第二个箭头在
幂等元上诱导双射，故第一个箭头也是如此。因此，再由该引理，
$(A, I)$ 是 Hensel 对。
\end{proof}

\begin{lemma}
\label{more-algebra-lemma-sum-henselian}
设 $A$ 为环，且 $(A, I)$ 与 $(A, I')$ 都是 Hensel 对。
则 $(A, I + I')$ 是 Hensel 对。
\end{lemma}

\begin{proof}
由引理 \ref{more-algebra-lemma-integral-over-henselian-pair}，对
$(A/I, (I' + I)/I)$ 是 Hensel 对。因此结论由引理
\ref{more-algebra-lemma-henselian-henselian-pair} 得出。
\end{proof}

\begin{lemma}
\label{more-algebra-lemma-product-henselian-pairs}
设 $J$ 为集合，且 $\{ (A_j, I_j)\}_{j \in J}$ 为一族对。
则 $(\prod_{j \in J} A_j, \prod_{j\in J} I_j)$ 是 Hensel 对，
当且仅当每个 $(A_j, I_j)$ 都是 Hensel 对。
\end{lemma}
 
\begin{proof}
对每个 $j \in J$，投影 $\prod_{j \in J} A_j \rightarrow A_j$
是整环映射。因此，若
$(\prod_{j \in J} A_j, \prod_{j\in J} I_j)$ 是 Hensel 对，
则引理 \ref{more-algebra-lemma-integral-over-henselian-pair} 证明每个
$(A_j, I_j)$ 都是 Hensel 对。

\medskip\noindent
反之，假设每个 $(A_j, I_j)$ 都是 Hensel 对。对任意
$x \in \prod_{j \in J} I_j$，元素 $1 + x$ 在
$\prod_{j \in J} A_j$ 中是单位元，因为由 Algebra, Lemma
\ref{algebra-lemma-contained-in-radical} 和定义
\ref{more-algebra-definition-henselian-pair}，这在每个分量上都成立。
再次由 Algebra, Lemma \ref{algebra-lemma-contained-in-radical}，
$\prod_{j \in J} I_j$ 包含于 $\prod_{j \in J} A_j$ 的 Jacobson 根中。
继续逐分量考察可知，对每个首一多项式
$f \in (\prod_{j \in J} A_j)[T]$，以及每个分解
$\overline{f} = g_0h_0$，其中首一多项式
$g_0, h_0 \in (\prod_{j \in J} A_j / \prod_{j \in J} I_j)[T] =
(\prod_{j \in J} A_j/I_j)[T]$ 在
$(\prod_{j \in J} A_j / \prod_{j \in J} I_j)[T]$ 中生成单位理想，
都存在 $(\prod_{j \in J} A_j)[T]$ 中的分解 $f = gh$，其中
$g$、$h$ 首一且分别约化为 $g_0$、$h_0$。综上，由定义
\ref{more-algebra-definition-henselian-pair}，
$(\prod_{j \in J} A_j, \prod_{j\in J} I_j)$ 是 Hensel 对。
\end{proof}

\begin{lemma}
\label{more-algebra-lemma-limits-henselian}
Hensel 性在对的极限下保持。更确切地说，设 $J$ 为预序集，并设
$(A_j, I_j)$ 是 $J$ 上 Hensel 对的逆向系统。则
$A = \lim A_j$ 配以理想 $I = \lim I_j$ 所得的对 $(A, I)$ 是
Hensel 对。
\end{lemma}

\begin{proof}
由 Categories, Lemma
\ref{categories-lemma-limits-products-equalizers}，只需考虑乘积和等化子。
对乘积，断言由引理 \ref{more-algebra-lemma-product-henselian-pairs} 得出。
因此，考虑等化子图
$$
\xymatrix{
(A, I) \ar[r] & (A', I')
\ar@<1ex>[r]^{\varphi} \ar@<-1ex>[r]_{\psi}
&
(A'', I'')  
}
$$
其中对 $(A', I')$ 与 $(A'', I'')$ 都是 Hensel 对。为验证
$(A, I)$ 也是 Hensel 对，我们使用引理
\ref{more-algebra-lemma-characterize-henselian-pair} 中 Gabber 的判别准则。
$1 + I$ 的每个元素都是 $A$ 中的单位元，因为由单位元之逆的唯一性，
这一点可在 $(A', I')$ 中检验。因此 $I$ 包含于 $A$ 的 Jacobson 根中，
见 Algebra, Lemma \ref{algebra-lemma-contained-in-radical}。
于是，令
$$
f(T) = T^{N - 1}(T - 1) + a_{N - 1} T^{N - 1} + \dotsb + a_1 T + a_0
$$
为 $A[T]$ 中的多项式，其中
$a_{N - 1}, \dotsc, a_0 \in I$ 且 $N \ge 1$。$f(T)$ 在 $A'[T]$
中的像有唯一根 $\alpha' \in 1 + I'$，其在 $A''[T]$ 中的进一步像也
同样如此。因此，由唯一性，$\varphi(\alpha') = \psi(\alpha')$；
也就是说，$\alpha'$ 在 $1 + I$ 中定义了 $f(T)$ 的一个根，正合所需。
\end{proof}

\begin{lemma}
\label{more-algebra-lemma-filtered-colimits-henselian}
Hensel 性在对的滤过余极限下保持。更确切地说，设 $J$ 为有向集，
并设 $(A_j, I_j)$ 是 $J$ 上 Hensel 对的系统。则
$A = \colim A_j$ 配以理想 $I = \colim I_j$ 所得的对 $(A, I)$
是 Hensel 对。
\end{lemma}

\begin{proof}
若 $u \in 1 + I$，则对某个 $j \in J$，$u$ 是某个
$u_j \in 1 + I_j$ 的像。由 Algebra, Lemma
\ref{algebra-lemma-contained-in-radical} 以及 $I_j$ 包含于 $A_j$ 的
Jacobson 根中这一假设，$u_j$ 在 $A_j$ 中可逆。因此 $u$ 在 $A$ 中
可逆。故由该引理，$I$ 包含于 $A$ 的 Jacobson 根中。

\medskip\noindent
设 $f \in A[T]$ 为首一多项式，并设 $\overline{f} = g_0 h_0$ 是一个
分解，其中 $g_0, h_0 \in A/I[T]$ 首一且生成 $A/I[T]$ 中的单位理想。
对某些 $g'_0, h'_0 \in A/I[T]$，写成
$1 = g_0 g'_0 + h_0 h'_0$。由于 $A = \colim A_j$ 与
$A/I = \colim A_j/I_j$ 都是滤过余极限，可以找到 $j \in J$、
$f_j \in A_j$ 以及分解 $\overline{f}_j = g_{j, 0} h_{j, 0}$，其中
$g_{j, 0}, h_{j, 0} \in A_j/I_j[T]$ 首一，并且对某些
$g'_{j, 0}, h'_{j, 0} \in A_j/I_j[T]$ 有
$1 = g_{j, 0} g'_{j, 0} + h_{j, 0} h'_{j, 0}$；这里
$f_j, g_{j, 0}, h_{j, 0}, g'_{j, 0}, h'_{j, 0}$ 分别映到
$f, g_0, h_0, g'_0, h'_0$。由于 $(A_j, I_j)$ 是 Hensel 对，
可把 $\overline{f}_j = g_{j, 0} h_{j, 0}$ 提升为 $A_j$ 上的分解；
取其在 $A$ 中的像，得到 $A$ 中的相应分解。因此 $(A, I)$ 是 Hensel 对。
\end{proof}

\begin{example}[Moret-Bailly]
\label{more-algebra-example-moret-bailly}
若余极限不是滤过的，则引理 \ref{more-algebra-lemma-filtered-colimits-henselian} 不成立。
例如，取 Hensel 对 $(\mathbf{Z}_p, (p))$ 与
$(\mathbf{Z}_p, (p))$ 的余积，则得到 $(A, pA)$，其中
$A = \mathbf{Z}_p \otimes_\mathbf{Z} \mathbf{Z}_p$。这不是 Hensel 对：
$A/pA = \mathbf{F}_p$，故若 $(A, pA)$ 是 Hensel 对，则 $A$ 必须是
局部环。然而，$\Spec(A)$ 不连通；例如，对奇素数 $p$，有非平凡幂等元
$$
(1/2 \otimes 1)
\left(1 \otimes 1 - (1 + p)^{-1}u \otimes u\right)
$$
其中 $u \in \mathbf{Z}_p$ 是 $1 + p$ 的一个平方根。略去一些细节。
\end{example}

\begin{lemma}
\label{more-algebra-lemma-largest-ideal-henselian}
设 $A$ 为环。存在最大的理想 $I \subset A$，使得 $(A, I)$ 是 Hensel 对。
\end{lemma}

\begin{proof}
结合引理 \ref{more-algebra-lemma-henselian-henselian-pair}、
\ref{more-algebra-lemma-sum-henselian} 和 \ref{more-algebra-lemma-filtered-colimits-henselian}。
\end{proof}

\begin{lemma}
\label{more-algebra-lemma-irreducible-henselian-pair-connected}
设 $(A, I)$ 为 Hensel 对。设 $\mathfrak p \subset A$ 为素理想。
则 $V(\mathfrak p + I)$ 连通。
\end{lemma}

\begin{proof}
由引理 \ref{more-algebra-lemma-integral-over-henselian-pair}，
$(A/\mathfrak p, I + \mathfrak p/\mathfrak p)$ 是 Hensel 对。
故只需证明：若 $(A, I)$ 是 Hensel 对且 $A$ 是整环，则
$\Spec(A/I) = V(I)$ 连通。否则，由 Algebra, Lemma
\ref{algebra-lemma-characterize-spec-connected}，$A/I$ 有非平凡幂等元。
由引理 \ref{more-algebra-lemma-characterize-henselian-pair}，这将蕴含 $A$ 有非平凡
幂等元，矛盾。
\end{proof}








\section{对的 Hensel 化}
\label{more-algebra-section-henselization}

\noindent
我们继续第 \ref{more-algebra-section-henselian-pairs} 节开始的讨论。

\begin{lemma}
\label{more-algebra-lemma-henselization}
包含函子
$$
\text{category of henselian pairs}
\longrightarrow
\text{category of pairs}
$$
有左伴随 $(A, I) \mapsto (A^h, I^h)$。
\end{lemma}

\begin{proof}
设 $(A, I)$ 为一个对。考虑由 \'etale 环映射 $A \to B$ 组成的范畴
$\mathcal{C}$，其中 $A/I \to B/IB$ 是同构。我们将证明范畴
$\mathcal{C}$ 是有向的，并且令
$A^h = \colim_{B \in \mathcal{C}} B$、$I^h = IA^h$ 给出所需伴随。

\medskip\noindent
先证明 $\mathcal{C}$ 是有向的
（Categories, Definition \ref{categories-definition-directed}）。
它非空，因为 $\text{id} : A \to A$ 是一个对象。若 $B$ 与 $B'$ 是
$\mathcal{C}$ 的两个对象，则 $B'' = B \otimes_A B'$ 是
$\mathcal{C}$ 的对象（使用 Algebra, Lemma \ref{algebra-lemma-etale}），
并且存在态射 $B \to B''$ 与 $B' \to B''$。设
$f, g : B \to B'$ 是 $\mathcal{C}$ 中两个对象之间的映射。则一个
余等化子是
$$
(B' \otimes_{f, B, g} B') \otimes_{(B' \otimes_A B')} B'
$$
由 Algebra, Lemmas \ref{algebra-lemma-etale} and
\ref{algebra-lemma-map-between-etale}，它在 $A$ 上 \'etale。
因此范畴 $\mathcal{C}$ 是有向的。

\medskip\noindent
由于对 $\mathcal{C}$ 的所有对象 $B$ 都有 $B/IB = A/I$，故
$A^h/I^h = A^h/IA^h = \colim B/IB = \colim A/I = A/I$。

\medskip\noindent
下面证明 $A^h = \colim_{B \in \mathcal{C}} B$ 配以
$I^h = IA^h$ 是 Hensel 对。为此，验证引理
\ref{more-algebra-lemma-characterize-henselian-pair} 的条件 (2)。具体地，设给定
\'etale 环映射 $A^h \to A'$ 以及 $A^h$-代数映射
$\sigma : A' \to A^h/I^h$。则存在 $B \in \mathcal{C}$ 及
\'etale 环映射 $B \to B'$，使得 $A' = B' \otimes_B A^h$，见
Algebra, Lemma \ref{algebra-lemma-etale}。由于 $A^h/I^h = A/I$，
映射 $\sigma$ 诱导 $A$-代数映射 $s : B' \to A/I$。于是作为
$A/I$-代数有 $B'/IB' = A/I \times C$，其中 $C$ 是由 $s$ 诱导的
映射 $B'/IB' \to A/I$ 的核。令 $g \in B'$ 映到
$(1, 0) \in A/I \times C$。则 $B \to B'_g$ 是 \'etale 的，且
$A/I \to B'_g/IB'_g$ 是同构；即 $B'_g$ 是 $\mathcal{C}$ 的对象。
因此得到典范映射 $B'_g \to A^h$，使得
$$
\vcenter{
\xymatrix{
B'_g \ar[r] & A^h \\
B \ar[u] \ar[ur]
}
}
\quad\text{且}\quad
\vcenter{
\xymatrix{
B' \ar[r] \ar[rrd]_s & B'_g \ar[r] & A^h \ar[d] \\
& & A/I
}
}
$$
交换。这诱导与 $\sigma$ 相容的映射
$A' = B' \otimes_B A^h \to A^h$，正合所需。

\medskip\noindent
设 $(A, I) \to (A', I')$ 为对的态射，且 $(A', I')$ 是 Hensel 对。
我们证明存在唯一分解 $A \to A^h \to A'$，从而完成证明。
事实上，对 $\mathcal{C}$ 中每个 $A \to B$，环映射
$A' \to B' = A' \otimes_A B$ 是 \'etale 的，并诱导同构
$A'/I' \to B'/I'B'$。故由引理
\ref{more-algebra-lemma-characterize-henselian-pair}，存在截面
$\sigma_B : B' \to A'$。给定 $\mathcal{C}$ 中的态射
$B_1 \to B_2$，我们断言下图交换：
$$
\xymatrix{
B'_1 \ar[rr] \ar[rd]_{\sigma_{B_1}} & &
B'_2 \ar[ld]^{\sigma_{B_2}} \\
& A'
}
$$
只需证明，对 $\mathcal{C}$ 中每个 $B$，截面 $\sigma_B$ 是唯一的
$A'$-代数映射 $B' \to A'$。由 Algebra, Lemma
\ref{algebra-lemma-diagonal-unramified}，对某个环 $R$ 有
$B' \otimes_{A'} B' = B' \times R$。在本情形中，由
$B'/I'B' = A'/I'$ 有 $R/I'R = 0$。因此，给定两个 $A'$-代数映射
$\sigma_B, \sigma_B' : B' \to A'$，则
$e = (\sigma_B \otimes \sigma_B')(0, 1) \in A'$ 是包含于 $I'$ 的
幂等元。由引理 \ref{more-algebra-lemma-idempotents-determined-modulo-radical}，
$e = 0$。故 $\sigma_B = \sigma_B'$，正合所需。
利用交换性，得到
$$
A^h = \colim_{B \in \mathcal{C}} B \to
\colim_{B \in \mathcal{C}} A' \otimes_A B \xrightarrow{\colim \sigma_B} A'
$$
正合所需。映射 $\sigma_B$ 的唯一性还保证该映射唯一。
故 $(A, I) \mapsto (A^h, I^h)$ 是所需伴随。
\end{proof}

\begin{lemma}
\label{more-algebra-lemma-henselization-flat}
设 $(A, I)$ 为一个对。令 $(A^h, I^h)$ 如引理
\ref{more-algebra-lemma-henselization} 所述。则 $A \to A^h$ 平坦，
$I^h = IA^h$，并且对所有 $n$，$A/I^n \to A^h/I^nA^h$ 都是同构。
\end{lemma}

\begin{proof}
在引理 \ref{more-algebra-lemma-henselization} 的证明中，我们已经看到，$A^h$ 是
满足 $A/I \to B/IB$ 为同构的 \'etale $A$-代数 $B$ 的滤过余极限，
并且 $I^h = IA^h$。由于 \'etale 环映射平坦
（Algebra, Lemma \ref{algebra-lemma-etale}），由 Algebra, Lemma
\ref{algebra-lemma-colimit-flat} 可知 $A \to A^h$ 平坦。
由于每个 $A \to B$ 都平坦，映射 $A/I^n \to B/I^nB$ 也都是同构
（例如由 Algebra, Lemma \ref{algebra-lemma-lift-basis}）。
取余极限，得到 $A/I^n = A^h/I^nA^h$，正合所需。
\end{proof}

\begin{lemma}
\label{more-algebra-lemma-henselization-local-ring}
\begin{slogan}
对的 Hensel 化与局部环的 Hensel 化相容。
\end{slogan}
引理 \ref{more-algebra-lemma-henselization} 的函子把局部环
$(A, \mathfrak m)$ 映为其 Hensel 化。
\end{lemma}

\begin{proof}
令 $(A^h, \mathfrak m^h)$ 为引理 \ref{more-algebra-lemma-henselization} 中构造的
对 $(A, \mathfrak m)$ 的 Hensel 化。由引理
\ref{more-algebra-lemma-henselization-flat}，$\mathfrak m^h = \mathfrak m A^h$ 是极大
理想；又因它包含于 Jacobson 根中，我们推出 $A^h$ 是以
$\mathfrak m^h$ 为极大理想的局部环。说完这些，有两种方法完成证明。

\medskip\noindent
第一种证明：注意，在 Algebra, Lemma
\ref{algebra-lemma-henselization} 的证明中作为余极限的构造，与引理
\ref{more-algebra-lemma-henselization} 中用于构造 $A^h$ 的余极限相同。
第二种证明：Hensel 化 $A \to S$ 与引理
\ref{more-algebra-lemma-henselization} 的 $A \to A^h$ 都是局部环同态；
$S$ 与 $A^h$ 都是 \'etale $A$-代数的滤过余极限；$S$ 与 $A^h$
都是 Hensel 局部环；并且 $S$ 与 $A^h$ 的剩余域都等于
$\kappa(\mathfrak m)$（第二种情形使用引理
\ref{more-algebra-lemma-henselization-flat}）。因此，由 Algebra, Lemma
\ref{algebra-lemma-uniqueness-henselian}，它们典范同构。
\end{proof}

\begin{lemma}
\label{more-algebra-lemma-henselization-Noetherian-pair}
\begin{slogan}
Noether 对的 Hensel 化仍是 Noether 对，并且具有相同的完备化。
\end{slogan}
设 $(A, I)$ 为一个对，且 $A$ 为 Noether 环。令 $(A^h, I^h)$ 如引理
\ref{more-algebra-lemma-henselization} 所述。则 $I$-进完备化之间的映射
$$
A^\wedge \to (A^h)^\wedge
$$
是同构。此外，$A^h$ 是 Noether 环，映射
$A \to A^h \to A^\wedge$ 都平坦，且
$A^h \to A^\wedge$ 忠实平坦。
\end{lemma}

\begin{proof}
第一个断言立即由引理 \ref{more-algebra-lemma-henselization-flat} 得出，事实上无需
假设 $A$ 是 Noether 环。在引理 \ref{more-algebra-lemma-henselization} 的证明中，
我们已经看到，$A^h$ 是满足 $A/I \to B/IB$ 为同构的 \'etale
$A$-代数 $B$ 的滤过余极限。对每个这样的 $A \to B$，诱导映射
$A^\wedge \to B^\wedge$ 是同构
（见引理 \ref{more-algebra-lemma-henselization-flat} 的证明）。
由 Algebra, Lemma \ref{algebra-lemma-completion-flat}，对每个 $B$，
环映射 $B \to A^\wedge = B^\wedge = (A^h)^\wedge$ 都平坦。
因此，由 Algebra, Lemma \ref{algebra-lemma-colimit-rings-flat}，
$A^h \to A^\wedge = (A^h)^\wedge$ 平坦。由于
$I^h = IA^h$ 包含于 $A^h$ 的 Jacobson 根中，并且
$A^h \to A^\wedge$ 诱导同构 $A^h/I^h \to A/I$，由
Algebra, Lemma \ref{algebra-lemma-ff}，$A^h \to A^\wedge$ 忠实平坦。
由 Algebra, Lemma \ref{algebra-lemma-completion-Noetherian-Noetherian}，
环 $A^\wedge$ 是 Noether 环。故由 Algebra, Lemma
\ref{algebra-lemma-descent-Noetherian}，$A^h$ 是 Noether 环。
\end{proof}

\begin{lemma}
\label{more-algebra-lemma-henselization-colimit}
设 $(A, I) = \colim (A_i, I_i)$ 是对的滤过余极限。引理
\ref{more-algebra-lemma-henselization} 的函子给出 $A^h = \colim A_i^h$ 和
$I^h = \colim I_i^h$。
\end{lemma}

\noindent
对非滤过余极限，此引理不成立，见例 \ref{more-algebra-example-moret-bailly}。

\begin{proof}
由 Categories, Lemma \ref{categories-lemma-adjoint-exact}，
$(A^h, I^h)$ 是系统 $(A_i^h, I_i^h)$ 在 Hensel 对范畴中的余极限。
因此，对 Hensel 对 $(B, J)$，有
$$
\Mor((A^h, I^h), (B, J)) =
\lim \Mor((A_i^h, I_i^h), (B, J)) =
\Mor(\colim (A_i^h, I_i^h), (B, J))
$$
这里余极限取在对的范畴中。由于该余极限滤过，在对的范畴中有
$\colim (A_i^h, I_i^h) = (\colim A_i^h, \colim I_i^h)$；略去细节。
再次利用余极限滤过，可知这是 Hensel 对（引理
\ref{more-algebra-lemma-filtered-colimits-henselian}）。故由 Yoneda 引理，
$(A^h, I^h) = (\colim A_i^h, \colim I_i^h)$。
\end{proof}

\begin{lemma}
\label{more-algebra-lemma-henselization-change-ideal}
\begin{slogan}
一个对的 Hensel 化只依赖于理想的根。
\end{slogan}
设 $A$ 为环，$I$ 和 $J$ 是其理想。若 $V(I) = V(J)$，则引理
\ref{more-algebra-lemma-henselization} 的函子对 $(A, I)$ 与 $(A, J)$ 产生同一个环。
\end{lemma}

\begin{proof}
令 $(A', IA')$ 为从对 $(A, I)$ 出发由引理
\ref{more-algebra-lemma-henselization} 产生的对，见引理
\ref{more-algebra-lemma-henselization-flat}。令 $(A'', JA'')$ 为从对 $(A, J)$ 出发由
引理 \ref{more-algebra-lemma-henselization} 产生的对。由引理
\ref{more-algebra-lemma-change-ideal-henselian-pair}，$(A', JA')$ 与
$(A'', IA'')$ 都是 Hensel 对。由该构造的泛性质，得到唯一的
$A$-代数映射 $A'' \to A'$ 与 $A' \to A''$。唯一性表明它们互为逆映射。
\end{proof}

\begin{lemma}
\label{more-algebra-lemma-henselization-integral}
\begin{slogan}
Hensel 化与整基变换交换。
\end{slogan}
设 $(A, I) \to (B, J)$ 为满足 $V(J) = V(IB)$ 的对的映射。
令 $(A^h , I^h) \to (B^h, J^h)$ 为 Hensel 化上诱导的映射
（引理 \ref{more-algebra-lemma-henselization}）。若 $A \to B$ 是整的，则诱导映射
$A^h \otimes_A B \to B^h$ 是同构。
\end{lemma}

\begin{proof}
由引理 \ref{more-algebra-lemma-henselization-change-ideal}，可假设 $J = IB$。
由引理 \ref{more-algebra-lemma-integral-over-henselian-pair}，对
$(A^h \otimes_A B, I^h(A^h \otimes_A B))$ 是 Hensel 对。
由 $(B^h, IB^h)$ 的泛性质，得到映射
$B^h \to A^h \otimes_A B$。我们略去证明该映射是引理中映射的逆。
\end{proof}


\begin{lemma}
\label{more-algebra-lemma-CRT-henselization}
设 $I_1, I_2, \dots, I_n \subset A$ 为环 $A$ 的理想。假设对
$i \not = j$、$1 \leq i, j \leq n$ 有 $I_i + I_j = A$。
以 $A^h$ 表示对 $(A, I_1 \cap \ldots \cap I_n)$ 的 Hensel 化，
以 $A_i^h$ 表示对 $(A, I_i)$ 的 Hensel 化。则自然态射
$$
A^h \to \prod\nolimits_{i = 1, \ldots, n} A^h_i
$$
是同构。
\end{lemma}

\begin{proof}
该映射来自构造的函子性。若 $n = 1$，结论立即成立。若 $n > 1$，
注意，由理想所满足的条件，
$(I_1 \cap \ldots \cap I_{n - 1}) + I_n = A$。因此，只要证明
$n = 2$ 的情形，引理对所有 $n$ 即由归纳法成立。

\medskip\noindent
考虑 $n = 2$。选取 $f_i \in I_i$，使得 $f_1 + f_2 = 1$。
回顾 $A^h$ 构造为如下 $A$-代数 $B$ 的余极限：$A \to B$ 是
\'etale 的，且 $A/I_1 \cap I_2 \to B/(I_1 \cap I_2)B$ 是同构。
见引理 \ref{more-algebra-lemma-henselization} 的证明。以 $Hens$ 表示这样的
$A$-代数 $B$ 所成的范畴。类似地，以 $Hens_i$、$i = 1, 2$ 表示如下
\'etale $A$-代数 $B$ 的范畴：$A/I_i \to B/I_iB$ 是同构。

\medskip\noindent
考虑由满足 $f_2$ 可逆的那些 $B$ 组成的子范畴
$Hens_1(f_2) \subset Hens_1$。这是余终子范畴，因为若 $B$ 属于
$Hens_1$，则 $B_{f_2}$ 也属于 $Hens_1$
（提示：$B_{f_2}/I_1B_{f_2} = (B/I_1B)_{f_2} = B/I_1B$，因为
$f_2$ 模 $I_1$ 是单位元）。类似地，子范畴
$Hens_2(f_1) \subset Hens_2$ 是余终的。因此，为计算 $A_1^h$ 与
$A_2^h$，可以使用这些较小的范畴，见 Categories, Lemma
\ref{categories-lemma-cofinal}。

\medskip\noindent
考虑函子
$$
Hens_1(f_2) \times Hens_2(f_1) \to Hens,\quad (B_1, B_2) \mapsto B_1 \times B_2
$$
该函子有定义，因为
\begin{align*}
(B_1 \times B_2)/(I_1 \cap I_2)(B_1 \times B_2)
& =
B_1/(I_1 \cap I_2)B_1 \times B_2/(I_1 \cap I_2)B_2 \\
& =
B_1/I_1B_1 \times B_2/I_2B_2 \\
& = A/I_1 \times A/I_2 \\
& = A/(I_1 \cap I_2)
\end{align*}
第二个等号成立，是因为 $f_2 \in I_2$ 在 $B_1$ 中可逆，从而
$(I_1 \cap I_2)B_1 = I_1B_1 \cap I_2B_1 = I_1B_1$
（这里还使用了 Algebra, Lemma
\ref{algebra-lemma-flat-intersect-ideals}）。最后一个等号由
Algebra, Lemma \ref{algebra-lemma-chinese-remainder} 得出。
要完成证明，只需证明所显示的函子余终（见上面引用的引理）。
Categories, Definition \ref{categories-definition-cofinal} 的第一个条件成立，
因为 $Hens$ 的对象 $B$ 映到 $B_{f_2} \times B_{f_1}$，而后者在该函子的
像中。第二个条件总是成立，因为若有
$B \to B_1 \times B_2$ 和 $B \to B'_1 \times B'_2$，其中
$B_1, B'_1 \in Hens_1(f_2)$ 且 $B_2, B'_2 \in Hens_2(f_1)$，则可令
$B''_1 = B_1 \otimes_B B'_1$，它属于 $Hens_1(f_2)$；类似地，令
$B''_2 = B_2 \otimes_B B''_2$，它属于 $Hens_2(f_1)$，从而找到适配于
下列交换图的映射 $B \to B''_1 \times B''_2$：
$$
\xymatrix{
& B \ar[ld] \ar[d] \ar[rd] \\
B_1 \times B_2 \ar[r] & B_1'' \times B_2'' & B_1' \times B_2' \ar[l]
}
$$
这就完成了证明。
\end{proof}








\section{提升与 Hensel 对}
\label{more-algebra-section-lifting-henselian-pairs}

\noindent
本节主要结合第 \ref{more-algebra-section-lifting} 节和第
\ref{more-algebra-section-henselian-pairs} 节的结果。

\begin{lemma}
\label{more-algebra-lemma-lift-finite-projective-module}
设 $(R, I)$ 为 Hensel 对。映射
$$
P \longrightarrow P/IP
$$
诱导有限投射 $R$-模的同构类集合与有限投射 $R/I$-模的同构类集合
之间的双射。特别地，任意有限投射 $R/I$-模都同构于某个有限投射
$R$-模 $P$ 的 $P/IP$。
\end{lemma}

\begin{proof}
先证明最后一个断言。令 $\overline{P}$ 为有限投射 $R/I$-模。
由引理 \ref{more-algebra-lemma-lift-projective-module}，可以找到某个在 $R$ 上
\'etale 的 $R'$ 上的有限投射模 $P'$，其中 $R/I = R'/IR'$，且
$P'/IP'$ 同构于 $\overline{P}$。由于 $(R, I)$ 是 Hensel 对，
\'etale 环映射 $R \to R'$ 有截面 $\tau : R' \to R$
（引理 \ref{more-algebra-lemma-characterize-henselian-pair}）。令
$P = P' \otimes_{R', \tau} R$，便推出 $P/IP$ 同构于
$\overline{P}$。当然，这说明引理陈述中的映射是满的。

\medskip\noindent
单射性。假设 $P_1$ 与 $P_2$ 是有限投射 $R$-模，并且作为
$R/I$-模有 $P_1/IP_1 \cong P_2/IP_2$。由于 $P_1$ 投射，可找到提升
该同构的 $R$-模映射 $u : P_1 \to P_2$。由 Nakayama 引理
（Algebra, Lemma \ref{algebra-lemma-NAK}），$u$ 是满射。
类似地，可找到满射 $v : P_2 \to P_1$。由 Algebra, Lemma
\ref{algebra-lemma-fun}，映射 $v \circ u$ 是同构，因此 $u$ 是同构。
\end{proof}

\begin{lemma}
\label{more-algebra-lemma-finite-etale-equivalence}
设 $(A, I)$ 为 Hensel 对。函子 $B \to B/IB$ 确定有限 \'etale
$A$-代数与有限 \'etale $A/I$-代数之间的等价。
\end{lemma}

\begin{proof}
设 $B, B'$ 为两个在 $A$ 上有限 \'etale 的 $A$-代数。
则 $B' \to B'' = B \otimes_A B'$ 也有限 \'etale
（Algebra, Lemmas \ref{algebra-lemma-etale} and
\ref{algebra-lemma-base-change-integral}）。现在，下列对象之间存在
$1$ 对 $1$ 的对应：
\begin{enumerate}
\item $A$-代数映射 $B \to B'$；
\item $B' \to B''$ 的截面；
\item $B''$ 的幂等元 $e$，使得 $B' \to B'' \to eB''$ 是同构。
\end{enumerate}
(2) 与 (3) 之间的双射把 $\sigma : B'' \to B'$ 映到 $e$，其中
$(1 - e)$ 是生成 $\sigma$ 之核的幂等元；它的存在性由
Algebra, Lemmas \ref{algebra-lemma-map-between-etale} and
\ref{algebra-lemma-surjective-flat-finitely-presented} 给出。
类似地，$A/I$-代数映射 $B/IB \to B'/IB'$ 与 $B''/IB''$ 的幂等元
$\overline{e}$ 之间也有对应，其中
$B'/IB' \to B''/IB'' \to \overline{e}(B''/IB'')$ 是同构。
然而，$B''/IB''$ 的每个幂等元 $\overline{e}$ 都唯一提升为
$B''$ 的幂等元 $e$（引理 \ref{more-algebra-lemma-characterize-henselian-pair}）。
此外，若 $B'/IB' \to \overline{e}(B''/IB'')$ 是同构，则由 Nakayama
引理（Algebra, Lemma \ref{algebra-lemma-NAK}），
$B' \to eB''$ 也是同构。由此可见，该函子是全忠实的。

\medskip\noindent
本质满性。设 $A/I \to C$ 为有限 \'etale 映射。由 Algebra, Lemma
\ref{algebra-lemma-lift-etale}，存在 \'etale 映射 $A \to B$，使得
$B/IB \cong C$。令 $B'$ 为 $A$ 在 $B$ 中的整闭包。由引理
\ref{more-algebra-lemma-helper-finite-type}，对某个环 $C'$ 有
$B'/IB' = C \times C'$，并且对某个映到
$(1, 0) \in C \times C'$ 的 $g \in B'$，有 $B'_g \cong B_g$。
由于幂等元可以提升（引理 \ref{more-algebra-lemma-characterize-henselian-pair}），
得到 $B' = B'_1 \times B'_2$，其中 $C = B'_1/IB'_1$ 且
$C' = B'_2/IB'_2$。$g$ 在 $B'_1$ 中的像可逆。于是
$B_g = B'_g = B'_1 \times (B_2)_g$，这蕴含 $A \to B'_1$ 是
\'etale 的。我们推出 $B'_1$ 在 $A$ 上有限 \'etale
（例如，由 Algebra, Lemma
\ref{algebra-lemma-characterize-finite-in-terms-of-integral}，整且
\'etale 蕴含有限 \'etale），证明完成。
\end{proof}

\begin{lemma}
\label{more-algebra-lemma-lift-smooth-henselian}
设 $R$ 为环，$S$ 为光滑 $R$-代数。假设 $A$ 是 $R$-代数，且
$(A,I)$ 是 Hensel 对。则任意 $R$-代数映射 $S \to A/I$ 都可提升为
$R$-代数映射 $S \to A$。
\end{lemma}

\begin{proof}
设 $\tau : S \to A/I$ 为 $R$-代数映射。注意，由 Algebra, Lemma
\ref{algebra-lemma-base-change-smooth}，$S \otimes_R A$ 是光滑
$A$-代数。因此，由引理 \ref{more-algebra-lemma-lift-section-smooth-morphism}，
可将诱导映射 $S \otimes_R A \to A/I$ 提升为 $A$-代数同态
$S \otimes_R A \to A'$，其中 $A \to A'$ 是 \'etale 的，并诱导同构
$A/I \to A'/IA'$。由于 $(A, I)$ 是 Hensel 对，存在 $A$-代数映射
$A' \to A$，见引理 \ref{more-algebra-lemma-characterize-henselian-pair}。
复合 $S \to S \otimes_R A \to A' \to A$ 即为所需提升。
\end{proof}

\begin{lemma}
\label{more-algebra-lemma-lim-finite-projective-gives-finite-projective}
设 $A = \lim A_n$ 是环的逆向系统 $(A_n)$ 的极限。假设给定
$A_n$-模 $M_n$ 以及 $A_{n + 1}$-模映射 $M_{n + 1} \to M_n$。假设：
\begin{enumerate}
\item 转移映射 $A_{n + 1} \to A_n$ 是满射且核局部幂零；
\item $M_1$ 是有限投射 $A_1$-模；
\item $M_n$ 是有限平坦 $A_n$-模；
\item 这些映射诱导同构
$M_{n + 1} \otimes_{A_{n + 1}} A_n \to M_n$。
\end{enumerate}
则 $M = \lim M_n$ 是有限投射 $A$-模，并且对所有 $n$，
$M \otimes_A A_n \to M_n$ 都是同构。
\end{lemma}

\begin{proof}
由引理 \ref{more-algebra-lemma-limit-henselian}，对
$(A, \Ker(A \to A_1))$ 是 Hensel 对。由引理
\ref{more-algebra-lemma-lift-finite-projective-module}，可选取有限投射 $A$-模 $P$
以及同构 $P \otimes_A A_1 \to M_1$。由于 $P$ 投射，可逐次将
$A$-模映射 $P \to M_1$ 提升为 $A$-模映射
$P \to M_2$、$P \to M_3$，依此类推。由此得到映射
$$
P \longrightarrow M
$$
由于 $P$ 有限投射，对某个 $m \geq 0$ 和某个 $A$-模 $Q$，可写成
$A^{\oplus m} = P \oplus Q$。由于 $A = \lim A_n$，得到
$P = \lim P \otimes_A A_n$。因此，为证明所显示的 $A$-模映射是同构，
只需证明映射 $P \otimes_A A_n \to M_n$ 是同构。由引理
\ref{more-algebra-lemma-lift-projective} 可知 $M_n$ 是有限投射模。由引理
\ref{more-algebra-lemma-isomorphic-finite-projective-lifts}，映射
$P \otimes_A A_n \to M_n$ 是同构。
\end{proof}






\section{绝对整闭包}
\label{more-algebra-section-absolute-integral-closure}

\noindent
下面给出我们的定义。

\begin{definition}
\label{more-algebra-definition-absolutely-integrally-closed}
若每个首一多项式 $f \in A[T]$ 都是一次因子的乘积，则称环 $A$
{\it 绝对整闭}。
\end{definition}

\noindent
请注意：$f$ 可能能够以许多不同方式写成一次因子的乘积。

\begin{lemma}
\label{more-algebra-lemma-absolutely-integrally-closed}
设 $A$ 为环。下列条件等价：
\begin{enumerate}
\item $A$ 绝对整闭；
\item 任意首一多项式 $f \in A[T]$ 在 $A$ 中都有根。
\end{enumerate}
\end{lemma}

\begin{proof}
略。
\end{proof}


\begin{lemma}
\label{more-algebra-lemma-absolutely-integrally-closed-quotient-localization}
设 $A$ 绝对整闭。
\begin{enumerate}
\item $A$ 的任意商环 $A/I$ 都绝对整闭。
\item 任意局部化 $S^{-1}A$ 都绝对整闭。
\end{enumerate}
\end{lemma}

\begin{proof}
略。
\end{proof}

\begin{lemma}
\label{more-algebra-lemma-integrally-closed-in-absolutely-integrally-closed}
设 $A$ 为环。设 $S \subset A$ 为由非零因子组成的乘法子集。
若 $S^{-1}A$ 绝对整闭，且 $A \subset S^{-1}A$ 在 $S^{-1}A$ 中整闭，
则 $A$ 绝对整闭。
\end{lemma}

\begin{proof}
略。
\end{proof}

\begin{lemma}
\label{more-algebra-lemma-normal-domain-absolutely-integrally-closed}
设 $A$ 为正规整环。则 $A$ 绝对整闭当且仅当其分式域代数闭。
\end{lemma}

\begin{proof}
注意，一个域代数闭当且仅当它作为环绝对整闭。因此该引理由引理
\ref{more-algebra-lemma-absolutely-integrally-closed-quotient-localization} 和
\ref{more-algebra-lemma-integrally-closed-in-absolutely-integrally-closed} 得出。
\end{proof}

\begin{lemma}
\label{more-algebra-lemma-construct-absolute-integral-closure}
对任意环 $A$，存在扩张 $A \subset B$，使得：
\begin{enumerate}
\item $B$ 是有限自由 $A$-代数的滤过余极限；
\item $B$ 作为 $A$-模自由；
\item $B$ 绝对整闭。
\end{enumerate}
\end{lemma}

\begin{proof}
令 $I$ 为 $A$ 上所有首一多项式所成的集合。对 $i \in I$，以
$x_i$ 表示一个变量，以 $P_i$ 表示变量 $x_i$ 中相应的首一多项式。
然后令
$$
F(A) = A[x_i; i \in I]/(P_i; i \in I)
$$
正如记号所表明的，$F$ 是从环范畴到自身的函子。注意
$A \subset F(A)$，$F(A)$ 作为 $A$-模自由，并且 $F(A)$ 是有限自由
$A$-代数的滤过余极限。然后取
$$
B = \colim F^n(A)
$$
其中转移映射是包含映射
$F^n(A) \subset F(F^n(A)) = F^{n + 1}(A)$。系数属于 $B$ 的任意
首一多项式实际上对某个 $n$ 具有属于 $F^n(A)$ 的系数，因而由构造，
它在 $F^{n + 1}(A)$ 中有解。由引理
\ref{more-algebra-lemma-absolutely-integrally-closed}，这蕴含 $B$ 绝对整闭。
其他性质的证明从略。
\end{proof}

\begin{lemma}
\label{more-algebra-lemma-absolutely-integrally-closed-strictly-henselian}
设 $A$ 绝对整闭。设 $\mathfrak p \subset A$ 为素理想。
则局部环 $A_\mathfrak p$ 是严格 Hensel 的。
\end{lemma}

\begin{proof}
由引理 \ref{more-algebra-lemma-absolutely-integrally-closed-quotient-localization}，
可假设 $A$ 是局部环，且 $\mathfrak p$ 是其极大理想。再由引理
\ref{more-algebra-lemma-absolutely-integrally-closed-quotient-localization}，剩余域代数闭。
每个首一多项式都完全分解为一次因子，故可直接应用
Algebra, Definition \ref{algebra-definition-henselian}。
\end{proof}

\begin{lemma}
\label{more-algebra-lemma-absolutely-integrally-closed-henselian-pair}
设 $A$ 绝对整闭。设 $I \subset A$ 为理想。若（且仅若）下列条件成立，
则 $(A, I)$ 是 Hensel 对：
\begin{enumerate}
\item $I$ 包含于 $A$ 的 Jacobson 根中；
\item $A \to A/I$ 在幂等元上诱导双射。
\end{enumerate}
\end{lemma}

\begin{proof}
设 $f \in A[T]$ 为首一多项式，并设
$f \bmod I = g_0 h_0$ 是 $A/I$ 上的一个分解，其中 $g_0$、$h_0$ 首一，
且 $g_0$ 与 $h_0$ 在 $A/I[T]$ 中生成单位理想。这意味着
$$
A/I[T]/(f) = A/I[T]/(g_0) \times A/I[T]/(h_0)
$$
以 $e \in A/I[T]/(f)$ 表示对应于右侧环中幂等元 $(1, 0)$ 的元素。
写成 $f = (T - a_1) \ldots (T - a_d)$，其中 $a_i \in A$。
对每个 $i \in \{1, \ldots, d\}$，得到 $A$-代数映射
$\varphi_i : A[T]/(f) \to A$、$T \mapsto a_i$，它诱导类似的
$A/I$-代数映射 $\overline{\varphi}_i : A/I[T]/(f) \to A/I$。
令 $e_i = \overline{\varphi}_i(e) \in A/I$。这些都是幂等元。
由假设 (2)，可将 $e_i$ 提升为 $A$ 中的幂等元。这意味着可将
$A = \prod A_j$ 写成环的有限乘积，使得在 $A_j/IA_j$ 中，每个
$e_i$ 都是 $0$ 或 $1$。略去一些细节。注意，$A_j$ 作为 $A$ 的
因子环绝对整闭。只需把 $f$ 在 $A_j/IA_j$ 上的分解提升到 $A_j$。
这将我们归结到下一段讨论的情形。

\medskip\noindent
假设对 $i = 1, \ldots, r$ 有 $e_i = 1$，而对
$i = r + 1, \ldots, d$ 有 $e_i = 0$。由
$(g_0, h_0) = A/I[T]$，存在 $k_0, l_0 \in A/I[T]$，使得
$g_0 k_0 + h_0 l_0 = 1$。可见 $e = h_0 l_0$ 且
$e_i = h_0(a_i) l_0(a_i)$。于是对 $i = 1, \dots ,r$，
$h_0(a_i)$ 是单位元。由于 $f(a_i) = 0$，有
$0 = h_0(a_i)g_0(a_i)$，从而对 $i = 1, \ldots, r$ 有
$g_0(a_i) = 0$。因此 $(T - a_1)$ 在 $A/I[T]$ 中整除 $g_0$；设
$g_0 = (T - a_1) g_0'$。令
$f' = (T - a_2) \ldots (T - a_d)$ 且 $h'_0 = h_0$。
对 $d$ 归纳，可将分解 $f' \bmod I = g'_0 h'_0$ 提升为 $A$ 上的分解
$f' = g' h'$；它给出分解 $f = (T - a_1) g' h'$，提升给定分解
$f \bmod I = g_0 h_0$，正合所需。
\end{proof}








\section{自伴随环}
\label{more-algebra-section-auto-ass}

\noindent
部分材料见 \cite{Autour}。

\begin{definition}
\label{more-algebra-definition-auto-ass}
若 $R$ 是局部环，且其极大理想 $\mathfrak m$ 弱伴随于 $R$，则称环
$R$ 是{\it 自伴随的}。
\end{definition}

\begin{lemma}
\label{more-algebra-lemma-auto-ass-implies-P}
自伴随环 $R$ 具有下列性质：(P) 每个真有限生成理想
$I \subset R$ 都有非零零化子。
\end{lemma}

\begin{proof}
由假设，存在非零元素 $x \in R$，使得对每个
$f \in \mathfrak m$ 都有 $f^n x = 0$。设 $I = (f_1, \ldots, f_r)$。
则 $x$ 属于 $R \to \bigoplus R_{f_i}$ 的核。因此，由 Algebra, Lemma
\ref{algebra-lemma-when-injective-covering}，存在非零 $y \in R$，使得对
所有 $i$ 都有 $f_i y = 0$。由于 $y \in \text{Ann}_R(I)$，证毕。
\end{proof}

\begin{lemma}
\label{more-algebra-lemma-P-universally-injective}
设 $R$ 为具有引理 \ref{more-algebra-lemma-auto-ass-implies-P} 中性质 (P) 的环。
设 $u : N \to M$ 为投射 $R$-模的同态。则 $u$ 普遍单射当且仅当
$u$ 是单射。
\end{lemma}

\begin{proof}
假设 $u$ 是单射。我们的目标是证明 $u$ 普遍单射。先选取模 $Q$，
使得 $N \oplus Q$ 自由。考察映射
$N \oplus Q \to M \oplus Q$ 可知，只需在 $N$ 自由时证明该引理。
在此情形中，$N$ 是有限自由 $R$-模的有向余极限。因此归结到
$N$ 为有限自由 $R$-模的情形，设 $N = R^{\oplus n}$。对 $n$ 归纳
证明引理。$n = 0$ 的情形是平凡的。

\medskip\noindent
设 $u : R^{\oplus n} \to M$ 为单射模映射，其中 $M$ 投射。
选取 $R$-模 $Q$，使得 $M \oplus Q$ 自由。以复合
$R^{\oplus n} \to M \to M \oplus Q$ 替换 $u$ 后，可假设 $M$ 自由。
于是可找到直和因子 $R^{\oplus m} \subset M$，使得
$u(R^{\oplus n}) \subset R^{\oplus m}$。故可假设
$M = R^{\oplus m}$。此时 $u$ 由矩阵 $A = (a_{ij})$ 给出，满足
$u(x_1, \ldots, x_n) = (\sum x_i a_{i1}, \ldots, \sum x_i a_{im})$。
由于 $u$ 是单射，特别地，若 $x \not = 0$，则
$u(x, 0, \ldots, 0) = (xa_{11}, xa_{12}, \ldots, xa_{1m}) \not = 0$；
又因 $R$ 具有性质 (P)，可知
$a_{11}R + a_{12}R + \ldots + a_{1m}R = R$。由此可见，
$R(a_{11}, \ldots, a_{1m}) \subset R^{\oplus m}$ 是
$R^{\oplus m}$ 的直和因子；特别地，
$R^{\oplus m}/R(a_{11}, \ldots, a_{1m})$ 是投射 $R$-模。
得到交换图
$$
\xymatrix{
0 \ar[r] &
R \ar[rr] \ar[d]^1 & & R^{\oplus n} \ar[r] \ar[d]^u &
R^{\oplus n - 1} \ar[r] \ar[d] & 0 \\
0 \ar[r] & R \ar[rr]^{(a_{11}, \ldots, a_{1m})} & &
R^{\oplus m} \ar[r] & R^{\oplus m}/R(a_{11}, \ldots, a_{1m}) \ar[r] & 0
}
$$
其两行均分裂正合。因此右侧竖直箭头是单射；应用归纳假设，得到
右侧竖直箭头普遍单射。由此中间的竖直箭头普遍单射。
\end{proof}

\begin{lemma}
\label{more-algebra-lemma-P-fPD-zero}
设 $R$ 为环。下列条件等价：
\begin{enumerate}
\item $R$ 具有引理 \ref{more-algebra-lemma-auto-ass-implies-P} 中的性质 (P)；
\item 投射 $R$-模的任意单射都是普遍单射；
\item 若 $u : N \to M$ 是单射，且 $N$、$M$ 是有限投射 $R$-模，
则 $\Coker(u)$ 是有限投射 $R$-模；
\item 若 $N \subset M$，且 $N$、$M$ 作为 $R$-模有限投射，
则 $N$ 是 $M$ 的直和因子；
\item 任意单射 $R \to R^{\oplus n}$ 都是分裂单射。
\end{enumerate}
\end{lemma}

\begin{proof}
(1) $\Rightarrow$ (2) 是引理 \ref{more-algebra-lemma-P-universally-injective}。
显然，(3) 与 (4) 等价。由 Algebra, Lemma
\ref{algebra-lemma-universally-exact-split}，有
(2) $\Rightarrow$ (3)、(4)。第 (5) 项是 (4) 的特殊情形。
假设 (5)。设 $I = (a_1, \ldots, a_n)$ 为 $R$ 的真有限生成理想。
由于 $I \not = R$，映射 $R \to R^{\oplus n}$、
$x \mapsto (xa_1, \ldots, xa_n)$ 不是分裂单射。因此它有非零核，
从而 $\text{Ann}_R(I) \not = 0$。故 (1) 成立。
\end{proof}

\begin{example}
\label{more-algebra-example-auto-ass-weird-flat}
即使引理 \ref{more-algebra-lemma-P-fPD-zero} 中的等价条件成立，也不总是每个自由
$R$-模之间的单射都是分裂单射。例如，设
$R = k[x_1, x_2, x_3, \ldots]/(x_i^2)$。这是一个自伴随环。考虑自由
$R$-模之间的映射
$$
u :
\bigoplus\nolimits_{i \geq 1} Re_i
\longrightarrow
\bigoplus\nolimits_{i \geq 1} Rf_i, \quad
e_i \longmapsto f_i - x_if_{i + 1}.
$$
对任意整数 $n$，$u$ 在 $\bigoplus_{i = 1, \ldots, n} Re_i$ 上的限制
是单射，因为像 $u(e_1), \ldots, u(e_n)$ 在 $R$ 上线性无关。
因此 $u$ 是单射，从而由该引理它普遍单射。由于
$u \otimes \text{id}_k$ 是双射，可知若 $u$ 是分裂单射，则 $u$ 将为
满射。但 $u$ 不是满射，因为 $f_1$ 的逆像将是元素
$$
\sum\nolimits_{i \geq 0} x_1 \ldots x_ie_{i + 1} =
e_1 + x_1e_2 + x_1x_2e_3 + \ldots
$$
而它不属于该直和。顺便指出，$\Coker(u)$ 是可数生成的平坦
$R$-模（因为 $u$ 普遍单射），但它不是投射模（因为 $u$ 不分裂），
因而不是 Mittag-Leffler 模（见 Algebra, Lemma
\ref{algebra-lemma-countgen-projective}）。
\end{example}

\noindent
在局部环为 Noether 环的情形，下列引理是 Algebra, Proposition
\ref{algebra-proposition-what-exact} 的特殊情形。

\begin{lemma}
\label{more-algebra-lemma-exact-length-1}
设 $(R, \mathfrak m)$ 为局部环。假设
$\varphi : R^m \to R^n$ 是有限自由模之间的映射。下列条件等价：
\begin{enumerate}
\item $\varphi$ 是单射；
\item $\varphi$ 的秩为 $m$，且 $I(\varphi)$ 在 $R$ 中的零化子为零。
\end{enumerate}
若 $R$ 为 Noether 环，则它们还等价于：
\begin{enumerate}
\item[(3)] $\varphi$ 的秩为 $m$，且 $I(\varphi) = R$，或它包含一个非零因子。
\end{enumerate}
这里 $\varphi$ 的秩与 $I(\varphi)$ 如 Algebra, Definition
\ref{algebra-definition-rank} 中所定义。
\end{lemma}

\begin{proof}
若 $\varphi$ 的某个矩阵系数不属于 $\mathfrak m$，则应用
Algebra, Lemma \ref{algebra-lemma-add-trivial-complex}，把 $\varphi$
写成 $1 : R \to R$ 与映射
$\varphi' : R^{m-1} \to R^{n-1}$ 的和。容易看出，对 $\varphi'$
成立的引理蕴含对 $\varphi$ 成立的引理。因此从一开始即可假设
$\varphi$ 的所有矩阵系数都属于 $\mathfrak m$。

\medskip\noindent
假设 $\varphi$ 是单射。可假设 $m > 0$。设
$\mathfrak q \in \text{WeakAss}(R)$，从而 $R_\mathfrak q$ 是
自伴随环。于是 $\varphi$ 诱导单射
$R_\mathfrak q^m \to R_\mathfrak q^n$；由引理
\ref{more-algebra-lemma-auto-ass-implies-P} 和 \ref{more-algebra-lemma-P-universally-injective}，
该映射普遍单射。因此
$\varphi : \kappa(\mathfrak q)^m \to \kappa(\mathfrak q)^n$ 是单射。
故 $\varphi \bmod \mathfrak q$ 的秩为 $m$，且
$I(\varphi \otimes \kappa(\mathfrak q))$ 不是零理想。由于 $m$ 是
$\varphi$ 可能具有的最大秩，我们推出 $\varphi$ 的秩也为 $m$
（矩阵的秩在基变换后只能下降）。因此
$I(\varphi) \cdot \kappa(\mathfrak q) =
I(\varphi \otimes \kappa(\mathfrak q))$ 非零。
所以 $I(\varphi)$ 不包含于 $\mathfrak q$。于是 $R$ 的弱伴随素理想
没有一个是 $R$-模 $\text{Ann}_R I(\varphi)$ 的弱伴随素理想。
故 $\text{Ann}_R I(\varphi)$ 没有弱伴随素理想，见 Algebra, Lemma
\ref{algebra-lemma-weakly-ass}。由 Algebra, Lemma
\ref{algebra-lemma-weakly-ass-zero}，$\text{Ann}_R I(\varphi)$ 为零。

\medskip\noindent
反之，假设 (2)。秩为 $m$ 蕴含 $n \geq m$。写成
$I(\varphi) = (f_1, \ldots, f_r)$；这是可能的，因为
$I(\varphi)$ 有限生成。由 Algebra, Lemma
\ref{algebra-lemma-matrix-left-inverse}，可以找到映射
$\psi_i : R^n \to R^m$，使得
$\psi \circ \varphi = f_i \text{id}_{R^m}$。因此
$\varphi(x) = 0$ 蕴含对 $i = 1, \ldots, r$ 有 $f_i x = 0$。
这蕴含 $x = 0$，故 $\varphi$ 是单射。

\medskip\noindent
关于 Noether 局部情形中 (1) 与 (3) 的等价性，参见
Algebra, Proposition \ref{algebra-proposition-what-exact}。
若环 $R$ 是 Noether 环但非局部环，读者可由局部情形推出它；
略去细节。另一种方法是使用伴随素理想重做上述论证，并使用下列事实：
伴随素理想只有有限多个；素理想回避成立；Noether 环的非零因子恰为
不包含于任何伴随素理想的元素。
\end{proof}

\begin{lemma}
\label{more-algebra-lemma-coker-injective-free}
设 $R$ 为环。假设 $\varphi : R^n \to R^n$ 是相同秩的有限自由模之间的
单射。则 $\Hom_R(\Coker(\varphi), R) = 0$。
\end{lemma}

\begin{proof}
令 $\varphi^t : R^n \to R^n$ 为 $\varphi$ 的转置。引理断言
$\varphi^t$ 是单射。使用引理 \ref{more-algebra-lemma-exact-length-1} 的记号可知，
$\varphi^t$ 的秩为 $n$，且 $I(\varphi) = I(\varphi^t)$。
因此结论由该引理中 (1) 与 (2) 的等价性得出。
\end{proof}








\section{平坦化分层}
\label{more-algebra-section-flattening}

\noindent
设 $R \to S$ 为环映射，$M$ 为 $S$-模。对任意 $R$-代数 $R'$，
可以考虑基变换 $S' = S \otimes_R R'$ 与
$M' = M \otimes_R R'$。若模 $M'$ 在 $R'$ 上平坦，则称
$R \to R'$ {\it 使 $M$ 平坦}。我们希望理解使 $M$ 平坦的环映射
$R \to R'$ 所成集合的结构。特别地，我们希望知道是否存在
{\it $M$ 的普遍平坦化 $R \to R_{univ}$}；也就是说，是否存在
使 $M$ 平坦的环映射 $R \to R_{univ}$，并且任何使 $M$ 平坦的环映射
$R \to R'$ 都通过 $R \to R_{univ}$ 分解。事实表明，这样的普遍解
通常不存在。

\medskip\noindent
我们将在 More on Flatness, Section \ref{flat-section-flattening} 中，
在概形论设定 $\mathcal{F}/X/S$ 下讨论{\it 普遍平坦化}与
{\it 平坦化分层}。若普遍平坦化 $R \to R_{univ}$ 存在，则概形态射
$\Spec(R_{univ}) \to \Spec(R)$ 是 $\Spec(S)$ 上拟凝聚模
$\widetilde{M}$ 的普遍平坦化。

\medskip\noindent
在本节及接下来的几节中，我们证明一些与此有关的基本代数事实。
最基本的结果也许是下面这个。

\begin{lemma}
\label{more-algebra-lemma-intersection-flat}
设 $R$ 为环，$M$ 为 $R$-模，$I_1$、$I_2$ 为 $R$ 的理想。若
$M/I_1M$ 在 $R/I_1$ 上平坦，且 $M/I_2M$ 在 $R/I_2$ 上平坦，
则 $M/(I_1 \cap I_2)M$ 在 $R/(I_1 \cap I_2)$ 上平坦。
\end{lemma}

\begin{proof}
以 $R/(I_1 \cap I_2)$ 替换 $R$，并以
$M/(I_1 \cap I_2)M$ 替换 $M$，可假设 $I_1 \cap I_2 = 0$。
设 $J \subset R$ 为理想。为证明 $M$ 在 $R$ 上平坦，必须证明
$J \otimes_R M \to M$ 是单射，见 Algebra, Lemma
\ref{algebra-lemma-flat}。由 $M/I_1M$ 在 $R/I_1$ 上的平坦性，映射
$$
J/(J \cap I_1) \otimes_R M =
(J + I_1)/I_1 \otimes_{R/I_1} M/I_1M
\longrightarrow M/I_1M
$$
是单射。由于 $0 \to (J \cap I_1) \to J \to J/(J \cap I_1) \to 0$
正合，得到交换图
$$
\xymatrix{
(J \cap I_1) \otimes_R M \ar[r] \ar[d] &
J \otimes_R M \ar[r] \ar[d] &
J/(J \cap I_1) \otimes_R M \ar[r] \ar[d] & 0 \\
M \ar@{=}[r] &
M \ar[r] &
M/I_1M
}
$$
因此，只需证明 $(J \cap I_1) \otimes_R M \to M$ 是单射。
由于 $I_1 \cap I_2 = 0$，理想 $J \cap I_1$ 同构地映到
一个理想 $J' \subset R/I_2$，并且
$(J \cap I_1) \otimes_R M = J' \otimes_{R/I_2} M/I_2M$。
由 $M/I_2M$ 在 $R/I_2$ 上的平坦性，映射
$J' \otimes_{R/I_2} M/I_2M \to M/I_2M$ 是单射；这显然蕴含
$(J \cap I_1) \otimes_R M \to M$ 是单射。
\end{proof}





\section{Artin 环上的平坦化}
\label{more-algebra-section-flattening-artinian}

\noindent
当基环是 Artin 局部环时，普遍平坦化存在。对任意模，它都存在。
因此，正如稍后将看到的，当基概形是 Artin 局部环的谱时，平坦化分层
存在。

\begin{lemma}
\label{more-algebra-lemma-flattening-artinian}
设 $R$ 为 Artin 环，$M$ 为 $R$-模。则存在最小理想
$I \subset R$，使得 $M/IM$ 在 $R/I$ 上平坦。
\end{lemma}

\begin{proof}
这直接由引理 \ref{more-algebra-lemma-intersection-flat} 与 Artin 性得出。
\end{proof}

\noindent
该理想具有下列泛性质。

\begin{lemma}
\label{more-algebra-lemma-flattening-artinian-universal-property}
设 $R$ 为 Artin 环，$M$ 为 $R$-模。设 $I \subset R$ 是使得
$M/IM$ 在 $R/I$ 上平坦的最小理想 $I \subset R$。则 $I$ 具有下列
泛性质：对每个环映射 $\varphi : R \to R'$，有
$$
R' \otimes_R M\text{ is flat over }R'
\Leftrightarrow
\text{有 }\varphi(I) = 0.
$$
\end{lemma}

\begin{proof}
注意，由引理 \ref{more-algebra-lemma-flattening-artinian}，$I$ 存在。
蕴含关系 $\Rightarrow$ 由 Algebra, Lemma
\ref{algebra-lemma-flat-base-change} 得出。设 $\varphi : R \to R'$ 满足
$M \otimes_R R'$ 在 $R'$ 上平坦。令 $J = \Ker(\varphi)$。
由 Algebra, Lemma
\ref{algebra-lemma-descent-flatness-injective-map-artinian-rings}，并利用
$R' \otimes_R M = R' \otimes_{R/J} M/JM$ 在 $R'$ 上平坦，推出
$M/JM$ 在 $R/J$ 上平坦。故 $I \subset J$，正合所需。
\end{proof}











\section{基的闭子集上的平坦化}
\label{more-algebra-section-flattening-local-base}

\noindent
设 $R \to S$ 为环映射，$I \subset R$ 为理想，$M$ 为 $S$-模。
下文考虑条件
\begin{equation}
\label{more-algebra-equation-flat-at-primes-over}
\forall \mathfrak q \in V(IS) \subset \Spec(S) :
M_{\mathfrak q}\text{ is flat over }R.
\end{equation}
在几何上，这意味着 $M$ 沿 $V(I)$ 在 $\Spec(S)$ 中的逆像在 $R$
上平坦。若 $R$ 与 $S$ 是 Noether 环，且 $M$ 是有限 $S$-模，则
(\ref{more-algebra-equation-flat-at-primes-over}) 等价于：对所有 $n \geq 1$，
$M/I^nM$ 在 $R/I^n$ 上平坦；见 Algebra, Lemma
\ref{algebra-lemma-flat-module-powers}。

\begin{lemma}
\label{more-algebra-lemma-base-change-flat-at-primes-over}
设 $R \to S$ 为环映射，$I \subset R$ 为理想，$M$ 为 $S$-模。
设 $R \to R'$ 为环映射，$IR' \subset I' \subset R'$ 为理想。
若 (\ref{more-algebra-equation-flat-at-primes-over}) 对
$(R \to S, I, M)$ 成立，则 (\ref{more-algebra-equation-flat-at-primes-over}) 对
$(R' \to S \otimes_R R', I', M \otimes_R R')$ 成立。
\end{lemma}

\begin{proof}
假设 (\ref{more-algebra-equation-flat-at-primes-over}) 对
$(R \to S, I \subset R, M)$ 成立。设
$I'(S \otimes_R R') \subset \mathfrak q'$ 是 $S \otimes_R R'$ 的
素理想。令 $\mathfrak q \subset S$ 为 $S$ 中对应的素理想。
则 $IS \subset \mathfrak q$。注意，
$(M \otimes_R R')_{\mathfrak q'}$ 是基变换
$M_{\mathfrak q} \otimes_R R'$ 的局部化。因此，作为平坦模的局部化，
$(M \otimes_R R')_{\mathfrak q'}$ 在 $R'$ 上平坦；见
Algebra, Lemmas \ref{algebra-lemma-flat-base-change} and
\ref{algebra-lemma-flat-localization}。
\end{proof}

\begin{lemma}
\label{more-algebra-lemma-flat-descent-flat-at-primes-over}
设 $R \to S$ 为环映射，$I \subset R$ 为理想，$M$ 为 $S$-模。
设 $R \to R'$ 为环映射，$IR' \subset I' \subset R'$ 为理想，并且：
\begin{enumerate}
\item 由 $\Spec(R') \to \Spec(R)$ 诱导的映射
$V(I') \to V(I)$ 是满射；
\item 对所有素理想 $\mathfrak p' \in V(I')$，
$R'_{\mathfrak p'}$ 在 $R$ 上平坦。
\end{enumerate}
若 (\ref{more-algebra-equation-flat-at-primes-over}) 对
$(R' \to S \otimes_R R', I', M \otimes_R R')$ 成立，则
(\ref{more-algebra-equation-flat-at-primes-over}) 对 $(R \to S, I, M)$ 成立。
\end{lemma}

\begin{proof}
假设 (\ref{more-algebra-equation-flat-at-primes-over}) 对
$(R' \to S \otimes_R R', IR', M \otimes_R R')$ 成立。
选取素理想 $IS \subset \mathfrak q \subset S$。令
$I \subset \mathfrak p \subset R$ 为 $R$ 中对应的素理想。
由假设，存在位于 $\mathfrak p$ 之上的 $R'$ 的素理想
$\mathfrak p' \in V(I')$，并且
$R_{\mathfrak p} \to R'_{\mathfrak p'}$ 平坦。选取素理想
$\overline{\mathfrak q}' \subset
\kappa(\mathfrak q) \otimes_{\kappa(\mathfrak p)} \kappa(\mathfrak p')$；
它对应于素理想 $\mathfrak q' \subset S \otimes_R R'$，后者位于
$\mathfrak q$ 与 $\mathfrak p'$ 之上。注意，
$(S \otimes_R R')_{\mathfrak q'}$ 是
$S_{\mathfrak q} \otimes_{R_{\mathfrak p}} R'_{\mathfrak p'}$ 的局部化。
由假设，模 $(M \otimes_R R')_{\mathfrak q'}$ 在
$R'_{\mathfrak p'}$ 上平坦。因此，Algebra, Lemma
\ref{algebra-lemma-base-change-flat-up-down} 蕴含
$M_{\mathfrak q}$ 在 $R_{\mathfrak p}$ 上平坦，正是所需结论。
\end{proof}

\begin{lemma}
\label{more-algebra-lemma-limit-preserving-flat-at-primes-over}
设 $R \to S$ 为有限表示环映射，$M$ 为有限表示 $S$-模。
设 $R' = \colim_{\lambda \in \Lambda} R_\lambda$ 是 $R$-代数的有向余极限。
设 $I_\lambda \subset R_\lambda$ 为理想，使得对所有
$\mu \geq \lambda$ 有 $I_\lambda R_\mu \subset I_\mu$，并令
$I' = \colim_\lambda I_\lambda$。若
(\ref{more-algebra-equation-flat-at-primes-over}) 对
$(R' \to S \otimes_R R', I', M \otimes_R R')$ 成立，则存在
$\lambda \in \Lambda$，使得 (\ref{more-algebra-equation-flat-at-primes-over}) 对
$(R_\lambda \to S \otimes_R R_\lambda, I_\lambda, M \otimes_R R_\lambda)$
成立。
\end{lemma}

\begin{proof}
记 $S_\lambda = S \otimes_R R_\lambda$、$S' = S \otimes_R R'$、
$M_\lambda = M \otimes_R R_\lambda$ 和 $M' = M \otimes_R R'$。
基变换 $S'$ 在 $R'$ 上有限表示，$M'$ 在 $S'$ 上有限表示；带下标
$\lambda$ 的版本也同样如此，见 Algebra, Lemma
\ref{algebra-lemma-base-change-finiteness}。由 Algebra, Theorem
\ref{algebra-theorem-openness-flatness}，集合
$$
U' = \{\mathfrak q' \in \Spec(S') \mid
M'_{\mathfrak q'}\text{ is flat over }R'\}
$$
在 $\Spec(S')$ 中为开集。注意，$V(I'S')$ 是拟紧空间，且由假设它
包含于 $U'$。因此存在有限多个 $g'_j \in S'$、
$j = 1, \ldots, m$，使得 $D(g'_j) \subset U'$，并且
$V(I'S') \subset \bigcup D(g'_j)$。特别地，
$(M')_{g'_j}$ 在 $R'$ 上平坦。

\medskip\noindent
我们将依次选取越来越大的 $\lambda \in \Lambda$。首先选得足够大，
使得可找到映到 $g'_j$ 的 $g_{j, \lambda} \in S_{\lambda}$。
包含关系 $V(I'S') \subset \bigcup D(g'_j)$ 意味着
$I'S' + (g'_1, \ldots, g'_m) = S'$；它可表示为
$1 = \sum z_sh_s + \sum f_jg'_j$，其中
$z_s \in I'$、$h_s, f_j \in S'$。增大 $\lambda$ 后，可假设这样的等式
在 $S_\lambda$ 中成立。因此可假设
$V(I_\lambda S_\lambda) \subset \bigcup D(g_{j, \lambda})$。
由 Algebra, Lemma \ref{algebra-lemma-flat-finite-presentation-limit-flat}，
对某个足够大的 $\lambda$，模 $(M_\lambda)_{g_{j, \lambda}}$ 在
$R_\lambda$ 上平坦。特别地，$M_\lambda$ 在位于理想 $I_\lambda$
之上的所有素理想处都在 $R_\lambda$ 上平坦。
\end{proof}









\section{源与基的闭子集上的平坦化}
\label{more-algebra-section-flattening-local-source-base}

\noindent
本节略微推广第 \ref{more-algebra-section-flattening-local-base} 节的讨论。
强烈建议读者先阅读并理解该节。

\begin{situation}
\label{more-algebra-situation-flattening-general}
设 $R \to S$ 为环映射，$J \subset S$ 为理想，$M$ 为 $S$-模。
\end{situation}

\noindent
在此情形中，给定 $R$-代数 $R'$ 与理想 $I' \subset R'$，令
$S' = S \otimes_R R'$、$M' = M \otimes_R R'$。我们考虑条件
\begin{equation}
\label{more-algebra-equation-flat-at-primes}
\forall \mathfrak q' \in V(I'S' + JS') \subset \Spec(S') :
M'_{\mathfrak q'}\text{ is flat over }R'.
\end{equation}
在几何上，这意味着 $M'$ 在 $V(I')$ 的逆像与 $V(J)$ 的逆像之交上
相对于 $R'$ 平坦。由于 $(R \to S, J, M)$ 已固定，条件
(\ref{more-algebra-equation-flat-at-primes}) 只依赖于对 $(R', I')$，其中 $R'$ 被视为
$R$-代数。

\begin{lemma}
\label{more-algebra-lemma-base-change-flat-at-primes}
在情形 \ref{more-algebra-situation-flattening-general} 中，设
$R' \to R''$ 为 $R$-代数映射。设 $I' \subset R'$ 以及
$I'R'' \subset I'' \subset R''$ 为理想。若
(\ref{more-algebra-equation-flat-at-primes}) 对 $(R', I')$ 成立，则
(\ref{more-algebra-equation-flat-at-primes}) 对
$(R'', I'')$ 成立。
\end{lemma}

\begin{proof}
假设 (\ref{more-algebra-equation-flat-at-primes}) 对 $(R', I')$ 成立。设
$I''S'' + JS'' \subset \mathfrak q''$ 是 $S''$ 的素理想。
令 $\mathfrak q' \subset S'$ 为 $S'$ 中对应的素理想。于是
$I'S' \subset \mathfrak q'$ 与 $JS' \subset \mathfrak q'$ 都成立，
因为相应条件对 $\mathfrak q''$ 成立。注意，
$(M'')_{\mathfrak q''}$ 是基变换
$M'_{\mathfrak q'} \otimes_R R''$ 的局部化。因此，作为平坦模的局部化，
$(M'')_{\mathfrak q''}$ 在 $R''$ 上平坦；见 Algebra, Lemmas
\ref{algebra-lemma-flat-base-change} and
\ref{algebra-lemma-flat-localization}。
\end{proof}

\begin{lemma}
\label{more-algebra-lemma-flat-descent-flat-at-primes}
在情形 \ref{more-algebra-situation-flattening-general} 中，设
$R' \to R''$ 为 $R$-代数映射。设 $I' \subset R'$ 以及
$I'R'' \subset I'' \subset R''$ 为理想。假设：
\begin{enumerate}
\item 由 $\Spec(R'') \to \Spec(R')$ 诱导的映射
$V(I'') \to V(I')$ 是满射；
\item 对所有素理想 $\mathfrak p'' \in V(I'')$，
$R''_{\mathfrak p''}$ 在 $R'$ 上平坦。
\end{enumerate}
若 (\ref{more-algebra-equation-flat-at-primes}) 对 $(R'', I'')$ 成立，则
(\ref{more-algebra-equation-flat-at-primes}) 对
$(R', I')$ 成立。
\end{lemma}

\begin{proof}
假设 (\ref{more-algebra-equation-flat-at-primes}) 对 $(R'', I'')$ 成立。
选取素理想 $I'S' + JS' \subset \mathfrak q' \subset S'$。
令 $I' \subset \mathfrak p' \subset R'$ 为 $R'$ 中对应的素理想。
由假设，存在位于 $\mathfrak p'$ 之上的 $R''$ 的素理想
$\mathfrak p'' \in V(I'')$，并且
$R'_{\mathfrak p'} \to R''_{\mathfrak p''}$ 平坦。选取素理想
$\overline{\mathfrak q}'' \subset
\kappa(\mathfrak q') \otimes_{\kappa(\mathfrak p')} \kappa(\mathfrak p'')$。
这对应于素理想 $\mathfrak q'' \subset S'' = S' \otimes_{R'} R''$，
后者位于 $\mathfrak q'$ 与 $\mathfrak p''$ 之上。特别地，
$I''S'' \subset \mathfrak q''$ 且 $JS'' \subset \mathfrak q''$。
注意，$(S' \otimes_{R'} R'')_{\mathfrak q''}$ 是
$S'_{\mathfrak q'} \otimes_{R'_{\mathfrak p'}} R''_{\mathfrak p''}$ 的
局部化。由假设，模 $(M' \otimes_{R'} R'')_{\mathfrak q''}$ 在
$R''_{\mathfrak p''}$ 上平坦。因此，Algebra, Lemma
\ref{algebra-lemma-base-change-flat-up-down} 蕴含
$M'_{\mathfrak q'}$ 在 $R'_{\mathfrak p'}$ 上平坦，正是所需结论。
\end{proof}

\begin{lemma}
\label{more-algebra-lemma-limit-preserving-flat-at-primes}
在情形 \ref{more-algebra-situation-flattening-general} 中，假设 $R \to S$ 本质有限表示，
且 $M$ 是有限表示 $S$-模。设
$R' = \colim_{\lambda \in \Lambda} R_\lambda$ 是 $R$-代数的有向余极限。
设 $I_\lambda \subset R_\lambda$ 为理想，使得对所有
$\mu \geq \lambda$ 有 $I_\lambda R_\mu \subset I_\mu$，并令
$I' = \colim_\lambda I_\lambda$。若 (\ref{more-algebra-equation-flat-at-primes}) 对
$(R', I')$ 成立，则存在 $\lambda \in \Lambda$，使得
(\ref{more-algebra-equation-flat-at-primes}) 对 $(R_\lambda, I_\lambda)$ 成立。
\end{lemma}

\begin{proof}
先证明 $R \to S$ 有限表示的情形，再说明一般情形需要作何修改。
记 $S_\lambda = S \otimes_R R_\lambda$、$S' = S \otimes_R R'$、
$M_\lambda = M \otimes_R R_\lambda$ 及 $M' = M \otimes_R R'$。
基变换 $S'$ 在 $R'$ 上有限表示，$M'$ 在 $S'$ 上有限表示；带下标
$\lambda$ 的版本也同样如此，见 Algebra, Lemma
\ref{algebra-lemma-base-change-finiteness}。由 Algebra, Theorem
\ref{algebra-theorem-openness-flatness}，集合
$$
U' = \{\mathfrak q' \in \Spec(S') \mid
M'_{\mathfrak q'}\text{ is flat over }R'\}
$$
在 $\Spec(S')$ 中为开集。注意，$V(I'S' + JS')$ 是拟紧空间，且由
假设它包含于 $U'$。因此存在有限多个
$g'_j \in S'$、$j = 1, \ldots, m$，使得 $D(g'_j) \subset U'$，且
$V(I'S' + JS') \subset \bigcup D(g'_j)$。特别地，
$(M')_{g'_j}$ 在 $R'$ 上平坦。

\medskip\noindent
我们将依次选取越来越大的 $\lambda \in \Lambda$。先选得足够大，
使得可以找到映到 $g'_j$ 的 $g_{j, \lambda} \in S_{\lambda}$。
包含关系 $V(I'S' + JS') \subset \bigcup D(g'_j)$ 意味着
$I'S' + JS' + (g'_1, \ldots, g'_m) = S'$；它可表示为
$$
1 = \sum y_tk_t + \sum z_sh_s + \sum f_jg'_j
$$
其中 $z_s \in I'$、$y_t \in J$、$k_t, h_s, f_j \in S'$。
增大 $\lambda$ 后，可假设这样的等式在 $S_\lambda$ 中成立。
因此可假设
$V(I_\lambda S_\lambda + J S_\lambda) \subset \bigcup D(g_{j, \lambda})$。
由 Algebra, Lemma \ref{algebra-lemma-flat-finite-presentation-limit-flat}，
对某个足够大的 $\lambda$，模 $(M_\lambda)_{g_{j, \lambda}}$ 在
$R_\lambda$ 上平坦。特别地，$M_\lambda$ 在对应于
$V(I_\lambda S_\lambda + J S_\lambda)$ 中各点的所有素理想处，都在
$R_\lambda$ 上平坦。

\medskip\noindent
若 $S$ 本质有限表示，则可写成 $S = \Sigma^{-1}C$，其中
$R \to C$ 有限表示，$\Sigma \subset C$ 是乘法子集。还可写成
$M = \Sigma^{-1}N$，其中 $N$ 是有限表示 $C$-模，见 Algebra, Lemma
\ref{algebra-lemma-construct-fp-module}。此时引入
$C_\lambda$、$C'$、$N_\lambda$、$N'$。在上述讨论中，我们得到开集
$U' \subset \Spec(C')$，在其上 $N'$ 相对于 $R'$ 平坦。条件
(\ref{more-algebra-equation-flat-at-primes}) 成立意味着 $V(I'S' + JS')$ 映入
$U'$，因为对素理想 $\mathfrak q' \subset S'$ 及其对应的素理想
$\mathfrak r' \subset C'$，有 $M'_{\mathfrak q'} = N'_{\mathfrak r'}$。
因此可找到 $g'_j \in C'$，使得 $\bigcup D(g'_j)$ 包含
$V(I'S' + JS')$ 的像。证明的其余部分与前面完全相同。
\end{proof}

\begin{lemma}
\label{more-algebra-lemma-flat-module-powers-variant}
在情形 \ref{more-algebra-situation-flattening-general} 中，设 $I \subset R$ 为理想。
假设：
\begin{enumerate}
\item $R$ 是 Noether 环；
\item $S$ 是 Noether 环；
\item $M$ 是有限 $S$-模；
\item 对每个 $n \geq 1$ 以及任意素理想
$\mathfrak q \in V(J + IS)$，模 $(M/I^n M)_{\mathfrak q}$ 在
$R/I^n$ 上平坦。
\end{enumerate}
则 (\ref{more-algebra-equation-flat-at-primes}) 对 $(R, I)$ 成立；即对每个素理想
$\mathfrak q \in V(J + IS)$，局部化 $M_{\mathfrak q}$ 在 $R$ 上平坦。
\end{lemma}

\begin{proof}
设 $\mathfrak q \in V(J + IS)$。将 Algebra, Lemma
\ref{algebra-lemma-flat-module-powers} 应用于 $R \to S_{\mathfrak q}$ 与
$M_{\mathfrak q}$，即可推出 $M_{\mathfrak q}$ 在 $R$ 上平坦。
\end{proof}










\section{Noether 完备局部环上的平坦化}
\label{more-algebra-section-flattening-Noetherian-complete-local}

\noindent
当基为完备 Noether 局部环 $R$，且给定模是在有限型 $R$-代数 $S$ 上
有限的模时，下面三个引理给出``局部''平坦化分层存在性的纯代数证明。

\begin{lemma}
\label{more-algebra-lemma-flattening-complete-local-noetherian}
设 $R \to S$ 为环映射，$M$ 为 $S$-模。假设：
\begin{enumerate}
\item $(R, \mathfrak m)$ 是完备 Noether 局部环；
\item $S$ 是 Noether 环；
\item $M$ 在 $S$ 上有限。
\end{enumerate}
则存在理想 $I \subset \mathfrak m$，使得：
\begin{enumerate}
\item 对 $S/IS$ 中位于 $\mathfrak m$ 之上的所有素理想
$\mathfrak q$，$(M/IM)_{\mathfrak q}$ 在 $R/I$ 上平坦；
\item 若 $J \subset R$ 是理想，并且对位于 $\mathfrak m$ 之上的所有
素理想 $\mathfrak q$，$(M/JM)_{\mathfrak q}$ 在 $R/J$ 上平坦，
则 $I \subset J$。
\end{enumerate}
换言之，$I$ 是使 (\ref{more-algebra-equation-flat-at-primes-over}) 对
$(\overline{R} \to \overline{S}, \overline{\mathfrak m}, \overline{M})$
成立的 $R$ 的最小理想，其中 $\overline{R} = R/I$、
$\overline{S} = S/IS$、$\overline{\mathfrak m} = \mathfrak m/I$ 且
$\overline{M} = M/IM$。
\end{lemma}

\begin{proof}
设 $J \subset R$ 为理想。将 Algebra, Lemma
\ref{algebra-lemma-flat-module-powers} 应用于环 $R/J$ 上的模 $M/JM$。
于是可见，对 $S/JS$ 中所有素理想 $\mathfrak q$，
$(M/JM)_{\mathfrak q}$ 在 $R/J$ 上平坦，当且仅当对所有
$n \geq 1$，$M/(J + \mathfrak m^n)M$ 在
$R/(J + \mathfrak m^n)$ 上平坦。下文使用这一观察。

\medskip\noindent
对每个 $n \geq 1$，局部环 $R/\mathfrak m^n$ 是 Artin 环。
故由引理 \ref{more-algebra-lemma-flattening-artinian}，存在最小理想
$I_n \supset \mathfrak m^n$，使得 $M/I_nM$ 在 $R/I_n$ 上平坦。
显然，$I_{n + 1} + \mathfrak m^n$ 包含 $I_n$；应用引理
\ref{more-algebra-lemma-intersection-flat}，可知
$I_n = I_{n + 1} + \mathfrak m^n$。由于
$R = \lim_n\ R/\mathfrak m^n$，可知
$I = \lim_n\ I_n/\mathfrak m^n$ 是 $R$ 中的理想，并且对所有
$n \geq 1$ 有 $I_n = I + \mathfrak m^n$。由证明开头的观察，
$I$ 满足 (1) 与 (2)。略去一些细节。
\end{proof}

\begin{lemma}
\label{more-algebra-lemma-flattening-complete-local-noetherian-property-by-finite-type}
采用引理 \ref{more-algebra-lemma-flattening-complete-local-noetherian} 中的记号
$R \to S$、$M$、$I$ 及假设。考虑局部环的局部同态
$\varphi : (R, \mathfrak m) \to (R', \mathfrak m')$，并假设
$R'$ 是 Noether 环。则下列条件等价：
\begin{enumerate}
\item (\ref{more-algebra-equation-flat-at-primes-over}) 对
$(R' \to S \otimes_R R', \mathfrak m', M \otimes_R R')$ 成立；
\item $\varphi(I) = 0$。
\end{enumerate}
\end{lemma}

\begin{proof}
(2) $\Rightarrow$ (1) 由引理
\ref{more-algebra-lemma-base-change-flat-at-primes-over} 得出。设
$\varphi : R \to R'$ 如引理所述并满足 (1)。必须证明
$\varphi(I) = 0$。这等价于 $\varphi(I)R' = 0$。由 Noether 局部环
$R'$ 中的 Artin--Rees（见 Algebra, Lemma
\ref{algebra-lemma-intersect-powers-ideal-module-zero}），这等价于：对所有
$n > 0$，有
$\varphi(I)R' + (\mathfrak m')^n = (\mathfrak m')^n$。
继而，这等价于：对每个 $n$，复合
$\varphi_n : R \to R' \to R'/(\mathfrak m')^n$ 零化 $I$。
由引理 \ref{more-algebra-lemma-base-change-flat-at-primes-over}，$\varphi$ 的假设 (1)
蕴含 $\varphi_n$ 的假设 (1)。这将我们归结到 $R'$ 为 Artin 局部环的
情形。

\medskip\noindent
假设 $R'$ 为 Artin 环。令 $J = \Ker(\varphi)$。必须证明
$I \subset J$。由引理 \ref{more-algebra-lemma-flattening-complete-local-noetherian}
中 $I$ 的构造，只需证明：对 $S/JS$ 中位于 $\mathfrak m$ 之上的每个
素理想 $\mathfrak q$，$(M/JM)_{\mathfrak q}$ 在 $R/J$ 上平坦。
由于 $R'$ 是 Artin 环，条件 (1) 表示 $M \otimes_R R'$ 在 $R'$ 上
平坦。由于 $R'$ 是 Artin 环，且 $R/J \to R'$ 是单射局部环映射，
$R/J$ 也是 Artin 环。因此，由 Algebra, Lemma
\ref{algebra-lemma-descent-flatness-injective-map-artinian-rings}，
$M \otimes_R R' = M/JM \otimes_{R/J} R'$ 在 $R'$ 上的平坦性蕴含
$M/JM$ 在 $R/J$ 上平坦。这就完成了证明。
\end{proof}

\begin{lemma}
\label{more-algebra-lemma-flattening-complete-local-universal-property}
采用引理 \ref{more-algebra-lemma-flattening-complete-local-noetherian} 中的记号
$R \to S$、$M$、$I$ 及假设，并另外假设 $R \to S$ 是有限型的。
则对任意局部环的局部同态
$\varphi : (R, \mathfrak m) \to (R', \mathfrak m')$，下列条件等价：
\begin{enumerate}
\item (\ref{more-algebra-equation-flat-at-primes-over}) 对
$(R' \to S \otimes_R R', \mathfrak m', M \otimes_R R')$ 成立；
\item $\varphi(I) = 0$。
\end{enumerate}
\end{lemma}

\begin{proof}
(2) $\Rightarrow$ (1) 由引理
\ref{more-algebra-lemma-base-change-flat-at-primes-over} 得出。设
$\varphi : R \to R'$ 如引理所述并满足 (1)。由于 $R$ 是 Noether 环，
$R \to S$ 有限表示，且 $M$ 是有限表示 $S$-模。把
$R' = \colim_\lambda R_\lambda$ 写成局部 $R$-子代数
$R_\lambda \subset R'$ 的有向余极限，其中极大理想
$\mathfrak m_\lambda = R_\lambda \cap \mathfrak m'$，且每个
$R_\lambda$ 在 $R$ 上本质有限型。由引理
\ref{more-algebra-lemma-limit-preserving-flat-at-primes-over}，对某个 $\lambda$，条件
(\ref{more-algebra-equation-flat-at-primes-over}) 对
$(R_\lambda \to S \otimes_R R_\lambda, \mathfrak m_\lambda,
M \otimes_R R_\lambda)$ 成立。因此，引理
\ref{more-algebra-lemma-flattening-complete-local-noetherian-property-by-finite-type}
可应用于环映射 $R \to R_\lambda$；可知 $I$ 在 $R_\lambda$ 中映为零，
从而更在 $R'$ 中映为零。
\end{proof}





\section{沿整映射的平坦性下降}
\label{more-algebra-section-descent-flatness-integral}

\noindent
先给出几个简单引理。

\begin{lemma}
\label{more-algebra-lemma-have-one-root}
设 $R$ 为环，$P(T)$ 为系数属于 $R$ 的首一多项式。设
$\alpha \in R$ 满足 $P(\alpha) = 0$。则对某个首一多项式
$Q(T) \in R[T]$，有 $P(T) = (T - \alpha)Q(T)$。
\end{lemma}

\begin{proof}
对 $P$ 的次数归纳。若 $\deg(P) = 1$，则
$P(T) = T - \alpha$，结论成立。若 $\deg(P) > 1$，则由多项式长除法，
可对某个次数 $< \deg(P)$ 的多项式 $Q \in R[T]$ 与某个 $r \in R$
写成 $P(T) = (T - \alpha)Q(T) + r$。由假设，
$0 = P(\alpha) = (\alpha - \alpha)Q(\alpha) + r = r$，故
$r = 0$，正合所需。
\end{proof}

\begin{lemma}
\label{more-algebra-lemma-adjoin-one-root}
设 $R$ 为环，$P(T)$ 为系数属于 $R$ 的首一多项式。存在有限自由环映射
$R \to R'$，使得对某个 $\alpha \in R'$ 及某个首一多项式
$Q(T) \in R'[T]$，有 $P(T) = (T - \alpha)Q(T)$。
\end{lemma}

\begin{proof}
写成 $P(T) = T^d + a_1T^{d - 1} + \ldots + a_0$。令
$R' = R[x]/(x^d + a_1x^{d - 1} + \ldots + a_0)$。令 $\alpha$ 等于
$x$ 的同余类。显然 $P(\alpha) = 0$。因此结论由引理
\ref{more-algebra-lemma-have-one-root} 得出。
\end{proof}

\begin{lemma}
\label{more-algebra-lemma-finite-split}
设 $R \to S$ 为有限环映射。存在有限自由环扩张 $R \subset R'$，
使得 $S \otimes_R R'$ 是下列形状的环的商环：
$$
R'[T_1, \ldots, T_n]/(P_1(T_1), \ldots, P_n(T_n))
$$
其中对某些 $\alpha_{ij} \in R'$，有
$P_i(T) = \prod_{j = 1, \ldots, d_i} (T - \alpha_{ij})$。
\end{lemma}

\begin{proof}
设 $x_1, \ldots, x_n \in S$ 是 $S$ 在 $R$ 上的一组生成元。
对每个 $i$，可选取首一多项式 $P_i(T) \in R[T]$，使得
$P_i(x_i) = 0$ 在 $S$ 中成立，见 Algebra, Lemma
\ref{algebra-lemma-finite-is-integral}。设 $\deg(P_i) = d_i$。
将引理 \ref{more-algebra-lemma-adjoin-one-root} 应用 $\sum d_i$ 次，得到有限自由环
扩张 $R \subset R'$，使得每个 $P_i$ 完全分裂：
$$
P_i(T) = \prod\nolimits_{j = 1, \ldots, d_i} (T - \alpha_{ij})
$$
其中某些 $\alpha_{ik} \in R'$。令
$R'[T_1, \ldots, T_n] \to S \otimes_R R'$ 为把 $T_i$ 映到
$x_i \otimes 1$ 的 $R'$-代数映射。由于它把 $P_i(T_i)$ 映为零，
便诱导所需满射。
\end{proof}

\begin{lemma}
\label{more-algebra-lemma-split-image}
设 $R$ 为环。令 $S = R[T_1, \ldots, T_n]/J$。假设 $J$ 包含形如
$P_i(T_i)$ 的元素，其中对某些 $\alpha_{ij} \in R$，有
$P_i(T) = \prod_{j = 1, \ldots, d_i} (T - \alpha_{ij})$。
对满足 $1 \leq k_i \leq d_i$ 的
$\underline{k} = (k_1, \ldots, k_n)$，考虑环映射
$$
\Phi_{\underline{k}} : R[T_1, \ldots, T_n] \to R,
\quad
T_i \longmapsto \alpha_{ik_i}
$$
令 $J_{\underline{k}} = \Phi_{\underline{k}}(J)$。则
$\Spec(S) \to \Spec(R)$ 的像等于 $V(\bigcap J_{\underline{k}})$。
\end{lemma}

\begin{proof}
该引理不证自明。提示：
$V(\bigcap J_{\underline{k}}) = \bigcup V(J_{\underline{k}})$。
\end{proof}

\noindent
下列结果归功于 Ferrand，见 \cite{Ferrand}。

\begin{lemma}
\label{more-algebra-lemma-descent-flatness-injective-finite-Noetherian-rings}
设 $R \to S$ 为 Noether 环的有限单射同态。设 $M$ 为 $R$-模。
若 $M \otimes_R S$ 是平坦 $S$-模，则 $M$ 是平坦 $R$-模。
\end{lemma}

\begin{proof}
设 $M$ 为满足 $M \otimes_R S$ 在 $S$ 上平坦的 $R$-模。
由 Algebra, Lemma \ref{algebra-lemma-flatness-descends}，为证明 $M$
平坦，可用任意忠实平坦环扩张替换 $R$。由引理
\ref{more-algebra-lemma-finite-split}，可找到有限局部自由环扩张 $R \subset R'$，
使得对某个理想 $J \subset R'[T_1, \ldots, T_n]$，有
$S' = S \otimes_R R' = R'[T_1, \ldots, T_n]/J$；该理想包含形如
$P_i(T_i)$ 的元素，其中对某些 $\alpha_{ij} \in R'$，有
$P_i(T) = \prod_{j = 1, \ldots, d_i} (T - \alpha_{ij})$。
注意，$R'$ 是 Noether 环，且 $R' \subset S'$ 是有限环扩张。
因此，可用 $R'$ 替换 $R$，并假设 $S$ 具有引理
\ref{more-algebra-lemma-split-image} 中的表示。注意，
$\Spec(S) \to \Spec(R)$ 是满射，见 Algebra, Lemma
\ref{algebra-lemma-integral-overring-surjective}。于是利用引理
\ref{more-algebra-lemma-split-image}，推出 $I = \bigcap J_{\underline{k}}$ 是满足
$V(I) = \Spec(R)$ 的理想。这意味着 $I \subset \sqrt{(0)}$；由于
$R$ 是 Noether 环，$I$ 幂零。映射 $\Phi_{\underline{k}}$ 诱导交换图
$$
\xymatrix{
S \ar[rr] & & R/J_{\underline{k}} \\
& R \ar[lu] \ar[ru]
}
$$
由此推出 $M/J_{\underline{k}}M$ 在 $R/J_{\underline{k}}$ 上平坦。
由引理 \ref{more-algebra-lemma-intersection-flat}，$M/IM$ 在 $R/I$ 上平坦。
最后应用 Algebra, Lemma \ref{algebra-lemma-lift-flatness}，推出 $M$
在 $R$ 上平坦。
\end{proof}

\begin{lemma}
\label{more-algebra-lemma-descent-flatness-injective-integral}
设 $R \to S$ 为单射整环映射。设 $M$ 为
$R[x_1, \ldots, x_n]$ 上的有限表示模。若 $M \otimes_R S$ 在 $S$ 上
平坦，则 $M$ 在 $R$ 上平坦。
\end{lemma}

\begin{proof}
选取表示
$$
R[x_1, \ldots, x_n]^{\oplus t} \to R[x_1, \ldots, x_n]^{\oplus r} \to M \to 0.
$$
设第一个映射由 $r \times t$ 矩阵 $T = (f_{ij})$ 给出，其中
$f_{ij} \in R[x_1, \ldots, x_n]$。写成
$f_{ij} = \sum f_{ij, I} x^I$，其中 $f_{ij, I} \in R$
（多重指标记号）。考虑交换图
$$
\xymatrix{
R \ar[r] & S \\
R_\lambda \ar[u] \ar[r] & S_\lambda \ar[u]
}
$$
其中 $R_\lambda$ 是 $R$ 的有限生成 $\mathbf{Z}$-子代数，包含所有
$f_{ij, I}$；$S_\lambda$ 是 $S$ 的有限
$R_\lambda$-子代数。令 $M_\lambda$ 为有限
$R_\lambda[x_1, \ldots, x_n]$-模，它由上述同样的表示定义，使用同一
矩阵 $T$，但现在把它视为 $R_\lambda[x_1, \ldots, x_n]$ 上的矩阵。
注意，$S$ 是 $S_\lambda$ 的有向余极限（略去细节）。由
Algebra, Lemma \ref{algebra-lemma-flat-finite-presentation-limit-flat}，
对某个 $\lambda$，模
$M_\lambda \otimes_{R_\lambda} S_\lambda$ 在 $S_\lambda$ 上平坦。
由引理 \ref{more-algebra-lemma-descent-flatness-injective-finite-Noetherian-rings}，
$M_\lambda$ 在 $R_\lambda$ 上平坦。由于
$M = M_\lambda \otimes_{R_\lambda} R$，结论由 Algebra, Lemma
\ref{algebra-lemma-flat-base-change} 得出。
\end{proof}

\begin{lemma}
\label{more-algebra-lemma-descent-projective-injective-finite-Noetherian-rings}
设 $R \to S$ 为 Noether 环的有限单射同态。设 $P$ 为 $R$-模。
若 $P \otimes_R S$ 是投射 $S$-模，则 $P$ 是投射 $R$-模。
\end{lemma}

\begin{proof}
设 $P$ 为满足 $P \otimes_R S$ 在 $S$ 上投射的 $R$-模。
由 Algebra, Theorem \ref{algebra-theorem-ffdescent-projectivity}，
为证明 $P$ 投射，可用任意忠实平坦环扩张替换 $R$。由引理
\ref{more-algebra-lemma-finite-split}，可找到有限局部自由环扩张 $R \subset R'$，
使得对某个理想 $J \subset R'[T_1, \ldots, T_n]$，有
$S' = S \otimes_R R' = R'[T_1, \ldots, T_n]/J$；该理想包含形如
$P_i(T_i)$ 的元素，其中对某些 $\alpha_{ij} \in R'$，有
$P_i(T) = \prod_{j = 1, \ldots, d_i} (T - \alpha_{ij})$。
注意，$R'$ 是 Noether 环，且 $R' \subset S'$ 是有限环扩张。
因此，可用 $R'$ 替换 $R$，并假设 $S$ 具有引理
\ref{more-algebra-lemma-split-image} 中的表示。注意，
$\Spec(S) \to \Spec(R)$ 是满射，见 Algebra, Lemma
\ref{algebra-lemma-integral-overring-surjective}。于是利用引理
\ref{more-algebra-lemma-split-image}，推出 $I = \bigcap J_{\underline{k}}$ 是满足
$V(I) = \Spec(R)$ 的理想。这意味着 $I \subset \sqrt{(0)}$；由于
$R$ 是 Noether 环，$I$ 幂零。映射 $\Phi_{\underline{k}}$ 诱导交换图
$$
\xymatrix{
S \ar[rr] & & R/J_{\underline{k}} \\
& R \ar[lu] \ar[ru]
}
$$
由此推出 $P/J_{\underline{k}}P$ 在 $R/J_{\underline{k}}$ 上投射。
由 Algebra, Lemma \ref{algebra-lemma-projective-modulo-two-ideals}，
$P/IP$ 在 $R/I$ 上投射。由引理
\ref{more-algebra-lemma-descent-flatness-injective-finite-Noetherian-rings}，$P$ 在
$R$ 上平坦；故可应用 Algebra, Lemma
\ref{algebra-lemma-lift-projective}，推出 $P$ 在 $R$ 上投射。
\end{proof}








\section{无挠模}
\label{more-algebra-section-torsion-flat}

\noindent
本节讨论无挠模及其与平坦性的关系（尤其是在一维环上）。

\begin{definition}
\label{more-algebra-definition-torsion}
设 $R$ 为整环，$M$ 为 $R$-模。
\begin{enumerate}
\item 若存在非零 $f \in R$ 使得 $fx = 0$，则称元素 $x \in M$
为{\it 挠元}；
\item 若 $M$ 的唯一挠元是 $0$，则称 $M$ {\it 无挠}；
\item 若 $M$ 的每个元素都是挠元，则称 $M$ 为{\it 挠模}。
\end{enumerate}
\end{definition}

\noindent
设 $R$ 为整环，$S = R \setminus \{0\}$ 为 $R$ 中非零元素所成的
乘法集。则 $R$-模 $M$ 无挠，当且仅当 $M \to S^{-1}M$ 是单射。
换言之，当且仅当映射 $M \to M \otimes_R K$ 是单射，其中
$K = S^{-1}R$ 是 $R$ 的分式域。

\begin{lemma}
\label{more-algebra-lemma-torsion}
设 $R$ 为整环，$M$ 为 $R$-模。$M$ 的挠元集合 $M_{tors}$ 是映射
$M \to M \otimes_R K$ 的核。因此 $M_{tors}$ 是 $M$ 的
$R$-子模。商模 $M/M_{tors}$ 无挠。
\end{lemma}

\begin{proof}
见上面的讨论。
\end{proof}

\begin{lemma}
\label{more-algebra-lemma-localize-torsion}
设 $R$ 为整环，$M$ 为无挠 $R$-模。对任意乘法集
$S \subset R$，模 $S^{-1}M$ 是无挠 $S^{-1}R$-模。
\end{lemma}

\begin{proof}
略。
\end{proof}

\begin{lemma}
\label{more-algebra-lemma-flat-pullback-torsion}
设 $R \to R'$ 为整环之间的平坦同态。若 $M$ 是无挠 $R$-模，
则 $M \otimes_R R'$ 是无挠 $R'$-模。
\end{lemma}

\begin{proof}
若 $M$ 无挠，则 $M \subset M \otimes_R K$ 是单射，其中 $K$ 是
$R$ 的分式域。由于 $R'$ 在 $R$ 上平坦，映射
$M \otimes_R R' \to (M \otimes_R K) \otimes_R R'$ 是单射。
由于 $M \otimes_R K$ 同构于若干个 $K$ 的直和，只需说明
$K \otimes_R R'$ 无挠。这成立，因为它是 $R'$ 的局部化。
\end{proof}

\begin{lemma}
\label{more-algebra-lemma-extension-torsion-free}
设 $R$ 为整环。设 $0 \to M \to M' \to M'' \to 0$ 为
$R$-模的短正合列。若 $M$ 与 $M''$ 无挠，则 $M'$ 无挠。
\end{lemma}

\begin{proof}
略。
\end{proof}

\begin{lemma}
\label{more-algebra-lemma-check-torsion}
设 $R$ 为整环，$M$ 为 $R$-模。则 $M$ 无挠，当且仅当对 $R$ 的所有
极大理想 $\mathfrak m$，$M_\mathfrak m$ 都是无挠
$R_\mathfrak m$-模。
\end{lemma}

\begin{proof}
略。提示：使用引理 \ref{more-algebra-lemma-localize-torsion} 与 Algebra, Lemma
\ref{algebra-lemma-characterize-zero-local}。
\end{proof}

\begin{lemma}
\label{more-algebra-lemma-finite-torsion-free-submodule-free}
设 $R$ 为整环，$M$ 为有限 $R$-模。则 $M$ 无挠，当且仅当 $M$ 是某个
有限自由模的子模。
\end{lemma}

\begin{proof}
若 $M$ 是 $R^{\oplus n}$ 的子模，则 $M$ 无挠。反之，假设 $M$ 无挠。
令 $K$ 为 $R$ 的分式域。则 $M \otimes_R K$ 是有限维 $K$-向量空间。
选取该向量空间的一组基 $e_1, \ldots, e_r$。令
$x_1, \ldots, x_n$ 为 $M$ 的生成元。对某些
$a_{ij}, b_{ij} \in R$ 且 $b_{ij} \not = 0$，写成
$x_i = \sum (a_{ij}/b_{ij}) e_j$。令 $b = \prod_{i, j} b_{ij}$。
由于 $M$ 无挠，映射 $M \to M \otimes_R K$ 是单射，并且其像包含于
$R^{\oplus r} = R e_1/b \oplus \ldots \oplus Re_r/b$。
\end{proof}

\begin{lemma}
\label{more-algebra-lemma-torsion-free-finite-noetherian-domain}
设 $R$ 为 Noether 整环，$M$ 为非零有限 $R$-模。下列条件等价：
\begin{enumerate}
\item $M$ 无挠；
\item $M$ 是有限自由模的子模；
\item $(0)$ 是 $M$ 的唯一伴随素理想；
\item $(0)$ 属于 $M$ 的支撑，且 $M$ 具有性质 $(S_1)$；
\item $(0)$ 属于 $M$ 的支撑，且 $M$ 没有嵌入伴随素理想。
\end{enumerate}
\end{lemma}

\begin{proof}
我们已在引理 \ref{more-algebra-lemma-finite-torsion-free-submodule-free} 中看到 (1)
与 (2) 等价，并在 Algebra, Lemma
\ref{algebra-lemma-criterion-no-embedded-primes} 中看到 (4) 与 (5) 等价。
(3) 与 (5) 的等价性由定义立即得出。无挠模的局部化无挠
（引理 \ref{more-algebra-lemma-localize-torsion}），因此显然 $M$ 没有不同于
$(0)$ 的伴随素理想。故 (1) 蕴含 (5)。反之，假设 (5)。
若 $M$ 有挠，则对 $R$ 的某个非零理想 $I$，存在嵌入
$R/I \subset M$。因此 $M$ 有不同于 $(0)$ 的伴随素理想
（见 Algebra, Lemmas \ref{algebra-lemma-ass} and
\ref{algebra-lemma-ass-zero}）。这是嵌入伴随素理想，与假设矛盾。
\end{proof}

\begin{lemma}
\label{more-algebra-lemma-flat-torsion-free}
设 $R$ 为整环。任意平坦 $R$-模都无挠。
\end{lemma}

\begin{proof}
若 $x \in R$ 非零，则 $x : R \to R$ 是单射；因此，若 $M$ 在 $R$ 上
平坦，则 $x : M \to M$ 是单射。故若 $M$ 在 $R$ 上平坦，则 $M$ 无挠。
\end{proof}

\begin{lemma}
\label{more-algebra-lemma-valuation-ring-torsion-free-flat}
设 $A$ 为赋值环。$A$-模 $M$ 在 $A$ 上平坦当且仅当 $M$ 无挠。
\end{lemma}

\begin{proof}
蕴含关系``平坦 $\Rightarrow$ 无挠''是引理
\ref{more-algebra-lemma-flat-torsion-free}。反之，假设 $M$ 无挠。由平坦性的方程判别
准则（见 Algebra, Lemma \ref{algebra-lemma-flat-eq}），必须证明 $M$
中的每个关系都是平凡的。为此，假设
$\sum_{i = 1, \ldots, n} a_i x_i = 0$，其中
$x_i \in M$ 且 $a_i \in A$。重新编号后，可假设对所有 $i$ 都有
$v(a_1) \leq v(a_i)$。故对某个 $a'_i \in A$，可写成
$a_i = a'_i a_1$。注意 $a'_1 = 1$。由于 $M$ 无挠，有
$x_1 = - \sum_{i \geq 2} a'_i x_i$。因此，若选取
$y_i = x_i$、$i = 2, \ldots, n$，则
$$
x_1 = \sum\nolimits_{j \geq 2} -a'_j y_j, \quad
x_i = y_i, (i \geq 2)\quad
0 = a_1 \cdot (-a'_j) + a_j \cdot 1 (j \geq 2)
$$
表明该关系平凡（明确地说，元素 $a_{ij}$ 定义如下：令
$a_{11} = 0$；对 $j > 1$，令 $a_{1j} = -a'_j$；对
$i, j \geq 2$，令 $a_{ij} = \delta_{ij}$）。
\end{proof}

\begin{lemma}
\label{more-algebra-lemma-dedekind-torsion-free-flat}
设 $A$ 为 Dedekind 整环（例如离散赋值环，或更一般的 PID）。
\begin{enumerate}
\item $A$-模平坦当且仅当它无挠。
\item 有限无挠 $A$-模有限局部自由。
\item 若 $A$ 是 PID，则有限无挠 $A$-模有限自由。
\end{enumerate}
\end{lemma}

\begin{proof}
（关于引理陈述中的括号说明，见 Algebra, Lemma
\ref{algebra-lemma-PID-dedekind}。）证明 (1)。由引理
\ref{more-algebra-lemma-check-torsion} 与 Algebra, Lemma
\ref{algebra-lemma-flat-localization}，只需对极大理想
$\mathfrak m \subset A$ 在 $A_\mathfrak m$ 上检验该陈述。
由于 $A_\mathfrak m$ 是离散赋值环
（Algebra, Lemma \ref{algebra-lemma-characterize-Dedekind}），
结论由引理 \ref{more-algebra-lemma-valuation-ring-torsion-free-flat} 得出。

\medskip\noindent
证明 (2)。它由 Algebra, Lemma \ref{algebra-lemma-finite-projective}
与 (1) 得出。

\medskip\noindent
证明 (3)。设 $A$ 为 PID，$M$ 为有限无挠模。由引理
\ref{more-algebra-lemma-finite-torsion-free-submodule-free}，对某个 $n$ 有
$M \subset A^{\oplus n}$。对 $n$ 归纳证明 $M$ 自由。
$n = 1$ 的情形恰好表述 $A$ 是 PID 这一事实。若 $n > 1$，令
$M' \subset A^{\oplus n - 1}$ 为其在 $A^{\oplus n}$ 后
$n - 1$ 个直和分量的投影之像。于是得到短正合列
$0 \to I \to M \to M' \to 0$，其中 $I$ 是 $M$ 与
$A^{\oplus n}$ 的第一个直和分量 $A$ 的交。由归纳法，$M$ 是有限自由
$A$-模的扩张，因而有限自由。
\end{proof}

\begin{lemma}
\label{more-algebra-lemma-hom-into-torsion-free}
设 $R$ 为整环，$M$、$N$ 为 $R$-模。若 $N$ 无挠，则
$\Hom_R(M, N)$ 也无挠。
\end{lemma}

\begin{proof}
选取满射 $\bigoplus_{i \in I} R \to M$。则
$\Hom_R(M, N) \subset \prod_{i \in I} N$。
\end{proof}





\section{模的秩}
\label{more-algebra-section-rank}

\noindent
下面给出我们的定义。

\begin{definition}
\label{more-algebra-definition-rank}
设 $R$ 为分式域为 $K$ 的整环。$R$-模 $M$ 的{\it 秩}为
$\dim_K(M \otimes_R K)$。
\end{definition}

\noindent
当两个定义都适用时，该定义与 Algebra, Definition
\ref{algebra-definition-locally-free} 中秩为 $r$ 的局部自由模概念不冲突。

\begin{lemma}
\label{more-algebra-lemma-torsion-does-not-matter-rank}
设 $R$ 为整环。若 $R$-模映射 $M \to M'$ 的核与余核都是挠模，
则 $M$ 的秩等于 $M'$ 的秩。
\end{lemma}

\begin{proof}
略。提示：若 $K$ 是 $R$ 的分式域，则诱导映射
$M \otimes_R K \to M' \otimes_R K$ 是同构。
\end{proof}

\begin{lemma}
\label{more-algebra-lemma-rank-additive}
设 $R$ 为整环。设 $0 \to M \to M' \to M'' \to 0$ 为 $R$-模的
短正合列。则 $M'$ 的秩等于 $M$ 与 $M''$ 的秩之和。
\end{lemma}

\begin{proof}
略。
\end{proof}

\begin{lemma}
\label{more-algebra-lemma-rank-tensor}
设 $R$ 为整环，$M$ 和 $N$ 为 $R$-模。
\begin{enumerate}
\item $M \otimes_R N$ 的秩等于 $M$ 与 $N$ 的秩之积。
\item 若 $M$ 是有限表示 $R$-模，则 $\Hom_R(M, N)$ 的秩等于
$M$ 与 $N$ 的秩之积。
\end{enumerate}
\end{lemma}

\begin{proof}
第 (1) 项由向量空间的相应事实得出。假设 $M$ 作为 $R$-模有限表示。
由 Algebra, Lemma \ref{algebra-lemma-hom-from-finitely-presented}，
作为 $K$-向量空间，$\Hom_R(M, N) \otimes_R K$ 同构于
$\Hom_K(M \otimes_R K, N \otimes_R K)$。又由于
$M \otimes_R K$ 有限维，结论由向量空间的相应事实得出。
\end{proof}

\begin{lemma}
\label{more-algebra-lemma-dominant-pullback-rank}
设 $R \subset R'$ 为整环扩张。若 $M$ 是 $R$-模，则 $M$ 在 $R$ 上的秩
等于 $M \otimes_R R'$ 在 $R'$ 上的秩。
\end{lemma}

\begin{proof}
这是因为向量空间的维数在基域扩张下不变。
\end{proof}

\begin{lemma}
\label{more-algebra-lemma-finite-module-rank}
设 $R$ 为整环，$M$ 为有限 $R$-模。则：
\begin{enumerate}
\item $M$ 有有限秩 $r \geq 0$；
\item 存在映射 $M \to R^{\oplus r}$，其核与余核都是挠模；
\item 存在 $f \in R$、$f \not = 0$，使得 $M_f$ 是秩为 $r$ 的自由模；
\item 存在单射 $R^{\oplus r} \to M$，其余核是挠模。
\end{enumerate}
\end{lemma}

\begin{proof}
由于 $M$ 有限，$M \otimes_R K$ 是 $R$ 的分式域 $K$ 上的有限维向量
空间。因此秩 $r$ 有限。

\medskip\noindent
选取 $M$ 的生成元 $x_1, \ldots, x_n$，并选取
$M \otimes_R K$ 的一组基 $e_1, \ldots, e_r$。对某些
$a_{ij} , b_{ij} \in R$ 且 $b_{ij} \not = 0$，可写成
$x_i \otimes 1 = \sum (a_{ij}/b_{ij}) e_j$。令
$b = \prod b_{ij}$，并以 $a_{ij}b/b_{ij}$ 替换 $a_{ij}$，可假设
$x_i \otimes 1 = \sum (a_{ij}/b) e_j$。注意，若 $M$ 中有关系
$\sum r_i x_i = 0$，则对所有 $j = 1, \ldots, r$，
$\sum r_ia_{ij} = 0$ 在 $R$ 中成立。

\medskip\noindent
证明 (2)。上面的最后一个观察说明，把 $x_i$ 映到
$\sum a_{ij} e_j$ 的映射 $M \to \bigoplus Re_j$ 良定义。
$e_1, \ldots, e_r$ 构成 $M \otimes_R K$ 的一组基这一事实说明，
该映射与 $K$ 张量后成为同构。因此这是 (2) 中所述的映射。

\medskip\noindent
证明 (3)。在 (2) 中构造的映射 $M \to R^{\oplus r}$ 的余核有限且有挠。
故可找到零化该余核的 $g \in R$、$g \not = 0$。换言之，映射
$M_g \to R_g^{\oplus r}$ 是满射。于是可选取分裂
$M_g = N \oplus R_g^{\oplus r}$。此时 $N$ 是有限 $R_g$-模，且
$N \otimes_{R_g} K = 0$。换言之，$N$ 是 $R_g$ 上的有限挠模。
故可找到零化 $N$ 的 $g' \in R$。令 $f= gg'$，便得到 (3)。

\medskip\noindent
为证明 (4)，选取 $M_f$ 的一组基 $y_1, \ldots, y_r$，其中 $f$
如 (3) 所述。然后对某些 $m_j \in M$ 与 $t_j \geq 0$ 写成
$y_j = m_j/f^{t_j}$。把第 $j$ 个基向量映到 $m_j$ 的映射
$R^{\oplus r} \to M$ 即如 (4) 所述。略去细节。
\end{proof}





\section{自反模}
\label{more-algebra-section-reflexive}

\noindent
下面给出我们的定义。

\begin{definition}
\label{more-algebra-definition-reflexive}
设 $R$ 为整环。若自然映射
$$
j : M \longrightarrow \Hom_R(\Hom_R(M, R), R)
$$
是同构，则称 $R$-模 $M$ {\it 自反}；该映射把 $m \in M$ 映为如下
映射：它把 $\varphi \in \Hom_R(M, R)$ 映到 $\varphi(m) \in R$。
\end{definition}

\noindent
可以对更一般的环作此定义，但即使上面的定义也有缺点。明智的做法是
限于 Noether 整环与有限无挠模，并（或许）对 $R$ 加上某些正则性条件
（例如 $R$ 正规）。

\begin{lemma}
\label{more-algebra-lemma-reflexive-torsion-free}
设 $R$ 为整环，$M$ 为 $R$-模。
\begin{enumerate}
\item 若 $M$ 自反，则 $M$ 无挠。
\item 若 $M$ 有限，则
$j : M \to \Hom_R(\Hom_R(M, R), R)$ 的核与余核都是挠模。
\item 若 $M$ 有限，则 $j$ 是单射当且仅当 $M$ 无挠。
\end{enumerate}
\end{lemma}

\begin{proof}
为证明 (1)，由引理 \ref{more-algebra-lemma-hom-into-torsion-free}，$M$ 的双对偶无挠。
假设 $M$ 有限。选取 $M$ 在 $R$ 上的生成元 $x_1, \ldots, x_n$。
设 $K$ 为 $R$ 的分式域。重新编号后，可以假设对某个
$0 \leq r \leq n$，元 $x_1, \ldots, x_r$ 在 $M \otimes_R K$
中的像构成它在 $K$ 上的一组基。对 $r < i \leq n$，写成
$x_i \otimes 1 = \sum_{1 \leq j \leq r} x_j \otimes k_{ij}$，其中
$k_{ij} \in K$。选取 $c \in R$，使 $a_{ij} = ck_{ij} \in R$，
并使 $cx_i =  \sum_{1 \leq j \leq r} a_{ij} x_j$ 在 $M$ 中成立
（略去一个小细节）。于是得到映射
$$
R^{\oplus r} \xrightarrow{\alpha} M \xrightarrow{\beta} R^{\oplus r}
\xrightarrow{\alpha} M
$$
其中 $\alpha \circ \beta$ 与 $\beta \circ \alpha$ 都是乘以 $c$。
这里 $\alpha$ 由 $x_1, \ldots, x_r$ 给出；$\beta$ 对
$1 \leq i \leq r$ 把 $x_i$ 映到第 $i$ 个基向量的 $c$ 倍，
对 $r < i \leq n$ 则映到向量 $(a_{i1}, \ldots, a_{ir})$。
考虑交换图
$$
\xymatrix{
R^{\oplus r} \ar[r] \ar[d]^{\text{id}} &
M \ar[r] \ar[d]^j &
R^{\oplus r} \ar[d]^{\text{id}} \ar[r] &
M \ar[d]^j \\
R^{\oplus r} \ar[r] &
\Hom_R(\Hom_R(M, R), R) \ar[r] &
R^{\oplus r} \ar[r] &
\Hom_R(\Hom_R(M, R), R)
}
$$
当然，底行中相继箭头的复合都等于乘以 $c$。此图表明
$j$ 的核与余核都被 $c$ 零化，这证明了 (2)。
对 (3)，注意若 $M$ 无挠，则 $\beta$ 是单射，从而 $j$ 是单射。
逆命题是清楚的，因为我们已经看到 $M$ 的双对偶无挠。
\end{proof}

\begin{lemma}
\label{more-algebra-lemma-cokernel-map-double-dual-dvr}
设 $R$ 为离散赋值环，$M$ 为有限 $R$-模。则映射
$j : M \to \Hom_R(\Hom_R(M, R), R)$ 是满射。
\end{lemma}

\begin{proof}
设 $M_{tors} \subset M$ 为挠子模。则
$\Hom_R(M, R) = \Hom_R(M/M_{tors}, R)$（对任意整环都成立）。
因此可以假设 $M$ 无挠。由引理
\ref{more-algebra-lemma-dedekind-torsion-free-flat}，$M$ 自由，引理便是清楚的。
\end{proof}

\begin{lemma}
\label{more-algebra-lemma-check-reflexive}
设 $R$ 为 Noether 整环，$M$ 为有限 $R$-模。下列条件等价：
\begin{enumerate}
\item $M$ 自反；
\item 对每个素理想 $\mathfrak p \subset R$，$M_\mathfrak p$ 都是自反
$R_\mathfrak p$-模；
\item 对 $R$ 的每个极大理想 $\mathfrak m$，$M_\mathfrak m$ 都是自反
$R_\mathfrak m$-模。
\end{enumerate}
\end{lemma}

\begin{proof}
映射 $j : M \to \Hom_R(\Hom_R(M, R), R)$ 在素理想
$\mathfrak p$ 处的局部化，就是 Noether 局部整环 $R_\mathfrak p$ 上的模
$M_\mathfrak p$ 所对应的映射。见 Algebra, Lemma
\ref{algebra-lemma-hom-from-finitely-presented}。因此由 Algebra, Lemma
\ref{algebra-lemma-characterize-zero-local} 即得本引理。
\end{proof}

\begin{lemma}
\label{more-algebra-lemma-sequence-reflexive}
设 $R$ 为 Noether 整环。设 $0 \to M \to M' \to M''$ 是有限
$R$-模的正合列。若 $M'$ 自反且 $M''$ 无挠，则 $M$ 自反。
\end{lemma}

\begin{proof}
下面不再说明地使用这一事实：对任意有限 $R$-模 $N$ 与 $N'$，
$\Hom_R(N, N')$ 都是有限 $R$-模；见 Algebra, Lemma
\ref{algebra-lemma-ext-noetherian}。取对偶得到序列
$$
\Hom_R(M, R) \leftarrow \Hom_R(M', R) \leftarrow \Hom_R(M'', R)
$$
再次取对偶，得到交换图
{\small
$$
\xymatrix{
\Hom_R(\Hom_R(M, R), R) \ar[r]_j &
\Hom_R(\Hom_R(M', R), R) \ar[r] &
\Hom_R(\Hom_R(M'', R), R) \\
M \ar[u] \ar[r] & M' \ar[u] \ar[r] & M'' \ar[u]
}
$$
}
我们并不知道顶行是否正合。但由假设，中间的竖直箭头是同构，
右边的竖直箭头是单射（引理 \ref{more-algebra-lemma-reflexive-torsion-free}）。
我们断言 $j$ 是单射。若承认该断言，作图追迹即知左边的竖直箭头是同构，
也就是说 $M$ 自反。

\medskip\noindent
断言的证明。考虑定义 $Q$ 的正合列
$\Hom_R(M',R)\to \Hom_R(M,R)\to Q \to 0$。应用
Algebra, Lemma \ref{algebra-lemma-hom-from-finitely-presented} 得
$$
\Hom_K(M'\otimes_R K,K) \to
\Hom_K(M\otimes_R K,K) \to
Q\otimes_R K\to 0
$$
但 $M \otimes_R K \to M' \otimes_R K$ 是向量空间之间的单射，
因而是分裂单射，所以 $Q \otimes_R K = 0$，即 $Q$ 为挠模。
于是得到正合列
$$
0 \to \Hom_R(Q,R) \to
\Hom_R(\Hom_R(M,R),R) \to
\Hom_R(\Hom_R(M',R),R)
$$
而因 $Q$ 为挠模，$\Hom_R(Q,R)=0$。
\end{proof}

\begin{lemma}
\label{more-algebra-lemma-characterize-reflexive}
设 $R$ 为 Noether 整环，$M$ 为有限 $R$-模。下列条件等价：
\begin{enumerate}
\item $M$ 自反；
\item 存在短正合列 $0 \to M \to F \to N \to 0$，其中 $F$ 有限自由，
$N$ 无挠。
\end{enumerate}
\end{lemma}

\begin{proof}
注意有限自由模自反。由引理 \ref{more-algebra-lemma-sequence-reflexive}，
(2) 推出 (1)。假设 $M$ 自反。选取一个表示
$R^{\oplus m} \to R^{\oplus n} \to \Hom_R(M, R) \to 0$。
取对偶后得到正合列
$$
0 \to \Hom_R(\Hom_R(M, R), R) \to R^{\oplus n} \to N \to 0
$$
其中 $N = \Im(R^{\oplus n} \to R^{\oplus m})$ 为无挠模。因为
$M = \Hom_R(\Hom_R(M, R), R)$，所以得到 (2) 中所述的正合列。
\end{proof}

\begin{lemma}
\label{more-algebra-lemma-flat-pullback-reflexive}
设 $R \to R'$ 为 Noether 整环之间的平坦同态。若 $M$ 是有限自反
$R$-模，则 $M \otimes_R R'$ 是有限自反 $R'$-模。
\end{lemma}

\begin{proof}
选取短正合列 $0 \to M \to F \to N \to 0$，其中 $F$ 有限自由，
$N$ 无挠；见引理 \ref{more-algebra-lemma-characterize-reflexive}。由于 $R \to R'$ 平坦，
我们得到短正合列
$0 \to M \otimes_R R' \to F \otimes_R R' \to N \otimes_R R' \to 0$，
其中 $F \otimes_R R'$ 有限自由，$N \otimes_R R'$ 无挠
（引理 \ref{more-algebra-lemma-flat-pullback-torsion}）。因此由引理
\ref{more-algebra-lemma-characterize-reflexive}，$M \otimes_R R'$ 自反。
\end{proof}

\begin{lemma}
\label{more-algebra-lemma-dual-reflexive}
设 $R$ 为 Noether 整环，$M$ 为有限 $R$-模，$N$ 为有限自反
$R$-模。则 $\Hom_R(M, N)$ 自反。
\end{lemma}

\begin{proof}
选取一个表示 $R^{\oplus m} \to R^{\oplus n} \to M \to 0$。
于是得到
$$
0 \to \Hom_R(M, N) \to N^{\oplus n} \to N' \to 0
$$
其中 $N' = \Im(N^{\oplus n} \to N^{\oplus m})$ 无挠。由引理
\ref{more-algebra-lemma-sequence-reflexive} 得到结论。
\end{proof}

\begin{definition}
\label{more-algebra-definition-reflexive-hull}
设 $R$ 为 Noether 整环，$M$ 为有限 $R$-模。称模
$M^{**} = \Hom_R(\Hom_R(M, R), R)$ 为 $M$ 的{\it 自反包}。
\end{definition}

\noindent
这一定义是合理的，因为由引理 \ref{more-algebra-lemma-dual-reflexive}，自反包自反。
对应 $M \mapsto M^{**}$ 是一个函子。若 $\varphi : M \to N$ 是到自反
$R$-模 $N$ 的 $R$-模映射，则 $\varphi$ 经过 $M$ 的自反包分解为
$M \to M^{**} \to N$。换言之，在 Noether 整环 $R$ 上，取自反包是包含函子
$$
\text{finite reflexive modules} \subset
\text{finite modules}
$$
的左伴随。

\begin{lemma}
\label{more-algebra-lemma-hom-into-depth}
设 $R$ 为 Noether 局部环，$M$、$N$ 为有限 $R$-模。
\begin{enumerate}
\item 若 $N$ 的深度 $\geq 1$，则 $\Hom_R(M, N)$ 的深度 $\geq 1$。
\item 若 $N$ 的深度 $\geq 2$，则 $\Hom_R(M, N)$ 的深度 $\geq 2$。
\end{enumerate}
\end{lemma}

\begin{proof}
选取一个表示 $R^{\oplus m} \to R^{\oplus n} \to M \to 0$。
取对偶后得到正合列
$$
0 \to \Hom_R(M, N) \to N^{\oplus n} \to N' \to 0
$$
其中 $N' = \Im(N^{\oplus n} \to N^{\oplus m})$。深度 $\geq 1$ 的模的子模深度
$\geq 1$；这直接由定义得出。因此 (1) 是清楚的。对 (2)，注意这里的
假设与上一观察说明 $N'$ 的深度 $\geq 1$。模 $N^{\oplus n}$ 的深度
$\geq 2$。由 Algebra, Lemma \ref{algebra-lemma-depth-in-ses}，可知
$\Hom_R(M, N)$ 的深度 $\geq 2$。
\end{proof}

\begin{lemma}
\label{more-algebra-lemma-hom-into-S2}
设 $R$ 为 Noether 环，$M$、$N$ 为有限 $R$-模。
\begin{enumerate}
\item 若 $N$ 具有性质 $(S_1)$，则 $\Hom_R(M, N)$ 具有性质 $(S_1)$。
\item 若 $N$ 具有性质 $(S_2)$，则 $\Hom_R(M, N)$ 具有性质 $(S_2)$。
\item 若 $R$ 是整环，$N$ 无挠且具有 $(S_2)$，则
$\Hom_R(M, N)$ 无挠且具有性质 $(S_2)$。
\end{enumerate}
\end{lemma}

\begin{proof}
对有限 $R$-模，在素理想处局部化与取 $\Hom_R$ 可交换
（Algebra, Lemma \ref{algebra-lemma-ext-noetherian}），因此 (1) 与 (2)
直接由引理 \ref{more-algebra-lemma-hom-into-depth} 得出。(3) 由 (2) 与引理
\ref{more-algebra-lemma-hom-into-torsion-free} 得出。
\end{proof}

\begin{lemma}
\label{more-algebra-lemma-check-injective-on-ass}
设 $R$ 为 Noether 环，$\varphi : M \to N$ 为 $R$-模的映射。假设对 $R$
的每个素理想 $\mathfrak p$，下列情形至少有一个发生：
\begin{enumerate}
\item $M_\mathfrak p \to N_\mathfrak p$ 是单射；或
\item $\mathfrak p \not \in \text{Ass}(M)$。
\end{enumerate}
则 $\varphi$ 是单射。
\end{lemma}

\begin{proof}
设 $\mathfrak p$ 为 $\Ker(\varphi)$ 的一个伴随素理想。则存在元
$x \in M_\mathfrak p$，它属于 $M_\mathfrak p \to N_\mathfrak p$ 的核，且其零化子为
$\mathfrak pR_\mathfrak p$（Algebra, Lemma
\ref{algebra-lemma-associated-primes-localize}）。在上述两种情形中这都不可能。
因此 $\text{Ass}(\Ker(\varphi)) = \emptyset$，再由 Algebra, Lemma
\ref{algebra-lemma-ass-zero} 得 $\Ker(\varphi) = 0$。
\end{proof}

\begin{lemma}
\label{more-algebra-lemma-check-isomorphism-via-depth-and-ass}
设 $R$ 为 Noether 环，$\varphi : M \to N$ 为 $R$-模的映射。假设 $M$ 有限，
且对 $R$ 的每个素理想 $\mathfrak p$，下列情形之一发生：
\begin{enumerate}
\item $M_\mathfrak p \to N_\mathfrak p$ 是同构；或
\item $\text{depth}(M_\mathfrak p) \geq 2$ 且
$\mathfrak p \not \in \text{Ass}(N)$。
\end{enumerate}
则 $\varphi$ 是同构。
\end{lemma}

\begin{proof}
由引理 \ref{more-algebra-lemma-check-injective-on-ass}，$\varphi$ 是单射。
设 $N' \subset N$ 是包含 $M$ 的像的有限生成 $R$-模。那么
$\text{Ass}(N_\mathfrak p) = \emptyset$ 推出
$\text{Ass}(N'_\mathfrak p) = \emptyset$。因此本引理的假设对 $M \to N'$ 成立。
为了证明 $\varphi$ 是同构，只须对每个这样的 $N' \subset N$ 证明同一结论。
因此可以假设 $N$ 是有限 $R$-模。在此情形下，
$\mathfrak p \not \in \text{Ass}(N) \Rightarrow
\text{depth}(N_\mathfrak p) \geq 1$；见 Algebra, Lemma
\ref{algebra-lemma-ideal-nonzerodivisor}。考虑定义 $Q$ 的短正合列
$$
0 \to M \to N \to Q \to 0
$$
检视各条件可知：在情形 (1) 中 $Q_\mathfrak p = 0$，而在情形 (2) 中，
由 Algebra, Lemma \ref{algebra-lemma-depth-in-ses} 有
$\text{depth}(Q_\mathfrak p) \geq 1$。这说明 $Q$ 没有任何伴随素理想，
因此由 Algebra, Lemma \ref{algebra-lemma-ass-zero} 得 $Q = 0$。
\end{proof}

\begin{lemma}
\label{more-algebra-lemma-isom-depth-2-torsion-free}
设 $R$ 为 Noether 整环，$\varphi : M \to N$ 为 $R$-模的映射。假设
$M$ 有限、$N$ 无挠，且对 $R$ 的每个素理想 $\mathfrak p$，下列情形之一发生：
\begin{enumerate}
\item $M_\mathfrak p \to N_\mathfrak p$ 是同构；或
\item $\text{depth}(M_\mathfrak p) \geq 2$。
\end{enumerate}
则 $\varphi$ 是同构。
\end{lemma}

\begin{proof}
这是引理 \ref{more-algebra-lemma-check-isomorphism-via-depth-and-ass} 的特殊情形。
\end{proof}

\begin{lemma}
\label{more-algebra-lemma-reflexive-depth-2}
设 $R$ 为 Noether 整环，$M$ 为有限 $R$-模。下列条件等价：
\begin{enumerate}
\item $M$ 自反；
\item 对 $R$ 的每个素理想 $\mathfrak p$，下列情形之一发生：
\begin{enumerate}
\item $M_\mathfrak p$ 是自反 $R_\mathfrak p$-模；或
\item $\text{depth}(M_\mathfrak p) \geq 2$。
\end{enumerate}
\end{enumerate}
\end{lemma}

\begin{proof}
若 (1) 成立，则由引理 \ref{more-algebra-lemma-check-reflexive}，对每个素理想
$\mathfrak p$，$M_\mathfrak p$ 都是自反模。因此 (1) $\Rightarrow$ (2)。
假设 (2)。令 $N = \Hom_R(\Hom_R(M, R), R)$，则对 $R$ 的每个素理想
$\mathfrak p$ 有
$$
N_\mathfrak p =
\Hom_{R_\mathfrak p}(
\Hom_{R_\mathfrak p}(M_\mathfrak p, R_\mathfrak p), R_\mathfrak p)
$$
见 Algebra, Lemma \ref{algebra-lemma-hom-from-finitely-presented}。
对映射 $j : M \to N$ 应用引理
\ref{more-algebra-lemma-isom-depth-2-torsion-free}。这是允许的，因为 $M$ 有限，而由引理
\ref{more-algebra-lemma-hom-into-torsion-free}，$N$ 无挠。在情形 (2)(a) 中，映射
$M_\mathfrak p \to N_\mathfrak p$ 是同构；在情形 (2)(b) 中，有
$\text{depth}(M_\mathfrak p) \geq 2$。
\end{proof}

\begin{lemma}
\label{more-algebra-lemma-reflexive-S2}
设 $R$ 为 Noether 整环，$M$ 为有限自反 $R$-模。设
$\mathfrak p \subset R$ 为素理想。
\begin{enumerate}
\item 若 $\text{depth}(R_\mathfrak p) \geq 2$，则
$\text{depth}(M_\mathfrak p) \geq 2$。
\item 若 $R$ 是 $(S_2)$，则 $M$ 是 $(S_2)$。
\end{enumerate}
\end{lemma}

\begin{proof}
由于形成自反包 $\Hom_R(\Hom_R(M, R), R)$ 与局部化可交换
（Algebra, Lemma \ref{algebra-lemma-hom-from-finitely-presented}），
(1) 由引理 \ref{more-algebra-lemma-hom-into-depth} 得出。(2) 直接由引理
\ref{more-algebra-lemma-hom-into-S2} 得出。
\end{proof}

\begin{example}
\label{more-algebra-example-ring-not-S2}
上下的结果使人猜想自反性与 $(S_2)$ 条件有关；下面的例子用来阻止过于乐观的猜想。
设 $k$ 为域。设 $R$ 是 $k[x, y]$ 由 $1, y, x^2, xy, x^3$ 生成的
$k$-子代数。则 $R$ 不是 $(S_2)$。所以把 $R$ 看作 $R$-模，它是一个自反但不是
$(S_2)$ 的 $R$-模之例。设 $M = k[x, y]$，并将其视为 $R$-模。
则 $M$ 是自反 $R$-模，因为
$$
\Hom_R(M, R) = \mathfrak m = (y, x^2, xy, x^3)
\quad\text{且}\quad
\Hom_R(\mathfrak m, R) = M
$$
且 $M$ 作为 $R$-模是 $(S_2)$（略去计算）。因此 $R$ 是一个 Noether 整环，
它拥有自反 $(S_2)$ 模，但 $R$ 本身不是 $(S_2)$。
\end{example}

\begin{lemma}
\label{more-algebra-lemma-reflexive-over-normal}
设 $R$ 为 Noether 正规整环，其分式域为 $K$。设 $M$ 为有限 $R$-模。
下列条件等价：
\begin{enumerate}
\item $M$ 自反；
\item $M$ 无挠且具有性质 $(S_2)$；
\item $M$ 无挠，且
$M = \bigcap_{\text{高度}(\mathfrak p) = 1} M_{\mathfrak p}$，
其中交取在 $M_K = M \otimes_R K$ 中。
\end{enumerate}
\end{lemma}

\begin{proof}
由 Algebra, Lemma \ref{algebra-lemma-criterion-normal}，$R$ 满足 $(R_1)$ 与 $(S_2)$。

\medskip\noindent
假设 (1)。由引理 \ref{more-algebra-lemma-reflexive-torsion-free}，$M$ 无挠；
由引理 \ref{more-algebra-lemma-reflexive-S2}，它满足 $(S_2)$。因此 (2) 成立。

\medskip\noindent
假设 (2)。依定义，
$M' = \bigcap_{\text{高度}(\mathfrak p) = 1} M_{\mathfrak p}$
是映射
$$
M_K
\longrightarrow
\bigoplus\nolimits_{\text{高度}(\mathfrak p) = 1} M_K/M_\mathfrak p
\subset
\prod\nolimits_{\text{高度}(\mathfrak p) = 1} M_K/M_\mathfrak p
$$
的核。注意，我们的映射确实如所示经过直和分解，因为给定
$a/b \in K$，至多只有有限多个高为 $1$ 的素理想 $\mathfrak p$ 满足
$b \in \mathfrak p$。设 $\mathfrak p_0$ 为一个高为 $1$ 的素理想。则除非两者相等，
否则 $(M_K/M_\mathfrak p)_{\mathfrak p_0} = 0$；
而在 $\mathfrak p = \mathfrak p_0$ 时，我们得到
$(M_K/M_\mathfrak p)_{\mathfrak p_0} = M_K/M_{\mathfrak p_0}$。
因此，由局部化的正合性以及局部化与直和可交换，有
$M'_{\mathfrak p_0} = M_{\mathfrak p_0}$。因为 $M$ 在高 $> 1$ 的素理想处深度
$\geq 2$，由引理 \ref{more-algebra-lemma-isom-depth-2-torsion-free}，$M \to M'$ 是同构。
因此 (3) 成立。

\medskip\noindent
假设 (3)。设 $\mathfrak p$ 为高 $1$ 的素理想。由 $(R_1)$，
$R_\mathfrak p$ 是离散赋值环。由引理 \ref{more-algebra-lemma-dedekind-torsion-free-flat}，
$M_\mathfrak p$ 有限自由，特别地它自反。因此映射 $M \to M^{**}$ 在所有高为
$1$ 的素理想 $\mathfrak p$ 处诱导出同构。于是条件
$M = \bigcap_{\text{高度}(\mathfrak p) = 1} M_{\mathfrak p}$
推出 $M = M^{**}$，从而 (1) 成立。
\end{proof}

\begin{lemma}
\label{more-algebra-lemma-describe-reflexive-hull}
设 $R$ 为 Noether 正规整环，$M$ 为有限 $R$-模。则 $M$ 的自反包是交
$$
M^{**} =
\bigcap\nolimits_{\text{高度}(\mathfrak p) = 1}
M_{\mathfrak p}/(M_\mathfrak p)_{tors} =
\bigcap\nolimits_{\text{高度}(\mathfrak p) = 1}
(M/M_{tors})_\mathfrak p
$$
，其交取在 $M \otimes_R K$ 中。
\end{lemma}

\begin{proof}
设 $\mathfrak p$ 为高 $1$ 的素理想。映射
$M_\mathfrak p \to M \otimes_R K$ 的核是 $M_\mathfrak p$ 的挠子模
$(M_\mathfrak p)_{tors}$。此外，
$(M/M_{tors})_\mathfrak p = M_\mathfrak p/(M_\mathfrak p)_{tors}$，
并且它是离散赋值环 $R_\mathfrak p$ 上的有限自由模
（引理 \ref{more-algebra-lemma-dedekind-torsion-free-flat}）。于是
$M_\mathfrak p/(M_\mathfrak p)_{tors} \to
(M_\mathfrak p)^{**} = (M^{**})_\mathfrak p$
是同构，因而本引理是引理 \ref{more-algebra-lemma-reflexive-over-normal} 的结论。
\end{proof}

\begin{lemma}
\label{more-algebra-lemma-integral-closure-reflexive}
设 $A$ 为 Noether 正规整环，分式域为 $K$。设 $L$ 是 $K$ 的有限扩张。
若 $A$ 在 $L$ 中的整闭包 $B$ 在 $A$ 上有限，则 $B$ 作为 $A$-模自反。
\end{lemma}

\begin{proof}
只须证明 $B = \bigcap B_\mathfrak p$，其中交遍历高为 $1$ 的素理想
$\mathfrak p \subset A$；见引理 \ref{more-algebra-lemma-reflexive-over-normal}。
设 $b \in \bigcap B_\mathfrak p$。设
$x^d + a_1x^{d - 1} + \ldots + a_d$ 是 $b$ 在 $K$ 上的最小多项式。
我们要证明 $a_i \in A$。由 Algebra, Lemma
\ref{algebra-lemma-minimal-polynomial-normal-domain}，对所有 $i$ 及所有高一素理想
$\mathfrak p$，有 $a_i \in A_\mathfrak p$。因此由 Algebra, Lemma
\ref{algebra-lemma-normal-domain-intersection-localizations-height-1}
得到所需结论（也可使用已引用的引理，因为 $A$ 是其自身上的自反模）。
\end{proof}








\section{内容理想}
\label{more-algebra-section-content}

\noindent
这个定义可能与你预期的不同。

\begin{definition}
\label{more-algebra-definition-content-ideal}
设 $A$ 为环，$M$ 为平坦 $A$-模，$x \in M$。若 $A$ 中所有满足
$x \in IM$ 的理想 $I$ 所成的集合有最小元，则称该最小元为
{\it $x$ 的内容理想}。
\end{definition}

\noindent
注意，因为 $M$ 在 $A$ 上平坦，对 $A$ 的一对理想 $I, I'$，有
$IM \cap I'M = (I \cap I')M$。这可由把正合列
$0 \to I \cap I' \to I \oplus I' \to I + I' \to 0$ 与 $M$ 张量而得。

\begin{lemma}
\label{more-algebra-lemma-content-finitely-generated}
设 $A$ 为环，$M$ 为平坦 $A$-模，$x \in M$。若 $x$ 的内容理想存在，
则它有限生成。
\end{lemma}

\begin{proof}
设 $x \in IM$。则可写成 $x = \sum_{i = 1, \ldots, n} f_i x_i$，其中
$f_i \in I$ 且 $x_i \in M$。因此 $x \in I'M$，其中
$I' = (f_1, \ldots, f_n)$。
\end{proof}

\begin{lemma}
\label{more-algebra-lemma-equal-content}
设 $(A, \mathfrak m)$ 为局部环。设 $u : M \to N$ 为平坦 $A$-模之间的映射，
且 $\overline{u} : M/\mathfrak m M \to N/\mathfrak m N$ 是单射。若 $x \in M$
的内容理想为 $I$，则 $u(x)$ 的内容理想也为 $I$。
\end{lemma}

\begin{proof}
显然 $u(x) \in IN$。若 $u(x) \in I'N$，则
$u(x) \in (I' \cap I)N$；见定义 \ref{more-algebra-definition-content-ideal} 后的讨论。
因此只须证明：若 $x \in I'N$ 且 $I' \subset I$、$I' \not =  I$，则
$u(x) \not \in I'N$。因为 $I/I'$ 是非零有限 $A$-模
（引理 \ref{more-algebra-lemma-content-finitely-generated}），由 Nakayama 引理
（Algebra, Lemma \ref{algebra-lemma-NAK}），存在 $A$-模的非零映射
$\chi : I/I' \to A/\mathfrak m$。由于 $I$ 是 $x$ 的内容理想，有
$x \not \in I''M$，其中 $I'' = \Ker(\chi)$。因而 $x$ 不在映射
$$
IM = I \otimes_A M \xrightarrow{\chi \otimes 1}
A/\mathfrak m \otimes M \cong M/\mathfrak m M
$$
的核中。应用关于 $\overline{u}$ 的假设，可知 $u(x)$ 在映射
$$
IN = I \otimes_A N \xrightarrow{\chi \otimes 1}
A/\mathfrak m \otimes N \cong N/\mathfrak m N
$$
下的像不为零，于是得到结论。
\end{proof}

\begin{lemma}
\label{more-algebra-lemma-content-exists-flat-Mittag-Leffler}
设 $A$ 为环，$M$ 为平坦 Mittag-Leffler 模。则 $M$ 的每个元都有内容理想。
\end{lemma}

\begin{proof}
这是 Algebra, Lemma \ref{algebra-lemma-flat-ML-criterion} 的特殊情形。
\end{proof}








\section{平坦性与有限性条件}
\label{more-algebra-section-flat-finite}

\noindent
本节讨论“平坦 $+$ 有限型 $\Rightarrow$ 有限表示”这一类型的一些推论。
我们将在平坦性一章重新讨论这一结果；见 More on Flatness, Section
\ref{flat-section-introduction}。这种类型的第一个结果已在 Algebra, Lemma
\ref{algebra-lemma-finite-flat-module-finitely-presented} 中证明。

\begin{lemma}
\label{more-algebra-lemma-flat-finite-type-finite-presentation-local-module}
设 $R$ 为环。设 $S = R[x_1, \ldots, x_n]$ 为 $R$ 上的多项式环。
设 $M$ 为 $S$-模。假设：
\begin{enumerate}
\item $R$ 有有限多个素理想 $\mathfrak p_1, \ldots, \mathfrak p_m$，使映射
$R \to \prod R_{\mathfrak p_j}$ 为单射；
\item $M$ 是有限 $S$-模；
\item $M$ 在 $R$ 上平坦；以及
\item 对 $R$ 的每个素理想 $\mathfrak p$，模 $M_{\mathfrak p}$ 在
$S_{\mathfrak p}$ 上有限表示。
\end{enumerate}
则 $M$ 在 $S$ 上有限表示。
\end{lemma}

\begin{proof}
选取 $M$ 作为 $S$-模的表示
$$
0 \to K \to S^{\oplus r} \to M \to 0
$$
设 $\mathfrak q$ 为 $S$ 的一个素理想，它位于 $R$ 的素理想
$\mathfrak p$ 之上。由假设，存在有限多个元 $k_1, \ldots, k_t \in K$，
使得若令 $K' = \sum Sk_j \subset K$，则
$K'_{\mathfrak p} = K_{\mathfrak p}$ 且对 $j = 1, \ldots, m$ 有
$K'_{\mathfrak p_j} = K_{\mathfrak p_j}$。令 $M' = S^{\oplus r}/K'$，特别可得
$M'_{\mathfrak q} = M_{\mathfrak q}$。由平坦性的开性，见 Algebra, Theorem
\ref{algebra-theorem-openness-flatness}，存在 $g \in S$、$g \not \in \mathfrak q$，
使 $M'_g$ 在 $R$ 上平坦。因此 $M'_g \to M_g$ 是平坦 $R$-模之间的满射。
考虑交换图
$$
\xymatrix{
M'_g \ar[r] \ar[d] & M_g \ar[d] \\
\prod (M'_g)_{\mathfrak p_j} \ar[r] & \prod (M_g)_{\mathfrak p_j}
}
$$
由 $k_1, \ldots, k_t$ 的选取，底部箭头是同构。由于
$R \to \prod R_{\mathfrak p_j}$ 是单射，而 $M'_g$ 在 $R$ 上平坦，
左边竖直箭头是单射。因此顶部水平箭头是单射，从而是同构。
这证明 $M_g$ 在 $S_g$ 上有限表示。应用 Algebra, Lemma
\ref{algebra-lemma-cover} 得到结论。
\end{proof}

\begin{lemma}
\label{more-algebra-lemma-flat-finite-type-finite-presentation-local}
设 $R \to S$ 为环同态。假设：
\begin{enumerate}
\item $R$ 有有限多个素理想 $\mathfrak p_1, \ldots, \mathfrak p_m$，使映射
$R \to \prod R_{\mathfrak p_j}$ 为单射；
\item $R \to S$ 为有限型；
\item $S$ 在 $R$ 上平坦；以及
\item 对 $R$ 的每个素理想 $\mathfrak p$，环 $S_{\mathfrak p}$ 在
$R_{\mathfrak p}$ 上有限表示。
\end{enumerate}
则 $S$ 在 $R$ 上有限表示。
\end{lemma}

\begin{proof}
由假设，$S$ 是 $R$ 上某个多项式环的商。因此结果直接由引理
\ref{more-algebra-lemma-flat-finite-type-finite-presentation-local-module} 得出。
\end{proof}

\begin{lemma}
\label{more-algebra-lemma-flat-graded-finite-type-finite-presentation-module}
设 $R$ 为环。设 $S = R[x_1, \ldots, x_n]$ 为 $R$ 上的分次多项式代数，
也就是说 $\deg(x_i) > 0$ 但不一定等于 $1$。设 $M$ 为分次 $S$-模。
假设：
\begin{enumerate}
\item $R$ 是局部环；
\item $M$ 是有限 $S$-模；以及
\item $M$ 在 $R$ 上平坦。
\end{enumerate}
则 $M$ 作为 $S$-模有限表示。
\end{lemma}

\begin{proof}
设 $M = \bigoplus M_d$ 为 $M$ 上的分次。选取 $M$ 的齐次生成元
$m_1, \ldots, m_r \in M$。设 $\deg(m_i) = d_i \in \mathbf{Z}$。
这给出表示
$$
0 \to K \to \bigoplus\nolimits_{i = 1, \ldots, r} S(-d_i) \to M \to 0
$$
它在每个次数 $d$ 上给出短正合列
$$
0 \to K_d \to \bigoplus\nolimits_{i = 1, \ldots, r} S_{d - d_i} \to
M_d \to 0.
$$
由假设，每个 $M_d$ 都是有限平坦 $R$-模。由 Algebra, Lemma
\ref{algebra-lemma-finite-flat-local}，每个 $M_d$ 都是有限自由 $R$-模。
因此每个 $K_d$ 都是有限 $R$-模。另外，由 Algebra, Lemma
\ref{algebra-lemma-flat-ses}，每个 $K_d$ 都在 $R$ 上平坦。再次由 Algebra, Lemma
\ref{algebra-lemma-finite-flat-local}，得出每个 $K_d$ 有限自由。

\medskip\noindent
设 $\mathfrak m$ 为 $R$ 的极大理想。由于 $M$ 在 $R$ 上平坦，上述短正合列与
$\kappa = \kappa(\mathfrak m)$ 张量后仍是短正合列。由于环
$S \otimes_R \kappa$ 是 Noether 环，存在齐次元
$k_1, \ldots, k_t \in K$，使其像 $\overline{k}_j$ 在
$S \otimes_R \kappa$ 上生成 $K \otimes_R \kappa$。设 $\deg(k_j) = e_j$。
因此对任意 $d$，映射
$$
\bigoplus\nolimits_{j = 1, \ldots, t} S_{d - e_j}
\longrightarrow
K_d
$$
与 $\kappa$ 张量后成为满射。由 Nakayama 引理
（Algebra, Lemma \ref{algebra-lemma-NAK}），该映射在 $R$ 上是满射。
因此 $K$ 由 $k_1, \ldots, k_t$ 在 $S$ 上生成，于是得证。
\end{proof}

\begin{lemma}
\label{more-algebra-lemma-flat-graded-finite-type-finite-presentation}
设 $R$ 为环。设 $S = \bigoplus_{n \geq 0} S_n$ 为分次 $R$-代数，
$M = \bigoplus_{d \in \mathbf{Z}} M_d$ 为分次 $S$-模。假设 $S$ 作为
$R$-代数有限生成，$S_0$ 是有限 $R$-代数，且存在有限多个素理想
$\mathfrak p_j$、$i = 1, \ldots, m$，使
$R \to \prod R_{\mathfrak p_j}$ 是单射。
\begin{enumerate}
\item 若 $S$ 在 $R$ 上平坦，则 $S$ 是有限表示 $R$-代数。
\item 若 $M$ 作为 $R$-模平坦，作为 $S$-模有限，则 $M$ 作为
$S$-模有限表示。
\end{enumerate}
\end{lemma}

\begin{proof}
由于 $S$ 作为 $R$-代数有限生成，它作为 $S_0$-代数也有限生成；
设由次数为 $d_1, \ldots, d_n > 0$ 的齐次元
$t_1, \ldots, t_n \in S$ 生成。令 $P = R[x_1, \ldots, x_n]$，并令
$\deg(x_i) = d_i$。由于 $S_0$ 是有限 $R$-模，环映射
$P \to S$、$x_i \to t_i$ 是有限的。为证明 (1)，只须证明 $S$ 是有限表示
$P$-模。为证明 (2)，只须证明 $M$ 是有限表示 $P$-模。因此，
只须证明如下情形：$S = P$ 是分次多项式环，$M$ 是有限 $S$-模且在
$R$ 上平坦，则 $M$ 作为 $S$-模有限表示。由引理
\ref{more-algebra-lemma-flat-graded-finite-type-finite-presentation-module}，对 $R$ 的每个素理想
$\mathfrak p$，$M_{\mathfrak p}$ 都是有限表示 $S_{\mathfrak p}$-模。
所以结果由引理 \ref{more-algebra-lemma-flat-finite-type-finite-presentation-local-module} 得出。
\end{proof}

\begin{remark}
\label{more-algebra-remark-when-does-condition-hold}
设 $R$ 为环。$R$ 何时满足引理
\ref{more-algebra-lemma-flat-finite-type-finite-presentation-local-module}、
\ref{more-algebra-lemma-flat-finite-type-finite-presentation-local} 与
\ref{more-algebra-lemma-flat-graded-finite-type-finite-presentation} 中提及的条件？
下列情形中该条件成立：
\begin{enumerate}
\item $R$ 是局部环；
\item $R$ 是 Noether 环；
\item $R$ 是整环；
\item $R$ 是只有有限多个极小素理想的既约环；或
\item $R$ 只有有限多个弱伴随素理想；见 Algebra, Lemma
\ref{algebra-lemma-zero-at-weakly-ass-zero}。
\end{enumerate}
因此这些引理在上述所有情形中都成立。
\end{remark}

\noindent
下一引理将在 More on Flatness, Proposition
\ref{flat-proposition-flat-finite-type-finite-presentation-domain} 中得到改进。

\begin{lemma}
\label{more-algebra-lemma-flat-finite-type-valuation-ring-finite-presentation}
\begin{reference}
\cite[Theorem 3]{Nagata-Finitely}
\end{reference}
设 $A$ 为赋值环。设 $A \to B$ 为有限型环映射，$M$ 为有限 $B$-模。
\begin{enumerate}
\item 若 $B$ 在 $A$ 上平坦，则 $B$ 是有限表示 $A$-代数。
\item 若 $M$ 作为 $A$-模平坦，则 $M$ 作为 $B$-模有限表示。
\end{enumerate}
\end{lemma}

\begin{proof}
我们将使用这一事实：$A$-模平坦当且仅当它无挠；见引理
\ref{more-algebra-lemma-valuation-ring-torsion-free-flat}。由 Algebra, Lemma
\ref{algebra-lemma-homogenize}，可找到分次 $A$-代数 $S$，其中 $S_0 = A$，
且它由有限多个次数为 $1$ 的元生成；还可找到一个元
$f \in S_1$ 和一个有限分次 $S$-模 $N$，使
$B \cong S_{(f)}$ 且 $M \cong N_{(f)}$。若 $M$ 无挠，则用 $N/N_{tors}$
作此替换后，可以取 $N$ 无挠；见引理 \ref{more-algebra-lemma-torsion}。
同样，若 $B$ 无挠，则用 $S/S_{tors}$ 替换后，可以取 $S$ 无挠。
因此在情形 (1) 中，可以应用引理
\ref{more-algebra-lemma-flat-graded-finite-type-finite-presentation}，得到 $S$ 是有限表示
$A$-代数，这推出 $B = S_{(f)}$ 是有限表示 $A$-代数。
为证明 (2)，可以先用分次多项式环替换 $S$，然后应用引理
\ref{more-algebra-lemma-flat-graded-finite-type-finite-presentation-module} 得到结论。
\end{proof}

\begin{lemma}
\label{more-algebra-lemma-valuation-ring-flat-essentially-finite-type}
设 $A$ 为赋值环。设 $A \to B$ 为本质有限型的局部同态，$M$ 为有限
$B$-模。
\begin{enumerate}
\item 若 $B$ 在 $A$ 上平坦，则 $B$ 在 $A$ 上本质有限表示。
\item 若 $M$ 作为 $A$-模平坦，则 $M$ 作为 $B$-模有限表示。
\end{enumerate}
\end{lemma}

\begin{proof}
由假设，可以把 $B$ 写成多项式代数
$P = A[x_1, \ldots, x_n]$ 在素理想 $\mathfrak q$ 处局部化的一个商。
在情形 (1) 中，我们把 $M = B$ 视为 $P_\mathfrak q$ 上的有限模；
在情形 (2) 中，我们把 $M$ 视为 $P_\mathfrak q$ 上的有限模。
在两种情形中，都要证明这是有限表示 $P_\mathfrak q$-模；
对情形 (2)，见 Algebra, Lemma
\ref{algebra-lemma-finitely-presented-over-subring}。

\medskip\noindent
选取表示 $0 \to K \to P_\mathfrak q^{\oplus r} \to M \to 0$；
这是可能的，因为 $M$ 在 $P_\mathfrak q$ 上有限。
令 $L = P^{\oplus r} \cap K$。则 $K = L_\mathfrak q$；见 Algebra, Lemma
\ref{algebra-lemma-submodule-localization}。此时 $N = P^{\oplus r}/L$ 是 $M$ 的子模，
从而由引理 \ref{more-algebra-lemma-valuation-ring-torsion-free-flat} 它平坦。
又因为 $N$ 是有限 $P$-模，由引理
\ref{more-algebra-lemma-flat-finite-type-valuation-ring-finite-presentation}，$N$ 作为 $P$-模有限表示。
由局部化的正合性（Algebra, Proposition
\ref{algebra-proposition-localization-exact}），可知 $N_\mathfrak q = M$，于是得证。
\end{proof}











\section{爆破与平坦性}
\label{more-algebra-section-blowup-flat}

\noindent
本节开始讨论如下形式的结果：“经过爆破后，严格变换变为平坦”。
这种类型的更多结果见 Divisors, Section
\ref{divisors-section-blowup-flat} 与 More on Flatness, Section
\ref{flat-section-blowup-flat}。

\begin{definition}
\label{more-algebra-definition-strict-transform}
设 $R$ 为环，$I \subset R$ 为理想，$a \in I$。设
$R[\frac{I}{a}]$ 为仿射爆破代数；见 Algebra, Definition
\ref{algebra-definition-blow-up}。设 $M$ 为 $R$-模。称 $R[\frac{I}{a}]$-模
$$
M' = \left(M \otimes_R R[\textstyle{\frac{I}{a}}]\right)/a\text{-幂挠子}
$$
为 {\it $M$ 沿 $R \to R[\frac{I}{a}]$ 的严格变换}。
\end{definition}

\noindent
下一结果是通过爆破实现平坦化的一个非常弱的版本，但有时已经很有用。

\begin{lemma}
\label{more-algebra-lemma-flatten-on-affine-blowup}
设 $(R, \mathfrak m)$ 为分式域为 $K$ 的局部整环。设 $S$ 为有限型
$R$-代数，$M$ 为有限 $S$-模。对每个支配 $R$ 的赋值环
$A \subset K$，存在理想 $I \subset \mathfrak m$ 与非零元 $a \in I$，使得：
\begin{enumerate}
\item $I$ 有限生成；
\item $A$ 在 $R[\frac{I}{a}]$ 上有中心；
\item $R \to R[\frac{I}{a}]$ 在 $\mathfrak m$ 处的纤维环不为零；以及
\item $S$ 沿 $R \to R[\frac{I}{a}]$ 的严格变换 $S_{I, a}$ 在 $R$ 上平坦且有限表示，
并且 $M$ 沿 $R \to R[\frac{I}{a}]$ 的严格变换 $M_{I, a}$ 在 $R$ 上平坦，
并在 $S_{I, a}$ 上有限表示。
\end{enumerate}
\end{lemma}

\begin{proof}
写成 $S = R[x_1, \ldots, x_n]/J$，并把 $N = S \oplus M$ 视为
$P = R[x_1, \ldots, x_n]$ 上的模。若能在 $S$ 是 $R$ 上多项式代数的情形证明本引理，
则可找到满足 (1)、(2)、(3) 的 $I, a$，使 $N$ 沿
$R \to R[\frac{I}{a}]$ 的严格变换 $N_{I, a}$ 在 $R$ 上平坦，并且作为
$P$ 的严格变换 $P_{I, a}$ 上的模有限表示。因为
$P_{I, a} = R[\frac{I}{a}][x_1, \ldots, x_n]$（略去一个小细节），
可知直和项 $S_{I, a} \subset N_{I, a}$ 在 $R$ 上平坦，并且作为
$R[\frac{I}{a}][x_1, \ldots, x_n]$-模有限表示。因此 $S_{I, a}$ 是有限表示
$R[\frac{I}{a}]$-代数。此外，直和项 $M_{I, a} \subset N_{I, a}$ 在 $R$ 上平坦，
并且作为 $P_{I, a}$-模有限表示，因而作为 $S_{I, a}$-模也有限表示；见
Algebra, Lemma \ref{algebra-lemma-finitely-presented-over-subring}。
这将我们化约到下一段讨论的情形。

\medskip\noindent
假设 $S = R[x_1, \ldots, x_n]$。选取表示
$$
0 \to K \to S^{\oplus r} \to M \to 0.
$$
设 $M_A$ 是 $M \otimes_R A$ 模去其挠子模得到的商；见引理
\ref{more-algebra-lemma-torsion}。则 $M_A$ 是 $S_A = A[x_1, \ldots, x_n]$ 上的有限模。
由引理 \ref{more-algebra-lemma-valuation-ring-torsion-free-flat}，$M_A$ 在 $A$ 上平坦。
由引理 \ref{more-algebra-lemma-flat-finite-type-valuation-ring-finite-presentation}，
$M_A$ 有限表示。因此存在有限多个元
$k_1, \ldots, k_t \in S_A^{\oplus r}$，它们作为 $S_A$-模生成表示
$S_A^{\oplus r} \to M_A$ 的核。对任意满足 (1)、(2)、(3) 的选取
$a \in I \subset \mathfrak m$，用 $M_{I, a}$ 表示 $M$ 沿
$R \to R[\frac{I}{a}]$ 的严格变换。它是
$S_{I, a} = R[\frac{I}{a}][x_1, \ldots, x_n]$ 上的有限模。
由 Algebra, Lemma \ref{algebra-lemma-valuation-ring-colimit-affine-blowups}，有
$A = \colim_{I, a} R[\frac{I}{a}]$。这推出
$S_A = \colim S_{I, a}$ 且
$$
\colim M \otimes_R R[\textstyle{\frac{I}{a}}] = M \otimes_R A
$$
选取 $I, a$ 以及提升 $k_1, \ldots, k_t \in S_{I, a}^{\oplus r}$。
由于 $M_A$ 是 $M \otimes_R A$ 模去挠得到的商，可知
$k_1, \ldots, k_t$ 在 $M \otimes_R A$ 中的像被某个非零元
$\alpha \in A$ 零化。用另一对替换 $I, a$（回忆该余极限是滤过的）后，
可以假设对某个非零 $x \in I^n$ 有 $\alpha = x/a^n$。
于是 $x k_1, \ldots, x k_t$ 在 $M \otimes_R A$ 中映到零。
再次用另一对替换 $I, a$ 后，可以假设对某个非零 $x \in R$，
$x k_1, \ldots, x k_t$ 在 $M \otimes_R R[\frac{I}{a}]$ 中映到零。
最后用 $xI, xa$ 替换 $I, a$，可以假设 $k_1, \ldots, k_t$ 在
$M \otimes_R R[\frac{I}{a}]$ 中映到被 $a$ 的某个幂零化的挠元。
对任意这样的对 $(I, a)$，令
$$
M'_{I, a} = S_{I, a}^{\oplus r}/ \sum S_{I, a}k_j.
$$
因为 $M_A = S_A^{\oplus r}/ \sum S_Ak_j$，可知
$M_A = \colim_{I, a} M'_{I, a}$。此时终于可以应用 Algebra, Lemma
\ref{algebra-lemma-flat-finite-presentation-limit-flat} (3)，得出对某个如上的对
$(I, a)$，$M'_{I, a}$ 平坦。该引理先验上不能应用于严格变换组成的系统
$$
M_{I, a} = (M \otimes_R R[\textstyle{\frac{I}{a}}])/a\text{-幂挠子}
$$
因为转移映射可能不满足该引理的假设。但现在，由于平坦推出无挠
（引理 \ref{more-algebra-lemma-flat-torsion-free}），并且 $M_{I, a}$ 是 $M'_{I, a}$ 的商
（因为我们已安排使元 $k_1, \ldots, k_t$ 在 $M_{I, a}$ 中映到零），
所商去的是被 $a$ 的某个幂零化的挠子模，所以对这样的对还有
$M'_{I, a} = M_{I, a}$，从而得证。
\end{proof}

\begin{lemma}
\label{more-algebra-lemma-blowup-fitting-ideal}
设 $R$ 为环，$M$ 为有限 $R$-模。设 $k \geq 0$ 且
$I = \text{Fit}_k(M)$。对每个 $a \in I$，令 $R' = R[\frac{I}{a}]$，则严格变换
$$
M' = (M \otimes_R R')/a\text{-幂挠子}
$$
满足 $\text{Fit}_k(M') = R'$。
\end{lemma}

\begin{proof}
首先注意 $\text{Fit}_k(M \otimes_R R') = IR' = aR'$。第一个等式由引理
\ref{more-algebra-lemma-fitting-ideal-basics} 的 (3) 得出，第二个等式由 Algebra, Lemma
\ref{algebra-lemma-affine-blowup} 得出。由引理
\ref{more-algebra-lemma-principal-fitting-ideal} 与局部化的正合性，可知对 $R'$ 的每个素理想
$\mathfrak p'$，$M'_{\mathfrak p'}$ 可由 $\leq k$ 个元生成。
例如由引理 \ref{more-algebra-lemma-fitting-ideal-generate-locally}，可得
$\text{Fit}_k(M') = R'$。
\end{proof}

\begin{lemma}
\label{more-algebra-lemma-blowup-fitting-ideal-locally-free}
设 $R$ 为环，$M$ 为有限 $R$-模。设 $k \geq 0$ 且
$I = \text{Fit}_k(M)$。假设对每个 $\mathfrak p \not \in V(I)$，
$M_\mathfrak p$ 都是秩为 $k$ 的自由模。则对每个 $a \in I$，令
$R' = R[\frac{I}{a}]$，严格变换
$$
M' = (M \otimes_R R')/a\text{-幂挠子}
$$
都局部自由且秩为 $k$。
\end{lemma}

\begin{proof}
由引理 \ref{more-algebra-lemma-blowup-fitting-ideal}，有
$\text{Fit}_k(M') = R'$。由引理 \ref{more-algebra-lemma-fitting-ideal-finite-locally-free}，
只须证明 $\text{Fit}_{k - 1}(M') = 0$。回忆
$R' \subset R'_a = R_a$；见 Algebra, Lemma \ref{algebra-lemma-affine-blowup}。
因此只须证明 $\text{Fit}_{k - 1}(M')$ 在 $R'_a = R_a$ 中映到零。
由于显然 $(M')_a = M_a$，而由引理 \ref{more-algebra-lemma-fitting-ideal-basics} 的 (3)，
形成 Fitting 理想与基变换可交换，问题归结为证明
$\text{Fit}_{k - 1}(M_a) = 0$。由假设，$M_a$ 是秩为 $k$ 的有限局部自由模
（见 Algebra, Lemma \ref{algebra-lemma-finite-projective}），故此结论由已引用的引理
\ref{more-algebra-lemma-fitting-ideal-finite-locally-free} 得出。
\end{proof}

\begin{lemma}
\label{more-algebra-lemma-blowup-module}
设 $R$ 为环，$M$ 为有限 $R$-模。设 $f \in R$ 为元，使 $M_f$ 是秩为
$r$ 的有限局部自由模。则存在有限生成理想 $I \subset R$，满足
$V(f) = V(I)$，且对所有 $a \in I$，令 $R' = R[\frac{I}{a}]$，
严格变换
$$
M' = (M \otimes_R R')/a\text{-幂挠子}
$$
局部自由且秩为 $r$。
\end{lemma}

\begin{proof}
选取满射 $R^{\oplus n} \to M$。选取有限子模
$K \subset \Ker(R^{\oplus n} \to M)$，使 $R^{\oplus n}/K \to M$ 在逆转 $f$ 后成为同构。
这是可能的，因为例如由 Algebra, Lemma
\ref{algebra-lemma-finite-projective}，$M_f$ 有限表示。
令 $M_1 = R^{\oplus n}/K$，并假设能对 $M_1$ 证明本引理。
设 $I \subset R$ 为相应的理想。则对 $a \in I$，映射
$$
M_1' = (M_1 \otimes_R R')/a\text{-幂挠子}
\longrightarrow
M' = (M \otimes_R R')/a\text{-幂挠子}
$$
是满射。在 $R'$ 中逆转 $a$ 后，它也是同构，因为
$R'_a = R_f$；见 Algebra, Lemma \ref{algebra-lemma-blowup-in-principal}。
但 $a$ 在 $M'_1$ 上是非零因子，因此所显示的映射是同构。
所以只须在 $M$ 是有限表示 $R$-模的情形证明本引理。

\medskip\noindent
假设 $M$ 是有限表示 $R$-模。则
$J = \text{Fit}_r(M) \subset S$ 是有限生成理想。我们断言 $I = fJ$ 可行。

\medskip\noindent
首先验证 $V(f) = V(I)$。包含关系 $V(f) \subset V(I)$ 是显然的。
反之，若 $f \not \in \mathfrak p$，则由引理
\ref{more-algebra-lemma-fitting-ideal-generate-locally}，$\mathfrak p$ 不是 $V(J)$ 的元。
因此 $\mathfrak p \not \in V(fJ) = V(I)$。

\medskip\noindent
设 $a \in I$，并令 $R' = R[\frac{I}{a}]$。可以对某个 $b \in J$ 写成
$a = fb$。由 Algebra, Lemmas \ref{algebra-lemma-affine-blowup} 与
\ref{algebra-lemma-blowup-add-principal}，有 $J R' = b R'$，且 $b$ 是 $R'$ 中的非零因子。
设 $\mathfrak p' \subset R' = R[\frac{I}{a}]$ 为素理想。则
$JR'_{\mathfrak p'}$ 由 $b$ 生成。由引理 \ref{more-algebra-lemma-principal-fitting-ideal}，
$M'_{\mathfrak p'}$ 可由 $r$ 个元生成。由于 $M'$ 有限，存在
$m_1, \ldots, m_r \in M'$ 与 $g \in R'$、$g \not \in \mathfrak p'$，使相应的映射
$(R')^{\oplus r} \to M'$ 在逆转 $g$ 后成为满射。

\medskip\noindent
最后，考虑理想 $J' = \text{Fit}_{k - 1}(M')$。注意 $J' R'_g$ 由
$m_1, \ldots, m_r$ 之间各关系的系数生成（Fitting 理想与基变换的相容性）。
因此只须证明 $J' = 0$；见引理 \ref{more-algebra-lemma-fitting-ideal-finite-locally-free}。
由于 $R'_a = R_f$（Algebra, Lemma \ref{algebra-lemma-blowup-in-principal}），而
$M'_a = M_f$ 自由且秩为 $r$，可知 $J'_a = 0$。因为 $a$ 是 $R'$ 中的非零因子，
所以 $J' = 0$，于是得证。
\end{proof}













\section{完备化与平坦性}
\label{more-algebra-section-completion-flat}

\noindent
本节讨论一个“大”平坦模的完备化何时平坦。

\begin{lemma}
\label{more-algebra-lemma-ui-completion-direct-sum-into-product}
设 $R$ 为环，$I \subset R$ 为理想，$A$ 为集合。假设 $R$ 是 Noether 环，
且对 $I$ 完备。存在从直和的 $I$-进完备化到直积的规范映射
$$
\left(\bigoplus\nolimits_{\alpha \in A} R\right)^\wedge
\longrightarrow
\prod\nolimits_{\alpha \in A} R
$$
，且该映射是普遍单射。
\end{lemma}

\begin{proof}
依定义，左边的元 $x$ 为 $x = (x_n)$，其中
$x_n = (x_{n, \alpha}) \in \bigoplus\nolimits_{\alpha \in A} R/I^n$，并满足
$x_{n, \alpha} = x_{n + 1, \alpha} \bmod I^n$。由于 $R = R^\wedge$，
对任意 $\alpha$ 存在 $y_\alpha \in R$，使
$x_{n, \alpha} = y_\alpha \bmod I^n$。注意，对每个 $n$，只有有限多个
$\alpha$ 使 $x_{n, \alpha}$ 不为零。反之，给定
$(y_\alpha) \in \prod_\alpha R$，若对每个 $n$，只有有限多个 $\alpha$ 使
$y_{\alpha} \bmod I^n$ 不为零，则它定义了左边的一个元。
所以可以把左边的元看作无限“收敛和”
$\sum_\alpha y_\alpha$，其中 $y_\alpha \in R$，且对每个 $n$，
只有有限多个 $y_\alpha$ 模 $I^n$ 不为零。所显示的映射把此元映到直积中的元
$(y_\alpha)$。特别地，该映射是单射。

\medskip\noindent
设 $Q$ 为有限 $R$-模。我们必须证明映射
$$
Q \otimes_R \left(\bigoplus\nolimits_{\alpha \in A} R\right)^\wedge
\longrightarrow
Q \otimes_R \left(\prod\nolimits_{\alpha \in A} R\right)
$$
是单射；见 Algebra, Theorem
\ref{algebra-theorem-universally-exact-criteria}。选取表示
$R^{\oplus k} \to R^{\oplus m} \to Q \to 0$，并用
$q_1, \ldots, q_m \in Q$ 表示 $Q$ 的相应生成元。由 Artin-Rees
（Algebra, Lemma \ref{algebra-lemma-Artin-Rees}），存在常数 $c$，使
$\Im(R^{\oplus k} \to R^{\oplus m}) \cap (I^N)^{\oplus m}
\subset \Im((I^{N - c})^{\oplus k} \to R^{\oplus m})$。考虑图
$$
\xymatrix{
\bigoplus_{l = 1}^k \left(\bigoplus\nolimits_{\alpha \in A} R\right)^\wedge
\ar[r] \ar[d] &
\bigoplus_{j = 1}^m \left(\bigoplus\nolimits_{\alpha \in A} R\right)^\wedge
\ar[r] \ar[d] &
Q \otimes_R \left(\bigoplus\nolimits_{\alpha \in A} R\right)^\wedge
\ar[r] \ar[d] &
0 \\
\bigoplus_{l = 1}^k \left(\prod\nolimits_{\alpha \in A} R\right)
\ar[r] &
\bigoplus_{j = 1}^m \left(\prod\nolimits_{\alpha \in A} R\right)
\ar[r] &
Q \otimes_R \left(\prod\nolimits_{\alpha \in A} R\right)
\ar[r] &
0
}
$$
其两行正合。选取元
$\sum_j \sum_\alpha y_{j, \alpha}$，它属于
$\bigoplus_{j = 1, \ldots, m}
\left(\bigoplus\nolimits_{\alpha \in A} R\right)^\wedge$。
若此元在模
$Q \otimes_R \left(\prod\nolimits_{\alpha \in A} R\right)$ 中映到零，
则特别地，对每个 $\alpha$，有
$\sum_j q_j \otimes y_{j, \alpha} = 0$ 在 $Q$ 中成立。
因此可找到元
$(z_{1, \alpha}, \ldots, z_{k, \alpha}) \in \bigoplus_{l = 1, \ldots, k} R$，
它映到
$(y_{1, \alpha}, \ldots, y_{m, \alpha}) \in \bigoplus_{j = 1, \ldots, m} R$。
此外，若对 $j = 1, \ldots, m$ 有 $y_{j, \alpha} \in I^{N_\alpha}$，
则可以假设对 $l = 1, \ldots, k$ 有
$z_{l, \alpha} \in I^{N_\alpha - c}$。因此和
$\sum_l \sum_\alpha z_{l, \alpha}$ 是“收敛的”，并定义元
$\bigoplus_{l = 1, \ldots, k}
\left(\bigoplus\nolimits_{\alpha \in A} R\right)^\wedge$，
该元映到我们开始所取的元 $\sum_j \sum_\alpha y_{j, \alpha}$。
因此右边的竖直箭头是单射，于是得证。
\end{proof}

\noindent
下一引理也可以由下面的引理 \ref{more-algebra-lemma-limit-flat} 推出。

\begin{lemma}
\label{more-algebra-lemma-completed-direct-sum-flat}
设 $R$ 为环，$I \subset R$ 为理想，$A$ 为集合。假设 $R$ 是 Noether 环。
完备化 $(\bigoplus\nolimits_{\alpha \in A} R)^\wedge$ 是平坦 $R$-模。
\end{lemma}

\begin{proof}
用 $R^\wedge$ 表示 $R$ 关于 $I$ 的完备化。由 Algebra, Lemma
\ref{algebra-lemma-completion-flat}，$R \to R^\wedge$ 平坦，所以只须证明
$(\bigoplus\nolimits_{\alpha \in A} R)^\wedge$ 是平坦 $R^\wedge$-模
（使用 Algebra, Lemma \ref{algebra-lemma-composition-flat}）。因为
$$
(\bigoplus\nolimits_{\alpha \in A} R)^\wedge
=
(\bigoplus\nolimits_{\alpha \in A} R^\wedge)^\wedge
$$
可以用 $R^\wedge$ 替换 $R$，并假设 $R$ 对 $I$ 完备
（见 Algebra, Lemma \ref{algebra-lemma-completion-complete}）。
在这种情形下，引理 \ref{more-algebra-lemma-ui-completion-direct-sum-into-product} 告诉我们，映射
$(\bigoplus\nolimits_{\alpha \in A} R)^\wedge \to \prod_{\alpha \in A} R$
是普遍单射。因此由 Algebra, Lemma \ref{algebra-lemma-ui-flat-domain}，
只须证明 $\prod_{\alpha \in A} R$ 平坦。由 Algebra, Proposition
\ref{algebra-proposition-characterize-coherent}（以及 Algebra, Lemma
\ref{algebra-lemma-Noetherian-coherent}），$\prod_{\alpha \in A} R$ 平坦。
\end{proof}

\begin{lemma}
\label{more-algebra-lemma-tor-strictly-pro-zero}
\begin{reference}
这是 \cite[Lemma 9.9]{quillenhomology}；注意，作者在陈述中遗漏了“严格”一词，
尽管显然原意包含该词。
\end{reference}
设 $A$ 为 Noether 环，$I$ 为 $A$ 的理想，$M$ 为有限 $A$-模。
对每个 $p > 0$，存在 $c > 0$，使对所有 $n \geq c$，映射
$\text{Tor}_p^A(M, A/I^n) \to \text{Tor}_p^A(M, A/I^{n - c})$ 为零。
\end{lemma}

\begin{proof}
先证明 $p = 1$ 的情形。选取短正合列
$0 \to K \to A^{\oplus t} \to M \to 0$。则
$\text{Tor}_1^A(M, A/I^n) = K \cap (I^n)^{\oplus t}/I^nK$。
由 Artin-Rees（Algebra, Lemma \ref{algebra-lemma-Artin-Rees}），
存在常数 $c \geq 0$，使对 $n \geq c$ 有
$K \cap (I^n)^{\oplus t} \subset I^{n - c}K$。这证明了 $p = 1$ 的结论。
对 $p > 1$，有
$\text{Tor}_p^A(M, A/I^n) = \text{Tor}^A_{p - 1}(K, A/I^n)$。
因此本引理由归纳得证。
\end{proof}

\begin{lemma}
\label{more-algebra-lemma-limit-flat}
设 $A$ 为 Noether 环，$I$ 为 $A$ 的理想。设 $(M_n)$ 为 $A$-模的逆系统，且：
\begin{enumerate}
\item $M_n$ 是平坦 $A/I^n$-模；
\item $M_{n + 1} \to M_n$ 是满射。
\end{enumerate}
则 $M = \lim M_n$ 是平坦 $A$-模，并且对每个有限 $A$-模 $Q$，有
$Q \otimes_A M = \lim Q \otimes_A M_n$。
\end{lemma}

\begin{proof}
首先对每个有限 $A$-模 $Q$ 证明
$Q \otimes_A M = \lim Q \otimes_A M_n$。选取由有限自由 $A$-模 $F_i$ 组成的解消
$F_2 \to F_1 \to F_0 \to Q \to 0$。则
$$
F_2 \otimes_A M_n \to F_1 \otimes_A M_n \to F_0 \otimes_A M_n
$$
是一个链复形，其 $0$ 次同调为 $Q \otimes_A M_n$，而其 $1$ 次同调为
$$
\text{Tor}_1^A(Q, M_n) = \text{Tor}_1^A(Q, A/I^n) \otimes_{A/I^n} M_n
$$
，因为 $M_n$ 在 $A/I^n$ 上平坦。由引理
\ref{more-algebra-lemma-tor-strictly-pro-zero}，该系统本质上是值为 $0$ 的常系统。
由 Homology, Lemma \ref{homology-lemma-apply-Mittag-Leffler-again}，得
$\lim Q \otimes_A M_n =
\Coker(\lim F_1 \otimes_A M_n \to \lim F_0 \otimes_A M_n)$。
由于 $F_i$ 有限自由，这等于
$\Coker(F_1 \otimes_A M \to F_0 \otimes_A M) = Q \otimes_A M$。

\medskip\noindent
其次，设 $Q \to Q'$ 是有限 $A$-模之间的单射。我们必须证明
$Q \otimes_A M \to Q' \otimes_A M$ 是单射
（Algebra, Lemma \ref{algebra-lemma-flat}）。由上面的结果，有
$$
\Ker(Q \otimes_A M \to Q' \otimes_A M) =
\Ker(\lim Q \otimes_A M_n \to \lim Q' \otimes_A M_n).
$$
对每个 $n$，有正合列
$$
\text{Tor}_1^A(Q', M_n) \to \text{Tor}_1^A(Q'', M_n) \to
Q \otimes_A M_n \to Q' \otimes_A M_n
$$
其中 $Q'' = \Coker(Q \to Q')$。上面已经看到，这些 Tor 的逆系统本质上是值为
$0$ 的常系统。由 Homology, Lemma
\ref{homology-lemma-apply-Mittag-Leffler-again}，最右边映射的逆极限是单射。
\end{proof}

\begin{lemma}
\label{more-algebra-lemma-flat-after-completion}
设 $R$ 为环，$I \subset R$ 为理想，$M$ 为 $R$-模。假设：
\begin{enumerate}
\item $I$ 有限生成；
\item $R/I$ 是 Noether 环；
\item $M/IM$ 在 $R/I$ 上平坦；
\item $\text{Tor}_1^R(M, R/I) = 0$。
\end{enumerate}
则 $I$-进完备化 $R^\wedge$ 是 Noether 环，且 $M^\wedge$ 在
$R^\wedge$ 上平坦。
\end{lemma}

\begin{proof}
由 Algebra, Lemma \ref{algebra-lemma-what-does-it-mean}，对所有 $n$，
模 $M/I^nM$ 在 $R/I^n$ 上平坦。由 Algebra, Lemma
\ref{algebra-lemma-hathat-finitely-generated}，有：(a) $R^\wedge$ 与 $M^\wedge$
都是 $I$-进完备的；(b) 对所有 $n$，有
$R/I^n = R^\wedge/I^nR^\wedge$。由 Algebra, Lemma
\ref{algebra-lemma-completion-Noetherian}，环 $R^\wedge$ 是 Noether 环。
应用引理 \ref{more-algebra-lemma-limit-flat}，可得
$M^\wedge = \lim M/I^nM$ 作为 $R^\wedge$-模平坦。
\end{proof}









\section{Koszul 复形}
\label{more-algebra-section-koszul}

\noindent
我们如下定义 Koszul 复形。

\begin{definition}
\label{more-algebra-definition-koszul}
设 $R$ 为环，$\varphi : E \to R$ 为 $R$-模映射。与 $\varphi$ 相关联的
{\it Koszul 复形} $K_\bullet(\varphi)$ 是如下定义的交换微分分次代数：
\begin{enumerate}
\item 其底分次代数是外代数
$K_\bullet(\varphi) = \wedge(E)$；
\item 微分 $d : K_\bullet(\varphi) \to K_\bullet(\varphi)$ 是唯一的导子，满足对所有
$e \in E = K_1(\varphi)$ 有 $d(e) = \varphi(e)$。
\end{enumerate}
\end{definition}

\noindent
明确地，若 $e_1 \wedge \ldots \wedge e_n$ 是 $K_\bullet(\varphi)$ 中次数为 $n$
的生成元之一，则
$$
d(e_1 \wedge \ldots \wedge e_n) =
\sum\nolimits_{i = 1, \ldots, n} (-1)^{i + 1}
\varphi(e_i)e_1 \wedge \ldots \wedge \widehat{e_i} \wedge \ldots \wedge e_n.
$$
容易验证，这在张量代数上给出良定义的导子，该导子零化
$e \otimes e$，从而经过外代数分解。

\medskip\noindent
我们常假设 $E$ 是有限自由模，设 $E = R^{\oplus n}$。
在这种情形下，映射 $\varphi$ 由一列元 $f_1, \ldots, f_n \in R$ 给出。

\begin{definition}
\label{more-algebra-definition-koszul-complex}
设 $R$ 为环，$f_1, \ldots, f_r \in R$。{\it $f_1, \ldots, f_r$ 上的 Koszul 复形}
是与映射 $(f_1, \ldots, f_r) : R^{\oplus r} \to R$ 相关联的 Koszul 复形。
记为 $K_\bullet(f_\bullet)$、$K_\bullet(f_1, \ldots, f_r)$、
$K_\bullet(R, f_1, \ldots, f_r)$ 或 $K_\bullet(R, f_\bullet)$。
\end{definition}

\noindent
当然，若 $E$ 有限局部自由，则 $K_\bullet(\varphi)$ 在 $\Spec(R)$ 上局部同构于
Koszul 复形 $K_\bullet(f_1, \ldots, f_r)$。该复形具有许多有趣的形式性质。

\begin{lemma}
\label{more-algebra-lemma-functorial}
设 $\varphi : E \to R$ 与 $\varphi' : E' \to R$ 为 $R$-模映射。设
$\psi : E \to E'$ 为满足 $\varphi' \circ \psi = \varphi$ 的 $R$-模映射。
则 $\psi$ 诱导微分分次代数同态
$K_\bullet(\varphi) \to K_\bullet(\varphi')$。
\end{lemma}

\begin{proof}
这直接由定义得出。
\end{proof}

\begin{lemma}
\label{more-algebra-lemma-change-basis}
设 $f_1, \ldots, f_r \in R$ 为一列元，$(x_{ij})$ 为系数在 $R$ 中的可逆
$r \times r$-矩阵。则复形 $K_\bullet(f_\bullet)$ 与
$$
K_\bullet(\sum x_{1j}f_j, \sum x_{2j}f_j, \ldots, \sum x_{rj}f_j)
$$
同构。
\end{lemma}

\begin{proof}
令 $g_i = \sum x_{ij}f_j$。矩阵 $(x_{ji})$ 给出同构
$x : R^{\oplus r} \to R^{\oplus r}$，使
$(g_1, \ldots, g_r) = (f_1, \ldots, f_r) \circ x$。
因此结论由引理 \ref{more-algebra-lemma-functorial} 中所述 Koszul 复形的函子性得出。
\end{proof}

\begin{lemma}
\label{more-algebra-lemma-homotopy-koszul-abstract}
设 $R$ 为环，$\varphi : E \to R$ 为 $R$-模映射。设 $e \in E$ 在 $R$ 中的像为
$f = \varphi(e)$。则作为 $K_\bullet(\varphi)$ 的自同态，有
$$
f = de + ed
$$
。
\end{lemma}

\begin{proof}
这是因为 $d(ea) = d(e)a - ed(a) = fa - ed(a)$。
\end{proof}

\begin{lemma}
\label{more-algebra-lemma-homotopy-koszul}
设 $R$ 为环，$f_1, \ldots, f_r \in R$ 为一列元。在 $K_\bullet(f_\bullet)$ 上乘以
$f_i$ 的映射零同伦，特别地，上同调模 $H_i(K_\bullet(f_\bullet))$ 被理想
$(f_1, \ldots, f_r)$ 零化。
\end{lemma}

\begin{proof}
这是引理 \ref{more-algebra-lemma-homotopy-koszul-abstract} 的特殊情形。
\end{proof}

\noindent
在 Derived Categories, Section \ref{derived-section-cones} 中，我们定义了上链复形态射的锥。
链复形态射 $f : A_\bullet \to B_\bullet$ 的锥记为 $C(f)_\bullet$。复形
$C(f)_\bullet$ 由下式给出：
$C(f)_n = B_n \oplus A_{n - 1}$，微分为
\begin{equation}
\label{more-algebra-equation-differential-cone}
d_{C(f), n} =
\left(
\begin{matrix}
d_{B, n} & f_{n - 1} \\
0 & -d_{A, n - 1}
\end{matrix}
\right)
\end{equation}
它带有规范复形态射 $i : B_\bullet \to C(f)_\bullet$ 与
$p : C(f)_\bullet \to A_\bullet[-1]$，它们由显然的映射
$B_n \to C(f)_n \to A_{n - 1}$ 诱导。

\begin{lemma}
\label{more-algebra-lemma-cone-koszul-abstract}
设 $R$ 为环，$\varphi : E \to R$ 为 $R$-模映射，$f \in R$。令
$E' = E \oplus R$，并定义 $\varphi' : E' \to R$：在 $E$ 上等于 $\varphi$，
在 $R$ 上等于乘以 $f$。复形 $K_\bullet(\varphi')$ 同构于复形映射
$$
f :
K_\bullet(\varphi)
\longrightarrow
K_\bullet(\varphi).
$$
的锥。
\end{lemma}

\begin{proof}
用 $e_0 \in E'$ 表示元 $1 \in R \subset R \oplus E$。由上面对锥的定义，有
$$
C(f)_n = K_n(\varphi) \oplus K_{n - 1}(\varphi) =
\wedge^n(E) \oplus \wedge^{n - 1}(E) = \wedge^n(E')
$$
其中在最后一个 $=$ 中，把 $(0, e_1 \wedge \ldots \wedge e_{n - 1})$
映到 $\wedge^n(E')$ 中的 $e_0 \wedge e_1 \wedge \ldots \wedge e_{n - 1}$。
计算表明此同构与微分相容。具体地，对第一个直和项中的元，这是清楚的，
因为 $\varphi'|_E = \varphi$，而 $d_{C(f)}$ 限制在第一个直和项上就是
$d_{K_\bullet(\varphi)}$。另一方面，若 $e_1 \wedge \ldots \wedge e_{n - 1}$
在第二个直和项中，则
$$
d_{C(f)}(0, e_1 \wedge \ldots \wedge e_{n - 1}) =
fe_1 \wedge \ldots \wedge e_{n - 1}
- d_{K_\bullet(\varphi)}(e_1 \wedge \ldots \wedge e_{n - 1})
$$
而另一方面
\begin{align*}
& d_{K_\bullet(\varphi')}(0, e_0 \wedge e_1 \wedge \ldots \wedge e_{n - 1}) \\
& =
\sum\nolimits_{i = 0, \ldots, n - 1}
(-1)^i \varphi'(e_i)e_0 \wedge \ldots \wedge \widehat{e_i}
\wedge \ldots \wedge e_{n - 1} \\
& =
fe_1 \wedge \ldots \wedge e_{n - 1} +
\sum\nolimits_{i = 1, \ldots, n - 1}
(-1)^i \varphi(e_i)e_0 \wedge \ldots \wedge \widehat{e_i}
\wedge \ldots \wedge e_{n - 1} \\
& =
fe_1 \wedge \ldots \wedge e_{n - 1} -
e_0 \left(\sum\nolimits_{i = 1, \ldots, n - 1}
(-1)^{i + 1} \varphi(e_i)e_1 \wedge \ldots \wedge \widehat{e_i}
\wedge \ldots \wedge e_{n - 1}\right)
\end{align*}
这正是上一计算结果的像。
\end{proof}

\begin{lemma}
\label{more-algebra-lemma-cone-koszul}
设 $R$ 为环，$f_1, \ldots, f_r$ 为 $R$ 中的一列元。复形
$K_\bullet(f_1, \ldots, f_r)$ 同构于复形映射
$$
f_r :
K_\bullet(f_1, \ldots, f_{r - 1})
\longrightarrow
K_\bullet(f_1, \ldots, f_{r - 1}).
$$
的锥。
\end{lemma}

\begin{proof}
这是引理 \ref{more-algebra-lemma-cone-koszul-abstract} 的特殊情形。
\end{proof}

\begin{lemma}
\label{more-algebra-lemma-cone-squared}
设 $R$ 为环，$A_\bullet$ 为 $R$-模的复形，$f, g \in R$。设
$C(f)_\bullet$ 为 $f : A_\bullet \to A_\bullet$ 的锥。类似定义
$C(g)_\bullet$ 与 $C(fg)_\bullet$。则 $C(fg)_\bullet$ 同伦等价于某个映射
$$
C(f)_\bullet[1] \longrightarrow C(g)_\bullet
$$
的锥。
\end{lemma}

\begin{proof}
先在 $A_\bullet$ 是仅在次数 $0$ 放置 $R$ 所得的复形时证明此结论。
在这种情形下，复形 $C(f)_\bullet$ 是
$$
\ldots \to 0 \to R \xrightarrow{f} R \to 0 \to \ldots
$$
其中 $R$ 位于（同调）次数 $1$ 与 $0$。我们使用的复形映射是
$$
\xymatrix{
0 \ar[r] \ar[d] &
0 \ar[r] \ar[d] &
R \ar[r]^f \ar[d]^1 &
R \ar[r] \ar[d] & 0 \ar[d] \\
0 \ar[r] & R \ar[r]^g & R \ar[r] & 0 \ar[r] & 0
}
$$
它的锥是由位于次数 $1$ 与 $0$ 的 $R^{\oplus 2}$ 组成的链复形，
其微分（\ref{more-algebra-equation-differential-cone}）为
$$
\left(
\begin{matrix}
g & 1 \\
0 & -f
\end{matrix}
\right) :
R^{\oplus 2} \longrightarrow R^{\oplus 2}
$$
为证明此链复形同伦等价于 $C(fg)_\bullet$，也就是等价于
$R \xrightarrow{fg} R$，考虑复形映射
$$
\xymatrix{
R \ar[d]_{(1, -g)} \ar[r]_{fg} & R \ar[d]^{(0, 1)} \\
R^{\oplus 2} \ar[r] & R^{\oplus 2}
}
\quad\quad
\xymatrix{
R^{\oplus 2} \ar[d]_{(1, 0)} \ar[r] & R^{\oplus 2} \ar[d]^{(f, 1)} \\
R \ar[r]^{fg} & R
}
$$
其记号的含义是显然的。这两个映射在一个方向上的复合是
$C(fg)_\bullet$ 上的恒等映射，但在另一方向上的复合不是恒等映射。
我们略去写出所需同伦的过程。

\medskip\noindent
为说明一般情形下结果成立，我们使用复形范畴上的函子
$K_\bullet \mapsto \text{Tot}(A_\bullet \otimes_R K_\bullet)$，它与同伦及锥相容。
然后把 $C(f)_\bullet$ 与 $C(g)_\bullet$ 写成双复形
$$
(R \xrightarrow{f} R) \otimes_R A_\bullet
\quad\text{且}\quad
(R \xrightarrow{g} R) \otimes_R A_\bullet
$$
的总复形，并由此从上面讨论的特殊情形推出结论。略去一些细节。
\end{proof}

\begin{lemma}
\label{more-algebra-lemma-koszul-mult-abstract}
设 $R$ 为环，$\varphi : E \to R$ 为 $R$-模映射，$f, g \in R$。令
$E' = E \oplus R$，并定义 $\varphi'_f, \varphi'_g, \varphi'_{fg} : E' \to R$：
在 $E$ 上都等于 $\varphi$，在 $R$ 上分别等于乘以 $f, g, fg$。
复形 $K_\bullet(\varphi'_{fg})$ 同伦等价于某个复形映射
$$
K_\bullet(\varphi'_f)[1]
\longrightarrow
K_\bullet(\varphi'_g).
$$
的锥。
\end{lemma}

\begin{proof}
由引理 \ref{more-algebra-lemma-cone-koszul-abstract}，复形 $K_\bullet(\varphi'_f)$
同构于在 $K_\bullet(\varphi)$ 上乘以 $f$ 的映射的锥；另外两种情形同理。
所以本引理由引理 \ref{more-algebra-lemma-cone-squared} 得出。
\end{proof}

\begin{lemma}
\label{more-algebra-lemma-koszul-mult}
设 $R$ 为环，$f_1, \ldots, f_{r - 1}$ 为 $R$ 中的一列元，$f, g \in R$。
复形 $K_\bullet(f_1, \ldots, f_{r - 1}, fg)$ 同伦等价于某个复形映射
$$
K_\bullet(f_1, \ldots, f_{r - 1}, f)[1]
\longrightarrow
K_\bullet(f_1, \ldots, f_{r - 1}, g)
$$
的锥。
\end{lemma}

\begin{proof}
这是引理 \ref{more-algebra-lemma-koszul-mult-abstract} 的特殊情形。
\end{proof}

\begin{lemma}
\label{more-algebra-lemma-join-sequences-koszul-complex}
设 $R$ 为环，$f_1, \ldots, f_r$、$g_1, \ldots, g_s$ 为 $R$ 中的元。
则存在 Koszul 复形的同构
$$
K_\bullet(R, f_1, \ldots, f_r, g_1, \ldots, g_s) =
\text{Tot}(K_\bullet(R, f_1, \ldots, f_r) \otimes_R
K_\bullet(R, g_1, \ldots, g_s)).
$$
\end{lemma}

\begin{proof}
略。提示：若 $K_\bullet(R, f_1, \ldots, f_r)$ 作为微分分次代数由
$x_1, \ldots, x_r$ 生成，且 $\text{d}(x_i) = f_i$，而
$K_\bullet(R, g_1, \ldots, g_s)$ 作为微分分次代数由
$y_1, \ldots, y_s$ 生成，且 $\text{d}(y_j) = g_j$，则可以把
$K_\bullet(R, f_1, \ldots, f_r, g_1, \ldots, g_s)$ 看作由元列
$x_1, \ldots, x_r, y_1, \ldots, y_s$ 生成的微分分次代数，其中
$\text{d}(x_i) = f_i$ 且 $\text{d}(y_j) = g_j$。
\end{proof}






\section{扩张交错 {\v C}ech 复形}
\label{more-algebra-section-alternating-cech}

\noindent
设 $R$ 为环，$f_1, \ldots, f_r \in R$。$R$ 的扩张交错 {\v C}ech 复形是上链复形
$$
R \to \bigoplus\nolimits_{i_0} R_{f_{i_0}} \to
\bigoplus\nolimits_{i_0 < i_1} R_{f_{i_0}f_{i_1}} \to
\ldots \to R_{f_1\ldots f_r}
$$
其中 $R$ 位于次数 $0$，项 $\bigoplus_{i_0} R_{f_{i_0}}$ 位于次数 $1$，依此类推。
各映射定义如下：
\begin{enumerate}
\item 映射 $R \to \bigoplus\nolimits_{i_0} R_{f_{i_0}}$ 由规范映射
$R \to R_{f_{i_0}}$ 给出。
\item 给定 $1 \leq i_0 < \ldots < i_{p + 1} \leq r$ 与
$0 \leq j \leq p + 1$，有规范局部化映射
$$
R_{f_{i_0} \ldots \hat f_{i_j} \ldots f_{i_{p + 1}}} \to
R_{f_{i_0} \ldots f_{i_{p + 1}}}
$$
\item 微分使用 (2) 中的规范映射，其符号为 $(-1)^j$。
\end{enumerate}
若 $M$ 是任意 $R$-模，则 $M$ 的扩张交错 {\v C}ech 复形是按同样方式构造的上链复形
$$
M \to \bigoplus\nolimits_{i_0} M_{f_{i_0}} \to
\bigoplus\nolimits_{i_0 < i_1} M_{f_{i_0}f_{i_1}} \to
\ldots \to M_{f_1\ldots f_r}
$$
其中像上面一样，$M$ 位于次数 $0$。

\begin{lemma}
\label{more-algebra-lemma-extended-alternating-is-complex}
上面定义的扩张交错 {\v C}ech 复形是 $R$-模的复形。
\end{lemma}

\begin{proof}
略。
\end{proof}

\begin{lemma}
\label{more-algebra-lemma-extended-alternating-form-module}
设 $R$ 为环，$f_1, \ldots, f_r \in R$，$M$ 为 $R$-模。$M$ 的扩张交错
{\v C}ech 复形是 $M$ 与 $R$ 的扩张交错 {\v C}ech 复形在 $R$ 上的张量积。
\end{lemma}

\begin{proof}
略。
\end{proof}

\begin{lemma}
\label{more-algebra-lemma-extended-alternating-base-change}
设 $R$ 为环，$f_1, \ldots, f_r \in R$，$M$ 为 $R$-模。设 $R \to S$ 为环映射，
用 $g_1, \ldots, g_r \in S$ 表示 $f_1, \ldots, f_r$ 的像，并令
$N = M \otimes_R S$。使用 $S$、$g_1, \ldots, g_r$ 与 $N$ 构造的扩张交错
{\v C}ech 复形，是 $M$ 的扩张交错 {\v C}ech 复形与 $S$ 在 $R$ 上的张量积。
\end{lemma}

\begin{proof}
略。
\end{proof}

\begin{lemma}
\label{more-algebra-lemma-extended-alternating-homotopy-zero}
设 $R$ 为环，$f_1, \ldots, f_r \in R$，$M$ 为 $R$-模。若存在
$i \in \{1, \ldots, r\}$ 使 $f_i$ 为单位元，则 $M$ 的扩张交错 {\v C}ech 复形与
$0$ 同伦等价。
\end{lemma}

\begin{proof}
我们将使用如下记号：$M$ 的扩张交错 {\v C}ech 复形中次数为
$p + 1$ 的上链 $x$ 为 $x = (x_{i_0 \ldots i_p})$，其中
$x_{i_0 \ldots i_p}$ 属于 $M_{f_{i_0} \ldots f_{i_p}}$。用此记号，有
$$
d(x)_{i_0 \ldots i_{p + 1}} =
\sum\nolimits_j (-1)^j x_{i_0 \ldots \hat i_j \ldots i_{p + 1}}
$$
作为同伦，使用映射
$$
h : \text{cochains of degree }p + 2 \to \text{cochains of degree }p + 1
$$
它由下列规则给出：
$$
h(x)_{i_0 \ldots i_p} = 0 \text{ 若 } i \in \{i_0, \ldots, i_p\}
\text{ 且 }
h(x)_{i_0 \ldots i_p} = 
(-1)^j x_{i_0 \ldots i_j i i_{j + 1} \ldots i_p} \text{ if not}
$$
在第二种情形中，$j$ 是满足 $i_j < i < i_{j + 1}$ 的唯一指标；
另外，因为 $f_i$ 是单位元，有等式
$$
M_{f_{i_0} \ldots f_{i_p}} =
M_{f_{i_0} \ldots f_{i_j} f_i f_{i_{j + 1}} \ldots f_{i_p}}
$$
借助此等式，可以把
$(-1)^j x_{i_0 \ldots i_j i i_{j + 1} \ldots i_p}$
理解为 $M_{f_{i_0} \ldots f_{i_p}}$ 的元。我们将通过计算证明
$d h + h d$ 等于恒等映射的负，从而完成证明。
为此，固定次数为 $p + 1$ 的上链 $x$，并令
$1 \leq i_0 < \ldots < i_p \leq r$。

\medskip\noindent
情形 I：$i \in \{i_0, \ldots, i_p\}$。设 $i = i_t$。则
$h(d(x))_{i_0 \ldots i_p} = 0$。另一方面，有
$$
d(h(x))_{i_0 \ldots i_p} =
\sum (-1)^j h(x)_{i_0 \ldots \hat i_j \ldots i_p} =
(-1)^t h(x)_{i_0 \ldots \hat i \ldots i_p} =
(-1)^t (-1)^{t - 1} x_{i_0 \ldots i_p}
$$
因此如所需，$(dh + hd)(x)_{i_0 \ldots i_p} = -x_{i_0 \ldots i_p}$。

\medskip\noindent
情形 II：$i \not \in \{i_0, \ldots, i_p\}$。设 $j$ 满足 $i_j < i < i_{j + 1}$。则
\begin{align*}
h(d(x))_{i_0 \ldots i_p}
& =
(-1)^j d(x)_{i_0 \ldots i_j i i_{j + 1} \ldots i_p} \\
& =
\sum\nolimits_{j' \leq j} (-1)^{j + j'}
x_{i_0 \ldots \hat i_{j'} \ldots i_j i i_{j + 1} \ldots i_p} -
x_{i_0 \ldots i_p} \\
&
+ \sum\nolimits_{j' > j} (-1)^{j + j' + 1}
x_{i_0 \ldots i_j i i_{j + 1} \ldots \hat i_{j'} \ldots i_p}
\end{align*}
另一方面，
\begin{align*}
d(h(x))_{i_0 \ldots i_p}
& =
\sum\nolimits_{j'} (-1)^{j'} h(x)_{i_0 \ldots \hat i_{j'} \ldots i_p} \\
& =
\sum\nolimits_{j' \leq j} (-1)^{j' + j - 1}
x_{i_0 \ldots \hat i_{j'} \ldots i_j i i_{j + 1} \ldots i_p} \\
& +
\sum\nolimits_{j' > j} (-1)^{j' + j}
x_{i_0 \ldots i_j i i_{j + 1} \ldots \hat i_{j'} \ldots i_p}
\end{align*}
将它们相加，如所需，得
$(dh + hd)(x)_{i_0 \ldots i_p} = - x_{i_0 \ldots i_p}$。
\end{proof}

\begin{lemma}
\label{more-algebra-lemma-extended-alternating-torsion}
设 $R$ 为环，$f_1, \ldots, f_r \in R$，$M$ 为 $R$-模。设 $H^q$ 为 $M$ 的扩张交错
{\v C}ech 复形的第 $q$ 个上同调模。则
\begin{enumerate}
\item 若 $q \not \in [0, r]$，则 $H^q = 0$；
\item 对 $x \in H^i$，存在 $n \geq 1$，使对 $i = 1, \ldots, r$ 有 $f_i^n x = 0$；
\item $H^q$ 的支撑包含于 $V(f_1, \ldots, f_r)$；
\item 若存在 $f \in (f_1, \ldots, f_r)$，它在 $M$ 上可逆地作用，则 $H^q = 0$。
\end{enumerate}
\end{lemma}

\begin{proof}
(1) 由扩张交错 {\v C}ech 复形在次数 $< 0$ 与 $> r$ 上为零这一事实得出。
为证明 (2)，只须证明对每个 $i$，存在 $n \geq 1$ 使 $f_i^n x = 0$。
为此，只须证明 $(H^q)_{f_i} = 0$。因为局部化正合，
$(H^q)_{f_i}$ 是 $M$ 的扩张交错复形在 $f_i$ 处局部化的第 $q$ 个上同调模。
由引理 \ref{more-algebra-lemma-extended-alternating-base-change}，该局部化是 $M_{f_i}$ 在 $R_{f_i}$ 上关于
$f_1, \ldots, f_r$ 在 $R_{f_i}$ 中的像所得的扩张交错 {\v C}ech 复形。
因此问题化约为：若 $f_i$ 可逆，则 $H^q$ 为零；这由引理
\ref{more-algebra-lemma-extended-alternating-homotopy-zero} 得出。(3) 由刚才证明的对所有 $i$
都有 $(H^q)_{f_i} = 0$ 得出。为证明 (4)，注意此时 $f$ 在 $H^q$ 上可逆地作用，
而由 (3)，$H^q$ 支撑在 $V(f)$ 上。这强制 $H^q$ 为零（略去一个小细节）。
\end{proof}

\begin{lemma}
\label{more-algebra-lemma-extended-alternating-Cech-is-colimit-koszul}
设 $R$ 为环，$f_1, \ldots, f_r \in R$。扩张交错 {\v C}ech 复形
$$
R \to \bigoplus\nolimits_{i_0} R_{f_{i_0}} \to
\bigoplus\nolimits_{i_0 < i_1} R_{f_{i_0}f_{i_1}} \to
\ldots \to R_{f_1\ldots f_r}
$$
是 Koszul 复形 $K(R, f_1^n, \ldots, f_r^n)$ 的余极限；精确陈述见证明。
\end{lemma}

\begin{proof}
我们强烈建议读者自行证明此结论。用 $K(R, f_1^n, \ldots, f_r^n)$
表示定义 \ref{more-algebra-definition-koszul-complex} 的 Koszul 复形，并将其视为位于次数
$0, \ldots, r$ 的上链复形。因此有
$$
K(R, f_1^n, \ldots, f_r^n) :
0 \to \wedge^r(R^{\oplus r}) \to
\wedge^{r - 1}(R^{\oplus r}) \to \ldots \to
R^{\oplus r} \to R \to 0
$$
其中项 $\wedge^r(R^{\oplus r})$ 位于次数 $0$。设
$e^n_1, \ldots, e^n_r$ 为 $R^{\oplus r}$ 的标准基。则元
$e^n_{j_1} \wedge \ldots \wedge e^n_{j_{r - p}}$，其中
$1 \leq j_1 < \ldots < j_{r - p} \leq r$，构成 Koszul 复形次数 $p$ 项的一组基。
此外，由 Section \ref{more-algebra-section-koszul} 中对 Koszul 复形的构造，注意
$$
d(e^n_{j_1} \wedge \ldots \wedge e^n_{j_{r - p}}) =
\sum (-1)^{a + 1} f_{j_a}^n e^n_{j_1} \wedge \ldots \wedge \hat e^n_{j_a} \wedge
\ldots \wedge e^n_{j_{r - p}}
$$
我们系统的转移映射
$$
K(R, f_1^n, \ldots, f_r^n) \to K(R, f_1^{n + 1}, \ldots, f_r^{n + 1})
$$
由如下规则给出：
$$
e^n_{j_1} \wedge \ldots \wedge e^n_{j_{r - p}}
\longmapsto
f_{i_0} \ldots f_{i_{p - 1}}
e^{n + 1}_{j_1} \wedge \ldots \wedge e^{n + 1}_{j_{r - p}}
$$
其中指标 $1 \leq i_0 < \ldots < i_{p - 1} \leq r$ 满足
$\{1, \ldots r\} =
\{i_0, \ldots, i_{p - 1}\} \amalg \{j_1, \ldots, j_{r - p}\}$。
我们略去说明此映射与微分相容的简短计算。注意，转移映射在次数
$0$ 上总是 $1$，在次数 $r$ 上等于 $f_1 \ldots f_r$。

\medskip\noindent
用 $K^p(R, f_1^n, \ldots, f_r^n)$ 表示 Koszul 复形的次数 $p$ 项。
注意，对任意 $f \in R$，有
$$
R_f = \colim (R \xrightarrow{f} R \xrightarrow{f} R \to \ldots )
$$
因此在次数 $p$ 上得到
$$
\colim K^p(R, f_1^n, \ldots f_r^n) =
\bigoplus\nolimits_{1 \leq i_0 < \ldots < i_{p - 1} \leq r}
R_{f_{i_0} \ldots f_{i_{p - 1}}}
$$
这里，上述 Koszul 复形的元
$e^n_{j_1} \wedge \ldots \wedge e^n_{j_{r - p}}$ 在余极限中映到直和项
$R_{f_{i_0} \ldots f_{i_{p - 1}}}$ 中的元
$(f_{i_0} \ldots f_{i_{p - 1}})^{-n}$，其中指标选取为满足
$\{1, \ldots r\} = \{i_0, \ldots, i_{p - 1}\}
\amalg \{j_1, \ldots, j_{r - p}\}$。因此此复形上的微分由下式给出：
$$
d(1\text{ 于 }R_{f_{i_0} \ldots f_{i_{p - 1}}}) =
\sum\nolimits_{i \not \in \{i_0, \ldots, i_{p - 1}\}}
(-1)^{i - t}\text{ 于 }
R_{f_{i_0} \ldots f_{i_t} f_i f_{i_{t + 1}} \ldots f_{i_{p - 1}}}
$$
因此，若考虑在次数 $p$ 上由映射
$$
\bigoplus\nolimits_{1 \leq i_0 < \ldots < i_{p - 1} \leq r}
R_{f_{i_0} \ldots f_{i_{p - 1}}}
\longrightarrow
\bigoplus\nolimits_{1 \leq i_0 < \ldots < i_{p - 1} \leq r}
R_{f_{i_0} \ldots f_{i_{p - 1}}}
$$
给出的复形映射，且该映射由规则
$$
1\text{ 于 }R_{f_{i_0} \ldots f_{i_{p - 1}}}
\longmapsto
(-1)^{i_0 + \ldots + i_{p - 1} + p}\text{ 于 }R_{f_{i_0} \ldots f_{i_{p - 1}}}
$$
确定，则得到从 $\colim K(R, f_1^n, \ldots, f_r^n)$ 到本节定义的扩张交错
{\v C}ech 复形的同构。我们略去对符号相符的验证。
\end{proof}








\section{Koszul 正则序列}
\label{more-algebra-section-koszul-regular}

\noindent
在阅读本节前，请先参看 Algebra, Sections
\ref{algebra-section-regular-sequences}、
\ref{algebra-section-quasi-regular} 与
\ref{algebra-section-depth}。

\begin{definition}
\label{more-algebra-definition-koszul-regular-sequence}
设 $R$ 为环，$r \geq 0$，$f_1, \ldots, f_r \in R$ 为一列元。设 $M$ 为 $R$-模。
称序列 $f_1, \ldots, f_r$：
\begin{enumerate}
\item 若对所有 $i \not = 0$ 有
$H_i(K_\bullet(f_1, \ldots, f_r) \otimes_R M) = 0$，则它是
{\it $M$-Koszul 正则的}；
\item 若 $H_1(K_\bullet(f_1, \ldots, f_r) \otimes_R M) = 0$，则它是
{\it $M$-$H_1$ 正则的}；
\item 若对所有 $i \not = 0$ 有 $H_i(K_\bullet(f_1, \ldots, f_r)) = 0$，
则它是{\it Koszul 正则的}；以及
\item 若 $H_1(K_\bullet(f_1, \ldots, f_r)) = 0$，则它是{\it $H_1$ 正则的}。
\end{enumerate}
\end{definition}

\noindent
我们将在引理 \ref{more-algebra-lemma-regular-koszul-regular}、
\ref{more-algebra-lemma-koszul-regular-H1-regular} 与
\ref{more-algebra-lemma-H1-regular-quasi-regular} 中看到，对环 $R$ 的元
$f_1, \ldots, f_r$，有下列推含：
\begin{align*}
f_1, \ldots, f_r\text{ is a regular sequence}
& \Rightarrow f_1, \ldots, f_r\text{ is a Koszul-regular sequence} \\
& \Rightarrow f_1, \ldots, f_r\text{ is an }H_1\text{-regular sequence} \\
& \Rightarrow f_1, \ldots, f_r\text{ is a quasi-regular sequence.}
\end{align*}
一般而言，这些推含均不可逆；但若 $R$ 是 Noether 局部环，且
$f_1, \ldots, f_r \in \mathfrak m_R$，则四个条件全部等价
（引理 \ref{more-algebra-lemma-noetherian-finite-all-equivalent}）。
若 $f = f_1 \in R$ 是长度为 $1$ 的序列，且 $f$ 不是 $R$ 的单位元，
则下列条件显然全部等价：
\begin{enumerate}
\item $f$ 是长度为一的正则序列；
\item $f$ 是长度为一的 Koszul 正则序列；以及
\item $f$ 是长度为一的 $H_1$ 正则序列。
\end{enumerate}
这些条件也显然推出 $f$ 是长度为一的拟正则序列。但确实存在长度为 $1$
而非正则序列的拟正则序列。具体地，令
$$
R = k[x, y_0, y_1, \ldots]/(xy_0, xy_1 - y_0, xy_2 - y_1, \ldots)
$$
并令 $f$ 为 $x$ 在 $R$ 中的像。则 $f$ 是零因子，但
$\bigoplus_{n \geq 0} (f^n)/(f^{n + 1}) \cong k[x]$ 是多项式环。

\begin{lemma}
\label{more-algebra-lemma-regular-koszul-regular}
设 $R$ 为环，$M$ 为 $R$-模，$f_1, \ldots, f_r \in R$，且对
$i = 1, \ldots, r$，乘以 $f_i$ 在 $M/(f_1, \ldots, f_{i - 1})M$ 上是单射。
则 $f_1, \ldots, f_r$ 是 $M$-Koszul 正则的。特别地，$M$-正则序列是
$M$-Koszul 正则的，且任意正则序列都是 Koszul 正则的。
\end{lemma}

\begin{proof}
设 $R$、$M$、$f_1, \ldots, f_r$ 如引理首句所述。若 $r = 1$，则
$f_1$ 是 $M$-Koszul 正则的这一结论是立即的。假设 $r > 1$。
因为 $f_1$ 在 $M$ 上是非零因子，我们得到复形的短正合列：
$$
0 \to K_\bullet(f_2, \ldots, f_r) \otimes M
\xrightarrow{f_1}
K_\bullet(f_2, \ldots, f_r) \otimes M \to
K_\bullet(\overline{f}_2, \ldots, \overline{f}_r) \otimes M/f_1M \to 0
$$
这里 $\overline{f}_i$ 是 $f_i$ 在 $R/(f_1)$ 中的像。由引理
\ref{more-algebra-lemma-cone-koszul}，复形 $K_\bullet(f_1, \ldots, f_r)$ 同构于在
$K_\bullet(f_2, \ldots, f_r)$ 上乘以 $f_1$ 的映射的锥。因此
$K_\bullet(R, f_1, \ldots, f_r) \otimes M$ 同构于第一个映射的锥。
从而 $K_\bullet(\overline{f}_2, \ldots, \overline{f}_r) \otimes M/f_1M$
与 $K_\bullet(f_1, \ldots, f_r) \otimes M$ 拟同构。因为
$R/(f_1)$、$M/f_1M$、$\overline{f}_2, \ldots, \overline{f}_r$ 满足引理的条件，
由归纳法，该复形在正次数上无同调。这完成第一个陈述的证明。
第二个陈述立即由第一个陈述得出。
\end{proof}

\begin{lemma}
\label{more-algebra-lemma-koszul-regular-H1-regular}
$M$-Koszul 正则序列是 $M$-$H_1$ 正则的。
Koszul 正则序列是 $H_1$ 正则的。
\end{lemma}

\begin{proof}
这直接由定义得出。
\end{proof}

\begin{lemma}
\label{more-algebra-lemma-mult-koszul-regular}
设 $f_1, \ldots, f_{r - 1} \in R$ 为一列元，$f, g \in R$，$M$ 为 $R$-模。
\begin{enumerate}
\item 若 $f_1, \ldots, f_{r - 1}, f$ 与 $f_1, \ldots, f_{r - 1}, g$
都是 $M$-$H_1$ 正则的，则 $f_1, \ldots, f_{r - 1}, fg$ 也是
$M$-$H_1$ 正则的。
\item 若 $f_1, \ldots, f_{r - 1}, f$ 与 $f_1, \ldots, f_{r - 1}, g$
都是 $M$-Koszul 正则的，则 $f_1, \ldots, f_{r - 1}, fg$ 也是
$M$-Koszul 正则的。
\end{enumerate}
\end{lemma}

\begin{proof}
由引理 \ref{more-algebra-lemma-koszul-mult}，对所有 $i$，有正合列
{\small
$$
H_i(K_\bullet(f_1, \ldots, f_{r - 1}, f) \otimes M) \to
H_i(K_\bullet(f_1, \ldots, f_{r - 1}, fg) \otimes M) \to
H_i(K_\bullet(f_1, \ldots, f_{r - 1}, g) \otimes M)
$$
}
\end{proof}

\begin{lemma}
\label{more-algebra-lemma-koszul-regular-flat-base-change}
设 $\varphi : R \to S$ 为平坦环映射，$f_1, \ldots, f_r \in R$。设 $M$ 为
$R$-模，并令 $N = M \otimes_R S$。
\begin{enumerate}
\item 若 $R$ 中的 $f_1, \ldots, f_r$ 是 $M$-$H_1$ 正则序列，则
$\varphi(f_1), \ldots, \varphi(f_r)$ 是 $S$ 中的 $N$-$H_1$ 正则序列。
\item 若 $f_1, \ldots, f_r$ 是 $R$ 中的 $M$-Koszul 正则序列，则
$\varphi(f_1), \ldots, \varphi(f_r)$ 是 $S$ 中的 $N$-Koszul 正则序列。
\end{enumerate}
\end{lemma}

\begin{proof}
这是因为
$K_\bullet(f_1, \ldots, f_r) \otimes_R S =
K_\bullet(\varphi(f_1), \ldots, \varphi(f_r))$，从而
$(K_\bullet(f_1, \ldots, f_r) \otimes_R M) \otimes_R S =
K_\bullet(\varphi(f_1), \ldots, \varphi(f_r)) \otimes_S N$。
\end{proof}

\begin{lemma}
\label{more-algebra-lemma-H1-regular-quasi-regular}
$M$-$H_1$ 正则序列是 $M$-拟正则的。
\end{lemma}

\begin{proof}
设 $R$ 为环，$M$ 为 $R$-模。设 $f_1, \ldots, f_r$ 为 $M$-$H_1$ 正则序列。
记 $J = (f_1, \ldots, f_r)$。该假设意味着有正合列
$$
\wedge^2(R^r) \otimes M \to R^{\oplus r} \otimes M \to JM \to 0
$$
其中第一个箭头由
$e_i \wedge e_j \otimes m \mapsto (f_ie_j - f_je_i) \otimes m$ 给出。
将此序列与 $R/J$ 张量，可知
$$
JM/J^2M = (R/J)^{\oplus r} \otimes_R M = (M/JM)^{\oplus r}
$$
是有限自由模。为完成证明，我们必须对每个 $n \geq 2$ 证明：
若
$$
\xi = \sum\nolimits_{|I| = n, I = (i_1, \ldots, i_r)}
m_I f_1^{i_1} \ldots f_r^{i_r} \in J^{n + 1}M
$$
则对所有 $I$ 有 $m_I \in JM$。在下一段中，我们对
$I = (0, \ldots, 0, n)$ 证明 $m_I \in JM$；最后一段将一般情形化归到此特殊情形。

\medskip\noindent
设 $I = (0, \ldots, 0, n)$，$\xi$ 如上所述。可以写成
$\xi = m_1 f_1 + \ldots + m_{r - 1}f_{r - 1} + m_I f_r^n$。
由于已假设 $\xi \in J^{n + 1}M$，也可写成
$\xi = \sum_{1 \leq i \leq j \leq r - 1} m_{ij}f_if_j +
\sum_{1 \leq i \leq r - 1}m'_i f_if_r^n + m'' f_r^{n + 1}$。
于是
$$
\begin{matrix}
(m_1 - m_{11}f_1 - m'_1f_r^n)f_1 + \\
(m_2 - m_{12}f_1 - m_{22}f_2 - m'_2f_r^n)f_2 + \\
\ldots + \\
(m_{r - 1} - m_{1 r - 1}f_1 - \ldots - m_{r - 1 r - 1}f_{r - 1}
- m'_{r - 1}f_r^n)f_{r - 1} + \\
(m_I - m'' f_r)f_r^n = 0
\end{matrix}
$$
由引理 \ref{more-algebra-lemma-mult-koszul-regular}，$f_1, \ldots, f_{r - 1}, f_r^n$
是 $M$-$H_1$ 正则的，因此
$m_I - m'' f_r$ 属于子模 $f_1M + \ldots + f_{r - 1}M + f_r^nM$。
从而 $m_I \in f_1M + \ldots + f_rM$。

\medskip\noindent
设 $S = R[x_1, x_2, \ldots, x_r, 1/x_r]$。环映射 $R \to S$ 是忠实平坦的，
因此 $f_1, \ldots, f_r$ 是 $S$ 中的 $M$-$H_1$ 正则序列；见引理
\ref{more-algebra-lemma-koszul-regular-flat-base-change}。由引理 \ref{more-algebra-lemma-change-basis}，可知
$$
g_1 = f_1 - \frac{x_1}{x_r} f_r,
\ \ldots,
\ g_{r - 1} = f_{r - 1} - \frac{x_{r - 1}}{x_r} f_r,
\ g_r = \frac{1}{x_r}f_r
$$
是 $S$ 中的 $M$-$H_1$ 正则序列。最后，注意我们的元 $\xi$ 可重写为
$$
\xi = \sum\nolimits_{|I| = n, I = (i_1, \ldots, i_r)}
m_I (g_1 + x_1 g_r)^{i_1} \ldots (g_{r - 1} + x_{r - 1} g_r)^{i_{r - 1}}
(x_rg_r)^{i_r}
$$
该表达式中 $g_r^n$ 的系数是
$$
\sum m_I x_1^{i_1} \ldots x_r^{i_r}
$$
由上一段讨论的情形，该和属于 $J(M \otimes_R S)$。因为单项式
$x_1^{i_1} \ldots x_r^{i_r}$ 构成 $S$ 在 $R$ 上的一组 $R$-基的一部分，所以如所需，
对所有 $I$ 有 $m_I \in J$。
\end{proof}

\noindent
对 Noether 局部环上的非零有限模，目前引入的所有各类正则序列彼此等价。

\begin{lemma}
\label{more-algebra-lemma-noetherian-finite-all-equivalent}
设 $(R, \mathfrak m)$ 为 Noether 局部环，$M$ 为非零有限 $R$-模，
$f_1, \ldots, f_r \in \mathfrak m$。下列条件等价：
\begin{enumerate}
\item $f_1, \ldots, f_r$ 是 $M$-正则序列；
\item $f_1, \ldots, f_r$ 是 $M$-Koszul 正则序列；
\item $f_1, \ldots, f_r$ 是 $M$-$H_1$ 正则序列；
\item $f_1, \ldots, f_r$ 是 $M$-拟正则序列。
\end{enumerate}
特别地，序列 $f_1, \ldots, f_r$ 是 $R$ 中的正则序列，当且仅当它是
Koszul 正则序列，当且仅当它是 $H_1$ 正则序列，当且仅当它是拟正则序列。
\end{lemma}

\begin{proof}
推含 (1) $\Rightarrow$ (2) 是引理 \ref{more-algebra-lemma-regular-koszul-regular}。
推含 (2) $\Rightarrow$ (3) 是引理 \ref{more-algebra-lemma-koszul-regular-H1-regular}。
推含 (3) $\Rightarrow$ (4) 是引理 \ref{more-algebra-lemma-H1-regular-quasi-regular}。
推含 (4) $\Rightarrow$ (1) 是 Algebra, Lemma
\ref{algebra-lemma-quasi-regular-regular}。
\end{proof}

\begin{lemma}
\label{more-algebra-lemma-H1-regular-in-quotient}
设 $A$ 为环，$I \subset A$ 为理想。设 $g_1, \ldots, g_m$ 是 $A$ 中的序列，
其在 $A/I$ 中的像是 $H_1$ 正则的。则
$I \cap (g_1, \ldots, g_m) = I(g_1, \ldots, g_m)$。
\end{lemma}

\begin{proof}
考虑复形的正合列
$$
0 \to I \otimes_A K_\bullet(A, g_1, \ldots, g_m)
\to K_\bullet(A, g_1, \ldots, g_m) \to
K_\bullet(A/I, g_1, \ldots, g_m) \to 0
$$
由假设，右边复形的 $H_1 = 0$，所以映射
$$
\Coker(I^{\oplus m} \to I)
\longrightarrow
\Coker(A^{\oplus m} \to A)
$$
是单射。这等价于本引理的断言。
\end{proof}

\begin{lemma}
\label{more-algebra-lemma-conormal-sequence-H1-regular}
设 $A$ 为环，$I \subset J \subset A$ 为理想。假设
$J/I \subset A/I$ 由一个 $H_1$ 正则序列生成。则 $I \cap J^2 = IJ$。
\end{lemma}

\begin{proof}
为证明此结论，选取 $g_1, \ldots, g_m \in J$，使其在 $A/I$ 中的像构成
生成 $J/I$ 的 $H_1$ 正则序列。特别地，$J = I + (g_1, \ldots, g_m)$。
设 $x \in I \cap J^2$。由于 $x \in J^2$，可写成
$$
x =
\sum a_{ij} g_ig_j +
\sum a_j g_j +
a
$$
其中 $a_{ij} \in A$、$a_j \in I$ 且 $a \in I^2$。于是
$\sum a_{ij}g_ig_j \in I \cap (g_1, \ldots, g_m)$，故由引理
\ref{more-algebra-lemma-H1-regular-in-quotient} 可知
$\sum a_{ij}g_ig_j \in I(g_1, \ldots, g_m)$。从而如所需，$x \in IJ$。
\end{proof}

\begin{lemma}
\label{more-algebra-lemma-join-quasi-regular-H1-regular}
设 $A$ 为环，$I$ 是由 $A$ 中的拟正则序列 $f_1, \ldots, f_n$ 生成的理想。
设 $g_1, \ldots, g_m \in A$，且其像
$\overline{g}_1, \ldots, \overline{g}_m$ 在 $A/I$ 中构成 $H_1$ 正则序列。
则 $f_1, \ldots, f_n, g_1, \ldots, g_m$ 是 $A$ 中的拟正则序列。
\end{lemma}

\begin{proof}
我们断言，对每个 $d$，$g_1, \ldots, g_m$ 在 $A/I^d$ 中构成
$H_1$ 正则序列。归纳地假设该结论在 $A/I^{d - 1}$ 中成立。我们有复形的短正合列
$$
0 \to K_\bullet(A, g_\bullet) \otimes_A I^{d - 1}/I^d
\to K_\bullet(A/I^d, g_\bullet) \to
K_\bullet(A/I^{d - 1}, g_\bullet) \to 0
$$
由于 $f_1, \ldots, f_n$ 拟正则，第一个复形是
$K_\bullet(A/I, g_1, \ldots, g_m)$ 的若干副本的直和，因而在次数 $1$ 上无同调。
由归纳假设，最后一个复形在次数 $1$ 上无同调，故中间复形也是如此。
特别地，对每个 $d \geq 1$，序列 $g_1, \ldots, g_m$ 在 $A/I^d$ 中拟正则；
见引理 \ref{more-algebra-lemma-H1-regular-quasi-regular}。
现在可以证明 $f_1, \ldots, f_n, g_1, \ldots, g_m$ 是 $A$ 中的拟正则序列。
事实上，令 $J = (f_1, \ldots, f_n, g_1, \ldots, g_m)$，并设（采用多重指标记号）
$$
\sum\nolimits_{|N| + |M| = d} a_{N, M} f^N g^M \in J^{d + 1}
$$
其中某些 $a_{N, M} \in A$。我们必须证明对所有 $N, M$ 都有
$a_{N, M} \in J$。令 $e \in \{0, 1, \ldots, d\}$。则
$$
\sum\nolimits_{|N| = d - e, \ |M| = e} a_{N, M} f^N g^M \in
(g_1, \ldots, g_m)^{e + 1} + I^{d - e + 1}
$$
因为 $g_1, \ldots, g_m$ 在 $A/I^{d - e + 1}$ 中是拟正则序列，所以对每个满足
$|M| = e$ 的 $M$，可推出
$$
\sum\nolimits_{|N| = d - e} a_{N, M} f^N \in
(g_1, \ldots, g_m) + I^{d - e + 1}
$$
在环 $A/I^{d - e + 1}$ 中将引理 \ref{more-algebra-lemma-H1-regular-in-quotient}
应用于 $I^{d - e}/I^{d - e + 1}$，可知
$\sum_{|N| = d - e} a_{N, M} f^N \in I^{d - e}(g_1, \ldots, g_m)$。
由于 $f_1, \ldots, f_n$ 在 $A$ 中拟正则，这又推出：对每个满足
$|N| = d - e$ 且 $|M| = e$ 的 $N, M$，都有 $a_{N, M} \in J$。
引理得证。
\end{proof}

\begin{lemma}
\label{more-algebra-lemma-join-H1-regular-sequences}
设 $A$ 为环，$I$ 是由 $A$ 中的 $H_1$ 正则序列
$f_1, \ldots, f_n$ 生成的理想。设 $g_1, \ldots, g_m \in A$，且其像
$\overline{g}_1, \ldots, \overline{g}_m$ 在 $A/I$ 中构成 $H_1$ 正则序列。
则 $f_1, \ldots, f_n, g_1, \ldots, g_m$ 是 $A$ 中的 $H_1$ 正则序列。
\end{lemma}

\begin{proof}
我们必须证明 $H_1(A, f_1, \ldots, f_n, g_1, \ldots, g_m) = 0$。
为此考虑交换图
$$
\xymatrix{
\wedge^2(A^{\oplus n + m}) \ar[r] \ar[d] &
A^{\oplus n + m} \ar[r] \ar[d] &
A \ar[r] \ar[d] & 0 \\
\wedge^2(A/I^{\oplus m}) \ar[r] &
A/I^{\oplus m} \ar[r] &
A/I \ar[r] & 0
}
$$
考虑映到 $A$ 中零元的一个元
$(a_1, \ldots, a_{n + m}) \in A^{\oplus n + m}$。
因为 $\overline{g}_1, \ldots, \overline{g}_m$ 在 $A/I$ 中构成
$H_1$ 正则序列，故
$(\overline{a}_{n + 1}, \ldots, \overline{a}_{n + m})$ 是
$\wedge^2(A/I^{\oplus m})$ 中某个元 $\overline{\alpha}$ 的像。
可将 $\overline{\alpha}$ 提升为一个元
$\alpha \in \wedge^2(A^{\oplus n + m})$，并从
$(a_1, \ldots, a_{n + m})$ 中减去它在 $A^{\oplus n + m}$ 中的像。
因此可以假设 $a_{n + 1}, \ldots, a_{n + m} \in I$。
由于 $I = (f_1, \ldots, f_n)$，可用顶部左侧水平箭头的像中诸元
$$
(0, \ldots, g_j, 0, \ldots, 0, f_i, 0, \ldots, 0)
$$
的线性组合修改 $(a_1, \ldots, a_{n + m})$，从而化归到
$a_{n + 1}, \ldots, a_{n + m}$ 均为零的情形。在此情形下，
$(a_1, \ldots, a_n, 0, \ldots, 0)$ 定义了
$H_1(A, f_1, \ldots, f_n)$ 的一个元，而我们已假设该模为零。
\end{proof}

\begin{lemma}
\label{more-algebra-lemma-truncate-H1-regular}
设 $A$ 为环，$f_1, \ldots, f_n, g_1, \ldots, g_m \in A$ 是
$H_1$ 正则序列。则 $\overline{g}_1, \ldots, \overline{g}_m$ 在
$A/(f_1, \ldots, f_n)$ 中的像构成 $H_1$ 正则序列。
\end{lemma}

\begin{proof}
令 $I = (f_1, \ldots, f_n)$。我们必须证明，$A/I$ 中的任一关系
$\sum_{j = 1, \ldots, m} \overline{a}_j \overline{g}_j$
都是平凡关系的线性组合。由于 $I = (f_1, \ldots, f_n)$，可将此关系提升为
$A$ 中的关系
$$
\sum\nolimits_{j = 1, \ldots, m} a_j g_j +
\sum\nolimits_{i = 1, \ldots, n} b_if_i = 0
$$
由假设，$A$ 中的这一关系是平凡关系的线性组合。取其在 $A/I$ 中的像，
便得到所需结论。
\end{proof}

\begin{lemma}
\label{more-algebra-lemma-join-koszul-regular-sequences}
设 $A$ 为环，$I$ 是由 $A$ 中的 Koszul 正则序列
$f_1, \ldots, f_n$ 生成的理想。设 $g_1, \ldots, g_m \in A$，且其像
$\overline{g}_1, \ldots, \overline{g}_m$ 在 $A/I$ 中构成 Koszul 正则序列。
则 $f_1, \ldots, f_n, g_1, \ldots, g_m$ 是 $A$ 中的 Koszul 正则序列。
\end{lemma}

\begin{proof}
我们的假设说明 $K_\bullet(A, f_1, \ldots, f_n)$ 是 $A/I$ 的有限自由解消，且
$K_\bullet(A/I, \overline{g}_1, \ldots, \overline{g}_m)$ 是
$A/(f_i, g_j)$ 在 $A/I$ 上的有限自由解消。于是
\begin{align*}
K_\bullet(A, f_1, \ldots, f_n, g_1, \ldots, g_m)
& = \text{Tot}(K_\bullet(A, f_1, \ldots, f_n) \otimes_A
K_\bullet(A, g_1, \ldots, g_m)) \\
& \cong A/I \otimes_A K_\bullet(A, g_1, \ldots, g_m) \\
& = K_\bullet(A/I, \overline{g}_1, \ldots, \overline{g}_m) \\
& \cong A/(f_i, g_j)
\end{align*}
第一个等式由引理 \ref{more-algebra-lemma-join-sequences-koszul-complex} 得出。
由于双复形
$K_\bullet(A, f_1, \ldots, f_n) \otimes_A K_\bullet(A, g_1, \ldots, g_m)$
的第 $q$ 行是 $A/I \otimes_A K_q(A, g_1, \ldots, g_m)$ 的解消，
第一个拟同构 $\cong$ 由 Homology, Lemma
\ref{homology-lemma-double-complex-gives-resolution}（的对偶）得出。
第二个等式是清楚的。最后一个拟同构由假设得出。故结论成立。
\end{proof}

\noindent
为得到下一个引理的结论，有必要同时假设
$f_1, \ldots, f_n$ 与 $f_1, \ldots, f_n, g_1, \ldots, g_m$
都是 Koszul 正则的。若删去 $f_1, \ldots, f_n$ Koszul 正则这一假设，
反例见 Examples, Lemma \ref{examples-lemma-strange-regular-sequence}。

\begin{lemma}
\label{more-algebra-lemma-truncate-koszul-regular}
设 $A$ 为环，$f_1, \ldots, f_n, g_1, \ldots, g_m \in A$。
若 $f_1, \ldots, f_n$ 与 $f_1, \ldots, f_n, g_1, \ldots, g_m$
都是 $A$ 中的 Koszul 正则序列，则
$\overline{g}_1, \ldots, \overline{g}_m$ 在 $A/(f_1, \ldots, f_n)$ 中
构成 Koszul 正则序列。
\end{lemma}

\begin{proof}
令 $I = (f_1, \ldots, f_n)$。我们的假设说明
$K_\bullet(A, f_1, \ldots, f_n)$ 是 $A/I$ 的有限自由解消，且
$K_\bullet(A, f_1, \ldots, f_n, g_1, \ldots, g_m)$ 是
$A/(f_i, g_j)$ 在 $A$ 上的有限自由解消。于是
\begin{align*}
A/(f_i, g_j) & \cong K_\bullet(A, f_1, \ldots, f_n, g_1, \ldots, g_m) \\
& = \text{Tot}(K_\bullet(A, f_1, \ldots, f_n) \otimes_A
K_\bullet(A, g_1, \ldots, g_m)) \\
& \cong A/I \otimes_A K_\bullet(A, g_1, \ldots, g_m) \\
& = K_\bullet(A/I, \overline{g}_1, \ldots, \overline{g}_m)
\end{align*}
第一个拟同构 $\cong$ 由假设得出。第一个等式由引理
\ref{more-algebra-lemma-join-sequences-koszul-complex} 得出。由于双复形
$K_\bullet(A, f_1, \ldots, f_n) \otimes_A K_\bullet(A, g_1, \ldots, g_m)$
的第 $q$ 行是 $A/I \otimes_A K_q(A, g_1, \ldots, g_m)$ 的解消，
第二个拟同构由 Homology, Lemma
\ref{homology-lemma-double-complex-gives-resolution}（的对偶）得出。
第二个等式是清楚的。故结论成立。
\end{proof}

\begin{lemma}
\label{more-algebra-lemma-independence-of-generators}
设 $R$ 为环，$I$ 是由 $f_1, \ldots, f_r \in R$ 生成的理想。
\begin{enumerate}
\item 若 $I$ 可由长度为 $r$ 的拟正则序列生成，则
$f_1, \ldots, f_r$ 是拟正则序列。
\item 若 $I$ 可由长度为 $r$ 的 $H_1$ 正则序列生成，则
$f_1, \ldots, f_r$ 是 $H_1$ 正则序列。
\item 若 $I$ 可由长度为 $r$ 的 Koszul 正则序列生成，则
$f_1, \ldots, f_r$ 是 Koszul 正则序列。
\end{enumerate}
\end{lemma}

\begin{proof}
若 $I$ 可由长度为 $r$ 的拟正则序列生成，则 $I/I^2$ 是秩为 $r$ 的自由
$R/I$-模。由于按假设 $f_1, \ldots, f_r$ 是一组生成元，可知其像
$\overline{f}_i$ 构成 $I/I^2$ 在 $R/I$ 上的一组基。因此
$f_1, \ldots, f_r$ 是拟正则序列，因为除 $I/I^2$ 的自由性外，这一切恰好意味着映射
$\text{Sym}^n_{R/I}(I/I^2) \to I^n/I^{n + 1}$ 均为同构。

\medskip\noindent
继续假设 $I$ 可由某个拟正则序列生成，记为 $g_1, \ldots, g_r$。
写成 $g_j = \sum a_{ij}f_i$。由上一段，$f_1, \ldots, f_r$ 拟正则，
故 $\det(a_{ij})$ 模 $I$ 可逆。矩阵 $a_{ij}$ 给出映射
$R^{\oplus r} \to R^{\oplus r}$，并诱导 Koszul 复形之间的映射
$\alpha : K_\bullet(R, f_1, \ldots, f_r) \to K_\bullet(R, g_1, \ldots, g_r)$；
见引理 \ref{more-algebra-lemma-functorial}。将 $\det(a_{ij})$ 可逆化后，该映射成为同构。
因为 $K_\bullet(R, f_1, \ldots, f_r)$ 与
$K_\bullet(R, g_1, \ldots, g_r)$ 的上同调模均被 $I$ 零化；见引理
\ref{more-algebra-lemma-homotopy-koszul}，可知 $\alpha$ 是拟同构。

\medskip\noindent
现在假设 $g_1, \ldots, g_r$ 是生成 $I$ 的 $H_1$ 正则序列。
由引理 \ref{more-algebra-lemma-H1-regular-quasi-regular}，$g_1, \ldots, g_r$ 是拟正则序列。
由上一段可知 $f_1, \ldots, f_r$ 是 $H_1$ 正则序列。
对 Koszul 正则序列同理。
\end{proof}

\begin{lemma}
\label{more-algebra-lemma-make-nonzero-divisor}
\begin{reference}
这是 \cite[Corollary]{McCoy} 的一个特殊情形。
\end{reference}
设 $R$ 为环，$a_1, \ldots, a_n \in R$ 为诸元，并且
$R \to R^{\oplus n}$，$x \mapsto (xa_1, \ldots, xa_n)$ 是单射。
则多项式环 $R[t_1, \ldots, t_n]$ 中的元 $\sum a_i t_i$ 是非零因子。
\end{lemma}

\begin{proof}
若某个 $a_i$ 是单位元，这恰是说，在 $R$ 上的多项式环中，任意形如
$t_1 + a_2 t_2 + \ldots + a_n t_n$ 的元都是非零因子。

\medskip\noindent
情形 I：$R$ 是 Noether 环。设 $\mathfrak q_j$，$j = 1, \ldots, m$，
为 $R$ 的伴随素理想。我们必须证明各映射
$$
\sum a_i t_i :
\text{Sym}^d(R^{\oplus n})
\longrightarrow
\text{Sym}^{d + 1}(R^{\oplus n})
$$
都是单射。由于 $\text{Sym}^d(R^{\oplus n})$ 是自由 $R$-模，
其伴随素理想为 $\mathfrak q_j$，$j = 1, \ldots, m$。
对每个 $j$，存在 $i = i(j)$ 使 $a_i \not \in \mathfrak q_j$，
这是因为存在 $x \in R$ 满足 $\mathfrak q_jx = 0$，而由假设，
对某个 $i$ 有 $a_i x \not = 0$。因此 $a_i$ 在 $R_{\mathfrak q_j}$ 中是单位元，
且该映射在 $\mathfrak q_j$ 处局部化后是单射。故该映射本身是单射；见
Algebra, Lemma \ref{algebra-lemma-zero-at-ass-zero}。

\medskip\noindent
情形 II：一般的 $R$。可以将 $R$ 写为 Noether 环 $R_\lambda$ 的并，且
$a_1, \ldots, a_n \in R_\lambda$。结论对每个 $R_\lambda$ 成立，
因而也对 $R$ 成立。
\end{proof}

\begin{lemma}
\label{more-algebra-lemma-Koszul-regular-flat-locally-regular}
设 $R$ 为环，$f_1, \ldots, f_n$ 是 $R$ 中的 Koszul 正则序列，且
$(f_1, \ldots, f_n) \not = R$。考虑忠实平坦的光滑环映射
$$
R \longrightarrow
S = R[\{t_{ij}\}_{i \leq j}, t_{11}^{-1}, t_{22}^{-1}, \ldots, t_{nn}^{-1}]
$$
对 $1 \leq i \leq n$，令
$$
g_i = \sum\nolimits_{i \leq j} t_{ij} f_j \in S.
$$
则 $g_1, \ldots, g_n$ 是 $S$ 中的正则序列，且
$(f_1, \ldots, f_n)S = (g_1, \ldots, g_n)$。
\end{lemma}

\begin{proof}
理想相等是显然的，因为矩阵
$$
\left(
\begin{matrix}
t_{11} & t_{12} & t_{13} & \ldots \\
0 & t_{22} & t_{23} & \ldots \\
0 & 0 & t_{33} & \ldots \\
\ldots & \ldots & \ldots & \ldots
\end{matrix}
\right)
$$
在 $S$ 中可逆。因为 $f_1, \ldots, f_n$ 是 Koszul 正则序列，可知
$R \to R^{\oplus n}$，$x \mapsto (xf_1, \ldots, xf_n)$ 的核为零
（因为它计算 $R$ 关于 $f_1, \ldots, f_n$ 的第 $n$ 个 Koszul 同调）。
因此由引理 \ref{more-algebra-lemma-make-nonzero-divisor} 可知
$g_1 = f_1 t_{11} + \ldots + f_n t_{1n}$ 在
$S' = R[t_{11}, t_{12}, \ldots, t_{1n}, t_{11}^{-1}]$ 中是非零因子。
由引理 \ref{more-algebra-lemma-koszul-regular-flat-base-change} 与
\ref{more-algebra-lemma-independence-of-generators}，可知
$g_1, f_2, \ldots, f_n$ 是 $S'$ 中的 Koszul 序列。
由引理 \ref{more-algebra-lemma-truncate-koszul-regular}，可推出
$\overline{f}_2, \ldots, \overline{f}_n$ 是 $S'/(g_1)$ 中的 Koszul 正则序列。
因此，对 $n$ 归纳可知，$g_2, \ldots, g_n$ 在
$S'/(g_1)[\{t_{ij}\}_{2 \leq i \leq j}, t_{22}^{-1}, \ldots, t_{nn}^{-1}]$
中的像 $\overline{g}_2, \ldots, \overline{g}_n$ 构成正则序列。
这又意味着 $g_1, \ldots, g_n$ 在 $S$ 中构成正则序列。
\end{proof}







\section{Koszul 正则序列的进一步结果}
\label{more-algebra-section-more-koszul-regular}

\noindent
我们继续第 \ref{more-algebra-section-koszul-regular} 节中的讨论。

\begin{lemma}
\label{more-algebra-lemma-vanishing-extended-alternating-koszul}
设 $R$ 为环，$f_1, \ldots, f_r \in R$ 是 Koszul 正则序列。则第
\ref{more-algebra-section-alternating-cech} 节中的增广交错 {\v C}ech 复形
$R \to \bigoplus\nolimits_{i_0} R_{f_{i_0}} \to
\bigoplus\nolimits_{i_0 < i_1} R_{f_{i_0}f_{i_1}} \to
\ldots \to R_{f_1\ldots f_r}$ 仅在次数 $r$ 上有上同调。
\end{lemma}

\begin{proof}
由引理 \ref{more-algebra-lemma-mult-koszul-regular} 与归纳法，对所有 $n \geq 1$，
序列 $f_1, \ldots, f_{r - 1}, f_r^n$ 都是 Koszul 正则的。
由引理 \ref{more-algebra-lemma-change-basis}，Koszul 正则序列的任一置换仍是
Koszul 正则序列。因此可将任意一个（或全部）$f_i$ 替换为其第 $n$ 次幂，
所得序列仍然 Koszul 正则。从而 $K_\bullet(R, f_1^n, \ldots, f_r^n)$
仅在同调次数 $0$ 上有非零上同调。由引理
\ref{more-algebra-lemma-extended-alternating-Cech-is-colimit-koszul}，这推出所需结论。
\end{proof}

\begin{lemma}
\label{more-algebra-lemma-blowup-regular-sequence}
设 $a, a_2, \ldots, a_r$ 是环 $R$ 中的 $H_1$ 正则序列
（例如 Koszul 正则序列或正则序列；见引理
\ref{more-algebra-lemma-regular-koszul-regular} 与
\ref{more-algebra-lemma-koszul-regular-H1-regular}）。令 $I = (a, a_2, \ldots, a_r)$，
则爆破代数 $R' = R[\frac{I}{a}]$ 同构于
$R'' = R[y_2, \ldots, y_r]/(a y_i - a_i)$。
\end{lemma}

\begin{proof}
由 Algebra, Lemma \ref{algebra-lemma-affine-blowup-quotient-description}，
只需证明 $R''$ 无 $a$-挠。

\medskip\noindent
我们断言 $a, ay_2 - a_2, \ldots, ay_n - a_r$ 在
$R[y_2, \ldots, y_r]$ 中是 $H_1$ 正则序列。事实上，用来定义
$a, ay_2 - a_2, \ldots, ay_n - a_r$ 上 Koszul 复形的映射
$$
(a, ay_2 - a_2, \ldots, ay_n - a_r) :
R[y_2, \ldots, y_r]^{\oplus r}
\longrightarrow
R[y_2, \ldots, y_r]
$$
通过同构
$$
R[y_2, \ldots, y_r]^{\oplus r}
\longrightarrow
R[y_2, \ldots, y_r]^{\oplus r}
$$
同构于用来定义 $a, a_2, \ldots, a_r$ 上 Koszul 复形的映射
$$
(a, a_2, \ldots, a_r) :
R[y_2, \ldots, y_r]^{\oplus r} \longrightarrow
R[y_2, \ldots, y_r]
$$
其中该同构将 $(b_1, \ldots, b_r)$ 映为
$(b_1 - b_2y_2 \ldots - b_ry_r, -b_2, \ldots, - b_r)$。
由引理 \ref{more-algebra-lemma-functorial}，这些 Koszul 复形同构。
将引理 \ref{more-algebra-lemma-koszul-regular-flat-base-change} 应用于平坦环映射
$R \to R[y_2, \ldots, y_r]$，便证明了断言。由引理
\ref{more-algebra-lemma-cone-koszul}，$a, ay_2 - a_2, \ldots, ay_n - a_r$
上的 Koszul 复形 $K$ 是映射 $a : L \to L$ 的锥，其中 $L$ 是
$ay_2 - a_2, \ldots, ay_n - a_r$ 上的 Koszul 复形。
由断言，$H_1(K) = 0$，故 $a : H_0(L) \to H_0(L)$ 是单射；换言之，
$R'' = R[y_2, \ldots, y_r]/(a y_i - a_i)$ 没有非零的 $a$-挠元，正合所需。
\end{proof}

\begin{lemma}
\label{more-algebra-lemma-base-change-H1-regular}
设 $A \to B$ 为环映射。设 $f_1, \ldots, f_r$ 是 $B$ 中的序列，且
$B/(f_1, \ldots, f_r)$ 在 $A$ 上平坦。设 $A \to A'$ 为环映射。
则典范映射
$$
H_1(K_\bullet(B, f_1, \ldots, f_r)) \otimes_A A'
\longrightarrow
H_1(K_\bullet(B', f'_1, \ldots, f'_r))
$$
是满射。这里 $B' = B \otimes_A A'$，而 $f_i' \in B'$ 是 $f_i$ 的像。
\end{lemma}

\begin{proof}
序列
$$
\wedge^2(B^{\oplus r}) \to B^{\oplus r} \to B \to B/J \to 0
$$
是 $A$-模的复形，其中 $B/J$ 在 $A$ 上平坦，并且在位置
$B^{\oplus r}$ 处的上同调群为
$H_1 = H_1(K_\bullet(B, f_1, \ldots, f_r))$。将它与 $A'$ 张量，得到复形
$$
\wedge^2((B')^{\oplus r}) \to (B')^{\oplus r} \to B' \to B'/J' \to 0
$$
该复形在 $B'$ 与 $B'/J'$ 处正合。为了计算它在 $(B')^{\oplus r}$ 处的上同调群
$$
H'_1 = H_1(K_\bullet(B', f'_1, \ldots, f'_r))
$$
。将上面的第一个序列拆成如下三个正合列：
$$
0 \to J \to B \to B/J \to 0
$$
$$
0 \to K \to B^{\oplus r} \to J \to 0
$$
$$
\wedge^2(B^{\oplus r}) \to K \to H_1 \to 0
$$
在 $A$ 上与 $A'$ 张量，得到正合列
$$
\begin{matrix}
0 \to J \otimes_A A' \to B \otimes_A A' \to (B/J) \otimes_A A' \to 0 \\
K \otimes_A A' \to B^{\oplus r} \otimes_A A' \to J \otimes_A A' \to 0 \\
\wedge^2(B^{\oplus r}) \otimes_A A' \to K \otimes_A A' \to H_1 \otimes_A A'
\to 0
\end{matrix}
$$
其中第一个序列正合，因为 $B/J$ 在 $A$ 上平坦；见 Algebra, Lemma
\ref{algebra-lemma-flat-tor-zero}。由此可知 $J' = J \otimes_A A' \subset B'$，
且 $K \otimes_A A' \to \Ker((B')^{\oplus r} \to B')$ 是满射。因此
\begin{align*}
H_1 \otimes_A A'
& =
\Coker\left(\wedge^2(B^{\oplus r}) \otimes_A A' \to K \otimes_A A'\right) \\
& \to
\Coker\left(
\wedge^2((B')^{\oplus r})  \to \Ker((B')^{\oplus r} \to B')
\right) = H'_1
\end{align*}
也是满射。
\end{proof}

\begin{lemma}
\label{more-algebra-lemma-relative-regular-immersion-algebra}
设 $A \to B$ 与 $A \to A'$ 为环映射，令 $B' = B \otimes_A A'$。
设 $f_1, \ldots, f_r \in B$。假设 $B/(f_1, \ldots, f_r)B$ 在 $A$ 上平坦。
\begin{enumerate}
\item 若 $f_1, \ldots, f_r$ 是拟正则序列，则其在 $B'$ 中的像是拟正则序列。
\item 若 $f_1, \ldots, f_r$ 是 $H_1$ 正则序列，则其在 $B'$ 中的像是
$H_1$ 正则序列。
\end{enumerate}
\end{lemma}

\begin{proof}
假设 $f_1, \ldots, f_r$ 拟正则。令 $J = (f_1, \ldots, f_r)$。
由假设，$J^n/J^{n + 1}$ 同构于 $B/J$ 的若干副本的直和，因而在 $A$ 上平坦。
由归纳法及 Algebra, Lemma \ref{algebra-lemma-flat-ses}，可知
$B/J^n$ 在 $A$ 上平坦。理想 $(J')^n$ 等于 $J^n \otimes_A A'$；见
Algebra, Lemma \ref{algebra-lemma-flat-tor-zero}。因此
$(J')^n/(J')^{n + 1} = J^n/J^{n + 1} \otimes_A A'$，这显然推出
$f_1, \ldots, f_r$ 在 $B'$ 中是拟正则序列。

\medskip\noindent
假设 $f_1, \ldots, f_r$ 是 $H_1$ 正则的。由引理
\ref{more-algebra-lemma-base-change-H1-regular}，Koszul 同调群
$H_1(K_\bullet(B, f_1, \ldots, f_r))$ 的消失推出
$H_1(K_\bullet(B', f'_1, \ldots, f'_r))$ 的消失，故结论成立。
\end{proof}

\begin{lemma}
\label{more-algebra-lemma-cut-by-koszul-flat}
设 $A' \to B'$ 为环映射，$I \subset A'$ 为理想。令
$A = A'/I$ 且 $B = B'/IB'$。设 $f'_1, \ldots, f'_r \in B'$。假设：
\begin{enumerate}
\item $A' \to B'$ 平坦且为有限表现的；
\item $I$ 是局部幂零的；
\item 诸像 $f_1, \ldots, f_r \in B$ 构成拟正则序列；
\item $B/(f_1, \ldots, f_r)$ 在 $A$ 上平坦。
\end{enumerate}
则 $B'/(f'_1, \ldots, f'_r)$ 在 $A'$ 上平坦。
\end{lemma}

\begin{proof}
令 $C' = B'/(f'_1, \ldots, f'_r)$。我们必须证明 $A' \to C'$ 平坦。
设 $\mathfrak r' \subset C'$ 是位于 $\mathfrak p' \subset A'$ 之上的素理想，
并令 $\mathfrak q' \subset B'$ 为 $\mathfrak r'$ 的逆像。
由 Algebra, Lemma \ref{algebra-lemma-flat-localization}，只需证明
$A'_{\mathfrak p'} \to C'_{\mathfrak q'}$ 平坦。
Algebra, Lemma \ref{algebra-lemma-grothendieck-regular-sequence-general}
说明，只需证明 $f'_1, \ldots, f'_r$ 映到
$$
B'_{\mathfrak q'}/\mathfrak p'B'_{\mathfrak q'} =
B_\mathfrak q/\mathfrak p B_\mathfrak q =
(B \otimes_A \kappa(\mathfrak p))_\mathfrak q
$$
中的一个正则序列；此处记号含义显然。已知 $f_1, \ldots, f_r$
在 $B$ 中是拟正则序列，且 $B/(f_1, \ldots, f_r)$ 在 $A$ 上平坦。
由引理 \ref{more-algebra-lemma-relative-regular-immersion-algebra}，$f'_1, \ldots, f'_r$
在 $B \otimes_A \kappa(\mathfrak p)$ 中的像
$\overline{f}_1, \ldots, \overline{f}_r$ 构成拟正则序列。
由于 $(B \otimes_A \kappa(\mathfrak p))_\mathfrak q$ 是 Noether 局部环，
由引理 \ref{more-algebra-lemma-noetherian-finite-all-equivalent} 即得结论。
\end{proof}

\begin{lemma}
\label{more-algebra-lemma-cut-by-koszul}
设 $A' \to B'$ 为环映射，$I \subset A'$ 为理想。令
$A = A'/I$ 且 $B = B'/IB'$。设 $f'_1, \ldots, f'_r \in B'$。假设：
\begin{enumerate}
\item $A' \to B'$ 平坦且为有限表现的（例如光滑）；
\item $I$ 是局部幂零的；
\item 诸像 $f_1, \ldots, f_r \in B$ 构成拟正则序列；
\item $B/(f_1, \ldots, f_r)$ 在 $A$ 上光滑。
\end{enumerate}
则 $B'/(f'_1, \ldots, f'_r)$ 在 $A'$ 上光滑。
\end{lemma}

\begin{proof}
令 $C' = B'/(f'_1, \ldots, f'_r)$ 且 $C = B/(f_1, \ldots, f_r)$。
则 $A' \to C'$ 为有限表现的。由引理 \ref{more-algebra-lemma-cut-by-koszul-flat}，
$A' \to C'$ 平坦。映射 $A' \to C'$ 的纤维环等于 $A \to C$ 的纤维环，
故由假设 (4) 它们光滑。因此由 Algebra, Lemma
\ref{algebra-lemma-flat-fibre-smooth}，$A' \to C'$ 光滑。
\end{proof}








\section{正则理想}
\label{more-algebra-section-ideals}

\noindent
我们将在 Divisors, Section \ref{divisors-section-regular-ideal-sheaves}
中以很大的一般性讨论正则理想层的概念。这里定义仿射情形中的相应概念，
即环中理想这一情形。

\begin{definition}
\label{more-algebra-definition-regular-ideal}
设 $R$ 为环，$I \subset R$ 为理想。
\begin{enumerate}
\item 若对每个 $\mathfrak p \in V(I)$，存在 $g \in R$、
$g \not \in \mathfrak p$ 以及正则序列
$f_1, \ldots, f_r \in R_g$，使 $I_g$ 由 $f_1, \ldots, f_r$ 生成，
则称 $I$ 为{\it 正则理想}。
\item 若对每个 $\mathfrak p \in V(I)$，存在 $g \in R$、
$g \not \in \mathfrak p$ 以及 Koszul 正则序列
$f_1, \ldots, f_r \in R_g$，使 $I_g$ 由 $f_1, \ldots, f_r$ 生成，
则称 $I$ 为{\it Koszul 正则理想}。
\item 若对每个 $\mathfrak p \in V(I)$，存在 $g \in R$、
$g \not \in \mathfrak p$ 以及 $H_1$ 正则序列
$f_1, \ldots, f_r \in R_g$，使 $I_g$ 由 $f_1, \ldots, f_r$ 生成，
则称 $I$ 为{\it $H_1$ 正则理想}。
\item 若对每个 $\mathfrak p \in V(I)$，存在 $g \in R$、
$g \not \in \mathfrak p$ 以及拟正则序列
$f_1, \ldots, f_r \in R_g$，使 $I_g$ 由 $f_1, \ldots, f_r$ 生成，
则称 $I$ 为{\it 拟正则理想}。
\end{enumerate}
\end{definition}

\noindent
显然，给定 $I \subset R$，有推含关系
\begin{align*}
I\text{ is a regular ideal}
& \Rightarrow
I\text{ is a Koszul-regular ideal} \\
& \Rightarrow
I\text{ is a }H_1\text{-regular ideal} \\
& \Rightarrow
I\text{ is a quasi-regular ideal}
\end{align*}
见引理 \ref{more-algebra-lemma-regular-koszul-regular}、
\ref{more-algebra-lemma-koszul-regular-H1-regular} 与
\ref{more-algebra-lemma-H1-regular-quasi-regular}。这样的理想总是有限生成的。

\begin{lemma}
\label{more-algebra-lemma-quasi-regular-ideal-finite}
拟正则理想是有限生成的。
\end{lemma}

\begin{proof}
设 $I \subset R$ 为拟正则理想。由于 $V(I)$ 拟紧，存在
$g_1, \ldots, g_m \in R$，使
$V(I) \subset D(g_1) \cup \ldots \cup D(g_m)$，并且 $I_{g_j}$ 由拟正则序列
$g_{j1}, \ldots, g_{jr_j} \in R_{g_j}$ 生成。对某些 $g'_{ij} \in I$，写成
$g_{ji} = g'_{ji}/g_j^{e_{ij}}$。写 $1 + x = \sum g_j h_j$，其中
$x \in I$；这是可能的，因为 $V(I) \subset D(g_1) \cup \ldots \cup D(g_m)$。
注意
$\Spec(R) = D(g_1) \cup \ldots \cup D(g_m) \bigcup D(x)$。
于是 $I$ 由诸元 $g'_{ij}$ 与 $x$ 生成，因为这些元在该覆盖的每一片上生成；
见 Algebra, Lemma \ref{algebra-lemma-cover}。
\end{proof}

\begin{lemma}
\label{more-algebra-lemma-quasi-regular-ideal-finite-projective}
设 $I \subset R$ 是环中的拟正则理想。则 $I/I^2$ 是有限投射 $R/I$-模。
\end{lemma}

\begin{proof}
这由 Algebra, Lemma \ref{algebra-lemma-finite-projective} 与定义得出。
\end{proof}

\noindent
我们证明 Koszul 正则理想、$H_1$ 正则理想和拟正则理想的平坦下降。

\begin{lemma}
\label{more-algebra-lemma-flat-descent-regular-ideal}
设 $A \to B$ 为忠实平坦环映射，$I \subset A$ 为理想。
若 $IB$ 是 $B$ 中的 Koszul 正则（相应地，$H_1$ 正则；拟正则）理想，
则 $I$ 是 $A$ 中的 Koszul 正则（相应地，$H_1$ 正则；拟正则）理想。
\end{lemma}

\begin{proof}
在整个证明中固定素理想 $\mathfrak p \supset I$。假设 $IB$ 拟正则。
由引理 \ref{more-algebra-lemma-quasi-regular-ideal-finite}，$IB$ 是有限模，故由
Algebra, Lemma \ref{algebra-lemma-descend-properties-modules}，
$I$ 是有限 $A$-模。因为 $A \to B$ 平坦，可知
$$
I/I^2 \otimes_{A/I} B/IB = I/I^2 \otimes_A B = IB/(IB)^2.
$$
由于 $IB$ 拟正则，$B/IB$-模 $IB/(IB)^2$ 有限局部自由。因此由
Algebra, Proposition \ref{algebra-proposition-ffdescent-finite-projectivity}，
$I/I^2$ 有限投射。特别地，对某个 $f \in A$、$f \not \in \mathfrak p$，
以 $A_f$ 替换 $A$ 后，可假设 $I/I^2$ 是秩为 $r$ 的自由模。
选取 $f_1, \ldots, f_r \in I$，使其给出 $I/I^2$ 的一组基。
由 Nakayama 引理（见 Algebra, Lemma \ref{algebra-lemma-NAK}），
再次进行上述替换 $A \leadsto A_f$ 后，可知 $I$ 由 $f_1, \ldots, f_r$ 生成。

\medskip\noindent
“拟正则”情形的证明。上面已经看到，$I/I^2$ 以 $r$ 个生成元
$f_1, \ldots, f_r$ 为自由基。为完成此情形的证明，必须证明对每个
$d \geq 2$，映射 $\text{Sym}^d(I/I^2) \to I^d/I^{d + 1}$ 都是同构。
这是清楚的，因为其忠实平坦基变换
$\text{Sym}^d(IB/(IB)^2) \to (IB)^d/(IB)^{d + 1}$ 由假设在 $B$ 上局部为同构。
略去细节。

\medskip\noindent
“$H_1$ 正则”与“Koszul 正则”情形的证明。考虑上面构造的生成 $I$ 的序列
$f_1, \ldots, f_r$。由引理 \ref{more-algebra-lemma-independence-of-generators}，
对任意 $g \in B$，只要 $IB$ 由某个 $H_1$ 正则或 Koszul 正则序列生成，
$f_1, \ldots, f_r$ 在 $B_g$ 中的像便是 $H_1$ 正则或 Koszul 正则序列。
因此 $K_\bullet(A, f_1, \ldots, f_r) \otimes_A B_g$ 的
$H_1$ 或 $H_i$（$i > 0$）消失。由于
$K_\bullet(B, f_1, \ldots, f_r) = K_\bullet(A, f_1, \ldots, f_r) \otimes_A B$
的同调被 $IB$ 零化（见引理 \ref{more-algebra-lemma-homotopy-koszul}），且
$V(IB) \subset \bigcup_{g\text{ 如上}} D(g)$，可知
$K_\bullet(A, f_1, \ldots, f_r) \otimes_A B$ 在次数 $1$ 或所有正次数上的同调消失。
再利用 $A \to B$ 忠实平坦，可知
$K_\bullet(A, f_1, \ldots, f_r)$ 也具有同样的性质。
\end{proof}

\begin{lemma}
\label{more-algebra-lemma-conormal-sequence-H1-regular-ideal}
设 $A$ 为环，$I \subset J \subset A$ 为理想。假设
$J/I \subset A/I$ 是 $H_1$ 正则理想。则 $I \cap J^2 = IJ$。
\end{lemma}

\begin{proof}
局部化后，这立即由引理 \ref{more-algebra-lemma-conormal-sequence-H1-regular} 得出。
\end{proof}







\section{局部完全交映射}
\label{more-algebra-section-lci}

\noindent
可以利用上面的材料，通过（有限）多项式代数的表现来定义环之间的局部完全交映射。

\begin{lemma}
\label{more-algebra-lemma-koszul-independence-presentation}
设 $A \to B$ 是有限型环映射。若对某个表现
$\alpha : A[x_1, \ldots, x_n] \to B$，其核 $I$ 是 Koszul 正则理想，
则对任一表现 $\beta : A[y_1, \ldots, y_m] \to B$，其核 $J$ 都是
Koszul 正则理想。
\end{lemma}

\begin{proof}
选取 $f_j \in A[x_1, \ldots, x_n]$ 使 $\alpha(f_j) = \beta(y_j)$，并选取
$g_i \in A[y_1, \ldots, y_m]$ 使 $\beta(g_i) = \alpha(x_i)$。
于是得到交换图
$$
\xymatrix{
A[x_1, \ldots, x_n, y_1, \ldots, y_m]
\ar[d]^{x_i \mapsto g_i} \ar[rr]_-{y_j \mapsto f_j} & &
A[x_1, \ldots, x_n] \ar[d] \\
A[y_1, \ldots, y_m] \ar[rr] & & B
}
$$
注意，$A[x_i, y_j] \to B$ 的核 $K$ 等于
$K = (I, y_j - f_j) = (J, x_i - g_i)$。特别地，由引理
\ref{more-algebra-lemma-quasi-regular-ideal-finite}，$I$ 有限生成，故
$J = K/(x_i - g_i)$ 也有限生成。

\medskip\noindent
选取素理想 $\mathfrak q \subset B$。由于
$I/I^2 \oplus B^{\oplus m} = J/J^2 \oplus B^{\oplus n}$
（Algebra, Lemma \ref{algebra-lemma-conormal-module}），可知
$$
\dim J/J^2 \otimes_B \kappa(\mathfrak q) + n =
\dim I/I^2 \otimes_B \kappa(\mathfrak q) + m.
$$
选取 $p_1, \ldots, p_t \in I$，使其映到
$I/I^2 \otimes \kappa(\mathfrak q) = I \otimes_{A[x_i]} \kappa(\mathfrak q)$
的一组基。选取 $q_1, \ldots, q_s \in J$，使其映到
$J/J^2 \otimes \kappa(\mathfrak q) = J \otimes_{A[y_j]} \kappa(\mathfrak q)$
的一组基。于是 $s + n = t + m$。由 Nakayama 引理，存在
$h \in A[x_i]$ 与 $h' \in A[y_j]$，两者都映到
$\kappa(\mathfrak q)$ 中的非零元，并且 $A[x_i, 1/h]$ 中
$I_h = (p_1, \ldots, p_t)$，$A[y_j, 1/h']$ 中
$J_{h'} = (q_1, \ldots, q_s)$。由于 $I$ Koszul 正则，还可假设
$I_h$ 由某个 Koszul 正则序列生成。该序列的长度必为
$t = \dim I/I^2 \otimes_B \kappa(\mathfrak q)$，故由引理
\ref{more-algebra-lemma-independence-of-generators}，$p_1, \ldots, p_t$ 是 Koszul 正则序列。
又因为 $y_1 - f_1, \ldots, y_m - f_m$ 是正则序列，由引理
\ref{more-algebra-lemma-join-koszul-regular-sequences} 可知
$$
y_1 - f_1, \ldots, y_m - f_m, p_1, \ldots, p_t
$$
是 $A[x_i, y_j, 1/h]$ 中的 Koszul 正则序列。该序列生成理想 $K_h$。
因此，理想 $K_{hh'}$ 由长度为 $m + t = n + s$ 的 Koszul 正则序列生成。
另一方面，它也由同样长度的序列
$$
x_1 - g_1, \ldots, x_n - g_n, q_1, \ldots, q_s
$$
生成，故由引理 \ref{more-algebra-lemma-independence-of-generators}，该序列也是 Koszul 正则的。
最后，由引理 \ref{more-algebra-lemma-truncate-koszul-regular}，可知
$q_1, \ldots, q_s$ 在
$$
A[x_i, y_j, 1/hh']/(x_1 - g_1, \ldots, x_n - g_n)
\cong A[y_j, 1/h'']
$$
中的像构成生成 $J_{h''}$ 的 Koszul 正则序列。由于 $h''$ 是 $hh'$ 的像，
它在 $\kappa(\mathfrak q)$ 中的像非零，故结论成立。
\end{proof}

\noindent
该引理使我们可以作如下定义。

\begin{definition}
\label{more-algebra-definition-local-complete-intersection}
若环映射 $A \to B$ 为有限型的，并且对某个（等价地，对任一）表现
$B = A[x_1, \ldots, x_n]/I$，理想 $I$ Koszul 正则，
则称它为{\it 局部完全交}。
\end{definition}

\noindent
这一概念是局部的。

\begin{lemma}
\label{more-algebra-lemma-lci-local}
设 $R \to S$ 为环映射，$g_1, \ldots, g_m \in S$ 生成单位理想。
若每个 $R \to S_{g_j}$ 都是局部完全交，则 $R \to S$ 也是局部完全交。
\end{lemma}

\begin{proof}
设 $S = R[x_1, \ldots, x_n]/I$ 为一个表现。选取
$h_j \in R[x_1, \ldots, x_n]$，使其在 $S$ 中映到 $g_j$。
则 $R[x_1, \ldots, x_n, x_{n + 1}]/(I, x_{n + 1}h_j - 1)$
是 $S_{g_j}$ 的表现。因此
$I_j = (I, x_{n + 1}h_j - 1)$ 是
$R[x_1, \ldots, x_n, x_{n + 1}]$ 中的 Koszul 正则理想。
选取素理想 $I \subset \mathfrak q \subset R[x_1, \ldots, x_n]$。
则对某个 $j$ 有 $h_j \not \in \mathfrak q$，且
$\mathfrak q_j = (\mathfrak q, x_{n + 1}h_j - 1)$ 是 $V(I_j)$ 中位于
$\mathfrak q$ 之上的素理想。选取 $f_1, \ldots, f_r \in I$，使其映到
$I/I^2 \otimes \kappa(\mathfrak q)$ 的一组基。于是序列
$x_{n + 1}h_j - 1, f_1, \ldots, f_r$ 的各元属于 $I_j$，并映到
$I_j \otimes \kappa(\mathfrak q_j)$ 的一组基；见 Algebra, Lemma
\ref{algebra-lemma-principal-localization-NL}。由 Nakayama 引理，存在
$h \in R[x_1, \ldots, x_n, x_{n + 1}]$，使 $(I_j)_h$ 由
$x_{n + 1}h_j - 1, f_1, \ldots, f_r$ 生成。还可假设 $(I_j)_h$
由某个长度为 $e$ 的 Koszul 正则序列生成。考察
$I_j \otimes \kappa(\mathfrak q_j)$ 的维数，可知 $e = r + 1$。
因此由引理 \ref{more-algebra-lemma-independence-of-generators}，
$x_{n + 1}h_j - 1, f_1, \ldots, f_r$ 是生成 $(I_j)_h$ 的 Koszul 正则序列，
其中某个 $h \in R[x_1, \ldots, x_n, x_{n + 1}]$ 满足
$h \not \in \mathfrak q_j$。由引理 \ref{more-algebra-lemma-truncate-koszul-regular}，
可知对某个 $h' \in R[x_1, \ldots, x_n]$、$h' \not \in \mathfrak q$，
$I_{h'}$ 由 Koszul 正则序列生成，正合所需。
\end{proof}

\begin{lemma}
\label{more-algebra-lemma-relative-global-complete-intersection-koszul}
设 $R$ 为环。若 $R[x_1, \ldots, x_n]/(f_1, \ldots, f_c)$
是相对全局完全交，则 $f_1, \ldots, f_c$ 是 Koszul 正则序列。
\end{lemma}

\begin{proof}
回忆同调群 $H_i(K_\bullet(f_\bullet))$ 被理想 $(f_1, \ldots, f_c)$ 零化。
因此只需证明：对所有包含 $(f_1, \ldots, f_c)$ 的素理想
$\mathfrak q \subset R[x_1, \ldots, x_n]$，
$H_i(K_\bullet(f_\bullet))_\mathfrak q$ 为零。这由 Algebra, Lemma
\ref{algebra-lemma-relative-global-complete-intersection-conormal}
以及正则序列是 Koszul 正则这一事实得出
（引理 \ref{more-algebra-lemma-regular-koszul-regular}）。
\end{proof}

\begin{lemma}
\label{more-algebra-lemma-syntomic-lci}
设 $R \to S$ 为环映射。下列条件等价：
\begin{enumerate}
\item $R \to S$ 是 syntomic 的
（Algebra, Definition \ref{algebra-definition-lci}）；以及
\item $R \to S$ 平坦且为局部完全交。
\end{enumerate}
\end{lemma}

\begin{proof}
假设 (1)。则按定义，$R \to S$ 平坦。由 Algebra, Lemma
\ref{algebra-lemma-syntomic} 与引理 \ref{more-algebra-lemma-lci-local}，
只需证明相对全局完全交同态是局部完全交同态；这正是引理
\ref{more-algebra-lemma-relative-global-complete-intersection-koszul}。

\medskip\noindent
假设 (2)。局部完全交为有限表现的，因为 Koszul 正则理想有限生成。
设 $R \to k$ 为到一个域的映射。只需证明
$S' = S \otimes_R k$ 是 $k$ 上的局部完全交；见 Algebra, Definition
\ref{algebra-definition-lci-field}。选取素理想 $\mathfrak q' \subset S'$。
写 $S = R[x_1, \ldots, x_n]/I$。则
$S' = k[x_1, \ldots, x_n]/I'$，其中
$I' \subset k[x_1, \ldots, x_n]$ 是 $I$ 的像。设
$\mathfrak p' \subset k[x_1, \ldots, x_n]$、$\mathfrak q \subset S$ 与
$\mathfrak p \subset R[x_1, \ldots, x_n]$ 为相应的素理想。
由定义 \ref{more-algebra-definition-regular-ideal}，存在
$g \in R[x_1, \ldots, x_n]$、$g \not \in \mathfrak p$ 与
$f_1, \ldots, f_r \in R[x_1, \ldots, x_n]_g$，它们构成生成 $I_g$ 的
Koszul 正则序列。由于 $S$、因而 $S_g$ 在 $R$ 上平坦，
这些元在 $k[x_1, \ldots, x_n]_g$ 中的像
$f'_1, \ldots, f'_r$ 构成生成 $I'_g$ 的 $H_1$ 正则序列；见引理
\ref{more-algebra-lemma-relative-regular-immersion-algebra}。因此，由引理
\ref{more-algebra-lemma-noetherian-finite-all-equivalent}，$f'_1, \ldots, f'_r$
在 $k[x_1, \ldots, x_n]_{\mathfrak p'}$ 中映到生成
$I'_{\mathfrak p'}$ 的正则序列。应用 Algebra, Lemma
\ref{algebra-lemma-lci}，可知对某个 $g' \in S$、
$g' \not \in \mathfrak q'$，$S'_{gg'}$ 是 $k$ 上的全局完全交，正合所需。
\end{proof}

\noindent
对局部完全交 $R \to S$，当 $n \geq 2$ 时有 $H_n(L_{S/R}) = 0$。
由于我们（尚）未定义完整的余切复形，无法陈述并证明这一点，
但可以推出它的一个后果。

\begin{lemma}
\label{more-algebra-lemma-transitive-lci-at-end}
设 $A \to B \to C$ 为环映射。假设 $B \to C$ 是局部完全交同态。
选取核为 $I$ 的表现 $\alpha : A[x_s, s \in S] \to B$，
以及核为 $J$ 的表现 $\beta : B[y_1, \ldots, y_m] \to C$。
设 $\gamma : A[x_s, y_t] \to C$ 为所诱导的 $C$ 的表现，其核为 $K$。
则得到典范交换图
$$
\xymatrix{
0 \ar[r] &
\Omega_{A[x_s]/A} \otimes C \ar[r] &
\Omega_{A[x_s, y_t]/A} \otimes C \ar[r] &
\Omega_{B[y_t]/B} \otimes C \ar[r] &
0 \\
0 \ar[r] &
I/I^2 \otimes C \ar[r] \ar[u] &
K/K^2 \ar[r] \ar[u] &
J/J^2 \ar[r] \ar[u] &
0
}
$$
其两行均正合。特别地，Algebra, Lemma
\ref{algebra-lemma-exact-sequence-NL} 中的六项正合列可在左侧补上一个零；
即序列
$$
0 \to H_1(\NL_{B/A} \otimes_B C) \to
H_1(L_{C/A}) \to
H_1(L_{C/B}) \to
\Omega_{B/A} \otimes_B C \to
\Omega_{C/A} \to
\Omega_{C/B} \to 0
$$
正合。
\end{lemma}

\begin{proof}
唯一需要证明的是映射 $I/I^2 \otimes C \to K/K^2$ 的单射性。
由假设，理想 $J$ Koszul 正则。因此由引理
\ref{more-algebra-lemma-conormal-sequence-H1-regular-ideal}，有
$IA[x_s, y_j] \cap K^2 = IK$。这意味着 $K/K^2 \to J/J^2$ 的核同构于
$IA[x_s, y_j]/IK$。由于张量积右正合，
$I/I^2 \otimes_A C = IA[x_s, y_j]/IK$，这便给出了所需的
$I/I^2 \otimes_A C \to K/K^2$ 的单射性。
\end{proof}

\begin{lemma}
\label{more-algebra-lemma-transitive-colimit-lci-at-end}
设 $A \to B \to C$ 为环映射。若 $B \to C$ 是局部完全交同态的滤过余极限，
则引理 \ref{more-algebra-lemma-transitive-lci-at-end} 的结论仍成立。
\end{lemma}

\begin{proof}
这由引理 \ref{more-algebra-lemma-transitive-lci-at-end} 与 Algebra, Lemma
\ref{algebra-lemma-colimits-NL} 得出。
\end{proof}

\begin{lemma}
\label{more-algebra-lemma-henselization-NL}
设 $A \to B$ 为局部环的局部同态。设 $A^h \to B^h$（相应地，
$A^{sh} \to B^{sh}$）为所诱导的 Hensel 化（相应地，严格 Hensel 化）之间的映射
（Algebra, Lemma \ref{algebra-lemma-henselian-functorial}，相应地，
Lemma \ref{algebra-lemma-strictly-henselian-functorial}）。则
$\NL_{B/A} \otimes_B B^h \to \NL_{B^h/A^h}$ 与
$\NL_{B/A} \otimes_B B^{sh} \to \NL_{B^{sh}/A^{sh}}$
在上同调群上诱导同构。
\end{lemma}

\begin{proof}
因为 $A^h$ 是 $A$ 上 étale 代数的滤过余极限，由 Algebra, Lemma
\ref{algebra-lemma-colimits-NL} 与 Algebra, Definition
\ref{algebra-definition-etale}，可知 $\NL_{A^h/A}$ 是非循环复形。
$B^h/B$ 也有同样性质。对 $A \to A^h \to B^h$ 使用 Jacobi--Zariski 序列
（Algebra, Lemma \ref{algebra-lemma-exact-sequence-NL}），可知
$\NL_{B^h/A} \to \NL_{B^h/A^h}$ 在上同调群上诱导同构。
此外，étale 环映射是局部完全交，因为它甚至是全局完全交；见
Algebra, Lemma \ref{algebra-lemma-etale-standard-smooth}。
由引理 \ref{more-algebra-lemma-transitive-colimit-lci-at-end}，得到与
$A \to B \to B^h$ 相伴的六项正合 Jacobi--Zariski 序列，它证明
$\NL_{B/A} \otimes_B B^h \to \NL_{B^h/A}$ 在上同调群上诱导同构。
这完成 Hensel 化之间映射的情形。严格 Hensel 化的情形以完全相同的方式证明。
\end{proof}







\section{Cartier 等式与几何正则性}
\label{more-algebra-section-cartier-equality}

\noindent
本节及下一节的参考文献是 \cite[Section 39]{MatCA}。
为顺利阅读本节，读者应熟悉朴素余切复形及其性质；见 Algebra, Section
\ref{algebra-section-netherlander}。

\begin{lemma}[Cartier 等式]
\label{more-algebra-lemma-cartier-equality}
设 $K/k$ 为有限生成域扩张。则 $\Omega_{K/k}$ 与 $H_1(L_{K/k})$
都是有限维的，并且
$\text{trdeg}_k(K) = \dim_K \Omega_{K/k} - \dim_K H_1(L_{K/k})$。
\end{lemma}

\begin{proof}
可以找到 $k$ 上的全局完全交
$A = k[x_1, \ldots, x_n]/(f_1, \ldots, f_c)$，使 $K$ 同构于 $A$ 的分式域；
见 Algebra, Lemma \ref{algebra-lemma-colimit-syntomic} 及其证明。
在此情形下，由 Algebra, Lemmas \ref{algebra-lemma-NL-homotopy} 与
\ref{algebra-lemma-localize-NL}，可知 $\NL_{K/k}$ 同伦等价于复形
$$
\bigoplus\nolimits_{j = 1, \ldots, c} K \longrightarrow
\bigoplus\nolimits_{i = 1, \ldots, n} K\text{d}x_i
$$
$K$ 在 $k$ 上的超越次数是 $A$ 的维数
（由 Algebra, Lemma \ref{algebra-lemma-dimension-prime-polynomial-ring}），
即 $n - c$，故结论成立。
\end{proof}

\begin{lemma}
\label{more-algebra-lemma-transitivity-gamma}
设 $M/L/K$ 为域扩张。则 Jacobi--Zariski 序列
$$
0 \to H_1(L_{L/K}) \otimes_L M \to
H_1(L_{M/K}) \to
H_1(L_{M/L}) \to
\Omega_{L/K} \otimes_L M \to
\Omega_{M/K} \to
\Omega_{M/L} \to 0
$$
正合。
\end{lemma}

\begin{proof}
结合引理 \ref{more-algebra-lemma-transitive-colimit-lci-at-end} 与 Algebra, Lemma
\ref{algebra-lemma-colimit-syntomic} 即得。
\end{proof}

\begin{lemma}
\label{more-algebra-lemma-gamma-commutative-diagram}
给定域的交换图
$$
\xymatrix{
K \ar[r] & K' \\
k \ar[u] \ar[r] & k' \ar[u]
}
$$
其中 $k'/k$ 与 $K'/K$ 是有限生成域扩张，则映射
$$
\alpha : \Omega_{K/k} \otimes_K K' \to \Omega_{K'/k'}
\quad\text{且}\quad
\beta : H_1(L_{K/k}) \otimes_K K' \to H_1(L_{K'/k'})
$$
的核与余核均为有限维的，并且
$$
\dim \Ker(\alpha) - \dim \Coker(\alpha)
-\dim \Ker(\beta) + \dim \Coker(\beta)
=
\text{trdeg}_k(k') - \text{trdeg}_K(K')
$$
\end{lemma}

\begin{proof}
$k \subset k' \subset K'$ 与 $k \subset K \subset K'$ 的
Jacobi--Zariski 序列分别为
$$
0 \to H_1(L_{k'/k}) \otimes K'  \to
H_1(L_{K'/k}) \to
H_1(L_{K'/k'}) \to
\Omega_{k'/k} \otimes K' \to
\Omega_{K'/k} \to
\Omega_{K'/k'} \to 0
$$
和
$$
0 \to H_1(L_{K/k}) \otimes K' \to
H_1(L_{K'/k}) \to
H_1(L_{K'/K}) \to
\Omega_{K/k} \otimes K' \to
\Omega_{K'/k} \to
\Omega_{K'/K} \to 0
$$
由引理 \ref{more-algebra-lemma-cartier-equality}，向量空间
$\Omega_{k'/k}$、$\Omega_{K'/K}$、$H_1(L_{K'/K})$ 与
$H_1(L_{k'/k})$ 均为有限维的，且其维数的交错和为
$\text{trdeg}_k(k') - \text{trdeg}_K(K')$。引理得证。
\end{proof}





\section{几何正则性}
\label{more-algebra-section-geometrically-regular}

\noindent
设 $k$ 为域，$(A, \mathfrak m, K)$ 为 Noether 局部 $k$-代数。
Jacobi--Zariski 序列（Algebra, Lemma
\ref{algebra-lemma-exact-sequence-NL}）是典范正合列
$$
H_1(L_{K/k}) \to \mathfrak m/\mathfrak m^2 \to \Omega_{A/k} \otimes_A K \to
\Omega_{K/k} \to 0
$$
这是因为由 Algebra, Lemma \ref{algebra-lemma-NL-surjection}，
$H_1(L_{K/A}) = \mathfrak m/\mathfrak m^2$。
我们将证明，该序列左端的正合性刻画了正则局部环 $A$ 是否在 $k$ 上几何正则。
在第 \ref{more-algebra-section-regular-fs} 节中，我们会把它与形式光滑的概念联系起来。

\begin{proposition}
\label{more-algebra-proposition-characterization-geometrically-regular}
设 $k$ 为特征 $p > 0$ 的域，$(A, \mathfrak m, K)$ 为 Noether 局部
$k$-代数。下列条件等价：
\begin{enumerate}
\item $A$ 在 $k$ 上几何正则；
\item 对所有在 $k$ 上有限的 $k \subset k' \subset k^{1/p}$，
环 $A \otimes_k k'$ 正则；
\item $A$ 正则，且典范映射
$H_1(L_{K/k}) \to \mathfrak m/\mathfrak m^2$ 是单射；以及
\item $A$ 正则，且映射
$\Omega_{k/\mathbf{F}_p} \otimes_k K \to \Omega_{A/\mathbf{F}_p} \otimes_A K$
是单射。
\end{enumerate}
\end{proposition}

\begin{proof}
(3) $\Rightarrow$ (1) 的证明。假设 (3)。设 $k'/k$ 为有限纯不可分扩张。
令 $A' = A \otimes_k k'$。这是极大理想为 $\mathfrak m'$ 的局部环。
令 $K' = A'/\mathfrak m'$。得到交换图
{\small
$$
\xymatrix{
0 \ar[r] &
H_1(L_{K/k}) \otimes K' \ar[r] \ar[d]_\beta &
\mathfrak m/\mathfrak m^2 \otimes K' \ar[r] \ar[d] &
\Omega_{A/k} \otimes_A K' \ar[r] \ar[d]_{\cong} &
\Omega_{K/k} \otimes K' \ar[r] \ar[d]_\alpha &
0
\\
 &
H_1(L_{K'/k'}) \ar[r] &
\mathfrak m'/(\mathfrak m')^2 \ar[r] &
\Omega_{A'/k'} \otimes_{A'} K' \ar[r] &
\Omega_{K'/k'} \ar[r] &
0
}
$$
}
其两行均正合。由微分模的基变换
（Algebra, Lemma \ref{algebra-lemma-differentials-base-change}），
第三个竖直箭头是同构。因此 $\alpha$ 是满射。由引理
\ref{more-algebra-lemma-gamma-commutative-diagram}，有
$$
\dim \Ker(\alpha) - \dim \Ker(\beta) + \dim \Coker(\beta) = 0
$$
（且这些维数均有限）。追图可得
$\dim \mathfrak m'/(\mathfrak m')^2 \leq \dim \mathfrak m/\mathfrak m^2$。
另一方面，因为 $A \to A'$ 有限平坦，由 Algebra, Lemma
\ref{algebra-lemma-dimension-base-fibre-total}，可知
$\dim(A) = \dim(A')$。故按定义，$A'$ 正则。

\medskip\noindent
(3) 与 (4) 的等价性。考虑交换图各行的 Jacobi--Zariski 序列
$$
\xymatrix{
\mathbf{F}_p \ar[r] & A \ar[r] & K \\
\mathbf{F}_p \ar[r] \ar[u] & k \ar[r] \ar[u] & K \ar[u]
}
$$
得到交换图
$$
\xymatrix{
0 \ar[r] &
\mathfrak m/\mathfrak m^2 \ar[r] &
\Omega_{A/\mathbf{F}_p} \otimes_A K \ar[r] &
\Omega_{K/\mathbf{F}_p} \ar[r] & 0 & \\
0 \ar[r] &
H_1(L_{K/k}) \ar[r] \ar[u] &
\Omega_{k/\mathbf{F}_p} \otimes_k K \ar[r] \ar[u] &
\Omega_{K/\mathbf{F}_p} \ar[r] \ar[u] &
\Omega_{K/k} \ar[r] \ar[u] &
0
}
$$
其两行均正合。这里使用了
$H_1(L_{K/A}) = \mathfrak m/\mathfrak m^2$，以及
$H_1(L_{K/\mathbf{F}_p}) = 0$；后者成立是因为 $K/\mathbf{F}_p$ 可分，
见 Algebra, Proposition
\ref{algebra-proposition-characterize-separable-field-extensions}。
因此显然，映射 $H_1(L_{K/k}) \to \mathfrak m/\mathfrak m^2$ 与
$\Omega_{k/\mathbf{F}_p} \otimes_k K \to \Omega_{A/\mathbf{F}_p} \otimes_A K$
的核具有相同维数。

\medskip\noindent
沿用 Faltings 的方法证明 (2) $\Rightarrow$ (4)；见
\cite{Faltings-einfacher}。设 $a_1, \ldots, a_n \in k$ 为诸元，使
$\text{d}a_1, \ldots, \text{d}a_n$ 在 $\Omega_{k/\mathbf{F}_p}$ 中线性无关。
考虑域扩张 $k' = k(a_1^{1/p}, \ldots, a_n^{1/p})$。由 Algebra, Lemma
\ref{algebra-lemma-size-extension-pth-roots}，可知
$k' = k[x_1, \ldots, x_n]/(x_1^p - a_1, \ldots, x_n^p - a_n)$。
特别地，可知 $k'/k$ 的朴素余切复形同伦于具有零微分的复形
$\bigoplus_{j = 1, \ldots, n} k' \rightarrow \bigoplus_{i = 1, \ldots, n} k'$，
因为在 $\Omega_{k[x_1, \ldots, x_n]/k}$ 中
$\text{d}(x_j^p - a_j) = 0$。如上令
$A' = A \otimes_k k'$ 与 $K' = A'/\mathfrak m'$。
由 Algebra, Lemma \ref{algebra-lemma-change-base-NL}，
$\NL_{A'/A}$ 同伦等价于具有零微分的复形
$\bigoplus_{j = 1, \ldots, n} A' \rightarrow \bigoplus_{i = 1, \ldots, n} A'$；
即 $H_1(L_{A'/A})$ 与 $\Omega_{A'/A}$ 都是秩为 $n$ 的自由模。
$\mathbf{F}_p \to A \to A'$ 的 Jacobi--Zariski 序列为
$$
H_1(L_{A'/A}) \to \Omega_{A/\mathbf{F}_p} \otimes_A A'
\to \Omega_{A'/\mathbf{F}_p} \to \Omega_{A'/A} \to 0
$$
利用核为 $(x_j^p - a_j)$ 的表现 $A[x_1, \ldots, x_n] \to A'$，
展开 Algebra, Lemma \ref{algebra-lemma-exact-sequence-NL} 中的映射，
可知 $H_1(L_{A'/A})$ 的第 $j$ 个基向量映到
$\Omega_{A/\mathbf{F}_p} \otimes A'$ 中的 $\text{d}a_j \otimes 1$。
因为 $\Omega_{A'/A}$ 自由（因而平坦），与 $K'$ 张量后得到正合列
$$
K'^{\oplus n} \to \Omega_{A/\mathbf{F}_p} \otimes_A K'
\xrightarrow{\beta} \Omega_{A'/\mathbf{F}_p} \otimes_{A'} K' \to
K'^{\oplus n} \to 0
$$
由此可知诸元 $\text{d}a_j \otimes 1$ 生成 $\Ker(\beta)$；
还需证明它们线性无关，即证明 $\dim(\Ker(\beta)) = n$。
考虑下列大图
$$
\xymatrix{
0 \ar[r] &
\mathfrak m'/(\mathfrak m')^2 \ar[r] &
\Omega_{A'/\mathbf{F}_p} \otimes K' \ar[r] &
\Omega_{K'/\mathbf{F}_p} \ar[r] & 0 \\
0 \ar[r] &
\mathfrak m/\mathfrak m^2 \otimes K' \ar[r] \ar[u]^\alpha &
\Omega_{A/\mathbf{F}_p} \otimes K' \ar[r] \ar[u]^\beta &
\Omega_{K/\mathbf{F}_p} \otimes K' \ar[r] \ar[u]^\gamma & 0
}
$$
由引理 \ref{more-algebra-lemma-cartier-equality} 以及
$\mathbf{F}_p \to K \to K'$ 的 Jacobi--Zariski 序列，
$\gamma$ 的核与余核具有相同的有限维数。由假设，$A'$ 正则
（且如上所见，其维数与 $A$ 相同），故 $\alpha$ 的核与余核具有相同维数。
从而 $\beta$ 的核与余核具有相同维数，这正是所需证明的。

\medskip\noindent
推含 (1) $\Rightarrow$ (2) 是平凡的。命题的证明完成。
\end{proof}

\begin{lemma}
\label{more-algebra-lemma-geometrically-regular-over-field}
设 $k$ 为特征 $p > 0$ 的域，$(A, \mathfrak m, K)$ 为 Noether 局部
$k$-代数。假设 $A$ 在 $k$ 上几何正则。设 $K/F/k$ 为有限生成子扩张。
设 $\varphi : k[y_1, \ldots, y_m] \to A$ 为 $k$-代数映射，使得
$y_i$ 在 $K$ 中映到 $F$ 的元，并且
$\text{d}y_1, \ldots, \text{d}y_m$ 映到 $\Omega_{F/k}$ 的一组基。
令 $\mathfrak p = \varphi^{-1}(\mathfrak m)$。则
$$
k[y_1, \ldots, y_m]_\mathfrak p \to A
$$
平坦，且 $A/\mathfrak pA$ 正则。
\end{lemma}

\begin{proof}
令 $A_0 = k[y_1, \ldots, y_m]_\mathfrak p$，其极大理想为
$\mathfrak m_0$，剩余域为 $K_0$。注意 $\Omega_{A_0/k}$ 是秩为 $m$ 的自由模，
且 $\Omega_{A_0/k} \otimes K_0 \to \Omega_{K_0/k}$ 是同构。
显然 $A_0$ 在 $k$ 上几何正则。因此
$H_1(L_{K_0/k}) \to \mathfrak m_0/\mathfrak m_0^2$ 是同构；见命题
\ref{more-algebra-proposition-characterization-geometrically-regular}。现在考虑
$$
\xymatrix{
H_1(L_{K_0/k}) \otimes K \ar[d] \ar[r] &
\mathfrak m_0/\mathfrak m_0^2 \otimes K \ar[d] \\
H_1(L_{K/k}) \ar[r] & \mathfrak m/\mathfrak m^2
}
$$
由引理 \ref{more-algebra-lemma-transitivity-gamma}，左侧竖直箭头是单射；
由命题 \ref{more-algebra-proposition-characterization-geometrically-regular}，
下方水平箭头是单射，故右侧竖直箭头也是单射。
因此 $A_0$ 中的一个正则参数系映到 $A$ 中正则参数系的一部分。
由 Algebra, Lemmas \ref{algebra-lemma-flat-over-regular} 与
\ref{algebra-lemma-regular-ring-CM} 即得结论。
\end{proof}













\section{拓扑环与拓扑模}
\label{more-algebra-section-topological-ring}

\noindent
我们简要讨论拓扑 Abel 群的一些性质。若 Abel 群 $M$ 带有一个拓扑，使加法
$M \times M \to M$、$(x, y) \mapsto x + y$ 与取逆
$M \to M$、$x \mapsto -x$ 连续，则称 $M$ 为{\it 拓扑 Abel 群}。
{\it 拓扑 Abel 群的同态}就是连续的 Abel 群同态。
交换拓扑群的范畴是加性范畴并且有核与余核，但不是 Abel 范畴
（因为公理 $\Im = \Coim$ 不成立）。若 $N \subset M$ 是子群，
我们也把 $N$ 与 $M/N$ 视为拓扑群：在 $N$ 上使用诱导拓扑，
在 $M/N$ 上使用商拓扑（即使 $M \to M/N$ 为商映射的拓扑）。
注意，若 $N \subset M$ 是开子群，则 $M/N$ 上的拓扑是离散的。

\medskip\noindent
若存在一个由子群组成的 $0$ 的邻域基，则称 $M$ 上的拓扑为{\it 线性的}。
此时这些子群也是开的。下面是一个例子。设 $I$ 为有向集，且
$G_i$ 是 $I$ 上（离散）Abel 群的逆系统。则
$$
G = \lim_{i \in I} G_i
$$
赋予逆极限拓扑后是线性拓扑化的；由 $\Ker(G \to G_i)$ 给出
$0$ 的一组邻域基。反之，设 $M$ 是线性拓扑化 Abel 群。
任取开子群的一个基本系统 $U_i \subset M$、$i \in I$
（即诸 $U_i$ 构成开邻域的基本系统，且每个 $U_i$ 都是 $M$ 的子群）。
令 $i \geq i' \Leftrightarrow U_i \subset U_{i'}$，可知 $I$ 是有向集。
得到线性拓扑化 Abel 群的同态
$$
c : M \longrightarrow \lim_{i \in I} M/U_i.
$$
显然，$M$（作为拓扑空间）{\it 分离}当且仅当 $c$ 是单射。
若 $c$ 是同构，则称 $M$ {\it 完备}\footnote{我们把分离性纳入完备性的定义，
因为希望完备群中的极限唯一。任意拓扑群都有一个完备性定义；除分离性问题外，
该定义在我们的特殊情形中与这里的定义一致。}。
留给读者验证，这一条件与上面所选的开子群基本系统
$\{U_i\}_{i \in I}$ 无关。事实上，拓扑 Abel 群
$M^\wedge = \lim_{i \in I} M/U_i$ 与这一选择无关，有时称为 $M$ 的
{\it 完备化}。任何上述 $G = \lim G_i$ 都是完备的；特别地，完备化
$M^\wedge$ 总是完备的。

\begin{definition}[拓扑环]
\label{more-algebra-definition-topological-ring}
\begin{reference}
\cite[Chapter 0, Sections 7.1 and 7.2]{EGA1}
\end{reference}
设 $R$ 为环，$M$ 为 $R$-模。
\begin{enumerate}
\item 若 $R$ 带有一个拓扑，使得加法与乘法作为映射
$R \times R \to R$ 均连续，其中 $R \times R$ 带乘积拓扑，
则称 $R$ 为{\it 拓扑环}。在此情形下，若 $M$ 带有一个拓扑，使加法
$M \times M \to M$ 与标量乘法 $R \times M \to M$ 连续，
则称 $M$ 为{\it 拓扑模}。
\item {\it 拓扑模的同态}就是连续的 $R$-模映射。
{\it 拓扑环的同态}是对给定拓扑连续的环同态。
\item 若 $0$ 有一个由子模组成的邻域基，则称 $M$ {\it 线性拓扑化}。
若 $0$ 有一个由理想组成的邻域基，则称 $R$ {\it 线性拓扑化}。
\item 若 $R$ 线性拓扑化，且 $I \subset R$ 是开理想，并且 $0$ 的每个邻域
都包含 $I^n$（对某个 $n$），则称 $I$ 为{\it 定义理想}。
\item 若 $R$ 线性拓扑化，则当 $R$ 有定义理想时，称 $R$ {\it 预容许}。
\item 若 $R$ 线性拓扑化、预容许且完备，则称 $R$ {\it 容许}\footnote{
依照我们的约定，这包括分离性。}。
\item 若 $R$ 线性拓扑化，且存在定义理想 $I$，使
$\{I^n\}_{n \geq 0}$ 构成 $0$ 的一组邻域基，则称 $R$ {\it 预 adic}。
\item 若 $R$ 线性拓扑化，则当 $R$ 预 adic 且完备时，称 $R$ {\it adic}。
\end{enumerate}
注意，（预）adic 拓扑环等价于这样一个（预）容许拓扑环：它有定义理想
$I$，且对所有 $n \geq 1$，$I^n$ 都是开的。
\end{definition}

\noindent
设 $R$ 为环，$M$ 为 $R$-模，$I \subset R$ 为理想。可以在 $R$ 上考虑
以 $\{I^n\}_{n \geq 0}$ 为 $0$ 的一组邻域基的线性拓扑。
该拓扑称为 $I$-adic 拓扑；$R$ 在 $I$-adic 拓扑下是预 adic 拓扑环\footnote{
因此，$I$-adic 拓扑有时称为 $I$-预 adic 拓扑。}。此外，在 $M$ 上以
$\{I^nM\}_{n \geq 0}$ 为 $0$ 的一组开邻域基的线性拓扑，使 $M$ 成为拓扑
$R$-模。这称为 $M$ 上的{\it $I$-adic 拓扑}。可知，$M$ 是 $I$-adic 完备的
（定义见 Algebra, Definition \ref{algebra-definition-complete}），当且仅当
$M$ 在 $I$-adic 拓扑下完备\footnote{可能出现如下情形：$I$-adic 完备化
$M^\wedge$ 并非 $I$-adic 完备，尽管 $M^\wedge$ 对极限拓扑总是完备的。
若 $I$ 有限生成，则 $M^\wedge$ 上的 $I$-adic 拓扑与极限拓扑一致；
见 Algebra, Lemma \ref{algebra-lemma-hathat-finitely-generated} 及其证明。}。
特别地，$R$ 是 $I$-adic 完备的，当且仅当 $R$ 在 $I$-adic 拓扑下是
adic 拓扑环。

\medskip\noindent
作为特殊情形，注意离散拓扑是 $0$-adic 拓扑，并且任意带离散拓扑的环都是 adic 的。

\begin{lemma}
\label{more-algebra-lemma-continuous}
设 $\varphi : R \to S$ 为环映射。设 $I \subset R$ 与 $J \subset S$ 为理想，
分别在 $R$ 与 $S$ 上赋予 $I$-adic 与 $J$-adic 拓扑。
则 $\varphi$ 是拓扑环同态，当且仅当对某个 $n \geq 1$ 有
$\varphi(I^n) \subset J$。
\end{lemma}

\begin{proof}
略。
\end{proof}

\begin{lemma}[Baire 纲定理]
\label{more-algebra-lemma-baire-category-complete-module}
设 $M$ 为拓扑 Abel 群。假设 $M$ 线性拓扑化、完备，并且有可数的
$0$ 的邻域基本系统。若 $U_n \subset M$、$n \geq 1$ 是开稠密子集，
则 $\bigcap_{n \geq 1} U_n$ 稠密。
\end{lemma}

\begin{proof}
设 $U_n$ 如引理陈述。以 $U_1 \cap \ldots \cap U_n$ 替换 $U_n$ 后，
可假设 $U_1 \supset U_2 \supset \ldots$。设
$M_n$、$n \in \mathbf{N}$ 为 $0$ 的邻域基本系统，可假设
$M_{n + 1} \subset M_n$。选取 $x \in M$。我们将证明：对每个
$k \geq 1$，存在 $y \in \bigcap_{n \geq 1} U_n$ 使 $x - y \in M_k$。

\medskip\noindent
如下构造 $y$。首先选取 $y_1 \in U_1$，使 $y_1 \in x + M_k$。
这是可能的，因为 $U_1$ 稠密，而 $x + M_k$ 是开的。
然后选取 $k_1 > k$，使 $y_1 + M_{k_1} \subset U_1$。
这是可能的，因为 $U_1$ 是开的。接着选取 $y_2 \in U_2$，使
$y_2 \in y_1 + M_{k_1}$。这是可能的，因为 $U_2$ 稠密，而
$y_2 + M_{k_1}$ 是开的。然后选取 $k_2 > k_1$，使
$y_2 + M_{k_2} \subset U_2$。这是可能的，因为 $U_2$ 是开的。

\medskip\noindent
如此继续，得到 $M$ 中收敛于 $y$ 的元的序列 $y_i$。
由构造，$x - y \in M_k$。由于
$$
y - y_i = (y_{i + 1} - y_i) + (y_{i + 2} - y_{i + 1}) + \ldots
$$
属于 $M_{k_i}$，故对所有 $i$，有
$y \in y_i + M_{k_i} \subset U_i$，正合所需。
\end{proof}

\begin{lemma}
\label{more-algebra-lemma-consequence-baire-complete-module}
在引理 \ref{more-algebra-lemma-baire-category-complete-module} 的相同假设下，
若对某些闭子群 $N_n$ 有 $M = \bigcup_{n \geq 1} N_n$，
则对某个 $n$，$N_n$ 是开的。
\end{lemma}

\begin{proof}
若不然，则对所有 $n$，$U_n = M \setminus N_n$ 都稠密，
这与引理 \ref{more-algebra-lemma-baire-category-complete-module} 矛盾。
\end{proof}

\begin{lemma}[开映射引理]
\label{more-algebra-lemma-open-mapping}
设 $u : N \to M$ 为拓扑化 Abel 群的连续映射。假设 $M$ 分离，
而 $N$ 完备、线性拓扑化，并且有可数的 $0$ 的邻域基本系统。
则下列两种情形恰有一种成立：
\begin{enumerate}
\item $u$ 是开映射；或者
\item 对某个开子群 $N' \subset N$，其像 $u(N')$ 在 $M$ 中无处稠密。
\end{enumerate}
\end{lemma}

\begin{proof}
设 $N_n$、$n \in \mathbf{N}$ 为 $0$ 的邻域基本系统。
可假设 $N_n$ 是子群，且 $N_{n + 1} \subset N_n$。若 (2) 不成立，
则对 $n = 1, 2, 3, \ldots$，$u(N_n)$ 的闭包 $M_n$ 是开子群。
因为 $u$ 连续，可知 $M_n$、$n \in \mathbf{N}$ 必构成 $M$ 中
$0$ 的开邻域基本系统。又因为 $M_n$ 是 $u(N_n)$ 的闭包，所以对所有
$n \geq 1$ 有
$$
u(N_n) + M_{n + 1} = M_n
$$
选取 $x_1 \in M_1$。于是可归纳选取 $y_i \in N_i$ 与
$x_{i + 1} \in M_{i + 1}$，使
$$
u(y_i) + x_{i + 1} = x_i
$$
由于 $N$ 完备，$N$ 中的元 $y = y_1 + y_2 + y_3 + \ldots$ 存在。
又因 $M$ 分离，便有 $x_1 = u(y)$。因此 $M_1 = u(N_1)$。
读者以完全相同的方式可证明对所有 $i \geq 2$，有
$M_i = u(N_i)$，从而 $u$ 是开映射。
\end{proof}









\section{拓扑环的形式光滑映射}
\label{more-algebra-section-formally-smooth}

\noindent
形式光滑性有一个适用于拓扑环同态的版本。

\begin{definition}
\label{more-algebra-definition-formally-smooth}
设 $R \to S$ 为拓扑环同态，且 $R$ 与 $S$ 线性拓扑化。
若对拓扑环同态的每一个实线交换图
$$
\xymatrix{
S \ar[r] \ar@{-->}[rd] & A/J \\
R \ar[r] \ar[u] & A \ar[u]
}
$$
其中 $A$ 为离散环，$J \subset A$ 是平方为零的理想，都存在使图交换的虚线箭头，
则称 $S$ 在 $R$ 上{\it 形式光滑}。
\end{definition}

\noindent
我们主要在给定理想 $\mathfrak m \subset R$ 与 $\mathfrak n \subset S$ 时使用此概念：
在 $R$ 上赋予 $\mathfrak m$-adic 拓扑，在 $S$ 上赋予 $\mathfrak n$-adic 拓扑。
$\varphi : R \to S$ 连续，当且仅当对某个 $m \geq 1$ 有
$\varphi(\mathfrak m^m) \subset  \mathfrak n$；见引理 \ref{more-algebra-lemma-continuous}。
事实表明，在此情形下只有 $S$ 上的拓扑是相关的。

\begin{lemma}
\label{more-algebra-lemma-formally-smooth}
设 $\varphi : R \to S$ 为环映射。
\begin{enumerate}
\item 若 $R \to S$ 在 Algebra, Definition
\ref{algebra-definition-formally-smooth} 的意义下形式光滑，
则对 $R$ 上任一线性拓扑和 $S$ 上任一预 adic 拓扑，只要
$R \to S$ 连续，$R \to S$ 就是形式光滑的。
\item 设 $\mathfrak n \subset S$ 与 $\mathfrak m \subset R$ 为理想，
并且 $\varphi$ 对 $R$ 上的 $\mathfrak m$-adic 拓扑和 $S$ 上的
$\mathfrak n$-adic 拓扑连续。则下列条件等价：
\begin{enumerate}
\item $\varphi$ 对 $R$ 上的 $\mathfrak m$-adic 拓扑和 $S$ 上的
$\mathfrak n$-adic 拓扑形式光滑；以及
\item $\varphi$ 对 $R$ 上的离散拓扑和 $S$ 上的 $\mathfrak n$-adic 拓扑形式光滑。
\end{enumerate}
\end{enumerate}
\end{lemma}

\begin{proof}
假设 $R \to S$ 在 Algebra, Definition
\ref{algebra-definition-formally-smooth} 的意义下形式光滑。
若 $S$ 有预 adic 拓扑，则存在理想 $\mathfrak n \subset S$，使 $S$
带有 $\mathfrak n$-adic 拓扑。设给定定义 \ref{more-algebra-definition-formally-smooth}
中的实线交换图。映射 $S \to A/J$ 的连续性意味着对某个 $k \geq 1$，
$\mathfrak n^k$ 映到 $A/J$ 中的零；见引理 \ref{more-algebra-lemma-continuous}。
由所假设的 $S$ 在 $R$ 上形式光滑，得到环映射 $\psi : S \to A$。
于是 $\psi(\mathfrak n^k) \subset J$，故因 $J^2 = 0$ 有
$\psi(\mathfrak n^{2k}) = 0$。所以由引理 \ref{more-algebra-lemma-continuous}，
$\psi$ 连续。这证明了 (1)。

\medskip\noindent
(2)(b) $\Rightarrow$ (2)(a) 的证明与 (1) 的证明相同。假设 (2)(a)。
设给定定义 \ref{more-algebra-definition-formally-smooth} 中的实线交换图，
但在 $R$ 上使用离散拓扑。因为 $\varphi$ 连续，可知对某个 $n \geq 1$，
$\varphi(\mathfrak m^n) \subset \mathfrak n$。因为 $S \to A/J$ 连续，
对某个 $k \geq 1$，$\mathfrak n^k$ 映到 $A/J$ 中的零。
故 $R \to A$ 将 $\mathfrak m^{nk}$ 映入 $J$，从而将
$\mathfrak m^{2nk}$ 映到 $A$ 中的零。因此 $R \to A$ 对
$\mathfrak m$-adic 拓扑连续，于是 (2)(a) 给出所需的虚线箭头。
\end{proof}

\begin{definition}
\label{more-algebra-definition-formally-smooth-adic}
设 $R \to S$ 为环映射，$\mathfrak n \subset S$ 为理想。
若引理 \ref{more-algebra-lemma-formally-smooth} 的等价条件 (2)(a) 与 (2)(b) 成立，
则称 $R \to S$ 对{\it $\mathfrak n$-adic 拓扑形式光滑}。
\end{definition}

\noindent
该性质由完备化继承。

\begin{lemma}
\label{more-algebra-lemma-formally-smooth-completion}
设 $(R, \mathfrak m)$ 与 $(S, \mathfrak n)$ 是带有限生成理想的环。
在 $R$ 与 $S$ 上分别赋予 $\mathfrak m$-adic 与 $\mathfrak n$-adic 拓扑。
设 $R \to S$ 为拓扑环同态。下列条件等价：
\begin{enumerate}
\item $R \to S$ 对 $\mathfrak n$-adic 拓扑形式光滑；
\item $R \to S^\wedge$ 对 $\mathfrak n^\wedge$-adic 拓扑形式光滑；
\item $R^\wedge \to S^\wedge$ 对 $\mathfrak n^\wedge$-adic 拓扑形式光滑。
\end{enumerate}
这里 $R^\wedge$ 与 $S^\wedge$ 分别是 $R$ 与 $S$ 的
$\mathfrak m$-adic 与 $\mathfrak n$-adic 完备化。
\end{lemma}

\begin{proof}
$\mathfrak m$ 有限生成这一假设推出：$R^\wedge$ 是
$\mathfrak mR^\wedge$-adic 完备的，$\mathfrak mR^\wedge = \mathfrak m^\wedge$，
并且 $R^\wedge/\mathfrak m^nR^\wedge = R/\mathfrak m^n$；见 Algebra, Lemma
\ref{algebra-lemma-hathat-finitely-generated} 及其证明。对
$(S, \mathfrak n)$ 同理。因此显然，在定义
\ref{more-algebra-definition-formally-smooth} 中，情形 (1)、(2) 与 (3) 的图一一对应。
\end{proof}

\noindent
使用 adic 环的优点是可以得到更强的提升性质。

\begin{lemma}
\label{more-algebra-lemma-lift-continuous}
设 $R \to S$ 为环映射，$\mathfrak n$ 为 $S$ 的理想。
假设 $R \to S$ 对 $\mathfrak n$-adic 拓扑形式光滑。
考虑拓扑环同态的实线交换图
$$
\xymatrix{
S \ar[r]_\psi \ar@{-->}[rd] & A/J \\
R \ar[r] \ar[u] & A \ar[u]
}
$$
其中 $A$ 是 adic 的，$A/J$ 是 $A$ 对闭理想 $J \subset A$ 的商
（作为拓扑环），并且对某个 $t \geq 1$，$J^t$ 包含于 $A$ 的一个定义理想。
则拓扑环范畴中存在使该图交换的虚线箭头。
\end{lemma}

\begin{proof}
设 $I \subset A$ 为定义理想，并使得对某个 $t$ 有 $I \supset J^t$。
则 $A = \lim A/I^n$；又因为假设 $J$ 闭，
$A/J = \lim A/J + I^n$。考虑离散 $R$-代数
$A_{n, m} = A/J^n + I^m$ 的下图：
$$
\xymatrix{
A/J^3 + I^3 \ar[r] \ar[d] &
A/J^2 + I^3 \ar[r] \ar[d] &
A/J + I^3 \ar[d] \\
A/J^3 + I^2 \ar[r] \ar[d] &
A/J^2 + I^2 \ar[r] \ar[d] &
A/J + I^2 \ar[d] \\
A/J^3 + I \ar[r] &
A/J^2 + I \ar[r] &
A/J + I
}
$$
注意，每个交换方块都定义了 $R$-代数的满射
$$
A_{n + 1, m + 1} \longrightarrow A_{n + 1, m} \times_{A_{n, m}} A_{n, m + 1}
$$
其核的平方为零。我们将归纳构造 $R$-代数映射
$\varphi_{n, m} : S \to A_{n, m}$。事实上，已有映射
$\varphi_{1, m} = \psi \bmod J + I^m$。注意，由于 $\psi$ 连续，
这些映射都连续。从选定的 $\varphi_{1, 1}$ 开始，利用 $S$ 在 $R$ 上的
形式光滑性，沿上图的底行向上提升，可归纳选取映射 $\varphi_{n, 1}$。
其余映射 $\varphi_{n, m}$ 按 $n + m$ 归纳构造。具体地，沿上面显示的满射
提升一对 $(\varphi_{n + 1, m}, \varphi_{n, m + 1})$，选取
$\varphi_{n + 1, m + 1}$（再次使用 $S$ 在 $R$ 上的形式光滑性）。
如此所得所有映射 $\varphi_{n, m}$ 都与该系统的转移映射相容。
因为 $J^t \subset I$，例如
$\varphi_n = \varphi_{nt, n} \bmod I^n$ 诱导映射 $S \to A/I^n$。
取极限 $\varphi = \lim \varphi_n$，得到映射
$S \to A = \lim A/I^n$。它与到 $A/J$ 的映射复合后等于 $\psi$，
因为已看到 $A/J = \lim A/J + I^n$。
最后证明 $\varphi$ 连续。由 $\psi$ 是拓扑环的态射这一假设及引理
\ref{more-algebra-lemma-continuous}，可知对某个 $r \geq 1$，
$\psi(\mathfrak n^r) \subset J + I/J$。因此
$\varphi(\mathfrak n^r) \subset J + I$，从而如所需，
$\varphi(\mathfrak n^{rt}) \subset I$。
\end{proof}

\begin{lemma}
\label{more-algebra-lemma-increase-ideal}
设 $R \to S$ 为环映射，$\mathfrak n \subset \mathfrak n' \subset S$ 为理想。
若 $R \to S$ 对 $\mathfrak n$-adic 拓扑形式光滑，
则 $R \to S$ 对 $\mathfrak n'$-adic 拓扑形式光滑。
\end{lemma}

\begin{proof}
略。
\end{proof}

\begin{lemma}
\label{more-algebra-lemma-compose-formally-smooth}
线性拓扑化环之间形式光滑连续同态的复合仍形式光滑。
\end{lemma}

\begin{proof}
略。（提示：这是完全形式的，由考察一个适当的图得出。）
\end{proof}

\begin{lemma}
\label{more-algebra-lemma-base-change-fs}
设 $R$、$S$ 为环，$\mathfrak n \subset S$ 为理想。
设 $R \to S$ 对 $\mathfrak n$-adic 拓扑形式光滑，并设
$R \to R'$ 为任意环映射。则
$R' \to S' = S \otimes_R R'$ 对
$\mathfrak n' = \mathfrak nS'$-adic 拓扑形式光滑。
\end{lemma}

\begin{proof}
给定定义 \ref{more-algebra-definition-formally-smooth} 中的实线图
$$
\xymatrix{
S \ar[r] \ar@{-->}[rrd] & S' \ar[r] \ar@{-->}[rd] & A/J \\
R  \ar[u] \ar[r] & R' \ar[r] \ar[u] & A \ar[u]
}
$$
则复合 $S \to S' \to A/J$ 连续。由假设，较长的虚线箭头存在。
由张量积的泛性质，得到较短的虚线箭头。
\end{proof}

\noindent
我们已经在 Algebra, Lemma \ref{algebra-lemma-descent-formally-smooth}
中见过沿忠实平坦环映射的形式光滑性下降。在当前拓扑环的语境中也有类似结果。
不过，这里只证明下面这个非常简单且易证、但已相当有用的版本。

\begin{lemma}
\label{more-algebra-lemma-descent-fs}
设 $R$、$S$ 为环，$\mathfrak n \subset S$ 为理想。设 $R \to R'$ 为环映射。
令 $S' = S \otimes_R R'$ 且 $\mathfrak n' = \mathfrak nS$。若
\begin{enumerate}
\item 映射 $R \to R'$ 将 $R$ 作为 $R'$ 的 $R$-模直和项嵌入；以及
\item $R' \to S'$ 对 $\mathfrak n'$-adic 拓扑形式光滑，
\end{enumerate}
则 $R \to S$ 对 $\mathfrak n$-adic 拓扑形式光滑。
\end{lemma}

\begin{proof}
给定定义 \ref{more-algebra-definition-formally-smooth} 中的实线图
$$
\xymatrix{
S \ar[r] & A/J \\
R \ar[u] \ar[r] & A \ar[u]
}
$$
令 $A' = A \otimes_R R'$，并令
$J' = \Im(J \otimes_R R' \to A')$。上图的基变换是具有连续箭头的图
$$
\xymatrix{
S' \ar[r] \ar@{-->}[rd]^{\psi'} & A'/J' \\
R' \ar[u] \ar[r] & A' \ar[u]
}
$$
由条件 (2)，得到虚线箭头 $\psi' : S' \to A'$。
利用条件 (1)，选取 $R$-模的直和分解 $R' = R \oplus C$。
（警告：$C$ 不是 $R'$ 中的理想。）于是
$A' = A \oplus A \otimes_R C$。令
$$
J'' = \Im(J \otimes_R C \to A \otimes_R C) \subset J' \subset A'.
$$
于是作为 $A$-模有 $J' = J \oplus J''$。将 $\psi'$ 与
$S \to S'$ 复合所得的 $\psi : S \to A'$，其像包含于
$A + J' = A \oplus J''$。然而，在环 $A + J' = A \oplus J''$ 中，
$A$-子模 $J''$ 是理想！（使用 $J^2 = 0$。）因此复合
$S \to A + J' \to (A + J')/J'' = A$ 就是所求箭头。
\end{proof}







\section{局部环的形式光滑映射}
\label{more-algebra-section-formally-smooth-local}

\noindent
对局部环之间的局部同态，可以限制需要检验提升性质的图。
请与 Algebra, Lemma \ref{algebra-lemma-smooth-test-artinian} 比较。

\begin{lemma}
\label{more-algebra-lemma-fs-local}
设 $(R, \mathfrak m) \to (S, \mathfrak n)$ 为局部环之间的局部同态。
下列条件等价：
\begin{enumerate}
\item $R \to S$ 对 $\mathfrak n$-adic 拓扑形式光滑；
\item 对局部环之间局部同态的每个实线交换图
$$
\xymatrix{
S \ar[r] \ar@{-->}[rd] & A/J \\
R \ar[r] \ar[u] & A \ar[u]
}
$$
其中 $J \subset A$ 是平方为零的理想，对某个 $n > 0$ 有
$\mathfrak m_A^n = 0$，并且 $S \to A/J$ 在剩余域上诱导同构，
都存在使该图交换的虚线箭头。
\end{enumerate}
若 $S$ 是 Noether 环，这些条件还等价于：
\begin{enumerate}
\item[(3)] 与 (2) 相同，但仅要求这样的图：附加地，$A \to A/J$ 是小扩张
（Algebra, Definition \ref{algebra-definition-small-extension}）。
\end{enumerate}
\end{lemma}

\begin{proof}
推含 (1) $\Rightarrow$ (2) 由定义得出。考虑定义
\ref{more-algebra-definition-formally-smooth} 中对 $R$ 上的 $\mathfrak m$-adic 拓扑及
$S$ 上的 $\mathfrak n$-adic 拓扑的一个图
$$
\xymatrix{
S \ar[r] \ar@{-->}[rd] & A/J \\
R \ar[r] \ar[u] & A \ar[u]
}
$$
选取 $m > 0$ 使 $\mathfrak n^m(A/J) = 0$
（由图中映射的连续性，这是可能的）。设 $A'$ 为 $A$ 的子环，
即 $S$ 在 $A/J$ 中的像的逆像。把 $J' = J$ 视为 $A'$ 中的理想。
则 $J'$ 是 $A'$ 中平方为零的理想，且 $A'/J'$ 是 $S/\mathfrak n^m$ 的商。
故 $A'$ 是局部环，并且 $\mathfrak m_{A'}^{2m} = 0$。于是得到 (2) 中的图
$$
\xymatrix{
S \ar[r] \ar@{-->}[rd] & A'/J' \\
R \ar[r] \ar[u] & A' \ar[u]
}
$$
若能在此图中构造虚线箭头，再与 $A' \to A$ 复合，便得到原图中的虚线箭头。
这样便看出 (2) 推出 (1)。

\medskip\noindent
假设 $S$ Noether。推含 (1) $\Rightarrow$ (3) 是立即的。
假设 (3)，并设给定 (2) 中的图。则对某个 $n > 0$ 有
$\mathfrak m_A^n J = 0$。考虑映射
$$
A \to A/\mathfrak m_A^{n - 1}J \to \ldots \to A/\mathfrak mJ \to A/J
$$
可知，只需在 $\mathfrak m_A J = 0$ 时给出提升。
假设 $\mathfrak m_A J = 0$，并令 $A' \subset A$ 为上面构造的环。
则 $A'/J'$ 是 Artin 的，因为它是 Artin 局部环
$S/\mathfrak n^m$ 的商。因此只需证明：给定性质 (3)，
在 (2) 中满足 $A/J$ Artin 且 $\mathfrak m_A J = 0$ 的图中，
可以找到虚线箭头。设 $\kappa$ 为 $A$、$A/J$ 与 $S$ 的共同剩余域。
由 (3)，若 $J_0 \subset J$ 是满足 $\dim_\kappa(J/J_0) = 1$ 的理想，
则可以给出虚线箭头 $S \to A/J_0$。取乘积得到
$$
S \longrightarrow \prod\nolimits_{J_0 \text{ 如上}} A/J_0
$$
显然，该箭头的像包含于这样的元组成的 $R$-子代数 $A'$：这些元映入小对角
$A/J \subset \prod_{J_0} A/J$。设 $J' \subset A'$ 为映到 $A/J$ 中零元的元。
则 $J'$ 是平方为零的理想，并且作为 $\kappa$-向量空间等于
$$
J' = \prod\nolimits_{J_0 \text{ 如上}} J/J_0
$$
因此映射 $J \to J'$ 是单射。由向量空间理论，可选取分裂
$J' = J \oplus M$。于是作为 $R$-代数有
$$
A' = A \oplus M
$$
故可将映射 $S \to A'$ 与投影 $A' \to A$ 复合，得到所需的虚线箭头，
从而完成引理的证明。
\end{proof}

\noindent
下列引理将在第 \ref{more-algebra-section-regular-fs} 节中得到改进。

\begin{lemma}
\label{more-algebra-lemma-fs-implies-regular}
设 $k$ 为域，$(A, \mathfrak m, K)$ 为 Noether 局部 $k$-代数。
若 $k \to A$ 对 $\mathfrak m$-adic 拓扑形式光滑，
则 $A$ 是正则局部环。
\end{lemma}

\begin{proof}
设 $k_0 \subset k$ 为素域。则 $k_0$ 完美，故 $k / k_0$ 可分，
从而由 Algebra, Lemma \ref{algebra-lemma-formally-smooth-extensions-easy}
形式光滑。由引理 \ref{more-algebra-lemma-formally-smooth} 与
\ref{more-algebra-lemma-compose-formally-smooth}，可知
$k_0 \to A$ 对 $A$ 上的 $\mathfrak m$-adic 拓扑形式光滑。
因此可假设 $k = \mathbf{Q}$ 或 $k = \mathbf{F}_p$。

\medskip\noindent
由 Algebra, Lemmas \ref{algebra-lemma-completion-faithfully-flat} 与
\ref{algebra-lemma-flat-under-regular}，只需证明完备化 $A^\wedge$ 正则。
由引理 \ref{more-algebra-lemma-formally-smooth-completion}，可用 $A^\wedge$ 替换 $A$。
因此可假设 $A$ 是 Noether 完备局部环。由 Cohen 结构定理
（Algebra, Theorem \ref{algebra-theorem-cohen-structure-theorem}），
存在映射 $K \to A$。由于 $k$ 是素域，$K \to A$ 是 $k$-代数映射。

\medskip\noindent
设 $x_1, \ldots, x_n \in \mathfrak m$ 为诸元，其像构成
$\mathfrak m/\mathfrak m^2$ 的一组基。令 $T = K[[X_1, \ldots, X_n]]$。
注意
$$
A/\mathfrak m^2 \cong K[x_1, \ldots, x_n]/(x_ix_j)
$$
以及
$$
T/\mathfrak m_T^2 \cong K[X_1, \ldots, X_n]/(X_iX_j).
$$
令 $A/\mathfrak m^2 \to T/m_T^2$ 为局部 $K$-代数同构，
它将 $x_i$ 的类映到 $X_i$ 的类。以
$f_1 : A \to T/\mathfrak m_T^2$ 表示该同构与商映射
$A \to A/\mathfrak m^2$ 的复合。$k \to A$ 对 $\mathfrak m$-adic 拓扑
形式光滑这一假设意味着，可以把 $f_1$ 提升为映射
$f_2 : A \to T/\mathfrak{m}_T^3$，再提升为映射
$f_3 : A \to T/\mathfrak{m}_T^4$，如此继续，对所有 $n \geq 1$。
警告：映射 $f_n$ 是连续 $k$-代数映射，但未必是 $K$-代数映射。
得到局部 $k$-代数的诱导映射
$f : A \to T = \lim T/\mathfrak m_T^n$。由我们对 $f_1$ 的选择，
$f$ 诱导同构 $\mathfrak m/\mathfrak m^2 \to \mathfrak m_T/\mathfrak m_T^2$，
故每个 $f_n$ 都是满射。因为 $A$ 完备，可推出 $f$ 是满射。
这意味着 $\dim(A) \geq \dim(T) = n$。故按定义，$A$ 正则。
（也可推出 $f$ 是同构。）
\end{proof}

\begin{lemma}
\label{more-algebra-lemma-lift-residue-field}
设 $k$ 为域，$(A, \mathfrak m, \kappa)$ 为完备局部 $k$-代数。
若 $\kappa/k$ 可分，则存在 $k$-代数映射 $\kappa \to A$，使
$\kappa \to A \to \kappa$ 是 $\text{id}_\kappa$。
\end{lemma}

\begin{proof}
由 Algebra, Proposition
\ref{algebra-proposition-characterize-separable-field-extensions}，
扩张 $\kappa/k$ 形式光滑。由引理 \ref{more-algebra-lemma-formally-smooth}，
$k \to \kappa$ 在定义 \ref{more-algebra-definition-formally-smooth} 的意义下形式光滑。
再由引理 \ref{more-algebra-lemma-lift-continuous} 得到 $\kappa \to A$。
\end{proof}

\begin{lemma}
\label{more-algebra-lemma-power-series-over-residue-field}
设 $k$ 为域，$(A, \mathfrak m, \kappa)$ 为完备局部 $k$-代数。
若 $\kappa/k$ 可分且 $A$ 正则，则存在 $k$-代数同构
$A \cong \kappa[[t_1, \ldots, t_d]]$。
\end{lemma}

\begin{proof}
按引理 \ref{more-algebra-lemma-lift-residue-field} 选取 $\kappa \to A$，
并应用 Algebra, Lemma
\ref{algebra-lemma-regular-complete-containing-coefficient-field}。
\end{proof}

\noindent
下列结果将在第 \ref{more-algebra-section-regular-fs} 节中得到改进。

\begin{lemma}
\label{more-algebra-lemma-regular-implies-fs}
设 $k$ 为域，$(A, \mathfrak m, K)$ 为正则局部 $k$-代数，
且 $K/k$ 可分。则 $k \to A$ 对 $\mathfrak m$-adic 拓扑形式光滑。
\end{lemma}

\begin{proof}
只需证明 $A$ 的完备化在 $k$ 上形式光滑；见引理
\ref{more-algebra-lemma-formally-smooth-completion}。因此可假设 $A$ 是完备局部正则
$k$-代数，其剩余域 $K$ 在 $k$ 上可分。由引理
\ref{more-algebra-lemma-power-series-over-residue-field}，可知
$A = K[[x_1, \ldots, x_n]]$。

\medskip\noindent
幂级数环 $K[[x_1, \ldots, x_n]]$ 在 $k$ 上形式光滑。
事实上，$K$ 在 $k$ 上形式光滑，而作为多项式代数，
$K[x_1, \ldots, x_n]$ 在 $K$ 上形式光滑。因此由 Algebra, Lemma
\ref{algebra-lemma-compose-formally-smooth}，
$K[x_1, \ldots, x_n]$ 在 $k$ 上形式光滑。
从而由引理 \ref{more-algebra-lemma-formally-smooth}，
$k \to K[x_1, \ldots, x_n]$ 对 $(x_1, \ldots, x_n)$-adic 拓扑形式光滑。
最后，由引理 \ref{more-algebra-lemma-formally-smooth-completion}，
$k \to K[[x_1, \ldots, x_n]]$ 对 $(x_1, \ldots, x_n)$-adic 拓扑形式光滑。
\end{proof}

\begin{lemma}
\label{more-algebra-lemma-formally-smooth-finite-type}
设 $A \to B$ 为有限型环映射，且 $A$ Noether。
设 $\mathfrak q \subset B$ 是位于 $\mathfrak p \subset A$ 之上的素理想。
下列条件等价：
\begin{enumerate}
\item $A \to B$ 在 $\mathfrak q$ 处光滑；以及
\item $A_\mathfrak p \to B_\mathfrak q$ 对 $\mathfrak q$-adic 拓扑形式光滑。
\end{enumerate}
\end{lemma}

\begin{proof}
推含 (2) $\Rightarrow$ (1) 由 Algebra, Lemma
\ref{algebra-lemma-smooth-test-artinian} 得出。反之，若 $A \to B$
在 $\mathfrak q$ 处光滑，则对某个 $g \in B$、
$g \not \in \mathfrak q$，$A \to B_g$ 光滑。于是由 Algebra, Proposition
\ref{algebra-proposition-smooth-formally-smooth}，$A \to B_g$ 形式光滑。
因此 $A_\mathfrak p \to B_\mathfrak q$ 形式光滑，因为局部化保持形式光滑
（例如，由 Algebra, Proposition
\ref{algebra-proposition-characterize-formally-smooth} 的判据，
以及余切复形对局部化表现良好这一事实；见 Algebra, Lemmas
\ref{algebra-lemma-NL-localize-bottom} 与
\ref{algebra-lemma-localize-NL}）。最后，由引理
\ref{more-algebra-lemma-formally-smooth}，$A_\mathfrak p \to B_\mathfrak q$
对 $\mathfrak q$-adic 拓扑形式光滑。
\end{proof}









\section{幂级数环的一些结果}
\label{more-algebra-section-power-series}

\noindent
Noether 局部环之间形式光滑映射的问题，往往可以归结为幂级数环之间
映射的问题。本节证明若干辅助引理，以便进行这类论证。

\begin{lemma}
\label{more-algebra-lemma-power-series-ring-over-Cohen-fs}
设 $K$ 是特征为 $0$ 的域，且 $A = K[[x_1, \ldots, x_n]]$。
设 $L$ 是特征为 $p > 0$ 的域，且 $B = L[[x_1, \ldots, x_n]]$。
设 $\Lambda$ 为 Cohen 环。令 $C = \Lambda[[x_1, \ldots, x_n]]$。
\begin{enumerate}
\item $\mathbf{Q} \to A$ 对 $\mathfrak m_A$-adic 拓扑形式光滑。
\item $\mathbf{F}_p \to B$ 对 $\mathfrak m_B$-adic 拓扑形式光滑。
\item $\mathbf{Z} \to C$ 对 $\mathfrak m_C$-adic 拓扑形式光滑。
\end{enumerate}
\end{lemma}

\begin{proof}
由幂级数环的泛性质，只需证明：
\begin{enumerate}
\item $\mathbf{Q} \to K$ 形式光滑。
\item $\mathbf{F}_p \to L$ 形式光滑。
\item $\mathbf{Z} \to \Lambda$ 对 $\mathfrak m_\Lambda$-adic 拓扑形式光滑。
\end{enumerate}
前两项即 Algebra, Proposition
\ref{algebra-proposition-characterize-separable-field-extensions}。
第三项由 Algebra, Lemma
\ref{algebra-lemma-cohen-ring-formally-smooth} 得出，因为对于如定义
\ref{more-algebra-definition-formally-smooth} 中的任意检验图，$p$ 的某个幂在
$A/J$ 中为零，因而 $p$ 的某个幂在 $A$ 中也为零。
\end{proof}

\begin{lemma}
\label{more-algebra-lemma-quotient-power-series-ring-over-Cohen}
设 $K$ 为域，且 $A = K[[x_1, \ldots, x_n]]$。
设 $\Lambda$ 为 Cohen 环，并令 $B = \Lambda[[x_1, \ldots, x_n]]$。
\begin{enumerate}
\item 若 $y_1, \ldots, y_n \in A$ 是一组正则参数系，则
$K[[y_1, \ldots, y_n]] \to A$ 是同构。
\item 若 $z_1, \ldots, z_r \in A$ 构成 $A$ 的一组正则参数系的一部分，
则 $r \leq n$，且
$A/(z_1, \ldots, z_r) \cong K[[y_1, \ldots, y_{n - r}]]$。
\item 若 $p, y_1, \ldots, y_n \in B$ 是一组正则参数系，则
$\Lambda[[y_1, \ldots, y_n]] \to B$ 是同构。
\item 若 $p, z_1, \ldots, z_r \in B$ 构成 $B$ 的一组正则参数系的一部分，
则 $r \leq n$，且
$B/(z_1, \ldots, z_r) \cong \Lambda[[y_1, \ldots, y_{n - r}]]$。
\end{enumerate}
\end{lemma}

\begin{proof}
证明 (1)。令 $A' = K[[y_1, \ldots, y_n]]$。显然，对每个 $n \geq 1$，
映射 $A' \to A$ 诱导同构
$A'/\mathfrak m_{A'}^n \to A/\mathfrak m_A^n$。
由于 $A$ 与 $A'$ 均完备，故 $A' \to A$ 是同构。
证明 (2)。将 $z_1, \ldots, z_r$ 扩充成 $A$ 的一组正则参数系
$z_1, \ldots, z_r, y_1, \ldots, y_{n - r}$。考虑映射
$A' = K[[z_1, \ldots, z_r, y_1, \ldots, y_{n - r}]] \to A$。
由 (1)，这是同构。又显然
$A'/(z_1, \ldots, z_r) \cong K[[y_1, \ldots, y_{n - r}]]$，
故 (2) 成立。(3) 与 (4) 的证明分别同 (1) 与 (2) 完全相同。
\end{proof}

\begin{lemma}
\label{more-algebra-lemma-embed-map-Noetherian-complete-local-rings}
设 $A \to B$ 为 Noether 完备局部环的局部同态。则存在交换图
$$
\xymatrix{
S \ar[r] & B \\
R \ar[u] \ar[r] & A \ar[u]
}
$$
并具有以下性质：
\begin{enumerate}
\item 水平箭头均为满射；
\item 若 $A/\mathfrak m_A$ 的特征为零，则 $S$ 与 $R$
都是域上的幂级数环；
\item 若 $A/\mathfrak m_A$ 的特征为 $p > 0$，则 $S$ 与 $R$
都是 Cohen 环上的幂级数环；以及
\item $R \to S$ 将 $R$ 的一组正则参数系映为 $S$
的一组正则参数系的一部分。
\end{enumerate}
特别地，$R \to S$ 平坦（见 Algebra, Lemma
\ref{algebra-lemma-flat-over-regular}），且其纤维
$S/\mathfrak m_R S$ 正则（见 Algebra, Lemma
\ref{algebra-lemma-regular-ring-CM}）。
\end{lemma}

\begin{proof}
利用 Cohen 结构定理（Algebra, Theorem
\ref{algebra-theorem-cohen-structure-theorem}），按引理所述选取满射
$S \to B$：若剩余特征为 $p > 0$，则取 $S$ 为 Cohen 环上的幂级数环；
否则取它为域上的幂级数环。设 $J \subset S$ 为 $S \to B$ 的核。
接着选取满射 $R = \Lambda[[x_1, \ldots, x_n]] \to A$：若 $A$
的剩余特征为 $p > 0$，则取 $\Lambda$ 为 Cohen 环；否则取
$\Lambda$ 为 $A$ 的剩余域。
我们将复合映射 $\Lambda[[x_1, \ldots, x_n]] \to A \to B$
提升为映射 $\varphi : R \to S$。这是可行的，因为
$\Lambda[[x_1, \ldots, x_n]]$ 对 $\mathbf{Z}$ 在
$\mathfrak m$-adic 拓扑下形式光滑（见引理
\ref{more-algebra-lemma-power-series-ring-over-Cohen-fs}）；应用引理
\ref{more-algebra-lemma-lift-continuous} 即得。
最后，以映射
$\varphi' : R = \Lambda[[x_1, \ldots, x_n]] \to S' = S[[y_1, \ldots, y_n]]$
替换 $\varphi$，其中 $\varphi'|_\Lambda = \varphi|_\Lambda$ 且
$\varphi'(x_i) = \varphi(x_i) + y_i$。同时，以将 $y_i$ 映为零的映射
$S' \to B$ 替换 $S \to B$。替换之后显然，$R$ 的一组正则参数系
被映为 $S'$ 中某个正则序列的一部分，结论得证。
\end{proof}

\noindent
下面这个引理应当有一个初等证明。

\begin{lemma}
\label{more-algebra-lemma-dominate-two-surjections}
设 $S \to R$ 与 $S' \to R$ 是完备 Noether 局部环的满射。
则 $S \times_R S'$ 是完备 Noether 局部环。
\end{lemma}

\begin{proof}
设 $k$ 为 $R$ 的剩余域。若 $k$ 的特征为 $p > 0$，则以 $\Lambda$
表示一个剩余域为 $k$ 的 Cohen 环
（Algebra, Definition \ref{algebra-definition-cohen-ring}；
Algebra, Lemma \ref{algebra-lemma-cohen-rings-exist}）。
若 $k$ 的特征为 $0$，则令 $\Lambda = k$。
选取满射 $\Lambda[[x_1, \ldots, x_n]] \to R$
（同 Cohen 结构定理中一样，见 Algebra, Theorem
\ref{algebra-theorem-cohen-structure-theorem}），并利用引理
\ref{more-algebra-lemma-power-series-ring-over-Cohen-fs} 与
\ref{more-algebra-lemma-lift-continuous} 将其提升为映射
% P02-E0167
$\Lambda[[x_1, \ldots, x_n]] \to S$ 以及
$\varphi : \Lambda[[x_1, \ldots, x_n]] \to S$ 和
$\varphi' : \Lambda[[x_1, \ldots, x_n]] \to S'$。
接着，选取 $f_1, \ldots, f_m \in S$ 生成 $S \to R$ 的核，并选取
$f'_1, \ldots, f'_{m'} \in S'$ 生成 $S' \to R$ 的核。于是映射
% P02-E0168
$$
\Lambda[[x_1, \ldots, x_n, y_1, \ldots, y_m, z_1, \ldots, z_{m'}]]
\longrightarrow S \times_R S,
$$
将 $x_i$ 映为 $(\varphi(x_i), \varphi'(x_i))$，
将 $y_j$ 映为 $(f_j, 0)$，并将
% P02-E0169
$z_{j'}$ 映为 $(0, f'_j)$；这个映射是满射。因此
$S \times_R S'$ 是某个完备局部环的商，故为完备环。
\end{proof}







\section{几何正则性与形式光滑性}
\label{more-algebra-section-regular-fs}

\noindent
本节综合前几节的结果，证明域上几何正则局部环的如下刻画。
随后，我们重新运用其中一些论证，证明 Noether 局部环之间在
$\mathfrak m$-adic 拓扑下形式光滑的映射的一种刻画。

\begin{theorem}
\label{more-algebra-theorem-regular-fs}
设 $k$ 为域。设 $(A, \mathfrak m, K)$ 为 Noether 局部
$k$-代数。若 $k$ 的特征为零，则下列条件等价：
\begin{enumerate}
\item $A$ 是正则局部环；以及
\item $k \to A$ 对 $\mathfrak m$-adic 拓扑形式光滑。
\end{enumerate}
若 $k$ 的特征为 $p > 0$，则下列条件等价：
\begin{enumerate}
\item $A$ 在 $k$ 上几何正则；
\item $k \to A$ 对 $\mathfrak m$-adic 拓扑形式光滑；
\item 对每个有限于 $k$ 的 $k \subset k' \subset k^{1/p}$，
环 $A \otimes_k k'$ 正则；
\item $A$ 正则，且典范映射
$H_1(L_{K/k}) \to \mathfrak m/\mathfrak m^2$ 为单射；以及
\item $A$ 正则，且映射
$\Omega_{k/\mathbf{F}_p} \otimes_k K \to \Omega_{A/\mathbf{F}_p} \otimes_A K$
为单射。
\end{enumerate}
\end{theorem}

\begin{proof}
若 $k$ 的特征为零，则 (1) 与 (2) 的等价性由引理
\ref{more-algebra-lemma-fs-implies-regular} 和 \ref{more-algebra-lemma-regular-implies-fs} 得出。

\medskip\noindent
若 $k$ 的特征为 $p > 0$，则由命题
\ref{more-algebra-proposition-characterization-geometrically-regular}，
(1)、(3)、(4) 与 (5) 等价。假设 (2) 成立。由引理
\ref{more-algebra-lemma-base-change-fs} 可见，
$k' \to A' = A \otimes_k k'$ 对
$\mathfrak m' = \mathfrak mA'$-adic 拓扑形式光滑。因此，若
$k \subset k'$ 为有限纯不可分扩张，则由引理
\ref{more-algebra-lemma-fs-implies-regular}，$A'$ 是正则局部环。故 (1) 成立。

\medskip\noindent
最后证明 (5) 推出 (2)。取定义 \ref{more-algebra-definition-formally-smooth}
中的如下实线图
$$
\xymatrix{
A \ar[r]_{\bar\psi} \ar@{-->}[rd] & B/J \\
k \ar[u]^i \ar[r]^\varphi & B \ar[u]_\pi
}
$$
由于 $J^2 = 0$，$J$ 具有典范的 $B/J$-模结构，并通过
$\bar\psi$ 具有 $A$-模结构。由于 $\bar\psi$ 对
$\mathfrak m$-adic 拓扑连续，故对某个 $n$ 有
$\mathfrak m^nJ = 0$。于是可以用 $B/J$-子模给 $J$ 作过滤
$0 \subset J_1 \subset J_2 \subset \ldots \subset J_n = J$，
使每个商 $J_{t + 1}/J_t$ 均被 $\mathfrak m$ 零化。
考察环映射序列
$B \to B/J_1 \to B/J_2 \to \ldots \to B/J$，
可见只需在 $J$ 被 $\mathfrak m$ 零化时，即 $J$ 是
$K$-向量空间时，证明虚线箭头存在。

\medskip\noindent
现设给定如上的图，并且 $J$ 被 $\mathfrak m$ 零化。
由引理 \ref{more-algebra-lemma-regular-implies-fs} 可见，
$\mathbf{F}_p \to A$ 对 $\mathfrak m$-adic 拓扑形式光滑。
因此存在环映射 $\psi : A \to B$，使得
$\pi \circ \psi = \bar \psi$。于是
$\psi \circ i, \varphi : k \to B$ 是两个与 $\pi$ 复合后相等的映射。
故 $D = \psi \circ i - \varphi : k \to J$ 是导子。
由 Algebra, Lemma \ref{algebra-lemma-universal-omega}，
可对某个 $k$-线性映射
$\xi : \Omega_{k/\mathbf{F}_p} \to J$ 写成
$D = \xi \circ \text{d}$。利用 $J$ 的 $K$-向量空间结构，
将 $\xi$ 延拓为 $K$-线性映射
$\xi' : \Omega_{k/\mathbf{F}_p} \otimes_k K \to J$。
利用 (5)，可以找到一个 $K$-线性映射
% P02-E0170
$\xi'' : \Omega_{A/\mathbf{F}_p} \otimes_A K$，其在
$\Omega_{k/\mathbf{F}_p} \otimes_k K$ 上的限制为 $\xi'$。记
$$
D' : A \xrightarrow{\text{d}} \Omega_{A/\mathbf{F}_p}
\to \Omega_{A/\mathbf{F}_p} \otimes_A K \xrightarrow{\xi''} J.
$$
最后令 $\psi' = \psi - D' : A \to B$。读者可验证，$\psi'$
是满足 $\pi \circ \psi' = \bar \psi$ 且
$\psi' \circ i = \varphi$ 的环映射，正如所需。
\end{proof}

\begin{example}
\label{more-algebra-example-fs}
设 $k$ 是特征为 $p > 0$ 的域。设 $a \in k$ 不是 $p$ 次幂。
一个剩余域在 $k$ 上纯不可分的几何正则局部 $k$-代数的标准例子是
$$
A = k[x, y]_{(x, y^p - a)}/(y^p - a - x)
$$
事实上，$A$ 是 $k$ 上某个光滑代数的局部化，故 $k \to A$
形式光滑，从而 $k \to A$ 对 $\mathfrak m$-adic 拓扑形式光滑。
下面是一个密切相关的例子。令
$k = \mathbf{F}_p(s)$ 且 $K = \mathbf{F}_p(t)^{perf}$。
我们断言环映射
$$
k \longrightarrow A = K[[x]],\quad s \longmapsto t + x
$$
对 $A$ 上的 $(x)$-adic 拓扑形式光滑。事实上，
$\Omega_{k/\mathbf{F}_p}$ 是以 $\text{d}s$ 为基的 $1$ 维向量空间。
它被映到 $\Omega_{A/\mathbf{F}_p}$ 中的元素
$\text{d}x + \text{d}t = \text{d}x$。
留给读者证明，$\Omega_{A/\mathbf{F}_p}$ 作为 $A$-模以
$\text{d}x$ 为基自由。因此，定理 \ref{more-algebra-theorem-regular-fs}
的条件 (5) 成立，从而 $k \to A$ 对 $(x)$-adic 拓扑形式光滑。
\end{example}

\begin{lemma}
\label{more-algebra-lemma-formally-smooth-flat}
设 $A \to B$ 为 Noether 局部环的局部同态。
假设 $A \to B$ 对 $\mathfrak m_B$-adic 拓扑形式光滑。
则 $A \to B$ 平坦。
\end{lemma}

\begin{proof}
% P02-E0171
由引理 \ref{more-algebra-lemma-formally-smooth-completion} 以及
Algebra, Lemma
\ref{algebra-lemma-completion-Noetherian-Noetherian}，
可以假设 $A$ 与 $B$ 是 Noether 完备局部环
（这里还用到 Algebra, Lemmas
\ref{algebra-lemma-flatness-descends-more-general} 和
\ref{algebra-lemma-completion-faithfully-flat}，以说明完备化之间的映射
平坦蕴含 $A \to B$ 平坦）。
选取引理 \ref{more-algebra-lemma-embed-map-Noetherian-complete-local-rings}
中的交换图
$$
\xymatrix{
S \ar[r] & B \\
R \ar[u] \ar[r] & A \ar[u]
}
$$
其中 $R \to S$ 平坦。设 $I \subset R$ 为 $R \to A$ 的核。
由于 $B$ 在 $A$ 上形式光滑，可见 $A$-代数映射
$$
S/IS \longrightarrow B
$$
有截面；见引理 \ref{more-algebra-lemma-lift-continuous}。
因此，$B$ 是平坦 $A$-模 $S/IS$ 的直和因子
（由平坦性的基变换，见 Algebra, Lemma
\ref{algebra-lemma-flat-base-change}），故为平坦模。
\end{proof}

\begin{lemma}
\label{more-algebra-lemma-formally-smooth-JZ}
设 $A \to B$ 为 Noether 局部环的局部同态。
假设 $A \to B$ 对 $\mathfrak m_B$-adic 拓扑形式光滑。
设 $K$ 为 $B$ 的剩余域。则 $A \to B \to K$ 的
Jacobi--Zariski 序列给出正合序列
$$
0 \to H_1(\NL_{K/A}) \to \mathfrak m_B/\mathfrak m_B^2
\to \Omega_{B/A} \otimes_B K \to \Omega_{K/A} \to 0
$$
\end{lemma}

\begin{proof}
由 Algebra, Lemma \ref{algebra-lemma-NL-surjection}，注意到
$\mathfrak m_B/\mathfrak m_B^2 = H_1(\NL_{K/B})$。
根据 Algebra, Lemma \ref{algebra-lemma-exact-sequence-NL}，
只需证明
$H_1(\NL_{K/A}) \to \mathfrak m_B/\mathfrak m_B^2$ 为单射。
设 $k$ 为 $A$ 的剩余域。$A \to k \to K$ 的
Jacobi--Zariski 序列给出 $\Omega_{K/A} = \Omega_{K/k}$ 以及正合序列
$$
\mathfrak m_A/\mathfrak m_A^2 \otimes_k K \to
H_1(\NL_{K/A}) \to
H_1(\NL_{K/k}) \to 0
$$
令 $\overline{B} = B \otimes_A k$。由于 $\overline{B}$ 正则，
理想 $\mathfrak m_{\overline{B}}$ 由一个正则序列生成。
对 $\mathfrak m_A B \subset \mathfrak m_B$ 应用引理
\ref{more-algebra-lemma-conormal-sequence-H1-regular} 和
\ref{more-algebra-lemma-noetherian-finite-all-equivalent}，得到
$\mathfrak m_A B / (\mathfrak m_AB \cap \mathfrak m_B^2) =
\mathfrak m_A B / \mathfrak m_A \mathfrak m_B$。
又由引理 \ref{more-algebra-lemma-formally-smooth-flat}，$A \to B$ 平坦，
所以上式等于 $\mathfrak m_A/\mathfrak m_A^2 \otimes_k K$。
因此得到短正合序列
$$
0 \to
\mathfrak m_A/\mathfrak m_A^2 \otimes_k K \to
\mathfrak m_B/\mathfrak m_B^2 \to
\mathfrak m_{\overline{B}}/\mathfrak m_{\overline{B}}^2 \to 0
$$
Jacobi--Zariski 序列的函子性表明，我们得到交换图
$$
\xymatrix{
&
\mathfrak m_A/\mathfrak m_A^2 \otimes_k K \ar[d] \ar[r] &
H_1(\NL_{K/A}) \ar[d] \ar[r] &
H_1(\NL_{K/k}) \ar[d] \ar[r] & 0 \\
0 \ar[r] &
\mathfrak m_A/\mathfrak m_A^2 \otimes_k K \ar[r] &
\mathfrak m_B/\mathfrak m_B^2 \ar[r] &
\mathfrak m_{\overline{B}}/\mathfrak m_{\overline{B}}^2 \ar[r] & 0
}
$$
由定理 \ref{more-algebra-theorem-regular-fs}，左侧竖直箭头为单射，因为由引理
\ref{more-algebra-lemma-base-change-fs}，$k \to \overline{B}$ 对
$\mathfrak m_{\overline{B}}$-adic 拓扑形式光滑。
蛇形引理即完成证明。
\end{proof}

\begin{proposition}
\label{more-algebra-proposition-fs-flat-fibre-fs}
设 $A \to B$ 为 Noether 局部环的局部同态。
设 $k$ 为 $A$ 的剩余域，并设
$\overline{B} = B \otimes_A k$ 为特殊纤维。下列条件等价：
\begin{enumerate}
\item $A \to B$ 平坦，且 $\overline{B}$ 在 $k$ 上几何正则；
\item $A \to B$ 平坦，且 $k \to \overline{B}$ 对
$\mathfrak m_{\overline{B}}$-adic 拓扑形式光滑；以及
\item $A \to B$ 对 $\mathfrak m_B$-adic 拓扑形式光滑。
\end{enumerate}
\end{proposition}

\begin{proof}
(1) 与 (2) 的等价性由定理 \ref{more-algebra-theorem-regular-fs} 得出。

\medskip\noindent
假设 (3)。由引理 \ref{more-algebra-lemma-formally-smooth-flat}，$A \to B$ 平坦。
由引理 \ref{more-algebra-lemma-base-change-fs}，$k \to \overline{B}$ 对
$\mathfrak m_{\overline{B}}$-adic 拓扑形式光滑。故 (2) 成立。

\medskip\noindent
假设 (2)。引理 \ref{more-algebra-lemma-formally-smooth-completion}
说明形式光滑性在完备化下保持。由 Algebra, Lemma
\ref{algebra-lemma-completion-faithfully-flat}，平坦性亦然。
因此，可分别用 $A$ 与 $B$ 的完备化替换它们，并假设
$A$ 与 $B$ 是 Noether 完备局部环。在这种情况下，选取引理
\ref{more-algebra-lemma-embed-map-Noetherian-complete-local-rings} 中的图
$$
\xymatrix{
S \ar[r] & B \\
R \ar[u] \ar[r] & A \ar[u]
}
$$
下文将使用此图的全部性质而不再一一说明。
固定 $R$ 的一组正则参数系 $t_1, \ldots, t_d$；当 $k$
的特征为 $p > 0$ 时，令 $t_1 = p$。
令 $\overline{S} = S \otimes_R k$。考察短正合序列
$$
0 \to J \to S \to B \to 0
$$
由于 $\overline{B}$ 与 $\overline{S}$ 均正则，
$\overline{S} \to \overline{B}$ 的核由元素
$\overline{x}_1, \ldots, \overline{x}_r$ 生成，且这些元素构成
$\overline{S}$ 的一组正则参数系的一部分；见 Algebra, Lemma
\ref{algebra-lemma-regular-quotient-regular}。
将这些元素提升为 $x_1, \ldots, x_r \in J$。于是
$t_1, \ldots, t_d, x_1, \ldots, x_r$ 构成 $S$
的一组正则参数系的一部分。因此，若 $k$ 的特征为零，则
$S/(x_1, \ldots, x_r)$ 是域上的幂级数环；若 $k$ 的特征为
$p > 0$，则它是 Cohen 环上的幂级数环；见引理
\ref{more-algebra-lemma-quotient-power-series-ring-over-Cohen}。
此外，$R \to S/(x_1, \ldots, x_r)$ 仍将
$t_1, \ldots, t_d$ 映为 $S/(x_1, \ldots, x_r)$
的一组正则参数系的一部分。换言之，可以用
$S/(x_1, \ldots, x_r)$ 替换 $S$，并假设我们有引理
\ref{more-algebra-lemma-embed-map-Noetherian-complete-local-rings} 中的图
$$
\xymatrix{
S \ar[r] & B \\
R \ar[u] \ar[r] & A \ar[u]
}
$$
而且 $\overline{S} = \overline{B}$。此时映射
$$
S \otimes_R A \longrightarrow B
$$
是同构，因为它是满射、在特殊纤维上是同构，并且源与目标均在
$A$ 上平坦（例如可用 Algebra, Lemma
\ref{algebra-lemma-mod-injective}；也可将短正合序列
$0 \to I \to S \otimes_R A \to B \to 0$ 在 $A$ 上与 $k$
张量，得到 $I \otimes_A k = 0$，从而由 Nakayama 引理有 $I = 0$）。
于是根据引理 \ref{more-algebra-lemma-base-change-fs}，只需证明
$R \to S$ 对 $\mathfrak m_S$-adic 拓扑形式光滑。
当然，由于 $\overline{S} = \overline{B}$，可知
$\overline{S}$ 在 $k = R/\mathfrak m_R$ 上形式光滑。

\medskip\noindent
选取元素 $y_1, \ldots, y_m \in S$，使
$t_1, \ldots, t_d, y_1, \ldots, y_m$ 是 $S$ 的一组正则参数系。
若 $k$ 的特征为零，则选取系数域 $K \subset S$；若 $k$
的特征为 $p > 0$，则选取剩余域为 $K$ 的 Cohen 环
$\Lambda \subset S$。此时，映射
$K[[t_1, \ldots, t_d, y_1, \ldots, y_m]] \to S$
（特征为零的情形）或
$\Lambda[[t_2, \ldots, t_d, y_1, \ldots, y_m]] \to S$
（特征为 $p > 0$ 的情形）是同构；见引理
\ref{more-algebra-lemma-quotient-power-series-ring-over-Cohen}。
从现在起，将 $S$ 视为上述幂级数环。

\medskip\noindent
余下的证明类似于定理 \ref{more-algebra-theorem-regular-fs} 证明中的论证。
取定义 \ref{more-algebra-definition-formally-smooth} 中的如下实线图
$$
\xymatrix{
S \ar[r]_{\bar\psi} \ar@{-->}[rd] & N/J \\
R \ar[u]^i \ar[r]^\varphi & N \ar[u]_\pi
}
$$
由于 $J^2 = 0$，$J$ 具有典范的 $N/J$-模结构，并通过
$\bar\psi$ 具有 $S$-模结构。由于 $\bar\psi$ 对
$\mathfrak m_S$-adic 拓扑连续，故对某个 $n$ 有
$\mathfrak m_S^nJ = 0$。于是可以用 $N/J$-子模给 $J$ 作过滤
$0 \subset J_1 \subset J_2 \subset \ldots \subset J_n = J$，
使每个商 $J_{t + 1}/J_t$ 均被 $\mathfrak m_S$ 零化。
考察环映射序列
$N \to N/J_1 \to N/J_2 \to \ldots \to N/J$，
可见只需在 $J$ 被 $\mathfrak m_S$ 零化时，即 $J$ 是
$K$-向量空间时，证明虚线箭头存在。

\medskip\noindent
现设给定如上的图，并且 $J$ 被 $\mathfrak m_S$ 零化。
由于 $\mathbf{Q} \to S$（特征为零的情形）或
$\mathbf{Z} \to S$（特征为 $p > 0$ 的情形）对
$\mathfrak m_S$-adic 拓扑形式光滑（见引理
\ref{more-algebra-lemma-power-series-ring-over-Cohen-fs}），故存在环映射
$\psi : S \to N$，使得 $\pi \circ \psi = \bar \psi$。
由于 $S$ 是某个子环上以 $t_1, \ldots, t_d$
（特征为零的情形）或以 $t_2, \ldots, t_d$
（特征为 $p > 0$ 的情形）为变量的幂级数环，故由幂级数环的泛性质，
可改变 $\psi$ 的选择，使 $\psi(t_i)$ 等于 $\varphi(t_i)$
（在特征为 $p$ 的情形，对 $t_1 = p$ 这是自动成立的）。
于是 $\psi \circ i$ 与 $\varphi : R \to N$ 是两个与 $\pi$
复合后相等、并且在 $t_1, \ldots, t_d$ 上相同的映射。
因此，$D = \psi \circ i - \varphi : R \to J$ 是将
$t_1, \ldots, t_d$ 零化的导子。
由 Algebra, Lemma \ref{algebra-lemma-universal-omega}，
可对某个将 $\text{d}t_1, \ldots, \text{d}t_d$
（依构造）和 $\mathfrak m_R \Omega_{R/\mathbf{Z}}$
（因为 $J$ 被 $\mathfrak m_R$ 零化）零化的 $R$-线性映射
$\xi : \Omega_{R/\mathbf{Z}} \to J$，写成
$D = \xi \circ \text{d}$。故 $\xi$ 分解为复合
$$
\Omega_{R/\mathbf{Z}} \to \Omega_{k/\mathbf{Z}} \xrightarrow{\xi'} J
$$
其中 $\xi'$ 为 $k$-线性。利用 $J$ 的 $K$-向量空间结构，
将 $\xi'$ 延拓为 $K$-线性映射
$$
\xi'' : \Omega_{k/\mathbf{Z}} \otimes_k K \longrightarrow J.
$$
由于 $\overline{S}/k$ 形式光滑，由定理
\ref{more-algebra-theorem-regular-fs} 可见
$$
\Omega_{k/\mathbf{Z}} \otimes_k K \to
\Omega_{\overline{S}/\mathbf{Z}} \otimes_S K
$$
为单射（在特征为零的情形也成立，因为此时甚至
$\Omega_{k/\mathbf{Z}} \to \Omega_{K/\mathbf{Z}}$
也是单射；见 Algebra, Proposition
\ref{algebra-proposition-characterize-separable-field-extensions}）。
因此，存在 $K$-线性映射
$\xi''' : \Omega_{\overline{S}/\mathbf{Z}} \otimes_S K \to J$，
其在 $\Omega_{k/\mathbf{Z}} \otimes_k K$ 上的限制为 $\xi''$。记
$$
D' : S \xrightarrow{\text{d}} \Omega_{S/\mathbf{Z}}
\to \Omega_{\overline{S}/\mathbf{Z}} \to
\Omega_{\overline{S}/\mathbf{Z}} \otimes_S K \xrightarrow{\xi'''} J.
$$
最后令 $\psi' = \psi - D' : S \to N$。读者可验证，$\psi'$
是满足 $\pi \circ \psi' = \bar \psi$ 且
$\psi' \circ i = \varphi$ 的环映射，正如所需。
\end{proof}

\noindent
作为上述结果的一个应用，我们证明形式光滑代数的形变没有障碍。

\begin{lemma}
\label{more-algebra-lemma-lift-fs}
设 $A$ 是剩余域为 $k$ 的 Noether 完备局部环。
设 $B$ 是 Noether 完备局部 $k$-代数。假设 $k \to B$
对 $\mathfrak m_B$-adic 拓扑形式光滑。则存在 Noether 完备局部环
$C$ 以及对 $\mathfrak m_C$-adic 拓扑形式光滑的局部同态
$A \to C$，使得 $C \otimes_A k \cong B$。
\end{lemma}

\begin{proof}
选取引理 \ref{more-algebra-lemma-embed-map-Noetherian-complete-local-rings} 中的图
$$
\xymatrix{
S \ar[r] & B \\
R \ar[u] \ar[r] & A \ar[u]
}
$$
设 $t_1, \ldots, t_d$ 为 $R$ 的一组正则参数系；当 $k$
的特征为 $p > 0$ 时，令 $t_1 = p$。
由于 $B$ 与 $\overline{S} = S \otimes_R k$ 均正则，
$\Ker(\overline{S} \to B)$ 由元素
$\overline{x}_1, \ldots, \overline{x}_r$ 生成，且这些元素构成
$\overline{S}$ 的一组正则参数系的一部分；见 Algebra, Lemma
\ref{algebra-lemma-regular-quotient-regular}。
将这些元素提升为 $x_1, \ldots, x_r \in S$。于是
$t_1, \ldots, t_d, x_1, \ldots, x_r$ 构成 $S$
的一组正则参数系的一部分。因此，若 $k$ 的特征为零，则
$S/(x_1, \ldots, x_r)$ 是域上的幂级数环；若 $k$ 的特征为
$p > 0$，则它是 Cohen 环上的幂级数环；见引理
\ref{more-algebra-lemma-quotient-power-series-ring-over-Cohen}。
此外，$R \to S/(x_1, \ldots, x_r)$ 仍将
$t_1, \ldots, t_d$ 映为 $S/(x_1, \ldots, x_r)$
的一组正则参数系的一部分。换言之，可以用
$S/(x_1, \ldots, x_r)$ 替换 $S$，并假设我们有引理
\ref{more-algebra-lemma-embed-map-Noetherian-complete-local-rings} 中的图
$$
\xymatrix{
S \ar[r] & B \\
R \ar[u] \ar[r] & A \ar[u]
}
$$
而且 $\overline{S} = B$。此时，由命题
\ref{more-algebra-proposition-fs-flat-fibre-fs}，$R \to S$ 对
$\mathfrak m_S$-adic 拓扑形式光滑。因此，由引理
\ref{more-algebra-lemma-base-change-fs}，基变换 $C = S \otimes_R A$
在 $A$ 上对 $\mathfrak m_C$-adic 拓扑形式光滑。
\end{proof}

\begin{remark}
\label{more-algebra-remark-what-does-it-mean}
引理 \ref{more-algebra-lemma-lift-fs} 的断言相当强。具体而言，设有 Noether
完备局部环的局部同态图
$$
\xymatrix{
& B \\
A \ar[r] & A' \ar[u]
}
$$
其中 $A \to A'$ 诱导剩余域同构
$k = A/\mathfrak m_A = A'/\mathfrak m_{A'}$，并且
$B \otimes_{A'} k$ 在 $k$ 上形式光滑。则可将其扩充为 Noether
完备局部环的局部同态交换图
$$
\xymatrix{
C \ar[r] & B \\
A \ar[r] \ar[u] & A' \ar[u]
}
$$
其中 $A \to C$ 对 $\mathfrak m_C$-adic 拓扑形式光滑，并且
$C \otimes_A k \cong B \otimes_{A'} k$。
事实上，根据引理 \ref{more-algebra-lemma-lift-fs}，选取 $A \to C$，
使其提升 $k$ 上的 $B \otimes_{A'} k$。由形式光滑性可得到箭头
$C \to B$；见引理 \ref{more-algebra-lemma-lift-continuous}。
以 $C \otimes_A^\wedge A'$ 表示 $C \otimes_A A'$ 关于理想
$C \otimes_A \mathfrak m_{A'}$ 的完备化。注意，
$C \otimes_A^\wedge A'$ 是 Noether 完备局部环
（见 Algebra, Lemma \ref{algebra-lemma-completion-Noetherian}），
并且在 $A'$ 上平坦（见 Algebra, Lemma
\ref{algebra-lemma-flat-module-powers}）。此外还有：
\begin{enumerate}
\item $C \otimes_A^\wedge A' \to B$ 是满射；
\item 若 $A \to A'$ 是满射，则 $C \to B$ 是满射；
\item 若 $A \to A'$ 有限，则 $C \to B$ 有限；以及
\item 若 $A' \to B$ 平坦，则 $C \otimes_A^\wedge A' \cong B$。
\end{enumerate}
事实上，由关于幂零理想的 Nakayama 引理
（见 Algebra, Lemma \ref{algebra-lemma-NAK}），从
$C \otimes_A k \cong B \otimes_{A'} k$ 可得，对每个 $n$，
$C \otimes_A A'/\mathfrak m_{A'}^n \to B/\mathfrak m_{A'}^nB$
均为满射。这证明了 (1)。(2) 与 (3) 由 (1) 得出。
(4) 由 Algebra, Lemma \ref{algebra-lemma-mod-injective} 得出。
\end{remark}




\section{正则环映射}
\label{more-algebra-section-regular}

\noindent
设 $k$ 为域。回忆：一个 Noether $k$-代数 $A$ 称为在 $k$ 上
{\it 几何正则}，当且仅当对于 $k$ 的每个有限纯不可分扩张 $k'$，
$A \otimes_k k'$ 都正则；见 Algebra, Definition
\ref{algebra-definition-geometrically-regular}。
此外，若此条件成立，则对每个有限生成域扩张 $k'/k$，
$A \otimes_k k'$ 都正则；见 Algebra, Lemma
\ref{algebra-lemma-geometrically-regular}。
我们在以下定义中使用这一概念。

\begin{definition}
\label{more-algebra-definition-regular}
环映射 $R \to \Lambda$ 称为{\it 正则}，如果它平坦，并且对于每个素理想
$\mathfrak p \subset R$，纤维环
$$
\Lambda \otimes_R \kappa(\mathfrak p) =
\Lambda_\mathfrak p/\mathfrak p\Lambda_\mathfrak p
$$
是 Noether 环，且在 $\kappa(\mathfrak p)$ 上几何正则。
\end{definition}

\noindent
若 $R \to \Lambda$ 是环映射且 $\Lambda$ Noether，
则纤维环总是 Noether 环。

\begin{lemma}[正则性是局部性质]
\label{more-algebra-lemma-regular-local}
设 $R \to \Lambda$ 为环映射，且 $\Lambda$ Noether。
下列条件等价：
\begin{enumerate}
\item $R \to \Lambda$ 正则；
\item 对每个在 $\mathfrak p \subset R$ 之上的
$\mathfrak q \subset \Lambda$，$R_\mathfrak p \to \Lambda_\mathfrak q$
正则；以及
\item 对每个在 $R$ 中的 $\mathfrak m$ 之上的极大理想
$\mathfrak m' \subset \Lambda$，
$R_\mathfrak m \to \Lambda_{\mathfrak m'}$ 正则。
\end{enumerate}
\end{lemma}

\begin{proof}
这是因为一个 Noether 环正则，当且仅当它的所有局部环都是正则局部环；
见 Algebra, Definition \ref{algebra-definition-regular}；而一个环映射平坦，
当且仅当所有诱导的局部环映射平坦；见 Algebra, Lemma
\ref{algebra-lemma-flat-localization}。
\end{proof}

\begin{lemma}[正则映射与基变换]
\label{more-algebra-lemma-regular-base-change}
设 $R \to \Lambda$ 为正则环映射。对于任意有限型环映射
$R \to R'$，基变换 $R' \to \Lambda \otimes_R R'$ 也正则。
\end{lemma}

\begin{proof}
平坦性在任意基变换下保持；见 Algebra, Lemma
\ref{algebra-lemma-flat-base-change}。
考察位于 $\mathfrak p \subset R$ 之上的素理想
$\mathfrak p' \subset R'$。由于 $R'$ 在 $R$ 上有限型，
剩余域扩张 $\kappa(\mathfrak p')/\kappa(\mathfrak p)$ 有限生成。
因此，纤维环
$$
(\Lambda \otimes_R R') \otimes_{R'} \kappa(\mathfrak p') =
\Lambda \otimes_R \kappa(\mathfrak p) \otimes_{\kappa(\mathfrak p)}
\kappa(\mathfrak p')
$$
是 Noether 环，这是由 Algebra, Lemma
\ref{algebra-lemma-Noetherian-field-extension} 以及对
$R \to \Lambda$ 的纤维环的假设得出的。
纤维的几何正则性由 Algebra, Lemma
\ref{algebra-lemma-geometrically-regular} 保持。
\end{proof}

\begin{lemma}[正则映射的复合]
\label{more-algebra-lemma-regular-composition}
设 $A \to B$ 与 $B \to C$ 为正则环映射。
若 $A \to C$ 的纤维环是 Noether 环，则 $A \to C$ 正则。
\end{lemma}

\begin{proof}
设 $\mathfrak p \subset A$ 为素理想。设
$\kappa(\mathfrak p) \subset k$ 为有限纯不可分扩张。
我们必须证明 $C \otimes_A k$ 正则。由引理
\ref{more-algebra-lemma-regular-base-change}，可以假设 $A = k$，
从而归结为证明 $C$ 正则。已知 $B$ 正则，且 $B \to C$
平坦并具有正则纤维。于是由 Algebra, Lemma
\ref{algebra-lemma-flat-over-regular-with-regular-fibre}，$C$ 正则。
略去若干细节。
\end{proof}

\begin{lemma}
\label{more-algebra-lemma-colimit-smooth-regular}
设 $R$ 为环。设 $(A_i, \varphi_{ii'})$ 为光滑 $R$-代数的有向系。
令 $\Lambda = \colim A_i$。若对所有 $\mathfrak p \subset R$，
纤维环 $\Lambda \otimes_R \kappa(\mathfrak p)$ 均为 Noether 环，
则 $R \to \Lambda$ 正则。
\end{lemma}

\begin{proof}
由 Algebra, Lemmas \ref{algebra-lemma-colimit-flat} 和
\ref{algebra-lemma-smooth-syntomic}，注意到 $\Lambda$ 在 $R$ 上平坦。
设 $\kappa(\mathfrak p) \subset k$ 为有限纯不可分扩张。注意到
$$
\Lambda \otimes_R \kappa(\mathfrak p) \otimes_{\kappa(\mathfrak p)} k =
\Lambda \otimes_R k = \colim A_i \otimes_R k
$$
是光滑 $k$-代数的余极限；见 Algebra, Lemma
\ref{algebra-lemma-base-change-smooth}。由 Algebra, Lemma
\ref{algebra-lemma-characterize-smooth-over-field}，光滑 $k$-代数的
每个局部环都正则，故由 Algebra, Lemma
\ref{algebra-lemma-colimit-regular}，$\Lambda \otimes_R k$
的所有局部环均正则。这就证明了引理。
\end{proof}

\noindent
下面考察域扩张何时定义正则环映射。

\begin{lemma}
\label{more-algebra-lemma-regular-field-extension}
设 $K/k$ 为域扩张。则 $k \to K$ 是正则环映射，当且仅当
$K$ 是 $k$ 的可分域扩张。
\end{lemma}

\begin{proof}
若 $k \to K$ 正则，则 $K$ 在 $k$ 上几何既约，因而由
Algebra, Proposition
\ref{algebra-proposition-characterize-separable-field-extensions}，
$K$ 在 $k$ 上可分。反之，若 $K/k$ 可分，则由 Algebra, Lemma
\ref{algebra-lemma-colimit-syntomic}，$K$ 是光滑 $k$-代数的余极限，
从而由引理 \ref{more-algebra-lemma-colimit-smooth-regular}，它是正则的。
\end{proof}

\begin{lemma}
\label{more-algebra-lemma-regular-permanence}
设 $A \to B \to C$ 为环映射。若 $A \to C$ 正则，并且
$B \to C$ 平坦且在谱上满射，则 $A \to B$ 正则。
\end{lemma}

\begin{proof}
由 Algebra, Lemma \ref{algebra-lemma-flat-permanence}，
$A \to B$ 平坦。设 $\mathfrak p \subset A$ 为素理想。
环映射
$B \otimes_A \kappa(\mathfrak p) \to C \otimes_A \kappa(\mathfrak p)$
平坦且在谱上满射。因此，由 Algebra, Lemma
\ref{algebra-lemma-geometrically-regular-descent}，
$B \otimes_A \kappa(\mathfrak p)$ 几何正则。
\end{proof}




\section{沿正则环映射上升的性质}
\label{more-algebra-section-ascending-properties}

\noindent
本节类似于 Algebra, Section
\ref{algebra-section-ascending-properties}，但这里环映射
$R \to S$ 是正则的。

\begin{lemma}
\label{more-algebra-lemma-reduced-goes-up}
设 $\varphi : R \to S$ 为环映射。假设
\begin{enumerate}
\item $\varphi$ 正则；
\item $S$ Noether；以及
\item $R$ Noether 且既约。
\end{enumerate}
则 $S$ 既约。
\end{lemma}

\begin{proof}
对于 Noether 环，既约等价于具有性质 $(S_1)$ 和 $(R_0)$；
见 Algebra, Lemma \ref{algebra-lemma-criterion-reduced}。
因此，可以应用 Algebra, Lemmas
\ref{algebra-lemma-Sk-goes-up} 和
\ref{algebra-lemma-Rk-goes-up}。
\end{proof}

\begin{lemma}
\label{more-algebra-lemma-normal-goes-up}
设 $\varphi : R \to S$ 为环映射。假设
\begin{enumerate}
\item $\varphi$ 正则；
\item $S$ Noether；以及
\item $R$ Noether 且正规。
\end{enumerate}
则 $S$ 正规。
\end{lemma}

\begin{proof}
对于 Noether 环，正规等价于具有性质 $(S_2)$ 和 $(R_1)$；
见 Algebra, Lemma \ref{algebra-lemma-criterion-normal}。
因此，可以应用 Algebra, Lemmas
\ref{algebra-lemma-Sk-goes-up} 和
\ref{algebra-lemma-Rk-goes-up}。
\end{proof}

\begin{lemma}
\label{more-algebra-lemma-regular-goes-up}
设 $\varphi : R \to S$ 为环映射。假设
\begin{enumerate}
\item $\varphi$ 正则；
\item $S$ Noether；以及
\item $R$ Noether 且正则。
\end{enumerate}
则 $S$ 正则。
\end{lemma}

\begin{proof}
对于 Noether 环，正则等价于对所有 $k$ 都具有性质 $(R_k)$。
因此，可以应用 Algebra, Lemma \ref{algebra-lemma-Rk-goes-up}。
\end{proof}

\begin{lemma}
\label{more-algebra-lemma-CM-goes-up}
设 $\varphi : R \to S$ 为环映射。假设
\begin{enumerate}
\item $\varphi$ 正则；
\item $S$ Noether；以及
\item $R$ Noether 且为 Cohen--Macaulay 环。
\end{enumerate}
则 $S$ 为 Cohen--Macaulay 环。
\end{lemma}

\begin{proof}
对于 Noether 环，成为 Cohen--Macaulay 环等价于对所有 $k$
都具有性质 $(S_k)$。因此，可以应用 Algebra, Lemma
\ref{algebra-lemma-Sk-goes-up}。
\end{proof}




\section{性质在完备化下的保持性}
\label{more-algebra-section-permanence-completion}

\noindent
给定 Noether 局部环 $(A, \mathfrak m)$，以 $A^\wedge$ 表示
$A$ 关于 $\mathfrak m$ 的完备化。下文将不再说明地使用如下事实：
$A^\wedge$ 是极大理想为
$\mathfrak m^\wedge = \mathfrak m A^\wedge$ 的 Noether
完备局部环，且 $A \to A^\wedge$ 忠实平坦。见 Algebra, Lemmas
\ref{algebra-lemma-completion-Noetherian-Noetherian}、
\ref{algebra-lemma-completion-complete} 和
\ref{algebra-lemma-completion-faithfully-flat}。

\begin{lemma}
\label{more-algebra-lemma-completion-dimension}
设 $A$ 为 Noether 局部环。则
$\dim(A) = \dim(A^\wedge)$。
\end{lemma}

\begin{proof}
由 Algebra, Lemma \ref{algebra-lemma-completion-complete}，
映射 $A \to A^\wedge$ 对 $n \geq 1$ 诱导同构
$A/\mathfrak m^n = A^\wedge/(\mathfrak m^\wedge)^n$。
由 Algebra, Lemma \ref{algebra-lemma-pushdown-module}，这蕴含
$$
\text{length}_A(A/\mathfrak m^n) =
\text{length}_{A^\wedge}(A^\wedge/(\mathfrak m^\wedge)^n)
$$
对所有 $n \geq 1$ 成立。因此 $d(A) = d(A^\wedge)$，
再由 Algebra, Proposition \ref{algebra-proposition-dimension} 得出结论。
另一种证明是使用 Algebra, Lemma
\ref{algebra-lemma-dimension-base-fibre-equals-total}。
\end{proof}

\begin{lemma}
\label{more-algebra-lemma-completion-depth}
设 $A$ 为 Noether 局部环。则
$\text{depth}(A) = \text{depth}(A^\wedge)$。
\end{lemma}

\begin{proof}
见 Algebra, Lemma \ref{algebra-lemma-apply-grothendieck}。
\end{proof}

\begin{lemma}
\label{more-algebra-lemma-completion-CM}
设 $A$ 为 Noether 局部环。则 $A$ 是 Cohen--Macaulay 环，
当且仅当 $A^\wedge$ 是 Cohen--Macaulay 环。
\end{lemma}

\begin{proof}
局部环 $A$ 是 Cohen--Macaulay 环，当且仅当
$\dim(A) = \text{depth}(A)$。由于这两个不变量都在完备化下保持
（引理 \ref{more-algebra-lemma-completion-dimension} 和
\ref{more-algebra-lemma-completion-depth}），断言得证。
\end{proof}

\begin{lemma}
\label{more-algebra-lemma-completion-regular}
设 $A$ 为 Noether 局部环。则 $A$ 正则，当且仅当
$A^\wedge$ 正则。
\end{lemma}

\begin{proof}
若 $A^\wedge$ 正则，则由 Algebra, Lemma
\ref{algebra-lemma-flat-under-regular}，$A$ 正则。
假设 $A$ 正则。设 $\mathfrak m$ 为 $A$ 的极大理想。则
$\dim_{\kappa(\mathfrak m)} \mathfrak m/\mathfrak m^2 =
\dim(A) = \dim(A^\wedge)$
（引理 \ref{more-algebra-lemma-completion-dimension}）。
另一方面，$\mathfrak mA^\wedge$ 是 $A^\wedge$ 的极大理想，
因而 $\mathfrak m_{A^\wedge}$ 至多由 $\dim(A^\wedge)$ 个元素生成。
故 $A^\wedge$ 正则。（也可以使用 Algebra, Lemma
\ref{algebra-lemma-flat-over-regular-with-regular-fibre}。）
\end{proof}

\begin{lemma}
\label{more-algebra-lemma-completion-dvr}
设 $A$ 为 Noether 局部环。则 $A$ 是离散赋值环，当且仅当
$A^\wedge$ 是离散赋值环。
\end{lemma}

\begin{proof}
这由引理 \ref{more-algebra-lemma-completion-dimension}、
\ref{more-algebra-lemma-completion-regular} 和 Algebra, Lemma
\ref{algebra-lemma-characterize-dvr} 得出。
\end{proof}

\begin{lemma}
\label{more-algebra-lemma-completion-reduced}
设 $A$ 为 Noether 局部环。
\begin{enumerate}
\item 若 $A^\wedge$ 既约，则 $A$ 也既约。
\item 一般而言，$A$ 既约并不蕴含 $A^\wedge$ 既约。
\item 若 $A$ 是 Nagata 环，则 $A$ 既约，当且仅当
$A^\wedge$ 既约。
\end{enumerate}
\end{lemma}

\begin{proof}
由于 $A \to A^\wedge$ 忠实平坦，由 Algebra, Lemma
\ref{algebra-lemma-descent-reduced} 得到 (1)。
关于 (2)，见 Algebra, Example
\ref{algebra-example-bad-dvr-char-p}
（特征为零时也有例子，见 Algebra, Remark
\ref{algebra-remark-resolution-dim-1}）。
关于 (3)，见 Algebra, Lemmas
\ref{algebra-lemma-local-nagata-domain-analytically-unramified} 和
\ref{algebra-lemma-analytically-unramified-easy}。
\end{proof}

\begin{lemma}
\label{more-algebra-lemma-completion-normal}
设 $A$ 为 Noether 局部环。若 $A^\wedge$ 正规，则 $A$ 也正规。
\end{lemma}

\begin{proof}
由于 $A \to A^\wedge$ 忠实平坦，这由 Algebra, Lemma
\ref{algebra-lemma-descent-normal} 得出。
\end{proof}

\begin{lemma}
\label{more-algebra-lemma-flat-completion}
设 $A \to B$ 为 Noether 局部环的局部同态。
则诱导的完备化映射 $A^\wedge \to B^\wedge$ 平坦，
当且仅当 $A \to B$ 平坦。
\end{lemma}

\begin{proof}
考察交换图
$$
\xymatrix{
A^\wedge \ar[r] & B^\wedge \\
A \ar[r] \ar[u] & B \ar[u]
}
$$
竖直箭头均忠实平坦。假设 $A^\wedge \to B^\wedge$ 平坦。
则 $A \to B^\wedge$ 平坦。因此，由 Algebra, Lemma
\ref{algebra-lemma-flatness-descends-more-general}，$B$ 在 $A$ 上平坦。

\medskip\noindent
假设 $A \to B$ 平坦。则 $A \to B^\wedge$ 平坦。
因此，对所有 $n \geq 1$，
$B^\wedge/\mathfrak m_A^n B^\wedge$ 在
$A/\mathfrak m_A^n$ 上平坦。注意，
$\mathfrak m_A^n A^\wedge$ 是 $A^\wedge$ 的极大理想
$\mathfrak m_A^\wedge$ 的第 $n$ 次幂，并且
$A/\mathfrak m_A^n = A^\wedge/(\mathfrak m_A^\wedge)^n$。
于是应用 Algebra, Lemma
\ref{algebra-lemma-flat-module-powers}
（取 $R = A^\wedge$、$I = \mathfrak m_A^\wedge$、
$S = B^\wedge$、$M = S$），可知 $B^\wedge$ 在
$A^\wedge$ 上平坦。
\end{proof}

\begin{lemma}
\label{more-algebra-lemma-flat-unramified}
设 $A \to B$ 为 Noether 局部环的平坦局部同态，并且
$\mathfrak m_A B = \mathfrak m_B$ 且
$\kappa(\mathfrak m_A) = \kappa(\mathfrak m_B)$。
则 $A \to B$ 在完备化上诱导同构 $A^\wedge \to B^\wedge$。
\end{lemma}

\begin{proof}
由 Algebra, Lemma \ref{algebra-lemma-finite-after-completion}，
$B^\wedge$ 是 $B$ 的 $\mathfrak m_A$-adic 完备化，并且
$A^\wedge \to B^\wedge$ 有限。由于 $A \to B$ 平坦，有
$\text{Tor}_1^A(B, \kappa(\mathfrak m_A)) = 0$。
因此，由引理 \ref{more-algebra-lemma-flat-after-completion}，
$B^\wedge$ 在 $A^\wedge$ 上平坦。进而由 Algebra, Lemma
\ref{algebra-lemma-finite-flat-local}，$B^\wedge$ 是自由
$A^\wedge$-模。根据我们的假设（以及 Algebra, Lemma
\ref{algebra-lemma-hathat-finitely-generated}），
$A^\wedge \to B^\wedge$ 诱导同构
$$
\kappa(\mathfrak m_A) = A^\wedge/\mathfrak m_A A^\wedge \to
B^\wedge/\mathfrak m_A B^\wedge = B^\wedge/\mathfrak m_B B^\wedge =
\kappa(\mathfrak m_B)
$$
故 $B^\wedge$ 是秩为 $1$ 的自由模。因此
$A^\wedge \to B^\wedge$ 是同构。
\end{proof}





\section{性质在 \'etale 映射下的保持性}
\label{more-algebra-section-permanence-etale}

\noindent
本节考察一个 \'etale 环映射 $\varphi : A \to B$，
研究 $A$ 的哪些性质由 $B$ 继承，以及 $B$ 在 $\mathfrak q$
处的局部环的哪些性质由 $A$ 在
$\mathfrak p = \varphi^{-1}(\mathfrak q)$ 处的局部环继承。
基本上，本节只是回顾并汇集先前的结果，不增加新内容。

\medskip\noindent
下文将不再说明地使用：\'etale 环映射平坦
（Algebra, Lemma \ref{algebra-lemma-etale}），且局部环的平坦局部同态
忠实平坦（Algebra, Lemma \ref{algebra-lemma-local-flat-ff}）。

\begin{lemma}
\label{more-algebra-lemma-Noetherian-etale-extension}
若 $A \to B$ 是 \'etale 环映射，且 $\mathfrak q$ 是 $B$
中位于 $\mathfrak p \subset A$ 之上的素理想，则
$A_{\mathfrak p}$ Noether，当且仅当 $B_{\mathfrak q}$ Noether。
\end{lemma}

\begin{proof}
由于 $A_\mathfrak p \to B_\mathfrak q$ 忠实平坦，
由 Algebra, Lemma \ref{algebra-lemma-descent-Noetherian} 可见，
$B_\mathfrak q$ Noether 蕴含 $A_\mathfrak p$ Noether。
反之，若 $A_\mathfrak p$ Noether，则 $B_\mathfrak q$ Noether，
因为它是某个有限型 $A_\mathfrak p$-代数的局部化。
\end{proof}

\begin{lemma}
\label{more-algebra-lemma-dimension-etale-extension}
若 $A \to B$ 是 \'etale 环映射，且 $\mathfrak q$ 是 $B$
中位于 $\mathfrak p \subset A$ 之上的素理想，则
$\dim(A_{\mathfrak p}) = \dim(B_{\mathfrak q})$。
\end{lemma}

\begin{proof}
事实上，由于 $A_{\mathfrak p} \to B_{\mathfrak q}$ 平坦，
下降定理成立，因而有不等式
$\dim(A_{\mathfrak p}) \leq \dim(B_{\mathfrak q})$；
见 Algebra, Lemma \ref{algebra-lemma-dimension-going-up}。
另一方面，设
$\mathfrak q_0 \subset \mathfrak q_1 \subset \ldots \subset \mathfrak q_n$
是 $B_{\mathfrak q}$ 中的一条素理想链。则对应的素理想序列
$\mathfrak p_0 \subset \mathfrak p_1 \subset \ldots \subset \mathfrak p_n$
（其中 $\mathfrak p_i = \mathfrak q_i \cap A_{\mathfrak p}$）
也确为一条链，即该序列中没有相等项。这是因为 \'etale 环映射拟有限
（见 Algebra, Lemma \ref{algebra-lemma-etale-quasi-finite}），
而拟有限环映射诱导的谱映射具有离散纤维（依定义）。
这说明 $\dim(A_{\mathfrak p}) \geq \dim(B_{\mathfrak q})$，正如所需。
\end{proof}

\begin{lemma}
\label{more-algebra-lemma-regular-etale-extension}
若 $A \to B$ 是 \'etale 环映射，且 $\mathfrak q$ 是 $B$
中位于 $\mathfrak p \subset A$ 之上的素理想，则
$A_{\mathfrak p}$ 正则，当且仅当 $B_{\mathfrak q}$ 正则。
\end{lemma}

\begin{proof}
由引理 \ref{more-algebra-lemma-Noetherian-etale-extension}，为证明等价性，
可以假设 $A_\mathfrak p$ 与 $B_\mathfrak q$ 均为 Noether 环。
设 $x_1, \ldots, x_t \in \mathfrak pA_\mathfrak p$
是一组极小生成元。由于
$A_\mathfrak p \to B_\mathfrak q$ 忠实平坦，它们在
$B_\mathfrak q$ 中的像 $y_1, \ldots, y_t$ 构成
$\mathfrak pB_\mathfrak q = \mathfrak q B_\mathfrak q$
的一组极小生成元
（Algebra, Lemma \ref{algebra-lemma-etale-at-prime}）。
依定义，$A_\mathfrak p$ 正则意味着
$t = \dim(A_\mathfrak p)$；对于 $B_\mathfrak q$ 亦然。
因此，引理由引理 \ref{more-algebra-lemma-dimension-etale-extension} 中的等式
$\dim(A_\mathfrak p) = \dim(B_\mathfrak q)$ 得出。
\end{proof}

\begin{lemma}
\label{more-algebra-lemma-Dedekind-etale-extension}
若 $A \to B$ 是 \'etale 环映射，且 $A$ 是 Dedekind 整环，
则 $B$ 是 Dedekind 整环的有限直积。特别地，对于
$\mathfrak q \subset B$ 为极大理想，局部化 $B_\mathfrak q$
是离散赋值环。
\end{lemma}

\begin{proof}
关于局部环的断言由引理
\ref{more-algebra-lemma-dimension-etale-extension}、
\ref{more-algebra-lemma-regular-etale-extension} 和 Algebra, Lemma
\ref{algebra-lemma-characterize-dvr} 得出。
由此，$B$ 是维数为 $1$ 的 Noether 正规环。
根据 Algebra, Lemma
\ref{algebra-lemma-characterize-reduced-ring-normal}，
$B$ 是维数为 $1$ 的正规整环的有限直积。
由 Algebra, Lemma \ref{algebra-lemma-characterize-Dedekind}，
这些整环都是 Dedekind 整环。
\end{proof}




\section{性质在 Hensel 化下的保持性}
\label{more-algebra-section-permanence-henselization}

\noindent
给定局部环 $R$，以 $R^h$ 和 $R^{sh}$ 分别表示 $R$ 的 Hensel 化
和严格 Hensel 化；见 Algebra, Definition
\ref{algebra-definition-henselization}。正如本节将证明的，$R$
的许多性质都会反映在 $R^h$ 与 $R^{sh}$ 中。

\begin{lemma}
\label{more-algebra-lemma-dumb-properties-henselization}
设 $(R, \mathfrak m, \kappa)$ 为局部环。则有：
\begin{enumerate}
\item $R \to R^h \to R^{sh}$ 是忠实平坦环映射；
\item $\mathfrak m R^h = \mathfrak m^h$ 且
$\mathfrak m R^{sh} = \mathfrak m^h R^{sh} = \mathfrak m^{sh}$；
\item 对所有 $n$，有 $R/\mathfrak m^n = R^h/\mathfrak m^nR^h$；
\item 存在元素 $x_i \in R^{sh}$，使得
$R^{sh}/\mathfrak m^nR^{sh}$ 是以
$x_i \bmod \mathfrak m^nR^{sh}$ 为基的自由
$R/\mathfrak m^n$-模。
\end{enumerate}
\end{lemma}

\begin{proof}
依构造，$R^h$ 是 \'etale $R$-代数的余极限；见 Algebra, Lemma
\ref{algebra-lemma-henselization}。由于 \'etale 环映射平坦
（Algebra, Lemma \ref{algebra-lemma-etale}），由 Algebra, Lemma
\ref{algebra-lemma-colimit-flat}，$R^h$ 在 $R$ 上平坦。
由于平坦局部环同态忠实平坦
（Algebra, Lemma \ref{algebra-lemma-local-flat-ff}），
$R \to R^h$ 忠实平坦。
环映射 $R^h \to R^{sh}$ 是有限 \'etale 环映射的余极限；
见 Algebra, Lemma \ref{algebra-lemma-strict-henselization} 的证明。
因此，与上述相同的论证表明 $R^h \to R^{sh}$ 忠实平坦。

\medskip\noindent
(2) 由 Algebra, Lemmas \ref{algebra-lemma-henselization} 和
\ref{algebra-lemma-strict-henselization} 得出。(3) 由
Algebra, Lemma \ref{algebra-lemma-local-artinian-basis-when-flat} 得出，
因为 $R/\mathfrak m \to R^h/\mathfrak mR^h$ 是同构，并且
$R/\mathfrak m^n \to R^h/\mathfrak m^nR^h$
是平坦环映射 $R \to R^h$ 的基变换，因而平坦
（Algebra, Lemma \ref{algebra-lemma-flat-base-change}）。
设 $\kappa^{sep}$ 为 $R^{sh}$ 的剩余域
（它是 $\kappa$ 的一个可分代数闭包）。选取 $x_i \in R^{sh}$，
使其像构成 $\kappa^{sep}$ 作为 $\kappa$-向量空间的一组基。
于是，完全如上地应用 Algebra, Lemma
\ref{algebra-lemma-local-artinian-basis-when-flat} 即得 (4)。
\end{proof}

\begin{lemma}
\label{more-algebra-lemma-henselization-formally-smooth}
设 $(R, \mathfrak m, \kappa)$ 为局部环。则
\begin{enumerate}
\item $R \to R^h$、$R^h \to R^{sh}$ 与 $R \to R^{sh}$
均为形式 \'etale 映射；
\item $R \to R^h$、$R^h \to R^{sh}$ 与 $R \to R^{sh}$
分别对 $\mathfrak m^h$、$\mathfrak m^{sh}$ 与
$\mathfrak m^{sh}$-拓扑形式光滑。
\end{enumerate}
\end{lemma}

\begin{proof}
(1) 由以下事实得出：依构造，$R^h$ 与 $R^{sh}$ 是 \'etale
代数的有向余极限；\'etale 代数是形式 \'etale 的
（Algebra, Lemma \ref{algebra-lemma-formally-etale-etale}）；
形式 \'etale 代数的余极限仍形式 \'etale
（Algebra, Lemma \ref{algebra-lemma-colimit-formally-etale}）。
(2) 由形式 \'etale 环映射形式光滑这一事实以及引理
\ref{more-algebra-lemma-formally-smooth} 得出。
\end{proof}

\begin{lemma}
\label{more-algebra-lemma-henselization-noetherian}
\begin{reference}
\cite[IV, Theorem 18.6.6 and Proposition 18.8.8]{EGA}
\end{reference}
设 $R$ 为局部环。下列条件等价：
\begin{enumerate}
\item $R$ Noether；
\item $R^h$ Noether；以及
\item $R^{sh}$ Noether。
\end{enumerate}
在这种情况下：
\begin{enumerate}
\item[(a)] $(R^h)^\wedge$ 与 $(R^{sh})^\wedge$ 是 Noether 完备局部环；
\item[(b)] $R^\wedge \to (R^h)^\wedge$ 是同构；
\item[(c)] $R^h \to (R^h)^\wedge$ 与
$R^{sh} \to (R^{sh})^\wedge$ 平坦；
\item[(d)] $R^\wedge \to (R^{sh})^\wedge$ 对
$\mathfrak m_{(R^{sh})^\wedge}$-adic 拓扑形式光滑；
\item[(e)] $(R^\wedge)^{sh} = R^\wedge \otimes_{R^h} R^{sh}$；以及
\item[(f)] $((R^\wedge)^{sh})^\wedge = (R^{sh})^\wedge$。
\end{enumerate}
\end{lemma}

\begin{proof}
由于 $R \to R^h \to R^{sh}$ 忠实平坦
（引理 \ref{more-algebra-lemma-dumb-properties-henselization}），由
Algebra, Lemma \ref{algebra-lemma-descent-Noetherian} 可见，
$R^h$ 或 $R^{sh}$ Noether 蕴含 $R$ Noether。
在余下的证明中，假设 $R$ Noether。

\medskip\noindent
由于 $\mathfrak m \subset R$ 有限生成，由引理
\ref{more-algebra-lemma-dumb-properties-henselization}，可知
$\mathfrak m^h = \mathfrak m R^h$ 与
$\mathfrak m^{sh} = \mathfrak mR^{sh}$ 有限生成。
因此，由 Algebra, Lemma
\ref{algebra-lemma-complete-local-ring-Noetherian}，
$(R^h)^\wedge$ 与 $(R^{sh})^\wedge$ 是 Noether 环。这证明了 (a)。

\medskip\noindent
注意，(b) 立即由引理
\ref{more-algebra-lemma-dumb-properties-henselization} 得出。
特别地，由 Algebra, Lemma
\ref{algebra-lemma-completion-faithfully-flat} 可见，
$(R^h)^\wedge$ 在 $R$ 上平坦。

\medskip\noindent
接着证明 $R^h \to (R^h)^\wedge$ 平坦。将
$R^h = \colim_i R_i$ 写成 \'etale $R$-代数的局部化的有向余极限。
根据 Algebra, Lemma \ref{algebra-lemma-colimit-rings-flat}，
若 $(R^h)^\wedge$ 在每个 $R_i$ 上平坦，则
$R^h \to (R^h)^\wedge$ 平坦。注意，依构造有
$R^h = R_i^h$。故由 (b)，
$R_i^\wedge = (R^h)^\wedge$ 在 $R_i$ 上平坦，正如所需。
为了完成 (c) 的证明，我们证明
$R^{sh} \to (R^{sh})^\wedge$ 平坦。为此，使用与上述相同的极限论证，
只需证明 $(R^{sh})^\wedge$ 在 $R$ 上平坦。由引理
\ref{more-algebra-lemma-dumb-properties-henselization} 可知，
$(R^{sh})^\wedge$ 是一个自由 $R$-模的完备化。
由引理 \ref{more-algebra-lemma-completed-direct-sum-flat}，它在 $R$ 上平坦，
正如所需。这完成了 (c) 的证明。

\medskip\noindent
至此已知 (c) 成立，且 $(R^h)^\wedge$ 与
$(R^{sh})^\wedge$ Noether。由 Algebra, Lemma
\ref{algebra-lemma-descent-Noetherian}，$R^h$ 与 $R^{sh}$ Noether。

\medskip\noindent
(d) 由引理 \ref{more-algebra-lemma-henselization-formally-smooth} 和
\ref{more-algebra-lemma-formally-smooth-completion} 得出。

\medskip\noindent
(e) 由 Algebra, Lemma \ref{algebra-lemma-sh-from-h-map} 以及
$R^\wedge$ 是 Hensel 环这一事实得出；后者见 Algebra, Lemma
\ref{algebra-lemma-complete-henselian}。

\medskip\noindent
证明 (f)。利用 (e)，存在映射
$R^{sh} \to (R^\wedge)^{sh}$；完备化后，它诱导映射
$(R^{sh})^\wedge \to ((R^\wedge)^{sh})^\wedge$。
再利用 (e)，存在映射 $R^\wedge \to (R^{sh})^\wedge$。
由于 $(R^{sh})^\wedge$ 是严格 Hensel 环（见上文），由
Algebra, Lemma \ref{algebra-lemma-strictly-henselian-functorial}，
此映射诱导映射 $(R^\wedge)^{sh} \to (R^{sh})^\wedge$。
完备化后得到映射
$((R^\wedge)^{sh})^\wedge \to (R^{sh})^\wedge$。
略去验证这两个映射互为逆映射的过程。
\end{proof}

\begin{lemma}
\label{more-algebra-lemma-henselization-reduced}
\begin{slogan}
既约性传递到（严格）Hensel 化。
\end{slogan}
设 $R$ 为局部环。下列条件等价：$R$ 既约，$R$ 的 Hensel 化
$R^h$ 既约，以及 $R$ 的严格 Hensel 化 $R^{sh}$ 既约。
\end{lemma}

\begin{proof}
环映射 $R \to R^h \to R^{sh}$ 忠实平坦。因此，由
Algebra, Lemma \ref{algebra-lemma-descent-reduced}
得到一个方向的推含。反之，假设 $R$ 既约。由于 $R^h$ 与
$R^{sh}$ 是 \'etale、因而光滑的 $R$-代数的滤过余极限，
结论由 Algebra, Lemma \ref{algebra-lemma-reduced-goes-up} 得出。
\end{proof}

\begin{lemma}
\label{more-algebra-lemma-henselization-nil}
设 $R$ 为局部环。以 $nil(R)$ 表示 $R$ 的幂零元素构成的理想。
则 $nil(R)R^h = nil(R^h)$ 且
$nil(R)R^{sh} = nil(R^{sh})$。
\end{lemma}

\begin{proof}
注意，$nil(R)$ 是由幂零元素组成且使商 $R/nil(R)$ 既约的最大理想。
由 Algebra, Lemma \ref{algebra-lemma-locally-nilpotent}，
$nil(R)R^h$ 由幂零元素组成。另注意，由 Algebra, Lemma
\ref{algebra-lemma-quotient-henselization}，
$R^h/nil(R) R^h$ 是 $R/nil(R)$ 的 Hensel 化。
因此，由引理 \ref{more-algebra-lemma-henselization-reduced}，
$R^h/nil(R)R^h$ 既约。故 $nil(R) R^h = nil(R^h)$，正如所需。
严格 Hensel 化的情形类似，但使用 Algebra, Lemma
\ref{algebra-lemma-quotient-strict-henselization}。
\end{proof}

\begin{lemma}
\label{more-algebra-lemma-henselization-normal}
设 $R$ 为局部环。下列条件等价：$R$ 是正规整环，$R$ 的 Hensel 化
$R^h$ 是正规整环，以及 $R$ 的严格 Hensel 化 $R^{sh}$ 是正规整环。
\end{lemma}

\begin{proof}
先说明，局部环正规当且仅当它是正规整环
（见 Algebra, Definition \ref{algebra-definition-ring-normal}）。
环映射 $R \to R^h \to R^{sh}$ 忠实平坦。因此，由
Algebra, Lemma \ref{algebra-lemma-descent-normal}
得到一个方向的推含。反之，假设 $R$ 正规。由于 $R^h$ 与
$R^{sh}$ 是 \'etale、因而光滑的 $R$-代数的滤过余极限，
结论由 Algebra, Lemmas \ref{algebra-lemma-normal-goes-up} 和
\ref{algebra-lemma-colimit-normal-ring} 得出。
\end{proof}

\begin{lemma}
\label{more-algebra-lemma-henselization-dimension}
对任意局部环 $R$，有
$\dim(R) = \dim(R^h) = \dim(R^{sh})$。
\end{lemma}

\begin{proof}
由于 $R \to R^{sh}$ 忠实平坦
（引理 \ref{more-algebra-lemma-dumb-properties-henselization}），由下降定理可知
$\dim(R^{sh}) \geq \dim(R)$；见 Algebra, Lemma
\ref{algebra-lemma-dimension-going-up}。
为证反向不等式，将 $R^{sh} = \colim R_i$ 写成局部环的有向余极限，
其中每个 $R_i$ 都是某个 \'etale $R$-代数的局部化。
若
$\mathfrak q_0 \subset \mathfrak q_1 \subset \ldots \subset \mathfrak q_n$
是 $R^{sh}$ 中的一条素理想链，则当 $i$ 充分大时，序列
$$
R_i \cap \mathfrak q_0 \subset
R_i \cap \mathfrak q_1 \subset \ldots \subset
R_i \cap \mathfrak q_n
$$
是 $R_i$ 中的一条素理想链。因此
$\dim(R^{sh}) \leq \sup_i \dim(R_i)$。
但由引理 \ref{more-algebra-lemma-dimension-etale-extension}，对每个 $i$
都有 $\dim(R_i) = \dim(R)$，结论得证。
\end{proof}

\begin{lemma}
\label{more-algebra-lemma-henselization-depth}
对任意 Noether 局部环 $R$，有
$\text{depth}(R) = \text{depth}(R^h) = \text{depth}(R^{sh})$。
\end{lemma}

\begin{proof}
由引理 \ref{more-algebra-lemma-henselization-noetherian}，已知
$R^h$ 与 $R^{sh}$ Noether。因此，引理由 Algebra, Lemma
\ref{algebra-lemma-apply-grothendieck} 得出。
\end{proof}

\begin{lemma}
\label{more-algebra-lemma-henselization-CM}
设 $R$ 为 Noether 局部环。下列条件等价：$R$ 是 Cohen--Macaulay 环，
$R$ 的 Hensel 化 $R^h$ 是 Cohen--Macaulay 环，以及 $R$
的严格 Hensel 化 $R^{sh}$ 是 Cohen--Macaulay 环。
\end{lemma}

\begin{proof}
由引理 \ref{more-algebra-lemma-henselization-noetherian}，已知
$R^h$ 与 $R^{sh}$ Noether，故该引理的陈述有意义。
再由引理 \ref{more-algebra-lemma-henselization-depth} 和
\ref{more-algebra-lemma-henselization-dimension}，有
$\text{depth}(R) = \text{depth}(R^h) = \text{depth}(R^{sh})$
以及
$\dim(R) = \dim(R^h) = \dim(R^{sh})$，结论随即得出。
\end{proof}

\begin{lemma}
\label{more-algebra-lemma-henselization-regular}
设 $R$ 为 Noether 局部环。下列条件等价：$R$ 是正则局部环，
$R$ 的 Hensel 化 $R^h$ 是正则局部环，以及 $R$
的严格 Hensel 化 $R^{sh}$ 是正则局部环。
\end{lemma}

\begin{proof}
由引理 \ref{more-algebra-lemma-henselization-noetherian}，已知
$R^h$ 与 $R^{sh}$ Noether，故该引理的陈述有意义。
设 $\mathfrak m$ 为 $R$ 的极大理想。
设 $x_1, \ldots, x_t \in \mathfrak m$ 是 $\mathfrak m$
的一组极小生成元，即它们在 $\mathfrak m/\mathfrak m^2$
中的像构成该空间在 $\kappa = R/\mathfrak m$ 上的一组基。
由于 $R \to R^h$ 与 $R \to R^{sh}$ 忠实平坦，故它们在
$R^h$ 中的像 $x_1^h, \ldots, x_t^h$ 以及在 $R^{sh}$
中的像 $x_1^{sh}, \ldots, x_t^{sh}$，分别构成
$\mathfrak m^h = \mathfrak mR^h$ 与
$\mathfrak m^{sh} = \mathfrak mR^{sh}$ 的一组极小生成元。
依定义，$R$ 正则意味着 $t = \dim(R)$；对于 $R^h$ 和 $R^{sh}$
亦然。因此，引理由引理 \ref{more-algebra-lemma-henselization-dimension} 中的维数等式
$\dim(R) = \dim(R^h) = \dim(R^{sh})$ 得出
% P02-E0174
\end{proof}

\begin{lemma}
\label{more-algebra-lemma-henselization-dvr}
设 $R$ 为 Noether 局部环。则 $R$ 是离散赋值环，当且仅当
$R^h$ 是离散赋值环，当且仅当 $R^{sh}$ 是离散赋值环。
\end{lemma}

\begin{proof}
这由引理 \ref{more-algebra-lemma-henselization-dimension}、
\ref{more-algebra-lemma-henselization-regular} 和 Algebra, Lemma
\ref{algebra-lemma-characterize-dvr} 得出。
\end{proof}

\begin{lemma}
\label{more-algebra-lemma-filtered-colimit-etale-noetherian-fibres}
设 $A$ 为环。设 $B$ 为 \'etale $A$-代数的滤过余极限。
设 $\mathfrak p$ 为 $A$ 的素理想。若 $B$ Noether，则位于
$\mathfrak p$ 之上的素理想只有有限多个
$\mathfrak q_1, \ldots, \mathfrak q_r$，并且
$B \otimes_A \kappa(\mathfrak p) = \prod \kappa(\mathfrak q_i)$；
各域扩张 $\kappa(\mathfrak q_i)/\kappa(\mathfrak p)$ 都是可分代数扩张。
\end{lemma}

\begin{proof}
将 $B$ 写成滤过余极限 $B = \colim B_i$，其中
$A \to B_i$ 为 \'etale 映射。一方面，
$B \otimes_A \kappa(\mathfrak p) = \colim B_i \otimes_A \kappa(\mathfrak p)$
是 \'etale $\kappa(\mathfrak p)$-代数的滤过余极限；另一方面，
它是 Noether 环。一个 \'etale $\kappa(\mathfrak p)$-代数是有限可分
域扩张的有限直积
（Algebra, Lemma \ref{algebra-lemma-etale-over-field}）。
因此，在诸代数 $B_i \otimes_A \kappa(\mathfrak p)$ 的素理想之间
没有非平凡的特殊化关系（这些素理想全都是极大且极小素理想），
从而在 $B \otimes_A \kappa(\mathfrak p)$ 的素理想之间也没有
非平凡的特殊化关系。因此，$B \otimes_A \kappa(\mathfrak p)$
既约，且只有有限多个素理想，并且它们全都是极小素理想。
% P02-E0172
所以，它是域的有限直积（使用 Algebra, Lemma
\ref{algebra-lemma-total-ring-fractions-no-embedded-points}
或 Algebra, Proposition
\ref{algebra-proposition-dimension-zero-ring}）。
这些域中的每一个都是有限可分扩张的余极限，故引理的最后一个断言成立。
\end{proof}

\begin{lemma}
\label{more-algebra-lemma-fibres-henselization}
设 $R$ 为 Noether 局部环。设 $\mathfrak p \subset R$ 为素理想。
则
$$
R^h \otimes_R \kappa(\mathfrak p) =
\prod\nolimits_{i = 1, \ldots, t} \kappa(\mathfrak q_i)
\quad\text{分别}\quad
R^{sh} \otimes_R \kappa(\mathfrak p) =
\prod\nolimits_{i = 1, \ldots, s} \kappa(\mathfrak r_i)
$$
其中 $\mathfrak q_1, \ldots, \mathfrak q_t$ 和
$\mathfrak r_1, \ldots, \mathfrak r_s$ 分别是 $R^h$ 与 $R^{sh}$
中位于 $\mathfrak p$ 之上的素理想。
% P02-E0173
此外，域扩张
$\kappa(\mathfrak q_i)/\kappa(\mathfrak p)$ 和
$\kappa(\mathfrak r_i)/\kappa(\mathfrak p)$ 分别是可分代数扩张。
\end{lemma}

\begin{proof}
这可以由更一般的引理
\ref{more-algebra-lemma-filtered-colimit-etale-noetherian-fibres} 推出，
其中使用 Hensel 化与严格 Hensel 化是 Noether 环
（如上所见）这一事实。不过，我们也给出如下直接证明。

\medskip\noindent
下文将不再说明地使用引理
\ref{more-algebra-lemma-dumb-properties-henselization} 和
\ref{more-algebra-lemma-henselization-noetherian} 的结果。
注意，$R^h/\mathfrak pR^h$ 与 $R^{sh}/\mathfrak pR^{sh}$
分别是 $R/\mathfrak p$ 的 Hensel 化与严格 Hensel 化；
分别见 Algebra, Lemma \ref{algebra-lemma-quotient-henselization}
与 Algebra, Lemma
\ref{algebra-lemma-quotient-strict-henselization}。
因此，可以用 $R/\mathfrak p$ 替换 $R$，并假设 $R$
是 Noether 局部整环且 $\mathfrak p = (0)$。
由于 $R^h$ 与 $R^{sh}$ Noether，它们分别只有有限多个极小素理想
$\mathfrak q_1, \ldots, \mathfrak q_t$ 和
$\mathfrak r_1, \ldots, \mathfrak r_s$。
由于 $R \to R^h$ 与 $R \to R^{sh}$ 平坦，由下降定理，
这些恰是位于 $\mathfrak p = (0)$ 之上的素理想。
最后，由于 $R$ 是整环，由引理
\ref{more-algebra-lemma-henselization-reduced}，$R^h$ 与 $R^{sh}$ 既约。
因此，$R^h \otimes_R \kappa(\mathfrak p)$ 与
$R^{sh} \otimes_R \kappa(\mathfrak p)$ 是只有有限多个素理想的
既约 Noether 环，并且所有这些素理想都是极小素理想
（因而也是极大理想）。于是这些环是 Artin 环，并且是它们在极大理想处
的局部化的直积；每个这样的局部化必为域
（见 Algebra, Proposition
\ref{algebra-proposition-dimension-zero-ring} 和 Algebra, Lemma
\ref{algebra-lemma-minimal-prime-reduced-ring}）。

\medskip\noindent
最后一个断言由如下事实得出：$R \to R^h$ 与
$R \to R^{sh}$ 是 \'etale 环映射的余极限，故诱导的剩余域扩张
是有限可分扩张的余极限；见 Algebra, Lemma
\ref{algebra-lemma-etale-at-prime}。
\end{proof}









\section{再论域扩张}
\label{more-algebra-section-p-bases}

\noindent
本节研究特征为 $p > 0$ 的域扩张的一些特殊现象。

\begin{definition}
\label{more-algebra-definition-p-basis}
设 $p$ 为素数。设 $k \to K$ 为特征 $p$ 的域的扩张。
以 $kK^p$ 表示 $k$ 与 $K^p$ 在 $K$ 中的合成域。
\begin{enumerate}
\item 子集 $\{x_i\} \subset K$ 称为{\it 在 $k$ 上 p-无关}，
如果满足 $0 \leq e_i < p$ 的元素
$x^E = \prod x_i^{e_i}$ 在 $kK^p$ 上线性无关。
\item $K$ 的子集 $\{x_i\}$ 称为{\it $K$ 在 $k$ 上的 p-基}，
如果诸元素 $x^E$ 构成 $K$ 在 $kK^p$ 上的一组基。
\end{enumerate}
\end{definition}

\noindent
这与 $\mathbf{F}_p$-代数的 $p$-基概念有关；我们将在后文讨论它
（此处插入未来的引用）。
% P02-E0175

\begin{lemma}
\label{more-algebra-lemma-p-basis}
设 $K/k$ 为域扩张。假设 $k$ 的特征为 $p > 0$。
设 $\{x_i\}$ 为 $K$ 的一个子集。下列条件等价：
\begin{enumerate}
\item 诸元素 $\{x_i\}$ 在 $k$ 上 $p$-无关；以及
\item 诸元素 $\text{d}x_i$ 在 $\Omega_{K/k}$ 中 $K$-线性无关。
\end{enumerate}
任意 $p$-无关族都可以扩充成 $K$ 在 $k$ 上的一组 $p$-基。
特别地，域 $K$ 在 $k$ 上有一组 $p$-基。
此外，下列条件等价：
\begin{enumerate}
\item[(a)] $\{x_i\}$ 是 $K$ 在 $k$ 上的一组 $p$-基；以及
\item[(b)] $\text{d}x_i$ 是 $K$-向量空间 $\Omega_{K/k}$ 的一组基。
\end{enumerate}
\end{lemma}

\begin{proof}
假设 (2)，并设 $\sum a_E x^E = 0$ 是一个线性关系，其中
$a_E \in k K^p$。设 $\theta_i : K \to K$ 为满足
$\theta_i(x_j) = \delta_{ij}$（Kronecker delta）的 $k$-导子。
注意，$K$ 的任意 $k$-导子都将 $kK^p$ 零化。
对给定关系应用 $\theta_i$，得到新的关系
$$
\sum\nolimits_{E, e_i > 0}
e_i a_E x_1^{e_1}\ldots x_i^{e_i - 1} \ldots x_n^{e_n} = 0
$$
因此，如果在某个 $a_E \not = 0$ 的关系中，选取总次数
$|E| = \sum e_i$ 最小的 $\sum a_E x^E$，就会得到矛盾。
故 (1) 成立。

\medskip\noindent
若 $\{x_i\}$ 是 $K$ 在 $k$ 上的一组 $p$-基，则
$K \cong kK^p[X_i]/(X_i^p - x_i^p)$。
由此可见，$\text{d}x_i$ 构成 $\Omega_{K/k}$ 在 $K$ 上的一组基。
故 (a) 推出 (b)。

\medskip\noindent
设 $\{x_i\}$ 是 $K$ 在 $k$ 上的一个 $p$-无关子集。
应用 Zorn 引理，可将其扩充成 $K$ 在 $k$ 上的一个极大 $p$-无关子集。
我们断言，$K$ 的任意极大 $p$-无关子集 $\{x_i\}$
都是 $K$ 在 $k$ 上的一组 $p$-基。该断言将说明 (1) 推出 (2)，
并确立 $p$-基的存在性。为证明该断言，设 $L$ 为
$K$ 中由 $kK^p$ 与诸 $x_i$ 生成的子域。必须证明 $L = K$。
若 $x \in K$ 但 $x \not \in L$，则 $x^p \in L$，且
% P02-E0176
$L(x) \cong L[z]/(z^p - x)$。因此，
$\{x_i\} \cup \{x\}$ 在 $k$ 上 $p$-无关，矛盾。

\medskip\noindent
最后，还须证明 (b) 推出 (a)。由 (1) 与 (2) 的等价性，
$\{x_i\}$ 是 $K$ 在 $k$ 上的极大 $p$-无关子集。
故由上述断言，它是一组 $p$-基。
\end{proof}

\begin{lemma}
\label{more-algebra-lemma-intersection-subfields-subspace}
设 $K/k$ 为域扩张。设 $\{K_\alpha\}_{\alpha \in A}$
为 $K$ 的一族子域，并具有以下性质：
\begin{enumerate}
\item 对所有 $\alpha \in A$，有 $k \subset K_\alpha$；
\item $k = \bigcap_{\alpha \in A} K_\alpha$；
\item 对任意 $\alpha, \alpha' \in A$，存在 $\alpha'' \in A$，
使得 $K_{\alpha''} \subset K_\alpha \cap K_{\alpha'}$。
\end{enumerate}
则对于 $n \geq 1$ 以及作为 $K$-向量空间的
$V \subset K^{\oplus n}$，有
$V \cap k^{\oplus n} \not = 0$，当且仅当对所有 $\alpha \in A$
都有 $V \cap K_\alpha^{\oplus n} \not = 0$。
\end{lemma}

\begin{proof}
对 $n$ 作归纳。$n = 1$ 的情形由假设得出。
假设已对 $K^{\oplus n - 1}$ 的子空间证明该结果。
设 $V \subset K^{\oplus n}$ 对所有 $\alpha \in A$
都与 $K_\alpha^{\oplus n}$ 的交非零。若
$V \cap 0 \oplus k^{\oplus n - 1}$ 非零，则结论成立。
故可假设它为零。由归纳假设，可以找到一个 $\alpha$，使得
$V \cap 0 \oplus K_\alpha^{\oplus n - 1}$ 为零。
取一个非零元素
$v = (x_1, \ldots, x_n) \in V \cap K_\alpha^{\oplus n}$。
由 $\alpha$ 的选取，$x_1$ 非零。以 $x_1^{-1}v$ 替换 $v$，
使得 $v = (1, x_2, \ldots, x_n)$。注意，若
$v' = (x_1', \ldots, x'_n) \in V \cap K_\alpha^{\oplus n}$，
则由 $\alpha$ 的选取有 $v' - x_1'v = 0$。
因此 $V \cap K_\alpha^{\oplus n} = K_\alpha v$。
若选取某个满足 $K_{\alpha'} \subset K_\alpha$ 的 $\alpha'$，
则必有 $v \in V \cap K_{\alpha'}^{\oplus n}$
（对 $\alpha'$ 应用同样的论证）。因此
$$
x_2, \ldots, x_n \in
\bigcap\nolimits_{\alpha' \in A, K_{\alpha'} \subset K_\alpha} K_{\alpha'}
$$
而由 (2) 与 (3)，右端等于 $k$。
\end{proof}

\begin{lemma}
\label{more-algebra-lemma-intersection-subfields}
设 $K$ 为特征 $p$ 的域。设 $\{K_\alpha\}_{\alpha \in A}$
为 $K$ 的一族子域，并具有以下性质：
\begin{enumerate}
\item 对所有 $\alpha \in A$，有 $K^p \subset K_\alpha$；
\item $K^p = \bigcap_{\alpha \in A} K_\alpha$；
\item 对任意 $\alpha, \alpha' \in A$，存在 $\alpha'' \in A$，
使得 $K_{\alpha''} \subset K_\alpha \cap K_{\alpha'}$。
\end{enumerate}
则
\begin{enumerate}
\item 映射
$\Omega_{K/\mathbf{F}_p} \to \Omega_{K/K_\alpha}$
的诸核之交为零；
\item 对任意有限扩张 $L/K$，有
$L^p = \bigcap_{\alpha \in A} L^pK_\alpha$。
\end{enumerate}
\end{lemma}

\begin{proof}
证明 (1)。选取 $K$ 在 $\mathbf{F}_p$ 上的一组 $p$-基
$\{x_i\}$。设
$\eta = \sum_{i \in I'} y_i \text{d}x_i$ 在每个 $\alpha \in A$
下都映为 $\Omega_{K/K_\alpha}$ 中的零，其中指标集 $I'$ 有限。
由引理 \ref{more-algebra-lemma-p-basis}，这意味着对每个 $\alpha$ 都存在关系
$$
\sum\nolimits_E a_{E, \alpha} x^E = 0,\quad a_{E, \alpha} \in K_\alpha
$$
其中 $E$ 遍历满足 $0 \leq e_i < p$ 的多重指标
$E = (e_i)_{i \in I'}$。另一方面，引理 \ref{more-algebra-lemma-p-basis}
保证不存在系数 $a_E \in K^p$ 的这种关系
$\sum a_E x^E = 0$。由引理
\ref{more-algebra-lemma-intersection-subfields-subspace}，这给出矛盾。

\medskip\noindent
证明 (2)。设有有限域扩张塔 $L/M/K$。
令 $M_\alpha = M^p K_\alpha$ 且
$L_\alpha = L^p K_\alpha = L^p M_\alpha$。
可以先证明 $M^p = \bigcap_{\alpha \in A} M_\alpha$，
再证明 $L^p = \bigcap_{\alpha \in A} L_\alpha$。
因此，只需对没有非平凡子域的单扩张证明 (2)。
首先，假设 $L = K(\theta)$ 在 $K$ 上可分。则 $L$
在 $K$ 上由 $\theta^p$ 生成，故可假设 $\theta \in L^p$。
在这种情况下，
$$
L^p = K^p \oplus K^p\theta \oplus \ldots K^p\theta^{d - 1}
\quad\text{且}\quad
L^pK_\alpha =
K_\alpha \oplus K_\alpha \theta \oplus \ldots K_\alpha\theta^{d - 1}
$$
其中 $d = [L : K]$。在这种情况下，结论显然。
另一种情形是 $L = K(\theta)$，其中
$\theta^p = t \in K$ 且 $t \not \in K^p$。此时
$$
L^p = K^p \oplus K^pt \oplus \ldots K^pt^{p - 1}
\quad\text{且}\quad
L^pK_\alpha =
K_\alpha \oplus K_\alpha t \oplus \ldots K_\alpha t^{p - 1}
$$
结论同样显然。
\end{proof}

\begin{lemma}
\label{more-algebra-lemma-power-series-ring-subfields}
设 $k$ 为特征 $p > 0$ 的域。设 $\{x_i\}_{i \in I}$
为 $k$ 的一组 $p$-基。设 $n, m \geq 0$。
设 $K$ 为
$A = k[[x_1, \ldots, x_n]][y_1, \ldots, y_m]$ 的分式域。
设 $J$ 为 $I$ 的有限子集。考察由 $k^p$ 以及满足
$i \in I \setminus J$ 的 $x_i$ 生成的子域
% P02-E0177
$k/k_J/k^p$。下列环
$$
A_J = k_J[[x_1^p, \ldots, x_n^p]][y_1^p, \ldots, y_m^p]
$$
的分式域 $K_J$ 构成 $K$ 的一族子域，并满足引理
\ref{more-algebra-lemma-intersection-subfields} 的条件。此外，每个环扩张
$A_J \subset A$ 都有限。
\end{lemma}

\begin{proof}
由于 $k/k_J$ 有限，由 Algebra, Lemma
\ref{algebra-lemma-finite-after-completion}，环扩张
% P02-E0178
$k_J[[x_1^p, \ldots, x_d^p]] \subset k[[x_1, \ldots, x_d]]$
有限。这蕴含 $A_J \to A$ 有限。

\medskip\noindent
下面验证引理 \ref{more-algebra-lemma-intersection-subfields} 的性质 (1)、(2)、(3)。
证明 (1)。对 $a \in A$，有 $a^p \in A_J$。故
$K^p \subset K_J$。证明 (2)。设
$f/g^p \in K$，其中 $f, g \in A$ 且 $g \not = 0$；
并假设它对 $J$ 的每个选择都属于 $K_J$。
暂时固定 $J$。由于 $f/g^p \in K_J$，可将其写成
$f/g^p = a/b^p$，其中 $a \in A_J$ 且 $b \in A$ 非零。
故 $b^p f \in A_J$。对任意 $A_J$-导子 $D : A \to A$，
有 $0 = D(b^pf) = b^p D(f)$；由于 $A$ 是整环，故 $D(f) = 0$。
取 $D = \partial_{x_i}$ 与 $D = \partial_{y_j}$，得到
% P02-E0179
$f \in k[[x_1^p, \ldots, x_n^p]][y_1^p, \ldots, y_d^p]$。
应用 $k_J$-导子 $\theta : k \to k$，同样得到 $f$
的所有系数都属于 $k_J$，即 $f \in A_J$。
由于显然
% P02-E0180
$A^p = \bigcap\nolimits_J A_J$，其中 $J$ 遍历引理中所述的所有子域，
故 $f \in A^p$，正如所需。
证明 (3)。这是显然的，因为
$K_{J \cup J'} \subset K_J \cap K_{J'}$。
\end{proof}





\section{奇异轨迹}
\label{more-algebra-section-singular-locus}

\noindent
设 $R$ 为 Noether 环。$X = \Spec(R)$ 的{\it 正则轨迹}
$\text{Reg}(X)$ 是由满足 $R_\mathfrak p$ 为正则局部环的素理想
$\mathfrak p$ 组成的集合。$X = \Spec(R)$ 的{\it 奇异轨迹}
$\text{Sing}(X)$ 是补集
$X \setminus \text{Reg}(X)$，即由满足 $R_\mathfrak p$
不是正则局部环的素理想 $\mathfrak p$ 组成的集合。
由 Algebra, Definition \ref{algebra-definition-regular}
之前的讨论可见，$\text{Reg}(X)$ 在一般化下稳定。
本节研究保证 $\text{Reg}(X)$ 为开集的条件。

\begin{definition}
\label{more-algebra-definition-J}
\begin{reference}
\cite[(32.B)]{MatCA}
\end{reference}
设 $R$ 为 Noether 环。设 $X = \Spec(R)$。
\begin{enumerate}
\item 若 $\text{Reg}(X)$ 包含一个非空开集，则称 $R$ 为 {\it J-0}；
\item 若 $\text{Reg}(X)$ 为开集，则称 $R$ 为 {\it J-1}；
\item 若任意有限型 $R$-代数都是 J-1 的，则称 $R$ 为 {\it J-2}。
\end{enumerate}
\end{definition}

\noindent
环 $\mathbf{Q}[x]/(x^2)$ 不满足 J-0，但满足 J-1。
另一方面，对于 Noether 整环以及更一般的非零既约 Noether 环，
J-1 蕴含 J-0，因为这种环在极小素理想处正则。
下面给出 J-1 性质的一种刻画。

\begin{lemma}
\label{more-algebra-lemma-J-1}
设 $R$ 为 Noether 环。设 $X = \Spec(R)$。
则 $R$ 为 J-1，当且仅当对所有
$\mathfrak p \in \text{Reg}(X)$，
$V(\mathfrak p) \cap \text{Reg}(X)$ 都包含
$V(\mathfrak p)$ 的一个非空开子集。
\end{lemma}

\begin{proof}
这由 Topology, Lemma
\ref{topology-lemma-characterize-open-Noetherian} 以及
$\text{Reg}(X)$ 在一般化下稳定这一事实得出；后者见
Algebra, Lemma
\ref{algebra-lemma-localization-of-regular-local-is-regular}。
\end{proof}

\begin{lemma}
\label{more-algebra-lemma-intersection-regular-with-closed}
设 $R$ 为 Noether 环。设 $X = \Spec(R)$。
假设对所有素理想 $\mathfrak p \subset R$，环
$R/\mathfrak p$ 都是 J-0 的。则 $R$ 是 J-1 的。
\end{lemma}

\begin{proof}
我们将证明引理 \ref{more-algebra-lemma-J-1} 的判据适用。
设 $\mathfrak p \in \text{Reg}(X)$ 为高度 $r$ 的素理想。
选取 $f_1, \ldots, f_r \in \mathfrak p$，使其像生成
$\mathfrak pR_\mathfrak p$。由于
$\mathfrak p \in \text{Reg}(X)$，可知
% P02-E0181
$f_1, \ldots, f_r$ 的像在 $R_\mathfrak p$ 中构成正则序列；
见 Algebra, Lemma \ref{algebra-lemma-regular-ring-CM}。
因此，由 Algebra, Lemma
\ref{algebra-lemma-regular-sequence-in-neighbourhood} 可见，
用某个满足 $g \not \in \mathfrak p$ 的 $g \in R$ 对 $R$
作 $R_g$ 替换后，序列 $f_1, \ldots, f_r$ 在 $R$ 中正则。
再作一次替换，还可假设 $f_1, \ldots, f_r$ 生成 $\mathfrak p$。
接着，设 $\mathfrak p \subset \mathfrak q$ 为素理想，并且
$(R/\mathfrak p)_\mathfrak q$ 是正则局部环。
由引理的假设，$V(\mathfrak p)$ 中存在一个由这类素理想组成的非空开子集，
故只需证明 $R_\mathfrak q$ 正则。
注意，$f_1, \ldots, f_r$ 在 $R_\mathfrak q$ 中构成正则序列，
且 $R_\mathfrak q/(f_1, \ldots, f_r)R_\mathfrak q$ 正则。
因此，由 Algebra, Lemma
\ref{algebra-lemma-regular-mod-x}，$R_\mathfrak q$ 正则。
\end{proof}

\begin{lemma}
\label{more-algebra-lemma-J-0-goes-down}
设 $R \to S$ 为环映射。假设：
\begin{enumerate}
\item $R$ 是 Noether 整环；
\item $R \to S$ 是单射且为有限型；以及
\item $S$ 是整环且为 J-0。
\end{enumerate}
则 $R$ 是 J-0。
\end{lemma}

\begin{proof}
用某个非零 $g \in S$ 对 $S$ 作 $S_g$ 替换后，可假设
$S$ 是正则环。由一般平坦性，还可假设 $R \to S$ 忠实平坦；
见 Algebra, Lemma
\ref{algebra-lemma-generic-flatness-Noetherian}。
于是由 Algebra, Lemma
\ref{algebra-lemma-descent-regular}，$R$ 正则。
\end{proof}

\begin{lemma}
\label{more-algebra-lemma-J-0-goes-up}
设 $R \to S$ 为环映射。假设：
\begin{enumerate}
\item $R$ 是 Noether 整环且为 J-0；
\item $R \to S$ 是单射且为有限型；
\item $S$ 是整环；以及
\item 所诱导的分式域扩张可分。
\end{enumerate}
则 $S$ 是 J-0。
\end{lemma}

\begin{proof}
可用一个主局部化替换 $R$，并假设 $R$ 是正则环。
由 Algebra, Lemma \ref{algebra-lemma-smooth-at-generic-point}，
环映射 $R \to S$ 在 $(0)$ 处光滑。
因此，用一个主局部化替换 $S$ 后，可假设 $S$ 在 $R$ 上光滑。
于是 $S$ 也正则；见 Algebra, Lemma
\ref{algebra-lemma-regular-goes-up}。
\end{proof}

\begin{lemma}
\label{more-algebra-lemma-J-2}
设 $R$ 为 Noether 环。下列条件等价：
\begin{enumerate}
\item $R$ 是 J-2；
\item 每个作为整环的有限型 $R$-代数都是 J-0；
\item 每个有限 $R$-代数都是 J-1；
\item 对每个素理想 $\mathfrak p$ 和每个有限纯不可分扩张
$L/\kappa(\mathfrak p)$，都存在一个有限 $R$-代数 $R'$，
它是 J-0 整环，且其分式域为 $L$。
\end{enumerate}
\end{lemma}

\begin{proof}
显然有推含 (1) $\Rightarrow$ (2) 和 (2) $\Rightarrow$ (4)。
回忆：一个 J-1 整环是 J-0 的。因此也有推含
(1) $\Rightarrow$ (3) 和 (3) $\Rightarrow$ (4)。

\medskip\noindent
设 $R \to S$ 为有限型环映射，尝试证明 $S$ 是 J-1。
由引理 \ref{more-algebra-lemma-intersection-regular-with-closed}，只需证明
对 $S$ 的每个素理想 $\mathfrak q$，$S/\mathfrak q$ 都是 J-0。
由此可见 (2) $\Rightarrow$ (1)。

\medskip\noindent
假设 (4)。我们将证明 (2)，从而完成证明。
设 $R \to S$ 为有限型环映射，且 $S$ 为整环。
令 $\mathfrak p = \Ker(R \to S)$。设 $K$ 为 $S$ 的分式域。
存在域的交换图
$$
\xymatrix{
K \ar[r] & K' \\
\kappa(\mathfrak p) \ar[u] \ar[r] & L \ar[u]
}
$$
其中水平箭头均为有限纯不可分域扩张，并且 $K'/L$ 可分；
见 Algebra, Lemma
\ref{algebra-lemma-make-separably-generated}。
按 (4) 选取 $R' \subset L$，并令 $S'$ 为映射
$S \otimes_R R' \to K'$ 的像。则 $S'$ 是分式域为 $K'$
的整环，故由引理 \ref{more-algebra-lemma-J-0-goes-up} 以及对 $R'$
的选取，$S'$ 是 J-0。再应用引理 \ref{more-algebra-lemma-J-0-goes-down}，
可知 $S$ 是 J-0，正如所需。
\end{proof}



\section{正则性与导子}
\label{more-algebra-section-regularity-derivations}

\noindent
设 $R \to S$ 为环映射。设 $D : R \to R$ 为导子。
若存在导子 $D' : S \to S$，使得交换图
$$
\xymatrix{
S \ar[r]_{D'} & S \\
R \ar[u] \ar[r]^D & R \ar[u]
}
$$
成立，则称 $D$ {\it 延拓到} $S$。

\begin{lemma}
\label{more-algebra-lemma-derivation-extends}
设 $R$ 为环。设 $D : R \to R$ 为导子。
\begin{enumerate}
\item 对任意理想 $I \subset R$，导子 $D$ 典范地延拓为
$I$-adic 完备化上的导子
$D^\wedge : R^\wedge \to R^\wedge$。
\item 对任意乘法子集 $S \subset R$，导子 $D$ 唯一地延拓到
$R$ 的局部化 $S^{-1}R$。
\end{enumerate}
若 $R \subset R'$ 是有限型环扩张，且对某个在 $R'$ 中为非零因子的
$g \in R$ 有 $R_g \cong R'_g$，则对某个 $N \geq 0$，
$g^ND$ 延拓到 $R'$。
\end{lemma}

\begin{proof}
证明 (1)。由 Leibniz 法则，对 $n \geq 2$ 有
$D(I^n) \subset I^{n - 1}$。因此 $D$ 诱导映射
$D_n : R/I^n \to R/I^{n - 1}$。取极限得到 $D^\wedge$。
略去验证 $D^\wedge$ 是导子的过程。

\medskip\noindent
证明 (2)。为将 $D$ 延拓到 $S^{-1}R$，只需令
$D(r/s) = D(r)/s - rD(s)/s^2$ 并验证各公理。

\medskip\noindent
证明最后一个断言。设 $x_1, \ldots, x_n \in R'$ 为
$R'$ 在 $R$ 上的一组生成元。选取 $N$，使得 $g^Nx_i \in R$。
考察 $g^{N + 1}D$。由 (2)，它延拓到 $R_g$。
此外，由 Leibniz 法则和上面对延拓的构造，有
$$
g^{N + 1}D(x_i) = g^{N + 1}D(g^{-N} g^Nx_i) = -Ng^Nx_iD(g) +
gD(g^Nx_i)
$$
且右端两项均属于 $R$。这蕴含
$$
g^{N + 1}D(x_1^{e_1} \ldots x_n^{e_n}) =
\sum e_i x_1^{e_1} \ldots x_i^{e_i - 1} \ldots x_n^{e_n} g^{N + 1}D(x_i)
$$
是 $R'$ 的元素。因此，$R'$ 的每个元素（它可写成诸 $x_i$
的单项式的和，系数在 $R$ 中）都被 $g^{N + 1}D$
映到 $R'$ 的元素，结论得证。
\end{proof}

\begin{lemma}
\label{more-algebra-lemma-quotient-regular}
\begin{slogan}
正确表述的超曲面 Jacobian 判据。
\end{slogan}
设 $R$ 为正则环。设 $f \in R$。假设存在导子
$D : R \to R$，使得 $D(f)$ 是 $R/(f)$ 的单位。
则 $R/(f)$ 正则。
\end{lemma}

\begin{proof}
只需在 $R$ 为具有极大理想 $\mathfrak m$、剩余域 $\kappa$
的局部环时证明。在这种情况下，只需证明
$f \not \in \mathfrak m^2$；见 Algebra, Lemma
\ref{algebra-lemma-regular-ring-CM}。
然而，若 $f \in \mathfrak m^2$，则由 Leibniz 法则，
$D(f) \in \mathfrak m$，矛盾。
\end{proof}

\begin{lemma}
\label{more-algebra-lemma-quotient-sequence-regular}
设 $(R, \mathfrak m, \kappa)$ 为正则局部环。设 $m \geq 1$。
设 $f_1, \ldots, f_m \in \mathfrak m$。假设存在导子
$D_1, \ldots, D_m : R \to R$，使得
$\det_{1 \leq i, j \leq m}(D_i(f_j))$ 是 $R$ 的单位。
则 $R/(f_1, \ldots, f_m)$ 正则，且
$f_1, \ldots, f_m$ 是正则序列。
\end{lemma}

\begin{proof}
只需证明 $f_1, \ldots, f_m$ 在
$\mathfrak m/\mathfrak m^2$ 中 $\kappa$-线性无关；
见 Algebra, Lemma \ref{algebra-lemma-regular-ring-CM}。
然而，若存在非平凡线性关系，则对某些 $a_i \in R$，且并非所有
$a_i \in \mathfrak m$，有
$\sum a_i f_i \in \mathfrak m^2$。
注意，由导子的 Leibniz 法则，
$D_i(\mathfrak m^2) \subset \mathfrak m$ 且
$D_i(a_j f_j) \equiv a_j D_i(f_j) \bmod \mathfrak m$。
因此，这将蕴含
$$
\sum a_j D_i(f_j) \in \mathfrak m
$$
与关于行列式的假设矛盾。
\end{proof}

\begin{lemma}
\label{more-algebra-lemma-degree-p-extension-regular}
设 $R$ 为正则环。设 $f \in R$。假设存在导子
$D : R \to R$，使得 $D(f)$ 是 $R$ 的单位。
则对任意整数 $n \geq 1$，$R[z]/(z^n - f)$ 正则。
更一般地，对任意 $p \in \mathbf{Z}[z]$，
$R[z]/(p(z) - f)$ 正则。
\end{lemma}

\begin{proof}
由 Algebra, Lemma \ref{algebra-lemma-regular-goes-up} 可见，
$R[z]$ 是正则环。对 $D$ 到 $R[z]$ 的延拓应用引理
\ref{more-algebra-lemma-quotient-regular}，其中将 $z$ 映为零。
这是可行的，因为 $D$ 将任意整系数多项式零化，并将 $f$ 映为单位。
\end{proof}

\begin{lemma}
\label{more-algebra-lemma-find-D}
设 $p$ 为素数。设 $B$ 为整环，并且在 $B$ 中 $p = 0$。
设 $f \in B$ 在 $B$ 的分式域中不是 $p$ 次幂。
若 $B$ 在某个 Noether 完备局部环上有限型，则存在导子
$D : B \to B$，使得 $D(f)$ 非零。
\end{lemma}

\begin{proof}
设 $R$ 为 Noether 完备局部环，并且存在有限型环映射
$R \to B$。当然，可以用它在 $B$ 中的像替换 $R$，
故可假设 $R$ 是特征 $p > 0$ 的整环
（同时仍是 Noether 完备局部环）。
由 Algebra, Lemma
\ref{algebra-lemma-complete-local-Noetherian-domain-finite-over-regular}，
可将 $R$ 写成某个域 $k$ 和整数 $n$ 所给的
$k[[x_1, \ldots, x_n]]$ 的有限扩张。因此，可用
$k[[x_1, \ldots, x_n]]$ 替换 $R$。
接着，使用 Algebra, Lemma
\ref{algebra-lemma-Noether-normalization-over-a-domain}，
将 $R \to B$ 分解为
$$
R \subset R[y_1, \ldots, y_d] \subset B' \subset B
$$
其中 $B'$ 在 $R[y_1, \ldots, y_d]$ 上有限，并且对某个非零
$g \in R$ 有 $B'_g \cong B_g$。注意，对某个充分大的整数 $N$，
$f' = g^{pN} f \in B'$。显然，$f'$ 在 $B'$ 的分式域中不是
$p$ 次幂。若能找到满足 $D'(f') \not = 0$ 的导子
$D' : B' \to B'$，则引理 \ref{more-algebra-lemma-derivation-extends}
保证对某个 $M > 0$，$D = g^MD'$ 延拓到 $B$。此时
$D(f) = g^MD'(f) = g^MD'(g^{-pN}f') = g^{M - pN}D'(f')$
非零。因此，只需在 $B$ 是
$A = k[[x_1, \ldots, x_n]][y_1, \ldots, y_m]$
的有限扩张时证明引理。

\medskip\noindent
假设 $B$ 是
$A = k[[x_1, \ldots, x_n]][y_1, \ldots, y_m]$
的有限扩张。以 $L$ 表示 $B$ 的分式域。
注意，$\text{d}f$ 在 $\Omega_{L/\mathbf{F}_p}$ 中非零；
见 Algebra, Lemma
\ref{algebra-lemma-derivative-zero-pth-power}。
应用引理 \ref{more-algebra-lemma-power-series-ring-subfields}，可找到
有限指数的子域 $k' \subset k$，使得令
$A' = k'[[x_1^p, \ldots, x_n^p]][y_1^p, \ldots, y_m^p]$ 后，
元素 $\text{d}f$ 在 $\Omega_{L/K'}$ 中的像非零，其中 $K'$
是 $A'$ 的分式域。因此，可选取满足 $D'(f) \not = 0$ 的
$K'$-导子 $D' : L \to L$。
依构造，$A' \subset A$ 与 $A \subset B$ 均有限，
故 $A' \subset B$ 有限。选取
$b_1, \ldots, b_t \in B$，使其生成 $B$ 作为 $A'$-模。
则对某些 $f_i, g_i \in B$ 且 $g_i \not = 0$，有
$D'(b_i) = f_i/g_i$。令 $D = g_1 \ldots g_t D'$ 即得结论。
\end{proof}

\begin{lemma}
\label{more-algebra-lemma-complete-Noetherian-domain-J-0}
设 $A$ 为 Noether 完备局部整环。则 $A$ 是 J-0。
\end{lemma}

\begin{proof}
由 Algebra, Lemma
\ref{algebra-lemma-complete-local-Noetherian-domain-finite-over-regular}，
可找到正则子环 $A_0 \subset A$，使 $A$ 在 $A_0$ 上有限。
所诱导的分式域扩张 $K/K_0$ 有限。若 $K/K_0$ 可分，则由引理
\ref{more-algebra-lemma-J-0-goes-up} 即得结论。若不可分，则 $A_0$ 与 $A$
的特征为 $p > 0$。对于任意子扩张 $K/M/K_0$，存在有限子扩张
$A_0 \subset B \subset A$，其分式域为 $M$。
因此，对 $[K : K_0]$ 作归纳，可假设存在
$A_0 \subset B \subset A$，使得 $B$ 是 J-0，且 $K/M$
没有非平凡子扩张。在这种情况下，若 $K/M$ 可分，则由引理
\ref{more-algebra-lemma-J-0-goes-up} 可知 $A$ 是 J-0。
若不可分，则对某些 $b_1, b_2 \in B$ 且 $b_2 \not = 0$，
有 $K = M[z]/(z^p - b_1/b_2)$，其中 $b_1/b_2$ 在 $M$
中不是 $p$ 次幂。选取非零 $a \in A$，使得 $az \in A$。
以 $b_2 a^p z$ 替换 $z$ 后，得到
$K = M[z]/(z^p - b)$，其中 $z \in A$ 且 $b \in B$
在 $M$ 中不是 $p$ 次幂。由引理 \ref{more-algebra-lemma-find-D}，
可找到满足 $D(b) \not = 0$ 的导子 $D : B \to B$。
应用引理 \ref{more-algebra-lemma-degree-p-extension-regular} 可见，
对于 $A$ 中位于 $B$ 的某个正则素理想之上、且不包含
$D(b)$ 的任意素理想 $\mathfrak p$，$A_\mathfrak p$ 正则。
由于 $B$ 是 J-0，故 $A$ 也是 J-0。
\end{proof}

\begin{proposition}
\label{more-algebra-proposition-ubiquity-J-2}
下列各类环都是 J-2：
\begin{enumerate}
\item 域；
\item Noether 完备局部环；
\item $\mathbf{Z}$；
\item 维数为 $1$ 的 Noether 局部环；
\item 维数为 $1$ 的 Nagata 环；
\item 分式域特征为零的 Dedekind 整环；
\item 上述任一类环的有限型环扩张。
\end{enumerate}
\end{proposition}

\begin{proof}
对于情形 (1)、(3)、(5) 与 (6)，可通过验证引理
\ref{more-algebra-lemma-J-2} 的条件 (4) 来证明。我们只处理 $R$
是维数为 $1$ 的 Nagata 环的情形。设
$\mathfrak p \subset R$ 为素理想，并设
$L/\kappa(\mathfrak p)$ 为有限纯不可分扩张。
若 $\mathfrak p \subset R$ 是极大理想，则 $R \to L$ 有限，
且 $L$ 是正则环，故已验证该条件。
若 $\mathfrak p \subset R$ 是极小素理想，则 Nagata 条件保证
% P02-E0182
$R$ 在 $L$ 中的整闭包 $R' \subset L$ 在 $R$ 上有限。
于是 $R'$ 是维数为 $1$ 的正规整环
（Algebra, Lemma \ref{algebra-lemma-integral-dim-up}），
因而正则（Algebra, Lemma
\ref{algebra-lemma-criterion-normal}），故在此情形也已验证该条件。

\medskip\noindent
对于情形 (2)，使用引理 \ref{more-algebra-lemma-J-2} 的条件 (3)。
设 $R$ 为 Noether 完备局部环。注意，若 $R \to R'$ 有限，
则 $R'$ 是 Noether 完备局部环的直积；见 Algebra, Lemma
\ref{algebra-lemma-quotient-complete-local}。
因此，由引理 \ref{more-algebra-lemma-intersection-regular-with-closed}，
只需证明作为整环的 Noether 完备局部环是 J-0；
这就是引理 \ref{more-algebra-lemma-complete-Noetherian-domain-J-0}。

\medskip\noindent
对于情形 (4)，也使用引理 \ref{more-algebra-lemma-J-2} 的条件 (3)。
事实上，若 $R$ 是维数为 $1$ 的 Noether 局部环，且
$R \to R'$ 有限，则 $\Spec(R')$ 是有限集。
由于正则轨迹在一般化下稳定，可见 $R'$ 是 J-1。
\end{proof}





\section{形式光滑性与正则性}
\label{more-algebra-section-fs-regular}

\noindent
本节标题指的是 Andr\'e 定理 \ref{more-algebra-theorem-fs-regular}。

\begin{lemma}
\label{more-algebra-lemma-lift-derivation-through-fs}
设 $A \to B$ 为 Noether 局部环的局部同态。
设 $D : A \to A$ 为导子。假设 $B$ 完备，并且
$A \to B$ 对 $\mathfrak m_B$-adic 拓扑形式光滑。
则存在 $D$ 的延拓 $D' : B \to B$。
\end{lemma}

\begin{proof}
以 $B[\epsilon] = B[x]/(x^2)$ 表示 $B$ 上的对偶数环。
考察环映射 $\psi : A \to B[\epsilon]$，
$a \mapsto a + \epsilon D(a)$。考察交换图
$$
\xymatrix{
B \ar[r]_1 & B \\
A \ar[u] \ar[r]^\psi & B[\epsilon] \ar[u]
}
$$
由引理 \ref{more-algebra-lemma-lift-continuous} 以及 $B/A$ 形式光滑的假设，
可找到适合该图的映射 $\varphi : B \to B[\epsilon]$。
写成 $\varphi(b) = b + \epsilon D'(b)$。于是
$D' : B \to B$ 就是所需的延拓。
\end{proof}

\begin{proposition}
\label{more-algebra-proposition-fs-regular}
设 $A \to B$ 为 Noether 完备局部环的局部同态。
设 $k$ 为 $A$ 的剩余域，并设
$\overline{B} = B \otimes_A k$ 为特殊纤维。
下列条件等价：
\begin{enumerate}
\item $A \to B$ 正则；
\item $A \to B$ 平坦，且 $\overline{B}$ 在 $k$ 上几何正则；
\item $A \to B$ 平坦，且 $k \to \overline{B}$ 对
$\mathfrak m_{\overline{B}}$-adic 拓扑形式光滑；以及
\item $A \to B$ 对 $\mathfrak m_B$-adic 拓扑形式光滑。
\end{enumerate}
\end{proposition}

\begin{proof}
我们已在命题 \ref{more-algebra-proposition-fs-flat-fibre-fs} 中证明
(2)、(3) 与 (4) 等价。显然 (1) 推出 (2)。
因此，假设等价条件 (2)、(3) 与 (4) 成立，并证明 (1)。

\medskip\noindent
设 $\mathfrak p$ 为 $A$ 的素理想。我们将证明
$B \otimes_A \kappa(\mathfrak p)$ 在
$\kappa(\mathfrak p)$ 上几何正则。
由引理 \ref{more-algebra-lemma-base-change-fs}，可以用
$A/\mathfrak p$ 替换 $A$，并用 $B/\mathfrak pB$ 替换 $B$。
故可假设 $A$ 是整环且 $\mathfrak p = (0)$。

\medskip\noindent
按 Algebra, Lemma
\ref{algebra-lemma-complete-local-Noetherian-domain-finite-over-regular}
选取 $A_0 \subset A$。下文将不再说明地使用该引理所述的全部性质。
由于 $A_0 \to A$ 在剩余域上诱导同构，且
$B/\mathfrak m_A B$ 在 $A/\mathfrak m_A$ 上几何正则，
可找到交换图
$$
\xymatrix{
C \ar[r] & B \\
A_0 \ar[r] \ar[u] & A \ar[u]
}
$$
其中 $A_0 \to C$ 对 $\mathfrak m_C$-adic 拓扑形式光滑，
且 $B = C \otimes_{A_0} A$；见注记
\ref{more-algebra-remark-what-does-it-mean}。
（由于 $A_0 \to A$ 有限，张量积中不需要完备化；见 Algebra, Lemma
\ref{algebra-lemma-completion-tensor}。）
因此，只需证明 $C \otimes_{A_0} K_0$ 是 $A_0$ 的分式域
$K_0$ 上的几何正则代数。

\medskip\noindent
上一段的要点是，可以假设
$A = k[[x_1, \ldots, x_n]]$，其中 $k$ 为域；或者
$A = \Lambda[[x_1, \ldots, x_n]]$，其中 $\Lambda$ 为 Cohen 环。
在这种情况下，由 Algebra, Lemma
\ref{algebra-lemma-flat-over-regular-with-regular-fibre}，
$B$ 是正则环。因此，$B \otimes_A K$ 也是正则环
（其中 $K$ 为 $A$ 的分式域）；若 $K$ 的特征为零，结论即得。

\medskip\noindent
因此，只剩下 $A = k[[x_1, \ldots, x_n]]$ 且 $k$
是特征 $p > 0$ 的域这一情形。
设 $L/K$ 为有限纯不可分域扩张。我们对 $[L : K]$ 作归纳，
证明 $B \otimes_A L$ 正则。基础情形为 $L = K$，已在上文证明。
设 $K \subset M \subset L$ 为子域，使得 $L$ 是 $M$
的次数为 $p$ 的扩张，并由添加某个元素 $f \in M$ 的 $p$ 次根得到。
设 $A'$ 是 $M$ 中分式域为 $M$ 的有限 $A$-子代数。
清除分母后，可以且确实假设 $f \in A'$。
令 $A'' = A'[z]/(z^p -f)$，并注意 $A' \subset A''$ 有限，
且 $A''$ 的分式域为 $L$。
由归纳假设，环
% P02-E0183
$B \otimes_A M$ 正则。我们有
$$
B \otimes_A L = B \otimes_A M[z]/(z^p - f)
$$
由引理 \ref{more-algebra-lemma-find-D}，存在导子 $D : A' \to A'$，
使得 $D(f) \not = 0$。由引理 \ref{more-algebra-lemma-descent-fs}，
$A' \to B \otimes_A A'$ 对 $\mathfrak m$-adic 拓扑形式光滑，
故可使用引理 \ref{more-algebra-lemma-lift-derivation-through-fs}
将 $D$ 延拓为导子
$D' : B \otimes_A A' \to B \otimes_A A'$。
注意，$D'(f) = D(f)$ 在 $B \otimes_A M$ 中为单位，
因为 $D(f)$ 在 $A' \subset M$ 中非零。
因此，由引理 \ref{more-algebra-lemma-degree-p-extension-regular}，
$B \otimes_A L$ 正则，结论得证。
\end{proof}

\begin{theorem}[Andr\'e]
\label{more-algebra-theorem-fs-regular}
\begin{reference}
推含 (4) $\Rightarrow$ (1) 是 \cite{Andre-smooth} 的主要结果。
\end{reference}
设 $A \to B$ 为 Noether 局部环的局部同态。
设 $k$ 为 $A$ 的剩余域，并设
$\overline{B} = B \otimes_A k$ 为特殊纤维。
假设 $A \to A^\wedge$ 正则\footnote{例如，若 $A$ 是 G-环或
（拟）优环；见第 \ref{more-algebra-section-G-ring} 节和第
\ref{more-algebra-section-excellent} 节。}。下列条件等价：
\begin{enumerate}
\item $A \to B$ 正则；
\item $A \to B$ 平坦，且 $\overline{B}$ 在 $k$ 上几何正则；
\item $A \to B$ 平坦，且 $k \to \overline{B}$ 对
$\mathfrak m_{\overline{B}}$-adic 拓扑形式光滑；以及
\item $A \to B$ 对 $\mathfrak m_B$-adic 拓扑形式光滑。
\end{enumerate}
\end{theorem}

\begin{proof}
我们已在命题 \ref{more-algebra-proposition-fs-flat-fibre-fs} 中证明
(2)、(3) 与 (4) 等价。显然 (1) 推出 (2)。
因此，假设 (4) 成立，并证明 (1)。

\medskip\noindent
由引理 \ref{more-algebra-lemma-formally-smooth-completion}，
$A^\wedge \to B^\wedge$ 形式光滑。由命题
\ref{more-algebra-proposition-fs-regular}，$A^\wedge \to B^\wedge$ 正则。
依假设，$A \to A^\wedge$ 正则，故由引理
\ref{more-algebra-lemma-regular-composition}，$A \to B^\wedge$ 正则。
由于 $B \to B^\wedge$ 忠实平坦，由引理
\ref{more-algebra-lemma-regular-permanence} 可知 $A \to B$ 正则。
\end{proof}




\section{G-环}
\label{more-algebra-section-G-ring}

\noindent
设 $A$ 为 Noether 局部环。在第
\ref{more-algebra-section-permanence-completion} 节中，
我们已经看到，$A$ 的某些性质（但并非所有性质）会反映在
$A$ 的完备化 $A^\wedge$ 上。为了进一步研究这一现象，
我们引入一些术语。对于 $A$ 的素理想 $\mathfrak q$，纤维环
$$
A^\wedge \otimes_A \kappa(\mathfrak q) =
(A^\wedge)_\mathfrak q/\mathfrak q(A^\wedge)_\mathfrak q =
(A/\mathfrak q)^\wedge \otimes_{A/q} \kappa(\mathfrak q)
$$
称为 $A$ 的一个{\it 形式纤维}。我们把形式纤维看成
$\kappa(\mathfrak q)$-代数。因此，$A \to A^\wedge$ 是正则环同态，
当且仅当所有形式纤维都是几何正则代数。

\begin{definition}
\label{more-algebra-definition-G-ring}
环 $R$ 称为{\it G-环}，如果 $R$ 是 Noether 环，并且对于
$R$ 的每个素理想 $\mathfrak p$，环同态
$R_\mathfrak p \to (R_\mathfrak p)^\wedge$ 都是正则的。
\end{definition}

\noindent
由上面的讨论可知，$R$ 是 G-环，当且仅当每个局部环
$R_\mathfrak p$ 都有几何正则的形式纤维。注意，若
$\mathbf{Q} \subset R$，则只需检验形式纤维正则。
以下引理给出 G-环条件的另一种表述。

\begin{lemma}
\label{more-algebra-lemma-check-G-ring-easy}
设 $R$ 为 Noether 环。那么，$R$ 是 G-环，当且仅当对于每一对素理想
$\mathfrak q \subset \mathfrak p \subset R$，代数
$$
(R/\mathfrak q)_\mathfrak p^\wedge \otimes_{R/\mathfrak q} \kappa(\mathfrak q)
$$
在 $\kappa(\mathfrak q)$ 上几何正则。
\end{lemma}

\begin{proof}
这由以下事实立即得到：作为 $\kappa(\mathfrak q)$-代数，
$$
R_\mathfrak p^\wedge \otimes_R \kappa(\mathfrak q) =
(R/\mathfrak q)_\mathfrak p^\wedge \otimes_{R/\mathfrak q} \kappa(\mathfrak q)
$$
\end{proof}

\begin{lemma}
\label{more-algebra-lemma-G-ring-goes-up-quasi-finite}
设 $R \to R'$ 为 Noether 环之间的有限型同态，并设
$$
\xymatrix{
\mathfrak q' \ar[r] & \mathfrak p' \ar[r] & R' \\
\mathfrak q \ar[r] \ar@{-}[u] &
\mathfrak p \ar[r] \ar@{-}[u] & R \ar[u]
}
$$
为素理想图。假设 $R \to R'$ 在 $\mathfrak p'$ 处拟有限。
\begin{enumerate}
\item 若形式纤维 $R_\mathfrak p^\wedge \otimes_R \kappa(\mathfrak q)$
在 $\kappa(\mathfrak q)$ 上几何正则，则形式纤维
$(R'_{\mathfrak p'})^\wedge \otimes_{R'} \kappa(\mathfrak q')$
在 $\kappa(\mathfrak q')$ 上几何正则。
\item 若 $R_\mathfrak p$ 的形式纤维均几何正则，
则 $R'_{\mathfrak p'}$ 的形式纤维均几何正则。
\item 若 $R \to R'$ 拟有限且 $R$ 是 G-环，则 $R'$ 是 G-环。
\end{enumerate}
\end{lemma}

\begin{proof}
显然，(1) $\Rightarrow$ (2) $\Rightarrow$ (3)。假设
$R_\mathfrak p^\wedge \otimes_R \kappa(\mathfrak q)$
在 $\kappa(\mathfrak q)$ 上几何正则。由 Algebra，引理
\ref{algebra-lemma-completion-at-quasi-finite-prime}，可知
$$
R_\mathfrak p^\wedge \otimes_R R'
=
(R'_{\mathfrak p'})^\wedge \times B
$$
对某个 $R_\mathfrak p^\wedge$-代数 $B$ 成立。因此，
$R'_{\mathfrak p'} \to (R'_{\mathfrak p'})^\wedge$ 是同态
$R_\mathfrak p \to R_\mathfrak p^\wedge$ 的某个基变换的一个因子。
于是 $(R'_{\mathfrak p'})^\wedge \otimes_{R'} \kappa(\mathfrak q')$
是
$$
R_\mathfrak p^\wedge \otimes_R R' \otimes_{R'} \kappa(\mathfrak q') =
R_\mathfrak p^\wedge \otimes_R \kappa(\mathfrak q)
\otimes_{\kappa(\mathfrak q)} \kappa(\mathfrak q').
$$
的一个因子。由于基域扩张保持几何正则性，结论成立；见
Algebra，引理 \ref{algebra-lemma-geometrically-regular}。
\end{proof}

\begin{lemma}
\label{more-algebra-lemma-check-G-ring}
设 $R$ 为 Noether 环。那么，$R$ 是 G-环，当且仅当对于每个有限自由环同态
$R \to S$，$S$ 的形式纤维都是正则环。
\end{lemma}

\begin{proof}
假设对于任意有限自由环同态 $R \to S$，环 $S$ 都有正则形式纤维。
设 $\mathfrak q \subset \mathfrak p \subset R$ 为素理想，并设
$\kappa(\mathfrak q) \subset L$ 为有限纯不可分扩张。要证明 $R$ 是 G-环，
只需证明
$$
R_\mathfrak p^\wedge \otimes_R \kappa(\mathfrak q)
\otimes_{\kappa(\mathfrak q)} L
$$
为正则环。选取有限自由扩张 $R \to R'$，使得
$\mathfrak q' = \mathfrak qR'$ 是素理想，并且
$\kappa(\mathfrak q')$ 在 $\kappa(\mathfrak q)$ 上同构于 $L$；见
Algebra，引理 \ref{algebra-lemma-finite-free-given-residue-field-extension}。
由 Algebra，引理 \ref{algebra-lemma-completion-finite-extension}，有
$$
R_\mathfrak p^\wedge \otimes_R R' = \prod (R'_{\mathfrak p_i'})^\wedge
$$
其中 $\mathfrak p_i'$ 是 $R'$ 中位于 $\mathfrak p$ 上方的素理想。因此，
$$
R_\mathfrak p^\wedge \otimes_R \kappa(\mathfrak q)
\otimes_{\kappa(\mathfrak q)} L =
R_\mathfrak p^\wedge \otimes_R R'
\otimes_{R'} \kappa(\mathfrak q')
=
\prod (R'_{\mathfrak p_i'})^\wedge
\otimes_{R'_{\mathfrak p'_i}} \kappa(\mathfrak q')
$$
依假设，右侧各环正则，故左侧的环也正则。因此 $R$ 是 G-环。
反向推论由引理 \ref{more-algebra-lemma-G-ring-goes-up-quasi-finite} 得到。
\end{proof}

\begin{lemma}
\label{more-algebra-lemma-helper-G-ring}
设 $k$ 为特征 $p$ 的域。令
$A = k[[x_1, \ldots, x_n]][y_1, \ldots, y_m]$，并以 $K$
表示 $A$ 的分式域。设 $\mathfrak p \subset A$ 为素理想。那么
$A_\mathfrak p^\wedge \otimes_A K$ 在 $K$ 上几何正则。
\end{lemma}

\begin{proof}
设 $L/K$ 为有限纯不可分域扩张。我们将对 $[L : K]$ 归纳证明
$A_\mathfrak p^\wedge \otimes L$ 正则。基始情形为 $L = K$：
由于 $A$ 正则，$A_\mathfrak p^\wedge$ 正则
（引理 \ref{more-algebra-lemma-completion-regular}），故其局部化
$A_\mathfrak p^\wedge \otimes K$ 正则。取中间域
$K \subset M \subset L$，使得 $L$ 是通过添加某个 $f \in M$
的一个 $p$ 次根而得到的 $M$ 的 $p$ 次扩张。令 $B$ 为 $M$
的一个有限 $A$-子代数，其分式域为 $M$。清去分母后，
我们可以并且确实假设 $f \in B$。置 $C = B[z]/(z^p -f)$，
并注意 $B \subset C$ 有限，且 $C$ 的分式域为 $L$。由于
$A \subset B \subset C$ 均为有限扩张且 $L/M/K$ 均为纯不可分扩张，
可知 $B$ 或 $C$ 的每个元素都有某个幂属于 $A$。因此，
存在唯一素理想 $\mathfrak r \subset B$，分别地，也存在唯一素理想
$\mathfrak q \subset C$ 位于 $\mathfrak p$ 上方。注意
$$
A_\mathfrak p^\wedge \otimes_A M = B_\mathfrak r^\wedge \otimes_B M
$$
见 Algebra，引理 \ref{algebra-lemma-completion-finite-extension}。
由归纳假设，该环正则。同理有
$$
A_\mathfrak p^\wedge \otimes_A L =
C_\mathfrak r^\wedge \otimes_C L =
B_\mathfrak r^\wedge \otimes_B M[z]/(z^p - f)
$$
最后一个等式成立，是因为 $C = B[z]/(z^p - f)$ 的完备化等于
$B_\mathfrak r^\wedge[z]/(z^p -f)$。由引理 \ref{more-algebra-lemma-find-D}，
存在导子 $D : B \to B$，使得 $D(f) \not = 0$。换言之，
$g = D(f)$ 是 $M$ 中的单位！由引理 \ref{more-algebra-lemma-derivation-extends}，
$D$ 依次通过局部化、完备化和局部化延拓为 $B_\mathfrak r$、
$B_\mathfrak r^\wedge$ 以及 $B_\mathfrak r^\wedge \otimes_B M$ 上的导子。
由于这些是延拓，我们最终得到 $B_\mathfrak r^\wedge \otimes_B M$
上的一个导子，它将 $f$ 映到 $g$，而 $g$ 是环
$B_\mathfrak r^\wedge \otimes_B M$ 的单位。
于是由引理 \ref{more-algebra-lemma-degree-p-extension-regular}，
$A_\mathfrak p^\wedge \otimes_A L$ 正则，结论得证。
\end{proof}

\begin{proposition}
\label{more-algebra-proposition-Noetherian-complete-G-ring}
Noether 完备局部环是 G-环。
\end{proposition}

\begin{proof}
设 $A$ 为 Noether 完备局部环。由引理
\ref{more-algebra-lemma-check-G-ring-easy}，只需检验 $B = A/\mathfrak q$
在 $B$ 的极小素理想 $(0)$ 上方具有几何正则的形式纤维。
因此，可以假设 $A$ 为整环，并且只需检验 $A$ 的极小素理想
$(0)$ 上方的形式纤维满足该条件。令 $K$ 为 $A$ 的分式域。

\medskip\noindent
可以选取子环 $A_0 \subset A$，它是正则完备局部环且
$A$ 在 $A_0$ 上有限；见 Algebra，引理
\ref{algebra-lemma-complete-local-Noetherian-domain-finite-over-regular}。
此外，可以假设 $A_0$ 是域或 Cohen 环上的幂级数环。
由引理 \ref{more-algebra-lemma-G-ring-goes-up-quasi-finite}，只需对 $A_0$ 证明结论。

\medskip\noindent
假设 $A$ 是域或 Cohen 环上的幂级数环。由于 $A$ 正则，
各局部化 $A_\mathfrak p$ 均正则（见 Algebra，定义
\ref{algebra-definition-regular} 及其前面的讨论）。因此，
各完备化 $A_\mathfrak p^\wedge$ 均正则；见引理
\ref{more-algebra-lemma-completion-regular}。于是纤维
$A_{\mathfrak p}^\wedge \otimes_A K$ 作为
$A_\mathfrak p^\wedge$ 的局部化也正则。所以当 $K$ 的特征为 $0$
时结论成立。正特征情形就是 $A = k[[x_1, \ldots, x_d]]$ 的情形，
它是引理 \ref{more-algebra-lemma-helper-G-ring} 的特殊情形。
\end{proof}

\begin{lemma}
\label{more-algebra-lemma-check-G-ring-maximal-ideals}
设 $R$ 为 Noether 环。那么，$R$ 是 G-环，当且仅当对于 $R$
的每个极大理想 $\mathfrak m$，$R_\mathfrak m$ 都有几何正则的形式纤维。
\end{lemma}

\begin{proof}
假设对于 $R$ 的每个极大理想 $\mathfrak m$，同态
$R_\mathfrak m \to R_\mathfrak m^\wedge$ 都正则。设 $\mathfrak p$
为 $R$ 的素理想，并选取极大理想 $\mathfrak p \subset \mathfrak m$。
由于 $R_\mathfrak m \to R_\mathfrak m^\wedge$ 忠实平坦，
可以在 $R_\mathfrak m^\wedge$ 中选取一个位于
$\mathfrak pR_\mathfrak m$ 上方的素理想 $\mathfrak p'$。
考虑交换图
$$
\xymatrix{
R_\mathfrak m^\wedge \ar[r] &
(R_\mathfrak m^\wedge)_{\mathfrak p'} \ar[r] &
(R_\mathfrak m^\wedge)_{\mathfrak p'}^\wedge
\\
R_\mathfrak m \ar[u] \ar[r] & R_\mathfrak p \ar[u] \ar[r] &
R_\mathfrak p^\wedge \ar[u]
}
$$
依假设，环同态 $R_\mathfrak m \to R_\mathfrak m^\wedge$ 正则。
由命题 \ref{more-algebra-proposition-Noetherian-complete-G-ring}，
$(R_\mathfrak m^\wedge)_{\mathfrak p'} \to
(R_\mathfrak m^\wedge)_{\mathfrak p'}^\wedge$ 正则。
局部化 $R_\mathfrak m^\wedge \to
(R_\mathfrak m^\wedge)_{\mathfrak p'}$ 正则。因此，由引理
\ref{more-algebra-lemma-regular-composition}，$R_\mathfrak m \to
(R_\mathfrak m^\wedge)_{\mathfrak p'}^\wedge$ 正则。
由于它分解经过局部化 $R_\mathfrak p$，环同态
$R_\mathfrak p \to (R_\mathfrak m^\wedge)_{\mathfrak p'}^\wedge$
也正则。因此，可以应用引理 \ref{more-algebra-lemma-regular-permanence} 得到
$R_\mathfrak p \to R_\mathfrak p^\wedge$ 正则。
\end{proof}

\begin{lemma}
\label{more-algebra-lemma-henselization-G-ring}
设 $R$ 为 Noether 局部 G-环。那么，Hensel 化 $R^h$
和严格 Hensel 化 $R^{sh}$ 都是 G-环。
\end{lemma}

\begin{proof}
我们将使用引理 \ref{more-algebra-lemma-check-G-ring-maximal-ideals} 的判据。
设 $\mathfrak q \subset R^h$ 为素理想，并置
$\mathfrak p = R \cap \mathfrak q$。置 $\mathfrak q_1 = \mathfrak q$，
并令 $\mathfrak q_2, \ldots, \mathfrak q_t$ 为 $R^h$ 中位于
$\mathfrak p$ 上方的其余素理想，于是
$R^h \otimes_R \kappa(\mathfrak p) =
\prod\nolimits_{i = 1, \ldots, t} \kappa(\mathfrak q_i)$；见引理
\ref{more-algebra-lemma-fibres-henselization}。利用 $(R^h)^\wedge = R^\wedge$
（引理 \ref{more-algebra-lemma-henselization-noetherian}），可得
$$
\prod\nolimits_{i = 1, \ldots, t}
(R^h)^\wedge \otimes_{R^h} \kappa(\mathfrak q_i) =
(R^h)^\wedge \otimes_{R^h} (R^h \otimes_R \kappa(\mathfrak p)) =
R^\wedge \otimes_R \kappa(\mathfrak p)
$$
依假设，$(R^h)^\wedge \otimes_{R^h} \kappa(\mathfrak q_i)$
在 $\kappa(\mathfrak p)$ 上几何正则。由于
$\kappa(\mathfrak q_i)$ 是 $\kappa(\mathfrak p)$ 的可分代数扩张，
由 Algebra，引理
\ref{algebra-lemma-geometrically-regular-over-separable-algebraic}，可知
$(R^h)^\wedge \otimes_{R^h} \kappa(\mathfrak q_i)$
在 $\kappa(\mathfrak q_i)$ 上几何正则。

\medskip\noindent
设 $\mathfrak r \subset R^{sh}$ 为素理想，并置
$\mathfrak p = R \cap \mathfrak r$。置 $\mathfrak r_1 = \mathfrak r$，
并令 $\mathfrak r_2, \ldots, \mathfrak r_s$ 为 $R^{sh}$ 中位于
$\mathfrak p$ 上方的其余素理想，于是
$R^{sh} \otimes_R \kappa(\mathfrak p) =
\prod\nolimits_{i = 1, \ldots, s} \kappa(\mathfrak r_i)$；见引理
\ref{more-algebra-lemma-fibres-henselization}。于是可得
$$
\prod\nolimits_{i = 1, \ldots, s}
(R^{sh})^\wedge \otimes_{R^{sh}} \kappa(\mathfrak r_i) =
(R^{sh})^\wedge \otimes_{R^{sh}} (R^{sh} \otimes_R \kappa(\mathfrak p)) =
(R^{sh})^\wedge \otimes_R \kappa(\mathfrak p)
$$
注意，由引理 \ref{more-algebra-lemma-henselization-noetherian}，
$R^\wedge \to (R^{sh})^\wedge$ 对
$\mathfrak m_{(R^{sh})^\wedge}$-adic 拓扑形式光滑。因此，由命题
\ref{more-algebra-proposition-fs-regular}，$R^\wedge \to (R^{sh})^\wedge$ 正则。
依假设，$R^\wedge \otimes_R \kappa(\mathfrak p)$ 在
$\kappa(\mathfrak p)$ 上正则，故由引理
\ref{more-algebra-lemma-regular-composition}，可知
$(R^{sh})^\wedge \otimes_{R^{sh}} \kappa(\mathfrak r_i)$
在 $\kappa(\mathfrak p)$ 上正则。由于 $\kappa(\mathfrak r_i)$
是 $\kappa(\mathfrak p)$ 的可分代数扩张，由 Algebra，引理
\ref{algebra-lemma-geometrically-regular-over-separable-algebraic}，可知
$(R^{sh})^\wedge \otimes_{R^{sh}} \kappa(\mathfrak r_i)$
在 $\kappa(\mathfrak r_i)$ 上几何正则。
\end{proof}

\begin{lemma}
\label{more-algebra-lemma-another-helper-G-ring}
设 $p$ 为素数。设 $A$ 为 Noether 完备局部整环，其分式域 $K$
的特征为 $p$。设 $\mathfrak q \subset A[x]$ 为位于 $A$
的极大理想上方的极大理想，并设
$(0) \not = \mathfrak r \subset \mathfrak q$ 为位于
$(0) \subset A$ 上方的素理想。那么，
$A[x]_\mathfrak q^\wedge \otimes_{A[x]} \kappa(\mathfrak r)$ 在
$\kappa(\mathfrak r)$ 上几何正则。
\end{lemma}

\begin{proof}
注意，$K \subset \kappa(\mathfrak r)$ 是有限扩张。因此，给定有限纯不可分扩张
$L/\kappa(\mathfrak r)$，存在 Noether 完备局部整环的有限扩张
$A \subset B$，使得 $\kappa(\mathfrak r) \otimes_A B$ 满射到 $L$。
具体地，取 $B \subset L$ 为有限 $A$-子代数，其分式域为 $L$。
以 $\mathfrak r' \subset B[x]$ 表示映射
$B[x] = A[x] \otimes_A B \to \kappa(\mathfrak r) \otimes_A B \to L$
的核，于是 $\kappa(\mathfrak r') = L$。那么
$$
A[x]_\mathfrak q^\wedge \otimes_{A[x]} L =
A[x]_\mathfrak q^\wedge \otimes_{A[x]} B[x]
\otimes_{B[x]} \kappa(\mathfrak r') =
\prod B[x]_{\mathfrak q_i}^\wedge \otimes_{B[x]} \kappa(\mathfrak r')
$$
其中 $\mathfrak q_1, \ldots, \mathfrak q_t$ 是 $B[x]$ 中位于
$\mathfrak q$ 上方的素理想；见 Algebra，引理
\ref{algebra-lemma-completion-finite-extension}。因此，只需证明各环
$B[x]_{\mathfrak q_i}^\wedge \otimes_{B[x]} \kappa(\mathfrak r')$
均正则。这就把问题化约为在 $K = \kappa(\mathfrak r)$
这一特殊情形下证明
$A[x]_\mathfrak q^\wedge \otimes_{A[x]} \kappa(\mathfrak r)$ 正则。

\medskip\noindent
假设 $K = \kappa(\mathfrak r)$。在此情形下，
$\mathfrak r K[x]$ 由某个 $f \in K$ 的 $x - f$ 生成，并且
$$
A[x]_\mathfrak q^\wedge \otimes_{A[x]} \kappa(\mathfrak r)
=
(A[x]_\mathfrak q^\wedge \otimes_A K)/(x - f)
$$
$A[x]$ 上的导子 $D = \text{d}/\text{d}x$ 延拓到 $K[x]$，
并将 $x - f$ 映到 $K[x]$ 的一个单位。此外，由引理
\ref{more-algebra-lemma-derivation-extends}，$D$ 延拓到
$A[x]_\mathfrak q^\wedge \otimes_A K$。由于
$A \to A[x]_\mathfrak q^\wedge$ 形式光滑（见引理
\ref{more-algebra-lemma-formally-smooth} 和
\ref{more-algebra-lemma-formally-smooth-completion}），由命题
\ref{more-algebra-proposition-fs-regular}，环
$A[x]_\mathfrak q^\wedge \otimes_A K$ 正则（在这一特殊情形下，
该命题证明中的论证显著简化）。最后应用引理
\ref{more-algebra-lemma-quotient-regular} 即得结论。
\end{proof}

\begin{proposition}
\label{more-algebra-proposition-finite-type-over-G-ring}
设 $R$ 为 G-环。若 $R \to S$ 本质有限型，则 $S$ 是 G-环。
\end{proposition}

\begin{proof}
由于 G-环性质是局部环的性质，显然 G-环的局部化仍为 G-环。
反之，若在每个素理想处的局部化都是 G-环，则该环是 G-环。
因此，只需对每个有限型 $R$-代数 $S$ 及 $S$ 的每个素理想
$\mathfrak q$ 证明 $S_\mathfrak q$ 是 G-环。将 $S$ 写成
$R[x_1, \ldots, x_n]$ 的商，由引理
\ref{more-algebra-lemma-G-ring-goes-up-quasi-finite} 可知，只需证明
$R[x_1, \ldots, x_n]$ 是 G-环。对 $n$ 归纳，只需证明
$R[x]$ 是 G-环。设 $\mathfrak q \subset R[x]$ 为极大理想。
由引理 \ref{more-algebra-lemma-check-G-ring-maximal-ideals}，只需证明
$$
R[x]_\mathfrak q \longrightarrow R[x]_\mathfrak q^\wedge
$$
正则。若 $\mathfrak q$ 位于 $\mathfrak p \subset R$ 上方，
则可以用 $R_\mathfrak p$ 代替 $R$。因此，可以假设 $R$
是极大理想为 $\mathfrak m$ 的 Noether 局部 G-环，并且
$\mathfrak q \subset R[x]$ 位于 $\mathfrak m$ 上方。注意，存在唯一的素理想
$\mathfrak q' \subset R^\wedge[x]$ 位于 $\mathfrak q$ 上方。考虑交换图
$$
\xymatrix{
R[x]_\mathfrak q^\wedge \ar[r] &
(R^\wedge[x]_{\mathfrak q'})^\wedge \\
R[x]_\mathfrak q \ar[r] \ar[u] & R^\wedge[x]_{\mathfrak q'} \ar[u]
}
$$
由于 $R$ 是 G-环，下方水平箭头正则（它是正则环同态
$R \to R^\wedge$ 的某个基变换的局部化）。假设可以证明右侧竖直箭头
正则。那么复合 $R[x]_\mathfrak q \to
(R^\wedge[x]_{\mathfrak q'})^\wedge$ 正则，从而由引理
\ref{more-algebra-lemma-regular-permanence}，左侧竖直箭头正则。
因此，可以假设 $R$ 是 Noether 完备局部环，并且 $\mathfrak q$
是位于 $R$ 的极大理想上方的素理想。

\medskip\noindent
设 $R$ 为 Noether 完备局部环，并设 $\mathfrak q \subset R[x]$
为位于 $R$ 的极大理想上方的极大理想。设
$\mathfrak r \subset \mathfrak q$ 为素理想。我们要证明
$R[x]_\mathfrak q^\wedge \otimes_{R[x]} \kappa(\mathfrak r)$
是 $\kappa(\mathfrak r)$ 上的几何正则代数。置
$\mathfrak p = R \cap \mathfrak r$。由引理
\ref{more-algebra-lemma-check-G-ring-easy}，可以用 $R/\mathfrak p$ 代替 $R$，
并用 $\mathfrak q$ 和 $\mathfrak r$ 在 $R/\mathfrak p[x]$ 中的像
代替它们。因此，可以假设 $R$ 是整环且 $\mathfrak r \cap R = (0)$。

\medskip\noindent
由 Algebra，引理
\ref{algebra-lemma-complete-local-Noetherian-domain-finite-over-regular}，
可以找到正则子环 $R_0 \subset R$，使得 $R$ 在 $R_0$ 上有限。
应用引理 \ref{more-algebra-lemma-G-ring-goes-up-quasi-finite} 可知，只需在上述假设之外
再假设 $R$ 是正则完备局部环，并证明
$R[x]_\mathfrak q^\wedge \otimes_{R[x]} \kappa(\mathfrak r)$
在 $\kappa(\mathfrak r)$ 上几何正则。

\medskip\noindent
现在 $R$ 是正则完备局部环，且有
$\mathfrak r \subset \mathfrak q \subset R[x]$、
$(0) = R \cap \mathfrak r$，并且 $\mathfrak q$ 是位于 $R$
的极大理想上方的极大理想。由于 $R$ 正则，由 Algebra，引理
\ref{algebra-lemma-regular-goes-up}，环 $R[x]$ 正则。
因此，局部化 $R[x]_\mathfrak q$ 正则。于是，完备化
$R[x]_\mathfrak q^\wedge$ 正则；见引理
\ref{more-algebra-lemma-completion-regular}。因此，纤维
$R[x]_{\mathfrak q}^\wedge \otimes_{R[x]} \kappa(\mathfrak r)$
作为 $R[x]_\mathfrak q^\wedge$ 的局部化也正则。所以，当 $R$
的分式域的特征为 $0$ 时，结论成立。

\medskip\noindent
若 $R$ 的特征为正，则 $R = k[[x_1, \ldots, x_n]]$。
在这种情况下，我们将论证分为两个子情形：
\begin{enumerate}
\item 情形 $\mathfrak r = (0)$。结论直接由引理
\ref{more-algebra-lemma-helper-G-ring} 得到。
\item 情形 $\mathfrak r \not = (0)$。这就是引理
\ref{more-algebra-lemma-another-helper-G-ring}。
\end{enumerate}
\end{proof}

\begin{remark}
\label{more-algebra-remark-G-does-not-survive-completion}
设 $R$ 为 G-环，并设 $I \subset R$ 为理想。一般而言，
$I$-adic 完备化 $R^\wedge$ 不一定是 G-环。Nishimura 在
\cite{Nishimura} 中给出了一个例子。该例子的推广以及某种意义上的澄清，
可见 \cite{Dumitrescu} 的最后一节。
\end{remark}

\begin{proposition}
\label{more-algebra-proposition-ubiquity-G-ring}
下列各类环都是 G-环：
\begin{enumerate}
\item 域；
\item Noether 完备局部环；
\item $\mathbf{Z}$；
\item 分式域特征为零的 Dedekind 整环；
\item 上述任一类环的有限型环扩张。
\end{enumerate}
\end{proposition}

\begin{proof}
对于域、$\mathbf{Z}$ 和特征为零的 Dedekind 整环，
结论立即由定义以及离散赋值环的完备化仍为离散赋值环这一事实得到。
由命题 \ref{more-algebra-proposition-Noetherian-complete-G-ring}，
Noether 完备局部环是 G-环。关于有限型扩张环的陈述就是命题
\ref{more-algebra-proposition-finite-type-over-G-ring}。
\end{proof}

\begin{lemma}
\label{more-algebra-lemma-henselian-local-limit-G-rings}
设 $(A, \mathfrak m)$ 为 Hensel 局部环。那么，$A$ 是一个由
Hensel 局部 G-环构成且转移同态均为局部同态的系统的滤过余极限。
\end{lemma}

\begin{proof}
将 $A = \colim A_i$ 写成有限型 $\mathbf{Z}$-代数的滤过余极限。
设 $\mathfrak p_i$ 为 $A_i$ 中位于 $\mathfrak m$ 下方的素理想。
可以用 $A_i$ 在 $\mathfrak p_i$ 处的局部化代替 $A_i$。
此时 $A_i$ 是 Noether 局部 G-环（命题
\ref{more-algebra-proposition-ubiquity-G-ring}）。由引理
\ref{more-algebra-lemma-henselization-colimit}，可知 $A = \colim A_i^h$。
由引理 \ref{more-algebra-lemma-henselization-G-ring}，各环 $A_i^h$ 都是 G-环。
\end{proof}

\begin{lemma}
\label{more-algebra-lemma-map-G-ring-to-completion-regular}
\begin{reference}
\cite[Theorem 79]{MatCA}
\end{reference}
设 $A$ 为 G-环。设 $I \subset A$ 为理想，并设 $A^\wedge$
为 $A$ 关于 $I$ 的完备化。那么 $A \to A^\wedge$ 正则。
\end{lemma}

\begin{proof}
由 Algebra，引理 \ref{algebra-lemma-completion-flat}，环同态
$A \to A^\wedge$ 平坦。由 Algebra，引理
\ref{algebra-lemma-completion-Noetherian-Noetherian}，环 $A^\wedge$
是 Noether 环。因此，只需检验引理 \ref{more-algebra-lemma-regular-local}
的第三个条件。设 $\mathfrak m' \subset A^\wedge$ 为位于
$\mathfrak m \subset A$ 上方的极大理想。由 Algebra，引理
\ref{algebra-lemma-radical-completion}，有
$IA^\wedge \subset \mathfrak m'$。由于 $A^\wedge/IA^\wedge = A/I$，
可知 $I \subset \mathfrak m$、
$\mathfrak m/I = \mathfrak m'/IA^\wedge$，并且
$A/\mathfrak m = A^\wedge/\mathfrak m'$。由于
$A^\wedge/\mathfrak m'$ 是域，可知 $\mathfrak m$ 也是极大理想。
于是，$A_\mathfrak m \to A^\wedge_{\mathfrak m'}$ 是 Noether
局部环之间的平坦局部同态，它诱导剩余域的同构，并满足
$\mathfrak m A^\wedge_{\mathfrak m'} =
\mathfrak m'A^\wedge_{\mathfrak m'}$。因此，由引理
\ref{more-algebra-lemma-flat-unramified}，它在完备局部环上诱导同构。
设 $(A_\mathfrak m)^\wedge$ 为 $A_\mathfrak m$ 关于其极大理想的完备化。
环同态
$$
(A^\wedge)_{\mathfrak m'} \to
((A^\wedge)_{\mathfrak m'})^\wedge = (A_\mathfrak m)^\wedge
$$
忠实平坦（Algebra，引理
\ref{algebra-lemma-completion-faithfully-flat}）。因此，可以对环同态
$$
A_\mathfrak m \to (A^\wedge)_{\mathfrak m'} \to (A_\mathfrak m)^\wedge
$$
应用引理 \ref{more-algebra-lemma-regular-permanence} 得到结论，因为 $A$ 是 G-环，
所以 $A_\mathfrak m \to (A_\mathfrak m)^\wedge$ 正则。
\end{proof}

\begin{lemma}
\label{more-algebra-lemma-henselization-pair-G-ring}
\begin{slogan}
G-环性质在沿理想的 Hensel 化下稳定
\end{slogan}
\begin{reference}
\cite[Theorem 5.3 i)]{Greco}
\end{reference}
设 $A$ 为 G-环。设 $I \subset A$ 为理想。设 $(A^h, I^h)$
为偶对 $(A, I)$ 的 Hensel 化；见引理 \ref{more-algebra-lemma-henselization}。
那么 $A^h$ 是 G-环。
\end{lemma}

\begin{proof}
设 $\mathfrak m^h \subset A^h$ 为极大理想。由引理
\ref{more-algebra-lemma-check-G-ring-maximal-ideals}，我们必须证明从
$A^h_{\mathfrak m^h}$ 到其完备化的同态具有几何正则的纤维。
设 $\mathfrak m$ 为 $\mathfrak m^h$ 在 $A$ 中的逆像。注意，
$I^h \subset \mathfrak m^h$，从而 $I \subset \mathfrak m$，
因为 $(A^h, I^h)$ 是 Hensel 偶对。回顾：$A^h$ 是 Noether 环，
$I^h = IA^h$，并且 $A \to A^h$ 在 $I$-adic 完备化上诱导同构；
见引理 \ref{more-algebra-lemma-henselization-Noetherian-pair}。于是，Noether
局部环之间的局部同态
$$
A_\mathfrak m \to A^h_{\mathfrak m^h}
$$
由引理 \ref{more-algebra-lemma-flat-unramified} 在极大理想完备化上诱导同构
（略去细节）。设 $\mathfrak q^h$ 为 $A^h_{\mathfrak m^h}$
中位于 $\mathfrak q \subset A_\mathfrak m$ 上方的素理想。
置 $\mathfrak q_1 = \mathfrak q^h$，并令
$\mathfrak q_2, \ldots, \mathfrak q_t$ 为 $A^h$ 中位于
$\mathfrak q$ 上方的其余素理想，于是
$A^h \otimes_A \kappa(\mathfrak q) =
\prod\nolimits_{i = 1, \ldots, t} \kappa(\mathfrak q_i)$；见引理
\ref{more-algebra-lemma-filtered-colimit-etale-noetherian-fibres}。利用上面讨论过的
$(A^h)_{\mathfrak m^h}^\wedge = (A_\mathfrak m)^\wedge$，可得
$$
\prod\nolimits_{i = 1, \ldots, t}
(A^h_{\mathfrak m^h})^\wedge \otimes_{A^h_{\mathfrak m^h}}
\kappa(\mathfrak q_i) =
(A^h_{\mathfrak m^h})^\wedge \otimes_{A^h_{\mathfrak m^h}}
(A^h_{\mathfrak m^h} \otimes_{A_{\mathfrak m}} \kappa(\mathfrak q)) =
(A_{\mathfrak m})^\wedge \otimes_{A_{\mathfrak m}} \kappa(\mathfrak q)
$$
因此，作为其中一个因子，环
$$
(A^h_{\mathfrak m^h})^\wedge \otimes_{A^h_{\mathfrak m^h}}
\kappa(\mathfrak q^h)
$$
依关于 $A$ 的假设在 $\kappa(\mathfrak q)$ 上几何正则。
由于 $\kappa(\mathfrak q^h)$ 是 $\kappa(\mathfrak q)$
的可分代数扩张，由 Algebra，引理
\ref{algebra-lemma-geometrically-regular-over-separable-algebraic}，可知
$$
(A^h_{\mathfrak m^h})^\wedge \otimes_{A^h_{\mathfrak m^h}}
\kappa(\mathfrak q^h)
$$
如所需地在 $\kappa(\mathfrak q^h)$ 上几何正则。
\end{proof}








\section{形式纤维的性质}
\label{more-algebra-section-properties-formal-fibres}

\noindent
在本节中，我们重新处理第 \ref{more-algebra-section-G-ring} 节中的一些论证，
以便能够准确讨论 Noether 环的形式纤维所具有的性质。

\medskip\noindent
设 $P$ 是域 $k$ 到 Noether 环 $R$ 的环同态 $k \to R$
的一种性质。对于 $B$ 为 Noether 环的环同态 $A \to B$，
若对 $A$ 的每个素理想 $\mathfrak q$，同态
$\kappa(\mathfrak q) \to B \otimes_A \kappa(\mathfrak q)$ 都具有 $P$，
则称 $P$ 对该环同态的纤维成立。以下将使用这些断言：
\begin{enumerate}
\item[(A)] 对有限生成域扩张 $k'/k$，有
$P(k \to R) \Rightarrow P(k' \to R \otimes_k k')$；
\item[(B)] $P(k \to R_\mathfrak p),\ \forall \mathfrak p \in \Spec(R)
\Leftrightarrow P(k \to R)$；
\item[(C)] 给定 Noether 环的平坦同态 $A \to B \to C$，
若 $A \to B$ 的纤维具有 $P$ 且 $B \to C$ 正则，
则 $A \to C$ 的纤维具有 $P$；
\item[(D)] 给定 Noether 环的平坦同态 $A \to B \to C$，
若 $A \to C$ 的纤维具有 $P$ 且 $B \to C$ 忠实平坦，
则 $A \to B$ 的纤维具有 $P$；
\item[(E)] 给定 $k \to k' \to R$，其中 $R$ 为 Noether 环，
若 $k'/k$ 是可分代数扩张且 $P(k \to R)$ 成立，
则 $P(k' \to R)$ 成立；以及
\item[(F)] 在此添加更多内容。
\end{enumerate}
给定 Noether 局部环 $A$，若 $P$ 对 $A \to A^\wedge$ 的纤维成立，
则称“{\it $A$ 的形式纤维具有 $P$}”。若 $R$ 是 Noether 环，
且对 $R$ 的所有素理想 $\mathfrak p$，$R_\mathfrak p$ 的形式纤维
都具有 $P$，则称{\it $R$ 是 $P$-环}。

\begin{lemma}
\label{more-algebra-lemma-check-P-ring-easy}
设 $R$ 为 Noether 环，并设 $P$ 为上述性质。那么，$R$ 是 $P$-环，
当且仅当对于每一对素理想 $\mathfrak q \subset \mathfrak p \subset R$，
$\kappa(\mathfrak q)$-代数
$$
(R/\mathfrak q)_\mathfrak p^\wedge \otimes_{R/\mathfrak q} \kappa(\mathfrak q)
$$
都具有性质 $P$。
\end{lemma}

\begin{proof}
这由以下事实立即得到：作为 $\kappa(\mathfrak q)$-代数，
$$
R_\mathfrak p^\wedge \otimes_R \kappa(\mathfrak q) =
(R/\mathfrak q)_\mathfrak p^\wedge \otimes_{R/\mathfrak q} \kappa(\mathfrak q)
$$
\end{proof}

\begin{lemma}
\label{more-algebra-lemma-P-local}
设 $R \to \Lambda$ 为 Noether 环同态。假设 $P$ 具有性质 (B)。
下列条件等价：
\begin{enumerate}
\item $R \to \Lambda$ 的纤维具有 $P$；
\item 对每个位于 $\mathfrak p \subset R$ 上方的素理想
$\mathfrak q \subset \Lambda$，$R_\mathfrak p \to \Lambda_\mathfrak q$
的纤维具有 $P$；以及
\item 对每个位于 $R$ 的 $\mathfrak m$ 上方的极大理想
$\mathfrak m' \subset \Lambda$，
$R_\mathfrak m \to \Lambda_{\mathfrak m'}$ 的纤维具有 $P$。
\end{enumerate}
\end{lemma}

\begin{proof}
设 $\mathfrak p \subset R$ 为素理想。那么，$\mathfrak p$ 上方的纤维是环
$\Lambda \otimes_R \kappa(\mathfrak p)$，其谱双射地映到
$\Spec(\Lambda)$ 中由位于 $\mathfrak p$ 上方的素理想
$\mathfrak q$ 构成的子集；见 Algebra，注
\ref{algebra-remark-fundamental-diagram}。对于这样的素理想
$\mathfrak q$，选取极大理想 $\mathfrak q \subset \mathfrak m'$，
并置 $\mathfrak m = R \cap \mathfrak m'$。那么
$\mathfrak p \subset \mathfrak m$，且有
$$
(\Lambda \otimes_R \kappa(\mathfrak p))_\mathfrak q \cong
(\Lambda_{\mathfrak m'} \otimes_{R_\mathfrak m}
\kappa(\mathfrak p))_\mathfrak q
$$
作为 $\kappa(\mathfrak q)$-代数。因此，(1)、(2) 和 (3) 等价，
因为由 (B)，可以在局部环上检验性质 $P$。
\end{proof}

\begin{lemma}
\label{more-algebra-lemma-P-ring-goes-up-quasi-finite}
设 $R \to R'$ 为 Noether 环之间的有限型同态，并设
$$
\xymatrix{
\mathfrak q' \ar[r] & \mathfrak p' \ar[r] & R' \\
\mathfrak q \ar[r] \ar@{-}[u] &
\mathfrak p \ar[r] \ar@{-}[u] & R \ar[u]
}
$$
为素理想图。假设 $R \to R'$ 在 $\mathfrak p'$ 处拟有限。
假设 $P$ 满足 (A) 和 (B)。
\begin{enumerate}
\item 若 $\kappa(\mathfrak q) \to
R_\mathfrak p^\wedge \otimes_R \kappa(\mathfrak q)$
具有 $P$，则
$\kappa(\mathfrak q') \to
(R'_{\mathfrak p'})^\wedge \otimes_{R'} \kappa(\mathfrak q')$
具有 $P$。
\item 若 $R_\mathfrak p$ 的形式纤维具有 $P$，
则 $R'_{\mathfrak p'}$ 的形式纤维具有 $P$。
\item 若 $R \to R'$ 拟有限且 $R$ 是 $P$-环，则 $R'$ 是 $P$-环。
\end{enumerate}
\end{lemma}

\begin{proof}
显然，(1) $\Rightarrow$ (2) $\Rightarrow$ (3)。假设 $P$ 对
$\kappa(\mathfrak q) \to R_\mathfrak p^\wedge \otimes_R \kappa(\mathfrak q)$
成立。由 Algebra，引理
\ref{algebra-lemma-completion-at-quasi-finite-prime}，可知
$$
R_\mathfrak p^\wedge \otimes_R R'
=
(R'_{\mathfrak p'})^\wedge \times B
$$
对某个 $R_\mathfrak p^\wedge$-代数 $B$ 成立。因此，
$R'_{\mathfrak p'} \to (R'_{\mathfrak p'})^\wedge$ 是同态
$R_\mathfrak p \to R_\mathfrak p^\wedge$ 的某个基变换的一个因子。
于是 $(R'_{\mathfrak p'})^\wedge \otimes_{R'} \kappa(\mathfrak q')$
是
$$
R_\mathfrak p^\wedge \otimes_R R' \otimes_{R'} \kappa(\mathfrak q') =
R_\mathfrak p^\wedge \otimes_R \kappa(\mathfrak q)
\otimes_{\kappa(\mathfrak q)} \kappa(\mathfrak q').
$$
的一个因子。因此，结论由关于 $P$ 的假设得到。
\end{proof}

\begin{lemma}
\label{more-algebra-lemma-check-P-ring-maximal-ideals}
设 $R$ 为 Noether 环。假设 $P$ 满足 (C) 和 (D)。那么，
$R$ 是 $P$-环，当且仅当对 $R$ 的每个极大理想 $\mathfrak m$，
$R_\mathfrak m$ 的形式纤维都具有 $P$。
\end{lemma}

\begin{proof}
假设对 $R$ 的所有极大理想 $\mathfrak m$，$R_\mathfrak m$
的形式纤维都具有 $P$。设 $\mathfrak p$ 为 $R$ 的素理想，
并选取极大理想 $\mathfrak p \subset \mathfrak m$。由于
$R_\mathfrak m \to R_\mathfrak m^\wedge$ 忠实平坦，
可以选取 $R_\mathfrak m^\wedge$ 中位于
$\mathfrak pR_\mathfrak m$ 上方的素理想 $\mathfrak p'$。
考虑交换图
$$
\xymatrix{
R_\mathfrak m^\wedge \ar[r] &
(R_\mathfrak m^\wedge)_{\mathfrak p'} \ar[r] &
(R_\mathfrak m^\wedge)_{\mathfrak p'}^\wedge
\\
R_\mathfrak m \ar[u] \ar[r] & R_\mathfrak p \ar[u] \ar[r] &
R_\mathfrak p^\wedge \ar[u]
}
$$
依假设，环同态 $R_\mathfrak m \to R_\mathfrak m^\wedge$ 的纤维
具有 $P$。由命题 \ref{more-algebra-proposition-Noetherian-complete-G-ring}，
$(R_\mathfrak m^\wedge)_{\mathfrak p'} \to
(R_\mathfrak m^\wedge)_{\mathfrak p'}^\wedge$ 正则。
局部化 $R_\mathfrak m^\wedge \to
(R_\mathfrak m^\wedge)_{\mathfrak p'}$ 正则。因此，由引理
\ref{more-algebra-lemma-regular-composition}，
$R_\mathfrak m^\wedge \to
(R_\mathfrak m^\wedge)_{\mathfrak p'}^\wedge$ 正则。
于是由 (C)，$R_\mathfrak m \to
(R_\mathfrak m^\wedge)_{\mathfrak p'}^\wedge$ 的纤维具有 $P$。
由于 $R_\mathfrak m \to
(R_\mathfrak m^\wedge)_{\mathfrak p'}^\wedge$ 分解经过局部化
$R_\mathfrak p$，同态 $R_\mathfrak p \to
(R_\mathfrak m^\wedge)_{\mathfrak p'}^\wedge$ 的纤维也具有 $P$。
因此，可以应用 (D) 得到
$R_\mathfrak p \to R_\mathfrak p^\wedge$ 的纤维具有 $P$。
\end{proof}

\begin{proposition}
\label{more-algebra-proposition-finite-type-over-P-ring}
设 $R$ 为 $P$-环，其中 $P$ 满足 (A)、(B)、(C) 和 (D)。
若 $R \to S$ 本质有限型，则 $S$ 是 $P$-环。
\end{proposition}

\begin{proof}
由于 $P$-环性质是局部环的性质，显然 $P$-环的局部化仍为 $P$-环。
反之，若在每个素理想处的局部化都是 $P$-环，则该环是 $P$-环。
因此，只需对每个有限型 $R$-代数 $S$ 及 $S$ 的每个素理想
$\mathfrak q$ 证明 $S_\mathfrak q$ 是 $P$-环。将 $S$ 写成
$R[x_1, \ldots, x_n]$ 的商，由引理
\ref{more-algebra-lemma-P-ring-goes-up-quasi-finite} 可知，只需证明
$R[x_1, \ldots, x_n]$ 是 $P$-环。对 $n$ 归纳，只需证明
$R[x]$ 是 $P$-环。设 $\mathfrak q \subset R[x]$ 为极大理想。
由引理 \ref{more-algebra-lemma-check-P-ring-maximal-ideals}，只需证明
$$
R[x]_\mathfrak q \longrightarrow R[x]_\mathfrak q^\wedge
$$
的纤维具有 $P$。若 $\mathfrak q$ 位于 $\mathfrak p \subset R$ 上方，
则可以用 $R_\mathfrak p$ 代替 $R$。因此，可以假设 $R$ 是极大理想为
$\mathfrak m$ 的 Noether 局部 $P$-环，并且
$\mathfrak q \subset R[x]$ 位于 $\mathfrak m$ 上方。注意，
存在唯一的素理想 $\mathfrak q' \subset R^\wedge[x]$
位于 $\mathfrak q$ 上方。考虑交换图
$$
\xymatrix{
R[x]_\mathfrak q^\wedge \ar[r] &
(R^\wedge[x]_{\mathfrak q'})^\wedge \\
R[x]_\mathfrak q \ar[r] \ar[u] & R^\wedge[x]_{\mathfrak q'} \ar[u]
}
$$
由于 $R$ 是 $P$-环，$R[x] \to R^\wedge[x]$ 的纤维具有 $P$：
它们是 $R \to R^\wedge$ 的纤维沿有限生成域扩张的基变换，
故可应用 (A)。于是，例如由引理 \ref{more-algebra-lemma-P-local}，
下方水平箭头的纤维具有 $P$。右侧竖直箭头正则，因为
$R^\wedge$ 是 G-环（命题 \ref{more-algebra-proposition-Noetherian-complete-G-ring}
和 \ref{more-algebra-proposition-finite-type-over-G-ring}）。因此，由 (C)，复合
$R[x]_\mathfrak q \to (R^\wedge[x]_{\mathfrak q'})^\wedge$
的纤维具有 $P$。从而由 (D)，左侧竖直箭头的纤维具有 $P$，
证明完成。
\end{proof}

\begin{lemma}
\label{more-algebra-lemma-map-P-ring-to-completion-P}
设 $A$ 为 $P$-环，其中 $P$ 满足 (B) 和 (D)。设 $I \subset A$
为理想，并设 $A^\wedge$ 为 $A$ 关于 $I$ 的完备化。那么
$A \to A^\wedge$ 的纤维具有 $P$。
\end{lemma}

\begin{proof}
由 Algebra，引理 \ref{algebra-lemma-completion-flat}，环同态
$A \to A^\wedge$ 平坦。由 Algebra，引理
\ref{algebra-lemma-completion-Noetherian-Noetherian}，环 $A^\wedge$
是 Noether 环。因此，只需检验引理 \ref{more-algebra-lemma-P-local}
的第三个条件。设 $\mathfrak m' \subset A^\wedge$ 为位于
$\mathfrak m \subset A$ 上方的极大理想。由 Algebra，引理
\ref{algebra-lemma-radical-completion}，有
$IA^\wedge \subset \mathfrak m'$。由于 $A^\wedge/IA^\wedge = A/I$，
可知 $I \subset \mathfrak m$、
$\mathfrak m/I = \mathfrak m'/IA^\wedge$，并且
$A/\mathfrak m = A^\wedge/\mathfrak m'$。由于
$A^\wedge/\mathfrak m'$ 是域，可知 $\mathfrak m$ 也是极大理想。
于是，$A_\mathfrak m \to A^\wedge_{\mathfrak m'}$ 是 Noether
局部环之间的平坦局部同态，它诱导剩余域的同构，并满足
$\mathfrak m A^\wedge_{\mathfrak m'} =
\mathfrak m'A^\wedge_{\mathfrak m'}$。因此，由引理
\ref{more-algebra-lemma-flat-unramified}，它在完备局部环上诱导同构。
设 $(A_\mathfrak m)^\wedge$ 为 $A_\mathfrak m$ 关于其极大理想的完备化。
环同态
$$
(A^\wedge)_{\mathfrak m'} \to
((A^\wedge)_{\mathfrak m'})^\wedge = (A_\mathfrak m)^\wedge
$$
忠实平坦（Algebra，引理
\ref{algebra-lemma-completion-faithfully-flat}）。因此，可以对环同态
$$
A_\mathfrak m \to (A^\wedge)_{\mathfrak m'} \to (A_\mathfrak m)^\wedge
$$
应用 (D) 得到结论，因为 $A$ 是 $P$-环，所以
$A_\mathfrak m \to (A_\mathfrak m)^\wedge$ 的纤维具有 $P$。
\end{proof}

\begin{lemma}
\label{more-algebra-lemma-henselization-pair-P-ring}
\begin{slogan}
环的 Hensel 化继承形式纤维的良好性质
\end{slogan}
设 $A$ 为 $P$-环，其中 $P$ 满足 (B)、(C)、(D) 和 (E)。
设 $I \subset A$ 为理想。设 $(A^h, I^h)$ 为偶对 $(A, I)$
的 Hensel 化；见引理 \ref{more-algebra-lemma-henselization}。那么 $A^h$ 是 $P$-环。
\end{lemma}

\begin{proof}
设 $\mathfrak m^h \subset A^h$ 为极大理想。由引理
\ref{more-algebra-lemma-check-P-ring-maximal-ideals}，我们必须证明
$A^h_{\mathfrak m^h} \to (A^h_{\mathfrak m^h})^\wedge$
的纤维具有 $P$。设 $\mathfrak m$ 为 $\mathfrak m^h$ 在 $A$
中的逆像。注意，$I^h \subset \mathfrak m^h$，从而
$I \subset \mathfrak m$，因为 $(A^h, I^h)$ 是 Hensel 偶对。
回顾：$A^h$ 是 Noether 环，$I^h = IA^h$，并且
$A \to A^h$ 在 $I$-adic 完备化上诱导同构；见引理
\ref{more-algebra-lemma-henselization-Noetherian-pair}。于是，Noether 局部环之间的局部同态
$$
A_\mathfrak m \to A^h_{\mathfrak m^h}
$$
由引理 \ref{more-algebra-lemma-flat-unramified} 在极大理想完备化上诱导同构
（略去细节）。设 $\mathfrak q^h$ 为 $A^h_{\mathfrak m^h}$
中位于 $\mathfrak q \subset A_\mathfrak m$ 上方的素理想。
置 $\mathfrak q_1 = \mathfrak q^h$，并令
$\mathfrak q_2, \ldots, \mathfrak q_t$ 为 $A^h$ 中位于
$\mathfrak q$ 上方的其余素理想，于是
$A^h \otimes_A \kappa(\mathfrak q) =
\prod\nolimits_{i = 1, \ldots, t} \kappa(\mathfrak q_i)$；见引理
\ref{more-algebra-lemma-filtered-colimit-etale-noetherian-fibres}。利用上面讨论过的
$(A^h)_{\mathfrak m^h}^\wedge = (A_\mathfrak m)^\wedge$，可得
$$
\prod\nolimits_{i = 1, \ldots, t}
(A^h_{\mathfrak m^h})^\wedge \otimes_{A^h_{\mathfrak m^h}}
\kappa(\mathfrak q_i) =
(A^h_{\mathfrak m^h})^\wedge \otimes_{A^h_{\mathfrak m^h}}
(A^h_{\mathfrak m^h} \otimes_{A_{\mathfrak m}} \kappa(\mathfrak q)) =
(A_{\mathfrak m})^\wedge \otimes_{A_{\mathfrak m}} \kappa(\mathfrak q)
$$
因此，考察局部环并使用 (B)，可知
$$
\kappa(\mathfrak q) \longrightarrow
(A^h_{\mathfrak m^h})^\wedge \otimes_{A^h_{\mathfrak m^h}}
\kappa(\mathfrak q^h)
$$
具有 $P$，因为依关于 $A$ 的假设，
$\kappa(\mathfrak q) \to
(A_\mathfrak m)^\wedge \otimes_{A_\mathfrak m} \kappa(\mathfrak q)$
也具有该性质。由于 $\kappa(\mathfrak q^h)/\kappa(\mathfrak q)$
是可分代数扩张，由 (E) 可知，如所需地
$\kappa(\mathfrak q^h) \to
(A^h_{\mathfrak m^h})^\wedge \otimes_{A^h_{\mathfrak m^h}}
\kappa(\mathfrak q^h)$ 具有 $P$。
\end{proof}

\begin{lemma}
\label{more-algebra-lemma-henselization-P-ring}
设 $R$ 为 Noether 局部 $P$-环，其中 $P$ 满足 (B)、(C)、(D) 和 (E)。
那么，Hensel 化 $R^h$ 和严格 Hensel 化 $R^{sh}$ 都是 $P$-环。
\end{lemma}

\begin{proof}
对于 Hensel 化，我们已在引理 \ref{more-algebra-lemma-henselization-pair-P-ring}
中证明了结论。要对严格 Hensel 化证明结论，由引理
\ref{more-algebra-lemma-check-P-ring-maximal-ideals}，只需证明 $R^{sh}$
的形式纤维具有 $P$。设 $\mathfrak r \subset R^{sh}$ 为素理想，
并置 $\mathfrak p = R \cap \mathfrak r$。置
$\mathfrak r_1 = \mathfrak r$，并令
$\mathfrak r_2, \ldots, \mathfrak r_s$ 为 $R^{sh}$ 中位于
$\mathfrak p$ 上方的其余素理想，于是
$R^{sh} \otimes_R \kappa(\mathfrak p) =
\prod\nolimits_{i = 1, \ldots, s} \kappa(\mathfrak r_i)$；见引理
\ref{more-algebra-lemma-fibres-henselization}。于是可得
$$
\prod\nolimits_{i = 1, \ldots, t}
(R^{sh})^\wedge \otimes_{R^{sh}} \kappa(\mathfrak r_i) =
(R^{sh})^\wedge \otimes_{R^{sh}} (R^{sh} \otimes_R \kappa(\mathfrak p)) =
(R^{sh})^\wedge \otimes_R \kappa(\mathfrak p)
$$
注意，由引理 \ref{more-algebra-lemma-henselization-noetherian}，
$R^\wedge \to (R^{sh})^\wedge$ 对
$\mathfrak m_{(R^{sh})^\wedge}$-adic 拓扑形式光滑。因此，由命题
\ref{more-algebra-proposition-fs-regular}，$R^\wedge \to (R^{sh})^\wedge$ 正则。
由 (C) 以及关于 $R$ 的假设，可知性质 $P$ 对
$\kappa(\mathfrak p) \to (R^{sh})^\wedge \otimes_R \kappa(\mathfrak p)$
成立。使用性质 (B)、上面的分解并考察局部环，可知性质 $P$ 对
$\kappa(\mathfrak p) \to
(R^{sh})^\wedge \otimes_{R^{sh}} \kappa(\mathfrak r)$ 成立。
由于 $\kappa(\mathfrak r)/\kappa(\mathfrak p)$ 是可分代数扩张，
由 (E)，可知 $P$ 对
$\kappa(\mathfrak r) \to
(R^{sh})^\wedge \otimes_{R^{sh}} \kappa(\mathfrak r)$ 成立。
\end{proof}

\begin{lemma}
\label{more-algebra-lemma-formal-fibres-reduced}
性质 (A)、(B)、(C)、(D) 和 (E) 对
$P(k \to R) =$“$R$ 在 $k$ 上几何既约”成立。
\end{lemma}

\begin{proof}
(A) 由几何既约代数的定义得到（Algebra，定义
\ref{algebra-definition-geometrically-reduced}）。(B) 也成立：
一个环既约，当且仅当它的所有局部环既约。
(C) 由引理 \ref{more-algebra-lemma-reduced-goes-up} 得到。
(D) 由 Algebra，引理 \ref{algebra-lemma-descent-reduced} 得到。
(E) 由 Algebra，引理
\ref{algebra-lemma-geometrically-reduced-over-separable-algebraic} 得到。
\end{proof}

\begin{lemma}
\label{more-algebra-lemma-formal-fibres-normal}
性质 (A)、(B)、(C)、(D) 和 (E) 对
$P(k \to R) =$“$R$ 在 $k$ 上几何正规”成立。
\end{lemma}

\begin{proof}
(A) 由几何正规代数的定义得到（Algebra，定义
\ref{algebra-definition-geometrically-normal}）。(B) 也成立：
一个环正规，当且仅当它的所有局部环正规。
(C) 由引理 \ref{more-algebra-lemma-normal-goes-up} 得到。
(D) 由 Algebra，引理 \ref{algebra-lemma-descent-normal} 得到。
(E) 由 Algebra，引理
\ref{algebra-lemma-geometrically-normal-over-separable-algebraic} 得到。
\end{proof}

\begin{lemma}
\label{more-algebra-lemma-formal-fibres-Sk}
固定 $n \geq 1$。性质 (A)、(B)、(C)、(D) 和 (E) 对
$P(k \to R) =$“$R$ 满足 $(S_n)$”成立。
\end{lemma}

\begin{proof}
设 $k \to R$ 为环同态，其中 $k$ 为域，$R$ 为 Noether 环。
设 $k'/k$ 为有限生成域扩张。由 Algebra，引理
\ref{algebra-lemma-tensor-fields-CM}，环同态
$R \to R \otimes_k k'$ 的纤维是 Cohen--Macaulay 的。
因此，可以对环同态 $R \to R \otimes_k k'$ 应用 Algebra，引理
\ref{algebra-lemma-Sk-goes-up}，得到：若 $R$ 满足 $(S_n)$，
则 $R \otimes_k k'$ 也如此。这证明了 (A)。
(B) 也成立：Noether 环满足 $(S_n)$，当且仅当其所有局部环都满足
$(S_n)$。(C) 由 Algebra，引理
\ref{algebra-lemma-Sk-goes-up} 得到，因为正则同态的纤维正则，
因而尤其是 Cohen--Macaulay 的。
(D) 由 Algebra，引理 \ref{algebra-lemma-descent-Sk} 得到。
(E) 是立即的，因为该条件不涉及基域。
\end{proof}

\begin{lemma}
\label{more-algebra-lemma-formal-fibres-CM}
性质 (A)、(B)、(C)、(D) 和 (E) 对
$P(k \to R) =$“$R$ 是 Cohen--Macaulay 环”成立。
\end{lemma}

\begin{proof}
这立即由引理 \ref{more-algebra-lemma-formal-fibres-Sk} 以及以下事实得到：
Noether 环是 Cohen--Macaulay 环，当且仅当它对所有 $n$
都满足条件 $(S_n)$。
\end{proof}

\begin{lemma}
\label{more-algebra-lemma-formal-fibres-Rk}
固定 $n \geq 0$。性质 (A)、(B)、(C)、(D) 和 (E) 对
$P(k \to R) =$“对所有有限扩张 $k'/k$，
$R \otimes_k k'$ 都满足 $(R_n)$”成立。
\end{lemma}

\begin{proof}
设 $k \to R$ 为环同态，其中 $k$ 为域，$R$ 为 Noether 环。
假设 $P(k \to R)$ 成立。设 $K/k$ 为有限生成域扩张。
由 Algebra，引理 \ref{algebra-lemma-make-separable}，可以找到交换图
$$
\xymatrix{
K \ar[r] & K' \\
k \ar[u] \ar[r] & k' \ar[u]
}
$$
其中 $k'/k$、$K'/K$ 是有限纯不可分域扩张，而 $K'/k'$ 可分。
由 Algebra，引理 \ref{algebra-lemma-localization-smooth-separable}，
存在光滑 $k'$-代数 $B$，使得 $K'$ 是 $B$ 的分式域。
现在可以作如下论证：
步骤 1：$R \otimes_k k'$ 满足 $(R_n)$，因为我们假设
$P$ 对 $k \to R$ 成立。
步骤 2：$R \otimes_k k' \to R \otimes_k k' \otimes_{k'} B$
是光滑环同态（Algebra，引理 \ref{algebra-lemma-base-change-smooth}），
由 Algebra，引理 \ref{algebra-lemma-Rk-goes-up} 可知
$R \otimes_k k' \otimes_{k'} B$ 满足 $(R_n)$
（并使用 Algebra，引理
\ref{algebra-lemma-characterize-smooth-over-field} 检验其假设成立）。
步骤 3：$R \otimes_k k' \otimes_{k'} K' = R \otimes_k K'$
满足 $(R_n)$，因为它是一个满足 $(R_n)$ 的环的局部化。
步骤 4：最后，由 $(R_n)$ 沿忠实平坦环同态
$K \otimes_k R \to K' \otimes_k R$ 的下降
（Algebra，引理 \ref{algebra-lemma-descent-Rk}），
$R \otimes_k K$ 满足 $(R_n)$。这证明了 (A)。
(B) 也成立：Noether 环满足 $(R_n)$，当且仅当其所有局部环都满足
$(R_n)$。(C) 由 Algebra，引理 \ref{algebra-lemma-Rk-goes-up}
得到，因为正则同态的纤维正则（略去一个小细节）。
(D) 由 Algebra，引理 \ref{algebra-lemma-descent-Rk} 得到
（略去一个小细节）。

\medskip\noindent
对于 (E)，设 $l/k$ 为域的可分代数扩张，并设 $l \to R$ 为环同态，
其中 $R$ 为 Noether 环。假设 $k \to R$ 具有 $P$。
我们必须证明 $l \to R$ 具有 $P$。设 $l'/l$ 为有限扩张。
首先注意，存在有限子扩张 $l/m/k$ 和有限扩张 $m'/m$，使得
$l' = l \otimes_m m'$。那么 $R \otimes_l l' = R \otimes_m m'$。
因此，只需证明 $m \to R$ 具有性质 $P$；换言之，可以假设
$l/k$ 有限。若 $l/k$ 有限，则 $l'/k$ 有限，并且有
$$
l' \otimes_l R = (l' \otimes_k R) \otimes_{l \otimes_k l} l
$$
它是具有性质 $(R_n)$ 的 Noether 环 $l' \otimes_k R$ 的一个局部化
（由 Algebra，引理 \ref{algebra-lemma-separable-algebraic-diagonal}），
其中该环具有此性质是因为我们假设 $P$ 对 $k \to R$ 成立。
这证明了 $l' \otimes_l R$ 如所需地具有性质 $(R_n)$。
\end{proof}








\section{优环}
\label{more-algebra-section-excellent}

\noindent
本节讨论 Grothendieck 的优环概念。关于 G-环、J-2 环和普遍链环的定义，
我们分别参见定义 \ref{more-algebra-definition-G-ring}、定义 \ref{more-algebra-definition-J}
以及 Algebra，定义 \ref{algebra-definition-universally-catenary}。

\begin{definition}
\label{more-algebra-definition-excellent}
设 $R$ 为环。
\begin{enumerate}
\item 若 $R$ 是 Noether 环、G-环且为 J-2，则称 $R$ 是{\it 拟优环}。
\item 若 $R$ 是拟优环且为普遍链环，则称 $R$ 是{\it 优环}。
\end{enumerate}
\end{definition}

\noindent
因此，若 Noether 环具有几何正则的形式纤维，并且其上任意有限型代数的
奇异集都闭，则它是拟优环。要使这样的环成为优环，还要求（局部地）
存在良好的维数函数。稍后我们将看到（第
\ref{more-algebra-section-formally-catenary} 节），普遍链性可以表述为对 $R$
的极大理想 $\mathfrak m$，同态
$R_\mathfrak m \to R_\mathfrak m^\wedge$ 所满足的条件。

\begin{lemma}
\label{more-algebra-lemma-finite-type-over-excellent}
（拟）优环上的有限型环的任意局部化都是（拟）优环。
\end{lemma}

\begin{proof}
对于有限型代数，这由 J-2 和普遍链性的定义得到。对于 G-环，见命题
\ref{more-algebra-proposition-finite-type-over-G-ring}。我们略去局部化保持（拟）优性的证明。
\end{proof}

\begin{proposition}
\label{more-algebra-proposition-ubiquity-excellent}
下列各类环都是优环：
\begin{enumerate}
\item 域；
\item Noether 完备局部环；
\item $\mathbf{Z}$；
\item 分式域特征为零的 Dedekind 整环；
\item 上述任一类环的有限型环扩张。
\end{enumerate}
\end{proposition}

\begin{proof}
由命题 \ref{more-algebra-proposition-ubiquity-G-ring} 和
\ref{more-algebra-proposition-ubiquity-J-2} 可知，这些环是 G-环且具有 J-2。
任意 Cohen--Macaulay 环都是普遍链环；见 Algebra，引理
\ref{algebra-lemma-CM-ring-catenary}。特别地，域、Dedekind 环以及更一般的
正则环都是普遍链环。借助 Cohen 结构定理，可知完备局部环是普遍链环；
见 Algebra，注
\ref{algebra-remark-Noetherian-complete-local-ring-universally-catenary}。
\end{proof}

\noindent
上面发展的材料对 Nagata 环有一些推论。

\begin{lemma}
\label{more-algebra-lemma-Nagata-local-ring}
设 $(A, \mathfrak m)$ 为 Noether 局部环。下列条件等价：
\begin{enumerate}
\item $A$ 是 Nagata 环；以及
\item $A$ 的形式纤维几何既约。
\end{enumerate}
\end{lemma}

\begin{proof}
假设 (2)。由 Algebra，引理
\ref{algebra-lemma-local-nagata-and-analytically-unramified}，
必须证明：若 $A \to B$ 有限、$B$ 为整环且
$\mathfrak m' \subset B$ 为极大理想，则
$B_{\mathfrak m'}$ 解析无分歧。结合引理
\ref{more-algebra-lemma-formal-fibres-reduced}、
\ref{more-algebra-lemma-check-P-ring-maximal-ideals} 以及命题
\ref{more-algebra-proposition-finite-type-over-P-ring}，可知
$B_{\mathfrak m'}$ 的形式纤维几何既约。特别地，若 $L$ 为 $B$
的分式域，则 $B_{\mathfrak m'}^\wedge \otimes_B L$ 既约。
由此可知 $B_{\mathfrak m'}^\wedge$ 既约；换言之，
$B_{\mathfrak m'}$ 解析无分歧。

\medskip\noindent
假设 (1)。设 $\mathfrak q \subset A$ 为素理想，并设
$K/\kappa(\mathfrak q)$ 为有限扩张。我们必须证明
$A^\wedge \otimes_A K$ 既约。取局部子环
$A/\mathfrak q \subset B \subset K$，使得它在 $A$ 上有限且分式域为 $K$。
为了构造 $B$，选取生成 $K$ 在 $\kappa(\mathfrak q)$ 上扩张的元素
$x_1, \ldots, x_n \in K$，它们满足首一多项式
$P_i(T) = T^{d_i} + a_{i, 1} T^{d_i - 1} + \ldots + a_{i, d_i} = 0$，
其中 $a_{i, j} \in \mathfrak m$。然后令 $B$ 为 $K$ 中由
$x_1, \ldots, x_n$ 生成的 $A$-子代数。（更多细节见 Algebra，引理
\ref{algebra-lemma-local-nagata-and-analytically-unramified} 的证明。）那么
$$
A^\wedge \otimes_A K =
(A^\wedge \otimes_A B)_\mathfrak q =
B^\wedge_\mathfrak q
$$
由 Algebra，引理
\ref{algebra-lemma-local-nagata-and-analytically-unramified}，
$B^\wedge$ 既约，故证明完成。
\end{proof}

\begin{lemma}
\label{more-algebra-lemma-quasi-excellent-nagata}
拟优环是 Nagata 环。
\end{lemma}

\begin{proof}
设 $R$ 为拟优环。利用有限型 $R$-代数仍为拟优环这一事实
（引理 \ref{more-algebra-lemma-finite-type-over-excellent}），由 Algebra，引理
\ref{algebra-lemma-check-universally-japanese}，只需证明任意拟优整环都是 N-1。
应用 Algebra，引理 \ref{algebra-lemma-characterize-N-1}
（并使用拟优环为 J-2 这一事实），问题化约为证明拟优局部整环 $R$ 是 N-1。
由于 $R \to R^\wedge$ 正则，由引理 \ref{more-algebra-lemma-reduced-goes-up}，
$R^\wedge$ 既约。换言之，$R$ 解析无分歧。因此，由 Algebra，引理
\ref{algebra-lemma-analytically-unramified-easy}，$R$ 是 N-1。
\end{proof}

\begin{lemma}
\label{more-algebra-lemma-completion-normal-local-ring}
设 $(A, \mathfrak m)$ 为 Noether 局部环。若 $A$ 正规，且 $A$
的形式纤维正规（例如，若 $A$ 为优环或拟优环），则 $A^\wedge$ 正规。
\end{lemma}

\begin{proof}
这立即由 Algebra，引理
\ref{algebra-lemma-normal-goes-up-noetherian} 得到。
\end{proof}





\section{模的 Abel 范畴}
\label{more-algebra-section-abelian-categories-modules}

\noindent
设 $R$ 为环。$R$-模范畴 $\text{Mod}_R$ 是 Abel 范畴。
以下是 $\text{Mod}_R$ 的一些 Abel 子范畴的例子（我们使用 Homology，
定义 \ref{homology-definition-serre-subcategory} 以及 Homology，引理
\ref{homology-lemma-characterize-serre-subcategory} 和
\ref{homology-lemma-characterize-weak-serre-subcategory} 中引入的术语）：
\begin{enumerate}
\item 相干 $R$-模范畴是 $\text{Mod}_R$ 的弱 Serre 子范畴。
这由 Algebra，引理 \ref{algebra-lemma-coherent} 得到。
\item 设 $S \subset R$ 为乘法子集。由所有满足下述条件的 $R$-模 $M$
构成的满子范畴是 $\text{Mod}_R$ 的 Serre 子范畴：对每个 $s \in S$，
相应乘法在 $M$ 上是同构。这由 Algebra，引理
\ref{algebra-lemma-localization-and-modules} 得到。
\item 设 $I \subset R$ 为有限生成理想。$I$-幂挠模的满子范畴是
$\text{Mod}_R$ 的 Serre 子范畴。见引理 \ref{more-algebra-lemma-I-power-torsion}。
\item 在一些文献中，{\it 挠模}定义为满足下述条件的模 $M$：
对每个 $x \in M$，存在非零因子 $f \in R$，使得 $fx = 0$。
挠模的满子范畴是 $\text{Mod}_R$ 的 Serre 子范畴。
\item 若 $R$ 不是 Noether 环，则有限生成 $R$-模范畴
$\text{Mod}^{fg}_R$ {\bf 不是} Abel 范畴。事实上，若
$I \subset R$ 是非有限生成理想，则映射 $R \to R/I$
在 $\text{Mod}^{fg}_R$ 中没有核。
\item 若 $R$ 是 Noether 环，则相干 $R$-模等同于有限生成
（即有限）$R$-模；见 Algebra，引理
\ref{algebra-lemma-Noetherian-coherent}、
\ref{algebra-lemma-coherent-ring} 和
\ref{algebra-lemma-Noetherian-finite-type-is-finite-presentation}。
因此，由上面的 (1)，$\text{Mod}^{fg}_R$ 是 Abel 范畴；事实上，
在这种情况下，范畴 $\text{Mod}_R^{fg}$ 是 $\text{Mod}_R$
的（强）Serre 子范畴。
\end{enumerate}




\section{内射 Abel 群}
\label{more-algebra-section-abelian-groups}

\noindent
本节证明 Abel 群范畴有足够多的内射对象。回顾：Abel 群 $M$
是{\it 可除的}，当且仅当对每个 $x \in M$ 和每个
$n \in \mathbf{N}$，存在 $y \in M$，使得 $n y = x$。

\begin{lemma}
\label{more-algebra-lemma-injective-abelian}
Abel 群 $J$ 是 Abel 群范畴中的内射对象，当且仅当 $J$ 可除。
\end{lemma}

\begin{proof}
假设 $J$ 不可除。那么，存在 $x \in J$ 和 $n \in \mathbf{N}$，
使得不存在满足 $n y = x$ 的 $y \in J$。于是，同态
$\mathbf{Z} \to J$，$m \mapsto mx$ 不能延拓到
$\frac{1}{n}\mathbf{Z} \supset \mathbf{Z}$。因此 $J$ 不是内射对象。

\medskip\noindent
设 $A \subset B$ 为 Abel 群。假设 $J$ 为可除 Abel 群。
设 $\varphi : A \to J$ 为同态。考虑所有同态
$\varphi' : A' \to J$ 的集合，其中 $A \subset A' \subset B$
且 $\varphi'|_A = \varphi$。定义
$(A', \varphi') \geq (A'', \varphi'')$，当且仅当
$A' \supset A''$ 且 $\varphi'|_{A''} = \varphi''$。
若 $(A_i, \varphi_i)_{i \in I}$ 是这类偶对的全序族，
则得到映射 $\bigcup_{i \in I} A_i \to J$：元素 $a \in A_i$
映到 $\varphi_i(a)$。因此，可应用 Zorn 引理。为得出结论，
必须证明：若偶对 $(A', \varphi')$ 极大，则 $A' = B$。
换言之，只需证明：给定任意子群 $A \subset B$、$A \not = B$
以及任意 $\varphi : A \to J$，可以找到
$\varphi' : A' \to J$，其中 $A \subset A' \subset B$，使得
(a) 包含 $A \subset A'$ 是严格的，且 (b) 同态 $\varphi'$
延拓 $\varphi$。

\medskip\noindent
为证明这一点，选取 $x \in B$、$x \not \in A$。
若不存在 $n\in \mathbf{N}$ 使得 $nx \in A$，则
$A \oplus \mathbf{Z} \cong A + \mathbf{Z}x$。于是，可以将
$\varphi$ 延拓到 $A' = A + \mathbf{Z}x$，例如在 $A$ 上使用
$\varphi$，并将 $x$ 映到零。若确实存在 $n \in \mathbf{N}$
使得 $nx \in A$，则取最小的此类整数 $n$。取 $z \in J$，
使得 $nz = \varphi(nx)$。定义同态
$\tilde\varphi : A \oplus \mathbf{Z} \to J$ 为
$(a, m) \mapsto \varphi(a) + mz$。由 $z$ 的选取，
$\tilde \varphi$ 的核包含映射 $A \oplus \mathbf{Z} \to B$，
$(a, m) \mapsto a + mx$ 的核。因此，$\tilde \varphi$
分解经过像 $A' = A + \mathbf{Z}x$，并由此延拓同态 $\varphi$。
\end{proof}

\noindent
可以用这个引理证明每个 Abel 群都能嵌入某个内射 Abel 群。
但这是下一节结果的一个特殊情形。






\section{内射模}
\label{more-algebra-section-injectives-modules}

\noindent
这里给出关于内射模的一些引理。

\begin{definition}
\label{more-algebra-definition-projective}
设 $R$ 为环。$R$-模 $J$ 是{\it 内射的}，当且仅当函子
$\Hom_R(-, J) : \text{Mod}_R \to \text{Mod}_R$ 是正合函子。
\end{definition}

\noindent
对任意 $R$-模 $M$，函子 $\Hom_R(- , M)$ 左正合；见 Algebra，引理
\ref{algebra-lemma-hom-exact}。因此，$J$ 为内射模这一条件实际表示：
给定 $R$-模的单射 $M \to M'$，映射
$\Hom_R(M', J) \to \Hom_R(M, J)$ 满射。

\medskip\noindent
在用 ${Ext}$-模重新表述这一点之前，我们先讨论
$\Ext^1_R(M, N)$ 与 Homology，第 \ref{homology-section-extensions} 节
中的扩张之间的关系。

\begin{lemma}
\label{more-algebra-lemma-relation-ext-ext}
设 $R$ 为环。设 $\mathcal{A}$ 为 $R$-模的 Abel 范畴。
存在典范同构
$\Ext_\mathcal{A}(M, N) = \Ext^1_R(M, N)$，它与 Algebra，引理
\ref{algebra-lemma-long-exact-seq-ext} 和
\ref{algebra-lemma-reverse-long-exact-seq-ext} 的长正合序列以及
Homology，引理 \ref{homology-lemma-six-term-sequence-ext}
的 $6$ 项正合序列相容。
\end{lemma}

\begin{proof}
略。
\end{proof}

\begin{lemma}
\label{more-algebra-lemma-characterize-injective}
设 $R$ 为环。设 $J$ 为 $R$-模。下列条件等价：
\begin{enumerate}
\item $J$ 是内射模；
\item 对每个 $R$-模 $M$，都有 $\Ext^1_R(M, J) = 0$。
\end{enumerate}
\end{lemma}

\begin{proof}
设 $0 \to M'' \to M' \to M \to 0$ 为 $R$-模的短正合序列。
考虑 Algebra，引理 \ref{algebra-lemma-reverse-long-exact-seq-ext}
中的长正合序列
$$
\begin{matrix}
0
\to \Hom_R(M, J)
\to \Hom_R(M', J)
\to \Hom_R(M'', J)
\\
\phantom{0\ }
\to \Ext^1_R(M, J)
\to \Ext^1_R(M', J)
\to \Ext^1_R(M'', J)
\to \ldots
\end{matrix}
$$
由此可知 (2) 推出 (1)。反之，若 $J$ 是内射模，
则由 Homology，引理 \ref{homology-lemma-characterize-injectives}
和引理 \ref{more-algebra-lemma-relation-ext-ext}，该 $\Ext$-群为零。
\end{proof}

\begin{lemma}
\label{more-algebra-lemma-characterize-injective-bis}
设 $R$ 为环。设 $J$ 为 $R$-模。下列条件等价：
\begin{enumerate}
\item $J$ 是内射模；
\item 对每个理想 $I \subset R$，都有 $\Ext^1_R(R/I, J) = 0$；以及
\item 对理想 $I \subset R$ 和模同态 $I \to J$，存在延拓 $R \to J$。
\end{enumerate}
\end{lemma}

\begin{proof}
若 $I \subset R$ 为理想，则短正合序列
$0 \to I \to R \to R/I \to 0$ 给出正合序列
$$
\Hom_R(R, J) \to
\Hom_R(I, J) \to
\Ext^1_R(R/I, J) \to 0
$$
这是由 Algebra，引理 \ref{algebra-lemma-reverse-long-exact-seq-ext}
以及如下事实得到的：$R$ 是投射模，故
$\Ext^1_R(R, J) = 0$（Algebra，引理
\ref{algebra-lemma-characterize-projective}）。因此，(2) 与 (3) 等价。
在本证明中，我们将证明 (1) $\Leftrightarrow$ (3)，这就是著名的
Baer 判据。

\medskip\noindent
假设 (1)。给定 (3) 中的模同态 $I \to J$，由于按定义映射
$\Hom_R(R, J) \to \Hom_R(I, J)$ 满射，便得到延拓 $R \to J$。

\medskip\noindent
假设 (3)。设 $M \subset N$ 为 $R$-模的包含。设
$\varphi : M \to J$ 为同态。我们将证明 $\varphi$ 延拓到 $N$，
这就完成引理的证明。考虑所有同态 $\varphi' : M' \to J$ 的集合，
其中 $M \subset M' \subset N$ 且 $\varphi'|_M = \varphi$。
定义 $(M', \varphi') \geq (M'', \varphi'')$，当且仅当
$M' \supset M''$ 且 $\varphi'|_{M''} = \varphi''$。
若 $(M_i, \varphi_i)_{i \in I}$ 是这类偶对的全序族，
则得到映射 $\bigcup_{i \in I} M_i \to J$：元素 $a \in M_i$
映到 $\varphi_i(a)$。因此，可应用 Zorn 引理。为得出结论，
必须证明：若偶对 $(M', \varphi')$ 极大，则 $M' = N$。
换言之，只需证明：给定任意子群 $M \subset N$、$M \not = N$
以及任意 $\varphi : M \to J$，可以找到
$\varphi' : M' \to J$，其中 $M \subset M' \subset N$，使得
(a) 包含 $M \subset M'$ 是严格的，且 (b) 同态 $\varphi'$
延拓 $\varphi$。

\medskip\noindent
为证明这一点，选取 $x \in N$、$x \not \in M$。令
$I = \{f \in R \mid fx \in M\}$。这是 $R$ 的理想。
定义同态 $\psi : I \to J$ 为 $f \mapsto \varphi(fx)$。
依假设 (3)，可以将其延拓为映射 $\tilde\psi : R \to J$。
由 $I$ 的选取，同态
$M \oplus R \to J$，$(y, f) \mapsto \varphi(y) + \tilde\psi(f)$
的核包含映射 $M \oplus R \to N$，$(y, f) \mapsto y + fx$ 的核。
因此，该同态分解经过像 $M' = M + Rx$，并如所需地延拓给定同态。
\end{proof}

\noindent
本节余下部分证明环 $R$ 上有足够多的内射模。首先注意，
$\mathbf{Q}/\mathbf{Z}$ 是内射 Abel 群。这由引理
\ref{more-algebra-lemma-injective-abelian} 得到。

\begin{definition}
\label{more-algebra-definition-simple-functors}
设 $R$ 为环。
\begin{enumerate}
\item 对任意 $R$ 上的 $R$-模 $M$，记
$M^\vee = \Hom(M, \mathbf{Q}/\mathbf{Z})$，并赋予其自然的
$R$-模结构。我们把{\it $M \mapsto M^\vee$} 看成从 $R$-模范畴
到自身的反变函子。
\item 对任意 $R$-模 $M$，记
$$
F(M) = \bigoplus\nolimits_{m \in M} R[m]
$$
为以 $m \in M$ 对应的元素 $[m]$ 为基的{\it 自由模}。
令 $F(M)\to M$，$\sum f_i [m_i] \mapsto \sum f_i m_i$
为 $R$-模的自然满射。我们把 $M \mapsto (F(M) \to M)$
看成从 $R$-模范畴到 $R$-模箭头范畴的函子。
\end{enumerate}
\end{definition}

\begin{lemma}
\label{more-algebra-lemma-vee-exact}
设 $R$ 为环。函子 $M \mapsto M^\vee$ 正合。
\end{lemma}

\begin{proof}
这是因为由引理 \ref{more-algebra-lemma-injective-abelian}，
$\mathbf{Q}/\mathbf{Z}$ 是内射 Abel 群。
\end{proof}

\noindent
存在由求值给出的典范映射 $ev : M \to (M^\vee)^\vee$：
给定 $x \in M$，令
$ev(x) \in (M^\vee)^\vee = \Hom(M^\vee, \mathbf{Q}/\mathbf{Z})$
为映射 $\varphi \mapsto \varphi(x)$。

\begin{lemma}
\label{more-algebra-lemma-ev-injective}
对任意 $R$-模 $M$，求值映射 $ev : M \to (M^\vee)^\vee$ 是单射。
\end{lemma}

\begin{proof}
可以利用 $\mathbf{Q}/\mathbf{Z}$ 是内射 Abel 群这一事实来检验。
具体地，若 $x \in M$ 非零，则令 $M' \subset M$ 为它生成的循环群。
存在非零映射 $M' \to \mathbf{Q}/\mathbf{Z}$，它必不消去 $x$。
该映射延拓为 $\varphi : M \to \mathbf{Q}/\mathbf{Z}$，于是
$ev(x)(\varphi) = \varphi(x) \not = 0$。
\end{proof}

\noindent
$R$-模的典范满射 $F(M) \to M$ 经过上述构造变为 $R$-模的典范单射
$$
(M^\vee)^\vee \longrightarrow (F(M^\vee))^\vee.
$$
置 $J(M) = (F(M^\vee))^\vee$。$ev$ 与这个显示映射的复合给出
$M \to J(M)$，并且该构造关于 $M$ 具有函子性。

\begin{lemma}
\label{more-algebra-lemma-JM-injective}
设 $R$ 为环。对每个 $R$-模 $M$，$R$-模 $J(M)$ 是内射模。
\end{lemma}

\begin{proof}
注意，作为 $R$-模，$J(M) \cong \prod_{\varphi \in M^\vee} R^\vee$。
由于内射模的积仍为内射模，只需证明 $R^\vee$ 是内射模。
为此，使用等式
$$
\Hom_R(N, R^\vee) =
\Hom_R(N, \Hom_{\mathbf{Z}}(R, \mathbf{Q}/\mathbf{Z})) =
N^\vee
$$
以及由引理 \ref{more-algebra-lemma-vee-exact} 得到的事实：
$(-)^\vee$ 是正合函子。
\end{proof}

\begin{lemma}
\label{more-algebra-lemma-injectives-modules}
设 $R$ 为环。上述构造定义了从 $R$-模范畴到 $R$-模箭头范畴的协变函子
$M \mapsto (M \to J(M))$，使得对每个模 $M$，输出
$M \to J(M)$ 都是将 $M$ 嵌入内射 $R$-模 $J(M)$ 的单射。
\end{lemma}

\begin{proof}
由上文立即得到。
\end{proof}

\noindent
特别地，对任意 $R$-模映射 $M \to N$，都有相伴的同态
$J(M) \to J(N)$，使下图交换：
$$
\xymatrix{
M \ar[d] \ar[r] & N \ar[d] \\
J(M) \ar[r] & J(N) }
$$
这正是我们一般希望得到的构造。在 Homology，第
\ref{homology-section-injectives} 节中，我们引入了表述这一点的术语：
这意味着 $R$-模范畴具有函子性内射嵌入。










\section{模的导出范畴}
\label{more-algebra-section-derived-modules}

\noindent
本节给出环上模的导出范畴的一些一般性结论。

\medskip\noindent
设 $A$ 为环。以 $\text{Mod}_A$ 表示 $A$-模范畴。
我们用符号 $K(A)$ 表示 $A$-模复形的同伦范畴；换言之，
作为范畴，置 $K(A) = K(\text{Mod}_A)$；见 Derived Categories，第
\ref{derived-section-homotopy} 节。有界版本为
$K^+(A)$、$K^-(A)$ 和 $K^b(A)$。按照 Derived Categories，第
\ref{derived-section-homotopy-triangulated} 节，我们把 $K(A)$
看成三角范畴。$A$ 的{\it 导出范畴}记作 $D(A)$，它是通过在
$K(A)$ 中逆转拟同构而得到的范畴；换言之，置
$D(A) = D(\text{Mod}_A)$；见 Derived Categories，第
\ref{derived-section-derived-categories} 节\footnote{另见 Injectives，注
\ref{injectives-remark-existence-D}。}。有界版本为
$D^+(A)$、$D^-(A)$ 和 $D^b(A)$。

\medskip\noindent
设 $A$ 为环。$A$-模范畴有积，并且积正合。由引理
\ref{more-algebra-lemma-injectives-modules}，$A$-模范畴有足够多的内射对象。
因此，每个 $A$-模复形都拟同构于某个 K-内射复形
（Derived Categories，引理
\ref{derived-lemma-enough-K-injectives-Ab4-star}）。由此，$D(A)$
有可数积（Derived Categories，引理 \ref{derived-lemma-products}），
事实上还有任意积（Injectives，引理
\ref{injectives-lemma-derived-products}）。这蕴含 $D(A)$
中对象的每个逆系统都有导出极限（在同构意义下良定义）；见
Derived Categories，第 \ref{derived-section-derived-limit} 节。

\begin{lemma}
\label{more-algebra-lemma-K-injective-flat}
设 $R \to S$ 为平坦环同态。若 $I^\bullet$ 是 $S$-模的 K-内射复形，
则作为 $R$-模复形，$I^\bullet$ 也是 K-内射的。
\end{lemma}

\begin{proof}
这是因为由 Algebra，引理 \ref{algebra-lemma-adjoint-tensor-restrict}
以及张量积以 $S$ 是正合的这一事实，有
$\Hom_{K(R)}(M^\bullet, I^\bullet) =
\Hom_{K(S)}(M^\bullet \otimes_R S, I^\bullet)$。
\end{proof}

\begin{lemma}
\label{more-algebra-lemma-K-injective-epimorphism}
设 $R \to S$ 为环的满态射。设 $I^\bullet$ 为 $S$-模复形。
若作为 $R$-模复形，$I^\bullet$ 是 K-内射的，则作为 $S$-模复形，
$I^\bullet$ 也是 K-内射的。
\end{lemma}

\begin{proof}
这是因为对任意 $S$-模复形 $N^\bullet$，有
$\Hom_{K(R)}(N^\bullet, I^\bullet) =
\Hom_{K(S)}(N^\bullet, I^\bullet)$；见 Algebra，引理
\ref{algebra-lemma-epimorphism-modules}。
\end{proof}

\begin{lemma}
\label{more-algebra-lemma-hom-K-injective}
设 $A \to B$ 为环同态。若 $I^\bullet$ 是 $A$-模的 K-内射复形，
则 $\Hom_A(B, I^\bullet)$ 是 $B$-模的 K-内射复形。
\end{lemma}

\begin{proof}
这是因为由 Algebra，引理 \ref{algebra-lemma-adjoint-hom-restrict}，有
$\Hom_{K(B)}(N^\bullet, \Hom_A(B, I^\bullet)) =
\Hom_{K(A)}(N^\bullet, I^\bullet)$。
\end{proof}










\section{Tor 的计算}
\label{more-algebra-section-computing-tor}

\noindent
设 $R$ 为环。以 $D(R)$ 表示 $R$-模 Abel 范畴
$\text{Mod}_R$ 的导出范畴。注意，$\text{Mod}_R$ 有足够多的投射对象，
因为每个自由 $R$-模都是投射模。因此，可以定义从
$\text{Mod}_R$ 到任意 Abel 范畴的任何加性函子的左导出函子。

\medskip\noindent
这特别适用于函子
$ - \otimes_R M : \text{Mod}_R \to \text{Mod}_R$，
其左导出函子是 Tor 函子 $\text{Tor}_i^R(-, M)$；见 Algebra，第
\ref{algebra-section-tor} 节。还有全左导出函子
\begin{equation}
\label{more-algebra-equation-derived-tensor-module}
-\otimes_R^{\mathbf{L}} M :
D^{-}(R)
\longrightarrow
D^{-}(R)
\end{equation}
记作 $-\otimes_R^{\mathbf{L}} M$。它的卫星函子是 Tor 模；换言之，
$$
H^{-p}(N \otimes_R^{\mathbf{L}} M) = \text{Tor}_p^R(N, M).
$$

\medskip\noindent
当考虑与 $R$-代数 $A$ 的张量积时，会出现一种特殊情形。
在此情形下，我们把 $- \otimes_R A$ 看成从 $\text{Mod}_R$
到 $\text{Mod}_A$ 的函子。因此，有全右导出函子
\begin{equation}
\label{more-algebra-equation-derived-tensor-algebra}
-\otimes_R^{\mathbf{L}} A :
D^{-}(R)
\longrightarrow
D^{-}(A)
\end{equation}
记作 $-\otimes_R^{\mathbf{L}} A$。它的卫星函子是 Tor 群；换言之，
$$
H^{-p}(N \otimes_R^{\mathbf{L}} A) = \text{Tor}_p^R(N, A).
$$
特别地，这些 Tor 群自然具有 $A$-模结构。

\medskip\noindent
在接下来的几节中，我们会把本节材料推广到无界复形。






\section{复形的张量积}
\label{more-algebra-section-symmetric-monoidal}

\noindent
设 $R$ 为环。$R$-模复形范畴 $\text{Comp}(R)$ 具有对称幺半结构。
具体地，设有两个 $R$-模复形 $L^\bullet$ 和 $M^\bullet$。
利用 Homology，例 \ref{homology-example-double-complex-as-tensor-product-of}
和 Homology，定义 \ref{homology-definition-associated-simple-complex}，
得到第三个 $R$-模复形
$$
\text{Tot}(L^\bullet \otimes_R M^\bullet)
$$
显然，该构造关于 $L^\bullet$ 和 $M^\bullet$ 都具有函子性。
结合约束取 Homology，注 \ref{homology-remark-triple-complex}
中从三重复形 $K^\bullet \otimes_R L^\bullet \otimes_R M^\bullet$
构造的复形典范同构
$$
\text{Tot}(\text{Tot}(K^\bullet \otimes_R L^\bullet) \otimes_R M^\bullet)
\longrightarrow
\text{Tot}(K^\bullet \otimes_R \text{Tot}(L^\bullet \otimes_R M^\bullet))
$$
交换约束是典范同构
$$
\text{Tot}(L^\bullet \otimes_R M^\bullet) \to
\text{Tot}(M^\bullet \otimes_R L^\bullet)
$$
它在直和项 $L^p \otimes_R M^q$ 上使用符号 $(-1)^{pq}$。
为了看出这是复形映射，对 $x \in L^p$ 和 $y \in M^q$ 计算：
$$
\text{d}(x \otimes y) =
\text{d}_L(x) \otimes y + (-1)^px \otimes \text{d}_M(y)
$$
按我们的规则，右侧映到
$$
(-1)^{(p + 1)q}y \otimes \text{d}_L(x) +
(-1)^{p + p(q + 1)} \text{d}_M(y) \otimes x
$$
另一方面，
$$
\text{d}((-1)^{pq}y \otimes x) =
(-1)^{pq} \text{d}_M(y) \otimes x +
(-1)^{pq + q} y \otimes \text{d}_L(x)
$$
直接观察即可看出这两个表达式如所需地相等。

\begin{lemma}
\label{more-algebra-lemma-symmetric-monoidal-cat-complexes}
设 $R$ 为环。在赋予函子
$(L^\bullet, M^\bullet) \mapsto \text{Tot}(L^\bullet \otimes_R M^\bullet)$
以及上述结合约束和交换约束之后，$R$-模复形范畴
$\text{Comp}(R)$ 是对称幺半范畴。
\end{lemma}

\begin{proof}
略。提示：取在次数 $0$ 处为 $R$、在其他次数处为零的复形作为单位
$\mathbf{1}$，并使用显然的同构
$\text{Tot}(\mathbf{1} \otimes_R M^\bullet) = M^\bullet$ 和
$\text{Tot}(K^\bullet \otimes_R \mathbf{1}) = K^\bullet$。
为证明该引理，需要检验若干图的交换性；见 Categories，定义
\ref{categories-definition-monoidal-category} 和
\ref{categories-definition-symmetric-monoidal-category}。
每种情形中的检验都很直接。
\end{proof}

\begin{lemma}
\label{more-algebra-lemma-derived-tor-homotopy}
设 $R$ 为环。设 $P^\bullet$ 为 $R$-模复形。设
$\alpha, \beta : L^\bullet \to M^\bullet$ 为复形的同伦映射。
那么，$\alpha$ 和 $\beta$ 诱导同伦映射
$$
\text{Tot}(\alpha \otimes \text{id}_P),
\text{Tot}(\beta \otimes \text{id}_P) :
\text{Tot}(L^\bullet \otimes_R P^\bullet)
\longrightarrow
\text{Tot}(M^\bullet \otimes_R P^\bullet).
$$
特别地，构造
$L^\bullet \mapsto \text{Tot}(L^\bullet \otimes_R P^\bullet)$
定义了复形同伦范畴的自函子。
\end{lemma}

\begin{proof}
设对某个由 $h^n : L^n \to M^{n - 1}$ 定义的同伦 $h$，有
$\alpha = \beta + dh + hd$。置
$$
H^n = \bigoplus\nolimits_{a + b = n} h^a \otimes \text{id}_{P^b} :
\bigoplus\nolimits_{a + b = n} L^a \otimes_R P^b
\longrightarrow
\bigoplus\nolimits_{a + b = n} M^{a - 1} \otimes_R P^b
$$
直接计算可知，作为映射
$\text{Tot}(L^\bullet \otimes_R P^\bullet) \to
\text{Tot}(M^\bullet \otimes_R P^\bullet)$，有
$$
\text{Tot}(\alpha \otimes \text{id}_P) =
\text{Tot}(\beta \otimes \text{id}_P) + dH + Hd
$$
\end{proof}

\begin{lemma}
\label{more-algebra-lemma-symmetric-monoidal-homotopy-category}
设 $R$ 为环。在赋予函子
$(L^\bullet, M^\bullet) \mapsto \text{Tot}(L^\bullet \otimes_R M^\bullet)$
以及上述结合约束和交换约束之后，$R$-模复形的同伦范畴 $K(R)$
是对称幺半范畴。
\end{lemma}

\begin{proof}
这由引理 \ref{more-algebra-lemma-symmetric-monoidal-cat-complexes} 和
\ref{more-algebra-lemma-derived-tor-homotopy} 得到。略去细节。
\end{proof}

\begin{lemma}
\label{more-algebra-lemma-derived-tor-exact}
设 $R$ 为环。设 $P^\bullet$ 为 $R$-模复形。函子
$$
K(R) \longrightarrow K(R), \quad
L^\bullet \longmapsto \text{Tot}(P^\bullet \otimes_R L^\bullet)
$$
以及
$$
K(R) \longrightarrow K(R), \quad
L^\bullet \longmapsto \text{Tot}(L^\bullet \otimes_R P^\bullet)
$$
都是三角范畴的正合函子。
\end{lemma}

\begin{proof}
这由 Derived Categories，注
\ref{derived-remark-double-complex-as-tensor-product-of} 得到。
\end{proof}




\section{导出张量积}
\label{more-algebra-section-derived-tensor-product}

\noindent
我们可以在更一般的情形下构造导出张量积。事实上，只要选取
K-平坦解消，就不需要有界性假设。

\begin{definition}
\label{more-algebra-definition-K-flat}
设 $R$ 为环。若对每个无环复形 $M^\bullet$，全复形
$\text{Tot}(M^\bullet \otimes_R K^\bullet)$ 都无环，
则称复形 $K^\bullet$ 是{\it K-平坦的}。
\end{definition}

\begin{lemma}
\label{more-algebra-lemma-K-flat-quasi-isomorphism}
设 $R$ 为环。设 $K^\bullet$ 为 K-平坦复形。那么，函子
$$
K(R) \longrightarrow K(R), \quad
L^\bullet \longmapsto \text{Tot}(L^\bullet \otimes_R K^\bullet)
$$
将拟同构变为拟同构。
\end{lemma}

\begin{proof}
这由引理 \ref{more-algebra-lemma-derived-tor-exact} 以及以下事实得到：
$K(R)$ 中的拟同构由其锥无环刻画。
\end{proof}

\begin{lemma}
\label{more-algebra-lemma-base-change-K-flat}
设 $R \to R'$ 为环同态。若 $K^\bullet$ 是 $R$-模的 K-平坦复形，
则 $K^\bullet \otimes_R R'$ 是 $R'$-模的 K-平坦复形。
\end{lemma}

\begin{proof}
这由定义以及以下事实得到：对任意 $R'$-模复形 $L^\bullet$，有
$(K^\bullet \otimes_R R') \otimes_{R'} L^\bullet =
K^\bullet \otimes_R L^\bullet$。
\end{proof}

\begin{lemma}
\label{more-algebra-lemma-tensor-product-K-flat}
设 $R$ 为环。若 $K^\bullet$、$L^\bullet$ 是 $R$-模的 K-平坦复形，
则 $\text{Tot}(K^\bullet \otimes_R L^\bullet)$ 是 $R$-模的 K-平坦复形。
\end{lemma}

\begin{proof}
这由同构
$$
\text{Tot}(M^\bullet \otimes_R \text{Tot}(K^\bullet \otimes_R L^\bullet))
=
\text{Tot}(\text{Tot}(M^\bullet \otimes_R K^\bullet) \otimes_R L^\bullet)
$$
和定义得到。
\end{proof}

\begin{lemma}
\label{more-algebra-lemma-K-flat-two-out-of-three}
设 $R$ 为环。设 $(K_1^\bullet, K_2^\bullet, K_3^\bullet)$
为 $K(R)$ 中的可区别三角。若 $K_i^\bullet$ 中三个有两个 K-平坦，
则第三个也 K-平坦。
\end{lemma}

\begin{proof}
这由引理 \ref{more-algebra-lemma-derived-tor-exact} 以及以下事实得到：
$K(R)$ 中的可区别三角若三个对象有两个无环，则第三个也无环。
\end{proof}

\begin{lemma}
\label{more-algebra-lemma-K-flat-two-out-of-three-ses}
设 $R$ 为环。设
$0 \to K_1^\bullet \to K_2^\bullet \to K_3^\bullet \to 0$
为复形的短正合序列。若对所有 $n \in \mathbf{Z}$，$K_3^n$
平坦，并且 $K_i^\bullet$ 中三个有两个 K-平坦，则第三个也 K-平坦。
\end{lemma}

\begin{proof}
设 $L^\bullet$ 为 $R$-模复形。那么
$$
0 \to
\text{Tot}(L^\bullet \otimes_R K_1^\bullet) \to
\text{Tot}(L^\bullet \otimes_R K_2^\bullet) \to
\text{Tot}(L^\bullet \otimes_R K_3^\bullet) \to 0
$$
是复形的短正合序列。事实上，对每个 $n, m$，模序列
$0 \to L^n \otimes_R K_1^m \to
L^n \otimes_R K_2^m \to
L^n \otimes_R K_3^m \to 0$ 正合，这是由 Algebra，引理
\ref{algebra-lemma-flat-tor-zero} 得到的，而上述复形序列是这些序列的直和。
因此，引理由此以及以下事实得到：复形的短正合序列若三个有两个无环，
则第三个也无环。
\end{proof}

\begin{lemma}
\label{more-algebra-lemma-derived-tor-quasi-isomorphism}
设 $R$ 为环。设 $P^\bullet$ 为由平坦 $R$-模构成的上有界复形。
那么 $P^\bullet$ 是 K-平坦的。
\end{lemma}

\begin{proof}
设 $L^\bullet$ 为无环 $R$-模复形。设
$\xi \in H^n(\text{Tot}(L^\bullet \otimes_R P^\bullet))$。
我们必须证明 $\xi = 0$。由于
$\text{Tot}^n(L^\bullet \otimes_R P^\bullet)$ 是以
$L^a \otimes_R P^b$ 为项的直和，可知对某个 $m \in \mathbf{Z}$，
$\xi$ 来自
$H^n(\text{Tot}(\tau_{\leq m}L^\bullet \otimes_R P^\bullet))$
中的元素。由于 $\tau_{\leq m}L^\bullet$ 也无环，可以用它代替
$L^\bullet$，即用 $\tau_{\leq m}L^\bullet$ 作此替换。因此，
可以假设 $L^\bullet$ 上有界。在这种情况下，
Homology，引理 \ref{homology-lemma-first-quadrant-ss} 的谱序列满足
$$
{}'E_1^{p, q} = H^p(L^\bullet \otimes_R P^q)
$$
由于 $P^q$ 平坦且 $L^\bullet$ 无环，该项为零。因此
$H^*(\text{Tot}(L^\bullet \otimes_R P^\bullet)) = 0$。
\end{proof}

\noindent
在下一个引理中，复形系统的余极限是指逐项余极限。

\begin{lemma}
\label{more-algebra-lemma-colimit-K-flat}
设 $R$ 为环。设
$K_1^\bullet \to K_2^\bullet \to \ldots$ 为 K-平坦复形的系统。
那么 $\colim_i K_i^\bullet$ 是 K-平坦的。更一般地，
K-平坦复形的任意滤过余极限都是 K-平坦的。
\end{lemma}

\begin{proof}
由于取的是逐项余极限，由 Algebra，引理
\ref{algebra-lemma-tensor-products-commute-with-limits}，有
$$
\colim_i \text{Tot}(M^\bullet \otimes_R K_i^\bullet) =
\text{Tot}(M^\bullet \otimes_R \colim_i K_i^\bullet)
$$
因此，引理由滤过余极限正合这一事实得到；见 Algebra，引理
\ref{algebra-lemma-directed-colimit-exact}。
\end{proof}

\begin{lemma}
\label{more-algebra-lemma-universally-acyclic-K-flat}
设 $R$ 为环。设 $K^\bullet$ 为 $R$-模复形。若对所有有限表示
$R$-模 $M$，$K^\bullet \otimes_R M$ 都无环，则 $K^\bullet$ 是 K-平坦的。
\end{lemma}

\begin{proof}
我们将反复使用张量积与余极限交换这一事实（Algebra，引理
\ref{algebra-lemma-tensor-products-commute-with-limits}）。因此，
对任意 $R$-模 $M$，$K^\bullet \otimes_R M$ 都无环，因为任意
$R$-模都是有限表示 $R$-模 $M$ 的滤过余极限；见 Algebra，引理
\ref{algebra-lemma-module-colimit-fp}。设 $M^\bullet$ 为无环
$R$-模复形。我们必须证明
$\text{Tot}(M^\bullet \otimes_R K^\bullet)$ 无环。由于
$M^\bullet = \colim \tau_{\leq n} M^\bullet$（逐项余极限），有
$$
\text{Tot}(M^\bullet \otimes_R K^\bullet) =
\colim \text{Tot}(\tau_{\leq n} M^\bullet \otimes_R K^\bullet)
$$
其中截断如 Homology，第 \ref{homology-section-truncations} 节所述。
由于滤过余极限正合（Algebra，引理
\ref{algebra-lemma-directed-colimit-exact}），可以用
$\tau_{\leq n}M^\bullet$ 代替 $M^\bullet$，并假设 $M^\bullet$ 上有界。
在上有界情形下，可以写成
$M^\bullet = \colim \sigma_{\geq -n} M^\bullet$，其中复形
$\sigma_{\geq -n} M^\bullet$ 有界，但可能不再无环。
按上述方式论证，问题化约为 $M^\bullet$ 为有界复形的情形。
最后，对有界复形 $M^a \to \ldots \to M^b$，可以对复形的长度
$b - a$ 作归纳。情形 $b - a = 1$ 已在上面处理。
当 $b - a > 1$ 时，考虑复形的分裂短正合序列
$$
0 \to \sigma_{\geq a + 1}M^\bullet \to M^\bullet \to M^a[-a] \to 0
$$
并应用引理 \ref{more-algebra-lemma-derived-tor-exact} 完成归纳步骤。
略去一些细节。
\end{proof}

\begin{lemma}
\label{more-algebra-lemma-K-flat-resolution}
设 $R$ 为环。对任意复形 $M^\bullet$，都存在各项为平坦 $R$-模的
K-平坦复形 $K^\bullet$，以及逐项满射的拟同构
$K^\bullet \to M^\bullet$。
\end{lemma}

\begin{proof}
令 $\mathcal{P} \subset \Ob(\text{Mod}_R)$ 为平坦 $R$-模的类。
由 Derived Categories，引理 \ref{derived-lemma-special-direct-system}，
存在系统 $K_1^\bullet \to K_2^\bullet \to \ldots$ 和交换图
$$
\xymatrix{
K_1^\bullet \ar[d] \ar[r] &
K_2^\bullet \ar[d] \ar[r] & \ldots \\
\tau_{\leq 1}M^\bullet \ar[r] &
\tau_{\leq 2}M^\bullet \ar[r] & \ldots
}
$$
并具有该引理列出的性质 (1)、(2)、(3)。这些性质蕴含每个复形
$K_i^\bullet$ 都是由平坦模构成的上有界复形。因此，由引理
\ref{more-algebra-lemma-derived-tor-quasi-isomorphism}，$K_i^\bullet$ 是 K-平坦的。
依构造，诱导映射 $\colim_i K_i^\bullet \to M^\bullet$
是拟同构且逐项满射。由引理 \ref{more-algebra-lemma-colimit-K-flat}，复形
$\colim_i K_i^\bullet$ 是 K-平坦的。各项 $\colim K_i^n$ 平坦，
因为平坦模的滤过余极限平坦；见 Algebra，引理
\ref{algebra-lemma-colimit-flat}。
\end{proof}

\begin{remark}
\label{more-algebra-remark-P-resolution}
事实上，可以得到比引理 \ref{more-algebra-lemma-K-flat-resolution} 更强的结果。
具体地，可以找到拟同构 $P^\bullet \to M^\bullet$，其中
$P^\bullet$ 是赋有子复形滤过
$$
0 = F_{-1}P^\bullet \subset F_0P^\bullet \subset
F_1P^\bullet \subset \ldots \subset P^\bullet
$$
的 $R$-模复形，并满足
\begin{enumerate}
\item $P^\bullet = \bigcup F_pP^\bullet$；
\item 包含 $F_iP^\bullet \to F_{i + 1}P^\bullet$ 是逐项分裂单射；
\item 商 $F_{i + 1}P^\bullet/F_iP^\bullet$ 同构于平移
$R[k]$ 的直和（作为复形，因此其微分为零）。
\end{enumerate}
这将在 Differential Graded Algebra，引理 \ref{dga-lemma-resolve}
中证明（也可以像引理 \ref{more-algebra-lemma-resolution-filtered-complex}
的证明那样论证）。此外，给定这样的复形，在 $D(R)$ 中得到可区别三角
$$
\bigoplus F_iP^\bullet \to \bigoplus F_iP^\bullet \to M^\bullet
\to \bigoplus F_iP^\bullet[1]
$$
利用它，有时可以把关于一般复形的陈述化约为关于 $R[k]$
的陈述（当然，这只在该陈述于取直和下保持时有效）。更精确地，
设 $T$ 为 $D(R)$ 中对象的一种性质。假设
\begin{enumerate}
\item 若 $K_i \in D(R)$、$i \in I$ 是一族对象，且对所有
$i \in I$ 都有 $T(K_i)$，则 $T(\bigoplus K_i)$；
\item 若 $K \to L \to M \to K[1]$ 是可区别三角，并且 $T$
对其中两个对象成立，则 $T$ 对第三个对象也成立；
\item 对所有 $k$，$T(R[k])$ 成立。
\end{enumerate}
那么 $T$ 对 $D(R)$ 的所有对象成立。
\end{remark}

\begin{lemma}
\label{more-algebra-lemma-derived-tor-quasi-isomorphism-other-side}
设 $R$ 为环。设
$\alpha : P^\bullet \to Q^\bullet$ 为 $R$-模的 K-平坦复形之间的拟同构。
对每个 $R$-模复形 $L^\bullet$，诱导映射
$$
\text{Tot}(\text{id}_L \otimes \alpha) :
\text{Tot}(L^\bullet \otimes_R P^\bullet)
\longrightarrow
\text{Tot}(L^\bullet \otimes_R Q^\bullet)
$$
是拟同构。
\end{lemma}

\begin{proof}
选取拟同构 $K^\bullet \to L^\bullet$，其中 $K^\bullet$
为 K-平坦复形；见引理 \ref{more-algebra-lemma-K-flat-resolution}。考虑交换图
$$
\xymatrix{
\text{Tot}(K^\bullet \otimes_R P^\bullet) \ar[r] \ar[d] &
\text{Tot}(K^\bullet \otimes_R Q^\bullet) \ar[d] \\
\text{Tot}(L^\bullet \otimes_R P^\bullet) \ar[r] &
\text{Tot}(L^\bullet \otimes_R Q^\bullet)
}
$$
由引理 \ref{more-algebra-lemma-K-flat-quasi-isomorphism}，两条竖直箭头和上方水平箭头
都是拟同构，故结论成立。
\end{proof}

\noindent
设 $R$ 为环。设 $M^\bullet$ 为 $D(R)$ 的对象。选取 K-平坦解消
$K^\bullet \to M^\bullet$；见引理 \ref{more-algebra-lemma-K-flat-resolution}。
由引理 \ref{more-algebra-lemma-derived-tor-homotopy} 和
\ref{more-algebra-lemma-derived-tor-exact}，得到三角范畴的正合函子
$$
K(R) \longrightarrow K(R), \quad
L^\bullet \longmapsto \text{Tot}(L^\bullet \otimes_R K^\bullet)
$$
由引理 \ref{more-algebra-lemma-K-flat-quasi-isomorphism}，该函子诱导函子
$D(R) \to D(R)$，这仅仅是因为 $D(R)$ 是 $K(R)$
关于拟同构的局部化。由于 $M^\bullet$ 的 $K$-平坦解消范畴是余滤过的，
并且有引理 \ref{more-algebra-lemma-derived-tor-quasi-isomorphism-other-side}，
所得函子（在同构意义下）不依赖于 $K^\bullet$ 的选取。

\begin{definition}
\label{more-algebra-definition-derived-tor}
设 $R$ 为环。设 $M^\bullet$ 为 $D(R)$ 的对象。{\it 导出张量积}
$$
- \otimes_R^{\mathbf{L}} M^\bullet : D(R) \longrightarrow D(R)
$$
是上述三角范畴的正合函子。
\end{definition}

\noindent
该函子延拓函子（\ref{more-algebra-equation-derived-tensor-module}）。
由我们的显式构造显然可见，只要 $L^\bullet$ 和 $M^\bullet$
都属于 $D(R)$，就存在同构（涉及符号的选取，见下文）
$$
M^\bullet \otimes_R^{\mathbf{L}} L^\bullet
\cong
L^\bullet \otimes_R^{\mathbf{L}} M^\bullet
$$
因此，在写 $M^\bullet \otimes_R^{\mathbf{L}} L^\bullet$ 时，
我们通常不区分使用哪个变量来定义与之作导出张量积的函子。

\begin{lemma}
\label{more-algebra-lemma-flip-douoble-tensor-product}
设 $R$ 为环。设 $K^\bullet, L^\bullet$ 为 $R$-模复形。
存在关于两个复形都具有函子性的典范同构
$$
K^\bullet \otimes_R^\mathbf{L} L^\bullet \longrightarrow
L^\bullet \otimes_R^\mathbf{L} K^\bullet
$$
它在映射 $K^p \otimes_R L^q \to L^q \otimes_R K^p$ 上使用符号
$(-1)^{pq}$（解释见证明）。
\end{lemma}

\begin{proof}
我们可以并且确实用 K-平坦复形 $K^\bullet$ 和 $L^\bullet$
代替原复形，然后使用第 \ref{more-algebra-section-symmetric-monoidal} 节中讨论的交换约束。
\end{proof}

\begin{lemma}
\label{more-algebra-lemma-triple-tensor-product}
设 $R$ 为环。设 $K^\bullet, L^\bullet, M^\bullet$ 为 $R$-模复形。
存在关于三个复形都具有函子性的典范同构
$$
(K^\bullet \otimes_R^\mathbf{L} L^\bullet) \otimes_R^\mathbf{L} M^\bullet
=
K^\bullet \otimes_R^\mathbf{L} (L^\bullet \otimes_R^\mathbf{L} M^\bullet)
$$
\end{lemma}

\begin{proof}
用 K-平坦复形代替这些复形，并使用第
\ref{more-algebra-section-symmetric-monoidal} 节中的结合约束。
\end{proof}

\begin{lemma}
\label{more-algebra-lemma-factor-through-K-flat}
设 $R$ 为环。设 $a : K^\bullet \to L^\bullet$ 为 $R$-模复形的映射。
若 $K^\bullet$ K-平坦，则存在复形 $N^\bullet$ 以及复形映射
$b : K^\bullet \to N^\bullet$ 和 $c : N^\bullet \to L^\bullet$，使得
\begin{enumerate}
\item $N^\bullet$ K-平坦；
\item $c$ 是拟同构；
\item $a$ 与 $c \circ b$ 同伦。
\end{enumerate}
若 $K^\bullet$ 的各项平坦，则可以选取 $N^\bullet$、$b$ 和 $c$，
使得 $N^\bullet$ 也具有同样的性质。
\end{lemma}

\begin{proof}
我们将使用同伦范畴 $K(R)$ 是三角范畴这一事实；见 Derived Categories，
命题 \ref{derived-proposition-homotopy-category-triangulated}。
选取可区别三角
$K^\bullet \to L^\bullet \to C^\bullet \to K^\bullet[1]$。
选取拟同构 $M^\bullet \to C^\bullet$，其中 $M^\bullet$
K-平坦且各项平坦；见引理 \ref{more-algebra-lemma-K-flat-resolution}。
由三角范畴的公理，可以把复合
$M^\bullet \to C^\bullet \to K^\bullet[1]$ 嵌入可区别三角
$K^\bullet \to N^\bullet \to M^\bullet \to K^\bullet[1]$。
由引理 \ref{more-algebra-lemma-K-flat-two-out-of-three}，$N^\bullet$ K-平坦。
再次使用三角范畴的公理，可以选取映射
$N^\bullet \to L^\bullet$，将其嵌入下列可区别三角的态射：
$$
\xymatrix{
K^\bullet \ar[r] \ar[d] &
N^\bullet \ar[r] \ar[d] &
M^\bullet \ar[r] \ar[d] &
K^\bullet[1] \ar[d] \\
K^\bullet \ar[r] &
L^\bullet \ar[r] &
C^\bullet \ar[r] &
K^\bullet[1]
}
$$
三条箭头中有两条是拟同构，故由这些可区别三角所伴随的上同调长正合序列，
第三条箭头 $N^\bullet \to L^\bullet$ 也是拟同构（也可以考察该图在
$D(R)$ 中的像，并使用 Derived Categories，引理
\ref{derived-lemma-third-isomorphism-triangle}）。这完成了 (1)、(2)、(3)
的证明。为证明最后一个断言，可以选取 $N^\bullet$，使得
$N^n \cong M^n \oplus K^n$；见 Derived Categories，引理
\ref{derived-lemma-improve-distinguished-triangle-homotopy}。
因此，若 $K^\bullet$ 的各项平坦，就得到所需的平坦性。
\end{proof}








\section{导出换环}
\label{more-algebra-section-derived-base-change}

\noindent
设 $R \to A$ 为环同态。设 $N^\bullet$ 为 $A$-模复形。
也可以使用 K-平坦解消，把函子
$$
- \otimes_R^{\mathbf{L}} N^\bullet : D(R) \to D(A)
$$
定义为函子 $K(R) \to K(A)$，
$M^\bullet \mapsto \text{Tot}(M^\bullet \otimes_R N^\bullet)$
的左导出函子。特别地，取 $N^\bullet = A[0]$，得到导出基变换函子
$$
- \otimes_R^{\mathbf{L}} A : D(R) \to D(A)
$$
它延拓函子（\ref{more-algebra-equation-derived-tensor-algebra}）。具体地，
对每个 $R$-模复形 $M^\bullet$，可以选取 K-平坦解消
$K^\bullet \to M^\bullet$，并置
$$
M^\bullet \otimes_R^{\mathbf{L}} N^\bullet =
\text{Tot}(K^\bullet \otimes_R N^\bullet).
$$
可以使用引理 \ref{more-algebra-lemma-K-flat-resolution} 和
\ref{more-algebra-lemma-derived-tor-quasi-isomorphism-other-side} 看出该定义良好。
不过，为了把所有细节都交代周全，使用一些一般理论或许更方便。

\begin{lemma}
\label{more-algebra-lemma-derived-base-change}
上述构造不依赖于选取，并定义三角范畴的正合函子
$- \otimes_R^\mathbf{L} N^\bullet : D(R) \to D(A)$。
对 $D(R)$ 中的 $E^\bullet$，存在函子性同构
$$
E^\bullet \otimes_R^\mathbf{L} N^\bullet =
(E^\bullet \otimes_R^\mathbf{L} A) \otimes_A^\mathbf{L} N^\bullet
$$
\end{lemma}

\begin{proof}
为了证明导出函子 $- \otimes_R^\mathbf{L} N^\bullet$ 的存在性，
使用 Derived Categories，第 \ref{derived-section-derived-functors} 节
中发展的一般理论。置 $\mathcal{D} = K(R)$ 和
$\mathcal{D}' = D(A)$。记 $F : \mathcal{D} \to \mathcal{D}'$
为由规则 $F(M^\bullet) = \text{Tot}(M^\bullet \otimes_R N^\bullet)$
定义的三角范畴正合函子。要证明 $F$ 的上述性质，使用引理
\ref{more-algebra-lemma-derived-tor-homotopy} 和 \ref{more-algebra-lemma-derived-tor-exact}。
令 $S$ 为 $\mathcal{D} = K(R)$ 中所有拟同构构成的集合。
这给出 Derived Categories，情形
\ref{derived-situation-derived-functor} 中的情况，因此可应用
Derived Categories，定义
\ref{derived-definition-right-derived-functor-defined}。
我们断言 $LF$ 处处有定义。由 Derived Categories，引理
\ref{derived-lemma-find-existence-computes} 可得这一点，其中取
$\mathcal{P} \subset \Ob(\mathcal{D})$ 为 K-平坦复形的集合：
(1) 由引理 \ref{more-algebra-lemma-K-flat-resolution} 得到，(2) 由引理
\ref{more-algebra-lemma-derived-tor-quasi-isomorphism-other-side} 得到。
因此，得到导出函子
$$
LF : D(R) = S^{-1}\mathcal{D} \longrightarrow \mathcal{D}' = D(A)
$$
见 Derived Categories，等式（\ref{derived-equation-everywhere}）。
最后，Derived Categories，引理
\ref{derived-lemma-find-existence-computes} 保证：当 $K^\bullet$
K-平坦时，
$LF(K^\bullet) = F(K^\bullet) =
\text{Tot}(K^\bullet \otimes_R N^\bullet)$；换言之，$LF$
确实按上述方式计算。此外，由引理 \ref{more-algebra-lemma-base-change-K-flat}，
复形 $K^\bullet \otimes_R A$ 是 $A$-模的 K-平坦复形。因此
$$
(K^\bullet \otimes_R^\mathbf{L} A) \otimes_A^\mathbf{L} N^\bullet =
\text{Tot}((K^\bullet \otimes_R A) \otimes_A N^\bullet) =
\text{Tot}(K^\bullet \otimes_A N^\bullet) =
K^\bullet \otimes_A^\mathbf{L} N^\bullet
$$
这证明了引理的最后一个陈述。
\end{proof}

\begin{lemma}
\label{more-algebra-lemma-functoriality-derived-base-change}
设 $R \to A$ 为环同态。设 $f : L^\bullet \to N^\bullet$
为 $A$-模复形的映射。那么 $f$ 诱导函子变换
$$
1 \otimes f :
- \otimes_A^\mathbf{L} L^\bullet
\longrightarrow
- \otimes_A^\mathbf{L} N^\bullet
$$
若 $f$ 是拟同构，则 $1 \otimes f$ 是函子的同构。
\end{lemma}

\begin{proof}
由于这些函子通过在 K-平坦复形 $K^\bullet$ 上取值来计算，
可以直接使用函子性
$$
\text{Tot}(K^\bullet \otimes_R L^\bullet) \to
\text{Tot}(K^\bullet \otimes_R N^\bullet)
$$
定义该变换。最后一个陈述由引理
\ref{more-algebra-lemma-K-flat-quasi-isomorphism} 得到。
\end{proof}

\begin{lemma}
\label{more-algebra-lemma-tensor-hom-adjoint}
设 $R \to A$ 为环同态。引理 \ref{more-algebra-lemma-derived-base-change}
的函子 $D(R) \to D(A)$，$E \mapsto E \otimes_R^\mathbf{L} A$，
是限制函子 $D(A) \to D(R)$ 的左伴随。
\end{lemma}

\begin{proof}
这由 Derived Categories，引理
\ref{derived-lemma-pre-derived-adjoint-functors-general}，
以及由 Algebra，引理 \ref{algebra-lemma-adjoint-tensor-restrict}
得到的事实，即 $- \otimes_R A$ 与限制互为伴随，得到。
\end{proof}

\begin{remark}[警告]
\label{more-algebra-remark-warning-compute-base-change}
设 $R \to A$ 为环同态，并设 $N$ 和 $N'$ 为 $A$-模。
以 $N_R$ 和 $N'_R$ 表示把 $N$ 和 $N'$ 限制为 $R$-模所得的模；
见 Algebra，第 \ref{algebra-section-base-change} 节。在这种情况下，
$D(A)$ 中的对象 $N_R \otimes_R^\mathbf{L} N'$ 和
$N \otimes_R^\mathbf{L} N'_R$ 一般不同构！换言之，必须仔细分辨
两侧中的哪一侧用于给出 $A$-模结构。

\medskip\noindent
作为具体例子，置 $R = k[x, y]$、$A = R/(xy)$、$N = R/(x)$
以及 $N' = A = R/(xy)$。解消
$0 \to R \xrightarrow{xy} R \to N'_R \to 0$ 表明在 $D(A)$ 中，
$N \otimes_R^\mathbf{L} N'_R = N[1] \oplus N$。解消
$0 \to R \xrightarrow{x} R \to N_R \to 0$ 表明
$N_R \otimes_R^\mathbf{L} N'$ 由复形 $A \xrightarrow{x} A$ 表示。
要看出这两个复形不同构，可以证明第二个复形在 $D(A)$ 中不同构于
其上同调群的直和，也可以证明第一个复形不是 $D(A)$ 的完美对象，
而第二个是。略去一些细节。
\end{remark}

\begin{lemma}
\label{more-algebra-lemma-double-base-change}
设 $A \to B \to C$ 为环同态。设 $N^\bullet$ 为 $B$-模复形，
$K^\bullet$ 为 $C$-模复形。下列函子的复合
$$
D(A) \xrightarrow{- \otimes_A^\mathbf{L} N^\bullet}
D(B) \xrightarrow{- \otimes_B^\mathbf{L} K^\bullet} D(C)
$$
是函子
$- \otimes_A^\mathbf{L} (N^\bullet \otimes_B^\mathbf{L} K^\bullet) :
D(A) \to D(C)$。如果 $M$、$N$、$K$ 分别是 $A$、$B$、$C$ 上的模，
则在 $D(C)$ 中有
$$
(M \otimes_A^\mathbf{L} N) \otimes_B^\mathbf{L} K =
M \otimes_A^\mathbf{L} (N \otimes_B^\mathbf{L} K) =
(M \otimes_A^\mathbf{L} C) \otimes_C^\mathbf{L} (N \otimes_B^\mathbf{L} K)
$$
还存在典范同构
$$
(M \otimes_A^\mathbf{L} N) \otimes_B^\mathbf{L} K \longrightarrow
(M \otimes_A^\mathbf{L} K) \otimes_C^\mathbf{L} (N \otimes_B^\mathbf{L} C)
$$
其中使用符号。对复形也有类似结果。
\end{lemma}

\begin{proof}
选取 $B$-模的 K-平坦复形 $P^\bullet$ 以及拟同构
$P^\bullet \to N^\bullet$（引理 \ref{more-algebra-lemma-K-flat-resolution}）。
设 $M^\bullet$ 为表示 $D(A)$ 中任意对象的 $A$-模 K-平坦复形。
那么，由引理 \ref{more-algebra-lemma-functoriality-derived-base-change}
应用于括号内的对象可知，
$$
(M^\bullet \otimes_A^\mathbf{L} P^\bullet) \otimes_B^\mathbf{L} K^\bullet
\longrightarrow
(M^\bullet \otimes_A^\mathbf{L} N^\bullet) \otimes_B^\mathbf{L} K^\bullet
$$
是同构。由引理 \ref{more-algebra-lemma-base-change-K-flat} 和
\ref{more-algebra-lemma-tensor-product-K-flat}，复形
$$
\text{Tot}(M^\bullet \otimes_A P^\bullet) =
\text{Tot}((M^\bullet \otimes_R A) \otimes_A P^\bullet
$$
作为 $B$-模复形是 K-平坦的，并且按构造表示 $D(B)$ 中的导出张量积。
因此，
$(M^\bullet \otimes_A^\mathbf{L} P^\bullet) \otimes_B^\mathbf{L} K^\bullet$
由 $C$-模复形
$$
\text{Tot}(\text{Tot}(M^\bullet \otimes_A P^\bullet)\otimes_B K^\bullet) =
\text{Tot}(M^\bullet \otimes_A \text{Tot}(P^\bullet \otimes_B K^\bullet))
$$
表示。该等式见 Homology，评注 \ref{homology-remark-triple-complex}。
沿原路返回可见，这等于
$$
M^\bullet \otimes_A^\mathbf{L} (P^\bullet \otimes_B^\mathbf{L} K^\bullet)
\longleftarrow
M^\bullet \otimes_A^\mathbf{L} (N^\bullet \otimes_B^\mathbf{L} K^\bullet)
$$
按照函子 $-\otimes_B^\mathbf{L} K^\bullet$ 的定义，该箭头是同构。
所有这些构造关于复形 $M^\bullet$ 都具有函子性，故得到所述函子同构。

\medskip\noindent
由上文，下式中的第一个等式成立：
$$
(M \otimes_A^\mathbf{L} N) \otimes_B^\mathbf{L} K =
M \otimes_A^\mathbf{L} (N \otimes_B^\mathbf{L} K) =
(M \otimes_A^\mathbf{L} C) \otimes_C^\mathbf{L} (N \otimes_B^\mathbf{L} K)
$$
第二个等式由引理 \ref{more-algebra-lemma-derived-base-change} 的最后一个陈述得到。
同一陈述还允许我们写成
$N \otimes_B^\mathbf{L} K = (N \otimes_B^\mathbf{L} C) \otimes_C^\mathbf{L} K$，
代入即得
\begin{align*}
(M \otimes_A^\mathbf{L} N) \otimes_B^\mathbf{L} K
& =
(M \otimes_A^\mathbf{L} C) \otimes_C^\mathbf{L}
((N \otimes_B^\mathbf{L} C) \otimes_C^\mathbf{L} K) \\
& =
(M \otimes_A^\mathbf{L} C) \otimes_C^\mathbf{L}
(K \otimes_C^\mathbf{L} (N \otimes_B^\mathbf{L} C)) \\
& =
((M \otimes_A^\mathbf{L} C) \otimes_C^\mathbf{L} K)
\otimes_C^\mathbf{L} (N \otimes_B^\mathbf{L} C)) \\
& =
(M \otimes_C^\mathbf{L} K)
\otimes_C^\mathbf{L} (N \otimes_B^\mathbf{L} C)
\end{align*}
这里使用了引理 \ref{more-algebra-lemma-flip-douoble-tensor-product}、
\ref{more-algebra-lemma-triple-tensor-product} 以及前面提到的引理。
\end{proof}







\section{Tor 无关性}
\label{more-algebra-section-tor-independence}

\noindent
考虑环的交换图
$$
\xymatrix{
A \ar[r] & A' \\
R \ar[r] \ar[u] & R' \ar[u]
}
$$
给定 $D(A)$ 的对象 $K$，可以考虑它到 $D(A')$ 的对象
$K \otimes_A^\mathbf{L} A'$ 的导出基变换。也可以先把 $K$ 限制为
$D(R)$ 的对象，再考虑它到 $D(R')$ 的导出基变换，记为
$K \otimes_R^\mathbf{L} R'$。我们断言，在 $D(R')$ 中存在函子性比较映射
\begin{equation}
\label{more-algebra-equation-comparison-map}
K \otimes_R^{\mathbf{L}} R'
\longrightarrow
K \otimes_A^{\mathbf{L}} A'
\end{equation}
为构造这个比较映射，选取表示 $K$ 的 $A$-模 K-平坦复形
$K^\bullet$。接着选取拟同构 $E^\bullet \to K^\bullet$，其中
$E^\bullet$ 是 $R$-模的 K-平坦复形。上述映射就是
$$
K \otimes_R^{\mathbf{L}} R' =
E^\bullet \otimes_R R'
\longrightarrow
K^\bullet \otimes_A A' =
K \otimes_A^{\mathbf{L}} A'
$$
一般而言，该映射不可能是同构。

\medskip\noindent
不过，我们经常遇到上述图是环的“基变换”图，即
$A' = A \otimes_R R'$ 的情形。在这种情况下，对任意 $A$-模 $M$，都有
$M \otimes_A A' = M \otimes_R R'$。因此，作为函子
$\text{Mod}_A \to \text{Mod}_{A'}$，$- \otimes_R R'$ 等于
$- \otimes_A A'$。一般而言，这个等式{\bf 并不延拓}到导出张量积。
换言之，比较映射不是同构。一个简单例子是取
$R = k[x]$、$A = R' = A' = k[x]/(x) = k$ 和 $K^\bullet = A[0]$。
显然，一个必要条件是对所有 $p > 0$ 都有
$\text{Tor}_p^R(A, R') = 0$。

\begin{definition}
\label{more-algebra-definition-tor-independent}
设 $R$ 为环。设 $A$、$B$ 为 $R$-代数。如果对所有 $p > 0$ 都有
$\text{Tor}_p^R(A, B) = 0$，则称 $A$ 与 $B$ {\it 在 $R$ 上 Tor 无关}。
\end{definition}

\begin{lemma}
\label{more-algebra-lemma-base-change-comparison}
如果 $A' = A \otimes_R R'$，且 $A$ 与 $R'$ 在 $R$ 上 Tor 无关，
则比较映射（\ref{more-algebra-equation-comparison-map}）是同构。
\end{lemma}

\begin{proof}
为证明这一点，选取 $R'$ 作为 $R$-模的自由解消
$F^\bullet \to R'$。由于 $A$ 与 $R'$ 在 $R$ 上 Tor 无关，
$F^\bullet \otimes_R A$ 是 $A'$ 在 $A$ 上的自由 $A$-模解消。
由上文对导出张量积的一般构造可见
$$
K^\bullet \otimes_A A' \cong
\text{Tot}(K^\bullet \otimes_A (F^\bullet \otimes_R A)) =
\text{Tot}(K^\bullet \otimes_R F^\bullet) \cong
\text{Tot}(E^\bullet \otimes_R F^\bullet) \cong
E^\bullet \otimes_R R'
$$
这正是所需结论。
\end{proof}

\begin{lemma}
\label{more-algebra-lemma-tor-independent-flat}
考虑环的交换图
$$
\xymatrix{
A' & R' \ar[r] \ar[l] & B' \\
A \ar[u] & R \ar[l] \ar[u] \ar[r] & B \ar[u]
}
$$
假设 $R'$ 在 $R$ 上平坦，$A'$ 在 $A \otimes_R R'$ 上平坦，
且 $B'$ 在 $R' \otimes_R B$ 上平坦。那么
$$
\text{Tor}_i^R(A, B) \otimes_{(A \otimes_R B)} (A' \otimes_{R'} B') =
\text{Tor}_i^{R'}(A', B')
$$
\end{lemma}

\begin{proof}
由 Algebra，第 \ref{algebra-section-functoriality-tor} 节，存在典范映射
$$
\text{Tor}_i^R(A, B) \longrightarrow
\text{Tor}_i^{R'}(A \otimes_R R', B \otimes_R R') \longrightarrow
\text{Tor}_i^{R'}(A', B')
$$
它们诱导从引理公式左侧到右侧的映射。

\medskip\noindent
取 $A$ 作为 $R$-模的自由解消 $F_\bullet \to A$。
那么 $F_\bullet \otimes_R R'$ 是 $A \otimes_R R'$ 的解消。
因而 $\text{Tor}_i^{R'}(A \otimes_R R', B \otimes_R R')$ 由
$F_\bullet \otimes_R B \otimes_R R'$ 计算。根据 $R'$ 在 $R$ 上平坦的假设，
它计算 $\text{Tor}_i^R(A, B) \otimes_R R'$。因此
$\text{Tor}_i^{R'}(A \otimes_R R', B \otimes_R R') =
\text{Tor}_i^R(A, B) \otimes_R R'$（这里只使用了 $R'$ 在 $R$ 上的平坦性）。

\medskip\noindent
由 Lazard 定理（Algebra，定理 \ref{algebra-theorem-lazard}），可以分别把
$A'$ 和 $B'$ 写成有限自由 $A \otimes_R R'$-模和
$B \otimes_R R'$-模的滤过余极限。记
$A' = \colim M_i$ 和 $B' = \colim N_j$。上面的结果给出
$$
\text{Tor}_i^{R'}(M_i, N_j) =
\text{Tor}_i^R(A, B) \otimes_{A \otimes_R B} (M_i \otimes_{R'} N_j)
$$
把一切用基写出即可看出这一点。取余极限便得到引理的结论。
\end{proof}

\begin{lemma}
\label{more-algebra-lemma-flat-base-change-tor-independent}
设 $R \to A$ 和 $R \to B$ 为环同态。设 $R \to R'$ 为环同态，
并置 $A' = A \otimes_R R'$ 和 $B' = B \otimes_R R'$。
如果 $A$ 与 $B$ 在 $R$ 上 Tor 无关，且 $R \to R'$ 平坦，
则 $A'$ 与 $B'$ 在 $R'$ 上 Tor 无关。
\end{lemma}

\begin{proof}
这立即由引理 \ref{more-algebra-lemma-tor-independent-flat} 和定义
\ref{more-algebra-definition-tor-independent} 得到。
\end{proof}

\begin{lemma}
\label{more-algebra-lemma-lemma-tor-independent-flat-compare}
假设同引理 \ref{more-algebra-lemma-tor-independent-flat}。对 $M \in D(A)$，
存在 $A' \otimes_{R'} B'$-模的典范同构
$$
H^i((M \otimes_A^\mathbf{L} A') \otimes_{R'}^\mathbf{L} B') =
H^i(M \otimes_R^\mathbf{L} B) \otimes_{(A \otimes_R B)} (A' \otimes_{R'} B')
$$
\end{lemma}

\begin{proof}
我们来阐明等式的两侧。在左侧，有函子
$D(A) \to D(A') \to D(R') \to D(B')$ 的复合，再与函子
$H^i : D(B') \to \text{Mod}_{B'}$ 复合。由于存在从 $A'$ 到
$D(B')$ 中对象 $(M \otimes_A^\mathbf{L} A') \otimes_{R'}^\mathbf{L} B'$
的自同态环的映射，可见左侧确实是 $A' \otimes_{R'} B'$-模。
由同样的论证可见，$H^i(M \otimes_R^\mathbf{L} B)$ 具有
$A \otimes_R B$-模结构。

\medskip\noindent
先在 $B' = R' \otimes_R B$ 的情形证明该结论。
在这种情况下，选取由自由 $R$-模组成的解消 $F^\bullet \to B$。
还选取表示 $M$ 的 $A$-模 K-平坦复形 $M^\bullet$。
于是左侧由下式表示：
\begin{align*}
H^i(\text{Tot}((M^\bullet \otimes_A A') \otimes_{R'} (R' \otimes_R F^\bullet)))
& =
H^i(\text{Tot}(M^\bullet \otimes_A A' \otimes_R F^\bullet)) \\
& =
H^i(\text{Tot}(M^\bullet \otimes_R F^\bullet) \otimes_A A') \\
& =
H^i(M \otimes_R^\mathbf{L} B) \otimes_A A'
\end{align*}
最后一个等式成立，是因为 $A \to A'$ 平坦。最后一个模正是所需的模，
因为本段假设 $B' = R' \otimes_R B$，故
$A' \otimes_{R'} B' = A' \otimes_R B$。

\medskip\noindent
现在考虑一般情形。假设 $B' \to B''$ 是平坦环同态。
容易看出
$$
H^i((M \otimes_A^\mathbf{L} A') \otimes_{R'}^\mathbf{L} B'') =
H^i((M \otimes_A^\mathbf{L} A') \otimes_{R'}^\mathbf{L} B')
\otimes_{B'} B''
$$
并且
$$
H^i(M \otimes_R^\mathbf{L} B) \otimes_{(A \otimes_R B)} (A' \otimes_{R'} B'')
=
\left(
H^i(M \otimes_R^\mathbf{L} B) \otimes_{(A \otimes_R B)} (A' \otimes_{R'} B')
\right) \otimes_{B'} B''
$$
因此，对 $B'$ 的结论推出对 $B''$ 的结论。上一段已经证明了
$R' \otimes_R B$ 的情形，故这推出一般结论。
\end{proof}

\begin{lemma}
\label{more-algebra-lemma-tor-independent}
设 $R$ 为环。设 $A$、$B$ 为 $R$-代数。下列条件等价：
\begin{enumerate}
\item $A$ 与 $B$ 在 $R$ 上 Tor 无关；
\item 对位于同一素理想 $\mathfrak r \subset R$ 之上的每一对素理想
$\mathfrak p \subset A$ 和 $\mathfrak q \subset B$，环 $A_\mathfrak p$
与 $B_\mathfrak q$ 在 $R_\mathfrak r$ 上 Tor 无关；
\item 对 $A \otimes_R B$ 的每个素理想 $\mathfrak s$，模
$$
\text{Tor}_i^R(A, B)_\mathfrak s =
\text{Tor}_i^{R_\mathfrak r}(A_\mathfrak p, B_\mathfrak q)_\mathfrak s
$$
为零（其中 $\mathfrak p = A \cap \mathfrak s$、
$\mathfrak q = B \cap \mathfrak s$ 且 $\mathfrak r = R \cap \mathfrak s$）。
\end{enumerate}
\end{lemma}

\begin{proof}
设 $\mathfrak s$ 为 (3) 中 $A \otimes_R B$ 的素理想。等式
$$
\text{Tor}_i^R(A, B)_\mathfrak s =
\text{Tor}_i^{R_\mathfrak r}(A_\mathfrak p, B_\mathfrak q)_\mathfrak s
$$
其中 $\mathfrak p = A \cap \mathfrak s$、$\mathfrak q = B \cap \mathfrak s$
且 $\mathfrak r = R \cap \mathfrak s$，由引理
\ref{more-algebra-lemma-tor-independent-flat} 得到。因此 (2) 推出 (3)。
由于模的消失性可以通过在素理想处局部化来检验（Algebra，引理
\ref{algebra-lemma-characterize-zero-local}），可知 (3) 推出 (1)。
为证明 (1) $\Rightarrow$ (2)，再次使用引理
\ref{more-algebra-lemma-tor-independent-flat} 给出的等式
$$
\text{Tor}_i^{R_\mathfrak r}(A_\mathfrak p, B_\mathfrak q) =
\text{Tor}_i^R(A, B) \otimes_{(A \otimes_R B)}
(A_\mathfrak p \otimes_{R_{\mathfrak r}} B_\mathfrak q)
$$
\end{proof}








\section{Tor 的谱序列}
\label{more-algebra-section-spectral-sequence-tor}


\noindent
本节汇集考察 Tor 函子时出现的各种谱序列。

\begin{example}
\label{more-algebra-example-cohomology-complex-tensored}
设 $R$ 为环。设 $K_\bullet$ 为 $R$-模链复形，并且当 $n \ll 0$ 时
$K_n = 0$。设 $M$ 为 $R$-模。选取由自由 $R$-模组成的 $M$ 的解消
$P_\bullet \to M$。得到双链复形 $K_\bullet \otimes_R P_\bullet$。
应用 Homology，第 \ref{homology-section-double-complex} 节的材料
（特别是 Homology，引理 \ref{homology-lemma-first-quadrant-ss}），
并将其改写为链复形的语言，得到两个收敛到
$H_*(K_\bullet \otimes_R^\mathbf{L} M)$ 的谱序列。一方面，存在
$E_2$-页为
$$
(E_2)_{i, j} = \text{Tor}^R_j(H_i(K_\bullet), M)
\Rightarrow
H_{i + j}(K_\bullet \otimes^{\mathbf{L}}_R M)
$$
的谱序列，其微分 $d_2$ 由映射
$\text{Tor}^R_j(H_i(K_\bullet), M) \to
\text{Tor}^R_{j - 2}(H_{i + 1}(K_\bullet), M)$ 给出。
另一方面，存在 $E_1$-页为
$$
(E_1)_{i, j} = \text{Tor}^R_j(K_i, M)
\Rightarrow
H_{i + j}(K_\bullet \otimes^{\mathbf{L}}_R M)
$$
的谱序列，其微分 $d_1$ 由 $K_i \to K_{i - 1}$ 诱导的映射
$\text{Tor}^R_j(K_i, M) \to \text{Tor}^R_j(K_{i - 1}, M)$ 给出。
\end{example}

\begin{example}
\label{more-algebra-example-tor-change-rings}
设 $R \to S$ 为环同态。设 $M$ 为 $R$-模，$N$ 为 $S$-模。
那么存在谱序列
$$
\text{Tor}^S_n(\text{Tor}^R_m(M, S), N) \Rightarrow
\text{Tor}^R_{n + m}(M, N).
$$
为构造它，选取 $M$ 的一个 $R$-自由解消 $P^\bullet$。于是
$$
M \otimes_R^{\mathbf{L}} N = P^\bullet \otimes_R N =
(P^\bullet \otimes_R S) \otimes_S N
$$
然后应用例 \ref{more-algebra-example-cohomology-complex-tensored} 的第一个谱序列。
\end{example}

\begin{example}
\label{more-algebra-example-tor-base-change}
考虑交换图
$$
\xymatrix{
B \ar[r] & B' = B \otimes_A A' \\
A \ar[r] \ar[u] & A' \ar[u]
}
$$
以及 $B$-模 $M, N$。置 $M' = M \otimes_A A' = M \otimes_B B'$ 和
$N' = N \otimes_A A' = N \otimes_B B'$。
{\it 假设 $A \to B$ 平坦，且 $M$ 和 $N$ 都在 $A$ 上平坦。}
那么存在谱序列
$$
\text{Tor}^A_i(\text{Tor}_j^B(M, N), A')
\Rightarrow
\text{Tor}^{B'}_{i + j}(M', N')
$$
理由如下。选取一个由自由 $B$-模组成的解消 $F_\bullet \to M$。
由于 $B$ 和 $M$ 都在 $A$ 上平坦，可见
$F_\bullet \otimes_A A'$ 是 $M'$ 的自由 $B'$-解消。
因此，群 $\text{Tor}^{B'}_n(M', N')$ 由复形
$$
(F_\bullet \otimes_A A') \otimes_{B'} N' =
(F_\bullet \otimes_B N) \otimes_A A' =
(F_\bullet \otimes_B N) \otimes^{\mathbf{L}}_A A'
$$
计算。最后一个等式成立，是因为 $N$ 在 $A$ 上平坦，所以
$F_\bullet \otimes_B N$ 是由平坦 $A$-模组成的复形。
于是应用例 \ref{more-algebra-example-cohomology-complex-tensored} 的谱序列便得到所述谱序列。
\end{example}

\begin{example}
\label{more-algebra-example-tor}
设 $K^\bullet, L^\bullet$ 为 $D^{-}(R)$ 的对象。那么存在谱序列
$$
E_2^{p, q} = H^p(K^\bullet \otimes_R^{\mathbf{L}} H^q(L^\bullet))
\Rightarrow H^{p + q}(K^\bullet \otimes_R^{\mathbf{L}} L^\bullet)
$$
以及另一个谱序列
$$
E_2^{p, q} =
H^p(H^q(K^\bullet) \otimes_R^{\mathbf{L}} L^\bullet)
\Rightarrow H^{p + q}(K^\bullet \otimes_R^{\mathbf{L}} L^\bullet)
$$
两个谱序列都有
$d_2^{p, q} : E_2^{p, q} \to E_2^{p + 2, q - 1}$。
用有上界的投射模复形代替 $K^\bullet$ 和 $L^\bullet$ 后，
这些谱序列就是 Homology，第 \ref{homology-section-double-complex} 节
讨论的、用于计算 $\text{Tot}(K^\bullet \otimes L^\bullet)$
上同调的两个谱序列。
\end{example}









\section{乘积与 Tor}
\label{more-algebra-section-products-tor}

\noindent
乘积映射最简单的例子来自下列情形。假设
$K^\bullet, L^\bullet \in D(R)$。那么存在映射
\begin{equation}
\label{more-algebra-equation-simple-tor-product}
H^i(K^\bullet) \otimes_R H^j(L^\bullet)
\longrightarrow
H^{i + j}(K^\bullet \otimes_R^{\mathbf{L}} L^\bullet)
\end{equation}
为定义这些映射，可以假设 $K^\bullet, L^\bullet$ 之一是
$R$-模的 K-平坦复形（例如由自由或投射 $R$-模组成的有上界复形）。
在这种情况下，$K^\bullet \otimes_R^{\mathbf{L}} L^\bullet$
由复形 $\text{Tot}(K^\bullet \otimes_R L^\bullet)$ 表示；见第
\ref{more-algebra-section-derived-tensor-product} 节（或第
\ref{more-algebra-section-computing-tor} 节）。接着，假设
$\xi \in H^i(K^\bullet)$ 且 $\zeta \in H^j(L^\bullet)$。
选取代表 $\xi$ 和 $\zeta$ 的
$k \in \Ker(K^i \to K^{i + 1})$ 和
$l \in \Ker(L^j \to L^{j + 1})$。定义
$$
\xi \cup \zeta =
\text{下列对象的类： }k \otimes l\text{ 于 }
H^{i + j}(\text{Tot}(K^\bullet \otimes_R L^\bullet)).
$$
该定义有意义，因为全复形上的微分 $\text{d}$ 的公式（见 Homology，
定义 \ref{homology-definition-associated-simple-complex}）表明
$k \otimes l$ 是余闭元。此外，如果对某个
$k'' \in K^{i - 1}$ 有 $k' = d_K(k'')$，则由于 $l$ 是余闭元，
$k' \otimes l = \text{d}(k'' \otimes l)$。类似地，改变代表
$\zeta$ 的 $l$ 的选取不会改变 $k \otimes l$ 的类。
同样显然，$\cup$ 是双线性的，因而把
$H^i(K^\bullet) \otimes_R H^j(L^\bullet)$ 的一般元素按
$$
\sum \xi_i \otimes \zeta_i \longmapsto \sum \xi_i \cup \zeta_i
$$
映到 $H^{i + j}(\text{Tot}(K^\bullet \otimes_R L^\bullet))$ 中。

\medskip\noindent
设 $R \to A$ 为环同态。设 $K^\bullet, L^\bullet \in D(R)$。
那么在 $D(A)$ 中存在典范等同
\begin{equation}
\label{more-algebra-equation-pullback-derived-tensor-product}
(K^\bullet \otimes_R^{\mathbf{L}} A)
\otimes_A^{\mathbf{L}}
(L^\bullet \otimes_R^{\mathbf{L}} A)
=
(K^\bullet \otimes_R^{\mathbf{L}} L^\bullet) \otimes_R^{\mathbf{L}} A
\end{equation}
其构造如下。首先，在 $R$ 上选取 K-平坦解消
$P^\bullet \to K^\bullet$ 和 $Q^\bullet \to L^\bullet$。
那么左侧由复形
$\text{Tot}((P^\bullet \otimes_R A) \otimes_A (Q^\bullet \otimes_R A))$
表示，右侧由复形
$\text{Tot}(P^\bullet \otimes_R Q^\bullet) \otimes_R A$ 表示。
这些复形典范同构。因此，上述构造诱导有时很有用的乘积
$$
\text{Tor}^R_n(K^\bullet, A) \otimes_A \text{Tor}^R_m(L^\bullet, A)
\longrightarrow
\text{Tor}_{n + m}^R(K^\bullet \otimes_R^\mathbf{L} L^\bullet, A)
$$

\medskip\noindent
设 $M$、$N$ 为 $R$-模。使用上述一般构造、典范映射
$M \otimes_R^\mathbf{L} N \to M \otimes_R N$ 以及
$\text{Tor}$ 的函子性，得到典范映射
\begin{equation}
\label{more-algebra-equation-tor-product}
\text{Tor}^R_n(M, A) \otimes_A \text{Tor}^R_m(N, A)
\longrightarrow \text{Tor}_{n + m}^R(M \otimes_R N, A)
\end{equation}
下面给出使用投射解消的直接构造。首先，在 $R$ 上选取投射解消
$$
P_\bullet \to M, \quad Q_\bullet \to N, \quad T_\bullet \to M \otimes_R N
$$
由 $\otimes_R$ 的右正合性，有
$H_0(\text{Tot}(P_\bullet \otimes_R Q_\bullet)) = M \otimes_R N$。
因此，Derived Categories，引理
\ref{derived-lemma-morphisms-lift-projective} 和
\ref{derived-lemma-morphisms-equal-up-to-homotopy-projective}
保证存在唯一的复形映射
$\mu : \text{Tot}(P_\bullet \otimes_R Q_\bullet) \to T_\bullet$，使得
$H_0(\mu) = \text{id}_{M \otimes_R N}$。这诱导 $D(A)$ 中的典范映射
\begin{align*}
(M \otimes_R^{\mathbf{L}} A) \otimes_A^{\mathbf{L}}
(N \otimes_R^{\mathbf{L}} A)
& =
\text{Tot}((P_\bullet \otimes_R A) \otimes_A (Q_\bullet \otimes_R A)) \\
& =
\text{Tot}(P_\bullet \otimes_R Q_\bullet) \otimes_R A \\
& \to
T_\bullet \otimes_R A \\
& = (M \otimes_R N) \otimes_R^{\mathbf{L}} A
\end{align*}
因此，为构造上面的乘积（\ref{more-algebra-equation-tor-product}），先在 $A$ 上使用
（\ref{more-algebra-equation-simple-tor-product}）构造
$$
\text{Tor}^R_n(M, A) \otimes_A \text{Tor}^R_m(N, A) \to
H^{-n-m}((M \otimes_R^{\mathbf{L}} A) \otimes_A^{\mathbf{L}}
(N \otimes_R^{\mathbf{L}} A))
$$
再与上面的显示映射复合，最终落入
$\text{Tor}_{n + m}^R(M \otimes_R N, A)$。

\medskip\noindent
上述构造的一个有趣特例是 $M = N = B$，其中 $B$ 是 $R$-代数。
在这种情况下，得到映射
$$
\text{Tor}_n^R(B, A) \otimes_A \text{Tor}_m^R(B, A)
\longrightarrow
\text{Tor}_{n + m}^R(B \otimes_R B, A)
\longrightarrow
\text{Tor}_{n + m}^R(B, A)
$$
第二条箭头由乘法映射 $B \otimes_R B \to B$ 通过
$\text{Tor}$ 的函子性诱导。换言之，在
$\text{Tor}^R_{\star}(B, A)$ 上得到一个 $A$-代数结构。
这个代数结构具有许多有趣的性质（结合性、分次交换性、
$B$-代数结构、在某些情形下的除幂等），我们将在别处讨论这些性质
（在此插入将来的引用）。

\begin{lemma}
\label{more-algebra-lemma-functoriality-product-tor}
设 $R$ 为环。设 $A, B, C$ 为 $R$-代数，并设 $B \to C$ 为
$R$-代数同态。那么诱导映射
$$
\text{Tor}^R_{\star}(B, A)
\longrightarrow
\text{Tor}^R_{\star}(C, A)
$$
是 $A$-代数同态。
\end{lemma}

\begin{proof}
略。提示：把所有复形显式写出并逐一定义即可证明。
\end{proof}








\section{K"unneth 谱序列}
\label{more-algebra-section-kunneth-spectral-sequence}

\noindent
设 $R$ 为环。设 $K^\bullet$ 和 $L^\bullet$ 为 $R$-模的滤过复形
（见 Homology，定义 \ref{homology-definition-filtered-complex}；
注意这里的滤过递减）。那么复形
$$
T^\bullet = \text{Tot}(K^\bullet \otimes_R L^\bullet)
$$
也有由公式
$$
F^nT^\bullet =
\Im\left(
\bigoplus\nolimits_{i + j = n}
\text{Tot}(F^iK^\bullet \otimes_R F^jL^\bullet) \to
\text{Tot}(K^\bullet \otimes_R L^\bullet)
\right)
$$
定义的递减滤过。在对这些滤过复形作出一些假设后，我们将确定由
Homology，第 \ref{homology-section-filtered-complex} 节的构造所得的谱序列。

\medskip\noindent
假设每个 $K^n$、$F^iK^n$、$\text{gr}^iK^n$、$L^m$、
$F^jL^m$ 和 $\text{gr}^jL^m$ 都是平坦 $R$-模。在这种情况下，模
$F^iK^n \otimes_R L^m$、$K^n \otimes_R F^jL^m$ 和
$F^iK^n \otimes_R F^jL^m$ 都是 $K^n \otimes_R L^m$ 的子模。
类似地，模 $\text{gr}^iK^n \otimes_R \text{gr}^jL^m$ 是
$K^n/F^{i + 1}K^n \otimes_R L^m/F^{j + 1}L^m$ 的子模。考虑映射
$$
F^nT^n \longrightarrow
\bigoplus\nolimits_{i + j = n}
\text{Tot}(K^\bullet / F^{i + 1}K^\bullet
\otimes_R L^\bullet / F^{j + 1}L^\bullet)
$$
由 $F^nT^\bullet$ 的定义，最终所得确为直和而非直积；这一点留给读者证明。
对于 $a + b = n$，上述箭头在 $F^nT^\bullet$ 的子复形
$\text{Tot}(F^aK^\bullet \otimes_R F^bL^\bullet)$ 上的限制映入
$i = a$ 且 $j = b$ 的直和项。此外，由平坦性假设，该像同构于
$\text{Tot}(\text{gr}^iK^\bullet \otimes_R \text{gr}^jL^\bullet)$，
而映射
$\text{Tot}(F^iK^\bullet \otimes_R F^jL^\bullet) \to
\text{Tot}(\text{gr}^iK^\bullet \otimes_R \text{gr}^jL^\bullet)$ 的核是
$\text{Tot}(F^{i + 1}K^\bullet \otimes_R F^jL^\bullet) +
\text{Tot}(F^iK^\bullet \otimes_R F^{j + 1}L^\bullet)$。
因此，在这些假设下有复形的短正合序列
$$
0 \to F^{n + 1}T^\bullet \to F^nT^\bullet \to
\bigoplus\nolimits_{i + j = n}
\text{Tot}(\text{gr}^iK^\bullet \otimes_R \text{gr}^jL^\bullet) \to 0
$$
换言之，这说明 $\text{gr}^nT^\bullet$ 是复形
$\text{Tot}(\text{gr}^iK^\bullet \otimes_R \text{gr}^jL^\bullet)$
在 $i + j = n$ 上的直和。

\medskip\noindent
进一步假设 $R$-模复形 $K^\bullet$、$F^iK^\bullet$、
$\text{gr}^iK^\bullet$、$L^\bullet$、$F^jL^\bullet$ 和
$\text{gr}^jL^\bullet$ 都 K-平坦。在这种情况下，Homology，
第 \ref{homology-section-filtered-complex} 节中与滤过复形 $T^\bullet$
相伴的谱序列具有项
$$
E_1^{p, q} =
\bigoplus\nolimits_{i + j = p}
H^{p + q}(\text{gr}^iK^\bullet \otimes_R^\mathbf{L} \text{gr}^jL^\bullet)
$$
其微分由短正合序列
$$
0 \to \text{gr}^{n + 1}T^\bullet \to
F^nT^\bullet/F^{n + 2}T^\bullet \to
\text{gr}^nT^\bullet \to 0
$$
诱导，解释见 Homology，引理
\ref{homology-lemma-spectral-sequence-filtered-complex-d1}。
特别地，读者可以证明 $E_1^{p, q}$ 的直和项
$H^{p + q}(\text{gr}^iK^\bullet \otimes_R^\mathbf{L} \text{gr}^jL^\bullet)$
映入 $E_1^{p + 1, q}$ 中的和
$$
H^{p + q + 1}(\text{gr}^{i + 1}K^\bullet \otimes_R^\mathbf{L}
\text{gr}^jL^\bullet) \oplus
H^{p + q + 1}(\text{gr}^iK^\bullet \otimes_R^\mathbf{L}
\text{gr}^{j + 1}L^\bullet)
$$

\begin{lemma}
\label{more-algebra-lemma-kunneth-converges}
在上述假设下，如果还满足
\begin{enumerate}
\item $K^\bullet$ 上的滤过有限，且 $L^\bullet$ 上的滤过有限；或者
\item 更一般地，下列条件成立：
\begin{enumerate}
\item 当 $i \gg 0$ 时，$F^iK^\bullet$ 无环；
\item 当 $i \ll 0$ 时，$F^iK^\bullet \to K^\bullet$ 是拟同构；
\item 当 $j \gg 0$ 时，$F^jL^\bullet$ 无环；并且
\item 当 $j \ll 0$ 时，$F^jL^\bullet \to L^\bullet$ 是拟同构。
\end{enumerate}
\end{enumerate}
那么该谱序列有界，每个
$H^n(T^\bullet) = H^n(K^\bullet \otimes_R^\mathbf{L} L^\bullet)$
上的伴随滤过有限，并且有收敛
$$
E_1^{p, q} =
\bigoplus\nolimits_{i + j = p}
H^{p + q}(\text{gr}^iK^\bullet \otimes_R^\mathbf{L} \text{gr}^jL^\bullet)
\Rightarrow
H^n(K^\bullet \otimes_R^\mathbf{L} L^\bullet)
$$
\end{lemma}

\begin{proof}
在情形 (1) 中，$T^\bullet$ 上的滤过有限，引理立即由 Homology，
引理 \ref{homology-lemma-biregular-ss-converges} 得到。
在情形 (2) 中，选取 $a < b$，使得当 $i, j > b$ 时
$F^iK^\bullet$ 和 $F^jL^\bullet$ 无环，并且当 $i, j < a$ 时
$F^iK^\bullet \to K^\bullet$ 和 $F^jL^\bullet \to L^\bullet$
是拟同构。我们断言，在这种情况下，当 $n > 2b$ 时复形
$F^nT^\bullet$ 无环，而当 $n < 2a - 1$ 时
$F^nT^\bullet \to T^\bullet$ 是拟同构。该断言表明可以应用
Homology，引理 \ref{homology-lemma-ss-converges-trivial}，故本引理成立。

\medskip\noindent
证明该断言。给定 $n$ 和整数 $i_1 \leq i_2$，考虑
$$
S^{n, i_1, i_2} =
\Im\left(
\bigoplus\nolimits_{i_1 \leq i \leq i_2}
\text{Tot}(F^iK^\bullet \otimes_R F^{n - i}L^\bullet)
\to
\text{Tot}(K^\bullet \otimes_R L^\bullet)
\right)
$$
把它看作 $\text{Tot}(K^\bullet \otimes_R L^\bullet)$ 的子复形。
注意 $F^nT^\bullet = \colim S^{n, i_1, i_2}$ 是滤过余极限。
因此，为证明该断言，只需证明：当 $n > 2b$ 时
$S^{n, i_1, i_2}$ 无环；并且当 $n < 2a - 1$ 时，对一组余终的
二元组 $i_1, i_2$，映射
$S^{n, i_1, i_2} \to \text{Tot}(K^\bullet \otimes L^\bullet)$
是拟同构。

\medskip\noindent
如果 $i_1 = i_2 = i$ 且 $n > 2b$，那么
$$
S^{n, i, i} =
\text{Tot}(F^iK^\bullet \otimes_R F^{n - i}L^\bullet)
$$
无环：如果 $i > b$，则由 $F^{n - i}L^\bullet$ K-平坦的假设，
该复形无环；如果 $i \leq b$，则
$n - i \geq n - b > b$，同理可得结论。使用平坦性假设，
读者可以证明存在复形的短正合序列
$$
0 \to
S^{n + 1, i_2 + 1, i_2 + 1}
\to
S^{n, i_1, i_2} \oplus S^{n, i_2 + 1, i_2 + 1}
\to
S^{n, i_1, i_2 + 1}
\to
0
$$
因此，对 $i_2 - i_1$ 归纳可见，当 $n > 2b$ 时所有复形
$S^{n, i_1, i_2}$ 都无环。

\medskip\noindent
假设 $n < 2a - 1$。注意
$$
S^{n, a - 1, a - 1} =
\text{Tot}(F^{a - 1}K^\bullet \otimes_R F^{n - a + 1}L^\bullet)
\to \text{Tot}(K^\bullet \otimes_R L^\bullet)
$$
是拟同构，因为 $a - 1 < a$ 且 $n - a + 1 < a$，并且所有复形都
K-平坦，所以等号两侧计算 $D(R)$ 的同一个对象。注意当
$i_2 \geq a - 1$ 且 $n < 2a - 1$ 时，映射
$$
S^{n + 1, i_2 + 1, i_2 + 1} \to S^{n, i_2 + 1, i_2 + 1}
$$
是拟同构，因为左侧和右侧都拟同构于
$\text{Tot}(F^{i_2 + 1}K^\bullet \otimes L^\bullet)$。
因此，使用上一段的复形短正合序列可知，对所有
$t \geq a - 1$，映射
$S^{n, a - 1, a - 1} \to S^{n, a - 1, t}$ 是拟同构。
另一方面，类似地存在短正合序列
$$
0 \to
S^{n + 1, i_1, i_1}
\to
S^{n, i_1 - 1, i_1 - 1} \oplus
S^{n, i_1, i_2}
\to
S^{n, i_1 - 1, i_2}
\to
0
$$
同上可知，当 $i_1 \leq a - 1$ 时，映射
$$
S^{n + 1, i_1, i_1} \to S^{n, i_1 - 1, i_1 - 1}
$$
是拟同构，因为左侧和右侧都拟同构于
$\text{Tot}(K^\bullet \otimes F^{n - i_1 + 1}L^\bullet)$。
因此，对所有 $s \leq a - 1$，映射
$S^{n, s, t} \to S^{n, a - 1, t}$ 是拟同构。综合可知，当
$s \leq a - 1$ 且 $t \geq a - 1$ 时，$S^{n, s, t}$
拟同构地映到 $\text{Tot}(K^\bullet \otimes_R L^\bullet)$。
这完成了引理的证明。
\end{proof}

\begin{lemma}
\label{more-algebra-lemma-resolution-filtered-complex}
设 $R$ 为环。设 $K^\bullet$ 为滤过复形。存在滤过复形的映射
$f : P^\bullet \to K^\bullet$，使得
\begin{enumerate}
\item 每个 $P^n$、$F^iP^n$、$\text{gr}^iP^n$ 都是自由 $R$-模；
\item $R$-模复形 $P^\bullet$、$F^iP^\bullet$ 和
$\text{gr}^iP^\bullet$ 都 K-平坦；
\item $f$ 诱导拟同构
$P^\bullet \to K^\bullet$、
$F^iP^\bullet \to F^iK^\bullet$ 和
$\text{gr}^iP^\bullet \to \text{gr}^iK^\bullet$。
\end{enumerate}
\end{lemma}

\begin{proof}
如果滤过复形 $L^\bullet$ 的每个 $L^n$、$F^iL^n$、
$\text{gr}^iL^n$ 都是自由 $R$-模，并且所有微分都为零，
则称该滤过复形 {\it 基本}。

\medskip\noindent
存在基本滤过复形 $P_0^\bullet$ 以及滤过复形的映射
$f_0 : P_0^\bullet \to K^\bullet$，使得 $f_0$ 和
$F^if_0$（$i \in \mathbf{Z}$）在上同调上都是满射。为看出这一点，置
$$
P_0^n = \bigoplus\nolimits_{z \in \Ker(d_K^n)} R \xi_z
\oplus
\bigoplus\nolimits_{j \in \mathbf{Z}}
\bigoplus\nolimits_{z \in \Ker(d_{F^jK}^n)} R \xi_z
$$
取零微分，并用 $f_0(\xi_z) = z$ 定义映射 $f_0$。至于滤过，置
$$
F^iP_0^n = \bigoplus\nolimits_{j \geq i \in \mathbf{Z}}
\bigoplus\nolimits_{z \in \Ker(d_{F^jK}^n)} R \xi_z
$$
所述构造确实给出断言中的 $f_0 : P_0^\bullet \to K^\bullet$；
验证留给读者。

\medskip\noindent
我们将对 $m \geq 0$ 归纳构造滤过复形的嵌入
$$
P_0^\bullet \subset \ldots P_m^\bullet \subset P_{m + 1}^\bullet
$$
以及映射 $f_m : P_m^\bullet \to K^\bullet$，使其具有下列性质：
\begin{enumerate}
\item 滤过复形 $P_{m + 1}^\bullet / P_m^\bullet$ 是基本的；
\item $H^n(f_m)$ 的核与 $H^n(F^if_m)$ 的核在
$H^n(P_{m + 1}^\bullet)$ 和 $H^n(F^iP_{m + 1}^\bullet)$ 中都映为零；
\item 映射 $f_{m + 1} : P_{m + 1}^\bullet \to K^\bullet$
延拓映射 $f_m$。
\end{enumerate}
为此，置
$$
\Omega_{n, m} = \Ker\left(\Ker(d^n_{P_m}) \to H^n(K^\bullet)\right),
\text{ 分别 }
\Omega_{n, i, m} = \Ker\left(\Ker(d^n_{F^iP_m}) \to H^n(F^iK^\bullet)\right),
$$
注意 $\Omega_{n, m}$ 满射到 $H^n(f_m)$ 的核，而
$\Omega_{n, i, m}$ 满射到 $H^n(F^if_n)$ 的核。
对每个 $z \in \Omega_{n, m}$（相应地，$z \in \Omega_{n, i, m}$），
选取 $y_z \in K^{n - 1}$（相应地，$y_z \in F^iK^{n - 1}$），
使得 $d_K(y_z) = f_m(z)$。然后置
$$
P^n_{m + 1} = P_m^n \oplus
\bigoplus\nolimits_{z \in \Omega_{n + 1, m}} R \eta_z
\oplus
\bigoplus\nolimits_{j \in \mathbf{Z}}
\bigoplus\nolimits_{z \in \Omega_{n + 1, j, m}} R \eta_z
$$
置 $d_{P_m}(\eta_z) = z$ 并置 $f_{m + 1}(\eta_z) = y_z$。
最后，置
$$
F^iP_{m + 1}^n = F^iP_m^n \oplus
\bigoplus\nolimits_{j \geq i}
\bigoplus\nolimits_{z \in \Omega_{n + 1, j, m}} R \eta_z
$$
所述构造确实给出断言中的 $P_m^\bullet \subset P_{m + 1}^\bullet$
和 $f_{m + 1} : P_{m + 1}^\bullet \to K^\bullet$；验证留给读者。

\medskip\noindent
此时只需取
$$
P^\bullet = \bigcup P_m^\bullet
$$
作为滤过复形，并取由 $\bigcup f_m$ 给出的映射
$f : P^\bullet \to K^\bullet$。

\medskip\noindent
引理陈述的 (1) 成立，因为作为滤过模的 $P^n$ 同构于
$P_0^n$ 与所有 $m \geq 0$ 的 $P_{m + 1}^n/P_m^n$ 的直和。
略去一个小细节。

\medskip\noindent
考虑陈述的 (3)。注意
$H^n(P^\bullet) = \colim H^n(P_m^\bullet)$。映射 $f_0$
在上同调上满射；因此每个 $f_m$ 也如此。根据构造，对每个 $m$，
嵌入 $P_m^\bullet \subset P_{m + 1}^\bullet$ 杀掉
$H^n(f_m)$ 的核。综合这些事实，读者容易推出 $H^n(f)$ 是同构。
$H^n(F^if_n)$ 的情形类似。于是 $H^n(\text{gr}^if)$ 也必为同构，
因为有短正合序列 $0 \to F^{i + 1} \to F^i \to \text{gr}^i \to 0$
（例如，把它看作滤过复形范畴上的函子序列）。略去一个小细节。

\medskip\noindent
考虑陈述的 (2)。为看出 $P^\bullet$ K-平坦，由引理
\ref{more-algebra-lemma-colimit-K-flat}，只需证明 $P_m^\bullet$ K-平坦。
由引理 \ref{more-algebra-lemma-K-flat-two-out-of-three-ses} 和归纳，只需注意微分为零、
各项自由的复形 K-平坦。相同论证说明，对所有
$i \in \mathbf{Z}$，$F^iP^\bullet$ 都 K-平坦。最后，再次应用引理
\ref{more-algebra-lemma-K-flat-two-out-of-three-ses} 可知
$\text{gr}^iP^\bullet$ K-平坦。
\end{proof}

\begin{proposition}
\label{more-algebra-proposition-kunneth-filtered}
设 $R$ 为环。设 $K^\bullet$ 和 $L^\bullet$ 为 $R$-模的滤过复形。
那么存在表示 $D(R)$ 中 $K^\bullet \otimes_R^\mathbf{L} L^\bullet$
的滤过复形 $T^\bullet$，使得相伴谱序列的 $E_1$-页为
$$
E_1^{p, q} =
\bigoplus\nolimits_{i + j = p}
H^{p + q}(\text{gr}^iK^\bullet \otimes_R^\mathbf{L} \text{gr}^jL^\bullet)
$$
如果
\begin{enumerate}
\item 当 $i \gg 0$ 时，$F^iK^\bullet$ 无环；
\item 当 $i \ll 0$ 时，$F^iK^\bullet \to K^\bullet$ 是拟同构；
\item 当 $j \gg 0$ 时，$F^jL^\bullet$ 无环；并且
\item 当 $j \ll 0$ 时，$F^jL^\bullet \to L^\bullet$ 是拟同构；
\end{enumerate}
则该谱序列有界，每个
$H^n(K^\bullet \otimes_R^\mathbf{L} L^\bullet)$ 上的伴随滤过有限，
并且该谱序列收敛。
\end{proposition}

\begin{proof}
选取引理 \ref{more-algebra-lemma-resolution-filtered-complex} 中的
$P^\bullet \to K^\bullet$ 和 $Q^\bullet \to L^\bullet$。
然后使用本节正文和引理 \ref{more-algebra-lemma-kunneth-converges} 所描述的、
滤过复形
$$
T^\bullet = \text{Tot}(P^\bullet \otimes_R Q^\bullet)
$$
的谱序列。
\end{proof}

\begin{lemma}[K"unneth 谱序列]
\label{more-algebra-lemma-kunneth-spectral-sequence}
设 $R$ 为环。设 $K$ 和 $L$ 为 $D^b(R)$ 的对象。
存在双分次有界谱序列 $\{E_r\}_{r \geq 2}$，其中
$$
E_2^{p, q} = \bigoplus\nolimits_{i + j = q} \text{Tor}^R_{-p}(H^i(K), H^j(L))
$$
且 $d_r$ 的双次数为 $(r, -r + 1)$，该谱序列收敛到
$H^{p + q}(K \otimes_R^\mathbf{L} L)$。
\end{lemma}

\begin{proof}
设 $K^\bullet$ 为表示 $K$ 的 $R$-模复形，设 $L^\bullet$
为表示 $L$ 的 $R$-模复形。置
$F^iK^\bullet = \tau_{\leq -i}K^\bullet$，并对 $L$ 作类似定义；
见 Homology，第 \ref{homology-section-truncations} 节。
应用命题 \ref{more-algebra-proposition-kunneth-filtered}，并注意
$\text{gr}^iK^\bullet = H^{-i}(K)[i]$，得到有界谱序列
$\{(E')_r\}_{r \geq 1}$，其
$$
(E')^{p, q}_1 =
\bigoplus\nolimits_{i + j = p} \text{Tor}^R_{-2p - q}(H^{-i}(K), H^{-j}(L))
$$
收敛到 $K \otimes_R^\mathbf{L} L$ 的上同调。按构造，该谱序列具有微分
$d_r^{p, q} : (E')_r^{p, q} \to (E')_r^{p + r, q - r + 1}$，
其中 $r \geq 1$。为得到引理中的谱序列，对 $r \geq 2$ 置
$$
E_r^{p, q} = (E')_{r - 1}^{-q, p + 2q}
$$
所得构造确实可行；验证留给读者。
\end{proof}

\begin{example}
\label{more-algebra-example-kunneth-dedekind}
如果 $R = \mathbf{Z}$，或更一般地，如果 $R$ 是 Dedekind 整环，
那么对所有 $R$-模 $M$ 和 $N$，当 $i \not \in \{0, 1\}$ 时都有
$\text{Tor}_i^R(M, N) = 0$。因此，引理
\ref{more-algebra-lemma-kunneth-spectral-sequence} 的谱序列在 $E_2$ 页退化，
并且对所有 $n \in \mathbf{Z}$ 都得到短正合序列
$$
0 \to \bigoplus_{i + j = n} H^i(K) \otimes_R H^j(L) \to
H^n(K \otimes_R^\mathbf{L} L) \to
\bigoplus_{i + j = n + 1} \text{Tor}^R_1(H^i(K), H^j(L)) \to 0
$$
\end{example}






\section{拟相干模，I}
\label{more-algebra-section-pseudo-coherent}

\noindent
设 $R$ 为环。回顾一下，如果存在满射 $R^{\oplus a} \to M$，
则称 $R$-模 $M$ 为有限型；如果存在表示
$R^{\oplus a_1} \to R^{\oplus a_0} \to M \to 0$，
则称其为有限表示。类似地，可以考虑存在长度为 $n$ 的解消的
$R$-模：
\begin{equation}
\label{more-algebra-equation-pseudo-coherent}
R^{\oplus a_n} \to R^{\oplus a_{n - 1}} \to \ldots \to R^{\oplus a_0} \to
M \to 0
\end{equation}
其中各项是有限自由 $R$-模。如果对每个 $n$ 都能找到这样的解消，
则称模是{\it 拟相干的}。下面给出形式定义。

\begin{definition}
\label{more-algebra-definition-pseudo-coherent}
设 $R$ 为环。记其导出范畴为 $D(R)$。设 $m \in \mathbf{Z}$。
\begin{enumerate}
\item 如果存在有限自由 $R$-模组成的有界复形 $E^\bullet$ 及映射
$\alpha : E^\bullet \to K^\bullet$，使得当 $i > m$ 时
$H^i(\alpha)$ 是同构、而 $H^m(\alpha)$ 是满射，则称 $D(R)$ 的对象
$K^\bullet$ 为{\it $m$-拟相干的}。
\item 如果对象 $K^\bullet$ 是 $D(R)$ 的对象且拟同构于有限自由 $R$-模组成的上有界复形，
则称其为{\it 拟相干的}。
\item 如果 $M[0]$ 是 $D(R)$ 中的 $m$-拟相干对象，则称 $R$-模 $M$
为{\it $m$-拟相干的}。
\item 如果 $M[0]$ 是 $D(R)$ 中的拟相干对象，则称 $R$-模 $M$ 为
{\it 拟相干的}\footnote{这与 \cite{Bourbaki-CA} 中对拟相干模的含义冲突。}。
\end{enumerate}
\end{definition}

\noindent
通常，我们也把这一术语用于 $R$-模复形。由于 $D(R)$ 中的任意映射
$E^\bullet \to K^\bullet$ 都由复形的实际映射表示（见 Derived Categories，
引理 \ref{derived-lemma-morphisms-from-projective-complex}），
所以没有歧义。
事实证明，$K^\bullet$ 拟相干当且仅当对所有 $m \in \mathbf{Z}$，
$K^\bullet$ 都是 $m$-拟相干的；见引理 \ref{more-algebra-lemma-pseudo-coherent}。
此外，如果环是 Noether 环，则该条件可以理解为上同调的有限生成条件；
见引理 \ref{more-algebra-lemma-Noetherian-pseudo-coherent}。先把它与上面的非正式讨论联系起来。

\begin{lemma}
\label{more-algebra-lemma-cone-pseudo-coherent}
设 $R$ 为环且 $m \in \mathbf{Z}$。设
$(K^\bullet, L^\bullet, M^\bullet, f, g, h)$ 为 $D(R)$ 中的可区别三角。
\begin{enumerate}
\item 如果 $K^\bullet$ 是 $(m + 1)$-拟相干的且 $L^\bullet$ 是
$m$-拟相干的，则 $M^\bullet$ 是 $m$-拟相干的。
\item 如果 $K^\bullet, M^\bullet$ 都是 $m$-拟相干的，则
$L^\bullet$ 是 $m$-拟相干的。
\item 如果 $L^\bullet$ 是 $(m + 1)$-拟相干的且 $M^\bullet$ 是
$m$-拟相干的，则 $K^\bullet$ 是 $(m + 1)$-拟相干的。
\end{enumerate}
\end{lemma}

\begin{proof}
证明 (1)。选取 $\alpha : P^\bullet \to K^\bullet$，其中
$P^\bullet$ 是有限自由模组成的有界复形，并且当 $i > m + 1$ 时
$H^i(\alpha)$ 是同构、在 $i = m + 1$ 时是满射。可以用
$\sigma_{\geq m + 1}P^\bullet$ 替代 $P^\bullet$，因而不妨设
$P^i = 0$（当 $i < m + 1$）。选取 $\beta : E^\bullet \to L^\bullet$，
其中 $E^\bullet$ 是有限自由模组成的有界复形，并且当 $i > m$ 时
$H^i(\beta)$ 是同构、在 $i = m$ 时是满射。由 Derived Categories，
引理 \ref{derived-lemma-lift-map}，可以找到映射
$\gamma : P^\bullet \to E^\bullet$，使图
$$
\xymatrix{
K^\bullet \ar[r] & L^\bullet \\
P^\bullet \ar[u] \ar[r]^\gamma & E^\bullet \ar[u]_\beta
}
$$
在 $D(R)$ 中交换。锥 $C(\gamma)^\bullet$ 是有限自由 $R$-模组成的有界复形，
且图的交换性意味着存在可区别三角的态射
$$
(P^\bullet, E^\bullet, C(\gamma)^\bullet)
\longrightarrow
(K^\bullet, L^\bullet, M^\bullet).
$$
由诱导的上同调长正合序列以及 Homology，引理
\ref{homology-lemma-four-lemma} 和 \ref{homology-lemma-five-lemma}，
可知 $C(\gamma)^\bullet \to M^\bullet$ 在次数 $> m$ 的上同调上诱导同构，
在次数 $m$ 上诱导满射。因此 $M^\bullet$ 是 $m$-拟相干的。

\medskip\noindent
(2) 和 (3) 由 (1) 旋转可区别三角得到。
\end{proof}

\begin{lemma}
\label{more-algebra-lemma-finite-cohomology}
设 $R$ 为环。设 $K^\bullet$ 为 $R$-模复形。设 $m \in \mathbf{Z}$。
\begin{enumerate}
\item 如果 $K^\bullet$ 是 $m$-拟相干的且当 $i > m$ 时
$H^i(K^\bullet) = 0$，则 $H^m(K^\bullet)$ 是有限型 $R$-模。
\item 如果 $K^\bullet$ 是 $m$-拟相干的且当 $i > m + 1$ 时
$H^i(K^\bullet) = 0$，则 $H^{m + 1}(K^\bullet)$ 是有限表示的
$R$-模。
\end{enumerate}
\end{lemma}

\begin{proof}
证明 (1)。选取有限投射 $R$-模组成的有界复形 $E^\bullet$ 以及映射
$\alpha : E^\bullet \to K^\bullet$，使其在次数 $> m$ 的上同调上诱导同构、
在次数 $m$ 上诱导满射。显然只需对 $E^\bullet$ 证明结论。
设 $n$ 为使 $E^n \not = 0$ 的最大整数。如果 $n = m$，结论显然。
如果 $n > m$，则因 $H^n(E^\bullet) = 0$，映射
$E^{n - 1} \to E^n$ 是满射。由于 $E^n$ 是有限投射的，可知
$E^{n - 1} = E' \oplus E^n$。因此只需对复形 $(E')^\bullet$ 证明，
它与 $E^\bullet$ 相同，只是第 $n - 1$ 次为 $E'$、第 $n$ 次为 $0$。
对 $n$ 归纳即可。

\medskip\noindent
证明 (2)。选取有限投射 $R$-模组成的有界复形 $E^\bullet$ 以及映射
$\alpha : E^\bullet \to K^\bullet$，使其在次数 $> m$ 的上同调上诱导同构、
在次数 $m$ 上诱导满射。与 (1) 的证明一样，可以化归到
$E^i = 0$（当 $i > m + 1$）的情形。于是
$$
H^{m + 1}(K^\bullet) \cong
H^{m + 1}(E^\bullet) = \Coker(E^m \to E^{m + 1})
$$
是有限表示的。
\end{proof}

\begin{lemma}
\label{more-algebra-lemma-n-pseudo-module}
设 $R$ 为环。设 $M$ 为 $R$-模。则
\begin{enumerate}
\item $M$ 是 $0$-拟相干的，当且仅当 $M$ 是有限 $R$-模；
\item $M$ 是 $(-1)$-拟相干的，当且仅当 $M$ 是有限表示的 $R$-模；
\item $M$ 是 $(-d)$-拟相干的，当且仅当存在长度为 $d$ 的解消
$$
R^{\oplus a_d} \to R^{\oplus a_{d - 1}} \to \ldots \to R^{\oplus a_0} \to
M \to 0
$$
并且
\item $M$ 是拟相干的，当且仅当存在由有限自由 $R$-模组成的无限解消
$$
\ldots \to R^{\oplus a_1} \to R^{\oplus a_0} \to M \to 0
$$
\end{enumerate}
\end{lemma}

\begin{proof}
如果 $M$ 是有限型（分别为有限表示），则由定义
\ref{more-algebra-definition-pseudo-coherent} 之前的讨论，$M$ 是 $0$-拟相干
（分别为 $(-1)$-拟相干）的。反之，如果 $M$ 是 $0$-拟相干的，
则 $M = H^0(M[0])$ 由引理 \ref{more-algebra-lemma-finite-cohomology} 为有限型。
如果 $M$ 是 $(-1)$-拟相干的，则它也是 $0$-拟相干的，因而为有限型。
选取满射 $R^{\oplus a} \to M$，记
$K = \Ker(R^{\oplus a} \to M)$。由引理 \ref{more-algebra-lemma-cone-pseudo-coherent}，
$K$ 是 $0$-拟相干的，因而为有限型，所以 $M$ 是有限表示的。

\medskip\noindent
第三、第四个断言用归纳和类似上述的论证即可证明（细节略）。
\end{proof}

\begin{lemma}
\label{more-algebra-lemma-pseudo-coherent}
设 $R$ 为环。设 $K^\bullet$ 为 $R$-模复形。下列条件等价：
\begin{enumerate}
\item $K^\bullet$ 是拟相干的；
\item 对每个 $m \in \mathbf{Z}$，$K^\bullet$ 都是 $m$-拟相干的；以及
\item $K^\bullet$ 拟同构于有限投射 $R$-模组成的上有界复形。
\end{enumerate}
如果 (1)、(2)、(3) 成立且当 $i > b$ 时
$H^i(K^\bullet) = 0$，则可以找到拟同构 $F^\bullet \to K^\bullet$，
其中 $F^i$ 是有限自由 $R$-模且当 $i > b$ 时 $F^i = 0$。
\end{lemma}

\begin{proof}
由于有限自由模是有限投射 $R$-模，故 (1) $\Rightarrow$ (3)。
反之，设 $P^\bullet$ 是有限投射 $R$-模组成的上有界复形。
设当 $i > n_0$ 时 $P^i = 0$。选取直和分解
$F^{n_0} = P^{n_0} \oplus C^{n_0}$，其中 $F^{n_0}$ 是有限自由
$R$-模，并归纳地对 $n \leq n_0$ 取
$$
F^{n - 1} = P^{n - 1} \oplus C^n \oplus C^{n - 1}
$$
其中 $F^{n_0}$ 是有限自由 $R$-模。令复形 $F^\bullet$ 的映射
$F^{n - 1} \to F^n$ 与 $P^{n - 1} \to P^n$ 一致，在
$C^n \to C^n$ 上诱导恒等映射，并且在 $C^{n - 1}$ 上为零。
映射 $F^\bullet \to P^\bullet$ 是拟同构（事实上是同伦等价），
所以 (3) 推出 (1)。

\medskip\noindent
假设 (1)。设 $E^\bullet$ 是有限自由 $R$-模组成的上有界复形，
且 $E^\bullet \to K^\bullet$ 是拟同构。由愚蠢截断得到的
诱导的映射 $\sigma_{\geq m}E^\bullet \to K^\bullet$ 由 $E^\bullet$ 到 $K^\bullet$ 的愚蠢截断得到，
可知 $K^\bullet$ 是 $m$-拟相干的。因此 (1) 推出 (2)。

\medskip\noindent
假设 (2)。由于 $K^\bullet$ 是 $0$-拟相干的，特别可知 $K^\bullet$ 上有界。
取整数 $b$ 使得当 $i > b$ 时 $H^i(K^\bullet) = 0$。我们将对
$n \in \mathbf{Z}$ 作下降归纳，构造当 $i \geq n$ 时的有限自由
$R$-模 $F^i$、微分 $d^i : F^i \to F^{i + 1}$（当 $i \geq n$）以及与微分相容的
映射 $\alpha : F^i \to K^i$，使得 (1) $H^i(\alpha)$ 在 $i > n$
时是同构、在 $i = n$ 时是满射；(2) 当 $i > b$ 时 $F^i = 0$。
图示为
$$
\xymatrix{
& F^n \ar[r] \ar[d]^\alpha & F^{n + 1} \ar[d]^\alpha \ar[r] & \ldots \\
K^{n - 1} \ar[r] & K^n \ar[r] & K^{n + 1} \ar[r] & \ldots
}
$$
基例为 $n = b + 1$，此时对所有 $i$ 可取 $F^i = 0$。

归纳步骤。令 $C^\bullet$ 为 $\alpha$ 的锥（Derived Categories，定义
\ref{derived-definition-cone}）。上同调长正合序列
$$
0 \to H^{n - 1}(K^\bullet) \to H^{n - 1}(C^\bullet) \to
H^n(F^\bullet) \to H^n(K^\bullet) \to H^n(C^\bullet) \to \ldots
$$
表明当 $i \geq n$ 时 $H^i(C^\bullet) = 0$。由引理
\ref{more-algebra-lemma-cone-pseudo-coherent} 可知 $C^\bullet$ 是
$(n - 1)$-拟相干的。由引理 \ref{more-algebra-lemma-finite-cohomology} 可知
$H^{n - 1}(C^\bullet)$ 是有限 $R$-模。
特别地，$H^n(F^\bullet) \to H^n(K^\bullet)$ 的核是有限 $R$-模。
选取有限自由 $R$-模 $F^{n - 1}$ 以及映射
$F^{n - 1} \to \Ker(F^n \to F^{n - 1})$，使 $F^{n - 1}$ 满射到
$H^n(F^\bullet) \to H^n(K^\bullet)$ 的核。
把复形延拓为
$$
\xymatrix{
& F^{n - 1} \ar[d]^{\alpha^{n - 1}} \ar[r] &
F^n \ar[r] \ar[d]^\alpha &
F^{n + 1} \ar[d]^\alpha \ar[r] & \ldots \\
\ldots \ar[r] &
K^{n - 1} \ar[r] & K^n \ar[r] & K^{n + 1} \ar[r] & \ldots
}
$$
注意 $\alpha^{n - 1}$ 存在，是因为由构造可知
$F^{n - 1} \to F^n \to K^n$ 的像包含于 $K^{n - 1} \to K^n$ 的像中。
再次记 $C^\bullet$ 为这个延拓的复形映射的锥。此时得到正合序列
$$
H^{n - 1}(F^\bullet) \to H^{n - 1}(K^\bullet) \to H^{n - 1}(C^\bullet) \to
H^n(F^\bullet) \cong H^n(K^\bullet) \to H^n(C^\bullet) \to \ldots
$$
换言之，$H^{n - 1}(F^\bullet) \to H^{n - 1}(K^\bullet)$ 的余核是有限
$R$-模，设其由 $\xi_1, \ldots, \xi_r \in \Ker(K^{n - 1} \to K^n)$
的类生成。然后把 $F^{n - 1}$ 替换为
$F^{n - 1} \oplus R^{\oplus r}$，使自由直和中的基元素映到 $F^n$ 中的零，
并映到 $K^{n - 1}$ 中的 $\xi_i$。这完成归纳步骤的证明。
\end{proof}

\begin{lemma}
\label{more-algebra-lemma-two-out-of-three-pseudo-coherent}
设 $R$ 为环。设 $(K^\bullet, L^\bullet, M^\bullet, f, g, h)$ 为
$D(R)$ 中的可区别三角。如果 $K^\bullet, L^\bullet, M^\bullet$ 中任意
两个拟相干，则第三个也拟相干。
\end{lemma}

\begin{proof}
结合引理 \ref{more-algebra-lemma-cone-pseudo-coherent} 和
\ref{more-algebra-lemma-pseudo-coherent} 即得。
\end{proof}

\begin{lemma}
\label{more-algebra-lemma-recognize-pseudo-coherent}
设 $R$ 为环。设 $K^\bullet$ 为 $R$-模复形。设 $m \in \mathbf{Z}$。
\begin{enumerate}
\item 如果对所有 $i \geq m$ 都有 $H^i(K^\bullet) = 0$，则
$K^\bullet$ 是 $m$-拟相干的。
\item 如果当 $i > m$ 时 $H^i(K^\bullet) = 0$，且 $H^m(K^\bullet)$
是有限 $R$-模，则 $K^\bullet$ 是 $m$-拟相干的。
\item 如果当 $i > m + 1$ 时 $H^i(K^\bullet) = 0$，且模
$H^{m + 1}(K^\bullet)$ 是有限表示的、$H^m(K^\bullet)$ 是有限型的，
则 $K^\bullet$ 是 $m$-拟相干的。
\end{enumerate}
\end{lemma}

\begin{proof}
只需证明 (3)。置 $M = H^{m + 1}(K^\bullet)$。注意
$\tau_{\geq m + 1}K^\bullet$ 拟同构于 $M[- m - 1]$。由引理
\ref{more-algebra-lemma-n-pseudo-module}，$M[- m - 1]$ 是 $m$-拟相干的。
由于有可区别三角
$$
(\tau_{\leq m}K^\bullet, K^\bullet, \tau_{\geq m + 1}K^\bullet)
$$
（Derived Categories，评注
\ref{derived-remark-truncation-distinguished-triangle}），由引理
\ref{more-algebra-lemma-cone-pseudo-coherent}，只需证明
$\tau_{\leq m}K^\bullet$ 拟相干。根据假设，
$H^m(\tau_{\leq m}K^\bullet)$ 是有限型 $R$-模。
因此可以找到有限自由 $R$-模 $E$ 及映射
$E \to \Ker(d_K^m)$，使复合
$E \to \Ker(d_K^m) \to H^m(\tau_{\leq m}K^\bullet)$ 满射。
于是 $E[-m] \to \tau_{\leq m}K^\bullet$ 见证
$\tau_{\leq m}K^\bullet$ 是 $m$-拟相干的事实。
\end{proof}

\begin{lemma}
\label{more-algebra-lemma-summands-pseudo-coherent}
设 $R$ 为环。设 $m \in \mathbf{Z}$。如果 $K^\bullet \oplus L^\bullet$
是 $m$-拟相干的（分别为拟相干的），则 $K^\bullet$ 和 $L^\bullet$
也是如此。
\end{lemma}

\begin{proof}
在本证明中省略上标 ${}^\bullet$。假设 $K \oplus L$ 是
$m$-拟相干的。显然 $K, L \in D^{-}(R)$。注意有可区别三角
$$
(K \oplus L, K \oplus L, L \oplus L[1]) =
(K, K, 0) \oplus (L, L, L \oplus L[1])
$$
见 Derived Categories，引理 \ref{derived-lemma-direct-sum-triangles}。
由引理 \ref{more-algebra-lemma-cone-pseudo-coherent}，$L \oplus L[1]$ 是
$m$-拟相干的。因此 $L[1] \oplus L[2]$ 也是 $m$-拟相干的。
归纳可知 $L[n] \oplus L[n + 1]$ 是 $m$-拟相干的。
由引理 \ref{more-algebra-lemma-recognize-pseudo-coherent}，当 $n$ 足够大时
$L[n]$ 是 $m$-拟相干的。然后利用可区别三角
$$
(L[n], L[n] \oplus L[n - 1], L[n - 1])
$$
向后工作，得到 $L[n], L[n - 1], \ldots, L$ 都是 $m$-拟相干的。
拟相干的情形由此及引理 \ref{more-algebra-lemma-pseudo-coherent} 得到。
\end{proof}

\begin{lemma}
\label{more-algebra-lemma-complex-pseudo-coherent-modules}
设 $R$ 为环。设 $m \in \mathbf{Z}$。设 $K^\bullet$ 为上有界的
$R$-模复形，且对所有 $i$，$K^i$ 是 $(m - i)$-拟相干的。
则 $K^\bullet$ 是 $m$-拟相干的。特别地，如果 $K^\bullet$ 是拟相干
$R$-模组成的上有界复形，则 $K^\bullet$ 拟相干。
\end{lemma}

\begin{proof}
可以用 $\sigma_{\geq m - 1}K^\bullet$ 替代 $K^\bullet$（例如），
因而不妨设 $K^\bullet$ 有界。由复形长度归纳，使用引理
\ref{more-algebra-lemma-cone-pseudo-coherent} 和愚蠢截断，复形 $K^\bullet$ 是
$m$-拟相干的，因为每个 $K^i[-i]$ 都是 $m$-拟相干的。
对于最后一个断言，只需证明对所有 $m \in \mathbf{Z}$，
$K^\bullet$ 都是 $m$-拟相干的；见引理 \ref{more-algebra-lemma-pseudo-coherent}。
这由第一部分得到。
\end{proof}

\begin{lemma}
\label{more-algebra-lemma-cohomology-pseudo-coherent}
设 $R$ 为环。设 $m \in \mathbf{Z}$。设 $K^\bullet \in D^{-}(R)$，使得对所有 $i$，
$H^i(K^\bullet)$ 都是 $(m - i)$-拟相干（分别为拟相干）的。
则 $K^\bullet$ 是 $m$-拟相干（分别为拟相干）的。
\end{lemma}

\begin{proof}
设 $K^\bullet$ 是 $D^{-}(R)$ 的对象，并且每个 $H^i(K^\bullet)$ 都是
$(m - i)$-拟相干的。令 $n$ 为使 $H^n(K^\bullet)$ 非零的最大整数。
我们对 $n$ 归纳证明引理。
如果 $n < m$，则由引理
\ref{more-algebra-lemma-recognize-pseudo-coherent}，$K^\bullet$ 是 $m$-拟相干的。
如果 $n \geq m$，则有可区别三角
$$
(\tau_{\leq n - 1}K^\bullet, K^\bullet, H^n(K^\bullet)[-n])
$$
（Derived Categories，评注
\ref{derived-remark-truncation-distinguished-triangle}）。由于按假设
$H^n(K^\bullet)[-n]$ 是 $m$-拟相干的，应用引理
\ref{more-algebra-lemma-cone-pseudo-coherent} 可知，只需证明
$\tau_{\leq n - 1}K^\bullet$ 是 $m$-拟相干的。由对 $n$ 的归纳即可得证。
（拟相干的情形由此及引理
\ref{more-algebra-lemma-pseudo-coherent} 得到。）
\end{proof}

\begin{lemma}
\label{more-algebra-lemma-finite-push-pseudo-coherent}
设 $A \to B$ 为环同态。假设 $B$ 作为 $A$-模是拟相干的。设 $K^\bullet$ 为
$B$-模复形。下列条件等价：
\begin{enumerate}
\item $K^\bullet$ 作为 $B$-模复形是 $m$-拟相干的；以及
\item $K^\bullet$ 作为 $A$-模复形是 $m$-拟相干的。
\end{enumerate}
拟相干性的情形也有同样的等价性。
\end{lemma}

\begin{proof}
假设 (1)。选取有限自由 $B$-模组成的有界复形 $E^\bullet$ 以及映射
$\alpha : E^\bullet \to K^\bullet$，使其在次数 $> m$ 的上同调上诱导同构、
在次数 $m$ 上诱导满射。考虑可区别三角
$(E^\bullet, K^\bullet, C(\alpha)^\bullet)$。由引理
\ref{more-algebra-lemma-recognize-pseudo-coherent}，$C(\alpha)^\bullet$ 作为 $A$-模复形是
$m$-拟相干的。因此只需证明 $E^\bullet$ 作为 $A$-模复形是拟相干的，
这由引理\ref{more-algebra-lemma-complex-pseudo-coherent-modules} 得到。
(1) $\Rightarrow$ (2) 的拟相干情形由此及引理
\ref{more-algebra-lemma-pseudo-coherent} 得到。

\medskip\noindent
假设 (2)。令 $n$ 为使 $H^n(K^\bullet) \not = 0$ 的最大整数。
我们对 $n - m$ 归纳证明 $K^\bullet$ 作为 $B$-模复形是 $m$-拟相干的。
$n < m$ 的情形由引理\ref{more-algebra-lemma-recognize-pseudo-coherent} 得到。
选取有限自由 $A$-模组成的有界复形 $E^\bullet$ 以及映射
$\alpha : E^\bullet \to K^\bullet$，使其在次数 $> m$ 的上同调上诱导同构、
在次数 $m$ 上诱导满射。考虑诱导的复形映射
$$
\alpha \otimes 1 : E^\bullet \otimes_A B \to K^\bullet.
$$
注意，由于按构造 $H^n(E) \to H^n(E^\bullet \otimes_A B) \to H^n(K^\bullet)$
是满射，并且由例\ref{more-algebra-example-tor} 的谱序列，当 $i > n$ 时
$H^i(E^\bullet \otimes_A B) = 0$，故 $C(\alpha \otimes 1)^\bullet$ 在次数
$\geq n$ 上无上同调。
另一方面，由于 $K^\bullet$ 和 $E^\bullet \otimes_A B$（见引理
\ref{more-algebra-lemma-complex-pseudo-coherent-modules}）都是如此，
故由引理\ref{more-algebra-lemma-cone-pseudo-coherent}，
$C(\alpha \otimes 1)^\bullet$ 作为 $A$-模复形是 $m$-拟相干的。
因此由归纳，$C(\alpha \otimes 1)^\bullet$ 作为 $B$-模复形是 $m$-拟相干的。
最后再次应用引理\ref{more-algebra-lemma-cone-pseudo-coherent}，可知 $K^\bullet$ 作为
$B$-模复形是 $m$-拟相干的（显然 $E^\bullet \otimes_A B$ 作为 $B$-模复形是
拟相干的）。(2) $\Rightarrow$ (1) 的拟相干情形由此及引理
\ref{more-algebra-lemma-pseudo-coherent} 得到。
\end{proof}

\begin{lemma}
\label{more-algebra-lemma-pull-pseudo-coherent}
设 $A \to B$ 为环同态。设 $K^\bullet$ 为 $A$-模的 $m$-拟相干（分别为拟相干）
复形。则 $K^\bullet \otimes_A^{\mathbf{L}} B$ 是 $B$-模的
$m$-拟相干（分别为拟相干）复形。
\end{lemma}

\begin{proof}
首先注意到引理的陈述有意义：$K^\bullet$ 上有界，因而
$K^\bullet \otimes_A^{\mathbf{L}} B$ 由式
（\ref{more-algebra-equation-derived-tensor-algebra}）定义。
在此基础上，选取有限自由 $A$-模组成的有界复形 $E^\bullet$ 以及
$\alpha : E^\bullet \to K^\bullet$，使得当 $i > m$ 时 $H^i(\alpha)$ 是同构、
当 $i = m$ 时是满射。于是锥 $C(\alpha)^\bullet$ 在次数
$\geq m$ 上无上同调。由于 $-\otimes_A^{\mathbf{L}} B$ 是正合函子，得到
复形的可区别三角
$$
(E^\bullet \otimes_A^{\mathbf{L}} B, K^\bullet \otimes_A^{\mathbf{L}} B,
C(\alpha)^\bullet \otimes_A^{\mathbf{L}} B)
$$
其对象都是 $B$-模复形。由 Derived Categories，引理
\ref{derived-lemma-negative-vanishing} 的对偶可知，当 $i \geq m$ 时
$H^i(C(\alpha)^\bullet \otimes_A^{\mathbf{L}} B) = 0$。
由于 $E^\bullet$ 是有限投射 $R$-模组成的复形，可知
$E^\bullet \otimes_A^{\mathbf{L}} B = E^\bullet \otimes_A B$，因而
$$
E^\bullet \otimes_A B
\longrightarrow
K^\bullet \otimes_A^{\mathbf{L}} B
$$
是 $B$-模复形之间的映射，见证了
$K^\bullet \otimes_A^{\mathbf{L}} B$ 是 $m$-拟相干的事实。
拟相干复形的情形由 $m$-拟相干复形的情形及引理
\ref{more-algebra-lemma-pseudo-coherent} 得到。
\end{proof}

\begin{lemma}
\label{more-algebra-lemma-flat-base-change-pseudo-coherent}
设 $A \to B$ 为平坦环同态。设 $M$ 为 $A$-模，且是 $m$-拟相干的
（分别为拟相干的）。则 $M \otimes_A B$ 是 $m$-拟相干的
（分别为拟相干的）$B$-模。
\end{lemma}

\begin{proof}
这是引理\ref{more-algebra-lemma-pull-pseudo-coherent} 以及
$M \otimes_A^{\mathbf{L}} B = M \otimes_A B$（因为 $B$ 在 $A$ 上平坦）这一事实的
直接推论。
\end{proof}

\noindent
下面的引理也可由更强的引理
\ref{more-algebra-lemma-flat-descent-pseudo-coherent} 得到。

\begin{lemma}
\label{more-algebra-lemma-glue-pseudo-coherent}
设 $R$ 为环。设 $f_1, \ldots, f_r \in R$ 为生成单位理想的元素。设
$m \in \mathbf{Z}$。设 $K^\bullet$ 为 $R$-模复形。如果对每个 $i$，复形
$K^\bullet \otimes_R R_{f_i}$ 是 $m$-拟相干的（分别为拟相干的），则
$K^\bullet$ 是 $m$-拟相干的（分别为拟相干的）。
\end{lemma}

\begin{proof}
以下不再特别说明，使用 $- \otimes_R R_{f_i}$ 是正合函子、因而有
$$
H^i(K^\bullet)_{f_i} =
H^i(K^\bullet) \otimes_R R_{f_i} = H^i(K^\bullet \otimes_R R_{f_i}).
$$
假设当 $i = 1, \ldots, r$ 时 $K^\bullet \otimes_R R_{f_i}$ 是 $m$-拟相干的。
令 $n \in \mathbf{Z}$ 为使某个 $i$ 的
$H^n(K^\bullet \otimes_R R_{f_i})$ 非零的最大整数。这特别推出当 $i > n$ 时
$H^i(K^\bullet) = 0$（并且 $H^n(K^\bullet) \not = 0$），见
Algebra，引理\ref{algebra-lemma-cover}。我们对 $n - m$ 归纳证明引理。
如果 $n < m$，则由引理\ref{more-algebra-lemma-recognize-pseudo-coherent}，引理成立。
如果 $n \geq m$，则对每个 $i$，由引理
\ref{more-algebra-lemma-finite-cohomology}，$H^n(K^\bullet)_{f_i}$ 是有限的
$R_{f_i}$-模。因此由 Algebra，引理\ref{algebra-lemma-cover}，
$H^n(K^\bullet)$ 是有限 $R$-模。选取有限自由 $R$-模 $E$ 以及满射
$E \to H^n(K^\bullet)$。由于 $E$ 是投射的，可将其提升为复形映射
$\alpha : E[-n] \to K^\bullet$。于是锥 $C(\alpha)^\bullet$ 在次数
$\geq n$ 上的上同调为零。另一方面，对每个 $i$，复形
$C(\alpha)^\bullet \otimes_R R_{f_i}$ 是 $m$-拟相干的，见引理
\ref{more-algebra-lemma-cone-pseudo-coherent}。因此由归纳，$C(\alpha)^\bullet$ 作为
$R$-模复形是 $m$-拟相干的。再次应用引理
\ref{more-algebra-lemma-cone-pseudo-coherent} 即得结论。
\end{proof}

\begin{lemma}
\label{more-algebra-lemma-flat-descent-pseudo-coherent}
设 $R$ 为环。设 $m \in \mathbf{Z}$。设 $K^\bullet$ 为 $R$-模复形。设
$R \to R'$ 为忠实平坦的环同态。如果复形
$K^\bullet \otimes_R R'$ 是 $m$-拟相干的（分别为拟相干的），则
$K^\bullet$ 是 $m$-拟相干的（分别为拟相干的）。
\end{lemma}

\begin{proof}
以下不再特别说明，使用 $- \otimes_R R'$ 是正合函子，因而有
$$
H^i(K^\bullet) \otimes_R R' = H^i(K^\bullet \otimes_R R').
$$
假设 $K^\bullet \otimes_R R'$ 是 $m$-拟相干的。令
$n \in \mathbf{Z}$ 为使 $H^n(K^\bullet)$ 非零的最大整数；于是 $n$ 也是使
$H^n(K^\bullet \otimes_R R')$ 非零的最大整数。
我们对 $n - m$ 归纳证明引理。如果 $n < m$，则由引理
\ref{more-algebra-lemma-recognize-pseudo-coherent}，引理成立。
如果 $n \geq m$，则 $H^n(K^\bullet) \otimes_R R'$ 是有限
$R'$-模，见引理\ref{more-algebra-lemma-finite-cohomology}。
因此由 Algebra，引理\ref{algebra-lemma-descend-properties-modules}，
$H^n(K^\bullet)$ 是有限 $R$-模。选取有限自由 $R$-模 $E$ 以及满射
$E \to H^n(K^\bullet)$。由于 $E$ 是投射的，可将其提升为复形映射
$\alpha : E[-n] \to K^\bullet$。于是锥 $C(\alpha)^\bullet$ 在次数
$\geq n$ 上的上同调为零。另一方面，复形
$C(\alpha)^\bullet \otimes_R R'$ 是 $m$-拟相干的，见引理
\ref{more-algebra-lemma-cone-pseudo-coherent}。
因此由归纳，$C(\alpha)^\bullet$ 作为 $R$-模复形是 $m$-拟相干的。
再次应用引理\ref{more-algebra-lemma-cone-pseudo-coherent} 即得结论。
\end{proof}

\begin{lemma}
\label{more-algebra-lemma-tensor-pseudo-coherent}
设 $R$ 为环。设 $K, L$ 为 $D(R)$ 中的对象。
\begin{enumerate}
\item 如果 $K$ 是 $n$-拟相干的且当 $i > a$ 时 $H^i(K) = 0$，并且
$L$ 是 $m$-拟相干的且当 $j > b$ 时 $H^j(L) = 0$，则
$K \otimes_R^\mathbf{L} L$ 是 $t$-拟相干的，其中
$t = \max(m + a, n + b)$。
\item 如果 $K$ 和 $L$ 都是拟相干的，则
$K \otimes_R^\mathbf{L} L$ 是拟相干的。
\end{enumerate}
\end{lemma}

\begin{proof}
证明 (1)。不妨设存在有限自由 $R$-模组成的有界复形 $K^\bullet$ 和
$L^\bullet$，以及映射 $\alpha : K^\bullet \to K$ 和
$\beta : L^\bullet \to L$，使得当 $i > n$ 时 $H^i(\alpha)$ 是同构、
当 $i = n$ 时是满射，并且当 $i > m$ 时 $H^i(\beta)$ 是同构、
当 $i = m$ 时是满射。于是映射
$$
\alpha \otimes^\mathbf{L} \beta :
\text{Tot}(K^\bullet \otimes_R L^\bullet)
\to K \otimes_R^\mathbf{L} L
$$
在次数 $i$ 的上同调上，当 $i > t$ 时诱导同构，在 $i = t$ 时诱导满射。
这是由 tors 的谱序列得到的（细节略）。(2) 由 (1) 及引理
\ref{more-algebra-lemma-pseudo-coherent} 得到。
\end{proof}

\begin{lemma}
\label{more-algebra-lemma-Noetherian-pseudo-coherent}
设 $R$ 为 Noether 环。则
\begin{enumerate}
\item $R$-模复形 $K^\bullet$ 是 $m$-拟相干的，当且仅当
$K^\bullet \in D^{-}(R)$ 且当 $i \geq m$ 时 $H^i(K^\bullet)$ 是有限
$R$-模。
\item $R$-模复形 $K^\bullet$ 是拟相干的，当且仅当
$K^\bullet \in D^{-}(R)$ 且对所有 $i$，$H^i(K^\bullet)$ 是有限 $R$-模。
\item $R$-模是拟相干的，当且仅当它是有限的。
\end{enumerate}
\end{lemma}

\begin{proof}
在 Algebra，引理\ref{algebra-lemma-resolution-by-finite-free} 中，
我们已经看到每个有限 $R$-模都是拟相干的。
另一方面，拟相干模是有限的，见引理\ref{more-algebra-lemma-n-pseudo-module}。
因此 (3) 成立。假设 $K^\bullet$ 是 $m$-拟相干复形。
则存在有限自由 $R$-模组成的有界复形 $E^\bullet$，使得当 $i > m$ 时
$H^i(K^\bullet)$ 与 $H^i(E^\bullet)$ 同构，并且
$H^m(K^\bullet)$ 是 $H^m(E^\bullet)$ 的商。因此显然对每个 $i \geq m$，
$H^i(K^\bullet)$ 都是有限模。(1) 的逆向蕴含由引理
\ref{more-algebra-lemma-cohomology-pseudo-coherent} 及 (3) 得到。
(2) 由 (1) 及引理\ref{more-algebra-lemma-pseudo-coherent} 得到。
\end{proof}

\begin{lemma}
\label{more-algebra-lemma-coherent-pseudo-coherent}
设 $R$ 为相容环（Algebra，定义
\ref{algebra-definition-coherent}）。设 $K \in D^-(R)$。下列条件等价：
\begin{enumerate}
\item $K$ 是 $m$-拟相干的；
\item $H^m(K)$ 是有限 $R$-模，且当 $i > m$ 时 $H^i(K)$ 是相容模；以及
\item $H^m(K)$ 是有限 $R$-模，且当 $i > m$ 时
$H^i(K)$ 是有限表示的。
\end{enumerate}
因此，$K$ 拟相干当且仅当对所有 $i$，$H^i(K)$ 是相容模。
\end{lemma}

\begin{proof}
回顾：$R$-模 $M$ 相容，当且仅当它是有限表示的（Algebra，引理
\ref{algebra-lemma-coherent-ring}）。这说明 (2) 与 (3) 等价。
在此情形下，若选取正合序列
$0 \to N \to R^{\oplus m} \to M \to 0$，则由 Algebra，引理
\ref{algebra-lemma-coherent}，$N$ 是相容模。因此重复对 $N$ 作此过程，得到解消
$$
\ldots \to R^{\oplus n} \to R^{\oplus m} \to M \to 0
$$
其各项都是有限自由 $R$-模。换言之，$M$ 是拟相干的。
(1) 与 (2) 的等价性由此以及引理
\ref{more-algebra-lemma-cohomology-pseudo-coherent} 和
\ref{more-algebra-lemma-n-pseudo-module} 得到。
最后一个断言由 (1) 与 (2) 的等价性以及引理
\ref{more-algebra-lemma-pseudo-coherent} 得到。
\end{proof}









\section{拟相干模，II}
\label{more-algebra-section-pseudo-coherent-bis}

\noindent
我们继续第\ref{more-algebra-section-pseudo-coherent}节开始的讨论。

\begin{lemma}
\label{more-algebra-lemma-pseudo-coherence-colimit-ext}
设 $R$ 为环。设 $M = \colim M_i$ 为 $R$-模的滤过余极限。
设 $K \in D(R)$ 为 $m$-拟相干的。则当 $n < -m$ 时
$\colim \Ext^n_R(K, M_i) = \Ext^n_R(K, M)$，并且
$\colim \Ext^{-m}_R(K, M_i) \to \Ext^{-m}_R(K, M)$ 是单射。
\end{lemma}

\begin{proof}
由定义，可以找到可区别三角
$$
E \to K \to L \to E[1]
$$
在 $D(R)$ 中，其中 $E$ 由有限自由 $R$-模组成的有界复形表示，并且当 $i \geq m$ 时
$H^i(L) = 0$。于是对任意 $R$-模 $N$，当 $n \leq -m$ 时
$\Ext^n_R(L, N) = 0$，见 Derived Categories，引理
\ref{derived-lemma-negative-exts}。
由与该可区别三角关联的 $\Ext$ 长正合序列可知，
$\Ext^n_R(K, N) \to \Ext^n_R(E, N)$ 在 $n < -m$ 时是同构，
在 $n = -m$ 时是单射。因此只需证明，当 $E$ 可由有限自由 $R$-模组成的
有界复形 $E^\bullet$ 表示时，函子 $M \mapsto \Ext_R^n(E, M)$ 与滤过余极限交换。
$\Ext^n_R(E, M)$ 由复形 $\Hom_R(E^\bullet, M)$ 计算，见 Derived Categories，
引理\ref{derived-lemma-morphisms-from-projective-complex}。
由于 $E^p$ 是有限自由的，函子 $M \mapsto \Hom_R(E^p, M)$ 与滤过余极限交换。
因此作为复形有
$\Hom_R(E^\bullet, M) = \colim \Hom_R(E^\bullet, M_i)$。
由于滤过余极限是正合的（Algebra，引理
\ref{algebra-lemma-directed-colimit-exact}），结论成立。
\end{proof}

\begin{lemma}
\label{more-algebra-lemma-characterize-pseudo-coherent-colimit-ext}
设 $R$ 为环。设 $K \in D^-(R)$。设 $m \in \mathbf{Z}$。
则 $K$ 是 $m$-拟相干的，当且仅当对任意 $R$-模的滤过余极限
$M = \colim M_i$，都有
$\colim \Ext^n_R(K, M_i) = \Ext^n_R(K, M)$（当 $n < -m$），并且
$\colim \Ext^{-m}_R(K, M_i) \to \Ext^{-m}_R(K, M)$ 是单射。
\end{lemma}

\begin{proof}
一个方向已在引理\ref{more-algebra-lemma-pseudo-coherence-colimit-ext} 中证明。
现在假设对任意 $R$-模的滤过余极限 $M = \colim M_i$，都有
$$
\colim \Ext^n_R(K, M_i) = \Ext^n_R(K, M)
$$
（当 $n < -m$），并且
$$
\colim \Ext^{-m}_R(K, M_i) \to \Ext^{-m}_R(K, M)
$$
是单射。
我们证明 $K$ 是 $m$-拟相干的。

\medskip\noindent
令 $t$ 为使 $H^t(K)$ 非零的最大整数。我们对 $t$ 归纳。
如果 $t < m$，则由引理\ref{more-algebra-lemma-recognize-pseudo-coherent}，$K$ 是
$m$-拟相干的。如果 $t \geq m$，由于
$\Hom_R(H^t(K), M) = \Ext^{-t}_R(K, M)$，可知对任意滤过余极限
$M = \colim M_i$，映射
$\colim \Hom_R(H^t(K), M_i) \to \Hom_R(H^t(K), M)$ 是单射。
由 Algebra，引理\ref{algebra-lemma-characterize-finite-module-hom}，
这推出 $H^t(K)$ 是有限 $R$-模。选取有限自由 $R$-模 $F$ 以及满射
$F \to H^t(K)$。可将其提升为 $D(R)$ 中的态射 $F[-t] \to K$，并选取可区别三角
$$
F[-t] \to K \to L \to F[-t + 1]
$$
在 $D(R)$ 中。于是当 $i \geq t$ 时 $H^i(L) = 0$。此外，与该可区别三角关联的
$\Ext$ 长正合序列表明，经过一个略去的简短论证，$L$ 继承我们对 $K$ 所作的假设。
由对 $t$ 的归纳，$L$ 是 $m$-拟相干的。因此由引理
\ref{more-algebra-lemma-cone-pseudo-coherent}，$K$ 是 $m$-拟相干的。
\end{proof}

\begin{lemma}
\label{more-algebra-lemma-pseudo-coherence-and-ext}
设 $R$ 为环。设 $L$、$M$、$N$ 为 $R$-模。
\begin{enumerate}
\item 如果 $M$ 是有限表示的且 $L$ 平坦，则典范映射
$\Hom_R(M, N) \otimes_R L \to \Hom_R(M, N \otimes_R L)$
是同构。
\item 如果 $M$ 是 $(-m)$-拟相干的且 $L$ 平坦，则典范映射
$\Ext^i_R(M, N) \otimes_R L \to \Ext^i_R(M, N \otimes_R L)$
在 $i < m$ 时是同构。
\end{enumerate}
\end{lemma}

\begin{proof}
选取各项为自由 $R$-模的解消 $F_\bullet \to M$，见 Algebra，引理
\ref{algebra-lemma-resolution-by-finite-free}。复形
$\Hom_R(F_\bullet, N)$ 计算 $\Ext^i_R(M, N)$，而复形
$\Hom_R(F_\bullet, N \otimes_R L)$ 计算 $\Ext^i_R(M, N \otimes_R L)$。
总有复形映射
$$
\Hom_R(F_\bullet, N) \otimes_R L
\longrightarrow
\Hom_R(F_\bullet, N \otimes_R L)
$$
它对所有 $i \geq 0$ 诱导典范映射
$\Ext^i_R(M, N) \otimes_R L \to \Ext^i_R(M, N \otimes_R L)$
（例如，在这些映射不依赖于解消 $F_\bullet$ 的选取这一意义下是典范的）。
如果 $L$ 平坦，则复形 $\Hom_R(F_\bullet, N) \otimes_R L$ 计算
$\Ext^i_R(M, N) \otimes_R L$，因为取上同调与张量 $L$ 可交换。

\medskip\noindent
在上述事实基础上，如果 $M$ 是 $(-m)$-拟相干的，则可以选取 $F_\bullet$，使得
当 $i = 0, \ldots, m$ 时 $F_i$ 是有限自由的。于是上面显示的复形映射在次数
$\leq m$ 上是同构，因而在次数 $< m$ 上对上同调群诱导同构。这就证明了 (2)。
如果 $M$ 是有限表示的，则由引理\ref{more-algebra-lemma-n-pseudo-module}，$M$ 是
$(-1)$-拟相干的；再由 $\Hom = \Ext^0$ 得到结论。
\end{proof}

\begin{lemma}
\label{more-algebra-lemma-pseudo-coherence-and-base-change-ext}
设 $R \to R'$ 为平坦环同态。设 $M$、$N$ 为 $R$-模。
\begin{enumerate}
\item 如果 $M$ 是有限表示的 $R$-模，则
$\Hom_R(M, N) \otimes_R R' = \Hom_{R'}(M \otimes_R R', N \otimes_R R')$。
\item 如果 $M$ 是 $(-m)$-拟相干的，则
$\Ext^i_R(M, N) \otimes_R R' = \Ext^i_{R'}(M \otimes_R R', N \otimes_R R')$
当 $i < m$ 时成立。
\end{enumerate}
特别地，如果 $R$ 是 Noether 环且 $M$ 是有限模，则上述等式对所有 $i$ 都成立。
\end{lemma}

\begin{proof}
由 Algebra，引理\ref{algebra-lemma-flat-base-change-ext}，有
$\Ext^i_{R'}(M \otimes_R R', N \otimes_R R') =
\Ext^i_R(M, N \otimes_R R')$。
结合引理\ref{more-algebra-lemma-pseudo-coherence-and-ext}，即得 (1) 和 (2)。
最后一个断言由此及引理\ref{more-algebra-lemma-Noetherian-pseudo-coherent} 得到。
\end{proof}

\begin{lemma}
\label{more-algebra-lemma-pseudo-coherent-tensor}
设 $R$ 为环。设 $K \in D^-(R)$。下列条件等价：
\begin{enumerate}
\item $K$ 是拟相干的；
\item 对 $R$-模的任意族 $(Q_{\alpha})_{\alpha \in A}$，典范映射
$$
\alpha :
K \otimes_R^\mathbf{L} \left( \prod\nolimits_\alpha Q_{\alpha} \right)
\longrightarrow
\prod\nolimits_\alpha (K \otimes_R^\mathbf{L} Q_{\alpha})
$$
是 $D(R)$ 中的同构；
\item 对每个 $R$-模 $Q$ 和每个集合 $A$，典范映射
$$
\beta : K \otimes_R^\mathbf{L} Q^A \longrightarrow (K \otimes_R^\mathbf{L} Q)^A
$$
是 $D(R)$ 中的同构；以及
\item 对每个集合 $A$，典范映射
$$
\gamma : K \otimes_R^\mathbf{L} R^A \longrightarrow K^A
$$
是 $D(R)$ 中的同构。
\end{enumerate}
给定 $m \in \mathbf{Z}$，下列条件等价：
\begin{enumerate}
\item[(a)] $K$ 是 $m$-拟相干的；
\item[(b)] 对 $R$-模的任意族 $(Q_{\alpha})_{\alpha \in A}$，并令 $\alpha$ 如上，
$H^i(\alpha)$ 当 $i > m$ 时是同构、当 $i = m$ 时是满射；
\item[(c)] 对每个 $R$-模 $Q$ 和每个集合 $A$，并令 $\beta$ 如上，
$H^i(\beta)$ 当 $i > m$ 时是同构、当 $i = m$ 时是满射；
\item[(d)] 对每个集合 $A$，并令 $\gamma$ 如上，
$H^i(\gamma)$ 当 $i > m$ 时是同构、当 $i = m$ 时是满射。
\end{enumerate}
\end{lemma}

\begin{proof}
如果 $K$ 拟相干，则 $K$ 可由有限自由 $R$-模组成的上有界复形表示。
于是导出张量积由与此复形张量计算。并且，$D(R)$ 中的乘积可通过对任意
代表复形的选取取乘积来给出。因此由模的相应事实，(1) 推出 (2)、(3)、(4)，见
Algebra，命题\ref{algebra-proposition-fp-tensor}。

\medskip\noindent
同理（使用张量积是右正合的），读者可证明 (a) 推出 (b)、(c) 和 (d)。

\medskip\noindent
假设 (4) 成立。要证明 $K$ 拟相干，只需证明对所有 $m$，$K$ 都是
$m$-拟相干的（引理\ref{more-algebra-lemma-pseudo-coherent}）。因此只需证明
(d) 推出 (a)。

\medskip\noindent
假设 (d)。令 $i$ 为使 $H^i(K)$ 非零的最大整数。如果 $i < m$，则结论成立。
否则，由 (d) 以及上面对 $D(R)$ 中乘积的描述，可知
$H^i(K) \otimes_R R^A \to H^i(K)^A$ 是满射。因此由 Algebra，命题
\ref{algebra-proposition-fg-tensor}，$H^i(K)$ 是有限生成 $R$-模。
于是可以选取一个复形 $L$，其仅由位于次数 $i$ 的有限自由模组成，并选取复形映射
$L \to K$，使 $H^i(L) \to H^i(K)$ 满射。特别地，$L$ 满足 (1)、(2)、(3)、(4)。
选取可区别三角
$$
L \to K \to M \to L[1]
$$
则 $H^j(M) = 0$（当 $j \geq i$）。另一方面，经过一个略去的简短论证，$M$
仍具有性质 (d)。由对 $i$ 的归纳，$M$ 是 $m$-拟相干的。因此由引理
\ref{more-algebra-lemma-cone-pseudo-coherent}，$K$ 是 $m$-拟相干的。
\end{proof}

\begin{lemma}
\label{more-algebra-lemma-detect-cohomology-pseudo-coherent}
设 $R$ 为环。设 $K \in D(R)$ 为拟相干对象。设 $i \in \mathbf{Z}$。
存在有限表示的 $R$-模 $M$ 以及 $D(R)$ 中的映射 $K \to M[-i]$，
该映射在 $H^i(K) \to M$ 上诱导单射。
\end{lemma}

\begin{proof}
由定义\ref{more-algebra-definition-pseudo-coherent}，可以用有限自由 $R$-模组成的复形
$P^\bullet$ 表示 $K$。置 $M = \Coker(P^{i - 1} \to P^i)$。
\end{proof}

\begin{lemma}
\label{more-algebra-lemma-detect-cohomology}
设 $A$ 为 Noether 环。设 $K \in D(A)$ 为拟相干对象，即
$K \in D^-(A)$ 且上同调模有限。设 $\mathfrak m$ 为 $A$ 的极大理想。
如果 $H^i(K)/\mathfrak m H^i(K) \not = 0$，则存在被 $\mathfrak m$ 的某个幂
湮灭的有限 $A$-模 $E$，以及映射 $K \to E[-i]$，该映射在 $H^i(K)$ 上非零。
\end{lemma}

\begin{proof}
（引理陈述中拟相干性的等价表述是引理
\ref{more-algebra-lemma-Noetherian-pseudo-coherent}。）
按引理\ref{more-algebra-lemma-detect-cohomology-pseudo-coherent}，选取
$K \to M[-i]$。由 Artin-Rees（Algebra，引理
\ref{algebra-lemma-Artin-Rees}），可找到 $n$ 使得
$H^i(K) \cap \mathfrak m^n M \subset \mathfrak m H^i(K)$。
取 $E = M/\mathfrak m^n M$。
\end{proof}








\section{Tor 维数}
\label{more-algebra-section-tor}

\noindent
我们不再用投射模作解消，而是考察平坦模的解消。这引出如下概念。

\begin{definition}
\label{more-algebra-definition-tor-amplitude}
设 $R$ 为环。记 $D(R)$ 为其导出范畴。设 $a, b \in \mathbf{Z}$。
\begin{enumerate}
\item 如果对所有 $R$-模 $M$ 以及所有 $i \not \in [a, b]$，都有
$H^i(K^\bullet \otimes_R^\mathbf{L} M) = 0$，则称 $D(R)$ 中的对象
$K^\bullet$ 具有 $[a, b]$ 内的{\it Tor-振幅}。
\item 如果某些 $a, b$ 使其具有 $[a, b]$ 内的 Tor-振幅，则称 $D(R)$ 中的对象
$K^\bullet$ 具有{\it 有限 Tor 维数}。
\item 如果 $R$-模 $M$ 作为 $D(R)$ 中的对象 $M[0]$ 具有
$[-d, 0]$ 内的 Tor-振幅，则称该模具有{\it Tor 维数 $\leq d$}。
\item 如果 $R$-模 $M$ 作为 $D(R)$ 中的对象 $M[0]$ 具有有限 Tor 维数，
则称该模具有{\it 有限 Tor 维数}。
\end{enumerate}
\end{definition}

\noindent
注意，如果 $K^\bullet$ 具有有限 Tor 维数，则
$K^\bullet \in D^b(R)$。

\begin{lemma}
\label{more-algebra-lemma-last-one-flat}
设 $R$ 为环。设 $K^\bullet$ 为由平坦 $R$-模组成的上有界复形，
且具有 $[a, b]$ 内的 Tor-振幅。则
$\Coker(d_K^{a - 1})$ 是平坦 $R$-模。
\end{lemma}

\begin{proof}
由于 $K^\bullet$ 是平坦模组成的上有界复形，可知
$K^\bullet \otimes_R M = K^\bullet \otimes_R^{\mathbf{L}} M$。
因此对每个 $R$-模 $M$，序列
$$
K^{a - 2} \otimes_R M \to K^{a - 1} \otimes_R M \to K^a \otimes_R M
$$
在中间项处正合。由于
$K^{a - 2} \to K^{a - 1} \to K^a \to \Coker(d_K^{a - 1}) \to 0$
是平坦解消，这推出
$\text{Tor}_1^R(\Coker(d_K^{a - 1}), M) = 0$
对所有 $R$-模 $M$ 都成立。这意味着
$\Coker(d_K^{a - 1})$ 是平坦的，见
Algebra，引理\ref{algebra-lemma-characterize-flat}。
\end{proof}

\begin{lemma}
\label{more-algebra-lemma-tor-amplitude}
设 $R$ 为环。设 $K^\bullet$ 为 $D(R)$ 中的对象。设 $a, b \in \mathbf{Z}$。
下列条件等价：
\begin{enumerate}
\item $K^\bullet$ 具有 $[a, b]$ 内的 Tor-振幅；
\item $K^\bullet$ 拟同构于平坦 $R$-模复形 $E^\bullet$，且当
$i \not \in [a, b]$ 时 $E^i = 0$。
\end{enumerate}
\end{lemma}

\begin{proof}
如果 (2) 成立，则可以计算
$K^\bullet \otimes_R^\mathbf{L} M = E^\bullet \otimes_R M$，故 (1) 显然成立。
假设 (1)。可以用一个满足当 $i > b$ 时 $K^i = 0$ 的投射解消替代
$K^\bullet$。见 Derived Categories，引理
\ref{derived-lemma-projective-resolutions-exist}。置
$E^\bullet = \tau_{\geq a}K^\bullet$。除 $E^a$ 平坦外一切都显然；其
平坦性立即由引理\ref{more-algebra-lemma-last-one-flat} 及定义得到。
\end{proof}

\begin{lemma}
\label{more-algebra-lemma-bounded-below-tor-amplitude}
设 $R$ 为环。设 $a \in \mathbf{Z}$，且设 $K$ 为 $D(R)$ 中的对象。
下列条件等价：
\begin{enumerate}
\item $K$ 具有 $[a, \infty]$ 内的 Tor-振幅；
\item $K$ 拟同构于一个 K-平坦复形 $E^\bullet$，其各项为平坦 $R$-模，且
$E^i = 0$ 当 $i \not \in [a, \infty]$ 时成立。
\end{enumerate}
\end{lemma}

\begin{proof}
(2) $\Rightarrow$ (1) 立即成立。设 (1) 成立。
首先选取一个各项平坦且表示 $K$ 的 K-平坦复形 $K^\bullet$，
见引理\ref{more-algebra-lemma-K-flat-resolution}。
对任意 $R$-模 $M$，复形
$$
K^{n - 1} \otimes_R M \to K^n \otimes_R M \to K^{n + 1} \otimes_R M
$$
的上同调计算 $H^n(K \otimes_R^\mathbf{L} M)$。当 $n < a$ 时它总为零。
因此将引理\ref{more-algebra-lemma-last-one-flat} 应用于复形
$\ldots \to K^{a - 1} \to K^a \to K^{a + 1}$，
可得 $N = \Coker(K^{a - 1} \to K^a)$ 是平坦 $R$-模。置
$$
E^\bullet = \tau_{\geq a}K^\bullet =
(\ldots \to 0 \to N \to K^{a + 1} \to \ldots )
$$
复形 $K^\bullet \to E^\bullet$ 的核 $L^\bullet$ 为复形
$$
L^\bullet = (\ldots \to K^{a - 1} \to I \to 0 \to \ldots)
$$
其中 $I \subset K^a$ 是 $K^{a - 1} \to K^a$ 的像。
由于有短正合列 $0 \to I \to K^a \to N \to 0$，
可知 $I$ 是平坦 $R$-模。因此 $L^\bullet$ 是由平坦模组成的上有界复形，
从而由引理\ref{more-algebra-lemma-derived-tor-quasi-isomorphism} 可知它是 K-平坦的。
由引理\ref{more-algebra-lemma-K-flat-two-out-of-three-ses}，$E^\bullet$ 也是 K-平坦的。
\end{proof}

\begin{lemma}
\label{more-algebra-lemma-cone-tor-amplitude}
设 $R$ 为环。
设 $(K^\bullet, L^\bullet, M^\bullet, f, g, h)$ 为 $D(R)$ 中的可区别三角。
设 $a, b \in \mathbf{Z}$。
\begin{enumerate}
\item 若 $K^\bullet$ 具有 $[a + 1, b + 1]$ 内的 Tor-振幅，且
$L^\bullet$ 具有 $[a, b]$ 内的 Tor-振幅，则 $M^\bullet$ 具有
$[a, b]$ 内的 Tor-振幅。
\item 若 $K^\bullet, M^\bullet$ 具有 $[a, b]$ 内的 Tor-振幅，则
$L^\bullet$ 具有 $[a, b]$ 内的 Tor-振幅。
\item 若 $L^\bullet$ 具有 $[a + 1, b + 1]$ 内的 Tor-振幅，且
$M^\bullet$ 具有 $[a, b]$ 内的 Tor-振幅，则
$K^\bullet$ 具有 $[a + 1, b + 1]$ 内的 Tor-振幅。
\end{enumerate}
\end{lemma}

\begin{proof}
略去证明。结论直接来自可区别三角所关联的上同调长正合列，以及
$- \otimes_R^{\mathbf{L}} M$ 保持可区别三角这一事实。
最容易证明的是 (2)，其余结论由平移从 (2) 得到。
\end{proof}

\begin{lemma}
\label{more-algebra-lemma-tor-dimension}
设 $R$ 为环。设 $M$ 为一个 $R$-模。
设 $d \geq 0$。下列条件等价：
\begin{enumerate}
\item $M$ 具有 Tor 维数 $\leq d$；
\item 存在解消
$$
0 \to F_d \to \ldots \to F_1 \to F_0 \to M \to 0
$$
其中 $F_i$ 为平坦 $R$-模。
\end{enumerate}
特别地，$R$-模的 Tor 维数为 $0$ 当且仅当它是平坦 $R$-模。
\end{lemma}

\begin{proof}
设 (2) 成立。取 $E^\bullet$ 为满足 $E^{-i} = F_i$ 的复形，
则它与 $M$ 拟同构。因此由引理\ref{more-algebra-lemma-tor-amplitude}，
$M$ 的 Tor 维数至多为 $d$。
反之，设 (1) 成立。令 $P^\bullet \to M$ 为 $M$ 的投射解消。
由引理\ref{more-algebra-lemma-last-one-flat} 可知，$\tau_{\geq -d}P^\bullet$
是 $M$ 的长度为 $d$ 的平坦解消，即 (2) 成立。
\end{proof}

\begin{lemma}
\label{more-algebra-lemma-summands-tor-amplitude}
设 $R$ 为环。设 $a, b \in \mathbf{Z}$。
若 $K^\bullet \oplus L^\bullet$ 具有 $[a, b]$ 内的 Tor 振幅，
则 $K^\bullet$ 和 $L^\bullet$ 也具有该 Tor 振幅。
\end{lemma}

\begin{proof}
这由 Tor 函子具有可加性立即得到。
\end{proof}

\begin{lemma}
\label{more-algebra-lemma-complex-finite-tor-dimension-modules}
设 $R$ 为环。设 $K^\bullet$ 为由 $R$-模组成的有界复形，
且对每个 $i$，$K^i$ 具有 $[a - i, b - i]$ 内的 Tor 振幅。
则 $K^\bullet$ 具有 $[a, b]$ 内的 Tor 振幅。特别地，
若 $K^\bullet$ 是由有限个具有有限 Tor 维数的 $R$-模组成的复形，
则 $K^\bullet$ 具有有限 Tor 维数。
\end{lemma}

\begin{proof}
对有限复形的长度作归纳即可；使用引理\ref{more-algebra-lemma-cone-tor-amplitude}
以及朴素截断。
\end{proof}

\begin{lemma}
\label{more-algebra-lemma-cohomology-tor-amplitude}
设 $R$ 为环。设 $a, b \in \mathbf{Z}$。设 $K^\bullet \in D^b(R)$，
使得对所有 $i$，$H^i(K^\bullet)$ 具有 $[a - i, b - i]$ 内的 Tor 振幅。
则 $K^\bullet$ 具有 $[a, b]$ 内的 Tor 振幅。特别地，
若 $K^\bullet \in D^b(R)$ 且其所有上同调群都具有有限 Tor
维数，则 $K^\bullet$ 具有有限 Tor 维数。
\end{lemma}

\begin{proof}
对有限复形的长度作归纳即可；使用引理\ref{more-algebra-lemma-cone-tor-amplitude}
以及典范截断。
\end{proof}

\begin{lemma}
\label{more-algebra-lemma-push-tor-amplitude}
设 $A \to B$ 为环映射。设 $K^\bullet$ 和 $L^\bullet$ 为由
$B$-模组成的复形。设 $a, b, c, d \in \mathbf{Z}$。若
\begin{enumerate}
\item $K^\bullet$ 作为 $B$-模复形具有 $[a, b]$ 内的 Tor 振幅；
\item $L^\bullet$ 作为 $A$-模复形具有 $[c, d]$ 内的 Tor 振幅，
\end{enumerate}
则 $K^\bullet \otimes^\mathbf{L}_B L^\bullet$ 作为 $A$-模复形
具有 $[a + c, b + d]$ 内的 Tor 振幅。
\end{lemma}

\begin{proof}
不妨设 $K^\bullet$ 是由平坦 $B$-模组成且当
$i \not \in [a, b]$ 时 $K^i = 0$ 的复形，见引理\ref{more-algebra-lemma-tor-amplitude}。
令 $M$ 为 $A$-模。选取自由解消 $F^\bullet \to M$。则
$$
(K^\bullet \otimes_B^\mathbf{L} L^\bullet) \otimes_A^{\mathbf{L}} M =
\text{Tot}(\text{Tot}(K^\bullet \otimes_B L^\bullet) \otimes_A F^\bullet) =
\text{Tot}(K^\bullet \otimes_B \text{Tot}(L^\bullet \otimes_A F^\bullet))
$$
关于第二个等式，见 Homology，引理\ref{homology-remark-triple-complex}。
由假设 (2)，复形
$\text{Tot}(L^\bullet \otimes_A F^\bullet)$
仅在次数 $[c, d]$ 中有非零上同调。因此，Homology，
引理\ref{homology-lemma-ss-double-complex} 关于双复形
$K^\bullet \otimes_B \text{Tot}(L^\bullet \otimes_A F^\bullet)$
的谱序列表明，
$(K^\bullet \otimes_B^\mathbf{L} L^\bullet) \otimes_A^{\mathbf{L}} M$
仅在次数 $[a + c, b + d]$ 中有非零上同调。
\end{proof}

\begin{lemma}
\label{more-algebra-lemma-flat-push-tor-amplitude}
设 $A \to B$ 为环映射。假设 $B$ 作为 $A$-模是平坦的。设
$K^\bullet$ 为由 $B$-模组成的复形。设 $a, b \in \mathbf{Z}$。若
$K^\bullet$ 作为 $B$-模复形具有 $[a, b]$ 内的 Tor 振幅，
则 $K^\bullet$ 作为 $A$-模复形具有 $[a, b]$ 内的 Tor 振幅。
\end{lemma}

\begin{proof}
这是引理\ref{more-algebra-lemma-push-tor-amplitude} 的特例，但也可直接如下看出。我们有
$K^\bullet \otimes_A^{\mathbf{L}} M =
K^\bullet \otimes_B^{\mathbf{L}} (M \otimes_A B)$，
因为 $K^\bullet$ 作为 $B$-模复形的任意投射解消都是
$K^\bullet$ 作为 $A$-模复形的平坦解消，并可用于计算
$K^\bullet \otimes_A^{\mathbf{L}} M$。
\end{proof}

\begin{lemma}
\label{more-algebra-lemma-finite-tor-dimension-push-tor-amplitude}
设 $A \to B$ 为环映射。假设 $B$ 作为 $A$-模具有 Tor 维数 $\leq d$。
设 $K^\bullet$ 为由 $B$-模组成的复形。设 $a, b \in \mathbf{Z}$。若
$K^\bullet$ 作为 $B$-模复形具有 $[a, b]$ 内的 Tor 振幅，
则 $K^\bullet$ 作为 $A$-模复形具有 $[a - d, b]$ 内的 Tor 振幅。
\end{lemma}

\begin{proof}
这是引理\ref{more-algebra-lemma-push-tor-amplitude} 的特例，但也可直接如下看出。
令 $M$ 为 $A$-模。选取自由解消 $F^\bullet \to M$。则
$$
K^\bullet \otimes_A^{\mathbf{L}} M =
\text{Tot}(K^\bullet \otimes_A F^\bullet) =
\text{Tot}(K^\bullet \otimes_B (F^\bullet \otimes_A B)) =
K^\bullet \otimes_B^{\mathbf{L}} (M \otimes_A^{\mathbf{L}} B).
$$
由 $B$ 作为 $A$-模所满足的假设可知，
$M \otimes_A^{\mathbf{L}} B$ 的上同调只可能位于
$-d, -d + 1, \ldots, 0$ 这些次数。由于 $K^\bullet$ 具有
$[a, b]$ 内的 Tor 振幅，由 Example\ref{more-algebra-example-tor} 中的谱序列可知，
$K^\bullet \otimes_B^{\mathbf{L}} (M \otimes_A^{\mathbf{L}} B)$
的上同调只可能位于 $[-d + a, b]$，正如所需。
\end{proof}

\begin{lemma}
\label{more-algebra-lemma-pull-tor-amplitude}
设 $A \to B$ 为环映射。
设 $a, b \in \mathbf{Z}$。
设 $K^\bullet$ 为具有 $[a, b]$ 内 Tor 振幅的 $A$-模复形。
则 $K^\bullet \otimes_A^{\mathbf{L}} B$ 作为 $B$-模复形
具有 $[a, b]$ 内的 Tor 振幅。
\end{lemma}

\begin{proof}
由引理\ref{more-algebra-lemma-tor-amplitude}，可找到拟同构
$E^\bullet \to K^\bullet$，其中 $E^\bullet$ 是由平坦 $A$-模组成且当
$i \not \in [a, b]$ 时 $E^i = 0$ 的复形。于是按构造，
$E^\bullet \otimes_A B$ 计算 $K^\bullet \otimes_A ^{\mathbf{L}} B$，
并且由 Algebra，引理\ref{algebra-lemma-flat-base-change}，
每个 $E^i \otimes_A B$ 都是平坦 $B$-模。
因此由引理\ref{more-algebra-lemma-tor-amplitude} 得到结论。
\end{proof}

\begin{lemma}
\label{more-algebra-lemma-flat-base-change-finite-tor-dimension}
设 $A \to B$ 为平坦环映射。设 $d \geq 0$。
设 $M$ 为 Tor 维数 $\leq d$ 的 $A$-模。
则 $M \otimes_A B$ 是 Tor 维数 $\leq d$ 的 $B$-模。
\end{lemma}

\begin{proof}
这是引理\ref{more-algebra-lemma-pull-tor-amplitude} 的直接推论，
以及 $M \otimes_A^{\mathbf{L}} B = M \otimes_A B$ 这一事实的结果，
其中使用了 $B$ 在 $A$ 上平坦。
\end{proof}

\begin{lemma}
\label{more-algebra-lemma-tor-amplitude-localization}
设 $A  \to B$ 为环映射。设 $K^\bullet$ 为由 $B$-模组成的复形。
设 $a, b \in \mathbf{Z}$。下列条件等价：
\begin{enumerate}
\item $K^\bullet$ 作为 $A$-模复形具有 $[a, b]$ 内的 Tor 振幅；
\item 对每个素理想 $\mathfrak q \subset B$，令
$\mathfrak p = A \cap \mathfrak q$，则 $K^\bullet_\mathfrak q$
作为 $A_\mathfrak p$-模复形具有 $[a, b]$ 内的 Tor 振幅；
\item 对每个极大理想 $\mathfrak m \subset B$，令
$\mathfrak p = A \cap \mathfrak m$，则 $K^\bullet_\mathfrak m$
作为 $A_\mathfrak p$-模复形具有 $[a, b]$ 内的 Tor 振幅。
\end{enumerate}
\end{lemma}

\begin{proof}
设 (3) 成立，令 $M$ 为 $A$-模。则
$H^i = H^i(K^\bullet \otimes_A^\mathbf{L} M)$ 是 $B$-模，且
$(H^i)_\mathfrak m =
H^i(K^\bullet_\mathfrak m \otimes_{A_\mathfrak p}^\mathbf{L} M_\mathfrak p)$。
因此由 Algebra，引理\ref{algebra-lemma-characterize-zero-local}，
当 $i \not \in [a, b]$ 时 $H^i = 0$。故
(3) $\Rightarrow$ (1)。我们略去 (1) $\Rightarrow$ (2) 以及
(2) $\Rightarrow$ (3) 的证明。
\end{proof}

\begin{lemma}
\label{more-algebra-lemma-glue-tor-amplitude}
设 $R$ 为环。设 $f_1, \ldots, f_r \in R$ 为生成单位理想的一组元素。
设 $a, b \in \mathbf{Z}$。设 $K^\bullet$ 为由 $R$-模组成的复形。
若对每个 $i$，复形 $K^\bullet \otimes_R R_{f_i}$ 具有
$[a, b]$ 内的 Tor 振幅，则 $K^\bullet$ 具有 $[a, b]$ 内的 Tor 振幅。
\end{lemma}

\begin{proof}
这立即由引理\ref{more-algebra-lemma-tor-amplitude-localization} 得到，但也可直接如下看出。
注意 $- \otimes_R R_{f_i}$ 是正合函子，因此
$$
H^i(K^\bullet)_{f_i} =
H^i(K^\bullet) \otimes_R R_{f_i} = H^i(K^\bullet \otimes_R R_{f_i}).
$$
类似地，对每个 $R$-模 $M$，有
$$
H^i(K^\bullet \otimes_R^{\mathbf{L}} M)_{f_i} =
H^i(K^\bullet \otimes_R^{\mathbf{L}} M) \otimes_R R_{f_i} =
H^i(K^\bullet \otimes_R R_{f_i} \otimes_{R_{f_i}}^{\mathbf{L}} M_{f_i}).
$$
因此结果来自如下事实：$R$-模 $N$ 为零当且仅当对每个 $i$，
$N_{f_i}$ 为零，见 Algebra，引理\ref{algebra-lemma-cover}。
\end{proof}

\begin{lemma}
\label{more-algebra-lemma-flat-descent-tor-amplitude}
设 $R$ 为环。设 $a, b \in \mathbf{Z}$。设 $K^\bullet$ 为由 $R$-模组成的复形。
设 $R \to R'$ 为忠实平坦环映射。若复形 $K^\bullet \otimes_R R'$
具有 $[a, b]$ 内的 Tor 振幅，则 $K^\bullet$ 具有 $[a, b]$ 内的 Tor 振幅。
\end{lemma}

\begin{proof}
令 $M$ 为 $R$-模。由于 $R \to R'$ 平坦，有
$$
(M \otimes_R^{\mathbf{L}} K^\bullet) \otimes_R R'
=
((M \otimes_R R') \otimes_{R'}^{\mathbf{L}} (K^\bullet \otimes_R R')
$$
并且取上同调与张量 $R'$ 可交换。
因此
$\text{Tor}_i^R(M, K^\bullet) \otimes_R R' =
\text{Tor}_i^{R'}(M \otimes_R R', K^\bullet \otimes_R R')$。
由于 $R \to R'$ 忠实平坦，当
$i \not \in [a, b]$ 时
$\text{Tor}_i^{R'}(M \otimes_R R', K^\bullet \otimes_R R')$ 的消灭
蕴含 $\text{Tor}_i^R(M, K^\bullet)$ 也消灭。
\end{proof}

\begin{lemma}
\label{more-algebra-lemma-no-change-tor-amplitude}
给定环映射 $R \to A \to B$，其中 $A \to B$ 忠实平坦，且
$K \in D(A)$，则 $K$ 在 $R$ 上的 Tor 振幅与
$K \otimes_A^\mathbf{L} B$ 在 $R$ 上的 Tor 振幅相同。
\end{lemma}

\begin{proof}
这是因为对任意 $R$-模 $M$，都有
$H^i(K \otimes_R^\mathbf{L} M) \otimes_A B =
H^i((K \otimes_A^\mathbf{L} B) \otimes_R^\mathbf{L} M)$
对所有 $i$ 成立。
具体地，用一个 $A$-模复形 $K^\bullet$ 表示 $K$，并选取自由解消
$F^\bullet \to M$。则有等式
$$
\text{Tot}(K^\bullet \otimes_A B \otimes_R F^\bullet) =
\text{Tot}(K^\bullet \otimes_R F^\bullet) \otimes_A B
$$
左端的上同调群为
$H^i((K \otimes_A^\mathbf{L} B) \otimes_R^\mathbf{L} M)$，
而右端得到
$H^i(K \otimes_R^\mathbf{L} M) \otimes_A B$。
\end{proof}

\begin{lemma}
\label{more-algebra-lemma-finite-gl-dim-tor-dimension}
设 $R$ 为有限整体维数为 $d$ 的环。则
\begin{enumerate}
\item 每个模都具有 Tor 维数 $\leq d$；
\item 若由 $R$-模组成的复形 $K^\bullet$ 仅在
$i \in [a, b]$ 时满足 $H^i(K^\bullet) \not = 0$，
则它具有 $[a - d, b]$ 内的 Tor 振幅；
\item 由 $R$-模组成的复形 $K^\bullet$ 具有有限 Tor 维数，当且仅当
$K^\bullet \in D^b(R)$。
\end{enumerate}
\end{lemma}

\begin{proof}
$R$ 的假设意味着每个模都有长度至多为 $d$ 的有限投射解消，
特别地，每个模的 Tor 维数 $\leq d$。第二个断言由
引理\ref{more-algebra-lemma-cohomology-tor-amplitude} 及定义得到。第三个断言是第二个的改述。
\end{proof}

\begin{lemma}
\label{more-algebra-lemma-check-tor-dimension-modulo-nilpotent}
设 $R' \to R$ 为满射环映射，其核为幂零理想。
设 $K' \in D(R')$，并置 $K = K' \otimes_{R'}^\mathbf{L} R$。
设 $a, b \in \mathbf{Z}$。
则 $K$ 具有 $[a, b]$ 内的 Tor 振幅，当且仅当 $K'$ 具有该振幅。
\end{lemma}

\begin{proof}
一个方向由引理\ref{more-algebra-lemma-pull-tor-amplitude} 得到。
反之，设 $K$ 具有 $[a, b]$ 内的 Tor 振幅，令 $M'$ 为 $R'$-模。
需证明 $K' \otimes_{R'}^\mathbf{L} M'$ 的非零上同调只出现在
区间 $[a, b]$ 内。

\medskip\noindent
令 $I = \Ker(R' \to R)$。则某个 $n$ 使 $I^n = 0$。若 $IM' = 0$，
则可将 $M'$ 视为 $R$-模，并如下论证：
$$
K' \otimes_{R'}^\mathbf{L} M' =
K' \otimes_{R'}^\mathbf{L} (R \otimes_R^\mathbf{L} M') =
(K' \otimes_{R'}^\mathbf{L} R) \otimes_R^\mathbf{L} M' =
K \otimes_R^\mathbf{L} M'
$$
由 $K$ 的假设，其上同调只在区间 $[a, b]$ 内非零。若
$I^{t + 1}M' = 0$，则考虑短正合列
$$
0 \to IM' \to M' \to M'/IM' \to 0
$$
对 $t$ 作归纳可知，$K' \otimes_{R'}^\mathbf{L} IM'$ 与
$K' \otimes_{R'}^\mathbf{L} M'/IM'$ 的上同调都只在区间 $[a, b]$ 内非零。
于是可区别三角
$$
K' \otimes_{R'}^\mathbf{L} IM' \to
K' \otimes_{R'}^\mathbf{L} M' \to
K' \otimes_{R'}^\mathbf{L} M'/IM' \to
(K' \otimes_{R'}^\mathbf{L} IM')[1]
$$
证明 $K' \otimes_{R'}^\mathbf{L} M'$ 也具有所需性质。
\end{proof}

\section{Ext 的谱序列}
\label{more-algebra-section-spectral-sequence-ext}

\noindent
本节汇集考察 Ext 函子时出现的各种谱序列。对于环 $R$ 的导出范畴
$D(R)$ 中的任意一对对象 $L$、$K$，记
$$
\Ext^n_R(L, K) = \Hom_{D(R)}(L, K[n])
$$
这遵循我们在 Derived Categories，第
\ref{derived-section-ext} 节中的一般约定。

\medskip\noindent
若 $M$ 为 $R$-模且 $K \in D^+(R)$，则存在谱序列
\begin{equation}
\label{more-algebra-equation-first-ss-ext}
E_2^{i, j} = \Ext_R^i(M, H^j(K)) \Rightarrow \Ext_R^{i + j}(M, K)
\end{equation}
若 $K$ 由下有界的 $R$-模复形 $K^\bullet$ 表示，则还有谱序列
\begin{equation}
\label{more-algebra-equation-second-ss-ext}
E_1^{i, j} = \Ext_R^j(M, K^i) \Rightarrow \Ext_R^{i + j}(M, K)
\end{equation}
这些谱序列来自将 Derived Categories，引理
\ref{derived-lemma-two-ss-complex-functor} 应用于函子
$\Hom_R(M, -)$。

\section{投射维数}
\label{more-algebra-section-projective-dimension}

\noindent
我们在 Algebra，定义\ref{algebra-definition-finite-proj-dim} 中定义了模的投射维数。

\begin{definition}
\label{more-algebra-definition-projective-dimension}
设 $R$ 为环。设 $K$ 为 $D(R)$ 中的对象。若 $K$ 可由投射模组成的
有界复形表示，则称 $K$ 具有{\it 有限投射维数}。若 $K$ 拟同构于复形
$$
\ldots \to 0 \to P^a \to P^{a + 1} \to \ldots \to
P^{b - 1} \to P^b \to 0 \to \ldots
$$
则称 $K$ 具有{\it $[a, b]$ 内的投射振幅}，其中对所有
$i \in \mathbf{Z}$，$P^i$ 都是投射 $R$-模。
\end{definition}

\noindent
显然，$K$ 具有有限投射维数，当且仅当对某些
$a, b \in \mathbf{Z}$，$K$ 具有 $[a, b]$ 内的投射振幅。
此外，若 $K$ 具有有限投射维数，则 $K$ 是有界的。下面给出一个
用于检测 $D(R)$ 中这类对象的引理。

\begin{lemma}
\label{more-algebra-lemma-projective-amplitude}
设 $R$ 为环。设 $K$ 为 $D(R)$ 中的对象。设 $a, b \in \mathbf{Z}$。
下列条件等价：
\begin{enumerate}
\item $K$ 具有 $[a, b]$ 内的投射振幅；
\item 对所有 $R$-模 $N$ 以及所有 $i \not \in [-b, -a]$，
$\Ext^i_R(K, N) = 0$；
\item 当 $n > b$ 时 $H^n(K) = 0$，且对所有 $R$-模 $N$ 以及所有
$i > -a$，都有 $\Ext^i_R(K, N) = 0$；
\item 当 $n \not \in [a - 1, b]$ 时 $H^n(K) = 0$，且对所有
$R$-模 $N$，都有 $\Ext^{-a + 1}_R(K, N) = 0$。
\end{enumerate}
\end{lemma}

\begin{proof}
设 (1) 成立。不妨设 $K$ 是复形
$$
\ldots \to 0 \to P^a \to P^{a + 1} \to \ldots \to
P^{b - 1} \to P^b \to 0 \to \ldots
$$
其中对所有 $i \in \mathbf{Z}$，$P^i$ 都是投射 $R$-模。
此时可用复形
$$
\ldots \to 0 \to \Hom_R(P^b, N) \to \ldots \to
\Hom_R(P^a, N) \to 0 \to \ldots
$$
计算 Ext 群，从而得到 (2)。

\medskip\noindent
设 (2) 成立。选取单射 $H^n(K) \to I$，其中 $I$ 是单射 $R$-模。
由于 $\Hom_R(-, I)$ 是正合函子，可知
$\Ext^{-n}(K, I) = \Hom_R(H^n(K), I)$。
特别地，当 $n > b$ 时 $H^n(K)$ 为零。因此 (2) 推出 (3)。

\medskip\noindent
与 (2) 推出 (3) 时相同的论证表明 (3) 推出 (4)。

\medskip\noindent
设 (4) 成立。(2) 推出 (3) 的相同论证表明
$H^{a - 1}(K) = 0$，即 $H^i(K) = 0$，除非 $i \in [a, b]$。
特别地，$K$ 上有界，因此可选取一个表示 $K$ 的复形
$P^\bullet$，使得对所有 $i \in \mathbf{Z}$，$P^i$ 为投射
（例如自由）模，且当 $i > b$ 时 $P^i = 0$。
见 Derived Categories，引理
\ref{derived-lemma-subcategory-left-resolution}。
令 $Q = \Coker(P^{a - 1} \to P^a)$。则 $K$ 拟同构于复形
$$
\ldots \to 0 \to Q \to P^{a + 1} \to \ldots \to P^b \to 0 \to \ldots
$$
因为当 $i < a$ 时 $H^i(K) = 0$。
记 $K' = (P^{a + 1} \to \ldots \to P^b)$ 为 $D(R)$ 中相应的对象。
我们得到可区别三角
$$
K' \to K \to Q[-a] \to K'[1]
$$
位于 $D(R)$ 中。因此对每个 $R$-模 $N$，有正合列
$$
\Ext^{-a}(K', N) \to \Ext^1(Q, N) \to \Ext^{1 - a}(K, N)
$$
由假设，右端项消失；由蕴含 (1) $\Rightarrow$ (2)，左端项也消失。
因此由 Algebra，引理\ref{algebra-lemma-characterize-projective}，
$Q$ 是投射 $R$-模。于是 (1) 成立，证明完毕。
\end{proof}

\begin{example}
\label{more-algebra-example-ext-not-bounded}
设 $k$ 为域，令 $R$ 为 $k$ 上的对偶数环，即 $R = k[x]/(x^2)$。
记 $\epsilon \in R$ 为 $x$ 的类。令 $M = R/(\epsilon)$。
则 $M$ 拟同构于复形
$$
R \xrightarrow{\epsilon} R \xrightarrow{\epsilon} R \to \ldots
$$
但按照 Algebra，定义\ref{algebra-definition-finite-proj-dim} 中的定义，
$M$ 不具有有限投射维数。
这说明在有限投射维数的定义中，我们考虑 $D(R)$ 的对象所对应的
在两个方向上均有界的投射模复形。
\end{example}

\section{内射维数}
\label{more-algebra-section-injective-dimension}

\noindent
本节是投射维数一节的对偶版本。

\begin{definition}
\label{more-algebra-definition-injective-dimension}
设 $R$ 为环。设 $K$ 为 $D(R)$ 中的对象。
若 $K$ 可由内射 $R$-模组成的有限复形表示，则称 $K$ 具有
{\it 有限内射维数}。若 $K$ 同构于复形
$$
\ldots \to 0 \to I^a \to I^{a + 1} \to \ldots \to
I^{b - 1} \to I^b \to 0 \to \ldots
$$
则称 $K$ 具有{\it $[a, b]$ 内的内射振幅}，其中对所有
$i \in \mathbf{Z}$，$I^i$ 都是内射 $R$-模。
\end{definition}

\noindent
显然，$K$ 具有有界内射维数，当且仅当对某些
$a, b \in \mathbf{Z}$，$K$ 具有 $[a, b]$ 内的内射振幅。
此外，若 $K$ 具有有界内射维数，则 $K$ 是有界的。下面给出必需的引理。

\begin{lemma}
\label{more-algebra-lemma-injective-amplitude}
设 $R$ 为环。设 $K$ 为 $D(R)$ 中的对象。设 $a, b \in \mathbf{Z}$。
下列条件等价：
\begin{enumerate}
\item $K$ 具有 $[a, b]$ 内的内射振幅；
\item 对所有 $R$-模 $N$ 以及所有 $i \not \in [a, b]$，
$\Ext^i_R(N, K) = 0$；
\item 对所有理想 $I \subset R$ 以及所有 $i \not \in [a, b]$，
$\Ext^i(R/I, K) = 0$。
\end{enumerate}
\end{lemma}

\begin{proof}
设 (1) 成立。不妨设 $K$ 是复形
$$
\ldots \to 0 \to I^a \to I^{a + 1} \to \ldots \to
I^{b - 1} \to I^b \to 0 \to \ldots
$$
其中对所有 $i \in \mathbf{Z}$，$I^i$ 都是内射 $R$-模。
此时可用复形
$$
\ldots \to 0 \to \Hom_R(N, I^a) \to \ldots \to
\Hom_R(N, I^b) \to 0 \to \ldots
$$
计算 Ext 群，从而得到 (2)。(2) 显然推出 (3)。

\medskip\noindent
设 (3) 成立。选取非零映射 $R \to H^n(K)$。由于
$\Hom_R(R, -)$ 是正合函子，可知
$\Ext^n_R(R, K) = \Hom_R(R, H^n(K)) = H^n(K)$。
由此可知，当 $n \not \in [a, b]$ 时 $H^n(K)$ 为零。
特别地，$K$ 下有界，因此可选取拟同构
$$
K \to I^\bullet
$$
其中对所有 $i \in \mathbf{Z}$，$I^i$ 为内射模，且当 $i < a$ 时
$I^i = 0$。见 Derived Categories，引理
\ref{derived-lemma-subcategory-right-resolution}。
令 $J = \Ker(I^b \to I^{b + 1})$。则 $K$ 拟同构于复形
$$
\ldots \to 0 \to I^a \to \ldots \to I^{b - 1} \to J \to 0 \to \ldots
$$
记 $K' = (I^a \to \ldots \to I^{b - 1})$ 为 $D(R)$ 中相应的对象。
我们得到可区别三角
$$
J[-b] \to K \to K' \to J[1 - b]
$$
位于 $D(R)$ 中。因此对每个理想 $I \subset R$，有正合列
$$
\Ext^b(R/I, K') \to \Ext^1(R/I, J) \to \Ext^{1 + b}(R/I, K)
$$
由假设右端项消失；由蕴含 (1) $\Rightarrow$ (2)，左端项也消失。
因此由引理\ref{more-algebra-lemma-characterize-injective-bis}，$J$ 是内射 $R$-模。
\end{proof}

\begin{example}
\label{more-algebra-example-finite-injective-finite-global-dimension}
设 $R$ 为 Dedekind 整环。则每个非零理想 $I$ 都是有限投射模，
见引理\ref{more-algebra-lemma-dedekind-torsion-free-flat}。
因此 $R/I$ 的投射维数为 $1$。由引理\ref{more-algebra-lemma-injective-amplitude}，
每个 $R$-模 $M$ 的内射维数都满足 $\leq 1$。
于是对任意一对 $R$-模 $M, N$，当 $i \geq 2$ 时
$\Ext^i_R(M, N) = 0$。
因此 $D^b(R)$ 中任意对象 $K$ 都同构于其上同调的直和：
$K \cong \bigoplus H^i(K)[-i]$，见 Derived Categories，
引理\ref{derived-lemma-ext-2-zero}。
\end{example}

\begin{example}
\label{more-algebra-example-ext-not-bounded-reversed}
设 $k$ 为域，令 $R$ 为 $k$ 上的对偶数环，即 $R = k[x]/(x^2)$。
记 $\epsilon \in R$ 为 $x$ 的类。令 $M = R/(\epsilon)$。
则 $M$ 拟同构于复形
$$
\ldots \to R \xrightarrow{\epsilon} R \xrightarrow{\epsilon} R
$$
且 $R$ 是内射 $R$-模。然而在这种情形下，通常不认为 $M$ 具有有限内射维数。
这说明在有界内射维数的定义中，我们考虑 $D(R)$ 的对象所对应的
在两个方向上均有界的内射模复形。
\end{example}

\begin{lemma}
\label{more-algebra-lemma-finite-injective-dimension}
设 $R$ 为环。设 $K \in D(R)$。
\begin{enumerate}
\item 若 $K$ 属于 $D^b(R)$ 且对所有 $i$，$H^i(K)$ 都具有有限内射维数，
则 $K$ 具有有限内射维数；
\item 若表示 $K$ 的 $K^\bullet$ 是由 $R$-模组成的有界复形，
且对所有 $i$，$K^i$ 都具有有限内射维数，则 $K$ 具有有限内射维数。
\end{enumerate}
\end{lemma}

\begin{proof}
略去证明。提示：将 Derived Categories，引理
\ref{derived-lemma-two-ss-complex-functor} 应用于函子
$F = \Hom_R(N, -)$，以计算 $\Ext^i_A(N, K)$，并使用引理
\ref{more-algebra-lemma-injective-amplitude} 的判据。
\end{proof}

\begin{lemma}
\label{more-algebra-lemma-finite-injective-dimension-Noetherian-radical}
设 $R$ 为 Noether 环。设 $I \subset R$ 为包含于 $R$ 的 Jacobson 根中的理想。
设 $K \in D^+(R)$ 且具有有限上同调模。则下列条件等价：
\begin{enumerate}
\item $K$ 具有有限内射维数；
\item 存在某个 $b$，使得对任意理想 $J \supset I$，当 $i > b$ 时
$\Ext^i_R(R/J, K) = 0$。
\end{enumerate}
\end{lemma}

\begin{proof}
(1) $\Rightarrow$ (2) 立即成立。设 (2) 成立。
设 $H^i(K) = 0$ 当 $i < a$。则对所有 $R$-模 $M$，当 $i < a$ 时
$\Ext^i(M, K) = 0$。因此只需证明：对任意有限 $R$-模 $M$，当
$i > b$ 时 $\Ext^i(M, K) = 0$，见引理\ref{more-algebra-lemma-injective-amplitude}。
由 Algebra，引理\ref{algebra-lemma-filter-Noetherian-module}，
模 $M$ 有有限滤过，其相继商形如 $R/\mathfrak p$，其中
$\mathfrak p$ 为素理想。若
$0 \to M_1 \to M \to M_2 \to 0$ 是短正合列，且对 $j = 1, 2$，
当 $i > b$ 时 $\Ext^i(M_j, K) = 0$，则当 $i > b$ 时
$\Ext^i(M, K) = 0$。因此不妨设 $M = R/\mathfrak p$。
若 $I \subset \mathfrak p$，则消灭性由假设得到。否则选取
$f \in I$，且 $f \not \in \mathfrak p$。考虑短正合列
$$
0 \to R/\mathfrak p \xrightarrow{f} R/\mathfrak p \to R/(\mathfrak p, f) \to 0
$$
$R$-模 $R/(\mathfrak p, f)$ 有有限滤过，其相继商为
$R/\mathfrak q$，且 $(\mathfrak p, f) \subset \mathfrak q$。
因此由 Noether 归纳及上述论证，可设消灭性对 $R/(\mathfrak p, f)$ 成立。
另一方面，模
$E^i = \Ext^i(R/\mathfrak p, K)$
由关于 $K$ 的假设（下有界且具有有限上同调模）、谱序列
(\ref{more-algebra-equation-first-ss-ext}) 以及 Algebra，引理
\ref{algebra-lemma-ext-noetherian}，都是有限模。
于是当 $i > b$ 时，$E^i$ 是满足 $E^i/fE^i = 0$ 的有限 $R$-模。
由 Nakayama 引理（Algebra，引理\ref{algebra-lemma-NAK}），$E^i$ 为零。
\end{proof}

\begin{lemma}
\label{more-algebra-lemma-finite-injective-dimension-Noetherian-local}
设 $(R, \mathfrak m, \kappa)$ 为局部 Noether 环。
设 $K \in D^+(R)$ 且具有有限上同调模。
则下列条件等价：
\begin{enumerate}
\item $K$ 具有有限内射维数；
\item 当 $i \gg 0$ 时 $\Ext^i_R(\kappa, K) = 0$。
\end{enumerate}
\end{lemma}

\begin{proof}
这是引理\ref{more-algebra-lemma-finite-injective-dimension-Noetherian-radical} 的特例。
\end{proof}

\section{接近投射的模}
\label{more-algebra-section-near-projective}

\noindent
文献中似乎有许多不同的“几乎投射模”定义。本节只讨论其中一种可能性。

\begin{lemma}
\label{more-algebra-lemma-ideal-factor-through-projective}
设 $R$ 为环。设 $M$、$N$ 为 $R$-模。
\begin{enumerate}
\item 给定 $R$-模映射 $\varphi : M \to N$，下列条件等价：
(a) $\varphi$ 经由某个投射 $R$-模分解；(b) $\varphi$ 经由某个自由 $R$-模分解。
\item 满足 (1) 中等价条件的所有 $\varphi : M \to N$ 构成
$\Hom_R(M, N)$ 的一个 $R$-子模。
\item 给定映射 $\psi : M' \to M$ 和 $\xi : N \to N'$，若
$\varphi : M \to N$ 满足 (1) 的等价条件，则
$\xi \circ \varphi \circ \psi : M' \to N'$ 也满足。
\end{enumerate}
\end{lemma}

\begin{proof}
(1)(a) 与 (1)(b) 的等价性由 Algebra，引理
\ref{algebra-lemma-characterize-projective} 得到。
若 $\varphi : M \to N$ 和 $\varphi' : M \to N$ 分别经由模
$P$ 和 $P'$ 分解，则 $\varphi + \varphi'$ 经由
$P \oplus P'$ 分解，而 $\lambda \varphi$ 对所有 $\lambda \in R$ 经由
$P$ 分解。这证明了 (2)。若 $\varphi : M \to N$ 经由 $P$ 分解，且
$\psi$、$\xi$ 如 (3) 所述，则 $\xi \circ \varphi \circ \psi$ 也经由
$P$ 分解。这证明了 (3)。
\end{proof}

\begin{lemma}
\label{more-algebra-lemma-factor-through-finite-projective}
设 $R$ 为环。设 $\varphi : M \to N$ 为 $R$-模映射。
若 $\varphi$ 经由投射模分解且 $M$ 是有限 $R$-模，则 $\varphi$
经由有限投射模分解。
\end{lemma}

\begin{proof}
由引理\ref{more-algebra-lemma-ideal-factor-through-projective}，可将
$\varphi = \tau \circ \sigma$ 分解，其中 $\sigma$ 的目标为某个集合
$I$ 所指标的 $\bigoplus_{i \in I} R$。选取 $M$ 的生成元
$x_1, \ldots, x_n$。写成 $\sigma(x_j) = (a_{ji})_{i \in I}$。
对每个 $j$，只有有限多个 $a_{ij}$ 非零。
因此 $\sigma$ 的像包含于有限自由 $R$-模中，结论成立。
\end{proof}

\noindent
设 $R$ 为环。注意，一个 $R$-模是投射的，当且仅当 $R$ 上的恒等映射
经由某个投射模分解。

\begin{lemma}
\label{more-algebra-lemma-near-projective}
设 $R$ 为环。设 $I \subset R$ 为理想。设 $M$ 为 $R$-模。
下列条件等价：
\begin{enumerate}
\item 对每个 $a \in I$，映射 $a : M \to M$ 经由投射 $R$-模分解；
\item 对每个 $a \in I$，映射 $a : M \to M$ 经由自由 $R$-模分解；以及
\item 对每个 $R$-模 $N$，$\Ext^1_R(M, N)$ 被 $I$ 湮灭。
\end{enumerate}
\end{lemma}

\begin{proof}
(1) 与 (2) 的等价性由引理\ref{more-algebra-lemma-ideal-factor-through-projective} 得到。
若 (1) 成立，则 (3) 成立，因为对任意 $N$ 和任意投射模 $P$，
$\Ext^1_R(P, N)$ 为零。
反之，设 (3) 成立。选取短正合列
$0 \to N \to P \to M \to 0$，其中 $P$ 投射（甚至可取自由）。
由假设，该列在 $\Ext^1_R(M, N)$ 中的相应元素被 $I$ 湮灭。
因此对每个 $a \in I$，映射 $a : M \to M$ 可经由满射 $P \to M$
分解，从而 (1) 成立。
\end{proof}

\noindent
为了方便讨论满足引理\ref{more-algebra-lemma-near-projective} 中等价条件的模，
我们为这一性质命名。

\begin{definition}
\label{more-algebra-definition-near-projective}
设 $R$ 为环。设 $I \subset R$ 为理想。设 $M$ 为 $R$-模。
若引理\ref{more-algebra-lemma-near-projective} 的等价条件成立，则称 $M$ 为
{\it $I$-投射模}\footnote{这是非标准记号。}
\end{definition}

\noindent
被 $I$ 湮灭的模都是 $I$-投射模。

\begin{lemma}
\label{more-algebra-lemma-torsion-near-projective}
设 $R$ 为环。设 $I \subset R$ 为理想。设 $M$ 为 $R$-模。
若 $M$ 被 $I$ 湮灭，则 $M$ 是 $I$-投射模。
\end{lemma}

\begin{proof}
这直接由定义以及零模是投射模这一事实得到。
\end{proof}

\begin{lemma}
\label{more-algebra-lemma-ses-near-projective}
设 $R$ 为环。设 $I \subset R$ 为理想。设
$$
0 \to K \to P \to M \to 0
$$
为 $R$-模的短正合列。
若 $M$ 是 $I$-投射模且 $P$ 是投射模，则 $K$ 是 $I$-投射模。
\end{lemma}

\begin{proof}
元素 $\text{id}_K \in \Hom_R(K, K)$ 映射到给定扩张在
$\Ext^1_R(M, K)$ 中的类。由假设，该类被每个 $a \in I$ 湮灭，
因此 $a : K \to K$ 经由 $K \to P$ 分解，结论成立。
\end{proof}

\begin{lemma}
\label{more-algebra-lemma-dual-near-projective}
设 $R$ 为环。设 $I \subset R$ 为理想。
若 $M$ 是有限的 $I$-投射 $R$-模，则
$M^\vee = \Hom_R(M, R)$ 是 $I$-投射模。
\end{lemma}

\begin{proof}
设 $M$ 有限且为 $I$-投射模。选取短正合列
$0 \to K \to R^{\oplus r} \to M \to 0$。
这给出单射 $M^\vee \to R^{\oplus r} = (R^{\oplus r})^\vee$。
由于该短正合列在 $\Ext^1_R(M, K)$ 中对应的扩张类被 $I$ 湮灭，
可知对每个 $a \in I$，存在映射 $M \to R^{\oplus r}$，使其与给定映射
$R^{\oplus r} \to M$ 的复合等于 $a : M \to M$。
取对偶可知，$a : M^\vee \to M^\vee$ 经由上面给出的
$M^\vee \to R^{\oplus r}$ 分解，结论成立。
\end{proof}

\section{Hom 复形}
\label{more-algebra-section-hom-complexes}

\noindent
设 $R$ 为环。设 $L^\bullet$ 和 $M^\bullet$ 为由 $R$-模组成的两个复形。
我们构造复形 $\Hom^\bullet(L^\bullet, M^\bullet)$。具体地，对每个 $n$，置
$$
\Hom^n(L^\bullet, M^\bullet) =
\prod\nolimits_{n = p + q} \Hom_R(L^{-q}, M^p)
$$
可以将 $\Hom^n$ 看作从 $L^\bullet$ 到 $M^\bullet$（视为分次模）的
所有齐次次数为 $n$ 的 $R$-线性映射所成的 $R$-模。
在此术语下，用如下规则定义微分：
$$
\text{d}(f) = \text{d}_M \circ f - (-1)^n f \circ \text{d}_L
$$
其中 $f \in \Hom^n(L^\bullet, M^\bullet)$。
我们略去验证 $\text{d}^2 = 0$。符号规则见第
\ref{more-algebra-section-sign-rules} 节。
此构造是 Differential Graded Algebra，例
\ref{dga-example-category-complexes} 的特例。
由构造立即得到
\begin{equation}
\label{more-algebra-equation-cohomology-hom-complex}
H^n(\Hom^\bullet(L^\bullet, M^\bullet)) = \Hom_{K(R)}(L^\bullet, M^\bullet[n])
\end{equation}
对所有 $n \in \mathbf{Z}$ 都成立。

\begin{lemma}
\label{more-algebra-lemma-compose}
设 $R$ 为环。给定由 $R$-模组成的复形
$K^\bullet, L^\bullet, M^\bullet$，存在典范同构
$$
\Hom^\bullet(K^\bullet, \Hom^\bullet(L^\bullet, M^\bullet))
=
\Hom^\bullet(\text{Tot}(K^\bullet \otimes_R L^\bullet), M^\bullet)
$$
其为由 $R$-模组成的复形的同构。
\end{lemma}

\begin{proof}
设左端有一个次数为 $n$ 的元素 $\alpha$。于是
$$
\alpha = (\alpha^{p, q}) \in
\prod\nolimits_{p + q = n} \Hom_R(K^{-q}, \Hom^p(L^\bullet, M^\bullet))
$$
每个 $\alpha^{p, q}$ 都是元素
$$
\alpha^{p, q} = (\alpha^{r, s, q}) \in
\prod\nolimits_{r + s + q = n} \Hom_R(K^{-q}, \Hom_R(L^{-s}, M^r))
$$
若作如下恒等识别：
\begin{equation}
\label{more-algebra-equation-identification}
\Hom_R(K^{-q}, \Hom_R(L^{-s}, M^r)) = \Hom_R(K^{-q} \otimes_R L^{-s}, M^r)
\end{equation}
则按符号规则有
\begin{align*}
\text{d}(\alpha^{r, s, q})
& =
\text{d}_{\Hom^\bullet(L^\bullet, M^\bullet)} \circ \alpha^{r, s, q}
- (-1)^n \alpha^{r, s, q} \circ \text{d}_K \\
& =
\text{d}_M \circ \alpha^{r, s, q}
- (-1)^{r + s} \alpha^{r, s, q} \circ \text{d}_L
- (-1)^{r + s + q} \alpha^{r, s, q} \circ \text{d}_K
\end{align*}
另一方面，若 $\beta$ 是右端次数为 $n$ 的元素，则
$$
\beta = (\beta^{r, s, q}) \in \prod\nolimits_{r + s + q = n}
\Hom_R(K^{-q} \otimes_R L^{-s}, M^r)
$$
并且由符号规则（Homology，定义
\ref{homology-definition-associated-simple-complex}）有
\begin{align*}
\text{d}(\beta^{r, s, q})
& =
\text{d}_M \circ \beta^{r, s, q}
- (-1)^n \beta^{r, s, q} \circ
\text{d}_{\text{Tot}(K^\bullet \otimes L^\bullet)} \\
& =
\text{d}_M \circ \beta^{r, s, q}
- (-1)^{r + s + q} \left(
\beta^{r, s, q} \circ \text{d}_K + (-1)^{-q} \beta^{r, s, q} \circ \text{d}_L
\right)
\end{align*}
因此，由恒等识别 (\ref{more-algebra-equation-identification}) 所诱导的映射确实是复形态射。
\end{proof}

\begin{remark}
\label{more-algebra-remark-internal-hom}
设 $R$ 为环。由 $R$-模复形构成的范畴 $\text{Comp}(R)$
是对称幺半范畴，其张量积由 $\text{Tot}(- \otimes_R -)$ 给出，见
引理\ref{more-algebra-lemma-symmetric-monoidal-cat-complexes}。
给定 $\text{Comp}(R)$ 中的 $L^\bullet$ 和 $M^\bullet$，
元素 $f \in \Hom^0(L^\bullet, M^\bullet)$ 定义复形映射
$f : L^\bullet \to M^\bullet$，当且仅当 $\text{d}(f) = 0$。
因此引理\ref{more-algebra-lemma-compose} 还告诉我们
$$
\Mor_{\text{Comp}(R)}(K^\bullet, \Hom^\bullet(L^\bullet, M^\bullet))
=
\Mor_{\text{Comp}(R)}(\text{Tot}(K^\bullet \otimes_R L^\bullet), M^\bullet)
$$
对 $\text{Comp}(R)$ 中的 $K^\bullet, L^\bullet, M^\bullet$ 具有函子性。
这意味着 $\Hom^\bullet( - , -)$ 是对称幺半范畴
$\text{Comp}(R)$ 的内 Hom，正如 Categories，
注\ref{categories-remark-internal-hom-monoidal} 所讨论的。
\end{remark}

\begin{lemma}
\label{more-algebra-lemma-composition}
设 $R$ 为环。给定由 $R$-模组成的复形
$K^\bullet, L^\bullet, M^\bullet$，存在典范复形态射
$$
\text{Tot}\left(
\Hom^\bullet(L^\bullet, M^\bullet) \otimes_R \Hom^\bullet(K^\bullet, L^\bullet)
\right)
\longrightarrow
\Hom^\bullet(K^\bullet, M^\bullet)
$$
其为由 $R$-模组成的复形的态射。
\end{lemma}

\begin{proof}
由注\ref{more-algebra-remark-internal-hom} 中的讨论，这样的典范映射的存在性
由 Categories，注\ref{categories-remark-internal-hom-monoidal} 得到。
我们也给出直接构造。

\medskip\noindent
左端的次数为 $n$ 的元素 $\alpha$ 为
$$
\alpha = (\alpha^{p, q}) \in
\bigoplus\nolimits_{p + q = n}
\Hom^p(L^\bullet, M^\bullet) \otimes_R \Hom^q(K^\bullet, L^\bullet)
$$
元素 $\alpha^{p, q}$ 是有限和
$\alpha^{p, q} = \sum \beta^p_i \otimes \gamma^q_i$，
其中
$$
\beta^p_i = (\beta^{r, s}_i)
\in \prod\nolimits_{r + s = p} \Hom_R(L^{-s}, M^r)
$$
以及
$$
\gamma^q_i = (\gamma^{u, v}_i)
\in \prod\nolimits_{u + v = q} \Hom_R(K^{-v}, L^u)
$$
该映射将 $\alpha$ 送到 $\delta = (\delta^{r, v})$，其中
$$
\delta^{r, v} =
\sum\nolimits_{i, s}  \beta^{r, s}_i
\circ \gamma^{-s, v}_i
\in \Hom_R(K^{-v}, M^r)
$$
给定 $r + v = n$ 时，此和是有限的，因为非零的
$\alpha^{p, q}$ 只有有限多个，从而非零的 $\beta^p_i$ 和
$\gamma^q_i$ 也只有有限多个。
按符号规则有
\begin{align*}
\text{d}(\alpha^{p, q})
& =
\text{d}_{\Hom^\bullet(L^\bullet, M^\bullet)}(\alpha^{p, q})
+ (-1)^p \text{d}_{\Hom^\bullet(K^\bullet, L^\bullet)}(\alpha^{p, q}) \\
& =
\sum \Big( \text{d}_M \circ \beta^p_i \circ \gamma^q_i
- (-1)^p \beta^p_i \circ \text{d}_L \circ \gamma^q_i \Big) \\
&
\quad + (-1)^p \sum \Big( \beta^p_i \circ \text{d}_L \circ \gamma^q_i
- (-1)^q \beta^p_i \circ \gamma^q_i \circ \text{d}_K \Big) \\
& =
\sum \Big( \text{d}_M \circ \beta^p_i \circ \gamma^q_i
-(-1)^n \beta^p_i \circ \gamma^q_i \circ \text{d}_K \Big)
\end{align*}
因此规则 $\alpha \mapsto \delta$ 与微分相容，结论得证。
\end{proof}


\begin{lemma}
\label{more-algebra-lemma-diagonal-better}
设 $R$ 为一个环。给定复形 $K^\bullet, L^\bullet, M^\bullet$
（它们都是 $R$-模复形），存在一个典范态射
$$
\text{Tot}(K^\bullet \otimes_R \Hom^\bullet(M^\bullet, L^\bullet))
\longrightarrow
\Hom^\bullet(M^\bullet, \text{Tot}(K^\bullet \otimes_R L^\bullet))
$$
它是 $R$-模复形之间的态射，并且关于这三个复形具有函子性。
\end{lemma}

\begin{proof}
由注\ref{more-algebra-remark-internal-hom} 中的讨论，这样一个典范映射的存在性
由 Categories，注\ref{categories-remark-internal-hom-monoidal} 得到。
我们也给出直接构造。

\medskip\noindent
设 $\alpha$ 是右端中次数为 $n$ 的元素。于是
$$
\alpha = (\alpha^{p, q}) \in \prod\nolimits_{p + q = n}
\Hom_R(M^{-q}, \text{Tot}^p(K^\bullet \otimes_R L^\bullet))
$$
每个 $\alpha^{p, q}$ 都是元素
$$
\alpha^{p, q} = (\alpha^{r, s, q}) \in
\Hom_R(M^{-q}, \bigoplus\nolimits_{r + s + q = n} K^r \otimes_R L^s)
$$
其中我们把 $\alpha^{r, s, q}$ 看作一族映射，使得对每个
$x \in M^{-q}$，只有有限多个 $\alpha^{r, s, q}(x)$ 非零。
根据符号规则有
\begin{align*}
\text{d}(\alpha^{r, s, q})
& =
\text{d}_{\text{Tot}(K^\bullet \otimes_R L^\bullet)} \circ \alpha^{r, s, q}
- (-1)^n \alpha^{r, s, q} \circ \text{d}_M \\
& =
\text{d}_K \circ \alpha^{r, s, q} + (-1)^r \text{d}_L \circ \alpha^{r, s, q}
- (-1)^n \alpha^{r, s, q} \circ \text{d}_M
\end{align*}
另一方面，若 $\beta$ 是左端中次数为 $n$ 的元素，则
$$
\beta = (\beta^{p, q}) \in
\bigoplus\nolimits_{p + q = n} K^p \otimes_R \Hom^q(M^\bullet, L^\bullet)
$$
我们可以写成 $\beta^{p, q} = \sum \gamma_i^p \otimes \delta_i^q$，其中
$\gamma_i^p \in K^p$，且
$$
\delta_i^q = (\delta_i^{r, s}) \in
\prod\nolimits_{r + s = q} \Hom_R(M^{-s}, L^r)
$$
根据符号规则有
\begin{align*}
\text{d}(\beta^{p, q})
& =
\text{d}_K(\beta^{p, q}) +
(-1)^p \text{d}_{\Hom^\bullet(M^\bullet, L^\bullet)}(\beta^{p, q}) \\
& =
\sum \text{d}_K(\gamma_i^p) \otimes \delta_i^q +
(-1)^p \sum \gamma_i^p \otimes
(\text{d}_L \circ \delta_i^q - (-1)^q \delta_i^q \circ \text{d}_M)
\end{align*}
我们将元素 $\beta$ 送到 $\alpha$，其中
$$
\alpha^{r, s, q} = c^{r, s, q}(\sum \gamma_i^r \otimes \delta_i^{s, q})
$$
其中 $c^{r, s, q} : K^r \otimes_R \Hom_R(M^{-q}, L^s) \to
\Hom_R(M^{-q}, K^r \otimes_R L^s)$ 是典范映射。
对于给定的 $\beta$ 和 $r$，只有有限多个 $\gamma_i^r$ 非零，
因此（对给定的 $r$）只有有限多个 $\alpha^{r, s, q}$
非零。于是这族映射满足上述条件，且该映射定义良好。
比较符号可知，这与微分相容。
\end{proof}

\begin{lemma}
\label{more-algebra-lemma-diagonal}
设 $R$ 为一个环。给定 $R$-模复形 $K^\bullet, L^\bullet$，存在一个典范态射
$$
K^\bullet
\longrightarrow
\Hom^\bullet(L^\bullet, \text{Tot}(K^\bullet \otimes_R L^\bullet))
$$
它是 $R$-模复形之间的态射，并且关于这两个复形具有函子性。
\end{lemma}

\begin{proof}
由注\ref{more-algebra-remark-internal-hom} 中的讨论，这样一个典范映射的存在性
由 Categories，注\ref{categories-remark-internal-hom-monoidal} 得到。
我们也给出直接构造。

\medskip\noindent
设 $\alpha$ 是右端中次数为 $n$ 的元素。于是
$$
\alpha = (\alpha^{p, q}) \in \prod\nolimits_{p + q = n}
\Hom_R(L^{-q}, \text{Tot}^p(K^\bullet \otimes_R L^\bullet))
$$
每个 $\alpha^{p, q}$ 都是元素
$$
\alpha^{p, q} = (\alpha^{r, s, q}) \in
\Hom_R(L^{-q}, \bigoplus\nolimits_{r + s + q = n} K^r \otimes_R L^s)
$$
其中我们把 $\alpha^{r, s, q}$ 看作一族映射，使得对每个
$x \in L^{-q}$，只有有限多个 $\alpha^{r, s, q}(x)$ 非零。
根据符号规则有
\begin{align*}
\text{d}(\alpha^{r, s, q})
& =
\text{d}_{\text{Tot}(K^\bullet \otimes_R L^\bullet)} \circ \alpha^{r, s, q}
- (-1)^n \alpha^{r, s, q} \circ \text{d}_L \\
& =
\text{d}_K \circ \alpha^{r, s, q} + (-1)^r \text{d}_L \circ \alpha^{r, s, q}
- (-1)^n \alpha^{r, s, q} \circ \text{d}_L
\end{align*}
现在将元素 $\beta \in K^n$ 送到 $\alpha$，令
$\alpha^{n, -q, q} = \beta \otimes \text{id}_{L^{-q}}$，
并在 $r \not = n$ 时令 $\alpha^{r, s, q} = 0$。
这是上述意义下的元素，因为对于固定的 $q$，只有一个
$\alpha^{r, s, q}$ 非零。对微分的描述表明这与微分相容。
\end{proof}


\begin{lemma}
\label{more-algebra-lemma-evaluate-and-more}
设 $R$ 为一个环。给定 $R$-模复形 $K^\bullet, L^\bullet, M^\bullet$，存在一个典范态射
$$
\text{Tot}(\Hom^\bullet(L^\bullet, M^\bullet) \otimes_R K^\bullet)
\longrightarrow
\Hom^\bullet(\Hom^\bullet(K^\bullet, L^\bullet), M^\bullet)
$$
它是 $R$-模复形之间的态射，并且关于这三个复形具有函子性。
\end{lemma}

\begin{proof}
由注\ref{more-algebra-remark-internal-hom} 中的讨论，这样一个典范映射的存在性
由 Categories，注\ref{categories-remark-internal-hom-monoidal} 得到。
我们也给出直接构造。

\medskip\noindent
考虑右端中次数为 $n$ 的元素 $\beta$。则
$$
\beta = (\beta^{p, s}) \in
\prod\nolimits_{p + s = n} \Hom_R(\Hom^{-s}(K^\bullet, L^\bullet), M^p)
$$
根据符号规则，
\begin{align*}
\text{d}(\beta^{p, s})
& =
\text{d}_M \circ \beta^{p, s}
- (-1)^n \beta^{p, s} \circ
\text{d}_{\Hom^\bullet(K^\bullet, L^\bullet)}
\end{align*}
我们可以如下描述最后一项：
$$
(\beta^{p, s} \circ \text{d}_{\Hom^\bullet(K^\bullet, L^\bullet)})(f) =
\beta^{p, s}(\text{d}_L \circ f - (-1)^{s + 1} f \circ \text{d}_K)
$$
若 $f \in \Hom^{-s - 1}(K^\bullet, L^\bullet)$。因此，在某种未具体说明的意义下，
$\text{d}(\beta^{p, s})$ 是三个项之和，其符号如下：
\begin{equation}
\label{more-algebra-equation-beta}
\text{d}(\beta^{p, s}) =
\text{d}_M(\beta^{p, s})
-(-1)^n\text{d}_L(\beta^{p, s}) +
(-1)^{p + 1}\text{d}_K(\beta^{p, s})
\end{equation}

\noindent
接下来考虑左端中次数为 $n$ 的元素 $\alpha$。
我们可以将它写成
$$
\alpha = (\alpha^{t, r}) \in
\bigoplus\nolimits_{t + r = n} \Hom^t(L^\bullet, M^\bullet) \otimes K^r
$$
每个 $\alpha^{t, r}$ 映到元素
$$
\alpha^{t, r} \mapsto (\alpha^{p, q, r}) \in
\prod\nolimits_{p + q = t} \Hom_R(L^{-q}, M^p) \otimes_R K^r
$$
根据符号规则，
\begin{align*}
\text{d}(\alpha^{p, q, r})
& =
\text{d}_{\Hom^\bullet(L^\bullet, M^\bullet)}(\alpha^{p, q, r})
+ (-1)^{p + q} \text{d}_K(\alpha^{p, q, r})
\end{align*}
若进一步写 $\alpha^{p, q, r} = \sum g_i^{p, q} \otimes k_i^r$，则
$$
\text{d}_{\Hom^\bullet(L^\bullet, M^\bullet)}(\alpha^{p, q, r}) =
\sum (\text{d}_M \circ g_i^{p, q}) \otimes k_i^r
- (-1)^{p + q} \sum (g_i^{p, q} \circ \text{d}_L) \otimes k_i^r
$$
因此，在某种未具体说明的意义下，$\text{d}(\alpha^{p, q, r})$
是三个项之和，其符号如下：
\begin{equation}
\label{more-algebra-equation-alpha}
\text{d}(\alpha^{p, q, r}) =
\text{d}_M(\alpha^{p, q, r})
-(-1)^{p + q}\text{d}_L(\alpha^{p, q, r}) +
(-1)^{p + q}\text{d}_K(\alpha^{p, q, r})
\end{equation}

\noindent
为了定义我们的映射，将使用典范映射
$$
c_{p, q, r} : 
\Hom_R(L^{-q}, M^p) \otimes_R K^r
\longrightarrow
\Hom_R(\Hom_R(K^r, L^{-q}), M^p)
$$
它将 $\varphi \otimes k$ 送到映射 $\psi \mapsto \varphi(\psi(k))$。
该映射关于三个变量均具有函子性。令 $s = q + r$，有包含
$$
\Hom_R(\Hom_R(K^r, L^{-q}), M^p) \subset
\Hom_R(\Hom^{-s}(K^\bullet, L^\bullet), M^p)
$$
这个包含来自投影
$\Hom^{-s}(K^\bullet, L^\bullet) \to \Hom_R(K^r, L^{-q})$。
由于 $\alpha^{p, q, r}$ 只对有限多个 $r$ 非零，
可知对给定的 $s$，满足 $q + r = s$ 的 $q, r$ 中只有有限多个。
因此可将 $\alpha$ 送到元素 $\beta$，其中
$$
\beta^{p, s} =
\sum\nolimits_{q + r = s} \epsilon_{p, q, r} c_{p, q, r}(\alpha^{p, q, r})
$$
% P02-E0247
其中上述和使用前面给出的包含映射，且
$\epsilon_{p, q, r} \in \{\pm 1\}$。比较
方程 (\ref{more-algebra-equation-beta}) 和 (\ref{more-algebra-equation-alpha}) 中的符号可知，
\begin{enumerate}
\item $\epsilon_{p, q, r} = \epsilon_{p + 1, q, r}$
\item $-(-1)^n\epsilon_{p, q, r} = -(-1)^{p + q}\epsilon_{p, q - 1, r}$
或等价地 $\epsilon_{p, q, r} = (-1)^r\epsilon_{p, q - 1, r}$
\item $(-1)^{p + 1}\epsilon_{p, q, r} = (-1)^{p + q}\epsilon_{p, q, r + 1}$
或等价地
$(-1)^{q + 1}\epsilon_{p, q, r} = \epsilon_{p, q, r + 1}$。
\end{enumerate}
一个好的解是取
% P02-E0248
$$
\epsilon_{p, r, s} = (-1)^{r + qr}
$$
这一符号的选取将在证明之后的备注中解释。
\end{proof}


\begin{remark}
\label{more-algebra-remark-sign-explanation}
现在说明，评价映射引理\ref{more-algebra-lemma-evaluate-and-more} 的直接构造中所使用的符号，
为何与由注\ref{more-algebra-remark-internal-hom} 以及 Categories，
注\ref{categories-remark-internal-hom-monoidal} 中的讨论所得到的构造相一致。
记
$- \otimes - = \text{Tot}(- \otimes_R -)$，并记
$hom(-, -) = \Hom^\bullet(-, -)$。由幺半范畴语言给出的构造要求在
$\text{Comp}(R)$ 中使用箭头
$$
hom(L^\bullet, M^\bullet) \otimes K^\bullet
\longrightarrow
hom(hom(K^\bullet, L^\bullet), M^\bullet)
$$
该箭头对应于箭头
$$
hom(L^\bullet, M^\bullet) \otimes K^\bullet \otimes hom(K^\bullet, L^\bullet)
\longrightarrow
M^\bullet
$$
后者先通过交换最后两个张量积的次序，再使用评价映射
$hom(K^\bullet, L^\bullet) \otimes K^\bullet \to L^\bullet$
和
% P02-E0249
$hom(L^\bullet, K^\bullet) \otimes L^\bullet \to M^\bullet$ 得到。
只有交换时会出现符号。具体而言，在同构
$$
K^\bullet \otimes hom(K^\bullet, L^\bullet)
\to
hom(K^\bullet, L^\bullet) \otimes K^\bullet
$$
中，在 $K^r \otimes_R \Hom_R(K^{-r'}, L^q)$ 上出现符号
$(-1)^{r(q + r')}$，见符号规则一节
\ref{more-algebra-section-sign-rules} 的条目
(\ref{more-algebra-item-flip-tensor-product})。
读者可以自行验证，由于我们用来描述进入内 Hom 的映射所采用的对应关系，
该符号只有在 $r = r'$ 时才有影响；此时得到
$(-1)^{r(q + r)} = (-1)^{r + qr}$，正如直接证明中所示。
\end{remark}


\section{符号规则}
\label{more-algebra-section-sign-rules}

\noindent
本节回顾目前使用的符号规则，并讨论它们的一些推论。
在环上的模复形范畴中讨论这些问题似乎也很合适，因为大多数有趣现象
已经在这一情形中出现。
我们诚挚希望读者不必使用本节中较为深奥的部分。

\medskip\noindent
在本节其余部分，固定一个环 $R$，并记
$M^\bullet$ 为一个 $R$-模复形，其微分为
$d^n_M : M^n \to M^{n + 1}$。
\begin{enumerate}
\item 第 $k$ 个平移复形 $M^\bullet[k]$ 的项为
$(M^\bullet[k])^n = M^{n + k}$，其微分为
$d_{M[k]}^n = (-1)^kd^{n + k}_M$，见
Homology，定义\ref{homology-definition-shift-cochain}。
\item 对于复形之间的映射 $f : M^\bullet \to N^\bullet$，我们定义
$f[k] : M^\bullet[k] \to N^\bullet[k]$，不引入符号，见
Homology，定义\ref{homology-definition-shift-cochain}。
\item 我们将 $H^n(M^\bullet[k])$ 与 $H^{n + k}(M^\bullet)$
视为相同，不引入符号，见
Homology，定义\ref{homology-definition-cohomology-shift}。
\item 复形短正合序列的边界映射按蛇引理定义，不引入符号，见
Homology，引理\ref{homology-lemma-long-exact-sequence-cochain}。
\item 与逐项分裂短正合序列
$0 \to K^\bullet \to L^\bullet \to M^\bullet \to 0$
相关的复形的三角为
$$
K^\bullet \to L^\bullet \to M^\bullet \to K^\bullet[1]
$$
其中，若 $s$ 和 $\pi$ 是相容的逐项分裂，则
$M^n \to K^{n + 1}$ 是映射 $\pi^{n + 1} \circ d^n_L \circ s^n$。
换言之，不引入符号。见
Derived Categories，定义\ref{derived-definition-distinguished-triangle}
及\ref{derived-definition-split-ses}。
\item 总复形 $\text{Tot}(M^\bullet \otimes_R N^\bullet)$
具有满足 Leibniz 规则的微分 $d$
$$
d(x \otimes y) = d(x) \otimes y + (-1)^{\deg(x)}x \otimes d(y)
$$
见 Homology，例
\ref{homology-example-double-complex-as-tensor-product-of} 和
Homology，定义\ref{homology-definition-associated-simple-complex}。
\item
\label{more-algebra-item-shift-tensor}
存在典范同构
$$
\text{Tot}(M^\bullet \otimes_R N^\bullet)[a + b] \to
\text{Tot}(M^\bullet[a] \otimes_R N^\bullet[b])
$$
其在直和项 $M^p \otimes_R N^q$ 上使用符号 $(-1)^{pb}$，
见 Homology，注\ref{homology-remark-shift-double-complex}。
通常考虑相应的平移映射会更方便：
$\text{Tot}(M^\bullet \otimes_R N^\bullet) \to
\text{Tot}(M^\bullet[a] \otimes_R N^\bullet[b])[-a - b]$。
\item 存在复形的典范同构
$$
\text{Tot}(
\text{Tot}(K^\bullet \otimes_R L^\bullet) \otimes_R M^\bullet) \to
\text{Tot}(K^\bullet \otimes_R \text{Tot}(L^\bullet \otimes_R M^\bullet))
$$
其定义不引入符号。见
\ref{more-algebra-section-symmetric-monoidal} 节。
\item
\label{more-algebra-item-flip-tensor-product}
存在典范同构
$$
\text{Tot}(L^\bullet \otimes_R M^\bullet) \to
\text{Tot}(M^\bullet \otimes_R L^\bullet)
$$
其在直和项 $L^p \otimes_R M^q$ 上使用符号 $(-1)^{pq}$。
见\ref{more-algebra-section-symmetric-monoidal} 节。
\end{enumerate}

在讨论关于 Hom 复形的符号约定之前，我们先依据上述约定
构造一个复形的对偶。

\begin{lemma}
\label{more-algebra-lemma-left-dual-module}
设 $R$ 为一个环，$M$ 为一个 $R$-模。设
$N, \eta, \epsilon$ 是 $R$-模的幺半范畴中 $M$ 的一个左对偶，见
Categories，定义\ref{categories-definition-dual}。则
\begin{enumerate}
\item $M$ 和 $N$ 都是有限投射 $R$-模；
\item 映射
$e : \Hom_R(M, R) \to N$，$\lambda \mapsto (\lambda \otimes 1)(\eta)$
是同构；
\item 对 $n \in N$ 和 $m \in M$ 有
$\epsilon(n, m) = e^{-1}(n)(m)$。
\end{enumerate}
\end{lemma}

\begin{proof}
假设意味着
$$
M \xrightarrow{\eta \otimes 1} M \otimes_R N \otimes_R M
\xrightarrow{1 \otimes \epsilon} M
\quad\text{且}\quad
N \xrightarrow{1 \otimes \eta} N \otimes_R M \otimes_R N
\xrightarrow{\epsilon \otimes 1} N
$$
均为恒等映射。可以选取有限自由模 $F$、一个 $R$-模映射
$F \to M$，以及 $\eta$ 的一个提升
$\tilde \eta : R \to F \otimes_R N$。得到交换图
$$
\xymatrix{
M \ar[rr]_-{\eta \otimes 1} \ar[rrd]_-{\tilde \eta \otimes 1} & &
M \otimes_R N \otimes_R M \ar[r]_-{1 \otimes \epsilon} &
M \\
& &
F \otimes_R N \otimes_R M \ar[u] \ar[r]^-{1 \otimes \epsilon} &
F \ar[u]
}
$$
这表明 $M$ 上的恒等映射经过一个有限自由模分解，因此
$M$ 是有限投射模。由对称性可知 $N$ 也是有限投射模。
这证明了 (1)。(2) 由 Categories，引理
\ref{categories-lemma-left-dual} 及其证明得到。
(3) 由该证明中的第一个等式得到。
\end{proof}

\begin{lemma}
\label{more-algebra-lemma-left-dual-complex}
设 $R$ 为一个环，$M^\bullet$ 为一个 $R$-模复形。
设 $N^\bullet, \eta, \epsilon$ 是 $R$-模复形的幺半范畴中
$M^\bullet$ 的一个左对偶。则
\begin{enumerate}
\item $M^\bullet$ 和 $N^\bullet$ 都有界；
\item $M^n$ 和 $N^n$ 都是有限投射 $R$-模；
\item 写成 $\epsilon = \sum \epsilon_n$，其中
$\epsilon_n : N^{-n} \otimes_R M^n \to R$，并写成
$\eta = \sum \eta_n$，其中
$\eta_n : R \to  M^n \otimes_R N^{-n}$，则
$(N^{-n}, \eta_n, \epsilon_n)$ 是引理
\ref{more-algebra-lemma-left-dual-module} 意义下 $M^n$ 的左对偶；
\item 微分 $d_N^n : N^n \to N^{n + 1}$ 等于映射
$$
N^n = \Hom_R(M^{-n}, R)
\xrightarrow{d_M^{-n - 1}}
\Hom_R(M^{-n - 1}, R) = N^{n + 1}
$$
的 $-(-1)^n$ 倍，其中等号是引理
\ref{more-algebra-lemma-left-dual-module} (2) 中给出的恒等。
\end{enumerate}
反之，给定由有限投射 $R$-模组成的有界复形 $M^\bullet$，
令 $N^n = \Hom_R(M^{-n}, R)$，微分如上；令
$\epsilon = \sum \epsilon_n$，其中
$\epsilon_n : N^{-n} \otimes_R M^n \to R$ 由评价给出；并令
$\eta = \sum \eta_n$，其中 $\eta_n : R \to M^n \otimes_R N^{-n}$
将 $1$ 映到
% P02-E0250
$\text{id}_{M_n}$，则得到 $R$-模复形的幺半范畴中
$M^\bullet$ 的一个左对偶。
\end{lemma}

\begin{proof}
由 Categories，定义\ref{categories-definition-dual}，
$(1 \otimes \epsilon) \circ (\eta \otimes 1) = \text{id}_{M^\bullet}$
且
$(\epsilon \otimes 1) \circ (1 \otimes \eta) = \text{id}_{N^\bullet}$，
因而立即有
$(1 \otimes \epsilon_n) \circ (\eta_n \otimes 1) = \text{id}_{M^n}$
以及
$(\epsilon_n \otimes 1) \circ (1 \otimes \eta_n) = \text{id}_{N^{-n}}$，
这证明了 (3)。由引理\ref{more-algebra-lemma-left-dual-module} 得到 (2)。
由于和 $\eta = \sum \eta_n$ 是有限的，得到 (1)。
由于 $\eta = \sum \eta_n$ 是复形映射
$R \to \text{Tot}(M^\bullet \otimes_R N^\bullet)$，由我们对
$\text{Tot}(M^\bullet \otimes_R N^\bullet)$ 的微分所选取的符号可知
$$
(d_M^{-n - 1} \otimes 1) \circ \eta_{-n - 1}  +
(-1)^n (1 \otimes d_N^{-n}) \circ \eta_{-n} = 0
$$
展开定义即证明了 (4)。要得到引理的最后一个陈述，只需将上述论证反向读取。
\end{proof}


\noindent
我们将使用引理\ref{more-algebra-lemma-left-dual-complex} 中对复形左对偶的描述，
% P02-E0251
作为 $\Hom$ 复形符号规则的动机。我们选择符号，使得
(\ref{more-algebra-item-compatible}) 成立。
我们继续讨论上述各种符号规则。
\begin{enumerate}
\setcounter{enumi}{9}
\item 给定复形 $K^\bullet$、$M^\bullet$，令
$\Hom^\bullet(M^\bullet, K^\bullet)$ 为如下复形，其项为
$$
\Hom^n(M^\bullet, K^\bullet) = \prod\nolimits_{n = p + q} \Hom_R(M^{-q}, K^p)
$$
微分由规则
$$
d(f) = d_K \circ f - (-1)^n f \circ d_M
$$
给出。
\item
\label{more-algebra-item-compatible}
上述选择满足：若 $M^\bullet$ 具有引理
\ref{more-algebra-lemma-left-dual-complex} 中的左对偶 $N^\bullet$，则存在典范同构
$$
\text{Tot}(K^\bullet \otimes_R N^\bullet)
\longrightarrow
\Hom^\bullet(M^\bullet, K^\bullet)
$$
其定义不引入符号，并通过
$N^q = \Hom_R(M^{-q}, R)$ 以及典范映射
$K^p \otimes_R \Hom_R(M^{-q}, R) \to \Hom_R(M^{-q}, K^p)$，
将直和项 $K^p \otimes_R N^q$ 送到直和项
$\Hom_R(M^{-q}, K^p)$。
\item 存在复合映射
$$
\text{Tot}(
\Hom^\bullet(L^\bullet, K^\bullet) \otimes_R
\Hom^\bullet(M^\bullet, L^\bullet))
\longrightarrow
\Hom^\bullet(M^\bullet, K^\bullet)
$$
其定义不引入符号，见引理\ref{more-algebra-lemma-composition}。
\item 存在典范同构
$$
\Hom^\bullet(K^\bullet, \Hom^\bullet(L^\bullet, M^\bullet))
=
\Hom^\bullet(\text{Tot}(K^\bullet \otimes_R L^\bullet), M^\bullet)
$$
其定义不引入符号，见引理\ref{more-algebra-lemma-compose}。
\item 存在典范映射
$$
\text{Tot}(K^\bullet \otimes_R \Hom^\bullet(M^\bullet, L^\bullet))
\longrightarrow
\Hom^\bullet(M^\bullet, \text{Tot}(K^\bullet \otimes_R L^\bullet))
$$
其定义不引入符号，见引理\ref{more-algebra-lemma-diagonal-better}。
\item 存在典范映射
$$
K^\bullet
\longrightarrow
\Hom^\bullet(L^\bullet, \text{Tot}(K^\bullet \otimes_R L^\bullet))
$$
其定义不引入符号，见引理\ref{more-algebra-lemma-diagonal}。
\item
由引理\ref{more-algebra-lemma-evaluate-and-more} 可得一个典范映射
% P02-E0252
$$
\text{Tot}(\Hom^\bullet(L^\bullet, M^\bullet) \otimes_R K^\bullet)
\longrightarrow
\Hom^\bullet(\Hom^\bullet(K^\bullet, L^\bullet), M^\bullet)
$$
它在 R-模
$\Hom_R(L^{-q}, M^p) \otimes_R K^r$ 上使用符号
$(-1)^{r + qr}$，其原因在备注\ref{more-algebra-remark-sign-explanation} 中解释。
\item
\label{more-algebra-item-evaluation}
取 $L^\bullet = M^\bullet$，并使用
$R \to \Hom^\bullet(M^\bullet, M^\bullet)$，上一项的映射变为评价映射
$$
ev : K^\bullet \longrightarrow
\Hom^\bullet(\Hom^\bullet(K^\bullet, M^\bullet), M^\bullet)
$$
它将 $x \in K^n$ 送到如下映射：将
$f \in \Hom^m(K^\bullet, M^\bullet)$ 送到 $(-1)^{nm}f(x)$。
\item
\label{more-algebra-item-shift-hom}
存在典范识别
$$
\Hom^\bullet(M^\bullet, K^\bullet)[a - b] \to
\Hom^\bullet(M^\bullet[b], K^\bullet[a])
$$
其中使用了符号。它定义为如下映射：相应的平移映射
$$
\Hom^\bullet(M^\bullet, K^\bullet) \to
\Hom^\bullet(M^\bullet[b], K^\bullet[a])[b - a]
$$
在模 $\Hom_R(M^{-q}, K^p)$ 上（其中 $p + q = n$）
使用符号 $(-1)^{nb}$。具体地，若
$f \in \Hom^n(M^\bullet, K^\bullet)$，则源上的
$$
d(f) = d_K \circ f - (-1)^n f \circ d_M
$$
而在目标上，$f$ 属于
$$
\left(\Hom^\bullet(M^\bullet[b], K^\bullet[a])[b - a]\right)^n =
\Hom^{n + b -a}(M^\bullet[b], K^\bullet[a])
$$
因此有
\begin{align*}
d(f) & =
(-1)^{b - a}
\left(d_{K[a]} \circ f - (-1)^{n + b - a} f \circ d_{M[b]}\right) \\
& =
(-1)^{b - a}
\left((-1)^a d_K \circ f - (-1)^{n + b - a} f \circ (-1)^b d_M \right) \\
& =
(-1)^b d_K \circ f - (-1)^{n + b} f \circ d_M
\end{align*}
由此可见，在次数 $n$ 中选取的符号 $(-1)^{nb}$
会针对这些微分产生一个复形映射。
\end{enumerate}


\section{导出 Hom}
\label{more-algebra-section-RHom}

\noindent
设 $R$ 为一个环。本节定义的导出 Hom 是一个函子
$$
D(R)^{opp} \times D(R) \longrightarrow D(R),\quad
(K, L) \longmapsto R\Hom_R(K, L)
$$
它是 $R$-模导出范畴中的内 Hom，亦即由公式
\begin{equation}
\label{more-algebra-equation-internal-hom}
\Hom_{D(R)}(K, R\Hom_R(L, M))
=
\Hom_{D(R)}(K \otimes_R^\mathbf{L} L, M)
\end{equation}
刻画。这里 $K, L, M$ 是 $D(R)$ 的对象。注意，根据 Yoneda
引理，这个公式在唯一同构意义下刻画这些对象。其构造如下。
选取表示 $M$ 的 $R$-模 K-内射复形 $I^\bullet$，选取表示 $L$
的复形 $L^\bullet$，并置
$$
R\Hom_R(L, M) = \Hom^\bullet(L^\bullet, I^\bullet)
$$
记号如 Hom 复形一节\ref{more-algebra-section-hom-complexes} 中所述。
这一构造的推广见 Differential Graded Algebra，
一节\ref{dga-section-restriction}。
由(\ref{more-algebra-equation-cohomology-hom-complex}) 以及
Derived Categories，引理\ref{derived-lemma-K-injective} 可知
% P02-E0253
\begin{equation}
\label{more-algebra-equation-h0-RHom}
H^n(R\Hom_R(L, M)) = \Hom_{D(R)}(L, M[n])
\end{equation}
对所有 $n \in \mathbf{Z}$ 成立。特别地，$D(R)$ 中的对象
$R\Hom_R(L, M)$ 是良定义的，即与 K-内射复形
$I^\bullet$ 的选取无关。

\begin{lemma}
\label{more-algebra-lemma-internal-hom}
设 $R$ 为一个环。设 $K, L, M$ 是 $D(R)$ 的对象。存在典范同构
$$
R\Hom_R(K, R\Hom_R(L, M)) = R\Hom_R(K \otimes_R^\mathbf{L} L, M)
$$
它位于 $D(R)$ 中，并且关于 $K, L, M$ 具有函子性；取 $H^0$
后恢复(\ref{more-algebra-equation-internal-hom})。
\end{lemma}

\begin{proof}
选取表示 $M$ 的 K-内射复形 $I^\bullet$，以及表示 $L$ 的
$R$-模 K-平坦复形 $L^\bullet$。对于任意 $R$-模复形 $K^\bullet$，有
$$
\Hom^\bullet(K^\bullet, \Hom^\bullet(L^\bullet, I^\bullet)) =
\Hom^\bullet(\text{Tot}(K^\bullet \otimes_R L^\bullet), I^\bullet)
$$
由引理\ref{more-algebra-lemma-compose} 得到。根据 $R\Hom$ 的定义，并且由于
$\text{Tot}(K^\bullet \otimes_R L^\bullet)$ 表示导出张量积，
即得本引理。
\end{proof}


\begin{lemma}
\label{more-algebra-lemma-RHom-out-of-projective}
设 $R$ 为一个环。设 $P^\bullet$ 是由投射 $R$-模组成的上有界复形，
设 $L^\bullet$ 是 $R$-模复形。则 $R\Hom_R(P^\bullet, L^\bullet)$
由复形 $\Hom^\bullet(P^\bullet, L^\bullet)$ 表示。
\end{lemma}

\begin{proof}
由(\ref{more-algebra-equation-cohomology-hom-complex}) 以及
Derived Categories，引理\ref{derived-lemma-morphisms-from-projective-complex}，
该复形的上同调群是“正确的”。因此，若选取拟同构
$L^\bullet \to I^\bullet$，其中 $I^\bullet$ 是 $R$-模的 K-内射复形，
则诱导映射
$$
\Hom^\bullet(P^\bullet, L^\bullet)
\longrightarrow
\Hom^\bullet(P^\bullet, I^\bullet)
$$
是拟同构。由于右端正是我们对
$R\Hom_R(P^\bullet, L^\bullet)$ 的定义，结论成立。
\end{proof}

\begin{lemma}
\label{more-algebra-lemma-internal-hom-evaluate}
设 $R$ 为一个环。设 $K, L, M$ 是 $D(R)$ 的对象。
存在典范态射
$$
R\Hom_R(L, M) \otimes_R^\mathbf{L} K
\longrightarrow
R\Hom_R(R\Hom_R(K, L), M)
$$
它位于 $D(R)$ 中，并且关于 $K$、$L$、$M$ 具有函子性。
\end{lemma}

\begin{proof}
选取表示 $M$ 的 K-内射复形 $I^\bullet$，
表示 $L$ 的 K-内射复形 $J^\bullet$，以及表示 $K$ 的
K-平坦复形 $K^\bullet$。该映射使用引理
\ref{more-algebra-lemma-evaluate-and-more} 的映射定义：
$$
\text{Tot}(\Hom^\bullet(J^\bullet, I^\bullet) \otimes_R K^\bullet)
\longrightarrow
\Hom^\bullet(\Hom^\bullet(K^\bullet, J^\bullet), I^\bullet)
$$
我们略去该映射关于 $D(R)$ 中三个对象均具有函子性的证明。
\end{proof}


\begin{lemma}
\label{more-algebra-lemma-internal-hom-composition}
设 $R$ 为一个环。给定 $D(R)$ 中的 $K, L, M$，存在典范态射
$$
R\Hom_R(L, M) \otimes_R^\mathbf{L} R\Hom_R(K, L) \longrightarrow R\Hom_R(K, M)
$$
它位于 $D(R)$ 中，并且关于 $K, L, M$ 具有函子性。
\end{lemma}

\begin{proof}
选取表示 $M$ 的 K-内射复形 $I^\bullet$，表示 $L$ 的
K-内射复形 $J^\bullet$，以及任意表示 $K$ 的 $R$-模复形 $K^\bullet$。
由引理\ref{more-algebra-lemma-composition}，存在复形映射
$$
\text{Tot}\left(
\Hom^\bullet(J^\bullet, I^\bullet) \otimes_R \Hom^\bullet(K^\bullet, J^\bullet)
\right)
\longrightarrow
\Hom^\bullet(K^\bullet, I^\bullet)
$$
$R$-模复形
$\Hom^\bullet(J^\bullet, I^\bullet)$、
$\Hom^\bullet(K^\bullet, J^\bullet)$ 和
$\Hom^\bullet(K^\bullet, I^\bullet)$ 分别表示
$R\Hom_R(L, M)$、$R\Hom_R(K, L)$ 和 $R\Hom_R(K, M)$。
若选取 K-平坦复形 $H^\bullet$ 以及拟同构
$H^\bullet \to \Hom^\bullet(K^\bullet, J^\bullet)$，则存在映射
$$
\text{Tot}\left(
\Hom^\bullet(J^\bullet, I^\bullet) \otimes_R H^\bullet
\right)
\longrightarrow
\text{Tot}\left(
\Hom^\bullet(J^\bullet, I^\bullet) \otimes_R \Hom^\bullet(K^\bullet, J^\bullet)
\right)
$$
其源表示
$R\Hom_R(L, M) \otimes_R^\mathbf{L} R\Hom_R(K, L)$。
复合上述两个显示的箭头即得到所需映射。我们略去证明该构造
具有函子性。
\end{proof}


\begin{lemma}
\label{more-algebra-lemma-internal-hom-diagonal-better}
设 $R$ 为一个环。给定 $D(R)$ 中的复形 $K, L, M$，
存在典范态射
$$
K \otimes_R^\mathbf{L} R\Hom_R(M, L)
\longrightarrow
R\Hom_R(M, K \otimes_R^\mathbf{L} L)
$$
它位于 $D(R)$ 中，并且关于 $K, L, M$ 具有函子性。
\end{lemma}

\begin{proof}
选取表示 $K$ 的 K-平坦复形 $K^\bullet$、表示 $L$ 的
K-内射复形 $I^\bullet$，以及任意表示 $M$ 的复形 $M^\bullet$。
选取拟同构
$\text{Tot}(K^\bullet \otimes_R I^\bullet) \to J^\bullet$，
其中 $J^\bullet$ 是 K-内射的。然后使用映射
$$
\text{Tot}\left(
K^\bullet \otimes_R \Hom^\bullet(M^\bullet, I^\bullet)
\right)
\to
\Hom^\bullet(M^\bullet, \text{Tot}(K^\bullet \otimes_R I^\bullet))
\to
\Hom^\bullet(M^\bullet, J^\bullet)
$$
其中第一个映射是引理\ref{more-algebra-lemma-diagonal-better} 中的映射。
\end{proof}

\begin{lemma}
\label{more-algebra-lemma-internal-hom-diagonal}
设 $R$ 为一个环。给定 $D(R)$ 中的复形 $K, L$，存在典范态射
$$
K \longrightarrow R\Hom_R(L, K \otimes_R^\mathbf{L} L)
$$
它位于 $D(R)$ 中，并且关于 $K$ 和 $L$ 具有函子性。
\end{lemma}

\begin{proof}
这是引理\ref{more-algebra-lemma-internal-hom-diagonal-better} 的特殊情形，
但我们也将直接证明。
选取表示 $K$ 的 K-平坦复形 $K^\bullet$ 和任意表示 $L$ 的复形
$L^\bullet$。选取拟同构
$\text{Tot}(K^\bullet \otimes_R L^\bullet) \to J^\bullet$，
其中 $J^\bullet$ 是 K-内射的。然后使用映射
$$
K^\bullet \to
\Hom^\bullet(L^\bullet, \text{Tot}(K^\bullet \otimes_R L^\bullet))
\to \Hom^\bullet(L^\bullet, J^\bullet)
$$
其中第一个映射是引理\ref{more-algebra-lemma-diagonal} 中的映射。
\end{proof}


\section{完美复形}
\label{more-algebra-section-perfect}

\noindent
完美复形是具有有限 Tor 维数的拟相干复形。
我们不会采用此作为定义，而是直接如下定义环上的完美复形。

\begin{definition}
\label{more-algebra-definition-perfect}
设 $R$ 为一个环。记 $D(R)$ 为 $R$-模阿贝尔范畴的导出范畴。
\begin{enumerate}
\item 若 $D(R)$ 中的对象 $K$ 拟同构于有限投射 $R$-模组成的有界复形，
则称其为{\it 完美}；
\item 若 $R$-模 $M$ 的 $M[0]$ 是 $D(R)$ 中的完美对象，
则称其为{\it 完美}。
\end{enumerate}
\end{definition}

\noindent
例如，在 Noether 环上，一个有限模是完美的，当且仅当它具有有限投射维数，见
引理\ref{more-algebra-lemma-perfect-module} 以及
Algebra，定义\ref{algebra-definition-finite-proj-dim}。

\begin{lemma}
\label{more-algebra-lemma-perfect}
设 $K^\bullet$ 是 $D(R)$ 中的对象。以下条件等价：
\begin{enumerate}
\item $K^\bullet$ 是完美的；
\item $K^\bullet$ 是拟相干的并且具有有限 Tor 维数。
\end{enumerate}
若 (1) 和 (2) 成立，且 $K^\bullet$ 的 Tor 振幅位于 $[a, b]$ 中，
则 $K^\bullet$ 拟同构于由有限投射 $R$-模组成的复形
$E^\bullet$，并且当 $i \not \in [a, b]$ 时 $E^i = 0$。
\end{lemma}

\begin{proof}
(1) 推出 (2) 是显然的，见引理
\ref{more-algebra-lemma-pseudo-coherent} 和\ref{more-algebra-lemma-tor-amplitude}。
设 (2) 成立且 $K^\bullet$ 的 Tor 振幅位于 $[a, b]$ 中。
特别地，当 $i > b$ 时 $H^i(K^\bullet) = 0$。
选取由有限自由 $R$-模组成的复形 $F^\bullet$，使得当 $i > b$ 时
$F^i = 0$，并取一个拟同构 $F^\bullet \to K^\bullet$
（引理\ref{more-algebra-lemma-pseudo-coherent}）。
置 $E^\bullet = \tau_{\geq a}F^\bullet$。注意，除 $E^a$ 是有限表示
$R$-模外，$E^i$ 都是有限自由的。由引理
\ref{more-algebra-lemma-last-one-flat}，$E^a$ 是平坦的。因此由
Algebra，引理\ref{algebra-lemma-finite-projective} 可知
$E^a$ 是有限投射的。
\end{proof}


\begin{lemma}
\label{more-algebra-lemma-perfect-module}
设 $M$ 是环 $R$ 上的模。以下条件等价：
\begin{enumerate}
\item $M$ 是完美模；
\item 存在解消
$$
0 \to F_d \to \ldots \to F_1 \to F_0 \to M \to 0
$$
其中每个 $F_i$ 都是有限投射 $R$-模。
\end{enumerate}
\end{lemma}

\begin{proof}
设 (2) 成立。令复形 $E^\bullet$ 满足 $E^{-i} = F_i$，
则其与 $M[0]$ 拟同构。因此 $M$ 是完美的。
反之，设 (1) 成立。由引理
\ref{more-algebra-lemma-perfect} 和\ref{more-algebra-lemma-n-pseudo-module}，可找到解消
$E^\bullet \to M$，其中 $E^{-i}$ 是有限自由 $R$-模。由
引理\ref{more-algebra-lemma-last-one-flat} 可知，对于充分大的某个 $d$，
% P02-E0254
$F_d = \Coker(E^{d - 1} \to E^d)$ 是平坦的。由
Algebra，引理\ref{algebra-lemma-finite-projective} 可知
$F_d$ 是有限投射的。因此
$$
0 \to F_d \to E^{-d+1} \to \ldots \to E^0 \to M \to 0
$$
是所需的解消。
\end{proof}

\begin{lemma}
\label{more-algebra-lemma-two-out-of-three-perfect}
设 $R$ 为一个环。设 $(K^\bullet, L^\bullet, M^\bullet, f, g, h)$
是 $D(R)$ 中的一个三角。若 $K^\bullet, L^\bullet, M^\bullet$
中任意两个是完美的，则第三个也是完美的。
\end{lemma}

\begin{proof}
结合引理\ref{more-algebra-lemma-perfect}、
\ref{more-algebra-lemma-two-out-of-three-pseudo-coherent} 以及
\ref{more-algebra-lemma-cone-tor-amplitude} 即得。
\end{proof}

\begin{lemma}
\label{more-algebra-lemma-summands-perfect}
设 $R$ 为一个环。若 $K^\bullet \oplus L^\bullet$ 是完美的，
则 $K^\bullet$ 和 $L^\bullet$ 也都是完美的。
\end{lemma}

\begin{proof}
由引理\ref{more-algebra-lemma-perfect}、
\ref{more-algebra-lemma-summands-pseudo-coherent} 以及
\ref{more-algebra-lemma-summands-tor-amplitude} 得到。
\end{proof}


\begin{lemma}
\label{more-algebra-lemma-complex-perfect-modules}
设 $R$ 为一个环。若 $K^\bullet$ 是由完美 $R$-模组成的有界复形，
则 $K^\bullet$ 是完美复形。
\end{lemma}

\begin{proof}
对有限复形的长度作归纳即可证明：使用引理
\ref{more-algebra-lemma-two-out-of-three-perfect} 以及朴素截断。
\end{proof}

\begin{lemma}
\label{more-algebra-lemma-cohomology-perfect}
设 $R$ 为一个环。若 $K^\bullet \in D^b(R)$ 且其所有上同调模都是完美的，
则 $K^\bullet$ 是完美的。
\end{lemma}

\begin{proof}
对有限复形的长度作归纳即可证明：使用引理
\ref{more-algebra-lemma-two-out-of-three-perfect} 以及典范截断。
\end{proof}

\begin{lemma}
\label{more-algebra-lemma-perfect-push-perfect}
设 $A \to B$ 为环映射。假设 $B$ 作为 $A$-模是完美的。
设 $K^\bullet$ 为完美的 $B$-模复形。则 $K^\bullet$
作为 $A$-模复形也是完美的。
\end{lemma}

\begin{proof}
使用引理\ref{more-algebra-lemma-perfect}，这可转化为关于拟相干模及
有限 Tor 维数模的相应结果。关于这些结果，见引理
\ref{more-algebra-lemma-finite-tor-dimension-push-tor-amplitude} 以及
\ref{more-algebra-lemma-finite-push-pseudo-coherent}。
\end{proof}

\begin{lemma}
\label{more-algebra-lemma-pull-perfect}
设 $A \to B$ 为环映射。
设 $K^\bullet$ 是由 $A$-模组成的完美复形。则
$K^\bullet \otimes_A^{\mathbf{L}} B$ 是由 $B$-模组成的完美复形。
\end{lemma}

\begin{proof}
使用引理\ref{more-algebra-lemma-perfect}，这可转化为关于拟相干模及
有限 Tor 维数模的相应结果。关于这些结果，见引理
\ref{more-algebra-lemma-pull-tor-amplitude} 以及
\ref{more-algebra-lemma-pull-pseudo-coherent}。
\end{proof}


\begin{lemma}
\label{more-algebra-lemma-flat-base-change-perfect}
设 $A \to B$ 为平坦环映射。设 $M$ 是完美 $A$-模。
则 $M \otimes_A B$ 是完美 $B$-模。
\end{lemma}

\begin{proof}
由引理\ref{more-algebra-lemma-perfect-module}，假设意味着 $M$ 具有由有限投射
% P02-E0255
$R$-模组成的有限解消 $F_\bullet$。由于 $A \to B$ 平坦，
复形 $F_\bullet \otimes_A B$ 是 $M \otimes_A B$ 的有限长度解消，
其各项为 $B$-上的有限投射模。因此 $M \otimes_A B$ 是完美的。
\end{proof}


\begin{lemma}
\label{more-algebra-lemma-tensor-perfect}
设 $R$ 为一个环。若 $K$ 和 $L$ 是 $D(R)$ 中的完美对象，
则 $K \otimes_R^\mathbf{L} L$ 也是完美对象。
\end{lemma}

\begin{proof}
可以如下使用定义证明。分别用由有限投射 $R$-模组成的有界复形
$K^\bullet$ 和 $L^\bullet$ 表示 $K$ 和 $L$。于是
$K \otimes_R^\mathbf{L} L$ 由有界复形
$\text{Tot}(K^\bullet \otimes_R L^\bullet)$ 表示。
该复形的项是模
% P02-E0256
$M^a \otimes_R L^b$ 的直和。由于 $M^a$ 和 $L^b$
是有限自由 $R$-模的直和因子，$M^a \otimes_R L^b$ 也是如此。因此
我们得出复形 $\text{Tot}(K^\bullet \otimes_R L^\bullet)$ 的各项
都是有限投射的。

\medskip\noindent
也可以使用引理\ref{more-algebra-lemma-perfect} 中对完美复形的刻画，
以及关于拟相干复形（引理\ref{more-algebra-lemma-tensor-pseudo-coherent}）
和 Tor 振幅（引理\ref{more-algebra-lemma-push-tor-amplitude}，取 $A = B = R$）
的相应引理来证明。
\end{proof}


\begin{lemma}
\label{more-algebra-lemma-glue-perfect}
设 $R$ 为一个环。设 $f_1, \ldots, f_r \in R$ 生成单位理想。
设 $K^\bullet$ 为 $R$-模复形。若对每个 $i$，复形
$K^\bullet \otimes_R R_{f_i}$ 都是完美的，则 $K^\bullet$ 是完美的。
\end{lemma}

\begin{proof}
使用引理\ref{more-algebra-lemma-perfect}，这可转化为关于拟相干模及
有限 Tor 维数模的相应结果。关于这些结果，见引理
\ref{more-algebra-lemma-glue-tor-amplitude} 以及
\ref{more-algebra-lemma-glue-pseudo-coherent}。
\end{proof}

\begin{lemma}
\label{more-algebra-lemma-flat-descent-perfect}
设 $R$ 为一个环。设 $K^\bullet$ 为 $R$-模复形。
设 $R \to R'$ 为忠实平坦环映射。若复形
$K^\bullet \otimes_R R'$ 是完美的，则 $K^\bullet$ 是完美的。
\end{lemma}

\begin{proof}
使用引理\ref{more-algebra-lemma-perfect}，这可转化为关于拟相干模及
有限 Tor 维数模的相应结果。关于这些结果，见引理
\ref{more-algebra-lemma-flat-descent-tor-amplitude} 以及
\ref{more-algebra-lemma-flat-descent-pseudo-coherent}。
\end{proof}


\begin{lemma}
\label{more-algebra-lemma-regular-perfect}
设 $R$ 为正则环。则
\begin{enumerate}
\item 一个 $R$-模是完美的，当且仅当它是有限 $R$-模；
\item $R$-模复形 $K^\bullet$ 是完美的，当且仅当
$K^\bullet \in D^b(R)$ 且每个 $H^i(K^\bullet)$ 都是有限 $R$-模。
\end{enumerate}
\end{lemma}

\begin{proof}
根据定义，任何完美 $R$-模都是有限的。反之，设 $M$ 是有限
$R$-模。选取解消
$$
\ldots \to F_2 \xrightarrow{d_2} F_1 \xrightarrow{d_1} F_0 \to M \to 0
$$
其中 $F_i$ 是有限自由 $R$-模
（Algebra，引理\ref{algebra-lemma-resolution-by-finite-free}）。
置 $M_i = \Ker(d_i)$。记
$U_i \subset \Spec(R)$ 为满足 $M_{i, \mathfrak p}$ 自由的素理想
$\mathfrak p$ 所成的集合；由 Algebra，引理
\ref{algebra-lemma-finitely-presented-localization-free}，$U_i$ 是开集。
% P02-E0257
有一个正合序列 $0 \to M_{i + 1} \to F_{i + 1} \to M_i \to 0$。
若 $\mathfrak p \in U_i$，则
$0 \to M_{i + 1, \mathfrak p} \to F_{i + 1, \mathfrak p} \to
M_{i, \mathfrak p} \to 0$ 分裂。因此 $M_{i + 1, \mathfrak p}$
是有限投射模，因而是自由模
（Algebra，引理\ref{algebra-lemma-finite-projective}）。
这说明 $U_i \subset U_{i + 1}$。我们断言
$\Spec(R) = \bigcup U_i$。事实上，对每个素理想
$\mathfrak p$，正则局部环 $R_\mathfrak p$ 由
Algebra，命题\ref{algebra-proposition-regular-finite-gl-dim} 知具有有限
整体维数。因此，当 $i \gg 0$ 时，例如由 Algebra，引理
\ref{algebra-lemma-independent-resolution}，$M_{i, \mathfrak p}$ 是有限投射的
（因而是自由的）。由于 $R$ 的谱是 Noether 的
（Algebra，引理\ref{algebra-lemma-Noetherian-topology}），我们得到
对某个 $n$ 有 $U_n = \Spec(R)$。于是由
Algebra，引理\ref{algebra-lemma-finite-projective}，$M_n$ 是投射
$R$-模。因此
$$
0 \to M_n \to F_n \to \ldots \to F_1 \to M \to 0
$$
是由有限投射模组成的有界解消，从而 $M$ 是完美的。
这证明了 (1)。

\medskip\noindent
设 $K^\bullet$ 是 $R$-模复形。若 $K^\bullet$ 完美，则它属于
$D^b(R)$，并且拟同构于由有限投射 $R$-模组成的有限复形，
所以每个 $H^i(K^\bullet)$ 都是有限 $R$-模（因为 $R$ 是
Noether 的）。反之，设 $K^\bullet$ 属于 $D^b(R)$，且每个
$H^i(K^\bullet)$ 都是有限 $R$-模。由 (1)，每个
$H^i(K^\bullet)$ 都是完美 $R$-模，因此由引理
\ref{more-algebra-lemma-cohomology-perfect}，$K^\bullet$ 是完美的。
\end{proof}
\begin{lemma}
\label{more-algebra-lemma-dual-perfect-complex}
设 $A$ 为一个环，设 $K \in D(A)$ 是完美对象。则
$K^\vee = R\Hom_A(K, A)$ 是一个完美复形，并且
$K \cong (K^\vee)^\vee$。存在函子性的同构
$$
L \otimes_A^\mathbf{L} K^\vee = R\Hom_A(K, L)
\quad\text{且}\quad
H^0(L \otimes_A^\mathbf{L} K^\vee) = \Ext_A^0(K, L)
$$
其中 $L \in D(A)$。
\end{lemma}

\begin{proof}
我们可以用有限投射 $A$-模组成的有界复形 $K^\bullet$ 表示 $K$。
由引理 \ref{more-algebra-lemma-RHom-out-of-projective}，对象 $K^\vee$ 由复形
$E^\bullet = \Hom^\bullet(K^\bullet, A)$ 表示。注意
$E^n = \Hom_A(K^{-n}, A)$ 同样是有限投射的。因此 $E^\bullet$
是由有限投射模组成的有界复形，从而 $K^\vee$ 是完美的。存在典范映射
$$
K^\bullet = \text{Tot}(\Hom^\bullet(A, A) \otimes_A K^\bullet)
\longrightarrow
\Hom^\bullet(\Hom^\bullet(K^\bullet, A), A)
$$
该映射在每个次数上（差一个符号）使用求值映射，参见引理
\ref{more-algebra-lemma-evaluate-and-more}。（符号规则参见
第 \ref{more-algebra-section-sign-rules} 节。）因此，由有限投射模的双对偶仍为自身，
该映射给出典范同构 $(K^\vee)^\vee \cong K$。

\medskip\noindent
第二个等式由第一个等式、引理 \ref{more-algebra-lemma-internal-hom} 以及
Derived Categories，见引理
\ref{derived-lemma-morphisms-from-projective-complex} 推出；还要结合
Ext 群的定义，参见 Derived Categories，第
\ref{derived-section-ext} 节。设 $L^\bullet$ 是表示 $L$ 的
$A$-模复形。根据第 \ref{more-algebra-section-sign-rules} 节的
第 \ref{more-algebra-item-compatible} 项，存在复形的典范同构
$$
\text{Tot}(L^\bullet \otimes_A E^\bullet)
\longrightarrow
\Hom^\bullet(K^\bullet, L^\bullet)
$$
这里使用了引理 \ref{more-algebra-lemma-left-dual-complex}，以说明在由 $A$-模组成的复形范畴中
$E^\bullet$ 是 $K^\bullet$ 的左对偶。于是得到第一个显示等式，证明完成。
\end{proof}

\begin{remark}
\label{more-algebra-remark-dual-perfect-complex}
设 $A$ 为一个环，设 $K \in D(A)$ 是完美对象。若 $K$ 的
tor-振幅为 $[a, b]$，则其对偶 $K^\vee$ 的 tor-振幅为 $[-b, -a]$。
这由引理 \ref{more-algebra-lemma-perfect} 的最后一个断言以及引理
\ref{more-algebra-lemma-dual-perfect-complex} 证明中对 $K^\vee$ 的描述得到。
\end{remark}

\begin{lemma}
\label{more-algebra-lemma-colim-and-lim-of-duals}
\begin{slogan}
完美对象系统的平凡对偶性。
\end{slogan}
设 $A$ 为一个环，设 $(K_n)_{n \in \mathbf{N}}$ 是 $D(A)$ 中完美对象
组成的系统。设 $K = \text{hocolim} K_n$ 是导出余极限
（Derived Categories，定义 \ref{derived-definition-derived-colimit}）。
则对 $D(A)$ 中任意对象 $E$，有
$$
R\Hom_A(K, E) = R\lim E \otimes^\mathbf{L}_A K_n^\vee
$$
其中 $(K_n^\vee)$ 是完美复形对偶组成的逆系统。
\end{lemma}

\begin{proof}
由引理 \ref{more-algebra-lemma-dual-perfect-complex}，有
$R\lim E \otimes^\mathbf{L}_A K_n^\vee =
R\lim R\Hom_A(K_n, E)$，
它嵌入如下区别三角
$$
R\lim R\Hom_A(K_n, E) \to
\prod R\Hom_A(K_n, E) \to
\prod R\Hom_A(K_n, E)
$$
由于 $K$ 同样嵌入区别三角
$\bigoplus K_n \to \bigoplus K_n \to K$，只需证明
$\prod R\Hom_A(K_n, E) = R\Hom_A(\bigoplus K_n, E)$。
这由（\ref{more-algebra-equation-internal-hom}）以及导出张量积与直和交换这一事实
形式地推出。
\end{proof}

\begin{lemma}
\label{more-algebra-lemma-colimit-perfect-complexes}
设 $R = \colim_{i \in I} R_i$ 是环的滤子余极限。
\begin{enumerate}
\item 给定 $D(R)$ 中的完美对象 $K$，存在 $i \in I$ 以及
$D(R_i)$ 中的完美对象 $K_i$，使得
$K \cong K_i \otimes_{R_i}^\mathbf{L} R$（在 $D(R)$ 中）。
\item 给定 $0 \in I$ 以及 $K_0, L_0 \in D(R_0)$，其中 $K_0$ 完美，
有
$$
\Hom_{D(R)}(K_0 \otimes_{R_0}^\mathbf{L} R, L_0 \otimes_{R_0}^\mathbf{L} R) =
\colim_{i \geq 0}
\Hom_{D(R_i)}(K_0 \otimes_{R_0}^\mathbf{L} R_i,
L_0 \otimes_{R_0}^\mathbf{L} R_i)
$$
\end{enumerate}
换言之，$R$ 上的完美复形三角范畴是 $R_i$ 上完美复形三角范畴的余极限。
\end{lemma}

\begin{proof}
我们将使用 Algebra，引理
\ref{algebra-lemma-module-map-property-in-colimit} 和
\ref{algebra-lemma-colimit-category-fp-modules} 的结果，以下不再特别说明。
这些引理特别表明，有限表示 $R$-模范畴是有限表示 $R_i$-模范畴的余极限。
由于有限投射模可以刻画为有限自由模的直和项
（Algebra，引理 \ref{algebra-lemma-finite-projective}），
有限投射模范畴也具有同样性质。这就由完美对象在 $D(R)$ 中的定义证明了
（1）。

\medskip\noindent
为证明（2），可用有限投射 $R_0$-模组成的有界复形 $K_0^\bullet$
表示 $K_0$。可用一个 K-平坦复形 $L_0^\bullet$ 表示 $L_0$
（引理 \ref{more-algebra-lemma-K-flat-resolution}）。于是
$$
\Hom_{D(R)}(K_0 \otimes_{R_0}^\mathbf{L} R, L_0 \otimes_{R_0}^\mathbf{L} R) =
\Hom_{K(R)}(K_0^\bullet \otimes_{R_0} R, L_0^\bullet \otimes_{R_0} R)
$$
由 Derived Categories，引理
\ref{derived-lemma-morphisms-from-projective-complex} 得到。对于 $\Hom$，将 $R$ 换成
$R_i$ 时也同样适用。由于右端只涉及有限多个项，并且
$$
\Hom_R(K_0^p \otimes_{R_0} R, L_0^q \otimes_{R_0} R) =
\colim_{i \geq 0}
\Hom_{R_i}(K_0^p \otimes_{R_0} R_i, L_0^q \otimes_{R_0} R_i)
$$
由证明开头所引的引理，再由于滤子余极限是正合的
（Algebra，引理 \ref{algebra-lemma-directed-colimit-exact}），
可得（2）。
\end{proof}
\section{复形的提升}
\label{more-algebra-section-lifting-complexes}

\noindent
设 $R$ 为一个环，设 $I \subset R$ 为一个理想。我们将考虑的提升问题如下。
假设给定 $D(R)$ 中的对象 $K$ 以及由 $R/I$-模组成的复形 $E^\bullet$，
且 $E^\bullet$ 在 $D(R)$ 中表示 $K \otimes_R^\mathbf{L} R/I$。
问题：是否存在一个由 $R$-模组成的复形 $P^\bullet$，提升 $E^\bullet$
并在 $D(R)$ 中表示 $K$？
一般来说这个问题的答案是否定的，但在良好情形下仍可有所作为。
我们先讨论无环复形的提升。

\begin{lemma}
\label{more-algebra-lemma-lift-acyclic-complex}
设 $R$ 为一个环，设 $I \subset R$ 为一个理想。设 $\mathcal{P}$
为一类 $R$-模。假设
\begin{enumerate}
\item 每个 $P \in \mathcal{P}$ 都是投射 $R$-模；
\item 若 $P_1 \in \mathcal{P}$ 且 $P_1 \oplus P_2 \in \mathcal{P}$，则
$P_2 \in \mathcal{P}$；
\item 若 $f : P_1 \to P_2$（其中 $P_1, P_2 \in \mathcal{P}$）模 $I$
满射，则 $f$ 是满射。
\end{enumerate}
则对任意项形如 $P/IP$（其中 $P \in \mathcal{P}$）的有界上无环复形
$E^\bullet$，存在一个项属于 $\mathcal{P}$ 的有界上无环复形
$P^\bullet$，提升 $E^\bullet$。
\end{lemma}

\begin{proof}
设 $E^i = 0$（当 $i > b$）。假设给定 $n$ 以及复形态射
$$
\xymatrix{
& & P^n \ar[r] \ar[d] & P^{n + 1} \ar[r] \ar[d] & \ldots \ar[r]
& P^b \ar[r] \ar[d] & 0 \ar[r] \ar[d] & \ldots \\
\ldots \ar[r] & E^{n - 1} \ar[r] &
E^n \ar[r] & E^{n + 1} \ar[r] & \ldots \ar[r] &
E^b \ar[r] & 0 \ar[r] & \ldots
}
$$
其中 $P^i \in \mathcal{P}$，且
$P^n \to P^{n + 1} \to \ldots \to P^b$ 在次数 $\geq n + 1$ 时无环，
并且竖直映射诱导同构 $P^i/IP^i \to E^i$。
在这种情况下，可以归纳地选取同构
$P^i = Z^i \oplus Z^{i + 1}$，使得映射 $P^i \to P^{i + 1}$
由
$Z^i \oplus Z^{i + 1} \to Z^{i + 1} \to Z^{i + 1} \oplus Z^{i + 2}$
给出。由性质（2）并作归纳论证可知 $Z^i \in \mathcal{P}$。
选取 $P^{n - 1} \in \mathcal{P}$ 以及同构
$P^{n - 1}/IP^{n - 1} \to E^{n - 1}$。由于 $P^{n - 1}$ 是投射的，且
$Z^n/IZ^n = \Im(E^{n - 1} \to E^n)$，
我们可以将映射 $P^{n - 1} \to E^{n - 1} \to E^n$ 提升为映射
$P^{n - 1} \to Z^n$。由性质（3)，映射 $P^{n - 1} \to Z^n$
是满射。因此，只需将 $P^{n - 1}$ 及刚构造的映射添加到 $P^n$ 的左侧，
就得到该图表的延拓。由于当 $n > b$ 时存在所需形式的图表，
对 $n$ 作归纳即可。
\end{proof}

\begin{lemma}
\label{more-algebra-lemma-lift-complex}
设 $R$ 为一个环，设 $I \subset R$ 为一个理想。设 $\mathcal{P}$
为一类 $R$-模。设 $K \in D(R)$，设 $E^\bullet$ 是由 $R/I$-模组成的
复形，表示 $K \otimes_R^\mathbf{L} R/I$。假设
\begin{enumerate}
\item 每个 $P \in \mathcal{P}$ 都是投射 $R$-模；
\item $P_1 \in \mathcal{P}$ 且 $P_1 \oplus P_2 \in \mathcal{P}$
当且仅当 $P_1, P_2 \in \mathcal{P}$；
\item 若 $f : P_1 \to P_2$（其中 $P_1, P_2 \in \mathcal{P}$）模 $I$
满射，则 $f$ 是满射；
\item $E^\bullet$ 有界上，且 $E^i$ 的形式为 $P/IP$，
其中 $P \in \mathcal{P}$；
\item $K$ 可由项属于 $\mathcal{P}$ 的有界上复形表示。
\end{enumerate}
则存在一个项属于 $\mathcal{P}$ 的有界上复形 $P^\bullet$，使得
$P^\bullet/IP^\bullet$ 与 $E^\bullet$ 同构，并在 $D(R)$ 中表示 $K$。
\end{lemma}

\begin{proof}
由假设（5），可用项属于 $\mathcal{P}$ 的有界上复形 $K^\bullet$ 表示 $K$。
于是 $K \otimes_R^\mathbf{L} R/I$ 由 $K^\bullet/IK^\bullet$ 表示。
由于根据（4），$E^\bullet$ 是由投射 $R/I$-模组成的有界上复形，
我们可以选取拟同构
$\delta : E^\bullet \to K^\bullet/IK^\bullet$
（Derived Categories，引理
\ref{derived-lemma-morphisms-from-projective-complex}）。
设 $C^\bullet$ 为 $\delta$ 的映射锥
（Derived Categories，定义 \ref{derived-definition-cone}）。
模 $C^i$ 是直和 $K^i/IK^i \oplus E^{i + 1}$，
因此由于（2）特别说明 $\mathcal{P}$ 对取直和封闭，
其形式为某个 $P/IP$（其中 $P \in \mathcal{P}$）。
% P02-E0258
由于 $C^\bullet$ 无环，可以应用引理
\ref{more-algebra-lemma-lift-acyclic-complex}，找到 $C^\bullet$ 的无环提升
$A^\bullet$。复形 $A^\bullet$ 有界上，且其项属于 $\mathcal{P}$。在
$$
\xymatrix{
K^\bullet \ar@{..>}[r] \ar[d] & A^\bullet \ar[d] \\
K^\bullet/IK^\bullet \ar[r] & C^\bullet \ar[r] & E^\bullet[1]
}
$$
中，根据 Derived Categories，引理
\ref{derived-lemma-morphisms-lift-projective}，可以找到使图表交换的虚线箭头。
下面将证明由（1）、（2）、（3）可知，对每个 $i$，映射
$K^i \to A^i$ 是直和项的包含映射。由性质（2），
$P^i = \Coker(K^i \to A^i)$ 属于 $\mathcal{P}$。
因此可取
$P^\bullet = \Coker(K^\bullet \to A^\bullet)[-1]$，从而得到结论。

\medskip\noindent
为完成证明，需证明如下命题：设 $f : P_1 \to P_2$，
其中 $P_1, P_2 \in \mathcal{P}$，且
$P_1/IP_1 \to P_2/IP_2$ 是带分裂的单射，其余核形式为某个
$P_3/IP_3$（其中 $P_3 \in \mathcal{P}$），则 $f$ 是带分裂的单射。
记 $E_i = P_i/IP_i$，则 $E_2 = E_1 \oplus E_3$。
由于 $P_2$ 是投射的，可以选取映射 $g : P_2 \to P_3$，
提升映射 $E_2 \to E_3$。由条件（3），映射 $g$ 是满射，
% P02-E0259
因而由于 $P_3$ 是投射的，它是可分裂的。置 $P_1' = \Ker(g)$，
并选取分裂 $P_2 = P'_1 \oplus P_3$。由（2），$P'_1 \in \mathcal{P}$。
我们不知道 $g \circ f = 0$，但可以考虑映射
$$
P_1 \xrightarrow{f} P_2 \xrightarrow{projection} P'_1
$$
模 $I$ 后的复合是同构。由于 $P'_1$ 是投射的，可以分裂
$P_1 = T \oplus P'_1$。若 $T = 0$，则完成证明，因为此时
$P_2 \to P'_1$ 是 $f$ 的分裂。由（2）可知 $T \in \mathcal{P}$。
模 $I$ 计算可得 $T/IT = 0$。由于 $0 \in \mathcal{P}$
（它是 $\mathcal{P}$ 中任意 $P$ 的直和项），可知映射 $0 \to T$
是满射，从而得到 $T = 0$，如所需。
\end{proof}
\begin{lemma}
\label{more-algebra-lemma-lift-complex-projectives}
设 $R$ 为一个环，设 $I \subset R$ 为一个理想。设 $E^\bullet$
为 $R/I$-模复形，设 $K$ 为 $D(R)$ 中的对象。假设
\begin{enumerate}
\item $E^\bullet$ 是由投射 $R/I$-模组成的有界上复形；
\item $K \otimes_R^\mathbf{L} R/I$ 在 $D(R/I)$ 中由
$E^\bullet$ 表示；
\item $I$ 是幂零理想。
\end{enumerate}
则存在一个由投射 $R$-模组成的有界上复形 $P^\bullet$，它在
$D(R)$ 中表示 $K$，使得 $P^\bullet \otimes_R R/I$
与 $E^\bullet$ 同构。
\end{lemma}

\begin{proof}
对所有投射 $R$-模组成的类 $\mathcal{P}$ 应用引理
\ref{more-algebra-lemma-lift-complex}。该引理的性质（1）和（2）是显然的。
性质（3）由 Nakayama 引理
（Algebra，引理 \ref{algebra-lemma-NAK}）得到。
性质（4）由投射 $R/I$-模可提升为投射 $R$-模这一事实得到，参见
Algebra，引理 \ref{algebra-lemma-lift-projective-module}。
为说明（5）成立，只需证明 $K$ 属于 $D^{-}(R)$。
由于 $E^\bullet$ 有界上，已知
$K \otimes_R^\mathbf{L} R/I$ 属于 $D^{-}(R/I)$；
结论由引理 \ref{more-algebra-lemma-check-tor-dimension-modulo-nilpotent} 得到。
\end{proof}

\begin{lemma}
\label{more-algebra-lemma-check-pseudo-coherent-modulo-nilpotent}
设 $R' \to R$ 是满的环映射，其核为幂零理想。设 $K' \in D(R')$，
并置 $K = K' \otimes_{R'}^\mathbf{L} R$。
则 $K$ 伪凝聚当且仅当 $K'$ 伪凝聚。
\end{lemma}

\begin{proof}
一个方向由引理 \ref{more-algebra-lemma-pull-pseudo-coherent} 得到。
对另一个方向，假设 $K$ 伪凝聚。由引理
\ref{more-algebra-lemma-pseudo-coherent}，可以用有限自由 $R$-模组成的有界上复形
$E^\bullet$ 表示 $K$。由引理
\ref{more-algebra-lemma-lift-complex-projectives}，可以用投射 $R'$-模组成的有界上复形
$P^\bullet$ 表示 $K'$，使得 $P^n \otimes_{R'} R = E^n$。
由 Nakayama 引理可知 $P^n$ 是有限自由的，从而 $K'$ 也是伪凝聚的。
\end{proof}

\begin{lemma}
\label{more-algebra-lemma-lift-complex-stably-frees}
设 $R$ 为一个环，设 $I \subset R$ 为一个理想。设 $E^\bullet$
为 $R/I$-模复形，设 $K$ 为 $D(R)$ 中的对象。假设
\begin{enumerate}
\item $E^\bullet$ 是由有限稳定自由 $R/I$-模组成的有界上复形；
\item $K \otimes_R^\mathbf{L} R/I$ 在 $D(R/I)$ 中由 $E^\bullet$ 表示；
% P02-E0260
\item $K^\bullet$ 是伪凝聚的；
\item $1 + I$ 中的每个元素都可逆。
\end{enumerate}
则存在一个由有限稳定自由 $R$-模组成的有界上复形 $P^\bullet$，
它在 $D(R)$ 中表示 $K$，使得 $P^\bullet \otimes_R R/I$
与 $E^\bullet$ 同构。此外，若 $E^i$ 自由，则 $P^i$ 自由。
\end{lemma}

\begin{proof}
对所有有限稳定自由 $R$-模组成的类 $\mathcal{P}$ 应用引理
\ref{more-algebra-lemma-lift-complex}。该引理的性质（1）是显然的。
性质（2）由引理 \ref{more-algebra-lemma-exact-category-stably-free} 得到。
性质（3）由 Nakayama 引理
（Algebra，引理 \ref{algebra-lemma-NAK}）得到。
性质（4）由有限稳定自由 $R/I$-模可提升为有限稳定自由 $R$-模这一事实
得到，参见引理 \ref{more-algebra-lemma-lift-stably-free}。
性质（5）成立，因为伪凝聚复形可由有限自由 $R$-模组成的有界上复形表示。
引理的最后一个断言由引理
\ref{more-algebra-lemma-isomorphic-finite-projective-lifts} 得到。
\end{proof}

\begin{lemma}
\label{more-algebra-lemma-lift-pseudo-coherent-from-residue-field}
设 $(R, \mathfrak m, \kappa)$ 是局部环，设 $K \in D(R)$ 伪凝聚。置
$d_i = \dim_\kappa H^i(K \otimes_R^\mathbf{L} \kappa)$。
则 $d_i < \infty$，并且对某个 $b \in \mathbf{Z}$，当 $i > b$ 时
$d_i = 0$。于是存在复形
$$
\ldots \to
R^{\oplus d_{b - 2}} \to
R^{\oplus d_{b - 1}} \to
R^{\oplus d_b} \to 0 \to \ldots
$$
在 $D(R)$ 中表示 $K$。此外，该复形在同构意义下唯一（！）。
\end{lemma}

\begin{proof}
注意，作为 $D(\kappa)$ 中的对象，
$K \otimes_R^\mathbf{L} \kappa$ 是伪凝聚的，参见引理
\ref{more-algebra-lemma-pull-pseudo-coherent}。因此其上同调空间是有限维的，
并在某个截断点以上消失。$D(\kappa)$ 的每个对象都在 $D(\kappa)$ 中
同构于一个微分为零的复形 $E^\bullet$。特别地，
$E^i \cong \kappa^{\oplus d_i}$ 是有限自由的。应用引理
\ref{more-algebra-lemma-lift-complex-stably-frees} 即得存在性。

\medskip\noindent
% P02-E0261
若有两个复形 $F^\bullet$ 和 $G^\bullet$，其中 $F^i$ 与 $G^i$
都是秩为 $d_i$ 的自由模，并且它们表示 $K$。
则可以选取复形映射 $\beta : F^\bullet \to G^\bullet$，表示同构
$F^\bullet \cong K \cong G^\bullet$，参见
Derived Categories，引理
\ref{derived-lemma-morphisms-from-projective-complex}。
所诱导的复形映射
$\beta \otimes 1 : F^\bullet \otimes_R^\mathbf{L} \kappa \to
G^\bullet \otimes_R^\mathbf{L} \kappa$
必为复形同构，因为
$F^\bullet \otimes_R^\mathbf{L} \kappa$ 和
$G^\bullet \otimes_R^\mathbf{L} \kappa$ 中的微分均为零。
因此 $\beta^i : F^i \to G^i$ 是有限自由 $R$-模之间的映射，
其模 $\mathfrak m$ 的约化是同构。故 $\beta^i$ 是同构，结论成立。
\end{proof}

\begin{lemma}
\label{more-algebra-lemma-lift-perfect-from-residue-field}
设 $R$ 为一个环，设 $\mathfrak p \subset R$ 为一个素理想，
设 $K \in D(R)$ 是完美的。置
$d_i =
\dim_{\kappa(\mathfrak p)} H^i(K \otimes_R^\mathbf{L} \kappa(\mathfrak p))$。
则 $d_i < \infty$，并且只有有限多个非零。
于是存在 $f \in R$、$f \not \in \mathfrak p$ 以及复形
$$
\ldots \to 0 \to R_f^{\oplus d_a} \to R_f^{\oplus d_{a + 1}} \to
\ldots \to
R_f^{\oplus d_{b - 1}} \to
R_f^{\oplus d_b} \to 0
\to \ldots
$$
在 $D(R_f)$ 中表示 $K \otimes_R^\mathbf{L} R_f$。
\end{lemma}

\begin{proof}
注意，作为 $D(\kappa(\mathfrak p))$ 中的对象，
$K \otimes_R^\mathbf{L} \kappa(\mathfrak p)$ 是完美的，参见引理
\ref{more-algebra-lemma-pull-perfect}。因此只有有限多个 $d_i$ 非零，且它们都是有限的。
应用引理 \ref{more-algebra-lemma-lift-pseudo-coherent-from-residue-field}，
可在局部环 $R_\mathfrak p$ 上得到一个具有所需形状、表示 $K$ 的复形。
对 $f \in R$、$f \not \in \mathfrak p$，有
$R_\mathfrak p = \colim R_f$
（Algebra，引理 \ref{algebra-lemma-localization-colimit}）。
由引理 \ref{more-algebra-lemma-colimit-perfect-complexes} 得到结论。
略去一些细节。
\end{proof}

\begin{lemma}
\label{more-algebra-lemma-compare-representatives-perfect}
设 $R$ 为一个环，设 $\mathfrak p \subset R$ 为一个素理想。设
$M^\bullet$ 与 $N^\bullet$ 是由有限投射 $R$-模组成的有界复形，
它们表示 $D(R)$ 中的同一个对象。则存在
$f \in R$、$f \not \in \mathfrak p$，使得存在复形的同构（！）
$$
M^\bullet_f \oplus P^\bullet \cong N^\bullet_f \oplus Q^\bullet
$$
% P02-E0262
其中 $P^\bullet$ 与 $Q^\bullet$ 是平凡复形的有限直和，即形如
$\ldots \to 0 \to R_f \xrightarrow{1} R_f \to 0 \to \ldots$
的复形（置于任意次数）。
\end{lemma}

\begin{proof}
若在局部化 $R_\mathfrak p$ 上有上述类型的同构，则利用
$R_\mathfrak p = \colim R_f$
（Algebra，引理 \ref{algebra-lemma-localization-colimit}），
可以将该同构下降为某个 $f$ 对应的 $R_f$ 上的同构。
因此可假设 $R$ 是局部环且 $\mathfrak p$ 是极大理想。
此时结论由表示完美对象的“极小”复形的唯一性得到，参见引理
\ref{more-algebra-lemma-lift-pseudo-coherent-from-residue-field}；还要使用任意复形都是
一个平凡复形与一个极小复形的直和这一事实
（Algebra，引理 \ref{algebra-lemma-add-trivial-complex}）。
\end{proof}
\begin{lemma}
\label{more-algebra-lemma-lift-complex-finite-projectives}
设 $R$ 为一个环，设 $I \subset R$ 为一个理想。设 $E^\bullet$
为 $R/I$-模复形，设 $K$ 为 $D(R)$ 中的对象。假设
\begin{enumerate}
\item $E^\bullet$ 是由有限投射 $R/I$-模组成的有界上复形；
\item $K \otimes_R^\mathbf{L} R/I$ 在 $D(R/I)$ 中由 $E^\bullet$ 表示；
\item $K$ 是伪凝聚的；
\item $(R, I)$ 是 Hensel 对。
\end{enumerate}
则存在一个由有限投射 $R$-模组成的有界上复形 $P^\bullet$，
它在 $D(R)$ 中表示 $K$，使得 $P^\bullet \otimes_R R/I$
与 $E^\bullet$ 同构。此外，若 $E^i$ 自由，则 $P^i$ 自由。
\end{lemma}

\begin{proof}
对所有有限投射 $R$-模组成的类 $\mathcal{P}$ 应用引理
\ref{more-algebra-lemma-lift-complex}。该引理的性质（1）和（2）是显然的。
性质（3）由 Nakayama 引理
（Algebra，引理 \ref{algebra-lemma-NAK}）得到。
性质（4）由有限投射 $R/I$-模可提升为有限投射 $R$-模这一事实得到，
参见引理 \ref{more-algebra-lemma-lift-finite-projective-module}。
性质（5）成立，因为伪凝聚复形可由有限自由 $R$-模组成的有界上复形表示。
因此引理 \ref{more-algebra-lemma-lift-complex} 可用，并得到所需的 $P^\bullet$。
引理的最后一个断言由引理
\ref{more-algebra-lemma-isomorphic-finite-projective-lifts} 得到。
\end{proof}



\section{复形的分裂}
\label{more-algebra-section-splitting}

\noindent
本节讨论使环的导出范畴中一个对象成为其截断之直和的条件。
我们的方法是在适当假设下使用下述引理，分裂典范区别三角
$$
\tau_{\leq i}K \to K \to \tau_{\geq i + 1}K \to (\tau_{\leq i}K)[1]
$$
（在 $D(R)$ 中），参见 Derived Categories，注
\ref{derived-remark-truncation-distinguished-triangle}。

\begin{lemma}
\label{more-algebra-lemma-splitting-unique}
设 $R$ 为一个环，设 $K$ 和 $L$ 是 $D(R)$ 中的对象。
假设 $L$ 的投射振幅为 $[a, b]$；例如，若 $L$ 是 Tor-振幅为
$[a, b]$ 的完美对象，则满足此条件。
\begin{enumerate}
\item 若对 $i \geq a$ 有 $H^i(K) = 0$，则
$\Hom_{D(R)}(L, K) = 0$。
\item 若对 $i \geq a + 1$ 有 $H^i(K) = 0$，则对任意给定的区别三角
$K \to M \to L \to K[1]$，存在 $D(R)$ 中的同构
$M \cong K \oplus L$，且与该区别三角中的映射相容。
\item 若对 $i \geq a$ 有 $H^i(K) = 0$，则（2）中的同构存在且唯一。
\end{enumerate}
\end{lemma}

\begin{proof}
$L$ 的投射振幅为 $[a, b]$ 这一假设意味着，我们可以用投射
$R$-模组成的复形 $L^\bullet$ 表示 $L$，且当
$i \not \in [a, b]$ 时 $L^i = 0$，参见定义
\ref{more-algebra-definition-projective-dimension}。
若 $L$ 是 Tor-振幅为 $[a, b]$ 的完美对象，则可以用有限投射
$R$-模组成的复形 $L^\bullet$ 表示 $L$，且当
$i \not \in [a, b]$ 时 $L^i = 0$，参见引理 \ref{more-algebra-lemma-perfect}。
若对 $i \geq a$ 有 $H^i(K) = 0$，则 $K$ 拟同构于
$\tau_{\leq a - 1}K$。因此可以用 $R$-模复形 $K^\bullet$ 表示 $K$，
且当 $i \geq a$ 时 $K^i = 0$。于是由 Derived Categories，引理
\ref{derived-lemma-morphisms-from-projective-complex} 得到
$$
\Hom_{D(R)}(L, K) = \Hom_{K(R)}(L^\bullet, K^\bullet) = 0
$$
这证明了（1）。在（2）的假设下，由（1）有
$\Hom_{D(R)}(L, K[1]) = 0$，所以根据 Derived Categories，引理
\ref{derived-lemma-split}，该区别三角是分裂的。
（3）中的唯一性由（1）得到。
\end{proof}

\begin{lemma}
\label{more-algebra-lemma-better-cut-complex-in-two}
设 $R$ 为一个环，设 $\mathfrak p \subset R$ 为一个素理想。
设 $K^\bullet$ 是伪凝聚 $R$-模复形。假设对某个
$i \in \mathbf{Z}$，映射
$$
H^i(K^\bullet) \otimes_R \kappa(\mathfrak p)
\longrightarrow
H^i(K^\bullet \otimes_R^{\mathbf{L}} \kappa(\mathfrak p))
$$
满射。则存在 $f \in R$、$f \not \in \mathfrak p$，使得
$\tau_{\geq i + 1}(K^\bullet \otimes_R R_f)$ 是 $D(R_f)$ 中
Tor-振幅为 $[i + 1, \infty]$ 的完美对象，并且在 $D(R_f)$ 中
存在典范同构
$$
K^\bullet \otimes_R R_f \cong
\tau_{\leq i}(K^\bullet \otimes_R R_f) \oplus
\tau_{\geq i + 1}(K^\bullet \otimes_R R_f)
$$
\end{lemma}

\begin{proof}
在本证明中，所有张量积均取在 $R$ 上，并记
$\kappa = \kappa(\mathfrak p)$。可假设 $K^\bullet$
是由有限自由 $R$-模组成的有界上复形。考察次数 $i$ 中的情形：
$$
\ldots \to K^{i - 1} \xrightarrow{d^{i - 1}} K^i \xrightarrow{d^i}
K^{i + 1} \to \ldots
$$
用下列公式定义 $0 \subset V \subset W \subset K^i \otimes \kappa$：
$$
V = \Im\left(
K^{i - 1} \otimes \kappa \to K^i \otimes \kappa
\right)
\quad\text{且}\quad
W = \Ker\left(
K^i \otimes \kappa \to K^{i + 1} \otimes \kappa
\right)
$$
置 $\dim(V) = r$、$\dim(W/V) = s$ 以及
$\dim(K^i \otimes \kappa/W) = t$。
可选取 $x_1, \ldots, x_r \in K^{i - 1}$，使它们在
$d^{i - 1}$ 下映为 $V$ 的一组基。由假设，可选取
$y_1, \ldots, y_s \in \Ker(d^i)$，其像构成 $W/V$ 的一组基。
最后，选取 $z_1, \ldots, z_t \in K^i$，其像构成
$K^i \otimes \kappa/W$ 的一组基。于是元素
$d^i(z_1), \ldots, d^i(z_t) \in K^{i + 1}$
在 $K^{i + 1} \otimes_R \kappa$ 中线性无关。
由 Algebra，引理 \ref{algebra-lemma-cokernel-flat}，将 $R$
替换为某个 $f \in R$、$f \not \in \mathfrak p$ 对应的 $R_f$
之后，可以假设
\begin{enumerate}
% P02-E0263
\item $d^i(x_a), y_b, z_c$ 是 $K^i$ 的一组 $R$-基；
\item $d^i(z_1), \ldots, d^i(z_t)$ 在 $K^{i + 1}$ 中
$R$-线性无关；
\item 商 $E^{i + 1} = K^{i + 1}/\sum Rd^i(z_c)$ 是有限投射的。
\end{enumerate}
由于 $d^i$ 将 $d^{i - 1}(x_a)$ 和 $y_b$ 映为零，
由条件（2）可知
$E^{i + 1} = \Coker(d^i : K^i \to K^{i + 1})$。
于是
$$
\tau_{\geq i + 1}K^\bullet =
(\ldots \to 0 \to E^{i + 1} \to K^{i + 2} \to \ldots)
$$
是由有限投射模组成的有界复形，其非零次数位于
$[i + 1, b]$，其中 $b$ 为某个整数。因此
$\tau_{\geq i + 1}K^\bullet$ 是振幅为 $[i + 1, b]$ 的完美对象。
由于 $\tau_{\leq i}K^\bullet$ 在次数 $> i$ 中没有上同调，
可以对区别三角
$$
\tau_{\leq i}K^\bullet \to K^\bullet \to \tau_{\geq i + 1}K^\bullet \to
(\tau_{\leq i}K^\bullet)[1]
$$
应用引理 \ref{more-algebra-lemma-splitting-unique}
（Derived Categories，注
\ref{derived-remark-truncation-distinguished-triangle}），得到结论。
\end{proof}
\begin{lemma}
\label{more-algebra-lemma-isolate-a-cohomology-group}
设 $R$ 为一个环，设 $\mathfrak p \subset R$ 为一个素理想。
设 $K^\bullet$ 是伪凝聚 $R$-模复形。假设对某个
$i \in \mathbf{Z}$，映射
$$
H^i(K^\bullet) \otimes_R \kappa(\mathfrak p)
\longrightarrow
H^i(K^\bullet \otimes_R^{\mathbf{L}} \kappa(\mathfrak p))
\quad\text{且}\quad
H^{i - 1}(K^\bullet) \otimes_R \kappa(\mathfrak p)
\longrightarrow
H^{i - 1}(K^\bullet \otimes_R^{\mathbf{L}} \kappa(\mathfrak p))
$$
均为满射。则存在 $f \in R$、$f \not \in \mathfrak p$，使得
\begin{enumerate}
\item $\tau_{\geq i + 1}(K^\bullet \otimes_R R_f)$ 是 $D(R_f)$ 中
Tor-振幅为 $[i + 1, \infty]$ 的完美对象；
\item $H^i(K^\bullet)_f$ 是有限自由 $R_f$-模；
\item 在 $D(R_f)$ 中存在典范直和分解
$$
K^\bullet \otimes_R R_f \cong
\tau_{\leq i - 1}(K^\bullet \otimes_R R_f) \oplus
H^i(K^\bullet)_f[-i] \oplus
\tau_{\geq i + 1}(K^\bullet \otimes_R R_f)
$$
。
\end{enumerate}
\end{lemma}

\begin{proof}
由引理 \ref{more-algebra-lemma-better-cut-complex-in-two} 得到（1），以及
$D(R_f)$ 中的分裂
$K^\bullet \otimes_R R_f = \tau_{\leq i}K^\bullet \otimes_R R_f \oplus
\tau_{\geq i + 1}K^\bullet \otimes_R R_f$。
再对 $\tau_{\leq i}K^\bullet \otimes_R R_f$ 应用一次引理
\ref{more-algebra-lemma-better-cut-complex-in-two}，可在适当选取 $f$ 后得到分裂
% P02-E0264
$\tau_{\leq i}K^\bullet \otimes_R R_f =
\tau_{\leq i - 1}K^\bullet \otimes_R R_f \oplus H^i(K^\bullet)_f$（在
$D(R_f)$ 中）；还得到 $H^i(K)_f$ 是平坦完美模，即有限投射模。
% P02-E0265
\end{proof}

\begin{lemma}
\label{more-algebra-lemma-cut-complex-in-two}
设 $R$ 为一个环，设 $\mathfrak p \subset R$ 为一个素理想，
设 $i \in \mathbf{Z}$。设 $K^\bullet$ 是伪凝聚 $R$-模复形，满足
$H^i(K^\bullet \otimes_R^{\mathbf{L}} \kappa(\mathfrak p)) = 0$。
则存在 $f \in R$、$f \not \in \mathfrak p$，以及在 $D(R_f)$ 中的
典范直和分解
$$
K^\bullet \otimes_R R_f =
\tau_{\geq i + 1}(K^\bullet \otimes_R R_f) \oplus
\tau_{\leq i - 1}(K^\bullet \otimes_R R_f)
$$
其中 $\tau_{\geq i + 1}(K^\bullet \otimes_R R_f)$ 是
Tor-振幅为 $[i + 1, \infty]$ 的完美复形。
\end{lemma}

\begin{proof}
这是引理 \ref{more-algebra-lemma-better-cut-complex-in-two} 的一个常用特殊情形。
直接证明如下。可假设 $K^\bullet$ 是由有限自由 $R$-模组成的有界上复形。
考察次数 $i$ 中的情形：
$$
\ldots \to K^{i - 2} \to R^{\oplus l}
\to R^{\oplus m} \to R^{\oplus n} \to K^{i + 2} \to \ldots
$$
设 $A$ 为对应于 $K^{i - 1} \to K^i$ 的 $m \times l$ 矩阵，
设 $B$ 为对应于 $K^i \to K^{i + 1}$ 的 $n \times m$ 矩阵。
假设说明 $A \bmod \mathfrak p$ 的秩为 $r$，而
$B \bmod \mathfrak p$ 的秩为 $m - r$。换言之，$A$ 有某个
$r \times r$ 子式 $a$ 不属于 $\mathfrak p$，而 $B$ 有某个
$(m - r) \times (m - r)$-子式 $b$ 不属于 $\mathfrak p$。
置 $f = ab$。则在逆转 $f$ 后，可以找到直和分解
$K^{i - 1} = R^{\oplus l - r} \oplus R^{\oplus r}$、
$K^i = R^{\oplus r} \oplus R^{\oplus m - r}$、
$K^{i + 1} = R^{\oplus m - r} \oplus R^{\oplus n - m + r}$，
使得模映射 $K^{i - 1} \to K^i$ 将
% P02-E0266
$R^{\oplus l - r}$ 映为零，并诱导 $R^{\oplus r}$ 到 $K^i$
相应直和项上的同构；且模映射 $K^i \to K^{i + 1}$ 将
$R^{\oplus r}$ 映为零，并诱导 $R^{\oplus m - r}$ 到
$K^{i + 1}$ 相应直和项上的同构。因此 $K^\bullet$ 与
$$
\ldots \to K^{i - 2} \to R^{\oplus l - r}
\to 0 \to R^{\oplus n - m + r} \to K^{i + 2} \to \ldots
$$
拟同构，一切显然。
\end{proof}

\begin{lemma}
\label{more-algebra-lemma-split-using-ext-injective}
设 $R$ 为一个环，设 $K \in D^-(R)$，设 $a \in \mathbf{Z}$。
假设对任意单射 $R$-模映射 $M \to M'$，映射
$\Ext^{-a}_R(K, M) \to \Ext^{-a}_R(K, M')$ 都是单射。
则存在唯一的直和分解
$K \cong \tau_{\leq a}K \oplus \tau_{\geq a + 1}K$，并且对某个
$b$，$\tau_{\geq a + 1}K$ 的投射振幅为 $[a + 1, b]$。
\end{lemma}

\begin{proof}
考虑 $D(R)$ 中的区别三角
$$
\tau_{\leq a}K \to K \to \tau_{\geq a + 1}K \to (\tau_{\leq a}K)[1]
$$
参见 Derived Categories，注
\ref{derived-remark-truncation-distinguished-triangle}。
注意
$\Ext^{-a}_R(\tau_{\leq a}K, M) = \Hom_R(H^a(K), M)$
以及
$\Ext^{-a - 1}_R(\tau_{\leq a}K, M) = 0$，参见
Derived Categories，引理 \ref{derived-lemma-negative-exts}。
因此 $\Ext$ 的长正合序列给出正合序列
$$
0 \to
\Ext^{-a}_R(\tau_{\geq a + 1}K, M) \to
\Ext^{-a}_R(K, M) \to
\Hom_R(H^a(K), M)
$$
它关于 $R$-模 $M$ 是函子性的。现在若 $I$ 是内射 $R$-模，则例如由
Derived Categories，引理
\ref{derived-lemma-compute-ext-resolutions}，有
$\Ext^{-a}_R(\tau_{\geq a + 1}K, I) = 0$。
由于每个模都可嵌入一个内射模，得到对每个 $R$-模 $M$ 都有
$\Ext^{-a}_R(\tau_{\geq a + 1}K, M) = 0$。
由引理 \ref{more-algebra-lemma-projective-amplitude} 可知，对某个 $b$，
$\tau_{\geq a + 1}K$ 的投射振幅为 $[a + 1, b]$
（此处使用了 $K$ 有界上这一条件）。
由引理 \ref{more-algebra-lemma-splitting-unique} 得到分裂。
\end{proof}

\begin{lemma}
\label{more-algebra-lemma-split-using-ext-zero}
设 $R$ 为一个环，设 $K \in D^-(R)$，设 $a \in \mathbf{Z}$。
假设对任意 $R$-模 $M$ 都有 $\Ext^{-a}_R(K, M) = 0$。
则存在唯一的直和分解
$K \cong \tau_{\leq a - 1}K \oplus \tau_{\geq a + 1}K$，
并且对某个 $b$，$\tau_{\geq a + 1}K$ 的投射振幅为 $[a + 1, b]$。
\end{lemma}

\begin{proof}
由引理 \ref{more-algebra-lemma-split-using-ext-injective}，有直和分解
$K \cong \tau_{\leq a}K \oplus \tau_{\geq a + 1}K$，
并且对某个 $b$，$\tau_{\geq a + 1}K$ 的投射振幅为 $[a + 1, b]$。
显然必有 $H^a(K) = 0$，从而在 $D(R)$ 中
$\tau_{\leq a}K = \tau_{\leq a - 1}K$。
\end{proof}













\section{识别完美复形}
\label{more-algebra-section-recognizing-perfect}

\noindent
下面给出一些可用来证明某些复形完美的引理。

\begin{lemma}
\label{more-algebra-lemma-lift-bounded-pseudo-coherent-to-perfect}
设 $R$ 为一个环，设 $\mathfrak p \subset R$ 为一个素理想。
设 $K$ 伪凝聚且有界下。置
$d_i = \dim_{\kappa(\mathfrak p)}
H^i(K \otimes_R^\mathbf{L} \kappa(\mathfrak p))$。
若存在 $a \in \mathbf{Z}$，使得当 $i < a$ 时 $d_i = 0$，
则存在 $f \in R$、$f \not \in \mathfrak p$ 以及复形
$$
\ldots \to 0 \to R_f^{\oplus d_a} \to R_f^{\oplus d_{a + 1}} \to
\ldots \to
R_f^{\oplus d_{b - 1}} \to
R_f^{\oplus d_b} \to 0
\to \ldots
$$
在 $D(R_f)$ 中表示 $K \otimes_R^\mathbf{L} R_f$。
特别地，$K \otimes_R^\mathbf{L} R_f$ 是完美的。
\end{lemma}

\begin{proof}
% P02-E0267
减小 $a$ 后，可以进一步假设对 $i < a$ 有
$H^i(K^\bullet) = 0$。由引理 \ref{more-algebra-lemma-cut-complex-in-two}，
将 $R$ 替换为某个 $f \in R$、$f \not \in \mathfrak p$
对应的 $R_f$ 后，可以在 $D(R)$ 中写成
$K^\bullet = \tau_{\leq a - 1}K^\bullet \oplus
\tau_{\geq a}K^\bullet$，其中 $\tau_{\geq a}K^\bullet$ 完美。
由于对 $i < a$ 有 $H^i(K^\bullet) = 0$，可知在 $D(R)$ 中
$\tau_{\leq a - 1}K^\bullet = 0$。因此 $K^\bullet$ 是完美的。
再应用引理 \ref{more-algebra-lemma-lift-perfect-from-residue-field} 即得结论。
\end{proof}
\begin{lemma}
\label{more-algebra-lemma-check-perfect-pointwise}
设 $R$ 为一个环，设 $a, b \in \mathbf{Z}$。
设 $K^\bullet$ 是伪凝聚 $R$-模复形。下列条件等价：
\begin{enumerate}
\item $K^\bullet$ 是 Tor-振幅为 $[a, b]$ 的完美对象；
\item 对每个素理想 $\mathfrak p$ 以及所有
$i \not \in [a, b]$，有
$H^i(K^\bullet \otimes_R^{\mathbf{L}} \kappa(\mathfrak p)) = 0$；
\item 对每个极大理想 $\mathfrak m$ 以及所有
$i \not \in [a, b]$，有
$H^i(K^\bullet \otimes_R^{\mathbf{L}} \kappa(\mathfrak m)) = 0$。
\end{enumerate}
\end{lemma}

\begin{proof}
略去蕴含（1）$\Rightarrow$（2）$\Rightarrow$（3）的证明。
假设（3）成立。设 $i \in \mathbf{Z}$ 且 $i \not \in [a, b]$。
由引理 \ref{more-algebra-lemma-cut-complex-in-two}，该假设说明对 $R$ 的所有极大理想，
都有 $H^i(K^\bullet)_{\mathfrak m} = 0$。因此
$H^i(K^\bullet) = 0$，参见 Algebra，引理
\ref{algebra-lemma-characterize-zero-local}。
此外，引理 \ref{more-algebra-lemma-cut-complex-in-two} 还说明，对每个极大理想
$\mathfrak m$，存在 $f \in R$、$f \not \in \mathfrak m$，
使得 $K^\bullet \otimes_R R_f$ 是 Tor-振幅为 $[a, b]$ 的完美对象。
由引理 \ref{more-algebra-lemma-glue-perfect} 和
\ref{more-algebra-lemma-glue-tor-amplitude} 得到结论。
\end{proof}

\begin{lemma}
\label{more-algebra-lemma-check-perfect-stalks}
设 $R$ 为一个环，设 $K^\bullet$ 是伪凝聚 $R$-模复形。
考虑下列条件：
\begin{enumerate}
\item $K^\bullet$ 是完美的；
\item 对每个素理想 $\mathfrak p$，复形
$K^\bullet \otimes_R R_{\mathfrak p}$ 是完美的；
\item 对每个极大理想 $\mathfrak m$，复形
$K^\bullet \otimes_R R_{\mathfrak m}$ 是完美的；
\item 对每个素理想 $\mathfrak p$，当所有 $i \ll 0$ 时，有
$H^i(K^\bullet \otimes_R^{\mathbf{L}} \kappa(\mathfrak p)) = 0$；
\item 对每个极大理想 $\mathfrak m$，当所有 $i \ll 0$ 时，有
$H^i(K^\bullet \otimes_R^{\mathbf{L}} \kappa(\mathfrak m)) = 0$。
\end{enumerate}
总有下列蕴含：
% P02-E0268
$$
(1) \Rightarrow (2) \Leftrightarrow (3) \Leftrightarrow (3)
\Leftrightarrow (4) \Leftrightarrow (5)
$$
若 $K^\bullet$ 有界下，则所有条件均等价。
\end{lemma}

\begin{proof}
由引理 \ref{more-algebra-lemma-pull-perfect} 可知（1）蕴含（2）。
（2）$\Rightarrow$（3）是显然的。由于每个素理想
$\mathfrak p$ 都包含于一个极大理想 $\mathfrak m$，对映射
$R_\mathfrak m \to R_\mathfrak p$ 应用引理
\ref{more-algebra-lemma-pull-perfect}，可知（3）蕴含（2）。
对剩余域映射
$R_\mathfrak p \to \kappa(\mathfrak p)$ 和
$R_\mathfrak m \to \kappa(\mathfrak m)$ 应用引理
\ref{more-algebra-lemma-pull-perfect}，可知（2）蕴含（4），而（3）蕴含（5）。

\medskip\noindent
假设 $R$ 是极大理想为 $\mathfrak m$、剩余域为 $\kappa$ 的局部环。
我们将证明，若对某个 $a$ 以及 $i < a$ 有
$H^i(K^\bullet \otimes^\mathbf{L} \kappa) = 0$，
则 $K$ 是完美的。
% P02-E0269
这将证明（4）蕴含（2）且（5）蕴含（3），从而得到引理的第一部分。
首先对 $i = a - 1$ 应用引理 \ref{more-algebra-lemma-cut-complex-in-two}，可知
在 $D(R)$ 中
$K^\bullet = \tau_{\leq a - 1}K^\bullet \oplus
\tau_{\geq a}K^\bullet$，其中 $\tau_{\geq a}K^\bullet$ 是
Tor-振幅包含于 $[a, \infty]$ 的完美对象。
为完成证明，需要说明 $\tau_{\leq a - 1}K$ 为零，
即其上同调群均为零。否则，设 $i$ 是使
$M = H^i(\tau_{\leq a - 1}K)$ 非零的最大指标。
由于 $\tau_{\leq a - 1}K^\bullet$ 是伪凝聚的，
$M$ 是有限 $R$-模（引理 \ref{more-algebra-lemma-finite-cohomology} 和
\ref{more-algebra-lemma-summands-pseudo-coherent}）。
因此由 Nakayama 引理（Algebra，引理 \ref{algebra-lemma-NAK}），
$M \otimes_R \kappa$ 非零。这说明
$$
H^i((\tau_{\leq a - 1}K^\bullet) \otimes_R^\mathbf{L} \kappa) =
H^i(K^\bullet \otimes_R^\mathbf{L} \kappa)
$$
非零，矛盾。

\medskip\noindent
假设等价条件（2）--（5）成立，且 $K^\bullet$ 有界下。
设对 $i < a$ 有 $H^i(K^\bullet) = 0$。选取 $R$ 的一个极大理想
$\mathfrak m$。只需证明存在 $f \in R$、$f \not \in \mathfrak m$，
使得 $K^\bullet \otimes_R^\mathbf{L} R_f$ 是完美的
（引理 \ref{more-algebra-lemma-glue-perfect} 和
Algebra，引理 \ref{algebra-lemma-quasi-compact}）。
这由引理 \ref{more-algebra-lemma-lift-bounded-pseudo-coherent-to-perfect} 得到。
\end{proof}

\begin{lemma}
\label{more-algebra-lemma-projective-amplitude-pseudo-coherent}
设 $R$ 为一个环，设 $K$ 是 $D(R)$ 中的伪凝聚对象。
设 $a, b \in \mathbf{Z}$。下列条件等价：
\begin{enumerate}
\item $K$ 的投射振幅为 $[a, b]$；
\item $K$ 是 Tor-振幅为 $[a, b]$ 的完美对象；
\item 对所有有限表示 $R$-模 $N$ 以及所有
$i \not \in [-b, -a]$，有 $\Ext^i_R(K, N) = 0$；
\item 对 $n > b$ 有 $H^n(K) = 0$，并且对所有有限表示
$R$-模 $N$ 以及所有 $i > -a$，有 $\Ext^i_R(K, N) = 0$；
\item 对 $n \not \in [a - 1, b]$ 有 $H^n(K) = 0$，并且对所有
有限表示 $R$-模 $N$，有 $\Ext^{-a + 1}_R(K, N) = 0$。
\end{enumerate}
\end{lemma}

\begin{proof}
由引理 \ref{more-algebra-lemma-perfect} 的最后一个断言可知（2）蕴含（1）。
若（1）成立，则 $K$ 可由投射模 $P^i$ 组成的复形表示，其中当
$i \not \in [a, b]$ 时 $P^i = 0$。
由于投射模是平坦模（它们是自由模的直和项），由引理
\ref{more-algebra-lemma-tor-amplitude} 可知 $K$ 的 Tor-振幅为 $[a, b]$。
于是由引理 \ref{more-algebra-lemma-perfect} 可知（2）成立。

\medskip\noindent
在条件（3）、（4）、（5）中，假设的 Ext 群
$\Ext^i_R(K, M)$ 对有限表示模 $M$ 的消失，等价于其对所有
$R$-模 $M$ 的消失；这由引理
\ref{more-algebra-lemma-pseudo-coherence-colimit-ext} 和
Algebra，引理 \ref{algebra-lemma-module-colimit-fp} 得到。
因此（1）、（3）、（4）、（5）的等价性由引理
\ref{more-algebra-lemma-projective-amplitude} 得到。
\end{proof}

\noindent
% P02-E0270
下面的引理可用于在多项式环
$B = A[x_1, \ldots, x_d]$ 上寻找完美复形。

\begin{lemma}
\label{more-algebra-lemma-perfect-over-polynomial-ring}
设 $A \to B$ 为一个环映射，设 $a, b \in \mathbf{Z}$，设
$d \geq 0$。设 $K^\bullet$ 是 $B$-模复形。假设
\begin{enumerate}
\item 环映射 $A \to B$ 平坦；
\item 对每个素理想 $\mathfrak p \subset A$，环
$B \otimes_A \kappa(\mathfrak p)$ 的整体维数有限且 $\leq d$；
\item 作为 $B$-模复形，$K^\bullet$ 是伪凝聚的；
\item 作为 $A$-模复形，$K^\bullet$ 的 Tor-振幅为 $[a, b]$。
\end{enumerate}
则作为 $B$-模复形，$K^\bullet$ 是完美的，且 Tor-振幅为
$[a - d, b]$。
\end{lemma}

\begin{proof}
可假设 $K^\bullet$ 是由有限自由 $B$-模组成的有界上复形。
特别地，作为 $A$-模复形，$K^\bullet$ 是平坦的，并且对任意
$A$-模 $M$ 有
$K^\bullet \otimes_A M = K^\bullet \otimes_A^{\mathbf{L}} M$。
对 $A$ 的每个素理想 $\mathfrak p$，复形
$$
K^\bullet \otimes_A \kappa(\mathfrak p)
$$
是由有限自由 $B \otimes_A \kappa(\mathfrak p)$-模组成的有界上复形，
且除 $i \in [a, b]$ 外所有 $H^i$ 均消失。
由于 $B \otimes_A \kappa(\mathfrak p)$ 的整体维数为 $d$，
由引理 \ref{more-algebra-lemma-finite-gl-dim-tor-dimension} 可知
$K^\bullet \otimes_A \kappa(\mathfrak p)$ 的 Tor-振幅为
$[a - d, b]$。设 $\mathfrak q$ 是 $B$ 的一个位于
$\mathfrak p$ 之上的素理想。由于
$K^\bullet \otimes_A \kappa(\mathfrak p)$ 是由自由
$B \otimes_A \kappa(\mathfrak p)$-模组成的有界上复形，有
\begin{align*}
K^\bullet \otimes_B^{\mathbf{L}} \kappa(\mathfrak q)
& = K^\bullet \otimes_B \kappa(\mathfrak q) \\
& = (K^\bullet \otimes_A \kappa(\mathfrak p))
\otimes_{B \otimes_A \kappa(\mathfrak p)} \kappa(\mathfrak q) \\
& = (K^\bullet \otimes_A \kappa(\mathfrak p))
\otimes^{\mathbf{L}}_{B \otimes_A \kappa(\mathfrak p)} \kappa(\mathfrak q)
\end{align*}
所以以上论证说明，当 $i \not \in [a - d, b]$ 时，
$H^i(K^\bullet \otimes_B^{\mathbf{L}} \kappa(\mathfrak q)) = 0$。
由引理 \ref{more-algebra-lemma-check-perfect-pointwise} 得到结论。
\end{proof}
\noindent
下述引理是引理 \ref{more-algebra-lemma-perfect-over-polynomial-ring} 的局部版本。
它可用于在正则局部环上寻找完美复形。

\begin{lemma}
\label{more-algebra-lemma-perfect-over-regular-local-ring}
设 $A \to B$ 是局部环同态，设 $a, b \in \mathbf{Z}$，设
$d \geq 0$。设 $K^\bullet$ 是 $B$-模复形。假设
\begin{enumerate}
\item 环映射 $A \to B$ 平坦；
\item 环 $B/\mathfrak m_AB$ 是维数为 $d$ 的正则环；
\item 作为 $B$-模复形，$K^\bullet$ 是伪凝聚的；
\item 作为 $A$-模复形，$K^\bullet$ 的 Tor-振幅为 $[a, b]$；
事实上，只需
$H^i(K^\bullet \otimes_A^\mathbf{L} \kappa(\mathfrak m_A))$
仅对 $i \in [a, b]$ 非零即可。
\end{enumerate}
则作为 $B$-模复形，$K^\bullet$ 是完美的，且 Tor-振幅为
$[a - d, b]$。
\end{lemma}

\begin{proof}
由（3），可假设 $K^\bullet$ 是由有限自由 $B$-模组成的有界上复形。
计算得
\begin{align*}
K^\bullet \otimes_B^{\mathbf{L}} \kappa(\mathfrak m_B)
& = K^\bullet \otimes_B \kappa(\mathfrak m_B) \\
& = (K^\bullet \otimes_A \kappa(\mathfrak m_A))
\otimes_{B/\mathfrak m_A B} \kappa(\mathfrak m_B) \\
& = (K^\bullet \otimes_A \kappa(\mathfrak m_A))
\otimes^{\mathbf{L}}_{B/\mathfrak m_A B} \kappa(\mathfrak m_B)
\end{align*}
% P02-E0271
第一个等式成立，是因为 $K^\bullet$ 是由平坦 $B$-模组成的有界上复形。
第二个等式由张量积的基本性质得到。第三个等式成立，是因为
$K^\bullet \otimes_A \kappa(\mathfrak m_A) =
K^\bullet/ \mathfrak m_A K^\bullet$
是由平坦 $B/\mathfrak m_A B$-模组成的有界上复形。
由（1），$K^\bullet$ 是由平坦 $A$-模组成的有界上复形；由假设（4），
复形 $K^\bullet \otimes_A \kappa(\mathfrak m_A)$ 的上同调模 $H^i$
仅对 $i \in [a, b]$ 非零。因此，例
\ref{more-algebra-example-cohomology-complex-tensored} 的谱序列，以及
$B/\mathfrak m_AB$ 的整体维数有限且为 $d$ 这一事实
（由（2）和 Algebra，命题
\ref{algebra-proposition-regular-finite-gl-dim}），说明当
$j \not \in [a - d, b]$ 时，
$H^j(K^\bullet \otimes_B^{\mathbf{L}} \kappa(\mathfrak m_B))$ 为零。
% P02-E0272
由引理 \ref{more-algebra-lemma-check-perfect-pointwise} 即完成证明。
\end{proof}





\section{完美复形的刻画}
\label{more-algebra-section-perfect-compact}

\noindent
本节证明，完美复形恰为环的导出范畴中的紧对象。先证明如下结论。

\begin{lemma}
\label{more-algebra-lemma-perfect-ring-classical-generator}
设 $R$ 为一个环。完美对象组成的全子范畴
$D_{perf}(R) \subset D(R)$ 是包含 $R = R[0]$ 的最小严格全、饱和三角子范畴。
换言之，$D_{perf}(R) = \langle R \rangle$。
特别地，$R$ 是 $D_{perf}(R)$ 的经典生成元。
\end{lemma}

\begin{proof}
关于该断言的含义，参见 Derived Categories，定义
\ref{derived-definition-saturated} 和
\ref{derived-definition-generators}。
引理 \ref{more-algebra-lemma-two-out-of-three-perfect} 和
\ref{more-algebra-lemma-summands-perfect} 已说明
$D_{perf}(R) \subset D(R)$ 是 $D(R)$ 的严格全、饱和三角子范畴。
当然，$R \in D_{perf}(R)$。

\medskip\noindent
回忆 $\langle R \rangle = \bigcup \langle R \rangle_n$。
为完成证明，将说明：若 $M \in D_{perf}(R)$ 由
$$
\ldots \to 0 \to M^a \to M^{a + 1} \to \ldots \to M^b \to 0 \to \ldots
$$
表示，其中 $M^i$ 为有限投射模，则
$M \in \langle R \rangle_{b - a + 1}$。对 $b - a$ 作归纳。
按定义，$\langle R \rangle_1$ 包含置于任意次数的每个有限投射
$R$-模；这处理了归纳的基础情形 $b - a = 0$。
一般地，考虑区别三角
% P02-E0273
$$
M_b[-b] \to M^\bullet \to \sigma_{\leq b - 1}M^\bullet \to M_b[-b + 1]
$$
由归纳假设，截断复形 $\sigma_{\leq b - 1}M^\bullet$ 属于
$\langle R \rangle_{b - a}$，而 $M_b[-b]$ 属于
$\langle R \rangle_1$。因此按定义，
$M^\bullet \in \langle R \rangle_{b - a + 1}$。
\end{proof}

\noindent
设 $R$ 为一个环。回忆 $D(R)$ 有直和，且它们由复形的直和直接给出，
参见 Derived Categories，引理 \ref{derived-lemma-direct-sums}。
本节各引理中将不再特别说明这一点。

\begin{lemma}
\label{more-algebra-lemma-commutes-with-countable-sums}
设 $R$ 为一个环。设 $K \in D(R)$ 是满足如下条件的对象：
对每个可数对象族 $E_n \in D(R)$，典范映射
$$
\bigoplus \Hom_{D(R)}(K, E_n) \longrightarrow \Hom_{D(R)}(K, \bigoplus E_n)
$$
是双射。则对任意以 $\mathbf{N}$ 为指标的复形系统 $L_n^\bullet$，映射
$$
\colim \Hom_{D(R)}(K, L^\bullet_n) \longrightarrow \Hom_{D(R)}(K, L^\bullet)
$$
是双射，其中 $L^\bullet$ 是逐项余极限，即对所有
$m \in \mathbf{Z}$ 有 $L^m = \colim L_n^m$。
\end{lemma}

\begin{proof}
考虑复形的短正合序列
$$
0 \to \bigoplus L_n^\bullet \to \bigoplus L_n^\bullet \to L^\bullet \to 0
$$
其中第一个映射在次数 $n$ 中由 $1 - t_n$ 给出，
$t_n : L_n^\bullet \to L_{n + 1}^\bullet$ 为转移映射。
由 Derived Categories，引理
\ref{derived-lemma-derived-canonical-delta-functor}，
这是 $D(R)$ 中的区别三角。应用同调函子
$\Hom_{D(R)}(K, -)$，参见 Derived Categories，引理
\ref{derived-lemma-representable-homological}。
% P02-E0274
于是得到长正合上同调序列
$$
\xymatrix{
& \ldots \ar[r] & \Hom_{D(R)}(K, \colim L^\bullet_n[-1]) \ar[lld] \\
\Hom_{D(R)}(K, \bigoplus L^\bullet_n) \ar[r] &
\Hom_{D(R)}(K, \bigoplus L^\bullet_n) \ar[r] &
\Hom_{D(R)}(K, \colim L^\bullet_n) \ar[lld] \\
\Hom_{D(R)}(K, \bigoplus L^\bullet_n[1]) \ar[r] & \ldots
}
$$
由于已假设 $\Hom_{D(R)}(K, \bigoplus L^\bullet_n)$ 等于
$\bigoplus \Hom_{D(R)}(K, L^\bullet_n)$，图表每一行的第一个映射都是单射
（这由该映射作为 $1 - t_n$ 所诱导映射之和的显式描述得到）。
因此 $\Hom_{D(R)}(K, \colim L^\bullet_n)$ 是上述图表中间一行
第一个映射的余核，这正是所需证明的。
\end{proof}
\noindent
下述命题将完美复形刻画为导出范畴的紧对象
（Derived Categories，定义 \ref{derived-definition-compact-object}），
它在多处出现。例如参见
\cite[命题 6.3 的证明]{Rickard}（这里处理有界情形）、
\cite[定理 2.4.3]{TT}（陈述并不完全一致），以及
\cite[命题 6.4]{Bokstedt-Neeman}（注意其极其糟糕的记号约定）。

\begin{proposition}
\label{more-algebra-proposition-perfect-is-compact}
设 $R$ 为一个环。对 $D(R)$ 中的对象 $K$，下列条件等价：
\begin{enumerate}
\item $K$ 是完美的；
\item $K$ 是 $D(R)$ 中的紧对象。
\end{enumerate}
\end{proposition}

\begin{proof}
假设 $K$ 是完美的，即 $K$ 拟同构于由有限投射模组成的有界复形
$P^\bullet$，参见定义 \ref{more-algebra-definition-perfect}。
若 $E_i$ 由复形 $E_i^\bullet$ 表示，则 $\bigoplus E_i$
由次数 $n$ 的项为 $\bigoplus E_i^n$ 的复形表示。
另一方面，对所有 $n$，由于 $P^n$ 是投射的，对每个 $R$-模复形
$K^\bullet$，有
$\Hom_{D(R)}(P^\bullet, K^\bullet) = \Hom_{K(R)}(P^\bullet, K^\bullet)$，
参见 Derived Categories，引理
\ref{derived-lemma-morphisms-from-projective-complex}。
因此，$\Hom_{D(R)}(P^\bullet, E^\bullet)$ 是下述复形的上同调：
$$
\prod \Hom_R(P^n, E^{n - 1}) \to
\prod \Hom_R(P^n, E^n) \to
\prod \Hom_R(P^n, E^{n + 1}).
$$
由于 $P^\bullet$ 有界，在上述复形中可以将 $\prod$ 号换成
$\bigoplus$ 号。由于每个 $P^n$ 都是有限 $R$-模，对所有
$n, m$ 有
$\Hom_R(P^n, \bigoplus_i E_i^m) = \bigoplus_i \Hom_R(P^n, E_i^m)$。
综合这些观察可知，Derived Categories，定义
\ref{derived-definition-compact-object} 中的映射是双射。

\medskip\noindent
反之，假设 $K$ 是紧的。用复形 $K^\bullet$ 表示 $K$，并考虑映射
$$
K^\bullet
\longrightarrow
\bigoplus\nolimits_{n \geq 0} \tau_{\geq n} K^\bullet
$$
其中使用了典范截断，参见 Homology，第
\ref{homology-section-truncations} 节。该映射有意义，因为在每个次数上，
右端的直和都是有限的。由假设，该映射经过一个有限直和分解。
因此，至少对某个 $n$，映射 $K \to \tau_{\geq n} K$ 为零；
换言之，$K$ 属于 $D^{-}(R)$。

\medskip\noindent
由于 $K \in D^{-}(R)$，且每个 $R$-模都是某个自由模的商，
可用自由 $R$-模组成的有界上复形 $K^\bullet$ 表示 $K$，参见
Derived Categories，引理
\ref{derived-lemma-subcategory-left-resolution}。注意
$$
K^\bullet = \bigcup\nolimits_{n \leq 0} \sigma_{\geq n}K^\bullet
$$
其中使用了硬截断，参见 Homology，第
\ref{homology-section-truncations} 节。因此由引理
\ref{more-algebra-lemma-commutes-with-countable-sums} 可知，$D(R)$ 中的
$1 : K^\bullet \to K^\bullet$ 经过
$\sigma_{\geq n}K^\bullet \to K^\bullet$ 分解。
于是对某个有界且各项为自由 $R$-模的复形 $L^\bullet$，
$D(R)$ 中的 $1 : K^\bullet \to K^\bullet$ 分解为
$$
K^\bullet \xrightarrow{\varphi} L^\bullet \xrightarrow{\psi} K^\bullet
$$
设当 $i \not \in [a, b]$ 时 $L^i = 0$。此后固定 $a, b$。
设 $c$ 是满足下列条件的最大整数 $\leq b + 1$：存在上述
$1_{K^\bullet}$ 的分解，且当 $i < c$ 时 $L^i$ 是有限自由的。
% P02-E0275
我们将归纳证明 $c = b + 1$。具体地，写成
$L^c = \bigoplus_{\lambda \in \Lambda} R$。
由于 $L^{c - 1}$ 是有限自由的，可以找到有限子集
$\Lambda' \subset \Lambda$，使得 $L^{c - 1} \to L^c$ 经过
$\bigoplus_{\lambda \in \Lambda'} R \subset L^c$ 分解。
考虑复形映射
$$
\pi :
L^\bullet
\longrightarrow
(\bigoplus\nolimits_{\lambda \in \Lambda \setminus \Lambda'} R)[-c]
$$
它在次数 $c$ 中由投影到对应于
$\Lambda \setminus \Lambda'$ 的因子给出。
由关于 $K$ 的假设，在必要时用更大的有限子集替换 $\Lambda'$ 后，
可以假设在 $D(R)$ 中 $\pi \circ \varphi = 0$。
设 $(L')^\bullet \subset L^\bullet$ 为 $\pi$ 的核。
由于 $\pi$ 满射，得到复形的短正合序列，进而得到 $D(R)$ 中的
区别三角（参见 Derived Categories，引理
\ref{derived-lemma-derived-canonical-delta-functor}）。
由于 $\Hom_{D(R)}(K, -)$ 是同调函子（参见
Derived Categories，引理
\ref{derived-lemma-representable-homological}），并且
$\pi \circ \varphi = 0$，可以找到 $D(R)$ 中的态射
$\varphi' : K^\bullet \to (L')^\bullet$，其与
$(L')^\bullet \to L^\bullet$ 的复合等于 $\varphi$。
令 $\psi'$ 等于 $\psi$ 与 $(L')^\bullet \to L^\bullet$ 的复合，
便得到新的分解。由于 $(L')^\bullet$ 除次数 $c$ 外与
$L^\bullet$ 一致，并且
$(L')^c = \bigoplus_{\lambda \in \Lambda'} R$，归纳步骤得证。

\medskip\noindent
上一段讨论的结论是，$1_K : K \to K$ 在 $D(R)$ 中分解为
$$
K \xrightarrow{\varphi} L \xrightarrow{\psi} K
$$
其中 $L$ 可由有限自由 $R$-模复形表示。特别地，$L$ 是完美的。
注意，$e = \varphi \circ \psi \in \text{End}_{D(R)}(L)$ 是幂等元。
由 Derived Categories，引理
\ref{derived-lemma-projectors-have-images-triangulated} 可知
$L = \Ker(e) \oplus \Ker(1 - e)$。
映射 $\varphi : K \to L$ 在 $D(R)$ 中诱导到
$\Ker(1 - e)$ 的同构。因此最后由引理
\ref{more-algebra-lemma-summands-perfect} 得到 $K$ 是完美的。
\end{proof}

\begin{lemma}
\label{more-algebra-lemma-perfect-modulo-nilpotent-ideal}
设 $R$ 为一个环，设 $I \subset R$ 为一个理想，设 $K$ 为
$D(R)$ 中的对象。假设
\begin{enumerate}
\item $K \otimes_R^\mathbf{L} R/I$ 在 $D(R/I)$ 中是完美的；
\item $I$ 是幂零理想。
\end{enumerate}
则 $K$ 在 $D(R)$ 中是完美的。
\end{lemma}

\begin{proof}
选取由有限投射 $R/I$-模组成的有限复形
$\overline{P}^\bullet$，表示
$K \otimes_R^\mathbf{L} R/I$，参见定义
\ref{more-algebra-definition-perfect}。由引理
\ref{more-algebra-lemma-lift-complex-projectives}，存在由投射 $R$-模组成的复形
$P^\bullet$，表示 $K$，使得
$\overline{P}^\bullet = P^\bullet/IP^\bullet$。
由 Nakayama 引理（Algebra，引理 \ref{algebra-lemma-NAK}）可知，
$P^\bullet$ 是由有限投射 $R$-模组成的有限复形。
\end{proof}

\begin{lemma}
\label{more-algebra-lemma-perfect-modulo-two-ideals}
设 $R$ 为一个环。设 $I, J \subset R$ 为理想。
设 $K$ 为 $D(R)$ 中的对象。假设
\begin{enumerate}
\item $K \otimes_R^\mathbf{L} R/I$ 在 $D(R/I)$ 中是完美的；
\item $K \otimes_R^\mathbf{L} R/J$ 在 $D(R/J)$ 中是完美的。
\end{enumerate}
则 $K \otimes_R^\mathbf{L} R/IJ$ 在 $D(R/IJ)$ 中是完美的。
\end{lemma}

\begin{proof}
% P02-E0276
显然，可以用 $R/IJ$ 替换 $R$，并用
$K \otimes_R^\mathbf{L} R/IJ$ 替换 $K$。此时
$R \to R/(I \cap J)$ 是核的平方为零的满射。因此由引理
\ref{more-algebra-lemma-perfect-modulo-nilpotent-ideal}，只需证明
$K \otimes_R^\mathbf{L} R/(I \cap J)$ 是完美的。
故可以假设 $I \cap J = 0$。

\medskip\noindent
我们在 $I \cap J = 0$ 的情形证明该引理。首先，可以用一个
K-平坦复形 $K^\bullet$ 表示 $K$，其中每个 $K^n$ 都是平坦的；
参见引理 \ref{more-algebra-lemma-K-flat-resolution}。于是有复形的短正合序列
$$
0 \to
K^\bullet \to K^\bullet/IK^\bullet \oplus K^\bullet/JK^\bullet \to
K^\bullet/(I + J)K^\bullet \to 0
$$
按导出张量积的构造，$K^\bullet/IK^\bullet$ 表示
$K \otimes^\mathbf{L}_R R/I$。对
$K^\bullet/JK^\bullet$ 和 $K^\bullet/(I + J)K^\bullet$ 亦然。
注意，$K^\bullet/(I + J)K^\bullet$ 是
$R/(I + J)$-模的完美复形；参见引理 \ref{more-algebra-lemma-pull-perfect}。
% P02-E0277
因此，复形 $K^\bullet/IK^\bullet$、
$K^\bullet/JK^\bullet$ 和 $K^\bullet/(I + J)K^\bullet$
只有有限多个非零上同调群（因为完美复形具有有限 Tor-振幅；
参见引理 \ref{more-algebra-lemma-perfect}）。
% P02-E0278
由与上面所示复形短正合序列相伴的长正合上同调序列，
可知 $K \in D^b(R)$。特别地，可以假设 $K^\bullet$
是自由 $R$-模组成的有界上复形（参见 Derived Categories，引理
\ref{derived-lemma-subcategory-left-resolution}）。

\medskip\noindent
现在用命题 \ref{more-algebra-proposition-perfect-is-compact} 的判据证明
$K$ 是完美的。
% P02-E0279
为此，令 $E_j \in D(R)$ 为由集合 $J$ 参数化的对象族。
选取逐项平坦且表示 $E_j$ 的复形 $E_j^\bullet$；例如参见引理
\ref{more-algebra-lemma-K-flat-resolution}。显然，
$$
0 \to
E_j^\bullet \to
E_j^\bullet/IE_j^\bullet \oplus E_j^\bullet/JE_j^\bullet \to
E_j^\bullet/(I + J)E_j^\bullet \to 0
$$
是复形的短正合序列。取直和，得到类似的短正合序列
$$
0 \to
\bigoplus E_j^\bullet \to
\bigoplus E_j^\bullet/IE_j^\bullet \oplus E_j^\bullet/JE_j^\bullet \to
\bigoplus E_j^\bullet/(I + J)E_j^\bullet \to 0
$$
（注意，$- \otimes_R R/I$ 与直和交换。）
该短正合序列在 $D(R)$ 中确定一个区别三角；参见
Derived Categories，引理
\ref{derived-lemma-derived-canonical-delta-functor}。
应用同调函子 $\Hom_{D(R)}(K, -)$（参见 Derived Categories，引理
\ref{derived-lemma-representable-homological}），得到交换图
% P02-E0280
$$
\xymatrix{
\bigoplus \Hom_{D(R)}(K^\bullet, E_j^\bullet/(I + J))[-1] \ar[r] \ar[d] &
\Hom_{D(R)}(K^\bullet, \bigoplus E_j^\bullet/(I + J))[-1] \ar[d] \\
\bigoplus \Hom_{D(R)}(K^\bullet, E_j^\bullet/I \oplus E_j^\bullet/J)[-1]
\ar[r] \ar[d] &
\Hom_{D(R)}(K^\bullet, \bigoplus E_j^\bullet/I \oplus E_j^\bullet/J)[-1]
\ar[d] \\
\bigoplus \Hom_{D(R)}(K^\bullet, E_j^\bullet) \ar[r] \ar[d] &
\Hom_{D(R)}(K^\bullet, \bigoplus E_j^\bullet) \ar[d] \\
\bigoplus \Hom_{D(R)}(K^\bullet, E_j^\bullet/I \oplus E_j^\bullet/J)
\ar[r] \ar[d] &
\Hom_{D(R)}(K^\bullet, \bigoplus E_j^\bullet/I \oplus E_j^\bullet/J)
\ar[d] \\
\bigoplus \Hom_{D(R)}(K^\bullet, E_j^\bullet/(I + J)) \ar[r] &
\Hom_{D(R)}(K^\bullet, \bigoplus E_j^\bullet/(I + J))
}
$$
其各列正合。显然，对任意 $R$-模复形 $E^\bullet$，有
\begin{align*}
\Hom_{D(R)}(K^\bullet, E^\bullet/I) & =
\Hom_{K(R)}(K^\bullet, E^\bullet/I) \\
& =
\Hom_{K(R/I)}(K^\bullet/IK^\bullet, E^\bullet/I) \\
& =
\Hom_{D(R/I)}(K^\bullet/IK^\bullet, E^\bullet/I)
\end{align*}
除以 $J$ 或 $I + J$ 时也同样成立；参见 Derived Categories，引理
\ref{derived-lemma-morphisms-from-projective-complex}。
% P02-E0281
因此，除中间的水平箭头可能例外以外，其余水平箭头都是同构，
因为复形 $K^\bullet/IK^\bullet$、$K^\bullet/JK^\bullet$、
$K^\bullet/(I + J)K^\bullet$ 分别是 $R/I$、$R/J$、
$R/(I + J)$-模的完美复形；参见命题
\ref{more-algebra-proposition-perfect-is-compact}。
由 $5$-引理（Homology，引理 \ref{homology-lemma-five-lemma}），
中间的映射是同构，再由命题 \ref{more-algebra-proposition-perfect-is-compact}
即得结论。
\end{proof}






\section{强生成元与正则环}
\label{more-algebra-section-strong-generators-regular}

\noindent
设 $R$ 为一个环。
% P02-E0282
用 $D(R)_c$ 表示 $D(R)$ 的饱和满三角子范畴。已经知道
$$
\langle R \rangle = D_{perf}(R) = D(R)_c
$$
参见引理 \ref{more-algebra-lemma-perfect-ring-classical-generator} 和命题
\ref{more-algebra-proposition-perfect-is-compact}。事实证明，若 $R$ 是正则的，
则 $R$ 是强生成元（Derived Categories，定义
\ref{derived-definition-generators}）。

\begin{lemma}
\label{more-algebra-lemma-ghost-lemma}
\begin{reference}
\cite{Kelly}
\end{reference}
设 $R$ 为一个环，设 $n \geq 1$，并设
$K \in \langle R \rangle_n$；记号见 Derived Categories，第
\ref{derived-section-generators} 节。考虑 $D(R)$ 中的映射
$$
K \xrightarrow{f_1} K_1 \xrightarrow{f_2} K_2
\xrightarrow{f_3} \ldots \xrightarrow{f_n} K_n
$$
若对所有 $i, j$ 都有 $H^i(f_j) = 0$，则
$f_n \circ \ldots \circ f_1 = 0$。
\end{lemma}

\begin{proof}
若 $n = 1$，则 $K$ 是 $D(R)$ 中某个有界复形
$P^\bullet$ 的直和项，该复形的各项是有限自由 $R$-模，且微分均为零。
因此，只需证明 $D(R)$ 中满足对所有 $i$ 都有 $H^i(f) = 0$ 的任意态射
$f : P^\bullet \to K_1$ 为零。由于 $P^\bullet$ 是有限直和
$P^\bullet = \bigoplus R[m_j]$，只需证明 $D(R)$ 中满足
$H^{-m}(g) = 0$ 的任意态射 $g : R[m] \to K_1$ 为零。
这由等式 $\Hom_{D(R)}(R[-m], K) = H^m(K)$ 得到。

\medskip\noindent
当 $n > 1$ 时，对 $n$ 作归纳。我们知道，$K$ 是 $D(R)$ 中某个对象
$P$ 的直和项，而该对象位于区别三角
$$
P' \xrightarrow{i} P \xrightarrow{p} P'' \to P'[1]
$$
中，其中 $P' \in \langle R \rangle_1$ 且
$P'' \in \langle R \rangle_{n - 1}$。与上面一样，可以用 $P$ 替换
$K$，并假设在 $D(R)$ 中有
$$
P \xrightarrow{f_1} K_1 \xrightarrow{f_2} K_2
\xrightarrow{f_3} \ldots \xrightarrow{f_n} K_n
$$
且 $f_j$ 在上同调上为零。由 $n = 1$ 的情形，复合
$f_1 \circ i$ 为零。因此由 Derived Categories，引理
\ref{derived-lemma-representable-homological}，可以找到态射
$h : P'' \to K_1$，使得 $f_1 = h \circ p$。注意
$f_2 \circ h$ 在上同调上为零。于是由归纳假设，
$f_n \circ \ldots \circ f_2 \circ h = 0$，从而
$f_n \circ \ldots \circ f_1 =
f_n \circ \ldots \circ f_2 \circ h \circ p = 0$，正合所需。
\end{proof}

\begin{lemma}
\label{more-algebra-lemma-not-regular-not-strong}
设 $R$ 为 Noether 环。若 $R$ 是 $D_{perf}(R)$ 的强生成元，
则 $R$ 是有限维正则环。
\end{lemma}

\begin{proof}
假设对某个 $n \geq 1$ 有
$D_{perf}(R) = \langle R \rangle_n$。对任意有限 $R$-模 $M$，
可以选取有限自由 $R$-模复形
% P02-E0283
$$
P = (
P^{-n - 1} \xrightarrow{d^{-n - 1}}
P^{-n} \xrightarrow{d^{-n}}
P^{-n + 1} \xrightarrow{d^1}
\ldots \xrightarrow{d^{-1}} P^0)
$$
使得当 $i = -n, \ldots, - 1$ 时有 $H^i(P) = 0$，并且
$M \cong \Coker(d^{-1})$。注意 $P$ 属于 $D_{perf}(R)$。
% P02-E0284
对任意 $R$-模 $N$，可以通过 $M$ 的有限自由解消 $P$ 计算
$\Ext^n_R(M, N)$；参见 Algebra，第 \ref{algebra-section-ext} 节，
并与 Derived Categories，第 \ref{derived-section-ext} 节比较。
特别地，上述序列定义元素
$$
\xi \in \Ext^n_R(\Coker(d^{-1}), \Coker(d^{-n - 1})) =
\Ext^n_R(M, \Coker(d^{-n - 1}))
$$
并且对 $\Ext^n_R(M, N)$ 中任意元素 $\overline{\xi}$，存在
$R$-模映射 $\varphi : \Coker(d^{-n - 1}) \to N$，使得
$\varphi$ 把 $\xi$ 映到 $\overline{\xi}$。
对 $j = 1, \ldots, n - 1$，考虑复形
$$
K_j = (\Coker(d^{-n - 1}) \to P^{-n + 1} \to \ldots \to P^{-j})
$$
其中 $\Coker(d^{-n - 1})$ 位于次数 $-n$，而 $P^t$ 位于次数 $t$。
还令 $K_n = \Coker(d^{-n - 1})[n]$。于是有映射
$$
P \to K_1 \to K_2 \to \ldots \to K_n
$$
它们在上同调上诱导零映射。由于
$P \in D_{perf}(R) = \langle R \rangle_n$，引理
\ref{more-algebra-lemma-ghost-lemma} 表明这些映射的复合在 $D(R)$ 中为零。
% P02-E0285
又由 Derived Categories，引理
\ref{derived-lemma-morphisms-from-projective-complex}，有
$\Hom_{D(R)}(P, K_n) = \Hom_{K(R)}(P, K_n)$，故 $\xi = 0$。
因此对所有 $R$-模 $N$ 都有 $\Ext^n_R(M, N) = 0$；参见上面的讨论。
由 Algebra，引理 \ref{algebra-lemma-projective-dimension-ext}，
$M$ 的投射维数 $\leq n - 1$。这对所有有限 $R$-模 $M$ 都成立，
故由 Algebra，引理 \ref{algebra-lemma-finite-gl-dim}，$R$ 具有有限整体维数。
最后由 Algebra，引理
\ref{algebra-lemma-finite-gl-dim-finite-dim-regular} 得到结论。
\end{proof}

\begin{lemma}
\label{more-algebra-lemma-ext-regular}
设 $R$ 为维数 $d < \infty$ 的 Noether 正则环，设
$K, L \in D^-(R)$。
% P02-E0286
假设存在一个 $k$，使得当 $i \leq k$ 时 $H^i(K) = 0$，
且当 $i \geq k - d + 1$ 时 $H^i(L) = 0$。则
$\Hom_{D(R)}(K, L) = 0$。
\end{lemma}

\begin{proof}
设 $K^\bullet$ 为表示 $K$ 的有界上复形，例如当
$i \geq n + 1$ 时 $K^i = 0$。用
$\tau_{\geq k + 1}K^\bullet$ 替换 $K^\bullet$ 后，可以假设当
$i \leq k$ 时 $K^i = 0$。于是利用区别三角
$$
K^n[-n] \to K^\bullet \to \sigma_{\leq n - 1}K^\bullet
$$
可知，只需对 $K^n[-n]$ 和 $\sigma_{\leq n - 1}K^\bullet$
证明该引理。对 $n$ 作归纳，可知只需处理下述情形：$K$ 由复形
$M[-m]$ 表示，其中 $M$ 为某个 $R$-模，且 $m \geq k + 1$。
由 Algebra，引理
\ref{algebra-lemma-finite-gl-dim-finite-dim-regular}，$R$ 的整体维数为
$d$，故 $M$ 有投射解消
$0 \to P_d \to \ldots \to P_0 \to M \to 0$。
令复形 $P^\bullet$ 在次数 $m - i$ 的项为 $P_i$，则它是表示
$M[-m]$ 的有界投射复形。另一方面，可以选取表示 $L$ 的复形
$L^\bullet$，使得当 $i \geq k - d  + 1$ 时 $L^i = 0$。
因此任意复形映射 $P^\bullet \to L^\bullet$ 都为零。
由 Derived Categories，引理
\ref{derived-lemma-morphisms-from-projective-complex} 即得结论。
\end{proof}

\begin{lemma}
\label{more-algebra-lemma-split-complex-regular}
设 $R$ 为维数 $1 \leq d < \infty$ 的 Noether 正则环。
设 $K \in D(R)$ 是完美的，并设 $k \in \mathbf{Z}$，使得当
$i = k - d + 2, \ldots, k$ 时有 $H^i(K) = 0$
（当 $d = 1$ 时此条件为空）。则
$K = \tau_{\leq k - d + 1}K \oplus \tau_{\geq k + 1}K$。
\end{lemma}

\begin{proof}
上同调消失表明存在区别三角
$$
\tau_{\leq k - d + 1}K \to K \to \tau_{\geq k + 1}K \to
(\tau_{\leq k - d + 1}K)[1]
$$
由 Derived Categories，引理 \ref{derived-lemma-split}，
只需证明第三个箭头为零。因此，只需证明
$\Hom_{D(R)}(\tau_{\geq k + 1}K, \allowbreak (\tau_{\leq k - d + 1}K)[1]) = 0$；
这由引理 \ref{more-algebra-lemma-ext-regular} 得到。
\end{proof}

\begin{lemma}
\label{more-algebra-lemma-regular-strong-generator}
设 $R$ 为有限维 Noether 正则环。则 $R$ 是完美对象所成满子范畴
$D_{perf}(R) \subset D(R)$ 的强生成元。
\end{lemma}

\begin{proof}
我们将用到：$D(R)$ 的对象 $K$ 是完美的，当且仅当 $K$ 有界且
各上同调模均为有限模；参见引理 \ref{more-algebra-lemma-regular-perfect}。
三角范畴的强生成元定义于 Derived Categories，定义
\ref{derived-definition-generators}。令 $d = \dim(R)$。

\medskip\noindent
设 $K \in D_{perf}(R)$。我们将证明
$K \in \langle R \rangle_{d + 1}$。由 Algebra，引理
\ref{algebra-lemma-finite-gl-dim-finite-dim-regular}，每个有限
$R$-模的投射维数都 $\leq d$。我们将对 $0 \leq i \leq d$ 归纳证明：
若对所有 $n \in \mathbf{Z}$，$H^n(K)$ 的投射维数都
$\leq i$，则 $K$ 属于 $\langle R \rangle_{i + 1}$。

\medskip\noindent
基础情形 $i = 0$。此时 $H^n(K)$ 是投射维数为 $0$ 的有限
$R$-模。换言之，每个上同调模都是投射 $R$-模。因此，对所有
$i > 0$ 及 $m, n \in \mathbf{Z}$，都有
$\Ext^i_R(H^n(K), H^m(K)) = 0$。由 Derived Categories，引理
\ref{derived-lemma-ext-2-zero-pre}，可知 $K$ 同构于其各上同调模之平移的直和。
由于每个上同调模都是有限投射 $R$-模，它是若干个 $R$ 副本之直和的直和项。
因此根据定义，$K$ 包含在 $\langle R \rangle_1$ 中。

\medskip\noindent
归纳步骤。假设断言对 $i < d$ 成立，并设
$K \in D_{perf}(R)$ 满足：对所有 $n \in \mathbf{Z}$，
$H^n(K)$ 的投射维数都 $\leq i + 1$。选取 $a \leq b$，使得当
$n \not \in [a, b]$ 时 $H^n(K)$ 为零。对每个 $n \in [a, b]$，
选取满射 $F^n \to H^n(K)$，其中 $F^n$ 是有限自由 $R$-模。
由于 $F^n$ 是投射的，可以在 $D(R)$ 中把
$F^n \to H^n(K)$ 提升为映射 $F^n[-n] \to K$
（略去一个小细节）。于是得到态射
$\bigoplus_{a \leq n \leq b} F^n[-n] \to K$，
它在上同调模上是满射。选取 $D(R)$ 中的区别三角
$$
K' \to \bigoplus\nolimits_{a \leq n \leq b} F^n[-n] \to K \to K'[1]
$$
当然，对象 $K'$ 有界，且各上同调模均为有限模。根据所作选择，
上同调长正合序列分解成短正合序列
$$
0 \to H^n(K') \to F^n \to H^n(K) \to 0
$$
由 Algebra，引理
\ref{algebra-lemma-exact-sequence-projective-dimension}，
$H^n(K')$ 的投射维数 $\leq \max(0, i)$。因此
$K' \in \langle R \rangle_{i + 1}$。根据定义，这意味着
$K$ 属于 $\langle R \rangle_{i + 1 + 1}$，正合所需。
\end{proof}

\begin{proposition}
\label{more-algebra-proposition-regular-strong-generator}
设 $R$ 为 Noether 环。下列条件等价：
\begin{enumerate}
\item $R$ 是有限维正则环；
\item $D_{perf}(R)$ 有强生成元；
\item $R$ 是 $D_{perf}(R)$ 的强生成元。
\end{enumerate}
\end{proposition}

\begin{proof}
这是引理 \ref{more-algebra-lemma-perfect-ring-classical-generator}、
\ref{more-algebra-lemma-not-regular-not-strong}、
\ref{more-algebra-lemma-regular-strong-generator} 以及 Derived Categories，引理
\ref{derived-lemma-classical-generator-strong-generator} 的形式推论。
\end{proof}







\section{相对有限表示模}
\label{more-algebra-section-relative-finite-presentation}

\noindent
设 $R$ 为一个环，设 $A \to B$ 为有限型 $R$-代数之间的有限同态，
并设 $M$ 为有限 $B$-模。在这种情形，下面的等价关系{\bf 并不成立}：
$$
M\text{ of finite presentation over }B
\Leftrightarrow
M\text{ of finite presentation over }A
$$
% P02-E0287
一个反例是 $R = k[x_1, x_2, x_3, \ldots]$、$A = R$、
$B = R/(x_i)$ 以及 $M = B$。为了“修正”这一点，
我们引入有限表示的相对概念。

\begin{lemma}
\label{more-algebra-lemma-relatively-finitely-presented}
设 $R \to A$ 为有限型环同态，并设 $M$ 为 $A$-模。
下列条件等价：
\begin{enumerate}
\item 对某个表示 $\alpha : R[x_1, \ldots, x_n] \to A$，
模 $M$ 是有限表示 $R[x_1, \ldots, x_n]$-模；
\item 对所有表示 $\alpha : R[x_1, \ldots, x_n] \to A$，
模 $M$ 是有限表示 $R[x_1, \ldots, x_n]$-模；
\item 对任意满射 $A' \to A$，其中 $A'$ 是有限表示 $R$-代数，
模 $M$ 作为 $A'$-模是有限表示的。
\end{enumerate}
在这种情形，$M$ 是有限表示 $A$-模。
\end{lemma}

\begin{proof}
% P02-E0288
设 $\alpha : R[x_1, \ldots, x_n] \to A$ 与
$\beta : R[y_1, \ldots, y_m] \to A$ 为两个表示。
选取 $f_j \in R[x_1, \ldots, x_n]$，使得
$\alpha(f_j) = \beta(y_j)$；并选取
$g_i \in R[y_1, \ldots, y_m]$，使得
$\beta(g_i) = \alpha(x_i)$。于是得到交换图
$$
\xymatrix{
R[x_1, \ldots, x_n, y_1, \ldots, y_m]
\ar[d]^{x_i \mapsto g_i} \ar[rr]_-{y_j \mapsto f_j} & &
R[x_1, \ldots, x_n] \ar[d] \\
R[y_1, \ldots, y_m] \ar[rr] & & A
}
$$
对 (1) 与 (2) 的等价性应用 Algebra，引理
\ref{algebra-lemma-finitely-presented-over-subring} 和
\ref{algebra-lemma-finite-finitely-presented-extension} 即可。
选取表示 $A' = R[x_1, \ldots, x_n]/(f_1, \ldots, f_m)$，
并应用 Algebra，引理
\ref{algebra-lemma-finite-finitely-presented-extension}，可知
$M$ 作为 $A'$-模有限表示，当且仅当 $M$ 作为
$R[x_1, \ldots, x_n]$-模有限表示；这证明了 (2) 与 (3) 等价。
\end{proof}

\begin{definition}
\label{more-algebra-definition-relatively-finitely-presented}
设 $R \to A$ 为有限型环同态，设 $M$ 为 $A$-模。
若引理 \ref{more-algebra-lemma-relatively-finitely-presented} 中的等价条件成立，
则称 $M$ 是{it 相对于 $R$ 有限表示的 $A$-模}。
\end{definition}

\noindent
注意，若 $R \to A$ 是有限表示的，则 $M$ 是相对于 $R$ 有限表示的
$A$-模，当且仅当 $M$ 是有限表示 $A$-模。同样显然，$A$ 作为
$A$-模相对于 $R$ 有限表示，当且仅当 $A$ 在 $R$ 上有限表示。
若 $R$ 是 Noether 的，该概念没有新内容。现在可以表述我们所寻求的结果。

\begin{lemma}
\label{more-algebra-lemma-finite-extension}
设 $R$ 为一个环，设 $A \to B$ 为有限型 $R$-代数之间的有限同态，
并设 $M$ 为 $B$-模。则 $M$ 作为 $A$-模相对于 $R$ 有限表示，
当且仅当 $M$ 作为 $B$-模相对于 $R$ 有限表示。
\end{lemma}

\begin{proof}
选取满射 $R[x_1, \ldots, x_n] \to A$。选取
$y_1, \ldots, y_m \in B$，使它们在 $A$ 上生成 $B$。
由于 $A \to B$ 是有限的，每个 $y_i$ 都满足一个系数在 $A$ 中的首一方程。
因此，可以找到首一多项式
$P_j(T) \in R[x_1, \ldots, x_n][T]$，使得在 $B$ 中
$P_j(y_j) = 0$。于是得到交换图
$$
\xymatrix{
R[x_1, \ldots, x_n] \ar[d] \ar[r] &
R[x_1, \ldots, x_n, y_1, \ldots, y_m]/(P_j(y_j)) \ar[d] \\
A \ar[r] & B
}
$$
上方箭头是有限且有限表示的环同态，故由 Algebra，引理
\ref{algebra-lemma-finite-finitely-presented-extension} 和定义即得结论。
\end{proof}

\noindent
由这个结果可知，相对概念是合理的，并且对于 $R$ 上有限型环之间的有限同态
表现良好。它还在局部化和基变换下稳定，并且能够良好地“粘合”。

\begin{lemma}
\label{more-algebra-lemma-localize-relative-finite-presentation}
设 $R$ 为一个环，$f \in R$ 为元素，$R_f \to A$ 为有限型环同态，
$g \in A$，并设 $M$ 为 $A$-模。
% P02-E0289
若 $M$ 相对于 $R_f$ 有限表示，则 $M_g$ 是相对于 $R$ 有限表示的
$A_g$-模。
\end{lemma}

\begin{proof}
选取表示 $R_f[x_1, \ldots, x_n] \to A$。写成
$R_f = R[x_0]/(fx_0 - 1)$。考虑表示
$R[x_0, x_1, \ldots, x_n, x_{n + 1}] \to A_g$；
它延拓给定映射，把 $x_0$ 映到 $1/f$ 的像，并把
$x_{n + 1}$ 映到 $1/g$。选取映到 $g$ 的
$g' \in R[x_0, x_1, \ldots, x_n]$（这是可能的）。假设
$$
R_f[x_1, \ldots, x_n]^{\oplus s} \to
R_f[x_1, \ldots, x_n]^{\oplus t} \to M \to 0
$$
是由矩阵 $(h_{ij})$ 给出的 $M$ 的表示。选取映到 $h_{ij}$ 的
$h'_{ij} \in R[x_0, x_1, \ldots, x_n]$。则
$$
R[x_0, x_1, \ldots, x_n, x_{n + 1}]^{\oplus s + 2t} \to
R[x_0, x_1, \ldots, x_n, x_{n + 1}]^{\oplus t} \to M_g \to 0
$$
% P02-E0290
是 $M_f$ 的一个表示。这里，定义该映射的
$t \times (s + 2t)$ 矩阵的第一个 $t \times s$ 分块为矩阵
$h'_{ij}$；
% P02-E0291
第二个 $t \times t$ 分块为 $(x_0f - )I_t$；第三个分块为
$(x_{n + 1}g' - 1)I_t$。
\end{proof}

\begin{lemma}
\label{more-algebra-lemma-base-change-relative-finite-presentation}
设 $R \to A$ 为有限型环同态，并设 $M$ 是相对于 $R$ 有限表示的
$A$-模。对任意环同态 $R \to R'$，$A \otimes_R R'$-模
% P02-E0292
$$
M \otimes_A A' = M \otimes_R R'
$$
相对于 $R'$ 有限表示。
\end{lemma}

\begin{proof}
选取满射 $R[x_1, \ldots, x_n] \to A$，并选取表示
$$
R[x_1, \ldots, x_n]^{\oplus s} \to
R[x_1, \ldots, x_n]^{\oplus t} \to M \to 0
$$
则
$$
R'[x_1, \ldots, x_n]^{\oplus s} \to
R'[x_1, \ldots, x_n]^{\oplus t} \to M \otimes_R R' \to 0
$$
是基变换的一个表示，结论得证。
\end{proof}

\begin{lemma}
\label{more-algebra-lemma-pull-relative-finite-presentation}
设 $R \to A$ 为有限型环同态，设 $M$ 是相对于 $R$ 有限表示的
$A$-模，并设 $A \to A'$ 为有限表示环同态。则 $A'$-模
$M \otimes_A A'$ 相对于 $R$ 有限表示。
\end{lemma}

\begin{proof}
选取满射 $R[x_1, \ldots, x_n] \to A$，并选取表示
$A' = A[y_1, \ldots, y_m]/(g_1, \ldots, g_l)$。
选取映到 $g_i$ 的
$g'_i \in R[x_1, \ldots, x_n, y_1, \ldots, y_m]$。设
$$
R[x_1, \ldots, x_n]^{\oplus s} \to
R[x_1, \ldots, x_n]^{\oplus t} \to M \to 0
$$
是由矩阵 $(h_{ij})$ 给出的 $M$ 的表示。则
% P02-E0293
$$
R[x_1, \ldots, x_n, y_1, \ldots, y_m]^{\oplus s + tl} \to
R[x_0, x_1, \ldots, x_n, y_1, \ldots, y_m]^{\oplus t} \to M \otimes_A A' \to 0
$$
是 $M \otimes_A A'$ 的一个表示。这里，定义该映射的
$t \times (s + lt)$ 矩阵的第一个 $t \times s$ 分块为矩阵
$h_{ij}$，随后是 $l$ 个大小为 $t \times t$ 的分块
$g'_iI_t$。
\end{proof}

\begin{lemma}
\label{more-algebra-lemma-composition-relative-finite-presentation}
设 $R \to A \to B$ 为有限型环同态，设 $M$ 为 $B$-模。
若 $M$ 相对于 $A$ 有限表示，且 $A$ 在 $R$ 上有限表示，
则 $M$ 相对于 $R$ 有限表示。
\end{lemma}

\begin{proof}
选取满射 $A[x_1, \ldots, x_n] \to B$，并选取由矩阵
$(h_{ij})$ 给出的表示
$$
A[x_1, \ldots, x_n]^{\oplus s} \to
A[x_1, \ldots, x_n]^{\oplus t} \to M \to 0
$$
再选取表示
$$
A = R[y_1, \ldots, y_m]/(g_1, \ldots, g_u).
$$
选取映到 $h_{ij}$ 的
$h'_{ij} \in R[y_1, \ldots, y_m, x_1, \ldots, x_n]$。
于是得到表示
$$
R[y_1, \ldots, y_m, x_1, \ldots, x_n]^{\oplus s + tu} \to
R[y_1, \ldots, y_m, x_1, \ldots, x_n]^{\oplus t} \to M \to 0
$$
其中 $t \times (s + tu)$ 矩阵的第一个 $t \times s$ 分块由
$h'_{ij}$ 构成，随后是 $u$ 个大小为 $t \times t$ 的分块
$g_iI_t$，$i = 1, \ldots, u$。
\end{proof}

\begin{lemma}
\label{more-algebra-lemma-glue-relative-finite-presentation}
设 $R \to A$ 为有限型环同态，设 $M$ 为 $A$-模，并设
$f_1, \ldots, f_r \in A$ 生成单位理想。下列条件等价：
\begin{enumerate}
\item 每个 $M_{f_i}$ 都相对于 $R$ 有限表示；
\item $M$ 相对于 $R$ 有限表示。
\end{enumerate}
\end{lemma}

\begin{proof}
(2) $\Rightarrow$ (1) 是引理
\ref{more-algebra-lemma-localize-relative-finite-presentation} 的内容。
假设 (1)。在 $A$ 中写成 $1 = \sum f_ig_i$。
% P02-E0294
选取满射
$R[x_1, \ldots, x_n, y_1, \ldots, y_r, z_1, \ldots, z_r] \to A$，
使得 $y_i$ 映到 $f_i$，而 $z_i$ 映到 $g_i$。于是可见存在满射
$$
P = R[x_1, \ldots, x_n, y_1, \ldots, y_r, z_1, \ldots, z_r]/(\sum y_iz_i - 1)
\longrightarrow
A.
$$
由引理 \ref{more-algebra-lemma-relatively-finitely-presented}，$M_{f_i}$ 是有限表示
$A_{f_i}$-模；故由 Algebra，引理 \ref{algebra-lemma-cover}，
$M$ 是有限表示 $A$-模。因此，$M$ 是有限 $P$-模（$P$ 如上）。
选取满射 $P^{\oplus t} \to M$。必须证明该映射的核 $K$ 是有限
$P$-模。由于 $P_{y_i}$ 满射到 $A_{f_i}$，由引理
\ref{more-algebra-lemma-relatively-finitely-presented} 和 Algebra，引理
\ref{algebra-lemma-extension} 可知，局部化 $K_{y_i}$ 是有限生成
$P_{y_i}$-模。选取元素
$k_{i, j} \in K$，$i = 1, \ldots, r$，$j = 1, \ldots, s_i$，
使得 $k_{i, j}$ 在 $K_{y_i}$ 中的像生成该模。令
$K' \subset K$ 为这些元素 $k_{i, j}$ 生成的 $P$-模。
则对所有 $i$，模 $K/K'$ 在 $y_i$ 处的局部化都为零。由于
$(y_1, \ldots, y_r) = P$，可知 $K/K' = 0$，正合所需。
\end{proof}

\begin{lemma}
\label{more-algebra-lemma-ses-relatively-finite-presentation}
设 $R \to A$ 为有限型环同态，并设
$0 \to M' \to M \to M'' \to 0$ 为 $A$-模的短正合序列。
\begin{enumerate}
\item 若 $M', M''$ 相对于 $R$ 有限表示，则 $M$ 也如此。
\item 若 $M'$ 是有限型 $A$-模，且 $M$ 相对于 $R$ 有限表示，
则 $M''$ 相对于 $R$ 有限表示。
\end{enumerate}
\end{lemma}

\begin{proof}
% P02-E0295
这立即由 Algebra，引理 \ref{algebra-lemma-extension} 得到。
\end{proof}

\begin{lemma}
\label{more-algebra-lemma-sum-relatively-finite-presentation}
设 $R \to A$ 为有限型环同态，设 $M, M'$ 为 $A$-模。
若 $M \oplus M'$ 相对于 $R$ 有限表示，则 $M$ 与 $M'$ 也如此。
\end{lemma}

\begin{proof}
略。
\end{proof}






\section{相对伪凝聚模}
\label{more-algebra-section-relative-pseudo-coherent}

\noindent
本节给出第 \ref{more-algebra-section-relative-finite-presentation} 节关于伪凝聚性的类似理论。

\begin{lemma}
\label{more-algebra-lemma-pull-push}
设 $R$ 为一个环，设 $K^\bullet$ 为 $R$-模复形。
考虑把 $x$ 映到零的 $R$-代数同态 $R[x] \to R$。则在 $D(R)$ 中
$$
K^\bullet \otimes_{R[x]}^{\mathbf{L}} R \cong K^\bullet \oplus K^\bullet[1]
$$
\end{lemma}

\begin{proof}
在 $R$ 上选取 K-平坦解消 $P^\bullet \to K^\bullet$，使得对所有
$n$，$P^n$ 都是平坦 $R$-模；参见引理
\ref{more-algebra-lemma-K-flat-resolution}。于是 $P^\bullet \otimes_R R[x]$ 是
$R[x]$-模的 K-平坦复形，且各项都是平坦 $R[x]$-模；参见引理
\ref{more-algebra-lemma-base-change-K-flat} 和 Algebra，引理
\ref{algebra-lemma-flat-base-change}。特别地，
$x : P^n \otimes_R R[x] \to P^n \otimes_R R[x]$ 是单射，
其余核同构于 $P^n$。因此
$$
P^\bullet \otimes_R R[x] \xrightarrow{x} P^\bullet \otimes_R R[x]
$$
是 $R[x]$-模的双复形，其相伴全复形拟同构于 $P^\bullet$，
从而拟同构于 $K^\bullet$。此外，例如由引理
\ref{more-algebra-lemma-tensor-product-K-flat} 或引理
\ref{more-algebra-lemma-K-flat-two-out-of-three}，该相伴全复形是
$R[x]$-模的 K-平坦复形。因此
\begin{align*}
K^\bullet \otimes_{R[x]}^{\mathbf{L}} R
& \cong
\text{Tot}(P^\bullet \otimes_R R[x] \xrightarrow{x} P^\bullet \otimes_R R[x])
\otimes_{R[x]} R =
\text{Tot}(P^\bullet \xrightarrow{0} P^\bullet) \\
& = P^\bullet \oplus P^\bullet[1] \cong K^\bullet \oplus K^\bullet[1]
\end{align*}
正合所需。
\end{proof}

\begin{lemma}
\label{more-algebra-lemma-add-variable-pseudo-coherent}
设 $R$ 为一个环，$K^\bullet$ 为 $R$-模复形，并设
$m \in \mathbf{Z}$。考虑把 $x$ 映到零的 $R$-代数同态
$R[x] \to R$。则 $K^\bullet$ 作为 $R$-模复形是
$m$-伪凝聚的，当且仅当 $K^\bullet$ 作为 $R[x]$-模复形是
$m$-伪凝聚的。
\end{lemma}

\begin{proof}
这是引理 \ref{more-algebra-lemma-finite-push-pseudo-coherent} 的特例。
下面也给出另一种证明。

\medskip\noindent
注意 $0 \to R[x] \to R[x] \to R \to 0$ 正合。
因此 $R$ 作为 $R[x]$-模是伪凝聚的。故该引理的一个方向由引理
\ref{more-algebra-lemma-finite-push-pseudo-coherent} 得到。
为证明另一方向，假设 $K^\bullet$ 作为 $R[x]$-模复形是
$m$-伪凝聚的。由引理 \ref{more-algebra-lemma-pull-pseudo-coherent}，
$K^\bullet \otimes^{\mathbf{L}}_{R[x]} R$ 作为 $R$-模复形是
$m$-伪凝聚的。由引理 \ref{more-algebra-lemma-pull-push}，
$K^\bullet \oplus K^\bullet[1]$ 作为 $R$-模复形是 $m$-伪凝聚的。
最后，由引理 \ref{more-algebra-lemma-summands-pseudo-coherent} 可知，
$K^\bullet$ 作为 $R$-模复形是 $m$-伪凝聚的。
\end{proof}

\begin{lemma}
\label{more-algebra-lemma-relatively-pseudo-coherent}
设 $R \to A$ 为有限型环同态，设 $K^\bullet$ 为 $A$-模复形，
并设 $m \in \mathbf{Z}$。下列条件等价：
\begin{enumerate}
\item 对某个表示 $\alpha : R[x_1, \ldots, x_n] \to A$，
复形 $K^\bullet$ 是 $R[x_1, \ldots, x_n]$-模的
$m$-伪凝聚复形；
\item 对所有表示 $\alpha : R[x_1, \ldots, x_n] \to A$，
复形 $K^\bullet$ 是 $R[x_1, \ldots, x_n]$-模的
$m$-伪凝聚复形。
\end{enumerate}
特别地，对伪凝聚性也有同样的等价性。
\end{lemma}

\begin{proof}
% P02-E0296 P02-E0297
设 $\alpha : R[x_1, \ldots, x_n] \to A$ 与
$\beta : R[y_1, \ldots, y_m] \to A$ 为两个表示。
选取 $f_j \in R[x_1, \ldots, x_n]$，使得
$\alpha(f_j) = \beta(y_j)$；并选取
$g_i \in R[y_1, \ldots, y_m]$，使得
$\beta(g_i) = \alpha(x_i)$。于是得到交换图
$$
\xymatrix{
R[x_1, \ldots, x_n, y_1, \ldots, y_m]
\ar[d]^{x_i \mapsto g_i} \ar[rr]_-{y_j \mapsto f_j} & &
R[x_1, \ldots, x_n] \ar[d] \\
R[y_1, \ldots, y_m] \ar[rr] & & A
}
$$
作坐标变换后，环同态
$R[x_1, \ldots, x_n, y_1, \ldots, y_m] \to R[x_1, \ldots, x_n]$
同构于把每个 $y_i$ 映到零的环同态。图中的左侧竖直映射亦然。
因此对变量数作归纳，该引理由引理
\ref{more-algebra-lemma-add-variable-pseudo-coherent} 得到。
伪凝聚情形则由此和引理 \ref{more-algebra-lemma-pseudo-coherent} 得到。
\end{proof}

\begin{definition}
\label{more-algebra-definition-relatively-pseudo-coherent}
设 $R \to A$ 为有限型环同态，设 $K^\bullet$ 为 $A$-模复形，
设 $M$ 为 $A$-模，并设 $m \in \mathbf{Z}$。
\begin{enumerate}
\item 若引理 \ref{more-algebra-lemma-relatively-pseudo-coherent} 中的等价条件成立，
则称 $K^\bullet$ {\it 相对于 $R$ 是 $m$-伪凝聚的}。
\item 若对所有 $m \in \mathbf{Z}$，$K^\bullet$ 相对于 $R$ 都是
$m$-伪凝聚的，则称 $K^\bullet$ {\it 相对于 $R$ 是伪凝聚的}。
\item 若 $M[0]$ 相对于 $R$ 是 $m$-伪凝聚的，则称 $M$
{\it 相对于 $R$ 是 $m$-伪凝聚的}。
\item 若 $M[0]$ 相对于 $R$ 是伪凝聚的，则称 $M$
{\it 相对于 $R$ 是伪凝聚的}。
\end{enumerate}
\end{definition}

\noindent
% P02-E0298
第 (2) 项意指：对任意满射 $R[y_1, \ldots, y_m] \to A$，
$K^\bullet$ 作为 $R[x_1, \ldots, x_n]$-模复形是伪凝聚的；
参见引理 \ref{more-algebra-lemma-pseudo-coherent}。该定义具有下述良好性质。

\begin{lemma}
\label{more-algebra-lemma-finite-extension-pseudo-coherent}
设 $R$ 为一个环，设 $A \to B$ 为有限型 $R$-代数之间的有限同态，
设 $m \in \mathbf{Z}$，并设 $K^\bullet$ 为 $B$-模复形。
则 $K^\bullet$ 相对于 $R$ 是 $m$-伪凝聚的（相应地，伪凝聚的），
当且仅当把 $K^\bullet$ 视为 $A$-模复形后，它相对于 $R$ 是
$m$-伪凝聚的（伪凝聚的）。
\end{lemma}

\begin{proof}
选取满射 $R[x_1, \ldots, x_n] \to A$。
% P02-E0299
选取 $y_1, \ldots, y_m \in B$，使它们在 $A$ 上生成 $B$。
由于 $A \to B$ 是有限的，每个 $y_i$ 都满足系数在 $A$ 中的首一方程。
因此，可以找到首一多项式
$P_j(T) \in R[x_1, \ldots, x_n][T]$，使得在 $B$ 中
$P_j(y_j) = 0$。于是得到交换图
$$
\xymatrix{
& R[x_1, \ldots, x_n, y_1, \ldots, y_m] \ar[d] \\
R[x_1, \ldots, x_n] \ar[d] \ar[r] &
R[x_1, \ldots, x_n, y_1, \ldots, y_m]/(P_j(y_j)) \ar[d] \\
A \ar[r] & B
}
$$
顶部水平箭头和右上方竖直箭头满足引理
\ref{more-algebra-lemma-finite-push-pseudo-coherent} 的假设。因此，
$K^\bullet$ 作为 $R[x_1, \ldots, x_n]$-模复形是
$m$-伪凝聚的（相应地，伪凝聚的），当且仅当 $K^\bullet$ 作为
$R[x_1, \ldots, x_n, y_1, \ldots, y_m]$-模复形是
$m$-伪凝聚的（相应地，伪凝聚的）。
\end{proof}

\begin{lemma}
\label{more-algebra-lemma-cone-relatively-pseudo-coherent}
设 $R$ 为一个环，设 $R \to A$ 为有限型环同态，设
$m \in \mathbf{Z}$，并设 $(K^\bullet, L^\bullet, M^\bullet, f, g, h)$
为 $D(A)$ 中的区别三角。
\begin{enumerate}
\item 若 $K^\bullet$ 相对于 $R$ 是 $(m + 1)$-伪凝聚的，且
$L^\bullet$ 相对于 $R$ 是 $m$-伪凝聚的，则 $M^\bullet$
相对于 $R$ 是 $m$-伪凝聚的。
\item 若 $K^\bullet, M^\bullet$ 相对于 $R$ 是 $m$-伪凝聚的，
则 $L^\bullet$ 相对于 $R$ 是 $m$-伪凝聚的。
\item 若 $L^\bullet$ 相对于 $R$ 是 $(m + 1)$-伪凝聚的，且
$M^\bullet$ 相对于 $R$ 是 $m$-伪凝聚的，则 $K^\bullet$
相对于 $R$ 是 $(m + 1)$-伪凝聚的。
\end{enumerate}
此外，若 $K^\bullet, L^\bullet, M^\bullet$ 中有两个相对于 $R$
是伪凝聚的，则第三个也如此。
% P02-E0300
\end{lemma}

\begin{proof}
% P02-E0301
这立即由引理 \ref{more-algebra-lemma-cone-pseudo-coherent} 和定义得到。
\end{proof}

\begin{lemma}
\label{more-algebra-lemma-rel-n-pseudo-module}
设 $R \to A$ 为有限型环同态，设 $M$ 为 $A$-模。则
\begin{enumerate}
\item $M$ 相对于 $R$ 是 $0$-伪凝聚的，当且仅当 $M$ 是有限型
$A$-模；
% P02-E0302
\item $M$ 相对于 $R$ 是 $(-1)$-伪凝聚的，当且仅当 $M$
相对于 $R$ 有限表示；
% P02-E0303
\item $M$ 相对于 $R$ 是 $(-d)$-伪凝聚的，当且仅当对每个满射
$R[x_1, \ldots, x_n] \to A$，都存在长度为 $d$ 的解消
$$
R[x_1, \ldots, x_n]^{\oplus a_d} \to R[x_1, \ldots, x_n]^{\oplus a_{d - 1}}
\to \ldots \to R[x_1, \ldots, x_n]^{\oplus a_0} \to M \to 0
$$
；
\item $M$ 相对于 $R$ 是伪凝聚的，当且仅当对每个表示
$R[x_1, \ldots, x_n] \to A$，都存在由有限自由
$R[x_1, \ldots, x_n]$-模组成的无穷解消
$$
\ldots \to R[x_1, \ldots, x_n]^{\oplus a_1} \to
R[x_1, \ldots, x_n]^{\oplus a_0} \to M \to 0
$$
。
\end{enumerate}
\end{lemma}

\begin{proof}
% P02-E0304
这立即由引理 \ref{more-algebra-lemma-n-pseudo-module} 和定义得到。
\end{proof}

\begin{lemma}
\label{more-algebra-lemma-summands-relative-pseudo-coherent}
设 $R \to A$ 为有限型环同态，设 $m \in \mathbf{Z}$，并设
$K^\bullet, L^\bullet \in D(A)$。若 $K^\bullet \oplus L^\bullet$
相对于 $R$ 是 $m$-伪凝聚的（相应地，伪凝聚的），则
$K^\bullet$ 与 $L^\bullet$ 也如此。
\end{lemma}

\begin{proof}
% P02-E0304
这立即由引理 \ref{more-algebra-lemma-summands-pseudo-coherent} 和定义得到。
\end{proof}

\begin{lemma}
\label{more-algebra-lemma-complex-relative-pseudo-coherent-modules}
设 $R \to A$ 为有限型环同态，设 $m \in \mathbf{Z}$，并设
$K^\bullet$ 为有界上 $A$-模复形，使得对所有 $i$，$K^i$
相对于 $R$ 是 $(m - i)$-伪凝聚的。则 $K^\bullet$ 相对于
$R$ 是 $m$-伪凝聚的。
% P02-E0305
特别地，若 $K^\bullet$ 是一个有界上复形，其各项都是相对于 $R$
伪凝聚的 $A$-模，则 $K^\bullet$ 相对于 $R$ 是伪凝聚的。
\end{lemma}

\begin{proof}
% P02-E0304
这立即由引理 \ref{more-algebra-lemma-complex-pseudo-coherent-modules} 和定义得到。
\end{proof}

\begin{lemma}
\label{more-algebra-lemma-cohomology-relative-pseudo-coherent}
设 $R \to A$ 为有限型环同态，设 $m \in \mathbf{Z}$，并设
$K^\bullet \in D^{-}(A)$，使得对所有 $i$，$H^i(K^\bullet)$
相对于 $R$ 是 $(m - i)$-伪凝聚的（相应地，伪凝聚的）。
则 $K^\bullet$ 相对于 $R$ 是 $m$-伪凝聚的
（相应地，伪凝聚的）。
\end{lemma}

\begin{proof}
% P02-E0304
这立即由引理 \ref{more-algebra-lemma-cohomology-pseudo-coherent} 和定义得到。
\end{proof}

\begin{lemma}
\label{more-algebra-lemma-localize-relative-pseudo-coherent}
% P02-E0306
设 $R$ 为一个环，$f \in R$ 为元素，$R_f \to A$ 为有限型环同态，
$g \in A$，并设 $K^\bullet$ 为 $A$-模复形。若 $K^\bullet$
相对于 $R_f$ 是 $m$-伪凝聚的（相应地，伪凝聚的），则
$K^\bullet \otimes_A A_g$ 相对于 $R$ 是 $m$-伪凝聚的
（相应地，伪凝聚的）。
\end{lemma}

\begin{proof}
首先证明 $K^\bullet$ 相对于 $R$ 是 $m$-伪凝聚的。
具体地，设 $R_f[x_1, \ldots, x_n] \to A$ 为满射。写成
$R_f = R[x_0]/(fx_0 - 1)$。则 $R[x_0, x_1, \ldots, x_n] \to A$
是满射，而 $R_f[x_1, \ldots, x_n]$ 作为
$R[x_0, \ldots, x_n]$-模是伪凝聚的。因此由引理
\ref{more-algebra-lemma-finite-push-pseudo-coherent}，$K^\bullet$ 作为
$R[x_0, x_1, \ldots, x_n]$-模复形是 $m$-伪凝聚的。

\medskip\noindent
选取映到 $g \in A$ 的元素
$g' \in R[x_0, x_1, \ldots, x_n]$。由引理
\ref{more-algebra-lemma-pull-pseudo-coherent}，可知
% P02-E0307
\begin{align*}
K^\bullet \otimes_{R[x_0, x_1, \ldots, x_n]}^{\mathbf{L}}
R[x_0, x_1, \ldots, x_n, \frac{1}{g'}] & =
K^\bullet \otimes_{R[x_0, x_1, \ldots, x_n]}
R[x_0, x_1, \ldots, x_n, \frac{1}{g'}] \\
& = K^\bullet \otimes_A A_f
\end{align*}
作为 $R[x_0, x_1, \ldots, x_n, \frac{1}{g'}]$-模复形是
$m$-伪凝聚的。
% P02-E0308
写成
$$
R[x_0, x_1, \ldots, x_n, \frac{1}{g'}] =
R[x_0, \ldots, x_n, x_{n + 1}]/(x_{n + 1}g' - 1).
$$
由于 $R[x_0, x_1, \ldots, x_n, \frac{1}{g'}]$ 作为
$R[x_0, \ldots, x_n, x_{n + 1}]$-模是伪凝聚的，由引理
\ref{more-algebra-lemma-finite-push-pseudo-coherent} 可知，
$K^\bullet \otimes_A A_g$ 作为
$R[x_0, \ldots, x_n, x_{n + 1}]$-模复形是 $m$-伪凝聚的，
正合所需。
\end{proof}

\begin{lemma}
\label{more-algebra-lemma-base-change-relative-pseudo-coherent}
设 $R \to A$ 为有限型环同态，设 $m \in \mathbf{Z}$，并设
$K^\bullet$ 为相对于 $R$ 是 $m$-伪凝聚的（相应地，伪凝聚的）
$A$-模复形。设 $R \to R'$ 为环同态，使得 $A$ 与 $R'$ 在 $R$
上 Tor 独立。令 $A' = A \otimes_R R'$。则
$K^\bullet \otimes_A^{\mathbf{L}} A'$ 相对于 $R'$ 是
$m$-伪凝聚的（相应地，伪凝聚的）。
\end{lemma}

\begin{proof}
选取满射 $R[x_1, \ldots, x_n] \to A$。注意
$$
K^\bullet \otimes_A^{\mathbf{L}} A' =
K^\bullet \otimes_R^{\mathbf{L}} R' =
K^\bullet \otimes_{R[x_1, \ldots, x_n]}^{\mathbf{L}} R'[x_1, \ldots, x_n]
$$
这是两次应用引理 \ref{more-algebra-lemma-base-change-comparison} 得到的。
于是由引理 \ref{more-algebra-lemma-pull-pseudo-coherent} 即得结论。
\end{proof}

\begin{lemma}
\label{more-algebra-lemma-pull-relative-pseudo-coherent}
设 $R \to A \to B$ 为有限型环同态，设 $m \in \mathbf{Z}$，
并设 $K^\bullet$ 为 $A$-模复形。假设 $B$ 作为 $B$-模相对于
$A$ 是伪凝聚的。若 $K^\bullet$ 相对于 $R$ 是 $m$-伪凝聚的
（相应地，伪凝聚的），则 $K^\bullet \otimes_A^{\mathbf{L}} B$
相对于 $R$ 是 $m$-伪凝聚的（相应地，伪凝聚的）。
\end{lemma}

\begin{proof}
% P02-E0309
选取满射 $A[y_1, \ldots, y_m] \to B$，并选取满射
$R[x_1, \ldots, x_n] \to A$。合起来得到满射
% P02-E0310
$R[x_1, \ldots, x_n, y_1, \ldots y_m] \to B$。
% P02-E0311
选取 $B$ 的解消 $E^\bullet \to B$，其中它是有限自由
$A[y_1, \ldots, y_n]$-模复形（由关于环同态 $A \to B$ 的假设，
这是可能的）。可以假设 $K^\bullet$ 是平坦 $A$-模组成的有界上复形。
于是
\begin{align*}
K^\bullet \otimes_A^{\mathbf{L}} B & =
\text{Tot}(K^\bullet \otimes_A B[0]) \\
& = \text{Tot}(K^\bullet \otimes_A A[y_1, \ldots, y_m]
\otimes_{A[y_1, \ldots, y_m]} B[0]) \\
& \cong
\text{Tot}\left(
(K^\bullet \otimes_A A[y_1, \ldots, y_m])
\otimes_{A[y_1, \ldots, y_m]} E^\bullet
\right) \\
& =
\text{Tot}(K^\bullet \otimes_A E^\bullet)
\end{align*}
在 $D(A[y_1, \ldots, y_m])$ 中成立。同构 $\cong$ 来自引理
\ref{more-algebra-lemma-derived-tor-quasi-isomorphism} 的一次应用。
因此，必须证明 $\text{Tot}(K^\bullet \otimes_A E^\bullet)$ 作为
$R[x_1, \ldots, x_n, y_1, \ldots y_m]$-模复形是 $m$-伪凝聚的。
注意，$\text{Tot}(K^\bullet \otimes_A E^\bullet)$ 有一个由子复形组成的滤过，
其逐次商为复形 $K^\bullet \otimes_A E^i[-i]$。当 $i \ll 0$ 时，
复形 $K^\bullet \otimes_A E^i[-i]$ 在次数 $\leq m$ 的上同调为零，
因而是 $m$-伪凝聚的（在任意环上皆然）。因此，应用引理
\ref{more-algebra-lemma-cone-relatively-pseudo-coherent} 并作归纳，只需证明
% P02-E0312
$K^\bullet \otimes_A E^i[-i]$ 对所有 $i$ 都相对于 $R$ 是伪凝聚的。
注意当 $i > 0$ 时 $E^i = 0$。又因 $E^i$ 是有限自由的，
这归结为证明 $K^\bullet \otimes_A A[y_1, \ldots, y_m]$
相对于 $R$ 是 $m$-伪凝聚的；例如这由引理
\ref{more-algebra-lemma-base-change-relative-pseudo-coherent} 得到。
\end{proof}

\begin{lemma}
\label{more-algebra-lemma-pull-relative-pseudo-coherent-module}
设 $R \to A \to B$ 为有限型环同态，设 $m \in \mathbf{Z}$，
并设 $M$ 为 $A$-模。假设 $B$ 在 $A$ 上平坦，且 $B$ 作为
$B$-模相对于 $A$ 是伪凝聚的。若 $M$ 相对于 $R$ 是
$m$-伪凝聚的（相应地，伪凝聚的），则 $M \otimes_A B$
相对于 $R$ 是 $m$-伪凝聚的（相应地，伪凝聚的）。
\end{lemma}

\begin{proof}
% P02-E0313
这立即由引理 \ref{more-algebra-lemma-pull-relative-pseudo-coherent} 得到。
\end{proof}

\begin{lemma}
\label{more-algebra-lemma-composition-relative-pseudo-coherent}
设 $R$ 为一个环，设 $A \to B$ 为有限型 $R$-代数之间的同态，
设 $m \in \mathbf{Z}$，并设 $K^\bullet$ 为 $B$-模复形。
假设 $A$ 相对于 $R$ 是伪凝聚的。则下列条件等价：
\begin{enumerate}
\item $K^\bullet$ 相对于 $A$ 是 $m$-伪凝聚的
（相应地，伪凝聚的）；
\item $K^\bullet$ 相对于 $R$ 是 $m$-伪凝聚的
（相应地，伪凝聚的）。
\end{enumerate}
\end{lemma}

\begin{proof}
选取满射 $R[x_1, \ldots, x_n] \to A$。
% P02-E0314
选取满射 $A[y_1, \ldots, y_m] \to B$。于是得到满射
$$
R[x_1, \ldots, x_n, y_1, \ldots, y_m] \to A[y_1, \ldots, y_m]
$$
它是 $R[x_1, \ldots, x_n] \to A$ 的平坦基变换。
根据假设，$A$ 是伪凝聚 $R[x_1, \ldots, x_n]$-模；故由引理
\ref{more-algebra-lemma-flat-base-change-pseudo-coherent}，可知
$A[y_1, \ldots, y_m]$ 在
$R[x_1, \ldots, x_n, y_1, \ldots, y_m]$ 上是伪凝聚的。
因此，该引理由引理 \ref{more-algebra-lemma-finite-push-pseudo-coherent} 和定义得到。
\end{proof}

\begin{lemma}
\label{more-algebra-lemma-glue-relative-pseudo-coherent}
设 $R \to A$ 为有限型环同态，设 $K^\bullet$ 为 $A$-模复形，
设 $m \in \mathbf{Z}$，并设 $f_1, \ldots, f_r \in A$ 生成单位理想。
下列条件等价：
\begin{enumerate}
\item 每个 $K^\bullet \otimes_A A_{f_i}$ 都相对于 $R$ 是
$m$-伪凝聚的；
\item $K^\bullet$ 相对于 $R$ 是 $m$-伪凝聚的。
\end{enumerate}
对于相对于 $R$ 的伪凝聚性，也有同样的等价性。
\end{lemma}

\begin{proof}
(2) $\Rightarrow$ (1) 是引理
\ref{more-algebra-lemma-localize-relative-pseudo-coherent} 的内容。
假设 (1)。在 $A$ 中写成 $1 = \sum f_ig_i$。
% P02-E0315
选取满射
$R[x_1, \ldots, x_n, y_1, \ldots, y_r, z_1, \ldots, z_r] \to A$，
使得 $y_i$ 映到 $f_i$，而 $z_i$ 映到 $g_i$。于是可见存在满射
$$
P = R[x_1, \ldots, x_n, y_1, \ldots, y_r, z_1, \ldots, z_r]/(\sum y_iz_i - 1)
\longrightarrow
A.
$$
注意，$P$ 作为
$R[x_1, \ldots, x_n, y_1, \ldots, y_r, z_1, \ldots, z_r]$-模是伪凝聚的，
而 $P[1/y_i]$ 作为
$R[x_1, \ldots, x_n, y_1, \ldots, y_r, z_1, \ldots, z_r, 1/y_i]$-模
是伪凝聚的。因此，由引理 \ref{more-algebra-lemma-finite-push-pseudo-coherent}，
对每个 $i$，$K^\bullet \otimes_A A_{f_i}$ 都是
$P[1/y_i]$-模的 $m$-伪凝聚复形。于是由引理
\ref{more-algebra-lemma-glue-pseudo-coherent}，可知 $K^\bullet$ 作为
% P02-E0316
$P$-模复形是伪凝聚的；再由引理
\ref{more-algebra-lemma-finite-push-pseudo-coherent}，可知 $K^\bullet$ 作为
$R[x_1, \ldots, x_n, y_1, \ldots, y_r, z_1, \ldots, z_r]$-模复形
是伪凝聚的。
\end{proof}

\begin{lemma}
\label{more-algebra-lemma-Noetherian-relative-pseudo-coherent}
设 $R$ 为 Noether 环，并设 $R \to A$ 为有限型环同态。则
\begin{enumerate}
\item $A$-模复形 $K^\bullet$ 相对于 $R$ 是 $m$-伪凝聚的，
当且仅当 $K^\bullet \in D^{-}(A)$，并且当 $i \geq m$ 时
$H^i(K^\bullet)$ 是有限 $A$-模。
\item $A$-模复形 $K^\bullet$ 相对于 $R$ 是伪凝聚的，
当且仅当 $K^\bullet \in D^{-}(A)$，并且对所有 $i$，
$H^i(K^\bullet)$ 都是有限 $A$-模。
\item 一个 $A$-模相对于 $R$ 是伪凝聚的，当且仅当它是有限的。
\end{enumerate}
\end{lemma}

\begin{proof}
% P02-E0317
这是引理 \ref{more-algebra-lemma-Noetherian-pseudo-coherent} 和定义的直接推论。
\end{proof}










\section{伪凝聚与完美环同态}
\label{more-algebra-section-pseudo-coherent-perfect-ring-map}

\noindent
可以如下定义这些类型的环同态。

\begin{definition}
\label{more-algebra-definition-pseudo-coherent-perfect}
设 $A \to B$ 为环同态。
\begin{enumerate}
\item 若 $A \to B$ 是有限型的，且 $B$ 作为 $B$-模相对于 $A$
是伪凝聚的，则称该同态为{\it 伪凝聚环同态}。
\item 若 $A \to B$ 是伪凝聚环同态，并且 $B$ 作为 $A$-模具有有限
Tor 维数，则称该同态为{\it 完美环同态}。
\end{enumerate}
\end{definition}

\noindent
该术语可能并不标准。由引理 \ref{more-algebra-lemma-rel-n-pseudo-module}，
可知 $A \to B$ 是伪凝聚的，当且仅当
$B = A[x_1, \ldots, x_n]/I$，且 $B$ 作为
$A[x_1, \ldots, x_n]$-模具有由有限自由
$A[x_1, \ldots, x_n]$-模组成的解消。完美环同态定义的动机来自引理
\ref{more-algebra-lemma-perfect}。
% P02-E0318
下述引理给出完美环同态的一个更有用且更直观的刻画。

\begin{lemma}
\label{more-algebra-lemma-perfect-ring-map}
环同态 $A \to B$ 是完美的，当且仅当
$B = A[x_1, \ldots, x_n]/I$，且 $B$ 作为
$A[x_1, \ldots, x_n]$-模具有由有限投射
$A[x_1, \ldots, x_n]$-模组成的有限解消。
\end{lemma}

\begin{proof}
若 $A \to B$ 是完美的，则 $B = A[x_1, \ldots, x_n]/I$，
$B$ 作为 $A[x_1, \ldots, x_n]$-模是伪凝聚的，并且作为
$A$-模具有有限 Tor 维数。因此，引理
\ref{more-algebra-lemma-perfect-over-polynomial-ring} 表明 $B$ 作为
% P02-E0319
$A[x_1, \ldots, x_n]$-模是完美的，即它具有由有限投射
$A[x_1, \ldots, x_n]$-模组成的有限解消
（引理 \ref{more-algebra-lemma-perfect-module}）。
反之，若 $B = A[x_1, \ldots, x_n]/I$，且 $B$ 作为
$A[x_1, \ldots, x_n]$-模具有由有限投射
$A[x_1, \ldots, x_n]$-模组成的有限解消，则 $B$ 作为
$A[x_1, \ldots, x_n]$-模是伪凝聚的，故 $A \to B$ 是伪凝聚的。
此外，给定的 $A[x_1, \ldots, x_n]$ 上解消是由平坦 $A$-模组成的
有限解消，故 $B$ 作为 $A$-模具有有限 Tor 维数。
\end{proof}

\noindent
前几节的许多结果都可以用这一术语重新表述。
% P02-E0320
关于概形态射的相应讨论，还可参见 More on Morphisms，第
\ref{more-morphisms-section-pseudo-coherent} 节和第
\ref{more-morphisms-section-perfect} 节。

\begin{lemma}
\label{more-algebra-lemma-Noetherian-pseudo-coherent-ring-map}
Noether 环之间的有限型环同态是伪凝聚的。
\end{lemma}

\begin{proof}
参见引理 \ref{more-algebra-lemma-Noetherian-relative-pseudo-coherent}。
\end{proof}

\begin{lemma}
\label{more-algebra-lemma-flat-finite-presentation-perfect}
平坦且有限表示的环同态是完美的。
\end{lemma}

\begin{proof}
设 $A \to B$ 为平坦且有限表示的环同态。显然，$B$ 具有有限
Tor 维数。由 Algebra，引理
\ref{algebra-lemma-flat-finite-presentation-limit-flat}，
存在有限型 $\mathbf{Z}$-代数 $A_0 \subset A$ 和平坦有限型环同态
$A_0 \to B_0$，使得 $B = B_0 \otimes_{A_0} A$。由引理
\ref{more-algebra-lemma-Noetherian-relative-pseudo-coherent}，可知
$A_0 \to B_0$ 是伪凝聚的。由于 $A_0 \to B_0$ 平坦，
$B_0$ 与 $A$ 在 $A_0$ 上 Tor 独立，故可以应用引理
\ref{more-algebra-lemma-base-change-relative-pseudo-coherent}，得到
$A \to B$ 是伪凝聚的。
\end{proof}

\begin{lemma}
\label{more-algebra-lemma-regular-perfect-ring-map}
设 $A \to B$ 为有限型环同态，其中 $A$ 是有限维正则环。
则 $A \to B$ 是完美的。
\end{lemma}

\begin{proof}
由 Algebra，引理
\ref{algebra-lemma-finite-gl-dim-finite-dim-regular}，关于 $A$ 的假设
意指 $A$ 具有有限整体维数。因此，每个模都具有有限 Tor 维数；
参见引理 \ref{more-algebra-lemma-finite-gl-dim-tor-dimension}，特别地 $B$ 也如此。
由引理 \ref{more-algebra-lemma-Noetherian-pseudo-coherent-ring-map}，该同态是伪凝聚的。
\end{proof}

\begin{lemma}
\label{more-algebra-lemma-lci-perfect}
局部完全交同态是完美的。
\end{lemma}

\begin{proof}
设 $A \to B$ 为局部完全交同态。由定义
\ref{more-algebra-definition-local-complete-intersection}，这意指
$B = A[x_1, \ldots, x_n]/I$，其中 $I$ 是
$A[x_1, \ldots, x_n]$ 中的 Koszul 理想。
由引理 \ref{more-algebra-lemma-perfect-ring-map} 和 \ref{more-algebra-lemma-perfect-module}，
只需证明 $I$ 是完美 $A[x_1, \ldots, x_n]$-模。
由引理 \ref{more-algebra-lemma-glue-perfect}，这是局部问题。因此由定义
\ref{more-algebra-definition-regular-ideal}，可以假设 $I$ 由一个 Koszul-正则序列生成。
这当然意指 $I$ 具有有限自由解消，结论得证。
\end{proof}

\begin{lemma}
\label{more-algebra-lemma-relative-pseudo-coherent-is-moot}
设 $R \to A$ 为伪凝聚环同态，并设 $K \in D(A)$。下列条件等价：
\begin{enumerate}
\item $K$ 相对于 $R$ 是 $m$-伪凝聚的（相应地，伪凝聚的）；
\item $K$ 在 $D(A)$ 中是 $m$-伪凝聚的（相应地，伪凝聚的）。
\end{enumerate}
\end{lemma}

\begin{proof}
这是引理 \ref{more-algebra-lemma-composition-relative-pseudo-coherent} 一个特例的重述。
\end{proof}

\begin{lemma}
\label{more-algebra-lemma-more-relative-pseudo-coherent-is-moot}
设 $R \to B \to A$ 为环同态，其中
$\varphi : B \to A$ 满射，并且 $R \to B$ 与 $R \to A$
平坦且有限表示。
% P02-E0321
对 $K \in D(A)$，用 $\varphi_*K \in D(B)$ 表示限制。下列条件等价：
\begin{enumerate}
\item $K$ 是伪凝聚的；
\item $K$ 相对于 $R$ 是伪凝聚的；
\item $K$ 相对于 $A$ 是伪凝聚的；
\item $\varphi_*K$ 是伪凝聚的；
\item $\varphi_*K$ 相对于 $R$ 是伪凝聚的。
\end{enumerate}
% P02-E0322
对于 $m$-伪凝聚性也有类似结论。
\end{lemma}

\begin{proof}
注意，$R \to A$ 和 $R \to B$ 是完美环同态
（引理 \ref{more-algebra-lemma-flat-finite-presentation-perfect}），
因而更是伪凝聚环同态。因此，由引理
\ref{more-algebra-lemma-relative-pseudo-coherent-is-moot}，有
(1) $\Leftrightarrow$ (2) 和 (4) $\Leftrightarrow$ (5)。

\medskip\noindent
利用 $A$ 相对于 $R$ 是伪凝聚的这一事实，应用引理
\ref{more-algebra-lemma-composition-relative-pseudo-coherent} 可知
(2) $\Leftrightarrow$ (3)。
% P02-E0323 P02-E0324
另一方面，由于 $A \to B$ 是满射，从定义
\ref{more-algebra-definition-relatively-pseudo-coherent} 直接可见
(3) 与 (4) 等价。
\end{proof}




\section{相对完美模}
\label{more-algebra-section-relatively-perfect}

\noindent
本节给出第 \ref{more-algebra-section-relative-pseudo-coherent} 节关于导出范畴中
完美对象的类似理论。
% P02-E0325
由于尚不确定一般情形下的正确概念是什么，我们只在有限的一般性下定义这一概念。
相关讨论见 Derived Categories of Schemes，评注
\ref{perfect-remark-discuss-rel-perfect}。

\begin{definition}
\label{more-algebra-definition-relatively-perfect}
设 $R \to A$ 为平坦有限表示环同态。若 $D(A)$ 的对象 $K$
是伪凝聚的（定义 \ref{more-algebra-definition-pseudo-coherent}），并且在
$R$ 上具有有限 Tor 维数（定义 \ref{more-algebra-definition-tor-amplitude}），
则称 $K$ 是{\it $R$-完美的}，或称它{\it 相对于 $R$ 完美}。
\end{definition}

\noindent
由引理 \ref{more-algebra-lemma-more-relative-pseudo-coherent-is-moot}，
要求 $K$ 相对于 $R$ 伪凝聚会得到同一个概念。
下面给出一些必备引理。

\begin{lemma}
\label{more-algebra-lemma-cone-relatively-perfect}
设 $R \to A$ 为平坦有限表示环同态。$D(A)$ 中的 $R$-完美对象组成
饱和的\footnote{Derived Categories，定义
\ref{derived-definition-saturated}。}三角严格满子范畴。
\end{lemma}

\begin{proof}
这由引理 \ref{more-algebra-lemma-cone-pseudo-coherent}、
\ref{more-algebra-lemma-summands-pseudo-coherent}、
\ref{more-algebra-lemma-cone-tor-amplitude} 和
\ref{more-algebra-lemma-summands-tor-amplitude} 得到。
\end{proof}

\begin{lemma}
\label{more-algebra-lemma-perfect-relatively-perfect}
设 $R \to A$ 为平坦有限表示环同态。$D(A)$ 中的完美对象是
$R$-完美的。若 $K, M \in D(A)$，$K$ 是完美的且 $M$ 是
$R$-完美的，则 $K \otimes_A^\mathbf{L} M$ 是 $R$-完美的。
\end{lemma}

\begin{proof}
在第二个断言中取 $M = A$，即得第一个断言。
第二个断言由引理 \ref{more-algebra-lemma-perfect}、
\ref{more-algebra-lemma-push-tor-amplitude} 和 \ref{more-algebra-lemma-tensor-pseudo-coherent} 得到。
\end{proof}

\begin{lemma}
\label{more-algebra-lemma-structure-relatively-perfect}
设 $R \to A$ 为平坦有限表示环同态，并设 $K \in D(A)$。
下列条件等价：
\begin{enumerate}
\item $K$ 是 $R$-完美的；
\item $K$ 同构于一个由 $R$-平坦有限表示 $A$-模组成的有限复形。
\end{enumerate}
\end{lemma}

\begin{proof}
为证明 (2) 蕴含 (1)，由引理 \ref{more-algebra-lemma-cone-relatively-perfect}，
只需证明一个 $R$-平坦有限表示 $A$-模 $M$ 定义 $D(A)$ 中的
$R$-完美对象。由于 $M$ 在 $R$ 上具有有限 Tor 维数，
只需证明 $M$ 是伪凝聚的。由 Algebra，引理
\ref{algebra-lemma-flat-finite-presentation-limit-flat}，
存在有限型 $\mathbf{Z}$-代数 $R_0 \subset R$、平坦有限型环同态
$R_0 \to A_0$，以及在 $R_0$ 上平坦的有限 $A_0$-模 $M_0$，
使得 $A = A_0 \otimes_{R_0} R$ 且
$M = M_0 \otimes_{R_0} R$。由引理
\ref{more-algebra-lemma-Noetherian-pseudo-coherent}，可知
% P02-E0326
$M_0$ 是伪凝聚 $A_0$-模。选取由有限自由 $A_0$-模
$P_0^n$ 组成的解消 $P_0^\bullet \to M_0$。由于 $A_0$ 在
$R_0$ 上平坦，这是一个平坦解消。由于 $M_0$ 在 $R_0$ 上平坦，
可知 $P^\bullet = P_0^\bullet \otimes_{R_0} R$ 仍然解消
$M = M_0 \otimes_{R_0} R$。（可以应用引理
\ref{more-algebra-lemma-base-change-comparison} 看出这一点。）
因此 $P^\bullet$ 是 $M$ 在 $A$ 上的有限自由解消，从而
$M$ 是伪凝聚的。

\medskip\noindent
假设 (1)。可以用有限自由 $A$-模组成的有界上复形
$P^\bullet$ 表示 $K$。假设把 $K$ 视为 $D(R)$ 中的对象时，
其 Tor 振幅包含在 $[a, b]$ 中。由引理
\ref{more-algebra-lemma-last-one-flat}，可知 $\tau_{\geq a}P^\bullet$
是表示 $K$ 的复形，其各项为 $R$-平坦有限表示 $A$-模。
\end{proof}

\begin{lemma}
\label{more-algebra-lemma-base-change-relatively-perfect}
设 $R \to A$ 为平坦有限表示环同态。设 $R \to R'$ 为环同态，
并令 $A' = A \otimes_R R'$。若 $K \in D(A)$ 是 $R$-完美的，
则 $K \otimes_A^\mathbf{L} A'$ 是 $R'$-完美的。
\end{lemma}

\begin{proof}
由引理 \ref{more-algebra-lemma-pull-pseudo-coherent}，可知
$K \otimes_A^\mathbf{L} A'$ 是伪凝聚的。
由引理 \ref{more-algebra-lemma-base-change-comparison}，可知
$K \otimes_A^\mathbf{L} A'$ 在 $D(R')$ 中等于
$K \otimes_R^\mathbf{L} R'$。然后可以应用引理
\ref{more-algebra-lemma-pull-tor-amplitude}，得到 $D(R')$ 中的
$K \otimes_R^\mathbf{L} R'$ 具有有限 Tor 维数。
\end{proof}

\begin{lemma}
\label{more-algebra-lemma-compute-RHom-relatively-perfect}
设 $R \to A$ 为平坦环同态，设 $K, L \in D(A)$，其中 $K$
是伪凝聚的，而 $L$ 在 $R$ 上具有有限 Tor 维数。可以选取：
\begin{enumerate}
\item 表示 $K$ 的有限自由 $A$-模有界上复形 $P^\bullet$；
\item 表示 $L$ 的 $R$-平坦 $A$-模有界复形 $F^\bullet$。
\end{enumerate}
在作出这些选择后，有：
\begin{enumerate}
\item[(a)] $E^\bullet = \Hom^\bullet(P^\bullet, F^\bullet)$
是表示 $R\Hom_A(K, L)$ 的 $R$-平坦 $A$-模有界下复形；
\item[(b)] 对任意环同态 $R \to R'$，令 $A' = A \otimes_R R'$，
则复形 $E^\bullet \otimes_R R'$ 表示
$R\Hom_{A'}(K \otimes_A^\mathbf{L} A',
L \otimes_A^\mathbf{L} A')$。
\end{enumerate}
若还假设 $R \to A$ 是有限表示的，且 $L$ 是 $R$-完美的，
则可以选取各 $F^p$ 为有限表示 $A$-模，因而各 $E^n$ 也是有限表示
$A$-模。
\end{lemma}

\begin{proof}
$P^\bullet$ 的存在性就是伪凝聚复形的定义。首先用自由 $A$-模组成的
有界上复形 $F^\bullet$ 表示 $L$（这是可能的，因为有界 Tor 维数
尤其蕴含有界）。再设把 $L$ 视为 $D(R)$ 中的对象时，其 Tor 振幅
包含在 $[a, b]$ 中。用 $\tau_{\geq a}F^\bullet$ 替换
$F^\bullet$ 后，得到 (2) 中的复形。这由引理
\ref{more-algebra-lemma-last-one-flat} 得到。

\medskip\noindent
(a) 的证明。
% P02-E0327
由于 $F^\bullet$ 有界且 $P^\bullet$ 有界上，可知当
$n \ll 0$ 时 $E^n = 0$，并且 $E^n$ 是有限（！）直和
$$
E^n = \bigoplus\nolimits_{p + q = n} \Hom_A(P^{-q}, F^p)
$$
由于 $P^{-q}$ 是有限自由的，这的确是 $R$-平坦 $A$-模。
$E^\bullet$ 表示 $R\Hom_A(K, L)$ 这一事实由引理
\ref{more-algebra-lemma-RHom-out-of-projective} 得到。

\medskip\noindent
(b) 的证明。
设 $R \to R'$ 为环同态，并令 $A' = A \otimes_R R'$。
由引理 \ref{more-algebra-lemma-base-change-comparison}，对象
$L \otimes_A^\mathbf{L} A'$ 由视为 $A'$-模复形的
$F^\bullet \otimes_R R'$ 表示（因为各 $F^p$ 在 $R$ 上平坦）。
对 $P^\bullet \otimes_R R'$ 亦然。与上面一样，
$R\Hom_{A'}(K \otimes_A^\mathbf{L} A',
L \otimes_A^\mathbf{L} A')$ 由
$$
\Hom^\bullet(P^\bullet \otimes_R R', F^\bullet \otimes_R R') =
E^\bullet \otimes_R R'
$$
表示。逐项考察该复形，并应用
$\Hom_{A'}(P^{-q} \otimes_R R', F^p \otimes_R R') =
\Hom_A(P^{-q}, F^p) \otimes_R R'$，即得该等式。
\end{proof}

\begin{lemma}
\label{more-algebra-lemma-colimit-relatively-perfect}
设 $R = \colim_{i \in I} R_i$ 为环的滤过余极限。
取 $0 \in I$，并设 $R_0 \to A_0$ 为平坦有限表示环同态。
对 $i \geq 0$，令 $A_i = R_i \otimes_{R_0} A_0$，并令
$A = R \otimes_{R_0} A_0$。
\begin{enumerate}
\item 给定 $D(A)$ 中的 $R$-完美对象 $K$，存在 $i \in I$
以及 $D(A_i)$ 中的 $R_i$-完美对象 $K_i$，使得在 $D(A)$ 中
$K \cong K_i \otimes_{A_i}^\mathbf{L} A$。
\item 给定 $K_0, L_0 \in D(A_0)$，其中 $K_0$ 伪凝聚，
且 $L_0$ 在 $R_0$ 上具有有限 Tor 维数，则有
$$
\Hom_{D(A)}(K_0 \otimes_{A_0}^\mathbf{L} A, L_0 \otimes_{A_0}^\mathbf{L} A) =
\colim_{i \geq 0}
\Hom_{D(A_i)}(K_0 \otimes_{A_0}^\mathbf{L} A_i,
L_0 \otimes_{A_0}^\mathbf{L} A_i)
$$
\end{enumerate}
特别地，$A$ 上 $R$-完美复形的三角范畴是各
$A_i$ 上 $R_i$-完美复形三角范畴的余极限。
\end{lemma}

\begin{proof}
由 Algebra，引理 \ref{algebra-lemma-colimit-category-fp-modules}，
有限表示 $A$-模范畴是各有限表示 $A_i$-模范畴的余极限。
在此基础上，由 Algebra，引理
\ref{algebra-lemma-flat-finite-presentation-limit-flat} 可知，
% P02-E0328
$R$-平坦有限表示 $A$-模范畴是各
$R_i$-平坦有限表示 $A_i$-模范畴的余极限。
因此，引理 \ref{more-algebra-lemma-structure-relatively-perfect} 中的刻画
证明了 (1)。

\medskip\noindent
为证明 (2)，按引理 \ref{more-algebra-lemma-compute-RHom-relatively-perfect}
选取表示 $K_0$ 的 $P_0^\bullet$ 和表示 $L_0$ 的
$F_0^\bullet$。于是
$E_0^\bullet = \Hom^\bullet(P_0^\bullet, F_0^\bullet)$ 满足
$$
H^0(E_0^\bullet \otimes_{R_0} R_i) =
\Hom_{D(A_i)}(K_0 \otimes_{A_0}^\mathbf{L} A_i,
L_0 \otimes_{A_0}^\mathbf{L} A_i)
$$
以及
$$
H^0(E_0^\bullet \otimes_{R_0} R) =
\Hom_{D(A)}(K_0 \otimes_{A_0}^\mathbf{L} A, L_0 \otimes_{A_0}^\mathbf{L} A)
$$
这是该引理所给出的。% P02-E0329
因此结论成立，因为张量积与余极限交换，并且滤过余极限是正合的
（Algebra，引理 \ref{algebra-lemma-directed-colimit-exact}）。
\end{proof}

\begin{lemma}
\label{more-algebra-lemma-thickening-relatively-perfect}
设 $R' \to A'$ 为平坦有限表示环同态。设 $R' \to R$
为满射环同态，其核为幂零理想。令 $A = A' \otimes_{R'} R$。
设 $K' \in D(A')$，并令
$K = K' \otimes_{A'}^\mathbf{L} A$，视为 $D(A)$ 中的对象。
若 $K$ 是 $R$-完美的，则 $K'$ 是 $R'$-完美的。
\end{lemma}

\begin{proof}
注意，$A' \to A$ 的核是幂零的，并且由 $R' \to A'$ 的平坦性，
有 $K = K' \otimes_{R'}^\mathbf{L} R$（见第
\ref{more-algebra-section-tor-independence} 节）。因此，结合引理
\ref{more-algebra-lemma-check-pseudo-coherent-modulo-nilpotent} 与
\ref{more-algebra-lemma-check-tor-dimension-modulo-nilpotent} 即得结论。
\end{proof}

\begin{lemma}
\label{more-algebra-lemma-lift-from-fibre-relatively-perfect}
设 $R$ 为环。设 $A = R[x_1, \ldots, x_d]/I$ 在 $R$ 上平坦且有限表示。
设素理想 $\mathfrak q \subset A$ 位于素理想
$\mathfrak p \subset R$ 之上。设 $K \in D(A)$ 伪凝聚，
并设 $a, b \in \mathbf{Z}$。若
$H^i(K_\mathfrak q \otimes_{R_\mathfrak p}^\mathbf{L} \kappa(\mathfrak p))$
仅在 $i \in [a, b]$ 时非零，则 $K_\mathfrak q$ 在 $R$ 上
具有 $[a - d, b]$ 内的 Tor 振幅。
\end{lemma}

\begin{proof}
由引理 \ref{more-algebra-lemma-more-relative-pseudo-coherent-is-moot}，
$K$ 作为 $R[x_1, \ldots, x_d]$-模复形是伪凝聚的。
因此可以假设 $A = R[x_1, \ldots, x_d]$。
对 $R_\mathfrak p \to A_\mathfrak q$ 和复形 $K_\mathfrak q$
应用引理 \ref{more-algebra-lemma-perfect-over-regular-local-ring}，并利用假设，
可知 $K_\mathfrak q$ 在 $D(A_\mathfrak q)$ 中是完美的，
且具有 $[a - d, b]$ 内的 Tor 振幅。由于
$R_\mathfrak p \to A_\mathfrak q$ 平坦，由引理
\ref{more-algebra-lemma-flat-push-tor-amplitude} 即得结论。
\end{proof}

\begin{lemma}
\label{more-algebra-lemma-bounded-on-fibres-relatively-perfect}
设 $R \to A$ 为平坦有限表示环同态，并设 $K \in D(A)$ 伪凝聚。
下列条件等价：
\begin{enumerate}
\item $K$ 是 $R$-完美的；
\item $K$ 下有界，并且对每个素理想 $\mathfrak p \subset R$，
对象 $K \otimes_R^\mathbf{L} \kappa(\mathfrak p)$ 下有界。
\end{enumerate}
\end{lemma}

\begin{proof}
注意，(1) 蕴含 (2)，因为按照定义，$R$-完美复形作为
$R$-模复形具有有界 Tor 维数。下面证明反向蕴含。

\medskip\noindent
写成 $A = R[x_1, \ldots, x_d]/I$。% P02-E0330
在 $D(R[x_1, \ldots, x_d])$ 中，用 $L$ 表示 $K$ 的限制。
由引理 \ref{more-algebra-lemma-more-relative-pseudo-coherent-is-moot}，
可知 $L$ 伪凝聚。由于 $L$ 与 $K$ 在 $D(R)$ 中具有相同的像，
故 $L$ 是 $R$-完美的，当且仅当 $K$ 是 $R$-完美的。
此外，$L \otimes_R^\mathbf{L} \kappa(\mathfrak p)$ 与
$K \otimes_R^\mathbf{L} \kappa(\mathfrak p)$ 是
$D(\kappa(\mathfrak p))$ 中相同的对象。这把问题归约到
$A = R[x_1, \ldots, x_d]$ 的情形。

\medskip\noindent
设 $A = R[x_1, \ldots, x_d]$，并设 $K$ 满足 (2)。
令素理想 $\mathfrak q \subset A$ 位于素理想
$\mathfrak p \subset R$ 之上。对
$R_\mathfrak p \to A_\mathfrak q$ 和复形 $K_\mathfrak q$
应用引理 \ref{more-algebra-lemma-perfect-over-regular-local-ring}，并利用假设，
可知 $K_\mathfrak q$ 在 $D(A_\mathfrak q)$ 中是完美的。
由于 $K$ 下有界，由引理 \ref{more-algebra-lemma-check-perfect-stalks}，
可知 $K$ 在 $D(A)$ 中是完美的。再由引理
\ref{more-algebra-lemma-perfect-relatively-perfect}，这蕴含 $K$ 是 $R$-完美的，
证明完成。
\end{proof}

\section{二项复形}
\label{more-algebra-section-two-term}

\noindent
本节证明若干关于模的二项复形的结果，这些结果将帮助我们理解
朴素余切复形所满足的条件。

\begin{lemma}
\label{more-algebra-lemma-ext-1-zero}
设 $R$ 为环。设 $K \in D(R)$，且对
$i \not \in \{-1, 0\}$ 有 $H^i(K) = 0$。下列条件等价：
\begin{enumerate}
\item $H^{-1}(K) = 0$，且 $H^0(K)$ 是投射模；
\item 对每个 $R$-模 $M$，有 $\Ext^1_R(K, M) = 0$。
\end{enumerate}
若 $R$ 是 Noether 环，并且对 $i = -1, 0$，$H^i(K)$ 是有限
$R$-模，则它们还等价于：
\begin{enumerate}
\item[(3)] 对每个有限 $R$-模 $M$，有 $\Ext^1_R(K, M) = 0$。
\end{enumerate}
\end{lemma}

\begin{proof}
(1) 与 (2) 的等价性由引理 \ref{more-algebra-lemma-projective-amplitude} 得到。
若 $R$ 是 Noether 环，并且对 $i = -1, 0$，$H^i(K)$ 是有限
$R$-模，则 $K$ 是伪凝聚的；见引理
\ref{more-algebra-lemma-Noetherian-pseudo-coherent}。因此，(1) 与 (3)
的等价性由引理 \ref{more-algebra-lemma-projective-amplitude-pseudo-coherent} 得到。
\end{proof}

\begin{remark}
\label{more-algebra-remark-smoothness-ext-1-zero}
下列两个陈述由引理 \ref{more-algebra-lemma-ext-1-zero}、Algebra，定义
\ref{algebra-definition-smooth} 和 Algebra，命题
\ref{algebra-proposition-characterize-formally-smooth} 得到。
\begin{enumerate}
\item 环同态 $A \to B$ 光滑，当且仅当 $A \to B$ 有限表示，且对每个
$B$-模 $N$，有 $\Ext^1_B(\NL_{B/A}, N) = 0$。
\item 环同态 $A \to B$ 形式光滑，当且仅当对每个
$B$-模 $N$，有 $\Ext^1_B(\NL_{B/A}, N) = 0$。
\end{enumerate}
\end{remark}

\begin{lemma}
\label{more-algebra-lemma-represent-two-term-complex}
设 $R$ 为环。设 $K$ 为 $D(R)$ 的对象，且对
$i \not \in \{-1, 0\}$ 有 $H^i(K) = 0$。则：
\begin{enumerate}
\item $K$ 可以由二项复形 $K^{-1} \to K^0$ 表示，其中
$K^0$ 是自由模；
\item 若 $R$ 是 Noether 环，并且对 $i = -1, 0$，$H^i(K)$
是有限 $R$-模，则 $K$ 可以由二项复形 $K^{-1} \to K^0$ 表示，
其中 $K^0$ 是有限自由模，且 $K^{-1}$ 有限。
\end{enumerate}
\end{lemma}

\begin{proof}
证明 (1)。假设 $K$ 由模复形 $M^\bullet$ 给出。首先可以用
$\tau_{\leq 0}M^\bullet$ 替换 $M^\bullet$。因此可以假设对
$i > 0$ 有 $M^i = 0$。% P02-E0331
接着，可以选取自由解消 $P^\bullet \to M^\bullet$，使得对
$i > 0$ 有 $P^i = 0$；见 Derived Categories，引理
\ref{derived-lemma-subcategory-left-resolution}。最后，可以令
$K^\bullet = \tau_{\geq -1}P^\bullet$。

\medskip\noindent
证明 (2)。假设 $R$ 是 Noether 环，并且对 $i = -1, 0$，
$H^i(K)$ 是有限 $R$-模。由引理 \ref{more-algebra-lemma-pseudo-coherent}，
可以选取拟同构 $F^\bullet \to M^\bullet$，使得对 $i > 0$
有 $F^i = 0$，且各 $F^i$ 有限自由。然后可以令
$K^\bullet = \tau_{\geq -1}F^\bullet$。
\end{proof}

\noindent
由下面的引理容易理解导出范畴中以朴素余切复形
$\NL_{B/A}$ 或 $\NL(\alpha)$ 为源的映射（见
Algebra，第 \ref{algebra-section-netherlander} 节）。

\begin{lemma}
\label{more-algebra-lemma-map-out-of-almost-free}
设 $R$ 为环。设 $M^\bullet$ 为 $R$ 上的模复形，对
$i > 0$ 有 $M^i = 0$，且 $M^0$ 是投射 $R$-模。
设 $K^\bullet$ 为另一个复形。
\begin{enumerate}
\item 假设对 $i \leq -2$ 有 $K^i = 0$。则
$\Hom_{D(R)}(M^\bullet, K^\bullet) = \Hom_{K(R)}(M^\bullet, K^\bullet)$。
\item 假设对 $i \not \in [-1, 0]$ 有 $K^i = 0$，且
$K^0$ 是投射 $R$-模。对复形映射
$a^\bullet : M^\bullet \to K^\bullet$，下列条件等价：
\begin{enumerate}
\item 对所有 $R$-模 $N$，$a^\bullet$ 诱导零映射
$\Ext^1_R(K^\bullet, N) \to \Ext^1_R(M^\bullet, N)$；
\item 存在映射 $h^0 : M^0 \to K^{-1}$，使得
$a^{-1} + h^0 \circ d^{-1}_K = 0$。
\end{enumerate}
\item 假设对 $i \leq -3$ 有 $K^i = 0$。设
$\alpha \in \Hom_{D(R)}(M^\bullet, K^\bullet)$。若 $\alpha$
与 $K^\bullet \to K^{-2}[2]$ 的复合来自 $R$-模映射
$a : M^{-2} \to K^{-2}$，且 $a \circ d_M^{-3} = 0$，则
$\alpha$ 可以由复形映射 $a^\bullet : M^\bullet \to K^\bullet$
表示，并且 $a^{-2} = a$。
\item % P02-E0332
在 (2) 中，若另一个复形映射
$(a')^\bullet : M^\bullet \to K^\bullet$ 表示 $\alpha$，且
$a = (a')^{-2}$，则对 $i = 0, -1$，存在
$h^i : M^i \to K^{i - 1}$，使得
$$
h^{-1} \circ d_M^{-2} = 0, \quad
(a')^{-1} = a^{-1} + d_K^{-2} \circ h^{-1} + h^0 \circ d_M^{-1},\quad
(a')^0 = a^0 + d_K^{-1} \circ h^0
$$
\end{enumerate}
\end{lemma}

\begin{proof}
令 $F^0 = M^0$。选取自由 $R$-模 $F^{-1}$ 以及满射
$F^{-1} \to M^{-1}$。选取自由 $R$-模 $F^{-2}$ 以及满射
$F^{-2} \to M^{-2} \times_{M^{-1}} F^{-1}$。如此继续，得到拟同构
$p^\bullet : F^\bullet \to M^\bullet$，它逐项满射，并且对所有
$i$，$F^i$ 都是投射的。

\medskip\noindent
证明 (1)。由 Derived Categories，引理
\ref{derived-lemma-morphisms-from-projective-complex}，有
$$
\Hom_{D(R)}(M^\bullet, K^\bullet) = \Hom_{K(R)}(F^\bullet, K^\bullet)
$$
若对 $i \leq -2$ 有 $K^i = 0$，则任意复形态射
$F^\bullet \to K^\bullet$ 均经由 $p^\bullet$ 分解。类似地，
任意同伦 $\{h^i : F^i \to K^{i - 1}\}$ 均经由
$p^\bullet$ 分解。因此 (1) 成立。

\medskip\noindent
证明 (2)。若 (2)(b) 成立，则 $a^\bullet$ 同伦于一个在次数
$-1$ 上为零的复形映射
$(a')^\bullet : M^\bullet \to K^\bullet$。
另一方面，设 $N \to I^\bullet$ 为内射解消。由 Derived Categories，
引理 \ref{derived-lemma-morphisms-into-injective-complex}，有
$$
\Ext^1_R(K^\bullet, N) = \Hom_{D(R)}(K^\bullet, I^\bullet[1]) =
\Hom_{K(R)}(K^\bullet, I^\bullet[1])
$$
设 $b^\bullet : K^\bullet \to I^\bullet[1]$ 为复形映射。
由于 $K^1 = 0$，映射 $b^0 : K^0 \to I^1$ 的像落在
$I^1 \to I^2$ 的核中，而该核是 $I^0 \to I^1$ 的像。
由于 $K^0$ 是投射的，可以把 $b^0$ 提升为映射
$h : K^0 \to I^0$。因此，$b^\bullet$ 同伦于复形映射
$(b')^\bullet$，且 $(b')^0 = 0$。由于对
$i \not \in [-1, 0]$ 有 $K^i = 0$，可知作为复形映射，
$(b')^\bullet \circ (a')^\bullet = 0$。故映射
$\Ext^1_R(K^\bullet, N) \to \Ext^1_R(M^\bullet, N)$ 为零。
由此可见，(2)(b) 蕴含 (2)(a)。反之，假设 (2)(a)。
可知 $\Ext^1_R(K^\bullet, K^{-1})$ 中的典范元素映到
$\Ext^1_R(M^\bullet, K^{-1})$ 中的零。利用 (1)，立即得到
(2)(b) 中的映射 $h^0$。

\medskip\noindent
证明 (3)。选取表示 $\alpha$ 的
$b^\bullet : F^\bullet \to K^\bullet$。$\alpha$ 与
$K^\bullet \to K^{-2}[2]$ 的复合由
$b^{-2} : F^{-2} \to K^{-2}$ 表示。由于它同伦于
$a \circ p^{-2} : F^{-2} \to M^{-2} \to K^{-2}$，存在映射
$h : F^{-1} \to K^{-2}$，使得
$b^{-2} = a \circ p^{-2} + h \circ d_F^{-2}$。
把 $h$ 看作从 $F^\bullet$ 到 $K^\bullet$ 的同伦，并据此调整
$b^\bullet$，可得 $b^{-2} = a \circ p^{-2}$。因此
$b^{-2}$ 经由 $p^{-2}$ 分解。由于 $F^0 = M^0$，
$p^{-2}$ 的核满射到 $p^{-1}$ 的核（例如，因为
$p^\bullet$ 的核是非周期复形，或通过追图可见）。因此
$b^{-1}$ 也必然经由 $p^{-1}$ 分解；由这些分解以及
$a^0 = b^0$，可知 (3) 成立。

\medskip\noindent
省略 (4) 的证明。提示：$a^\bullet \circ p^\bullet$ 与
$(a')^\bullet \circ p^\bullet$ 之间存在同伦，并且与前面一样论证，
该同伦经由 $p^\bullet$ 分解。
\end{proof}

\noindent
设 $A \to B$ 为有限表示环同态。给定理想 $I \subset B$，
可以考虑如下条件：
\begin{enumerate}
\item[(*)] 对所有 $B$-模 $N$，$\Ext^1_B(\NL_{B/A}, N)$ 被 $I$ 零化。
\end{enumerate}
这一条件是概念“$A \to B$ 的奇异轨迹在概形论意义下包含于
$V(I)$ 中”的一种可能的精确数学表述。请将其与评注
\ref{more-algebra-remark-smoothness-ext-1-zero} 及下面的引理作比较。

\begin{lemma}
\label{more-algebra-lemma-ext-1-annihilated-definite}
设 $R$ 为环，$I \subset R$ 为理想。设 $K \in D(R)$。
假设对 $i \not \in \{-1, 0\}$ 有 $H^i(K) = 0$。下列条件等价：
\begin{enumerate}
\item 对所有 $R$-模 $N$，$\Ext^1_R(K, N)$ 被 $I$ 零化；
\item $K$ 可以由复形 $K^{-1} \to K^0$ 表示，其中 $K^0$ 自由，
并且对任意 $a \in I$，映射 $a : K^{-1} \to K^{-1}$ 经由
$d_K^{-1} : K^{-1} \to K^0$ 分解；
\item 每当 $K$ 由二项复形 $K^{-1} \to K^0$ 表示且
$K^0$ 投射时，对任意 $a \in I$，映射
$a : K^{-1} \to K^{-1}$ 经由
$d_K^{-1} : K^{-1} \to K^0$ 分解。
\end{enumerate}
若 $R$ 是 Noether 环，并且对 $i = -1, 0$，$H^i(K)$ 是有限
$R$-模，则它们还等价于：
\begin{enumerate}
\item[(4)] 对每个有限 $R$-模 $N$，$\Ext^1_R(K, N)$ 被 $I$ 零化；
\item[(5)] $K$ 可以由复形 $K^{-1} \to K^0$ 表示，其中
$K^0$ 有限自由且 $K^{-1}$ 有限，并且对任意 $a \in I$，映射
$a : K^{-1} \to K^{-1}$ 经由
$d_K^{-1} : K^{-1} \to K^0$ 分解。
\end{enumerate}
\end{lemma}

\begin{proof}

\medskip\noindent
假设 (1)，并设 $K^{-1} \to K^0$ 是表示 $K$ 的二项复形，
其中 $K^0$ 投射。以下将直接使用引理
\ref{more-algebra-lemma-map-out-of-almost-free} 给出的以 $K^\bullet$
为源的 $D(R)$ 中映射之描述，不再另行说明。取
$N = K^{-1}$，并考虑 $\Ext^1_R(K, N)$ 中由
$\text{id}_{K^{-1}} : K^{-1} \to K^{-1}$ 给出的元素 $\xi$。
% P02-E0333
由于它被 $a \in I$ 零化，得到下列交换图中的虚线箭头：
$$
\xymatrix{
K^{-1} \ar[d]_a \ar[r]_{d_K^{-1}} & K^0 \ar@{..>}[ld]^h \\
K^{-1}
}
$$
这证明了 (3)。由引理 \ref{more-algebra-lemma-represent-two-term-complex} 的
(1)，(3) 蕴含 (2)。假设 $K^\bullet$ 如 (2) 所述，且 $N$
为任意 $R$-模。$\Ext^1_R(K, N)$ 的任意元素 $\xi$ 都由映射
$\varphi : K^{-1} \to N$ 的类给出。于是，对 $a \in I$，
由假设可以选取上图中的映射 $h$，从而
$$
a\varphi = \varphi \circ a = \varphi \circ h \circ \text{d}_K^{-1}
$$
这证明 $a \xi$ 在 $\Ext^1_R(K, N)$ 中为零。
因此 (1)、(2) 和 (3) 等价。

\medskip\noindent
假设 $R$ 是 Noether 环，并且对 $i = -1, 0$，$H^i(K)$ 是有限
$R$-模。由引理 \ref{more-algebra-lemma-represent-two-term-complex} 的 (2)，
(3) 蕴含 (5)。显然 (5) 蕴含 (2)，而 (1) 显然蕴含 (4)。
因此，为完成证明，只需证明 (4) 蕴含其余任一条件。
由引理 \ref{more-algebra-lemma-represent-two-term-complex} 的 (2)，
设 $K^{-1} \to K^0$ 是表示 $K$ 的复形，其中 $K^0$
有限自由且 $K^{-1}$ 有限。证明 (2) $\Rightarrow$ (1) 时所用的论证表明，
若 $\Ext^1_R(K, K^{-1})$ 被 $I$ 零化，则 (1) 成立。
由此可见 (4) 蕴含 (1)，证明完成。
\end{proof}

\begin{lemma}
\label{more-algebra-lemma-two-term-base-change}
设 $R$ 为环。设 $K$ 为 $D(R)$ 的对象，且对
$i \not \in \{-1, 0\}$ 有 $H^i(K) = 0$。设
$K^{-1} \to K^0$ 是表示 $K$ 的 $R$-模二项复形，并且
$K^0$ 是平坦 $R$-模（例如投射模或自由模）。设
$R \to R'$ 为环同态。则复形 $K^\bullet \otimes_R R'$
表示 $\tau_{\geq -1}(K \otimes_R^\mathbf{L} R')$。
\end{lemma}

\begin{proof}
在 $D(R)$ 中有区分三角
$$
K^0 \to K^\bullet \to K^{-1}[1] \to K^0[1]
$$
它确定区分三角之间的映射
$$
\xymatrix{
K^0 \otimes_R^\mathbf{L} R' \ar[d] \ar[r] &
K^\bullet \otimes_R^\mathbf{L} R' \ar[r] \ar[d] &
K^{-1} \otimes_R^\mathbf{L} R'[1] \ar[r] \ar[d] &
K^0 \otimes_R^\mathbf{L} R'[1] \ar[d] \\
K^0 \otimes_R R' \ar[r] &
K^\bullet \otimes_R R' \ar[r] &
K^{-1} \otimes_R R'[1] \ar[r] &
K^0 \otimes_R R'[1]
}
$$
由于 $K^0$ 平坦，左、右两个竖直箭头都是同构。映射
$K^{-1} \otimes_R^\mathbf{L} R' \to K^{-1} \otimes_R R'$
在次数 $0$ 的上同调上是同构，故得结论。
\end{proof}

\begin{lemma}
\label{more-algebra-lemma-base-change-property-ext-1-annihilated}
设 $I$ 为环 $R$ 的理想。设 $K$ 为 $D(R)$ 的对象，且对
$i \not \in \{-1, 0\}$ 有 $H^i(K) = 0$。设
$R \to R'$ 为环同态。若 $K$ 关于 $(R, I)$ 满足引理
\ref{more-algebra-lemma-ext-1-annihilated-definite} 的等价条件 (1)、(2) 和 (3)，
则 $\tau_{\geq -1}(K \otimes_R^\mathbf{L} R')$ 关于 $(R', IR')$
满足引理 \ref{more-algebra-lemma-ext-1-annihilated-definite} 的等价条件
(1)、(2) 和 (3)。
\end{lemma}

\begin{proof}
可以假设 $K$ 由二项复形 $K^{-1} \to K^0$ 表示，其中 $K^0$
自由，并且对任意 $a \in I$，映射 $a : K^{-1} \to K^{-1}$
等于某个映射 $h_a : K^0 \to K^{-1}$ 所给出的复合
$h_a \circ d_K^{-1}$。由引理 \ref{more-algebra-lemma-two-term-base-change}，
可知 $\tau_{\geq -1}(K \otimes_R^\mathbf{L} R')$ 由
$K^\bullet \otimes_R R'$ 表示。于是对每个 $a \in I$，显然有
$a \otimes 1 : K^{-1} \otimes_R R' \to K^{-1} \otimes_R R'$
等于 $(h_a \otimes 1) \circ (d_K^{-1} \otimes 1)$。
经由 $\text{d}_K^{-1} \otimes 1$ 分解的映射
$K^{-1} \otimes_R R' \to K^{-1} \otimes_R R'$ 全体组成一个
$R'$-模，故得结论。
\end{proof}

\begin{lemma}
\label{more-algebra-lemma-two-term-surjection-map-zero}
设 $R$ 为环。设 $\alpha : K \to K'$ 为 $D(R)$ 中的态射。假设：
\begin{enumerate}
\item 对 $i \not \in \{-1, 0\}$，有 $H^i(K) = H^i(K') = 0$；
\item $H^0(\alpha)$ 是同构，且 $H^{-1}(\alpha)$ 满射。
\end{enumerate}
对任意 $f \in R$，若 $f : K \to K$ 为 $0$，则
$f : K' \to K'$ 为 $0$。
\end{lemma}

\begin{proof}
令 $M = \Ker(H^{-1}(\alpha))$。于是 $\alpha$ 嵌入区分三角
$$
M[1] \to K \to K' \to M[2]
$$
由假设，$K \to K' \xrightarrow{f} K'$ 为零，故
$f : K' \to K'$ 经由某个映射 $M[2] \to K'$ 分解。然而
$\Hom(M[2], K') = 0$；例如可由 Derived Categories，引理
\ref{derived-lemma-negative-exts} 得到。
\end{proof}

\begin{lemma}
\label{more-algebra-lemma-surjection-property-ext-1-annihilated}
设 $I$ 为环 $R$ 的理想。设 $\alpha : K \to K'$ 为
$D(R)$ 中的态射。假设：
\begin{enumerate}
\item 对 $i \not \in \{-1, 0\}$，有 $H^i(K) = H^i(K') = 0$；
\item $H^0(\alpha)$ 是同构，且 $H^{-1}(\alpha)$ 满射。
\end{enumerate}
若 $K$ 满足引理 \ref{more-algebra-lemma-ext-1-annihilated-definite} 的等价条件
(1)、(2) 和 (3)，则 $K'$ 也满足这些条件。
\end{lemma}

\begin{proof}
令 $M = \Ker(H^{-1}(\alpha))$。于是 $\alpha$ 嵌入区分三角
$$
M[1] \to K \to K' \to M[2]
$$
对任意 $R$-模 $N$，它确定正合序列
$$
\Ext^0_R(M[1], N) \to
\Ext^1_R(K', N) \to
\Ext^1_R(K, N)
$$
由于 $\Ext^0_R(M[1], N) = \Ext^{-1}_R(M, N) = 0$，可知
$\Ext^1_R(K', N)$ 是 $\Ext^1_R(K, N)$ 的子模。因此，若
$\Ext^1_R(K, N)$ 被 $I$ 零化，则 $\Ext^1_R(K', N)$ 也如此。
\end{proof}

\begin{lemma}
\label{more-algebra-lemma-ext-1-annihilated}
设 $R$ 为环，% P02-E0334
并设 $I \subset R$ 为理想。设 $K \in D(R)$，且对
$i \not \in \{-1, 0\}$ 有 $H^i(K) = 0$。下列条件等价：
\begin{enumerate}
\item 存在 $c \geq 0$，使得对 $K$ 和理想 $I^c$，
引理 \ref{more-algebra-lemma-ext-1-annihilated-definite} 的等价条件
(1)、(2)、(3) 成立；
\item 存在 $c \geq 0$，使得 (a) $I^c$ 零化 $H^{-1}(K)$，
且 (b) $H^0(K)$ 是 $I^c$-投射模（见第
\ref{more-algebra-section-near-projective} 节）。
\end{enumerate}
若 $R$ 是 Noether 环，并且对 $i = -1, 0$，$H^i(K)$ 是有限
$R$-模，则它们还等价于：
\begin{enumerate}
\item[(3)] 存在 $c \geq 0$，使得对 $K$ 和理想 $I^c$，
引理 \ref{more-algebra-lemma-ext-1-annihilated-definite} 的等价条件 (4)、(5) 成立；
\item[(4)] $H^{-1}(K)$ 是 $I$-幂挠的，并且存在
$f_1, \ldots, f_s \in R$，满足
$V(f_1, \ldots, f_s) \subset V(I)$，使得各局部化
$H^0(K)_{f_i}$ 是投射 $R_{f_i}$-模；
\item[(5)] $H^{-1}(K)$ 是 $I$-幂挠的，并且存在
$f_1, \ldots, f_s \in I$，满足
$V(f_1, \ldots, f_s) = V(I)$，使得各局部化
$H^0(K)_{f_i}$ 是投射 $R_{f_i}$-模；
\item[(6)] $H^{-1}(K)$ 是 $I$-幂挠的，并且对任意满足
$V(f_1, \ldots, f_s) = V(I)$ 的
$f_1, \ldots, f_s \in I$，各局部化 $H^0(K)_{f_i}$
都是投射 $R_{f_i}$-模。
\end{enumerate}
\end{lemma}

\begin{proof}
区分三角
$H^{-1}(K)[1] \to K \to H^0(K)[0] \to H^{-1}(K)[2]$
确定正合序列
$$
0 \to \Ext^1_R(H^0(K), N) \to \Ext^1_R(K, N) \to \Hom_R(H^{-1}(K), N) \to
\Ext^2_R(H^0(K), N)
$$
因此，(2) 蕴含对每个 $R$-模 $N$，$I^{2c}$ 零化
$\Ext^1_R(K, N)$。假设 (1)，立即可知 $H^0(K)$ 是
$I^c$-投射的。另一方面，可以选取某个内射 $R$-模 $N$
以及内射映射 $H^{-1}(K) \to N$。由于上述序列中的
$\Ext^2$ 消失，该映射是 $\Ext^1_R(K, N)$ 中某个元素的像，
故可知 $H^{-1}(K)$ 被 $I^c$ 零化。

\medskip\noindent
假设 $R$ 是 Noether 环，并且对 $i = -1, 0$，$H^i(K)$ 是有限
$R$-模。由引理 \ref{more-algebra-lemma-ext-1-annihilated-definite}，
可知 (3) 等价于 (1) 和 (2)。此外，若 (3) 成立，则对
$f \in I$，$f$ 在 $H^0(K)$ 上的乘法经由投射模分解；
% P02-E0335
这蕴含 $H^0(K)_f$ 是某个投射 $R_f$-模的直和项，因而自身也是投射
$R_f$-模。因此等价条件 (1)、(2) 和 (3) 蕴含 (6)。
当然，(6) 蕴含 (5)，而 (5) 显然蕴含 (4)。

\medskip\noindent
假设 (4)。由于 $H^{-1}(K)$ 是有限 $R$-模且为 $I$-幂挠，
可知对某个 $c_1 \geq 0$，$I^{c_1}$ 零化 $H^{-1}(K)$。
选取短正合序列
$$
0 \to M \to R^{\oplus r} \to H^0(K) \to 0
$$
它确定元素 $\xi \in \Ext^1_R(H^0(K), M)$。对任意 $f \in I$，
由引理 \ref{more-algebra-lemma-pseudo-coherence-and-base-change-ext}，有
$\Ext^1_R(H^0(K), M)_f = \Ext^1_{R_f}(H^0(K)_f, M_f)$。
因此，若 $H^0(K)_f$ 投射，则 $f$ 的某个幂零化 $\xi$。
故对某个 $c_2 \geq 0$，$(f_1, \ldots, f_s)^{c_2}$ 零化 $\xi$。
由于 $V(f_1, \ldots, f_s) \subset V(I)$，有
% P02-E0336
$\sqrt{I} \subset (f_1, \ldots, f_s)$
（Algebra，引理 \ref{algebra-lemma-Zariski-topology}）。
由于 $R$ 是 Noether 环，对某个 $c_3 \geq 0$，有
$I^{c_3} \subset (f_1, \ldots, f_s)$
（Algebra，引理 \ref{algebra-lemma-Noetherian-power}）。
% P02-E0337
因此 $I^{c2c3}$ 零化 $\xi$。这又说明 $H^0(K)$ 是
$I^{c_2c_3}$-投射的（因为对零化 $\xi$ 的 $a \in I$，
% P02-E0338
乘法经由 $R^{\oplus r}$ 分解）。因此取
$c = \max(c_1, c_2c_3)$，可知 (2) 成立。
\end{proof}

\begin{lemma}
\label{more-algebra-lemma-zero-in-derived}
设 $R$ 为环。对 $j = 1, 2, 3$，设 $K_j \in D(R)$，且对
$i \not \in \{-1, 0\}$ 有 $H^i(K_j) = 0$。设
$\varphi : K_1 \to K_2$ 和 $\psi : K_2 \to K_3$ 为
$D(R)$ 中的映射。若 $H^0(\varphi) = 0$ 且
$H^{-1}(\psi) = 0$，则
% P02-E0339
$\varphi \circ \psi = 0$。
\end{lemma}

\begin{proof}
应用 Derived Categories，引理
\ref{derived-lemma-trick-vanishing-composition}，可知
% P02-E0339
$\varphi \circ \psi$ 经由 $\tau_{\leq -2}K_2 = 0$ 分解。
\end{proof}

\begin{lemma}
\label{more-algebra-lemma-silly}
设 $R$ 为环。设 $K \in D(R)$ 由形如
$R^{\oplus n} \to R^{\oplus n}$ 的二项复形给出。
% P02-E0340
用 $A \in \text{Mat}(n \times n, R)$ 表示微分的矩阵。
% P02-E0341
则 $\det(a) : K \to K$ 在 $D(R)$ 中为零。
\end{lemma}

\begin{proof}
省略。这是一个很好的练习。
\end{proof}

\section{朴素余切复形}
\label{more-algebra-section-netherlander}

\noindent
本节继续 Algebra，第 \ref{algebra-section-netherlander} 节中开始的讨论。
先讨论基变换。第一个引理表明，把朴素余切复形与环扩张作朴素张量积，
并不像人们可能以为的那样朴素。

\begin{lemma}
\label{more-algebra-lemma-tensor-NL}
设 $R \to S$ 和 $S \to S'$ 为环同态。典范映射
$\NL_{S/R} \otimes_S^\mathbf{L} S' \to \NL_{S/R} \otimes_S S'$
在 $D(S')$ 中诱导同构
$\tau_{\geq -1}(\NL_{S/R} \otimes_S^\mathbf{L} S') \to \NL_{S/R} \otimes_S S'$。
类似地，给定 $S$ 在 $R$ 上的一个表示 $\alpha$，典范映射
$\NL(\alpha) \otimes_S^\mathbf{L} S' \to \NL(\alpha) \otimes_S S'$
在 $D(S')$ 中诱导同构
$\tau_{\geq -1}(\NL(\alpha) \otimes_S^\mathbf{L} S') \to
\NL(\alpha) \otimes_S S'$。
\end{lemma}

\begin{proof}
这是引理 \ref{more-algebra-lemma-two-term-base-change} 的特例。% P02-E0342
\end{proof}

\begin{lemma}
\label{more-algebra-lemma-base-change-NL}
设 $R \to S$ 和 $R \to R'$ 为环同态。设
$\alpha : P \to S$ 为 $S$ 在 $R$ 上的一个表示。则
$\alpha' : P \otimes_R R' \to S \otimes_R R'$ 是
$S' = S \otimes_R R'$ 在 $R'$ 上的一个表示。典范映射
% P02-E0343
$$
NL(\alpha) \otimes_S S' \to \NL(\alpha')
$$
在 $H^0$ 上是同构，在 $H^{-1}$ 上是满射。特别地，典范映射
$$
\NL_{S/R} \otimes_S S' \to \NL_{S'/R'}
$$
在 $H^0$ 上是同构，在 $H^{-1}$ 上是满射。
\end{lemma}

\begin{proof}
% P02-E0344
记 $I = \Ker(P \to S)$。记 $P' = P \otimes_R R'$，并记
$I' = \Ker(P' \to S')$。假设 $P$ 是以 $j \in J$ 为指标的
$x_j$ 上的多项式代数。引理中显示的映射变为
$$
\xymatrix{
\bigoplus_{j \in J} S' \text{d}x_j \ar[r] &
\bigoplus_{j \in J} S' \text{d}x_j \\
I/I^2 \otimes_S S' \ar[r] \ar[u] &
I'/(I')^2 \ar[u]
}
$$
其中左列为 $\NL(\alpha) \otimes_S S'$，右列为
$\NL(\alpha')$。由张量积的右正合性，可知
$I \otimes_R R' \to I'$ 满射。因此底部箭头满射。
这证明了引理的第一个陈述。关于
$\NL_{S/R} \otimes_S S' \to \NL_{S'/R'}$ 的陈述随之成立，
因为这些复形分别同伦于 $\NL(\alpha) \otimes_S S'$ 和
$\NL(\alpha')$。
\end{proof}

\begin{lemma}
\label{more-algebra-lemma-base-change-NL-flat}
考虑环的余笛卡尔图
$$
\xymatrix{
B \ar[r] & B' \\
A \ar[r] \ar[u] & A' \ar[u]
}
$$
若 $B$ 在 $A$ 上平坦，则典范映射
$\NL_{B/A} \otimes_B B' \to \NL_{B'/A'}$ 是拟同构。
若还假设 $\NL_{B/A}$ 的 Tor 振幅位于 $[-1, 0]$ 中，则
$\NL_{B/A} \otimes_B^\mathbf{L} B' \to \NL_{B'/A'}$
也是拟同构。
\end{lemma}

\begin{proof}
按 Algebra，第 \ref{algebra-section-netherlander} 节选取表示
$\alpha : P \to B$。令 $I = \Ker(\alpha)$，置
$P' = P \otimes_A A'$，并用
% P02-E0345
$\alpha' : P' \to B'$ 表示 $B'$ 在 $A'$ 上的相应表示。
由于 $B$ 在 $A$ 上平坦，可知 $I' = \Ker(\alpha')$ 等于
$I \otimes_A A'$。因此
$$
I'/(I')^2 = \Coker(I^2 \otimes_A A' \to I \otimes_A A') =
I/I^2 \otimes_A A' = I/I^2 \otimes_B B'
$$
由于两边具有相同的基，有
$\Omega_{P'/A'} = \Omega_{P/A} \otimes_A A'$。于是
$\Omega_{P'/A'} \otimes_{P'} B' = \Omega_{P/A} \otimes_P B \otimes_B B'$。
这证明 $\NL(\alpha) \otimes_B B' \to \NL(\alpha')$ 是复形同构，
从而第一个陈述成立。

\medskip\noindent
作为 $B$-模复形，有
$$
\NL(\alpha) = I/I^2 \longrightarrow \Omega_{P/A} \otimes_P B
$$
其中 $I/I^2$ 置于次数 $-1$。由于次数 $0$ 的项自由，
由引理 \ref{more-algebra-lemma-last-one-flat}，该复形的 Tor 振幅位于
$[-1, 0]$ 中，当且仅当 $I/I^2$ 是平坦 $B$-模。
若此条件成立，则
$\NL(\alpha) \otimes_B^\mathbf{L} B' =
\NL(\alpha) \otimes_B B'$，因而得到第二个陈述。
\end{proof}

\begin{lemma}
\label{more-algebra-lemma-lci-NL}
设 $A \to B$ 是定义 \ref{more-algebra-definition-local-complete-intersection}
意义下的局部完全交。则 $\NL_{B/A}$ 是 $D(B)$ 中 Tor 振幅位于
$[-1, 0]$ 的完美对象。
\end{lemma}

\begin{proof}
写成 $B = A[x_1, \ldots, x_n]/I$。则 $\NL_{B/A}$ 由
$B$-模复形
$$
I/I^2 \longrightarrow \bigoplus B \text{d}x_i
$$
表示，其中 $I/I^2$ 置于次数 $-1$。由于次数 $0$ 的项有限自由，
由引理 \ref{more-algebra-lemma-last-one-flat}，该复形的 Tor 振幅位于
$[-1, 0]$ 中，当且仅当 $I/I^2$ 是平坦 $B$-模。
按照定义，$I$ 是 Koszul 正则理想，因而是拟正则理想；见第
\ref{more-algebra-section-ideals} 节。因此 $I/I^2$ 是有限投射 $B$-模
（引理 \ref{more-algebra-lemma-quasi-regular-ideal-finite-projective}），
从而 $\NL_{B/A}$ 既是完美的，又具有 $[-1, 0]$ 内的 Tor 振幅。
\end{proof}

\begin{lemma}
\label{more-algebra-lemma-base-change-lci-bis}
考虑环的余笛卡尔图
$$
\xymatrix{
B \ar[r] & B' \\
A \ar[r] \ar[u] & A' \ar[u]
}
$$
若 $A \to B$ 与 $A' \to B'$ 都是定义
\ref{more-algebra-definition-local-complete-intersection} 意义下的局部完全交，
则映射
$H^{-1}(\NL_{B/A} \otimes_B B') \to H^{-1}(\NL_{B'/A'})$
的核是有限投射 $B'$-模。
\end{lemma}

\begin{proof}
由引理 \ref{more-algebra-lemma-lci-NL}，复形 $\NL_{B/A}$ 和
$\NL_{B'/A'}$ 都是 Tor 振幅位于 $[-1, 0]$ 的完美复形。
结合引理 \ref{more-algebra-lemma-tensor-NL}、\ref{more-algebra-lemma-pull-perfect}
和 \ref{more-algebra-lemma-pull-tor-amplitude}，有
$\NL_{B/A} \otimes_B B' = \NL_{B/A} \otimes_B^\mathbf{L} B'$，
并且该复形也是 Tor 振幅位于 $[-1, 0]$ 的完美复形。
在 $D(B')$ 中选取区分三角
$$
C \to \NL_{B/A} \otimes_B B'  \to \NL_{B'/A'} \to C[1]
$$
由引理 \ref{more-algebra-lemma-two-out-of-three-perfect} 和
\ref{more-algebra-lemma-cone-tor-amplitude}，可知 $C$ 是 Tor 振幅位于
$[-1, 1]$ 的完美复形。由引理 \ref{more-algebra-lemma-base-change-NL}，
复形 $C$ 仅有一个非零上同调模，即引理中的模，位于次数 $-1$。
该模有限表示（引理 \ref{more-algebra-lemma-n-pseudo-module}）且平坦
（引理 \ref{more-algebra-lemma-tor-dimension}）。因此，由 Algebra，引理
\ref{algebra-lemma-finite-projective}，它是有限投射的。
\end{proof}

\section{Abel 群的 Rlim}
\label{more-algebra-section-Rlim}

\noindent
我们简要讨论 Abel 群上的 $R\lim$。本节用
% P02-E0346
$\textit{Ab}(\mathbf{N})$ 表示 Abel 群逆系组成的 Abel 范畴。
这一记号与位点上 Abel 群层的记号相容，因为 Abel 群逆系等同于范畴
$\mathbf{N}$ 上的群层（当 $i \leq j$ 时，恰有一个态射
$i \to j$）；见评注 \ref{more-algebra-remark-Rlim-cohomology}。
本节的许多论证重复了构造位点上 Abel 群层的上同调机制时所用的论证。

\begin{lemma}
\label{more-algebra-lemma-compute-Rlim}
函子 $\lim : \textit{Ab}(\mathbf{N}) \to \textit{Ab}$ 具有右导出函子
\begin{equation}
\label{more-algebra-equation-Rlim}
R\lim : D(\textit{Ab}(\mathbf{N})) \longrightarrow D(\textit{Ab})
\end{equation}
按惯例，置 $R^p\lim(K) = H^p(R\lim(K))$。此外：
\begin{enumerate}
\item 对 $\textit{Ab}(\mathbf{N})$ 中任意 $(A_n)$，当
$p > 1$ 时有 $R^p\lim A_n = 0$；
\item $D(\textit{Ab})$ 的对象 $R\lim A_n$ 由复形
$$
\prod A_n \to \prod A_n,\quad (x_n) \mapsto (x_n - f_{n + 1}(x_{n + 1}))
$$
表示，该复形位于次数 $0$ 和 $1$；
\item 若 $(A_n)$ 是 ML 的，则 $R^1\lim A_n = 0$，即
$(A_n)$ 对 $\lim$ 是右非周期的；
\item 每个 $K^\bullet \in D(\textit{Ab}(\mathbf{N}))$
都拟同构于一个逐项对 $\lim$ 右非周期的复形；
\item 若每个 $K^p = (K^p_n)$ 都对 $\lim$ 右非周期，即
% P02-E0347
$R^1\lim_n K^p_n = 0$，则 $R\lim K$ 由次数 $p$ 的项为
$\lim_n K_n^p$ 的复形表示。
\end{enumerate}
\end{lemma}

\begin{proof}
设 $(A_n)$ 为任意逆系。定义逆系 $(B_n)$，其中
$$
B_n = A_n \oplus A_{n - 1} \oplus \ldots \oplus A_1
$$
转移映射由投影给出。映射 $A_n \to B_n$ 由
$(1, f_n, f_{n - 1} \circ f_n, \ldots, f_2 \circ \ldots \circ f_n)$
给出，其中 $f_i : A_i \to A_{i - 1}$ 是转移映射。
由此可见，每个逆系都是某个 ML 系统的子对象
（Homology，第 \ref{homology-section-inverse-systems} 节）。
由 Derived Categories，引理
\ref{derived-lemma-subcategory-right-acyclics} 和 Homology，引理
\ref{homology-lemma-Mittag-Leffler}，可知每个 ML 系统都对
$\lim$ 右非周期，即 (3) 成立。这已经蕴含
% P02-E0348
$RF$ 定义在 $D^+(\textit{Ab}(\mathbf{N}))$ 上；见
Derived Categories，命题 \ref{derived-proposition-enough-acyclics}。
对 $n > 1$，置
$C_n = A_{n - 1} \oplus \ldots \oplus A_1$，并置 $C_1 = 0$，
转移映射同样由投影给出。于是有逆系的短正合序列
$0 \to (A_n) \to (B_n) \to (C_n) \to 0$，其中
$B_n \to C_n$ 由
$(x_i) \mapsto (x_i - f_{i + 1}(x_{i + 1}))$ 给出。
由于 $(C_n)$ 也是 ML 的，由上面引用的命题可知 (2) 成立，
而这也蕴含 (1)。最后，由 Derived Categories，引理
\ref{derived-lemma-unbounded-right-derived}，可知 $R\lim$
事实上定义在整个 $D(\textit{Ab}(\mathbf{N}))$ 上。
事实上，Derived Categories，引理
\ref{derived-lemma-unbounded-right-derived} 的证明正是通过证明
(4) 和 (5) 来进行的。
\end{proof}

\begin{lemma}
\label{more-algebra-lemma-six-term-Rlim}
设
$$
0 \to (A_i) \to (B_i) \to (C_i) \to 0
$$
为 Abel 群逆系的短正合序列。则有相伴的 $6$ 项正合序列
$0 \to \lim A_i \to \lim B_i \to \lim C_i \to
R^1\lim A_i \to R^1\lim B_i \to R^1\lim C_i \to 0$。
\end{lemma}

\begin{proof}
这由引理 \ref{more-algebra-lemma-compute-Rlim} 中的消失性得到。% P02-E0349
\end{proof}

\noindent
下面是 Homology，引理
\ref{homology-lemma-apply-Mittag-Leffler-again} 的“正确”表述。

\begin{lemma}
\label{more-algebra-lemma-apply-Mittag-Leffler-again}
设
$$
(A^{-2}_n \to A^{-1}_n \to A^0_n \to A^1_n)
$$
为 Abel 群复形的逆系，并用
% P02-E0350
$A^{-2} \to A^{-1} \to A^0 \to A^1$ 表示其极限。
用 $(H_n^{-1})$、$(H_n^0)$ 表示上同调的逆系，并用
$H^{-1}$、$H^0$ 表示复形
$A^{-2} \to A^{-1} \to A^0 \to A^1$ 的上同调。若：
\begin{enumerate}
\item $(A^{-2}_n)$ 和 $(A^{-1}_n)$ 的 $R^1\lim$ 均消失；
\item $(H^{-1}_n)$ 的 $R^1\lim$ 消失；
\end{enumerate}
则 $H^0 = \lim H_n^0$。
\end{lemma}

\begin{proof}
设 $K \in D(\textit{Ab}(\mathbf{N}))$ 为由下述复形系表示的对象：
其第 $n$ 个分量是复形
$A^{-2}_n \to A^{-1}_n \to A^0_n \to A^1_n$。
我们将使用 Derived Categories，引理
\ref{derived-lemma-two-ss-complex-functor} 的两个谱序列来计算
$H^0(R\lim K)$\footnote{为使用这些谱序列，必须证明
$\textit{Ab}(\mathbf{N})$ 有足够多的内射对象。% P02-E0351
Abel 群逆系 $(I_n)$ 是内射的，当且仅当每个 $I_n$
都是内射 Abel 群，并且转移映射是分裂满射。每个系统都嵌入某个
这样的系统中。细节省略。}。第一个谱序列的 $E_1$-页为
$$
\begin{matrix}
0 & 0 & R^1\lim A^0_n  & R^1\lim A^1_n \\
A^{-2} & A^{-1} & A^0  & A^1
\end{matrix}
$$
其微分是横向的，且所有高阶微分均为零。% P02-E0352
第二个谱序列的 $E_2$ 页为
$$
\begin{matrix}
R^1\lim H^{-2}_n & 0 & R^1\lim H^0_n & R^1 \lim H^1_n \\
\lim H^{-2}_n &
\lim H^{-1}_n &
\lim H^0_n &
\lim H^1_n
\end{matrix}
$$
并在此页退化。结论随即得到。
\end{proof}

\begin{lemma}
\label{more-algebra-lemma-pro-isomorphism-iso-Rlim}
设 $(A_n)$ 和 $(B_n)$ 为 Abel 群逆系。pro-系的态射
$\varphi : (A_n) \to (B_n)$ 确定映射
$\lim A_n \to \lim B_n$ 和 $R^1\lim A_n \to R^1\lim B_n$。
若 $\varphi$ 是 pro-同构，则这些映射都是同构。
\end{lemma}

\begin{proof}
关于 pro-系态射的讨论，请见 Categories，例
\ref{categories-example-pro-morphism-inverse-systems}。
映射 $\varphi$ 由 $1 \leq m_1 < m_2 < m_3 < \ldots$ 以及与转移映射
相容的映射 $\varphi_n : A_{m_n} \to B_n$ 给出。置
$\varphi', \varphi'' : \prod A_n \to \prod B_n$ 分别等于
$\varphi'((a_n)) = (\varphi_n(a_{m_n}))$ 以及
$$
\varphi''((a_n)) =
(\varphi_n(a_{m_n} + f(a_{m_n + 1}) + \ldots + f(a_{m_{n + 1} - 1})))
$$
其中每次出现的 $f$ 都表示逆系 $(A_n)$ 的一个适当转移映射。
% P02-E0353
于是交换图
$$
\xymatrix{
\prod A_n \ar[r]_\delta \ar[d]^{\varphi'} &
\prod A_n \ar[d]^{\varphi''} \\
\prod B_n \ar[r]^\delta & \prod B_n
}
$$
中的横向箭头如引理 \ref{more-algebra-lemma-compute-Rlim} 所述。
由此得到所需的映射。该构造具有如下函子性：如果给定逆系
$(C_n)$、整数 $1 \leq m'_1 < m'_2 < m'_3 < \ldots$，
以及与转移映射相容的映射 $\psi_n : B_{m'_n} \to C_n$，
% P02-E0354
则
$(\psi \circ \varphi)' = \psi' \circ \varphi'$ 且
$(\psi \circ \varphi)'' = \psi'' \circ \varphi''$，
其中复合 $\psi \circ \varphi$ 指整数
$1 \leq m_{m'_1} < m_{m'_2} < \ldots$ 以及映射
$\psi_n \circ \varphi_{m'_n} : A_{m_{m'_n}} \to C_n$。
下面证明：若 $B_n = A_n$ 且 $\varphi_n$ 是转移映射，
则所得映射为恒等映射。这将同时证明该构造与
$\varphi$ 的代表元 $(m_n, \varphi_n)$ 的选择无关，并证明引理的
最后一个陈述。

\medskip\noindent
因此，令 $B_n = A_n$，并令 $\varphi_n$ 为转移映射。
设 $(a_n) \in \lim A_n$ 为一个元素。则
$\varphi'((a_n))$ 在次数 $n$ 的分量是 $a_{m_n}$ 的像，
而它等于 $a_n$。这证明了关于 $\lim A_n$ 的陈述。
设 $\xi \in R^1\lim A_n$ 是 $\prod A_n$ 中元素 $(a_n)$ 的类。
考虑元素 $(b_n)$，其中
$$
b_i = a_i + f(a_{i + 1}) + \ldots + f(a_{m_n - 1})
$$
若 $m_{n - 1} < i < m_n$，并且若对某个 $n$ 有 $i = m_n$，
则取 $0$。于是
$(a'_n) = (a_n) - \delta((b_n))$ 定义 $R^1\lim A_n$ 中相同的类，
且计算表明，除非对某个 $n$ 有 $i = m_n$，否则 $a'_i = 0$。
因此可以并且确实假设，除非对某个 $n$ 有 $i = m_n$，否则
$a_i = 0$。注意，
$$
(0, \ldots, -f(a_{m_j}), 0, \ldots, 0, a_{m_j}, 0, \ldots)
$$
在位置 $j$ 和 $m_j$ 上非零，它是
$$
c_j = (0, \ldots, 0, f(a_{m_j}), f(a_{m_j}), \ldots, a_{m_j}, 0, \ldots)
$$
在 $\delta$ 下的像，其中非零项位于位置
$j + 1, \ldots, m_j$。和 $c = \sum c_j$ 在 $\prod A_n$ 中有意义。
回忆 $\varphi''((a_n)) = (a_{m_1}, a_{m_2}, \ldots)$，可知
$$
(a_n) - \varphi''((a_n)) = \delta(c)
$$
证明完成。
\end{proof}

\begin{lemma}
\label{more-algebra-lemma-map-into-Rlim}
设 $\mathcal{D}$ 为三角范畴，$(K_n)$ 为 $\mathcal{D}$ 中对象的逆系，
并设 $K$ 为系统 $(K_n)$ 的一个导出极限。则对
$\mathcal{D}$ 中每个 $L$，有短正合序列
$$
0 \to R^1\lim \Hom_\mathcal{D}(L, K_n[-1]) \to
\Hom_\mathcal{D}(L, K) \to
\lim \Hom_\mathcal{D}(L, K_n) \to 0
$$
\end{lemma}

\begin{proof}
这由 Derived Categories，定义
\ref{derived-definition-derived-limit}、引理
\ref{derived-lemma-representable-homological}，以及上面引理
\ref{more-algebra-lemma-compute-Rlim} 中对 $\lim$ 和 $R^1\lim$ 的描述得到。
\end{proof}

\begin{lemma}
\label{more-algebra-lemma-pro-isomorphism-Rlim}
设 $\mathcal{D}$ 为三角范畴，$(K_n)$ 与 $(M_n)$ 为
$\mathcal{D}$ 中对象的逆系，其导出极限分别为 $K$ 和 $M$。
设 $a : (K_n) \to (M_n)$ 为 pro-对象的 pro-同构。
则可以用 $a$ 构造一个（非典范）同构 $K \to M$。
\end{lemma}

\begin{proof}
由 Derived Categories，评注
\ref{derived-remark-functoriality-derived-limit}，得到一个嵌入定义区分三角之态射
$$
\xymatrix{
K \ar[r] \ar[d] & \prod K_n \ar[r] \ar[d]^{a'} &
\prod K_n \ar[d]^{a''} \\
M \ar[r] & \prod M_n \ar[r] & \prod M_n
}
$$
的箭头 $K \to M$。因此，对 $\mathcal{D}$ 的每个对象 $L$，
得到短正合序列之间的映射
$$
\xymatrix{
0 \ar[r] &
R^1\lim \Hom_\mathcal{D}(L, K_n[-1]) \ar[r] \ar[d] &
\Hom_\mathcal{D}(L, K) \ar[r] \ar[d] &
\lim \Hom_\mathcal{D}(L, K_n) \ar[r] \ar[d] &
0 \\
0 \ar[r] &
R^1\lim \Hom_\mathcal{D}(L, M_n[-1]) \ar[r] &
\Hom_\mathcal{D}(L, M) \ar[r] &
\lim \Hom_\mathcal{D}(L, M_n) \ar[r] &
0
}
$$
参见引理 \ref{more-algebra-lemma-map-into-Rlim} 及其证明。
由引理 \ref{more-algebra-lemma-pro-isomorphism-iso-Rlim}，左、右两个竖直箭头
都是同构\footnote{若 pro-同构 $a$ 由映射
$a_n : K_{m_n} \to M_n$ 给出，则使用相应的映射
$\Hom_\mathcal{D}(L, K_{m_n}) \to \Hom_\mathcal{D}(L, M_n)$，
% P02-E0355
得到 $\Hom$ 上的 pro-同构。因此，引理
\ref{more-algebra-lemma-pro-isomorphism-iso-Rlim} 的证明中构造
$\lim$ 与 $R^1\lim$ 上箭头的方法，与 Derived Categories，评注
\ref{derived-remark-functoriality-derived-limit} 中构造
$a'$ 和 $a''$ 的方法一致。}。因此中间箭头是同构。
由 Yoneda 引理可知，映射 $K \to M$ 是同构。
\end{proof}

\begin{lemma}
\label{more-algebra-lemma-map-from-hocolim}
设 $\mathcal{D}$ 为三角范畴，$(K_n)$ 为 $\mathcal{D}$ 中的对象系，
并设 $K$ 为系统 $(K_n)$ 的一个导出余极限。则对
$\mathcal{D}$ 中每个 $L$，有短正合序列
$$
0 \to R^1\lim \Hom_\mathcal{D}(K_n, L[-1]) \to
\Hom_\mathcal{D}(K, L) \to
\lim \Hom_\mathcal{D}(K_n, L) \to 0
$$
\end{lemma}

\begin{proof}
这由 Derived Categories，定义
\ref{derived-definition-derived-colimit}、引理
\ref{derived-lemma-representable-homological}，以及上面引理
\ref{more-algebra-lemma-compute-Rlim} 中对 $\lim$ 和 $R^1\lim$ 的描述得到。
\end{proof}

\begin{remark}[作为上同调的 Rlim]
\label{more-algebra-remark-Rlim-cohomology}
考虑对象为自然数的范畴 $\mathbf{N}$；若 $j \geq i$，
则从 $i \to j$ 恰有一个箭头。赋予 $\mathbf{N}$ 混沌拓扑
（Sites，例 \ref{sites-example-indiscrete}），于是层
$\mathcal{F}$ 等同于 $\mathbf{N}$ 上的集合逆系
$$
\mathcal{F}_1 \leftarrow \mathcal{F}_2 \leftarrow \mathcal{F}_3
\leftarrow \ldots
$$
注意 $\Gamma(\mathbf{N}, \mathcal{F}) = \lim \mathcal{F}_n$。
对 Abel 群逆系 $\mathcal{F}_n$，有
$$
R^p\lim \mathcal{F}_n = H^p(\mathbf{N}, \mathcal{F})
$$
因为两边都是函子
$\mathcal{F} \mapsto \lim \mathcal{F}_n = H^0(\mathbf{N}, \mathcal{F})$
的高阶右导出函子。因此，$R\lim$ 的存在性也可由
Cohomology on Sites，第
\ref{sites-cohomology-section-cohomology-sheaves} 节和第
\ref{sites-cohomology-section-unbounded} 节的一般理论得到。
\end{remark}

\noindent
下一个引理中的积既可以看作复形的逐项积，也可以看作导出范畴
$D(\textit{Ab})$ 中的积；见 Derived Categories，引理
\ref{derived-lemma-products}。

\begin{lemma}
\label{more-algebra-lemma-distinguished-triangle-Rlim}
设 $K = (K_n^\bullet)$ 为 $D(\textit{Ab}(\mathbf{N}))$ 的对象。
在 $D(\textit{Ab})$ 中存在典范区分三角
$$
R\lim K \to \prod\nolimits_n K_n^\bullet \to \prod\nolimits_n K_n^\bullet
\to R\lim K[1]
$$
换言之，$R\lim K$ 是 $D(\textit{Ab})$ 中逆系
$(K_n^\bullet)$ 的一个导出极限；见 Derived Categories，定义
\ref{derived-definition-derived-limit}。
\end{lemma}

\begin{proof}
假设对每个 $p$，逆系 $(K_n^p)$ 都对 $\lim$ 右非周期。
由引理 \ref{more-algebra-lemma-compute-Rlim}，对每个 $p$ 都得到短正合序列
$$
0 \to \lim_n K^p_n \to \prod\nolimits_n K^p_n \to \prod\nolimits_n K^p_n \to 0
$$
由于由 $\lim_n K^p_n$ 组成的复形按引理
\ref{more-algebra-lemma-compute-Rlim} 计算 $R\lim K$，故该引理在此情形成立。

\medskip\noindent
接着，设 $K = (K_n^\bullet)$ 为一般对象。由引理
\ref{more-algebra-lemma-compute-Rlim}，在 $D(\textit{Ab}(\mathbf{N}))$
中存在拟同构 $K \to L$，使得对每个 $p$，$(L_n^p)$ 非周期。
由于 $\textit{Ab}$ 中积是正合的，$\prod K_n^\bullet$
拟同构于 $\prod L_n^\bullet$；因而上面已证明的关于 $L$ 的结论
蕴含关于 $K$ 的结论。
\end{proof}

\begin{lemma}
\label{more-algebra-lemma-break-long-exact-sequence}
沿用引理 \ref{more-algebra-lemma-distinguished-triangle-Rlim} 的记号，
与该区分三角相伴的上同调长正合序列分解为短正合序列
$$
0 \to R^1\lim_n H^{p - 1}(K_n^\bullet) \to
H^p(R\lim K) \to
\lim_n H^p(K_n^\bullet) \to 0
$$
\end{lemma}

\begin{proof}
该区分三角的长正合序列为
$$
\ldots \to H^p(R\lim K) \to \prod\nolimits_n H^p(K_n^\bullet)
\to \prod\nolimits_n H^p(K_n^\bullet) \to
H^{p + 1}(R\lim K) \to \ldots
$$
根据引理中给出的显式描述，中间映射的核是
$\lim_n H^p(K_n^\bullet)$。由引理
\ref{more-algebra-lemma-compute-Rlim}，该映射的余核是
$R^1\lim_n H^p(K_n^\bullet)$。
\end{proof}

\noindent
{\bf 警告。}$D(\textit{Ab}(\mathbf{N}))$ 的对象是 Abel 群逆系的复形。
也可以把它看作复形的逆系 $(K_n^\bullet)$。然而，这与
$D(\textit{Ab})$ 的对象之逆系{\bf 并不}相同；下一个引理与评注说明二者的差别。

\begin{lemma}
\label{more-algebra-lemma-lift-to-system-complexes-Ab}
设 $(K_n)$ 为 $D(\textit{Ab})$ 的对象之逆系。则存在
$D(\textit{Ab}(\mathbf{N}))$ 的对象 $M = (M_n^\bullet)$，以及
$D(\textit{Ab})$ 中的同构 $M_n^\bullet \to K_n$，使得图
$$
\xymatrix{
M_{n + 1}^\bullet \ar[d] \ar[r] &
M_n^\bullet \ar[d] \\
K_{n + 1} \ar[r] & K_n
}
$$
在 $D(\textit{Ab})$ 中交换。
\end{lemma}

\begin{proof}
具体地，设 $M_1^\bullet$ 为表示 $K_1$ 的 Abel 群复形。
假设已经构造
$M_e^\bullet \to M_{e - 1}^\bullet \to \ldots \to M_1^\bullet$
以及映射 $\psi_i : M_i^\bullet \to K_i$，使得对所有
$n < e$，引理陈述中的图都交换。然后考虑
% P02-E0356
$$
\xymatrix{
& M_n^\bullet \ar[d]^{\psi_n} \\
K_{n + 1} \ar[r] & K_n
}
$$
作为 $D(\textit{Ab})$ 中的图。根据 $D(\textit{Ab})$ 中态射的定义，
可以找到 Abel 群复形 $M_{n + 1}^\bullet$、
$D(\textit{Ab})$ 中的同构 $M_{n + 1}^\bullet \to K_{n + 1}$，
以及复形态射 $M_{n + 1}^\bullet \to M_n^\bullet$；它在
$D(\textit{Ab})$ 中表示复合
$$
K_{n + 1} \to K_n \xrightarrow{\psi_n^{-1}} M_n^\bullet
$$
。
故由归纳法，引理成立。
\end{proof}

\begin{remark}
\label{more-algebra-remark-compare-derived-limit}
设 $(K_n)$ 为 $D(\textit{Ab})$ 的对象之逆系，并设
$K = R\lim K_n$ 为该系统的一个导出极限（见 Derived Categories，
第 \ref{derived-section-derived-limit} 节）。这样的导出极限存在，
因为 $D(\textit{Ab})$ 具有可数积
（Derived Categories，引理 \ref{derived-lemma-products}）。
由引理 \ref{more-algebra-lemma-lift-to-system-complexes-Ab}，也可以把
$(K_n)$ 提升为 $D(\textit{Ab}(\mathbf{N}))$ 的对象 $M$。
于是 $K \cong R\lim M$，其中 $R\lim$ 是函子
(\ref{more-algebra-equation-Rlim})；这是因为由引理
\ref{more-algebra-lemma-distinguished-triangle-Rlim}，$R\lim M$ 也是系统
$(K_n)$ 的导出极限。因此，尽管系统 $(K_n)$ 的提升 $M$
可能有许多同构类，$R\lim M$ 的同构型却与选择无关，
因为它同构于该系统的导出极限 $K = R\lim K_n$。
因此，可以把本节证明的关于 $R\lim$ 的结果应用于导出极限。
例如，对每个 $p \in \mathbf{Z}$，都有典范短正合序列
$$
0 \to R^1\lim H^{p - 1}(K_n) \to H^p(K) \to \lim H^p(K_n) \to 0
$$
因为可以对 $M$ 应用引理 \ref{more-algebra-lemma-break-long-exact-sequence}。
% P02-E0357
不借助 $M$ 的存在性，也可以直接看出这一点：把引理
\ref{more-algebra-lemma-break-long-exact-sequence} 证明中的论证应用于（定义性的）
区分三角 $K \to \prod K_n \to \prod K_n \to K[1]$ 即可。
\end{remark}

\begin{lemma}
\label{more-algebra-lemma-Rlim-pro-equal}
设 $E \to D$ 为 $D(\textit{Ab}(\mathbf{N}))$ 的态射。
设 $(E_n)$、相应地 $(D_n)$ 为分别与 $E$、相应地 $D$ 关联的
$D(\textit{Ab})$ 对象之系统。若
$(E_n) \to (D_n)$ 是 pro-对象的同构，则
$R\lim E \to R\lim D$ 是 $D(\textit{Ab})$ 中的同构。
\end{lemma}

\begin{proof}
该假设特别蕴含 pro-对象 $H^p(E_n)$ 与 $H^p(D_n)$ 同构。
结论由引理 \ref{more-algebra-lemma-break-long-exact-sequence} 的短正合序列
以及引理 \ref{more-algebra-lemma-pro-isomorphism-iso-Rlim} 得到。
\end{proof}

\begin{lemma}[Emmanouil]
\label{more-algebra-lemma-emmanouil}
\begin{reference}
取自 \cite{Emmanouil}。
\end{reference}
设 $(A_n)$ 为 Abel 群逆系。下列条件等价：
\begin{enumerate}
\item $(A_n)$ 是 Mittag-Leffler 的；
\item $R^1\lim A_n = 0$，并且
$\bigoplus_{i \in \mathbf{N}} (A_n)$ 也满足相同结论。
\end{enumerate}
\end{lemma}

\begin{proof}
置 $B = \bigoplus_{i \in \mathbf{N}} (A_n)$，因而
$B = (B_n)$，其中 $B_n = \bigoplus_{i \in \mathbf{N}} A_n$。
若 $(A_n)$ 是 ML 的，则 $B$ 是 ML 的，故由引理
\ref{more-algebra-lemma-compute-Rlim}，有 $R^1\lim A_n = 0$ 且
$R^1\lim B_n = 0$。

\medskip\noindent
反之，假设 $(A_n)$ 不是 ML 的。于是可以选取 $m$、整数列
$m < m_1 < m_2 < \ldots$，以及元素 $x_i \in A_{m_i}$，
使得其像 $y_i \in A_m$ 不在 $A_{m_i + 1} \to A_m$ 的像中。
我们将用元素 $x_i$ 和 $y_i$ 以两种方式证明
$R^1\lim B_n \not = 0$。这将完成引理的证明。

\medskip\noindent
第一种证明。
置 $C = (C_n)$，其中 $C_n = \prod_{i \in \mathbf{N}} A_n$。
有典范单射 $B_n \to C_n$，其余核为 $Q_n$。置 $Q = (Q_n)$。
可以并且确实把 $Q_n$ 的元素 $q_n$ 看作元素序列
$q_n = (q_{n, 1}, q_{n, 2}, \ldots)$，其中 $q_{n, i} \in A_n$，
并对尾部为零的序列取模（换言之，差仅在有限个位置的序列被等同）。
有逆系的短正合序列
$$
0 \to (B_n) \to (C_n) \to (Q_n) \to 0
$$
考虑由下式给出的元素 $q_n \in Q_n$：
$$
q_{n, i} =
\left\{
\begin{matrix}
\text{image of }x_i &\text{若}& m_i \geq n \\
0 & \text{否则}
\end{matrix}
\right.
$$
显然，$q_{n + 1}$ 映到 $q_n$。因此得到
$q = (q_n) \in \lim Q_n$。另一方面，我们断言 $q$ 不在
$\lim C_n \to \lim Q_n$ 的像中。事实上，假设 $c = (c_n)$
映到 $q$。可以写成 $c_n = (c_{n, i})$；由于对
$n' \geq n$ 有 $c_{n', i} \mapsto c_{n, i}$，可知对所有
$n, i, n' \geq n$，有
$c_{n, i} \in \Im(C_{n'} \to C_n)$。特别地，
% P02-E0359
$c_{m, i}$ 在 $A_m$ 中的像属于
$\Im(A_{m_i + 1} \to A_m)$，故不可能等于 $y_i$。
因此 $c_m$ 与 $q_m = (y_1, y_2, y_3, \ldots)$ 在无穷多个位置上不同，
这是矛盾。考察上同调长正合序列
$$
0 \to \lim B_n \to \lim C_n \to \lim Q_n \to R^1\lim B_n
$$
可知最后一个群非零，如所需。

\medskip\noindent
第二种证明。对 $n' \geq n$，用
% P02-E0358
$A_{n, n'} = \Im(A_{n'} \to A_n)$ 表示其像。于是
$y_i \in A_m$，且 $y_i \not \in A_{m, m_i + 1}$。
设 $\xi = (\xi_n) \in \prod B_n$ 为如下元素：
除非 $n = m_i$，否则 $\xi_n = 0$；而
$\xi_{m_i} = (0, \ldots, 0, x_i, 0, \ldots)$，其中 $x_i$
置于第 $i$ 个直和项中。我们断言，$\xi$ 不在引理
\ref{more-algebra-lemma-compute-Rlim} 的映射 $\prod B_n \to \prod B_n$ 的像中。
这表明 $R^1\lim B_n$ 非零，并完成证明。事实上，假设 $\xi$
是 $\eta = (z_1, z_2, \ldots)$ 的像，其中
$z_n = \sum z_{n, i} \in \bigoplus_i A_n$。注意
$x_i = z_{m_i, i} \bmod A_{m_i, m_i + 1}$。
于是 $z_{m_i - 1, i}$ 是 $z_{m_i, i}$ 在
$A_{m_i} \to A_{m_i - 1}$ 下的像，依此类推，故
$z_{m, i}$ 是 $z_{m_i, i}$ 在 $A_{m_i} \to A_m$ 下的像。
因此，$z_{m, i}$ 模 $A_{m, m_i + 1}$ 同余于 $y_i$。
特别地，$z_{m, i} \not = 0$。但
$\sum z_{m, i} \in \bigoplus_i A_m$，所以只有有限多个
$z_{m, i}$ 可以非零，矛盾。
\end{proof}

\begin{lemma}
\label{more-algebra-lemma-Mittag-Leffler}
设
$$
0 \to (A_i) \to (B_i) \to (C_i) \to 0
$$
为 Abel 群逆系的短正合序列。若 $(A_i)$ 和 $(C_i)$ 是 ML 的，
则 $(B_i)$ 也是 ML 的。
\end{lemma}

\begin{proof}
这由引理 \ref{more-algebra-lemma-emmanouil}、无穷直和保持正合这一事实，
以及与 $R\lim$ 相伴的上同调长正合序列得到。
\end{proof}

\begin{lemma}
\label{more-algebra-lemma-Rlim-zero-of-direct-sums}
设 $(A_n)$ 为 Abel 群逆系。下列条件等价：
\begin{enumerate}
\item $(A_n)$ 作为 pro-对象为零；
\item $\lim A_n = 0$ 且 $R^1\lim A_n = 0$，并且
$\bigoplus_{i \in \mathbf{N}} (A_n)$ 也满足相同结论。
\end{enumerate}
\end{lemma}

\begin{proof}
由引理 \ref{more-algebra-lemma-pro-isomorphism-iso-Rlim}，(1) 蕴含 (2)。
假设 (2)。由引理 \ref{more-algebra-lemma-emmanouil}，$(A_n)$ 是 ML 的。
对 $m \geq n$，令 $A_{n, m} = \Im(A_m \to A_n)$，于是
$A_n = A_{n, n} \supset A_{n, n + 1} \supset \ldots$。
注意，$(A_n)$ 作为 pro-对象为零，当且仅当对每个 $n$，
存在 $m \geq n$ 使得 $A_{n, m} = 0$。还要注意，$(A_n)$ 是 ML 的，
当且仅当对每个 $n$，存在 $m_n \geq n$，使得
$A_{n, m} = A_{n, m + 1} = \ldots$。在 ML 情形，显然
$\lim A_n = 0$ 蕴含 $A_{n, m_n} = 0$，因为映射
% P02-E0360
$A_{n + 1, m_{n + 1}} \to A_{n, m}$ 都是满射。
证明完成。
\end{proof}

\section{模的 Rlim}
\label{more-algebra-section-Rlim-modules}

\noindent
我们简要讨论模上的 $R\lim$。本节的许多论证重复了构造环化位点上
模的上同调机制时所用的论证。

\medskip\noindent
设 $(A_n)$ 为环的逆系。我们用
% P02-E0361
$\textit{Mod}(\mathbf{N}, (A_n))$ 表示 Abel 群逆系 $(M_n)$ 的范畴，
其中每个 $M_n$ 都赋予 $A_n$-模结构，% P02-E0362
且转移映射 $M_{n + 1} \to M_n$ 是 $A_{n + 1}$-模映射。
这是一个 Abel 范畴。置 $A = \lim A_n$。给定
$\textit{Mod}(\mathbf{N}, (A_n))$ 的对象 $(M_n)$，极限
$\lim M_n$ 是 $A$-模。

\begin{lemma}
\label{more-algebra-lemma-compute-Rlim-modules}
在上述情形中，% P02-E0363
函子 $\lim : \textit{Mod}(\mathbf{N}, (A_n)) \to \text{Mod}_A$
具有右导出函子
$$
R\lim :
D(\textit{Mod}(\mathbf{N}, (A_n)))
\longrightarrow
D(A)
$$
按惯例，置 $R^p\lim(K) = H^p(R\lim(K))$。此外：
\begin{enumerate}
\item 对 $\textit{Mod}(\mathbf{N}, (A_n))$ 中任意 $(M_n)$，
当 $p > 1$ 时有 $R^p\lim M_n = 0$；
\item $D(\text{Mod}_A)$ 的对象 $R\lim M_n$ 由复形
$$
\prod M_n \to \prod M_n,\quad
(x_n) \mapsto (x_n - f_{n + 1}(x_{n + 1}))
$$
表示，该复形位于次数 $0$ 和 $1$；
\item 若 $(M_n)$ 是 ML 的，则 $R^1\lim M_n = 0$，即
$(M_n)$ 对 $\lim$ 右非周期；
\item 每个 $K^\bullet \in D(\textit{Mod}(\mathbf{N}, (A_n)))$
都拟同构于一个逐项对 $\lim$ 右非周期的复形；
\item 若每个 $K^p = (K^p_n)$ 都对 $\lim$ 右非周期，即
% P02-E0364
$R^1\lim_n K^p_n = 0$，则 $R\lim K$ 由次数 $p$ 的项为
$\lim_n K_n^p$ 的复形表示。
\end{enumerate}
\end{lemma}

\begin{proof}
其证明与引理 \ref{more-algebra-lemma-compute-Rlim} 的证明逐字相同。
\end{proof}

\begin{remark}
\label{more-algebra-remark-Rlim-cohomology-modules}
本评注续接评注 \ref{more-algebra-remark-Rlim-cohomology}。
$\mathbf{N}$ 上的环层正是环的逆系 $(A_n)$。
$(A_n)$ 上的模层正好等同于上面定义的范畴
$\textit{Mod}(\mathbf{N}, (A_n))$ 的对象。引理
\ref{more-algebra-lemma-compute-Rlim-modules} 的导出函子 $R\lim$
就是从模的导出范畴到结构层整体截面环上模的导出范畴的
$R\Gamma(\mathbf{N}, -)$。一般而言，群和模的上同调相同；见
Cohomology on Sites，引理
\ref{sites-cohomology-lemma-cohomology-modules-abelian-agree}。
\end{remark}

\noindent
下一个引理中的积既可以看作复形的逐项积，也可以看作导出范畴
$D(A)$ 中的积；见 Derived Categories，引理
\ref{derived-lemma-products}。

\begin{lemma}
\label{more-algebra-lemma-distinguished-triangle-Rlim-modules}
设 $K = (K_n^\bullet)$ 为
$D(\textit{Mod}(\mathbf{N}, (A_n)))$ 的对象。
在 $D(A)$ 中存在典范区分三角
$$
R\lim K \to \prod\nolimits_n K_n^\bullet \to \prod\nolimits_n K_n^\bullet
\to R\lim K[1]
$$
换言之，$R\lim K$ 是 $D(A)$ 中逆系 $(K_n^\bullet)$ 的一个导出极限；
见 Derived Categories，定义 \ref{derived-definition-derived-limit}。
\end{lemma}

\begin{proof}
其证明与引理 \ref{more-algebra-lemma-distinguished-triangle-Rlim} 的证明完全相同，
但使用引理 \ref{more-algebra-lemma-compute-Rlim-modules} 代替
% P02-E0365
引理 \ref{more-algebra-lemma-compute-Rlim}。
\end{proof}

\begin{lemma}
\label{more-algebra-lemma-break-long-exact-sequence-modules}
沿用引理 \ref{more-algebra-lemma-distinguished-triangle-Rlim-modules} 的记号，
与该区分三角相伴的上同调长正合序列分解为 $A$-模的短正合序列
$$
0 \to R^1\lim_n H^{p - 1}(K_n^\bullet) \to
H^p(R\lim K) \to
\lim_n H^p(K_n^\bullet) \to 0
$$
\end{lemma}

\begin{proof}
其证明与引理 \ref{more-algebra-lemma-break-long-exact-sequence} 的证明完全相同，
但使用引理 \ref{more-algebra-lemma-compute-Rlim-modules} 代替
% P02-E0365
引理 \ref{more-algebra-lemma-compute-Rlim}。
\end{proof}

\noindent
{\bf 警告。}与 Abel 群的情形一样，
$D(\textit{Mod}(\mathbf{N}, (A_n)))$ 的对象
$M = (M_n^\bullet)$ 是模复形的逆系，这与各导出范畴中对象的逆系
{\bf 并不}相同。在下一个引理中，我们证明导出范畴中对象的逆系
总能提升为 $D(\textit{Mod}(\mathbf{N}, (A_n)))$ 的对象。

\begin{lemma}
\label{more-algebra-lemma-lift-to-system-complexes}
设 $(A_n)$ 为环的逆系。假设给定：
\begin{enumerate}
\item 对每个 $n$，给定 $D(A_n)$ 的对象 $K_n$；
\item 对每个 $n$，给定 $D(A_{n + 1})$ 中的映射
$\varphi_n : K_{n + 1} \to K_n$，其中通过
$A_{n + 1} \to A_n$ 作限制，把 $K_n$ 看作
$D(A_{n + 1})$ 的对象。
\end{enumerate}
则存在对象
$M = (M_n^\bullet) \in D(\textit{Mod}(\mathbf{N}, (A_n)))$
以及 $D(A_n)$ 中的同构 $\psi_n : M_n^\bullet \to K_n$，
使得图
$$
\xymatrix{
M_{n + 1}^\bullet \ar[d]_{\psi_{n + 1}} \ar[r] & M_n^\bullet \ar[d]^{\psi_n} \\
K_{n + 1} \ar[r]^{\varphi_n} & K_n
}
$$
在 $D(A_{n + 1})$ 中交换。
\end{lemma}

\begin{proof}
我们详细写出证明。对 $A_n$-模 $T$，用
$T_{A_{n + 1}}$ 表示视为 $A_{n + 1}$-模的同一个模。
% P02-E0366
假设 $K_n^\bullet$ 是表示 $K_n$ 的 $A_n$-模复形。
则 $K_{n, A_{n + 1}}^\bullet$ 是同一个复形，但视为
$A_{n + 1}$-模复形。根据导出范畴的构造，映射
% P02-E0367
$\psi_n$ 可以写成
$$
\psi_n = \tau_n \circ \sigma_n^{-1}
$$
其中 $\sigma_n : L_{n + 1}^\bullet \to K_{n + 1}^\bullet$
是 $A_{n + 1}$-模复形的拟同构，而
$\tau_n : L_{n + 1}^\bullet \to K_{n, A_{n + 1}}^\bullet$
是 $A_{n + 1}$-模复形的映射。

\medskip\noindent
现在归纳构造复形 $M_n^\bullet$。作为基础情形，令
$M_1^\bullet = K_1^\bullet$。假设已经构造
$M_e^\bullet \to M_{e - 1}^\bullet \to \ldots \to M_1^\bullet$
以及复形映射 $\psi_i : M_i^\bullet \to K_i^\bullet$，使得对所有
$n < e$，上面的图
$$
\xymatrix{
M_{n + 1}^\bullet \ar[d]_{\psi_{n + 1}} \ar[rr] & &
M_{n, A_{n + 1}}^\bullet \ar[d]^{\psi_{n, A_{n + 1}}} \\
K_{n + 1}^\bullet & L_{n + 1}^\bullet \ar[l]_{\sigma_n} \ar[r]^{\tau_n} &
K_{n, A_{n + 1}}^\bullet
}
$$
在 $D(A_{n + 1})$ 中交换。然后考虑 $D(A_{e + 1})$ 中的图
$$
\xymatrix{
& & M_{e, A_{e + 1}}^\bullet \ar[d]^{\psi_{e, A_{e + 1}}} \\
K_{e + 1}^\bullet & L_{e + 1}^\bullet \ar[r]^{\tau_e} \ar[l]_{\sigma_e} &
K_{e, A_{e + 1}}^\bullet
}
$$
由于 $\psi_e$ 是拟同构，$\psi_{e, A_{e + 1}}$ 也是拟同构。
根据 $D(A_{e + 1})$ 中态射的定义，可以找到
$A_{e + 1}$-模复形的拟同构
$\psi_{e + 1} : M_{e + 1}^\bullet \to K_{e + 1}^\bullet$，
使得存在 $A_{e + 1}$-模的复形态射
$M_{e + 1}^\bullet \to M_{e, A_{e + 1}}^\bullet$，它在
$D(A_{e + 1})$ 中表示复合
$\psi_{e, A_{e + 1}}^{-1} \circ \tau_e \circ \sigma_e^{-1}$。
故由归纳法，引理成立。
\end{proof}

\begin{remark}
\label{more-algebra-remark-how-unique}
沿用引理 \ref{more-algebra-lemma-lift-to-system-complexes} 的假设。% P02-E0368
先验地，$D(\textit{Mod}(\mathbf{N}, (A_n)))$ 中有许多同构类的对象
$M$ 能给出引理中的系统 $(K_n, \varphi_n)$。对每个这样的 $M$，
可以考察复形 $R\lim M \in D(A)$，其中 $A = \lim A_n$。
由引理 \ref{more-algebra-lemma-distinguished-triangle-Rlim-modules}，
$R\lim M$ 是 $D(A)$ 中逆系 $(K_n)$ 的导出极限。
因此，$R\lim M$ 在 $D(A)$ 中的同构类与构造 $M$ 时所作的选择无关。
特别地，可以把本节证明的关于 $R\lim$ 的结果应用于
$D(A)$ 中逆系的导出极限。例如，对每个 $p \in \mathbf{Z}$，
都有典范短正合序列
$$
0 \to R^1\lim H^{p - 1}(K_n) \to H^p(R\lim K_n) \to \lim H^p(K_n) \to 0
$$
因为可以对 $M$ 应用引理
\ref{more-algebra-lemma-distinguished-triangle-Rlim-modules}。
% P02-E0369
不借助 $M$ 的存在性，也可以直接看出这一点：把引理
\ref{more-algebra-lemma-distinguished-triangle-Rlim-modules} 证明中的论证应用于
导出极限的（定义性）区分三角
$R\lim K_n \to \prod K_n \to \prod K_n \to (R\lim K_n)[1]$
即可。
\end{remark}

\begin{lemma}
\label{more-algebra-lemma-get-ML-system}
设 $(A_n)$ 为环的逆系。每个
$K \in D(\textit{Mod}(\mathbf{N}, (A_n)))$
都可以由复形系 $(M_n^\bullet)$ 表示，使得所有转移映射
$M_{n + 1}^\bullet \to M_n^\bullet$ 都满射。
\end{lemma}

\begin{proof}
设 $K$ 由系统 $(K_n^\bullet)$ 表示。置
$M_1^\bullet = K_1^\bullet$。假设已经构造复形的满射
$M_n^\bullet \to M_{n - 1}^\bullet \to \ldots \to M_1^\bullet$
以及同伦等价 $\psi_e : K_e^\bullet \to M_e^\bullet$，使得对所有
$e < n$，图
$$
\xymatrix{
K_{e + 1}^\bullet \ar[d] \ar[r] & K_e^\bullet \ar[d] \\
M_{e + 1}^\bullet \ar[r] & M_e^\bullet
}
$$
交换。然后考虑图
$$
\xymatrix{
K_{n + 1}^\bullet \ar[r] & K_n^\bullet \ar[d] \\
& M_n^\bullet
}
$$
由 Derived Categories，引理 \ref{derived-lemma-make-surjective}，
可以把复合 $K_{n + 1}^\bullet \to M_n^\bullet$ 分解为
$K_{n + 1}^\bullet \to M_{n + 1}^\bullet \to M_n^\bullet$，
使得第一个箭头是同伦等价，第二个箭头是逐项分裂满射。
由此并结合归纳法，引理成立。
\end{proof}

\begin{lemma}
\label{more-algebra-lemma-get-K-flat-system}
设 $(A_n)$ 为环的逆系。每个
$K \in D(\textit{Mod}(\mathbf{N}, (A_n)))$
都可以由复形系 $(K_n^\bullet)$ 表示，使得每个
$K_n^\bullet$ 都是 K-平坦的。
\end{lemma}

\begin{proof}
首先使用引理 \ref{more-algebra-lemma-get-ML-system}，以复形系
$(M_n^\bullet)$ 表示 $K$，使得所有转移映射
$M_{n + 1}^\bullet \to M_n^\bullet$ 都满射。
接着，取拟同构 $K_1^\bullet \to M_1^\bullet$，其中
$K_1^\bullet$ 是 K-平坦 $A_1$-模复形
（引理 \ref{more-algebra-lemma-K-flat-resolution}）。
假设已经构造
$K_n^\bullet \to K_{n - 1}^\bullet \to \ldots \to K_1^\bullet$
以及复形映射 $\psi_e : K_e^\bullet \to M_e^\bullet$，使得对所有
$e < n$，图
$$
\xymatrix{
K_{e + 1}^\bullet \ar[d] \ar[r] & K_e^\bullet \ar[d] \\
M_{e + 1}^\bullet \ar[r] & M_e^\bullet
}
$$
交换。然后考虑 $D(A_{n + 1})$ 中的图
$$
\xymatrix{
C^\bullet \ar@{..>}[d] \ar@{..>}[r] & K_n^\bullet \ar[d]^{\psi_n} \\
M_{n + 1}^\bullet \ar[r]^{\varphi_n} & M_n^\bullet
}
$$
由于 $M_{n + 1}^\bullet \to M_n^\bullet$ 逐项满射，
置于左上角、各项为
$$
C^p = M_{n + 1}^p \times_{M_n^p} K_n^p
$$
的复形 $C^\bullet$ 拟同构于 $M_{n + 1}^\bullet$（细节省略）。
选取拟同构 $K_{n + 1}^\bullet \to C^\bullet$，其中
$K_{n +1}^\bullet$ K-平坦。故由归纳法，引理成立。
\end{proof}

\begin{lemma}
\label{more-algebra-lemma-derived-tensor-product-systems}
设 $(A_n)$ 为环的逆系。给定
$K, L \in D(\textit{Mod}(\mathbf{N}, (A_n)))$，在
% P02-E0370
$D(\mathbf{N}, (A_n))$ 中有典范导出张量积
$K \otimes^\mathbf{L} L$，它与到 $D(A_n)$ 的映射相容。
该构造关于 $K$ 和 $L$ 对称，并且在每个变量上都是三角范畴的正合函子。
\end{lemma}

\begin{proof}
为 $K$ 选取代表元 $(K_n^\bullet)$，使得每个 $K_n^\bullet$
都是 K-平坦复形（引理 \ref{more-algebra-lemma-get-K-flat-system}）。
然后，对 $L$ 的任意代表元选择 $(L_n^\bullet)$，可以把
$K \otimes^\mathbf{L} L$ 定义为由复形系
$$
(\text{Tot}(K_n^\bullet \otimes_{A_n} L_n^\bullet))
$$
表示的对象。由引理 \ref{more-algebra-lemma-K-flat-quasi-isomorphism}
和 \ref{more-algebra-lemma-derived-tor-quasi-isomorphism-other-side}，
这在两个变量上都是良定义的。它与到 $D(A_n)$ 的映射相容是显然的。
正合性与引理 \ref{more-algebra-lemma-derived-tor-exact} 中完全相同。
\end{proof}

\begin{remark}
\label{more-algebra-remark-constructing-tensor-with-limits-functorially}
设 $A$ 为环，$(E_n)$ 为 $D(A)$ 的对象之逆系。上面已经看到，
导出极限 $R\lim E_n$ 存在。因此，对 $D(A)$ 的每个对象 $K$，
导出极限 $R\lim( K \otimes_A^\mathbf{L} E_n )$ 也存在。
事实上，可以关于 $K$ 函子性地构造这些导出极限，并得到三角范畴的正合函子
$$
R\lim(- \otimes_A^\mathbf{L} E_n) : D(A) \longrightarrow D(A)
$$
具体地，首先把 $(E_n)$ 提升为 $D(\mathbf{N}, A)$ 的对象 $E$；
见引理 \ref{more-algebra-lemma-lift-to-system-complexes}。
（该函子将依赖于这一提升的选择。）接着，注意有“对角”或“常值”函子
$$
\Delta : D(A) \longrightarrow D(\mathbf{N}, A)
$$
它把复形 $K^\bullet$ 映到值为 $K^\bullet$ 的常值复形逆系。
然后直接定义
$$
R\lim(K \otimes_A^\mathbf{L} E_n) = R\lim(\Delta(K)\otimes^\mathbf{L} E)
$$
右边使用引理 \ref{more-algebra-lemma-compute-Rlim-modules} 的函子 $R\lim$
以及引理 \ref{more-algebra-lemma-derived-tensor-product-systems} 的函子
$- \otimes^\mathbf{L} -$。
\end{remark}

\begin{lemma}
\label{more-algebra-lemma-tensor-Rlim-exact}
设 $A$ 为环。设 $E \to D \to F \to E[1]$ 为
$D(\mathbf{N}, A)$ 的区分三角。设 $(E_n)$、相应地 $(D_n)$、
相应地 $(F_n)$ 为分别与 $E$、相应地 $D$、相应地 $F$
关联的 $D(A)$ 对象之系统。则对每个 $K \in D(A)$，
都有典范区分三角
$$
R\lim (K \otimes^\mathbf{L}_A E_n) \to
R\lim (K \otimes^\mathbf{L}_A D_n) \to
R\lim (K \otimes^\mathbf{L}_A F_n) \to
R\lim (K \otimes^\mathbf{L}_A E_n)[1]
$$
位于 $D(A)$ 中；记号如评注
\ref{more-algebra-remark-constructing-tensor-with-limits-functorially}。
\end{lemma}

\begin{proof}
这由评注 \ref{more-algebra-remark-constructing-tensor-with-limits-functorially}
中的构造，以及 $\Delta : D(A) \to D(\mathbf{N}, A)$、
$- \otimes^\mathbf{L} -$ 和 $R\lim$ 都是三角范畴的正合函子这一事实
直接得到。
\end{proof}

\begin{lemma}
\label{more-algebra-lemma-tensor-Rlim-pro-equal}
设 $A$ 为环。设 $E \to D$ 为 $D(\mathbf{N}, A)$ 的态射。
设 $(E_n)$、相应地 $(D_n)$ 为分别与 $E$、相应地 $D$
关联的 $D(A)$ 对象之系统。若 $(E_n) \to (D_n)$ 是 pro-对象的同构，
则对每个 $K \in D(A)$，相应映射
$$
R\lim (K \otimes^\mathbf{L}_A E_n)
\longrightarrow
R\lim (K \otimes^\mathbf{L}_A D_n)
$$
是 $D(A)$ 中的同构（记号如评注
\ref{more-algebra-remark-constructing-tensor-with-limits-functorially}）。
\end{lemma}

\begin{proof}
这由定义和引理 \ref{more-algebra-lemma-Rlim-pro-equal} 得到。% P02-E0371
\end{proof}

\section{挠模}
\label{more-algebra-section-torsion}

\noindent
本节中的“挠模”指支撑在仿射概形 $\Spec(R)$ 的给定闭子集
$V(I)$ 上的模。这不同于、但类似于整环上的挠模概念
（定义 \ref{more-algebra-definition-torsion}）。

\begin{definition}
\label{more-algebra-definition-f-power-torsion}
设 $R$ 为环，$M$ 为 $R$-模。
\begin{enumerate}
\item 设 $I \subset R$ 为理想。若对每个 $m \in M$，存在
$n > 0$ 使得 $I^n m = 0$，则称 $M$ 为{\it $I$-幂挠模}。
\item 设 $f \in R$。若对每个 $m \in M$，存在 $n > 0$
使得 $f^n m = 0$，则称 $M$ 为{\it $f$-幂挠模}。
\end{enumerate}
\end{definition}

\noindent
因此，$f$-幂挠模等同于取 $I = (f)$ 时的 $I$-幂挠模。
对 $R$-模 $M$，将使用记号
$$
M[I^n] = \{m \in M \mid I^nm = 0\}
$$
以及
$$
M[I^\infty] = \bigcup M[I^n]
$$
因此，$M$ 是 $I$-幂挠的，当且仅当
$M = M[I^\infty]$，当且仅当 $M = \bigcup M[I^n]$。

\begin{lemma}
\label{more-algebra-lemma-I-power-torsion-presentation}
设 $R$ 为环，$I$ 为 $R$ 的理想，并设 $M$ 为 $I$-幂挠模。
则 $M$ 具有解消
$$
\ldots \to K_2 \to K_1 \to K_0 \to M \to 0
$$
其中每个 $K_i$ 都是若干 $R/I^n$ 副本的直和，而 $n$ 可变。
% P02-E0372
\end{lemma}

\begin{proof}
有典范满射
$$
\oplus_{m \in M} R/I^{n_m} \to M \to 0
$$
其中 $n_m$ 是满足 $I^{n_m} \cdot m = 0$ 的最小正整数。
上述满射的核也是 $I$-幂挠模。归纳进行，即构造出所需的
$M$ 的解消。
\end{proof}

\begin{lemma}
\label{more-algebra-lemma-torsion-free}
设 $R$ 为环，$I$ 为 $R$ 的理想。对任意 $R$-模 $M$，置
$M[I^n] = \{m \in M \mid I^nm = 0\}$。若 $I$ 有限生成，
则下列条件等价：
\begin{enumerate}
\item $M[I] = 0$；
\item 对所有 $n \geq 1$，有 $M[I^n] = 0$；
\item 若 $I = (f_1, \ldots, f_t)$，则映射
$M \to \bigoplus M_{f_i}$ 是单射。
\end{enumerate}
\end{lemma}

\begin{proof}
这由 Algebra，引理 \ref{algebra-lemma-when-injective-covering} 得到。
\end{proof}

\begin{lemma}
\label{more-algebra-lemma-divide-by-torsion}
设 $R$ 为环，$I$ 为 $R$ 的有限生成理想。
\begin{enumerate}
\item 对任意 $R$-模 $M$，有 $(M/M[I^\infty])[I] = 0$。
\item $I$-幂挠模的扩张仍是 $I$-幂挠的。
\end{enumerate}
\end{lemma}

\begin{proof}
设 $m \in M$。若 $m$ 映到 $(M/M[I^\infty])[I]$ 的元素，
则 $Im \subset M[I^\infty]$。写成 $I = (f_1, \ldots, f_t)$。
于是 $f_i m \in M[I^\infty]$，即对某个 $n_i > 0$，
有 $I^{n_i}f_i m = 0$。因此 $I^Nm = 0$，其中
$N = \sum n_i + 2$。故 $m$ 在 $(M/M[I^\infty])$ 中映到零，
这证明了引理的第一个陈述。

\medskip\noindent
对第二个陈述，设
$0 \to M' \to M \to M'' \to 0$ 为模的短正合序列，其中
$M'$ 和 $M''$ 都是 $I$-幂挠模。则
$M[I^\infty] \supset M'$，因而 $M/M[I^\infty]$ 是 $M''$
的商，所以是 $I$-幂挠的。结合第一个陈述和引理
\ref{more-algebra-lemma-torsion-free}，可知它为零。
\end{proof}

\begin{lemma}
\label{more-algebra-lemma-I-power-torsion}
设 $I$ 为环 $R$ 的有限生成理想。$I$-幂挠模组成 Abel 范畴
$\text{Mod}_R$ 的一个 Serre 子范畴；见 Homology，定义
\ref{homology-definition-serre-subcategory}。
\end{lemma}

\begin{proof}
显然，$I$-幂挠模的子模与商模都是 $I$-幂挠的。% P02-E0373
此外，由引理 \ref{more-algebra-lemma-divide-by-torsion}，两个 $I$-幂挠模的扩张
也是 $I$-幂挠的。因此，由 Homology，引理
\ref{homology-lemma-characterize-serre-subcategory}，引理的陈述成立。
% P02-E0374
\end{proof}

\begin{lemma}
\label{more-algebra-lemma-local-cohomology-closed}
设 $R$ 为环，$I \subset R$ 为有限生成理想。子范畴
$I^\infty\text{-torsion} \subset \text{Mod}_R$
仅依赖于闭子集 $Z = V(I) \subset \Spec(R)$。
事实上，$R$-模 $M$ 是 $I$-幂挠的，当且仅当其支撑包含于 $Z$。
\end{lemma}

\begin{proof}
设 $M$ 为 $R$-模，$x \in M$。若 $x \in M[I^\infty]$，
则对所有 $f \in I$，$x$ 在 $M_f$ 中映到零。因此，对所有
$\mathfrak p \not \supset I$，$x$ 在 $M_\mathfrak p$ 中映到零。
反之，若对所有 $\mathfrak p \not \supset I$，$x$ 在
$M_\mathfrak p$ 中映到零，则对所有 $f \in I$，$x$ 在
$M_f$ 中映到零。因此，若 $I = (f_1, \ldots, f_r)$，
则对某个 $n_i \geq 1$，有 $f_i^{n_i}x = 0$。
于是 $x \in M[I^{\sum n_i}]$。故 $M[I^\infty]$ 是映射
$M \to \prod_{\mathfrak p \not \in Z} M_\mathfrak p$ 的核。
引理的第二个陈述成立，并且它蕴含第一个陈述。
\end{proof}

\noindent
下面两个引理或许应放在别处。

\begin{lemma}
\label{more-algebra-lemma-bounded-above-mod-I}
设 $R$ 为环，$I \subset R$ 为理想。设 $K$ 为 $D(R)$ 的对象，
且当 $i > 0$ 时，有 $H^i(K \otimes_R^\mathbf{L} R/I) = 0$。则：
\begin{enumerate}
\item 对所有 $n \geq 1$ 和 $i > 0$，有
$H^i(K \otimes_R^\mathbf{L} R/I^n) = 0$；
\item 对任意 $I$-幂挠 $R$-模 $N$ 和 $i > 0$，有
$H^i(K \otimes_R^\mathbf{L} N) = 0$；
\item 对任意 $M \in D^b(R)$，若其上同调模 $H^i(M)$
都是 $I$-幂挠的，并且在 $i > 0$ 时为 $0$，则在 $i > 0$ 时有
$H^i(K \otimes_R^\mathbf{L} M) = 0$。
\end{enumerate}
\end{lemma}

\begin{proof}
证明 (2)。可以写成 $N = \bigcup N[I^n]$。由于张量积与余极限交换，
有
$K \otimes_R^\mathbf{L} N = \text{hocolim}_n K \otimes_R^\mathbf{L} N[I^n]$
（细节省略；提示：用 K-平坦复形表示 $K$ 并直接计算）。
因此可以假设 $N$ 被 $I^n$ 零化。考虑 $R$-代数
$R' = R/I^n \oplus N$，其中 $N$ 是平方为零的理想。
只需证明 $D(R')$ 的对象 $K' = K \otimes_R^\mathbf{L} R'$
在正次数上的上同调消失。存在 $R$-代数满射
$R' \to R/I$，其核 $J$ 为幂零的
% P02-E0375
（核中任意 $n$ 个元素的乘积为零）。有
$$
K \otimes_R^\mathbf{L} R/I =
(K \otimes_R^\mathbf{L} R') \otimes_{R'}^\mathbf{L} R/I =
K' \otimes_{R'}^\mathbf{L} R/I
$$
由假设，复形 $K \otimes_R^\mathbf{L} R/I$ 的 Tor 振幅位于
$[-\infty, 0]$ 中。因此，由引理
\ref{more-algebra-lemma-check-tor-dimension-modulo-nilpotent} 得到结论。
% P02-E0376

\medskip\noindent
(1) 从 (2) 平凡地得到。(3) 由 (2)、对 $M$ 的非零上同调模个数的归纳，
以及 Derived Categories，评注
\ref{derived-remark-truncation-distinguished-triangle} 的截断区分三角得到。
细节省略。
\end{proof}

\begin{lemma}
\label{more-algebra-lemma-derived-vanishing-mod-I}
设 $R$ 为环，$I \subset R$ 为理想。设 $K$ 为 $D(R)$ 的对象，
且 $K \otimes_R^\mathbf{L} R/I = 0$ 在 $D(R)$ 中成立。则：
\begin{enumerate}
\item 对所有 $n \geq 1$，有 $K \otimes_R^\mathbf{L} R/I^n = 0$；
\item 对任意 $I$-幂挠 $R$-模 $N$，有
$K \otimes_R^\mathbf{L} N = 0$；
\item 对任意 $M \in D^b(R)$，若其上同调模都是 $I$-幂挠的，
则 $K \otimes_R^\mathbf{L} M = 0$。
\end{enumerate}
\end{lemma}

\begin{proof}
这是引理 \ref{more-algebra-lemma-bounded-above-mod-I} 的推论（省略一些细节）。
% P02-E0377
\end{proof}

\begin{lemma}
\label{more-algebra-lemma-torsion-completion}
\begin{reference}
\cite[Lemme~1]{Beauville-Laszlo} 的轻微推广。
\end{reference}
设 $R \to R'$ 为环同态。设 $I \subset R$ 为理想，使得对
$n > 0$，$R/I^n \to R'/I^nR'$ 是同构。对任意 $I$-幂挠
$R$-模 $M$，映射 $M \to M \otimes_R R'$ 是同构。
例如，若 $I$ 有限生成，且 $R^\wedge$ 是 $R$ 关于 $I$ 的完备化，
则有 $M \cong M \otimes_R R^\wedge$。
\end{lemma}

\begin{proof}
若 $M$ 被 $I^n$ 零化，则
$$
M \otimes_R R' \cong
M \otimes_{R/I^n} R'/I^n R' \cong
M \otimes_{R/I^n} R/I^n \cong M.
$$
若 $M$ 是 $I$-幂挠的，则 $M = \bigcup M[I^n]$。
由于张量积与直极限交换
（Algebra，引理 \ref{algebra-lemma-tensor-products-commute-with-limits}），
得到所需同构。最后一个陈述由 Algebra，引理
\ref{algebra-lemma-hathat-finitely-generated} 可知是第一个陈述的特例。
\end{proof}

\section{模范畴的形式粘合}
\label{more-algebra-section-formal-glueing}

\noindent
固定一个 Noether 概形 $X$，以及一个闭子概形 $Z$，其补集为 $U$。
我们的目标是说明，如何从 $X$ 沿 $Z$ 的形式完备化上的凝聚层，
以及 $U$ 上在交叠处满足适当相容性的凝聚层，（唯一地）构造出
$X$ 上的凝聚层。我们先仅用交换代数完成此事（本节），随后在
代数空间的框架下说明它
（Pushouts of Spaces，第 \ref{spaces-pushouts-section-formal-glueing} 节）。

\medskip\noindent
以下文献讨论了本节的部分内容：
\cite[Section 2]{ArtinII}、
\cite[Appendix]{Ferrand-Raynaud}、
\cite{Beauville-Laszlo}、
\cite{MB} 以及
\cite[Section 4.6]{dJ-crystalline}。

\begin{lemma}
\label{more-algebra-lemma-characterize-flatness-on-torsion}
设 $\varphi : R \to S$ 为环映射，$I \subset R$ 为理想。
下列条件等价：
\begin{enumerate}
\item $\varphi$ 平坦，且 $R/I \to S/IS$ 忠实平坦；
\item $\varphi$ 平坦，且映射
$\Spec(S/IS) \to \Spec(R/I)$ 满射；
\item $\varphi$ 平坦，且基变换函子
$M \mapsto M \otimes_R S$ 在被 $I$ 零化的模上是忠实的；
\item $\varphi$ 平坦，且基变换函子
$M \mapsto M \otimes_R S$ 在 $I$-幂挠模上是忠实的。
\end{enumerate}
\end{lemma}

\begin{proof}
若 $R \to S$ 平坦，则对每个 $n$，$R/I^n \to S/I^nS$ 都平坦；见
Algebra，引理 \ref{algebra-lemma-flat-base-change}。
故由 Algebra，引理 \ref{algebra-lemma-ff-rings}，(1) 与 (2) 等价。
将 $I$-挠 $R$-模等同于 $R/I$-模，并利用对被 $I$ 零化的
$R$-模 $M$ 有
$$
M \otimes_R S = M \otimes_{R/I} S/IS
$$
再由 Algebra，引理 \ref{algebra-lemma-easy-ff}，可得 (1) 与 (3) 等价。
(4) $\Rightarrow$ (3) 是立即的。假设 (3)。上面已经看到
$R/I^n \to S/I^nS$ 平坦；由假设，它在谱上诱导满射，因为
$\Spec(R/I^n) = \Spec(R/I)$，对 $S$ 亦然。故基变换函子在被
$I^n$ 零化的模上是忠实的。任何 $I$-幂挠模 $M$ 都是并
$M = \bigcup M_n$，其中 $M_n$ 被 $I^n$ 零化，因此基变换函子
在所有 $I$-幂挠模组成的范畴上是忠实的
（因为张量积与余极限交换）。
% P02-E0378
\end{proof}

\begin{lemma}
\label{more-algebra-lemma-neighbourhood-isomorphism}
假设 $(\varphi : R \to S, I)$ 满足引理
\ref{more-algebra-lemma-characterize-flatness-on-torsion} 的等价条件。
下列条件等价：
\begin{enumerate}
\item 对任意 $I$-幂挠模 $M$，自然映射
$M \to M \otimes_R S$ 是同构；
\item $R/I \to S/IS$ 是同构。
\end{enumerate}
\end{lemma}

\begin{proof}
(1) $\Rightarrow$ (2) 是立即的。假设 (2)。先假设 $M$ 被 $I$ 零化。
此时 $M$ 是 $R/I$-模，故有同构
$$
M \otimes_R S = M \otimes_{R/I} S/IS = M \otimes_{R/I} R/I = M
$$
从而结论成立。接着归纳证明：对任意被 $I^n$ 零化的模 $M$，
$M \to M \otimes_R S$ 是同构。
% P02-E0379
假设归纳假设对 $n$ 成立，并设 $M$ 被 $I^{n + 1}$ 零化。于是有短正合列
$$
0 \to I^nM \to M \to M/I^nM \to 0
$$
由于 $R \to S$ 平坦，这给出短正合列
$$
0 \to I^nM \otimes_R S \to M \otimes_R S \to M/I^nM \otimes_R S \to 0
$$
对 $M' = I^nM$ 和 $M'' = M/I^nM$，规范映射由归纳假设是同构，
故对 $M$ 亦然。最后，设 $M$ 为一般的 $I$-幂挠模。则
$M = \bigcup M_n$，其中 $M_n$ 被 $I^n$ 零化；再利用张量积与
余极限交换即可得出结论。
\end{proof}

\begin{lemma}
\label{more-algebra-lemma-neighbourhood-equivalence}
设 $\varphi : R \to S$ 为平坦环映射，$I \subset R$ 为有限生成理想，
且 $R/I \to S/IS$ 是同构。则：
\begin{enumerate}
\item 对任意 $R$-模 $M$，映射 $M \to M \otimes_R S$ 在
$I$-幂挠子模上诱导同构
$M[I^\infty] \to (M \otimes_R S)[(IS)^\infty]$；
\item 若 $M$ 或 $N$ 是 $I$-幂挠的，则自然映射
$$
\Hom_R(M, N) \longrightarrow \Hom_S(M \otimes_R S, N \otimes_R S)
$$
是同构；
\item 基变换函子 $M \mapsto M \otimes_R S$ 给出 $I$-幂挠模范畴与
$IS$-幂挠模范畴之间的范畴等价。
\end{enumerate}
\end{lemma}

\begin{proof}
引理 \ref{more-algebra-lemma-characterize-flatness-on-torsion} 与引理
\ref{more-algebra-lemma-neighbourhood-isomorphism} 的等价条件都成立。
下文将不再说明而直接使用这些条件。
先证明 (1)。设 $M$ 为任意 $R$-模。置
$M' = M/M[I^\infty]$，并考虑正合列
$$
0 \to M[I^\infty] \to M \to M' \to 0
$$
由于 $M[I^\infty] = M[I^\infty] \otimes_R S$，只需证明
$(M' \otimes_R S)[(IS)^\infty] = 0$。
写成 $I = (f_1, \ldots, f_t)$。由引理
\ref{more-algebra-lemma-divide-by-torsion}，有 $M'[I^\infty] = 0$。
故对每个 $n > 0$，映射
$$
M' \longrightarrow \bigoplus\nolimits_{i = 1, \ldots t} M',
\quad
x \longmapsto (f_1^n x, \ldots, f_t^n x)
$$
是单射。由于 $S$ 在 $R$ 上平坦，相应的映射
$M' \otimes_R S \to \bigoplus_{i = 1, \ldots t} M' \otimes_R S$
也为单射。这意味着 $(M' \otimes_R S)[I^n] = 0$，如所需。

\medskip\noindent
接着证明 (2)。若 $N$ 是 $I$-幂挠的，则
$N \otimes_R S = N$，且由 Algebra，引理
\ref{algebra-lemma-adjoint-tensor-restrict}，(2) 中所显示的映射是同构。
若 $M$ 是 $I$-幂挠的，则任意映射 $M \to N$ 的像都经由
$N[I^\infty]$ 分解，而任意映射
$M \otimes_R S \to N \otimes_R S$ 的像都经由
$(N \otimes_R S)[(IS)^\infty]$ 分解。因此在此情形，(1) 保证我们
可以用 $N[I^\infty]$ 替换 $N$；结论随即归结为刚讨论过的
$N$ 是 $I$-幂挠的情形。

\medskip\noindent
最后证明 (3)。由 (2)，该函子是全忠实的。关于本质满，只需注意：
对任意 $IS$-幂挠 $S$-模 $N$，自然映射 $N \otimes_R S \to N$ 是同构。
\end{proof}

\begin{lemma}
\label{more-algebra-lemma-map-identifies-koszul-and-cech-complexes}
设 $\varphi : R \to S$ 为平坦环映射，$I \subset R$ 为有限生成理想，
且 $R/I \to S/IS$ 是同构。对任意满足
$V(f_1, \ldots, f_r) = V(I)$ 的 $f_1, \ldots, f_r \in R$，有：
\begin{enumerate}
\item Koszul 复形的映射
$K(R, f_1, \ldots, f_r) \to K(S, f_1, \ldots, f_r)$ 是拟同构；
\item 扩充交错 {\v C}ech 复形的映射
$$
\xymatrix{
R \to \prod_{i_0} R_{f_{i_0}} \to \prod_{i_0 < i_1} R_{f_{i_0}f_{i_1}}
\to \ldots \to R_{f_1\ldots f_r} \ar[d] \\
S \to \prod_{i_0} S_{f_{i_0}} \to \prod_{i_0 < i_1} S_{f_{i_0}f_{i_1}}
\to \ldots \to S_{f_1\ldots f_r}
}
$$
是拟同构。
\end{enumerate}
\end{lemma}

\begin{proof}
在两种情形下都有一个由 $R$ 模组成的复形 $K_\bullet$，并且要证明
$K_\bullet \to K_\bullet \otimes_R S$ 是拟同构。
% P02-E0380
由引理 \ref{more-algebra-lemma-neighbourhood-isomorphism} 及 $R \to S$ 的平坦性，
只要 $K$ 的所有同调群都是 $I$-幂挠的，结论即成立。
对 Koszul 复形，这由引理 \ref{more-algebra-lemma-homotopy-koszul} 成立；
对扩充交错 {\v C}ech 复形，这由引理
\ref{more-algebra-lemma-extended-alternating-torsion} 成立。
\end{proof}

\begin{lemma}
\label{more-algebra-lemma-naive-Koszul-complex}
设 $R$ 为环，$I = (f_1, \ldots, f_n)$ 为 $R$ 的有限生成理想。
设 $M$ 为由元素 $e_1, \ldots, e_n$ 生成、以关系
$f_i e_j - f_j e_i = 0$ 为约束的 $R$-模。存在短正合列
$$
0 \to K \to M \to I \to 0
$$
其中 $K$ 被 $I$ 零化。
\end{lemma}

\begin{proof}
这只是 Koszul 复形的一个截断。映射 $M \to I$ 由规则
$e_i \mapsto f_i$ 决定。若 $m = \sum a_i e_i$ 属于 $M \to I$ 的核，
即 $\sum a_i f_i = 0$，则
$f_j m = \sum f_j a_i e_i = (\sum f_i a_i) e_j = 0$。
\end{proof}

\begin{lemma}
\label{more-algebra-lemma-explicit-ext}
设 $R$ 为环，$I = (f_1, \ldots, f_n)$ 为 $R$ 的有限生成理想。
对任意 $R$-模 $N$，置
$$
H_1(N, f_\bullet) =
\frac{\{(x_1, \ldots, x_n) \in N^{\oplus n} \mid f_i x_j = f_j x_i \}}
{\{f_1x, \ldots, f_nx) \mid x \in N\}}
$$
% P02-E0381
对任意 $R$-模 $N$，存在规范短正合列
% P02-E0382
$$
0 \to \Ext_R(R/I, N) \to H_1(N, f_\bullet) \to \Hom_R(K, N)
$$
其中 $K$ 如引理 \ref{more-algebra-lemma-naive-Koszul-complex} 中所述。
\end{lemma}

\begin{proof}
上面的记号表示 $\text{Mod}_R$ 中的 $\Ext$-群，其定义见
Homology，第 \ref{homology-section-extensions} 节。
它们记作 $\Ext_R(M, N)$。利用短正合列
$0 \to I \to R \to R/I \to 0$ 所对应的
Homology，引理 \ref{homology-lemma-six-term-sequence-ext} 的长正合列，
以及 $\Ext_R(R, N) = 0$，可得
$$
\Ext_R(R/I, N) =
\Coker(N \longrightarrow \Hom(I, N))
$$
利用引理 \ref{more-algebra-lemma-naive-Koszul-complex} 的短正合列，可得到复形
$$
N \to \Hom(M, N) \to \Hom_R(K, N)
$$
其中间同调规范同构于 $\Ext_R(R/I, N)$。由于第一个映射的余核
规范同构于 $H_1(N, f_\bullet)$，引理得证。
\end{proof}

\begin{lemma}
\label{more-algebra-lemma-koszul-homology-annihilated}
设 $R$ 为环，$I = (f_1, \ldots, f_n)$ 为 $R$ 的有限生成理想。
对任意 $R$-模 $N$，引理 \ref{more-algebra-lemma-explicit-ext} 中定义的 Koszul 同调群
$H_1(N, f_\bullet)$ 被 $I$ 零化。
\end{lemma}

\begin{proof}
设 $(x_1, \ldots, x_n) \in N^{\oplus n}$，且 $f_i x_j = f_j x_i$。
则 $f_i(x_1, \ldots, x_n) = (f_i x_i, \ldots, f_i x_n)$。
% P02-E0383
换言之，$f_i$ 零化 $H_1(N, f_\bullet)$。
\end{proof}

\noindent
我们可以加强引理 \ref{more-algebra-lemma-neighbourhood-equivalence} 的全忠实性：
若 $\Ext$-群的源是 $I$-幂挠的，则传到 $S$ 同样不改变该群。
下一个引理的导出版本见 Dualizing Complexes，引理
\ref{dualizing-lemma-neighbourhood-extensions}。

\begin{lemma}
\label{more-algebra-lemma-neighbourhood-extensions}
假设 $\varphi : R \to S$ 为平坦环映射，$I \subset R$ 为有限生成理想，
且 $R/I \to S/IS$ 是同构。设 $M$、$N$ 为 $R$-模，并假设 $M$ 是
$I$-幂挠的。给定短正合列
% P02-E0384
$$
0 \to N \otimes_R S \to \tilde E \to M \otimes_R S \to 0
$$
则存在交换图
$$
\xymatrix{
0 \ar[r] &
N \ar[r] \ar[d] &
E \ar[r] \ar[d] &
M \ar[r] \ar[d] &
0 \\
0 \ar[r] &
N \otimes_R S \ar[r] &
\tilde E \ar[r] &
M \otimes_R S \ar[r] &
0
}
$$
且两行均正合。
\end{lemma}

\begin{proof}
因为 $M$ 是 $I$-幂挠的，由引理
\ref{more-algebra-lemma-neighbourhood-isomorphism} 可知 $M \otimes_R S = M$。
下文将不再说明而使用这一等同。由于 $R \to S$ 平坦，基变换函子正合，
从而得到 $\Ext$-群的函子性映射
$$
\Ext_R(M, N)
\longrightarrow
\Ext_S(M \otimes_R S, N \otimes_R S),
$$
参见 Homology，引理 \ref{homology-lemma-exact-functor-ext}。
引理的断言是：当 $M$ 是 $I$-幂挠的时，此映射满射。
事实上，我们将证明它是同构。由引理
\ref{more-algebra-lemma-I-power-torsion-presentation}，可以找到满射 $M' \to M$，
其中 $M'$ 是形如 $R/I^n$ 的模的直和。利用 Homology，引理
\ref{homology-lemma-six-term-sequence-ext} 的长正合列以及引理
\ref{more-algebra-lemma-neighbourhood-equivalence}，可知只需对 $M'$ 证明该引理。
利用 $\Ext$ 与直和的相容性（略去细节），可约化到
$M = R/I^n$ 的情形，其中 $n$ 为某个整数。

\medskip\noindent
设 $f_1, \ldots, f_t$ 为 $I^n$ 的生成元。由引理
\ref{more-algebra-lemma-explicit-ext}，有交换图
$$
\xymatrix{
0 \ar[r] &
\Ext_R(R/I^n, N) \ar[r] \ar[d] &
H_1(N, f_\bullet) \ar[r] \ar[d] &
\Hom_R(K, N) \ar[d] \\
0 \ar[r] &
\Ext_S(S/I^nS, N \otimes S) \ar[r] &
H_1(N \otimes S, f_\bullet) \ar[r] &
\Hom_S(K \otimes S, N \otimes S)
}
$$
% P02-E0385
两行正合，其中 $K$ 如引理
\ref{more-algebra-lemma-naive-Koszul-complex} 中所述。
因此只需证明右侧两个竖直箭头是同构。由于 $K$ 被 $I^n$ 零化，
由引理 \ref{more-algebra-lemma-neighbourhood-equivalence}，有
$\Hom_R(K, N) = \Hom_S(K \otimes_R S, N \otimes_R S)$。
由于 $R \to S$ 平坦，有
$H_1(N, f_\bullet) \otimes_R S = H_1(N \otimes_R S, f_\bullet)$。
由引理 \ref{more-algebra-lemma-koszul-homology-annihilated}，
$H_1(N, f_\bullet)$ 被 $I^n$ 零化，故由引理
\ref{more-algebra-lemma-neighbourhood-isomorphism}，
$H_1(N, f_\bullet) \otimes_R S = H_1(N, f_\bullet)$。
\end{proof}

\noindent
设 $R \to S$ 为环映射。设 $f_1, \ldots, f_t \in R$，并置
$I = (f_1, \ldots, f_t)$。于是对任意 $R$-模 $M$，可定义复形
\begin{equation}
\label{more-algebra-equation-glueing-complex}
0 \to M \xrightarrow{\alpha}
M \otimes_R S \times \prod M_{f_i} \xrightarrow{\beta}
\prod (M \otimes_R S)_{f_i}
\times
\prod M_{f_if_j}
\end{equation}
其中 $\alpha(m) = (m \otimes 1, m/1, \ldots, m/1)$，且
$$
\beta(m', m_1, \ldots, m_t) =
((m'/1 - m_1 \otimes 1, \ldots, m'/1 - m_t \otimes 1),
(m_1 - m_2, \ldots, m_{t - 1} - m_t).
$$
% P02-E0386
% P02-E0387
我们希望知道此复形何时正合。

\begin{lemma}
\label{more-algebra-lemma-recover-module-from-glueing-data}
假设 $\varphi : R \to S$ 为平坦环映射，且
$I = (f_1, \ldots, f_t) \subset R$ 是满足
$R/I \to S/IS$ 为同构的理想。设 $M$ 为 $R$-模。
则复形 (\ref{more-algebra-equation-glueing-complex}) 正合。
\end{lemma}

\begin{proof}
第一种证明。用
$\check{\mathcal{C}}_R \to \check{\mathcal{C}}_S$ 表示引理
\ref{more-algebra-lemma-map-identifies-koszul-and-cech-complexes} 中扩充交错
{\v C}ech 复形的拟同构。由于这些复形有界且各项平坦，
$M \otimes_R \check{\mathcal{C}}_R \to M \otimes_R \check{\mathcal{C}}_S$
也是拟同构（引理 \ref{more-algebra-lemma-derived-tor-quasi-isomorphism} 与
\ref{more-algebra-lemma-derived-tor-quasi-isomorphism-other-side}）。
复形 (\ref{more-algebra-equation-glueing-complex}) 是映射
$M \otimes_R \check{\mathcal{C}}_R \to M \otimes_R \check{\mathcal{C}}_S$
的锥的一个截断，结论即得。

\medskip\noindent
第二种计算性证明。
设 $m \in M$。若 $\alpha(m) = 0$，则由引理
\ref{more-algebra-lemma-torsion-free}，$m \in M[I^\infty]$。选取 $n$ 使
$I^n m = 0$，并考虑映射 $\varphi : R/I^n \to M$。
若 $m \otimes 1 = 0$，则 $\varphi \otimes 1_S = 0$，故
$\varphi = 0$（见引理 \ref{more-algebra-lemma-neighbourhood-equivalence}），
从而 $m = 0$。这说明 $\alpha$ 是单射。

\medskip\noindent
设 $(m', m'_1, \ldots, m'_t) \in \Ker(\beta)$。
对某个 $n > 0$ 及 $m_i \in M$，写成 $m'_i = m_i/f_i^n$。
必要时增大 $n$，可以假设
% P02-E0388
$f_i^n m' = m_i \otimes 1$ 在 $M \otimes_R S$ 中成立，且
$f_j^nm_i - f_i^nm_j = 0$ 在 $M$ 中成立。特别地，
$(m_1, \ldots, m_t)$ 定义了
$H_1(M, (f_1^n, \ldots, f_t^n))$ 的一个元素 $\xi$。
由于 $H_1(M, (f_1^n, \ldots, f_t^n))$ 被 $I^{tn + 1}$ 零化
（见引理 \ref{more-algebra-lemma-koszul-homology-annihilated}），且
$R \to S$ 平坦，由引理
\ref{more-algebra-lemma-neighbourhood-isomorphism} 可得
$$
H_1(M, (f_1^n, \ldots, f_t^n)) =
H_1(M, (f_1^n, \ldots, f_t^n)) \otimes_R S =
H_1(M \otimes_R S, (f_1^n, \ldots, f_t^n))
$$
% P02-E0389
$m'$ 的存在性蕴含 $\xi$ 在最后一个群中的像为零，即元素
$\xi$ 为零。因此存在 $m \in M$，使得 $m_i = f_i^n m$。于是对某个
$m'' \in (M \otimes_R S)[(IS)^\infty]$，有
$(m', m'_1, \ldots, m'_t) - \alpha(m)
= (m'', 0, \ldots, 0)$。由引理
\ref{more-algebra-lemma-neighbourhood-equivalence}，可知
$m'' \in M[I^\infty]$，结论即得。
\end{proof}

\begin{remark}
\label{more-algebra-remark-glueing-data}
本评注定义粘合数据的一个范畴。设 $R \to S$ 为环映射。
设 $f_1, \ldots, f_t \in R$，并置 $I = (f_1, \ldots, f_t)$。
范畴 $\text{Glue}(R \to S, f_1, \ldots, f_t)$ 定义如下：
\begin{enumerate}
\item 其对象为系统 $(M', M_i, \alpha_i, \alpha_{ij})$，其中
$M'$ 是 $S$-模，$M_i$ 是 $R_{f_i}$-模，
$\alpha_i : (M')_{f_i} \to M_i \otimes_R S$ 是同构，并且
$\alpha_{ij} : (M_i)_{f_j} \to (M_j)_{f_i}$ 是满足下列条件的同构：
\begin{enumerate}
\item 作为 $(M')_{f_if_j} \to (M_j)_{f_i}$ 的映射，
$\alpha_{ij} \circ \alpha_i = \alpha_j$；
\item 作为 $(M_i)_{f_jf_k} \to (M_k)_{f_if_j}$ 的映射，
$\alpha_{jk} \circ \alpha_{ij} = \alpha_{ik}$（上链条件）。
\end{enumerate}
\item 态射
$(M', M_i, \alpha_i, \alpha_{ij}) \to (N', N_i, \beta_i, \beta_{ij})$
由映射 $\varphi' : M' \to N'$ 与 $\varphi_i : M_i \to N_i$ 给出，
并与给定映射 $\alpha_i, \beta_i, \alpha_{ij}, \beta_{ij}$ 相容。
\end{enumerate}
有规范函子
$$
\text{Can} : \text{Mod}_R
\longrightarrow
\text{Glue}(R \to S, f_1, \ldots, f_t),
\quad
M \longmapsto (M \otimes_R S, M_{f_i}, \text{can}_i, \text{can}_{ij})
$$
其中 $\text{can}_i : (M \otimes_R S)_{f_i} \to M_{f_i} \otimes_R S$
与 $\text{can}_{ij} : (M_{f_i})_{f_j} \to (M_{f_j})_{f_i}$
是规范同构。对范畴
$\text{Glue}(R \to S, f_1, \ldots, f_t)$ 的任意对象
$\mathbf{M} = (M', M_i, \alpha_i, \alpha_{ij})$，定义
$$
H^0(\mathbf{M}) =
\{(m', m_i) \mid \alpha_i(m') = m_i \otimes 1, \alpha_{ij}(m_i) = m_j\}
$$
换言之，它由与 (\ref{more-algebra-equation-glueing-complex}) 类似的正合列
$$
0 \to H^0(\mathbf{M}) \to
M' \times \prod M_i \to
\prod M'_{f_i}
\times
\prod (M_i)_{f_j}
$$
定义。把 $H^0(\mathbf{M})$ 看作 $R$-模。这样也得到函子
$$
H^0 :
\text{Glue}(R \to S, f_1, \ldots, f_t)
\longrightarrow
\text{Mod}_R
$$
接下来的目标是说明：函子 $\text{Can}$ 与 $H^0$ 有时互为拟逆。
\end{remark}

\begin{lemma}
\label{more-algebra-lemma-H0-Can-adjoint}
在评注 \ref{more-algebra-remark-glueing-data} 中，函子
$H^0 : \text{Glue}(R \to S, f_1, \ldots, f_t) \to \text{Mod}_R$
是函子
$\text{Can} : \text{Mod}_R \to \text{Glue}(R \to S, f_1, \ldots, f_t)$
的右伴随。
\end{lemma}

\begin{proof}
设 $\mathbf{M} = (M', M_i, \alpha_i, \alpha_{ij})$ 是
$\text{Glue}(R \to S, f_1, \ldots, f_t)$ 的对象。对任意 $R$-模 $N$，
有映射
$$
\Hom_{\text{Glue}(R \to S, f_1, \ldots, f_t)}(\text{Can}(N), \mathbf{M})
\to
\Hom_R(N, H^0(\mathbf{M}))
$$
它把 $\psi$ 映到 $H^0(\psi)$ 与显然映射
$N \to H^0(\text{Can}(N))$ 的复合。由构造，所显示的映射对
$N = R$ 是同构（即使一般而言 $R \to H^0(\text{Can}(R))$
不是同构）。范畴 $\text{Glue}(R \to S, f_1, \ldots, f_t)$
有直和与余核。函子 $\text{Can}$ 与直和及余核交换。
将 $N$ 写成自由 $R$-模之间一个映射的余核，由这些观察可知
所显示的映射是双射。略去细节。
\end{proof}

\begin{lemma}
\label{more-algebra-lemma-H0-inverse}
假设 $\varphi : R \to S$ 为平坦环映射，且
$I = (f_1, \ldots, f_t) \subset R$ 是满足
$R/I \to S/IS$ 为同构的理想。则函子 $H^0$ 是评注
\ref{more-algebra-remark-glueing-data} 中函子 $\text{Can}$ 的左拟逆。
\end{lemma}

\begin{proof}
这是引理 \ref{more-algebra-lemma-recover-module-from-glueing-data} 的改述。
\end{proof}

\begin{lemma}
\label{more-algebra-lemma-exact}
假设 $\varphi : R \to S$ 为平坦环映射，并设
$I = (f_1, \ldots, f_t) \subset R$ 为理想。
则 $\text{Glue}(R \to S, f_1, \ldots, f_t)$ 是 Abel 范畴，
且函子 $\text{Can}$ 正合并与任意余极限交换。
\end{lemma}

\begin{proof}
给定范畴 $\text{Glue}(R \to S, f_1, \ldots, f_t)$ 中的态射
$$
(\varphi', \varphi_i) :
(M', M_i, \alpha_i, \alpha_{ij})
\to
(N', N_i, \beta_i, \beta_{ij})
$$
其核存在，且等于对象
$(\Ker(\varphi'), \Ker(\varphi_i), \alpha_i, \alpha_{ij})$；
其余核也存在，且等于对象
$(\Coker(\varphi'), \Coker(\varphi_i), \beta_i, \beta_{ij})$。
这是因为 $R \to S$ 平坦，故取核或余核与 $- \otimes_R S$ 交换。
略去细节。正合性来自 $R_{f_i}$ 与 $S$ 在 $R$ 上平坦；
与余极限交换则来自张量积与余极限交换。
\end{proof}

\begin{lemma}
\label{more-algebra-lemma-equivalence-I-unit}
设 $\varphi : R \to S$ 为平坦环映射，且
$(f_1, \ldots, f_t) = R$。则 $\text{Can}$ 与 $H^0$ 是互为拟逆的
范畴等价
$$
\text{Mod}_R = \text{Glue}(R \to S, f_1, \ldots, f_t)
$$
\end{lemma}

\begin{proof}
考虑 $\text{Glue}(R \to S, f_1, \ldots, f_t)$ 的对象
$\mathbf{M} = (M', M_i, \alpha_i, \alpha_{ij})$。由 Algebra，引理
\ref{algebra-lemma-glue-modules}，存在唯一的模 $M$ 及同构
$M_{f_i} \to M_i$，它们恢复粘合数据 $\alpha_{ij}$。
于是 $M'$ 与 $M \otimes_R S$ 都是 $S$-模，并且在各 $f_i$ 处局部化后
都恢复模 $M_i \otimes_R S$。因此有规范同构
$M \otimes_R S \to M'$。这说明 $\mathbf{M}$ 属于 $\text{Can}$ 的本质像。
结合引理 \ref{more-algebra-lemma-H0-inverse} 即得本引理。
\end{proof}

\begin{lemma}
\label{more-algebra-lemma-base-change-glue}
设 $\varphi : R \to S$ 为平坦环映射，并设
$I = (f_1, \ldots, f_t)$ 为理想。设 $R \to R'$ 为平坦环映射，并置
$S' = S \otimes_R R'$。
% P02-E0390
则得到如下范畴与函子的交换图：
$$
\xymatrix{
\text{Mod}_R \ar[r]_-{\text{Can}} \ar[d]_{-\otimes_R R'} &
\text{Glue}(R \to S, f_1, \ldots, f_t) \ar[r]_-{H^0} \ar[d]^{-\otimes_R R'} &
\text{Mod}_R \ar[d]^{-\otimes_R R'} \\
\text{Mod}_{R'} \ar[r]^-{\text{Can}} &
\text{Glue}(R' \to S', f_1, \ldots, f_t) \ar[r]^-{H^0} &
\text{Mod}_{R'}
}
$$
\end{lemma}

\begin{proof}
略。
\end{proof}

\begin{proposition}
\label{more-algebra-proposition-equivalence}
假设 $\varphi : R \to S$ 为平坦环映射，且
$I = (f_1, \ldots, f_t) \subset R$ 是满足
$R/I \to S/IS$ 为同构的理想。则 $\text{Can}$ 与 $H^0$
是互为拟逆的范畴等价
$$
\text{Mod}_R = \text{Glue}(R \to S, f_1, \ldots, f_t)
$$
\end{proposition}

\begin{proof}
我们已经看到 $H^0 \circ \text{Can}$ 同构于恒等函子；见引理
\ref{more-algebra-lemma-H0-inverse}。考虑
$\text{Glue}(R \to S, f_1, \ldots, f_t)$ 的对象
$\mathbf{M} = (M', M_i, \alpha_i, \alpha_{ij})$。
得到自然态射
$$
\Psi :
(H^0(\mathbf{M}) \otimes_R S, H^0(\mathbf{M})_{f_i},
\text{can}_i, \text{can}_{ij})
\longrightarrow
(M', M_i, \alpha_i, \alpha_{ij}).
$$
具体而言，按定义，$H^0(\mathbf{M})$ 带有相容的 $R$-模映射
$H^0(\mathbf{M}) \to M'$ 与 $H^0(\mathbf{M}) \to M_i$。
必须证明此映射是同构。

\medskip\noindent
选取指标 $i$ 并置 $R' = R_{f_i}$。结合引理
\ref{more-algebra-lemma-base-change-glue} 与 \ref{more-algebra-lemma-equivalence-I-unit}，
可知 $\Psi \otimes_R R'$ 是同构。因此 $\Psi$ 的核与余核分别是
形如 $(K, 0, 0, 0)$ 与 $(Q, 0, 0, 0)$ 的系统。注意
$H^0((K, 0, 0, 0)) = K$，函子 $H^0$ 左正合，并且由构造
$H^0(\Psi)$ 是双射。因此 $K = 0$，即 $\Psi$ 的核为零。

\medskip\noindent
上述结论给出短正合列
$$
0 \to H^0(\mathbf{M}) \otimes_R S \to M' \to Q \to 0
$$
并且 $M_i = H^0(\mathbf{M})_{f_i}$。可以把 $Q$ 看作 $I$-幂挠
$R$-模，因而 $Q = Q \otimes_R S$。由引理
\ref{more-algebra-lemma-neighbourhood-extensions}，存在交换图
$$
\xymatrix{
0 \ar[r] &
H^0(\mathbf{M}) \ar[r] \ar[d] &
E \ar[r] \ar[d] &
Q \ar[r] \ar[d] &
0 \\
0 \ar[r] &
H^0(\mathbf{M}) \otimes_R S \ar[r] &
M' \ar[r] &
Q \ar[r] &
0
}
$$
且两行正合。这显然在范畴
$\text{Glue}(R \to S, f_1, \ldots, f_t)$ 中确定同构
$\text{Can}(E) \to (M', M_i, \alpha_i, \alpha_{ij})$，结论即得。
（当然，事后可知 $Q = 0$。）
\end{proof}

\begin{lemma}
\label{more-algebra-lemma-application-formal-glueing}
设 $\varphi : R \to S$ 为平坦环映射，$I \subset R$ 为有限生成理想，
且 $R/I \to S/IS$ 是同构。
% P02-E0391
\begin{enumerate}
\item 给定 $R$-模 $N$、$S$-模 $M'$ 以及核和余核均为 $I$-幂挠的
$S$-模映射 $\varphi : M' \to N \otimes_R S$，则存在 $R$-模映射
$\psi : M \to N$ 及与 $\varphi$、$\psi$ 相容的同构
$M \otimes_R S = M'$。
\item 给定 $R$-模 $M$、$S$-模 $N'$ 以及核和余核均为 $I$-幂挠的
$S$-模映射 $\varphi : M \otimes_R S \to N'$，则存在 $R$-模映射
$\psi : M \to N$ 及与 $\varphi$、$\psi$ 相容的同构
$N \otimes_R S = N'$。
\end{enumerate}
在两种情形下均有 $\Ker(\varphi) \cong \Ker(\psi)$ 及
$\Coker(\varphi) \cong \Coker(\psi)$。
\end{lemma}

\begin{proof}
(1) 的证明。设 $I = (f_1, \ldots, f_t)$。显然，局部化
$\varphi_{f_i}$ 是同构。因此
$(M', N_{f_i}, \varphi_{f_i}, can_{ij})$ 是
$\text{Glue}(R \to S, f_1, \ldots, f_t)$ 的对象；见评注
\ref{more-algebra-remark-glueing-data}。
% P02-E0392
由命题 \ref{more-algebra-proposition-equivalence}，存在 $R$-模 $M$，使得
$M' = M \otimes_R S$ 及 $N_{f_i} = M_{f_i}$ 与同构
$\varphi_{f_i}$、$can_{ij}$ 相容。有
$$
(M \otimes_R S, M_{f_i}, can_i, can_{ij}) =
(M', N_{f_i}, \varphi_{f_i}, can_{ij})
\to
(N \otimes_R S, N_{f_i}, can_i, can_{ij})
$$
这一 $\text{Glue}(R \to S, f_1, \ldots, f_t)$ 中的态射，
其第一个分量使用 $\varphi$。由命题
\ref{more-algebra-proposition-equivalence} 的范畴等价，它对应于 $R$-模映射
$\psi : M \to N$。$M \to N$ 的基变换与同构
$M' \cong M \otimes_R S$ 的复合是 $\varphi$；换言之，
$M \to N$ 与 $\varphi$ 相容。

\medskip\noindent
(2) 的证明。这只是上述论证的对偶。具体而言，局部化
$\varphi_{f_i}$ 是同构。因此
$(N', M_{f_i}, \varphi_{f_i}^{-1}, can_{ij})$ 是
$\text{Glue}(R \to S, f_1, \ldots, f_t)$ 的对象；见评注
\ref{more-algebra-remark-glueing-data}。由命题
\ref{more-algebra-proposition-equivalence}，存在 $R$-模 $N$，使得
$N' = N \otimes_R S$ 及 $N_{f_i} = M_{f_i}$ 与同构
$\varphi_{f_i}^{-1}$、$can_{ij}$ 相容。有态射
$$
(M \otimes_R S, M_{f_i}, can_i, can_{ij}) \to
(N', M_{f_i}, \varphi_{f_i}, can_{ij}) =
(N \otimes_R S, N_{f_i}, can_i, can_{ij})
$$
% P02-E0393
这一 $\text{Glue}(R \to S, f_1, \ldots, f_t)$ 中的态射，
其第一个分量使用 $\varphi$。由命题
\ref{more-algebra-proposition-equivalence} 的范畴等价，它对应于 $R$-模映射
$\psi : M \to N$。$M \to N$ 的基变换与同构
$N' \cong N \otimes_R S$ 的复合是 $\varphi$；换言之，
$M \to N$ 与 $\varphi$ 相容。

\medskip\noindent
最后一个陈述例如可由引理
\ref{more-algebra-lemma-neighbourhood-equivalence} 得到。
\end{proof}

\noindent
接下来将命题 \ref{more-algebra-proposition-equivalence} 特化，以得到更便于使用的形式。
具体而言，若 $I = (f)$ 是主理想，则
$\text{Glue}(R \to S, f)$ 的对象不过是三元组
$(M', M_1, \alpha_1)$，而且{\it 无需}检验上链条件！

\begin{theorem}
\label{more-algebra-theorem-formal-glueing}
设 $R$ 为环，$f \in R$。设 $\varphi : R \to S$ 为平坦环映射，
并诱导同构 $R/fR \to S/fS$。则函子
$$
\text{Mod}_R
\longrightarrow
\text{Mod}_S \times_{\text{Mod}_{S_f}} \text{Mod}_{R_f},
\quad
M
\longmapsto
(M \otimes_R S, M_f, \text{can})
$$
是等价。
\end{theorem}

\begin{proof}
箭头右侧的范畴是三元组 $(M', M_1, \alpha_1)$ 组成的范畴，其中
$M'$ 是 $S$-模，$M_1$ 是 $R_f$-模，并且
$\alpha_1 : M'_f \to M_1 \otimes_R S$ 是 $S_f$-同构；见
Categories，例 \ref{categories-example-2-fibre-product-categories}。
% P02-E0394
故该定理是命题 \ref{more-algebra-proposition-equivalence} 的特例。
\end{proof}

\noindent
定理 \ref{more-algebra-theorem-formal-glueing} 的一个有用特例是：$R$ 为 Noether 环，
而 $S$ 是 $R$ 关于元素 $f$ 的完备化。完备化映射 $R \to S$ 平坦；
当 $M$ 有限生成时，函子 $M \mapsto M \otimes_R S$ 可等同于
$f$-进完备化函子。为更精确地陈述这一点，用
$\text{Mod}^{fg}_R$ 表示有限生成 $R$-模的范畴。

\begin{proposition}
\label{more-algebra-proposition-formal-glueing}
设 $R$ 为 Noether 环，$f \in R$ 为元素，并设 $R^\wedge$ 为
$R$ 的 $f$-进完备化。则函子
$M \mapsto (M^\wedge, M_f, \text{can})$ 定义等价
$$
\text{Mod}^{fg}_R
\longrightarrow
\text{Mod}^{fg}_{R^\wedge}
\times_{\text{Mod}^{fg}_{(R^\wedge)_f}}
\text{Mod}^{fg}_{R_f}
$$
\end{proposition}

\begin{proof}
由 Algebra，引理 \ref{algebra-lemma-completion-flat}，环映射
$R \to R^\wedge$ 平坦。显然 $R/fR = R^\wedge/fR^\wedge$。
由 Algebra，引理 \ref{algebra-lemma-completion-tensor}，有限
$R$-模 $M$ 的完备化等于 $M \otimes_R R^\wedge$。
故命题所显示的函子等于定理 \ref{more-algebra-theorem-formal-glueing} 中的函子，
特别地它全忠实。设 $(M_1, M_2, \psi)$ 为右侧的对象。
由定理 \ref{more-algebra-theorem-formal-glueing}，存在 $R$-模 $M$，使得
$M_1 = M \otimes_R R^\wedge$ 及 $M_2 = M_f$。
由于 $R \to R^\wedge \times R_f$ 忠实平坦，由 Algebra，引理
\ref{algebra-lemma-cover} 可知 $M$ 有限生成，即
$M \in \text{Mod}^{fg}_R$。命题得证。
\end{proof}

\begin{remark}
\label{more-algebra-remark-formal-glueing-algebras}
命题 \ref{more-algebra-proposition-equivalence}、定理
\ref{more-algebra-theorem-formal-glueing} 及命题
\ref{more-algebra-proposition-formal-glueing} 的等价保持模的性质。例如，若
$M$ 对应于 $\mathbf{M} = (M', M_i, \alpha_i, \alpha_{ij})$，
则 $M$ 在 $R$ 上有限、有限表现、平坦或投射，当且仅当
$M'$ 与 $M_i$ 分别在 $S$ 与 $R_{f_i}$ 上具有相应性质。
这是因为 $R \to S \times \prod R_{f_i}$ 忠实平坦，并且这些性质
可沿忠实平坦映射下降与上升；见 Algebra，引理
\ref{algebra-lemma-descend-properties-modules} 及定理
\ref{algebra-theorem-ffdescent-projectivity}。
% P02-E0395
这些函子还保持两侧的 $\otimes$-结构。因此，它定义由
% P02-E0396
偶 $(\text{Mod}_R, \otimes)$ 构造的各类范畴之间的等价，
例如代数范畴之间的等价。
\end{remark}

\begin{remark}
\label{more-algebra-remark-topological-analogue}
给定微分流形 $X$ 及紧闭子流形 $Z$，其补集为 $U$，指定 $X$ 上的层
等价于指定 $U$ 上的层、$Z$ 在 $X$ 中某个未指定的管状邻域 $T$ 上的层，
以及这两个层在 $T \cap U$ 上所得层之间的同构。
% P02-E0397
代数几何中并不存在这种意义下的管状邻域，但命题
\ref{more-algebra-proposition-equivalence}、定理
\ref{more-algebra-theorem-formal-glueing} 及命题
\ref{more-algebra-proposition-formal-glueing} 等结果允许我们改用形式邻域。
\end{remark}

\section{Beauville--Laszlo 定理}
\label{more-algebra-section-beauville-laszlo}

\noindent
设 $R$ 为环，$f$ 为 $R$ 的元素。用
$R^\wedge = \lim R/f^n R$ 表示 $R$ 的 $f$-进完备化。
% P02-E0398
本节讨论并稍作推广 Beauville 与 Laszlo 的一个定理；见
\cite{Beauville-Laszlo}。该定理断言：在适当条件下，可以把
$R^\wedge$ 上的模与 $R_f$ 上的模，沿着它们基扩张到
$(R^\wedge)_f$ 后所得模之间的一个同构“粘合起来”，从而构造
$R$ 上的模。

\medskip\noindent
\cite{Beauville-Laszlo} 假设 $f$ 在 $R$ 与 $M$ 上都是非零因子。
事实上，只需假设
$$
R[f^\infty] \longrightarrow R^\wedge[f^\infty]
$$
是双射，并且
$$
M[f^\infty] \longrightarrow M \otimes_R R^\wedge
$$
是单射。这一优化部分受到 \cite[\S 1.3]{Kedlaya-Liu-I}
为非 Archimedes 解析空间理论而引入的另一种粘合方法的启发。

\medskip\noindent
事实上，我们将在更一般的环映射
$$
R \longrightarrow R'
$$
的框架下建立 Beauville--Laszlo 定理；它对每个 $n > 0$ 都诱导同构
$R/f^nR \to R'/f^nR'$，并诱导同构
$R[f^\infty] \to R'[f^\infty]$。这个框架更适合整体化，而且不能从
$R'$ 是 $R$ 的完备化的情形形式地推出。例如，条件
$R[f^\infty] \to R'[f^\infty]$ 为双射并不蕴含
$R[f^\infty] \to R^\wedge[f^\infty]$ 为双射。

\medskip\noindent
本节所证明的 Beauville 与 Laszlo 定理可看作定理
\ref{more-algebra-theorem-formal-glueing} 的非平坦版本；当
$R' = R^\wedge$ 时，也可看作命题
\ref{more-algebra-proposition-formal-glueing} 的非 Noether 版本。
它与平坦下降的比较见评注 \ref{more-algebra-remark-not-descent}。

\medskip\noindent
还可以建立更强的结果（例如不对 $M$ 施加限制），但为此必须在
导出范畴层面工作。详见 \cite[\S 5]{Bhatt-Algebraize}。

\begin{lemma}
\label{more-algebra-lemma-same-quotients}
设 $R$ 为环，$f \in R$。对每个正整数 $n$，映射
$R/f^nR \to R^\wedge/f^n R^\wedge$ 是同构。
\end{lemma}

\begin{proof}
这是 Algebra，引理
\ref{algebra-lemma-hathat-finitely-generated} 的特例。
\end{proof}

\noindent
我们将使用第 \ref{more-algebra-section-torsion} 节引入的记号。因此，对
$R$-模 $M$，用 $M[f^n]$ 表示 $M$ 中被 $f^n$ 零化的子模，并置
% P02-E0399
$$
M[f^\infty] = \bigcup\nolimits_{n = 1}^\infty M[f^n] = \Ker(M \to M_f).
$$
若 $M = M[f^\infty]$，则称 $M$ 为 $f$-幂挠模。

\begin{lemma}
\label{more-algebra-lemma-BL-faithful}
设 $R$ 为环，$f \in R$，且 $R \to R'$ 为对 $n > 0$ 诱导同构
$R/f^nR \to R'/f^nR'$ 的环映射。$R$-模 $R' \oplus R_f$ 是忠实的：
对每个非零 $R$-模 $M$，模 $M \otimes_R (R' \oplus R_f)$ 也非零。
例如，若 $M$ 非零，则 $M \otimes_R (R^\wedge \oplus R_f)$ 非零。
\end{lemma}

\noindent
然而，映射 $M \to M \otimes_R (R' \oplus R_f)$ 未必是单射；
见例 \ref{more-algebra-example-not-glueing-pair}。

\begin{proof}
若 $M \neq 0$ 但 $M \otimes_R R_f = 0$，则 $M$ 是 $f$-幂挠的。
由引理 \ref{more-algebra-lemma-torsion-completion}，
$M \otimes_R R' \cong M \neq 0$。由引理
\ref{more-algebra-lemma-same-quotients}，最后一个陈述是第一个陈述的特例。
\end{proof}

\begin{lemma}
\label{more-algebra-lemma-cover-spec}
设 $R$ 为环，$f \in R$，且 $R \to R'$ 为诱导同构
$R/fR \to R'/fR'$ 的环映射。映射
$\Spec(R') \amalg \Spec(R_f) \to \Spec(R)$ 是满射。
例如，映射
$\Spec(R^\wedge) \amalg \Spec(R_f) \to \Spec(R)$ 是满射。
\end{lemma}

\begin{proof}
回忆 $\Spec(R) = V(f) \amalg D(f)$，其中
$V(f) = \Spec(R/fR)$ 且 $D(f) = \Spec(R_f)$；见
Algebra，第 \ref{algebra-section-spectrum-ring} 节，尤其是引理
\ref{algebra-lemma-spec-closed} 与
\ref{algebra-lemma-standard-open}。由于映射 $R \to R/fR$
经由 $R'$ 分解，引理成立。由引理
\ref{more-algebra-lemma-same-quotients}，最后一个陈述是第一个陈述的特例。
\end{proof}

\begin{lemma}
\label{more-algebra-lemma-faithful-descent}
\begin{reference}
\cite[Lemme~2(a)]{Beauville-Laszlo} 的轻微推广。
\end{reference}
设 $R$ 为环，$f \in R$，且 $R \to R'$ 为对 $n > 0$ 诱导同构
$R/f^nR \to R'/f^nR'$ 的环映射。$R$-模 $M$ 有限生成，当且仅当
（$R' \oplus R_f$）-模 $M \otimes_R (R' \oplus R_f)$ 有限生成。
例如，若 $M \otimes_R (R^\wedge \oplus R_f)$ 作为
$R^\wedge \oplus R_f$ 上的模有限生成，则 $M$ 是有限生成 $R$-模。
\end{lemma}

\begin{proof}
“仅当”方向是清楚的，故假设
$M \otimes_R (R' \oplus R_f)$ 有限生成。将每个生成元写成纯张量之和，
可知 $M \otimes_R (R' \oplus R_f)$ 有一个由 $M$ 中元素组成的
有限生成集。换言之，存在从有限自由 $R$-模到 $M$ 的态射，其余核
在与 $R' \oplus R_f$ 作张量后被零化；对该余核应用引理
\ref{more-algebra-lemma-BL-faithful}，可推出 $M$ 有限生成。
% P02-E0400
由引理 \ref{more-algebra-lemma-same-quotients}，最后一个陈述是第一个陈述的特例。
\end{proof}

\begin{remark}
\label{more-algebra-remark-not-descent}
虽然 $R \to R_f$ 总是平坦的，但除非 $R$ 是 Noether 环，
$R \to R^\wedge$ 通常并不平坦（见 Algebra，引理
\ref{algebra-lemma-completion-flat} 以及 Examples，第
\ref{examples-section-nonflat} 节中的讨论）。因此，一般不能把
Descent，第 \ref{descent-section-descent-modules} 节所讨论的忠实平坦下降
应用于态射 $R \to R^\wedge \oplus R_f$。此外，即使在 Noether 情形，
该态射的通常下降数据定义也涉及环
$R^\wedge \otimes_R R^\wedge$，而本节将避免考虑此环。
\end{remark}

\noindent
{\bf 粘合对。}设 $R \to R'$ 为对 $n > 0$ 诱导同构
$R/f^nR \to R'/f^nR'$ 的环映射。考虑序列
\begin{equation}
\label{more-algebra-equation-BL-cech-re}
0 \to R \to R' \oplus R_f \to R'_f \to 0,
\end{equation}
其中右侧映射是两个规范同态之差。若此序列正合，则称
$(R \to R', f)$ 为一个\emph{粘合对}。若
$(R \to R^\wedge, f)$ 是粘合对，则称 $(R, f)$ 为
\emph{粘合对}；由引理 \ref{more-algebra-lemma-same-quotients}，这一定义有意义。
因此，$(R, f)$ 是粘合对，当且仅当序列
\begin{equation}
\label{more-algebra-equation-BL-cech}
0 \to R \to R^\wedge \oplus R_f \to (R^\wedge)_f \to 0,
\end{equation}
正合。

\begin{lemma}
\label{more-algebra-lemma-same-f-torsion}
设 $R$ 为环，$f \in R$，且 $R \to R'$ 为对 $n > 0$ 诱导同构
$R/f^nR \to R'/f^nR'$ 的环映射。序列
(\ref{more-algebra-equation-BL-cech-re})：
\begin{enumerate}
\item 在右端正合；
\item 在左端正合，当且仅当
$R[f^\infty] \to R'[f^\infty]$ 是单射；
\item 在中间正合，当且仅当
$R[f^\infty] \to R'[f^\infty]$ 是满射。
\end{enumerate}
特别地，$(R \to R', f)$ 是粘合对，当且仅当
$R[f^\infty] \to R'[f^\infty]$ 是双射。例如，$(R, f)$ 是粘合对，
当且仅当 $R[f^\infty] \to R^\wedge[f^\infty]$ 是双射。
\end{lemma}

\begin{proof}
设 $x \in R'_f$。写成 $x = x'/f^n$，其中 $x' \in R'$。
再写成 $x' = x'' + f^n y$，其中 $x'' \in R$ 且 $y \in R'$。
于是 $(y, -x''/f^n)$ 映到 $x$，故 (1) 成立。

\medskip\noindent
(2) 来自等式 $\Ker(R \to R_f) = R[f^\infty]$。

\medskip\noindent
若序列在中间正合，则形如 $(x, 0)$、其中
$x \in R'[f^\infty]$ 的元素都在第一个箭头的像中。
这蕴含 $R[f^\infty] \to R'[f^\infty]$ 是满射。
反之，假设 $R[f^\infty] \to R'[f^\infty]$ 满射。设
$(x, y)$ 是映到右端零的中间元素。对某个 $y' \in R$，写成
$y = y'/f^n$。于是 $f^n x - y'$ 在 $R'$ 中被 $f$ 的某个幂零化。
由假设，可对某个 $z \in R[f^\infty]$ 写成
$f^nx - y' = z$。于是 $y = y''/f^n$，其中
$y'' = y' + z$ 属于 $R \to R/f^nR$ 的核。因此对某个
$y''' \in R$，$y$ 可表示为 $y'''/1$。此时
$x - y'''$ 属于 $R'[f^\infty]$，故
$x - y''' = z' \in R[f^\infty]$。于是
$(x, y'''/1) = (y''' + z', (y''' + z')/1)$，如所需。

\medskip\noindent
由引理 \ref{more-algebra-lemma-same-quotients}，引理的最后一个陈述是倒数第二个
陈述的特例。
\end{proof}

\begin{remark}
\label{more-algebra-remark-BL-special-case}
假设 $f$ 是非零因子。则 Algebra，引理
\ref{algebra-lemma-completion-differ-by-torsion} 说明
$f$ 在 $R^\wedge$ 中是非零因子。故 $(R, f)$ 是粘合对。
\end{remark}

\begin{remark}
\label{more-algebra-remark-noetherian-case}
若 $R \to R^\wedge$ 平坦，则对每个正整数 $n$，将序列
$0 \to R[f^n] \to R \to R$ 与 $R^\wedge$ 作张量，得到序列
$0 \to R[f^n] \otimes_R R^\wedge \to R^\wedge \to R^\wedge$。
结合引理 \ref{more-algebra-lemma-torsion-completion}，可得
$R[f^n] \to R^\wedge[f^n]$ 是同构。故 $(R, f)$ 是粘合对。
特别地，当 $R$ 为 Noether 环时这成立；见 Algebra，引理
\ref{algebra-lemma-completion-flat}。
\end{remark}

\begin{example}
\label{more-algebra-example-not-glueing-pair}
设 $k$ 为域，并置
$$
R = k[f, T_1, T_2, \ldots]/(fT_1, fT_2 - T_1, fT_3 - T_2, \ldots).
$$
则 $(R, f)$ 不是粘合对，因为映射
$R[f^\infty] \to R^\wedge[f^\infty]$ 不是单射：$T_1$ 的像在
$R^\wedge$ 中可被 $f$ 整除。对于
$$
R = k[f, T_1, T_2, \ldots]/(fT_1, f^2T_2, \ldots),
$$
映射 $R[f^\infty] \to R^\wedge[f^\infty]$ 不是满射，因为元素
$T_1 + fT_2 + f^2 T_3 + \ldots$ 不在像中。特别地，由评注
\ref{more-algebra-remark-noetherian-case}，这二者都是
$R \to R^\wedge$ 不平坦的例子。
\end{example}

\noindent
{\bf 可粘合模。}
设 $R \to R'$ 为对 $n > 0$ 诱导同构
$R/f^nR \to R'/f^nR'$ 的环映射。对任意 $R$-模 $M$，可将
(\ref{more-algebra-equation-BL-cech-re}) 与 $M$ 作张量，得到序列
\begin{equation}
\label{more-algebra-equation-BL-cech-mod-re}
0 \to M \to (M \otimes_R R') \oplus (M \otimes_R R_f) \to
M \otimes_R R'_f \to 0
\end{equation}
注意 $M \otimes_R R_f = M_f$，并且
$M \otimes_R R'_f = (M \otimes_R R')_f$。
若此序列正合，则称 $M$ \emph{对 $(R \to R', f)$ 可粘合}。
若 $R$ 为环且 $f \in R$，并且 $M$ 对
$(R \to R^\wedge, f)$ 可粘合，则称 $R$-模\emph{可粘合}。
% P02-E0401
因此，$M$ 可粘合，当且仅当序列
\begin{equation}
\label{more-algebra-equation-BL-cech-mod}
0 \to M \to (M \otimes_R R^\wedge) \oplus (M \otimes_R R_f) \to
M \otimes_R (R^\wedge)_f \to 0
\end{equation}
正合。

\begin{lemma}
\label{more-algebra-lemma-same-f-torsion-module}
设 $R$ 为环，$f \in R$，且 $R \to R'$ 为对 $n > 0$ 诱导同构
$R/f^nR \to R'/f^nR'$ 的环映射。序列
(\ref{more-algebra-equation-BL-cech-mod-re})：
\begin{enumerate}
\item 在右端正合；
\item 在左端正合，当且仅当
$M[f^\infty] \to (M \otimes_R R')[f^\infty]$ 是单射；
\item 在中间正合，当且仅当
$M[f^\infty] \to (M \otimes_R R')[f^\infty]$ 是满射。
\end{enumerate}
因此，$M$ 对 $(R \to R', f)$ 可粘合，当且仅当
$M[f^\infty] \to (M \otimes_R R')[f^\infty]$ 是双射。
若 $(R \to R', f)$ 是粘合对，则 $M$ 对 $(R \to R', f)$ 可粘合，
当且仅当 $M[f^\infty] \to (M \otimes_R R')[f^\infty]$ 是单射。
例如，若 $(R, f)$ 是粘合对，则 $M$ 可粘合，当且仅当
$M[f^\infty] \to (M \otimes_R R^\wedge)[f^\infty]$ 是单射。
\end{lemma}

\begin{proof}
下文将不再说明而使用引理 \ref{more-algebra-lemma-same-f-torsion} 的结果。
函子 $M \otimes_R -$ 右正合
（Algebra，引理 \ref{algebra-lemma-tensor-product-exact}），故得到 (1)。

\medskip\noindent
$M \to M \otimes_R R_f = M_f$ 的核为 $M[f^\infty]$，故 (2) 成立。

\medskip\noindent
若序列在中间正合，则形如 $(x, 0)$、其中
$x \in (M \otimes_R R')[f^\infty]$ 的元素都在第一个箭头的像中。
这蕴含
$M[f^\infty] \to (M \otimes_R R')[f^\infty]$ 是满射。
反之，假设
$M[f^\infty] \to (M \otimes_R R')[f^\infty]$ 是满射。
设 $(x, y)$ 为映到右端零的中间元素。
对某个 $y' \in M$，写成 $y = y'/f^n$。于是
$f^n x - y'$ 在 $M \otimes_R R'$ 中被 $f$ 的某个幂零化。
由假设，可对某个 $z \in M[f^\infty]$ 写成
$f^nx - y' = z$。于是 $y = y''/f^n$，其中
$y'' = y' + z$ 属于 $M \to M/f^nM$ 的核。因此对某个
$y''' \in M$，$y$ 可表示为 $y'''/1$。此时
$x - y'''$ 属于 $(M \otimes_R R')[f^\infty]$，故
$x - y''' = z' \in M[f^\infty]$。于是
$(x, y'''/1) = (y''' + z', (y''' + z')/1)$，如所需。

\medskip\noindent
若 $(R \to R', f)$ 是粘合对，则对任意 $M$，序列
(\ref{more-algebra-equation-BL-cech-mod-re}) 由 Algebra，引理
\ref{algebra-lemma-tensor-product-exact} 在中间正合。
这给出引理的倒数第二个陈述。最后一个陈述由此以及如下事实得到：
$(R, f)$ 是粘合对，当且仅当
$(R \to R^\wedge, f)$ 是粘合对。
\end{proof}

\begin{remark}
\label{more-algebra-remark-glueable}
设 $(R \to R', f)$ 为粘合对，$M$ 为 $R$-模。以下观察可用来判定
$M$ 是否对 $(R \to R', f)$ 可粘合。
\begin{enumerate}
\item 由引理 \ref{more-algebra-lemma-same-f-torsion-module}，$M$ 对
$(R \to R^\wedge, f)$ 可粘合，当且仅当
$M[f^\infty] \to M \otimes_R R^\wedge$ 是单射。
若 $M[f] \to M^\wedge$ 是单射，即
$M[f] \cap \bigcap_{n = 1}^\infty f^n M = 0$，则该条件成立。
\item 若 $\text{Tor}_1^R(M, R'_f) = 0$，则 $M$ 对
$(R \to R', f)$ 可粘合
（使用 Algebra，引理 \ref{algebra-lemma-long-exact-sequence-tor}）。
这等价于 $\text{Tor}_1^R(M, R')$ 是 $f$-幂挠的。
特别地，任意平坦 $R$-模都对 $(R \to R', f)$ 可粘合。
\item 若 $R \to R'$ 平坦，则对每个 $R$-模都有
$\text{Tor}_1^R(M, R') = 0$，故每个 $R$-模都对
$(R \to R', f)$ 可粘合。特别地，当 $R$ 为 Noether 环且
$R' = R^\wedge$ 时，该条件成立；见 Algebra，引理
\ref{algebra-lemma-completion-flat}
% P02-E0402
\end{enumerate}
\end{remark}

\begin{example}[不可粘合模]
\label{more-algebra-example-not-glueable-module}
% P02-E0403
\begin{reference}
\cite[\S 4, Remarques]{Beauville-Laszlo}
\end{reference}
设 $R$ 为 $\mathbf{R}$ 上 $0$ 处的 $C^\infty$ 函数芽之环。
设 $f \in R$ 为函数 $f(x) = x$。则 $f$ 在 $R$ 中是非零因子，
故 $(R, f)$ 是粘合对，且 $R^\wedge \cong \mathbf{R}[[x]]$。
设 $\varphi \in R$ 为函数
$\varphi(x) = \text{exp}(-1/x^2)$。
% P02-E0404
则 $\varphi$ 的 Taylor 级数为零，故
$\varphi \in \Ker(R \to R^\wedge)$。由于当 $x \neq 0$ 时
$\varphi(x) \neq 0$，可知 $\varphi$ 在 $R$ 中是非零因子。
函数 $\varphi/f$ 的 Taylor 级数也为零，故它在
$M = R/\varphi R$ 中的像是 $M[f]$ 的非零元素，而该元素映到
$$
M \otimes_R R^\wedge = R^\wedge/\varphi R^\wedge = R^\wedge
$$
中的零。因此 $M$ 不可粘合。
\end{example}

\noindent
接下来计算一些 Tor 群。

\begin{lemma}
\label{more-algebra-lemma-first-tor}
设 $(R \to R', f)$ 为粘合对。则对每个 $n > 0$，
$\text{Tor}^R_1(R', f^n R) = 0$。
\end{lemma}

\begin{proof}
由正合列 $0 \to R[f^n] \to R \to f^n R \to 0$，只需检验
$R[f^n] \otimes_R R' \to R'$ 是单射。由引理
\ref{more-algebra-lemma-torsion-completion}，有
$R[f^n] \otimes_R R' = R[f^n]$；又由引理
\ref{more-algebra-lemma-same-f-torsion} 及 $(R \to R', f)$ 是粘合对，
$R[f^n] \to R'$ 是单射。
\end{proof}

\begin{lemma}
\label{more-algebra-lemma-first-tor-total}
设 $(R \to R',f)$ 为粘合对。
则 $\text{Tor}^R_1(R', R/R[f^\infty]) = 0$。
\end{lemma}

\begin{proof}
有 $R/R[f^\infty] = \colim R/R[f^n] = \colim f^nR$。
由于 Tor 群的形成与滤过余极限交换
（Algebra，引理 \ref{algebra-lemma-tor-commutes-filtered-colimits}），
可应用引理 \ref{more-algebra-lemma-first-tor}。
\end{proof}

\begin{lemma}
\label{more-algebra-lemma-BL3}
\begin{reference}
\cite[Lemme 3(a)]{Beauville-Laszlo} 的轻微推广。
\end{reference}
设 $(R \to R', f)$ 为粘合对。对每个 $R$-模 $M$，有
$\text{Tor}^R_1(R', \Coker(M \to M_f)) = 0$。
\end{lemma}

\begin{proof}
置 $\overline{M} = M/M[f^\infty]$。则
$\Coker(M \to M_f) \cong \Coker(\overline{M} \to \overline{M}_f)$，
故不妨假设 $f$ 在 $M$ 上是非零因子。此时 $M \subset M_f$，且
$M_f/M = \colim M/f^nM$，其中转移映射由乘以 $f$ 给出。
由于 Tor 群的形成与余极限交换
（Algebra，引理 \ref{algebra-lemma-tor-commutes-filtered-colimits}），
只需证明 $\text{Tor}^R_1(R', M/f^n M) = 0$。

\medskip\noindent
先处理 $M = R/R[f^\infty]$ 的情形。由引理
\ref{more-algebra-lemma-same-f-torsion}，有
$M \otimes_R R' = R'/R'[f^\infty]$。从短正合列
$0 \to M \to M \to M/f^nM \to 0$，由 Algebra，引理
\ref{algebra-lemma-long-exact-sequence-tor} 得到正合列
$$
\xymatrix{
\text{Tor}_1^R(R', R/R[f^\infty]) \ar[r] &
\text{Tor}_1^R(R', M/f^n M) \ar[r] &
R'/R'[f^\infty] \ar[dll]_{f^n} \\
R'/R'[f^\infty] \ar[r] &
(R'/R'[f^\infty])/(f^n
(R'/R'[f^\infty])) \ar[r] & 0
}
$$
这里对角箭头是单射。由引理 \ref{more-algebra-lemma-first-tor-total}，第一个群
$\text{Tor}_1^R(R', R/R[f^\infty])$ 为零，故推出
$\text{Tor}_1^R(R', M/f^nM) = 0$，如所需。

\medskip\noindent
为处理一般情形，选取满射 $F \to M$，其中 $F$ 是自由
$R/R[f^\infty]$-模，并构造正合列
$$
0 \to N \to F/f^n F \to M/f^n M \to 0.
$$
由引理 \ref{more-algebra-lemma-torsion-completion}，与 $R'$ 作张量后此序列不变，
因而仍正合。由上一段，
$\text{Tor}^R_1(R', F/f^n F) = 0$，故推出
$\text{Tor}^R_1(R', M/f^n M) = 0$，如所需。
\end{proof}

\noindent
设 $(R \to R', f)$ 为粘合对。这意味着对 $n > 0$，
$R/f^nR \to R'/f^nR'$ 是同构，且序列
$$
0 \to R \to R' \oplus R_f \to R_f' \to 0
$$
正合。考虑评注 \ref{more-algebra-remark-glueing-data} 中引入的范畴
$\text{Glue}(R \to R', f)$。我们称
$\text{Glue}(R \to R', f)$ 的对象 $(M', M_1, \alpha_1)$ 为一个
\emph{粘合数据}。它由一个 $R'$-模 $M'$、一个 $R_f$-模 $M_1$
及同构 $\alpha_1 : (M')_f \to M_1 \otimes_R R'$ 组成。
有显然的函子
$$
\text{Can} : \text{Mod}_R \longrightarrow \text{Glue}(R \to R', f),\quad
M \longmapsto (M \otimes_R R', M_f, \text{can}),
$$
并且反方向有函子
$$
H^0 : \text{Glue}(R \to R', f) \longrightarrow \text{Mod}_R,\quad
(M', M_1, \alpha_1) \longmapsto \Ker(M' \oplus M_1 \to (M')_f)
$$
其精确定义见评注 \ref{more-algebra-remark-glueing-data}。

\begin{theorem}
\label{more-algebra-theorem-BL-glueing}
\begin{reference}
\cite{Beauville-Laszlo} 主定理的轻微推广。
\end{reference}
设 $(R \to R',f)$ 为粘合对。函子
$\text{Can} : \text{Mod}_R \longrightarrow \text{Glue}(R \to R', f)$
确定如下两个范畴之间的等价：对 $(R \to R', f)$ 可粘合的
$R$-模范畴，与粘合数据的范畴
$\text{Glue}(R \to R', f)$。
% P02-E0407
\end{theorem}

\begin{proof}
设 $(M', M_1, \alpha_1)$ 为粘合数据。我们将证明
$M = H^0((M', M_1, \alpha_1))$ 是对 $(R \to R', f)$ 可粘合的模，
且 $(M', M_1, \alpha_1) \cong \text{Can}(M)$。

\medskip\noindent
先检验函子 $H^0$ 定义中所用的映射
$\text{d} : M' \oplus M_1 \to (M')_f$ 是满射。注意，
$(x, y) \in M' \oplus M_1$ 映到 $(M')_f$ 中的
$\text{d}(x, y) = x/1 - \alpha_1^{-1}(y \otimes 1)$。
若 $z \in (M')_f$，则可写成
$\alpha_1(z) = \sum y_i \otimes g_i$，其中 $g_i \in R'$ 且
$y_i \in M_1$。对某个 $y'_i \in M'$ 与 $n \geq 0$，写成
$\alpha_1^{-1}(y_i \otimes 1) = y_i'/f^n$
（可以对所有 $i$ 选取同一个 $n$）。再写成
$g_i = a_i + f^n b_i$，其中 $a_i \in R$ 且 $b_i \in R'$。
于是取 $y = \sum a_i y_i \in M_1$ 及
$x = \sum b_i y'_i \in M'$，即有
$\text{d}(x, -y) = z$，如所需。

\medskip\noindent
由于 $M = H^0((M', M_1, \alpha_1)) = \Ker(\text{d})$，
得到 $R$-模的正合列
\begin{equation}
\label{more-algebra-equation-define-M}
0 \to M \to M' \oplus M_1 \to (M')_f \to 0.
\end{equation}
我们将证明映射 $M \to M'$ 与 $M \to M_1$ 诱导同构
$M \otimes_R R' \to M'$ 与 $M \otimes_R R_f \to M_1$。
这将蕴含 $M$ 对 $(R \to R', f)$ 可粘合，且
$\text{Can}(M) \cong (M', M_1, \alpha_1)$，如所需。

\medskip\noindent
由于 $f$ 在 $M_1$ 上是非零因子，有
$M[f^\infty] \cong M'[f^\infty]$。这给出正合列
\begin{equation}
\label{more-algebra-equation-exact-mod-torsion}
0 \to M/M[f^\infty] \to M_1 \to (M')_f/M' \to 0.
\end{equation}
由于 $R \to R_f$ 平坦，可将该正合列与 $R_f$ 作张量，推出
$M \otimes_R R_f = (M/M[f^\infty]) \otimes_R R_f \to M_1$
是同构。

\medskip\noindent
由引理 \ref{more-algebra-lemma-BL3}，有
$\text{Tor}_1^R(R', \Coker(M' \to (M')_f)) = 0$。
因此，序列 (\ref{more-algebra-equation-exact-mod-torsion}) 与 $R'$ 在 $R$ 上
作张量后仍正合。利用 $\alpha_1$ 与引理
\ref{more-algebra-lemma-torsion-completion}，所得正合列可写成
\begin{equation}
\label{more-algebra-equation-mod-torsion-sequence}
0 \to (M/M[f^\infty]) \otimes_R R' \to
(M')_f \to (M')_f/M' \to 0
\end{equation}
这给出同构
$(M/M[f^\infty]) \otimes_R R' \cong M'/M'[f^\infty]$。
由此，在图
$$
\xymatrix{
& M[f^\infty] \otimes_R R' \ar[r] \ar[d] &
M \otimes_R R'  \ar[r] \ar[d] &
(M/M[f^\infty]) \otimes_R R' \ar[r] \ar[d] & 0 \\
0 \ar[r] &
M'[f^\infty] \ar[r] &
M' \ar[r] &
M'/M'[f^\infty] \ar[r] & 0,
}
$$
中，第三个竖直箭头是同构。两行均正合；由引理
\ref{more-algebra-lemma-torsion-completion} 及
$M[f^\infty] = M'[f^\infty]$，第一个竖直箭头是同构。
五引理因而蕴含 $M \otimes_R R' \to M'$ 是同构。

\medskip\noindent
以上说明 $\text{Can}$ 本质满，且函子 $H^0$ 的像落在可粘合模范畴中。
由于对可粘合模，(\ref{more-algebra-equation-BL-cech-mod-re}) 正合，所以在可粘合
模范畴上有 $H^0 \circ \text{Can} = \text{id}$。结合引理
\ref{more-algebra-lemma-H0-Can-adjoint} 与 Categories，引理
\ref{categories-lemma-adjoint-fully-faithful}，这蕴含
$\text{Can}$ 全忠实。证明完成。
\end{proof}

\begin{remark}
\label{more-algebra-remark-what-you-get-for-general-modules}
设 $(R \to R', f)$ 为粘合对。设 $M$ 为不一定对
$(R \to R', f)$ 可粘合的 $R$-模。置
$M' = M \otimes_R R'$ 及 $M_1 = M_f$，得到粘合数据
$\text{Can}(M) = (M', M_1, \text{can})$。则
$\tilde M = H^0(M', M_1, \text{can})$ 是对
$(R \to R', f)$ 可粘合的 $R$-模，并且规范映射
$M \to \tilde M$ 给出同构
$M \otimes_R R' \to \tilde M \otimes_R R'$ 及
$M_f \to \tilde M_f$；见定理 \ref{more-algebra-theorem-BL-glueing}。
由序列
$$
M \to (M \otimes_R R' )\oplus M_f \to M \otimes_R (R')_f  \to 0
$$
及
$$
0 \to \tilde M \to (\tilde M \otimes_R R') \oplus \tilde M_f \to
\tilde M \otimes_R  (R')_f \to 0
$$
的正合性，可知映射 $M \to \tilde M$ 是满射。
\end{remark}

\noindent
回忆：粘合对 $(R \to R', f)$ 上的平坦 $R$-模是可粘合的
（评注 \ref{more-algebra-remark-glueable}）。因此，下一个引理说明定理
\ref{more-algebra-theorem-BL-glueing} 确定如下两个范畴之间的等价：
平坦 $R$-模的范畴，与粘合数据 $(M', M_1, \alpha_1)$ 的范畴，
其中 $M'$、$M_1$ 分别在 $R'$、$R_f$ 上平坦。

\begin{lemma}
\label{more-algebra-lemma-BL-flat}
设 $(R \to R', f)$ 为粘合对。设 $M$ 为不一定对
$(R \to R', f)$ 可粘合的 $R$-模。则 $M$ 在 $R$ 上平坦，
当且仅当 $M \otimes_R R'$ 在 $R'$ 上平坦，且
$M_f$ 在 $R_f$ 上平坦。
\end{lemma}

\begin{proof}
引理的一个方向来自 Algebra，引理
\ref{algebra-lemma-flat-base-change}。为证明另一个方向，假设
$M \otimes_R R'$ 在 $R'$ 上平坦，且 $M_f$ 在 $R_f$ 上平坦。
设 $\tilde M$ 如评注 \ref{more-algebra-remark-what-you-get-for-general-modules} 中所述。
若 $\tilde M$ 在 $R$ 上平坦，则对短正合列
$0 \to \Ker(M \to \tilde M) \to M \to \tilde M \to 0$
应用 Algebra，引理 \ref{algebra-lemma-flat-tor-zero}，可知
$\Ker(M \to \tilde M) \otimes_R (R' \oplus R_f)$ 为零。
故由引理 \ref{more-algebra-lemma-BL-faithful}，$M = \tilde M$，从而得到结论。
换言之，可以用 $\tilde M$ 替换 $M$，并假设 $M$ 对
$(R \to R', f)$ 可粘合。设 $N$ 为另一个 $R$-模。
由 Algebra，引理 \ref{algebra-lemma-characterize-flat}，只需证明
$\text{Tor}_1^R(M, N) = 0$。

\medskip\noindent
短正合列 $0 \to R \to R' \oplus R_f \to (R')_f \to 0$
与 $N$ 所对应的 Tor 长正合列给出正合列
$$
0 \to \text{Tor}_1^R(R', N) \to \text{Tor}_1^R((R')_f, N)
$$
以及对 $i \geq 2$ 的同构
$\text{Tor}_i^R(R', N) = \text{Tor}_i^R((R')_f, N)$。
由于
$\text{Tor}_i^R((R')_f, N) = \text{Tor}_i^R(R', N)_f$，
可知 $f$ 在 $\text{Tor}_1^R(R', N)$ 上是非零因子，并且在
$i \geq 2$ 时于 $\text{Tor}_i^R(R', N)$ 上可逆。
由于 $M \otimes_R R'$ 在 $R'$ 上平坦，由例
\ref{more-algebra-example-tor-change-rings} 的谱序列，有
$$
\text{Tor}_i^R(M \otimes_R R', N) =
(M \otimes_R R') \otimes_{R'} \text{Tor}_i^R(R', N)
$$
将 $M \otimes_R R'$ 写成有限自由 $R'$-模的滤过余极限
（Algebra，定理 \ref{algebra-theorem-lazard}），可知 $f$ 在
$\text{Tor}_1^R(M \otimes_R R', N)$ 上是非零因子，并在
$\text{Tor}_i^R(M \otimes_R R', N)$ 上可逆。
% P02-E0405
接着考虑由 $M$ 可粘合而来的正合列
$0 \to M \to M \otimes_R R' \oplus M_f \to M \otimes_R (R')_f \to 0$
及其所对应的 $\text{Tor}$ 长正合列。相关部分为
$$
\xymatrix{
\text{Tor}_1^R(M, N) \ar[r] &
\text{Tor}_1^R(M \otimes_R R', N) \ar[r] &
\text{Tor}_1^R(M \otimes_R (R')_f, N) \\
& \text{Tor}_2^R(M \otimes_R R', N) \ar[r] &
\text{Tor}_2^R(M \otimes_R (R')_f, N) \ar[llu]
}
$$
由上面对 $f$ 在
$\text{Tor}_i^R(M \otimes_R R', N)$ 上作用的讨论，可得
$\text{Tor}_1^R(M, N) = 0$。
% P02-E0406
\end{proof}

\noindent
注意，对于“有限生成”这一性质，我们已经在引理
\ref{more-algebra-lemma-faithful-descent} 中看到下一个引理的结论。

\begin{lemma}
\label{more-algebra-lemma-BL-properties}
设 $(R \to R', f)$ 为粘合对。设 $M$ 为不一定对
$(R \to R', f)$ 可粘合的 $R$-模。则 $M$ 是有限投射 $R$-模，
当且仅当 $M \otimes_R R'$ 在 $R'$ 上有限投射，且
$M_f$ 在 $R_f$ 上有限投射。
\end{lemma}

\begin{proof}
假设 $M \otimes_R R'$ 是有限投射 $R'$-模，且 $M_f$ 是有限投射
$R_f$-模。我们的任务是证明 $M$ 在 $R$ 上有限投射。下文将不再说明
而使用 Algebra，引理 \ref{algebra-lemma-finite-projective}。
由引理 \ref{more-algebra-lemma-BL-flat}，$M$ 平坦。由引理
\ref{more-algebra-lemma-faithful-descent}，$M$ 有限。选取短正合列
$0 \to K \to R^{\oplus n} \to M \to 0$。由于有限投射模是有限表现的，
且该序列与 $R'$ 作张量后（由 Algebra，引理
\ref{algebra-lemma-flat-tor-zero}）以及与 $R_f$ 作张量后仍正合，
可知 $K \otimes_R R'$ 与 $K_f$ 是有限模。利用上面的引理，
可知 $K$ 有限生成。因此 $M$ 有限表现，从而有限投射。
\end{proof}

\begin{remark}
\label{more-algebra-remark-compare-BL}
\cite{Beauville-Laszlo} 假设 $f$ 在 $R$ 中是非零因子，并且
$R' = R^\wedge$；由引理 \ref{more-algebra-lemma-same-f-torsion}，这给出粘合对。
即使在该框架下，定理 \ref{more-algebra-theorem-BL-glueing} 仍给出新内容：
\cite{Beauville-Laszlo} 的结果只适用于 $f$ 为非零因子的模
（因而按我们的定义可粘合；见引理
\ref{more-algebra-lemma-same-f-torsion-module}）。引理
\ref{more-algebra-lemma-BL-properties} 也轻微推广了
\cite{Beauville-Laszlo} 的结果：不仅可允许 $M$ 具有非零的
$f$-幂挠，甚至不要求它可粘合。
\end{remark}

\section{导出完备化}
\label{more-algebra-section-derived-completion}

\noindent
本节内容的一些参考文献为
\cite{Dwyer-Greenlees}、\cite{Greenlees-May}、\cite{PSY}、\cite{dag12}
（尤其是第 4 章）。我们的论述遵循 \cite{BS}。
本节对于挠模的类似物（或“对偶”）见 Dualizing Complexes，第
\ref{dualizing-section-local-cohomology} 节。具有挠同调的复形之导出范畴
与导出完备复形之间的关系见 Dualizing Complexes，第
\ref{dualizing-section-torsion-and-complete} 节。

\medskip\noindent
设 $K \in D(A)$，$f \in A$。用 $T(K, f)$ 表示系
% P02-E0408
$$
\ldots \to K \xrightarrow{f} K \xrightarrow{f} K
$$
在 $D(A)$ 中的一个导出极限。

\begin{lemma}
\label{more-algebra-lemma-hom-from-Af}
设 $A$ 为环，$f \in A$，$K \in D(A)$。下列条件等价：
\begin{enumerate}
\item 对所有 $n$，$\Ext^n_A(A_f, K) = 0$；
\item 对 $D(A_f)$ 中所有 $E$，$\Hom_{D(A)}(E, K) = 0$；
\item $T(K, f) = 0$；
\item 对每个 $p \in \mathbf{Z}$，$T(H^p(K), f) = 0$；
\item 对每个 $p \in \mathbf{Z}$，都有
$\Hom_A(A_f, H^p(K)) = 0$ 及 $\Ext^1_A(A_f, H^p(K)) = 0$；
\item $R\Hom_A(A_f, K) = 0$；
\item 映射 $\prod_{n \geq 0} K \to \prod_{n \geq 0} K$，
$(x_0, x_1, \ldots) \mapsto (x_0 - fx_1, x_1 - fx_2, \ldots)$
是 $D(A)$ 中的同构；
\item 此处添加更多条件。
% P02-E0409
\end{enumerate}
\end{lemma}

\begin{proof}
显然 (2) 蕴含 (1)，且 (1) 等价于 (6)。假设 (1)。
设 $I^\bullet$ 为表示 $K$ 的 K-内射 $A$-模复形。
条件 (1) 表示 $\Hom_A(A_f, I^\bullet)$ 非周期。
设 $M^\bullet$ 为表示 $E$ 的 $A_f$-模复形。则由 Algebra，引理
\ref{algebra-lemma-adjoint-hom-restrict}，
$$
\Hom_{D(A)}(E, K) =
\Hom_{K(A)}(M^\bullet, I^\bullet) =
\Hom_{K(A_f)}(M^\bullet, \Hom_A(A_f, I^\bullet))
$$
由引理 \ref{more-algebra-lemma-hom-K-injective}，
$\Hom_A(A_f, I^\bullet)$ 是 K-内射 $A_f$-模复形；
% P02-E0410
它非周期这一事实蕴含它同伦等价于零
（Derived Categories，引理 \ref{derived-lemma-K-injective}）。
故得到 (2)。

\medskip\noindent
$A$-模 $A_f$ 的一个自由消解为
$$
0 \to \bigoplus\nolimits_{n \in \mathbf{N}} A \to
\bigoplus\nolimits_{n \in \mathbf{N}} A
\to A_f \to 0
$$
其中第一个映射把 $(a_0, a_1, a_2, \ldots)$ 映到
% P02-E0411
$(a_0, a_1 - fa_0, a_2 - fa_1, \ldots)$，第二个映射把
$(a_0, a_1, a_2, \ldots)$ 映到
$a_0 + a_1/f + a_2/f^2 + \ldots$。应用
$\Hom_A(-, I^\bullet)$，得到
$$
0 \to \Hom_A(A_f, I^\bullet) \to \prod I^\bullet \to \prod I^\bullet \to 0
$$
由于 $\prod I^\bullet$ 表示 $\prod_{n \geq 0} K$，这证明
(1) 与 (7) 等价。另一方面，由 Derived Categories，第
\ref{derived-section-derived-limit} 节中导出极限的构造，
所显示的正合列说明对象 $T(K, f)$ 是 $D(A)$ 中
$R\Hom_A(A_f, K)$ 的一个代表。因此 (1) 与 (3) 等价。
% P02-E0412

\medskip\noindent
有谱序列
$$
E_2^{p, q} = \Ext^p_A(A_f, H^q(K)) \Rightarrow
\Ext^{p + q}_A(A_f, K)
$$
见等式 (\ref{more-algebra-equation-first-ss-ext})\footnote{若 $K$ 不有界下，
则不能使用所给参考。此时，选择一个表示 $K$ 的复形 $K^\bullet$
以及一个两项自由消解
$0 \to F^{-1} \to F^0 \to A_f \to 0$（见上文）。
% P02-E0413
于是得到各项为 $\Hom_A(F^{-q}, K^p)$ 的双复形，其总化计算
$R\Hom(A_f, K)$。与此双复形相联系的第二个谱序列
（Homology，第 \ref{homology-section-double-complex} 节）
给出所需结论。略去细节。}。
因为 $A_f$ 具有长度为 $1$ 的投射 $A$-模消解（见上文），
$E_2$ 页只有两个非零列，故此谱序列在 $E_2$ 处退化。
因此得到短正合列
$$
0 \to \Ext^1_A(A_f, H^{p - 1}(K)) \to
\Ext^p_A(A_f, K) \to
\Hom_A(A_f, H^p(K)) \to 0
$$
这证明 (4) 与 (5) 均等价于 (1)。
% P02-E0414
\end{proof}

\begin{lemma}
\label{more-algebra-lemma-ideal-of-elements-complete-wrt}
设 $A$ 为环，$K \in D(A)$。使 $T(K, f) = 0$ 的所有
$f \in A$ 组成的集合 $I$ 是 $A$ 的根理想。
\end{lemma}

\begin{proof}
下文将不再说明而使用引理 \ref{more-algebra-lemma-hom-from-Af} 的结果。
若 $f \in I$ 且 $g \in A$，则 $A_{gf}$ 是 $A_f$-模，故对所有
$n$ 有 $\Ext^n_A(A_{gf}, K) = 0$，从而 $gf \in I$。
设 $f, g \in I$。则有短正合列
$$
0 \to A_{f + g} \to A_{f(f + g)} \oplus A_{g(f + g)} \to A_{gf(f + g)} \to 0
$$
因为 $f, g$ 在 $A_{f + g}$ 中生成单位理想。这来自 Algebra，引理
\ref{algebra-lemma-standard-covering} 以及末箭头为满射这一简单事实。
由 $\Ext$ 的长正合列及对所有 $n$ 均有
$\Ext^n_A(A_{f(f + g)}, K)$、
$\Ext^n_A(A_{g(f + g)}, K)$ 与
$\Ext^n_A(A_{gf(f + g)}, K)$ 为零，可推出对所有 $n$，
$\Ext^n_A(A_{f + g}, K)$ 为零。最后，若对某个 $n > 0$ 有
$f^n \in I$，则 $f \in I$，因为
$T(K, f) = T(K, f^n)$，也可因为 $A_f \cong A_{f^n}$。
\end{proof}

\begin{lemma}
\label{more-algebra-lemma-complete-derived-complete}
设 $A$ 为环，$I \subset A$ 为理想，$M$ 为 $A$-模。
\begin{enumerate}
\item 若 $M$ 是 $I$-进完备的，则对所有 $f \in I$，
$T(M, f) = 0$。
\item 反之，若对所有 $f \in I$ 有 $T(M, f) = 0$，且 $I$ 有限生成，
则 $M \to \lim M/I^nM$ 是满射。
\end{enumerate}
\end{lemma}

\begin{proof}
(1) 的证明。假设 $M$ 是 $I$-进完备的。由引理
\ref{more-algebra-lemma-hom-from-Af}，只需证明
$\Ext^1_A(A_f, M) = 0$ 及 $\Hom_A(A_f, M) = 0$。
由于 $M = \lim M/I^nM$，并且
$\Hom_A(A_f, M/I^nM) = 0$，可得
$\Hom_A(A_f, M) = 0$。设有扩张
$$
0 \to M \to E \to A_f \to 0
$$
对 $n \geq 0$，选取映到 $1/f^n$ 的 $e_n \in E$。
对 $n \geq 0$，置 $\delta_n = fe_{n + 1} - e_n \in M$。
将 $e_n$ 替换为
$$
e'_n = e_n + \delta_n + f\delta_{n + 1} + f^2 \delta_{n + 2} + \ldots
$$
由于 $M$ 关于 $I$ 完备且 $f \in I$，此无穷和存在。
简单计算表明 $fe'_{n + 1} = e'_n$。因此，将 $1/f^n$ 映到
$e'_n$，即得到该扩张的分裂。

\medskip\noindent
(2) 的证明。假设 $I = (f_1, \ldots, f_r)$，并且对
$i = 1, \ldots, r$ 有 $T(M, f_i) = 0$。由 Algebra，引理
\ref{algebra-lemma-when-surjective-to-completion}，可以假设
$I = (f)$ 且 $T(M, f) = 0$。设对 $n \geq 0$ 有 $x_n \in M$。
考虑扩张
$$
0 \to M \to E \to A_f \to 0
$$
其中
$$
E = M \oplus \bigoplus Ae_n\Big/\langle x_n - fe_{n + 1} + e_n\rangle
$$
并把 $e_n$ 映到 $A_f$ 中的 $1/f^n$（见上文）。
由假设及引理 \ref{more-algebra-lemma-hom-from-Af}，此扩张分裂，故得到元素
$x + e_0$，它在 $E$ 中生成 $A_f$ 的一个副本。于是
$$
x + e_0 = x - x_0 + fe_1 = x - x_0 - f x_1 + f^2 e_2 = \ldots
$$
由蛇引理，$M/f^nM = E/f^nE$，故模 $f^nM$ 有
$x = x_0 + fx_1 + \ldots + f^{n - 1}x_{n - 1}$。
换言之，映射 $M \to \lim M/f^nM$ 满射，如所需。
\end{proof}

\noindent
受上述结果启发，作如下定义。

\begin{definition}
\label{more-algebra-definition-derived-complete}
设 $A$ 为环，$K \in D(A)$，$I \subset A$ 为理想。若对每个
$f \in I$ 都有 $T(K, f) = 0$，则称 $K$
{\it 关于 $I$ 导出完备}。若 $M$ 是 $A$-模，且
$M[0] \in D(A)$ 关于 $I$ 导出完备，则称 $M$
{\it 关于 $I$ 导出完备}。
\end{definition}

\noindent
由导出完备对象组成的全子范畴
$D_{comp}(A) = D_{comp}(A, I) \subset D(A)$ 是严格全、饱和的三角
子范畴；见 Derived Categories，定义
\ref{derived-definition-triangulated-subcategory} 与
\ref{derived-definition-saturated}。由引理
\ref{more-algebra-lemma-ideal-of-elements-complete-wrt}，子范畴
$D_{comp}(A, I)$ 只依赖于 $I$ 的根理想 $\sqrt{I}$；换言之，
它只依赖于 $\Spec(A)$ 的闭子集 $Z = V(I)$。
子范畴 $D_{comp}(A, I)$ 在 $D(A)$ 中的乘积与同伦极限下保持，
但一般不在可数直和下保持。若理想 $I$ 的选择由语境清楚，
我们常简称 $M$ 为导出完备模。

\begin{proposition}
\label{more-algebra-proposition-derived-complete-modules}
设 $I \subset A$ 为环 $A$ 的有限生成理想，$M$ 为 $A$-模。
下列条件等价：
\begin{enumerate}
\item $M$ 是 $I$-进完备的；
\item $M$ 关于 $I$ 导出完备，且 $\bigcap I^nM = 0$。
\end{enumerate}
\end{proposition}

\begin{proof}
这由引理 \ref{more-algebra-lemma-complete-derived-complete} 的结果立即得到。
\end{proof}

\noindent
下一个引理说明导出完备模的范畴 $\mathcal{C}$ 是 Abel 范畴。
$\mathcal{C}$ 并非 Grothendieck Abel 范畴；见 Examples，第
\ref{examples-section-derived-complete-modules} 节。

\begin{lemma}
\label{more-algebra-lemma-serre-subcategory}
设 $I$ 为环 $A$ 的理想。
\begin{enumerate}
\item 导出完备 $A$-模组成 $\text{Mod}_A$ 的弱 Serre 子范畴
$\mathcal{C}$。
\item $D_\mathcal{C}(A) \subset D(A)$ 是导出完备对象的全子范畴。
\end{enumerate}
\end{lemma}

\begin{proof}
(2) 由引理 \ref{more-algebra-lemma-hom-from-Af} 及定义立即得到。对 (1)，设
$M \to N$ 为导出完备模之间的映射。用 $K = (M \to N)$ 表示
% P02-E0415
$D(A)$ 中相应的对象。选取 $f \in I$。对所有 $n$，
$\Ext_A^n(A_f, K)$ 为零，因为对所有 $n$，
$\Ext_A^n(A_f, M)$ 与 $\Ext_A^n(A_f, N)$ 都为零。
故 $K$ 导出完备。由 (2)，$\Ker(M \to N)$ 与
$\Coker(M \to N)$ 是 $\mathcal{C}$ 的对象。
最后，设 $0 \to M_1 \to M_2 \to M_3 \to 0$ 是 $A$-模的短正合列，
且 $M_1$、$M_3$ 导出完备。则由 $\Ext$ 的长正合列，
$M_2$ 导出完备。因此，由 Homology，引理
\ref{homology-lemma-characterize-weak-serre-subcategory}，
$\mathcal{C}$ 是弱 Serre 子范畴。
\end{proof}

\noindent
我们将在引理 \ref{more-algebra-lemma-derived-complete-zero-bis} 中推广下一个引理。

\begin{lemma}
\label{more-algebra-lemma-derived-complete-zero}
设 $I$ 为环 $A$ 的有限生成理想，$M$ 为导出完备 $A$-模。
若 $M/IM = 0$，则 $M = 0$。
\end{lemma}

\begin{proof}
假设 $M/IM$ 为零。设 $I = (f_1, \ldots, f_r)$。
令 $i < r$ 为使
$N = M/(f_1, \ldots, f_i)M$ 非零的最大整数。
若这样的 $i$ 不存在，则 $M = 0$，这正是要证明的。
由这一选择及引理 \ref{more-algebra-lemma-serre-subcategory}，$N$ 作为
导出完备模之间一个映射的余核是导出完备的。又由 $i$ 的选择，
$f_{i + 1} : N \to N$ 是满射。故
$$
\lim (\ldots \to N \xrightarrow{f_{i + 1}} N \xrightarrow{f_{i + 1}} N)
$$
非零，这与 $N$ 的导出完备性矛盾。
\end{proof}

\noindent
若环是 $I$-进完备的，则可得到丰富的导出完备复形。

\begin{lemma}
\label{more-algebra-lemma-pseudo-coherent-is-derived-complete}
设 $A$ 为环，$I \subset A$ 为理想。若 $A$ 导出完备
（例如，$I$-进完备），则 $D(A)$ 的任意伪凝聚对象都导出完备。
\end{lemma}

\begin{proof}
（引理 \ref{more-algebra-lemma-complete-derived-complete} 解释了引理的括号内陈述。）
设 $K$ 为 $D(A)$ 的伪凝聚对象。按定义，这意味着 $K$ 可由有限自由
$A$-模组成的有界上复形 $K^\bullet$ 表示。由于 $A$ 导出完备，
由引理 \ref{more-algebra-lemma-serre-subcategory} 的 (1)，对所有 $n$，
$H^n(K)$ 导出完备。再由同一引理的 (2)，这蕴含 $K$ 导出完备。
\end{proof}

\begin{lemma}
\label{more-algebra-lemma-double-localize}
设 $A$ 为环，$f, g \in A$。则对 $K \in D(A)$，有
$R\Hom_A(A_f, R\Hom_A(A_g, K)) = R\Hom_A(A_{fg}, K)$。
\end{lemma}

\begin{proof}
这由引理 \ref{more-algebra-lemma-internal-hom} 得出。
\end{proof}

\begin{lemma}
\label{more-algebra-lemma-derived-completion}
\begin{slogan}
沿有限生成理想的导出完备化存在，并且可用 {\v C}ech 方法计算。
\end{slogan}
设 $I$ 为环 $A$ 的有限生成理想。
包含函子 $D_{comp}(A, I) \to D(A)$ 有左伴随；也就是说，给定
$D(A)$ 的任意对象 $K$，存在从 $K$ 到 $D(A)$ 的一个导出完备对象
% P02-E0416
$K \to K^\wedge$ 的映射，使得只要 $E$ 是 $D(A)$ 的导出完备对象，映射
$$
\Hom_{D(A)}(K^\wedge, E) \longrightarrow \Hom_{D(A)}(K, E)
$$
就是双射。事实上，若 $I$ 由 $f_1, \ldots, f_r \in A$ 生成，则有
$$
K^\wedge = R\Hom\left((A \to \prod\nolimits_{i_0} A_{f_{i_0}} \to
\prod\nolimits_{i_0 < i_1} A_{f_{i_0}f_{i_1}}
\to \ldots \to A_{f_1\ldots f_r}), K\right)
$$
且该等式关于 $K$ 函子性成立。
\end{lemma}

\begin{proof}
用引理中最后一个展示公式定义 $K^\wedge$。存在复形映射
$$
(A \to \prod\nolimits_{i_0} A_{f_{i_0}} \to
\prod\nolimits_{i_0 < i_1} A_{f_{i_0}f_{i_1}} \to
\ldots \to A_{f_1\ldots f_r}) \longrightarrow A
$$
它诱导映射 $K \to K^\wedge$。只需证明 $K^\wedge$ 导出完备，并且
若 $K$ 导出完备，则 $K \to K^\wedge$ 是同构\footnote{具体而言，
若 $E$ 导出完备且 $a : K \to E$ 是映射，则交换图
$$
\xymatrix{
K \ar[r]_a \ar[d] & E \ar[d]^{\cong} \\
K^\wedge \ar[r]^{a^\wedge} & E^\wedge
}
$$
表明到 $E$ 的任意映射都经 $K^\wedge$ 分解。选取一个杰出三角形
$K \to K^\wedge \to C$ 并施以 ${}^\wedge$。于是得到
$C^\wedge = 0$。由上文，这意味着对导出完备的 $E$，
$\Hom(C, E) = 0$，从而按所需证明了
$\Hom(K, E) = \Hom(K^\wedge, E)$。}。

\medskip\noindent
设 $f \in A$。由引理 \ref{more-algebra-lemma-double-localize}，对象
$R\Hom_A(A_f, K^\wedge)$ 等于
$$
R\Hom\left((A_f \to \prod\nolimits_{i_0} A_{ff_{i_0}} \to
\prod\nolimits_{i_0 < i_1} A_{ff_{i_0}f_{i_1}} \to
\ldots \to A_{ff_1\ldots f_r}), K\right)
$$
若 $f \in I$，则 $f_1, \ldots, f_r$ 在 $A_f$ 中生成单位理想，
故增广交错 {\v C}ech 复形
$$
A_f \to \prod\nolimits_{i_0} A_{ff_{i_0}} \to
\prod\nolimits_{i_0 < i_1} A_{ff_{i_0}f_{i_1}} \to
\ldots \to A_{ff_1\ldots f_r}
$$
由引理 \ref{more-algebra-lemma-extended-alternating-torsion} 在 $D(A)$ 中为零。
（事实上，若对某个 $i$ 有 $f = f_i$，则由引理
\ref{more-algebra-lemma-extended-alternating-homotopy-zero}，此复形零伦；
这是我们所需的唯一情形。）
因此 $R\Hom_A(A_f, K^\wedge) = 0$，再由引理
\ref{more-algebra-lemma-hom-from-Af} 得出 $K^\wedge$ 导出完备。

\medskip\noindent
反之，若 $K$ 导出完备，则对所有
$f = f_{i_0} \ldots f_{i_p}$、$p \geq 0$，
$R\Hom_A(A_f, K)$ 都为零。因此 $K \to K^\wedge$ 在
$D(A)$ 中是同构。
\end{proof}

\begin{remark}
\label{more-algebra-remark-derived-completion}
设 $A$ 为环，$I \subset A$ 为有限生成理想。由引理
\ref{more-algebra-lemma-derived-completion} 存在的包含函子
$D_{comp}(A, I) \to D(A)$ 的左伴随称为{\it 导出完备化}。
为表明这一点，我们会说“令 $K^\wedge$ 为 $K$ 的导出完备化”。
请记住，伴随的单位是函子性映射 $K \to K^\wedge$。
\end{remark}

\begin{lemma}
\label{more-algebra-lemma-derived-completion-vanishes}
设 $A$ 为环，$I \subset A$ 为有限生成理想。设 $K^\bullet$ 为
$A$-模复形，且对某个 $f \in I$，映射
$f : K^\bullet \to K^\bullet$ 是同构，也就是说，$K^\bullet$
是 $A_f$-模复形。则 $K^\bullet$ 的导出完备化为零。
\end{lemma}

\begin{proof}
事实上，在这种情形下，对任意导出完备复形 $L$，
% P02-E0417
$R\Hom_A(K, L)$ 都为零，见引理 \ref{more-algebra-lemma-hom-from-Af}。
因此，由引理 \ref{more-algebra-lemma-derived-completion} 中的泛性质，
$K^\wedge$ 为零。
\end{proof}

\begin{lemma}
\label{more-algebra-lemma-completion-RHom}
设 $A$ 为环，$I \subset A$ 为有限生成理想。设
$K, L \in D(A)$。则
$$
R\Hom_A(K, L)^\wedge = R\Hom_A(K, L^\wedge) = R\Hom_A(K^\wedge, L^\wedge)
$$
\end{lemma}

\begin{proof}
由引理 \ref{more-algebra-lemma-derived-completion}，导出完备化由某个
$C \in D(A)$ 的 $R\Hom_A(C, -)$ 给出。于是
\begin{align*}
R\Hom_A(C, R\Hom_A(K, L))
& =
R\Hom_A(C \otimes_A^\mathbf{L} K, L) \\
& =
R\Hom_A(K, R\Hom_A(C, L))
\end{align*}
这由引理 \ref{more-algebra-lemma-internal-hom} 得出。它证明第一个等式。
映射 $K \to K^\wedge$ 诱导映射
$$
R\Hom_A(K^\wedge, L^\wedge) \to R\Hom_A(K, L^\wedge)
$$
按照导出完备化作为包含函子的左伴随这一良定义，它在 $D(A)$
中是同构。
\end{proof}

\begin{lemma}
\label{more-algebra-lemma-naive-derived-completion}
设 $A$ 为环，$I \subset A$ 为理想。设 $(K_n)$ 为 $D(A)$ 中对象的
逆系，使得对所有 $f \in I$ 与 $n$，都存在
$e = e(n, f)$，使 $f^e$ 在 $K_n$ 上为零。则对 $K \in D(A)$，
对象 $K' = R\lim (K \otimes_A^\mathbf{L} K_n)$ 关于 $I$ 导出完备。
\end{lemma}

\begin{proof}
由于导出完备对象的范畴在 $R\lim$ 下保持，只需证明每个
$K \otimes_A^\mathbf{L} K_n$ 都导出完备。由假设，对所有
$f \in I$，都存在 $e$，使 $f^e$ 在
$K \otimes_A^\mathbf{L} K_n$ 上为零。这当然蕴含
$T(K \otimes_A^\mathbf{L} K_n, f) = 0$，结论即得。
\end{proof}

\begin{situation}
\label{more-algebra-situation-koszul}
设 $A$ 为环，$I = (f_1, \ldots, f_r) \subset A$。令
$K_n^\bullet = K_\bullet(A, f_1^n, \ldots, f_r^n)$ 为关于
$f_1^n, \ldots, f_r^n$ 的 Koszul 复形，并把它看作位于次数
$-r, -r + 1, \ldots, 0$ 的上链复形。利用引理
\ref{more-algebra-lemma-functorial} 的函子性，得到逆系
$$
\ldots \to K_3^\bullet \to K_2^\bullet \to K_1^\bullet
$$
它与逆系
$H^0(K_n^\bullet) = A/(f_1^n, \ldots, f_r^n)$
以及映射 $A \to K_n^\bullet$ 相容。
\end{situation}

\noindent
% P02-E0418
下文讨论的一个关键特征在于：若 $m > n$，则映射
$$
K_m^{-p} = \wedge^p(A^{\oplus r}) \to \wedge^p(A^{\oplus r}) = K_n^{-p}
$$
在基元素 $e_{i_1} \wedge \ldots \wedge e_{i_p}$ 上由乘以
$f_{i_1}^{m - n} \ldots f_{i_p}^{m - n}$ 给出。

\begin{lemma}
\label{more-algebra-lemma-koszul-derived-completion-complete}
在情形 \ref{more-algebra-situation-koszul} 中。对 $K \in D(A)$，对象
$K' = R\lim (K \otimes_A^\mathbf{L} K_n^\bullet)$ 关于 $I$ 导出完备。
\end{lemma}

\begin{proof}
这是引理 \ref{more-algebra-lemma-naive-derived-completion} 的特殊情形，因为
$f_i^n$ 在 $K_n^\bullet$ 上的作用是一个自同态，而由引理
\ref{more-algebra-lemma-homotopy-koszul}，该自同态零伦。
\end{proof}

\begin{lemma}
\label{more-algebra-lemma-characterize-derived-complete-Koszul}
在情形 \ref{more-algebra-situation-koszul} 中。设 $K \in D(A)$。
以下条件等价：
\begin{enumerate}
\item $K$ 关于 $I$ 导出完备；以及
\item 典范映射 $K \to R\lim (K \otimes_A^\mathbf{L} K_n^\bullet)$
是 $D(A)$ 中的同构。
\end{enumerate}
\end{lemma}

\begin{proof}
若 (2) 成立，则由引理
\ref{more-algebra-lemma-koszul-derived-completion-complete}，$K$ 关于 $I$
导出完备。反之，假设 $K$ 关于 $I$ 导出完备。考虑由愚蠢截断
（Homology，第 \ref{homology-section-truncations} 节）给出的滤过
$$
K_n^\bullet \supset
\sigma_{\geq -r + 1}K_n^\bullet \supset
\sigma_{\geq -r + 2}K_n^\bullet \supset \ldots \supset
\sigma_{\geq -1}K_n^\bullet \supset
\sigma_{\geq 0}K_n^\bullet = A
$$
。因为构造 $R\lim(K \otimes E)$ 关于第二个变元是正合的
（引理 \ref{more-algebra-lemma-tensor-Rlim-exact}），所以只需证明
$$
R\lim \left(
K \otimes_A^\mathbf{L}
(\sigma_{\geq p}K_n^\bullet/ \sigma_{\geq p + 1}K_n^\bullet)
\right) = 0
$$
对 $p < 0$ 成立。上面对 Koszul 复形的显式描述表明
$$
R\lim \left(
K \otimes_A^\mathbf{L}
(\sigma_{\geq p}K_n^\bullet/ \sigma_{\geq p + 1}K_n^\bullet)
\right) =
\bigoplus\nolimits_{i_1, \ldots, i_{-p}}
T(K, f_{i_1}\ldots f_{i_{-p}})
$$
，而由关于 $K$ 的假设，这在 $p < 0$ 时为零。
\end{proof}

\begin{lemma}
\label{more-algebra-lemma-derived-completion-koszul}
在情形 \ref{more-algebra-situation-koszul} 中。
把 $K \in D(A)$ 送到导出极限
$K' = R\lim( K \otimes_A^\mathbf{L} K_n^\bullet )$ 的函子，是引理
\ref{more-algebra-lemma-derived-completion} 中所构造的包含函子
$D_{comp}(A) \to D(A)$ 的左伴随。
\end{lemma}

\begin{proof}[第一种证明]
赋值 $K \leadsto K'$ 是函子，而且由引理
\ref{more-algebra-lemma-koszul-derived-completion-complete}，$K'$ 关于 $I$
导出完备。由一个形式论证（略去），只需证明当 $K$ 关于 $I$
导出完备时，$K \to K'$ 是同构。这正是引理
\ref{more-algebra-lemma-characterize-derived-complete-Koszul}。
\end{proof}

\begin{proof}[第二种证明]
% P02-E0419
用 $K \mapsto K^\wedge$ 表示引理
\ref{more-algebra-lemma-derived-completion} 中构造的伴随。由该引理，
$$
K^\wedge = R\Hom\left((A \to \prod\nolimits_{i_0} A_{f_{i_0}} \to
\prod\nolimits_{i_0 < i_1} A_{f_{i_0}f_{i_1}}
\to \ldots \to A_{f_1\ldots f_r}), K\right)
$$
在引理 \ref{more-algebra-lemma-extended-alternating-Cech-is-colimit-koszul} 中，
我们已经看到增广交错 {\v C}ech 复形
$$
A \to \prod\nolimits_{i_0} A_{f_{i_0}} \to
\prod\nolimits_{i_0 < i_1} A_{f_{i_0}f_{i_1}}
\to \ldots \to A_{f_1\ldots f_r}
$$
是位于次数 $0, \ldots, r$ 的 Koszul 复形
$K^n = K(A, f_1^n, \ldots, f_r^n)$ 的余极限。注意，$K^n$
是由有限自由 $A$-模组成的有限链复形，其对偶（如引理
\ref{more-algebra-lemma-dual-perfect-complex} 中所述）为
$R\Hom_A(K^n, A) = K_n$，其中 $K_n$ 是位于次数
$-r, \ldots, 0$ 的 Koszul 上链复形（同往常一样）。
因此只需证明
$$
R\Hom_A(\text{hocolim} K^n, K) =  R\lim (K \otimes_A^\mathbf{L} K_n)
$$
这由引理 \ref{more-algebra-lemma-colim-and-lim-of-duals} 得出。
\end{proof}

\begin{lemma}
\label{more-algebra-lemma-derived-complete-bounded}
设 $I = (f_1, \ldots, f_r)$ 为环 $A$ 的有限生成理想。
设 $K$ 为 $D(A)$ 的导出完备对象。以下条件等价：
\begin{enumerate}
\item 对 $i > 0$，有 $H^i(K) = 0$；
\item 对 $i > 0$，有 $H^i(K \otimes_A^\mathbf{L} A/I) = 0$；
\item 对 $i > 0$，有 $H^i(K \otimes_A^\mathbf{L} K_1^\bullet) = 0$，
其中 $K_1^\bullet$ 如情形 \ref{more-algebra-situation-koszul} 中所述。
\end{enumerate}
\end{lemma}

\begin{proof}
蕴含关系 (1) $\Rightarrow$ (2) 总成立。蕴含关系
(2) $\Rightarrow$ (3) 由引理 \ref{more-algebra-lemma-bounded-above-mod-I} 得出。
假设 (3)。对 $s = 0, \ldots, r$，考虑复形
$$
K(s) =  K \otimes_A^\mathbf{L} K_\bullet(A, f_1, \ldots, f_s)
$$
，其中 $K_\bullet(A, f_1, \ldots, f_s)$ 是置于上同调次数
$-s, \ldots, 0$ 的 Koszul 复形。由引理
\ref{more-algebra-lemma-cone-koszul}，我们有一个
% P02-E0420
杰出三角形
$$
K(s - 1) \xrightarrow{f_s} K(s - 1) \to K(s) \to K(s - 1)[1]
$$
。由于 $K(0) = K$ 导出完备，向上归纳得出每个 $K(s)$ 都是
$D(A)$ 的导出完备对象。我们将向下归纳证明 $K(s)$ 在次数
$> 0$ 没有非零上同调。事实上，由假设 (3)，这对 $s = r$
成立。若它对 $K(s)$ 成立且 $s > 0$，则对 $i > 0$，
$f_s : H^i(K(s - 1)) \to H^i(K(s - 1))$ 是满射；再由引理
\ref{more-algebra-lemma-serre-subcategory} 与
\ref{more-algebra-lemma-derived-complete-zero}，得出 $H^i(K(s - 1))$ 为零。
\end{proof}

\begin{lemma}
\label{more-algebra-lemma-derived-complete-zero-bis}
\begin{slogan}
导出 Nakayama
\end{slogan}
\begin{reference}
一个相关结果是 \cite[Proposition 6.5]{Dwyer-Greenlees}。
例如，在 Bhatt 于 2018 年秋季在哥伦比亚大学所作关于棱柱上同调的
第三讲中，可在第 2 节性质 (2) 找到导出 Nakayama 引理。Leonid
Positselski 在 \url{https://mathoverflow.net/a/331501} 中提出了一份证明。
不过，我们采用评论中 Anonymous 建议的证明。
\end{reference}
设 $I$ 为环 $A$ 的有限生成理想。设 $K$ 为 $D(A)$ 的导出完备对象。
若 $K \otimes_A^\mathbf{L} A/I = 0$，则 $K = 0$。
\end{lemma}

\begin{proof}
这是引理 \ref{more-algebra-lemma-derived-complete-bounded} 的一个特殊情形，
但我们也给出一个直接证明。选取 $I$ 的生成元
$f_1, \ldots, f_r$。
% P02-E0421
用 $K_n$ 表示 $A$ 上关于 $f_1^n, \ldots, f_r^n$ 的 Koszul 复形。
回忆 $K_n$ 有界，并且 $K_n$ 的上同调模被
$f_1^n, \ldots, f_r^n$ 消去，因而被 $I^{nr}$ 消去。
由引理 \ref{more-algebra-lemma-derived-vanishing-mod-I}，得到
$K \otimes_A^\mathbf{L} K_n = 0$。由于 $K$ 导出完备，由引理
\ref{more-algebra-lemma-derived-completion-koszul} 有
$K = R\lim K \otimes_A^\mathbf{L} K_n = 0$，如所需。
\end{proof}

\noindent
作为与 Koszul 复形之关系的一个应用，我们得到导出完备化具有
有限上同调维数。

\begin{lemma}
\label{more-algebra-lemma-derived-completion-finite-cohomological-dimension}
设 $A$ 为环，$I \subset A$ 为可由 $r$ 个元素生成的理想。
则导出完备化具有有限上同调维数：
\begin{enumerate}
\item 设 $K \to L$ 为 $D(A)$ 中的态射，使得
$H^i(K) \to H^i(L)$ 在 $i \geq 1$ 时是同构，在 $i = 0$ 时是满射。
则 $H^i(K^\wedge) \to H^i(L^\wedge)$ 在 $i \geq 1$ 时是同构，
在 $i = 0$ 时是满射。
\item 设 $K \to L$ 为 $D(A)$ 的态射，使得
$H^i(K) \to H^i(L)$ 在 $i \leq -1$ 时是同构，在 $i = 0$ 时是单射。
则 $H^i(K^\wedge) \to H^i(L^\wedge)$ 在 $i \leq -r - 1$ 时是同构，
在 $i = -r$ 时是单射。
\end{enumerate}
\end{lemma}

\begin{proof}
设 $I$ 由 $f_1, \ldots, f_r$ 生成。对任意 $K \in D(A)$，
由引理 \ref{more-algebra-lemma-derived-completion-koszul} 有
$K^\wedge = R\lim K \otimes_A^\mathbf{L} K_n$，其中 $K_n$ 是关于
$f_1^n, \ldots, f_r^n$ 的 Koszul 复形。因此，由引理
\ref{more-algebra-lemma-break-long-exact-sequence-modules} 得到短正合列
$$
0 \to R^1\lim H^{i - 1}(K \otimes_A^\mathbf{L} K_n)
\to H^i(K^\wedge) \to \lim H^i(K \otimes_A^\mathbf{L} K_n) \to 0
$$
。

\medskip\noindent
(1) 的证明。选取杰出三角形 $K \to L \to C \to K[1]$。
则对 $i \geq 0$，有 $H^i(C) = 0$。由于 $K_n$ 位于次数
$\leq 0$，可见对 $i \geq 0$，有
$H^i(C \otimes_A^\mathbf{L} K_n) = 0$，并且
$H^{-1}(C \otimes_A^\mathbf{L} K_n) =
H^{-1}(C) \otimes_A A/(f_1^n, \ldots, f_r^n)$ 是转移映射均为满射的
一个系。上面的展示等式表明对 $i \geq 0$，
$H^i(C^\wedge) = 0$。应用杰出三角形
$K^\wedge \to L^\wedge \to C^\wedge \to K^\wedge[1]$，即得 (1)。

\medskip\noindent
(2) 的证明。选取杰出三角形 $K \to L \to C \to K[1]$。
则对 $i < 0$，有 $H^i(C) = 0$。由于 $K_n$ 位于次数
$-r, \ldots, 0$，可见对 $i < -r$，有
$H^i(C \otimes_A^\mathbf{L} K_n) = 0$。上面的展示等式表明
% P02-E0422
对 $i < r$，有 $H^i(C^\wedge) = 0$。应用杰出三角形
$K^\wedge \to L^\wedge \to C^\wedge \to K^\wedge[1]$，即得 (2)。
\end{proof}

\begin{lemma}
\label{more-algebra-lemma-derived-completion-spectral-sequence}
设 $A$ 为环，$I \subset A$ 为有限生成理想。设 $K^\bullet$ 为
$A$-模的滤过复形。存在双分次导出完备 $A$-模的典范谱序列
$(E_r, \text{d}_r)_{r \geq 1}$，其中
% P02-E0423
$d_r$ 的双次数为 $(r, -r + 1)$，并且
$$
E_1^{p, q} = H^{p + q}((\text{gr}^pK^\bullet)^\wedge)
$$
若每个 $K^n$ 上的滤过有限，则该谱序列有界并收敛于
$H^*((K^\bullet)^\wedge)$。
\end{lemma}

\begin{proof}
由引理 \ref{more-algebra-lemma-derived-completion}，导出完备化由某个
$C \in D^b(A)$ 的 $R\Hom_A(C, -)$ 给出。由引理
\ref{more-algebra-lemma-derived-completion-finite-cohomological-dimension} 和
\ref{more-algebra-lemma-projective-amplitude}，可见 $C$ 具有有限投射维数。
因此可选取一个表示 $C$ 的有界投射模复形 $P^\bullet$。于是
$$
M^\bullet = \Hom^\bullet(P^\bullet, K^\bullet)
$$
是一个表示 $(K^\bullet)^\wedge$ 的 $A$-模复形。它带有由
$F^pM^\bullet = \Hom^\bullet(P^\bullet, F^pK^\bullet)$ 给出的滤过。
可见 $F^pM^\bullet$ 表示 $(F^pK^\bullet)^\wedge$，因而
% P02-E0424
$\text{gr}^pM^\bullet$ 表示 $(\text{gr}K^\bullet)^\wedge$。
于是取滤过复形 $M^\bullet$ 的谱序列，便得到我们的谱序列；
见 Homology，第 \ref{homology-section-filtered-complex} 节。
若每个 $K^n$ 上的滤过有限，则因为 $P^\bullet$ 是有界复形，
每个 $M^n$ 上的滤过也有限。因此，最后一个陈述由 Homology，
引理 \ref{homology-lemma-biregular-ss-converges} 得出。
\end{proof}

\begin{example}
\label{more-algebra-example-derived-completion-spectral-sequence}
设 $A$ 为环，$I \subset A$ 为有限生成理想。设 $K^\bullet$ 为
$A$-模复形。可以在引理
\ref{more-algebra-lemma-derived-completion-spectral-sequence} 中取
$F^pK^\bullet = \tau_{\leq -p}K^\bullet$。于是得到有界谱序列
$$
E_1^{p, q} = H^{p + q}(H^{-p}(K^\bullet)^\wedge[p]) =
H^{2p + q}(H^{-p}(K^\bullet)^\wedge)
$$
，它收敛于 $H^{p + q}((K^\bullet)^\wedge)$。重新编号
$p = -j$ 与 $q = i + 2j$ 后，得到：对任意 $K \in D(A)$，
存在双分次导出完备模的有界谱序列
$(E'_r, d'_r)_{r \geq 2}$，其中 $d'_r$ 的双次数为
$(r, -r + 1)$，满足
$$
(E'_2)^{i, j} = H^i(H^j(K)^\wedge)
$$
，并收敛于 $H^{i + j}(K^\wedge)$。
\end{example}

\begin{lemma}
\label{more-algebra-lemma-restriction-derived-complete}
设 $A \to B$ 为环映射，$I \subset A$ 为理想。$D_{comp}(A, I)$
在限制函子 $D(B) \to D(A)$ 下的原像是 $D_{comp}(B, IB)$。
\end{lemma}

\begin{proof}
利用引理 \ref{more-algebra-lemma-ideal-of-elements-complete-wrt} 可见，
$L \in D(B)$ 属于 $D_{comp}(B, IB)$，当且仅当对每个
% P02-E0425
局部截面 $f \in I$，$T(L, f)$ 为零。注意，$T(L, f)$ 的上同调
在阿贝尔群范畴中计算，所以无论把 $f$ 看作 $A$ 的元素，还是取
$f$ 在 $B$ 中的像，都没有区别。由此及导出完备对象的定义，
立即得到该引理。
\end{proof}

\begin{lemma}
\label{more-algebra-lemma-restriction-derived-complete-equivalence}
设 $A \to B$ 为环映射，$I \subset A$ 为有限生成理想。
若 $A \to B$ 平坦且 $A/I \cong B/IB$，则限制函子
$D(B) \to D(A)$ 诱导等价
$D_{comp}(B, IB) \to D_{comp}(A, I)$。
\end{lemma}

\begin{proof}
选取 $I$ 的生成元 $f_1, \ldots, f_r$。
% P02-E0426
用 $\check{\mathcal{C}}^\bullet_A \to \check{\mathcal{C}}^\bullet_B$
表示引理 \ref{more-algebra-lemma-map-identifies-koszul-and-cech-complexes} 中
增广交错 {\v C}ech 复形的拟同构。设
$K \in D_{comp}(A, I)$。设 $I^\bullet$ 为表示 $K$ 的
K-内射 $A$-模复形。由于对所有 $f \in I$ 与
$n \in \mathbf{Z}$，$\Ext^n_A(A_f, K)$ 和
$\Ext^n_A(B_f, K)$ 均为零（引理 \ref{more-algebra-lemma-hom-from-Af}），
我们得出 $\check{\mathcal{C}}^\bullet_A \to A$ 与
$\check{\mathcal{C}}^\bullet_B \to B$ 诱导拟同构
$$
I^\bullet = \Hom_A(A, I^\bullet) \longrightarrow
\text{Tot}(\Hom_A(\check{\mathcal{C}}^\bullet_A, I^\bullet))
$$
以及
$$
\Hom_A(B, I^\bullet) \longrightarrow
\text{Tot}(\Hom_A(\check{\mathcal{C}}^\bullet_B, I^\bullet))
$$
。略去一些细节。
由于 $\check{\mathcal{C}}^\bullet_A \to \check{\mathcal{C}}^\bullet_B$
是拟同构且 $I^\bullet$ 是 K-内射的，得出
$\Hom_A(B, I^\bullet) \to I^\bullet$ 是拟同构。
因为复形 $\Hom_A(B, I^\bullet)$ 是 $B$-模复形，得出 $K$
属于限制映射的像，也就是说，
% P02-E0427
该函子本质满。

\medskip\noindent
事实上，该论证表明
$F : D_{comp}(A, I) \to D_{comp}(B, IB)$，
$K \mapsto \Hom_A(B, I^\bullet)$ 是限制函子的左逆。最后，假设
$L \in D_{comp}(B, IB)$。用 $B$-模的 K-内射复形 $J^\bullet$
表示 $L$。则 $J^\bullet$ 作为 $A$-模复形也是 K-内射的
（引理 \ref{more-algebra-lemma-K-injective-flat}），因而
$F(\text{下列对象的限制： }L) = \Hom_A(B, J^\bullet)$。
存在 $B$-模复形的映射
$J^\bullet \to \Hom_A(B, J^\bullet)$，它与
$\Hom_A(B, J^\bullet) \to J^\bullet$ 的复合是恒等映射。
我们得出 $F$ 也是限制函子的右逆，证明完成。
\end{proof}

\section{导出完备模的范畴}
\label{more-algebra-section-derived-complete-modules}

\noindent
设 $A$ 为环，$I$ 为理想。
% P02-E0428
用 $\mathcal{C}$ 表示导出完备模的范畴，见定义
\ref{more-algebra-definition-derived-complete}。本节讨论此范畴的一些性质。
在 Examples，第 \ref{examples-section-derived-complete-modules} 节中，
我们说明 $\mathcal{C}$ 一般不是 Grothendieck 阿贝尔范畴。

\medskip\noindent
由引理 \ref{more-algebra-lemma-serre-subcategory}，范畴 $\mathcal{C}$ 是阿贝尔的，
且包含函子 $\mathcal{C} \to \text{Mod}_A$ 正合。

\medskip\noindent
由于 $D_{comp}(A) \subset D(A)$ 在乘积下封闭
（见定义 \ref{more-algebra-definition-derived-complete} 后的讨论），并且
$D(A)$ 中的乘积在复形层面计算，可见 $\mathcal{C}$ 有乘积，
且它们与 $\text{Mod}_A$ 中的乘积一致。因此，$\mathcal{C}$
事实上有任意极限，包含函子
$\mathcal{C} \to \text{Mod}_A$ 与它们交换；见 Categories，
引理 \ref{categories-lemma-limits-products-equalizers}。

\medskip\noindent
假设 $I$ 有限生成。用 ${}^\wedge : D(A) \to D(A)$ 表示引理
\ref{more-algebra-lemma-derived-completion} 的导出完备化函子。我们来证明函子
$$
\text{Mod}_A \longrightarrow \mathcal{C},\quad
M \longmapsto H^0(M^\wedge)
$$
是包含函子 $\mathcal{C} \to \text{Mod}_A$ 的左伴随。
例如，由引理
\ref{more-algebra-lemma-derived-completion-finite-cohomological-dimension}，
对 $i > 0$，有 $H^i(M^\wedge) = 0$。因此，若 $N$ 是导出完备
$A$-模，则有
\begin{align*}
\Hom_\mathcal{C}(H^0(M^\wedge), N)
& =
\Hom_{D_{comp}(A)}(M^\wedge, N)\\
& =
\Hom_{D(A)}(M, N) \\
& =
\Hom_A(M, N)
\end{align*}
，如所需。

\medskip\noindent
设 $T$ 为预序集，$t \mapsto M_t$ 为导出完备 $A$-模的一个系，
也就是 $\mathcal{C}$ 中以 $T$ 为指标的一个系；见 Categories，
第 \ref{categories-section-posets-limits} 节。
% P02-E0429
用 $\colim_{t \in T} M_t$ 表示该系在 $\text{Mod}_A$ 中的余极限。
由上文形式地得到
$$
H^0((\colim_{t \in T} M_t)^\wedge)
$$
是该系在 $\mathcal{C}$ 中的余极限。这样便看出
$\mathcal{C}$ 有所有余极限。一般而言，包含函子
$\mathcal{C} \to \text{Mod}_A$ 不与余极限交换；见 Examples，
第 \ref{examples-section-derived-complete-modules} 节。

\begin{lemma}
\label{more-algebra-lemma-derived-complete-modules}
设 $A$ 为环，$I \subset A$ 为理想。导出完备模的范畴
$\mathcal{C}$ 是阿贝尔的并有任意极限，且包含函子
$F : \mathcal{C} \to \text{Mod}_A$ 正合并与极限交换。
若 $I$ 有限生成，则 $\mathcal{C}$ 有任意余极限，且
% P02-E0430
$F$ 有左伴随
\end{lemma}

\begin{proof}
这是对上文讨论的总结。
\end{proof}

\section{主理想的导出完备化}
\label{more-algebra-section-derived-completion-principal}

\noindent
本节讨论理想由单个元素生成时导出完备化的情形。

\begin{lemma}
\label{more-algebra-lemma-lift-universally}
设 $A$ 为环，$f \in A$。若存在整数 $c \geq 1$，使
$A[f^c] = A[f^{c + 1}] = A[f^{c + 2}] = \ldots$
（例如，当 $A$ 为 Noether 环时），则对所有 $n \geq 1$，
在 $D(A)$ 中存在映射
$$
(A \xrightarrow{f^n} A) \longrightarrow A/(f^n),
\quad\text{且}\quad
A/(f^{n + c}) \longrightarrow (A \xrightarrow{f^n} A)
$$
，它们诱导 $D(A)$ 中 pro-对象 $\{A/(f^n)\}$ 与
$\{(f^n : A \to A)\}$ 的同构。
\end{lemma}

\begin{proof}
第一个展示箭头是显然的。可以用图
$$
\xymatrix{
A/A[f^c] \ar[r]_-{f^{n + c}} \ar[d]_{f^c} & A \ar[d]^1 \\
A \ar[r]^{f^n} & A
}
$$
定义引理中的第二个箭头。由于上方水平箭头是单射，上行中的复形
拟同构于 $A/f^{n + c}A$。我们略去为证明关于 pro 对象的陈述而需
进行的复合计算。
\end{proof}

\begin{lemma}
\label{more-algebra-lemma-when-does-it-work}
设 $A$ 为环且 $f \in A$。置 $I = (f)$。在此情形下，有朴素导出完备化
$K \mapsto K' = R\lim (K \otimes_A^\mathbf{L} A/f^nA)$，以及引理
\ref{more-algebra-lemma-derived-completion-koszul} 的导出完备化
$$
K \mapsto K^\wedge = R\lim (K \otimes_A^\mathbf{L} (A \xrightarrow{f^n} A))
$$
。函子的自然变换 $K^\wedge \to K'$ 是同构，当且仅当 $A$ 的
$f$-幂挠有界。
\end{lemma}

\begin{proof}
若 $f$-幂挠有界，则由引理 \ref{more-algebra-lemma-lift-universally}，pro-对象
$\{(f^n : A \to A)\}$ 与 $\{A/f^nA\}$ 同构。因此，由引理
\ref{more-algebra-lemma-Rlim-pro-equal}，这些函子同构。反之，由引理
\ref{more-algebra-lemma-tensor-Rlim-exact} 可见，该条件恰好是
$$
R\lim (K \otimes_A^\mathbf{L} A[f^n])
$$
对所有 $K \in D(A)$ 都为零。这里，系 $(A[f^n])$ 的映射由乘以
$f$ 给出。取 $K = A$ 以及 $K = \bigoplus_{i \in \mathbf{N}} A$，
由引理 \ref{more-algebra-lemma-Rlim-zero-of-direct-sums} 可见，这蕴含
$(A[f^n])$ 作为 pro-对象为零，也就是说，对某个 $n$，
$f^{n - 1}A[f^n] = 0$，即 $A[f^{n - 1}] = A[f^n]$，也即
$f$-幂挠有界。
\end{proof}

\begin{example}
\label{more-algebra-example-derived-complete-modules}
设 $A$ 为环，$f \in A$ 为非零因子。应当牢记的一个例子是
$A = \mathbf{Z}_p$ 与 $f = p$。设 $M$ 为 $A$-模。断言：
$M$ 关于 $f$ 导出完备，当且仅当存在短正合列
$$
0 \to K \to L \to M \to 0
$$
，其中 $K, L$ 是 $f$-进完备模，且其 $f$-挠为零。具体而言，
% P02-E0431
若存在这样的短正合列，则
$$
M \otimes_A^\mathbf{L} (A \xrightarrow{f^n} A) = (K/f^nK \to L/f^nL)
$$
，因为 $f$ 在 $K$ 和 $L$ 上都是非零因子；于是得出
$R\lim (M \otimes_A^\mathbf{L} (A \xrightarrow{f^n} A))$
拟同构于 $K \to L$，也就是 $M$。由引理
\ref{more-algebra-lemma-characterize-derived-complete-Koszul}，这表明 $M$
导出完备。

反之，假设 $M$ 导出完备。选取满射 $F \to M$，其中 $F$ 是自由
$A$-模。由于 $f$ 在 $F$ 上是非零因子，$F$ 的导出完备化为
$L = \lim F/f^nF$。注意 $L$ 无 $f$-挠：若 $x_n \in F$ 的
$(x_n)$ 表示 $L$ 的一个元素 $\xi$ 且 $f\xi = 0$，则对某些
$z_n, y_n \in F$，有 $x_n = x_{n + 1} + f^nz_n$ 与
$fx_n = f^ny_n$。于是
$f^n y_n = fx_n =
fx_{n + 1} + f^{n + 1}z_n = f^{n + 1}y_{n + 1} + f^{n + 1}z_n$；
由于 $f$ 在 $F$ 上是非零因子，可见 $y_n \in fF$，这蕴含
$x_n \in f^nF$，即 $\xi = 0$。由于 $L$ 是导出完备化，
泛性质给出分解 $F \to M$ 的映射 $L \to M$。令
$K = \Ker(L \to M)$ 为其核。仍然有 $K$ 无 $f$-挠，故 $K$
的导出完备化为 $\lim K/f^nK$。另一方面，$M$ 与 $L$
都导出完备，所以由引理 \ref{more-algebra-lemma-serre-subcategory}，$K$
也导出完备。因此 $K = \lim K/f^nK$，断言得证。
\end{example}

\begin{example}
\label{more-algebra-example-derived-complete-not-complete}
设 $p$ 为素数。考虑将 $x$ 送到 $py$ 的多项式代数映射
$\mathbf{Z}_p[x] \to \mathbf{Z}_p[y]$。考虑诱导的（通常）
$p$-进完备化映射的余核
$M = \Coker(\mathbf{Z}_p[x]^\wedge \to \mathbf{Z}_p[y]^\wedge)$。
由命题 \ref{more-algebra-proposition-derived-complete-modules} 和引理
\ref{more-algebra-lemma-serre-subcategory}，$M$ 是导出完备
$\mathbf{Z}_p$-模；另见例 \ref{more-algebra-example-derived-complete-modules}
中的讨论。然而，$M$ 不是 $p$-进完备的，因为
$1 + py + p^2 y^2 + \ldots$ 映到 $M$ 的一个非零元素，而该元素
属于 $\bigcap p^nM$。
\end{example}

\begin{example}
\label{more-algebra-example-spectral-sequence-principal}
设 $A$ 为环且 $f \in A$。
% P02-E0432
用 $K \mapsto K^\wedge$ 表示关于 $(f)$ 的导出完备化。设 $M$
为 $A$-模。利用
$$
M^\wedge = R\lim (M \xrightarrow{f^n} M)
$$
（引理 \ref{more-algebra-lemma-derived-completion-koszul}）以及引理
\ref{more-algebra-lemma-break-long-exact-sequence-modules}，得到
$$
H^{-1}(M^\wedge) = \lim M[f^n] = T_f(M)
$$
，即{\it $M$ 的 $f$-进 Tate 模}。这里，映射
$M[f^n] \to M[f^{n - 1}]$ 由乘以 $f$ 给出。于是有描述
$H^0(M^\wedge)$ 的短正合列
$$
0 \to R^1\lim M[f^n] \to H^0(M^\wedge) \to \lim M/f^n M \to 0
$$
。因为转移映射是满射（引理 \ref{more-algebra-lemma-compute-Rlim-modules}），
有 $H^1(M^\wedge) = R^1\lim M/f^nM = 0$。由于平凡的理由，
$M^\wedge$ 的所有其他上同调也为零。最后，对
$K \in D(A)$ 与 $p \in \mathbf{Z}$，有短正合列
$$
0 \to H^0(H^p(K)^\wedge) \to H^p(K^\wedge) \to T_f(H^{p + 1}(K)) \to 0
$$
。这由例 \ref{more-algebra-example-derived-completion-spectral-sequence} 的谱序列
得出，因为它在 $E_2$ 退化（只有 $i = -1, 0$ 给出非零项）；
下一个引理给出更多信息。
\end{example}

\begin{lemma}
\label{more-algebra-lemma-compare-completion-derived-completion}
设 $A$ 为环，$f \in A$，$K$ 为 $D(A)$ 的对象。
% P02-E0433
置 $K_n = K \otimes_A^\mathbf{L} (A \xrightarrow{f^n} A)$。
对所有 $p \in \mathbf{Z}$，存在交换图
$$
\xymatrix{
& 0 & 0 \\
0 \ar[r] &
\widehat{H^p(K)} \ar[r] \ar[u] &
\lim H^p(K_n) \ar[r] \ar[u] &
T_f(H^{p + 1}(K)) \ar[r] &
0 \\
0 \ar[r] &
H^0(H^p(K)^\wedge) \ar[r] \ar[u] &
H^p(K^\wedge) \ar[r] \ar[u] &
T_f(H^{p + 1}(K)) \ar[r] \ar@{=}[u] &
0 \\
&
R^1\lim H^p(K)[f^n] \ar[u] \ar[r]^\cong &
R^1\lim H^{p - 1}(K_n) \ar[u] \\
& 0 \ar[u] & 0 \ar[u]
}
$$
，其行与列均正合，其中
$\widehat{H^p(K)} = \lim H^p(K)/f^nH^p(K)$ 是通常的 $f$-进完备化。
左侧竖直短正合列与中间水平短正合列取自例
% P02-E0434
\ref{more-algebra-example-spectral-sequence-principal}。中间竖直短正合列来自引理
\ref{more-algebra-lemma-break-long-exact-sequence-modules}。
\end{lemma}

\begin{proof}
为构造上方水平短正合列，注意我们有如下逆系的短正合列
$$
0 \to H^p(K)/f^nH^p(K) \to H^p(K_n) \to H^{p + 1}(K)[f^n] \to 0
$$
，它来自把 $K_n$ 构造为 $f^n : K \to K$ 的锥的移位。
取这些逆极限，便得到上方水平短正合列；见 Homology，引理
\ref{homology-lemma-Mittag-Leffler}。

\medskip\noindent
我们来证明有引理中的交换图。考虑映射
$L = \tau_{\leq p}K \to K$。置
$L_n = L \otimes_A^\mathbf{L} (A \xrightarrow{f^n} A)$，
得到逆系的映射 $(L_n) \to (K_n)$，它诱导短正合列的映射
$$
\xymatrix{
0 & 0 \\
\lim H^p(L_n) \ar[r] \ar[u] &
\lim H^p(K_n) \ar[u] \\
H^p(L^\wedge) \ar[r] \ar[u] &
H^p(K^\wedge) \ar[u] \\
R^1\lim H^{p - 1}(L_n) \ar[r] \ar[u] &
R^1\lim H^{p - 1}(K_n) \ar[u] \\
0 \ar[u] & 0 \ar[u]
}
$$
。由于对 $i > p$ 有 $H^i(L) = 0$，且
$H^p(L) = H^p(K)$，利用引理陈述中的参考结果进行计算可得
$H^p(L^\wedge) = H^0(H^p(K)^\wedge)$，以及
$H^p(L_n) = H^p(K)/f^nH^p(K)$。另一方面，有
$H^{p - 1}(L_n) = H^{p - 1}(K_n)$，所以得到引理陈述中所示的
同构，因为已经知道
$H^0(H^p(K)^\wedge) \to \widehat{H^p(K)}$ 的核等于
$R^1\lim H^p(K)[f^n]$。若用证明第一段的构造定义顶行，
我们略去验证图中最右方正方形交换。
\end{proof}

\begin{remark}
\label{more-algebra-remark-ML-condition}
沿用引理 \ref{more-algebra-lemma-compare-completion-derived-completion} 的记号，
我们还看到，逆系 $H^p(K_n)$ 具有 ML 性质，当且仅当逆系
$H^{p + 1}(K)[f^n]$ 具有 ML 性质。这由逆系的短正合列
$0 \to H^p(K)/f^nH^p(K) \to H^p(K_n) \to H^{p + 1}(K)[f^n] \to 0$
（见该引理的证明），结合 Homology，引理
\ref{homology-lemma-Mittag-Leffler} 与引理
\ref{more-algebra-lemma-Mittag-Leffler} 得出。
\end{remark}

\begin{lemma}[Bhatt]
\label{more-algebra-lemma-torsion-and-derived-complete}
设 $I$ 为环 $A$ 中的有限生成理想。设 $M$ 为导出完备 $A$-模。
若 $M$ 是 $I$-幂挠模，则对某个 $n$ 有 $I^nM = 0$。
\end{lemma}

\begin{proof}
设 $I = (f_1, \ldots, f_r)$。只需证明对每个 $i$，都存在
$n_i$，使 $f_i^{n_i}M = 0$。因此可假设 $I = (f)$ 为主理想。
设 $B = \mathbf{Z}[x] \to A$ 为将 $x$ 送到 $f$ 的环映射。
由引理 \ref{more-algebra-lemma-restriction-derived-complete}，$M$ 作为
$B$-模关于理想 $(x)$ 导出完备。用 $B$ 代替 $A$ 后，可假设
$f$ 是 $A$ 中的非零因子。

\medskip\noindent
假设 $I = (f)$，其中 $f \in A$ 为非零因子。根据例
\ref{more-algebra-example-derived-complete-modules}，存在短正合列
$$
0 \to K \xrightarrow{u} L \to M \to 0
$$
，其中 $K$ 与 $L$ 是 $I$-进完备 $A$-模，且其 $f$-挠为零
\footnote{对该证明，只需说明存在序列
$K \xrightarrow{u} L \to M \to 0$，其中 $K$ 与 $L$ 是
$I$-进完备 $A$-模。为此可选取一个表示
$F_1 \to F_0 \to M \to 0$，其中 $F_i$ 自由，然后分别令
$K$ 与 $L$ 等于 $F_1$ 与 $F_0$ 的 $f$-进完备化。事实上，
由于 $f$ 是非零因子，这些完备化就是导出完备化，且该序列
仍然正合。}。把 $K$ 与 $L$ 看作带 $I$-进拓扑的拓扑模。
则 $u$ 连续。令
$$
L_n = \{x \in L \mid f^n x \in u(K)\}
$$
。由于 $M$ 是 $f$-幂挠的，可见 $L = \bigcup L_n$。
设 $N_n$ 为 $L_n$ 在 $L$ 中的闭包。由引理
\ref{more-algebra-lemma-consequence-baire-complete-module}，对某个 $n$，
$N_n$ 在 $L$ 中开。固定这样的 $n$。由于
$f^{n + m} : L \to L$ 是连续开映射，且
$f^{n + m} L_n \subset u(f^m K)$，得出对所有 $m \geq 1$，
$u(f^mK)$ 的闭包都是开的。因此，由引理
\ref{more-algebra-lemma-open-mapping}，得出 $u$ 是开映射。故对某个 $t$，
$f^tL \subset \Im(u)$，从而 $f^t$ 消去 $M$，如所需。
\end{proof}

\begin{lemma}
\label{more-algebra-lemma-kernel-to-completion-square-zero}
设 $f \in A$ 为环中的元素，置 $J = \bigcap f^nA$。设 $M$
为关于 $f$ 导出完备的 $A$-模。则 $JM' = 0$，其中
$M' = \Ker(M \to \lim M/f^nM)$。特别地，若 $A$ 导出完备，
则 $J$ 是平方零理想。
\end{lemma}

\begin{proof}
取 $x \in M'$ 与 $g \in J$。对每个 $n \geq 1$，可写
$x = f^n x_n$。由于 $g$ 属于 $f^nA$，可见 $M'$ 中的元素
$y_n = gx_n$ 与 $x_n$ 的选择无关。特别地，可取
$x_n = fx_{n + 1}$，于是得到 $y_n = fy_{n + 1}$。因此得到映射
$A_f \to M$，它把 $1/f^n$ 送到 $y_n$。由于 $M$ 导出完备，
此映射必须为零（引理 \ref{more-algebra-lemma-hom-from-Af}）
% P02-E0435
，因而对所有 $n$ 都有 $y_n = 0$。由于
$gx = gfx_1 = fy_1$，证明完成。
\end{proof}

\begin{lemma}
\label{more-algebra-lemma-derived-complete-henselian}
设 $A$ 为关于理想 $I$ 导出完备的环。则 $(A, I)$ 是 Hensel 对。
\end{lemma}

\begin{proof}
设 $f \in I$。由引理 \ref{more-algebra-lemma-largest-ideal-henselian}，
只需证明 $(A, fA)$ 是 Hensel 对。注意 $A$ 关于 $fA$ 导出完备
（这立即由定义 \ref{more-algebra-definition-derived-complete} 得出）。
由引理 \ref{more-algebra-lemma-complete-derived-complete}，从 $A$ 到
$A$ 的 $f$-进完备化 $A'$ 的映射是满射。由引理
\ref{more-algebra-lemma-complete-henselian}，对 $(A', fA')$ 是 Hensel 对。
因此，只需证明 $(A, \bigcap f^nA)$ 是 Hensel 对；见引理
\ref{more-algebra-lemma-henselian-henselian-pair}。这由引理
\ref{more-algebra-lemma-kernel-to-completion-square-zero} 与
\ref{more-algebra-lemma-locally-nilpotent-henselian} 得出。
\end{proof}

\begin{lemma}
\label{more-algebra-lemma-kernel-to-completion-nilpotent}
设 $A$ 为关于理想 $I$ 导出完备的环。置 $J = \bigcap I^n$。
若 $I$ 可由 $r$ 个元素生成，则 $J^N = 0$，其中 $N = 2^r$。
\end{lemma}

\begin{proof}
当 $r = 1$ 时，这是引理
\ref{more-algebra-lemma-kernel-to-completion-square-zero}。设
$I = (f_1, \ldots, f_r)$，其中 $r > 1$。由引理
\ref{more-algebra-lemma-serre-subcategory}，环 $A_t = A/f_r^tA$ 关于 $I$
导出完备，因而更关于
$I_t = (f_1, \ldots, f_{r - 1})A_t$ 导出完备。注意
$A \to A_t$ 把 $J$ 送入 $J_t = \bigcap I_t^n$。由归纳，
$J_t^{N/2} = 0$，其中 $N = 2^r$。理想
$\bigcap \Ker(A \to A_t) = \bigcap f_r^t A$ 由 $r = 1$
的情形是平方零的。这就完成证明。
\end{proof}

\begin{lemma}
\label{more-algebra-lemma-reduced-derived-complete-complete}
设 $A$ 为关于有限生成理想 $I$ 导出完备的既约环。
则 $A$ 是 $I$-进完备的。
\end{lemma}

\begin{proof}
% P02-E0436
这由引理 \ref{more-algebra-lemma-kernel-to-completion-nilpotent} 与命题
\ref{more-algebra-proposition-derived-complete-modules} 得出。
\end{proof}

\section{Noether 环的导出完备化}
\label{more-algebra-section-derived-completion-noetherian}

\noindent
设 $A$ 为环，$I \subset A$ 为理想。对任意 $K \in D(A)$，
可考虑导出极限
$$
K' = R\lim (K \otimes_A^\mathbf{L} A/I^n)
$$
。这是关于 $K$ 的函子，见评注
\ref{more-algebra-remark-constructing-tensor-with-limits-functorially}。映射系
$A \to A/I^n$ 诱导映射 $K \to K'$，且 $K'$ 关于 $I$
导出完备（引理 \ref{more-algebra-lemma-naive-derived-completion}）。
一般而言，这一“朴素”导出完备化构造并不等于引理
\ref{more-algebra-lemma-derived-completion} 的伴随。例如，若
$A = \mathbf{Z}_p \oplus \mathbf{Q}_p/\mathbf{Z}_p$，且第二个直和项
为平方零理想，$K = A[0]$，$I = (p)$，则朴素导出完备化给出
$\mathbf{Z}_p[0]$，而引理 \ref{more-algebra-lemma-derived-completion} 的构造给出
$K^\wedge \cong \mathbf{Z}_p[1] \oplus \mathbf{Z}_p[0]$
（略去计算）。当 $I$ 由单个元素生成时，引理
\ref{more-algebra-lemma-when-does-it-work} 刻画了这两个函子何时一致。

\medskip\noindent
% P02-E0437
本节的主要目标是证明：若 $A$ 为 Noether 环，则朴素导出完备化
等于导出完备化。

\begin{lemma}
\label{more-algebra-lemma-sequence-Koszul-complexes}
在情形 \ref{more-algebra-situation-koszul} 中。若 $A$ 为 Noether 环，则
$D(A)$ 的 pro-对象 $\{K_n^\bullet\}$ 与
$\{A/(f_1^n, \ldots, f_r^n)\}$ 同构
\footnote{特别地，对每个 $n$，都存在 $m \geq n$，使
$K_m^\bullet \to K_n^\bullet$ 经映射
$K_m^\bullet \to A/(f_1^m, \ldots, f_r^m)$ 分解。}。
\end{lemma}

\begin{proof}
% P02-E0438
我们有杰出三角形的一个逆系
$$
\tau_{\leq -1}K_n^\bullet \to K_n^\bullet \to
A/(f_1^m, \ldots, f_r^m) \to
(\tau_{\leq -1}K_n^\bullet)[1]
$$
，见 Derived Categories，评注
\ref{derived-remark-truncation-distinguished-triangle}。由 Derived
Categories，引理 \ref{derived-lemma-pro-isomorphism}，只需证明逆系
$\tau_{\leq -1}K_n^\bullet$ 是 pro-零的。回忆 $K_n^\bullet$
仅在满足 $-r \leq i \leq 0$ 的次数 $i$ 上有非零项。因此，由
Derived Categories，引理
\ref{derived-lemma-essentially-constant-cohomology}，只需证明对
$p \leq -1$，$H^p(K_n^\bullet)$ 是 pro-零的。换言之，对每个
$n \in \mathbf{N}$，必须证明存在 $m \geq n$，使
$H^p(K_m^\bullet) \to H^p(K_n^\bullet)$ 为零。由于 $A$
为 Noether 环，可见
$$
H^p(K_n^\bullet) =
\frac{\Ker(K_n^p \to K_n^{p + 1})}{\Im(K_n^{p - 1} \to K_n^p)}
$$
是有限 $A$-模。此外，当 $p < 0$ 时，映射
$K_m^p \to K_n^p$ 由一个对角矩阵给出，其各项属于理想
$(f_1^{m - n}, \ldots, f_r^{m - n})$。注意
$H^p(K_n^\bullet)$ 被
$J = (f_1^n, \ldots, f_r^n)$ 消去，见引理
\ref{more-algebra-lemma-homotopy-koszul}。现在，当 $m - n \geq tn$ 时，
有 $(f_1^{m - n}, \ldots, f_r^{m - n}) \subset J^t$。因此，把
Algebra，引理 \ref{algebra-lemma-Artin-Rees}（Artin--Rees）
应用于理想 $J$ 和模 $M = K_n^p$ 及其子模
% P02-E0439
$N = \Ker(K_n^p \to K_n^{p + 1})$，则当 $m$ 充分大时，
$K_m^p \to K_n^p$ 的像与
$\Ker(K_n^p \to K_n^{p + 1})$ 的交包含于
$J \Ker(K_n^p \to K_n^{p + 1})$。对这样的 $m$，得到零映射。
\end{proof}

\begin{proposition}
\label{more-algebra-proposition-noetherian-naive-completion-is-completion}
设 $A$ 为 Noether 环，$I \subset A$ 为理想。把
$K \in D(A)$ 送到导出极限
$K' = R\lim( K \otimes_A^\mathbf{L} A/I^n )$ 的函子，是引理
\ref{more-algebra-lemma-derived-completion} 中所构造的包含函子
$D_{comp}(A) \to D(A)$ 的左伴随。
\end{proposition}

\begin{proof}
设 $(f_1, \ldots, f_r) = I$，并令 $K_n^\bullet$ 为关于
$f_1^n, \ldots, f_r^n$ 的 Koszul 复形。由引理
\ref{more-algebra-lemma-derived-completion-koszul}，只需证明
$$
R\lim (K \otimes_A^\mathbf{L} K_n^\bullet) =
R\lim (K \otimes_A^\mathbf{L} A/(f_1^n, \ldots, f_r^n) ) =
R\lim (K \otimes_A^\mathbf{L} A/I^n ).
$$
由引理 \ref{more-algebra-lemma-sequence-Koszul-complexes}，$D(A)$ 的 pro-对象
$\{K_n^\bullet\}$ 与 $\{A/(f_1^n, \ldots, f_r^n)\}$ 同构。
显然，pro-对象 $\{A/(f_1^n, \ldots, f_r^n)\}$ 与
$\{A/I^n\}$ 同构。因此，由引理
\ref{more-algebra-lemma-tensor-Rlim-pro-equal}，从左到右的映射是同构。
\end{proof}

\begin{lemma}
\label{more-algebra-lemma-noetherian-calculate}
设 $I$ 为 Noether 环 $A$ 的理想，$M$ 为 $A$-模，其导出完备化为
$M^\wedge$。则存在短正合列
$$
0 \to R^1\lim \text{Tor}_{i + 1}^A(M, A/I^n) \to
H^{-i}(M^\wedge) \to \lim \text{Tor}_i^A(M, A/I^n) \to 0
$$
。对 $M \in D^-(A)$ 也有类似结果。
\end{lemma}

\begin{proof}
% P02-E0440
这是命题 \ref{more-algebra-proposition-noetherian-naive-completion-is-completion}
与引理 \ref{more-algebra-lemma-break-long-exact-sequence-modules} 的直接推论。
\end{proof}

\noindent
作为上述命题的一个应用，我们在 Noether 情形下为伪凝聚复形
识别导出完备化。

\begin{lemma}
\label{more-algebra-lemma-derived-completion-pseudo-coherent}
设 $A$ 为 Noether 环，$I \subset A$ 为理想。设 $K$ 为 $D(A)$
的对象，使得对所有
% P02-E0441
$n \in \mathbf{Z}$，$H^n(K)$ 都是有限 $A$-模。则导出完备化的
上同调模 $H^n(K^\wedge)$ 是上同调模 $H^n(K)$
的 $I$-进完备化。
\end{lemma}

\begin{proof}
由引理 \ref{more-algebra-lemma-Noetherian-pseudo-coherent}，对所有 $m$，
复形 $\tau_{\leq m}K$ 都是伪凝聚的。因此
$\tau_{\leq m}K$ 可由有限自由 $A$-模的有界上复形
$P^\bullet$ 表示。于是
$\tau_{\leq m}K \otimes_A^\mathbf{L} A/I^n = P^\bullet/I^nP^\bullet$。
故
$(\tau_{\leq m}K)^\wedge = R\lim P^\bullet/I^nP^\bullet$
（命题 \ref{more-algebra-proposition-noetherian-naive-completion-is-completion}）；
又因为 $R\lim$ 由逐项的 $\lim$ 给出
（引理 \ref{more-algebra-lemma-compute-Rlim-modules}），且 $I$-进完备化在有限
$A$-模上是正合函子
（Algebra，引理 \ref{algebra-lemma-completion-flat}），得出该结果
对 $\tau_{\leq m}K$ 成立。由于导出完备化具有有限上同调维数，
该结果因而对 $K$ 成立；见引理
\ref{more-algebra-lemma-derived-completion-finite-cohomological-dimension}。
\end{proof}

\begin{lemma}
\label{more-algebra-lemma-derived-complete-finite}
设 $I$ 为 Noether 环 $A$ 的理想。设 $M$ 为导出完备 $A$-模。
若 $M/IM$ 是有限 $A/I$-模，则
$M = \lim M/I^nM$，且 $M$ 是有限 $A^\wedge$-模。
\end{lemma}

\begin{proof}
假设 $M/IM$ 有限。选取 $x_1, \ldots, x_t \in M$，使其像生成
$M/IM$。得到映射 $A^{\oplus t} \to M$，它把第 $i$ 个基向量
送到 $x_i$。由命题
\ref{more-algebra-proposition-noetherian-naive-completion-is-completion}，$A$
的导出完备化为 $A^\wedge = \lim A/I^n$。由于 $M$ 导出完备，
可见我们的映射经一个映射
$q : (A^\wedge)^{\oplus t} \to M$ 分解。由引理
\ref{more-algebra-lemma-derived-complete-zero}，模 $\Coker(q)$ 为零。
因此 $M$ 是有限 $A^\wedge$-模。由于 $A^\wedge$ 是 Noether 环，
且关于 $IA^\wedge$ 完备，得出 $M$ 是 $I$-进完备的
（使用 Algebra，引理
\ref{algebra-lemma-completion-Noetherian}、
\ref{algebra-lemma-when-finite-module-complete-over-complete-ring} 以及
\ref{algebra-lemma-Artin-Rees}）。
\end{proof}

\begin{lemma}
\label{more-algebra-lemma-when-derived-completion-is-completion}
设 $I$ 为 Noether 环 $A$ 中的理想。
\begin{enumerate}
\item 若 $M$ 是有限 $A$-模且 $N$ 是平坦 $A$-模，则
$M \otimes_A N$ 的导出 $I$-进完备化就是
$M \otimes_A N$ 的通常 $I$-进完备化。
\item 若 $M$ 是有限 $A$-模且 $f \in A$，则 $M_f$ 的导出
$I$-进完备化就是 $M_f$ 的通常 $I$-进完备化。
\end{enumerate}
\end{lemma}

\begin{proof}
对 $A$-模 $M$，
% P02-E0442
用 $M^\wedge$ 表示导出完备化，用 $\lim M/I^nM$ 表示通常完备化。
假设 $M$ 有限。对 $i > 0$，系
$\text{Tor}^A_i(M, A/I^n)$ 是 pro-零的，见引理
\ref{more-algebra-lemma-tor-strictly-pro-zero}。因为 $N$ 平坦，
$\text{Tor}_i^A(M \otimes_A N, A/I^n) =
\text{Tor}_i^A(M, A/I^n) \otimes_A N$，所以系
$\text{Tor}^A_i(M \otimes_A N, A/I^n)$ 也如此。
由引理 \ref{more-algebra-lemma-noetherian-calculate}，得出
$R\lim (M \otimes_A N) \otimes_A^\mathbf{L} A/I^n$
仅在次数 $0$ 有上同调，且该上同调由通常完备化
$\lim M \otimes_A N/ I^n(M \otimes_A N)$ 给出。这证明 (1)。
(2) 由 (1) 以及 $M_f = M \otimes_A A_f$ 得出。
\end{proof}

\begin{lemma}
\label{more-algebra-lemma-derived-completion-tensor-finite}
设 $I$ 为 Noether 环 $A$ 中的理想。令
${}^\wedge$ 表示关于 $I$ 的导出完备化。设 $K \in D^-(A)$。
\begin{enumerate}
\item 若 $M$ 是有限 $A$-模，则
$(K \otimes_A^\mathbf{L} M)^\wedge = K^\wedge \otimes_A^\mathbf{L} M$。
\item 若 $L \in D(A)$ 是伪凝聚的，则
$(K \otimes_A^\mathbf{L} L)^\wedge = K^\wedge \otimes_A^\mathbf{L} L$。
\end{enumerate}
\end{lemma}

\begin{proof}
设 $L$ 如 (2) 中所述。可以用自由 $A$-模的有界上复形
$P^\bullet$ 表示 $K$，并用有限自由 $A$-模的有界上复形
$F^\bullet$ 表示 $L$。由于
$\text{Tot}(P^\bullet \otimes_A F^\bullet)$ 表示
$K \otimes_A^\mathbf{L} L$，可见
$(K \otimes_A^\mathbf{L} L)^\wedge$ 由
$$
\text{Tot}((P^\bullet)^\wedge \otimes_A F^\bullet)
$$
表示，其中 $(P^\bullet)^\wedge$ 的各项是
% P02-E0443
通常的 $=$ 导出完备化 $(P^n)^\wedge$；例如见命题
\ref{more-algebra-proposition-noetherian-naive-completion-is-completion} 与引理
\ref{more-algebra-lemma-when-derived-completion-is-completion}。这证明 (2)。
(1) 是 (2) 的特殊情形。
\end{proof}

\section{Berthelot 与 Ogus 引入的算子}
\label{more-algebra-section-eta}

\noindent
本节讨论 \cite[Section 8]{Berthelot-Ogus} 中引入并在
\cite[Section 6]{BMS} 中推广的一个构造。我们强烈建议读者查看
讨论这一概念的原始论文。

\medskip\noindent
设 $A$ 为环，$f \in A$ 为非零因子。若 $M$ 为 $A$-模，则由引理
% P02-E0444
\ref{more-algebra-lemma-torsion-free}，以下条件等价：
\begin{enumerate}
\item $f$ 在 $M$ 上是非零因子；
\item $M[f] = 0$；
\item 对所有 $n \geq 1$，有 $M[f^n] = 0$；以及
\item 映射 $M \to M_f$ 是单射。
\end{enumerate}
若这些等价条件成立，则（在本节中）我们称{\it $M$ 无 $f$-挠}。
在这种情况下，
% P02-E0445
用 $f^iM \subset M_f$ 表示由形如 $f^ix$（其中 $x \in M$）
的元素组成的子模。当然，$f^iM$ 作为 $A$-模同构于 $M$。
设 $M^\bullet$ 为由无 $f$-挠 $A$-模组成的复形，其微分为
$d^i : M^i \to M^{i + 1}$。在这种情形下，定义
$\eta_fM^\bullet$ 为具有各项
$$
(\eta_fM)^i = \{x \in f^iM^i \mid d^i(x) \in f^{i + 1}M^{i + 1}\}
$$
以及由 $d^i$ 诱导的微分的复形。注意 $\eta_fM^\bullet$
是另一个由无 $f$-挠 $A$-模组成的复形。若
$a^\bullet : M^\bullet \to N^\bullet$ 是无 $f$-挠 $A$-模复形的
映射，则由映射 $f^iM^i \to f^iN^i$ 诱导复形映射
$$
\eta_fa^\bullet : \eta_fM^\bullet \longrightarrow \eta_fN^\bullet
$$
。读者可检验，我们在无 $f$-挠 $A$-模复形的范畴上得到一个自函子。
若 $a^\bullet, b^\bullet : M^\bullet \to N^\bullet$ 是无
$f$-挠 $A$-模复形的两个映射，而
$h = \{h^i : M^i \to N^{i - 1}\}$ 是 $a^\bullet$ 与 $b^\bullet$
之间的伦，则定义 $\eta_fh$ 为映射族
$(\eta_fh)^i : (\eta_fM)^i \to (\eta_fN)^{i - 1}$，它把 $x$
送到 $h^i(x)$；这是有意义的，因为 $x \in f^iM^i$ 蕴含
$h^i(x) \in f^iN^{i - 1}$，而后者当然包含于
$(\eta_fN)^{i - 1}$。读者可检验，$\eta_fh$ 是
$\eta_fa^\bullet$ 与 $\eta_fb^\bullet$ 之间的伦。总之，可见在无
$f$-挠 $A$-模这一加性范畴的伦范畴
（Derived Categories，第 \ref{derived-section-homotopy} 节）上，
我们得到函子
$$
\eta_f :
K(f\text{-torsion free }A\text{-模})
\longrightarrow
K(f\text{-torsion free }A\text{-模})
$$
。$\eta_f$ 在任何意义下都不是三角范畴的正合函子；见例
\ref{more-algebra-example-eta-not-distinguished}。

\begin{example}
\label{more-algebra-example-eta-not-distinguished}
设 $A$ 为环，$f \in A$ 为非零因子。考虑函子
$\eta_f : K(f\text{-torsion free }A\text{-模})
\to K(f\text{-torsion free }A\text{-模})$。设
$M^\bullet$ 为无 $f$-挠 $A$-模的复形。乘以 $f$ 定义同构
$\eta_f(M^\bullet[1]) \to (\eta_fM^\bullet)[1]$，所以在这一意义下
$\eta_f$ 与移位相容。然而，考虑图
$$
\xymatrix{
A \ar[r]_f &
A \ar[r]_1 &
A \ar[r] &
0 \\
0 \ar[r] \ar[u] &
0 \ar[r] \ar[u] &
A \ar[r]^{-1} \ar[u]^f &
A \ar[u]
}
$$
。把每一列看作无 $f$-挠 $A$-模的复形，其中上方的模在次数 $1$，
其下的模在次数 $0$。于是，此图在
$K(f\text{-torsion free }A\text{-模})$ 中给出一个杰出三角形，
其三角结构如 Derived Categories，第
\ref{derived-section-homotopy-triangulated} 节中所述。具体而言，
第三个复形是前两个复形之间映射的锥。然而，对每一列施以
$\eta_f$，得到
$$
\xymatrix{
fA \ar[r]_f &
fA \ar[r]_1 &
fA \ar[r] &
0 \\
0 \ar[r] \ar[u] &
0 \ar[r] \ar[u] &
A \ar[r]^{-1} \ar[u]^f &
A \ar[u]
}
$$
。然而，第三个复形是非周期的，甚至零伦。因此，若这是杰出三角形，
则第一个箭头必须是伦范畴中的同构；除非 $f$ 是单位，否则这不成立。
\end{example}

\begin{lemma}
\label{more-algebra-lemma-eta-first-property}
设 $A$ 为环，$f \in A$ 为非零因子。设 $M^\bullet$ 为无
$f$-挠 $A$-模的复形。存在由乘以 $f^i$ 给出的典范同构
$$
f^i : H^i(M^\bullet)/H^i(M^\bullet)[f] \longrightarrow H^i(\eta_fM^\bullet)
$$
。
\end{lemma}

\begin{proof}
注意
$\Ker(d^i : (\eta_fM)^i \to (\eta_fM)^{i + 1})$ 等于
$\Ker(d^i : f^iM^i \to f^iM^{i + 1}) =
f^i\Ker(d^i : M^i \to M^{i + 1})$。因此，
% P02-E0446
把 $z \in \Ker(d^i : M^i \to M^{i + 1})$ 的类送到 $f^iz$ 的类，
得到满射
$f^i : H^i(M^\bullet) \to H^i(\eta_fM^\bullet)$。若由此在
$H^i(\eta_fM^\bullet)$ 中得到零类，则对某个
$y \in M^{i - 1}$，有
$f^i z = d^{i - 1}(f^{i - 1}y)$。因为 $f$ 在所涉及的所有模上
都是非零因子，这意味着 $f z = d^{i - 1}(y)$，而这恰好意味着
$z$ 的类是 $f$-挠的，如所需。
\end{proof}

\begin{lemma}
\label{more-algebra-lemma-eta-qis}
设 $A$ 为环，$f \in A$ 为非零因子。若
$M^\bullet \to N^\bullet$ 是无 $f$-挠 $A$-模复形的拟同构，
则诱导映射 $\eta_fM^\bullet \to \eta_fN^\bullet$ 也是拟同构。
\end{lemma}

\begin{proof}
这是因为引理 \ref{more-algebra-lemma-eta-first-property} 的同构与复形映射相容。
\end{proof}

\begin{lemma}
\label{more-algebra-lemma-Leta}
设 $A$ 为环，$f \in A$ 为非零因子。存在加性函子
\footnote{注意，此函子不正合，也就是说，它不把杰出三角形变换为
杰出三角形。见例 \ref{more-algebra-example-eta-not-distinguished}。}
$L\eta_f : D(A) \to D(A)$，使得若 $M \in D(A)$ 由无
$f$-挠 $A$-模复形 $M^\bullet$ 表示，则
$L\eta_fM = \eta_fM^\bullet$；对态射也类似。
\end{lemma}

\begin{proof}
% P02-E0447
用 $\mathcal{T} \subset \text{Mod}_A$ 表示无 $f$-挠 $A$-模的
全子范畴。相应地，有 $K(\mathcal{T})$ 作为 $K(A)$ 的全三角
子范畴的包含
$$
K(\mathcal{T}) \quad\subset\quad K(\text{Mod}_A) = K(A)
$$
。设 $S \subset \text{Arrows}(K(\mathcal{T}))$ 为拟同构。
我们将应用 Derived Categories，引理
\ref{derived-lemma-localization-subcategory} 证明映射
$$
S^{-1}K(\mathcal{T}) \longrightarrow D(A)
$$
是三角范畴的等价。该引理表明，只需证明：给定 $A$-模复形
$M^\bullet$，存在拟同构 $K^\bullet \to M^\bullet$，其中
$K^\bullet$ 是无 $f$-挠模的复形。由引理
\ref{more-algebra-lemma-K-flat-resolution}，可找到拟同构
$K^\bullet \to M^\bullet$，使复形 $K^\bullet$ 是 K-平坦的
（我们不会用到这一点），且由平坦 $A$-模 $K^i$ 组成。
特别地，对所有 $i$，$f$ 在 $K^i$ 上都是非零因子，如所需。

\medskip\noindent
准备工作完成后，可以定义 $L\eta_f$。事实上，由本节开头的讨论，
我们已有良定义函子
$$
K(\mathcal{T}) \xrightarrow{\eta_f} K(\mathcal{T}) \to K(A) \to D(A)
$$
，而根据引理 \ref{more-algebra-lemma-eta-qis}，它把拟同构送到拟同构。因此，
% P02-E0448
由 Categories，引理
\ref{categories-lemma-properties-left-localization}，此函子经
$S^{-1}K(\mathcal{T}) = D(A)$ 分解。
\end{proof}

\begin{remark}
\label{more-algebra-remark-eta-BZ}
设 $A$ 为环，$f \in A$ 为非零因子。设 $M^\bullet$ 为无
$f$-挠 $A$-模的复形。对每个 $i$，置
$\overline{M}^i = M^i/fM^i$。
% P02-E0449
用 $B^i \subset Z^i \subset \overline{M}^i$ 表示复形
$\overline{M}^\bullet = M^\bullet \otimes_A A/fA$ 的微分的
上边界与上循环。我们断言存在行正合的交换图
$$
\xymatrix{
0 \ar[r] &
B^{i + 1} \ar[r] \ar@{=}[d] &
B^{i + 1} \oplus B^i \ar[r] \ar[d]^{s, s'} &
B^i \ar[r] \ar[d] & 0 \\
0 \ar[r] &
B^{i + 1} \ar[r]^-s &
(\eta_fM)^i /f(\eta_fM)^i \ar[r]^-t &
Z^i \ar[r] & 0
}
$$
。这些映射构造如下：
\begin{enumerate}
\item 若 $x \in (\eta_fM)^i$，则 $x = f^ix'$，且
$d^i(x') = 0$ 在 $\overline{M}^{i + 1}$ 中成立。因此，可把 $x$
送到 $x'$ 的类来定义映射 $t$。
\item 若 $y \in M^{i + 1}$ 在
$B^{i + 1} \subset \overline{M}^{i + 1}$ 中的类为
$\overline{y}$，则可对某些 $y' \in M^{i + 1}$ 与 $x \in M^i$
写成 $y = fy' + d^i(x)$。因此，可把 $\overline{y}$ 送到
$(\eta_fM)^i /f(\eta_fM)^i$ 中 $f^{i + 1}x$ 的类来定义映射
$s$；我们略去该映射良定义的验证。
\item 若 $x \in M^i$ 在 $B^i \subset \overline{M}^i$ 中的类为
$\overline{x}$，则可对某些 $x' \in M^i$ 与
$z \in M^{i - 1}$ 写成
$x = fx' + d^{i - 1}(z)$。定义映射 $s'$，把 $\overline{x}$
送到 $(\eta_fM)^i/f(\eta_fM)^i$ 中 $f^i d^{i - 1}(z)$ 的类。
这是良定义的，因为若 $fx' + d^{i - 1}(z) = 0$，则
$f^ix'$ 属于 $(\eta_fM)^i$，因而
$f^id^{i - 1}(z)$ 属于 $f(\eta_fM)^i$。
\end{enumerate}
我们略去验证展示图的下行是模的短正合列。由这些构造立即可见，
我们有交换图
$$
\xymatrix{
B^{i + 1} \oplus B^i \ar[d]^{s, s'} \ar[r] &
B^{i + 2} \oplus B^{i + 1} \ar[d]^{s, s'} \\
(\eta_fM)^i /f(\eta_fM)^i \ar[r] &
(\eta_fM)^{i + 1} /f(\eta_fM)^{i + 1}
}
$$
，其中上方水平箭头由源与靶中直和项 $B^{i + 1}$ 的同一给出。
换言之，我们找到了
$\eta_fM^\bullet / f(\eta_fM^\bullet) =
\eta_fM^\bullet \otimes_A A/fA$ 的一个非周期子复形，而以此子复形
取商所得复形的各项 $Z^i/B^i$，是复形
$\overline{M}^\bullet = M^\bullet \otimes_A A/fA$ 的上同调模。
\end{remark}

\noindent
为以更典范的方式解释评注 \ref{more-algebra-remark-eta-BZ} 中观察到的现象，
我们将构造 Bockstein 算子。设 $A$ 为环，$f \in A$ 为非零因子。
设 $M^\bullet$ 为无 $f$-挠 $A$-模的复形。对每个
$i \in \mathbf{Z}$，都有交换图（张量积均在 $A$ 上取）
$$
\xymatrix{
0 \ar[r] &
M^\bullet \otimes f^{i + 1}A \ar[r] \ar[d] &
M^\bullet \otimes f^iA \ar[r] \ar[d] &
M^\bullet \otimes f^iA/f^{i + 1}A \ar[r] \ar@{=}[d] &
0 \\
0 \ar[r] &
M^\bullet \otimes f^{i + 1}A/f^{i + 2}A \ar[r] &
M^\bullet \otimes f^iA/f^{i + 2}A \ar[r] &
M^\bullet \otimes f^iA/f^{i + 1}A \ar[r] &
0
}
$$
，其两行都是复形的短正合列。当然，通过适当地乘以 $f$ 的幂，
不同 $i$ 的这些短正合列都彼此同构。特别地，下方序列的上同调
长正合列对所有 $i \in \mathbf{Z}$ 确定{\it Bockstein 算子}
$$
\beta = \beta^i : H^i(M^\bullet \otimes f^iA/f^{i + 1}A) \to
H^{i + 1}(M^\bullet \otimes f^{i + 1}A/f^{i + 2}A)
$$
。为供后文使用，我们在此记录：由上面的交换图，Bockstein 算子
具有分解
\begin{equation}
\label{more-algebra-equation-factorization-bockstein}
\vcenter{
\xymatrix{
H^i(M^\bullet \otimes f^iA/f^{i + 1}A)
\ar[r]_\delta \ar[rd]_\beta &
H^{i + 1}(M^\bullet \otimes f^{i + 1}A) \ar[d] \\
&
H^{i + 1}(M^\bullet \otimes f^{i + 1}A/f^{i + 2}A)
}
}
\end{equation}
，其中 $\delta$ 是来自上述交换图顶行的边界算子。
我们来证明得到一个复形
\begin{equation}
\label{more-algebra-equation-complex-bocksteins}
H^\bullet(M^\bullet/f) =
\left[
\begin{matrix}
\ldots \\
\downarrow \\
H^{i - 1}(M^\bullet \otimes f^{i - 1}A/f^iA) \\
\downarrow \beta \\
H^i(M^\bullet \otimes f^iA/f^{i + 1}A) \\
\downarrow \beta \\
H^{i + 1}(M^\bullet \otimes f^{i + 1}A/f^{i + 2}A) \\
\downarrow \\
\ldots
\end{matrix}
\right]
\end{equation}
，也就是说，$\beta \circ \beta = 0$
\footnote{另一种方法是论证 $\beta$ 作为复形
$(M^\bullet)_f$ 关于子复形 $f^iM^\bullet$ 的滤过之谱序列的
微分出现。还有一个证明更强结论的论证：先考虑
$M^\bullet = A$ 的情形。此时短正合列
$0 \to f^{i + 1}A/f^{i + 2}A \to f^iA/f^{i + 2}A \to
f^iA/f^{i + 1}A \to 0$ 在 $D(A)$ 中定义映射
$\beta^i : f^iA/f^{i + 1}A \to f^{i + 1}A/f^{i + 2}A[1]$。
然后可计算（论证与正文类似）出复合
$f^iA/f^{i + 1}A \to f^{i + 1}A/f^{i + 2}A[1] \to
f^{i + 2}A/f^{i + 3}A[2]$ 在 $D(A)$ 中为零。由于关于
$M^\bullet$ 的各项无 $f$-挠这一假设，有
$M^\bullet \otimes f^iA/f^{i + 1}A =
M^\bullet \otimes^\mathbf{L} f^iA/f^{i + 1}A$，所以得出
% P02-E0450
复合
$$
(M^\bullet \otimes f^iA/f^{i + 1}A)
\to
(M^\bullet \otimes f^{i + 1}A/f^{i + 2}A)[1]
\to
(M^\bullet \otimes f^{i + 2}A/f^{i + 3}A)[2]
$$
在 $D(A)$ 中也为零。}。事实上，利用分解
(\ref{more-algebra-equation-factorization-bockstein})，只需证明
$$
H^{i + 1}(M^\bullet \otimes f^{i + 1}A)
\to
H^{i + 1}(M^\bullet \otimes f^{i + 1}A/f^{i + 2}A)
\xrightarrow{\beta^{i + 1}}
H^{i + 2}(M^\bullet \otimes f^{i + 2}A/f^{i + 3}A)
$$
为零。这成立是因为 $\beta^{i + 1}$ 的核由可提升到
$H^{i + 1}(M^\bullet \otimes f^{i + 1}A/f^{i + 3}A)$ 的
上同调类组成，而第一个映射像中的类当然可以提升！

\begin{lemma}
\label{more-algebra-lemma-eta-second-property}
设 $A$ 为环，$f \in A$ 为非零因子。设 $M^\bullet$ 为无
$f$-挠 $A$-模的复形。存在复形的典范映射
$$
\eta_fM^\bullet \otimes_A A/fA
\longrightarrow
H^\bullet(M^\bullet/f)
$$
，它是拟同构，其中右端是复形
(\ref{more-algebra-equation-complex-bocksteins})。
\end{lemma}

\begin{proof}
设 $x \in (\eta_fM)^i$。则
% P02-E0451
$x = f^ix' \in f^iM$，且
$d^i(x) = f^{i + 1}y \in f^{i + 1}M^{i + 1}$。于是 $d^i$
把 $x' \otimes f^i$ 映到
$M^{i + 1} \otimes f^iA/f^{i + 1}A$ 中的零。本证明中的所有
张量积都在 $A$ 上取。因此，可把 $x$ 映到
$H^i(M^\bullet \otimes f^iA/f^{i + 1}A)$ 中
$x' \otimes f^i$ 的类。显然，此规则定义 $A/fA$-模的映射
$$
(\eta_fM)^i \otimes A/fA
\longrightarrow
H^i(M^\bullet \otimes f^iA/f^{i + 1}A)
$$
。注意，在上述情形中，可把 $x' \otimes f^i$ 看作
$M^i \otimes f^iA/f^{i + 2}A$ 的元素，其微分为
$d^i(x' \otimes f^i) = y \otimes f^{i + 1}$。由上面对
$\beta$ 的构造，得到
$\beta(x' \otimes f^i) = y \otimes f^{i + 1}$，因而我们的映射
与微分相容，也就是说，得到复形的映射。

\medskip\noindent
为完成证明，注意上一段给出的构造与评注 \ref{more-algebra-remark-eta-BZ}
中讨论的映射
$(\eta_fM)^i \otimes A/fA \to Z^i/B^i$ 一致。由于已经看到这些
映射的核是 $\eta_fM^\bullet \otimes A/fA$ 的非周期子复形，
引理得证。
\end{proof}

\begin{lemma}
\label{more-algebra-lemma-vanishing-beta}
设 $A$ 为环，$f \in A$ 为非零因子。设 $M^\bullet$ 为无
$f$-挠 $A$-模的复形。对 $i \in \mathbf{Z}$，以下条件等价：
\begin{enumerate}
\item $\Ker(d^i \bmod f^2)$ 满射到 $\Ker(d^i \bmod f)$；
\item 映射
$\beta : H^i(M^\bullet \otimes_A f^iA/f^{i + 1}A) \to
H^{i + 1}(M^\bullet \otimes_A f^{i + 1}A/f^{i + 2}A)$ 为零。
\end{enumerate}
条件 $H^{i + 1}(M^\bullet)[f] = 0$ 蕴含这些等价条件。
\end{lemma}

\begin{proof}
(1) 与 (2) 的等价性，由 $\beta$ 定义为一个复形短正合列的
上同调边界映射，以及该短正合列同构于复形短正合列
$0 \to fM^\bullet/f^2M^\bullet \to M^\bullet/f^2M^\bullet \to
M^\bullet/fM^\bullet \to 0$ 得出。若 $\beta \not = 0$，则由分解
(\ref{more-algebra-equation-factorization-bockstein})，
$H^{i + 1}(M^\bullet)[f] \not = 0$。
\end{proof}

\begin{lemma}
\label{more-algebra-lemma-eta-vanishing-beta}
设 $A$ 为环，$f \in A$ 为非零因子。设 $M^\bullet$ 为无
$f$-挠 $A$-模的复形。若 $\Ker(d^i \bmod f^2)$ 满射到
$\Ker(d^i \bmod f)$，则典范映射
$$
(1, d^i) :
(\eta_fM)^i / f(\eta_fM)^i \longrightarrow
f^iM^i/f^{i + 1}M^i \oplus f^{i + 1}M^{i + 1}/f^{i + 2}M^{i + 1}
$$
把左端同一为右端的若干子模的直和。
\end{lemma}

\begin{proof}
沿用评注 \ref{more-algebra-remark-eta-BZ} 的记号，定义映射
$t^{-1} : Z^i \to (\eta_fM)^i / f(\eta_fM)^i$。具体而言，
对满足 $d^i(x) = f^2y$ 的 $x \in M^i$，把 $x$ 在 $Z^i$ 中的类
送到 $(\eta_fM)^i / f(\eta_fM)^i$ 中 $f^ix$ 的类。我们略去验证
这是良定义的；引理的假设恰好意味着此运算的定义域是整个
$Z^i$。于是 $t \circ t^{-1} = \text{id}_{Z^i}$。故 $t^{-1}$
定义了评注 \ref{more-algebra-remark-eta-BZ} 中短正合列的一个分裂，所得直和分解
$$
(\eta_fM)^i / f(\eta_fM)^i = Z^i \oplus B^{i + 1}
$$
与引理中展示的映射相容。
\end{proof}

\begin{lemma}
\label{more-algebra-lemma-eta-third-property}
设 $A$ 为环，$f, g \in A$ 为非零因子。设 $M^\bullet$ 为
$A$-模复形，使 $fg$ 在所有 $M^i$ 上都是非零因子。则
$\eta_f\eta_gM^\bullet = \eta_{fg}M^\bullet$。
\end{lemma}

\begin{proof}
该陈述意味着，在次数 $i$，我们得到局部化
$M^i_{fg} = (M^i_g)_f$ 的同一子模。略去细节。
\end{proof}

\begin{lemma}
\label{more-algebra-lemma-eta-flat-base-change}
设 $A$ 为环，$f \in A$ 为非零因子。设 $A \to B$ 为平坦环映射，
并设
% P02-E0452
$g \in B$ 为 $f$ 的像。设 $M^\bullet$ 为无 $f$-挠 $A$-模的复形。
则 $g$ 是非零因子，$M^\bullet \otimes_A B$ 是无 $g$-挠模的复形，
并且
$\eta_fM^\bullet \otimes_A B = \eta_g(M^\bullet \otimes_A B)$。
\end{lemma}

\begin{proof}
略。
\end{proof}

\section{完美复形与 eta 算子}
\label{more-algebra-section-perfect-eta}

\noindent
本节将作若干代数准备，以用于我们版本的 Macpherson 图构造；见
《关于平坦性的进一步结果》第
\ref{flat-section-blowup-complexes-III} 节。我们将使用第
\ref{more-algebra-section-eta} 节引入的 $\eta_f$ 算子。

\medskip\noindent
设 $A$ 为环，$f \in A$ 为非零因子。设 $M^\bullet$ 为有限自由
$A$-模的有界复形。对每个 $i$，令 $r_i$ 为 $M^i$ 的秩，并置
$$
I_i(M^\bullet, f) = \text{ideal generated by the }
r_i \times r_i\text{-minors of }
(f, d^i) : M^i \to M^i \oplus M^{i + 1}
$$
。注意 $f^{r_i} \in I_i(M^\bullet, f)$。

\begin{lemma}
\label{more-algebra-lemma-ideal-well-defined}
设 $A$ 为环，$f \in A$ 为非零因子。设 $M^\bullet$ 与
$N^\bullet$ 为有限自由 $A$-模的两个有界复形，它们代表
$D(A)$ 中同一个对象。则
$$
f^m I_i(M^\bullet, f) = f^n I_i(N^\bullet, f)
$$
是 $A$ 的理想，其中整数 $n, m \geq 0$ 满足
$$
m + \sum\nolimits_{j \geq i} (-1)^{j - i}rk(M^j) =
n + \sum\nolimits_{j \geq i} (-1)^{j - i}rk(N^j)
$$
。
\end{lemma}

\begin{proof}
只需在 $A$ 的每个素理想处局部化后证明该等式。因此，由引理
\ref{more-algebra-lemma-compare-representatives-perfect} 以及我们略去的归纳论证，
% P02-E0453
可以假设 $N^\bullet = M^\bullet \oplus Q^\bullet$，其中
$Q^\bullet$ 为某个平凡复形，即
$$
Q^\bullet = \ldots \to 0 \to A \xrightarrow{1} A \to 0 \to \ldots
$$
，两个 $A$ 分别置于次数 $j$ 与 $j + 1$。若
$j \not = i - 1, i, i + 1$，则显然
$I_i(M^\bullet, f) = I_i(N^\bullet, f)$ 且 $m = n$，从而得到所需
等式。若 $j = i + 1$，则映射
$$
(f, d^i) : M^i \to M^i \oplus M^{i + 1}
\quad\text{且}\quad
(f, d^i, 0) : M^i \to M^i \oplus M^{i + 1} \oplus A
$$
具有相同的非零子式，故此时同样有
$I_i(M^\bullet, f) = I_i(N^\bullet, f)$ 且 $m = n$。若 $j = i$，
则 $I_i(M^\bullet, f)$ 是由映射
$$
(f, d^i) : M^i \to M^i \oplus M^{i + 1}
$$
的 $r_i \times r_i$ 子式生成的理想，而 $I_i(N^\bullet, f)$
是由映射
$$
(f \oplus f, d^i \oplus 1) :
(M^i \oplus A) \to (M^i \oplus A) \oplus (M^{i + 1} \oplus A)
$$
的 $(r_i + 1) \times (r_i + 1)$ 子式生成的理想。适当选取坐标后，
可见第二个映射的矩阵具有分块形式
% P02-E0454
$$
T =
\left(
\begin{matrix}
T_1 & 0 \\
0 & T_2
\end{matrix}
\right),
\quad
T_1 = \text{matrix of first map},
\quad
T_2 =
\left(
\begin{matrix}
f \\
1
\end{matrix}
\right)
$$
。沿用引理 \ref{more-algebra-lemma-ideals-generated-by-minors} 的记号，
$I_0(T_2) = A$，$I_1(T_2) = A$，且对 $p \geq 2$ 有
$I_p(T_2) = 0$，因而
$I_{r_i + 1}(T) = I_{r_i + 1}(T_1) + I_{r_i}(T_1) =
I_{r_i}(T_1)$；这意味着
$I_i(M^\bullet, f) = I_i(N^\bullet, f)$。又有 $m = n$，故
$j = i$ 的情形得证。

最后，设 $j = i - 1$。此时可见 $m = n + 1$，故须证明
$fI_i(M^\bullet, f) = I_i(N^\bullet, f)$。在此情形下，
$I_i(M^\bullet, f)$ 是由映射
$$
(f, d^i) : M^i \to M^i \oplus M^{i + 1}
$$
的 $r_i \times r_i$ 子式生成的理想，而 $I_i(N^\bullet, f)$
是由映射
$$
(f \oplus f, d^i) : (M^i \oplus A) \to (M^i \oplus A) \oplus M^{i + 1}
$$
的 $(r_i + 1) \times (r_i + 1)$ 子式生成的理想。适当选取坐标后，
可见第二个映射的矩阵具有分块形式
% P02-E0454
$$
T =
\left(
\begin{matrix}
T_1 & 0 \\
0 & T_2
\end{matrix}
\right),
\quad
T_1 = \text{matrix of first map},
\quad
T_2 =
\left(
\begin{matrix}
f
\end{matrix}
\right)
$$
。与上面同样论证，便得到
$fI_i(M^\bullet, f) = I_i(N^\bullet, f)$。
\end{proof}

\begin{lemma}
\label{more-algebra-lemma-eta-change-unit}
设 $f \in A$ 为环 $A$ 的非零因子，$u \in A$ 为单位。设
$M^\bullet$ 为有限自由 $A$-模的有界复形。则
$I_i(M^\bullet, f) = I_i(M^\bullet, uf)$。
\end{lemma}

\begin{proof}
略。
\end{proof}

\begin{lemma}
\label{more-algebra-lemma-eta-base-change-pre}
设 $A \to B$ 为环映射，$f \in A$ 为非零因子。设 $M^\bullet$
为有限自由 $A$-模的有界复形。假设 $f$ 在 $B$ 中的像 $g$
是非零因子。则
$I_i(M^\bullet, f)B = I_i(M^\bullet \otimes_A B, g)$。
\end{lemma}

\begin{proof}
映射 $(f, d^i) : M^i \to M^i \oplus M^{i + 1}$ 的各子式映到
相应映射
$(g, d^i) : M^i \otimes_A B \to M^i \otimes_A B \oplus M^{i + 1} \otimes_A B$
的对应子式。
\end{proof}

\begin{lemma}
\label{more-algebra-lemma-eta-cohomology-free}
设 $A$ 为环，$\mathfrak p \subset A$ 为素理想，且 $f \in A$
为非零因子。设 $M^\bullet$ 为有限自由 $A$-模的有界复形。若对所有
$i$，$H^i(M^\bullet)_\mathfrak p$ 都自由，则对所有 $i$，
$I_i(M^\bullet, f)_\mathfrak p$ 都是主理想；事实上，它由 $f$
的一个幂生成。
\end{lemma}

\begin{proof}
由引理 \ref{more-algebra-lemma-eta-base-change-pre}，可以假设 $A$ 为极大理想
$\mathfrak p$ 下的局部环。由引理 \ref{more-algebra-lemma-ideal-well-defined}，
还可用一个拟同构复形替换 $M^\bullet$。根据上同调模自由这一假设，
可见 $M^\bullet$ 拟同构于这样的复形：其次数 $i$ 项是
$H^i(M^\bullet)$，且微分全为零；例如见《导出范畴》引理
\ref{derived-lemma-ext-2-zero-pre}。换言之，可以假设复形
$M^\bullet$ 的所有微分都为零。此时显然
$I_i(M^\bullet, f) = (f^{r_i})$ 是主理想。
\end{proof}

\begin{lemma}
\label{more-algebra-lemma-eta-locally-free}
设 $A$ 为环，$f \in A$ 为非零因子。设 $M^\bullet$ 为有限自由
$A$-模的有界复形。假设 $I_i(M^\bullet, f)$ 是主理想。则
$(\eta_fM)^i$ 是秩为 $r_i$ 的局部自由模，且映射
$(1, d^i) : (\eta_fM)^i \to f^iM^i \oplus f^{i + 1}M^{i + 1}$
是一个直和项的嵌入。
\end{lemma}

\begin{proof}
选取 $I_i(M^\bullet, f)$ 的一个生成元 $g$。由于
$f^{r_i} \in I_i(M^\bullet, f)$，可见 $g$ 整除 $f$ 的某个幂。
特别地，$g$ 是 $A$ 中的非零因子。映射
$(f, d^i) : M^i \to M^i \oplus M^{i + 1}$ 的
$r_i \times r_i$ 子式生成理想 $I_i(M^\bullet, f)$，而
$(f, d^i)$ 的 $(r_i + 1) \times (r_i + 1)$ 子式为零：这可在
$f$ 处局部化后检验，因为局部化后该映射的秩等于 $r_i$。考虑满射
$$
M^i \oplus M^{i + 1}
\longrightarrow
Q = \Coker(f, d^i)/g\text{-torsion}
$$
。由引理 \ref{more-algebra-lemma-fitting-ideals-and-pd1}，模 $Q$ 是秩为
$r_{i + 1}$ 的有限局部自由模。因此 $Q$ 无 $f$-挠，从而
$(f, d^i)$ 的余核模去 $f$-幂挠后同样为 $Q$。

\medskip\noindent
考虑有限自由 $A$-模复形
$$
0 \to f^{i + 1}M^i \xrightarrow{1, d^i} f^iM^i \oplus f^{i + 1}M^{i + 1}
\xrightarrow{d^i, -1} f^iM^{i + 1} \to 0
$$
；它在 $f$ 处局部化后成为分裂正合复形。映射
$(1, d^i) : f^{i + 1}M^i \to f^iM^i \oplus f^{i + 1}M^{i + 1}$
同构于上面研究过的映射
$(f, d^i) : M^i \to M^i \oplus M^{i + 1}$。因此像
$$
Q' = \Im(f^iM^i \oplus f^{i + 1}M^{i + 1} \xrightarrow{d^i, -1} f^iM^{i + 1})
$$
同构于 $Q$，特别地是投射模。
% P02-E0455
另一方面，由第 \ref{more-algebra-section-eta} 节中 $\eta_f$ 的构造，
单射
$(1, d^i) : (\eta_fM)^i \to f^iM^i \oplus f^{i + 1}M^{i + 1}$
的像是 $(d^i, -1)$ 的核。于是得到同构
$(\eta_fM)^i \oplus Q' = f^iM^i \oplus f^{i + 1}M^{i + 1}$，
并且可见
% P02-E0456
$\eta_fM^i$ 确为秩 $r_i$ 的有限局部自由模，而 $(1, d^i)$
是一个直和项的嵌入。
\end{proof}

\begin{lemma}
\label{more-algebra-lemma-eta-base-change}
设 $A \to B$ 为环映射，$f \in A$ 为非零因子。设 $M^\bullet$
为有限自由 $A$-模的有界复形。假设 $f$ 在 $B$ 中的像 $g$
是非零因子，并且对所有 $i \in \mathbf{Z}$，
$I_i(M^\bullet, f)$ 都是主理想。则有典范同构
$\eta_fM^\bullet \otimes_A B = \eta_g(M^\bullet \otimes_A B)$。
\end{lemma}

\begin{proof}
置 $N^i = M^i \otimes_A B$。注意，作为 $(N^i)_g$ 的子模，
$f^iM^i \otimes_A B = g^iN^i$。映射
% P02-E0457
$$
(\eta_fM)^i \otimes_A B \to g^iN^i \otimes g^{i + 1}N^{i + 1}
\quad\text{且}\quad
(\eta_gN)^i \to g^iN^i \otimes g^{i + 1}N^{i + 1}
$$
由引理 \ref{more-algebra-lemma-eta-locally-free} 都是直和项的嵌入。它们的像在
$g$ 处局部化后相同，结论随即得出。
\end{proof}

\begin{lemma}
\label{more-algebra-lemma-ideal-direct-summand}
设 $A$ 为环，$M$、$N_1$、$N_2$ 为有限投射 $A$-模。设
$s : M \to N_1 \oplus N_2$ 为分裂单射。存在具有以下性质的有限生成
理想 $J \subset A$：环映射 $A \to B$ 经过 $A/J$ 分解，当且仅当
$s \otimes \text{id}_B$ 把 $M \otimes_A B$ 同一为
$N_1 \otimes_A B \oplus N_2 \otimes_A B$ 的若干子模的直和。
\end{lemma}

\begin{proof}
选取 $s$ 的一个分裂
$\pi : N_1 \oplus N_2 \to M$。
% P02-E0458
以 $q_i : N_1 \oplus N_2 \to N_1 \oplus N_2$ 表示到 $N_i$
上的投影算子。置 $p_i = \pi \circ q_i \circ s$。注意
$p_1 + p_2 = \text{id}_M$。我们断言，$M$ 是
$N_1 \oplus N_2$ 的若干子模的直和，当且仅当 $p_1$ 与 $p_2$
是正交投影算子。因此，$J$ 是 $A$ 的满足下列条件的最小理想：
$p_1 \circ p_1 - p_1$、$p_2 \circ p_2 - p_2$、
$p_1 \circ p_2$ 与
$p_2 \circ p_1$ 都包含在
$J \otimes_A \text{End}_A(M)$ 中。略去一些细节。
\end{proof}

\noindent
设 $A$ 为环，$f \in A$ 为非零因子。设 $M^\bullet$ 为有限自由
$A$-模的有界复形。假设对所有 $i \in \mathbf{Z}$，理想
$I_i(M^\bullet, f)$ 都是主理想。于是，由引理
\ref{more-algebra-lemma-eta-locally-free}，映射
$$
(1, d^i) : (\eta_fM)^i / f(\eta_fM)^i \longrightarrow
f^iM^i/f^{i + 1}M^i \oplus f^{i + 1}M^{i + 1}/f^{i + 2}M^{i + 1}
$$
是分裂单射。
% P02-E0459
以 $J_i(M^\bullet, f) \subset A/fA$ 表示引理
\ref{more-algebra-lemma-ideal-direct-summand} 中对应于上面所示分裂单射
$(1, d^i)$ 的有限生成理想。

\begin{lemma}
\label{more-algebra-lemma-eta-vanishing-beta-plus-pre}
设 $A$ 为环，$f \in A$ 为非零因子。设 $M^\bullet$ 与
$N^\bullet$ 为有限自由 $A$-模的两个有界复形，它们代表
$D(A)$ 中同一个对象。假设对所有 $i \in \mathbf{Z}$，
$I_i(M^\bullet, f)$ 都是主理想。则作为 $A/fA$ 的理想，
$J_i(M^\bullet, f) = J_i(N^\bullet, f)$。
\end{lemma}

\begin{proof}
注意，
% P02-E0460
$I_i(M^\bullet, f)$ 是主理想这一事实，由引理
\ref{more-algebra-lemma-ideal-well-defined} 蕴含 $I_i(M^\bullet, f)$ 是主理想，
故该陈述有意义。如同引理 \ref{more-algebra-lemma-ideal-well-defined} 的证明，
可以假设 $N^\bullet = M^\bullet \oplus Q^\bullet$，其中
$Q^\bullet$ 为某个平凡复形，即
$$
Q^\bullet = \ldots \to 0 \to A \xrightarrow{1} A \to 0 \to \ldots
$$
，两个 $A$ 分别置于次数 $j$ 与 $j + 1$。由于 $\eta_f$ 与直和相容，
可见映射
$$
(1, d^i) : (\eta_fN)^i / f(\eta_fN)^i \longrightarrow
f^iN^i/f^{i + 1}N^i \oplus f^{i + 1}N^{i + 1}/f^{i + 2}N^{i + 1}
$$
是 $M^\bullet$ 的相应映射与 $Q^\bullet$ 的相应映射的直和。由定义
所论理想的泛性质，得到
$J_i(N^\bullet, f) = J_i(M^\bullet, f) + J_i(Q^\bullet, f)$。
因此，只需证明对所有 $i$ 都有 $J_i(Q^\bullet, f) = 0$。这是一个
我们略去的计算。
\end{proof}

\begin{lemma}
\label{more-algebra-lemma-eta-vanishing-beta-plus}
设 $A$ 为环，$f \in A$ 为非零因子。设 $M^\bullet$ 为有限自由
$A$-模的有界复形。假设对所有 $i \in \mathbf{Z}$，
$I_i(M^\bullet, f)$ 都是主理想。考虑 $A/fA$ 的理想
$J(M^\bullet, f) = \sum_i J_i(M^\bullet, f)$。考虑素理想集合
\begin{align*}
E
& =
\{f \in \mathfrak p \subset A \mid
\Ker(d^i \bmod f^2)_\mathfrak p \text{ surjects onto }
\Ker(d^i \bmod f)_\mathfrak p \text{ 对所有 }i \in \mathbf{Z}\} \\
& =
\{f \in \mathfrak p \subset A \mid
\text{the localizations }\beta_\mathfrak p
\text{ of the Bockstein operators are zero}\}
\end{align*}
。则有：
\begin{enumerate}
\item $J(M^\bullet, f)$ 是有限生成的；
\item $A/fA \to C = (A/fA)/J(M^\bullet, f)$ 是有限表示满射；
\item 对 $\mathfrak p \in E$，有
$J(M^\bullet, f)_\mathfrak p = 0$；
\item 若 $f \in \mathfrak p$，且对所有 $i \in \mathbf{Z}$，
$H^i(M^\bullet)_\mathfrak p$ 都自由，则 $\mathfrak p \in E$；以及
\item $\eta_f M^\bullet \otimes_A C$ 的各上同调模都是有限局部自由
$C$-模。
\end{enumerate}
\end{lemma}

\begin{proof}
$E$ 定义中两个集合的相等性由引理 \ref{more-algebra-lemma-vanishing-beta} 得出；
此外，该引理的最后一个陈述蕴含 (4)。

\medskip\noindent
(1) 成立，因为理想 $J_i(M^\bullet, f)$ 都有限生成，而且
$M^\bullet$ 有界，故几乎所有 $i$ 对应的 $J_i(M^\bullet, f)$ 都为零。
(2) 只是 (1) 的改述。

\medskip\noindent
证明 (3)。由引理 \ref{more-algebra-lemma-eta-locally-free}，对所有 $i$，
$(\eta_fM)^i$ 都是秩为 $r_i$ 的有限局部自由模。考虑映射
$$
(1, d^i) :
(\eta_fM)^i / f(\eta_fM)^i
\longrightarrow
f^iM^i/f^{i + 1}M^i \oplus f^{i + 1}M^{i + 1}/f^{i + 2}M^{i + 1}
$$
。取 $\mathfrak p \in E$。由引理
\ref{more-algebra-lemma-eta-vanishing-beta} 以及各模 $(\eta_fM)^i$ 的局部自由性，
可以写成
$$
\left((\eta_fM)^i / f(\eta_fM)^i\right)_\mathfrak p =
(A/fA)_\mathfrak p^{\oplus m_i} \oplus (A/fA)_\mathfrak p^{\oplus n_i}
$$
，并与上面的箭头 $(1, d^i)$ 相容。由理想
$J_i(M^\bullet, f)$ 的泛性质，得到
$J_i(M^\bullet, f)_\mathfrak p = 0$。因此，
% P02-E0461
$I_\mathfrak p = fA_\mathfrak p$ 对 $\mathfrak p \in E$ 成立。

\medskip\noindent


\medskip\noindent
证明 (5)。注意，$\eta_fM^\bullet$ 上的微分嵌入交换图
% P02-E0462
$$
\xymatrix{
(\eta_fM)^i \ar[d] \ar[r] &
f^iM^i \oplus f^{i + 1}M^{i + 1}
\ar[d]^{\left(
\begin{matrix}
0 & 1 \\
0 & 0
\end{matrix}
\right)} \\
(\eta_fM)^{i + 1} \ar[r] &
f^{i + 1}M^i \oplus f^{i + 2}M^{i + 2}
}
$$
。由构造，在与 $C$ 张量后，左边各模是右边各直和项的若干直和项
之直和。考察图
% P02-E0462
$$
\xymatrix{
(\eta_fM)^i \otimes_A C \ar[d] \ar@{=}[r] &
K^i \oplus L^i \ar[r] \ar[d] &
f^iM^i \otimes_A C \oplus f^{i + 1}M^{i + 1} \otimes_A C
\ar[d]^{\left(
\begin{matrix}
0 & 1 \\
0 & 0
\end{matrix}
\right)} \\
(\eta_fM)^{i + 1} \otimes_A C \ar@{=}[r] &
K^{i + 1} \oplus L^{i + 1} \ar[r] &
f^{i + 1}M^i \otimes_A C \oplus f^{i + 2}M^{i + 2} \otimes_A C
}
$$
，其中各水平箭头既与直和分解相容，又是直和项的嵌入。由此，微分把
$L^i$ 同一为 $K^{i + 1}$ 的一个直和项，因而
$\eta_fM^\bullet \otimes_A C$ 在次数 $i$ 的上同调为模
% P02-E0463
$K^{i + 1}/L^i$；它按要求是有限投射模。
\end{proof}

\section{复形的极限}
\label{more-algebra-section-limits}

\noindent
本节讨论一个复形具有“形式变形”并对其取极限时会发生什么。我们将
考虑两种情形：
\begin{enumerate}
\item 有环的一个逆系统的极限 $A = \lim A_n$，其过渡映射为满射，
核局部幂零；并有对象 $K_n \in D(A_n)$，它们按
$K_n = K_{n + 1} \otimes_{A_{n + 1}}^\mathbf{L} A_n$ 的意义相容；或者
\item 有环 $A$、理想 $I$ 以及对象 $K_n \in D(A/I^n)$，它们按
$K_n = K_{n + 1} \otimes_{A/I^{n + 1}}^\mathbf{L} A/I^n$ 的意义相容。
\end{enumerate}
在附加假设下，可以证明 $K = R\lim K_n$ 重现该系统，即
$K_n = K \otimes_A^\mathbf{L} A_n$ 或
$K_n = K \otimes_A^\mathbf{L} A/I^n$。

\begin{lemma}
\label{more-algebra-lemma-Rlim-pseudo-coherent-gives-pseudo-coherent}
设 $A = \lim A_n$ 为环的一个逆系统 $(A_n)$ 的极限。给定
$K_n \in D(A_n)$ 以及 $D(A_{n + 1})$ 中的映射
$K_{n + 1} \to K_n$。假设：
\begin{enumerate}
\item 过渡映射 $A_{n + 1} \to A_n$ 是满射且核局部幂零；
\item 或者所有 $K_n$ 都伪凝聚，或者存在整数 $n_0$，使得
$K_{n_0}$ 伪凝聚，且当 $n \geq n_0$ 时，
$A_{n + 1} \to A_n$ 的核是幂零理想；
\item 这些映射诱导同构
$K_{n + 1} \otimes_{A_{n + 1}}^\mathbf{L} A_n \to K_n$。
\end{enumerate}
则 $K = R\lim K_n$ 是 $D(A)$ 的伪凝聚对象，并且对所有 $n$，
$K \otimes_A^\mathbf{L} A_n \to K_n$ 都是同构。
\end{lemma}

\begin{proof}
由假设，可找到一个代表 $K_1$ 的有限自由 $A_1$-模上有界复形
$P_1^\bullet$；见定义 \ref{more-algebra-definition-pseudo-coherent}。由引理
\ref{more-algebra-lemma-check-pseudo-coherent-modulo-nilpotent}，得到对所有 $n$，
$K_n$ 都伪凝聚。于是，由引理 \ref{more-algebra-lemma-lift-complex-stably-frees}，
对 $n > 1$ 归纳可找到代表 $K_n$ 的有限自由 $A_n$-模复形
$P_n^\bullet$，以及代表映射 $K_n \to K_{n - 1}$ 的映射
$P_n^\bullet \to P_{n - 1}^\bullet$；它们诱导复形同构（！）
$P_n^\bullet \otimes_{A_n} A_{n - 1} \to P_{n - 1}^\bullet$。
因此，由引理 \ref{more-algebra-lemma-compute-Rlim-modules} 与评注
\ref{more-algebra-remark-how-unique}，
$K = R\lim K_n$ 由 $P^\bullet = \lim P_n^\bullet$ 代表。
对每个 $n$，$P_n^i$ 是有限自由 $A_n$-模，又因
$A = \lim A_n$，可见对每个 $i$，$P^i$ 都是与 $P_1^i$ 同秩的
有限自由模。这意味着 $K$ 伪凝聚。又有
$K \otimes_A^\mathbf{L} A_n$ 由
$P^\bullet \otimes_A A_n = P_n^\bullet$ 代表，从而证明最后一个断言。
\end{proof}

\begin{lemma}
\label{more-algebra-lemma-Rlim-pseudo-coherent-gives-complete-pseudo-coherent}
设 $A$ 为环，$I \subset A$ 为理想。给定 $K_n \in D(A/I^n)$
以及 $D(A/I^{n + 1})$ 中的映射 $K_{n + 1} \to K_n$。假设：
\begin{enumerate}
\item $A$ 是 $I$-进完备的；
\item $K_1$ 伪凝聚；且
\item 这些映射诱导同构
$K_{n + 1} \otimes_{A/I^{n + 1}}^\mathbf{L} A/I^n \to K_n$。
\end{enumerate}
则 $K = R\lim K_n$ 是 $D(A)$ 的伪凝聚导出完备对象，并且对所有
$n$，$K \otimes_A^\mathbf{L} A/I^n \to K_n$ 都是同构。
\end{lemma}

\begin{proof}
由引理 \ref{more-algebra-lemma-Rlim-pseudo-coherent-gives-pseudo-coherent}，已经知道
$K$ 伪凝聚，并且对所有 $n$，
$K \otimes_A^\mathbf{L} A/I^n \to K_n$ 都是同构。最后，由引理
\ref{more-algebra-lemma-naive-derived-completion}，$K$ 导出完备。
\end{proof}

\begin{lemma}
\label{more-algebra-lemma-Rlim-perfect-gives-perfect}
\begin{reference}
\cite[Lemma 4.2]{Bhatt-Algebraize}
\end{reference}
设 $A = \lim A_n$ 为环的一个逆系统 $(A_n)$ 的极限。给定
$K_n \in D(A_n)$ 以及 $D(A_{n + 1})$ 中的映射
$K_{n + 1} \to K_n$。假设：
\begin{enumerate}
\item 过渡映射 $A_{n + 1} \to A_n$ 是满射且核局部幂零；
\item 或者所有 $K_n$ 都完美，或者存在整数 $n_0$，使得
$K_{n_0}$ 完美，且当 $n \geq n_0$ 时，
$A_{n + 1} \to A_n$ 的核为幂零理想；且
\item 这些映射诱导同构
$K_{n + 1} \otimes_{A_{n + 1}}^\mathbf{L} A_n \to K_n$。
\end{enumerate}
则 $K = R\lim K_n$ 是 $D(A)$ 的完美对象，并且对所有 $n$，
$K \otimes_A^\mathbf{L} A_n \to K_n$ 都是同构。
\end{lemma}

\begin{proof}
由引理 \ref{more-algebra-lemma-Rlim-pseudo-coherent-gives-pseudo-coherent}，已经知道
$K$ 伪凝聚，并且对所有 $n$，
$K \otimes_A^\mathbf{L} A_n \to K_n$ 都是同构。考虑一个满射
$A \to \kappa$，其核为极大理想 $\mathfrak m$。由《代数》引理
\ref{algebra-lemma-locally-nilpotent-unit}，$A$ 中任何在 $A_1$ 中
映为单位的元素在 $A$ 中都是单位；因而，由《代数》引理
\ref{algebra-lemma-contained-in-radical}，
$\Ker(A \to A_1)$ 包含在 $A$ 的 Jacobson 根中。故
$A \to \kappa$ 分解为 $A \to A_1 \to \kappa$。于是
$$
K \otimes_A^\mathbf{L} \kappa =
K \otimes_A^\mathbf{L} A_1 \otimes_{A_1}^\mathbf{L} \kappa =
K_1 \otimes_{A_1}^\mathbf{L} \kappa
$$
。
% P02-E0464
注意，由假设 (2)，$K_1$ 完美。因此，由引理
\ref{more-algebra-lemma-perfect}，$K_1$ 具有有限 Tor 维数。故存在
$a, b \in \mathbf{Z}$，使得对所有 $i \not \in [a, b]$，
$H^i(K \otimes_A^\mathbf{L} \kappa) = 0$。由引理
\ref{more-algebra-lemma-check-perfect-pointwise}，得到 $K$ 完美。
\end{proof}

\begin{lemma}
\label{more-algebra-lemma-Rlim-perfect-gives-complete}
设 $A$ 为环，$I \subset A$ 为理想。给定 $K_n \in D(A/I^n)$
以及 $D(A/I^{n + 1})$ 中的映射 $K_{n + 1} \to K_n$。假设：
\begin{enumerate}
\item $A$ 是 $I$-进完备的；
\item $K_1$ 是完美对象；且
\item 这些映射诱导同构
$K_{n + 1} \otimes_{A/I^{n + 1}}^\mathbf{L} A/I^n \to K_n$。
\end{enumerate}
则 $K = R\lim K_n$ 是 $D(A)$ 的完美导出完备对象，并且对所有
$n$，$K \otimes_A^\mathbf{L} A/I^n \to K_n$ 都是同构。
\end{lemma}

\begin{proof}
合并引理 \ref{more-algebra-lemma-Rlim-perfect-gives-perfect} 与
\ref{more-algebra-lemma-Rlim-pseudo-coherent-gives-complete-pseudo-coherent}
（后者给出导出完备性）。
\end{proof}

\noindent
我们不知道下一个引理对无界复形是否成立。

\begin{lemma}
\label{more-algebra-lemma-Rlim-gives-complete}
设 $A$ 为环，$I \subset A$ 为理想。给定 $K_n \in D(A/I^n)$
以及 $D(A/I^{n + 1})$ 中的映射 $K_{n + 1} \to K_n$。若
\begin{enumerate}
\item $A$ 是 Noether 环；
\item $K_1$ 上有界；且
\item 这些映射诱导同构
$K_{n + 1} \otimes_{A/I^{n + 1}}^\mathbf{L} A/I^n \to K_n$，
\end{enumerate}
则 $K = R\lim K_n$ 是 $D^-(A)$ 的导出完备对象，并且对所有 $n$，
$K \otimes_A^\mathbf{L} A/I^n \to K_n$ 都是同构。
\end{lemma}

\begin{proof}
由引理 \ref{more-algebra-lemma-naive-derived-completion}，$D(A)$ 的对象 $K$
导出完备。

\medskip\noindent
假设当 $i > b$ 时 $H^i(K_1) = 0$。则可找到代表 $K_1$ 的自由
$A/I$-模复形 $P_1^\bullet$，且当 $i > b$ 时 $P_1^i = 0$。由引理
\ref{more-algebra-lemma-lift-complex-projectives}，对 $n > 1$ 归纳可找到代表
$K_n$ 的自由 $A/I^n$-模复形 $P_n^\bullet$，以及代表映射
$K_n \to K_{n - 1}$ 的映射
$P_n^\bullet \to P_{n - 1}^\bullet$；它们诱导复形同构（！）
$P_n^\bullet/I^{n - 1}P_n^\bullet \to P_{n - 1}^\bullet$。

\medskip\noindent
于是得到如下情形：$R\lim K_n$ 由
$P^\bullet = \lim P_n^\bullet$ 代表；见引理
\ref{more-algebra-lemma-compute-Rlim-modules} 与评注 \ref{more-algebra-remark-how-unique}。
各复形 $P_n^\bullet$ 是平坦 $A/I^n$-模的一致上有界复形，且过渡映射
逐项满射。于是，由引理 \ref{more-algebra-lemma-limit-flat}，$P^\bullet$
是平坦 $A$-模的上有界复形。因此
$K \otimes_A^\mathbf{L} A/I^t$ 由
$P^\bullet \otimes_A A/I^t$ 代表。由引理 \ref{more-algebra-lemma-limit-flat}，
逐项有
$P^\bullet \otimes_A A/I^t = \lim P_n^\bullet \otimes_A A/I^t$。
根据对 $P_n^\bullet$ 的选取，当 $n \geq t$ 时，过渡映射
$P_{n + 1}^\bullet \otimes_A A/I^t \to P_n^\bullet \otimes_A A/I^t$
都是同构，因而
$\lim P_n^\bullet \otimes_A A/I^t =
P_t^\bullet \otimes_A A/I^t = P_t^\bullet$。由于 $P_t^\bullet$
代表 $K_t$，可见
$K \otimes_A^\mathbf{L} A/I^t \to K_t$ 是同构。
\end{proof}

\noindent
下面给出另一类结果。

\begin{lemma}[Koll\'ar-Kov\'acs]
\label{more-algebra-lemma-kollar-kovacs}
\begin{reference}
% P02-E0465
Kovacs 于 23/02/2018 发来的电子邮件。
\end{reference}
设 $I$ 为 Noether 环 $A$ 的理想，$K \in D(A)$。置
$K_n = K \otimes_A^\mathbf{L} A/I^n$。假设对所有
$i \in \mathbf{Z}$ 都有：
\begin{enumerate}
\item $H^i(K)$ 是有限 $A$-模；且
\item 系统 $H^i(K_n)$ 满足 Mittag-Leffler 性质。
\end{enumerate}
则对所有 $i \in \mathbf{Z}$，
$\lim H^i(K)/I^nH^i(K)$ 等于 $\lim H^i(K_n)$。
\end{lemma}

\begin{proof}
回忆，$K^\wedge = R\lim K_n$ 是 $K$ 的导出完备化；见命题
\ref{more-algebra-proposition-noetherian-naive-completion-is-completion}。由引理
\ref{more-algebra-lemma-derived-completion-pseudo-coherent}，有
$H^i(K^\wedge) = \lim H^i(K)/I^nH^i(K)$。由引理
\ref{more-algebra-lemma-break-long-exact-sequence-modules}，得到短正合列
% P02-E0464
$$
0 \to R^1\lim H^{i - 1}(K_n) \to H^i(K^\wedge) \to \lim H^i(K_n) \to 0
$$
。Mittag-Leffler 性质保证左端项为零（引理
\ref{more-algebra-lemma-compute-Rlim-modules}），故引理得证。
\end{proof}

\section{若干求值映射}
\label{more-algebra-section-evaluation}

\noindent
本节证明，对适当类型的复形，某些典范的 $R\Hom$ 映射是同构。

\begin{lemma}
\label{more-algebra-lemma-internal-hom-evaluate-isomorphism}
设 $R$ 为环，$K, L, M$ 为 $D(R)$ 的对象。
% P02-E0466
引理 \ref{more-algebra-lemma-internal-hom-evaluate} 的映射
$$
R\Hom_R(L, M) \otimes_R^\mathbf{L} K \longrightarrow R\Hom_R(R\Hom_R(K, L), M)
$$
在以下两种情形中是同构：
\begin{enumerate}
\item $K$ 完美；或者
\item $K$ 伪凝聚，$L \in D^+(R)$，并且
% P02-E0467
$M$ 具有有限内射维数。
\end{enumerate}
\end{lemma}

\begin{proof}
选取代表 $M$ 的 K-内射复形 $I^\bullet$、代表 $L$ 的 K-内射复形
$J^\bullet$，以及代表 $K$ 的有限投射模上有界复形 $K^\bullet$。
考虑引理 \ref{more-algebra-lemma-evaluate-and-more} 的复形映射
$$
\text{Tot}(\Hom^\bullet(J^\bullet, I^\bullet) \otimes_R K^\bullet)
\longrightarrow
\Hom^\bullet(\Hom^\bullet(K^\bullet, J^\bullet), I^\bullet)
$$
。注意
$$
\left(\prod\nolimits_{p + r = t} \Hom_R(J^{-r}, I^p)\right) \otimes_R K^s =
\prod\nolimits_{p + r = t} \Hom_R(J^{-r}, I^p) \otimes_R K^s
$$
，因为 $K^s$ 有限投射。该映射由映射
$$
c_{p, r, s} :
\Hom_R(J^{-r}, I^p) \otimes_R K^s
\longrightarrow
\Hom_R(\Hom_R(K^s, J^{-r}), I^p)
$$
给出；由于 $K^s$ 有限投射，它们都是同构。对于左端次数 $n$
的任意元素 $\alpha = (\alpha^{p, r, s})$，只有有限多个 $s$ 的值
使 $\alpha^{p, r, s}$ 非零（这里对某些满足
$n = p + r + s$ 的 $p, r$ 而言）。因此，若右端元素
$\beta = (\beta^{p, r, s})$ 也被迫满足同一消失条件，则该映射是同构。
若 $K^\bullet$ 是有限投射模的有界复形，这一点显然。另一方面，若可取
$I^\bullet$ 有界而 $J^\bullet$ 下有界，则当 $p$ 落在某个固定范围之外、
$s \gg 0$ 或 $r \gg 0$ 时，$\beta^{p, r, s}$ 为零。因此，在满足
$n = p + r + s$ 且 $\beta^{p, r, s}$ 非零的解中，只出现有限多个
$s$ 的值。
\end{proof}

\begin{lemma}
\label{more-algebra-lemma-internal-hom-evaluate-isomorphism-technical}
设 $R$ 为环，$K, L, M$ 为 $D(R)$ 的对象。
% P02-E0466
若以下三个条件成立，则引理 \ref{more-algebra-lemma-internal-hom-evaluate} 的映射
$$
R\Hom_R(L, M) \otimes_R^\mathbf{L} K \longrightarrow R\Hom_R(R\Hom_R(K, L), M)
$$
是同构：
\begin{enumerate}
\item $L, M$ 具有有限内射维数；
\item $R\Hom_R(L, M)$ 具有有限 Tor 维数；
% P02-E0468
\item 对每个 $n \in \mathbf{Z}$，截断 $\tau_{\leq n}K$ 都伪凝聚。
\end{enumerate}
\end{lemma}

\begin{proof}
选取整数 $n$，并考虑区分三角
% P02-E0469
$$
\tau_{\leq n}K \to K \to \tau_{\geq n + 1}K \to \tau_{\leq n}K[1]
$$
；见《导出范畴》评注
\ref{derived-remark-truncation-distinguished-triangle}。由假设 (3)
与引理 \ref{more-algebra-lemma-internal-hom-evaluate-isomorphism}，该映射对
$\tau_{\leq n}K$ 是同构。因此，只需证明
$$
R\Hom_R(L, M) \otimes_R^\mathbf{L} \tau_{\geq n + 1}K
\quad\text{且}\quad
R\Hom_R(R\Hom_R(\tau_{\geq n + 1}K, L), M)
$$
的上同调在次数 $\leq n - c$ 中消失，其中 $c$ 为某个整数。这立即由
假设 (2) 与 (1) 得出。
\end{proof}

\begin{lemma}
\label{more-algebra-lemma-internal-hom-evaluate-tensor-isomorphism}
设 $R$ 为环，$K, L, M$ 为 $D(R)$ 的对象。引理
\ref{more-algebra-lemma-internal-hom-diagonal-better} 的映射
$$
K \otimes_R^\mathbf{L} R\Hom_R(M, L) \longrightarrow
R\Hom_R(M, K \otimes_R^\mathbf{L} L)
$$
在以下情形中是同构：
\begin{enumerate}
\item $M$ 完美；或者
\item $K$ 完美；或者
\item $M$ 伪凝聚，$L \in D^+(R)$，并且 $K$ 的 Tor 振幅位于
% P02-E0470
$[a, \infty]$ 中。
\end{enumerate}
\end{lemma}

\begin{proof}
先证明 $M$ 完美的情形。注意，该箭头的两端都把 $M$ 中的区分三角
变为区分三角，并与直和交换。因此，只需检验 $M = R[n]$ 时该断言
成立；见《导出范畴》评注
\ref{derived-remark-check-on-generator} 与引理
\ref{more-algebra-lemma-perfect-ring-classical-generator}。在此情形下，结果显然。

\medskip\noindent
再证明 $K$ 完美的情形。论证与前一种情形相同。

\medskip\noindent
最后证明 (3)。可以用 $R$-模下有界复形 $K^\bullet$ 与
$L^\bullet$ 分别代表 $K$ 与 $L$。可以假设 $K^\bullet$ 是由平坦
$R$-模组成的 K-平坦复形；见引理
\ref{more-algebra-lemma-bounded-below-tor-amplitude}。可以用有限自由 $R$-模的
上有界复形 $M^\bullet$ 代表 $M$；见定义
\ref{more-algebra-definition-pseudo-coherent}。于是左端对象由
% P02-E0469
$$
\text{Tot}(K^\bullet \otimes_R \Hom^\bullet(M^\bullet, L^\bullet))
$$
代表，而右端对象由
$$
\Hom^\bullet(M^\bullet, \text{Tot}(K^\bullet \otimes_R L^\bullet))
$$
代表。这使用引理 \ref{more-algebra-lemma-RHom-out-of-projective}。两个复形在次数
$n$ 的模都是
$$
\bigoplus\nolimits_{p + q + r = n} K^p \otimes \Hom_R(M^{-r}, L^q) =
\bigoplus\nolimits_{p + q + r = n} \Hom_R(M^{-r}, K^p \otimes_R L^q)
$$
，因为 $M^{-r}$ 有限自由
% P02-E0471
（这些也都是有限直和）。引理
\ref{more-algebra-lemma-internal-hom-diagonal-better} 中定义的映射来自引理
\ref{more-algebra-lemma-diagonal-better} 中定义的复形映射；后者使用这些模之间的
典范同构。
\end{proof}

\begin{lemma}
\label{more-algebra-lemma-hom-complex-K-flat}
设 $R$ 为环，$P^\bullet$ 为投射 $R$-模的上有界复形，
$K^\bullet$ 为 $R$-模的 K-平坦复形。若 $P^\bullet$ 是 $D(R)$
的完美对象，则 $\Hom^\bullet(P^\bullet, K^\bullet)$ 是 K-平坦的，
并代表 $R\Hom_R(P^\bullet, K^\bullet)$。
\end{lemma}

\begin{proof}
最后一个陈述即引理 \ref{more-algebra-lemma-RHom-out-of-projective}。由于
$P^\bullet$ 代表完美对象，存在有限投射 $R$-模的有限复形
$F^\bullet$，使得 $P^\bullet$ 与 $F^\bullet$ 在 $D(R)$ 中同构；
见定义 \ref{more-algebra-definition-perfect}。于是，由《导出范畴》引理
\ref{derived-lemma-morphisms-from-projective-complex}，
$P^\bullet$ 与 $F^\bullet$ 伦等价。因此
$\Hom^\bullet(P^\bullet, K^\bullet)$ 与
$\Hom^\bullet(F^\bullet, K^\bullet)$ 伦等价。故前者 K-平坦当且仅当
后者 K-平坦（这由定义 \ref{more-algebra-definition-K-flat} 与引理
\ref{more-algebra-lemma-derived-tor-homotopy} 得出）。显然，
% P02-E0469
$$
\Hom^\bullet(F^\bullet, K^\bullet) =
\text{Tot}(E^\bullet \otimes_R K^\bullet)
$$
，其中 $E^\bullet$ 是 $F^\bullet$ 的对偶复形，其项为
$E^n = \Hom_R(F^{-n}, R)$；见引理
\ref{more-algebra-lemma-dual-perfect-complex} 及其证明。由于 $E^\bullet$
是投射模的有界复形，由引理
\ref{more-algebra-lemma-derived-tor-quasi-isomorphism}，它是 K-平坦的。最后由引理
\ref{more-algebra-lemma-tensor-product-K-flat} 得出结论。
\end{proof}

\section{导出 Hom 的基变换}
\label{more-algebra-section-base-change-RHom}

\noindent
关于这一点，我们已经在引理
\ref{more-algebra-lemma-pseudo-coherence-and-base-change-ext} 与《代数》第
\ref{algebra-section-functoriality-ext} 节中看到一些内容。

\begin{lemma}
\label{more-algebra-lemma-upgrade-adjoint-tensor-RHom}
设 $R \to R'$ 为环映射。对 $K \in D(R)$ 与 $M \in D(R')$，
存在典范同构
$$
R\Hom_R(K, M) = R\Hom_{R'}(K \otimes_R^\mathbf{L} R', M)
$$
。
\end{lemma}

\begin{proof}
选取代表 $M$ 的 $R'$-模 K-内射复形 $J^\bullet$。选取拟同构
$J^\bullet \to I^\bullet$，其中 $I^\bullet$ 是 $R$-模的 K-内射
复形。选取代表 $K$ 的 $R$-模 K-平坦复形 $K^\bullet$。考虑映射
% P02-E0472
$$
\Hom^\bullet(K^\bullet \otimes_R R', J^\bullet)
\longrightarrow
\Hom^\bullet(K^\bullet, I^\bullet)
$$
。次数 $n$ 项上的映射由映射
$$
\prod\nolimits_{n = p + q} \Hom_{R'}(K^{-q} \otimes_R R', J^p)
\longrightarrow
\prod\nolimits_{n = p + q} \Hom_R(K^{-q}, I^p)
$$
给出；后者来自与 $K^{-q} \to K^{-q} \otimes_R R'$ 预复合并与
$J^p \to I^p$ 后复合。为完成证明，只需证明在上同调群上得到同构：
% P02-E0472
$$
\Hom_{D(R)}(K, M) = \Hom_{D(R')}(K \otimes_R^\mathbf{L} R', M)
$$
。这成立是因为，由引理 \ref{more-algebra-lemma-tensor-hom-adjoint}，基变换
$- \otimes_R^\mathbf{L} R' : D(R) \to D(R')$ 左伴随于限制函子
$D(R') \to D(R)$。
\end{proof}

\noindent
设 $R \to R'$ 为环映射。有基变换映射
\begin{equation}
\label{more-algebra-equation-base-change-RHom}
R\Hom_R(K, M) \otimes_R^\mathbf{L} R'
\longrightarrow
R\Hom_{R'}(K \otimes_R^\mathbf{L} R', M \otimes_R^\mathbf{L} R')
\end{equation}
，它位于 $D(R')$ 中，并对 $K, M \in D(R)$ 函子性。事实上，由
$- \otimes_R^\mathbf{L} R' : D(R) \to D(R')$ 与限制函子
$D(R') \to D(R)$ 之间的伴随性，这等价于映射
$$
R\Hom_R(K, M)
\longrightarrow
R\Hom_{R'}(K \otimes_R^\mathbf{L} R', M \otimes_R^\mathbf{L} R') =
R\Hom_R(K, M \otimes_R^\mathbf{L} R')
$$
（等号由引理 \ref{more-algebra-lemma-upgrade-adjoint-tensor-RHom} 给出）；为此可使用
典范映射 $M \to M \otimes_R^\mathbf{L} R'$（伴随的单位）。

\begin{lemma}
\label{more-algebra-lemma-base-change-RHom}
设 $R \to R'$ 为环映射，$K, M \in D(R)$。映射
(\ref{more-algebra-equation-base-change-RHom})
$$
R\Hom_R(K, M) \otimes_R^\mathbf{L} R'
\longrightarrow
R\Hom_{R'}(K \otimes_R^\mathbf{L} R', M \otimes_R^\mathbf{L} R')
$$
在以下情形中是 $D(R')$ 中的同构：
\begin{enumerate}
\item $K$ 完美；
\item $R'$ 作为 $R$-模完美；
\item $R \to R'$ 平坦，$K$ 伪凝聚，且 $M \in D^{+}(R)$；或者
% P02-E0473
\item $R'$ 作为 $R$-模具有有限 Tor 维数，$K$ 伪凝聚，且
$M \in D^{+}(R)$。
\end{enumerate}
\end{lemma}

\begin{proof}
可以在应用限制函子 $D(R') \to D(R)$ 后检验该映射是同构。应用此
函子后，我们的映射成为
% P02-E0474
$$
R\Hom_R(K, L) \otimes_R^\mathbf{L} R'
\longrightarrow
R\Hom_R(K, L \otimes_R^\mathbf{L} R')
$$
，即引理 \ref{more-algebra-lemma-internal-hom-diagonal-better} 的映射。为匹配左右
两端，见该引理上方的讨论；特别地，这里使用引理
\ref{more-algebra-lemma-upgrade-adjoint-tensor-RHom}。于是由引理
\ref{more-algebra-lemma-internal-hom-evaluate-tensor-isomorphism} 得出结论。
\end{proof}

\section{模的系统}
\label{more-algebra-section-systems}

\noindent
设 $I$ 为 Noether 环 $A$ 的理想。本节将补充我们关于有限 $A$-模与
有限 $A/I^n$-模系统之间关系的认识。

\begin{lemma}
\label{more-algebra-lemma-consequence-Artin-Rees}
设 $I$ 为 Noether 环 $A$ 的理想。设
$
K \xrightarrow{\alpha} L \xrightarrow{\beta} M
$
为有限 $A$-模的复形。置 $H = \Ker(\beta)/\Im(\alpha)$。对
$n \geq 0$，设
$$
K/I^nK \xrightarrow{\alpha_n} L/I^nL \xrightarrow{\beta_n} M/I^nM
$$
为诱导复形。置 $H_n = \Ker(\beta_n)/\Im(\alpha_n)$。则存在典范
$A$-模映射，给出交换图
$$
\xymatrix{
& & & H \ar[lld] \ar[ld] \ar[d] \\
\ldots \ar[r] & H_3 \ar[r] & H_2 \ar[r] & H_1
}
$$
。此外，存在 $c > 0$，并且当 $n \geq c$ 时存在典范 $A$-模映射
$H_n \to H/I^{n - c}H$，使复合
$$
H/I^n H \to H_n \to  H/I^{n - c}H
\quad\text{且}\quad
H_n \to H/I^{n - c}H \to H_{n - c}
$$
都是典范映射。再者：
\begin{enumerate}
\item $(H_n)$ 与 $(H/I^nH)$ 作为 $\text{Mod}_A$ 的 pro-对象同构；
\item $\lim H_n = \lim H/I^n H$；
\item 逆系统 $(H_n)$ 满足 Mittag-Leffler 性质；
\item $H_{n + c} \to H_n$ 的像等于 $H \to H_n$ 的像；
\item 复合
$I^cH_n \to H_n \to H/I^{n - c}H \to H_n/I^{n - c}H_n$
是嵌入 $I^cH_n \to H_n$ 后接商映射
$H_n \to H_n/I^{n - c}H_n$；且
% P02-E0478
\item $H/I^nH \to H_n$ 的核与余核被 $I^c$ 零化。
\end{enumerate}
\end{lemma}

\begin{proof}
注意，
% P02-E0475
$H_n = \beta^{-1}(I^nM)/\Im(\alpha) + I^nL$。对 $n \geq 2$，有
$\beta^{-1}(I^nM) \subset \beta^{-1}(I^{n - 1}M)$ 以及
$\Im(\alpha) + I^nL \subset \Im(\alpha) + I^{n - 1}L$。因此得到典范
映射 $H_n \to H_{n - 1}$。类似地，
$\Ker(\beta) \subset \beta^{-1}(I^nM)$ 且
$\Im(\alpha) \subset \Im(\alpha) + I^nL$，从而产生典范映射
$H \to H_n$。我们略去该图交换的验证。

\medskip\noindent
由 Artin-Rees，可以选取 $c_1, c_2 \geq 0$，使得当
$n \geq c_1$ 时
$\beta^{-1}(I^nM) \subset \Ker(\beta) + I^{n - c_1}L$，且当
$n \geq c_2$ 时
$\Ker(\beta) \cap I^nL \subset I^{n - c_2}\Ker(\beta)$；见《代数》
引理 \ref{algebra-lemma-map-AR} 与
\ref{algebra-lemma-Artin-Rees}。置 $c = c_1 + c_2$。

\medskip\noindent
设 $n \geq c$。如下定义
$\psi_n : H_n \to H/I^{n - c}H$。取 $x \in H_n$，并选取代表 $x$ 的
$y \in \beta^{-1}(I^nM)$。写成 $y = z + w$，其中
$z \in \Ker(\beta)$ 且 $w \in I^{n - c_1}L$（由 $c_1$ 的选取可行）。
定义 $\psi_n(x)$ 为 $z$ 在 $H/I^{n - c}H$ 中的类。为证明这是良定义
的，假设对 $x$ 有另一组选取 $y', z', w'$，记号含义同上。此时
$y' - y \in \Im(\alpha) + I^nL$；写
$y' - y = \alpha(v) + u$，其中 $v \in K$ 且 $u \in I^nL$。于是
$$
y' = z' + w' = \alpha(v) + u + z + w
\Rightarrow
z' = z + \alpha(v) + u + w - w'
$$
。因为 $\beta(z' - z - \alpha(v)) = 0$，得到
$u + w - w' \in \Ker(\beta) \cap I^{n - c_1}L$；由 $c_2$ 的选取，
后者包含于
$I^{n - c_1 - c_2}\Ker(\beta) = I^{n - c}\Ker(\beta)$。故按要求，
$z'$ 与 $z$ 在 $H/I^{n - c}H$ 中有相同的像。

\medskip\noindent
复合 $H/I^n H \to H_n \to  H/I^{n - c}H$ 是典范映射。事实上，若
$z \in \Ker(\beta)$ 代表
% P02-E0475
$H/I^nH = \Ker(\beta)/\Im(\alpha) + I^n\Ker(\beta)$ 中的元素 $x$，
则由上面的构造，$x$ 在所得映射下映到 $z$ 在
$H/I^{n - c}H$ 中的类。

\medskip\noindent
考虑复合 $H_n \to H/I^{n - c}H \to H_{n - c}$。给定上面构造
$\psi_n$ 时的 $x, y, z, w$，可见 $x$ 映到 $z$ 在 $H_{n - c}$ 中的
% P02-E0476
类。另一方面，证明第一段的典范映射 $H_n \to H_{n - c}$ 把 $x$
送到 $y$ 的类。因此须证明
$y - z \in \Im(\alpha) + I^{n - c}L$；这成立，因为
$y - z = w \in I^{n - c_1}L \subset I^{n - c}L$。

\medskip\noindent
(1)--(4) 是刚才所证内容的形式推论。具体地，(1) 由这些映射的存在性
以及《范畴》评注 \ref{categories-remark-pro-category} 中 pro-对象
态射的定义得出。(2) 成立，因为同构的 pro-对象具有同构的极限。
(3) 立即由 (4) 得出。(4) 由典范映射 $H_{n + c} \to H_n$ 的分解
$H_{n + c} \to H/I^nH \to H_n$ 得出。

\medskip\noindent
证明 (5)。设 $x \in I^cH_n$。写成 $x = \sum f_i x_i$，其中
$x_i \in H_n$ 且 $f_i \in I^c$。对 $x_i$，按 $\psi_n$ 的构造选取
$y_i, z_i, w_i$。则计算 $x$ 的 $\psi_n$ 时，可取
$y = \sum f_iy_i$、$z = \sum f_i z_i$ 与 $w = \sum f_i w_i$，从而
可见 $\psi_n(x)$ 由 $z$ 的类给出。它在
$H_n/I^{n - c}H_n$ 中的像等于 $y$ 的类，因为
$w = \sum f_i w_i$ 位于 $I^nL$ 中。这证明 (5)。

\medskip\noindent
证明 (6)。设 $y \in \Ker(\beta)$，其在 $H$ 中的类为 $x$。若 $x$
在 $H_n$ 中映为零，则 $y \in I^nL + \Im(\alpha)$。故对某个
$v \in K$，有
$y - \alpha(v) \in \Ker(\beta) \cap I^nL$。于是
$y - \alpha(v) \in I^{n - c_2}\Ker(\beta)$，从而 $y$ 在 $H/I^nH$
中的类被 $I^{c_2}$ 零化。最后，设 $x \in H_n$ 是
$y \in \beta^{-1}(I^nM)$ 的类。仍如上写成 $y = z + w$，其中
$z \in \Ker(\beta)$ 且 $w \in I^{n - c_1}L$。显然，若
$f \in I^{c_1}$，则
% P02-E0477
$fx$ 是 $fy + fw \equiv fy$ 模 $\Im(\alpha) + I^nL$ 的类，因而
按要求，$fx$ 是 $H$ 中 $fy$ 的类之像。
\end{proof}

\begin{lemma}
\label{more-algebra-lemma-kollar-kovacs-pseudo-coherent}
\begin{reference}
% P02-E0479
Kovacs 于 23/02/2018 发来的电子邮件。
\end{reference}
设 $I$ 为 Noether 环 $A$ 的理想。设 $K \in D(A)$ 伪凝聚。置
$K_n = K \otimes_A^\mathbf{L} A/I^n$。则对所有
$i \in \mathbf{Z}$，系统 $H^i(K_n)$ 满足 Mittag-Leffler 性质，并且
$\lim H^i(K)/I^nH^i(K)$ 等于 $\lim H^i(K_n)$。
\end{lemma}

\begin{proof}
可以用有限自由 $A$-模的上有界复形 $P^\bullet$ 代表 $K$。于是
$K_n$ 由 $P^\bullet/I^nP^\bullet$ 代表。因此，
% P02-E0480
由引理 \ref{more-algebra-lemma-consequence-Artin-Rees} 得到 Mittag-Leffler 性质。
最后一个陈述随即由引理 \ref{more-algebra-lemma-kollar-kovacs} 得出。
\end{proof}

\begin{lemma}
\label{more-algebra-lemma-derived-completion-plain-completion}
设 $A$ 为 Noether 环，$I \subset A$ 为理想。设 $M^\bullet$ 为有限
$A$-模的有界复形。映射的逆系统
% P02-E0481
$$
M^\bullet \otimes_A^\mathbf{L} A/I^n \longrightarrow M^\bullet/I^nM^\bullet
$$
定义 $D(A)$ 的 pro-对象的一个同构。
\end{lemma}

\begin{proof}
设 $I = (f_1, \ldots, f_r)$。令 $K_n \in D(A)$ 为由
$f_1^n, \ldots, f_r^n$ 上的 Koszul 复形所代表的对象。回忆，存在
映射 $K_n \to A/I^n$，它们诱导逆系统的 pro-同构；见引理
\ref{more-algebra-lemma-sequence-Koszul-complexes}。因此，只需证明
$$
M^\bullet \otimes_A^\mathbf{L} K_n \longrightarrow M^\bullet/I^nM^\bullet
$$
定义 $D(A)$ 的 pro-对象的一个同构。由于 $K_n$ 由位于次数
$-r, \ldots, 0$ 的有限自由 $A$-模复形代表，存在
$a, b \in \mathbf{Z}$，使得对所有 $n$，所示箭头的源与靶的上同调
都在 $[a, b]$ 之外消失。因此，可应用《导出范畴》引理
\ref{derived-lemma-pro-isomorphism-bis}，得到只需证明：对任意
$i \in \mathbf{Z}$，映射
$$
H^i(M^\bullet \otimes_A^\mathbf{L} A/I^n) \to H^i(M^\bullet/I^nM^\bullet)
$$
定义 $A$-模的 pro-系统的同构。为此，选取拟同构
$P^\bullet \to M^\bullet$，其中 $P^\bullet$ 是有限自由 $A$-模的
上有界复形。上述箭头由映射
% P02-E0481
$$
H^i(P^\bullet/I^nP^\bullet) \to H^i(M^\bullet/I^nM^\bullet)
$$
给出。由引理 \ref{more-algebra-lemma-consequence-Artin-Rees}，它们定义
pro-系统的同构。具体而言，该引理表明两者都同构于 pro-系统
$H^i/I^nH^i$，其中
$H^i = H^i(M^\bullet) = H^i(P^\bullet)$。
\end{proof}

\begin{lemma}
\label{more-algebra-lemma-hom-systems-ML}
设 $A$ 为 Noether 环，$I \subset A$ 为理想。设 $M$、$N$ 为有限
$A$-模。置 $M_n = M/I^nM$ 与 $N_n = N/I^nN$。则：
\begin{enumerate}
\item 系统 $(\Hom_A(M_n, N_n))$ 与
$(\text{Isom}_A(M_n, N_n))$ 满足 Mittag-Leffler 性质；
\item 存在 $c \geq 0$，使得对所有 $n$，映射
$$
\Hom_A(M, N)/I^n\Hom_A(M, N) \to \Hom_A(M_n, N_n)
$$
的核与余核被 $I^c$ 零化；
\item 有
$\lim \Hom_A(M_n, N_n) =\Hom_A(M, N)^\wedge =
\Hom_{A^\wedge}(M^\wedge, N^\wedge)$；且
% P02-E0482
\item $\lim \text{Isom}_A(M_n, N_n) =
\text{Isom}_{A^\wedge}(M^\wedge, N^\wedge)$。
\end{enumerate}
这里 ${}^\wedge$ 表示通常的 $I$-进完备化。
\end{lemma}

\begin{proof}
注意 $\Hom_A(M_n, N_n) = \Hom_A(M, N_n)$。选取表示
$$
A^{\oplus t} \to A^{\oplus s} \to M \to 0
$$
。应用
% P02-E0483
右正合函子 $\Hom_A(-, N)$，得到复形
$$
0 \xrightarrow{\alpha} N^{\oplus s} \xrightarrow{\beta} N^{\oplus t}
$$
，其在中间项的上同调为 $\Hom_A(M, N)$；并且对 $n \geq 0$，
复形
$$
0 \xrightarrow{\alpha_n} N_n^{\oplus s} \xrightarrow{\beta_n} N_n^{\oplus t}
$$
的上同调为 $\Hom_A(M_n, N_n)$。对这个 $A$、$I$、$\alpha$ 与
$\beta$，令 $c \geq 0$ 如引理
\ref{more-algebra-lemma-consequence-Artin-Rees} 所给。由该引理 (3)，得到系统
$(\Hom_A(M_n, N_n))$ 的 Mittag-Leffler 性质。映射
$\Hom_A(M, N)/I^n\Hom_A(M, N) \to \Hom_A(M_n, N_n)$ 的核与余核
% P02-E0484
由该引理 (6) 被 $I^c$ 零化。由该引理 (2)，得到
$\lim \Hom_A(M_n, N_n) = \Hom_A(M, N)^\wedge$。等式
% P02-E0481
$$
\Hom_{A^\wedge}(M^\wedge, N^\wedge) = \lim \Hom_A(M_n, N_n)
$$
形式地由 $M^\wedge = \lim M_n$、$M_n = M^\wedge/I^nM^\wedge$
以及 $N$ 的相应事实得出；见《代数》引理
\ref{algebra-lemma-completion-complete}。

\medskip\noindent
同构的结果来自对 $(\Hom(M_n, N_n))$ 与 $(\Hom(N_n, M_n))$ 都应用
同态的情形，以及如下事实：若 $n > m > 0$，且有映射
$\alpha : M_n \to N_n$ 与 $\beta : N_n \to M_n$，它们诱导
% P02-E0485
同构 $M_m \to N_m$ 与 $N_m \to M_m$，则 $\alpha$ 与 $\beta$
是同构。事实上，由 Nakayama 引理（《代数》引理
\ref{algebra-lemma-NAK}），$\alpha \circ \beta$ 是满射，因而由
《代数》引理 \ref{algebra-lemma-fun}，$\alpha \circ \beta$ 是同构。
\end{proof}

\begin{lemma}
\label{more-algebra-lemma-isomorphic-completions}
设 $A$ 为 Noether 环，$I \subset A$ 为理想。设 $M$、$N$ 为有限
$A$-模。置 $M_n = M/I^nM$ 与 $N_n = N/I^nN$。若对所有 $n$，
$M_n \cong N_n$，则作为 $A^\wedge$-模，
$M^\wedge \cong N^\wedge$。
\end{lemma}

\begin{proof}
由引理 \ref{more-algebra-lemma-hom-systems-ML}，系统
$(\text{Isom}_A(M_n, N_n))$ 满足 Mittag-Leffler 性质。由假设，
每个集合 $\text{Isom}_A(M_n, N_n)$ 都非空，故
$\lim \text{Isom}_A(M_n, N_n)$ 非空。由于
$\lim \text{Isom}_A(M_n, N_n) =
\text{Isom}_{A^\wedge}(M^\wedge, N^\wedge)$，得到一个同构。
\end{proof}

\begin{remark}
\label{more-algebra-remark-weird-systems}
设 $I$ 为 Noether 环 $A$ 的理想。对 $n \geq 1$，置
$A_n = A/I^n$。考虑如下范畴：
\begin{enumerate}
\item 对象是序列 $\{E_n\}_{n \geq 1}$，其中 $E_n$ 是有限
$A_n$-模。
\item 态射 $\{E_n\} \to \{E'_n\}$ 由映射
$$
\varphi_n : I^cE_n \longrightarrow E'_n/E'_n[I^c]
\quad\text{对 }n \geq c
$$
给出，其中 $E'_n[I^c]$ 是挠子模（第 \ref{more-algebra-section-torsion} 节）；
并对以下等价关系取商：称 $(c, \varphi_n)$ 与
$(c + 1, \overline{\varphi}_n)$ 相同，其中
$\overline{\varphi}_n : I^{c + 1}E_n \longrightarrow E'_n/E'_n[I^{c + 1}]$
是诱导映射。
\end{enumerate}
态射 $(c, \varphi_n) : \{E_n\} \to \{E'_n\}$ 与
$(c', \varphi'_n) : \{E'_n\} \to \{E''_n\}$ 的复合由显然的复合
% P02-E0486
$$
I^{c + c'}E_n \to I^{c'}E'_n/E'_n[I^{c}] \to E''_n/E''_n[I^{c + c'}]
$$
定义，其中 $n \geq c + c'$。我们略去它确为一个范畴的验证。
\end{remark}

\begin{lemma}
\label{more-algebra-lemma-iso}
评注 \ref{more-algebra-remark-weird-systems} 的范畴中的态射
$(c, \varphi_n)$ 是同构，当且仅当存在 $c' \geq 0$，使得对所有
$n \gg 0$，$\Ker(\varphi_n)$ 与 $\Coker(\varphi_n)$ 都是
$I^{c'}$-挠的。
\end{lemma}

\begin{proof}
可以并且确实假设 $c' \geq c$，且对所有 $n$，
$\Ker(\varphi_n)$ 与 $\Coker(\varphi_n)$ 都是 $I^{c'}$-挠的。对
$n \geq c'$ 及 $x \in I^{c'}E'_n$，由于
$\Coker(\varphi_n)$ 被 $I^{c'}$ 零化，可以选取 $y \in I^cE_n$，
使得 $x = \varphi_n(y) \bmod E'_n[I^c]$。定义 $\psi_n(x)$ 为 $y$
在 $E_n/E_n[I^{c'}]$ 中的类。若另有选取 $y' \in I^cE_n$ 满足
$x = \varphi_n(y') \bmod E'_n[I^c]$，则差 $y - y'$ 在
$E'_n/E'_n[I^c]$ 中映为零，因而它在 $I^cE_n$ 中被 $I^{c'}$
零化。因此，映射
$\psi_n : I^{c'}E'_n \to E_n/E_n[I^{c'}]$ 良定义。我们略去验证
$(c', \psi_n)$ 是该范畴中 $(c, \varphi_n)$ 的逆。
\end{proof}

\begin{lemma}
\label{more-algebra-lemma-dejong-kollar-kovacs}
\begin{reference}
% P02-E0479
Janos Kollar、Sandor Kovacs 与 Johan de Jong 于 23/02/2018
之间的电子邮件通信。
\end{reference}
设 $I$ 为 Noether 环 $A$ 的理想。设 $M$ 与 $N$ 为有限 $A$-模。
写 $A_n = A/I^n$、$M_n = M/I^nM$ 以及 $N_n = N/I^nN$。对每个
$i \geq 0$，对象
$$
\{\Ext^i_A(M, N)/I^n\Ext^i_A(M, N)\}_{n \geq 1}
\quad\text{且}\quad
\{\Ext^i_{A_n}(M_n, N_n)\}_{n \geq 1}
$$
在评注 \ref{more-algebra-remark-weird-systems} 的范畴 $\mathcal{C}$ 中同构。
\end{lemma}

\begin{proof}
选取短正合列
$$
0 \to K \to A^{\oplus r} \to M \to 0
$$
，并置 $K_n = K/I^nK$。对 $n \geq 1$，定义
$K(n) = \Ker(A_n^{\oplus r} \to M_n)$，于是有正合列
$$
0 \to K(n) \to A_n^{\oplus r} \to M_n \to 0
$$
以及满射 $K_n \to K(n)$。事实上，由引理
\ref{more-algebra-lemma-consequence-Artin-Rees}，存在 $c \geq 0$ 与“几乎互逆”
的映射 $K(n) \to K_n/I^{n - c}K_n$。由于
$I^{n - c}K_n \subset K_n[I^c]$，
% P02-E0487
这些映射见证了系统 $\{K(n)\}_{n \geq 1}$ 与
$\{K_n\}_{n \geq 1}$ 在 $\mathcal{C}$ 中同构。

\medskip\noindent
我们断言系统
$$
\{\Ext^i_{A_n}(K(n), N_n)\}_{n \geq 1}
\quad\text{且}\quad
\{\Ext^i_{A_n}(K_n, N_n)\}_{n \geq 1}
$$
在范畴 $\mathcal{C}$ 中同构。事实上，满射 $K_n \to K(n)$ 的核
被 $I^c$ 零化，因此确定映射
$$
\Ext^i_{A_n}(K(n), N_n) \to \Ext^i_{A_n}(K_n, N_n)
$$
，其核与余核被 $I^c$ 零化。于是该断言由引理
\ref{more-algebra-lemma-iso} 得出。

\medskip\noindent
当 $i \geq 2$ 时，有同构
$$
\Ext^{i - 1}_A(K, N) = \Ext^i_A(M, N)
\quad\text{且}\quad
\Ext^{i - 1}_{A_n}(K(n), N_n) = \Ext^i_{A_n}(M_n, N_n)
$$
。由此可见，只需对 $i = 0, 1$ 证明引理。

\medskip\noindent
对 $i = 0, 1$，考虑交换图
{\footnotesize
$$
\xymatrix{
0 \ar[r] &
\Hom(M, N) \ar[r] \ar[dd] &
N^{\oplus r} \ar[r]_-\varphi \ar[dd] &
\Hom(K, N) \ar[r] \ar[d] &
\Ext^1(M, N) \ar[r] &
0 \\
& & &
\Hom(K_n, N_n)
\\
0 \ar[r] &
\Hom(M_n, N_n) \ar[r] &
N_n^{\oplus r} \ar[r] &
\Hom(K(n), N_n) \ar[r] \ar[u] &
\Ext^1(M_n, N_n) \ar[r] &
0
}
$$
}
。由引理 \ref{more-algebra-lemma-hom-systems-ML}，可见映射
$$
\Hom(M, N)/I^n \Hom(M, N) \to \Hom(M_n, N_n)
$$
与
% P02-E0488
$$
\Hom(K, N)/I^n \Hom(K, N) \to \Hom(K_n, N_n)
$$
的核与余核，被某个
$c \geq 0$ 所对应且与 $n$ 无关的 $I^c$ 零化。上面已经看到，在必要时增大 $c$ 后，
单射 $\Hom(K(n), N_n) \to \Hom(K_n, N_n)$ 的余核被 $I^c$ 零化。
对这样的 $c$，得到嵌入图中的映射
$$
\delta_n : I^c\Hom(K, N)/I^n\Hom(K, N) \to \Hom(K(n), N_n)
$$
（略去精确定式）。必要时增大 $c$ 后，$\delta_n$ 的核与余核被
$I^c$ 零化，因为已知
$$
\Hom(K, N)/I^n \Hom(K, N) \to \Hom(K_n, N_n)
$$
及
$$
\Hom(K(n), N_n) \to \Hom(K_n, N_n)
$$
具有同一性质。然后利用实线图
的交换性
{\scriptsize
$$
\xymatrix{
\varphi^{-1}(I^c\Hom(K, N)) \ar[r]_-\varphi \ar[d] &
I^c\Hom(K, N)/I^n\Hom(K, N) \ar[r] \ar[d]^{\delta_n} &
I^c\Ext^1(M, N)/I^n\Ext^1(M, N) \ar[r] \ar@{..>}[d] & 0 \\
N_n^{\oplus r} \ar[r] &
\Hom(K(n), N_n) \ar[r] &
\Ext^1(M_n, N_n) \ar[r] & 0
}
$$
}
来定义虚线箭头。直接追图（略去）表明，在最后一次必要地增大 $c$
之后，虚线箭头的核与余核
% P02-E0489
被 $I^c$ 零化。
\end{proof}

\begin{remark}
\label{more-algebra-remark-awkward}
引理 \ref{more-algebra-lemma-dejong-kollar-kovacs} 陈述中的别扭，部分源于对变化的
$n$，模 $\Ext^i_{A_n}(M_n, N_n)$ 之间没有明显映射这一事实。由该
引理可以推出：存在 $c \geq 0$，使得当 $m \gg n \gg 0$ 时，有
（典范）映射
% P02-E0490
$$
I^c\Ext^i_{A_n}(M_m, N_m)/I^n\Ext^i_{A_n}(M_m, N_m) \to
\Ext^i_{A_n}(M_n, N_n)/\Ext^i_{A_n}(M_n, N_n)[I^c]
$$
，其核与余核被 $I^c$ 零化。这正是我们得到模的系统的（弱）意义。
\end{remark}

\begin{example}
\label{more-algebra-example-has-to-be-awkward}
设 $k$ 为域，$A = k[[x, y]]/(xy)$。稍微滥用记号，
% P02-E0491
以 $x$ 与 $y$ 表示 $x$ 与 $y$ 在 $A$ 中的像。设 $I = (x)$，
$M = A/(y)$。有自由分解
% P02-E0481
$$
\ldots \to
A \xrightarrow{y}
A \xrightarrow{x}
A \xrightarrow{y} A \to M \to 0
$$
。由此得到
$$
\Ext^2_A(M, N) = N[y]/xN
$$
，其中 $N[y] = \Ker(y : N \to N)$。以
$A_n = A/I^n$、$M_n = M/I^nM$ 以及 $N_n = N/I^nN$ 表示相应对象。
对每个 $n$，有自由分解
$$
\ldots \to
A_n^{\oplus 2} \xrightarrow{y, x^{n - 1}}
A_n \xrightarrow{x}
A_n \xrightarrow{y} A_n \to M_n \to 0
$$
。由此得到
$$
\Ext^2_{A_n}(M_n, N_n) = (N_n[y] \cap N_n[x^{n - 1}])/xN_n
$$
，其中 $N_n[y] = \Ker(y : N_n \to N_n)$，且
% P02-E0492
$N[x^{n - 1}] = \Ker(x^{n - 1} : N_n \to N_n)$。取
$N = A/(y)$。则
$$
\Ext^2_A(M, N) = N[y]/xN = N/xN \cong k
$$
，但
$$
\Ext^2_{A_n}(M_n, N_n) = (N_n[y] \cap N_n[x^{n - 1}])/xN_n =
N_n[x^{n - 1}]/xN_n = 0
$$
对所有
% P02-E0493
$r$ 都成立，因为 $N_n = k[x]/(x^n)$，且序列
$$
N_n \xrightarrow{x} N_n \xrightarrow{x^{n - 1}} N_n
$$
正合。因此，为得到引理 \ref{more-algebra-lemma-dejong-kollar-kovacs} 那样的结果，
必须忽略某种 $I$-幂挠。
\end{example}

\begin{lemma}
\label{more-algebra-lemma-not-awkward}
\begin{reference}
% P02-E0479
Janos Kollar、Sandor Kovacs 与 Johan de Jong 于 23/02/2018
之间的电子邮件通信。
\end{reference}
设 $A \to B$ 为 Noether 环的平坦同态，$I \subset A$ 为理想。
% P02-E0494
设 $M, N$ 为 $A$-模。置 $B_n = B/I^nB$、
$M_n = M/I^nM$、$N_n = N/I^nN$。若 $M$ 在 $A$ 上平坦，则对所有
$i \in \mathbf{Z}$ 有
$$
\lim \Ext^i_B(M, N)/I^n \Ext^i_B(M, N) =
\lim \Ext^i_{B_n}(M_n, N_n)
$$
。
\end{lemma}

\begin{proof}
选取有限自由 $B$-模 $P_i$ 构成的分解
% P02-E0495
$$
\ldots \to P_2 \to P_1 \to P_0 \to M \to 0
$$
。置 $P_{i, n} = P_i/I^nP_i$。由于 $M$ 与 $B$ 在 $A$ 上平坦，
序列
$$
\ldots \to P_{2, n} \to P_{1, n} \to P_{0, n} \to M_n \to 0
$$
正合。一方面，复形
$$
\Hom_B(P_0, N) \to \Hom_B(P_1, N) \to \Hom_B(P_2, N) \to \ldots
$$
计算各模 $\Ext^i_B(M, N)$；另一方面，复形
$$
\Hom_{B_n}(P_{0, n}, N_n) \to \Hom_{B_n}(P_{1, n}, N_n) \to
\Hom_{B_n}(P_{2, n}, N_n) \to \ldots
$$
计算各模 $\Ext^i_{B_n}(M_n, N_n)$。由于
$$
\Hom_{B_n}(P_{i, n}, N_n) = \Hom_B(P_i, N)/I^n \Hom_B(P_i, N)
$$
，由引理 \ref{more-algebra-lemma-consequence-Artin-Rees} 的 (2) 得到结果。
\end{proof}

\section{模的系统（续）}
\label{more-algebra-section-systems-bis}

\noindent
设 $I$ 为 Noether 环 $A$ 的理想。在第 \ref{more-algebra-section-systems} 节中，
我们考察了有限 $A$-模 $M$ 的 $M/I^nM$ 型系统。本节改为考察系统
$I^nM$。

\begin{lemma}
\label{more-algebra-lemma-consequence-Artin-Rees-bis}
设 $I$ 为 Noether 环 $A$ 的理想。设
$
K \xrightarrow{\alpha} L \xrightarrow{\beta} M
$
为有限 $A$-模的复形。置 $H = \Ker(\beta)/\Im(\alpha)$。对
$n \geq 0$，设
$$
I^nK \xrightarrow{\alpha_n} I^nL \xrightarrow{\beta_n} I^nM
$$
为诱导复形。置 $H_n = \Ker(\beta_n)/\Im(\alpha_n)$。则存在典范
$A$-模映射
$$
\ldots \to H_3 \to H_2 \to H_1 \to H
$$
。存在 $c > 0$，使得当 $n \geq c$ 时，$H_n \to H$ 的像包含在
$I^{n - c}H$ 中；并存在典范 $A$-模映射
$I^nH \to H_{n - c}$，使复合
$$
I^n H \to H_{n - c} \to  I^{n - 2c}H
\quad\text{且}\quad
H_n \to I^{n - c}H \to H_{n - 2c}
$$
都是典范映射。特别地，逆系统 $(H_n)$ 与 $(I^nH)$ 作为
$\text{Mod}_A$ 的 pro-对象同构。
\end{lemma}

\begin{proof}
有
% P02-E0496
$H_n = \Ker(\beta) \cap I^nL/\alpha(I^nK)$。由于
$\Ker(\beta) \cap I^nL \subset \Ker(\beta) \cap I^{n - 1}L$ 且
$\alpha(I^nK) \subset \alpha(I^{n - 1}K)$，得到映射
$H_n \to H_{n - 1}$。映射 $H_1 \to H$ 的情形类似。

\medskip\noindent
由 Artin-Rees，可以选取 $c_1, c_2 \geq 0$，使得当
$n \geq c_1$ 时
$\Im(\alpha) \cap I^nL \subset \alpha(I^{n - c_1}K)$，且当
$n \geq c_2$ 时
$\Ker(\beta) \cap I^nL \subset I^{n - c_2}\Ker(\beta)$；见《代数》
引理 \ref{algebra-lemma-map-AR} 与
\ref{algebra-lemma-Artin-Rees}。置 $c = c_1 + c_2$。

\medskip\noindent
由 $c \geq c_2$ 的选取立即得到：当 $n \geq c$ 时，
$H_n \to H$ 的像包含在 $I^{n - c}H$ 中。

\medskip\noindent
设 $n \geq c$。如下定义 $\psi_n : I^nH \to H_{n - c}$。取
$x \in I^nH$，并选取代表 $x$ 的 $y \in I^n\Ker(\beta)$。定义
$\psi_n(x)$ 为 $y$ 在 $H_{n - c}$ 中的类。为证明其良定义，假设
$y'$ 是 $x$ 的另一个上述代表。则
$y' - y \in \Im(\alpha)$。由 $c \geq c_1$ 的选取，得到
$y' - y \in \alpha(I^{n - c}K)$，这意味着 $y$ 与 $y'$ 代表
$H_{n - c}$ 的同一个元素。因此 $\psi_n$ 良定义。

\medskip\noindent
关于复合 $I^n H \to H_{n - c} \to  I^{n - 2c}H$ 与
$H_n \to I^{n - c}H \to H_{n - 2c}$ 的陈述立即由定义得出。
\end{proof}

\begin{lemma}
\label{more-algebra-lemma-ext-factors}
设 $A$ 为 Noether 环，$I \subset A$ 为理想。设 $M$、$N$ 为
$A$-模，其中 $M$ 有限。对每个 $p > 0$，存在 $c \geq 0$，使得当
$n \geq c$ 时，映射
$\Ext_A^p(M, N) \to \Ext_A^p(I^nM, N)$ 经过映射
$\Ext^p_A(I^nM, I^{n - c}N) \to \Ext_A^p(I^nM, N)$ 分解。
\end{lemma}

\begin{proof}
对 $p = 0$，若 $\varphi : M \to N$ 是 $A$-线性映射，则对
% P02-E0497
$f_i \in A$ 与 $m_i \in M$，有
$\varphi(\sum f_i m_i) = \sum f_i \varphi(m_i)$。因此，对所有 $n$，
$\varphi$ 诱导映射 $I^nM \to I^nN$，结论对 $c = 0$ 成立。

\medskip\noindent
选取短正合列 $0 \to K \to A^{\oplus t} \to M \to 0$。对每个 $n$，
选取短正合列 $0 \to L_n \to A^{\oplus s_n} \to I^nM \to 0$。
显然可以构造短正合列之间的映射
$$
\xymatrix{
0 \ar[r] &
L_n \ar[r] \ar[d] &
A^{\oplus s_n} \ar[r] \ar[d] &
I^nM \ar[r] \ar[d] & 0 \\
0 \ar[r] &
K \ar[r] &
A^{\oplus t} \ar[r] &
M \ar[r] & 0
}
$$
，使得 $A^{\oplus s_n} \to A^{\oplus t}$ 的像位于
$(I^n)^{\oplus t}$ 中。由 Artin-Rees（《代数》引理
\ref{algebra-lemma-Artin-Rees}），存在 $c \geq 0$，使得当
$n \geq c$ 时，$L_n \to K$ 经过 $I^{n - c}K$ 分解。

\medskip\noindent
对 $p = 1$，上述选取诱导实线交换图
{\footnotesize
$$
\xymatrix{
\Hom_A(A^{\oplus s_n}, N) \ar[r] &
\Hom_A(L_n, N) \ar[r] &
\Ext_A^1(I^nM, N) \ar[r] & 0 \\
\Hom_A((I^n)^{\oplus t}, I^{n - c}N) \ar[r] \ar[u] &
\Hom_A(K \cap (I^n)^{\oplus t}, I^{n - c}N) \ar[r] \ar[u] &
\Ext_A^1(I^nM, I^{n - c}N) \ar[u] \\
\Hom_A(A^{\oplus t}, N) \ar[r] \ar[u] &
\Hom_A(K, N) \ar[r] \ar[u] &
\Ext_A^1(M, N) \ar@{..>}[u] \ar[r] & 0
}
$$
}
，其水平箭头正合。下方中间的竖直箭头来自
$K \cap (I^n)^{\oplus t} \subset I^{n - c}K$；因此，按第一段的论证，
% P02-E0498
任意 $A$-线性映射 $K \to N$ 都诱导 $A$-线性映射
$(I^n)^{\oplus t} \to I^{n - c}N$。于是按要求得到虚线箭头。

\medskip\noindent
对 $p > 1$，得到交换图
$$
\xymatrix{
\Ext^{p - 1}_A(I^{n - c}K, N) \ar[r] &
\Ext^{p - 1}_A(L_n, N) \ar[r] &
\Ext_A^p(I^nM, N) \\
\Ext^{p - 1}_A(K, N) \ar[rr] \ar[u] & &
\Ext_A^p(M, N) \ar[u]
}
$$
，其下方水平箭头是同构。对 $p$ 归纳，存在某个 $c' \geq 0$，
使得对所有 $n \geq c + c'$，左侧竖直映射经过
$\Ext^{p - 1}_A(I^{n - c}K, I^{n - c - c'}N)$ 分解。使用复合
$\Ext^{p - 1}_A(I^{n - c}K, I^{n - c - c'}N) \to
\Ext^{p - 1}_A(L_n, I^{n - c - c'}N) \to
\Ext^p_A(I^nM, I^{n - c - c'}N)$，便得到所需分解（对
$n \geq c + c'$，并以 $c + c'$ 替换 $c$）。
\end{proof}

\begin{lemma}
\label{more-algebra-lemma-ext-annihilated}
设 $A$ 为 Noether 环，$I \subset A$ 为理想。设 $M$、$N$ 为
$A$-模，其中 $M$ 有限且 $N$ 被 $I$ 的某个幂零化。对每个
$p > 0$，存在 $n$，使映射
$\Ext_A^p(M, N) \to \Ext_A^p(I^nM, N)$ 为零。
\end{lemma}

\begin{proof}
这立即由引理 \ref{more-algebra-lemma-ext-factors} 以及存在 $m > 0$ 使
$I^mN = 0$ 的事实得出。
\end{proof}

\begin{lemma}
\label{more-algebra-lemma-ext-induced-toplogy}
设 $A$ 为 Noether 环，$I \subset A$ 为理想。设 $K \in D(A)$
伪凝聚，$M$ 为有限 $A$-模。对每个 $p \in \mathbf{Z}$，存在
% P02-E0499
$c$，使得当 $n \geq c$ 时，
$\Ext_A^p(K, I^nM) \to \Ext_A^p(K, M)$ 的像包含在
$I^{n - c}\Ext_A^p(K, M)$ 中。
\end{lemma}

\begin{proof}
选取代表 $K$ 的有限自由 $A$-模上有界复形 $P^\bullet$。则
$\Ext_A^p(K, M)$ 是
% P02-E0500
$$
\Hom_A(F^{-p + 1}, M) \xrightarrow{a}
\Hom_A(F^{-p}, M) \xrightarrow{b}
\Hom_A(F^{-p - 1}, M)
$$
的上同调，而 $\Ext_A^p(K, I^nM)$ 由把这些有限 $A$-模替换为它们的
$I^n$ 倍所得复形计算。因此，
% P02-E0501
结果由引理 \ref{more-algebra-lemma-consequence-Artin-Rees-bis} 得出
（事实上还有更强的结论）。
\end{proof}

\noindent
在情形 \ref{more-algebra-situation-koszul} 中，我们定义复形 $I_n^\bullet$，使得
在 $A$-模复形模伦的三角范畴 $K(A)$ 中有区分三角
% P02-E0502
$$
I_n^\bullet \to A \to K_n^\bullet \to I_n^\bullet[1]
$$
。具体地，置 $I_n^\bullet = \sigma_{\leq -1}K_n^\bullet[-1]$。
存在逐项分裂的复形短正合列
$$
0 \to A \to K_n^\bullet \to I_n^\bullet[1] \to 0
$$
，它们由 $K(A)$ 上三角结构的定义给出区分三角。旋转这些三角便得到
上面所需的区分三角。注意，$I_n^0 = A^{\oplus r} \to A$ 在第
$i$ 个因子上由乘以 $f_i^n$ 给出。因此，$I_n^\bullet \to A$
分解为
$$
I_n^\bullet \to (f_1^n, \ldots, f_r^n) \to A
$$
。事实上，有短正合列
$$
0 \to H^{-1}(K_n^\bullet) \to H^0(I_n^\bullet) \to (f_1^n, \ldots, f_r^n) \to 0
$$
，并且对每个 $i < 0$，有
$H^i(I_n^\bullet) = H^{i - 1}(K_n^\bullet)$。映射
$K_{n + 1}^\bullet \to K_n^\bullet$ 诱导映射
$I_{n + 1}^\bullet \to I_n^\bullet$，并得到交换图
$$
\xymatrix{
\ldots \ar[r] &
I_3^\bullet \ar[d] \ar[r] &
I_2^\bullet \ar[d] \ar[r] &
I_1^\bullet \ar[d] \\
\ldots \ar[r] &
(f_1^3, \ldots, f_r^3) \ar[r] &
(f_1^2, \ldots, f_r^2) \ar[r] &
(f_1, \ldots, f_r)
}
$$
，它位于 $K(A)$ 中。

\begin{lemma}
\label{more-algebra-lemma-sequence-powers-pro-bounded}
在情形 \ref{more-algebra-situation-koszul} 中，假设 $A$ 为 Noether 环。沿用上述
记号，逆系统 $(I^n)$ 在 $D(A)$ 中 pro-同构于逆系统
$(I_n^\bullet)$。
\end{lemma}

\begin{proof}
容易证明，逆系统 $I^n$ 在 $A$-模范畴中 pro-同构于逆系统
$(f_1^n, \ldots, f_r^n)$。考虑区分三角的逆系统
$$
I_n^\bullet \to (f_1^n, \ldots, f_r^n) \to C_n^\bullet \to I_n^\bullet[1]
$$
，其中 $C_n^\bullet$ 是第一个箭头的锥。由《导出范畴》引理
\ref{derived-lemma-pro-isomorphism}，只需证明逆系统
$C_n^\bullet$ 是 pro-零的。复形 $I_n^\bullet$ 仅在满足
$-r + 1 \leq i \leq 0$ 的次数 $i$ 上有非零项，故
$C_n^\bullet$ 也类似地有界。因此，由《导出范畴》引理
\ref{derived-lemma-essentially-constant-cohomology}，只需证明
$H^p(C_n^\bullet)$ 是 pro-零的。由上面的讨论，当 $p \leq -1$ 时，
$H^p(C_n^\bullet) = H^p(K_n^\bullet)$；当 $p \geq 0$ 时，
$H^p(C_n^\bullet) = 0$。在引理
\ref{more-algebra-lemma-sequence-Koszul-complexes} 的证明中，已经证明各逆系统
$H^p(K_n^\bullet)$ 是 pro-零的（也正是在这里使用了 $A$ 为
Noether 环这一假设）。
\end{proof}

\begin{lemma}
\label{more-algebra-lemma-tensoring-Deligne-system}
设 $A$ 为 Noether 环，$I \subset A$ 为理想。设 $M^\bullet$ 为有限
$A$-模的有界复形。映射的逆系统
$$
I^n \otimes_A^\mathbf{L} M^\bullet \longrightarrow I^nM^\bullet
$$
定义 $D(A)$ 的 pro-对象的一个同构。
\end{lemma}

\begin{proof}
选取 $I$ 的生成元 $f_1, \ldots, f_r \in I$。逆系统 $I^n$ 在
$A$-模范畴中 pro-同构于逆系统
$(f_1^n, \ldots, f_r^n)$。沿用引理
\ref{more-algebra-lemma-sequence-powers-pro-bounded} 的记号，可见只需证明映射的
逆系统
$$
I_n^\bullet \otimes_A^\mathbf{L} M^\bullet
\longrightarrow
(f_1^n, \ldots, f_r^n)M^\bullet
$$
定义 $D(A)$ 的 pro-对象的一个同构。设 $a \leq b$ 满足：若
$i \not \in [a, b]$，则 $M^i = 0$。于是上述箭头的源与靶仅在次数
$[-r + a, b]$ 中有上同调。因此，只需证明：对任意
$p \in \mathbf{Z}$，映射的逆系统
$$
H^p(I_n^\bullet \otimes_A^\mathbf{L} M^\bullet)
\longrightarrow
H^p((f_1^n, \ldots, f_r^n)M^\bullet)
$$
定义 $A$-模的 pro-对象的一个同构；见《导出范畴》引理
\ref{derived-lemma-pro-isomorphism-bis}。利用
$I_n^\bullet \otimes_A^\mathbf{L} M^\bullet$ 与
$I^n \otimes_A^\mathbf{L} M^\bullet$ 之间的 pro-同构，以及
$(f_1^n, \ldots, f_r^n)M^\bullet$ 与 $I^nM^\bullet$ 之间的
pro-同构，这等价于证明映射的逆系统
$$
H^p(I^n \otimes_A^\mathbf{L} M^\bullet)
\longrightarrow
H^p(I^nM^\bullet)
$$
定义 $A$-模的 pro-对象的一个同构。
% P02-E0502
选取有限自由 $A$-模的上有界复形 $P^\bullet$ 及拟同构
$P^\bullet \to M^\bullet$。于是，只需证明映射的逆系统
$$
H^p(I^nP^\bullet)
\longrightarrow
H^p(I^nM^\bullet)
$$
是 pro-同构。这由引理
\ref{more-algebra-lemma-consequence-Artin-Rees-bis} 以及
$H^p(P^\bullet) = H^p(M^\bullet)$ 得出。
\end{proof}

\begin{lemma}
\label{more-algebra-lemma-factor-through-derived-tensor-product}
设 $A$ 为 Noether 环，$I \subset A$ 为理想，$M$ 为有限 $A$-模。
存在整数 $n > 0$，使得 $I^nM \to M$ 在 $D(A)$ 中经过映射
$I \otimes_A^\mathbf{L} M \to M$ 分解。
\end{lemma}

\begin{proof}
这由引理 \ref{more-algebra-lemma-tensoring-Deligne-system} 得出。
% P02-E0503
也可如下直接看出。考虑区分三角
% P02-E0502
$$
I \otimes_A^\mathbf{L} M \to M \to A/I \otimes_A^\mathbf{L} M \to
I \otimes_A^\mathbf{L} M[1]
$$
。由三角范畴的公理，只需证明对某个 $n$，
$I^nM \to A/I \otimes_A^\mathbf{L} M$ 在 $D(A)$ 中为零。选取 $I$
的生成元 $f_1, \ldots, f_r$，令
$K = K_\bullet(A, f_1, \ldots, f_r)$ 为 Koszul 复形，并考虑商映射
的分解 $A \to K \to A/I$。于是只需证明：对某个 $n > 0$，
$I^nM \to K \otimes_A M$ 在 $D(A)$ 中为零。假设已找到
$n > 0$，使 $I^nM \to K \otimes_A M$ 在 $D(A)$ 中经过
$\tau_{\geq t}(K \otimes_A M)$ 分解。将其进一步经过
$\tau_{\geq t + 1}(K \otimes_A M)$ 分解的障碍，是
$\Ext^t(I^nM, H_t(K \otimes_A M))$ 的一个元素。有限 $A$-模
$H_t(K \otimes_A M)$ 被 $I$ 零化。由引理
\ref{more-algebra-lemma-ext-annihilated}，增大 $n$ 后可假设该障碍元素为零。
有限次重复此过程，便找到 $n$，使得 $I^nM \to K \otimes_A M$
在 $D(A)$ 中经过
$0 = \tau_{\geq r + 1}(K \otimes_A M)$ 分解，证明完成。
\end{proof}

\section{杂项}
\label{more-algebra-section-misc}

\noindent
下面给出一些不适合归入其他章节的结果。

\begin{lemma}
\label{more-algebra-lemma-ext-annihilated-into}
设 $A$ 为 Noether 环，$I \subset A$ 为理想，$K \in D(A)$ 伪凝聚，
$a \in \mathbf{Z}$。假设对每个有限 $A$-模 $M$，当 $i \geq a$ 时，
模 $\Ext^i_A(K, M)$ 都是 $I$-幂挠的。则对 $i \geq a$ 及有限模
$M$，系统 $\Ext^i_A(K, M/I^nM)$ 本质常值，其值为
% P02-E0504
$$
\Ext^i_A(K, M) = \lim \Ext^i_A(K, M/I^nM)
$$
。
\end{lemma}

\begin{proof}
设 $M$ 为有限 $A$-模。由于 $K$ 伪凝聚，$\Ext^i_A(K, M)$ 是有限
$A$-模。因此，当 $i \geq a$ 时，它被某个 $I^t$（$t \geq 0$）
零化。由引理 \ref{more-algebra-lemma-ext-induced-toplogy}，对某个 $n > 0$，
$\Ext^i_A(K, I^nM) \to \Ext^i_A(K, M)$ 的像为零。短正合列
$0 \to I^nM \to M \to M/I^n M \to 0$ 给出长正合列
$$
\Ext^i_A(K, I^nM) \to \Ext^i_A(K, M) \to
\Ext^i_A(K, M/I^nM) \to \Ext^{i + 1}_A(K, I^nM)
$$
。由刚才的论证（将其应用于有限 $A$-模 $I^mM$），系统
$\Ext^i_A(K, I^nM)$ 与 $\Ext^{i + 1}_A(K, I^nM)$ 都本质常值，
其值为 $0$。追图可知，$\Ext^i_A(K, M/I^nM)$ 本质常值，其值为
$\Ext^i_A(K, M)$。
\end{proof}

\begin{lemma}
\label{more-algebra-lemma-tor-annihilated}
设 $A$ 为 Noether 环，$I \subset A$ 为理想，$M$ 为有限 $A$-模，
$N$ 为被 $I$ 零化的 $A$-模。存在整数 $n > 0$，使得对所有
$p \geq 0$，映射
$\text{Tor}^A_p(I^nM, N) \to \text{Tor}^A_p(M, N)$ 为零。
\end{lemma}

\begin{proof}
由引理 \ref{more-algebra-lemma-factor-through-derived-tensor-product}，映射
$I^nM \to M$ 可分解为
$I^nM \to M \otimes_A^\mathbf{L} I \to M$。我们断言复合
$$
I^nM \otimes_A^\mathbf{L} N \to
(M \otimes_A^\mathbf{L} I) \otimes_A^\mathbf{L} N
\to M \otimes_A^\mathbf{L} N
$$
为零。事实上，图
$$
\xymatrix{
(M \otimes_A^\mathbf{L} I) \otimes_A^\mathbf{L} N \ar[rr] \ar[rd] & &
M \otimes_A^\mathbf{L} (I \otimes_A^\mathbf{L} N) \ar[ld] \\
& M \otimes_A^\mathbf{L} N
}
$$
交换（略去细节），而映射 $I \otimes_A^\mathbf{L} N \to N$
为零，因为 $N$ 被 $I$ 零化。
\end{proof}

\begin{lemma}
\label{more-algebra-lemma-pseudo-coherent-tensor-limit}
设 $R$ 为环，$K \in D(R)$ 伪凝聚，$(M_n)$ 为 $R$-模的逆系统。则
$R\lim K \otimes_R^\mathbf{L} M_n = K \otimes_R^\mathbf{L} R\lim M_n$。
\end{lemma}

\begin{proof}
考虑定义性的区分三角
$$
R\lim M_n \to \prod M_n \to \prod M_n \to R\lim M_n[1]
$$
，并应用引理 \ref{more-algebra-lemma-pseudo-coherent-tensor}。
\end{proof}

\begin{lemma}
\label{more-algebra-lemma-additivity-of-pd}
设 $R$ 为 Noether 局部环，$I \subset R$ 为理想，$E$ 为非零有限
$R/I$-模。若 $R/I$ 具有有限投射维数，且 $E$ 作为 $R/I$-模具有
有限投射维数，则 $E$ 作为 $R$-模具有有限投射维数，并且
$$
\text{pd}_R(E) = \text{pd}_R(R/I) + \text{pd}_{R/I}(E)
$$
。
\end{lemma}

\begin{proof}
我们将使用如下事实：对有限模而言，在 $R$（相应地，在 $R/I$）上
具有有限投射维数，等价于它是完美模；见定义
\ref{more-algebra-definition-perfect} 后的讨论。由引理
\ref{more-algebra-lemma-cohomology-perfect}，$E$ 在 $R$ 上具有有限投射维数。
因此可以应用 Auslander-Buchsbaum（《代数》命题
\ref{algebra-proposition-Auslander-Buchsbaum}），得到
$$
\text{pd}_R(E) + \text{depth}(E) = \text{depth}(R),\quad
\text{pd}_{R/I}(E) + \text{depth}(E) = \text{depth}(R/I),
$$
以及
$$
\text{pd}_R(R/I) + \text{depth}(R/I) = \text{depth}(R)
$$
。注意，第一个等式中的 $E$ 的深度按 $R$-模计算，第二个等式中
则按 $R/I$-模计算。不过这两个深度相同（这很显然，也可由《代数》
引理 \ref{algebra-lemma-depth-goes-down-finite} 得出）。证明完毕。
\end{proof}

\begin{lemma}
\label{more-algebra-lemma-enlarge}
设 $A \to B$ 为环映射。存在基数
$\kappa = \kappa(A \to B)$，具有如下性质。设 $M^\bullet$ 为
$A$-模复形，$N^\bullet$ 为 $B$-模复形。设
$a : M^\bullet \to N^\bullet$ 为 $A$-模复形的映射，并且它在
$D(B)$ 中诱导同构
$M^\bullet \otimes_A^\mathbf{L} B \to N^\bullet$。设
$M_1^\bullet \subset M^\bullet$ 为 $A$-模子复形，
$N_1^\bullet \subset N^\bullet$ 为 $B$-模子复形，并满足
$a(M_1^\bullet) \subset N_1^\bullet$。则存在子复形
$$
M_1^\bullet \subset M_2^\bullet \subset M^\bullet
\quad\text{且}\quad
N_1^\bullet \subset N_2^\bullet \subset N^\bullet
$$
，使得 $a(M_2^\bullet) \subset N_2^\bullet$，且具有以下性质：
\begin{enumerate}
\item $\Ker(H^i(M_1^\bullet \otimes_A^\mathbf{L} B) \to H^i(N_1^\bullet))$
在 $H^i(M_2^\bullet \otimes_A^\mathbf{L} B)$ 中的像为零；
\item
% P02-E0505
$\Im(H^i(N_1^\bullet) \to H^i(N_2^\bullet))$ 包含在
$\Im(H^i(M_2^\bullet \otimes_A^\mathbf{L} B) \to H^2(N_2^\bullet))$
中；
\item $|\bigcup M_2^i \cup \bigcup N_2^i| \leq
\max(\kappa, |\bigcup M_1^i \cup \bigcup N_1^i|)$。
\end{enumerate}
\end{lemma}

\begin{proof}
置 $\kappa = \max(|A|, |B|, \aleph_0)$。记
$|M^\bullet| = |\bigcup M^i|$，其他复形也采用类似记号。于是，对
引理陈述中使用的量，有
% P02-E0506
$$
\max(\kappa, |\bigcup M_1^i \cup \bigcup N_1^i|) =
\max(\kappa, |M_1^\bullet|, |M_2^\bullet|)
$$
。下文将不再特别说明地使用此等由基数算术得到的观察。

\medskip\noindent
首先说明存在足够多的“小”子复形。对有限子集族
$E = \{E^i\}$ 与 $F = \{F^i\}$ 的每一对，其中
$E^i \subset M^i$、$i \in \mathbf{Z}$，且
$F^i \subset N^i$、$i \in \mathbf{Z}$，令
$$
M_1^\bullet \subset M_1(E, F)^\bullet \subset M^\bullet
\quad\text{且}\quad
N_1^\bullet \subset N_1(E, F)^\bullet \subset N^\bullet
$$
为分别由 $A$-模和 $B$-模构成的最小子复形，使得
$a(M_1(E, F)^\bullet) \subset N_1(E, F)^\bullet$，并且
% P02-E0507
$E^i \subset M_1(E, F)^i$ 与 $F^i \subset M_2(E, F)^i$。容易看出
% P02-E0508
$$
|M_1(E, F)^\bullet| \leq \max(\kappa, |M_1^\bullet|)
\quad\text{且}\quad
|M_2(E, F)^\bullet| \leq \max(\kappa, |M_2^\bullet|)
$$
。略去细节。显然有
$$
M^\bullet = \colim_{(E, F)} M_1(E, F)^\bullet
\quad\text{且}\quad
N^\bullet = \colim_{(E, F)} N_1(E, F)^\bullet
$$
，并且这些余极限是（逐项的）滤过余极限。

\medskip\noindent
存在由自由 $A$-模 $F_i$ 构成的解消
% P02-E0509
$\ldots \to F^{-1} \to F^0 \to B$，满足 $|F_i| \leq \kappa$
（略去细节）。$M_1^\bullet \otimes_A^\mathbf{L} B$ 的上同调模由
$\text{Tot}(M_1^\bullet \otimes_A F^\bullet)$ 计算。因此
$|H^i(M_1^\bullet \otimes_A^\mathbf{L} B)| \leq
\max(\kappa, |M_1^\bullet|)$。

\medskip\noindent
设 $i \in \mathbf{Z}$，并设
$\xi \in H^i(M_1^\bullet \otimes_A^\mathbf{L} B)$ 是一个在
$H^i(N_1^\bullet)$ 中映为零的元素。则 $\xi$ 在
$H^i(N^\bullet)$ 中映为零，从而 $\xi$ 在
$H^i(M^\bullet \otimes_A^\mathbf{L} B)$ 中也映为零。由于导出张量积
与滤过余极限交换，可以找到上述有限子集族 $E_\xi$ 与 $F_\xi$，
使得 $\xi$ 在
$H^i(M_1(E_\xi, F_\xi)^\bullet \otimes_A^\mathbf{L} B)$ 中映为零。

\medskip\noindent
设 $i \in \mathbf{Z}$，并设 $\eta \in H^i(N_1^\bullet)$。则
$\eta$ 在 $H^i(N^\bullet)$ 中的像属于
$H^i(M^\bullet \otimes_A^\mathbf{L} B) \to H^i(N^\bullet)$ 的像。
因此，与前面一样，可以找到上述有限子集族 $E_\eta$ 与 $F_\eta$，
使得 $\eta$ 映到
% P02-E0510
$H^i(N_1(E_\eta, F_\eta)$ 的一个元素，而该元素属于映射
$H^i(M_1(E_\eta, F_\eta)^\bullet \otimes_A^\mathbf{L} B) \to
H^i(N_1(E_\eta, F_\eta)$ 的像。

\medskip\noindent
现在只需定义
$$
M_2^\bullet =
\sum\nolimits_\xi M_1(E_\xi, F_\xi)^\bullet +
\sum\nolimits_\eta M_1(E_\eta, F_\eta)^\bullet
$$
，其中求和遍历前两段中的 $\xi$ 与 $\eta$，并在 $M^\bullet$ 内
取和。类似地，置
$$
N_2^\bullet =
\sum\nolimits_\xi N_1(E_\xi, F_\xi)^\bullet +
\sum\nolimits_\eta N_1(E_\eta, F_\eta)^\bullet
$$
，其中在 $N^\bullet$ 内取和。由构造，这些选取满足性质 (1) 与
(2)。性质 (3) 的界也随之成立，因为 $\xi$ 与 $\eta$ 所成集合的
基数至多为 $\max(\kappa, |M_1^\bullet|, |N_1^\bullet|)$。
\end{proof}

\begin{lemma}
\label{more-algebra-lemma-f-bounded}
设 $R$ 为环，$f \in R$，$M \in D(R)$，并设 $C$ 为
$f : M \to M$ 的锥。若当 $i < 0$ 时 $H^i(M)_f = 0$，且当
$i < -1$ 时 $H^i(C) = 0$，则当 $i < 0$ 时 $H^i(M) = 0$。
\end{lemma}

\begin{proof}
记 $M_f = M \otimes_R^\mathbf{L} R_f$，并选取区分三角
$F \to M \to M_f$。由假设，当 $i < 0$ 时，
$H^i(M_f) = H^i(M)_f = 0$。因此，只需证明当 $i < 0$ 时
$H^i(F) = 0$。注意，对所有 $i \in \mathbf{Z}$，$H^i(F)$ 都是
$f$-幂挠的。另一方面，由于 $f : M_f \to M_f$ 是同构，可见
$C$ 同构于 $f : F \to F$ 的锥（使用《导出范畴》命题
\ref{derived-proposition-9}）。现在，若 $H^i(F) \not = 0$，则
$f : H^i(F) \to H^i(F)$ 的核非零，这意味着
$H^{i - 1}(C)$ 非零。由假设，当 $i - 1 < -1$ 时不可能发生这种
情形，证明完毕。
\end{proof}

\begin{lemma}
\label{more-algebra-lemma-f-flat}
设 $R$ 为环，$f \in R$，$M \in D(R)$。假设
\begin{enumerate}
\item 当 $i > 0$ 时，$H^i(M) = 0$；
\item $M \otimes_R^\mathbf{L} R_f$ 同构于置于次数 $0$ 的平坦
$R_f$-模；
\item $M \otimes_R^\mathbf{L} R/fR$ 同构于置于次数 $0$ 的平坦
$R/fR$-模。
\end{enumerate}
则 $M$ 同构于置于次数 $0$ 的平坦 $R$-模。
\end{lemma}

\begin{proof}
设 $K$ 为 $R$-模。只需证明
$N = M \otimes_R^\mathbf{L} K$ 同构于置于次数 $0$ 的
$R$-模；见第 \ref{more-algebra-section-tor} 节。由假设 (1)，$N$ 仅在
$\leq 0$ 的次数上可能有非零上同调。因此，由引理
\ref{more-algebra-lemma-f-bounded}，只需证明当 $i < 0$ 时
$H^i(N)_f = 0$，并证明当 $i < -1$ 时 $H^i(C) = 0$；这里
$C$ 是 $f : N \to N$ 的锥。对于第一点，注意
% P02-E0504
$$
H^i(N)_f = H^i(N \otimes_R^\mathbf{L} R_f) =
H^i(M \otimes_R^\mathbf{L} K \otimes_R^\mathbf{L} R_f ) =
H^i((M \otimes_R^\mathbf{L} R_f) \otimes_{R_f}^\mathbf{L}
(K \otimes_R^\mathbf{L} R_f))
$$
。由假设 (2)，除非 $i = 0$，否则它为零；此时得到
$H^0(M)_f \otimes_R K$。对于第二点，设 $C'$ 为
$f : K \to K$ 的锥。则
$$
C = M \otimes_R^\mathbf{L} C'
$$
。复形 $C'$ 仅在次数 $0$ 与 $-1$ 上有非零上同调，分别等于
$K/fK$ 与 $K[f]$；换言之，存在区分三角
$$
(K[f])[1] \to C' \to K/fK
$$
。另一方面，对任意被 $f$ 零化的 $R$-模 $K'$，有
$$
M \otimes_R^\mathbf{L} K' =
M \otimes_R^\mathbf{L} R/fR \otimes_{R/fR}^\mathbf{L} K'
$$
。由假设 (3)，它等于置于次数 $0$ 的模
$H^0(M \otimes_R^\mathbf{L} R/fR) \otimes_{R/fR} K'$。
结合以上结论，$C$ 仅在次数 $0$ 与 $-1$ 上有非零上同调，证明完成。
\end{proof}

\section{双复形的若干技巧}
\label{more-algebra-section-tricks}

\noindent
本节继续《同调》第
\ref{homology-section-double-complexes-abelian-groups} 节的讨论。

\begin{lemma}
\label{more-algebra-lemma-vanishing-coh-prod-totalization}
设 $A_0^\bullet \to A_1^\bullet \to A_2^\bullet \to \ldots$
为阿贝尔群复形所成的复形。假设对所有 $p \geq 0$，有
$H^{-p}(A_p^\bullet) = 0$。置 $A^{p, q} = A_p^q$，并把
$A^{\bullet, \bullet}$ 视为双复形。则
$H^0(\text{Tot}_\pi(A^{\bullet, \bullet})) = 0$。
\end{lemma}

\begin{proof}
记 $f_p : A_p^\bullet \to A_{p + 1}^\bullet$ 为给定的复形映射。
回忆 $\text{Tot}_\pi(A^{\bullet, \bullet})$ 上的微分在次数 $n$
的元素上由
$$
\prod\nolimits_{p + q = n} (f^q_p + (-1)^p\text{d}^q_{A_p^\bullet})
$$
给出。设 $\xi \in H^0(\text{Tot}_\pi(A^{\bullet, \bullet}))$
为上同调类。我们将证明 $\xi$ 为零。用上循环
% P02-E0511
$x = (x_p) \in \prod A^{p, -p}$ 的类表示 $\xi$。由
$\text{d}(x) = 0$，得到
$\text{d}_{A_0^\bullet}(x_0) = 0$。由于
$H^0(A_0^\bullet) = 0$，存在 $y_{-1} \in A^{0, -1}$，满足
$\text{d}_{A_0^\bullet}(y_{-1}) = x_0$。于是
$\text{d}_{A_1^\bullet}(x_1 + f_0(y_{-1})) = 0$。由于
$H^{-1}(A_1^\bullet) = 0$，可以找到 $y_{-2} \in A^{1, -2}$，
使得 $-\text{d}_{A_1^\bullet}(y_{-2}) = x_1 + f_0(y_{-1})$。
归纳可找到 $y_{-p - 1} \in A^{p, -p - 1}$，使得
% P02-E0513
$$
(-1)^p\text{d}_{A_p^\bullet}(y_{-p - 1}) = x_p + f_{p - 1}(y_{-p})
$$
。这意味着 $\text{d}(y) = x$，其中 $y = (y_{-p - 1})$。
\end{proof}

\begin{lemma}
\label{more-algebra-lemma-prod-qis-gives-qis}
设
$$
(A_0^\bullet \to A_1^\bullet \to A_2^\bullet \to \ldots)
\longrightarrow
(B_0^\bullet \to B_1^\bullet \to B_2^\bullet \to \ldots)
$$
为两个阿贝尔群复形所成复形之间的映射。置
$A^{p, q} = A_p^q$、$B^{p, q} = B_p^q$，得到双复形。设
$\text{Tot}_\pi(A^{\bullet, \bullet})$ 与
$\text{Tot}_\pi(B^{\bullet, \bullet})$ 为这些双复形相伴的乘积
全复形。若每个 $A_p^\bullet \to B_p^\bullet$ 都是拟同构，则
$\text{Tot}_\pi(A^{\bullet, \bullet}) \to
\text{Tot}_\pi(B^{\bullet, \bullet})$ 是拟同构。
\end{lemma}

\begin{proof}
回忆 $\text{Tot}_\pi(A^{\bullet, \bullet})$ 在次数 $n$ 上由
% P02-E0512
$\prod_{p + q = n} A^{p, q} = \prod_{p + 1 = n} A^q_p$ 给出。
设 $C_p^\bullet$ 为映射 $A_p^\bullet \to B_p^\bullet$ 的锥；
见《导出范畴》第 \ref{derived-section-cones} 节。由锥构造的函子性，
得到复形所成的复形
$$
C_0^\bullet \to C_1^\bullet \to C_2^\bullet \to \ldots
$$
。可见 $\text{Tot}_\pi(C^{\bullet, \bullet})$ 在次数 $n$ 上由
$$
\prod_{p + q = n} C^{p, q} = \prod_{p + q = n} C^q_p =
\prod_{p + q = n} (B^q_p \oplus A^{q + 1}_p) =
\prod_{p + q = n} B^q_p \oplus \prod_{p + q = n} A^{q + 1}_p
$$
给出。因此，$\text{Tot}_\pi(C^{\bullet, \bullet})$ 是映射
$\text{Tot}_\pi(A^{\bullet, \bullet}) \to
\text{Tot}_\pi(B^{\bullet, \bullet})$ 的锥（
% P02-E0514
略去验证微分一致的过程）。所以只需证明：若每个
$A_p^\bullet$ 都无环，则 $\text{Tot}_\pi(A^{\bullet, \bullet})$
无环。这由引理 \ref{more-algebra-lemma-vanishing-coh-prod-totalization} 得出。
\end{proof}

\section{弱 \'{e}tale 环映射}
\label{more-algebra-section-weakly-etale}

\noindent
本节的大多数结果来自 Olivier 的论文 \cite{Olivier-AF}。另见相关论文
\cite{Ferrand-epi}。

\begin{definition}
\label{more-algebra-definition-weakly-etale}
若每个 $A$-模都在 $A$ 上平坦，则称环 $A$ {\it 绝对平坦}。若
$A \to B$ 与 $B \otimes_A B \to B$ 都平坦，则称环映射
$A \to B$ {\it 弱 \'{e}tale}或{\it 绝对平坦}。
\end{definition}

\noindent
绝对平坦环有时称为 von Neumann 正则环（这一名称常用于非交换环的
语境）。局部化是弱 \'{e}tale 环映射。\'{e}tale 环映射是弱
\'{e}tale 的。下面是一个简单而关键的性质。

\begin{lemma}
\label{more-algebra-lemma-key}
设 $A \to B$ 为环映射，并且 $B \otimes_A B \to B$ 平坦。设 $N$
为 $B$-模。若 $N$ 作为 $A$-模平坦，则 $N$ 作为 $B$-模平坦。
\end{lemma}

\begin{proof}
假设 $N$ 作为 $A$-模平坦。
% P02-E0515
则函子
$$
\text{Mod}_B \longrightarrow \text{Mod}_{B \otimes_A B},\quad
N' \mapsto N \otimes_A N'
$$
正合。由于 $B \otimes_A B \to B$ 平坦，可知函子
$$
\text{Mod}_B \longrightarrow \text{Mod}_B,\quad
N' \mapsto (N \otimes_A N') \otimes_{B \otimes_A B} B = N \otimes_B N'
$$
正合，因此 $N$ 在 $B$ 上平坦。
\end{proof}

\begin{definition}
\label{more-algebra-definition-weak-dimension}
设 $A$ 为环，$d \geq 0$ 为整数。若每个 $A$-模的 Tor 维数
$\leq d$，则称 $A$ 的{\it 弱维数 $\leq d$}。
\end{definition}

\begin{lemma}
\label{more-algebra-lemma-weak-dimension-goes-up}
设 $A \to B$ 为弱 \'{e}tale 环映射。若 $A$ 的弱维数至多为 $d$，
则 $B$ 也如此。
\end{lemma}

\begin{proof}
设 $N$ 为 $B$-模。若 $d = 0$，则 $N$ 作为 $A$-模平坦，从而由
引理 \ref{more-algebra-lemma-key}，它作为 $B$-模平坦。假设 $d > 0$。选取
自由 $B$-模解消 $F_\bullet \to N$。我们的假设蕴含
$K = \Im(F_d \to F_{d - 1})$ 是 $A$-平坦的；见引理
\ref{more-algebra-lemma-last-one-flat}。故由引理 \ref{more-algebra-lemma-key}，它是
$B$-平坦的。因此
$0 \to K \to F_{d - 1} \to \ldots \to F_0 \to N \to 0$
是长度为 $d$ 的平坦解消，可见 $N$ 的 Tor 维数至多为 $d$。
\end{proof}

\begin{lemma}
\label{more-algebra-lemma-absolutely-flat}
设 $A$ 为环。下列条件等价：
\begin{enumerate}
\item $A$ 的弱维数 $\leq 0$；
\item $A$ 绝对平坦；
\item $A$ 既约，且每个素理想都是极大理想；
\item $A$ 的每个局部环都是域。
\end{enumerate}
\end{lemma}

\begin{proof}
(1) 与 (2) 的等价性由定义立即得到。

\medskip\noindent
证明 (1) $\Rightarrow$ (3)。
假设 $A$ 绝对平坦。则 $A$ 的每个理想都是纯理想；见《代数》定义
\ref{algebra-definition-pure-ideal}。因此，由《代数》引理
\ref{algebra-lemma-finitely-generated-pure-ideal}，每个有限生成理想
都由一个幂等元生成。若 $f \in A$，则对某个幂等元 $e \in A$，
有 $(f) = (e)$，且 $D(f) = D(e)$ 既开又闭（《代数》引理
\ref{algebra-lemma-idempotent-spec}）。由《代数》引理
\ref{algebra-lemma-ring-with-only-minimal-primes}，这蕴含 $A$ 的每个
素理想都是极大理想。此外，若 $f$ 幂零，则 $e = 0$，故 $f = 0$。
因此 $A$ 既约。

\medskip\noindent
证明 (3) $\Rightarrow$ (4)。若 $A$ 既约且每个素理想都是极大理想，
则由《代数》引理 \ref{algebra-lemma-minimal-prime-reduced-ring}，
每个局部环都是域。

\medskip\noindent
证明 (4) $\Rightarrow$ (2)。若 $A$ 的每个局部环都是域，则由
《代数》引理 \ref{algebra-lemma-flat-localization}，每个 $A$-模
都平坦。
\end{proof}

\begin{lemma}
\label{more-algebra-lemma-product-fields-absolutely-flat}
域的乘积是绝对平坦环。
\end{lemma}

\begin{proof}
设 $K_i$ 为一族域。若 $f = (f_i) \in \prod K_i$，则 $f$ 生成的理想
与幂等元 $e = (e_i)$ 生成的理想相同，其中根据 $f_i$ 是否为 $0$，
$e_i = 0, 1$。因此 $D(f) = D(e)$ 既开又闭；结论由引理
\ref{more-algebra-lemma-absolutely-flat} 与《代数》引理
\ref{algebra-lemma-ring-with-only-minimal-primes} 得出。
\end{proof}

\begin{lemma}
\label{more-algebra-lemma-base-change-weakly-etale}
设 $A \to B$ 与 $A \to A'$ 为环映射，设
$B' = B \otimes_A A'$ 为 $B$ 的基变换。
\begin{enumerate}
\item 若 $B \otimes_A B \to B$ 平坦，则
$B' \otimes_{A'} B' \to B'$ 平坦。
\item 若 $A \to B$ 弱 \'{e}tale，则 $A' \to B'$ 弱 \'{e}tale。
\end{enumerate}
\end{lemma}

\begin{proof}
假设 $B \otimes_A B \to B$ 平坦。环映射
$B' \otimes_{A'} B' \to B'$ 是 $B \otimes_A B \to B$ 沿
$A \to A'$ 的基变换。因此，由《代数》引理
\ref{algebra-lemma-flat-base-change}，它是平坦的。这证明了 (1)。
(2) 由 (1) 以及刚才使用的平坦环映射的基变换仍平坦这一事实得出。
\end{proof}

\begin{lemma}
\label{more-algebra-lemma-absolutely-flat-over-absolutely-flat}
设 $A \to B$ 为环映射，并且 $B \otimes_A B \to B$ 平坦。
\begin{enumerate}
\item 若 $A$ 是绝对平坦环，则 $B$ 也如此。
\item 若 $A$ 既约且 $A \to B$ 弱 \'{e}tale，则 $B$ 既约。
\end{enumerate}
\end{lemma}

\begin{proof}
(1) 立即由引理 \ref{more-algebra-lemma-key} 与定义得出。若 $A$ 既约，则存在
单射
$$
A \to A' = \prod_{\mathfrak p \subset A\text{ minimal}} A_\mathfrak p
$$
，把 $A$ 嵌入一个绝对平坦环中（《代数》引理
\ref{algebra-lemma-reduced-ring-sub-product-fields} 与引理
\ref{more-algebra-lemma-product-fields-absolutely-flat}）。若 $A \to B$ 平坦，
则诱导映射 $B \to B' = B \otimes_A A'$ 也是单射。由引理
\ref{more-algebra-lemma-base-change-weakly-etale}，环映射 $A' \to B'$ 弱
\'{e}tale。由 (1)，$B'$ 绝对平坦。再由引理
\ref{more-algebra-lemma-absolutely-flat}，环 $B'$ 既约。因此 $B$ 既约。
\end{proof}

\begin{lemma}
\label{more-algebra-lemma-composition-weakly-etale}
设 $A \to B$ 与 $B \to C$ 为环映射。
\begin{enumerate}
\item 若 $B \otimes_A B \to B$ 与 $C \otimes_B C \to C$ 平坦，则
$C \otimes_A C \to C$ 平坦。
\item 若 $A \to B$ 与 $B \to C$ 弱 \'{e}tale，则 $A \to C$
弱 \'{e}tale。
\end{enumerate}
\end{lemma}

\begin{proof}
(1) 由乘法映射的分解
$$
C \otimes_A C \longrightarrow C \otimes_B C \longrightarrow C
$$
、等式
$$
C \otimes_B C = (C \otimes_A C) \otimes_{B \otimes_A B} B,
$$
、平坦映射的基变换仍平坦以及平坦环映射的复合仍平坦这些事实得出。
见《代数》引理 \ref{algebra-lemma-flat-base-change} 与
\ref{algebra-lemma-composition-flat}。(2) 由 (1) 以及刚才使用的
平坦环映射的复合仍平坦这一事实得出。
\end{proof}

\begin{lemma}
\label{more-algebra-lemma-go-down}
设 $A \to B \to C$ 为环映射。
\begin{enumerate}
\item 若 $B \to C$ 忠实平坦，且 $C \otimes_A C \to C$ 平坦，则
$B \otimes_A B \to B$ 平坦。
\item 若 $B \to C$ 忠实平坦，且 $A \to C$ 弱 \'{e}tale，则
$A \to B$ 弱 \'{e}tale。
\end{enumerate}
\end{lemma}

\begin{proof}
假设 $B \to C$ 忠实平坦且 $C \otimes_A C \to C$ 平坦。考虑交换图
% P02-E0522
$$
\xymatrix{
C \otimes_A C \ar[r] & C \\
B \otimes_A B \ar[r] \ar[u] & B \ar[u]
}
$$
。竖直箭头平坦，上方水平箭头平坦。因此 $C$ 作为
$B \otimes_A B$-模平坦。映射 $B \to C$ 忠实平坦，且
$C = B \otimes_B C$。故由《代数》引理
\ref{algebra-lemma-flatness-descends-more-general}，$B$ 作为
$B \otimes_A B$-模平坦。这证明了 (1)。(2) 由 (1) 以及如下事实
得出：若 $A \to C$ 平坦且 $B \to C$ 忠实平坦，则 $A \to B$
平坦（《代数》引理
\ref{algebra-lemma-flatness-descends-more-general}）。
\end{proof}

\begin{lemma}
\label{more-algebra-lemma-weakly-etale-permanence}
设 $A$ 为环。设 $B \to C$ 为弱 \'{e}tale $A$-代数之间的
$A$-代数映射。则 $B \to C$ 弱 \'{e}tale。
\end{lemma}

\begin{proof}
由引理 \ref{more-algebra-lemma-key}，环映射 $B \to C$ 平坦。环映射
$C \otimes_A C \to C \otimes_B C$ 满射，因而是上满射。于是，
% P02-E0516
由引理 \ref{more-algebra-lemma-key}，由于 $C$ 在 $C \otimes_A C$ 上平坦，
$C$ 在 $C \otimes_B C$ 上也平坦。
\end{proof}

\begin{lemma}
\label{more-algebra-lemma-formally-unramified}
设 $A \to B$ 为环映射，并且 $B \otimes_A B \to B$ 平坦。则
$\Omega_{B/A} = 0$，即 $B$ 在 $A$ 上形式非分歧。
\end{lemma}

\begin{proof}
设 $I \subset B \otimes_A B$ 为平坦满射
$B \otimes_A B \to B$ 的核。则 $I$ 是纯理想（《代数》定义
\ref{algebra-definition-pure-ideal}），故 $I^2 = I$（《代数》引理
\ref{algebra-lemma-pure}）。由于 $\Omega_{B/A} = I/I^2$
（《代数》引理 \ref{algebra-lemma-differentials-diagonal}），得到
所述消失。由《代数》引理
\ref{algebra-lemma-characterize-formally-unramified}，这意味着 $B$
在 $A$ 上形式非分歧。
\end{proof}

\begin{lemma}
\label{more-algebra-lemma-weakly-etale-finite-type}
设 $A \to B$ 为环映射，并且 $B \otimes_A B \to B$ 平坦。
\begin{enumerate}
\item 若 $A \to B$ 为有限型，则 $A \to B$ 非分歧。
\item 若 $A \to B$ 有限表现且平坦，则 $A \to B$ 是 \'{e}tale 的。
\end{enumerate}
特别地，有限表现的弱 \'{e}tale 环映射是 \'{e}tale 的。
\end{lemma}

\begin{proof}
(1) 由引理 \ref{more-algebra-lemma-formally-unramified} 与《代数》定义
\ref{algebra-definition-unramified} 得出。(2) 由 (1) 与《代数》引理
\ref{algebra-lemma-etale-flat-unramified-finite-presentation} 得出。
\end{proof}

\begin{lemma}
\label{more-algebra-lemma-when-weakly-etale}
设 $A \to B$ 为环映射。在以下每种情形中，$A \to B$ 都弱
\'{e}tale：
\begin{enumerate}
\item $B = S^{-1}A$ 是 $A$ 的局部化；
\item $A \to B$ 是 \'{e}tale 的；
\item $B$ 是弱 \'{e}tale $A$-代数的滤过余极限。
\end{enumerate}
\end{lemma}

\begin{proof}
\'{e}tale 环映射平坦，而且作为 \'{e}tale $A$-代数之间的映射，
$B \otimes_A B \to B$ 也是 \'{e}tale 的（《代数》引理
\ref{algebra-lemma-map-between-etale}）。这证明了 (2)。

\medskip\noindent
设 $B_i$ 为弱 \'{e}tale $A$-代数的有向系统。由《代数》引理
\ref{algebra-lemma-colimit-flat}，$B = \colim B_i$ 在 $A$ 上平坦。
注意，由引理 \ref{more-algebra-lemma-weakly-etale-permanence}，过渡映射
$B_i \to B_{i'}$ 平坦。因此，对每个 $i$，$B$ 在 $B_i$ 上平坦；
由《代数》引理 \ref{algebra-lemma-composition-flat}，可见 $B$
在 $B_i \otimes_A B_i$ 上平坦。故由《代数》引理
\ref{algebra-lemma-colimit-rings-flat}，$B$ 在
$B \otimes_A B = \colim B_i \otimes_A B_i$ 上平坦。

\medskip\noindent
(1) 可以直接证明，也可由 (2) 与 (3) 合并得到。
\end{proof}

\begin{lemma}
\label{more-algebra-lemma-absolutely-flat-fields}
设 $L/K$ 为域扩张。若 $L \otimes_K L \to L$ 平坦，则 $L$ 是
$K$ 的代数可分扩张。
\end{lemma}

\begin{proof}
由引理 \ref{more-algebra-lemma-go-down} 可见，
% P02-E0517
对任意中间域 $K \subset L' \subset L$，映射
$L' \otimes_K L' \to L'$ 都平坦。因此可假设 $L$ 是 $K$ 的有限生成
域扩张。在这种情形下，$L/K$ 形式非分歧（引理
\ref{more-algebra-lemma-formally-unramified}）这一事实蕴含 $L/K$ 有限可分；
见《代数》引理
\ref{algebra-lemma-characterize-separable-algebraic-field-extensions}。
\end{proof}

\begin{lemma}
\label{more-algebra-lemma-absolutely-flat-over-field}
设 $B$ 为域 $K$ 上的代数。下列条件等价：
\begin{enumerate}
\item $B \otimes_K B \to B$ 平坦；
\item $K \to B$ 弱 \'{e}tale；
\item $B$ 是 \'{e}tale $K$-代数的滤过余极限。
\end{enumerate}
此外，$B$ 的每个有限生成 $K$-子代数都在 $K$ 上 \'{e}tale。
\end{lemma}

\begin{proof}
(1) 与 (2) 等价，因为每个 $K$-代数都在 $K$ 上平坦。由引理
\ref{more-algebra-lemma-when-weakly-etale}，(3) 蕴含 (1) 与 (2)。
% P02-E0518

\medskip\noindent
假设 (1) 与 (2) 成立。我们将证明 (3) 以及引理中的有限性陈述。
% P02-E0519
域是绝对平坦环，故由引理
\ref{more-algebra-lemma-absolutely-flat-over-absolutely-flat}，$B$ 是绝对平坦环。
因此 $B$ 既约，且每个局部环都是域；见引理
\ref{more-algebra-lemma-absolutely-flat}。

\medskip\noindent
设 $\mathfrak q \subset B$ 为素理想。环映射
$B \to B_\mathfrak q$ 弱 \'{e}tale，故
$B_\mathfrak q$ 在 $K$ 上弱 \'{e}tale（引理
\ref{more-algebra-lemma-composition-weakly-etale}）。因此，由引理
\ref{more-algebra-lemma-absolutely-flat-fields}，$B_\mathfrak q$ 是 $K$ 的
可分代数扩张。

\medskip\noindent
设 $K \subset A \subset B$ 为有限生成 $K$-子代数。
% P02-E0520
我们将证明 $A$ 在 $K$ 上 \'{e}tale，这将完成引理的证明。
于是，每个极小素理想 $\mathfrak p \subset A$ 都是 $B$ 的
某个素理想 $\mathfrak q$ 的像；见《代数》引理
\ref{algebra-lemma-injective-minimal-primes-in-image}。因此，作为
$B_\mathfrak q = \kappa(\mathfrak q)$ 的子域，
$\kappa(\mathfrak p)$ 是 $K$ 的可分代数扩张。故
$\Spec(A)$ 的每个泛点都是闭点（《代数》引理
\ref{algebra-lemma-finite-residue-extension-closed}），从而
$\dim(A) = 0$。于是 $A$ 是其各局部环的乘积，例如可由《代数》命题
\ref{algebra-proposition-dimension-zero-ring} 得到。此外，由于 $A$
既约，所有局部环都等于各自的剩余域，而这些剩余域都是 $K$ 的有限
可分扩张。由《代数》引理 \ref{algebra-lemma-etale-over-field}，
这意味着 $A$ 在 $K$ 上 \'{e}tale，并完成证明。
\end{proof}

\begin{lemma}
\label{more-algebra-lemma-weakly-etale-residue-field-extensions}
设 $A \to B$ 为环映射。若 $A \to B$ 弱 \'{e}tale，则
$A \to B$ 诱导可分代数的剩余域扩张。
\end{lemma}

\begin{proof}
设 $\mathfrak p$ 为 $A$ 的素理想。由引理
\ref{more-algebra-lemma-base-change-weakly-etale}，
$\kappa(\mathfrak p) \to B \otimes_A \kappa(\mathfrak p)$ 弱
\'{e}tale。因此，由引理 \ref{more-algebra-lemma-absolutely-flat-over-field}，
$B \otimes_A \kappa(\mathfrak p)$ 是 \'{e}tale
$\kappa(\mathfrak p)$-代数的滤过余极限。故对位于
$\mathfrak p$ 上方的 $\mathfrak q \subset B$，由《代数》引理
\ref{algebra-lemma-etale-over-field}，扩张
$\kappa(\mathfrak q)/\kappa(\mathfrak p)$ 是有限可分扩张的滤过余极限。
\end{proof}

\begin{lemma}
\label{more-algebra-lemma-weak-dimension-at-most-1}
设 $A$ 为环。下列条件等价：
\begin{enumerate}
\item $A$ 的弱维数 $\leq 1$；
\item $A$ 的每个理想都平坦；
\item $A$ 的每个有限生成理想都平坦；
\item 平坦 $A$-模的每个子模都平坦；
\item $A$ 的每个局部环都是赋值环。
\end{enumerate}
\end{lemma}

\begin{proof}
若 $A$ 的弱维数 $\leq 1$，则解消
$0 \to I \to A \to A/I \to 0$ 与引理 \ref{more-algebra-lemma-last-one-flat}
表明每个理想 $I$ 都平坦。因此 (1) $\Rightarrow$ (2)。

\medskip\noindent
假设 (4)。设 $M$ 为 $A$-模。选取满射 $F \to M$，其中 $F$ 是自由
$A$-模。由假设，$\Ker(F \to M)$ 平坦；由引理
\ref{more-algebra-lemma-tor-dimension}，可见 $M$ 的 Tor 维数 $\leq 1$。
因此 (4) $\Rightarrow$ (1)。

\medskip\noindent
每个理想都是其中各有限生成理想的并。因此，由《代数》引理
\ref{algebra-lemma-colimit-flat}，(3) 蕴含 (2)。从而
(3) $\Leftrightarrow$ (2)。

\medskip\noindent
假设 (2)。设 $N \subset M$，其中 $M$ 为平坦 $A$-模。我们将证明
$N$ 平坦。由《代数》定理 \ref{algebra-theorem-lazard}，可写
$M = \colim M_i$，其中每个 $M_i$ 都是有限自由模。令
$N_i \subset M_i$ 为 $N$ 的逆像，可见 $N = \colim N_i$。由
% P02-E0521
《代数》引理 \ref{algebra-lemma-colimit-flat}，只需证明 $N_i$
平坦，于是归约到 $M = R^{\oplus n}$ 的情形。
% P02-E0523
在这种情形下，模 $N$ 具有由子模 $R^{\oplus j} \cap N$ 构成的
有限滤过，其逐次商都是理想。由 (2)，这些理想平坦；故由《代数》
引理 \ref{algebra-lemma-flat-ses}，$N$ 平坦。因此
(2) $\Rightarrow$ (4)。

\medskip\noindent
假设 $A$ 满足 (1)，并设 $\mathfrak p \subset A$ 为素理想。由引理
\ref{more-algebra-lemma-when-weakly-etale} 与
\ref{more-algebra-lemma-weak-dimension-goes-up}，可见 $A_\mathfrak p$ 满足 (1)。
下面证明：若 $A$ 是满足 (3) 的局部环，则 $A$ 是赋值环。设
$f \in \mathfrak m$ 为非零元素。则 $(f)$ 是由一个元素生成的
非零平坦模。故由《代数》引理
\ref{algebra-lemma-finite-flat-local}，它是自由 $A$-模。于是 $f$
是非零因子，且 $A$ 是整环。若 $I \subset A$ 是有限生成理想，
则同理可见 $I$ 是有限自由 $A$-模，因而（考察秩）是秩 $1$ 的自由模，
且 $I$ 是主理想。由《代数》引理
\ref{algebra-lemma-characterize-valuation-ring}，$A$ 是赋值环。
因此 (1) $\Rightarrow$ (5)。

\medskip\noindent
假设 (5)。设 $I \subset A$ 为有限生成理想。则
$I_\mathfrak p \subset A_\mathfrak p$ 是赋值环中的有限生成理想，
故为主理想（《代数》引理
\ref{algebra-lemma-characterize-valuation-ring}），从而平坦。因此，
由《代数》引理 \ref{algebra-lemma-flat-localization}，$I$ 平坦。
所以 (5) $\Rightarrow$ (3)。引理得证。
\end{proof}

\begin{lemma}
\label{more-algebra-lemma-product-weak-dimension-at-most-1}
设 $J$ 为集合。对每个 $j \in J$，设 $A_j$ 为分式域为 $K_j$ 的
赋值环。置 $A = \prod A_j$ 与 $K = \prod K_j$。则 $A$ 的弱维数
至多为 $1$，并且 $A \to K$ 是局部化。
\end{lemma}

\begin{proof}
设 $I \subset A$ 为有限生成理想。由引理
\ref{more-algebra-lemma-weak-dimension-at-most-1}，只需证明 $I$ 是平坦
$A$-模。设 $I_j \subset A_j$ 为 $I$ 的像。注意
$I_j = I \otimes_A A_j$，故由《代数》命题
\ref{algebra-proposition-fg-tensor}，$I \to \prod I_j$ 满射。因此
$I = \prod I_j$。由于 $A_j$ 是赋值环，理想 $I_j$ 由单个元素生成
（《代数》引理 \ref{algebra-lemma-characterize-valuation-ring}）。
设 $I_j = (f_j)$。则 $I$ 由元素 $f = (f_j)$ 生成。设 $e \in A$
为如下幂等元：根据 $f_j$ 是否为 $0$，它在 $A_j$ 中的分量为
$0$ 或 $1$。于是对某个非零因子 $g \in A$，有 $f = g e$：
取 $g = (g_j)$，其中若 $f_j = 0$，则 $g_j = 1$；否则
% P02-E0524
$g_j = f_j$。因此，作为模有 $I \cong (e)$。由于 $(e)$ 是 $A$
的直和项，故 $I$ 平坦。最后一个陈述成立，因为
$K = S^{-1}A$，其中 $S = \prod (A_j \setminus \{0\})$。
\end{proof}

\begin{lemma}
\label{more-algebra-lemma-product-found-valuation-rings}
设 $A$ 为分式域为 $K$ 的正规整环。存在环的笛卡尔图
$$
\xymatrix{
A \ar[d] \ar[r] & K \ar[d] \\
V \ar[r] & L
}
$$
，其中 $V$ 的弱维数至多为 $1$，且 $V \to L$ 是平坦、单射的
环满态射。
\end{lemma}

\begin{proof}
对每个 $x \in K$、$x \not \in A$，选取《代数》引理
% P02-E0525
\ref{algebra-lemma-find-valuation-rings} 中的 $V_x \subset K$。
置 $V = \prod_{x \in K \setminus A} V_x$ 与
$L = \prod_{x \in K \setminus A} K$。由引理
\ref{more-algebra-lemma-product-weak-dimension-at-most-1}，环 $V$ 的弱维数至多为
$1$；同一引理也表明 $V \to L$ 是局部化。局部化平坦且为满态射；
见《代数》引理 \ref{algebra-lemma-flat-localization} 与
\ref{algebra-lemma-epimorphism-local}。
\end{proof}

\begin{lemma}
\label{more-algebra-lemma-weak-dimension-at-most-1-integrally-closed}
设 $A$ 为弱维数至多为 $1$ 的环。若 $A \to B$ 是平坦、单射的
环满态射，则 $A$ 在 $B$ 中整闭。
\end{lemma}

\begin{proof}
设 $x \in B$ 在 $A$ 上整。令 $A' = A[x] \subset B$。由《代数》引理
\ref{algebra-lemma-characterize-finite-in-terms-of-integral}，
$A'$ 是 $A$ 的有限环扩张。由《代数》引理
\ref{algebra-lemma-finite-epimorphism-surjective}，为证明
$A = A'$，只需证明 $A \to A'$ 是满态射。由关于 $A$ 的假设、
$B$ 在 $A$ 上平坦这一事实以及引理
\ref{more-algebra-lemma-weak-dimension-at-most-1}，$A'$ 在 $A$ 上平坦。因此复合
$$
A' \otimes_A A' \to B \otimes_A A' \to B \otimes_A B \to B
$$
是单射，即 $A' \otimes_A A' \cong A'$，引理得证。
\end{proof}

\begin{lemma}
\label{more-algebra-lemma-normality-goes-up}
设 $A$ 为分式域为 $K$ 的正规整环。设 $A \to B$ 弱
\'{e}tale。则 $B$ 在 $B \otimes_A K$ 中整闭。
\end{lemma}

\begin{proof}
选取引理 \ref{more-algebra-lemma-product-found-valuation-rings} 中的图。由于
$A \to B$ 平坦，基变换给出环的笛卡尔图
$$
\xymatrix{
B \ar[d] \ar[r] & B \otimes_A K \ar[d] \\
B \otimes_A V \ar[r] & B \otimes_A L
}
$$
。注意，$V \to B \otimes_A V$ 弱 \'{e}tale（引理
\ref{more-algebra-lemma-base-change-weakly-etale}），故由引理
\ref{more-algebra-lemma-weak-dimension-goes-up}，$B \otimes_A V$ 的弱维数至多
为 $1$。又注意，$B \otimes_A V \to B \otimes_A L$ 是某个
平坦、单射的环满态射的平坦基变换，因而仍是平坦、单射的环满态射
（《代数》引理 \ref{algebra-lemma-flat-base-change} 与
\ref{algebra-lemma-base-change-epimorphism}）。由引理
\ref{more-algebra-lemma-weak-dimension-at-most-1-integrally-closed}，
$B \otimes_A V$ 在 $B \otimes_A L$ 中整闭。由该图的笛卡尔性质，
可知 $B$ 在 $B \otimes_A K$ 中整闭。
\end{proof}

\begin{lemma}
\label{more-algebra-lemma-integral-over-henselian}
设 $A \to B$ 为环同态。假设
\begin{enumerate}
\item $A$ 是 Hensel 局部环；
\item $A \to B$ 是整的；
\item $B$ 是整环。
\end{enumerate}
则 $B$ 是 Hensel 局部环，且 $A \to B$ 是局部同态。若 $A$
严格 Hensel，则 $B$ 是严格 Hensel 局部环，而且剩余域扩张
$\kappa(\mathfrak m_B)/\kappa(\mathfrak m_A)$ 纯不可分。
\end{lemma}

\begin{proof}
把 $B$ 写成有限 $A$-子代数的滤过余极限
$B = \colim B_i$。
% P02-E0520
若对每个 $B_i$ 都证明这些结论，则 $B$ 的结论随之成立；见《代数》
引理 \ref{algebra-lemma-colimit-henselian}。若 $A \to B$ 有限，则
由《代数》引理 \ref{algebra-lemma-finite-over-henselian}，$B$
是若干 Hensel 局部环的乘积。由于 $B$ 是整环，可见 $B$ 是局部环。
由 $A \to B$ 的上升定理（《代数》引理
\ref{algebra-lemma-integral-going-up}），$B$ 的极大理想位于 $A$
的极大理想上方。若 $A$ 严格 Hensel，则代数域扩张
$\kappa(\mathfrak m_B)/\kappa(\mathfrak m_A)$ 必为纯不可分。
于是当然 $\kappa(\mathfrak m_B)$ 可分代数闭，且 $B$ 严格 Hensel。
\end{proof}

\begin{theorem}[Olivier]
\label{more-algebra-theorem-olivier}
设 $A \to B$ 为局部环之间的局部同态。若 $A$ 严格 Hensel，且
$A \to B$ 弱 \'{e}tale，则 $A = B$。
\end{theorem}

\begin{proof}
我们将证明：对所有 $\mathfrak p \subset A$，存在唯一的
$\mathfrak q \subset B$ 位于 $\mathfrak p$ 上方，并且
$\kappa(\mathfrak p) = \kappa(\mathfrak q)$。这蕴含
$B \otimes_A B \to B$ 在谱上双射，而且该映射满射且平坦。因此，
例如由《代数》引理
\ref{algebra-lemma-pure-open-closed-specializations} 中对纯理想的
描述，它是同构。故 $A \to B$ 是忠实平坦的环满态射。由《代数》
引理 \ref{algebra-lemma-faithfully-flat-epimorphism}，得到
$A = B$。

\medskip\noindent
注意，由引理 \ref{more-algebra-lemma-base-change-weakly-etale} 与
\ref{more-algebra-lemma-absolutely-flat-over-field}，纤维环
$B \otimes_A \kappa(\mathfrak p)$ 是
$\kappa(\mathfrak p)$ 的 \'{e}tale 扩张的余极限。因此，若有不止
一个素理想位于 $\mathfrak p$ 上方，或对某个 $\mathfrak q$ 有
$\kappa(\mathfrak p) \not = \kappa(\mathfrak q)$，则对某个（可分）
代数域扩张 $L/\kappa(\mathfrak p)$，$B \otimes_A L$ 具有非平凡幂等元。

\medskip\noindent
设 $L/\kappa(\mathfrak p)$ 为代数域扩张。令 $A' \subset L$ 为
$A/\mathfrak p$ 在 $L$ 中的整闭包。由引理
\ref{more-algebra-lemma-integral-over-henselian}，$A'$ 是严格 Hensel 局部环，
其剩余域是 $A$ 的剩余域的纯不可分扩张。故由《代数》引理
\ref{algebra-lemma-local-tensor-with-integral}，
$B \otimes_A A'$ 是局部环。另一方面，由引理
\ref{more-algebra-lemma-normality-goes-up}，$B \otimes_A A'$ 在
$B \otimes_A L$ 中整闭。由于 $B \otimes_A A'$ 是局部环，可知
$B \otimes_A L$ 没有非平凡幂等元，这正是所要证明的。
\end{proof}

\section{域上的弱 \'{e}tale 代数}
\label{more-algebra-section-weakly-etale-over-field}

\noindent
若 $K$ 为域，则代数 $B$ 在 $K$ 上弱 \'{e}tale，当且仅当它是
\'{e}tale $K$-代数的滤过余极限。这就是引理
\ref{more-algebra-lemma-absolutely-flat-over-field}。

\begin{lemma}
\label{more-algebra-lemma-class-weakly-etale-over-field}
设 $K$ 为域。若 $B$ 在 $K$ 上弱 \'{e}tale，则
\begin{enumerate}
\item $B$ 既约；
\item $B$ 在 $K$ 上整；
\item $B$ 的任意有限生成 $K$-子代数都是 $K$ 的有限可分扩张的
有限乘积；
\item $B$ 是域，当且仅当 $B$ 没有非平凡幂等元；在这种情形下，
它是 $K$ 的可分代数扩张；
\item
% P02-E0526
$B$ 的任意 $K$-子代数或其商代数都在 $K$ 上弱 \'{e}tale；
\item 若 $B'$ 在 $K$ 上弱 \'{e}tale，则
$B \otimes_K B'$ 在 $K$ 上弱 \'{e}tale。
\end{enumerate}
\end{lemma}

\begin{proof}
(1) 由引理 \ref{more-algebra-lemma-absolutely-flat-over-absolutely-flat} 得出，当然
它也由 (3) 得出。(3) 由引理
\ref{more-algebra-lemma-absolutely-flat-over-field} 以及如下事实得出：
\'{e}tale $K$-代数是 $K$ 的有限可分扩张的有限乘积；见《代数》引理
\ref{algebra-lemma-etale-over-field}。(3) 蕴含 (2)。(4) 由 (3)
得出，因为域的乘积是域，当且仅当它没有非平凡幂等元。

\medskip\noindent
若 $S \subset B$ 是子代数，则它是其各有限生成子代数的滤过
余极限；由上文，这些子代数都在 $K$ 上 \'{e}tale，故由引理
\ref{more-algebra-lemma-absolutely-flat-over-field}，$S$ 在 $K$ 上弱
\'{e}tale。若 $B \to Q$ 是商代数，则 $Q$ 是有限可分扩张
$L_i/K$ 的有限乘积 $\prod_{i \in I} L_i$ 的 $K$-代数商的滤过
余极限。这种商具有形式 $\prod_{i \in J} L_i$，其中
$J \subset I$；故以同一论证，商代数的结论也成立。

\medskip\noindent
关于张量积的陈述可以类似证明，也可合并引理
\ref{more-algebra-lemma-base-change-weakly-etale} 与
\ref{more-algebra-lemma-composition-weakly-etale} 得到。
\end{proof}

\begin{lemma}
\label{more-algebra-lemma-max-weakly-etale-subalgebra}
设 $K$ 为域，$A$ 为 $K$-代数。存在极大的弱 \'{e}tale
$K$-子代数 $B_{max} \subset A$。
\end{lemma}

\begin{proof}
设 $B_1, B_2 \subset A$ 为弱 \'{e}tale $K$-子代数。则
$B_1 \otimes_K B_2$ 在 $K$ 上弱 \'{e}tale，而且
$B_1 \otimes_K B_2 \to A$ 的像也如此（引理
\ref{more-algebra-lemma-class-weakly-etale-over-field}）。因此，弱 \'{e}tale
$K$-子代数 $B \subset A$ 所成集合 $\mathcal{B}$ 是有向的；由引理
\ref{more-algebra-lemma-when-weakly-etale}，余极限
$B_{max} = \colim_{B \in \mathcal{B}} B$ 是弱 \'{e}tale $K$-代数。
因此，$B_{max} \to A$ 的像在 $K$ 上弱 \'{e}tale（引用前一个引理）。
于是该像属于 $\mathcal{B}$，故 $\mathcal{B}$ 有极大元（而该像
与 $B_{max}$ 相同）。
\end{proof}

\begin{lemma}
\label{more-algebra-lemma-properties-of-max-weakly-etale-subalgebra}
设 $K$ 为域。对 $K$-代数 $A$，以 $B_{max}(A)$ 表示引理
% P02-E0527
\ref{more-algebra-lemma-max-weakly-etale-subalgebra} 中 $A$ 的极大弱
\'{e}tale $K$-子代数。则
\begin{enumerate}
\item 任意 $K$-代数映射 $A' \to A$ 都诱导 $K$-代数映射
$B_{max}(A') \to B_{max}(A)$；
\item 若 $A' \subset A$，则
$B_{max}(A') = B_{max}(A) \cap A'$；
\item 若 $A = \colim A_i$ 为滤过余极限，则
$B_{max}(A) = \colim B_{max}(A_i)$；
\item 映射 $B_{max}(A) \to B_{max}(A_{red})$ 是同构；
\item $B_{max}(A_1 \times \ldots \times A_n) =
B_{max}(A_1) \times \ldots \times B_{max}(A_n)$；
\item 若 $A$ 没有非平凡幂等元，则 $B_{max}(A)$ 是域，并且是
$K$ 的可分代数扩张；
\item
% P02-E0528
在此添加更多内容。
\end{enumerate}
\end{lemma}

\begin{proof}
证明 (1)。这是因为，由引理
\ref{more-algebra-lemma-class-weakly-etale-over-field}，
$B_{max}(A') \to A$ 的像在 $K$ 上弱 \'{e}tale。

\medskip\noindent
证明 (2)。由 (1)，有 $B_{max}(A') \subset B_{max}(A)$。反之，由引理
\ref{more-algebra-lemma-class-weakly-etale-over-field}，
$B_{max}(A) \cap A'$ 是弱 \'{e}tale $K$-代数，因而包含在
$B_{max}(A')$ 中。

\medskip\noindent
证明 (3)。由 (1)，存在映射 $\colim B_{max}(A_i) \to A$；由于系统
滤过且 $B_{max}(A_i) \subset A_i$，该映射单射。由引理
\ref{more-algebra-lemma-when-weakly-etale}，余极限
$\colim B_{max}(A_i)$ 在 $K$ 上弱 \'{e}tale。因此得到单射
$\colim B_{max}(A_i) \to B_{max}(A)$。设
$a \in B_{max}(A)$。则 $a$ 生成有限表现 $K$-子代数
$B \subset B_{max}(A)$。由《代数》引理
\ref{algebra-lemma-characterize-finite-presentation}，存在 $i$ 以及
$K$-代数映射 $f : B \to A_i$，它提升给定映射 $B \to A$。由引理
\ref{more-algebra-lemma-class-weakly-etale-over-field}，$B$ 弱 \'{e}tale，
故 $f(B) \subset B_{max}(A_i)$；于是 $a$ 属于
$\colim B_{max}(A_i) \to B_{max}(A)$ 的像。

\medskip\noindent
证明 (4)。把 $B_{max}(A_{red}) = \colim B_i$ 写成
\'{e}tale $K$-代数的滤过余极限（引理
\ref{more-algebra-lemma-absolutely-flat-over-field}）。由《代数》引理
\ref{algebra-lemma-smooth-strong-lift}，对每个 $i$，存在
$K$-代数映射 $f_i : B_i \to A$，提升给定映射
$B_i \to A_{red}$。因此，
% P02-E0529
典范映射 $B_{max}(A_{red}) \to B_{max}(A)$ 满射。其核由幂零元组成，
因而为零，因为
% P02-E0530
$B_{max}(A_{red})$ 既约（引理
\ref{more-algebra-lemma-class-weakly-etale-over-field}）。

\medskip\noindent
证明 (5)。略。

\medskip\noindent
证明 (6)。这由引理 \ref{more-algebra-lemma-class-weakly-etale-over-field}
的 (4) 得出。
\end{proof}

\begin{lemma}
\label{more-algebra-lemma-change-fields-max-weakly-etale-subalgebra}
设 $L/K$ 为域扩张，$A$ 为 $K$-代数。设 $B \subset A$ 为引理
\ref{more-algebra-lemma-max-weakly-etale-subalgebra} 中 $A$ 的极大弱
\'{e}tale $K$-子代数。则 $B \otimes_K L$ 是
$A \otimes_K L$ 的极大弱 \'{e}tale $L$-子代数。
\end{lemma}

\begin{proof}
对 $K$ 上的代数 $A$，以 $B_{max}(A/K)$ 表示 $A$ 的极大弱
\'{e}tale $K$-子代数。类似地，若 $A'$ 是 $L$-代数，则以
$B_{max}(A'/L)$ 表示 $A'$ 的极大弱 \'{e}tale $L$-子代数。由于
$B_{max}(A/K) \otimes_K L$ 在 $L$ 上弱 \'{e}tale（引理
\ref{more-algebra-lemma-base-change-weakly-etale}），并且
$B_{max}(A/K) \otimes_K L \subset A \otimes_K L$，得到典范单射
% P02-E0533
$$
B_{max}(A/K) \otimes_K L \to B_{max}((A \otimes_K L)/L)
$$
。引理断言该映射是同构。

\medskip\noindent
为对 $L$ 与当前 $K$-代数 $A$ 证明引理，只需对 $L$ 的任意域扩张
$L'$ 证明引理。事实上，有分解
$$
B_{max}(A/K) \otimes_K L' \to
B_{max}((A \otimes_K L)/L) \otimes_L L' \to
B_{max}((A \otimes_K L')/L')
$$
，故若
$B_{max}(A/K) \otimes_K L \to B_{max}((A \otimes_K L)/L)$
不满射，则该复合不可能满射。因此可假设 $L$ 代数闭。

\medskip\noindent
归约到有限型 $K$-代数。可以把 $A$ 写成其各有限型
$K$-子代数的滤过余极限。
% P02-E0531
利用引理 \ref{more-algebra-lemma-properties-of-max-weakly-etale-subalgebra}，
可见只需对有限型 $K$-代数证明引理。

\medskip\noindent
假设 $A$ 是有限型 $K$-代数。由于
$A \to A_{red}$ 的核幂零，
$A \otimes_K L \to A_{red} \otimes_K L$ 的核也幂零。于是
$$
B_{max}((A \otimes_K L)/L) \to
B_{max}((A_{red} \otimes_K L)/L)
$$
单射，因为其核幂零，而弱 \'{e}tale $L$-代数
$B_{max}((A \otimes_K L)/L)$ 既约（引理
\ref{more-algebra-lemma-class-weakly-etale-over-field}）。由引理
\ref{more-algebra-lemma-properties-of-max-weakly-etale-subalgebra}，
$B_{max}(A/K) = B_{max}(A_{red}/K)$，故只需对 $A_{red}$ 证明引理。

\medskip\noindent
假设 $A$ 是既约有限型 $K$-代数。令 $Q = Q(A)$ 为 $A$ 的全商环。
则 $A \subset Q$，且
% P02-E0532
$A \otimes_K L \subset Q \otimes_A L$，从而
$$
B_{max}(A/K) = A \cap B_{max}(Q/K)
$$
以及
$$
B_{max}((A \otimes_K L)/L) =
(A \otimes_K L) \cap B_{max}((Q \otimes_K L)/L)
$$
；这是引理 \ref{more-algebra-lemma-properties-of-max-weakly-etale-subalgebra}
的结论。由于 $-\otimes_K L$ 是正合函子，若对 $Q$ 证明了结论，
则 $A$ 的结论随之成立。由于 $Q$ 是域的有限乘积（《代数》引理
\ref{algebra-lemma-total-ring-fractions-no-embedded-points}、
\ref{algebra-lemma-minimal-prime-reduced-ring}、
\ref{algebra-lemma-Noetherian-irreducible-components} 与
\ref{algebra-lemma-Noetherian-permanence}），而 $B_{max}$ 与乘积交换
（引理 \ref{more-algebra-lemma-properties-of-max-weakly-etale-subalgebra}），
只需在 $A$ 为域时证明引理。

\medskip\noindent
假设 $A$ 为域。利用证明第三段中的论证，可归约到 $A$ 在 $K$ 上
有限生成的情形。（事实上，我们归约到域的方式产生了 $K$ 的
有限生成域扩张。）

\medskip\noindent
假设 $A$ 是 $K$ 的有限生成域扩张。由引理
\ref{more-algebra-lemma-properties-of-max-weakly-etale-subalgebra} 的 (6)，
$K' = B_{max}(A/K)$ 是 $K$ 的可分代数域扩张。因此，$K'$ 是 $K$
的有限可分域扩张，并且由《代数》引理
\ref{algebra-lemma-make-geometrically-irreducible}，$A$ 在 $K'$
上几何不可约。由于 $L$ 代数闭且 $K'/K$ 有限可分，可见
$$
K' \otimes_K L \to \prod\nolimits_{\sigma \in \Hom_K(K', L)} L,\quad
\alpha \otimes \beta \mapsto (\sigma(\alpha)\beta)_\sigma
$$
是同构（《域》引理
\ref{fields-lemma-finite-separable-tensor-alg-closed}）。因此
$$
A \otimes_K L = A \otimes_{K'} (K' \otimes_K L) =
\prod\nolimits_{\sigma \in \Hom_K(K', L)} A \otimes_{K', \sigma} L
$$
。由于 $A$ 在 $K'$ 上几何不可约，可见
$A \otimes_{K', \sigma} L$ 有唯一的极小素理想。由于 $L$ 代数闭，
由引理 \ref{more-algebra-lemma-properties-of-max-weakly-etale-subalgebra} 的 (6)，
$B_{max}((A \otimes_{K', \sigma} L)/L) = L$，因为这个 $L$-代数是
$L$ 的代数域扩张。于是，$A \otimes_K L$ 的极大弱
\'{e}tale $K' \otimes_K L$-子代数是 $K' \otimes_K L$，因为可以
像上面一样把这些子代数分解为乘积。故包含
$K' \otimes_K L \subset B_{max}((A \otimes_K L)/L)$ 是等式：
由引理 \ref{more-algebra-lemma-weakly-etale-permanence}，环映射
$K' \otimes_K L \to B_{max}((A \otimes_K L)/L)$ 弱 \'{e}tale。
\end{proof}
\section{局部不可约性}
\label{more-algebra-section-unibranch}

\noindent
下述定义看来是普遍接受的定义。为理解它，注意若 $A \subset B$
是局部整环的整扩张，则由上升定理，$A \to B$ 是局部环同态
（《代数》引理 \ref{algebra-lemma-integral-going-up}）。

\begin{definition}
\label{more-algebra-definition-unibranch}
\begin{reference}
\cite[Chapter 0 (23.2.1)]{EGA4}
\end{reference}
设 $A$ 为局部环。若约化环 $A_{red}$ 是整环，并且 $A_{red}$
在其分式域中的整闭包 $A'$ 是局部环，则称 $A$ 为{\it 单支的}。
若 $A$ 单支，并且 $A'$ 的剩余域是 $A$ 的剩余域的纯不可分扩张，
则称 $A$ 为{\it 几何单支的}。
\end{definition}

\noindent
设 $A$ 为局部环。以下是等价的表述：
\begin{enumerate}
\item $A$ 单支，当且仅当 $A$ 有唯一的极小素理想 $\mathfrak p$，
并且 $A/\mathfrak p$ 在其分式域中的整闭包是局部环；
\item $A$ 几何单支，当且仅当 $A$ 有唯一的极小素理想
$\mathfrak p$，并且 $A/\mathfrak p$ 在其分式域中的整闭包是局部环，
而且该局部环的剩余域是 $A$ 的剩余域的纯不可分扩张。
\end{enumerate}
正规局部环是几何单支的（这由定义 \ref{more-algebra-definition-unibranch}
与《代数》定义 \ref{algebra-definition-ring-normal} 得出）。
引理 \ref{more-algebra-lemma-unibranch} 与 \ref{more-algebra-lemma-geometrically-unibranch}
表明，研究（几何）单支这一性质是合理的。

\begin{lemma}
\label{more-algebra-lemma-branches}
设 $A$ 为局部环，并假设 $A$ 只有有限多个极小素理想。令 $A'$
为 $A$ 在 $A_{red}$ 的全商环中的整闭包，令 $A^h$ 为 $A$
的 Hensel 化。考虑映射
$$
\Spec(A') \leftarrow \Spec((A')^h) \rightarrow \Spec(A^h)
$$
，其中 $(A')^h = A' \otimes_A A^h$。则：
\begin{enumerate}
\item 左侧箭头在极大理想上双射；
\item 右侧箭头在极小素理想上双射；
\item $(A')^h$ 的每个极小素理想都包含于唯一的极大理想中，
且每个极大理想恰好包含一个极小素理想。
\end{enumerate}
\end{lemma}

\begin{proof}
令 $I \subset A$ 为幂零元构成的理想。由《代数》引理
\ref{algebra-lemma-quotient-henselization}，有
$(A/I)^h = A^h/IA^h$。环 $A$、$A^h$、$A'$ 与 $(A')^h$ 的谱分别与
$A/I$、$A^h/IA^h$、$A'$ 及
$(A')^h = A' \otimes_{A/I} A^h/IA^h$ 的谱相同。因此可以用
$A_{red} = A/I$ 替换 $A$，并假设 $A$ 既约。于是 $A \subset A'$；
以下将不再特别说明这一点。

\medskip\noindent
证明 (1)。由于 $A'$ 在 $A$ 上整，$(A')^h$ 在 $A^h$ 上整。
由上升定理（《代数》引理 \ref{algebra-lemma-integral-going-up}），
$A'$ 与 $(A')^h$ 的每个极大理想分别位于 $A$ 的极大理想
$\mathfrak m$ 与 $A^h$ 的极大理想 $\mathfrak m^h$ 之上。
因此 (1) 由同构
$$
(A')^h \otimes_{A^h} \kappa^h =
A' \otimes_A A^h \otimes_{A^h} \kappa^h = A' \otimes_A \kappa
$$
得出，因为 $A \to A^h$ 所诱导的剩余域扩张
$\kappa^h/\kappa$ 是平凡的。下文将用到：所显示的环在一个域上整，
故其谱为 profinite 空间；见《代数》引理
\ref{algebra-lemma-integral-over-field} 与
\ref{algebra-lemma-ring-with-only-minimal-primes}。
% P02-E0534

\medskip\noindent
证明 (3)。环 $A'$ 是正规环；事实上，它是正规整环的有限乘积，
见《代数》引理
\ref{algebra-lemma-characterize-reduced-ring-normal}。由于 $A^h$
是 \'etale $A$-代数的滤过余极限，$(A')^h$ 是 \'etale
$A'$-代数的滤过余极限；故由《代数》引理
\ref{algebra-lemma-normal-goes-up} 与
\ref{algebra-lemma-colimit-normal-ring}，$(A')^h$ 是正规环。
% P02-E0535
因此 $(A')^h$ 的每个局部环都是正规整环，从而每个极大理想都包含
唯一的极小素理想。对 $A^h \to (A')^h$ 应用引理
\ref{more-algebra-lemma-integral-over-henselian-pair}，可见
$((A')^h, \mathfrak m(A')^h)$ 是 Hensel 对。若
$\mathfrak q \subset (A')^h$ 是极小素理想（或任意素理想），则由
引理 \ref{more-algebra-lemma-irreducible-henselian-pair-connected}，
$V(\mathfrak q)$ 与 $V(\mathfrak m (A')^h)$ 的交连通。又因
$V(\mathfrak m (A')^h) = \Spec((A')^h \otimes \kappa^h)$ 是
profinite 空间，可见存在唯一的极大理想包含 $\mathfrak q$。
% P02-E0536

\medskip\noindent
证明 (2)。$A'$ 的极小素理想恰为位于 $A$ 的极小素理想之上的素理想
（由构造可知）。由于 $A' \to (A')^h$ 平坦，由下降定理
（《代数》引理 \ref{algebra-lemma-flat-going-down}），$(A')^h$
的每个极小素理想都位于 $A'$ 的某个极小素理想之上。反之，
$(A')^h$ 中位于 $A'$ 的极小素理想之上的任意素理想都是极小的，
因为 $(A')^h$ 是 $A'$ 上 \'etale、因而拟有限的代数的滤过余极限
（略去一个小细节）。由此，$(A')^h$ 的极小素理想恰为位于 $A$
的极小素理想之上的那些素理想。类似地，$A^h$ 的极小素理想恰为
位于 $A$ 的极小素理想之上的素理想。由构造，
$A' \otimes_A Q(A) = Q(A)$，其中 $Q(A)$ 是既约局部环 $A$
的全商环。当然，$Q(A)$ 是 $A$ 的各极小素理想的剩余域的有限乘积。
因此
% P02-E0537
$$
(A')^h \otimes_A Q(A) = A^h \otimes_A A' \otimes_A Q(A) = A^h \otimes_A Q(A)
$$
。上述讨论表明，左侧环的谱是 $(A')^h$ 的极小素理想集合，
右侧环的谱是 $A^h$ 的极小素理想集合。这就完成了证明。
% P02-E0538
\end{proof}

\begin{lemma}
\label{more-algebra-lemma-unibranch}
\begin{reference}
\cite[Chapter IV Proposition 18.6.12]{EGA4}
\end{reference}
设 $A$ 为局部环，$A^h$ 为 $A$ 的 Hensel 化。以下条件等价：
\begin{enumerate}
\item $A$ 单支；
\item $A^h$ 有唯一的极小素理想。
\end{enumerate}
\end{lemma}

\begin{proof}
这由引理 \ref{more-algebra-lemma-branches} 得出，不过我们也给出直接证明。
以 $\mathfrak m$ 表示环 $A$ 的极大理想。回忆剩余域
$\kappa = A/\mathfrak m$ 与 $A^h$ 的剩余域相同。
% P02-E0539

\medskip\noindent
假设 (2) 成立。令 $\mathfrak p^h$ 为 $A^h$ 的唯一极小素理想。
$A \to A^h$ 的平坦性蕴含
$\mathfrak p = A \cap \mathfrak p^h$ 是 $A$ 的唯一极小素理想
（由下降定理，见《代数》引理
\ref{algebra-lemma-flat-going-down}）。又因
$A^h/\mathfrak pA^h = (A/\mathfrak p)^h$（见《代数》引理
\ref{algebra-lemma-quotient-henselization}）由引理
\ref{more-algebra-lemma-henselization-reduced} 是既约的，所以
$\mathfrak p^h = \mathfrak pA^h$。令 $A'$ 为
$A/\mathfrak p$ 在其分式域中的整闭包。我们须证明 $A'$ 是局部环。
由于 $A \to A'$ 是整映射，$A'$ 的每个极大理想都位于
$\mathfrak m$ 之上（由整环映射的上升定理，见《代数》引理
\ref{algebra-lemma-integral-going-up}）。若 $A'$ 不是局部环，则可取
两个不同的极大理想 $\mathfrak m_1$、$\mathfrak m_2$。选取元素
$f_1, f_2 \in A'$，使得 $f_i \in \mathfrak m_i$ 且
$f_i \not \in \mathfrak m_{3 - i}$。我们得到有限子代数
% P02-E0540
$B = A/\mathfrak p[f_1, f_2] \subset A'$，其极大理想
$B \cap \mathfrak m_i$（$i = 1, 2$）彼此不同。注意包含关系
$$
A/\mathfrak p \subset B \subset \kappa(\mathfrak p)
$$
与平坦环映射 $A \to A^h$ 张量后给出包含关系
$$
A^h/\mathfrak p^h \subset
B \otimes_A A^h \subset
\kappa(\mathfrak p) \otimes_A A^h \subset
\kappa(\mathfrak p^h)
$$
；最后一个包含成立，是因为
$\kappa(\mathfrak p) \otimes_A A^h =
\kappa(\mathfrak p) \otimes_{A/\mathfrak p} A^h/\mathfrak p^h$
是整环 $A^h/\mathfrak p^h$ 的局部化。注意
$B \otimes_A \kappa$ 至少有两个极大理想，因为
$B/\mathfrak mB$ 有两个极大理想。因此，由于 $A^h$ 是 Hensel 环，
$B \otimes_A A^h$ 是 $\geq 2$ 个局部环的乘积；见《代数》引理
\ref{algebra-lemma-mop-up}。但我们刚刚看到
$B \otimes_A A^h$ 是一个整环的子环，由此得到矛盾。

\medskip\noindent
假设 (1) 成立。令 $\mathfrak p \subset A$ 为唯一的极小素理想，
并令 $A'$ 为 $A/\mathfrak p$ 在其分式域中的整闭包。设
$A \to B$ 为局部环之间的局部映射，它诱导剩余域的同构，并且是
某个 \'etale $A$-代数的局部化。特别地，$\mathfrak m_B$ 是包含
$\mathfrak m B$ 的唯一素理想。于是 $B' = A' \otimes_A B$
在 $B$ 上整，而 $A \to A'$ 是局部映射这一假设蕴含 $B'$ 是局部环
（《代数》引理 \ref{algebra-lemma-local-tensor-with-integral}）。
另一方面，$A' \to B'$ 是 \'etale 环映射的局部化，故 $B'$ 正规；
见《代数》引理 \ref{algebra-lemma-normal-goes-up}。因此 $B'$
是（局部）正规整环。最后，有
% P02-E0537
$$
B/\mathfrak pB \subset B \otimes_A \kappa(\mathfrak p)
= B' \otimes_{A'} (\text{fraction field of }A') \subset
\text{fraction field of }B'
$$
。故 $B/\mathfrak pB$ 是整环，从而 $B$ 有唯一的极小素理想
（因为由 $A \to B$ 的平坦性，这些极小素理想全都必须位于
$\mathfrak p$ 之上）。由于 $A^h$ 是这些局部环 $B$ 的滤过余极限，
$A^h$ 有唯一的极小素理想。具体而言，若某些非幂零元
$f, g$ 在 $A^h$ 中满足 $fg = 0$，则可找到一个如上的 $B$
同时包含 $f$ 与 $g$，这会导致矛盾。
\end{proof}

\begin{lemma}
\label{more-algebra-lemma-geometric-branches}
设 $(A, \mathfrak m, \kappa)$ 为局部环，并假设 $A$ 只有有限多个
极小素理想。令 $A'$ 为 $A$ 在 $A_{red}$ 的全商环中的整闭包。
选取 $\kappa$ 的代数闭包 $\overline{\kappa}$，并以
$\kappa^{sep} \subset \overline{\kappa}$ 表示 $\kappa$
的可分代数闭包。令 $A^{sh}$ 为 $A$ 关于 $\kappa^{sep}$
的严格 Hensel 化。考虑映射
$$
\Spec(A') \xleftarrow{c} \Spec((A')^{sh}) \xrightarrow{e} \Spec(A^{sh})
$$
，其中 $(A')^{sh} = A' \otimes_A A^{sh}$。则：
\begin{enumerate}
\item 对极大理想 $\mathfrak m' \subset A'$，剩余域
$\kappa'$ 是 $\kappa$ 的代数扩张，且 $c$ 在 $\mathfrak m'$
上的纤维可典范地与
$\Hom_\kappa(\kappa', \overline{\kappa})$ 等同；
\item 右侧箭头在极小素理想上双射；
\item $(A')^{sh}$ 的每个极小素理想都包含于唯一的极大理想中，
且每个极大理想包含唯一的极小素理想。
\end{enumerate}
\end{lemma}

\begin{proof}
本证明几乎与引理 \ref{more-algebra-lemma-branches} 的证明完全相同。令
$I \subset A$ 为幂零元构成的理想。由《代数》引理
\ref{algebra-lemma-quotient-henselization}，有
$(A/I)^{sh} = A^{sh}/IA^{sh}$。环
% P02-E0541
$A$、$A^{sh}$、$A'$ 与 $(A')^h$ 的谱分别与
$A/I$、$A^{sh}/IA^{sh}$、$A'$ 及
$(A')^{sh} = A' \otimes_{A/I} A^{sh}/IA^{sh}$ 的谱相同。
因此可以用 $A_{red} = A/I$ 替换 $A$，并假设 $A$ 既约。于是
$A \subset A'$；以下将不再特别说明这一点。

\medskip\noindent
证明 (1)。域扩张 $\kappa'/\kappa$ 是代数的，因为 $A'$
在 $A$ 上整。由于 $A'$ 在 $A$ 上整，$(A')^{sh}$ 在
$A^{sh}$ 上整。由上升定理（《代数》引理
\ref{algebra-lemma-integral-going-up}），$A'$ 与 $(A')^{sh}$
的每个极大理想分别位于 $A$ 的极大理想 $\mathfrak m$ 与
% P02-E0542
$A^h$ 的极大理想 $\mathfrak m^{sh}$ 之上。我们有
$$
(A')^{sh} \otimes_{A^{sh}} \kappa^{sep} =
A' \otimes_A A^h \otimes_{A^h} \kappa^{sep} =
(A' \otimes_A \kappa) \otimes_{\kappa} \kappa^{sep}
$$
，因为 $A^{sh}$ 的剩余域是 $\kappa^{sep}$。因此 $c$
在 $\mathfrak m'$ 上的纤维是
$\kappa' \otimes_\kappa \kappa^{sep}$ 的谱。由双射
% P02-E0537
$$
\Hom_\kappa(\kappa', \overline{\kappa}) \to
\Spec(\kappa' \otimes_\kappa \kappa^{sep}),\quad
\sigma \mapsto \Ker(
\sigma \otimes 1 : \kappa' \otimes_\kappa \kappa^{sep} \to \overline{\kappa}
)
$$
，可知 (1) 成立。下文将用到：所显示的环在一个域上整，故其谱为
profinite 空间；见《代数》引理
\ref{algebra-lemma-integral-over-field} 与
\ref{algebra-lemma-ring-with-only-minimal-primes}。

\medskip\noindent
证明 (3)。环 $A'$ 是正规环；事实上，它是正规整环的有限乘积，
见《代数》引理
\ref{algebra-lemma-characterize-reduced-ring-normal}。由于
$A^{sh}$ 是 \'etale $A$-代数的滤过余极限，$(A')^{sh}$ 是
\'etale $A'$-代数的滤过余极限；故由《代数》引理
\ref{algebra-lemma-normal-goes-up} 与
\ref{algebra-lemma-colimit-normal-ring}，$(A')^{sh}$ 是正规环。
% P02-E0535
因此 $(A')^{sh}$ 的每个局部环都是正规整环，从而每个极大理想都包含
唯一的极小素理想。
对 $A^{sh} \to (A')^{sh}$ 应用引理
\ref{more-algebra-lemma-integral-over-henselian-pair}，可见
$((A')^{sh}, \mathfrak m(A')^{sh})$ 是 Hensel 对。若
$\mathfrak q \subset (A')^{sh}$ 是极小素理想（或任意素理想），
则由引理 \ref{more-algebra-lemma-irreducible-henselian-pair-connected}，
$V(\mathfrak q)$ 与 $V(\mathfrak m (A')^{sh})$ 的交连通。又因
% P02-E0543
$V(\mathfrak m (A')^{sh}) = \Spec((A')^{sh} \otimes \kappa^{sh})$
是 profinite 空间，可见存在唯一的极大理想包含 $\mathfrak q$。

\medskip\noindent
证明 (2)。$A'$ 的极小素理想恰为位于 $A$ 的极小素理想之上的素理想
（由构造可知）。由于 $A' \to (A')^{sh}$ 平坦，由下降定理
（《代数》引理 \ref{algebra-lemma-flat-going-down}），
$(A')^{sh}$ 的每个极小素理想都位于 $A'$ 的某个极小素理想之上。
反之，$(A')^{sh}$ 中位于 $A'$ 的极小素理想之上的任意素理想都是
极小的，因为 $(A')^{sh}$ 是 $A'$ 上 \'etale、因而拟有限的代数的
滤过余极限（略去一个小细节）。由此，$(A')^{sh}$ 的极小素理想
恰为位于 $A$ 的极小素理想之上的那些素理想。类似地，$A^{sh}$
的极小素理想恰为位于 $A$ 的极小素理想之上的素理想。由构造，
$A' \otimes_A Q(A) = Q(A)$，其中 $Q(A)$ 是既约局部环 $A$
的全商环。当然，$Q(A)$ 是 $A$ 的各极小素理想的剩余域的有限乘积。
因此
% P02-E0537
$$
(A')^{sh} \otimes_A Q(A) =
A^{sh} \otimes_A A' \otimes_A Q(A) = A^{sh} \otimes_A Q(A)
$$
。上述讨论表明，左侧环的谱是 $(A')^{sh}$ 的极小素理想集合，
右侧环的谱是 $A^{sh}$ 的极小素理想集合。这就完成了证明。
% P02-E0538
\end{proof}

\begin{lemma}
\label{more-algebra-lemma-geometrically-unibranch}
\begin{reference}
\cite[Lemma 2.2]{Etale-coverings} and
\cite[Chapter IV Proposition 18.8.15]{EGA4}
\end{reference}
设 $A$ 为局部环，$A^{sh}$ 为 $A$ 的一个严格 Hensel 化。
以下条件等价：
\begin{enumerate}
\item $A$ 几何单支；
\item $A^{sh}$ 有唯一的极小素理想。
\end{enumerate}
\end{lemma}

\begin{proof}
这由引理 \ref{more-algebra-lemma-geometric-branches} 得出，不过我们也给出一个
直接证明；它与引理 \ref{more-algebra-lemma-unibranch} 的直接证明几乎完全相同。
以 $\mathfrak m$ 表示环 $A$ 的极大理想，以
$\kappa$、$\kappa^{sh}$ 分别表示 $A$、$A^{sh}$ 的剩余域。
% P02-E0539

\medskip\noindent
假设 (2) 成立。令 $\mathfrak p^{sh}$ 为 $A^{sh}$ 的唯一极小素理想。
$A \to A^{sh}$ 的平坦性蕴含
$\mathfrak p = A \cap \mathfrak p^{sh}$ 是 $A$ 的唯一极小素理想
（由下降定理，见《代数》引理
\ref{algebra-lemma-flat-going-down}）。又因
$A^{sh}/\mathfrak pA^{sh} = (A/\mathfrak p)^{sh}$
（见《代数》引理
\ref{algebra-lemma-quotient-strict-henselization}）由引理
\ref{more-algebra-lemma-henselization-reduced} 是既约的，所以
$\mathfrak p^{sh} = \mathfrak pA^{sh}$。令 $A'$ 为
$A/\mathfrak p$ 在其分式域中的整闭包。我们须证明 $A'$ 是局部环，
且其剩余域是 $\kappa$ 的纯不可分扩张。由于 $A \to A'$ 是整映射，
$A'$ 的每个极大理想都位于 $\mathfrak m$ 之上（由整环映射的
上升定理，见《代数》引理
\ref{algebra-lemma-integral-going-up}）。若 $A'$ 不是局部环，则可取
两个不同的极大理想 $\mathfrak m_1$、$\mathfrak m_2$。选取元素
$f_1, f_2 \in A'$，使得 $f_i \in \mathfrak m_i, f_i \not \in \mathfrak m_{3 - i}$，
便得到有限子代数
$B = A[f_1, f_2] \subset A'$，其极大理想
$B \cap \mathfrak m_i$（$i = 1, 2$）彼此不同。
% P02-E0544
若 $A'$ 是以 $\mathfrak m'$ 为极大理想的局部环，但
$A/\mathfrak m \subset A'/\mathfrak m'$ 不是纯不可分扩张，
则可取 $f \in A'$，使它在 $A'/\mathfrak m'$ 中的像生成
$A/\mathfrak m$ 的一个有限但非纯不可分扩张；由此得到有限局部
子代数 $B = A[f] \subset A'$，其剩余域不是
$A/\mathfrak m$ 的纯不可分扩张。注意包含关系
$$
A/\mathfrak p \subset B \subset \kappa(\mathfrak p)
$$
与平坦环映射 $A \to A^{sh}$ 张量后给出包含关系
$$
A^{sh}/\mathfrak p^{sh} \subset
B \otimes_A A^{sh} \subset
\kappa(\mathfrak p) \otimes_A A^{sh} \subset
\kappa(\mathfrak p^{sh})
$$
；最后一个包含成立，是因为
$\kappa(\mathfrak p) \otimes_A A^{sh} =
\kappa(\mathfrak p) \otimes_{A/\mathfrak p} A^{sh}/\mathfrak p^{sh}$
是整环 $A^{sh}/\mathfrak p^{sh}$ 的局部化。注意
$B \otimes_A \kappa^{sh}$ 至少有两个极大理想，因为
$B/\mathfrak mB$ 或者有两个极大理想，或者只有一个极大理想而其
剩余域不是 $\kappa$ 的纯不可分扩张，并且 $\kappa^{sh}$
是可分代数闭的。因此，由于 $A^{sh}$ 是严格 Hensel 环，
$B \otimes_A A^{sh}$ 是 $\geq 2$ 个局部环的乘积；见《代数》引理
\ref{algebra-lemma-mop-up-strictly-henselian}。但我们刚刚看到
$B \otimes_A A^{sh}$ 是一个整环的子环，由此得到矛盾。

\medskip\noindent
假设 (1) 成立。令 $\mathfrak p \subset A$ 为唯一的极小素理想，
并令 $A'$ 为 $A/\mathfrak p$ 在其分式域中的整闭包。设
$A \to B$ 为局部环之间的局部映射，它是某个
\'etale $A$-代数的局部化。特别地，$\mathfrak m_B$ 是包含
$\mathfrak m_AB$ 的唯一素理想。于是 $B' = A' \otimes_A B$
在 $B$ 上整，而 $A \to A'$ 是局部映射且诱导纯不可分剩余域扩张
这一假设蕴含 $B'$ 是局部环（《代数》引理
\ref{algebra-lemma-local-tensor-with-integral}）。另一方面，
$A' \to B'$ 是 \'etale 环映射的局部化，故 $B'$ 正规；
见《代数》引理 \ref{algebra-lemma-normal-goes-up}。因此 $B'$
是（局部）正规整环。最后，有
% P02-E0537
$$
B/\mathfrak pB \subset B \otimes_A \kappa(\mathfrak p)
= B' \otimes_{A'} (\text{fraction field of }A') \subset
\text{fraction field of }B'
$$
。故 $B/\mathfrak pB$ 是整环，从而 $B$ 有唯一的极小素理想
（因为由 $A \to B$ 的平坦性，这些极小素理想全都必须位于
$\mathfrak p$ 之上）。由于 $A^{sh}$ 是这些局部环 $B$ 的滤过余极限，
$A^{sh}$ 有唯一的极小素理想。具体而言，若某些非幂零元
$f, g$ 在 $A^{sh}$ 中满足 $fg = 0$，则可找到一个如上的 $B$
同时包含 $f$ 与 $g$，这会导致矛盾。
\end{proof}

\begin{definition}
\label{more-algebra-definition-number-of-branches}
设 $A$ 为局部环，其 Hensel 化与严格 Hensel 化分别为
$A^h$、$A^{sh}$。若 $A^h$ 的极小素理想数有限，则定义
{\it $A$ 的分支数}为这一数目；否则定义为 $\infty$。
若 $A^{sh}$ 的极小素理想数有限，则定义
{\it $A$ 的几何分支数}为这一数目；否则定义为 $\infty$。
\end{definition}

\noindent
下面具体说明它与定义 \ref{more-algebra-definition-unibranch} 的关系。

\begin{lemma}
\label{more-algebra-lemma-number-of-branches-1}
设 $(A, \mathfrak m, \kappa)$ 为局部环。
\begin{enumerate}
\item 若 $A$ 有无限多个极小素理想，则 $A$ 的（几何）分支数为
$\infty$。
\item $A$ 的分支数为 $1$，当且仅当 $A$ 单支。
\item $A$ 的几何分支数为 $1$，当且仅当 $A$ 几何单支。
\end{enumerate}
假设 $A$ 只有有限多个极小素理想，并令 $A'$ 为 $A$ 在
$A_{red}$ 的全商环中的整闭包。则：
\begin{enumerate}
\item[(4)] $A$ 的分支数等于 $A'$ 的极大理想 $\mathfrak m'$ 的数目；
\item[(5)] 为得到 $A$ 的几何分支数，须按
$\kappa(\mathfrak m')/\kappa$ 的可分次数给出的重数，计数 $A'$
的每个极大理想 $\mathfrak m'$。
\end{enumerate}
\end{lemma}

\begin{proof}
本引理由各定义、引理 \ref{more-algebra-lemma-branches}、引理
\ref{more-algebra-lemma-geometric-branches} 以及《域》引理
\ref{fields-lemma-separable-degree} 立即得出。
\end{proof}

\begin{lemma}
\label{more-algebra-lemma-invariance-number-branches-smooth}
设 $A \to B$ 为局部环之间的局部同态，并且是某个光滑环映射的
局部化。
\begin{enumerate}
\item $A$ 的几何分支数等于 $B$ 的几何分支数。
\item 若 $A \to B$ 诱导剩余域的纯不可分扩张，则 $A$ 的分支数
等于 $B$ 的分支数。
\end{enumerate}
\end{lemma}

\begin{proof}
我们将用到：光滑环映射平坦（《代数》引理
\ref{algebra-lemma-smooth-syntomic}）；局部化平坦（《代数》引理
\ref{algebra-lemma-flat-localization}）；平坦环映射的复合平坦
（《代数》引理 \ref{algebra-lemma-composition-flat}）；平坦环映射
经基变换后仍平坦（《代数》引理
\ref{algebra-lemma-flat-base-change}）；平坦局部同态忠实平坦
（《代数》引理 \ref{algebra-lemma-local-flat-ff}）；（严格）
Hensel 化平坦（引理 \ref{more-algebra-lemma-dumb-properties-henselization}）；
以及平坦环映射的下降定理（《代数》引理
\ref{algebra-lemma-flat-going-down}）。
% P02-E0545

\medskip\noindent
证明 (2)。令 $A^h$、$B^h$ 分别为 $A$、$B$ 的 Hensel 化。
则 $B^h$ 是 $A^h \otimes_A B$ 在位于 $\mathfrak m_B$
之上的唯一极大理想处的 Hensel 化；见《代数》引理
\ref{algebra-lemma-henselian-functorial-improve}。因此可以并且确实
假设 $A$ 是 Hensel 环。由于 $A \to B \to B^h$ 平坦，
$B^h$ 的每个极小素理想都位于 $A$ 的某个极小素理想之上；
又由于 $A \to B^h$ 忠实平坦，$A$ 的每个极小素理想之上确有
$B^h$ 的某个极小素理想；两处都使用平坦环映射的下降定理。
因此只须证明：给定极小素理想 $\mathfrak p \subset A$，
$B^h$ 中至多有一个极小素理想位于 $\mathfrak p$ 之上。
用 $A/\mathfrak p$ 替换 $A$、用 $B/\mathfrak p B$ 替换 $B$ 后，
可假设 $A$ 为整环；此时由《代数》引理
\ref{algebra-lemma-quotient-henselization}，$A$ 仍是 Hensel 环。
% P02-E0546
由引理 \ref{more-algebra-lemma-unibranch}，$A$ 在其分式域中的整闭包
$A'$ 是局部整环。当然，$A'$ 是正规整环。由《代数》引理
\ref{algebra-lemma-normal-goes-up}，
$A' \otimes_A B^h$ 是正规环（该引理直接给出
$A' \otimes_A B$ 的结论；为推到 $A' \otimes_A B^h$，
使用 $B^h$ 是 \'etale $B$-代数的余极限以及《代数》引理
\ref{algebra-lemma-colimit-normal-ring}）。由《代数》引理
\ref{algebra-lemma-local-tensor-with-integral}，
$A' \otimes_A B^h$ 是局部环（这里使用了关于 $A$ 与 $B$
的剩余域的假设）。因此 $A' \otimes_A B^h$ 是局部正规环，
从而是局部整环。由 $A \to B^h$ 的平坦性，
$B^h \subset A' \otimes_A B^h$，故 $B^h$ 是整环，如所求。

\medskip\noindent
证明 (1)。令 $A^{sh}$、$B^{sh}$ 分别为 $A$、$B$ 的严格
Hensel 化。则 $B^{sh}$ 是
% P02-E0547
$A^h \otimes_A B$ 在某个同时位于 $\mathfrak m_B$
与 $\mathfrak m_{A^h}$ 之上的极大理想处的严格 Hensel 化；
见《代数》引理
\ref{algebra-lemma-strictly-henselian-functorial-improve}。
因此可以并且确实假设 $A$ 是严格 Hensel 环。由于
$A \to B \to B^{sh}$ 平坦，$B^{sh}$ 的每个极小素理想都位于
$A$ 的某个极小素理想之上；又由于 $A \to B^{sh}$ 忠实平坦，
$A$ 的每个极小素理想之上确有 $B^{sh}$ 的某个极小素理想；
两处都使用平坦环映射的下降定理。因此只须证明：给定极小素理想
$\mathfrak p \subset A$，$B^{sh}$ 中至多有一个极小素理想位于
$\mathfrak p$ 之上。用 $A/\mathfrak p$ 替换 $A$、用
$B/\mathfrak p B$ 替换 $B$ 后，可假设 $A$ 为整环；此时由
《代数》引理 \ref{algebra-lemma-quotient-strict-henselization}，
$A$ 仍是严格 Hensel 环。由引理
\ref{more-algebra-lemma-geometrically-unibranch}，$A$ 在其分式域中的整闭包
$A'$ 是局部整环，而且其剩余域是 $A$ 的剩余域的纯不可分扩张。
当然，$A'$ 是正规整环。由《代数》引理
\ref{algebra-lemma-normal-goes-up}，
$A' \otimes_A B^{sh}$ 是正规环（该引理直接给出
$A' \otimes_A B$ 的结论；为推到 $A' \otimes_A B^{sh}$，
使用 $B^{sh}$ 是 \'etale $B$-代数的余极限以及《代数》引理
\ref{algebra-lemma-colimit-normal-ring}）。由《代数》引理
\ref{algebra-lemma-local-tensor-with-integral}，
$A' \otimes_A B^{sh}$ 是局部环（因为 $A \subset A'$
诱导纯不可分剩余域扩张）。因此 $A' \otimes_A B^{sh}$
是局部正规环，从而是局部整环。由 $A \to B^{sh}$ 的平坦性，
$B^{sh} \subset A' \otimes_A B^{sh}$，故 $B^{sh}$
是整环，如所求。
\end{proof}
\section{分支杂论}
\label{more-algebra-section-branches}

\noindent
本节给出一些与第 \ref{more-algebra-section-unibranch} 节所定义的局部环分支有关的
结果。

\begin{lemma}
\label{more-algebra-lemma-ideal-inverse-nonzero}
设 $A$ 与 $B$ 为整环，$A \to B$ 为环映射。进一步假设
$A \to B$ 至少具有以下性质之一：
% P02-E0548
\begin{enumerate}
\item 它是某个 \'etale 环映射的局部化；
\item 它平坦，并且是某个非分歧环映射的局部化；
\item 它平坦，并且是某个拟有限环映射的局部化；
\item 它平坦，并且是某个整环映射的局部化；
\item 它平坦，并且 $\Spec(B) \to \Spec(A)$ 的各纤维中的点之间
不存在非平凡特殊化；
\item $\Spec(B) \to \Spec(A)$ 把泛点映到泛点，并且各纤维中的点
之间不存在非平凡特殊化；或者
\item $\Spec(B)$ 中恰有一个点映到 $\Spec(A)$ 的泛点。
\end{enumerate}
则对 $B$ 的每个非零理想 $J$，$A \cap J$ 都非零。
\end{lemma}

\begin{proof}
先证明情形 (7)。分别令 $K$ 与 $L$ 为 $A$ 与 $B$ 的分式域。
由《代数》引理
\ref{algebra-lemma-minimal-prime-image-minimal-prime}，
$\Spec(B)$ 中映到泛点
$(0) \in \Spec(A)$ 的唯一点是 $(0) \in \Spec(B)$。由此，
$B \otimes_A K$ 是仅有一个素理想且该素理想的剩余域为 $L$
的环（事实上它等于 $L$，但我们不需要这一点）。选取非零的
$b \in J$。则 $b$ 在 $L$ 中的像是单位。因此 $b$ 在
$B \otimes_A K$ 中的像是单位（《代数》引理
\ref{algebra-lemma-surjective-on-spec-units}）。由于
$B \otimes_A K = \colim_{f \in A \setminus \{0\}} B_f$，
存在非零的 $f \in A$，使 $b$ 在 $B_f$ 中的像为单位。
这意味着对某个 $b' \in B$ 与 $n \geq 1$，有
$b b' = f^n$。故 $f^n \in A \cap J$，如所求。

\medskip\noindent
在证明的其余部分，我们说明另外每项假设都蕴含 (7)。
% P02-E0549
在假设 (1)--(5) 下，环映射 $A \to B$ 平坦，因而
$A \to B$ 单射（因为由《代数》引理
\ref{algebra-lemma-local-flat-ff}，平坦局部同态忠实平坦）。
所以 $\Spec(B)$ 的泛点映到 $\Spec(A)$ 的泛点。现在，若
$\Spec(B) \to \Spec(A)$ 的各纤维中的点之间不存在非平凡特殊化，
则 $\Spec(B)$ 的这个泛点当然必须是映到 $\Spec(A)$ 泛点的唯一点。
故 (6) 蕴含 (7)。最后，为完成证明，我们说明在情形 (1)--(5) 中，
$\Spec(B) \to \Spec(A)$ 的各纤维中的点之间不存在非平凡特殊化。
具体而言，整映射情形见《代数》引理
\ref{algebra-lemma-integral-no-inclusion}，拟有限情形见《代数》定义
\ref{algebra-definition-quasi-finite}；并使用非分歧环映射与
\'etale 环映射都是拟有限的这一事实（《代数》引理
\ref{algebra-lemma-unramified-quasi-finite} 与
\ref{algebra-lemma-etale-quasi-finite}）。
\end{proof}

\begin{lemma}
\label{more-algebra-lemma-local-unramified-extension-unibranch-domain-is-etale}
设 $A \to B$ 为环映射，$\mathfrak q \subset B$ 为位于素理想
$\mathfrak p \subset A$ 之上的素理想。假设：
\begin{enumerate}
\item $A$ 是整环；
\item $A_\mathfrak p$ 几何单支；
\item $A \to B$ 在 $\mathfrak q$ 处非分歧；
\item $A_\mathfrak p \to B_\mathfrak q$ 单射。
\end{enumerate}
则存在 $g \in B$、$g \not \in \mathfrak q$，使得
$B_g$ 在 $A$ 上 \'etale。
\end{lemma}

\begin{proof}
由《代数》命题
\ref{algebra-proposition-unramified-locally-standard}，用一个
主局部化替换 $B$ 后，可以找到标准 \'etale 环映射
$A \to B'$ 与满射 $B' \to B$。以 $\mathfrak q' \subset B'$
表示 $\mathfrak q$ 的逆像。
% P02-E0550
我们将证明，必要时再用一个主局部化替换 $B'$ 后，
$B' \to B$ 是单射。

\medskip\noindent
本段将问题归约到 $B'$ 是整环的情形。因为 $A$ 是整环，
由《代数》引理 \ref{more-algebra-lemma-reduced-goes-up}，环 $B'$ 既约。
令 $K$ 为 $A$ 的分式域。则 $B' \otimes_A K$ 在一个域上
\'etale，因而是域的有限乘积；见《代数》引理
\ref{algebra-lemma-etale-over-field}。由于 $A \to B'$ 是
\'etale 的（因而平坦），$B'$ 的极小素理想位于
$(0) \subset A$ 之上（由平坦环映射的下降定理）。
% P02-E0551
由此 $B'$ 只有有限多个极小素理想，记作
$\mathfrak r_1, \ldots, \mathfrak r_r \subset B'$。由于
$A_\mathfrak p$ 几何单支且 $A \to B'$ 是 \'etale 的，
由引理 \ref{more-algebra-lemma-invariance-number-branches-smooth} 与
\ref{more-algebra-lemma-number-of-branches-1}，环 $B'_{\mathfrak q'}$
是整环。因此恰对一个 $i = i_0$，有
$\mathfrak q' \supset \mathfrak r_i$。由《代数》引理
\ref{algebra-lemma-silly}，可以选取 $g' \in B'$，使
$g' \not \in \mathfrak r_{i_0}$，但对 $i \not = i_0$ 有
$g' \in \mathfrak r_i$。分别用 $B'_{g'}$ 与 $B_{g'}$
替换 $B'$ 与 $B$ 后，$B'$ 成为整环。

\medskip\noindent
假设 $B'$ 是整环；特别地，$B' \subset B'_{\mathfrak q'}$。
若 $B' \to B$ 不单射，则
$J = \Ker(B'_{\mathfrak q'} \to B_\mathfrak q)$ 非零。
对 $A_\mathfrak p \to B'_{\mathfrak q'}$ 应用引理
\ref{more-algebra-lemma-ideal-inverse-nonzero}，得到非零元素
$a \in A_\mathfrak p$，它在 $B_\mathfrak q$ 中的像为零；
这与假设 (4) 矛盾。证明完成。
\end{proof}

\begin{lemma}
\label{more-algebra-lemma-unramified-extension-unibranch-domain-is-etale}
\begin{reference}
Generalization of \cite[Expose I, Theorem 9.5 part (ii)]{SGA1}
\end{reference}
设 $(A, \mathfrak m)$ 为几何单支局部整环。设
$A \to B$ 为局部环之间的单射局部同态，并且本质有限型。
若 $\mathfrak m B$ 是 $B$ 的极大理想，且所诱导的剩余域扩张可分，
则 $A \to B$ 是某个 \'etale 环映射的局部化。
\end{lemma}

\begin{proof}
可以写成 $B = C_\mathfrak q$，其中 $A \to C$ 是有限型环映射，
$\mathfrak q \subset C$ 是位于 $\mathfrak m$ 之上的素理想。
由《代数》引理 \ref{algebra-lemma-characterize-unramified}，
环映射 $A \to C$ 在 $\mathfrak q$ 处非分歧。由《代数》命题
\ref{algebra-proposition-unramified-locally-standard}，用一个
主局部化替换 $C$ 后，可以找到标准 \'etale 环映射
$A \to C'$ 与满射 $C' \to C$。以 $\mathfrak q' \subset C'$
表示 $\mathfrak q$ 的逆像，并置 $B' = C'_{\mathfrak q'}$。
则 $B' \to B$ 满射。只须证明 $B' \to B$ 也单射。

\medskip\noindent
由于 $A$ 是整环，由《代数》引理
\ref{more-algebra-lemma-reduced-goes-up}，环 $C'$ 与 $B'$ 既约。由于
$A$ 几何单支，由引理
\ref{more-algebra-lemma-invariance-number-branches-smooth} 与
\ref{more-algebra-lemma-number-of-branches-1}，环 $B'$ 是整环。
% P02-E0552
若 $B' \to B$ 不单射，则由引理
\ref{more-algebra-lemma-ideal-inverse-nonzero}，$A \cap \Ker(B' \to B)$ 非零，
这与 $A \to B$ 单射的假设矛盾。
\end{proof}

\begin{lemma}
\label{more-algebra-lemma-minimal-primes-tensor-strictly-henselian}
设 $k$ 为代数闭域。设 $A$、$B$ 为剩余域等于 $k$ 的严格
Hensel 局部 $k$-代数。令 $C$ 为 $A \otimes_k B$ 在极大理想
$\mathfrak m_A \otimes_k B + A \otimes_k \mathfrak m_B$
处的严格 Hensel 化。则 $C$ 的极小素理想与 $A$、$B$
的极小素理想对按 $1$ 对 $1$ 的方式对应。
% P02-E0553
\end{lemma}

\begin{proof}
首先注意，$C$ 的极小素理想 $\mathfrak r$ 在 $A$ 中的像是
极小素理想 $\mathfrak p$，在 $B$ 中的像是极小素理想
$\mathfrak q$，因为环映射 $A \to C$ 与 $B \to C$ 平坦
（由平坦环映射的下降定理，《代数》引理
\ref{algebra-lemma-flat-going-down}）。
% P02-E0554
因此只须证明
% P02-E0555
$(A/\mathfrak p \otimes_k B/\mathfrak q)_{
\mathfrak m_A \otimes_k B + A \otimes_k \mathfrak m_B}$
的严格 Hensel 化有唯一的极小素理想。由《代数》引理
\ref{algebra-lemma-quotient-strict-henselization}，环
$A/\mathfrak p$、$B/\mathfrak q$ 严格 Hensel。因此可假设
$A$ 与 $B$ 是严格 Hensel 局部整环；我们的目标是证明 $C$
有唯一的极小素理想。由引理
\ref{more-algebra-lemma-geometrically-unibranch}，$A$ 在其分式域中的整闭包
$A'$ 是剩余域为 $k$ 的正规局部整环。$B$ 在其分式域中的整闭包
$B'$ 也一样。
% P02-E0556
由《代数》引理
\ref{algebra-lemma-geometrically-normal-tensor-normal}，
$A' \otimes_k B'$ 是正规环。因此其局部化
% P02-E0558
$$
R = (A' \otimes_k B')_{
\mathfrak m_{A'} \otimes_k B' + A' \otimes_k \mathfrak m_{B'}}
$$
是正规局部整环。注意 $A \otimes_k B \to A' \otimes_k B'$
是整映射（因而上升定理成立——《代数》引理
\ref{algebra-lemma-integral-going-up}），并且
% P02-E0557
$\mathfrak m_{A'} \otimes_k B' + A' \otimes_k \mathfrak m_{B'}$
是 $A' \otimes_k B'$ 中位于
$\mathfrak m_A \otimes_k B + A \otimes_k \mathfrak m_B$
之上的唯一极大理想。因此，由《代数》引理
\ref{algebra-lemma-unique-prime-over-localize-below}，
% P02-E0558
$$
R = (A' \otimes_k B')_{
\mathfrak m_A \otimes_k B + A \otimes_k \mathfrak m_B}
$$
。由此，映射
% P02-E0558
$$
(A \otimes_k B)_{
\mathfrak m_A \otimes_k B + A \otimes_k \mathfrak m_B}
\longrightarrow
R
$$
是整映射。因此，$R$ 是
$(A \otimes_k B)_{
\mathfrak m_A \otimes_k B + A \otimes_k \mathfrak m_B}$
在其分式域中的整闭包；再次应用引理
\ref{more-algebra-lemma-geometrically-unibranch}，可知 $C$ 有唯一的素理想。
% P02-E0559
\end{proof}
\section{完备化的分支}
\label{more-algebra-section-branches-completion}

\noindent
设 $(A, \mathfrak m)$ 为 Noether 局部环。考虑映射
$A \to A^h \to A^\wedge$。一般而言，映射
$A^h \to A^\wedge$ 未必诱导极小素理想之间的双射；见《例》一书
第 \ref{examples-section-bad} 节。换言之，$A$ 的分支数
（定义见 \ref{more-algebra-definition-number-of-branches}）可能不同于
$A^\wedge$ 的分支数。不过，在某些条件下两者相同；例如，
$A$ 的维数为 $1$ 时即如此。

\begin{lemma}
\label{more-algebra-lemma-nr-branches-completion}
设 $(A, \mathfrak m)$ 为 Noether 局部环。
\begin{enumerate}
\item 映射 $A^h \to A^\wedge$ 定义了从 $A^\wedge$ 的极小素理想
到 $A^h$ 的极小素理想的满射。
\item $A$ 的分支数不大于 $A^\wedge$ 的分支数。
\item $A$ 的几何分支数不大于 $A^\wedge$ 的几何分支数。
\end{enumerate}
\end{lemma}

\begin{proof}
由引理 \ref{more-algebra-lemma-henselization-noetherian}，映射
$A^h \to A^\wedge$ 平坦且单射。结合下降定理（《代数》引理
\ref{algebra-lemma-flat-going-down}）与《代数》引理
\ref{algebra-lemma-injective-minimal-primes-in-image}，可得 (1)。
(2) 由此、定义 \ref{more-algebra-definition-number-of-branches} 以及
$A^\wedge$ 是 Hensel 环这一事实得出（《代数》引理
\ref{algebra-lemma-complete-henselian}）。由引理
\ref{more-algebra-lemma-henselization-noetherian}，有
$(A^\wedge)^{sh} = A^{sh} \otimes_{A^h} A^\wedge$。因此，使用平坦
单射映射 $A^{sh} \to (A^\wedge)^{sh}$ 重复上述论证，即得 (3)。
\end{proof}

\begin{lemma}
\label{more-algebra-lemma-equal-nr-branches-completion}
设 $(A, \mathfrak m)$ 为 Noether 局部环。$A$ 的分支数等于
$A^\wedge$ 的分支数，当且仅当对 Hensel 化的每个极小素理想
$\mathfrak q \subset A^h$，$\sqrt{\mathfrak qA^\wedge}$ 都是素理想。
\end{lemma}

\begin{proof}
这由引理 \ref{more-algebra-lemma-nr-branches-completion} 以及如下事实得出：
由《代数》引理
\ref{algebra-lemma-Noetherian-irreducible-components} 及
$A^h$ 与 $A^\wedge$ 都是 Noether 环这一事实（引理
\ref{more-algebra-lemma-henselization-noetherian}），$A$ 与 $A^\wedge$
都只有有限多个分支。
% P02-E0560
\end{proof}

\noindent
下面是一个简单的黏合引理。

\begin{lemma}
\label{more-algebra-lemma-glueing-sum-components-open}
设 $A$ 为环，$I$ 为有限生成理想。设 $A \to C$ 为环映射，
使得对每个 $f \in I$，环映射 $A_f \to C_f$ 都是关于某个幂等元
的局部化。则存在满射 $A \to C'$，使得对每个 $f \in I$，
$A_f \to (C \times C')_f$ 都是同构。
\end{lemma}

\begin{proof}
选取 $I$ 的一组生成元 $f_1, \ldots, f_r$。写成
$$
C_{f_i} = (A_{f_i})_{e_i}
$$
，其中 $e_i \in A_{f_i}$ 为某个幂等元。对某个
$a_i \in A$ 与 $n \geq 0$，写 $e_i = a_i/f_i^n$；对所有
$i = 1, \ldots, r$ 可以使用同一个 $n$。取适当的
$m \gg 0$，用 $f_i^ma_i$ 替换 $a_i$、用 $n + m$ 替换 $n$
之后，可以假设对所有 $i$ 都有 $a_i^2 = f_i^n a_i$。由于
$e_i$ 在 $C_{f_if_j} = (A_{f_if_j})_{e_j} = A_{f_if_ja_j}$
中的像为 $1$，可见对某个 $N$ 有
$$
(f_if_ja_j)^N(f_j^n a_i  - f_i^na_j) = 0
$$
（对所有序对 $i, j$ 可以选取同一个 $N$）。利用
$a_j^2 = f_j^na_j$，得到
% P02-E0561
$$
f_i^{N + n} f_j^{N + nN} a_j = f_i^N f_j^{N + n} a_ia_j^N
$$
。把 $n$ 增大为 $n + N + nN$ 并用 $f_i^{N + nN}a_i$
替换 $a_i$ 后，可见对所有序对 $i, j$，$f_i^n a_j$
都属于由 $a_i$ 生成的理想。令
$C' = A/(a_1, \ldots, a_r)$。则
$$
C'_{f_i} = A_{f_i}/(a_i) = A_{f_i}/(e_i)
$$
，因为反演 $f_i$ 后，$a_j$ 属于由 $a_i$ 生成的理想。对环 $B$
的幂等元 $e$，有 $B = B_e \times B/(e)$，故当 $f$ 等于
$f_1, \ldots, f_r$ 之一时，引理的结论成立。利用函数的黏合，
具体使用《代数》引理 \ref{algebra-lemma-cover}，可知结论对所有
$f \in I$ 都成立。事实上，对 $f \in I$，元素
$f_1, \ldots, f_r$ 在 $A_f$ 中生成单位理想，故
$A_f \to (C \times C')_f$ 是同构，当且仅当在分别关于
$f_1, \ldots, f_r$ 局部化后它是同构。
\end{proof}

\noindent
可以利用引理 \ref{more-algebra-lemma-quotient-by-idempotent}，从形式完备化的
给定有限型扩张构造有限型扩张。我们将在《形式空间的代数化》引理
\ref{restricted-lemma-approximate-by-etale-over-complement}
中推广这一引理。

\begin{lemma}
\label{more-algebra-lemma-quotient-by-idempotent}
设 $A$ 为 Noether 环，$I$ 为理想，$B$ 为有限型 $A$-代数。
设 $B^\wedge \to C$ 为核是 $J$ 的满射，其中 $B^\wedge$
是 $I$-进完备化。若对某个 $c \geq 0$，$J/J^2$ 被 $I^c$ 零化，
则 $C$ 同构于某个有限型 $A$-代数的完备化。
\end{lemma}

\begin{proof}
令 $f \in I$。由于 $B^\wedge$ 是 Noether 环（《代数》引理
\ref{algebra-lemma-completion-Noetherian-Noetherian}），
$J$ 是有限生成理想。因此，由《代数》引理
\ref{algebra-lemma-ideal-is-squared-union-connected} 可知
$$
C_f = ((B^\wedge)_f)_e
$$
，其中 $e \in (B^\wedge)_f$ 是某个幂等元。由引理
\ref{more-algebra-lemma-glueing-sum-components-open}，可以找到满射
$B^\wedge \to C'$，使得反演任意 $f \in I$ 后，
$B^\wedge \to C \times C'$ 都成为同构。注意
$C \times C'$ 是有限 $B^\wedge$-代数。

\medskip\noindent
选取生成元 $f_1, \ldots, f_r \in I$。以
$\alpha_i : (C \times C')_{f_i} \to B_{f_i} \otimes_B B^\wedge$
表示上面得到的 $(B^\wedge)_{f_i}$-代数同构的逆，以
$\alpha_{ij} : (B_{f_i})_{f_j} \to (B_{f_j})_{f_i}$
表示显然的 $B$-代数同构。
% P02-E0562
考虑注记 \ref{more-algebra-remark-glueing-data} 中所引入的范畴
$\text{Glue}(B \to B^\wedge, f_1, \ldots, f_r)$ 的对象
$$
(C \times C', B_{f_i}, \alpha_i, \alpha_{ij})
$$
。我们略去条件 (1)(a) 与 (1)(b) 的验证。因为
$B \to B^\wedge$ 是平坦映射（《代数》引理
\ref{algebra-lemma-completion-flat}），并诱导同构
$B/IB \to B^\wedge/IB^\wedge$，所以可以应用命题
\ref{more-algebra-proposition-equivalence} 与注记
\ref{more-algebra-remark-formal-glueing-algebras}。由此，对某个有限
$B$-代数 $D$，$C \times C'$ 同构于 $D \otimes_B B^\wedge$。
于是 $D/ID \cong C/IC \times C'/IC'$。令
$\overline{e} \in D/ID$ 为对应于因子 $C/IC$ 的幂等元。由引理
\ref{more-algebra-lemma-lift-idempotent-upstairs}，存在 \'etale 环映射
$B \to B'$，它诱导同构 $B/IB \to B'/IB'$，且
$D' = D \otimes_B B'$ 包含提升 $\overline{e}$ 的幂等元 $e$。
由于 $C \times C'$ 是 $I$-进完备的，对
$(C \times C', IC \times IC')$ 是 Hensel 对（引理
\ref{more-algebra-lemma-complete-henselian}）。因此，可以把映射
$B \to C \times C'$ 分解经过 $B'$（见引理
\ref{more-algebra-lemma-characterize-henselian-pair}）。这样做以后，可以用
$B'$ 替换 $B$、用 $D'$ 替换 $D$（因为这不改变 $I$-进完备化）。
% P02-E0563
于是
$D = D_e \times D_{1 - e} = D/(1 - e) \times D/(e)$
是有限型 $A$-代数的乘积；第一个因子的完备化是 $C$，第二个因子的
完备化是 $C'$。
\end{proof}

\begin{lemma}
\label{more-algebra-lemma-one-dimensional-formal-branch}
设 $(A, \mathfrak m)$ 为 Noether 局部环，其 Hensel 化为 $A^h$。
设 $\mathfrak q \subset A^\wedge$ 为极小素理想，且
$\dim(A^\wedge/\mathfrak q) = 1$。则存在 $A^h$ 的极小素理想
$\mathfrak q^h$，使得
$\mathfrak q = \sqrt{\mathfrak q^hA^\wedge}$。
\end{lemma}

\begin{proof}
由于 $A$ 与 $A^h$ 的完备化相同，可以假设 $A$ 是 Hensel 环
（引理 \ref{more-algebra-lemma-henselization-noetherian}）。我们将对映射
$A^\wedge \to A^\wedge/J$ 应用引理
\ref{more-algebra-lemma-quotient-by-idempotent}，其中
$J = \Ker(A^\wedge \to (A^\wedge)_{\mathfrak q})$。由于
$\dim((A^\wedge)_\mathfrak q) = 0$，对某个 $n$ 有
$\mathfrak q^n \subset J$。所以 $J/J^2$ 被 $\mathfrak q^n$
零化。另一方面，因 $J_\mathfrak q = 0$，有
$(J/J^2)_\mathfrak q = 0$。故 $\mathfrak m$ 是 $J/J^2$
唯一的伴随素理想，从而 $\mathfrak m$ 的某个幂零化 $J/J^2$。
于是该引理适用，得到
$A^\wedge/J = C^\wedge$，其中 $C$ 是某个有限型 $A$-代数。

\medskip\noindent
由于 $A^\wedge/J$ 具有相应性质，有
$C/\mathfrak m C = A/\mathfrak m$。因此
$\mathfrak m_C = \mathfrak m C$ 是极大理想，且
$A \to C$ 在 $\mathfrak m_C$ 处非分歧（《代数》引理
\ref{algebra-lemma-characterize-unramified}）。用一个主局部化替换
$C$ 后，可以假设 $C$ 是某个 \'etale $A$-代数 $B$ 的商；
见《代数》命题
\ref{algebra-proposition-unramified-locally-standard}。不过，
由于 $A \to C_{\mathfrak m_C}$ 的剩余域扩张平凡且 $A$
是 Hensel 环，再作一次局部化后可知 $B = A$。因此对某个理想
$I \subset A$，有 $C = A/I$，从而
$J = IA^\wedge$（因为在当前情形下完备化正合，见《代数》引理
\ref{algebra-lemma-completion-flat}），并且
$I = J \cap A$（由 $A \to A^\wedge$ 的平坦性）。由于
$\mathfrak q^n \subset J \subset \mathfrak q$，可知
$\mathfrak p = \mathfrak q \cap A$ 满足
$\mathfrak p^n \subset I \subset \mathfrak p$。于是
$\sqrt{\mathfrak p A^\wedge} = \mathfrak q$，证明完成。
\end{proof}

\begin{lemma}
\label{more-algebra-lemma-completion-disconnected}
设 $(A, \mathfrak m)$ 为 Noether 局部环。$A^\wedge$ 的去心谱不连通，
当且仅当 $A^h$ 的去心谱不连通。
\end{lemma}

\begin{proof}
由于 $A$ 与 $A^h$ 的完备化相同，可以假设 $A$ 是 Hensel 环
（引理 \ref{more-algebra-lemma-henselization-noetherian}）。

\medskip\noindent
由于 $A \to A^\wedge$ 忠实平坦（见刚才给出的参考文献），
从 $A^\wedge$ 的去心谱到 $A$ 的去心谱的映射是满射
（见《代数》引理 \ref{algebra-lemma-ff-rings}）。因此，若
$A$ 的去心谱不连通，则 $A^\wedge$ 的去心谱也不连通。

\medskip\noindent
假设 $A^\wedge$ 的去心谱不连通。这意味着
% P02-E0564
$$
\Spec(A^\wedge) \setminus \{\mathfrak m^\wedge\} = Z \amalg Z'
$$
，其中 $Z$ 与 $Z'$ 闭。令
% P02-E0565
$\overline{Z}, \overline{Z}' \subset \Spec(A^\wedge)$ 为其闭包。
设对某些理想 $J, J' \subset A^\wedge$ 有
$\overline{Z} = V(J)$、$\overline{Z}' = V(J')$。则
$V(J + J') = \{\mathfrak m^\wedge\}$ 且
$V(JJ') = \Spec(A^\wedge)$。第一个等式意味着
$\mathfrak m^\wedge = \sqrt{J + J'}$，从而对某个
$e \geq 1$ 有
$(\mathfrak m^\wedge)^e \subset J + J'$。第二个等式蕴含
$JJ'$ 的每个元素都幂零，故对某个 $n \geq 1$ 有
$(JJ')^n = 0$。合在一起，这意味着 $J^n/J^{2n}$ 被
$J^n$ 与 $(J')^n$ 零化，因而被
$(\mathfrak m^\wedge)^{2en}$ 零化。
% P02-E0566
于是可以应用引理 \ref{more-algebra-lemma-quotient-by-idempotent}，得到有限型
$A$-代数 $C$ 与同构 $A^\wedge/J^n = C^\wedge$。

\medskip\noindent
证明的其余部分与引理
\ref{more-algebra-lemma-one-dimensional-formal-branch} 证明的第二部分完全相同；
当然，该引理正是本引理的一个特殊情形！由于 $A^\wedge/J^n$
具有相应性质，有 $C/\mathfrak m C = A/\mathfrak m$。因此
$\mathfrak m_C = \mathfrak m C$ 是极大理想，且
$A \to C$ 在 $\mathfrak m_C$ 处非分歧（《代数》引理
\ref{algebra-lemma-characterize-unramified}）。用一个主局部化替换
$C$ 后，可以假设 $C$ 是某个 \'etale $A$-代数 $B$ 的商；
见《代数》命题
\ref{algebra-proposition-unramified-locally-standard}。不过，
由于 $A \to C_{\mathfrak m_C}$ 的剩余域扩张平凡且 $A$
是 Hensel 环，再作一次局部化后可知 $B = A$。因此对某个理想
$I \subset A$，有 $C = A/I$，从而
$J^n = IA^\wedge$（因为在当前情形下完备化正合，见《代数》引理
\ref{algebra-lemma-completion-flat}），并且
$I = J^n \cap A$（由 $A \to A^\wedge$ 的平坦性）。
由对称性，$I' = (J')^n \cap A$ 满足 $(J')^n = I'A^\wedge$。
于是 $\mathfrak m^e \subset I + I'$ 且 $II' = 0$，从而
$V(I)$ 与 $V(I')$ 是闭子概形，并给出 $A$ 的去心谱所需的不交并分解。
% P02-E0567
\end{proof}

\begin{lemma}
\label{more-algebra-lemma-one-dimensional-number-of-branches}
设 $(A, \mathfrak m)$ 为维数 $1$ 的 Noether 局部环。则
$A$ 与 $A^\wedge$ 的（几何）分支数相同。
\end{lemma}

\begin{proof}
为说明分支数的结论，结合引理
\ref{more-algebra-lemma-nr-branches-completion}、
\ref{more-algebra-lemma-equal-nr-branches-completion} 与
\ref{more-algebra-lemma-one-dimensional-formal-branch}，并使用
$A^\wedge$ 的维数为一这一事实；见引理
\ref{more-algebra-lemma-completion-dimension}。为说明几何分支数的结论，
使用刚证明的分支数结论、严格 Hensel 化不改变维数这一事实
（引理 \ref{more-algebra-lemma-henselization-dimension}），以及由引理
\ref{more-algebra-lemma-henselization-noetherian} 得到的
$(A^{sh})^\wedge = ((A^\wedge)^{sh})^\wedge$。
\end{proof}
\begin{lemma}
\label{more-algebra-lemma-geometrically-normal-formal-fibres-number-of-branches}
\begin{reference}
\cite[Theorem 2.3]{Beddani}
\end{reference}
设 $(A, \mathfrak m)$ 为 Noether 局部环。若 $A$ 的形式纤维
几何正规（例如 $A$ 优秀或拟优秀），则 $A$ 是 Nagata 环，
并且 $A$ 与 $A^\wedge$ 的（几何）分支数相同。
\end{lemma}

\begin{proof}
由于正规环既约，由引理 \ref{more-algebra-lemma-Nagata-local-ring}，
$A$ 是 Nagata 环。在证明的其余部分，我们将使用引理
\ref{more-algebra-lemma-formal-fibres-normal}、命题
\ref{more-algebra-proposition-finite-type-over-P-ring} 与引理
\ref{more-algebra-lemma-check-P-ring-maximal-ideals}。这些结果告诉我们，
$A$ 是 P-环，其中 $P(k \to R) = $``$R$ is geometrically
normal over $k$''，而且任意（本质）有限型 $A$-代数也具有同样性质。

\medskip\noindent
令 $\mathfrak q \subset A$ 为极小素理想。则
$A^\wedge/\mathfrak q A^\wedge = (A/\mathfrak q)^\wedge$
且 $A^h/\mathfrak qA^h = (A/\mathfrak q)^h$
（《代数》引理 \ref{algebra-lemma-quotient-henselization}）。
因此，$A$ 的分支数是各环 $A/\mathfrak q$ 的分支数之和；
$A^\wedge$ 也同样。这样便归约到 $A$ 是整环的情形。

\medskip\noindent
假设 $A$ 是整环。令 $A'$ 为 $A$ 在 $A$ 的分式域 $K$ 中的整闭包。
由于 $A$ 是 Nagata 环，$A \to A'$ 有限。回忆 $A$ 的分支数等于
$A'$ 的极大理想 $\mathfrak m'$ 的数目（引理
\ref{more-algebra-lemma-branches}）。还回忆
% P02-E0568
$$
(A')^\wedge = A' \otimes_A A^\wedge =
\prod\nolimits_{\mathfrak m' \subset A'} (A'_{\mathfrak m'})^\wedge
$$
，见《代数》引理
\ref{algebra-lemma-completion-finite-extension}。因为
$A'_{\mathfrak m'}$ 是形式纤维几何正规的局部环，由引理
\ref{more-algebra-lemma-completion-normal-local-ring}，
$(A'_{\mathfrak m'})^\wedge$ 正规。因此
$A' \otimes_A A^\wedge$ 的极小素理想与上述分解中的因子按
$1$ 对 $1$ 的方式对应。由 $A \to A^\wedge$ 的平坦性，有
% P02-E0568
$$
A^\wedge \subset A' \otimes_A A^\wedge \subset K \otimes_A A^\wedge
$$
。由于左侧环与右侧环具有同一个极小素理想集合，中间的环也如此
（略去一个小细节），这就完成了证明。
% P02-E0569

\medskip\noindent
为说明几何分支数的结论，使用刚证明的分支数结论、
$A^{sh}$ 的形式纤维几何正规这一事实（引理
\ref{more-algebra-lemma-formal-fibres-normal} 与
\ref{more-algebra-lemma-henselization-P-ring}），以及由引理
\ref{more-algebra-lemma-henselization-noetherian} 得到的
$(A^{sh})^\wedge = ((A^\wedge)^{sh})^\wedge$。
\end{proof}
\section{形式链环}
\label{more-algebra-section-formally-catenary}

\noindent
本节证明 Ratliff 的定理 \cite{Ratliff}：Noether 局部环是普遍链环，
当且仅当它是形式链环。

\begin{definition}
\label{more-algebra-definition-formally-catenary}
若对每个极小素理想 $\mathfrak p \subset A$，
$A^\wedge/\mathfrak p A^\wedge$ 的谱都等维，则称 Noether
局部环 $A$ 为{\it 形式链环}。
\end{definition}

\noindent
设 $A$ 为形式链 Noether 局部环。由 Ratliff 的结果（命题
\ref{more-algebra-proposition-ratliff}），$A$ 的任意商也形式链
（因为普遍链环这一类对取商稳定）。于是，对 $A$ 的每个素理想
$\mathfrak p$，$A^\wedge/\mathfrak p A^\wedge$ 的谱都等维。

\begin{lemma}
\label{more-algebra-lemma-not-formally-catenary}
设 $(A, \mathfrak m)$ 为非形式链的 Noether 局部环。则 $A$
不是普遍链环。
\end{lemma}

\begin{proof}
由假设，存在极小素理想 $\mathfrak p \subset A$，使
$A^\wedge /\mathfrak p A^\wedge$ 的谱不等维。用
$A/\mathfrak p$ 替换 $A$ 后，可以假设 $A$ 是整环，且
$A^\wedge$ 的谱不等维。令 $\mathfrak q$ 为 $A^\wedge$
的一个极小素理想，使 $d = \dim(A^\wedge/\mathfrak q)$ 最小；
于是 $0 < d < \dim(A)$。我们对 $d$ 作归纳以证明本引理。

\medskip\noindent
先考虑 $d = 1$ 的情形。此时
$\dim(A^\wedge_\mathfrak q) = 0$。故
$A^\wedge_\mathfrak q$ 是 Artin 局部环，从而对某个 $n > 0$，
理想 $J = \mathfrak q^n$ 在 $A^\wedge_\mathfrak q$ 中的像为零。
由此，$\mathfrak m$ 是 $J/J^2$ 唯一的伴随素理想，所以对某个
$m > 0$，$\mathfrak m^m$ 零化 $J/J^2$。因此，可以利用引理
\ref{more-algebra-lemma-quotient-by-idempotent} 找到有限型的 $A \to B$，
使得 $B^\wedge \cong A^\wedge/J$。
% P02-E0570
于是 $\mathfrak m_B = \sqrt{\mathfrak mB}$ 是极大理想，
其剩余域与 $\mathfrak m$ 的剩余域相同，且 $B^\wedge$
是 $\mathfrak m_B$-进完备化（《代数》引理
\ref{algebra-lemma-finite-after-completion}）。于是
$$
\dim(B_{\mathfrak m_B}) = \dim(B^\wedge) = 1 = d.
$$
由于有分解 $A \to B \to A^\wedge/J$，$\mathfrak q/J$
的逆像是素理想 $\mathfrak q' \subset \mathfrak m_B$，它位于
$A$ 的 $(0)$ 之上。因此，若 $A$ 是普遍链环，则维数公式
（《代数》引理 \ref{algebra-lemma-dimension-formula}）将给出
\begin{align*}
\dim(B_{\mathfrak m_B})
& \geq
\dim((B/\mathfrak q')_{\mathfrak m_B}) \\
& =
\dim(A) + \text{trdeg}_A(B/\mathfrak q') -
\text{trdeg}_{\kappa(\mathfrak m)}(\kappa(\mathfrak m_B)) \\
& =
\dim(A) + \text{trdeg}_A(B/\mathfrak q')
\end{align*}
。这个矛盾完成了 $d = 1$ 情形的论证。
% P02-E0571

\medskip\noindent
假设 $d > 1$。令 $Z \subset \Spec(A^\wedge)$ 为除
$V(\mathfrak q)$ 外其余不可约分支之并。令
$\mathfrak r_1, \ldots, \mathfrak r_m \subset A^\wedge$
为对应于 $V(\mathfrak q) \cap Z$ 中维数 $> 0$ 的不可约分支的
素理想。利用素理想回避（《代数》引理
\ref{algebra-lemma-silly}），选取
$f \in \mathfrak m$、$f \not \in A \cap \mathfrak r_j$。
则 $\dim(A/fA) = \dim(A) - 1$，且 $V(\mathfrak q, f)$
有某个维数为 $d - 1$ 的不可约分支。因此 $A/fA$ 不是形式链环，
且不变量 $d$ 已减小。由归纳假设，$A/fA$ 不是普遍链环，
故 $A$ 也不是普遍链环。
\end{proof}

\begin{lemma}
\label{more-algebra-lemma-flat-under-catenary-equidimensional}
设 $A \to B$ 为局部 Noether 环之间的平坦局部环映射。假设
$B$ 是链环，且 $\Spec(B)$ 等维。则：
% P02-E0572
\begin{enumerate}
\item 对所有 $\mathfrak p \subset A$，
$\Spec(B/\mathfrak p B)$ 都等维；
\item $A$ 是链环且 $\Spec(A)$ 等维。
\end{enumerate}
\end{lemma}

\begin{proof}
令 $\mathfrak p \subset A$ 为素理想，令 $\mathfrak q \subset B$
为 $\mathfrak pB$ 上的极小素理想。由 $A \to B$ 的下降定理
（《代数》引理 \ref{algebra-lemma-flat-going-down}），
$\mathfrak q \cap A = \mathfrak p$。故
$A_\mathfrak p \to B_\mathfrak q$ 是特殊纤维维数为 $0$
的平坦局部环映射，因此
$$
\dim(A_\mathfrak p) = \dim(B_\mathfrak q) = \dim(B) - \dim(B/\mathfrak q)
$$
（《代数》引理
\ref{algebra-lemma-dimension-base-fibre-equals-total}）。
第二个等式成立，是因为 $\Spec(B)$ 等维且 $B$ 是链环。
% P02-E0573
所以 $\dim(B/\mathfrak q)$ 与 $\mathfrak q$ 的选择无关，从而
$\Spec(B/\mathfrak p B)$ 等维，其维数为
$\dim(B) - \dim(A_\mathfrak p)$。另一方面，由平坦性
（见刚引用的引理），有
$\dim(B/\mathfrak p B) = \dim(A/\mathfrak p) + \dim(B/\mathfrak m_A B)$
以及
$\dim(B) = \dim(A) + \dim(B/\mathfrak m_A B)$，故得到
% P02-E0574
$$
\dim(A_\mathfrak p) = \dim(A) - \dim(A/\mathfrak p)
$$
，这对 $A$ 中所有 $\mathfrak p$ 都成立。把它应用于 $A$
的所有极小素理想，可见 $A$ 等维。若
$\mathfrak p \subset \mathfrak p'$ 是中间没有素理想的严格包含，
则可把上述结论应用于 $A/\mathfrak p$ 中的素理想
$\mathfrak p'/\mathfrak p$，因为
$A/\mathfrak p \to B/\mathfrak p B$ 平坦且
$\Spec(B/\mathfrak p B)$ 等维，从而得到
% P02-E0575
$$
1 = \dim((A/\mathfrak p)_{\mathfrak p'}) =
\dim(A/\mathfrak p) - \dim(A/\mathfrak p')
$$
。所以 $\mathfrak p \mapsto \dim(A/\mathfrak p)$ 是维数函数，
由此 $A$ 是链环。
\end{proof}

\begin{lemma}
\label{more-algebra-lemma-formally-catenary}
设 $A$ 为形式链 Noether 局部环。则 $A$ 是普遍链环。
\end{lemma}

\begin{proof}
可以用 $A/\mathfrak p$ 替换 $A$，其中 $\mathfrak p$ 是 $A$
的极小素理想；见《代数》引理
\ref{algebra-lemma-catenary-check-irreducible}。因此，可以假设
$A^\wedge$ 的谱等维。只须证明 $A$ 上本质有限型的每个局部环都是
链环（例如见《代数》引理
\ref{algebra-lemma-catenary-check-local}）。所以只须证明
$A[x_1, \ldots, x_n]_\mathfrak m$ 是链环，其中
$\mathfrak m \subset A[x_1, \ldots, x_n]$ 是位于
$\mathfrak m_A$ 之上的极大理想；见《代数》引理
\ref{algebra-lemma-localization-at-closed-point-special-fibre}
（以及《代数》引理 \ref{algebra-lemma-quotient-catenary} 与
\ref{algebra-lemma-localization-catenary}）。令
$\mathfrak m' \subset A^\wedge[x_1, \ldots, x_n]$ 为位于
$\mathfrak m$ 之上的唯一极大理想。则
% P02-E0576
$$
A[x_1, \ldots, x_n]_\mathfrak m \to A^\wedge[x_1, \ldots, x_n]_{\mathfrak m'}
$$
是局部平坦映射（《代数》引理
\ref{algebra-lemma-completion-flat}）。因此，由引理
\ref{more-algebra-lemma-flat-under-catenary-equidimensional}，只须证明右侧的环
是谱等维的链环。它是链环，因为完备局部环都是普遍链环
（《代数》注记
\ref{algebra-remark-Noetherian-complete-local-ring-universally-catenary}）。
任取 $A^\wedge[x_1, \ldots, x_n]_{\mathfrak m'}$ 的极小素理想
$\mathfrak q$。则对 $A^\wedge$ 的某个极小素理想 $\mathfrak p$，
有
$\mathfrak q = \mathfrak p A^\wedge[x_1, \ldots, x_n]_{\mathfrak m'}$
（略去一个小细节）。所以
% P02-E0577
$$
\dim(A^\wedge[x_1, \ldots, x_n]_{\mathfrak m'}/\mathfrak q) =
\dim(A^\wedge/\mathfrak p) + n = \dim(A^\wedge) + n
$$
；第一个等式由《代数》引理
\ref{algebra-lemma-dimension-base-fibre-equals-total} 得到，
第二个等式成立是因为 $A^\wedge$ 的谱等维。这就完成了证明。
\end{proof}

\begin{proposition}[Ratliff]
\label{more-algebra-proposition-ratliff}
\begin{reference}
\cite{Ratliff}
\end{reference}
Noether 局部环是普遍链环，当且仅当它是形式链环。
\end{proposition}

\begin{proof}
结合引理 \ref{more-algebra-lemma-not-formally-catenary} 与
\ref{more-algebra-lemma-formally-catenary}。
\end{proof}

\begin{lemma}
\label{more-algebra-lemma-geometrically-normal-fibres-universally-catenary}
\begin{reference}
\cite[Corollary 2.3]{Heinzer-Rotthaus-Wiegand}
\end{reference}
设 $(A, \mathfrak m)$ 为形式纤维几何正规的 Noether 局部环。则：
\begin{enumerate}
\item $A^h$ 是普遍链环；
\item 若 $A$ 单支（例如为正规环时），则 $A$ 是普遍链环。
\end{enumerate}
\end{lemma}

\begin{proof}
由引理
\ref{more-algebra-lemma-geometrically-normal-formal-fibres-number-of-branches}，
$A$ 与 $A^\wedge$ 的分支数相同，故引理
\ref{more-algebra-lemma-equal-nr-branches-completion} 适用。于是，对任意极小素理想
$\mathfrak q \subset A^h$，$A^\wedge/\mathfrak q A^\wedge$
有唯一的极小素理想。因此，$A^h$ 依定义是形式链环，进而由命题
\ref{more-algebra-proposition-ratliff} 是普遍链环。若 $A$ 单支，则 $A^h$
有唯一的极小素理想，故 $A^\wedge$ 有唯一的极小素理想，
于是 $A$ 是形式链环；同理可得结论。
\end{proof}
\section{群作用与整闭包}
\label{more-algebra-section-group-actions-integral}

\noindent
在某种意义下，本节是《代数》第
\ref{algebra-section-going-down-integral-over-normal} 节的延续。
在《基本群》第 \ref{pione-section-group-actions-integral} 节中还可找到
更多类似材料。
% P02-E0578

\begin{lemma}
\label{more-algebra-lemma-pol-lifting}
设 $\varphi : A \to B$ 为环满射，$G$ 为作用在
$\varphi : A \to B$ 上的 $n$ 阶有限群。若 $b \in B^G$，
则存在首一多项式 $P \in A^G[T]$，它在 $B^G[T]$ 中的像为
$(T - b)^n$。
\end{lemma}

\begin{proof}
选取提升 $b$ 的 $a \in A$，并置
$P = \prod_{\sigma \in G} (T - \sigma(a))$。
\end{proof}

\begin{lemma}
\label{more-algebra-lemma-invariants-modulo}
设 $R$ 为环，$G$ 为作用在 $R$ 上的有限群。设 $I \subset R$
为理想，使得对所有 $\sigma \in G$ 都有 $\sigma(I) \subset I$。
则 $R^G/I^G \subset (R/I)^G$ 是整环扩张；它在谱上诱导同胚，
并诱导剩余域的纯不可分扩张。
\end{lemma}

\begin{proof}
由于 $I^G = R^G \cap I$，该映射显然单射。引理
\ref{more-algebra-lemma-pol-lifting} 表明《代数》引理
\ref{algebra-lemma-universally-bijective} 适用。
\end{proof}

\begin{lemma}
\label{more-algebra-lemma-pol-in-ideal}
设 $G$ 为作用在环 $R$ 上的 $n$ 阶有限群，$J \subset R^G$
为理想。对 $x \in JR$，有
$\prod_{\sigma \in G} (T - \sigma(x)) =
T^n + a_1 T^{n - 1} + \ldots + a_n$，其中 $a_i \in J$。
\end{lemma}

\begin{proof}
注意该多项式确实首一，且系数属于 $R^G$。可以写
$x = f_1 b_1 + \ldots + f_m b_m$，其中 $f_j \in J$、
$b_j \in R$。因此，对 $m$ 作归纳，可以假设
$x = y - fb$，其中 $f \in J$、$b \in R$ 且 $y \in JR$，
并且结论对 $y$ 成立。于是
$$
\prod\nolimits_{\sigma \in G} (T - \sigma(x)) =
\prod\nolimits_{\sigma \in G} (T - \sigma(y) + f\sigma(b)) =
\prod\nolimits_{\sigma \in G} (T - \sigma(y)) +
\sum_{i = 1, \ldots, n} f^i a_i
$$
，其中
% P02-E0579
$$
a_i =
\sum\nolimits_{S \subset G,\ |S| = i}
\prod\nolimits_{\sigma \in S} \sigma(b)
\prod\nolimits_{\sigma \not \in S} (T - \sigma(y))
$$
。我们略去的一个计算表明 $a_i \in R^G$（提示：所给表达式是
对称的）。因此，引理陈述中关于 $x$ 的多项式模 $J$
同余于关于 $y$ 的多项式，这就完成了归纳步骤。
\end{proof}

\begin{lemma}
\label{more-algebra-lemma-invariants-modulo-bis}
设 $R$ 为环，$G$ 为作用在 $R$ 上的 $n$ 阶有限群，
$J \subset R^G$ 为理想。则 $R^G/J \to (R/JR)^G$ 是环映射，
并且：
% P02-E0580
\begin{enumerate}
\item 对 $b \in (R/JR)^G$，存在首一多项式
$P \in R^G/J[T]$，它在 $(R/JR)^G[T]$ 中的像为 $(T - b)^n$；
\item 对 $a \in \Ker(R^G/J \to (R/JR)^G)$，在
$R^G/J[T]$ 中有 $(T - a)^n = T^n$。
\end{enumerate}
特别地，$R^G/J \to (R/JR)^G$ 是整环映射；它在谱上诱导同胚，
并诱导剩余域的纯不可分扩张。
\end{lemma}

\begin{proof}
(1) 由取 $I = JR$ 的引理 \ref{more-algebra-lemma-pol-lifting} 得出。
% P02-E0581
若 $a$ 如 (2) 中所述，则 $a$ 是某个
$x \in R^G \cap JR$ 的像。因此
$(T - x)^n = \prod_{\sigma \in G} (T - \sigma(x))$
由引理 \ref{more-algebra-lemma-pol-in-ideal} 模 $J$ 同余于 $T^n$。
这就证明了 (2)。为得到最后的陈述，可以应用《代数》引理
\ref{algebra-lemma-universally-bijective}。
\end{proof}

\begin{remark}
\label{more-algebra-remark-when-iso}
在引理 \ref{more-algebra-lemma-invariants-modulo-bis} 中可见，若 $n$ 在 $R$
中可逆，则映射 $R^G/J \to (R/JR)^G$ 是同构。
\end{remark}

\begin{lemma}
\label{more-algebra-lemma-functor-invariants-tensor}
设 $R$ 为环，$G$ 为作用在 $R$ 上的 $n$ 阶有限群，
$A$ 为 $R^G$-代数。
\begin{enumerate}
\item 对 $b \in (A \otimes_{R^G} R)^G$，存在首一多项式
$P \in A[T]$，它在 $(A \otimes_{R^G} R)^G[T]$ 中的像为
$(T - b)^n$；
\item 对 $a \in \Ker(A \to (A \otimes_{R^G} R)^G)$，在
$A[T]$ 中有 $(T - a)^n = T^n$。
\end{enumerate}
\end{lemma}

\begin{proof}
选取满射 $E \to A$，其中 $E$ 是 $R^G$ 上的多项式代数。由于
$E$ 是自由 $R^G$-模，有
$(E \otimes_{R^G} R)^G = E$。以
$J = \Ker(E \to A)$ 表示其核。
% P02-E0582
由于张量积右正合，$A \otimes_{R^G} R$ 是
$E \otimes_{R^G} R$ 关于 $J$ 生成的理想的商。这样可见，
本引理是引理 \ref{more-algebra-lemma-invariants-modulo-bis} 的特殊情形。
\end{proof}

\begin{lemma}
\label{more-algebra-lemma-base-change-invariants}
设 $R$ 为环，$G$ 为作用在 $R$ 上的有限群，$R^G \to A$
为环映射。映射
$$
A \to (A \otimes_{R^G} R)^G
$$
在 $R^G \to A$ 平坦时是同构。一般而言，该映射是整的，在谱上
诱导同胚，并诱导纯不可分的剩余域扩张。
\end{lemma}

\begin{proof}
为证明第一项陈述，考虑正合列
$0 \to R^G \to R \to \bigoplus_{\sigma \in G} R$，其中第二个映射
把 $x$ 送到 $(\sigma(x) - x)_{\sigma \in G}$。若
$R^G \to A$ 平坦，与 $A$ 张量后该序列仍正合。因此，$A$
是 $(A \otimes_{R^G} R)^G$ 中的 $G$-不变量。
% P02-E0583

\medskip\noindent
第二项陈述由引理 \ref{more-algebra-lemma-functor-invariants-tensor} 与《代数》引理
\ref{algebra-lemma-universally-bijective} 得出。
\end{proof}

\begin{lemma}
\label{more-algebra-lemma-one-orbit}
设 $G$ 为作用在环 $R$ 上的有限群。对位于 $R^G$ 中同一个素理想
之上的任意两个素理想 $\mathfrak q, \mathfrak q' \subset R$，
存在 $\sigma \in G$，使得
$\sigma(\mathfrak q) = \mathfrak q'$。
\end{lemma}

\begin{proof}
扩张 $R^G \subset R$ 是整的，因为每个 $x \in R$ 都是
$R^G[T]$ 中首一多项式
$\prod_{\sigma \in G}(T - \sigma(x))$ 的根。因此，位于给定素理想
$\mathfrak p$ 之上的素理想之间不存在包含关系（《代数》引理
\ref{algebra-lemma-integral-no-inclusion}）。若本引理不成立，
则可选取 $x \in \mathfrak q'$，使得对所有 $\sigma \in G$
都有 $x \not \in \sigma(\mathfrak q)$；见《代数》引理
\ref{algebra-lemma-silly}。于是
$y = \prod_{\sigma \in G} \sigma(x)$ 属于 $R^G$，并属于
$\mathfrak p = R^G \cap \mathfrak q'$。另一方面，对所有
$\sigma$ 都有 $x \not \in \sigma(\mathfrak q)$，意味着对所有
$\sigma$ 都有 $\sigma(x) \not \in \mathfrak q$。由于
$\mathfrak q$ 是素理想，$y \not \in \mathfrak q$。
这不可能，因为 $y \in \mathfrak p \subset \mathfrak q$。
\end{proof}

\begin{lemma}
\label{more-algebra-lemma-one-orbit-geometric}
设 $G$ 为作用在环 $R$ 上的有限群。设素理想
$\mathfrak q \subset R$ 位于 $\mathfrak p \subset R^G$ 之上。
则 $\kappa(\mathfrak q)/\kappa(\mathfrak p)$ 是代数正规扩张，
并且映射
$$
D = \{\sigma \in G \mid \sigma(\mathfrak q) = \mathfrak q\}
\longrightarrow
\text{Aut}(\kappa(\mathfrak q)/\kappa(\mathfrak p))
$$
满射\footnote{回忆我们只在 Galois 扩张的情形使用记号
$\text{Gal}$。}。
\end{lemma}

\begin{proof}
取 $A = (R^G)_\mathfrak p$ 与 $B = A \otimes_{R^G} R$。由于局部化
平坦，由引理 \ref{more-algebra-lemma-base-change-invariants}，
$A = B^G$。注意 $\mathfrak pA$ 与 $\mathfrak qB$ 是素理想，
$D$ 是 $\mathfrak qB$ 的稳定子，并且
$\kappa(\mathfrak p) = \kappa(\mathfrak pA)$、
$\kappa(\mathfrak q) = \kappa(\mathfrak qB)$。因此可以用 $B$
替换 $R$，并假设 $\mathfrak p$ 是极大理想。由于
$R^G \subset R$ 是整环扩张，$R$ 的极大理想恰为位于
$\mathfrak p$ 之上的素理想（这由《代数》引理
\ref{algebra-lemma-integral-no-inclusion} 与
\ref{algebra-lemma-integral-going-up} 得出）。由引理
\ref{more-algebra-lemma-one-orbit}，它们只有有限多个，记作
$\mathfrak q = \mathfrak q_1, \mathfrak q_2, \ldots, \mathfrak q_m$，
且构成 $G$ 的一个轨道。由中国剩余定理（《代数》引理
\ref{algebra-lemma-chinese-remainder}），映射
$R \to \prod_{j = 1, \ldots, m} R/\mathfrak q_j$ 满射。

\medskip\noindent
首先证明该扩张正规。取元素
$\alpha \in \kappa(\mathfrak q)$。我们须证明 $\alpha$
在 $\kappa(\mathfrak p)$ 上的极小多项式 $P$ 完全分裂。
由上述结论，可选取
$a \in \mathfrak q_2 \cap \ldots \cap \mathfrak q_m$，使它在
$\kappa(\mathfrak q)$ 中的像为 $\alpha$。考虑 $R^G[T]$
中的多项式 $Q = \prod_{\sigma \in G} (T - \sigma(a))$。
$Q$ 在 $R[T]$ 中的像完全分裂为一次因子，故其在
$\kappa(\mathfrak q)[T]$ 中的像也如此。由于 $P$ 整除 $Q$
在 $\kappa(\mathfrak p)[T]$ 中的像，可知 $P$ 在
$\kappa(\mathfrak q)$ 上完全分裂为一次因子，如所求。

\medskip\noindent
由于 $\kappa(\mathfrak q)/\kappa(\mathfrak p)$ 正规，可以假设
$\kappa(\mathfrak q) = \kappa_1 \otimes_{\kappa(\mathfrak p)} \kappa_2$，
其中 $\kappa_1/\kappa(\mathfrak p)$ 纯不可分，而
$\kappa_2/\kappa(\mathfrak p)$ 为 Galois 扩张；见《域》引理
\ref{fields-lemma-normal-case}。选取 $\alpha \in \kappa_2$：
若该扩张有限，使其在 $\kappa(\mathfrak p)$ 上生成
$\kappa_2$；若该扩张无限，使它生成次数 $> |G|$ 的子域
（以得到矛盾）。
% P02-E0584
由《域》引理 \ref{fields-lemma-primitive-element}，这是可能的。
如前段一样选取 $a$、$P$ 与 $Q$。若
$\alpha' \in \kappa_2$ 是 $\alpha$ 在 $\kappa(\mathfrak p)$ 上的
Galois 共轭元，则
% P02-E0585
$P$ 整除 $P$ 在 $\kappa(\mathfrak p)[T]$ 中的像这一事实表明，
存在 $\sigma \in G$，使得 $\sigma(a)$ 映到 $\alpha'$。
由 $a$ 的选取方式（它在其他极大理想处消失），这蕴含
$\sigma \in D$，且 $\sigma$ 在
$\text{Aut}(\kappa(\mathfrak q)/\kappa(\mathfrak p))$
中的像把 $\alpha$ 映到 $\alpha'$。因此得到所需的满射；
若 $\alpha$ 在 $\kappa(\mathfrak p)$ 上的次数 $> |G|$，
则得到所需的荒谬结论。
% P02-E0586
\end{proof}

\begin{lemma}
\label{more-algebra-lemma-one-orbit-geometric-galois}
设 $A$ 为分式域是 $K$ 的正规整环，$L/K$ 为（可能无限的）
Galois 扩张。令 $G = \text{Gal}(L/K)$，并令 $B$ 为 $A$
在 $L$ 中的整闭包。
\begin{enumerate}
\item 对位于 $A$ 中同一个素理想之上的任意两个素理想
$\mathfrak q, \mathfrak q' \subset B$，存在
$\sigma \in G$，使得 $\sigma(\mathfrak q) = \mathfrak q'$。
\item 设素理想 $\mathfrak q \subset B$ 位于
$\mathfrak p \subset A$ 之上。则
$\kappa(\mathfrak q)/\kappa(\mathfrak p)$ 是代数正规扩张，
并且映射
$$
D = \{\sigma \in G \mid \sigma(\mathfrak q) = \mathfrak q\}
\longrightarrow
\text{Aut}(\kappa(\mathfrak q)/\kappa(\mathfrak p))
$$
满射。
\end{enumerate}
\end{lemma}

\begin{proof}
证明 (1)。考虑序对 $(M, \sigma)$，其中
$K \subset M \subset L$ 是子域，$M/K$ 为 Galois 扩张，
$\sigma \in \text{Gal}(M/K)$，并且
$\sigma(\mathfrak q \cap M) = \mathfrak q' \cap M$。
定义 $(M', \sigma') \geq (M, \sigma)$，当且仅当
$M \subset M'$ 且 $\sigma'|_M = \sigma$。注意
$(K, \text{id}_K)$ 是这样的序对，因为 $A$ 是正规整环，所以
$A = K \cap B$。这些序对构成的集合满足 Zorn 引理的假设，
故存在极大序对 $(M, \sigma)$。若 $M \not = L$，则可以找到
$M \subset M' \subset L$，使得 $M'/M$ 非平凡且有限，并且
$M'/K$ 为 Galois 扩张（《域》引理
\ref{fields-lemma-normal-closure-inside-normal}）。选取
$\sigma' \in \text{Gal}(M'/K)$，使它在 $M$ 上的限制为
$\sigma$（《域》引理 \ref{fields-lemma-galois-infinite}）。
于是素理想 $\sigma'(\mathfrak q \cap M')$ 与
$\mathfrak q' \cap M'$ 在 $B \cap M$ 中的收缩相同。由于
$B \cap M = (B \cap M')^{\text{Gal}(M'/M)}$，可以应用引理
\ref{more-algebra-lemma-one-orbit}，找到 $\tau \in \text{Gal}(M'/M)$，使得
$\tau(\sigma'(\mathfrak q \cap M')) = \mathfrak q' \cap M'$。
于是 $(M', \tau \circ \sigma') > (M, \sigma)$，这与
$(M, \sigma)$ 的极大性矛盾。

\medskip\noindent
(2) 的证明方式与 (1) 完全相同。下面写出细节。选取
$\overline{\sigma} \in
\text{Aut}(\kappa(\mathfrak q)/\kappa(\mathfrak p))$。
考虑序对 $(M, \sigma)$，其中 $K \subset M \subset L$ 是子域，
$M/K$ 为 Galois 扩张，$\sigma \in \text{Gal}(M/K)$，
$\sigma(\mathfrak q \cap M) = \mathfrak q \cap M$，并且图
$$
\xymatrix{
\kappa(\mathfrak q \cap M) \ar[r] \ar[d]_\sigma &
\kappa(\mathfrak q) \ar[d]_{\overline{\sigma}} \\
\kappa(\mathfrak q \cap M) \ar[r] & \kappa(\mathfrak q)
}
$$
交换。定义 $(M', \sigma') \geq (M, \sigma)$，当且仅当
$M \subset M'$ 且 $\sigma'|_M = \sigma$。与上面一样，
$(K, \text{id}_K)$ 是这样的序对。这些序对构成的集合满足
Zorn 引理的假设，故存在极大序对 $(M, \sigma)$。若
$M \not = L$，则可以找到 $M \subset M' \subset L$，使得
$M'/M$ 有限，并且 $M'/K$ 为 Galois 扩张（《域》引理
\ref{fields-lemma-normal-closure-inside-normal}）。选取
$\sigma' \in \text{Gal}(M'/K)$，使它在 $M$ 上的限制为
$\sigma$（《域》引理 \ref{fields-lemma-galois-infinite}）。
于是素理想 $\sigma'(\mathfrak q \cap M')$ 与
$\mathfrak q \cap M'$ 在 $B \cap M$ 中的收缩相同。像第一段那样
调整 $\sigma'$ 的选择后，可以假设
$\sigma'(\mathfrak q \cap M') = \mathfrak q \cap M'$。
于是 $\sigma'$ 与 $\overline{\sigma}$ 定义映射
$\kappa(\mathfrak q \cap M') \to \kappa(\mathfrak q)$，
且它们在 $\kappa(\mathfrak q \cap M)$ 上相同。由于
$B \cap M = (B \cap M')^{\text{Gal}(M'/M)}$，可以应用引理
\ref{more-algebra-lemma-one-orbit-geometric}，找到
$\tau \in \text{Gal}(M'/M)$，使得
$\tau(\mathfrak q \cap M') = \mathfrak q \cap M'$，并使
% P02-E0587
$\tau \circ \sigma$ 与 $\overline{\sigma}$ 在
$\kappa(\mathfrak q \cap M')$ 上诱导同一个映射。这里有一个小细节：
该引理首先保证
$\kappa(\mathfrak q \cap M')/\kappa(\mathfrak q \cap M)$ 正规，
由此可知两个映射之差是该扩张的自同构（《域》引理
\ref{fields-lemma-normal-embeddings-differ-by-aut}），再对它应用该引理
即可得到 $\tau$。因此
$(M', \tau \circ \sigma') > (M, \sigma)$，这与
$(M, \sigma)$ 的极大性矛盾。
\end{proof}

\begin{lemma}
\label{more-algebra-lemma-one-orbit-geometric-galois-compare}
设 $A$ 为分式域是 $K$ 的正规整环，
$M/L/K$ 为 $K$ 的（可能无限的）Galois 扩张塔。令
$H = \text{Gal}(M/K)$、$G = \text{Gal}(L/K)$，并分别令
$C$ 与 $B$ 为 $A$ 在 $M$ 与 $L$ 中的整闭包。
% P02-E0588
设 $\mathfrak r \subset C$，并令
$\mathfrak q = B \cap \mathfrak r$。置
$D_\mathfrak r = \{\tau \in H \mid \tau(\mathfrak r) = \mathfrak r\}$
以及
$I_\mathfrak r = \{\tau \in D_\mathfrak r \mid
\tau \bmod \mathfrak r = \text{id}_{\kappa(\mathfrak r)}\}$，
并类似地定义 $D_\mathfrak q$ 与 $I_\mathfrak q$。
在映射 $H \to G$ 下，所诱导的映射
$D_\mathfrak r \to D_\mathfrak q$ 与
$I_\mathfrak r \to I_\mathfrak q$ 都满射。
\end{lemma}

\begin{proof}
令 $\sigma \in D_\mathfrak q$。选取映到 $\sigma$ 的
$\tau \in H$；由《域》引理
\ref{fields-lemma-galois-infinite}，这是可能的。于是
$\tau(\mathfrak r)$ 与 $\mathfrak r$ 都位于 $\mathfrak q$ 之上。
故由引理 \ref{more-algebra-lemma-one-orbit-geometric-galois}，存在
$\sigma' \in \text{Gal}(M/L)$，使得
$\sigma'(\tau(\mathfrak r)) = \mathfrak r$。于是
$\sigma'\tau \in D_\mathfrak r$ 映到 $\sigma$。
惰性群的情形完全相同，只须利用到自同构群的满射性。
\end{proof}
\section{离散赋值环的扩张}
\label{more-algebra-section-discrete-valuation-rings}

\noindent
在本节及随后几节中，我们使用以下定义。

\begin{definition}
\label{more-algebra-definition-extension-discrete-valuation-rings}
若 $A$ 与 $B$ 都是离散赋值环，并且 $A \to B$ 单射且为局部映射，
则称 $A \to B$ 或 $A \subset B$ 为{\it 离散赋值环的扩张}。
特别地，若 $\pi_A$ 与 $\pi_B$ 分别为 $A$ 与 $B$ 的一致化元，
则对某个 $e \geq 1$ 以及 $B$ 的某个单位 $u$，有
$\pi_A = u \pi_B^e$。整数 $e$ 与一致化元的选择无关，因为它也是
使下式成立的唯一整数 $\geq 1$：
% P02-E0589
$$
\mathfrak m_A B = \mathfrak m_B^e
$$
。整数 $e$ 称为 $B$ 在 $A$ 上的{\it 分歧指数}。若 $e = 1$，
则称 $B$ 在 $A$ 上{\it 弱非分歧}。若剩余域扩张
$\kappa_A = A/\mathfrak m_A \subset \kappa_B = B/\mathfrak m_B$
有限，则置 $f = [\kappa_B : \kappa_A]$，并称它为扩张
$A \subset B$ 的{\it 剩余次数}。
\end{definition}

\noindent
注意，我们不要求分式域扩张有限。

\begin{lemma}
\label{more-algebra-lemma-inequality}
设 $A \subset B$ 为离散赋值环的扩张，其分式域为
$K \subset L$。若扩张 $L/K$ 有限，则剩余域扩张有限，并且
$ef \leq [L : K]$。
\end{lemma}

\begin{proof}
% P02-E0595
剩余域扩张的有限性见《代数》引理
\ref{algebra-lemma-finite-extension-residue-fields-dimension-1}。
不等式由《代数》引理 \ref{algebra-lemma-finite-length} 与
\ref{algebra-lemma-pushdown-module} 得出。
\end{proof}

\begin{lemma}
\label{more-algebra-lemma-multiplicative-e-f}
设 $A \subset B \subset C$ 为离散赋值环的扩张。则 $B/A$
与 $C/B$ 的分歧指数之积等于 $C/A$ 的分歧指数，即
% P02-E0596
$e_{C/A} = e_{B/A} e_{C/B}$。若各剩余次数有限，它们也满足同样的
乘法性。
\end{lemma}

\begin{proof}
这由定义及《域》引理
\ref{fields-lemma-multiplicativity-degrees} 立即得出。
\end{proof}

\begin{lemma}
\label{more-algebra-lemma-ramification-index-a-power-of-p}
设 $A \subset B$ 为离散赋值环的扩张，所诱导的域扩张为
$K \subset L$。若 $K$ 的特征为 $p > 0$，且 $L$ 在 $K$
上纯不可分，则分歧指数 $e$ 是 $p$ 的幂。
\end{lemma}

\begin{proof}
对某个 $u \in B^*$，写 $\pi_A = u \pi_B^e$。另一方面，
对某个 $p$ 的幂 $q$，有 $\pi_B^q \in K$。对某个
$v \in A^*$ 与 $k \in \mathbf{Z}$，写
$\pi_B^q = v \pi_A^k$。于是
$\pi_A^q = u^q \pi_B^{qe} = u^q v^e \pi_A^{ke}$。
在 $B$ 中取赋值，得到 $ke = q$。
\end{proof}

\noindent
在下一个引理中，我们讨论离散赋值环的扩张 $A \subset B$
为“非分歧”的含义，即分歧指数为 $1$，且剩余域扩张可分
（可能非代数）。不过，我们不能直接使用“非分歧”一词，因为已经有
非分歧环映射这一概念；见《代数》第
\ref{algebra-section-unramified} 节。

\begin{lemma}
\label{more-algebra-lemma-extension-dvrs-formally-smooth}
设 $A \subset B$ 为离散赋值环的扩张。以下条件等价：
\begin{enumerate}
\item $A \to B$ 关于 $\mathfrak m_B$-进拓扑形式光滑；
\item $A \to B$ 弱非分歧，且 $\kappa_B/\kappa_A$
是可分域扩张。
\end{enumerate}
\end{lemma}

\begin{proof}
这由命题 \ref{more-algebra-proposition-fs-flat-fibre-fs} 与《代数》命题
\ref{algebra-proposition-characterize-separable-field-extensions}
得出。
\end{proof}

\begin{remark}
\label{more-algebra-remark-finite-separable-extension}
设 $A$ 为分式域是 $K$ 的离散赋值环，$L/K$ 为有限可分域扩张，
$B \subset L$ 为 $A$ 在 $L$ 中的整闭包。图为：
% P02-E0590
$$
\xymatrix{
B \ar[r] & L \\
A \ar[u] \ar[r] & K \ar[u]
}
$$
。由《代数》引理
\ref{algebra-lemma-Noetherian-normal-domain-finite-separable-extension}，
环扩张 $A \subset B$ 有限，故 $B$ 是 Noether 环。由《代数》引理
\ref{algebra-lemma-integral-sub-dim-equal}，$B$ 的维数为 $1$，
故 $B$ 是 Dedekind 整环；见《代数》引理
\ref{algebra-lemma-characterize-Dedekind}。令
$\mathfrak m_1, \ldots, \mathfrak m_n$ 为 $B$ 的极大理想
（即位于 $\mathfrak m_A$ 之上的素理想）。我们得到离散赋值环的扩张
% P02-E0590
$$
A \subset B_{\mathfrak m_i}
$$
，从而得到分歧指数 $e_i$ 与剩余次数 $f_i$。由《代数》引理
\ref{algebra-lemma-finite-extension-dim-1} 对 $A$ 中的一致化元的应用，
有
% P02-E0590
$$
[L : K] = \sum\nolimits_{i = 1, \ldots, n} e_i f_i
$$
。注意，若 $A$ 是 Hensel 环，则 $n = 1$（由《代数》引理
\ref{algebra-lemma-finite-over-henselian}）；例如 $A$ 完备时即如此。
\end{remark}

\begin{definition}
\label{more-algebra-definition-types-of-extensions}
设 $A$ 为分式域是 $K$ 的离散赋值环，$L/K$ 为有限可分扩张。
沿用注记 \ref{more-algebra-remark-finite-separable-extension} 中的 $B$ 与
$\mathfrak m_i$（$i = 1, \ldots, n$），称扩张 $L/K$：
\begin{enumerate}
\item {\it 关于 $A$ 非分歧}，若对所有 $i$ 都有 $e_i = 1$，
且扩张 $\kappa(\mathfrak m_i)/\kappa_A$ 可分；
\item {\it 关于 $A$ 驯分歧}，若 $\kappa_A$ 的特征为 $0$，或者
$\kappa_A$ 的特征为 $p > 0$、各域扩张
$\kappa(\mathfrak m_i)/\kappa_A$ 可分且各分歧指数 $e_i$
与 $p$ 互素；
\item {\it 关于 $A$ 完全分歧}，若 $n = 1$，且剩余域扩张
$\kappa(\mathfrak m_1)/\kappa_A$ 平凡。
\end{enumerate}
若从上下文可明确离散赋值环 $A$，则有时简称 $L/K$ 非分歧、
完全分歧或驯分歧。
\end{definition}

\noindent
对非分歧扩张，有以下基本引理。

\begin{lemma}
\label{more-algebra-lemma-permanence-unramified}
设 $A$ 为分式域是 $K$ 的离散赋值环。
\begin{enumerate}
\item 若 $M/L/K$ 是有限可分扩张，且 $M$ 关于 $A$ 非分歧，
则 $L$ 关于 $A$ 非分歧。
\item 若 $L/K$ 是关于 $A$ 非分歧的有限可分扩张，则存在包含
$L$ 且关于 $A$ 非分歧的 Galois 扩张 $M/K$。
\item 若 $L_1/K$、$L_2/K$ 是关于 $A$ 非分歧的有限可分扩张，
则存在包含 $L_1$ 与 $L_2$ 且关于 $A$ 非分歧的有限可分扩张
$L/K$。
% P02-E0591
\end{enumerate}
\end{lemma}

\begin{proof}
以下不再特别说明地使用注记
\ref{more-algebra-remark-finite-separable-extension} 中讨论的结果。

\medskip\noindent
证明 (1)。令 $C/B/A$ 分别为 $A$ 在 $M/L/K$ 中的整闭包。
% P02-E0592
由于 $C$ 是 $B$ 的有限环扩张，
$\Spec(C) \to \Spec(B)$ 满射。因此，对每个极大理想
$\mathfrak m \subset B$，都有一个位于 $\mathfrak m$ 之上的极大理想
$\mathfrak m' \subset C$。由分歧指数的乘法性（引理
\ref{more-algebra-lemma-multiplicative-e-f}）与假设，$B_\mathfrak m$
在 $A$ 上的分歧指数为 $1$。由于
$\kappa(\mathfrak m')/\kappa_A$ 有限可分，
$\kappa(\mathfrak m)/\kappa_A$ 也如此。

\medskip\noindent
证明 (2)。令 $M$ 为 $L$ 在 $K$ 上的正规闭包；见《域》定义
\ref{fields-definition-normal-closure}。由《域》引理
\ref{fields-lemma-normal-closure-galois}，$M/K$ 是 Galois 扩张。
另一方面，存在 $K$-代数的满射
% P02-E0593
$$
L \otimes_K \ldots \otimes_K L \longrightarrow M
$$
；见《域》引理
\ref{fields-lemma-normal-closure-tensor-product}。令 $B$ 为
$A$ 在 $L$ 中的整闭包，如注记
\ref{more-algebra-remark-finite-separable-extension} 所述。$L$ 关于 $A$
非分歧这一条件恰好意味着 $A \to B$ 是 \'etale 环映射；
见《代数》引理 \ref{algebra-lemma-characterize-etale}。
由 \'etale 环映射的恒久性质，
$$
B \otimes_A \ldots \otimes_A B
$$
在 $A$ 上 \'etale；见《代数》引理
\ref{algebra-lemma-etale}。因此，所显示的环是 Dedekind 整环的
乘积；见引理 \ref{more-algebra-lemma-Dedekind-etale-extension}。由此，
$M$ 是某个在 $A$ 上有限 \'etale 的 Dedekind 整环的分式域。
这意味着 $M$ 如所求地关于 $A$ 非分歧。

\medskip\noindent
证明 (3)。令 $B_i \subset L_i$ 分别为 $A$ 在相应域中的整闭包。
% P02-E0594
采用与上面相同的论证，可知
$B_1 \otimes_A B_2$ 在 $A$ 上有限 \'etale。略去细节。
\end{proof}

\begin{lemma}
\label{more-algebra-lemma-composition-unramified}
设 $A$ 为分式域是 $K$ 的离散赋值环，$M/L/K$ 为有限可分扩张。
令 $B$ 为 $A$ 在 $L$ 中的整闭包。若 $L/K$ 关于 $A$ 非分歧，
并且对 $B$ 的每个极大理想 $\mathfrak m$，$M/L$
都关于 $B_\mathfrak m$ 非分歧，则 $M/K$ 关于 $A$ 非分歧。
\end{lemma}

\begin{proof}
令 $C$ 为 $A$ 在 $M$ 中的整闭包。$C$ 的每个极大理想
$\mathfrak m'$ 都位于 $B$ 的某个极大理想 $\mathfrak m$ 之上。
于是，本引理由分歧指数的乘法性（引理
\ref{more-algebra-lemma-multiplicative-e-f}）以及有限域扩张塔
$\kappa(\mathfrak m')/\kappa(\mathfrak m)/\kappa_A$ 得出。
\end{proof}
\section{Galois 扩张与分歧}
\label{more-algebra-section-ramification}

\noindent
在 Galois 扩张的情形，可以进一步展开第
\ref{more-algebra-section-discrete-valuation-rings} 节的讨论。

\begin{lemma}
\label{more-algebra-lemma-galois}
设 $A$ 为分式域是 $K$ 的离散赋值环，$L/K$ 为 Galois 群是
$G$ 的有限 Galois 扩张。则 $G$ 作用在注记
\ref{more-algebra-remark-finite-separable-extension} 的环 $B$ 上，并传递地作用在
$B$ 的极大理想集合上。
\end{lemma}

\begin{proof}
注意 $A = B^G$，因为 $A$ 在 $K$ 中整闭且 $K = L^G$。
故本引理是引理 \ref{more-algebra-lemma-one-orbit} 的特殊情形。
\end{proof}

\begin{lemma}
\label{more-algebra-lemma-galois-conclusion}
设 $A$ 为分式域是 $K$ 的离散赋值环，$L/K$ 为有限 Galois 扩张。
则存在 $e \geq 1$ 与 $f \geq 1$，使得对所有 $i$ 都有
$e_i = e$ 与 $f_i = f$（记号同注记
\ref{more-algebra-remark-finite-separable-extension}）。特别地，
$[L : K] = n e f$。
\end{lemma}

\begin{proof}
这是引理 \ref{more-algebra-lemma-galois} 与各定义的直接推论。
\end{proof}

\begin{definition}
\label{more-algebra-definition-decomposition-inertia}
设 $A$ 为分式域是 $K$ 的离散赋值环，$L/K$ 为 Galois 群是
$G$ 的有限 Galois 扩张。令 $B$ 为 $A$ 在 $L$ 中的整闭包，
$\mathfrak m \subset B$ 为极大理想。
\begin{enumerate}
\item $\mathfrak m$ 的{\it 分解群}是子群
$D = \{\sigma \in G \mid \sigma(\mathfrak m) = \mathfrak m\}$。
\item $\mathfrak m$ 的{\it 惰性群}是映射
$D \to \text{Aut}(\kappa(\mathfrak m)/\kappa_A)$ 的核 $I$。
\end{enumerate}
\end{definition}

\noindent
注意，域 $\kappa(\mathfrak m)$ 在 $\kappa_A$ 上可能不可分。
特别地，域扩张 $\kappa(\mathfrak m)/\kappa_A$ 未必是 Galois 扩张。
若 $\kappa_A$ 完美，则它是 Galois 扩张。

\begin{lemma}
\label{more-algebra-lemma-galois-galois}
设 $A$ 为分式域是 $K$、剩余域是 $\kappa$ 的离散赋值环，
$L/K$ 为 Galois 群是 $G$ 的有限 Galois 扩张。令 $B$ 为
$A$ 在 $L$ 中的整闭包，$\mathfrak m$ 为 $B$ 的极大理想。则：
\begin{enumerate}
\item 域扩张 $\kappa(\mathfrak m)/\kappa$ 正规；
\item $D \to \text{Aut}(\kappa(\mathfrak m)/\kappa)$ 满射。
\end{enumerate}
若对某个（等价地，每个）极大理想 $\mathfrak m \subset B$，
域扩张 $\kappa(\mathfrak m)/\kappa$ 可分，则：
\begin{enumerate}
\item[(3)] $\kappa(\mathfrak m)/\kappa$ 是 Galois 扩张；
\item[(4)] $D \to \text{Gal}(\kappa(\mathfrak m)/\kappa)$ 满射。
\end{enumerate}
这里 $D \subset G$ 是 $\mathfrak m$ 的分解群。
\end{lemma}

\begin{proof}
注意 $A = B^G$，因为 $A$ 在 $K$ 中整闭且 $K = L^G$。故
(1) 与 (2) 由引理 \ref{more-algebra-lemma-one-orbit-geometric} 得出。
引理中的“等价地，每个”部分由引理 \ref{more-algebra-lemma-galois} 得出。
假设 $\kappa(\mathfrak m)/\kappa$ 可分，则 (3) 与 (4)
立即由 (1) 与 (2) 得出。
\end{proof}

\begin{lemma}
\label{more-algebra-lemma-galois-inertia}
设 $A$ 为分式域是 $K$ 的离散赋值环，$L/K$ 为 Galois 群是
$G$ 的有限 Galois 扩张。令 $B$ 为 $A$ 在 $L$ 中的整闭包，
$\mathfrak m \subset B$ 为极大理想。$\mathfrak m$ 的惰性群
$I$ 位于典范正合列
$$
1 \to P \to I \to I_t \to 1
$$
中，并满足：
\begin{enumerate}
\item 若 $D$ 是分解群，则
$$
P = \{\sigma \in I \mid \sigma|_{\mathfrak m/\mathfrak m^2} =
\text{id}_{\mathfrak m/\mathfrak m^2}\}
=
\{\sigma \in D \mid \sigma \text{ acts trivially on }\text{Gr}_\mathfrak m(B)\}
$$
；
\item $P$ 是 $D$ 的正规子群；
\item 若 $\kappa_A$ 的特征为 $p > 0$，则 $P$ 是 $p$-群；
若 $\kappa_A$ 的特征为零，则 $P = \{1\}$；
\item $I_t$ 是循环群，其阶为整数 $e$ 的与 $p$ 互素部分；
% P02-E0597
\item 存在典范同构
$\theta : I_t \to \mu_e(\kappa(\mathfrak m))$；
\item 若 $\kappa(m)$ 在 $A$ 的剩余域上可分，则
% P02-E0598
$P =
\{\sigma \in D \mid \sigma|_{B/\mathfrak m^2} = \text{id}_{B/\mathfrak m^2}\}$。
\end{enumerate}
这里 $e$ 是引理 \ref{more-algebra-lemma-galois-conclusion} 中的整数。
\end{lemma}

\begin{proof}
回忆 $|G| = [L : K] = nef$；见引理
\ref{more-algebra-lemma-galois-conclusion}。由于 $G$ 传递地作用在 $B$
的极大理想集合
$\{\mathfrak m_1, \ldots, \mathfrak m_n\}$ 上
（引理 \ref{more-algebra-lemma-galois}），且 $D$ 是一个元素的稳定子，
可见 $|D| = ef$。由引理 \ref{more-algebra-lemma-galois-galois}，有
$$
ef = |D| = |I| \cdot |\text{Aut}(\kappa(\mathfrak m)/\kappa)|
$$
，其中 $\kappa$ 是 $A$ 的剩余域。由于
$\kappa(\mathfrak m)$ 在 $\kappa$ 上正规，
$\text{Aut}(\kappa(\mathfrak m)/\kappa)$ 的阶与 $f$
相差 $p$ 的一个幂（见《域》引理
\ref{fields-lemma-normal-and-automorphisms} 以及《域》定义
\ref{fields-definition-insep-degree} 后的讨论）。因此，$|I|$
的与 $p$ 互素部分等于 $e$ 的与 $p$ 互素部分。

\medskip\noindent
置 $C = B_\mathfrak m$。则 $I$ 在 $A$ 上作用于 $C$，并且
平凡作用于 $C$ 的剩余域。令 $\pi_A \in A$ 与 $\pi_C \in C$
为一致化元。对 $C$ 中某个单位 $u$，写
$\pi_A = u \pi_C^e$。对 $\sigma \in I$，对 $C$ 中某个单位
$\theta_\sigma$，写
$\sigma(\pi_C) = \theta_\sigma \pi_C$。于是
% P02-E0599
$$
\pi_A = \sigma(\pi_A) = \sigma(u) (\theta_\sigma \pi_C)^e
= \sigma(u) \theta_\sigma^e \pi_C^e = \frac{\sigma(u)}{u} \theta_\sigma^e \pi_A
$$
。由于 $\sigma \in I$，有
$\sigma(u) \equiv u \bmod \mathfrak m_C$，故
$\theta_\sigma$ 在 $\kappa_C = \kappa(\mathfrak m)$ 中的像
$\overline{\theta}_\sigma$ 是 $e$ 次单位根。我们得到映射
\begin{equation}
\label{more-algebra-equation-inertia-character}
\theta : I \longrightarrow \mu_e(\kappa(\mathfrak m)),\quad
\sigma \mapsto \overline{\theta}_\sigma
\end{equation}
。我们断言 $\theta$ 是群同态，且与一致化元 $\pi_C$ 的选择无关。
事实上，若 $\tau$ 是 $I$ 的另一个元素，则
$\tau(\sigma(\pi_C)) = \tau(\theta_\sigma \pi_C) =
\tau(\theta_\sigma) \theta_\tau \pi_C$，故
$\theta_{\tau \sigma} = \tau(\theta_\sigma) \theta_\tau$。由于
$\tau \in I$，可知
$\overline{\theta}_{\tau \sigma} =
\overline{\theta}_\sigma \overline{\theta}_\tau$。若 $\pi'_C$
是另一个一致化元，则对 $C$ 中某个单位 $w$，有
$\pi'_C = w \pi_C$，并且
$\sigma(\pi'_C) = w^{-1}\sigma(w)\theta_\sigma \pi'_C$。
因此 $\theta'_\sigma = w^{-1}\sigma(w)\theta_\sigma$，
从而 $\theta'_\sigma$ 与 $\theta_\sigma$ 如前一样映到剩余域中的
同一个元素。

\medskip\noindent
由于 $\kappa(\mathfrak m)$ 的特征为 $p$，群
$\mu_e(\kappa(\mathfrak m))$ 是循环群，其阶至多为 $e$
的与 $p$ 互素部分（见《域》第
\ref{fields-section-roots-of-1} 节）。
% P02-E0597

\medskip\noindent
令 $P = \Ker(\theta)$。由构造，$P$ 的元素恰为 $I$
中在 $\mathfrak m/\mathfrak m^2 =
\text{Gr}_\mathfrak m^1(B)$ 上平凡作用的元素，即
$P = \{\sigma \in I \mid \sigma|_{\mathfrak m/\mathfrak m^2} =
\text{id}_{\mathfrak m/\mathfrak m^2}\}$。此外，$I$ 由 $D$
中在 $\kappa(\mathfrak m) = B/\mathfrak m$ 上平凡作用的元素构成。
由于分次环 $\text{Gr}_\mathfrak m(B)$ 由
$\text{Gr}^0_\mathfrak m(B) = B/\mathfrak m$ 上的
$\text{Gr}^1_\mathfrak m(B) = \mathfrak m/\mathfrak m^2$ 生成，
可知
$P = \{\sigma \in D \mid \sigma \text{ acts trivially on }
\text{Gr}_\mathfrak m(B)\}$。所以 (1) 成立。由此也得到 (2)，
因为 $P$ 是同态
$D \to \text{Aut}(\text{Gr}_\mathfrak m(B))$ 的核。
若能证明 (3)，则 (4) 与 (5) 也随之成立：$I_t$
将同构于 $\mu_e(\kappa(\mathfrak m))$，因为上述论证表明
$|I_t| \geq |\mu_e(\kappa(\mathfrak m))|$。

\medskip\noindent
因此，只须证明 $\theta$ 的核 $P$ 是 $p$-群。令 $\sigma$
为该核的非平凡元素。则对所有 $i$，$\sigma - \text{id}$
把 $\mathfrak m_C^i$ 映入 $\mathfrak m_C^{i + 1}$。
令 $m$ 为 $\sigma$ 的阶。选取 $c \in C$，使得
$\sigma(c) \not = c$。则对某个 $i$，有
$\sigma(c) - c \in \mathfrak m_C^i$ 但
$\sigma(c) - c \not \in \mathfrak m_C^{i + 1}$，且
\begin{align*}
0
& =
\sigma^m(c) - c \\
& =
\sigma^m(c) - \sigma^{m - 1}(c) + \ldots + \sigma(c) - c \\
& =
\sum\nolimits_{j = 0, \ldots, m - 1} \sigma^j(\sigma(c) - c) \\
& \equiv
m(\sigma(c) - c) \bmod \mathfrak m_C^{i + 1}
\end{align*}
。由此 $p | m$（若 $p = 1$，则 $m = 0$）。因此，
$\theta$ 的核中每个元素的阶都被 $p$ 整除，即
$\Ker(\theta)$ 是 $p$-群。
% P02-E0600

\medskip\noindent
证明 (6)。假设 $\kappa(\mathfrak m)/\kappa$ 可分。此时，由引理
\ref{more-algebra-lemma-galois-galois}，$f = |D/I|$ 且 $I$ 的阶为 $e$。
若 $e = 1$，则 $P = I = \{\text{id}\}$，结论成立。若
$e > 1$，则 $B/\mathfrak m^2 = C/\pi_C^2C$ 是
$\kappa$-代数，也是 $\kappa(\mathfrak m)$ 的一个小扩张。
由于 $\kappa(\mathfrak m)$ 在 $\kappa$ 上形式
\'etale（《代数》引理
\ref{algebra-lemma-characterize-separable-algebraic-field-extensions}），
存在唯一的 $\kappa$-代数映射
$\kappa(\mathfrak m) \to B/\mathfrak m^2$，它是
$B/\mathfrak m^2 \to \kappa(\mathfrak m)$ 的右逆。换言之，
存在 $\kappa$-代数的 $D$-等变同构
$B/\mathfrak m^2 = \mathfrak m/\mathfrak m^2 \oplus \kappa(\mathfrak m)$。
这立即表明，在此情形下
$P =
\{\sigma \in D \mid \sigma|_{B/\mathfrak m^2} = \text{id}_{B/\mathfrak m^2}\}$。
\end{proof}

\begin{example}
\label{more-algebra-example-counter-description-P}
若没有关于 $\kappa(\mathfrak m)$ 的假设，引理
\ref{more-algebra-lemma-galois-inertia} 的 (6) 中的等式不成立。一个反例是由
$B = A[x]/(x^2 - t)$ 给出的 $A = Z_2[t]_{(2)}$ 的扩张。
事实上，此时 $\sigma(x) = -x = x - 2x$，而
$2x$ 不属于 $\mathfrak m^2$。
\end{example}

\begin{definition}
\label{more-algebra-definition-wild-inertia}
沿用引理 \ref{more-algebra-lemma-galois-inertia} 的假设与记号。
% P02-E0601
\begin{enumerate}
\item $\mathfrak m$ 的{\it 狂惰性群}是子群 $P$。
\item $\mathfrak m$ 的{\it 驯惰性群}是商 $I \to I_t$。
\end{enumerate}
以 $\theta : I \to \mu_e(\kappa(\mathfrak m))$ 表示满射
(\ref{more-algebra-equation-inertia-character})；其核为 $P$，并诱导同构
$I_t \to \mu_e(\kappa(\mathfrak m))$。
% P02-E0602
\end{definition}

\begin{lemma}
\label{more-algebra-lemma-inertia-character}
沿用引理 \ref{more-algebra-lemma-galois-inertia} 的假设与记号。惰性特征标
$\theta : I \to \mu_e(\kappa(\mathfrak m))$ 满足：
% P02-E0603
$$
\theta(\tau \sigma \tau^{-1}) = \tau(\theta(\sigma))
$$
，其中 $\tau \in D$ 且 $\sigma \in I$。
\end{lemma}

\begin{proof}
该公式有意义，因为 $I$ 是 $D$ 的正规子群，并且 $\tau$
通过例如引理 \ref{more-algebra-lemma-galois-galois} 中讨论的映射
$D \to \text{Aut}(\kappa(\mathfrak m))$ 作用在
$\kappa(\mathfrak m)$ 上。回忆 $\theta$ 的构造。选取
$B_\mathfrak m$ 的一致化元 $\pi$，并对 $\sigma \in I$ 写
$\sigma(\pi) = \theta_\sigma \pi$。则 $\theta(\sigma)$ 是
$\theta_\sigma$ 在剩余域中的像
$\overline{\theta}_\sigma$。对任意 $\tau \in D$，可以对某个单位
$\theta_\tau$ 写成 $\tau(\pi) = \theta_\tau \pi$。于是
$\theta_{\tau^{-1}} = \tau^{-1}(\theta_\tau^{-1})$。计算得到
\begin{align*}
\theta_{\tau \sigma \tau^{-1}}
& =
\tau(\sigma(\tau^{-1}(\pi)))/\pi \\
& =
\tau(\sigma(\tau^{-1}(\theta_\tau^{-1}) \pi))/\pi \\
& =
\tau(\sigma(\tau^{-1}(\theta_\tau^{-1})) \theta_\sigma \pi)/\pi \\
& =
\tau(\sigma(\tau^{-1}(\theta_\tau^{-1}))) \tau(\theta_\sigma) \theta_\tau
\end{align*}
。不过，由于 $\sigma$ 模 $\pi$ 平凡作用，乘积
$\tau(\sigma(\tau^{-1}(\theta_\tau^{-1}))) \theta_\tau$
在剩余域中的像为 $1$。这就证明了本引理。
% P02-E0604
\end{proof}

\noindent
我们将在《基本群》引理
\ref{pione-lemma-inertial-invariants-unramified} 中推广下一个引理。

\begin{lemma}
\label{more-algebra-lemma-inertial-invariants-unramified}
设 $A$ 为分式域是 $K$ 的离散赋值环，$L/K$ 为有限 Galois 扩张。
设 $\mathfrak m \subset B$ 为 $A$ 在 $L$ 中的整闭包的一个极大理想，
$I \subset G$ 为 $\mathfrak m$ 的惰性群。则 $B^I$ 是 $A$
在 $L^I$ 中的整闭包，且
$A \to (B^I)_{B^I \cap \mathfrak m}$ 是 \'etale 映射。
\end{lemma}

\begin{proof}
写 $B' = B^I$。由定义可知，$B' = B^I$ 是 $A$ 在 $L^I$
中的整闭包。写
$\mathfrak m' = B^I \cap \mathfrak m = B' \cap \mathfrak m \subset B'$。
由引理 \ref{more-algebra-lemma-one-orbit}，$\mathfrak m$ 是 $B$
中位于 $\mathfrak m'$ 之上的唯一素理想。由于 $I$
平凡作用在 $\kappa(\mathfrak m)$ 上，由引理
\ref{more-algebra-lemma-invariants-modulo} 可见扩张
$\kappa(\mathfrak m)/\kappa(\mathfrak m')$ 纯不可分
（一个也许更简单的替代办法是应用引理
\ref{more-algebra-lemma-one-orbit-geometric} 的结果）。由于 $D/I$
忠实作用在 $\kappa(\mathfrak m')$ 上，可知 $D/I$
忠实作用在 $\kappa(\mathfrak m)$ 上。当然，$A$ 的剩余域
$\kappa$ 的元素在此作用下不变。由 Galois 理论，
$[\kappa(\mathfrak m') : \kappa] \geq |D/I|$；
见《域》引理 \ref{fields-lemma-galois-over-fixed-field}。

\medskip\noindent
令 $\pi$ 为 $A$ 的一致化元。由于
$\text{Norm}_{L/K}(\pi) = \pi^{[L : K]}$，由《代数》引理
\ref{algebra-lemma-finite-extension-dim-1} 可知
% P02-E0605
$$
|G| = [L : K] = [L : K]\ \text{ord}_A(\pi) =
|G/D|\ [\kappa(\mathfrak m) : \kappa]\ \text{ord}_{B_\mathfrak m}(\pi)
$$
；这是因为 $B$ 有 $n = |G/D|$ 个极大理想，且它们在 $G$
下全都共轭；见注记 \ref{more-algebra-remark-finite-separable-extension}
与引理 \ref{more-algebra-lemma-galois}。对次数为 $|I|$ 的有限扩张
$L/L^I$ 应用同样的论证，得到
% P02-E0606
$$
|I|\ \text{ord}_{B'_{\mathfrak m'}}(\pi) =
[\kappa(\mathfrak m) : \kappa(\mathfrak m')]\ \text{ord}_{B_\mathfrak m}(\pi)
$$
。由此
% P02-E0605
$$
\text{ord}_{B'_{\mathfrak m'}}(\pi) =
\frac{|D/I|}{[\kappa(\mathfrak m') : \kappa]}
$$
。由于左侧是正整数，而由上述结论右侧 $\leq 1$，可知必有等号，
即 $\text{ord}_{B'_{\mathfrak m'}}(\pi) = 1$，且
$\kappa(\mathfrak m')/\kappa$ 的次数为 $|D/I|$。因此
% P02-E0607
$\pi B'_{\mathfrak m'} = \mathfrak m' B_\mathfrak m'$，且
$\kappa(\mathfrak m')$ 在 $\kappa$ 上是 Galois 扩张，其
Galois 群为 $D/I$；特别地，它可分，见《域》引理
\ref{fields-lemma-finite-Galois}。由《代数》引理
\ref{algebra-lemma-characterize-etale}，
$A \to B'_{\mathfrak m'}$ 是 \'etale 映射，如所求。
\end{proof}

\begin{remark}
\label{more-algebra-remark-tower-of-rings}
设 $A$ 为分式域是 $K$ 的离散赋值环，$L/K$ 为有限 Galois 扩张。
设 $\mathfrak m \subset B$ 为 $A$ 在 $L$ 中的整闭包的一个极大理想。
令
$$
P \subset I \subset D \subset G
$$
依次为 $\mathfrak m$ 的狂惰性群、惰性群与分解群。
% P02-E0608
考虑图
$$
\xymatrix{
\mathfrak m \ar@{-}[d] \ar@{-}[r] &
\mathfrak m^P \ar@{-}[d] \ar@{-}[r] &
\mathfrak m^I \ar@{-}[d] \ar@{-}[r] &
\mathfrak m^D \ar@{-}[d] \ar@{-}[r] &
A \cap \mathfrak m \ar@{-}[d] \\
B & B^P \ar[l] & B^I \ar[l] & B^D \ar[l] & A \ar[l]
}
$$
。注意，$B^P, B^I, B^D$ 分别是 $A$ 在域
$L^P$、$L^I$、$L^D$ 中的整闭包。因此还可看出，$B^P$
是 $B^I$ 在 $L^P$ 中的整闭包，依此类推。注意
$\mathfrak m^P = \mathfrak m \cap B^P$、
$\mathfrak m^I = \mathfrak m \cap B^I$ 且
$\mathfrak m^D = \mathfrak m \cap B^D$。所以，图的顶行对应于
$\mathfrak m \in \Spec(B)$ 在所诱导的谱映射下的各个像。
综上，有：
\begin{enumerate}
\item 扩张 $L^I/L^D$ 是 Galois 扩张，其群为 $D/I$；
\item 扩张 $L^P/L^I$ 是 Galois 扩张，其群为 $I_t = I/P$；
\item 扩张 $L^P/L^D$ 是 Galois 扩张，其群为 $D/P$；
\item $\mathfrak m^I$ 是 $B^I$ 中位于 $\mathfrak m^D$
之上的唯一素理想；
\item $\mathfrak m^P$ 是 $B^P$ 中位于 $\mathfrak m^I$
之上的唯一素理想；
\item $\mathfrak m$ 是 $B$ 中位于 $\mathfrak m^P$ 之上的唯一素理想；
\item $\mathfrak m^P$ 是 $B^P$ 中位于 $\mathfrak m^D$
之上的唯一素理想；
\item $\mathfrak m$ 是 $B$ 中位于 $\mathfrak m^I$ 之上的唯一素理想；
\item $\mathfrak m$ 是 $B$ 中位于 $\mathfrak m^D$ 之上的唯一素理想；
\item $A \to B^D_{\mathfrak m^D}$ 是 \'etale 映射，并诱导平凡的
剩余域扩张；
\item $B^D_{\mathfrak m^D} \to B^I_{\mathfrak m^I}$ 是
\'etale 映射，并诱导 Galois 群为 $D/I$ 的 Galois 剩余域扩张；
\item $A \to B^I_{\mathfrak m^I}$ 是 \'etale 映射；
\item $B^I_{\mathfrak m^I} \to B^P_{\mathfrak m^P}$
的分歧指数为与 $p$ 互素的 $|I/P|$，并诱导平凡的剩余域扩张；
\item $B^D_{\mathfrak m^D} \to B^P_{\mathfrak m^P}$
的分歧指数为与 $p$ 互素的 $|I/P|$，并诱导可分剩余域扩张；
\item $A \to B^P_{\mathfrak m^P}$ 的分歧指数为与 $p$
互素的 $|I/P|$，并诱导可分剩余域扩张。
% P02-E0609
\end{enumerate}
(1)、(2) 与 (3) 由 Galois 理论（《域》第
\ref{fields-section-galois-theory} 节）及引理
\ref{more-algebra-lemma-galois-inertia} 立即得出。(4)--(9) 由引理
\ref{more-algebra-lemma-galois} 显然成立。(12) 是引理
\ref{more-algebra-lemma-inertial-invariants-unramified}。由于有分解
$A \to B^D_{\mathfrak m^D} \to B^I_{\mathfrak m^I}$，
可随之得到 (10) 与 (11) 中的 \'etale 性。(10) 中的剩余域扩张
必须平凡，因为它可分，并且如引理 \ref{more-algebra-lemma-galois-galois}
所示，$D/I$ 映满
$\text{Aut}(\kappa(\mathfrak m)/\kappa_A)$。同一论证给出 (11)
中关于剩余域的陈述。为证明 (13)、(14) 与 (15)，只须证明 (13)。
由上述结论，扩张 $L^P/L^I$ 是 Galois 扩张，其 Galois 群循环且
阶与 $p$ 互素；$\mathfrak m^P$ 是位于 $\mathfrak m^I$
之上的唯一素理想，且 $I/P$ 在剩余域上的作用平凡。因此，可以对
该扩张及离散赋值环 $B^I_{\mathfrak m^I}$ 应用引理
\ref{more-algebra-lemma-galois-inertia}，得到 (13)。
\end{remark}

\begin{lemma}
\label{more-algebra-lemma-compare-inertia}
设 $A$ 为分式域是 $K$ 的离散赋值环，$M/L/K$ 为扩张塔，其中
$M/K$ 与 $L/K$ 都是有限 Galois 扩张。令 $C$、$B$ 分别为
$A$ 在 $M$、$L$ 中的整闭包。
% P02-E0610
设 $\mathfrak m' \subset C$ 为极大理想，并置
$\mathfrak m = \mathfrak m' \cap B$。令
$$
P \subset I \subset D \subset \text{Gal}(L/K)
\quad\text{且}\quad
P' \subset I' \subset D' \subset \text{Gal}(M/K)
$$
依次为 $\mathfrak m$ 与 $\mathfrak m'$ 的狂惰性群、惰性群及
分解群。则典范满射
$\text{Gal}(M/K) \to \text{Gal}(L/K)$ 诱导满射
$P' \to P$、$I' \to I$ 与 $D' \to D$。此外，它们组成交换图
$$
\vcenter{
\xymatrix{
D' \ar[r] \ar[d] &
\text{Aut}(\kappa(\mathfrak m')/\kappa_A) \ar[d] \\
D \ar[r] &
\text{Aut}(\kappa(\mathfrak m)/\kappa_A)
}
}
\quad\text{且}\quad
\vcenter{
\xymatrix{
I' \ar[r]_-{\theta'} \ar[d] &
\mu_{e'}(\kappa(\mathfrak m')) \ar[d]^{(-)^{e'/e}} \\
I \ar[r]^-\theta &
\mu_e(\kappa(\mathfrak m))
}
}
$$
，其中 $e'$ 与 $e$ 分别是
$A \to C_{\mathfrak m'}$ 与 $A \to B_\mathfrak m$ 的分歧指数。
\end{lemma}

\begin{proof}
在映射 $\text{Gal}(M/K) \to \text{Gal}(L/K)$ 下，各群
$P', I', D'$ 分别映入 $P, I, D$，这由这些群的定义立即得出。
第一个图的交换性是清楚的（注意，由于
$\kappa(\mathfrak m)/\kappa_A$ 正规，
$\kappa(\mathfrak m')$ 在 $\kappa_A$ 上的每个自同构确实诱导
$\kappa(\mathfrak m)$ 在 $\kappa_A$ 上的自同构，故得到第一个图
中的右侧竖直箭头；见引理 \ref{more-algebra-lemma-galois-galois} 与《域》引理
\ref{fields-lemma-lift-maps}）。

\medskip\noindent
由引理 \ref{more-algebra-lemma-one-orbit-geometric-galois-compare}，映射
$I' \to I$ 与 $D' \to D$ 满射。$P' \to P$ 的满射性由
$P'$ 与 $P$ 分别是 $I'$ 与 $I$ 的 Sylow p-子群得出。

\medskip\noindent
为证明第二个图交换，选取 $C_{\mathfrak m'}$ 的一致化元
$\pi'$ 与 $B_\mathfrak m$ 的一致化元 $\pi$。则对
$C_{\mathfrak m'}$ 中某个单位 $c'$，有
$\pi = c' (\pi')^{e'/e}$。对 $\sigma' \in I'$，其像
$\sigma \in I$ 就是 $\sigma'$ 在 $L$ 上的限制。对某个单位
$c \in C_{\mathfrak m'}$ 写
$\sigma'(\pi') = c \pi'$，并对 $B_\mathfrak m$ 中某个单位 $b$
写 $\sigma(\pi) = b \pi$。于是
$\sigma'(\pi) = b \pi$，并得到
% P02-E0611
$$
b \pi = \sigma'(\pi) = \sigma'(c' (\pi')^{e'/e}) =
\sigma'(c') c^{e'/e} (\pi')^{e'/e} =
\frac{\sigma'(c')}{c'} c^{e'/e} \pi
$$
。由于 $\sigma' \in I'$，$b$ 与 $c^{e'/e}$ 在剩余域中的像相同，
这就证明了所需结论。
\end{proof}

\begin{remark}
\label{more-algebra-remark-canonical-inertia-character}
为把惰性特征标
$\theta : I \to \mu_e(\kappa(\mathfrak m))$ 用于无限 Galois 扩张，
对它作缩放是方便的。沿用引理 \ref{more-algebra-lemma-galois-inertia}
与定义 \ref{more-algebra-definition-wild-inertia} 中的
$A, K, L, B, \mathfrak m, G, P, I, D, e, \theta$。
则 $e = q |I_t|$，其中 $q$ 是 $\kappa(\mathfrak m)$ 的特征
$p$ 的一个幂（若该特征为正），或在特征为零时为 $1$。
% P02-E0612
注意，因 $\kappa(\mathfrak m)$ 的特征为 $p$，有
$\mu_e(\kappa(\mathfrak m)) = \mu_{|I_t|}(\kappa(\mathfrak m))$。
考虑映射
% P02-E0613
$$
\theta_{can} = q\theta : I \longrightarrow \mu_{|I_t|}(\kappa(\mathfrak m))
$$
。该映射诱导同构
$\theta_{can} : I_t \to \mu_{|I_t|}(\kappa(\mathfrak m))$。
由引理 \ref{more-algebra-lemma-inertia-character}，对 $\tau \in D$ 与
$\sigma \in I$，有
$\theta_{can}(\tau \sigma \tau^{-1}) = \tau(\theta_{can}(\sigma))$。
最后，若 $M/L$ 是使 $M/K$ 成为 Galois 扩张的扩张，且
$\mathfrak m'$ 是 $A$ 在 $M$ 中的整闭包中位于
$\mathfrak m$ 之上的素理想，则得到交换图
% P02-E0613
$$
\xymatrix{
I' \ar[r]_-{\theta'_{can}} \ar[d] &
\mu_{|I'_t|}(\kappa(\mathfrak m')) \ar[d]^{(-)^{|I'_t|/|I_t|}} \\
I \ar[r]^-{\theta_{can}} &
\mu_{|I_t|}(\kappa(\mathfrak m))
}
$$
；这是引理 \ref{more-algebra-lemma-compare-inertia} 的结论。
\end{remark}

\section{Krasner 引理}
\label{more-algebra-section-krasner}

\noindent
下面给出离散赋值域情形下的 Krasner 引理。

\begin{lemma}[Krasner 引理]
\label{more-algebra-lemma-krasner}
设 $A$ 是维数为 $1$ 的完备局部整环。设 $P(t) \in A[t]$
是系数在 $A$ 中的多项式。设 $\alpha \in A$ 是 $P$ 的根，
但不是导数 $P' = \text{d}P/\text{d}t$ 的根。对每个
$c \geq 0$，存在整数 $n$，使得对任意系数均属于
$\mathfrak m_A^n$ 的 $Q \in A[t]$，多项式 $P + Q$ 都有根
$\beta \in A$，并且 $\beta - \alpha \in \mathfrak m_A^c$。
\end{lemma}

\begin{proof}
选取非零元 $\pi \in \mathfrak m$。
% P02-E0614
由于 $A$ 的维数为 $1$，有 $\mathfrak m = \sqrt{(\pi)}$。
由假设，可对某个 $m \geq 0$ 与 $a \in A$ 写成
$P'(\alpha)^{-1} = \pi^{-m} a$。不妨设
$c \geq m + 1$。选取 $n$，使得
$\mathfrak m_A^n \subset (\pi^{c + m})$。任取命题中的
$Q$。为后面使用，注意可写成
% P02-E0618
$$
P(x + y) = P(x) + P'(x)y + R(x, y)y^2
$$
，其中某个 $R(x, y) \in A[x, y]$。我们将用归纳法证明，
可以找到序列 $\alpha_m, \alpha_{m + 1}, \alpha_{m + 2}, \ldots$，
使得
% P02-E0615
\begin{enumerate}
\item $\alpha_k \equiv \alpha \bmod \pi^c$，
\item $\alpha_{k + 1} - \alpha_k \in (\pi^k)$，并且
\item $(P + Q)(\alpha_k) \in (\pi^{m + k})$。
\end{enumerate}
置 $\beta = \lim \alpha_k$ 即可完成证明。

\medskip\noindent
初始情形。由于 $Q$ 的系数均属于 $(\pi^{c + m})$，有
$(P + Q)(\alpha) \in (\pi^{c + m})$。因此
$\alpha_m = \alpha$ 可行。这个选择保证对所有
$k \geq m$ 都有 $\alpha_k \equiv \alpha \bmod \pi^c$。
% P02-E0615

\medskip\noindent
归纳步骤。给定 $\alpha_k$，对某个 $\delta \in (\pi^k)$ 写
$\alpha_{k + 1} = \alpha_k + \delta$。于是
% P02-E0618
$$
(P + Q)(\alpha_{k + 1}) =
P(\alpha_k + \delta) + Q(\alpha_k + \delta)
$$
。由于 $Q$ 的系数均属于 $(\pi^{c + m})$，可知
$Q(\alpha_k + \delta) \equiv Q(\alpha_k) \bmod \pi^{c + m + k}$。
另一方面，
% P02-E0618
$$
P(\alpha_k + \delta) =
P(\alpha_k) + P'(\alpha_k)\delta + R(\alpha_k, \delta)\delta^2
$$
。注意，由于 $\alpha_k \equiv \alpha \bmod \pi^{m + 1}$，
有 $P'(\alpha_k) \equiv P'(\alpha) \bmod (\pi^{m + 1})$。
因此
% P02-E0618
$$
P(\alpha_k + \delta) \equiv P(\alpha_k) + P'(\alpha) \delta
\bmod \pi^{k + m + 1}
$$
。重新合并两项，得到
% P02-E0618
$$
(P + Q)(\alpha_{k + 1}) \equiv (P + Q)(\alpha_k) + P'(\alpha) \delta
\bmod \pi^{k + m + 1}
$$
。因此可取
$\delta = -P'(\alpha)^{-1} (P + Q)(\alpha_k) =  - \pi^{-m} a (P + Q)(\alpha_k)$；
由归纳假设，它属于 $(\pi^k)$。
\end{proof}

\begin{lemma}
\label{more-algebra-lemma-approximate-separable-extension}
设 $A$ 是分式域为 $K$ 的离散赋值环。设 $A^\wedge$ 是
$A$ 的完备化，其分式域为 $K^\wedge$。若 $M/K^\wedge$
是有限可分扩张，则存在有限可分扩张 $L/K$，使得
$M = K^\wedge \otimes_K L$。
\end{lemma}

\begin{proof}
注意，$A^\wedge$ 也是离散赋值环（由引理
\ref{more-algebra-lemma-completion-regular} 与
\ref{more-algebra-lemma-completion-dimension}）。特别地，$A^\wedge$
是整环。以下证明更一般地适用于满足下述条件的 Noether 局部环
$A$：其 $A^\wedge$ 是维数为 $1$ 的局部整环。

\medskip\noindent
取在 $K^\wedge$ 上生成 $M$ 的元素 $\theta \in M$。
（本原元定理。）设 $P(t) \in K^\wedge[t]$ 是
$\theta$ 在 $K^\wedge$ 上的极小多项式。取非零元
$\pi \in \mathfrak m_A$。把 $\theta$ 替换为
$\pi^n\theta$ 后，可设 $P(t)$ 的系数均属于
$A^\wedge$。令
$B = A^\wedge[\theta] = A^\wedge[t]/(P(t))$。注意，
$B$ 是维数为 $1$ 的完备局部整环，因为它在 $A$ 上有限且包含于
$M$。
% P02-E0616
由于 $M$ 在 $K$ 上可分，元素 $\theta$ 不是 $P$ 的导数的根。
% P02-E0617
对任意整数 $n$，可找到首一多项式 $P_1 \in A[t]$，使
$P - P_1$ 的系数均属于 $\pi^nA^\wedge[t]$。由 Krasner 引理
（引理 \ref{more-algebra-lemma-krasner}），当 $n$ 足够大时，$P_1$
在 $B$ 中有根 $\beta$。此外，可设（若 $n$ 取得足够大）
$\theta - \beta \in \pi B$。考虑把 $t$ 映到 $\beta$ 的
$A^\wedge$-代数映射
$\Phi : A^\wedge[t]/(P_1) \to B$。由于
$B = \pi B + \sum_{i < \deg(P)} A^\wedge \theta^i$，
由 Nakayama 引理，映射 $\Phi$ 是满射。因
$\deg(P_1) = \deg(P)$，可知 $\Phi$ 是同构。于是环扩张
$L = K[t]/(P_1(t))$ 满足 $K^\wedge \otimes_K L \cong M$。
这蕴含 $L$ 是域，证明完成。
\end{proof}

\begin{definition}
\label{more-algebra-definition-mixed}
设 $A$ 是离散赋值环。若 $A$ 的剩余域的特征为
$p > 0$，而 $A$ 的分式域的特征为 $0$，则称
$A$ 具有{\it 混合特征}。在这种情形下，得到离散赋值环扩张
$\mathbf{Z}_{(p)} \subset A$；$A$ 的{\it 绝对分歧指数}
是这个扩张的分歧指数。
\end{definition}
\section{Abhyankar 引理与驯分歧}
\label{more-algebra-section-abhyankar-tame}

\noindent
本节证明我们认为对离散赋值环而言最一般形式的 Abhyankar 引理。
随后应用它证明离散赋值环分式域的驯分歧扩张的一些结果。

\begin{remark}
\label{more-algebra-remark-construction}
设 $A \to B$ 是分式域分别为 $K \subset L$ 的离散赋值环扩张。
设 $K_1/K$ 是有限域扩张。设 $A_1 \subset K_1$ 是
$A$ 在 $K_1$ 中的整闭包。另一方面，令
$L_1 = (L \otimes_K K_1)_{red}$。则 $L_1$ 是由 $L$
的有限域扩张组成的非空有限乘积。设 $B_1$ 是 $B$ 在
$L_1$ 中的整闭包。得到相容的交换图
% P02-E0619
$$
\vcenter{
\xymatrix{
L \ar[r] & L_1 \\
K \ar[u] \ar[r] & K_1 \ar[u]
}
}
\quad\text{且}\quad
\vcenter{
\xymatrix{
B \ar[r] & B_1 \\
A \ar[u] \ar[r] & A_1 \ar[u]
}
}
$$
。在这种情形下，有如下结论：
\begin{enumerate}
\item 由 Algebra，引理
\ref{algebra-lemma-integral-closure-Dedekind}，环 $A_1$
是 Dedekind 整环，而 $B_1$ 是 Dedekind 整环的有限乘积。
\item 注意，$L \otimes_K K_1 = (B \otimes_A A_1)_\pi$，
其中 $\pi \in A$ 是一致化元，并且 $\pi$ 是
$B \otimes_A A_1$ 上的非零因子。因此环映射
$B \otimes_A A_1 \to B_1$ 是整的，其核由幂零元组成。
故 $\Spec(B_1) \to \Spec(B \otimes_A A_1)$ 在谱上是满射
（Algebra，引理
\ref{algebra-lemma-integral-overring-surjective}）。
映射 $\Spec(B \otimes_A A_1) \to \Spec(A_1)$ 是满射，因为
$A_1/\mathfrak m_A A_1 \to
B/\mathfrak m_AB \otimes_{\kappa_A} A_1/\mathfrak m_A A_1$
是单射环映射，且 $A_1/\mathfrak m_A A_1$ 是 Artin 环。
于是 $\Spec(B_1) \to \Spec(A_1)$ 是满射。
\item 设 $\mathfrak m_i$（$i = 1, \ldots n$，其中 $n \geq 1$）
是 $A_1$ 的极大理想。
% P02-E0626
对每个 $i = 1, \ldots, n$，设
$\mathfrak m_{ij}$（$j = 1, \ldots, m_i$，其中 $m_i \geq 1$）
是 $B_1$ 中位于 $\mathfrak m_i$ 之上的极大理想。得到离散赋值环扩张图
% P02-E0619
$$
\xymatrix{
B \ar[r] & (B_1)_{\mathfrak m_{ij}} \\
A \ar[u] \ar[r] & (A_1)_{\mathfrak m_i} \ar[u]
}
$$
。
\item 若 $A$ 是 Hensel 环（例如完备），则 $A_1$ 是离散赋值环，
即 $n = 1$。事实上，$A_1$ 是 $A$ 的有限扩张的并，且这些扩张均为整环，
故由 Algebra，引理
\ref{algebra-lemma-finite-over-henselian}，它们都是局部环。
\item 若 $B$ 是 Hensel 环（例如完备），则 $B_1$
是离散赋值环的乘积，即对 $i = 1, \ldots, n$ 有 $m_i = 1$。
\item 若 $K \subset K_1$ 是纯不可分扩张，则 $A_1$ 与 $B_1$
都是离散赋值环，即 $n = 1$ 且 $m_1 = 1$。这是因为对每个
$b \in B_1$，$b$ 的某个 $p$-幂属于 $B$，故 $B_1$
只能有一个极大理想。
% P02-E0620
\item 若 $K \subset K_1$ 是有限可分扩张，则
$L_1 = L \otimes_K K_1$，而且它也是有限可分扩张的有限乘积。
因此由 Algebra，引理
\ref{algebra-lemma-Noetherian-normal-domain-finite-separable-extension}，
$A \subset A_1$ 与 $B \subset B_1$ 都是有限的。
\item 若 $A$ 是 Nagata 环，则 $A \subset A_1$ 是有限的。
\item 若 $B$ 是 Nagata 环，则 $B \subset B_1$ 是有限的。
\end{enumerate}
\end{remark}

\begin{lemma}
\label{more-algebra-lemma-pull-root-uniformizer}
设 $A$ 是一致化元为 $\pi$ 的离散赋值环。设 $n \geq 2$。
令 $K_1 = K[\pi^{1/n}]$。则
\begin{enumerate}
\item $[K_1 : K] = n$；
\item $A$ 在 $K_1$ 中的整闭包 $A_1$ 是环 $A[\pi^{1/n}]$；
\item $A_1$ 是离散赋值环；
\item $A_1$ 在 $A$ 上的分歧指数为 $n$；
\item $K_1$ 关于 $A$ 完全分歧；并且
\item 若 $n$ 与 $A$ 的剩余特征互素，则 $K_1/K$ 驯分歧，
且 $K_1/K$ 的任意子扩张都由 $\pi^{1/d}$ 生成，其中 $d$
是 $n$ 的某个因子。
\end{enumerate}
\end{lemma}

\begin{proof}
考虑环 $A' = A[x]/(x^n - \pi)$，并用 $\pi'$ 表示 $x$ 在
$A_1$ 中的像。
% P02-E0621
% P02-E0622
由于 $A'$ 作为 $A$-模是秩为 $n$ 的有限自由模，元素
$\pi$ 在 $A'$ 中是非零因子，因而 $\pi'$ 也是非零因子。
此外，$A'/\pi'A'$ 同构于域 $A/\pi A$。故 $A'$ 是以
$\pi'$ 为一致化元的离散赋值环。因此，$A'$ 的分式域 $K'$
是次数为 $n$ 的扩张。显然，$K'$ 由在 $K$ 上添加 $\pi$
的一个 $n$ 次根 $\pi' = \pi^{1/n}$ 得到，即
$K' \cong K_1$。
% P02-E0623
由于 $A'$ 正规，可知 $A'$ 是 $A$ 在 $K'$ 中的整闭包，
即 $A' = A_1$。这证明了 (1)--(5)。

\medskip\noindent
由定义 \ref{more-algebra-definition-types-of-extensions} 及 (3)、(4)，
若 $n$ 在剩余域中可逆，则 $K_1/K$ 驯分歧的断言是直接的。
为证明最后一个陈述，只须证明每个 $n$ 次单位根
$\zeta \in K'$ 都属于 $K$；见 Fields，引理
\ref{fields-lemma-subfields-kummer}。显然
$\zeta \in (A')^*$。可写 $\zeta = a + b$，其中
$a \in A$ 且
$b \in A \pi' \oplus \ldots \oplus A (\pi')^{n - 1}$。
若 $b$ 非零，则可写 $b = u(\pi')^t$，其中
$t \geq 1$，且 $b$ 是 $A'$ 的单位。
% P02-E0624
于是
% P02-E0619
$$
1 = \zeta^n \equiv a^n + n a^{n - 1} u (\pi')^t \bmod (\pi')^{t + 1}
$$
。由于 $n$ 与 $a$ 都是 $A$ 的单位，这将推出
$b = u(\pi')^t$ 模 $(\pi')^{t + 1}$ 同余于
$(1 - a^n)/(n a^{n - 1}) \in A$。直和分解
$A' = A \oplus \ldots \oplus A(\pi')^{n - 1}$ 表明这是不可能的，
因为乘以 $(\pi')^{t + 1}$ 的像是相应的直和，而它与
$A(\pi')^t$ 的交为 $(\pi A)(\pi')^t$。故 $b = 0$，
并如愿得到 $\zeta \in A$。
\end{proof}

\begin{lemma}
\label{more-algebra-lemma-formally-smooth-goes-up}
设 $A \to B$ 是分式域为 $K \subset L$ 的离散赋值环扩张。
假设 $A \to B$ 在 $\mathfrak m_B$-进拓扑下形式光滑。
则对任意有限扩张 $K_1/K$，有
$L_1 = L \otimes_K K_1$、$B_1 = B \otimes_A A_1$，而且每个扩张
$(A_1)_{\mathfrak m_i} \subset (B_1)_{\mathfrak m_{ij}}$
（见评注 \ref{more-algebra-remark-construction}）在
$\mathfrak m_{ij}$-进拓扑下形式光滑。
\end{lemma}

\begin{proof}
我们将使用引理 \ref{more-algebra-lemma-extension-dvrs-formally-smooth} 的等价条件，
不再另行说明。设 $\pi \in A$ 与
$\pi_i \in (A_1)_{\mathfrak m_i}$ 分别为一致化元。
由于 $\kappa_A \subset \kappa_B$ 可分，环
% P02-E0619
$$
(B \otimes_A (A_1)_{\mathfrak m_i})/\pi_i (B \otimes_A (A_1)_{\mathfrak m_i}) =
B/\pi B \otimes_{A/\pi A} (A_1)_{\mathfrak m_i}/\pi_i (A_1)_{\mathfrak m_i}
$$
是若干域的乘积，其中每个域在 $\kappa_{\mathfrak m_i}$ 上可分。
因此 $B \otimes_A (A_1)_{\mathfrak m_i}$ 中的元素 $\pi_i$
是非零因子，且对它取商得到域的乘积。由此
$B \otimes_A A_1$ 是 Dedekind 整环，特别地是既约的。
% P02-E0625
因此 $B \otimes_A A_1 \subset B_1$ 是等式。
\end{proof}


\noindent
下面的引理是我们针对离散赋值环所用的 Abhyankar 引理版本。
注意，并不假设 $\kappa_B/\kappa_A$ 是代数域扩张。

\begin{lemma}[Abhyankar 引理]
\label{more-algebra-lemma-abhyankar}
设 $A \subset B$ 是离散赋值环扩张。假设 $A$ 的剩余特征为
$0$；或者其剩余特征为 $p$，分歧指数 $e$ 与 $p$ 互素，且
$\kappa_B/\kappa_A$ 是可分域扩张。设 $K_1/K$ 是有限扩张。
使用评注 \ref{more-algebra-remark-construction} 的记号，并假设对某个 $i$，
$e$ 整除 $A \subset (A_1)_{\mathfrak m_i}$ 的分歧指数。
则对所有 $j = 1, \ldots, m_i$，扩张
$(A_1)_{\mathfrak m_i} \subset (B_1)_{\mathfrak m_{ij}}$
在 $\mathfrak m_{ij}$-进拓扑下形式光滑。
\end{lemma}

\begin{proof}
设 $\pi \in A$ 是一致化元。设 $\pi_1$ 是
$(A_1)_{\mathfrak m_i}$ 的一致化元。写成
$\pi = u \pi_1^{e_1}$，其中 $u$ 是
$(A_1)_{\mathfrak m_i}$ 的单位，而 $e_1$ 是
$A \subset (A_1)_{\mathfrak m_i}$ 的分歧指数。

\medskip\noindent
断言：可设 $u$ 是 $K_1$ 中的 $e$ 次幂。事实上，设
$K_2$ 是通过添加 $x^e = u$ 的一个根而得到的 $K_1$ 的扩张；
于是 $K_2$ 是 $K_1[x]/(x^e - u)$ 的一个因子。由关于 $e$
的假设，$K_2/K_1$ 是有限可分扩张，因而 $A_1 \subset A_2$
是有限的。由于
$(A_1)_{\mathfrak m_i} \to (A_1)_{\mathfrak m_i}[x]/(x^e - u)$
是有限 \'etale 的（因为 $e$ 与剩余特征互素且 $u$ 是单位），
可知 $(A_2)_{\mathfrak m_i}$ 是
$(A_1)_{\mathfrak m_i}$ 的某个有限 \'etale 扩张的因子，
故其本身在 $(A_1)_{\mathfrak m_i}$ 上有限 \'etale。
% P02-E0627
同样的推理表明，$B_1 \subset B_2$ 诱导有限 \'etale 扩张
$(B_1)_{\mathfrak m_{ij}} \subset (B_2)_{\mathfrak m_{ij}}$。
% P02-E0627
选取极大理想 $\mathfrak m'_{ij} \subset B_2$，使其位于
$\mathfrak m_{ij} \subset B_1$ 之上（当然可能不止一个），并考虑
% P02-E0630
$$
\xymatrix{
(B_1)_{\mathfrak m_{ij}} \ar[r] & (B_2)_{\mathfrak m'_{ij}} \\
(A_1)_{\mathfrak m_i} \ar[u] \ar[r] &
(A_2)_{\mathfrak m'_i} \ar[u]
}
$$
，其中 $\mathfrak m'_i \subset A_2$ 是其像。
% P02-E0628
现在，水平箭头的分歧指数为 $1$，并诱导有限可分剩余域扩张。
因此，使用引理 \ref{more-algebra-lemma-extension-dvrs-formally-smooth} 的等价条件，
只须证明右侧竖直箭头在 $\mathfrak m'_{ij}$-进拓扑下形式光滑。
由于 $u$ 在 $K_2$ 中有一个 $e$ 次根，断言得证。
% P02-E0629

\medskip\noindent
假设 $u$ 在 $K_1$ 中有一个 $e$ 次根。由于 $e | e_1$，且
$u$ 在 $K_1$ 中有一个 $e$ 次根，可知对某个
$\theta \in K_1$ 有 $\pi = \theta^e$。
% P02-E0629
设 $K'_1 = K[\theta] \subset K_1$ 是由 $\theta$ 生成的子域。
由引理 \ref{more-algebra-lemma-pull-root-uniformizer}，$A$ 在
$K[\theta]$ 中的整闭包 $A'_1$ 是离散赋值环
$A'_1 = A[\theta]$，它在 $A$ 上的分歧指数为 $e$。
若能对扩张 $K'_1/K$ 证明本引理，则对图
% P02-E0630
$$
\xymatrix{
(B'_1)_{B'_1 \cap \mathfrak m_{ij}} \ar[r] & (B_1)_{\mathfrak m_{ij}} \\
A'_1 \ar[u] \ar[r] &
(A_1)_{\mathfrak m_i} \ar[u]
}
$$
应用引理 \ref{more-algebra-lemma-formally-smooth-goes-up}，即可对所有
$j = 1, \ldots, m_i$ 得到结论。这把问题约化为下一段讨论的情形。

\medskip\noindent
假设 $K_1 = K[\pi^{1/e}]$，并置 $\theta = \pi^{1/e}$。
设 $\pi_B$ 是 $B$ 的一致化元，并对 $B$ 的某个单位 $w$
写成 $\pi = w \pi_B^e$。于是可见
$L_1 = L \otimes_K K_1$ 通过添加 $\theta / \pi_B$ 得到，
后者是单位 $w$ 的一个 $e$ 次根。因此 $B \subset B_1$
是有限 \'etale 的。于是对任意极大理想
$\mathfrak m \subset B_1$，考虑交换图
% P02-E0630
$$
\xymatrix{
B \ar[r]_1 & (B_1)_{\mathfrak m} \\
A \ar[u]^e \ar[r]^e & A_1 \ar[u]_{e_\mathfrak m}
}
$$
。箭头旁的数是分歧指数。由分歧指数的乘法性
（引理 \ref{more-algebra-lemma-multiplicative-e-f}），得到
$e_\mathfrak m = 1$。考察剩余域扩张，可知
$\kappa(\mathfrak m)$ 是 $\kappa_B$ 的有限可分扩张，
而后者在 $\kappa_A$ 上可分。因此
$\kappa(\mathfrak m)$ 在 $\kappa_A$ 上可分；后者
等于 $A_1$ 的剩余域。于是由引理
\ref{more-algebra-lemma-extension-dvrs-formally-smooth} 得到所需结论。
\end{proof}

\begin{lemma}
\label{more-algebra-lemma-composition-tame}
设 $A$ 是分式域为 $K$ 的离散赋值环。设 $M/L/K$
是有限可分扩张。设 $B$ 是 $A$ 在 $L$ 中的整闭包。
若 $L/K$ 关于 $A$ 驯分歧，并且对 $B$ 的每个极大理想
$\mathfrak m$，$M/L$ 关于 $B_\mathfrak m$ 驯分歧，
则 $M/K$ 关于 $A$ 驯分歧。
\end{lemma}

\begin{proof}
设 $C$ 是 $A$ 在 $M$ 中的整闭包。$C$ 的每个极大理想
$\mathfrak m'$ 都位于 $B$ 的某个极大理想 $\mathfrak m$ 之上。
于是本引理由分歧指数的乘法性（引理
\ref{more-algebra-lemma-multiplicative-e-f}）以及有限域扩张塔
$\kappa(\mathfrak m')/\kappa(\mathfrak m)/\kappa_A$ 得出。
\end{proof}

\begin{lemma}
\label{more-algebra-lemma-subextension-tame}
设 $A$ 是分式域为 $K$ 的离散赋值环。若 $M/L/K$
是有限可分扩张，且 $M$ 关于 $A$ 驯分歧，则 $L$
关于 $A$ 驯分歧。
\end{lemma}

\begin{proof}
我们将使用评注 \ref{more-algebra-remark-finite-separable-extension} 中讨论的结果，
不再另行说明。设 $C/B/A$ 是 $A$ 在 $M/L/K$ 中的整闭包。
% P02-E0631
由于 $C$ 是 $B$ 的有限环扩张，可知
$\Spec(C) \to \Spec(B)$ 是满射。因此，对每个极大理想
$\mathfrak m \subset B$，都存在极大理想
$\mathfrak m' \subset C$ 位于 $\mathfrak m$ 之上。
% P02-E0632
由分歧指数的乘法性（引理
\ref{more-algebra-lemma-multiplicative-e-f}）及假设，得到
$B_\mathfrak m$ 在 $A$ 上的分歧指数与剩余特征互素。
由于 $\kappa(\mathfrak m')/\kappa_A$ 有限可分，
$\kappa(\mathfrak m)/\kappa_A$ 也具有同一性质。
\end{proof}


\begin{lemma}
\label{more-algebra-lemma-characterize-tame}
设 $A$ 是分式域为 $K$ 的离散赋值环。设
$\pi \in A$ 是一致化元。设 $L/K$ 是有限可分扩张。
下列条件等价：
\begin{enumerate}
\item $L$ 关于 $A$ 驯分歧；
\item 存在一个在 $\kappa_A$ 中可逆的 $e \geq 1$，以及一个关于
$A' = A[\pi^{1/e}]$ 非分歧的扩张
$L'/K' = K[\pi^{1/e}]$，使得 $L$ 包含于 $L'$；并且
\item 存在一个在 $\kappa_A$ 中可逆的 $e_0 \geq 1$，使得对每个
在 $\kappa_A$ 中可逆的 $d \geq 1$，(2) 对
$e = de_0$ 成立。
\end{enumerate}
\end{lemma}

\begin{proof}
注意，$A'$ 是分式域为 $K'$ 的离散赋值环；见引理
\ref{more-algebra-lemma-pull-root-uniformizer}。当然，$A'$ 在 $A$ 上的分歧指数为
$e$。因此，若 (2) 成立，则由引理 \ref{more-algebra-lemma-composition-tame}，
$L'$ 关于 $A$ 驯分歧。再由引理
\ref{more-algebra-lemma-subextension-tame}，$L$ 关于 $A$ 驯分歧。

\medskip\noindent
蕴含 (3) $\Rightarrow$ (2) 是直接的。

\medskip\noindent
假设 (1) 成立。设 $B$ 是 $A$ 在 $L$ 中的整闭包，并设
$\mathfrak m_1, \ldots, \mathfrak m_n$ 是它的极大理想。
用 $e_i$ 表示 $A \to B_{\mathfrak m_i}$ 的分歧指数。
% P02-E0633
设 $e_0$ 是 $e_1, \ldots, e_r$ 的最小公倍数。
% P02-E0634
由假设 (1)，它在 $\kappa_A$ 中可逆。按 (3) 令
$e = de_0$。置 $A' = A[\pi^{1/e}]$。则
$A \to A'$ 是分式域为 $K' = K[\pi^{1/e}]$ 的离散赋值环扩张；
见引理 \ref{more-algebra-lemma-pull-root-uniformizer}。选取乘积分解
% P02-E0636
$$
L \otimes_K K' = \prod L'_j
$$
，其中 $L'_j$ 是域。设 $B'_j$ 是 $A$ 在 $L'_j$ 中的整闭包。
设 $\mathfrak m_{ijk}$ 是 $B'_j$ 中位于 $\mathfrak m_i$
之上的极大理想。注意，$(B'_j)_{\mathfrak m_i}$ 是
$B_{\mathfrak m_i}$ 在 $L'_j$ 中的整闭包。
% P02-E0635
把 Abhyankar 引理（引理 \ref{more-algebra-lemma-abhyankar}）应用于
$A \subset B_{\mathfrak m_i}$ 与扩张 $K'/K$，可知
$A' \to (B'_j)_{\mathfrak m_{ijk}}$ 在
$\mathfrak m_{ijk}$-进拓扑下形式光滑。这蕴含分歧指数为
$1$ 且剩余域扩张可分（引理
\ref{more-algebra-lemma-extension-dvrs-formally-smooth}）。
由此可见，$L'_j$ 关于 $A'$ 非分歧。证明完成：任取某个
$j$，令 $L' = L'_j$。
\end{proof}


\begin{lemma}
\label{more-algebra-lemma-permanence-tame}
设 $A$ 是分式域为 $K$ 的离散赋值环。
\begin{enumerate}
\item 若 $L/K$ 是关于 $A$ 驯分歧的有限可分扩张，则存在
包含 $L$ 且关于 $A$ 驯分歧的 Galois 扩张 $M/K$。
\item 若 $L_1/K$、$L_2/K$ 是关于 $A$ 驯分歧的有限可分扩张，
则存在包含 $L_1$ 与 $L_2$ 且关于 $A$ 驯分歧的有限可分扩张
$L/K$。
\end{enumerate}
\end{lemma}

\begin{proof}
证明 (2)。选取一致化元 $\pi \in A$。可以选取在
$\kappa_A$ 中可逆的整数 $e$，以及关于
$A' = A[\pi^{1/e}]$ 非分歧的扩张
$L_i'/K' = K[\pi^{1/e}]$，使得作为 $K$ 的扩张有
$L'_i/L_i$；见引理 \ref{more-algebra-lemma-characterize-tame}。
% P02-E0637
由引理 \ref{more-algebra-lemma-permanence-unramified}，可以找到关于 $A'$
非分歧的扩张 $L'/K'$，使得对 $i = 1, 2$，
$L'_i/K$ 同构于 $L'/K'$ 的一个子扩张。
% P02-E0638
由于 $L'/K$ 驯分歧，这完成了 (3) 的证明
（使用上面同一个引理）。
% P02-E0639

\medskip\noindent
证明 (1)。首先可以把 $L$ 替换为更大的扩张，并假设
$L$ 是 $K' = K[\pi^{1/e}]$ 的扩张，关于
$A' = A[\pi^{1/e}]$ 非分歧，其中 $e$ 在 $\kappa_A$
中可逆；见引理 \ref{more-algebra-lemma-characterize-tame}。
设 $M$ 是 $L$ 在 $K$ 上的正规闭包；见 Fields，定义
\ref{fields-definition-normal-closure}。由 Fields，引理
\ref{fields-lemma-normal-closure-galois}，$M/K$ 是 Galois 扩张。
另一方面，有 $K$-代数满射
% P02-E0640
% P02-E0642
$$
L \otimes_K \ldots \otimes_K L \longrightarrow M
$$
；见 Fields，引理
\ref{fields-lemma-normal-closure-tensor-product}。
如评注 \ref{more-algebra-remark-finite-separable-extension}，设 $B$ 是
$A$ 在 $L$ 中的整闭包。$L$ 关于
$A' = A[\pi^{1/e}]$ 非分歧这一条件，恰好意味着
$A' \to B$ 是 \'etale 环映射；见 Algebra，引理
\ref{algebra-lemma-characterize-etale}。
断言：
% P02-E0640
% P02-E0642
$$
K' \otimes_K \ldots \otimes_K K' = \prod K'_i
$$
是关于 $A$ 驯分歧的域扩张 $K'_i/K$ 的乘积。于是，若
$A'_i$ 是 $A$ 在 $K'_i$ 中的整闭包，则
% P02-E0640
% P02-E0642
$$
\prod A'_i \otimes_{(A' \otimes_A \ldots \otimes_A A')}
(B \otimes_A \ldots \otimes_A B)
$$
在 $\prod A'_i$ 上有限 \'etale，因而是 Dedekind 整环的乘积
（引理 \ref{more-algebra-lemma-Dedekind-etale-extension}）。
可知 $M$ 是这些 Dedekind 整环之一的分式域，而该整环在
$A'_i$ 上有限 \'etale，其中 $i$ 为某个指标。因此 $M/K'_i$ 关于
$A'_i$ 的每个极大理想都非分歧，从而由引理
\ref{more-algebra-lemma-composition-tame}，$M/K$ 驯分歧。

\medskip\noindent
还需证明该断言。为此写成
$A' = A[x]/(x^e - \pi)$，并且可见
% P02-E0641
% P02-E0640
% P02-E0642
$$
A' \otimes_A \ldots \otimes_A A' =
A'[x_1, \ldots, x_r]/(x_1^e - \pi, \ldots, x_r^e - \pi)
$$
该环的正规化当然包含元素 $y_i = x_i/x_1$（其中
$i = 2, \ldots, r$），它们满足关系 $y_i^e - 1 = 0$，并得到
% P02-E0642
$$
A[x_1, y_2, \ldots, y_r]/(x_1^e - \pi, y_2^e - 1, \ldots, y_r^e - 1) =
A'[y_2, \ldots, y_r]/(y_2^e - 1, \ldots, y_r^e - 1)
$$
。由于 $e$ 在 $A'$ 中可逆，该环在 $A'$ 上有限 \'etale。
因此它是 Dedekind 整环的乘积，每个因子在 $A'$ 上非分歧，
如愿得证（若有疑惑，参见上面给出的参考文献）。
\end{proof}


\begin{lemma}
\label{more-algebra-lemma-tame-goes-up}
设 $A \subset B$ 是离散赋值环扩张。用 $L/K$ 表示相应的分式域扩张。
设 $K'/K$ 是有限可分扩张。则
% P02-E0644
$$
K' \otimes_K L = \prod L'_i
$$
是域的有限乘积，并且下列结论成立：
\begin{enumerate}
\item 若 $K'$ 关于 $A$ 非分歧，则每个 $L'_i$ 关于 $B$ 非分歧。
\item 若 $K'$ 关于 $A$ 驯分歧，则每个 $L'_i$ 关于 $B$ 驯分歧。
\end{enumerate}
\end{lemma}

\begin{proof}
代数 $K' \otimes_K L$ 是 $L$ 上的有限 \'etale 代数，
故为域的有限乘积。设 $A'$ 是 $A$ 在 $K'$ 中的整闭包。

\medskip\noindent
在情形 (1) 中，环映射 $A \to A'$ 有限 \'etale。因此
$B' = B \otimes_A A'$ 在 $B$ 上有限 \'etale，而且是
Dedekind 整环的有限乘积（引理
\ref{more-algebra-lemma-Dedekind-etale-extension}）。故 $B'$ 是
$B$ 在 $K' \otimes_K L$ 中的整闭包。立即可知每个
$L'_i$ 关于 $B$ 非分歧。

\medskip\noindent
选取一致化元 $\pi \in A$。为证明 (2)，可把 $K'$ 替换为
关于 $A$ 驯分歧的更大扩张（略去细节；提示：使用引理
\ref{more-algebra-lemma-subextension-tame}）。
% P02-E0643
因此，由引理 \ref{more-algebra-lemma-characterize-tame}，可设存在某个在
$\kappa_A$ 中可逆的 $e \geq 1$，使得 $K'$ 包含
$K[\pi^{1/e}]$，并且 $K'$ 关于 $A[\pi^{1/e}]$ 非分歧。
选取乘积分解
% P02-E0644
$$
K[\pi^{1/e}] \otimes_K L = \prod L_{e, j}
$$
。对每个 $i$，存在 $j_i$，使得 $L'_i/L_{e, j_i}$
是有限可分扩张。设 $B_{e, j}$ 是 $B$ 在 $L_{e, j}$
中的整闭包。把 (1) 应用于 $K'/K[\pi^{1/e}]$ 与
$A[\pi^{1/e}] \subset (B_{e, j_i})_\mathfrak m$，可知
对每个极大理想 $\mathfrak m \subset B_{e, j_i}$，
$L'_i$ 关于 $(B_{e, j_i})_\mathfrak m$ 非分歧。
因此，由引理 \ref{more-algebra-lemma-composition-tame}，只须证明
$L_{e, j}$ 关于 $B$ 驯分歧。

\medskip\noindent
在 $B$ 中选取一致化元 $\theta$。写成
$\pi = u \theta^t$，其中 $u$ 是 $B$ 的单位且 $t \geq 1$。
于是
% P02-E0644
$$
A[\pi^{1/e}] \otimes_A B = B[x]/(x^e - u \theta^t)
\subset B[y, z]/(y^{e'} - \theta, z^e - u)
$$
，其中 $e' = e/\gcd(e, t)$。该映射把 $x$ 映到
$z y^{t/\gcd(e, t)}$。由于右侧是 Dedekind 整环的乘积，
其中每个因子在 $B$ 上驯分歧，证明完成（略去细节）。
\end{proof}


\section{消除分歧}
\label{more-algebra-section-eliminating-ramification}

\noindent
本节讨论 Helmut Epp 的一个结果，见 \cite{Epp}。我们强烈建议读者阅读原文。
我们的方法略有不同，因为我们试图用同一种方法处理混合特征与等特征情形。
相关结果还可参见
\cite{Ponomarev}、\cite{Ponomarev-Abhyankar}、\cite{Kuhlmann}
与 \cite{ZK}。

\medskip\noindent
设 $A \subset B$ 是分式域为 $K \subset L$ 的离散赋值环扩张。
本节的目标是找到有限扩张 $K_1/K$，使得在
% P02-E0645
$$
\vcenter{
\xymatrix{
L \ar[r] & L_1 \\
K \ar[u] \ar[r] & K_1 \ar[u]
}
}
\quad\text{且}\quad
\vcenter{
\xymatrix{
B \ar[r] & B_1 \ar[r] & (B_1)_{\mathfrak m_{ij}} \\
A \ar[u] \ar[r] & A_1 \ar[r] \ar[u] & (A_1)_{\mathfrak m_i} \ar[u]
}
}
$$
如评注 \ref{more-algebra-remark-construction} 所示时，扩张
$(A_1)_{\mathfrak m_i} \subset (B_1)_{\mathfrak m_{ij}}$
全都弱非分歧，乃至在相应进拓扑下形式光滑。
最简单（但非平凡）的例子是 Abhyankar 引理，见引理
\ref{more-algebra-lemma-abhyankar}。

\begin{definition}
\label{more-algebra-definition-solution}
设 $A \to B$ 是分式域为 $K \subset L$ 的离散赋值环扩张。
\begin{enumerate}
\item 若评注 \ref{more-algebra-remark-construction} 的所有扩张
$(A_1)_{\mathfrak m_i} \subset (B_1)_{\mathfrak m_{ij}}$
都弱非分歧，则称有限域扩张 $K_1/K$ 是
$A \subset B$ 的一个{\it 弱解}。
\item 若评注 \ref{more-algebra-remark-construction} 的每个扩张
$(A_1)_{\mathfrak m_i} \subset (B_1)_{\mathfrak m_{ij}}$
都在 $\mathfrak m_{ij}$-进拓扑下形式光滑，则称有限域扩张
$K_1/K$ 是 $A \subset B$ 的一个{\it 解}。
\end{enumerate}
若解 $K_1/K$ 满足 $K_1/K$ 可分，则称它是一个{\it 可分的解}。
\end{definition}

\noindent
一般而言，弱解或解未必存在；\cite{Epp} 中有一个例子。
在对剩余域扩张施加温和假设后，我们将依照 \cite{Epp} 在定理
\ref{more-algebra-theorem-epp} 中证明弱解的存在性。下一节将从中推出几何上有意义的
情形中解的存在性，有时还可推出可分的解的存在性；见命题
\ref{more-algebra-proposition-epp-essentially-finite-type} 与引理
\ref{more-algebra-lemma-epp-essentially-finite-type-separable}。
不过，下面的例子表明，一般而言，即使为了得到弱解也需要不可分扩张。

\begin{example}
\label{more-algebra-example-inseparable-necessary}
设 $k$ 是特征为 $p > 0$ 的完美域。令 $A = k[[x]]$ 且
$K = k((x))$。令 $B = A[x^{1/p}]$。$A \to B$ 的任意弱解
$K_1/K$ 都不可分（而 $K$ 的任意有限不可分扩张都是解）。
略去证明。
\end{example}

\noindent
解在进一步扩张下稳定，见引理 \ref{more-algebra-lemma-solution-goes-up}。
弱解未必具有这一性质。弱解在某种意义下对完全分歧扩张稳定，见引理
\ref{more-algebra-lemma-weakly-unramified-goes-up-along-totally-ramified}。

\begin{lemma}
\label{more-algebra-lemma-weakly-unramified-goes-up-along-totally-ramified}
设 $A \to B$ 是分式域为 $K \subset L$ 的离散赋值环扩张。
假设 $A \to B$ 弱非分歧。则对任意关于 $A$ 完全分歧的有限可分扩张
$K_1/K$，$L_1 = L \otimes_K K_1$ 是域，$A_1$ 与
$B_1 = B \otimes_A A_1$ 是离散赋值环，而且扩张
$A_1 \subset B_1$（见评注 \ref{more-algebra-remark-construction}）弱非分歧。
\end{lemma}

\begin{proof}
设 $\pi \in A$ 与 $\pi_1 \in A_1$ 是一致化元。由于
$K_1/K$ 关于 $A$ 完全分歧，对 $A_1$ 中的某个单位 $u_1$ 有
$\pi_1^e = u_1 \pi$。因此 $A_1$ 在 $A$ 上由 $\pi_1$ 生成，
而 $\pi_1$ 在 $K$ 上的极小多项式 $P(t)$ 具有形式
% P02-E0645
$$
P(t) = t^e + a_{e - 1} t^{e - 1} + \ldots + a_0
$$
，其中 $a_i \in (\pi)$，且对 $A$ 的某个单位 $u$ 有
$a_0 = u\pi$。还要注意，$e = [K_1 : K]$。由于
$A \to B$ 弱非分歧，$\pi$ 是 $B$ 的一致化元，故
$B_1 = B[t]/(P(t))$ 是离散赋值环，其一致化元是 $t$ 的类。
引理于是显然。
\end{proof}

\begin{lemma}
\label{more-algebra-lemma-solutions-go-down}
设 $A \to B \to C$ 是分式域为
$K \subset L \subset M$ 的离散赋值环扩张。设 $K_1/K$ 是有限扩张。
\begin{enumerate}
\item 若 $K_1$ 是 $A \to C$ 的弱解或解，则 $K_1$ 是
$A \to B$ 的弱解或解。
\item 若 $K_1$ 是 $A \to B$ 的弱解或解，且
$L_1 = (L \otimes_K K_1)_{red}$ 是若干域的乘积，其中每个域都是
$B \to C$ 的弱解或解，则 $K_1$ 是 $A \to C$ 的弱解或解。
\end{enumerate}
\end{lemma}

\begin{proof}
令 $L_1 = (L \otimes_K K_1)_{red}$ 与
$M_1 = (M \otimes_K K_1)_{red}$，并设
$B_1 \subset L_1$ 与 $C_1 \subset M_1$ 是 $B$ 与 $C$ 的整闭包。
% P02-E0646
注意，$M_1 = (M \otimes_L L_1)_{red}$，而 $L_1$ 是由
$L$ 的有限扩张组成的非空有限乘积。因此环映射
$B_1 \to C_1$ 是评注 \ref{more-algebra-remark-construction} 中所讨论形式的环映射的
有限乘积。特别地，映射 $\Spec(C_1) \to \Spec(B_1)$ 是满射。
选取极大理想 $\mathfrak m \subset C_1$，并考虑离散赋值环扩张
% P02-E0645
$$
(A_1)_{A_1 \cap \mathfrak m} \to
(B_1)_{B_1 \cap \mathfrak m} \to
(C_1)_\mathfrak m
$$
。若复合映射弱非分歧，则映射
$(A_1)_{A_1 \cap \mathfrak m} \to
(B_1)_{B_1 \cap \mathfrak m}$ 也弱非分歧。
若剩余域扩张
$\kappa_{A_1 \cap \mathfrak m} \to \kappa_\mathfrak m$
可分，则子扩张
$\kappa_{A_1 \cap \mathfrak m} \to
\kappa_{B_1 \cap \mathfrak m}$ 也可分。结合引理
\ref{more-algebra-lemma-extension-dvrs-formally-smooth}，这证明了 (1)。
类似论证适用于 (2)。
\end{proof}

\begin{lemma}
\label{more-algebra-lemma-solution-after-strict-henselization}
设 $A \to B$ 是离散赋值环扩张。存在离散赋值环扩张的交换图
% P02-E0645
$$
\xymatrix{
B \ar[r] & B' \\
A \ar[r] \ar[u] & A' \ar[u]
}
$$
，使得
\begin{enumerate}
\item 分式域扩张 $K'/K$ 与 $L'/L$ 可分代数；
\item $A'$ 与 $B'$ 的剩余域分别是 $A$ 与 $B$ 的剩余域的
可分代数闭包；并且
\item 若 $A' \to B'$ 存在解、弱解或可分的解，则
$A \to B$ 存在解、弱解或可分的解。
\end{enumerate}
\end{lemma}

\begin{proof}
由 Algebra，引理
\ref{algebra-lemma-colimit-finite-etale-given-residue-field}，
存在扩张 $A \subset A'$，它是有限 \'etale 扩张的滤过余极限，
且 $A'$ 的剩余域是 $A$ 的剩余域的可分代数闭包。
于是 $A \subset A'$ 是离散赋值环扩张，所诱导的分式域扩张
$K'/K$ 可分代数。

\medskip\noindent
设 $B \subset B'$ 是 $B$ 的严格 Hensel 化。则
$B \subset B'$ 是离散赋值环扩张，其分式域扩张可分代数。
由 Algebra，引理
\ref{algebra-lemma-strictly-henselian-functorial-prepare}，
存在引理陈述中的交换图。引理的 (1)、(2) 是清楚的。

\medskip\noindent
设 $K'_1/K'$ 是 $A' \to B'$ 的弱解或解。由于 $A'$ 是余极限，
可找到有限 \'etale 扩张 $A \subset A_1'$，以及 $A_1'$
的分式域 $F$ 的有限扩张 $K_1$，使得
$K'_1 = K' \otimes_F K_1$。由于 $A \subset A_1'$ 有限
\'etale 且 $B'$ 严格 Hensel，$B' \otimes_A A_1'$ 是与
$B'$ 同构的环的有限乘积。因此
% P02-E0645
$$
L' \otimes_K K_1 = L' \otimes_K F \otimes_F K_1
$$
是与 $L' \otimes_{K'} K'_1$ 同构的环的有限乘积。
由此可见，$K_1/K$ 是 $A \to B'$ 的弱解或解。
故由引理 \ref{more-algebra-lemma-solutions-go-down}，它也是
$A \to B$ 的弱解或解。
% P02-E0647
\end{proof}

\begin{lemma}
\label{more-algebra-lemma-galois-relative}
设 $A \to B$ 是分式域为 $K \subset L$ 的离散赋值环扩张。设
$K_1/K$ 是正规扩张，并令 $G = \text{Aut}(K_1/K)$。则 $G$
作用在评注 \ref{more-algebra-remark-construction} 的环 $K_1$、$L_1$、$A_1$
与 $B_1$ 上，并且在 $B_1$ 的极大理想集合上传递作用。
% P02-E0648
\end{lemma}

\begin{proof}
除最后一个断言外，一切都是清楚的。若该作用有两个或更多轨道，则可找到
$b \in B_1$，它在一个轨道的所有极大理想处为零，而在另一个轨道的所有
极大理想处的剩余类都是 $1$。于是
$b' = \prod_{\sigma \in G} \sigma(b)$ 是
$B_1 \subset L_1 = (L \otimes_K K_1)_{red}$ 的一个 $G$-不变元素；
它属于 $B_1$ 的某些极大理想，却不属于 $B_1$ 的所有极大理想。将它提升为
$L \otimes_K K_1$ 的一个元素并取充分高次幂，可得到
$L \otimes_K K_1$ 的一个 $G$-不变元素 $b''$，它对某个 $N > 0$
映到 $(b')^N$；事实上，只有在特征为 $p > 0$ 时才需要这样做，而在这种情形，
取充分大的 $p$ 次幂 $q$ 给出典范映射
$(L \otimes_K K_1)_{red} \to L \otimes_K K_1$。
由于 $K = (K_1)^G$，可得 $b'' \in L$。又因为 $b''$ 映到 $B_1$
的一个元素，故 $b'' \in B$（因为 $B$ 正规）。一方面，由于 $b'$ 属于
$B_1$ 的某个极大理想，必有 $b'' \in \mathfrak m_B$；另一方面，由于
$b'$ 不属于 $B_1$ 的所有极大理想，又必有
$b'' \not \in \mathfrak m_B$。这一矛盾完成了引理的证明。
\end{proof}

\begin{lemma}
\label{more-algebra-lemma-make-degree-q-extension}
设 $A$ 是以 $\pi$ 为一致化元的离散赋值环。若 $A$ 的剩余特征为
$p > 0$，则对每个 $n > 1$ 和 $p$ 的幂 $q$，都存在关于 $A$
完全分歧的次数为 $q$ 的可分扩张 $L/K$，使得 $A$ 在 $L$ 中的整闭包
$B$ 的分歧指数为 $q$，并且有一致化元 $\pi_B$，满足
$\pi_B^q = \pi + \pi^n b$ 和
$\pi_B^q = \pi + (\pi_B)^{nq}b'$，其中某些 $b, b' \in B$。
\end{lemma}

\begin{proof}
若 $K$ 的特征为零，则可取由 $\pi_B^q = \pi$ 给出的扩张，见引理
\ref{more-algebra-lemma-pull-root-uniformizer}。若 $K$ 的特征为 $p > 0$，则可取由
$z^q - \pi^n z = \pi^{1 - q}$ 给出的 $K$ 的扩张。事实上，令
$y = \pi z$，便有 $y^q - \pi^{n + q - 1} y = \pi$。取
$\pi_B = y$ 即得所求结论。
\end{proof}

\begin{lemma}
\label{more-algebra-lemma-pre-purely-inseparable-case}
% P02-E0649
设 $A$ 是离散赋值环。假设剩余域 $\kappa_A$ 的特征为 $p > 0$，并且
$a \in A$ 的剩余类在 $\kappa_A$ 中不是 $p$ 次幂。则 $a$ 在 $K$
中也不是 $p$ 次幂，而且 $A$ 在 $K[a^{1/p}]$ 中的整闭包是环
$A[a^{1/p}]$；这是 $A$ 上弱非分歧的离散赋值环。
\end{lemma}

\begin{proof}
本引理不证自明。
\end{proof}

\begin{lemma}
\label{more-algebra-lemma-purely-inseparable-case}
% P02-E0650
设 $A \subset B \subset C$ 是分式域为 $K \subset L \subset M$
的离散赋值环扩张，并设 $\pi \in A$ 是一致化元。假设
\begin{enumerate}
\item $B$ 是 Nagata 环；
\item $A \subset B$ 弱非分歧；
\item $M$ 是 $L$ 的次数为 $p$ 的纯不可分扩张。
\end{enumerate}
则下列情形之一成立：
\begin{enumerate}
\item $A \to C$ 弱非分歧；
\item $C = B[\pi^{1/p}]$；
\item 存在关于 $A$ 完全分歧的次数为 $p$ 的可分扩张 $K_1/K$，使得
$L_1 = L \otimes_K K_1$ 与 $M_1 = M \otimes_K K_1$ 都是域，
并且整闭包之间的映射 $A_1 \to B_1 \to C_1$ 都是离散赋值环的
弱非分歧扩张。
\end{enumerate}
\end{lemma}

\begin{proof}
设 $e$ 是 $C$ 在 $B$ 上的分歧指数。若 $e = 1$，则结论成立。
否则，由引理 \ref{more-algebra-lemma-inequality} 与
\ref{more-algebra-lemma-ramification-index-a-power-of-p} 可知 $e = p$。
这又蕴含 $B$ 与 $C$ 的剩余域相同。选取 $C$ 的一致化元 $\pi_C$，并写成
$\pi_C^p = u \pi$，其中 $u$ 是 $C$ 的某个单位。由于
$\pi_C^p \in L$，可知 $u \in B^*$。此外，$M = L[\pi_C]$。

\medskip\noindent
假设存在整数 $m \geq 0$，使得
$$
u = \sum\nolimits_{0 \leq i < m} b_i^p \pi^i + b \pi^m
$$
，其中 $b_i \in B$，且 $b \in B$ 在 $\kappa_B$ 中的像不是
$p$ 次幂。按引理 \ref{more-algebra-lemma-make-degree-q-extension} 选取扩张
$K_1/K$，其中取 $n = m + 2$；并把 $A$ 在 $K_1$ 中的整闭包
$A_1$ 的一致化元记作 $\pi'$，使得对某个 $a \in A_1$ 有
$\pi = (\pi')^p + (\pi')^{np} a$。设 $B_1$ 是 $B$ 在
$L \otimes_K K_1$ 中的整闭包。由引理
\ref{more-algebra-lemma-weakly-unramified-goes-up-along-totally-ramified}，
$A_1 \to B_1$ 弱非分歧。在 $B_1$ 中有
$$
u \pi =
\left(\sum\nolimits_{0 \leq i < m} b_i (\pi')^{i + 1}\right)^p +
b (\pi')^{(m + 1)p} + (\pi')^{np} b_1
$$
，其中某个 $b_1 \in B_1$（计算从略）。由此可知，$M_1$ 是由在
$L_1$ 上添加下式的一个 $p$ 次根而得到的：
$$
b + (\pi')^{n - m - 1} b_1
$$
。由于 $B_1$ 的剩余域等于 $B$ 的剩余域，由引理
\ref{more-algebra-lemma-pre-purely-inseparable-case} 可知 $M_1/L_1$ 的次数为
$p$，并且 $B_1$ 的整闭包 $C_1$ 在 $B_1$ 上弱非分歧。因此在此情形
得出结论。

\medskip\noindent
% P02-E0651
若不存在前一段中的整数 $m$，则 $u$ 在 $B_1$ 的 $\pi$-进完备化中
是 $p$ 次幂。由于 $B$ 是 Nagata 环，由 Algebra，引理
\ref{algebra-lemma-nagata-pth-roots}，这意味着 $u$ 在 $B_1$ 中是
$p$ 次幂。因此，引理陈述中的第二种情形成立。
\end{proof}

\begin{lemma}
\label{more-algebra-lemma-cohen}
设 $A$ 是由素数 $p$ 零化的局部环，且其极大理想是幂零的。存在环映射
$\sigma : \kappa_A \to A$，它是剩余映射 $A \to \kappa_A$ 的截面。
若 $A \to A'$ 是局部环之间的局部同态，则只要扩张
$\kappa_{A'}/\kappa_A$ 可分，就可以选取与 $\sigma$ 相容的类似环映射
$\sigma' : \kappa_{A'} \to A'$。
\end{lemma}

\begin{proof}
由 Algebra，命题
\ref{algebra-proposition-characterize-separable-field-extensions}，
可分扩张形式光滑。因此，由 $\mathbf{F}_p \to \kappa_A$ 可分可得
$\sigma$ 的存在性。与 $\sigma$ 相容的 $\sigma'$ 的存在性同理。
\end{proof}

\begin{lemma}
\label{more-algebra-lemma-pre-characteristic-p-case}
设 $A$ 是分式域 $K$ 的特征为 $p > 0$ 的离散赋值环。设
$\xi \in K$，并设 $L$ 是在 $K$ 上添加 $z^p - z = \xi$ 的一个根
而得到的扩张。则 $L/K$ 是 Galois 扩张，并且下列情形之一成立：
\begin{enumerate}
\item $L = K$；
\item $L/K$ 关于 $A$ 非分歧，次数为 $p$；
\item $L/K$ 关于 $A$ 完全分歧，分歧指数为 $p$；
\item $A$ 在 $L$ 中的整闭包 $B$ 是离散赋值环，$A \subset B$
弱非分歧，而且 $A \to B$ 诱导次数为 $p$ 的纯不可分剩余域扩张。
\end{enumerate}
设 $\pi$ 是 $A$ 的一致化元。则有以下蕴含：
\begin{enumerate}
\item[(A)] 若 $\xi \in A$，则处于情形 (1) 或 (2)。
\item[(B)] 若 $\xi = \pi^{-n}a$，其中 $n > 0$ 不被 $p$ 整除，且
$a$ 是 $A$ 中的单位，则处于情形 (3)。
\item[(C)] 若 $\xi = \pi^{-n} a$，其中 $n > 0$ 被 $p$ 整除，且
$a$ 在 $\kappa_A$ 中的像不是 $p$ 次幂，则处于情形 (4)。
\end{enumerate}
\end{lemma}

\begin{proof}
由 Fields，章节 \ref{fields-section-Artin-Schreier} 中的讨论，
该扩张是阶数整除 $p$ 的 Galois 扩张。由章节
\ref{more-algebra-section-ramification} 中的讨论立即可知，我们处于引理所列的
(1)--(4) 中的某种情形。

\medskip\noindent
情形 (A)。此时 $A \to A[x]/(x^p - x - \xi)$ 是有限
\'etale 环扩张，故我们处于情形 (1) 或 (2)。

\medskip\noindent
情形 (B)。写成 $\xi = \pi^{-n}a$，其中 $p$ 不整除 $n$。设
$B \subset L$ 是 $A$ 在 $L$ 中的整闭包。若对某个极大理想
$\mathfrak m$ 有 $C = B_\mathfrak m$，则显然
$p \text{ord}_C(z) = -n \text{ord}_C(\pi)$。特别地，
$A \subset C$ 的分歧指数被 $p$ 整除。于是该指数等于 $p$，且
$B = C$。

\medskip\noindent
情形 (C)。令 $k = n/p$。于是可将方程改写为
$$
(\pi^kz)^p - \pi^{n - k} (\pi^kz) = a
$$
。由于 $A[y]/(y^p - \pi^{n - k}y - a)$ 是 $A$ 上弱非分歧的
离散赋值环，引理得证。
\end{proof}

\begin{lemma}
\label{more-algebra-lemma-characteristic-p-case}
设 $A \subset B \subset C$ 是分式域为 $K \subset L \subset M$
的离散赋值环扩张。假设
\begin{enumerate}
\item $A \subset B$ 弱非分歧；
\item $K$ 的特征为 $p$；
\item $M$ 是 $L$ 的次数为 $p$ 的 Galois 扩张；并且
\item $\kappa_A = \bigcap_{n \geq 1} \kappa_B^{p^n}$。
\end{enumerate}
则存在关于 $A$ 完全分歧的有限 Galois 扩张 $K_1/K$，它是
$A \to C$ 的弱解。
\end{lemma}

\begin{proof}
由于 $L$ 的特征为 $p$，我们知道 $M$ 是 $L$ 的 Artin--Schreier
扩张（Fields，引理 \ref{fields-lemma-Artin-Schreier}）。因此可选取
$z \in M$、$z \not \in L$，使得 $\xi = z^p - z \in L$。选取
$n \geq 0$，使得 $\pi^n\xi \in B$。在所有这些选取中，选取使
$n$ 最小的 $z$。若 $n = 0$，则 $M/L$ 关于 $B$ 非分歧
（引理 \ref{more-algebra-lemma-pre-characteristic-p-case}），结论成立。因此
$n > 0$。

\medskip\noindent
假设 (4) 蕴含 $\kappa_A$ 是完美域。因而可按引理
\ref{more-algebra-lemma-cohen} 选取相容的环映射
% P02-E0652
$\overline{\sigma} : \kappa_A \to A/\pi^n A$ 与
$\overline{\sigma} : \kappa_B \to B/\pi^n B$。把其中第二个映射提升为
集合映射 $\sigma : \kappa_B \to B$\footnote{若 $B$ 完备，则可选取
$\sigma$ 为环映射。若 $A$ 也完备且 $\sigma$ 是环映射，则
$\sigma$ 把 $\kappa_A$ 映入 $A$。}。于是可写成
$$
\xi = \sum\nolimits_{i = n, \ldots, 1} \sigma(\lambda_i) \pi^{-i} + b
$$
，其中某些 $\lambda_i \in \kappa_B$ 与 $b \in B$。令
$$
I = \{i \in \{n, \ldots, 1\} \mid \lambda_i \in \kappa_A\}
$$
以及
% P02-E0653
$$
J = \{j \in \{n, \ldots, 1\} \mid \lambda_i \not \in \kappa_A\}
$$
。我们将对有限集 $J$ 的基数作归纳。

\medskip\noindent
情形 $J = \emptyset$。此时，对每个 $i \in \{n, \ldots, 1\}$，
由 $\sigma$ 的选取，可写
$\sigma(\lambda_i) = a_i + \pi^n b_i$，其中某些 $a_i \in A$ 与
$b_i \in B$。因此，对某些 $a \in A$ 与 $b \in B$ 有
$\xi = \pi^{-n} a + b$。若 $p | n$，则可写
$a = a_0^p + \pi a_1$，其中某些 $a_0, a_1 \in A$
（因为 $A$ 的剩余域完美）。计算得
$$
(z - \pi^{-n/p}a_0)^p - (z - \pi^{-n/p}a_0) =
\pi^{-(n - 1)}(a_1 + \pi^{n - 1 - n/p}a_0) + b'
$$
，其中某个 $b' \in B$。这与 $n$ 的最小性矛盾。因此 $p$ 不整除
$n$。考虑由 $w^p - w = \pi^{-n}a$ 给出的 $K$ 的次数为 $p$
的扩张 $K_1$。由引理 \ref{more-algebra-lemma-pre-characteristic-p-case}，该扩张是
Galois 扩张，且关于 $A$ 完全分歧。因此
$L_1 = L \otimes_K K_1$ 是域，而且 $A_1 \subset B_1$ 弱非分歧
（引理 \ref{more-algebra-lemma-weakly-unramified-goes-up-along-totally-ramified}）。
由引理 \ref{more-algebra-lemma-pre-characteristic-p-case}，环
$M_1 = M \otimes_K K_1$ 或者是 $L_1$ 的 $p$ 个副本的乘积
（此时结论成立），或者是 $L_1$ 的次数为 $p$ 的域扩张。此外，在第二种
情形中，或者 $C_1$ 在 $B_1$ 上弱非分歧（此时结论成立），或者
$M_1/L_1$ 的次数为 $p$、是 Galois 扩张，并且关于 $B_1$ 完全分歧。
在最后这种情形中，扩张 $M_1/L_1$ 由元素 $z - w$ 生成，并且
$$
(z - w)^p - (z - w) = z^p - z - (w^p - w) = b
$$
，其中 $b \in B$（见上文）。再用一次引理
\ref{more-algebra-lemma-pre-characteristic-p-case}，扩张 $M_1/L_1$ 关于 $B_1$
非分歧，故可断定 $K_1$ 是 $A \to C$ 的弱解。从现在起假设
$J \not = \emptyset$。

\medskip\noindent
假设存在 $j', j \in J$ 与某个 $r > 0$，使得 $j' = p^r j$。
把所选的 $z$ 改为
$$
z' = z -
(\sigma(\lambda_j) \pi^{-j} + \sigma(\lambda_j^p) \pi^{-pj} + \ldots +
\sigma(\lambda_j^{p^{r - 1}}) \pi^{-p^{r - 1}j})
$$
。于是 $\xi$ 按下式变为 $\xi' = (z')^p - (z')$：
$$
\xi' =
\xi - \sigma(\lambda_j) \pi^{-j} + \sigma(\lambda_j^{p^r}) \pi^{-j'}
+ \text{something in }B
$$
。像先前一样写成
$\xi' = \sum\nolimits_{i = n, \ldots, 1} \sigma(\lambda'_i) \pi^{-i} + b'$，
便得到：当 $i \not = j, j'$ 时有 $\lambda'_i = \lambda_i$，并且
$\lambda'_j = 0$。因此集合 $J$ 变小了。由对 $J$ 的基数所作归纳，
可以假设不存在这样的指标对 $j, j'$。（请注意，在此过程中可能回到
上面已经处理过的 $J = \emptyset$ 情形。）

\medskip\noindent
对 $j \in J$，写
$\lambda_j = \mu_j^{p^{r_j}}$，其中某个 $r_j \geq 0$，而且
$\mu_j \in \kappa_B$ 不是 $p$ 次幂。由假设 (4)，这是可能的。
令 $j \in J$ 是使 $j p^{-r_j}$ 最大的唯一指标（由前一段的结论，
该指标唯一）。选取
% P02-E0654
$r > \max(r_j + 1)$，并使得对 $j \in J$ 有
$j p^{r - r_j} > n$。选取关于 $A$ 完全分歧、次数为 $p^r$ 的可分扩张
$K_1/K$，使得相应离散赋值环 $A_1 \subset K_1$ 有一致化元 $\pi'$，
且对某个 $a \in A_1$ 有
$(\pi')^{p^r} = \pi + \pi^{n + 1}a$
（引理 \ref{more-algebra-lemma-make-degree-q-extension}）。注意，
$L_1 = L \otimes_K K_1$ 是域，而且 $L_1/L$ 关于 $B$ 完全分歧
（引理 \ref{more-algebra-lemma-weakly-unramified-goes-up-along-totally-ramified}）。
在整闭包 $B_1$ 中计算得
$$
\xi = \sum\nolimits_{i \in I} \sigma(\lambda_i) (\pi')^{-i p^r} +
\sum\nolimits_{j \in J} \sigma(\mu_j)^{p^{r_j}} (\pi')^{-j p^r} + b_1
$$
，其中某个 $b_1 \in B_1$。注意，对 $i \in I$，
% P02-E0655
$\sigma(\lambda_i)$ 模 $\pi^n$ 是 $q$ 次幂，换言之模
$(\pi')^{n p^r}$ 也有上述性质。因此可将上式改写为
$$
\xi = \sum\nolimits_{i \in I} x_i^{p^r} (\pi')^{-i p^r} +
\sum\nolimits_{j \in J} \sigma(\mu_j)^{p^{r_j}} (\pi')^{- j p^r}
+ b_1
$$
。如前一段那样，把所选的 $z$ 改为
% P02-E0656
\begin{align*}
z' & = z \\
& -
\sum\nolimits_{i \in I}
\left(x_i (\pi')^{-i} + \ldots + x_i^{p^{r - 1}} (\pi')^{-i p^{r - 1}}\right)
\\
& -
\sum\nolimits_{j \in J}
\left(
\sigma(\mu_j) (\pi')^{- j p^{r - r_j}}
+ \ldots +
\sigma(\mu_j)^{p^{r_j - 1}} (\pi')^{- j p^{r - 1}}
\right)
\end{align*}
，便得到
$$
(z')^p - z' =
\sum\nolimits_{i \in I} x_i (\pi')^{-i} +
\sum\nolimits_{j \in J} \sigma(\mu_j) (\pi')^{- j p^{r - r_j}} + b_1'
$$
，其中某个 $b'_1 \in B_1$。由于恰有一个 $j$ 使得
$j p^{r - r_j}$ 最大，而且 $j p^{r - r_j}$ 大于每个 $i \in I$
并被 $p$ 整除，可知 $M_1 / L_1$ 属于引理
\ref{more-algebra-lemma-pre-characteristic-p-case} 的情形 (C)。证明完成。
\end{proof}

\begin{lemma}
\label{more-algebra-lemma-prepare}
设 $A$ 是包含本原 $p$ 次单位根 $\zeta$ 的环。令 $w = 1 - \zeta$。则
$$
P(z) = \frac{(1 + wz)^p - 1}{w^p} =
z^p - z + \sum\nolimits_{0 < i < p} a_i z^i
$$
是 $A[z]$ 的元素，并且事实上 $a_i \in (w)$。此外，在多项式环
$A[z_1, z_2]$ 中有
$$
P(z_1 + z_2 + w z_1 z_2) = P(z_1) + P(z_2) + w^p P(z_1) P(z_2)
$$
。
\end{lemma}

\begin{proof}
只需在
$$
A = \mathbf{Z}[\zeta] = \mathbf{Z}[x]/(x^{p - 1} + \ldots + x + 1)
$$
是圆分域的整数环时证明。多项式恒等式
$t^p - 1 = (t - 1)(t - \zeta) \ldots (t - \zeta^{p - 1})$
（比较两边的根即可证明）表明
$t^{p - 1} + \ldots + t + 1 = (t - \zeta) \ldots (t - \zeta^{p - 1})$。
代入 $t = 1$，得到
$p = (1 - \zeta)(1 - \zeta^2) \ldots (1 - \zeta^{p - 1})$。
极大理想 $(p, w) = (w)$ 是 $A$ 中位于 $p$ 之上的唯一素理想
（因为特征为 $p$ 的域没有非平凡的 $1$ 的 $p$ 次根）。因此对某个单位
$u$ 有 $p = u w^{p - 1}$。这蕴含
$$
a_i = \frac{1}{p} {p \choose i} u w^{i - 1}
$$
对 $p > i > 1$ 成立，而且 $- 1 + a_1 = pw/w^p = u$。由于
$P(-1) = 0$，可知模 $(w)$ 有 $0 = (-1)^p - u$。因此
$a_1 \in (w)$，第一部分的证明完成。第二部分由从略的直接计算得到。
% P02-E0657
\end{proof}

\begin{lemma}
% P02-E0658
\label{more-algebra-lemma-extension-defined-by-nice-polynial}
设 $A$ 是包含 $1$ 的一个本原 $p$ 次根的混合特征 $(0, p)$ 离散赋值环。
设 $P(t) \in A[t]$ 是引理 \ref{more-algebra-lemma-prepare} 的多项式。设
$\xi \in K$，并设 $L$ 是在 $K$ 上添加 $P(z) = \xi$ 的一个根
而得到的扩张。则 $L/K$ 是 Galois 扩张，并且下列情形之一成立：
\begin{enumerate}
\item $L = K$；
\item $L/K$ 关于 $A$ 非分歧，次数为 $p$；
\item $L/K$ 关于 $A$ 完全分歧，分歧指数为 $p$；
\item $A$ 在 $L$ 中的整闭包 $B$ 是离散赋值环，$A \subset B$
弱非分歧，而且 $A \to B$ 诱导次数为 $p$ 的纯不可分剩余域扩张。
\end{enumerate}
设 $\pi$ 是 $A$ 的一致化元。则有以下蕴含：
\begin{enumerate}
\item[(A)] 若 $\xi \in A$，则处于情形 (1) 或 (2)。
\item[(B)] 若 $\xi = \pi^{-n}a$，其中 $n > 0$ 不被 $p$ 整除，且
$a$ 是 $A$ 中的单位，则处于情形 (3)。
\item[(C)] 若 $\xi = \pi^{-n} a$，其中 $n > 0$ 被 $p$ 整除，且
$a$ 在 $\kappa_A$ 中的像不是 $p$ 次幂，则处于情形 (4)。
\end{enumerate}
\end{lemma}

\begin{proof}
% P02-E0659
添加 $P(z) = \xi$ 的一个根等价于添加
$y^p = w^p(1 + \xi)$ 的一个根。由于 $K$ 包含 $1$ 的一个本原
$p$ 次根，由 Fields，章节 \ref{fields-section-Kummer} 中的讨论，
该扩张是阶数整除 $p$ 的 Galois 扩张。由章节
\ref{more-algebra-section-ramification} 中的讨论立即可知，我们处于引理所列的
(1)--(4) 中的某种情形。

\medskip\noindent
情形 (A)。此时 $A \to A[x]/(P(x) - \xi)$ 是有限
\'etale 环扩张，故我们处于情形 (1) 或 (2)。

\medskip\noindent
情形 (B)。写成 $\xi = \pi^{-n}a$，其中 $p$ 不整除 $n$。设
$B \subset L$ 是 $A$ 在 $L$ 中的整闭包。若对某个极大理想
$\mathfrak m$ 有 $C = B_\mathfrak m$，则显然
$p \text{ord}_C(z) = -n \text{ord}_C(\pi)$。特别地，
$A \subset C$ 的分歧指数被 $p$ 整除。于是该指数等于 $p$，且
$B = C$。

\medskip\noindent
情形 (C)。令 $k = n/p$。于是可将方程改写为
$$
(\pi^kz)^p - \pi^{n - k} (\pi^kz) + \sum a_i \pi^{n - ik} (\pi^kz)^i = a
$$
。由于
% P02-E0660
$A[y]/(y^p - \pi^{n - k}y - \sum  a_i \pi^{n - ik} y^i - a)$
是 $A$ 上弱非分歧的离散赋值环，引理得证。
\end{proof}

\noindent
设 $A$ 是包含 $1$ 的一个本原 $p$ 次根的混合特征 $(0, p)$
离散赋值环。设 $w \in A$ 与 $P(t) \in A[t]$ 如引理
\ref{more-algebra-lemma-prepare} 中所述，并设 $L$ 是 $K$ 的有限扩张。若 $L$
是在 $K$ 上添加方程 $P(z) = \xi$ 的一个根而得到的次数为 $p$
的扩张，其中 $\xi \in K$ 且 $w^p \xi \in \mathfrak m_A$，
则称 $L/K$ 是一个\textit{有限层级的次数为 $p$ 的扩张}。

\medskip\noindent
由于下面这个直接的引理，这一定义与本节的讨论有关。

\begin{lemma}
\label{more-algebra-lemma-make-finite-level}
设 $A \subset B \subset C$ 是分式域为 $K \subset L \subset M$
的离散赋值环扩张。假设
\begin{enumerate}
\item $A$ 具有混合特征 $(0, p)$；
\item $A \subset B$ 弱非分歧；
\item $B$ 包含 $1$ 的一个本原 $p$ 次根；并且
\item $M/L$ 是次数为 $p$ 的 Galois 扩张。
\end{enumerate}
则存在关于 $A$ 完全分歧的有限 Galois 扩张 $K_1/K$，它或者是
$A \to C$ 的弱解，或者使得 $M_1/L_1$ 是有限层级的次数为 $p$
的扩张。
\end{lemma}

\begin{proof}
设 $\pi \in A$ 是一致化元。由 Kummer 理论（Fields，引理
\ref{fields-lemma-Kummer}），$M$ 是从 $L$ 添加某个 $b \in L$
的方程 $y^p = b$ 的根而得到的。

\medskip\noindent
若 $\text{ord}_B(b)$ 与 $p$ 互素，则选取关于 $A$ 完全分歧的
次数为 $p$ 的可分扩张 $K_1/K$（例如使用引理
\ref{more-algebra-lemma-make-degree-q-extension}）。设 $A_1$ 是 $A$ 在 $K_1$
中的整闭包。由引理
\ref{more-algebra-lemma-weakly-unramified-goes-up-along-totally-ramified}，
$B$ 在 $L_1 = L \otimes_K K_1$ 中的整闭包 $B_1$ 是 $A_1$ 上
弱非分歧的离散赋值环。若 $K_1/K$ 不是 $A \to C$ 的弱解，则
$C$ 在 $M_1 = M \otimes_K K_1$ 中的整闭包 $C_1$ 是离散赋值环，
而 $B_1 \to C_1$ 的分歧指数为 $p$。在这种情形中，域 $M_1$
由在 $L_1$ 上添加 $b$ 的 $p$ 次根而得到，而
$\text{ord}_{B_1}(b)$ 被 $p$ 整除。用 $A_1$ 替换 $A$ 等对象，
便可假设 $b = \pi^n u$，其中 $u \in B$ 是单位，且 $n$ 被 $p$
整除。当然，在此情形，扩张 $M$ 是由在 $L$ 上添加某个单位的
$p$ 次根而得到的。

\medskip\noindent
假设 $M$ 是从 $L$ 添加方程 $y^p = u$ 的根而得到的，其中 $u$
是 $B$ 的某个单位。若 $u$ 在 $\kappa_B$ 中的剩余类不是 $p$ 次幂，
则 $B \subset C$ 弱非分歧（引理
\ref{more-algebra-lemma-pre-purely-inseparable-case}），结论成立。否则，可用
$y/v$ 替换所选的 $y$，其中 $v^p$ 与 $u$ 在 $\kappa_B$ 中的像相同。
作此替换后有
$$
y^p = 1 + \pi b
$$
，其中某个 $b \in B$。于是由 $z = (y - 1)/w$ 可知
$P(z) = \pi b/ w^p$。因此可见，该扩张是有限层级的次数为 $p$
的扩张，其中 $\xi = \pi b / w^p$。
\end{proof}

\noindent
设 $A$ 是包含 $1$ 的一个本原 $p$ 次根的混合特征 $(0, p)$
离散赋值环。设 $w \in A$ 与 $P(t) \in A[t]$ 如引理
\ref{more-algebra-lemma-prepare} 中所述。设 $L$ 是 $K$ 的有限层级的次数为 $p$
的扩张。选取在 $K$ 上生成 $L$ 的 $z \in L$，使得
$\xi = P(z) \in K$。选取 $A$ 的一致化元 $\pi$，并写成
$w = u \pi^{e_1}$，其中整数 $e_1 = \text{ord}_A(w)$，而 $u \in A$
是单位。最后，选取 $n \geq 0$，使得
$$
\pi^n \xi \in A
$$
。当 $z$ 取遍所有满足 $\xi = P(z) \in K$ 且生成扩张 $L/K$
的元素时，量 $n/e_1$ 所取的最小值称为 $L/K$ 的\textit{层级}。
% P02-E0661

\medskip\noindent
作几点说明。由于该扩张具有有限层级，我们知道可选取 $z$ 使得
$n < pe_1$。因此层级是包含在 $[0, p)$ 中的有理数。若层级为零，
则由引理 \ref{more-algebra-lemma-extension-defined-by-nice-polynial}，$L/K$ 关于
$A$ 非分歧。我们的下一个目标是降低层级。

\begin{lemma}
\label{more-algebra-lemma-lowering-the-level}
设 $A \subset B \subset C$ 是分式域为 $K \subset L \subset M$
的离散赋值环扩张。假设
\begin{enumerate}
\item $A$ 具有混合特征 $(0, p)$；
\item $A \subset B$ 弱非分歧；
\item $B$ 包含 $1$ 的一个本原 $p$ 次根；
\item $M/L$ 是有限层级 $l > 0$ 的次数为 $p$ 的扩张；
\item $\kappa_A = \bigcap_{n \geq 1} \kappa_B^{p^n}$。
\end{enumerate}
则存在 $K$ 的关于 $A$ 完全分歧的有限可分扩张 $K_1$，使得
$K_1$ 或者是 $A \to C$ 的弱解，或者扩张 $M_1/L_1$ 是层级
$\leq \max(0, l - 1, 2l - p)$ 的次数为 $p$ 的有限层级扩张。
\end{lemma}

\begin{proof}
设 $\pi \in A$ 是一致化元。设 $w \in B$ 与 $P \in B[t]$ 如引理
\ref{more-algebra-lemma-prepare} 中所述（对 $B$ 而言）。令
$e_1 = \text{ord}_B(w)$，于是 $w$ 与 $\pi^{e_1}$ 在 $B$ 中相伴。
选取在 $L$ 上生成 $M$ 的 $z \in M$，使得 $\xi = P(z) \in K$，
并按照 $M$ 在 $L$ 上的层级定义选取 $n$，使得
$\pi^n\xi \in B$；即 $l = n/e_1$。

\medskip\noindent
本引理的证明与引理 \ref{more-algebra-lemma-characteristic-p-case} 的证明完全类似。
为说明其中的情形，注意到，对任意满足
% P02-E0663
$\pi^{-n} P(z) \in B$ 的 $z \in L$（利用 $z$ 的赋值最坏为
$-n/p$ 以及多项式 $P$ 的形状），都有
\begin{equation}
\label{more-algebra-equation-first-congruence}
P(z) \equiv z^p - z \bmod \pi^{-n + e_1}B
\end{equation}
。此外，对 $\xi_1, \xi_2 \in \pi^{-n}B$ 有
\begin{equation}
\label{more-algebra-equation-second-congruence}
\xi_1 + \xi_2 + w^p \xi_1 \xi_2 \equiv \xi_1 + \xi_2 \bmod \pi^{-2n + pe_1}B 
\end{equation}
。最后，注意
% P02-E0664
$n - e_1 = (l - 1)/e_1$ 与
$-2n + pe_1 = -(2l - p)e_1$。写
$m = n - e_1 \max(0, l - 1, 2l - p)$。上述事实表明，在
$\pi^{-n}B / \pi^{-n + m}B$ 中计算时，多项式 $P$ 的表现与多项式
$z^p -  z$ 完全相同。这解释了引理为何成立，不过下面仍给出细节。

\medskip\noindent
% P02-E0662
假设 (4) 蕴含 $\kappa_A$ 是完美域。注意 $m \leq e_1$，所以
$A/\pi^m$ 被 $w$ 零化，从而被 $p$ 零化。因此可按引理
\ref{more-algebra-lemma-cohen} 选取相容的环映射
$\overline{\sigma} : \kappa_A \to A/\pi^mA$ 与
$\overline{\sigma} : \kappa_B \to B/\pi^mB$。把其中第二个映射提升为
集合映射 $\sigma : \kappa_B \to B$。于是可写成
% P02-E0665
$$
\xi =
\sum\nolimits_{i = n, \ldots, n - m + 1} \sigma(\lambda_i) \pi^{-i} +
\pi^{-n + m)} b
$$
，其中某些 $\lambda_i \in \kappa_B$ 与 $b \in B$。令
$$
I = \{i \in \{n, \ldots, n - m + 1\} \mid \lambda_i \in \kappa_A\}
$$
以及
% P02-E0666
$$
J = \{j \in \{n, \ldots, n - m + 1\} \mid \lambda_i \not \in \kappa_A\}
$$
。我们将对有限集 $J$ 的基数作归纳。

\medskip\noindent
情形 $J = \emptyset$。此时，对每个
$i \in \{n, \ldots, n - m + 1\}$，由 $\overline{\sigma}$ 的选取，
% P02-E0667
可写 $\sigma(\lambda_i) = a_i + \pi^{n - m}b_i$，其中某些
$a_i \in A$ 与 $b_i \in B$。因此，对某些 $a \in A$ 与 $b \in B$
有 $\xi = \pi^{-n} a + \pi^{-n + m} b$。若 $p | n$，则可写
$a = a_0^p + \pi a_1$，其中某些 $a_0, a_1 \in A$
（因为 $A$ 的剩余域完美）。令 $z_1 = - \pi^{-n/p} a_0$。注意
$P(z_1) \in \pi^{-n}B$，而且 $z + z_1 + w z z_1$ 是在 $L$ 上
生成 $M$ 的元素（注意，由于 $n < pe_1$，有
$wz_1 \not = -1$）。此外，由引理 \ref{more-algebra-lemma-prepare} 有
$$
P(z + z_1 + w z z_1) = P(z) + P(z_1) + w^p P(z) P(z_1) \in K
$$
，而由方程 (\ref{more-algebra-equation-first-congruence}) 与
(\ref{more-algebra-equation-second-congruence}) 有
$$
P(z) + P(z_1) + w^p P(z) P(z_1)
\equiv
\xi + z_1^p - z_1 \bmod \pi^{-n + m}B
$$
。这里某个 $b' \in B$。
% P02-E0668
这与 $n$ 的最小性矛盾！因此 $p$ 不整除 $n$。考虑由
$P(y) = -\pi^{-n}a$ 给出的 $K$ 的次数为 $p$ 的扩张 $K_1$。
由引理 \ref{more-algebra-lemma-extension-defined-by-nice-polynial}，该扩张可分，
且关于 $A$ 完全分歧。因此 $L_1 = L \otimes_K K_1$ 是域，且
$A_1 \subset B_1$ 弱非分歧
（引理 \ref{more-algebra-lemma-weakly-unramified-goes-up-along-totally-ramified}）。
由引理 \ref{more-algebra-lemma-extension-defined-by-nice-polynial}，环
$M_1 = M \otimes_K K_1$ 或者是 $L_1$ 的 $p$ 个副本的乘积
（此时结论成立），或者是 $L_1$ 的次数为 $p$ 的域扩张。此外，在第二种
情形中，或者 $C_1$ 在 $B_1$ 上弱非分歧（此时结论成立），或者
$M_1/L_1$ 的次数为 $p$、是 Galois 扩张并且关于 $B_1$ 完全分歧。
在最后一种情形，扩张 $M_1/L_1$ 由元素 $z + y + wzy$ 生成，
而且 $P(z + y + wzy) \in L_1$，并且
\begin{align*}
P(z + y + wzy)
& = P(z) + P(y) + w^p P(z) P(y) \\
& \equiv
\xi - \pi^{-n}a \bmod \pi^{-n + m}B_1 \\
& \equiv
0 \bmod \pi^{-n + m}B_1
\end{align*}
，方式与上文完全相同。由 $m$ 的选取，这恰好意味着 $M_1/L_1$
的层级至多为 $\max(0, l - 1, 2l - p)$。从现在起假设
$J \not = \emptyset$。

\medskip\noindent
假设存在 $j', j \in J$ 与某个 $r > 0$，使得 $j' = p^r j$。
令
$$
z_1 = - \sigma(\lambda_j) \pi^{-j} - \sigma(\lambda_j^p) \pi^{-pj} -
\ldots - \sigma(\lambda_j^{p^{r - 1}}) \pi^{-p^{r - 1}j}
$$
，并把 $z$ 改为 $z' = z + z_1 + wzz_1$。注意，$z' \in M$
在 $L$ 上生成 $M$，并且有
% P02-E0669
$\xi' = P(z') = P(z) + P(z_1) + wP(z)P(z_1) \in L$，且
$$
\xi' \equiv
\xi - \sigma(\lambda_j) \pi^{-j} + \sigma(\lambda_j^{p^r}) \pi^{-j'}
\bmod \pi^{-n + m}B
$$
；这是像上面一样使用方程 (\ref{more-algebra-equation-first-congruence}) 与
(\ref{more-algebra-equation-second-congruence}) 得到的。像先前一样写成
$$
\xi' = \sum\nolimits_{i = n, \ldots, n - m + 1} \sigma(\lambda'_i) \pi^{-i}
+ \pi^{-n + m}b'
$$
，便得到：当 $i \not = j, j'$ 时有 $\lambda'_i = \lambda_i$，并且
$\lambda'_j = 0$。因此集合 $J$ 变小了。由对 $J$ 的基数所作归纳，
可以假设 $J$ 中不存在使 $j'/j$ 为 $p$ 的幂的指标对 $j, j'$。
（请注意，在此过程中可能回到上面已经处理过的
$J = \emptyset$ 情形。）

\medskip\noindent
对 $j \in J$，写
$\lambda_j = \mu_j^{p^{r_j}}$，其中某个 $r_j \geq 0$，而且
$\mu_j \in \kappa_B$ 不是 $p$ 次幂。由假设 (4)，这是可能的。
令 $j \in J$ 是使 $j p^{-r_j}$ 最大的唯一指标（由前一段的结论，
该指标唯一）。选取
% P02-E0670
$r > \max(r_j + 1)$，并使得对 $j \in J$ 有
$j p^{r - r_j} > n$。令 $K_1/K$ 是由
$(\pi')^{p^r} = \pi$ 定义的、关于 $A$ 完全分歧的次数为 $p^r$
的扩张。注意，$\pi'$ 是相应离散赋值环 $A_1 \subset K_1$
的一致化元。还要注意，$L_1 = L \otimes_K K_1$ 是域，而且
$L_1/L$ 关于 $B$ 完全分歧
（引理 \ref{more-algebra-lemma-weakly-unramified-goes-up-along-totally-ramified}）。
在整闭包 $B_1$ 中计算得
$$
\xi = \sum\nolimits_{i \in I} \sigma(\lambda_i) (\pi')^{-i p^r} +
\sum\nolimits_{j \in J} \sigma(\mu_j)^{p^{r_j}} (\pi')^{-j p^r} +
\pi^{-n + m} b_1
$$
，其中某个 $b_1 \in B_1$。注意，对 $i \in I$，
% P02-E0671
$\sigma(\lambda_i)$ 模 $\pi^m$ 是 $q$ 次幂，换言之模
$(\pi')^{m p^r}$ 也有上述性质。因此可将上式改写为
$$
\xi = \sum\nolimits_{i \in I} x_i^{p^r} (\pi')^{-i p^r} +
\sum\nolimits_{j \in J} \sigma(\mu_j)^{p^{r_j}} (\pi')^{- j p^r}
+ \pi^{-n + m}b_1
$$
。与前一段中的选取类似，令
% P02-E0672
% P02-E0673
\begin{align*}
z_1 & - \sum\nolimits_{i \in I}
\left(x_i (\pi')^{-i} + \ldots + x_i^{p^{r - 1}} (\pi')^{-i p^{r - 1}}\right)
\\
& - \sum\nolimits_{j \in J}
\left(
\sigma(\mu_j) (\pi')^{- j p^{r - r_j}}
+ \ldots +
\sigma(\mu_j)^{p^{r_j - 1}} (\pi')^{- j p^{r - 1}}
\right)
\end{align*}
，并把所选的 $z$ 改为 $z' = z + z_1 + wzz_1$。于是 $z'$ 在
$L_1$ 上生成 $M_1$，并且
$\xi' = P(z') = P(z) + P(z_1) + w^p P(z) P(z_1) \in L_1$，
而计算表明
$$
\xi' \equiv
\sum\nolimits_{i \in I} x_i (\pi')^{-i} +
\sum\nolimits_{j \in J} \sigma(\mu_j) (\pi')^{- j p^{r - r_j}} +
(\pi')^{(-n + m)p^r}b'_1
$$
，其中某个 $b'_1 \in B_1$。恰有一个 $j$ 使得
$j p^{r - r_j}$ 最大，而且 $j p^{r - r_j}$ 大于每个 $i \in I$。
若 $j p^{r - r_j} \leq (n - m)p^r$，则扩张 $M_1/L_1$ 的层级小于
% P02-E0674
$\max(0, l - 1, 2l - p)$。否则，由于 $p$ 整除
$j p^{r - r_j}$，可知 $M_1 / L_1$ 属于引理
\ref{more-algebra-lemma-extension-defined-by-nice-polynial} 的情形 (C)。证明完成。
\end{proof}

\begin{lemma}
\label{more-algebra-lemma-special-case}
设 $A \subset B \subset C$ 是分式域为 $K \subset L \subset M$
的离散赋值环扩张。假设
\begin{enumerate}
\item $A$ 的剩余域 $k$ 是特征为 $p > 0$ 的代数闭域；
\item $A$ 与 $B$ 完备；
\item $A \to B$ 弱非分歧；
\item $M$ 是 $L$ 的有限扩张；
\item $k = \bigcap\nolimits_{n \geq 1} \kappa_B^{p^n}$。
\end{enumerate}
则存在 $A \to C$ 的弱解，即有限扩张 $K_1/K$。
\end{lemma}

\begin{proof}
设 $M'$ 是 $L$ 的任意有限扩张，并考虑 $B$ 在 $M'$ 中的整闭包
$C'$。由 Algebra，引理
\ref{algebra-lemma-Noetherian-complete-local-Nagata}，$B$ 是 Nagata
环，故 $C'$ 在 $B$ 上有限。此外，$C'$ 是离散赋值环，见评注
\ref{more-algebra-remark-construction} 中的讨论。又因为 $C'$ 作为 $B$-模完备，
所以它作为离散赋值环完备，见 Algebra，章节
\ref{algebra-section-completion}。特别地，这表明 $C$ 是 $B$ 在 $M$
中的整闭包（由赋值环是支配关系下的极大元这一定义）。

\medskip\noindent
设 $M \subset M'$ 是有限扩张，并设 $C' \subset M'$ 是如上所述的
$B$ 的整闭包。由引理 \ref{more-algebra-lemma-solutions-go-down}，只需对
$A \to B \to C'$ 证明结论。因此，由 Fields，引理
\ref{fields-lemma-normal-closure}，可以假设 $M/L$ 正规。

\medskip\noindent
若 $M / L$ 正规，则可找到有限扩张链
$$
L = L^0 \subset L^1 \subset L^2 \subset \ldots \subset L^r = M
$$
，使得每个扩张 $L^{j + 1}/L^j$ 都属于以下类型之一：
\begin{enumerate}
\item[(a)] 次数为 $p$ 的纯不可分扩张；
\item[(b)] 关于 $B^j$ 完全分歧、次数为 $p$ 的 Galois 扩张；
\item[(c)] 关于 $B^j$ 完全分歧、阶数与 $p$ 互素的循环 Galois 扩张；
\item[(d)] 关于 $B^j$ 非分歧的 Galois 扩张。
\end{enumerate}
这里 $B^j$ 是 $B$ 在 $L^j$ 中的整闭包。事实上，由于 $M/L$
正规，可将它写成一个 Galois 扩张与一个纯不可分扩张的合成
（Fields，引理 \ref{fields-lemma-normal-case}）。对纯不可分扩张，
滤过的存在性是清楚的。在 Galois 情形，注意 $G$ 是“该”分解群，
并令 $I \subset G$ 为惯性群。一方面，由引理
\ref{more-algebra-lemma-galois-inertia}，$I$ 可解；另一方面，由引理
\ref{more-algebra-lemma-inertial-invariants-unramified}，扩张 $M^I/L$ 关于 $B$
非分歧。这就证明了所述滤过的存在性。

\medskip\noindent
我们将对整数 $r$ 作归纳。假设可找到有限扩张 $K_1/K$，它是
$A \to B^1$ 的弱解，其中 $B^1$ 是 $B$ 在 $L^1$ 中的整闭包。
设 $K'_1$ 是 $K_1/K$ 的正规闭包
（Fields，引理 \ref{fields-lemma-normal-closure}）。由于 $A$ 完备且
$A$ 的剩余域代数闭，可知 $K'_1/K_1$ 关于 $A_1$ 可分并完全分歧
（略去一些细节）。因此，由引理
\ref{more-algebra-lemma-weakly-unramified-goes-up-along-totally-ramified}，
$K'_1/K$ 也是 $A \to B^1$ 的弱解。换言之，我们可以并且确实假设
$K_1$ 是 $K$ 的正规扩张。在此之后，考虑序列
$$
L^0_1 = (L^0 \otimes_K K_1)_{red} \subset
L^1_1 = (L^1 \otimes_K K_1)_{red} \subset \ldots \subset
L^r_1 = (L^r \otimes_K K_1)_{red}
$$
以及相应的整闭包 $B^i_1$。注意，$C_1 = B^r_1$ 是若干离散赋值环
的乘积；由引理 \ref{more-algebra-lemma-galois-relative}，
$G = \text{Aut}(K_1/K)$ 传递地置换这些离散赋值环。特别地，所有离散
赋值环扩张 $A_1 \to (C_1)_\mathfrak m$ 都彼此同构，而其中一个的
弱解就是所有这些扩张的弱解。可将归纳假设用于序列
$$
A_1 \to (B^1_1)_{B^1_1 \cap \mathfrak m} \to
(B^2_1)_{B^2_1 \cap \mathfrak m} \to
\ldots \to
(B^r_1)_{B^r_1 \cap \mathfrak m} =
(C_1)_\mathfrak m
$$
，得到 $A_1 \to (C_1)_\mathfrak m$ 的弱解 $K_2/K_1$。于是，由前述
讨论，扩张 $K_2/K$ 是 $A \to C$ 的弱解。注意，归纳假设确实适用：
由 $K_1$ 的选取，环映射
$A_1 \to (B^1_1)_{B^1_1 \cap \mathfrak m}$ 弱非分歧；而分式域扩张
的序列中的每一项仍具有上述 (a)、(b)、(c) 或 (d) 中的一种性质。
此外，注意对任意有限扩张 $\kappa_B \subset \kappa$，仍有
$k = \bigcap \kappa^{p^n}$。

\medskip\noindent
因此，一切都归结为：当域扩张 $M/L$ 具有 (a)、(b)、(c) 或 (d)
中的一种性质时，为 $A \subset C$ 找到弱解。

\medskip\noindent
情形 (d)。此情形平凡，因为此时 $B \to C$ 已经非分歧。

\medskip\noindent
情形 (c)。设 $M/L$ 是阶数为 $n$ 且其阶数与 $p$ 互素的循环扩张。
因为 $M/L$ 关于 $B$ 完全分歧，所以 $B \subset C$ 的分歧指数为
$n$，从而 $A \subset C$ 的分歧指数也是 $n$。选取一致化元
$\pi \in A$，并令 $K_1 = K[\pi^{1/n}]$。由 Abhyankar 引理
（引理 \ref{more-algebra-lemma-abhyankar}），$K_1/K$ 是 $A \subset C$ 的解。

\medskip\noindent
情形 (b)。把此情形分成混合特征情形与等特征情形。在等特征情形，
结论是引理 \ref{more-algebra-lemma-characteristic-p-case}。在混合特征情形，
先用引理 \ref{more-algebra-lemma-make-finite-level}，以一个有限扩张替换 $K$，
从而达到 $M/L$ 是有限层级的次数为 $p$ 的扩张这一情形。此时层级是
有理数 $l \in [0, p)$，见引理 \ref{more-algebra-lemma-lowering-the-level} 之前的
讨论。若层级为 $0$，则 $B \to C$ 弱非分歧，结论成立。否则，由引理
\ref{more-algebra-lemma-lowering-the-level}，可以用一个有限扩张替换域 $K$，得到
% P02-E0675
层级为 $l' \leq \max(0, l - 1, 2l - p)$ 的新情形。若
$l = p - \epsilon$ 且 $\epsilon < 1$，则
$l' \leq p - 2\epsilon$。因此，经过有限次替换，就得到层级
$\leq p - 1$ 的情形。再作至多 $p - 1$ 次这样的替换，便达到层级为零
的情形。

\medskip\noindent
情形 (a) 是引理 \ref{more-algebra-lemma-purely-inseparable-case}。这是唯一可能需要
$K$ 的纯不可分扩张的情形；也就是说，在该引理陈述的情形 (2) 中，
通过添加元素 $\pi$ 的一个 $p$ 次幂即可取胜。引理证明完成。
% P02-E0676
\end{proof}

\noindent
至此，我们已经收集了证明本节主要结果所需的所有引理。

\begin{theorem}[Epp]
\label{more-algebra-theorem-epp}
设 $A \subset B$ 是分式域为 $K \subset L$ 的离散赋值环扩张。
若 $\kappa_A$ 的特征为 $p > 0$，则假设
$$
\bigcap\nolimits_{n \geq 1} \kappa_B^{p^n}
$$
中的每个元素在 $\kappa_A$ 上都可分且代数。于是存在有限扩张
$K_1/K$，它是定义 \ref{more-algebra-definition-solution} 所定义的 $A \to B$
的弱解。
\end{theorem}

\begin{proof}
若 $\kappa_A$ 的特征为零，或者剩余特征为 $p$、分歧指数与 $p$
互素且剩余域扩张可分，则结论由 Abhyankar 引理
（引理 \ref{more-algebra-lemma-abhyankar}）得到。事实上，设分歧指数为 $e$。
选取一致化元 $\pi \in A$。令 $K_1/K$ 为添加 $\pi$ 的一个 $e$
次根而得到的扩张。由引理 \ref{more-algebra-lemma-pull-root-uniformizer}，
$A$ 在 $K_1$ 中的整闭包 $A_1$ 是在 $A$ 上具有分歧指数的离散赋值环。
% P02-E0677
因此，由引理 \ref{more-algebra-lemma-abhyankar}，对 $B_1$ 的每个极大理想
$\mathfrak m$，$A_1 \to (B_1)_\mathfrak m$ 在
$\mathfrak m$-进拓扑中形式光滑；从而更有，这些都是离散赋值环的
弱非分歧扩张。

\medskip\noindent
从现在起，设 $p$ 是素数，并假设 $\kappa_A$ 的特征为 $p$。首先使用
引理 \ref{more-algebra-lemma-solution-after-strict-henselization}，约化到 $A$ 与
$B$ 都有可分闭剩余域的情形。由于此过程分别用其可分代数闭包替换
$\kappa_A$ 与 $\kappa_B$，从定理的条件可知，我们得到
$$
\kappa_A \supset \bigcap\nolimits_{n \geq 1} \kappa_B^{p^n}
$$
。

\medskip\noindent
设 $\pi \in A$ 是一致化元，并设 $A^\wedge$ 与 $B^\wedge$ 分别为
$A$ 与 $B$ 的完备化。我们有离散赋值环扩张的交换图
$$
\xymatrix{
B \ar[r] & B^\wedge \\
A \ar[u] \ar[r] & A^\wedge \ar[u]
}
$$
。设 $K^\wedge$ 是 $A^\wedge$ 的分式域。假设可找到有限扩张
$M/K^\wedge$，使它 (a) 是 $A^\wedge \to B^\wedge$ 的弱解，并且
(b) 是一个可分扩张与一个通过添加 $\pi$ 的某个 $p$-幂次根而得到的
扩张的合成。于是由引理 \ref{more-algebra-lemma-approximate-separable-extension}，
可找到有限扩张 $K_1/K$，使得
$K^\wedge \otimes_K K_1 = M$。令 $A_1$（相应地，$A_1^\wedge$）
为 $A$（相应地，$A^\wedge$）在 $K_1$（相应地，$M$）中的整闭包。
由于 $A \to A^\wedge$ 在 $\mathfrak m^\wedge$-进拓扑中形式光滑
（引理 \ref{more-algebra-lemma-extension-dvrs-formally-smooth}），可知
$A_1 \to A_1^\wedge$ 在 $\mathfrak m_1^\wedge$-进拓扑中形式光滑
（引理 \ref{more-algebra-lemma-formally-smooth-goes-up}；由评注
\ref{more-algebra-remark-construction} 中的讨论，$A_1$ 与 $A_1^\wedge$ 是离散赋值环）。
由引理 \ref{more-algebra-lemma-solutions-go-down} 的 (2) 可得，$K_1/K$ 是
$A \to B^\wedge$ 的弱解。再用引理 \ref{more-algebra-lemma-solutions-go-down}
的 (1)，可知 $K_1/K$ 是 $A \to B$ 的弱解。

\medskip\noindent
因此，可以假设 $A$ 与 $B$ 是完备离散赋值环，具有特征为 $p$ 的
可分闭剩余域，并且
$\kappa_A \supset \bigcap\nolimits_{n \geq 1} \kappa_B^{p^n}$。
此外，给定一致化元 $\pi \in A$，我们必须找到 $A \to B$ 的弱解，
它是一个可分扩张与一个通过取 $\pi$ 的 $p$-幂次根而得到的域的合成。
注意，若 $A$ 具有混合特征，则第二个条件自动成立。

\medskip\noindent
令 $k = \bigcap\nolimits_{n \geq 1} \kappa_B^{p^n}$。注意，$k$
是特征为 $p$ 的代数闭域。若 $A$ 具有混合特征，则令 $\Lambda$ 为
$k$ 的 Cohen 环；在等特征情形，令 $\Lambda = k[[t]]$。可以选取环映射
$\Lambda \to A$，它在等特征情形把 $t$ 映到 $\pi$。在等特征情形，
这由 Cohen 结构定理（Algebra，定理
\ref{algebra-theorem-cohen-structure-theorem}）得到；在混合特征情形，
这由 $\mathbf{Z}_p \to \Lambda$ 在进拓扑中形式光滑得到
（引理 \ref{more-algebra-lemma-extension-dvrs-formally-smooth} 与
\ref{more-algebra-lemma-lift-continuous}）。应用引理 \ref{more-algebra-lemma-solutions-go-down}，
可知只需证明 $\Lambda \to B$ 存在弱解，并要求在等特征 $p$ 情形，
该弱解是一个可分扩张与一个通过取 $t$ 的 $p$-幂次根而得到的域的合成。
然而，由于在等特征情形 $\Lambda = k[[t]]$，且 $k((t))$ 的任意扩张
都是这样的合成，现在可以去掉这一要求！

\medskip\noindent
这样便达到如下情形：$A$ 与 $B$ 完备；$A$ 的剩余域 $k$ 是特征为
$p > 0$ 的代数闭域；有 $k = \bigcap \kappa_B^{p^n}$；并且在混合特征
情形，$p$ 是 $A$ 的一致化元（即 $A$ 是 $k$ 的 Cohen 环）。若 $A$
具有混合特征，则选取 $\kappa_B$ 的 Cohen 环 $\Lambda$；在等特征情形，
令 $\Lambda = \kappa_B[[t]]$。像上面一样论证，可以选取环映射
$A \to \Lambda$，它提升 $k \to \kappa_B$ 并把一致化元映到一致化元。
由于 $k \subset \kappa_B$ 可分，环映射 $A \to \Lambda$ 在进拓扑中
形式光滑（引理 \ref{more-algebra-lemma-extension-dvrs-formally-smooth}）。因此可找到
环映射 $\Lambda \to B$，使得复合映射
$A \to \Lambda \to B$ 是给定的环映射 $A \to B$
（见引理 \ref{more-algebra-lemma-lift-continuous}）。由于 $\Lambda$ 与 $B$ 是具有
相同剩余域的完备离散赋值环，由 Algebra，引理
\ref{algebra-lemma-finite-over-complete-ring}，$B$ 在 $\Lambda$ 上有限。
这就约化到引理 \ref{more-algebra-lemma-special-case} 中讨论的特殊情形。
\end{proof}

\section{消除分歧，II}
\label{more-algebra-section-eliminating-ramification-bis}

\noindent
本节使用章节 \ref{more-algebra-section-eliminating-ramification} 的结果，在若干情形
得到（可分的）解。

\begin{lemma}
\label{more-algebra-lemma-solution-goes-up}
设 $A \to B$ 是分式域为 $K \subset L$ 的离散赋值环扩张。若
$K_1/K$ 是 $A \subset B$ 的解，则对任意有限扩张 $K_2/K_1$，
扩张 $K_2/K$ 都是 $A \subset B$ 的解。
\end{lemma}

\begin{proof}
这由引理 \ref{more-algebra-lemma-formally-smooth-goes-up} 得到。细节从略。
\end{proof}

\begin{lemma}
\label{more-algebra-lemma-nagata-goes-down}
设 $A \subset B$ 是离散赋值环扩张。若 $B$ 是 Nagata 环，且分式域
扩张 $L/K$ 可分，则 $A$ 是 Nagata 环。
\end{lemma}

\begin{proof}
离散赋值环是 Nagata 环当且仅当它是 N-2 的。设 $K_1/K$ 是有限纯不
可分域扩张。我们必须证明 $A$ 在 $K_1$ 中的整闭包 $A_1$ 在 $A$
上有限，见 Algebra，引理 \ref{algebra-lemma-domain-char-p-N-1-2}。
由于 $L/K$ 可分而 $K_1/K$ 纯不可分，代数 $L \otimes_K K_1$ 是域
（由 Algebra，引理
\ref{algebra-lemma-separable-extension-preserves-reducedness} 与
\ref{algebra-lemma-radicial-integral-bijective}）。设 $B_1$ 是 $B$ 在
$L \otimes_K K_1$ 中的整闭包。由于 $B$ 是 Nagata 环，$B_1$ 在
$B$ 上有限。由 $B \otimes_A A_1 \subset B_1$ 以及 $B$ 是
Noether 环，可知 $B \otimes_A A_1$ 在 $B$ 上有限。由于
$A \to B$ 忠实平坦，这蕴含 $A_1$ 在 $A$ 上有限，见 Algebra，
引理 \ref{algebra-lemma-descend-properties-modules}。
\end{proof}

\begin{lemma}
\label{more-algebra-lemma-construct-extension}
设 $A' \subset A$ 是环扩张，并设 $f \in A'$。假设 (a) $A$ 在
$A'$ 上有限，(b) $f$ 是 $A$ 上的非零因子，并且 (c)
$A'_f = A_f$。则存在整数 $n_0 > 0$，使得对所有 $n \geq n_0$，
以下结论成立：给定环 $B'$、非零因子 $g \in B'$ 以及满足
$\varphi'(f) \equiv g$ 的同构
$\varphi' : A'/f^n A' \to B'/g^n B'$，存在有限扩张
$B' \subset B$ 以及与 $\varphi'$ 相容的同构
$\varphi : A/fA \to B/gB$。
\end{lemma}

\begin{proof}
% P02-E0678
由于 $A$ 在 $A'$ 上有限且 $A'_f = A_f$，可选取 $t > 0$，使得
$f^t A \subset A'$。令 $n_0 = 2t$。给定引理陈述中的
$n, B', g, \varphi'$，以 $N \subset B'$ 表示由满足
$b \bmod g^nB' \in \varphi'(f^tA)$ 的 $b \in B'$ 组成的集合。
令 $B = g^{-t}N$。由于 $f^tA' \subset f^tA$，且 $\varphi'$ 把
$f$ 映到 $g$，有 $g^tB' \subset N$，故 $B' \subset B$。由于
$f^tA \cdot f^tA \subset f^t \cdot f^tA$，且 $\varphi'$ 把 $f$
映到 $g$，可知 $N \cdot N \subset g^t N$。因此在 $B$ 上得到一个
延拓 $B'$ 之乘法的乘法。我们有 $A'/f^nA'$-模的同构
% P02-E0679
$$
A/f^tA' \xrightarrow{f^t} f^tA/f^nA' \xrightarrow{\varphi'}
g^tB/g^nB' \xrightarrow{g^{-t}} B/g^tB'
$$
，其中右边的模结构用 $\varphi'$ 定义。由于 $A/f^tA'$ 是有限
$A'$-模，可得 $B/g^tB'$ 是有限 $B'$-模，从而 $B' \to B$ 有限。
最后，留给读者验证：所显示的模同构把 $fA$ 映入 $gB$，并诱导与
$\varphi'$ 相容的环同构 $\varphi : A/fA \to B/gB$（它甚至诱导同构
$A/f^tA \to B/g^tB$，但我们不需要这一点）。
\end{proof}

\begin{remark}
\label{more-algebra-remark-functoriality-construct-extension}
引理 \ref{more-algebra-lemma-construct-extension} 中的构造满足以下“函子性”。假设有
交换图
$$
\xymatrix{
A'_2 \ar[r] & A_2 \\
A'_1 \ar[r] \ar[u] & A_1 \ar[u]
}
$$
，其中水平箭头单射。假设给定元素 $f \in A'_1$，使得
$(A'_1 \subset A_1, f)$ 与 $(A'_2 \subset A_2, f)$ 满足引理
\ref{more-algebra-lemma-construct-extension} 的性质 (a)、(b)、(c)。设
$n_{0, 1}$ 与 $n_{0, 2}$ 是引理对这两个情形所给出的整数。最后，
设 $B'_1 \to B'_2$ 是环映射，设
% P02-E0680
$g \in B'_1$ 是 $B_1$ 与 $B_2$ 上的非零因子，设
$n \geq \max(n_{0, 1}, n_{0, 2})$，并给定交换图
% P02-E0681
$$
\xymatrix{
A'_2/f^nA'_2 \ar[r]_{\varphi'_2} & B'_2/g^nB'_2 \\
A'_1/f^nA'_1 \ar[r]^{\varphi'_1} \ar[u] & B'_2/g^nB'_2 \ar[u]
}
$$
，其中水平箭头是同构，且 $\varphi'_1(f) \equiv g$。于是得到交换图
% P02-E0682
$$
\vcenter{
\xymatrix{
B'_2 \ar[r] & B_2 \\
B'_1 \ar[r] \ar[u] & B_1 \ar[u]
}
}
\quad\text{且}\quad
\vcenter{
\xymatrix{
A_2/fA_2 \ar[r]_{\varphi_2} & B_2/gB_2 \\
A_1/fA_1 \ar[r]^{\varphi_1} \ar[u] & B_2/gB_2 \ar[u]
}
}
$$
，其中 $(B'_1 \subset B_1, \varphi_1)$ 与
$(B'_2 \subset B_2, \varphi_2)$ 按引理
\ref{more-algebra-lemma-construct-extension} 的证明构造。详细验证从略。
\end{remark}

\begin{lemma}
\label{more-algebra-lemma-approximate-solution}
设 $p$ 是素数。设 $A \subset B$ 是分式域扩张为 $L/K$ 的离散赋值环
扩张，并设 $K_2/K_1/K$ 是有限域扩张塔。假设
\begin{enumerate}
\item $K$ 的特征为 $p$；
\item $L/K$ 可分；
\item $B$ 是 Nagata 环；
\item $K_2$ 是 $A \subset B$ 的解；
\item $K_2/K_1$ 是次数为 $p$ 的纯不可分扩张。
\end{enumerate}
则存在可分扩张 $K_3/K_1$，它是 $A \subset B$ 的解。
\end{lemma}

\begin{proof}
使用评注 \ref{more-algebra-remark-construction} 中的记号；我们将使用那里所作的
全部观察。由于 $L/K$ 可分，代数 $L_1 = L \otimes_K K_1$ 约化
（Algebra，引理
\ref{algebra-lemma-separable-extension-preserves-reducedness}）。
由于 $B$ 是 Nagata 环，环扩张 $B \subset B_1$ 有限，其中 $B_1$
是 $B$ 在 $L_1$ 中的整闭包，并且 $B_1$ 是 Nagata 环。类似地，由
引理 \ref{more-algebra-lemma-nagata-goes-down}，环 $A$ 是 Nagata 环，因而
$A \subset A_1$ 有限，且 $A_1$ 也是 Nagata 环。此外，对 $K_2$
也有同样断言；即 $L_2 = L \otimes_K K_2$ 约化，环扩张
$A_1 \subset A_2$ 与 $B_1 \subset B_2$ 有限，其中
$A_2$（相应地，$B_2$）是 $A$（相应地，$B$）在 $K_2$
（相应地，$L_2$）中的整闭包。

\medskip\noindent
设 $\pi \in A$ 是一致化元。注意，$\pi$ 是 $K_1$、$K_2$、$A_1$、
$A_2$、$L_1$、$L_2$、$B_1$ 与 $B_2$ 上的非零因子，并且
$K_1 = (A_1)_\pi$、$K_2 = (A_2)_\pi$、$L_1 = (B_1)_\pi$ 且
$L_2 = (B_2)_\pi$。由 Fields，引理
\ref{fields-lemma-finite-purely-inseparable}，可写
$K_2 = K_1(\alpha)$，其中 $\alpha^p = a_1 \in K_1$。把 $\alpha$
乘以 $\pi$ 的某个幂后，可以并且确实假设 $a_1 \in A_1$。在证明的
其余部分，为方便起见写成
$K_2 = K_1[x]/(x^p - a_1)$ 与 $L_2 = L_1[x]/(x^p - a_1)$。
考虑环扩张
$$
A'_2 = A_1[x]/(x^p - a_1) \subset A_2
\quad\text{且}\quad
B'_2 = B_1[x]/(x^p - a_1) \subset B_2
$$
。可将引理 \ref{more-algebra-lemma-construct-extension} 分别用于
$A'_2 \subset A_2$、$f = \pi^2$ 以及
$B'_2 \subset B_2$、$f = \pi^2$。选取对这两者都适用的充分大整数
$n$。

\medskip\noindent
考虑代数
% P02-E0685
$$
K_3 = K_1[x]/(x^p - \pi^{2n} x - a_1)
\quad\text{且}\quad
L_3 = L_1[x]/(x^p - \pi^{2n} x - a_1)
$$
。注意，$K_3/K_1$ 与 $L_3/L_1$ 是次数为 $p$ 的有限
\'etale 代数扩张。考虑子环
% P02-E0683
$$
A'_3 = A_1[x]/(x^p - \pi^n x - a_1)
\quad\text{且}\quad
B'_3 = B_1[x]/(x^p - \pi^n x - a_1)
$$
；它们分别是
% P02-E0684
$K_3 = (A'_2)_\pi$ 与 $L_3 = (B'_3)_\pi$ 的子环。我们将构造交换图
$$
\xymatrix{
B'_2/\pi^{2n} B'_2 \ar[r]_{\psi'} & B'_3/\pi^{2n} B'_3 \\
A'_2/\pi^{2n} A'_2 \ar[r]^{\varphi'} \ar[u] & A'_3/\pi^{2n} A'_3 \ar[u]
}
$$
。事实上，$\varphi'$ 是把 $x$ 的类映到 $x$ 的类的唯一 $A_1$-代数
同构。类似地，$\psi'$ 是把 $x$ 的类映到 $x$ 的类的唯一 $B_1$-代数
同构。由 $n$ 的选取，通过引理 \ref{more-algebra-lemma-construct-extension} 与评注
\ref{more-algebra-remark-functoriality-construct-extension}，得到有限环扩张
$A'_3 \subset A_3$ 与 $B'_3 \subset B_3$，使得
$A'_3 \to B'_3$ 延拓为环映射 $A_3 \to B_3$，并得到交换图
$$
\xymatrix{
B_2/\pi^2 B_2 \ar[r]_\psi & B_3/\pi^2B_3 \\
A_2/\pi^2 A_2 \ar[r]^\varphi \ar[u] & A_3/\pi^2A_3 \ar[u]
}
$$
，它具有所引用结果中断言的所有性质（特别地，$\varphi$ 与 $\psi$
是同构）。

\medskip\noindent
有了所有这些数据，就可以完成证明。首先注意，$A_3$ 与 $B_3$ 是
若干 Dedekind 整环的有限乘积，并且 $\pi$ 属于它们的所有极大理想。
事实上，若 $\mathfrak p \subset A_3$ 是极大理想，则由于
$A \to A_3$ 有限，有 $\pi \in \mathfrak p$。于是
$\mathfrak p/\pi^2 A_3$ 通过 $\varphi$ 对应于
$A_2 / \pi^2A_2$ 中的极大理想；后者是主理想，因为 $A_2$ 是若干
Dedekind 整环的有限乘积。可得 $\mathfrak p/\pi^2 A_3$ 是主理想，
因而由 Nakayama 引理，$\mathfrak p (A_3)_\mathfrak p$ 是主理想。
同样的论证适用于 $B_3$。由此可得，$A_3$ 是 $A$ 在 $K_3$ 中的
整闭包，而 $B_3$ 是 $B$ 在 $L_3$ 中的整闭包。设
$\mathfrak q \subset B_3$ 是位于 $\mathfrak p \subset A_3$ 之上的
极大理想。为完成证明，必须证明
$(A_3)_\mathfrak p \to (B_3)_\mathfrak q$ 在
$\mathfrak q$-进拓扑中形式光滑。由引理
\ref{more-algebra-lemma-extension-dvrs-formally-smooth} 的判据，只需证明
$\mathfrak p (B_3)_\mathfrak q = \mathfrak q (B_3)_\mathfrak q$，并且
域扩张 $\kappa(\mathfrak q)/\kappa(\mathfrak p)$ 可分。这是真的，
因为两个断言都可以通过观察环映射
$A_3/\pi^2 A_3 \to B_3/\pi^2 B_3$ 来检验，而该环映射同构于环映射
$A_2/\pi^2 A_2 \to B_2/\pi^2 B_2$；由于假设 $K_2$ 是
$A \subset B$ 的解，相应断言对后一个映射成立。略去一些细节。
\end{proof}

\begin{lemma}
\label{more-algebra-lemma-separable-solution-separable-solution}
设 $A \subset B$ 是离散赋值环扩张。假设
\begin{enumerate}
\item 分式域扩张 $L/K$ 可分；
\item $B$ 是 Nagata 环；并且
\item $A \subset B$ 存在解。
\end{enumerate}
则 $A \subset B$ 存在可分的解。
\end{lemma}

\begin{proof}
若 $K$ 的特征为零，引理平凡；因此可以并且确实假设 $K$ 的特征为
$p > 0$。

\medskip\noindent
设 $K_2/K$ 是 $A \to B$ 的解。我们将对 $K_2/K$ 的不可分次数
$[K_2 : K]_i$（Fields，定义 \ref{fields-definition-insep-degree}）
作归纳。若 $[K_2 : K]_i = 1$，则 $K_2$ 在 $K$ 上可分，结论成立。
否则，存在子域塔 $K_2/K_1/K$，使得 $K_2/K_1$ 是次数为 $p$ 的
纯不可分扩张（Fields，引理 \ref{fields-lemma-separable-first} 与
\ref{fields-lemma-finite-purely-inseparable}）。由引理
\ref{more-algebra-lemma-approximate-solution}，存在可分扩张 $K_3/K_1$，它是
$A \subset B$ 的解。于是
$[K_3 : K]_i = [K_1 : K]_i = [K_2 : K]_i/p$
（Fields，引理 \ref{fields-lemma-multiplicativity-all-degrees}）更小，
由归纳得出结论。
\end{proof}

\begin{lemma}
\label{more-algebra-lemma-big-extension-is-ok}
设 $A \to B$ 是分式域为 $K \subset L$ 的离散赋值环扩张。假设 $B$
在 $A$ 上本质有限型。设 $K'/K$ 是代数域扩张，使得 $A$ 在 $K'$
中的整闭包 $A'$ 是 Noether 环。则 $B$ 在
$L' = (L \otimes_K K')_{red}$ 中的整闭包 $B'$ 也是 Noether 环。
此外，映射 $\Spec(B') \to \Spec(A')$ 满射，并且相应的剩余域扩张
是有限生成域扩张。
\end{lemma}

\begin{proof}
设 $A \to C$ 是有限型环映射，使得 $B$ 是 $C$ 在素理想
$\mathfrak p$ 处的局部化。则 $C' = C \otimes_A A'$ 是有限型
$A'$-代数，特别地是 Noether 环。由于 $A \to A'$ 整，
$C \to C'$ 也整。因此 $B = C_\mathfrak p \subset C'_\mathfrak p$
也整。由此，$C'_\mathfrak p$ 的维数为 $1$
（Algebra，引理 \ref{algebra-lemma-integral-sub-dim-equal}）。当然，
$C'_\mathfrak p$ 是 Noether 环。设
$\mathfrak q_1, \ldots, \mathfrak q_n$ 是 $C'_\mathfrak p$ 的极小
素理想。设 $B'_i$ 是 $B = C_\mathfrak p$ 的整闭包；或者由上文等价地，
% P02-E0686
是 $C'_\mathfrak p$ 在 $C'_{\mathfrak p'}/\mathfrak q_i$ 的分式域
中的整闭包。由 Krull--Akizuki 定理（将 Algebra，引理
\ref{algebra-lemma-krull-akizuki} 用于 $C'_\mathfrak p$ 在其极大理想
处的有限多个局部化），每个 $B'_i$ 都是 Noether 环。此外，由
Algebra，引理
\ref{algebra-lemma-finite-extension-residue-fields-dimension-1}，
$C'_\mathfrak p \to B'_i$ 中的剩余域扩张有限。最后注意，
$B' = \prod B'_i$ 是 $B$ 在
$L' = (L \otimes_K K')_{red}$ 中的整闭包。
\end{proof}

\begin{proposition}
\label{more-algebra-proposition-epp-essentially-finite-type}
\begin{reference}
特殊情形见 \cite[Lemma 2.13]{alterations}。
\end{reference}
设 $A \to B$ 是分式域为 $K \subset L$ 的离散赋值环扩张。若 $B$
在 $A$ 上本质有限型，则存在有限扩张 $K_1/K$，它是定义
\ref{more-algebra-definition-solution} 所定义的 $A \to B$ 的解。
\end{proposition}

\begin{proof}
注意，若 $A$ 的剩余域完美，则弱解就是解，见引理
\ref{more-algebra-lemma-extension-dvrs-formally-smooth}。因此，若 $A$ 的剩余特征为
$0$，该命题立即由定理 \ref{more-algebra-theorem-epp} 得到（事实上，无需假设
$A \to B$ 本质有限型）。若 $A$ 的剩余特征为 $p > 0$，也将从 Epp
定理推出该命题。

\medskip\noindent
设 $x_i \in A$（$i \in I$）是一族元素，它们映到 $A$ 的剩余域
$\kappa$ 的一个 $p$-基。令
% P02-E0688
$$
A' = \bigcup\nolimits_{n \geq 1} A[t_{i, n}]/(t_{i, n}^{p^n} - x_i)
$$
，其中转移映射把 $t_{i, n + 1}$ 映到 $t_{i, n}^p$。注意，$A'$
是 $A$ 上离散赋值环的弱非分歧有限扩张的滤过余极限。因此 $A'$ 是
离散赋值环，且 $A \to A'$ 弱非分歧。由构造，剩余域
$\kappa' = A'/\mathfrak m_A A'$ 是 $\kappa$ 的完美化。

\medskip\noindent
设 $K'$ 是 $A'$ 的分式域。可将引理 \ref{more-algebra-lemma-big-extension-is-ok}
用于扩张 $K'/K$。因此 $B'$ 是若干 Dedekind 整环的有限乘积。设
$\mathfrak m_1, \ldots, \mathfrak m_n$ 是 $B'$ 的极大理想。使用 Epp
定理（定理 \ref{more-algebra-theorem-epp}），对每个扩张
$A' \subset B'_{\mathfrak m_i}$ 找到弱解 $K'_i/K'$。由于 $A'$
的剩余域完美，这些实际上都是解。设 $K'_1/K'$ 是包含每个 $K'_i$
的有限扩张。由引理 \ref{more-algebra-lemma-solution-goes-up}，$K'_1/K'$ 仍是每个
$A' \subset B'_{\mathfrak m_i}$ 的解。

\medskip\noindent
设 $A'_1$ 是 $A$ 在 $K'_1$ 中的整闭包。注意，把评注
\ref{more-algebra-remark-construction} 中的讨论用于 $K' \subset K'_1$，可知
$A'_1$ 是 Dedekind 整环。因此，引理 \ref{more-algebra-lemma-big-extension-is-ok}
适用于 $K'_1/K$。从而 $B$ 在
$L'_1 = (L \otimes_K K'_1)_{red}$ 中的整闭包 $B'_1$ 是 Dedekind
整环；并且由于 $K'_1/K'$ 是每个 $A' \subset B'_{\mathfrak m_i}$
的解，可知对每个极大理想 $\mathfrak m \subset B'_1$，
% P02-E0687
$(A'_1)_{A'_1 \cap \mathfrak m} \to (B'_1)_{\mathfrak m}$ 在
$\mathfrak m$-进拓扑中形式光滑。

\medskip\noindent
由构造，域 $K'_1$ 是 $K$ 的有限扩张的滤过余极限。写
$K'_1 = \colim_{i \in I} K_i$。对每个 $i$，令 $A_i$（相应地，
$B_i$）为 $A$（相应地，$B$）在 $K_i$（相应地，
$L_i = (L \otimes_K K_i)_{red}$）中的整闭包。于是显然
$$
A'_1 =  \colim A_i\quad\text{且}\quad B'_1 = \colim B_i
$$
。由于环映射 $A_i \to A'_1$ 与 $B_i \to B'_1$ 是单射整环映射，
而 $A'_1$ 与 $B'_1$ 的谱有限，可知对所有充分大的 $i$，环映射
$A_i \to A'_1$ 与 $B_i \to B'_1$ 在谱上双射。一旦如此，对所有充分
大的 $i$，映射 $A_i \to A'_1$ 与 $B_i \to B'_1$ 将弱非分歧
（只要一致化元已在像中）。由分歧指数的乘法性可知，对这样的 $i$，
$A_i \to B_i$ 在 $B_i$ 的所有极大理想局部化上诱导弱非分歧映射。
再稍微增大 $i$，可知
$$
B_i \otimes_{A_i} A'_1 \longrightarrow B'_1
$$
在剩余域上诱导满射（因为由引理 \ref{more-algebra-lemma-big-extension-is-ok}，
$B'_1$ 的剩余域在 $A'_1$ 的相应剩余域上有限生成）。在 $B'_1$
的一个选定极大理想之下的极大理想处，剩余域的图为
$$
\xymatrix{
\kappa_{B_i} \ar[r] &
\kappa_{B_{i'}} \ar[r] &
 & \ldots &
\kappa_{B'_1} \\
\kappa_{A_i} \ar[r] \ar[u] &
\kappa_{A_{i'}} \ar[r] \ar[u] &
 & \ldots &
\kappa_{A'_1} \ar[u]
}
$$
。因此，$\kappa_{B_i}$ 是 $\kappa_{A_i}$ 的有限生成扩张，使得
$\kappa_{B_i}$ 与 $\kappa_{A'_1}$ 在 $\kappa_{B'_1}$ 中的合成在
$\kappa_{A'_1}$ 上可分。此事已经在某个有限阶段发生：例如，设
$\kappa_{B'_1}$ 在 $\kappa_{A'_1}(x_1, \ldots, x_n)$ 上有限可分；
只需增大 $i$，使得 $x_1, \ldots, x_n$ 属于 $\kappa_{B_i}$，并使
所有生成元都满足系数在 $\kappa_{A_i}(x_1, \ldots, x_n)$ 中的可分
多项式方程。这意味着，对 $B_i$ 的所有极大理想 $\mathfrak m$，
$A_i \to (B_i)_\mathfrak m$ 在 $\mathfrak m$-进拓扑中形式光滑，
证明完成。
\end{proof}

\begin{lemma}
\label{more-algebra-lemma-epp-essentially-finite-type-separable}
设 $A \to B$ 是分式域为 $K \subset L$ 的离散赋值环扩张。假设
\begin{enumerate}
\item $B$ 在 $A$ 上本质有限型；
\item $A$ 或 $B$ 是 Nagata 环；并且
\item $L/K$ 可分。
\end{enumerate}
则 $A \to B$ 存在可分的解（定义 \ref{more-algebra-definition-solution}）。
\end{lemma}

\begin{proof}
注意，若 $A$ 是 Nagata 环，则 $B$ 也是 Nagata 环
（Algebra，引理 \ref{algebra-lemma-nagata-localize} 与命题
\ref{algebra-proposition-nagata-universally-japanese}）。因此，结合命题
\ref{more-algebra-proposition-epp-essentially-finite-type} 与引理
\ref{more-algebra-lemma-separable-solution-separable-solution} 即得本引理。
\end{proof}

\section{环的 Picard 群}
\label{more-algebra-section-picard}

\noindent
首先如下定义可逆模。

\begin{definition}
\label{more-algebra-definition-invertible}
设 $R$ 是环。若函子
$$
\text{Mod}_R \longrightarrow \text{Mod}_R,\quad
N \longmapsto M \otimes_R N
$$
是范畴等价，则称 $R$-模 $M$ \textit{可逆}。若可逆 $R$-模作为
$R$-模同构于 $R$，则称它\textit{平凡}。
\end{definition}

\begin{lemma}
\label{more-algebra-lemma-invertible}
设 $R$ 是环，并设 $M$ 是 $R$-模。下列条件等价：
\begin{enumerate}
\item $M$ 是秩为 $1$ 的有限局部自由模；
\item $M$ 可逆；
\item 存在 $R$-模 $N$，使得 $M \otimes_R N \cong R$。
\end{enumerate}
此外，在这种情形，(3) 中的模 $N$ 同构于 $\Hom_R(M, R)$。
\end{lemma}

\begin{proof}
假设 (1)。考虑模 $N = \Hom_R(M, R)$ 与求值映射
$M \otimes_R N = M \otimes_R \Hom_R(M, R) \to R$。若
$f \in R$ 使得 $M_f \cong R_f$，则求值映射在 $f$ 处局部化后成为
同构（细节从略）。因此，由 Algebra，引理 \ref{algebra-lemma-cover}，
求值映射是同构。故 (1) $\Rightarrow$ (3)。

\medskip\noindent
假设 (3)。则函子 $K \mapsto K \otimes_R N$ 是函子
$K \mapsto K \otimes_R M$ 的拟逆。因此 (3) $\Rightarrow$ (2)。
反之，若 (2) 成立，则 $K \mapsto K \otimes_R M$ 本质满，从而 (3)
成立。

\medskip\noindent
假设等价条件 (2) 与 (3) 成立。以
$\psi : M \otimes_R N \to R$ 表示 (3) 中的同构。选取元素
$\xi = \sum_{i = 1, \ldots, n} x_i \otimes y_i$，使得
$\psi(\xi) = 1$。考虑同构
$$
M \to M \otimes_R M \otimes_R N \to M
$$
，其中第一个箭头把 $x$ 映到 $\sum x_i \otimes x \otimes y_i$，
第二个箭头把 $x \otimes x' \otimes y$ 映到
$\psi(x' \otimes y)x$。由此可得，
$x \mapsto \sum \psi(x \otimes y_i)x_i$ 是 $M$ 的自同构。该自同构
分解为
$$
M \to R^{\oplus n} \to M
$$
，其中第一个箭头由
$x \mapsto (\psi(x \otimes y_1), \ldots, \psi(x \otimes y_n))$
给出，第二个箭头由
$(a_1, \ldots, a_n) \mapsto \sum a_i x_i$ 给出。这样便可断定 $M$
是有限自由 $R$-模的直和项。这意味着 $M$ 有限局部自由
（Algebra，引理 \ref{algebra-lemma-finite-projective}）。由对称性，
$N$ 也如此；又因为 $M \otimes_R N \cong R$，可知 $M$ 与 $N$
的秩都必须为 $1$。
\end{proof}

\noindent
这些模的同构类集合通常称为 $R$ 的\textit{类群}或\textit{Picard 群}。
群结构由如下规则确定：把可逆模 $L$ 与 $L'$ 的同构类映到
$L \otimes_R L'$ 的同构类。可逆模 $L$ 的逆元是模
$$
L^{\otimes -1} = \Hom_R(L, R),
$$
因为如引理 \ref{more-algebra-lemma-invertible} 的证明所示，求值映射
$L \otimes_R L^{\otimes -1} \to R$ 是同构。以 $\Pic(R)$ 表示 $R$
的 Picard 群。

\begin{lemma}
\label{more-algebra-lemma-UFD-Pic-trivial}
设 $R$ 是 UFD。则 $\Pic(R)$ 平凡。
\end{lemma}

\begin{proof}
设 $L$ 是可逆 $R$-模。由引理 \ref{more-algebra-lemma-invertible}，$L$ 是有限局部
自由 $R$-模。特别地，$L$ 无挠并且在 $R$ 上有限。从对偶可逆模中
选取非零元素 $\varphi \in \Hom_R(L, R)$。则
$I = \varphi(L) \subset R$ 是作为模可逆的理想。选取非零
$f \in I$，并设
$$
f = u p_1^{e_1} \ldots p_r^{e_r}
$$
是它分解成两两不同的素元 $p_i$ 的分解。由于 $L$ 有限局部自由，
存在 $a_i \in R$、$a_i \not \in (p_i)$，使得对某个
$g_i \in R_{a_i}$ 有 $I_{a_i} = (g_i)$。于是 $p_i$ 在 UFD
$R_{a_i}$ 中仍是素元，并且可对某个不被 $p_i$ 整除的
$g'_i \in R_{a_i}$ 写成 $g_i = p_i^{c_i} g'_i$。由于
$f \in I_{a_i}$，可知 $e_i \geq c_i$。我们断言 $I$ 由
$h = p_1^{c_1} \ldots p_r^{c_r}$ 生成；这将完成证明。

\medskip\noindent
为证明该断言，只需对任意使 $I_a$ 为主理想的 $a \in R$ 证明
$I_a$ 由 $h$ 生成（Algebra，引理 \ref{algebra-lemma-cover}）。设
$I_a = (g)$。令 $J \subset \{1, \ldots, r\}$ 为使 $p_i$ 在 $R_a$
中是非单位（从而是素元）的 $i$ 的集合。由于 $f \in I_a = (g)$，
得到素因子分解
% P02-E0689
$g = v \prod_{i \in J} p_j^{b_j}$，其中 $v$ 是单位，且
$b_j \leq e_j$。对每个 $j \in J$，有
$I_{aa_j} = g R_{aa_j} = g_j R_{aa_j}$；换言之，$g$ 与 $g_j$
在 $R_{aa_j}$ 中的像相伴。由分解的唯一性，这蕴含 $b_j = c_j$，
证明完成。
\end{proof}

\section{行列式}
\label{more-algebra-section-determinants}

\noindent
设 $R$ 是环，并设 $M$ 是有限投射 $R$-模。存在乘积分解
$R = R_0 \times \ldots \times R_t$，使得在 $M$ 的相应分解
$M = M_0 \times \ldots \times M_t$ 中，$M_i$ 是 $R_i$ 上秩为 $i$
的有限局部自由模。这由 Algebra，引理
\ref{algebra-lemma-finite-projective}（用以看出秩局部常值），以及
Algebra，引理 \ref{algebra-lemma-disjoint-decomposition} 与
\ref{algebra-lemma-disjoint-implies-product}（用以把 $R$ 分解为乘积）
得到。在这种情形，定义 $R$-模
$$
\det(M) = \wedge^0_{R_0}(M_0) \times \ldots \times \wedge^t_{R_t}(M_t)
$$
。这是秩为 $1$ 的有限局部自由模，因为每一项都是秩为 $1$ 的有限局部
自由模。若 $\varphi : M \to N$ 是有限投射 $R$-模的同构，则得到秩为
$1$ 的局部自由模的典范同构
$$
\det(\varphi) : \det(M) \longrightarrow \det(N)
$$
。更一般地，若对 $R$ 的所有素理想 $\mathfrak p$，自由模
$M_\mathfrak p$ 与 $N_\mathfrak p$ 的秩相同，则任意 $R$-模同态
$\varphi : M \to N$ 都诱导 $R$-模映射
$\det(\varphi) : \det(M) \to \det(N)$。最后，若 $M = N$，则
$\det(\varphi) : \det(M) \to \det(M)$ 是可逆 $R$-模的自同态。
由于对可逆 $R$-模有 $R = \Hom_R(L, L)$，可以并且确实把
$\det(\varphi)$ 看作 $R$ 的元素。这样便得到\textit{行列式}
$$
\det : \Hom_R(M, M) \longrightarrow R
$$
；这是乘法映射。

\begin{remark}
\label{more-algebra-remark-determinant-as-socle}
设 $R$ 是环，并设 $M$ 是有限投射 $R$-模。可以考虑 $R$ 上由 $M$
生成的分次交换 $R$-代数，即外代数 $\wedge^*_R(M)$。关于
$\det(M)$ 的一个公式是：$\det(M) \subset \wedge^*_R(M)$ 是
$M \subset \wedge^*_R(M)$ 的零化子。这个公式有时很有用，因为它不
涉及把 $R$ 分解成乘积。当然，要证明它具有所需性质，或者必须把 $R$
分解成乘积（如上），或者必须观察 $R$ 的各素理想处的局部化。
\end{remark}

\noindent
接下来考虑给定有限投射模的短正合列时，行列式会怎样。
% P02-E0690

\begin{lemma}
\label{more-algebra-lemma-det-ses}
设 $R$ 是环。设
$$
0 \to M' \to M \to M'' \to 0
$$
是有限投射 $R$-模的短正合列。则存在典范同构
$$
\gamma : \det(M') \otimes \det(M'') \longrightarrow \det(M)
$$
\end{lemma}

\begin{proof}[第一种证明]
第一种证明。把 $R$ 分解成环 $R_{ij}$ 的乘积，使得
$M' = \prod M'_{ij}$ 与 $M'' = \prod M''_{ij}$，其中 $M'_{ij}$
的秩为 $i$，而 $M''_{ij}$ 的秩为 $j$。当然，此时
$M = \prod M_{ij}$，并且 $M_{ij}$ 的秩为 $i + j$。这把我们约化到
$M'$ 与 $M''$ 具有常秩的情形，设其秩分别为 $i$ 与 $j$。在这种情形，
必须构造典范映射
$$
\wedge^i(M') \otimes \wedge^j(M'') \longrightarrow \wedge^{i + j}(M)
$$
。为此，选取 $M'$ 中的 $m'_1, \ldots, m'_i$ 与 $M''$ 中的
$m''_1, \ldots, m''_j$。以 $m_1, \ldots, m_i \in M$ 表示
$m'_1, \ldots, m'_i$ 的像，并以
$m_{i + 1}, \ldots , m_{i + j} \in M$ 表示映到 $M''$ 中
$m''_1, \ldots, m''_j$ 的元素。规则为
$$
m'_1 \wedge \ldots \wedge m'_i \otimes
m''_1 \wedge \ldots \wedge m''_j
\longmapsto
m_1 \wedge \ldots \wedge m_{i + j}
$$
。关于它良定义且是同构的详细证明从略。
\end{proof}

\begin{proof}[第二种证明]
我们将使用评注 \ref{more-algebra-remark-determinant-as-socle} 中给出的
$\det(M)$、$\det(M')$ 与 $\det(M'')$ 的描述。考虑 $R$-代数映射
$\wedge^*_R(M') \to \wedge^*_R(M)$ 与
$\wedge^*_R(M) \to \wedge^*_R(M'')$。第一个单射，第二个满射。取元素
$x' \in \det(M') \subset \wedge^*_R(M')$ 与元素
$x'' \in \det(M'') \subset \wedge^*_R(M'')$。选取映到 $x''$
的元素 $y'' \in \wedge^*(M)$，并令
$$
\gamma(x' \otimes x'') = x' \wedge y'' \in \det(M) \subset \wedge^*_R(M)
$$
。通过观察各素理想处的局部化，读者容易验证它良定义且是同构。
% P02-E0691
此外，这一构造给出与第一种证明中构造相同的映射。
\end{proof}

\begin{lemma}
\label{more-algebra-lemma-det-ses-functorial}
设 $R$ 是环。设
$$
\xymatrix{
0 \ar[r] &
M' \ar[r] \ar[d]^u &
M \ar[r] \ar[d]^v &
M'' \ar[r] \ar[d]^w &
0 \\
0 \ar[r] &
K' \ar[r] &
K \ar[r] &
K'' \ar[r] &
0
}
$$
是有限投射 $R$-模的交换图，其中竖直箭头都是同构。则得到同构的交换图
$$
\xymatrix{
\det(M') \otimes \det(M'') \ar[r]_-\gamma \ar[d]_{\det(u) \otimes \det(w)} &
\det(M) \ar[d]^{\det(v)} \\
\det(K') \otimes \det(K'') \ar[r]^-\gamma & \det(K)
}
$$
，其中水平箭头是引理 \ref{more-algebra-lemma-det-ses} 中构造的箭头。
\end{lemma}

\begin{proof}
从略。提示：使用引理 \ref{more-algebra-lemma-det-ses} 中映射 $\gamma$ 的第二种构造。
\end{proof}

\begin{lemma}
\label{more-algebra-lemma-det-filtration}
设 $R$ 是环。设
$$
K \subset L \subset M
$$
是 $R$-模，使得 $K$、$L/K$ 与 $M/L$ 是有限投射 $R$-模。则图
$$
\xymatrix{
\det(K) \otimes \det(L/K) \otimes \det(M/L) \ar[r] \ar[d] &
\det(L) \otimes \det(M/L) \ar[d] \\
\det(K) \otimes \det(M/K) \ar[r] &
\det(M)
}
$$
交换，其中映射是引理 \ref{more-algebra-lemma-det-ses} 中的映射。
\end{lemma}

\begin{proof}
从略。提示：在 $R$ 的某个素理想处局部化后，可以假设
$K \subset L \subset M$ 同构于
$R^{\oplus a} \subset R^{\oplus a + b}  \subset R^{\oplus a + b + c}$；
在这种情形，结论是显然的计算。
\end{proof}

\begin{lemma}
\label{more-algebra-lemma-det-direct-sum}
设 $R$ 是环，并设 $M'$ 与 $M''$ 是两个有限投射 $R$-模。则图
$$
\xymatrix{
\det(M') \otimes \det(M'') \ar[r]
\ar[d]_{\epsilon \cdot (\text{switch tensors})} &
\det(M' \oplus M'') \ar[d]^{\det(\text{switch summands})} \\
\det(M'') \otimes \det(M') \ar[r] &
\det(M'' \oplus M')
}
$$
交换，其中 $\epsilon = \det( -\text{id}_{M' \otimes M''}) \in R^*$，
水平箭头是引理 \ref{more-algebra-lemma-det-ses} 中的箭头。
\end{lemma}

\begin{proof}
从略。
\end{proof}

\begin{lemma}
\label{more-algebra-lemma-det-switch}
设 $R$ 是环，并设 $M$、$N$ 是有限投射 $R$-模。设
$a : M \to N$ 与 $b : N \to M$ 是 $R$-线性映射。则作为 $R$ 的元素，
有
$$
\det(\text{id} + a \circ b) = \det(\text{id} + b \circ a)
$$
。
\end{lemma}

\begin{proof}
只需在以素理想局部化的环替换 $R$ 后证明该断言。因此可假设 $R$
局部，而 $M$ 与 $N$ 有限自由。在这种情形，必须证明通常的矩阵行列式
之间的等式
$$
\det(I_n + AB) = \det(I_m + BA)
$$
，其中 $A$ 的大小为 $n \times m$，而 $B$ 的大小为 $m \times n$。
这约化到环
$R = \mathbf{Z}[a_{ij}, b_{ji}; 1 \leq i \leq n, 1 \leq j \leq m]$
的情形，其中
% P02-E0693
$a_{ij}$ 与 $b_{ij}$ 是变量以及矩阵 $A$ 与 $B$ 的元素。取分式域后，
这约化到特征为零的域的情形。在特征为零时，存在一个普遍多项式，
它以所述矩阵各次幂的迹表示大小为 $\leq N$ 的矩阵的行列式。因此只需对
所有 $k \geq 1$ 证明
% P02-E0692
$$
\text{Trace}((I_n + AB)^k) = \text{Trace}((I_m + BA)^k)
$$
。展开后可知，只需对所有 $k \geq 0$ 证明
$\text{Trace}((AB)^k) = \text{Trace}((BA)^k)$。当 $k = 1$ 时，
这是熟知的事实 $\text{Trace}(AB) = \text{Trace}(BA)$。当 $k > 1$
时，把 $(AB)^k = A(BA)^{k - 1}B$ 与
$(BA)^k = (BA)^{k - 1} A B$ 写出，便由该事实得到结论。
\end{proof}

\noindent
回顾我们在 Algebra，章节 \ref{algebra-section-K-groups} 中定义了群
$K_0(R)$：它是以有限投射 $R$-模的同构类为生成元的自由群，模去关系
$[M'] + [M''] = [M' \oplus M'']$。

\begin{lemma}
\label{more-algebra-lemma-det}
设 $R$ 是环。存在映射
$$
\det : K_0(R) \longrightarrow \Pic(R)
$$
；若 $M$ 是秩为 $n$ 的有限局部自由模，则它把 $[M]$ 映到可逆模
$\wedge^n(M)$ 的类。
\end{lemma}

\begin{proof}
这立即由上述构造得到；特别地，用引理 \ref{more-algebra-lemma-det-ses} 可看出这些
关系被映到 $0$。
\end{proof}

\section{完美复形与 K-群}
\label{more-algebra-section-perfect-K-group}

\noindent
我们迅速证明：环 $R$ 的完美复形导出范畴的第零 K-群，与 Algebra，
章节 \ref{algebra-section-K-groups} 中定义的 $K_0(R)$ 相同。

\begin{lemma}
\label{more-algebra-lemma-perfect-to-K-group}
设 $R$ 是环。存在映射
$$
c : \text{perfect complexes over }R \longrightarrow K_0(R)
$$
，它具有以下性质：
\begin{enumerate}
\item 对完美复形 $K$ 有 $c(K[n]) = (-1)^nc(K)$；
\item 若 $K \to L \to M \to K[1]$ 是完美复形的可区别三角，则
$c(L) = c(K) + c(M)$；
\item 若 $K$ 由有限投射模组成的有限复形 $M^\bullet$ 表示，则
% P02-E0694
$c(K) = \sum (-1)^i[M_i]$。
\end{enumerate}
\end{lemma}

\begin{proof}
设 $K$ 是 $D(R)$ 的完美对象。由定义，可用有限投射 $R$-模的有限复形
$M^\bullet$ 表示 $K$。通过令
$$
c(K) = \sum (-1)^n[M^n]
$$
在 $K_0(R)$ 中定义 $c$。当然，必须证明它良定义；但一旦良定义，(1)
与 (3) 就立即成立。暂时把映射 $c$ 看作定义在有限投射 $R$-模的复形上。

\medskip\noindent
假设 $L^\bullet \to M^\bullet$ 是有限投射 $R$-模的有限复形之间的
满射。设 $K^\bullet$ 为其核。于是得到 $R$-模的短正合列
$$
0 \to K^n \to L^n \to M^n \to 0
$$
；这些列分裂，因为 $M^n$ 投射。因此 $K^\bullet$ 也是有限投射
$R$-模的有限复形，并且在 $K_0(R)$ 中有
$c(L^\bullet) = c(K^\bullet) + c(M^\bullet)$。

\medskip\noindent
假设给定有限投射 $R$-模的非循环有限复形 $M^\bullet$。设当
$n \not \in [a, b]$ 时 $M^n = 0$。则可把 $M^\bullet$ 分解成短正合列
% P02-E0695
$$
\begin{matrix}
0 \to M^a \to M^{a + 1} \to N^{a + 1} \to 0, \\
0 \to N^{a + 1} \to M^{a + 2} \to N^{a + 3} \to 0, \\
\ldots \\
0 \to N^{b - 3} \to M^{b - 2} \to N^{b - 2} \to 0, \\
0 \to N^{b - 2} \to M^{b - 1} \to M^b \to 0
\end{matrix}
$$
。作降序归纳，可知 $N^{b - 2}, \ldots, N^{a + 1}$ 是有限投射
$R$-模，这些列分裂正合，并且
% P02-E0696
$$
c(M^\bullet) = \sum (-1)[M^n] = \sum (-1)^n([N^{n - 1}] + [N^n]) = 0
$$
。因此，我们的构造在非循环复形上给出零。

\medskip\noindent
由前两段的结果形式地可知，$c$ 良定义并满足 (2)。事实上，假设有限投射
$R$-模的有限复形 $M^\bullet$ 与 $L^\bullet$ 表示 $D(R)$ 的同一对象。
由 Derived Categories，引理
\ref{derived-lemma-morphisms-from-projective-complex}，可以用复形映射
$f : M^\bullet \to L^\bullet$ 表示该同构。由 Derived Categories，
定义 \ref{derived-definition-cone}，得到复形的短正合列
% P02-E0697
$$
0 \to L^\bullet \to C(f)^\bullet \to K^\bullet[1] \to 0
$$
。由于 $f$ 是拟同构，锥 $C(f)^\bullet$ 非循环（例如，这由 Derived
Categories，章节 \ref{derived-section-canonical-delta-functor} 中的讨论
得到）。因此
$$
0 = c(C(f)^\bullet) = c(L^\bullet) + c(K^\bullet[1]) =
c(L^\bullet) - c(K^\bullet)
$$
，如所需。类似的 (2) 之证明从略。
\end{proof}

\noindent
下面的引理表明 $K_0(R)$ 等于 $K_0(D_{perf}(R))$。

\begin{lemma}
\label{more-algebra-lemma-perfect-to-K-group-universal}
设 $R$ 是环。设 $D_{perf}(R)$ 是完美对象的导出范畴，见引理
\ref{more-algebra-lemma-perfect-ring-classical-generator}。引理
\ref{more-algebra-lemma-perfect-to-K-group} 的映射 $c$ 给出同构
$K_0(D_{perf}(R)) = K_0(R)$。
\end{lemma}

\begin{proof}
由 $K_0(D_{perf}(R))$ 的定义（Derived Categories，定义
\ref{derived-definition-K-zero}），$c$ 诱导同态
$K_0(D_{perf}(R)) \to K_0(R)$。

\medskip\noindent
给定 $R$ 上的有限投射模 $M$，以 $M[0]$ 表示如下 $R$ 上的完美复形：
$M$ 位于次数 $0$，其他次数均为零。给定有限投射模的短正合列
$0 \to M \to M' \to M'' \to 0$，得到可区别三角
$M[0] \to M'[0] \to M''[0] \to M[1]$，见 Derived Categories，
章节 \ref{derived-section-canonical-delta-functor}。这表明，把 $[M]$
映到 $[M[0]]$ 得到映射 $K_0(R) \to K_0(D_{perf}(R))$；请原谅这糟糕的
记号。

\medskip\noindent
显然，$K_0(R) \to K_0(D_{perf}(R)) \to K_0(R)$ 是恒等映射。另一方面，
若 $M^\bullet$ 是有限投射 $R$-模的有界复形，则“粗截断”的可区别三角
的存在性（见 Homology，章节 \ref{homology-section-truncations}）
% P02-E0698
$$
\sigma_{\geq n}M^\bullet \to \sigma_{\geq n - 1}M^\bullet \to
M^{n - 1}[-n + 1] \to (\sigma_{\geq n}M^\bullet)[1]
$$
与归纳表明
$$
[M^\bullet] = \sum (-1)^i[M^i[0]]
$$
在 $K_0(D_{perf}(R))$ 中成立（再次为该记号致歉）。因此映射
$K_0(R) \to K_0(D_{perf}(R))$ 满射，证明完成。
\end{proof}



\section{有限长度模的自同态的行列式}
\label{more-algebra-section-determinants-finite-length}

\noindent
设 $(R, \mathfrak m, \kappa)$ 是局部环。考虑如下二元组 $(M, \varphi)$
所成的范畴：前一分量是有限长度 $R$-模，而
$\varphi : M \to M$ 是一个自同态。该范畴是 Abel 范畴，并且每个对象
既是 Artin 的也是 Noether 的。相关定义见 Homology，章节
\ref{homology-section-jordan-holder}。

\medskip\noindent
若 $(M, \varphi)$ 是该范畴的单对象，则 $M$ 被 $\mathfrak m$ 零化；
否则
% P02-E0699
$(\mathfrak m M, \varphi|_{\mathfrak m M})$
会是一个非平凡子对象。此外，$\dim_\kappa(M) = \text{length}_R(M)$
是有限的。因此可以用线性代数把行列式和迹
$$
\det\nolimits_\kappa(\varphi),\quad \text{Trace}_\kappa(\varphi)
$$
定义为 $\kappa$ 的元素。在此情形下，$\varphi$ 的特征多项式也可类似定义。

\medskip\noindent
由 Homology，引理 \ref{homology-lemma-finite-length}，对上述范畴的任意对象
$(M, \varphi)$，存在由 $\varphi$-稳定子模组成的有限滤过
$$
0 \subset M_1 \subset \ldots \subset M_n = M
$$
，使得 $(M_i/M_{i - 1}, \varphi_i)$ 是该范畴的单对象；这里
$\varphi_i : M_i/M_{i - 1} \to M_i/M_{i - 1}$
是诱导映射。定义 $(M, \varphi)$ 在 $\kappa$ 上的{\it 行列式}为
$$
\det\nolimits_\kappa(\varphi) = \prod \det\nolimits_\kappa(\varphi_i)
$$
，其中 $\det_\kappa(\varphi_i)$ 如上一段所定义。定义
$(M, \varphi)$ 在 $\kappa$ 上的{\it 迹}为
$$
\text{Trace}_\kappa(\varphi) = \sum \text{Trace}_\kappa(\varphi_i)
$$
，其中 $\text{Trace}_\kappa(\varphi_i)$ 如上一段所定义。类似地，
可将 $\varphi$ 在 $\kappa$ 上的特征多项式定义为上一段所定义的
$\varphi_i$ 的特征多项式之积。由 Jordan--H\"older 定理
（Homology，引理 \ref{homology-lemma-jordan-holder}），这些定义良好。

\begin{lemma}
\label{more-algebra-lemma-ses}
设 $(R, \mathfrak m, \kappa)$ 是局部环。设
$0 \to (M, \varphi) \to (M', \varphi') \to (M'', \varphi'') \to 0$
是上述范畴中的短正合列。则
$$
\det\nolimits_\kappa(\varphi') =
\det\nolimits_\kappa(\varphi)\det\nolimits_\kappa(\varphi''),\quad
\text{Trace}_\kappa(\varphi') = \text{Trace}_\kappa(\varphi) + 
\text{Trace}_\kappa(\varphi'')
$$
。此外，$\varphi'$ 在 $\kappa$ 上的特征多项式等于
$\varphi$ 与 $\varphi''$ 的特征多项式之积。
\end{lemma}

\begin{proof}
留作练习。
\end{proof}

\begin{lemma}
\label{more-algebra-lemma-multiplication}
设 $(R, \mathfrak m, \kappa) \to (R', \mathfrak m', \kappa')$
是局部环的局部同态，并假设 $\kappa'/\kappa$ 是有限扩张。设
$u \in R'$。则对任意有限长度 $R'$-模 $M'$，有
$$
\det\nolimits_\kappa(u : M' \to M') =
\text{Norm}_{\kappa'/\kappa}(u \bmod \mathfrak m')^m
$$
，其中 $m = \text{length}_{R'}(M')$。
\end{lemma}

\begin{proof}
注意到
$\text{length}_R(M') = \text{length}_{R'}(M') [\kappa' : \kappa]$，
所以命题有意义。若 $M' = \kappa'$，则等式由范数的定义成立，因为范数
就是由乘以 $u$ 给出的线性算子的行列式。一般情形下，选取适当滤过并应用
引理 \ref{more-algebra-lemma-ses} 的乘法性，即可归约到这一情形。略去若干细节。
\end{proof}

\begin{lemma}
\label{more-algebra-lemma-flat-base-change-det}
设 $(R, \mathfrak m, \kappa) \to (R', \mathfrak m', \kappa')$
是局部环的平坦局部同态，并且
$m = \text{length}_{R'}(R'/\mathfrak mR') < \infty$。
对任意上述 $(M, \varphi)$，元素 $\det_\kappa(\varphi)^m$
在 $\kappa'$ 中的像等于
$\det_{\kappa'}(\varphi \otimes 1 : M \otimes_R R' \to M \otimes_R R')$。
\end{lemma}

\begin{proof}
$R \to R'$ 的平坦性保证引理 \ref{more-algebra-lemma-ses} 中的短正合列经基变换后
仍为 $R'$ 上的短正合列。因此，由引理 \ref{more-algebra-lemma-ses} 的乘法性，
可以假设 $(M, \varphi)$ 是本范畴的单对象（见本节引言）。
在单对象情形下，$M$ 被 $\mathfrak m$ 零化。选取滤过
$$
0 \subset I_1 \subset I_2 \subset \ldots \subset I_{m - 1} \subset
R'/\mathfrak mR'
$$
，使其各相继商作为 $R'$-模都同构于 $\kappa'$。于是得到滤过
$$
0 \subset
M \otimes_\kappa I_1 \subset
M \otimes_\kappa I_2 \subset
\ldots \subset
M \otimes_\kappa I_{m - 1} \subset
M \otimes_\kappa R'/\mathfrak mR' = M \otimes_R R'
$$
，其各相继商都同构于 $M \otimes_\kappa \kappa'$。并且这些子模在
$\varphi \otimes 1$ 下不变。由引理 \ref{more-algebra-lemma-ses} 得
$$
\det\nolimits_{\kappa'}(\varphi \otimes 1 : M \otimes_R R' \to M \otimes_R R')
=
\det\nolimits_{\kappa'}(\varphi \otimes 1 :
M \otimes_\kappa \kappa' \to M \otimes_\kappa \kappa')^m =
\det\nolimits_\kappa(\varphi)^m
$$
。最后一个等式来自线性映射的行列式与域扩张的相容性。这就证明了引理。
\end{proof}







\section{正则局部环是唯一分解整环}
\label{more-algebra-section-regular-local-UFD}

\noindent
我们证明节标题中所述的结果。

\begin{lemma}
\label{more-algebra-lemma-regular-local-Pic-zero}
设 $R$ 是正则局部环，且 $f \in R$。则 $\Pic(R_f) = 0$。
\end{lemma}

\begin{proof}
设 $L$ 是可逆 $R_f$-模。特别地，$L$ 是有限 $R_f$-模。存在有限
$R$-模 $M$ 使得 $M_f \cong L$，见 Algebra，引理
\ref{algebra-lemma-construct-fp-module}。由 Algebra，命题
\ref{algebra-proposition-regular-finite-gl-dim}，$M$ 在 $R$ 上有有限自由
消解 $F_\bullet$。由此 $L$ 拟同构于自由 $R_f$-模的有限复形。因此，
由引理 \ref{more-algebra-lemma-perfect-to-K-group}，对某个 $n \in \mathbf{Z}$，
% P02-E0700
$K_0(R_f)$ 中有 $[L_f] = n[R_f]$。应用引理 \ref{more-algebra-lemma-det} 的映射，
可知 $L$ 平凡。
\end{proof}

\begin{lemma}
\label{more-algebra-lemma-regular-local-UFD}
正则局部环是唯一分解整环。
\end{lemma}

\begin{proof}
回忆正则局部环是整环，见 Algebra，引理
\ref{algebra-lemma-regular-domain}。我们对正则局部环 $R$ 的维数归纳，
证明唯一分解性质。若 $\dim(R) = 0$，则 $R$ 是域，因而尤其是唯一分解
整环。假设 $\dim(R) > 0$。取
$x \in \mathfrak m$，$x \not \in \mathfrak m^2$。由 Algebra，引理
\ref{algebra-lemma-regular-ring-CM}，$R/(x)$ 正则；再由 Algebra，
引理 \ref{algebra-lemma-regular-domain}，它是整环，故 $x$ 是素元。
设 $\mathfrak p \subset R$ 是高为 $1$ 的素理想。我们必须证明
$\mathfrak p$ 是主理想，见 Algebra，引理
\ref{algebra-lemma-characterize-UFD-height-1}。可以假设
$x \not \in \mathfrak p$，因为若 $x \in \mathfrak p$，则
$\mathfrak p = (x)$，结论已得。
对每个非极大素理想 $\mathfrak q \subset R$，局部环
$R_\mathfrak q$ 都是正则局部环，见 Algebra，引理
\ref{algebra-lemma-localization-of-regular-local-is-regular}。
由归纳假设，$\mathfrak pR_\mathfrak q$ 是主理想。特别地，
$R_x$-模 $\mathfrak p_x = \mathfrak pR_x \subset R_x$ 是有限呈示 $R_x$-模，
并且在任意素理想处的局部化都是秩 $1$ 自由模。由 Algebra，引理
\ref{algebra-lemma-finite-projective}，$\mathfrak p_x$ 是可逆
$R_x$-模。由引理 \ref{more-algebra-lemma-regular-local-Pic-zero}，对某个
$y \in R_x$ 有 $\mathfrak p_x = (y)$。对某个
$f \in \mathfrak p$ 与 $e \in \mathbf{Z}$，可以写成
$y = x^e f$。把 $f = a_1 \ldots a_r$ 分解为 $R$ 的不可约元之积
（Algebra，引理 \ref{algebra-lemma-factorization-exists}）。
由于 $\mathfrak p$ 是素理想，对某个 $i$ 有
$a_i \in \mathfrak p$。由于 $\mathfrak p_x = (y)$ 是素理想，并且在
$R_x$ 中 $a_i | y$，故 $\mathfrak p_x$ 在 $R_x$ 中由 $a_i$ 生成，
即 $a_i$ 在 $R_x$ 中的像是素元。由于 $x$ 是素元，由 Algebra，引理
\ref{algebra-lemma-invert-prime-elements} 可知 $a_i$ 在 $R$ 中是素元。
又因 $(a_i) \subset \mathfrak p$ 且 $\mathfrak p$ 的高为 $1$，
所以 $(a_i) = \mathfrak p$，如所需。
\end{proof}

\begin{lemma}
\label{more-algebra-lemma-picard-group-generic-fibre-regular}
设 $R$ 是分式域为 $K$、剩余域为 $\kappa$ 的赋值环。设
$R \to A$ 是环同态，使得
\begin{enumerate}
\item $A$ 是局部环，且 $R \to A$ 是局部同态；
\item $A$ 在 $R$ 上平坦且本质有限型；
% P02-E0701
\item $A \otimes_R \kappa$ 正则。
\end{enumerate}
则 $\Pic(A \otimes_R K) = 0$。
\end{lemma}

\begin{proof}
设 $L$ 是可逆 $A \otimes_R K$-模。特别地，$L$ 是有限模。
存在有限 $A$-模 $M$ 使得 $M \otimes_R K \cong L$，见 Algebra，引理
\ref{algebra-lemma-construct-fp-module}。可以假设 $M$ 作为 $R$-模无挠。
因此 $M$ 作为 $R$-模平坦（引理
\ref{more-algebra-lemma-valuation-ring-torsion-free-flat}）。由引理
\ref{more-algebra-lemma-flat-finite-type-valuation-ring-finite-presentation}，
$M$ 作为 $A$-模是有限呈示的，并且 $A$ 作为 $R$-代数本质有限呈示。
由引理 \ref{more-algebra-lemma-structure-relatively-perfect}，$M$ 相对于 $R$ 是完美的；
特别地，$M$ 作为 $A$-模是伪凝聚的。由引理
\ref{more-algebra-lemma-perfect-over-regular-local-ring}，$M$ 是完美的，故
$M$ 在 $A$ 上有有限自由消解 $F_\bullet$。由此 $L$ 拟同构于自由
$A \otimes_R K$-模的有限复形。因此，由引理
\ref{more-algebra-lemma-perfect-to-K-group}，对某个 $n \in \mathbf{Z}$，
$K_0(A \otimes_R K)$ 中有
$[L] = n[A \otimes_R K]$。应用引理 \ref{more-algebra-lemma-det} 的映射，可知
$L$ 平凡。
\end{proof}
\section{复形的行列式}
\label{more-algebra-section-determinants-complexes}

\noindent
在章节 \ref{more-algebra-section-perfect-K-group} 中，我们已经看到，对环 $R$ 上的完美
复形 $K$，可以关联一个可逆 $R$-模的同构类，即 $\Pic(R)$ 的一个元素。
事实上，与章节 \ref{more-algebra-section-determinants} 类似，存在函子
$$
\det :
\left\{
\begin{matrix}
\text{category of perfect complexes} \\
\text{morphisms are isomorphisms}
\end{matrix}
\right\}
\longrightarrow
\left\{
\begin{matrix}
\text{category of invertible modules} \\
\text{morphisms are isomorphisms}
\end{matrix}
\right\}
$$
。此外，若 $(L, F)$ 是 $R$ 的滤过导出范畴 $DF(R)$ 中的对象，其滤过
有限且各分次部分都是完美复形，则存在典范同构
$\det(\text{gr}L) \to \det(L)$。原始论述见 \cite{determinant}。
% P02-E0702
我们将在以后补入这些材料（插入未来的引用）。

\medskip\noindent
眼下，我们对 $D(R)$ 中 tor-幅度位于 $[-1, 0]$ 的完美对象 $L$
给出一种特设构造。这样的对象可由复形
$$
L^\bullet = \ldots \to 0 \to L^{-1} \to L^0 \to 0 \to \ldots
$$
表示，其中 $L^{-1}$ 与 $L^0$ 是有限投射 $R$-模，见引理
\ref{more-algebra-lemma-perfect}。在此情形下，置
$$
\det(L^\bullet) = \det(L^0) \otimes_R \det(L^{-1})^{\otimes -1} =
\Hom_R(\det(L^{-1}), \det(L^0))
$$
。若这种形式的复形对 $R$ 的所有素理想都有
$L^{-1}_\mathfrak p$ 与 $L^0_\mathfrak p$ 秩相同，则称其{\it 秩为 $0$}。
若 $L^\bullet$ 的秩为 $0$，则章节 \ref{more-algebra-section-determinants} 已表明，
存在典范元素
$$
\delta(L^\bullet) \in \det(L^\bullet)
$$
；它就是 $d : L^{-1} \to L^0$ 的
% P02-E0703
行列式。注意，$\delta(L^\bullet)$ 是 $\det(L^\bullet)$ 的平凡化，
当且仅当 $L^\bullet$ 非循环。

\medskip\noindent
考虑复形映射 $a^\bullet : K^\bullet \to L^\bullet$，并假设
\begin{enumerate}
\item $a^\bullet$ 是拟同构；
\item 对所有 $n$，$a^n : K^n \to L^n$ 都满射；
\item $K^n$、$L^n$ 都是有限投射 $R$-模，并且仅当
$n \in \{-1, 0\}$ 时可能非零。
\end{enumerate}
在此情形下，我们将构造同构
$$
\det(a^\bullet) : \det(K^\bullet) \longrightarrow \det(L^\bullet)
$$
。利用正合列 $0 \to \Ker(a^i) \to K^i \to L^i \to 0$，
由引理 \ref{more-algebra-lemma-det-ses}，对 $i = -1, 0$ 得到同构
$$
\gamma^i : \det(\Ker(a^i)) \otimes \det(L^i) \to \det(K^i)
$$
。由于 $a^\bullet$ 是拟同构，复形 $\Ker(a^\bullet)$ 非循环且秩为 $0$。
因此，上述典范元素 $\delta(\Ker(a^\bullet))$ 是可逆 $R$-模
$\det(\Ker(a^\bullet))$ 的一个平凡化。把
$\det(a^\bullet) : \det(K^\bullet) \to \det(L^\bullet)$ 定义为使图表
$$
\xymatrix{
\det(K^\bullet)
\ar[rr]_{\det(a^\bullet)}
& &
\det(L^\bullet)
\ar[ld]^{\delta(\Ker(a^\bullet))}
\\
& \det(L^\bullet) \otimes \det(\Ker(a^\bullet))
\ar[lu]^{\gamma^0 \otimes (\gamma^{-1})^{\otimes -1}}
}
$$
交换的唯一同构。

\begin{lemma}
\label{more-algebra-lemma-canonical-element-well-defined}
设 $R$ 是环。设 $a^\bullet : K^\bullet \to L^\bullet$ 是满足上述
(1)、(2)、(3) 的 $R$-模复形映射。若 $L^\bullet$ 秩为 $0$，则
$\det(a^\bullet)$ 把典范元素 $\delta(K^\bullet)$ 映到
$\delta(L^\bullet)$。
\end{lemma}

\begin{proof}
记 $M^i = \Ker(a^i)$。于是有短正合列的映射
$$
\xymatrix{
0 \ar[r] &
M^{-1} \ar[r] \ar[d]_{d_M} &
K^{-1} \ar[r] \ar[d]_{d_K} &
L^{-1} \ar[r] \ar[d]_{d_L} &
0 \\
0 \ar[r] &
M^0 \ar[r] &
K^0 \ar[r] &
L^0 \ar[r] &
0
}
$$
。由引理 \ref{more-algebra-lemma-det-ses-functorial}，作为映射的
$\det(d_K)$ 对应于 $\det(d_M) \otimes \det(d_L)$。展开各定义即得所需
等式。
\end{proof}

\begin{lemma}
\label{more-algebra-lemma-homotopic-surjections}
设 $R$ 是环。设 $a^\bullet : K^\bullet \to L^\bullet$ 是满足上述
(1)、(2)、(3) 的 $R$-模复形映射。设 $h : K^0 \to L^{-1}$ 是映射，
使得 $b^0 = a^0 + d \circ h$ 与
$b^{-1} = a^{-1} + h \circ d$ 都满射。则作为
$\det(K^\bullet) \to \det(L^\bullet)$ 的映射，有
$\det(a^\bullet) = \det(b^\bullet)$。
\end{lemma}

\begin{proof}
假设存在映射 $\tilde h : K^0 \to K^{-1}$，使得
$h = a^{-1} \circ \tilde h$，并且
$k^0 = \text{id} + d \circ \tilde h : K^0 \to K^0$ 与
% P02-E0704
$k^1 = \text{id} + \tilde h \circ d : K^{-1} \to K^{-1}$ 都是同构。
于是得到复形的交换图
$$
\xymatrix{
0 \ar[r] &
\Ker(b^\bullet) \ar[r] \ar[d]_{c^\bullet} &
K^\bullet \ar[r]_{b^\bullet}
\ar[d]_{k^\bullet} &
L^\bullet \ar[r] \ar[d]^{\text{id}} &
0 \\
0 \ar[r] &
\Ker(a^\bullet) \ar[r] &
K^\bullet \ar[r]^{a^\bullet} &
L^\bullet \ar[r] &
0
}
$$
，其中 $c^\bullet$ 是诱导的核之间的同构。利用引理
\ref{more-algebra-lemma-det-ses-functorial} 可知
$$
\xymatrix{
\det(\Ker(b^i)) \otimes \det(L^i) \ar[r] \ar[d]_{\det(c^i) \otimes 1} &
\det(K^i) \ar[d]^{\det(k^i)} \\
\det(\Ker(a^i)) \otimes \det(L^i) \ar[r] &
\det(K^i)
}
$$
交换。由于
% P02-E0705
$\det(c^\bullet)$ 把 $\det(\Ker(a^\bullet))$ 的典范平凡化映到
$\Ker(b^\bullet)$ 的典范平凡化
（引理 \ref{more-algebra-lemma-canonical-element-well-defined}），可知当且仅当
% P02-E0706
$$
\det(k^0) = \det(k^{-1})
$$
在 $R$ 中成立时便可得到结论；而该等式由引理
\ref{more-algebra-lemma-det-switch} 得到。

\medskip\noindent
假设存在直和因子 $U \subset K^{-1}$，使得
$a^{-1}|_U : U \to L^{-1}$ 与 $b^{-1}|_U : U \to L^{-1}$
都是同构。把 $\tilde h$ 定义为 $h$ 与 $a^{-1}|_U$ 的逆之复合。
我们断言 $\tilde h$ 是证明第一段中所需的映射。事实上，由构造有
$h = a^{-1} \circ \tilde h$。
为证明 $k^{-1} : K^{-1} \to K^{-1}$ 是同构，只需证明它满射
（Algebra，引理 \ref{algebra-lemma-fun}）。设 $u \in U$。可以选取
$u' \in U$ 使得 $b^{-1}(u') = a^{-1}(u)$。于是
$u = k^{-1}(u')$。事实上，$u$ 与 $k^{-1}(u')$ 都在 $U$ 中，并且由计算
\footnote{$a^{-1}(k^{-1}(u')) =
a^{-1}(u') + a^{-1}(\tilde h(d(u'))) =
a^{-1}(u') + h(d(u')) = b^{-1}(u') = a^{-1}(u)$}
有 $a^{-1}(u) = a^{-1}(k^{-1}(u'))$。由于 $a^{-1}|_U$ 是同构，
得到该等式。因此 $U \subset \Im(k^{-1})$。另一方面，若
$x \in \Ker(a^{-1})$，则 $x = k^{-1}(x) \bmod U$。由于
$K^{-1} = \Ker(a^{-1}) + U$，可知 $k^{-1}$ 满射。
最后证明 $k^0 : K^0 \to K^0$ 满射。首先，由
$a^0 \circ k^0 = b^0$ 可知 $a^0 \circ k^0$ 满射。若
$x \in \Ker(a^0)$，则对某个 $y \in \Ker(a^{-1})$ 有 $x = d(y)$。
由上可对某个 $z \in K^{-1}$ 写成 $y = k^{-1}(z)$。于是
$x = k^0(d(z))$，结论得证。

\medskip\noindent
证明的最后一步。只需找到上一段中的 $U$，但这未必总能直接做到。不过，
为证明两个 $R$-模映射相等，只需在 $R$ 的各素理想处局部化后证明。
因此可以假设 $R$ 是局部环。于是问题变为：假设
$$
\alpha, \beta : R^{\oplus n} \longrightarrow R^{\oplus m}
$$
是两个满 $R$-线性映射，要找一个直和因子
$U \subset R^{\oplus n}$，使得 $\alpha|_U$ 与 $\beta|_U$
都是同构。若 $R$ 是域，这可由线性代数做到。一般地，先在剩余域上取一个
解，再把它提升为局部环 $R$ 上的解。略去若干细节。
\end{proof}
\begin{lemma}
\label{more-algebra-lemma-compose-surjections}
设 $R$ 是环。设 $a^\bullet : K^\bullet \to L^\bullet$ 与
$b^\bullet : L^\bullet \to M^\bullet$ 是满足上述 (1)、(2)、(3) 的
$R$-模复形映射。则
$\det(b^\bullet) \circ \det(a^\bullet) = 
\det(b^\bullet \circ a^\bullet)$；这里两边作为
% P02-E0707
$\det(M^\bullet) \to \det(K^\bullet)$ 的映射相等。
\end{lemma}

\begin{proof}
从略。提示：直接应用引理
\ref{more-algebra-lemma-det-ses}、\ref{more-algebra-lemma-det-ses-functorial}
与 \ref{more-algebra-lemma-det-filtration}。
\end{proof}

\begin{lemma}
\label{more-algebra-lemma-determinant-two-term-complexes}
设 $R$ 是环。上述构造确定函子
$$
\det :
\left\{
\begin{matrix}
\text{category of perfect complexes} \\
\text{with tor amplitude in }[-1, 0] \\
\text{morphisms are isomorphisms}
\end{matrix}
\right\}
\longrightarrow
\left\{
\begin{matrix}
\text{category of invertible modules} \\
\text{morphisms are isomorphisms}
\end{matrix}
\right\}
$$
。此外，给定 $D(R)$ 中秩为 $0$ 且 tor-幅度位于 $[-1, 0]$ 的完美对象
$L$，存在典范元素 $\delta(L) \in \det(L)$，使得对 $D(R)$ 中任意同构
$a : L \to K$ 都有 $\det(a)(\delta(L)) = \delta(K)$。
\end{lemma}

\begin{proof}
由引理 \ref{more-algebra-lemma-perfect}，源范畴的每个对象都可由复形
$$
L^\bullet = \ldots \to 0 \to L^{-1} \to L^0 \to 0 \to \ldots
$$
表示，其中 $L^{-1}$ 与 $L^0$ 是有限投射 $R$-模。暂称这种复形为良好
复形。由 Derived Categories，引理
\ref{derived-lemma-morphisms-from-projective-complex}，导出范畴中良好
复形之间的态射就是复形映射的链同伦类。因此，可以使用良好复形来工作，
并可使用上文研究的行列式
$\det(L^\bullet) = \det(L^0) \otimes \det(L^{-1})^{\otimes -1}$。

\medskip\noindent
设 $a^\bullet : L^\bullet \to K^\bullet$ 是良好复形之间的态射，并且
它在 $D(R)$ 中是同构，即为拟同构。若图表
$$
\xymatrix{
L^\bullet \ar[rr]_{a^\bullet} & & K^\bullet \\
& M^\bullet \ar[lu]^{b^\bullet} \ar[ru]_{c^\bullet}
}
$$
在链同伦意义下交换，且 $b^\bullet$ 与 $c^\bullet$ 满足上述
(1)、(2)、(3)，则称其为良好图表。一旦有这样的图表，就可以有意义地定义
$$
\det(a^\bullet) = \det(c^\bullet) \circ \det(b^\bullet)^{-1}
$$
，其中 $\det(c^\bullet)$ 与 $\det(b^\bullet)$ 是上文构造的同构。
我们将证明良好图表总是存在，并且所得映射 $\det(a^\bullet)$ 与良好图表
的选择无关。

\medskip\noindent
证明良好复形的拟同构 $a^\bullet : L^\bullet \to K^\bullet$
存在良好图表。选取满射 $p : R^{\oplus n} \to K^{-1}$。然后考虑新的
良好复形
$$
M^\bullet = \ldots \to 0 \to
L^{-1} \oplus R^{\oplus n} \xrightarrow{d \oplus 1}
L^0 \oplus R^{\oplus n} \to 0 \to \ldots
$$
，并取投影映射 $b^\bullet : M^\bullet \to L^\bullet$；映射
$c^\bullet : M^\bullet \to K^\bullet$ 在次数 $-1$ 使用
$a^{-1} \oplus p$，在次数 $0$ 使用 $a^0 \oplus d \circ p$。
映射 $b^\bullet : M^\bullet \to L^\bullet$ 与
$c^\bullet : M^\bullet \to K^\bullet$ 满足上述 (1)、(2)、(3)，
因而得到良好图表。

\medskip\noindent
假设有良好图表
$$
\xymatrix{
L^\bullet \ar[rr]_{\text{id}^\bullet} & & L^\bullet \\
& M^\bullet \ar[lu]^{b^\bullet} \ar[ru]_{c^\bullet}
}
$$
。由引理 \ref{more-algebra-lemma-homotopic-surjections}，
$\det(c^\bullet) = \det(b^\bullet)$。因此，
$\det(\text{id}^\bullet) = \text{id}$ 与良好图表的选择无关。

\medskip\noindent
在证明一般的无关性之前，先考虑复合。设有良好复形的拟同构
$L_1^\bullet \to L_2^\bullet$ 与 $L_2^\bullet \to L_3^\bullet$，
以及良好图表
$$
\vcenter{
\xymatrix{
L_1^\bullet \ar[rr] & &
L_2^\bullet \\
& M_{12}^\bullet \ar[lu] \ar[ru]
}
}
\quad\text{且}\quad
\vcenter{
\xymatrix{
L_2^\bullet \ar[rr] & &
L_3^\bullet \\
& M_{23}^\bullet \ar[lu] \ar[ru]
}
}
$$
。可以把它扩充为图表
$$
\xymatrix{
L_1^\bullet \ar[rr] & &
L_2^\bullet \ar[rr] & &
L_3^\bullet \\
& M_{12}^\bullet \ar[lu] \ar[ru]
& & M_{23}^\bullet \ar[lu] \ar[ru] \\
& &
M_{123}^\bullet \ar[lu] \ar[ru]
}
$$
，其中 $M_{123}^\bullet \to M_{12}^\bullet$ 与
$M_{123}^\bullet \to M_{23}^\bullet$ 具有性质 (1)、(2)、(3)，
且图中的方形交换：只需取
$M_{123}^n = M_{12}^n \times_{L_2^n} M_{23}^n$。于是引理
\ref{more-algebra-lemma-compose-surjections} 表明
$$
\xymatrix{
\det(L_2^\bullet) &
\det(M_{23}^\bullet) \ar[l] \\
\det(M_{12}^\bullet) \ar[u] &
\det(M_{123}^\bullet) \ar[l] \ar[u]
}
$$
交换。追图表明，由使用 $M_{12}^\bullet$ 与 $M_{23}^\bullet$ 的两个
良好图表所关联的映射之复合
$\det(L_1^\bullet) \to \det(L_2^\bullet) \to \det(L_3^\bullet)$，
等于良好图表
$$
\xymatrix{
L_1^\bullet \ar[rr] & & L_3^\bullet \\
& M_{123}^\bullet \ar[lu] \ar[ru]
}
$$
所关联的映射。因此，如果能够证明这些映射与选择无关，那么复合律也成立，
从而得到所需函子。

\medskip\noindent
无关性。给定良好复形的拟同构
$a^\bullet : L^\bullet \to K^\bullet$。选取逆拟同构
$b^\bullet : K^\bullet \to L^\bullet$。令
% P02-E0708
$L_1^\bullet = L$、$L_2^\bullet = K^\bullet$ 且
$L_3^\bullet = L^\bullet$，可以固定对 $b^\bullet$ 所选的良好图表，
并改变 $a^\bullet$ 的良好图表。于是前几段的结果表明，无论作何选择，
其复合总是 $\det(L^\bullet)$ 上的恒等映射。这显然证明了与这些选择无关。

\medskip\noindent
关于典范元素的陈述立即由引理
\ref{more-algebra-lemma-canonical-element-well-defined} 及我们的构造得到。
\end{proof}
\section{赋值环的扩张}
\label{more-algebra-section-valuation-rings}

\noindent
本节是章节 \ref{more-algebra-section-discrete-valuation-rings} 对一般赋值环的对应版本。

\begin{definition}
\label{more-algebra-definition-extension-valuation-rings}
若 $A$ 与 $B$ 都是赋值环，并且 $A \to B$ 是单射局部同态，
则称 $A \to B$ 或 $A \subset B$ 是{\it 赋值环的扩张}。
这样的扩张诱导交换图
$$
\xymatrix{
A \setminus \{0\} \ar[r] \ar[d]_v & B \setminus \{0\} \ar[d]^v \\
\Gamma_A \ar[r] & \Gamma_B
}
$$
，其中 $\Gamma_A$ 与 $\Gamma_B$ 是值群。若下方的水平箭头是双射，
则称 $B$ 在 $A$ 上{\it 弱非分歧}。若剩余域扩张
$\kappa_A = A/\mathfrak m_A \subset \kappa_B = B/\mathfrak m_B$
有限，则置 $f = [\kappa_B : \kappa_A]$，并称其为扩张
$A \subset B$ 的{\it 剩余次数}。
\end{definition}

\noindent
注意，$\Gamma_A \to \Gamma_B$ 是单射，因为在映射 $A \to B$ 下，
% P02-E0709
$A$ 的单位元是 $B$ 的单位元的逆。还要注意，我们并不要求分式域扩张有限。

\begin{lemma}
\label{more-algebra-lemma-inequality-general}
设 $A \subset B$ 是赋值环的扩张，分式域为 $K \subset L$。
若扩张 $L/K$ 有限，则剩余域扩张有限，$\Gamma_A$ 在 $\Gamma_B$
中的指数有限，并且
$$
[\Gamma_B : \Gamma_A] [\kappa_B : \kappa_A] \leq [L : K].
$$
\end{lemma}

\begin{proof}
设 $b_1, \ldots, b_n \in B$ 是单位元，并且它们在 $\kappa_B$ 中的像
在 $\kappa_A$ 上线性无关。设 $c_1, \ldots, c_m \in B$ 是非零元素，
并且它们在 $\Gamma_B/\Gamma_A$ 中的像两两不同。我们断言
$b_i c_j$ 在 $L$ 中于 $K$ 上线性无关。也就是说，断言一个
$a_{ij} \in K$ 不全为零的和
$$
\sum a_{ij} b_i c_j
$$
不可能为零。选取 $(i_0, j_0)$ 使
$v(a_{i_0j_0}b_{i_0}c_{j_0})$ 最小。把 $a_{ij}$ 替换为
$a_{ij}/a_{i_0j_0}$，从而 $a_{i_0 j_0} = 1$。令
$$
P = \{(i, j) \mid
v(a_{ij}b_ic_j) = v(a_{i_0j_0}b_{i_0}c_{j_0}) \}
$$
。由对 $c_1, \ldots, c_m$ 的选择，$(i, j) \in P$ 蕴含
$j = j_0$。因此若 $(i, j) \in P$，则
$v(a_{ij}) = v(a_{i_0j_0}) = 0$，即 $a_{ij}$ 是单位元。
由对 $b_1, \ldots, b_n$ 的选择，
$$
\sum\nolimits_{(i, j) \in P} a_{ij}b_i
$$
是 $B$ 中的单位元。因此
$\sum\nolimits_{(i, j) \in P} a_{ij}b_ic_j$ 的赋值为
$v(c_{j_0}) = v(a_{i_0j_0}b_{i_0}c_{j_0})$。第一个陈列和中
$(i, j) \not \in P$ 的各项具有严格更大的赋值，故该和不可能为零，
从而证明了引理。
\end{proof}

\begin{lemma}
\label{more-algebra-lemma-valuation-ring-purely-inseparable}
设 $A$ 是赋值环，其分式域 $K$ 的特征为 $p > 0$。设 $L/K$
是纯不可分扩张。则 $A$ 在 $L$ 中的整闭包 $B$ 是分式域为 $L$ 的赋值环，
并且 $A \subset B$ 是赋值环的扩张。
\end{lemma}

\begin{proof}
从略。提示：例如使用 Algebra，引理
\ref{algebra-lemma-x-or-x-inverse-valuation-ring} 与
\ref{algebra-lemma-integral-overring-surjective}。
\end{proof}

\begin{lemma}
\label{more-algebra-lemma-extension-normal-domains-and-roots}
设 $A \to B$ 是 Noether 局部正规整环的平坦局部同态。设
$f \in A$、$h \in B$，并且对某个 $n > 1$ 以及 $B$ 的某个单位元
$w$ 有 $f = w h^n$。假设对每个高为 $1$ 的素理想
$\mathfrak p \subset A$，都存在位于 $\mathfrak p$ 上方的高为 $1$
的素理想 $\mathfrak q \subset B$，使扩张
$A_\mathfrak p \subset B_\mathfrak q$ 弱非分歧。则对某个
$g \in A$ 及 $A$ 的某个单位元 $u$，有 $f = u g^n$。
\end{lemma}

\begin{proof}
$A$ 与 $B$ 在高为 $1$ 的素理想处的局部环都是离散赋值环
（Algebra，引理 \ref{algebra-lemma-characterize-dvr}）。因此该假设
是有意义的（通过定义
\ref{more-algebra-definition-extension-discrete-valuation-rings}）。
设 $\mathfrak p_1, \ldots, \mathfrak p_r$ 是 $A$ 中在 $f$ 上极小的
素理想。由 Algebra，引理 \ref{algebra-lemma-minimal-over-1}，
它们的高都是 $1$。对每个 $i$，设
$\mathfrak q_{i, j} \subset B$（$j = 1, \ldots, r_i$）是 $B$ 中
位于 $\mathfrak p_i$ 上方的高为 $1$ 的素理想。给它们编号，使得
$A_{\mathfrak p_i} \to B_{\mathfrak q_{i, 1}}$ 弱非分歧。
由于 $f$ 在 $B_{\mathfrak q_{i, 1}}$ 中映为一个 $n$ 次幂乘以单位元，
$f$ 在 $A_{\mathfrak p_i}$ 中的赋值 $v_i$ 被 $n$ 整除。写成
$v_i = n w_i$，其中 $w_i \geq 0$。考虑定义理想 $I$ 的正合列
$$
0 \to I \to A \to
\prod\nolimits_{i = 1, \ldots, r}
A_{\mathfrak p_i}/\mathfrak p_i^{w_i}A_{\mathfrak p_i}
$$
。应用正合函子 $- \otimes_A B$，得到正合列
$$
0 \to I \otimes_A B \to B \to
\prod\nolimits_{i = 1, \ldots, r}
(A_{\mathfrak p_i}/\mathfrak p_i^{w_i}A_{\mathfrak p_i}) \otimes_A B
$$
。固定 $i$。我们断言典范映射
$$
(A_{\mathfrak p_i}/\mathfrak p_i^{w_i}A_{\mathfrak p_i}) \otimes_A B
\to
\prod\nolimits_{j = 1, \ldots, r_i}
B_{\mathfrak q_{i, j}}/\mathfrak q_{i, j}^{e_{i, j}w_i}B_{\mathfrak q_{i, j}}
$$
是单射。这里 $e_{i, j}$ 是
$A_{\mathfrak p_i} \to B_{\mathfrak q_{i, j}}$ 的分歧指数。
该断言说，$\mathfrak p_i^{w_i}B_{\mathfrak p_i}$ 等于
$B_{\mathfrak p_i}$ 中在 $\mathfrak q_{i, j}$ 处的赋值
$\geq e_{i, j}w_i$ 的元素 $b$ 所成的集合。选取主理想
$\mathfrak p_i^{w_i}$ 的生成元
$a \in A_{\mathfrak p_i}$。则 $a$ 在 $\mathfrak q_{i, j}$ 处的赋值
等于 $e_{i, j}w_i$。由于 $B_{\mathfrak p_i}$ 是正规整环，其高为一的
素理想恰为 $\mathfrak q_{i, j}$（$j = 1, \ldots, r_i$），
对上述 $b$，由 Algebra，引理
\ref{algebra-lemma-normal-domain-intersection-localizations-height-1}
可知 $b/a \in B_{\mathfrak p_i}$。断言得证。

\medskip\noindent
该断言与上面的第二个正合列合在一起，确定正合列
$$
0 \to I \otimes_A B \to B \to
\prod\nolimits_{i = 1, \ldots, r}
\prod\nolimits_{j = 1, \ldots, r_i}
B_{\mathfrak q_{i, j}}/\mathfrak q_{i, j}^{e_{i, j}w_i}B_{\mathfrak q_{i, j}}
$$
。由此，$I \otimes_A B$ 就是 $B$ 中在
$\mathfrak q_{i, j}$ 处赋值 $\geq e_{i, j}w_i$ 的元素 $h'$
所成的集合。由于在 $B$ 中 $f = wh^n$，可知 $h$ 在
$\mathfrak q_{i, j}$ 处的赋值为 $e_{i, j}w_i$。于是由 Algebra，引理
\ref{algebra-lemma-normal-domain-intersection-localizations-height-1}，
$h'/h \in B$。因此 $I \otimes_A B$ 是秩为 $1$ 的自由 $B$-模
（由 $h$ 生成）。所以 $I$ 是秩为 $1$ 的自由 $A$-模，见 Algebra，
引理 \ref{algebra-lemma-finite-projective-descends}。设
$g \in I$ 是生成元。则 $g$ 与 $h$ 在 $B$ 中相差一个单位元。
倒推可知，$g$ 在 $A_{\mathfrak p_i}$ 中的赋值为
$w_i = v_i/n$。所以 $g^n$ 与 $f$ 在 $A$ 中相差一个单位元
（由 Algebra，引理
\ref{algebra-lemma-normal-domain-intersection-localizations-height-1}），
如所需。
\end{proof}

\begin{lemma}
\label{more-algebra-lemma-etale-extension-valuation-ring}
设 $A$ 是赋值环。设 $A \to B$ 是 \'etale 环映射，并设
$\mathfrak m \subset B$ 是位于 $A$ 的极大理想上方的素理想。
则 $A \subset B_\mathfrak m$ 是弱非分歧的赋值环扩张。
\end{lemma}

\begin{proof}
由引理 \ref{more-algebra-lemma-weak-dimension-at-most-1}，环 $A$ 的弱维数
$\leq 1$。再由引理 \ref{more-algebra-lemma-weak-dimension-goes-up} 与
\ref{more-algebra-lemma-when-weakly-etale}，$B$ 的弱维数 $\leq 1$。
% P02-E0710
故由引理 \ref{more-algebra-lemma-weak-dimension-at-most-1}，局部环
$B_\mathfrak m$ 是赋值环。由于扩张 $A \subset B_\mathfrak m$
诱导有限的分式域扩张，可知 $\Gamma_A$ 在 $B_{\mathfrak m}$
的值群中的指数有限。因此，对每个 $h \in B_\mathfrak m$，存在
$n > 0$、元素 $f \in A$ 以及单位元 $w \in B_\mathfrak m$，
使得在 $B_\mathfrak m$ 中 $f = w h^n$。我们将证明，这蕴含对某个
$g \in A$ 及单位元 $u \in A$ 有 $f = ug^n$；这将表明
$A$ 与 $B_\mathfrak m$ 的值群相同，从而得到引理的断言。

\medskip\noindent
把 $A = \colim A_i$ 写成它的各个在 $\mathbf{Z}$ 上本质有限型的局部
子环之余极限。由于 $A$ 是正规整环（Algebra，引理
\ref{algebra-lemma-valuation-ring-normal}），可以假设每个
$A_i$ 都正规（这里使用如下事实：取正规化后，局部环仍在
$\mathbf{Z}$ 上本质有限型；见 Algebra，命题
\ref{algebra-proposition-ubiquity-nagata}）。对某个 $i$，
可以找到 \'etale 扩张 $A_i \to B_i$，使得
$B = A \otimes_{A_i} B_i$，见 Algebra，引理
\ref{algebra-lemma-etale}。令 $\mathfrak m_i$ 为
$B_i$ 与 $\mathfrak m$ 的交。于是可以对环映射
$A_i \to (B_i)_{\mathfrak m_i}$ 应用引理
\ref{more-algebra-lemma-extension-normal-domains-and-roots} 得到结论。
该引理的假设成立，因为：
\begin{enumerate}
\item $A_i$ 与 $(B_i)_{\mathfrak m_i}$ 都是 Noether 的，因为它们在
$\mathbf{Z}$ 上本质有限型；
\item $A_i \to (B_i)_{\mathfrak m_i}$ 平坦，因为 $A_i \to B_i$
是 \'etale 的；
\item $B_i$ 正规，因为 $A_i \to B_i$ 是 \'etale 的，见 Algebra，
引理 \ref{algebra-lemma-normal-goes-up}；
\item 对 $A_i$ 的每个高为 $1$ 的素理想，都存在
$(B_i)_{\mathfrak m_i}$ 中位于其上方的高为 $1$ 的素理想；这由
Algebra，引理 \ref{algebra-lemma-finite-in-codim-1} 以及
$\Spec((B_i)_{\mathfrak m_i}) \to \Spec(A_i)$ 满射得到；
\item 对每个位于素理想 $\mathfrak p$ 上方的素理想 $\mathfrak q$，
诱导扩张 $(A_i)_\mathfrak p \to (B_i)_\mathfrak q$ 非分歧，因为
$A_i \to B_i$ 是 \'etale 的。
\end{enumerate}
这就完成了引理的证明。
\end{proof}

\begin{lemma}
\label{more-algebra-lemma-henselization-valuation-ring}
设 $A$ 是赋值环。设 $A^h$、$A^{sh}$ 分别是它的 Hensel 化与严格
Hensel 化。则
$$
A \subset A^h \subset A^{sh}
$$
是诱导值群双射的赋值环扩张，即弱非分歧扩张。
\end{lemma}

\begin{proof}
把 $A^h = \colim (B_i)_{\mathfrak q_i}$ 写成如下形式：$A \to B_i$
是 \'etale 的，而 $\mathfrak q_i \subset B_i$ 是位于
$\mathfrak m_A$ 上方的素理想；见 Algebra，引理
\ref{algebra-lemma-henselization-different}。引理
\ref{more-algebra-lemma-etale-extension-valuation-ring} 告诉我们，
$(B_i)_{\mathfrak q_i}$ 是赋值环，并且诱导映射
$$
(A \setminus \{0\})/A^* \longrightarrow
((B_i)_{\mathfrak q_i} \setminus \{0\}) / (B_i)_{\mathfrak q_i}^*
$$
是双射。由 Algebra，引理
\ref{algebra-lemma-colimit-valuation-rings}，$A^h$ 是赋值环。
而且，$(A \setminus \{0\})/A^* \to (A^h \setminus \{0\})/(A^h)^*$
也是双射。这证明了包含 $A \subset A^h$ 的情形。
为证明 $A \subset A^{sh}$，可以使用完全相同的论证，只需把
Algebra，引理 \ref{algebra-lemma-henselization-different} 换成
Algebra，引理 \ref{algebra-lemma-strict-henselization-different}。
由于 $A^{sh} = (A^h)^{sh}$，这也证明了包含
$A^h \subset A^{sh}$ 的所有断言。
\end{proof}
\section{主理想整环上模的结构}
\label{more-algebra-section-structure-modules}

\noindent
我们稍作推广（遵循 Warfield 的论文 \cite{Warfield-Purity} 与
\cite{Warfield-Decomposition}），使证明也适用于赋值环。

\begin{lemma}
\label{more-algebra-lemma-characterize-PD-modules}
\begin{reference}
\cite[Corollary 1]{Warfield-Purity}
\end{reference}
设 $P$ 是环 $R$ 上的模。下列条件等价：
\begin{enumerate}
\item $P$ 是形如 $R/fR$（其中 $f \in R$ 变化）的模之直和的直和因子。
\item 对每个满足如下条件的 $R$-模短正合列
$0 \to A \to B \to C \to 0$：对所有 $f \in R$ 都有
$fA = A \cap fB$，映射 $\Hom_R(P, B) \to \Hom_R(P, C)$ 满射。
\end{enumerate}
\end{lemma}

\begin{proof}
设 $0 \to A \to B \to C \to 0$ 是 (2) 中的正合列。为证明 (1) 蕴含
(2)，只需对每个 $f \in R$ 证明
$\Hom_R(R/fR, B) \to \Hom_R(R/fR, C)$ 满射。设
$\psi : R/fR \to C$ 是映射，并设 $\psi(1)$ 是 $b \in B$ 的像。
则 $fb \in A$。因此存在 $a \in A$ 使 $fa = fb$。于是
$f(b - a) = 0$，故得到把 $1$ 映到 $b - a$ 的态射
$\varphi : R/fR \to B$，它提升 $\psi$。

\medskip\noindent
反过来，假设 (2) 成立。设 $I$ 是所有二元组 $(f, \varphi)$ 所成的集合，
其中 $f \in R$ 且 $\varphi : R/fR \to P$。对 $i \in I$，以
$(f_i, \varphi_i)$ 表示对应的二元组。考虑映射
$$
B = \bigoplus\nolimits_{i \in I} R/f_iR \longrightarrow P
$$
，它把分量 $R/f_iR$ 中的元素 $r$ 映到 $P$ 中的
$\varphi_i(r)$。令 $A = \Ker(B \to P)$。那么，只要序列
$$
0 \to A \to B \to P \to 0
$$
是 (2) 中的正合列，(1) 就成立。为此，设 $f \in R$，且
$a \in A$ 在 $B$ 中等于 $f b$。写
$b = (r_i)_{i \in I}$，其中几乎所有 $r_i = 0$。于是 $P$ 中有
$$
f\sum \varphi_i(r_i) = 0
$$
。故存在 $i_0 \in I$，使得 $f_{i_0} = f$ 且
$\varphi_{i_0}(1) = \sum \varphi_i(r_i)$。设
$x_{i_0} \in R/f_{i_0}R$ 是 $1$ 的类。则
$$
a' = (r_i)_{i \in I} - (0, \ldots, 0, x_{i_0}, 0, \ldots )
$$
属于 $A$，并且 $fa' = a$，如所需。
\end{proof}

\begin{lemma}[广义赋值环]
\label{more-algebra-lemma-generalized-valuation-ring}
\begin{reference}
\cite{Warfield-Decomposition}
\end{reference}
设 $R$ 是非零环。下列条件等价：
\begin{enumerate}
\item 对 $a, b \in R$，或者 $a$ 整除 $b$，或者 $b$ 整除 $a$。
\item 每个有限生成理想都是主理想，且 $R$ 是局部环。
\item $R$ 的理想集在包含关系下全序。
\end{enumerate}
特别地，当 $R$ 是赋值环时这些条件成立。
\end{lemma}

\begin{proof}
假设 (2)，并取 $a, b \in R$。则 $(a, b) = (c)$。若 $c = 0$，
则 $a = b = 0$，故 $a$ 整除 $b$。假设 $c \not = 0$。写成
$c = ua + vb$、$a = wc$ 及 $b = zc$。于是
$c(1 - uw - vz) = 0$。由于 $R$ 是局部环，这蕴含
$1 - uw - vz \in \mathfrak m$。所以 $w$ 或 $z$ 中有一个是单位元，
从而 $a$ 整除 $b$ 或 $b$ 整除 $a$。故 (2) 蕴含 (1)。

\medskip\noindent
假设 (1)。若 $R$ 有两个极大理想 $\mathfrak m_i$，则可以选取
$a \in \mathfrak m_1$ 且 $a \not \in \mathfrak m_2$，以及
$b \in \mathfrak m_2$ 且 $b \not \in \mathfrak m_1$。此时
$a$ 不整除 $b$，而 $b$ 也不整除 $a$。因此 $R$ 有唯一极大理想，
从而是局部环。由条件 (1) 与归纳容易推出每个有限生成理想都是主理想。
所以 (1) 蕴含 (2)。

\medskip\noindent
(1) 与 (3) 等价的证明直接。最后一个陈述是 Algebra，引理
\ref{algebra-lemma-valuation-ring-x-or-x-inverse}。
\end{proof}

\begin{lemma}
\label{more-algebra-lemma-generalized-valuation-ring-modules}
\begin{reference}
\cite[Theorem 1]{Warfield-Decomposition}
\end{reference}
设环 $R$ 满足引理 \ref{more-algebra-lemma-generalized-valuation-ring} 中的等价条件。
则每个有限呈示 $R$-模都同构于有限个形如 $R/fR$ 的模之直和。
\end{lemma}

\begin{proof}
设 $M$ 是有限呈示 $R$-模。下文将不加说明地使用引理
\ref{more-algebra-lemma-generalized-valuation-ring} 中 $R$ 的全部等价性质。
记 $\mathfrak m \subset R$ 为极大理想，$\kappa = R/\mathfrak m$
为剩余域。设 $I \subset R$ 是 $M$ 的零化子。选取有限维
$\kappa$-向量空间 $M/\mathfrak m M$ 的一组基
$y_1, \ldots, y_n$。我们对 $n$ 归纳。

\medskip\noindent
由 Nakayama 引理，任意一组提升 $M/\mathfrak m M$ 中元素
$y_1, \ldots, y_n$ 的 $x_1, \ldots, x_n \in M$ 都生成 $M$，见
Algebra，引理 \ref{algebra-lemma-NAK}。这立即证明归纳的基础情形
$n = 0$。

\medskip\noindent
我们断言，存在指标 $i$，使得对任意映到 $y_i$ 的
$x_i \in M$，$x_i$ 的零化子都是 $I$。否则，可以选取
$x_1, \ldots, x_n$，使对所有 $i$ 都有
$I_i = \text{Ann}(x_i) \not = I$。但对所有 $i$ 都有
$I \subset I_i$；理想全序蕴含对 $i = 1, \ldots, n$，
$I_i$ 严格大于 $I$，再次使用全序便得到
$\text{Ann}(M) = I_1 \cap \ldots \cap I_n$ 大于 $I$，矛盾。
重新编号后，可以假设 $y_1$ 具有该性质：对任意提升 $y_1$ 的
$x_1 \in M$，$x_1$ 的零化子都是 $I$。

\medskip\noindent
置 $A = Rx_1 \subset M$。考虑正合列
$0 \to A \to M \to M/A \to 0$。由于 $A$ 有限，
$M/A$ 是生成元更少的有限呈示 $R$-模
（Algebra，引理 \ref{algebra-lemma-extension}）。因此由归纳，
$M/A \cong \bigoplus_{j = 1, \ldots, m} R/f_jR$。另一方面，我们断言
$A \to M$ 满足如下性质：若 $f \in R$，则
$fA = A \cap fM$。包含 $fA \subset A \cap fM$ 显然。反之，
若 $x \in A \cap fM$，则对某些 $g \in R$ 与 $y \in M$ 有
$x = gx_1 = f y$。若 $f$ 整除 $g$，则如所需有 $x \in fA$。
否则，可对某个 $h \in \mathfrak m$ 写成 $f = hg$。由上一段，
元素 $x'_1 = x_1 - hy$ 的零化子是 $I$。因此 $g \in I$，从而
$x = 0$，如所需。该断言与引理
\ref{more-algebra-lemma-characterize-PD-modules} 蕴含序列
$0 \to A \to M \to M/A \to 0$ 分裂，故
$M \cong A \oplus \bigoplus_{j = 1, \ldots, m} R/f_jR$。
于是 $A = R/I$ 是有限呈示的（因为它是 $M$ 的直和因子），所以
$I$ 有限生成，进而为主理想。证明完成。
\end{proof}

\begin{lemma}
\label{more-algebra-lemma-warfield}
\begin{reference}
\cite[Theorem 3]{Warfield-Decomposition}
\end{reference}
设 $R$ 是环，并且 $R$ 在每个极大理想处的局部环都满足引理
\ref{more-algebra-lemma-generalized-valuation-ring} 中的等价条件。则每个有限呈示
$R$-模都是某个有限直和的直和因子；该直和的各项形如
$R/fR$，其中 $f$ 在 $R$ 中变化。
\end{lemma}

\begin{proof}
设 $M$ 是有限呈示 $R$-模。先证明 $M$ 是形如 $R/fR$ 的模之直和的
直和因子，最后再证明该直和可以取为有限直和。设
$$
0 \to A \to B \to C \to 0
$$
是 $R$-模短正合列，并且对所有 $f \in R$ 都有
$fA = A \cap fB$。由引理 \ref{more-algebra-lemma-characterize-PD-modules}，
必须证明 $\Hom_R(M, B) \to \Hom_R(M, C)$ 满射。只需在各极大理想
$\mathfrak m$ 处局部化后证明这一点，见 Algebra，引理
\ref{algebra-lemma-characterize-zero-local}。注意，局部化序列
$0 \to A_\mathfrak m \to B_\mathfrak m \to C_\mathfrak m \to 0$
满足如下条件：对所有 $f \in R_\mathfrak m$，都有
$fA_\mathfrak m = A_\mathfrak m \cap fB_\mathfrak m$（因为可以写
$f = uf'$，其中 $u \in R_\mathfrak m$ 是单位元而 $f' \in R$，
并且局部化正合）。由于 $M$ 有限呈示，由 Algebra，引理
\ref{algebra-lemma-hom-from-finitely-presented}，
$$
\Hom_R(M, B)_\mathfrak m = \Hom_{R_\mathfrak m}(M_\mathfrak m, B_\mathfrak m)
\quad\text{且}\quad
\Hom_R(M, C)_\mathfrak m = \Hom_{R_\mathfrak m}(M_\mathfrak m, C_\mathfrak m)
$$
。模 $M_\mathfrak m$ 是有限呈示 $R_\mathfrak m$-模。由引理
\ref{more-algebra-lemma-generalized-valuation-ring-modules}，$M_\mathfrak m$
是形如 $R_\mathfrak m/fR_\mathfrak m$ 的模之直和。因此，由引理
\ref{more-algebra-lemma-characterize-PD-modules}，局部化后的映射满射。

\medskip\noindent
此时已知 $M$ 是 $\bigoplus_{i \in I} R/f_i R$ 的直和因子。
考虑映射 $M \to \bigoplus_{i \in I} R/f_i R$。由于 $M$ 是有限
$R$-模，对某个有限子集 $I' \subset I$，其像包含于
$\bigoplus_{i \in I'} R/f_i R$。证明完成。
\end{proof}

\begin{definition}
\label{more-algebra-definition-bezout}
设 $R$ 是整环。
\begin{enumerate}
\item 若 $R$ 的每个有限生成理想都是主理想，则称 $R$ 为
{\it B\'ezout 整环}。
\item 若对所有 $n , m \geq 1$ 以及每个 $n \times m$ 矩阵 $A$，
都存在大小分别为 $n \times n, m \times m$ 的可逆矩阵 $U, V$，使得
$$
U A V =
\left(
\begin{matrix}
f_1 & 0 & 0 & \ldots \\
0 & f_2 & 0 & \ldots \\
0 & 0 & f_3 & \ldots \\
\ldots & \ldots & \ldots & \ldots
\end{matrix}
\right)
$$
，其中 $f_1, \ldots, f_{\min(n, m)} \in R$ 且
$f_1 | f_2 | \ldots$，则称 $R$ 为{\it 初等因子整环}。
\end{enumerate}
\end{definition}

\noindent
每个 B\'ezout 整环 $R$ 是否都是初等因子整环，看来仍是一个未决问题。
这等价于问每个有限呈示 $R$-模是否都是循环模的直和。反向蕴含成立。

\begin{lemma}
\label{more-algebra-lemma-elementary-divisor-is-bezout}
初等因子整环是 B\'ezout 整环。
\end{lemma}

\begin{proof}
设 $a, b \in R$ 非零。考虑 $1 \times 2$ 矩阵 $A = (a\ b)$。
则 $u(a\ b)V = (f\ 0)$，其中 $u \in R$ 可逆，而
$V = (g_{ij})$ 是可逆 $2 \times 2$ 矩阵。于是
$f = u a g_{11} + u b g_{2 1}$ 且 $(g_{11}, g_{2 1}) = R$。
所以 $(a, b) = (f)$。归纳论证（从略）进而表明 $R$ 的每个有限生成
理想都由一个元素生成。
\end{proof}

\begin{lemma}
\label{more-algebra-lemma-localize-bezout}
B\'ezout 整环的局部化是 B\'ezout 整环。B\'ezout 整环的每个局部环
都是赋值环。局部整环是 B\'ezout 整环，当且仅当它是赋值环。
\end{lemma}

\begin{proof}
关于局部化的陈述之证明从略。最后一个陈述是 Algebra，引理
\ref{algebra-lemma-characterize-valuation-ring}。第二个陈述由另外两个推出。
\end{proof}

\begin{lemma}
\label{more-algebra-lemma-split-off-free-part}
设 $R$ 是 B\'ezout 整环。
\begin{enumerate}
\item 自由模的每个有限子模都是有限自由模。
\item 每个有限呈示 $R$-模 $M$ 都是一个有限自由模与一个挠模的直和，
其中该挠模 $M_{tors}$ 是形如
$\bigoplus_{i = 1, \ldots, n} R/f_iR$ 的模之直和因子，且
$f_1, \ldots, f_n \in R$ 非零。
\end{enumerate}
\end{lemma}

\begin{proof}
证明 (1)。设 $M \subset F$ 是自由模 $F$ 的有限子模。由于 $M$ 有限，
可以假设 $F$ 是有限自由模（细节从略）。设
$F = R^{\oplus n}$。对 $n$ 归纳。若 $n = 1$，则 $M$ 是有限生成理想，
因 $R$ 是 B\'ezout 整环而为主理想。若 $n > 1$，考虑 $M$ 在投影
$R^{\oplus n} \to R$（投到最后一个分量）下的像 $I$。
若 $I = (0)$，则 $M \subset R^{\oplus n - 1}$，由归纳得证。
若 $I \not = 0$，则 $I = (f) \cong R$。因此
$M \cong R \oplus \Ker(M \to I)$，再次由归纳得证。

\medskip\noindent
设 $M$ 是有限呈示 $R$-模。由于
% P02-E0711
$R$ 的局部化是极大理想处的赋值环（引理
\ref{more-algebra-lemma-localize-bezout}），可以应用引理 \ref{more-algebra-lemma-warfield}。
因此 $M$ 是形如
$R^{\oplus r} \oplus \bigoplus_{i = 1, \ldots, n} R/f_iR$
的模之直和因子，其中 $f_i \not = 0$。由于取挠子模是函子，可知
$M_{tors}$ 是 $\bigoplus_{i = 1, \ldots, n} R/f_iR$ 的直和因子，
而 $M/M_{tors}$ 是 $R^{\oplus r}$ 的直和因子。由证明的第一部分，
$M/M_{tors}$ 是有限自由模。因此
$M \cong M_{tors} \oplus M/M_{tors}$，如所需。
\end{proof}

\begin{lemma}
\label{more-algebra-lemma-modules-PID}
设 $R$ 是主理想整环。每个有限 $R$-模 $M$
% P02-E0712
都属于同构于形如
$$
R^{\oplus r} \oplus \bigoplus\nolimits_{i = 1, \ldots, n} R/f_iR
$$
的模，其中 $r, n \geq 0$，且 $f_1, \ldots, f_n \in R$ 非零。
\end{lemma}

\begin{proof}
主理想整环是 Noether B\'ezout 环。由引理
\ref{more-algebra-lemma-split-off-free-part}，只需在 $M$ 为挠模时证明结论。
由于 $M$ 有限，这意味着 $M$ 的零化子非零。设对某个非零
$f \in R$ 有 $fM = 0$。于是可把 $M$ 视为 $R/fR$-模。
由于 $R/fR$ 是维数为 $0$ 的 Noether 环（略去一个小细节），
% P02-E0713
$R/fR = \prod R_j$ 是 Artin 局部环 $R_i$ 的有限乘积
（Algebra，命题 \ref{algebra-proposition-dimension-zero-ring}）。
每个 $R_i$ 都是局部环及主理想整环的商，因而是引理
\ref{more-algebra-lemma-generalized-valuation-ring} 意义下的广义赋值环
（略去一个小细节）。写 $M = \prod M_j$，其中
$M_j = e_j M$，而 $e_j \in R/fR$ 是对应于因子 $R_j$ 的幂等元。
由引理 \ref{more-algebra-lemma-generalized-valuation-ring-modules}，
$M_j = \bigoplus_{i = 1, \ldots, n_j} R_j/\overline{f}_{ji}R_j$，
其中 $\overline{f}_{ji} \in R_j$。选取提升 $f_{ji} \in R$，并选取
% P02-E0714
$g_{ji} \in R$ 使 $(g_{ji}) = (f_j, f_{ji})$。于是作为 $R$-模，
$$
M \cong \bigoplus R/g_{ji}R
$$
，证明完成。
\end{proof}

\noindent
% P02-E0715
还可证明主理想整环是一个初等因子整环（在此插入未来的引用）；
一种方法是证明类似于下列引理的结果。

\begin{lemma}
\label{more-algebra-lemma-unimodular-vector}
设 $R$ 是 B\'ezout 整环。设 $n \geq 1$，并设
$f_1, \ldots, f_n \in R$ 生成单位理想。则存在 $R$ 中的可逆
$n \times n$ 矩阵，其第一行为 $f_1 \ldots f_n$。
\end{lemma}

\begin{proof}
这由引理 \ref{more-algebra-lemma-split-off-free-part} 得到；也可如下直接证明。
对 $n$ 归纳。$n = 1$ 时结论成立。假设 $n > 1$。重新编号后可以假设
$f_1 \not = 0$。选取 $f \in R$，使
$(f) = (f_1, \ldots, f_{n - 1})$。设 $A$ 是
$(n - 1) \times (n - 1)$ 矩阵，其第一行为
$f_1/f, \ldots, f_{n - 1}/f$。选取 $a, b \in R$，使
$af - bf_n = 1$；这是可能的，因为
$1 \in (f_1, \ldots, f_n) = (f, f_n)$。则一个解为矩阵
$$
\left(
\begin{matrix}
f & 0 & \ldots & 0 & f_n \\
0 & 1 & \ldots & 0 & 0 \\
  &   & \ldots \\
0 & 0 & \ldots & 1 & 0 \\
b & 0 & \ldots & 0 & a
\end{matrix}
\right)
\left(
\begin{matrix}
  & & & 0 \\
  & A \\
  & & & 0 \\
0 & \ldots & 0 & 1
\end{matrix}
\right)
$$
。注意，左侧矩阵的行列式为 $1$，故可逆。
\end{proof}
\section{主根理想}
\label{more-algebra-section-principal-radical-ideals}

\noindent
% P02-E0716
本节证明，一个链环 Noether 局部正规整环存在非平凡的主根理想。
该结果可见于 \cite{Artin-Lipman}。

\begin{lemma}
\label{more-algebra-lemma-polypoly}
设 $(R,\mathfrak m)$ 是一维 Noether 局部环，并设
$x\in\mathfrak m$ 不属于 $R$ 的任何极小素理想。则
\begin{enumerate}
\item 函数 $P : n \mapsto \text{length}_R(R/x^n R)$ 对
$n \geq 0$ 满足 $P(n) \leq n P(1)$；
\item 若 $x$ 是非零因子，则对 $n \geq 0$ 有 $P(n) = nP(1)$。
\end{enumerate}
\end{lemma}

\begin{proof}
由于 $\dim(R) = 1$，有 $\dim(R/x^n R) = 0$，故对每个 $n$，
$\text{length}_R(R/x^n R)$ 都有限
（Algebra，引理 \ref{algebra-lemma-support-point}）。
我们对 $n$ 归纳以证明引理。由于 $x^0 R = R$，有
$P(0) = \text{length}_R(R/x^0R) = \text{length}_R 0 = 0$。
$n = 1$ 时陈述也成立。现设 $n \geq 2$，并假设陈述对 $n - 1$
成立。下列序列正合：
$$
R/x^{n-1}R \xrightarrow{x} R/x^nR \to R/xR \to 0
$$
，其中 $x$ 表示乘以 $x$ 的映射。由于长度可加
（Algebra，引理 \ref{algebra-lemma-length-additive}），有
$P(n) \leq P(n - 1) + P(1)$。由归纳假设，
$P(n - 1) \leq (n - 1)P(1)$，所以 $P(n) \leq nP(1)$。
这证明了归纳步骤。

\medskip\noindent
若 $x$ 是非零因子，则上面的陈列正合列在左端也正合。因此对所有
$n \geq 1$，有 $P(n) = P(n - 1) + P(1)$。
\end{proof}

\begin{lemma}
\label{more-algebra-lemma-minprimespoly}
设 $(R, \mathfrak m)$ 是维数为 $1$ 的 Noether 局部环。设
$x \in \mathfrak m$ 不属于 $R$ 的任何极小素理想。设 $t$ 是 $R$
的极小素理想数目。则 $t \leq \text{length}_R(R/xR)$。
\end{lemma}

\begin{proof}
设 $\mathfrak p_1, \ldots, \mathfrak p_t$ 是 $R$ 的极小素理想。置
$R' = R/\sqrt{0} = R/(\bigcap_{i = 1}^t \mathfrak p_i)$。
我们断言，只需对 $R'$ 证明引理。事实上，显然 $R'$ 也有 $t$ 个极小
素理想，并且
$\text{length}_{R'}(R'/xR') = \text{length}_R(R'/xR')$
小于 $\text{length}_R(R/xR)$，因为存在满射
$R/xR \to R'/xR'$。因此可以假设 $R$ 既约。

\medskip\noindent
假设 $R$ 既约，其极小素理想为
$\mathfrak p_1, \ldots, \mathfrak p_t$。这意味着存在正合列
$$
0 \to R \to
\prod\nolimits_{i = 1}^t R/\mathfrak p_i \to Q \to 0
$$
。这里 $Q$ 是第一个映射的余核。写
$M = \prod_{i = 1}^t R/\mathfrak p_i$。在 $\mathfrak p_j$ 处局部化，
可见
$$
R_{\mathfrak p_j} \to M_{\mathfrak p_j} =
\left(\prod\nolimits_{i=1}^t R/\mathfrak p_i\right)_{\mathfrak p_j} =
(R/\mathfrak p_j)_{\mathfrak p_j}
$$
满射。因此对所有 $j$ 都有 $Q_{\mathfrak p_j} = 0$。由于
$\mathfrak m$ 是 $R$ 中唯一不同于各 $\mathfrak p_i$ 的素理想，
得到 $\text{Supp}(Q) = \{\mathfrak m\}$。所以 $Q$ 长度有限
（Algebra，引理 \ref{algebra-lemma-support-point}）。由于
$\text{Supp}(Q) = \{\mathfrak m\}$，可以选取 $n \gg 0$，使
$x^n$ 在 $Q$ 上作用为 $0$
（Algebra，引理
\ref{algebra-lemma-Noetherian-power-ideal-kills-module}）。
现考虑图表
$$
\xymatrix{
0 \ar[r] & R \ar[r] \ar[d]^-{x^n} & M
\ar[r] \ar[d]^-{x^n} & Q \ar[r] \ar[d]^-{x^n} & 0 \\
0 \ar[r] & R \ar[r] & M \ar[r] & Q \ar[r] & 0
}
$$
，其中竖直映射均为乘以 $x^n$。因为 $x$ 不属于任何
$\mathfrak p_i$，该映射在 $R$ 与 $M$ 上都是单射。由蛇引理
（Algebra，引理 \ref{algebra-lemma-snake}），下列序列正合：
$$
0 \to Q \to R/x^nR \to M/x^nM \to Q \to 0
$$
。所以当 $n$ 充分大时，
$\text{length}_R(R/x^nR) = \text{length}_R(M/x^nM)$。
写 $R_i = R/\mathfrak p_i$，则
$\text{length}(M/x^nM) =
\sum_{i = 1}^t \text{length}_R(R_i/x^nR_i)$。
应用引理 \ref{more-algebra-lemma-polypoly}，并利用 $x$ 在 $R$ 及 $R_i$ 上都是
非零因子，得到
$$
n \text{length}_R(R/xR) =
\sum\nolimits_{i = 1}^t n \text{length}_{R_i}(R_i/x R_i)
$$
。由于 $\text{length}_{R_i}(R_i/x R_i) \geq 1$，引理得证。
\end{proof}

\begin{lemma}
\label{more-algebra-lemma-sopexists}
设 $(R,\mathfrak m)$ 是维数 $d > 1$ 的 Noether 局部环，
设 $f \in \mathfrak m$ 不属于 $R$ 的任何极小素理想，并设
$k\in\mathbf{N}$。则存在
$g_1, \ldots, g_{d - 1} \in \mathfrak m^k$，使得
$f, g_1, \ldots, g_{d - 1}$ 是一组参数系。
\end{lemma}

\begin{proof}
由 Algebra，引理 \ref{algebra-lemma-one-equation}，
$\dim(R/fR) = d - 1$。在 $R/fR$ 中选取参数系
$\overline{g}_1, \ldots, \overline{g}_{d - 1}$
（Algebra，命题 \ref{algebra-proposition-dimension}），并在 $R$ 中
取提升 $g_1, \ldots, g_{d - 1}$。容易看出
$f, g_1, \ldots, g_{d - 1}$ 是 $R$ 的参数系。
于是 $f, g_1^k, \ldots, g_{d - 1}^k$ 也是参数系，证明完成。
\end{proof}

\begin{lemma}
\label{more-algebra-lemma-syspar}
设 $(R,\mathfrak m)$ 是二维 Noether 局部环，并设
$f \in \mathfrak m$ 不属于 $R$ 的任何极小素理想。则存在
$g \in \mathfrak m$ 与 $N \in \mathbf{N}$，使得
\begin{enumerate}
\item[(a)] $f,g$ 构成 $R$ 的一组参数系。
\item[(b)] 若 $h \in \mathfrak m^N$，则 $f + h, g$ 是参数系，并且
$\text{length}_R (R/(f, g)) = \text{length}_R(R/(f + h, g))$。
\end{enumerate}
\end{lemma}

\begin{proof}
由引理 \ref{more-algebra-lemma-sopexists}，存在 $g \in \mathfrak m$，使
$f, g$ 是 $R$ 的参数系。于是
$\mathfrak m = \sqrt{(f, g)}$。因此存在 $n$，使
$\mathfrak m^n \subset (f, g)$，见 Algebra，引理
\ref{algebra-lemma-Noetherian-power}。我们断言 $N = n + 1$ 可用。
事实上，设 $h \in \mathfrak m^N$。由 $N$ 的选择，可以写
$h = af + bg$，其中 $a, b \in \mathfrak m$。于是
$$
(f + h, g) = (f + af + bg, g) = ((1 + a)f, g) = (f, g)
$$
，因为 $1 + a$ 是 $R$ 中的单位元。这证明了长度相等，并证明
$f + h, g$ 是参数系。
\end{proof}

\begin{lemma}
\label{more-algebra-lemma-radical-element}
设 $R$ 是维数为 $2$ 的 Noether 局部正规整环。设
$\mathfrak p_1, \ldots, \mathfrak p_r$ 是两两不同的高为 $1$
的素理想。则存在非零元素
$f \in \mathfrak p_1 \cap \ldots \cap \mathfrak p_r$，使得
$R/fR$ 既约。
\end{lemma}

\begin{proof}
取非零元素
$f \in \mathfrak p_1 \cap \ldots \cap \mathfrak p_r$。
我们将稍微修改 $f$，得到生成根理想的元素。$R$ 在每个高为一
的素理想 $\mathfrak p$ 处的局部化 $R_\mathfrak p$ 都是离散赋值环，
见 Algebra，引理 \ref{algebra-lemma-characterize-dvr} 或
Algebra，引理 \ref{algebra-lemma-criterion-normal}。记
$\text{ord}_\mathfrak p(f)$ 为 $f$ 在 $R_{\mathfrak p}$ 中的相应
赋值。设 $\mathfrak q_1, \ldots, \mathfrak q_s$ 是包含 $f$
的所有不同的高为一的素理想。对每个 $j$，写
$\text{ord}_{\mathfrak q_j}(f) = m_j \geq 1$。定义
$\text{div}(f) = \sum_{j = 1}^s m_j\mathfrak q_j$，把它看作高为一
的素理想以整数为系数的形式线性组合。留意以后会用到：每个
$\mathfrak p_i$ 都出现在某个 $\mathfrak q_j$ 中。
$R/fR$ 既约，当且仅当对 $j = 1, \ldots, s$ 都有 $m_j = 1$。
事实上，若 $m_j$ 为 $1$，则
% P02-E0717
$(R/fR)\mathfrak q_j$ 既约，并且
$R/fR \subset \prod (R/fR)_{\mathfrak q_j}$，因为
% P02-E0718
$\mathfrak q_1, \ldots, \mathfrak q_j$ 是 $R/fR$ 的伴随素理想；
见 Algebra，引理 \ref{algebra-lemma-zero-at-ass-zero} 与
\ref{algebra-lemma-normal-domain-intersection-localizations-height-1}。

\medskip\noindent
选定引理 \ref{more-algebra-lemma-syspar} 中的 $g$ 与 $N$。对非零
$y \in R$，以 $t(y)$ 表示 $y$ 上的极小素理想数目。由于 $R$
是正规整环，这些素理想的高为一，并且与 $R/yR$ 的极小素理想形成
$1$ 对 $1$ 的对应
（Algebra，引理 \ref{algebra-lemma-minimal-over-1} 与
\ref{algebra-lemma-normal-domain-intersection-localizations-height-1}）。
例如，$t(f) = s$ 就是出现在 $\text{div}(f)$ 中的素理想
$\mathfrak q_j$ 的数目。设 $h \in \mathfrak m^N$。由引理
\ref{more-algebra-lemma-minprimespoly}，
\begin{align*}
t(f + h) & \leq \text{length}_{R/(f + h)}(R/(f + h, g)) \\
& = \text{length}_R(R/(f + h, g)) \\
& = \text{length}_R(R/(f, g))
\end{align*}
；第一个等式见 Algebra，引理
\ref{algebra-lemma-length-independent}。因此 $t(f + h)$ 有一个与
$h \in \mathfrak m^N$ 无关的上界。

\medskip\noindent
由上述有界性，可以选取
$h \in \mathfrak m^N \cap \mathfrak p_1 \cap \ldots \cap \mathfrak p_r$，
使 $t(f + h)$ 在这样的 $h$ 中达到最大值。置 $f' = f + h$。
给定
$h' \in \mathfrak m^N \cap \mathfrak p_1 \cap \ldots \cap \mathfrak p_r$，
有 $t(f' + h') \leq t(f + h)$。因此，把 $f$ 换成 $f'$ 后，可以假设
对每个
$h \in \mathfrak m^N \cap \mathfrak p_1 \cap \ldots \cap \mathfrak p_r$
都有 $t(f + h) \leq s$。

\medskip\noindent
接着，假设能找到 $h \in \mathfrak m^N$，使对每个 $j$ 都有
$\text{ord}_{\mathfrak q_j}(h) \geq 1$，并且
$\text{ord}_{\mathfrak q_j}(h) = 1 \Leftrightarrow m_j > 1$。
注意，
$h \in \mathfrak m^N \cap \mathfrak p_1 \cap \ldots \cap \mathfrak p_r$。
由赋值的初等性质，对每个 $j$ 都有
$\text{ord}_{\mathfrak q_j}(f + h) = 1$。因此
$$
\text{div}(f + h) = \sum\nolimits_{j = 1}^s \mathfrak q_j +
\sum\nolimits_{k = 1}^v e_k \mathfrak r_k
$$
，其中 $\mathfrak r_1, \ldots, \mathfrak r_v$ 是两两不同的高为一
的素理想，且 $e_k \geq 1$。但是
$s = t(f) \geq t(f + h)$，所以 $v = 0$，我们找到了所需元素。

\medskip\noindent
现在选取满足上述条件的 $h$。由素理想回避
（Algebra，引理 \ref{algebra-lemma-silly}），对每个
$1 \leq j \leq s$，可以找到 $a_j \in \mathfrak q_j$，使
$a_j \not \in \mathfrak q_{j'}$（$j' \not = j$），且
$a_j \not \in \mathfrak q_j^{(2)}$。这里
$\mathfrak q_j^{(2)} = \{x \in R \mid \text{ord}_{\mathfrak q_j}(x) \geq 2\}$
是 $\mathfrak q_j$ 的第二符号幂。然后取
$$
h = \prod\nolimits_{m_j = 1} a_j^2 \times
\prod\nolimits_{m_j > 1} a_j
$$
。显然 $h$ 满足上面对赋值施加的条件。若
$h \not \in \mathfrak m^N$，则把它乘以 $\mathfrak m^N$ 中一个对所有
$j$ 都不属于 $\mathfrak q_j$ 的元素。
\end{proof}

\begin{lemma}
\label{more-algebra-lemma-divides-radical}
设 $(A, \mathfrak m, \kappa)$ 是维数为 $2$ 的 Noether 局部正规整环。
若 $a \in \mathfrak m$ 非零，则存在 $c \in A$，使得 $A/cA$
既约，并且对某个 $n$，$a$ 整除 $c^n$。
\end{lemma}

\begin{proof}
沿用引理 \ref{more-algebra-lemma-radical-element} 证明中的记号，写
$\text{div}(a) = \sum_{i = 1, \ldots, r} n_i \mathfrak p_i$。
由引理 \ref{more-algebra-lemma-radical-element}，选取
$c \in \mathfrak p_1 \cap \ldots \cap \mathfrak p_r$，使 $A/cA$
既约。对 $n \geq \max(n_i)$，
$-\text{div}(a) + \text{div}(c^n)$ 是有效除子（所有系数非负）。
因此，由 Algebra，引理
\ref{algebra-lemma-normal-domain-intersection-localizations-height-1}，
$c^n/a \in A$。
\end{proof}

\noindent
本节余下部分证明维数 $> 2$ 时的结论。

\begin{lemma}
\label{more-algebra-lemma-multiplicity}
设 $(R, \mathfrak m)$ 是维数为 $d$ 的 Noether 局部环，
$g_1, \ldots, g_d$ 是参数系，并令 $I = (g_1, \ldots, g_d)$。
若数值多项式
$n \mapsto \text{length}_R(R/I^{n+1})$ 的首项系数为 $e_I/d!$，则
$e_I \leq \text{length}_R(R/I)$。
\end{lemma}

\begin{proof}
由 Algebra，命题
\ref{algebra-proposition-hilbert-function-polynomial}，该函数是数值多项式。
由 Algebra，命题 \ref{algebra-proposition-dimension}，它的次数为 $d$。
若 $d = 0$，结论显然。若 $d = 1$，结论就是引理
\ref{more-algebra-lemma-polypoly}。为证明一般情形，注意存在满射
$$
\bigoplus\nolimits_{i_1, \ldots, i_d \geq 0,\ \sum i_j = n} R/I
\longrightarrow
I^n/I^{n + 1}
$$
，它把对应于 $i_1, \ldots, i_d$ 的基元素送到
$I^n/I^{n + 1}$ 中 $g_1^{i_1} \ldots g_d^{i_d}$ 的类。因此
$$
\text{length}_R(R/I^{n + 1}) - \text{length}_R(R/I^n)
\leq  \text{length}_R(R/I) {n + d - 1 \choose d - 1}
$$
。由于 $d \geq 2$，左边的数值多项式次数为 $d - 1$，首项系数为
$e_I / (d - 1)!$。右边的多项式次数为 $d - 1$，首项系数为
$\text{length}_R(R/I)/ (d - 1)!$。这就证明了引理。
\end{proof}

\begin{lemma}
\label{more-algebra-lemma-minprimespolyhigher}
设 $(R, \mathfrak m)$ 是维数为 $d$ 的 Noether 局部环，$t$ 是 $R$
中维数为 $d$ 的极小素理想数目，并设 $(g_1,\ldots,g_d)$ 是参数系。则
% P02-E0719
$t \leq \text{length}_R(R/(g_1,\ldots,g_n))$。
\end{lemma}

\begin{proof}
若 $d = 0$，引理显然。若 $d = 1$，引理就是引理
\ref{more-algebra-lemma-minprimespoly}。因此可以假设 $d > 1$。设
$\mathfrak p_1, \ldots, \mathfrak p_s$ 是 $R$ 的极小素理想，其中前
$t$ 个的维数为 $d$，并记
% P02-E0719
$I = (g_1, \ldots, g_n)$。完全仿照引理
\ref{more-algebra-lemma-minprimespoly} 的证明，可以假设 $R$ 既约。

\medskip\noindent
% P02-E0720
假设 $R$ 既约，其极小素理想为
$\mathfrak p_1, \ldots, \mathfrak p_t$。这意味着存在正合列
$$
0 \to R \to
\prod\nolimits_{i = 1}^t R/\mathfrak p_i \to Q \to 0
$$
。这里 $Q$ 是第一个映射的余核。记
$M = \prod_{i = 1}^t R/\mathfrak p_i$。在 $\mathfrak p_j$ 处局部化，得到
$$
R_{\mathfrak p_j} \to M_{\mathfrak p_j} =
\left(\prod\nolimits_{i=1}^t R/\mathfrak p_i\right)_{\mathfrak p_j} =
(R/\mathfrak p_j)_{\mathfrak p_j}
$$
满射。因此对所有 $j$ 都有 $Q_{\mathfrak p_j} = 0$。所以 $R$ 的高为
$0$ 的素理想都不在 $Q$ 的支集内。由此，数值多项式
$n \mapsto \text{length}_R(Q/I^nQ)$ 的次数等于
$\dim(\text{Supp}(Q)) < d$；见 Algebra，引理
\ref{algebra-lemma-support-dimension-d}。由 Algebra，引理
\ref{algebra-lemma-hilbert-ses-chi}（它适用是因为 $R$ 不是有限长度），
多项式
$$
n \longmapsto
\text{length}_R(M/I^nM) - \text{length}_R(R/I^n) - \text{length}_R(Q/I^nQ)
$$
的次数 $< d$。由于 $M = \prod R/\mathfrak p_i$，并且
% P02-E0721
$n \to \text{length}_R(R/\mathfrak p_i + I^n)$ 对
$i = 1, \ldots, t$ 是次数恰好为(!) $d$ 的数值多项式（由 Algebra，引理
\ref{algebra-lemma-support-dimension-d}），可见
$n \mapsto \text{length}_R(M/I^nM)$ 的首项系数至少为 $t/d!$。
于是由引理 \ref{more-algebra-lemma-multiplicity} 得到结论。
\end{proof}

\begin{lemma}
\label{more-algebra-lemma-sysparhigher}
设 $(R, \mathfrak m)$ 是维数为 $d$ 的 Noether 局部环，且
$f \in \mathfrak m$ 不属于 $R$ 的任何极小素理想。则存在元素
$g_1, \ldots, g_{d - 1} \in \mathfrak m$ 与 $N \in \mathbf{N}$，使得
\begin{enumerate}
\item $f, g_1, \ldots, g_{d - 1}$ 构成 $R$ 的参数系；
\item 若 $h \in \mathfrak m^N$，则
$f + h, g_1, \ldots, g_{d - 1}$ 是参数系，并且
$\text{length}_R R/(f, g_1, \ldots, g_{d-1}) =
\text{length}_R R/(f + h, g_1, \ldots, g_{d-1})$。
\end{enumerate}
\end{lemma}

\begin{proof}
由引理 \ref{more-algebra-lemma-sopexists}，存在
$g_1, \ldots, g_{d - 1} \in \mathfrak m$，使
$f, g_1, \ldots, g_{d - 1}$ 是 $R$ 的参数系。于是
$\mathfrak m = \sqrt{(f, g_1, \ldots, g_{d - 1})}$。因此存在 $n$，使
% P02-E0722
$\mathfrak m^n \subset (f, g)$；见 Algebra，引理
\ref{algebra-lemma-Noetherian-power}。我们断言 $N = n + 1$ 可用。
事实上，设 $h \in \mathfrak m^N$。由 $N$ 的选择，可以写
$h = af + \sum b_ig_i$，其中 $a, b_i \in \mathfrak m$。于是
\begin{align*}
(f + h, g_1, \ldots, g_{d - 1})
& =
(f + af + \sum b_ig_i, g_1, \ldots, g_{d - 1}) \\
& =
((1 + a)f, g_1, \ldots, g_{d - 1}) \\
& =
(f, g_1, \ldots, g_{d - 1})
\end{align*}
，因为 $1 + a$ 是 $R$ 中的单位元。这证明了长度相等，并证明
$f + h, g_1, \ldots, g_{d - 1}$ 是参数系。
\end{proof}

\begin{proposition}
\label{more-algebra-proposition-propdimd}
\begin{reference}
\cite[Lemma 3.14]{Artin-Lipman} 在不假设环是链环的情形下给出了这一结果
\end{reference}
设 $R$ 是链环 Noether 局部正规整环。设 $J \subset R$ 是根理想。
则存在非零元素 $f \in J$，使得 $R/fR$ 既约。
\end{proposition}

\begin{proof}
证明与引理 \ref{more-algebra-lemma-radical-element} 的证明相同，只需以引理
\ref{more-algebra-lemma-minprimespolyhigher} 代替引理 \ref{more-algebra-lemma-minprimespoly}，并以
引理 \ref{more-algebra-lemma-sysparhigher} 代替引理 \ref{more-algebra-lemma-syspar}。
可以使用引理 \ref{more-algebra-lemma-minprimespolyhigher}，因为 $R$ 是链整环，所以
$R$ 中每个高为一的素理想的维数都是 $d - 1$，从而
$R/(f + h)$ 的谱是等维的。为方便读者，下面写出细节。

\medskip\noindent
取非零元素 $f \in J$。我们将稍微修改 $f$，得到生成根理想的元素。
$R$ 在每个高为一的素理想 $\mathfrak p$ 处的局部化
$R_\mathfrak p$ 都是离散赋值环；见 Algebra，引理
\ref{algebra-lemma-characterize-dvr} 或 Algebra，引理
\ref{algebra-lemma-criterion-normal}。记
$\text{ord}_\mathfrak p(f)$ 为 $f$ 在 $R_{\mathfrak p}$ 中的相应赋值。
设 $\mathfrak q_1, \ldots, \mathfrak q_s$ 是包含 $f$ 的所有不同的
高为一的素理想。对每个 $j$，写
$\text{ord}_{\mathfrak q_j}(f) = m_j \geq 1$。定义
$\text{div}(f) = \sum_{j = 1}^s m_j\mathfrak q_j$，把它看作高为一
的素理想以整数为系数的形式线性组合。$R/fR$ 既约，当且仅当
对 $j = 1, \ldots, s$ 都有 $m_j = 1$。事实上，若 $m_j$ 为 $1$，则
% P02-E0717
$(R/fR)\mathfrak q_j$ 既约，并且
$R/fR \subset \prod (R/fR)_{\mathfrak q_j}$，因为
% P02-E0718
$\mathfrak q_1, \ldots, \mathfrak q_j$ 是 $R/fR$ 的伴随素理想；见
Algebra，引理 \ref{algebra-lemma-zero-at-ass-zero} 与
\ref{algebra-lemma-normal-domain-intersection-localizations-height-1}。

\medskip\noindent
% P02-E0723
选定引理 \ref{more-algebra-lemma-sysparhigher} 中的
$g_2, \ldots, g_{d - 1}$ 与 $N$。对非零 $y \in R$，以 $t(y)$ 表示
$y$ 上的极小素理想数目。由于 $R$ 是正规整环，这些素理想的高为一，
并且与 $R/yR$ 的极小素理想形成 $1$ 对 $1$ 的对应（Algebra，引理
\ref{algebra-lemma-minimal-over-1} 与
\ref{algebra-lemma-normal-domain-intersection-localizations-height-1}）。
例如，$t(f) = s$ 就是出现在 $\text{div}(f)$ 中的素理想
$\mathfrak q_j$ 的数目。设 $h \in \mathfrak m^N$。因为 $R$ 是链环，
对 $R$ 的每个高为一的素理想 $\mathfrak p$，都有
% P02-E0724
$\dim(R/\mathfrak p) = d$。因此由引理
\ref{more-algebra-lemma-minprimespolyhigher}，
\begin{align*}
t(f + h) & \leq \text{length}_{R/(f + h)}(R/(f + h, g_1, \ldots, g_{d - 1})) \\
& = \text{length}_R(R/(f + h, g_1, \ldots, g_{d - 1})) \\
& = \text{length}_R(R/(f, g_1, \ldots, g_{d - 1}))
\end{align*}
；第一个等式见 Algebra，引理
\ref{algebra-lemma-length-independent}。因此 $t(f + h)$ 有一个与
$h \in \mathfrak m^N$ 无关的上界。

\medskip\noindent
由上述有界性，可以选取 $h \in \mathfrak m^N \cap J$，使
$t(f + h)$ 在这样的 $h$ 中达到最大值。置 $f' = f + h$。
给定 $h' \in \mathfrak m^N \cap J$，有
$t(f' + h') \leq t(f + h)$。因此，把 $f$ 换成 $f'$ 后，可以假设
对每个 $h \in \mathfrak m^N \cap J$ 都有 $t(f + h) \leq s$。

\medskip\noindent
接着，假设能找到 $h \in \mathfrak m^N \cap J$，使对每个 $j$ 都有
$\text{ord}_{\mathfrak q_j}(h) \geq 1$，并且
$\text{ord}_{\mathfrak q_j}(h) = 1 \Leftrightarrow m_j > 1$。
由赋值的初等性质，对每个 $j$ 都有
$\text{ord}_{\mathfrak q_j}(f + h) = 1$。因此
$$
\text{div}(f + h) = \sum\nolimits_{j = 1}^s \mathfrak q_j +
\sum\nolimits_{k = 1}^v e_k \mathfrak r_k
$$
，其中 $\mathfrak r_1, \ldots, \mathfrak r_v$ 是两两不同的高为一
的素理想，且 $e_k \geq 1$。但是
$s = t(f) \geq t(f + h)$，所以 $v = 0$，我们找到了所需元素。

\medskip\noindent
现在选取满足上述条件的 $h$。由素理想回避（Algebra，引理
\ref{algebra-lemma-silly}），对每个 $1 \leq j \leq s$，可以找到
$a_j \in \mathfrak q_j \cap J$，使
$a_j \not \in \mathfrak q_{j'}$（$j' \not = j$）。接着可以选取
% P02-E0725
$b_j \in J \cap \mathfrak q_1 \cap \ldots \cap q_s$，使
$b_j \not \in \mathfrak q_j^{(2)}$。这里
$\mathfrak q_j^{(2)} = \{x \in R \mid \text{ord}_{\mathfrak q_j}(x) \geq 2\}$
是 $\mathfrak q_j$ 的第二符号幂。素理想回避可以应用，因为理想
% P02-E0725
$J' = J \cap \mathfrak q_1 \cap \ldots \cap q_s$ 是根理想，故
$R/J'$ 既约，故 $(R/J')_{\mathfrak q_j}$ 既约，故 $J'$ 包含一个元素
$x$，满足 $\text{ord}_{\mathfrak q_j}(x) = 1$；因此
$J' \not \subset \mathfrak q_j^{(2)}$。于是元素
$$
c = \sum\nolimits_{j = 1, \ldots, s}
b_j \times \prod\nolimits_{j' \not = j} a_{j'}
$$
属于 $J$，且由赋值的初等性质，对所有 $j = 1, \ldots, s$ 都有
$\text{ord}_{\mathfrak q_j}(c) = 1$。最后令
$$
h = c \times \prod\nolimits_{m_j = 1} a_j \times y
$$
，其中 $y \in \mathfrak m^N$ 是对所有 $j$ 都不属于 $\mathfrak q_j$
的元素。
\end{proof}

\section{导出范畴中的可逆对象}
\label{more-algebra-section-invertible-D-or-R}

\noindent
我们刻画环的导出范畴中的可逆对象。

\begin{lemma}
\label{more-algebra-lemma-symmetric-monoidal-derived}
设 $R$ 是环。$R$ 的导出范畴 $D(R)$ 是对称幺半范畴；其张量积由导出
张量积给出，结合约束与交换约束如第 \ref{more-algebra-section-sign-rules} 节所述。
\end{lemma}

\begin{proof}
略。提示：结合约束是引理 \ref{more-algebra-lemma-triple-tensor-product} 的同构，交换
约束是引理 \ref{more-algebra-lemma-flip-douoble-tensor-product} 的同构。说明这些以后，
各图表的交换性由 $R$-模复形范畴的相应结果得到；见第
\ref{more-algebra-section-symmetric-monoidal} 节。
\end{proof}

\noindent
因此，我们知道 $D(R)$ 的对象有（左）对偶或可逆是什么意思。在具体确定
这等价于什么之前，需要一个简单引理。

\begin{lemma}
\label{more-algebra-lemma-complex-bounded-above-free-colim-bounded-finite-free}
设 $R$ 是环，$F^\bullet$ 是自由 $R$-模组成的上有界复形。给定成对数据
$(n_i, f_i)$，$i = 1, \ldots, N$，其中 $n_i \in \mathbf{Z}$ 且
$f_i \in F^{n_i}$，则存在包含所有 $f_i$ 的子复形
$G^\bullet \subset F^\bullet$，它有界且由有限自由 $R$-模组成。
\end{lemma}

\begin{proof}
对 $a = \min(n_i; i = 1, \ldots, N)$ 作降归纳。若对 $n \geq a$ 都有
$F^n = 0$，则取 $G^\bullet$ 为零复形即可。一般而言，重新编号后，可以
假设存在
% P02-E0726
$1 \leq r \leq N$，使 $n_1 = \ldots = n_r = a$，并且对 $i > r$
有 $n_i > a$。为 $F^a$ 选择一组基 $b_j, j \in J$。可以选取有限子集
$J' \subset J$，使对 $i = 1, \ldots, r$ 都有
$f_i \in \bigoplus_{j \in J'} Rb_j$。为 $F^{a + 1}$ 选择一组基
$c_k, k \in K$。可以选取有限子集 $K' \subset K$，使对 $j \in J'$
都有 $\text{d}_F^a(b_j) \in \bigoplus_{k \in K'} Rc_k$。于是可以应用
归纳假设，找到子复形 $H^\bullet \subset F^\bullet$，它包含
$c_k \in F^{a + 1}$（$k \in K'$）以及 $f_i \in F^{n_i}$（$i > r$）。
令 $G^\bullet$ 在次数 $> a$ 时等于 $H^\bullet$，在次数 $a$ 时等于
$\bigoplus_{j \in J'} Rb_j$。
\end{proof}

\begin{lemma}
\label{more-algebra-lemma-have-dual-derived}
设 $R$ 是环，$M$ 是 $D(R)$ 的对象。以下条件等价：
\begin{enumerate}
\item $M$ 在 $D(R)$ 中有 Categories，定义
\ref{categories-definition-dual} 意义下的左对偶；
\item $M$ 是 $D(R)$ 的完美对象。
\end{enumerate}
此外，在这种情形下，$M$ 的左对偶是引理
\ref{more-algebra-lemma-dual-perfect-complex} 的对象 $M^\vee$。
\end{lemma}

\begin{proof}
若 $M$ 完美，则可以用有限投射 $R$-模组成的有界复形 $M^\bullet$ 表示
$M$。此时，由引理 \ref{more-algebra-lemma-left-dual-complex}，$M^\bullet$ 在复形
范畴中有左对偶，因而它当然也是 $D(R)$ 中的左对偶。

\medskip\noindent
假设 (1)。设 $N$，
$\eta : R \to M \otimes_R^\mathbf{L} N$，以及
% P02-E0727
$\epsilon : M \otimes_R^\mathbf{L} N \to R$ 构成 Categories，定义
\ref{categories-definition-dual} 意义下的左对偶。选择表示 $M$ 的复形
$M^\bullet$。选择一个各项平坦、表示 $N$ 的
% P02-E0728
K-平坦复形 $N^\bullet$；见引理 \ref{more-algebra-lemma-K-flat-resolution}。于是
$\eta$ 由复形映射
$$
\eta : R \longrightarrow \text{Tot}(M^\bullet \otimes_R N^\bullet)
$$
给出。可以把 $1$ 的像写成有限和
$$
\eta(1) = \sum\nolimits_n \sum\nolimits_i m_{n, i} \otimes n_{-n, i}
$$
，其中 $m_{n, i} \in M^n$ 且 $n_{-n, i} \in N^{-n}$。令
$K^\bullet \subset M^\bullet$ 为所有元素 $m_{n, i}$ 与
$\text{d}(m_{n, i})$ 生成的子复形。由 $N^\bullet$ 的选择，有
$\text{Tot}(K^\bullet \otimes_R N^\bullet) \subset
\text{Tot}(M^\bullet \otimes_R N^\bullet)$，而根据上面的选择，
$\eta(1)$ 位于这个子复形中。以 $K$ 表示 $K^\bullet$ 所表示的
$D(R)$ 的对象。于是 $\eta$ 通过映射
$\tilde \eta : R \longrightarrow K \otimes_R^\mathbf{L} N$ 分解。
由于
$(1 \otimes \epsilon) \circ (\eta \otimes 1) = \text{id}_M$，
由交换图表
$$
\xymatrix{
M \ar[rr]_-{\eta \otimes 1} \ar[rrd]_{\tilde \eta \otimes 1} & &
M \otimes_R^\mathbf{L} N \otimes_R^\mathbf{L} M
\ar[r]_-{1 \otimes \epsilon} &
M \\
& &
K \otimes_R^\mathbf{L} N \otimes_R^\mathbf{L} M
\ar[u] \ar[r]^-{1 \otimes \epsilon} &
K \ar[u]
}
$$
可知 $M$ 上的恒等映射通过 $K$ 分解。由于 $K$ 上有界，得到
$M \in D^-(R)$。因此，可以用自由 $R$-模组成的上有界复形
$M^\bullet$ 表示 $M$；例如见 Derived Categories，引理
\ref{derived-lemma-subcategory-left-resolution}。像以前一样写
$\eta(1) = \sum\nolimits_n \sum\nolimits_i m_{n, i} \otimes n_{-n, i}$。
由引理 \ref{more-algebra-lemma-complex-bounded-above-free-colim-bounded-finite-free}，
可以找到包含所有元素 $m_{n, i}$ 的子复形
$K^\bullet \subset M^\bullet$，它有界且由有限自由 $R$-模组成。
与上面一样，$M$ 上的恒等映射通过 $K$ 分解。由于 $K$ 完美，得到
$M$ 也完美；见引理 \ref{more-algebra-lemma-summands-perfect}。
\end{proof}

\begin{lemma}
\label{more-algebra-lemma-invertible-derived}
设 $R$ 是环，$M$ 是 $D(R)$ 的对象。以下条件等价：
\begin{enumerate}
\item $M$ 在 $D(R)$ 中可逆；见 Categories，定义
\ref{categories-definition-invertible}；
\item 对每个素理想 $\mathfrak p \subset R$，存在
$f \in R$，$f \not \in \mathfrak p$，使对某个 $n \in \mathbf{Z}$
有 $M_f \cong R_f[-n]$。
\end{enumerate}
此外，在这种情形下：
\begin{enumerate}
\item[(a)] $M$ 是 $D(R)$ 的完美对象；
\item[(b)] 在 $D(R)$ 中，$M = \bigoplus H^n(M)[-n]$；
\item[(c)] 每个 $H^n(M)$ 都是有限投射 $R$-模；
\item[(d)] 可以写成 $R = \prod_{a \leq n \leq b} R_n$，使
$H^n(M)$ 对应于可逆 $R_n$-模。
\end{enumerate}
\end{lemma}

\begin{proof}
假设 (2)。考虑对象 $R\Hom_R(M, R)$ 与复合映射
$$
R\Hom(M, R) \otimes_R^\mathbf{L} M \to R
$$
。在局部检查可知这是同构；细节从略。由于 $D(R)$ 是对称幺半范畴，
$M$ 可逆。

\medskip\noindent
假设 (1)。注意，幺半范畴中的可逆对象有左对偶，即它的逆。因此由引理
\ref{more-algebra-lemma-have-dual-derived}，$M$ 完美。考虑素理想
$\mathfrak p \subset R$，其剩余域为 $\kappa$。于是
$M \otimes_R^\mathbf{L} \kappa$ 是 $D(\kappa)$ 的可逆对象。显然，这
意味着 $\dim H^i(M \otimes_R^\mathbf{L} \kappa)$ 恰好对一个 $i$
非零，并且此时等于 $1$。由引理
\ref{more-algebra-lemma-lift-perfect-from-residue-field}，这给出 (2)。

\medskip\noindent
在上面的证明中已经看到 (a) 成立。令 $U_n \subset \Spec(R)$ 为所有形如
$D(f)$ 且满足 $M_f \cong R_f[-n]$ 的开集之并。显然，若
$n \not = n'$，则 $U_n \cap U_{n'} = \emptyset$。若 $M$ 的 tor-幅度在
$[a, b]$ 内，则当 $n \not \in [a, b]$ 时 $U_n = \emptyset$。因此得到
(d) 中的乘积分解 $R = \prod_{a \leq n \leq b} R_n$，使 $U_n$ 对应于
$\Spec(R_n)$；见 Algebra，引理
\ref{algebra-lemma-disjoint-implies-product}。由于
$D(R) = \prod_{a \leq n \leq b} D(R_n)$，模范畴也同样，
% P02-E0729
(b)、(c) 与 (d) 立即成立。
\end{proof}

\section{分裂出一个自由模}
\label{more-algebra-section-serre-splitting}

\noindent
本节中的论证归功于 Serre；见 \cite{Serre-projective}。

\begin{situation}
\label{more-algebra-situation-splitting}
这里 $R$ 是环，$M$ 是有限表现 $R$-模。以 $\Omega \subset \Spec(R)$
表示所有闭点组成的集合，并赋予诱导拓扑。对 $x \in \Omega$，以
$M(x) = M/xM$ 表示 $M$ 在 $x$ 处的纤维。这是 $x$ 处剩余域
$\kappa(x)$ 上的有限维向量空间。给定 $s \in M$，以 $s(x)$ 表示
$s$ 在 $M(x)$ 中的像。
\end{situation}

\begin{lemma}
\label{more-algebra-lemma-which-elements-split}
在情形 \ref{more-algebra-situation-splitting} 中，设 $x \in \Omega$。存在
$\kappa(x)$-向量空间的典范短正合列
$$
0 \to B(x) \to M(x) \to V(x) \to 0
$$
，它具有以下性质：
% P02-E0730
对 $s_1, \ldots, s_r \in M$，以下条件等价：
\begin{enumerate}
\item 存在 $f \in R$，$f \not \in x$，使映射
$s_1, \ldots, s_r : R^{\oplus r} \to M$ 在逆转 $f$ 后成为一个
直和项的嵌入；
\item $s_1(x), \ldots, s_r(x)$ 在 $V(x)$ 中的像线性无关。
\end{enumerate}
\end{lemma}

\begin{proof}
把 $B(x) \subset M(x)$ 定义为映射
$$
\Hom_R(M, R) \to \Hom_{\kappa(x)}(M(x), \kappa(x))
$$
之像的正交补，并令 $V(x) = M(x)/B(x)$。于是每个 $R$-线性映射
$\varphi : M \to R$ 都诱导映射
$\overline{\varphi} : V(x) \to \kappa(x)$；反之，每个
$\kappa(x)$-线性映射 $\lambda : V(x) \to \kappa(x)$ 都等于某个
$\varphi$ 的 $\overline{\varphi}$。设 $s_1, \ldots, s_r \in M$。

\medskip\noindent
假设 $s_1, \ldots, s_r$ 在 $V(x)$ 中的像线性无关。于是可以找到
$\varphi_1, \ldots, \varphi_r \in \Hom_R(M, R)$，使
$\varphi_i(s_j)$ 在 $\kappa(x)$ 中的像为
$\delta_{ij}$\footnote{Kronecker 符号。}。因此复合映射
$$
R^{\oplus r}
\xrightarrow{s_1, \ldots, s_r}
M
\xrightarrow{\varphi_1, \ldots, \varphi_r}
R^{\oplus r}
$$
的矩阵有一个行列式 $f \in R$，它在 $\kappa(x)$ 中的像为 $1$
% P02-E0731
。显然，这说明 $s_1, \ldots, s_r : R^{\oplus r} \to M$ 在逆转
$f$ 后是一个直和项的嵌入。

\medskip\noindent
反之，假设存在 $f \in R$，$f \not \in x$，使
$s_1, \ldots, s_r : R^{\oplus r} \to M$ 在逆转 $f$ 后是一个
直和项的嵌入。于是可以找到 $R_f$-线性映射
$\varphi_i : M_f \to R_f$，使
$\varphi_i(s_j) = \delta_{ij} \in R_f$。由 Algebra，引理
\ref{algebra-lemma-hom-from-finitely-presented}，有
$\Hom_R(M, R)_f = \Hom_{R_f}(M_f, R_f)$，故可以找到 $n \geq 0$ 与
$\varphi'_i \in \Hom_R(M, R)$，使
$\varphi'_i(s_j) = f^n\delta_{ij} \in R$。由此，
$s_1, \ldots, s_r$ 在 $V(x)$ 中的像线性无关，因为
% P02-E0732
$\overline{\varphi}'_i(s_j) = f^n\delta_{ij}$。
\end{proof}

\noindent
在情形 \ref{more-algebra-situation-splitting} 中，给定
$s_1, \ldots, s_r \in M$，以
$Z(s_1, \ldots, s_r) \subset \Omega$ 表示所有满足如下条件的
$x \in \Omega$ 所成的集合：$s_1(x), \ldots, s_r(x)$ 在 $V(x)$ 中
的像线性相关。由引理，这是 $\Omega$ 的闭子集。

\begin{lemma}
\label{more-algebra-lemma-choose-values}
在情形 \ref{more-algebra-situation-splitting} 中，设
$x_1, \ldots, x_n \in \Omega$ 两两不同，并设
$v_i \in V(x_i)$。则存在 $s \in M$，使对
$i = 1, \ldots, n$，$s(x_i)$ 的像为 $v_i$。
\end{lemma}

\begin{proof}
由于 $x_i$ 是 $R$ 的极大理想，可以应用 Algebra，引理
\ref{algebra-lemma-chinese-remainder}，得到
$M(x_1) \oplus \ldots \oplus M(x_n)$ 是 $M$ 的商。
\end{proof}

\begin{proposition}
\label{more-algebra-proposition-splitting}
\begin{reference}
\cite[Theorem 2]{Serre-projective}
\end{reference}
在情形 \ref{more-algebra-situation-splitting} 中，假设 $\Omega$ 是 Noether
拓扑空间。设 $s_1, \ldots, s_h \in M$。设
$Z(s_1, \ldots, s_h) \subset F \subset \Omega$ 是闭集。设
$x_1, \ldots, x_n \in F$ 两两不同，并设 $v_i  \in V(x_i)$。
设 $k \geq 0$ 是整数，满足
$$
(*)\quad h + k \leq \dim_{\kappa(x)} V(x)\text{ 对所有 }x \in \Omega
$$
。则存在 $s \in M$ 与闭集 $F' \subset \Omega$，使得：
\begin{enumerate}
\item[(a)] $s(x_i)$ 的像为 $v_i$；
\item[(b)] $Z(s_1, \ldots, s_h, s) \subset F \cup F'$；
\item[(c)] $F'$ 的每个不可约分支在 $\Omega$ 中的余维都
$\geq k$。
\end{enumerate}
\end{proposition}

\begin{proof}
我们指出，余维定义于 Topology，第
\ref{topology-section-catenary-spaces} 节；还将使用 Topology，第
\ref{topology-section-noetherian} 节所收录的若干关于 Noether
拓扑空间的结果。

\medskip\noindent
对 $k$ 作归纳。若 $k = 0$，则按引理
\ref{more-algebra-lemma-choose-values} 选取 $s \in M$，并取 $F' = \Omega$。

\medskip\noindent
假设 $k > 0$。由归纳假设，可以选取 $u \in M$ 与闭集
$G \subset \Omega$，它们对
$s_1, \ldots, s_h$、$F$、$x_1, \ldots, x_n$、
$v_1, \ldots, v_n$ 以及 $k - 1$ 满足 (a)、(b)、(c)。

\medskip\noindent
设 $G = G_1 \cup \ldots \cup G_m$ 是 $G$ 的不可约分支分解。若
$G_j \subset F$，则可从列表中去掉它。因此可以假设，对
$j = 1, \ldots, m$，$G_j$ 不包含于 $F$。对
$j = 1, \ldots, m$，选择 $y_j \in G_j$，使
$y_j \not \in F$，且当 $j' \not = j$ 时
$y_j \not \in G_{j'}$。这是可行的，因为 $G$ 的不可约分支之间没有
包含关系。选取 $w_j \in V(y_j)$，使它不属于
$s_1(y_j), \ldots, s_h(y_j)$ 诸像张成的空间；这是可行的，因为
$h + k \leq \dim V(y_j)$ 且 $k > 0$。

\medskip\noindent
把归纳假设应用于 $h + 1$ 个截面
$s_1, \ldots, s_h, u$，闭集 $F \cup G$，诸点
$x_1, \ldots, x_n, y_1, \ldots, y_m \in F \cup G$，诸元素
$0 \in V(x_i)$ 与 $w_j \in V(y_j)$，以及整数 $k - 1$。注意，我们
把 $h$ 增加了 $1$，把 $k$ 减少了 $1$，所以命题的假设 $(*)$
仍然成立。这给出 $t \in M$ 与闭集 $H \subset \Omega$；它们对
$s_1, \ldots, s_h, u$、$F \cup G$、
$x_1, \ldots, x_n, y_1, \ldots, y_m$、
$0, \ldots, 0, w_1, \ldots, w_m$ 以及 $k - 1$ 满足 (a)、(b)、(c)。

\medskip\noindent
设 $H_1, \ldots, H_p \subset H$ 是 $H$ 中不包含于
$F \cup G$ 的不可约分支。与前面一样，选取 $z_l \in H_l$，使
$z_l \not \in F \cup G$，且当 $l' \not = l$ 时
$z_l \not \in H_{l'}$。利用 Algebra，引理
\ref{algebra-lemma-chinese-remainder}，可以选取 $f \in R$，使
$f(y_j) = 1$（$j = 1, \ldots, m$），且
$f(z_l) = 0$（$l = 1, \ldots, p$）。断言：元素
$s = u + f t$ 可用。

\medskip\noindent
首先，因为 $t(x_i) = 0$，$s(x_i)$ 的值与 $u(x_i)$ 相同，所以
$s(x_i)$ 的像为 $v_i$。这证明了 (a)。为完成证明，只需说明
$Z(s_1, \ldots, s_h, s)$ 中每个不包含于 $F$ 的不可约分支 $Z$
在 $\Omega$ 中的余维都 $\geq k$。事实上，这时可以令 $F'$ 等于
这些分支之并，从而得到 (b) 与 (c)。下面说明
$Z(s_1, \ldots, s_h, s)$ 不存在余维 $\leq k - 1$ 的不可约分支
$Z$：
\begin{enumerate}
\item 注意，
% P02-E0733
$Z(s_1, \ldots, s_h, s) \subset Z(s_1, \ldots, s_h, u, t) = F \cup H$
，因为 $s = u + ft$。故 $Z \subset H$。
\item $H$ 的不可约分支的余维都 $\geq k - 1$。因为 $Z$ 的余维
$\leq k - 1$，所以 $Z$ 等于 $H$ 的某个不可约分支。因此，对某个
$l \in \{1, \ldots, p\}$ 有 $Z = H_l$，或者对某个
% P02-E0734
$j \in \{1, \ldots m\}$ 有 $Z = G_j$。
\item 但 $Z = G_j$ 不可能，因为
$s_1(y_j), \ldots, s_h(y_j)$ 在 $V(y_j)$ 中的像线性无关，并且
$s(y_j) = u(y_j) + f(y_j) t(y_j) = u(y_j) + t(y_j)$ 的像形如
% P02-E0735
$$
\text{linear combination images of }s_i(y_j) + w_j
$$
；由 $w_j$ 的选择，它与 $s_1(y_j), \ldots, s_h(y_j)$ 在
$V(y_j)$ 中的像线性无关。
\item 同样，$Z = H_l$ 也不可能。事实上，
$s_1(z_l), \ldots, s_h(z_l)$ 在 $V(z_l)$ 中的像线性无关，而
$s(z_l) = u(z_l) + f(z_l) t(z_l) = u(z_l)$ 的像在 $V(z_l)$ 中
与它们线性无关，因为 $z_l \not \in F \cup G$。
\end{enumerate}
证明完毕。
\end{proof}

\begin{theorem}
\label{more-algebra-theorem-splitting}
\begin{reference}
\cite[Theorem 1]{Serre-projective}
\end{reference}
设环 $R$ 的极大谱 $\Omega \subset \Spec(R)$ 是维数
$d < \infty$ 的 Noether 拓扑空间。设 $M$ 是有限表现 $R$-模，并且
对所有 $\mathfrak m \in \Omega$，$R_\mathfrak m$-模
$M_\mathfrak m$ 都有秩 $> d$ 的自由直和项。则
$M \cong R \oplus M'$。
\end{theorem}

\begin{proof}
设对 $\mathfrak m \in \Omega$，$R_\mathfrak m^{\oplus r}$ 是
$M_\mathfrak m$ 的直和项。由 Algebra，引理
\ref{algebra-lemma-localization-colimit} 与
\ref{algebra-lemma-colimit-category-fp-modules}，对某个
$f \in R$，$f \not \in \mathfrak m$，$R_f^{\oplus r}$ 是
$M_f$ 的直和项。因此，该假设意味着对所有 $x \in \Omega$，
$\dim V(x) > d$，其中 $V(x)$ 如引理
\ref{more-algebra-lemma-which-elements-split} 所定义。把命题
\ref{more-algebra-proposition-splitting} 应用于 $F = \emptyset$、$h = 0$ 且没有
$s_i$、$n = 0$ 且没有 $x_i, v_i$，以及 $k = d + 1$，得到
$s \in M$ 与 $F' \subset \Omega$，使 $F'$ 的每个不可约分支的余维
都 $\geq d + 1$，且 $Z(s) \subset F'$。由于
$d = \dim(\Omega)$，这迫使 $F' = \emptyset$。因此
$s : R \to M$ 在每个极大理想处都是一个直和项的嵌入。由此，$s$
是普遍单射；见 Algebra，引理
\ref{algebra-lemma-universally-injective-check-stalks}。于是，由
Algebra，引理 \ref{algebra-lemma-universally-exact-split}，$s$ 是
分裂单射。
\end{proof}

\section{大投射模是自由模}
\label{more-algebra-section-big-projective-free}

\noindent
本节讨论 \cite{Bass} 的一个结果；建议读者参阅原论文。我们的论证将使用
\cite{Akasaki} 与 \cite{Hinohara} 中给出的略为简化的证明。

\begin{lemma}[Eilenberg 引理]
\label{more-algebra-lemma-eilenberg-swindle}
\begin{slogan}
Eilenberg 骗术
\end{slogan}
\begin{history}
在 \cite{Bass} 中可以读到：“……这是一个优雅的小骗术，Eilenberg
多年前便已注意到它，而它很可能是从 Barry Mazur 的额头上迸发出来的。”
\end{history}
\begin{reference}
\cite[Eilenberg's lemma]{Bass}
\end{reference}
若 $P \oplus Q \cong F$，其中 $F$ 是非有限生成的自由模，则
$P \oplus F \cong F$。
\end{lemma}

\begin{proof}
$$
F \cong F \oplus F \oplus \ldots \cong
P \oplus Q \oplus P \oplus Q \oplus \ldots \cong
P \oplus F \oplus F \oplus \ldots \cong P \oplus F
$$
\end{proof}

\begin{lemma}
\label{more-algebra-lemma-projective-plus-free-is-free}
设 $R$ 是环，$P$ 是投射模。存在自由模 $F$，使 $P \oplus F$ 自由。
\end{lemma}

\begin{proof}
由于 $P$ 投射，对某个模 $Q$，$F_0 = P \oplus Q$ 是自由模。令
$F = \bigoplus_{n \geq 1} F_0$。由引理
\ref{more-algebra-lemma-eilenberg-swindle}，$P \oplus F \cong F$。
\end{proof}

\begin{lemma}
\label{more-algebra-lemma-element-projective}
设 $R$ 是环，$P$ 是投射模。设 $s \in P$。则存在有限自由模 $F$ 与
有限自由直和项 $K \subset F \oplus P$，使
$(0, s) \in K$。
\end{lemma}

\begin{proof}
由引理 \ref{more-algebra-lemma-projective-plus-free-is-free}，可以找到一个（可能无限的）
自由模 $F$，使 $F \oplus P$ 自由。于是 $(0, s)$ 当然包含在某个
有限自由直和项 $K \subset F \oplus P$ 中。进而，$K$ 包含在
$F' \oplus P$ 中，其中 $F' \subset F$ 是有限自由直和项。
\end{proof}

\begin{lemma}
\label{more-algebra-lemma-trick-to-find-good-element}
设 $R$ 是环，其 Jacobson 根为 $J$，且 $R/J$ 是 Noether 环。设 $P$
是投射 $R$-模，并且对 $R$ 的每个极大理想 $\mathfrak m$，
$P_\mathfrak m$ 都有无限秩。设 $s \in P$ 且 $M \subset P$，满足
$Rs + M = P$。则可以找到 $m \in M$，使 $R(s + m)$ 是 $P$ 的自由
直和项。
\end{lemma}

\begin{proof}
该陈述有意义，因为由 Algebra，定理
\ref{algebra-theorem-projective-free-over-local-ring}，
$P_\mathfrak m$ 自由。

\medskip\noindent
以 $M' \subset P/JP$ 表示 $M$ 的像，以 $s' \in P/JP$ 表示 $s$ 的像。
注意，
% P02-E0736
$R/J s' + M' = P/JP$。假设可以找到 $m' \in M'$，使
$R/J(s' + m')$ 是 $M'$ 的自由直和项。
% P02-E0737
选择给出分裂的 $\varphi' : P/JP \to R/J$，即在 $R/J$ 中有
$\varphi'(s' + m') = 1$。由于 $P$ 是投射 $R$-模，可以找到
$\varphi'$ 的提升 $\varphi : P \to R$。选取映到 $m'$ 的
$m \in M$。于是 $\varphi(s + m) \in R$ 模 $J$ 同余于 $1$，因而是
$R$ 中的单位元（Algebra，引理
\ref{algebra-lemma-contained-in-radical}）。所以 $R(s + m)$ 是 $P$
的自由直和项。这把问题约化为下一段讨论的情形。

\medskip\noindent
假设 $R$ 是 Noether 环。设 $m \in M$，并设
$\varphi_1, \ldots, \varphi_n : P \to R$ 是 $R$-线性映射。以
$$
Z(s + m, \varphi_1, \ldots, \varphi_n) \subset \Spec(R)
$$
表示 $\varphi_1(s + m), \ldots, \varphi_n(s + m) \in R$ 的零点集。

\medskip\noindent
假设 $\mathfrak m$ 是 $R$ 的极大理想，且
$\mathfrak m \in Z(s + m, \varphi_1, \ldots, \varphi_n)$。令
$K = M \cap \bigcap \Ker(\varphi_i)$。我们断言，映射
$$
K/\mathfrak mK \to P/\mathfrak m P
$$
的像有无限维。事实上，商 $P/K$ 是有限 $R$-模，因为它同构于
$P/M \oplus R^{\oplus n}$ 的一个子模。因此，该映射的核是
$\text{Tor}_1^R(P/K, \kappa(\mathfrak m))$ 的一个商；由 Algebra，
引理 \ref{algebra-lemma-tor-noetherian}，后者有限。结合
$P/\mathfrak mP$ 有无限维这一事实，便得到断言。因此可以找到
$t \in K$，它映到向量空间 $P/\mathfrak mP$ 的元素
$\overline{t}$，而该元素与 $s + m$ 的像
$\overline{s + m}$ 线性无关。由线性代数，可以找到 $R$-线性映射
% P02-E0738
$\overline{\varphi} : P \to \kappa(\mathfrak m)$，使
$\overline{\varphi}(\overline{t}) = 1$ 且
$\overline{\varphi}(\overline{s + m}) = 0$。由于 $P$ 投射，可以找到
提升 $\overline{\varphi}$ 的 $R$-线性映射 $\varphi : P \to R$。
于是零点集
$Z(s + m + t, \varphi_1, \ldots, \varphi_n, \varphi)$ 包含于
$Z(s + m, \varphi_1, \ldots, \varphi_n)$，但不包含
$\mathfrak m$；换言之，它严格小于
$Z(s + m, \varphi_1, \ldots, \varphi_n)$。

\medskip\noindent
由于 $\Spec(R)$ 是 Noether 拓扑空间，由上述论证可知，可以找到
$m \in M$ 与 $\varphi_1, \ldots, \varphi_n : P \to R$，使闭子集
$Z(s + m, \varphi_1, \ldots, \varphi_n)$ 不包含 $\Spec(R)$ 的任何
闭点。因此 $Z(s + m, \varphi_1, \ldots, \varphi_n) = \emptyset$。
所以可以找到 $r_1, \ldots, r_n \in R$，使
$\sum r_i\varphi_i(s + m) = 1$。于是
$$
R \xrightarrow{s + m} P \xrightarrow{\sum r_i \varphi_i} R
$$
就是所需的分裂。
\end{proof}

\begin{lemma}
\label{more-algebra-lemma-element-in-free-summand}
设 $R$ 是环，其 Jacobson 根为 $J$，且 $R/J$ 是 Noether 环。设 $P$
是投射 $R$-模，并且对 $R$ 的每个极大理想 $\mathfrak m$，
$P_\mathfrak m$ 都有无限秩。设 $s \in P$。则可以找到有限稳定自由
直和项 $M \subset P$，使 $s \in M$。
\end{lemma}

\begin{proof}
由引理 \ref{more-algebra-lemma-element-projective}，可以找到有限自由模 $F$ 与
有限自由直和项 $K \subset F \oplus P$，使
$(0, s) \in K$。对 $F$ 的秩作归纳，把问题约化为下一段讨论的情形。

\medskip\noindent
假设存在有限稳定自由直和项 $K \subset R \oplus P$，使
$(0, s) \in K$。选择 $K$ 的补 $K'$，即
$R \oplus P = K \oplus K'$。投影
% P02-E0739
$\pi : R \oplus P \to K'$ 是满射；因此由引理
\ref{more-algebra-lemma-trick-to-find-good-element}，可以找到 $p \in P$，使
$\pi(1, p) \in K'$ 生成自由直和项。相应地，写
$K' = R\pi(1, p) \oplus K''$。于是
$$
R \oplus P = K \oplus K' = K \oplus R\pi(1, p) \oplus K''
$$
。投影 $\pi' : P \to K''$ 是满射\footnote{事实上，若
$k'' \in K''$，则把 $k''$ 看作 $K'$ 的元素，可以写成
$k'' = \lambda \pi(1, 0) + \pi(0, q)$，其中 $\lambda \in R$ 且
$q \in P$。这意味着
$k'' = \lambda \pi(1, p) + \pi(0, q - \lambda p)$。进而，这意味着
$q - \lambda p$ 在复合映射
$P \to R \oplus P \xrightarrow{\pi} K' \to K''$ 下映到 $k''$，因为
$K' \to K''$ 零化 $\pi(1, p)$。}，因而它分裂
% P02-E0740
（因为 $K''$ 投射）。因此 $\Ker(\pi') \subset P$ 是包含 $s$ 的
直和项。最后，由构造有同构
$$
R \oplus \Ker(\pi') \cong K \oplus R\pi(1, p)
$$
；所以，由于 $K$ 有限且稳定自由，$\Ker(\pi')$ 也如此。
\end{proof}

\begin{theorem}
\label{more-algebra-theorem-countable-free}
\begin{reference}
\cite[Theorem 3.1]{Bass} 的交换情形
\end{reference}
设 $R$ 是环，其 Jacobson 根为 $J$，且 $R/J$ 是 Noether 环。设 $P$
是可数生成投射 $R$-模，并且对 $R$ 的每个极大理想
$\mathfrak m$，$P_\mathfrak m$ 都有无限秩。则 $P$ 自由。
\end{theorem}

\begin{proof}
先证明 $P$ 是有限稳定自由模的可数直和。设
$x_1, x_2, \ldots$ 是 $P$ 的一组可数生成元。归纳构造 $P$ 的有限
稳定自由直和项 $F_1, F_2, \ldots$，使对所有 $n$，
$F_1 \oplus \ldots \oplus F_n$ 都是包含
$x_1, \ldots, x_n$ 的 $P$ 的直和项。事实上，给定具有所需性质的
$F_1, \ldots, F_n$，写
$$
P = F_1 \oplus \ldots \oplus F_n \oplus P'
$$
，并令 $s \in P'$ 是 $x_{n + 1}$ 的像。由引理
\ref{more-algebra-lemma-element-in-free-summand}，可以找到包含 $s$ 的有限
稳定自由直和项 $F_{n + 1} \subset P'$。于是
$P = \bigoplus_{i = 1}^{\infty} F_i$。

\medskip\noindent
假设 $P$ 是非零有限稳定自由模的无限直和
$P = \bigoplus_{i = 1}^{\infty} F_i$。下面按如下方式使用各 $F_i$
的稳定自由性：每个 $F_i$ 的秩都是常数（且为正）。因此，对 $R$
的每个极大理想 $\mathfrak m$，$P_\mathfrak m$ 都是可数无限秩的
自由模。把引理 \ref{more-algebra-lemma-trick-to-find-good-element} 应用于
$s = 0$ 与 $M = P$，可以找到 $t_1 \in P$，使 $Rt_1$ 是 $P$ 的
自由直和项。于是对某个 $n_1 > n_0 = 0$，
$t_1$ 包含在 $F_1 \oplus \ldots \oplus F_{n_1}$ 中。把同一论证
应用于 $\bigoplus_{n > n_1} F_n$，得到 $n_1 < n_2$ 与
$t_2 \in F_{n_1 + 1} \oplus \ldots \oplus F_{n_2}$，后者生成一个
自由直和项。继续如此，得到无限秩自由直和项
$$
\bigoplus\nolimits_{i \geq 1} t_i :
\bigoplus\nolimits_{i \geq 1} R
\longrightarrow
\bigoplus\nolimits_{i \geq 1}
\bigoplus\nolimits_{n_i \geq n > n_{i - 1}} F_n = P
$$
。因此，对某个可数秩自由 $R$-模 $F$，有
$P \cong Q \oplus F$。由于 $Q$ 可数生成，对某个模 $Q'$，有
$Q \oplus Q' \cong F$。于是 Eilenberg 骗术（引理
\ref{more-algebra-lemma-eilenberg-swindle}）说明 $Q \oplus F \cong F$，从而
$P$ 自由。
\end{proof}
% Generated mechanically from the frozen smoothing.tex target.
% Do not edit this generated file; edit the bound locale source instead.

% OK, start here.
%

\setcounter{chapter}{15}
\chapter{环同态的光滑化}



\phantomsection
\label{smoothing-section-phantom}


\section{引言}
\label{smoothing-section-introduction}

\noindent
本章的主要结论如下：
$$
\fbox{诺特环之间的正则同态是光滑同态的滤过余极限。}
$$
该定理由 Popescu 证明，见 \cite{popescu-letter}。
Richard Swan 给出了 Popescu 证明的一份易读阐述，见 \cite{swan}；
他采用了 Andr\'e 的笔记以及 Ogoma 的一篇论文，见 \cite{Ogoma}。

\medskip\noindent
我们的叙述沿循 Swan 的处理，不过会先证明一个中间结果，
以便在稍为简单的情形中展开论证。下面概述全章结构。
首先解决如下“提升问题”：光滑代数的滤过余极限的平坦无穷小形变，
仍是光滑代数的滤过余极限。这个结果实质上说明，只须对约化诺特环之间的
同态证明主定理。随后证明两个十分巧妙的引理，分别称为“提升引理”和
“消奇异化引理”。我们将说明，合用这两个引理后，主定理可归结为证明：
域 $k$ 上几何正则的诺特代数 $\Lambda$ 是光滑 $k$-代数的滤过余极限。
接着讨论 Popescu--Ogoma--Swan--Andr\'e 证明所需的局部技巧。
最后三节给出证明。

\medskip\noindent
最后列出若干参考文献。设 $A$ 为 Hensel 诺特局部环。
若对任意 $f_1, \ldots, f_m \in A[x_1, \ldots, x_n]$，方程组
$f_1 = 0, \ldots, f_m = 0$ 在 $A$ 的完备化中有解当且仅当它在 $A$ 中有解，
则称 $A$ 具有{\it 逼近性质}。这一定义由 Artin 提出。
Artin 首先对解析方程组证明了逼近性质，见
\cite{Artin-Analytic-Approximation}。
在 \cite{Artin-Algebraic-Approximation} 中，Artin 对本质有限型于
优秀离散赋值环的局部环证明了逼近性质。
Artin 猜想（见 \cite{Artin-Algebraic-Approximation} 第 26 页），
每个优秀 Hensel 局部环都应具有逼近性质。

\medskip\noindent
后来人们提出猜想：诺特环的每个正则同态都是光滑代数的滤过余极限（例如见
\cite{Raynaud-Rennes}, \cite{popescu-global}, \cite{Artin-power-series},
\cite{Artin-Denef}）。我们不确定这一猜想应归于谁。\footnote{在
\cite{Artin-power-series}、\cite{Artin-Denef} 和 \cite{popescu-global}
中表述的问题或猜想更强；\cite{Cipu} 证明了它与原版本等价。}
它与逼近性质的联系如下：若 $A$ 如上，并且 $A \to A^\wedge$ 是光滑代数的
余极限，则逼近性质成立（此处待补引用）。此外，若 $A$ 是优秀局部环，
则主定理适用于同态 $A \to A^\wedge$，因为优秀局部环的条件之一是其
形式纤维几何正则。优秀局部环由 Grothendieck 定义，该定义于 1965 年刊印。

\medskip\noindent
\cite{Artin-power-series} 证明了：对任意本质有限型于域的局部环 $R$，
$R \to R^\wedge$ 都是光滑代数的滤过余极限。
\cite{Rotthaus-Artin} 证明了：对任意本质有限型于优秀离散赋值环的局部环
$R$，$R \to R^\wedge$ 都是光滑代数的滤过余极限。
最后，如上所述，主定理在
\cite{popescu-GND}, \cite{popescu-GNDA}, \cite{popescu-letter}, 
\cite{Ogoma} 和 \cite{swan} 中得到证明。

\medskip\noindent
反过来，利用上述若干结果，\cite{Rotthaus-excellent} 证明了：
任何具有逼近性质的诺特局部环都是优秀环。

\medskip\noindent
论文 \cite{Spivakovsky} 给出了主定理的另一种处理方式，但似乎颇难读懂
（例如，\cite[Lemma 5.2]{Spivakovsky} 看来是对 \cite[Lemma 3]{Elkik}
的不正确改述）。关于这些材料还有一篇 Bourbaki 报告，见 \cite{Teissier}。












\section{奇异理想}
\label{smoothing-section-singular-ideal}

\noindent
设 $R \to A$ 为环同态。$A$ 相对于 $R$ 的奇异理想，是 $A$ 中定义态射
$\Spec(A) \to \Spec(R)$ 的奇异轨迹的根理想。其严格定义如下。

\begin{definition}
\label{smoothing-definition-singular-ideal}
设 $R \to A$ 为环同态。$A$ 相对于 $R$ 的{\it 奇异理想}记作 $H_{A/R}$，
它是满足下式的唯一根理想 $H_{A/R} \subset A$：
$$
V(H_{A/R}) = \{\mathfrak q \in \Spec(A) \mid R \to A
\text{ 在 }\mathfrak q\text{ 处不光滑}\}.
$$
\end{definition}

\noindent
这一定义合理，因为 $R \to A$ 光滑的素理想所成集合是开集，见
《代数》定义 \ref{algebra-definition-smooth-at-prime}。
为了找出奇异理想的一组显式生成元，我们先证明以下引理。

\begin{lemma}
\label{smoothing-lemma-find-strictly-standard}
设 $R$ 为环，$A = R[x_1, \ldots, x_n]/(f_1, \ldots, f_m)$，
$\mathfrak q \subset A$ 为素理想。假设 $R \to A$ 在 $\mathfrak q$ 处光滑。
则存在 $a \in A$ 且 $a \not \in \mathfrak q$、整数 $c$ 且
$0 \leq c \leq \min(n, m)$，以及基数均为 $c$ 的子集
$U \subset \{1, \ldots, n\}$、$V \subset \{1, \ldots, m\}$，使得
$$
a = a' \det(\partial f_j/\partial x_i)_{j \in V, i \in U}
$$
对某个 $a' \in A$ 成立，并且
$$
a f_\ell \in (f_j, j \in V) + (f_1, \ldots, f_m)^2
$$
对所有 $\ell \in \{1, \ldots, m\}$ 都成立。
\end{lemma}

\begin{proof}
令 $I = (f_1, \ldots, f_m)$。于是 $A$ 在 $R$ 上的朴素余切复形同伦等价于
$I/I^2 \to \bigoplus A\text{d}x_i$，见
《代数》引理 \ref{algebra-lemma-NL-homotopy}。
我们将用到朴素余切复形的构造与局部化可交换；见
《代数》第 \ref{algebra-section-netherlander} 节，特别是
《代数》引理 \ref{algebra-lemma-localize-NL}。
由《代数》定义 \ref{algebra-definition-smooth} 和
\ref{algebra-definition-smooth-at-prime} 可知，对某个
$a \in A$、$a \not \in \mathfrak q$，同态
$(I/I^2)_a \to \bigoplus A_a\text{d}x_i$ 是分裂单射。
重排 $x_1, \ldots, x_n$ 与 $f_1, \ldots, f_m$ 的编号后，可以假设
$f_1, \ldots, f_c$ 构成向量空间
$I/I^2 \otimes_A \kappa(\mathfrak q)$ 的一组基，并且
$\text{d}x_{c + 1}, \ldots, \text{d}x_n$ 的像构成
$\Omega_{A/R} \otimes_A \kappa(\mathfrak q)$ 的一组基。
因此，用 $aa'$ 代替 $a$（其中 $a' \in A$ 且
$a' \not \in \mathfrak q$）后，可以假设 $f_1, \ldots, f_c$
构成 $(I/I^2)_a$ 的一组基，并且
$\text{d}x_{c + 1}, \ldots, \text{d}x_n$ 的像构成
$(\Omega_{A/R})_a$ 的一组基。在此情形中，当整数 $N$ 充分大时，
$a^N$ 满足引理的条件（取 $U = V = \{1, \ldots, c\}$）。
\end{proof}

\noindent
我们将使用 $A$ 中相对于 $R$ 的{\it 严格标准}元这一概念。
这里的概念比 Swan 的论文 \cite{swan} 中所用的概念稍弱。
我们还把上一个引理中所得类型的元素定义为{\it 初等标准}元。
引理 \ref{smoothing-lemma-compare-standard} 将比较这些不同类型的元素。

\begin{definition}
\label{smoothing-definition-strictly-standard}
设 $R \to A$ 为有限表示的环同态，$a \in A$。若存在一个表示
$A = R[x_1, \ldots, x_n]/(f_1, \ldots, f_m)$
以及满足 $0 \leq c \leq \min(n, m)$ 的整数 $c$，使得
\begin{equation}
\label{smoothing-equation-elementary-standard-one}
a = a' \det(\partial f_j/\partial x_i)_{i, j = 1, \ldots, c}
\end{equation}
对某个 $a' \in A$ 成立，并且
\begin{equation}
\label{smoothing-equation-elementary-standard-two}
a f_{c + j} \in (f_1, \ldots, f_c) + (f_1, \ldots, f_m)^2
\end{equation}
对 $j = 1, \ldots, m - c$ 成立，则称 $a$ 是
{\it $A$ 中相对于 $R$ 的初等标准元}。若存在一个表示
$A = R[x_1, \ldots, x_n]/(f_1, \ldots, f_m)$
以及满足 $0 \leq c \leq \min(n, m)$ 的整数 $c$，使得
\begin{equation}
\label{smoothing-equation-strictly-standard-one}
a = \sum\nolimits_{I \subset \{1, \ldots, n\},\ |I| = c}
a_I \det(\partial f_j/\partial x_i)_{j = 1, \ldots, c,\ i \in I}
\end{equation}
对某些 $a_I \in A$ 成立，并且
\begin{equation}
\label{smoothing-equation-strictly-standard-two}
a f_{c + j} \in (f_1, \ldots, f_c) + (f_1, \ldots, f_m)^2
\end{equation}
对 $j = 1, \ldots, m - c$ 成立，则称 $a$ 是
{\it $A$ 中相对于 $R$ 的严格标准元}。
\end{definition}

\noindent
以下引理可用来理解式 (\ref{smoothing-equation-strictly-standard-one}) 所蕴含的结论。

\begin{lemma}
\label{smoothing-lemma-parse-equation-strictly-standard-one}
设 $R$ 为环，$A = R[x_1, \ldots, x_n]/(f_1, \ldots, f_m)$，并记
$I = (f_1, \ldots, f_m)$。设 $a \in A$。则式
(\ref{smoothing-equation-strictly-standard-one}) 蕴含存在一个 $A$-线性同态
$\psi : \bigoplus\nolimits_{i = 1, \ldots, n} A \text{d}x_i \to A^{\oplus c}$
使得复合
$$
A^{\oplus c} \xrightarrow{(f_1, \ldots, f_c)}
I/I^2 \xrightarrow{f \mapsto \text{d}f}
\bigoplus\nolimits_{i = 1, \ldots, n} A \text{d}x_i
\xrightarrow{\psi}
A^{\oplus c}
$$
就是乘以 $a$ 的同态。反之，若这样的 $\psi$ 存在，则 $a^c$ 满足
(\ref{smoothing-equation-strictly-standard-one})。
\end{lemma}

\begin{proof}
这是《代数》引理 \ref{algebra-lemma-matrix-left-inverse} 的一个特例。
\end{proof}

\begin{lemma}[Elkik]
\label{smoothing-lemma-elkik}
设 $R \to A$ 为有限表示的环同态。奇异理想 $H_{A/R}$ 既是由
$A$ 中相对于 $R$ 的严格标准元生成的理想之根，也是由
$A$ 中相对于 $R$ 的初等标准元生成的理想之根。
\end{lemma}

\begin{proof}
假设 $a$ 是 $A$ 中相对于 $R$ 的严格标准元。我们断言 $A_a$ 在 $R$ 上光滑，
这便证明了 $a \in H_{A/R}$。事实上，令
$A = R[x_1, \ldots, x_n]/(f_1, \ldots, f_m)$、$c$ 和 $a' \in A$
如定义 \ref{smoothing-definition-strictly-standard} 所述。记
$I = (f_1, \ldots, f_m)$，则 $A$ 在 $R$ 上的朴素余切复形为
$I/I^2 \to \bigoplus A\text{d}x_i$。
条件 (\ref{smoothing-equation-strictly-standard-two}) 蕴含 $(I/I^2)_a$
由 $f_1, \ldots, f_c$ 的类生成。条件
(\ref{smoothing-equation-strictly-standard-one}) 蕴含微分同态
$(I/I^2)_a \to \bigoplus A_a\text{d}x_i$ 有左逆，见
引理 \ref{smoothing-lemma-parse-equation-strictly-standard-one}。
因此，由定义和《代数》引理 \ref{algebra-lemma-localize-NL}，
$R \to A_a$ 光滑。

\medskip\noindent
令 $H_e, H_s \subset A$ 分别为由 $A$ 中相对于 $R$ 的初等标准元、
严格标准元生成的理想之根。由定义和刚才的证明，
$H_e \subset H_s \subset H_{A/R}$。包含关系 $H_{A/R} \subset H_e$
由引理 \ref{smoothing-lemma-find-strictly-standard} 得出。
\end{proof}

\begin{example}
\label{smoothing-example-not-quasi-compact}
有限表示环同态的光滑点集未必是拟紧开集。例如，令
$R = k[x, y_1, y_2, y_3, \ldots]/(xy_i)$ 且 $A = R/(x)$。
则 $R \to A$ 的光滑轨迹为 $\bigcup D(y_i)$，它不是拟紧的。
\end{example}

\begin{lemma}
\label{smoothing-lemma-strictly-standard-base-change}
设 $R \to A$ 为有限表示的环同态，$R \to R'$ 为环同态。
若 $a \in A$ 是 $A$ 中相对于 $R$ 的初等标准元（相应地，严格标准元），
则 $a \otimes 1$ 是 $A \otimes_R R'$ 中相对于 $R'$ 的初等标准元
（相应地，严格标准元）。
\end{lemma}

\begin{proof}
若 $A = R[x_1, \ldots, x_n]/(f_1, \ldots, f_m)$ 是 $A$ 在 $R$ 上的一个表示，
则
$A \otimes_R R' = R'[x_1, \ldots, x_n]/(f'_1, \ldots, f'_m)$
是 $A \otimes_R R'$ 在 $R'$ 上的一个表示，其中 $f'_j$ 是 $f_j$
在 $R'[x_1, \ldots, x_n]$ 中的像。因此结论由定义立即得出。
\end{proof}

\begin{lemma}
\label{smoothing-lemma-final-solve}
设 $R \to A \to \Lambda$ 为环同态，且 $A$ 是有限表示 $R$-代数。
假设 $H_{A/R} \Lambda = \Lambda$。则存在分解
$A \to B \to \Lambda$，其中 $B$ 在 $R$ 上光滑。
\end{lemma}

\begin{proof}
选取 $f_1, \ldots, f_r \in H_{A/R}$ 和
$\lambda_1, \ldots, \lambda_r \in \Lambda$，使得在 $\Lambda$ 中
$\sum f_i\lambda_i = 1$。令
$B = A[x_1, \ldots, x_r]/(f_1x_1 + \ldots + f_rx_r - 1)$
并以 $x_i \mapsto \lambda_i$ 定义 $B \to \Lambda$。
为验证 $B$ 在 $R$ 上光滑，注意由 $H_{A/R}$ 的定义，$A_{f_i}$
在 $R$ 上光滑，而 $B_{f_i}$ 在 $A_{f_i}$ 上光滑。细节从略。
\end{proof}





\section{代数的表示}
\label{smoothing-section-presentations}

\noindent
本节的若干结果由 Elkik 证明。注意，下一个引理中的代数 $C$ 是 $A$ 上的
对称代数。此外，若 $R$ 是诺特环，则 $C$ 是有限表示 $R$-代数。

\begin{lemma}
\label{smoothing-lemma-improve-presentation}
设 $R$ 为环，$A$ 为有限表示 $R$-代数。存在具有缩回的有限型
$R$-代数同态 $A \to C$，且满足以下两个性质：
\begin{enumerate}
\item 对每个使 $R \to A_a$ 为局部完全交的 $a \in A$
（见《代数详论》定义
\ref{more-algebra-definition-local-complete-intersection}），
$C_a$ 在 $A_a$ 上光滑，并且有一个表示
$C_a = R[y_1, \ldots, y_m]/J$，其中 $J/J^2$ 是自由 $C_a$-模；
\item 对每个使 $A_a$ 在 $R$ 上光滑的 $a \in A$，模
$\Omega_{C_a/R}$ 是自由 $C_a$-模。
\end{enumerate}
\end{lemma}

\begin{proof}
选取一个表示 $A = R[x_1, \ldots, x_n]/I$，并记
$I = (f_1, \ldots, f_m)$。用短正合列
$$
0 \to K \to A^{\oplus m} \to I/I^2 \to 0
$$
定义 $A$-模 $K$；其中中间项的第 $j$ 个基向量 $e_j$ 映到右端 $f_j$ 的类。
令
$$
C = \text{Sym}^*_A(I/I^2).
$$
缩回就是投影到 $C$ 的 $0$ 次部分。存在满射
$R[x_1, \ldots, x_n, y_1, \ldots, y_m] \to C$，把 $y_j$ 映到
$I/I^2$ 中 $f_j$ 的类。此同态的核 $J$ 由 $f_1, \ldots, f_m$，
以及所有形如 $\sum h_j y_j$ 的元素生成；这里
$h_j \in R[x_1, \ldots, x_n]$，且 $\sum h_j e_j$ 定义 $K$ 的一个元素。
把《代数》引理 \ref{algebra-lemma-exact-sequence-NL} 应用于
$R \to A \to C$ 及上述表示，并使用《代数详论》引理
\ref{more-algebra-lemma-cotangent-complex-symmetric-algebra}，得到短正合列
\begin{equation}
\label{smoothing-equation-sequence}
I/I^2 \otimes_A C \to J/J^2 \to K \otimes_A C \to 0
\end{equation}
的 $C$-模。设 $h \in R[x_1, \ldots, x_n]$ 的像为 $a \in A$。
对局部化环采用如下表示：
$$
A_a = R[x_0, x_1, \ldots, x_n]/I'
\quad\text{且}\quad
C_a = R[x_0, x_1, \ldots, x_n, y_1, \ldots, y_m]/J'
$$
其中 $I' = (hx_0 - 1, I)$，$J' = (hx_0 - 1, J)$。因此，作为
$A_a$-模有 $I'/(I')^2 = A_a \oplus (I/I^2)_a$，作为 $C_a$-模有
$J'/(J')^2 = C_a \oplus (J/J^2)_a$。于是得到
\begin{equation}
\label{smoothing-equation-sequence-localized}
C_a \oplus I/I^2 \otimes_A C_a \to
C_a \oplus (J/J^2)_a \to
K \otimes_A C_a \to 0
\end{equation}
这正是《代数》引理 \ref{algebra-lemma-exact-sequence-NL} 对应于
$R \to A_a \to C_a$ 及上述表示的正合列。

\medskip\noindent
接下来，假设 $a \in A$ 使得 $A_a$ 是 $R$ 上的局部完全交。
则 $(I/I^2)_a$ 是有限投射 $A_a$-模，见《代数详论》引理
\ref{more-algebra-lemma-quasi-regular-ideal-finite-projective}。
因此 $K_a \oplus (I/I^2)_a \cong A_a^{\oplus m}$ 是自由模；特别地，
$K_a$ 也是有限投射模。由《代数详论》引理
\ref{more-algebra-lemma-transitive-lci-at-end}，正合列
(\ref{smoothing-equation-sequence-localized}) 在左端也正合。因此
$$
J'/(J')^2 \cong
C_a \oplus I/I^2 \otimes_A C_a \oplus K \otimes_A C_a \cong
C_a^{\oplus m + 1}
$$
这就证明了 (1)。最后，再假设 $A_a$ 在 $R$ 上光滑。同一表示表明，
$\Omega_{C_a/R}$ 是同态
$$
J'/(J')^2 \longrightarrow
\bigoplus\nolimits_i C_a\text{d}x_i \oplus \bigoplus\nolimits_j C_a\text{d}y_j
$$
的余核。在上述分解中，$J'/(J')^2$ 的直和项 $C_a$ 对应于
$hx_0 - 1$，因而同构地映到直和项 $C_a\text{d}x_0$。
$J'/(J')^2$ 的直和项 $I/I^2 \otimes_A C_a$ 单射到
$\bigoplus_{i = 1, \ldots, n} C_a\text{d}x_i$，其商为
$\Omega_{A_a/R} \otimes_{A_a} C_a$。直和项 $K \otimes_A C_a$
单射到 $\bigoplus_{j \geq 1} C_a\text{d}y_j$，其商同构于
$I/I^2 \otimes_A C_a$。所以，上一个显示同态的余核是模
$I/I^2 \otimes_A C_a \oplus \Omega_{A_a/R} \otimes_{A_a} C_a$。
由光滑环同态的定义，$(I/I^2)_a \oplus \Omega_{A_a/R}$ 是自由模，
从而 (2) 成立。
\end{proof}

\noindent
对 Hensel 对上的光滑环同态，以下命题由 Elkik 在 \cite{Elkik} 中证明。
关于光滑环同态的表述可见 \cite{Arabia}；该文还证明了光滑代数之间的
环同态可以提升。

\begin{proposition}
\label{smoothing-proposition-lift-smooth}
\begin{slogan}
光滑代数和 syntomic 代数可沿满射提升
\end{slogan}
设 $R \to R_0$ 是核为 $I$ 的满环同态。
\begin{enumerate}
\item 若 $R_0 \to A_0$ 是 syntomic 环同态，则存在 syntomic 环同态
$R \to A$，使得 $A/IA \cong A_0$。
\item 若 $R_0 \to A_0$ 是光滑环同态，则存在光滑环同态
$R \to A$，使得 $A/IA \cong A_0$。
\end{enumerate}
\end{proposition}

\begin{proof}
假设 $R_0 \to A_0$ 是 syntomic 的；特别地，它是局部完全交
（见《代数详论》引理 \ref{more-algebra-lemma-syntomic-lci}）。
选取表示 $A_0 = R_0[x_1, \ldots, x_n]/J_0$，并令
$C_0 = \text{Sym}^*_{A_0}(J_0/J_0^2)$。注意，$J_0/J_0^2$ 是有限投射
$A_0$-模（见《代数》引理
\ref{algebra-lemma-syntomic-presentation-ideal-mod-squares}）。
由引理 \ref{smoothing-lemma-improve-presentation}，环同态 $A_0 \to C_0$ 光滑，
并且可找到表示 $C_0 = R_0[y_1, \ldots, y_m]/K_0$，使得
$K_0/K_0^2$ 是自由 $C_0$-模。由《代数》引理
\ref{algebra-lemma-huber}，可以假设
$C_0 = R_0[y_1, \ldots, y_m]/(\overline{f}_1, \ldots, \overline{f}_c)$
其中 $\overline{f}_1, \ldots, \overline{f}_c$ 的像构成 $K_0/K_0^2$
在 $C_0$ 上的一组基。选取 $\overline{f}_1, \ldots, \overline{f}_c$
的提升 $f_1, \ldots, f_c \in R[y_1, \ldots, y_c]$，并令
$$
C = R[y_1, \ldots, y_m]/(f_1, \ldots, f_c)
$$
由构造，$C_0 = C/IC$。由《代数》引理
\ref{algebra-lemma-localize-relative-complete-intersection}，
用 $C_g$ 代替 $C$ 后，可以假设 $C$ 是 $R$ 上的相对整体完全交。
因此存在有限投射 $A_0$-模 $P_0$，使
$C_0 = \text{Sym}^*_{A_0}(P_0)$ 同构于某个 syntomic $R$-代数
$C$ 的 $C/IC$。

\medskip\noindent
选取整数 $n$ 以及直和分解
$A_0^{\oplus n} = P_0 \oplus Q_0$.
由《代数详论》引理 \ref{more-algebra-lemma-lift-projective-module}，
可找到平展环同态 $C \to C'$，它诱导同构 $C/IC \to C'/IC'$；
还可找到有限投射 $C'$-模 $Q$，使 $Q/IQ$ 同构于
$Q_0 \otimes_{A_0} C/IC$。于是 $D = \text{Sym}_{C'}^*(Q)$ 是光滑
$C'$-代数（见《代数详论》引理
\ref{more-algebra-lemma-symmetric-algebra-smooth}）。图示如下：
$$
\xymatrix{
R \ar[d] \ar[rr] & &
C \ar[r] \ar[d] &
C' \ar[r] \ar[d] &
D \ar[d] \\
R/I \ar[r] &
A_0 \ar[r] &
C/IC \ar[r]^{\cong} &
C'/IC' \ar[r] &
D/ID
}
$$
注意，$Q$ 的选取给出
\begin{align*}
D/ID & =
\text{Sym}_{C/IC}^*(Q_0 \otimes_{A_0} C/IC) \\
& =
\text{Sym}_{A_0}^*(Q_0) \otimes_{A_0} C/IC \\
& =
\text{Sym}_{A_0}^*(Q_0) \otimes_{A_0}
\text{Sym}_{A_0}^*(P_0) \\
& =
\text{Sym}_{A_0}^*(Q_0 \oplus P_0) \\
& =
\text{Sym}_{A_0}^*(A_0^{\oplus n}) \\
& =
A_0[x_1, \ldots, x_n]
\end{align*}
选取 $f_1, \ldots, f_n \in D$，使它们在
$D/ID = A_0[x_1, \ldots, x_n]$ 中分别映到 $x_1, \ldots, x_n$。
令 $A = D/(f_1, \ldots, f_n)$。注意 $A_0 = A/IA$。
我们断言，$R \to A$ 在 $V(IA)$ 的某个邻域上是 syntomic 的。
若断言成立，就可找到映到 $1 \in A_0$ 的 $f \in A$，使 $A_f$
在 $R$ 上 syntomic；于是 (1) 的证明完成。

\medskip\noindent
证明上述断言。注意，$R \to D$ 是 syntomic 环同态 $R \to C$、
平展环同态 $C \to C'$ 与光滑环同态 $C' \to D$ 的复合，因而是
syntomic 的（见《代数》引理
\ref{algebra-lemma-composition-syntomic} 以及
\ref{algebra-lemma-smooth-syntomic}）。这个问题在 $\Spec(D)$ 上是局部的，
故可假设 $D$ 是相对整体完全交（见《代数》引理
\ref{algebra-lemma-syntomic}）。设
$D = R[y_1, \ldots, y_m]/(g_1, \ldots, g_s)$。
令 $f'_1, \ldots, f'_n \in R[y_1, \ldots, y_m]$ 为
$f_1, \ldots, f_n$ 的提升。应用《代数》引理
\ref{algebra-lemma-localize-relative-complete-intersection} 即得断言。

\medskip\noindent
证明 (2)。由于光滑环同态是 syntomic 的，可找到 syntomic 环同态
$R \to A$，使得 $A_0 = A/IA$。由假设，$R \to A$ 的纤维在
$V(I)$ 中的素理想上光滑，故 $R \to A$ 在 $V(IA)$ 的一个开邻域上光滑
（见《代数》引理 \ref{algebra-lemma-flat-fibre-smooth}）。
因此，以一个局部化代替 $A$，便得到所需结论。
\end{proof}

\noindent
我们知道，任意 syntomic 环同态 $R \to A$ 在局部都是相对整体完全交，见
《代数》引理 \ref{algebra-lemma-syntomic}。下一个引理说明，
$\Spec(A)$ 上的向量丛是相对整体完全交。

\begin{lemma}
\label{smoothing-lemma-syntomic-complete-intersection}
设 $R \to A$ 为 syntomic 环同态。则存在具有缩回的光滑 $R$-代数同态
$A \to C$，使得 $C$ 是 $R$ 上的相对整体完全交，即
$$
C \cong R[x_1, \ldots, x_n]/(f_1, \ldots, f_c)
$$
在 $R$ 上平坦，且所有纤维的维数均为 $n - c$。
\end{lemma}

\begin{proof}
应用引理 \ref{smoothing-lemma-improve-presentation} 得到 $A \to C$。
由《代数》引理 \ref{algebra-lemma-huber}，可写成
$C = R[x_1, \ldots, x_n]/(f_1, \ldots, f_c)$，其中 $f_i$ 的像构成
$J/J^2$ 的一组基。环同态 $R \to C$ 是一个 syntomic 环同态与一个光滑
环同态的复合，因而是 syntomic 的（所以是平坦的）。由《代数》引理
\ref{algebra-lemma-lci}，纤维的维数为 $n - c$
（这些纤维是局部完全交，故该引理适用）。
\end{proof}

\begin{lemma}
\label{smoothing-lemma-smooth-standard-smooth}
设 $R \to A$ 为光滑环同态。则存在具有缩回的光滑 $R$-代数同态
$A \to B$，使得 $B$ 在 $R$ 上标准光滑，即
$$
B \cong R[x_1, \ldots, x_n]/(f_1, \ldots, f_c)
$$
并且 $\det(\partial f_j/\partial x_i)_{i, j = 1, \ldots, c}$
在 $B$ 中可逆。
\end{lemma}

\begin{proof}
应用引理 \ref{smoothing-lemma-syntomic-complete-intersection}，得到具有缩回的光滑
$R$-代数同态 $A \to C$，使得
$C = R[x_1, \ldots, x_n]/(f_1, \ldots, f_c)$ 是 $R$ 上的相对整体完全交。
由于 $C$ 在 $R$ 上光滑，存在短正合列
$$
0 \to
\bigoplus\nolimits_{j = 1, \ldots, c} C f_j \to
\bigoplus\nolimits_{i = 1, \ldots, n} C\text{d}x_i \to
\Omega_{C/R} \to 0
$$
由于 $\Omega_{C/R}$ 是投射 $C$-模，此正合列分裂。选取第一个同态的左逆
$t$。记
$t(\text{d}x_i) = \sum c_{ij} f_j$
于是 $\sum_i \frac{\partial f_j}{\partial x_i} c_{i\ell} = \delta_{j\ell}$
（Kronecker delta）。令
$$
B' = C[y_1, \ldots, y_c] =
R[x_1, \ldots, x_n, y_1, \ldots, y_c]/(f_1, \ldots, f_c)
$$
$R$-代数同态 $C \to B'$ 有一个由 $y_j \mapsto 0$ 给出的缩回。
我们断言，同态
$$
R[z_1, \ldots, z_n] \longrightarrow B',\quad
z_i \longmapsto x_i - \sum\nolimits_j c_{ij} y_j
$$
在 $\Spec(C) \to \Spec(B')$ 的像中每一点处都是平展的。
在 $\Omega_{B'/R[z_1, \ldots, z_n]}$ 中有
$$
0 =
\text{d}f_j - \sum\nolimits_i \frac{\partial f_j}{\partial x_i} \text{d}z_i
\equiv
\sum\nolimits_{i, \ell}
\frac{\partial f_j}{\partial x_i} c_{i\ell} \text{d}y_\ell
\equiv
\text{d}y_j \bmod (y_1, \ldots, y_c)\Omega_{B'/R[z_1, \ldots, z_n]}
$$
又因为模去
$\sum B'\text{d}y_j + (y_1, \ldots, y_c)\Omega_{B'/R[z_1, \ldots, z_n]}$
后有 $0 = \text{d}z_i = \text{d}x_i$，故
$$
\Omega_{B'/R[z_1, \ldots, z_n]}/
(y_1, \ldots, y_c)\Omega_{B'/R[z_1, \ldots, z_n]} = 0.
$$
由于 $\Omega_{B'/R[z_1, \ldots, z_n]}$ 是有限 $B'$-模，Nakayama 引理
给出某个 $g \in 1 + (y_1, \ldots, y_c)$，使得
$(\Omega_{B'/R[z_1, \ldots, z_n]})_g = 0$。这证明了
$R[z_1, \ldots, z_n] \to B'_g$ 非分歧，见《代数》定义
\ref{algebra-definition-unramified}。对任意以域 $k$ 为目标的环同态
$R \to k$，得到两个维数为 $n$ 的光滑 $k$-代数之间的非分歧环同态
$k[z_1, \ldots, z_n] \to (B'_g) \otimes_R k$。由《代数》引理
\ref{algebra-lemma-CM-over-regular-flat} 和
\ref{algebra-lemma-characterize-smooth-kbar}，
$k[z_1, \ldots, z_n] \to (B'_g) \otimes_R k$ 是平坦的。
再由逐纤维平坦性判别法（crit\`ere de platitude par fibre；见《代数》引理
\ref{algebra-lemma-criterion-flatness-fibre}），可知
$R[z_1, \ldots, z_n] \to B'_g$ 平坦。最后，《代数》引理
\ref{algebra-lemma-characterize-etale} 蕴含
$R[z_1, \ldots, z_n] \to B'_g$ 平展。令 $B = B'_g$。
注意 $C \to B$ 光滑且有缩回，故 $A \to B$ 也光滑且有缩回。
此外，$R[z_1, \ldots, z_n] \to B$ 平展。由《代数》引理
\ref{algebra-lemma-etale-standard-smooth}，可写成
$$
B = R[z_1, \ldots, z_n, w_1, \ldots, w_m]/(g_1, \ldots, g_m)
$$
其中 $\det(\partial g_j/\partial w_i)$ 在 $B$ 中可逆。引理得证。
\end{proof}

\begin{lemma}
\label{smoothing-lemma-colimit-standard-smooth}
设 $R \to \Lambda$ 为环同态。若 $\Lambda$ 是光滑 $R$-代数的滤过余极限，
则 $\Lambda$ 是标准光滑 $R$-代数的滤过余极限。
\end{lemma}

\begin{proof}
设 $A \to \Lambda$ 为 $R$-代数同态，其中 $A$ 是有限表示 $R$-代数。
按照《代数》引理 \ref{algebra-lemma-when-colimit}，只需把此同态经由某个
标准光滑代数分解。我们已知它可以分解为 $A \to B \to \Lambda$，其中
$B$ 在 $R$ 上光滑。由引理 \ref{smoothing-lemma-smooth-standard-smooth}，选取具有缩回
$C \to B$ 的 $R$-代数同态 $B \to C$，使 $C$ 在 $R$ 上标准光滑。
所求分解便是 $A \to B \to C \to B \to \Lambda$。
\end{proof}

\begin{lemma}
\label{smoothing-lemma-standard-smooth-include-generators}
设 $R \to A$ 为标准光滑环同态，$E \subset A$ 为有限子集，且 $|E| = n$。
则存在表示
$A = R[x_1, \ldots, x_{n + m}]/(f_1, \ldots, f_c)$，其中 $c \geq n$，
使 $\det(\partial f_j/\partial x_i)_{i, j = 1, \ldots, c}$ 在 $A$ 中可逆，
并且 $E$ 恰为 $x_1, \ldots, x_n$ 的同余类所成集合。
\end{lemma}

\begin{proof}
选取表示 $A = R[y_1, \ldots, y_m]/(g_1, \ldots, g_d)$，使得
$\det(\partial g_j/\partial y_i)_{i, j = 1, \ldots, d}$ 的像在 $A$ 中可逆。
选定枚举 $E = \{a_1, \ldots, a_n\}$，并选取
$h_i \in R[y_1, \ldots, y_m]$，使其在 $A$ 中的像为 $a_i$。
考虑表示
$$
A = R[x_1, \ldots, x_n, y_1, \ldots, y_m]/
(x_1 - h_1, \ldots, x_n - h_n, g_1, \ldots, g_d)
$$
并令 $c = n + d$。
\end{proof}

\begin{lemma}
\label{smoothing-lemma-compare-standard}
设 $R \to A$ 为有限表示环同态，$a \in A$。考虑关于 $a$ 的下列条件：
\begin{enumerate}
\item $A_a$ 在 $R$ 上光滑；
\item $A_a$ 在 $R$ 上光滑，且 $\Omega_{A_a/R}$ 稳定自由；
\item $A_a$ 在 $R$ 上光滑，且 $\Omega_{A_a/R}$ 自由；
\item $A_a$ 在 $R$ 上标准光滑；
\item $a$ 是 $A$ 中相对于 $R$ 的严格标准元；
\item $a$ 是 $A$ 中相对于 $R$ 的初等标准元。
\end{enumerate}
则有
\begin{enumerate}
\item[(a)] (4) $\Rightarrow$ (3) $\Rightarrow$ (2) $\Rightarrow$ (1),
\item[(b)] (6) $\Rightarrow$ (5),
\item[(c)] (6) $\Rightarrow$ (4),
\item[(d)] (5) $\Rightarrow$ (2),
\item[(e)] (2) $\Rightarrow$ 当 $e \geq e_0$ 时，$a^e$ 是 $A$ 中相对于
$R$ 的严格标准元；
\item[(f)] (4) $\Rightarrow$ 当 $e \geq e_0$ 时，$a^e$ 是 $A$ 中相对于
$R$ 的初等标准元。
\end{enumerate}
\end{lemma}

\begin{proof}
由定义和《代数》引理 \ref{algebra-lemma-standard-smooth}，(a) 显然。
由定义 \ref{smoothing-definition-strictly-standard}，(b) 显然。

\medskip\noindent
证明 (c)。选取表示 $A = R[x_1, \ldots, x_n]/(f_1, \ldots, f_m)$，
使 (\ref{smoothing-equation-elementary-standard-one}) 和
(\ref{smoothing-equation-elementary-standard-two}) 成立。
选取映到 $a$ 的 $h \in R[x_1, \ldots, x_n]$。于是
$$
A_a = R[x_0, x_1, \ldots, x_n]/(x_0h - 1, f_1, \ldots, f_m).
$$
记 $J = (x_0h - 1, f_1, \ldots, f_m)$。由
(\ref{smoothing-equation-elementary-standard-two})，$A_a$-模 $J/J^2$ 由
$x_0h - 1, f_1, \ldots, f_c$ 生成。因此，与《代数》引理
\ref{algebra-lemma-huber} 的证明相同，可以选取 $g \in 1 + J$，使得
$$
A_a = R[x_0, \ldots, x_n, x_{n + 1}]/
(x_0h - 1, f_1, \ldots, f_c, gx_{n + 1} - 1).
$$
此时，(\ref{smoothing-equation-elementary-standard-one}) 蕴含 $R \to A_a$
标准光滑（用坐标 $x_0, x_1, \ldots, x_c, x_{n + 1}$ 求偏导）。

\medskip\noindent
证明 (d)。选取表示 $A = R[x_1, \ldots, x_n]/(f_1, \ldots, f_m)$，
使 (\ref{smoothing-equation-strictly-standard-one}) 和
(\ref{smoothing-equation-strictly-standard-two}) 成立。记 $I = (f_1, \ldots, f_m)$。
由引理 \ref{smoothing-lemma-elkik}，我们已经知道 $A_a$ 在 $R$ 上光滑。
由引理 \ref{smoothing-lemma-parse-equation-strictly-standard-one}，
$(I/I^2)_a$ 是以 $f_1, \ldots, f_c$ 为基的自由模，并同构地映到
$\bigoplus A_a \text{d}x_i$ 的一个直和项。由于
$\Omega_{A_a/R} = (\Omega_{A/R})_a$ 是同态
$(I/I^2)_a \to \bigoplus A_a \text{d}x_i$ 的余核，故它稳定自由。

\medskip\noindent
证明 (e)。选取表示 $A = R[x_1, \ldots, x_n]/I$，其中 $I$ 有限生成。
由假设，存在短正合列
$$
0 \to (I/I^2)_a \to \bigoplus\nolimits_{i = 1, \ldots, n} A_a\text{d}x_i \to
\Omega_{A_a/R} \to 0
$$
且它分裂正合。因此 $(I/I^2)_a \oplus \Omega_{A_a/R}$ 是自由
$A_a$-模。由于 $\Omega_{A_a/R}$ 稳定自由，$(I/I^2)_a$ 也稳定自由。
于是，对某个 $r$，用表示
$A = R[x_1, \ldots, x_n, x_{n + 1}, \ldots, x_{n + r}]/J$ 代替上述表示，
其中 $J = (I, x_{n + 1}, \ldots, x_{n + r})$，便可使 $(J/J^2)_a$
为（有限）自由模。选取 $f_1, \ldots, f_c \in J$，使其像构成
$(J/J^2)_a$ 的一组基，再把它们扩充为一列生成元
$f_1, \ldots, f_m \in J$。考虑表示
$A = R[x_1, \ldots, x_{n + r}]/(f_1, \ldots, f_m)$。
由构造，当 $e$ 充分大时，$a^e$ 满足
(\ref{smoothing-equation-strictly-standard-two})。此外，由于
$(J/J^2)_a \to \bigoplus\nolimits_{i = 1, \ldots, n + r} A_a\text{d}x_i$
是分裂单射，可以找到一个 $A_a$-线性左逆。用基
$f_1, \ldots, f_c$ 写出此左逆并清去分母，得到线性同态
$\psi_0 : A^{\oplus n + r} \to A^{\oplus c}$，使得
$$
A^{\oplus c} \xrightarrow{(f_1, \ldots, f_c)}
J/J^2 \xrightarrow{f \mapsto \text{d}f}
\bigoplus\nolimits_{i = 1, \ldots, n + r} A \text{d}x_i
\xrightarrow{\psi_0}
A^{\oplus c}
$$
是乘以 $a^{e_0}$ 的同态，其中 $e_0 \geq 1$。由引理
\ref{smoothing-lemma-parse-equation-strictly-standard-one}，
(\ref{smoothing-equation-strictly-standard-one}) 对 $a^{ce_0}$ 成立，
因而对所有满足 $e \geq ce_0$ 的 $a^e$ 成立。

\medskip\noindent
证明 (f)。选取表示
$A_a = R[x_1, \ldots, x_n]/(f_1, \ldots, f_c)$，使得
$\det(\partial f_j/\partial x_i)_{i, j = 1, \ldots, c}$ 在 $A_a$ 中可逆。
由引理 \ref{smoothing-lemma-standard-smooth-include-generators}，可以假设对某个
$m < n$，元素 $x_1, \ldots, x_m$ 的类分别对应于 $a_i/1$，其中
$a_1, \ldots, a_m \in A$ 是 $A$ 在 $R$ 上的一组生成元。
对 $m < i \leq n$ 用 $a^Nx_i$ 代替 $x_i$ 后，可以假设 $x_i$ 的类
为某个 $a_i \in A$ 所给出的 $a_i/1 \in A_a$。考虑环同态
$$
\Psi : R[x_1, \ldots, x_n] \longrightarrow A,\quad
x_i \longmapsto a_i.
$$
这是满环同态。用 $a^Nf_j$ 代替 $f_j$ 后，可以假设
$f_j \in R[x_1, \ldots, x_n]$ 且 $\Psi(f_j) = 0$
（毕竟在 $A_a$ 中 $f_j(a_1/1, \ldots, a_n/1) = 0$）。令
$J = \Ker(\Psi)$。则 $A = R[x_1, \ldots, x_n]/J$ 是一个表示，
并且 $f_1, \ldots, f_c \in J$ 满足：$(J/J^2)_a$ 以
$f_1, \ldots, f_c$ 为自由基，而且
$\det(\partial f_j/\partial x_i)_{i, j = 1, \ldots, c}$ 的像在 $A_a$
中可逆。因此，对所有充分大的 $e$，$a^e$ 满足
(\ref{smoothing-equation-elementary-standard-one}) 和
(\ref{smoothing-equation-elementary-standard-two})，正是所需结论。
\end{proof}






\section{插曲：N\'eron 消奇异化}
\label{smoothing-section-neron}

\noindent
我们暂且中断对 Popescu 定理一般情形的攻克，转而研究一个较容易但已十分
有趣的情形，即 $R \to \Lambda$ 是离散赋值环之间的同态。
对此的讨论见 \cite[Section 4]{Artin-Algebraic-Approximation}。

\begin{situation}
\label{smoothing-situation-neron}
设 $R \subset \Lambda$ 是分歧指数为 $1$ 的离散赋值环扩张
（见《代数详论》定义
\ref{more-algebra-definition-extension-discrete-valuation-rings}）。
给定分解
$$
R \to A \xrightarrow{\varphi} \Lambda
$$
其中 $R \to A$ 平坦且为有限型。令 $\mathfrak q = \Ker(\varphi)$，
$\mathfrak p = \varphi^{-1}(\mathfrak m_\Lambda)$。
\end{situation}

\noindent
在情形 \ref{smoothing-situation-neron} 中，令 $\pi \in R$ 为一致化元。
回顾一下，$A$ 在 $R$ 上平坦意味着 $\pi$ 是 $A$ 上的非零因子
（见《代数详论》引理
\ref{more-algebra-lemma-valuation-ring-torsion-free-flat}）。
由关于 $R \subset \Lambda$ 的假设，$\pi$ 在 $\Lambda$ 中的像也是一致化元。
由于 $\pi \in \mathfrak p$，可以考虑 N\'eron 仿射吹起代数
（见《代数》第 \ref{algebra-section-blow-up} 节）
$$
\varphi' :
A' = A[\textstyle{\frac{\mathfrak p}{\pi}}]
\longrightarrow
\Lambda
$$
它带有到 $\Lambda$ 的诱导同态，把 $a/\pi^n$（其中
$a \in \mathfrak p^n$）映到 $\Lambda$ 中的 $\pi^{-n}\varphi(a)$。
以 $\mathfrak q' \subset \mathfrak p' \subset A'$ 记 $A'$ 中相应的素理想。
注意，$A'$ 的同构类不依赖于一致化元的选取。重复此构造得到序列
$$
A \to A' \to A'' \to \ldots \to \Lambda
$$

\begin{lemma}
\label{smoothing-lemma-neron-functorial}
在情形 \ref{smoothing-situation-neron} 中，N\'eron 吹起在下述意义下具有函子性：
\begin{enumerate}
\item 若 $a \in A$ 且 $a \not \in \mathfrak p$，则 $A_a$ 的 N\'eron 吹起
为 $A'_a$；
\item 若 $B \to A$ 是核为 $I$ 的平坦有限型 $R$-代数之间的满同态，
则 $A'$ 是 $B'/IB'$ 对其 $\pi$-幂挠子模取商所得的环。
\end{enumerate}
\end{lemma}

\begin{proof}
(1) 和 (2) 都是《代数》引理
\ref{algebra-lemma-blowup-base-change} 的特例。事实上，只要有
$A_1 \to A_2 \to \Lambda$ 且 $\mathfrak p_1 A_2 = \mathfrak p_2$，
$A_2'$ 就是 $A_1' \otimes_{A_1} A_2$ 对其 $\pi$-幂挠子模取商所得的环。
\end{proof}

\begin{lemma}
\label{smoothing-lemma-neron-blowup-smooth}
在情形 \ref{smoothing-situation-neron} 中，假设 $R \to A$ 在 $\mathfrak p$ 处光滑，
且 $R/\pi R \subset \Lambda/\pi \Lambda$ 是可分域扩张。
则 $R \to A'$ 在 $\mathfrak p'$ 处光滑，并且存在短正合列
$$
0 \to
\Omega_{A/R} \otimes_A A'_{\mathfrak p'} \to
\Omega_{A'/R, \mathfrak p'} \to
(A'/\pi A')_{\mathfrak p'}^{\oplus c} \to 0
$$
其中 $c = \dim((A/\pi A)_\mathfrak p)$。
\end{lemma}

\begin{proof}
由引理 \ref{smoothing-lemma-neron-functorial}，可以用 $A$ 关于某个不属于
$\mathfrak p$ 的元素的局部化代替 $A$；下文将不再特别说明这一点。
记 $\kappa = R/\pi R$。由于光滑性在基变换下保持
（见《代数》引理 \ref{algebra-lemma-base-change-smooth}），
$A/\pi A$ 在 $\mathfrak p$ 处是光滑 $\kappa$-代数。因此
$(A/\pi A)_\mathfrak p$ 是正则局部环
（见《代数》引理 \ref{algebra-lemma-characterize-smooth-over-field}）。
选取 $g_1, \ldots, g_c \in \mathfrak p$，使其像构成
$(A/\pi A)_\mathfrak p$ 的一组正则参数。必要时局部化 $A$ 后，有
$\mathfrak p = (\pi, g_1, \ldots, g_c)$。注意
$\pi, g_1, \ldots, g_c$ 是 $A_\mathfrak p$ 中的正则序列
（首先 $\pi$ 是非零因子，再对序列的其余部分应用《代数》引理
\ref{algebra-lemma-regular-ring-CM}）。再次局部化 $A$ 后，可以假设
$\pi, g_1, \ldots, g_c$ 是 $A$ 中的正则序列
（见《代数》引理 \ref{algebra-lemma-regular-sequence-in-neighbourhood}）。
由此
$$
A' = A[y_1, \ldots, y_c]/(\pi y_1 - g_1, \ldots, \pi y_c - g_c) =
A[y_1, \ldots, y_c]/I
$$
由《代数详论》引理 \ref{more-algebra-lemma-blowup-regular-sequence} 得出。
下文将直接使用以朴素余切复形表述的光滑性定义
（见《代数》定义 \ref{algebra-definition-smooth}）以及《代数》引理
\ref{algebra-lemma-smooth-at-point} 的判别准则，不再另行说明。
《代数》引理 \ref{algebra-lemma-exact-sequence-NL} 对
$R \to A[y_1, \ldots, y_c] \to A'$ 给出的正合列为
$$
0 \to H_1(\NL_{A'/R}) \to I/I^2 \to
\Omega_{A/R} \otimes_A A' \oplus
\bigoplus\nolimits_{i = 1, \ldots, c} A' \text{d}y_i \to
\Omega_{A'/R} \to 0
$$
其中 $I/I^2$ 中 $\pi y_i - g_i$ 的类映到下一项中的
$- \text{d}g_i + \pi \text{d}y_i$。这里用《代数》引理
\ref{algebra-lemma-NL-surjection} 计算
$\NL_{A'/A[y_1, \ldots, y_c]}$，并使用了
$R \to A[y_1, \ldots, y_c]$ 光滑这一事实。因此
$H_1(\NL_{A[y_1, \ldots, y_c]/R}) = 0$，而且
$\Omega_{A[y_1, \ldots, y_c]/R}$ 是有限投射（因而平坦）
$A[y_1, \ldots, y_c]$-模；事实上，它是
$\Omega_{A/R} \otimes_A A[y_1, \ldots, y_c]$ 与以 $\text{d}y_i$
为基的自由模之直和。为完成证明，只需说明
$\text{d}g_1, \ldots, \text{d}g_c$ 是有限自由模
$\Omega_{A/R, \mathfrak p}$ 的一组基的一部分。事实上，这将说明
$(I/I^2)_\mathfrak p$ 是以 $\pi y_i - g_i$ 为基的自由模，
正合列中间同态在 $\mathfrak p$ 处的局部化是单射，因而
$H_1(\NL_{A'/R})_\mathfrak p = 0$，并且余核
$\Omega_{A'/R, \mathfrak p}$ 有限自由。为此，只需证明
$\text{d}g_i$ 的像在
$\Omega_{A/R, \mathfrak p}/\pi = \Omega_{(A/\pi A)/\kappa, \mathfrak p}$
中在 $\kappa(\mathfrak p)$ 上线性无关（该等式由《代数》引理
\ref{algebra-lemma-differentials-base-change} 给出）。由于
$\kappa \subset \kappa(\mathfrak p) \subset \Lambda/\pi \Lambda$，
$\kappa(\mathfrak p)$ 在 $\kappa$ 上可分
（见《代数》定义 \ref{algebra-definition-separable-field-extension}）。
所需的线性无关性由《代数》引理
\ref{algebra-lemma-computation-differential} 得出。
\end{proof}

\begin{lemma}
\label{smoothing-lemma-neron-when-smooth}
在情形 \ref{smoothing-situation-neron} 中，假设 $R \to A$ 在 $\mathfrak q$ 处光滑，
并有核为 $I$ 的满 $R$-代数同态 $B \to A$。假设 $R \to B$ 在
$\mathfrak p_B = (B \to A)^{-1}\mathfrak p$ 处光滑。若同态
$$
I/I^2 \otimes_A \Lambda \to \Omega_{B/R} \otimes_B \Lambda
$$
的余核为自由 $\Lambda$-模，则 $R \to A$ 在 $\mathfrak p$ 处光滑。
\end{lemma}

\begin{proof}
同态 $I/I^2 \to \Omega_{B/R} \otimes_B A$ 的余核为 $\Omega_{A/R}$，
见《代数》引理 \ref{algebra-lemma-differential-seq}。令
$d = \dim_\mathfrak q(A/R)$ 为 $R \to A$ 在 $\mathfrak q$ 处的相对维数，
即 $\Spec(A[1/\pi])$ 在 $\mathfrak q$ 处的维数；见《代数》定义
\ref{algebra-definition-relative-dimension}。于是
$\Omega_{A/R, \mathfrak q}$ 是秩为 $d$ 的自由 $A_\mathfrak q$-模
（见《代数》引理 \ref{algebra-lemma-characterize-smooth-over-field}）。
因此，若引理的假设成立，则 $\Omega_{A/R} \otimes_A \Lambda$
是秩为 $d$ 的自由模。从而 $\Omega_{A/R} \otimes_A \kappa(\mathfrak p)$
的维数为 $d$（与 $\Lambda/\pi \Lambda$ 张量后即可看出）。
由于 $R \to A$ 平坦，且 $\mathfrak p$ 是 $\mathfrak q$ 的特化，
《代数》引理 \ref{algebra-lemma-dimension-fibres-bounded-open-upstairs}
给出 $\dim_\mathfrak p(A/R) \geq d$。最后，由《代数》引理
\ref{algebra-lemma-flat-fibre-smooth} 和
\ref{algebra-lemma-characterize-smooth-over-field}，
$R \to A$ 在 $\mathfrak p$ 处光滑。
\end{proof}

\begin{lemma}
\label{smoothing-lemma-neron-desingularization}
在情形 \ref{smoothing-situation-neron} 中，假设 $R \to A$ 在 $\mathfrak q$ 处光滑，
且 $R/\pi R \subset \Lambda/\pi \Lambda$ 是可分域扩张。
则经过有限次仿射 N\'eron 吹起后，代数 $A$ 在 $\mathfrak p$ 处变为
光滑 $R$-代数。
\end{lemma}

\begin{proof}
选取 $R$-代数 $B$ 和满同态 $B \to A$。令
$\mathfrak p_B = (B \to A)^{-1}(\mathfrak p)$，并以 $r$ 记
$R \to B$ 在 $\mathfrak p_B$ 处的相对维数。选取 $B$，使得
$R \to B$ 在 $\mathfrak p_B$ 处光滑。例如，可取 $B$ 为 $R$ 上有
$r$ 个变量的多项式代数。考虑引理 \ref{smoothing-lemma-neron-when-smooth} 中的复形
$$
I/I^2 \otimes_A \Lambda \longrightarrow \Omega_{B/R} \otimes_B \Lambda
$$
由有限 $\Lambda$-模的结构
（见《代数详论》引理 \ref{more-algebra-lemma-modules-PID}），
其余核形如
$$
\Lambda^{\oplus d} \oplus
\bigoplus\nolimits_{i = 1, \ldots, n} \Lambda/\pi^{e_i} \Lambda
$$
其中 $d \geq 0$、$n \geq 0$ 且 $e_i \geq 1$。注意，$d$ 是 $A/R$
在 $\mathfrak q$ 处的相对维数
（见《代数》引理 \ref{algebra-lemma-characterize-smooth-over-field}）。
若缺陷 $e = \sum_{i = 1, \ldots, n} e_i$ 为零，
则引理 \ref{smoothing-lemma-neron-when-smooth} 已给出结论。

\medskip\noindent
下面考察施行 N\'eron 吹起后发生的情形。回顾一下，$A'$ 是
$B'/IB'$ 对其 $\pi$-幂挠子模取商所得的环
（见引理 \ref{smoothing-lemma-neron-functorial}），而 $R \to B'$ 在
$\mathfrak p_{B'}$ 处光滑（见引理 \ref{smoothing-lemma-neron-blowup-smooth}）。
因此吹起后得到完全相同的设定。图示如下：
$$
\xymatrix{
0 \ar[r] & I' \ar[r] & B' \ar[r] & A' \ar[r] & 0 \\
0 \ar[r] & I \ar[u] \ar[r] & B \ar[u] \ar[r] & A \ar[r] \ar[u] & 0
}
$$
由于 $I \subset \mathfrak p_B$，同态 $I \to I'$ 经由 $\pi I'$ 分解。
考察所诱导的复形同态，得到
$$
\xymatrix{
I'/(I')^2 \otimes_{A'} \Lambda \ar[r] &
\Omega_{B'/R} \otimes_{B'} \Lambda \ar@{=}[r] & M' \\
I/I^2 \otimes_A \Lambda \ar[r] \ar[u] &
\Omega_{B/R} \otimes_B \Lambda \ar[u] \ar@{=}[r] & M
}
$$
此时 $M \subset M'$ 是有限自由 $\Lambda$-模，且商 $M'/M$ 被 $\pi$ 零化，
见引理 \ref{smoothing-lemma-neron-blowup-smooth}。令 $N \subset M$ 与
$N' \subset M'$ 分别为水平同态的像，并记 $Q = M/N$、$Q' = M'/N'$。
得到交换图
$$
\xymatrix{
0 \ar[r] &
N' \ar[r] &
M' \ar[r] &
Q' \ar[r] &
0 \\
0 \ar[r] &
N \ar[r] \ar[u] &
M \ar[r] \ar[u] &
Q \ar[r] \ar[u] &
0
}
$$
其中 $N \subset N'$ 是秩为 $r - d$ 的自由 $\Lambda$-模。
由于 $I$ 映入 $\pi I'$，有 $N \subset \pi N'$。

\medskip\noindent
令 $K = \Lambda_\pi$ 为 $\Lambda$ 的分式域。存在交换图
$$
\xymatrix{
0 \ar[r] &
N' \ar[r] &
N'_K \cap M' \ar[r] &
Q'_{tor} \ar[r] &
0 \\
0 \ar[r] &
N \ar[r] \ar[u] &
N_K \cap M \ar[r] \ar[u] &
Q_{tor} \ar[r] \ar[u] &
0
}
$$
其两行都是短正合列。这表明缺陷的变化量为
$$
e - e' =
\text{length}(Q_{tor}) - \text{length}(Q'_{tor})
=
\text{length}(N'/N) - \text{length}(N'_K \cap M' / N_K \cap M)
$$
由于 $M'/M$ 被 $\pi$ 零化，$N'_K \cap M' / N_K \cap M$ 也被
$\pi$ 零化，故其长度至多为 $\dim_K(N_K)$。又因
$N \subset \pi N'$，有 $\text{length}(N'/N) \ge \dim_K(N_K)$，
而等号成立当且仅当 $N = \pi N'$。

\medskip\noindent
为完成证明，需要说明：当 $A$ 在 $\mathfrak p$ 处不光滑时，
$N$ 严格小于 $\pi N'$；这正是 N\'eron 论证中的关键计算。
为此，考虑正合列
$$
I/I^2 \otimes_B \kappa(\mathfrak p_B)
\to \Omega_{B/R} \otimes_B \kappa(\mathfrak p_B)
\to \Omega_{A/R} \otimes_A \kappa(\mathfrak p) \to 0
$$
（由《代数》引理 \ref{algebra-lemma-differential-seq} 得出）。
由于 $R \to A$ 在 $\mathfrak p$ 处不光滑，
$\Omega_{A/R} \otimes_A \kappa(\mathfrak p)$ 的维数 $s$ 大于 $d$。
另一方面，第一个同态经由单射
$$
\mathfrak p B_\mathfrak p/\mathfrak p^2 B_\mathfrak p
\to \Omega_{B/R} \otimes_B \kappa(\mathfrak p_B)
$$
分解；该单射来自《代数》引理
\ref{algebra-lemma-computation-differential}。注意，由关于
$R/\pi R \subset \Lambda/\pi \Lambda$ 的假设，
$\kappa(\mathfrak p)$ 在 $k$ 上可分。因此可以找到生成元
$g_1, \ldots, g_t \in I$，使得当 $j > r - s$ 时
$g_j \in \mathfrak p^2$。于是当 $j > r - s$ 时，$g_j$ 在 $A'$ 中的像
属于 $\pi^2 I'$。由于 $r - s < r - d$，$N$ 的至少一个极小生成元
在 $N'$ 中变得可被 $\pi^2$ 整除。因此缺陷 $e$ 至少减少 $1$，证明完成。
\end{proof}

\noindent
若 $R \to \Lambda$ 是离散赋值环扩张，则 $R \to \Lambda$ 正则当且仅当：
(a) 分歧指数为 $1$；(b) 分式域扩张可分；并且
(c) $R/\mathfrak m_R \subset \Lambda/\mathfrak m_\Lambda$ 可分。
因此，以下结果是定理 \ref{smoothing-theorem-popescu} 中一般 N\'eron 消奇异化的特例。

\begin{lemma}
\label{smoothing-lemma-neron-colimit}
\begin{slogan}
DVR 的非分歧扩张是 ind-smooth 的，亦即 N\'eron 消奇异化
\end{slogan}
设 $R \subset \Lambda$ 是分歧指数为 $1$ 的离散赋值环扩张，
并且它所诱导的剩余域扩张和分式域扩张均可分。
则 $\Lambda$ 是光滑 $R$-代数的滤过余极限。
\end{lemma}

\begin{proof}
由《代数》引理 \ref{algebra-lemma-when-colimit}，只需证明情形
\ref{smoothing-situation-neron} 中的任意 $R \to A \to \Lambda$ 都可分解为
$A \to B \to \Lambda$，其中 $B$ 是光滑 $R$-代数。用 $A$ 在
$\Lambda$ 中的像代替 $A$ 后，可以假设 $A$ 是整环，其分式域 $K$
是 $\Lambda$ 的分式域的子域。特别地，由假设，$A$ 在 $R$ 的分式域上可分。
于是由《代数》引理 \ref{algebra-lemma-smooth-at-generic-point}，
$R \to A$ 在 $\mathfrak q = (0)$ 处光滑。经有限次 N\'eron 吹起后，
可以假设 $R \to A$ 在 $\mathfrak p$ 处光滑，见引理
\ref{smoothing-lemma-neron-desingularization}。随后，对某个
$a \in A$、$a \not \in \mathfrak p$，用其局部化代替 $A$，即可使
$A$ 在 $R$ 上光滑。引理得证。
\end{proof}












\section{提升问题}
\label{smoothing-section-lifting}

\noindent
本节的目标是证明（命题 \ref{smoothing-proposition-lift}）：由光滑代数的滤过余极限
构成的代数类在平坦无穷小形变下封闭。证明是初等的，只用到第
\ref{smoothing-section-presentations} 节关于光滑代数表示的结果。

\begin{lemma}
\label{smoothing-lemma-lift-once}
设 $R \to \Lambda$ 为环同态，$I \subset R$ 为理想。假设
\begin{enumerate}
\item $I^2 = 0$；
\item $\Lambda/I\Lambda$ 是光滑 $R/I$-代数的滤过余极限。
\end{enumerate}
设 $\varphi : A \to \Lambda$ 为 $R$-代数同态，其中 $A$ 是有限表示
$R$-代数。则存在分解
$$
A \to B/J \to \Lambda
$$
其中 $B$ 是光滑 $R$-代数，而 $J \subset IB$ 是有限生成理想。
\end{lemma}

\begin{proof}
选取分解
$$
A/IA \to \bar B \to \Lambda/I\Lambda
$$
其中 $\bar B$ 在 $R/I$ 上标准光滑；由假设和引理
\ref{smoothing-lemma-colimit-standard-smooth}，这样的分解存在。写成
$$
\bar B = A/IA[t_1, \ldots, t_r]/(\bar g_1, \ldots, \bar g_s)
$$
并设 $\bar B \to \Lambda/I\Lambda$ 把 $t_i$ 映到 $\lambda_i$ 模
$I\Lambda$ 的类。选取 $\bar g_1, \ldots, \bar g_s$ 的提升
$g_1, \ldots, g_s \in A[t_1, \ldots, t_r]$。写成
$\varphi(g_i)(\lambda_1, \ldots, \lambda_r) =
\sum \epsilon_{ij} \mu_{ij}$
其中某些 $\epsilon_{ij} \in I$，$\mu_{ij} \in \Lambda$。定义
$$
A' = A[t_1, \ldots, t_r, \delta_{i, j}]/
(g_i - \sum \epsilon_{ij} \delta_{ij})
$$
并考虑同态
$$
A' \longrightarrow \Lambda,\quad
a \longmapsto \varphi(a),\quad
t_i \longmapsto \lambda_i,\quad
\delta_{ij} \longmapsto \mu_{ij}
$$
有
$$
A'/IA' = A/IA[t_1, \ldots, t_r]/(\bar g_1, \ldots, \bar g_s)[\delta_{ij}]
\cong \bar B[\delta_{ij}]
$$
由于 $\bar B$ 标准光滑，这是一个标准光滑 $R/I$-代数。选取表示
$A'/IA' = R/I[x_1, \ldots, x_n]/(\bar f_1, \ldots, \bar f_c)$，使得
$\det(\partial \bar f_j/\partial x_i)_{i, j = 1, \ldots, c}$ 在
$A'/IA'$ 中可逆。选取 $\bar f_1, \ldots, \bar f_c$ 的提升
$f_1, \ldots, f_c \in R[x_1, \ldots, x_n]$。于是
$$
B = R[x_1, \ldots, x_n, x_{n + 1}]/
(f_1, \ldots, f_c,
x_{n + 1}\det(\partial f_j/\partial x_i)_{i, j = 1, \ldots, c} - 1)
$$
在 $R$ 上光滑。由于光滑环同态形式光滑
（见《代数》命题 \ref{algebra-proposition-smooth-formally-smooth}），
存在模 $I$ 后为同构的 $R$-代数同态 $B \to A'$。再由 Nakayama 引理
（《代数》引理 \ref{algebra-lemma-NAK}），$B \to A'$ 是满射。
因此 $A' = B/J$，其中 $J \subset IB$ 有限生成
（见《代数》引理 \ref{algebra-lemma-finite-presentation-independent}）。
\end{proof}

\begin{lemma}
\label{smoothing-lemma-lift-twice}
设 $R \to \Lambda$ 为环同态，$I \subset R$ 为理想。假设
\begin{enumerate}
\item $I^2 = 0$；
\item $\Lambda/I\Lambda$ 是光滑 $R/I$-代数的滤过余极限；
\item $R \to \Lambda$ 平坦。
\end{enumerate}
设 $\varphi : B \to \Lambda$ 为 $R$-代数同态，其中 $B$ 在 $R$ 上光滑。
设 $J \subset IB$ 为有限生成理想，且 $\varphi(J) = 0$。
则存在 $R$-代数同态
$$
B \xrightarrow{\alpha} B' \xrightarrow{\beta} \Lambda
$$
使得 $B'$ 在 $R$ 上光滑、$\alpha(J) = 0$，并且
$\beta \circ \alpha = \varphi$。
\end{lemma}

\begin{proof}
只要能对 $J = (h)$ 的情形证明引理，便可对 $J$ 的生成元个数作归纳。
具体地，假设 $J$ 可由 $n$ 个元素 $h_1, \ldots, h_n$ 生成，且引理对
$J$ 由 $n - 1$ 个元素生成的一切情形成立。先应用 $n = 1$ 的情形，
得到 $B \to B' \to \Lambda$，其中第一个同态把 $h_n$ 映为零。
再令 $J'$ 为 $B'$ 中由 $h_1, \ldots, h_{n - 1}$ 的像生成的理想，
应用 $n - 1$ 的情形得到 $B' \to B'' \to \Lambda$。
容易验证 $B \to B'' \to \Lambda$ 满足要求。

\medskip\noindent
假设 $J = (h)$，并把 $h$ 写成 $h = \sum \epsilon_i b_i$，其中
$\epsilon_i \in I$、$b_i \in B$。注意
$0 = \varphi(h) = \sum \epsilon_i \varphi(b_i)$。由于 $\Lambda$ 在
$R$ 上平坦，平坦性的方程判别准则
（见《代数》引理 \ref{algebra-lemma-flat-eq}）表明，存在
$\lambda_j \in \Lambda$（$j = 1, \ldots, m$）和 $a_{ij} \in R$，使得
$\varphi(b_i) = \sum_j a_{ij} \lambda_j$ 且
$\sum_i \epsilon_i a_{ij} = 0$。令
$$
C = B[x_1, \ldots, x_m]/(b_i - \sum a_{ij} x_j)
$$
并以 $\varphi$ 和 $x_j \mapsto \lambda_j$ 定义 $C \to \Lambda$。
按引理 \ref{smoothing-lemma-lift-once} 选取分解
$$
C \to B'/J' \to \Lambda
$$
由于 $B$ 在 $R$ 上光滑，可以把同态 $B \to C \to B'/J'$ 提升为
$\alpha : B \to B'$。于是 $\varphi = \beta \circ \alpha$。
为完成证明，只需验证 $\alpha(h) = 0$。事实上，$\alpha$ 是
$B \to C \to B'/J'$ 的提升这一事实蕴含
$$
\alpha(b_i) = \sum a_{ij} \xi_j + \theta_i
$$
其中某些 $\xi_j \in B'$，$\theta_i \in J' \subset IB'$。因此
$$
\alpha(h) = \alpha(\sum \epsilon_i b_i) =
\sum \epsilon_i a_{ij} \xi_j + \sum \epsilon_i \theta_i = 0
$$
这是因为上述关系成立且 $I^2 = 0$。
\end{proof}

\begin{proposition}
\label{smoothing-proposition-lift}
\begin{slogan}
代数的 ind-光滑性在无穷小形变下保持
\end{slogan}
设 $R \to \Lambda$ 为环同态，$I \subset R$ 为理想。假设
\begin{enumerate}
\item $I$ 是幂零理想；
\item $\Lambda/I\Lambda$ 是光滑 $R/I$-代数的滤过余极限；
\item $R \to \Lambda$ 平坦。
\end{enumerate}
则 $\Lambda$ 是光滑 $R$-代数的滤过余极限。
\end{proposition}

\begin{proof}
对某个 $n$ 有 $I^n = 0$，故对 $n$ 作归纳可知，只需考虑 $I^2 = 0$
的情形。设 $\varphi : A \to \Lambda$ 为 $R$-代数同态，其中 $A$ 是
有限表示 $R$-代数。由《代数》引理 \ref{algebra-lemma-when-colimit}，
需要找到分解 $A \to B \to \Lambda$，其中 $B$ 在 $R$ 上光滑。
由引理 \ref{smoothing-lemma-lift-once}，可以假设 $A = B/J$，其中 $B$ 在 $R$
上光滑，而 $J \subset IB$ 是有限生成理想。由引理
\ref{smoothing-lemma-lift-twice}，可以找到交换图
$$
\xymatrix{
B \ar[rr]_\alpha \ar[rd]_\varphi & & B' \ar[ld]^\beta \\
& \Lambda
}
$$
其中各同态均为 $R$-代数同态，$B'$ 在 $R$ 上光滑，且
$\alpha(J) = 0$。所以 $\alpha$ 分解为 $B \to A \to B'$，证明完成。
\end{proof}






\section{提升引理}
\label{smoothing-section-lifting-lemma}

\noindent
下面是一个极其巧妙的引理。

\begin{lemma}
\label{smoothing-lemma-lifting}
设 $R$ 为诺特环，$\Lambda$ 为 $R$-代数。设 $\pi \in R$，并假设
$\text{Ann}_R(\pi) = \text{Ann}_R(\pi^2)$ 且
$\text{Ann}_\Lambda(\pi) = \text{Ann}_\Lambda(\pi^2)$。
假设有 $R$-代数同态
$R/\pi^2R \to \bar C \to \Lambda/\pi^2\Lambda$，其中 $\bar C$ 有限表示。
则存在 $R$-代数同态 $D \to \Lambda$ 及交换图
$$
\xymatrix{
R/\pi^2R \ar[r] \ar[d] &
\bar C \ar[r] \ar[d] &
\Lambda/\pi^2\Lambda \ar[d] \\
R/\pi R \ar[r] &
D/\pi D \ar[r] &
\Lambda/\pi \Lambda
}
$$
满足以下性质：
\begin{enumerate}
\item[(a)] $D$ 有限表示；
\item[(b)] 对任意满足 $\pi \not \in \mathfrak q$ 的素理想 $\mathfrak q$，
$R \to D$ 在 $\mathfrak q$ 处光滑；
\item[(c)] 对任意满足 $\pi \in \mathfrak q$、且位于 $\bar C$ 中一个
$R/\pi^2 R \to \bar C$ 光滑的素理想之上的素理想 $\mathfrak q$，
$R \to D$ 在 $\mathfrak q$ 处光滑；
\item[(d)] 对任意位于 $\bar C$ 中一个 $R/\pi^2R \to \bar C$ 光滑的
素理想之上的素理想，$\bar C/\pi \bar C \to D/\pi D$ 在该处光滑。
\end{enumerate}
\end{lemma}

\begin{proof}
选取表示
$$
\bar C = R[x_1, \ldots, x_n]/(f_1, \ldots, f_m)
$$
再记 $I = (f_1, \ldots, f_m)$，并以 $\bar I$ 记 $I$ 在
$R/\pi^2R[x_1, \ldots, x_n]$ 中的像。由于 $R$ 是诺特环，$\bar C$
也是诺特环。因此 $R/\pi^2 R \to \bar C$ 的光滑轨迹拟紧，见
《拓扑》引理 \ref{topology-lemma-Noetherian}。应用引理
\ref{smoothing-lemma-find-strictly-standard}，可以选取有限个元素
$a_1, \ldots, a_r \in R[x_1, \ldots, x_n]$，使得
\begin{enumerate}
\item 开子空间 $\Spec(\bar C_{a_k}) \subset \Spec(\bar C)$ 的并覆盖
$R/\pi^2 R \to \bar C$ 的光滑轨迹；
\item 对每个 $k = 1, \ldots, r$，存在有限子集
$E_k \subset \{1, \ldots, m\}$，使得 $(\bar I/\bar I^2)_{a_k}$
以 $f_j$（$j \in E_k$）的类为自由基。
\end{enumerate}
令 $I_k = (f_j, j \in E_k) \subset I$，并以 $\bar I_k$ 记 $I_k$ 在
$R/\pi^2R[x_1, \ldots, x_n]$ 中的像。由 (2) 和 Nakayama 引理，
对某个 $b'_k \in \bar I_{a_k}$，$(\bar I/\bar I_k)_{a_k}$ 被
$1 + b'_k$ 零化。设 $b'_k$ 是某个 $b_k \in I$ 和某个整数 $N$ 所给出的
$b_k/(a_k)^N$ 的像。用 $a_kb_k$ 代替 $a_k$ 后，得到
\begin{enumerate}
\item[(3)] $(\bar I_k)_{a_k} = (\bar I)_{a_k}$.
\end{enumerate}
因此，必要时用 $a_k$ 的一个高次幂代替 $a_k$ 后，可以写成
\begin{enumerate}
\item[(4)]
$a_k f_\ell = \sum\nolimits_{j \in E_k} h_{k, \ell}^jf_j + \pi^2 g_{k, \ell}$
\end{enumerate}
其中 $\ell \in \{1, \ldots, m\}$ 任意，而某些
$h_{i, \ell}^j, g_{i, \ell} \in R[x_1, \ldots, x_n]$。
若 $\ell \in E_k$，则取 $h_{k, \ell}^j = a_k\delta_{\ell, j}$
（Kronecker delta）和 $g_{k, \ell} = 0$。令
$$
D = R[x_1, \ldots, x_n, z_1, \ldots, z_m]/
(f_j - \pi z_j, p_{k, \ell}).
$$
这里 $j \in \{1, \ldots, m\}$、$k \in \{1, \ldots, r\}$、
$\ell \in \{1, \ldots, m\}$，并且
$$
p_{k, \ell} = a_k z_\ell - \sum\nolimits_{j \in E_k} h_{k, \ell}^j z_j
- \pi g_{k, \ell}.
$$
注意，由上述选取，当 $\ell \in E_k$ 时 $p_{k, \ell} = 0$。

\medskip\noindent
同态 $R \to D$ 取为给定同态。设 $\bar C \to \Lambda/\pi^2\Lambda$
把 $x_i$ 映到 $\lambda_i$ 模 $\pi^2$ 的类。对元素
$f \in R[x_1, \ldots, x_n]$，以 $f(\lambda) \in \Lambda$ 记把
$\lambda_i$ 代入 $x_i$ 后的结果。于是对某些 $\mu_j \in \Lambda$，有
$f_j(\lambda) = \pi^2 \mu_j$。按 $x_i \mapsto \lambda_i$ 和
$z_j \mapsto \pi\mu_j$ 定义 $D \to \Lambda$。这是良定义的，因为
\begin{align*}
p_{k, \ell} & \mapsto
a_k(\lambda) \pi \mu_\ell -
\sum\nolimits_{j \in E_k} h_{k, \ell}^j(\lambda) \pi \mu_j
- \pi g_{k, \ell}(\lambda) \\
& =
\pi\left(a_k(\lambda) \mu_\ell -
\sum\nolimits_{j \in E_k} h_{k, \ell}^j(\lambda) \mu_j
- g_{k, \ell}(\lambda)\right)
\end{align*}
在上述 (4) 中代入 $x_i = \lambda_i$，可见括号内的表达式被 $\pi^2$
零化；由于假设 $\text{Ann}_\Lambda(\pi) = \text{Ann}_\Lambda(\pi^2)$，
它也被 $\pi$ 零化。同态 $\bar C \to D/\pi D$ 由
$x_i \mapsto x_i$ 决定（显然良定义）。因此，只须证明 (b)、(c) 和 (d)。

\medskip\noindent
利用 (4) 得到以下关键等式：
\begin{align*}
\pi p_{k, \ell} & =
\pi a_k z_\ell - \sum\nolimits_{j \in E_k} \pi h_{k, \ell}^jz_j
- \pi^2 g_{k, \ell} \\
& =
- a_k (f_\ell - \pi z_\ell) + a_k f_\ell +
\sum\nolimits_{j \in E_k} h_{k, \ell}^j (f_j - \pi z_j) -
\sum\nolimits_{j \in E_k} h_{k, \ell}^j f_j - \pi^2 g_{k, \ell} \\
& =
-a_k(f_\ell - \pi z_\ell) +
\sum\nolimits_{j \in E_k} h_{k, \ell}^j(f_j - \pi z_j)
\end{align*}
最后所得表达式属于由 $f_j - \pi z_j$ 生成的理想。特别地，
$D[1/\pi]$ 同构于
$R[1/\pi][x_1, \ldots, x_n, z_1, \ldots, z_m]/(f_j - \pi z_j)$，
后者又同构于 $R[1/\pi][x_1, \ldots, x_n]$，因而在 $R$ 上光滑。
这证明了 (b)。

\medskip\noindent
固定 $k \in \{1, \ldots, r\}$，考虑环
$$
D_k = R[x_1, \ldots, x_n, z_1, \ldots, z_m]/
(f_j - \pi z_j, j \in E_k, p_{k, \ell})
$$
由于当 $\ell \in E_k$ 时 $p_{k, \ell}$ 为零，方程的数目为
$m = |E_k| + (m - |E_k|)$。还要注意
\begin{align*}
(D_k/\pi D_k)_{a_k}
& =
R/\pi R[x_1, \ldots, x_n, 1/a_k, z_1, \ldots, z_m]/
(f_j, j \in E_k, p_{k, \ell}) \\
& =
(\bar C/\pi \bar C)_{a_k}[z_1, \ldots, z_m]/
(a_kz_\ell - \sum\nolimits_{j \in E_k} h_{k, \ell}^j z_j) \\
& \cong
(\bar C/\pi \bar C)_{a_k}[z_j, j \in E_k]
\end{align*}
特别地，$(D_k/\pi D_k)_{a_k}$ 在 $(\bar C/\pi \bar C)_{a_k}$ 上光滑。
由我们对 $a_k$ 的选择，$(\bar C/\pi \bar C)_{a_k}$ 在
$R/\pi R$ 上光滑且相对维数为 $n - |E_k|$，见 (2)。因此，若素理想
$\mathfrak q_k \subset D_k$ 含有 $\pi$ 且位于
$\Spec(\bar C_{a_k})$ 之上，则 $R \to D_k$ 的纤维环
在 $\mathfrak q_k$ 处光滑且维数为 $n$。于是，由上述方程数目的计算，
$R \to D_k$ 在 $\mathfrak q_k$ 处为 syntomic，见《代数》
引理 \ref{algebra-lemma-localize-relative-complete-intersection}。
因而 $R \to D_k$ 在 $\mathfrak q_k$ 处光滑，见《代数》
引理 \ref{algebra-lemma-flat-fibre-smooth}。

\medskip\noindent
为完成证明，设 $\mathfrak q \subset D$ 是一个含有 $\pi$ 的素理想，
它位于某个使 $R/\pi^2 R \to \bar C$ 光滑的素理想之上。
由 (1)，对某个 $k$ 有 $a_k \not \in \mathfrak q$。
我们将证明满射 $D_k \to D$ 在 $\mathfrak q$ 处诱导局部环的同构。
由上一段，我们已经知道环同态
$\bar C/\pi \bar C \to D_k/\pi D_k$ 和 $R \to D_k$
在相应的素理想 $\mathfrak q_k$ 处光滑。这将证明 (c) 和 (d)，
从而完成整个证明。

\medskip\noindent
首先注意，对任意 $\ell$，上面证明的等式
$\pi p_{k, \ell} = -a_k(f_\ell - \pi z_\ell) +
\sum_{j \in E_k} h_{k, \ell}^j (f_j - \pi z_j)$ 表明
$f_\ell - \pi z_\ell$ 在 $(D_k)_{a_k}$ 中的像为零，因而它在
$(D_k)_{\mathfrak q_k}$ 中的像尤其为零。
关系式 (4) 蕴含在 $I/I^2$ 中有 $a_k f_\ell =
\sum_{j \in E_k} h_{k, \ell}^j f_j$。
由于 $(\bar I_k/\bar I_k^2)_{a_k}$ 是以
$f_j$（$j \in E_k$）为基的自由模，可知
$$
a_{k'} h_{k, \ell}^j -
\sum\nolimits_{j' \in E_{k'}} h_{k', \ell}^{j'} h_{k, j'}^j
$$
对每个 $k,k',\ell$ 和 $j \in E_k$，这个元素在 $\bar C_{a_k}$ 中为零。
因此，可以取充分大的整数 $N$，使得
$$
a_k^N\left(
a_{k'} h_{k, \ell}^j -
\sum\nolimits_{j' \in E_{k'}} h_{k', \ell}^{j'} h_{k, j'}^j
\right)
$$
属于 $I_k + \pi^2R[x_1, \ldots, x_n]$。模 $\pi$ 计算可得
\begin{align*}
&
a_kp_{k', \ell} - a_{k'}p_{k, \ell} + \sum h_{k', \ell}^{j'} p_{k, j'}
\\
&
=
- a_k \sum h_{k', \ell}^{j'} z_{j'}
+ a_{k'} \sum h_{k, \ell}^j z_j
+ \sum h_{k', \ell}^{j'} a_k z_{j'}
- \sum \sum h_{k', \ell}^{j'} h_{k, j'}^j z_j \\
&
=
\sum \left(
a_{k'} h_{k, \ell}^j
- \sum h_{k', \ell}^{j'} h_{k, j'}^j
\right) z_j
\end{align*}
其中采用爱因斯坦求和约定。结合前述结果可知，在
$R[x_1, \ldots, x_n, z_1, \ldots, z_m]$ 中，
$a_k^{N + 1} p_{k', \ell}$ 属于由 $I_k$ 和 $\pi$ 生成的理想。
因此 $p_{k', \ell}$ 的像属于 $\pi (D_k)_{a_k}$。另一方面，等式
$$
\pi p_{k', \ell} =
-a_{k'} (f_\ell - \pi z_\ell) +
\sum\nolimits_{j' \in E_{k'}} h_{k', \ell}^{j'}(f_{j'} - \pi z_{j'})
$$
表明 $\pi p_{k', \ell}$ 在 $(D_k)_{a_k}$ 中为零。
我们已经假设 $\text{Ann}_R(\pi) = \text{Ann}_R(\pi^2)$，
而 $(D_k)_{\mathfrak q_k}$ 在 $R$ 上光滑，因而在 $R$ 上平坦，故有
$\text{Ann}_{(D_k)_{\mathfrak q_k}}(\pi) =
\text{Ann}_{(D_k)_{\mathfrak q_k}}(\pi^2)$.
于是 $p_{k', \ell}$ 的像也为零，从而
$D_{\mathfrak q} = (D_k)_{\mathfrak q_k}$，证毕。
\end{proof}






\section{消奇异化引理}
\label{smoothing-section-desingularization-lemma}

\noindent
下面是另一个极为巧妙的引理。

\begin{lemma}
\label{smoothing-lemma-desingularize}
设 $R$ 为诺特环，$\Lambda$ 为 $R$-代数。取 $\pi \in R$，
并假设 $\text{Ann}_\Lambda(\pi) = \text{Ann}_\Lambda(\pi^2)$。
设 $A \to \Lambda$ 为 $R$-代数同态，其中 $A$ 为有限表示。假设
\begin{enumerate}
\item $\pi$ 的像在 $A$ 中相对于 $R$ 是严格标准元；
\item 存在与到 $\Lambda/\pi^4 \Lambda$ 的同态相容的截面
$\rho : A/\pi^4 A \to R/\pi^4 R$。
\end{enumerate}
则可以找到 $R$-代数同态 $A \to B \to \Lambda$，其中 $B$ 为有限表示，
并且 $\mathfrak a B \subset H_{B/R}$，这里
$\mathfrak a = \text{Ann}_R(\text{Ann}_R(\pi^2)/\text{Ann}_R(\pi))$。
\end{lemma}

\begin{proof}
选取一个表示
$$
A = R[x_1, \ldots, x_n]/(f_1, \ldots, f_m)
$$
以及 $0 \leq c \leq \min(n, m)$，使得
(\ref{smoothing-equation-strictly-standard-one}) 对 $\pi$ 成立，并且
\begin{equation}
\label{smoothing-equation-star}
\pi f_{c + j} \in (f_1, \ldots, f_c) + (f_1, \ldots, f_m)^2
\end{equation}
对 $j = 1, \ldots, m - c$ 成立。设 $\rho$ 把 $x_i$ 映到
$r_i \in R$ 的类。把 $x_i$ 替换为 $x_i-r_i$ 后，可以假设
$\rho(x_i) = 0$ 于 $R/\pi^4 R$ 中。这蕴含
$f_j(0) \in \pi^4R$，并且 $A \to \Lambda$ 把 $x_i$
映到某个 $\pi^4\lambda_i$，其中 $\lambda_i \in \Lambda$。写成
$$
f_j = f_j(0) + \sum\nolimits_{i = 1, \ldots, n} r_{ji} x_i + \text{高阶项}
$$
于是 $\partial f_j/\partial x_i$ 的常数项为 $r_{ji}$。
把 $\rho$ 应用于关于 $\pi$ 的 (\ref{smoothing-equation-strictly-standard-one})，可得
$$
\pi = \sum\nolimits_{I \subset \{1, \ldots, n\},\ |I| = c}
r_I \det(r_{ji})_{j = 1, \ldots, c,\ i \in I} \bmod \pi^4R
$$
其中某些 $r_I \in R$。因此有
$$
u\pi = \sum\nolimits_{I \subset \{1, \ldots, n\},\ |I| = c}
r_I \det(r_{ji})_{j = 1, \ldots, c,\ i \in I}
$$
其中某个 $u \in 1 + \pi^3R$。由《代数》
引理 \ref{algebra-lemma-matrix-left-inverse}，存在一个
$n \times c$ 矩阵 $(s_{ik})$，使得
$$
u\pi \delta_{jk} = \sum\nolimits_{i = 1, \ldots, n} r_{ji}s_{ik}\quad
\text{对所有 } j, k = 1, \ldots, c
$$
（Kronecker 符号）。引入辅助变量
$v_1, \ldots, v_c, w_1, \ldots, w_n$，并置
$$
h_i = x_i - \pi^2 \sum\nolimits_{j = 1, \ldots c} s_{ij} v_j - \pi^3 w_i
$$
下文将不再另行说明地使用等式
$$
R[x_1, \ldots, x_n, v_1, \ldots, v_c, w_1, \ldots, w_n]/
(h_1, \ldots, h_n) = R[v_1, \ldots, v_c, w_1, \ldots, w_n]
$$
。在
$R[x_1, \ldots, x_n, v_1, \ldots, v_c, w_1, \ldots, w_n]/
(h_1, \ldots, h_n)$ 中有
\begin{align*}
f_j & = f_j(x_1 - h_1, \ldots, x_n - h_n) \\
& =
\pi^2 \sum\nolimits_{k = 1}^c
\left(\sum\nolimits_{i = 1}^n r_{ji} s_{ik}\right) v_k
+
\pi^3 \sum\nolimits_{i = 1}^n r_{ji}w_i \bmod \pi^4 \\
& =
\pi^3 v_j + \pi^3 \sum\nolimits_{i = 1}^n r_{ji}w_i \bmod \pi^4
\end{align*}
，其中 $1 \leq j \leq c$。因此，可以选取元素
$g_j \in R[v_1, \ldots, v_c, w_1, \ldots, w_n]$，使得
$g_j = v_j + \sum r_{ji}w_i \bmod \pi$，并且在 $R$-代数
$R[x_1, \ldots, x_n, v_1, \ldots, v_c, w_1, \ldots, w_n]/
(h_1, \ldots, h_n)$ 中有 $f_j = \pi^3 g_j$。置
$$
B = R[x_1, \ldots, x_n, v_1, \ldots, v_c, w_1, \ldots, w_n]/
(f_1, \ldots, f_m, h_1, \ldots, h_n, g_1, \ldots, g_c).
$$
$A \to B$ 的同态是显然的。定义 $B \to \Lambda$，令
$x_i \mapsto \pi^4\lambda_i$、$v_i \mapsto 0$ 且
$w_i \mapsto \pi \lambda_i$。于是，元素 $f_j$ 和 $h_i$ 显然都映到
$\Lambda$ 中的零。此外，$g_i$ 映到 $\pi\Lambda$ 中的某个元素 $t$，
且 $\pi^3t = 0$（因为模掉由诸 $h$ 生成的理想后，
$f_i = \pi^3 g_i$）。所以，假设
$\text{Ann}_\Lambda(\pi) = \text{Ann}_\Lambda(\pi^2)$ 蕴含 $t=0$。
因而，只需证明关于光滑性的断言即可。

\medskip\noindent
注意，$B_\pi \cong A_\pi[v_1, \ldots, v_c]$，因为方程
$f_i=0$ 蕴含 $g_i=0$。又因假设 $\pi$ 在 $A$ 中相对于 $R$ 是严格标准元，
$A_\pi$ 在 $R$ 上光滑，见引理 \ref{smoothing-lemma-elkik}，故
$B_\pi$ 在 $R$ 上光滑。

\medskip\noindent
置 $B' = R[v_1, \ldots, v_c, w_1, \ldots, w_n]/(g_1, \ldots, g_c)$。
由于 $g_i = v_i + \sum r_{ji}w_i \bmod \pi$，可知
$B'/\pi B' = R/\pi R[w_1, \ldots, w_n]$。因此，由《代数》引理
\ref{algebra-lemma-localize-relative-complete-intersection} 和
\ref{algebra-lemma-flat-fibre-smooth}
，$R \to B'$ 在 $V(\pi)$ 的每一点都光滑且相对维数为 $n$
（第一个引理表明它在这些素理想处为 syntomic，因而特别是平坦的；
随后第二个引理表明它是光滑的）。

\medskip\noindent
设 $\mathfrak q \subset B$ 是一个满足 $\pi \in \mathfrak q$ 的素理想，
并且对某个 $r \in \mathfrak a$ 有 $r \not \in \mathfrak q$。
记 $\mathfrak q' = B' \cap \mathfrak q$。我们断言，满射 $B' \to B$
诱导局部环的同构 $(B')_{\mathfrak q'} \to B_\mathfrak q$；这将完成
引理的证明。注意，$B_\mathfrak q$ 是 $(B')_{\mathfrak q'}$ 对由
$f_{c + j}$（$j = 1, \ldots, m - c$）生成之理想的商。我们作两个观察：
第一，$f_{c+j}$ 在 $(B')_{\mathfrak q'}$ 中的像可被 $\pi^2$ 整除；
第二，由 (\ref{smoothing-equation-star})，$\pi f_{c+j}$ 在
$(B')_{\mathfrak q'}$ 中的像可以写成
$\sum b_{j_1 j_2} f_{c + j_1}f_{c + j_2}$。因此，每个
$\pi f_{c+j}$ 的像都属于由元素 $\pi^2 f_{c+j'}$ 生成的理想。
由于 $(B')_{\mathfrak q'}$ 是诺特局部环，故其中
$\pi f_{c+j}=0$，见《代数》
引理 \ref{algebra-lemma-intersect-powers-ideal-module-zero}。
又因 $R \to (B')_{\mathfrak q'}$ 平坦，可知
$$
\left(\text{Ann}_R(\pi^2)/\text{Ann}_R(\pi)\right)
\otimes_R (B')_{\mathfrak q'}
=
\text{Ann}_{(B')_{\mathfrak q'}}(\pi^2)/\text{Ann}_{(B')_{\mathfrak q'}}(\pi)
$$
由于 $r \in \mathfrak a$ 在 $(B')_{\mathfrak q'}$ 中可逆，
这个模必为零。因此，$f_{c+j}$ 在 $(B')_{\mathfrak q'}$ 中的像为零，
正合所需。
\end{proof}

\begin{lemma}
\label{smoothing-lemma-desingularize-strictly-standard}
设 $R$ 为诺特环，$\Lambda$ 为 $R$-代数。取 $\pi \in R$，并假设
$\text{Ann}_R(\pi) = \text{Ann}_R(\pi^2)$ 且
$\text{Ann}_\Lambda(\pi) = \text{Ann}_\Lambda(\pi^2)$。
设 $A \to \Lambda$ 和 $D \to \Lambda$ 为 $R$-代数同态，
其中 $A$ 和 $D$ 均为有限表示。假设
\begin{enumerate}
\item $\pi$ 在 $A$ 中相对于 $R$ 是严格标准元；
\item 存在与到 $\Lambda/\pi^4\Lambda$ 的同态相容的 $R$-代数同态
$A/\pi^4 A \to D/\pi^4 D$。
\end{enumerate}
则可找到 $R$-代数同态 $B \to \Lambda$，其中 $B$ 为有限表示，
以及与到 $\Lambda$ 的同态相容的 $R$-代数同态 $A \to B$ 和 $D \to B$，
使得 $H_{D/R}B \subset H_{B/D}$ 且 $H_{D/R}B \subset H_{B/R}$。
\end{lemma}

\begin{proof}
把引理 \ref{smoothing-lemma-desingularize} 应用于
$$
D \longrightarrow A \otimes_R D \longrightarrow \Lambda
$$
以及 $\pi$ 在 $D$ 中的像。由引理
\ref{smoothing-lemma-strictly-standard-base-change}，$\pi$ 在 $A \otimes_R D$
中相对于 $D$ 是严格标准元。取 (2) 中同态所诱导的同态作为截面
$\rho : (A \otimes_R D)/\pi^4 (A \otimes_R D) \to D/\pi^4 D$。
于是引理 \ref{smoothing-lemma-desingularize} 适用，并给出分解
$A \otimes_R D \to B \to \Lambda$，其中 $B$ 为有限表示，且
$\mathfrak a B \subset H_{B/D}$，这里
$$
\mathfrak a = \text{Ann}_D(\text{Ann}_D(\pi^2)/\text{Ann}_D(\pi)).
$$
对 $D$ 的任意素理想 $\mathfrak q$，若 $D_\mathfrak q$ 在 $R$ 上平坦，
则有
$\text{Ann}_{D_\mathfrak q}(\pi^2)/\text{Ann}_{D_\mathfrak q}(\pi) = 0$
这是因为元素的零化子与平坦基变换可交换，且我们假设了
$\text{Ann}_R(\pi) = \text{Ann}_R(\pi^2)$。由于 $D$ 是诺特环，
$\text{Ann}_D(\pi^2)/\text{Ann}_D(\pi)$ 是有限 $D$-模，
故取其零化子与局部化可交换。因此 $\mathfrak a \not \subset \mathfrak q$，
从而 $D \to B$ 在 $B$ 中位于 $\mathfrak q$ 之上的任意素理想处光滑。
凡使 $R \to D$ 光滑的 $D$ 中素理想 $\mathfrak q$，都有
$D_\mathfrak q$ 在 $R$ 上平坦，所以 $H_{D/R}B \subset H_{B/D}$。
最后的包含关系 $H_{D/R}B \subset H_{B/R}$ 源于光滑环同态的复合仍光滑
（《代数》，引理 \ref{algebra-lemma-compose-smooth}）。
\end{proof}

\begin{lemma}
\label{smoothing-lemma-desingularize-lifting-apply}
设 $R$ 为诺特环，$\Lambda$ 为 $R$-代数。取 $\pi \in R$，并假设
$\text{Ann}_R(\pi) = \text{Ann}_R(\pi^2)$ 且
$\text{Ann}_\Lambda(\pi) = \text{Ann}_\Lambda(\pi^2)$。
设 $A \to \Lambda$ 为 $R$-代数同态，其中 $A$ 为有限表示，
并假设 $\pi$ 在 $A$ 中相对于 $R$ 是严格标准元。设
$$
A/\pi^8A \to \bar C \to \Lambda/\pi^8\Lambda
$$
是一个分解，其中 $\bar C$ 为有限表示。则可以找到分解
$A \to B \to \Lambda$，其中 $B$ 为有限表示，
$R_\pi \to B_\pi$ 光滑，并且
$$
H_{\bar C/(R/\pi^8 R)} \cdot \Lambda/\pi^8\Lambda
\subset
\sqrt{H_{B/R} \Lambda} \bmod \pi^8\Lambda.
$$
\end{lemma}

\begin{proof}
应用引理 \ref{smoothing-lemma-lifting}，得到 $R \to D \to \Lambda$ 以及分解
$\bar C/\pi^4\bar C \to D/\pi^4 D \to \Lambda/\pi^4\Lambda$
，使 $R \to D$ 在任何不含 $\pi$ 的素理想处都光滑，
并且在任何位于 $\bar C/\pi^4\bar C$ 中使
$R/\pi^8 R \to \bar C$ 光滑之素理想上方的素理想处也光滑。
由引理 \ref{smoothing-lemma-desingularize-strictly-standard}，可以找到有限表示的
$R$-代数 $B$，以及分解 $A \to B \to \Lambda$ 和
$D \to B \to \Lambda$，使得 $H_{D/R}B \subset H_{B/R}$。
略去验证这确实给出了本引理所述问题的解。
\end{proof}












\section{预备：约化到基域}
\label{smoothing-section-reduction}

\noindent
本节应用前几节的引理，证明只需在基环为域时证明主结果，
见引理 \ref{smoothing-lemma-reduce-to-field}。

\begin{situation}
\label{smoothing-situation-global}
这里 $R \to \Lambda$ 是诺特环之间的正则环同态。
\end{situation}

\noindent
设 $R \to \Lambda$ 如情形 \ref{smoothing-situation-global} 所述。
若 $\Lambda$ 是光滑 $R$-代数的滤过余极限，则称
{\it $R \to \Lambda$ 满足 PT}。

\begin{lemma}
\label{smoothing-lemma-product}
设 $R_i \to \Lambda_i$（$i=1,2$）如情形
\ref{smoothing-situation-global} 所述。若 $R_i \to \Lambda_i$（$i=1,2$）
均满足 PT，则 $R_1 \times R_2 \to \Lambda_1 \times \Lambda_2$ 满足 PT。
\end{lemma}

\begin{proof}
略。提示：滤过余极限的乘积仍是滤过余极限。
\end{proof}

\begin{lemma}
\label{smoothing-lemma-delocalize-base}
设 $R \to A \to \Lambda$ 为环同态，其中 $A$ 在 $R$ 上为有限表示；
设 $S \subset R$ 为乘法集。给定分解
$S^{-1}A \to B' \to S^{-1}\Lambda$，其中 $B'$ 在 $S^{-1}R$ 上光滑。
则可找到分解 $A \to B \to \Lambda$，使某个 $s \in S$ 的像
在 $B$ 中相对于 $R$ 是初等标准元
（定义 \ref{smoothing-definition-strictly-standard}）。
\end{lemma}

\begin{proof}
先对 $S^{-1}R \to B'$ 应用引理
\ref{smoothing-lemma-smooth-standard-smooth}，从而可假设 $B'$ 在 $S^{-1}R$
上标准光滑。写
$A = R[x_1, \ldots, x_n]/(g_1, \ldots, g_t)$，并设 $x_i$
在 $\Lambda$ 中映到 $\lambda_i$。可以写成
$B' = S^{-1}R[x_1, \ldots, x_{n + m}]/(f_1, \ldots, f_c)$
，其中某个 $c \geq n$，使得
$\det(\partial f_j/\partial x_i)_{i, j = 1, \ldots, c}$
在 $B'$ 中可逆，并且 $A \to B'$ 由 $x_i \mapsto x_i$ 给出，
见引理 \ref{smoothing-lemma-standard-smooth-include-generators}。
将 $i>n$ 的 $x_i$ 乘以 $S$ 中某个元素，并相应修改方程 $f_j$ 后，
可假设对 $i>n$，$B' \to S^{-1}\Lambda$ 把 $x_i$ 映到某个
$\lambda_i/1$，其中 $\lambda_i \in \Lambda$。选取关系式
$$
1 =
a_0 \det(\partial f_j/\partial x_i)_{i, j = 1, \ldots, c}
+
\sum\nolimits_{j = 1, \ldots, c} a_jf_j
$$
，其中某些 $a_j \in S^{-1}R[x_1, \ldots, x_{n + m}]$。
由于 $S$ 的每个元素在 $B'$ 中都可逆，消去分母后，可以假设
$f_j, a_j \in R[x_1, \ldots, x_{n + m}]$，并且
$$
s_0 = a_0 \det(\partial f_j/\partial x_i)_{i, j = 1, \ldots, c}
+
\sum\nolimits_{j = 1, \ldots, c} a_jf_j
$$
，其中某个 $s_0 \in S$。由于 $g_j$ 在
$S^{-1}R[x_1, \ldots, x_{n + m}]/(f_1, \ldots, x_c)$
中的像为零，可以找到元素 $s_j \in S$，使得 $s_jg_j=0$ 于
$R[x_1, \ldots, x_{n + m}]/(f_1, \ldots, f_c)$.
由于 $f_j$ 在 $S^{-1}\Lambda$ 中的像为零，可以找到 $s'_j \in S$，
使 $s'_j f_j(\lambda_1, \ldots, \lambda_{n + m}) = 0$ 于
$\Lambda$ 中。考虑环
$$
B = R[x_1, \ldots, x_{n + m}]/
(s'_1f_1, \ldots, s'_cf_c, g_1, \ldots, g_t)
$$
以及分解 $A \to B \to \Lambda$，其中 $B \to \Lambda$ 由
$x_i \mapsto \lambda_i$ 给出。我们断言
$s = s_0s_1 \ldots s_ts'_1 \ldots s'_c$ 在 $B$ 中相对于 $R$
是初等标准元，这将完成证明。事实上，
$s_j g_j \in (f_1, \ldots, f_c)$，因而
$sg_j \in (s'_1f_1, \ldots, s'_cf_c)$。最后有
$$
a_0\det(\partial s'_jf_j/\partial x_i)_{i, j = 1, \ldots, c}
+
\sum\nolimits_{j = 1, \ldots, c}
(s'_1 \ldots \hat{s'_j} \ldots s'_c) a_j s'_jf_j
=
s_0s'_1\ldots s'_c
$$
，而它整除 $s$，正合所需。
\end{proof}

\begin{lemma}
\label{smoothing-lemma-reduce-to-field}
\begin{slogan}
证明 Popescu 逼近可约化到域上代数的情形
\end{slogan}
若在情形 \ref{smoothing-situation-global} 中每当 $R$ 是域时 PT 都成立，
则 PT 一般地成立。
\end{lemma}

\begin{proof}
假设在情形 \ref{smoothing-situation-global} 中，只要 $R$ 是域，PT 就成立。
任取情形 \ref{smoothing-situation-global} 中的 $R \to \Lambda$。注意，
$R/I \to \Lambda/I\Lambda$ 仍是诺特环之间的正则环同态，
见《代数详论》引理 \ref{more-algebra-lemma-regular-base-change}。
考虑理想的集合
$$
\mathcal{I} = \{I \subset R \mid R/I \to \Lambda/I\Lambda
\text{不满足 PT}\}
$$
我们必须证明 $\mathcal{I}$ 为空。若它非空，则因 $R$ 是诺特环，
其中存在极大元。把 $R$ 替换为 $R/I$，把 $\Lambda$ 替换为
$\Lambda/I$，便得到这样一种情形：对 $R$ 的任意非零理想 $I$，
$R/I \to \Lambda/I\Lambda$ 都满足 PT。特别地，应用命题
\ref{smoothing-proposition-lift} 可知 $R$ 是既约环。

\medskip\noindent
设 $A \to \Lambda$ 为 $R$-代数同态，其中 $A$ 为有限表示。
必须找到分解 $A \to B \to \Lambda$，使 $B$ 在 $R$ 上光滑，
见《代数》引理 \ref{algebra-lemma-when-colimit}。

\medskip\noindent
设 $S \subset R$ 为非零因子的集合，并考虑 $R$ 的全分式环
$Q=S^{-1}R$。已知 $Q = K_1 \times \ldots \times K_n$ 是若干域的乘积，
见《代数》引理 \ref{algebra-lemma-total-ring-fractions-no-embedded-points}
和 \ref{algebra-lemma-Noetherian-irreducible-components}。
由引理 \ref{smoothing-lemma-product} 及我们的假设，环同态
$S^{-1}R \to S^{-1}\Lambda$ 满足 PT。因此，可以找到分解
$S^{-1}A \to B' \to S^{-1}\Lambda$，其中 $B'$ 在 $S^{-1}R$ 上光滑。

\medskip\noindent
应用引理 \ref{smoothing-lemma-delocalize-base}，得到分解
$A \to B \to \Lambda$，使某个 $\pi \in S$ 在 $B$ 中相对于 $R$
是初等标准元。把 $A$ 替换为 $B$ 后，可假设 $\pi$ 在 $A$ 中是
初等标准元，因而是严格标准元。已知
$R/\pi^8R \to \Lambda/\pi^8\Lambda$ 满足 PT，所以可以找到分解
$R/\pi^8 R \to A/\pi^8A \to \bar C \to \Lambda/\pi^8\Lambda$
，其中 $R/\pi^8 R \to \bar C$ 光滑。由引理 \ref{smoothing-lemma-lifting}，
可以找到 $R$-代数同态 $D \to \Lambda$，其中 $D$ 在 $R$ 上光滑，
以及分解
$R/\pi^4 R \to A/\pi^4A \to D/\pi^4D \to \Lambda/\pi^4\Lambda$.
由引理 \ref{smoothing-lemma-desingularize-strictly-standard}，可以找到
$A \to B \to \Lambda$，其中 $B$ 在 $R$ 上光滑，从而完成证明。
\end{proof}










\section{局部技巧}
\label{smoothing-section-local}


\begin{situation}
\label{smoothing-situation-local}
给定诺特环 $R$、$R$-代数同态 $A \to \Lambda$，以及素理想
$\mathfrak q \subset \Lambda$。假设 $A$ 在 $R$ 上为有限表示。
在此情形下，记 $\mathfrak h_A = \sqrt{H_{A/R} \Lambda}$。
\end{situation}

\noindent
设 $R \to A \to \Lambda \supset \mathfrak q$ 如情形
\ref{smoothing-situation-local} 所述。若存在分解 $A \to B \to \Lambda$，
其中 $B$ 为有限表示，且
$\mathfrak h_A \subset \mathfrak h_B \not \subset \mathfrak q$，
则称 {\it $R \to A \to \Lambda \supset \mathfrak q$ 可消解}。
此时称分解 $A \to B \to \Lambda$ 为
{\it $R \to A \to \Lambda \supset \mathfrak q$ 的一个消解}。

\begin{lemma}
\label{smoothing-lemma-lift-solution}
设 $R \to A \to \Lambda \supset \mathfrak q$ 如情形
\ref{smoothing-situation-local} 所述。设 $r \geq 1$，并设
$\pi_1, \ldots, \pi_r \in R$ 的像属于 $\mathfrak q$。假设
\begin{enumerate}
\item 对 $i = 1, \ldots, r$ 有
$$
\text{Ann}_{R/(\pi_1^8, \ldots, \pi_{i - 1}^8)R}(\pi_i)
=
\text{Ann}_{R/(\pi_1^8, \ldots, \pi_{i - 1}^8)R}(\pi_i^2)
$$
且
$$
\text{Ann}_{\Lambda/(\pi_1^8, \ldots, \pi_{i - 1}^8)\Lambda}(\pi_i)
=
\text{Ann}_{\Lambda/(\pi_1^8, \ldots, \pi_{i - 1}^8)\Lambda}(\pi_i^2)
$$
\item 对 $i = 1, \ldots, r$，元素 $\pi_i$ 的像在 $A$ 中相对于 $R$
是严格标准元。
\end{enumerate}
那么，若
$$
R/(\pi_1^8, \ldots, \pi_r^8)R \to A/(\pi_1^8, \ldots, \pi_r^8)A
\to \Lambda/(\pi_1^8, \ldots, \pi_r^8)\Lambda \supset
\mathfrak q/(\pi_1^8, \ldots, \pi_r^8)\Lambda
$$
可消解，则 $R \to A \to \Lambda \supset \mathfrak q$ 也可消解。
\end{lemma}

\begin{proof}
对 $r$ 作归纳证明。

\medskip\noindent
先考虑 $r=1$。此时假设存在分解
$A/\pi_1^8 \to \bar C \to \Lambda/\pi_1^8$，它消解了模
$\pi_1^8$ 的情形。条件 (1) 和 (2) 正是应用引理
\ref{smoothing-lemma-desingularize-lifting-apply} 所需的假设。
因此，可以把消解 $\bar C$“提升”为
$R \to A \to \Lambda \supset \mathfrak q$ 的消解。

\medskip\noindent
再考虑 $r>1$。此时，把关于 $r-1$ 的归纳假设应用于情形
$R/\pi_1^8 \to A/\pi_1^8 \to \Lambda/\pi_1^8
\supset \mathfrak q/\pi_1^8\Lambda$.
注意，引理 \ref{smoothing-lemma-strictly-standard-base-change} 表明性质 (2)
在这里得到保持。
\end{proof}

\begin{lemma}
\label{smoothing-lemma-delocalize-weak}
\begin{reference}
\cite[Lemma 12.2]{swan} 或
\cite[Lemma 2]{popescu-GND}
\end{reference}
设 $R \to A \to \Lambda \supset \mathfrak q$ 如情形
\ref{smoothing-situation-local} 所述，并令 $\mathfrak p=R\cap\mathfrak q$。
假设 $\mathfrak q$ 是包含 $\mathfrak h_A$ 的极小素理想，并且
$R_\mathfrak p \to A_\mathfrak p \to \Lambda_\mathfrak q
\supset \mathfrak q\Lambda_\mathfrak q$ 可消解。
则存在分解 $A \to C \to \Lambda$，其中 $C$ 为有限表示，且
$H_{C/R} \Lambda \not \subset \mathfrak q$。
\end{lemma}

\begin{proof}
设 $A_\mathfrak p \to C \to \Lambda_\mathfrak q$ 是
$R_\mathfrak p \to A_\mathfrak p \to \Lambda_\mathfrak q
\supset \mathfrak q\Lambda_\mathfrak q$ 的一个消解。由
$\mathfrak q$ 是包含 $\mathfrak h_A$ 的极小素理想这一假设，
这意味着 $H_{C/R_\mathfrak p} \Lambda_\mathfrak q = \Lambda_\mathfrak q$。
由引理 \ref{smoothing-lemma-final-solve}，可假设 $C$ 在 $R_\mathfrak p$ 上光滑；
再由引理 \ref{smoothing-lemma-smooth-standard-smooth}，可假设 $C$ 在
$R_\mathfrak p$ 上标准光滑。写
$A = R[x_1, \ldots, x_n]/(g_1, \ldots, g_t)$，并设
$A \to \Lambda$ 由 $x_i \mapsto \lambda_i$ 给出。写
$C = R_\mathfrak p[x_1, \ldots, x_{n + m}]/(f_1, \ldots, f_c)$，
其中某个 $c \geq n$，使 $A \to C$ 把 $x_i$ 映到 $x_i$，并且
$\det(\partial f_j/\partial x_i)_{i, j = 1, \ldots, c}$
在 $C$ 中可逆，见引理 \ref{smoothing-lemma-standard-smooth-include-generators}。
消去分母后，可假设 $f_1,\ldots,f_c$ 属于
$R[x_1, \ldots, x_{n + m}]$。当然，
$\det(\partial f_j/\partial x_i)_{i, j = 1, \ldots, c}$
在 $R[x_1, \ldots, x_{n + m}]/(f_1, \ldots, f_c)$ 中未必可逆，
但将某个 $s_0 \in R$（$s_0 \not \in \mathfrak p$）取逆后便可逆。
由于 $g_j$ 经 $R[x_1, \ldots, x_n] \to A \to C$ 映到零，
可以找到 $s_j \in R$，其中 $s_j \not \in \mathfrak p$，使
$s_jg_j$ 在 $R[x_1, \ldots, x_{n + m}]/(f_1, \ldots, f_c)$ 中为零。
将 $f_j$ 写成 $f_j = F_j(x_1, \ldots, x_{n + m}, 1)$，其中多项式
$F_j \in R[x_1, \ldots, x_n, X_{n + 1}, \ldots, X_{n + m + 1}]$
关于 $X_{n + 1}, \ldots, X_{n + m + 1}$ 齐次。选取
$\lambda_{n + i} \in \Lambda$（$i = 1, \ldots, m + 1$），其中
$\lambda_{n + m + 1} \not \in \mathfrak q$，使 $x_{n+i}$ 在
$\Lambda_\mathfrak q$ 中映到
$\lambda_{n + i}/\lambda_{n + m + 1}$。于是
\begin{align*}
F_j(\lambda_1, \ldots, \lambda_{n + m + 1})
& =
(\lambda_{n + m + 1})^{\deg(F_j)} F_j(\lambda_1, \ldots, \lambda_n,
\frac{\lambda_{n + 1}}{\lambda_{n + m + 1}}, \ldots,
\frac{\lambda_{n + m}}{\lambda_{n + m + 1}}, 1) \\
& =
(\lambda_{n + m + 1})^{\deg(F_j)} f_j(\lambda_1, \ldots, \lambda_n,
\frac{\lambda_{n + 1}}{\lambda_{n + m + 1}}, \ldots,
\frac{\lambda_{n + m}}{\lambda_{n + m + 1}}) \\
& = 0
\end{align*}
在 $\Lambda_\mathfrak q$ 中成立。因此，可以找到
$\lambda_0 \in \Lambda$，其中 $\lambda_0 \not \in \mathfrak q$，使得
$\lambda_0 F_j(\lambda_1, \ldots, \lambda_{n + m + 1}) = 0$
在 $\Lambda$ 中成立。现在置 $B$ 为
$$
R[x_0, \ldots, x_{n + m + 1}]/
(g_1, \ldots, g_t, x_0F_1(x_1, \ldots, x_{n + m + 1}), \ldots,
x_0F_c(x_1, \ldots, x_{n + m + 1}))
$$
，并通过 $x_i \mapsto \lambda_i$ 将其映到 $\Lambda$。
令 $b$ 为 $x_0x_{n+m+1}s_0s_1\ldots s_t$ 在 $B$ 中的像。
则 $B_b$ 同构于
$$
R_{s_0s_1 \ldots s_t}[x_0, x_1, \ldots, x_{n + m + 1}, 1/x_0x_{n + m + 1}]/
(f_1, \ldots, f_c)
$$
，后者依构造在 $R$ 上光滑。由于 $b$ 的像不属于 $\mathfrak q$，
结论得证。
\end{proof}

\begin{lemma}
\label{smoothing-lemma-delocalize-height-zero}
设 $R \to A \to \Lambda \supset \mathfrak q$ 如情形
\ref{smoothing-situation-local} 所述，并令 $\mathfrak p=R\cap\mathfrak q$。假设
\begin{enumerate}
\item $\mathfrak q$ 是包含 $\mathfrak h_A$ 的极小素理想；
\item $R_\mathfrak p \to A_\mathfrak p \to \Lambda_\mathfrak q
\supset \mathfrak q\Lambda_\mathfrak q$ 可消解；
\item $\dim(\Lambda_\mathfrak q) = 0$.
\end{enumerate}
则 $R \to A \to \Lambda \supset \mathfrak q$ 可消解。
\end{lemma}

\begin{proof}
由 (3)，环 $\Lambda_\mathfrak q$ 是 Artin 局部环，因而
$\mathfrak q\Lambda_\mathfrak q$ 幂零。故对某个 $N>0$，有
$(\mathfrak h_A)^N \Lambda_\mathfrak q = 0$。于是存在
$\lambda \in \Lambda$，其中 $\lambda \not \in \mathfrak q$，使
$\lambda(\mathfrak h_A)^N=0$ 于 $\Lambda$ 中。设
$H_{A/R}=(a_1,\ldots,a_r)$，于是 $\lambda a_i^N=0$ 于 $\Lambda$ 中。
由引理 \ref{smoothing-lemma-delocalize-weak}，可以找到分解
$A \to C \to \Lambda$，其中 $C$ 为有限表示，且
$\mathfrak h_C \not \subset \mathfrak q$。写
$C=A[x_1,\ldots,x_n]/(f_1,\ldots,f_m)$。置
$$
B = A[x_1, \ldots, x_n, y_1, \ldots, y_r, z, t_{ij}]/
(f_j - \sum y_i t_{ij}, zy_i)
$$
其中 $t_{ij}$ 是 $rm$ 个变量。令 $t_{ij}$ 全部等于零，便得到同态
$B \to C[y_i,z]/(y_iz)$。同态 $B \to \Lambda$ 是复合
$B \to C[y_i,z]/(y_iz) \to \Lambda$，其中
$C[y_i,z]/(y_iz) \to \Lambda$ 延拓给定同态 $C \to \Lambda$，
把 $z$ 映到 $\lambda$，并把 $y_i$ 映到 $a_i^N$ 在 $\Lambda$ 中的像。

\medskip\noindent
我们断言，$B$ 是 $R \to A \to \Lambda \supset \mathfrak q$ 的一个解。
首先，$B_z$ 同构于 $C[y_1, \ldots, y_r,z,z^{-1}]$，因而是光滑的。
另一方面，
$B_{y_\ell} \cong A[x_i, y_i, y_\ell^{-1}, t_{ij}, i \not = \ell]$
在 $A$ 上光滑。因此，$z$ 以及各个 $a_\ell y_\ell$
（光滑同态的复合仍光滑）全都属于 $H_{B/R}$。引理得证。
\end{proof}




\section{可分剩余域}
\label{smoothing-section-separable}

\noindent
本节说明当剩余域扩张可分时如何求解一个局部问题。

\begin{lemma}[Ogoma]
\label{smoothing-lemma-ogoma}
设 $A$ 为诺特环，$M$ 为有限 $A$-模，$S\subset A$ 为乘法集。
若 $\pi\in A$ 且
$\Ker(\pi : S^{-1}M \to S^{-1}M) =
\Ker(\pi^2 : S^{-1}M \to S^{-1}M)$
，则存在 $s\in S$，使得对任意 $n>0$ 都有
$\Ker(s^n\pi : M \to M) = \Ker((s^n\pi)^2 : M \to M)$.
\end{lemma}

\begin{proof}
令 $K=\Ker(\pi:M\to M)$、
$K'=\{m\in M\mid \pi^2m=0\text{ 于 }S^{-1}M\text{ 中}\}$，
并令 $Q=K'/K$。由假设，$S^{-1}Q=0$。由于 $A$ 是诺特环，
$Q$ 是有限 $A$-模。因此可以找到 $s\in S$ 使 $s$ 零化 $Q$；
这个 $s$ 即合要求。
\end{proof}

\begin{lemma}
\label{smoothing-lemma-find-sequence}
设 $\Lambda$ 为诺特环，$I\subset\Lambda$ 为理想，
且 $I\subset\mathfrak q$，其中 $\mathfrak q$ 为素理想。
设 $n,e$ 为正整数。假设
$\mathfrak q^n\Lambda_\mathfrak q\subset I\Lambda_\mathfrak q$，
且 $\Lambda_\mathfrak q$ 是维数为 $d$ 的正则局部环。
则存在 $n>0$ 以及 $\pi_1,\ldots,\pi_d\in\Lambda$，使得
\begin{enumerate}
\item $(\pi_1, \ldots, \pi_d)\Lambda_\mathfrak q =
\mathfrak q\Lambda_\mathfrak q$；
\item $\pi_1^n, \ldots, \pi_d^n \in I$；
\item 对 $i=1,\ldots,d$ 有
$$
\text{Ann}_{\Lambda/(\pi_1^e, \ldots, \pi_{i - 1}^e)\Lambda}(\pi_i) =
\text{Ann}_{\Lambda/(\pi_1^e, \ldots, \pi_{i - 1}^e)\Lambda}(\pi_i^2).
$$
\end{enumerate}
\end{lemma}

\begin{proof}
置 $S=\Lambda\setminus\mathfrak q$，于是
$\Lambda_\mathfrak q=S^{-1}\Lambda$。先选取满足 (1) 的
$\pi_1,\ldots,\pi_d$；这因 $\Lambda_\mathfrak q$ 正则而可行。
由假设，$\pi_i^n\in I\Lambda_\mathfrak q$，故可找到
$s_1,\ldots,s_d\in S$，使 $s_i\pi_i^n\in I$。以
$s_i\pi_i$ 替换 $\pi_i$ 后便得到 (2)。注意，继续乘以 $S$ 中元素
仍保持 (1) 和 (2)。假设对某个 $t\in\{0,\ldots,d\}$，
(3) 对 $i=1,\ldots,t$ 成立。注意，
$\pi_1,\ldots,\pi_d$ 是 $S^{-1}\Lambda$ 中的正则序列，
见《代数》引理 \ref{algebra-lemma-regular-ring-CM}。
特别地，由《代数》引理 \ref{algebra-lemma-regular-sequence-powers}，
$\pi_1^e,\ldots,\pi_t^e,\pi_{t+1}$ 是
$S^{-1}\Lambda=\Lambda_\mathfrak q$ 中的正则序列。因此
$$
\text{Ann}_{S^{-1}\Lambda/(\pi_1^e, \ldots, \pi_{i - 1}^e)}(\pi_i) =
\text{Ann}_{S^{-1}\Lambda/(\pi_1^e, \ldots, \pi_{i - 1}^e)}(\pi_i^2).
$$
由引理 \ref{smoothing-lemma-ogoma}，取某个 $s\in S$，以
$s\pi_{t+1}$ 替换 $\pi_{t+1}$ 后，便可使 (3) 对 $i=t+1$ 成立。
对 $t$ 归纳，即得到满足 (1)、(2) 和 (3) 的序列。
\end{proof}

\begin{lemma}
\label{smoothing-lemma-resolve-special}
设 $k \to A \to \Lambda \supset \mathfrak q$ 如情形
\ref{smoothing-situation-local} 所述，其中
\begin{enumerate}
\item $k$ 是域；
\item $\Lambda$ 是诺特环；
\item $\mathfrak q$ 是包含 $\mathfrak h_A$ 的极小素理想；
\item $\Lambda_\mathfrak q$ 是正则局部环；
\item 域扩张 $\kappa(\mathfrak q)/k$ 可分。
\end{enumerate}
则 $k \to A \to \Lambda \supset \mathfrak q$ 可消解。
\end{lemma}

\begin{proof}
置 $d=\dim\Lambda_\mathfrak q$ 及 $R=k[x_1,\ldots,x_d]$。
选取 $n>0$，使
$\mathfrak q^n\Lambda_\mathfrak q\subset\mathfrak h_A\Lambda_\mathfrak q$；
这因 $\mathfrak q$ 是包含 $\mathfrak h_A$ 的极小素理想而可行。
选取 $H_{A/R}$ 的生成元 $a_1,\ldots,a_r$。置
$$
B = A[x_1, \ldots, x_d, z_{ij}]/(x_i^n - \sum z_{ij}a_j)
$$
每个 $B_{a_j}$ 都在 $R$ 上光滑：它是
$A_{a_j}[x_1,\ldots,x_d]$ 上的多项式代数，而 $A_{a_j}$ 在 $k$ 上光滑。
因此 $B_{x_i}$ 在 $R$ 上光滑。令 $B\to C$ 为引理
\ref{smoothing-lemma-improve-presentation} 中构造的 $R$-代数同态；该构造还给出
$R$-代数缩回 $C\to B$，特别给出与上述图表相容的同态
$C\to\Lambda$。依构造，$C_{x_i}$ 是光滑 $R$-代数，且
$\Omega_{C_{x_i}/R}$ 自由。因此，可以找到 $c>0$，使 $x_i^c$
在 $C/R$ 中是严格标准元，见引理 \ref{smoothing-lemma-compare-standard}。
现在按照引理 \ref{smoothing-lemma-find-sequence} 选取
$\pi_1,\ldots,\pi_d\in\Lambda$，
其中 $n=n$、$e=8c$、$\mathfrak q=\mathfrak q$ 且
$I=\mathfrak h_A$。写
$\pi_i^n=\sum\lambda_{ij}a_j$，其中某些 $\pi_{ij}\in\Lambda$。
存在同态 $B\to\Lambda$，由 $x_i\mapsto\pi_i$ 及
$z_{ij}\mapsto\lambda_{ij}$ 给出。置 $R=k[x_1,\ldots,x_d]$。图表为
$$
\xymatrix{
R \ar[r] & B \ar[rd] \\
k \ar[u] \ar[r] & A \ar[u] \ar[r] & \Lambda
}
$$
现在把引理 \ref{smoothing-lemma-lift-solution} 应用于
$R \to C \to \Lambda \supset \mathfrak q$ 以及 $R$ 中的元素序列
$x_1^c,\ldots,x_d^c$。假设 (2) 显然成立。对 $R$ 而言，
直接检验可知假设 (1) 成立；对 $\Lambda$ 而言，它由我们对
$\pi_1,\ldots,\pi_d$ 的选择成立。（注意，若
$\text{Ann}_\Lambda(\pi)=\text{Ann}_\Lambda(\pi^2)$，则对所有
$c>0$ 都有 $\text{Ann}_\Lambda(\pi)=\text{Ann}_\Lambda(\pi^c)$。）
因此，只需消解
$$
R/(x_1^e, \ldots, x_d^e) \to C/(x_1^e, \ldots, x_d^e) \to
\Lambda/(\pi_1^e, \ldots, \pi_d^e) \supset
\mathfrak q/(\pi_1^e, \ldots, \pi_d^e)
$$
，其中 $e=8c$。由引理 \ref{smoothing-lemma-delocalize-height-zero}，
只需在 $\mathfrak q$ 处局部化后消解它。由于
$x_1,\ldots,x_d$ 在 $\Lambda_\mathfrak q$ 中映成正则序列，
$R_\mathfrak p\to\Lambda_\mathfrak q$ 平坦，见《代数》
引理 \ref{algebra-lemma-flat-over-regular}。所以
$$
R_\mathfrak p/(x_1^e, \ldots, x_d^e) \to
\Lambda_\mathfrak q/(\pi_1^e, \ldots, \pi_d^e)
$$
是 Artin 局部环之间的平坦环同态。此外，由假设，该同态在剩余域上
诱导可分域扩张。因此，由《代数》引理
\ref{algebra-lemma-colimit-syntomic} 和命题 \ref{smoothing-proposition-lift}，
这个同态是光滑代数的滤过余极限。所需解的存在性由《代数》
引理 \ref{algebra-lemma-when-colimit} 得出。
\end{proof}






\section{不可分剩余域}
\label{smoothing-section-inseparable}

\noindent
本节说明当剩余域扩张不可分时如何求解一个局部问题。

\begin{lemma}
\label{smoothing-lemma-helper}
设 $k$ 为特征 $p>0$ 的域，
$(\Lambda,\mathfrak m,K)$ 为 Artin 局部 $k$-代数。
假设 $\dim H_1(L_{K/k})<\infty$。则 $\Lambda$ 是 Artin 局部
$k$-代数 $A$ 的滤过余极限，其中每个同态 $A\to\Lambda$ 都平坦，
$\mathfrak m_A\Lambda=\mathfrak m$，且 $A$ 在 $k$ 上本质有限型。
\end{lemma}

\begin{proof}
注意，$A\to\Lambda$ 的平坦性蕴含该同态为单射，故本引理实际说明
$\Lambda$ 是这类子环 $A\subset\Lambda$ 的有向并。令 $n$ 为满足
$\mathfrak m^n=0$ 的最小整数。对 $n$ 归纳证明本引理。
$n=1$ 的情形显然，因为任一域扩张都是有限生成域扩张的并。

\medskip\noindent
选取生成 $\mathfrak m$ 的元素
$\lambda_1,\ldots,\lambda_d\in\mathfrak m$。由于 $K$ 在
$\mathbf{F}_p$ 上形式光滑（见《代数》引理
\ref{algebra-lemma-formally-smooth-extensions-easy}），可以找到环同态
$\sigma:K\to\Lambda$，它是商同态 $\Lambda\to K$ 的截面。
一般而言，$\sigma$ {\bf 不是} $k$-代数同态。给定 $\sigma$，定义
$$
\Psi_\sigma : K[x_1, \ldots, x_d] \longrightarrow \Lambda
$$
：在 $K$ 的元素上使用 $\sigma$，并令 $x_i\mapsto\lambda_i$。
断言：存在 $\sigma:K\to\Lambda$ 以及在 $k$ 上有限生成的子域
$k\subset F\subset K$，使 $k$ 在 $\Lambda$ 中的像包含于
$\Psi_\sigma(F[x_1,\ldots,x_d])$。

\medskip\noindent
对满足 $\mathfrak m^n=0$ 的最小整数 $n$ 作归纳，以证明该断言。
$n=1$ 时显然。若 $n>1$，置 $I=\mathfrak m^{n-1}$ 及
$\Lambda'=\Lambda/I$。由归纳假设，可认为已给定
$\sigma':K\to\Lambda'$ 以及有限生成的 $k\subset F'\subset K$，
使 $k\to\Lambda\to\Lambda'$ 的像包含于
$A'=\Psi_{\sigma'}(F'[x_1,\ldots,x_d])$。记
$\tau':k\to A'$ 为所诱导的同态。选取 $\sigma'$ 的一个提升
$\sigma:K\to\Lambda$；这可由前述 $K/\mathbf{F}_p$ 的形式光滑性做到。
留意，稍后可用某个导子 $D:K\to I$ 把 $\sigma$ 改为 $\sigma+D$。
置 $A=F[x_1,\ldots,x_d]/(x_1,\ldots,x_d)^n$。于是
$\Psi_\sigma$ 诱导环同态 $\Psi_\sigma:A\to\Lambda$；它与商同态
$\Lambda\to\Lambda'$ 的复合诱导一个核为幂零理想的满射 $A\to A'$。
选取 $\tau'$ 的提升 $\tau:k\to A$；这可由
$k/\mathbf{F}_p$ 的形式光滑性做到。于是得到两个从 $k$ 到
$\Lambda$ 的同态，即 $\Psi_\sigma\circ\tau:k\to\Lambda$ 和给定同态
$i:k\to\Lambda$。它们模 $I$ 相同，故其差为导子
$\theta=i-\Psi_\sigma\circ\tau:k\to I$。注意，若把 $\sigma$
改为 $\sigma+D$，则 $\theta$ 相应改为 $\theta-D|_k$。

\medskip\noindent
在 $k$ 中选取一族元素 $\{y_j\}_{j\in J}$，使其微分
$\text{d}y_j$ 构成 $\Omega_{k/\mathbf{F}_p}$ 的一组基。
对 $\mathbf{F}_p\subset k\subset K$，Jacobi--Zariski 序列为
$$
0 \to H_1(L_{K/k}) \to \Omega_{k/\mathbf{F}_p} \otimes K \to
\Omega_{K/\mathbf{F}_p} \to \Omega_{K/k} \to 0
$$
由于 $\dim H_1(L_{K/k})<\infty$，可以找到有限子集
$J_0\subset J$，使第一个同态的像包含于
$\bigoplus_{j\in J_0}K\text{d}y_j$。因此，元素
$\text{d}y_j$（$j\in J\setminus J_0$）在
$\Omega_{K/\mathbf{F}_p}$ 中的像在 $K$ 上线性无关。故可选取
$D:K\to I$，使 $\theta-D|_k=\xi\circ\text{d}$，其中 $\xi$ 是复合
$$
\Omega_{k/\mathbf{F}_p} = \bigoplus\nolimits_{j \in J} k \text{d}y_j
\longrightarrow \bigoplus\nolimits_{j \in J_0} k \text{d}y_j
\longrightarrow I
$$
对 $j\in J_0$，令 $f_j=\xi(\text{d}y_j)\in I$。按上面所述，
把 $\sigma$ 改为 $\sigma+D$。于是，对 $a\in k$，若在
$\Omega_{k/\mathbf{F}_p}$ 中有
$\text{d}a=\sum a_j\text{d}y_j$，则
$\theta(a)=\sum_{j\in J_0}a_jf_j$。注意，$I$ 由
$\lambda_1,\ldots,\lambda_d$ 中总次数
$|E|=\sum e_i=n-1$ 的单项式
$\lambda^E=\lambda_1^{e_1}\ldots\lambda_d^{e_d}$ 生成。写
$f_j=\sum_Ec_{j,E}\lambda^E$，其中 $c_{j,E}\in K$。
把 $F'$ 替换为 $F=F'(c_{j,E})$，断言即得证。

\medskip\noindent
按断言选取 $\sigma$ 和 $F$。$\Psi_\sigma$ 的核由有限多个多项式
$g_1,\ldots,g_t\in K[x_1,\ldots,x_d]$ 生成。向 $F$ 中添加有限多个
元素将其扩大后，可假设这些多项式的系数都在 $F$ 中。此时同态
$A = F[x_1, \ldots, x_d]/(g_1, \ldots, g_t) \to
K[x_1, \ldots, x_d]/(g_1, \ldots, g_t) = \Lambda$ 显然平坦。
由该断言，$A$ 是 $\Lambda$ 的 $k$-子代数。由于 $K$ 是这些子域
$F$ 的滤过并，$\Lambda$ 显然是这些代数的滤过余极限。最后，
由《代数》引理
\ref{algebra-lemma-essentially-of-finite-type-into-artinian-local}，
这些代数在 $k$ 上本质有限型。
\end{proof}

\begin{lemma}
\label{smoothing-lemma-solution-modulo}
设 $k$ 为特征 $p>0$ 的域，$\Lambda$ 为几何正则的诺特
$k$-代数，$\mathfrak q\subset\Lambda$ 为素理想。设 $n\geq1$ 为整数，
且 $E\subset\Lambda_\mathfrak q/\mathfrak q^n\Lambda_\mathfrak q$
为有限子集。则可以找到 $m\geq0$ 以及
$\varphi:k[y_1,\ldots,y_m]\to\Lambda$，使得
\begin{enumerate}
\item 令 $\mathfrak p=\varphi^{-1}(\mathfrak q)$，则
$\mathfrak q\Lambda_\mathfrak q=\mathfrak p\Lambda_\mathfrak q$，且
$k[y_1,\ldots,y_m]_\mathfrak p\to\Lambda_\mathfrak q$ 平坦；
\item 存在由 Artin 局部环同态组成的分解
$$
k[y_1, \ldots, y_m]_\mathfrak p/\mathfrak p^n k[y_1, \ldots, y_m]_\mathfrak p
\to D \to
\Lambda_\mathfrak q/\mathfrak q^n\Lambda_\mathfrak q
$$
，其中第一个箭头本质光滑，第二个箭头平坦；
\item 模 $\mathfrak q^n\Lambda_\mathfrak q$ 意义下，$E$ 包含于 $D$。
\end{enumerate}
\end{lemma}

\begin{proof}
置 $\bar\Lambda=\Lambda_\mathfrak q/\mathfrak q^n\Lambda_\mathfrak q$。
由《代数详论》命题
\ref{more-algebra-proposition-characterization-geometrically-regular}.
，有 $\dim H_1(L_{\kappa(\mathfrak q)/k})<\infty$。
按照引理 \ref{smoothing-lemma-helper}，选取含 $E$ 的 $A\subset\bar\Lambda$，
使 $A$ 是在 $k$ 上本质有限型的 Artin 局部环，$A\to\bar\Lambda$
平坦，且 $\mathfrak m_A$ 生成 $\bar\Lambda$ 的极大理想。
记剩余域 $F=A/\mathfrak m_A$，于是 $k\subset F\subset K$。
选取 $\lambda_1,\ldots,\lambda_t\in\Lambda$，它们在
$\bar\Lambda$ 中映到 $A$ 的元素，并且
$\text{d}\lambda_1,\ldots,\text{d}\lambda_t$ 的像构成
$\Omega_{F/k}$ 的一组基。考虑把 $y_j$ 映到 $\lambda_j$ 的同态
$\varphi':k[y_1,\ldots,y_t]\to\Lambda$，并置
$\mathfrak p'=(\varphi')^{-1}(\mathfrak q)$。由《代数详论》引理
\ref{more-algebra-lemma-geometrically-regular-over-field}
，环同态 $k[y_1,\ldots,y_t]_{\mathfrak p'}\to\Lambda_\mathfrak q$
平坦，且 $\Lambda_\mathfrak q/\mathfrak p'\Lambda_\mathfrak q$ 正则。
因此，可以再选取元素 $\lambda_{t+1},\ldots,\lambda_m\in\Lambda$；
它们在 $\bar\Lambda$ 中映入 $A$，且在
$\Lambda_\mathfrak q/\mathfrak p'\Lambda_\mathfrak q$ 中映成一组
正则参数系。由此得到满足性质 (1) 的
$\varphi:k[y_1,\ldots,y_m]\to\Lambda$，并且
$k[y_1, \ldots, y_m]_\mathfrak p/\mathfrak p^n k[y_1, \ldots, y_m]_\mathfrak p
\to \bar\Lambda$
经由 $A$ 分解。因此，
$k[y_1, \ldots, y_m]_\mathfrak p/\mathfrak p^n k[y_1, \ldots, y_m]_\mathfrak p
\to A$ 由《代数》引理
\ref{algebra-lemma-flatness-descends-more-general} 可知是平坦的。
依构造，剩余域扩张 $F/\kappa(\mathfrak p)$ 有限生成，且
$\Omega_{F/\kappa(\mathfrak p)}=0$。故由《代数详论》引理
\ref{more-algebra-lemma-cartier-equality}，它是有限可分扩张。因此
$k[y_1, \ldots, y_m]_\mathfrak p/\mathfrak p^n k[y_1, \ldots, y_m]_\mathfrak p
\to A$
由《代数》引理
\ref{algebra-lemma-essentially-of-finite-type-into-artinian-local}
可知是有限的。最后，由《代数》引理
\ref{algebra-lemma-characterize-etale}，它是平展的。
平展环同态当然本质光滑，故结论得证。
\end{proof}

\begin{lemma}
\label{smoothing-lemma-enlarge-solution-modulo}
设 $\varphi:k[y_1,\ldots,y_m]\to\Lambda$、$n$、$\mathfrak q$、
$\mathfrak p$ 以及
$$
k[y_1, \ldots, y_m]_\mathfrak p/\mathfrak p^n \to
D \to \Lambda_\mathfrak q/\mathfrak q^n \Lambda_\mathfrak q
$$
如引理 \ref{smoothing-lemma-solution-modulo} 所述。则对任意
$\lambda\in\Lambda\setminus\mathfrak q$，存在整数 $q>0$ 以及分解
$$
k[y_1, \ldots, y_m]_\mathfrak p/\mathfrak p^n \to
D \to D' \to \Lambda_\mathfrak q/\mathfrak q^n \Lambda_\mathfrak q
$$
，使 $D\to D'$ 是 Artin 局部环之间的本质光滑同态，
最后一个箭头平坦，并且 $\lambda^q$ 属于 $D'$。
\end{lemma}

\begin{proof}
置 $\bar\Lambda=\Lambda_\mathfrak q/\mathfrak q^n\Lambda_\mathfrak q$，
并令 $\bar\lambda$ 为 $\lambda$ 在 $\bar\Lambda$ 中的像，
$\alpha\in\kappa(\mathfrak q)$ 为 $\lambda$ 在剩余域中的像。
设 $k\subset F\subset\kappa(\mathfrak q)$ 为 $D$ 的剩余域。
若 $\alpha\in F$，则可找到 $x\in D$，使
$x\bar\lambda=1\bmod\mathfrak q$。若 $q$ 可被 $p$ 整除，则
$(x\bar\lambda)^q=1\bmod(\mathfrak q)^q$，故
$\bar\lambda^q\in D$。若 $\alpha$ 在 $F$ 上超越，则可取
$D'=(D[\bar\lambda])_\mathfrak m$，即把由 $D$ 和 $\bar\lambda$
生成的子环在
$\mathfrak m=D[\bar\lambda]\cap\mathfrak q\bar\Lambda$ 处局部化。
这是可行的，因为此时 $D[\bar\lambda]$ 实际上是 $D$ 上的多项式代数。
最后，若 $\lambda\bmod\mathfrak q$ 在 $F$ 上代数，则可以找到
$p$ 的幂 $q$，使 $\alpha^q$ 在 $F$ 上可分代数，
见《域》一章第 \ref{fields-section-algebraic} 节。
注意，$D$ 和 $\bar\Lambda$ 都是 Hensel 局部环，见《代数》引理
\ref{algebra-lemma-local-dimension-zero-henselian}。令 $D\to D'$ 为
有限平展扩张，其剩余域扩张为 $F(\alpha^q)/F$，见《代数》引理
\ref{algebra-lemma-henselian-cat-finite-etale}。由于 $\bar\Lambda$
是 Hensel 环，且其剩余域包含 $F(\alpha^q)$，可以找到分解
$D'\to\bar\Lambda$。由论证的第一部分可知，对某个 $q'>0$，有
$\bar\lambda^{qq'}\in D'$。
\end{proof}

\begin{lemma}
\label{smoothing-lemma-resolve-general}
设 $k \to A \to \Lambda \supset \mathfrak q$ 如情形
\ref{smoothing-situation-local} 所述，其中
\begin{enumerate}
\item $k$ 是特征 $p>0$ 的域；
\item $\Lambda$ 是诺特环，且在 $k$ 上几何正则；
\item $\mathfrak q$ 是包含 $\mathfrak h_A$ 的极小素理想。
\end{enumerate}
则 $k \to A \to \Lambda \supset \mathfrak q$ 可消解。
\end{lemma}

\begin{proof}
按给定顺序完成以下步骤即可证明本引理。下面将逐一说明这些步骤。
\begin{enumerate}
\item
\label{smoothing-item-power}
选取整数 $N>0$，使
$\mathfrak q^N\Lambda_\mathfrak q \subset H_{A/k}\Lambda_\mathfrak q$.
\item
\label{smoothing-item-generators}
选取理想 $H_{A/R}$ 的生成元 $a_1,\ldots,a_t\in A$。
\item
\label{smoothing-item-dimension}
置 $d=\dim(\Lambda_\mathfrak q)$。
\item
\label{smoothing-item-standardizer}
置 $B=A[x_1,\ldots,x_d,z_{ij}]/(x_i^{2N}-\sum z_{ij}a_j)$。
\item
\label{smoothing-item-elkik}
把 $B$ 看作 $k[x_1,\ldots,x_d]$-代数，并令 $B\to C$ 如引理
\ref{smoothing-lemma-improve-presentation} 所述。我们还得到截面 $C\to B$。
\item
\label{smoothing-item-strictly-standard}
选取 $c>0$，使每个 $x_i^c$ 在 $C$ 中相对于
$k[x_1,\ldots,x_d]$ 都是严格标准元。
\item
\label{smoothing-item-set-n}
置 $n=N+dc$ 及 $e=8c$。
\item
\label{smoothing-item-set-E}
令 $E\subset\Lambda_\mathfrak q/\mathfrak q^n\Lambda_\mathfrak q$
为 $A$ 作为 $k$-代数的一组生成元的像。
\item
\label{smoothing-item-NP}
选取整数 $m$、$k$-代数同态
$\varphi : k[y_1, \ldots, y_m] \to \Lambda$
以及由 Artin 局部环组成的分解
$$
k[y_1, \ldots, y_m]_\mathfrak p/\mathfrak p^n k[y_1, \ldots, y_m]_\mathfrak p
\to D \to
\Lambda_\mathfrak q/\mathfrak q^n\Lambda_\mathfrak q
$$
，使第一个箭头本质光滑，第二个箭头平坦，$E$ 包含于 $D$；并且令
$\mathfrak p=\varphi^{-1}(\mathfrak q)$ 时，同态
$k[y_1,\ldots,y_m]_\mathfrak p\to\Lambda_\mathfrak q$ 平坦，且
$\mathfrak p\Lambda_\mathfrak q=\mathfrak q\Lambda_\mathfrak q$。
\item
\label{smoothing-item-choose-pii}
选取 $\pi_1,\ldots,\pi_d\in\mathfrak p$，它们映成
$k[y_1,\ldots,y_m]_\mathfrak p$ 的一组正则参数系。
\item
\label{smoothing-item-set-R}
令 $R=k[y_1,\ldots,y_m,t_1,\ldots,t_m]$ 且 $\gamma_i=\pi_it_i$。
\item
\label{smoothing-item-modify-pii}
必要时修改 $\pi_i$ 的选择，使得对 $i=1,\ldots,d$ 有
$$
\text{Ann}_{R/(\gamma_1^e, \ldots, \gamma_{i - 1}^e)R}(\gamma_i)
=
\text{Ann}_{R/(\gamma_1^e, \ldots, \gamma_{i - 1}^e)R}(\gamma_i^2)
$$
\item
\label{smoothing-item-choose-deltai}
存在 $\delta_1,\ldots,\delta_d\in\Lambda$，其中
$\delta_i \not \in \mathfrak q$，以及分解
$D \to D' \to \Lambda_\mathfrak q/\mathfrak q^n\Lambda_\mathfrak q$
，其中 $D'$ 是 Artin 局部环，$D\to D'$ 本质光滑，且同态
$D'\to\Lambda_\mathfrak q/\mathfrak q^n\Lambda_\mathfrak q$ 平坦；
令 $\pi_i'=\delta_i\pi_i$，则对 $i=1,\ldots,d$ 有
\begin{enumerate}
\item 在 $\Lambda$ 中，$(\pi_i')^{2N}=\sum a_j\lambda_{ij}$，其中
$\lambda_{ij}\bmod\mathfrak q^n\Lambda_\mathfrak q$ 是 $D'$ 的元素；
\item $\text{Ann}_{\Lambda/({\pi'}_1^e, \ldots, {\pi'}_{i - 1}^e)}({\pi'}_i) =
\text{Ann}_{\Lambda/({\pi'}_1^e, \ldots, {\pi'}_{i - 1}^e)}({\pi'}_i^2)$；
\item $\delta_i\bmod\mathfrak q^n\Lambda_\mathfrak q$ 是 $D'$ 的元素。
\end{enumerate}
\item
\label{smoothing-item-map-B-Lambda}
令 $x_i\mapsto\pi_i'$、$z_{ij}\mapsto\lambda_{ij}$，以此定义
$B\to\Lambda$。把 $B\to\Lambda$ 与缩回 $C\to B$ 复合，
定义 $C\to\Lambda$。
\item
\label{smoothing-item-set-map-R}
定义 $R\to\Lambda$：在 $k[y_1,\ldots,y_m]$ 上取 $\varphi$，
并令 $t_i\mapsto\delta_i$。再引入同态
$$
k[x_1, \ldots, x_d]
\longrightarrow
R = k[y_1, \ldots, y_m, t_1, \ldots, t_d]
$$
，令 $x_i\mapsto\gamma_i=\pi_it_i$。
\item
\label{smoothing-item-first-solve}
只需消解
$$
R
\to
C \otimes_{k[x_1, \ldots, x_d]} R
\to
\Lambda \supset \mathfrak q
$$
\item
\label{smoothing-item-set-I}
置 $I=(\gamma_1^e,\ldots,\gamma_d^e)\subset R$。
\item
\label{smoothing-item-second-resolve}
只需消解
$$
R/I
\to
C \otimes_{k[x_1, \ldots, x_d]} R/I
\to
\Lambda/I\Lambda \supset \mathfrak q/I\Lambda
$$
\item
\label{smoothing-item-set-r}
记 $\mathfrak q$ 在
$R=k[y_1,\ldots,y_m,t_1,\ldots,t_d]$ 中的逆像为
$\mathfrak r\subset R$。
\item
\label{smoothing-item-third-resolve}
只需消解
$$
(R/I)_\mathfrak r \to
C \otimes_{k[x_1, \ldots, x_d]} (R/I)_\mathfrak r \to
\Lambda_\mathfrak q/I\Lambda_\mathfrak q
\supset
\mathfrak q\Lambda_\mathfrak q/I\Lambda_\mathfrak q
$$
\item
\label{smoothing-item-fourth-resolve}
在 $k[y_1,\ldots,y_m]$ 中置 $J=(\pi_1^e,\ldots,\pi_d^e)$。
\item
\label{smoothing-item-fifth-resolve}
只需消解
$$
(R/JR)_\mathfrak p \to
C \otimes_{k[x_1, \ldots, x_d]} (R/JR)_\mathfrak p \to
\Lambda_\mathfrak q/J\Lambda_\mathfrak q
\supset
\mathfrak q\Lambda_\mathfrak q/J\Lambda_\mathfrak q
$$
\item
\label{smoothing-item-sixth-resolve}
只需消解
$$
(R/\mathfrak p^nR)_\mathfrak p \to
C \otimes_{k[x_1, \ldots, x_d]} (R/\mathfrak p^nR)_\mathfrak p \to
\Lambda_\mathfrak q/\mathfrak q^n\Lambda_\mathfrak q
\supset
\mathfrak q\Lambda_\mathfrak q/\mathfrak q^n\Lambda_\mathfrak q
$$
\item
\label{smoothing-item-seventh-resolve}
只需消解
$$
(R/\mathfrak p^nR)_\mathfrak p \to
B \otimes_{k[x_1, \ldots, x_d]} (R/\mathfrak p^nR)_\mathfrak p \to
\Lambda_\mathfrak q/\mathfrak q^n\Lambda_\mathfrak q
\supset
\mathfrak q\Lambda_\mathfrak q/\mathfrak q^n\Lambda_\mathfrak q
$$
\item
\label{smoothing-item-eighth-resolve}
借助给定同态
$k[y_1, \ldots, y_m]_\mathfrak p/\mathfrak p^n k[y_1, \ldots, y_m]_\mathfrak p
\to D'$ 以及赋值 $t_i\mapsto t_i$，赋予环
$D'[t_1,\ldots,t_d]$ 一个
$R_\mathfrak p/\mathfrak p^nR_\mathfrak p$-代数结构。只需找到分解
$$
B \otimes_{k[x_1, \ldots, x_d]} (R/\mathfrak p^nR)_\mathfrak p
\to D'[t_1, \ldots, t_d] \to
\Lambda_\mathfrak q/\mathfrak q^n\Lambda_\mathfrak q
$$
，其中第二个箭头把 $t_i$ 映到 $\delta_i$，并诱导给定同态
$D'\to\Lambda_\mathfrak q/\mathfrak q^n\Lambda_\mathfrak q$。
\item
\label{smoothing-item-done}
由上面对 $D'$ 的选择，这样的分解存在。
\end{enumerate}
下面逐步给出论证；仅引入记号的步骤不再说明。

\medskip\noindent
关于 (\ref{smoothing-item-power})。这是可行的，因为 $\mathfrak q$ 是包含
$\mathfrak h_A=\sqrt{H_{A/k}\Lambda}$ 的极小素理想。

\medskip\noindent
关于 (\ref{smoothing-item-strictly-standard})。注意，$A_{a_i}$ 在 $k$ 上光滑。
因此，$B_{a_j}$ 同构于 $A_{a_j}[x_1,\ldots,x_d]$ 上的多项式代数，
从而在 $k[x_1,\ldots,x_d]$ 上光滑。于是 $B_{x_i}$ 在
$k[x_1,\ldots,x_d]$ 上光滑。由引理
\ref{smoothing-lemma-improve-presentation}，$C_{x_i}$ 在
$k[x_1,\ldots,x_d]$ 上光滑，且其微分模有限自由。因此，由引理
\ref{smoothing-lemma-compare-standard}，$x_i$ 的某个幂在 $C$ 中相对于
$k[x_1,\ldots,x_n]$ 是严格标准元。

\medskip\noindent
关于 (\ref{smoothing-item-NP})。这由应用引理 \ref{smoothing-lemma-solution-modulo} 得出。

\medskip\noindent
关于 (\ref{smoothing-item-choose-pii})。由于
$k[y_1, \ldots, y_m]_\mathfrak p \to \Lambda_\mathfrak q$
依构造是平坦的，且
$\mathfrak p\Lambda_\mathfrak q=\mathfrak q\Lambda_\mathfrak q$，
由《代数》引理 \ref{algebra-lemma-dimension-base-fibre-equals-total}，
有 $\dim(k[y_1,\ldots,y_m]_\mathfrak p)=d$。因而可以找到
$\pi_1,\ldots,\pi_d\in\Lambda$，它们在 $\Lambda_\mathfrak q$
中映成一组正则参数系。

\medskip\noindent
关于 (\ref{smoothing-item-modify-pii})。由《代数》引理
\ref{algebra-lemma-regular-ring-CM}，序列
$\pi_1,\ldots,\pi_d$ 的任意排列都是
$k[y_1,\ldots,y_m]_\mathfrak p$ 中的正则序列。因此
$\gamma_1=\pi_1t_1,\ldots,\gamma_d=\pi_dt_d$ 是
$R_\mathfrak p=k[y_1,\ldots,y_m]_\mathfrak p[t_1,\ldots,t_d]$
中的正则序列，见《代数》引理
\ref{algebra-lemma-regular-sequence-in-polynomial-ring}。
令 $S=k[y_1,\ldots,y_m]\setminus\mathfrak p$，于是
$R_\mathfrak p=S^{-1}R$。注意，即使把各个 $\pi_i$ 乘以 $S$ 中元素，
$\pi_1,\ldots,\pi_d$ 和 $\gamma_1,\ldots,\gamma_d$ 仍是正则序列。
假设
$$
\text{Ann}_{R/(\gamma_1^e, \ldots, \gamma_{i - 1}^e)R}(\gamma_i)
=
\text{Ann}_{R/(\gamma_1^e, \ldots, \gamma_{i - 1}^e)R}(\gamma_i^2)
$$
对某个 $t\in\{0,\ldots,d\}$ 及 $i=1,\ldots,t$ 成立。由《代数》
引理 \ref{algebra-lemma-regular-sequence-powers}，
$\gamma_1^e,\ldots,\gamma_t^e,\gamma_{t+1}$ 是 $S^{-1}R$
中的正则序列，故
$$
\text{Ann}_{S^{-1}R/(\gamma_1^e, \ldots, \gamma_{i - 1}^e)}(\gamma_i) =
\text{Ann}_{S^{-1}R/(\gamma_1^e, \ldots, \gamma_{i - 1}^e)}(\gamma_i^2).
$$
因此，
$$
\text{Ann}_{R/(\gamma_1^e, \ldots, \gamma_t^e)R}(\gamma_{t + 1})
=
\text{Ann}_{R/(\gamma_1^e, \ldots, \gamma_t^e)R}(\gamma_{t + 1}^2)
$$
由引理 \ref{smoothing-lemma-ogoma}，取某个 $s\in S$ 并以
$s\pi_{t+1}$ 替换 $\pi_{t+1}$ 后，上式成立。对 $t$ 归纳，
即得到所需序列。

\medskip\noindent
关于 (\ref{smoothing-item-choose-deltai})。令 $S=\Lambda\setminus\mathfrak q$，
于是 $\Lambda_\mathfrak q=S^{-1}\Lambda$；并置
$\bar\Lambda=\Lambda_\mathfrak q/\mathfrak q^n\Lambda_\mathfrak q$。
假设已有 $t\in\{0,\ldots,d\}$、$\delta_1,\ldots,\delta_t\in S$，
以及 (\ref{smoothing-item-choose-deltai}) 中的分解
$D\to D'\to\bar\Lambda$，使 (a)、(b)、(c) 对
$i=1,\ldots,t$ 成立。由 (\ref{smoothing-item-power})，
$\mathfrak q^N\Lambda_\mathfrak q\subset H_{A/k}\Lambda_\mathfrak q$，
所以 $\pi_{t+1}^N\in H_{A/k}\Lambda_\mathfrak q$，进而
$\pi_{t+1}^N\in H_{A/k}\bar\Lambda$。由于
$D'\to\bar\Lambda$ 忠实平坦，由《代数》引理
\ref{algebra-lemma-faithfully-flat-universally-injective} 又有
$\pi_{t+1}^N\in H_{A/k}D'$。回忆 $H_{A/k}=(a_1,\ldots,a_t)$。
设在 $D'$ 中 $\pi_{t+1}^N=\sum a_jd_j$，并选取
$c_j\in\Lambda_\mathfrak q$ 提升 $d_j\in D'$。则
$\pi_{t+1}^N=\sum c_ja_j+\epsilon$，其中
$\epsilon\in\mathfrak q^n\Lambda_\mathfrak q\subset
\mathfrak q^{n-N}H_{A/k}\Lambda_\mathfrak q$。对某些
$c'_j\in\mathfrak q^{n-N}\Lambda_\mathfrak q$，写
$\epsilon=\sum a_jc'_j$。于是
$\pi_{t+1}^{2N}=\sum(\pi_{t+1}^Nc_j+\pi_{t+1}^Nc'_j)a_j$。
注意，$\pi_{t+1}^Nc'_j$ 在 $\bar\Lambda$ 中映到零；这个虽显然却
关键的观察将确保稍后 (a) 成立。
现在选取 $s\in S$，使存在 $\mu_{t + 1j}\in\Lambda$，一方面在
$S^{-1}\Lambda$ 中有
$\pi_{t+1}^Nc_j+\pi_{t+1}^Nc'_j=\mu_{t + 1j}/s^{2N}$，
另一方面在 $\Lambda$ 中有
$(s\pi_{t+1})^{2N}=\sum\mu_{t + 1j}a_j$（略去一个小细节）。
还可把 $s$ 替换为它的某个幂并扩大 $D'$，使 $s$ 的像属于 $D'$。
在这些选择下，$\mu_{t + 1j}$ 映到 $s^{2N}d_j$，后者属于 $D'$。
注意，由我们对 $\varphi$ 的选择，$\pi_1,\ldots,\pi_d$ 是
$S^{-1}\Lambda$ 中的一组正则参数序列。因此，由《代数》引理
\ref{algebra-lemma-regular-ring-CM}，它们构成
$\Lambda_\mathfrak q$ 中的正则序列。再由《代数》引理
\ref{algebra-lemma-regular-sequence-powers}，
${\pi'}_1^e,\ldots,{\pi'}_t^e,s\pi_{t+1}$ 是
$S^{-1}\Lambda$ 中的正则序列。所以
$$
\text{Ann}_{S^{-1}\Lambda/({\pi'}_1^e, \ldots, {\pi'}_t^e)}(s\pi_{t + 1}) =
\text{Ann}_{S^{-1}\Lambda/({\pi'}_1^e, \ldots, {\pi'}_t^e)}((s\pi_{t + 1})^2).
$$
因此，可以应用引理 \ref{smoothing-lemma-ogoma} 找到 $s'\in S$，使得
$$
\text{Ann}_{\Lambda/({\pi'}_1^e, \ldots, {\pi'}_t^e)}((s')^qs\pi_{t + 1})
=
\text{Ann}_{\Lambda/({\pi'}_1^e, \ldots, {\pi'}_t^e)}(((s')^qs\pi_{t + 1})^2).
$$
对任意 $q>0$ 都成立。由引理 \ref{smoothing-lemma-enlarge-solution-modulo}，
可以选择 $q$ 并扩大 $D'$，使 $(s')^q$ 的像属于 $D'$。
置 $\delta_{t+1}=(s')^qs$，便得 (a)、(b)、(c) 对
$i=1,\ldots,t+1$ 成立。对于 (a)，注意可取
$\lambda_{t + 1j}=(s')^{2Nq}\mu_{t + 1j}$。对 $t$ 归纳即得结论。

\medskip\noindent
关于 (\ref{smoothing-item-first-solve})。依构造，
$H_{(C\otimes_{k[x_1,\ldots,x_d]}R)/R}\Lambda$ 的根包含
$\mathfrak h_A$。事实上，元素 $a_j\in H_{A/k}$ 映到
$H_{B/k[x_1,\ldots,x_n]}$ 的元素，继而映到
$H_{C/k[x_1,\ldots,x_n]}$ 的元素，所以 $a_j\otimes1$ 映到
$H_{C\otimes_{k[x_1,\ldots,x_d]}R/R}$ 的元素。此外，若有
$C\otimes_{k[x_1,\ldots,x_n]}R\to T\to\Lambda$，它是下列情形的解：
$$
R
\to
C \otimes_{k[x_1, \ldots, x_d]} R
\to
\Lambda \supset \mathfrak q
$$
，则因 $R$ 在 $k$ 上光滑，有 $H_{T/R}\subset H_{T/k}$。
所以 $T$ 也是原情形 $k\to A\to\Lambda\supset\mathfrak q$ 的解。

\medskip\noindent
关于 (\ref{smoothing-item-second-resolve})。把引理 \ref{smoothing-lemma-lift-solution}
应用于
$R \to C \otimes_{k[x_1, \ldots, x_d]} R
\to \Lambda \supset \mathfrak q$ 以及元素序列
$\gamma_1^c,\ldots,\gamma_d^c$ 即得。注意，由于 $x_i^c$
在 $C$ 中相对于 $k[x_1,\ldots,x_d]$ 是严格标准元，引理
\ref{smoothing-lemma-strictly-standard-base-change} 表明 $\gamma_i^c$ 在
$C\otimes_{k[x_1,\ldots,x_d]}R$ 中相对于 $R$ 是严格标准元。
引理 \ref{smoothing-lemma-lift-solution} 的另一个假设由步骤
(\ref{smoothing-item-modify-pii}) 和 (\ref{smoothing-item-choose-deltai}) 成立。

\medskip\noindent
关于 (\ref{smoothing-item-third-resolve})。把引理
\ref{smoothing-lemma-delocalize-height-zero} 应用于
(\ref{smoothing-item-second-resolve}) 中的情形。以下论证中的目标环均为
Artin 局部环，因此我们所求的是经由源环上某个光滑代数 $T$ 的分解。

\medskip\noindent
关于 (\ref{smoothing-item-fifth-resolve})。假设
$C\otimes_{k[x_1,\ldots,x_d]}(R/JR)_\mathfrak p\to
T\to\Lambda_\mathfrak q/J\Lambda_\mathfrak q$ 是下列情形的解：
$$
(R/JR)_\mathfrak p \to
C \otimes_{k[x_1, \ldots, x_d]} (R/JR)_\mathfrak p \to
\Lambda_\mathfrak q/J\Lambda_\mathfrak q
\supset
\mathfrak q\Lambda_\mathfrak q/J\Lambda_\mathfrak q
$$
则
$C\otimes_{k[x_1,\ldots,x_d]}(R/I)_\mathfrak r\to T_\mathfrak r\to
\Lambda_\mathfrak q/I\Lambda_\mathfrak q$
是 (\ref{smoothing-item-third-resolve}) 中情形的解。

\medskip\noindent
关于 (\ref{smoothing-item-sixth-resolve})。$n=N+dc$ 足够大，因而
$\mathfrak p^nk[y_1,\ldots,y_m]_\mathfrak p\subset J_\mathfrak p$
且 $\mathfrak q^n\Lambda_\mathfrak q\subset J\Lambda_\mathfrak q$。
因此，若有
$C \otimes_{k[x_1, \ldots, x_d]} (R/\mathfrak p^nR)_\mathfrak p \to
T \to \Lambda_\mathfrak q/\mathfrak q^n\Lambda_\mathfrak q$
作为 (\ref{smoothing-item-fifth-resolve} 的一个解，则可取 $T/JT$ 作为
(\ref{smoothing-item-sixth-resolve}) 的解。

\medskip\noindent
关于 (\ref{smoothing-item-seventh-resolve})。这是因为在 $R$-代数范畴中
已有截面 $C\to B$。

\medskip\noindent
关于 (\ref{smoothing-item-eighth-resolve})。这是因为 $D'$ 在 Artin 局部环
$k[y_1, \ldots, y_m]_\mathfrak p/\mathfrak p^n k[y_1, \ldots, y_m]_\mathfrak p$
上本质光滑，并且
$$
R_\mathfrak p/\mathfrak p^nR_\mathfrak p =
k[y_1, \ldots, y_m]_\mathfrak p/
\mathfrak p^n k[y_1, \ldots, y_m]_\mathfrak p[t_1, \ldots, t_d].
$$
因此，$D'[t_1,\ldots,t_d]$ 是光滑
$R_\mathfrak p/\mathfrak p^nR_\mathfrak p$-代数的滤过余极限，并且
$B \otimes_{k[x_1, \ldots, x_d]} (R_\mathfrak p/\mathfrak p^nR_\mathfrak p)$
经由其中一个代数分解。

\medskip\noindent
关于 (\ref{smoothing-item-done})。证明的最后一个微妙之处在于，不能直接使用
$B\to D'$ 这个把 $x_i$ 映到 $\pi_i'$ 在 $D'$ 中的像、把
$z_{ij}$ 映到 $\lambda_{ij}$ 在 $D'$ 中之像的同态，因为需要图表
$$
\xymatrix{
B \ar[r] & D'[t_1, \ldots, t_d] \\
k[x_1, \ldots, x_d] \ar[r] \ar[u] &
R_\mathfrak p/\mathfrak p^nR_\mathfrak p \ar[u]
}
$$
可交换，并且需要复合
$B \to D'[t_1, \ldots, t_d] \to
\Lambda_\mathfrak q/\mathfrak q^n\Lambda_\mathfrak q$
等于 (\ref{smoothing-item-map-B-Lambda}) 中的同态。这就要求把 $x_i$ 映到
$\pi_it_i$ 在 $D'[t_1,\ldots,t_d]$ 中的像。因此，把 $z_{ij}$ 映到
$\lambda_{ij}t_i^{2N}/\delta_i^{2N}$ 在
$D'[t_1,\ldots,t_d]$ 中的像，一切便都清楚了。
\end{proof}







\section{主定理}
\label{smoothing-section-main}

\noindent
本节汇总前面的论证。

\begin{theorem}[Popescu]
\label{smoothing-theorem-popescu}
任意诺特环之间的正则同态都是光滑环同态的滤过余极限。
\end{theorem}

\begin{proof}
由引理 \ref{smoothing-lemma-reduce-to-field}，只需对 $k\to\Lambda$ 证明结论，
其中 $\Lambda$ 是在 $k$ 上几何正则的诺特环。设
$k\to A\to\Lambda$ 是一个分解，其中 $A$ 为有限型 $k$-代数。
只需构造分解 $A\to B\to\Lambda$，其中 $B$ 为有限型且
$\mathfrak h_B=\Lambda$，见引理 \ref{smoothing-lemma-final-solve}。
因此，可以对理想 $\mathfrak h_A$ 作诺特归纳。选取
$\mathfrak q\supset\mathfrak h_A$，使 $\mathfrak q$ 是包含
$\mathfrak h_A$ 的极小素理想。现在只需消解
$k\to A\to\Lambda\supset\mathfrak q$
（其定义见情形 \ref{smoothing-situation-local} 后的正文）。若 $k$ 的特征为零，
结论由引理 \ref{smoothing-lemma-resolve-special} 得出；若 $k$ 的特征为
$p>0$，结论由引理 \ref{smoothing-lemma-resolve-general} 得出。
\end{proof}




\section{G-环的逼近性质}
\label{smoothing-section-approximation-G-rings}

\noindent
设 $R$ 为诺特局部环。此时，$R$ 是 G-环当且仅当环同态
$R\to R^\wedge$ 正则，见《代数详论》引理
\ref{more-algebra-lemma-check-G-ring-maximal-ideals}。在这种情况下，
$R$ 的 Hensel 化 $R^h$ 及严格 Hensel 化 $R^{sh}$ 都是 G-环，
见《代数详论》引理 \ref{more-algebra-lemma-henselization-G-ring}。
此外，在一个域、完备局部环、$\mathbf{Z}$ 或特征零 Dedekind 环上
本质有限型的任意代数都是 G-环，见《代数详论》命题
\ref{more-algebra-proposition-ubiquity-G-ring}。因此，有大量环可应用
下面的结果。

\medskip\noindent
设 $R$ 为环，$f_1,\ldots,f_m\in R[x_1,\ldots,x_n]$，且 $S$ 为
$R$-代数。在这种情形下，称向量 $(a_1,\ldots,a_n)\in S^n$
是{\it $S$ 中的解}，当且仅当
$$
f_j(a_1, \ldots, a_n) = 0 \text{ 于 } S \text{ 中，对 }
j = 1, \ldots, m
$$
代数几何中的一个重要问题自然是判定多项式方程组何时有解。
下面的定理说明，在诺特局部环的完备化中有解，往往足以推出
在该环的 Hensel 化中有解。

\begin{theorem}
\label{smoothing-theorem-approximation-property}
设 $R$ 为诺特局部环，且
$f_1,\ldots,f_m\in R[x_1,\ldots,x_n]$。假设
$(a_1,\ldots,a_n)\in(R^\wedge)^n$ 是 $R^\wedge$ 中的解。
若 $R$ 是 Hensel G-环，则对每个整数 $N$，都存在 $R$ 中的解
$(b_1,\ldots,b_n)\in R^n$，使
$a_i-b_i\in\mathfrak m^NR^\wedge$。
\end{theorem}

\begin{proof}
选取 $c_i\in R$，使 $a_i-c_i\in\mathfrak m^N$。选取生成元
$\mathfrak m^N=(d_1,\ldots,d_M)$，并写成
$a_i=c_i+\sum a_{i,l}d_l$。考虑多项式环 $R[x_{i,l}]$ 及元素
$$
g_j = f_j(c_1 + \sum x_{1, l} d_l , \ldots, c_n + \sum x_{n, l} d_{n, l})
\in R[x_{i, l}]
$$
方程组 $g_j=0$ 有解 $(a_{i,l})$。假设可以证明它在 $R$ 中有解
$(b_{i,l})$，则 $b_i=c_i+\sum b_{i,l}d_l$ 是 $f_j=0$ 的解，
且模 $\mathfrak m^N$ 与 $a_i$ 同余。因此，只需证明在 $R^\wedge$
上可解蕴含在 $R$ 上可解。

\medskip\noindent
令 $A\subset R^\wedge$ 为由 $a_1,\ldots,a_n$ 生成的 $R$-子代数。
我们假设了 $R$ 是 G-环，即 $R\to R^\wedge$ 正则，因此存在分解
$$
A \to B \to R^\wedge
$$
，其中 $B$ 在 $R$ 上光滑，见定理 \ref{smoothing-theorem-popescu}。
记剩余域 $\kappa=R/\mathfrak m$。它也是 $R^\wedge$ 的剩余域，
所以得到交换图
$$
\xymatrix{
B \ar[rd] \ar@{..>}[r] & R' \ar@{..>}[d] \\
R \ar[r] \ar[u] & \kappa
}
$$
由于竖直箭头光滑，《代数详论》引理
\ref{more-algebra-lemma-lift-section-smooth-morphism} 蕴含存在平展环同态
$R\to R'$，它诱导同构 $R/\mathfrak m\to R'/\mathfrak mR'$，
并存在使上图交换的 $R$-代数同态 $B\to R'$。由于 $R$ 是 Hensel 环，
$R\to R'$ 有截面，见《代数》引理
\ref{algebra-lemma-characterize-henselian}。令 $b_i\in R$ 为 $a_i$
经环同态 $A\to B\to R'\to R$ 所得的像。由于这些全是
$R$-代数同态，$(b_1,\ldots,b_n)$ 是 $R$ 中的解。
\end{proof}

\noindent
给定诺特局部环 $(R,\mathfrak m)$、平展环同态 $R\to R'$，
以及位于 $\mathfrak m$ 之上且满足
$\kappa(\mathfrak m)=\kappa(\mathfrak m')$ 的极大理想
$\mathfrak m'\subset R'$，则有包含关系
$$
R \subset R_{\mathfrak m'} \subset R^h \subset R^\wedge,
$$
，见《代数》引理 \ref{algebra-lemma-henselian-functorial-prepare}
和《代数详论》引理 \ref{more-algebra-lemma-henselization-noetherian}。

\begin{theorem}
\label{smoothing-theorem-approximation-property-variant}
设 $R$ 为诺特局部环，且
$f_1,\ldots,f_m\in R[x_1,\ldots,x_n]$。假设
$(a_1,\ldots,a_n)\in(R^\wedge)^n$ 是一个解。若 $R$ 是 G-环，
则对每个整数 $N$，都存在
\begin{enumerate}
\item 平展环同态 $R\to R'$；
\item 位于 $\mathfrak m$ 之上的极大理想 $\mathfrak m'\subset R'$；
\item $R'$ 中的解 $(b_1,\ldots,b_n)\in(R')^n$，
\end{enumerate}
使得 $\kappa(\mathfrak m)=\kappa(\mathfrak m')$ 且
$a_i-b_i\in(\mathfrak m')^NR^\wedge$。
\end{theorem}

\begin{proof}
可以由定理 \ref{smoothing-theorem-approximation-property} 推出本定理：
由《代数详论》引理 \ref{more-algebra-lemma-henselization-G-ring}，
Hensel 化 $R^h$ 是 G-环；再把 $R^h$ 写成平展扩张 $R'$ 的有向余极限即可。
不过，我们改为在当前情形下重做上一条定理的证明。

\medskip\noindent
选取 $c_i\in R$，使 $a_i-c_i\in\mathfrak m^N$。选取生成元
$\mathfrak m^N=(d_1,\ldots,d_M)$，并写成
$a_i=c_i+\sum a_{i,l}d_l$。考虑多项式环 $R[x_{i,l}]$ 及元素
$$
g_j = f_j(c_1 + \sum x_{1, l} d_l , \ldots, c_n + \sum x_{n, l} d_{n, l})
\in R[x_{i, l}]
$$
方程组 $g_j=0$ 有解 $(a_{i,l})$。假设可以证明：对某个平展环同态
$R\to R'$，存在极大理想 $\mathfrak m'$ 满足
$\kappa(\mathfrak m)=\kappa(\mathfrak m')$，且 $g_j=0$ 在 $R'$
中有解 $(b_{i,l})$。则 $b_i=c_i+\sum b_{i,l}d_l$ 是 $f_j=0$
的解，并且模 $(\mathfrak m')^N$ 与 $a_i$ 同余。因此，只需证明
$R^\wedge$ 上可解蕴含在某个平展环扩张上可解，而且该扩张在
$\mathfrak m$ 上方的某个素理想处诱导平凡剩余域扩张。

\medskip\noindent
令 $A\subset R^\wedge$ 为由 $a_1,\ldots,a_n$ 生成的 $R$-子代数。
我们假设了 $R$ 是 G-环，即 $R\to R^\wedge$ 正则，故存在分解
$$
A \to B \to R^\wedge
$$
，其中 $B$ 在 $R$ 上光滑，见定理 \ref{smoothing-theorem-popescu}。
记剩余域 $\kappa=R/\mathfrak m$。它也是 $R^\wedge$ 的剩余域，
所以得到交换图
$$
\xymatrix{
B \ar[rd] \ar@{..>}[r] & R' \ar@{..>}[d] \\
R \ar[r] \ar[u] & \kappa
}
$$
由于竖直箭头光滑，《代数详论》引理
\ref{more-algebra-lemma-lift-section-smooth-morphism} 蕴含存在平展环同态
$R\to R'$，它诱导同构 $R/\mathfrak m\to R'/\mathfrak mR'$，
并存在使上图交换的 $R$-代数同态 $B\to R'$。令 $b_i\in R'$
为 $a_i$ 经环同态 $A\to B\to R'$ 所得的像。由于这些全是
$R$-代数同态，$(b_1,\ldots,b_n)$ 是 $R'$ 中的解。
\end{proof}

\begin{example}
\label{smoothing-example-describe-henselian}
设 $(R,\mathfrak m)$ 为诺特局部环，其 Hensel 化为 $R^h$。
完备化之间的同态 $R^\wedge\to(R^h)^\wedge$ 是同构，见《代数详论》
引理 \ref{more-algebra-lemma-henselization-noetherian}。由于 $R^h$
也是诺特环（同上），可将 $R^h$ 看作其完备化的子环
（因为完备化忠实平坦）。由此，可把 $R^h$ 视为 $R^\wedge$ 的子环。

\medskip\noindent
来考察 $R^\wedge$ 中哪些元素属于 $R^h$。为简单起见，假设 $R$
为整环，其分式域为 $K$。显然，$R^h$ 的每个元素 $f$ 都在 $R$
上代数；这意味着存在形如
$a_nf^n+\ldots+a_1f+a_0=0$ 的方程，其中 $a_i\in R$、$n>0$ 且
$a_n \not = 0$。

\medskip\noindent
反之，假设 $f\in R^\wedge$、$n\in\mathbf{N}$ 且
$a_0,\ldots,a_n\in R$，其中 $a_n \not = 0$，并满足
$a_nf^n+\ldots+a_1f+a_0=0$。若 $R$ 是 G-环，则对每个 $N>0$，
都存在 $g\in R^h$，使
$a_ng^n+\ldots+a_1g+a_0=0$ 且
$f-g\in\mathfrak m^NR^\wedge$，见定理
\ref{smoothing-theorem-approximation-property-variant}。我们希望推出当
$N\gg0$ 时 $f=g$。若非如此，就会在 $R^h$ 中找到 $P(T)$ 的无穷多个
根 $g$。这是不可能的，因为 (1) $R^h\subset R^h\otimes_RK$，
且 (2) $R^h\otimes_RK$ 是 $K$ 的有限多个域扩张的乘积。事实上，
$R\to K$ 为单射而 $R\to R^h$ 平坦，故
$R^h\to R^h\otimes_RK$ 为单射；(2) 则由《代数详论》引理
\ref{more-algebra-lemma-fibres-henselization} 得出。

\medskip\noindent
结论：若 $R$ 是分式域为 $K$ 的诺特局部整环且为 G-环，
则 $R^h\subset R^\wedge$ 恰为所有在 $K$ 上代数的元素所成的集合。
\end{example}

\noindent
下面是本节主定理的另一个变体。

\begin{lemma}
\label{smoothing-lemma-approximation-property-variant}
设 $R$ 为诺特环，$\mathfrak p\subset R$ 为素理想，且
$f_1,\ldots,f_m\in R[x_1,\ldots,x_n]$。假设
$(a_1,\ldots,a_n)\in((R_\mathfrak p)^\wedge)^n$ 是一个解。
若 $R_\mathfrak p$ 是 G-环，则对每个整数 $N$，都存在
\begin{enumerate}
\item 平展环同态 $R\to R'$；
\item 位于 $\mathfrak p$ 之上的素理想 $\mathfrak p'\subset R'$；
\item $R'$ 中的解 $(b_1,\ldots,b_n)\in(R')^n$，
\end{enumerate}
使得 $\kappa(\mathfrak p)=\kappa(\mathfrak p')$ 且
$a_i-b_i\in(\mathfrak p')^N(R'_{\mathfrak p'})^\wedge$。
\end{lemma}

\begin{proof}
由定理 \ref{smoothing-theorem-approximation-property-variant}，可以在某个
$R_\mathfrak p$ 上平展的环 $R''$ 中找到解
$(b'_1,\ldots,b'_n)$；并有位于 $\mathfrak p$ 之上的素理想
$\mathfrak p''$，满足 $\kappa(\mathfrak p)=\kappa(\mathfrak p'')$ 且
$a_i-b'_i\in(\mathfrak p'')^N(R''_{\mathfrak p''})^\wedge$。
对某个平展 $R$-代数 $R'$，可以写成
$R''=R'\otimes_RR_\mathfrak p$（见《代数》引理
\ref{algebra-lemma-etale}）。必要时把 $R'$ 替换为一个主局部化，
便可假设 $(b'_1,\ldots,b'_n)$ 来自 $R'$ 中的解
$(b_1,\ldots,b_n)$。置 $\mathfrak p'=R'\cap\mathfrak p''$，则
$R''_{\mathfrak p''}=R'_{\mathfrak p'}$，证明完成。
\end{proof}



\section{Hensel 对的逼近}
\label{smoothing-section-approximation-pairs}

\noindent
可以把第 \ref{smoothing-section-approximation-G-rings} 节的讨论推广到 Hensel 对。
Hensel 对的定义见《代数详论》第
\ref{more-algebra-section-henselian-pairs} 节。

\begin{lemma}
\label{smoothing-lemma-henselian-pair}
\begin{slogan}
Hensel 对的逼近。
\end{slogan}
设 $(A,I)$ 为 Hensel 对，其中 $A$ 是诺特环。令 $A^\wedge$
为 $A$ 的 $I$-进完备化。假设下列条件至少有一个成立：
\begin{enumerate}
\item $A\to A^\wedge$ 是正则环同态；
\item $A$ 是诺特 G-环；
\item $(A,I)$ 是某个对 $(B,J)$ 的 Hensel 化
（《代数详论》，引理 \ref{more-algebra-lemma-henselization}），
其中 $B$ 是诺特 G-环。
\end{enumerate}
给定 $f_1,\ldots,f_m\in A[x_1,\ldots,x_n]$ 以及
$\hat{a}_1,\ldots,\hat{a}_n\in A^\wedge$，满足对
$j=1,\ldots,m$ 有 $f_j(\hat{a}_1,\ldots,\hat{a}_n)=0$。
则对每个 $N\geq1$，都存在 $a_1,\ldots,a_n\in A$，使
$\hat{a}_i-a_i\in I^N$，且对 $j=1,\ldots,m$ 有
$f_j(a_1,\ldots,a_n)=0$。
\end{lemma}

\begin{proof}
由《代数详论》引理
\ref{more-algebra-lemma-henselization-pair-G-ring}，(3) 蕴含 (2)；
由《代数详论》引理
\ref{more-algebra-lemma-map-G-ring-to-completion-regular}，(2) 蕴含 (1)。
因此，只需在 $A\to A^\wedge$ 为正则环同态的情形证明本引理。

\medskip\noindent
令 $\hat{a}_1,\ldots,\hat{a}_n$ 如引理陈述所述。由定理
\ref{smoothing-theorem-popescu}，可以找到分解 $A\to B\to A^\wedge$，
其中 $A\to P$ 光滑，并有 $b_1,\ldots,b_n\in B$ 满足
$f_j(b_1,\ldots,b_n)=0$ 于 $B$ 中。记 $\sigma$ 为复合
$B\to A^\wedge\to A/I^N$。由《代数详论》引理
\ref{more-algebra-lemma-lift-section-smooth-morphism}，可以找到平展环同态
$A\to A'$，它诱导同构 $A/I^N\to A'/I^NA'$，并有提升 $\sigma$ 的
$A$-代数同态 $\tilde\sigma:B\to A'$。由于 $(A,I)$ 是 Hensel 对，
存在 $A$-代数同态 $\chi:A'\to A$，见《代数详论》引理
\ref{more-algebra-lemma-characterize-henselian-pair}。
置 $a_i=\chi(\tilde\sigma(b_i))$，便得到一个解。
\end{proof}
% Generated mechanically from the frozen modules.tex target.
% Do not edit this generated file; edit the bound locale source instead.

% OK, start here.
%

\setcounter{chapter}{16}
\chapter{模层}



\phantomsection
\label{modules-section-phantom}


\section{引言}
\label{modules-section-introduction}

\noindent
本章阐述模层的基本概念。其中尤其包括阿贝尔群层的情形，因为它们可视为
$\underline{\mathbf{Z}}$-模层。基本参考文献为 \cite{FAC}、\cite{EGA}
和 \cite{SGA4}。

\medskip\noindent
关于环拓扑斯上的模层，我们在另一章中加以讨论（见《站点上的模》
第 \ref{sites-modules-section-introduction} 节），不过那里大部分内容只是重述本章的讨论。





\section{退化情形}
\label{modules-section-pathology}

\noindent
环空间是由拓扑空间 $X$ 与环层 $\mathcal{O}$ 组成的序对。定义允许
$\mathcal{O} = 0$。此时模范畴只有一个对象（即 $0$）。它仍是阿贝尔范畴等，
但略显退化。类似地，层 $\mathcal{O}$ 在 $X$ 的某些开子集上也可能为零，等等。

\medskip\noindent
考察局部环空间时不会出现这种情形（我们稍后将这样做）。








\section{模层的阿贝尔范畴}
\label{modules-section-kernels}

\noindent
设 $(X, \mathcal{O}_X)$ 为环空间，见《层》定义
\ref{sheaves-definition-ringed-space}。设 $\mathcal{F}$、$\mathcal{G}$
为 $\mathcal{O}_X$-模层，见《层》定义
\ref{sheaves-definition-sheaf-modules}。设
$\varphi, \psi : \mathcal{F} \to \mathcal{G}$ 为 $\mathcal{O}_X$-模层的态射。
定义 $\varphi + \psi : \mathcal{F} \to \mathcal{G}$：在每个开集
$U \subset X$ 上，它是由 $\varphi$、$\psi$ 诱导的两个映射之和。
显然，这仍是 $\mathcal{O}_X$-模层的态射。也显然，$\mathcal{O}_X$-模层态射的
复合关于这一加法是双线性的。因此 $\textit{Mod}(\mathcal{O}_X)$
是预加性范畴，见《同调》定义 \ref{homology-definition-preadditive}。

\medskip\noindent
记 $0$ 为在每个开集 $U \subset X$ 上恒取值 $\{0\}$ 的
$\mathcal{O}_X$-模层。显然，它既是 $\textit{Mod}(\mathcal{O}_X)$ 的终对象，
也是其始对象。给定 $\mathcal{O}_X$-模层态射
$\varphi : \mathcal{F} \to \mathcal{G}$，下列条件等价：
(a) $\varphi$ 为零；(b) $\varphi$ 分解通过 $0$；
(c) $\varphi$ 在每个开集 $U$ 的截面上为零；
(d) 对所有 $x \in X$，都有 $\varphi_x = 0$。见《层》引理
\ref{sheaves-lemma-points-exactness}。

\medskip\noindent
此外，给定两个 $\mathcal{O}_X$-模层 $\mathcal{F}$、$\mathcal{G}$，
可将直和定义为
$$
\mathcal{F} \oplus \mathcal{G} = \mathcal{F} \times \mathcal{G}
$$
并配以《同调》定义 \ref{homology-definition-direct-sum} 中显然的映射
$(i, j, p, q)$。因此 $\textit{Mod}(\mathcal{O}_X)$ 是加性范畴，见
《同调》定义 \ref{homology-definition-additive-category}。

\medskip\noindent
设 $\varphi : \mathcal{F} \to \mathcal{G}$ 为 $\mathcal{O}_X$-模层态射。
可将 $\Ker(\varphi)$ 定义为 $\mathcal{F}$ 的子层，其截面为
$$
\Ker(\varphi)(U) =
\{ s \in \mathcal{F}(U) \mid \varphi(s) = 0 \text{ 于 } \mathcal{G}(U)\}
$$
其中 $U$ 遍历 $X$ 的所有开集。容易看出，这确实是 $\mathcal{O}_X$-模层范畴中的核。
换言之，态射 $\alpha : \mathcal{H} \to \mathcal{F}$ 分解通过
$\Ker(\varphi)$，当且仅当 $\varphi \circ \alpha = 0$。此外，在茎的层次上有
$\Ker(\varphi)_x = \Ker(\varphi_x)$。

\medskip\noindent
另一方面，将 $\Coker(\varphi)$ 定义为下述 $\mathcal{O}_X$-模预层所关联的
$\mathcal{O}_X$-模层：
$$
U
\longmapsto
\Coker(\mathcal{F}(U)\to \mathcal{G}(U)) =
\mathcal{G}(U)/\varphi(\mathcal{F}(U)).
$$
由于取茎与层化可交换，见《层》引理
\ref{sheaves-lemma-stalk-sheafification}，可得
$\Coker(\varphi)_x = \Coker(\varphi_x)$。因此映射
$\mathcal{G} \to \Coker(\varphi)$ 是满射（作为集合层的态射），见《层》第
\ref{sheaves-section-exactness-points} 节。为证明这是余核，注意若
$\beta : \mathcal{G} \to \mathcal{H}$ 是满足
$\beta \circ \varphi = 0$ 的 $\mathcal{O}_X$-模层态射，则对每个开集
$U \subset X$，$\beta$ 都诱导一个从
$\mathcal{G}(U)/\varphi(\mathcal{F}(U))$ 到 $\mathcal{H}(U)$ 的映射。
由层化的泛性质（见《层》引理
\ref{sheaves-lemma-sheafification-presheaf-modules}），得到典范映射
$\Coker(\varphi) \to \mathcal{H}$，使得原来的 $\beta$ 等于复合
$\mathcal{G} \to \Coker(\varphi) \to \mathcal{H}$。由上述满射性，
态射 $\Coker(\varphi) \to \mathcal{H}$ 是唯一的。

\begin{lemma}
\label{modules-lemma-abelian}
设 $(X, \mathcal{O}_X)$ 为环空间。范畴 $\textit{Mod}(\mathcal{O}_X)$
是阿贝尔范畴。此外，一个复形
$$
\mathcal{F} \to \mathcal{G} \to \mathcal{H}
$$
在 $\mathcal{G}$ 处正合，当且仅当对所有 $x \in X$，复形
$$
\mathcal{F}_x \to \mathcal{G}_x \to \mathcal{H}_x
$$
在 $\mathcal{G}_x$ 处正合。
\end{lemma}

\begin{proof}
依《同调》定义 \ref{homology-definition-abelian-category}，只需证明像与余像一致。
依《层》引理 \ref{sheaves-lemma-points-exactness}，只需证明像与余像在每个
$x \in X$ 处的茎相同。由上述核与余核的构造，这些茎正是
$\mathcal{O}_{X, x}$-模范畴中的余像与像。结论遂由环上模的范畴是阿贝尔范畴这一事实得到。
\end{proof}

\noindent
事实上，范畴 $\textit{Mod}(\mathcal{O}_X)$ 还有许多其他性质。以下给出两种构造。
\begin{enumerate}
\item 给定任意集合 $I$，并对每个 $i \in I$ 给定一个 $\mathcal{O}_X$-模层，
可作乘积
$$
\prod\nolimits_{i \in I} \mathcal{F}_i
$$
它是在每个开集 $U$ 上取诸模 $\mathcal{F}_i(U)$ 之积的层。这也就是范畴积，
如《范畴论》定义 \ref{categories-definition-product} 所述。
\item 给定任意集合 $I$，并对每个 $i \in I$ 给定一个 $\mathcal{O}_X$-模层，
可作直和
$$
\bigoplus\nolimits_{i \in I} \mathcal{F}_i
$$
它是在每个开集 $U$ 上取诸模 $\mathcal{F}_i(U)$ 之直和的预层的{\it 层化}。
这也就是范畴余积，如《范畴论》定义
\ref{categories-definition-coproduct} 所述；由层化的泛性质即可看出这一点。
\end{enumerate}
由这些构造可知，$\textit{Mod}(\mathcal{O}_X)$ 中所有极限与余极限都存在。

\begin{lemma}
\label{modules-lemma-limits-colimits}
设 $(X, \mathcal{O}_X)$ 为环空间。
\begin{enumerate}
\item $\textit{Mod}(\mathcal{O}_X)$ 中所有极限都存在。极限与
$\mathcal{O}_X$-模预层的相应极限相同（即极限与在开集上取截面可交换）。
\item $\textit{Mod}(\mathcal{O}_X)$ 中所有余极限都存在。余极限是预层范畴中
相应余极限的层化。取余极限与取茎可交换。
\item 滤过余极限是正合的。
\item 有限直和与 $\mathcal{O}_X$-模预层的相应有限直和相同。
\end{enumerate}
\end{lemma}

\begin{proof}
由于 $\textit{Mod}(\mathcal{O}_X)$ 是阿贝尔范畴（引理 \ref{modules-lemma-abelian}），
它具有所有有限极限与有限余极限（《同调》引理
\ref{homology-lemma-colimit-abelian-category}）。因此，由积和余积的存在及其上述描述，
再结合《范畴论》引理 \ref{categories-lemma-limits-products-equalizers} 与
\ref{categories-lemma-colimits-coproducts-coequalizers}，便得到极限与余极限的存在性及其描述。
由于层化与取茎可交换，可见取余极限与取茎可交换。第 (3) 项的含义如下：
给定有向集 $I$ 上一族 $\mathcal{O}_X$-模层的正合列
$0 \to \mathcal{F}_i \to \mathcal{G}_i \to \mathcal{H}_i \to 0$，则序列
$0 \to \colim \mathcal{F}_i \to \colim \mathcal{G}_i \to
\colim \mathcal{H}_i \to 0$ 也正合。由于可在茎上检验正合性（引理
\ref{modules-lemma-abelian}），这由模的相应情形，即《代数》引理
\ref{algebra-lemma-directed-colimit-exact} 得出。第 (4) 项的证明从略。
\end{proof}

\noindent
极限与余极限的存在，使我们可以像《范畴论》第
\ref{categories-section-exact-functor} 节那样，以极限和余极限来考察定义在
$\mathcal{O}$-模层范畴上的函子的正合性。关于用短正合列描述正合性，见《同调》引理
\ref{homology-lemma-exact-functor}。

\begin{lemma}
\label{modules-lemma-exactness-pushforward-pullback}
设 $f : (X, \mathcal{O}_X) \to (Y, \mathcal{O}_Y)$ 为环空间态射。
\begin{enumerate}
\item 函子
$f_* : \textit{Mod}(\mathcal{O}_X) \to \textit{Mod}(\mathcal{O}_Y)$
是左正合的。事实上，它与所有极限可交换。
\item 函子
$f^* : \textit{Mod}(\mathcal{O}_Y) \to \textit{Mod}(\mathcal{O}_X)$
是右正合的。事实上，它与所有余极限可交换。
\item 阿贝尔群层上的逆像函子
$f^{-1} : \textit{Ab}(Y) \to \textit{Ab}(X)$ 是正合的。
\end{enumerate}
\end{lemma}

\begin{proof}
第 (1)、(2) 项成立，是因为 $(f^*, f_*)$ 是一对伴随函子，见《层》引理
\ref{sheaves-lemma-adjoint-pullback-pushforward-modules} 和《范畴论》第
\ref{categories-section-adjoint} 节。第 (3) 项由可在茎上检验正合性（引理
\ref{modules-lemma-abelian}）以及逆像之茎的描述得出，见《层》引理
\ref{sheaves-lemma-pullback-abelian-stalk}。
\end{proof}

\begin{lemma}
\label{modules-lemma-j-shriek-exact}
设 $j : U \to X$ 为拓扑空间的开浸入。函子
$j_! : \textit{Ab}(U) \to \textit{Ab}(X)$ 是正合的。
\end{lemma}

\begin{proof}
这由《层》引理 \ref{sheaves-lemma-j-shriek-abelian} 给出的茎的描述得出。
\end{proof}

\begin{lemma}
\label{modules-lemma-section-direct-sum-quasi-compact}
设 $(X, \mathcal{O}_X)$ 为环空间，$I$ 为集合。对 $i \in I$，设
$\mathcal{F}_i$ 为 $\mathcal{O}_X$-模层。若 $U \subset X$ 为拟紧开集，则映射
$$
\bigoplus\nolimits_{i \in I} \mathcal{F}_i(U)
\longrightarrow
\left(\bigoplus\nolimits_{i \in I} \mathcal{F}_i\right)(U)
$$
是双射。
\end{lemma}

\begin{proof}
若 $s$ 是右端的一个元素，则存在开覆盖
$U = \bigcup_{j \in J} U_j$，使得 $s|_{U_j}$ 是有限和
$\sum_{i \in I_j} s_{ji}$，其中 $s_{ji} \in \mathcal{F}_i(U_j)$。
由于 $U$ 拟紧，可设此覆盖为有限覆盖，即 $J$ 有限。于是
$I' = \bigcup_{j \in J} I_j$ 是 $I$ 的有限子集。显然，$s$ 是子层
$\bigoplus_{i \in I'} \mathcal{F}_i$ 的一个截面。有限直和无需层化，故结论由上面的
引理 \ref{modules-lemma-limits-colimits} 得出。
\end{proof}



\section{模层的截面}
\label{modules-section-sections}

\noindent
设 $(X, \mathcal{O}_X)$ 为环空间，$\mathcal{F}$ 为 $\mathcal{O}_X$-模层，
并设 $s \in \Gamma(X, \mathcal{F}) = \mathcal{F}(X)$ 为整体截面。
存在唯一的{\it $\mathcal{O}_X$-模层态射}
$$
\mathcal{O}_X \longrightarrow \mathcal{F}, \ f \longmapsto fs
$$
称为{\it 与 $s$ 关联的态射}。上式表示：$\mathcal{O}_X$ 的一个局部截面 $f$，
即某个开集 $U$ 上的截面 $f$，被映到 $f$ 与 $s$ 在 $U$ 上的限制之乘积。
反过来，任意态射 $\varphi : \mathcal{O}_X \to \mathcal{F}$ 都给出截面
$s = \varphi(1)$，并且 $\varphi$ 正是与 $s$ 关联的态射。

\begin{definition}
\label{modules-definition-globally-generated}
设 $(X, \mathcal{O}_X)$ 为环空间，$\mathcal{F}$ 为 $\mathcal{O}_X$-模层。
若存在集合 $I$ 以及整体截面
$s_i \in \Gamma(X, \mathcal{F})$（$i \in I$），使得映射
$$
\bigoplus\nolimits_{i \in I}
\mathcal{O}_X \longrightarrow \mathcal{F}
$$
为满射，其中它在与 $i$ 对应的直和项上就是与 $s_i$ 关联的态射，
则称 $\mathcal{F}$ {\it 由整体截面生成}。在这种情形下，称诸截面 $s_i$
{\it 生成} $\mathcal{F}$。
\end{definition}

\noindent
我们经常沿用《层》第 \ref{sheaves-section-stalks} 节引入的符号滥用：
若 $s$ 是定义在点 $x \in X$ 的某个开邻域上的 $\mathcal{F}$ 的局部截面，
则以 $s_x$，甚至仍以 $s$，表示 $s$ 在茎 $\mathcal{F}_x$ 中的像。

\begin{lemma}
\label{modules-lemma-globally-generated}
设 $(X, \mathcal{O}_X)$ 为环空间，$\mathcal{F}$ 为 $\mathcal{O}_X$-模层，
$I$ 为集合，并设 $s_i \in \Gamma(X, \mathcal{F})$（$i \in I$）为整体截面。
诸截面 $s_i$ 生成 $\mathcal{F}$，当且仅当对所有 $x\in X$，诸元素
$s_{i, x} \in \mathcal{F}_x$ 生成 $\mathcal{O}_{X, x}$-模
$\mathcal{F}_x$。
\end{lemma}

\begin{proof}
略。
\end{proof}

\begin{lemma}
\label{modules-lemma-tensor-product-globally-generated}
\begin{slogan}
整体生成的模层之张量积仍是整体生成的。
\end{slogan}
设 $(X, \mathcal{O}_X)$ 为环空间，$\mathcal{F}$、$\mathcal{G}$ 为
$\mathcal{O}_X$-模层。若 $\mathcal{F}$ 与 $\mathcal{G}$ 由整体截面生成，
则 $\mathcal{F} \otimes_{\mathcal{O}_X} \mathcal{G}$ 亦然。
\end{lemma}

\begin{proof}
略。
\end{proof}

\begin{lemma}
\label{modules-lemma-generated-by-local-sections}
设 $(X, \mathcal{O}_X)$ 为环空间，$\mathcal{F}$ 为 $\mathcal{O}_X$-模层，
$I$ 为集合。设 $s_i$（$i \in I$）为 $\mathcal{F}$ 的一族局部截面，
即对某些开集 $U_i \subset X$ 有 $s_i \in \mathcal{F}(U_i)$。
存在唯一的最小 $\mathcal{O}_X$-模子层 $\mathcal{G}$，使得每个 $s_i$
都对应于 $\mathcal{G}$ 的一个局部截面。
\end{lemma}

\begin{proof}
考虑由下述规则定义的 $\mathcal{O}_X$-模子预层：
$$
U
\longmapsto
\{
\text{形如 } \sum\nolimits_{i \in J} f_i (s_i|_U) \text{，其中 }
J \text{ 有限，} U \subset U_i \text{ 若 } i\in J \text{，且 }
f_i \in \mathcal{O}_X(U)
\}
$$
令 $\mathcal{G}$ 为此子预层的层化。依《层》引理
\ref{sheaves-lemma-sheafify-epi-mono}，它是 $\mathcal{F}$ 的子层。
由于所有上述有限和显然都必须属于 $\mathcal{G}$，这正是所求的最小子层。
\end{proof}

\begin{definition}
\label{modules-definition-generated-by-local-sections}
设 $(X, \mathcal{O}_X)$ 为环空间，$\mathcal{F}$ 为 $\mathcal{O}_X$-模层。
给定集合 $I$ 以及 $\mathcal{F}$ 的局部截面 $s_i$（$i \in I$），称上述
引理 \ref{modules-lemma-generated-by-local-sections} 中的子层 $\mathcal{G}$ 为
{\it 由诸 $s_i$ 生成的子层}。
\end{definition}

\begin{lemma}
\label{modules-lemma-generated-by-local-sections-stalk}
设 $(X, \mathcal{O}_X)$ 为环空间，$\mathcal{F}$ 为 $\mathcal{O}_X$-模层。
给定集合 $I$ 以及 $\mathcal{F}$ 的局部截面 $s_i$（$i \in I$）。
设 $\mathcal{G}$ 为由诸 $s_i$ 生成的子层，并设 $x\in X$。则
$\mathcal{G}_x$ 是 $\mathcal{F}_x$ 中由所有 $s_{i, x}$ 生成的
$\mathcal{O}_{X, x}$-子模，其中 $i$ 遍历使 $s_i$ 在 $x$ 处有定义的指标。
\end{lemma}

\begin{proof}
这由引理 \ref{modules-lemma-generated-by-local-sections} 的证明中 $\mathcal{G}$ 的构造立即可见。
\end{proof}










\section{模层与截面的支集}
\label{modules-section-support}

\begin{definition}
\label{modules-definition-support}
设 $(X, \mathcal{O}_X)$ 为环空间，$\mathcal{F}$ 为 $\mathcal{O}_X$-模层。
\begin{enumerate}
\item $\mathcal{F}$ 的{\it 支集}是所有满足 $\mathcal{F}_x \not = 0$
的点 $x \in X$ 组成的集合。
\item 以 $\text{Supp}(\mathcal{F})$ 表示 $\mathcal{F}$ 的支集。
\item 设 $s \in \Gamma(X, \mathcal{F})$ 为整体截面。$s$ 的{\it 支集}
是所有使 $s$ 的像 $s_x \in \mathcal{F}_x$ 非零的点 $x \in X$ 组成的集合。
\end{enumerate}
\end{definition}

\noindent
由此当然也定义了局部截面的支集，因为局部截面就是 $\mathcal{F}$ 的某个限制的整体截面。

\begin{lemma}
\label{modules-lemma-support-section-closed}
设 $(X, \mathcal{O}_X)$ 为环空间，$\mathcal{F}$ 为 $\mathcal{O}_X$-模层，
$U \subset X$ 为开集。
\begin{enumerate}
\item $s \in \mathcal{F}(U)$ 的支集在 $U$ 中是闭集。
\item $fs$ 的支集包含于 $f \in \mathcal{O}_X(U)$ 与
$s \in \mathcal{F}(U)$ 的支集之交。
\item $s + s'$ 的支集包含于 $s, s' \in \mathcal{F}(U)$ 的支集之并。
\item $\mathcal{F}$ 的支集是 $\mathcal{F}$ 的所有局部截面的支集之并。
\item 若 $\varphi : \mathcal{F} \to \mathcal{G}$ 为 $\mathcal{O}_X$-模层态射，
则 $\varphi(s)$ 的支集包含于 $s \in \mathcal{F}(U)$ 的支集。
\end{enumerate}
\end{lemma}

\begin{proof}
这是因为，若 $s_x = 0$，则由茎的定义，$s$ 在 $x$ 的某个开邻域上为零。
$f$ 的情形同理。细节从略。
\end{proof}

\noindent
一般而言，模层的支集不一定闭。例如，可在带通常阿基米德拓扑的
$\mathbf{R}$ 上取一个阿贝尔群层，它是无穷多个非零摩天楼层的直和，
其中每个摩天楼层都支撑在 $\mathbf{R}$ 的单点 $p_i$ 上。此时支集就是点集
$\{p_i\}$，而这个集合未必闭。

\medskip\noindent
另一个例子是开浸入 $j : U = (0 , \infty) \to \mathbf{R} = X$ 及阿贝尔群层
$j_!\underline{\mathbf{Z}}_U$。由《层》第
\ref{sheaves-section-open-immersions} 节，此层的支集恰为 $U$。

\begin{lemma}
\label{modules-lemma-support-sheaf-rings-closed}
设 $X$ 为拓扑空间。环层的支集是闭集。
\end{lemma}

\begin{proof}
这是因为，依我们的约定，一个环为 $0$ 当且仅当 $1 = 0$；因此环层的支集
就是其单位截面的支集。
\end{proof}






\section{闭浸入与阿贝尔群层}
\label{modules-section-closed-immersions}

\noindent
回忆，我们将拓扑空间 $X$ 上的阿贝尔群层视为
$\underline{\mathbf{Z}}_X$-模层。因此，关于模层的任何结果和定义都可用于阿贝尔群层。

\begin{lemma}
\label{modules-lemma-i-star-exact}
设 $X$ 为拓扑空间，$Z \subset X$ 为闭子集。以 $i : Z \to X$ 表示包含映射。函子
$$
i_* : \textit{Ab}(Z) \longrightarrow \textit{Ab}(X)
$$
是正合且全忠实的，其本质像恰由支集包含于 $Z$ 的阿贝尔群层组成。
函子 $i^{-1}$ 是 $i_*$ 的左逆。
\end{lemma}

\begin{proof}
正合性由《层》引理 \ref{sheaves-lemma-stalks-closed-pushforward}
中的茎的描述及引理 \ref{modules-lemma-abelian} 得出。其余断言已在《层》引理
\ref{sheaves-lemma-equivalence-categories-closed-abelian} 中证明。
\end{proof}

\noindent
设 $\mathcal{F}$ 为拓扑空间 $X$ 上的阿贝尔群层。给定闭子集 $Z$，
$\mathcal{F}$ 有一个典范阿贝尔子层，它恰由支集包含于 $Z$ 的截面组成。
准确陈述如下。

\begin{remark}
\label{modules-remark-sections-support-in-closed}
设 $X$ 为拓扑空间，$Z \subset X$ 为闭子集，$\mathcal{F}$ 为 $X$ 上的阿贝尔群层。
对开集 $U \subset X$，定义
$$
\mathcal{H}_Z(\mathcal{F})(U) =
\{s \in \mathcal{F}(U) \mid
\text{截面 }s\text{ 的支集包含于 }Z \cap U\}.
$$
则 $\mathcal{H}_Z(\mathcal{F})$ 是 $\mathcal{F}$ 的阿贝尔子层，
并且是 $\mathcal{F}$ 中支集包含于 $Z$ 的最大阿贝尔子层。由引理
\ref{modules-lemma-i-star-exact}，可将（并且我们确实将）
$\mathcal{H}_Z(\mathcal{F})$ 视为 $Z$ 上的阿贝尔群层。如此得到左正合函子
$$
\textit{Ab}(X) \longrightarrow \textit{Ab}(Z),\quad
\mathcal{F} \longmapsto \mathcal{H}_Z(\mathcal{F})
\text{视为 }Z\text{ 上的阿贝尔群层}
$$
上述所有断言均直接由引理 \ref{modules-lemma-support-section-closed} 得出。
\end{remark}

\noindent
这里正好可以证明，函子 $i_*$ 在阿贝尔群层上有右伴随。

\begin{lemma}
\label{modules-lemma-i-star-right-adjoint}
设 $i : Z \to X$ 为闭子集到拓扑空间 $X$ 的包含映射。注
\ref{modules-remark-sections-support-in-closed} 中的函子
$\textit{Ab}(X) \to \textit{Ab}(Z)$，
$\mathcal{F} \mapsto \mathcal{H}_Z(\mathcal{F})$，是
$i_* : \textit{Ab}(Z) \to \textit{Ab}(X)$ 的右伴随。特别地，
$i_*$ 与任意余极限可交换。
\end{lemma}

\begin{proof}
只需证明，对 $X$ 上任意阿贝尔群层 $\mathcal{F}$ 以及 $Z$ 上任意阿贝尔群层
$\mathcal{G}$，都有
$$
\Hom_{\textit{Ab}(X)}(i_*\mathcal{G}, \mathcal{F}) =
\Hom_{\textit{Ab}(Z)}(\mathcal{G}, \mathcal{H}_Z(\mathcal{F}))
$$
这是显然的，因为 $i_*\mathcal{G}$ 的任意截面的支集终究都包含于 $Z$。
细节从略。
\end{proof}

\begin{remark}
\label{modules-remark-i-star-right-adjoint}
在《层》注 \ref{sheaves-remark-i-star-not-exact} 中，我们证明了：
$i_*$ 作为集合层范畴上的函子没有右伴随，原因很简单，它不正合。
不过，与此非常接近的说法确实成立。事实上，函子 $i_*$ 在带基点集合层上正合；
对带基点集合层可以定义支集包含于 $Z$ 的截面，而 $\mathcal{H}_Z$ 也有意义，
并且是 $i_*$ 的右伴随。
\end{remark}









\section{一个典范正合列}
\label{modules-section-canonical-exact-sequence}

\noindent
我们为这个正合列单列一节。

\begin{lemma}
\label{modules-lemma-canonical-exact-sequence}
设 $X$ 为拓扑空间，$U \subset X$ 为开子集，其补集为 $Z \subset X$。
以 $j : U \to X$ 表示开浸入，以 $i : Z \to X$ 表示闭浸入。
对 $X$ 上任意阿贝尔群层 $\mathcal{F}$，伴随态射
$j_{!}j^{-1}\mathcal{F} \to \mathcal{F}$ 与
$\mathcal{F} \to i_*i^{-1}\mathcal{F}$ 给出阿贝尔群层的短正合列
$$
0 \to j_{!}j^{-1}\mathcal{F} \to \mathcal{F} \to i_*i^{-1}\mathcal{F} \to 0
$$
。对 $X$ 上阿贝尔群层的任意态射
$\varphi : \mathcal{F} \to \mathcal{G}$，得到如下短正合列之间的态射：
$$
\xymatrix{
0 \ar[r] &
j_{!}j^{-1}\mathcal{F} \ar[r] \ar[d] &
\mathcal{F} \ar[r] \ar[d] &
i_*i^{-1}\mathcal{F} \ar[r] \ar[d] &
0 \\
0 \ar[r] &
j_{!}j^{-1}\mathcal{G} \ar[r] &
\mathcal{G} \ar[r] &
i_*i^{-1}\mathcal{G} \ar[r] &
0
}
$$
\end{lemma}

\begin{proof}
短正合列的函子性立即来自伴随态射的自然性。可在茎上检验正合性（引理
\ref{modules-lemma-abelian}）。关于这里所涉及的茎的描述，见《层》引理
\ref{sheaves-lemma-j-shriek-abelian} 与
\ref{sheaves-lemma-stalks-closed-pushforward}。
\end{proof}








\section{局部由截面生成的模层}
\label{modules-section-locally-generated}

\noindent
设 $(X, \mathcal{O}_X)$ 为环空间。在本节及下一节中，我们常把层限制到开子空间
$U \subset X$，见《层》第 \ref{sheaves-section-open-immersions} 节。
特别地，我们常以 $(U, \mathcal{O}_U)$ 表示这个开子空间，而不用更准确的记号
$(U, \mathcal{O}_X|_U)$，见《层》定义 \ref{sheaves-definition-restriction}。

\medskip\noindent
考虑开浸入 $j : U = (0 , \infty) \to \mathbf{R} = X$ 以及阿贝尔群层
$j_!\underline{\mathbf{Z}}_U$。由《层》第
\ref{sheaves-section-open-immersions} 节，$j_!\underline{\mathbf{Z}}_U$
在 $x = 0$ 处的茎为 $0$。事实上，此层在任意包含 $0$ 的开区间上的截面都为
$0$。因此，不存在点 $0$ 的开邻域，使得此层在其上可由截面生成。

\begin{definition}
\label{modules-definition-locally-generated}
设 $(X, \mathcal{O}_X)$ 为环空间，$\mathcal{F}$ 为 $\mathcal{O}_X$-模层。
若对每个 $x \in X$ 都存在 $x$ 的开邻域 $U$，使得 $\mathcal{F}|_U$
作为 $\mathcal{O}_U$-模层是整体生成的，则称 $\mathcal{F}$
{\it 局部由截面生成}。
\end{definition}

\noindent
换言之，存在集合 $I$，且对每个 $i$ 存在截面
$s_i \in \mathcal{F}(U)$，使得相应态射
$$
\bigoplus\nolimits_{i \in I} \mathcal{O}_U
\longrightarrow
\mathcal{F}|_U
$$
为满射。

\begin{lemma}
\label{modules-lemma-pullback-locally-generated}
设 $f : (X, \mathcal{O}_X) \to (Y, \mathcal{O}_Y)$ 为环空间态射。
若 $\mathcal{G}$ 局部由截面生成，则其拉回 $f^*\mathcal{G}$ 亦然。
\end{lemma}

\begin{proof}
给定 $Y$ 的开子空间 $V$，可考虑如下环空间交换图：
$$
\xymatrix{
(f^{-1}V, \mathcal{O}_{f^{-1}V}) \ar[r]_{j'} \ar[d]_{f'} &
(X, \mathcal{O}_X) \ar[d]^f \\
(V, \mathcal{O}_V) \ar[r]^j &
(Y, \mathcal{O}_Y)
}
$$
我们知道 $f^*\mathcal{G}|_{f^{-1}V} \cong (f')^*(\mathcal{G}|_V)$，
见《层》引理 \ref{sheaves-lemma-push-pull-composition-modules}。
因此可设 $\mathcal{G}$ 是整体生成的。

\medskip\noindent
由引理 \ref{modules-lemma-exactness-pushforward-pullback}，我们已经知道 $f^*$
与所有余极限可交换，并且是右正合的。因此，若有满射
$$
\bigoplus\nolimits_{i \in I}
\mathcal{O}_Y
\to
\mathcal{G}
\to
0
$$
，则施用 $f^*$ 后得到满射
$$
\bigoplus\nolimits_{i \in I}
\mathcal{O}_X
\to
f^*\mathcal{G}
\to
0.
$$
这就证明了引理。
\end{proof}












\section{有限型模层}
\label{modules-section-finite-type}

\begin{definition}
\label{modules-definition-finite-type}
设 $(X, \mathcal{O}_X)$ 为环空间，$\mathcal{F}$ 为 $\mathcal{O}_X$-模层。
若对每个 $x \in X$ 都存在开邻域 $U$，使得 $\mathcal{F}|_U$
由有限多个截面生成，则称 $\mathcal{F}$ 是{\it 有限型的}。
\end{definition}

\begin{lemma}
\label{modules-lemma-pullback-finite-type}
设 $f : (X, \mathcal{O}_X) \to (Y, \mathcal{O}_Y)$ 为环空间态射。
有限型 $\mathcal{O}_Y$-模层的拉回 $f^*\mathcal{G}$ 是有限型
$\mathcal{O}_X$-模层。
\end{lemma}

\begin{proof}
按照引理 \ref{modules-lemma-pullback-locally-generated} 的证明，可设 $\mathcal{G}$
由有限多个整体截面生成。由引理
\ref{modules-lemma-exactness-pushforward-pullback}，我们已经知道 $f^*$
与所有余极限可交换，并且是右正合的。因此，若有满射
$$
\bigoplus\nolimits_{i = 1, \ldots, n}
\mathcal{O}_Y
\to
\mathcal{G}
\to
0
$$
，则施用 $f^*$ 后得到满射
$$
\bigoplus\nolimits_{i = 1, \ldots, n}
\mathcal{O}_X
\to
f^*\mathcal{G}
\to
0.
$$
这就证明了引理。
\end{proof}

\begin{lemma}
\label{modules-lemma-extension-finite-type}
设 $X$ 为环空间。有限型 $\mathcal{O}_X$-模层之间态射的像是有限型的。设
$0 \to \mathcal{F}_1 \to \mathcal{F}_2 \to \mathcal{F}_3 \to 0$
为 $\mathcal{O}_X$-模层的短正合列。若 $\mathcal{F}_1$ 与
$\mathcal{F}_3$ 是有限型的，则 $\mathcal{F}_2$ 也是有限型的。
\end{lemma}

\begin{proof}
关于像的断言是平凡的。关于短正合列的断言来自以下事实：
$\mathcal{F}_3$ 的截面局部地可提升为 $\mathcal{F}_2$ 的截面，
并且环上模范畴中有相应结论（例如将其用于各个茎）。
\end{proof}

\begin{lemma}
\label{modules-lemma-finite-type-surjective-on-stalk}
设 $X$ 为环空间，$\varphi : \mathcal{G} \to \mathcal{F}$ 为
$\mathcal{O}_X$-模层同态，且 $x \in X$。假设 $\mathcal{F}$ 是有限型的，
并且茎上的映射 $\varphi_x : \mathcal{G}_x \to \mathcal{F}_x$ 为满射。
则存在开邻域 $x \in U \subset X$，使得 $\varphi|_U$ 为满射。
\end{lemma}

\begin{proof}
取 $x$ 的开邻域 $U \subset X$，使得 $\mathcal{F}$ 在 $U$ 上由
$s_1, \ldots, s_n \in \mathcal{F}(U)$ 生成。由 $\varphi_x$ 的满射性，
缩小 $U$ 后可设 $s_i = \varphi(t_i)$，其中
$t_i \in \mathcal{G}(U)$。于是这个 $U$ 满足要求。
\end{proof}

\begin{lemma}
\label{modules-lemma-finite-type-stalk-zero}
设 $X$ 为环空间，$\mathcal{F}$ 为 $\mathcal{O}_X$-模层，且 $x \in X$。
假设 $\mathcal{F}$ 是有限型的并且 $\mathcal{F}_x = 0$。则存在开邻域
$x \in U \subset X$，使得 $\mathcal{F}|_U$ 为零。
\end{lemma}

\begin{proof}
这是把引理 \ref{modules-lemma-finite-type-surjective-on-stalk}
用于态射 $0 \to \mathcal{F}$ 所得的特殊情形。
\end{proof}

\begin{lemma}
\label{modules-lemma-support-finite-type-closed}
\begin{slogan}
在任意环空间上，有限型模层都具有闭支集。
\end{slogan}
设 $(X, \mathcal{O}_X)$ 为环空间，$\mathcal{F}$ 为 $\mathcal{O}_X$-模层。
若 $\mathcal{F}$ 是有限型的，则 $\mathcal{F}$ 的支集是闭集。
\end{lemma}

\begin{proof}
这只是引理 \ref{modules-lemma-finite-type-stalk-zero} 的改述。
\end{proof}

\begin{lemma}
\label{modules-lemma-finite-type-quasi-compact-colimit}
设 $X$ 为环空间，$I$ 为预序集，并设 $(\mathcal{F}_i, f_{ii'})$
为 $I$ 上由 $\mathcal{O}_X$-模层组成的系统（见《范畴论》第
\ref{categories-section-posets-limits} 节）。令
$\mathcal{F} = \colim \mathcal{F}_i$ 为其余极限。假设
(a) $I$ 有向；(b) $\mathcal{F}$ 是有限型 $\mathcal{O}_X$-模层；
(c) $X$ 拟紧。则存在 $i$，使得 $\mathcal{F}_i \to \mathcal{F}$ 为满射。
若转移映射 $f_{ii'}$ 都是单射，则可得对某个 $i \in I$ 有
$\mathcal{F} = \mathcal{F}_i$。
\end{lemma}

\begin{proof}
设 $x \in X$。存在 $x$ 的开邻域 $U \subset X$ 及有限多个截面
$s_j \in \mathcal{F}(U)$（$j = 1, \ldots, m$），使得
$s_1, \ldots, s_m$ 作为 $\mathcal{O}_U$-模层生成 $\mathcal{F}$。
必要时把 $U$ 缩小为 $x$ 的更小开邻域，可设每个 $s_j$ 都来自某个
$i \in I$ 所对应的 $\mathcal{F}_i$ 的截面。于是，由于 $X$ 拟紧，
可找到有限开覆盖 $X = \bigcup_{j = 1, \ldots, m} U_j$；对每个 $j$，
可找到指标 $i_j$ 及有限多个截面
$s_{jl} \in \mathcal{F}_{i_j}(U_j)$，使它们的像生成 $\mathcal{F}$
在 $U_j$ 上的限制。显然，任取不小于所有 $i_j$ 的指标 $i \in I$，
引理的结论都成立。
\end{proof}

\begin{lemma}
\label{modules-lemma-set-isomorphism-classes-finite-type-modules}
设 $X$ 为环空间。存在一组有限型 $\mathcal{O}_X$-模层
$\{\mathcal{F}_i\}_{i \in I}$，使得 $X$ 上每个有限型
$\mathcal{O}_X$-模层都恰与其中一个 $\mathcal{F}_i$ 同构。
\end{lemma}

\begin{proof}
对每个开覆盖 $\mathcal{U} : X = \bigcup U_j$，考虑这样的
$\mathcal{O}_X$-模层 $\mathcal{F}$：每个限制 $\mathcal{F}|_{U_j}$
都是某个 $r_j \geq 0$ 所对应的 $\mathcal{O}_{U_j}^{\oplus r_j}$ 的商。
这些层由子层 $\mathcal{K}_j \subset \mathcal{O}_{U_j}^{\oplus r_j}$
及粘合数据参数化：
$$
\varphi_{jj'} :
\mathcal{O}_{U_j \cap U_{j'}}^{\oplus r_j}/
(\mathcal{K}_j|_{U_j \cap U_{j'}})
\longrightarrow
\mathcal{O}_{U_j \cap U_{j'}}^{\oplus r_{j'}}/
(\mathcal{K}_{j'}|_{U_j \cap U_{j'}})
$$
见《层》第 \ref{sheaves-section-glueing-sheaves} 节。注意，所有粘合数据组成一个集合。
所有满足映射 $J \to \mathcal{P}(X)$，$j \mapsto U_j$ 为单射的覆盖
$\mathcal{U} : X = \bigcup_{j \in J} U_i$ 也组成一个集合。因此，
由上述商粘合而成的所有 $\mathcal{O}_X$-模层组成一个集合 $\mathcal{I}$。
按定义，每个有限型 $\mathcal{O}_X$-模层都与 $\mathcal{I}$ 中的某个元素同构。
从 $\mathcal{I}$ 内每个同构类中选取一个元素，便得到所求的层集（这里使用选择公理）。
\end{proof}













\section{拟凝聚模层}
\label{modules-section-quasi-coherent}

\noindent
本节引入拟凝聚 $\mathcal{O}_X$-模层这一抽象概念。它在代数几何中非常有用，
因为概形上的拟凝聚模层在任意仿射开集上都有良好的描述。不过我们提醒读者：
在一般的（局部）环空间情形下，这一概念的行为很差。拟凝聚层范畴一般不是
阿贝尔范畴；拟凝聚层的无穷直和也未必拟凝聚；等等。

\begin{definition}
\label{modules-definition-quasi-coherent}
设 $(X, \mathcal{O}_X)$ 为环空间，$\mathcal{F}$ 为 $\mathcal{O}_X$-模层。
若对每个点 $x \in X$ 都存在开邻域 $x\in U \subset X$，使得
$\mathcal{F}|_U$ 同构于某个如下态射的余核：
$$
\bigoplus\nolimits_{j \in J}
\mathcal{O}_U
\longrightarrow
\bigoplus\nolimits_{i \in I}
\mathcal{O}_U
$$
则称 $\mathcal{F}$ 为{\it 拟凝聚 $\mathcal{O}_X$-模层}。
拟凝聚 $\mathcal{O}_X$-模层的范畴记为 $\QCoh(\mathcal{O}_X)$。
\end{definition}

\noindent
此定义是说，$X$ 可由这样的开集 $U$ 覆盖：$\mathcal{F}|_U$ 具有如下形式的
{\it 表示}：
$$
\bigoplus\nolimits_{j \in J}
\mathcal{O}_U
\longrightarrow
\bigoplus\nolimits_{i \in I}
\mathcal{O}_U
\longrightarrow
\mathcal{F}|_U
\longrightarrow
0.
$$
这里“表示”是指所展示的序列正合。换言之：
\begin{enumerate}
\item 对 $X$ 的每个点 $x$，都存在一个开邻域 $U$，使得
$\mathcal{F}|_U$ 由整体截面生成；
\item 适当选取这些截面后，相应满射的核也由整体截面生成。
\end{enumerate}

\begin{lemma}
\label{modules-lemma-direct-sum-quasi-coherent}
设 $(X, \mathcal{O}_X)$ 为环空间。两个拟凝聚 $\mathcal{O}_X$-模层的直和
仍是拟凝聚 $\mathcal{O}_X$-模层。
\end{lemma}

\begin{proof}
略。
\end{proof}

\begin{remark}
\label{modules-remark-infinite-direct-sum-quasi-coherent-not}
警告：拟凝聚 $\mathcal{O}_X$-模层的无穷直和一般并非拟凝聚的。
关于拟凝聚模层更为奇特的行为，见例 \ref{modules-example-quasi-coherent}。
\end{remark}

\begin{lemma}
\label{modules-lemma-pullback-quasi-coherent}
设 $f : (X, \mathcal{O}_X) \to (Y, \mathcal{O}_Y)$ 为环空间态射。
拟凝聚 $\mathcal{O}_Y$-模层的拉回 $f^*\mathcal{G}$ 是拟凝聚的。
\end{lemma}

\begin{proof}
按照引理 \ref{modules-lemma-pullback-locally-generated} 的证明，可设 $\mathcal{G}$
有一个由若干份 $\mathcal{O}_Y$ 的直和给出的整体表示。由引理
\ref{modules-lemma-exactness-pushforward-pullback}，我们已经知道 $f^*$
与所有余极限可交换，并且是右正合的。因此，若有正合列
$$
\bigoplus\nolimits_{j \in J}
\mathcal{O}_Y
\longrightarrow
\bigoplus\nolimits_{i \in I}
\mathcal{O}_Y
\longrightarrow
\mathcal{G}
\longrightarrow
0
$$
，则施用 $f^*$ 后得到正合列
$$
\bigoplus\nolimits_{j \in J}
\mathcal{O}_X
\longrightarrow
\bigoplus\nolimits_{i \in I}
\mathcal{O}_X
\longrightarrow
f^*\mathcal{G}
\longrightarrow
0.
$$
这就证明了引理。
\end{proof}

\noindent
这给出了大量拟凝聚层的例子。

\begin{lemma}
\label{modules-lemma-construct-quasi-coherent-sheaves}
设 $(X, \mathcal{O}_X)$ 为环空间，
$\alpha : R \to \Gamma(X, \mathcal{O}_X)$ 为从环 $R$ 到 $X$ 上整体截面环的
环同态，并设 $M$ 为 $R$-模。下列三种构造给出典范同构的
$\mathcal{O}_X$-模层：
\begin{enumerate}
\item 令 $\pi : (X, \mathcal{O}_X) \longrightarrow (\{*\}, R)$
为如下环空间态射：$\pi : X \to \{*\}$ 是唯一映射，且 $\pi$-映射
$\pi^\sharp$ 是给定同态 $\alpha : R \to \Gamma(X, \mathcal{O}_X)$。
置 $\mathcal{F}_1 = \pi^*M$。
\item 选取一个表示
$\bigoplus_{j \in J} R \to \bigoplus_{i \in I} R \to M \to 0$.
Set
$$
\mathcal{F}_2 = \Coker\left(
\bigoplus\nolimits_{j \in J} \mathcal{O}_X
\to
\bigoplus\nolimits_{i \in I} \mathcal{O}_X
\right).
$$
这里，在与 $j \in J$ 对应的分量 $\mathcal{O}_X$ 上，此态射由截面
$\sum_i \alpha(r_{ij})$ 给出，其中 $r_{ij}$ 是 $M$ 的表示中该态射的矩阵系数。
\item 令 $\mathcal{F}_3$ 为预层
$U \mapsto \mathcal{O}_X(U) \otimes_R M$ 所关联的层，其中映射
$R \to \mathcal{O}_X(U)$ 是 $\alpha$ 与限制映射
$\mathcal{O}_X(X) \to \mathcal{O}_X(U)$ 的复合。
\end{enumerate}
此构造具有下列性质：
\begin{enumerate}
\item 所得 $\mathcal{O}_X$-模层
$\mathcal{F}_M = \mathcal{F}_1 = \mathcal{F}_2 = \mathcal{F}_3$
是拟凝聚的。
\item 此构造给出从 $R$-模范畴到 $X$ 上拟凝聚层范畴的函子，
并且它与任意余极限可交换。
\item 对任意 $x \in X$，有
$\mathcal{F}_{M, x} = \mathcal{O}_{X, x} \otimes_R M$，且该等式关于 $M$ 是函子性的。
\item 给定任意 $\mathcal{O}_X$-模层 $\mathcal{G}$，有
$$
\Mor_{\mathcal{O}_X}(\mathcal{F}_M, \mathcal{G})
=
\Hom_R(M, \Gamma(X, \mathcal{G}))
$$
其中 $\Gamma(X, \mathcal{G})$ 上的 $R$-模结构由其
$\Gamma(X, \mathcal{O}_X)$-模结构经 $\alpha$ 得到。
\end{enumerate}
\end{lemma}

\begin{proof}
由于 $\pi^*$ 按定义就是第 (3) 项中预层的层化，故 $\mathcal{F}_1$
与 $\mathcal{F}_3$ 同构，见《层》第
\ref{sheaves-section-ringed-spaces-functoriality-modules} 节。
第 (2) 项与第 (1) 项的构造同构，是因为函子 $\pi^*$ 右正合，所以
$\pi^*(\bigoplus_{j \in J} R) \to \pi^*(\bigoplus_{i \in I} R) \to
\pi^*M \to 0$ 正合；$\pi^*$ 与任意直和可交换，见引理
\ref{modules-lemma-exactness-pushforward-pullback}；最后还要用到
$\pi^*(R) = \mathcal{O}_X$。

\medskip\noindent
断言 (1) 由构造 (2) 显然成立。断言 (2) 由 $\pi^*$ 具有上述性质显然成立。
断言 (3) 来自拉回层之茎的描述，见《层》引理
\ref{sheaves-lemma-stalk-pullback-modules}。断言 (4) 来自 $\pi_*$ 与
$\pi^*$ 的伴随性。
\end{proof}

\begin{definition}
\label{modules-definition-sheaf-associated}
在引理 \ref{modules-lemma-construct-quasi-coherent-sheaves} 的情形下，称
$\mathcal{F}_M$ 为{\it 与模 $M$ 及环同态 $\alpha$ 关联的层}。
若 $R = \Gamma(X, \mathcal{O}_X)$ 且 $\alpha = \text{id}_R$，
则简称 $\mathcal{F}_M$ 为{\it 与模 $M$ 关联的层}。
\end{definition}


\begin{lemma}
\label{modules-lemma-restrict-quasi-coherent}
设 $(X, \mathcal{O}_X)$ 为环空间，并置
$R = \Gamma(X, \mathcal{O}_X)$。设 $M$ 为 $R$-模，$\mathcal{F}_M$
为与 $M$ 关联的拟凝聚 $\mathcal{O}_X$-模层。若
$g : (Y, \mathcal{O}_Y) \to (X, \mathcal{O}_X)$ 为环空间态射，
则 $g^*\mathcal{F}_M$ 是与 $\Gamma(Y, \mathcal{O}_Y)$-模
$\Gamma(Y, \mathcal{O}_Y) \otimes_R M$ 关联的层。
\end{lemma}

\begin{proof}
由引理 \ref{modules-lemma-construct-quasi-coherent-sheaves} 中把
$\mathcal{F}_M$ 描述为 $\pi^*M$ 的第一种描述，以及如下环空间交换图，
即可得到断言：
$$
\xymatrix{
(Y, \mathcal{O}_Y) \ar[r]_-\pi \ar[d]_g &
(\{*\}, \Gamma(Y, \mathcal{O}_Y)) \ar[d]^{\text{由 }g^\sharp\text{ 诱导}} \\
(X, \mathcal{O}_X) \ar[r]^-\pi &
(\{*\}, \Gamma(X, \mathcal{O}_X))
}
$$
（还要使用《层》引理 \ref{sheaves-lemma-push-pull-composition-modules}。）
\end{proof}

\begin{lemma}
\label{modules-lemma-quasi-coherent-module}
设 $(X, \mathcal{O}_X)$ 为环空间，$x \in X$ 为一点。假设 $x$ 有拟紧邻域基。
任取拟凝聚 $\mathcal{O}_X$-模层 $\mathcal{F}$。则存在 $x$ 的开邻域 $U$，
使得 $\mathcal{F}|_U$ 同构于 $(U, \mathcal{O}_U)$ 上与某个
$\Gamma(U, \mathcal{O}_U)$-模 $M$ 关联的模层 $\mathcal{F}_M$。
\end{lemma}

\begin{proof}
首先，可用 $x$ 的一个开邻域替换 $X$，并设 $\mathcal{F}$ 同构于如下态射的余核：
$$
\Psi :
\bigoplus\nolimits_{j \in J}
\mathcal{O}_X
\longrightarrow
\bigoplus\nolimits_{i \in I}
\mathcal{O}_X.
$$
问题在于，这一态射未必由一个“矩阵”给出，因为直和的整体截面模一般不同于
各整体截面模的直和。

\medskip\noindent
取 $x$ 的拟紧邻域 $x \in E \subset X$（注意，$E$ 未必开），再取包含于
$E$ 的开邻域 $x \in U \subset E$。接着仿照引理
\ref{modules-lemma-section-direct-sum-quasi-compact} 的证明。对每个 $j \in J$，
以 $s_j \in \Gamma(X, \bigoplus\nolimits_{i \in I} \mathcal{O}_X)$
表示与 $j$ 对应的直和项 $\mathcal{O}_X$ 中截面 $1$ 的像。
存在有限族开集 $U_{jk}$（$k \in K_j$），使得
$E \subset \bigcup_{k \in K_j} U_{jk}$，并且每个限制
$s_j|_{U_{jk}}$ 都是有限和 $\sum_{i \in I_{jk}} f_{jki}$，其中
$I_{jk} \subset I$，且 $f_{jki}$ 属于与 $i \in I$ 对应的直和项
$\mathcal{O}_X$。置 $I_j = \bigcup_{k \in K_j} I_{jk}$；这是有限集。
由于 $U \subset E \subset \bigcup_{k \in K_j} U_{jk}$，截面 $s_j|_U$
是有限直和 $\bigoplus_{i \in I_j} \mathcal{O}_X$ 的截面。由引理
\ref{modules-lemma-limits-colimits}，实际上 $s_j|_U$ 是和
$\sum_{i \in I_j} f_{ij}$，其中
$f_{ij} \in \mathcal{O}_X(U) = \Gamma(U, \mathcal{O}_U)$。

\medskip\noindent
现在可将模 $M$ 定义为如下映射的余核：
$$
\bigoplus\nolimits_{j \in J}
\Gamma(U, \mathcal{O}_U)
\longrightarrow
\bigoplus\nolimits_{i \in I}
\Gamma(U, \mathcal{O}_U)
$$
其矩阵由 $(f_{ij})$ 给出。由引理
\ref{modules-lemma-construct-quasi-coherent-sheaves} 的构造 (2)，可见
$\mathcal{F}_M$ 与 $\mathcal{F}|_U$ 有相同的表示，因而
$\mathcal{F}_M \cong \mathcal{F}|_U$。
\end{proof}

\begin{example}
\label{modules-example-quasi-coherent}
令 $X$ 为可数多个实直线副本 $L_1, L_2, L_3, \ldots$ 在 $0$ 处全部粘合
而成的空间；$0$ 的一个邻域基为 $\{U_n\}_{n \in \mathbf{N}}$，其中
$U_n \cap L_i = (-1/n, 1/n)$。令 $\mathcal{O}_X$ 为实值连续函数层。
设 $f : \mathbf{R} \to \mathbf{R}$ 为连续函数，它在 $(-1, 1)$ 上恒为零，
而在 $(-\infty, -2) \cup (2, \infty)$ 上恒为 $1$。以 $f_n$ 表示 $X$
上的连续函数，它在每条 $L_j = \mathbf{R}$ 上都等于
$x \mapsto f(nx)$。令 $1_{L_j}$ 为 $L_j$ 的示性函数。考虑映射
$$
\bigoplus\nolimits_{j \in \mathbf{N}}
\mathcal{O}_X
\longrightarrow
\bigoplus\nolimits_{j, i \in \mathbf{N}}
\mathcal{O}_X, \quad
e_j \longmapsto \sum\nolimits_{i \in \mathbf{N}} f_j 1_{L_i} e_{ij}
$$
，其中记号含义显然。这一定义有意义，因为 $f_j$ 在 $0$ 的一个邻域上为零，
故该和局部有限。在 $U_n$ 上，当 $j > 2n$ 时，$e_j$ 的像不能写成
$\sum g_{ij} e_{ij}$ 这样的有限线性组合，其中 $g_{ij}$ 连续。因此，
不存在 $0 \in X$ 的邻域，使所展示的映射像上面引理
\ref{modules-lemma-quasi-coherent-module} 的证明中那样由一个“矩阵”给出。

\medskip\noindent
注意，$\bigoplus\nolimits_{j \in \mathbf{N}} \mathcal{O}_X$
是与以 $e_j$ 为基的自由模关联的层；另一个直和亦然。因此可见，
与模关联的层之间的态射一般即使在 $X$ 上局部地也未必来自模之间的态射。
类似地，应当存在环空间 $X$ 与拟凝聚 $\mathcal{O}_X$-模层 $\mathcal{F}$，
使得 $\mathcal{F}$ 局部地也不具有 $\mathcal{F}_M$ 的形式。
（若找到这样的例子，请发电子邮件告知。）此外，应当还存在局部紧空间 $X$
及态射 $\mathcal{F}_M \to \mathcal{F}_N$，它们局部地也不来自模之间的态射
（引理 \ref{modules-lemma-quasi-coherent-module} 的证明表明，若 $N$ 自由，
这种情形不会发生）。
\end{example}












\section{有限表示模层}
\label{modules-section-finite-presentation}

\noindent
定义如下。

\begin{definition}
\label{modules-definition-finite-presentation}
设 $(X, \mathcal{O}_X)$ 为环空间，$\mathcal{F}$ 为 $\mathcal{O}_X$-模层。
若对每个点 $x \in X$，都存在开邻域 $x\in U \subset X$ 以及
$n, m \in \mathbf{N}$，使得 $\mathcal{F}|_U$ 同构于如下态射的余核：
$$
\bigoplus\nolimits_{j = 1, \ldots, m}
\mathcal{O}_U
\longrightarrow
\bigoplus\nolimits_{i = 1, \ldots, n}
\mathcal{O}_U
$$
则称 $\mathcal{F}$ 是{\it 有限表示的}。
\end{definition}

\noindent
这就是说，$X$ 可由这样的开集 $U$ 覆盖：$\mathcal{F}|_U$ 有如下形式的
{\it 表示}：
$$
\bigoplus\nolimits_{j = 1, \ldots, m}
\mathcal{O}_U
\longrightarrow
\bigoplus\nolimits_{i = 1, \ldots, n}
\mathcal{O}_U
\to
\mathcal{F}|_U
\to
0.
$$
这里“表示”是指所展示的序列正合。换言之：
\begin{enumerate}
\item 对 $X$ 的每个点 $x$，都存在一个开邻域 $U$，使得
$\mathcal{F}|_U$ 由有限多个整体截面生成；
\item 适当选取这些截面后，相应满射的核也由有限多个整体截面生成。
\end{enumerate}

\begin{lemma}
\label{modules-lemma-finite-presentation-quasi-coherent}
设 $(X, \mathcal{O}_X)$ 为环空间。任意有限表示
$\mathcal{O}_X$-模层都是拟凝聚的。
\end{lemma}

\begin{proof}
由定义立即得到。
\end{proof}

\begin{lemma}
\label{modules-lemma-cokernel-finite-finite-presentation}
设 $(X,\mathcal{O}_X)$ 为环空间，$\mathcal{F}$ 为有限表示
$\mathcal{O}_X$-模层，并设
$\varphi : \mathcal{G} \to \mathcal{F}$ 为 $\mathcal{O}_X$-模层态射。
若 $\mathcal{G}$ 是有限型的，则 $\Coker(\varphi)$ 是有限表示的。
\end{lemma}

\begin{proof}
在 $X$ 上局部地，可写成
$\mathcal{F} = \mathcal{O}_X^{\oplus n} / \mathcal{M}$，其中
$\mathcal{M} \subset \mathcal{O}_X^{\oplus n}$ 是有限型
$\mathcal{O}_X$-子模层。因此
$\Im(\varphi) = \mathcal{N} / \mathcal{M}$，其中
$\mathcal{N} \subset \mathcal{O}_X^{\oplus n}$ 是包含 $\mathcal{M}$ 的
$\mathcal{O}_X$-子模层。由于 $\mathcal{G} \to \Im(\varphi)$ 为满射且
$\mathcal{G}$ 是有限型的，$\mathcal{O}_X$-模层 $\Im(\varphi)$ 是有限型的。
由引理 \ref{modules-lemma-extension-finite-type}，$\mathcal{N}$ 是有限型的。因此
$\Coker(\varphi) = \mathcal{O}_X^{\oplus n} / \mathcal{N}$ 是有限表示的。
\end{proof}

\begin{lemma}
\label{modules-lemma-kernel-surjection-finite-free-onto-finite-presentation}
设 $(X, \mathcal{O}_X)$ 为环空间，$\mathcal{F}$ 为有限表示
$\mathcal{O}_X$-模层。
\begin{enumerate}
\item 若 $\psi : \mathcal{O}_X^{\oplus r} \to \mathcal{F}$ 为满射，
则 $\Ker(\psi)$ 是有限型的。
\item 若 $\theta : \mathcal{G} \to \mathcal{F}$ 为满射且
$\mathcal{G}$ 是有限型的，则 $\Ker(\theta)$ 是有限型的。
\end{enumerate}
\end{lemma}

\begin{proof}
先证 (1)。设 $x \in X$。取 $x$ 的开邻域 $U \subset X$，使得存在表示
$$
\mathcal{O}_U^{\oplus m}
\xrightarrow{\chi}
\mathcal{O}_U^{\oplus n}
\xrightarrow{\varphi}
\mathcal{F}|_U
\to
0.
$$
令 $e_k$ 为生成 $\mathcal{O}_X^{\oplus r}$ 第 $k$ 个直和因子的截面。
对每个 $k = 1, \ldots, r$，把 $U$ 缩小为 $x$ 的一个较小邻域后，
可将 $\psi(e_k)$ 提升为 $U$ 上 $\mathcal{O}_U^{\oplus n}$ 的截面
$\tilde e_k$。由此得到层态射
$\alpha : \mathcal{O}_U^{\oplus r} \to \mathcal{O}_U^{\oplus n}$，
满足 $\varphi \circ \alpha = \psi$。类似地，缩小 $U$ 后可找到态射
$\beta : \mathcal{O}_U^{\oplus n} \to \mathcal{O}_U^{\oplus r}$，
满足 $\psi \circ \beta = \varphi$。于是映射
$$
\mathcal{O}_U^{\oplus m} \oplus
\mathcal{O}_U^{\oplus r}
\xrightarrow{\beta \circ \chi, 1 - \beta \circ \alpha}
\mathcal{O}_U^{\oplus r}
$$
满射到 $\psi$ 的核。

\medskip\noindent
为证 (2)，可局部选取满射
$\eta : \mathcal{O}_X^{\oplus r} \to \mathcal{G}$。由 (1)，
$\Ker(\theta \circ \eta)$ 是有限型的。由于
$\Ker(\theta) = \eta(\Ker(\theta \circ \eta))$，结论成立。
\end{proof}

\begin{lemma}
\label{modules-lemma-pullback-finite-presentation}
设 $f : (X, \mathcal{O}_X) \to (Y, \mathcal{O}_Y)$ 为环空间态射。
有限表示模层的拉回 $f^*\mathcal{G}$ 仍是有限表示的。
\end{lemma}

\begin{proof}
与引理 \ref{modules-lemma-pullback-quasi-coherent} 的证明完全相同，
只是指标集都是有限的。
\end{proof}

\begin{lemma}
\label{modules-lemma-quasi-coherent-limit-finite-presentation}
设 $(X, \mathcal{O}_X)$ 为环空间，置 $R = \Gamma(X, \mathcal{O}_X)$，
并设 $M$ 为 $R$-模。与 $M$ 关联的 $\mathcal{O}_X$-模层
$\mathcal{F}_M$ 是有限表示 $\mathcal{O}_X$-模层的有向余极限。
\end{lemma}

\begin{proof}
这立即由引理 \ref{modules-lemma-construct-quasi-coherent-sheaves}
以及任意模都是有限表示模的有向余极限这一事实得到，见《代数》引理
\ref{algebra-lemma-module-colimit-fp}。
\end{proof}

\begin{lemma}
\label{modules-lemma-finite-presentation-stalk-free}
设 $(X, \mathcal{O}_X)$ 为环空间，$\mathcal{F}$ 为有限表示
$\mathcal{O}_X$-模层。设 $x \in X$ 且
$\mathcal{F}_x \cong \mathcal{O}_{X, x}^{\oplus r}$。则存在 $x$ 的开邻域
$U$，使得 $\mathcal{F}|_U \cong \mathcal{O}_U^{\oplus r}$。
\end{lemma}

\begin{proof}
选取 $s_1, \ldots, s_r \in \mathcal{F}_x$，使它们经该同构映到
$\mathcal{O}_{X, x}^{\oplus r}$ 的一组基。取 $x$ 的开邻域 $U$，
使每个 $s_i$ 都提升为 $s_i \in \mathcal{F}(U)$。缩小 $U$ 后，
诱导态射 $\psi : \mathcal{O}_U^{\oplus r} \to \mathcal{F}|_U$ 为满射
（引理 \ref{modules-lemma-finite-type-surjective-on-stalk}）。由引理
\ref{modules-lemma-kernel-surjection-finite-free-onto-finite-presentation}，
$\Ker(\psi)$ 是有限型的。于是 $\Ker(\psi)_x = 0$ 意味着再次缩小 $U$ 后
$\Ker(\psi)$ 变为零（引理 \ref{modules-lemma-finite-type-stalk-zero}）。
\end{proof}





\section{凝聚模层}
\label{modules-section-coherent}

\noindent
本节的参考文献为 \cite{FAC}。

\medskip\noindent
环空间 $X$ 上的凝聚层范畴比拟凝聚层范畴更为合理：无论 $X$ 如何，
它至少都是 $\textit{Mod}(\mathcal{O}_X)$ 的阿贝尔子范畴。另一方面，
在一般环空间的情形下，凝聚模层的拉回“几乎从不”凝聚。

\begin{definition}
\label{modules-definition-coherent}
设 $(X, \mathcal{O}_X)$ 为环空间，$\mathcal{F}$ 为 $\mathcal{O}_X$-模层。
若下列两个条件成立，则称 $\mathcal{F}$ 为{\it 凝聚 $\mathcal{O}_X$-模层}：
\begin{enumerate}
\item $\mathcal{F}$ 是有限型的；
\item 对每个开集 $U \subset X$ 及每个有限族
$s_i \in \mathcal{F}(U)$（$i = 1, \ldots, n$），相应态射
$\bigoplus_{i = 1, \ldots, n} \mathcal{O}_U \to \mathcal{F}|_U$
的核都是有限型的。
\end{enumerate}
凝聚 $\mathcal{O}_X$-模层的范畴记为 $\textit{Coh}(\mathcal{O}_X)$。
\end{definition}

\begin{lemma}
\label{modules-lemma-coherent-finite-presentation}
设 $(X, \mathcal{O}_X)$ 为环空间。任意凝聚 $\mathcal{O}_X$-模层都是
有限表示的，因而是拟凝聚的。
\end{lemma}

\begin{proof}
设 $\mathcal{F}$ 为 $X$ 上的凝聚层。取一点 $x \in X$。由凝聚定义的 (1)，
可找到开邻域 $U$ 及 $\mathcal{F}$ 在 $U$ 上的截面
$s_i$（$i = 1, \ldots, n$），使得
$\Psi : \bigoplus_{i = 1, \ldots, n} \mathcal{O}_U \to \mathcal{F}$
为满射。由凝聚定义的 (2)，可找到开邻域 $x \in V \subset U$ 以及
$\bigoplus_{i = 1, \ldots, n} \mathcal{O}_V$ 的截面
$t_1, \ldots, t_m$，它们生成 $\Psi|_V$ 的核。于是，在 $V$ 上得到表示
$$
\bigoplus\nolimits_{j = 1, \ldots, m}
\mathcal{O}_V
\longrightarrow
\bigoplus\nolimits_{i = 1, \ldots, n}
\mathcal{O}_V
\to
\mathcal{F}|_V
\to
0
$$
，如所求。
\end{proof}

\begin{example}
\label{modules-example-coherent-not-Noetherian}
假设 $X$ 是一个点。此时上述定义给出了环上模的一个概念。
“凝聚”的定义意味着什么？它与诺特性密切相关，但二者并不相同。
事实上，环 $R = \mathbf{C}[x_1, x_2, x_3, \ldots]$ 作为自身上的模是凝聚的，
却不是诺特的。更多讨论见《代数》第 \ref{algebra-section-coherent} 节。
\end{example}

\begin{lemma}
\label{modules-lemma-coherent-abelian}
设 $(X, \mathcal{O}_X)$ 为环空间。
\begin{enumerate}
\item 凝聚层的任意有限型子层都是凝聚的。
\item 设 $\varphi : \mathcal{F} \to \mathcal{G}$ 为从有限型层
$\mathcal{F}$ 到凝聚层 $\mathcal{G}$ 的态射。则 $\Ker(\varphi)$ 是有限型的。
\item 设 $\varphi : \mathcal{F} \to \mathcal{G}$ 为凝聚
$\mathcal{O}_X$-模层之间的态射。则 $\Ker(\varphi)$ 与
$\Coker(\varphi)$ 都是凝聚的。
\item 给定 $\mathcal{O}_X$-模层的短正合列
$0 \to \mathcal{F}_1 \to \mathcal{F}_2 \to \mathcal{F}_3 \to 0$
，若其中两个模层凝聚，则第三个也凝聚。
\item 范畴 $\textit{Coh}(\mathcal{O}_X)$ 是
$\textit{Mod}(\mathcal{O}_X)$ 的弱 Serre 子范畴。特别地，凝聚模层范畴是
阿贝尔范畴，且包含函子
$\textit{Coh}(\mathcal{O}_X) \to \textit{Mod}(\mathcal{O}_X)$
是正合的。
\end{enumerate}
\end{lemma}

\begin{proof}
定义 \ref{modules-definition-coherent} 的条件 (2) 对凝聚层的任意子层都成立。
因此得到 (1)。

\medskip\noindent
假设 (2) 的条件成立。来证明 $\Ker(\varphi)$ 是有限型的。取 $x \in X$，
并取 $x$ 在 $X$ 中的开邻域 $U$，使得 $\mathcal{F}|_U$ 由
$s_1, \ldots, s_n$ 生成。由定义 \ref{modules-definition-coherent}，诱导态射
$\bigoplus_{i = 1}^n \mathcal{O}_U \to \mathcal{G}$,
$e_i \mapsto \varphi(s_i)$ 的核 $\mathcal{K}$ 是有限型的。因此，
$\Ker(\varphi)$ 作为复合
$\mathcal{K} \to \bigoplus_{i = 1}^n \mathcal{O}_U \to \mathcal{F}$
的像，是有限型的。

\medskip\noindent
假设 (3) 的条件成立。由 (2)，$\varphi$ 的核是有限型的；再由 (1)，它是凝聚的。

\medskip\noindent
在相同条件下，来证明 $\Coker(\varphi)$ 凝聚。由于 $\mathcal{G}$ 是有限型的，
$\Coker(\varphi)$ 也是有限型的。设 $U \subset X$ 为开集，且
$\overline{s}_i \in \Coker(\varphi)(U)$（$i = 1, \ldots, n$）为截面。
只需证明相应态射
$\overline{\Psi} : \bigoplus_{i = 1}^n \mathcal{O}_U \to \Coker(\varphi)$
的核是有限型的。存在 $U$ 的一个开覆盖，使得在其中每个开集上，
所有 $\overline{s}_i$ 都提升为 $\mathcal{G}$ 的截面 $s_i$。
因此可设这在整个 $U$ 上成立。还可设存在 $\Im(\varphi)$ 在 $U$ 上的截面
$t_j$（$j = 1, \ldots, m$），它们在 $U$ 上生成 $\Im(\varphi)$。
令 $\Phi : \bigoplus_{j = 1}^m \mathcal{O}_U \to \Im(\varphi)$
由 $t_j$ 定义，并令
$\Psi :
\bigoplus_{j = 1}^m \mathcal{O}_U \oplus
\bigoplus_{i = 1}^n \mathcal{O}_U \to \mathcal{G}$
由 $t_j$ 与 $s_i$ 定义。考虑如下交换图：
$$
\xymatrix{
0 \ar[r] &
\bigoplus_{j = 1}^m \mathcal{O}_U \ar[d]_\Phi \ar[r] &
\bigoplus_{j = 1}^m \mathcal{O}_U \oplus
\bigoplus_{i = 1}^n \mathcal{O}_U \ar[d]_\Psi \ar[r] &
\bigoplus_{i = 1}^n \mathcal{O}_U \ar[d]_{\overline{\Psi}} \ar[r] &
0 \\
0 \ar[r] &
\Im(\varphi) \ar[r] &
\mathcal{G} \ar[r] &
\Coker(\varphi) \ar[r] &
0
}
$$
由蛇形引理得到正合列
$\Ker(\Psi) \to \Ker(\overline{\Psi}) \to 0$。由于 $\Ker(\Psi)$
是有限型模层，可见 $\Ker(\overline{\Psi})$ 是有限型的。

\medskip\noindent
证明 (4)。设
$0 \to \mathcal{F}_1 \to \mathcal{F}_2 \to \mathcal{F}_3 \to 0$
为 $\mathcal{O}_X$-模层的短正合列。由 (3)，只需证明：若
$\mathcal{F}_1$ 与 $\mathcal{F}_3$ 凝聚，则 $\mathcal{F}_2$ 也凝聚。
由引理 \ref{modules-lemma-extension-finite-type}，$\mathcal{F}_2$ 是有限型的。
设 $s_1, \ldots, s_n$ 为 $\mathcal{F}_2$ 的有限多个局部截面，
它们定义在 $X$ 的同一开集 $U$ 上。只需证明它们之间的关系模层
$\mathcal{K}$ 是有限型的。考虑如下交换图：
$$
\xymatrix{
0 \ar[r] &
0 \ar[r] \ar[d] &
\bigoplus_{i = 1}^{n} \mathcal{O}_U \ar[r] \ar[d] &
\bigoplus_{i = 1}^{n} \mathcal{O}_U \ar[r] \ar[d] &
0 \\
0 \ar[r] &
\mathcal{F}_1 \ar[r] &
\mathcal{F}_2 \ar[r] &
\mathcal{F}_3 \ar[r] &
0
}
$$
，其中记号含义显然。由蛇形引理得到短正合列
$0 \to \mathcal{K} \to \mathcal{K}_3 \to \mathcal{F}_1$
，其中 $\mathcal{K}_3$ 是诸截面 $s_i$ 在 $\mathcal{F}_3$ 中的像之间的关系模层。
由于 $\mathcal{F}_1$ 凝聚，可见 $\mathcal{K}$ 是从有限型模层到凝聚模层的
某个态射之核，故由 (2)，它是有限型的。

\medskip\noindent
证明 (5)。这是因为 (3) 与 (4) 表明《同调》引理
\ref{homology-lemma-characterize-weak-serre-subcategory} 适用。
\end{proof}

\begin{lemma}
\label{modules-lemma-coherent-structure-sheaf}
设 $(X, \mathcal{O}_X)$ 为环空间，$\mathcal{F}$ 为 $\mathcal{O}_X$-模层。
假设 $\mathcal{O}_X$ 作为 $\mathcal{O}_X$-模层是凝聚的。则
$\mathcal{F}$ 凝聚，当且仅当它是有限表示的。
\end{lemma}

\begin{proof}
略。
\end{proof}

\begin{lemma}
\label{modules-lemma-finite-type-to-coherent-injective-on-stalk}
设 $X$ 为环空间，$\varphi : \mathcal{G} \to \mathcal{F}$ 为
$\mathcal{O}_X$-模层同态，且 $x \in X$。假设 $\mathcal{G}$ 是有限型的，
$\mathcal{F}$ 凝聚，并且茎上的映射
$\varphi_x : \mathcal{G}_x \to \mathcal{F}_x$ 为单射。则存在开邻域
$x \in U \subset X$，使得 $\varphi|_U$ 为单射。
\end{lemma}

\begin{proof}
以 $\mathcal{K} \subset \mathcal{G}$ 表示 $\varphi$ 的核。由引理
\ref{modules-lemma-coherent-abelian}，$\mathcal{K}$ 是有限型
$\mathcal{O}_X$-模层。假设给出 $\mathcal{K}_x = 0$。由引理
\ref{modules-lemma-finite-type-stalk-zero}，存在 $x$ 的开邻域 $U$，使得
$\mathcal{K}|_U = 0$。于是这个 $U$ 满足要求。
\end{proof}













\section{环空间的闭浸入}
\label{modules-section-closed-immersion}

\noindent
何时称环空间态射
$i : (Z, \mathcal{O}_Z) \to (X, \mathcal{O}_X)$ 为闭浸入？

\medskip\noindent
受正规拓扑空间的闭浸入（配以连续函子层）或微分流形的闭浸入
（配以可微函数层）这些例子的启发，至少作如下假设似乎是自然的：
\begin{enumerate}
\item 映射 $i$ 是拓扑空间的闭浸入。
\item 相应态射 $\mathcal{O}_X \to i_*\mathcal{O}_Z$ 为满射。
以 $\mathcal{I}$ 表示其核。
\end{enumerate}
仅这些条件就已蕴含若干令人满意的结果。例如，我们将证明
$\mathcal{O}_Z$-模层范畴等价于被 $\mathcal{I}$ 零化的
$\mathcal{O}_X$-模层范畴，从而推广第 \ref{modules-section-closed-immersions} 节
关于阿贝尔群层的结果。

\medskip\noindent
不过，在 Stacks Project 中，我们选取的定义保证：若 $i$ 是闭浸入且
$(X, \mathcal{O}_X)$ 是概形，则 $(Z, \mathcal{O}_Z)$ 也是概形。
此外，在这种情形下，我们希望 $i_*$ 与 $i^*$ 给出拟凝聚
$\mathcal{O}_Z$-模层范畴和被 $\mathcal{I}$ 零化的拟凝聚
$\mathcal{O}_X$-模层范畴之间的等价。一个最低条件是
$i_*\mathcal{O}_Z$ 为拟凝聚 $\mathcal{O}_X$-模层。保证这一点的一种好办法
是要求 $\mathcal{I}$ 局部由截面生成。可以把这个条件理解为：
“$(Z, \mathcal{O}_Z)$ 在 $(X, \mathcal{O}_X)$ 上局部地由令某些正则函数
$f_i$（即 $\mathcal{O}_X$ 的局部截面）等于零来定义”。这引出如下定义。

\begin{definition}
\label{modules-definition-closed-immersion}
环空间的{\it 闭浸入}\footnote{这是非标准记法；见上面的讨论。}是态射
$i : (Z, \mathcal{O}_Z) \to (X, \mathcal{O}_X)$
，并具有下列性质：
\begin{enumerate}
\item 映射 $i$ 是拓扑空间的闭浸入。
\item 相应态射 $\mathcal{O}_X \to i_*\mathcal{O}_Z$ 为满射。
以 $\mathcal{I}$ 表示其核。
\item $\mathcal{O}_X$-模层 $\mathcal{I}$ 局部由截面生成。
\end{enumerate}
\end{definition}

\noindent
事实上，此定义仍不能保证拟凝聚 $\mathcal{O}_Z$-模层的 $i_*$
是拟凝聚 $\mathcal{O}_X$-模层。问题在于，尚不清楚如何把拟凝聚
$\mathcal{O}_Z$-模层的局部表示转化为其直像的局部表示。不过下述结论是平凡的。

\begin{lemma}
\label{modules-lemma-i-star-quasi-coherent}
设 $i : (Z, \mathcal{O}_Z) \to (X, \mathcal{O}_X)$ 为环空间的闭浸入，
$\mathcal{F}$ 为拟凝聚 $\mathcal{O}_Z$-模层。则在 $X$ 上局部地，
$i_*\mathcal{F}$ 是拟凝聚 $\mathcal{O}_X$-模层之间某个态射的余核。
\end{lemma}

\begin{proof}
由定义，$i_*\mathcal{O}_Z$ 是拟凝聚的。并且在 $Z$ 上局部地，
层 $\mathcal{F}$ 是若干份 $\mathcal{O}_Z$ 的直和之间某个态射的余核。
此外，同一个拟凝聚层的任意多份副本之直和是拟凝聚的。最后，
$i_*$ 与任意余极限可交换，见引理 \ref{modules-lemma-i-star-right-adjoint}。
略去一些细节。
\end{proof}

\begin{lemma}
\label{modules-lemma-i-star-reflects-finite-type}
设 $i : (Z, \mathcal{O}_Z) \to (X, \mathcal{O}_X)$ 为环空间态射。
假设 $i$ 是从 $Z$ 到 $X$ 的某个闭子集上的同胚，并且
$\mathcal{O}_X \to i_*\mathcal{O}_Z$ 为满射。设 $\mathcal{F}$ 为
$\mathcal{O}_Z$-模层。则 $i_*\mathcal{F}$ 是有限型的，当且仅当
$\mathcal{F}$ 是有限型的。
\end{lemma}

\begin{proof}
假设 $\mathcal{F}$ 是有限型的。取 $x \in X$。若 $x \not \in Z$，
则 $i_*\mathcal{F}$ 在 $x$ 的某个邻域上为零，因而在该邻域上有限生成。
若 $x = i(z)$，则取开邻域 $z \in V \subset Z$ 及截面
$s_1, \ldots, s_n \in \mathcal{F}(V)$，使它们在 $V$ 上生成
$\mathcal{F}$。对某个开集 $U \subset X$ 写成 $V = Z \cap U$。
注意，$U$ 是 $x$ 的邻域。显然，诸 $s_i$ 给出 $i_*\mathcal{F}$ 在
$U$ 上的截面 $s_i$。所得态射
$$
\bigoplus\nolimits_{i = 1, \ldots, n} \mathcal{O}_U
\longrightarrow
i_*\mathcal{F}|_U
$$
为满射，这可由考察其在茎上的作用看出（这里使用
$\mathcal{O}_X \to i_*\mathcal{O}_Z$ 为满射）。因此
$i_*\mathcal{F}$ 是有限型的。

\medskip\noindent
反过来，假设 $i_*\mathcal{F}$ 是有限型的。取 $z \in Z$，置
$x = i(z)$。依假设，存在 $x$ 的开邻域 $U \subset X$ 及截面
$s_1, \ldots, s_n \in (i_*\mathcal{F})(U)$，它们在 $U$ 上生成
$i_*\mathcal{F}$。置 $V = Z \cap U$。由 $i_*$ 的定义，诸截面 $s_i$
对应于 $\mathcal{F}$ 在 $V$ 上的截面 $s_i$。所得态射
$$
\bigoplus\nolimits_{i = 1, \ldots, n} \mathcal{O}_V
\longrightarrow
\mathcal{F}|_V
$$
为满射，这可由考察其在茎上的作用看出。因此 $\mathcal{F}$ 是有限型的。
\end{proof}

\begin{lemma}
\label{modules-lemma-i-star-equivalence}
设 $i : (Z, \mathcal{O}_Z) \to (X, \mathcal{O}_X)$ 为环空间态射。
假设 $i$ 是从 $Z$ 到 $X$ 的某个闭子集上的同胚，并且
$i^\sharp : \mathcal{O}_X \to i_*\mathcal{O}_Z$ 为满射。
以 $\mathcal{I} \subset \mathcal{O}_X$ 表示 $i^\sharp$ 的核。函子
$$
i_* :
\textit{Mod}(\mathcal{O}_Z)
\longrightarrow
\textit{Mod}(\mathcal{O}_X)
$$
是正合且全忠实的，其本质像由满足 $\mathcal{I}\mathcal{G} = 0$ 的
$\mathcal{O}_X$-模层 $\mathcal{G}$ 组成。
\end{lemma}

\begin{proof}
我们断言，对 $\mathcal{O}_Z$-模层 $\mathcal{F}$，典范态射
$$
i^*i_*\mathcal{F} \longrightarrow \mathcal{F}
$$
是同构。可在茎上检验这一点。设 $z \in Z$ 且 $x = i(z)$。有
$$
(i^*i_*\mathcal{F})_z =
(i_*\mathcal{F})_x \otimes_{\mathcal{O}_{X, x}} \mathcal{O}_{Z, z} =
\mathcal{F}_z \otimes_{\mathcal{O}_{X, x}} \mathcal{O}_{Z, z} =
\mathcal{F}_z
$$
，这里使用《层》引理 \ref{sheaves-lemma-stalk-pullback-modules}、
$\mathcal{O}_{Z, z}$ 是 $\mathcal{O}_{X, x}$ 的商这一事实，以及《层》引理
\ref{sheaves-lemma-stalks-closed-pushforward}。由此可知 $i_*$ 全忠实。

\medskip\noindent
设 $\mathcal{G}$ 为满足 $\mathcal{I}\mathcal{G} = 0$ 的
$\mathcal{O}_X$-模层。我们将证明典范态射
$$
\mathcal{G} \longrightarrow i_*i^*\mathcal{G}
$$
是同构。这就证明了 $\mathcal{G} = i_*\mathcal{F}$，其中
$\mathcal{F} = i^*\mathcal{G}$，从而完成证明。检验所展示的态射在茎上诱导同构。
若 $x \in X$ 且 $x \not \in i(Z)$，则 $\mathcal{G}_x = 0$，
因为此时 $\mathcal{I}_x = \mathcal{O}_{X, x}$。同上，由《层》引理
\ref{sheaves-lemma-stalks-closed-pushforward}，有
$(i_*i^*\mathcal{G})_x = 0$。另一方面，若 $x \in Z$，则得到态射
$$
\mathcal{G}_x
\longrightarrow
\mathcal{G}_x \otimes_{\mathcal{O}_{X, x}} \mathcal{O}_{Z, x}
$$
，这里使用《层》引理 \ref{sheaves-lemma-stalk-pullback-modules} 与
\ref{sheaves-lemma-stalks-closed-pushforward}。此态射是同构，因为
$\mathcal{O}_{Z, x} = \mathcal{O}_{X, x}/\mathcal{I}_x$，
且依假设 $\mathcal{G}_x$ 被 $\mathcal{I}_x$ 零化。
\end{proof}

\begin{remark}
\label{modules-remark-sections-support-in-closed-modules}
设 $(X, \mathcal{O}_X)$ 为环空间，$Z \subset X$ 为闭子集。对
$\mathcal{O}_X$-模层 $\mathcal{F}$，可以考虑{\it 支集包含于 $Z$ 的截面子模层}，
记为 $\mathcal{H}_Z(\mathcal{F})$，由下式定义：
$$
\mathcal{H}_Z(\mathcal{F})(U) =
\{s \in \mathcal{F}(U) \mid \text{Supp}(s) \subset U \cap Z\}
$$
注意，$\mathcal{H}_Z(\mathcal{F})(U)$ 是 $\mathcal{O}_X(U)$-模；
换言之，$\mathcal{H}_Z(\mathcal{F})$ 是 $\mathcal{O}_X$-模层。
由构造，它是 $\mathcal{F}$ 中支集包含于 $Z$ 的最大
$\mathcal{O}_X$-子模层。将引理 \ref{modules-lemma-i-star-equivalence}
用于环空间态射 $(Z, \mathcal{O}_X|_Z) \to (X, \mathcal{O}_X)$，
可将（并且我们确实将）$\mathcal{H}_Z(\mathcal{F})$ 视为 $Z$ 上的
$\mathcal{O}_X|_Z$-模层。由此得到函子
$$
\textit{Mod}(\mathcal{O}_X) \longrightarrow \textit{Mod}(\mathcal{O}_X|_Z),
\quad
\mathcal{F} \longmapsto \mathcal{H}_Z(\mathcal{F})
\text{视为一个 }\mathcal{O}_X|_Z\text{-模层，其底空间为 }Z
$$
此函子左正合，但一般并不正合。上述所有断言均直接由引理
\ref{modules-lemma-support-section-closed} 得出。显然，此构造与注
\ref{modules-remark-sections-support-in-closed} 中的构造相容。
\end{remark}

\begin{lemma}
\label{modules-lemma-adjoint-section-with-support}
设 $(X, \mathcal{O}_X)$ 为环空间，$i : Z \to X$ 为闭子集的包含映射。注
\ref{modules-remark-sections-support-in-closed-modules} 中的函子
$\mathcal{H}_Z : \textit{Mod}(\mathcal{O}_X) \to
\textit{Mod}(\mathcal{O}_X|_Z)$ 是
$i_* : \textit{Mod}(\mathcal{O}_X|_Z) \to \textit{Mod}(\mathcal{O}_X)$
的右伴随。
\end{lemma}

\begin{proof}
只需证明，对任意 $\mathcal{O}_X$-模层 $\mathcal{F}$ 以及任意
$\mathcal{O}_X|_Z$-模层 $\mathcal{G}$，都有
$$
\Hom_{\mathcal{O}_X|_Z}(\mathcal{G}, \mathcal{H}_Z(\mathcal{F})) =
\Hom_{\mathcal{O}_X}(i_*\mathcal{G}, \mathcal{F})
$$
这是显然的，因为 $i_*\mathcal{G}$ 的任意截面的支集终究都包含于 $Z$。
细节从略。
\end{proof}










\section{局部自由层}
\label{modules-section-locally-free}

\noindent
设 $(X, \mathcal{O}_X)$ 为环空间。我们的约定允许（某些）茎
$\mathcal{O}_{X, x}$ 为零环。这意味着在定义局部自由层的秩时必须稍加谨慎。

\begin{definition}
\label{modules-definition-locally-free}
设 $(X, \mathcal{O}_X)$ 为环空间，$\mathcal{F}$ 为 $\mathcal{O}_X$-模层。
\begin{enumerate}
\item 若对每个点 $x \in X$，都存在集合 $I$ 及开邻域
$x \in U \subset X$，使得 $\mathcal{F}|_U$ 作为
$\mathcal{O}_X|_U$-模层同构于
$\bigoplus_{i \in I} \mathcal{O}_X|_U$，则称 $\mathcal{F}$ {\it 局部自由}。
\item 若可将指标集 $I$ 选为有限集，则称 $\mathcal{F}$ {\it 有限局部自由}。
\item 若可将指标集 $I$ 选为基数 $r$ 的集合，则称 $\mathcal{F}$
{\it 有限局部自由且秩为 $r$}。
\end{enumerate}
\end{definition}

\noindent
（有限）局部自由层的有限直和仍是（有限）局部自由的。不过，
局部自由层的无穷直和未必局部自由。

\begin{lemma}
\label{modules-lemma-locally-free-quasi-coherent}
设 $(X, \mathcal{O}_X)$ 为环空间，$\mathcal{F}$ 为 $\mathcal{O}_X$-模层。
若 $\mathcal{F}$ 局部自由，则它拟凝聚。
\end{lemma}

\begin{proof}
略。
\end{proof}

\begin{lemma}
\label{modules-lemma-pullback-locally-free}
设 $f : (X, \mathcal{O}_X) \to (Y, \mathcal{O}_Y)$ 为环空间态射。
若 $\mathcal{G}$ 是局部自由 $\mathcal{O}_Y$-模层，则
$f^*\mathcal{G}$ 是局部自由 $\mathcal{O}_X$-模层。
\end{lemma}

\begin{proof}
略。
\end{proof}

\begin{lemma}
\label{modules-lemma-rank}
设 $(X, \mathcal{O}_X)$ 为环空间。假设 $\mathcal{O}_X$ 的支集为 $X$，
即 $\mathcal{O}_X$ 的所有茎都是非零环。设 $\mathcal{F}$ 为局部自由
$\mathcal{O}_X$-模层。存在局部常值函数
$$
\text{rank}_\mathcal{F} :
X \longrightarrow \{0, 1, 2, \ldots\}\cup\{\infty\}
$$
，使得对任意点 $x \in X$，凡集合 $I$ 满足 $\mathcal{F}$ 在 $x$ 的某个邻域上
同构于 $\bigoplus_{i\in I} \mathcal{O}_X$，其基数都等于
$\text{rank}_\mathcal{F}(x)$。
\end{lemma}

\begin{proof}
在引理的假设下，$I$ 的基数可由非零环 $\mathcal{O}_{X, x}$ 上自由模
$\mathcal{F}_x$ 的秩读出，并且它在 $x$ 的某个邻域上为常数。
\end{proof}

\begin{lemma}
\label{modules-lemma-map-finite-locally-free}
设 $(X, \mathcal{O}_X)$ 为环空间，$r \geq 0$。设
$\varphi : \mathcal{F} \to \mathcal{G}$ 为秩 $r$ 的有限局部自由
$\mathcal{O}_X$-模层之间的态射。则 $\varphi$ 是同构，当且仅当
$\varphi$ 为满射。
\end{lemma}

\begin{proof}
假设 $\varphi$ 为满射。取 $x \in X$。存在 $x$ 的开邻域 $U$，使得
$\mathcal{F}|_U$ 与 $\mathcal{G}|_U$ 都同构于
$\mathcal{O}_U^{\oplus r}$。选取 $\mathcal{G}|_U$ 的自由生成元的提升，
得到态射 $\psi : \mathcal{G}|_U \to \mathcal{F}|_U$，满足
$\varphi|_U \circ \psi = \text{id}$。因此，由 $\varphi$ 诱导的映射
$\Gamma(U, \mathcal{F}) \to \Gamma(U, \mathcal{G})$ 为满射。
由于 $\Gamma(U, \mathcal{F})$ 与 $\Gamma(U, \mathcal{G})$
作为 $\Gamma(U, \mathcal{O}_U)$-模都同构于
$\Gamma(U, \mathcal{O}_U)^{\oplus r}$，可应用《代数》引理
\ref{algebra-lemma-fun}，得到
$\Gamma(U, \mathcal{F}) \to \Gamma(U, \mathcal{G})$ 为单射。
证明完成。
\end{proof}

\begin{lemma}
\label{modules-lemma-direct-summand-of-locally-free-is-locally-free}
设 $(X, \mathcal{O}_X)$ 为环空间。若所有茎 $\mathcal{O}_{X, x}$
都是局部环，则有限局部自由 $\mathcal{O}_X$-模层的任意直和项都是有限局部自由的。
\end{lemma}

\begin{proof}
假设 $\mathcal{F}$ 是有限局部自由 $\mathcal{O}_X$-模层
$\mathcal{H}$ 的直和项。设 $x \in X$。则 $\mathcal{H}_x$ 是有限自由
$\mathcal{O}_{X, x}$-模。由于 $\mathcal{O}_{X, x}$ 是局部环，
对某个 $r$ 有 $\mathcal{F}_x \cong \mathcal{O}_{X, x}^{\oplus r}$，
见《代数》引理 \ref{algebra-lemma-finite-projective}。由引理
\ref{modules-lemma-finite-presentation-stalk-free}，$\mathcal{F}$ 在 $x$ 的某个开邻域上
是秩 $r$ 的自由模层。（注意，$\mathcal{F}$ 作为 $\mathcal{H}$ 的直和项是有限表示的。）
\end{proof}





\section{双线性映射}
\label{modules-section-bilinear}


\noindent
设 $(X, \mathcal{O}_X)$ 为环空间，$\mathcal{F}$、$\mathcal{G}$、
$\mathcal{H}$ 为 $\mathcal{O}_X$-模层。$\mathcal{O}_X$-模层的
{\it 双线性映射} $f : \mathcal{F} \times \mathcal{G} \to \mathcal{H}$
是所示的集合层态射，并且对每个开集 $U \subset X$，诱导映射
$$
\mathcal{F}(U) \times \mathcal{G}(U) \to \mathcal{H}(U)
$$
都是模的 $\mathcal{O}_X(U)$-双线性映射。等价地，可以模仿模的双线性映射的
通常公理，要求集合层的某些态射图交换。例如，公理
$f(x + y, z) = f(x, z) + f(y, z)$ 由下图的交换性表示：
$$
\xymatrix{
\mathcal{F} \times \mathcal{F} \times \mathcal{G}
\ar[rrr]_{(f \circ \text{pr}_{13}, f \circ \text{pr}_{23})}
\ar[d]_{(+ \circ \text{pr}_{12}, \text{pr}_3)} & & &
\mathcal{H} \times \mathcal{H}
\ar[d]^{+} \\
\mathcal{F} \times \mathcal{G} \ar[rrr]^f & & &
\mathcal{H}
}
$$
另一种刻画如下：若
$f : \mathcal{F} \times \mathcal{G} \to \mathcal{H}$ 是集合层态射，
并且它在 $X$ 的每个点的茎上都诱导模的双线性映射，则 $f$ 是模层的双线性映射。
这是因为，可以通过在茎上检验来判断局部截面是否相等。

\medskip\noindent
以 $\text{Mor}( - , - )$ 表示 $X$ 上集合层范畴中的态射。双线性映射还有如下刻画：
集合层态射 $f : \mathcal{F} \times \mathcal{G} \to \mathcal{H}$
是双线性的，如果对任意集合层 $\mathcal{S}$，规则
$$
\text{Mor}(\mathcal{S}, \mathcal{F}) \times
\text{Mor}(\mathcal{S}, \mathcal{G}) \to
\text{Mor}(\mathcal{S}, \mathcal{H}),\quad
(a, b) \mapsto f \circ (a \times b)
$$
都是环 $\text{Mor}(\mathcal{S}, \mathcal{O}_X)$ 上模的双线性映射。
我们通常不采用这一观点，因为从局部截面集的角度思考更容易，而且两种观点显然等价。

\medskip\noindent
最后，还可用另一种方式表述此定义：$\mathcal{O}_X$ 是集合层范畴中的环对象，
而 $\mathcal{F}$、$\mathcal{G}$、$\mathcal{H}$ 是该环对象上的模对象。
于是，在任意范畴中，都可对环对象上的模对象定义双线性映射。
为了表述何为环对象、何为环对象上的模对象，以及这类对象在范畴中的双线性映射，
假设该范畴具有有限积会很方便（但并非绝对必要）；集合层范畴确实具有这一性质。







\section{张量积}
\label{modules-section-tensor-product}

\noindent
我们已经在《层》第 \ref{sheaves-section-presheaves-modules} 节与第
\ref{sheaves-section-sheafification-presheaves-modules} 节中，
在换环的情形下简要讨论过张量积。现在将其推广到模层的张量积。

\medskip\noindent
设 $(X, \mathcal{O}_X)$ 为环空间，$\mathcal{F}$、$\mathcal{G}$
为 $\mathcal{O}_X$-模层。首先定义{\it 张量积预层}
$$
\mathcal{F} \otimes_{p, \mathcal{O}_X} \mathcal{G}
$$
：它给每个开集 $U \subset X$ 指派 $\mathcal{O}_X(U)$-模
$\mathcal{F}(U) \otimes_{\mathcal{O}_X(U)} \mathcal{G}(U)$。
然后把{\it 张量积层}定义为上述预层的层化：
$$
\mathcal{F} \otimes_{\mathcal{O}_X} \mathcal{G}
=
(\mathcal{F} \otimes_{p, \mathcal{O}_X} \mathcal{G})^\#
$$
它可刻画为这样的 $\mathcal{O}_X$-模层：对任意第三个
$\mathcal{O}_X$-模层 $\mathcal{H}$，都有
$$
\Hom_{\mathcal{O}_X}
(\mathcal{F} \otimes_{\mathcal{O}_X} \mathcal{G}, \mathcal{H})
=
\text{Bilin}_{\mathcal{O}_X}(\mathcal{F} \times \mathcal{G}, \mathcal{H}).
$$
这里右端表示第 \ref{modules-section-bilinear} 节所定义的
$\mathcal{O}_X$-模层双线性映射的集合。

\medskip\noindent
环 $R$ 上模 $M,N$ 的张量积满足对称性，即
$M \otimes_R N = N \otimes_R M$；因此模层的张量积也满足同样的性质，即有
$$
\mathcal{F} \otimes_{\mathcal{O}_X} \mathcal{G}
=
\mathcal{G} \otimes_{\mathcal{O}_X} \mathcal{F}
$$
，且关于 $\mathcal{F}$、$\mathcal{G}$ 是函子性的。由于模的张量积满足结合性，
还得到典范的函子性同构
$$
(\mathcal{F} \otimes_{\mathcal{O}_X} \mathcal{G})
\otimes_{\mathcal{O}_X} \mathcal{H}
=
\mathcal{F} \otimes_{\mathcal{O}_X}
(\mathcal{G} \otimes_{\mathcal{O}_X} \mathcal{H})
$$
，它关于 $\mathcal{F}$、$\mathcal{G}$、$\mathcal{H}$ 是函子性的。

\begin{lemma}
\label{modules-lemma-stalk-tensor-product}
设 $(X, \mathcal{O}_X)$ 为环空间，$\mathcal{F}$、$\mathcal{G}$
为 $\mathcal{O}_X$-模层，且 $x \in X$。存在 $\mathcal{O}_{X, x}$-模的典范同构
$$
(\mathcal{F} \otimes_{\mathcal{O}_X} \mathcal{G})_x
=
\mathcal{F}_x \otimes_{\mathcal{O}_{X, x}} \mathcal{G}_x
$$
，它关于 $\mathcal{F}$ 与 $\mathcal{G}$ 是函子性的。
\end{lemma}

\begin{proof}
略。
\end{proof}

\begin{lemma}
\label{modules-lemma-tensor-product-sheafification}
设 $(X, \mathcal{O}_X)$ 为环空间。设 $\mathcal{F}'$、$\mathcal{G}'$
为 $\mathcal{O}_X$-模预层，其层化分别为 $\mathcal{F}$、$\mathcal{G}$。则
$\mathcal{F} \otimes_{\mathcal{O}_X} \mathcal{G} =
(\mathcal{F}' \otimes_{p, \mathcal{O}_X} \mathcal{G}')^\#$.
\end{lemma}

\begin{proof}
略。
\end{proof}


\begin{lemma}
\label{modules-lemma-tensor-product-exact}
设 $(X, \mathcal{O}_X)$ 为环空间，$\mathcal{G}$ 为 $\mathcal{O}_X$-模层。若
$\mathcal{F}_1
\to \mathcal{F}_2
\to \mathcal{F}_3
\to 0$
是 $\mathcal{O}_X$-模层的正合列，则诱导序列
$$
\mathcal{F}_1 \otimes_{\mathcal{O}_X} \mathcal{G} \to
\mathcal{F}_2 \otimes_{\mathcal{O}_X} \mathcal{G} \to
\mathcal{F}_3 \otimes_{\mathcal{O}_X} \mathcal{G} \to
0
$$
正合。
\end{lemma}

\begin{proof}
这由以下三点得出：可在茎上检验正合性（引理 \ref{modules-lemma-abelian}）；
茎的描述（引理 \ref{modules-lemma-stalk-tensor-product}）；以及模的张量积的相应结果
（《代数》引理 \ref{algebra-lemma-tensor-product-exact}）。
\end{proof}

\begin{lemma}
\label{modules-lemma-tensor-product-pullback}
设 $f : (X, \mathcal{O}_X) \to (Y, \mathcal{O}_Y)$ 为环空间态射，
$\mathcal{F}$、$\mathcal{G}$ 为 $\mathcal{O}_Y$-模层。则
$f^*(\mathcal{F} \otimes_{\mathcal{O}_Y} \mathcal{G})
= f^*\mathcal{F} \otimes_{\mathcal{O}_X} f^*\mathcal{G}$
，且关于 $\mathcal{F}$、$\mathcal{G}$ 是函子性的。
\end{lemma}

\begin{proof}
略。
\end{proof}

\begin{lemma}
\label{modules-lemma-tensor-commute-colimits}
设 $(X, \mathcal{O}_X)$ 为环空间。对任意 $\mathcal{O}_X$-模层
$\mathcal{F}$，函子
$$
\textit{Mod}(\mathcal{O}_X) \longrightarrow \textit{Mod}(\mathcal{O}_X)
, \quad
\mathcal{G} \longmapsto \mathcal{F} \otimes_{\mathcal{O}_X} \mathcal{G}
$$
与任意余极限可交换。
\end{lemma}

\begin{proof}
设 $I$ 为预序集，$\{\mathcal{G}_i\}$ 为 $I$ 上的系统，并置
$\mathcal{G} = \colim_i \mathcal{G}_i$。回忆，$\mathcal{G}$ 是与预层
$\mathcal{G}' : U \mapsto \colim_i \mathcal{G}_i(U)$ 关联的层，见《层》第
\ref{sheaves-section-limits-sheaves} 节。由引理
\ref{modules-lemma-tensor-product-sheafification}，张量积
$\mathcal{F} \otimes_{\mathcal{O}_X} \mathcal{G}$ 是如下预层的层化：
$$
U \longmapsto
\mathcal{F}(U) \otimes_{\mathcal{O}_X(U)} \colim_i \mathcal{G}_i(U) =
\colim_i \mathcal{F}(U) \otimes_{\mathcal{O}_X(U)} \mathcal{G}_i(U)
$$
其中等号来自《代数》引理
\ref{algebra-lemma-tensor-products-commute-with-limits}。因此，引理由
$\textit{Mod}(\mathcal{O}_X)$ 中余极限的描述得出，见引理
\ref{modules-lemma-limits-colimits}。
\end{proof}

\begin{lemma}
\label{modules-lemma-tensor-product-permanence}
设 $(X, \mathcal{O}_X)$ 为环空间，$\mathcal{F}$、$\mathcal{G}$
为 $\mathcal{O}_X$-模层。
\begin{enumerate}
\item 若 $\mathcal{F}$、$\mathcal{G}$ 局部由截面生成，则
$\mathcal{F} \otimes_{\mathcal{O}_X} \mathcal{G}$ 亦然。
\item 若 $\mathcal{F}$、$\mathcal{G}$ 是有限型的，则
$\mathcal{F} \otimes_{\mathcal{O}_X} \mathcal{G}$ 亦然。
\item 若 $\mathcal{F}$、$\mathcal{G}$ 拟凝聚，则
$\mathcal{F} \otimes_{\mathcal{O}_X} \mathcal{G}$ 亦然。
\item 若 $\mathcal{F}$、$\mathcal{G}$ 是有限表示的，则
$\mathcal{F} \otimes_{\mathcal{O}_X} \mathcal{G}$ 亦然。
\item 若 $\mathcal{F}$ 是有限表示的且 $\mathcal{G}$ 凝聚，则
$\mathcal{F} \otimes_{\mathcal{O}_X} \mathcal{G}$ 凝聚。
\item 若 $\mathcal{F}$、$\mathcal{G}$ 凝聚，则
$\mathcal{F} \otimes_{\mathcal{O}_X} \mathcal{G}$ 亦然。
\item 若 $\mathcal{F}$、$\mathcal{G}$ 局部自由，则
$\mathcal{F} \otimes_{\mathcal{O}_X} \mathcal{G}$ 亦然。
\end{enumerate}
\end{lemma}

\begin{proof}
首先证明局部自由 $\mathcal{O}_X$-模层的张量积局部自由。只需证明
$(\bigoplus_{i \in I} \mathcal{O}_X) \otimes_{\mathcal{O}_X}
(\bigoplus_{j \in J} \mathcal{O}_X) \cong
\bigoplus_{(i, j) \in I \times J} \mathcal{O}_X$.
层 $\bigoplus_{i \in I} \mathcal{O}_X$ 是与预层
$U \mapsto \bigoplus_{i \in I} \mathcal{O}_X(U)$ 关联的层。因此，
该张量积是与如下预层关联的层：
$$
U \longmapsto
(\bigoplus\nolimits_{i \in I} \mathcal{O}_X(U))
\otimes_{\mathcal{O}_X(U)}
(\bigoplus\nolimits_{j \in J} \mathcal{O}_X(U)).
$$
由于对任意环 $R$ 都有
$(\bigoplus_{i \in I} R) \otimes_R (\bigoplus_{j \in J} R) =
\bigoplus_{(i, j) \in I \times J} R$，遂得到所需结论。

\medskip\noindent
若 $\mathcal{F}_2 \to \mathcal{F}_1 \to \mathcal{F} \to 0$ 正合，
则由引理 \ref{modules-lemma-tensor-product-exact}，复形
$\mathcal{F}_2 \otimes_{\mathcal{O}_X} \mathcal{G} \to
\mathcal{F}_1 \otimes_{\mathcal{O}_X} \mathcal{G} \to
\mathcal{F} \otimes_{\mathcal{O}_X} \mathcal{G} \to 0$
正合。由此可证明 (5)。事实上，在这种情形下，局部存在上述正合列，且
$\mathcal{F}_i$（$i = 1, 2$）有限自由。因此两项
$\mathcal{F}_2 \otimes_{\mathcal{O}_X} \mathcal{G}$ 与
$\mathcal{F}_1 \otimes_{\mathcal{O}_X} \mathcal{G}$
同构于 $\mathcal{G}$ 的有限直和（例如由引理
\ref{modules-lemma-tensor-commute-colimits}）。由于有限直和是凝聚层，
这两项凝聚，因而该态射的余核也凝聚，见引理 \ref{modules-lemma-coherent-abelian}。

\medskip\noindent
若此外 $\mathcal{G}_2 \to \mathcal{G}_1 \to \mathcal{G} \to 0$ 也正合，
则可见
$$
\mathcal{F}_2 \otimes_{\mathcal{O}_X} \mathcal{G}_1
\oplus
\mathcal{F}_1 \otimes_{\mathcal{O}_X} \mathcal{G}_2
\to
\mathcal{F}_1 \otimes_{\mathcal{O}_X} \mathcal{G}_1
\to
\mathcal{F} \otimes_{\mathcal{O}_X} \mathcal{G}
\to 0
$$
正合。由此例如可证明 (3)。事实上，假设意味着局部可找到上述表示，
其中 $\mathcal{F}_i$、$\mathcal{G}_i$ 都是自由
$\mathcal{O}_X$-模层。因此，所展示的表示也是用自由层给出的张量积表示。

\medskip\noindent
其余断言的证明从略。
\end{proof}





\section{平坦模层}
\label{modules-section-flat}

\noindent
可以完全仿照环上模的情形来定义平坦模层。

\begin{definition}
\label{modules-definition-flat}
设 $(X, \mathcal{O}_X)$ 为环空间。若函子
$$
\textit{Mod}(\mathcal{O}_X)
\longrightarrow
\textit{Mod}(\mathcal{O}_X), \quad
\mathcal{G} \mapsto \mathcal{G} \otimes_{\mathcal{O}_X} \mathcal{F}
$$
是正合的，则称 $\mathcal{O}_X$-模层 $\mathcal{F}$ {\it 平坦}。
\end{definition}

\noindent
平坦性可以通过考察茎来刻画。

\begin{lemma}
\label{modules-lemma-flat-stalks-flat}
设 $(X, \mathcal{O}_X)$ 为环空间。$\mathcal{O}_X$-模层 $\mathcal{F}$
平坦，当且仅当对所有 $x \in X$，茎 $\mathcal{F}_x$ 都是平坦
$\mathcal{O}_{X, x}$-模。
\end{lemma}

\begin{proof}
假设对所有 $x \in X$，$\mathcal{F}_x$ 都是平坦
$\mathcal{O}_{X, x}$-模。此时，若
$\mathcal{G} \to \mathcal{H} \to \mathcal{K}$ 正合，则
$\mathcal{G} \otimes_{\mathcal{O}_X} \mathcal{F} \to
\mathcal{H} \otimes_{\mathcal{O}_X} \mathcal{F} \to
\mathcal{K} \otimes_{\mathcal{O}_X} \mathcal{F}$
也正合，因为可以在茎上检验正合性，且张量积与取茎可交换，见引理
\ref{modules-lemma-stalk-tensor-product}。反过来，假设 $\mathcal{F}$ 平坦，
并设 $x \in X$。考虑摩天楼层 $i_{x, *} M$，其中 $M$ 为
$\mathcal{O}_{X, x}$-模。注意
$$
M \otimes_{\mathcal{O}_{X, x}} \mathcal{F}_x =
\left(i_{x, *} M \otimes_{\mathcal{O}_X} \mathcal{F}\right)_x
$$
，这再次由引理 \ref{modules-lemma-stalk-tensor-product} 得到。由于 $i_{x, *}$ 正合，
$\mathcal{F}$ 平坦意味着
$M \mapsto M \otimes_{\mathcal{O}_{X, x}} \mathcal{F}_x$ 正合。
因此 $\mathcal{F}_x$ 是平坦 $\mathcal{O}_{X, x}$-模。
\end{proof}

\noindent
因此下述定义是合理的。

\begin{definition}
\label{modules-definition-flat-at-point}
设 $(X, \mathcal{O}_X)$ 为环空间，$x \in X$。若 $\mathcal{F}_x$
是平坦 $\mathcal{O}_{X, x}$-模，则称 $\mathcal{O}_X$-模层
$\mathcal{F}$ {\it 在 $x$ 处平坦}。
\end{definition}

\noindent
因此，$\mathcal{F}$ 是平坦 $\mathcal{O}_X$-模层，当且仅当它在每一点处平坦。

\begin{lemma}
\label{modules-lemma-base-change-flat}
设 $f : (X, \mathcal{O}_X) \to (Y, \mathcal{O}_Y)$ 为环空间态射。
若 $\mathcal{G}$ 是平坦 $\mathcal{O}_Y$-模层，则
$f^*\mathcal{G}$ 是平坦 $\mathcal{O}_X$-模层。
\end{lemma}

\begin{proof}
结合引理 \ref{modules-lemma-flat-stalks-flat}、《层》引理
\ref{sheaves-lemma-stalk-pullback-modules} 及《代数》引理
\ref{algebra-lemma-flat-base-change} 即得。
\end{proof}

\begin{lemma}
\label{modules-lemma-colimits-flat}
设 $(X, \mathcal{O}_X)$ 为环空间。平坦 $\mathcal{O}_X$-模层的滤过余极限
是平坦的；平坦 $\mathcal{O}_X$-模层的直和也是平坦的。
\end{lemma}

\begin{proof}
这由引理 \ref{modules-lemma-tensor-commute-colimits}、引理
\ref{modules-lemma-stalk-tensor-product}、《代数》引理
\ref{algebra-lemma-directed-colimit-exact}，以及可以在茎上检验正合性这一事实得出。
\end{proof}

\begin{lemma}
\label{modules-lemma-j-shriek-flat}
设 $(X, \mathcal{O}_X)$ 为环空间，$U \subset X$ 为开集。
层 $j_{U!}\mathcal{O}_U$ 是平坦 $\mathcal{O}_X$-模层。
\end{lemma}

\begin{proof}
$j_{U!}\mathcal{O}_U$ 的茎或为零，或等于 $\mathcal{O}_{X, x}$。
应用引理 \ref{modules-lemma-flat-stalks-flat} 即得。
\end{proof}

\begin{lemma}
\label{modules-lemma-module-quotient-flat}
设 $(X, \mathcal{O}_X)$ 为环空间。
\begin{enumerate}
\item 任意 $\mathcal{O}_X$-模层都是某个直和
$\bigoplus j_{U_i!}\mathcal{O}_{U_i}$ 的商。
\item 任意 $\mathcal{O}_X$-模层都是某个平坦 $\mathcal{O}_X$-模层的商。
\end{enumerate}
\end{lemma}

\begin{proof}
设 $\mathcal{F}$ 为 $\mathcal{O}_X$-模层。对每个开集 $U \subset X$
及每个 $s \in \mathcal{F}(U)$，都得到态射
$j_{U!}\mathcal{O}_U \to \mathcal{F}$，即态射
$\mathcal{O}_U \to \mathcal{F}|_U$，$1 \mapsto s$ 的伴随态射。显然，态射
$$
\bigoplus\nolimits_{(U, s)} j_{U!}\mathcal{O}_U
\longrightarrow
\mathcal{F}
$$
为满射；结合引理 \ref{modules-lemma-colimits-flat} 与
\ref{modules-lemma-j-shriek-flat} 可知其源平坦。
\end{proof}

\begin{lemma}
\label{modules-lemma-flat-tor-zero}
设 $(X, \mathcal{O}_X)$ 为环空间。设
$$
0 \to \mathcal{F}'' \to \mathcal{F}' \to \mathcal{F} \to 0
$$
为 $\mathcal{O}_X$-模层的短正合列。假设 $\mathcal{F}$ 平坦。
则对任意 $\mathcal{O}_X$-模层 $\mathcal{G}$，序列
$$
0 \to
\mathcal{F}'' \otimes_\mathcal{O} \mathcal{G} \to
\mathcal{F}' \otimes_\mathcal{O} \mathcal{G} \to
\mathcal{F} \otimes_\mathcal{O} \mathcal{G} \to 0
$$
正合。
\end{lemma}

\begin{proof}
对每个 $x \in X$，$\mathcal{F}_x$ 都是平坦 $\mathcal{O}_{X, x}$-模，
且可在茎上检验正合性；故结论由《代数》引理
\ref{algebra-lemma-flat-tor-zero} 得出。
\end{proof}

\begin{lemma}
\label{modules-lemma-flat-ses}
\begin{slogan}
环空间上平坦模层的满态射之核以及平坦模层的扩张仍是平坦的。
\end{slogan}
设 $(X, \mathcal{O}_X)$ 为环空间。设
$$
0 \to
\mathcal{F}_2 \to
\mathcal{F}_1 \to
\mathcal{F}_0 \to 0
$$
为 $\mathcal{O}_X$-模层的短正合列。
\begin{enumerate}
\item 若 $\mathcal{F}_2$ 与 $\mathcal{F}_0$ 平坦，则 $\mathcal{F}_1$ 也平坦。
\item 若 $\mathcal{F}_1$ 与 $\mathcal{F}_0$ 平坦，则 $\mathcal{F}_2$ 也平坦。
\end{enumerate}
\end{lemma}

\begin{proof}
由于正合性与平坦性都可在茎的层次上检验，结论由《代数》引理
\ref{algebra-lemma-flat-ses} 得出。
\end{proof}

\begin{lemma}
\label{modules-lemma-flat-resolution-of-flat}
设 $(X, \mathcal{O}_X)$ 为环空间。设
$$
\ldots \to
\mathcal{F}_2 \to
\mathcal{F}_1 \to
\mathcal{F}_0 \to
\mathcal{Q} \to 0
$$
为 $\mathcal{O}_X$-模层的正合复形。若 $\mathcal{Q}$ 与所有
$\mathcal{F}_i$ 都是平坦 $\mathcal{O}_X$-模层，则对任意
$\mathcal{O}_X$-模层 $\mathcal{G}$，复形
$$
\ldots \to
\mathcal{F}_2 \otimes_{\mathcal{O}_X} \mathcal{G} \to
\mathcal{F}_1 \otimes_{\mathcal{O}_X} \mathcal{G} \to
\mathcal{F}_0 \otimes_{\mathcal{O}_X} \mathcal{G} \to
\mathcal{Q} \otimes_{\mathcal{O}_X} \mathcal{G} \to 0
$$
也正合。
\end{lemma}

\begin{proof}
由引理 \ref{modules-lemma-flat-tor-zero}，把该复形拆成短正合列，并用引理
\ref{modules-lemma-flat-ses} 归纳证明
$\Im(\mathcal{F}_{i + 1} \to \mathcal{F}_i)$ 平坦，即得结论。
\end{proof}

\noindent
下述引理给出平坦性的方程判据（《代数》引理
\ref{algebra-lemma-flat-eq}）的一个方向。

\begin{lemma}
\label{modules-lemma-flat-eq}
设 $(X, \mathcal{O}_X)$ 为环空间，$\mathcal{F}$ 为平坦
$\mathcal{O}_X$-模层，$U \subset X$ 为开集，并设
$$
\mathcal{O}_U \xrightarrow{(f_1, \ldots, f_n)}
\mathcal{O}_U^{\oplus n} \xrightarrow{(s_1, \ldots, s_n)}
\mathcal{F}|_U
$$
为 $\mathcal{O}_U$-模层的复形。对每个 $x \in U$，都存在 $x$ 的开邻域
$V \subset U$ 以及 $(s_1, \ldots, s_n)|_V$ 的分解
$$
\mathcal{O}_V^{\oplus n}
\xrightarrow{A}
\mathcal{O}_V^{\oplus m} \xrightarrow{(t_1, \ldots, t_m)}
\mathcal{F}|_V
$$
，使得 $A \circ (f_1, \ldots, f_n)|_V = 0$。
\end{lemma}

\begin{proof}
令 $\mathcal{I} \subset \mathcal{O}_U$ 为由 $f_1, \ldots, f_n$
生成的理想层。则 $\sum f_i \otimes s_i$ 是
$\mathcal{I} \otimes_{\mathcal{O}_U} \mathcal{F}|_U$ 的截面，
并且它在 $\mathcal{F}|_U$ 中的像为零。由于 $\mathcal{F}|_U$ 平坦，态射
$\mathcal{I} \otimes_{\mathcal{O}_U} \mathcal{F}|_U \to \mathcal{F}|_U$
为单射。由于 $\mathcal{I} \otimes_{\mathcal{O}_U} \mathcal{F}|_U$
是与预层张量积关联的层，可见存在 $x$ 的开邻域 $V \subset U$，使得
$\sum f_i|_V \otimes s_i|_V$ 在
$\mathcal{I}(V) \otimes_{\mathcal{O}(V)} \mathcal{F}(V)$ 中为零。
利用《代数》引理 \ref{algebra-lemma-relations} 展开定义，得到
$t_1, \ldots, t_m \in \mathcal{F}(V)$ 及 $a_{ij} \in \mathcal{O}(V)$，
使得 $\sum a_{ij}f_i|_V = 0$ 且 $s_i|_V = \sum a_{ij}t_j$。
\end{proof}














\section{对偶}
\label{modules-section-duals}

\noindent
设 $(X, \mathcal{O}_X)$ 为环空间。赋予第
\ref{modules-section-tensor-product} 节所构造张量积的
$\mathcal{O}_X$-模层范畴是对称幺半范畴。对于
$\mathcal{O}_X$-模层 $\mathcal{F}$，下列条件等价：
\begin{enumerate}
\item $\mathcal{F}$ 在 $\mathcal{O}_X$-模层的幺半范畴中有左对偶；
\item $\mathcal{F}$ 局部为有限自由 $\mathcal{O}_X$-模层的直和项；
\item 作为 $\mathcal{O}_X$-模层，$\mathcal{F}$ 有限表示且平坦。
\end{enumerate}
本节的例 \ref{modules-example-dual} 及引理
\ref{modules-lemma-left-dual-module}、\ref{modules-lemma-flat-locally-finite-presentation}
证明了这一点。

\begin{example}
\label{modules-example-dual}
设 $(X, \mathcal{O}_X)$ 为环空间。设 $\mathcal{F}$ 为
$\mathcal{O}_X$-模层，且局部为有限自由
$\mathcal{O}_X$-模层的直和项。则态射
$$
\mathcal{F} \otimes_{\mathcal{O}_X}
\SheafHom_{\mathcal{O}_X}(\mathcal{F}, \mathcal{O}_X)
\longrightarrow
\SheafHom_{\mathcal{O}_X}(\mathcal{F}, \mathcal{F})
$$
是同构。事实上，这是一个局部问题；当 $\mathcal{F}$ 有限自由时结论成立，
而且若结论对某个模层成立，则对它的任一直和项也成立。记
$$
\eta :
\mathcal{O}_X
\longrightarrow
\mathcal{F} \otimes_{\mathcal{O}_X}
\SheafHom_{\mathcal{O}_X}(\mathcal{F}, \mathcal{O}_X)
$$
为把 $1$ 映到上述同构之下与
$\text{id}_\mathcal{F}$ 对应的截面的态射。记
$$
\epsilon : 
\SheafHom_{\mathcal{O}_X}(\mathcal{F}, \mathcal{O}_X)
\otimes_{\mathcal{O}_X} \mathcal{F}
\longrightarrow
\mathcal{O}_X
$$
为求值态射。于是
$\SheafHom_{\mathcal{O}_X}(\mathcal{F}, \mathcal{O}_X), \eta, \epsilon$
是《范畴论》定义 \ref{categories-definition-dual}
意义下 $\mathcal{F}$ 的左对偶。我们略去下列等式的验证：
$(1 \otimes \epsilon) \circ (\eta \otimes 1) = \text{id}_\mathcal{F}$
以及
$(\epsilon \otimes 1) \circ (1 \otimes \eta) =
\text{id}_{\SheafHom_{\mathcal{O}_X}(\mathcal{F}, \mathcal{O}_X)}$.
\end{example}

\begin{lemma}
\label{modules-lemma-left-dual-module}
设 $(X, \mathcal{O}_X)$ 为环空间，$\mathcal{F}$ 为
$\mathcal{O}_X$-模层。设 $\mathcal{G}, \eta, \epsilon$
是 $\mathcal{O}_X$-模层的幺半范畴中 $\mathcal{F}$ 的左对偶，参见
《范畴论》定义 \ref{categories-definition-dual}。则
\begin{enumerate}
\item $\mathcal{F}$ 局部为有限自由 $\mathcal{O}_X$-模层的直和项；
\item 态射
$e : \SheafHom_{\mathcal{O}_X}(\mathcal{F}, \mathcal{O}_X) \to \mathcal{G}$
把局部截面 $\lambda$ 映到 $(\lambda \otimes 1)(\eta)$，且它是同构；
\item 对 $\mathcal{F}$ 与 $\mathcal{G}$ 的局部截面 $f$ 与 $g$，有
$\epsilon(f, g) = e^{-1}(g)(f)$。
\end{enumerate}
\end{lemma}

\begin{proof}
这些假设意味着
$$
\mathcal{F} \xrightarrow{\eta \otimes 1}
\mathcal{F} \otimes_{\mathcal{O}_X} \mathcal{G}
\otimes_{\mathcal{O}_X} \mathcal{F}
\xrightarrow{1 \otimes \epsilon} \mathcal{F}
\quad\text{且}\quad
\mathcal{G} \xrightarrow{1 \otimes \eta}
\mathcal{G} \otimes_{\mathcal{O}_X} \mathcal{F}
\otimes_{\mathcal{O}_X} \mathcal{G}
\xrightarrow{\epsilon \otimes 1} \mathcal{G}
$$
均为恒等态射。设 $x \in X$。可以找到 $x$ 的开邻域 $U$，
以及 $\mathcal{F}$ 与 $\mathcal{G}$ 在 $U$ 上的有限多个截面
$f_1, \ldots, f_n$ 与 $g_1, \ldots, g_n$，使得
$\eta(1) = \sum f_i g_i$。记
$$
\mathcal{O}_U^{\oplus n} \to \mathcal{F}|_U
$$
为把第 $i$ 个基向量映到 $f_i$ 的态射。于是态射 $\eta|_U$
可经由态射
$\tilde \eta : \mathcal{O}_U \to
\mathcal{O}_U^{\oplus n} \otimes_{\mathcal{O}_U} \mathcal{G}|_U$
分解。由此得到交换图
$$
\xymatrix{
\mathcal{F}|_U
\ar[rr]_-{\eta \otimes 1} \ar[rrd]_-{\tilde \eta \otimes 1} & &
\mathcal{F}|_U \otimes \mathcal{G}|_U \otimes \mathcal{F}|_U
\ar[r]_-{1 \otimes \epsilon} &
\mathcal{F}|_U \\
& &
\mathcal{O}_U^{\oplus n} \otimes \mathcal{G}|_U \otimes \mathcal{F}|_U
\ar[u] \ar[r]^-{1 \otimes \epsilon} &
\mathcal{O}_U^{\oplus n} \ar[u]
}
$$
这表明 $\mathcal{F}$ 上的恒等态射在 $X$ 上局部地经由有限自由模层分解，
从而证明了 (1)。由《范畴论》引理
\ref{categories-lemma-left-dual} 及其证明得到 (2)。
由本证明中的第一个等式得到 (3)。也可以由左对偶的唯一性
（《范畴论》注 \ref{categories-remark-left-dual-adjoint}）
以及例 \ref{modules-example-dual} 中左对偶的构造推出 (2) 与 (3)。
\end{proof}

\begin{lemma}
\label{modules-lemma-flat-locally-finite-presentation}
设 $(X, \mathcal{O}_X)$ 为环空间。设 $\mathcal{F}$ 为有限表示的平坦
$\mathcal{O}_X$-模层。则 $\mathcal{F}$ 局部为有限自由
$\mathcal{O}_X$-模层的直和项。
\end{lemma}

\begin{proof}
以一个开覆盖的成员替换 $X$ 后，可以假设存在表示
$$
\mathcal{O}_X^{\oplus r} \to
\mathcal{O}_X^{\oplus n} \to \mathcal{F} \to 0
$$
设 $x \in X$。由引理 \ref{modules-lemma-flat-eq}，把 $X$
缩小为 $x$ 的某个开邻域后，可以假设存在分解
$$
\mathcal{O}_X^{\oplus n} \to
\mathcal{O}_X^{\oplus n_1} \to \mathcal{F}
$$
使得复合
$\mathcal{O}_X^{\oplus r} \to \mathcal{O}_X^{\oplus n} \to
\mathcal{O}_X^{\oplus n_1}$
零化 $\mathcal{O}_X^{\oplus r}$ 的第一个直和项。
再重复此论证 $r - 1$ 次，得到分解
$$
\mathcal{O}_X^{\oplus n} \to
\mathcal{O}_X^{\oplus n_r} \to \mathcal{F}
$$
使得复合
$\mathcal{O}_X^{\oplus r} \to \mathcal{O}_X^{\oplus n}
\to \mathcal{O}_X^{\oplus n_r}$ 为零。这意味着满射
$\mathcal{O}_X^{\oplus n_r} \to \mathcal{F}$ 有截面，故得证。
\end{proof}







\section{可构造集合层}
\label{modules-section-constructible}

\noindent
设 $X$ 为拓扑空间。给定集合 $S$，回顾 $\underline{S}$
或 $\underline{S}_X$ 表示取值为 $S$ 的常值层，参见
《层》定义 \ref{sheaves-definition-constant-sheaf}。
设 $U \subset X$ 为拓扑空间 $X$ 的开子集。我们以 $j_U$
表示包含态射，并以 $j_{U!} : \Sh(U) \to \Sh(X)$ 表示
《层》第 \ref{sheaves-section-open-immersions} 节所述的以空集延拓。

\begin{lemma}
\label{modules-lemma-surjection}
设 $X$ 为拓扑空间，$\mathcal{B}$ 为 $X$ 的一个拓扑基，
$\mathcal{F}$ 为 $X$ 上的集合层。则存在集合 $I$，并且对每个
$i \in I$ 存在 $U_i \in \mathcal{B}$ 及有限集 $S_i$，使得存在满射
$\coprod_{i \in I} j_{U_i!}\underline{S_i} \to \mathcal{F}$。
\end{lemma}

\begin{proof}
设 $S$ 为单点集。我们将以 $S_i = S$ 证明结论。对每个 $x \in X$
及元素 $s \in \mathcal{F}_x$，可选取 $U(x, s) \in \mathcal{B}$
与 $s(x, s) \in \mathcal{F}(U(x, s))$，使后者在
$\mathcal{F}_x$ 中映到 $s$。由《层》引理
\ref{sheaves-lemma-j-shriek}，截面 $s(x, s)$ 对应于层态射
$j_{U(x, s)!}\underline{S} \to \mathcal{F}$。于是
$$
\coprod\nolimits_{(x, s)} j_{U(x, s)!}\underline{S} \to \mathcal{F}
$$
在茎上满射，因而是满射。
\end{proof}

\begin{lemma}
\label{modules-lemma-filtered-colimit-constructibles}
设 $X$ 为拓扑空间，$\mathcal{B}$ 为 $X$ 的一个拓扑基，并假设每个
$U \in \mathcal{B}$ 都拟紧。则 $X$ 上的每个集合层都是如下形式的层的
滤过余极限：
\begin{equation}
\label{modules-equation-towards-constructible-sets}
\text{余等化子}\left(
\xymatrix{
\coprod\nolimits_{b = 1, \ldots, m} j_{V_b!}\underline{S_b}
\ar@<1ex>[r] \ar@<-1ex>[r] &
\coprod\nolimits_{a = 1, \ldots, n} j_{U_a!}\underline{S_a}
}
\right)
\end{equation}
其中 $U_a,V_b \in \mathcal{B}$，而 $S_a,S_b$ 为有限集。
\end{lemma}

\begin{proof}
由引理 \ref{modules-lemma-surjection}，每个集合层 $\mathcal{F}$
都是某个满射的靶，该满射的源 $\mathcal{F}_0$ 是形如
$j_{U!}\underline{S}$ 的层的余积，其中 $U \in \mathcal{B}$
且 $S$ 有限。把这一结论应用于
$\mathcal{F}_0 \times_\mathcal{F} \mathcal{F}_0$，可见
$\mathcal{F}$ 是一对态射的余等化子：
$$
\xymatrix{
\coprod\nolimits_{b \in B} j_{V_b!}\underline{S_b}
\ar@<1ex>[r] \ar@<-1ex>[r] &
\coprod\nolimits_{a \in A} j_{U_a!}\underline{S_a}
}
$$
这里 $A,B$ 为某些指标集，$V_b,U_a \in \mathcal{B}$，
而 $S_a,S_b$ 有限。对每个有限子集 $B' \subset B$，存在有限子集
$A' \subset A$，使得两态射都把 $b \in B'$ 上的余积映入
$a \in A'$ 上的余积。事实上，可以把右端视为具有单射转移态射的滤过余极限。
因此，在拟紧开集 $V_b$（$b \in B'$）上取截面与此余积可交换，参见
《层》引理 \ref{sheaves-lemma-directed-colimits-sections}。
所以我们的层是有限余积之间这些态射的余核的余极限。
\end{proof}

\begin{lemma}
\label{modules-lemma-constructible-comes-from-finite}
设 $X$ 为谱拓扑空间，$\mathcal{B}$ 为 $X$ 的拟紧开子集之集，
$\mathcal{F}$ 为式 (\ref{modules-equation-towards-constructible-sets})
所示的集合层。则存在到某个有限清醒拓扑空间 $Y$ 的连续谱映射
$f : X \to Y$，以及 $Y$ 上具有有限茎的集合层 $\mathcal{G}$，
使得 $f^{-1}\mathcal{G} \cong \mathcal{F}$。
\end{lemma}

\begin{proof}
可以把 $X = \lim X_i$ 写成有限清醒空间的有向逆极限，参见
《拓扑学》引理
\ref{topology-lemma-spectral-inverse-limit-finite-sober-spaces}。
转移态射 $X_{i'} \to X_i$ 自然都是谱映射，故由《拓扑学》引理
\ref{topology-lemma-directed-inverse-limit-spectral-spaces}，
态射 $p_i : X \to X_i$ 都是谱映射。对某个 $i$，可找到
$X_i$ 的开集 $U_{a, i}$ 与 $V_{b, i}$，其逆像分别为
$U_a$ 与 $V_b$，参见《拓扑学》引理
\ref{topology-lemma-descend-opens}。以 $\mathcal{F}$ 为余等化子的两个态射
$$
\beta, \gamma :
\coprod\nolimits_{b \in B} j_{V_b!}\underline{S_b}
\longrightarrow
\coprod\nolimits_{a \in A} j_{U_a!}\underline{S_a}
$$
经伴随对应于两族集合映射
$$
\beta_b, \gamma_b :
S_b
\longrightarrow
\Gamma(V_b, \coprod\nolimits_{a \in A} j_{U_a!}\underline{S_a}), \quad
b \in B
$$
。注意
$p_i^{-1}(j_{U_{a, i}!}\underline{S_a}) = j_{U_a!}\underline{S_a}$
，且 $(X_{i'} \to X_i)^{-1}(j_{U_{a, i}!}\underline{S_a}) =
j_{U_{a, i'}!}\underline{S_a}$。由《层》引理
\ref{sheaves-lemma-compute-pullback-to-limit}（并利用 $S_b$
与 $B$ 均为有限集），增大 $i$ 后可找到映射
$$
\beta_{b, i}, \gamma_{b, i} :
S_b
\longrightarrow
\Gamma(V_{b, i}, \coprod\nolimits_{a \in A} j_{U_{a, i}!}\underline{S_a})
, \quad b \in B
$$
，它们经 $p_i$ 拉回后给出映射 $\beta_b$ 与 $\gamma_b$。
这些映射又依次对应于 $X_i$ 上的层态射
$$
\beta_i, \gamma_i :
\coprod\nolimits_{b \in B} j_{V_{b, i}!}\underline{S_b}
\longrightarrow
\coprod\nolimits_{a \in A} j_{U_{a, i}!}\underline{S_a}
$$
。于是可取 $Y = X_i$ 以及
$$
\mathcal{G} =
\text{余等化子}\left(
\xymatrix{
\coprod\nolimits_{b = 1, \ldots, m} j_{V_{b, i}!}\underline{S_b}
\ar@<1ex>[r] \ar@<-1ex>[r] &
\coprod\nolimits_{a = 1, \ldots, n} j_{U_{a, i}!}\underline{S_a}
}
\right)
$$
我们略去若干细节。
\end{proof}

\begin{lemma}
\label{modules-lemma-constructible-in-constant}
设 $X$ 为谱拓扑空间，$\mathcal{B}$ 为 $X$ 的拟紧开子集之集，
$\mathcal{F}$ 为式 (\ref{modules-equation-towards-constructible-sets})
所示的集合层。则存在有限多个可构造闭子集
$Z_1, \ldots, Z_n \subset X$ 及有限集 $S_i$，使得
$\mathcal{F}$ 同构于 $\prod (Z_i \to X)_*\underline{S_i}$ 的一个子层。
\end{lemma}

\begin{proof}
由引理 \ref{modules-lemma-constructible-comes-from-finite}，可归约到有限清醒拓扑空间
以及具有有限茎的层这一情形。此时
$\mathcal{F} \subset \prod_{x \in X} i_{x, *}\mathcal{F}_x$，
其中 $i_x : \{x\} \to X$ 为嵌入。我们略去
$i_{x, *}\mathcal{F}_x$ 是 $\overline{\{x\}}$ 上常值层的证明。
\end{proof}










\section{环空间的平坦态射}
\label{modules-section-flat-morphisms}

\noindent
下述逐点定义由上面的引理 \ref{modules-lemma-flat-stalks-flat}
与定义 \ref{modules-definition-flat-at-point} 所启发。

\begin{definition}
\label{modules-definition-flat-morphism}
设 $f : X \to Y$ 为环空间态射，$x \in X$。若环同态
$\mathcal{O}_{Y, f(x)} \to \mathcal{O}_{X, x}$ 平坦，则称
$f$ {\it 在 $x$ 处平坦}。若 $f$ 在每个 $x \in X$ 处都平坦，
则称 $f$ {\it 平坦}。
\end{definition}

\noindent
考虑环层态射
$f^\sharp : f^{-1}\mathcal{O}_Y \to \mathcal{O}_X$。可见它在 $x$
处的茎就是环同态
$f^\sharp_x : \mathcal{O}_{Y, f(x)} \to \mathcal{O}_{X, x}$。
因此，$f$ 在 $x$ 处平坦，当且仅当作为
$f^{-1}\mathcal{O}_Y$-模层的 $\mathcal{O}_X$ 在 $x$ 处平坦；
$f$ 平坦，当且仅当 $\mathcal{O}_X$ 作为
$f^{-1}\mathcal{O}_Y$-模层平坦。平坦态射的一个非常特殊的情形是开浸入。

\begin{lemma}
\label{modules-lemma-pullback-flat}
设 $f : X \to Y$ 为环空间的平坦态射。则拉回函子
$f^* : \textit{Mod}(\mathcal{O}_Y) \to \textit{Mod}(\mathcal{O}_X)$
正合。
\end{lemma}

\begin{proof}
函子 $f^*$ 是正合函子
$f^{-1} : \textit{Mod}(\mathcal{O}_Y) \to \textit{Mod}(f^{-1}\mathcal{O}_Y)$
与换环函子
$$
\textit{Mod}(f^{-1}\mathcal{O}_Y) \to \textit{Mod}(\mathcal{O}_X), \quad
\mathcal{F} \longmapsto
\mathcal{F} \otimes_{f^{-1}\mathcal{O}_Y} \mathcal{O}_X.
$$
的复合。因此结论由定义 \ref{modules-definition-flat-morphism} 后的讨论得到。
\end{proof}

\begin{definition}
\label{modules-definition-flat-module}
设 $f : (X, \mathcal{O}_X) \to (Y, \mathcal{O}_Y)$ 为环空间态射，
$\mathcal{F}$ 为 $\mathcal{O}_X$-模层。
\begin{enumerate}
\item 若茎 $\mathcal{F}_x$ 是平坦 $\mathcal{O}_{Y, f(x)}$-模，
则称 $\mathcal{F}$ {\it 在点 $x \in X$ 处在 $Y$ 上平坦}。
\item 若 $\mathcal{F}$ 在 $X$ 的每一点 $x$ 处都在 $Y$ 上平坦，
则称 $\mathcal{F}$ {\it 在 $Y$ 上平坦}。
\end{enumerate}
\end{definition}

\noindent
根据此定义可见，$\mathcal{F}$ 在 $x$ 处在 $Y$ 上平坦，当且仅当
$\mathcal{F}$ 作为 $f^{-1}\mathcal{O}_Y$-模层在 $x$ 处平坦，因为由
《层》引理 \ref{sheaves-lemma-stalk-pullback}，有
$(f^{-1}\mathcal{O}_Y)_x = \mathcal{O}_{Y, f(x)}$。

\begin{lemma}
\label{modules-lemma-pullback-tensor-flat-module}
设 $f : X \to Y$ 为环空间态射，$\mathcal{F}$ 为在 $Y$ 上平坦的
$\mathcal{O}_X$-模层。则函子
$$
\textit{Mod}(\mathcal{O}_Y) \to \textit{Mod}(\mathcal{O}_X),\quad
\mathcal{G} \longmapsto f^*\mathcal{G} \otimes_{\mathcal{O}_X} \mathcal{F}
$$
正合。
\end{lemma}

\begin{proof}
这是因为
$f^*\mathcal{G} \otimes_{\mathcal{O}_X} \mathcal{F} =
f^{-1}\mathcal{G} \otimes_{f^{-1}\mathcal{O}_Y} \mathcal{F}$,
函子 $f^{-1}$ 正合，而 $\mathcal{F}$ 是平坦
$f^{-1}\mathcal{O}_Y$-模层。
\end{proof}














\section{对称幂与外幂}
\label{modules-section-symmetric-exterior}

\noindent
设 $(X, \mathcal{O}_X)$ 为环空间，$\mathcal{F}$ 为
$\mathcal{O}_X$-模层。定义 $\mathcal{F}$ 的{\it 张量代数}为
非交换 $\mathcal{O}_X$-代数层
$$
\text{T}(\mathcal{F})
=
\text{T}_{\mathcal{O}_X}(\mathcal{F})
= \bigoplus\nolimits_{n \geq 0} \text{T}^n(\mathcal{F}).
$$
这里 $\text{T}^0(\mathcal{F}) = \mathcal{O}_X$，
$\text{T}^1(\mathcal{F}) = \mathcal{F}$
，且当 $n \geq 2$ 时有
$$
\text{T}^n(\mathcal{F}) =
\mathcal{F} \otimes_{\mathcal{O}_X} \ldots \otimes_{\mathcal{O}_X} \mathcal{F}
\ \ (n\text{ 个因子})
$$
定义 $\wedge(\mathcal{F})$ 为 $\text{T}(\mathcal{F})$
关于如下双边理想的商：该理想由 $\text{T}^2(\mathcal{F})$
中形如 $s \otimes s$ 的局部截面生成，其中 $s$ 是
$\mathcal{F}$ 的局部截面。这称为 $\mathcal{F}$ 的{\it 外代数}。
类似地，定义 $\text{Sym}(\mathcal{F})$ 为
$\text{T}(\mathcal{F})$ 关于如下双边理想的商：该理想由
$\text{T}^2(\mathcal{F})$ 中形如
$s \otimes t - t \otimes s$ 的局部截面生成。

\medskip\noindent
代数 $\text{T}(\mathcal{F})$、$\wedge(\mathcal{F})$ 与
$\text{Sym}(\mathcal{F})$ 是《微分分次层》定义
\ref{sdga-definition-ga} 意义下的分次 $\mathcal{O}_X$-代数。
此外，$\text{Sym}(\mathcal{F})$ 交换，而
$\wedge(\mathcal{F})$ 分次交换。

\begin{lemma}
\label{modules-lemma-local-tensor-algebra}
在上述情形中，层 $\wedge^n\mathcal{F}$ 是预层
$$
U \longmapsto \wedge^n_{\mathcal{O}_X(U)}(\mathcal{F}(U)).
$$
的层化，参见《代数》第 \ref{algebra-section-tensor-algebra} 节。
类似地，层 $\text{Sym}^n\mathcal{F}$ 是预层
$$
U \longmapsto \text{Sym}^n_{\mathcal{O}_X(U)}(\mathcal{F}(U)).
$$
的层化。
\end{lemma}

\begin{proof}
略。用这种方式定义 $\text{Sym}(\mathcal{F})$
与 $\wedge(\mathcal{F})$，可能比上述方法更高效。
\end{proof}

\begin{lemma}
\label{modules-lemma-stalk-tensor-algebra}
在上述情形中，设 $x \in X$。存在 $\mathcal{O}_{X, x}$-模的典范同构
$\text{T}(\mathcal{F})_x = \text{T}(\mathcal{F}_x)$、
$\text{Sym}(\mathcal{F})_x = \text{Sym}(\mathcal{F}_x)$ 与
$\wedge(\mathcal{F})_x = \wedge(\mathcal{F}_x)$.
\end{lemma}

\begin{proof}
由上面的引理 \ref{modules-lemma-local-tensor-algebra} 与
《代数》引理 \ref{algebra-lemma-colimit-tensor-algebra} 立即得到。
\end{proof}

\begin{lemma}
\label{modules-lemma-pullback-tensor-algebra}
设 $f : (X, \mathcal{O}_X) \to (Y, \mathcal{O}_Y)$ 为环空间态射，
$\mathcal{F}$ 为 $\mathcal{O}_Y$-模层。则
$f^*\text{T}(\mathcal{F}) = \text{T}(f^*\mathcal{F})$；
与 $\mathcal{F}$ 关联的外代数与对称代数也有类似结论。
\end{lemma}

\begin{proof}
略。
\end{proof}

\begin{lemma}
\label{modules-lemma-presentation-sym-exterior}
设 $(X, \mathcal{O}_X)$ 为环空间。设
$\mathcal{F}_2 \to \mathcal{F}_1 \to \mathcal{F} \to 0$
为 $\mathcal{O}_X$-模层的正合列。对每个 $n \geq 1$，存在正合列
$$
\mathcal{F}_2 \otimes_{\mathcal{O}_X} \text{Sym}^{n - 1}(\mathcal{F}_1)
\to
\text{Sym}^n(\mathcal{F}_1)
\to
\text{Sym}^n(\mathcal{F})
\to
0
$$
以及类似的正合列
$$
\mathcal{F}_2 \otimes_{\mathcal{O}_X} \wedge^{n - 1}(\mathcal{F}_1)
\to
\wedge^n(\mathcal{F}_1)
\to
\wedge^n(\mathcal{F})
\to
0
$$
\end{lemma}

\begin{proof}
参见《代数》引理 \ref{algebra-lemma-presentation-sym-exterior}。
\end{proof}

\begin{lemma}
\label{modules-lemma-tensor-algebra-permanence}
设 $(X, \mathcal{O}_X)$ 为环空间，$\mathcal{F}$ 为
$\mathcal{O}_X$-模层。
\begin{enumerate}
\item 若 $\mathcal{F}$ 局部由截面生成，则每个
$\text{T}^n(\mathcal{F})$、$\wedge^n(\mathcal{F})$
与 $\text{Sym}^n(\mathcal{F})$ 也如此。
\item 若 $\mathcal{F}$ 为有限型，则每个
$\text{T}^n(\mathcal{F})$、$\wedge^n(\mathcal{F})$
与 $\text{Sym}^n(\mathcal{F})$ 也为有限型。
\item 若 $\mathcal{F}$ 有限表示，则每个
$\text{T}^n(\mathcal{F})$、$\wedge^n(\mathcal{F})$
与 $\text{Sym}^n(\mathcal{F})$ 也有限表示。
\item 若 $\mathcal{F}$ 凝聚，则对 $n > 0$，每个
$\text{T}^n(\mathcal{F})$、$\wedge^n(\mathcal{F})$
与 $\text{Sym}^n(\mathcal{F})$ 都凝聚。
\item 若 $\mathcal{F}$ 拟凝聚，则每个
$\text{T}^n(\mathcal{F})$、$\wedge^n(\mathcal{F})$
与 $\text{Sym}^n(\mathcal{F})$ 也拟凝聚。
\item 若 $\mathcal{F}$ 局部自由，则每个
$\text{T}^n(\mathcal{F})$、$\wedge^n(\mathcal{F})$
与 $\text{Sym}^n(\mathcal{F})$ 也局部自由。
\end{enumerate}
\end{lemma}

\begin{proof}
关于 $\text{T}^n(\mathcal{F})$ 的这些陈述由引理
\ref{modules-lemma-tensor-product-permanence} 得到。

\medskip\noindent
陈述 (1) 与 (2) 源于 $\wedge^n(\mathcal{F})$ 和
$\text{Sym}^n(\mathcal{F})$ 都是 $\text{T}^n(\mathcal{F})$ 的商。

\medskip\noindent
陈述 (6) 由《代数》引理
\ref{algebra-lemma-free-tensor-algebra} 得到。

\medskip\noindent
对 (3) 与 (5)，我们将使用上面的引理
\ref{modules-lemma-presentation-sym-exterior}。局部选取表示
$\mathcal{F}_2 \to \mathcal{F}_1 \to \mathcal{F} \to 0$
，其中 $\mathcal{F}_i$ 自由或有限自由，并应用该引理，可见
$\text{Sym}^n(\mathcal{F})$、$\wedge^n(\mathcal{F})$
有类似表示；这里使用了 (6) 与引理
\ref{modules-lemma-tensor-product-permanence}。

\medskip\noindent
为证明 (4)，我们将使用《代数》引理
\ref{algebra-lemma-present-sym-wedge}。可以在 $X$ 上局部化，并假设
$\mathcal{F}$ 由有限族整体截面 $(s_i)_{i \in I}$ 生成。
上述引理结合引理 \ref{modules-lemma-local-tensor-algebra} 表明，
当 $n \geq 2$ 时存在正合列
$$
\bigoplus\nolimits_{j \in J}
\text{T}^{n - 2}(\mathcal{F})
\to
\text{T}^n(\mathcal{F})
\to
\text{Sym}^n(\mathcal{F})
\to
0
$$
其中指标集 $J$ 有限。已知 $\text{T}^{n - 2}(\mathcal{F})$
有限生成，故第一个箭头的像是 $\text{T}^n(\mathcal{F})$
的凝聚子层，参见引理 \ref{modules-lemma-coherent-abelian}。
再由同一引理可知 $\text{Sym}^n(\mathcal{F})$ 凝聚。
\end{proof}

\begin{lemma}
\label{modules-lemma-whole-tensor-algebra-permanence}
设 $(X, \mathcal{O}_X)$ 为环空间，$\mathcal{F}$ 为
$\mathcal{O}_X$-模层。
\begin{enumerate}
\item 若 $\mathcal{F}$ 拟凝聚，则
$\text{T}(\mathcal{F})$、$\wedge(\mathcal{F})$
与 $\text{Sym}(\mathcal{F})$ 也都拟凝聚。
\item 若 $\mathcal{F}$ 局部自由，则
$\text{T}(\mathcal{F})$、$\wedge(\mathcal{F})$
与 $\text{Sym}(\mathcal{F})$ 也都局部自由。
\end{enumerate}
\end{lemma}

\begin{proof}
局部自由模层的无限直和 $\bigoplus \mathcal{G}_i$ 未必局部自由，
拟凝聚模层的无限直和也未必拟凝聚。问题在于，给定点 $x \in X$，
使 $\mathcal{G}_i$ 变为自由（相应地，具有适当表示）的 $x$ 的开邻域
$U_i$，其交集未必仍为 $x$ 的开邻域。然而，在引理
\ref{modules-lemma-tensor-algebra-permanence} 的证明中我们看到，
一旦为 $\mathcal{F}$ 选定适当的开邻域，该开邻域便同时适用于每个层
$\text{T}^n(\mathcal{F})$、$\wedge^n(\mathcal{F})$
及 $\text{Sym}^n(\mathcal{F})$。引理得证。
\end{proof}







\section{内 Hom}
\label{modules-section-internal-hom}

\noindent
设 $(X, \mathcal{O}_X)$ 为环空间，$\mathcal{F},\mathcal{G}$
为 $\mathcal{O}_X$-模层。考虑规则
$$
U \longmapsto \Hom_{\mathcal{O}_X|_U}(\mathcal{F}|_U, \mathcal{G}|_U).
$$
由《层》第 \ref{sheaves-section-glueing-sheaves} 节的讨论，
这是一个阿贝尔群层。此外，给定元素
$\varphi \in \Hom_{\mathcal{O}_X|_U}(\mathcal{F}|_U, \mathcal{G}|_U)$
与截面 $f \in \mathcal{O}_X(U)$，可以定义
$f\varphi \in \Hom_{\mathcal{O}_X|_U}(\mathcal{F}|_U, \mathcal{G}|_U)$
：或者前复合 $\mathcal{F}|_U$ 上的乘 $f$ 态射，或者后复合
$\mathcal{G}|_U$ 上的乘 $f$ 态射（两者给出相同结果）。
因此实际上得到一个 $\mathcal{O}_X$-模层，记为
$\SheafHom_{\mathcal{O}_X}(\mathcal{F}, \mathcal{G})$。
存在典范的“求值”态射
$$
\mathcal{F}
\otimes_{\mathcal{O}_X}
\SheafHom_{\mathcal{O}_X}(\mathcal{F}, \mathcal{G})
\longrightarrow
\mathcal{G}.
$$
对每个 $x \in X$，还存在典范态射
$$
\SheafHom_{\mathcal{O}_X}(\mathcal{F}, \mathcal{G})_x
\to
\Hom_{\mathcal{O}_{X, x}}(\mathcal{F}_x, \mathcal{G}_x)
$$
，但它很少是同构。

\begin{lemma}
\label{modules-lemma-internal-hom}
设 $(X, \mathcal{O}_X)$ 为环空间，$\mathcal{F},\mathcal{G},\mathcal{H}$
为 $\mathcal{O}_X$-模层。存在典范同构
$$
\SheafHom_{\mathcal{O}_X}
(\mathcal{F} \otimes_{\mathcal{O}_X} \mathcal{G}, \mathcal{H})
\longrightarrow
\SheafHom_{\mathcal{O}_X}
(\mathcal{F}, \SheafHom_{\mathcal{O}_X}(\mathcal{G}, \mathcal{H}))
$$
，它对三个变量都具有函子性（三处都是层 Hom）。特别地，给出态射
$\mathcal{F} \otimes_{\mathcal{O}_X} \mathcal{G} \to \mathcal{H}$
等价于给出态射
$\mathcal{F} \to \SheafHom_{\mathcal{O}_X}(\mathcal{G}, \mathcal{H})$。
\end{lemma}

\begin{proof}
这是《代数》引理
\ref{algebra-lemma-hom-from-tensor-product} 的类似结果。
证明相同，故略。
\end{proof}

\begin{lemma}
\label{modules-lemma-internal-hom-exact}
设 $(X, \mathcal{O}_X)$ 为环空间，$\mathcal{F},\mathcal{G}$
为 $\mathcal{O}_X$-模层。
\begin{enumerate}
\item 若 $\mathcal{F}_2 \to \mathcal{F}_1 \to \mathcal{F} \to 0$
是 $\mathcal{O}_X$-模层的正合列，则
$$
0 \to
\SheafHom_{\mathcal{O}_X}(\mathcal{F}, \mathcal{G}) \to
\SheafHom_{\mathcal{O}_X}(\mathcal{F}_1, \mathcal{G}) \to
\SheafHom_{\mathcal{O}_X}(\mathcal{F}_2, \mathcal{G})
$$
正合。
\item 若 $0 \to \mathcal{G} \to \mathcal{G}_1 \to \mathcal{G}_2$
是 $\mathcal{O}_X$-模层的正合列，则
$$
0 \to
\SheafHom_{\mathcal{O}_X}(\mathcal{F}, \mathcal{G}) \to
\SheafHom_{\mathcal{O}_X}(\mathcal{F}, \mathcal{G}_1) \to
\SheafHom_{\mathcal{O}_X}(\mathcal{F}, \mathcal{G}_2)
$$
正合。
\end{enumerate}
\end{lemma}

\begin{proof}
设 $\mathcal{F}_2 \to \mathcal{F}_1 \to \mathcal{F} \to 0$
如 (1) 所述。对每个开集 $U \subset X$，序列
$$
0 \to \Hom_{\mathcal{O}_U}(\mathcal{F}|_U, \mathcal{G}|_U) \to
\Hom_{\mathcal{O}_U}(\mathcal{F}_1|_U, \mathcal{G}|_U) \to
\Hom_{\mathcal{O}_U}(\mathcal{F}_2|_U, \mathcal{G}|_U)
$$
由《同调》引理 \ref{homology-lemma-check-exactness} 正合。
这意味着对 (1) 中的层序列在 $U$ 上取截面会得到阿贝尔群的正合列。
因此按照定义，(1) 中的序列正合。(2) 的证明完全相同。
\end{proof}

\begin{lemma}
\label{modules-lemma-adjoint-tensor-restrict}
设 $X$ 为拓扑空间，$\mathcal{O}_1 \to \mathcal{O}_2$
为环层同态。则有
$$
\Hom_{\mathcal{O}_1}(\mathcal{F}_{\mathcal{O}_1}, \mathcal{G}) =
\Hom_{\mathcal{O}_2}(\mathcal{F},
\SheafHom_{\mathcal{O}_1}(\mathcal{O}_2, \mathcal{G}))
$$
，它关于 $\mathcal{F} \in \textit{Mod}(\mathcal{O}_2)$
与 $\mathcal{G} \in \textit{Mod}(\mathcal{O}_1)$ 双函子地成立。
\end{lemma}

\begin{proof}
略。这是《代数》引理 \ref{algebra-lemma-adjoint-hom-restrict}
的类似结果，证明完全相同。
\end{proof}

\begin{lemma}
\label{modules-lemma-stalk-internal-hom}
设 $(X, \mathcal{O}_X)$ 为环空间，$\mathcal{F},\mathcal{G}$
为 $\mathcal{O}_X$-模层。若 $\mathcal{F}$ 为有限型，则典范态射
$$
\SheafHom_{\mathcal{O}_X}(\mathcal{F}, \mathcal{G})_x
\to
\Hom_{\mathcal{O}_{X, x}}(\mathcal{F}_x, \mathcal{G}_x)
$$
为单射。若 $\mathcal{F}$ 有限表示，则该典范态射为同构。
\end{lemma}

\begin{proof}
该态射把
$\SheafHom_{\mathcal{O}_X}(\mathcal{F}, \mathcal{G})_x$
中 $(U,\varphi)$ 的等价类映到在 $x$ 处的诱导茎态射
$\varphi_x : \mathcal{F}_x \to \mathcal{G}_x$，其中
$x \in U \subset X$ 为开集，且
$\varphi \in \Hom_{\mathcal{O}_U}(\mathcal{F}|_U, \mathcal{G}|_U)$。

\medskip\noindent
假设 $\mathcal{F}$ 为有限型。取该态射核中某元素 $\sigma$
的代表 $(U,\varphi)$，即 $\varphi_x=0$。必要时缩小 $U$，
选取生成 $\mathcal{F}|_U$ 的截面
$s^1, \ldots, s^n \in \mathcal{F}(U)$。由于
$\varphi_x(s^i_x)=0$，且所涉截面只有有限多个，可以找到
$x$ 的开邻域 $V \subset U$，使得对所有 $i=1,\ldots,n$ 都有
$\varphi_V(s^i|_V)=0$。由于 $s^i|_V$（$i=1,\ldots,n$）
生成 $\mathcal{F}|_V$，这意味着 $\varphi|_V=0$。
又因 $(U,\varphi)$ 与 $(V,\varphi|_V)$ 等价，故
$\sigma=0$，于是该态射为单射。

\medskip\noindent
接着假设 $\mathcal{F}$ 有限表示。在 $X$ 上局部化后，可以假设
$\mathcal{F}$ 有表示
$$
\bigoplus\nolimits_{j = 1, \ldots, m}
\mathcal{O}_X
\longrightarrow
\bigoplus\nolimits_{i = 1, \ldots, n}
\mathcal{O}_X
\to
\mathcal{F}
\to
0.
$$
由引理 \ref{modules-lemma-internal-hom-exact}，这给出正合列
$
0 \to
\SheafHom_{\mathcal{O}_X}(\mathcal{F}, \mathcal{G}) \to
\bigoplus\nolimits_{i = 1, \ldots, n} \mathcal{G}
\longrightarrow
\bigoplus\nolimits_{j = 1, \ldots, m} \mathcal{G}.
$
取茎得到正合列
$
0 \to
\SheafHom_{\mathcal{O}_X}(\mathcal{F}, \mathcal{G})_x \to
\bigoplus\nolimits_{i = 1, \ldots, n} \mathcal{G}_x
\longrightarrow
\bigoplus\nolimits_{j = 1, \ldots, m} \mathcal{G}_x
$
。结论随即成立，因为 $\mathcal{F}_x$ 位于正合列
$
\bigoplus\nolimits_{j = 1, \ldots, m}
\mathcal{O}_{X, x}
\longrightarrow
\bigoplus\nolimits_{i = 1, \ldots, n}
\mathcal{O}_{X, x}
\to
\mathcal{F}_x
\to
0
$
中，而该列诱导正合列
$
0 \to
\Hom_{\mathcal{O}_{X, x}}(\mathcal{F}_x, \mathcal{G}_x) \to
\bigoplus\nolimits_{i = 1, \ldots, n} \mathcal{G}_x
\longrightarrow
\bigoplus\nolimits_{j = 1, \ldots, m} \mathcal{G}_x
$
，它与上面的正合列相同。
\end{proof}

\begin{lemma}
\label{modules-lemma-pullback-internal-hom}
设 $f : (X, \mathcal{O}_X) \to (Y, \mathcal{O}_Y)$ 为环空间态射，
$\mathcal{F},\mathcal{G}$ 为 $\mathcal{O}_Y$-模层。
若 $\mathcal{F}$ 有限表示且 $f$ 平坦，则典范态射
$$
f^*\SheafHom_{\mathcal{O}_Y}(\mathcal{F}, \mathcal{G})
\longrightarrow
\SheafHom_{\mathcal{O}_X}(f^*\mathcal{F}, f^*\mathcal{G})
$$
为同构。
\end{lemma}

\begin{proof}
注意 $f^*\mathcal{F}$ 也有限表示
（引理 \ref{modules-lemma-pullback-finite-presentation}）。
设 $x \in X$ 映到 $y \in Y$。考察 $x$ 处的茎，由引理
\ref{modules-lemma-stalk-internal-hom} 与《代数详论》引理
\ref{more-algebra-lemma-pseudo-coherence-and-base-change-ext}，
可知在此情形中 $\Hom$ 与经
$\mathcal{O}_{Y,y}\to\mathcal{O}_{X,x}$ 的基变换可交换，
故得到同构。第二种证明：使用引理
\ref{modules-lemma-stalk-internal-hom} 证明中完全相同的论证。
\end{proof}

\begin{lemma}
\label{modules-lemma-internal-hom-locally-kernel-direct-sum}
设 $(X,\mathcal{O}_X)$ 为环空间，$\mathcal{F},\mathcal{G}$
为 $\mathcal{O}_X$-模层。若 $\mathcal{F}$ 有限表示，则层
$\SheafHom_{\mathcal{O}_X}(\mathcal{F}, \mathcal{G})$
局部是由 $\mathcal{G}$ 的有限直和之间某态射的核。
特别地，若 $\mathcal{G}$ 凝聚，则
$\SheafHom_{\mathcal{O}_X}(\mathcal{F}, \mathcal{G})$ 也凝聚。
\end{lemma}

\begin{proof}
第一个断言已见于引理 \ref{modules-lemma-stalk-internal-hom} 的证明。
关于凝聚层的结论随即由引理 \ref{modules-lemma-coherent-abelian} 得到。
\end{proof}

\begin{lemma}
\label{modules-lemma-hom-finite-presentation-colimit}
设 $X$ 为环空间，$\mathcal{F}$ 为有限表示的
$\mathcal{O}_X$-模层。设
$\mathcal{G} = \colim_{\lambda \in \Lambda} \mathcal{G}_\lambda$
为 $\mathcal{O}_X$-模层的滤过余极限。则典范态射
$$
\colim_\lambda \SheafHom_{\mathcal{O}_X}(\mathcal{F}, \mathcal{G}_\lambda)
\longrightarrow
\SheafHom_{\mathcal{O}_X}(\mathcal{F}, \mathcal{G})
$$
为同构。
\end{lemma}

\begin{proof}
模层的余极限与限制到开集可交换，参见《层》第
\ref{sheaves-section-limits-sheaves} 节。因此可以假设
$\mathcal{F}$ 有整体表示
$$
\bigoplus\nolimits_{j = 1, \ldots, m}
\mathcal{O}_X
\longrightarrow
\bigoplus\nolimits_{i = 1, \ldots, n}
\mathcal{O}_X
\to
\mathcal{F}
\to
0
$$
函子 $\SheafHom_{\mathcal{O}_X}(-,-)$ 对任一变量都与有限直和可交换，
而 $\SheafHom_{\mathcal{O}_X}(\mathcal{O}_X,-)$ 是恒等函子。
由此及引理 \ref{modules-lemma-internal-hom-exact}，得到正合列
$$
0 \to
\SheafHom_{\mathcal{O}_X}(\mathcal{F}, \mathcal{G}) \to
\bigoplus\nolimits_{i = 1, \ldots, n} \mathcal{G} \to
\bigoplus\nolimits_{j = 1, \ldots, m} \mathcal{G}
$$
由于滤过余极限在 $\textit{Mod}(\mathcal{O}_X)$ 中正合，
下列交换图的顶行也正合：
$$
\small
\xymatrix{
0 \ar[r] &
\colim_\lambda \SheafHom_{\mathcal{O}_X}(\mathcal{F}, \mathcal{G}_\lambda)
\ar[r] \ar[d] &
\colim_\lambda \bigoplus\nolimits_{i = 1, \ldots, n} \mathcal{G}_\lambda
\ar[r] \ar[d] &
\colim_\lambda \bigoplus\nolimits_{j = 1, \ldots, m} \mathcal{G}_\lambda
\ar[d] \\
0 \ar[r] &
\SheafHom_{\mathcal{O}_X}(\mathcal{F}, \mathcal{G}) \ar[r] &
\bigoplus\nolimits_{i = 1, \ldots, n} \mathcal{G} \ar[r] &
\bigoplus\nolimits_{j = 1, \ldots, m} \mathcal{G}
}
$$
右侧两个竖直箭头都是同构，故得结论。
\end{proof}

\begin{lemma}
\label{modules-lemma-finite-presentation-quasi-compact-colimit}
设 $X$ 为环空间，$I$ 为预序集，且
$(\mathcal{F}_i,\varphi_{ii'})$ 为 $I$ 上由
$\mathcal{O}_X$-模层组成的系统（参见《范畴论》第
\ref{categories-section-posets-limits} 节）。假设
\begin{enumerate}
\item $I$ 有向；
\item $\mathcal{G}$ 是有限表示的 $\mathcal{O}_X$-模层；
\item $X$ 有一个共尾的开覆盖系统
$\mathcal{U} : X = \bigcup_{j\in J} U_j$，
，其中 $J$ 有限，且对所有 $j,j'\in J$，
$U_j\cap U_{j'}$ 都拟紧。
\end{enumerate}
则有
$$
\colim_i \Hom_X(\mathcal{G}, \mathcal{F}_i)
=
\Hom_X(\mathcal{G}, \colim_i \mathcal{F}_i).
$$
\end{lemma}

\begin{proof}
置
$\mathcal{H} = \SheafHom_{\mathcal{O}_X}(\mathcal{G}, \colim \mathcal{F}_i)$
及 $\mathcal{H}_i = \SheafHom_{\mathcal{O}_X}(\mathcal{G}, \mathcal{F}_i)$。
回顾按照构造有
$$
\Hom_X(\mathcal{G}, \mathcal{F}) = \Gamma(X, \mathcal{H})
\quad\text{且}\quad
\Hom_X(\mathcal{G}, \mathcal{F}_i) = \Gamma(X, \mathcal{H}_i)
$$
。由引理 \ref{modules-lemma-hom-finite-presentation-colimit}，有
$\mathcal{H}=\colim\mathcal{H}_i$。因此本引理由《层》引理
\ref{sheaves-lemma-directed-colimits-sections} 得到。
\end{proof}

\begin{remark}
\label{modules-remark-condition-necessary}
在上述引理中，除 $X$ 拟紧之外还需要某种条件。参见
《层》例 \ref{sheaves-example-conditions-needed-colimit}。
\end{remark}




\section{模层的零化子}
\label{modules-section-annihilator}

\noindent
设 $(X,\mathcal{O}_X)$ 为环空间，$\mathcal{F}$ 为
$\mathcal{O}_X$-模层。存在典范的 $\mathcal{O}_X$-模层态射
$$
\mathcal{O}_X
\longrightarrow
\SheafHom_{\mathcal{O}_X}(\mathcal{F}, \mathcal{F})
$$
，它把局部截面 $f\in\mathcal{O}_X(U)$ 映到由乘 $f$ 给出的态射
$f:\mathcal{F}|_U\to\mathcal{F}|_U$。

\begin{definition}
\label{modules-definition-annihilator-sheaf}
设 $(X,\mathcal{O}_X)$ 为环空间，$\mathcal{F}$ 为
$\mathcal{O}_X$-模层。$\mathcal{F}$ 的{\it 零化子}，记作
$\text{Ann}_{\mathcal{O}_X}(\mathcal{F})$，是上述态射
$\mathcal{O}_X \to \SheafHom_{\mathcal{O}_X}(\mathcal{F}, \mathcal{F})$
的核。
\end{definition}

\noindent
对每个 $x\in X$，有 $\mathcal{O}_{X,x}$ 的理想之间的包含：
\begin{equation}
\label{modules-equation-inclusion-of-annihilator-ideals}
(\text{Ann}_{\mathcal{O}_X}(\mathcal{F}))_x
\subset
\text{Ann}_{\mathcal{O}_{X, x}}(\mathcal{F}_x)
\end{equation}
因为 $\text{Ann}_{\mathcal{O}_X}(\mathcal{F})$ 的任一截面，
都会零化 $\mathcal{F}$ 在其定义域中各点处的茎。下述简单情形中，
(\ref{modules-equation-inclusion-of-annihilator-ideals}) 成为等式。

\begin{lemma}
\label{modules-lemma-finite-type-annihilator-stalk}
设 $(X,\mathcal{O}_X)$ 为环空间，$\mathcal{F}$ 为
$\mathcal{O}_X$-模层。若 $\mathcal{F}$ 为有限型，则
$(\text{Ann}_{\mathcal{O}_X}(\mathcal{F}))_x =
\text{Ann}_{\mathcal{O}_{X, x}}(\mathcal{F}_x)$.
\end{lemma}

\begin{proof}
由引理 \ref{modules-lemma-stalk-internal-hom}，态射
$$
\SheafHom_{\mathcal{O}_X}(\mathcal{F}, \mathcal{F})_x \longrightarrow
\Hom_{\mathcal{O}_{X,x}}(\mathcal{F}_x, \mathcal{F}_x)
$$
为单射。因此，若 $\mathcal{O}_X$ 在 $x$ 的开邻域 $U$ 上的截面
$f$ 在 $\mathcal{F}_x$ 上作用为零，则存在开集
$U\supset V\ni x$，使得 $f$ 在 $\mathcal{F}|_V$ 上作用为零。
故包含式 (\ref{modules-equation-inclusion-of-annihilator-ideals}) 是等式。
\end{proof}

\begin{lemma}
\label{modules-lemma-induced-quotient-ring-sheaf-module-structure}
设 $(X,\mathcal{O}_X)$ 为环空间，$\mathcal{F}$ 为
$\mathcal{O}_X$-模层，$\mathcal{I}\subset\mathcal{O}_X$ 为理想层。若
$\mathcal{I} \subset \text{Ann}_{\mathcal{O}_X}(\mathcal{F})$,
则 $\mathcal{F}$ 有自然的 $\mathcal{O}_X/\mathcal{I}$-模层结构，
且它在茎上与通常的交换代数构造一致。
\end{lemma}

\begin{proof}
对包含态射 $\mathcal{I}\to\mathcal{O}_X$ 的余核应用泛性质，
得到 $\mathcal{O}_X$-模层的交换图
$$
\xymatrix{
\mathcal{O}_X \ar[r] \ar[d] &
\SheafHom_{\mathcal{O}_X}(\mathcal{F}, \mathcal{F}) \\
\mathcal{O}_X/\mathcal{I} \ar@{-->}[ur]
}
$$
。由引理 \ref{modules-lemma-internal-hom}，所得态射
$\mathcal{O}_X/\mathcal{I} \to
\SheafHom_{\mathcal{O}_X}(\mathcal{F}, \mathcal{F})$
对应于 $\mathcal{O}_X$-模层态射
$$
\mathcal{O}_X/\mathcal{I} \otimes_{\mathcal{O}_X} \mathcal{F}
\longrightarrow
\mathcal{F}
$$
，这意味着 $\mathcal{F}$ 上有与给定 $\mathcal{O}_X$-模层结构相容的
$\mathcal{O}_X/\mathcal{I}$-模层结构。我们略去关于茎的断言的验证。
\end{proof}

\begin{lemma}
\label{modules-lemma-coherent-annihilator}
设 $(X,\mathcal{O}_X)$ 为环空间。若
$\mathcal{O}_X$ 与 $\mathcal{F}$ 凝聚，则
$\text{Ann}_{\mathcal{O}_X}(\mathcal{F})$ 也凝聚。
\end{lemma}

\begin{proof}
由于 $\text{Ann}_{\mathcal{O}_X}(\mathcal{F})$ 是
$\mathcal{O}_X \to \SheafHom_{\mathcal{O}_X}(\mathcal{F}, \mathcal{F})$
的核，由引理 \ref{modules-lemma-coherent-abelian}，只需证明
$\SheafHom_{\mathcal{O}_X}(\mathcal{F}, \mathcal{F})$ 凝聚。
这由引理 \ref{modules-lemma-internal-hom-locally-kernel-direct-sum}
以及如下事实得到：$\mathcal{F}$ 凝聚，因而特别有限表示
（引理 \ref{modules-lemma-coherent-finite-presentation}）。
\end{proof}





\section{Koszul 复形}
\label{modules-section-koszul-complex}

\noindent
建议先阅读《代数详论》第 \ref{more-algebra-section-koszul} 节中
关于 Koszul 复形的内容。我们如下定义 $\mathcal{O}_X$-模层范畴中的
Koszul 复形。

\begin{definition}
\label{modules-definition-koszul}
设 $X$ 为环空间，$\varphi:\mathcal{E}\to\mathcal{O}_X$
为 $\mathcal{O}_X$-模层态射。与 $\varphi$ 关联的
{\it Koszul 复形} $K_\bullet(\varphi)$ 是如下定义的交换微分分次代数层：
\begin{enumerate}
\item 其底层分次代数是外代数
$K_\bullet(\varphi)=\wedge(\mathcal{E})$；
\item 微分 $d:K_\bullet(\varphi)\to K_\bullet(\varphi)$
是满足如下条件的唯一导子：对
$\mathcal{E}=K_1(\varphi)$ 的所有局部截面 $e$，
$d(e)=\varphi(e)$。
\end{enumerate}
\end{definition}

\noindent
具体地，若 $e_1\wedge\ldots\wedge e_n$ 是
$\mathcal{E}$ 的局部截面的外积，则
$$
d(e_1 \wedge \ldots \wedge e_n) =
\sum\nolimits_{i = 1, \ldots, n} (-1)^{i + 1}
\varphi(e_i)e_1 \wedge \ldots \wedge \widehat{e_i} \wedge \ldots \wedge e_n.
$$
易见这在张量代数上给出良定义的导子；它零化
$e\wedge e$，因而经由外代数分解。

\begin{definition}
\label{modules-definition-koszul-complex}
设 $X$ 为环空间，且
$f_1,\ldots,f_n\in\Gamma(X,\mathcal{O}_X)$。
{\it $f_1,\ldots,f_n$ 上的 Koszul 复形}是与态射
$(f_1, \ldots, f_n) : \mathcal{O}_X^{\oplus n} \to \mathcal{O}_X$.
关联的 Koszul 复形。记作
$K_\bullet(\mathcal{O}_X,f_1,\ldots,f_n)$
或 $K_\bullet(\mathcal{O}_X,f_\bullet)$。
\end{definition}

\noindent
当然，给定 $\mathcal{O}_X$-模层态射
$\varphi:\mathcal{E}\to\mathcal{O}_X$，若 $\mathcal{E}$
有限局部自由，则 $K_\bullet(\varphi)$ 在 $X$ 上局部同构于某个
Koszul 复形 $K_\bullet(\mathcal{O}_X,f_1,\ldots,f_n)$。



\section{可逆模层}
\label{modules-section-invertible}

\noindent
与环上模的情形类似（《代数详论》第
\ref{more-algebra-section-picard} 节），有下述定义。

\begin{definition}
\label{modules-definition-invertible}
设 $(X,\mathcal{O}_X)$ 为环空间。{\it 可逆
$\mathcal{O}_X$-模层}是满足如下条件的
$\mathcal{O}_X$-模层 $\mathcal{L}$：函子
$$
\textit{Mod}(\mathcal{O}_X) \longrightarrow \textit{Mod}(\mathcal{O}_X),\quad
\mathcal{F} \longmapsto \mathcal{L} \otimes_{\mathcal{O}_X} \mathcal{F}
$$
是范畴等价。若 $\mathcal{L}$ 作为 $\mathcal{O}_X$-模层
同构于 $\mathcal{O}_X$，则称 $\mathcal{L}$ {\it 平凡}。
\end{definition}

\noindent
下面的引理 \ref{modules-lemma-invertible-is-locally-free-rank-1}
说明它与秩 $1$ 局部自由模层之间的关系。

\begin{lemma}
\label{modules-lemma-invertible}
设 $(X,\mathcal{O}_X)$ 为环空间，$\mathcal{L}$ 为
$\mathcal{O}_X$-模层。下列条件等价：
\begin{enumerate}
\item $\mathcal{L}$ 可逆；
\item 存在 $\mathcal{O}_X$-模层 $\mathcal{N}$，使得
$\mathcal{L} \otimes_{\mathcal{O}_X} \mathcal{N} \cong \mathcal{O}_X$.
\end{enumerate}
在此情形中，$\mathcal{L}$ 局部为有限自由
$\mathcal{O}_X$-模层的直和项，且 (2) 中的模层 $\mathcal{N}$ 同构于
$\SheafHom_{\mathcal{O}_X}(\mathcal{L}, \mathcal{O}_X)$.
\end{lemma}

\begin{proof}
假设 (1) 成立。则函子 $-\otimes_{\mathcal{O}_X}\mathcal{L}$
本质满，故存在 (2) 所述的 $\mathcal{O}_X$-模层 $\mathcal{N}$。
若 (2) 成立，则函子 $-\otimes_{\mathcal{O}_X}\mathcal{N}$
是函子 $-\otimes_{\mathcal{O}_X}\mathcal{L}$ 的拟逆，故 (1) 成立。

\medskip\noindent
假设 (1) 与 (2) 成立。记
$\psi : \mathcal{L} \otimes_{\mathcal{O}_X} \mathcal{N} \to \mathcal{O}_X$
为给定同构。设 $x\in X$。选取开邻域 $U$、整数 $n\geq1$
以及截面 $s_i\in\mathcal{L}(U)$、$t_i\in\mathcal{N}(U)$，
使得 $\psi(\sum s_i\otimes t_i)=1$。考虑同构
$$
\mathcal{L}|_U \to
\mathcal{L}|_U \otimes_{\mathcal{O}_U}
\mathcal{L}|_U \otimes_{\mathcal{O}_U} \mathcal{N}|_U \to \mathcal{L}|_U
$$
，其中第一个箭头把 $s$ 映到 $\sum s_i\otimes s\otimes t_i$，
第二个箭头把 $s\otimes s'\otimes t$ 映到 $\psi(s'\otimes t)s$。
可知 $s\mapsto\sum\psi(s\otimes t_i)s_i$ 是
$\mathcal{L}|_U$ 的自同构。此自同构分解为
$$
\mathcal{L}|_U \to \mathcal{O}_U^{\oplus n} \to \mathcal{L}|_U
$$
，其中第一个箭头由
$s \mapsto (\psi(s \otimes t_1), \ldots, \psi(s \otimes t_n))$
给出，第二个箭头由 $(a_1,\ldots,a_n)\mapsto\sum a_i s_i$ 给出。
因此 $\mathcal{L}|_U$ 是有限自由 $\mathcal{O}_U$-模层的直和项。

\medskip\noindent
假设 (1) 与 (2) 成立。考虑求值态射
$$
\mathcal{L} \otimes_{\mathcal{O}_X}
\SheafHom_{\mathcal{O}_X}(\mathcal{L}, \mathcal{O}_X)
\longrightarrow \mathcal{O}_X
$$
为完成引理的证明，我们将通过检验它在茎上诱导同构来证明它是同构。
设 $x\in X$。由上一段已知 $\mathcal{L}$ 是有限表示的
$\mathcal{O}_X$-模层，故由引理 \ref{modules-lemma-stalk-internal-hom}，
只需证明
$$
\mathcal{L}_x \otimes_{\mathcal{O}_{X, x}}
\Hom_{\mathcal{O}_{X, x}}(\mathcal{L}_x, \mathcal{O}_{X, x})
\longrightarrow \mathcal{O}_{X, x}
$$
为同构。由于
$\mathcal{L}_x \otimes_{\mathcal{O}_{X, x}} \mathcal{N}_x =
(\mathcal{L} \otimes_{\mathcal{O}_X} \mathcal{N})_x =
\mathcal{O}_{X, x}$（引理 \ref{modules-lemma-stalk-tensor-product}）
，所需结论由《代数详论》引理
\ref{more-algebra-lemma-invertible} 得到。
\end{proof}

\begin{lemma}
\label{modules-lemma-pullback-invertible}
设 $f:(X,\mathcal{O}_X)\to(Y,\mathcal{O}_Y)$ 为环空间态射。
可逆 $\mathcal{O}_Y$-模层的拉回 $f^*\mathcal{L}$ 可逆。
\end{lemma}

\begin{proof}
由引理 \ref{modules-lemma-invertible}，存在 $\mathcal{O}_Y$-模层
$\mathcal{N}$，使得
$\mathcal{L} \otimes_{\mathcal{O}_Y} \mathcal{N} \cong \mathcal{O}_Y$.
拉回后，由引理 \ref{modules-lemma-tensor-product-pullback} 得到
$f^*\mathcal{L} \otimes_{\mathcal{O}_X} f^*\mathcal{N} \cong \mathcal{O}_X$
。因此由引理 \ref{modules-lemma-invertible}，$f^*\mathcal{L}$ 可逆。
\end{proof}

\begin{lemma}
\label{modules-lemma-invertible-is-locally-free-rank-1}
设 $(X,\mathcal{O}_X)$ 为环空间。任一秩 $1$ 局部自由
$\mathcal{O}_X$-模层都是可逆的。若所有茎
$\mathcal{O}_{X,x}$ 都是局部环，则逆命题也成立
（但一般情形并非如此）。
\end{lemma}

\begin{proof}
括号中的断言可由如下例子看出：考虑单点空间 $X$，其环层
$\mathcal{O}_X$ 由环 $R$ 给出。此时可逆
$\mathcal{O}_X$-模层对应于可逆 $R$-模，因此只要
$\Pic(R)$ 不是平凡群，就得到一个例子。

\medskip\noindent
假设 $\mathcal{L}$ 是秩 $1$ 局部自由模层，考虑求值态射
$$
\mathcal{L} \otimes_{\mathcal{O}_X}
\SheafHom_{\mathcal{O}_X}(\mathcal{L}, \mathcal{O}_X)
\longrightarrow \mathcal{O}_X
$$
在使 $\mathcal{L}$ 平凡化的开覆盖上考察，可见此态射为同构。
故由引理 \ref{modules-lemma-invertible}，$\mathcal{L}$ 可逆。

\medskip\noindent
假设所有茎 $\mathcal{O}_{X,x}$ 都是局部环，且 $\mathcal{L}$ 可逆。
在引理 \ref{modules-lemma-invertible} 的证明中已看到，对所有 $x\in X$，
$\mathcal{L}_x$ 都是可逆 $\mathcal{O}_{X,x}$-模。
由于 $\mathcal{O}_{X,x}$ 是局部环，故
$\mathcal{L}_x \cong \mathcal{O}_{X, x}$
（《代数详论》第 \ref{more-algebra-section-picard} 节）。
由引理 \ref{modules-lemma-invertible}，$\mathcal{L}$ 有限表示，
故由引理 \ref{modules-lemma-finite-presentation-stalk-free}，
$\mathcal{L}$ 是秩 $1$ 局部自由模层。
\end{proof}

\begin{lemma}
\label{modules-lemma-constructions-invertible}
设 $(X,\mathcal{O}_X)$ 为环空间。
\begin{enumerate}
\item 若 $\mathcal{L},\mathcal{N}$ 是可逆
$\mathcal{O}_X$-模层，则
$\mathcal{L} \otimes_{\mathcal{O}_X} \mathcal{N}$.
也可逆。
\item 若 $\mathcal{L}$ 是可逆 $\mathcal{O}_X$-模层，则
$\SheafHom_{\mathcal{O}_X}(\mathcal{L},\mathcal{O}_X)$ 也可逆，且求值态射
$\mathcal{L} \otimes_{\mathcal{O}_X}
\SheafHom_{\mathcal{O}_X}(\mathcal{L}, \mathcal{O}_X) \to \mathcal{O}_X$
为同构。
\end{enumerate}
\end{lemma}

\begin{proof}
(1) 由定义立即得到，(2) 由引理 \ref{modules-lemma-invertible} 及其证明得到。
\end{proof}

\begin{definition}
\label{modules-definition-powers}
设 $(X,\mathcal{O}_X)$ 为环空间。给定 $X$ 上的可逆层
$\mathcal{L}$ 及 $n\in\mathbf{Z}$，定义 $\mathcal{L}$ 的第
$n$ 个{\it 张量幂} $\mathcal{L}^{\otimes n}$ 为
$\mathcal{O}_X$ 在恰好应用 $n$ 次如下等价后的像：
$\mathcal{F} \mapsto\mathcal{F} \otimes_{\mathcal{O}_X} \mathcal{L}$
\end{definition}

\noindent
这对负的 $n$ 也有意义，因为我们已把可逆 $\mathcal{O}_X$-模层
定义为张量作用给出等价的模层。更具体地，有
$$
\mathcal{L}^{\otimes n} =
\left\{
\begin{matrix}
\mathcal{O}_X & \text{若} & n = 0 \\
\SheafHom_{\mathcal{O}_X}(\mathcal{L}, \mathcal{O}_X) & \text{若} & n = -1\\
\mathcal{L} \otimes_{\mathcal{O}_X} \ldots \otimes_{\mathcal{O}_X} \mathcal{L}
& \text{若} & n > 0 \\
\mathcal{L}^{\otimes -1} \otimes_{\mathcal{O}_X} \ldots
\otimes_{\mathcal{O}_X} \mathcal{L}^{\otimes -1}
& \text{若} & n < -1
\end{matrix}
\right.
$$
，参见引理 \ref{modules-lemma-constructions-invertible}。根据此定义，
有典范同构
$\mathcal{L}^{\otimes n} \otimes_{\mathcal{O}_X}
\mathcal{L}^{\otimes m} \to
\mathcal{L}^{\otimes n + m}$，且这些同构满足交换性约束与结合性约束
（略去表述）。

\medskip\noindent
设 $(X,\mathcal{O}_X)$ 为环空间。可在
$\bigoplus\Gamma(X,\mathcal{L}^{\otimes n})$ 上定义
$\mathbf{Z}$-分次环结构：把
$s\in\Gamma(X,\mathcal{L}^{\otimes n})$ 与
$t\in\Gamma(X,\mathcal{L}^{\otimes m})$ 映到
$\Gamma(X,\mathcal{L}^{\otimes n+m})$ 中与 $s\otimes t$ 对应的截面。
我们略去它定义含幺元的交换结合环的验证。然而，按照《代数》第
\ref{algebra-section-graded} 节中的约定，分次环在负次数中没有非零元。
这引出下述定义。

\begin{definition}
\label{modules-definition-gamma-star}
设 $(X,\mathcal{O}_X)$ 为环空间。给定 $X$ 上的可逆层
$\mathcal{L}$，定义{\it 关联分次环}为
$$
\Gamma_*(X, \mathcal{L})
=
\bigoplus\nolimits_{n \geq 0} \Gamma(X, \mathcal{L}^{\otimes n})
$$
给定 $\mathcal{O}_X$-模层 $\mathcal{F}$，置
$$
\Gamma_*(X, \mathcal{L}, \mathcal{F})
=
\bigoplus\nolimits_{n \in \mathbf{Z}} \Gamma(X,
\mathcal{F} \otimes_{\mathcal{O}_X} \mathcal{L}^{\otimes n})
$$
，将其视为分次 $\Gamma_*(X,\mathcal{L})$-模。
\end{definition}

\noindent
我们常简写为 $\Gamma_*(\mathcal{L})$ 与 $\Gamma_*(\mathcal{F})$
（但若 $\mathcal{F}$ 可逆，这一记法有歧义）。
$\Gamma_*(\mathcal{L})$ 在 $\Gamma_*(\mathcal{F})$ 上的乘法
用上述同构定义。若 $\gamma:\mathcal{F}\to\mathcal{G}$
为 $\mathcal{O}_X$-模层态射，则得到
$\Gamma_*(\mathcal{L})$-模同态
$\gamma:\Gamma_*(\mathcal{F})\to\Gamma_*(\mathcal{G})$。
若 $\alpha:\mathcal{L}\to\mathcal{N}$ 是可逆
$\mathcal{O}_X$-模层之间的态射，则得到分次环同态
$\Gamma_*(\mathcal{L})\to\Gamma_*(\mathcal{N})$。
若 $f:(Y,\mathcal{O}_Y)\to(X,\mathcal{O}_X)$ 为环空间态射，
且 $\mathcal{L}$ 在 $X$ 上可逆，则在 $Y$ 上得到可逆层
$f^*\mathcal{L}$（引理 \ref{modules-lemma-pullback-invertible}）
以及诱导的分次环同态
$$
f^* :
\Gamma_*(X, \mathcal{L})
\longrightarrow
\Gamma_*(Y, f^*\mathcal{L})
$$
此外，上述构造之间还存在一些相容性；我们略去其表述。

\begin{lemma}
\label{modules-lemma-pic-set}
设 $(X,\mathcal{O}_X)$ 为环空间。存在一族可逆模层
$\{\mathcal{L}_i\}_{i\in I}$，使得 $X$ 上的每个可逆模层
恰好同构于其中一个 $\mathcal{L}_i$。
\end{lemma}

\begin{proof}
回顾任一可逆 $\mathcal{O}_X$-模层局部都是有限自由
$\mathcal{O}_X$-模层的直和项，参见引理 \ref{modules-lemma-invertible}。
对每个开覆盖 $\mathcal{U}:X=\bigcup_{j\in J}U_j$
及映射 $r:J\to\mathbf{N}$，考虑满足如下条件的
$\mathcal{O}_X$-模层 $\mathcal{F}$：
$\mathcal{F}_j=\mathcal{F}|_{U_j}$ 是
$\mathcal{O}_{U_j}^{\oplus r(j)}$ 的直和项。
$\mathcal{F}_j$ 的同构类之全体是集合，因为
$\Hom_{\mathcal{O}_U}(\mathcal{O}_U^{\oplus r}, \mathcal{O}_U^{\oplus r})$
是集合。层 $\mathcal{F}$ 由粘合 $\mathcal{F}_j$ 得到，参见《层》第
\ref{sheaves-section-glueing-sheaves} 节。注意所有粘合数据之全体构成集合。
所有覆盖 $\mathcal{U} : X = \bigcup_{j \in J} U_i$ 之全体也构成集合，
其中 $J \to \mathcal{P}(X)$、$j \mapsto U_j$ 为单射。对每个覆盖，映射
$r:J\to\mathbf{N}$ 之全体为集合。因此所有 $\mathcal{F}$ 之全体构成集合。
\end{proof}

\noindent
粗略地说，此引理断言可逆层的同构类之全体构成集合。
引理 \ref{modules-lemma-constructions-invertible} 断言张量积在该集合上定义阿贝尔群结构。

\begin{definition}
\label{modules-definition-pic}
设 $(X,\mathcal{O}_X)$ 为环空间。$X$ 的 {\it Picard 群}
$\Pic(X)$ 是如下阿贝尔群：其元素为可逆
$\mathcal{O}_X$-模层的同构类，加法对应于张量积。
\end{definition}

\begin{lemma}
\label{modules-lemma-s-open}
\begin{slogan}
线丛的（局部）平凡化等价于（局部）非消没截面。
\end{slogan}
设 $X$ 为环空间。假设每个茎 $\mathcal{O}_{X,x}$ 都是
以 $\mathfrak m_x$ 为极大理想的局部环。设 $\mathcal{L}$
为可逆 $\mathcal{O}_X$-模层。对任一截面
$s\in\Gamma(X,\mathcal{L})$，集合
$$
X_s = \{x \in X \mid \text{像 }s \not\in \mathfrak m_x\mathcal{L}_x\}
$$
在 $X$ 中开。态射 $s:\mathcal{O}_{X_s}\to\mathcal{L}|_{X_s}$
是同构，且存在 $\mathcal{L}^{\otimes-1}$ 在 $X_s$ 上的截面 $s'$，
使得 $s'(s|_{X_s})=1$。
\end{lemma}

\begin{proof}
设 $x\in X_s$。由引理 \ref{modules-lemma-constructions-invertible}，有同构
$$
\mathcal{L}_x \otimes_{\mathcal{O}_{X, x}} (\mathcal{L}^{\otimes -1})_x
\longrightarrow
\mathcal{O}_{X, x}
$$
。$\mathcal{L}_x$ 与 $(\mathcal{L}^{\otimes-1})_x$
都是秩 $1$ 自由 $\mathcal{O}_{X,x}$-模。由《代数》中的
Nakayama 引理 \ref{algebra-lemma-NAK}，可知 $s_x$
是 $\mathcal{L}_x$ 的一个基。因此存在基元
$t_x\in(\mathcal{L}^{\otimes-1})_x$，使得
$s_x\otimes t_x$ 映到 $1$。选取 $x$ 的开邻域 $U$，
使 $t_x$ 来自 $\mathcal{L}^{\otimes-1}$ 在 $U$ 上的截面 $t$，
且 $s\otimes t$ 映到 $1\in\mathcal{O}_X(U)$。
显然，对每个 $x'\in U$，$s$ 都生成模 $\mathcal{L}_{x'}$。
故 $U\subset X_s$，这证明 $X_s$ 为开集。此外，上面在 $U$
上构造的截面 $t$ 是唯一的，因而它们粘合成引理中的截面 $s'$。
\end{proof}

\begin{remark}
\label{modules-remark-pullback-s-open}
给定局部环空间态射 $f:Y\to X$（参见《概形》定义
\ref{schemes-definition-locally-ringed-space}），逆像
$f^{-1}(X_s)$ 等于 $Y_{f^*s}$，其中
$f^*s\in\Gamma(Y,f^*\mathcal{L})$ 是 $s$ 的拉回。
\end{remark}


\section{秩与行列式}
\label{modules-section-rank-and-det}

\noindent
设 $(X,\mathcal{O}_X)$ 为环空间。考虑有限局部自由
$\mathcal{O}_X$-模层的范畴 $\textit{Vect}(X)$。这是正合范畴
（参见《内射对象》注 \ref{injectives-remark-embed-exact-category}），
其可容许满态射为满射，而可容许单态射为满射的核。此外，
$\textit{Vect}(X)$ 的对象同构类构成集合（略去证明）。
因此可以构造第零 Grothendieck $K$-群
$K_0(\textit{Vect}(X))$。具体地，此时
$K_0(\textit{Vect}(X))$ 是由有限局部自由
$\mathcal{O}_X$-模层 $\mathcal{E}$ 所对应的 $[\mathcal{E}]$
生成、并服从如下关系的阿贝尔群：
$$
[\mathcal{E}] = [\mathcal{E}'] + [\mathcal{E}'']
$$
，只要存在有限局部自由 $\mathcal{O}_X$-模层的短正合列
$0 \to \mathcal{E}' \to \mathcal{E} \to \mathcal{E}'' \to 0$
。

\medskip\noindent
{\bf 秩。}假设所有茎 $\mathcal{O}_{X,x}$ 都是非零环。
给定有限局部自由 $\mathcal{O}_X$-模层 $\mathcal{E}$，
其{\it 秩}是局部常值函数
$$
\text{rank}_\mathcal{E} : X \longrightarrow \mathbf{Z}_{\geq 0},\quad
x \longmapsto \text{rank}_{\mathcal{O}_{X, x}} \mathcal{E}_x
$$
，参见引理 \ref{modules-lemma-rank}。由局部自由模层的定义，
函数 $\text{rank}_\mathcal{E}$ 局部常值。若
$0 \to \mathcal{E}' \to \mathcal{E} \to \mathcal{E}'' \to 0$
是有限局部自由 $\mathcal{O}_X$-模层的短正合列，则
$\text{rank}_\mathcal{E} = \text{rank}_{\mathcal{E}'} +
\text{rank}_{\mathcal{E}''}$,
。因此秩定义同态
$$
K_0(\textit{Vect}(X)) \longrightarrow \text{Map}_{cont}(X, \mathbf{Z}),\quad
[\mathcal{E}] \longmapsto \text{rank}_\mathcal{E}
$$

\medskip\noindent
{\bf 行列式。}给定有限局部自由 $\mathcal{O}_X$-模层
$\mathcal{E}$，得到不交并分解
$$
X = X_0 \amalg X_1 \amalg X_2 \amalg \ldots
$$
，其中 $X_i$ 开且闭，而 $\mathcal{E}$ 在 $X_i$ 上为秩 $i$
的有限局部自由模层（这恰好等价于
$\text{rank}_\mathcal{E}$ 局部常值）。在此情形中，
定义 $\det(\mathcal{E})$ 为 $X$ 上如下可逆层：对每个
$i\geq0$，它在 $X_i$ 上等于 $\wedge^i(\mathcal{E}|_{X_i})$。
由于上述分解不交，无需检验任何粘合条件。由下面的引理
\ref{modules-lemma-det-ses}，这定义阿贝尔群同态
$$
\det : K_0(\textit{Vect}(X)) \longrightarrow \Pic(X),\quad
[\mathcal{E}] \longmapsto \det(\mathcal{E})
$$
。由此得到的 $\Pic(X)$ 中元素都是秩 $1$ 局部自由的
（推广见引理之后）。

\begin{lemma}
\label{modules-lemma-det-ses}
设 $X$ 为环空间。设
$0 \to \mathcal{E}' \to \mathcal{E} \to \mathcal{E}'' \to 0$
为有限局部自由 $\mathcal{O}_X$-模层的短正合列。则存在典范同构
$$
\det(\mathcal{E}') \otimes_{\mathcal{O}_X}\det(\mathcal{E}'')
\longrightarrow
\det(\mathcal{E})
$$
，它是 $\mathcal{O}_X$-模层同构。
\end{lemma}

\begin{proof}
可把 $X$ 分解为不交的开闭子集，使得 $\mathcal{E}'$
与 $\mathcal{E}''$ 在每个子集上都具有常秩。于是可归约到
$\mathcal{E}'$ 与 $\mathcal{E}''$ 分别具有常秩
$r'$ 与 $r''$ 的情形。在此情形中如下定义
$$
\wedge^{r'}(\mathcal{E}') \otimes_{\mathcal{O}_X} \wedge^{r''}(\mathcal{E}'')
\longrightarrow
\wedge^{r' + r''}(\mathcal{E})
$$
。给定 $\mathcal{E}'$ 的局部截面 $s'_1,\ldots,s'_{r'}$
及 $\mathcal{E}''$ 的局部截面 $s''_1,\ldots,s''_{r''}$，
我们把
$$
s'_1 \wedge \ldots \wedge s'_{r'} \otimes
s''_1 \wedge \ldots \wedge s''_{r''}
\quad\text{映到}\quad
s'_1 \wedge \ldots \wedge s'_{r'} \wedge
\tilde s''_1 \wedge \ldots \wedge \tilde s''_{r''}
$$
，其中 $\tilde s''_i$ 是截面 $s''_i$ 到 $\mathcal{E}$
的某个局部提升。我们略去细节。
\end{proof}

\noindent
设 $(X,\mathcal{O}_X)$ 为环空间。除了有限局部自由
$\mathcal{O}_X$-模层，还可以考察如下 $\mathcal{O}_X$-模层
$\mathcal{F}$：它在 $X$ 上局部为有限自由
$\mathcal{O}_X$-模层的直和项。这等价于要求 $\mathcal{F}$
是有限表示的平坦 $\mathcal{O}_X$-模层，参见引理
\ref{modules-lemma-flat-locally-finite-presentation}。若所有茎
$\mathcal{O}_{X,x}$ 都是局部环，则这样的 $\mathcal{F}$
有限局部自由，参见引理
\ref{modules-lemma-direct-summand-of-locally-free-is-locally-free}。
然而一般情形并非如此；例如，$X$ 可以是一个点，
$\Gamma(X,\mathcal{O}_X)$ 可以是两个非零环的积 $A\times B$，
而 $\mathcal{F}$ 可以对应于 $A\times0$。因此对这样的模层，
秩函数没有定义。不过仍可定义 $\det(\mathcal{F})$，而且它是定义
\ref{modules-definition-invertible} 意义下的可逆 $\mathcal{O}_X$-模层
（未必是秩 $1$ 局部自由模层）。当 $\mathcal{F}$ 有限局部自由时，
我们的构造与上述构造一致。我们建议读者跳过本节余下部分。

\begin{lemma}
\label{modules-lemma-determinant-as-socle}
设 $(X,\mathcal{O}_X)$ 为环空间，$\mathcal{F}$ 为平坦且有限表示的
$\mathcal{O}_X$-模层。记
$$
\det(\mathcal{F}) \subset
\wedge^*_{\mathcal{O}_X}(\mathcal{F})
$$
为 $\mathcal{F}\subset\wedge^*_{\mathcal{O}_X}(\mathcal{F})$ 的零化子。
则 $\det(\mathcal{F})$ 是可逆 $\mathcal{O}_X$-模层。
\end{lemma}

\begin{proof}
为证明此结论，可以在 $X$ 上局部工作。因此可以假设
$\mathcal{F}$ 是有限自由模层的直和项，参见引理
\ref{modules-lemma-flat-locally-finite-presentation}。设
$\mathcal{F} \oplus \mathcal{G} = \mathcal{O}_X^{\oplus n}$.
置 $R=\mathcal{O}_X(X)$。于是
$\mathcal{F}(X) \oplus \mathcal{G}(X) = R^{\oplus n}$
，相应地，对所有开集 $U\subset X$ 有
$\mathcal{F}(U) \oplus \mathcal{G}(U) = \mathcal{O}_X(U)^{\oplus n}$
。可知 $\mathcal{F}=\mathcal{F}_M$，如引理
\ref{modules-lemma-construct-quasi-coherent-sheaves} 所述，其中
$M=\mathcal{F}(X)$ 是有限投射 $R$-模。换言之，
$\mathcal{F}(U)=M\otimes_R\mathcal{O}_X(U)$。这蕴含
$\det(M) \otimes_R \mathcal{O}_X(U) = \det(\mathcal{F}(U))$
对所有开集 $U\subset X$ 成立，其中 $\det$ 如《代数详论》第
\ref{more-algebra-section-determinants} 节所述。由《代数详论》注
\ref{more-algebra-remark-determinant-as-socle}，可见
$$
\det(M) \otimes_R \mathcal{O}_X(U) =
\det(\mathcal{F}(U)) \subset
\wedge^*_{\mathcal{O}_X(U)}(\mathcal{F}(U))
$$
是 $\mathcal{F}(U)$ 的零化子。可知引理陈述中所定义的
$\det(\mathcal{F})$ 等于 $\mathcal{F}_{\det(M)}$。
略去一些细节；必须注意，不能把零化子定义为“在开集上对截面取零化子”
所得预层的层化。因此 $\det(\mathcal{F})$ 是一个可逆模的拉回，故得结论。
\end{proof}
















\section{环层的局部化}
\label{modules-section-localizing-sheaves-rings}

\noindent
设 $X$ 为拓扑空间，$\mathcal{O}_X$ 为环预层，
$\mathcal{S}\subset\mathcal{O}_X$ 为包含于 $\mathcal{O}_X$
的集合预层。假设对每个开集 $U\subset X$，集合
$\mathcal{S}(U)\subset\mathcal{O}_X(U)$ 都是乘法子集，参见
《代数》定义 \ref{algebra-definition-multiplicative-subset}。
此时可以考虑环预层
$$
\mathcal{S}^{-1}\mathcal{O}_X :
U \longmapsto \mathcal{S}(U)^{-1}\mathcal{O}_X(U).
$$
若 $V\subset U$ 是 $X$ 的开集，则限制映射把截面
$f/s$（$f\in\mathcal{O}_X(U)$、$s\in\mathcal{S}(U)$）
映到 $(f|_V)/(s|_V)$。

\begin{lemma}
\label{modules-lemma-simple-invert}
设 $X$ 为拓扑空间，$\mathcal{O}_X$ 为环预层，
$\mathcal{S}\subset\mathcal{O}_X$ 为包含于 $\mathcal{O}_X$
的集合预层。假设对每个开集 $U\subset X$，集合
$\mathcal{S}(U)\subset\mathcal{O}_X(U)$ 都是乘法子集。
\begin{enumerate}
\item 存在环预层态射
$\mathcal{O}_X \to \mathcal{S}^{-1}\mathcal{O}_X$
，使得 $\mathcal{S}$ 的每个局部截面都映到
$\mathcal{O}_X$ 的可逆截面。
\item 对任一环预层同态 $\mathcal{O}_X\to\mathcal{A}$，
若 $\mathcal{S}$ 的每个局部截面都映到 $\mathcal{A}$ 的可逆截面，
则存在唯一分解
$\mathcal{S}^{-1}\mathcal{O}_X \to \mathcal{A}$.
\item 对任意 $x\in X$，有
$$
(\mathcal{S}^{-1}\mathcal{O}_X)_x = \mathcal{S}_x^{-1} \mathcal{O}_{X, x}.
$$
\item 层化 $(\mathcal{S}^{-1}\mathcal{O}_X)^\#$ 是环层，
并带有环层态射
$(\mathcal{O}_X)^\# \to (\mathcal{S}^{-1}\mathcal{O}_X)^\#$
；它对于 $(\mathcal{O}_X)^\#$ 到如下环层的态射具有泛性质：
$\mathcal{S}$ 的每个局部截面都映到可逆截面。
\item 对任意 $x\in X$，有
$$
(\mathcal{S}^{-1}\mathcal{O}_X)^\#_x = \mathcal{S}_x^{-1} \mathcal{O}_{X, x}.
$$
\end{enumerate}
\end{lemma}

\begin{proof}
略。
\end{proof}

\noindent
设 $X$ 为拓扑空间，$\mathcal{O}_X$ 为环预层，
$\mathcal{S}\subset\mathcal{O}_X$ 为包含于 $\mathcal{O}_X$
的集合预层。假设对每个开集 $U\subset X$，集合
$\mathcal{S}(U)\subset\mathcal{O}_X(U)$ 都是乘法子集。
设 $\mathcal{F}$ 为 $\mathcal{O}_X$-模预层。此时可以考虑
$\mathcal{S}^{-1}\mathcal{O}_X$-模预层
$$
\mathcal{S}^{-1}\mathcal{F} :
U \longmapsto \mathcal{S}(U)^{-1}\mathcal{F}(U).
$$
若 $V\subset U$ 是 $X$ 的开集，则限制映射把截面
$t/s$（$t\in\mathcal{F}(U)$、$s\in\mathcal{S}(U)$）
映到 $(t|_V)/(s|_V)$。

\begin{lemma}
\label{modules-lemma-simple-invert-module}
设 $X$ 为拓扑空间，$\mathcal{O}_X$ 为环预层，
$\mathcal{S}\subset\mathcal{O}_X$ 为包含于 $\mathcal{O}_X$
的集合预层。假设对每个开集 $U\subset X$，集合
$\mathcal{S}(U)\subset\mathcal{O}_X(U)$ 都是乘法子集。
对任一 $\mathcal{O}_X$-模预层 $\mathcal{F}$，有
$$
\mathcal{S}^{-1}\mathcal{F}
=
\mathcal{S}^{-1}\mathcal{O}_X \otimes_{p, \mathcal{O}_X} \mathcal{F}
$$
（记法见《层》第 \ref{sheaves-section-presheaves-modules} 节）；
若 $\mathcal{F}$ 与 $\mathcal{O}_X$ 都是层，则
$$
(\mathcal{S}^{-1}\mathcal{F})^\#
=
(\mathcal{S}^{-1}\mathcal{O}_X)^\# \otimes_{\mathcal{O}_X} \mathcal{F}
$$
（记法见《层》第
\ref{sheaves-section-sheafification-presheaves-modules} 节）。
\end{lemma}

\begin{proof}
略。
\end{proof}








\section{微分模}
\label{modules-section-differentials}

\noindent
本节简述如何为环空间态射定义相对微分模。建议读者先看交换代数章节中的
相应部分（《代数》第 \ref{algebra-section-differentials} 节）。

\begin{definition}
\label{modules-definition-derivation}
设 $X$ 为拓扑空间，$\varphi:\mathcal{O}_1\to\mathcal{O}_2$
为环层同态，$\mathcal{F}$ 为 $\mathcal{O}_2$-模层。
取值于 $\mathcal{F}$ 的一个 {\it $\mathcal{O}_1$-导子}，
更准确地说一个 {\it $\varphi$-导子}，是满足如下条件的态射
$D:\mathcal{O}_2\to\mathcal{F}$：它可加，零化
$\mathcal{O}_1\to\mathcal{O}_2$ 的像，并满足 {\it Leibniz 法则}
$$
D(ab) = aD(b) + D(a)b
$$
，对 $\mathcal{O}_2$ 的所有局部截面 $a,b$（在二者同时有定义之处）
都成立。以
$\text{Der}_{\mathcal{O}_1}(\mathcal{O}_2,\mathcal{F})$
表示取值于 $\mathcal{F}$ 的 $\varphi$-导子之集。
\end{definition}

\noindent
这是《代数》定义 \ref{algebra-definition-derivation}
的层论类似物。给定定义中所述的导子
$D:\mathcal{O}_2\to\mathcal{F}$，整体截面上的态射
$$
D : \Gamma(X, \mathcal{O}_2) \longrightarrow \Gamma(X, \mathcal{F})
$$
是代数定义意义下的 $\Gamma(X,\mathcal{O}_1)$-导子。
注意，若 $\alpha:\mathcal{F}\to\mathcal{G}$ 为
$\mathcal{O}_2$-模层态射，则存在诱导态射
$$
\text{Der}_{\mathcal{O}_1}(\mathcal{O}_2, \mathcal{F})
\longrightarrow
\text{Der}_{\mathcal{O}_1}(\mathcal{O}_2, \mathcal{G})
$$
，由规则 $D\mapsto\alpha\circ D$ 给出。换言之，我们得到一个函子。

\begin{lemma}
\label{modules-lemma-universal-module}
设 $X$ 为拓扑空间，$\varphi:\mathcal{O}_1\to\mathcal{O}_2$
为环层同态。函子
$$
\textit{Mod}(\mathcal{O}_2) \longrightarrow \textit{Ab}, \quad
\mathcal{F} \longmapsto \text{Der}_{\mathcal{O}_1}(\mathcal{O}_2, \mathcal{F})
$$
可表。
\end{lemma}

\begin{proof}
证明与代数中的类似陈述完全相同。在本证明中，对 $X$ 上任一集合层
$\mathcal{F}$，以 $\mathcal{O}_2[\mathcal{F}]$ 表示预层
$U\mapsto\mathcal{O}_2(U)[\mathcal{F}(U)]$ 的层化，其中后者表示以集合
$\mathcal{F}(U)$ 为基的自由 $\mathcal{O}_2(U)$-模。
对 $s\in\mathcal{F}(U)$，以 $[s]$ 表示
$\mathcal{O}_2[\mathcal{F}]$ 在 $U$ 上的相应截面。
若 $\mathcal{F}$ 是 $\mathcal{O}_2$-模层，则存在典范态射
$$
c : \mathcal{O}_2[\mathcal{F}] \longrightarrow \mathcal{F}
$$
，它在预层层面由规则
$\sum f_s[s]\mapsto\sum f_s s$ 给出。我们将用简写
$[s]\mapsto s$ 描述此态射，下面的其他态射也类似。考虑
$\mathcal{O}_2$-模层态射
\begin{equation}
\label{modules-equation-define-module-differentials}
\begin{matrix}
\mathcal{O}_2[\mathcal{O}_2 \times \mathcal{O}_2] \oplus
\mathcal{O}_2[\mathcal{O}_2 \times \mathcal{O}_2] \oplus
\mathcal{O}_2[\mathcal{O}_1] &
\longrightarrow &
\mathcal{O}_2[\mathcal{O}_2] \\
[(a, b)] \oplus [(f, g)] \oplus [h] & \longmapsto & [a + b] - [a] - [b] + \\
& & [fg] - g[f] - f[g] + \\
& & [\varphi(h)]
\end{matrix}
\end{equation}
，这里使用上述简写。置 $\Omega_{\mathcal{O}_2/\mathcal{O}_1}$
等于此态射的余核。显然存在集合层态射
$$
\text{d} : \mathcal{O}_2 \longrightarrow \Omega_{\mathcal{O}_2/\mathcal{O}_1}
$$
，它把局部截面 $f$ 映到 $[f]$ 在
$\Omega_{\mathcal{O}_2/\mathcal{O}_1}$ 中的像。按照构造，
$\text{d}$ 是 $\mathcal{O}_1$-导子。接着，设 $\mathcal{F}$
为 $\mathcal{O}_2$-模层，$D:\mathcal{O}_2\to\mathcal{F}$
为 $\mathcal{O}_1$-导子。考虑把 $[g]$ 映到 $D(g)$ 的
$\mathcal{O}_2$-线性态射
$\mathcal{O}_2[\mathcal{O}_2]\to\mathcal{F}$。
由导子的定义，此态射零化态射
(\ref{modules-equation-define-module-differentials}) 像中的截面，因而定义态射
$$
\alpha_D : \Omega_{\mathcal{O}_2/\mathcal{O}_1} \longrightarrow \mathcal{F}
$$
显然 $D=\alpha_D\circ\text{d}$，引理得证。
\end{proof}

\begin{definition}
\label{modules-definition-module-differentials}
设 $X$ 为拓扑空间，$\varphi:\mathcal{O}_1\to\mathcal{O}_2$
为 $X$ 上的环层同态。$\varphi$ 的{\it 微分模}是表示函子
$\mathcal{F} \mapsto \text{Der}_{\mathcal{O}_1}(\mathcal{O}_2, \mathcal{F})$
的对象；其存在性由引理 \ref{modules-lemma-universal-module} 保证。
它记作 $\Omega_{\mathcal{O}_2/\mathcal{O}_1}$，而{\it 泛
$\varphi$-导子}记作
$\text{d} : \mathcal{O}_2 \to \Omega_{\mathcal{O}_2/\mathcal{O}_1}$.
\end{definition}

\noindent
注意 $\Omega_{\mathcal{O}_2/\mathcal{O}_1}$ 是
$\mathcal{O}_2$-模层态射
(\ref{modules-equation-define-module-differentials}) 的余核。此外，态射
$\text{d}$ 由如下规则描述：$\text{d}f$ 是局部截面 $[f]$ 的像。

\begin{lemma}
\label{modules-lemma-differentials-sheafify}
设 $X$ 为拓扑空间，$\varphi:\mathcal{O}_1\to\mathcal{O}_2$
为 $X$ 上的环层同态。则
$\Omega_{\mathcal{O}_2/\mathcal{O}_1}$ 是与预层
$U\mapsto\Omega_{\mathcal{O}_2(U)/\mathcal{O}_1(U)}$ 关联的层。
\end{lemma}

\begin{proof}
考虑态射 (\ref{modules-equation-define-module-differentials})。
存在类似的预层态射，其在开集 $U$ 上的值为
$$
\mathcal{O}_2(U)[\mathcal{O}_2(U) \times \mathcal{O}_2(U)] \oplus
\mathcal{O}_2(U)[\mathcal{O}_2(U) \times \mathcal{O}_2(U)] \oplus
\mathcal{O}_2(U)[\mathcal{O}_1(U)]
\longrightarrow
\mathcal{O}_2(U)[\mathcal{O}_2(U)]
$$
由《代数》定义 \ref{algebra-definition-differentials}
中微分模的构造，此态射的余核在 $U$ 上的值为
$\Omega_{\mathcal{O}_2(U)/\mathcal{O}_1(U)}$。另一方面，
(\ref{modules-equation-define-module-differentials}) 中的各层是上述预层的层化。
由于层化正合，结论随即成立。
\end{proof}

\begin{lemma}
\label{modules-lemma-localize-differentials}
设 $X$ 为拓扑空间，$\varphi:\mathcal{O}_1\to\mathcal{O}_2$
为环层同态。对开集 $U\subset X$，存在典范同构
$$
\Omega_{\mathcal{O}_2/\mathcal{O}_1}|_U =
\Omega_{(\mathcal{O}_2|_U)/(\mathcal{O}_1|_U)}
$$
，它与泛导子相容。
\end{lemma}

\begin{proof}
这是因为 $\Omega_{\mathcal{O}_2/\mathcal{O}_1}$
是态射 (\ref{modules-equation-define-module-differentials}) 的余核。
\end{proof}

\begin{lemma}
\label{modules-lemma-pullback-differentials}
设 $f:Y\to X$ 为拓扑空间的连续映射，
$\varphi:\mathcal{O}_1\to\mathcal{O}_2$ 为 $X$ 上的环层同态。
则存在典范等同
$f^{-1}\Omega_{\mathcal{O}_2/\mathcal{O}_1} =
\Omega_{f^{-1}\mathcal{O}_2/f^{-1}\mathcal{O}_1}$
，它与泛导子相容。
\end{lemma}

\begin{proof}
这是因为层 $\Omega_{\mathcal{O}_2/\mathcal{O}_1}$
是态射 (\ref{modules-equation-define-module-differentials}) 的余核，
$\Omega_{f^{-1}\mathcal{O}_2/f^{-1}\mathcal{O}_1}$ 也有类似陈述，
而函子 $f^{-1}$ 正合，并且
$f^{-1}(\mathcal{O}_2[\mathcal{O}_2]) =
f^{-1}\mathcal{O}_2[f^{-1}\mathcal{O}_2]$,
$f^{-1}(\mathcal{O}_2[\mathcal{O}_2 \times \mathcal{O}_2]) =
f^{-1}\mathcal{O}_2[f^{-1}\mathcal{O}_2 \times f^{-1}\mathcal{O}_2]$，且
$f^{-1}(\mathcal{O}_2[\mathcal{O}_1]) =
f^{-1}\mathcal{O}_2[f^{-1}\mathcal{O}_1]$.
\end{proof}

\begin{lemma}
\label{modules-lemma-stalk-module-differentials}
设 $X$ 为拓扑空间，$\mathcal{O}_1\to\mathcal{O}_2$
为 $X$ 上的环层同态，$x\in X$。则有
$\Omega_{\mathcal{O}_2/\mathcal{O}_1, x} =
\Omega_{\mathcal{O}_{2, x}/\mathcal{O}_{1, x}}$.
\end{lemma}

\begin{proof}
这是把引理 \ref{modules-lemma-pullback-differentials} 应用于包含映射
$\{x\}\to X$ 所得的特殊情形。另一种证明使用引理
\ref{modules-lemma-differentials-sheafify}、《层》引理
\ref{sheaves-lemma-stalk-sheafification} 及《代数》引理
\ref{algebra-lemma-colimit-differentials}。
\end{proof}

\begin{lemma}
\label{modules-lemma-functoriality-differentials}
设 $X$ 为拓扑空间。设
$$
\xymatrix{
\mathcal{O}_2 \ar[r]_\varphi & \mathcal{O}_2' \\
\mathcal{O}_1 \ar[r] \ar[u] & \mathcal{O}'_1 \ar[u]
}
$$
为 $X$ 上环层的交换图。态射
$\mathcal{O}_2 \to \mathcal{O}'_2$ 与态射
$\text{d} : \mathcal{O}'_2 \to \Omega_{\mathcal{O}'_2/\mathcal{O}'_1}$
的复合是 $\mathcal{O}_1$-导子。因此得到典范的
$\mathcal{O}_2$-模层态射
$\Omega_{\mathcal{O}_2/\mathcal{O}_1} \to
\Omega_{\mathcal{O}'_2/\mathcal{O}'_1}$.
。它由下述性质唯一刻画：
$\text{d}(f) \mapsto \text{d}(\varphi(f))$
对 $\mathcal{O}_2$ 的任一局部截面 $f$ 成立。这样，
$\Omega_{-/-}$ 成为环层箭头范畴上的函子。
\end{lemma}

\begin{proof}
本引理由其表述即得。
\end{proof}

\begin{lemma}
\label{modules-lemma-differential-seq}
在引理 \ref{modules-lemma-functoriality-differentials} 中，假设
$\mathcal{O}_2\to\mathcal{O}'_2$ 为满射，其核为
$\mathcal{I}\subset\mathcal{O}_2$，并假设
$\mathcal{O}_1=\mathcal{O}'_1$。则存在典范的
$\mathcal{O}'_2$-模层正合列
$$
\mathcal{I}/\mathcal{I}^2
\longrightarrow
\Omega_{\mathcal{O}_2/\mathcal{O}_1} \otimes_{\mathcal{O}_2} \mathcal{O}'_2
\longrightarrow
\Omega_{\mathcal{O}'_2/\mathcal{O}_1}
\longrightarrow
0
$$
最左侧态射由如下规则刻画：$\mathcal{I}$ 的局部截面 $f$
映到 $\text{d}f\otimes1$。
\end{lemma}

\begin{proof}
对 $\mathcal{I}$ 的局部截面 $f$，以 $\overline f$ 表示
$f$ 在 $\mathcal{I}/\mathcal{I}^2$ 中的像。为证明态射
$\overline f\mapsto\text{d}f\otimes1$ 良定义，只需检验：
若 $f_1,f_2$ 是 $\mathcal{I}$ 的局部截面，则
$\text{d}f_1f_2\otimes1=0$。这由 Leibniz 法则立即得到：
$\text{d} f_1f_2 \otimes 1 =
(f_1 \text{d}f_2 + f_2 \text{d} f_1 )\otimes 1 =
\text{d}f_2 \otimes f_1 + \text{d}f_1 \otimes f_2 = 0$.
类似计算表明，此态射是
$\mathcal{O}'_2=\mathcal{O}_2/\mathcal{I}$-线性的。右侧态射来自引理
\ref{modules-lemma-functoriality-differentials}。为证明该序列正合，
可在茎上检验（引理 \ref{modules-lemma-abelian}）。由引理
\ref{modules-lemma-stalk-module-differentials}，这归结为《代数》引理
\ref{algebra-lemma-differential-seq}。
\end{proof}

\begin{definition}
\label{modules-definition-differentials}
设 $(f,f^\sharp):(X,\mathcal{O}_X)\to(S,\mathcal{O}_S)$
为环空间态射。
\begin{enumerate}
\item 设 $\mathcal{F}$ 为 $\mathcal{O}_X$-模层。取值于
$\mathcal{F}$ 的 {\it $S$-导子}是
$f^{-1}\mathcal{O}_S$-导子，更准确地说，是定义
\ref{modules-definition-derivation} 意义下的 $f^\sharp$-导子。
以 $\text{Der}_S(\mathcal{O}_X,\mathcal{F})$
表示取值于 $\mathcal{F}$ 的 $S$-导子之集。
\item {\it $X$ 在 $S$ 上的微分层 $\Omega_{X/S}$} 是带有其泛
$S$-导子 $\text{d}_{X/S}:\mathcal{O}_X\to\Omega_{X/S}$
的微分模 $\Omega_{\mathcal{O}_X/f^{-1}\mathcal{O}_S}$。
\end{enumerate}
\end{definition}

\noindent
下面给出导子自然出现的一种特殊情形。

\begin{lemma}
\label{modules-lemma-double-structure-gives-derivation}
设 $(f,f^\sharp):(X,\mathcal{O}_X)\to(S,\mathcal{O}_S)$
为环空间态射。考虑短正合列
$$
0 \to \mathcal{I} \to \mathcal{A} \to \mathcal{O}_X \to 0
$$
这里 $\mathcal{A}$ 是 $f^{-1}\mathcal{O}_S$-代数层，
$\pi:\mathcal{A}\to\mathcal{O}_X$ 是
$f^{-1}\mathcal{O}_S$-代数层的满射，而
$\mathcal{I}=\Ker(\pi)$ 是其核。假设 $\mathcal{I}$
是 $\mathcal{A}$ 中平方为零的理想层。因此 $\mathcal{I}$
自然具有 $\mathcal{O}_X$-模层结构。$\pi$ 的截面
$s:\mathcal{O}_X\to\mathcal{A}$ 是满足
$\pi\circ s=\text{id}$ 的 $f^{-1}\mathcal{O}_S$-代数层态射。
给定 $\pi$ 的任一截面 $s:\mathcal{O}_X\to\mathcal{A}$
与任一 $S$-导子 $D:\mathcal{O}_X\to\mathcal{I}$，态射
$$
s + D : \mathcal{O}_X \to \mathcal{A}
$$
是 $\pi$ 的截面，并且每个截面 $s'$ 都能唯一地写成
$s+D$，其中 $D$ 是 $S$-导子。
\end{lemma}

\begin{proof}
回顾 $\mathcal{I}$ 上的 $\mathcal{O}_X$-模层结构由
$h\tau=\tilde h\tau$（在 $\mathcal{A}$ 中相乘）给出，其中
$h$ 是 $\mathcal{O}_X$ 的局部截面，$\tilde h$ 是 $h$
到 $\mathcal{A}$ 的局部截面的某个局部提升，而 $\tau$
是 $\mathcal{I}$ 的局部截面。特别地，给定 $s$，可取
$\tilde h=s(h)$。为验证 $s+D$ 是环层同态，计算
\begin{eqnarray*}
(s + D)(ab) & = & s(ab) + D(ab) \\
& = & s(a)s(b) + aD(b) + D(a)b \\
& = & s(a) s(b) + s(a)D(b) + D(a)s(b) \\
& = & (s(a) + D(a))(s(b) + D(b))
\end{eqnarray*}
，这里使用了 Leibniz 法则。以同样方式，由于 $D$ 是
$S$-导子，可证明 $s+D$ 是 $f^{-1}\mathcal{O}_S$-代数层态射。
反之，给定 $s'$，置 $D=s'-s$。略去细节。
\end{proof}

\begin{lemma}
\label{modules-lemma-functoriality-differentials-ringed-spaces}
设
$$
\xymatrix{
X' \ar[d]_{h'} \ar[r]_f & X \ar[d]^h \\
S' \ar[r]^g & S
}
$$
为环空间的交换图。
\begin{enumerate}
\item 典范态射 $\mathcal{O}_X\to f_*\mathcal{O}_{X'}$ 与
$f_*\text{d}_{X'/S'}:f_*\mathcal{O}_{X'}\to f_*\Omega_{X'/S'}$
的复合是 $S$-导子，并由此得到典范的 $\mathcal{O}_X$-模层态射
$\Omega_{X/S} \to f_*\Omega_{X'/S'}$.
\item 交换图
$$
\xymatrix{
f^{-1}\mathcal{O}_X \ar[r] & \mathcal{O}_{X'} \\
f^{-1}h^{-1}\mathcal{O}_S \ar[u] \ar[r] & (h')^{-1}\mathcal{O}_{S'} \ar[u]
}
$$
由引理 \ref{modules-lemma-pullback-differentials} 与
\ref{modules-lemma-functoriality-differentials} 诱导典范态射
$f^{-1}\Omega_{X/S}\to\Omega_{X'/S'}$。
\end{enumerate}
这两个态射（通过 $f_*$ 与 $f^*$ 的伴随性，并通过
$f^*\Omega_{X/S} =
f^{-1}\Omega_{X/S} \otimes_{f^{-1}\mathcal{O}_X} \mathcal{O}_{X'}$ 以及
《层》引理 \ref{sheaves-lemma-adjointness-tensor-restrict}）
对应于同一个 $\mathcal{O}_{X'}$-模层同态
$$
c_f : f^*\Omega_{X/S} \longrightarrow \Omega_{X'/S'}
$$
，它由如下性质唯一刻画：对 $\mathcal{O}_X$ 的任一局部截面
$a$，$f^*\text{d}_{X/S}(a)$ 映到
$\text{d}_{X'/S'}(f^*a)$。
\end{lemma}

\begin{proof}
略。
\end{proof}

\begin{lemma}
\label{modules-lemma-check-functoriality-differentials}
设
$$
\xymatrix{
X'' \ar[d] \ar[r]_g & X' \ar[d] \ar[r]_f & X \ar[d] \\
S'' \ar[r] & S' \ar[r] & S
}
$$
为环空间的交换图。使用引理
\ref{modules-lemma-functoriality-differentials-ringed-spaces} 中的记法，有
$$
c_{f \circ g} = c_g \circ g^* c_f
$$
，这是从 $(f\circ g)^*\Omega_{X/S}$ 到 $\Omega_{X''/S''}$ 的态射等式。
\end{lemma}

\begin{proof}
略。
\end{proof}

\begin{lemma}
\label{modules-lemma-triangle-differentials}
设 $f:X\to Y$、$g:Y\to S$ 为环空间态射。则存在典范正合列
$$
f^*\Omega_{Y/S} \to \Omega_{X/S} \to \Omega_{X/Y} \to 0
$$
，其中各态射由应用引理
\ref{modules-lemma-functoriality-differentials-ringed-spaces} 得到。
\end{lemma}

\begin{proof}
在 $x\in X$ 处取诱导的茎态射，并使用引理
\ref{modules-lemma-stalk-module-differentials}，得到序列
$$
\mathcal{O}_{X, x} \otimes_{\mathcal{O}_{Y, f(x)}}
\Omega_{\mathcal{O}_{Y, f(x)}/\mathcal{O}_{S, g(f(x))}} \to
\Omega_{\mathcal{O}_{X, x}/\mathcal{O}_{S, g(f(x))}} \to
\Omega_{\mathcal{O}_{X, x}/\mathcal{O}_{Y, f(x)}} \to 0
$$
只需说明该序列中的态射与《代数》引理
\ref{algebra-lemma-exact-sequence-differentials} 中的态射相同。
这是因为引理 \ref{modules-lemma-functoriality-differentials-ringed-spaces}
对这些态射的刻画，并且通过《层》引理
\ref{sheaves-lemma-stalk-pullback-modules} 的同构，对任一
$\mathcal{O}_Y$-模层的局部截面 $s$，有
$(f^*s)_x=s_x\otimes1$。
\end{proof}


















\section{有限阶微分算子}
\label{modules-section-differential-operators}

\noindent
本节介绍有限阶微分算子。建议读者先看交换代数章节中的相应部分
（《代数》第 \ref{algebra-section-differential-operators} 节）。

\begin{definition}
\label{modules-definition-differential-operators}
设 $X$ 为拓扑空间，$\varphi:\mathcal{O}_1\to\mathcal{O}_2$
为 $X$ 上的环层同态，$k\geq0$ 为整数，而
$\mathcal{F},\mathcal{G}$ 为 $\mathcal{O}_2$-模层。
{\it $k$ 阶微分算子 $D:\mathcal{F}\to\mathcal{G}$}是满足如下条件的
$\mathcal{O}_1$-线性态射：对 $\mathcal{O}_2$ 的所有局部截面 $g$，
态射 $s\mapsto D(gs)-gD(s)$ 是 $k-1$ 阶微分算子。
对于起始情形 $k=0$，定义 $0$ 阶微分算子为
$\mathcal{O}_2$-线性态射。
\end{definition}

\noindent
若 $D:\mathcal{F}\to\mathcal{G}$ 是 $k$ 阶微分算子，则对
$\mathcal{O}_2$ 的所有局部截面 $g$，态射 $gD$ 也是 $k$ 阶微分算子。
两个 $k$ 阶微分算子之和仍是 $k$ 阶微分算子。因此所有这类算子之集
$$
\text{Diff}^k(\mathcal{F}, \mathcal{G}) =
\text{Diff}^k_{\mathcal{O}_2/\mathcal{O}_1}(\mathcal{F}, \mathcal{G})
$$
是 $\Gamma(X,\mathcal{O}_2)$-模。我们有
$$
\text{Diff}^0(\mathcal{F}, \mathcal{G}) \subset
\text{Diff}^1(\mathcal{F}, \mathcal{G}) \subset
\text{Diff}^2(\mathcal{F}, \mathcal{G}) \subset \ldots
$$
把开集 $U\subset X$ 映到 $k$ 阶微分算子
$D:\mathcal{F}|_U\to\mathcal{G}|_U$ 所成模的规则，
是 $X$ 上的 $\mathcal{O}_2$-模层。由此得到微分算子层
（若以后需要，我们会在此补上定义）。

\begin{lemma}
\label{modules-lemma-composition-differential-operators}
设 $X$ 为拓扑空间，$\mathcal{O}_1\to\mathcal{O}_2$
为 $X$ 上的环层态射，$\mathcal{E},\mathcal{F},\mathcal{G}$
为 $\mathcal{O}_2$-模层。若
$D:\mathcal{E}\to\mathcal{F}$ 与
$D':\mathcal{F}\to\mathcal{G}$ 分别为 $k$ 阶与 $k'$ 阶微分算子，
则 $D'\circ D$ 是 $k+k'$ 阶微分算子。
\end{lemma}

\begin{proof}
设 $g$ 为 $\mathcal{O}_2$ 的局部截面。于是把
$\mathcal{E}$ 的局部截面 $x$ 映到
$$
D'(D(gx)) - gD'(D(x)) = D'(D(gx)) - D'(gD(x)) + D'(gD(x)) - gD'(D(x))
$$
的态射，是两个较低阶微分算子的复合之和。因此对 $k+k'$ 归纳即得本引理。
\end{proof}

\begin{lemma}
\label{modules-lemma-module-principal-parts}
设 $X$ 为拓扑空间，$\mathcal{O}_1\to\mathcal{O}_2$
为 $X$ 上的环层态射，$\mathcal{F}$ 为 $\mathcal{O}_2$-模层，
$k\geq0$。存在 $\mathcal{O}_2$-模层
$\mathcal{P}^k_{\mathcal{O}_2/\mathcal{O}_1}(\mathcal{F})$
及典范同构
$$
\text{Diff}^k_{\mathcal{O}_2/\mathcal{O}_1}(\mathcal{F}, \mathcal{G}) =
\Hom_{\mathcal{O}_2}(
\mathcal{P}^k_{\mathcal{O}_2/\mathcal{O}_1}(\mathcal{F}), \mathcal{G})
$$
，它关于 $\mathcal{O}_2$-模层 $\mathcal{G}$ 具有函子性。
\end{lemma}

\begin{proof}
存在性由一般范畴论论证得到（此处以后补引用），但我们也给出直接构造，
因为此构造将在后续证明中有用。我们将自由使用引理
\ref{modules-lemma-universal-module} 的证明中引入的记法。
给定任一微分算子 $D:\mathcal{F}\to\mathcal{G}$，得到
$\mathcal{O}_2$-线性态射
$L_D : \mathcal{O}_2[\mathcal{F}] \to \mathcal{G}$
，把 $[m]$ 映到 $D(m)$。若 $D$ 为 $0$ 阶，则 $L_D$
零化局部截面
$$
[m + m'] - [m] - [m'],\quad
g_0[m] - [g_0m]
$$
，其中 $g_0$ 是 $\mathcal{O}_2$ 的局部截面，而 $m,m'$
是 $\mathcal{F}$ 的局部截面。若 $D$ 为 $1$ 阶，则 $L_D$
零化局部截面
$$
[m + m' - [m] - [m'],\quad
f[m] - [fm], \quad
g_0g_1[m] - g_0[g_1m] - g_1[g_0m] + [g_1g_0m]
$$
，其中 $f$ 是 $\mathcal{O}_1$ 的局部截面，
$g_0,g_1$ 是 $\mathcal{O}_2$ 的局部截面，而 $m,m'$
是 $\mathcal{F}$ 的局部截面。若 $D$ 为 $k$ 阶，则 $L_D$
零化局部截面 $[m+m']-[m]-[m']$、$f[m]-[fm]$ 以及局部截面
$$
g_0g_1\ldots g_k[m] - \sum g_0 \ldots \hat g_i \ldots g_k[g_im] + \ldots
+(-1)^{k + 1}[g_0\ldots g_km]
$$
反之，若 $L:\mathcal{O}_2[\mathcal{F}]\to\mathcal{G}$
是零化前一句所列全部局部截面的 $\mathcal{O}_2$-线性态射，
则 $m\mapsto L([m])$ 是 $k$ 阶微分算子。因此
$\mathcal{P}^k_{\mathcal{O}_2/\mathcal{O}_1}(\mathcal{F})$
是 $\mathcal{O}_2[\mathcal{F}]$ 关于由这些局部截面生成的
$\mathcal{O}_2$-子模层的商。
\end{proof}

\begin{definition}
\label{modules-definition-module-principal-parts}
设 $X$ 是一个拓扑空间，
$\mathcal{O}_1 \to \mathcal{O}_2$ 是 $X$ 上环层之间的一个同态，
$\mathcal{F}$ 是一个 $\mathcal{O}_2$-模层。
称引理 \ref{modules-lemma-module-principal-parts} 中构造的模层
$\mathcal{P}^k_{\mathcal{O}_2/\mathcal{O}_1}(\mathcal{F})$
为 $\mathcal{F}$ 的 {\it $k$ 阶主部模层}。
\end{definition}

\noindent
注意，包含关系
$$
\text{Diff}^0(\mathcal{F}, \mathcal{G}) \subset
\text{Diff}^1(\mathcal{F}, \mathcal{G}) \subset
\text{Diff}^2(\mathcal{F}, \mathcal{G}) \subset \ldots
$$
通过 Yoneda 引理（《范畴论》，引理 \ref{categories-lemma-yoneda}）
对应于满射
$$
\ldots \to \mathcal{P}^2_{\mathcal{O}_2/\mathcal{O}_1}(\mathcal{F})
\to \mathcal{P}^1_{\mathcal{O}_2/\mathcal{O}_1}(\mathcal{F})
\to \mathcal{P}^0_{\mathcal{O}_2/\mathcal{O}_1}(\mathcal{F}) = \mathcal{F}
$$

\begin{lemma}
\label{modules-lemma-differential-operators-sheafify}
设 $X$ 是一个拓扑空间，$\mathcal{O}_1 \to \mathcal{O}_2$
是 $X$ 上环预层之间的一个同态，$\mathcal{F}$ 是一个
$\mathcal{O}_2$-模预层。则
$\mathcal{P}^k_{\mathcal{O}_2^\#/\mathcal{O}_1^\#}(\mathcal{F}^\#)$
是预层
$U \mapsto P^k_{\mathcal{O}_2(U)/\mathcal{O}_1(U)}(\mathcal{F}(U))$.
的伴随层。
\end{lemma}

\begin{proof}
这可以完全仿照引理 \ref{modules-lemma-differentials-sheafify} 中微分层的情形证明。
也许更自然的办法是直接利用引理
\ref{modules-lemma-module-principal-parts} 的泛性质来得到该等式。
细节从略。
\end{proof}

\begin{lemma}
\label{modules-lemma-sequence-of-principal-parts}
设 $X$ 是一个拓扑空间，$\mathcal{O}_1 \to \mathcal{O}_2$
是 $X$ 上环层之间的一个同态，$\mathcal{F}$ 是一个
$\mathcal{O}_2$-模层。存在典范短正合列
$$
0 \to
\Omega_{\mathcal{O}_2/\mathcal{O}_1} \otimes_{\mathcal{O}_2} \mathcal{F} \to
\mathcal{P}^1_{\mathcal{O}_2/\mathcal{O}_1}(\mathcal{F}) \to
\mathcal{F} \to 0
$$
它对 $\mathcal{F}$ 具有函子性，称为{\it 主部序列}。
\end{lemma}

\begin{proof}
这由交换代数版本（《代数》，引理
\ref{algebra-lemma-sequence-of-principal-parts}）以及引理
\ref{modules-lemma-differentials-sheafify} 和
\ref{modules-lemma-differential-operators-sheafify} 推出。
\end{proof}

\begin{remark}
\label{modules-remark-functoriality-principal-parts}
设 $X$ 是一个拓扑空间，并给定 $X$ 上环层的交换图
$$
\xymatrix{
\mathcal{B} \ar[r] & \mathcal{B}' \\
\mathcal{A} \ar[u] \ar[r] & \mathcal{A}' \ar[u]
}
$$
以及一个 $\mathcal{B}$-模层 $\mathcal{F}$、
一个 $\mathcal{B}'$-模层 $\mathcal{F}'$ 和一个
$\mathcal{B}$-线性同态 $\mathcal{F} \to \mathcal{F}'$。
于是得到一族相容的模同态
$$
\xymatrix{
\ldots \ar[r] &
\mathcal{P}^2_{\mathcal{B}'/\mathcal{A}'}(\mathcal{F}') \ar[r] &
\mathcal{P}^1_{\mathcal{B}'/\mathcal{A}'}(\mathcal{F}') \ar[r] &
\mathcal{P}^0_{\mathcal{B}'/\mathcal{A}'}(\mathcal{F}') \\
\ldots \ar[r] &
\mathcal{P}^2_{\mathcal{B}/\mathcal{A}}(\mathcal{F}) \ar[r] \ar[u] &
\mathcal{P}^1_{\mathcal{B}/\mathcal{A}}(\mathcal{F}) \ar[r] \ar[u] &
\mathcal{P}^0_{\mathcal{B}/\mathcal{A}}(\mathcal{F}) \ar[u]
}
$$
这些同态与此类同态的进一步复合相容。最容易的验证方法是使用引理
\ref{modules-lemma-module-principal-parts} 的证明中以（局部）生成元和关系
给出的模层 $\mathcal{P}^k_{\mathcal{B}/\mathcal{A}}(\mathcal{M})$
的描述；也可以直接从这些模层的泛性质看出这一点。此外，这些同态与引理
\ref{modules-lemma-sequence-of-principal-parts} 的短正合列相容。
\end{remark}

\noindent
下面将定义推广到环空间的态射。

\begin{definition}
\label{modules-definition-relative-differential-operators}
设 $(f, f^\sharp) : (X, \mathcal{O}_X) \to (S, \mathcal{O}_S)$
是环空间的一个态射，$\mathcal{F}$ 和 $\mathcal{G}$ 是
$\mathcal{O}_X$-模层，$k \geq 0$ 是一个整数。
所谓 $X/S$ 上的一个{\it $k$ 阶微分算子}，是指关于
$f^\sharp : f^{-1}\mathcal{O}_S \to \mathcal{O}_X$
的微分算子 $D : \mathcal{F} \to \mathcal{G}$。
用 $\text{Diff}^k_{X/S}(\mathcal{F}, \mathcal{G})$
表示这些微分算子的集合。
\end{definition}












\section{de Rham 复形}
\label{modules-section-de-rham-complex}

\noindent
本节是《代数》第 \ref{algebra-section-de-rham-complex} 节对环空间态射的类似版本。
我们建议读者先阅读该节。

\medskip\noindent
设 $X$ 是一个拓扑空间，$\mathcal{A} \to \mathcal{B}$ 是环层的一个同态。
用
$\text{d} : \mathcal{B} \to \Omega_{\mathcal{B}/\mathcal{A}}$
表示第 \ref{modules-section-differentials} 节中构造的微分模层及其泛
$\mathcal{A}$-导子。对 $i \geq 0$，令
$$
\Omega_{\mathcal{B}/\mathcal{A}}^i =
\wedge^i_\mathcal{B}(\Omega_{\mathcal{B}/\mathcal{A}})
$$
为第 \ref{modules-section-symmetric-exterior} 节意义下的第 $i$ 个外幂。

\begin{definition}
\label{modules-definition-de-rham-complex}
在上述情形中，$\mathcal{B}$ 在 $\mathcal{A}$ 上的
{\it de Rham 复形}是唯一的 $\mathcal{A}$-模层复形
$$
\Omega_{\mathcal{B}/\mathcal{A}}^0 \to
\Omega_{\mathcal{B}/\mathcal{A}}^1 \to
\Omega_{\mathcal{B}/\mathcal{A}}^2 \to \ldots
$$
其 $0$ 次微分由
$\text{d} : \mathcal{B} \to \Omega_{\mathcal{B}/\mathcal{A}}$
给出，而更高次数的微分具有以下性质：
\begin{equation}
\label{modules-equation-rule}
\text{d}\left(b_0\text{d}b_1 \wedge \ldots \wedge \text{d}b_p\right) =
\text{d}b_0 \wedge \text{d}b_1 \wedge \ldots \wedge \text{d}b_p
\end{equation}
其中 $b_0, \ldots, b_p \in \mathcal{B}(U)$ 是同一个开集
$U \subset X$ 上的截面。
\end{definition}

\noindent
我们可以重复《代数》第 \ref{algebra-section-de-rham-complex} 节中繁复的论证来构造此复形。
不过，回忆 $\Omega_{\mathcal{B}/\mathcal{A}}$ 是预层
$U \mapsto \Omega_{\mathcal{B}(U)/\mathcal{A}(U)}$ 的层化，见引理
\ref{modules-lemma-differentials-sheafify}。因此，
$\Omega_{\mathcal{B}/\mathcal{A}}^i$ 是预层
$U \mapsto \Omega^i_{\mathcal{B}(U)/\mathcal{A}(U)}$ 的层化，见引理
\ref{modules-lemma-local-tensor-algebra}。所以，可以把 de Rham 复形定义为下述规则的层化：
$$
U \longmapsto
\Omega^\bullet_{\mathcal{B}(U)/\mathcal{A}(U)}
$$

\begin{lemma}
\label{modules-lemma-pullback-de-rham-complex}
设 $f : Y \to X$ 是拓扑空间的一个连续映射，
$\mathcal{A} \to \mathcal{B}$ 是 $X$ 上环层的一个同态。
则 de Rham 复形之间存在典范等同
$f^{-1}\Omega^\bullet_{\mathcal{B}/\mathcal{A}} =
\Omega^\bullet_{f^{-1}\mathcal{B}/f^{-1}\mathcal{A}}$
。
\end{lemma}

\begin{proof}
从略。提示：与引理 \ref{modules-lemma-pullback-differentials} 比较。
\end{proof}

\begin{lemma}
\label{modules-lemma-differentials-de-rham-complex-order-1}
设 $X$ 是一个拓扑空间，$\mathcal{A} \to \mathcal{B}$
是 $X$ 上环层的一个同态。微分
$\text{d} : \Omega^i_{\mathcal{B}/\mathcal{A}}  \to
\Omega^{i + 1}_{\mathcal{B}/\mathcal{A}}$
是一阶微分算子。
\end{lemma}

\begin{proof}
根据上面把 de Rham 复形构造为规则
$U \mapsto \Omega^\bullet_{\mathcal{B}(U)/\mathcal{A}(U)}$ 的层化，
这由《代数》引理
\ref{algebra-lemma-differentials-de-rham-complex-order-1} 推出。
\end{proof}

\noindent
设 $X$ 是一个拓扑空间，并设
$$
\xymatrix{
\mathcal{B} \ar[r] & \mathcal{B}' \\
\mathcal{A} \ar[r] \ar[u] & \mathcal{A}' \ar[u]
}
$$
是 $X$ 上环层的一个交换图。存在 de Rham 复形的自然同态
$$
\Omega^\bullet_{\mathcal{B}/\mathcal{A}} \longrightarrow
\Omega^\bullet_{\mathcal{B}'/\mathcal{A}'}
$$
具体而言，在 $0$ 次它是同态 $\mathcal{B} \to \mathcal{B}'$；
在 $1$ 次它是第 \ref{modules-section-differentials} 节中构造的同态
$\Omega_{\mathcal{B}/\mathcal{A}} \to \Omega_{\mathcal{B}'/\mathcal{A}'}$
；而对 $p \geq 2$，它是诱导同态
$\Omega^p_{\mathcal{B}/\mathcal{A}} =
\wedge^p_\mathcal{B}(\Omega_{\mathcal{B}/\mathcal{A}}) \to
\wedge^p_{\mathcal{B}'}(\Omega_{\mathcal{B}'/\mathcal{A}'}) =
\Omega^p_{\mathcal{B}'/\mathcal{A}'}$.
它与微分的相容性由公式 (\ref{modules-equation-rule}) 对微分的刻画得到。

\begin{definition}
\label{modules-definition-de-rham-complex-morphism-ringed-spaces}
设 $f : (X, \mathcal{O}_X) \to (Y, \mathcal{O}_Y)$ 是环空间的一个态射。
$f$ 的 de Rham 复形，或称 $X$ 在 $Y$ 上的 {\it de Rham 复形}，是复形
$$
\Omega^\bullet_{X/Y} = \Omega^\bullet_{\mathcal{O}_X/f^{-1}\mathcal{O}_Y}
$$
\end{definition}

\noindent
考虑环空间的交换图
$$
\xymatrix{
X' \ar[d]_{h'} \ar[r]_f & X \ar[d]^h \\
S' \ar[r]^g & S
}
$$
于是得到 de Rham 复形的典范同态
$$
\Omega^\bullet_{X/S} \to f_*\Omega^\bullet_{X'/S'}
$$
事实上，$X'$ 上环层的交换图
$$
\xymatrix{
f^{-1}\mathcal{O}_X \ar[r] & \mathcal{O}_{X'} \\
f^{-1}h^{-1}\mathcal{O}_S \ar[u] \ar[r] & (h')^{-1}\mathcal{O}_{S'} \ar[u]
}
$$
产生复形同态（见上文）
$$
f^{-1}\Omega^\bullet_{X/S} =
\Omega^\bullet_{f^{-1}\mathcal{O}_X/f^{-1}h^{-1}\mathcal{O}_S}
\longrightarrow
\Omega^\bullet_{\mathcal{O}_{X'}/(h')^{-1}\mathcal{O}_{S'}} =
\Omega^\bullet_{X'/S'}
$$
（第一个等式使用引理 \ref{modules-lemma-pullback-de-rham-complex}），
然后应用伴随性即可。

\begin{lemma}
\label{modules-lemma-differentials-relative-de-rham-complex-order-1}
设 $f : X \to Y$ 是环空间的一个态射。微分
$\text{d} : \Omega^i_{X/Y}  \to \Omega^{i + 1}_{X/Y}$
是 $X/Y$ 上的一阶微分算子。
\end{lemma}

\begin{proof}
这立即由引理 \ref{modules-lemma-differentials-de-rham-complex-order-1} 和定义得到。
\end{proof}





















\section{朴素余切复形}
\label{modules-section-netherlander}

\noindent
本节是《代数》第 \ref{algebra-section-netherlander} 节对环空间态射的类似版本。
我们建议读者先阅读该节。

\medskip\noindent
设 $X$ 是一个拓扑空间，$\mathcal{A} \to \mathcal{B}$ 是环层的一个同态。
在本节中，对 $X$ 上任意集合层 $\mathcal{E}$，用
$\mathcal{A}[\mathcal{E}]$ 表示预层
$U \mapsto \mathcal{A}(U)[\mathcal{E}(U)]$ 的层化。这里，
$\mathcal{A}(U)[\mathcal{E}(U)]$
表示 $\mathcal{A}(U)$ 上以 $\mathcal{E}(U)$ 的元素为变量的多项式代数。
对 $e \in \mathcal{E}(U)$，用
$[e] \in \mathcal{A}(U)[\mathcal{E}(U)]$ 表示与 $e$ 对应的变量。
存在 $\mathcal{A}$-代数的典范满射
\begin{equation}
\label{modules-equation-canonical-presentation}
\mathcal{A}[\mathcal{B}] \longrightarrow \mathcal{B},\quad [b] \longmapsto b
\end{equation}
把它的核记作 $\mathcal{I} \subset \mathcal{A}[\mathcal{B}]$。
容易看出，$\mathcal{I}$ 由局部截面
$[b][b'] - [bb']$ 和 $[a] - a$ 生成。根据引理
\ref{modules-lemma-differential-seq}，存在典范同态
\begin{equation}
\label{modules-equation-naive-cotangent-complex}
\mathcal{I}/\mathcal{I}^2
\longrightarrow
\Omega_{\mathcal{A}[\mathcal{B}]/\mathcal{A}}
\otimes_{\mathcal{A}[\mathcal{B}]} \mathcal{B}
\end{equation}
其余核典范同构于 $\Omega_{\mathcal{B}/\mathcal{A}}$。

\begin{definition}
\label{modules-definition-naive-cotangent-complex}
设 $X$ 是一个拓扑空间，$\mathcal{A} \to \mathcal{B}$ 是环层的一个同态。
{\it 朴素余切复形} $\NL_{\mathcal{B}/\mathcal{A}}$ 是链复形
(\ref{modules-equation-naive-cotangent-complex})：
$$
\NL_{\mathcal{B}/\mathcal{A}} =
\left(\mathcal{I}/\mathcal{I}^2
\longrightarrow
\Omega_{\mathcal{A}[\mathcal{B}]/\mathcal{A}}
\otimes_{\mathcal{A}[\mathcal{B}]} \mathcal{B}\right)
$$
其中 $\mathcal{I}/\mathcal{I}^2$ 位于 $-1$ 次，而
$\Omega_{\mathcal{A}[\mathcal{B}]/\mathcal{A}}
\otimes_{\mathcal{A}[\mathcal{B}]} \mathcal{B}$
位于 $0$ 次。
\end{definition}

\noindent
此构造具有类似于引理 \ref{modules-lemma-functoriality-differentials}
中对微分模层所讨论的函子性。具体而言，给定交换图
\begin{equation}
\label{modules-equation-commutative-square-sheaves}
\vcenter{
\xymatrix{
\mathcal{B} \ar[r] & \mathcal{B}' \\
\mathcal{A} \ar[u] \ar[r] & \mathcal{A}' \ar[u]
}
}
\end{equation}
作为 $X$ 上环层的图，存在典范的 $\mathcal{B}$-线性复形同态
$$
\NL_{\mathcal{B}/\mathcal{A}} \longrightarrow \NL_{\mathcal{B}'/\mathcal{A}'}
$$
事实上，交换图中的同态给出典范同态
$\mathcal{A}[\mathcal{B}] \to \mathcal{A}'[\mathcal{B}']$
，它把 $\mathcal{I}$ 映入
$\mathcal{I}' = \Ker(\mathcal{A}'[\mathcal{B}'] \to \mathcal{B}')$.
由此得到同态
$\mathcal{I}/\mathcal{I}^2 \to \mathcal{I}'/(\mathcal{I}')^2$
以及微分模层之间的同态，二者共同给出所需的朴素余切复形之间的同态。
该同态在如下意义下与复合相容：给定交换图
$$
\xymatrix{
\mathcal{B} \ar[r] &
\mathcal{B}' \ar[r] &
\mathcal{B}'' \\
\mathcal{A} \ar[u] \ar[r] &
\mathcal{A}' \ar[u] \ar[r] &
\mathcal{A}'' \ar[u]
}
$$
作为环层的图，则复合
$$
\NL_{\mathcal{B}/\mathcal{A}} \longrightarrow
\NL_{\mathcal{B}'/\mathcal{A}'} \longrightarrow
\NL_{\mathcal{B}''/\mathcal{A}''}
$$
就是外部矩形所对应的同态。

\medskip\noindent
可以选取 $\mathcal{B}$ 作为 $\mathcal{A}$ 上多项式代数之商的另一个表示，
而仍在 $D(\mathcal{B})$ 中得到同一个对象。为说明这一点，设
$\mathcal{E}$ 是 $X$ 上的一个集合层，
$\alpha : \mathcal{E} \to \mathcal{B}$ 是集合层的一个映射。
于是得到一个 $\mathcal{A}$-代数同态
$\mathcal{A}[\mathcal{E}] \to \mathcal{B}$。若此同态满射，也就是说，
若 $\alpha(\mathcal{E})$ 作为 $\mathcal{A}$-代数生成 $\mathcal{B}$，则令
$$
\NL(\alpha) = \left(
\mathcal{J}/\mathcal{J}^2
\longrightarrow
\Omega_{\mathcal{A}[\mathcal{E}]/\mathcal{A}}
\otimes_{\mathcal{A}[\mathcal{E}]} \mathcal{B}\right)
$$
其中 $\mathcal{J} \subset \mathcal{A}[\mathcal{E}]$ 是满射
$\mathcal{A}[\mathcal{E}] \to \mathcal{B}$ 的核。结论如下。

\begin{lemma}
\label{modules-lemma-NL-up-to-qis}
在上述情形中，$D(\mathcal{B})$ 中存在典范同构
$\NL(\alpha) = \NL_{\mathcal{B}/\mathcal{A}}$。
\end{lemma}

\begin{proof}
注意 $\NL_{\mathcal{B}/\mathcal{A}} = \NL(\text{id}_\mathcal{B})$。
因此，只需证明：给定如上的两个映射
$\alpha_i : \mathcal{E}_i \to \mathcal{B}$，在 $D(\mathcal{B})$ 中存在典范拟同构
$\NL(\alpha_1) = \NL(\alpha_2)$。为此，令
$\mathcal{E} = \mathcal{E}_1 \amalg \mathcal{E}_2$ 且
$\alpha = \alpha_1 \amalg \alpha_2 : \mathcal{E} \to \mathcal{B}$.
令
$\mathcal{J}_i = \Ker(\mathcal{A}[\mathcal{E}_i] \to \mathcal{B})$
以及
$\mathcal{J} = \Ker(\mathcal{A}[\mathcal{E}] \to \mathcal{B})$.
我们得到把 $\mathcal{J}_i$ 映入 $\mathcal{J}$ 的同态
$\mathcal{A}[\mathcal{E}_i] \to \mathcal{A}[\mathcal{E}]$。
因而得到典范复形同态
$$
\NL(\alpha_i) \longrightarrow \NL(\alpha)
$$
，只需证明这些同态是拟同构。为此，只需在茎上检验（引理
\ref{modules-lemma-abelian}）。若 $x \in X$，则 $\NL(\alpha)$ 的茎是
《代数》第 \ref{algebra-section-netherlander} 节中与映射
$\alpha_x : \mathcal{E}_x \to \mathcal{B}_x$ 所给表示
$\mathcal{A}_x[\mathcal{E}_x] \to \mathcal{B}_x$ 相伴的复形
$\NL(\alpha_x)$。（省略若干细节；用引理
\ref{modules-lemma-stalk-module-differentials} 验证形成微分与取茎相容。）
于是由《代数》引理 \ref{algebra-lemma-NL-homotopy} 得到结论。
\end{proof}

\begin{lemma}
\label{modules-lemma-pullback-NL}
设 $f : X \to Y$ 是拓扑空间的一个连续映射，
$\mathcal{A} \to \mathcal{B}$ 是 $Y$ 上环层的一个同态。则
$f^{-1}\NL_{\mathcal{B}/\mathcal{A}} =
\NL_{f^{-1}\mathcal{B}/f^{-1}\mathcal{A}}$。
\end{lemma}

\begin{proof}
从略。提示：使用引理 \ref{modules-lemma-pullback-differentials}。
\end{proof}

\begin{lemma}
\label{modules-lemma-stalk-NL}
设 $X$ 是一个拓扑空间，$\mathcal{A} \to \mathcal{B}$
是 $X$ 上环层的一个同态，并设 $x \in X$。则
$\NL_{\mathcal{B}/\mathcal{A}, x} =
\NL_{\mathcal{B}_x/\mathcal{A}_x}$。
\end{lemma}

\begin{proof}
这是引理 \ref{modules-lemma-pullback-NL} 对包含映射 $\{x\} \to X$ 的特殊情形。
\end{proof}

\begin{lemma}
\label{modules-lemma-exact-sequence-NL}
设 $X$ 是一个拓扑空间，并设
$\mathcal{A} \to \mathcal{B} \to \mathcal{C}$
是环层的同态。令 $C$ 为复形同态
$\NL_{\mathcal{C}/\mathcal{A}} \to \NL_{\mathcal{C}/\mathcal{B}}$
的锥（《导出范畴》，定义 \ref{derived-definition-cone}）。
存在 $\mathcal{C}$-模层复形的典范同态
$$
c :
\NL_{\mathcal{B}/\mathcal{A}} \otimes_\mathcal{B} \mathcal{C}
\longrightarrow
C[-1]
$$
，它给出上同调层的典范六项正合列
$$
\xymatrix{
H^0(\NL_{\mathcal{B}/\mathcal{A}} \otimes_\mathcal{B} \mathcal{C}) \ar[r] &
H^0(\NL_{\mathcal{C}/\mathcal{A}}) \ar[r] &
H^0(\NL_{\mathcal{C}/\mathcal{B}}) \ar[r] &
0 \\
H^{-1}(\NL_{\mathcal{B}/\mathcal{A}} \otimes_\mathcal{B} \mathcal{C}) \ar[r] &
H^{-1}(\NL_{\mathcal{C}/\mathcal{A}}) \ar[r] &
H^{-1}(\NL_{\mathcal{C}/\mathcal{B}}) \ar[llu]
}
$$
。
\end{lemma}

\begin{proof}
为给出同态 $c$，必须给出同态
$c_1 : \NL_{\mathcal{B}/\mathcal{A}} \otimes_\mathcal{B} \mathcal{C}
\to \NL_{\mathcal{C}/\mathcal{A}}$，以及复合
$$
\NL_{\mathcal{B}/\mathcal{A}} \otimes_\mathcal{B} \mathcal{C} \to
\NL_{\mathcal{C}/\mathcal{A}} \to \NL_{\mathcal{C}/\mathcal{B}}
$$
与零同态之间的一个显式同伦，见《导出范畴》引理
\ref{derived-lemma-map-from-cone}。对 $c_1$，对显然的图使用上述函子性。
对该同伦，使用同态
$$
\NL_{\mathcal{B}/\mathcal{A}}^0 \otimes_\mathcal{B} \mathcal{C}
\longrightarrow
\NL_{\mathcal{C}/\mathcal{B}}^{-1},\quad
\text{d}[b] \otimes 1 \longmapsto [\varphi(b)] - b[1]
$$
其中 $\varphi : \mathcal{B} \to \mathcal{C}$ 是给定同态。请与《代数》评注
\ref{algebra-remark-composition-homotopy-equivalent-to-zero} 比较。
为得到关于上同调层的结论，只需证明 $H^0(c)$ 是同构且
$H^{-1}(c)$ 是满射。为此可以考察茎，见引理
\ref{modules-lemma-stalk-NL}；然后使用交换代数中的相应结论，见《代数》引理
\ref{algebra-lemma-exact-sequence-NL}。省略若干细节。
\end{proof}

\noindent
环空间态射的余切复形用上面定义的余切复形来定义。

\begin{definition}
\label{modules-definition-cotangent-complex-morphism-ringed-topoi}
环空间态射 $f : (X, \mathcal{O}_X) \to (Y, \mathcal{O}_Y)$ 的
{\it 朴素余切复形} $\NL_f = \NL_{X/Y}$ 是
$\NL_{\mathcal{O}_X/f^{-1}\mathcal{O}_Y}$。
\end{definition}

\noindent
给定环空间的交换图
$$
\xymatrix{
X' \ar[r]_g \ar[d]_{f'} & X \ar[d]^f \\
Y' \ar[r]^h & Y
}
$$
，存在典范同态 $c : g^*\NL_{X/Y} \to \NL_{X'/Y'}$。具体而言，它是同态
$$
g^*\NL_{X/Y} =
\mathcal{O}_{X'} \otimes_{g^{-1}\mathcal{O}_X}
\NL_{g^{-1}\mathcal{O}_X/g^{-1}f^{-1}\mathcal{O}_Y}
\longrightarrow
\NL_{\mathcal{O}_{X'}/(f')^{-1}\mathcal{O}_{Y'}} = \NL_{X'/Y'}
$$
，其中箭头来自环层的交换图
$$
\xymatrix{
g^{-1}\mathcal{O}_X \ar[r]_{g^\sharp} &
\mathcal{O}_{X'} \\
g^{-1}f^{-1}\mathcal{O}_Y
\ar[r]^{g^{-1}h^\sharp}
\ar[u]^{g^{-1}f^\sharp} &
(f')^{-1}\mathcal{O}_{Y'} \ar[u]_{(f')^\sharp}
}
$$
，如上面的 (\ref{modules-equation-commutative-square-sheaves})。再给定一个这样的图
$$
\xymatrix{
X'' \ar[r]_{g'} \ar[d] & X' \ar[d] \\
Y'' \ar[r] & Y'
}
$$
，则 $(g')^*c$ 与同态
$c' : (g')^*\NL_{X'/Y'} \to \NL_{X''/Y''}$
的复合是同态 $(g \circ g')^*\NL_{X''/Y''} \to \NL_{X/Y}$。

\begin{lemma}
\label{modules-lemma-exact-sequence-NL-ringed-topoi}
设 $f : X \to Y$ 和 $g : Y \to Z$ 是环空间的态射。
令 $C$ 为 $\mathcal{O}_X$-模层复形同态
$\NL_{X/Z} \to \NL_{X/Y}$ 的锥。存在典范同态
$$
f^*\NL_{Y/Z} \to C[-1]
$$
，它给出上同调层的典范六项正合列
$$
\xymatrix{
H^0(f^*\NL_{Y/Z}) \ar[r] &
H^0(\NL_{X/Z}) \ar[r] &
H^0(\NL_{X/Y}) \ar[r] &
0 \\
H^{-1}(f^*\NL_{Y/Z}) \ar[r] &
H^{-1}(\NL_{X/Z}) \ar[r] &
H^{-1}(\NL_{X/Y}) \ar[llu]
}
$$
。
\end{lemma}

\begin{proof}
考虑环层的同态
$$
(g \circ f)^{-1}\mathcal{O}_Z \to f^{-1}\mathcal{O}_Y \to \mathcal{O}_X
$$
并应用引理 \ref{modules-lemma-exact-sequence-NL}。
\end{proof}
% Generated mechanically from the frozen sites-modules.tex target.
% Do not edit this generated file; edit the bound locale source instead.

% OK, start here.
%

\setcounter{chapter}{17}
\chapter{站点上的模}



\phantomsection
\label{sites-modules-section-phantom}


\section{引言}
\label{sites-modules-section-introduction}

\noindent
本章阐述环拓扑斯或环站点上的模层之基本概念。我们先给出阿贝尔层的若干
基本事实，随后引入环站点与环拓扑斯。我们讨论 $\mathcal{O}$-模（预）层的
一些最基本概念；这些内容类似于空间上的层一章中关于 $\mathcal{O}$-模
（预）层的材料。完成这些准备后，我们重述《模层》中相当一部分讨论
（见《模层》第 \ref{modules-section-introduction} 节）。基本参考文献为
\cite{FAC}、\cite{EGA} 和 \cite{SGA4}。






\section{阿贝尔预层}
\label{sites-modules-section-abelian-pre-sheaves}

\noindent
设 $\mathcal{C}$ 为范畴。阿贝尔预层已在《站点》第
\ref{sites-section-presheaves}、\ref{sites-section-sheaves} 节中引入，
并在《站点》第 \ref{sites-section-sheaves-algebraic-structures} 节中作了
进一步讨论。我们沿用后一处的约定：将阿贝尔预层看作集合预层，其每个截面
集合上都带有与限制映射相容的加法法则。回顾 $\mathcal{C}$ 上阿贝尔预层的
范畴记为 $\textit{PAb}(\mathcal{C})$。

\medskip\noindent
依《同调》定义 \ref{homology-definition-abelian-category}，范畴
$\textit{PAb}(\mathcal{C})$ 是阿贝尔范畴。给定预层映射
$\varphi : \mathcal{G}_1 \to \mathcal{G}_2$，$\varphi$ 的核是阿贝尔预层
$U \mapsto \Ker(\mathcal{G}_1(U) \to \mathcal{G}_2(U))$，
而 $\varphi$ 的余核是预层
$U \mapsto \Coker(\mathcal{G}_1(U) \to \mathcal{G}_2(U))$。
由于阿贝尔群范畴是阿贝尔范畴，并且这一点对每个 $U$ 都成立，故
$\Coim = \Im$。阿贝尔预层列
$$
\mathcal{G}_1 \longrightarrow
\mathcal{G}_2 \longrightarrow
\mathcal{G}_3
$$
正合，当且仅当对所有 $U \in \Ob(\mathcal{C})$，阿贝尔群列
$\mathcal{G}_1(U) \to \mathcal{G}_2(U) \to \mathcal{G}_3(U)$
正合。验证留给读者。

\begin{lemma}
\label{sites-modules-lemma-limits-colimits-abelian-presheaves}
设 $\mathcal{C}$ 为范畴。
\begin{enumerate}
\item $\textit{PAb}(\mathcal{C})$ 中一切极限与余极限均存在。
\item 一切极限与余极限均与在 $\mathcal{C}$ 的对象上取截面可交换。
\end{enumerate}
\end{lemma}

\begin{proof}
设 $\mathcal{I} \to \textit{PAb}(\mathcal{C})$，$i \mapsto \mathcal{F}_i$
为图。我们可直接按下列法则定义阿贝尔预层 $L$ 与 $C$：
$$
L : U \longmapsto \lim_i \mathcal{F}_i(U)
$$
以及
$$
C : U \longmapsto \colim_i \mathcal{F}_i(U).
$$
在每个 $U$ 上的截面群中使用相应映射，显然得到阿贝尔预层映射
$L \to \mathcal{F}_i$ 与 $\mathcal{F}_i \to C$。容易验证，配备这些映射的
$L$ 与 $C$ 分别是该图在 $\textit{PAb}(\mathcal{C})$ 中的极限与余极限。
这证明了 (1) 与 (2)。细节从略。
\end{proof}


\section{阿贝尔层}
\label{sites-modules-section-abelian-sheaves}

\noindent
设 $\mathcal{C}$ 为站点。$\mathcal{C}$ 上阿贝尔层的范畴记为
$\textit{Ab}(\mathcal{C})$。它是 $\textit{PAb}(\mathcal{C})$ 的全子范畴，
由底层集合预层为层的那些阿贝尔预层组成。《站点》第
\ref{sites-section-sheaves-algebraic-structures} 节的性质
($\alpha$)--($\zeta$) 成立；见《站点》命题
\ref{sites-proposition-functoriality-algebraic-structures-topoi}。
特别地，包含函子
$\textit{Ab}(\mathcal{C}) \to \textit{PAb}(\mathcal{C})$
具有左伴随，即层化函子 $\mathcal{G} \mapsto \mathcal{G}^\#$。

\medskip\noindent
我们建议读者在草稿纸上自行证明下述引理，而不是阅读其证明。

\begin{lemma}
\label{sites-modules-lemma-abelian-abelian}
设 $\mathcal{C}$ 为站点，$\varphi : \mathcal{F} \to \mathcal{G}$ 为
$\mathcal{C}$ 上阿贝尔层的态射。
\begin{enumerate}
\item 范畴 $\textit{Ab}(\mathcal{C})$ 是阿贝尔范畴。
\item $\varphi$ 的核 $\Ker(\varphi)$ 与将 $\varphi$ 视为预层态射时的核相同。
\item 态射 $\varphi$ 是单射（《同调》定义
\ref{homology-definition-injective-surjective}），当且仅当它作为预层映射是单射
（《站点》定义 \ref{sites-definition-presheaves-injective-surjective}），又当且仅当
它作为层映射是单射（《站点》定义
\ref{sites-definition-sheaves-injective-surjective}）。
\item $\varphi$ 的余核 $\Coker(\varphi)$ 是将 $\varphi$ 视为预层态射时所得
余核的层化。
\item 态射 $\varphi$ 是满射（《同调》定义
\ref{homology-definition-injective-surjective}），当且仅当它作为层映射是满射
（《站点》定义 \ref{sites-definition-sheaves-injective-surjective}）。
\item 阿贝尔层复形
$$
\mathcal{F} \to \mathcal{G} \to \mathcal{H}
$$
在 $\mathcal{G}$ 处正合，当且仅当对每个 $U \in \Ob(\mathcal{C})$ 及每个
在 $\mathcal{H}(U)$ 中映为零的 $s \in \mathcal{G}(U)$，都存在
$\mathcal{C}$ 中的覆盖 $\{U_i \to U\}_{i \in I}$，使每个 $s|_{U_i}$
都属于 $\mathcal{F}(U_i) \to \mathcal{G}(U_i)$ 的像。
\end{enumerate}
\end{lemma}

\begin{proof}
我们断言，《同调》引理 \ref{homology-lemma-adjoint-get-abelian} 可用于范畴
$\mathcal{A} = \textit{Ab}(\mathcal{C})$、
$\mathcal{B} = \textit{PAb}(\mathcal{C})$ 以及函子
$a : \mathcal{A} \to \mathcal{B}$（包含）与
$b : \mathcal{B} \to \mathcal{A}$（层化）。下面核验《同调》引理
\ref{homology-lemma-adjoint-get-abelian} 的假设。
假设 (1) 要求 $\mathcal{A}$、$\mathcal{B}$ 为加性范畴，$a$、$b$ 为加性
函子，且 $a$ 是 $b$ 的右伴随。前两点显然，而伴随性就是《站点》第
\ref{sites-section-sheaves-algebraic-structures} 节的 ($\epsilon$)。假设 (2)
要求 $\textit{PAb}(\mathcal{C})$ 为阿贝尔范畴——这已在第
\ref{sites-modules-section-abelian-pre-sheaves} 节中看到——并要求层化为左正合函子；后者
即《站点》第 \ref{sites-section-sheaves-algebraic-structures} 节的 ($\zeta$)。
最后一个假设是 $ba \cong \text{id}_\mathcal{A}$，即《站点》第
\ref{sites-section-sheaves-algebraic-structures} 节的 ($\delta$)。
故《同调》引理 \ref{homology-lemma-adjoint-get-abelian} 适用，从而
$\textit{Ab}(\mathcal{C})$ 是阿贝尔范畴。

\medskip\noindent
《同调》引理 \ref{homology-lemma-adjoint-get-abelian} 的证明表明，
$\Ker(\varphi)$ 与 $\Coker(\varphi)$ 分别等于将 $\varphi$ 视为阿贝尔预层
态射时所得核与余核的层化。这证明了 (4)。由于核是一个等化子（因而是极限），
且层化与有限极限可交换，遂得 (2)。

\medskip\noindent
断言 (2) 蕴含 (3)。由我们对层化的描述，断言 (4) 蕴含 (5)。断言 (6) 中
对正合性的刻画来自 (2)、(5)，以及该列正合当且仅当
$\Im(\mathcal{F} \to \mathcal{G}) =
\Ker(\mathcal{G} \to \mathcal{H})$.
\end{proof}

\noindent
该引理第 (6) 项还可表述为：阿贝尔层列
$$
\mathcal{F}_1 \longrightarrow
\mathcal{F}_2 \longrightarrow
\mathcal{F}_3
$$
正合，当且仅当预层
$U \mapsto \Im(\mathcal{F}_1(U) \to \mathcal{F}_2(U))$
的层化等于 $\mathcal{F}_2 \to \mathcal{F}_3$ 的核。

\begin{lemma}
\label{sites-modules-lemma-limits-colimits-abelian-sheaves}
设 $\mathcal{C}$ 为站点。
\begin{enumerate}
\item $\textit{Ab}(\mathcal{C})$ 中一切极限与余极限均存在。
\item 极限与 $\mathcal{C}$ 上阿贝尔预层的相应极限相同（即与在
$\mathcal{C}$ 的对象上取截面可交换）。
\item 有限直和与 $\mathcal{C}$ 上阿贝尔预层范畴中的相应有限直和相同。
\item 余极限是阿贝尔预层范畴中相应余极限的层化。
\item 滤过余极限正合。
\end{enumerate}
\end{lemma}

\begin{proof}
由引理 \ref{sites-modules-lemma-limits-colimits-abelian-presheaves}，阿贝尔预层的极限与
余极限存在，并可描述为在对象上的截面层面取极限与余极限。

\medskip\noindent
设 $\mathcal{I} \to \textit{Ab}(\mathcal{C})$，$i \mapsto \mathcal{F}_i$
为图，并令 $\lim_i \mathcal{F}_i$ 为该图在阿贝尔预层范畴中的极限。由
《站点》引理 \ref{sites-lemma-limit-sheaf}，它是阿贝尔层。此时容易看出，
$\lim_i \mathcal{F}_i$ 也是该图在 $\textit{Ab}(\mathcal{C})$ 中的极限。
这证明了极限存在，且 (2) 成立。

\medskip\noindent
由《范畴论》引理 \ref{categories-lemma-adjoint-exact}，并利用层化是包含函子的
左伴随，可知 $\colim_i \mathcal{F}$ 存在，且它是
$\textit{PAb}(\mathcal{C})$ 中余极限的层化。这证明了余极限存在，且 (4) 成立。

\medskip\noindent
在任何阿贝尔范畴中，有限直和与有限直积相同。因此 (3) 来自 (2)。

\medskip\noindent
证明 (5)。该断言意指：给定由有向集 $I$ 参数化的阿贝尔层正合列系统
$0 \to \mathcal{F}_i \to \mathcal{G}_i \to \mathcal{H}_i \to 0$，列
$0 \to \colim \mathcal{F}_i \to \colim \mathcal{G}_i \to
\colim \mathcal{H}_i \to 0$ 也正合。利用《同调》引理
\ref{homology-lemma-check-exactness} 与余极限定义作形式论证，可知列
$\colim \mathcal{F}_i \to \colim \mathcal{G}_i \to \colim \mathcal{H}_i \to 0$
正合。注意，$\colim \mathcal{F}_i \to \colim \mathcal{G}_i$ 是预层余极限
映射的层化；由于每个 $\mathcal{F}_i \to \mathcal{G}_i$ 都是单射，后一映射
也是单射。层化正合，故得结论。
\end{proof}







\section{自由阿贝尔预层}
\label{sites-modules-section-free-abelian-presheaf}

\noindent
为下述定义准备符号，我们约定将集合 $S$ 上的自由阿贝尔群记为\footnote{在其他
章节中，符号 $\mathbf{Z}[S]$ 有时表示以 $S$ 为变量集的 $\mathbf{Z}$ 上多项式环。}
$\mathbf{Z}[S] = \bigoplus_{s \in S} \mathbf{Z}$。其刻画性质为
$$
\Mor_{\textit{Ab}}(\mathbf{Z}[S], A)
=
\Mor_{\textit{Sets}}(S, A)
$$
换言之，构造 $S \mapsto \mathbf{Z}[S]$ 是遗忘函子
$\textit{Ab} \to \textit{Sets}$ 的左伴随。

\begin{definition}
\label{sites-modules-definition-free-abelian-presheaf-on}
设 $\mathcal{C}$ 为范畴，$\mathcal{G}$ 为集合预层。$\mathcal{G}$ 上的
{\it 自由阿贝尔预层} $\mathbf{Z}_\mathcal{G}$ 是由下列法则定义的阿贝尔预层：
$$
U \longmapsto \mathbf{Z}[\mathcal{G}(U)].
$$
在 $\mathcal{G} = h_X$ 是与 $\mathcal{C}$ 的对象 $X$ 相伴的可表预层这一
特殊情形，我们采用符号 $\mathbf{Z}_X = \mathbf{Z}_{h_X}$。换言之，
$$
\mathbf{Z}_X(U) = \mathbf{Z}[\Mor_\mathcal{C}(U, X)].
$$
\end{definition}

\noindent
此构造关于预层 $\mathcal{G}$ 显然具有函子性。事实上，它是遗忘函子
$\textit{PAb}(\mathcal{C}) \to \textit{PSh}(\mathcal{C})$ 的左伴随。
准确陈述如下。

\begin{lemma}
\label{sites-modules-lemma-obvious-adjointness}
设 $\mathcal{C}$ 为范畴，$\mathcal{G}$、$\mathcal{F}$ 为集合预层，
$\mathcal{A}$ 为阿贝尔预层，$U$ 为 $\mathcal{C}$ 的对象。则
\begin{align*}
\Mor_{\textit{PSh}(\mathcal{C})}(h_U, \mathcal{F})
& =
\mathcal{F}(U), \\
\Mor_{\textit{PAb}(\mathcal{C})}(\mathbf{Z}_\mathcal{G}, \mathcal{A})
& =
\Mor_{\textit{PSh}(\mathcal{C})}(\mathcal{G}, \mathcal{A}), \\
\Mor_{\textit{PAb}(\mathcal{C})}(\mathbf{Z}_U, \mathcal{A})
& =
\mathcal{A}(U).
\end{align*}
所有这些等式均具有函子性。
\end{lemma}

\begin{proof}
从略。
\end{proof}

\begin{lemma}
\label{sites-modules-lemma-coproduct-sum-free-abelian-presheaf}
设 $\mathcal{C}$ 为范畴，$I$ 为集合。对每个 $i \in I$，设
$\mathcal{G}_i$ 为集合预层。则
$$
\mathbf{Z}_{\coprod_i \mathcal{G}_i}
=
\bigoplus\nolimits_{i \in I} \mathbf{Z}_{\mathcal{G}_i}
$$
在 $\textit{PAb}(\mathcal{C})$ 中成立。
\end{lemma}

\begin{proof}
从略。
\end{proof}



\section{自由阿贝尔层}
\label{sites-modules-section-free-abelian-sheaf}

\noindent
下面给出集合层上的自由阿贝尔层之概念。

\begin{definition}
\label{sites-modules-definition-free-abelian-sheaf-on}
设 $\mathcal{C}$ 为站点，$\mathcal{G}$ 为集合预层。$\mathcal{G}$ 上的
{\it 自由阿贝尔层} $\mathbf{Z}_\mathcal{G}^\#$ 是阿贝尔层
$\mathbf{Z}_\mathcal{G}^\#$，即 $\mathcal{G}$ 上自由阿贝尔预层的层化。
在 $\mathcal{G} = h_X$ 是与 $\mathcal{C}$ 的对象 $X$ 相伴的可表预层这一
特殊情形，我们采用符号 $\mathbf{Z}_X^\#$。
\end{definition}

\noindent
此构造关于预层 $\mathcal{G}$ 显然具有函子性。事实上，它给出遗忘函子
$\textit{Ab}(\mathcal{C}) \to \Sh(\mathcal{C})$ 的左伴随。准确陈述如下。

\begin{lemma}
\label{sites-modules-lemma-obvious-adjointness-sheaves}
设 $\mathcal{C}$ 为站点，$\mathcal{G}$、$\mathcal{F}$ 为集合层，
$\mathcal{A}$ 为阿贝尔层，$U$ 为 $\mathcal{C}$ 的对象。则
\begin{align*}
\Mor_{\Sh(\mathcal{C})}(h_U^\#, \mathcal{F})
& =
\mathcal{F}(U), \\
\Mor_{\textit{Ab}(\mathcal{C})}(\mathbf{Z}_\mathcal{G}^\#,
\mathcal{A})
& =
\Mor_{\Sh(\mathcal{C})}(\mathcal{G}, \mathcal{A}), \\
\Mor_{\textit{Ab}(\mathcal{C})}(\mathbf{Z}_U^\#, \mathcal{A})
& =
\mathcal{A}(U).
\end{align*}
所有这些等式均具有函子性。
\end{lemma}

\begin{proof}
从略。
\end{proof}

\begin{lemma}
\label{sites-modules-lemma-may-sheafify-before-abelianize}
设 $\mathcal{C}$ 为站点，$\mathcal{G}$ 为集合预层。则
$\mathbf{Z}_\mathcal{G}^\# = (\mathbf{Z}_{\mathcal{G}^\#})^\#$。
\end{lemma}

\begin{proof}
从略。
\end{proof}








\section{环站点}
\label{sites-modules-section-ringed-sites}

\noindent
本章主要研究环站点上的模层，因此需要定义这一概念。

\begin{definition}
\label{sites-modules-definition-ringed-site}
环站点。
\begin{enumerate}
\item {\it 环站点}是偶 $(\mathcal{C}, \mathcal{O})$，其中
$\mathcal{C}$ 为站点，$\mathcal{O}$ 为 $\mathcal{C}$ 上的环层。层
$\mathcal{O}$ 称为该环站点的{\it 结构层}。
\item 设 $(\mathcal{C}, \mathcal{O})$、$(\mathcal{C}', \mathcal{O}')$
为环站点。一个{\it 环站点态射}
$$
(f, f^\sharp) :
(\mathcal{C}, \mathcal{O})
\longrightarrow
(\mathcal{C}', \mathcal{O}')
$$
由站点态射 $f : \mathcal{C} \to \mathcal{C}'$（见《站点》定义
\ref{sites-definition-morphism-sites}）以及环层映射
$f^\sharp : f^{-1}\mathcal{O}' \to \mathcal{O}$ 给出；由伴随性，后者等价于
环层映射 $f^\sharp : \mathcal{O}' \to f_*\mathcal{O}$。
\item 设
$(f, f^\sharp) :
(\mathcal{C}_1, \mathcal{O}_1) \to (\mathcal{C}_2, \mathcal{O}_2)$ 以及
$(g, g^\sharp) :
(\mathcal{C}_2, \mathcal{O}_2) \to (\mathcal{C}_3, \mathcal{O}_3)$
为环站点态射。按下列法则定义{\it 环站点态射的复合}：
$$
(g, g^\sharp) \circ (f, f^\sharp) = (g \circ f, f^\sharp \circ g^\sharp).
$$
这里使用《站点》定义 \ref{sites-definition-composition-morphisms-sites}
中定义的站点态射之复合，而 $f^\sharp \circ g^\sharp$ 表示环层态射
$$
\mathcal{O}_3 \xrightarrow{g^\sharp} g_*\mathcal{O}_2
\xrightarrow{g_*f^\sharp} g_*f_*\mathcal{O}_1 = (g \circ f)_*\mathcal{O}_1
$$
\end{enumerate}
\end{definition}






\section{环拓扑斯}
\label{sites-modules-section-ringed-topoi}

\noindent
环拓扑斯本身就是环站点；区别只在于，环拓扑斯态射的概念不同于环站点态射。

\begin{definition}
\label{sites-modules-definition-ringed-topos}
环拓扑斯。
\begin{enumerate}
\item {\it 环拓扑斯}是偶
$(\Sh(\mathcal{C}), \mathcal{O})$
，其中 $\mathcal{C}$ 为站点，$\mathcal{O}$ 为 $\mathcal{C}$ 上的环层。
层 $\mathcal{O}$ 称为该环拓扑斯的{\it 结构层}。
\item 设 $(\Sh(\mathcal{C}), \mathcal{O})$、
$(\Sh(\mathcal{C}'), \mathcal{O}')$ 为环拓扑斯。一个{\it 环拓扑斯态射}
$$
(f, f^\sharp) :
(\Sh(\mathcal{C}), \mathcal{O})
\longrightarrow
(\Sh(\mathcal{C}'), \mathcal{O}')
$$
由拓扑斯态射 $f : \Sh(\mathcal{C}) \to \Sh(\mathcal{C}')$（见《站点》定义
\ref{sites-definition-topos}）以及环层映射
$f^\sharp : f^{-1}\mathcal{O}' \to \mathcal{O}$ 给出；由伴随性，后者等价于
环层映射 $f^\sharp : \mathcal{O}' \to f_*\mathcal{O}$。
\item 设
$(f, f^\sharp) :
(\Sh(\mathcal{C}_1), \mathcal{O}_1)
\to (\Sh(\mathcal{C}_2), \mathcal{O}_2)$ 以及
$(g, g^\sharp) :
(\Sh(\mathcal{C}_2), \mathcal{O}_2) \to
(\Sh(\mathcal{C}_3), \mathcal{O}_3)$
为环拓扑斯态射。按下列法则定义{\it 环拓扑斯态射的复合}：
$$
(g, g^\sharp) \circ (f, f^\sharp) = (g \circ f, f^\sharp \circ g^\sharp).
$$
这里使用《站点》定义 \ref{sites-definition-topos} 中定义的拓扑斯态射之复合，
而 $f^\sharp \circ g^\sharp$ 表示环层态射
$$
\mathcal{O}_3 \xrightarrow{g^\sharp} g_*\mathcal{O}_2
\xrightarrow{g_*f^\sharp} g_*f_*\mathcal{O}_1 = (g \circ f)_*\mathcal{O}_1
$$
\end{enumerate}
\end{definition}

\noindent
每个环拓扑斯态射都是一个环拓扑斯等价与一个由环站点态射诱导的环拓扑斯态射
之复合。准确陈述如下。

\begin{lemma}
\label{sites-modules-lemma-morphism-ringed-topoi-comes-from-morphism-ringed-sites}
设 $(f, f^\sharp) :
(\Sh(\mathcal{C}), \mathcal{O}_\mathcal{C})
\to (\Sh(\mathcal{D}), \mathcal{O}_\mathcal{D})$
为环拓扑斯态射。存在分解
$$
\xymatrix{
(\Sh(\mathcal{C}), \mathcal{O}_\mathcal{C})
\ar[rr]_{(f, f^\sharp)}
\ar[d]_{(g, g^\sharp)}
& &
(\Sh(\mathcal{D}), \mathcal{O}_\mathcal{D}) \ar[d]^{(e, e^\sharp)}
\\
(\Sh(\mathcal{C}'), \mathcal{O}_{\mathcal{C}'})
\ar[rr]^{(h, h^\sharp)} & &
(\Sh(\mathcal{D}'), \mathcal{O}_{\mathcal{D}'})
}
$$
其中
\begin{enumerate}
\item $g : \Sh(\mathcal{C}) \to \Sh(\mathcal{C}')$
是由特殊余连续函子 $\mathcal{C} \to \mathcal{C}'$ 诱导的拓扑斯等价
（见《站点》定义 \ref{sites-definition-special-cocontinuous-functor}），
\item $e : \Sh(\mathcal{D}) \to \Sh(\mathcal{D}')$
是由特殊余连续函子 $\mathcal{D} \to \mathcal{D}'$ 诱导的拓扑斯等价
（见《站点》定义 \ref{sites-definition-special-cocontinuous-functor}），
\item $\mathcal{O}_{\mathcal{C}'} = g_*\mathcal{O}_\mathcal{C}$
且 $g^\sharp$ 为显然映射，
\item $\mathcal{O}_{\mathcal{D}'} = e_*\mathcal{O}_\mathcal{D}$
且 $e^\sharp$ 为显然映射，
\item 站点 $\mathcal{C}'$ 与 $\mathcal{D}'$ 具有终对象与纤维积（即一切有限极限），
\item $h$ 是由与一切有限极限可交换的连续函子
$u : \mathcal{D}' \to \mathcal{C}'$ 诱导的站点态射（即它满足《站点》命题
\ref{sites-proposition-get-morphism} 的假设），以及
\item 给定 $\mathcal{C}$ 上任意一族层 $\mathcal{F}_i$（相应地，
$\mathcal{D}$ 上任意一族层 $\mathcal{G}_j$），可假设它们分别是
$\mathcal{C}'$（相应地，$\mathcal{D}'$）上的可表层。
\end{enumerate}
此外，若 $(f, f^\sharp)$ 是环拓扑斯等价，则可选择该图使得
$\mathcal{C}' = \mathcal{D}'$,
$\mathcal{O}_{\mathcal{C}'} = \mathcal{O}_{\mathcal{D}'}$
且 $(h, h^\sharp)$ 为恒等态射。
\end{lemma}

\begin{proof}
这由《站点》引理 \ref{sites-lemma-morphism-topoi-comes-from-morphism-sites}
以及《站点》注
\ref{sites-remark-morphism-topoi-comes-from-morphism-sites}、
\ref{sites-remark-equivalence-topoi-comes-from-morphism-sites} 得出；只需同时
携带相应环层。部分细节从略。
\end{proof}







\section{环拓扑斯的 2-态射}
\label{sites-modules-section-2-category}

\noindent
本节简要讨论环拓扑斯的 $2$-态射这一概念。

\begin{definition}
\label{sites-modules-definition-2-morphism-ringed-topoi}
设
$f, g :
(\Sh(\mathcal{C}), \mathcal{O}_\mathcal{C})
\to
(\Sh(\mathcal{D}), \mathcal{O}_\mathcal{D})$
为两个环拓扑斯态射。一个{\it 从 $f$ 到 $g$ 的 2-态射}由函子变换
$t : f_* \to g_*$ 给出，使得
$$
\xymatrix{
& \mathcal{O}_\mathcal{D}
\ar[ld]_{f^\sharp}
\ar[rd]^{g^\sharp} \\
f_*\mathcal{O}_\mathcal{C} \ar[rr]^t & &
g_*\mathcal{O}_\mathcal{C}
}
$$
可交换。
\end{definition}

\noindent
在图示中，我们有时将 $t$ 表示如下：
$$
\xymatrix{
(\Sh(\mathcal{C}), \mathcal{O}_\mathcal{C})
\rrtwocell^f_g{t}
&
&
(\Sh(\mathcal{D}), \mathcal{O}_\mathcal{D})
}
$$
如《站点》第 \ref{sites-section-2-category} 节所述，给出 2-态射
$t : f_* \to g_*$ 等价于给出 $t : g^{-1} \to f^{-1}$（通常仍以同一符号
表示），使得图
$$
\xymatrix{
f^{-1}\mathcal{O}_\mathcal{D}
\ar[rd]_{f^\sharp}  & &
g^{-1}\mathcal{O}_\mathcal{D} \ar[ll]^t \ar[ld]^{g^\sharp} \\
& \mathcal{O}_\mathcal{C}
}
$$
可交换。同样依《站点》第 \ref{sites-section-2-category} 节，按该处所述定义
水平复合与垂直复合后，严格 2-范畴的公理成立。










\section{模预层}
\label{sites-modules-section-presheaves-modules}

\noindent
设 $\mathcal{C}$ 为范畴，$\mathcal{O}$ 为 $\mathcal{C}$ 上的环预层。我们尚未
定义 $\mathcal{O}$-模预层，现立即给出定义。

\begin{definition}
\label{sites-modules-definition-presheaf-modules}
设 $\mathcal{C}$ 为范畴，$\mathcal{O}$ 为 $\mathcal{C}$ 上的环预层。
\begin{enumerate}
\item 一个{\it $\mathcal{O}$-模预层}由阿贝尔预层 $\mathcal{F}$ 以及集合
预层映射
$$
\mathcal{O} \times \mathcal{F} \longrightarrow \mathcal{F}
$$
给出，并要求对 $\mathcal{C}$ 的每个对象 $U$，映射
$\mathcal{O}(U) \times \mathcal{F}(U) \to \mathcal{F}(U)$
在阿贝尔群 $\mathcal{F}(U)$ 上定义 $\mathcal{O}(U)$-模结构。
\item 一个{\it $\mathcal{O}$-模预层态射
$\varphi : \mathcal{F} \to \mathcal{G}$}是阿贝尔预层态射
$\varphi : \mathcal{F} \to \mathcal{G}$，使图
$$
\xymatrix{
\mathcal{O} \times \mathcal{F} \ar[r] \ar[d]_{\text{id} \times \varphi} &
\mathcal{F} \ar[d]^{\varphi} \\
\mathcal{O} \times \mathcal{G} \ar[r] &
\mathcal{G}
}
$$
可交换。
\item 上述 $\mathcal{O}$-模态射组成的集合记为
$\Hom_\mathcal{O}(\mathcal{F}, \mathcal{G})$。
\item $\mathcal{O}$-模预层的范畴记为 $\textit{PMod}(\mathcal{O})$。
\end{enumerate}
\end{definition}

\noindent
设 $\mathcal{O}_1 \to \mathcal{O}_2$ 是范畴 $\mathcal{C}$ 上环预层的态射。
若 $\mathcal{F}$ 是 $\mathcal{O}_2$-模预层，则借助复合
$$
\mathcal{O}_1 \times \mathcal{F}
\to
\mathcal{O}_2 \times \mathcal{F}
\to
\mathcal{F}.
$$
可将 $\mathcal{F}$ 视为 $\mathcal{O}_1$-模预层。为表明标量限制，我们有时将
它记为 $\mathcal{F}_{\mathcal{O}_1}$，并称之为 $\mathcal{F}$ 的
{\it 标量限制}。由此得到标量限制函子
$$
\textit{PMod}(\mathcal{O}_2)
\longrightarrow
\textit{PMod}(\mathcal{O}_1)
$$

\medskip\noindent
另一方面，给定 $\mathcal{O}_1$-模预层 $\mathcal{G}$，可按下列法则构造
$\mathcal{O}_2$-模预层
$\mathcal{O}_2 \otimes_{p, \mathcal{O}_1} \mathcal{G}$：
$$
U \longmapsto
\left(\mathcal{O}_2 \otimes_{p, \mathcal{O}_1} \mathcal{G}\right)(U)
=
\mathcal{O}_2(U) \otimes_{\mathcal{O}_1(U)} \mathcal{G}(U)
$$
其中 $U \in \Ob(\mathcal{C})$，限制映射取显然的映射。下标 $p$ 表示
“预层”，而非“点”。该预层称为张量积预层。由此得到{\it 换环}函子
$$
\textit{PMod}(\mathcal{O}_1)
\longrightarrow
\textit{PMod}(\mathcal{O}_2)
$$

\begin{lemma}
\label{sites-modules-lemma-adjointness-tensor-restrict-presheaves}
沿用上述 $\mathcal{C}$、$\mathcal{O}_1 \to \mathcal{O}_2$、$\mathcal{F}$ 与
$\mathcal{G}$，存在典范双射
$$
\Hom_{\mathcal{O}_1}(\mathcal{G}, \mathcal{F}_{\mathcal{O}_1})
=
\Hom_{\mathcal{O}_2}(
\mathcal{O}_2 \otimes_{p, \mathcal{O}_1} \mathcal{G},
\mathcal{F}
)
$$
换言之，上述标量限制函子与换环函子彼此伴随。
\end{lemma}

\begin{proof}
这由下述事实得出：对环同态 $A \to B$，标量限制函子与换环函子彼此伴随。
\end{proof}


\section{模层}
\label{sites-modules-section-sheaves-modules}

\begin{definition}
\label{sites-modules-definition-sheaf-modules}
设 $\mathcal{C}$ 为站点，$\mathcal{O}$ 为 $\mathcal{C}$ 上的环层。
\begin{enumerate}
\item 一个{\it $\mathcal{O}$-模层}是 $\mathcal{O}$-模预层 $\mathcal{F}$
（见定义 \ref{sites-modules-definition-presheaf-modules}），并且其底层阿贝尔群预层
$\mathcal{F}$ 是层。
\item 一个{\it $\mathcal{O}$-模层态射}就是 $\mathcal{O}$-模预层态射。
\item 给定 $\mathcal{O}$-模层 $\mathcal{F}$ 与 $\mathcal{G}$，以
$\Hom_\mathcal{O}(\mathcal{F}, \mathcal{G})$ 表示 $\mathcal{O}$-模层态射
组成的集合。
\item $\mathcal{O}$-模层的范畴记为 $\textit{Mod}(\mathcal{O})$。
\end{enumerate}
\end{definition}

\noindent
即使 $\mathcal{O}$ 只是环预层，这一定义也勉强说得通；不过我们不知道它在
何种例子中有用，并且当 $\mathcal{O}$ 不是环层时，将避免使用
“$\mathcal{O}$-模层”这一术语。



\section{模预层的层化}
\label{sites-modules-section-sheafification-presheaves-modules}

\begin{lemma}
\label{sites-modules-lemma-sheafification-presheaf-modules}
设 $\mathcal{C}$ 为站点，$\mathcal{O}$ 为 $\mathcal{C}$ 上的环预层，
$\mathcal{F}$ 为 $\mathcal{O}$-模预层。设 $\mathcal{O}^\#$ 为将
$\mathcal{O}$ 视为环预层时的层化；见《站点》第
\ref{sites-section-sheaves-algebraic-structures} 节。设 $\mathcal{F}^\#$ 为将
$\mathcal{F}$ 视为阿贝尔群预层时的层化。存在唯一的集合层映射
$$
\mathcal{O}^\# \times \mathcal{F}^\#
\longrightarrow
\mathcal{F}^\#
$$
使得图
$$
\xymatrix{
\mathcal{O} \times \mathcal{F} \ar[r] \ar[d] &
\mathcal{F} \ar[d] \\
\mathcal{O}^\# \times \mathcal{F}^\# \ar[r] &
\mathcal{F}^\#
}
$$
可交换，并使 $\mathcal{F}^\#$ 成为 $\mathcal{O}^\#$-模层。此外，若
$\mathcal{G}$ 是 $\mathcal{O}^\#$-模层，则任意 $\mathcal{O}$-模预层态射
$\mathcal{F} \to \mathcal{G}$（其靶是 $\mathcal{G}$ 限制为
$\mathcal{O}$-模后的预层）都唯一分解为
$\mathcal{F} \to \mathcal{F}^\# \to \mathcal{G}$，其中
$\mathcal{F}^\# \to \mathcal{G}$ 是 $\mathcal{O}^\#$-模态射。
\end{lemma}

\begin{proof}
从略。
\end{proof}

\noindent
这实际上意味着函子
$i : \textit{Mod}(\mathcal{O}^\#) \to \textit{PMod}(\mathcal{O})$
（结合标量限制与将层包含进预层）和该引理中的层化函子
${}^\# : \textit{PMod}(\mathcal{O}) \to \textit{Mod}(\mathcal{O}^\#)$
彼此伴随。公式为
$$
\Mor_{\textit{PMod}(\mathcal{O})}(\mathcal{F}, i\mathcal{G})
=
\Mor_{\textit{Mod}(\mathcal{O}^\#)}(\mathcal{F}^\#, \mathcal{G})
$$
一个重要情形是 $\mathcal{O}$ 已经为环层。此时公式成为
$$
\Mor_{\textit{PMod}(\mathcal{O})}(\mathcal{F}, i\mathcal{G})
=
\Mor_{\textit{Mod}(\mathcal{O})}(\mathcal{F}^\#, \mathcal{G})
$$
因为此时 $\mathcal{O} = \mathcal{O}^\#$。

\begin{lemma}
\label{sites-modules-lemma-sheafification-exact}
设 $\mathcal{C}$ 为站点，$\mathcal{O}$ 为 $\mathcal{C}$ 上的环预层。层化函子
$$
\textit{PMod}(\mathcal{O}) \longrightarrow \textit{Mod}(\mathcal{O}^\#), \quad
\mathcal{F} \longmapsto \mathcal{F}^\#
$$
正合。
\end{lemma}

\begin{proof}
这是因为层化 $\textit{PAb}(\mathcal{C}) \to \textit{Ab}(\mathcal{C})$
具有这一性质；见第 \ref{sites-modules-section-abelian-sheaves} 节的讨论。
\end{proof}

\noindent
设 $\mathcal{C}$ 为站点，$\mathcal{O}_1 \to \mathcal{O}_2$ 为
$\mathcal{C}$ 上环层的态射。在第 \ref{sites-modules-section-presheaves-modules} 节中，
我们定义了此情形所伴随的模预层标量限制函子与换环函子。

\medskip\noindent
若 $\mathcal{F}$ 是 $\mathcal{O}_2$-模层，则 $\mathcal{F}$ 的标量限制
$\mathcal{F}_{\mathcal{O}_1}$ 显然为 $\mathcal{O}_1$-模层。由此得到标量
限制函子
$$
\textit{Mod}(\mathcal{O}_2)
\longrightarrow
\textit{Mod}(\mathcal{O}_1)
$$

\medskip\noindent
另一方面，给定 $\mathcal{O}_1$-模层 $\mathcal{G}$，$\mathcal{O}_2$-模预层
$\mathcal{O}_2 \otimes_{p, \mathcal{O}_1} \mathcal{G}$ 一般不是层。因此按公式
定义{\it 张量积层} $\mathcal{O}_2 \otimes_{\mathcal{O}_1} \mathcal{G}$：
$$
\mathcal{O}_2 \otimes_{\mathcal{O}_1} \mathcal{G}
=
(\mathcal{O}_2 \otimes_{p, \mathcal{O}_1} \mathcal{G})^\#
$$
即取上述预层构造的层化。由此得到{\it 换环}函子
$$
\textit{Mod}(\mathcal{O}_1)
\longrightarrow
\textit{Mod}(\mathcal{O}_2)
$$

\begin{lemma}
\label{sites-modules-lemma-adjointness-tensor-restrict}
沿用上述 $X$、$\mathcal{O}_1$、$\mathcal{O}_2$、$\mathcal{F}$ 与
$\mathcal{G}$，存在典范双射
$$
\Hom_{\mathcal{O}_1}(\mathcal{G}, \mathcal{F}_{\mathcal{O}_1})
=
\Hom_{\mathcal{O}_2}(
\mathcal{O}_2 \otimes_{\mathcal{O}_1} \mathcal{G},
\mathcal{F}
)
$$
换言之，标量限制函子与换环函子彼此伴随。
\end{lemma}

\begin{proof}
这由引理 \ref{sites-modules-lemma-adjointness-tensor-restrict-presheaves} 以及下述事实得出：
$\Hom_{\mathcal{O}_2}(
\mathcal{O}_2 \otimes_{\mathcal{O}_1} \mathcal{G},
\mathcal{F}
)
=
\Hom_{\mathcal{O}_2}(
\mathcal{O}_2 \otimes_{p, \mathcal{O}_1} \mathcal{G},
\mathcal{F}
)$
因为 $\mathcal{F}$ 是层。
\end{proof}

\begin{lemma}
\label{sites-modules-lemma-epimorphism-modules}
设 $\mathcal{C}$ 为站点，$\mathcal{O} \to \mathcal{O}'$ 为环层的满态射，
$\mathcal{G}_1, \mathcal{G}_2$ 为 $\mathcal{O}'$-模层。则
$$
\Hom_{\mathcal{O}'}(\mathcal{G}_1, \mathcal{G}_2) =
\Hom_\mathcal{O}(\mathcal{G}_1, \mathcal{G}_2).
$$
换言之，标量限制函子
$\textit{Mod}(\mathcal{O}') \to \textit{Mod}(\mathcal{O})$ 充分忠实。
\end{lemma}

\begin{proof}
这是《代数》引理 \ref{algebra-lemma-epimorphism-modules} 的层版本，证明完全相同。
\end{proof}




\section{拓扑斯态射与模层}
\label{sites-modules-section-sheaves-modules-functorial}

\noindent
这里的一切材料都完全直接。我们以拓扑斯态射的情形陈述各项结果；当然，
它们对站点态射也成立。

\begin{lemma}
\label{sites-modules-lemma-pushforward-module}
设 $\mathcal{C}$、$\mathcal{D}$ 为站点，
$f : \Sh(\mathcal{C}) \to \Sh(\mathcal{D})$ 为拓扑斯态射，
$\mathcal{O}$ 为 $\mathcal{C}$ 上的环层，$\mathcal{F}$ 为
$\mathcal{O}$-模层。存在自然的集合层映射
$$
f_*\mathcal{O} \times f_*\mathcal{F}
\longrightarrow
f_*\mathcal{F}
$$
使 $f_*\mathcal{F}$ 成为 $f_*\mathcal{O}$-模层。此构造关于
$\mathcal{F}$ 具有函子性。
\end{lemma}

\begin{proof}
以 $\mu : \mathcal{O} \times \mathcal{F} \to \mathcal{F}$ 表示乘法映射。
回顾 $f_*$（作用于集合层时）左正合，因而与直积可交换。因此
$f_*\mu$ 正是所示映射。这证明了引理。
\end{proof}

\begin{lemma}
\label{sites-modules-lemma-pullback-module}
设 $\mathcal{C}$、$\mathcal{D}$ 为站点，
$f : \Sh(\mathcal{C}) \to \Sh(\mathcal{D})$ 为拓扑斯态射，
$\mathcal{O}$ 为 $\mathcal{D}$ 上的环层，$\mathcal{G}$ 为
$\mathcal{O}$-模层。存在自然的集合层映射
$$
f^{-1}\mathcal{O} \times f^{-1}\mathcal{G}
\longrightarrow
f^{-1}\mathcal{G}
$$
使 $f^{-1}\mathcal{G}$ 成为 $f^{-1}\mathcal{O}$-模层。此构造关于
$\mathcal{G}$ 具有函子性。
\end{lemma}

\begin{proof}
以 $\mu : \mathcal{O} \times \mathcal{G} \to \mathcal{G}$ 表示乘法映射。
回顾 $f^{-1}$（作用于集合层时）正合，因而与直积可交换。因此
$f^{-1}\mu$ 正是所示映射。这证明了引理。
\end{proof}

\begin{lemma}
\label{sites-modules-lemma-adjoint-push-pull-modules}
设 $\mathcal{C}$、$\mathcal{D}$ 为站点，
$f : \Sh(\mathcal{C}) \to \Sh(\mathcal{D})$ 为拓扑斯态射，
$\mathcal{O}$ 为 $\mathcal{D}$ 上的环层，$\mathcal{G}$ 为
$\mathcal{O}$-模层，$\mathcal{F}$ 为 $f^{-1}\mathcal{O}$-模层。则
$$
\Mor_{\textit{Mod}(f^{-1}\mathcal{O})}(f^{-1}\mathcal{G}, \mathcal{F})
=
\Mor_{\textit{Mod}(\mathcal{O})}(\mathcal{G}, f_*\mathcal{F}).
$$
这里使用引理 \ref{sites-modules-lemma-pullback-module} 与
\ref{sites-modules-lemma-pushforward-module}，并通过
$\mathcal{O} \to f_*f^{-1}\mathcal{O}$ 作标量限制，将
$f_*\mathcal{F}$ 视为 $\mathcal{O}$-模层。
\end{lemma}

\begin{proof}
首先注意到
$$
\Mor_{\textit{Ab}(\mathcal{C})}(f^{-1}\mathcal{G}, \mathcal{F})
=
\Mor_{\textit{Ab}(\mathcal{D})}(\mathcal{G}, f_*\mathcal{F}).
$$
，这是《站点》命题
\ref{sites-proposition-functoriality-algebraic-structures-topoi}。
设 $\alpha : f^{-1}\mathcal{G} \to \mathcal{F}$ 与
$\beta : \mathcal{G} \to f_*\mathcal{F}$ 是在上述公式下对应的阿贝尔层
态射。我们须证明，$\alpha$ 是 $f^{-1}\mathcal{O}$-线性的，当且仅当
$\beta$ 是 $\mathcal{O}$-线性的。例如，设 $\alpha$ 是
$f^{-1}\mathcal{O}$-线性的，则 $f_*\alpha$ 显然是
$f_*f^{-1}\mathcal{O}$-线性的，因而（由于标量限制是函子）也是
$\mathcal{O}$-线性的。故只需证明伴随映射
$\mathcal{G} \to f_*f^{-1}\mathcal{G}$ 是 $\mathcal{O}$-线性的。利用
$f_*$ 与 $f^{-1}$（作用于集合层时）都与直积可交换，这归结为证明图
$$
\xymatrix{
\mathcal{O} \times \mathcal{G} \ar[r] \ar[d] &
f_*f^{-1}(\mathcal{O} \times \mathcal{G}) \ar[d] \\
\mathcal{G} \ar[r] & f_*f^{-1}\mathcal{G}
}
$$
可交换。这是因为伴随映射
$\text{id}_{\Sh(\mathcal{D})} \to f_*f^{-1}$ 是函子变换。蕴含关系
$\beta$ 线性 $\Rightarrow$ $\alpha$ 线性的证明从略。
\end{proof}

\begin{lemma}
\label{sites-modules-lemma-adjoint-pull-push-modules}
设 $\mathcal{C}$、$\mathcal{D}$ 为站点，
$f : \Sh(\mathcal{C}) \to \Sh(\mathcal{D})$ 为拓扑斯态射，
$\mathcal{O}$ 为 $\mathcal{C}$ 上的环层，$\mathcal{F}$ 为
$\mathcal{O}$-模层，$\mathcal{G}$ 为 $f_*\mathcal{O}$-模层。则
$$
\Mor_{\textit{Mod}(\mathcal{O})}(
\mathcal{O} \otimes_{f^{-1}f_*\mathcal{O}} f^{-1}\mathcal{G}, \mathcal{F})
=
\Mor_{\textit{Mod}(f_*\mathcal{O})}(\mathcal{G}, f_*\mathcal{F}).
$$
这里使用引理 \ref{sites-modules-lemma-pullback-module} 与
\ref{sites-modules-lemma-pushforward-module}，并在张量积定义中使用典范映射
$f^{-1}f_*\mathcal{O} \to \mathcal{O}$。
\end{lemma}

\begin{proof}
注意到
$$
\Mor_{\textit{Mod}(\mathcal{O})}(
\mathcal{O} \otimes_{f^{-1}f_*\mathcal{O}} f^{-1}\mathcal{G}, \mathcal{F})
=
\Mor_{\textit{Mod}(f^{-1}f_*\mathcal{O})}(
f^{-1}\mathcal{G}, \mathcal{F}_{f^{-1}f_*\mathcal{O}})
$$
，这是引理 \ref{sites-modules-lemma-adjointness-tensor-restrict}。于是结论由引理
\ref{sites-modules-lemma-adjoint-push-pull-modules} 得出。
\end{proof}






\section{环拓扑斯态射与模}
\label{sites-modules-section-functoriality-modules}

\noindent
现在已经引入足够多的符号，可以定义模沿环拓扑斯态射的拉回与直像。

\begin{definition}
\label{sites-modules-definition-pushforward}
设
$(f, f^\sharp) :
(\Sh(\mathcal{C}), \mathcal{O}_\mathcal{C})
\to
(\Sh(\mathcal{D}), \mathcal{O}_\mathcal{D})$
为环拓扑斯态射或环站点态射。
\begin{enumerate}
\item 设 $\mathcal{F}$ 为 $\mathcal{O}_\mathcal{C}$-模层。定义
$\mathcal{F}$ 的{\it 直像}为如下 $\mathcal{O}_\mathcal{D}$-模层：作为阿贝尔
群层，它等于 $f_*\mathcal{F}$；其模结构由下述模结构经
$f^\sharp : \mathcal{O}_\mathcal{D} \to f_*\mathcal{O}_\mathcal{C}$
作标量限制而得：
$$
f_*\mathcal{O}_\mathcal{C} \times f_*\mathcal{F}
\longrightarrow
f_*\mathcal{F}
$$
这里使用引理 \ref{sites-modules-lemma-pushforward-module}。
\item 设 $\mathcal{G}$ 为 $\mathcal{O}_\mathcal{D}$-模层。定义其
{\it 拉回} $f^*\mathcal{G}$ 为由下式定义的
$\mathcal{O}_\mathcal{C}$-模层：
$$
f^*\mathcal{G}
=
\mathcal{O}_\mathcal{C} \otimes_{f^{-1}\mathcal{O}_\mathcal{D}}
f^{-1}\mathcal{G}
$$
其中环同态 $f^{-1}\mathcal{O}_\mathcal{D} \to \mathcal{O}_\mathcal{C}$
为 $f^\sharp$，模结构由引理 \ref{sites-modules-lemma-pullback-module} 给出。
\end{enumerate}
\end{definition}

\noindent
这样便定义了函子
\begin{eqnarray*}
f_* : \textit{Mod}(\mathcal{O}_\mathcal{C})
& \longrightarrow &
\textit{Mod}(\mathcal{O}_\mathcal{D}) \\
f^* : \textit{Mod}(\mathcal{O}_\mathcal{D})
& \longrightarrow &
\textit{Mod}(\mathcal{O}_\mathcal{C})
\end{eqnarray*}
关于这两个函子的最后一个结果是：它们确如预期彼此伴随。

\begin{lemma}
\label{sites-modules-lemma-adjoint-pullback-pushforward-modules}
设
$(f, f^\sharp) :
(\Sh(\mathcal{C}), \mathcal{O}_\mathcal{C})
\to
(\Sh(\mathcal{D}), \mathcal{O}_\mathcal{D})$
为环拓扑斯态射或环站点态射。设 $\mathcal{F}$ 为
$\mathcal{O}_\mathcal{C}$-模层，$\mathcal{G}$ 为
$\mathcal{O}_\mathcal{D}$-模层。存在典范双射
$$
\Hom_{\mathcal{O}_\mathcal{C}}(f^*\mathcal{G}, \mathcal{F})
=
\Hom_{\mathcal{O}_\mathcal{D}}(\mathcal{G}, f_*\mathcal{F}).
$$
换言之，函子 $f^*$ 是 $f_*$ 的左伴随。
\end{lemma}

\begin{proof}
这由前面的工作得出：
\begin{eqnarray*}
\Hom_{\mathcal{O}_\mathcal{C}}(f^*\mathcal{G}, \mathcal{F})
& = &
\Mor_{\textit{Mod}(\mathcal{O}_\mathcal{C})}(
\mathcal{O}_\mathcal{C}
\otimes_{f^{-1}\mathcal{O}_\mathcal{D}} f^{-1}\mathcal{G},
\mathcal{F}) \\
& = &
\Mor_{\textit{Mod}(f^{-1}\mathcal{O}_\mathcal{D})}(
f^{-1}\mathcal{G}, \mathcal{F}_{f^{-1}\mathcal{O}_\mathcal{D}}) \\
& = &
\Hom_{\mathcal{O}_\mathcal{D}}(\mathcal{G}, f_*\mathcal{F}).
\end{eqnarray*}
这里使用引理 \ref{sites-modules-lemma-adjointness-tensor-restrict} 与
\ref{sites-modules-lemma-adjoint-push-pull-modules}。
\end{proof}

\begin{lemma}
\label{sites-modules-lemma-push-pull-composition-modules}
设 $(f, f^\sharp) :
(\Sh(\mathcal{C}_1), \mathcal{O}_1)
\to (\Sh(\mathcal{C}_2), \mathcal{O}_2)$ 以及
$(g, g^\sharp) :
(\Sh(\mathcal{C}_2), \mathcal{O}_2) \to
(\Sh(\mathcal{C}_3), \mathcal{O}_3)$
为环拓扑斯态射。存在函子的典范同构
$(g \circ f)_* \cong g_* \circ f_*$ 以及
$(g \circ f)^* \cong f^* \circ g^*$.
\end{lemma}

\begin{proof}
这由定义即得。
\end{proof}





\section{模层的阿贝尔范畴}
\label{sites-modules-section-kernels}

\noindent
设 $(\Sh(\mathcal{C}), \mathcal{O})$ 为环拓扑斯，$\mathcal{F}$、
$\mathcal{G}$ 为 $\mathcal{O}$-模层；见《层》定义
\ref{sheaves-definition-sheaf-modules}。设
$\varphi, \psi : \mathcal{F} \to \mathcal{G}$ 为 $\mathcal{O}$-模层态射。
定义 $\varphi + \psi : \mathcal{F} \to \mathcal{G}$ 为将 $\varphi$ 与
$\psi$ 视作阿贝尔层态射时的和。这显然仍为 $\mathcal{O}$-模映射。也显然，
$\mathcal{O}$-模映射的复合关于此加法是双线性的。因此
$\textit{Mod}(\mathcal{O})$ 为预加性范畴；见《同调》定义
\ref{homology-definition-preadditive}。

\medskip\noindent
以 $0$ 表示这样的 $\mathcal{O}$-模层：对 $\mathcal{C}$ 的每个对象 $U$，
其值恒为 $\{0\}$。它显然既是 $\textit{Mod}(\mathcal{O})$ 的终对象，
也是始对象。给定 $\mathcal{O}$-模态射
$\varphi : \mathcal{F} \to \mathcal{G}$，下列条件等价：
(a) $\varphi$ 为零；(b) $\varphi$ 经过 $0$ 分解；
(c) $\varphi$ 在每个对象 $U$ 上的截面中为零。

\medskip\noindent
此外，给定一对 $\mathcal{O}$-模层 $\mathcal{F}$、$\mathcal{G}$，可定义直和为
$$
\mathcal{F} \oplus \mathcal{G} = \mathcal{F} \times \mathcal{G}
$$
，并取《同调》定义 \ref{homology-definition-direct-sum} 中显然的映射
$(i, j, p, q)$。因此 $\textit{Mod}(\mathcal{O})$ 是加性范畴；见《同调》定义
\ref{homology-definition-additive-category}。

\medskip\noindent
设 $\varphi : \mathcal{F} \to \mathcal{G}$ 为 $\mathcal{O}$-模态射。
可定义 $\Ker(\varphi)$ 为将 $\varphi$ 视作阿贝尔层映射时的核。由第
\ref{sites-modules-section-abelian-sheaves} 节，这是 $\mathcal{F}$ 的子层，其截面为
$$
\Ker(\varphi)(U) =
\{ s \in \mathcal{F}(U) \mid \varphi(s) = 0 \text{ 于 } \mathcal{G}(U)\}
$$
，其中 $U$ 遍历 $\mathcal{C}$ 的所有对象。容易看出，这确为
$\mathcal{O}$-模范畴中的核。换言之，态射
$\alpha : \mathcal{H} \to \mathcal{F}$ 经过 $\Ker(\varphi)$ 分解，当且仅当
$\varphi \circ \alpha = 0$。

\medskip\noindent
类似地，定义 $\Coker(\varphi)$ 为将 $\varphi$ 视作阿贝尔层映射时的余核。
存在唯一的乘法映射
$$
\mathcal{O} \times \Coker(\varphi) \longrightarrow \Coker(\varphi)
$$
使映射 $\mathcal{G} \to \Coker(\varphi)$ 成为 $\mathcal{O}$-模态射（验证从略）。
映射 $\mathcal{G} \to \Coker(\varphi)$ 是满射（将其视为集合层映射；见第
\ref{sites-modules-section-abelian-sheaves} 节）。为证明 $\Coker(\varphi)$ 是
$\textit{Mod}(\mathcal{O})$ 中的余核，注意：若
$\beta : \mathcal{G} \to \mathcal{H}$ 是满足
$\beta \circ \varphi = 0$ 的 $\mathcal{O}$-模态射，则对 $\mathcal{C}$ 的每个
对象 $U$，$\beta$ 诱导映射
$\mathcal{G}(U)/\varphi(\mathcal{F}(U)) \to \mathcal{H}(U)$。由层化的泛性质
（见《层》引理 \ref{sheaves-lemma-sheafification-presheaf-modules}），得到典范
映射 $\Coker(\varphi) \to \mathcal{H}$，使原来的 $\beta$ 等于复合
$\mathcal{G} \to \Coker(\varphi) \to \mathcal{H}$。由上述满射性，态射
$\Coker(\varphi) \to \mathcal{H}$ 唯一。

\begin{lemma}
\label{sites-modules-lemma-abelian}
设 $(\Sh(\mathcal{C}), \mathcal{O})$ 为环拓扑斯。范畴
$\textit{Mod}(\mathcal{O})$ 是阿贝尔范畴。遗忘函子
$\textit{Mod}(\mathcal{O}) \to \textit{Ab}(\mathcal{C})$
正合。因此，$\mathcal{O}$-模的核、余核及正合性分别对应于阿贝尔层的相应概念。
\end{lemma}

\begin{proof}
上面已经看到，$\textit{Mod}(\mathcal{O})$ 是具有核与余核的加性范畴，且
$\textit{Mod}(\mathcal{O}) \to \textit{Ab}(\mathcal{C})$ 保持核与余核。
依《同调》定义 \ref{homology-definition-abelian-category}，只须证明像与余像
一致。这显然成立，因为它在 $\textit{Ab}(\mathcal{C})$ 中成立。引理得证。
\end{proof}

\begin{lemma}
\label{sites-modules-lemma-limits-colimits}
设 $(\Sh(\mathcal{C}), \mathcal{O})$ 为环拓扑斯。
$\textit{Mod}(\mathcal{O})$ 中一切极限与余极限均存在，且遗忘函子
$\textit{Mod}(\mathcal{O}) \to \textit{Ab}(\mathcal{C})$
与它们可交换。此外，滤过余极限正合。
\end{lemma}

\begin{proof}
最后一个断言来自第一个断言，因为由引理
\ref{sites-modules-lemma-limits-colimits-abelian-sheaves}，滤过余极限在
$\textit{Ab}(\mathcal{C})$ 中正合。设
$\mathcal{I} \to \textit{Mod}(\mathcal{C})$，$i \mapsto \mathcal{F}_i$
为图，并设 $\lim_i \mathcal{F}_i$ 为该图在
$\textit{Ab}(\mathcal{C})$ 中的极限。由引理
\ref{sites-modules-lemma-limits-colimits-abelian-sheaves} 对此极限的描述，立即可见存在乘法
$$
\mathcal{O} \times \lim_i \mathcal{F}_i
\longrightarrow
\lim_i \mathcal{F}_i
$$
使 $\lim_i \mathcal{F}_i$ 成为 $\mathcal{O}$-模层。容易看出，这是该图在
$\textit{Mod}(\mathcal{C})$ 中的极限。设 $\colim_i \mathcal{F}_i$ 为该图在
$\textit{PAb}(\mathcal{C})$ 中的余极限。由引理
\ref{sites-modules-lemma-limits-colimits-abelian-presheaves} 的证明中对此余极限的描述，
立即可见存在乘法
$$
\mathcal{O} \times \colim_i \mathcal{F}_i
\longrightarrow
\colim_i \mathcal{F}_i
$$
使 $\colim_i \mathcal{F}_i$ 成为 $\mathcal{O}$-模预层。施行层化，得到
$\mathcal{O}$-模层 $(\colim_i \mathcal{F}_i)^\#$；见引理
\ref{sites-modules-lemma-sheafification-presheaf-modules}。容易看出，
$(\colim_i \mathcal{F}_i)^\#$ 是该图在
$\textit{Mod}(\mathcal{O})$ 中的余极限；又由引理
\ref{sites-modules-lemma-limits-colimits-abelian-sheaves}，遗忘 $\mathcal{O}$-模结构后，
它就是 $\textit{Ab}(\mathcal{C})$ 中的余极限。
\end{proof}

\noindent
极限与余极限的存在，使我们可以像《范畴论》第
\ref{categories-section-exact-functor} 节那样，用极限与余极限讨论定义在
$\mathcal{O}$-模范畴上的函子的正合性。关于用短正合列描述正合性，见《同调》
引理 \ref{homology-lemma-exact-functor}。

\begin{lemma}
\label{sites-modules-lemma-exactness-pushforward-pullback}
设 $f : (\Sh(\mathcal{C}), \mathcal{O}_\mathcal{C})
\to (\Sh(\mathcal{D}), \mathcal{O}_\mathcal{D})$
为环拓扑斯态射。
\begin{enumerate}
\item 函子 $f_*$ 左正合；事实上，它与一切极限可交换。
\item 函子 $f^*$ 右正合；事实上，它与一切余极限可交换。
\end{enumerate}
\end{lemma}

\begin{proof}
这是因为 $(f^*, f_*)$ 是一对伴随函子；见引理
\ref{sites-modules-lemma-adjoint-pullback-pushforward-modules}。另见《范畴论》第
\ref{categories-section-adjoint} 节。
\end{proof}

\begin{lemma}
\label{sites-modules-lemma-check-exactness-stalks}
设 $\mathcal{C}$ 为站点。若 $\{p_i\}_{i \in I}$ 是保守点族，则可在各点
$p_i$（$i \in I$）处的茎上检验阿贝尔层列的正合性。若 $\mathcal{C}$ 有
足够多的点，则阿贝尔层列的正合性可在茎上检验。
\end{lemma}

\begin{proof}
这由《站点》引理 \ref{sites-lemma-exactness-stalks} 立即得出。
\end{proof}





\section{直像的正合性}
\label{sites-modules-section-pushforward}

\noindent
下面给出若干关于直像之正合性质的技术性引理。

\begin{lemma}
\label{sites-modules-lemma-reflect-surjections}
设 $f : \Sh(\mathcal{C}) \to \Sh(\mathcal{D})$ 为拓扑斯态射。下列条件等价：
\begin{enumerate}
\item 对 $\textit{Ab}(\mathcal{C})$ 中所有 $\mathcal{F}$，
$f^{-1}f_*\mathcal{F} \to \mathcal{F}$ 均为满射；以及
\item $f_* : \textit{Ab}(\mathcal{C}) \to \textit{Ab}(\mathcal{D})$
反映满射。
\end{enumerate}
在此情形下，函子
$f_* : \textit{Ab}(\mathcal{C}) \to \textit{Ab}(\mathcal{D})$
忠实。
\end{lemma}

\begin{proof}
假设 (1)。设 $a : \mathcal{F} \to \mathcal{F}'$ 为 $\mathcal{C}$ 上阿贝尔层
的映射，且 $f_*a$ 为满射。由于 $f^{-1}$ 正合，可知
$f^{-1}f_*a : f^{-1}f_*\mathcal{F} \to f^{-1}f_*\mathcal{F}'$
为满射。结合 (1)，可知 $a$ 为满射。这表明 (2) 成立。

\medskip\noindent
假设 (2)。设 $\mathcal{F}$ 为 $\mathcal{C}$ 上阿贝尔层。我们须证明映射
$f^{-1}f_*\mathcal{F} \to \mathcal{F}$ 为满射。由 (2)，只需证明
$f_*f^{-1}f_*\mathcal{F} \to f_*\mathcal{F}$ 为满射。这确实成立，因为存在
作为单侧逆的典范映射 $f_*\mathcal{F} \to f_*f^{-1}f_*\mathcal{F}$。

\medskip\noindent
最后一个断言的证明从略。
\end{proof}

\begin{lemma}
\label{sites-modules-lemma-exactness}
设 $f : \Sh(\mathcal{C}) \to \Sh(\mathcal{D})$ 为拓扑斯态射。假设下列性质
中至少一个成立：
\begin{enumerate}
\item $f_*$ 将集合层的满射变为满射；
\item $f_*$ 将阿贝尔层的满射变为满射；
\item $f_*$ 与集合层上的余等化子可交换；
\item $f_*$ 与集合层上的推出可交换。
\end{enumerate}
则 $f_* : \textit{Ab}(\mathcal{C}) \to \textit{Ab}(\mathcal{D})$ 正合。
\end{lemma}

\begin{proof}
由于 $f_* : \textit{Ab}(\mathcal{C}) \to \textit{Ab}(\mathcal{D})$
是右伴随，我们已经知道，它将 $\mathcal{C}$ 上阿贝尔层的短正合列
$0 \to \mathcal{F}_1 \to \mathcal{F}_2 \to \mathcal{F}_3 \to 0$
变为正合列
$$
0 \to f_*\mathcal{F}_1 \to f_*\mathcal{F}_2 \to f_*\mathcal{F}_3
$$
；见《范畴论》第 \ref{categories-section-exact-functor}、
\ref{categories-section-adjoint} 节以及《同调》第
\ref{homology-section-functors} 节。因此只需证明映射
$f_*\mathcal{F}_2 \to f_*\mathcal{F}_3$ 为满射。若 (1) 或 (2) 成立，
这由定义即得。由《站点》引理 \ref{sites-lemma-exactness-properties}，(3) 或
(4) 都形式地蕴含 (1)，故这些情形也完成了证明。
\end{proof}

\begin{lemma}
\label{sites-modules-lemma-morphism-ringed-sites-almost-cocontinuous}
设 $f : \mathcal{D} \to \mathcal{C}$ 为连续函子
$u : \mathcal{C} \to \mathcal{D}$ 所伴随的站点态射。假设 $u$ 几乎余连续。则
\begin{enumerate}
\item $f_* : \textit{Ab}(\mathcal{D}) \to \textit{Ab}(\mathcal{C})$ 正合。
\item 若给定 $f^\sharp : f^{-1}\mathcal{O}_\mathcal{C} \to
\mathcal{O}_\mathcal{D}$，使 $f$ 成为环站点态射，则
$f_* : \textit{Mod}(\mathcal{O}_\mathcal{D}) \to
\textit{Mod}(\mathcal{O}_\mathcal{C})$ 正合。
\end{enumerate}
\end{lemma}

\begin{proof}
由引理 \ref{sites-modules-lemma-limits-colimits}，第 (2) 项来自第 (1) 项。第 (1) 项来自
《站点》引理 \ref{sites-lemma-morphism-of-sites-almost-cocontinuous} 与
\ref{sites-lemma-exactness-properties}。
\end{proof}





\section{下叹号函子的正合性}
\label{sites-modules-section-exactness-lower-shriek}

\noindent
设 $u : \mathcal{C} \to \mathcal{D}$ 为站点之间的函子。假设
\begin{enumerate}
\item[(a)] $u$ 余连续；且
\item[(b)] $u$ 连续。
\end{enumerate}
设 $g : \Sh(\mathcal{C}) \to \Sh(\mathcal{D})$ 为 $u$ 所伴随的拓扑斯态射；
见《站点》引理 \ref{sites-lemma-cocontinuous-morphism-topoi}。回顾
$g^{-1} = u^p$，即 $g^{-1}$ 由简单公式
$(g^{-1}\mathcal{G})(U) = \mathcal{G}(u(U))$ 给出；见《站点》引理
\ref{sites-lemma-when-shriek}。我们要证明
$g^{-1} : \textit{Ab}(\mathcal{D}) \to \textit{Ab}(\mathcal{C})$
具有左伴随 $g_!$。由《站点》引理 \ref{sites-lemma-when-shriek}，函子
$g^{Sh}_! = (u_p\ )^\#$ 是集合层上的左伴随。此外，我们知道
$g^{Sh}_!\mathcal{F}$ 是与预层
$$
V \longmapsto \colim_{V \to u(U)} \mathcal{F}(U)
$$
相伴的层，其中余极限取遍 $(\mathcal{I}_V^u)^{opp}$，并在集合范畴中计算。
因此下述定义很自然。

\begin{definition}
\label{sites-modules-definition-g-shriek}
设 $u : \mathcal{C} \to \mathcal{D}$ 满足上述 (a)、(b)。对
$\mathcal{F} \in \textit{PAb}(\mathcal{C})$，定义
{\it $g_{p!}\mathcal{F}$} 为预层
$$
V \longmapsto \colim_{V \to u(U)} \mathcal{F}(U)
$$
，其中取遍 $(\mathcal{I}_V^u)^{opp}$ 的余极限在 $\textit{Ab}$ 中计算。
对 $\mathcal{F} \in \textit{PAb}(\mathcal{C})$，令
{\it $g_!\mathcal{F} = (g_{p!}\mathcal{F})^\#$}。
\end{definition}

\noindent
之所以如此明确，是因为函子 $g^{Sh}_!$ 与 $g_!$ 不同。每当同时使用二者时，
都必须谨慎地区分它们。

\begin{lemma}
\label{sites-modules-lemma-g-shriek-adjoint}
函子 $g_{p!}$ 是 $u^p$ 的左伴随；函子 $g_!$ 是 $g^{-1}$ 的左伴随。换言之，公式
\begin{align*}
\Mor_{\textit{PAb}(\mathcal{C})}(\mathcal{F}, u^p\mathcal{G})
& =
\Mor_{\textit{PAb}(\mathcal{D})}(g_{p!}\mathcal{F}, \mathcal{G}), \\
\Mor_{\textit{Ab}(\mathcal{C})}(\mathcal{F}, g^{-1}\mathcal{G})
& =
\Mor_{\textit{Ab}(\mathcal{D})}(g_!\mathcal{F}, \mathcal{G})
\end{align*}
关于 $\mathcal{F}$ 与 $\mathcal{G}$ 双函子性地成立。
\end{lemma}

\begin{proof}
第二个公式形式地来自第一个公式，因为若 $\mathcal{F}$ 与 $\mathcal{G}$
为阿贝尔层，则
\begin{align*}
\Mor_{\textit{Ab}(\mathcal{C})}(\mathcal{F}, g^{-1}\mathcal{G})
& =
\Mor_{\textit{PAb}(\mathcal{D})}(g_{p!}\mathcal{F}, \mathcal{G}) \\
& =
\Mor_{\textit{Ab}(\mathcal{D})}(g_!\mathcal{F}, \mathcal{G})
\end{align*}
，这是层化的泛性质。

\medskip\noindent
为证明第一个公式，设 $\mathcal{F}$、$\mathcal{G}$ 为阿贝尔预层。我们将从
左端群到右端群构造映射，并略去验证这些映射互为逆映射。

\medskip\noindent
注意，存在阿贝尔预层的典范映射
$\mathcal{F} \to u^pg_{p!}\mathcal{F}$；它在 $U$ 上的截面中是自然映射
$\mathcal{F}(U) \to \colim_{u(U) \to u(U')} \mathcal{F}(U')$；见《站点》
引理 \ref{sites-lemma-recover}。给定映射
$\alpha : g_{p!}\mathcal{F} \to \mathcal{G}$，得到
$u^p\alpha : u^pg_{p!}\mathcal{F} \to u^p\mathcal{G}$，再将它与映射
$\mathcal{F} \to u^pg_{p!}\mathcal{F}$ 前复合。

\medskip\noindent
注意，存在阿贝尔预层的典范映射
$g_{p!}u^p\mathcal{G} \to \mathcal{G}$；它在 $V$ 上的截面中是自然映射
$\colim_{V \to u(U)} \mathcal{G}(u(U)) \to \mathcal{G}(V)$。它把与
$t : V \to u(U)$ 对应的直和项中的截面 $s \in u(U)$ 映为
$t^*s \in \mathcal{G}(V)$。因此，给定映射
$\beta : \mathcal{F} \to u^p\mathcal{G}$，得到映射
$g_{p!}\beta : g_{p!}\mathcal{F} \to g_{p!}u^p\mathcal{G}$，再将它与上述映射
$g_{p!}u^p\mathcal{G} \to \mathcal{G}$ 后复合。
\end{proof}

\begin{lemma}
\label{sites-modules-lemma-exactness-lower-shriek}
设 $\mathcal{C}$ 与 $\mathcal{D}$ 为站点，$u : \mathcal{C} \to \mathcal{D}$
为函子。假设
\begin{enumerate}
\item[(a)] $u$ 余连续；
\item[(b)] $u$ 连续；且
\item[(c)] $\mathcal{C}$ 中存在纤维积与等化子，并且 $u$ 与它们可交换。
\end{enumerate}
在此情形下，函子
$g_! : \textit{Ab}(\mathcal{C}) \to \textit{Ab}(\mathcal{D})$
正合。
\end{lemma}

\begin{proof}
比较《站点》引理 \ref{sites-lemma-preserve-equalizers}。假设 (a)、(b)、(c)。
我们已经知道 $g_!$ 右正合，因为它是左伴随；见《范畴论》引理
\ref{categories-lemma-exact-adjoint} 与《同调》第
\ref{homology-section-functors} 节。我们有 $g_! = (g_{p!}\ )^\#$。须证明
$g_!$ 将阿贝尔层的单射映为阿贝尔预层的单射。回顾阿贝尔预层的层化正合；
见引理 \ref{sites-modules-lemma-limits-colimits-abelian-sheaves}。因此，只需证明
$g_{p!}$ 将阿贝尔预层的单射映为阿贝尔预层的单射。为此，只需《站点》第
\ref{sites-section-functoriality-PSh} 节中取遍范畴
$(\mathcal{I}_V^u)^{opp}$ 的余极限把图之间的单射映为单射。这由《站点》引理
\ref{sites-lemma-almost-directed} 与《代数》引理
\ref{algebra-lemma-almost-directed-colimit-exact} 得出。
\end{proof}

\begin{lemma}
\label{sites-modules-lemma-back-and-forth}
设 $\mathcal{C}$ 与 $\mathcal{D}$ 为站点，$u : \mathcal{C} \to \mathcal{D}$
为函子。假设
\begin{enumerate}
\item[(a)] $u$ 余连续；
\item[(b)] $u$ 连续；且
\item[(c)] $u$ 充分忠实。
\end{enumerate}
对上述 $g_!, g^{-1}, g_*$，典范映射
$\mathcal{F} \to g^{-1}g_!\mathcal{F}$ 与
$g^{-1}g_*\mathcal{F} \to \mathcal{F}$
对 $\mathcal{C}$ 上一切阿贝尔层 $\mathcal{F}$ 均为同构。
\end{lemma}

\begin{proof}
由《站点》引理 \ref{sites-lemma-back-and-forth}，以及阿贝尔层的逆像与直像
同其底层集合层的逆像与直像一致这一事实，映射
$g^{-1}g_*\mathcal{F} \to \mathcal{F}$ 为同构。

\medskip\noindent
取 $U \in \Ob(\mathcal{C})$。我们将证明
$g^{-1}g_!\mathcal{F}(U) = \mathcal{F}(U)$。先注意到
$g^{-1}g_!\mathcal{F}(U) = g_!\mathcal{F}(u(U))$。因此只需证明
$g_!\mathcal{F}(u(U)) = \mathcal{F}(U)$。我们知道，$g_!\mathcal{F}$ 是与
预层 $g_{p!}\mathcal{F}$ 相伴的阿贝尔层，而后者由法则
$$
V \longmapsto \colim_{V \to u(U')} \mathcal{F}(U')
$$
定义，其中余极限在 $\textit{Ab}$ 中计算。若 $V = u(U)$，则由于 $u$ 充分
忠实，此余极限取遍 $U \to U'$。故得到
$g_{p!}\mathcal{F}(u(U) = \mathcal{F}(U)$。由于 $u$ 既余连续又连续，
$\mathcal{D}$ 中 $u(U)$ 的任意覆盖都可由 $\mathcal{D}$ 的覆盖（！）
$\{u(U_i) \to u(U)\}$ 加细，其中 $\{U_i \to U\}$ 是 $\mathcal{C}$ 中的覆盖。
这还蕴含 $(g_{p!}\mathcal{F})^+(u(U)) = \mathcal{F}(U)$，因为在定义
$(g_{p!}\mathcal{F})^+$ 于 $u(U)$ 的值的余极限中，可限制到上述覆盖
$\{u(U_i) \to u(U)\}$ 所成的共终系统。于是对 $\mathcal{C}$ 的所有对象 $U$，
同样有 $(g_{p!}\mathcal{F})^+(u(U)) = \mathcal{F}(U)$。再重复一次此论证，
对 $\mathcal{C}$ 的所有对象 $U$ 得到
$(g_{p!}\mathcal{F})^\#(u(U)) = \mathcal{F}(U)$。这给出所需等式
$g^{-1}g_!\mathcal{F} = \mathcal{F}$。
\end{proof}

\begin{remark}
\label{sites-modules-remark-no-extension}
一般而言，当 $g$ 是环拓扑斯态射（的一部分）时，函子 $g_!$ 不能延伸到模
范畴。事实上，给定任意环同态 $A \to B$，函子
$M \mapsto B \otimes_A M$ 具有右伴随（标量限制），但一般没有左伴随
（因为左伴随的存在将蕴含 $A \to B$ 平坦）。下面将在第
\ref{sites-modules-section-localize} 节看到：对站点的局部化，可以在模层上定义 $j_!$。
我们将在第 \ref{sites-modules-section-lower-shriek-modules} 节以更大的一般性讨论此事。
\end{remark}

\begin{lemma}
\label{sites-modules-lemma-have-left-adjoint}
设 $\mathcal{C}$ 与 $\mathcal{D}$ 为站点，
$g : \Sh(\mathcal{C}) \to \Sh(\mathcal{D})$ 为连续且余连续的函子
$u : \mathcal{C} \to \mathcal{D}$ 所伴随的拓扑斯态射。
\begin{enumerate}
\item 若 $u$ 具有左伴随 $w$，则 $g_!$ 与 $g_!^{\Sh}$ 在底层集合层上
一致，且 $g_!$ 正合。
\item 若此外 $w$ 余连续，则 $g_! = h^{-1}$ 且 $g^{-1} = h_*$，其中
$h : \Sh(\mathcal{D}) \to \Sh(\mathcal{C})$ 是 $w$ 所伴随的拓扑斯态射。
\end{enumerate}
\end{lemma}

\begin{proof}
本引理对应于《站点》引理 \ref{sites-lemma-have-left-adjoint}。由《站点》引理
\ref{sites-lemma-adjoint-functors}，范畴 $\mathcal{I}_V^u$ 具有始对象。因此，
取遍 $(\mathcal{I}_V^u)^{opp}$ 的阿贝尔群系统之余极限的底层集合，就是底层
集合的余极限。结合定义 \ref{sites-modules-definition-g-shriek} 中对 $g_!$ 的构造，遂得
$g_!^{\Sh}$ 与 $g_!$ 一致。正合性及 (2) 立即来自《站点》引理
\ref{sites-lemma-have-left-adjoint} 中的相应断言。
\end{proof}





\section{模的整体类型}
\label{sites-modules-section-global}

\begin{definition}
\label{sites-modules-definition-global}
设 $(\Sh(\mathcal{C}), \mathcal{O})$ 为环拓扑斯，$\mathcal{F}$ 为
$\mathcal{O}$-模层。
\begin{enumerate}
\item 若 $\mathcal{F}$ 作为 $\mathcal{O}$-模层同构于形如
$\bigoplus_{i \in I} \mathcal{O}$ 的层，则称 $\mathcal{F}$ 为
{\it 自由 $\mathcal{O}$-模层}。
\item 若 $\mathcal{F}$ 作为 $\mathcal{O}$-模层同构于形如
$\bigoplus_{i \in I} \mathcal{O}$ 的层，其中指标集 $I$ 有限，则称
$\mathcal{F}$ {\it 有限自由}。
\item 若存在从自由 $\mathcal{O}$-模层到 $\mathcal{F}$ 的满射
$$
\bigoplus\nolimits_{i \in I} \mathcal{O} \longrightarrow \mathcal{F}
$$
，则称 $\mathcal{F}$ {\it 由整体截面生成}。
\item 给定 $r \geq 0$，若存在满射
$\mathcal{O}^{\oplus r} \to \mathcal{F}$，则称 $\mathcal{F}$
{\it 由 $r$ 个整体截面生成}。
\item 若对某个 $r \geq 0$，$\mathcal{F}$ 由 $r$ 个整体截面生成，则称
$\mathcal{F}$ {\it 由有限多个整体截面生成}。
\item 若存在 $\mathcal{O}$-模层正合列
$$
\bigoplus\nolimits_{j \in J} \mathcal{O} \longrightarrow
\bigoplus\nolimits_{i \in I} \mathcal{O} \longrightarrow
\mathcal{F} \longrightarrow 0
$$
，则称 $\mathcal{F}$ 具有{\it 整体表示}。
\item 若存在 $\mathcal{O}$-模层正合列
$$
\bigoplus\nolimits_{j \in J} \mathcal{O} \longrightarrow
\bigoplus\nolimits_{i \in I} \mathcal{O} \longrightarrow
\mathcal{F} \longrightarrow 0
$$
，其中 $I$ 与 $J$ 为有限集，则称 $\mathcal{F}$ 具有{\it 整体有限表示}。
\end{enumerate}
\end{definition}

\noindent
注意，对任意集合 $I$，直和 $\bigoplus_{i \in I} \mathcal{O}$ 存在
（引理 \ref{sites-modules-lemma-limits-colimits}），并且它是预层
$U \mapsto \bigoplus_{i \in I} \mathcal{O}(U)$ 的层化。该模层称为
{\it 集合 $I$ 上的自由 $\mathcal{O}$-模层}。

\begin{lemma}
\label{sites-modules-lemma-global-pullback}
设
$(f, f^\sharp) :
(\Sh(\mathcal{C}), \mathcal{O}_\mathcal{C})
\to
(\Sh(\mathcal{D}), \mathcal{O}_\mathcal{D})$
为环拓扑斯态射，$\mathcal{F}$ 为 $\mathcal{O}_\mathcal{D}$-模层。
\begin{enumerate}
\item 若 $\mathcal{F}$ 自由，则 $f^*\mathcal{F}$ 自由。
\item 若 $\mathcal{F}$ 有限自由，则 $f^*\mathcal{F}$ 有限自由。
\item 若 $\mathcal{F}$ 由整体截面生成，则 $f^*\mathcal{F}$ 由整体截面生成。
\item 给定 $r \geq 0$，若 $\mathcal{F}$ 由 $r$ 个整体截面生成，则
$f^*\mathcal{F}$ 由 $r$ 个整体截面生成。
\item 若 $\mathcal{F}$ 由有限多个整体截面生成，则 $f^*\mathcal{F}$ 亦然。
\item 若 $\mathcal{F}$ 具有整体表示，则 $f^*\mathcal{F}$ 亦然。
\item 若 $\mathcal{F}$ 具有整体有限表示，则 $f^*\mathcal{F}$ 亦然。
\end{enumerate}
\end{lemma}

\begin{proof}
这是因为 $f^*$ 与任意余极限可交换（引理
\ref{sites-modules-lemma-exactness-pushforward-pullback}），且
$f^*\mathcal{O}_\mathcal{D} = \mathcal{O}_\mathcal{C}$。
\end{proof}






\section{模的内蕴性质}
\label{sites-modules-section-intrinsic}

\noindent
设 $\mathcal{P}$ 为环拓扑斯上的模层之性质。若每当
$(f, f^\sharp) :
(\Sh(\mathcal{C}'), \mathcal{O}')
\to
(\Sh(\mathcal{C}), \mathcal{O})$
是环拓扑斯等价时，都有
$\mathcal{P}(\mathcal{F}) \Leftrightarrow \mathcal{P}(f^*\mathcal{F})$
，则称 $\mathcal{P}$ 为{\it 内蕴性质}。例如，自由性是内蕴的。事实上，
集合 $I$ 上的自由 $\mathcal{O}$-模层由下述性质刻画：
$$
\Mor_{\textit{Mod}(\mathcal{O})}(
\bigoplus\nolimits_{i \in I} \mathcal{O},
\mathcal{F})
=
\prod\nolimits_{i \in I} \Mor_{\Sh(\mathcal{C})}(\{*\},
\mathcal{F})
$$
其中 $\mathcal{F}$ 遍历 $\textit{Mod}(\mathcal{O})$。也可以利用引理
\ref{sites-modules-lemma-global-pullback} 看出自由性是内蕴的。事实上，出于同一理由，
定义 \ref{sites-modules-definition-global} 中的每个性质都是内蕴的。接下来如何定义
$\mathcal{O}$-模的其他内蕴性质？

\medskip\noindent
引理 \ref{sites-modules-lemma-morphism-ringed-topoi-comes-from-morphism-ringed-sites}
的要点如下：假设要定义拓扑斯上 $\mathcal{O}$-模的内蕴性质 $\mathcal{P}$，
可按以下步骤进行：
\begin{enumerate}
\item 给定任意站点 $\mathcal{C}$、$\mathcal{C}$ 上任意环层
$\mathcal{O}$ 及任意 $\mathcal{O}$-模层 $\mathcal{F}$，定义相应性质
$\mathcal{P}(\mathcal{C}, \mathcal{O}, \mathcal{F})$。
\item 对任意一对站点 $\mathcal{C}$、$\mathcal{C}'$，任意特殊余连续函子
$u : \mathcal{C} \to \mathcal{C}'$，$\mathcal{C}$ 上任意环层
$\mathcal{O}$ 及任意 $\mathcal{O}$-模层 $\mathcal{F}$，证明
$$
\mathcal{P}(\mathcal{C}, \mathcal{O}, \mathcal{F})
\Leftrightarrow
\mathcal{P}(\mathcal{C}', g_*\mathcal{O}, g_*\mathcal{F})
$$
，其中 $g : \Sh(\mathcal{C}) \to \Sh(\mathcal{C}')$ 是 $u$ 所伴随的
拓扑斯等价。
\end{enumerate}
在此情形下，给定任意环拓扑斯 $(\Sh(\mathcal{C}), \mathcal{O})$ 与任意
$\mathcal{O}$-模层 $\mathcal{F}$，若
$\mathcal{P}(\mathcal{C}, \mathcal{O}, \mathcal{F})$ 成立，就简称
$\mathcal{F}$ 具有性质 $\mathcal{P}$。引理
\ref{sites-modules-lemma-morphism-ringed-topoi-comes-from-morphism-ringed-sites} 与上述 (2)
共同保证此定义良定。

\medskip\noindent
此外，同一引理 \ref{sites-modules-lemma-morphism-ringed-topoi-comes-from-morphism-ringed-sites}
还保证，若进一步有
\begin{enumerate}
\item[(3)] 对任意环站点态射
$(f, f^\sharp) :
(\mathcal{C}, \mathcal{O}_\mathcal{C})
\to
(\mathcal{D}, \mathcal{O}_\mathcal{D})$
，其中 $f$ 由满足《站点》命题 \ref{sites-proposition-get-morphism} 之假设的
函子 $u : \mathcal{D} \to \mathcal{C}$ 给出，并且对任意
$\mathcal{O}_\mathcal{D}$-模层 $\mathcal{G}$，都有
$$
\mathcal{P}(\mathcal{D}, \mathcal{O}_\mathcal{D}, \mathcal{F})
\Rightarrow
\mathcal{P}(\mathcal{C}, \mathcal{O}_\mathcal{C}, f^*\mathcal{F})
$$
\end{enumerate}
，则 $\mathcal{P}$ 在模沿任意环拓扑斯态射拉回时得以保持。

\medskip\noindent
以下各节将用此方法看出：局部自由、局部由截面生成、局部由 $r$ 个截面生成、
有限型、有限表示、拟凝聚与凝聚，都是模的内蕴性质。

\medskip\noindent
与其采用此处概述的繁琐过程，也许为这些概念寻找内蕴定义更令人满意。另一方面，
在许多要应用这些定义的几何情形中，给定的是一个具体环站点与一个具体模层；
因此，已有适合这套语言的定义也很方便。




\section{环站点的局部化}
\label{sites-modules-section-localize}

\noindent
设 $(\mathcal{C}, \mathcal{O})$ 为环站点，$U \in \Ob(\mathcal{C})$。
我们说明《站点》第 \ref{sites-section-localize} 节中各结果在此情形下的对应物。

\medskip\noindent
以 $\mathcal{O}_U = j_U^{-1}\mathcal{O}$ 表示 $\mathcal{O}$ 到站点
$\mathcal{C}/U$ 的限制。它由简单法则
$\mathcal{O}_U(V/U) = \mathcal{O}(V)$ 描述。采用此符号后，局部化态射
$j_U$ 成为环拓扑斯态射
$$
(j_U, j_U^\sharp) :
(\Sh(\mathcal{C}/U), \mathcal{O}_U)
\longrightarrow
(\Sh(\mathcal{C}), \mathcal{O})
$$
；也就是说，取 $j_U^\sharp : j_U^{-1}\mathcal{O} \to \mathcal{O}_U$
为恒等映射。此外，模的直像与拉回具有下述描述。

\begin{definition}
\label{sites-modules-definition-localize-ringed-site}
设 $(\mathcal{C}, \mathcal{O})$ 为环站点，$U \in \Ob(\mathcal{C})$。
\begin{enumerate}
\item 环站点 $(\mathcal{C}/U, \mathcal{O}_U)$ 称为
{\it 环站点 $(\mathcal{C}, \mathcal{O})$ 在对象 $U$ 处的局部化}。
\item 环拓扑斯态射
$(j_U, j_U^\sharp) :
(\Sh(\mathcal{C}/U), \mathcal{O}_U)
\to
(\Sh(\mathcal{C}), \mathcal{O})$
称为{\it 局部化态射}。
\item 函子
$j_{U*} : \textit{Mod}(\mathcal{O}_U) \to \textit{Mod}(\mathcal{O})$
称为{\it 直像函子}。
\item 对 $\mathcal{C}$ 上的 $\mathcal{O}$-模层 $\mathcal{F}$，层
$j_U^*\mathcal{F}$ 称为{\it $\mathcal{F}$ 到 $\mathcal{C}/U$ 的限制}。
有时将它记为 $\mathcal{F}|_{\mathcal{C}/U}$，甚至简记为
$\mathcal{F}|_U$。它由简单法则
$j_U^*(\mathcal{F})(X/U) = \mathcal{F}(X)$ 描述。
\item 限制函子的左伴随
$j_{U!} : \textit{Mod}(\mathcal{O}_U) \to \textit{Mod}(\mathcal{O})$
称为{\it 以零延拓}。由引理 \ref{sites-modules-lemma-extension-by-zero} 与
\ref{sites-modules-lemma-extension-by-zero-exact}，它存在且正合。
\end{enumerate}
\end{definition}

\noindent
与拓扑情形相同（见《层》第 \ref{sheaves-section-open-immersions} 节），
以零延拓函子 $j_{U!}$ 不同于在集合层上定义的以空集延拓 $j_{U!}$。
下面的引理定义以零延拓。

\begin{lemma}
\label{sites-modules-lemma-extension-by-zero}
设 $(\mathcal{C}, \mathcal{O})$ 为环站点，$U \in \Ob(\mathcal{C})$。限制函子
$j_U^* : \textit{Mod}(\mathcal{O}) \to \textit{Mod}(\mathcal{O}_U)$
具有左伴随
$j_{U!} : \textit{Mod}(\mathcal{O}_U) \to \textit{Mod}(\mathcal{O})$.
因此
$$
\Mor_{\textit{Mod}(\mathcal{O}_U)}(\mathcal{G}, j_U^*\mathcal{F})
=
\Mor_{\textit{Mod}(\mathcal{O})}(j_{U!}\mathcal{G}, \mathcal{F})
$$
对 $\mathcal{F} \in \Ob(\textit{Mod}(\mathcal{O}))$ 与
$\mathcal{G} \in \Ob(\textit{Mod}(\mathcal{O}_U))$ 成立。此外，
$\mathcal{G}$ 的以零延拓 $j_{U!}\mathcal{G}$ 是与预层
$$
V
\longmapsto
\bigoplus\nolimits_{\varphi \in \Mor_\mathcal{C}(V, U)}
\mathcal{G}(V \xrightarrow{\varphi} U)
$$
相伴的层，其中限制映射与 $\mathcal{O}$-模结构均取显然者。
\end{lemma}

\begin{proof}
预层上的 $\mathcal{O}$-模结构定义如下。若 $f \in \mathcal{O}(V)$ 且
$s \in \mathcal{G}(V \xrightarrow{\varphi} U)$，则定义
$f \cdot s = fs$，其中
$f \in \mathcal{O}_U(\varphi : V \to U) = \mathcal{O}(V)$
（因为 $\mathcal{O}_U$ 是 $\mathcal{O}$ 到 $\mathcal{C}/U$ 的限制）。

\medskip\noindent
类似地，设 $\alpha : \mathcal{G} \to \mathcal{F}|_U$ 为
$\mathcal{O}_U$-模态射。此时可定义从引理中的预层到 $\mathcal{F}$ 的映射
$$
\bigoplus\nolimits_{\varphi \in \Mor_\mathcal{C}(V, U)}
\mathcal{G}(V \xrightarrow{\varphi} U)
\longrightarrow
\mathcal{F}(V)
$$
，法则为：将 $s \in \mathcal{G}(V \xrightarrow{\varphi} U)$ 映为
$\alpha(s) \in \mathcal{F}(V)$。这显然为 $\mathcal{O}$-线性的，因而由模层化
的性质（引理 \ref{sites-modules-lemma-sheafification-presheaf-modules}）诱导
$\mathcal{O}$-模态射 $\alpha' : j_{U!}\mathcal{G} \to \mathcal{F}$。

\medskip\noindent
反之，设 $\beta : j_{U!}\mathcal{G} \to \mathcal{F}$ 为
$\mathcal{O}$-模映射。回顾《站点》第 \ref{sites-section-localize} 节：在集合
层上存在以空集延拓
$j^{Sh}_{U!} : \Sh(\mathcal{C}/U) \to \Sh(\mathcal{C})$
，它是 $j_U^{-1}$ 的左伴随。此外，$j^{Sh}_{U!}\mathcal{G}$ 是与预层
$$
V
\longmapsto
\coprod\nolimits_{\varphi \in \Mor_\mathcal{C}(V, U)}
\mathcal{G}(V \xrightarrow{\varphi} U)
$$
相伴的层。因此存在集合层的自然映射
$j^{Sh}_{U!}\mathcal{G} \to j_{U!}\mathcal{G}$。将 $\beta$ 与此映射前复合，
得到集合层映射 $j^{Sh}_{U!}\mathcal{G} \to \mathcal{F}$；由伴随性，它对应于
集合层映射 $\beta' : \mathcal{G} \to \mathcal{F}|_U$。我们断言
$\beta'$ 是 $\mathcal{O}_U$-线性的。事实上，设
$\varphi : V \to U$ 为 $\mathcal{C}/U$ 的对象，
$s, s' \in \mathcal{G}(\varphi : V \to U)$，且
$f \in \mathcal{O}(V) = \mathcal{O}_U(\varphi : V \to U)$。由上述讨论，
$\mathcal{F}|_U(\varphi : V \to U)$ 中的 $\beta'(s + s')$、相应地
$\beta'(fs)$，分别对应于 $\mathcal{F}(V)$ 中的 $\beta(s + s')$、
相应地 $\beta(fs)$。由于 $\beta$ 是同态，结论成立。

\medskip\noindent
为完成引理的证明，须证明构造 $\alpha \mapsto \alpha'$ 与
$\beta \mapsto \beta'$ 互为逆构造。验证从略。
\end{proof}

\noindent
注意，在定义 \ref{sites-modules-definition-localize-ringed-site} 的情形中，对每个
$\mathcal{O}$-模 $\mathcal{F}$，都有
\begin{equation}
\label{sites-modules-equation-map-lower-shriek-OU-into-module}
\Hom_\mathcal{O}(j_{U!}\mathcal{O}_U, \mathcal{F}) =
\Hom_{\mathcal{O}_U}(\mathcal{O}_U, j_U^*\mathcal{F}) =
\mathcal{F}(U)
\end{equation}
。事实上，第一个等号来自 $j_{U!}$ 与 $j_U^*$ 的伴随性，第二个等号则来自
{\small $\Hom_{\mathcal{O}_U}(\mathcal{O}_U, j_U^*\mathcal{F}) =
j_U^*\mathcal{F}(U/U)$}；又有
{\small $j_U^*\mathcal{F}(U/U) = \mathcal{F}|_U(U/U) = \mathcal{F}(U)$}。

\begin{lemma}
\label{sites-modules-lemma-extension-by-zero-exact}
设 $(\mathcal{C}, \mathcal{O})$ 是环站点，且 $U \in \Ob(\mathcal{C})$。函子
$j_{U!} : \textit{Mod}(\mathcal{O}_U) \to \textit{Mod}(\mathcal{O})$
是正合的。
\end{lemma}

\begin{proof}
由于 $j_{U!}$ 是 $j_U^*$ 的左伴随，它是右正合的（见
《范畴》，引理 \ref{categories-lemma-exact-adjoint}，以及
《同调代数》，第 \ref{homology-section-functors} 节）。
因此，只需证明：若 $\mathcal{G}_1 \to \mathcal{G}_2$ 是
$\mathcal{O}_U$-模的单射，则
$j_{U!}\mathcal{G}_1 \to j_{U!}\mathcal{G}_2$ 也是单射。
在对象 $V$ 上的预层截面映射（如引理 \ref{sites-modules-lemma-extension-by-zero} 中所述）为
$$
\bigoplus\nolimits_{\varphi \in \Mor_\mathcal{C}(V, U)}
\mathcal{G}_1(V \xrightarrow{\varphi} U)
\longrightarrow
\bigoplus\nolimits_{\varphi \in \Mor_\mathcal{C}(V, U)}
\mathcal{G}_2(V \xrightarrow{\varphi} U)
$$
，由假设它是单射。根据引理 \ref{sites-modules-lemma-sheafification-exact}，层化是正合的，
故 $j_{U!}\mathcal{G}_1 \to j_{U!}\mathcal{G}_2$ 是单射，结论成立。
\end{proof}

\begin{lemma}
\label{sites-modules-lemma-j-shriek-reflects-exactness}
设 $(\mathcal{C}, \mathcal{O})$ 是环站点，且 $U \in \Ob(\mathcal{C})$。
$\mathcal{O}_U$-模的复形
$\mathcal{G}_1 \to \mathcal{G}_2 \to \mathcal{G}_3$ 正合，当且仅当
$j_{U!}\mathcal{G}_1 \to j_{U!}\mathcal{G}_2 \to j_{U!}\mathcal{G}_3$
作为 $\mathcal{O}$-模列正合。
\end{lemma}

\begin{proof}
由引理 \ref{sites-modules-lemma-extension-by-zero-exact}，我们已经知道 $j_{U!}$ 正合。
因此，只需证明
$j_{U!} :  \textit{Mod}(\mathcal{O}_U) \to \textit{Mod}(\mathcal{O})$
反映单射和满射。

\medskip\noindent
对 $\textit{Mod}(\mathcal{O}_U)$ 中的每个 $\mathcal{G}$，有该伴随的单位
$\mathcal{G} \to j_U^*j_{U!}\mathcal{G}$。我们断言此映射是层的单射。
事实上，考察引理 \ref{sites-modules-lemma-extension-by-zero} 的构造可知，此映射是如下规则的层化：
它将 $\mathcal{C}/U$ 的对象 $V/U$ 送到单射
$$
\mathcal{G}(V/U) \longrightarrow
\bigoplus\nolimits_{\varphi \in \Mor_\mathcal{C}(V, U)}
\mathcal{G}(V \xrightarrow{\varphi} U)
$$
，即嵌入到对应于结构态射 $V \to U$ 的直和项中。由于层化正合，断言成立。
略去若干细节。

\medskip\noindent
若 $\mathcal{G} \to \mathcal{G}'$ 是 $\mathcal{O}_U$-模映射，且
$j_{U!}\mathcal{G} \to j_{U!}\mathcal{G}'$ 是单射，则
$j_U^*j_{U!}\mathcal{G} \to j_U^*j_{U!}\mathcal{G}'$ 是单射
（限制是正合的），从而 $\mathcal{G} \to j_U^*j_{U!}\mathcal{G}'$
是单射，进而 $\mathcal{G} \to \mathcal{G}'$ 是单射。
所以 $j_{U!}$ 反映单射。

\medskip\noindent
设 $a : \mathcal{G} \to \mathcal{G}'$ 是 $\mathcal{O}_U$-模映射，且
$j_{U!}\mathcal{G} \to j_{U!}\mathcal{G}'$ 是满射。令 $\mathcal{H}$ 为
$a$ 的余核。由于 $j_{U!}$ 正合，有 $j_{U!}\mathcal{H} = 0$。
由上可知，映射 $\mathcal{H} \to j^*_U j_{U!}\mathcal{H}$ 是单射，
故 $\mathcal{H} = 0$，正合所需。
\end{proof}

\begin{lemma}
\label{sites-modules-lemma-relocalize}
设 $(\mathcal{C}, \mathcal{O})$ 是环站点，且 $f : V \to U$ 是
$\mathcal{C}$ 的态射。则存在环拓扑斯的交换图
$$
\xymatrix{
(\Sh(\mathcal{C}/V), \mathcal{O}_V)
\ar[rd]_{(j_V, j_V^\sharp)} \ar[rr]_{(j, j^\sharp)} & &
(\Sh(\mathcal{C}/U), \mathcal{O}_U)
\ar[ld]^{(j_U, j_U^\sharp)} \\
& (\Sh(\mathcal{C}), \mathcal{O}) &
}
$$
。这里 $(j, j^\sharp)$ 是与环站点
$(\mathcal{C}/V, \mathcal{O}_V)$ 的对象 $V/U$ 相伴的局部化态射。
\end{lemma}

\begin{proof}
只需检验 $j_V^\sharp = j^\sharp \circ j^{-1}(j_U^\sharp)$，
其余结论直接来自《站点》，引理 \ref{sites-lemma-relocalize} 和
公式 (\ref{sites-equation-relocalize})。略去该等式的验证。
\end{proof}

\begin{remark}
\label{sites-modules-remark-localize-shriek-equal}
在引理 \ref{sites-modules-lemma-extension-by-zero} 的情形中，图
$$
\xymatrix{
\textit{Mod}(\mathcal{O}_U) \ar[r]_{j_{U!}} \ar[d]_{forget} &
\textit{Mod}(\mathcal{O}_\mathcal{C}) \ar[d]^{forget} \\
\textit{Ab}(\mathcal{C}/U) \ar[r]^{j^{Ab}_{U!}} &
\textit{Ab}(\mathcal{C})
}
$$
交换。这直接来自该引理中函子 $j_{U!}$ 的显式描述。
\end{remark}

\begin{remark}
\label{sites-modules-remark-localize-presheaves}
局部化与模预层；见《站点》，注记 \ref{sites-remark-localize-presheaves}。
设 $\mathcal{C}$ 是范畴，$\mathcal{O}$ 是环预层，且 $U$ 是
$\mathcal{C}$ 的对象。严格地说，我们尚未对 $\mathcal{O}$-模预层定义
函子 $j_U^*$、$j_{U*}$ 和 $j_{U!}$。不过，可以把预层看成
$\mathcal{C}$ 上混沌拓扑的层（见《站点》，例 \ref{sites-example-indiscrete}）。
因此也得到函子
$$
j_U^* :
\textit{PMod}(\mathcal{O})
\longrightarrow
\textit{PMod}(\mathcal{O}_U)
$$
以及函子
$$
j_{U*}, j_{U!} :
\textit{PMod}(\mathcal{O}_U)
\longrightarrow
\textit{PMod}(\mathcal{O})
$$
，它们分别是 $j_U^*$ 的右伴随和左伴随。考察引理
\ref{sites-modules-lemma-extension-by-zero} 的证明可知，$j_{U!}\mathcal{G}$ 是预层
$$
V \longmapsto
\bigoplus\nolimits_{\varphi \in \Mor_\mathcal{C}(V, U)}
\mathcal{G}(V \xrightarrow{\varphi} U)
$$
。此外，函子 $j_{U!}$ 正合（对离散拓扑应用引理
\ref{sites-modules-lemma-extension-by-zero-exact}）。再者，若 $\mathcal{C}$ 实际上是站点，
且 $\mathcal{O}$ 实际上是环层，则图
$$
\xymatrix{
\textit{Mod}(\mathcal{O}_U) \ar[r]_{j_{U!}} \ar[d]_{forget} &
\textit{Mod}(\mathcal{O}) \\
\textit{PMod}(\mathcal{O}_U) \ar[r]^{j_{U!}} &
\textit{PMod}(\mathcal{O}) \ar[u]_{(\ )^\#}
}
$$
交换。
\end{remark}

\begin{lemma}
\label{sites-modules-lemma-restrict-back}
设 $\mathcal{C}$ 是站点，$U \in \Ob(\mathcal{C})$。假设对
$\mathcal{C}$ 中每个 $X$，至多有一个从 $X$ 到 $U$ 的态射。
设 $\mathcal{F}$ 是 $\mathcal{C}/U$ 上的阿贝尔层。则典范映射
$\mathcal{F} \to j_U^{-1}j_{U!}\mathcal{F}$ 和
$j_U^{-1}j_{U*}\mathcal{F} \to \mathcal{F}$ 都是同构。
\end{lemma}

\begin{proof}
这是引理 \ref{sites-modules-lemma-back-and-forth} 的特例，因为关于 $U$ 的假设等价于
局部化函子 $\mathcal{C}/U \to \mathcal{C}$ 是全忠实的。
\end{proof}









\section{环站点态射的局部化}
\label{sites-modules-section-localize-morphisms}

\noindent
本节是《站点》，第 \ref{sites-section-localize-morphisms} 节的类似版本。

\begin{lemma}
\label{sites-modules-lemma-localize-morphism-ringed-sites}
设
$(f, f^\sharp) :
(\mathcal{C}, \mathcal{O})
\longrightarrow
(\mathcal{D}, \mathcal{O}')$
是环站点态射，其中 $f$ 由连续函子
$u : \mathcal{D} \to \mathcal{C}$ 给出。设 $V$ 是 $\mathcal{D}$ 的对象，
并令 $U = u(V)$。则存在典范的环层映射 $(f')^\sharp$，使《站点》，
引理 \ref{sites-lemma-localize-morphism} 的图成为环拓扑斯的交换图
$$
\xymatrix{
(\Sh(\mathcal{C}/U), \mathcal{O}_U)
\ar[rr]_{(j_U, j_U^\sharp)} \ar[d]_{(f', (f')^\sharp)} & &
(\Sh(\mathcal{C}), \mathcal{O})
\ar[d]^{(f, f^\sharp)} \\
(\Sh(\mathcal{D}/V), \mathcal{O}'_V)
\ar[rr]^{(j_V, j_V^\sharp)} & &
(\Sh(\mathcal{D}), \mathcal{O}').
}
$$
。此外，在此情形中有 $f'_*j_U^{-1} = j_V^{-1}f_*$ 和
$f'_*j_U^* = j_V^*f_*$。
\end{lemma}

\begin{proof}
只需取 $(f')^\sharp$ 为
$$
(f')^{-1}\mathcal{O}'_V =
(f')^{-1}j_V^{-1}\mathcal{O}' =
j_U^{-1}f^{-1}\mathcal{O}' \xrightarrow{j_U^{-1}f^\sharp}
j_U^{-1}\mathcal{O} = \mathcal{O}_U
$$
，其余结论都来自《站点》，引理 \ref{sites-lemma-localize-morphism}。
（注意，若 $j$ 是局部化态射，则在模层上 $j^{-1} = j^*$，
因而第一个函子等式蕴含第二个。）
\end{proof}

\begin{lemma}
\label{sites-modules-lemma-relocalize-morphism-ringed-sites}
设
$(f, f^\sharp) :
(\mathcal{C}, \mathcal{O})
\longrightarrow
(\mathcal{D}, \mathcal{O}')$
是环站点态射，其中 $f$ 由连续函子
$u : \mathcal{D} \to \mathcal{C}$ 给出。设
$V \in \Ob(\mathcal{D})$、$U \in \Ob(\mathcal{C})$，且
$c : U \to u(V)$ 是 $\mathcal{C}$ 的态射。则存在环拓扑斯的交换图
$$
\xymatrix{
(\Sh(\mathcal{C}/U), \mathcal{O}_U)
\ar[rr]_{(j_U, j_U^\sharp)} \ar[d]_{(f_c, f_c^\sharp)} & &
(\Sh(\mathcal{C}), \mathcal{O}) \ar[d]^{(f, f^\sharp)} \\
(\Sh(\mathcal{D}/V), \mathcal{O}'_V)
\ar[rr]^{(j_V, j_V^\sharp)} & &
(\Sh(\mathcal{D}), \mathcal{O}').
}
$$
。态射 $(f_c, f_c^\sharp)$ 等于如下态射
$$
(f', (f')^\sharp) :
(\Sh(\mathcal{C}/u(V)), \mathcal{O}_{u(V)})
\longrightarrow
(\Sh(\mathcal{D}/V), \mathcal{O}'_V)
$$
（见引理 \ref{sites-modules-lemma-localize-morphism-ringed-sites}）与如下态射
$$
(j, j^\sharp) :
(\Sh(\mathcal{C}/U), \mathcal{O}_U)
\to
(\Sh(\mathcal{C}/u(V)), \mathcal{O}_{u(V)})
$$
（见引理 \ref{sites-modules-lemma-relocalize}）的复合。给定任意态射
$b : V' \to V$、$a : U' \to U$ 和 $c' : U' \to u(V')$，使得
$$
\xymatrix{
U' \ar[r]_-{c'} \ar[d]_a & u(V') \ar[d]^{u(b)} \\
U \ar[r]^-c & u(V)
}
$$
交换，则下列环拓扑斯之图
$$
\xymatrix{
(\Sh(\mathcal{C}/U'), \mathcal{O}_{U'})
\ar[rr]_{(j_{U'/U}, j_{U'/U}^\sharp)} \ar[d]_{(f_{c'}, f_{c'}^\sharp)} & &
(\Sh(\mathcal{C}/U), \mathcal{O}_U)
\ar[d]^{(f_c, f_c^\sharp)} \\
(\Sh(\mathcal{D}/V'), \mathcal{O}'_{V'})
\ar[rr]^{(j_{V'/V}, j_{V'/V}^\sharp)} & &
(\Sh(\mathcal{D}/V), \mathcal{O}'_{V'})
}
$$
交换。
\end{lemma}

\begin{proof}
在拓扑斯态射的层面，这就是《站点》，引理
\ref{sites-lemma-relocalize-morphism}。为检验这些图作为环拓扑斯态射之图交换，
使用引理 \ref{sites-modules-lemma-relocalize} 和
\ref{sites-modules-lemma-localize-morphism-ringed-sites}，完全仿照《站点》，引理
\ref{sites-lemma-relocalize-morphism} 的证明即可。
\end{proof}
















\section{环拓扑斯的局部化}
\label{sites-modules-section-localize-ringed-topoi}

\noindent
本节是在环拓扑斯的设定下，对《站点》，第
\ref{sites-section-localize-topoi} 节所作的类似论述。

\begin{lemma}
\label{sites-modules-lemma-localize-ringed-topos}
设 $(\Sh(\mathcal{C}), \mathcal{O})$ 是环拓扑斯，且
$\mathcal{F} \in \Sh(\mathcal{C})$ 是层。对 $\mathcal{C}$ 上的层
$\mathcal{H}$，以 $\mathcal{H}_\mathcal{F}$ 表示层
$\mathcal{H} \times \mathcal{F}$，并将它视为范畴
$\Sh(\mathcal{C})/\mathcal{F}$ 的对象。则
$(\Sh(\mathcal{C})/\mathcal{F}, \mathcal{O}_\mathcal{F})$
是环拓扑斯，且存在环拓扑斯的典范态射
$$
(j_\mathcal{F}, j_\mathcal{F}^\sharp) :
(\Sh(\mathcal{C})/\mathcal{F}, \mathcal{O}_\mathcal{F})
\longrightarrow
(\Sh(\mathcal{C}), \mathcal{O})
$$
，它是第 \ref{sites-modules-section-localize} 节意义下的局部化，并且
\begin{enumerate}
\item 函子 $j_\mathcal{F}^{-1}$ 是函子
$\mathcal{H} \mapsto \mathcal{H}_\mathcal{F}$,
\item 函子 $j_\mathcal{F}^*$ 是函子
$\mathcal{H} \mapsto \mathcal{H}_\mathcal{F}$,
\item 集合层上的函子 $j_{\mathcal{F}!}$ 是遗忘函子
$\mathcal{G}/\mathcal{F} \mapsto \mathcal{G}$，
\item 模层上的函子 $j_{\mathcal{F}!}$ 将 $\mathcal{O}_\mathcal{F}$-模
$\varphi : \mathcal{G} \to \mathcal{F}$ 送到如下预层之层化所得的
$\mathcal{O}$-模：
$$
V \longmapsto
\bigoplus\nolimits_{s \in \mathcal{F}(V)}
\{\sigma \in \mathcal{G}(V) \mid \varphi(\sigma) = s \}
$$
其中 $V \in \Ob(\mathcal{C})$。
\end{enumerate}
\end{lemma}

\begin{proof}
由《站点》，引理 \ref{sites-lemma-localize-topos} 可知，
$\Sh(\mathcal{C})/\mathcal{F}$ 是拓扑斯，且 (1) 和 (3) 成立。
特别地，这说明
$j_\mathcal{F}^{-1}\mathcal{O} = \mathcal{O}_\mathcal{F}$
，并说明 $\mathcal{O}_\mathcal{F}$ 是环层。因此可取映射
$j_\mathcal{F}^\sharp$ 为恒等映射，特别地，(2) 成立。
此外，《站点》，引理 \ref{sites-lemma-localize-topos} 的证明说明，
可以假设 $\mathcal{C}$ 是具有所有有限极限和次典范拓扑的站点，且对
$\mathcal{C}$ 的某个对象 $U$ 有 $\mathcal{F} = h_U$。
此时 (4) 来自引理 \ref{sites-modules-lemma-extension-by-zero} 中对
$j_{\mathcal{F}!}$ 的描述。另一种方法是直接证明 (4) 中描述的函子是
$j_\mathcal{F}^*$ 的左伴随。
\end{proof}

\begin{definition}
\label{sites-modules-definition-localize-ringed-topos}
设 $(\Sh(\mathcal{C}), \mathcal{O})$ 是环拓扑斯，且
$\mathcal{F} \in \Sh(\mathcal{C})$。
\begin{enumerate}
\item 环拓扑斯
$(\Sh(\mathcal{C})/\mathcal{F}, \mathcal{O}_\mathcal{F})$
称为{\it 环拓扑斯 $(\Sh(\mathcal{C}), \mathcal{O})$ 在
$\mathcal{F}$ 处的局部化}。
\item 引理 \ref{sites-modules-lemma-localize-ringed-topos} 的环拓扑斯态射
$(j_\mathcal{F}, j_\mathcal{F}^\sharp) :
(\Sh(\mathcal{C})/\mathcal{F}, \mathcal{O}_\mathcal{F})
\to
(\Sh(\mathcal{C}), \mathcal{O})$
称为{\it 局部化态射}。
\end{enumerate}
\end{definition}

\noindent
我们延续《站点》一章中确立的做法：每当有意义时，检验拓扑斯上的局部化构造
与站点上的局部化构造相容。

\begin{lemma}
\label{sites-modules-lemma-localize-compare}
设 $(\Sh(\mathcal{C}), \mathcal{O})$ 和
$\mathcal{F} \in \Sh(\mathcal{C})$ 如引理
\ref{sites-modules-lemma-localize-ringed-topos} 中所述。若对 $\mathcal{C}$ 的某个对象
$U$ 有 $\mathcal{F} = h_U^\#$，则经由《站点》，引理
\ref{sites-lemma-essential-image-j-shriek} 的等同
$\Sh(\mathcal{C}/U) = \Sh(\mathcal{C})/h_U^\#$，有
\begin{enumerate}
\item 典范地有 $\mathcal{O}_U = \mathcal{O}_\mathcal{F}$；
\item 在这些等同下，
$(j_\mathcal{F}, j_\mathcal{F}^\sharp) = (j_U, j_U^\sharp)$。
\end{enumerate}
\end{lemma}

\begin{proof}
关于底层拓扑斯的断言就是《站点》，引理
\ref{sites-lemma-localize-compare}。注意，$\mathcal{O}_U$ 是
$\mathcal{O}$ 的限制；由《站点》，引理
\ref{sites-lemma-compute-j-shriek-restrict}，在《站点》，引理
\ref{sites-lemma-essential-image-j-shriek} 的等价下，它对应于
$\mathcal{O} \times h_U^\#$。由 $\mathcal{O}_\mathcal{F}$ 的定义即得 (1)。
余下只须证明在此等同下 $j_\mathcal{F}^\sharp = j_U^\sharp$。验证从略。
\end{proof}

\noindent
局部化在以下两种意义下具有函子性：可以对一个局部化再次“局部化”（见引理
\ref{sites-modules-lemma-relocalize-ringed-topos}）；也可以给定环拓扑斯态射，在上方拓扑斯中
对下方拓扑斯上某层的逆像作局部化，从而得到环拓扑斯的交换图（见引理
\ref{sites-modules-lemma-localize-morphism-ringed-topoi}）。

\begin{lemma}
\label{sites-modules-lemma-relocalize-ringed-topos}
设 $(\Sh(\mathcal{C}), \mathcal{O})$ 是环拓扑斯。若
$s : \mathcal{G} \to \mathcal{F}$ 是 $\mathcal{C}$ 上的层态射，
则存在环拓扑斯态射的自然交换图
$$
\xymatrix{
(\Sh(\mathcal{C})/\mathcal{G}, \mathcal{O}_\mathcal{G})
\ar[rd]_{(j_\mathcal{G}, j_\mathcal{G}^\sharp)} \ar[rr]_{(j, j^\sharp)} & &
(\Sh(\mathcal{C})/\mathcal{F}, \mathcal{O}_\mathcal{F})
\ar[ld]^{(j_\mathcal{F}, j_\mathcal{F}^\sharp)} \\
& (\Sh(\mathcal{C}), \mathcal{O}) &
}
$$
，其中 $(j, j^\sharp)$ 是环拓扑斯
$(\Sh(\mathcal{C})/\mathcal{F}, \mathcal{O}_\mathcal{F})$
在对象 $\mathcal{G}/\mathcal{F}$ 处的局部化态射。
\end{lemma}

\begin{proof}
除等式
$j_\mathcal{G}^\sharp = j^\sharp \circ j^{-1}(j_\mathcal{F}^\sharp)$
外，所有断言都来自《站点》，引理 \ref{sites-lemma-relocalize-topos}。
该等式的验证从略。
\end{proof}

\begin{lemma}
\label{sites-modules-lemma-relocalize-compare}
设 $(\Sh(\mathcal{C}), \mathcal{O})$ 和
$s : \mathcal{G} \to \mathcal{F}$ 如引理
\ref{sites-modules-lemma-relocalize-ringed-topos} 中所述。若存在 $\mathcal{C}$ 的态射
$f : V \to U$，使得 $\mathcal{G} = h_V^\#$、
$\mathcal{F} = h_U^\#$，且 $s$ 由 $f$ 诱导，则引理
\ref{sites-modules-lemma-relocalize} 与引理
\ref{sites-modules-lemma-relocalize-ringed-topos} 的图经由等同
$(j_\mathcal{F}, j_\mathcal{F}^\sharp) = (j_U, j_U^\sharp)$
和
$(j_\mathcal{G}, j_\mathcal{G}^\sharp) = (j_V, j_V^\sharp)$
（见引理 \ref{sites-modules-lemma-localize-compare}）彼此一致。
\end{lemma}

\begin{proof}
除两个映射 $j^\sharp$ 一致这一断言外，所有断言都来自《站点》，引理
\ref{sites-lemma-relocalize-compare}。在两种情形中，映射 $j^\sharp$
都只是恒等映射，故该断言成立。略去若干细节。
\end{proof}










\section{环拓扑斯态射的局部化}
\label{sites-modules-section-localize-morphisms-ringed-topoi}

\noindent
本节是《站点》，第 \ref{sites-section-localize-morphisms-topoi} 节的类似版本。

\begin{lemma}
\label{sites-modules-lemma-localize-morphism-ringed-topoi}
设
$$
f :
(\Sh(\mathcal{C}), \mathcal{O})
\longrightarrow
(\Sh(\mathcal{D}), \mathcal{O}')
$$
是环拓扑斯态射。设 $\mathcal{G}$ 是 $\mathcal{D}$ 上的层，并令
$\mathcal{F} = f^{-1}\mathcal{G}$。则存在环拓扑斯的交换图
$$
\xymatrix{
(\Sh(\mathcal{C})/\mathcal{F}, \mathcal{O}_\mathcal{F})
\ar[rr]_{(j_\mathcal{F}, j_\mathcal{F}^\sharp)}
\ar[d]_{(f', (f')^\sharp)} & &
(\Sh(\mathcal{C}), \mathcal{O}) \ar[d]^{(f, f^\sharp)} \\
(\Sh(\mathcal{D})/\mathcal{G}, \mathcal{O}'_\mathcal{G})
\ar[rr]^{(j_\mathcal{G}, j_\mathcal{G}^\sharp)} & &
(\Sh(\mathcal{D}), \mathcal{O}')
}
$$
。有 $f'_*j_\mathcal{F}^{-1} = j_\mathcal{G}^{-1}f_*$ 和
$f'_*j_\mathcal{F}^* = j_\mathcal{G}^*f_*$。此外，态射 $f'$ 由规则
$$
(f')^{-1}(\mathcal{H} \xrightarrow{\varphi} \mathcal{G})
=
(f^{-1}\mathcal{H} \xrightarrow{f^{-1}\varphi} \mathcal{F}).
$$
\end{lemma}

\begin{proof}
由《站点》，引理 \ref{sites-lemma-localize-morphism-topoi}，得到底层拓扑斯之图、
等式 $f'_*j_\mathcal{F}^{-1} = j_\mathcal{G}^{-1}f_*$，以及
对 $(f')^{-1}$ 的描述。为定义 $(f')^\sharp$，使用映射
$$
(f')^\sharp :
\mathcal{O}'_\mathcal{G} =
j_\mathcal{G}^{-1} \mathcal{O}'
\xrightarrow{j_\mathcal{G}^{-1}f^\sharp}
j_\mathcal{G}^{-1} f_*\mathcal{O} =
f'_* j_\mathcal{F}^{-1}\mathcal{O} =
f'_* \mathcal{O}_\mathcal{F}
$$
，或等价地使用映射
$$
(f')^\sharp :
(f')^{-1}\mathcal{O}'_\mathcal{G} =
(f')^{-1}j_\mathcal{G}^{-1} \mathcal{O}' =
j_\mathcal{F}^{-1}f^{-1}\mathcal{O}'
\xrightarrow{j_\mathcal{F}^{-1}f^\sharp}
j_\mathcal{F}^{-1} \mathcal{O} =
\mathcal{O}_\mathcal{F}.
$$
。略去这两个映射确实互为伴随映射的验证。$(f')^\sharp$ 的第二种构造说明，
该图在环拓扑斯的 $2$-范畴中交换（因为映射
$j_\mathcal{F}^\sharp$ 和 $j_\mathcal{G}^\sharp$ 都是恒等映射）。
最后，等式 $f'_*j_\mathcal{F}^* = j_\mathcal{G}^*f_*$ 来自等式
$f'_*j_\mathcal{F}^{-1} = j_\mathcal{G}^{-1}f_*$
以及模层与集合层的拉回一致这一事实；见引理
\ref{sites-modules-lemma-localize-ringed-topos}。
\end{proof}

\begin{lemma}
\label{sites-modules-lemma-localize-morphism-compare}
设
$$
f :
(\Sh(\mathcal{C}), \mathcal{O})
\longrightarrow
(\Sh(\mathcal{D}), \mathcal{O}')
$$
是环拓扑斯态射。设 $\mathcal{G}$ 是 $\mathcal{D}$ 上的层，并令
$\mathcal{F} = f^{-1}\mathcal{G}$。若 $f$ 由连续函子
$u : \mathcal{D} \to \mathcal{C}$ 给出，且 $\mathcal{G} = h_V^\#$，
则引理 \ref{sites-modules-lemma-localize-morphism-ringed-sites} 与引理
\ref{sites-modules-lemma-localize-morphism-ringed-topoi} 的交换图经由引理
\ref{sites-modules-lemma-localize-compare} 的等同彼此一致。
\end{lemma}

\begin{proof}
在拓扑斯态射的层面，这就是《站点》，引理
\ref{sites-lemma-localize-morphism-compare}。这在环拓扑斯态射的层面也成立，
因为引理 \ref{sites-modules-lemma-localize-morphism-ringed-sites} 与引理
\ref{sites-modules-lemma-localize-morphism-ringed-topoi} 的证明中定义
$(f')^\sharp$ 的公式一致。
\end{proof}

\begin{lemma}
\label{sites-modules-lemma-relocalize-morphism-ringed-topoi}
设
$(f, f^\sharp) :
(\Sh(\mathcal{C}), \mathcal{O})
\to
(\Sh(\mathcal{D}), \mathcal{O}')$
是环拓扑斯态射。设 $\mathcal{G}$ 是 $\mathcal{D}$ 上的层，
$\mathcal{F}$ 是 $\mathcal{C}$ 上的层，且
$s : \mathcal{F} \to f^{-1}\mathcal{G}$ 是层态射。则存在环拓扑斯的交换图
$$
\xymatrix{
(\Sh(\mathcal{C})/\mathcal{F}, \mathcal{O}_\mathcal{F})
\ar[rr]_{(j_\mathcal{F}, j_\mathcal{F}^\sharp)}
\ar[d]_{(f_c, f_c^\sharp)} & &
(\Sh(\mathcal{C}), \mathcal{O})
\ar[d]^{(f, f^\sharp)} \\
(\Sh(\mathcal{D})/\mathcal{G}, \mathcal{O}'_\mathcal{G})
\ar[rr]^{(j_\mathcal{G}, j_\mathcal{G}^\sharp)} & &
(\Sh(\mathcal{D}), \mathcal{O}').
}
$$
。态射 $(f_s, f_s^\sharp)$ 等于如下态射
$$
(f', (f')^\sharp) :
(\Sh(\mathcal{C})/f^{-1}\mathcal{G}, \mathcal{O}_{f^{-1}\mathcal{G}})
\longrightarrow
(\Sh(\mathcal{D})/{\mathcal{G}}, \mathcal{O}'_\mathcal{G})
$$
（见引理 \ref{sites-modules-lemma-localize-morphism-ringed-topoi}）与如下态射
$$
(j, j^\sharp) :
(\Sh(\mathcal{C})/\mathcal{F}, \mathcal{O}_\mathcal{F})
\to
(\Sh(\mathcal{C})/f^{-1}\mathcal{G}, \mathcal{O}_{f^{-1}\mathcal{G}})
$$
（见引理 \ref{sites-modules-lemma-relocalize-ringed-topos}）的复合。给定任意态射
$b : \mathcal{G}' \to \mathcal{G}$、$a : \mathcal{F}' \to \mathcal{F}$ 和
$s' : \mathcal{F}' \to f^{-1}\mathcal{G}'$，使得
$$
\xymatrix{
\mathcal{F}' \ar[r]_-{s'} \ar[d]_a &
f^{-1}\mathcal{G}' \ar[d]^{f^{-1}b} \\
\mathcal{F} \ar[r]^-s &
f^{-1}\mathcal{G}
}
$$
交换，则下列环拓扑斯之图
$$
\xymatrix{
(\Sh(\mathcal{C})/\mathcal{F}', \mathcal{O}_{\mathcal{F}'})
\ar[rr]_{(j_{\mathcal{F}'/\mathcal{F}}, j_{\mathcal{F}'/\mathcal{F}}^\sharp)}
\ar[d]_{(f_{s'}, f_{s'}^\sharp)} & &
(\Sh(\mathcal{C})/\mathcal{F}, \mathcal{O}_\mathcal{F})
\ar[d]^{(f_s, f_s^\sharp)} \\
(\Sh(\mathcal{D})/\mathcal{G}', \mathcal{O}'_{\mathcal{G}'})
\ar[rr]^{(j_{\mathcal{G}'/\mathcal{G}}, j_{\mathcal{G}'/\mathcal{G}}^\sharp)}
& &
(\Sh(\mathcal{D})/\mathcal{G}, \mathcal{O}'_{\mathcal{G}'})
}
$$
交换。
\end{lemma}

\begin{proof}
在拓扑斯态射的层面，这就是《站点》，引理
\ref{sites-lemma-relocalize-morphism-topoi}。为检验这些图作为环拓扑斯态射之图
交换，使用引理 \ref{sites-modules-lemma-relocalize-ringed-topos} 和
\ref{sites-modules-lemma-localize-morphism-ringed-topoi} 的交换图即可。
\end{proof}

\begin{lemma}
\label{sites-modules-lemma-relocalize-morphism-compare}
设
$(f, f^\sharp) :
(\Sh(\mathcal{C}), \mathcal{O})
\to
(\Sh(\mathcal{D}), \mathcal{O}')$,
$s : \mathcal{F} \to f^{-1}\mathcal{G}$ 如引理
\ref{sites-modules-lemma-relocalize-morphism-ringed-topoi} 中所述。若 $f$ 由连续函子
$u : \mathcal{D} \to \mathcal{C}$
给出，且 $\mathcal{G} = h_V^\#$、$\mathcal{F} = h_U^\#$，而 $s$
来自态射 $c : U \to u(V)$，则引理
\ref{sites-modules-lemma-relocalize-morphism-ringed-sites} 与引理
\ref{sites-modules-lemma-relocalize-morphism-ringed-topoi} 的交换图经由引理
\ref{sites-modules-lemma-localize-compare} 的等同彼此一致。
\end{lemma}

\begin{proof}
这是引理 \ref{sites-modules-lemma-relocalize-compare} 和
\ref{sites-modules-lemma-localize-morphism-compare} 的形式推论。
\end{proof}

















\section{模的局部类型}
\label{sites-modules-section-local}

\noindent
按照第 \ref{sites-modules-section-intrinsic} 节说明的一般策略，我们先定义环站点上模层的
局部类型，随后立即证明这些类型是内蕴的，因而对环拓扑斯上的模层也有意义。

\begin{definition}
\label{sites-modules-definition-site-local}
设 $(\mathcal{C}, \mathcal{O})$ 是环站点，且 $\mathcal{F}$ 是
$\mathcal{O}$-模层。我们将自由使用定义 \ref{sites-modules-definition-global} 中定义的概念。
\begin{enumerate}
\item 称 $\mathcal{F}$ 是{\it 局部自由的}，如果对 $\mathcal{C}$ 的每个对象
$U$，存在 $\mathcal{C}$ 的覆盖 $\{U_i \to U\}_{i \in I}$，使每个限制
$\mathcal{F}|_{\mathcal{C}/U_i}$ 都是自由 $\mathcal{O}_{U_i}$-模。
\item 称 $\mathcal{F}$ 是{\it 有限局部自由的}，如果对 $\mathcal{C}$ 的每个对象
$U$，存在 $\mathcal{C}$ 的覆盖 $\{U_i \to U\}_{i \in I}$，使每个限制
$\mathcal{F}|_{\mathcal{C}/U_i}$ 都是有限自由 $\mathcal{O}_{U_i}$-模。
\item 称 $\mathcal{F}$ 是{\it 局部由截面生成的}，如果对 $\mathcal{C}$ 的每个
对象 $U$，存在覆盖 $\{U_i \to U\}_{i \in I}$，使每个限制
$\mathcal{F}|_{\mathcal{C}/U_i}$ 都是由整体截面生成的
$\mathcal{O}_{U_i}$-模。
\item 给定 $r \geq 0$，称 $\mathcal{F}$ 是{\it 局部由 $r$ 个截面生成的}，
如果对 $\mathcal{C}$ 的每个对象 $U$，存在覆盖
$\{U_i \to U\}_{i \in I}$，使每个限制
$\mathcal{F}|_{\mathcal{C}/U_i}$ 都是由 $r$ 个整体截面生成的
$\mathcal{O}_{U_i}$-模。
\item 称 $\mathcal{F}$ 是{\it 有限型的}，如果对 $\mathcal{C}$ 的每个对象
$U$，存在覆盖 $\{U_i \to U\}_{i \in I}$，使每个限制
$\mathcal{F}|_{\mathcal{C}/U_i}$ 都是由有限多个整体截面生成的
$\mathcal{O}_{U_i}$-模。
\item 称 $\mathcal{F}$ 是{\it 拟凝聚的}，如果对 $\mathcal{C}$ 的每个对象
$U$，存在覆盖 $\{U_i \to U\}_{i \in I}$，使每个限制
$\mathcal{F}|_{\mathcal{C}/U_i}$ 都是具有整体表示的
$\mathcal{O}_{U_i}$-模。
\item 称 $\mathcal{F}$ 是{\it 有限表示的}，如果对 $\mathcal{C}$ 的每个对象
$U$，存在覆盖 $\{U_i \to U\}_{i \in I}$，使每个限制
$\mathcal{F}|_{\mathcal{C}/U_i}$ 都是具有有限整体表示的
$\mathcal{O}_{U_i}$-模。
\item 称 $\mathcal{F}$ 是{\it 凝聚的}，当且仅当 $\mathcal{F}$ 是有限型的，
并且对 $\mathcal{C}$ 的每个对象 $U$ 及任意
$s_1, \ldots, s_n \in \mathcal{F}(U)$，映射
$\bigoplus_{i = 1, \ldots, n} \mathcal{O}_U \to \mathcal{F}|_U$
的核在 $(\mathcal{C}/U, \mathcal{O}_U)$ 上是有限型的。
\end{enumerate}
\end{definition}

\begin{lemma}
\label{sites-modules-lemma-special-locally-free}
定义 \ref{sites-modules-definition-site-local} 的性质 (1)--(8) 中任意一个都是内蕴的
（见第 \ref{sites-modules-section-intrinsic} 节的讨论）。
\end{lemma}

\begin{proof}
设 $\mathcal{C}$、$\mathcal{D}$ 是站点，且
$u : \mathcal{C} \to \mathcal{D}$ 是特殊余连续函子。设
$\mathcal{O}$ 是 $\mathcal{C}$ 上的环层，$\mathcal{F}$ 是
$\mathcal{C}$ 上的 $\mathcal{O}$-模层。设
$g : \Sh(\mathcal{C}) \to \Sh(\mathcal{D})$ 是与 $u$ 相伴的拓扑斯等价。
令 $\mathcal{O}' = g_*\mathcal{O}$，并令
$g^\sharp : \mathcal{O}' \to g_*\mathcal{O}$ 为恒等映射。最后，令
$\mathcal{F}' = g_*\mathcal{F}$。设 $\mathcal{P}_l$ 是定义
\ref{sites-modules-definition-site-local} 所列性质 (1)--(7) 之一（凝聚情形将在证明末尾讨论）。
以 $\mathcal{P}_g$ 表示定义 \ref{sites-modules-definition-global} 中相应的性质。
我们已经看到 $\mathcal{P}_g$ 是内蕴的。须证明
$\mathcal{P}_l(\mathcal{C}, \mathcal{O}, \mathcal{F})$
成立，当且仅当
$\mathcal{P}_l(\mathcal{D}, \mathcal{O}', \mathcal{F}')$
成立。

\medskip\noindent
假设 $\mathcal{F}$ 具有性质 $\mathcal{P}_l$。设 $V$ 是 $\mathcal{D}$ 的对象。
特殊余连续函子的性质之一是：站点 $\mathcal{D}$ 中存在覆盖
$\{u(U_i) \to V\}_{i \in I}$。由假设，对每个 $i$，在 $\mathcal{C}$ 中
存在覆盖 $\{U_{ij} \to U_i\}_{j \in J_i}$，使每个限制
$\mathcal{F}|_{U_{ij}}$ 都具有性质 $\mathcal{P}_g$。由《站点》，引理
\ref{sites-lemma-localize-special-cocontinuous}，有环拓扑斯的交换图
$$
\xymatrix{
(\Sh(\mathcal{C}/U_{ij}), \mathcal{O}_{U_{ij}}) \ar[r] \ar[d] &
(\Sh(\mathcal{C}), \mathcal{O}) \ar[d] \\
(\Sh(\mathcal{D}/u(U_{ij})), \mathcal{O}'_{u(U_{ij})}) \ar[r] &
(\Sh(\mathcal{D}), \mathcal{O}')
}
$$
，其中竖直箭头都是等价。因此 $\mathcal{F}'|_{u(U_{ij})}$ 也具有性质
$\mathcal{P}_g$。此外，$\{u(U_{ij}) \to V\}_{i \in I, j \in J_i}$
是站点 $\mathcal{D}$ 的覆盖，故 $\mathcal{F}'$ 具有性质
$\mathcal{P}_l$。

\medskip\noindent
假设 $\mathcal{F}'$ 具有性质 $\mathcal{P}_l$。设 $U$ 是
$\mathcal{C}$ 的对象。由假设，存在覆盖
$\{V_i \to u(U)\}_{i \in I}$，使 $\mathcal{F}'|_{V_i}$ 具有性质
$\mathcal{P}_g$。由于 $u$ 余连续，可以用族
$\{u(U_j) \to u(U)\}_{j \in J}$ 加细此覆盖，其中
$\{U_j \to U\}_{j \in J}$ 是 $\mathcal{C}$ 中的覆盖。设此加细由
$\alpha : J \to I$ 和 $u(U_j) \to V_{\alpha(j)}$ 给出。
限制具有传递性，即
$(\mathcal{F}'|_{V_{\alpha(j)}})|_{u(U_j)} = \mathcal{F}'|_{u(U_j)}$.
因此由引理 \ref{sites-modules-lemma-global-pullback} 可知，
$\mathcal{F}'|_{u(U_j)}$ 具有性质 $\mathcal{P}_g$。于是图
$$
\xymatrix{
(\Sh(\mathcal{C}/U_j), \mathcal{O}_{U_j}) \ar[r] \ar[d] &
(\Sh(\mathcal{C}), \mathcal{O}) \ar[d] \\
(\Sh(\mathcal{D}/u(U_j)), \mathcal{O}'_{u(U_j)})
\ar[r] &
(\Sh(\mathcal{D}), \mathcal{O}')
}
$$
（其中竖直箭头都是等价）说明 $\mathcal{F}|_{U_j}$ 也具有性质
$\mathcal{P}_g$。故 $\mathcal{F}$ 具有性质 $\mathcal{P}_l$，正合所需。

\medskip\noindent
最后证明 $\mathcal{P}_l = coherent$ 的情形\footnote{这里的操作略显繁琐，
我们建议读者跳过证明的这一部分。}。假设 $\mathcal{F}$ 凝聚。
于是 $\mathcal{F}$ 是有限型的，由证明的第一部分，$\mathcal{F}'$ 也是有限型的。
设 $V$ 是 $\mathcal{D}$ 的对象，且
$s_1, \ldots, s_n \in \mathcal{F}'(V)$。须证明映射
$\bigoplus_{j = 1, \ldots, n} \mathcal{O}_V \to \mathcal{F}'|_V$ 的核
$\mathcal{K}'$
在 $\mathcal{D}/V$ 上是有限型的。这意味着须证明：对任意 $V'/V$，
存在覆盖 $\{V'_i \to V'\}$，使 $\mathcal{F}'|_{V'_i}$ 由有限多个截面生成。
以 $V'$ 替换 $V$（并把截面 $s_j$ 限制到 $V'$），可约化到
$V' = V$ 的情形。由于 $u$ 是特殊余连续函子，站点 $\mathcal{D}$ 中
存在覆盖 $\{u(U_i) \to V\}_{i \in I}$。利用拓扑斯同构
$\Sh(\mathcal{C}/U_i) = \Sh(\mathcal{D}/u(U_i))$
可知，$\mathcal{K}'|_{u(U_i)}$ 对应于映射
$\bigoplus_{j = 1, \ldots, n} \mathcal{O}_{U_i} \to \mathcal{F}|_{U_i}$.
的核 $\mathcal{K}_i$。由于 $\mathcal{F}$ 凝聚，$\mathcal{K}_i$ 是有限型的。
再次由证明的第一部分可知，$\mathcal{K}|_{u(U_i)}$ 是有限型的。因此存在覆盖
$\{V_{il} \to u(U_i)\}$，使 $\mathcal{K}|_{V_{il}}$ 由有限多个整体截面生成。
由于 $\{V_{il} \to V\}$ 是 $\mathcal{D}$ 的覆盖，得到
$\mathcal{K}$ 是有限型的，正合所需。

\medskip\noindent
假设 $\mathcal{F}'$ 凝聚。于是 $\mathcal{F}'$ 是有限型的，由证明的第一部分，
$\mathcal{F}$ 也是有限型的。设 $U$ 是 $\mathcal{C}$ 的对象，且
$s_1, \ldots, s_n \in \mathcal{F}(U)$。须证明映射
$\bigoplus_{j = 1, \ldots, n} \mathcal{O}_U \to \mathcal{F}|_U$ 的核
$\mathcal{K}$
在 $\mathcal{C}/U$ 上是有限型的。利用拓扑斯同构
$\Sh(\mathcal{C}/U) = \Sh(\mathcal{D}/u(U))$
可知，$\mathcal{K}|_U$ 对应于映射
$\bigoplus_{j = 1, \ldots, n} \mathcal{O}_{u(U)} \to \mathcal{F}'|_{u(U)}$
的核 $\mathcal{K}'$。
由于 $\mathcal{F}'$ 凝聚，$\mathcal{K}'$ 是有限型的。再次由证明的第一部分，
得到 $\mathcal{K}$ 是有限型的。
\end{proof}

\noindent
因此，从现在起可以不指定站点，直接使用定义 \ref{sites-modules-definition-site-local}
中定义的 $\mathcal{O}$-模性质。

\begin{lemma}
\label{sites-modules-lemma-local-final-object}
设 $(\Sh(\mathcal{C}), \mathcal{O})$ 是环拓扑斯，且 $\mathcal{F}$ 是
$\mathcal{O}$-模。假设站点 $\mathcal{C}$ 有终对象 $X$。则
\begin{enumerate}
\item 下列条件等价：
\begin{enumerate}
\item $\mathcal{F}$ 局部自由；
\item $\mathcal{C}$ 中存在覆盖 $\{X_i \to X\}$，使每个限制
$\mathcal{F}|_{\mathcal{C}/X_i}$ 都是局部自由 $\mathcal{O}_{X_i}$-模；
\item $\mathcal{C}$ 中存在覆盖 $\{X_i \to X\}$，使每个限制
$\mathcal{F}|_{\mathcal{C}/X_i}$ 都是自由 $\mathcal{O}_{X_i}$-模。
\end{enumerate}
\item 下列条件等价：
\begin{enumerate}
\item $\mathcal{F}$ 有限局部自由；
\item $\mathcal{C}$ 中存在覆盖 $\{X_i \to X\}$，使每个限制
$\mathcal{F}|_{\mathcal{C}/X_i}$ 都是有限局部自由
$\mathcal{O}_{X_i}$-模；
\item $\mathcal{C}$ 中存在覆盖 $\{X_i \to X\}$，使每个限制
$\mathcal{F}|_{\mathcal{C}/X_i}$ 都是有限自由 $\mathcal{O}_{X_i}$-模。
\end{enumerate}
\item 下列条件等价：
\begin{enumerate}
\item $\mathcal{F}$ 局部由截面生成；
\item $\mathcal{C}$ 中存在覆盖 $\{X_i \to X\}$，使每个限制
$\mathcal{F}|_{\mathcal{C}/X_i}$ 都是局部由截面生成的
$\mathcal{O}_{X_i}$-模；
\item $\mathcal{C}$ 中存在覆盖 $\{X_i \to X\}$，使每个限制
$\mathcal{F}|_{\mathcal{C}/X_i}$ 都是由整体截面生成的
$\mathcal{O}_{X_i}$-模。
\end{enumerate}
\item 给定 $r \geq 0$，下列条件等价：
\begin{enumerate}
\item $\mathcal{F}$ 局部由 $r$ 个截面生成；
\item $\mathcal{C}$ 中存在覆盖 $\{X_i \to X\}$，使每个限制
$\mathcal{F}|_{\mathcal{C}/X_i}$ 都是局部由 $r$ 个截面生成的
$\mathcal{O}_{X_i}$-模；
\item $\mathcal{C}$ 中存在覆盖 $\{X_i \to X\}$，使每个限制
$\mathcal{F}|_{\mathcal{C}/X_i}$ 都是由 $r$ 个整体截面生成的
$\mathcal{O}_{X_i}$-模。
\end{enumerate}
\item 下列条件等价：
\begin{enumerate}
\item $\mathcal{F}$ 是有限型的；
\item $\mathcal{C}$ 中存在覆盖 $\{X_i \to X\}$，使每个限制
$\mathcal{F}|_{\mathcal{C}/X_i}$ 都是有限型 $\mathcal{O}_{X_i}$-模；
\item $\mathcal{C}$ 中存在覆盖 $\{X_i \to X\}$，使每个限制
$\mathcal{F}|_{\mathcal{C}/X_i}$ 都是由有限多个整体截面生成的
$\mathcal{O}_{X_i}$-模。
\end{enumerate}
\item 下列条件等价：
\begin{enumerate}
\item $\mathcal{F}$ 拟凝聚；
\item $\mathcal{C}$ 中存在覆盖 $\{X_i \to X\}$，使每个限制
$\mathcal{F}|_{\mathcal{C}/X_i}$ 都是拟凝聚 $\mathcal{O}_{X_i}$-模；
\item $\mathcal{C}$ 中存在覆盖 $\{X_i \to X\}$，使每个限制
$\mathcal{F}|_{\mathcal{C}/X_i}$ 都是具有整体表示的
$\mathcal{O}_{X_i}$-模。
\end{enumerate}
\item 下列条件等价：
\begin{enumerate}
\item $\mathcal{F}$ 是有限表示的；
\item $\mathcal{C}$ 中存在覆盖 $\{X_i \to X\}$，使每个限制
$\mathcal{F}|_{\mathcal{C}/X_i}$ 都是有限表示 $\mathcal{O}_{X_i}$-模；
\item $\mathcal{C}$ 中存在覆盖 $\{X_i \to X\}$，使每个限制
$\mathcal{F}|_{\mathcal{C}/X_i}$ 都是具有有限整体表示的
$\mathcal{O}_{X_i}$-模。
\end{enumerate}
\item 下列条件等价：
\begin{enumerate}
\item $\mathcal{F}$ 凝聚；
\item $\mathcal{C}$ 中存在覆盖 $\{X_i \to X\}$，使每个限制
$\mathcal{F}|_{\mathcal{C}/X_i}$ 都是凝聚 $\mathcal{O}_{X_i}$-模。
\end{enumerate}
\end{enumerate}
\end{lemma}

\begin{proof}
在每种情形中都有 (a) $\Rightarrow (b)$。对情形 (1)--(6)，由站点的公理 (2)
（见《站点》，定义 \ref{sites-definition-site}）和模的局部类型定义，
条件 (b) 蕴含条件 (c)。设 $\{X_i \to X\}$ 是覆盖。则对
$\mathcal{C}$ 的每个对象 $U$，得到诱导覆盖
$\{X_i \times_X U \to U\}$。此外，由引理 \ref{sites-modules-lemma-global-pullback}，
(c) 中 $\mathcal{F}|_{\mathcal{C}/X_i}$ 的整体性质蕴含
$\mathcal{F}|_{\mathcal{C}/X_i \times_X U}$ 的相应整体性质，
因而按定义该层具有性质 (a)。略去情形 (7) 中 (b) $\Rightarrow$ (a) 的证明。
\end{proof}

\begin{lemma}
\label{sites-modules-lemma-local-pullback}
设
$(f, f^\sharp) :
(\Sh(\mathcal{C}), \mathcal{O}_\mathcal{C})
\to
(\Sh(\mathcal{D}), \mathcal{O}_\mathcal{D})$
是环拓扑斯态射，且 $\mathcal{F}$ 是 $\mathcal{O}_\mathcal{D}$-模。
\begin{enumerate}
\item 若 $\mathcal{F}$ 局部自由，则 $f^*\mathcal{F}$ 局部自由。
\item 若 $\mathcal{F}$ 有限局部自由，则 $f^*\mathcal{F}$ 有限局部自由。
\item 若 $\mathcal{F}$ 局部由截面生成，则 $f^*\mathcal{F}$ 局部由截面生成。
\item 若 $\mathcal{F}$ 局部由 $r$ 个截面生成，则
$f^*\mathcal{F}$ 局部由 $r$ 个截面生成。
\item 若 $\mathcal{F}$ 是有限型的，则 $f^*\mathcal{F}$ 是有限型的。
\item 若 $\mathcal{F}$ 拟凝聚，则 $f^*\mathcal{F}$ 拟凝聚。
\item 若 $\mathcal{F}$ 是有限表示的，则 $f^*\mathcal{F}$ 是有限表示的。
\end{enumerate}
\end{lemma}

\begin{proof}
按照第 \ref{sites-modules-section-intrinsic} 节的讨论，只须对环站点态射
$(f, f^\sharp) :
(\mathcal{C}, \mathcal{O}_\mathcal{C})
\to
(\mathcal{D}, \mathcal{O}_\mathcal{D})$
检验拉回保持这些性质，其中 $f$ 由具有所有有限极限的站点之间的左正合连续函子
$u : \mathcal{D} \to \mathcal{C}$ 给出。设 $\mathcal{G}$ 是
$\mathcal{O}_\mathcal{D}$-模层，并具有定义
\ref{sites-modules-definition-site-local} 的性质 (1)--(6) 之一。我们知道
$\mathcal{D}$ 有终对象 $Y$，且 $X = u(Y)$ 是 $\mathcal{C}$ 的终对象。
由假设，存在覆盖 $\{Y_i \to Y\}$，使
$\mathcal{G}|_{\mathcal{D}/Y_i}$ 具有相应的整体性质。令
$X_i = u(Y_i)$，于是 $\{X_i \to X\}$ 是 $\mathcal{C}$ 中的覆盖。
由《站点》，引理 \ref{sites-lemma-localize-morphism-strong}，得到环站点态射的交换图
$$
\xymatrix{
(\mathcal{C}/X_i, \mathcal{O}_\mathcal{C}|_{X_i}) \ar[r] \ar[d] &
(\mathcal{C}, \mathcal{O}_\mathcal{C}) \ar[d] \\
(\mathcal{D}/Y_i, \mathcal{O}_\mathcal{D}|_{Y_i}) \ar[r] &
(\mathcal{D}, \mathcal{O}_\mathcal{D})
}
$$
。因此由引理 \ref{sites-modules-lemma-global-pullback}，
$f^*\mathcal{G}|_{X_i}$ 具有相应的整体性质。故由引理
\ref{sites-modules-lemma-local-final-object}，得到 $\mathcal{G}$ 具有起初所考虑的局部性质。
\end{proof}







\section{模的局部类型的基本结果}
\label{sites-modules-section-basics}

\noindent
下面给出与上述定义有关的基本引理。

\begin{lemma}
\label{sites-modules-lemma-cokernel-finite-finite-presentation}
设 $(\mathcal{C},\mathcal{O})$ 是环站点，$\mathcal{F}$ 是有限表示
$\mathcal{O}$-模，且
$\varphi : \mathcal{G} \to \mathcal{F}$ 是 $\mathcal{O}$-模态射。
若 $\mathcal{G}$ 是有限型的，则 $\Coker(\varphi)$ 是有限表示的。
\end{lemma}

\begin{proof}
从略。提示：见《模层》，引理
\ref{modules-lemma-cokernel-finite-finite-presentation}。
\end{proof}

\begin{lemma}
\label{sites-modules-lemma-kernel-surjection-finite-onto-finite-presentation}
设 $(\mathcal{C}, \mathcal{O})$ 是环站点。设
$\theta : \mathcal{G} \to \mathcal{F}$ 是 $\mathcal{O}$-模的满射，
其中 $\mathcal{F}$ 是有限表示的，$\mathcal{G}$ 是有限型的。
则 $\Ker(\theta)$ 是有限型的。
\end{lemma}

\begin{proof}
从略。提示：见《模层》，引理
\ref{modules-lemma-kernel-surjection-finite-free-onto-finite-presentation}。
\end{proof}

\begin{lemma}
\label{sites-modules-lemma-chaotic-quasi-coherent}
设 $\mathcal{C}$ 是范畴，并赋予混沌拓扑而视为站点；见《站点》，例
\ref{sites-example-indiscrete}。设 $\mathcal{O}$ 是 $\mathcal{C}$ 上的环层，
$\mathcal{F}$ 是 $\mathcal{O}$-模层。则 $\mathcal{F}$ 拟凝聚，当且仅当
对 $\mathcal{C}$ 中所有 $U \to V$，典范映射
$$
\mathcal{F}(V) \otimes_{\mathcal{O}(V)} \mathcal{O}(U)
\longrightarrow
\mathcal{F}(U)
$$
是同构。
\end{lemma}

\begin{proof}
假设 $\mathcal{F}$ 拟凝聚，并设 $U \to V$ 是 $\mathcal{C}$ 的态射。
由于 $V$ 的每个覆盖都由一个同构给出，定义
\ref{sites-modules-definition-site-local} 说明存在表示
$$
\bigoplus\nolimits_{j \in J} \mathcal{O}_V \longrightarrow
\bigoplus\nolimits_{i \in I} \mathcal{O}_V \longrightarrow
\mathcal{F}|_{\mathcal{C}/V} \longrightarrow 0
$$
。由于 $\mathcal{C}$ 上的拓扑是混沌拓扑，在 $\mathcal{C}$ 的任意对象上
取截面是正合的。因此得到 $\mathcal{O}(V)$-模 $\mathcal{F}(V)$ 的表示
$$
\bigoplus\nolimits_{j \in J} \mathcal{O}(V) \longrightarrow
\bigoplus\nolimits_{i \in I} \mathcal{O}(V) \longrightarrow
\mathcal{F}(V) \longrightarrow 0
$$
，对 $\mathcal{F}(U)$ 亦然。这立即说明引理陈述中的映射是同构。

\medskip\noindent
反之，假设对 $\mathcal{C}$ 中每个态射 $U \to V$，引理陈述中的映射都是同构。
固定 $V$，并选取 $\mathcal{O}(V)$-模 $\mathcal{F}(V)$ 的表示
$$
\bigoplus\nolimits_{j \in J} \mathcal{O}(V) \longrightarrow
\bigoplus\nolimits_{i \in I} \mathcal{O}(V) \longrightarrow
\mathcal{F}(V) \longrightarrow 0
$$
。则关于 $\mathcal{F}$ 的假设恰好意味着相应的列
$$
\bigoplus\nolimits_{j \in J} \mathcal{O}_V \longrightarrow
\bigoplus\nolimits_{i \in I} \mathcal{O}_V \longrightarrow
\mathcal{F}|_{\mathcal{C}/V} \longrightarrow 0
$$
正合，故 $\mathcal{F}$ 拟凝聚。
\end{proof}

\begin{lemma}
\label{sites-modules-lemma-chaotic-flat-quasi-coherent}
设 $\mathcal{C}$ 是范畴，并赋予混沌拓扑而视为站点；见《站点》，例
\ref{sites-example-indiscrete}。设 $\mathcal{O}$ 是 $\mathcal{C}$ 上的环层。
假设对 $\mathcal{C}$ 中所有 $U \to V$，限制映射
$\mathcal{O}(V) \to \mathcal{O}(U)$ 都是平坦环同态。
则拟凝聚 $\mathcal{O}$-模构成 $\textit{Mod}(\mathcal{O})$
的弱 Serre 子范畴。
\end{lemma}

\begin{proof}
检验弱 Serre 子范畴的定义；见《同调代数》，定义
\ref{homology-definition-serre-subcategory}。为此使用引理
\ref{sites-modules-lemma-chaotic-quasi-coherent} 给出的拟凝聚模刻画。考虑正合列
$$
\mathcal{F}_0 \to \mathcal{F}_1 \to
\mathcal{F}_2 \to \mathcal{F}_3 \to
\mathcal{F}_4
$$
，它在 $\textit{Mod}(\mathcal{O})$ 中，且 $\mathcal{F}_0$、
$\mathcal{F}_1$、$\mathcal{F}_3$ 和 $\mathcal{F}_4$ 拟凝聚。
设 $U \to V$ 是 $\mathcal{C}$ 的态射，并考虑交换图
$$
\vcenter{\hbox{\resizebox{0.98\textwidth}{!}{$
\xymatrix{
\mathcal{F}_0(V) \otimes_{\mathcal{O}(V)} \mathcal{O}(U) \ar[r] \ar[d] &
\mathcal{F}_1(V) \otimes_{\mathcal{O}(V)} \mathcal{O}(U) \ar[r] \ar[d] &
\mathcal{F}_2(V) \otimes_{\mathcal{O}(V)} \mathcal{O}(U) \ar[r] \ar[d] &
\mathcal{F}_3(V) \otimes_{\mathcal{O}(V)} \mathcal{O}(U) \ar[r] \ar[d] &
\mathcal{F}_4(V) \otimes_{\mathcal{O}(V)} \mathcal{O}(U) \ar[d] \\
\mathcal{F}_0(U) \ar[r] &
\mathcal{F}_1(U) \ar[r] &
\mathcal{F}_2(U) \ar[r] &
\mathcal{F}_3(U) \ar[r] &
\mathcal{F}_4(U)
}
$}}}
$$
由假设，指标为 $0$、$1$、$3$、$4$ 的竖直箭头都是同构。
由于 $\mathcal{C}$ 上的拓扑是混沌拓扑，在 $\mathcal{C}$ 的对象上取截面正合，
故下行正合。又因 $\mathcal{O}(V) \to \mathcal{O}(U)$ 平坦，上行也正合。
因此由五引理（《同调代数》，引理 \ref{homology-lemma-five-lemma}），
中间箭头是同构。
\end{proof}






\section{环拓扑斯的闭浸入}
\label{sites-modules-section-closed-immersion}

\noindent
何时称环拓扑斯态射
$i : (\Sh(\mathcal{C}), \mathcal{O}) \to (\Sh(\mathcal{D}), \mathcal{O}')$
为闭浸入？类比《模层》，第 \ref{modules-section-closed-immersion} 节的讨论，
至少作如下假设似乎是自然的：
\begin{enumerate}
\item 函子 $i$ 是拓扑斯的闭浸入（《站点》，定义
\ref{sites-definition-immersion-topoi}）。
\item 相伴映射 $\mathcal{O}' \to i_*\mathcal{O}$ 是满射。
\end{enumerate}
这些条件已经蕴含若干令人满意的结果，本节将讨论它们。不过，由于存在多种可用的
定义，不实际定义环拓扑斯闭浸入这一概念似乎更为审慎。

\begin{lemma}
\label{sites-modules-lemma-i-star-equivalence}
设 $i : (\Sh(\mathcal{C}), \mathcal{O}) \to (\Sh(\mathcal{D}), \mathcal{O}')$
是环拓扑斯态射。假设 $i$ 是拓扑斯的闭浸入，且
$i^\sharp : \mathcal{O}' \to i_*\mathcal{O}$ 是满射。以
$\mathcal{I} \subset \mathcal{O}'$ 表示 $i^\sharp$ 的核。函子
$$
i_* :
\textit{Mod}(\mathcal{O})
\longrightarrow
\textit{Mod}(\mathcal{O}')
$$
正合且全忠实，其本质像恰为满足
$\mathcal{I}\mathcal{G} = 0$ 的 $\mathcal{O}'$-模 $\mathcal{G}$。
\end{lemma}

\begin{proof}
由引理 \ref{sites-modules-lemma-exactness} 和《站点》，引理
\ref{sites-lemma-closed-immersion} 可知 $i_*$ 正合。由 $i_*$ 在集合层上全忠实，
以及 $i^\sharp$ 是满射，得到 $i_*$ 作为函子
$\textit{Mod}(\mathcal{O}) \to \textit{Mod}(\mathcal{O}')$ 也是全忠实的。
事实上，设
$\alpha : i_*\mathcal{F}_1 \to i_*\mathcal{F}_2$ 是
$\mathcal{O}'$-模映射。由 $i_*$ 的全忠实性，得到集合层映射
$\beta : \mathcal{F}_1 \to \mathcal{F}_2$。为证明 $\beta$ 是模映射，
须证明
$$
\xymatrix{
\mathcal{O} \times \mathcal{F}_1 \ar[r] \ar[d] &
\mathcal{F}_1 \ar[d] \\
\mathcal{O} \times \mathcal{F}_2 \ar[r] &
\mathcal{F}_2
}
$$
交换。只须在应用 $i_*$ 后证明交换性。考虑
$$
\xymatrix{
\mathcal{O}' \times i_*\mathcal{F}_1 \ar[r] \ar[d] &
i_*\mathcal{O} \times i_*\mathcal{F}_1 \ar[r] \ar[d] &
i_*\mathcal{F}_1 \ar[d] \\
\mathcal{O}' \times i_*\mathcal{F}_2 \ar[r] &
i_*\mathcal{O} \times i_*\mathcal{F}_2 \ar[r] &
i_*\mathcal{F}_2
}
$$
外矩形交换。由于 $i^\sharp$ 是满射，结论成立。

\medskip\noindent
为完成证明，还须证明关于 $i_*$ 本质像的断言。显然，对任意
$\mathcal{O}$-模 $\mathcal{F}$，$i_*\mathcal{F}$ 被 $\mathcal{I}$ 零化。
反之，设 $\mathcal{G}$ 是满足 $\mathcal{I}\mathcal{G} = 0$ 的
$\mathcal{O}'$-模。由闭子拓扑斯的定义，$\mathcal{D}$ 的终对象存在一个子层
$\mathcal{U}$，使 $i_*$ 在集合层上的本质像恰为如下集合层
$\mathcal{H}$ 所成的类：映射
$\mathcal{H} \times \mathcal{U} \to \mathcal{U}$ 是同构。
特别地，$i_*\mathcal{O} \times \mathcal{U} = \mathcal{U}$。这蕴含
$\mathcal{I} \times \mathcal{U} = \mathcal{O} \times \mathcal{U}$.
因此模 $\mathcal{G}$ 满足
$\mathcal{G} \times \mathcal{U} = \{0\} \times \mathcal{U} = \mathcal{U}$
（因为零模同构于集合层的终对象）。从而存在 $\mathcal{C}$ 上的集合层
$\mathcal{F}$，使 $i_*\mathcal{F} = \mathcal{G}$。由于 $i_*$ 在集合层上
全忠实，为定义加法 $\mathcal{F} \times \mathcal{F} \to \mathcal{F}$ 和
乘法 $\mathcal{O} \times \mathcal{F} \to \mathcal{F}$，只须使用加法
$$
\mathcal{G} \times \mathcal{G} \longrightarrow \mathcal{G}
$$
（这是由 $\mathcal{G}$ 的 $\mathcal{O}'$-模结构给定的）以及乘法
$$
i_*\mathcal{O} \times \mathcal{G} \to \mathcal{G}
$$
。该乘法由 $\mathcal{O}'$ 的乘法给出，因为按假设它零化
$\mathcal{I}$，且 $i_*\mathcal{O} = \mathcal{O}'/\mathcal{I}$。
依构造，$\mathcal{G}$ 同构于由此构造的 $\mathcal{O}$-模
$\mathcal{F}$ 的直像。
\end{proof}












\section{张量积}
\label{sites-modules-section-tensor-product}

\noindent
在第 \ref{sites-modules-section-presheaves-modules} 节和第
\ref{sites-modules-section-sheafification-presheaves-modules} 节中，我们用张量积构造定义了
换环函子。这一构造也可以用来定义 $\mathcal{O}$-模预层的张量积。
确切地，设 $\mathcal{C}$ 是范畴，$\mathcal{O}$ 是环预层，且
$\mathcal{F}$、$\mathcal{G}$ 是 $\mathcal{O}$-模预层。此时定义
$\mathcal{F} \otimes_{p, \mathcal{O}} \mathcal{G}$ 为预层
$$
U
\longmapsto
(\mathcal{F} \otimes_{p, \mathcal{O}} \mathcal{G})(U)
=
\mathcal{F}(U) \otimes_{\mathcal{O}(U)} \mathcal{G}(U)
$$
。若 $\mathcal{C}$ 是站点，$\mathcal{O}$ 是环层，且
$\mathcal{F}$、$\mathcal{G}$ 是 $\mathcal{O}$-模层，则定义
$$
\mathcal{F} \otimes_\mathcal{O} \mathcal{G}
=
(\mathcal{F} \otimes_{p, \mathcal{O}} \mathcal{G})^\#
$$
为与预层 $\mathcal{F} \otimes_{p, \mathcal{O}} \mathcal{G}$ 相伴的
$\mathcal{O}$-模层。

\medskip\noindent
下面列出一些此后将直接使用而不再说明的公式：
$$
(\mathcal{F}
\otimes_{p, \mathcal{O}} \mathcal{G})
\otimes_{p, \mathcal{O}} \mathcal{H}
=
\mathcal{F}
\otimes_{p, \mathcal{O}} (\mathcal{G}
\otimes_{p, \mathcal{O}} \mathcal{H}),
$$
，对层也有类似公式。若 $\mathcal{O}_1 \to \mathcal{O}_2$ 是环预层映射，则
$$
(\mathcal{F} \otimes_{p, \mathcal{O}_1} \mathcal{G})
\otimes_{p, \mathcal{O}_1} \mathcal{O}_2 =
(\mathcal{F} \otimes_{p, \mathcal{O}_1} \mathcal{O}_2)
\otimes_{p, \mathcal{O}_2}
(\mathcal{G} \otimes_{p, \mathcal{O}_1} \mathcal{O}_2),
$$
，对层也有类似公式。这些公式来自相应的代数公式和层化。

\begin{lemma}
\label{sites-modules-lemma-sheafification-tensor}
设 $\mathcal{C}$ 是站点，$\mathcal{O}$ 是环预层，且
$\mathcal{F}$、$\mathcal{G}$ 是 $\mathcal{O}$-模预层。则
$\mathcal{F}^\# \otimes_{\mathcal{O}^\#} \mathcal{G}^\#$
等于
$(\mathcal{F} \otimes_{p, \mathcal{O}} \mathcal{G})^\#$.
\end{lemma}

\begin{proof}
从略。提示：使用下文由双线性映射给出的张量积刻画，以及层化的泛性质。
\end{proof}

\noindent
设 $\mathcal{C}$ 是站点，$\mathcal{O}$ 是环层，且
$\mathcal{F}$、$\mathcal{G}$、$\mathcal{H}$ 是 $\mathcal{O}$-模层。
此时定义
$$
\text{Bilin}_\mathcal{O}(\mathcal{F} \times \mathcal{G}, \mathcal{H})
=
\{\varphi \in
\Mor_{\Sh(\mathcal{C})}(
\mathcal{F} \times \mathcal{G}, \mathcal{H}) \mid
\varphi \text{ 是 }\mathcal{O}\text{-双线性的}\}.
$$
由此定义，有
$$
\Hom_\mathcal{O}
(\mathcal{F} \otimes_\mathcal{O} \mathcal{G}, \mathcal{H})
=
\text{Bilin}_\mathcal{O}(\mathcal{F} \times \mathcal{G}, \mathcal{H}).
$$
换言之，$\mathcal{F} \otimes_\mathcal{O} \mathcal{G}$ 表示如下函子：
它将 $\mathcal{H}$ 送到双线性映射
$\mathcal{F} \times \mathcal{G} \to \mathcal{H}$ 的集合。
特别地，由于在环拓扑斯上两个模之间的双线性映射概念有意义，模层的张量积
对拓扑斯而言是内蕴的（比较第 \ref{sites-modules-section-intrinsic} 节的讨论）。
事实上，有下述结果。

\begin{lemma}
\label{sites-modules-lemma-tensor-product-pullback}
设 $f : (\Sh(\mathcal{C}), \mathcal{O}_\mathcal{C})
\to (\Sh(\mathcal{D}), \mathcal{O}_\mathcal{D})$ 是环拓扑斯态射，且
$\mathcal{F}$、$\mathcal{G}$ 是 $\mathcal{O}_\mathcal{D}$-模。则
$f^*(\mathcal{F} \otimes_{\mathcal{O}_\mathcal{D}} \mathcal{G})
= f^*\mathcal{F} \otimes_{\mathcal{O}_\mathcal{C}} f^*\mathcal{G}$
，且关于 $\mathcal{F}$、$\mathcal{G}$ 具有函子性。
\end{lemma}

\begin{proof}
对 $\mathcal{O}_\mathcal{C}$-模层 $\mathcal{H}$，有
\begin{align*}
\Hom_{\mathcal{O}_\mathcal{C}}(
f^*(\mathcal{F} \otimes_\mathcal{O} \mathcal{G}), \mathcal{H})
& =
\Hom_{\mathcal{O}_\mathcal{D}}(
\mathcal{F} \otimes_\mathcal{O} \mathcal{G}, f_*\mathcal{H}) \\
& =
\text{Bilin}_{\mathcal{O}_\mathcal{D}}(
\mathcal{F} \times \mathcal{G}, f_*\mathcal{H}) \\
& =
\text{Bilin}_{f^{-1}\mathcal{O}_\mathcal{D}}(
f^{-1}\mathcal{F} \times f^{-1}\mathcal{G}, \mathcal{H}) \\
& =
\Hom_{f^{-1}\mathcal{O}_\mathcal{D}}(
f^{-1}\mathcal{F} \otimes_{f^{-1}\mathcal{O}_\mathcal{D}} f^{-1}\mathcal{G},
\mathcal{H}) \\
& =
\Hom_{\mathcal{O}_\mathcal{C}}(
f^*\mathcal{F} \otimes_{f^*\mathcal{O}_\mathcal{D}} f^*\mathcal{G},
\mathcal{H})
\end{align*}
这一串等式中有实质内容的是第三个“$=$”。它直接来自定义以及前面各节所讨论的
$f_*$ 与 $f^{-1}$ 的伴随性。
\end{proof}

\begin{lemma}
\label{sites-modules-lemma-tensor-product-permanence}
设 $(\mathcal{C}, \mathcal{O})$ 是环站点，且 $\mathcal{F}$、
$\mathcal{G}$ 是 $\mathcal{O}$-模层。
\begin{enumerate}
\item 若 $\mathcal{F}$、$\mathcal{G}$ 局部自由，则
$\mathcal{F} \otimes_\mathcal{O} \mathcal{G}$ 也局部自由。
\item 若 $\mathcal{F}$、$\mathcal{G}$ 有限局部自由，则
$\mathcal{F} \otimes_\mathcal{O} \mathcal{G}$ 也有限局部自由。
\item 若 $\mathcal{F}$、$\mathcal{G}$ 局部由截面生成，则
$\mathcal{F} \otimes_\mathcal{O} \mathcal{G}$ 也局部由截面生成。
\item 若 $\mathcal{F}$、$\mathcal{G}$ 是有限型的，则
$\mathcal{F} \otimes_\mathcal{O} \mathcal{G}$ 也是有限型的。
\item 若 $\mathcal{F}$、$\mathcal{G}$ 拟凝聚，则
$\mathcal{F} \otimes_\mathcal{O} \mathcal{G}$ 也拟凝聚。
\item 若 $\mathcal{F}$、$\mathcal{G}$ 是有限表示的，则
$\mathcal{F} \otimes_\mathcal{O} \mathcal{G}$ 也是有限表示的。
\item 若 $\mathcal{F}$ 是有限表示的，且 $\mathcal{G}$ 凝聚，则
$\mathcal{F} \otimes_\mathcal{O} \mathcal{G}$ 凝聚。
\item 若 $\mathcal{F}$、$\mathcal{G}$ 凝聚，则
$\mathcal{F} \otimes_\mathcal{O} \mathcal{G}$ 也凝聚。
\end{enumerate}
\end{lemma}

\begin{proof}
从略。提示：比较《模层》，引理 \ref{modules-lemma-tensor-product-permanence}。
\end{proof}



\section{内 Hom}
\label{sites-modules-section-internal-hom}

\noindent
设 $\mathcal{C}$ 是范畴，$\mathcal{O}$ 是环预层，且
$\mathcal{F}$、$\mathcal{G}$ 是 $\mathcal{O}$-模预层。考虑规则
$$
U \longmapsto \Hom_{\mathcal{O}_U}(\mathcal{F}|_U, \mathcal{G}|_U).
$$
。对 $\mathcal{C}$ 中的 $\varphi : V \to U$，定义限制映射
$$
\Hom_{\mathcal{O}_U}(\mathcal{F}|_U, \mathcal{G}|_U)
\longrightarrow
\Hom_{\mathcal{O}_V}(\mathcal{F}|_V, \mathcal{G}|_V)
$$
，即经由再局部化态射 $j : \mathcal{C}/V \to \mathcal{C}/U$ 作限制；
见《站点》，引理 \ref{sites-lemma-relocalize}。由此定义预层
$\SheafHom_\mathcal{O}(\mathcal{F}, \mathcal{G})$。此外，给定元素
$\varphi \in \Hom_{\mathcal{O}|_U}(\mathcal{F}|_U, \mathcal{G}|_U)$
和截面 $f \in \mathcal{O}(U)$，可以定义
$f\varphi \in \Hom_{\mathcal{O}|_U}(\mathcal{F}|_U, \mathcal{G}|_U)$
，方法是与 $\mathcal{F}|_U$ 上乘以 $f$ 的映射前复合，或与
$\mathcal{G}|_U$ 上乘以 $f$ 的映射后复合（两者结果相同）。
因此实际得到 $\mathcal{O}$-模预层。存在典范的“求值”态射
$$
\mathcal{F}
\otimes_{p, \mathcal{O}}
\SheafHom_\mathcal{O}(\mathcal{F}, \mathcal{G})
\longrightarrow
\mathcal{G}.
$$

\begin{lemma}
\label{sites-modules-lemma-internal-hom}
若 $\mathcal{C}$ 是站点，$\mathcal{O}$ 是环层，$\mathcal{F}$ 是
$\mathcal{O}$-模预层，且 $\mathcal{G}$ 是 $\mathcal{O}$-模层，则
$\SheafHom_\mathcal{O}(\mathcal{F}, \mathcal{G})$
是 $\mathcal{O}$-模层。
\end{lemma}

\begin{proof}
从略。提示：先注意
$\SheafHom_\mathcal{O}(\mathcal{F}, \mathcal{G})
= \SheafHom_\mathcal{O}(\mathcal{F}^\#, \mathcal{G})$，这把问题约化到
$\mathcal{F}$ 和 $\mathcal{G}$ 都是层的情形。关于集合层的结果见
《站点》，引理 \ref{sites-lemma-glue-maps}。
\end{proof}

\begin{lemma}
\label{sites-modules-lemma-internal-hom-restriction}
设 $(\mathcal{C}, \mathcal{O})$ 是环站点，且 $\mathcal{F}$、
$\mathcal{G}$ 是 $\mathcal{O}$-模层。则对 $U \in \Ob(\mathcal{C})$，
$\SheafHom_\mathcal{O}(\mathcal{F}, \mathcal{G})$ 的形成与到 $U$ 的限制交换。
\end{lemma}

\begin{proof}
由定义立即可得。
\end{proof}

\begin{remark}
\label{sites-modules-remark-pullback-internal-hom}
设 $f : (\mathcal{C}, \mathcal{O}_\mathcal{C}) \to
(\mathcal{D}, \mathcal{O}_\mathcal{D})$ 是环站点态射，且
$\mathcal{F}$、$\mathcal{G}$ 是 $\mathcal{O}_\mathcal{D}$-模层。
存在典范映射
$$
f^*\SheafHom_{\mathcal{O}_\mathcal{D}}(\mathcal{F}, \mathcal{G})
\longrightarrow
\SheafHom_{\mathcal{O}_\mathcal{C}}(f^*\mathcal{F}, f^*\mathcal{G})
$$
。事实上，此映射与如下映射伴随：
$$
\SheafHom_{\mathcal{O}_\mathcal{D}}(\mathcal{F}, \mathcal{G})
\longrightarrow
f_*\SheafHom_{\mathcal{O}_\mathcal{C}}(f^*\mathcal{F}, f^*\mathcal{G})
$$
该映射定义如下。设 $f$ 由连续函子 $u : \mathcal{D} \to \mathcal{C}$ 给出。
对 $V \in \Ob(\mathcal{D})$ 上的截面，使用映射
\begin{align*}
\Gamma(V, \SheafHom_{\mathcal{O}_\mathcal{D}}(\mathcal{F}, \mathcal{G}))
& =
\Hom_{\mathcal{O}_V}(\mathcal{F}|_V, \mathcal{G}|_V) \\
& \longrightarrow
\Hom_{\mathcal{O}_{u(V)}}(f^*\mathcal{F}|_{u(V)}, f^*\mathcal{G}|_{u(V)}) \\
& =
\Gamma(u(V), \SheafHom_{\mathcal{O}_\mathcal{C}}(f^*\mathcal{F},
f^*\mathcal{G})) \\
& =
\Gamma(V, f_*\SheafHom_{\mathcal{O}_\mathcal{C}}(f^*\mathcal{F},
f^*\mathcal{G}))
\end{align*}
，其中箭头使用由 $f$ 诱导的态射
$(\mathcal{C}/u(V), \mathcal{O}_{u(V)}) \to
(\mathcal{D}/V, \mathcal{O}_V)$ 的拉回。
\end{remark}

\noindent
在引理 \ref{sites-modules-lemma-internal-hom} 的情形中，“求值”态射通过模层的张量积分解：
$$
\mathcal{F}
\otimes_\mathcal{O}
\SheafHom_\mathcal{O}(\mathcal{F}, \mathcal{G})
\longrightarrow
\mathcal{G}.
$$

\begin{lemma}
\label{sites-modules-lemma-internal-hom-commute-limits}
内 Hom 与（余）极限。设 $\mathcal{C}$ 是范畴，且 $\mathcal{O}$ 是环预层。
\begin{enumerate}
\item 对任意 $\mathcal{O}$-模预层 $\mathcal{F}$，函子
$$
\textit{PMod}(\mathcal{O}) \longrightarrow \textit{PMod}(\mathcal{O})
, \quad
\mathcal{G} \longmapsto \SheafHom_\mathcal{O}(\mathcal{F}, \mathcal{G})
$$
与任意极限交换。
\item 对任意 $\mathcal{O}$-模预层 $\mathcal{G}$，函子
$$
\textit{PMod}(\mathcal{O}) \longrightarrow \textit{PMod}(\mathcal{O})^{opp}
, \quad
\mathcal{F} \longmapsto \SheafHom_\mathcal{O}(\mathcal{F}, \mathcal{G})
$$
与任意余极限交换；公式为
$$
\SheafHom_\mathcal{O}(\colim_i \mathcal{F}_i, \mathcal{G})
=
\lim_i \SheafHom_\mathcal{O}(\mathcal{F}_i, \mathcal{G}).
$$
\end{enumerate}
再假设 $\mathcal{C}$ 是站点，且 $\mathcal{O}$ 是环层。
\begin{enumerate}
\item[(3)] 对任意 $\mathcal{O}$-模层 $\mathcal{F}$，函子
$$
\textit{Mod}(\mathcal{O}) \longrightarrow \textit{Mod}(\mathcal{O})
, \quad
\mathcal{G} \longmapsto \SheafHom_\mathcal{O}(\mathcal{F}, \mathcal{G})
$$
与任意极限交换。
\item[(4)] 对任意 $\mathcal{O}$-模层 $\mathcal{G}$，函子
$$
\textit{Mod}(\mathcal{O}) \longrightarrow \textit{Mod}(\mathcal{O})^{opp}
, \quad
\mathcal{F} \longmapsto \SheafHom_\mathcal{O}(\mathcal{F}, \mathcal{G})
$$
与任意余极限交换；公式为
$$
\SheafHom_\mathcal{O}(\colim_i \mathcal{F}_i, \mathcal{G})
=
\lim_i \SheafHom_\mathcal{O}(\mathcal{F}_i, \mathcal{G}).
$$
\end{enumerate}
\end{lemma}

\begin{proof}
设 $\mathcal{I} \to \textit{PMod}(\mathcal{O})$、
$i \mapsto \mathcal{G}_i$ 是图，且 $U$ 是范畴 $\mathcal{C}$ 的对象。
由于 $j_U^*$ 既是左伴随又是右伴随，有
$\lim_i j_U^*\mathcal{G}_i = j_U^* \lim_i \mathcal{G}_i$.
因此
\begin{align*}
\SheafHom_\mathcal{O}(\mathcal{F}, \lim_i \mathcal{G}_i)(U)
& =
\Hom_{\mathcal{O}_U}(\mathcal{F}|_U, \lim_i \mathcal{G}_i|_U) \\
& =
\lim_i \Hom_{\mathcal{O}_U}(\mathcal{F}|_U, \mathcal{G}_i|_U) \\
& = \lim_i \SheafHom_\mathcal{O}(\mathcal{F}, \mathcal{G}_i)(U)
\end{align*}
，其中用到了极限的定义。这证明了 (1)。(2) 的证明完全相同。
(3) 来自 (1)，因为层图的极限与预层范畴中的极限相同。
最后，(4) 成立，因为公式中有
$$
\Mor_{\textit{Mod}(\mathcal{O})}(
\colim_i \mathcal{F}_i, \mathcal{G})
=
\Mor_{\textit{PMod}(\mathcal{O})}(
\colim^{PSh}_i \mathcal{F}_i, \mathcal{G})
$$
，这是因为余极限 $\colim_i \mathcal{F}_i$ 是
$\textit{PMod}(\mathcal{O})$ 中余极限
$\colim^{PSh}_i \mathcal{F}_i$ 的层化。因此 (4) 来自 (2)
（并再次使用上述关于极限的说明）。
\end{proof}

\begin{lemma}
\label{sites-modules-lemma-internal-hom-exact}
设 $(\mathcal{C}, \mathcal{O})$ 是环站点，且 $\mathcal{F}$、
$\mathcal{G}$ 是 $\mathcal{O}$-模。
\begin{enumerate}
\item 若 $\mathcal{F}_2 \to \mathcal{F}_1 \to \mathcal{F} \to 0$
是 $\mathcal{O}$-模的正合列，则
$$
0 \to
\SheafHom_\mathcal{O}(\mathcal{F}, \mathcal{G}) \to
\SheafHom_\mathcal{O}(\mathcal{F}_1, \mathcal{G}) \to
\SheafHom_\mathcal{O}(\mathcal{F}_2, \mathcal{G})
$$
正合。
\item 若 $0 \to \mathcal{G} \to \mathcal{G}_1 \to \mathcal{G}_2$
是 $\mathcal{O}$-模的正合列，则
$$
0 \to
\SheafHom_\mathcal{O}(\mathcal{F}, \mathcal{G}) \to
\SheafHom_\mathcal{O}(\mathcal{F}, \mathcal{G}_1) \to
\SheafHom_\mathcal{O}(\mathcal{F}, \mathcal{G}_2)
$$
正合。
\end{enumerate}
\end{lemma}

\begin{proof}
由引理 \ref{sites-modules-lemma-internal-hom-commute-limits} 和《同调代数》，引理
\ref{homology-lemma-exact-functor} 可得。
\end{proof}

\begin{lemma}
\label{sites-modules-lemma-internal-hom-adjoint-tensor}
设 $\mathcal{C}$ 是范畴，且 $\mathcal{O}$ 是环预层。
\begin{enumerate}
\item 设 $\mathcal{F}$、$\mathcal{G}$、$\mathcal{H}$ 是
$\mathcal{O}$-模预层。存在典范同构
$$
\SheafHom_\mathcal{O}
(\mathcal{F} \otimes_{p, \mathcal{O}} \mathcal{G}, \mathcal{H})
\longrightarrow
\SheafHom_\mathcal{O}
(\mathcal{F}, \SheafHom_\mathcal{O}(\mathcal{G}, \mathcal{H}))
$$
，它关于三个变元均具有函子性（三处都取层 Hom）。特别地，
$$
\Mor_{\textit{PMod}(\mathcal{O})}(
\mathcal{F} \otimes_{p, \mathcal{O}} \mathcal{G}, \mathcal{H})
=
\Mor_{\textit{PMod}(\mathcal{O})}(
\mathcal{F}, \SheafHom_\mathcal{O}(\mathcal{G}, \mathcal{H}))
$$
\item
假设 $\mathcal{C}$ 是站点，$\mathcal{O}$ 是环层，且
$\mathcal{F}$、$\mathcal{G}$、$\mathcal{H}$ 是 $\mathcal{O}$-模层。
存在典范同构
$$
\SheafHom_\mathcal{O}
(\mathcal{F} \otimes_\mathcal{O} \mathcal{G}, \mathcal{H})
\longrightarrow
\SheafHom_\mathcal{O}
(\mathcal{F}, \SheafHom_\mathcal{O}(\mathcal{G}, \mathcal{H}))
$$
，它关于三个变元均具有函子性（三处都取层 Hom）。特别地，
$$
\Mor_{\textit{Mod}(\mathcal{O})}(
\mathcal{F} \otimes_\mathcal{O} \mathcal{G}, \mathcal{H})
=
\Mor_{\textit{Mod}(\mathcal{O})}(
\mathcal{F}, \SheafHom_\mathcal{O}(\mathcal{G}, \mathcal{H}))
$$
\end{enumerate}
\end{lemma}

\begin{proof}
这是《代数》，引理 \ref{algebra-lemma-hom-from-tensor-product} 的类似版本。
证明相同，从略。
\end{proof}

\begin{lemma}
\label{sites-modules-lemma-tensor-commute-colimits}
张量积与余极限。设 $\mathcal{C}$ 是范畴，且 $\mathcal{O}$ 是环预层。
\begin{enumerate}
\item 对任意 $\mathcal{O}$-模预层 $\mathcal{F}$，函子
$$
\textit{PMod}(\mathcal{O}) \longrightarrow \textit{PMod}(\mathcal{O})
, \quad
\mathcal{G} \longmapsto \mathcal{F} \otimes_{p, \mathcal{O}} \mathcal{G}
$$
与任意余极限交换。
\item
假设 $\mathcal{C}$ 是站点，且 $\mathcal{O}$ 是环层。对任意
$\mathcal{O}$-模层 $\mathcal{F}$，函子
$$
\textit{Mod}(\mathcal{O}) \longrightarrow \textit{Mod}(\mathcal{O})
, \quad
\mathcal{G} \longmapsto \mathcal{F} \otimes_\mathcal{O} \mathcal{G}
$$
与任意余极限交换。
\end{enumerate}
\end{lemma}

\begin{proof}
这是因为由引理 \ref{sites-modules-lemma-internal-hom-adjoint-tensor}，张量积与内 Hom 伴随。
见《范畴》，引理 \ref{categories-lemma-adjoint-exact}。
\end{proof}

\begin{lemma}
\label{sites-modules-lemma-adjoint-hom-restrict}
设 $\mathcal{C}$ 是范畴，相应地是站点；设
$\mathcal{O} \to \mathcal{O}'$ 是环预层映射，相应地是环层映射。则
$$
\Hom_\mathcal{O}(\mathcal{G}, \mathcal{F}) =
\Hom_{\mathcal{O}'}(\mathcal{G},
\SheafHom_\mathcal{O}(\mathcal{O}', \mathcal{F}))
$$
，其中 $\mathcal{G}$ 是任意 $\mathcal{O}'$-模，$\mathcal{F}$ 是任意
$\mathcal{O}$-模。
\end{lemma}

\begin{proof}
这是《代数》，引理 \ref{algebra-lemma-adjoint-hom-restrict} 的类似版本。
证明相同，从略。
\end{proof}

\begin{lemma}
\label{sites-modules-lemma-j-shriek-and-tensor}
设 $(\mathcal{C}, \mathcal{O})$ 是环站点，且
$U \in \Ob(\mathcal{C})$。对
$\textit{Mod}(\mathcal{O}_U)$ 中的 $\mathcal{G}$ 和
$\textit{Mod}(\mathcal{O})$ 中的 $\mathcal{F}$，有
$j_{U!}\mathcal{G} \otimes_\mathcal{O} \mathcal{F} =
j_{U!}(\mathcal{G} \otimes_{\mathcal{O}_U} \mathcal{F}|_U)$.
\end{lemma}

\begin{proof}
设 $\mathcal{H}$ 是 $\textit{Mod}(\mathcal{O})$ 的对象。则
\begin{align*}
\Hom_\mathcal{O}(j_{U!}(\mathcal{G} \otimes_{\mathcal{O}_U} \mathcal{F}|_U),
\mathcal{H})
& =
\Hom_{\mathcal{O}_U}(\mathcal{G} \otimes_{\mathcal{O}_U} \mathcal{F}|_U,
\mathcal{H}|_U) \\
& =
\Hom_{\mathcal{O}_U}(\mathcal{G},
\SheafHom_{\mathcal{O}_U}(\mathcal{F}|_U, \mathcal{H}|_U)) \\
& =
\Hom_{\mathcal{O}_U}(\mathcal{G},
\SheafHom_\mathcal{O}(\mathcal{F}, \mathcal{H})|_U) \\
& =
\Hom_\mathcal{O}(j_{U!}\mathcal{G},
\SheafHom_\mathcal{O}(\mathcal{F}, \mathcal{H})) \\
& =
\Hom_\mathcal{O}(j_{U!}\mathcal{G} \otimes_\mathcal{O} \mathcal{F},
\mathcal{H})
\end{align*}
第一个等式成立是因为 $j_{U!}$ 是模限制的左伴随；第二个等式由引理
\ref{sites-modules-lemma-internal-hom-adjoint-tensor} 得到；第三个等式由引理
\ref{sites-modules-lemma-internal-hom-restriction} 得到；第四个等式成立是因为
$j_{U!}$ 是模限制的左伴随；第五个等式由引理
\ref{sites-modules-lemma-internal-hom-adjoint-tensor} 得到。再由 Yoneda 引理即得本引理。
\end{proof}

\begin{remark}
\label{sites-modules-remark-j-shriek-tensor}
设 $\mathcal{C}$ 是站点。设 $\mathcal{F}$ 是 $\mathcal{C}$ 上的集合层，并考虑局部化态射
$j : \Sh(\mathcal{C})/\mathcal{F} \to \Sh(\mathcal{C})$.
见《站点》，定义 \ref{sites-definition-localize-topos}。
我们断言：(a) $j_!\mathbf{Z} = \mathbf{Z}_\mathcal{F}^\#$；
(b) 对 $\mathcal{C}$ 上任意阿贝尔层 $\mathcal{H}$，都有
$j_!(j^{-1}\mathcal{H}) = j_!\mathbf{Z} \otimes_\mathbf{Z} \mathcal{H}$。
令 $\mathcal{G}$ 为 $\mathcal{C}$ 上的阿贝尔层。由 Yoneda 引理，(a) 成立，因为
$$
\Hom(j_!\mathbf{Z}, \mathcal{G}) =
\Hom(\mathbf{Z}, j^{-1}\mathcal{G}) =
\Hom(\mathbf{Z}_\mathcal{F}^\#, \mathcal{G})
$$
，其中第二个等式成立，是因为等式两边都给出从 $\mathcal{F}$ 到
$\mathcal{G}$ 的映射集，并将其视为阿贝尔群。对于 (b)，使用 Yoneda 引理以及
\begin{align*}
\Hom(j_!(j^{-1}\mathcal{H}), \mathcal{G})
& =
\Hom(j^{-1}\mathcal{H}, j^{-1}\mathcal{G}) \\
& =
\Hom(\mathbf{Z}, \SheafHom(j^{-1}\mathcal{H}, j^{-1}\mathcal{G})) \\
& =
\Hom(\mathbf{Z}, j^{-1}\SheafHom(\mathcal{H}, \mathcal{G})) \\
& =
\Hom(j_!\mathbf{Z}, \SheafHom(\mathcal{H}, \mathcal{G})) \\
& =
\Hom(j_!\mathbf{Z} \otimes_\mathbf{Z} \mathcal{H}, \mathcal{G})
\end{align*}
即可。这里用到了伴随性、取 $\SheafHom$ 与局部化交换这一事实，以及引理
\ref{sites-modules-lemma-internal-hom-adjoint-tensor}。
\end{remark}

\begin{lemma}
\label{sites-modules-lemma-hom-finite-presentation-colimit}
设 $(\mathcal{C}, \mathcal{O})$ 是环站点，且 $\mathcal{F}$ 是有限表示
$\mathcal{O}$-模。设
$\mathcal{G} = \colim_{\lambda \in \Lambda} \mathcal{G}_\lambda$
是 $\mathcal{O}$-模的滤过余极限。则典范映射
$$
\colim_\lambda \SheafHom_\mathcal{O}(\mathcal{F}, \mathcal{G}_\lambda)
\longrightarrow
\SheafHom_\mathcal{O}(\mathcal{F}, \mathcal{G})
$$
是同构。
\end{lemma}

\begin{proof}
只需证明对 $\mathcal{C}$ 中每个 $U$ 作限制后，该箭头是同构。模层的余极限和
内 Hom 都与到 $U$ 的限制交换；例如见引理
\ref{sites-modules-lemma-exactness-pushforward-pullback} 和
\ref{sites-modules-lemma-internal-hom-restriction}。固定 $U$。给定覆盖
$\{U_i \to U\}_{i \in I}$，只需证明到每个 $U_i$ 的限制是同构。因此可以假设
$\mathcal{F}$ 有整体表示
$$
\bigoplus\nolimits_{j = 1, \ldots, m}
\mathcal{O}
\longrightarrow
\bigoplus\nolimits_{i = 1, \ldots, n}
\mathcal{O}
\to
\mathcal{F}
\to
0
$$
。函子 $\SheafHom_\mathcal{O}(-, -)$ 对任一变元都与有限直和交换，且
$\SheafHom_\mathcal{O}(\mathcal{O}, -)$ 是恒等函子。由此以及引理
\ref{sites-modules-lemma-internal-hom-exact}，得到正合列
$$
0 \to
\SheafHom_\mathcal{O}(\mathcal{F}, \mathcal{G}) \to
\bigoplus\nolimits_{i = 1, \ldots, n} \mathcal{G} \to
\bigoplus\nolimits_{j = 1, \ldots, m} \mathcal{G}
$$
。由引理 \ref{sites-modules-lemma-limits-colimits}，滤过余极限在
$\textit{Mod}(\mathcal{O})$ 中正合，所以下列交换图的上行也正合：
{\small
$$
\xymatrix{
0 \ar[r] &
\colim_\lambda \SheafHom_\mathcal{O}(\mathcal{F}, \mathcal{G}_\lambda)
\ar[r] \ar[d] &
\colim_\lambda \bigoplus\nolimits_{i = 1, \ldots, n} \mathcal{G}_\lambda
\ar[r] \ar[d] &
\colim_\lambda \bigoplus\nolimits_{j = 1, \ldots, m} \mathcal{G}_\lambda
\ar[d] \\
0 \ar[r] &
\SheafHom_\mathcal{O}(\mathcal{F}, \mathcal{G}) \ar[r] &
\bigoplus\nolimits_{i = 1, \ldots, n} \mathcal{G} \ar[r] &
\bigoplus\nolimits_{j = 1, \ldots, m} \mathcal{G}
}
$$
}
右侧两个竖直箭头都是同构，故结论成立。
\end{proof}

\begin{lemma}
\label{sites-modules-lemma-finite-presentation-quasi-compact-colimit}
设 $(\mathcal{C}, \mathcal{O})$ 是环站点。设
$\mathcal{G} = \colim_{\lambda \in \Lambda} \mathcal{G}_\lambda$
是 $\mathcal{O}$-模的滤过余极限，且 $\mathcal{F}$ 是有限表示
$\mathcal{O}$-模。若站点 $\mathcal{C}$ 满足《站点》，引理
\ref{sites-lemma-directed-colimits-global-sections} 第 (4) 部分的假设，则有
$$
\colim_\lambda \Hom_\mathcal{O}(\mathcal{F}, \mathcal{G}_\lambda)
=
\Hom_\mathcal{O}(\mathcal{F}, \mathcal{G}).
$$
另见《站点》，注 \ref{sites-remark-stronger-conditions}。
\end{lemma}

\begin{proof}
置 $\mathcal{H} =
\SheafHom_\mathcal{O}(\mathcal{F}, \colim \mathcal{G}_\lambda)$
以及 $\mathcal{H}_\lambda =
\SheafHom_\mathcal{O}(\mathcal{F}, \mathcal{G}_\lambda)$。由构造可知
$$
\Hom_\mathcal{O}(\mathcal{F}, \mathcal{G}) = \Gamma(\mathcal{C}, \mathcal{H})
\quad\text{且}\quad
\Hom_\mathcal{O}(\mathcal{F}, \mathcal{G}_\lambda) =
\Gamma(\mathcal{C}, \mathcal{H}_\lambda)
$$
。由引理 \ref{sites-modules-lemma-hom-finite-presentation-colimit}，有
$\mathcal{H} = \colim \mathcal{H}_\lambda$。故本引理由《站点》，引理
\ref{sites-lemma-directed-colimits-global-sections} 得出。
\end{proof}









\section{平坦模}
\label{sites-modules-section-flat}

\noindent
可以与环上模的情形完全相同地定义平坦模。

\begin{definition}
\label{sites-modules-definition-flat}
设 $\mathcal{C}$ 是范畴，且 $\mathcal{O}$ 是环预层。
\begin{enumerate}
\item 如果函子
$$
\textit{PMod}(\mathcal{O})
\longrightarrow
\textit{PMod}(\mathcal{O}), \quad
\mathcal{G} \mapsto \mathcal{G} \otimes_{p, \mathcal{O}} \mathcal{F}
$$
正合，则称 $\mathcal{O}$-模预层 $\mathcal{F}$ 是{\it 平坦的}。
\item 如果 $\mathcal{O}'$ 作为 $\mathcal{O}$-模预层是平坦的，则称环预层映射
$\mathcal{O} \to \mathcal{O}'$ 是{\it 平坦的}。
\item 假设 $\mathcal{C}$ 是站点，$\mathcal{O}$ 是环层，且
$\mathcal{F}$ 是 $\mathcal{O}$-模层。如果函子
$$
\textit{Mod}(\mathcal{O})
\longrightarrow
\textit{Mod}(\mathcal{O}), \quad
\mathcal{G} \mapsto \mathcal{G} \otimes_\mathcal{O} \mathcal{F}
$$
正合，则称 $\mathcal{F}$ 是{\it 平坦的}。
\item 如果 $\mathcal{O}'$ 作为 $\mathcal{O}$-模层是平坦的，则称站点上的环层映射
$\mathcal{O} \to \mathcal{O}'$ 是{\it 平坦的}。
\end{enumerate}
\end{definition}

\noindent
平坦模和平坦环映射的概念是内蕴的（第 \ref{sites-modules-section-intrinsic} 节）。

\begin{lemma}
\label{sites-modules-lemma-flatness-presheaves}
设 $\mathcal{C}$ 是范畴，$\mathcal{O}$ 是环预层，且
$\mathcal{F}$ 是 $\mathcal{O}$-模预层。如果每个 $\mathcal{F}(U)$ 都是平坦
$\mathcal{O}(U)$-模，则 $\mathcal{F}$ 是平坦的。
\end{lemma}

\begin{proof}
由定义立即得到。
\end{proof}

\begin{lemma}
\label{sites-modules-lemma-flatness-sheafification}
设 $\mathcal{C}$ 是站点，$\mathcal{O}$ 是环预层，且
$\mathcal{F}$ 是 $\mathcal{O}$-模预层。如果 $\mathcal{F}$ 是平坦
$\mathcal{O}$-模，则 $\mathcal{F}^\#$ 是平坦 $\mathcal{O}^\#$-模。
\end{lemma}

\begin{proof}
略。（提示：层化是正合的。）
\end{proof}

\begin{lemma}
\label{sites-modules-lemma-flatness-sheafification-refined}
设 $\mathcal{C}$ 是站点，$\mathcal{O}$ 是环预层，且
$\mathcal{F}$ 是 $\mathcal{O}$-模预层。假设 $\mathcal{C}$ 的每个对象 $U$
都有覆盖 $\{U_i \to U\}_{i \in I}$，使得 $\mathcal{F}(U_i)$ 是平坦
$\mathcal{O}(U_i)$-模。则 $\mathcal{F}^\#$ 是平坦 $\mathcal{O}^\#$-模。
\end{lemma}

\begin{proof}
设 $\mathcal{G} \subset \mathcal{G}'$ 是 $\mathcal{O}^\#$-模的包含。需要证明
$$
\mathcal{G} \otimes_{\mathcal{O}^\#} \mathcal{F}^\#
\longrightarrow
\mathcal{G}' \otimes_{\mathcal{O}^\#} \mathcal{F}^\#
$$
是单射。由引理 \ref{sites-modules-lemma-sheafification-tensor}，该箭头的源是预层
$\mathcal{G} \otimes_{p, \mathcal{O}} \mathcal{F}$ 的层化，靶也类似。
如果 $U$ 是 $\mathcal{C}$ 的对象，且 $\mathcal{F}(U)$ 是平坦
$\mathcal{O}(U)$-模，则
$$
(\mathcal{G} \otimes_{p, \mathcal{O}} \mathcal{F})(U) =
\mathcal{G}(U) \otimes_{\mathcal{O}(U)} \mathcal{F}(U)
\longrightarrow
\mathcal{G}'(U) \otimes_{\mathcal{O}(U)} \mathcal{F}(U) =
(\mathcal{G}' \otimes_{p, \mathcal{O}} \mathcal{F})(U)
$$
是单射。因此归结为证明：给定 $\mathcal{C}$ 上预层映射
$f : \mathcal{H} \to \mathcal{H}'$，如果 $\mathcal{C}$ 中每个 $U$ 都有覆盖
$\{U_i \to U\}_{i \in I}$，使得
$\mathcal{H}(U_i) \to \mathcal{H}'(U_i)$ 是单射，则 $f^\#$ 是单射。
将此留给读者作为练习。
\end{proof}

\begin{lemma}
\label{sites-modules-lemma-colimits-flat}
余极限与张量积。
\begin{enumerate}
\item 平坦模预层的滤过余极限是平坦的。平坦模预层的直和是平坦的。
\item 平坦模层的滤过余极限是平坦的。平坦模层的直和是平坦的。
\end{enumerate}
\end{lemma}

\begin{proof}
考察对象上的截面，由引理 \ref{sites-modules-lemma-tensor-commute-colimits} 和《代数》，引理
\ref{algebra-lemma-directed-colimit-exact} 即得 (1)。对于 (2)，使用引理
\ref{sites-modules-lemma-tensor-commute-colimits} 以及正合复形的滤过余极限仍为正合复形这一事实
（这里用到层化正合且与余极限交换）。略去若干细节。
\end{proof}

\begin{lemma}
\label{sites-modules-lemma-restriction-flat}
设 $(\mathcal{C}, \mathcal{O})$ 是环站点，且 $U$ 是 $\mathcal{C}$ 的对象。
如果 $\mathcal{F}$ 是平坦 $\mathcal{O}$-模，则 $\mathcal{F}|_U$ 是平坦
$\mathcal{O}_U$-模。
\end{lemma}

\begin{proof}
设 $\mathcal{G}_1 \to \mathcal{G}_2 \to \mathcal{G}_3$ 是
$\mathcal{O}_U$-模的正合复形。由于 $j_{U!}$ 正合（引理
\ref{sites-modules-lemma-extension-by-zero-exact}），且 $\mathcal{F}$ 作为
$\mathcal{O}$-模是平坦的，故由下列模组成的复形正合：
$$
j_{U!}(\mathcal{G}_i \otimes_{\mathcal{O}_U} \mathcal{F}|_U) =
j_{U!}\mathcal{G}_i \otimes_\mathcal{O} \mathcal{F}
$$
（引理 \ref{sites-modules-lemma-j-shriek-and-tensor}）。由引理
\ref{sites-modules-lemma-j-shriek-reflects-exactness} 可知
$\mathcal{G}_1 \otimes_{\mathcal{O}_U} \mathcal{F}|_U \to
\mathcal{G}_2 \otimes_{\mathcal{O}_U} \mathcal{F}|_U \to
\mathcal{G}_3 \otimes_{\mathcal{O}_U} \mathcal{F}|_U$
正合。
\end{proof}

\begin{lemma}
\label{sites-modules-lemma-j-shriek-flat}
设 $\mathcal{C}$ 是范畴，$\mathcal{O}$ 是环预层，且 $U$ 是
$\mathcal{C}$ 的对象。考虑函子 $j_U : \mathcal{C}/U \to \mathcal{C}$。
\begin{enumerate}
\item $\mathcal{O}$-模预层 $j_{U!}\mathcal{O}_U$ 是平坦的（见注
\ref{sites-modules-remark-localize-presheaves}）。
\item 如果 $\mathcal{C}$ 是站点且 $\mathcal{O}$ 是环层，则
$j_{U!}\mathcal{O}_U$ 是平坦 $\mathcal{O}$-模层。
\end{enumerate}
\end{lemma}

\begin{proof}
证明 (1)。由注 \ref{sites-modules-remark-localize-presheaves} 中的讨论可知
$$
j_{U!}\mathcal{O}_U(V)
=
\bigoplus\nolimits_{\varphi \in \Mor_\mathcal{C}(V, U)}
\mathcal{O}(V)
$$
，它是平坦 $\mathcal{O}(V)$-模。因此 (1) 由引理
\ref{sites-modules-lemma-flatness-presheaves} 得到。又因为
$j_{U!}\mathcal{O}_U = (j_{U!}\mathcal{O}_U)^\#$（第一个 $j_{U!}$
作用于层，第二个作用于预层），故由引理
\ref{sites-modules-lemma-flatness-sheafification} 得到 (2)。
\end{proof}

\begin{lemma}
\label{sites-modules-lemma-module-quotient-flat}
设 $\mathcal{C}$ 是范畴，且 $\mathcal{O}$ 是环预层。
\begin{enumerate}
\item 任意 $\mathcal{O}$-模预层都是某个直和
$\bigoplus j_{U_i!}\mathcal{O}_{U_i}$ 的商。
\item 任意 $\mathcal{O}$-模预层都是某个平坦 $\mathcal{O}$-模预层的商。
\item 如果 $\mathcal{C}$ 是站点且 $\mathcal{O}$ 是环层，则任意
$\mathcal{O}$-模层都是某个直和
$\bigoplus j_{U_i!}\mathcal{O}_{U_i}$ 的商。
\item 如果 $\mathcal{C}$ 是站点且 $\mathcal{O}$ 是环层，则任意
$\mathcal{O}$-模层都是某个平坦 $\mathcal{O}$-模层的商。
\end{enumerate}
\end{lemma}

\begin{proof}
证明 (1)。对 $\mathcal{C}$ 的每个对象 $U$ 以及每个
$s \in \mathcal{F}(U)$，得到态射 $j_{U!}\mathcal{O}_U \to \mathcal{F}$，
即态射 $\mathcal{O}_U \to \mathcal{F}|_U$，$1 \mapsto s$ 的伴随态射。
显然，映射
$$
\bigoplus\nolimits_{(U, s)} j_{U!}\mathcal{O}_U
\longrightarrow
\mathcal{F}
$$
是满射。结合引理 \ref{sites-modules-lemma-colimits-flat} 和
\ref{sites-modules-lemma-j-shriek-flat} 可知其源是平坦的，这就证明了 (2)。层的情形可通过
层化得出，也可重复同一论证。
\end{proof}

\begin{lemma}
\label{sites-modules-lemma-flat-tor-zero}
设 $\mathcal{C}$ 是范畴，$\mathcal{O}$ 是环预层。设
$$
0 \to \mathcal{F}'' \to \mathcal{F}' \to \mathcal{F} \to 0
$$
是 $\mathcal{O}$-模预层的短正合列，且 $\mathcal{G}$ 是
$\mathcal{O}$-模预层。
\begin{enumerate}
\item 如果 $\mathcal{F}$ 是平坦模预层，则序列
$$
0 \to
\mathcal{F}'' \otimes_{p, \mathcal{O}} \mathcal{G} \to
\mathcal{F}' \otimes_{p, \mathcal{O}} \mathcal{G} \to
\mathcal{F} \otimes_{p, \mathcal{O}} \mathcal{G} \to 0
$$
正合。
\item 如果 $\mathcal{C}$ 是站点，$\mathcal{O}$、$\mathcal{F}$、
$\mathcal{F}'$、$\mathcal{F}''$ 和 $\mathcal{G}$ 都是层，并且
$\mathcal{F}$ 作为模层是平坦的，则序列
$$
0 \to
\mathcal{F}'' \otimes_\mathcal{O} \mathcal{G} \to
\mathcal{F}' \otimes_\mathcal{O} \mathcal{G} \to
\mathcal{F} \otimes_\mathcal{O} \mathcal{G} \to 0
$$
正合。
\end{enumerate}
\end{lemma}

\begin{proof}
选择满射到 $\mathcal{G}$ 的平坦 $\mathcal{O}$-模预层 $\mathcal{G}'$；
由引理 \ref{sites-modules-lemma-module-quotient-flat}，这是可能的。令
$\mathcal{G}'' = \Ker(\mathcal{G}' \to \mathcal{G})$。对下图应用蛇引理即可：
$$
\begin{matrix}
 & & 0 & & 0 & & 0 & & \\
 & & \uparrow & & \uparrow & & \uparrow & & \\
 & & \mathcal{F}'' \otimes_{p, \mathcal{O}} \mathcal{G} & \to &
     \mathcal{F}' \otimes_{p, \mathcal{O}} \mathcal{G} & \to &
     \mathcal{F} \otimes_{p, \mathcal{O}} \mathcal{G} & \to & 0 \\
 & & \uparrow & & \uparrow & & \uparrow & & \\
0 & \to & \mathcal{F}'' \otimes_{p, \mathcal{O}} \mathcal{G}' & \to &
          \mathcal{F}' \otimes_{p, \mathcal{O}} \mathcal{G}' & \to &
          \mathcal{F} \otimes_{p, \mathcal{O}} \mathcal{G}' & \to & 0 \\
 & & \uparrow & & \uparrow & & \uparrow & & \\
 & & \mathcal{F}'' \otimes_{p, \mathcal{O}} \mathcal{G}'' & \to &
     \mathcal{F}' \otimes_{p, \mathcal{O}} \mathcal{G}'' & \to &
     \mathcal{F} \otimes_{p, \mathcal{O}} \mathcal{G}'' & \to & 0 \\
 & & & & & & \uparrow & & \\
 & & & & & & 0 & &
\end{matrix}
$$
其中各行各列均正合。中间一行正合，是因为与平坦模 $\mathcal{G}'$
作张量积是正合的。层的情形证明完全相同。
\end{proof}

\begin{lemma}
\label{sites-modules-lemma-flat-ses}
设 $\mathcal{C}$ 是范畴，$\mathcal{O}$ 是环预层。设
$$
0 \to
\mathcal{F}_2 \to
\mathcal{F}_1 \to
\mathcal{F}_0 \to 0
$$
是 $\mathcal{O}$-模预层的短正合列。
\begin{enumerate}
\item 如果 $\mathcal{F}_2$ 和 $\mathcal{F}_0$ 平坦，则
$\mathcal{F}_1$ 也平坦。
\item 如果 $\mathcal{F}_1$ 和 $\mathcal{F}_0$ 平坦，则
$\mathcal{F}_2$ 也平坦。
\end{enumerate}
如果 $\mathcal{C}$ 是站点且 $\mathcal{O}$ 是环层，则同一结论在
$\textit{Mod}(\mathcal{O})$ 中成立。
\end{lemma}

\begin{proof}
设 $\mathcal{G}^\bullet$ 是 $\mathcal{O}$-模预层的任意正合复形。假设
$\mathcal{F}_0$ 平坦。由引理 \ref{sites-modules-lemma-flat-tor-zero} 可知
$$
0 \to
\mathcal{G}^\bullet \otimes_{p, \mathcal{O}} \mathcal{F}_2 \to
\mathcal{G}^\bullet \otimes_{p, \mathcal{O}} \mathcal{F}_1 \to
\mathcal{G}^\bullet \otimes_{p, \mathcal{O}} \mathcal{F}_0 \to 0
$$
是 $\mathcal{O}$-模预层复形的短正合列。因此 (1) 和 (2) 由蛇引理得出。
模层的情形以相同方式证明。
\end{proof}

\begin{lemma}
\label{sites-modules-lemma-flat-resolution-of-flat}
设 $\mathcal{C}$ 是范畴，$\mathcal{O}$ 是环预层。设
$$
\ldots \to
\mathcal{F}_2 \to
\mathcal{F}_1 \to
\mathcal{F}_0 \to
\mathcal{Q} \to 0
$$
是 $\mathcal{O}$-模预层的正合复形。如果 $\mathcal{Q}$ 和所有
$\mathcal{F}_i$ 都是平坦 $\mathcal{O}$-模，则对任意
$\mathcal{O}$-模预层 $\mathcal{G}$，复形
$$
\ldots \to
\mathcal{F}_2 \otimes_{p, \mathcal{O}} \mathcal{G} \to
\mathcal{F}_1 \otimes_{p, \mathcal{O}} \mathcal{G} \to
\mathcal{F}_0 \otimes_{p, \mathcal{O}} \mathcal{G} \to
\mathcal{Q} \otimes_{p, \mathcal{O}} \mathcal{G} \to 0
$$
也正合。如果 $\mathcal{C}$ 是站点且 $\mathcal{O}$ 是环层，则同一结论在
$\textit{Mod}(\mathcal{O})$ 中成立。
\end{lemma}

\begin{proof}
本结论由引理 \ref{sites-modules-lemma-flat-tor-zero} 得出：将复形拆成短正合列，并利用引理
\ref{sites-modules-lemma-flat-ses} 归纳证明
$\Im(\mathcal{F}_{i + 1} \to \mathcal{F}_i)$ 平坦。
\end{proof}

\begin{lemma}
\label{sites-modules-lemma-tensor-flats}
设 $(\mathcal{C}, \mathcal{O})$ 是环站点。如果 $\mathcal{G}$ 和
$\mathcal{F}$ 是平坦 $\mathcal{O}$-模，则
$\mathcal{G} \otimes_\mathcal{O} \mathcal{F}$ 是平坦 $\mathcal{O}$-模。
\end{lemma}

\begin{proof}
这是因为
$$
(\mathcal{G} \otimes_\mathcal{O} \mathcal{F}) \otimes_\mathcal{O} \mathcal{H}
=
\mathcal{G} \otimes_\mathcal{O} (\mathcal{F} \otimes_\mathcal{O} \mathcal{H})
$$
，而正合函子的复合仍然正合。
\end{proof}

\begin{lemma}
\label{sites-modules-lemma-flat-change-of-rings}
设 $\mathcal{O}_1 \to \mathcal{O}_2$ 是站点 $\mathcal{C}$ 上的环层映射。
如果 $\mathcal{G}$ 是平坦 $\mathcal{O}_1$-模，则
$\mathcal{G} \otimes_{\mathcal{O}_1} \mathcal{O}_2$
是平坦 $\mathcal{O}_2$-模。
\end{lemma}

\begin{proof}
这是因为
$$
(\mathcal{G} \otimes_{\mathcal{O}_1} \mathcal{O}_2)
\otimes_{\mathcal{O}_2} \mathcal{H}
=
\mathcal{G} \otimes_{\mathcal{O}_1} \mathcal{F}
$$
（例如作为阿贝尔群层）。
\end{proof}

\noindent
下述引理是平坦性的方程判别准则（《代数》，引理
\ref{algebra-lemma-flat-eq}）的类似版本。

\begin{lemma}
\label{sites-modules-lemma-flat-eq}
设 $(\mathcal{C}, \mathcal{O})$ 是环站点，且 $\mathcal{F}$ 是
$\mathcal{O}$-模。下列条件等价：
\begin{enumerate}
\item $\mathcal{F}$ 是平坦 $\mathcal{O}$-模。
\item 设 $U$ 是 $\mathcal{C}$ 的对象，并设
$$
\mathcal{O}_U \xrightarrow{(f_1, \ldots, f_n)}
\mathcal{O}_U^{\oplus n} \xrightarrow{(s_1, \ldots, s_n)}
\mathcal{F}|_U
$$
是 $\mathcal{O}_U$-模复形。则存在覆盖 $\{U_i \to U\}$，且对每个 $i$
存在分解
$$
\mathcal{O}_{U_i}^{\oplus n}
\xrightarrow{B_i}
\mathcal{O}_{U_i}^{\oplus l_i} \xrightarrow{(t_{i1}, \ldots, t_{il_i})}
\mathcal{F}|_{U_i}
$$
，它分解 $(s_1, \ldots, s_n)|_{U_i}$，并且
$B_i \circ (f_1, \ldots, f_n)|_{U_i} = 0$。
\item 设 $U$ 是 $\mathcal{C}$ 的对象，并设
$$
\mathcal{O}_U^{\oplus m} \xrightarrow{A}
\mathcal{O}_U^{\oplus n} \xrightarrow{(s_1, \ldots, s_n)}
\mathcal{F}|_U
$$
是 $\mathcal{O}_U$-模复形。则存在覆盖 $\{U_i \to U\}$，且对每个 $i$
存在分解
$$
\mathcal{O}_{U_i}^{\oplus n}
\xrightarrow{B_i}
\mathcal{O}_{U_i}^{\oplus l_i} \xrightarrow{(t_{i1}, \ldots, t_{il_i})}
\mathcal{F}|_{U_i}
$$
，它分解 $(s_1, \ldots, s_n)|_{U_i}$，并且
$B_i \circ A|_{U_i} = 0$。
\end{enumerate}
\end{lemma}

\begin{proof}
假设 (1)。设 $\mathcal{I} \subset \mathcal{O}_U$ 是由
$f_1, \ldots, f_n$ 生成的理想层。则 $\sum f_j \otimes s_j$ 是
$\mathcal{I} \otimes_{\mathcal{O}_U} \mathcal{F}|_U$ 的截面，并在
$\mathcal{F}|_U$ 中映为零。因为 $\mathcal{F}|_U$ 平坦（引理
\ref{sites-modules-lemma-restriction-flat}），映射
$\mathcal{I} \otimes_{\mathcal{O}_U} \mathcal{F}|_U \to \mathcal{F}|_U$
是单射。由于 $\mathcal{I} \otimes_{\mathcal{O}_U} \mathcal{F}|_U$
是预层张量积的关联层，故存在覆盖 $\{U_i \to U\}$，使得
$\sum f_j|_{U_i} \otimes s_j|_{U_i}$ 在
$\mathcal{I}(U_i) \otimes_{\mathcal{O}(U_i)} \mathcal{F}(U_i)$.
中为零。利用《代数》，引理 \ref{algebra-lemma-relations} 展开定义，可找到
$t_{i1}, \ldots, t_{i l_i} \in \mathcal{F}(U_i)$ 和
$a_{ijk} \in \mathcal{O}(U_i)$
，使得 $\sum_j a_{ijk}f_j|_{U_i} = 0$ 且
$s_j|_{U_i} = \sum_k a_{ijk}t_{ik}$。故 (2) 成立。

\medskip\noindent
假设 (2)。给定 (3) 中的 $U$、$n$、$m$、$A$ 以及
$s_1, \ldots, s_n$。注意 $A$ 有 $m$ 列。对 $m$ 归纳证明 (3) 的断言。
在初始情形 $m = 0$，可取覆盖 $\{U \to U\}$，并使用分解
$\mathcal{O}_U^n \to \mathcal{O}_U^n \to \mathcal{F}|_U$，其中第一个映射是
恒等映射，第二个是 $(s_1, \ldots, s_n)$。假设 $m > 0$。令
$(f_1, \ldots, f_n)$ 为 $A$ 的最后一列，并应用 (2)。这给出新图
$$
\mathcal{O}_{U_i}^{\oplus m} \xrightarrow{B_i \circ A|_{U_i}}
\mathcal{O}_{U_i}^{\oplus l_i} \xrightarrow{(t_{i1}, \ldots, t_{il_i})}
\mathcal{F}|_{U_i}
$$
，但 $A_i = B_i \circ A|_{U_i}$ 的最后一列为零。因此可将归纳假设用于
$U_i$、$l_i$、$m - 1$、由 $A_i$ 前 $m - 1$ 列组成的矩阵以及
$t_{i1}, \ldots, t_{il_i}$，得到覆盖 $\{U_{ij} \to U_j\}$ 和分解
$$
\mathcal{O}_{U_{ij}}^{\oplus l_i}
\xrightarrow{C_{ij}}
\mathcal{O}_{U_{ij}}^{\oplus k_{ij}}
\xrightarrow{(v_{ij1}, \ldots, v_{ij k_{ij}})}
\mathcal{F}|_{U_{ij}}
$$
，它分解 $(t_{i1}, \ldots, t_{il_i})|_{U_{ij}}$，并且
$C_i \circ B_i|_{U_{ij}} \circ A|_{U_{ij}} = 0$。于是
$\{U_{ij} \to U\}$ 是覆盖，使用
$B_{ij} = C_i \circ B_i|_{U_{ij}}$ 和 $v_{ija}$ 即得到所需分解。
由此可见 (2) 蕴含 (3)。

\medskip\noindent
假设 (3)。设 $\mathcal{G} \to \mathcal{H}$ 是 $\mathcal{O}$-模的单同态。
需要证明
$\mathcal{G} \otimes_\mathcal{O} \mathcal{F} \to
\mathcal{H} \otimes_\mathcal{O} \mathcal{F}$
是单射。设 $U$ 是 $\mathcal{C}$ 的对象，并设
$s \in (\mathcal{G} \otimes_\mathcal{O} \mathcal{F})(U)$ 是在
$\mathcal{H} \otimes_\mathcal{O} \mathcal{F}$ 中映为零的截面。需要证明
$s$ 为零。由于
$\mathcal{G} \otimes_\mathcal{O} \mathcal{F}$
是层，只需找到 $\mathcal{C}$ 的覆盖 $\{U_i \to U\}_{i \in I}$，使得对所有
$i \in I$，$s|_{U_i}$ 为零。因此总可用某个覆盖的成员替换 $U$。特别地，
由于 $\mathcal{G} \otimes_\mathcal{O} \mathcal{F}$ 是
$\mathcal{G} \otimes_{p, \mathcal{O}} \mathcal{F}$ 的层化，可以假设 $s$ 是某个
$s' \in 
\mathcal{G}(U) \otimes_{\mathcal{O}(U)} \mathcal{F}(U)$.
的像。对 $\mathcal{H} \otimes_\mathcal{O} \mathcal{F}$ 作类似论证，可以假设
$s'$ 在 $\mathcal{H}(U) \otimes_{\mathcal{O}(U)} \mathcal{F}(U)$ 中映为零。
将 $\mathcal{F}(U) = \colim M_\alpha$ 写成有限表示
$\mathcal{O}(U)$-模 $M_\alpha$ 的滤过余极限（《代数》，引理
\ref{algebra-lemma-module-colimit-fp}）。由于张量积与滤过余极限交换
（《代数》，引理 \ref{algebra-lemma-tensor-products-commute-with-limits}），
可以选择 $\alpha$，使得 $s'$ 来自某个
$s'' \in \mathcal{G}(U) \otimes_{\mathcal{O}(U)} M_\alpha$，且 $s''$ 在
$\mathcal{H}(U) \otimes_{\mathcal{O}(U)} M_\alpha$ 中映为零。
固定 $\alpha$ 和 $s''$，并选择一个表示
$$
\mathcal{O}(U)^{\oplus m} \xrightarrow{A} \mathcal{O}(U)^{\oplus n}
\to M_\alpha \to 0
$$
。对相应的 $\mathcal{O}_U$-模复形应用 (3)：
$$
\mathcal{O}_U^{\oplus m} \xrightarrow{A}
\mathcal{O}_U^{\oplus n} \xrightarrow{(s_1, \ldots, s_n)}
\mathcal{F}|_U
$$
。用覆盖 $U_i$ 的成员替换 $U$ 后，可知映射
$$
M_\alpha \to \mathcal{F}(U)
$$
分解通过某个自由模 $\mathcal{O}(U)^{\oplus l}$。由于
$\mathcal{G}(U) \to \mathcal{H}(U)$ 是单射，可知
$$
\mathcal{G}(U) \otimes_{\mathcal{O}(U)} \mathcal{O}(U)^{\oplus l}
\to
\mathcal{H}(U) \otimes_{\mathcal{O}(U)} \mathcal{O}(U)^{\oplus l}
$$
也是单射。因此既然 $s''$ 在右侧模中映为零，它在左侧模中也映为零，
即 $s$ 如所需为零。
\end{proof}

\begin{lemma}
\label{sites-modules-lemma-flat-over-thickening}
设 $\mathcal{C}$ 是站点。设 $\mathcal{O}' \to \mathcal{O}$ 是环层的满射，
其核 $\mathcal{I}$ 是平方为零的理想。设 $\mathcal{F}'$ 是
$\mathcal{O}'$-模，并置
$\mathcal{F} = \mathcal{F}'/\mathcal{I}\mathcal{F}'$。下列条件等价：
\begin{enumerate}
\item $\mathcal{F}'$ 是平坦 $\mathcal{O}'$-模；
\item $\mathcal{F}$ 是平坦 $\mathcal{O}$-模，并且
$\mathcal{I} \otimes_\mathcal{O} \mathcal{F} \to \mathcal{F}'$
是单射。
\end{enumerate}
\end{lemma}

\begin{proof}
如果 (1) 成立，则由引理 \ref{sites-modules-lemma-flat-change-of-rings}，
$\mathcal{F} = \mathcal{F}' \otimes_{\mathcal{O}'} \mathcal{O}$ 在
$\mathcal{O}$ 上平坦。对正合列
$0 \to \mathcal{I} \to \mathcal{O}' \to \mathcal{O} \to 0$
应用 $- \otimes_{\mathcal{O}'} \mathcal{F}'$，可知映射
$\mathcal{I} \otimes_\mathcal{O} \mathcal{F} \to \mathcal{F}'$
是单射；见引理 \ref{sites-modules-lemma-flat-tor-zero}。假设 (2)。在以下证明中将不再说明地使用：
对任意被 $\mathcal{I}$ 零化的 $\mathcal{O}'$-模 $\mathcal{K}$，有
$\mathcal{K} \otimes_{\mathcal{O}'} \mathcal{F}' =
\mathcal{K} \otimes_\mathcal{O} \mathcal{F}$
。设 $\alpha : \mathcal{G}' \to \mathcal{H}'$ 是
$\mathcal{O}'$-模的单射。分别令 $\mathcal{G} \subset \mathcal{G}'$、
$\mathcal{H} \subset \mathcal{H}'$ 为由被 $\mathcal{I}$ 零化的截面组成的子层。
考虑交换图
$$
\xymatrix{
\mathcal{G} \otimes_{\mathcal{O}'} \mathcal{F}' \ar[r] \ar[d] &
\mathcal{G}' \otimes_{\mathcal{O}'} \mathcal{F}' \ar[r] \ar[d] &
\mathcal{G}'/\mathcal{G} \otimes_{\mathcal{O}'} \mathcal{F}' \ar[r] \ar[d] &
0 \\
\mathcal{H} \otimes_{\mathcal{O}'} \mathcal{F}' \ar[r] &
\mathcal{H}' \otimes_{\mathcal{O}'} \mathcal{F}' \ar[r] &
\mathcal{H}'/\mathcal{H} \otimes_{\mathcal{O}'} \mathcal{F}' \ar[r] & 0
}
$$
注意，$\mathcal{G}'/\mathcal{G}$ 和 $\mathcal{H}'/\mathcal{H}$
都被 $\mathcal{I}$ 零化，且
$\mathcal{G}'/\mathcal{G} \to \mathcal{H}'/\mathcal{H}$ 是单射。
由于 $\mathcal{F}$ 在 $\mathcal{O}$ 上平坦，右侧竖直箭头是单射；左侧竖直箭头
也如此。因此中间的竖直箭头是单射，故 $\mathcal{F}'$ 平坦。
\end{proof}

\begin{lemma}
\label{sites-modules-lemma-neighbourhood-isomorphism}
设 $\mathcal{C}$ 是站点。设 $\mathcal{O} \to \mathcal{O}'$
是环层的平坦同态。设
$\mathcal{I} \subset \mathcal{O}$
是理想层，使得诱导映射
$\mathcal{O}/\mathcal{I} \to \mathcal{O}'/\mathcal{I}\mathcal{O}'$
是同构。对任意被某个 $\mathcal{I}^n$（$n \geq 0$）零化的
$\mathcal{O}$-模 $\mathcal{F}$，映射
$\text{id} \otimes 1 :
\mathcal{F} \to \mathcal{F} \otimes_\mathcal{O} \mathcal{O}'$
是同构。
\end{lemma}

\begin{proof}
略。提示：见《代数详论》，引理
\ref{more-algebra-lemma-neighbourhood-isomorphism}。
\end{proof}





\section{对偶}
\label{sites-modules-section-duals}

\noindent
设 $(\mathcal{C}, \mathcal{O})$ 是环站点。赋予第
\ref{sites-modules-section-tensor-product} 节构造的张量积后，$\mathcal{O}$-模范畴成为
对称幺半范畴。对 $\mathcal{O}$-模 $\mathcal{F}$，下列条件等价：
\begin{enumerate}
\item $\mathcal{F}$ 在 $\mathcal{O}$-模幺半范畴中有左对偶；
\item 对 $\mathcal{C}$ 的每个对象 $U$，存在覆盖 $\{U_i \to U\}$，使得
$\mathcal{F}|_{U_i}$ 是某个有限自由 $\mathcal{O}|_{U_i}$-模的直和项；
\item $\mathcal{F}$ 是有限表示且平坦的 $\mathcal{O}$-模。
\end{enumerate}
这将在本节的例 \ref{sites-modules-example-dual} 以及引理
\ref{sites-modules-lemma-left-dual-module} 和
\ref{sites-modules-lemma-flat-locally-finite-presentation} 中证明。

\begin{example}
\label{sites-modules-example-dual}
设 $(\mathcal{C}, \mathcal{O})$ 是环站点。设 $\mathcal{F}$ 是
$\mathcal{O}$-模，并且对 $\mathcal{C}$ 的每个对象 $U$，都存在覆盖
$\{U_i \to U\}$，使得 $\mathcal{F}|_{U_i}$ 是某个有限自由
$\mathcal{O}|_{U_i}$-模的直和项。则映射
$$
\mathcal{F} \otimes_\mathcal{O}
\SheafHom_\mathcal{O}(\mathcal{F}, \mathcal{O})
\longrightarrow
\SheafHom_\mathcal{O}(\mathcal{F}, \mathcal{F})
$$
是同构。事实上，这是局部问题；当 $\mathcal{F}$ 有限自由时成立，并且若对某个模
成立，则对它的任一直和项也成立（略去细节）。记
$$
\eta :
\mathcal{O}
\longrightarrow
\mathcal{F} \otimes_\mathcal{O}
\SheafHom_\mathcal{O}(\mathcal{F}, \mathcal{O})
$$
为把 $1$ 送到上述同构下对应于 $\text{id}_\mathcal{F}$ 的截面的映射。记
$$
\epsilon : 
\SheafHom_\mathcal{O}(\mathcal{F}, \mathcal{O})
\otimes_\mathcal{O} \mathcal{F}
\longrightarrow
\mathcal{O}
$$
为求值映射。于是可见
$\SheafHom_\mathcal{O}(\mathcal{F}, \mathcal{O}), \eta, \epsilon$
是《范畴》，定义 \ref{categories-definition-dual} 意义下
$\mathcal{F}$ 的左对偶。略去下列等式的验证：
$(1 \otimes \epsilon) \circ (\eta \otimes 1) = \text{id}_\mathcal{F}$
以及
$(\epsilon \otimes 1) \circ (1 \otimes \eta) =
\text{id}_{\SheafHom_\mathcal{O}(\mathcal{F}, \mathcal{O})}$.
\end{example}

\begin{lemma}
\label{sites-modules-lemma-left-dual-module}
设 $(\mathcal{C}, \mathcal{O})$ 是环站点，且 $\mathcal{F}$ 是
$\mathcal{O}$-模。设 $\mathcal{G}, \eta, \epsilon$ 是
$\mathcal{F}$ 在 $\mathcal{O}$-模幺半范畴中的左对偶；见《范畴》，定义
\ref{categories-definition-dual}。则
\begin{enumerate}
\item 对 $\mathcal{C}$ 的每个对象 $U$，存在覆盖 $\{U_i \to U\}$，使得
$\mathcal{F}|_{U_i}$ 是某个有限自由 $\mathcal{O}|_{U_i}$-模的直和项；
\item 映射
$e : \SheafHom_\mathcal{O}(\mathcal{F}, \mathcal{O}) \to \mathcal{G}$
把局部截面 $\lambda$ 送到 $(\lambda \otimes 1)(\eta)$，它是同构；
\item 对 $\mathcal{F}$ 和 $\mathcal{G}$ 的局部截面 $f$ 和 $g$，有
$\epsilon(f, g) = e^{-1}(g)(f)$。
\end{enumerate}
\end{lemma}

\begin{proof}
这些假设意味着
$$
\mathcal{F} \xrightarrow{\eta \otimes 1}
\mathcal{F} \otimes_\mathcal{O} \mathcal{G}
\otimes_\mathcal{O} \mathcal{F}
\xrightarrow{1 \otimes \epsilon} \mathcal{F}
\quad\text{且}\quad
\mathcal{G} \xrightarrow{1 \otimes \eta}
\mathcal{G} \otimes_\mathcal{O} \mathcal{F}
\otimes_\mathcal{O} \mathcal{G}
\xrightarrow{\epsilon \otimes 1} \mathcal{G}
$$
都是恒等映射。设 $U$ 是 $\mathcal{C}$ 的对象。用 $U$ 的某个覆盖的成员替换
$U$ 后，可以找到 $\mathcal{F}$ 和 $\mathcal{G}$ 在 $U$ 上的有限多个截面
$f_1, \ldots, f_n$ 和 $g_1, \ldots, g_n$，使得
$\eta(1) = \sum f_i g_i$。记
$$
\mathcal{O}_U^{\oplus n} \to \mathcal{F}|_U
$$
为把第 $i$ 个基向量送到 $f_i$ 的映射。于是映射 $\eta|_U$ 分解通过映射
$\tilde \eta : \mathcal{O}_U \to
\mathcal{O}_U^{\oplus n} \otimes_{\mathcal{O}_U} \mathcal{G}|_U$。
得到交换图
$$
\xymatrix{
\mathcal{F}|_U
\ar[rr]_-{\eta \otimes 1} \ar[rrd]_-{\tilde \eta \otimes 1} & &
\mathcal{F}|_U \otimes \mathcal{G}|_U \otimes \mathcal{F}|_U
\ar[r]_-{1 \otimes \epsilon} &
\mathcal{F}|_U \\
& &
\mathcal{O}_U^{\oplus n} \otimes \mathcal{G}|_U \otimes \mathcal{F}|_U
\ar[u] \ar[r]^-{1 \otimes \epsilon} &
\mathcal{O}_U^{\oplus n} \ar[u]
}
$$
这表明 $\mathcal{F}|_U$ 上的恒等映射分解通过有限自由
$\mathcal{O}_U$-模，从而证明了 (1)。(2) 由《范畴》，引理
\ref{categories-lemma-left-dual} 及其证明得到。(3) 由本证明中的第一个等式得到。
也可以从左对偶的唯一性（《范畴》，注
\ref{categories-remark-left-dual-adjoint}）以及例
\ref{sites-modules-example-dual} 中左对偶的构造推出 (2) 和 (3)。
\end{proof}

\begin{lemma}
\label{sites-modules-lemma-flat-locally-finite-presentation}
设 $(\mathcal{C}, \mathcal{O})$ 是环站点。设 $\mathcal{F}$ 局部有限表示且
平坦。则给定 $\mathcal{C}$ 的对象 $U$，存在覆盖 $\{U_i \to U\}$，使得
$\mathcal{F}|_{U_i}$ 是某个有限自由 $\mathcal{O}_{U_i}$-模的直和项。
\end{lemma}

\begin{proof}
选择 $\mathcal{C}$ 的对象 $U$。用某个覆盖的成员替换 $U$ 后，可以假设存在表示
$$
\mathcal{O}_U^{\oplus r} \to
\mathcal{O}_U^{\oplus n} \to \mathcal{F}|_U \to 0
$$
。由引理 \ref{sites-modules-lemma-flat-eq}，再次用某个覆盖的成员替换 $U$ 后，可以假设存在分解
$$
\mathcal{O}_U^{\oplus n} \to
\mathcal{O}_U^{\oplus n_1} \to \mathcal{F}|_U
$$
，使得复合
$\mathcal{O}_U^{\oplus r} \to \mathcal{O}_U^{\oplus n}
\to \mathcal{O}_U^{\oplus n_r}$ 为零。这意味着满射
$\mathcal{O}_U^{\oplus n_r} \to \mathcal{F}|_U$ 有截面，结论成立。
\end{proof}











\section{走向可构造模层}
\label{sites-modules-section-constructible}

\noindent
回顾一下，站点的拟紧对象粗略说来是其每个覆盖都可由有限覆盖加细的对象
（实际定义略为复杂；见《站点》，第 \ref{sites-section-quasi-compact} 节）。
事实表明，如果站点的每个对象都有由拟紧对象组成的覆盖，那么当 $U$ 拟紧时，
模层 $j_!\mathcal{O}_U$ 构成全体模层范畴的一组特别良好的生成元。

\begin{lemma}
\label{sites-modules-lemma-covering-gives-surjection}
设 $(\mathcal{C}, \mathcal{O})$ 是环站点，且 $\{U_i \to U\}$ 是
$\mathcal{C}$ 的覆盖。则序列
$$
\bigoplus j_{U_i \times_U U_j!}\mathcal{O}_{U_i \times_U U_j} \to
\bigoplus j_{U_i!}\mathcal{O}_{U_i} \to j_!\mathcal{O}_U \to 0
$$
正合。
\end{lemma}

\begin{proof}
对任意 $\mathcal{O}$-模 $\mathcal{F}$，函子
$\Hom_\mathcal{O}(-, \mathcal{F})$ 将上述序列变为正合列
$0 \to \mathcal{F}(U) \to \prod \mathcal{F}(U_i) \to
\prod \mathcal{F}(U_i \times_U U_j)$；见
(\ref{sites-modules-equation-map-lower-shriek-OU-into-module})。再由《同调代数》，引理
\ref{homology-lemma-check-exactness} 即得本引理。
\end{proof}

\begin{lemma}
\label{sites-modules-lemma-silly-quasi-compact}
设 $(\mathcal{C}, \mathcal{O})$ 是环站点。设
$\mathcal{U} = \{U_i \to U\}_{i \in I}$ 是 $\mathcal{C}$ 的覆盖。
如果 $U$ 拟紧，则存在有限子集 $I' \subset I$，使得序列
$$
\bigoplus\nolimits_{i, i' \in I'}
j_{U_i \times_U U_{i'}!}\mathcal{O}_{U_i \times_U U_{i'}} \to
\bigoplus\nolimits_{i \in I'}
j_{U_i!}\mathcal{O}_{U_i} \to
j_!\mathcal{O}_U \to 0
$$
正合。
\end{lemma}

\begin{proof}
本引理由引理 \ref{sites-modules-lemma-covering-gives-surjection} 立即得到，只要 $U$ 满足
《站点》，引理 \ref{sites-lemma-conclude-quasi-compact} 的条件 (3)。
我们建议读者跳过一般情形的证明。由定义，存在覆盖
$\mathcal{V} = \{V_j \to U\}_{j \in J}$ 以及具有固定靶的映射族的态射
$\mathcal{V} \to \mathcal{U}$，它由 $\text{id} : U \to U$、
$\alpha : J \to I$ 和
$f_j : V_j \to U_{\alpha(j)}$
给出（见《站点》，定义 \ref{sites-definition-morphism-coverings}），并且
$\alpha$ 的像 $I' \subset I$ 有限。由《同调代数》，引理
\ref{homology-lemma-check-exactness}，只需证明对任意 $\mathcal{O}$-模层
$\mathcal{F}$，函子 $\Hom_\mathcal{O}(-, \mathcal{F})$ 将本引理的序列变为
正合列。由
(\ref{sites-modules-equation-map-lower-shriek-OU-into-module})
得到通常的序列
$$
0 \to
\mathcal{F}(U) \to
\prod\nolimits_{i \in I'} \mathcal{F}(U_i) \to
\prod\nolimits_{i, i' \in I'} \mathcal{F}(U_i \times_U U_{i'})
$$
。对由覆盖 $\mathcal{V}$ 加细的映射族
$\{U_i \to U\}_{i \in I'}$ 应用《站点》，引理
\ref{sites-lemma-compare-sheaf-condition}，可知这是正合列。
\end{proof}

\begin{lemma}
\label{sites-modules-lemma-sections-over-quasi-compact}
设 $\mathcal{C}$ 是站点，且 $W$ 是 $\mathcal{C}$ 的拟紧对象。
\begin{enumerate}
\item 函子 $\Sh(\mathcal{C}) \to \textit{Sets}$，
$\mathcal{F} \mapsto \mathcal{F}(W)$ 与余积交换。
\item 设 $\mathcal{O}$ 是 $\mathcal{C}$ 上的环层。函子
$\textit{Mod}(\mathcal{O}) \to \textit{Ab}$,
$\mathcal{F} \mapsto \mathcal{F}(W)$
与直和交换。
\end{enumerate}
\end{lemma}

\begin{proof}
证明 (1)。由《站点》，引理 \ref{sites-lemma-directed-colimits-sections}，
取 $W$ 上截面与转移映射为单射的滤过余极限交换。设 $\mathcal{F}_i$
是由集合 $I$ 标号的集合层族。则 $\coprod \mathcal{F}_i$ 是余积
$\mathcal{F}_E = \coprod_{i \in E} \mathcal{F}_i$ 在有限子集
$E \subset I$ 的偏序集上的滤过余极限。由于转移映射为单射，结论成立。

\medskip\noindent
证明 (2)。设 $\mathcal{F}_i$ 是由集合 $I$ 标号的 $\mathcal{O}$-模层族。
则 $\bigoplus \mathcal{F}_i$ 是直和
$\mathcal{F}_E = \bigoplus_{i \in E} \mathcal{F}_i$ 在有限子集
$E \subset I$ 的偏序集上的滤过余极限。阿贝尔层的滤过余极限可在集合层范畴中
计算。此外，对 $E \subset E'$，转移映射
$\mathcal{F}_E \to \mathcal{F}_{E'}$ 是单射（因为层化正合，且在底层预层上
单射性显然）。因此由《站点》，引理
\ref{sites-lemma-directed-colimits-sections}，只需证明有限指标集的情形。
有限情形已由引理 \ref{sites-modules-lemma-limits-colimits-abelian-sheaves} 处理
（它在 $\mathcal{C}$ 的任意对象上都成立）。
\end{proof}

\begin{lemma}
\label{sites-modules-lemma-quasi-compact-hom-from}
设 $(\mathcal{C}, \mathcal{O})$ 是环站点，且 $U$ 是 $\mathcal{C}$ 的拟紧对象。
则函子 $\Hom_\mathcal{O}(j_!\mathcal{O}_U, -)$ 与直和交换。
\end{lemma}

\begin{proof}
这是因为由 (\ref{sites-modules-equation-map-lower-shriek-OU-into-module})，有
$\Hom_\mathcal{O}(j_!\mathcal{O}_U, \mathcal{F}) = \mathcal{F}(U)$
，而由引理 \ref{sites-modules-lemma-sections-over-quasi-compact}，函子
$\mathcal{F} \mapsto \mathcal{F}(U)$ 与直和交换。
\end{proof}

\noindent
为陈述下文尽可能精确的结果，引入一些记号。

\begin{situation}
\label{sites-modules-situation-quasi-compact-objects}
设 $\mathcal{C}$ 是站点，并设
$\mathcal{B} \subset \text{Ob}(\mathcal{C})$ 是对象集。考虑下列条件：
\begin{enumerate}
\item
\label{sites-modules-item-enough}
$\mathcal{C}$ 的每个对象都有由 $\mathcal{B}$ 中元素组成的覆盖。
\item
\label{sites-modules-item-enough-qc}
每个 $U \in \mathcal{B}$ 都是拟紧的
（《站点》，第 \ref{sites-section-quasi-compact} 节）。
\item
\label{sites-modules-item-enough-qc-qs}
对满足 $U_i, U \in \mathcal{B}$ 的覆盖 $\{U_i \to U\}$，纤维积
$U_i \times_U U_j$ 是拟紧的。
\end{enumerate}
\end{situation}

\begin{lemma}
\label{sites-modules-lemma-module-quotient-direct-sum}
在情形 \ref{sites-modules-situation-quasi-compact-objects} 中，假设
(\ref{sites-modules-item-enough}) 成立。
\begin{enumerate}
\item 每个集合层都是某个满射的靶，该满射的源是余积
$\coprod h_{U_i}^\#$，其中 $U_i$ 属于 $\mathcal{B}$。
\item 如果 $\mathcal{O}$ 是环层，则每个 $\mathcal{O}$-模都是某个直和
$\bigoplus\nolimits j_{U_i!}\mathcal{O}_{U_i}$ 的商，其中 $U_i$
属于 $\mathcal{B}$。
\end{enumerate}
\end{lemma}

\begin{proof}
第 (1) 部分由《站点》，引理
\ref{sites-lemma-sheaf-coequalizer-representable} 和
\ref{sites-lemma-covering-surjective-after-sheafification} 得出。
第 (2) 部分由引理 \ref{sites-modules-lemma-module-quotient-flat} 和
\ref{sites-modules-lemma-covering-gives-surjection} 得出。
\end{proof}

\begin{lemma}
\label{sites-modules-lemma-module-filtered-colimit-constructibles}
在情形 \ref{sites-modules-situation-quasi-compact-objects} 中，假设
(\ref{sites-modules-item-enough}) 和 (\ref{sites-modules-item-enough-qc}) 成立。
\begin{enumerate}
\item 每个集合层都是下列形式的层的滤过余极限：
\begin{equation}
\label{sites-modules-equation-towards-constructible-sets}
\text{余等化子}\left(
\xymatrix{
\coprod\nolimits_{j = 1, \ldots, m} h_{V_j}^\#
\ar@<1ex>[r] \ar@<-1ex>[r] &
\coprod\nolimits_{i = 1, \ldots, n} h_{U_i}^\#
}
\right)
\end{equation}
，其中 $U_i$ 和 $V_j$ 属于 $\mathcal{B}$。
\item 如果 $\mathcal{O}$ 是环层，则每个 $\mathcal{O}$-模都是下列形式的层的
滤过余极限：
\begin{equation}
\label{sites-modules-equation-towards-constructible}
\Coker\left(
\bigoplus\nolimits_{j = 1, \ldots, m} j_{V_j!}\mathcal{O}_{V_j}
\longrightarrow
\bigoplus\nolimits_{i = 1, \ldots, n} j_{U_i!}\mathcal{O}_{U_i}
\right)
\end{equation}
，其中 $U_i$ 和 $V_j$ 属于 $\mathcal{B}$。
\end{enumerate}
\end{lemma}

\begin{proof}
证明 (1)。由引理 \ref{sites-modules-lemma-module-quotient-direct-sum}，每个集合层
$\mathcal{F}$ 都是某个满射的靶，其源 $\mathcal{F}_0$ 是形如
$h_U^\#$（$U \in \mathcal{B}$）的层的余积。把这一点应用于
$\mathcal{F}_0 \times_\mathcal{F} \mathcal{F}_0$，可知 $\mathcal{F}$ 是一对映射
$$
\xymatrix{
\coprod\nolimits_{j \in J} h_{V_j}^\#
\ar@<1ex>[r] \ar@<-1ex>[r] &
\coprod\nolimits_{i \in I} h_{U_i}^\#
}
$$
的余等化子，其中 $I$、$J$ 是某些指标集，且 $V_j$ 和 $U_i$ 属于
$\mathcal{B}$。对每个有限子集 $J' \subset J$，存在有限子集
$I' \subset I$，使得 $j \in J'$ 上的余积经两个映射都映入
$i \in I'$ 上的余积；见《站点》，引理
\ref{sites-lemma-directed-colimits-sections}。（略去细节；提示：无限余积是有限子余积的
滤过余极限。）因此我们的层是这些有限余积之间映射的余核的余极限。

\medskip\noindent
证明 (2)。由引理 \ref{sites-modules-lemma-module-quotient-direct-sum}，每个模都是形如
$j_{U!}\mathcal{O}_U$（$U \in \mathcal{B}$）的模的某个直和的商。因此每个模
都是一个余核
$$
\Coker\left(
\bigoplus\nolimits_{j \in J} j_{V_j!}\mathcal{O}_{V_j}
\longrightarrow
\bigoplus\nolimits_{i \in I} j_{U_i!}\mathcal{O}_{U_i}
\right)
$$
，其中 $I$、$J$ 是某些指标集，且 $V_j$ 和 $U_i$ 属于 $\mathcal{B}$。
对每个有限子集 $J' \subset J$，存在有限子集 $I' \subset I$，使得
$j \in J'$ 上的直和映入 $i \in I'$ 上的直和；见引理
\ref{sites-modules-lemma-quasi-compact-hom-from}。因此我们的模是这些有限直和之间映射的余核的
余极限。
\end{proof}

\begin{lemma}
\label{sites-modules-lemma-cokernel-map-towards-constructibles}
在情形 \ref{sites-modules-situation-quasi-compact-objects} 中，假设
(\ref{sites-modules-item-enough}) 和 (\ref{sites-modules-item-enough-qc}) 成立。设 $\mathcal{O}$ 是环层。
则形如 (\ref{sites-modules-equation-towards-constructible}) 的模之间映射的余核，仍是形如
(\ref{sites-modules-equation-towards-constructible}) 的模。
\end{lemma}

\begin{proof}
设 $\mathcal{F} = \Coker(\bigoplus j_{V_j!}\mathcal{O}_{V_j} \to
\bigoplus j_{U_i!}\mathcal{O}_{U_i})$
，如 (\ref{sites-modules-equation-towards-constructible}) 所示。只需证明，当
$W \in \mathcal{B}$ 时，映射
$\varphi : j_{W!}\mathcal{O}_W \to \mathcal{F}$ 的余核仍是同类型的模。
映射 $\varphi$ 对应于 $s \in \mathcal{F}(W)$。由于
$\bigoplus j_{U_i!}\mathcal{O}_{U_i} \to \mathcal{F}$ 是满射，由
(\ref{sites-modules-item-enough}) 可选择满足 $W_k \in \mathcal{B}$ 的覆盖
$\{W_k \to W\}_{k \in K}$，使得 $s|_{W_k}$ 是
$\bigoplus j_{U_i!}\mathcal{O}_{U_i})$ 的某个截面 $s_k$ 的像。
由 (\ref{sites-modules-item-enough-qc})，对象 $W$ 拟紧。由引理
\ref{sites-modules-lemma-silly-quasi-compact}，存在有限子集 $K' \subset K$，使得
$\bigoplus_{k \in K'} j_{W_k!}\mathcal{O}_{W_k} \to j_{W!}\mathcal{O}_W$
是满射。由此可知 $\Coker(\varphi)$ 等于
$$
\Coker\left(
\bigoplus\nolimits_{k \in K'} j_{W_k!}\mathcal{O}_{W_k} \oplus
\bigoplus j_{V_j!}\mathcal{O}_{V_j}
\longrightarrow
\bigoplus j_{U_i!}\mathcal{O}_{U_i}
\right)
$$
，其中映射 $\bigoplus_{k \in K'} j_{W_k!}\mathcal{O}_{W_k}
\to \bigoplus j_{U_i!}\mathcal{O}_{U_i}$ 对应于
$\sum_{k \in K'} s_k$。证明完毕。
\end{proof}

\begin{lemma}
\label{sites-modules-lemma-change-presentation-towards-constructibles}
在情形 \ref{sites-modules-situation-quasi-compact-objects} 中，假设
(\ref{sites-modules-item-enough})、(\ref{sites-modules-item-enough-qc}) 和
(\ref{sites-modules-item-enough-qc-qs}) 成立。设 $\mathcal{O}$ 是环层。假设给定映射
$$
\bigoplus\nolimits_{j = 1, \ldots, m} j_{V_j!}\mathcal{O}_{V_j}
\longrightarrow
\bigoplus\nolimits_{i = 1, \ldots, n} j_{U_i!}\mathcal{O}_{U_i}
$$
，其中 $U_i$ 和 $V_j$ 属于 $\mathcal{B}$；并给定满足
$U_{ik} \in \mathcal{B}$ 的覆盖
$\{U_{ik} \to U_i\}_{k \in K_i}$。则存在有限子集
$K'_i \subset K_i$、由 $W_l \in \mathcal{B}$ 组成的有限集 $L$，以及交换图
$$
\xymatrix{
\bigoplus_{l \in L} j_{W_l!}\mathcal{O}_{W_l} \ar[d] \ar[r] &
\bigoplus_{i = 1, \ldots, n} \bigoplus_{k \in K'_i}
j_{U_{ik}!}\mathcal{O}_{U_{ik}} \ar[d] \\
\bigoplus_{j = 1, \ldots, m} j_{V_j!}\mathcal{O}_{V_j} \ar[r] &
\bigoplus_{i = 1, \ldots, n} j_{U_i!}\mathcal{O}_{U_i}
}
$$
，它在横向映射的余核上诱导同构。
\end{lemma}

\begin{proof}
由于 $U_i$ 拟紧，可以按引理 \ref{sites-modules-lemma-silly-quasi-compact} 选择有限子集
$K'_i \subset K_i$。又因为
$\bigoplus_{i = 1, \ldots, n}
\bigoplus_{k \in K'_i} j_{U_{ik}!}\mathcal{O}_{U_{ik}} \to
\bigoplus_{i = 1, \ldots, n} j_{U_i!}\mathcal{O}_{U_i}$ 是满射，
可以找到满足 $V_{jm} \in \mathcal{B}$ 的覆盖
$\{V_{jm} \to V_j\}_{m \in M_j}$，使得存在交换图
$$
\xymatrix{
\bigoplus_{j = 1, \ldots, m} \bigoplus_{m \in M_j}
j_{V_{jm}!}\mathcal{O}_{V_{jm}} \ar[d] \ar[r] &
\bigoplus_{i = 1, \ldots n} \bigoplus_{k \in K'_i}
j_{U_{ik}!}\mathcal{O}_{U_{ik}} \ar[d] \\
\bigoplus_{j = 1, \ldots, m} j_{V_j!}\mathcal{O}_{V_j} \ar[r] &
\bigoplus_{i = 1, \ldots, n} j_{U_i!}\mathcal{O}_{U_i}
}
$$
。由于 $V_j$ 拟紧，可以按引理 \ref{sites-modules-lemma-silly-quasi-compact}
选择有限子集 $M'_j \subset M_j$。置
$$
L = \left(\coprod\nolimits_{i = 1, \ldots, n} K'_i \times K'_i \right)
\coprod
\left(\coprod\nolimits_{j = 1, \ldots, m} M'_j\right)
$$
；对 $l = (k, k') \in K'_i \times K'_i \subset L$，置
$W_l = U_{ik} \times_{U_i} U_{ik'}$；对
$l = m \in M'_j \subset L$，置 $W_l = V_{jm}$。对映射族
$\{U_{ik} \to U_i\}_{k \in K'_i}$，已有引理
\ref{sites-modules-lemma-silly-quasi-compact} 的正合列，因此得到本引理陈述中的图
（略去细节），但尚不清楚 $W_l \in \mathcal{B}$。不过，对所有
$l \in L$，$W_l$ 都拟紧，所以再次应用引理
\ref{sites-modules-lemma-silly-quasi-compact}，找到满足 $W_{lt} \in \mathcal{B}$ 的有限映射族
$\{W_{lt} \to W_l\}_{t \in T_l}$，使得
$\bigoplus j_{W_{lt}!}\mathcal{O}_{W_{lt}} \to j_{W_l!}\mathcal{O}_{W_l}$
是满射。然后用 $\coprod_{l \in L} T_l$ 替换 $L$，结论即清楚。
\end{proof}

\begin{lemma}
\label{sites-modules-lemma-extension-towards-constructibles}
在情形 \ref{sites-modules-situation-quasi-compact-objects} 中，假设
(\ref{sites-modules-item-enough})、(\ref{sites-modules-item-enough-qc}) 和
(\ref{sites-modules-item-enough-qc-qs}) 成立。设 $\mathcal{O}$ 是环层。
则形如 (\ref{sites-modules-equation-towards-constructible}) 的模的扩张仍是形如
(\ref{sites-modules-equation-towards-constructible}) 的模。
\end{lemma}

\begin{proof}
设 $0 \to \mathcal{F}_1 \to \mathcal{F}_2 \to \mathcal{F}_3 \to 0$
是 $\mathcal{O}$-模的短正合列，其中 $\mathcal{F}_1$ 和 $\mathcal{F}_3$
形如 (\ref{sites-modules-equation-towards-constructible})。选择表示
$$
\bigoplus A_{V_j} \to \bigoplus A_{U_i} \to \mathcal{F}_1 \to 0
\quad\text{且}\quad
\bigoplus A_{T_j} \to \bigoplus A_{W_i} \to \mathcal{F}_3 \to 0
$$
。在此证明中直和总是有限的，并且对 $U \in \mathcal{B}$，记
$A_U = j_{U!}\mathcal{O}_U$。由于 $\mathcal{F}_2 \to \mathcal{F}_3$
是满射，可选择满足 $W_{ik} \in \mathcal{B}$ 的覆盖
$\{W_{ik} \to W_i\}$，使得 $A_{W_{ik}} \to \mathcal{F}_3$ 可提升为映射
$A_{W_{ik}} \to \mathcal{F}_2$。由引理
\ref{sites-modules-lemma-change-presentation-towards-constructibles}，可以用集合
$\{W_{ik}\}$ 的某个有限子集替换集合 $\{W_i\}$，并假设映射
$\bigoplus A_{W_i} \to \mathcal{F}_3$ 可提升为到 $\mathcal{F}_2$ 的映射。
考虑核
$$
\mathcal{K}_2 = \Ker(\bigoplus A_{U_i} \oplus \bigoplus A_{W_i}
\longrightarrow \mathcal{F}_2)
$$
。由蛇引理，该核满射到
$\mathcal{K}_3 = \Ker(\bigoplus A_{W_i} \to \mathcal{F}_3)$.
。因此按上述论证，在用 $\mathcal{B}$ 中元素的有限族替换每个 $T_j$ 后
（由引理 \ref{sites-modules-lemma-silly-quasi-compact}，这是允许的），可以假设存在映射
$\bigoplus A_{T_j} \to \mathcal{K}_2$，提升给定映射
$\bigoplus A_{T_j} \to \mathcal{K}_3$。于是
$\bigoplus A_{V_j} \oplus \bigoplus A_{T_j} \to \mathcal{K}_2$
是满射，证明完毕。
\end{proof}

\begin{lemma}
\label{sites-modules-lemma-towards-constructible-when-serre-subcategory}
在情形 \ref{sites-modules-situation-quasi-compact-objects} 中，假设
(\ref{sites-modules-item-enough})、(\ref{sites-modules-item-enough-qc}) 和
(\ref{sites-modules-item-enough-qc-qs}) 成立。设 $\mathcal{O}$ 是环层，并设
$\mathcal{A} \subset \textit{Mod}(\mathcal{O})$ 是由同构于形如
(\ref{sites-modules-equation-towards-constructible}) 的余核的模组成的全子范畴。
如果每个形如
$$
\bigoplus\nolimits_{j = 1, \ldots, m} j_{V_j!}\mathcal{O}_{V_j}
\longrightarrow
\bigoplus\nolimits_{i = 1, \ldots, n} j_{U_i!}\mathcal{O}_{U_i}
$$
的 $\mathcal{O}$-模映射的核都属于 $\mathcal{A}$，其中 $U_i$ 和 $V_j$
属于 $\mathcal{B}$，则 $\mathcal{A}$ 是
$\textit{Mod}(\mathcal{O})$ 的弱 Serre 子范畴。
\end{lemma}

\begin{proof}
使用《同调代数》，引理
\ref{homology-lemma-characterize-weak-serre-subcategory} 的判别准则。
由引理 \ref{sites-modules-lemma-cokernel-map-towards-constructibles} 和
\ref{sites-modules-lemma-extension-towards-constructibles} 的结果，只需证明
$\mathcal{A}$ 中对象之间的映射 $\mathcal{F} \to \mathcal{G}$ 的核属于
$\mathcal{A}$。为此选择表示
$$
\bigoplus A_{V_j} \to \bigoplus A_{U_i} \to \mathcal{F} \to 0
\quad\text{且}\quad
\bigoplus A_{T_j} \to \bigoplus A_{W_i} \to \mathcal{G} \to 0
$$
。在此证明中直和总是有限的，并且对 $U \in \mathcal{B}$，记
$A_U = j_{U!}\mathcal{O}_U$。使用引理
\ref{sites-modules-lemma-covering-gives-surjection} 和
\ref{sites-modules-lemma-change-presentation-towards-constructibles}，并像引理
\ref{sites-modules-lemma-extension-towards-constructibles} 的证明那样论证，可以假设映射
$\mathcal{F} \to \mathcal{G}$ 提升为表示之间的映射
$$
\xymatrix{
\bigoplus A_{V_j} \ar[r] \ar[d] &
\bigoplus A_{U_i} \ar[r] \ar[d] &
\mathcal{F} \ar[r] \ar[d] & 0 \\
\bigoplus A_{T_j} \ar[r] &
\bigoplus A_{W_i} \ar[r] &
\mathcal{G} \ar[r] & 0
}
$$
。于是可见
$$
\Ker(\mathcal{F} \to \mathcal{G}) =
\Coker\left(\bigoplus A_{V_j} \to
\Ker\left(
\bigoplus A_{T_j} \oplus \bigoplus A_{U_i} \to \bigoplus A_{W_i}\right)\right)
$$
，本引理由假设以及引理
\ref{sites-modules-lemma-cokernel-map-towards-constructibles} 得出。
\end{proof}








\section{平坦态射}
\label{sites-modules-section-flat-morphisms}

\begin{definition}
\label{sites-modules-definition-flat-morphism}
设
$(f, f^\sharp) :
(\Sh(\mathcal{C}), \mathcal{O})
\longrightarrow
(\Sh(\mathcal{C}'), \mathcal{O}')$
是环拓扑斯的态射。如果环映射
$f^\sharp : f^{-1}\mathcal{O}' \to \mathcal{O}$ 平坦，则称
$(f, f^\sharp)$ 是{\it 平坦的}。如果环站点的态射所关联的环拓扑斯态射平坦，
则称该环站点态射是{\it 平坦的}。
\end{definition}

\begin{lemma}
\label{sites-modules-lemma-flat-pullback-exact}
设 $f : \Sh(\mathcal{C}) \to \Sh(\mathcal{C}')$ 是环拓扑斯的态射。则
$$
f^{-1} : \textit{Ab}(\mathcal{C}') \longrightarrow \textit{Ab}(\mathcal{C}),
\quad
\mathcal{F} \longmapsto f^{-1}\mathcal{F}
$$
正合。如果
$(f, f^\sharp) :
(\Sh(\mathcal{C}), \mathcal{O})
\to
(\Sh(\mathcal{C}'), \mathcal{O}')$
是环拓扑斯的平坦态射，则
$$
f^* : \textit{Mod}(\mathcal{O}') \longrightarrow \textit{Mod}(\mathcal{O}),
\quad
\mathcal{F} \longmapsto f^*\mathcal{F}
$$
正合。
\end{lemma}

\begin{proof}
给定 $\mathcal{C}'$ 上的阿贝尔层 $\mathcal{G}$，$f^{-1}\mathcal{G}$ 的底层
集合层与 $\mathcal{G}$ 的底层集合层的 $f^{-1}$ 相同；见《站点》，第
\ref{sites-section-sheaves-algebraic-structures} 节。因此 $f^{-1}$ 对集合层的
正合性（拓扑斯态射定义中的要求；见《站点》，定义
\ref{sites-definition-topos}）蕴含 $f^{-1}$ 作为阿贝尔层上的函子正合。

\medskip\noindent
对于模层的陈述，回顾 $f^*\mathcal{F}$ 定义为张量积
$f^{-1}\mathcal{F} \otimes_{f^{-1}\mathcal{O}', f^\sharp} \mathcal{O}$。
因此 $f^*$ 是两个正合函子的复合。
\end{proof}

\begin{definition}
\label{sites-modules-definition-flat-module}
设 $f : (\Sh(\mathcal{C}), \mathcal{O}) \to
(\Sh(\mathcal{D}), \mathcal{O}')$ 是环拓扑斯的态射，且
$\mathcal{F}$ 是 $\mathcal{O}$-模层。如果 $\mathcal{F}$ 作为
$f^{-1}\mathcal{O}'$-模平坦，则称 $\mathcal{F}$
{\it 在 $(\Sh(\mathcal{D}), \mathcal{O}')$ 上平坦}。
\end{definition}

\noindent
这与对环空间态射定义的概念相容；见《模层》，定义
\ref{modules-definition-flat-module} 及其后的讨论。

\begin{lemma}
\label{sites-modules-lemma-pullback-internal-hom}
设 $f : (\mathcal{C}, \mathcal{O}_\mathcal{C}) \to
(\mathcal{D}, \mathcal{O}_\mathcal{D})$ 是环站点的态射，且
$\mathcal{F}$、$\mathcal{G}$ 是 $\mathcal{O}_\mathcal{D}$-模。
如果 $\mathcal{F}$ 有限表示且 $f$ 平坦，则注
\ref{sites-modules-remark-pullback-internal-hom} 中的典范映射
$$
f^*\SheafHom_{\mathcal{O}_\mathcal{D}}(\mathcal{F}, \mathcal{G})
\longrightarrow
\SheafHom_{\mathcal{O}_\mathcal{C}}(f^*\mathcal{F}, f^*\mathcal{G})
$$
是同构。
\end{lemma}

\begin{proof}
设 $f$ 由连续函子 $u : \mathcal{D} \to \mathcal{C}$ 给出。需要证明对任意
$U \in \Ob(\mathcal{C})$，该映射到 $\mathcal{C}/U$ 的限制是同构。
可以用 $U$ 的某个覆盖的成员替换 $U$。因此由《站点》，引理
\ref{sites-lemma-morphism-of-sites-covering}，可以假设对某个
$V \in \Ob(\mathcal{C})$ 存在态射 $U \to u(V)$。当然，此时可用 $u(V)$
替换 $U$。又因为 $u$ 连续，可以用某个覆盖替换 $V$，并假设存在表示
$\mathcal{O}_V^{\oplus m} \to \mathcal{O}_V^{\oplus n} \to
\mathcal{F}|_V \to 0$
于 $\mathcal{D}/V$ 上。由于 $\SheafHom$ 的形成与局部化交换（引理
\ref{sites-modules-lemma-internal-hom-restriction}），可以用由 $f$ 诱导的态射
$(\mathcal{C}/u(V), \mathcal{O}_{u(V)}) \to
(\mathcal{D}/V, \mathcal{O}_V)$ 替换 $f$。因此归结到 $\mathcal{F}$ 有整体表示
$\mathcal{O}_\mathcal{D}^{\oplus m} \to \mathcal{O}_\mathcal{D}^{\oplus n} \to
\mathcal{F} \to 0$ 的情形。因为 $f$ 平坦且
$f^*\mathcal{O}_\mathcal{D} = \mathcal{O}_\mathcal{C}$，得到相应的表示
$\mathcal{O}_\mathcal{C}^{\oplus m} \to \mathcal{O}_\mathcal{C}^{\oplus n} \to
f^*\mathcal{F} \to 0$；见引理 \ref{sites-modules-lemma-flat-pullback-exact}。
利用 $\SheafHom$ 对第一个变元与有限直和交换、
$\SheafHom_{\mathcal{O}_\mathcal{C}}(\mathcal{O}_\mathcal{C}, -)$ 和
$\SheafHom_{\mathcal{O}_\mathcal{D}}(\mathcal{O}_\mathcal{D}, -)$
都是恒等函子，以及注 \ref{sites-modules-remark-pullback-internal-hom} 中构造的函子性，
得到交换图
$$
\xymatrix{
0 \ar[r] &
f^*\SheafHom_{\mathcal{O}_\mathcal{D}}(\mathcal{F}, \mathcal{G})
\ar[d] \ar[r] &
f^*\mathcal{G}^{\oplus n} \ar[d] \ar[r] &
f^*\mathcal{G}^{\oplus n} \ar[d] \\
0 \ar[r] &
\SheafHom_{\mathcal{O}_\mathcal{C}}(f^*\mathcal{F}, f^*\mathcal{G})
\ar[r] &
f^*\mathcal{G}^{\oplus n} \ar[r] &
f^*\mathcal{G}^{\oplus n}
}
$$
，其中右侧两个竖直箭头是同构。由引理
\ref{sites-modules-lemma-internal-hom-exact}，两行都正合。再由五引理即得结论。
\end{proof}







\section{可逆模}
\label{sites-modules-section-invertible}

\noindent
定义如下。

\begin{definition}
\label{sites-modules-definition-invertible-sheaf}
设 $(\mathcal{C}, \mathcal{O})$ 为环站点。
\begin{enumerate}
\item 若对 $\mathcal{C}$ 的每个对象 $U$，都存在 $U$ 的一个覆盖
$\{U_i \to U\}$，使得作为 $\mathcal{O}_{U_i}$-模，
$\mathcal{F}|_{U_i}$ 同构于 $\mathcal{O}_{U_i}^{\oplus r}$，
则称有限局部自由 $\mathcal{O}$-模 $\mathcal{F}$ 的{\it 秩为 $r$}。
\item 若函子
$$
\textit{Mod}(\mathcal{O}) \longrightarrow \textit{Mod}(\mathcal{O}),\quad
\mathcal{F}  \longmapsto \mathcal{F} \otimes_\mathcal{O} \mathcal{L}
$$
是等价，则称 $\mathcal{O}$-模 $\mathcal{L}$ {\it 可逆}。
\item 层 {\it $\mathcal{O}^*$} 是 $\mathcal{O}$ 的子层，按如下规则定义：
$$
U \longmapsto \mathcal{O}^*(U) = \{f \in \mathcal{O}(U) \mid
\exists g \in \mathcal{O}(U)\text{ 使得 }fg = 1\}
$$
它是以乘法为群运算的阿贝尔群层。
\end{enumerate}
\end{definition}

\noindent
下面的引理 \ref{sites-modules-lemma-invertible-is-locally-free-rank-1}
说明了它与秩为 $1$ 的局部自由模之间的关系。

\begin{lemma}
\label{sites-modules-lemma-invertible}
设 $(\mathcal{C}, \mathcal{O})$ 为环站点，$\mathcal{L}$
为 $\mathcal{O}$-模。下列条件等价：
\begin{enumerate}
\item $\mathcal{L}$ 可逆；
\item 存在 $\mathcal{O}$-模 $\mathcal{N}$，使得
$\mathcal{L} \otimes_\mathcal{O} \mathcal{N} \cong \mathcal{O}$.
\end{enumerate}
在此情形下，还有：
\begin{enumerate}
\item[(a)] $\mathcal{L}$ 是有限表示的平坦 $\mathcal{O}$-模；
\item[(b)] 对 $\mathcal{C}$ 的每个对象 $U$，都存在一个覆盖
$\{U_i \to U\}$，使得 $\mathcal{L}|_{U_i}$
是某个有限自由模的直和因子；
\item[(c)] (2) 中的模 $\mathcal{N}$ 同构于
$\SheafHom_\mathcal{O}(\mathcal{L}, \mathcal{O})$.
\end{enumerate}
\end{lemma}

\begin{proof}
假设 (1) 成立。此时函子 $- \otimes_\mathcal{O} \mathcal{L}$
本质满，故存在 (2) 中所述的 $\mathcal{O}$-模
$\mathcal{N}$。反之，若 (2) 成立，则函子
$- \otimes_\mathcal{O} \mathcal{N}$ 是函子
$- \otimes_\mathcal{O} \mathcal{L}$ 的拟逆，因而 (1) 成立。

\medskip\noindent
假设 (1) 与 (2) 成立。由于 $- \otimes_\mathcal{O} \mathcal{L}$
是等价，它是正合函子，故 $\mathcal{L}$ 平坦。记给定的同构为
$\psi : \mathcal{L} \otimes_\mathcal{O} \mathcal{N} \to \mathcal{O}$。
设 $U$ 为 $\mathcal{C}$ 的一个对象。我们将证明：在 $U$
的某个覆盖的各成员上，$\mathcal{L}$ 的限制是自由模的直和因子；
这显然蕴含 $\mathcal{L}$ 是有限表示的。根据 $\otimes$ 的构造，
将 $U$ 替换为某个覆盖的各成员后，可以假设存在整数
$n \geq 1$ 以及截面 $x_i \in \mathcal{L}(U)$、
$y_i \in \mathcal{N}(U)$，使得 $\psi(\sum x_i \otimes y_i) = 1$。
考虑同构
$$
\mathcal{L}|_U \to
\mathcal{L}|_U \otimes_{\mathcal{O}_U}
\mathcal{L}|_U \otimes_{\mathcal{O}_U} \mathcal{N}|_U \to \mathcal{L}|_U
$$
其中第一个箭头将 $x$ 映到 $\sum x_i \otimes x \otimes y_i$，
第二个箭头将 $x \otimes x' \otimes y$ 映到 $\psi(x' \otimes y)x$。
由此可知，$x \mapsto \sum \psi(x \otimes y_i)x_i$
是 $\mathcal{L}|_U$ 的自同构。该自同构分解为
$$
\mathcal{L}|_U \to \mathcal{O}_U^{\oplus n} \to \mathcal{L}|_U
$$
其中第一个箭头由
$x \mapsto (\psi(x \otimes y_1), \ldots, \psi(x \otimes y_n))$
给出，第二个箭头由 $(a_1, \ldots, a_n) \mapsto \sum a_i x_i$
给出。因此 $\mathcal{L}|_U$ 是某个有限自由
$\mathcal{O}_U$-模的直和因子。

\medskip\noindent
仍假设 (1) 与 (2) 成立。考虑求值映射
$$
\mathcal{L} \otimes_\mathcal{O}
\SheafHom_\mathcal{O}(\mathcal{L}, \mathcal{O}_X)
\longrightarrow \mathcal{O}_X
$$
为完成引理的证明，我们将证明它是同构。由引理
\ref{sites-modules-lemma-internal-hom-adjoint-tensor}，有
$$
\Hom_\mathcal{O}(\mathcal{O}, \mathcal{O}) =
\Hom_\mathcal{O}
(\mathcal{N} \otimes_\mathcal{O} \mathcal{L}, \mathcal{O})
\longrightarrow
\Hom_\mathcal{O}
(\mathcal{N}, \SheafHom_\mathcal{O}(\mathcal{L}, \mathcal{O}))
$$
$1$ 的像给出态射
$\mathcal{N} \to \SheafHom_\mathcal{O}(\mathcal{L}, \mathcal{O})$。
与 $\mathcal{L}$ 作张量积，得到
$$
\mathcal{O} = \mathcal{L} \otimes_\mathcal{O} \mathcal{N}
\longrightarrow
\mathcal{L} \otimes_\mathcal{O} \SheafHom_\mathcal{O}(\mathcal{L}, \mathcal{O})
$$
此映射是求值映射的逆；计算从略。
\end{proof}

\begin{lemma}
\label{sites-modules-lemma-pullback-invertible}
设 $f : (\Sh(\mathcal{C}), \mathcal{O}_\mathcal{C}) \to
(\Sh(\mathcal{D}), \mathcal{O}_\mathcal{D})$ 为环拓扑斯态射。
可逆 $\mathcal{O}_\mathcal{D}$-模
$\mathcal{L}$ 的拉回 $f^*\mathcal{L}$ 是可逆的。
\end{lemma}

\begin{proof}
由引理 \ref{sites-modules-lemma-invertible}，存在
$\mathcal{O}_\mathcal{D}$-模 $\mathcal{N}$，使得
$\mathcal{L} \otimes_{\mathcal{O}_\mathcal{D}} \mathcal{N} \cong
\mathcal{O}_\mathcal{D}$。作拉回并应用引理
\ref{sites-modules-lemma-tensor-product-pullback}，得到
$f^*\mathcal{L} \otimes_{\mathcal{O}_\mathcal{C}} f^*\mathcal{N} \cong
\mathcal{O}_\mathcal{C}$。
因此由引理 \ref{sites-modules-lemma-invertible}，$f^*\mathcal{L}$ 可逆。
\end{proof}

\begin{lemma}
\label{sites-modules-lemma-constructions-invertible}
设 $(\mathcal{C}, \mathcal{O})$ 为环站点。
\begin{enumerate}
\item 若 $\mathcal{L}$、$\mathcal{N}$ 是可逆
$\mathcal{O}$-模，则
$\mathcal{L} \otimes_\mathcal{O} \mathcal{N}$.
也可逆。
\item 若 $\mathcal{L}$ 是可逆 $\mathcal{O}$-模，则
$\SheafHom_\mathcal{O}(\mathcal{L}, \mathcal{O})$ 也可逆，且求值映射
$\mathcal{L} \otimes_\mathcal{O}
\SheafHom_\mathcal{O}(\mathcal{L}, \mathcal{O}) \to \mathcal{O}$
是同构。
\end{enumerate}
\end{lemma}

\begin{proof}
(1) 由定义即得，(2) 由引理 \ref{sites-modules-lemma-invertible} 及其证明即得。
\end{proof}

\begin{lemma}
\label{sites-modules-lemma-pic-set}
设 $(\mathcal{C}, \mathcal{O})$ 为环站点。存在一族可逆模
$\{\mathcal{L}_i\}_{i \in I}$，使得
$(\mathcal{C}, \mathcal{O})$ 上的每个可逆模恰与其中一个
$\mathcal{L}_i$ 同构。
\end{lemma}

\begin{proof}
从略；见《模层》，引理 \ref{modules-lemma-pic-set}。
\end{proof}

\noindent
引理 \ref{sites-modules-lemma-pic-set} 表明，可逆层的同构类全体构成一个集合。
引理 \ref{sites-modules-lemma-constructions-invertible} 表明，张量积在此集合上
定义了一个阿贝尔群结构，其中 $\mathcal{L}$ 的逆元由
$\SheafHom_\mathcal{O}(\mathcal{L}, \mathcal{O})$ 给出。

\medskip\noindent
事实上，给定可逆 $\mathcal{O}$-模 $\mathcal{L}$
及 $n \in \mathbf{Z}$，将等价
$\mathcal{F} \mapsto \mathcal{F} \otimes_\mathcal{O} \mathcal{L}$
对 $\mathcal{O}$ 迭代应用恰好 $n$ 次所得的像，定义为
$\mathcal{L}$ 的第 $n$ 个{\it 张量幂}
$\mathcal{L}^{\otimes n}$。因为我们把可逆 $\mathcal{O}$-模
定义为与其作张量积给出等价的模，所以该定义对负整数 $n$
也有意义。更明确地，
$$
\mathcal{L}^{\otimes n} =
\left\{
\begin{matrix}
\mathcal{O} & \text{若} & n = 0 \\
\SheafHom_\mathcal{O}(\mathcal{L}, \mathcal{O}) & \text{若} & n = -1\\
\mathcal{L} \otimes_\mathcal{O} \ldots \otimes_\mathcal{O} \mathcal{L}
& \text{若} & n > 0 \\
\mathcal{L}^{\otimes -1} \otimes_\mathcal{O} \ldots
\otimes_\mathcal{O} \mathcal{L}^{\otimes -1}
& \text{若} & n < -1
\end{matrix}
\right.
$$
见引理 \ref{sites-modules-lemma-constructions-invertible}。
按此定义，我们有典范同构
$\mathcal{L}^{\otimes n} \otimes_\mathcal{O}
\mathcal{L}^{\otimes m} \to
\mathcal{L}^{\otimes n + m}$，并且这些同构满足交换性约束与结合性约束
（表述从略）。

\begin{definition}
\label{sites-modules-definition-pic}
设 $(\mathcal{C}, \mathcal{O})$ 为环站点。该环站点的
{\it Picard 群} $\Pic(\mathcal{O})$ 是以可逆
$\mathcal{O}$-模的同构类为元素、加法对应于张量积的阿贝尔群。
\end{definition}









\section{微分模}
\label{sites-modules-section-differentials}

\noindent
本节简要说明如何为环拓扑斯态射定义相对微分模。建议读者参阅
交换代数一章中的相应一节
（《代数》，第 \ref{algebra-section-differentials} 节）。

\begin{definition}
\label{sites-modules-definition-derivation}
设 $\mathcal{C}$ 为站点，$\varphi : \mathcal{O}_1 \to \mathcal{O}_2$
为环层同态，$\mathcal{F}$ 为 $\mathcal{O}_2$-模。取值于
$\mathcal{F}$ 的{\it $\mathcal{O}_1$-导子}，更准确地说是
{\it $\varphi$-导子}，是一个映射
$D : \mathcal{O}_2 \to \mathcal{F}$，它可加，零化
$\mathcal{O}_1 \to \mathcal{O}_2$ 的像，并满足{\it Leibniz 法则}
$$
D(ab) = aD(b) + D(a)b
$$
其中 $a,b$ 是 $\mathcal{O}_2$ 的任意两个局部截面
（在二者均有定义之处）。记
$\text{Der}_{\mathcal{O}_1}(\mathcal{O}_2, \mathcal{F})$
为取值于 $\mathcal{F}$ 的 $\varphi$-导子所成的集合。
\end{definition}

\noindent
这是《代数》，定义 \ref{sites-modules-definition-derivation} 的层论版本。
给定定义中所述的导子 $D : \mathcal{O}_2 \to \mathcal{F}$，
整体截面上的映射
$$
D : \Gamma(\mathcal{O}_2) \longrightarrow \Gamma(\mathcal{F})
$$
显然是代数定义意义下的 $\Gamma(\mathcal{O}_1)$-导子。注意，若
$\alpha : \mathcal{F} \to \mathcal{G}$ 是
$\mathcal{O}_2$-模的映射，则有诱导映射
$$
\text{Der}_{\mathcal{O}_1}(\mathcal{O}_2, \mathcal{F})
\longrightarrow
\text{Der}_{\mathcal{O}_1}(\mathcal{O}_2, \mathcal{G})
$$
由规则 $D \mapsto \alpha \circ D$ 给出。换言之，我们得到一个函子。

\begin{lemma}
\label{sites-modules-lemma-universal-module}
设 $\mathcal{C}$ 为站点，$\varphi : \mathcal{O}_1 \to \mathcal{O}_2$
为环层同态。函子
$$
\textit{Mod}(\mathcal{O}_2) \longrightarrow \textit{Ab}, \quad
\mathcal{F} \longmapsto \text{Der}_{\mathcal{O}_1}(\mathcal{O}_2, \mathcal{F})
$$
可表。
\end{lemma}

\begin{proof}
证明与代数中的相应命题完全相同。为叙述该证明，对
$\mathcal{C}$ 上任意集合层 $\mathcal{F}$，记
$\mathcal{O}_2[\mathcal{F}]$ 为预层
$U \mapsto \mathcal{O}_2(U)[\mathcal{F}(U)]$ 的层化，其中右端表示
由集合 $\mathcal{F}(U)$ 生成的自由 $\mathcal{O}_2(U)$-模。
对 $s \in \mathcal{F}(U)$，以 $[s]$ 表示
$\mathcal{O}_2[\mathcal{F}]$ 在 $U$ 上的相应截面。若
$\mathcal{F}$ 是 $\mathcal{O}_2$-模层，则有典范映射
$$
c : \mathcal{O}_2[\mathcal{F}] \longrightarrow \mathcal{F}
$$
它在预层层次上由规则
$\sum f_s[s] \mapsto \sum f_s s$ 给出。以下以简写
$[s] \mapsto s$ 描述此映射，其他映射亦同。考虑
$\mathcal{O}_2$-模的映射
\begin{equation}
\label{sites-modules-equation-define-module-differentials}
\begin{matrix}
\mathcal{O}_2[\mathcal{O}_2 \times \mathcal{O}_2] \oplus
\mathcal{O}_2[\mathcal{O}_2 \times \mathcal{O}_2] \oplus
\mathcal{O}_2[\mathcal{O}_1] &
\longrightarrow &
\mathcal{O}_2[\mathcal{O}_2] \\
[(a, b)] \oplus [(f, g)] \oplus [h] & \longmapsto & [a + b] - [a] - [b] + \\
& & [fg] - g[f] - f[g] + \\
& & [\varphi(h)]
\end{matrix}
\end{equation}
其中采用上述简写。令 $\Omega_{\mathcal{O}_2/\mathcal{O}_1}$
为此映射的余核。显然存在集合层的映射
$$
\text{d} : \mathcal{O}_2 \longrightarrow \Omega_{\mathcal{O}_2/\mathcal{O}_1}
$$
它将局部截面 $f$ 映到 $[f]$ 在
$\Omega_{\mathcal{O}_2/\mathcal{O}_1}$ 中的像。由构造，
$\text{d}$ 是 $\mathcal{O}_1$-导子。再设 $\mathcal{F}$
为 $\mathcal{O}_2$-模层，且
$D : \mathcal{O}_2 \to \mathcal{F}$ 为 $\mathcal{O}_1$-导子。
考虑将 $[g]$ 映到 $D(g)$ 的 $\mathcal{O}_2$-线性映射
$\mathcal{O}_2[\mathcal{O}_2] \to \mathcal{F}$。
由导子的定义，该映射零化映射
(\ref{sites-modules-equation-define-module-differentials}) 像中的截面，故定义了映射
$$
\alpha_D : \Omega_{\mathcal{O}_2/\mathcal{O}_1} \longrightarrow \mathcal{F}
$$
显然 $D = \alpha_D \circ \text{d}$，故引理得证。
\end{proof}

\begin{definition}
\label{sites-modules-definition-module-differentials}
设 $\mathcal{C}$ 为站点，$\varphi : \mathcal{O}_1 \to \mathcal{O}_2$
为环层同态。环映射 $\varphi$ 的{\it 微分模}是表示函子
$\mathcal{F} \mapsto \text{Der}_{\mathcal{O}_1}(\mathcal{O}_2, \mathcal{F})$
的对象；由引理 \ref{sites-modules-lemma-universal-module}，该对象存在。
记它为 $\Omega_{\mathcal{O}_2/\mathcal{O}_1}$，并将{\it 泛
$\varphi$-导子}记为
$\text{d} : \mathcal{O}_2 \to \Omega_{\mathcal{O}_2/\mathcal{O}_1}$.
\end{definition}

\noindent
由于此模及该导子构成表示一个函子的泛对象，这一概念显然是内蕴的
（即不依赖于环拓扑斯所基于的站点的选择；见第
\ref{sites-modules-section-intrinsic} 节）。注意，
$\Omega_{\mathcal{O}_2/\mathcal{O}_1}$ 是
$\mathcal{O}_2$-模映射
(\ref{sites-modules-equation-define-module-differentials}) 的余核。此外，映射
$\text{d}$ 由如下规则描述：$\text{d}f$ 是局部截面 $[f]$ 的像。

\begin{lemma}
\label{sites-modules-lemma-differentials-sheafify}
设 $\mathcal{C}$ 为站点，$\varphi : \mathcal{O}_1 \to \mathcal{O}_2$
为环预层同态。则
$\Omega_{\mathcal{O}_2^\#/\mathcal{O}_1^\#}$ 是预层
$U \mapsto \Omega_{\mathcal{O}_2(U)/\mathcal{O}_1(U)}$ 的关联层。
\end{lemma}

\begin{proof}
考虑映射 (\ref{sites-modules-equation-define-module-differentials})。
存在一个类似的预层映射，它在 $U \in \Ob(\mathcal{C})$ 上的取值为
$$
\mathcal{O}_2(U)[\mathcal{O}_2(U) \times \mathcal{O}_2(U)] \oplus
\mathcal{O}_2(U)[\mathcal{O}_2(U) \times \mathcal{O}_2(U)] \oplus
\mathcal{O}_2(U)[\mathcal{O}_1(U)]
\longrightarrow
\mathcal{O}_2(U)[\mathcal{O}_2(U)]
$$
由《代数》，定义 \ref{algebra-definition-differentials}
中微分模的构造，此映射的余核在 $U$ 上的取值为
$\Omega_{\mathcal{O}_2(U)/\mathcal{O}_1(U)}$。另一方面，
(\ref{sites-modules-equation-define-module-differentials}) 中的各层正是上述
各预层的层化。由于层化正合，结论随即成立。
\end{proof}

\begin{lemma}
\label{sites-modules-lemma-pullback-differentials}
设 $f : \Sh(\mathcal{D}) \to \Sh(\mathcal{C})$ 为拓扑斯态射，
$\varphi : \mathcal{O}_1 \to \mathcal{O}_2$
为 $\mathcal{C}$ 上的环层同态。则有与泛导子相容的典范等同
$f^{-1}\Omega_{\mathcal{O}_2/\mathcal{O}_1} =
\Omega_{f^{-1}\mathcal{O}_2/f^{-1}\mathcal{O}_1}$。
\end{lemma}

\begin{proof}
这是因为层 $\Omega_{\mathcal{O}_2/\mathcal{O}_1}$ 是映射
(\ref{sites-modules-equation-define-module-differentials}) 的余核，
$\Omega_{f^{-1}\mathcal{O}_2/f^{-1}\mathcal{O}_1}$ 也有类似描述，
函子 $f^{-1}$ 正合，并且
$f^{-1}(\mathcal{O}_2[\mathcal{O}_2]) =
f^{-1}\mathcal{O}_2[f^{-1}\mathcal{O}_2]$,
$f^{-1}(\mathcal{O}_2[\mathcal{O}_2 \times \mathcal{O}_2]) =
f^{-1}\mathcal{O}_2[f^{-1}\mathcal{O}_2 \times f^{-1}\mathcal{O}_2]$, and
$f^{-1}(\mathcal{O}_2[\mathcal{O}_1]) =
f^{-1}\mathcal{O}_2[f^{-1}\mathcal{O}_1]$.
\end{proof}

\begin{lemma}
\label{sites-modules-lemma-localize-differentials}
设 $\mathcal{C}$ 为站点，$\varphi : \mathcal{O}_1 \to \mathcal{O}_2$
为环层同态。对 $\mathcal{C}$ 的任意对象 $U$，有与泛导子相容的
典范同构
$$
\Omega_{\mathcal{O}_2/\mathcal{O}_1}|_U =
\Omega_{(\mathcal{O}_2|_U)/(\mathcal{O}_1|_U)}
$$
\end{lemma}

\begin{proof}
这是引理 \ref{sites-modules-lemma-pullback-differentials} 的特例。
\end{proof}

\begin{lemma}
\label{sites-modules-lemma-functoriality-differentials}
设 $\mathcal{C}$ 为站点。设
$$
\xymatrix{
\mathcal{O}_2 \ar[r]_\varphi & \mathcal{O}_2' \\
\mathcal{O}_1 \ar[r] \ar[u] & \mathcal{O}'_1 \ar[u]
}
$$
为 $\mathcal{C}$ 上环层的交换图。将映射
$\mathcal{O}_2 \to \mathcal{O}'_2$ 与映射
$\text{d} : \mathcal{O}'_2 \to \Omega_{\mathcal{O}'_2/\mathcal{O}'_1}$
复合，得到一个 $\mathcal{O}_1$-导子。因此得到典范
$\mathcal{O}_2$-模映射
$\Omega_{\mathcal{O}_2/\mathcal{O}_1} \to
\Omega_{\mathcal{O}'_2/\mathcal{O}'_1}$.
它由下述性质唯一刻画：对 $\mathcal{O}_2$ 的任意局部截面 $f$，
$\text{d}(f)$ 映到 $\text{d}(\varphi(f))$。
由此，$\Omega_{-/-}$ 成为环层箭头范畴上的函子。
\end{lemma}

\begin{proof}
此引理由其陈述即得。
\end{proof}

\begin{lemma}
\label{sites-modules-lemma-differential-seq}
在引理 \ref{sites-modules-lemma-functoriality-differentials} 中，假设
$\mathcal{O}_2 \to \mathcal{O}'_2$ 满射，其核为
$\mathcal{I} \subset \mathcal{O}_2$，并假设
$\mathcal{O}_1 = \mathcal{O}'_1$。则有典范的
$\mathcal{O}'_2$-模正合列
$$
\mathcal{I}/\mathcal{I}^2
\longrightarrow
\Omega_{\mathcal{O}_2/\mathcal{O}_1} \otimes_{\mathcal{O}_2} \mathcal{O}'_2
\longrightarrow
\Omega_{\mathcal{O}'_2/\mathcal{O}_1}
\longrightarrow
0
$$
最左侧映射由下述规则刻画：$\mathcal{I}$ 的局部截面
$f$ 映到 $\text{d}f \otimes 1$。
\end{lemma}

\begin{proof}
对 $\mathcal{I}$ 的局部截面 $f$，记 $\overline{f}$ 为
$f$ 在 $\mathcal{I}/\mathcal{I}^2$ 中的像。为证明映射
$\overline{f} \mapsto \text{d}f \otimes 1$ 良定义，只需验证：
若 $f_1,f_2$ 是 $\mathcal{I}$ 的局部截面，则
$\text{d} f_1f_2 \otimes 1 = 0$。这由 Leibniz 法则立即得到：
$\text{d} f_1f_2 \otimes 1 =
(f_1 \text{d}f_2 + f_2 \text{d} f_1 )\otimes 1 =
\text{d}f_2 \otimes f_1 + \text{d}f_2 \otimes f_1 = 0$.
类似的计算表明，此映射是
$\mathcal{O}'_2 = \mathcal{O}_2/\mathcal{I}$-线性的。右侧映射是
引理 \ref{sites-modules-lemma-functoriality-differentials} 中的映射。

\medskip\noindent
为证明该列正合，作如下论证。令
$\mathcal{O}''_2 \subset \mathcal{O}'_2$ 为
$\mathcal{O}_1$-代数预层，它在 $U$ 上的取值是
$\mathcal{O}_2(U) \to \mathcal{O}'_2(U)$ 的像。由
《代数》，引理 \ref{algebra-lemma-differential-seq}，对
$\mathcal{C}$ 的每个对象 $U$，下列序列正合：
$$
\mathcal{I}(U)/\mathcal{I}(U)^2
\longrightarrow
\Omega_{\mathcal{O}_2(U)/\mathcal{O}_1(U)}
\otimes_{\mathcal{O}_2(U)} \mathcal{O}''_2(U)
\longrightarrow
\Omega_{\mathcal{O}''_2(U)/\mathcal{O}_1(U)}
\longrightarrow
0
$$
由于层化正合，这给出一个
$(\mathcal{O}'_2)^\#$-模层的正合列。再由引理
\ref{sites-modules-lemma-differentials-sheafify} 及
$(\mathcal{O}''_2)^\# = \mathcal{O}'_2$，即得结论。
\end{proof}

\noindent
下面给出导子自然出现的一种特殊情形。

\begin{lemma}
\label{sites-modules-lemma-double-structure-gives-derivation}
设 $\mathcal{C}$ 为站点，$\varphi : \mathcal{O}_1 \to \mathcal{O}_2$
为环层同态。考虑短正合列
$$
0 \to \mathcal{F} \to \mathcal{A} \to \mathcal{O}_2 \to 0
$$
这里 $\mathcal{A}$ 是 $\mathcal{O}_1$-代数层，
$\pi : \mathcal{A} \to \mathcal{O}_2$ 是
$\mathcal{O}_1$-代数层的满射，且
$\mathcal{F} = \Ker(\pi)$ 是其核。假设 $\mathcal{F}$
是 $\mathcal{A}$ 中平方为零的理想层。因此
$\mathcal{F}$ 自然带有 $\mathcal{O}_2$-模结构。
$\pi$ 的截面 $s : \mathcal{O}_2 \to \mathcal{A}$
是满足 $\pi \circ s = \text{id}$ 的 $\mathcal{O}_1$-代数映射。
给定 $\pi$ 的任意截面 $s : \mathcal{O}_2 \to \mathcal{F}$
以及任意 $\varphi$-导子 $D : \mathcal{O}_1 \to \mathcal{F}$，映射
$$
s + D : \mathcal{O}_1 \to \mathcal{A}
$$
是 $\pi$ 的截面，并且每个截面 $s'$ 都唯一地写成
$s + D$，其中 $D$ 是 $\varphi$-导子。
\end{lemma}

\begin{proof}
回忆，$\mathcal{F}$ 上的 $\mathcal{O}_2$-模结构由
$h \tau = \tilde h \tau$ 给出（右侧是在 $\mathcal{A}$ 中的乘法）；
这里 $h$ 是 $\mathcal{O}_2$ 的局部截面，
$\tilde h$ 是 $h$ 到 $\mathcal{A}$ 的局部截面的一个局部提升，
而 $\tau$ 是 $\mathcal{F}$ 的局部截面。特别地，给定 $s$ 后，
可以取 $\tilde h = s(h)$。为验证 $s + D$ 是环层同态，计算
\begin{eqnarray*}
(s + D)(ab) & = & s(ab) + D(ab) \\
& = & s(a)s(b) + aD(b) + D(a)b \\
& = & s(a) s(b) + s(a)D(b) + D(a)s(b) \\
& = & (s(a) + D(a))(s(b) + D(b))
\end{eqnarray*}
其中使用了 Leibniz 法则。类似地，由于 $D$ 是
$\mathcal{O}_1$-导子，可证 $s + D$ 是
$\mathcal{O}_1$-代数映射。反之，给定 $s'$，令
$D = s' - s$ 即可。细节从略。
\end{proof}

\begin{definition}
\label{sites-modules-definition-sheaf-differentials}
设 $X = (\Sh(\mathcal{C}), \mathcal{O})$ 与
$Y = (\Sh(\mathcal{C}'), \mathcal{O}')$ 为环拓扑斯，并设
$(f, f^\sharp) : X \to Y$ 为环拓扑斯态射。在此情形下：
\begin{enumerate}
\item 对 $\mathcal{O}$-模层 $\mathcal{F}$，{\it $Y$-导子}
$D : \mathcal{O} \to \mathcal{F}$ 就是 $f^\sharp$-导子；
\item {\it $X$ 在 $Y$ 上的微分层 $\Omega_{X/Y}$}
是如下映射的微分模：
$f^\sharp : f^{-1}\mathcal{O}' \to \mathcal{O}$,
见定义 \ref{sites-modules-definition-module-differentials}。
\end{enumerate}
因此 $\Omega_{X/Y}$ 配备一个{\it 泛 $Y$-导子}
$\text{d}_{X/Y} : \mathcal{O} \longrightarrow \Omega_{X/Y}$。
有时也写作 $\Omega_{X/Y} = \Omega_f$。
\end{definition}

\noindent
回忆 $f^\sharp : f^{-1}\mathcal{O}' \to \mathcal{O}$，故此定义有意义。

\begin{lemma}
\label{sites-modules-lemma-functoriality-differentials-ringed-topoi}
设
$X = (\Sh(\mathcal{C}_X), \mathcal{O}_X)$,
$Y = (\Sh(\mathcal{C}_Y), \mathcal{O}_Y)$,
$X' = (\Sh(\mathcal{C}_{X'}), \mathcal{O}_{X'})$, and
$Y' = (\Sh(\mathcal{C}_{Y'}), \mathcal{O}_{Y'})$ 为环拓扑斯。设
$$
\xymatrix{
X' \ar[d] \ar[r]_f & X \ar[d] \\
Y' \ar[r] & Y
}
$$
为环拓扑斯态射的交换图。将映射
$f^\sharp : \mathcal{O}_X \to f_*\mathcal{O}_{X'}$ 与映射
$f_*\text{d}_{X'/Y'} : f_*\mathcal{O}_{X'} \to f_*\Omega_{X'/Y'}$
复合，得到一个 $Y$-导子。因此得到典范
$\mathcal{O}_X$-模映射
$\Omega_{X/Y} \to f_*\Omega_{X'/Y'}$；再由
$f_*$ 与 $f^*$ 的伴随性，得到典范的
$\mathcal{O}_{X'}$-模同态
$$
c_f : f^*\Omega_{X/Y} \longrightarrow \Omega_{X'/Y'}.
$$
它由下述性质唯一刻画：对 $\mathcal{O}_X$ 的任意局部截面 $t$，
$f^*\text{d}_{X/Y}(t)$ 映到 $\text{d}_{X'/Y'}(f^* t)$。
\end{lemma}

\begin{proof}
除最后一个断言外均显然。下面解释该断言的含义。设
$U \in \Ob(\mathcal{C}_X)$ 且 $t \in \mathcal{O}_X(U)$；
这正是说 $t$ 是 $\mathcal{O}_X$ 的局部截面。可以把 $t$
视为集合层的映射
$t : h_U^\# \to \mathcal{O}_X$。于是有
$f^{-1}t : f^{-1}h_U^\# \to f^{-1}\mathcal{O}_X$。
所谓 $f^*t$，是指复合
$$
\xymatrix{
f^{-1}h_U^\# \ar[rr]^{f^{-1}t} \ar@/^4ex/[rrrr]^{f^*t} & &
f^{-1}\mathcal{O}_X \ar[rr]^{f^\sharp} & &
\mathcal{O}_{X'}
}
$$
注意 $\text{d}_{X/Y}(t) \in \Omega_{X/Y}(U)$。因此可以把
$\text{d}_{X/Y}(t)$ 视为映射
$\text{d}_{X/Y}(t) : h_U^\# \to \Omega_{X/Y}$。于是
$f^{-1}\text{d}_{X/Y}(t) : f^{-1}h_U^\# \to f^{-1}\Omega_{X/Y}$。
所谓 $f^*\text{d}_{X/Y}(t)$，是指复合
$$
\xymatrix{
f^{-1}h_U^\#
\ar[rr]^{f^{-1}\text{d}_{X/Y}(t)}
\ar@/^4ex/[rrrr]^{f^*\text{d}_{X/Y}(t)} & &
f^{-1}\Omega_{X/Y} \ar[rr]^{1 \otimes \text{id}} & &
f^*\Omega_{X/Y}
}
$$
于是，引理的断言是指
$$
c_f \circ f^*t = f^*\text{d}_{X/Y}(t)
$$
作为从 $f^{-1}h_U^\#$ 到 $\Omega_{X'/Y'}$ 的映射相等。
略去验证引理中定义的 $c_f$ 满足此性质。（提示：第一个映射
$c'_f : \Omega_{X/Y} \to f_*\Omega_{X'/Y'}$ 满足
$c'_f(\text{d}_{X/Y}(t)) = f_*\text{d}_{X'/Y'}(f^\sharp(t))$，
这是 $f_*\Omega_{X'/Y'}$ 在 $U$ 上的截面等式；须利用伴随性
将其化为上述等式。）此性质唯一刻画 $c_f$，因为
$f^*\text{d}_{X/Y}(t)$ 的像生成 $\mathcal{O}_{X'}$-模
$f^*\Omega_{X/Y}$；这又是因为局部截面
$\text{d}_{X/Y}(t)$ 生成 $\mathcal{O}_X$-模 $\Omega_{X/Y}$。
\end{proof}





\section{有限阶微分算子}
\label{sites-modules-section-differential-operators}

\noindent
本节介绍有限阶微分算子。建议读者参阅交换代数一章中的相应一节
（《代数》，第 \ref{algebra-section-differential-operators} 节）。

\begin{definition}
\label{sites-modules-definition-differential-operators}
设 $\mathcal{C}$ 为站点，$\varphi : \mathcal{O}_1 \to \mathcal{O}_2$
为环层同态，$k \geq 0$ 为整数，并设 $\mathcal{F}$、$\mathcal{G}$
为 $\mathcal{O}_2$-模层。所谓{\it 阶数为 $k$ 的微分算子
$D : \mathcal{F} \to \mathcal{G}$}，是指一个
$\mathcal{O}_1$-线性映射，满足：对 $\mathcal{O}_2$
的任意局部截面 $g$，映射 $s \mapsto D(gs) - gD(s)$
是阶数为 $k - 1$ 的微分算子。作为起始情形，当 $k = 0$ 时，
将阶数为 $0$ 的微分算子定义为 $\mathcal{O}_2$-线性映射。
\end{definition}

\noindent
若 $D : \mathcal{F} \to \mathcal{G}$ 是阶数为 $k$ 的微分算子，
则对 $\mathcal{O}_2$ 的任意局部截面 $g$，映射 $gD$
也是阶数为 $k$ 的微分算子。两个阶数为 $k$ 的微分算子之和
仍是同阶微分算子。因此所有这些算子所成的集合
$$
\text{Diff}^k(\mathcal{F}, \mathcal{G}) =
\text{Diff}^k_{\mathcal{O}_2/\mathcal{O}_1}(\mathcal{F}, \mathcal{G})
$$
是 $\Gamma(\mathcal{C}, \mathcal{O}_2)$-模。并且
$$
\text{Diff}^0(\mathcal{F}, \mathcal{G}) \subset
\text{Diff}^1(\mathcal{F}, \mathcal{G}) \subset
\text{Diff}^2(\mathcal{F}, \mathcal{G}) \subset \ldots
$$
把 $U \in \Ob(\mathcal{C})$ 映到阶数为 $k$ 的微分算子
$D : \mathcal{F}|_U \to \mathcal{G}|_U$ 所成模的规则，
给出站点 $\mathcal{C}$ 上的一个 $\mathcal{O}_2$-模层。
由此得到微分算子层（若以后需要，我们将在此补充定义）。

\begin{lemma}
\label{sites-modules-lemma-composition-differential-operators}
设 $\mathcal{C}$ 为站点，
$\mathcal{O}_1 \to \mathcal{O}_2$ 为环层映射，并设
$\mathcal{E}, \mathcal{F}, \mathcal{G}$ 为
$\mathcal{O}_2$-模层。若
$D : \mathcal{E} \to \mathcal{F}$ 与
$D' : \mathcal{F} \to \mathcal{G}$ 分别是阶数为 $k$ 与 $k'$
的微分算子，则 $D' \circ D$ 是阶数为 $k + k'$ 的微分算子。
\end{lemma}

\begin{proof}
设 $g$ 为 $\mathcal{O}_2$ 的局部截面。将 $\mathcal{E}$
的局部截面 $x$ 映到
$$
D'(D(gx)) - gD'(D(x)) = D'(D(gx)) - D'(gD(x)) + D'(gD(x)) - gD'(D(x))
$$
的映射，是两个较低阶微分算子复合的和。因此对 $k + k'$
作归纳即可证明引理。
\end{proof}

\begin{lemma}
\label{sites-modules-lemma-module-principal-parts}
设 $\mathcal{C}$ 为站点，
$\mathcal{O}_1 \to \mathcal{O}_2$ 为环层映射，
$\mathcal{F}$ 为 $\mathcal{O}_2$-模层，且 $k \geq 0$。
则存在 $\mathcal{O}_2$-模层
$\mathcal{P}^k_{\mathcal{O}_2/\mathcal{O}_1}(\mathcal{F})$
及典范同构
$$
\text{Diff}^k_{\mathcal{O}_2/\mathcal{O}_1}(\mathcal{F}, \mathcal{G}) =
\Hom_{\mathcal{O}_2}(
\mathcal{P}^k_{\mathcal{O}_2/\mathcal{O}_1}(\mathcal{F}), \mathcal{G})
$$
它关于 $\mathcal{O}_2$-模 $\mathcal{G}$ 函子性成立。
\end{lemma}

\begin{proof}
存在性可由一般范畴论论证得到（此处将来插入参考文献），但我们还将
给出一个直接构造，因为该构造会用于后续证明。以下自由使用引理
\ref{sites-modules-lemma-universal-module} 的证明中引入的记号。
给定任意微分算子 $D : \mathcal{F} \to \mathcal{G}$，
得到 $\mathcal{O}_2$-线性映射
$L_D : \mathcal{O}_2[\mathcal{F}] \to \mathcal{G}$
将 $[m]$ 映到 $D(m)$。若 $D$ 的阶数为 $0$，则
$L_D$ 零化局部截面
$$
[m + m'] - [m] - [m'],\quad
g_0[m] - [g_0m]
$$
其中 $g_0$ 是 $\mathcal{O}_2$ 的局部截面，而 $m,m'$
是 $\mathcal{F}$ 的局部截面。若 $D$ 的阶数为 $1$，则
$L_D$ 零化局部截面
$$
[m + m'] - [m] - [m'],\quad
f[m] - [fm], \quad
g_0g_1[m] - g_0[g_1m] - g_1[g_0m] + [g_1g_0m]
$$
其中 $f$ 是 $\mathcal{O}_1$ 的局部截面，
$g_0,g_1$ 是 $\mathcal{O}_2$ 的局部截面，而
$m,m'$ 是 $\mathcal{F}$ 的局部截面。若 $D$ 的阶数为 $k$，
则 $L_D$ 零化局部截面
$[m + m'] - [m] - [m']$、$f[m] - [fm]$ 以及
$$
g_0g_1\ldots g_k[m] - \sum g_0 \ldots \hat g_i \ldots g_k[g_im] + \ldots
+(-1)^{k + 1}[g_0\ldots g_km]
$$
反之，若 $L : \mathcal{O}_2[\mathcal{F}] \to \mathcal{G}$
是零化上句所列全部局部截面的 $\mathcal{O}_2$-线性映射，
则 $m \mapsto L([m])$ 是阶数为 $k$ 的微分算子。因此
$\mathcal{P}^k_{\mathcal{O}_2/\mathcal{O}_1}(\mathcal{F})$
是 $\mathcal{O}_2[\mathcal{F}]$ 对由这些局部截面生成的
$\mathcal{O}_2$-子模的商。
\end{proof}

\begin{definition}
\label{sites-modules-definition-module-principal-parts}
设 $\mathcal{C}$ 为站点，
$\mathcal{O}_1 \to \mathcal{O}_2$ 为环层映射，
$\mathcal{F}$ 为 $\mathcal{O}_2$-模层。引理
\ref{sites-modules-lemma-module-principal-parts} 中构造的模
$\mathcal{P}^k_{\mathcal{O}_2/\mathcal{O}_1}(\mathcal{F})$
称为 $\mathcal{F}$ 的{\it $k$ 阶主部模}。
\end{definition}

\noindent
注意，包含关系
$$
\text{Diff}^0(\mathcal{F}, \mathcal{G}) \subset
\text{Diff}^1(\mathcal{F}, \mathcal{G}) \subset
\text{Diff}^2(\mathcal{F}, \mathcal{G}) \subset \ldots
$$
由 Yoneda 引理（《范畴》，引理 \ref{categories-lemma-yoneda}）
对应于满射
$$
\ldots \to \mathcal{P}^2_{\mathcal{O}_2/\mathcal{O}_1}(\mathcal{F})
\to \mathcal{P}^1_{\mathcal{O}_2/\mathcal{O}_1}(\mathcal{F})
\to \mathcal{P}^0_{\mathcal{O}_2/\mathcal{O}_1}(\mathcal{F}) = \mathcal{F}
$$

\begin{lemma}
\label{sites-modules-lemma-differential-operators-sheafify}
设 $\mathcal{C}$ 为站点，$\mathcal{O}_1 \to \mathcal{O}_2$
为环预层同态，并设 $\mathcal{F}$ 为
$\mathcal{O}_2$-模预层。则
$\mathcal{P}^k_{\mathcal{O}_2^\#/\mathcal{O}_1^\#}(\mathcal{F}^\#)$
是预层
$U \mapsto P^k_{\mathcal{O}_2(U)/\mathcal{O}_1(U)}(\mathcal{F}(U))$
的关联层。
\end{lemma}

\begin{proof}
这可以用与引理 \ref{sites-modules-lemma-differentials-sheafify}
中微分层完全相同的方法证明。或许更自然的方法是直接利用引理
\ref{sites-modules-lemma-module-principal-parts} 的泛性质看出该等式。细节从略。
\end{proof}

\begin{lemma}
\label{sites-modules-lemma-sequence-of-principal-parts}
设 $\mathcal{C}$ 为站点，$\mathcal{O}_1 \to \mathcal{O}_2$
为环层同态，并设 $\mathcal{F}$ 为 $\mathcal{O}_2$-模层。
有典范短正合列
$$
0 \to
\Omega_{\mathcal{O}_2/\mathcal{O}_1} \otimes_{\mathcal{O}_2} \mathcal{F} \to
\mathcal{P}^1_{\mathcal{O}_2/\mathcal{O}_1}(\mathcal{F}) \to
\mathcal{F} \to 0
$$
它关于 $\mathcal{F}$ 函子性成立，称为{\it 主部序列}。
\end{lemma}

\begin{proof}
由交换代数版本（《代数》，引理
\ref{algebra-lemma-sequence-of-principal-parts}）及引理
\ref{sites-modules-lemma-differentials-sheafify}、\ref{sites-modules-lemma-differential-operators-sheafify}
即得。
\end{proof}

\begin{remark}
\label{sites-modules-remark-functoriality-principal-parts}
设 $\mathcal{C}$ 为站点，并给定环层的交换图
$$
\xymatrix{
\mathcal{B} \ar[r] & \mathcal{B}' \\
\mathcal{A} \ar[u] \ar[r] & \mathcal{A}' \ar[u]
}
$$
以及 $\mathcal{B}$-模 $\mathcal{F}$、
$\mathcal{B}'$-模 $\mathcal{F}'$ 和
$\mathcal{B}$-线性映射 $\mathcal{F} \to \mathcal{F}'$。
则得到相容的模映射系
$$
\xymatrix{
\ldots \ar[r] &
\mathcal{P}^2_{\mathcal{B}'/\mathcal{A}'}(\mathcal{F}') \ar[r] &
\mathcal{P}^1_{\mathcal{B}'/\mathcal{A}'}(\mathcal{F}') \ar[r] &
\mathcal{P}^0_{\mathcal{B}'/\mathcal{A}'}(\mathcal{F}') \\
\ldots \ar[r] &
\mathcal{P}^2_{\mathcal{B}/\mathcal{A}}(\mathcal{F}) \ar[r] \ar[u] &
\mathcal{P}^1_{\mathcal{B}/\mathcal{A}}(\mathcal{F}) \ar[r] \ar[u] &
\mathcal{P}^0_{\mathcal{B}/\mathcal{A}}(\mathcal{F}) \ar[u]
}
$$
这些映射与此类映射的进一步复合相容。最容易的验证方式，是使用
引理 \ref{sites-modules-lemma-module-principal-parts} 的证明中对模
$\mathcal{P}^k_{\mathcal{B}/\mathcal{A}}(\mathcal{M})$
用（局部）生成元与关系给出的描述；也可以直接从这些模的泛性质看出。
此外，这些映射与引理
\ref{sites-modules-lemma-sequence-of-principal-parts} 的短正合列相容。
\end{remark}








\section{朴素余切复形}
\label{sites-modules-section-netherlander}

\noindent
本节是《代数》，第 \ref{algebra-section-netherlander} 节
及《模》，第 \ref{modules-section-netherlander} 节的对应版本。
建议读者先阅读这些章节。

\medskip\noindent
设 $\mathcal{C}$ 为站点，$\mathcal{A} \to \mathcal{B}$
为 $\mathcal{C}$ 上的环层同态。本节中，对 $\mathcal{C}$
上的任意集合层 $\mathcal{E}$，记
$\mathcal{A}[\mathcal{E}]$ 为预层
$U \mapsto \mathcal{A}(U)[\mathcal{E}(U)]$ 的层化。这里
$\mathcal{A}(U)[\mathcal{E}(U)]$
表示 $\mathcal{A}(U)$ 上的多项式代数，其变量与
$\mathcal{E}(U)$ 的元素一一对应。以
$[e] \in \mathcal{A}(U)[\mathcal{E}(U)]$ 表示对应于
$e \in \mathcal{E}(U)$ 的变量。有典范的
$\mathcal{A}$-代数满射
\begin{equation}
\label{sites-modules-equation-canonical-presentation}
\mathcal{A}[\mathcal{B}] \longrightarrow \mathcal{B},\quad [b] \longmapsto b
\end{equation}
记其核为 $\mathcal{I} \subset \mathcal{A}[\mathcal{B}]$。
容易看出，$\mathcal{I}$ 由局部截面
$[b][b'] - [bb']$ 与 $[a] - a$ 生成。由引理
\ref{sites-modules-lemma-differential-seq}，有典范映射
\begin{equation}
\label{sites-modules-equation-naive-cotangent-complex}
\mathcal{I}/\mathcal{I}^2
\longrightarrow
\Omega_{\mathcal{A}[\mathcal{B}]/\mathcal{A}}
\otimes_{\mathcal{A}[\mathcal{B}]} \mathcal{B}
\end{equation}
其余核典范同构于 $\Omega_{\mathcal{B}/\mathcal{A}}$。

\begin{definition}
\label{sites-modules-definition-naive-cotangent-complex}
设 $\mathcal{C}$ 为站点，$\mathcal{A} \to \mathcal{B}$
为 $\mathcal{C}$ 上的环层同态。{\it 朴素余切复形}
$\NL_{\mathcal{B}/\mathcal{A}}$ 是链复形
(\ref{sites-modules-equation-naive-cotangent-complex})
$$
\NL_{\mathcal{B}/\mathcal{A}} =
\left(\mathcal{I}/\mathcal{I}^2
\longrightarrow
\Omega_{\mathcal{A}[\mathcal{B}]/\mathcal{A}}
\otimes_{\mathcal{A}[\mathcal{B}]} \mathcal{B}\right)
$$
其中 $\mathcal{I}/\mathcal{I}^2$ 置于次数 $-1$，而
$\Omega_{\mathcal{A}[\mathcal{B}]/\mathcal{A}}
\otimes_{\mathcal{A}[\mathcal{B}]} \mathcal{B}$
置于次数 $0$。
\end{definition}

\noindent
此构造满足与引理 \ref{sites-modules-lemma-functoriality-differentials}
中微分模的函子性类似的函子性。具体地，给定交换图
\begin{equation}
\label{sites-modules-equation-commutative-square-sheaves}
\vcenter{
\xymatrix{
\mathcal{B} \ar[r] & \mathcal{B}' \\
\mathcal{A} \ar[u] \ar[r] & \mathcal{A}' \ar[u]
}
}
\end{equation}
它是 $\mathcal{C}$ 上环层的图，则有典范的
$\mathcal{B}$-线性复形映射
$$
\NL_{\mathcal{B}/\mathcal{A}} \longrightarrow \NL_{\mathcal{B}'/\mathcal{A}'}
$$
事实上，交换图中的映射诱导典范映射
$\mathcal{A}[\mathcal{B}] \to \mathcal{A}'[\mathcal{B}']$
它把 $\mathcal{I}$ 映入
$\mathcal{I}' = \Ker(\mathcal{A}'[\mathcal{B}'] \to \mathcal{B}')$。
于是得到映射
$\mathcal{I}/\mathcal{I}^2 \to \mathcal{I}'/(\mathcal{I}')^2$
及微分模之间的映射；二者共同给出朴素余切复形之间所需的映射。

\medskip\noindent
也可以选取另一个把 $\mathcal{B}$ 表示为
$\mathcal{A}$ 上多项式代数之商的表示，而仍在 $D(\mathcal{B})$
中得到同一对象。为说明这一点，设 $\mathcal{E}$ 是
$\mathcal{C}$ 上的集合层，且
$\alpha : \mathcal{E} \to \mathcal{B}$ 是集合层的映射。
于是得到 $\mathcal{A}$-代数同态
$\mathcal{A}[\mathcal{E}] \to \mathcal{B}$。假设该映射满射，
并令 $\mathcal{J} \subset \mathcal{A}[\mathcal{E}]$ 为其核。置
$$
\NL(\alpha) = \left(
\mathcal{J}/\mathcal{J}^2
\longrightarrow
\Omega_{\mathcal{A}[\mathcal{E}]/\mathcal{A}}
\otimes_{\mathcal{A}[\mathcal{E}]} \mathcal{B}\right)
$$
结论如下。

\begin{lemma}
\label{sites-modules-lemma-NL-up-to-qis}
在上述情形下，$D(\mathcal{B})$ 中有典范同构
$\NL(\alpha) = \NL_{\mathcal{B}/\mathcal{A}}$。
\end{lemma}

\begin{proof}
注意 $\NL_{\mathcal{B}/\mathcal{A}} = \NL(\text{id}_\mathcal{B})$。
因此只需证明：给定上述两个映射
$\alpha_i : \mathcal{E}_i \to \mathcal{B}$，在 $D(\mathcal{B})$
中存在典范拟同构
$\NL(\alpha_1) = \NL(\alpha_2)$。为此，置
$\mathcal{E} = \mathcal{E}_1 \amalg \mathcal{E}_2$ 且
$\alpha = \alpha_1 \amalg \alpha_2 : \mathcal{E} \to \mathcal{B}$.
再置
$\mathcal{J}_i = \Ker(\mathcal{A}[\mathcal{E}_i] \to \mathcal{B})$
以及
$\mathcal{J} = \Ker(\mathcal{A}[\mathcal{E}] \to \mathcal{B})$.
得到把 $\mathcal{J}_i$ 映入 $\mathcal{J}$ 的映射
$\mathcal{A}[\mathcal{E}_i] \to \mathcal{A}[\mathcal{E}]$。
从而得到典范复形映射
$$
\NL(\alpha_i) \longrightarrow \NL(\alpha)
$$
只需证明这些映射是拟同构。为此作如下论证。首先，由引理
\ref{sites-modules-lemma-differential-seq}，
$H^0(\NL(\alpha_i)) = \Omega_{\mathcal{B}/\mathcal{A}}$ 且
$H^0(\NL(\alpha)) = \Omega_{\mathcal{B}/\mathcal{A}}$；
故该映射在次数 $0$ 的上同调层上是同构。类似地，我们断言
$H^{-1}(\NL(\alpha_i))$ 与 $H^{-1}(\NL(\alpha))$ 均为预层
$U \mapsto H_1(L_{\mathcal{B}(U)/\mathcal{A}(U)})$ 的关联层，其中
$H_1(L_{-/-})$ 如《代数》，定义
\ref{algebra-definition-naive-cotangent-complex}。若此断言成立，
证明便完成。

\medskip\noindent
证明该断言。设 $\alpha : \mathcal{E} \to \mathcal{B}$
如上。令 $\mathcal{B}' \subset \mathcal{B}$ 为
$\mathcal{A}$-代数子预层，它在 $U$ 上的取值是
$\mathcal{A}(U)[\mathcal{E}(U)] \to \mathcal{B}(U)$ 的像。
令 $\mathcal{I}'$ 为预层，它在 $U$ 上的取值是
$\mathcal{A}(U)[\mathcal{E}(U)] \to \mathcal{B}(U)$ 的核。
于是 $\mathcal{I}$ 是 $\mathcal{I}'$ 的层化，而
$\mathcal{B}$ 是 $\mathcal{B}'$ 的层化。类似地，
$H^{-1}(\NL(\alpha))$ 是下列预层的层化：
$$
U \longmapsto
\Ker(\mathcal{I}'(U)/\mathcal{I}'(U)^2 \to
\Omega_{\mathcal{A}(U)[\mathcal{E}(U)]/\mathcal{A}(U)}
\otimes_{\mathcal{A}(U)[\mathcal{E}(U)]} \mathcal{B}'(U))
$$
这是引理 \ref{sites-modules-lemma-differentials-sheafify} 的结论。
由《代数》，引理 \ref{algebra-lemma-NL-homotopy}，可知
$H^{-1}(\NL(\alpha))$ 是预层
$U \mapsto H_1(L_{\mathcal{B}'(U)/\mathcal{A}(U)})$ 的关联层。
因此只需证明映射
$H_1(L_{\mathcal{B}'(U)/\mathcal{A}(U)}) \to
H_1(L_{\mathcal{B}(U)/\mathcal{A}(U)})$ 在层化后诱导同构
$\mathcal{H}'_1 \to \mathcal{H}_1$。

\medskip\noindent
证明 $\mathcal{H}'_1 \to \mathcal{H}_1$ 的单射性。设
$f \in H_1(L_{\mathcal{B}'(U)/\mathcal{A}(U)})$ 在
$\mathcal{H}_1(U)$ 中映到零。须证：$f$ 在
$\mathcal{H}'_1(U)$ 中映到零。该假设意味着存在覆盖
$\{U_i \to U\}$，使得对所有 $i$，$f$ 在
$H_1(L_{\mathcal{B}(U_i)/\mathcal{A}(U_i)})$ 中映到零。
以 $U_i$ 替换 $U$，即可归约到 $f$ 在
$H_1(L_{\mathcal{B}(U)/\mathcal{A}(U)})$ 中映到零的情形。
由《代数》，引理 \ref{algebra-lemma-colimits-NL}，可找到有限生成子代数
$\mathcal{B}'(U) \subset B \subset \mathcal{B}(U)$，使得
$f$ 在 $H_1(L_{B/\mathcal{A}(U)})$ 中映到零。由于
$\mathcal{B} = (\mathcal{B}')^\#$，可找到覆盖
$\{U_i \to U\}$，使得 $B \to \mathcal{B}(U_i)$ 通过
$\mathcal{B}'(U_i)$ 分解。因此 $f$ 在
$H_1(L_{\mathcal{B}'(U_i)/\mathcal{A}(U_i)})$ 中映到零，如所需。

\medskip\noindent
$\mathcal{H}'_1 \to \mathcal{H}_1$ 的满射性完全类似地证明。
\end{proof}

\begin{lemma}
\label{sites-modules-lemma-pullback-NL}
设 $f : \Sh(\mathcal{C}) \to \Sh(\mathcal{D})$ 为拓扑斯态射，
$\mathcal{A} \to \mathcal{B}$ 为 $\mathcal{D}$ 上的环层同态。
则 $f^{-1}\NL_{\mathcal{B}/\mathcal{A}} =
\NL_{f^{-1}\mathcal{B}/f^{-1}\mathcal{A}}$.
\end{lemma}

\begin{proof}
从略。提示：使用引理 \ref{sites-modules-lemma-pullback-differentials}。
\end{proof}

\noindent
环拓扑斯态射的余切复形用上面定义的余切复形来定义。

\begin{definition}
\label{sites-modules-definition-cotangent-complex-morphism-ringed-topoi}
设 $X = (\Sh(\mathcal{C}), \mathcal{O})$ 与
$Y = (\Sh(\mathcal{C}'), \mathcal{O}')$ 为环拓扑斯，并设
$(f, f^\sharp) : X \to Y$ 为环拓扑斯态射。给定环拓扑斯态射的
{\it 朴素余切复形} $\NL_f = \NL_{X/Y}$ 定义为
$\NL_{\mathcal{O}/f^{-1}\mathcal{O}'}$。
有时也写成 $\NL_{X/Y} = \NL_{\mathcal{O}/\mathcal{O}'}$。
\end{definition}









\section{模的茎}
\label{sites-modules-section-stalks}

\noindent
在点处取茎时必须稍加小心，因为定义茎的余极限（见《站点》，公式
\ref{sites-equation-stalk}）未必是滤过的\footnote{当然，在几乎所有
自然出现的情形中，该余极限都是滤过的，因而本节的部分讨论可以简化。}。
另一方面，由站点之点的定义，茎函子正合且与任意余极限交换。
换言之，它的行为恰如该余极限是滤过的一样。

\begin{lemma}
\label{sites-modules-lemma-stalk-exact}
设 $\mathcal{C}$ 为站点，$p$ 为 $\mathcal{C}$ 的一个点。
\begin{enumerate}
\item 对 $\mathcal{C}$ 上任意集合预层，有
$(\mathcal{F}^\#)_p = \mathcal{F}_p$。
\item 茎函子
$\Sh(\mathcal{C}) \to \textit{Sets}$,
$\mathcal{F} \mapsto \mathcal{F}_p$ 正合（见《范畴》，定义
\ref{categories-definition-exact}），并与任意余极限交换。
\item 茎函子
$\textit{PSh}(\mathcal{C}) \to \textit{Sets}$,
$\mathcal{F} \mapsto \mathcal{F}_p$ 正合（见《范畴》，定义
\ref{categories-definition-exact}），并与任意余极限交换。
\end{enumerate}
\end{lemma}

\begin{proof}
由《站点》，引理 \ref{sites-lemma-point-pushforward-sheaf}
得到 (1)。由《站点》，引理
\ref{sites-lemma-adjoint-point-push-stalk} 可知
$\textit{PSh}(\mathcal{C}) \to \textit{Sets}$,
$\mathcal{F} \mapsto \mathcal{F}_p$ 是左伴随；由《站点》，引理
\ref{sites-lemma-point-pushforward-sheaf}，对
$\Sh(\mathcal{C}) \to \textit{Sets}$,
$\mathcal{F} \mapsto \mathcal{F}_p$ 也有相同结论。
故茎函子与任意余极限交换（见《范畴》，引理
\ref{categories-lemma-adjoint-exact}）。由站点之点的定义
（见《站点》，定义 \ref{sites-definition-point}），
$\Sh(\mathcal{C}) \to \textit{Sets}$,
$\mathcal{F} \mapsto \mathcal{F}_p$
正合。由于层化正合（《站点》，引理
\ref{sites-lemma-sheafification-exact}），可知
$\textit{PSh}(\mathcal{C}) \to \textit{Sets}$,
$\mathcal{F} \mapsto \mathcal{F}_p$
也正合。
\end{proof}

\noindent
特别地，由于预层上的茎函子
$\mathcal{F} \mapsto \mathcal{F}_p$ 与所有有限极限和余极限交换，
可以应用《站点》，命题
\ref{sites-proposition-functoriality-algebraic-structures-topoi}
证明中的论证。其结论是：若 $\mathcal{F}$ 是《站点》命题
\ref{sites-proposition-functoriality-algebraic-structures-topoi} 所列代数结构的
（预）层，则茎 $\mathcal{F}_p$ 自然成为同类代数结构。下面在
$\mathcal{F}$ 为阿贝尔预层时详细说明。此时加法映射
$+ : \mathcal{F} \times \mathcal{F} \to \mathcal{F}$ 诱导映射
$$
+ :
\mathcal{F}_p \times \mathcal{F}_p
=
(\mathcal{F} \times \mathcal{F})_p
\longrightarrow
\mathcal{F}_p
$$
其中等号用到了集合预层上的茎函子与有限极限交换。这在茎
$\mathcal{F}_p$ 上定义群结构。由此得到茎函子
$$
\textit{PAb}(\mathcal{C}) \longrightarrow \textit{Ab}, \quad
\mathcal{F} \longmapsto \mathcal{F}_p
$$
按构造，$\mathcal{F}_p$ 的底集合就是其底集合预层的茎。
再与包含
$\textit{Ab}(\mathcal{C}) \subset \textit{PAb}(\mathcal{C})$
预复合，也就定义了阿贝尔群层的茎函子。

\begin{lemma}
\label{sites-modules-lemma-stalk-exact-abelian}
设 $\mathcal{C}$ 为站点，$p$ 为 $\mathcal{C}$ 的一个点。
\begin{enumerate}
\item 函子
$\textit{Ab}(\mathcal{C}) \to \textit{Ab}$,
$\mathcal{F} \mapsto \mathcal{F}_p$ 正合。
\item 茎函子
$\textit{PAb}(\mathcal{C}) \to \textit{Ab}$,
$\mathcal{F}  \mapsto  \mathcal{F}_p$
正合。
\item 对 $\mathcal{F} \in \Ob(\textit{PAb}(\mathcal{C}))$，
有 $\mathcal{F}_p = \mathcal{F}^\#_p$。
\end{enumerate}
\end{lemma}

\begin{proof}
由引理 \ref{sites-modules-lemma-stalk-exact} 的结论和上述茎函子的构造形式地得到。
\end{proof}

\noindent
下面转向模层的情形。设 $(\mathcal{C}, \mathcal{O})$ 为环站点
（对以下讨论，只需 $\mathcal{O}$ 是环预层），
$\mathcal{F}$ 为 $\mathcal{O}$-模预层，$p$ 为
$\mathcal{C}$ 的一个点。此时得到映射
$$
\cdot :
\mathcal{O}_p \times \mathcal{O}_p
=
(\mathcal{O} \times \mathcal{O})_p
\longrightarrow
\mathcal{O}_p
$$
它是乘法映射的茎；以及
$$
\cdot :
\mathcal{O}_p \times \mathcal{F}_p
=
(\mathcal{O} \times \mathcal{F})_p
\longrightarrow
\mathcal{F}_p
$$
它也是乘法映射的茎。略去验证这在 $\mathcal{O}_p$
上定义环结构，并在 $\mathcal{F}_p$ 上定义
$\mathcal{O}_p$-模结构。由此得到函子
$$
\textit{PMod}(\mathcal{O}) \longrightarrow \textit{Mod}(\mathcal{O}_p), \quad
\mathcal{F} \longmapsto \mathcal{F}_p
$$
按构造，$\mathcal{F}_p$ 的底集合就是其底集合预层的茎。
再与包含
$\textit{Mod}(\mathcal{O}) \subset \textit{PMod}(\mathcal{O})$
预复合，也就定义了 $\mathcal{O}$-模层的茎函子。

\begin{lemma}
\label{sites-modules-lemma-stalk-exact-modules}
设 $(\mathcal{C}, \mathcal{O})$ 为环站点，
$p$ 为 $\mathcal{C}$ 的一个点。
\begin{enumerate}
\item 函子
$\textit{Mod}(\mathcal{O}) \to \textit{Mod}(\mathcal{O}_p)$,
$\mathcal{F} \mapsto \mathcal{F}_p$ 正合。
\item 茎函子
$\textit{PMod}(\mathcal{O}) \to \textit{Mod}(\mathcal{O}_p)$,
$\mathcal{F} \mapsto \mathcal{F}_p$
正合。
\item 对 $\mathcal{F} \in \Ob(\textit{PMod}(\mathcal{O}))$，
有 $\mathcal{F}_p = \mathcal{F}^\#_p$。
\end{enumerate}
\end{lemma}

\begin{proof}
由引理 \ref{sites-modules-lemma-stalk-exact-abelian} 的结论、
上述茎函子的构造及引理 \ref{sites-modules-lemma-abelian} 形式地得到。
\end{proof}

\begin{lemma}
\label{sites-modules-lemma-pullback-stalk}
设
$(f, f^\sharp) :
(\Sh(\mathcal{C}), \mathcal{O}_\mathcal{C})
\to
(\Sh(\mathcal{D}), \mathcal{O}_\mathcal{D})$
为环拓扑斯或环站点的态射。设 $p$ 为 $\mathcal{C}$
或 $\Sh(\mathcal{C})$ 的一个点，并置 $q = f \circ p$。则对任意
$\mathcal{O}_\mathcal{D}$-模 $\mathcal{F}$，有
$$
(f^*\mathcal{F})_p =
\mathcal{F}_q \otimes_{\mathcal{O}_{\mathcal{D}, q}}
\mathcal{O}_{\mathcal{C}, p}
$$
\end{lemma}

\begin{proof}
由定义，
$$
f^*\mathcal{F} =
f^{-1}\mathcal{F} \otimes_{f^{-1}\mathcal{O}_\mathcal{D}}
\mathcal{O}_\mathcal{C}
$$
由于在 $p$ 处取茎（即应用 $p^{-1}$）由引理
\ref{sites-modules-lemma-tensor-product-pullback} 与 $\otimes$ 交换，
再用《站点》，引理 \ref{sites-lemma-point-morphism-sites}
或引理 \ref{sites-lemma-point-morphism-topoi}
中所述的 $p$ 处拉回之茎与 $q$ 处茎的关系，即得结论。
\end{proof}






\section{摩天楼层}
\label{sites-modules-section-skyscraper}

\noindent
设 $p$ 为站点 $\mathcal{C}$ 或拓扑斯
$\Sh(\mathcal{C})$ 的一个点。本节研究把阿贝尔群 $A$
映到摩天楼层 $p_*A$ 的函子的正合性。首先回忆，
$p_* : \textit{Sets} \to \Sh(\mathcal{C})$ 具有许多正合性；
见《站点》，引理 \ref{sites-lemma-stalk-skyscraper} 及
\ref{sites-lemma-skyscraper-functor-exact}。

\begin{lemma}
\label{sites-modules-lemma-skyscraper-exact}
设 $\mathcal{C}$ 为站点，$p$ 为 $\mathcal{C}$
或其关联拓扑斯的一个点。
\begin{enumerate}
\item 函子 $p_* : \textit{Ab} \to \textit{Ab}(\mathcal{C})$,
$A \mapsto p_*A$ 正合。
\item 有函子性的直和分解
$$
p^{-1}p_*A = A \oplus I(A)
$$
其中 $A \in \Ob(\textit{Ab})$。
\end{enumerate}
\end{lemma}

\begin{proof}
由《站点》，引理 \ref{sites-lemma-stalk-skyscraper}，有函子性映射
$A \to p^{-1}p_*A \to A$，其复合等于 $\text{id}_A$。
故得到 (2) 中的函子性直和分解，其中 $I(A)$ 是伴随映射
$p^{-1}p_*A \to A$ 的核。由引理
\ref{sites-modules-lemma-exactness-pushforward-pullback}，函子 $p_*$ 左正合。
由《站点》，引理 \ref{sites-lemma-skyscraper-functor-exact}，
函子 $p_*$ 把满射变为满射。因此 (1) 成立。
\end{proof}

\noindent
为对模层作相同处理，设 $p$ 为环拓扑斯
$(\Sh(\mathcal{C}), \mathcal{O})$ 的一个点。回忆
$p^{-1}$ 就是茎函子。因此可以把 $p$ 看成环拓扑斯态射
$$
(p, \text{id}_{\mathcal{O}_p}) :
(\Sh(pt), \mathcal{O}_p)
\longrightarrow
(\Sh(\mathcal{C}), \mathcal{O}).
$$
于是得到拉回函子
$p^* : \textit{Mod}(\mathcal{O}) \to \textit{Mod}(\mathcal{O}_p)$
它等于茎函子，已在引理 \ref{sites-modules-lemma-stalk-exact-modules} 中讨论。
本节考虑函子
$p_* : \textit{Mod}(\mathcal{O}_p) \to \textit{Mod}(\mathcal{O})$.

\begin{lemma}
\label{sites-modules-lemma-skyscraper-modules-exact}
设 $(\Sh(\mathcal{C}), \mathcal{O})$ 为环拓扑斯，
$p$ 为拓扑斯 $\Sh(\mathcal{C})$ 的一个点。
\begin{enumerate}
\item 函子
$p_* : \textit{Mod}(\mathcal{O}_p) \to \textit{Mod}(\mathcal{O})$,
$M \mapsto p_*M$ 正合。
\item 典范满射 $p^{-1}p_*M \to M$ 是 $\mathcal{O}_p$-线性的。
\item 引理 \ref{sites-modules-lemma-skyscraper-exact} 的函子性直和分解
$p^{-1}p_*M = M \oplus I(M)$
一般而言{\bf 不是} $\mathcal{O}_p$-线性的。
\end{enumerate}
\end{lemma}

\begin{proof}
(1) 及 (2) 中的满射性立即由引理 \ref{sites-modules-lemma-skyscraper-exact}
中阿贝尔层的相应结论得到。由于
$p^{-1}\mathcal{O} = \mathcal{O}_p$，有 $p^{-1} = p^*$；
故 $p^{-1}p_*M \to M$ 就是模的伴随中的余单位
$p^*p_*M \to M$，因而是线性的。

\medskip\noindent
证明 (3)。设 $G$ 为群。考虑拓扑斯
$G\textit{-Sets} = \Sh(\mathcal{T}_G)$
及点 $p : \textit{Sets} \to G\textit{-Sets}$。见《站点》，第
\ref{sites-section-example-sheaf-G-sets} 节及例
\ref{sites-example-point-G-sets}。这里 $p^{-1}$
是遗忘 $G$-作用的函子，而 $p_*$ 是该遗忘函子的右伴随，
将 $M$ 映到 $\text{Map}(G, M)$。直和分解中的映射为
$$
M \to \text{Map}(G, M) \to M
$$
其中第一个映射把 $m \in M$ 映到值恒为 $m$ 的常值映射，
第二个映射是在 $G$ 的单位元 $1$ 处取值。再设 $R$
为配备 $G$-作用的环。这在 $\mathcal{T}_G$ 上确定环层
$\mathcal{O}$。$\mathcal{O}$-模范畴就是配备与 $R$
上的作用相容的 $G$-作用的 $R$-模 $M$ 所成范畴。
$\text{Map}(G, M)$ 上的 $R$-模结构由
$$
( r f ) (\sigma) = \sigma(r) f(\sigma)
$$
给出，其中 $r \in R$ 且 $f \in \text{Map}(G, M)$。这是因为它是
与在 $1$ 处取值相容的唯一 $G$-不变 $R$-模结构。容易看出，
一般而言 $M \to \text{Map}(G, M)$ 的像并不是 $R$-子模
（例如取 $M = R$ 并假设 $G$-作用非平凡），证明完成。
\end{proof}

\begin{example}
\label{sites-modules-example-ring-with-group-action}
设 $G$ 为群。考虑站点 $\mathcal{T}_G$ 及其点 $p$；
见《站点》，例 \ref{sites-example-point-G-sets}。设
$R$ 为配备 $G$-作用的环，它对应于 $\mathcal{T}_G$
上的环层 $\mathcal{O}$。则在遗忘 $G$-作用后，
$\mathcal{O}_p = R$。在此情形下
$p^{-1}p_*M = \text{Map}(G, M)$，且
$I(M) = \{f : G \to M \mid f(1_G) = 0\}$，而
$M \to \text{Map}(G, M)$ 把 $m \in M$ 映到值恒为 $m$ 的常值函数。
\end{example}




\section{局部化与点}
\label{sites-modules-section-localize-points}

\begin{lemma}
\label{sites-modules-lemma-stalk-j-shriek}
设 $(\mathcal{C}, \mathcal{O})$ 为环站点，
$p$ 为 $\mathcal{C}$ 的一个点，$U$ 为 $\mathcal{C}$ 的对象。
对 $\textit{Mod}(\mathcal{O}_U)$ 中的 $\mathcal{G}$，有
$$
(j_{U!}\mathcal{G})_p =
\bigoplus\nolimits_q \mathcal{G}_q
$$
其中余积遍历位于 $p$ 上方的 $\mathcal{C}/U$ 的各点 $q$；
见《站点》，引理 \ref{sites-lemma-points-above-point}。
\end{lemma}

\begin{proof}
使用引理 \ref{sites-modules-lemma-extension-by-zero} 中的描述：
$j_{U!}\mathcal{G}$ 是预层
$V \mapsto
\bigoplus\nolimits_{\varphi \in \Mor_\mathcal{C}(V, U)}
\mathcal{G}(V/_\varphi U)$
的关联层。由引理 \ref{sites-modules-lemma-stalk-exact-modules}，
$j_{U!}\mathcal{G}$ 在 $p$ 处的茎等于此预层的茎。
设 $u : \mathcal{C} \to \textit{Sets}$ 为对应于 $p$ 的函子
（见《站点》，第 \ref{sites-section-points} 节）。于是
$$
(j_{U!}\mathcal{G})_p = \colim_{(V, y)}
\bigoplus\nolimits_{\varphi : V \to U} \mathcal{G}(V/_\varphi U)
$$
其中余极限在阿贝尔群范畴中取。对该余极限中出现的四元组
$(V, y, \varphi, s)$，可以赋予
$x = u(\varphi)(y) \in u(U)$。因此得到
$$
(j_{U!}\mathcal{G})_p =
\bigoplus\nolimits_{x \in u(U)}
\colim_{(\varphi : V \to U, y), \ u(\varphi)(y) = x} \mathcal{G}(V/_\varphi U).
$$
由《站点》，引理 \ref{sites-lemma-points-above-point}
中对位于 $x$ 上方的点 $q$ 的描述，这等于引理中的表达式。
\end{proof}

\begin{remark}
\label{sites-modules-remark-not-pushforward}
警告：引理 \ref{sites-modules-lemma-stalk-j-shriek} 的结论对
$j_{U, *}$ 没有类似物。
\end{remark}












\section{平坦模的拉回}
\label{sites-modules-section-pullback-flat}

\noindent
平坦模沿环拓扑斯态射的拉回仍然平坦。证明这一点略有困难。

\begin{lemma}
\label{sites-modules-lemma-pullback-flat}
\begin{reference}
\cite[Expos\'e V, Corollary 1.7.1]{SGA4}
\end{reference}
设
$(f, f^\sharp) :
(\Sh(\mathcal{C}), \mathcal{O}_\mathcal{C})
\to
(\Sh(\mathcal{D}), \mathcal{O}_\mathcal{D})$
为环拓扑斯或环站点的态射。若 $\mathcal{F}$ 是平坦
$\mathcal{O}_\mathcal{D}$-模，则 $f^*\mathcal{F}$
是平坦 $\mathcal{O}_\mathcal{C}$-模。
\end{lemma}

\begin{proof}
选取引理
\ref{sites-modules-lemma-morphism-ringed-topoi-comes-from-morphism-ringed-sites}
中所述的图。回忆，成为平坦模是内蕴性质（见第
\ref{sites-modules-section-intrinsic} 节及定义 \ref{sites-modules-definition-flat}）。
故只需对态射 $(h, h^\sharp) :
(\Sh(\mathcal{C}'), \mathcal{O}_{\mathcal{C}'})
\to
(\Sh(\mathcal{D}'), \mathcal{O}_{\mathcal{D}'})$ 证明引理。
换言之，可以假设站点 $\mathcal{C}$ 与 $\mathcal{D}$
具有所有有限极限，并且 $f$ 是由与有限极限交换的连续函子
$u : \mathcal{D} \to \mathcal{C}$ 诱导的站点态射。

\medskip\noindent
回忆 $f^*\mathcal{F} =
\mathcal{O}_\mathcal{C} \otimes_{f^{-1}\mathcal{O}_\mathcal{D}}
f^{-1}\mathcal{F}$（定义 \ref{sites-modules-definition-pushforward}）。
由引理 \ref{sites-modules-lemma-flat-change-of-rings}，只需证明
$f^{-1}\mathcal{F}$ 是平坦
$f^{-1}\mathcal{O}_\mathcal{D}$-模。结合上一段，问题归约到下一段的情形。

\medskip\noindent
假设 $\mathcal{C}$ 与 $\mathcal{D}$ 是具有所有有限极限的站点，
且 $u : \mathcal{D} \to \mathcal{C}$ 是与有限极限交换的连续函子。
设 $\mathcal{O}$ 为 $\mathcal{D}$ 上的环层，
$\mathcal{F}$ 为平坦 $\mathcal{O}$-模。则 $u$
定义站点态射 $f : \mathcal{C} \to \mathcal{D}$
（《站点》，命题 \ref{sites-proposition-get-morphism}）。
须证：$f^{-1}\mathcal{F}$ 是平坦
$f^{-1}\mathcal{O}$-模。设 $U$ 为 $\mathcal{C}$ 的对象，并设
$$
f^{-1}\mathcal{O}|_U \xrightarrow{(f_1, \ldots, f_n)}
f^{-1}\mathcal{O}|_U^{\oplus n} \xrightarrow{(s_1, \ldots, s_n)}
f^{-1}\mathcal{F}|_U
$$
为 $f^{-1}\mathcal{O}|_U$-模的复形。目标是在 $U$
的某个覆盖的各成员上，构造 $(s_1, \ldots, s_n)$ 的分解，
如引理 \ref{sites-modules-lemma-flat-eq} 的 (2) 所述。考虑元素
$s_a \in f^{-1}\mathcal{F}(U)$ 与
$f_a \in f^{-1}\mathcal{O}(U)$。由于
$f^{-1}\mathcal{F}$（相应地 $f^{-1}\mathcal{O}$）是
$u_p\mathcal{F}$ 的层化，将 $U$ 替换为某个覆盖的各成员后，
可以假设 $s_a$ 是某个 $s'_a \in u_p\mathcal{F}(U)$ 的像，
而 $f_a$ 是某个 $f'_a \in u_p\mathcal{O}(U)$ 的像。再次将
$U$ 替换为某个覆盖的各成员后，可以假设
$\sum f'_as'_a$ 在 $u_p\mathcal{F}(U)$ 中为零。回忆范畴
$(\mathcal{I}_U^u)^{opp}$ 是有向的（《站点》，引理
\ref{sites-lemma-directed}），并且
$u_p\mathcal{F}(U) = \colim_{(\mathcal{I}_U^u)^{opp}} \mathcal{F}(V)$，
$u_p\mathcal{O}(U) = \colim_{(\mathcal{I}_U^u)^{opp}} \mathcal{O}(V)$。
因此可以假设存在
$(V, \phi) \in \Ob(\mathcal{I}_U^u)$，其中 $V$ 是
$\mathcal{D}$ 的对象，$\phi$ 是 $\mathcal{D}$ 中的态射
$\phi : U \to u(V)$，并且存在元素
$s''_a \in \mathcal{F}(V)$ 与 $f''_a \in \mathcal{O}(V)$，
它们在 $u_p\mathcal{F}(U)$ 与 $u_p\mathcal{O}(U)$ 中的像
分别等于 $s'_a$ 与 $f'_a$，且
$\sum f''_a s''_a = 0$ 在 $\mathcal{F}(V)$ 中成立。
于是得到复形
$$
\mathcal{O}|_V \xrightarrow{(f''_1, \ldots, f''_n)}
\mathcal{O}|_V^{\oplus n} \xrightarrow{(s''_1, \ldots, s''_n)}
\mathcal{F}|_V
$$
应用引理 \ref{sites-modules-lemma-flat-eq} 的另一方向，可知存在
$\mathcal{D}$ 中的覆盖 $\{V_i \to V\}$，并且对每个 $i$
都有分解
$$
\mathcal{O}|_{V_i}^{\oplus n}
\xrightarrow{B''_i}
\mathcal{O}|_{V_i}^{\oplus l_i} \xrightarrow{(t''_{i1}, \ldots, t''_{il_i})}
\mathcal{F}|_{V_i}
$$
它是 $(s''_1, \ldots, s''_n)|_{V_i}$ 的分解，且
$B_i \circ (f''_1, \ldots, f''_n)|_{V_i} = 0$.
置 $U_i = U \times_{\phi, u(V)} u(V_i)$，记
$B_i \in \text{Mat}(l_i \times n, f^{-1}\mathcal{O}(U_i))$
为 $B''_i$ 的像，并记
$t_{ij} \in f^{-1}\mathcal{F}(U_i)$ 为 $t''_{ij}$ 的像。
于是得到分解
$$
f^{-1}\mathcal{O}|_{U_i}^{\oplus n}
\xrightarrow{B_i}
f^{-1}\mathcal{O}|_{U_i}^{\oplus l_i}
\xrightarrow{(t_{i1}, \ldots, t_{il_i})}
\mathcal{F}|_{U_i}
$$
它是 $(s_1, \ldots, s_n)|_{U_i}$ 的分解，且
$B_i \circ (f_1, \ldots, f_n)|_{U_i} = 0$.
证明完成。
\end{proof}

\begin{lemma}
\label{sites-modules-lemma-stalk-flat}
设 $(\mathcal{C}, \mathcal{O})$ 为环站点，
$p$ 为 $\mathcal{C}$ 的一个点。若 $\mathcal{F}$
是平坦 $\mathcal{O}$-模，则 $\mathcal{F}_p$
是平坦 $\mathcal{O}_p$-模。
\end{lemma}

\begin{proof}
第 \ref{sites-modules-section-skyscraper} 节表明，可以把 $p$ 看成环拓扑斯态射
$$
(p, \text{id}_{\mathcal{O}_p}) :
(\Sh(pt), \mathcal{O}_p)
\longrightarrow
(\Sh(\mathcal{C}), \mathcal{O}).
$$
使得拉回函子
$p^* : \textit{Mod}(\mathcal{O}) \to \textit{Mod}(\mathcal{O}_p)$
等于茎函子。因此结论由引理 \ref{sites-modules-lemma-pullback-flat} 得到。
\end{proof}

\begin{lemma}
\label{sites-modules-lemma-check-flat-stalks}
设 $(\mathcal{C}, \mathcal{O})$ 为环站点，
$\mathcal{F}$ 为 $\mathcal{O}$-模层，并设
$\{p_i\}_{i \in I}$ 为 $\mathcal{C}$ 的一族保守点。
则 $\mathcal{F}$ 平坦，当且仅当对所有 $i \in I$，
$\mathcal{F}_{p_i}$ 是平坦 $\mathcal{O}_{p_i}$-模。
\end{lemma}

\begin{proof}
一个方向由引理 \ref{sites-modules-lemma-stalk-flat} 得到。反之，使用引理
\ref{sites-modules-lemma-tensor-product-pullback}（在 $p$ 处取茎由
$p^{-1}$ 给出）中的等式
$(\mathcal{F} \otimes_\mathcal{O} \mathcal{G})_p =
\mathcal{F}_p \otimes_{\mathcal{O}_p} \mathcal{G}_p$，
再应用引理 \ref{sites-modules-lemma-check-exactness-stalks}。
\end{proof}

\begin{lemma}
\label{sites-modules-lemma-pullback-ses}
设
$f : (\Sh(\mathcal{C}'), \mathcal{O}') \to (\Sh(\mathcal{C}), \mathcal{O})$
为环拓扑斯态射。设
$0 \to \mathcal{F} \to \mathcal{G} \to \mathcal{H} \to 0$
为 $\mathcal{O}$-模的短正合列，且 $\mathcal{H}$
是平坦 $\mathcal{O}$-模。则序列
$0 \to f^*\mathcal{F} \to f^*\mathcal{G} \to f^*\mathcal{H} \to 0$
也正合。
\end{lemma}

\begin{proof}
由于 $f^{-1}$ 正合，有 $f^{-1}\mathcal{O}$-模的短正合列
$0 \to f^{-1}\mathcal{F} \to f^{-1}\mathcal{G} \to f^{-1}\mathcal{H} \to 0$
由引理 \ref{sites-modules-lemma-pullback-flat}，
$f^{-1}\mathcal{O}$-模 $f^{-1}\mathcal{H}$ 平坦。由引理
\ref{sites-modules-lemma-flat-tor-zero}，这蕴含该序列在
$f^{-1}\mathcal{O}$ 上与 $\mathcal{O}'$ 作张量积后仍正合。由于
$f^*\mathcal{F} = f^{-1}\mathcal{F} \otimes_{f^{-1}\mathcal{O}} \mathcal{O}'$
，而 $\mathcal{G}$ 与 $\mathcal{H}$ 也同样，故得结论。
\end{proof}





\section{局部环拓扑斯}
\label{sites-modules-section-locally-ringed}


\noindent
本节的参考文献为 \cite[Expos\'e IV, Exercice 13.9]{SGA4}。

\begin{lemma}
\label{sites-modules-lemma-locally-ringed}
设 $(\mathcal{C}, \mathcal{O})$ 为环站点。下列条件等价：
\begin{enumerate}
\item 对 $\mathcal{C}$ 的每个对象 $U$ 及
$f \in \mathcal{O}(U)$，存在覆盖 $\{U_j \to U\}$，使得对每个
$j$，$f|_{U_j}$ 或 $(1 - f)|_{U_j}$ 可逆。
\item 若 $U \in \Ob(\mathcal{C})$、$n \geq 1$，且
$f_1, \ldots, f_n \in \mathcal{O}(U)$ 在 $\mathcal{O}(U)$
中生成单位理想，则存在覆盖 $\{U_j \to U\}$，使得对每个 $j$
都存在 $i$，使 $f_i|_{U_j}$ 可逆。
\item 集合层的映射
$$
(\mathcal{O} \times \mathcal{O})
\amalg
(\mathcal{O} \times \mathcal{O})
\longrightarrow
\mathcal{O} \times \mathcal{O}
$$
把第一分支中的 $(f, a)$ 映到 $(f, af)$，并把第二分支中的
$(f, b)$ 映到 $(f, b(1 - f))$；此映射是满射。
\end{enumerate}
\end{lemma}

\begin{proof}
(2) 显然蕴含 (1)。为证明 (1) 蕴含 (2)，对 $n$ 作归纳。
第一个非平凡情形是 $n = 2$（$n = 1$ 显然）。此时存在
$a_1, a_2 \in \mathcal{O}(U)$，使
$a_1f_1 + a_2f_2 = 1$。由假设，可找到覆盖
$\{U_j \to U\}$，使得对每个 $j$，
$a_1f_1|_{U_j}$ 或 $a_2f_2|_{U_j}$ 可逆。因此
$f_1|_{U_j}$ 或 $f_2|_{U_j}$ 可逆，如所需。当 $n > 2$ 时，
存在 $a_1, \ldots, a_n \in \mathcal{O}(U)$，使
$a_1f_1 + \ldots + a_nf_n = 1$。由 $n = 2$ 的情形，存在覆盖
$\{U_j \to U\}_{j \in J}$，使得对每个 $j$，
$f_n|_{U_j}$ 或
$a_1f_1 + \ldots + a_{n - 1}f_{n - 1}|_{U_j}$ 可逆。
设第一种情形恰发生于 $j \in J_n$，并置
$J' = J \setminus J_n$。由归纳假设，对每个 $j \in J'$
可找到覆盖 $\{U_{jk} \to U_j\}_{k \in K_j}$，使得对每个
$k \in K_j$ 都存在 $i \in \{1, \ldots, n - 1\}$，使
$f_i|_{U_{jk}}$ 可逆。由站点公理，态射族
$\{U_j \to U\}_{j \in J_n} \cup \{U_{jk} \to U\}_{j \in J', k \in K_j}$
是具有所需性质的覆盖。

\medskip\noindent
假设 (1)。为证 (3) 中的映射满射，设 $(f, c)$ 是
$\mathcal{O} \times \mathcal{O}$ 在 $U$ 上的截面。由假设，
存在覆盖 $\{U_j \to U\}$，使得对每个 $j$，$f$ 或 $1-f$
的限制可逆。第一种情形取
$a = c|_{U_j} (f|_{U_j})^{-1}$，第二种情形取
$b = c|_{U_j} (1 - f|_{U_j})^{-1}$。故在每个覆盖成员上，
$(f,c)$ 都属于该映射的像。反之，假设 (3) 成立。对任意
$U$ 及 $f \in \mathcal{O}(U)$，存在 $U$ 的覆盖
$\{U_j \to U\}$，使截面 $(f,1)|_{U_j}$ 属于 (3)
中的映射在 $U_j$ 上截面的像。这恰好意味着 $f$ 或 $1-f$
在 $U_j$ 上的限制可逆，故 (1) 成立。
\end{proof}

\begin{lemma}
\label{sites-modules-lemma-locally-ringed-stalk}
设 $(\mathcal{C}, \mathcal{O})$ 为环站点。考虑下列条件：
\begin{enumerate}
\item 对 $\mathcal{C}$ 的每个对象 $U$ 及
$f \in \mathcal{O}(U)$，存在覆盖 $\{U_j \to U\}$，使得对每个
$j$，$f|_{U_j}$ 或 $(1-f)|_{U_j}$ 可逆。
\item 对 $\mathcal{C}$ 的每个点 $p$，茎 $\mathcal{O}_p$
是零环或局部环。
\end{enumerate}
总有 (1) $\Rightarrow$ (2)。若 $\mathcal{C}$ 有足够多的点，
则 (1) 与 (2) 等价。
\end{lemma}

\begin{proof}
假设 (1)。设 $p$ 是由函子
$u : \mathcal{C} \to \textit{Sets}$ 给出的 $\mathcal{C}$ 的点，
并取 $f_p \in \mathcal{O}_p$。由于 $\mathcal{O}_p$
由《站点》，公式 (\ref{sites-equation-stalk}) 计算，可以用三元组
$(U,x,f)$ 表示 $f_p$，其中 $x \in U(U)$ 且
$f \in \mathcal{O}(U)$。由假设，存在覆盖
$\{U_i \to U\}$，使得对每个 $i$，$f$ 或 $1-f$
在 $U_i$ 上可逆。由于 $u$ 定义站点的点，对某个 $i$
存在 $x_i \in u(U_i)$ 映到 $x \in u(U)$。由《站点》，公式
(\ref{sites-equation-stalk}) 附近的讨论，
$(U,x,f)$ 与 $(U_i,x_i,f|_{U_i})$ 定义
$\mathcal{O}_p$ 中的同一元素。因此 $f_p$ 或 $1-f_p$ 可逆。
于是环 $\mathcal{O}_p$ 的每个元素 $a$ 都满足
$a$ 或 $1-a$ 可逆。这意味着 $\mathcal{O}_p$ 是零环或局部环；
见《代数》，引理 \ref{algebra-lemma-characterize-local-ring}。

\medskip\noindent
假设 (2)，并假设 $\mathcal{C}$ 有足够多的点。考虑集合层的映射
$$
\mathcal{O} \times \mathcal{O} \amalg \mathcal{O} \times \mathcal{O}
\longrightarrow
\mathcal{O} \times \mathcal{O}
$$
即引理 \ref{sites-modules-lemma-locally-ringed} 的 (3) 中的映射。
对任意局部环 $R$，相应映射
$(R \times R) \amalg (R \times R) \to R \times R$
是满射；例如见《代数》，引理
\ref{algebra-lemma-characterize-local-ring}。由于每个
$\mathcal{O}_p$ 是局部环或零环，该映射在茎上满射。
由 $\mathcal{C}$ 有足够多的点的假设，该映射满射，结论得证。
\end{proof}

\noindent
《模》，第 \ref{modules-section-pathology} 节指出，在环化空间
$(X,\mathcal{O}_X)$ 中可能存在一个开子空间，使结构层在其上为零。
为排除这种情形，可以要求截面 $1$ 与 $0$ 在空间 $X$
的每个茎中取不同值。在环拓扑斯与环站点的语境下，该条件是
\begin{equation}
\label{sites-modules-equation-one-is-never-zero}
\emptyset^\# \longrightarrow
\text{Equalizer}(0, 1 : * \longrightarrow \mathcal{O})
\end{equation}
是层的同构。这里 $*$ 是单元素层，而 $\emptyset^\#$ 是“空层”，
即分别为层范畴的终对象与始对象；见《站点》，例
\ref{sites-example-singleton-sheaf} 及第
\ref{sites-section-almost-cocontinuous} 节。换言之，该条件是：
只要 $U \in \Ob(\mathcal{C})$ 在层论意义下不为空，
就有 $1,0 \in \mathcal{O}(U)$ 不相等。下面陈述必需的引理。

\begin{lemma}
\label{sites-modules-lemma-ringed-stalk-not-zero}
设 $(\mathcal{C}, \mathcal{O})$ 为环站点。考虑下列命题：
\begin{enumerate}
\item (\ref{sites-modules-equation-one-is-never-zero}) 是同构；
\item 对 $\mathcal{C}$ 的每个点 $p$，茎 $\mathcal{O}_p$ 不是零环。
\end{enumerate}
总有 (1) $\Rightarrow$ (2)；若 $\mathcal{C}$ 有足够多的点，
则 (1) $\Leftrightarrow$ (2)。
\end{lemma}

\begin{proof}
从略。
\end{proof}

\noindent
引理 \ref{sites-modules-lemma-locally-ringed}、
\ref{sites-modules-lemma-locally-ringed-stalk} 及
\ref{sites-modules-lemma-ringed-stalk-not-zero} 引出下列定义。

\begin{definition}
\label{sites-modules-definition-locally-ringed}
若 (\ref{sites-modules-equation-one-is-never-zero}) 是同构，且引理
\ref{sites-modules-lemma-locally-ringed} 中的等价性质成立，则称环站点
$(\mathcal{C}, \mathcal{O})$ 为{\it 局部环站点}。
\end{definition}

\noindent
\cite[Expos\'e IV, Exercice 13.9]{SGA4} 中缺少
(\ref{sites-modules-equation-one-is-never-zero}) 为同构这一条件，因而其中的
局部环站点与局部环拓扑斯概念略有不同。由于我们的动机来自
局部环化空间的概念，故加入了这一条件（见上文解释）。

\begin{lemma}
\label{sites-modules-lemma-locally-ringed-intrinsic}
成为局部环站点是内蕴性质。更准确地：
\begin{enumerate}
\item 若 $f : \Sh(\mathcal{C}') \to \Sh(\mathcal{C})$
为拓扑斯态射，且 $(\mathcal{C}, \mathcal{O})$
为局部环站点，则 $(\mathcal{C}', f^{-1}\mathcal{O})$
为局部环站点；
\item 若
$(f, f^\sharp) : (\Sh(\mathcal{C}'), \mathcal{O}')
\to (\Sh(\mathcal{C}), \mathcal{O})$
为环拓扑斯的等价，则 $(\mathcal{C}, \mathcal{O})$
局部环，当且仅当
$(\mathcal{C}', \mathcal{O}')$
局部环。
\end{enumerate}
\end{lemma}

\begin{proof}
(2) 显然由 (1) 得到。为证明 (1)，注意 $f^{-1}$ 正合，故
$f^{-1}* = *$、$f^{-1}\emptyset^\# = \emptyset^\#$，
且 $f^{-1}$ 与积及等化子交换，并把同构与满射分别变为同构与满射。
因此 $f^{-1}$ 把同构 (\ref{sites-modules-equation-one-is-never-zero})
变为 $f^{-1}\mathcal{O}$ 的相应同构，并把引理
\ref{sites-modules-lemma-locally-ringed} 的 (3) 中的满射变为
$f^{-1}\mathcal{O}$ 的相应满射。
\end{proof}

\noindent
事实上，引理 \ref{sites-modules-lemma-locally-ringed-intrinsic} 的 (2)
是《概形》，引理 \ref{schemes-lemma-isomorphism-locally-ringed}
的对应版本。它保证下列定义有意义。

\begin{definition}
\label{sites-modules-definition-locally-ringed-topos}
若环拓扑斯 $(\Sh(\mathcal{C}), \mathcal{O})$ 所基于的环站点
$(\mathcal{C}, \mathcal{O})$ 局部环，则称该环拓扑斯
{\it 局部环}。
\end{definition}

\noindent
下面给出局部环性质的一个推论。

\begin{lemma}
\label{sites-modules-lemma-invertible-is-locally-free-rank-1}
设 $(\Sh(\mathcal{C}), \mathcal{O})$ 为环拓扑斯。
任何秩为 $1$ 的局部自由 $\mathcal{O}$-模都可逆。
若 $(\mathcal{C}, \mathcal{O})$ 局部环，则逆命题也成立
（一般情形下则不成立）。
\end{lemma}

\begin{proof}
假设 $\mathcal{L}$ 是秩为 $1$ 的局部自由模，并考虑求值映射
$$
\mathcal{L} \otimes_\mathcal{O}
\SheafHom_\mathcal{O}(\mathcal{L}, \mathcal{O})
\longrightarrow \mathcal{O}
$$
对 $\mathcal{C}$ 的任意对象 $U$，限制到使 $\mathcal{L}$
平凡化的某个覆盖的各成员上，可见此映射是同构（细节从略）。
因此由引理 \ref{sites-modules-lemma-invertible}，$\mathcal{L}$ 可逆。

\medskip\noindent
假设 $(\Sh(\mathcal{C}), \mathcal{O})$ 局部环，并设
$U$ 为 $\mathcal{C}$ 的对象。引理 \ref{sites-modules-lemma-invertible}
的证明表明，存在覆盖 $\{U_i \to U\}$，使
$\mathcal{L}|_{\mathcal{C}/U_i}$ 是某个有限自由
$\mathcal{O}_{U_i}$-模的直和因子。以 $U_i$ 替换 $U$ 后，设
$p : \mathcal{O}_U^{\oplus r} \to \mathcal{O}_U^{\oplus r}$
为投影算子，其像同构于 $\mathcal{L}|_{\mathcal{C}/U}$。
则 $p$ 对应于矩阵
$$
P = (p_{ij}) \in \text{Mat}(r \times r, \mathcal{O}(U))
$$
它是投影矩阵：$P^2 = P$。置 $A = \mathcal{O}(U)$，
于是 $P \in \text{Mat}(r \times r, A)$。由《代数》，引理
\ref{algebra-lemma-finite-projective}，$P$ 的像是
$A$ 上的有限局部自由模 $M$。故存在生成单位理想的
$f_1,\ldots,f_t \in A$，使 $M_{f_i}$ 有限自由。由引理
\ref{sites-modules-lemma-locally-ringed}，将 $U$ 替换为某个开覆盖的各成员后，
可以假设 $M$ 自由。这意味着 $\mathcal{L}|_U$ 自由
（细节从略）。当然，由于 $\mathcal{L}$ 可逆，只有当
$\mathcal{L}|_U$ 的秩为 $1$ 时才可能如此，证明完成。
\end{proof}

\noindent
下面说明局部环化空间之间的态射意味着什么。为此需要下列引理。

\begin{lemma}
\label{sites-modules-lemma-locally-ringed-morphism}
设
$(f, f^\sharp) : (\Sh(\mathcal{C}), \mathcal{O}_\mathcal{C})
\to (\Sh(\mathcal{D}), \mathcal{O}_\mathcal{D})$
为环拓扑斯态射。考虑下列条件：
\begin{enumerate}
\item 层的图
$$
\xymatrix{
f^{-1}(\mathcal{O}^*_\mathcal{D}) \ar[r]_-{f^\sharp} \ar[d] &
\mathcal{O}^*_\mathcal{C} \ar[d] \\
f^{-1}(\mathcal{O}_\mathcal{D}) \ar[r]^-{f^\sharp} &
\mathcal{O}_\mathcal{C}
}
$$
是笛卡尔图。
\item 对 $\mathcal{C}$ 的任意点 $p$，置 $q = f \circ p$，则集合的图
$$
\xymatrix{
\mathcal{O}^*_{\mathcal{D}, q} \ar[r] \ar[d] &
\mathcal{O}^*_{\mathcal{C}, p} \ar[d] \\
\mathcal{O}_{\mathcal{D}, q} \ar[r] &
\mathcal{O}_{\mathcal{C}, p}
}
$$
是笛卡尔图。
\end{enumerate}
总有 (1) $\Rightarrow$ (2)。若 $\mathcal{C}$ 有足够多的点，
则 (1) 与 (2) 等价。若
$(\Sh(\mathcal{C}), \mathcal{O}_\mathcal{C})$
与
$(\Sh(\mathcal{D}), \mathcal{O}_\mathcal{D})$
均为局部环拓扑斯，则 (2) 等价于
\begin{enumerate}
\item[(3)] 对 $\mathcal{C}$ 的任意点 $p$，置 $q = f \circ p$，
则环映射
$\mathcal{O}_{\mathcal{D}, q} \to \mathcal{O}_{\mathcal{C}, p}$
是局部环同态。
\end{enumerate}
事实上，若对一族保守点有性质 (2) 或 (3)，则 (1) 成立。
\end{lemma}

\begin{proof}
展开各定义即可证明。
\end{proof}

\begin{definition}
\label{sites-modules-definition-morphism-locally-ringed-topoi}
设 $(f, f^\sharp) : (\Sh(\mathcal{C}), \mathcal{O}_\mathcal{C})
\to (\Sh(\mathcal{D}), \mathcal{O}_\mathcal{D})$
为环拓扑斯态射。假设
$(\Sh(\mathcal{C}), \mathcal{O}_\mathcal{C})$
与
$(\Sh(\mathcal{D}), \mathcal{O}_\mathcal{D})$
均为局部环拓扑斯。若且唯若层的图
$$
\xymatrix{
f^{-1}(\mathcal{O}^*_\mathcal{D}) \ar[r]_-{f^\sharp} \ar[d] &
\mathcal{O}^*_\mathcal{C} \ar[d] \\
f^{-1}(\mathcal{O}_\mathcal{D}) \ar[r]^-{f^\sharp} &
\mathcal{O}_\mathcal{C}
}
$$
是笛卡尔图（见引理 \ref{sites-modules-lemma-locally-ringed-morphism}），
则称 $(f,f^\sharp)$ 为{\it 局部环拓扑斯态射}。
若 $(f,f^\sharp)$ 是环站点态射，并且其关联的环拓扑斯态射
是局部环拓扑斯态射，则称其为{\it 局部环站点态射}。
\end{definition}

\noindent
显然，局部环拓扑斯之间的环拓扑斯同构自动是局部环拓扑斯同构。

\begin{lemma}
\label{sites-modules-lemma-composition-morphisms-locally-ringed-topoi}
设
$(f, f^\sharp) :
(\Sh(\mathcal{C}_1), \mathcal{O}_1)
\to (\Sh(\mathcal{C}_2), \mathcal{O}_2)$ and
$(g, g^\sharp) :
(\Sh(\mathcal{C}_2), \mathcal{O}_2) \to
(\Sh(\mathcal{C}_3), \mathcal{O}_3)$
为局部环拓扑斯态射。则复合
$(g, g^\sharp) \circ (f, f^\sharp)$（见定义
\ref{sites-modules-definition-ringed-topos}）也是局部环拓扑斯态射。
\end{lemma}

\begin{proof}
从略。
\end{proof}

\begin{lemma}
\label{sites-modules-lemma-locally-ringed-intrinsic-morphism}
设 $f : \Sh(\mathcal{C}') \to \Sh(\mathcal{C})$
为拓扑斯态射。若 $\mathcal{O}$ 是 $\mathcal{C}$ 上的环层，则
$$
f^{-1}(\mathcal{O}^*) = (f^{-1}\mathcal{O})^*.
$$
特别地，若 $\mathcal{O}$ 使 $\mathcal{C}$ 成为局部环站点，
则置 $f^\sharp = \text{id}$ 后，环拓扑斯态射
$$
(f, f^\sharp) :
(\Sh(\mathcal{C}'), f^{-1}\mathcal{O})
\to
(\Sh(\mathcal{C}, \mathcal{O})
$$
是局部环拓扑斯态射。
\end{lemma}

\begin{proof}
注意，图
$$
\xymatrix{
\mathcal{O}^* \ar[rr] \ar[d]_{u \mapsto (u, u^{-1})} & &
{*} \ar[d]^{1} \\
\mathcal{O} \times \mathcal{O} \ar[rr]^-{(a, b) \mapsto ab} & &
\mathcal{O}
}
$$
是笛卡尔图。由于 $f^{-1}$ 正合，图
$$
\xymatrix{
f^{-1}(\mathcal{O}^*)
\ar[d]_{u \mapsto (u, u^{-1})} \ar[rr] & &
{*} \ar[d]^{1} \\
f^{-1}\mathcal{O} \times f^{-1}\mathcal{O} \ar[rr]^-{(a, b) \mapsto ab} & &
f^{-1}\mathcal{O}
}
$$
也是笛卡尔图，这蕴含第一个断言。对第二个断言，注意由引理
\ref{sites-modules-lemma-locally-ringed-intrinsic}，
$(\mathcal{C}', f^{-1}\mathcal{O})$ 是局部环站点，故该断言有意义。
此时第一部分蕴含该态射是局部环拓扑斯态射。
\end{proof}

\begin{lemma}
\label{sites-modules-lemma-localize-morphism-locally-ringed-topoi}
局部环站点与拓扑斯的局部化。
\begin{enumerate}
\item 设 $(\mathcal{C}, \mathcal{O})$ 为局部环站点，
$U$ 为 $\mathcal{C}$ 的对象。则局部化
$(\mathcal{C}/U, \mathcal{O}_U)$ 是局部环站点，且局部化态射
$$
(j_U, j_U^\sharp) :
(\Sh(\mathcal{C}/U), \mathcal{O}_U)
\to
(\Sh(\mathcal{C}), \mathcal{O})
$$
是局部环拓扑斯态射。
\item 设 $(\mathcal{C}, \mathcal{O})$ 为局部环站点，
$f : V \to U$ 为 $\mathcal{C}$ 中的态射。则引理
\ref{sites-modules-lemma-relocalize} 中的态射
$$
(j, j^\sharp) :
(\Sh(\mathcal{C}/V), \mathcal{O}_V)
\to
(\Sh(\mathcal{C}/U), \mathcal{O}_U)
$$
是局部环拓扑斯态射。
\item 设
$(f, f^\sharp) :
(\mathcal{C}, \mathcal{O})
\longrightarrow
(\mathcal{D}, \mathcal{O}')$
为局部环站点态射，其中 $f$ 由连续函子
$u : \mathcal{D} \to \mathcal{C}$ 给出。设 $V$ 为
$\mathcal{D}$ 的对象，并令 $U = u(V)$。则引理
\ref{sites-modules-lemma-localize-morphism-ringed-sites} 中的态射
$$
(f', (f')^\sharp) :
(\Sh(\mathcal{C}/U), \mathcal{O}_U)
\to
(\Sh(\mathcal{D}/V), \mathcal{O}'_V)
$$
是局部环站点态射。
\item 设
$(f, f^\sharp) :
(\mathcal{C}, \mathcal{O})
\longrightarrow
(\mathcal{D}, \mathcal{O}')$
为局部环站点态射，其中 $f$ 由连续函子
$u : \mathcal{D} \to \mathcal{C}$ 给出。设
$V \in \Ob(\mathcal{D})$、$U \in \Ob(\mathcal{C})$，且
$c : U \to u(V)$。则引理
\ref{sites-modules-lemma-relocalize-morphism-ringed-sites} 中的态射
$$
(f_c, (f_c)^\sharp) :
(\Sh(\mathcal{C}/U), \mathcal{O}_U)
\to
(\Sh(\mathcal{D}/V), \mathcal{O}'_V)
$$
是局部环拓扑斯态射。
\item 设 $(\Sh(\mathcal{C}), \mathcal{O})$ 为局部环拓扑斯，
$\mathcal{F}$ 为 $\mathcal{C}$ 上的层。则局部化
$(\Sh(\mathcal{C})/\mathcal{F}, \mathcal{O}_\mathcal{F})$
是局部环拓扑斯，且局部化态射
$$
(j_\mathcal{F}, j_\mathcal{F}^\sharp) :
(\Sh(\mathcal{C})/\mathcal{F}, \mathcal{O}_\mathcal{F})
\to
(\Sh(\mathcal{C}), \mathcal{O})
$$
是局部环拓扑斯态射。
\item 设 $(\Sh(\mathcal{C}), \mathcal{O})$ 为局部环拓扑斯，
$s : \mathcal{G} \to \mathcal{F}$ 为 $\mathcal{C}$ 上层的映射。
则引理 \ref{sites-modules-lemma-relocalize-ringed-topos} 中的态射
$$
(j, j^\sharp) :
(\Sh(\mathcal{C})/\mathcal{G}, \mathcal{O}_\mathcal{G})
\longrightarrow
(\Sh(\mathcal{C})/\mathcal{F}, \mathcal{O}_\mathcal{F})
$$
是局部环拓扑斯态射。
\item 设
$f :
(\Sh(\mathcal{C}), \mathcal{O})
\longrightarrow
(\Sh(\mathcal{D}), \mathcal{O}')$
为局部环拓扑斯态射，$\mathcal{G}$ 为 $\mathcal{D}$ 上的层，
并置 $\mathcal{F} = f^{-1}\mathcal{G}$。则引理
\ref{sites-modules-lemma-localize-morphism-ringed-topoi} 中的态射
$$
(f', (f')^\sharp) :
(\Sh(\mathcal{C})/\mathcal{F}, \mathcal{O}_\mathcal{F})
\longrightarrow
(\Sh(\mathcal{D})/\mathcal{G}, \mathcal{O}'_\mathcal{G})
$$
是局部环拓扑斯态射。
\item 设
$f :
(\Sh(\mathcal{C}), \mathcal{O})
\longrightarrow
(\Sh(\mathcal{D}), \mathcal{O}')$
为局部环拓扑斯态射，$\mathcal{G}$ 为 $\mathcal{D}$ 上的层，
$\mathcal{F}$ 为 $\mathcal{C}$ 上的层，并设
$s : \mathcal{F} \to f^{-1}\mathcal{G}$ 为层态射。则引理
\ref{sites-modules-lemma-relocalize-morphism-ringed-topoi} 中的态射
$$
(f_s, (f_s)^\sharp) :
(\Sh(\mathcal{C})/\mathcal{F}, \mathcal{O}_\mathcal{F})
\longrightarrow
(\Sh(\mathcal{D})/\mathcal{G}, \mathcal{O}'_\mathcal{G})
$$
是局部环拓扑斯态射。
\end{enumerate}
\end{lemma}

\begin{proof}
(1) 显然，因为 $\mathcal{O}_U$ 就是 $\mathcal{O}$ 的限制，
故可应用引理 \ref{sites-modules-lemma-locally-ringed-intrinsic} 与
\ref{sites-modules-lemma-locally-ringed-intrinsic-morphism}。(2) 也显然，因为态射
$(j,j^\sharp)$ 实际上是局部环站点的局部化，故可应用 (1)。
(3) 同样显然，因为 $(f')^\sharp$ 就是 $f^\sharp$
在拓扑斯 $\Sh(\mathcal{C})/\mathcal{F}$ 上的限制；见引理
\ref{sites-modules-lemma-localize-morphism-ringed-topoi} 的证明
（因此定义 \ref{sites-modules-definition-morphism-locally-ringed-topoi}
中态射 $f'$ 的图，就是态射 $f$ 的相应图的限制，而限制函子正合）。
结合 (2) 与 (3)，形式地得到 (4)。按引理
\ref{sites-modules-lemma-morphism-ringed-topoi-comes-from-morphism-ringed-sites}
扩大站点，并利用引理 \ref{sites-modules-lemma-localize-compare}、
\ref{sites-modules-lemma-relocalize-compare}、
\ref{sites-modules-lemma-localize-morphism-compare} 及
\ref{sites-modules-lemma-relocalize-morphism-compare} 的比较，把一切翻译成
站点及站点态射的语言，即由 (1)、(2)、(3)、(4) 分别得到
(5)、(6)、(7)、(8)。（也可以直接使用证明 (1)、(2)、(3)、(4)
时相同的论证来证明 (5)、(6)、(7)、(8)。）
\end{proof}






\section{模的下叹号函子}
\label{sites-modules-section-lower-shriek-modules}

\noindent
本节把第 \ref{sites-modules-section-exactness-lower-shriek} 节讨论的
$g_!$ 构造推广到模层的情形。

\begin{lemma}
\label{sites-modules-lemma-lower-shriek-modules}
设 $u : \mathcal{C} \to \mathcal{D}$ 为站点之间既连续又余连续的函子，
记 $g : \Sh(\mathcal{C}) \to \Sh(\mathcal{D})$
为关联的拓扑斯态射。设 $\mathcal{O}_\mathcal{D}$
为 $\mathcal{D}$ 上的环层，并置
$\mathcal{O}_\mathcal{C} = g^{-1}\mathcal{O}_\mathcal{D}$。
于是 $g$ 成为满足 $g^* = g^{-1}$ 的环拓扑斯态射。
此时存在函子
$$
g_! :
\textit{Mod}(\mathcal{O}_\mathcal{C})
\longrightarrow
\textit{Mod}(\mathcal{O}_\mathcal{D})
$$
它是 $g^*$ 的左伴随。
\end{lemma}

\begin{proof}
设 $U$ 为 $\mathcal{C}$ 的对象。对任意
$\mathcal{O}_\mathcal{D}$-模 $\mathcal{G}$，有
\begin{align*}
\Hom_{\mathcal{O}_\mathcal{C}}(j_{U!}\mathcal{O}_U, g^{-1}\mathcal{G})
& =
g^{-1}\mathcal{G}(U) \\
& =
\mathcal{G}(u(U)) \\
& =
\Hom_{\mathcal{O}_\mathcal{D}}(j_{u(U)!}\mathcal{O}_{u(U)}, \mathcal{G})
\end{align*}
因为 $g^{-1}$ 由限制描述；见《站点》，引理
\ref{sites-lemma-when-shriek}。对形如
$j_{U!}\mathcal{O}_U$ 的模的直和，当然有类似公式。
由《同调》，引理 \ref{homology-lemma-partially-defined-adjoint}
及引理 \ref{sites-modules-lemma-module-quotient-flat}，可知 $g_!$ 存在。
\end{proof}

\begin{remark}
\label{sites-modules-remark-when-shriek-equal}
警告！设 $u : \mathcal{C} \to \mathcal{D}$、$g$、
$\mathcal{O}_\mathcal{D}$ 与 $\mathcal{O}_\mathcal{C}$
如引理 \ref{sites-modules-lemma-lower-shriek-modules}。一般而言，下图{\bf 不}交换：
$$
\xymatrix{
\textit{Mod}(\mathcal{O}_\mathcal{C}) \ar[r]_{g_!} \ar[d]_{forget} &
\textit{Mod}(\mathcal{O}_\mathcal{D}) \ar[d]^{forget} \\
\textit{Ab}(\mathcal{C}) \ar[r]^{g^{Ab}_!} &
\textit{Ab}(\mathcal{D})
}
$$
（这里 $g^{Ab}_!$ 是引理 \ref{sites-modules-lemma-g-shriek-adjoint}
中的函子。）有函子变换
$$
g_!^{Ab} \circ forget \longrightarrow forget \circ g_!
$$
由引理 \ref{sites-modules-lemma-lower-shriek-modules} 的证明可知，
该变换为同构，当且仅当对 $\mathcal{C}$ 的每个对象 $U$，
$g^{Ab}_!j_{U!}\mathcal{O}_U \to g_!j_{U!}\mathcal{O}_U$
都是同构。由于
$g_!j_{U!}\mathcal{O}_U = j_{u(U)!}\mathcal{O}_{u(U)}$，
这又等价于
$$
g^{Ab}_!j_{U!}\mathcal{O}_U \longrightarrow j_{u(U)!}\mathcal{O}_{u(U)}
$$
对 $\mathcal{C}$ 的每个对象 $U$ 都是同构。注意，对这样的 $U$
有交换图
$$
\xymatrix{
\mathcal{C}/U \ar[r]_-{u'} \ar[d]_{j_U} & \mathcal{D}/u(U) \ar[d]^{j_{u(U)}} \\
\mathcal{C} \ar[r]^u & \mathcal{D}
}
$$
其中各箭头是站点的余连续函子；见《站点》，引理
\ref{sites-lemma-localize-cocontinuous}。因此
$g^{Ab}_!j_{U!} = j_{u(U)!}(g')^{Ab}_!$，其中
$g' : \Sh(\mathcal{C}/U) \to \Sh(\mathcal{D}/u(U))$
是由余连续函子 $u'$ 诱导的拓扑斯态射。于是，若对
$\mathcal{C}$ 的每个对象 $U$，典范映射
\begin{equation}
\label{sites-modules-equation-compare-on-localizations}
(g')^{Ab}_!\mathcal{O}_U \longrightarrow \mathcal{O}_{u(U)}
\end{equation}
都是同构，则 $g_! = g^{Ab}_!$。
\end{remark}

\noindent
下面两个结论的性质略有不同。

\begin{lemma}
\label{sites-modules-lemma-special-square-cocontinuous}
假设给定环拓扑斯的交换图
$$
\xymatrix{
(\Sh(\mathcal{C}'), \mathcal{O}_{\mathcal{C}'})
\ar[r]_{(g', (g')^\sharp)} \ar[d]_{(f', (f')^\sharp)} &
(\Sh(\mathcal{C}), \mathcal{O}_\mathcal{C}) \ar[d]^{(f, f^\sharp)} \\
(\Sh(\mathcal{D}'), \mathcal{O}_{\mathcal{D}'}) \ar[r]^{(g, g^\sharp)} &
(\Sh(\mathcal{D}), \mathcal{O}_\mathcal{D})
}
$$
并假设：
\begin{enumerate}
\item $f$、$f'$、$g$、$g'$ 分别对应于余连续函子
$u$、$u'$、$v$、$v'$，如《站点》，引理
\ref{sites-lemma-cocontinuous-morphism-topoi}；
\item $v \circ u' = u \circ v'$,
\item $v$ 与 $v'$ 既连续又余连续；
\item 对 $\mathcal{D}'$ 的任意对象 $V'$，由 $v$ 给出的函子
${}^{u'}_{V'}\mathcal{I} \to {}^{\ \ \ u}_{v(V')}\mathcal{I}$
是余终的；
\item $g^{-1}\mathcal{O}_{\mathcal{D}} = \mathcal{O}_{\mathcal{D}'}$
且 $(g')^{-1}\mathcal{O}_{\mathcal{C}} = \mathcal{O}_{\mathcal{C}'}$。
\end{enumerate}
则在模上有 $f'_* \circ (g')^* = g^* \circ f_*$ 及
$g'_! \circ (f')^{-1} = f^{-1} \circ g_!$。
\end{lemma}

\begin{proof}
由条件 (5)，
$(g')^*\mathcal{F} = (g')^{-1}\mathcal{F}$ 且
$g^*\mathcal{G} = g^{-1}\mathcal{G}$。因此第一个等式立即由
《站点》，引理 \ref{sites-lemma-special-square-cocontinuous}
中的相应等式得到。由引理 \ref{sites-modules-lemma-lower-shriek-modules}，
$g^*$ 与 $(g')^*$ 的左伴随函子 $g_!$ 与 $g'_!$ 存在，
故第二个等式由伴随函子的唯一性得到。
\end{proof}

\begin{lemma}
\label{sites-modules-lemma-special-square-continuous}
考虑环拓扑斯的交换图
$$
\xymatrix{
(\Sh(\mathcal{C}'), \mathcal{O}_{\mathcal{C}'})
\ar[r]_{(g', (g')^\sharp)} \ar[d]_{(f', (f')^\sharp)} &
(\Sh(\mathcal{C}), \mathcal{O}_\mathcal{C}) \ar[d]^{(f, f^\sharp)} \\
(\Sh(\mathcal{D}'), \mathcal{O}_{\mathcal{D}'}) \ar[r]^{(g, g^\sharp)} &
(\Sh(\mathcal{D}), \mathcal{O}_\mathcal{D})
}
$$
并假设有函子
$$
\xymatrix{
\mathcal{C}' \ar[r]_{v'} &
\mathcal{C} \\
\mathcal{D}' \ar[r]^v \ar[u]^{u'} &
\mathcal{D} \ar[u]_u
}
$$
使得（采用《站点》，第 \ref{sites-section-morphism-sites} 节及
\ref{sites-section-cocontinuous-morphism-topoi} 节中的记号）：
\begin{enumerate}
\item $u$ 与 $u'$ 连续，并诱导态射 $f$ 与 $f'$；
\item $v$ 与 $v'$ 余连续，并诱导态射 $g$ 与 $g'$；
\item $u \circ v = v' \circ u'$,
\item $v$ 与 $v'$ 既连续又余连续；
\item $g^{-1}\mathcal{O}_{\mathcal{D}} = \mathcal{O}_{\mathcal{D}'}$
且 $(g')^{-1}\mathcal{O}_{\mathcal{C}} = \mathcal{O}_{\mathcal{C}'}$。
\end{enumerate}
则在模上有 $f'_* \circ (g')^* = g^* \circ f_*$ 及
$g'_! \circ (f')^{-1} = f^{-1} \circ g_!$。
\end{lemma}

\begin{proof}
由条件 (5)，
$(g')^*\mathcal{F} = (g')^{-1}\mathcal{F}$ 且
$g^*\mathcal{G} = g^{-1}\mathcal{G}$。因此第一个等式立即由
《站点》，引理 \ref{sites-lemma-special-square-continuous}
中的相应等式得到。由引理 \ref{sites-modules-lemma-lower-shriek-modules}，
$g^*$ 与 $(g')^*$ 的左伴随函子 $g_!$ 与 $g'_!$ 存在，
故第二个等式由伴随函子的唯一性得到。
\end{proof}







\section{常值层}
\label{sites-modules-section-constant}

\noindent
设 $E$ 为集合，$\mathcal{C}$ 为站点。取值为 $E$ 的常值层记为
$\underline{E}$；见《站点》，例 \ref{sites-example-constant-sheaf}。
若 $E$ 是阿贝尔群、环、模等，则 $\underline{E}$
分别是阿贝尔群层、环层、模层等。

\begin{lemma}
\label{sites-modules-lemma-constant-exact}
设 $\mathcal{C}$ 为站点。若
$0 \to A \to B \to C \to 0$ 是阿贝尔群的短正合列，则
$0 \to \underline{A} \to \underline{B} \to \underline{C} \to 0$
是阿贝尔层的正合列；事实上，它作为阿贝尔预层的序列也正合。
\end{lemma}

\begin{proof}
由于层化正合，显然
$0 \to \underline{A} \to \underline{B} \to \underline{C} \to 0$
是阿贝尔层的正合列。因此
$0 \to \underline{A} \to \underline{B} \to \underline{C}$
是阿贝尔预层的正合列。为证
$\underline{B} \to \underline{C}$ 满射，选取一个集合论截面
$s : C \to B$。这诱导集合层的截面
$\underline{s} : \underline{C} \to \underline{B}$，
它是满射 $\underline{B} \to \underline{C}$ 的左逆。
\end{proof}

\begin{lemma}
\label{sites-modules-lemma-tensor-with-finitely-presented}
设 $\mathcal{C}$ 为站点，$\Lambda$ 为环，而 $M$ 与 $Q$
为 $\Lambda$-模。若 $Q$ 是有限表示的 $\Lambda$-模，则对所有
$U \in \Ob(\mathcal{C})$，有
$\underline{M \otimes_\Lambda Q}(U) = \underline{M}(U) \otimes_\Lambda Q$
。
\end{lemma}

\begin{proof}
选取表示
$\Lambda^{\oplus m} \to \Lambda^{\oplus n} \to Q \to 0$。
它给出正合列
$M^{\oplus m} \to M^{\oplus n} \to M \otimes Q \to 0$.
由引理 \ref{sites-modules-lemma-constant-exact}，得到正合列
$$
\underline{M}(U)^{\oplus m} \to
\underline{M}(U)^{\oplus n} \to
\underline{M \otimes Q}(U) \to 0
$$
这证明了引理。（注意，在 $U$ 上取截面总与有限直和交换，
但不一定与任意直和交换。）
\end{proof}

\begin{lemma}
\label{sites-modules-lemma-flat-sections}
设 $\mathcal{C}$ 为站点，$\Lambda$ 为凝聚环，
$M$ 为平坦 $\Lambda$-模。对 $U \in \Ob(\mathcal{C})$，
模 $\underline{M}(U)$ 是平坦 $\Lambda$-模。
\end{lemma}

\begin{proof}
设 $I \subset \Lambda$ 为有限生成理想。由《代数》，引理
\ref{algebra-lemma-flat}，只需证明
$\underline{M}(U) \otimes_\Lambda I \to \underline{M}(U)$
单射。由于 $\Lambda$ 凝聚，$I$ 作为 $\Lambda$-模有限表示。
由引理 \ref{sites-modules-lemma-tensor-with-finitely-presented}，
$\underline{M}(U) \otimes I = \underline{M \otimes I}$。
由于 $M$ 平坦，映射 $M \otimes I \to M$ 单射，因而
$\underline{M \otimes I} \to \underline{M}$ 单射。
\end{proof}

\begin{lemma}
\label{sites-modules-lemma-completion-flat}
设 $\mathcal{C}$ 为站点，$\Lambda$ 为 Noether 环，
$I \subset \Lambda$ 为理想。层
$\underline{\Lambda}^\wedge = \lim \underline{\Lambda/I^n}$
是平坦 $\underline{\Lambda}$-代数。此外有典范等同
$$
\underline{\Lambda}/I\underline{\Lambda} =
\underline{\Lambda}/\underline{I} =
\underline{\Lambda}^\wedge/I\underline{\Lambda}^\wedge =
\underline{\Lambda}^\wedge/\underline{I} \cdot \underline{\Lambda}^\wedge =
\underline{\Lambda}^\wedge/\underline{I}^\wedge =
\underline{\Lambda/I}
$$
其中 $\underline{I}^\wedge = \lim \underline{I/I^n}$。
\end{lemma}

\begin{proof}
为证明 $\underline{\Lambda}^\wedge$ 平坦，只需证明对每个
$U \in \Ob(\mathcal{C})$，$\underline{\Lambda}^\wedge(U)$
作为 $\Lambda$-模平坦；见引理 \ref{sites-modules-lemma-flatness-presheaves}
及 \ref{sites-modules-lemma-flatness-sheafification}。由引理
\ref{sites-modules-lemma-flat-sections}，
$$
\underline{\Lambda}^\wedge(U) = \lim \underline{\Lambda/I^n}(U)
$$
是平坦 $\Lambda/I^n$-模系统的极限。由引理
\ref{sites-modules-lemma-constant-exact}，转移映射均满射。再由
《更多代数》，引理 \ref{more-algebra-lemma-limit-flat} 即得结论。

\medskip\noindent
为证明这些等式，注意由引理
\ref{sites-modules-lemma-tensor-with-finitely-presented}，
$\underline{\Lambda}(U)/I\underline{\Lambda}(U) = \underline{\Lambda/I}(U)$
。因此
$\underline{\Lambda}/I\underline{\Lambda} =
\underline{\Lambda}/\underline{I} = \underline{\Lambda/I}$。
短正合列系统
$$
0 \to \underline{I/I^n}(U) \to \underline{\Lambda/I^n}(U) \to
\underline{\Lambda/I}(U) \to 0
$$
具有满的转移映射，故给出短正合列
$$
0 \to \lim \underline{I/I^n}(U) \to \lim \underline{\Lambda/I^n}(U) \to
\lim \underline{\Lambda/I}(U) \to 0
$$
见《同调》，引理 \ref{homology-lemma-Mittag-Leffler}。因此
$\underline{\Lambda}^\wedge/\underline{I}^\wedge =
\underline{\Lambda/I}$。由于
$$
I \underline{\Lambda}^\wedge \subset
\underline{I} \cdot \underline{\Lambda}^\wedge \subset
\underline{I}^\wedge
$$
只需证明对所有 $U$，
$I \underline{\Lambda}^\wedge(U) = \underline{I}^\wedge(U)$
。选取生成元 $I = (f_1, \ldots, f_r)$。对每个 $n$，得到短正合列
$$
0 \to K_n/(I^n)^{\oplus r} \to (\Lambda/I^n)^{\oplus r}
\xrightarrow{(f_1, \ldots, f_r)}
I/I^{n + 1} \to 0
$$
其中
$K_n =
\{(x_1, \ldots, x_r) \in \Lambda^{\oplus r} \mid \sum x_i f_i \in I^{n + 1}\}$.
由此得到短正合列
$$
0 \to \underline{K_n/(I^n)^{\oplus r}}(U) \to
\underline{(\Lambda/I^n)^{\oplus r}}(U) \to
\underline{I/I^{n + 1}}(U) \to 0
$$
计算表明 $K_n = K + (I^n)^{\oplus r}$，故转移映射
$K_{n + 1}/(I^{n + 1})^{\oplus r} \to K_n/(I^n)^{\oplus r}$
满射。因此左侧模的系统具有满的转移映射，从而更满足 ML 条件。
由此可知
$(f_1, \ldots, f_r) :
(\underline{\Lambda}^\wedge)^{\oplus r}(U) \to \underline{I}^\wedge(U)$
由《同调》，引理 \ref{homology-lemma-Mittag-Leffler} 是满射，
这正是所需结论。
\end{proof}

\begin{lemma}
\label{sites-modules-lemma-locally-constant-finite-type}
设 $\mathcal{C}$ 为站点，$\Lambda$ 为环，$M$ 为 $\Lambda$-模。
假设 $\Sh(\mathcal{C})$ 不是空拓扑斯。则
\begin{enumerate}
\item $\underline{M}$ 是有限型 $\underline{\Lambda}$-模层，
当且仅当 $M$ 是有限 $\Lambda$-模；
\item $\underline{M}$ 是有限表示的 $\underline{\Lambda}$-模层，
当且仅当 $M$ 是有限表示的 $\Lambda$-模。
\end{enumerate}
\end{lemma}

\begin{proof}
证明 (1)。若 $M$ 由 $x_1,\ldots,x_r$ 生成，则这些元素定义
$\underline{M}$ 的整体截面并生成该层，故 $\underline{M}$ 有限型。
反之，假设 $\underline{M}$ 有限型。取一个在层论意义下不为空的
对象 $U \in \mathcal{C}$（《站点》，定义
\ref{sites-definition-empty}）。因为假设
$\Sh(\mathcal{C})$ 不是空拓扑斯，这样的对象存在。于是存在覆盖
$\{U_i \to U\}$ 及有限多个截面
$s_{ij} \in \underline{M}(U_i)$，它们生成
$\underline{M}|_{U_i}$。加细覆盖后，可以假设 $s_{ij}$
来自 $M$ 的元素 $x_{ij}$。于是 $x_{ij}$ 定义
$\underline{M}$ 的整体截面，它们在 $U$ 上的限制生成
$\underline{M}$。

\medskip\noindent
假设存在 $M$ 的元素 $x_1,\ldots,x_r$，它们定义
$\underline{M}$ 的整体截面，并作为
$\underline{\Lambda}$-模层生成 $\underline{M}$。我们证明
$x_1,\ldots,x_r$ 作为 $\Lambda$-模生成 $M$。取 $x \in M$。
可找到覆盖 $\{U_i \to U\}_{i \in I}$ 及
$f_{i,j} \in \underline{\Lambda}(U_i)$，使
$x|_{U_i} = \sum f_{i,j}x_j|_{U_i}$。加细覆盖后，可以假设
$f_{i,j} \in \Lambda$。由于 $U$ 在层论意义下不为空，
至少存在一个 $i \in I$，使 $U_i$ 在层论意义下不为空。
此时映射 $M \to \underline{M}(U_i)$ 单射（细节从略）。
因此 $M$ 中有 $x = \sum f_{i,j}x_j$，如所需。

\medskip\noindent
证明 (2)。假设 $\underline{M}$ 是有限表示的
$\underline{\Lambda}$-模。由 (1)，$M$ 有限型。选取
$M$ 作为 $\Lambda$-模的一组生成元 $x_1,\ldots,x_r$。
这确定短正合列
$0 \to K \to \Lambda^{\oplus r} \to M \to 0$，
由引理 \ref{sites-modules-lemma-constant-exact}，它变为短正合列
$$
0 \to \underline{K} \to \underline{\Lambda}^{\oplus r} \to \underline{M} \to 0
$$
由引理 \ref{sites-modules-lemma-kernel-surjection-finite-onto-finite-presentation}，
$\underline{K}$ 有限型。故由 (1)，$K$ 是有限
$\Lambda$-模。因此 $M$ 是有限表示的 $\Lambda$-模。
\end{proof}




\section{局部常值层}
\label{sites-modules-section-locally-constant}

\noindent
一般定义如下。

\begin{definition}
\label{sites-modules-definition-locally-constant}
设 $\mathcal{C}$ 为站点，$\mathcal{F}$ 为集合、群、阿贝尔群、环、
固定环 $\Lambda$ 上的模等的层。
\begin{enumerate}
\item 若 $\mathcal{F}$ 作为集合、群、阿贝尔群、环、固定环
$\Lambda$ 上的模等的层，同构于第 \ref{sites-modules-section-constant} 节所述的
某个常值层 $\underline{E}$，则称 $\mathcal{F}$ 是相应类型的
{\it 常值层}。
\item 若对 $\mathcal{C}$ 的每个对象 $U$，都存在覆盖
$\{U_i \to U\}$，使 $\mathcal{F}|_{U_i}$ 为常值层，
则称 $\mathcal{F}$ {\it 局部常值}。
\item 若 $\mathcal{F}$ 是集合层或群层，且其常值取值是有限集合
或有限群，则称 $\mathcal{F}$ {\it 有限局部常值}。
\end{enumerate}
\end{definition}

\begin{lemma}
\label{sites-modules-lemma-pullback-locally-constant}
设 $f : \Sh(\mathcal{C}) \to \Sh(\mathcal{D})$ 为拓扑斯态射。
若 $\mathcal{G}$ 是 $\mathcal{D}$ 上集合、群、阿贝尔群、环、
固定环 $\Lambda$ 上的模等的局部常值层，则
$f^{-1}\mathcal{G}$ 是 $\mathcal{C}$ 上相应类型的局部常值层。
\end{lemma}

\begin{proof}
从略。
\end{proof}

\begin{lemma}
\label{sites-modules-lemma-morphism-locally-constant}
设 $\mathcal{C}$ 为具有终对象 $X$ 的站点。
\begin{enumerate}
\item 设 $\varphi : \mathcal{F} \to \mathcal{G}$ 为
$\mathcal{C}$ 上局部常值集合层的映射。若 $\mathcal{F}$
有限局部常值，则存在覆盖 $\{U_i \to X\}$，使
$\varphi|_{U_i}$ 是由集合映射诱导的常值层映射。
\item 设 $\varphi : \mathcal{F} \to \mathcal{G}$ 为
$\mathcal{C}$ 上局部常值阿贝尔群层的映射。若 $\mathcal{F}$
有限局部常值，则存在覆盖 $\{U_i \to X\}$，使
$\varphi|_{U_i}$ 是由阿贝尔群同态诱导的常值阿贝尔层映射。
\item 设 $\Lambda$ 为环，
$\varphi : \mathcal{F} \to \mathcal{G}$ 为
$\mathcal{C}$ 上局部常值 $\Lambda$-模层的映射。若
$\mathcal{F}$ 有限型，则存在覆盖 $\{U_i \to X\}$，使
$\varphi|_{U_i}$ 是由 $\Lambda$-模同态诱导的常值
$\Lambda$-模层映射。
\end{enumerate}
\end{lemma}

\begin{proof}
证明从略。
\end{proof}

\begin{lemma}
\label{sites-modules-lemma-locally-constant}
设 $\mathcal{C}$ 为站点，$\Lambda$ 为环，
$M,N$ 为 $\Lambda$-模，而 $\mathcal{F},\mathcal{G}$
为局部常值 $\Lambda$-模层。
\begin{enumerate}
\item 若 $M$ 有限表示，则
$$
\underline{\Hom_\Lambda(M, N)} =
\SheafHom_{\underline{\Lambda}}(\underline{M}, \underline{N})
$$
\item 若 $M$ 与 $N$ 都有限表示，则
$$
\underline{\text{Isom}_\Lambda(M, N)} =
\mathit{Isom}_{\underline{\Lambda}}(\underline{M}, \underline{N})
$$
\item 若 $\mathcal{F}$ 有限表示，则
$\SheafHom_{\underline{\Lambda}}(\mathcal{F}, \mathcal{G})$
是局部常值 $\Lambda$-模层。
\item 若 $\mathcal{F}$ 与 $\mathcal{G}$ 都有限表示，则
$\mathit{Isom}_{\underline{\Lambda}}(\mathcal{F}, \mathcal{G})$
是局部常值集合层。
\end{enumerate}
\end{lemma}

\begin{proof}
证明 (1)。置 $E = \Hom_\Lambda(M, N)$。须证典范映射
$$
\underline{E}
\longrightarrow
\SheafHom_{\underline{\Lambda}}(\underline{M}, \underline{N})
$$
是同构。模 $M$ 有表示
$\Lambda^{\oplus s} \to \Lambda^{\oplus t} \to M \to 0$.
于是 $E$ 位于正合列
$$
0 \to E \to \Hom_\Lambda(\Lambda^{\oplus t}, N) \to
\Hom_\Lambda(\Lambda^{\oplus s}, N)
$$
中；类似地有
$$
0 \to
\SheafHom_{\underline{\Lambda}}(\underline{M}, \underline{N}) \to
\SheafHom_{\underline{\Lambda}}(\underline{\Lambda^{\oplus t}}, \underline{N})
\to
\SheafHom_{\underline{\Lambda}}(\underline{\Lambda^{\oplus s}}, \underline{N})
$$
这把问题归约到 $M$ 为有限自由模的情形，而此时结论显然。

\medskip\noindent
证明 (3)。该问题在 $\mathcal{C}$ 上是局部的，故可假设
$\mathcal{F} = \underline{M}$ 且 $\mathcal{G} = \underline{N}$，
其中 $M,N$ 为某些 $\Lambda$-模。由引理
\ref{sites-modules-lemma-locally-constant-finite-type}，模 $M$ 有限表示。
故结论由 (1) 得到。

\medskip\noindent
(2) 与 (4) 由 (1)、(3) 及下述事实得到：
$\mathit{Isom}$ 可视为
$\SheafHom_{\underline{\Lambda}}(\mathcal{F}, \mathcal{G})$
的截面所成的子层，其中这些截面在
$\SheafHom_{\underline{\Lambda}}(\mathcal{G}, \mathcal{F})$
中有逆。
\end{proof}

\begin{lemma}
\label{sites-modules-lemma-kernel-finite-locally-constant}
设 $\mathcal{C}$ 为站点。
\begin{enumerate}
\item 有限局部常值集合层的范畴在 $\Sh(\mathcal{C})$
中对有限极限与有限余极限封闭。
\item 有限局部常值阿贝尔层的范畴是
$\textit{Ab}(\mathcal{C})$ 的弱 Serre 子范畴。
\item 设 $\Lambda$ 为 Noether 环。$\mathcal{C}$ 上有限型局部常值
$\Lambda$-模层的范畴是
$\textit{Mod}(\mathcal{C}, \Lambda)$ 的弱 Serre 子范畴。
\end{enumerate}
\end{lemma}

\begin{proof}
证明 (1)。可以在 $\mathcal{C}$ 上局部地工作。因此由引理
\ref{sites-modules-lemma-morphism-locally-constant}，可以假设给定有限集合的有限图，
而我们的层图是其关联的常值层图。此时只须在集合范畴中取极限或余极限，
再取关联的常值层。略去部分细节。

\medskip\noindent
为证明 (2) 与 (3)，使用《同调》，引理
\ref{homology-lemma-characterize-weak-serre-subcategory} 的判据。
核与余核的存在性用与上面相同的方法论证。当然，(3)
中使用 Noether 环，是为了保证有限 $\Lambda$-模之间映射的核仍是
有限 $\Lambda$-模。为证明该范畴对扩张封闭
（在 $\Lambda$-模层的情形），假设给定 $\mathcal{C}$ 上
$\Lambda$-模层的扩张
$$
0 \to \mathcal{F} \to \mathcal{E} \to \mathcal{G} \to 0
$$
其中 $\mathcal{F},\mathcal{G}$ 有限型且局部常值。在
$\mathcal{C}$ 上局部化后，可以假设 $\mathcal{F}$ 与
$\mathcal{G}$ 都是常值层，即有
$$
0 \to \underline{M} \to \mathcal{E} \to \underline{N} \to 0
$$
其中 $M,N$ 为某些 $\Lambda$-模。选取 $N$ 的生成元
$y_1,\ldots,y_m$，得到 $\Lambda$-模的短正合列
$0 \to K \to \Lambda^{\oplus m} \to N \to 0$。
进一步局部化后，可以假设 $y_j$ 提升为 $\mathcal{E}$ 的截面
$s_j$。于是 $\mathcal{E}$ 是下图所示的推出：
$$
\xymatrix{
0 \ar[r] &
\underline{K} \ar[d] \ar[r] &
\underline{\Lambda^{\oplus m}} \ar[d] \ar[r] &
\underline{N} \ar[d] \ar[r] & 0 \\
0 \ar[r] &
\underline{M} \ar[r] &
\mathcal{E} \ar[r] &
\underline{N} \ar[r] & 0
}
$$
再次应用引理 \ref{sites-modules-lemma-morphism-locally-constant}
（并使用 $\Lambda$ 为 Noether 环从而 $K$ 是有限
$\Lambda$-模这一事实），可见映射
$\underline{K} \to \underline{M}$ 局部常值，故得结论。
\end{proof}

\begin{lemma}
\label{sites-modules-lemma-tensor-product-locally-constant}
设 $\mathcal{C}$ 为站点，$\Lambda$ 为环。$\mathcal{C}$
上两个局部常值 $\Lambda$-模层的张量积仍是局部常值
$\Lambda$-模层。
\end{lemma}

\begin{proof}
从略。
\end{proof}











\section{环层的局部化}
\label{sites-modules-section-localizing-sheaves-rings}

\noindent
设 $(\mathcal{C}, \mathcal{O})$ 为环站点，
$\mathcal{S} \subset \mathcal{O}$ 为集合子预层，满足：
对所有 $U \in \Ob(\mathcal{C})$，集合
$\mathcal{S}(U) \subset \mathcal{O}(U)$ 是乘法子集；见《代数》，
定义 \ref{algebra-definition-multiplicative-subset}。此时可以考虑环预层
$$
\mathcal{S}^{-1}\mathcal{O} :
U \longmapsto \mathcal{S}(U)^{-1}\mathcal{O}(U).
$$
对 $\mathcal{C}$ 中的 $V \to U$，限制映射把截面
$f/s$（其中 $f \in \mathcal{O}(U)$、$s \in \mathcal{S}(U)$）
映到 $(f|_V)/(s|_V)$。

\begin{lemma}
\label{sites-modules-lemma-simple-invert}
在上述情形下，到层化的映射
$$
\mathcal{O} \longrightarrow (\mathcal{S}^{-1}\mathcal{O})^\#
$$
是环层同态，并具有如下泛性质：对任意环层同态
$\mathcal{O} \to \mathcal{A}$，若 $\mathcal{S}$
的每个局部截面都映到 $\mathcal{A}$ 的可逆截面，则存在唯一分解
$(\mathcal{S}^{-1}\mathcal{O})^\# \to \mathcal{A}$。
\end{lemma}

\begin{proof}
从略。
\end{proof}

\noindent
设 $(\mathcal{C}, \mathcal{O})$ 为环站点，
$\mathcal{S} \subset \mathcal{O}$ 为集合子预层，满足：
对所有 $U \in \mathcal{C}$，集合
$\mathcal{S}(U) \subset \mathcal{O}(U)$ 是乘法子集。
设 $\mathcal{F}$ 为 $\mathcal{O}$-模层。此时可以考虑
$\mathcal{S}^{-1}\mathcal{O}$-模预层
$$
\mathcal{S}^{-1}\mathcal{F} :
U \longmapsto \mathcal{S}(U)^{-1}\mathcal{F}(U).
$$
若 $V \to U$ 是 $\mathcal{C}$ 中的态射，则限制映射把截面
$t/s$（其中 $t \in \mathcal{F}(U)$、$s \in \mathcal{S}(U)$）
映到 $(t|_V)/(s|_V)$。于是 $\mathcal{S}^{-1}\mathcal{F}$
是 $\mathcal{S}^{-1}\mathcal{O}$-模预层。

\begin{lemma}
\label{sites-modules-lemma-simple-invert-module}
在上述情形下，到层化的映射
$$
\mathcal{F} \longrightarrow (\mathcal{S}^{-1}\mathcal{F})^\#
$$
具有如下泛性质：对任意 $\mathcal{O}$-模同态
$\mathcal{F} \to \mathcal{G}$，若 $\mathcal{S}$
的每个局部截面都在 $\mathcal{G}$ 上可逆地作用，则存在唯一分解
$(\mathcal{S}^{-1}\mathcal{F})^\# \to \mathcal{G}$。
此外有
$$
(\mathcal{S}^{-1}\mathcal{F})^\#
=
(\mathcal{S}^{-1}\mathcal{O})^\# \otimes_\mathcal{O} \mathcal{F}
$$
作为 $(\mathcal{S}^{-1}\mathcal{O})^\#$-模层。
\end{lemma}

\begin{proof}
从略。
\end{proof}









\section{带基点集合层}
\label{sites-modules-section-pointed}

\noindent
本节汇集一些关于带基点集合层的事实；此前我们只在阿贝尔层的情形下
提到过这些事实。

\medskip\noindent
带基点集合是一个二元组 $(S, 0)$，其中 $S$ 是集合，$0 \in S$
是 $S$ 的一个元素。带基点集合的态射
$(S, 0) \to (S', 0')$ 就是把 $0$ 映到 $0'$ 的集合映射
$S \to S'$。我们将稍微滥用记号，把“设 $S$ 为带基点集合”
理解为 $S$ 已配备一个指定元素 $0 \in S$。带基点集合层
等同于配备“基点标记” $0 : * \to \mathcal{F}$ 的集合层
$\mathcal{F}$，其中 $*$ 是终层
（《站点》，例 \ref{sites-example-singleton-sheaf}）。

\medskip\noindent
给定一个站点态射或拓扑斯态射，在带基点集合层的范畴上有直像函子
和拉回函子；见《站点》，第
\ref{sites-section-sheaves-algebraic-structures} 节。
它们由底层集合层分别取直像与拉回并赋予适当的基点标记而构造
（这里使用了终层的拉回仍是终层这一事实）。

\medskip\noindent
设 $u : \mathcal{C} \to \mathcal{D}$ 为站点之间的连续且余连续函子。
设 $g : \Sh(\mathcal{C}) \to \Sh(\mathcal{D})$ 为与 $u$ 相关联的
拓扑斯态射；见《站点》，引理
\ref{sites-lemma-cocontinuous-morphism-topoi}。则在带基点集合层上，
$g^{-1}$ 也有左伴随 $g_!$。这个函子的构造与第
\ref{sites-modules-section-exactness-lower-shriek} 节中阿贝尔层上的 $g_!$
的构造完全类似。

\medskip\noindent
类似地，若 $j : \mathcal{C}/U \to \mathcal{C}$ 如第
\ref{sites-modules-section-localize} 节所述，则带基点集合层上的函子 $j^{-1}$
有左伴随 $j_!$。

\medskip\noindent
若日后需要这些事实与构造，我们将在此精确陈述并证明相应的引理。
% Generated mechanically from the frozen injectives.tex target.
% Do not edit this generated file; edit the bound locale source instead.

% OK, start here.
%

\setcounter{chapter}{18}
\chapter{内射对象}



\phantomsection
\label{injectives-section-phantom}


\section{引言}
\label{injectives-section-introduction}

\noindent
在后续各章中，我们将利用内射对象与 K-内射复形的存在性，研究环站点上
模层的上同调。本章说明如何构造内射对象与 K-内射复形：先讨论站点上的模，
随后更一般地讨论 Grothendieck 阿贝尔范畴。

\medskip\noindent
我们注意到，阿贝尔群范畴以及环上模范畴均有足够多的内射对象；见《代数详论》，
第
\ref{more-algebra-section-abelian-groups} 与
\ref{more-algebra-section-injectives-modules} 节。





\section{模的 Baer 论证}
\label{injectives-section-baer}

\noindent
还有一种更具集合论色彩的方法，可以证明任意 $R$-模 $M$ 都能嵌入一个内射模。
该方法借助推出的超限余极限来构造内射模。虽然这种方法比《代数详论》第
\ref{more-algebra-section-injectives-modules} 节中的方法更为抽象且更复杂，
但它也更为一般。这个方法似乎源自 Baer，后来 Cartan 与 Eilenberg 在
\cite{Cartan-Eilenberg} 中、Grothendieck 在 \cite{Tohoku} 中重新讨论了它。
Grothendieck 在后一文中用它证明许多其他阿贝尔范畴也有足够多的内射对象。
稍后我们将回到一般情形（第 \ref{injectives-section-grothendieck-categories} 节）。

\medskip\noindent
先作几条集合论方面的说明。设 $\{B_{\beta}\}_{\beta \in \alpha}$ 是某范畴
$\mathcal{C}$ 中由序数 $\alpha$ 指标的对象归纳系统。假设
$\colim_{\beta \in \alpha} B_\beta$ 在 $\mathcal{C}$ 中存在。
若 $A$ 是 $\mathcal{C}$ 的对象，则有自然映射
\begin{equation}
\label{injectives-equation-compare}
\colim_{\beta \in \alpha} \Mor_\mathcal{C}(A, B_\beta)
\longrightarrow
\Mor_\mathcal{C}(A, \colim_{\beta \in \alpha} B_\beta).
\end{equation}
这是因为，若给定某个 $\beta$ 以及映射 $A \to B_\beta$，将它与
$B_\beta \to \colim_{\beta \in \alpha} B_\alpha$ 复合，便自然得到从
$A$ 到该余极限的映射。注意，左侧是集合的余极限！一般而言，
(\ref{injectives-equation-compare}) 既不是单射，也不是满射。

\begin{example}
\label{injectives-example-not-surjective}
考虑集合范畴。令 $A = \mathbf{N}$，并令
$B_n = \{1, \ldots, n\}$ 为由自然数指标的归纳系统，其中当
$n \leq m$ 时，$B_n \to B_m$ 是显然映射。此时
$\colim B_n = \mathbf{N}$，因而存在映射 $A \to \colim B_n$，
但对任意 $m$，它都不能分解为 $A \to B_m$。因此，
$\colim \Mor(A, B_n) \to \Mor(A, \colim B_n)$
不是满射。
\end{example}

\begin{example}
\label{injectives-example-not-injective}
下面给出该映射不是单射的例子。令
$B_n = \mathbf{N}/\{1,  2, \ldots, n\}$，即把 $\mathbf{N}$ 的前
$n$ 个元素压缩为一个元素所得的商集。当 $n \leq m$ 时有自然映射
$B_n \to B_m$，故 $\{B_n\}$ 构成 $\mathbf{N}$ 上的集合系统。
容易看出 $\colim B_n = \{*\}$，即单点集。因此，对任意集合 $A$，
$\Mor(A, \colim B_n)$ 都是单元素集。然而，
$\colim \Mor(A , B_n)$ {\bf 不是}单元素集。考虑两族映射
$A \to B_n$：第一族就是自然投影
$\mathbf{N} \to \mathbf{N}/\{1, 2, \ldots, n\}$；第二族把整个
$A$ 映到 $1$ 的等价类。这两族映射在每一步都不同，因而在
$\colim \Mor(A, B_n)$ 中也不同，但它们诱导出相同的映射
$A \to \colim B_n$。
\end{example}

\noindent
不过，若映射从有限集出发，则 (\ref{injectives-equation-compare}) 总是同构。

\begin{lemma}
\label{injectives-lemma-out-of-finite}
在 (\ref{injectives-equation-compare}) 中，若 $\mathcal{C}$ 是集合范畴且
$A$ 是{\it 有限集}，则该映射是双射。
\end{lemma}

\begin{proof}
设 $f : A \to \colim B_\beta$。$f$ 的像有限；设它包含余极限中的元素
$c_1, \ldots, c_r \in \colim B_\beta$。由余极限的定义，当
$\beta \in \alpha$ 足够大时，这些元素都来自 $B_\beta$ 中的某些元素。
因此可以定义 $\widetilde{f} : A \to B_\beta$，在某个有限阶段提升
$f$。这证明了 (\ref{injectives-equation-compare}) 是满射。

接下来，假设两个映射 $f : A \to B_\gamma$ 与
$f' : A \to B_{\gamma'}$ 定义了相同的映射
$A \to \colim B_\beta$。于是 $A$ 的每个元素在余极限中的像都相同。
由集合余极限的定义，存在 $\beta \geq \gamma, \gamma'$，使得在
$f$ 与 $f'$ 下，$A$ 的有限多个元素在 $B_\beta$ 中分别映到相同的点。
这证明了 (\ref{injectives-equation-compare}) 是单射。
\end{proof}

\noindent
该引理最有趣的情形是 $\alpha = \omega$，也就是说，系统
$\{B_\beta\}$ 是例 \ref{injectives-example-not-surjective} 与
\ref{injectives-example-not-injective} 中那样由自然数指标的系统
$\{B_n\}_{n \in \mathbf{N}}$。核心思想是，相对于长复合链
$B_1 \to B_2 \to \ldots$，$A$ 是“小”的，因而它必定分解通过某个
有限阶段。该引理的更一般版本见《集合》，引理
\ref{sets-lemma-map-from-set-lifts}。下面把它推广到模范畴。

\begin{definition}
\label{injectives-definition-small}
设 $\mathcal{C}$ 为范畴，$I \subset \text{Arrows}(\mathcal{C})$，
$\alpha$ 为序数。若对任意以 $\alpha$ 为指标且转移映射均属于 $I$
的系统 $\{B_\beta\}$，映射 (\ref{injectives-equation-compare}) 都是同构，
则称 $\mathcal{C}$ 的对象 $A$ {\it 相对于 $I$ 为 $\alpha$-小的}。
\end{definition}

\noindent
本节以下固定一个交换环 $R$，并只考虑 $R$-模范畴。我们还令 $I$
为单射的集合，即 $R$-模范畴中的{\it 单态射}集合。在这种情形下，
对定义中任意系统 $\{B_\beta\}$，每个映射
$$
B_\beta \to \colim_{\beta \in \alpha} B_\beta
$$
都是单射。因此映射 (\ref{injectives-equation-compare}) 是{\it 单射}。事实上，
可以把各 $B_\beta$ 视为模
$B = \colim_{\beta \in \alpha} B_\beta$ 的子模，此时
$B = \bigcup_{\beta \in \alpha} B_\beta$。若把 $B_\alpha$
与它在余极限中的像认同，这并非记号的滥用。下面证明：对“足够大”的序数
$\alpha$，模总是小的。

\begin{proposition}
\label{injectives-proposition-modules-are-small}
设 $R$ 为环，$M$ 为 $R$-模。设 $\kappa$ 为 $M$ 的子模集合的基数。
若序数 $\alpha$ 的共尾数大于 $\kappa$，则 $M$ 相对于单射为
$\alpha$-小的。
\end{proposition}

\begin{proof}
证明是直接的，不过先考虑一个特殊情形。若 $M$ 有限，则断言是：对任意由极限序数
指标、转移映射均为单射的归纳系统 $\{B_\beta\}$，任意映射
$M \to \colim B_\beta$ 都分解通过某个 $B_\beta$。这已在引理
\ref{injectives-lemma-out-of-finite} 中证明。

\medskip\noindent
现在证明一般情形。只需证明映射 (\ref{injectives-equation-compare}) 是满射。
设 $f : M \to \colim B_\beta$ 为映射。把 $B_\beta$ 看作余极限
$B = \bigcup_\beta B_\beta$ 的子对象，并考虑 $M$ 的子对象
$\{f^{-1}(B_\beta)\}$。若其中某个子对象（比如
$f^{-1}(B_\beta)$）等于 $M$，则该映射分解通过 $B_\beta$。

\medskip\noindent
反设所有 $f^{-1}(B_\beta)$ 都是 $M$ 的真子对象。但我们知道
$$
\bigcup f^{-1}(B_\beta) = f^{-1}\left(\bigcup B_\beta\right) = M.
$$
由假设，在各 $f^{-1}(B_\alpha)$ 中至多出现 $\kappa$ 个不同的 $M$
的子对象。因此可以找到基数至多为 $\kappa$ 的子集
$S \subset \alpha$，使得当 $\beta'$ 遍历 $S$ 时，
$f^{-1}(B_{\beta'})$ 恰好遍历\emph{所有} $f^{-1}(B_\alpha)$。

\medskip\noindent
由于 $\alpha$ 的共尾数大于 $\kappa$，$S$ 有一个上界
$\widetilde{\alpha} < \alpha$。特别地，对所有 $\beta' \in S$，
$f^{-1}(B_{\beta'})$ 都包含于
$f^{-1}(B_{\widetilde{\alpha}})$。因此
$f^{-1}(B_{\widetilde{\alpha}}) = M$，从而映射 $f$ 分解通过
$B_{\widetilde{\alpha}}$。
\end{proof}

\noindent
由此引理可以推出大量内射对象的存在性。先回顾 Baer 判别法。

\begin{lemma}[Baer 判别法]
\label{injectives-lemma-criterion-baer}
\begin{reference}
\cite[Theorem 1]{Baer}
\end{reference}
设 $R$ 为环。$R$-模 $Q$ 是内射模，当且仅当在每个交换图
$$
\xymatrix{
\mathfrak{a} \ar[d] \ar[r] &  Q \\
R \ar@{-->}[ru]
}
$$
中虚线箭头都存在，其中 $\mathfrak{a} \subset R$ 为理想。
\end{lemma}

\begin{proof}
这就是《代数详论》，引理
\ref{more-algebra-lemma-characterize-injective-bis} 中 (1) 与 (3)
的等价性。请注意，那里给出的证明是初等的（既不使用 $\text{Ext}$ 群，
也不使用 $R$-模范畴中内射对象或投射对象的存在性）。
\end{proof}

\noindent
若 $M$ 是 $R$-模，则一般可能有一个如引理
\ref{injectives-lemma-criterion-baer} 中那样尚缺一箭头的图。在其中可以形成\emph{推出}
$$
\xymatrix{
\mathfrak{a} \ar[d]  \ar[r] &  M \ar[d] \\
R \ar[r] &  R \oplus_{\mathfrak{a}} M.
}
$$
这里的竖直映射是单射，且该图交换。要点是：只要把 $M$ 扩大为更大的模
$R \oplus_{\mathfrak{a}} M$，就\emph{可以}把
$\mathfrak{a} \to M$ 延拓到 $R$。

\medskip\noindent
Baer 论证的关键是超限多次地重复这个过程。为此，对给定的 $R$-模 $M$，
先定义如下（巨大）推出
\begin{equation}
\label{injectives-equation-huge-diagram}
\vcenter{
\xymatrix{
\bigoplus_{\mathfrak a}
\bigoplus_{\varphi \in \Hom_R(\mathfrak a, M)}
\mathfrak{a} \ar[r] \ar[d] & M \ar[d] \\
\bigoplus_{\mathfrak a}
\bigoplus_{\varphi \in \Hom_R(\mathfrak a, M)}
R \ar[r] &  \mathbf{M}(M).
}
}
\end{equation}
上方水平箭头把与 $\varphi$ 对应的直和项中元素
$a \in \mathfrak a$ 映到 $\varphi(a) \in M$。左侧竖直箭头则把与
$\varphi$ 对应的直和项中的 $a \in \mathfrak a$ 直接映到与
$\varphi$ 对应的直和项中的 $a \in R$。该构造的基本性质表述于下一个引理。

\begin{lemma}
\label{injectives-lemma-construction}
设 $R$ 为环。
\begin{enumerate}
\item 构造 $M \mapsto (M \to \mathbf{M}(M))$ 关于 $M$ 具有函子性。
\item 映射 $M \to \mathbf{M}(M)$ 是单射。
\item 对任意理想 $\mathfrak{a}$ 以及任意 $R$-模同态
$\varphi : \mathfrak a \to M$，存在 $R$-模同态
$\varphi' : R \to \mathbf{M}(M)$，使得
$$
\xymatrix{
\mathfrak{a} \ar[d] \ar[r]_\varphi &  M \ar[d] \\
R \ar[r]^{\varphi'} & \mathbf{M}(M)
}
$$
交换。
\end{enumerate}
\end{lemma}

\begin{proof}
(2) 与 (3) 由构造立即得到。为证明 (1)，设
$\chi : M \to N$ 为 $R$-模同态。我们断言存在典范交换图
$$
\xymatrix{
\bigoplus_{\mathfrak a}
\bigoplus_{\varphi \in \Hom_R(\mathfrak a, M)}
\mathfrak{a} \ar[r] \ar[d] \ar[rrd] & M \ar[rrd]^\chi \\
\bigoplus_{\mathfrak a}
\bigoplus_{\varphi \in \Hom_R(\mathfrak a, M)}
R \ar[rrd] & &
\bigoplus_{\mathfrak a}
\bigoplus_{\psi \in \Hom_R(\mathfrak a, N)}
\mathfrak{a} \ar[r] \ar[d] & N \\
& & \bigoplus_{\mathfrak a}
\bigoplus_{\psi \in \Hom_R(\mathfrak a, N)}
R
}
$$
它诱导出所需映射 $\mathbf{M}(M) \to \mathbf{M}(N)$。
中间那条东南向箭头通过 $\text{id}_{\mathfrak a}$，把与 $\varphi$
对应的直和项 $\mathfrak a$ 映到与
$\psi = \chi \circ \varphi$ 对应的直和项 $\mathfrak a$。
下方东南向箭头同理。略去细节。
\end{proof}

\noindent
现在的思路是超限多次地应用函子 $\mathbf{M}$。对每个序数 $\alpha$，
我们将在 $R$-模范畴上定义函子 $\mathbf{M}_\alpha$，并同时给出自然单射
$N \to \mathbf{M}_\alpha(N)$。用超限递归完成此事。首先，令
$\mathbf{M}_1 = \mathbf{M}$，即上述函子。现在设给定序数 $\alpha$，
并且对所有 $\alpha' < \alpha$ 都已定义 $\mathbf{M}_{\alpha'}$。
若 $\alpha$ 有直接前驱 $\widetilde{\alpha}$，则令
$$
\mathbf{M}_\alpha = \mathbf{M} \circ \mathbf{M}_{\widetilde{\alpha}}.
$$
否则，即若 $\alpha$ 是极限序数，则令
$$
\mathbf{M}_{\alpha}(N) =
\colim_{\alpha' < \alpha} \mathbf{M}_{\alpha'}(N).
$$
显然（例如通过归纳），各 $\mathbf{M}_{\alpha}(N)$ 构成序数上的归纳系统，
故这个定义是合理的。

\begin{theorem}
\label{injectives-theorem-baer-grothendieck}
设 $\kappa$ 为 $R$ 的理想集合的基数，$\alpha$ 为共尾数大于
$\kappa$ 的序数。则 $\mathbf{M}_\alpha(N)$ 是内射 $R$-模，且
$N \to \mathbf{M}_\alpha(N)$ 是函子性的内射嵌入。
\end{theorem}

\begin{proof}
由 Baer 判别法（引理 \ref{injectives-lemma-criterion-baer}），只需证明：
若 $\mathfrak{a} \subset R$ 为理想，则任意映射
$f : \mathfrak{a} \to \mathbf{M}_\alpha(N)$ 都可延拓为
$R \to \mathbf{M}_\alpha(N)$。由于 $\alpha$ 是极限序数，我们有
$$
\mathbf{M}_{\alpha}(N) =
\colim_{\beta < \alpha} \mathbf{M}_{\beta}(N),
$$
因此由命题 \ref{injectives-proposition-modules-are-small} 得到
$$
\Hom_R(\mathfrak{a}, \mathbf{M}_{\alpha}(N)) =
\colim_{\beta < \alpha} \Hom_R(\mathfrak a, \mathbf{M}_{\beta}(N)).
$$
特别地，存在某个 $\beta' < \alpha$，使得 $f$ 分解通过子模
$\mathbf{M}_{\beta'}(N)$，即
$$
f : \mathfrak{a} \to \mathbf{M}_{\beta'}(N) \to
\mathbf{M}_{\alpha}(N).
$$
然而，由函子 $\mathbf{M}$ 的基本性质（见引理
\ref{injectives-lemma-construction} 的 (3)），映射
$\mathfrak{a} \to \mathbf{M}_{\beta'}(N)$ 可以延拓为
$$
R \to \mathbf{M}( \mathbf{M}_{\beta'}(N)) =
\mathbf{M}_{\beta' + 1}(N),
$$
而最后一个对象嵌入 $\mathbf{M}_{\alpha}(N)$（因为 $\alpha$
是极限序数，所以 $\beta' + 1 < \alpha$）。特别地，$f$ 可以延拓到
$\mathbf{M}_{\alpha}(N)$。
\end{proof}




\section{$G$-模}
\label{injectives-section-G-modules}

\noindent
稍后我们将看到（《微分分次代数》，第
\ref{dga-section-modules-noncommutative} 节），代数上模的范畴具有函子性的
内射嵌入。该构造与《代数详论》，第
\ref{more-algebra-section-injectives-modules} 节中的构造完全相同。

\begin{lemma}
\label{injectives-lemma-G-modules}
设 $G$ 为拓扑群，$R$ 为环。$R\text{-}G$-模的范畴
$\text{Mod}_{R, G}$（见《平展上同调》，定义
\ref{etale-cohomology-definition-G-module-continuous}）具有函子性的内射嵌入。
特别地，离散 $G$-模范畴也具有此性质。
\end{lemma}

\begin{proof}
由引理之前的说明，范畴 $\text{Mod}_{R[G]}$ 具有函子性的内射嵌入。
考虑遗忘函子
$v : \text{Mod}_{R, G} \to \text{Mod}_{R[G]}$。
该函子全忠实，把单射变为单射，并且具有右伴随，即
$$
u : M \mapsto u(M) = \{x \in M \mid x\text{ 的稳定子是开子群}\}
$$
由于 $v(M) = 0 \Rightarrow M = 0$，由《同调代数》，引理
\ref{homology-lemma-adjoint-functorial-injectives} 得到结论。
\end{proof}



\section{空间上的阿贝尔层}
\label{injectives-section-abelian-sheaves-space}


\begin{lemma}
\label{injectives-lemma-abelian-sheaves-space}
设 $X$ 为拓扑空间。$X$ 上阿贝尔层的范畴有足够多的内射对象。
事实上，它具有函子性的内射嵌入。
\end{lemma}

\begin{proof}
对阿贝尔群 $A$，记 $j : A \to J(A)$ 为《代数详论》，第
\ref{more-algebra-section-injectives-modules} 节中构造的函子性内射嵌入。
设 $\mathcal{F}$ 为 $X$ 上的阿贝尔层。由《层论》，例
\ref{sheaves-example-sheaf-product-pointwise}，赋值
$$
\mathcal{I} : U \mapsto
\mathcal{I}(U) = \prod\nolimits_{x\in U} J(\mathcal{F}_x)
$$
是阿贝尔层。存在典范映射 $\mathcal{F} \to \mathcal{I}$，它把
$s \in \mathcal{F}(U)$ 映到 $\prod_{x \in U} j(s_x)$，其中
$s_x \in \mathcal{F}_x$ 表示 $s$ 在 $x$ 处的芽。这个映射是单射；
例如见《层论》，引理 \ref{sheaves-lemma-sheaf-subset-stalks}。

\medskip\noindent
还需证明如下断言：给定规则 $x \mapsto I_x$，它为每个点 $x \in X$
指定一个内射阿贝尔群，则层
$\mathcal{I} : U \mapsto \prod_{x \in U} I_x$ 是内射的。注意
$$
\mathcal{I} = \prod\nolimits_{x \in X} i_{x, *}I_x
$$
是摩天楼层 $i_{x, *}I_x$ 的积（记号见《层论》，第
\ref{sheaves-section-skyscraper-sheaves} 节）。我们有
$$
\Mor_{\textit{Ab}}(\mathcal{F}_x, I_x)
=
\Mor_{\textit{Ab}(X)}(\mathcal{F}, i_{x, *}I_x).
$$
见《层论》，引理 \ref{sheaves-lemma-stalk-skyscraper-adjoint}。因此显然
每个 $i_{x, *}I_x$ 都是内射的。于是 $\mathcal{I}$ 的内射性由
《同调代数》，引理 \ref{homology-lemma-product-injectives} 得出。
\end{proof}









\section{环空间上的模层}
\label{injectives-section-sheaves-modules-space}


\begin{lemma}
\label{injectives-lemma-sheaves-modules-space}
设 $(X, \mathcal{O}_X)$ 为环空间，见《层论》，第
\ref{sheaves-section-ringed-spaces} 节。$X$ 上 $\mathcal{O}_X$-模层的范畴
有足够多的内射对象。事实上，它具有函子性的内射嵌入。
\end{lemma}

\begin{proof}
对任意环 $R$ 和任意 $R$-模 $M$，记 $j : M \to J_R(M)$ 为
《代数详论》，第 \ref{more-algebra-section-injectives-modules} 节中构造的
函子性内射嵌入。设 $\mathcal{F}$ 为 $X$ 上的 $\mathcal{O}_X$-模层。
由《层论》，例 \ref{sheaves-example-sheaf-product-pointwise} 和
\ref{sheaves-example-sheaf-product-pointwise-algebraic-structure}，赋值
$$
\mathcal{I} : U \mapsto
\mathcal{I}(U) = \prod\nolimits_{x\in U} J_{\mathcal{O}_{X, x}}(\mathcal{F}_x)
$$
是阿贝尔层。存在典范映射 $\mathcal{F} \to \mathcal{I}$，它把
$s \in \mathcal{F}(U)$ 映到 $\prod_{x \in U} j(s_x)$，其中
$s_x \in \mathcal{F}_x$ 表示 $s$ 在 $x$ 处的芽。这个映射是单射；
例如见《层论》，引理 \ref{sheaves-lemma-sheaf-subset-stalks}。

\medskip\noindent
还需证明如下断言：给定规则 $x \mapsto I_x$，它为每个点 $x \in X$
指定一个内射 $\mathcal{O}_{X, x}$-模，则层
$\mathcal{I} : U \mapsto \prod_{x \in U} I_x$ 是内射的。注意
$$
\mathcal{I} = \prod\nolimits_{x \in X} i_{x, *}I_x
$$
是摩天楼层 $i_{x, *}I_x$ 的积（记号见《层论》，第
\ref{sheaves-section-skyscraper-sheaves} 节）。我们有
$$
\Hom_{\mathcal{O}_{X, x}}(\mathcal{F}_x, I_x)
=
\Hom_{\mathcal{O}_X}(\mathcal{F}, i_{x, *}I_x).
$$
见《层论》，引理 \ref{sheaves-lemma-stalk-skyscraper-adjoint}。因此显然
每个 $i_{x, *}I_x$ 都是内射 $\mathcal{O}_X$-模（见《同调代数》，引理
\ref{homology-lemma-adjoint-preserve-injectives}，或直接论证）。于是
$\mathcal{I}$ 的内射性由《同调代数》，引理
\ref{homology-lemma-product-injectives} 得出。
\end{proof}













\section{范畴上的阿贝尔预层}
\label{injectives-section-injectives-presheaves}

\noindent
设 $\mathcal{C}$ 为范畴。回忆，这意味着 $\Ob(\mathcal{C})$ 是集合。
一方面，考虑 $\mathcal{C}$ 上的阿贝尔预层，见《站点》，第
\ref{sites-section-presheaves} 节。另一方面，考虑由
$\Ob(\mathcal{C})$ 的元素作指标的阿贝尔群族；换言之，考虑对象集为
$\Ob(\mathcal{C})$ 的离散范畴上的预层。我们把这个离散范畴简称为
$\Ob(\mathcal{C})$。存在自然函子
$$
i : \Ob(\mathcal{C}) \longrightarrow \mathcal{C}
$$
因而存在自然的限制函子（或遗忘函子）
$$
v = i^p :
\textit{PAb}(\mathcal{C})
\longrightarrow
\textit{PAb}(\Ob(\mathcal{C}))
$$
比较《站点》，第 \ref{sites-section-functoriality-PSh} 节。我们用 $B$
表示 $\mathcal{C}$ 上的预层，用 $A$ 表示 $\Ob(\mathcal{C})$ 上的预层。

\medskip\noindent
还有两个函子，即 $i_p$ 和 ${}_pi$；它们给 $\Ob(\mathcal{C})$ 上的
阿贝尔预层指派 $\mathcal{C}$ 上的阿贝尔预层，见《站点》，第
\ref{sites-section-functoriality-PSh} 节和第
\ref{sites-section-more-functoriality-PSh} 节。这里我们将使用
$u = {}_pi$；在当前情形下，它定义如下：
$$
uA(U) = \prod\nolimits_{U' \to U} A(U').
$$
因此，一个元素是族 $(a_\phi)_\phi$，其中 $\phi$ 遍历
$\mathcal{C}$ 中所有以 $U$ 为目标的态射。$uA$ 上对应于
$g : V \to U$ 的限制映射把元素 $(a_\phi)_\phi$ 映到元素
$(a_{g \circ \psi})_\psi$。

\medskip\noindent
存在典范满射 $vuA \to A$ 和典范单射 $B \to uvB$。留给读者证明
$$
\Mor_{\textit{PAb}(\mathcal{C})}(B, uA)
=
\Mor_{\textit{PAb}(\Ob(\mathcal{C}))}(vB, A).
$$
在这个简单情形下成立；一般情形见《站点》，第
\ref{sites-section-functoriality-PSh} 节。因此，$(u, v)$ 是一对伴随函子的
例子，见《范畴》，第 \ref{categories-section-adjoint} 节。

\medskip\noindent
此时可以列出关于上述情形的如下事实。
\begin{enumerate}
\item 函子 $u$ 和 $v$ 是正合的。这由上面对这些函子的显式描述得出。
\item 特别地，函子 $v$ 把单射变为单射。
\item 范畴 $\textit{PAb}(\Ob(\mathcal{C}))$ 有足够多的内射对象。
\item 事实上，存在如《同调代数》，定义
\ref{homology-definition-functorial-injective-embedding} 中所述的函子性内射嵌入
$A \mapsto \big(A \to J(A)\big)$。具体地，可以取 $J(A)$ 为预层
$U\mapsto J(A(U))$，其中 $J(-)$ 是《代数详论》，第
\ref{more-algebra-section-injectives-modules} 节中对环 $\mathbf{Z}$ 构造的函子。
\end{enumerate}
综合以上各点，得到把 $\textit{PAb}(\mathcal{C})$ 的对象 $B$ 嵌入内射对象的
如下方法：$B \to uJ(vB)$。见《同调代数》，引理
\ref{homology-lemma-adjoint-functorial-injectives}。

\begin{proposition}
\label{injectives-proposition-presheaves-injectives}
对范畴上的阿贝尔预层，存在函子性的内射嵌入。
\end{proposition}

\begin{proof}
见以上讨论。
\end{proof}












\section{站点上的阿贝尔层}
\label{injectives-section-injectives-sheaves}

\noindent
设 $\mathcal{C}$ 为站点。本节证明 $\mathcal{C}$ 上的阿贝尔层有足够多的
内射对象。

\medskip\noindent
记
$i : \textit{Ab}(\mathcal{C}) \longrightarrow \textit{PAb}(\mathcal{C})$
为从阿贝尔层到阿贝尔预层的遗忘函子。令
${}^\# : \textit{PAb}(\mathcal{C}) \longrightarrow \textit{Ab}(\mathcal{C})$
表示层化函子。回忆，${}^\#$ 是 $i$ 的左伴随，${}^\#$ 是正合的，并且对
任意阿贝尔层 $\mathcal{F}$ 都有 $i\mathcal{F}^\# = \mathcal{F}$。最后，
记 $\mathcal{G} \to J(\mathcal{G})$ 为第
\ref{injectives-section-injectives-presheaves} 节中得到的嵌入内射预层的典范嵌入。

\medskip\noindent
对 $\textit{Ab}(\mathcal{C})$ 中的任意层 $\mathcal{F}$ 和任意序数
$\beta$，通过超限递归定义层 $J_\beta(\mathcal{F})$。置
$J_0(\mathcal{F}) = \mathcal{F}$，并定义
$J_1(\mathcal{F}) = J(i\mathcal{F})^\#$。对典范映射
$i\mathcal{F} \to J(i\mathcal{F})$ 作层化，得到函子性映射
$$
\mathcal{F} \longrightarrow J_1(\mathcal{F})
$$
由于 $\#$ 正合，该映射是单射。置
$J_{\alpha + 1}(\mathcal{F}) = J_1(J_\alpha(\mathcal{F}))$。于是存在典范单射
$J_\alpha(\mathcal{F}) \to J_{\alpha + 1}(\mathcal{F})$。对极限序数
$\beta$，定义
$$
J_\beta(\mathcal{F}) = \colim_{\alpha < \beta} J_\alpha(\mathcal{F}).
$$
注意，这是有向余极限。因此，对任意序数 $\alpha < \beta$，都有单射
$J_\alpha(\mathcal{F}) \to J_\beta(\mathcal{F})$。

\begin{lemma}
\label{injectives-lemma-map-into-next-one}
采用上述记号。设 $\mathcal{G}_1 \to \mathcal{G}_2$ 是
$\mathcal{C}$ 上阿贝尔层的单射。设 $\alpha$ 为序数，并设
$\mathcal{G}_1 \to J_\alpha(\mathcal{F})$ 为层的态射。存在态射
$\mathcal{G}_2 \to J_{\alpha + 1}(\mathcal{F})$，使下图交换：
$$
\xymatrix{
\mathcal{G}_1 \ar[d] \ar[r] & \mathcal{G}_2 \ar[d] \\
J_{\alpha}(\mathcal{F}) \ar[r] & J_{\alpha + 1}(\mathcal{F}) }
$$
\end{lemma}

\begin{proof}
这是因为映射 $i\mathcal{G}_1 \to i\mathcal{G}_2$ 是单射，因而
$i\mathcal{G}_1 \to iJ_\alpha(\mathcal{F})$ 可以延拓为
$i\mathcal{G}_2 \to J(iJ_\alpha(\mathcal{F}))$；应用层化函子后，便得到
所需映射。
\end{proof}

\noindent
这个引理表明，在某种意义下，系统 $\{J_{\alpha}(\mathcal{F})\}$ 是
$\mathcal{F}$ 的内射嵌入。当然，我们不能对所有 $\alpha$ 取极限，因为
它们构成类而非集合。不过，结合下一个引理，现有思路是不必对所有单射
$\mathcal{G}_1 \to \mathcal{G}_2$ 检验内射性。

\begin{lemma}
\label{injectives-lemma-map-into-smaller}
设 $\mathcal{G}_i$（$i\in I$）是 $\mathcal{C}$ 上阿贝尔层的一个集合。
存在序数 $\beta$，使得对任意层 $\mathcal{F}$、任意 $i\in I$ 及任意映射
$\varphi : \mathcal{G}_i \to J_\beta(\mathcal{F})$，都存在
$\alpha < \beta$，使 $ \varphi $ 通过 $J_\alpha(\mathcal{F})$ 分解。
\end{lemma}

\begin{proof}
取所有 $\mathcal{G}_i$ 的直和，即可将问题约化到单个层
$\mathcal{G}$ 的情形。

\medskip\noindent
考虑集合
$$
S = \coprod\nolimits_{U \in \Ob(\mathcal{C})} \mathcal{G}(U).
$$
以及
$$
T_\beta
=
\coprod\nolimits_{U \in \Ob(\mathcal{C})} J_\beta(\mathcal{F})(U)
$$
集合 $T_\beta$ 之间的转移映射都是单射。如果 $\beta$ 的共尾数足够大，
则 $T_\beta = \colim_{\alpha < \beta} T_\alpha$，见《站点》，引理
\ref{sites-lemma-colimit-over-ordinal-sections}。态射
$\mathcal{G} \to J_\beta(\mathcal{F})$ 通过
$J_\alpha(\mathcal{F})$ 分解，当且仅当相应映射 $S \to T_\beta$ 通过
$T_\alpha$ 分解。由《集合论》，引理
\ref{sets-lemma-map-from-set-lifts}，如果 $\beta$ 的共尾数大于 $S$ 的基数，
则本引理的结论成立。因此，本引理由存在共尾数任意大的序数这一事实得出；
见《集合论》，命题 \ref{sets-proposition-exist-ordinals-large-cofinality}。
\end{proof}

\noindent
回忆，对 $\mathcal{C}$ 的对象 $X$，记 $\mathbf{Z}_X$ 为阿贝尔群预层
$\Gamma(U, \mathbf{Z}_X) = \oplus_{U \to X} \mathbf{Z}$；见
《站点上的模》，第 \ref{sites-modules-section-free-abelian-presheaf} 节。
与这个预层相伴的层记为 $\mathbf{Z}_X^\#$；见《站点上的模》，第
\ref{sites-modules-section-free-abelian-sheaf} 节。它可由如下性质刻画：
\begin{equation}
\label{injectives-equation-free-sheaf-on}
\Mor_{\textit{Ab}(\mathcal{C})}(\mathbf{Z}_X^\#, \mathcal{G})
=
\mathcal{G}(X)
\end{equation}
其中，左端的元素 $\varphi$ 映到右端的
$\varphi(1 \cdot \text{id}_X)$。可以用这些层刻画内射阿贝尔层。

\begin{lemma}
\label{injectives-lemma-characterize-injectives}
设 $\mathcal{J}$ 是具有如下性质的阿贝尔群层：对所有
$X\in \Ob(\mathcal{C})$、任意阿贝尔子层
$\mathcal{S} \subset \mathbf{Z}_X^\#$ 和任意态射
$\varphi : \mathcal{S} \to \mathcal{J}$，都存在延拓 $\varphi$ 的态射
$\mathbf{Z}_X^\# \to \mathcal{J}$。则 $\mathcal{J}$ 是内射阿贝尔群层。
\end{lemma}

\begin{proof}
设 $\mathcal{F} \to \mathcal{G}$ 为阿贝尔层的单射，并设
$\varphi : \mathcal{F} \to \mathcal{J}$ 为态射。仿照《代数详论》，引理
\ref{more-algebra-lemma-injective-abelian} 的证明，只需证明：如果
$\mathcal{F} \not = \mathcal{G}$，则可以找到阿贝尔层
$\mathcal{F}'$，满足 $\mathcal{F} \subset \mathcal{F}' \subset \mathcal{G}$，
使得 (a) 包含 $\mathcal{F} \subset \mathcal{F}'$ 是严格的，并且
(b) $\varphi$ 可以延拓到 $\mathcal{F}'$。为寻找 $\mathcal{F}'$，取
$\mathcal{C}$ 的对象 $X$，使包含
$\mathcal{F}(X) \subset \mathcal{G}(X)$ 是严格的。选取
$s \in \mathcal{G}(X)$，且 $s \not \in \mathcal{F}(X)$。令
$\psi : \mathbf{Z}_X^\# \to \mathcal{G}$ 为通过
(\ref{injectives-equation-free-sheaf-on}) 与截面 $s$ 对应的态射，并置
$\mathcal{S} = \psi^{-1}(\mathcal{F})$。由假设，态射
$$
\mathcal{S} \xrightarrow{\psi} \mathcal{F} \xrightarrow{\varphi} \mathcal{J}
$$
可以延拓为态射 $\varphi' : \mathbf{Z}_X^\# \to \mathcal{J}$。注意，
$\varphi'$ 零化 $\psi$ 的核（因为 $\varphi$ 具有这一性质）。因此，
$\varphi'$ 诱导态射 $\varphi'' : \Im(\psi) \to \mathcal{J}$；按构造，
它在交 $\mathcal{F} \cap \Im(\psi)$ 上与 $\varphi$ 一致。于是可以粘合
$\varphi$ 和 $\varphi''$，得到 $\varphi$ 到严格更大的子层
$\mathcal{F}' = \mathcal{F} + \Im(\psi)$ 上的延拓。
\end{proof}

\begin{theorem}
\label{injectives-theorem-sheaves-injectives}
站点上的阿贝尔群层范畴有足够多的内射对象。事实上，存在函子性的内射嵌入；
见《同调代数》，定义
\ref{homology-definition-functorial-injective-embedding}。
\end{theorem}

\begin{proof}
取一组阿贝尔层 $\mathcal{G}_i$（$i \in I$），使每个
$\mathbf{Z}_X^\#$ 的每个子层都作为某个 $\mathcal{G}_i$ 出现。对这个集合
应用引理 \ref{injectives-lemma-map-into-smaller}，得到序数 $\beta$。我们断言，对任意
阿贝尔群层 $\mathcal{F}$，映射
$\mathcal{F} \to J_\beta(\mathcal{F})$ 都是把 $\mathcal{F}$ 嵌入某个
内射对象的单射。注意，按构造，赋值
$\mathcal{F} \mapsto \big(\mathcal{F} \to J_\beta(\mathcal{F})\big)$
确实是函子性的。

\medskip\noindent
该断言的证明如下。由引理 \ref{injectives-lemma-characterize-injectives}，只需把从某个
$\mathbf{Z}_X^\#$ 的子层 $\mathcal{G}$ 出发的任意态射
$\gamma : \mathcal{G} \to J_\beta(\mathcal{F})$ 延拓到整个
$\mathbf{Z}_X^\#$。由引理 \ref{injectives-lemma-map-into-smaller}，映射 $\gamma$
可提升到某个 $J_\alpha(\mathcal{F})$，其中 $\alpha < \beta$。最后，
应用引理 \ref{injectives-lemma-map-into-next-one}，得到所需的 $\gamma$ 的延拓，
它取值于 $J_{\alpha + 1}(\mathcal{F}) \to J_\beta(\mathcal{F})$。
\end{proof}






\section{环站点上的模}
\label{injectives-section-sheaves-modules}

\noindent
设 $\mathcal{C}$ 为站点，$\mathcal{O}$ 为 $\mathcal{C}$ 上的环层。
仿照《代数详论》，第 \ref{more-algebra-section-injectives-modules} 节，
我们尝试证明有足够多的内射 $\mathcal{O}$-模。首先选取内射嵌入
$$
\bigoplus\nolimits_{U, \mathcal{I}}
j_{U!}\mathcal{O}_U/\mathcal{I}
\longrightarrow
\mathcal{J}
$$
其中 $\mathcal{J}$ 是内射阿贝尔层（其存在性由上一节得出）。这里的直和
遍历 $\mathcal{C}$ 的所有对象 $U$ 及所有 $\mathcal{O}$-子模
$\mathcal{I} \subset j_{U!}\mathcal{O}_U$。关于局部化函子
$j_U : \mathcal{C}/U \to \mathcal{C}$ 的限制函子和以 $0$ 延拓函子，
见《站点上的模》，第 \ref{sites-modules-section-localize} 节。

\medskip\noindent
对任意 $\mathcal{O}$-模层 $\mathcal{F}$，记
$$
\mathcal{F}^\vee
=
\SheafHom(\mathcal{F}, \mathcal{J})
$$
并赋予其自然的 $\mathcal{O}$-模结构。
在此插入关于内 Hom 的未来参考文献。
我们还需要 $\mathcal{O}$-模层的一个典范平坦消解。可作如下构造：
对任意 $\mathcal{O}$-模 $\mathcal{F}$，记
$$
F(\mathcal{F})
=
\bigoplus\nolimits_{U \in \Ob(\mathcal{C}), s \in \mathcal{F}(U)}
j_{U!}\mathcal{O}_U.
$$
这是平坦 $\mathcal{O}$-模层，并带有典范满射
$F(\mathcal{F}) \to \mathcal{F}$；见《站点上的模》，引理
\ref{sites-modules-lemma-module-quotient-flat}。此外，构造
$\mathcal{F} \mapsto F(\mathcal{F})$ 关于 $\mathcal{F}$ 是函子性的。

\begin{lemma}
\label{injectives-lemma-vee-exact-sheaves}
函子 $\mathcal{F} \mapsto \mathcal{F}^\vee$ 是正合的。
\end{lemma}

\begin{proof}
这是因为 $\mathcal{J}$ 是内射阿贝尔层。
\end{proof}

\noindent
存在由求值给出的典范映射
$ev : \mathcal{F} \to (\mathcal{F}^\vee)^\vee$：给定
$x \in \mathcal{F}(U)$，令
$ev(x) \in (\mathcal{F}^\vee)^\vee =
\SheafHom(\mathcal{F}^\vee, \mathcal{J})$ 为映射
$\varphi \mapsto \varphi(x)$。

\begin{lemma}
\label{injectives-lemma-ev-injective-sheaves}
对任意 $\mathcal{O}$-模 $\mathcal{F}$，求值映射
$ev : \mathcal{F} \to (\mathcal{F}^\vee)^\vee$ 是单射。
\end{lemma}

\begin{proof}
这可由 $\mathcal{J}$ 的定义检验。具体地，如果
$s \in \mathcal{F}(U)$ 非零，则令
$j_{U!}\mathcal{O}_U \to \mathcal{F}$ 为它通过伴随对应的
$\mathcal{O}$-模映射，并令 $\mathcal{I}$ 为这个映射的核。存在不零化
$s$ 的非零映射
$\mathcal{F} \supset j_{U!}\mathcal{O}_U/\mathcal{I}
\to \mathcal{J}$。由于 $\mathcal{J}$ 是内射 $\mathcal{O}$-模，该映射可延拓为
$\varphi : \mathcal{F} \to \mathcal{J}$。于是
$ev(s)(\varphi) = \varphi(s) \not = 0$，这正是所要证明的。
\end{proof}

\noindent
上述 $\mathcal{O}$-模的典范满射
$F(\mathcal{F}) \to \mathcal{F}$ 变为如下 $\mathcal{O}$-模的典范单射：
$$
(\mathcal{F}^\vee)^\vee \longrightarrow (F(\mathcal{F}^\vee))^\vee.
$$
置 $J(\mathcal{F}) = (F(\mathcal{F}^\vee))^\vee$。将 $ev$ 与这个
显示映射复合，得到关于 $\mathcal{F}$ 函子性的映射
$\mathcal{F} \to J(\mathcal{F})$。

\begin{lemma}
\label{injectives-lemma-JM-injective-sheaves}
设 $\mathcal{O}$ 为环层。对每个 $\mathcal{O}$-模 $\mathcal{F}$，
$\mathcal{O}$-模 $J(\mathcal{F})$ 是内射的。
\end{lemma}

\begin{proof}
需要证明函子 $\Hom_\mathcal{O}(\mathcal{G}, J(\mathcal{F}))$ 是正合的。
注意
\begin{eqnarray*}
\Hom_\mathcal{O}(\mathcal{G}, J(\mathcal{F}))
& = &
\Hom_\mathcal{O}(\mathcal{G}, (F(\mathcal{F}^\vee))^\vee) \\
& = &
\Hom_\mathcal{O}
(\mathcal{G}, \SheafHom(F(\mathcal{F}^\vee), \mathcal{J})) \\
& = &
\Hom(\mathcal{G} \otimes_\mathcal{O} F(\mathcal{F}^\vee), \mathcal{J})
\end{eqnarray*}
因此，所需结论由 $F(\mathcal{F}^\vee)$ 平坦且 $\mathcal{J}$ 内射这一事实得出。
\end{proof}

\begin{theorem}
\label{injectives-theorem-sheaves-modules-injectives}
设 $\mathcal{C}$ 为站点，$\mathcal{O}$ 为 $\mathcal{C}$ 上的环层。
站点上的 $\mathcal{O}$-模层范畴有足够多的内射对象。事实上，存在函子性的
内射嵌入；见《同调代数》，定义
\ref{homology-definition-functorial-injective-embedding}。
\end{theorem}

\begin{proof}
由本节的讨论即得。
\end{proof}

\begin{proposition}
\label{injectives-proposition-presheaves-modules}
设 $\mathcal{C}$ 为范畴，$\mathcal{O}$ 为 $\mathcal{C}$ 上的环预层。
$\mathcal{O}$-模预层的范畴 $\textit{PMod}(\mathcal{O})$ 具有函子性的
内射嵌入。
\end{proposition}

\begin{proof}
可以沿用第 \ref{injectives-section-injectives-presheaves} 节的讨论证明此命题，
但这里改用上面的定理。赋予 $\mathcal{C}$ 如下站点结构：对象 $U$ 的覆盖
集合由所有单元素集 $\{f : V \to U\}$ 构成，其中 $f$ 是同构。略去它确实
定义站点的验证。对这个拓扑，层与预层相同（略去证明）。故上面的定理适用。
\end{proof}








\section{阿贝尔范畴的嵌入}
\label{injectives-section-embedding}

\noindent
本节证明，阿贝尔范畴可嵌入某个有足够多点的站点上的阿贝尔层范畴。
该站点由下一个引理描述。

\begin{lemma}
\label{injectives-lemma-site-abelian-category}
设 $\mathcal{A}$ 为阿贝尔范畴。令
$$
\text{Cov} = \{\{f : V \to U\} \mid f\text{ 为满射}\}.
$$
则 $(\mathcal{A}, \text{Cov})$ 是站点；见《站点》，定义
\ref{sites-definition-site}。
\end{lemma}

\begin{proof}
注意，按我们关于范畴的约定，$\Ob(\mathcal{A})$ 是集合。同构是满射态射，
满射态射的复合仍为满射态射，而 $\mathcal{A}$ 中满射态射的基变换仍为
满射态射；见《同调代数》，引理
\ref{homology-lemma-epimorphism-universal-abelian-category}。
\end{proof}

\noindent
设 $\mathcal{A}$ 为预加性范畴。在此情形下，Yoneda 嵌入
$\mathcal{A} \to \textit{PSh}(\mathcal{A})$，$X \mapsto h_X$，通过函子
$\mathcal{A} \to \textit{PAb}(\mathcal{A})$ 分解。

\begin{lemma}
\label{injectives-lemma-embedding}
设 $\mathcal{A}$ 为阿贝尔范畴，并令
$\mathcal{C} = (\mathcal{A}, \text{Cov})$ 为引理
\ref{injectives-lemma-site-abelian-category} 中定义的站点。则 $X \mapsto h_X$ 定义
全忠实正合函子
$$
\mathcal{A} \longrightarrow \textit{Ab}(\mathcal{C}).
$$
此外，站点 $\mathcal{C}$ 有足够多的点。
\end{lemma}

\begin{proof}
设 $f : V \to U$ 是 $\mathcal{A}$ 的满射态射，并令
$K = \Ker(f)$。回忆，
$V \times_U V = \Ker((f, -f) : V \oplus V \to U)$；见《同调代数》，
例 \ref{homology-example-fibre-product-pushouts}。特别地，存在单射
$K \oplus K \to V \times_U V$。令
$p, q : V \times_U V \to V$ 为两个投影态射。注意，
$p - q : V \times_U V \to V$ 满足 $f \circ (p  - q) = 0$。因此
$p - q$ 通过 $K \to V$ 分解。将这个态射记为
$c : V \times_U V \to K$。由于复合
$K \oplus K \to V \times_U V \to K$ 是满射，故 $c$ 是满射。于是
$$
V \times_U V \xrightarrow{p - q} V \to U \to 0
$$
是 $\mathcal{A}$ 中的正合列。因此，对 $\mathcal{A}$ 的对象 $X$，序列
$$
0 \to
\Hom_\mathcal{A}(U, X) \to
\Hom_\mathcal{A}(V, X) \to
\Hom_\mathcal{A}(V \times_U V, X)
$$
是阿贝尔群的正合列；见《同调代数》，引理
\ref{homology-lemma-check-exactness}。这意味着 $h_X$ 满足
$\mathcal{C}$ 上的层条件。

\medskip\noindent
由《范畴》，引理 \ref{categories-lemma-yoneda}，该函子全忠实。
由《同调代数》，引理 \ref{homology-lemma-check-exactness}，它是阿贝尔
范畴之间的左正合函子。为证明它右正合，设 $X \to Y$ 是
$\mathcal{A}$ 的满射态射。设 $U$ 为 $\mathcal{A}$ 的对象，并设
$s \in h_Y(U) = \Mor_\mathcal{A}(U, Y)$ 为 $h_Y$ 在 $U$ 上的截面。
由《同调代数》，引理
\ref{homology-lemma-epimorphism-universal-abelian-category}，投影
$U \times_Y X \to U$ 是满射。因此，
$\{V = U \times_Y X \to U\}$ 是 $U$ 的覆盖，并且 $s|_V$ 可提升为
$h_X$ 的截面。这证明 $h_X \to h_Y$ 是阿贝尔层的满射；见《站点》，
引理 \ref{sites-lemma-mono-epi-sheaves}。

\medskip\noindent
由《站点》，命题 \ref{sites-proposition-criterion-points}，站点
$\mathcal{C}$ 有足够多的点。
\end{proof}

\begin{remark}
\label{injectives-remark-embedding}
Freyd--Mitchell 嵌入定理断言，从任意阿贝尔范畴 $\mathcal{A}$ 到某个环上
模范畴存在全忠实正合函子。引理 \ref{injectives-lemma-embedding} 的结论没有这么强。
不过，这个结果适合 Stacks 项目，因为无论如何我们都必须详细理解站点上的
阿贝尔群层。此外，在 $\mathcal{C}$ 上阿贝尔层的范畴中可以进行“追图”；
例如，可以使用对象上的截面，也可以利用 $\mathcal{C}$ 有足够多的点而在茎
的层面工作。关于如何由引理 \ref{injectives-lemma-embedding} 推出 Freyd--Mitchell
嵌入定理，见注 \ref{injectives-remark-embedding-freyd}。
\end{remark}

\begin{remark}
\label{injectives-remark-embedding-big}
如果 $\mathcal{A}$ 是“大”阿贝尔范畴，即 $\mathcal{A}$ 的对象构成类，
则引理 \ref{injectives-lemma-embedding} 不适用。在这种情形下，给定任意对象集
$E \subset \Ob(\mathcal{A})$，存在阿贝尔全子范畴
$\mathcal{A}' \subset \mathcal{A}$，使得 $\Ob(\mathcal{A}')$ 是集合且
$E \subset \Ob(\mathcal{A}')$。于是可以对 $\mathcal{A}'$ 应用引理
\ref{injectives-lemma-embedding}。由此可证明，依赖追图的结果在 $\mathcal{A}$ 中成立。
\end{remark}

\begin{remark}
\label{injectives-remark-embedding-freyd}
设 $\mathcal{C}$ 为站点。注意，$\textit{Ab}(\mathcal{C})$ 有足够多的
内射对象；见定理 \ref{injectives-theorem-sheaves-injectives}。（当
$\mathcal{C}$ 有足够多的点时，这是直接的：如果 $I$ 是内射
$\mathbf{Z}$-模而 $p$ 是一个点，则 $p_*I$ 是内射层。）此外，
$\textit{Ab}(\mathcal{C})$ 有余生成元（略去细节）。因此，引理
\ref{injectives-lemma-embedding} 证明存在全忠实正合嵌入
$\mathcal{A} \to \mathcal{B}$，其中 $\mathcal{B}$ 有余生成元和足够多的
内射对象。把此结论应用于 $\mathcal{A}^{opp}$，得到全忠实正合函子
$i : \mathcal{A} \to \mathcal{D} = \mathcal{B}^{opp}$，其中
$\mathcal{D}$ 有足够多的投射对象和一个生成元。因此，$\mathcal{D}$ 有
投射生成元 $P$。置 $R = \Mor_\mathcal{D}(P, P)$。则
$$
\mathcal{A} \longrightarrow \text{Mod}_R, \quad
X \longmapsto \Hom_\mathcal{D}(P, X).
$$
可以检验这是全忠实正合函子。换言之，我们重新得到上面注
\ref{injectives-remark-embedding} 中提到的 Freyd--Mitchell 定理。
\end{remark}

\begin{remark}
\label{injectives-remark-embed-exact-category}
证明引理 \ref{injectives-lemma-site-abelian-category} 和
\ref{injectives-lemma-embedding} 的论证也适用于{\it 正合范畴}；见
\cite[Appendix A]{Buhler} 和 \cite[1.1.4]{BBD}。这里作简要回顾；
若 Stacks 项目以后需要，我们再补充更多细节。

\medskip\noindent
设 $\mathcal{A}$ 为加性范畴。一个{\it 核—余核对}是
$\mathcal{A}$ 的态射对 $(i, p)$，其中 $i : A \to B$、
$p : B \to C$，且 $i$ 是 $p$ 的核、$p$ 是 $i$ 的余核。给定核—余核对的
集合 $\mathcal{E}$，如果对某个态射 $p$ 有
$(i, p) \in \mathcal{E}$，则称 $i : A \to B$ 为{\it 容许单态射}。
类似地，如果对某个态射 $i$ 有 $(i, p) \in \mathcal{E}$，则称态射
$p : B \to C$ 为{\it 容许满态射}。如果下列公理成立，则称
$(\mathcal{A}, \mathcal{E})$ 为{\it 正合范畴}：
\begin{enumerate}
\item $\mathcal{E}$ 对核—余核对的同构封闭；
\item 对任意对象 $A$，态射 $1_A$ 既是容许满态射又是容许单态射；
\item 容许单态射在复合下稳定；
\item 容许满态射在复合下稳定；
\item 容许单态射 $i : A \to B$ 沿任意态射 $A \to A'$ 的推出存在，
且诱导态射 $i' : A' \to B'$ 是容许单态射；以及
\item 容许满态射 $p : B \to C$ 沿任意态射 $C' \to C$ 的基变换存在，
且诱导态射 $p' : B' \to C'$ 是容许满态射。
\end{enumerate}
给定这样的结构，令 $\mathcal{C} = (\mathcal{A}, \text{Cov})$，其中覆盖
（即 $\text{Cov}$ 的元素）由容许满态射给出。上述公立即刻蕴含这是站点。
考虑函子
$$
F : \mathcal{A} \longrightarrow \textit{Ab}(\mathcal{C}), \quad
X \longmapsto h_X
$$
正如引理 \ref{injectives-lemma-embedding} 中一样。可以证明，这个函子全忠实、正合，
并且反映正合性。此外，$F$ 的本质像中对象的任意扩张仍在 $F$ 的本质像中。
\end{remark}






\section{Grothendieck 的 AB 条件}
\label{injectives-section-grothendieck-conditions}

\noindent
本节及接下来的几节主要与“大”阿贝尔范畴有关，即《范畴》，注
\ref{categories-remark-big-categories} 中列出的那些范畴。一个值得牢记的
例子是环站点上的模层范畴。

\medskip\noindent
Grothendieck 在论文 \cite{Tohoku} 中以极大的一般性证明了内射对象的存在性。
他用下列条件挑出具有特殊性质的阿贝尔范畴。

\begin{definition}
\label{injectives-definition-grothendieck-conditions}
设 $\mathcal{A}$ 为阿贝尔范畴。列出如下条件：
\begin{enumerate}
\item[AB3] $\mathcal{A}$ 有直和；
\item[AB4] $\mathcal{A}$ 满足 AB3，并且直和正合；
\item[AB5] $\mathcal{A}$ 满足 AB3，并且滤过余极限正合。
\end{enumerate}
下面是对偶概念：
\begin{enumerate}
\item[AB3*] $\mathcal{A}$ 有积；
\item[AB4*] $\mathcal{A}$ 满足 AB3*，并且积正合；
\item[AB5*] $\mathcal{A}$ 满足 AB3*，并且余滤过极限正合。
\end{enumerate}
如果对 $\mathcal{A}$ 中每个 $N \subset M$、$N \not = M$，都存在不通过
$N$ 分解的态射 $U \to M$，则称 $\mathcal{A}$ 的对象 $U$ 为{\it 生成元}。
如果 $\mathcal{A}$ 满足 AB5 且有生成元，则称其为
{\it Grothendieck 阿贝尔范畴}。
\end{definition}

\noindent
讨论：阿贝尔范畴中的直和就是余积。如果阿贝尔范畴有直和（即满足 AB3），
则它有余极限；见《范畴》，引理
\ref{categories-lemma-colimits-coproducts-coequalizers}。类似地，如果
$\mathcal{A}$ 满足 AB3*，则它有极限；见《范畴》，引理
\ref{categories-lemma-limits-products-equalizers}。直和的正合性意指如下性质：
给定指标集 $I$ 和短正合列
$$
0 \to A_i \to B_i \to C_i \to 0,\quad i \in I
$$
于 $\mathcal{A}$ 中，则序列
$$
0 \to
\bigoplus\nolimits_{i \in I} A_i \to
\bigoplus\nolimits_{i \in I} B_i \to
\bigoplus\nolimits_{i \in I} C_i \to 0
$$
也正合。不假设 AB4 时，一般只能断言这个序列右端正合（即若直和存在，
取直和是右正合函子）。类似地，滤过余极限的正合性意指如下性质：给定有向集
$I$ 及 $\mathcal{A}$ 中定义在 $I$ 上的一族短正合列
$$
0 \to A_i \to B_i \to C_i \to 0
$$
则序列
$$
0 \to
\colim_{i \in I} A_i \to
\colim_{i \in I} B_i \to
\colim_{i \in I} C_i \to 0
$$
也正合。不假设 AB5 时，一般只能断言这个序列右端正合（即若余极限存在，
取余极限是右正合函子）。AB4* 和 AB5* 有类似解释。



\section{Grothendieck 范畴中的内射对象}
\label{injectives-section-grothendieck-categories}

\noindent
生成元的存在性蕴含：给定 Grothendieck 阿贝尔范畴 $\mathcal{A}$ 的对象
$M$，其子对象构成集合。（对一般的“大”阿贝尔范畴，这未必成立。）

\begin{lemma}
\label{injectives-lemma-set-of-subobjects}
设 $\mathcal{A}$ 为有生成元 $U$ 的阿贝尔范畴，$X$ 为
$\mathcal{A}$ 的对象。如果 $\kappa$ 是 $\Mor(U, X)$ 的基数，则
\begin{enumerate}
\item 不存在由大于 $\kappa$ 的基数作指标的 $X$ 的严格递增（或严格递减）
子对象链；
\item 如果 $\alpha$ 是共尾数 $> \kappa$ 的序数，则由 $\alpha$ 作指标的
$X$ 的任意递增（或递减）子对象列最终为常值；
\item $X$ 的子对象集的基数 $\leq 2^\kappa$。
\end{enumerate}
\end{lemma}

\begin{proof}
对 (1)，设 $\kappa' > \kappa$ 为基数，并设 $X_i$（$i \in \kappa'$）
严格递增。于是，对每个 $i$，取 $\phi_i \in \Mor(U, X)$，使
$\phi_i$ 通过 $X_{i + 1}$ 分解但不通过 $X_i$ 分解。则态射
$\phi_i$ 两两不同，这与 $\kappa$ 的定义矛盾。

\medskip\noindent
(2) 由共尾数的定义和 (1) 得出。

\medskip\noindent
证明 (3)。对任意子对象 $Y \subset X$，定义
$S_Y \in \mathcal{P}(\Mor(U, X))$（幂集）为
$S_Y = \{\phi \in \Mor(U,X) : \phi)\text{ 通过 }Y\text{ 分解}\}$.
则 $Y = Y'$ 当且仅当 $S_Y = S_{Y'}$。因此，子对象集的基数至多是这个
幂集的基数。
\end{proof}

\noindent
由引理 \ref{injectives-lemma-set-of-subobjects}，下述定义是有意义的。

\begin{definition}
\label{injectives-definition-size}
设 $\mathcal{A}$ 为 Grothendieck 阿贝尔范畴，$M$ 为 $\mathcal{A}$ 的对象。
$M$ 的{\it 大小} $|M|$ 是 $M$ 的子对象集的基数。
\end{definition}

\begin{lemma}
\label{injectives-lemma-size-goes-down}
设 $\mathcal{A}$ 为 Grothendieck 阿贝尔范畴。如果
$0 \to M' \to M \to M'' \to 0$ 是 $\mathcal{A}$ 中的短正合列，则
$|M'|, |M''| \leq |M|$。
\end{lemma}

\begin{proof}
由定义立即得出。
\end{proof}

\begin{lemma}
\label{injectives-lemma-set-iso-classes-bounded-size}
设 $\mathcal{A}$ 为有生成元 $U$ 的 Grothendieck 阿贝尔范畴。
\begin{enumerate}
\item 如果 $|M| \leq \kappa$，则 $M$ 是至多 $\kappa$ 个 $U$ 的副本的
直和的商；
\item 对每个基数 $\kappa$，满足 $|M| \leq \kappa$ 的对象 $M$ 的同构类
构成集合。
\end{enumerate}
\end{lemma}

\begin{proof}
对 (1)，为每个真子对象 $M' \subset M$ 选取态射
$\varphi_{M'} : U \to M$，使其像不包含在 $M'$ 中。于是
$\bigoplus_{M' \subset M} \varphi_{M'} : \bigoplus_{M' \subset M} U \to M$
是满射。显然，(1) 蕴含 (2)。
\end{proof}

\begin{proposition}
\label{injectives-proposition-objects-are-small}
设 $\mathcal{A}$ 为 Grothendieck 阿贝尔范畴，$M$ 为
$\mathcal{A}$ 的对象，并令 $\kappa = |M|$。如果 $\alpha$ 是共尾数大于
$\kappa$ 的序数，则 $M$ 相对于单射为 $\alpha$-小的。
\end{proposition}

\begin{proof}
比较命题 \ref{injectives-proposition-modules-are-small}。只需证明映射
(\ref{injectives-equation-compare}) 是满射。设
$f : M \to \colim B_\beta$ 为映射。考虑 $M$ 的子对象
$\{f^{-1}(B_\beta)\}$，其中把 $B_\beta$ 看作余极限
$B = \bigcup_\beta B_\beta$ 的子对象。如果其中某个子对象，例如
$f^{-1}(B_\beta)$，等于 $M$，则该映射通过 $B_\beta$ 分解。

\medskip\noindent
由于 $\mathcal{A}$ 满足 AB5，我们有
$$
\colim f^{-1}(B_\beta) = f^{-1}\left(\colim B_\beta\right) = M.
$$
由假设，在 $f^{-1}(B_\alpha)$ 中出现的 $M$ 的不同子对象至多有
$\kappa$ 个。因此，可以找到基数至多为 $\kappa$ 的子集
$S \subset \alpha$，使得当 $\beta'$ 遍历 $S$ 时，
$f^{-1}(B_{\beta'})$ 遍历\emph{所有} $f^{-1}(B_\alpha)$。

\medskip\noindent
由于 $\alpha$ 的共尾数大于 $\kappa$，$S$ 有上界
$\widetilde{\alpha} < \alpha$。特别地，所有
$f^{-1}(B_{\beta'})$（$\beta' \in S$）都包含在
$f^{-1}(B_{\widetilde{\alpha}})$ 中。因此
$f^{-1}(B_{\widetilde{\alpha}}) = M$。特别地，映射 $f$ 通过
$B_{\widetilde{\alpha}}$ 分解。
\end{proof}

\begin{lemma}
\label{injectives-lemma-characterize-injective}
\begin{slogan}
要检验一个对象是内射的，只需对生成元的子对象检验提升性质。
\end{slogan}
设 $\mathcal{A}$ 为有生成元 $U$ 的 Grothendieck 阿贝尔范畴。
$\mathcal{A}$ 的对象 $I$ 是内射的，当且仅当在每个交换图
$$
\xymatrix{
M \ar[d] \ar[r] &  I \\
U \ar@{-->}[ru]
}
$$
中，只要 $M \subset U$ 是子对象，虚线箭头就存在。
\end{lemma}

\begin{proof}
模的情形见引理 \ref{injectives-lemma-criterion-baer}。选取单射
$A \subset B$ 和态射 $\varphi : A \to I$。考虑由如下对
$(A', \varphi')$ 构成的集合 $S$：$A \subset A' \subset B$ 为子对象，
$\varphi' : A' \to I$ 为 $\varphi$ 的延拓。以显然的方式在这个集合上
定义偏序。选取全序子集 $T \subset S$。则
$$
A' = \colim_{t \in T} A_t \xrightarrow{\colim_{t \in T} \varphi_t} I
$$
是上界。因此，由 Zorn 引理，集合 $S$ 有极大元 $(A', \varphi')$。
我们断言 $A' = B$。否则，选取不通过 $A'$ 分解的态射
$\psi : U \to B$。置 $N = A' \cap \psi(U)$ 和
$M = \psi^{-1}(N)$。则映射
$$
M \to N \to A' \xrightarrow{\varphi'} I
$$
可延拓为态射 $\chi : U \to I$。由于
$\chi|_{\Ker(\psi)} = 0$，$\chi$ 分解为
$$
U \to \Im(\psi) \xrightarrow{\varphi''} I
$$
由于 $\varphi'$ 和 $\varphi''$ 在 $N = A' \cap \Im(\psi)$ 上一致，
它们合起来定义态射 $A' + \Im(\psi) \to I$，这与 $A'$ 的极大性矛盾。
\end{proof}

\begin{theorem}
\label{injectives-theorem-injective-embedding-grothendieck}
设 $\mathcal{A}$ 为 Grothendieck 阿贝尔范畴。则 $\mathcal{A}$ 具有
函子性的内射嵌入。
\end{theorem}

\begin{proof}
比较定理 \ref{injectives-theorem-baer-grothendieck} 的证明。选取
$\mathcal{A}$ 的生成元 $U$。对对象 $M$，用如下推出图定义
$\mathbf{M}(M)$：
$$
\xymatrix{
\bigoplus_{N \subset U}
\bigoplus_{\varphi \in \Hom(N, M)}
N \ar[r] \ar[d] & M \ar[d] \\
\bigoplus_{N \subset U}
\bigoplus_{\varphi \in \Hom(N, M)}
U \ar[r] &  \mathbf{M}(M).
}
$$
注意，$M \mapsto \mathbf{M}(M)$ 是函子，且存在函子性单射
$M \to \mathbf{M}(M)$。通过超限归纳，对每个序数 $\alpha$ 定义函子
$\mathbf{M}_\alpha(M)$。具体地，置 $\mathbf{M}_0(M) = M$；给定
$\mathbf{M}_\alpha(M)$，置
$\mathbf{M}_{\alpha + 1}(M) = \mathbf{M}(\mathbf{M}_\alpha(M))$。
对极限序数 $\beta$，置
$$
\mathbf{M}_\beta(M) = \colim_{\alpha < \beta} \mathbf{M}_\alpha(M).
$$
最后，选取任意共尾数大于 $|U|$ 的序数 $\alpha$。这样的序数存在；见
《集合论》，命题 \ref{sets-proposition-exist-ordinals-large-cofinality}。
我们断言，$M \to \mathbf{M}_\alpha(M)$ 是所需的函子性内射嵌入。
事实上，如果 $N \subset U$ 是子对象，且
$\varphi : N \to \mathbf{M}_\alpha(M)$ 是态射，则由命题
\ref{injectives-proposition-objects-are-small}，$\varphi$ 通过某个
$\mathbf{M}_{\alpha'}(M)$ 分解，其中 $\alpha' < \alpha$。由
$\mathbf{M}(-)$ 的构造，$\varphi$ 可延拓为从 $U$ 到
$\mathbf{M}_{\alpha' + 1}(M)$ 的态射，因而也延拓到
$\mathbf{M}_\alpha(M)$。由引理 \ref{injectives-lemma-characterize-injective}，
得出 $\mathbf{M}_\alpha(M)$ 是内射的。
\end{proof}













\section{Grothendieck 范畴中的 K-内射对象}
\label{injectives-section-K-injective}

\noindent
本节材料取自 Serpé 的论文 \cite{serpe}。该论文把 Spaltenstein 的
\cite{Spaltenstein} 中的一些结果推广到一般 Grothendieck 阿贝尔范畴。
我们的引理 \ref{injectives-lemma-characterize-K-injective} 只隐含在 Serpé 的论文中。
这里的方法是仿照 Grothendieck 对定理
\ref{injectives-theorem-injective-embedding-grothendieck} 的证明。

\begin{lemma}
\label{injectives-lemma-surjection-bounded-size}
设 $\mathcal{A}$ 为有生成元 $U$ 的 Grothendieck 阿贝尔范畴。令 $c$ 为
基数上的函数，定义为
$c(\kappa) = |\bigoplus_{\alpha \in \kappa} U|$。如果
$\pi : M \to N$ 是满射，则存在满射到 $N$ 的子对象 $M' \subset M$，
且 $|M'| \leq c(|N|)$。
\end{lemma}

\begin{proof}
对每个真子对象 $N' \subset N$，选取态射
$\varphi_{N'} : U \to M$，使 $U \to M \to N$ 不通过 $N'$ 分解。置
$$
M' = \Im\left(
\bigoplus\nolimits_{N' \subset N} \varphi_{N'} :
\bigoplus\nolimits_{N' \subset N} U \longrightarrow M\right)
$$
则 $M'$ 满足要求。
\end{proof}

\begin{lemma}
\label{injectives-lemma-acyclic-quotient-complexes-bounded-size}
设 $\mathcal{A}$ 为 Grothendieck 阿贝尔范畴。存在基数 $\kappa$，使得对
任意无环复形 $M^\bullet$，都有
\begin{enumerate}
\item 如果 $M^\bullet$ 非零，则存在非零子复形 $N^\bullet$，它有上界、
无环，且 $|N^n| \leq \kappa$；
\item 存在复形的满射
$$
\bigoplus\nolimits_{i \in I} M_i^\bullet \longrightarrow M^\bullet
$$
其中 $M_i^\bullet$ 有上界、无环，且 $|M_i^n| \leq \kappa$。
\end{enumerate}
\end{lemma}

\begin{proof}
选取 $\mathcal{A}$ 的生成元 $U$。记 $c$ 为引理
\ref{injectives-lemma-surjection-bounded-size} 中的函数。置
$\kappa = \sup \{c^n(|U|), n = 1, 2, 3, \ldots\}$。设
$n \in \mathbf{Z}$，并设 $\psi : U \to M^n$ 为态射。为证明 (1) 和 (2)，
只需证明存在子复形 $N^\bullet \subset M^\bullet$，它有上界、无环且
$|N^m| \leq \kappa$，并使 $\psi$ 通过 $N^n$ 分解。为此，置
$N^n = \Im(\psi)$，
$N^{n + 1} = \Im(U \to M^n \to M^{n + 1})$，并对
$m \geq n + 2$ 置 $N^m = 0$。设对某个 $k \leq n$，我们已经对所有
$n > m \geq k$ 构造了 $N^m \subset M^m$，使得
\begin{enumerate}
\item 对所有 $m \geq k$，$\text{d}(N^m) \subset N^{m + 1}$；
\item 对所有 $m \geq k + 1$，
$\Im(N^{m - 1} \to N^m) = \Ker(N^m \to N^{m + 1})$；以及
\item 对所有 $m \geq k$，
$|N^m| \leq c^{\max\{n - m, 0\}}(|U|)$。
\end{enumerate}
由于 $M^\bullet$ 无环，子对象
$\text{d}^{-1}(\Ker(N^k \to N^{k + 1})) \subset M^{k - 1}$ 满射到
$\Ker(N^k \to N^{k + 1})$。因此，由引理
\ref{injectives-lemma-surjection-bounded-size}，可以选取满射到
$\Ker(N^k \to N^{k + 1})$ 的 $N^{k - 1} \subset M^{k - 1}$，且
$|N^{k - 1}| \leq c^{n - k + 1}(|U|)$。对 $k$ 归纳即完成证明。
\end{proof}

\begin{lemma}
\label{injectives-lemma-characterize-K-injective}
设 $\mathcal{A}$ 为 Grothendieck 阿贝尔范畴，$\kappa$ 为引理
\ref{injectives-lemma-acyclic-quotient-complexes-bounded-size} 中的基数。设
$I^\bullet$ 是满足下列条件的复形：
\begin{enumerate}
\item 每个 $I^j$ 都是内射的；并且
\item 对每个有上界的无环复形 $M^\bullet$，如果
$|M^n| \leq \kappa$，则
$\Hom_{K(\mathcal{A})}(M^\bullet, I^\bullet) = 0$。
\end{enumerate}
则 $I^\bullet$ 是 K-内射复形。
\end{lemma}

\begin{proof}
设 $M^\bullet$ 为无环复形。按如下方式对序数 $\alpha$ 归纳构造无环子复形
$K_\alpha^\bullet \subset M^\bullet$。对 $\alpha = 0$，置
$K_0^\bullet = 0$。对 $\alpha > 0$，作如下处理：
\begin{enumerate}
\item 如果 $\alpha = \beta + 1$ 且 $K_\beta^\bullet = M^\bullet$，
则取 $K_\alpha^\bullet = K_\beta^\bullet$；
\item 如果 $\alpha = \beta + 1$ 且
$K_\beta^\bullet \not = M^\bullet$，则
$M^\bullet/K_\beta^\bullet$ 是非零无环复形。按引理
\ref{injectives-lemma-acyclic-quotient-complexes-bounded-size} 选取子复形
$N_\alpha^\bullet \subset M^\bullet/K_\beta^\bullet$。最后，令
$K_\alpha^\bullet \subset M^\bullet$ 为 $N_\alpha^\bullet$ 的逆像；
\item 如果 $\alpha$ 是极限序数，则置
$K_\beta^\bullet = \colim K_\alpha^\bullet$.
\end{enumerate}
显然，对充分大的序数 $\alpha$，有
$M^\bullet = K_\alpha^\bullet$。我们将对 $\alpha$ 作超限归纳，证明
$$
\Hom_{K(\mathcal{A})}(K_\alpha^\bullet, I^\bullet)
$$
为零。对 $\alpha = 0$，由于 $K_0^\bullet$ 为零，结论成立。设结论对
$\beta$ 成立，且 $\alpha = \beta + 1$。在上面列表的情形 (1) 中，
结论显然；在情形 (2) 中，存在复形的短正合列
$$
0 \to K_\beta^\bullet \to K_\alpha^\bullet \to N_\alpha^\bullet \to 0
$$
由于 $I^\bullet$ 的每个分量都是内射的，我们得到正合列
$$
\Hom_{K(\mathcal{A})}(K_\beta^\bullet, I^\bullet) \leftarrow
\Hom_{K(\mathcal{A})}(K_\alpha^\bullet, I^\bullet) \leftarrow
\Hom_{K(\mathcal{A})}(N_\alpha^\bullet, I^\bullet)
$$
由归纳假设，左端项为零；由关于 $I^\bullet$ 的假设，右端项为零。因此，
中间群也为零。最后，设 $\alpha$ 为极限序数。则
$$
\Hom^\bullet(K_\alpha^\bullet, I^\bullet) =
\lim_{\beta < \alpha} \Hom^\bullet(K_\beta^\bullet, I^\bullet)
$$
其中记号见《代数详论》，第
\ref{more-algebra-section-hom-complexes} 节。由《代数详论》，等式
(\ref{more-algebra-equation-cohomology-hom-complex})，这些复形计算
$K(\mathcal{A})$ 中的态射。注意，因为每个 $I^j$ 都是内射的，系统中的
转移映射都是满射。此外，对极限序数 $\alpha$，极限等于其值（见上面的
显示公式）。因此，可应用《同调代数》，引理
\ref{homology-lemma-ML-over-ordinals} 得到结论。
\end{proof}

\begin{lemma}
\label{injectives-lemma-functorial-homotopies}
设 $\mathcal{A}$ 为 Grothendieck 阿贝尔范畴，并设
$(K_i^\bullet)_{i \in I}$ 为一组无环复形。存在函子
$M^\bullet \mapsto \mathbf{M}^\bullet(M^\bullet)$ 和自然变换
$j_{M^\bullet} : M^\bullet \to \mathbf{M}^\bullet(M^\bullet)$，使得
\begin{enumerate}
\item $j_{M^\bullet}$ 是（逐项）单射的拟同构；并且
\item 对每个 $i \in I$ 和 $w : K_i^\bullet \to M^\bullet$，态射
$j_{M^\bullet} \circ w$ 零伦。
\end{enumerate}
\end{lemma}

\begin{proof}
对每个 $i \in I$，选取复形的（逐项）单射
$K_i^\bullet \to L_i^\bullet$，它零伦，且 $L_i^\bullet$ 拟同构于零。
例如，可取 $L_i^\bullet$ 为 $K_i^\bullet$ 的恒等态射的锥。用如下推出图
定义 $\mathbf{M}^\bullet(M^\bullet)$：
$$
\xymatrix{
\bigoplus_{i \in I}
\bigoplus_{w : K_i^\bullet \to M^\bullet}
K_i^\bullet \ar[r] \ar[d] & M^\bullet \ar[d] \\
\bigoplus_{i \in I}
\bigoplus_{w : K_i^\bullet \to M^\bullet}
L_i^\bullet \ar[r] &  \mathbf{M}^\bullet(M^\bullet).
}
$$
于是 $M^\bullet \mapsto \mathbf{M}^\bullet(M^\bullet)$ 是函子。右侧竖直箭头
定义函子性单射 $j_{M^\bullet}$。$j_{M^\bullet}$ 的余核同构于映射
$K_i^\bullet \to L_i^\bullet$ 的余核的直和，因而无环。因此，
$j_{M^\bullet}$ 是拟同构。(2) 由构造成立。
\end{proof}

\begin{lemma}
\label{injectives-lemma-functorial-injective}
设 $\mathcal{A}$ 为 Grothendieck 阿贝尔范畴。存在函子
$M^\bullet \mapsto \mathbf{N}^\bullet(M^\bullet)$ 和自然变换
$j_{M^\bullet} : M^\bullet \to \mathbf{N}^\bullet(M^\bullet)$，使得
\begin{enumerate}
\item $j_{M^\bullet}$ 是（逐项）单射的拟同构；并且
\item 对每个 $n \in \mathbf{Z}$，映射
$M^n \to \mathbf{N}^n(M^\bullet)$ 通过子对象
$I^n \subset \mathbf{N}^n(M^\bullet)$ 分解，其中 $I^n$ 是
$\mathcal{A}$ 的内射对象。
\end{enumerate}
\end{lemma}

\begin{proof}
选取函子性内射嵌入 $i_M : M \to I(M)$；见定理
\ref{injectives-theorem-injective-embedding-grothendieck}。对每个复形
$M^\bullet$，记 $J^\bullet(M^\bullet)$ 为各项
$J^n(M^\bullet) = I(M^n) \oplus I(M^{n + 1})$、微分为
$$
d_{J^\bullet(M^\bullet)} =
\left(
\begin{matrix}
0 & 1 \\
0 & 0
\end{matrix}
\right)
$$
的复形。存在复形的典范单射
$u_{M^\bullet} : M^\bullet \to J^\bullet(M^\bullet)$：通过映射
$i_{M^n} : M^n \to I(M^n)$ 和
$i_{M^{n + 1}} \circ d : M^n \to M^{n + 1} \to I(M^{n + 1})$，
把 $M^n$ 映到 $I(M^n) \oplus I(M^{n + 1})$。因此，得到关于
$M^\bullet$ 函子性的复形短正合列
$$
0 \to M^\bullet \xrightarrow{u_{M^\bullet}}
J^\bullet(M^\bullet) \xrightarrow{v_{M^\bullet}}
Q^\bullet(M^\bullet) \to 0
$$
置
$$
\mathbf{N}^\bullet(M^\bullet) = C(v_{M^\bullet})^\bullet[-1].
$$
注意
$$
\mathbf{N}^n(M^\bullet) = Q^{n - 1}(M^\bullet) \oplus J^n(M^\bullet)
$$
其微分为
$$
\left(
\begin{matrix}
- d^{n - 1}_{Q^\bullet(M^\bullet)} & - v^n_{M^\bullet} \\
0 & d^n_{J^\bullet(M)}
\end{matrix}
\right)
$$
因此，由 $u$ 诱导复形映射
$j_{M^\bullet} : M^\bullet \to \mathbf{N}^\bullet(M^\bullet)$。按构造，
它是单射并通过内射子对象分解。考察与上述复形短正合列相伴的上同调长正合列，
可证明映射 $j_{M^\bullet}$ 是拟同构。
\end{proof}

\begin{theorem}
\label{injectives-theorem-K-injective-embedding-grothendieck}
\begin{slogan}
Grothendieck 阿贝尔范畴中 K-内射复形的存在性。
\end{slogan}
设 $\mathcal{A}$ 为 Grothendieck 阿贝尔范畴。对每个复形
$M^\bullet$，存在拟同构 $M^\bullet \to I^\bullet$，使得对所有 $n$，
$M^n \to I^n$ 是单射，$I^n$ 是 $\mathcal{A}$ 的内射对象，且
$I^\bullet$ 是 K-内射复形。此外，该构造关于 $M^\bullet$ 是函子性的。
\end{theorem}

\begin{proof}
比较定理 \ref{injectives-theorem-baer-grothendieck} 和定理
\ref{injectives-theorem-injective-embedding-grothendieck} 的证明。选取引理
\ref{injectives-lemma-acyclic-quotient-complexes-bounded-size} 和
\ref{injectives-lemma-characterize-K-injective} 中的基数 $\kappa$。选取一组有上界的
无环复形 $(K_i^\bullet)_{i \in I}$，使每个满足
$|K^n| \leq \kappa$ 的有上界无环复形 $K^\bullet$ 都同构于某个
$K_i^\bullet$。由引理 \ref{injectives-lemma-set-iso-classes-bounded-size}，这是可能的。
记 $\mathbf{M}^\bullet(-)$ 为引理
\ref{injectives-lemma-functorial-homotopies} 中构造的函子，记
$\mathbf{N}^\bullet(-)$ 为引理 \ref{injectives-lemma-functorial-injective} 中构造的
函子。这两个函子都带有单射变换
$\text{id} \to \mathbf{M}$ 和 $\text{id} \to \mathbf{N}$。

\medskip\noindent
通过超限递归，定义一列函子 $\mathbf{T}_\alpha(-)$ 及相应变换
$\text{id} \to \mathbf{T}_\alpha$。具体地，置
$\mathbf{T}_0(M^\bullet) = M^\bullet$。给定 $\mathbf{T}_\alpha$ 后，置
$$
\mathbf{T}_{\alpha + 1}(M^\bullet) =
\mathbf{N}^\bullet(\mathbf{M}^\bullet(\mathbf{T}_\alpha(M^\bullet)))
$$
如果 $\beta$ 是极限序数，则置
$$
\mathbf{T}_\beta(M^\bullet) =
\colim_{\alpha < \beta} \mathbf{T}_\alpha(M^\bullet)
$$
这个系统的转移映射是单射拟同构。由 AB5，余极限仍拟同构于
$M^\bullet$。我们断言，如果 $\alpha$ 的共尾数大于
$\max(\kappa, |U|)$，其中 $U$ 是 $\mathcal{A}$ 的生成元，则
$M^\bullet \to \mathbf{T}_\alpha(M^\bullet)$ 满足要求。事实上，只需检验
引理 \ref{injectives-lemma-characterize-K-injective} 的条件 (1) 和 (2)。

\medskip\noindent
对 (1)，使用引理 \ref{injectives-lemma-characterize-injective} 的判据。设
$M \subset U$，且对某个 $n \in \mathbf{Z}$，
$\varphi : M \to \mathbf{T}^n_\alpha(M^\bullet)$ 为态射。由命题
\ref{injectives-proposition-objects-are-small}，$\varphi$ 通过某个
$\mathbf{T}^n_{\alpha'}(M^\bullet)$ 分解，其中
$\alpha' < \alpha$。特别地，由函子 $\mathbf{N}^\bullet(-)$ 的构造，
$\varphi$ 通过 $\mathcal{A}$ 的内射对象分解。这表明 $\varphi$ 可提升为
$U$ 上的态射。

\medskip\noindent
对 (2)，设 $w : K^\bullet  \to \mathbf{T}_\alpha(M^\bullet)$ 为复形态射，
其中 $K^\bullet$ 是满足 $|K^n| \leq \kappa$ 的有上界无环复形。则对某个
$i \in I$，有 $K^\bullet \cong K_i^\bullet$。此外，再次由命题
\ref{injectives-proposition-objects-are-small}，$w$ 通过某个
$\mathbf{T}^n_{\alpha'}(M^\bullet)$ 分解，其中
$\alpha' < \alpha$。特别地，由函子 $\mathbf{M}^\bullet(-)$ 的构造，
$w$ 零伦。证明完成。
\end{proof}







\section{Grothendieck 阿贝尔范畴的补充说明}
\label{injectives-section-additional-Grothendieck}

\noindent
本节给出一些关于 Grothendieck 阿贝尔范畴的流传已久的结果。

\begin{lemma}
\label{injectives-lemma-grothendieck-brown}
设 $\mathcal{A}$ 为 Grothendieck 阿贝尔范畴，并设
$F : \mathcal{A}^{opp} \to \textit{Sets}$ 为函子。则 $F$ 可表当且仅当
$F$ 与余极限交换，即
$$
F(\colim_i N_i) = \lim F(N_i)
$$
对任意图式 $\mathcal{I} \to \mathcal{A}$，$i \in \mathcal{I}$，均有上式。
\end{lemma}

\begin{proof}
若 $F$ 可表，则由余极限的定义，它与余极限交换。

\medskip\noindent
现设 $F$ 与余极限交换。于是
$F(M \oplus N) = F(M) \times F(N)$，可据此在 $F(M)$ 上定义群结构。
从而得到加性且右正合的函子 $F : \mathcal{A} \to \textit{Ab}$；也就是说，
它把短正合列 $0 \to K \to L \to M \to 0$ 变为正合列
$F(K) \leftarrow F(L) \leftarrow F(M) \leftarrow 0$（比较《同调代数》
第 \ref{homology-section-functors} 节）。

\medskip\noindent
设 $U$ 为 $\mathcal{A}$ 的生成元。置
$A = \bigoplus_{s \in F(U)} U$，并令
$s_{univ} = (s)_{s \in F(U)} \in F(A) = \prod_{s \in F(U)} F(U)$。
设 $A' \subset A$ 为使 $s_{univ}$ 在其上限制为零的最大子对象。
由于 $\mathcal{A}$ 是 Grothendieck 范畴且 $F$ 与余极限交换，这一对象存在。
又因 $F$ 与余极限交换，存在唯一元素
$\overline{s}_{univ} \in F(A/A')$，其在 $F(A)$ 中的像为 $s_{univ}$。
我们断言 $A/A'$ 表示 $F$；换言之，Yoneda 映射
$$
\overline{s}_{univ} : h_{A/A'} \longrightarrow F
$$
是同构。设 $M \in \Ob(\mathcal{A})$ 且 $s \in F(M)$。考虑满射
$$
c_M :
A_M = \bigoplus\nolimits_{\varphi \in \Hom_\mathcal{A}(U, M)} U
\longrightarrow
M.
$$
它给出 $F(c_M)(s) = (s_\varphi) \in \prod_\varphi F(U)$。考虑映射
$$
\psi :
A_M = \bigoplus\nolimits_{\varphi \in \Hom_\mathcal{A}(U, M)} U
\longrightarrow
\bigoplus\nolimits_{s \in F(U)} U = A
$$
它借助 $U$ 上的恒等映射，将对应于 $\varphi$ 的直和项映至对应于
$s_\varphi$ 的直和项。由构造，$s_{univ}$ 映至
$(s_\varphi)_\varphi$。换言之，下图右方块交换：
$$
\xymatrix{
A' \ar[r] &
A \ar@{..>}[r]_{s_{univ}} & F \\
K \ar[r] \ar[u]^{?} & A_M \ar[u]^\psi \ar[r] &
M \ar@{..>}[u]_s
}
$$
令 $K = \Ker(A_M \to M)$。由于 $s$ 在 $K$ 上限制为零，由 $A'$ 的定义有
$\psi(K) \subset A'$。因此诱导出态射 $M \to A/A'$。此构造给出映射
$h_{A/A'}(M) \to F(M)$ 的逆（略去细节）。
\end{proof}

\begin{lemma}
\label{injectives-lemma-grothendieck-products}
Grothendieck 阿贝尔范畴满足 Ab3*。
\end{lemma}

\begin{proof}
设 $M_i$（$i \in I$）为以集合 $I$ 为指标的 $\mathcal{A}$ 中对象族。
函子 $F = \prod_{i \in I} h_{M_i}$ 与余极限交换，故引理
\ref{injectives-lemma-grothendieck-brown} 适用。
\end{proof}

\begin{remark}
\label{injectives-remark-existence-D}
在导出范畴一章中，我们始终处理“小”阿贝尔范畴（Stacks 项目惯例如此）。
对“大”阿贝尔范畴 $\mathcal{A}$，导出范畴 $D(\mathcal{A})$ 是否存在
并不显然，因为导出范畴中的态射是否构成集合并不显然。一般而言这并不成立，
见《例》引理 \ref{examples-lemma-big-abelian-category}。不过，若
$\mathcal{A}$ 是 Grothendieck 阿贝尔范畴，且给定
$K^\bullet, L^\bullet \in K(\mathcal{A})$，则由定理
\ref{injectives-theorem-K-injective-embedding-grothendieck}，存在到 K-内射复形
$I^\bullet$ 的拟同构 $L^\bullet \to I^\bullet$，而《导出范畴》引理
\ref{derived-lemma-K-injective} 表明
$$
\Hom_{D(\mathcal{A})}(K^\bullet, L^\bullet) =
\Hom_{K(\mathcal{A})}(K^\bullet, I^\bullet)
$$
这是一个集合。Grothendieck 阿贝尔范畴的例子包括环上的模范畴；更一般地，
还包括环站点上的模层范畴。
\end{remark}

\begin{lemma}
\label{injectives-lemma-derived-products}
设 $\mathcal{A}$ 为 Grothendieck 阿贝尔范畴。则
\begin{enumerate}
\item $D(\mathcal{A})$ 既有直和也有积；
\item 直和可由任意复形的逐项直和得到；
\item 积可由 K-内射复形的逐项积得到。
\end{enumerate}
\end{lemma}

\begin{proof}
设 $K^\bullet_i$（$i \in I$）为以集合 $I$ 为指标的
$D(\mathcal{A})$ 中对象族。我们断言逐项直和
$\bigoplus_{i \in I} K^\bullet_i$ 是 $D(\mathcal{A})$ 中的直和。
事实上，设 $I^\bullet$ 为 K-内射复形，则
\begin{align*}
\Hom_{D(\mathcal{A})}(\bigoplus\nolimits_{i \in I} K^\bullet_i, I^\bullet)
& =
\Hom_{K(\mathcal{A})}(\bigoplus\nolimits_{i \in I} K^\bullet_i, I^\bullet) \\
& =
\prod\nolimits_{i \in I} \Hom_{K(\mathcal{A})}(K^\bullet_i, I^\bullet) \\
& =
\prod\nolimits_{i \in I} \Hom_{D(\mathcal{A})}(K^\bullet_i, I^\bullet)
\end{align*}
正合所需。由定理 \ref{injectives-theorem-K-injective-embedding-grothendieck}，
任意复形均可由 K-内射复形表示，故这已经足够。为构造积，对每个 $i$
选取 K-内射消解 $K_i^\bullet \to I_i^\bullet$。我们断言
$\prod_{i \in I} I_i^\bullet$ 是 $D(\mathcal{A})$ 中的积。
这由《导出范畴》引理 \ref{derived-lemma-product-K-injective} 得出。
\end{proof}

\begin{remark}
\label{injectives-remark-direct-sum-product-derived}
设 $R$ 为环，并设 $M_n$（$n \in \mathbf{Z}$）为 $R$-模。记
$E_n = M_n[-n] \in D(R)$。我们断言 $E = \bigoplus M_n[-n]$
{\it 同时}是 $D(R)$ 中对象 $E_n$ 的直和与积。要看它是直和，只须查看引理
\ref{injectives-lemma-derived-products} 的证明。要看它是积，取内射消解
$M_n \to I_n^\bullet$。由引理 \ref{injectives-lemma-derived-products} 的证明，有
$$
\prod E_n = \prod I_n^\bullet[-n]
$$
在 $D(R)$ 中成立。由于 $\text{Mod}_R$ 中的积正合，
$\prod I_n^\bullet[-n]$ 与 $E$ 拟同构。更一般地，对满足 Ab4* 的
Grothendieck 阿贝尔范畴 $\mathcal{A}$，此论证在 $D(\mathcal{A})$ 中同样成立。
\end{remark}

\begin{lemma}
\label{injectives-lemma-RF-commutes-with-Rlim}
设 $F : \mathcal{A} \to \mathcal{B}$ 为阿贝尔范畴间的加性函子。假设
\begin{enumerate}
\item $\mathcal{A}$ 是 Grothendieck 阿贝尔范畴；
\item $\mathcal{B}$ 的可数积正合；
\item $F$ 与可数积交换。
\end{enumerate}
则 $RF : D(\mathcal{A}) \to D(\mathcal{B})$ 与导出极限交换。
\end{lemma}

\begin{proof}
注意，因 $\mathcal{A}$ 有足够多的 K-内射对象，$RF$ 存在（定理
\ref{injectives-theorem-K-injective-embedding-grothendieck} 与《导出范畴》引理
\ref{derived-lemma-K-injective-defined}）。断言的含义是：若
$K = R\lim K_n$，则 $RF(K) = R\lim RF(K_n)$；记号见《导出范畴》定义
\ref{derived-definition-derived-limit}。由于 $RF$ 是三角范畴间的正合函子，
只须证明 $RF$ 与 $D(\mathcal{A})$ 中对象的可数积交换。引理
\ref{injectives-lemma-derived-products} 的证明表明，$D(\mathcal{A})$ 中的积通过
K-内射复形的积计算，而且 K-内射复形的积仍为 K-内射复形。此外，
《导出范畴》引理 \ref{derived-lemma-products} 表明，$D(\mathcal{B})$
中的积通过逐项积计算。$RF$ 由在 K-内射代表上施用 $F$ 计算，而我们假设
$F$ 与可数积交换，故引理成立。
\end{proof}

\noindent
以下引理可看作 Cartan--Eilenberg 消解存在性（《导出范畴》第
\ref{derived-section-cartan-eilenberg} 节）的某种推广。

\begin{lemma}
\label{injectives-lemma-K-injective-embedding-filtration}
设 $\mathcal{A}$ 为 Grothendieck 阿贝尔范畴，$K^\bullet$ 为
$\mathcal{A}$ 的滤过复形，见《同调代数》定义
\ref{homology-definition-filtered-complex}。则存在 $\mathcal{A}$
的滤过复形态射 $j : K^\bullet \to J^\bullet$，使得
\begin{enumerate}
\item $J^n$、$F^pJ^n$、$J^n/F^pJ^n$ 与 $F^pJ^n/F^{p'}J^n$
均为 $\mathcal{A}$ 的内射对象；
\item $J^\bullet$、$F^pJ^\bullet$、$J^\bullet/F^pJ^\bullet$ 与
$F^pJ^\bullet/F^{p'}J^\bullet$ 均为 K-内射复形；
\item $j$ 诱导拟同构
$K^\bullet \to J^\bullet$,
$F^pK^\bullet \to F^pJ^\bullet$,
$K^\bullet/F^pK^\bullet \to J^\bullet/F^pJ^\bullet$，以及
$F^pK^\bullet/F^{p'}K^\bullet \to F^pJ^\bullet/F^{p'}J^\bullet$.
\end{enumerate}
\end{lemma}

\begin{proof}
由定理 \ref{injectives-theorem-K-injective-embedding-grothendieck}，得到拟同构
$i : K^\bullet \to I^\bullet$ 与
$i^p : F^pK^\bullet \to I^{p, \bullet}$，以及交换图
$$
\vcenter{
\xymatrix{
K^\bullet \ar[d]_i & F^pK^\bullet \ar[l] \ar[d]_{i^p} \\
I^\bullet & I^{p, \bullet} \ar[l]_{\alpha^p}
}
}
\quad\text{且}\quad
\vcenter{
\xymatrix{
F^{p'}K^\bullet \ar[d]_{i^{p'}} &
F^pK^\bullet \ar[l] \ar[d]_{i^p} \\
I^{p', \bullet} &
I^{p, \bullet} \ar[l]_{\alpha^{p p'}}
}
}
\quad\text{其中 }p' \leq p
$$
并满足 $\alpha^p \circ \alpha^{p' p} = \alpha^{p'}$ 及
$\alpha^{p'p''} \circ \alpha^{pp'} = \alpha^{pp''}$。问题在于映射
$\alpha^p : I^{p, \bullet} \to I^\bullet$ 未必是单射。对每个 $p$，
选取到无环 K-内射复形 $J^{p, \bullet}$ 的单射
$t^p : I^{p, \bullet} \to J^{p, \bullet}$，其中
$J^{p, \bullet}$ 各项均为 $\mathcal{A}$ 的内射对象（先映至恒等映射的锥，
再使用该定理）。选取复形映射
$s^p : I^\bullet \to J^{p, \bullet}$，使下图交换：
$$
\xymatrix{
K^\bullet \ar[d]_i & F^pK^\bullet \ar[l] \ar[d]_{i^p} \\
I^\bullet \ar[rd]_{s^p} & I^{p, \bullet} \ar[d]^{t^p} \\
& J^{p, \bullet}
}
$$
这是可能的：因 $J^{p, \bullet}$ 无环且 K-内射，复合
$F^pK^\bullet \to J^{p, \bullet}$ 零伦（《导出范畴》引理
\ref{derived-lemma-K-injective}）。对象 $J^{p, n - 1}$ 内射，且
$F^pK^n \to K^n \to I^n$ 为单态射，故可将给出该同伦的映射
$F^pK^n \to J^{p, n - 1}$ 延拓为映射
$h^n : I^n \to J^{p, n - 1}$。然后置
$s^p = h \circ \text{d} + \text{d} \circ h$。
（警告：并不一定有 $t^p = s^p \circ \alpha^p$，故下文须注意不要使用此式。）

\medskip\noindent
考虑
$$
J^\bullet = I^\bullet \times \prod\nolimits_p J^{p, \bullet}
$$
由于 $D(\mathcal{A})$ 中的积由 K-内射复形的积给出（引理
\ref{injectives-lemma-derived-products}），且 $J^{p, \bullet}$ 在
$D(\mathcal{A})$ 中同构于 $0$，可见
$J^\bullet \to I^\bullet$ 是 $D(\mathcal{A})$ 中的同构。考虑映射
$$
j = i \times (s^p \circ i)_{p \in \mathbf{Z}} :
K^\bullet
\longrightarrow
I^\bullet \times \prod\nolimits_p J^{p, \bullet} = J^\bullet
$$
由上述说明，这是拟同构；它也是单射。对 $p \in \mathbf{Z}$，令
$F^pJ^\bullet \subset J^\bullet$ 为
$$
\Im\left(
\alpha^p \times (t^{p'} \circ \alpha^{pp'})_{p' \leq p} :
I^{p, \bullet}
\to
I^\bullet \times \prod\nolimits_{p' \leq p} J^{p', \bullet}
\right)
\times \prod\nolimits_{p' > p} J^{p', \bullet}
$$
此复形同构于复形
$I^{p, \bullet} \times \prod_{p' > p} J^{p, \bullet}$
，因为 $\alpha^{pp} = \text{id}$ 且 $t^p$ 为单射。因此
$F^pJ^\bullet$ 与 $I^{p, \bullet}$ 拟同构（论证同上）。由上图的交换性，
有 $j(F^pK^\bullet) \subset F^pJ^\bullet$。由刚才所述，对应的复形映射
$F^pK^\bullet \to F^pJ^\bullet$ 是拟同构。最后，为证明
$F^{p + 1}J^\bullet \subset F^pJ^\bullet$
，使用 $\alpha^{p + 1p} \circ \alpha^{pp'} = \alpha^{p + 1p'}$
以及本证明第一段中第一个陈列图的交换性。

\medskip\noindent
我们断言 $j : K^\bullet \to J^\bullet$ 正是引理所求。事实上，
$F^pJ^n$ 同构于 $I^{p, n} \times \prod_{p' > p} J^{p', n}$，而内射对象的积
仍内射，故 $F^pJ^n$ 是 $\mathcal{A}$ 的内射对象。于是单态射
$F^pJ^n \to J^n$ 分裂，所以商 $J^n/F^pJ^n$ 作为内射对象 $J^n$
的直和因子也是内射的。对 $F^pJ^n/F^{p'}J^n$ 同理。特别地，这意味着
$0 \to F^pJ^\bullet \to J^\bullet \to J^\bullet/F^pJ^\bullet \to 0$
是逐项分裂的复形短正合列，故按约定在 $K(\mathcal{A})$ 中定义一个特异三角。
由于 $J^\bullet$ 与 $F^pJ^\bullet$ 均为 K-内射复形，由《导出范畴》引理
\ref{derived-lemma-triangle-K-injective}，$J^\bullet/F^pJ^\bullet$
也如此。类似论证表明 $F^pJ^\bullet/F^{p'}J^\bullet$ 为 K-内射复形。
由构造，$j : K^\bullet \to J^\bullet$ 以及诱导映射
$F^pK^\bullet \to F^pJ^\bullet$ 均为拟同构。利用所涉及复形的上同调长正合列，
可知 $K^\bullet/F^pK^\bullet \to J^\bullet/F^pJ^\bullet$ 以及
$F^pK^\bullet/F^{p'}K^\bullet \to F^pJ^\bullet/F^{p'}J^\bullet$
也具有相同性质。
\end{proof}

\begin{remark}
\label{injectives-remark-ext-into-filtered-complex}
设 $\mathcal{A}$ 为 Grothendieck 阿贝尔范畴，$K^\bullet$ 为
$\mathcal{A}$ 的滤过复形，见《同调代数》定义
\ref{homology-definition-filtered-complex}。为简化记号，分别以
$K$、$F^pK$、$\text{gr}^pK$ 表示由 $K^\bullet$、$F^pK^\bullet$、
$\text{gr}^pK^\bullet$ 表示的 $D(\mathcal{A})$ 中对象。设
$M \in D(\mathcal{A})$。利用引理
\ref{injectives-lemma-K-injective-embedding-filtration}，可构造由 $\mathcal{A}$
中双分次对象组成的谱序列 $(E_r, d_r)_{r \geq 1}$，其中 $d_r$
的双次数为 $(r, -r + 1)$，并且
$$
E_1^{p, q} = \Ext^{p + q}(M, \text{gr}^pK)
$$
若对每个 $n$ 均有
$$
\Ext^n(M, F^pK) = 0 \text{ 当 } p \gg 0
\quad\text{且}\quad
\Ext^n(M, F^pK) = \Ext^n(M, K) \text{ 当 } p \ll 0
$$
，则此谱序列有界并收敛于 $\Ext^{p + q}(M, K)$。事实上，任取表示 $M$
的复形 $M^\bullet$，选取引理中的 $j : K^\bullet \to J^\bullet$，并考虑复形
$$
\Hom^\bullet(M^\bullet, J^\bullet)
$$
，其定义与《代数详论》第 \ref{more-algebra-section-hom-complexes} 节完全相同。
置 $F^p\Hom^\bullet(M^\bullet, J^\bullet) =
\Hom^\bullet(M^\bullet, F^pJ^\bullet)$，便得到一个滤过复形。
《同调代数》第 \ref{homology-section-filtered-complex} 节的谱序列具有上述
微分与各项；细节从略。有界性与收敛性由《同调代数》引理
\ref{homology-lemma-ss-converges-trivial} 得出。
\end{remark}

\begin{remark}
\label{injectives-remark-spectral-sequences-ext}
设 $\mathcal{A}$ 为 Grothendieck 阿贝尔范畴，$M,K$ 为
$D(\mathcal{A})$ 的对象。任取表示 $K$ 的复形 $K^\bullet$，利用滤过
$F^pK^\bullet = \tau_{\leq -p}K^\bullet$ 以及评注
\ref{injectives-remark-ext-into-filtered-complex} 中的讨论，可得到满足
$$
E_1^{p, q} = \Ext^{2p + q}(M, H^{-p}(K))
$$
的谱序列。它不依赖于表示 $K$ 的复形 $K^\bullet$ 的选择。重新编号
$p = -j$、$q = i + 2j$ 后，得到谱序列
$(E'_r, d'_r)_{r \geq 2}$，其中 $d'_r$ 的双次数为 $(r, -r + 1)$，并且
$$
(E'_2)^{i, j} = \Ext^i(M, H^j(K))
$$
若 $M \in D^-(\mathcal{A})$ 且 $K \in D^+(\mathcal{A})$，则
$E_r$ 与 $E'_r$ 均有界并收敛于 $\Ext^{p + q}(M, K)$。若使用滤过
$F^pK^\bullet = \sigma_{\geq p}K^\bullet$，则得到
$$
E_1^{p, q} = \Ext^q(M, K^p)
$$
若 $M \in D^-(\mathcal{A})$ 且 $K^\bullet$ 下有界，则此谱序列有界并收敛于
$\Ext^{p + q}(M, K)$。
\end{remark}

\begin{remark}
\label{injectives-remark-ext-from-filtered-complex}
设 $\mathcal{A}$ 为 Grothendieck 阿贝尔范畴，$K \in D(\mathcal{A})$。
设 $M^\bullet$ 为 $\mathcal{A}$ 的滤过复形，见《同调代数》定义
\ref{homology-definition-filtered-complex}。为简化记号，分别以
$M$、$M/F^pM$、$\text{gr}^pM$ 表示由 $M^\bullet$、
$M^\bullet/F^pM^\bullet$、$\text{gr}^pM^\bullet$ 表示的
$D(\mathcal{A})$ 中对象。与评注 \ref{injectives-remark-ext-into-filtered-complex}
对偶地，可构造由 $\mathcal{A}$ 中双分次对象组成的谱序列
$(E_r, d_r)_{r \geq 1}$，其中 $d_r$ 的双次数为 $(r, -r + 1)$，并且
$$
E_1^{p, q} = \Ext^{p + q}(\text{gr}^{-p}M, K)
$$
若对每个 $n$ 均有
$$
\Ext^n(M/F^pM, K) = 0 \text{ 当 } p \ll 0
\quad\text{且}\quad
\Ext^n(M/F^pM, K) = \Ext^n(M, K) \text{ 当 } p \gg 0
$$
，则此谱序列有界并收敛于 $\Ext^{p + q}(M, K)$。事实上，选取表示 $K$
且各项内射的 K-内射复形 $I^\bullet$，见定理
\ref{injectives-theorem-K-injective-embedding-grothendieck}。考虑复形
$$
\Hom^\bullet(M^\bullet, I^\bullet)
$$
，其定义与《代数详论》第 \ref{more-algebra-section-hom-complexes} 节完全相同。
置
$$
F^p\Hom^\bullet(M^\bullet, I^\bullet) =
\Hom^\bullet(M^\bullet/F^{-p + 1}M^\bullet, I^\bullet)
$$
，便得到一个滤过复形（注意滤过中的符号与移位）。《同调代数》第
\ref{homology-section-filtered-complex} 节的谱序列具有上述微分与各项；
细节从略。有界性与收敛性由《同调代数》引理
\ref{homology-lemma-ss-converges-trivial} 得出。
\end{remark}

\begin{remark}
\label{injectives-remark-spectral-sequences-ext-variant}
设 $\mathcal{A}$ 为 Grothendieck 阿贝尔范畴，$M,K$ 为
$D(\mathcal{A})$ 的对象。任取表示 $M$ 的复形 $M^\bullet$，利用滤过
$F^pM^\bullet = \tau_{\leq -p}M^\bullet$ 以及评注
\ref{injectives-remark-ext-into-filtered-complex} 中的讨论，可得到满足
$$
E_1^{p, q} = \Ext^{2p + q}(H^p(M), K)
$$
的谱序列。它不依赖于表示 $M$ 的复形 $M^\bullet$ 的选择。重新编号
$p = -j$、$q = i + 2j$ 后，得到谱序列
$(E'_r, d'_r)_{r \geq 2}$，其中 $d'_r$ 的双次数为
$(r, -r + 1)$，并且
$$
(E'_2)^{i, j} = \Ext^i(H^{-j}(M), K)
$$
若 $M \in D^-(\mathcal{A})$ 且 $K \in D^+(\mathcal{A})$，则
$E_r$ 与 $E'_r$ 均有界并收敛于 $\Ext^{p + q}(M, K)$。若使用滤过
$F^pM^\bullet = \sigma_{\geq p}M^\bullet$，则得到
$$
E_1^{p, q} = \Ext^q(M^{-p}, K)
$$
若 $K \in D^+(\mathcal{A})$ 且 $M^\bullet$ 上有界，则此谱序列有界并收敛于
$\Ext^{p + q}(M, K)$。
\end{remark}

\begin{lemma}
\label{injectives-lemma-represent-by-filtered-complex}
设 $\mathcal{A}$ 为 Grothendieck 阿贝尔范畴。给定对象
$E \in D(\mathcal{A})$、以 $\mathbf{Z}$ 为指标的 $D(\mathcal{A})$
中对象的逆系统 $\{E^i\}_{i \in \mathbf{Z}}$，以及相容的映射系
$E^i \to E$。图示为
$$
\ldots \to E^{i + 1} \to E^i \to E^{i - 1} \to \ldots \to E
$$
则存在 $\mathcal{A}$ 的滤过复形 $K^\bullet$（《同调代数》定义
\ref{homology-definition-filtered-complex}），使 $K^\bullet$ 表示 $E$，
而 $F^iK^\bullet$ 以与给定映射相容的方式表示 $E^i$。
\end{lemma}

\begin{proof}
由定理 \ref{injectives-theorem-K-injective-embedding-grothendieck}，可选取表示 $E$
的 K-内射复形 $I^\bullet$，其每一项 $I^n$ 均为 $\mathcal{A}$ 的内射对象。
选取表示 $E^0$ 的复形 $G^{0, \bullet}$，再选取表示 $E^0 \to E$
的复形映射 $\varphi^0 : G^{0, \bullet} \to I^\bullet$。对 $i > 0$，
归纳地用复形映射 $\delta : G^{i, \bullet} \to G^{i - 1, \bullet}$
表示 $E^i \to E^{i - 1}$，并置
$\varphi^i = \delta \circ \varphi^{i - 1}$。对 $i < 0$，归纳地用逐项单射的
复形映射
$\delta : G^{i + 1, \bullet} \to G^{i, \bullet}$
表示 $E^{i + 1} \to E^i$（例如可用《导出范畴》引理
\ref{derived-lemma-make-injective}）。我们断言可找到复形映射
$\varphi^i : G^{i, \bullet} \to I^\bullet$
，它表示映射 $E^i \to E$，并使下图交换：
$$
\xymatrix{
G^{i + 1, \bullet} \ar[r]_\delta \ar[d]_{\varphi^{i + 1}} &
G^{i, \bullet} \ar[ld]^{\varphi^i} \\
I^\bullet
}
$$
中。事实上，先任取表示映射 $E^i \to E$ 的复形映射
$\varphi : G^{i, \bullet} \to I^\bullet$。于是
$\varphi \circ \delta$ 与 $\varphi^{i + 1}$ 通过某个同伦
$h^p : G^{i + 1, p} \to I^{p - 1}$ 同伦。由于 $I^\bullet$ 的各项内射，
且 $\delta$ 逐项单射，可将 $h^p$ 延拓为
$(h')^p : G^{i, p} \to I^{p - 1}$。置
$\varphi^i = \varphi + h' \circ d + d \circ h'$，便得到所断言的映射。

\medskip\noindent
接着，对每个 $i$ 选取逐项单射的复形映射
$a^i : G^{i, \bullet} \to J^{i, \bullet}$，其中 $J^{i, \bullet}$
无环且 K-内射，而 $J^{i, p}$ 为 $\mathcal{A}$ 的内射对象。为此，先将
$G^{i, \bullet}$ 映至恒等映射的锥，再应用上述定理。与上面相同地论证，
可找到复形映射 $\delta' : J^{i, \bullet} \to J^{i - 1, \bullet}$，使图
$$
\xymatrix{
G^{i, \bullet} \ar[r]_\delta \ar[d]_{a^i} &
G^{i - 1, \bullet} \ar[d]^{a^{i - 1}} \\
J^{i, \bullet} \ar[r]^{\delta'} &
J^{i - 1, \bullet}
}
$$
交换。（也可利用锥的函子性以及该定理中的函子性得到这一点。）然后考虑映射
$$
\xymatrix{
G^{i + 1, \bullet} \times \prod\nolimits_{p > i + 1} J^{p, \bullet}
\ar[r] \ar[rd] &
G^{i, \bullet} \times \prod\nolimits_{p > i} J^{p, \bullet}
\ar[r] \ar[d] &
G^{i - 1, \bullet} \times \prod\nolimits_{p > i - 1} J^{p, \bullet}
\ar[ld] \\
& I^\bullet \times \prod\nolimits_p J^{p, \bullet}
}
$$
此处 $J^{p, \bullet}$ 上的箭头是显然的箭头（恒等映射或零映射）。在因子
$G^{i, \bullet}$ 上，使用
$\delta : G^{i, \bullet} \to G^{i - 1, \bullet}$、映射
$\varphi^i : G^{i, \bullet} \to I^\bullet$、当 $p > i$ 时的零映射
$0 : G^{i, \bullet} \to J^{p, \bullet}$、当 $p = i$ 时的映射
$a^i : G^{i, \bullet} \to J^{p, \bullet}$，以及当 $p < i$ 时的映射
$(\delta')^{i - p} \circ a^i = a^p \circ \delta^{i - p} :
G^{i, \bullet} \to J^{p, \bullet}$。略去图中所有箭头均逐项单射的验证。
这样便得到一个滤过复形。由于 $D(\mathcal{A})$ 中的积由 K-内射复形的积给出
（引理 \ref{injectives-lemma-derived-products}），且 $J^{p, \bullet}$ 在
$D(\mathcal{A})$ 中为零，可知此图在导出范畴中表示给定的图。证明完成。
\end{proof}

\begin{lemma}
\label{injectives-lemma-represent-by-filtered-complex-bis}
在引理 \ref{injectives-lemma-represent-by-filtered-complex} 的情形中，进一步假设有第二个
逆系统 $\{(E')^i\}_{i \in \mathbf{Z}}$ 以及相容的映射系
$(E')^i \to E$。则存在 $\mathcal{A}$ 的双滤过复形 $K^\bullet$，使
$K^\bullet$ 表示 $E$，$F^iK^\bullet$ 表示 $E^i$，而
$(F')^iK^\bullet$ 以与给定映射相容的方式表示 $(E')^i$。
\end{lemma}

\begin{proof}
先利用该引理选取 $K^\bullet$ 与 $F$，再为
$\{(E')^i\}_{i \in \mathbf{Z}}$ 及映射 $(E')^i \to E$ 选取适用的
$(K')^\bullet$ 与 $F'$。利用引理
\ref{injectives-lemma-K-injective-embedding-filtration}，可假设 $K^\bullet$
为 K-内射复形。于是可选取对应于给定认同
$(K')^\bullet \cong E \cong K^\bullet$ 的复形映射
$(K')^\bullet \to K^\bullet$。还可选取逐项单射的映射
$(K')^\bullet \to J^\bullet$，其中 $J^\bullet$ 无环且 K-内射。
（为此先将 $(K')^\bullet$ 映至恒等映射的锥，再应用定理
\ref{injectives-theorem-K-injective-embedding-grothendieck}。）于是
$(K')^\bullet \to K^\bullet \times J^\bullet$ 与
$K^\bullet \to K^\bullet \times J^\bullet$ 均逐项单射且为拟同构
（由引理 \ref{injectives-lemma-derived-products}，该积表示 $E$）。只须取
$K^\bullet$ 与 $(K')^\bullet$ 上的滤过在这些映射下的像，即得结论。
\end{proof}





\section{Gabriel--Popescu 定理}
\label{injectives-section-gabriel-popescu}

\noindent
本节讨论 \cite{GP} 的主定理。证明方法遵循 Jacob Lurie 的一份讲义以及
Akhil Mathew 的另一份讲义；两者又都遵循 \cite{Kuhn} 中 Kuhn 的陈述。
另见 \cite{Takeuchi}。

\medskip\noindent
设 $\mathcal{A}$ 为 Grothendieck 阿贝尔范畴，$U$ 为 $\mathcal{A}$ 的生成元，
见定义 \ref{injectives-definition-grothendieck-conditions}。令
$R = \Hom_\mathcal{A}(U, U)$。考虑由
$$
G(A) = \Hom_\mathcal{A}(U, A)
$$
给出的函子 $G : \mathcal{A} \to \text{Mod}_R$，其中 $G(A)$ 带有典范的右
$R$-模结构。

\begin{lemma}
\label{injectives-lemma-gabriel-popescu-left-adjoint}
上述函子 $G$ 有左伴随函子 $F : \text{Mod}_R \to \mathcal{A}$。
\end{lemma}

\begin{proof}
我们给出此引理的两个证明。

\medskip\noindent
第一个证明使用伴随函子定理，见《范畴》定理
\ref{categories-theorem-adjoint-functor}。注意
$G : \mathcal{A} \to \text{Mod}_R$ 左正合，并把积映为积，故 $G$ 与极限交换。
为验证该定理中的集合论条件，设 $M$ 为 $\text{Mod}_R$ 的对象。选取充分大的
基数 $\kappa$，并以 $E$ 表示 $\mathcal{A}$ 中对象的一个集合，使每个满足
$|A| \leq \kappa$ 的对象 $A$ 均同构于 $E$ 的某个元素。由引理
\ref{injectives-lemma-set-iso-classes-bounded-size}，这是可能的。置
$I = \coprod_{A \in E} \Hom_R(M, G(A))$。把元素 $i \in I$ 看作一对
$(A_i, f_i)$。最后，任取 $\mathcal{A}$ 的对象 $A$ 与映射
$f : M \to G(A)$。我们把
$\Im(f) \subset G(A) = \Hom_\mathcal{A}(U, A)$ 中的元素看作映射
$u : U \to A$。置
$$
A' = \Im(\bigoplus\nolimits_{u \in \Im(f)} U \xrightarrow{u} A)
$$
由于 $G$ 左正合，$G(A') \subset G(A)$ 包含 $\Im(f)$，并得到使 $f$ 分解的
$f' : M \to G(A')$。另一方面，对象 $A'$ 是至多 $|M|$ 份 $U$ 的直和的商。
因此，若 $\kappa = |\bigoplus_{|M|} U|$，则 $(A', f')$ 同构于 $E$ 的某个元素
$(A_i, f_i)$，从而 $f$ 按所需方式分解为
$M \xrightarrow{f_i} G(A_i) \to G(A)$。

\medskip\noindent
第二个证明将构造 $F$，并表明在某种意义下“$F(M) = M \otimes_R U$”。
具体地，对任意 $R$-模 $M$，可选取消解
$$
\bigoplus\nolimits_{j \in J} R \to
\bigoplus\nolimits_{i \in I} R \to
M \to 0
$$
然后用相应的正合列定义 $F(M)$：
$$
\bigoplus\nolimits_{j \in J} U \to
\bigoplus\nolimits_{i \in I} U \to
F(M) \to 0
$$
此构造不依赖于消解的选择，并且是函子性的；细节从略。对
$\mathcal{A}$ 中任意 $A$，得到正合列
$$
0 \to \Hom_\mathcal{A}(F(M), A) \to
\prod\nolimits_{i \in I} G(A) \to
\prod\nolimits_{j \in J} G(A)
$$
，它同构于正合列
$$
0 \to \Hom_R(M, G(A)) \to
\Hom_R(\bigoplus\nolimits_{i \in I} R, G(A)) \to
\Hom_R(\bigoplus\nolimits_{j \in J} R, G(A))
$$
。这表明 $F$ 是 $G$ 的左伴随函子。
\end{proof}

\begin{lemma}
\label{injectives-lemma-F-G-monos}
设 $f : M \to G(A)$ 为 $\text{Mod}_R$ 中的单射，则伴随映射
$f' : F(M) \to A$ 也是单射。
\end{lemma}

\begin{proof}
选取映射 $R^{\oplus n} \to M$，并考虑相应的映射
$U^{\oplus n} \to F(M)$。考虑映射 $v : U \to U^{\oplus n}$，使复合
$U \to U^{\oplus n} \to F(M) \to A$ 为 $0$。于是箭头
$v : U \to U^{\oplus n}$ 是 $R^{\oplus n}$ 中的元素 $v$，且它在
$G(A)$ 中映为零。由于 $f$ 为单射，$v$ 在 $M$ 中映为零；由引理
\ref{injectives-lemma-gabriel-popescu-left-adjoint} 的证明中 $F(M)$ 的构造，这意味着
$U \to U^{\oplus n} \to F(M)$ 为零。由于 $U$ 是生成元，可得
$$
\Ker(U^{\oplus n} \to F(M) \to A) = \Ker(U^{\oplus n} \to F(M))
$$
为完成证明，选取满射 $\bigoplus_{i \in I} R \to M$，并考虑相应的满射
$$
\pi : \bigoplus\nolimits_{i \in I} U \longrightarrow F(M)
$$
要证明 $f'$ 为单射，只须证明作为 $\bigoplus_{i \in I} U$ 的子对象有
$\Ker(\pi) = \Ker(f' \circ \pi)$。现在可将 $\bigoplus_{i \in I} U$
写成其子对象 $\bigoplus_{i \in I'} U$ 的滤过余极限，其中
$I' \subset I$ 遍历有限子集。由于 $\mathcal{A}$ 满足 AB5，滤过余极限正合，
故
$$
\Ker(\pi) =
\colim_{I' \subset I\text{ 有限}}
\left(\bigoplus\nolimits_{i \in I'} U\right)
\bigcap \Ker(\pi)
$$
并且
$$
\Ker(f' \circ \pi) =
\colim_{I' \subset I\text{ 有限}}
\left(\bigoplus\nolimits_{i \in I'} U\right)
\bigcap \Ker(f' \circ \pi)
$$
。由上面第一个陈列等式，对每个 $I'$ 两者均相等，故得到所需等式。
\end{proof}

\begin{theorem}
\label{injectives-theorem-gabriel-popescu}
设 $\mathcal{A}$ 为 Grothendieck 阿贝尔范畴。则存在一个（非交换）环 $R$
以及函子 $G : \mathcal{A} \to \text{Mod}_R$ 与
$F : \text{Mod}_R \to \mathcal{A}$，使得
\begin{enumerate}
\item $F$ 是 $G$ 的左伴随函子；
\item $G$ 为全忠实函子；
\item $F$ 正合。
\end{enumerate}
而且，这些函子正是上面构造的函子。
\end{theorem}

\begin{proof}
先证明 $G$ 全忠实，等价地，证明 $F \circ G \to \text{id}$ 是同构；见《范畴》
引理 \ref{categories-lemma-adjoint-fully-faithful}。首先，给定对象 $A$，
映射 $F(G(A)) \to A$ 为满射，因为由构造，每个映射 $U \to A$ 均通过
$F(G(A))$ 分解。另一方面，映射 $F(G(A)) \to A$ 伴随于映射
$\text{id} : G(A) \to G(A)$，故由引理 \ref{injectives-lemma-F-G-monos} 它为单射。

\medskip\noindent
由于 $F$ 是左伴随函子，它右正合。$\text{Mod}_R$ 有足够多的投射对象，
故要证明 $F$ 正合，只须证明第一左导出函子 $L_1F$ 为零。为对某个 $R$-模
$M$ 证明 $L_1F(M) = 0$，选取 $R$-模的正合列
$0 \to K \to P \to M \to 0$，其中 $P$ 自由。只须证明
$F(K) \to F(P)$ 为单射。可将此列写成正合列
$0 \to K_i \to P_i \to M_i \to 0$ 的滤过余极限，其中 $P_i$ 为有限自由
$R$-模：只须以这种方式写出 $P$，并置
$K_i = K \cap P_i$ 与 $M_i = \Im(P_i \to M)$。由于 $F$ 是左伴随函子，
它与余极限交换；又因 $\mathcal{A}$ 为 Grothendieck 阿贝尔范畴，若每个
$F(K_i) \to F(P_i)$ 都是单射，则 $F(K) \to F(P)$ 也是单射。因此只须在
$K \subset P = R^{\oplus n}$ 时验证 $F(K) \to F(P)$ 为单射。
此时应用引理 \ref{injectives-lemma-F-G-monos}，便知
$F(K) \to U^{\oplus n}$ 为单射。
\end{proof}

\begin{lemma}
\label{injectives-lemma-gabriel-popescu}
\begin{reference}
\cite[Corollary 4.1]{serpe}
\end{reference}
设 $\mathcal{A}$ 为 Grothendieck 阿贝尔范畴，$R$、$F$、$G$ 如
Gabriel--Popescu 定理（定理 \ref{injectives-theorem-gabriel-popescu}）中所述。
则得到导出函子
$$
RG : D(\mathcal{A}) \to D(\text{Mod}_R)
\quad\text{且}\quad
F : D(\text{Mod}_R) \to D(\mathcal{A})
$$
，其中 $F$ 是 $RG$ 的左伴随函子，$RG$ 全忠实，并且
$F \circ RG = \text{id}$。
\end{lemma}

\begin{proof}
这些函子的存在性与伴随性由定理 \ref{injectives-theorem-gabriel-popescu}、
\ref{injectives-theorem-K-injective-embedding-grothendieck} 以及《导出范畴》引理
\ref{derived-lemma-K-injective-defined}、
\ref{derived-lemma-right-derived-exact-functor}、
\ref{derived-lemma-derived-adjoint-functors} 得出。等式
$F \circ RG = \text{id}$ 的理由如下：在 $D(\mathcal{A})$ 的对象上，
可通过对适当的代表复形 $I^\bullet$（例如 K-内射代表）施用 $G$ 来计算
$RG$；而由 $F \circ G = \text{id}$，有
$F(G(I^\bullet)) = I^\bullet$。再由《范畴》引理
\ref{categories-lemma-adjoint-fully-faithful}，可知 $RG$ 全忠实。
\end{proof}





\section{Brown 可表性与 Grothendieck 阿贝尔范畴}
\label{injectives-section-brown}

\noindent
本节快速证明一个关于 Grothendieck 阿贝尔范畴之导出范畴的可表性定理。
读者应先阅读《导出范畴》第 \ref{derived-section-brown} 节中紧生成三角范畴
的情形。之后，与其阅读本节，不如查阅文献中这类更一般的结果；例如见
\cite{Franke}、\cite{Neeman}、\cite{Krause}，或参看《导出范畴》第
\ref{derived-section-brown-bis} 节。

\begin{lemma}
\label{injectives-lemma-brown}
设 $\mathcal{A}$ 为 Grothendieck 阿贝尔范畴，并设
$H : D(\mathcal{A}) \to \textit{Ab}$ 为把直和变为积的反变上同调函子。
则 $H$ 可表。
\end{lemma}

\begin{proof}
设 $R,F,G,RG$ 如引理 \ref{injectives-lemma-gabriel-popescu} 中所述，并考虑函子
$H \circ F : D(\text{Mod}_R) \to \textit{Ab}$。由于 $F$ 是左伴随函子，
它把直和映为直和，故 $H \circ F$ 把直和变为积。另一方面，导出范畴
$D(\text{Mod}_R)$ 由单个紧对象 $R$ 生成。由《导出范畴》引理
\ref{derived-lemma-brown}，$H \circ F$ 可表；设其由
$L \in D(\text{Mod}_R)$ 表示。选取特异三角
$$
M \to L \to RG(F(L)) \to M[1]
$$
于 $D(\text{Mod}_R)$ 中。因 $F \circ RG = \text{id}$，有 $F(M) = 0$。
从而 $H(F(M)) = 0$，进而 $\Hom(M, L) = 0$。因此
$L \to RG(F(L))$ 是一个直和因子的嵌入；见《导出范畴》引理
\ref{derived-lemma-split}。对 $D(\mathcal{A})$ 中的 $A$，得到
\begin{align*}
H(A)
& =
H(F(RG(A)) \\
& =
\Hom(RG(A), L) \\
& \to
\Hom(RG(A), RG(F(L))) \\
& =
\Hom(F(RG(A)), F(L)) \\
& =
\Hom(A, F(L))
\end{align*}
，其中箭头有一个关于 $A$ 函子性的左逆。换言之，$H$ 是某个可表函子的直和
因子。由于 $D(\mathcal{A})$ 是 Karoubi 完备的（《导出范畴》引理
\ref{derived-lemma-projectors-have-images-triangulated}），结论成立。
\end{proof}

\begin{proposition}
\label{injectives-proposition-brown}
设 $\mathcal{A}$ 为 Grothendieck 阿贝尔范畴，$\mathcal{D}$ 为三角范畴。
设 $F : D(\mathcal{A}) \to \mathcal{D}$ 为把直和映为直和的三角范畴间正合
函子。则 $F$ 有正合的右伴随函子。
\end{proposition}

\begin{proof}
对 $\mathcal{D}$ 的对象 $Y$，考虑反变函子
$$
D(\mathcal{A}) \to \textit{Ab},\quad
W \mapsto \Hom_\mathcal{D}(F(W), Y)
$$
由于 $F$ 正合，这是上同调函子；又因 $F$ 把直和映为直和，该反变函子把直和
变为积。
故由引理 \ref{injectives-lemma-brown}，存在 $D(\mathcal{A})$ 的对象 $X$，使
$\Hom_{D(\mathcal{A})}(W, X) = \Hom_\mathcal{D}(F(W), Y)$。
伴随函子的存在性由《范畴》引理 \ref{categories-lemma-adjoint-exists} 得出；
其正合性由《导出范畴》引理 \ref{derived-lemma-adjoint-is-exact} 得出。
\end{proof}
% Generated mechanically from the frozen cohomology.tex target.
% Do not edit this generated file; edit the bound locale source instead.

% OK, start here.
%

\setcounter{chapter}{19}
\chapter{层上同调}



\phantomsection
\label{cohomology-section-phantom}


\section{引言}
\label{cohomology-section-introduction}

\noindent
本章讨论拓扑空间上层的上同调若干专题。我们主要在环层上的模这一一般框架下
工作，并研究环空间的态射。关于站点上的层会出现什么情形，见《站点上的
上同调》第 \ref{sites-cohomology-section-introduction} 节。基本参考文献为
\cite{Godement} 与 \cite{Iversen}。





\section{层上同调}
\label{cohomology-section-cohomology-sheaves}

\noindent
设 $X$ 为拓扑空间，$\mathcal{F}$ 为阿贝尔层。我们已知 $X$ 上阿贝尔层的
范畴有足够多的内射对象；见《内射对象》引理
\ref{injectives-lemma-abelian-sheaves-space}。因此可以选取内射消解
$\mathcal{F}[0] \to \mathcal{I}^\bullet$。按照通常约定，定义
\begin{equation}
\label{cohomology-equation-cohomology}
H^i(X, \mathcal{F}) = H^i(\Gamma(X, \mathcal{I}^\bullet))
\end{equation}
并称之为阿贝尔层 $\mathcal{F}$ 的{\it 第 $i$ 个上同调群}。
函子族 $H^i(X, -)$ 构成从 $\textit{Ab}(X)$ 到 $\textit{Ab}$ 的
泛 $\delta$-函子。

\medskip\noindent
设 $f : X \to Y$ 为拓扑空间的连续映射。沿用上面的
$\mathcal{F}[0] \to \mathcal{I}^\bullet$，定义
\begin{equation}
\label{cohomology-equation-higher-direct-image}
R^if_*\mathcal{F} = H^i(f_*\mathcal{I}^\bullet)
\end{equation}
并称之为 $\mathcal{F}$ 的{\it 第 $i$ 个高阶直像}。
函子族 $R^if_*$ 构成从 $\textit{Ab}(X)$ 到 $\textit{Ab}(Y)$ 的
泛 $\delta$-函子。

\medskip\noindent
设 $(X, \mathcal{O}_X)$ 为环空间，$\mathcal{F}$ 为
$\mathcal{O}_X$-模。我们已知 $X$ 上 $\mathcal{O}_X$-模的范畴有足够多的
内射对象；见《内射对象》引理
\ref{injectives-lemma-sheaves-modules-space}。因此可以选取内射消解
$\mathcal{F}[0] \to \mathcal{I}^\bullet$。按照通常约定，定义
\begin{equation}
\label{cohomology-equation-cohomology-modules}
H^i(X, \mathcal{F}) = H^i(\Gamma(X, \mathcal{I}^\bullet))
\end{equation}
并称之为 $\mathcal{F}$ 的{\it 第 $i$ 个上同调群}。
函子族 $H^i(X, -)$ 构成从 $\textit{Mod}(\mathcal{O}_X)$ 到
$\text{Mod}_{\mathcal{O}_X(X)}$ 的泛 $\delta$-函子。

\medskip\noindent
设 $f : (X, \mathcal{O}_X) \to (Y, \mathcal{O}_Y)$ 为环空间的态射。
沿用上面的 $\mathcal{F}[0] \to \mathcal{I}^\bullet$，定义
\begin{equation}
\label{cohomology-equation-higher-direct-image-modules}
R^if_*\mathcal{F} = H^i(f_*\mathcal{I}^\bullet)
\end{equation}
并称之为 $\mathcal{F}$ 的{\it 第 $i$ 个高阶直像}。
函子族 $R^if_*$ 构成从 $\textit{Mod}(\mathcal{O}_X)$ 到
$\textit{Mod}(\mathcal{O}_Y)$ 的泛 $\delta$-函子。





\section{导出函子}
\label{cohomology-section-derived-functors}

\noindent
下面简要说明如何利用消解函子得到右导出函子。关于无界导出函子，见
第 \ref{cohomology-section-unbounded} 节。

\medskip\noindent
设 $(X, \mathcal{O}_X)$ 为环空间。范畴
$\textit{Mod}(\mathcal{O}_X)$ 是阿贝尔范畴；见《模层》引理
\ref{modules-lemma-abelian}。本章记
$$
K(\mathcal{O}_X) = K(\textit{Mod}(\mathcal{O}_X))
\quad
\text{以及}
\quad
D(\mathcal{O}_X) = D(\textit{Mod}(\mathcal{O}_X)).
$$
对于《导出范畴》定义 \ref{derived-definition-complexes-notation} 与定义
\ref{derived-definition-unbounded-derived-category} 中引入的三角范畴，
其有界版本采用类似记号。由《导出范畴》注
\ref{derived-remark-big-abelian-category}，存在消解函子
$$
j = j_X :
K^{+}(\textit{Mod}(\mathcal{O}_X))
\longrightarrow
K^{+}(\mathcal{I})
$$
其中 $\mathcal{I}$ 是 $\textit{Mod}(\mathcal{O}_X)$ 中由内射层构成的
严格满加性子范畴。对任意取值于阿贝尔范畴 $\mathcal{B}$ 的左正合函子
$F : \textit{Mod}(\mathcal{O}_X) \to \mathcal{B}$
，以 $RF$ 表示《导出范畴》第 \ref{derived-section-right-derived-functor} 节
所述、利用刚才的消解函子 $j_X$ 构造的右导出函子：
\begin{equation}
\label{cohomology-equation-RF}
RF = F \circ j_X' : D^{+}(X) \longrightarrow D^{+}(\mathcal{B})
\end{equation}
记号见《导出范畴》引理 \ref{derived-lemma-right-derived-functor}。
注意，视具体情形，可把 $RF$ 看作定义于
$\textit{Mod}(\mathcal{O}_X)$、
$\text{Comp}^{+}(\textit{Mod}(\mathcal{O}_X))$,
$K^{+}(X)$ 或 $D^{+}(X)$ 上。按照《导出范畴》定义
\ref{derived-definition-higher-derived-functors}，得到第 $i$ 个右导出函子
\begin{equation}
\label{cohomology-equation-RFi}
R^iF = H^i \circ RF : \textit{Mod}(\mathcal{O}_X) \longrightarrow \mathcal{B}
\end{equation}
从而 $R^0F = F$，且 $\{R^iF, \delta\}_{i \geq 0}$ 是泛
$\delta$-函子；见《导出范畴》引理
\ref{derived-lemma-higher-derived-functors}。

\medskip\noindent
下面给出这一构造的两个特殊情形。给定环 $R$，记
$K(R) = K(\text{Mod}_R)$ 与 $D(R) = D(\text{Mod}_R)$，有界版本也采用
类似记号。对任意开集 $U \subset X$，有左正合函子
$
\Gamma(U, -) :
\textit{Mod}(\mathcal{O}_X)
\longrightarrow
\text{Mod}_{\mathcal{O}_X(U)}
$
，由上述讨论得到
\begin{equation}
\label{cohomology-equation-total-derived-cohomology}
R\Gamma(U, -) :
D^{+}(X)
\longrightarrow
D^{+}(\mathcal{O}_X(U))
\end{equation}
。置 $H^i(U, -) = R^i\Gamma(U, -)$。当 $U = X$ 时，便恢复
(\ref{cohomology-equation-cohomology-modules})。若 $f : X \to Y$ 是环空间的态射，
则有左正合函子
$
f_* :
\textit{Mod}(\mathcal{O}_X)
\longrightarrow
\textit{Mod}(\mathcal{O}_Y)
$
，由此得到{\it 导出直像}
\begin{equation}
\label{cohomology-equation-total-derived-direct-image}
Rf_* :
D^{+}(X)
\longrightarrow
D^{+}(Y)
\end{equation}
。$Rf_*\mathcal{F}^\bullet$ 的第 $i$ 个上同调层记为
$R^if_*\mathcal{F}^\bullet$，并依照
(\ref{cohomology-equation-higher-direct-image-modules}) 称为第 $i$ 个{\it 高阶直像}。
上述两个显示式中的函子都是导出范畴之间的正合函子。

\medskip\noindent
{\bf 记号的滥用：}当函子 $Rf_*$ 或任何其他导出函子作用于 $X$ 上的层
$\mathcal{F}$ 或层复形时，约定先把 $\mathcal{F}$ 替换为一个适当的消解。
为方便进行这类操作，对给定对象
$\mathcal{F}^\bullet \in D(\mathcal{O}_X)$，若存在拟同构
$\mathcal{F}^\bullet \to \mathcal{I}^\bullet$，其中
$\mathcal{I}^\bullet$ 是由 $\textit{Mod}(\mathcal{O}_X)$ 中内射对象组成的
下有界复形，则称 $\mathcal{I}^\bullet$ 在导出范畴中{\it 表示
$\mathcal{F}^\bullet$}。同理，“设
$\alpha : \mathcal{F}^\bullet \to \mathcal{G}^\bullet$ 为
$D(\mathcal{O}_X)$ 的态射”并不意味着 $\alpha$ 由复形态射表示。若所指确为
复形态射，我们会明确说明。









\section{第一上同调与挠子}
\label{cohomology-section-h1-torsors}

\begin{definition}
\label{cohomology-definition-torsor}
设 $X$ 为拓扑空间，$\mathcal{G}$ 为 $X$ 上的（未必交换的）群层。
所谓{\it 伪挠子}，更准确地说，所谓{\it 伪 $\mathcal{G}$-挠子}，是指
$X$ 上的集合层 $\mathcal{F}$，配备作用
$\mathcal{G} \times \mathcal{F} \to \mathcal{F}$，使得
\begin{enumerate}
\item 只要 $\mathcal{F}(U)$ 非空，作用
$\mathcal{G}(U) \times \mathcal{F}(U) \to \mathcal{F}(U)$
便是单可迁的；
\end{enumerate}
伪 $\mathcal{G}$-挠子的态射 $\mathcal{F} \to \mathcal{F}'$ 是与
$\mathcal{G}$-作用相容的集合层态射。所谓{\it 挠子}，更准确地说，所谓
{\it $\mathcal{G}$-挠子}，是还满足下列条件的伪 $\mathcal{G}$-挠子：
\begin{enumerate}
\item[(2)] 对每个 $x \in X$，茎 $\mathcal{F}_x$ 非空。
\end{enumerate}
$\mathcal{G}$-挠子的态射是伪 $\mathcal{G}$-挠子的态射。
{\it 平凡 $\mathcal{G}$-挠子}是配备显然的左 $\mathcal{G}$-作用的层
$\mathcal{G}$。
\end{definition}

\noindent
显然，挠子之间的态射自动是同构。

\begin{lemma}
\label{cohomology-lemma-trivial-torsor}
设 $X$ 为拓扑空间，$\mathcal{G}$ 为 $X$ 上的（未必交换的）群层。
$\mathcal{G}$-挠子 $\mathcal{F}$ 平凡，当且仅当
$\mathcal{F}(X) \not = \emptyset$。
\end{lemma}

\begin{proof}
从略。
\end{proof}

\begin{lemma}
\label{cohomology-lemma-torsors-h1}
设 $X$ 为拓扑空间，$\mathcal{H}$ 为 $X$ 上的阿贝尔层。
$\mathcal{H}$-挠子的同构类集合与 $H^1(X, \mathcal{H})$ 之间有典范双射。
\end{lemma}

\begin{proof}
设 $\mathcal{F}$ 为 $\mathcal{H}$-挠子。考虑由 $\mathcal{F}$ 生成的
自由阿贝尔层 $\mathbf{Z}[\mathcal{F}]$。它是下述规则的层化：对开集
$U \subset X$，取所有有限形式和 $\sum n_i[s_i]$，其中
$n_i \in \mathbf{Z}$ 且 $s_i \in \mathcal{F}(U)$。存在自然映射
$$
\sigma : \mathbf{Z}[\mathcal{F}] \longrightarrow \underline{\mathbf{Z}}
$$
把局部截面 $\sum n_i[s_i]$ 映到 $\sum n_i$。$\sigma$ 的核由形如
$[s] - [s']$ 的局部截面生成。存在典范映射
$a : \Ker(\sigma) \to \mathcal{H}$
，把 $[s] - [s']$ 映到满足 $h \cdot s' = s$ 的 $\mathcal{H}$ 的局部截面
$h$。考虑推出图
$$
\xymatrix{
0 \ar[r] &
\Ker(\sigma) \ar[r] \ar[d]^a &
\mathbf{Z}[\mathcal{F}] \ar[r] \ar[d] &
\underline{\mathbf{Z}} \ar[r] \ar[d] &
0 \\
0 \ar[r] &
\mathcal{H} \ar[r] &
\mathcal{E} \ar[r] &
\underline{\mathbf{Z}} \ar[r] &
0
}
$$
其中 $\mathcal{E}$ 是通过推出得到的扩张。对下方短正合列应用相应的上同调
长正合列，并把连接同态作用于
$1 \in H^0(X, \underline{\mathbf{Z}})$，得到元素
$\xi = \xi_\mathcal{F} \in H^1(X, \mathcal{H})$
。

\medskip\noindent
反过来，给定 $\xi \in H^1(X, \mathcal{H})$，可按如下方式为它配上一个
挠子。选取 $\mathcal{H}$ 到内射阿贝尔层 $\mathcal{I}$ 中的嵌入
$\mathcal{H} \to \mathcal{I}$。置
$\mathcal{Q} = \mathcal{I}/\mathcal{H}$，于是有短正合列
$$
\xymatrix{
0 \ar[r] &
\mathcal{H} \ar[r] &
\mathcal{I} \ar[r] &
\mathcal{Q} \ar[r] &
0
}
$$
因为 $H^1(X, \mathcal{I}) = 0$（见《导出范畴》引理
\ref{derived-lemma-higher-derived-functors}），元素 $\xi$ 是某个整体截面
$q \in H^0(X, \mathcal{Q})$ 的像。令 $\mathcal{F} \subset \mathcal{I}$
为由在层 $\mathcal{Q}$ 中映到 $q$ 的截面组成的子层（视为集合层）。容易验证
$\mathcal{F}$ 是挠子。

\medskip\noindent
关于上述两个构造互为逆构造的验证从略。
\end{proof}







\section{第一上同调与扩张}
\label{cohomology-section-h1-extensions}

\begin{lemma}
\label{cohomology-lemma-h1-extensions}
设 $(X, \mathcal{O}_X)$ 为环空间，$\mathcal{F}$ 为
$\mathcal{O}_X$-模层。存在典范双射
$$
\Ext^1_{\textit{Mod}(\mathcal{O}_X)}(\mathcal{O}_X, \mathcal{F})
\longrightarrow
H^1(X, \mathcal{F})
$$
，它把扩张
$$
0 \to \mathcal{F} \to \mathcal{E} \to \mathcal{O}_X \to 0
$$
映到 $1 \in \Gamma(X, \mathcal{O}_X)$ 在 $H^1(X, \mathcal{F})$ 中的像。
\end{lemma}

\begin{proof}
下面构造引理中映射的逆映射。设 $\xi \in H^1(X, \mathcal{F})$。
选取嵌入 $\mathcal{F} \subset \mathcal{I}$，其中 $\mathcal{I}$ 是
$\textit{Mod}(\mathcal{O}_X)$ 中的内射对象。置
$\mathcal{Q} = \mathcal{I}/\mathcal{F}$。由上同调长正合列，$\xi$ 是截面
$\tilde \xi \in \Gamma(X, \mathcal{Q}) =
\Hom_{\mathcal{O}_X}(\mathcal{O}_X, \mathcal{Q})$ 的像。现在取拉回
$$
\xymatrix{
0 \ar[r] &
\mathcal{F} \ar[r] \ar@{=}[d] &
\mathcal{E} \ar[r] \ar[d] &
\mathcal{O}_X \ar[r] \ar[d]^{\tilde \xi} &
0 \\
0 \ar[r] &
\mathcal{F} \ar[r] &
\mathcal{I} \ar[r] &
\mathcal{Q} \ar[r] &
0
}
$$
；见《同调》第 \ref{homology-section-extensions} 节。
\end{proof}








\section{第一上同调与可逆模层}
\label{cohomology-section-invertible-sheaves}

\noindent
环空间的 Picard 群定义见《模层》第 \ref{modules-section-invertible} 节。

\begin{lemma}
\label{cohomology-lemma-h1-invertible}
设 $(X, \mathcal{O}_X)$ 为环空间。若所有茎 $\mathcal{O}_{X, x}$ 都是
局部环，则存在阿贝尔群的典范同构
$$
H^1(X, \mathcal{O}_X^*) = \Pic(X).
$$
\end{lemma}

\begin{proof}
设 $\mathcal{L}$ 为可逆 $\mathcal{O}_X$-模。考虑由下述规则定义的预层
$\mathcal{L}^*$：
$$
U \longmapsto \{s \in \mathcal{L}(U)
\text{ 使得 } \mathcal{O}_U \xrightarrow{s \cdot -} \mathcal{L}_U
\text{ 是同构}\}
$$
该预层满足层条件。此外，若 $f \in \mathcal{O}_X^*(U)$ 且
$s \in \mathcal{L}^*(U)$，则显然 $fs \in \mathcal{L}^*(U)$。同理，若
$s, s' \in \mathcal{L}^*(U)$，则存在唯一的
$f \in \mathcal{O}_X^*(U)$ 使 $fs = s'$。再者，由《模层》引理
\ref{modules-lemma-invertible-is-locally-free-rank-1}，层
$\mathcal{L}^*$ 局部有截面。换言之，$\mathcal{L}^*$ 是
$\mathcal{O}_X^*$-挠子。因此得到映射
$$
\begin{matrix}
\text{环空间 }(X, \mathcal{O}_X)\text{ 上的可逆模层} \\
\text{不计同构}
\end{matrix}
\longrightarrow
\begin{matrix}
\mathcal{O}_X^*\text{-挠子} \\
\text{不计同构}
\end{matrix}
$$
关于它是阿贝尔群同态的验证从略。由引理 \ref{cohomology-lemma-torsors-h1}，右侧与
$H^1(X, \mathcal{O}_X^*)$ 典范双射。因此只需证明此映射既单又满。

\medskip\noindent
单射性。若挠子 $\mathcal{L}^*$ 平凡，则由引理
\ref{cohomology-lemma-trivial-torsor}，$\mathcal{L}^*$ 有整体截面。这恰好意味着
$\mathcal{L} \cong \mathcal{O}_X$ 是 $\Pic(X)$ 中的单位元。

\medskip\noindent
满射性。设 $\mathcal{F}$ 为 $\mathcal{O}_X^*$-挠子。考虑集合预层
$$
\mathcal{L}_1 : U \longmapsto
(\mathcal{F}(U) \times \mathcal{O}_X(U))/\mathcal{O}_X^*(U)
$$
，其中 $f \in \mathcal{O}_X^*(U)$ 在 $(s, g)$ 上的作用为
$(fs, f^{-1}g)$。规定
$(s, g) + (s', g') = (s, g + (s'/s)g')$，其中 $s'/s$ 是满足
$fs = s'$ 的 $\mathcal{O}_X^*$ 的局部截面 $f$；并对
$\mathcal{O}_X$ 的局部截面 $h$ 规定 $h(s, g) = (s, hg)$。如此
$\mathcal{L}_1$ 成为 $\mathcal{O}_X$-模预层。关于其层化
$\mathcal{L} = \mathcal{L}_1^\#$ 是可逆 $\mathcal{O}_X$-模，且相应的
$\mathcal{O}_X^*$-挠子 $\mathcal{L}^*$ 同构于 $\mathcal{F}$ 的验证从略。
\end{proof}













\section{上同调的局部性}
\label{cohomology-section-locality}

\noindent
下面的引理说明，在开集上定义层 $\mathcal{F}$ 的上同调不会产生歧义。

\begin{lemma}
\label{cohomology-lemma-cohomology-of-open}
设 $X$ 是环化空间。
设 $U \subset X$ 是开子空间。
\begin{enumerate}
\item 若 $\mathcal{I}$ 是内射 $\mathcal{O}_X$-模，
则 $\mathcal{I}|_U$ 是内射 $\mathcal{O}_U$-模。
\item 对任意 $\mathcal{O}_X$-模层 $\mathcal{F}$，有
$H^p(U, \mathcal{F}) = H^p(U, \mathcal{F}|_U)$.
\end{enumerate}
\end{lemma}

\begin{proof}
以 $j : U \to X$ 表示该开浸入。
回忆限制到 $U$ 的函子 $j^{-1}$ 是延零函子 $j_!$ 的右伴随，见
《层》引理 \ref{sheaves-lemma-j-shriek-modules}。
此外，$j_!$ 是正合的。因此 (1) 由
《同调代数》引理 \ref{homology-lemma-adjoint-preserve-injectives}
推出。

\medskip\noindent
按定义，
$H^p(U, \mathcal{F}) = H^p(\Gamma(U, \mathcal{I}^\bullet))$，
其中 $\mathcal{F} \to \mathcal{I}^\bullet$ 是
$\textit{Mod}(\mathcal{O}_X)$ 中的内射消解。
由上可知，$\mathcal{F}|_U \to \mathcal{I}^\bullet|_U$
是 $\textit{Mod}(\mathcal{O}_U)$ 中的内射消解。
因此 $H^p(U, \mathcal{F}|_U)$ 等于
$H^p(\Gamma(U, \mathcal{I}^\bullet|_U))$。
当然，对 $X$ 上任意层 $\mathcal{F}$，都有
$\Gamma(U, \mathcal{F}) = \Gamma(U, \mathcal{F}|_U)$。
故 (2) 中的等式成立。
\end{proof}

\noindent
设 $X$ 是环化空间，$\mathcal{F}$ 是 $\mathcal{O}_X$-模层，
并设 $U \subset V \subset X$ 是开子集。
则存在典范的{\it 限制映射}
\begin{equation}
\label{cohomology-equation-restriction-mapping}
H^n(V, \mathcal{F})
\longrightarrow
H^n(U, \mathcal{F}), \quad
\xi \longmapsto \xi|_U
\end{equation}
且它关于 $\mathcal{F}$ 是函子性的。事实上，任取内射消解
$\mathcal{F} \to \mathcal{I}^\bullet$。层 $\mathcal{I}^p$
的限制映射给出复形态射
$$
\Gamma(V, \mathcal{I}^\bullet)
\longrightarrow
\Gamma(U, \mathcal{I}^\bullet)
$$
左端是表示 $R\Gamma(V, \mathcal{F})$ 的复形，
右端是表示 $R\Gamma(U, \mathcal{F})$ 的复形。
施用函子 $H^n$ 即得到上同调群之间的映射。
如上所示，我们用记号 $\xi \mapsto \xi|_U$ 表示此映射。
因此规则 $U \mapsto H^n(U, \mathcal{F})$ 构成
$\mathcal{O}_X$-模预层。该预层通常记作
$\underline{H}^n(\mathcal{F})$。
我们将在引理 \ref{cohomology-lemma-include} 中给出它的另一种解释。

\begin{lemma}
\label{cohomology-lemma-kill-cohomology-class-on-covering}
设 $X$ 是环化空间，$\mathcal{F}$ 是 $\mathcal{O}_X$-模层，
$U \subset X$ 是开子空间。设 $n > 0$ 且
$\xi \in H^n(U, \mathcal{F})$。则存在开覆盖
$U = \bigcup_{i\in I} U_i$，使得对所有 $i \in I$，
均有 $\xi|_{U_i} = 0$。
\end{lemma}

\begin{proof}
设 $\mathcal{F} \to \mathcal{I}^\bullet$ 是内射消解。于是
$$
H^n(U, \mathcal{F}) =
\frac{\Ker(\mathcal{I}^n(U) \to \mathcal{I}^{n + 1}(U))}
{\Im(\mathcal{I}^{n - 1}(U) \to \mathcal{I}^n(U))}.
$$
取表示上述商中该上同调类的元素
$\tilde \xi \in \mathcal{I}^n(U)$。
由于 $\mathcal{I}^\bullet$ 是 $\mathcal{F}$ 的内射消解且 $n > 0$，
复形 $\mathcal{I}^\bullet$ 在次数 $n$ 处正合。因此，作为层有
$\Im(\mathcal{I}^{n - 1} \to \mathcal{I}^n) =
\Ker(\mathcal{I}^n \to \mathcal{I}^{n + 1})$。
由于 $\tilde \xi$ 是核层在 $U$ 上的截面，存在开覆盖
$U = \bigcup_{i \in I} U_i$，使得
$\tilde \xi|_{U_i}$ 是某个截面
$\xi_i \in \mathcal{I}^{n - 1}(U_i)$ 在 $d$ 下的像。
按限制 $\xi|_{U_i}$ 的定义，它对应于
$\tilde \xi|_{U_i}$ 的类，故结论成立。
\end{proof}

\begin{lemma}
\label{cohomology-lemma-describe-higher-direct-images}
设 $f : X \to Y$ 是环化空间的态射，
$\mathcal{F}$ 是 $\mathcal{O}_X$-模。
层 $R^if_*\mathcal{F}$ 是与预层
$$
V \longmapsto H^i(f^{-1}(V), \mathcal{F})
$$
相伴的层，其中限制映射如
Equation (\ref{cohomology-equation-restriction-mapping}) 所示。
对 $R^if_*$ 施用于下有界复形 $\mathcal{F}^\bullet$，也有类似陈述。
\end{lemma}

\begin{proof}
设 $\mathcal{F} \to \mathcal{I}^\bullet$ 是内射消解。
按定义，$R^if_*\mathcal{F}$ 是复形
$$
f_*\mathcal{I}^0 \to f_*\mathcal{I}^1 \to f_*\mathcal{I}^2 \to \ldots
$$
的第 $i$ 个上同调层。由 $\mathcal{O}_Y$-模上阿贝尔范畴结构的定义，
该上同调层是与预层
$$
V
\longmapsto
\frac{\Ker(f_*\mathcal{I}^i(V) \to f_*\mathcal{I}^{i + 1}(V))}
{\Im(f_*\mathcal{I}^{i - 1}(V) \to f_*\mathcal{I}^i(V))}
$$
相伴的层，而此商显然等于
$$
\frac{\Ker(\mathcal{I}^i(f^{-1}(V)) \to \mathcal{I}^{i + 1}(f^{-1}(V)))}
{\Im(\mathcal{I}^{i - 1}(f^{-1}(V)) \to \mathcal{I}^i(f^{-1}(V)))}
$$
后者等于 $H^i(f^{-1}(V), \mathcal{F})$，故结论成立。
\end{proof}

\begin{lemma}
\label{cohomology-lemma-localize-higher-direct-images}
设 $f : X \to Y$ 是环化空间的态射，
$\mathcal{F}$ 是 $\mathcal{O}_X$-模，且
$V \subset Y$ 是开子空间。以
$g : f^{-1}(V) \to V$ 表示 $f$ 的限制。则有
$$
R^pg_*(\mathcal{F}|_{f^{-1}(V)}) = (R^pf_*\mathcal{F})|_V
$$
对导出直像 $Rf_*\mathcal{F}^\bullet$ 也有类似陈述，
其中 $\mathcal{F}^\bullet$ 是 $\mathcal{O}_X$-模的下有界复形。
\end{lemma}

\begin{proof}
第一种证明。施用引理 \ref{cohomology-lemma-describe-higher-direct-images}
和 \ref{cohomology-lemma-cohomology-of-open} 即得所示等式。
第二种证明。取内射消解 $\mathcal{F} \to \mathcal{I}^\bullet$，
并利用
$\mathcal{F}|_{f^{-1}(V)} \to \mathcal{I}^\bullet|_{f^{-1}(V)}$
也是内射消解。
\end{proof}

\begin{remark}
\label{cohomology-remark-daniel}
下面给出上述引理 \ref{cohomology-lemma-kill-cohomology-class-on-covering}
和 \ref{cohomology-lemma-describe-higher-direct-images} 的另一种证明方法。
设 $(X, \mathcal{O}_X)$ 是环化空间，
$i_X : \textit{Mod}(\mathcal{O}_X) \to \textit{PMod}(\mathcal{O}_X)$
是包含函子，$\#$ 是层化函子。回忆 $i_X$ 左正合而 $\#$ 正合。
\begin{enumerate}
\item 先证明下面的引理 \ref{cohomology-lemma-include}；它说明 $i_X$ 的右导出函子为
$R^pi_X\mathcal{F} = \underline{H}^p(\mathcal{F})$.
也可如下证明：当 $p = 0$ 时等式显然成立。
$(R^pi_X)_{p \geq 0}$ 与 $(\underline{H}^p)_{p \geq 0}$
都是在内射对象上消失的 $\delta$-函子，故二者都是泛的，从而彼此同构。见
《同调代数》第 \ref{homology-section-cohomological-delta-functor} 节。
\item 引理 \ref{cohomology-lemma-kill-cohomology-class-on-covering}
可重述为：对任意 $\mathcal{O}_X$-模层 $\mathcal{F}$ 及 $p > 0$，
有 $(\underline{H}^p(\mathcal{F}))^\# = 0$。
为说明这一点，利用 ${}^\#$ 正合可得
$$
(\underline{H}^p(\mathcal{F}))^\# =
(R^pi_X\mathcal{F})^\# =
R^p(\# \circ i_X)(\mathcal{F}) = 0
$$
因为 $\# \circ i_X$ 是恒等函子。
\item 设 $f : X \to Y$ 是环化空间的态射，
$\mathcal{F}$ 是 $\mathcal{O}_X$-模。预层
$V \mapsto H^p(f^{-1}V, \mathcal{F})$ 等于
$R^p (i_Y \circ f_*)\mathcal{F}$。
与上面 (1) 的论证一样，注意二者都给出泛 $\delta$-函子即可证明此点。
因此引理 \ref{cohomology-lemma-describe-higher-direct-images} 说明
$R^p f_* \mathcal{F}= (R^p (i_Y \circ f_*)\mathcal{F})^\#$。
再次利用 $\#$ 正合且 $\# \circ i_Y$ 是恒等函子，得到
$$
R^p f_* \mathcal{F} =
R^p(\# \circ i_Y \circ f_*)\mathcal{F} =
(R^p (i_Y \circ f_*)\mathcal{F})^\#
$$
如所需。
\end{enumerate}
\end{remark}






\section{Mayer--Vietoris 序列}
\label{cohomology-section-mayer-vietoris}

\noindent
下文将构造从 {\v C}ech 上同调到层上同调的谱序列，见引理 \ref{cohomology-lemma-cech-spectral-sequence}。
该谱序列的一个特例是 Mayer--Vietoris 长正合序列。
由于它是谱序列的一种基础、实用而易于理解的变体，
我们在此单独讨论。

\begin{lemma}
\label{cohomology-lemma-injective-restriction-surjective}
\begin{slogan}
内射层是松弛层。
\end{slogan}
设 $X$ 是环化空间，$U' \subset U \subset X$ 是开子空间。
对任意内射 $\mathcal{O}_X$-模 $\mathcal{I}$，
限制映射 $\mathcal{I}(U) \to \mathcal{I}(U')$ 是满射。
\end{lemma}

\begin{proof}
设 $j : U \to X$ 和 $j' : U' \to X$ 为开浸入。
回忆 $j_!\mathcal{O}_U$ 是
$\mathcal{O}_U = \mathcal{O}_X|_U$ 的延零，见
《层》第 \ref{sheaves-section-open-immersions} 节。
由于 $j_!$ 是限制函子的左伴随，对任意 $\mathcal{O}_X$-模层
$\mathcal{F}$，有
$$
\Hom_{\mathcal{O}_X}(j_!\mathcal{O}_U, \mathcal{F})
=
\Hom_{\mathcal{O}_U}(\mathcal{O}_U, \mathcal{F}|_U)
=
\mathcal{F}(U)
$$
见 《层》引理 \ref{sheaves-lemma-j-shriek-modules}。
类似地，层 $j'_!\mathcal{O}_{U'}$ 表示函子
$\mathcal{F} \mapsto \mathcal{F}(U')$。
此外有一个显然的典范 $\mathcal{O}_X$-模映射
$$
j'_!\mathcal{O}_{U'} \longrightarrow j_!\mathcal{O}_U
$$
由 Yoneda 引理（《范畴》引理 \ref{categories-lemma-yoneda}），
它对应于限制映射
$\mathcal{F}(U) \to \mathcal{F}(U')$。
根据层 $j'_!\mathcal{O}_{U'}$ 和 $j_!\mathcal{O}_U$
之茎的描述可知，上述映射是单射（见上引引理）。
因此，若 $\mathcal{I}$ 是内射 $\mathcal{O}_X$-模，则映射
$$
\Hom_{\mathcal{O}_X}(j_!\mathcal{O}_U, \mathcal{I})
\longrightarrow
\Hom_{\mathcal{O}_X}(j'_!\mathcal{O}_{U'}, \mathcal{I})
$$
是满射，见
《同调代数》引理 \ref{homology-lemma-characterize-injectives}。
综合以上事实即得引理。
\end{proof}

\begin{lemma}[Mayer-Vietoris]
\label{cohomology-lemma-mayer-vietoris}
设 $X$ 是环化空间，并设 $X = U \cup V$ 是两个开子集之并。
对每个 $\mathcal{O}_X$-模 $\mathcal{F}$，存在上同调长正合序列
$$
0 \to
H^0(X, \mathcal{F}) \to
H^0(U, \mathcal{F}) \oplus H^0(V, \mathcal{F}) \to
H^0(U \cap V, \mathcal{F}) \to
H^1(X, \mathcal{F}) \to \ldots
$$
该长正合序列关于 $\mathcal{F}$ 是函子性的。
\end{lemma}

\begin{proof}
层条件说明，对任意阿贝尔层 $\mathcal{F}$，映射
$(1, -1) : \mathcal{F}(U) \oplus \mathcal{F}(V) \to \mathcal{F}(U \cap V)$
的核等于 $\mathcal{F}(X)$ 在第一个映射下的像。
上面的引理 \ref{cohomology-lemma-injective-restriction-surjective} 说明，
只要 $\mathcal{I}$ 是内射 $\mathcal{O}_X$-模，映射
$(1, -1) : \mathcal{I}(U) \oplus \mathcal{I}(V) \to \mathcal{I}(U \cap V)$
就是满射。因此，若
$\mathcal{F} \to \mathcal{I}^\bullet$ 是 $\mathcal{F}$ 的内射消解，
便得到复形的短正合序列
$$
0 \to
\mathcal{I}^\bullet(X) \to
\mathcal{I}^\bullet(U) \oplus \mathcal{I}^\bullet(V) \to
\mathcal{I}^\bullet(U \cap V) \to
0.
$$
取上同调即得结论（使用
《同调代数》引理 \ref{homology-lemma-long-exact-sequence-cochain}）。
该序列的函子性证明从略。
\end{proof}

\begin{lemma}[相对 Mayer--Vietoris]
\label{cohomology-lemma-relative-mayer-vietoris}
设 $f : X \to Y$ 是环化空间的态射，并设
$X = U \cup V$ 是两个开子集之并。记
$a = f|_U : U \to Y$、$b = f|_V : V \to Y$，以及
$c = f|_{U \cap V} : U \cap V \to Y$。
对每个 $\mathcal{O}_X$-模 $\mathcal{F}$，存在长正合序列
$$
0 \to
f_*\mathcal{F} \to
a_*(\mathcal{F}|_U) \oplus b_*(\mathcal{F}|_V) \to
c_*(\mathcal{F}|_{U \cap V}) \to
R^1f_*\mathcal{F} \to \ldots
$$
该长正合序列关于 $\mathcal{F}$ 是函子性的。
\end{lemma}

\begin{proof}
设 $\mathcal{F} \to \mathcal{I}^\bullet$ 是 $\mathcal{F}$ 的内射消解。
我们断言，存在复形的短正合序列
$$
0 \to
f_*\mathcal{I}^\bullet \to
a_*\mathcal{I}^\bullet|_U \oplus b_*\mathcal{I}^\bullet|_V \to
c_*\mathcal{I}^\bullet|_{U \cap V} \to
0.
$$
事实上，对任意开集 $W \subset Y$ 和任意 $n \geq 0$，
相应的 $W$ 上截面群序列
$$
0 \to
\mathcal{I}^n(f^{-1}(W)) \to
\mathcal{I}^n(U \cap f^{-1}(W))
\oplus \mathcal{I}^n(V \cap f^{-1}(W)) \to
\mathcal{I}^n(U \cap V \cap f^{-1}(W)) \to
0
$$
已在引理 \ref{cohomology-lemma-mayer-vietoris} 的证明中证为短正合。
取上同调层，并利用 $\mathcal{I}^\bullet|_U$
是 $\mathcal{F}|_U$ 的内射消解这一事实，即得引理。
对 $\mathcal{I}^\bullet|_V$、
$\mathcal{I}^\bullet|_{U \cap V}$ 亦然，见引理 \ref{cohomology-lemma-cohomology-of-open}。
\end{proof}




















\section{{\v C}ech 复形与 {\v C}ech 上同调}
\label{cohomology-section-cech}

\noindent
设 $X$ 是拓扑空间，
$\mathcal{U} : U = \bigcup_{i \in I} U_i$ 是开覆盖，见
《拓扑》基本概念（\ref{topology-item-covering}）。
按照惯例，以
$U_{i_0\ldots i_p} = U_{i_0} \cap \ldots \cap U_{i_p}$
表示 $\mathcal{U}$ 中成员的 $(p + 1)$ 重交。
设 $\mathcal{F}$ 是 $X$ 上的阿贝尔预层。置
$$
\check{\mathcal{C}}^p(\mathcal{U}, \mathcal{F})
=
\prod\nolimits_{(i_0, \ldots, i_p) \in I^{p + 1}}
\mathcal{F}(U_{i_0\ldots i_p}).
$$
这是一个阿贝尔群。对
$s \in \check{\mathcal{C}}^p(\mathcal{U}, \mathcal{F})$，
以 $s_{i_0\ldots i_p}$ 表示它在
$\mathcal{F}(U_{i_0\ldots i_p})$ 中的分量。
注意，若 $s \in \check{\mathcal{C}}^1(\mathcal{U}, \mathcal{F})$
且 $i, j \in I$，则 $s_{ij}$ 与 $s_{ji}$
都是 $\mathcal{F}(U_i \cap U_j)$ 的元素，但我们并未规定它们之间的关系。
换言之，这里{\it 不}使用交错上链（将在
第 \ref{cohomology-section-alternating-cech} 节 中定义）。定义
$$
d : \check{\mathcal{C}}^p(\mathcal{U}, \mathcal{F})
\longrightarrow
\check{\mathcal{C}}^{p + 1}(\mathcal{U}, \mathcal{F})
$$
为
\begin{equation}
\label{cohomology-equation-d-cech}
d(s)_{i_0\ldots i_{p + 1}}
=
\sum\nolimits_{j = 0}^{p + 1}
(-1)^j
s_{i_0\ldots \hat i_j \ldots i_{p + 1}}|_{U_{i_0\ldots i_{p + 1}}}
\end{equation}
直接可见 $d \circ d = 0$。换言之，
$\check{\mathcal{C}}^\bullet(\mathcal{U}, \mathcal{F})$ 是复形。

\begin{definition}
\label{cohomology-definition-cech-complex}
设 $X$ 是拓扑空间，
$\mathcal{U} : U = \bigcup_{i \in I} U_i$ 是开覆盖，
$\mathcal{F}$ 是 $X$ 上的阿贝尔预层。
复形 $\check{\mathcal{C}}^\bullet(\mathcal{U}, \mathcal{F})$
称为与 $\mathcal{F}$ 和开覆盖 $\mathcal{U}$ 相伴的
{\it {\v C}ech 复形}。其上同调群
$H^i(\check{\mathcal{C}}^\bullet(\mathcal{U}, \mathcal{F}))$
称为与 $\mathcal{F}$ 和覆盖 $\mathcal{U}$ 相伴的
{\it {\v C}ech 上同调群}，记作
$\check H^i(\mathcal{U}, \mathcal{F})$。
\end{definition}

\begin{lemma}
\label{cohomology-lemma-cech-h0}
设 $X$ 是拓扑空间，$\mathcal{F}$ 是 $X$ 上的阿贝尔预层。
下列条件等价：
\begin{enumerate}
\item $\mathcal{F}$ 是阿贝尔层；
\item 对每个开覆盖 $\mathcal{U} : U = \bigcup_{i \in I} U_i$，
自然映射
$$
\mathcal{F}(U) \to \check{H}^0(\mathcal{U}, \mathcal{F})
$$
是双射。
\end{enumerate}
\end{lemma}

\begin{proof}
这是因为层条件恰好说明，对每个开覆盖，
$\mathcal{F}(U) \to \check{H}^0(\mathcal{U}, \mathcal{F})$
都是双射。
\end{proof}

\begin{lemma}
\label{cohomology-lemma-cech-trivial}
设 $X$ 是拓扑空间，$\mathcal{F}$ 是 $X$ 上的阿贝尔预层，
$\mathcal{U} : U = \bigcup_{i \in I} U_i$ 是开覆盖。
若对某个 $i \in I$ 有 $U_i = U$，则增广 {\v C}ech 复形
$$
\mathcal{F}(U) \to \check{\mathcal{C}}^\bullet(\mathcal{U}, \mathcal{F})
$$
由将 $\mathcal{F}(U)$ 置于次数 $-1$ 并取微分为从
$\mathcal{F}(U)$ 到
$\check{\mathcal{C}}^0(\mathcal{U}, \mathcal{F})$
的典范映射而得；该复形链同伦等价于零复形。
\end{lemma}

\begin{proof}
固定满足 $U = U_i$ 的元素 $i \in I$。注意，当 $i_j = i$ 时，
$U_{i_0 \ldots i_p} = U_{i_0 \ldots \hat i_j \ldots i_p}$。
定义同伦
$$
h :
\prod\nolimits_{i_0 \ldots i_{p + 1}} \mathcal{F}(U_{i_0 \ldots i_{p + 1}})
\longrightarrow
\prod\nolimits_{i_0 \ldots i_p} \mathcal{F}(U_{i_0 \ldots i_p})
$$
如下：
$$
h(s)_{i_0 \ldots i_p} = s_{i i_0 \ldots i_p}
$$
换言之，
$h : \prod_{i_0} \mathcal{F}(U_{i_0}) \to \mathcal{F}(U)$
是到因子 $\mathcal{F}(U_i) = \mathcal{F}(U)$ 的投影；
一般地，映射 $h$ 等于到下列因子的投影：
$\mathcal{F}(U_{i i_1 \ldots i_{p + 1}}) =
\mathcal{F}(U_{i_1 \ldots i_{p + 1}})$.
计算得
\begin{align*}
(dh + hd)(s)_{i_0 \ldots i_p}
& =
\sum\nolimits_{j = 0}^p
(-1)^j
h(s)_{i_0 \ldots \hat i_j \ldots i_p}
+
d(s)_{i i_0 \ldots i_p}\\
& =
\sum\nolimits_{j = 0}^p
(-1)^j
s_{i i_0 \ldots \hat i_j \ldots i_p}
+
s_{i_0 \ldots i_p}
+
\sum\nolimits_{j = 0}^p
(-1)^{j + 1}
s_{i i_0 \ldots \hat i_j \ldots i_p} \\
& =
s_{i_0 \ldots i_p}
\end{align*}
这证明恒等映射同伦于零，如所需。
\end{proof}






\section{{\v C}ech 上同调作为预层上的函子}
\label{cohomology-section-cech-functor}

\noindent
警告：本节几乎全程处理预层与预层范畴；
在层与层范畴的情形，这些结论完全错误！

\medskip\noindent
设 $X$ 是环化空间，
$\mathcal{U} : U = \bigcup_{i \in I} U_i$ 是开覆盖，
$\mathcal{F}$ 是 $\mathcal{O}_X$-模预层。
只需将 $\mathcal{F}$ 看作阿贝尔群预层，便得到它的 {\v C}ech 复形
$\check{\mathcal{C}}^\bullet(\mathcal{U}, \mathcal{F})$
。不过，每一项
$\check{\mathcal{C}}^p(\mathcal{U}, \mathcal{F})$
都自然具有 $\mathcal{O}_X(U)$-模结构，且微分由
$\mathcal{O}_X(U)$-模映射给出。此外，构造
$$
\mathcal{F} \longmapsto \check{\mathcal{C}}^\bullet(\mathcal{U}, \mathcal{F})
$$
显然关于 $\mathcal{F}$ 是函子性的。事实上，它是函子
\begin{equation}
\label{cohomology-equation-cech-functor}
\check{\mathcal{C}}^\bullet(\mathcal{U}, -) :
\textit{PMod}(\mathcal{O}_X)
\longrightarrow
\text{Comp}^{+}(\text{Mod}_{\mathcal{O}_X(U)})
\end{equation}
；记号见
《导出范畴》定义 \ref{derived-definition-complexes-notation}。
回忆阿贝尔范畴中下有界复形所成的范畴仍是阿贝尔范畴，见
《同调代数》引理 \ref{homology-lemma-cat-cochain-abelian}。

\begin{lemma}
\label{cohomology-lemma-cech-exact-presheaves}
Equation (\ref{cohomology-equation-cech-functor}) 给出的函子是正合函子
（见 《同调代数》引理 \ref{homology-lemma-exact-functor}）。
\end{lemma}

\begin{proof}
对任意开集 $W \subset U$，函子
$\mathcal{F} \mapsto \mathcal{F}(W)$
是从 $\textit{PMod}(\mathcal{O}_X)$ 到
$\text{Mod}_{\mathcal{O}_X(U)}$ 的加性正合函子。
复形的各项
$\check{\mathcal{C}}^p(\mathcal{U}, \mathcal{F})$
是这些正合函子的积，因而正合。
此外，复形序列正合当且仅当它在每个给定次数的项所成序列正合。
故引理成立。
\end{proof}

\begin{lemma}
\label{cohomology-lemma-cech-cohomology-delta-functor-presheaves}
设 $X$ 是环化空间，
$\mathcal{U} : U = \bigcup_{i \in I} U_i$ 是开覆盖。
诸函子 $\mathcal{F} \mapsto \check{H}^n(\mathcal{U}, \mathcal{F})$
组成从 $\mathcal{O}_X$-模预层的阿贝尔范畴到
$\mathcal{O}_X(U)$-模范畴的 $\delta$-函子
（见 《同调代数》定义 \ref{homology-definition-cohomological-delta-functor}）。
\end{lemma}

\begin{proof}
由引理 \ref{cohomology-lemma-cech-exact-presheaves}，
$\mathcal{O}_X$-模预层的短正合序列
$0 \to \mathcal{F}_1 \to \mathcal{F}_2 \to \mathcal{F}_3 \to 0$
被送到 $\mathcal{O}_X(U)$-模复形的短正合序列。因此可用
《同调代数》引理 \ref{homology-lemma-long-exact-sequence-cochain}
得到连接映射
$\delta_{\mathcal{F}_1 \to \mathcal{F}_2 \to \mathcal{F}_3} :
\check{H}^n(\mathcal{U}, \mathcal{F}_3) \to
\check{H}^{n + 1}(\mathcal{U}, \mathcal{F}_1)$
及相应的长正合序列。
这些映射与预层短正合序列之间的映射相容，其验证从略。
\end{proof}


\noindent
在下面引理的表述中，我们使用相对于开浸入
$j : U \to X$ 的模预层延零函子 $j_{p!}$。
见 《层》第 \ref{sheaves-section-open-immersions} 节。
对任意开集 $W \subset X$ 及任意
$\mathcal{O}_X|_U$-模预层 $\mathcal{G}$，有
$$
(j_{p!}\mathcal{G})(W) =
\left\{
\begin{matrix}
\mathcal{G}(W) & \text{若 } W \subset U \\
0 & \text{else.}
\end{matrix}
\right.
$$
此外，函子 $j_{p!}$ 是限制函子的左伴随，见
《层》引理 \ref{sheaves-lemma-j-shriek-modules}。
特别地，有公式
$$
\Hom_{\mathcal{O}_X}(j_{p!}\mathcal{O}_U, \mathcal{F})
=
\Hom_{\mathcal{O}_U}(\mathcal{O}_U, \mathcal{F}|_U)
=
\mathcal{F}(U).
$$
由于函子 $\mathcal{F} \mapsto \mathcal{F}(U)$
在预层范畴上正合，预层 $j_{p!}\mathcal{O}_U$
是范畴 $\textit{PMod}(\mathcal{O}_X)$ 中的投射对象，见
《同调代数》引理 \ref{homology-lemma-characterize-projectives}。

\medskip\noindent
注意，若给定开子集 $U \subset V \subset X$
及相应的开浸入 $j_U,j_V$，便有典范映射
$(j_U)_{p!}\mathcal{O}_U \to (j_V)_{p!}\mathcal{O}_V$。
它在任意开集 $W \subset U$ 上的截面上为恒等映射，在其他情形为 $0$。
在恒等式
$\Hom_{\mathcal{O}_X}((j_U)_{p!}\mathcal{O}_U, (j_V)_{p!}\mathcal{O}_V) =
(j_V)_{p!}\mathcal{O}_V(U) = \mathcal{O}_V(U)$
之下，它对应于元素 $1 \in \mathcal{O}_V(U)$。

\begin{lemma}
\label{cohomology-lemma-cech-map-into}
设 $X$ 是环化空间，
$\mathcal{U} : U = \bigcup_{i \in I} U_i$ 是开覆盖。
以 $j_{i_0\ldots i_p} : U_{i_0 \ldots i_p} \to X$
表示开浸入。考虑 $\mathcal{O}_X$-模预层的链复形
$K(\mathcal{U})_\bullet$
$$
\ldots
\to
\bigoplus_{i_0i_1i_2} (j_{i_0i_1i_2})_{p!}\mathcal{O}_{U_{i_0i_1i_2}}
\to
\bigoplus_{i_0i_1} (j_{i_0i_1})_{p!}\mathcal{O}_{U_{i_0i_1}}
\to
\bigoplus_{i_0} (j_{i_0})_{p!}\mathcal{O}_{U_{i_0}}
\to 0 \to \ldots
$$
，其中最后一个非零项置于次数 $0$，而映射
$$
(j_{i_0\ldots i_{p + 1}})_{p!}\mathcal{O}_{U_{i_0\ldots i_{p + 1}}}
\longrightarrow
(j_{i_0\ldots \hat i_j \ldots i_{p + 1}})_{p!}
\mathcal{O}_{U_{i_0\ldots \hat i_j \ldots i_{p + 1}}}
$$
由典范映射的 $(-1)^j$ 倍给出。则存在同构
$$
\Hom_{\mathcal{O}_X}(K(\mathcal{U})_\bullet, \mathcal{F})
=
\check{\mathcal{C}}^\bullet(\mathcal{U}, \mathcal{F})
$$
且它关于
$\mathcal{F} \in \Ob(\textit{PMod}(\mathcal{O}_X))$ 是函子性的。
\end{lemma}

\begin{proof}
在引理之前的讨论中已经看到
$$
\Hom_{\mathcal{O}_X}(
(j_{i_0\ldots i_p})_{p!}\mathcal{O}_{U_{i_0\ldots i_p}},
\mathcal{F})
=
\mathcal{F}(U_{i_0\ldots i_p}).
$$
因此，直和
$$
\bigoplus\nolimits_{i_0 \ldots i_p}
(j_{i_0 \ldots i_p})_{p!}\mathcal{O}_{U_{i_0 \ldots i_p}}
$$
确实表示函子
$$
\mathcal{F}
\longmapsto
\prod\nolimits_{i_0\ldots i_p} \mathcal{F}(U_{i_0\ldots i_p}).
$$
由 《范畴》Yoneda 引理 \ref{categories-lemma-yoneda}，
可知存在具有上述各项的复形 $K(\mathcal{U})_\bullet$。
直接可见其映射正如引理所述。
\end{proof}

\begin{lemma}
\label{cohomology-lemma-homology-complex}
设 $X$ 是环化空间，
$\mathcal{U} : U = \bigcup_{i \in I} U_i$ 是开覆盖。
令 $\mathcal{O}_\mathcal{U} \subset \mathcal{O}_X$
为映射 $\bigoplus j_{p!}\mathcal{O}_{U_i} \to \mathcal{O}_X$
的像预层。上面引理 \ref{cohomology-lemma-cech-map-into}
中的预层链复形 $K(\mathcal{U})_\bullet$ 具有同调预层
$$
H_i(K(\mathcal{U})_\bullet) =
\left\{
\begin{matrix}
0 & \text{若} & i \not = 0 \\
\mathcal{O}_\mathcal{U} & \text{若} & i = 0
\end{matrix}
\right.
$$
\end{lemma}

\begin{proof}
考虑如下增广复形 $K^{ext}_\bullet$：将
$\mathcal{O}_\mathcal{U}$ 置于次数 $-1$，并取显然的映射
$K(\mathcal{U})_0 =
\bigoplus_{i_0} (j_{i_0})_{p!}\mathcal{O}_{U_{i_0}} \to
\mathcal{O}_\mathcal{U}$.
只需证明，对任意开集 $W \subset X$，取该增广复形在 $W$ 上的截面
所得复形无环。事实上，我们断言对每个 $W \subset X$，
复形 $K^{ext}_\bullet(W)$ 链同伦等价于零复形。
写作 $I = I_1 \amalg I_2$，其中
$W \subset U_i$ 当且仅当 $i \in I_1$。

\medskip\noindent
若 $I_1 = \emptyset$，则复形 $K^{ext}_\bullet(W) = 0$，无需证明。

\medskip\noindent
若 $I_1 \not = \emptyset$，则
$\mathcal{O}_\mathcal{U}(W) = \mathcal{O}_X(W)$，并且
$$
K^{ext}_p(W) =
\bigoplus\nolimits_{i_0 \ldots i_p \in I_1} \mathcal{O}_X(W).
$$
这是因为预层
$(j_{i_0 \ldots i_p})_{p!}\mathcal{O}_{U_{i_0 \ldots i_p}}$
有上述简单描述。此外，复形 $K^{ext}_\bullet(W)$ 的微分为
$$
d(s)_{i_0 \ldots i_p} =
\sum\nolimits_{j = 0, \ldots, p + 1} \sum\nolimits_{i \in I_1}
(-1)^j s_{i_0 \ldots i_{j - 1} i i_j \ldots i_p}.
$$
由于元素 $s$ 具有有限支撑，该和是有限的。
固定元素 $i_{\text{fix}} \in I_1$。定义映射
$$
h : K^{ext}_p(W) \longrightarrow K^{ext}_{p + 1}(W)
$$
如下：
$$
h(s)_{i_0 \ldots i_{p + 1}} =
\left\{
\begin{matrix}
0 & \text{若} & i_0 \not = i_{\text{fix}} \\
s_{i_1 \ldots i_{p + 1}} & \text{若} & i_0 = i_{\text{fix}}
\end{matrix}
\right.
$$
我们将使用简记
$h(s)_{i_0 \ldots i_{p + 1}} = (i_0 = i_{\text{fix}}) s_{i_1 \ldots i_p}$
。于是计算得
\begin{eqnarray*}
& & (dh + hd)(s)_{i_0 \ldots i_p} \\
& = &
\sum_j \sum_{i \in I_1} (-1)^j h(s)_{i_0 \ldots i_{j - 1} i i_j \ldots i_p}
+
(i = i_0) d(s)_{i_1 \ldots i_p} \\
& = &
s_{i_0 \ldots i_p} +
\sum_{j \geq 1}\sum_{i \in I_1}
(-1)^j (i_0 = i_{\text{fix}}) s_{i_1 \ldots i_{j - 1} i i_j \ldots i_p}
+
(i_0 = i_{\text{fix}}) d(s)_{i_1 \ldots i_p}
\end{eqnarray*}
它等于 $s_{i_0 \ldots i_p}$，如所需。
\end{proof}

\begin{lemma}
\label{cohomology-lemma-cech-cohomology-derived-presheaves}
设 $X$ 是环化空间，
$\mathcal{U} : U = \bigcup_{i \in I} U_i$
是 $U \subset X$ 的开覆盖。
{\v C}ech 上同调函子 $\check{H}^p(\mathcal{U}, -)$
作为 $\delta$-函子，典范同构于函子
$$
\check{H}^0(\mathcal{U}, -) :
\textit{PMod}(\mathcal{O}_X)
\longrightarrow
\text{Mod}_{\mathcal{O}_X(U)}.
$$
的右导出函子。此外，存在函子性拟同构
$$
\check{\mathcal{C}}^\bullet(\mathcal{U}, \mathcal{F})
\longrightarrow
R\check{H}^0(\mathcal{U}, \mathcal{F})
$$
，其中右端表示左正合函子
$$
R\check{H}^0(\mathcal{U}, -) :
D^{+}(\textit{PMod}(\mathcal{O}_X))
\longrightarrow
D^{+}(\mathcal{O}_X(U))
$$
$\check{H}^0(\mathcal{U}, -)$ 的右导出函子。
\end{lemma}

\begin{proof}
注意，$\mathcal{O}_X$-模预层的范畴有足够多的内射对象，见
《内射对象》Proposition \ref{injectives-proposition-presheaves-modules}。
又注意到 $\check{H}^0(\mathcal{U}, -)$
是从 $\mathcal{O}_X$-模预层范畴到
$\mathcal{O}_X(U)$-模范畴的左正合函子。
因此导出函子和右导出函子存在，见
《导出范畴》第 \ref{derived-section-right-derived-functor} 节。

\medskip\noindent
设 $\mathcal{I}$ 是内射 $\mathcal{O}_X$-模预层。
此时函子 $\Hom_{\mathcal{O}_X}(-, \mathcal{I})$
在 $\textit{PMod}(\mathcal{O}_X)$ 上正合。由引理 \ref{cohomology-lemma-cech-map-into}，有
$$
\Hom_{\mathcal{O}_X}(K(\mathcal{U})_\bullet, \mathcal{I})
=
\check{\mathcal{C}}^\bullet(\mathcal{U}, \mathcal{I}).
$$
由引理 \ref{cohomology-lemma-homology-complex}，
$K(\mathcal{U})_\bullet$ 拟同构于
$\mathcal{O}_\mathcal{U}[0]$。因此，由上述取到
$\mathcal{I}$ 的 Hom 函子的正合性可知，对所有 $i > 0$，
$\check{H}^i(\mathcal{U}, \mathcal{I}) = 0$。
故 $\delta$-函子 $(\check{H}^n, \delta)$
（见引理 \ref{cohomology-lemma-cech-cohomology-delta-functor-presheaves}）
满足
《同调代数》引理 \ref{homology-lemma-efface-implies-universal},
的假设，因而是泛 $\delta$-函子。

\medskip\noindent
由 《导出范畴》引理 \ref{derived-lemma-higher-derived-functors}，
序列 $R^i\check{H}^0(\mathcal{U}, -)$ 也组成泛 $\delta$-函子。
由泛 $\delta$-函子的唯一性，见
《同调代数》引理 \ref{homology-lemma-uniqueness-universal-delta-functor}
，得到
$R^i\check{H}^0(\mathcal{U}, -) = \check{H}^i(\mathcal{U}, -)$.
这对大多数应用已经足够，建议读者略过证明的其余部分。

\medskip\noindent
设 $\mathcal{F}$ 是任意 $\mathcal{O}_X$-模预层。
在范畴 $\textit{PMod}(\mathcal{O}_X)$ 中取内射消解
$\mathcal{F} \to \mathcal{I}^\bullet$。
考虑各项为
$\check{\mathcal{C}}^p(\mathcal{U}, \mathcal{I}^q)$
的双复形
$\check{\mathcal{C}}^\bullet(\mathcal{U}, \mathcal{I}^\bullet)$，
以及相伴的全复形
$\text{Tot}(\check{\mathcal{C}}^\bullet(\mathcal{U}, \mathcal{I}^\bullet))$,
见 《同调代数》定义 \ref{homology-definition-associated-simple-complex}。
由诸映射
$\check{\mathcal{C}}^p(\mathcal{U}, \mathcal{F})
\to \check{\mathcal{C}}^p(\mathcal{U}, \mathcal{I}^0)$
得到复形映射
$$
\check{\mathcal{C}}^\bullet(\mathcal{U}, \mathcal{F})
\longrightarrow
\text{Tot}(\check{\mathcal{C}}^\bullet(\mathcal{U}, \mathcal{I}^\bullet))
$$
，并且由诸映射
$\check{H}^0(\mathcal{U}, \mathcal{I}^q) \to
\check{\mathcal{C}}^0(\mathcal{U}, \mathcal{I}^q)$
得到复形映射
$$
\check{H}^0(\mathcal{U}, \mathcal{I}^\bullet)
\longrightarrow
\text{Tot}(\check{\mathcal{C}}^\bullet(\mathcal{U}, \mathcal{I}^\bullet))
$$
。施用 《同调代数》引理 \ref{homology-lemma-double-complex-gives-resolution}
可知，这两个映射都是拟同构。
事实上，由于作为函子的 {\v C}ech 复形是正合的
（引理 \ref{cohomology-lemma-cech-exact-presheaves}），
该双复形的各列在正次数处正合；
又由于刚才看到内射预层 $\mathcal{I}^q$
的高阶 {\v C}ech 上同调群为零，其各行在正次数处也正合。
拟同构在 $D^{+}(\mathcal{O}_X(U))$ 中可逆，
故得到引理中最后一个显示态射。
该态射的函子性验证从略。
\end{proof}





\section{{\v C}ech 上同调与层上同调}
\label{cohomology-section-cech-cohomology-cohomology}

\begin{lemma}
\label{cohomology-lemma-injective-trivial-cech}
设 $X$ 是环化空间，
$\mathcal{U} : U = \bigcup_{i \in I} U_i$ 是开覆盖，
$\mathcal{I}$ 是内射 $\mathcal{O}_X$-模。则
$$
\check{H}^p(\mathcal{U}, \mathcal{I}) =
\left\{
\begin{matrix}
\mathcal{I}(U) & \text{若} & p = 0 \\
0 & \text{若} & p > 0
\end{matrix}
\right.
$$
\end{lemma}

\begin{proof}
内射 $\mathcal{O}_X$-模作为范畴
$\textit{PMod}(\mathcal{O}_X)$ 中的对象仍是内射的
（例如，这是因为层化函子是包含函子的正合左伴随；使用
《同调代数》引理 \ref{homology-lemma-adjoint-preserve-injectives}）。
因此施用引理 \ref{cohomology-lemma-cech-cohomology-derived-presheaves}
（或其证明）即得结论。
\end{proof}

\begin{lemma}
\label{cohomology-lemma-cech-cohomology}
设 $X$ 是环化空间，
$\mathcal{U} : U = \bigcup_{i \in I} U_i$ 是开覆盖。
存在函子
$\textit{Mod}(\mathcal{O}_X) \to D^{+}(\mathcal{O}_X(U))$
之间的变换
$$
\check{\mathcal{C}}^\bullet(\mathcal{U}, -)
\longrightarrow
R\Gamma(U, -)
$$
。特别地，对遍历
$\textit{Mod}(\mathcal{O}_X)$ 的 $\mathcal{F}$，
它给出典范映射
$\check{H}^p(\mathcal{U}, \mathcal{F}) \to H^p(U, \mathcal{F})$。
\end{lemma}

\begin{proof}
设 $\mathcal{F}$ 是 $\mathcal{O}_X$-模，并取内射消解
$\mathcal{F} \to \mathcal{I}^\bullet$。
考虑各项为
$\check{\mathcal{C}}^p(\mathcal{U}, \mathcal{I}^q)$
的双复形
$\check{\mathcal{C}}^\bullet(\mathcal{U}, \mathcal{I}^\bullet)$。
由诸映射
$\mathcal{I}^q(U) \to \check{H}^0(\mathcal{U}, \mathcal{I}^q)$
得到复形映射
$$
\alpha :
\Gamma(U, \mathcal{I}^\bullet)
\longrightarrow
\text{Tot}(\check{\mathcal{C}}^\bullet(\mathcal{U}, \mathcal{I}^\bullet))
$$
，并由映射 $\mathcal{F} \to \mathcal{I}^0$
得到复形映射
$$
\beta :
\check{\mathcal{C}}^\bullet(\mathcal{U}, \mathcal{F})
\longrightarrow
\text{Tot}(\check{\mathcal{C}}^\bullet(\mathcal{U}, \mathcal{I}^\bullet))
$$
。施用
《同调代数》引理 \ref{homology-lemma-double-complex-gives-resolution}
可知 $\alpha$ 是拟同构。
事实上，引理 \ref{cohomology-lemma-injective-trivial-cech} 说明，
双复形
$\check{\mathcal{C}}^\bullet(\mathcal{U}, \mathcal{I}^\bullet)$
的第 $q$ 行是 $\Gamma(U, \mathcal{I}^q)$ 的消解。
故 $\alpha$ 在 $D^{+}(\mathcal{O}_X(U))$ 中可逆，
而引理中的变换是先施用 $\beta$ 再施用 $\alpha$ 的逆所得的复合。
其函子性验证从略。
\end{proof}

\begin{lemma}
\label{cohomology-lemma-cech-h1}
设 $X$ 是拓扑空间，$\mathcal{H}$ 是 $X$ 上的阿贝尔层，
$\mathcal{U} : X = \bigcup_{i \in I} U_i$ 是开覆盖。映射
$$
\check{H}^1(\mathcal{U}, \mathcal{H}) \longrightarrow H^1(X, \mathcal{H})
$$
是单射，并通过引理 \ref{cohomology-lemma-torsors-h1} 的双射，
将 $\check{H}^1(\mathcal{U}, \mathcal{H})$
等同于如下 $\mathcal{H}$-挠子的同构类集合：
这些挠子限制到每个 $U_i$ 上均为平凡挠子。
\end{lemma}

\begin{proof}
为此我们构造逆映射。设 $\mathcal{F}$ 是一个
$\mathcal{H}$-挠子，其在 $U_i$ 上的限制平凡。
由引理 \ref{cohomology-lemma-trivial-torsor}，这意味着存在截面
$s_i \in \mathcal{F}(U_i)$。在 $U_{i_0} \cap U_{i_1}$ 上，
存在唯一的 $\mathcal{H}$ 的截面 $s_{i_0i_1}$，使得
$s_{i_0i_1} \cdot s_{i_0}|_{U_{i_0} \cap U_{i_1}} =
s_{i_1}|_{U_{i_0} \cap U_{i_1}}$。
直接计算可见，$s_{i_0i_1}$ 是 {\v C}ech 上循环，
且其类是良定义的（即不依赖于截面 $s_i$ 的选取）。
逆映射把 $\mathcal{F}$ 的同构类送到上循环
$(s_{i_0i_1})$ 的上同调类。
该映射确为逆映射的验证从略。
\end{proof}

\begin{lemma}
\label{cohomology-lemma-include}
设 $X$ 是环化空间。考虑函子
$i : \textit{Mod}(\mathcal{O}_X) \to \textit{PMod}(\mathcal{O}_X)$.
这是左正合函子，其右导出函子为
$$
R^pi(\mathcal{F}) = \underline{H}^p(\mathcal{F}) :
U \longmapsto H^p(U, \mathcal{F})
$$
，见 第 \ref{cohomology-section-locality} 节 中的讨论。
\end{lemma}

\begin{proof}
显然 $i$ 左正合。取内射消解
$\mathcal{F} \to \mathcal{I}^\bullet$。
按定义，$R^pi$ 是复形 $\mathcal{I}^\bullet$ 的第 $p$ 个
上同调{\it 预层}。换言之，
$R^pi(\mathcal{F})$ 在开集 $U$ 上的截面为
$$
\frac{\Ker(\mathcal{I}^p(U) \to \mathcal{I}^{p + 1}(U))}
{\Im(\mathcal{I}^{p - 1}(U) \to \mathcal{I}^p(U))}.
$$
这正是 $H^p(U, \mathcal{F})$ 的定义。
\end{proof}

\begin{lemma}
\label{cohomology-lemma-cech-spectral-sequence}
设 $X$ 是环化空间，
$\mathcal{U} : U = \bigcup_{i \in I} U_i$ 是开覆盖。
对任意 $\mathcal{O}_X$-模层 $\mathcal{F}$，
存在谱序列 $(E_r, d_r)_{r \geq 0}$，满足
$$
E_2^{p, q} = \check{H}^p(\mathcal{U}, \underline{H}^q(\mathcal{F}))
$$
并收敛到 $H^{p + q}(U, \mathcal{F})$。
该谱序列关于 $\mathcal{F}$ 是函子性的。
\end{lemma}

\begin{proof}
这是函子
$$
i :  \textit{Mod}(\mathcal{O}_X) \to \textit{PMod}(\mathcal{O}_X)
\quad\text{且}\quad
\check{H}^0(\mathcal{U}, - ) : \textit{PMod}(\mathcal{O}_X)
\to \text{Mod}_{\mathcal{O}_X(U)}.
$$
的 Grothendieck 谱序列（见
《导出范畴》引理 \ref{derived-lemma-grothendieck-spectral-sequence})
）。事实上，由引理 \ref{cohomology-lemma-cech-h0}，
$\check{H}^0(\mathcal{U}, i(\mathcal{F})) = \mathcal{F}(U)$。
由引理 \ref{cohomology-lemma-injective-trivial-cech}，
$i(\mathcal{I})$ 是 {\v C}ech 无环的。
又由引理 \ref{cohomology-lemma-cech-cohomology-derived-presheaves}，
在 $\textit{PMod}(\mathcal{O}_X)$ 上有函子等式
$\check{H}^p(\mathcal{U}, -) = R^p\check{H}^0(\mathcal{U}, -)$。
综合这些事实即得引理。
\end{proof}

\begin{lemma}
\label{cohomology-lemma-cech-spectral-sequence-application}
设 $X$ 是环化空间，
$\mathcal{U} : U = \bigcup_{i \in I} U_i$ 是开覆盖，
$\mathcal{F}$ 是 $\mathcal{O}_X$-模。
假设对所有 $i > 0$、所有 $p \geq 0$
以及所有 $i_0, \ldots, i_p \in I$，均有
$H^i(U_{i_0 \ldots i_p}, \mathcal{F}) = 0$。
则作为 $\mathcal{O}_X(U)$-模，有
$\check{H}^p(\mathcal{U}, \mathcal{F}) = H^p(U, \mathcal{F})$。
\end{lemma}

\begin{proof}
使用引理 \ref{cohomology-lemma-cech-spectral-sequence} 的谱序列。
这些假设意味着，对所有满足 $q \not = 0$ 的 $(p,q)$，
均有 $E_2^{p, q} = 0$。因此谱序列在 $E_2$ 页退化，结论随之成立。
\end{proof}

\begin{lemma}
\label{cohomology-lemma-ses-cech-h1}
设 $X$ 是环化空间。设
$$
0 \to \mathcal{F} \to \mathcal{G} \to \mathcal{H} \to 0
$$
是 $\mathcal{O}_X$-模的短正合序列，
$U \subset X$ 是开子集。
若存在 $U$ 的开覆盖 $\mathcal{U}$ 所成的共尾系，
使得 $\check{H}^1(\mathcal{U}, \mathcal{F}) = 0$，
则映射 $\mathcal{G}(U) \to \mathcal{H}(U)$ 是满射。
\end{lemma}

\begin{proof}
取元素 $s \in \mathcal{H}(U)$。选择开覆盖
$\mathcal{U} : U = \bigcup_{i \in I} U_i$，使得
(a) $\check{H}^1(\mathcal{U}, \mathcal{F}) = 0$，且
(b) $s|_{U_i}$ 是某截面 $s_i \in \mathcal{G}(U_i)$ 的像。
显然能找到满足 (b) 的 $\mathcal{U}$；
由引理的假设，可找到同时满足 (a) 与 (b) 的 $\mathcal{U}$。
考虑截面
$$
s_{i_0i_1} = s_{i_1}|_{U_{i_0i_1}} - s_{i_0}|_{U_{i_0i_1}}.
$$
由于 $s_i$ 提升 $s$，可知
$s_{i_0i_1} \in \mathcal{F}(U_{i_0i_1})$。
由 $\check{H}^1(\mathcal{U}, \mathcal{F})$ 的消失性，
可找到截面 $t_i \in \mathcal{F}(U_i)$，使得
$$
s_{i_0i_1} = t_{i_1}|_{U_{i_0i_1}} - t_{i_0}|_{U_{i_0i_1}}.
$$
于是截面 $s_i-t_i$ 显然满足层条件，因而黏合成
$\mathcal{G}$ 在 $U$ 上的截面，且它映到 $s$。结论成立。
\end{proof}

\begin{lemma}
\label{cohomology-lemma-cech-vanish}
\begin{slogan}
若阿贝尔层对所有开覆盖的高阶 {\v C}ech 上同调均消失，
则其高阶层上同调消失。
\end{slogan}
设 $X$ 是环化空间，$\mathcal{F}$ 是 $\mathcal{O}_X$-模，且
$$
\check{H}^p(\mathcal{U}, \mathcal{F}) = 0
$$
对所有 $p > 0$ 以及 $X$ 中任意开集的任意开覆盖
$\mathcal{U} : U = \bigcup_{i \in I} U_i$ 均成立。
则对所有 $p > 0$ 及任意开集 $U \subset X$，有
$H^p(U, \mathcal{F}) = 0$。
\end{lemma}

\begin{proof}
设 $\mathcal{F}$ 是满足引理假设的层。
我们称之为“$\mathcal{F}$ 对任意开覆盖的高阶
{\v C}ech 上同调消失”。
选取到内射 $\mathcal{O}_X$-模的嵌入
$\mathcal{F} \to \mathcal{I}$。
由引理 \ref{cohomology-lemma-injective-trivial-cech}，
$\mathcal{I}$ 对任意开覆盖的高阶 {\v C}ech 上同调消失。
令 $\mathcal{Q} = \mathcal{I}/\mathcal{F}$，于是有短正合序列
$$
0 \to \mathcal{F} \to \mathcal{I} \to \mathcal{Q} \to 0.
$$
由引理 \ref{cohomology-lemma-ses-cech-h1} 及我们的假设，
该序列事实上作为预层序列也是正合的！
特别地，对任意开覆盖 $\mathcal{U}$，
都有 {\v C}ech 上同调群的长正合序列，例如见引理 \ref{cohomology-lemma-cech-cohomology-delta-functor-presheaves}
。这说明 $\mathcal{Q}$ 也是对所有开覆盖具有消失高阶
{\v C}ech 上同调的 $\mathcal{O}_X$-模。

\medskip\noindent
接着，对任意开集 $U \subset X$，考虑上同调长正合序列
$$
\xymatrix{
0 \ar[r] &
H^0(U, \mathcal{F}) \ar[r] &
H^0(U, \mathcal{I}) \ar[r] &
H^0(U, \mathcal{Q}) \ar[lld] \\
&
H^1(U, \mathcal{F}) \ar[r] &
H^1(U, \mathcal{I}) \ar[r] &
H^1(U, \mathcal{Q}) \ar[lld] \\
&
\ldots & \ldots & \ldots \\
}
$$
。由于 $\mathcal{I}$ 内射，对 $n > 0$ 有
$H^n(U, \mathcal{I}) = 0$（见
《导出范畴》引理 \ref{derived-lemma-higher-derived-functors}）。
由上可知，$H^0(U, \mathcal{I}) \to H^0(U, \mathcal{Q})$
是满射，因而 $H^1(U, \mathcal{F}) = 0$。
由于 $\mathcal{F}$ 是具有消失高阶 {\v C}ech 上同调的任意
$\mathcal{O}_X$-模，而 $\mathcal{Q}$ 也是这类层（见上），
同样得到 $H^1(U, \mathcal{Q}) = 0$。
长正合序列进而给出 $H^2(U, \mathcal{F}) = 0$，如此继续即可。
\end{proof}

\begin{lemma}
\label{cohomology-lemma-cech-vanish-basis}
（引理 \ref{cohomology-lemma-cech-vanish} 的变体。）
设 $X$ 是环化空间，$\mathcal{B}$ 是 $X$ 的拓扑基，
$\mathcal{F}$ 是 $\mathcal{O}_X$-模。
假设存在开覆盖的集合 $\text{Cov}$，满足下列性质：
\begin{enumerate}
\item 对每个 $\mathcal{U} \in \text{Cov}$，若
$\mathcal{U} : U = \bigcup_{i \in I} U_i$，则
$U,U_i \in \mathcal{B}$，且每个
$U_{i_0 \ldots i_p} \in \mathcal{B}$。
\item 对每个 $U \in \mathcal{B}$，$\text{Cov}$ 中出现的
$U$ 的开覆盖构成 $U$ 的开覆盖的共尾系。
\item 对每个 $\mathcal{U} \in \text{Cov}$ 及所有 $p > 0$，有
$\check{H}^p(\mathcal{U}, \mathcal{F}) = 0$。
\end{enumerate}
则对所有 $p > 0$ 及任意 $U \in \mathcal{B}$，有
$H^p(U, \mathcal{F}) = 0$。
\end{lemma}

\begin{proof}
设 $\mathcal{F}$ 与 $\text{Cov}$ 如引理所述。
我们称之为“$\mathcal{F}$ 对任意
$\mathcal{U} \in \text{Cov}$ 的高阶 {\v C}ech 上同调消失”。
选取到内射 $\mathcal{O}_X$-模的嵌入
$\mathcal{F} \to \mathcal{I}$。
由引理 \ref{cohomology-lemma-injective-trivial-cech}，
$\mathcal{I}$ 对任意 $\mathcal{U} \in \text{Cov}$
的高阶 {\v C}ech 上同调消失。
令 $\mathcal{Q} = \mathcal{I}/\mathcal{F}$，于是有短正合序列
$$
0 \to \mathcal{F} \to \mathcal{I} \to \mathcal{Q} \to 0.
$$
由引理 \ref{cohomology-lemma-ses-cech-h1} 及假设 (2)，
对每个 $U \in \mathcal{B}$，该序列给出正合序列
$$
0 \to \mathcal{F}(U) \to \mathcal{I}(U) \to \mathcal{Q}(U) \to 0.
$$
因此，对任意 $\mathcal{U} \in \text{Cov}$，
得到 {\v C}ech 复形的短正合序列
$$
0 \to
\check{\mathcal{C}}^\bullet(\mathcal{U}, \mathcal{F}) \to
\check{\mathcal{C}}^\bullet(\mathcal{U}, \mathcal{I}) \to
\check{\mathcal{C}}^\bullet(\mathcal{U}, \mathcal{Q}) \to 0
$$
，因为由假设 (1)，{\v C}ech 复形的每一项都是
$\mathcal{B}$ 中元素上取值的积。
特别地，对任意开覆盖 $\mathcal{U} \in \text{Cov}$，
都有 {\v C}ech 上同调群的长正合序列。
这说明 $\mathcal{Q}$ 也是对所有
$\mathcal{U} \in \text{Cov}$ 具有消失高阶 {\v C}ech 上同调的
$\mathcal{O}_X$-模。

\medskip\noindent
接着，对任意 $U \in \mathcal{B}$，考虑上同调长正合序列
$$
\xymatrix{
0 \ar[r] &
H^0(U, \mathcal{F}) \ar[r] &
H^0(U, \mathcal{I}) \ar[r] &
H^0(U, \mathcal{Q}) \ar[lld] \\
&
H^1(U, \mathcal{F}) \ar[r] &
H^1(U, \mathcal{I}) \ar[r] &
H^1(U, \mathcal{Q}) \ar[lld] \\
&
\ldots & \ldots & \ldots \\
}
$$
。由于 $\mathcal{I}$ 内射，对 $n > 0$ 有
$H^n(U, \mathcal{I}) = 0$（见
《导出范畴》引理 \ref{derived-lemma-higher-derived-functors}）。
由上可知，$H^0(U, \mathcal{I}) \to H^0(U, \mathcal{Q})$
是满射，因而 $H^1(U, \mathcal{F}) = 0$。
由于 $\mathcal{F}$ 是对所有
$\mathcal{U} \in \text{Cov}$ 具有消失高阶 {\v C}ech 上同调的任意
$\mathcal{O}_X$-模，而 $\mathcal{Q}$ 也是这类层（见上），
同样得到 $H^1(U, \mathcal{Q}) = 0$。
长正合序列进而给出 $H^2(U, \mathcal{F}) = 0$，如此继续即可。
\end{proof}

\begin{lemma}
\label{cohomology-lemma-pushforward-injective}
设 $f : X \to Y$ 是环化空间的态射，
$\mathcal{I}$ 是内射 $\mathcal{O}_X$-模。则
\begin{enumerate}
\item $\check{H}^p(\mathcal{V}, f_*\mathcal{I}) = 0$
对所有 $p > 0$ 及 $Y$ 的任意开覆盖
$\mathcal{V} : V = \bigcup_{j \in J} V_j$ 成立；
\item 对所有 $p > 0$ 及每个开集 $V \subset Y$，
有 $H^p(V, f_*\mathcal{I}) = 0$。
\end{enumerate}
换言之，对任意开集 $V \subset Y$，
$f_*\mathcal{I}$ 关于 $\Gamma(V, -)$ 是右无环的
（见 《导出范畴》定义 \ref{derived-definition-derived-functor}）。
\end{lemma}

\begin{proof}
置
$\mathcal{U} : f^{-1}(V) = \bigcup_{j \in J} f^{-1}(V_j)$。
这是 $X$ 的开覆盖，并且
$$
\check{\mathcal{C}}^\bullet(\mathcal{V}, f_*\mathcal{I}) =
\check{\mathcal{C}}^\bullet(\mathcal{U}, \mathcal{I}).
$$
这是因为
$$
f_*\mathcal{I}(V_{j_0 \ldots j_p})
= \mathcal{I}(f^{-1}(V_{j_0 \ldots j_p})) =
\mathcal{I}(f^{-1}(V_{j_0}) \cap \ldots \cap f^{-1}(V_{j_p}))
= \mathcal{I}(U_{j_0 \ldots j_p}).
$$
因此引理的第一条由引理 \ref{cohomology-lemma-injective-trivial-cech} 推出；
第二条由第一条及引理 \ref{cohomology-lemma-cech-vanish} 推出。
\end{proof}

\noindent
下面的引理特别说明，
$f_* : \textit{Ab}(X) \to \textit{Ab}(Y)$
把内射阿贝尔层送到内射阿贝尔层。

\begin{lemma}
\label{cohomology-lemma-pushforward-injective-flat}
设 $f : X \to Y$ 是环化空间的态射，并假设 $f$ 平坦。
则对任意内射 $\mathcal{O}_X$-模 $\mathcal{I}$，
$f_*\mathcal{I}$ 是内射 $\mathcal{O}_Y$-模。
\end{lemma}

\begin{proof}
此时函子 $f^*$ 把单射送到单射
（《模层》引理 \ref{modules-lemma-pullback-flat}）。
故结论由
《同调代数》引理 \ref{homology-lemma-adjoint-preserve-injectives}
推出。
\end{proof}

\begin{lemma}
\label{cohomology-lemma-cohomology-products}
设 $(X, \mathcal{O}_X)$ 是环化空间，$I$ 是集合。
对 $i \in I$，设 $\mathcal{F}_i$ 是 $\mathcal{O}_X$-模，
并设 $U \subset X$ 是开集。典范映射
$$
H^p(U, \prod\nolimits_{i \in I} \mathcal{F}_i)
\longrightarrow
\prod\nolimits_{i \in I} H^p(U, \mathcal{F}_i)
$$
当 $p = 0$ 时是同构，当 $p = 1$ 时是单射。
\end{lemma}

\begin{proof}
当 $p = 0$ 时的陈述成立，因为层的积等于其底层预层的积，见
《层》第 \ref{sheaves-section-limits-sheaves} 节。
下面证明 $p = 1$ 的情形。置
$\mathcal{F} = \prod \mathcal{F}_i$。
设 $\xi \in H^1(U, \mathcal{F})$ 在
$\prod H^1(U, \mathcal{F}_i)$ 中的像为零。
由上同调的局部性，见引理 \ref{cohomology-lemma-kill-cohomology-class-on-covering}，
存在开覆盖 $\mathcal{U} : U = \bigcup U_j$，
使得对所有 $j$ 有 $\xi|_{U_j} = 0$。
由引理 \ref{cohomology-lemma-cech-h1}，这意味着 $\xi$ 来自某元素
$\check \xi \in \check H^1(\mathcal{U}, \mathcal{F})$。
由于对所有 $i$，映射
$\check H^1(\mathcal{U}, \mathcal{F}_i) \to H^1(U, \mathcal{F}_i)$
均为单射（引理 \ref{cohomology-lemma-cech-h1}），
而 $\xi$ 在 $\prod H^1(U, \mathcal{F}_i)$ 中的像为零，
故其像 $\check \xi_i$ 在
$\check H^1(\mathcal{U}, \mathcal{F}_i)$ 中为零。
另一方面，由于 $\mathcal{F} = \prod \mathcal{F}_i$，
$\check{\mathcal{C}}^\bullet(\mathcal{U}, \mathcal{F})$
是诸复形
$\check{\mathcal{C}}^\bullet(\mathcal{U}, \mathcal{F}_i)$
的积。因此由
《同调代数》引理 \ref{homology-lemma-product-abelian-groups-exact}，
得到 $\check \xi = 0$，如所需。
\end{proof}







\section{松弛层}
\label{cohomology-section-flasque}

\noindent
定义如下。

\begin{definition}
\label{cohomology-definition-flasque}
设 $X$ 是拓扑空间。若对 $X$ 中任意开集
$U \subset V$，限制映射
$\mathcal{F}(V) \to \mathcal{F}(U)$ 都是满射，
则称集合预层 $\mathcal{F}$ 是{\it 松弛的}（亦称 flabby）。
\end{definition}

\noindent
对于阿贝尔层，以及当 $X$ 为环化空间时的模层，
我们也使用这一术语。
显然，只需假设对每个开集 $U \subset X$，
限制映射 $\mathcal{F}(X) \to \mathcal{F}(U)$ 是满射。

\begin{lemma}
\label{cohomology-lemma-injective-flasque}
设 $(X, \mathcal{O}_X)$ 是环化空间。
则任意内射 $\mathcal{O}_X$-模都是松弛的。
\end{lemma}

\begin{proof}
这是引理 \ref{cohomology-lemma-injective-restriction-surjective} 的改述。
\end{proof}

\begin{lemma}
\label{cohomology-lemma-flasque-acyclic}
设 $(X, \mathcal{O}_X)$ 是环化空间。
任意松弛 $\mathcal{O}_X$-模关于 $R\Gamma(X,-)$ 无环，
并且对 $X$ 的任意开集 $U$，关于 $R\Gamma(U,-)$ 也无环。
\end{lemma}

\begin{proof}
我们使用
《导出范畴》引理 \ref{derived-lemma-subcategory-right-acyclics}.
来证明。由于每个内射模都是松弛的，任意 $\mathcal{O}_X$-模
都可嵌入一个松弛模，见
《内射对象》引理 \ref{injectives-lemma-abelian-sheaves-space}。
因此只需证明：给定短正合序列
$$
0 \to \mathcal{F} \to \mathcal{G} \to \mathcal{H} \to 0
$$
，其中 $\mathcal{F}$、$\mathcal{G}$ 松弛，则 $\mathcal{H}$
松弛，且在 $X$ 的任意开集上取截面后，该序列仍短正合。
事实上，后一陈述蕴含前一陈述。
设 $U \subset X$ 是开子空间，$s \in \mathcal{H}(U)$。
我们证明 $s$ 可提升为 $\mathcal{G}$ 在 $U$ 上的截面。
考虑所有二元组 $(V,t)$ 所成的集合 $T$，其中
$V \subset U$ 开，且 $t \in \mathcal{G}(V)$ 是在
$\mathcal{H}$ 中映到 $s|_V$ 的截面。
规定
$(V,t) \leq (V',t')$ 当且仅当 $V \subset V'$ 且 $t'|_V=t$，
从而在 $T$ 上定义偏序。
若 $(V_\alpha,t_\alpha)$（$\alpha\in A$）是 $T$ 的全序子集，
则 $V=\bigcup V_\alpha$ 开；由 $\mathcal{G}$ 的层条件，
存在唯一的 $t\in\mathcal{G}(V)$，其在 $V_\alpha$ 上限制为
$t_\alpha$。故由 Zorn 引理，$T$ 中存在极大元 $(V,t)$。
我们证明 $V=U$，从而结束证明。
任取 $x\in U$。可找到 $x$ 的小开邻域 $W\subset U$
以及 $t'\in\mathcal{G}(W)$，使其在 $\mathcal{H}$ 中映到 $s|_W$。
于是 $t'|_{W\cap V}-t|_{W\cap V}$ 在 $\mathcal{H}$ 中映到零，
故来自某个截面 $r'\in\mathcal{F}(W\cap V)$。
由于 $\mathcal{F}$ 松弛，可找到截面
$r\in\mathcal{F}(W)$，其在 $W\cap V$ 上限制为 $r'$。
用 $r$ 的像修正 $t'$ 后，可假设 $t$ 与 $t'$
在 $W\cap V$ 上限制为同一截面。
由 $\mathcal{G}$ 的层条件，可找到 $\mathcal{G}$
在 $W\cup V$ 上的截面 $\tilde t$，其限制分别为 $t$ 和 $t'$。
由 $(V,t)$ 的极大性，$V\cup W=V$。故 $x\in V$，结论成立。
\end{proof}

\noindent
下面的引理对松弛预层并不成立。

\begin{lemma}
\label{cohomology-lemma-flasque-acyclic-cech}
设 $(X, \mathcal{O}_X)$ 是环化空间，
$\mathcal{F}$ 是 $\mathcal{O}_X$-模层，
$\mathcal{U}:U=\bigcup U_i$ 是开覆盖。
若 $\mathcal{F}$ 松弛，则对 $p>0$，
$\check{H}^p(\mathcal{U},\mathcal{F})=0$。
\end{lemma}

\begin{proof}
引理 \ref{cohomology-lemma-cech-spectral-sequence} 的陈述中所用的预层
$\underline{H}^q(\mathcal{F})$ 由引理 \ref{cohomology-lemma-flasque-acyclic} 可知为零。
再由引理 \ref{cohomology-lemma-flasque-acyclic}，得到
$\check{H}^p(U,\mathcal{F})=H^p(U,\mathcal{F})=0$。
\end{proof}

\begin{lemma}
\label{cohomology-lemma-flasque-acyclic-pushforward}
设 $f:(X,\mathcal{O}_X)\to(Y,\mathcal{O}_Y)$ 是环化空间的态射，
$\mathcal{F}$ 是 $\mathcal{O}_X$-模层。
若 $\mathcal{F}$ 松弛，则对 $p>0$，
$R^pf_*\mathcal{F}=0$。
\end{lemma}

\begin{proof}
这立即由引理 \ref{cohomology-lemma-describe-higher-direct-images}
及引理 \ref{cohomology-lemma-flasque-acyclic} 推出。
\end{proof}

\noindent
对于有限覆盖，下面的引理可用初等归纳法证明；
可与 \cite{FOAG} 中关于 {\v C}ech 上同调的讨论比较。

\begin{lemma}
\label{cohomology-lemma-vanishing-ravi}
设 $X$ 是拓扑空间，$\mathcal{F}$ 是 $X$ 上的阿贝尔层，
$\mathcal{U}:U=\bigcup_{i\in I}U_i$ 是开覆盖。
假设当 $U'$ 是形如 $U_{i_0\ldots i_p}$ 的开集的任意并时，
限制映射 $\mathcal{F}(U)\to\mathcal{F}(U')$ 都是满射。
则对 $p>0$，$\check{H}^p(\mathcal{U},\mathcal{F})$ 消失。
\end{lemma}

\begin{proof}
令 $Y$ 为 $I$ 的所有非空子集所成的集合。
我们用字母 $A,B,C,\ldots$ 表示 $Y$ 的元素，即 $I$ 的非空子集。
对有限非空子集 $J\subset I$，令
$$
V_J = \{A \in Y \mid J \subset A\}
$$
这意味着 $V_{\{i\}}=\{A\in Y\mid i\in A\}$，并且
$V_J = \bigcap_{j \in J} V_{\{j\}}$.
此时 $V_J\subset V_K$ 当且仅当 $J\supset K$。
$Y$ 上存在唯一的拓扑，使得诸子集 $V_J$ 构成拓扑基。
任意开集形如
$$
V = \bigcup\nolimits_{t \in T} V_{J_t}
$$
，其中 $(J_t)_{t\in T}$ 是有限子集族。
若 $J_t\subset J_{t'}$，则从该族中移去 $J_{t'}$ 不会改变 $V$。
因此可假设诸 $J_t$ 之间不存在包含关系。
此时 $V$ 的极小元素正是诸集合 $A=J_t$。
故可由开集 $V$ 恢复该族 $(J_t)_{t\in T}$。

\medskip\noindent
我们可以完全描述 $Y$ 中的开覆盖。首先，由于
$A\in Y$ 是 $I$ 的非空子集，有
$$
Y = \bigcup\nolimits_{i \in I} V_{\{i\}}
$$
为理解其他覆盖，设 $V$ 如上，并设开集 $V_s\subset Y$
对应于族 $(J_{s,t})_{t\in T_s}$。则
$$
V = \bigcup\nolimits_{s \in S} V_s
$$
当且仅当对每个 $t\in T$，存在 $s\in S$ 及
$t_s\in T_s$，使得 $J_t=J_{s,t_s}$。
事实上，由于族 $(J_t)_{t\in T}$ 是极小的，
极小元素 $A=J_t$ 必属于某个 $V_s$，故对某个
$t_s\in T_s$ 有 $A\in V_{J_{t_s}}$。
但 $A$ 在 $V_s$ 中也极小，故 $J_{t_s}=J_t$。

\medskip\noindent
接着，把 $Y$ 的开集集合映到 $X$ 的开集集合。
具体地，把 $Y$ 送到 $U$，并在开集 $V_J$ 上使用规则
$$
V_J \mapsto U_J = \bigcap\nolimits_{i \in J} U_i
$$
，再用规则
$$
V = \bigcup\nolimits_{t \in T} V_{J_t}
\mapsto
\bigcup\nolimits_{t \in T} U_{J_t}
$$
将其延拓到任意开集 $V$。
上述对 $Y$ 的开覆盖的分类说明，该规则把开覆盖送到开覆盖。
因此可如下得到 $Y$ 上的阿贝尔层 $\mathcal{G}$：
置 $\mathcal{G}(Y)=\mathcal{F}(U)$，并对
$V=\bigcup\nolimits_{t\in T}V_{J_t}$ 置
$$
\mathcal{G}(V) = \mathcal{F}\left(\bigcup\nolimits_{t \in T} U_{J_t}\right)
$$
，其限制映射取自 $\mathcal{F}$ 的限制映射。

\medskip\noindent
有了这些准备，即可证明引理。$Y$ 有开覆盖
$\mathcal{V} : Y = \bigcup_{i \in I} V_{\{i\}}$。
按构造，有 {\v C}ech 复形的等式
$$
\check{C}^\bullet(\mathcal{V}, \mathcal{G}) =
\check{C}^\bullet(\mathcal{U}, \mathcal{F})
$$
。由引理陈述中关于 $\mathcal{F}$ 的假设，
层 $\mathcal{G}$ 在 $Y$ 上松弛，故所需消失性由引理 \ref{cohomology-lemma-flasque-acyclic-cech} 推出。
\end{proof}




\section{Leray 谱序列}
\label{cohomology-section-Leray}

\begin{lemma}
\label{cohomology-lemma-before-Leray}
设 $f:X\to Y$ 是环化空间的态射。存在交换图
$$
\xymatrix{
D^{+}(X) \ar[rr]_-{R\Gamma(X, -)} \ar[d]_{Rf_*} & &
D^{+}(\mathcal{O}_X(X)) \ar[d]^{\text{restriction}} \\
D^{+}(Y) \ar[rr]^-{R\Gamma(Y, -)} & &
D^{+}(\mathcal{O}_Y(Y))
}
$$
更一般地，对任意开集 $V\subset Y$ 及 $U=f^{-1}(V)$，存在交换图
$$
\xymatrix{
D^{+}(X) \ar[rr]_-{R\Gamma(U, -)} \ar[d]_{Rf_*} & &
D^{+}(\mathcal{O}_X(U)) \ar[d]^{\text{restriction}} \\
D^{+}(Y) \ar[rr]^-{R\Gamma(V, -)} & &
D^{+}(\mathcal{O}_Y(V))
}
$$
更多说明见注记 \ref{cohomology-remark-elucidate-lemma}。
\end{lemma}

\begin{proof}
令
$\Gamma_{res} : \textit{Mod}(\mathcal{O}_X) \to \text{Mod}_{\mathcal{O}_Y(Y)}$
为如下函子：它把 $\mathcal{O}_X$-模 $\mathcal{F}$ 送到
$\mathcal{F}$ 的全局截面，并通过映射
$f^\sharp:\mathcal{O}_Y(Y)\to\mathcal{O}_X(X)$
将其视为 $\mathcal{O}_Y(Y)$-模。令
$restriction : \text{Mod}_{\mathcal{O}_X(X)} \to \text{Mod}_{\mathcal{O}_Y(Y)}$
为 $f^\sharp:\mathcal{O}_Y(Y)\to\mathcal{O}_X(X)$
诱导的限制标量函子。注意 $restriction$ 正合，故其右导出函子
只需逐项施用限制标量函子即可计算，见
《导出范畴》引理 \ref{derived-lemma-right-derived-exact-functor}.
显然
$$
\Gamma_{res}
=
restriction \circ \Gamma(X, -)
=
\Gamma(Y, -) \circ f_*
$$
我们断言
《导出范畴》引理 \ref{derived-lemma-compose-derived-functors}
适用于这两个复合。对第一个复合，这由上述讨论显然成立。
对第二个复合，它由引理 \ref{cohomology-lemma-pushforward-injective} 推出；该引理说明，
内射 $\mathcal{O}_X$-模被送到 $Y$ 上关于
$\Gamma(Y,-)$ 无环的层。
\end{proof}

\begin{remark}
\label{cohomology-remark-elucidate-lemma}
下面具体说明引理 \ref{cohomology-lemma-before-Leray} 的含义。
它说明，给定
$f : X \to Y$ 与 $\mathcal{F} \in \textit{Mod}(\mathcal{O}_X)$
以及内射消解 $\mathcal{F}\to\mathcal{I}^\bullet$，有
$$
\begin{matrix}
R\Gamma(X, \mathcal{F}) & \text{由下式表示} &
\Gamma(X, \mathcal{I}^\bullet) \\
Rf_*\mathcal{F} & \text{由下式表示} & f_*\mathcal{I}^\bullet \\
R\Gamma(Y, Rf_*\mathcal{F}) & \text{由下式表示} &
\Gamma(Y, f_*\mathcal{I}^\bullet)
\end{matrix}
$$
，最后一个事实来自 Leray 无环性引理
(《导出范畴》引理 \ref{derived-lemma-leray-acyclicity})
及引理 \ref{cohomology-lemma-pushforward-injective}。
最后，将此与显然的等式
$$
\Gamma(X, \mathcal{I}^\bullet)
=
\Gamma(Y, f_*\mathcal{I}^\bullet).
$$
结合，即得到引理中图表的交换性。
\end{remark}

\begin{lemma}
\label{cohomology-lemma-modules-abelian}
设 $X$ 是环化空间，$\mathcal{F}$ 是 $\mathcal{O}_X$-模。
\begin{enumerate}
\item 对开集 $U\subset X$，把 $\mathcal{F}$ 视为
$\mathcal{O}_X$-模或视为阿贝尔层所计算出的上同调群
$H^i(U,\mathcal{F})$ 相同。
\item 设 $f:X\to Y$ 是环化空间的态射。
把 $\mathcal{F}$ 视为 $\mathcal{O}_X$-模或视为阿贝尔层
所计算出的高阶直像 $R^if_*\mathcal{F}$ 相同。
\end{enumerate}
对 $\mathcal{O}_X$-模的下有界复形也有类似陈述。
\end{lemma}

\begin{proof}
考虑环化空间的态射
$(X, \mathcal{O}_X) \to (X, \underline{\mathbf{Z}}_X)$
；它在底层拓扑空间上为恒等映射，在结构层上由环层的唯一映射
$\underline{\mathbf{Z}}_X\to\mathcal{O}_X$ 给出。
设 $\mathcal{F}$ 是 $\mathcal{O}_X$-模。
以 $\mathcal{F}_{ab}$ 表示将同一层视为
$\underline{\mathbf{Z}}_X$-模，即视为阿贝尔群层。
设 $\mathcal{F}\to\mathcal{I}^\bullet$ 是内射消解。
由注记 \ref{cohomology-remark-elucidate-lemma}，
$\Gamma(X,\mathcal{I}^\bullet)$ 同时计算
$R\Gamma(X,\mathcal{F})$ 与
$R\Gamma(X,\mathcal{F}_{ab})$。这证明 (1)。

\medskip\noindent
为证明 (2)，使用 (1) 及引理 \ref{cohomology-lemma-describe-higher-direct-images}，结论立即成立。
\end{proof}

\begin{lemma}[Leray 谱序列]
\label{cohomology-lemma-Leray}
设 $f:X\to Y$ 是环化空间的态射，
$\mathcal{F}^\bullet$ 是 $\mathcal{O}_X$-模的下有界复形。
存在谱序列
$$
E_2^{p, q} = H^p(Y, R^qf_*(\mathcal{F}^\bullet))
$$
，它收敛到 $H^{p+q}(X,\mathcal{F}^\bullet)$。
\end{lemma}

\begin{proof}
这就是由函子复合
$\Gamma_{res}=\Gamma(Y,-)\circ f_*$ 所得的 Grothendieck 谱序列
《导出范畴》引理 \ref{derived-lemma-grothendieck-spectral-sequence}
，其中 $\Gamma_{res}$ 如引理 \ref{cohomology-lemma-before-Leray} 的证明中所定义。
关于
《导出范畴》引理 \ref{derived-lemma-grothendieck-spectral-sequence}
的假设为何满足，可参见引理 \ref{cohomology-lemma-before-Leray} 的证明或注记 \ref{cohomology-remark-elucidate-lemma}。
\end{proof}

\begin{remark}
\label{cohomology-remark-Leray-ss-more-structure}
按引理 \ref{cohomology-lemma-Leray} 中的证明方式，
Leray 谱序列是 $\Gamma(Y,\mathcal{O}_Y)$-模的谱序列。
不过很容易看出，它事实上是 $\Gamma(X,\mathcal{O}_X)$-模的谱序列。
例如，$f$ 给出环化空间的态射
$f' :  (X, \mathcal{O}_X) \to (Y, f_*\mathcal{O}_X)$.
由引理 \ref{cohomology-lemma-modules-abelian}，对 $\mathcal{O}_X$-模
$\mathcal{F}$ 而言，与 $f$ 相伴的 Leray 谱序列之各项
$E_r^{p,q}$ 至少在 $r\geq2$ 时，与 $f'$ 所给出的各项相同。
事实上，二者都等于把 $\mathcal{F}$ 视为阿贝尔层所得的
Leray 谱序列之各项。又因
$(f_*\mathcal{O}_X)(Y)=\mathcal{O}_X(X)$，结论成立。
Leray 谱序列往往还带有额外结构。
\end{remark}

\begin{lemma}
\label{cohomology-lemma-apply-Leray}
设 $f:X\to Y$ 是环化空间的态射，
$\mathcal{F}$ 是 $\mathcal{O}_X$-模。
\begin{enumerate}
\item 若对 $q>0$ 有 $R^qf_*\mathcal{F}=0$，则对所有 $p$，
$H^p(X,\mathcal{F})=H^p(Y,f_*\mathcal{F})$。
\item 若对所有 $q$ 及 $p>0$ 有
$H^p(Y,R^qf_*\mathcal{F})=0$，则对所有 $q$，
$H^q(X,\mathcal{F})=H^0(Y,R^qf_*\mathcal{F})$。
\end{enumerate}
\end{lemma}

\begin{proof}
这两个简单条件都迫使 Leray 谱序列在 $E_2$ 页退化。
也可以不使用谱序列直接证明这些事实；
这是层上同调方面的一个好练习。
\end{proof}

\begin{lemma}
\label{cohomology-lemma-higher-direct-images-compose}
\begin{slogan}
复合函子的全导出函子等于诸全导出函子的复合。
\end{slogan}
设 $f:X\to Y$ 及 $g:Y\to Z$ 是环化空间的态射。
此时，作为从 $D^{+}(X)$ 到 $D^{+}(Z)$ 的函子，有
$Rg_*\circ Rf_*=R(g\circ f)_*$。
\end{lemma}

\begin{proof}
施用
《导出范畴》引理 \ref{derived-lemma-compose-derived-functors}.
显然 $g_*\circ f_*=(g\circ f)_*$，见
《层》引理 \ref{sheaves-lemma-pushforward-composition}。
只需证明 $f_*\mathcal{I}$ 关于 $g_*$ 无环。
这由引理 \ref{cohomology-lemma-pushforward-injective}
及引理 \ref{cohomology-lemma-describe-higher-direct-images}
中对高阶直像 $R^ig_*$ 的描述推出。
\end{proof}

\begin{lemma}[相对 Leray 谱序列]
\label{cohomology-lemma-relative-Leray}
设 $f:X\to Y$ 及 $g:Y\to Z$ 是环化空间的态射，
$\mathcal{F}$ 是 $\mathcal{O}_X$-模。
存在谱序列，其
$$
E_2^{p, q} = R^pg_*(R^qf_*\mathcal{F})
$$
收敛到 $R^{p+q}(g\circ f)_*\mathcal{F}$。
该谱序列关于 $\mathcal{F}$ 是函子性的，
并且对 $\mathcal{O}_X$-模的下有界复形也有相应版本。
\end{lemma}

\begin{proof}
这是函子复合的 Grothendieck 谱序列，由引理 \ref{cohomology-lemma-higher-direct-images-compose} 及
《导出范畴》引理 \ref{derived-lemma-grothendieck-spectral-sequence} 推出。
\end{proof}














\section{上同调的函子性}
\label{cohomology-section-functoriality}

\begin{lemma}
\label{cohomology-lemma-functoriality}
设 $f:X\to Y$ 是环化空间的态射。
设 $\mathcal{G}^\bullet$、相应地 $\mathcal{F}^\bullet$
分别是 $\mathcal{O}_Y$-模、相应地 $\mathcal{O}_X$-模的下有界复形。
设
$\varphi : \mathcal{G}^\bullet \to f_*\mathcal{F}^\bullet$
是复形态射。则在 $D^{+}(Y)$ 中存在典范态射
$$
\mathcal{G}^\bullet
\longrightarrow
Rf_*(\mathcal{F}^\bullet)
$$
。此外，该构造关于三元组
$(\mathcal{G}^\bullet,\mathcal{F}^\bullet,\varphi)$ 是函子性的。
\end{lemma}

\begin{proof}
取内射消解
$\mathcal{F}^\bullet\to\mathcal{I}^\bullet$。
按定义，$Rf_*(\mathcal{F}^\bullet)$ 在
$K^{+}(\mathcal{O}_Y)$ 中由 $f_*\mathcal{I}^\bullet$ 表示。
复合
$$
\mathcal{G}^\bullet \to f_*\mathcal{F}^\bullet \to f_*\mathcal{I}^\bullet
$$
是 $K^{+}(Y)$ 中的态射；施用局部化函子
$j_Y:K^{+}(Y)\to D^{+}(Y)$ 后，它变为引理中的态射。
\end{proof}

\noindent
设 $f:X\to Y$ 是环化空间的态射，
$\mathcal{G}$ 是 $\mathcal{O}_Y$-模，
$\mathcal{F}$ 是 $\mathcal{O}_X$-模。
回忆，从 $\mathcal{G}$ 到 $\mathcal{F}$ 的 $f$-映射
$\varphi$ 是映射
$\varphi:\mathcal{G}\to f_*\mathcal{F}$；
等价地，它是映射
$\varphi:f^*\mathcal{G}\to\mathcal{F}$。见
《层》定义 \ref{sheaves-definition-f-map}。
这样的 $f$-映射给出 $D^{+}(\mathcal{O}_Y(Y))$ 中的复形态射
\begin{equation}
\label{cohomology-equation-functorial-derived}
\varphi :
R\Gamma(Y, \mathcal{G})
\longrightarrow
R\Gamma(X, \mathcal{F})
\end{equation}
。具体地，使用引理 \ref{cohomology-lemma-functoriality}
给出的 $D^{+}(Y)$ 中的态射
$\mathcal{G}\to Rf_*\mathcal{F}$，并施用 $R\Gamma(Y,-)$。
由引理 \ref{cohomology-lemma-before-Leray}，有
$R\Gamma(X, \mathcal{F}) = R\Gamma(Y, Rf_*\mathcal{F})$
，故得到所示箭头。下面的注记 \ref{cohomology-remark-explain-arrow} 将完整展开这一构造。
特别地，它给出上同调上的映射
\begin{equation}
\label{cohomology-equation-functorial}
\varphi : H^i(Y, \mathcal{G}) \longrightarrow H^i(X, \mathcal{F}).
\end{equation}

\begin{remark}
\label{cohomology-remark-explain-arrow}
设 $f:X\to Y$ 是环化空间的态射，
$\mathcal{G}$ 是 $\mathcal{O}_Y$-模，
$\mathcal{F}$ 是 $\mathcal{O}_X$-模，
$\varphi$ 是从 $\mathcal{G}$ 到 $\mathcal{F}$ 的 $f$-映射。
选取由内射 $\mathcal{O}_X$-模复形给出的消解
$\mathcal{F}\to\mathcal{I}^\bullet$。
再选取由内射 $\mathcal{O}_Y$-模复形给出的消解
$\mathcal{G}\to\mathcal{J}^\bullet$ 及
$f_*\mathcal{I}^\bullet\to(\mathcal{J}')^\bullet$。
由
《导出范畴》引理 \ref{derived-lemma-morphisms-lift}
，存在复形映射 $\beta$，使得图
\begin{equation}
\label{cohomology-equation-choice}
\xymatrix{
\mathcal{G} \ar[d] \ar[r] &
f_*\mathcal{F} \ar[r] &
f_*\mathcal{I}^\bullet \ar[d] \\
\mathcal{J}^\bullet \ar[rr]^\beta & &
(\mathcal{J}')^\bullet
}
\end{equation}
交换。施用全局截面函子，得到图
$$
\xymatrix{
 & & \Gamma(Y, f_*\mathcal{I}^\bullet) \ar[d]_{qis} \ar@{=}[r] &
\Gamma(X, \mathcal{I}^\bullet) \\
\Gamma(Y, \mathcal{J}^\bullet) \ar[rr]^\beta & &
\Gamma(Y, (\mathcal{J}')^\bullet) &
}
$$
。左下方的复形表示 $R\Gamma(Y,\mathcal{G})$，
右上方的复形表示 $R\Gamma(X,\mathcal{F})$。
由引理 \ref{cohomology-lemma-before-Leray}，竖直箭头是拟同构；
施用局部化函子
$K^{+}(\mathcal{O}_Y(Y))\to D^{+}(\mathcal{O}_Y(Y))$
后它变得可逆。
箭头 (\ref{cohomology-equation-functorial-derived})
由水平映射与竖直映射之逆的复合给出。
\end{remark}





\section{加细与 {\v C}ech 上同调}
\label{cohomology-section-refinements-cech}

\noindent
设 $(X,\mathcal{O}_X)$ 是环化空间，并设
$\mathcal{U} : X = \bigcup_{i \in I} U_i$ 及
$\mathcal{V} : X = \bigcup_{j \in J} V_j$ 是开覆盖。
假设 $\mathcal{U}$ 是 $\mathcal{V}$ 的加细。
选择映射 $c:I\to J$，使得对所有 $i\in I$，
$U_i\subset V_{c(i)}$。它诱导 {\v C}ech 复形之间的映射
$$
\gamma :
\check{\mathcal{C}}^\bullet(\mathcal{V}, \mathcal{F})
\longrightarrow
\check{\mathcal{C}}^\bullet(\mathcal{U}, \mathcal{F}),
\quad
(\xi_{j_0 \ldots j_p})
\longmapsto
(\xi_{c(i_0) \ldots c(i_p)}|_{U_{i_0 \ldots i_p}})
$$
，且关于 $\mathcal{O}_X$-模层 $\mathcal{F}$ 是函子性的。
假设 $c':I\to J$ 是另一个映射，使得对所有 $i\in I$，
$U_i\subset V_{c'(i)}$。则相应的映射 $\gamma$ 与 $\gamma'$
彼此同伦。具体地，
$\gamma - \gamma' = \text{d} \circ h + h \circ \text{d}$
，其中
$h : \check{\mathcal{C}}^{p + 1}(\mathcal{V}, \mathcal{F}) \to
\check{\mathcal{C}}^p(\mathcal{U}, \mathcal{F})$
由规则
$$
h(\alpha)_{i_0 \ldots i_p} =
\sum\nolimits_{a = 0}^{p}
(-1)^a
\alpha_{c(i_0)\ldots c(i_a) c'(i_a) \ldots c'(i_p)}
$$
给出。验证此公式的计算从略；更一般情形的证明见
(\ref{cohomology-equation-transformation}) 之后的讨论。
特别地，{\v C}ech 上同调群上的映射与 $c$ 的选择无关。
此外，若 $\mathcal{W}:X=\bigcup_{k\in K}W_k$ 是第三个开覆盖，
且 $\mathcal{V}$ 是 $\mathcal{W}$ 的加细，则显然映射之复合
$$
\check{\mathcal{C}}^\bullet(\mathcal{W}, \mathcal{F})
\longrightarrow
\check{\mathcal{C}}^\bullet(\mathcal{V}, \mathcal{F})
\longrightarrow
\check{\mathcal{C}}^\bullet(\mathcal{U}, \mathcal{F})
$$
（分别与映射 $I\to J$ 及 $J\to K$ 相伴）
等于与复合 $I\to K$ 相伴的映射。
特别地，可定义 {\v C}ech 上同调群
$$
\check{H}^p(X, \mathcal{F}) =
\colim_\mathcal{U} \check{H}^p(\mathcal{U}, \mathcal{F})
$$
，其中余极限取遍按加细预序排列的 $X$ 的所有开覆盖。

\medskip\noindent
上述映射 $\gamma$ 与到层上同调的映射相容。换言之，复合
$$
\check{H}^p(\mathcal{V}, \mathcal{F}) \to
\check{H}^p(\mathcal{U}, \mathcal{F})
\xrightarrow{\text{引理 \ref{cohomology-lemma-cech-cohomology}}}
H^p(X, \mathcal{F})
$$
正是引理 \ref{cohomology-lemma-cech-cohomology}
中从第一个群到层上同调的典范映射。
下面的引理将在稍一般的情形中证明这一点。
由此得到良定义映射
\begin{equation}
\label{cohomology-equation-cech-to-cohomology}
\check{H}^p(X, \mathcal{F}) =
\colim_\mathcal{U} \check{H}^p(\mathcal{U}, \mathcal{F})
\longrightarrow
H^p(X, \mathcal{F})
\end{equation}
，它从 {\v C}ech 上同调映到层上同调。

\begin{lemma}
\label{cohomology-lemma-functoriality-cech}
设 $f:X\to Y$ 是环化空间的态射，
$\varphi:f^*\mathcal{G}\to\mathcal{F}$
是从 $\mathcal{O}_Y$-模 $\mathcal{G}$ 到
$\mathcal{O}_X$-模 $\mathcal{F}$ 的 $f$-映射。
设 $\mathcal{U}:X=\bigcup_{i\in I}U_i$ 及
$\mathcal{V}:Y=\bigcup_{j\in J}V_j$ 是开覆盖。
假设 $\mathcal{U}$ 是
$f^{-1}\mathcal{V}:X=\bigcup_{j\in J}f^{-1}(V_j)$ 的加细。
此时在 $D^{+}(\mathcal{O}_X(X))$ 中存在交换图
$$
\xymatrix{
\check{\mathcal{C}}^\bullet(\mathcal{U}, \mathcal{F}) \ar[r] &
R\Gamma(X, \mathcal{F}) \\
\check{\mathcal{C}}^\bullet(\mathcal{V}, \mathcal{G}) \ar[r]
\ar[u]^\gamma &
R\Gamma(Y, \mathcal{G}) \ar[u]
}
$$
，其中水平箭头由引理 \ref{cohomology-lemma-cech-cohomology} 给出，
右侧竖直箭头由
(\ref{cohomology-equation-functorial-derived}) 给出。
特别地，得到上同调群的交换图
$$
\xymatrix{
\check{H}^p(\mathcal{U}, \mathcal{F}) \ar[r] &
H^p(X, \mathcal{F}) \\
\check{H}^p(\mathcal{V}, \mathcal{G}) \ar[r]
\ar[u]^\gamma &
H^p(Y, \mathcal{G}) \ar[u]
}
$$
，其中右侧竖直箭头是 (\ref{cohomology-equation-functorial})。
\end{lemma}

\begin{proof}
先定义左侧竖直箭头。选择映射 $c:I\to J$，使得对所有
$i\in I$，有 $U_i\subset f^{-1}(V_{c(i)})$。
在次数 $p$，用规则
$$
\gamma(s)_{i_0 \ldots i_p} = \varphi(s)_{c(i_0) \ldots c(i_p)}
$$
定义该映射。这是有意义的，因为由假设，
$\varphi$ 确实诱导映射
$\mathcal{G}(V_{c(i_0) \ldots c(i_p)}) \to \mathcal{F}(U_{i_0 \ldots i_p})$
。显然，这定义了复形态射。选取 $X$ 上的内射消解
$\mathcal{F} \to \mathcal{I}^\bullet$
以及 $Y$ 上的内射消解
$\mathcal{G}\to J^\bullet$。
按照引理 \ref{cohomology-lemma-cech-cohomology} 的证明，
引入双复形 $A^{\bullet,\bullet}$ 与 $B^{\bullet,\bullet}$，其各项为
$$
B^{p, q} = \check{\mathcal{C}}^p(\mathcal{V}, \mathcal{J}^q)
\quad
\text{以及}
\quad
A^{p, q} = \check{\mathcal{C}}^p(\mathcal{U}, \mathcal{I}^q).
$$
如上面的注记 \ref{cohomology-remark-explain-arrow}，还选择 $Y$ 上的内射消解
$f_*\mathcal{I} \to (\mathcal{J}')^\bullet$
及复形态射
$\beta:\mathcal{J}\to(\mathcal{J}')^\bullet$，
使 (\ref{cohomology-equation-choice}) 交换。
再引入双复形 $(B')^{\bullet,\bullet}$ 与
$(B'')^{\bullet,\bullet}$，其各项为
$$
(B')^{p, q} = \check{\mathcal{C}}^p(\mathcal{V}, (\mathcal{J}')^q)
\quad
\text{以及}
\quad
(B'')^{p, q} = \check{\mathcal{C}}^p(\mathcal{V}, f_*\mathcal{I}^q).
$$
注意，存在从 $f_*\mathcal{I}^\bullet$ 到
$\mathcal{I}^\bullet$ 的复形 $f$-映射。
因此显然，上述同一规则定义双复形态射
$$
\gamma : (B'')^{\bullet, \bullet} \longrightarrow A^{\bullet, \bullet}.
$$
考虑复形图
$$
\xymatrix{
\check{\mathcal{C}}^\bullet(\mathcal{U}, \mathcal{F})
\ar[r] &
\text{Tot}(A^{\bullet, \bullet}) & & &
\Gamma(X, \mathcal{I}^\bullet) \ar[lll]^{qis}
\ar@{=}[ddl]\\
\check{\mathcal{C}}^\bullet(\mathcal{V}, \mathcal{G})
\ar[r] \ar[u]^\gamma &
\text{Tot}(B^{\bullet, \bullet}) \ar[r]^\beta &
\text{Tot}((B')^{\bullet, \bullet}) &
\text{Tot}((B'')^{\bullet, \bullet}) \ar[l] \ar[llu]_{s\gamma} \\
& \Gamma(Y, \mathcal{J}^\bullet) \ar[u]^{qis} \ar[r]^\beta &
\Gamma(Y, (\mathcal{J}')^\bullet) \ar[u] &
\Gamma(Y, f_*\mathcal{I}^\bullet) \ar[u] \ar[l]_{qis}
}
$$
。以 $\text{Tot}(A^{\bullet,\bullet})$ 及
$\text{Tot}(B^{\bullet,\bullet})$ 为目标的两个水平箭头，
正是引理 \ref{cohomology-lemma-cech-cohomology} 中解释的箭头。
左上方的图形（五边形）交换，这直接来自
(\ref{cohomology-equation-choice}) 的交换性。
下面两个正方形显然交换。
由定义也立即可知，右上方的图形（正方形）交换。
现在由定义及下述事实得到引理：沿图的外边界，
无论从上方还是下方由
$\check{\mathcal{C}}^\bullet(\mathcal{V},\mathcal{G})$
走到 $\Gamma(X,\mathcal{I}^\bullet)$，所得结果相同
（途中须将所遇到的拟同构取逆）。
\end{proof}





\section{Hausdorff 拟紧空间上的上同调}
\label{cohomology-section-cohomology-LC}

\noindent
对这样的空间，{\v C}ech 上同调与层上同调一致。

\begin{lemma}
\label{cohomology-lemma-cech-always}
设 $X$ 是拓扑空间，$\mathcal{F}$ 是阿贝尔层。
则 (\ref{cohomology-equation-cech-to-cohomology}) 中定义的映射
$\check{H}^1(X,\mathcal{F})\to H^1(X,\mathcal{F})$ 是同构。
\end{lemma}

\begin{proof}
设 $\mathcal{U}$ 是 $X$ 的开覆盖。
由引理 \ref{cohomology-lemma-cech-spectral-sequence}，存在正合序列
$$
0 \to \check{H}^1(\mathcal{U}, \mathcal{F}) \to H^1(X, \mathcal{F})
\to \check{H}^0(\mathcal{U}, \underline{H}^1(\mathcal{F}))
$$
故该映射是单射。为证明满射性，只需证明：
用 $\mathcal{U}$ 的某个加细替换 $\mathcal{U}$ 后，
$\check{H}^0(\mathcal{U},\underline{H}^1(\mathcal{F}))$
中的任意元素都映到零。
这由定义及下述事实立即得到：
$\underline{H}^1(\mathcal{F})$ 是阿贝尔群预层，
由上同调的局部性，其层化为零；见引理 \ref{cohomology-lemma-kill-cohomology-class-on-covering}。
\end{proof}

\begin{lemma}
\label{cohomology-lemma-cech-Hausdorff-quasi-compact}
设 $X$ 是 Hausdorff 拟紧拓扑空间，
$\mathcal{F}$ 是 $X$ 上的阿贝尔层。
则对所有 $n$，(\ref{cohomology-equation-cech-to-cohomology}) 中定义的映射
$\check{H}^n(X,\mathcal{F})\to H^n(X,\mathcal{F})$ 是同构。
\end{lemma}

\begin{proof}
我们已经知道，当 $n=0,1$ 时，
$\check{H}^n(X,-)\to H^n(X,-)$ 是函子同构，见引理 \ref{cohomology-lemma-cech-always}。
诸函子 $H^n(X,-)$ 组成泛 $\delta$-函子，见
《导出范畴》引理 \ref{derived-lemma-higher-derived-functors}。
若证明 $\check{H}^n(X,-)$ 组成泛 $\delta$-函子，
且 $\check{H}^n(X,-)\to H^n(X,-)$ 与连接映射相容，
则由泛 $\delta$-函子的唯一性，该映射自动是同构，见
《同调代数》引理 \ref{homology-lemma-uniqueness-universal-delta-functor}。

\medskip\noindent
设
$0\to\mathcal{F}\to\mathcal{G}\to\mathcal{H}\to0$
是 $X$ 上阿贝尔层的短正合序列，
$\mathcal{U}:X=\bigcup_{i\in I}U_i$ 是开覆盖。
它给出复形的复形
$$
0 \to \check{\mathcal{C}}^\bullet(\mathcal{U}, \mathcal{F}) \to
\check{\mathcal{C}}^\bullet(\mathcal{U}, \mathcal{G}) \to
\check{\mathcal{C}}^\bullet(\mathcal{U}, \mathcal{H}) \to 0
$$
，该复形一般在右端不正合。此序列定义映射
$$
\check{H}^n(\mathcal{U}, \mathcal{F}) \to
\check{H}^n(\mathcal{U}, \mathcal{G}) \to
\check{H}^n(\mathcal{U}, \mathcal{H})
$$
，但不足以定义连接算子
$\delta : \check{H}^n(\mathcal{U}, \mathcal{H}) \to
\check{H}^{n + 1}(\mathcal{U}, \mathcal{F})$。
事实上，这样的算子一般不存在。不过，给定元素
$\overline{h}\in\check{H}^n(\mathcal{U},\mathcal{H})$，
它是上循环 $h=(h_{i_0\ldots i_n})$ 的上同调类；
我们可以选择开覆盖
$$
U_{i_0 \ldots i_n} = \bigcup W_{i_0 \ldots i_n, k}
$$
，使得
$h_{i_0\ldots i_n}|_{W_{i_0\ldots i_n,k}}$
可提升为 $\mathcal{G}$ 在 $W_{i_0\ldots i_n,k}$ 上的截面。
由 《拓扑》引理 \ref{topology-lemma-refine-covering}
（正是在此使用 $X$ 为 Hausdorff 拟紧空间的假设），
可选择开覆盖 $\mathcal{V}:X=\bigcup_{j\in J}V_j$
及 $\alpha:J\to I$，使得
$V_j\subset U_{\alpha(j)}$（故它是加细），并且对所有
$j_0,\ldots,j_n\in J$，存在 $k$ 使得
$V_{j_0 \ldots j_n} \subset W_{\alpha(j_0) \ldots \alpha(j_n), k}$.
于是得到复形之间的映射
$$
\xymatrix{
0 \ar[r] &
\check{\mathcal{C}}^\bullet(\mathcal{U}, \mathcal{F}) \ar[d] \ar[r] &
\check{\mathcal{C}}^\bullet(\mathcal{U}, \mathcal{G}) \ar[d] \ar[r] &
\check{\mathcal{C}}^\bullet(\mathcal{U}, \mathcal{H}) \ar[d] \ar[r] &
0 \\
0 \ar[r] &
\check{\mathcal{C}}^\bullet(\mathcal{V}, \mathcal{F}) \ar[r] &
\check{\mathcal{C}}^\bullet(\mathcal{V}, \mathcal{G}) \ar[r] &
\check{\mathcal{C}}^\bullet(\mathcal{V}, \mathcal{H}) \ar[r] &
0
}
$$
事实上，竖直箭头正是定义 {\v C}ech 上同调群之间转移映射
所用的复形映射。我们对加细的选择说明，可以选择
$$
g_{j_0 \ldots j_n} \in
\mathcal{G}(V_{j_0 \ldots j_n}),\quad
g_{j_0 \ldots j_n} \longmapsto
h_{\alpha(j_0) \ldots \alpha(j_n)}|_{V_{j_0 \ldots j_n}}
$$
上链 $g=(g_{j_0\ldots j_n})$ 一般不是上循环，
但其 {\v C}ech 边界 $\text{d}(g)$ 在
$\check{\mathcal{C}}^{n+1}(\mathcal{V},\mathcal{H})$
中映到零（由上面的交换图及 $h$ 为上循环这一事实）。
故 $\text{d}(g)$ 是
$\check{\mathcal{C}}^\bullet(\mathcal{V},\mathcal{F})$ 中的上循环。
这使我们可以定义
$$
\delta(\overline{h}) = \text{下式中 }\text{d}(g)\text{ 的类：}
\check{H}^{n + 1}(\mathcal{V}, \mathcal{F})
$$
现在，给定元素 $\xi\in\check{H}^n(X,\mathcal{G})$，
选择开覆盖 $\mathcal{U}$ 及元素
$\overline{h}\in\check{H}^n(\mathcal{U},\mathcal{G})$，
它在定义 {\v C}ech 上同调的余极限中映到 $\xi$。
再如上选择 $\mathcal{V}$ 与 $g$，并定义
$\delta(\xi)$ 为 $\delta(\overline{h})$
在 $\check{H}^n(X,\mathcal{F})$ 中的像。
此时需验证许多性质，但它们都是直接的。
例如，需验证该构造与
$\mathcal{U},\overline{h},\mathcal{V},\alpha:J\to I,g$
的选择无关。
$\text{d}(g)$ 的类与提升 $g_{i_0\ldots i_n}$ 的选择无关，
因为二者之差将是上边缘。
它与 $\alpha$ 的选择无关\footnote{这项验证很重要：
$\alpha$ 的非唯一性正是阻止我们对 $X$ 的所有开覆盖取
{\v C}ech 复形的余极限，从而得到计算 {\v C}ech 上同调的
复形短正合序列的唯一障碍。}，
因为 $\alpha$ 的不同选择在上图中确定彼此同伦的竖直复形映射，
见 第 \ref{cohomology-section-refinements-cech} 节。
对其他选择，我们使用如下事实：
给定 $X$ 的有限多个开覆盖，总能找到同时加细它们的开覆盖。
还需验证加性，方法相同。
最后，需验证映射
$\check{H}^n(X,-)\to H^n(X,-)$ 与连接映射相容。
为此选择内射消解
$$
\xymatrix{
0 \ar[r] &
\mathcal{F} \ar[r] \ar[d] &
\mathcal{G} \ar[r] \ar[d] &
\mathcal{H} \ar[r] \ar[d] &
0 \\
0 \ar[r] &
\mathcal{I}_1^\bullet \ar[r] &
\mathcal{I}_2^\bullet \ar[r] &
\mathcal{I}_3^\bullet \ar[r] &
0
}
$$
，如 《导出范畴》引理 \ref{derived-lemma-injective-resolution-ses} 中所述。
这给出交换图
$$
\xymatrix{
0 \ar[r] &
\check{\mathcal{C}}^\bullet(\mathcal{U}, \mathcal{F}) \ar[r] \ar[d] &
\check{\mathcal{C}}^\bullet(\mathcal{U}, \mathcal{F}) \ar[r] \ar[d] &
\check{\mathcal{C}}^\bullet(\mathcal{U}, \mathcal{F}) \ar[r] \ar[d] &
0 \\
0 \ar[r] &
\text{Tot}(\check{\mathcal{C}}^\bullet(\mathcal{U}, \mathcal{I}_1^\bullet))
\ar[r] &
\text{Tot}(\check{\mathcal{C}}^\bullet(\mathcal{U}, \mathcal{I}_2^\bullet))
\ar[r] &
\text{Tot}(\check{\mathcal{C}}^\bullet(\mathcal{U}, \mathcal{I}_3^\bullet))
\ar[r] &
0
}
$$
这里 $\mathcal{U}$ 是如上的开覆盖，
而竖直映射正是定义映射
$\check{H}^n(\mathcal{U},-)\to H^n(X,-)$ 时所用的映射，见引理 \ref{cohomology-lemma-cech-cohomology}。
下面的复形正合，因为内射对象的复形序列逐项分裂正合。
因此，上同调中的连接映射由下方正合序列的通常程序计算，见
《同调代数》引理 \ref{homology-lemma-long-exact-sequence-cochain}。
传递到定义 {\v C}ech 上同调连接映射的加细
$\mathcal{V}$ 后，同样如此。
故这些连接映射一致，因为它们使用相同的构造
（只要前者定义在某给定开覆盖的 {\v C}ech 上同调元素上）。
至此，我们完成了关于 {\v C}ech 上同调上
$\delta$-函子结构之构造，以及为何该结构与通常上同调上给定的
$\delta$-函子结构相容的讨论。

\medskip\noindent
最后，施用引理 \ref{cohomology-lemma-injective-trivial-cech}
可知内射层的高阶 {\v C}ech 上同调为零。
因此由
《同调代数》引理 \ref{homology-lemma-efface-implies-universal}，
{\v C}ech 上同调是泛 $\delta$-函子。
\end{proof}

\begin{lemma}
\label{cohomology-lemma-cohomology-of-closed}
\begin{reference}
\cite[Expose V bis, 4.1.3]{SGA4}
\end{reference}
设 $X$ 是拓扑空间，$Z\subset X$ 是拟紧子集，
且 $Z$ 中任意两点在 $X$ 中都有不交的开邻域。
对 $X$ 上每个阿贝尔层 $\mathcal{F}$，典范映射
$$
\colim H^p(U, \mathcal{F})
\longrightarrow
H^p(Z, \mathcal{F}|_Z)
$$
是同构，其中余极限取遍 $Z$ 在 $X$ 中的开邻域 $U$。
\end{lemma}

\begin{proof}
先证明 $p=0$ 的情形。单射性由 $\mathcal{F}|_Z$ 的定义推出，
并且一般都成立（对任意拓扑空间 $X$ 的任意子集）。
再设 $s\in H^0(Z,\mathcal{F}|_Z)$。
可找到开集 $U_i\subset X$，使得
$Z\subset\bigcup U_i$，且
$s|_{Z\cap U_i}$ 来自某个 $s_i\in\mathcal{F}(U_i)$。
于是存在开集 $W_{ij}\subset U_i\cap U_j$，满足
$W_{ij}\cap Z=U_i\cap U_j\cap Z$，且
$s_i|_{W_{ij}}=s_j|_{W_{ij}}$。
施用 《拓扑》引理
\ref{topology-lemma-lift-covering-of-quasi-compact-hausdorff-subset}，
可找到 $X$ 的开集 $V_i$，使得
$V_i\subset U_i$ 且 $V_i\cap V_j\subset W_{ij}$。
因此，诸 $s_i|_{V_i}$ 黏合成 $\mathcal{F}$
在 $Z$ 的开邻域 $\bigcup V_i$ 上的截面。

\medskip\noindent
为完成证明，只需说明：若 $\mathcal{I}$ 是 $X$ 上的内射阿贝尔层，
则对 $p>0$ 有 $H^p(Z,\mathcal{I}|_Z)=0$。
这可用短正合序列与维数转移推出；细节从略。
因此，设对某个 $p>0$，
$\overline{\xi}\in H^p(Z,\mathcal{I}|_Z)$。
由引理 \ref{cohomology-lemma-cech-Hausdorff-quasi-compact}，
元素 $\overline{\xi}$ 来自
$\check{H}^p(\mathcal{V},\mathcal{I}|_Z)$，
其中 $\mathcal{V}:Z=\bigcup V_i$ 是 $Z$ 的某个开覆盖。
设 $\overline{\xi}$ 是
$\check{\mathcal{C}}^p(\mathcal{V},\mathcal{I}|_Z)$
中上循环 $\xi=(\xi_{i_0\ldots i_p})$ 之类的像。

\medskip\noindent
令 $\mathcal{I}'\subset\mathcal{I}|_Z$ 为由下列规则定义的子预层：
$$
\mathcal{I}'(V) =
\{s \in \mathcal{I}|_Z(V) \mid
\exists (U, t),\ U \subset X\text{ 开},
\ t \in \mathcal{I}(U),\ V = Z \cap U,\ s = t|_{Z \cap U} \}
$$
则 $\mathcal{I}|_Z$ 是 $\mathcal{I}'$ 的层化。
因此，对每个 $(p+1)$-元组 $i_0\ldots i_p$，
可找到开覆盖
$V_{i_0\ldots i_p}=\bigcup W_{i_0\ldots i_p,k}$，
使得
$\xi_{i_0\ldots i_p}|_{W_{i_0\ldots i_p,k}}$
是 $\mathcal{I}'$ 的截面。
施用 《拓扑》引理 \ref{topology-lemma-refine-covering}，
在加细 $\mathcal{V}$ 后，可假设每个
$\xi_{i_0\ldots i_p}$ 都是预层 $\mathcal{I}'$ 的截面。

\medskip\noindent
将 $V_i$ 写成 $Z\cap U_i$，其中 $U_i\subset X$ 开。
由于 $\mathcal{I}$ 松弛
（引理 \ref{cohomology-lemma-injective-flasque}），
且对每个 $(p+1)$-元组 $i_0\ldots i_p$，
$\xi_{i_0\ldots i_p}$ 都是 $\mathcal{I}'$ 的截面，
可选择截面
$s_{i_0\ldots i_p}\in\mathcal{I}(U_{i_0\ldots i_p})$，
它在
$V_{i_0\ldots i_p}=Z\cap U_{i_0\ldots i_p}$
上限制为 $\xi_{i_0\ldots i_p}$。
（还可以再次施用
《拓扑》引理
\ref{topology-lemma-lift-covering-of-quasi-compact-hausdorff-subset}，
从而避免使用内射层松弛这一事实。）
令 $s=(s_{i_0\ldots i_p})$ 为开覆盖
$U=\bigcup U_i$ 的相应上链。
由于 $\text{d}(\xi)=0$，截面
$\text{d}(s)_{i_0\ldots i_{p+1}}$
在 $Z\cap U_{i_0\ldots i_{p+1}}$ 上限制为零。
故由证明开头的讨论，存在开子集
$W_{i_0\ldots i_{p+1}}\subset U_{i_0\ldots i_{p+1}}$，
满足
$Z\cap W_{i_0\ldots i_{p+1}}
=Z\cap U_{i_0\ldots i_{p+1}}$，且
$\text{d}(s)_{i_0\ldots i_{p+1}}|_{W_{i_0\ldots i_{p+1}}}=0$。
由 《拓扑》引理
\ref{topology-lemma-lift-covering-of-quasi-compact-hausdorff-subset}
，可找到 $U'_i\subset U_i$，使得
$Z\subset\bigcup U'_i$，且
$U'_{i_0\ldots i_{p+1}}\subset W_{i_0\ldots i_{p+1}}$。
令 $s'=(s'_{i_0\ldots i_p})$，其中
$s'_{i_0 \ldots i_p} = s_{i_0 \ldots i_p}|_{U'_{i_0 \ldots i_p}}$
。则 $s'$ 是 $\mathcal{I}$ 关于开覆盖
$U'=\bigcup U'_i$ 的上循环，其中 $U'$ 是 $Z$ 的开邻域。
由于 $\mathcal{I}$ 的高阶 {\v C}ech 上同调群为零
（引理 \ref{cohomology-lemma-injective-trivial-cech}），
$s'$ 是上边缘。
因此，$\xi$ 在开覆盖
$Z=\bigcup Z\cap U'_i$ 的 {\v C}ech 复形中的像是上边缘，证明完成。
\end{proof}






\section{基变换映射}
\label{cohomology-section-base-change-map}

\noindent
在若干情形下，我们需要知道如何构造基变换映射。由于尚未讨论
导出拉回，这里只讨论沿环空间的平坦态射作基变换的情形。
在陈述结果之前，先讨论导出范畴上的平坦拉回。具体而言，设
$g : X \to Y$ 为环空间的平坦态射。由《模层》引理
\ref{modules-lemma-pullback-flat}，函子
$g^* : \textit{Mod}(\mathcal{O}_Y) \to
\textit{Mod}(\mathcal{O}_X)$ 正合。因此它有导出函子
$$
g^* : D^{+}(Y) \to D^{+}(X)
$$
其计算只需拉回 $D^{+}(Y)$ 中给定对象的一个代表，见《导出范畴》
引理 \ref{derived-lemma-right-derived-exact-functor}。因此，如上所示，
我们将此函子记为 $g^*$ 而非 $Lg^*$。

\begin{lemma}
\label{cohomology-lemma-base-change-map-flat-case}
设
$$
\xymatrix{
X' \ar[r]_{g'} \ar[d]_{f'} &
X \ar[d]^f \\
S' \ar[r]^g &
S
}
$$
为环空间的交换图。设 $\mathcal{F}^\bullet$ 为
$\mathcal{O}_X$-模层的下有界复形。假设 $g$ 与 $g'$ 均平坦。
则在 $D^{+}(S')$ 中存在典范基变换映射
$$
g^*Rf_*\mathcal{F}^\bullet
\longrightarrow
R(f')_*(g')^*\mathcal{F}^\bullet
$$
\end{lemma}

\begin{proof}
取内射消解 $\mathcal{F}^\bullet \to \mathcal{I}^\bullet$
与 $(g')^*\mathcal{F}^\bullet \to \mathcal{J}^\bullet$。
由引理 \ref{cohomology-lemma-pushforward-injective-flat}，
$(g')_*\mathcal{J}^\bullet$ 是表示
$R(g')_*(g')^*\mathcal{F}^\bullet$ 的内射复形。因此，由《导出范畴》
引理 \ref{derived-lemma-morphisms-lift}
与 \ref{derived-lemma-morphisms-equal-up-to-homotopy}，图
$$
\xymatrix{
(g')_*(g')^*\mathcal{F}^\bullet \ar[r] &
(g')_*\mathcal{J}^\bullet \\
\mathcal{F}^\bullet \ar[u]^{adjunction} \ar[r] &
\mathcal{I}^\bullet \ar[u]_\beta
}
$$
中的箭头 $\beta$ 存在且在同伦意义下唯一。取到 $S$ 上的直像，得到
$$
f_*\beta :
f_*\mathcal{I}^\bullet
\longrightarrow
f_*(g')_*\mathcal{J}^\bullet
=
g_*(f')_*\mathcal{J}^\bullet
$$
由 $g^*$ 与 $g_*$ 的伴随性，得到复形态射
$g^*f_*\mathcal{I}^\bullet \to (f')_*\mathcal{J}^\bullet$。
注意，此映射在同伦意义下唯一，因为整个过程中唯一的选择是映射
$\beta$，而且所有构造均在复形层面完成。
\end{proof}

\begin{remark}
\label{cohomology-remark-correct-version-base-change-map}
基变换映射的“正确”版本是映射
$$
Lg^* Rf_* \mathcal{F}^\bullet
\longrightarrow
R(f')_* L(g')^*\mathcal{F}^\bullet.
$$
该映射的构造涉及无界复形，见评注 \ref{cohomology-remark-base-change}。
\end{remark}





\section{拓扑中的真基变换}
\label{cohomology-section-proper-base-change}

\noindent
本节证明拓扑中真基变换定理的一个非常一般的版本。它说明
高阶直像 $R^pf_*$ 的茎可在纤维上计算。

\begin{lemma}
\label{cohomology-lemma-proper-base-change}
设 $f : (X, \mathcal{O}_X) \to (Y, \mathcal{O}_Y)$ 为环空间态射，
且 $y \in Y$。假设
\begin{enumerate}
\item $f$ 是闭映射，
\item $f$ 是分离的，并且
\item $f^{-1}(y)$ 拟紧。
\end{enumerate}
则对 $D^+(\mathcal{O}_X)$ 中的 $E$，在
$D^+(\mathcal{O}_{Y, y})$ 中有
$(Rf_*E)_y = R\Gamma(f^{-1}(y), E|_{f^{-1}(y)})$。
\end{lemma}

\begin{proof}
引理 \ref{cohomology-lemma-base-change-map-flat-case} 的基变换映射给出典范映射
$(Rf_*E)_y \to R\Gamma(f^{-1}(y), E|_{f^{-1}(y)})$。
为证明此映射为同构，以内射对象的下有界复形
$\mathcal{I}^\bullet$ 表示 $E$。置 $Z = f^{-1}(\{y\})$。引理 \ref{cohomology-lemma-cohomology-of-closed} 的假设成立，见《拓扑》
引理 \ref{topology-lemma-separated}。因此各限制
$\mathcal{I}^n|_Z$ 对 $\Gamma(Z, -)$ 无环。于是
$R\Gamma(Z, E|_Z)$ 由复形
$\Gamma(Z, \mathcal{I}^\bullet|_Z)$ 表示，见《导出范畴》
引理 \ref{derived-lemma-leray-acyclicity}。换言之，我们须证明映射
$$
\colim_V \mathcal{I}^\bullet(f^{-1}(V))
\longrightarrow
\Gamma(Z, \mathcal{I}^\bullet|_Z)
$$
为同构。利用引理 \ref{cohomology-lemma-cohomology-of-closed} 可知，只需证明
$Z = f^{-1}(\{y\})$ 的开邻域 $f^{-1}(V)$ 所成的集合族在全部开邻域
组成的系统中共尾。若 $f^{-1}(\{y\}) \subset U$ 是一个开邻域，
则由于 $f$ 为闭映射，集合
$V = Y \setminus f(X \setminus U)$ 是 $y$ 的开邻域，且
$f^{-1}(V) \subset U$。这就证明了引理。
\end{proof}

\begin{theorem}[真基变换]
\label{cohomology-theorem-proper-base-change}
\begin{reference}
\cite[Expose V bis, 4.1.1]{SGA4}
\end{reference}
考虑拓扑空间的笛卡尔方块
$$
\xymatrix{
X' = Y' \times_Y X \ar[d]_{f'} \ar[r]_-{g'} & X \ar[d]^f \\
Y' \ar[r]^g & Y
}
$$
假设 $f$ 为真映射。设 $E$ 为 $D^+(X)$ 中的对象。则引理
\ref{cohomology-lemma-base-change-map-flat-case} 的基变换映射
$$
g^{-1}Rf_*E \longrightarrow Rf'_*(g')^{-1}E
$$
是 $D^+(Y')$ 中的同构。
\end{theorem}

\begin{proof}
设 $y' \in Y'$ 为像是 $y \in Y$ 的点。只需证明基变换映射在
$y'$ 处的茎上诱导同构。由于 $f$ 为真映射，可知 $f'$ 也是真映射，
$f$ 与 $f'$ 的纤维均拟紧，并且 $f$ 与 $f'$ 均为闭映射；见《拓扑》
定理 \ref{topology-theorem-characterize-proper} 与引理
\ref{topology-lemma-base-change-separated}。因而可两次应用引理
\ref{cohomology-lemma-proper-base-change}，得到
$$
(Rf'_*(g')^{-1}E)_{y'} = R\Gamma((f')^{-1}(y'), (g')^{-1}E|_{(f')^{-1}(y')})
$$
以及
$$
(Rf_*E)_y = R\Gamma(f^{-1}(y), E|_{f^{-1}(y)})
$$
纤维间的诱导映射 $(f')^{-1}(y') \to f^{-1}(y)$ 是拓扑空间的同胚，
而 $E|_{f^{-1}(y)}$ 的拉回为
$(g')^{-1}E|_{(f')^{-1}(y')}$。所需结论随即成立。
\end{proof}

\begin{lemma}[集合层的真基变换]
\label{cohomology-lemma-proper-base-change-sheaves-of-sets}
考虑拓扑空间的笛卡尔方块
$$
\xymatrix{
X' \ar[d]_{f'} \ar[r]_-{g'} & X \ar[d]^f \\
Y' \ar[r]^g & Y
}
$$
假设 $f$ 为真映射。则对 $X$ 上的任意集合层 $\mathcal{F}$，有
$g^{-1}f_*\mathcal{F} = f'_*(g')^{-1}\mathcal{F}$。
\end{lemma}

\begin{proof}
完全仿照定理 \ref{cohomology-theorem-proper-base-change} 的证明，可知只需证明
$(f_*\mathcal{F})_y = \Gamma(X_y, \mathcal{F}|_{X_y})$。
再仿照引理 \ref{cohomology-lemma-proper-base-change} 的论证，将此归结为集合层情形下引理 \ref{cohomology-lemma-cohomology-of-closed} 的 $p = 0$ 情形。引理
\ref{cohomology-lemma-cohomology-of-closed} 证明的第一部分同样适用于集合层，
从而完成证明。略去若干细节。
\end{proof}



\section{上同调与余极限}
\label{cohomology-section-limits}

\noindent
设 $X$ 为环空间。设 $(\mathcal{F}_i, \varphi_{ii'})$ 为由有向集
$I$ 指标的 $\mathcal{O}_X$-模层系统，见《范畴》第
\ref{categories-section-posets-limits} 节。由于对每个 $i$ 都有典范映射
$\mathcal{F}_i \to \colim_i \mathcal{F}_i$，我们对每个 $p \geq 0$
得到典范映射
$$
\colim_i H^p(X, \mathcal{F}_i)
\longrightarrow
H^p(X, \colim_i \mathcal{F}_i)
$$
对每个开集 $U \subset X$ 当然也有类似映射。一般而言，这些映射并非
同构，即使 $p = 0$ 也是如此。本节推广《层》引理
\ref{sheaves-lemma-directed-colimits-sections} 的结果。另见《模层》
引理 \ref{modules-lemma-finite-presentation-quasi-compact-colimit}
（取特殊情形 $\mathcal{G} = \mathcal{O}_X$）。

\begin{lemma}
\label{cohomology-lemma-quasi-separated-cohomology-colimit}
设 $X$ 为环空间。假设 $X$ 的底拓扑空间具有下列性质：
\begin{enumerate}
\item 存在由拟紧开子集组成的拓扑基，并且
\item 任意两个拟紧开集的交仍拟紧。
\end{enumerate}
则对 $\mathcal{O}_X$-模层的任意有向系统
$(\mathcal{F}_i, \varphi_{ii'})$ 以及任意拟紧开集 $U \subset X$，
典范映射
$$
\colim_i H^q(U, \mathcal{F}_i)
\longrightarrow
H^q(U, \colim_i \mathcal{F}_i)
$$
对每个 $q \geq 0$ 都是同构。
\end{lemma}

\begin{proof}
本证明的关键是同时对所有拟紧开集 $U \subset X$ 论证。由《层》引理
\ref{sheaves-lemma-directed-colimits-sections}（结合《拓扑》引理
\ref{topology-lemma-topology-quasi-separated-scheme}），当 $q = 0$
时结论对任意拟紧开集 $U \subset X$ 成立。假设已经对所有
$q \leq q_0$ 证明了结论；下面证明 $q = q_0 + 1$ 的情形。

\medskip\noindent
按我们关于有向系统的约定，指标集 $I$ 是有向的，而 $I$ 上的任意
$\mathcal{O}_X$-模层系统 $(\mathcal{F}_i, \varphi_{ii'})$ 都是有向系统。
由《内射对象》引理 \ref{injectives-lemma-sheaves-modules-space}，
$\mathcal{O}_X$-模层范畴有函子性内射嵌入。因此，对任意系统
$(\mathcal{F}_i, \varphi_{ii'})$，存在系统
$(\mathcal{I}_i, \varphi_{ii'})$，其中每个 $\mathcal{I}_i$ 都是内射
$\mathcal{O}_X$-模层，并且存在由内射 $\mathcal{O}_X$-模同态
$\mathcal{F}_i \to \mathcal{I}_i$ 给出的系统态射。记其余核为
$\mathcal{Q}_i$，于是有短正合列
$$
0 \to
\mathcal{F}_i \to
\mathcal{I}_i \to
\mathcal{Q}_i \to 0.
$$
我们断言序列
$$
0 \to
\colim_i \mathcal{F}_i \to
\colim_i \mathcal{I}_i \to
\colim_i \mathcal{Q}_i \to 0.
$$
也是 $\mathcal{O}_X$-模层的短正合列。可在茎上检验这一点。由《层》第
\ref{sheaves-section-limits-presheaves} 节与第
\ref{sheaves-section-limits-sheaves} 节，取茎与余极限可交换。由于
阿贝尔群短正合列的有向余极限仍短正合（见《代数》引理
\ref{algebra-lemma-directed-colimit-exact}），断言成立。我们还断言：
对每个拟紧开集 $U \subset X$ 及每个 $q \geq 1$，有
$H^q(U, \colim_i \mathcal{I}_i) = 0$。暂且接受这一断言，考虑图
$$
\xymatrix{
\colim_i H^{q_0}(U, \mathcal{I}_i) \ar[d] \ar[r] &
\colim_i H^{q_0}(U, \mathcal{Q}_i) \ar[d] \ar[r] &
\colim_i H^{q_0 + 1}(U, \mathcal{F}_i) \ar[d] \ar[r] &
0 \ar[d] \\
H^{q_0}(U, \colim_i \mathcal{I}_i) \ar[r] &
H^{q_0}(U, \colim_i \mathcal{Q}_i) \ar[r] &
H^{q_0 + 1}(U, \colim_i \mathcal{F}_i) \ar[r] &
0
}
$$
右下角的零来自上述断言，右上角的零则来自各层
$\mathcal{I}_i$ 的内射性。应用《代数》引理
\ref{algebra-lemma-directed-colimit-exact} 可知上行正合。因此由蛇形引理
得到 $q = q_0 + 1$ 的结论。

\medskip\noindent
只剩证明该断言。我们将使用引理 \ref{cohomology-lemma-cech-vanish-basis}。
令 $\mathcal{B}$ 为 $X$ 的全部拟紧开子集所成的集合。按假设，
这是 $X$ 的拓扑基。令 $\text{Cov}$ 为有限开覆盖
$\mathcal{U} : U = \bigcup_{j = 1, \ldots, m} U_j$ 所成的集合，其中
$U$ 与各 $U_j$ 均为 $X$ 中的拟紧开集。由 $q = 0$ 的结论，对
$\mathcal{U} \in \text{Cov}$ 有
$$
\check{\mathcal{C}}^\bullet(\mathcal{U}, \colim_i \mathcal{I}_i)
=
\colim_i \check{\mathcal{C}}^\bullet(\mathcal{U}, \mathcal{I}_i)
$$
因为所有多重交 $U_{j_0 \ldots j_p}$ 都拟紧。由引理
\ref{cohomology-lemma-injective-trivial-cech}，{\v C}ech 复形余极限中的每个复形
在次数 $\geq 1$ 都无环。因此由《代数》引理
\ref{algebra-lemma-directed-colimit-exact}，{\v C}ech 复形
$\check{\mathcal{C}}^\bullet(\mathcal{U}, \colim_i \mathcal{I}_i)$
在次数 $\geq 1$ 也无环。换言之，对所有 $p \geq 1$，有
$\check{H}^p(\mathcal{U},  \colim_i \mathcal{I}_i) = 0$。于是满足引理
\ref{cohomology-lemma-cech-vanish-basis} 的假设，断言得证。
\end{proof}

\begin{lemma}
\label{cohomology-lemma-higher-direct-image-colimit}
设 $f : X \to Y$ 为拓扑空间的连续映射。设
$(\mathcal{F}_i, \varphi_{ii'})$ 为 $X$ 上的阿贝尔层系统。置
$\mathcal{F} = \colim \mathcal{F}_i$。设 $p \geq 0$ 为整数。假设使得
$H^p(f^{-1}(V), \mathcal{F}) = \colim H^p(f^{-1}(V), \mathcal{F}_i)$
成立的开集 $V \subset Y$ 组成 $Y$ 的一个拓扑基。则
$R^pf_*\mathcal{F} = \colim R^pf_*\mathcal{F}_i$.
\end{lemma}

\begin{proof}
回忆 $R^pf_*\mathcal{F}$ 是预层 $\mathcal{G}$ 的层化，其中
$\mathcal{G}$ 将 $V$ 映为 $H^p(f^{-1}(V), \mathcal{F})$；见引理
\ref{cohomology-lemma-describe-higher-direct-images}。类似地，
$R^pf_*\mathcal{F}_i$ 是预层 $\mathcal{G}_i$ 的层化，其中
$\mathcal{G}_i$ 将 $V$ 映为 $H^p(f^{-1}(V), \mathcal{F}_i)$。
回忆层化函子是从层到预层的包含函子的左伴随，见《层》第
\ref{sheaves-section-sheafification} 节。因此层化与余极限可交换，
见《范畴》引理 \ref{categories-lemma-adjoint-exact}。故只需证明预层态射
（其中余极限在预层范畴中取）
$$
\colim \mathcal{G}_i \longrightarrow \mathcal{G}
$$
在层化后诱导同构。为此，只需证明预层 $\mathcal{G}$ 与
$\colim \mathcal{G}_i$ 在 $Y$ 的一个拓扑基上一致。事实上，此时两者
层化后的茎一致，因为这些茎可直接由预层在该基的元素上的取值计算。
所需的一致性恰为引理的假设。
\end{proof}

\noindent
下面给出《层》引理 \ref{sheaves-lemma-descend-opens} 的上同调版本。
设 $X$ 为谱拓扑空间，它写成谱拓扑空间 $X_i$ 的余滤过极限，相应图的
转移态射均为谱态射，如《拓扑》引理
\ref{topology-lemma-directed-inverse-limit-spectral-spaces} 所述。假设给定
\begin{enumerate}
\item 对每个 $i \in \Ob(\mathcal{I})$，给定 $X_i$ 上的阿贝尔层
$\mathcal{F}_i$；
\item 对每个 $a : j \to i$，给定阿贝尔层的一个 $f_a$-映射
$\varphi_a : \mathcal{F}_i \to \mathcal{F}_j$（见《层》定义
\ref{sheaves-definition-f-map}）。
\end{enumerate}
并且只要 $c = a \circ b$，就有
$\varphi_c = \varphi_b \circ \varphi_a$。在 $X$ 上置
$\mathcal{F} = \colim p_i^{-1}\mathcal{F}_i$。

\begin{lemma}
\label{cohomology-lemma-colimit}
在上述情形下，设 $i \in \Ob(\mathcal{I})$，并设
$U_i \subset X_i$ 为拟紧开集。则
$$
\colim_{a : j \to i} H^p(f_a^{-1}(U_i), \mathcal{F}_j) =
H^p(p_i^{-1}(U_i), \mathcal{F})
$$
对所有 $p \geq 0$ 成立。特别地，有
$H^p(X, \mathcal{F}) = \colim H^p(X_i, \mathcal{F}_i)$.
\end{lemma}

\begin{proof}
当 $p = 0$ 时，这就是《层》引理 \ref{sheaves-lemma-descend-opens}。

\medskip\noindent
本段证明可以找到系统态射
$(\gamma_i) : (\mathcal{F}_i, \varphi_a) \to (\mathcal{G}_i, \psi_a)$，
其中 $\mathcal{G}_i$ 为内射阿贝尔层，且 $\gamma_i$ 为单射。对每个
$i$，选取单射 $\mathcal{F}_i \to \mathcal{I}_i$，其中
$\mathcal{I}_i$ 是 $X_i$ 上的内射阿贝尔层。于是可考虑映射族
$$
\gamma_i :
\mathcal{F}_i
\longrightarrow
\prod\nolimits_{b : k \to i} f_{b, *}\mathcal{I}_k = \mathcal{G}_i
$$
其分量映射是映射
$f_b^{-1}\mathcal{F}_i \to \mathcal{F}_k \to \mathcal{I}_k$
的伴随映射。对 $\mathcal{I}$ 中的 $a : j \to i$，有典范映射
$$
\psi_a : f_a^{-1}\mathcal{G}_i \to \mathcal{G}_j
$$
其对应于 $b : k \to j$ 的分量是典范映射
$f_b^{-1}f_{a \circ b, *}\mathcal{I}_k \to f_{b, *}\mathcal{I}_k$。
因而得到系统的单射
$\{\gamma_i\} : \{\mathcal{F}_i, \varphi_a) \to (\mathcal{G}_i, \psi_a)$
。注意 $\mathcal{G}_i$ 是 $X_i$ 上的内射阿贝尔群层，见引理
\ref{cohomology-lemma-pushforward-injective-flat} 与《同调代数》引理
\ref{homology-lemma-product-injectives}。构造至此完成。

\medskip\noindent
完全仿照引理 \ref{cohomology-lemma-quasi-separated-cohomology-colimit} 的证明，
可知只需证明当 $p > 0$ 时
$H^p(X, \colim f_i^{-1}\mathcal{G}_i) = 0$。

\medskip\noindent
置 $\mathcal{G} = \colim f_i^{-1}\mathcal{G}_i$。为证明
$\mathcal{G}$ 在 $X$ 的每个拟紧开集上的上同调消失，只需证明：
对任意由有限个拟紧开集覆盖 $X$ 的一个拟紧开集的开覆盖
$\mathcal{U}$，$\mathcal{G}$ 的 {\v C}ech 上同调为零；见引理
\ref{cohomology-lemma-cech-vanish-basis}。由《拓扑》引理
\ref{topology-lemma-descend-opens}，这样的覆盖是某个 $X_i$ 上同类覆盖
$\mathcal{U}_i$ 在 $p_i$ 下的原像。由 $p = 0$ 的情形，有
$$
\check{\mathcal{C}}^\bullet(\mathcal{U}, \mathcal{G}) =
\colim_{a : j \to i}
\check{\mathcal{C}}^\bullet(f_a^{-1}(\mathcal{U}_i), \mathcal{G}_j)
$$
右端是复形的滤过余极限，而由引理
\ref{cohomology-lemma-injective-trivial-cech}，其中每个复形在正次数都无环。
因此由《代数》引理 \ref{algebra-lemma-directed-colimit-exact} 得到结论。
\end{proof}









\section{诺特拓扑空间上的消失}
\label{cohomology-section-vanishing-Noetherian}

\noindent
本节旨在证明 Grothendieck 的一个定理，即命题
\ref{cohomology-proposition-vanishing-Noetherian}。见 \cite{Tohoku}。

\begin{lemma}
\label{cohomology-lemma-cohomology-and-closed-immersions}
设 $i : Z \to X$ 为拓扑空间的闭浸入。
对 $Z$ 上任意阿贝尔层 $\mathcal{F}$，有
$H^p(Z, \mathcal{F}) = H^p(X, i_*\mathcal{F})$.
\end{lemma}

\begin{proof}
这是因为 $i_*$ 正合（见《模层》引理
\ref{modules-lemma-i-star-exact}），因而作为函子有
$R^pi_* = 0$（《导出范畴》引理
\ref{derived-lemma-right-derived-exact-functor}）。
故可应用引理 \ref{cohomology-lemma-apply-Leray}。
\end{proof}

\begin{lemma}
\label{cohomology-lemma-irreducible-constant-cohomology-zero}
设 $X$ 为不可约拓扑空间。则对任意阿贝尔群 $A$
及所有 $p > 0$，有 $H^p(X, \underline{A}) = 0$。
\end{lemma}

\begin{proof}
回忆 $\underline{A}$ 是《层》定义
\ref{sheaves-definition-constant-sheaf} 中所定义的常值层。
由于 $X$ 不可约，任一非空开集 $U$ 都不可约，因而尤其连通。
故对任意非空开集 $U \subset X$，有 $\underline{A}(U) = A$。
又有 $\underline{A}(\emptyset) = 0$。因此 $\underline{A}$
是 $X$ 上的松弛阿贝尔层，所述消失性由引理
\ref{cohomology-lemma-flasque-acyclic} 得到。
\end{proof}

\begin{lemma}
\label{cohomology-lemma-subsheaf-of-constant-sheaf}
\begin{reference}
\cite[Page 168]{Tohoku}.
\end{reference}
设 $X$ 为一个拓扑空间，且任意两个拟紧开集之交仍拟紧。设
$\mathcal{F} \subset \underline{\mathbf{Z}}$
为由拟紧开集上的有限多个截面生成的子层。则存在由阿贝尔子层组成的有限滤过
$$
(0) = \mathcal{F}_0 \subset \mathcal{F}_1 \subset \ldots \subset
\mathcal{F}_n = \mathcal{F}
$$
使得对每个 $0 < i \leq n$，都存在短正合列
$$
0 \to j'_!\underline{\mathbf{Z}}_V \to j_!\underline{\mathbf{Z}}_U \to
\mathcal{F}_i/\mathcal{F}_{i - 1} \to 0
$$
其中 $j : U \to X$ 与 $j' : V \to X$ 是拟紧开集到 $X$ 的包含映射。
\end{lemma}

\begin{proof}
设 $\mathcal{F}$ 由拟紧开集 $U_1, \ldots, U_t$ 上的截面
$s_1, \ldots, s_t$ 生成。由于 $U_i$ 拟紧且 $s_i$ 是到
$\mathbf{Z}$ 的局部常值函数，必要时可将 $U_i$ 换成其某个有限开闭分解的各部分，
从而不妨设 $s_i$ 为常值截面。
设 $s_i = n_i$，其中 $n_i \in \mathbf{Z}$。若 $n_i = 0$，
当然可从表中删去 $(U_i, n_i)$。必要时变号后，还可假设 $n_i > 0$。
其次，对任意子集 $I \subset \{1, \ldots, t\}$，可把
$\bigcap_{i \in I} U_i$ 与 $\gcd(n_i, i \in I)$ 加入表中。
这样处理后，表 $(U_1, n_1), \ldots, (U_t, n_t)$ 具有如下性质：
对 $x \in X$，令 $I_x = \{i \in \{1, \ldots, t\} \mid x \in U_i\}$。
则对某个 $i \in I_x$，$\gcd(n_i, i \in I_x)$ 等于 $n_i$。

\medskip\noindent
取 $\mathcal{F}_0 = (0)$，并令 $\mathcal{F}_n$ 由所有满足
$n_i \leq n$ 的 $i$ 所对应的 $U_i$ 上截面 $n_i$ 生成，以此作为所需滤过。
显然当 $n \gg 0$ 时 $\mathcal{F}_n = \mathcal{F}$。此外，商
$\mathcal{F}_n/\mathcal{F}_{n - 1}$ 由
$U = \bigcup_{n_i \leq n} U_i$ 上的截面 $n$ 生成，而映射
$j_!\underline{\mathbf{Z}}_U \to \mathcal{F}_n/\mathcal{F}_{n - 1}$
的核由 $V = \bigcup_{n_i \leq n - 1} U_i$ 上的截面 $n$ 生成。
因而得到引理陈述中的短正合列。
\end{proof}

\begin{lemma}
\label{cohomology-lemma-vanishing-generated-one-section}
\begin{reference}
这是 \cite[Proposition 3.6.1]{Tohoku} 的一个特例。
\end{reference}
设 $X$ 为拓扑空间，$d \geq 0$ 为整数。假设
\begin{enumerate}
\item $X$ 拟紧；
\item 拟紧开集构成 $X$ 的一组基；
\item 两个拟紧开集之交拟紧；
\item 对所有 $p > d$ 以及任意拟紧开集的包含映射
$j : U \to X$，有 $H^p(X, j_!\underline{\mathbf{Z}}_U) = 0$。
\end{enumerate}
则对 $X$ 上任意阿贝尔层 $\mathcal{F}$ 及所有 $p > d$，
有 $H^p(X, \mathcal{F}) = 0$。
\end{lemma}

\begin{proof}
令 $S = \coprod_{U \subset X} \mathcal{F}(U)$，其中 $U$ 遍历
$X$ 的拟紧开集。对任一有限子集
$A = \{s_1, \ldots, s_n\} \subset S$，令 $\mathcal{F}_A$
为由所有 $s_i$ 生成的 $\mathcal{F}$ 的子层（见《模层》定义
\ref{modules-definition-generated-by-local-sections}）。
注意，若 $A \subset A'$，则 $\mathcal{F}_A \subset \mathcal{F}_{A'}$。
因此 $\{\mathcal{F}_A\}$ 构成以 $S$ 的有限子集所成有向偏序集为指标的系统。
由《模层》引理 \ref{modules-lemma-generated-by-local-sections-stalk}，
观察各茎即知
$$
\colim_A \mathcal{F}_A = \mathcal{F}
$$
由引理 \ref{cohomology-lemma-quasi-separated-cohomology-colimit}，有
$$
H^p(X, \mathcal{F}) =
\colim_A H^p(X, \mathcal{F}_A)
$$
因此只需对阿贝尔层 $\mathcal{F}_A$ 证明所述消失性。换言之，
只需处理 $\mathcal{F}$ 由 $X$ 的拟紧开集上的有限多个局部截面生成的情形。

\medskip\noindent
设 $\mathcal{F}$ 由局部截面 $s_1, \ldots, s_n$ 生成，
并令 $\mathcal{F}' \subset \mathcal{F}$ 为由
$s_1, \ldots, s_{n - 1}$ 生成的子层。于是有短正合列
$$
0 \to \mathcal{F}' \to \mathcal{F} \to \mathcal{F}/\mathcal{F}' \to 0
$$
由上同调长正合列可知，只需对由少于 $n$ 个局部截面生成的阿贝尔层
$\mathcal{F}'$ 与 $\mathcal{F}/\mathcal{F}'$ 证明消失性。
故只需对至多由一个局部截面生成的层证明消失性。这样的层恰好是
$j_!\underline{\mathbf{Z}}_U$ 的商，其中 $U$ 是 $X$ 的拟紧开集。

\medskip\noindent
现设有短正合列
$$
0 \to \mathcal{K} \to j_!\underline{\mathbf{Z}}_U \to \mathcal{F} \to 0
$$
其中 $U$ 是 $X$ 的拟紧开集。只需证明当 $q \geq d + 1$ 时
$H^q(X, \mathcal{K})$ 为零。与上面一样，可将 $\mathcal{K}$
写成由拟紧开集上的有限多个截面生成的子层 $\mathcal{K}'$ 的滤过余极限。
于是 $\mathcal{F}$ 是各层
$j_!\underline{\mathbf{Z}}_U/\mathcal{K}'$ 的滤过余极限。
由此约化到 $\mathcal{K}$ 由拟紧开集上的有限多个截面生成的情形。
注意 $\mathcal{K}$ 是 $\underline{\mathbf{Z}}_X$ 的子层。
故由引理 \ref{cohomology-lemma-subsheaf-of-constant-sheaf}，$\mathcal{K}$
存在有限滤过，其逐次商 $\mathcal{Q}$ 嵌入短正合列
$$
0 \to j''_!\underline{\mathbf{Z}}_W \to
j'_!\underline{\mathbf{Z}}_V \to \mathcal{Q} \to 0
$$
其中 $j'' : W \to X$ 与 $j' : V \to X$ 是拟紧开集的包含映射。
因此，由本引理关于 $j''_!\underline{\mathbf{Z}}_W$
与 $j'_!\underline{\mathbf{Z}}_V$ 的上同调群消失的假设，
当 $p > d$ 时有 $H^p(X, \mathcal{Q}) = 0$。
回到 $\mathcal{K}$，沿上述滤过运用上同调长正合列归纳，
也得到所需的消失性。
\end{proof}

\begin{example}
\label{cohomology-example-datta}
令 $X = \mathbf{N}$，赋予它如下拓扑：开集为 $\emptyset$、$X$
以及 $n \geq 1$ 时的 $U_n = \{i \mid i \leq n\}$。
$X$ 上的阿贝尔层 $\mathcal{F}$ 等同于阿贝尔群逆系统
$A_n = \mathcal{F}(U_n)$，并且
$\Gamma(X, \mathcal{F}) = \lim A_n$。由于逆极限函子在逆系统范畴上不正合，
可知存在 $H^1$ 非零的阿贝尔层。
最后，读者可验证：若 $j : U = U_n \to X$ 为包含映射，则
$p \geq 1$ 时 $H^p(X, j_!\mathbf{Z}_U) = 0$。
因此，$X$ 是一个在 $d = 0$ 时满足引理
\ref{cohomology-lemma-vanishing-generated-one-section} 的条件 (2)、(3)、(4)
却不满足结论的空间。
\end{example}

\begin{lemma}
\label{cohomology-lemma-subsheaf-irreducible}
设 $X$ 为不可约拓扑空间，并设
$\mathcal{H} \subset \underline{\mathbf{Z}}$ 为该常值层的阿贝尔子层。
则存在非空开集 $U \subset X$，使得对某个 $d \in \mathbf{Z}$ 有
$\mathcal{H}|_U = \underline{d\mathbf{Z}}_U$。
\end{lemma}

\begin{proof}
回忆，对 $X$ 的任意非空开集 $V$，有
$\underline{\mathbf{Z}}(V) = \mathbf{Z}$（见引理
\ref{cohomology-lemma-irreducible-constant-cohomology-zero} 的证明）。
若 $\mathcal{H} = 0$，则取 $d = 0$ 即得引理。
若 $\mathcal{H} \not = 0$，则存在非空开集 $U \subset X$，
使得 $\mathcal{H}(U) \not = 0$。设对某个 $n \geq 1$ 有
$\mathcal{H}(U) = n\mathbf{Z}$。于是
$\underline{n\mathbf{Z}}_U
\subset \mathcal{H}|_U \subset
\underline{\mathbf{Z}}_U$。若第一个包含严格，则可找到非空开集
$U' \subset U$ 与整数 $1 \leq n' < n$，使得
$\underline{n'\mathbf{Z}}_{U'}
\subset \mathcal{H}|_{U'} \subset
\underline{\mathbf{Z}}_{U'}$。
此过程必在有限步后停止，因而得到引理。
\end{proof}

\begin{proposition}[Grothendieck]
\label{cohomology-proposition-vanishing-Noetherian}
\begin{reference}
\cite[Theorem 3.6.5]{Tohoku}.
\end{reference}
设 $X$ 为诺特拓扑空间。若 $\dim(X) \leq d$，则对 $X$ 上任意阿贝尔层
$\mathcal{F}$ 及所有 $p > d$，有 $H^p(X, \mathcal{F}) = 0$。
\end{proposition}

\begin{proof}
对 $d$ 归纳证明此引理。固定 $d$，并假设对所有维数 $< d$
的诺特拓扑空间，此引理均成立。

\medskip\noindent
设 $\mathcal{F}$ 为 $X$ 上的阿贝尔层。设 $U \subset X$ 为开集，
并以 $Z \subset X$ 表示其闭补集。以
$j : U \to X$ 与 $i : Z \to X$ 表示包含映射。于是有短正合列
$$
0 \to j_{!}j^*\mathcal{F} \to \mathcal{F} \to i_*i^*\mathcal{F} \to 0
$$
见《模层》引理 \ref{modules-lemma-canonical-exact-sequence}。
注意 $j_!j^*\mathcal{F}$ 支撑在 $U$ 的拓扑闭包 $Z'$ 上；
换言之，它形如 $i'_*\mathcal{F}'$，其中 $\mathcal{F}'$
是 $Z'$ 上的某个阿贝尔层，而 $i' : Z' \to X$ 是包含映射。

\medskip\noindent
利用这一点可约化到 $X$ 不可约的情形。事实上，由《拓扑》引理
\ref{topology-lemma-Noetherian}，$X$ 只有有限多个不可约分支。
若 $X$ 有多于一个不可约分支，取 $X$ 的一个不可约分支
$Z \subset X$，并令 $U = X \setminus Z$。由上述结论及上同调长正合列，
只需证明当 $p > d$ 时
$H^p(X, i_*i^*\mathcal{F})$ 与 $H^p(X, i'_*\mathcal{F}')$ 消失。
由引理 \ref{cohomology-lemma-cohomology-and-closed-immersions}，又只需证明当
$p > d$ 时 $H^p(Z, i^*\mathcal{F})$ 与
$H^p(Z', \mathcal{F}')$ 消失。由于 $Z'$ 和 $Z$
的不可约分支都更少，确实约化到了 $X$ 不可约的情形。

\medskip\noindent
若 $d = 0$ 且 $X$ 不可约，则 $X$ 是 $X$ 的唯一非空开子集。
因此每个层都是常值层，高阶上同调群消失（例如由引理
\ref{cohomology-lemma-irreducible-constant-cohomology-zero}）。

\medskip\noindent
设 $X$ 不可约且维数为 $d > 0$。由引理
\ref{cohomology-lemma-vanishing-generated-one-section}，可约化到
$\mathcal{F} = j_!\underline{\mathbf{Z}}_U$ 的情形，其中
$U \subset X$ 为某个开集。此时考察短正合列
$$
0 \to j_!(\underline{\mathbf{Z}}_U) \to
\underline{\mathbf{Z}}_X \to i_*\underline{\mathbf{Z}}_Z \to 0
$$
其中 $Z = X \setminus U$。由引理
\ref{cohomology-lemma-irreducible-constant-cohomology-zero}，对所有 $p \geq 1$，
$H^p(X, \underline{\mathbf{Z}}_X)$ 消失。由归纳假设，
$H^p(X, i_*\underline{\mathbf{Z}}_Z) = H^p(Z, \underline{\mathbf{Z}}_Z) = 0$
对 $p \geq d$ 成立。故由上同调长正合列即得结论。
\end{proof}





\section{支撑在闭子集上的上同调}
\label{cohomology-section-cohomology-support}

\noindent
本节只讨论最基本的内容；相关讨论将在第
\ref{cohomology-section-cohomology-support-bis} 节继续。

\medskip\noindent
设 $X$ 为拓扑空间，$Z \subset X$ 为闭子集，并设
$\mathcal{F}$ 为 $X$ 上的阿贝尔层。令
$$
\Gamma_Z(X, \mathcal{F}) =
\{s \in \mathcal{F}(X) \mid \text{Supp}(s) \subset Z\}
$$
为支集包含在 $Z$ 中的截面所成子集。截面的支集定义见
《模层》定义 \ref{modules-definition-support}。《模层》引理
\ref{modules-lemma-support-section-closed} 蕴含
$\Gamma_Z(X, \mathcal{F})$ 是 $\Gamma(X, \mathcal{F})$ 的子群。
同一引理还保证赋值
$\mathcal{F} \mapsto \Gamma_Z(X, \mathcal{F})$
关于 $\mathcal{F}$ 构成函子。此函子左正合，但一般并不正合。

\medskip\noindent
由于阿贝尔层范畴有足够多的内射对象（《内射对象》引理
\ref{injectives-lemma-abelian-sheaves-space}），可得到右导出函子
$$
R\Gamma_Z(X, -) : D^+(X) \longrightarrow D^+(\textit{Ab})
$$
见《导出范畴》引理
\ref{derived-lemma-enough-injectives-right-derived}。
为计算 $R\Gamma_Z(X, -)$ 在对象 $K$ 上的值，用内射阿贝尔层的下有界复形
$\mathcal{I}^\bullet$ 表示 $K$，再取
$\Gamma_Z(X, \mathcal{I}^\bullet)$；见《导出范畴》引理
\ref{derived-lemma-injective-acyclic}。阿贝尔层 $\mathcal{F}$
支撑在 $Z$ 上的上同调群定义为
$H^q_Z(X, \mathcal{F}) = R^q\Gamma_Z(X, \mathcal{F})$。

\medskip\noindent
设 $\mathcal{I}$ 为 $X$ 上的内射阿贝尔层，并令
$U = X \setminus Z$。则限制映射
$\mathcal{I}(X) \to \mathcal{I}(U)$ 是满射（引理
\ref{cohomology-lemma-injective-restriction-surjective}），其核为
$\Gamma_Z(X, \mathcal{I})$。由此立即得到：对 $K \in D^+(X)$，
$D^+(\textit{Ab})$ 中存在特异三角
$$
R\Gamma_Z(X, K) \to R\Gamma(X, K) \to R\Gamma(U, K) \to R\Gamma_Z(X, K)[1]
$$
因而得到上同调长正合列
$$
\ldots \to H^i_Z(X, K) \to H^i(X, K) \to H^i(U, K) \to
H^{i + 1}_Z(X, K) \to \ldots
$$
它对 $D^+(X)$ 中任意 $K$ 都成立。

\medskip\noindent
对 $X$ 上的阿贝尔层 $\mathcal{F}$，可考虑{\it 支撑在 $Z$ 上的截面子层}，
记为 $\mathcal{H}_Z(\mathcal{F})$，其定义为
$$
\mathcal{H}_Z(\mathcal{F})(U) =
\{s \in \mathcal{F}(U) \mid \text{Supp}(s) \subset U \cap Z\} =
\Gamma_{Z \cap U}(U, \mathcal{F}|_U)
$$
利用《模层》引理 \ref{modules-lemma-i-star-exact} 中的等价，
可将 $\mathcal{H}_Z(\mathcal{F})$ 看作 $Z$ 上的阿贝尔层；见《模层》注
\ref{modules-remark-sections-support-in-closed}。由此得到函子
$$
\textit{Ab}(X) \longrightarrow \textit{Ab}(Z),\quad
\mathcal{F} \longmapsto
\mathcal{H}_Z(\mathcal{F})\text{（看作 }Z\text{ 上的层）}
$$
此函子左正合，但一般并不正合。与上面完全相同，可得到右导出函子
$$
R\mathcal{H}_Z : D^+(X) \longrightarrow D^+(Z)
$$
并置 $\mathcal{H}^q_Z(\mathcal{F}) = R^q\mathcal{H}_Z(\mathcal{F})$，
于是 $\mathcal{H}^0_Z(\mathcal{F}) = \mathcal{H}_Z(\mathcal{F})$。

\medskip\noindent
注意，对任意阿贝尔层 $\mathcal{F}$，有
$\Gamma_Z(X, \mathcal{F}) = \Gamma(Z, \mathcal{H}_Z(\mathcal{F}))$
由下面的引理 \ref{cohomology-lemma-sections-with-support-acyclic}，
函子 $\mathcal{H}_Z$ 将内射阿贝尔层变为关于
$\Gamma(Z, -)$ 右无环的层。因此由《导出范畴》引理
\ref{derived-lemma-grothendieck-spectral-sequence}，得到收敛的
Grothendieck 谱序列
$$
E_2^{p, q} = H^p(Z, \mathcal{H}^q_Z(K)) \Rightarrow H^{p + q}_Z(X, K)
$$
，它关于 $D^+(X)$ 中的 $K$ 具有函子性。

\begin{lemma}
\label{cohomology-lemma-sections-with-support-acyclic}
设 $i : Z \to X$ 为闭子集的包含映射，并设 $\mathcal{I}$
为 $X$ 上的内射阿贝尔层。则
$\mathcal{H}_Z(\mathcal{I})$ 是 $Z$ 上的内射阿贝尔层。
\end{lemma}

\begin{proof}
由于 $\mathcal{H}_Z(-)$ 是正合函子 $i_*$ 的右伴随，
结论由《同调代数》引理
\ref{homology-lemma-adjoint-preserve-injectives} 得到。
另见《模层》引理 \ref{modules-lemma-i-star-exact} 与
\ref{modules-lemma-i-star-right-adjoint}。
\end{proof}




\section{谱拓扑空间上的上同调}
\label{cohomology-section-spectral}

\noindent
谱拓扑空间上同调的一个关键结果是引理 \ref{cohomology-lemma-colimit}。
粗略地说，它断言上同调与谱拓扑空间范畴中的余滤过极限可交换；
该范畴定义见《拓扑》定义 \ref{topology-definition-spectral-space}。
应用这一结果，可如下得到引理 \ref{cohomology-lemma-cohomology-of-closed}
与 \ref{cohomology-lemma-proper-base-change} 的类似结论。

\begin{lemma}
\label{cohomology-lemma-cohomology-of-neighbourhoods-of-closed}
设 $X$ 为谱拓扑空间，$\mathcal{F}$ 为 $X$ 上的阿贝尔层，
$E \subset X$ 为拟紧子集。令 $W \subset X$ 为 $X$
中所有特化到 $E$ 中某点的点所成集合。
\begin{enumerate}
\item $H^p(W, \mathcal{F}|_W) = \colim H^p(U, \mathcal{F})$
，其中余极限遍历 $E$ 的拟紧开邻域；
\item $H^p(W \setminus E, \mathcal{F}|_{W \setminus E}) =
\colim H^p(U \setminus E, \mathcal{F}|_{U \setminus E})$
，只要 $E$ 是可构造子集。
\end{enumerate}
\end{lemma}

\begin{proof}
由《拓扑》引理 \ref{topology-lemma-make-spectral-space}，
$W = \lim U$，其中极限遍历包含 $E$ 的拟紧开集。由《拓扑》引理
\ref{topology-lemma-spectral-sub}，每个 $U$ 都是谱拓扑空间。
故可应用引理 \ref{cohomology-lemma-colimit}，得到 (1)。
对 (2) 作相同证明，只需改用《拓扑》引理
\ref{topology-lemma-make-spectral-space-minus}。
\end{proof}

\begin{lemma}
\label{cohomology-lemma-proper-base-change-spectral}
设 $f : X \to Y$ 为谱拓扑空间之间的谱态射，$y \in Y$。
令 $E \subset Y$ 为所有特化到 $y$ 的点所成集合，并设
$\mathcal{F}$ 为 $X$ 上的阿贝尔层。则
$(R^pf_*\mathcal{F})_y = H^p(f^{-1}(E), \mathcal{F}|_{f^{-1}(E)})$。
\end{lemma}

\begin{proof}
注意 $E = \bigcap V$，其中 $V$ 遍历 $Y$ 中 $y$ 的拟紧开邻域。
故 $f^{-1}(E) = \bigcap f^{-1}(V)$，这又蕴含作为拓扑空间有
$f^{-1}(E) = \lim f^{-1}(V)$。由于 $f$ 是谱态射，每个
$f^{-1}(V)$ 也都是谱拓扑空间（《拓扑》引理
\ref{topology-lemma-spectral-sub}）。因此 $f^{-1}(E)$ 是谱拓扑空间，
并且由引理 \ref{cohomology-lemma-colimit} 有
$$
H^p(f^{-1}(E), \mathcal{F}|_{f^{-1}(E)}) =
\colim H^p(f^{-1}(V), \mathcal{F})
$$
另一方面，$R^pf_*\mathcal{F}$ 在 $y$ 处的茎正由右端的余极限给出。
\end{proof}

\begin{lemma}
\label{cohomology-lemma-vanishing-for-profinite}
设 $X$ 为预有限拓扑空间。则对任意阿贝尔层 $\mathcal{F}$
及所有 $q > 0$，有 $H^q(X, \mathcal{F}) = 0$。
\end{lemma}

\begin{proof}
$X$ 的任意开覆盖都可由一个以开集为各部分的有限不交并分解加细；
见《拓扑》引理
\ref{topology-lemma-profinite-refine-open-covering}。
因此，若 $\mathcal{F} \to \mathcal{G}$ 是 $X$ 上阿贝尔层的满射，
则 $\mathcal{F}(X) \to \mathcal{G}(X)$ 为满射。换言之，整体截面函子正合。
故其高阶导出函子均为零；见《导出范畴》引理
\ref{derived-lemma-right-derived-exact-functor}。
\end{proof}

\noindent
下面关于上同调消失的结果改进了 Grothendieck 的结果
（命题 \ref{cohomology-proposition-vanishing-Noetherian}），可见于
\cite{Scheiderer}。

\begin{proposition}
\label{cohomology-proposition-cohomological-dimension-spectral}
\begin{reference}
第 (1) 部分是 \cite{Scheiderer} 的主定理。
\end{reference}
设 $X$ 为 Krull 维数为 $d$ 的谱拓扑空间，$\mathcal{F}$
为 $X$ 上的阿贝尔层。
\begin{enumerate}
\item 当 $q > d$ 时，$H^q(X, \mathcal{F}) = 0$；
\item 对任意拟紧开集 $U \subset X$，映射
$H^d(X, \mathcal{F}) \to H^d(U, \mathcal{F})$ 为满射；
\item 对任意可构造闭子集 $Z \subset X$，当 $q > d$ 时，
$H^q_Z(X, \mathcal{F}) = 0$。
\end{enumerate}
\end{proposition}

\begin{proof}
对 $d$ 归纳证明。

\medskip\noindent
若 $d = 0$，则 $X$ 是预有限空间；见《拓扑》引理
\ref{topology-lemma-characterize-profinite-spectral}。
故 (1) 由引理 \ref{cohomology-lemma-vanishing-for-profinite} 成立。
若 $U \subset X$ 是拟紧开集，则作为 Hausdorff 空间的拟紧子集，
$U$ 也是闭集。因此作为拓扑空间有
$X = U \amalg (X \setminus U)$，从而 (2) 成立。
对 (3) 中的 $Z$，考察长正合列
$$
H^{q - 1}(X, \mathcal{F}) \to
H^{q - 1}(X \setminus Z, \mathcal{F}) \to
H^q_Z(X, \mathcal{F}) \to H^q(X, \mathcal{F})
$$
由于 $X$ 与 $U = X \setminus Z$ 都是预有限空间（确切地说，
因 $Z$ 可构造，$U$ 拟紧），再结合 (2) 与 (1)，便得到支撑在
$Z$ 上的上同调群所需的消失性。

\medskip\noindent
以下作归纳步骤。设 $d \geq 1$，并假设本命题对所有维数 $< d$
的谱拓扑空间成立。先证明 $X$ 的 (2)。设 $U$ 为拟紧开集，
$\xi \in H^d(U, \mathcal{F})$。置 $Z = X \setminus U$，
并令 $W \subset X$ 为所有特化到 $Z$ 的点所成集合。
由引理 \ref{cohomology-lemma-cohomology-of-neighbourhoods-of-closed}，有
$$
H^d(W \setminus Z, \mathcal{F}|_{W \setminus Z}) =
\colim_{Z \subset V} H^d(V \setminus Z, \mathcal{F})
$$
其中余极限遍历 $Z$ 在 $X$ 中的拟紧开邻域 $V$。
由《拓扑》引理 \ref{topology-lemma-make-spectral-space}，
$W \setminus Z$ 是谱拓扑空间。由于 $W$ 的每一点都特化到 $Z$
中的某点，$W \setminus Z$ 是 Krull 维数 $< d$ 的谱拓扑空间。
由归纳假设，$\xi$ 在
$H^d(W \setminus Z, \mathcal{F}|_{W \setminus Z})$ 中的像为零。
由上述显示式，存在拟紧开集 $Z \subset V \subset X$，使得
$\xi|_{V \setminus Z} = 0$。由于 $V \setminus Z = V \cap U$，
对覆盖 $X = U \cup V$ 应用 Mayer--Vietoris（引理
\ref{cohomology-lemma-mayer-vietoris}），可知存在
$\tilde \xi \in H^d(X, \mathcal{F})$，其在 $U$ 上限制为 $\xi$，
在 $V$ 上限制为零。换言之，(2) 成立。

\medskip\noindent
假设 (2) 成立，证明 (1)。选取内射消解
$\mathcal{F} \to \mathcal{I}^\bullet$，并置
$$
\mathcal{G} = \Im(\mathcal{I}^{d - 1} \to \mathcal{I}^d) =
\Ker(\mathcal{I}^d \to \mathcal{I}^{d + 1})
$$
对拟紧开集 $U \subset X$，有如下正合列之间的映射：
$$
\xymatrix{
\mathcal{I}^{d - 1}(X) \ar[r] \ar[d] &
\mathcal{G}(X) \ar[r] \ar[d] &
H^d(X, \mathcal{F}) \ar[d] \ar[r] & 0 \\
\mathcal{I}^{d - 1}(U) \ar[r] &
\mathcal{G}(U) \ar[r] &
H^d(U, \mathcal{F}) \ar[r] & 0
}
$$
由引理 \ref{cohomology-lemma-injective-flasque} 及 $d \geq 1$，
层 $\mathcal{I}^{d - 1}$ 松弛。由 (2)，右侧竖直箭头为满射。
作图追踪可知映射 $\mathcal{G}(X) \to \mathcal{G}(U)$ 为满射。
由引理 \ref{cohomology-lemma-vanishing-ravi}，对拟紧开集由拟紧开集组成的任意有限覆盖
$\mathcal{U} : U = U_1 \cup \ldots \cup U_n$ 及任意 $q > 0$，有
$\check{H}^q(\mathcal{U}, \mathcal{G}) = 0$。
应用引理 \ref{cohomology-lemma-cech-vanish-basis}，可知对 $X$
的所有拟紧开集 $U$ 及所有 $q > 0$，有
$H^q(U, \mathcal{G}) = 0$。由 Leray 无环性引理
（《导出范畴》引理 \ref{derived-lemma-leray-acyclicity}），得到
$$
H^q(X, \mathcal{F}) =
H^q\left(
\Gamma(X, \mathcal{I}^0) \to \ldots \to
\Gamma(X, \mathcal{I}^{d - 1}) \to \Gamma(X, \mathcal{G})
\right)
$$
特别地，当 $q > d$ 时该上同调群消失。

\medskip\noindent
最后证明 (3)。对 (3) 中的 $Z$，考察长正合列
$$
H^{q - 1}(X, \mathcal{F}) \to
H^{q - 1}(X \setminus Z, \mathcal{F}) \to
H^q_Z(X, \mathcal{F}) \to H^q(X, \mathcal{F})
$$
由于 $X$ 与 $U = X \setminus Z$ 都是维数 $\leq d$ 的谱拓扑空间
（《拓扑》引理 \ref{topology-lemma-spectral-sub}），再结合 (2) 与 (1)，
便得到所需的消失性。
\end{proof}













\section{交替 {\v C}ech 复形}
\label{cohomology-section-alternating-cech}

\noindent
本节比较 {\v C}ech 复形、交替 {\v C}ech 复形以及若干相关复形。

\medskip\noindent
设 $X$ 为拓扑空间，$\mathcal{U} : U = \bigcup_{i \in I} U_i$
为开覆盖。对 $p \geq 0$，置
$$
\check{\mathcal{C}}_{alt}^p(\mathcal{U}, \mathcal{F})
=
\left\{
\begin{matrix}
s \in  \check{\mathcal{C}}^p(\mathcal{U}, \mathcal{F})
\text{，满足 }
s_{i_0 \ldots i_p} = 0 \text{，若 } i_n = i_m \text{ 对某些 } n \not = m \text{ 成立}\\
\text{且}
s_{i_0\ldots i_n \ldots i_m \ldots i_p}
=
-s_{i_0\ldots i_m \ldots i_n \ldots i_p}
\text{总成立。}
\end{matrix}
\right\}
$$
我们略去如下验证：方程 (\ref{cohomology-equation-d-cech}) 中的微分 $d$ 将
$\check{\mathcal{C}}^p_{alt}(\mathcal{U}, \mathcal{F})$ 映入
$\check{\mathcal{C}}^{p + 1}_{alt}(\mathcal{U}, \mathcal{F})$。

\begin{definition}
\label{cohomology-definition-alternating-cech-complex}
设 $X$ 为拓扑空间，$\mathcal{U} : U = \bigcup_{i \in I} U_i$
为开覆盖，$\mathcal{F}$ 为 $X$ 上的阿贝尔预层。复形
$\check{\mathcal{C}}_{alt}^\bullet(\mathcal{U}, \mathcal{F})$
称为与 $\mathcal{F}$ 及开覆盖 $\mathcal{U}$ 相伴的
{\it 交替 {\v C}ech 复形}。
\end{definition}

\noindent
于是存在复形的典范态射
$$
\check{\mathcal{C}}_{alt}^\bullet(\mathcal{U}, \mathcal{F})
\longrightarrow
\check{\mathcal{C}}^\bullet(\mathcal{U}, \mathcal{F})
$$
即交替 {\v C}ech 复形到通常 {\v C}ech 复形的包含映射。

\medskip\noindent
设覆盖 $\mathcal{U} : U = \bigcup_{i \in I} U_i$ 的指标集 $I$
带有全序 $<$。此时置
$$
\check{\mathcal{C}}_{ord}^p(\mathcal{U}, \mathcal{F})
=
\prod\nolimits_{(i_0, \ldots, i_p) \in I^{p + 1}, i_0 < \ldots < i_p}
\mathcal{F}(U_{i_0\ldots i_p}).
$$
这是一个阿贝尔群。对
$s \in \check{\mathcal{C}}_{ord}^p(\mathcal{U}, \mathcal{F})$，
以 $s_{i_0\ldots i_p}$ 表示它在
$\mathcal{F}(U_{i_0\ldots i_p})$ 中的值。定义
$$
d : \check{\mathcal{C}}_{ord}^p(\mathcal{U}, \mathcal{F})
\longrightarrow
\check{\mathcal{C}}_{ord}^{p + 1}(\mathcal{U}, \mathcal{F})
$$
如下：
$$
d(s)_{i_0\ldots i_{p + 1}}
=
\sum\nolimits_{j = 0}^{p + 1}
(-1)^j
s_{i_0\ldots \hat i_j \ldots i_{p + 1}}|_{U_{i_0\ldots i_{p + 1}}}
$$
其中 $i_0 < \ldots < i_{p + 1}$。注意此公式与方程
(\ref{cohomology-equation-d-cech}) 完全相同。直接可见 $d \circ d = 0$；
换言之，$\check{\mathcal{C}}_{ord}^\bullet(\mathcal{U}, \mathcal{F})$
是复形。

\begin{definition}
\label{cohomology-definition-ordered-cech-complex}
设 $X$ 为拓扑空间，$\mathcal{U} : U = \bigcup_{i \in I} U_i$
为开覆盖，并在 $I$ 上给定全序。设 $\mathcal{F}$ 为 $X$ 上的阿贝尔预层。
复形 $\check{\mathcal{C}}_{ord}^\bullet(\mathcal{U}, \mathcal{F})$
称为与 $\mathcal{F}$、开覆盖 $\mathcal{U}$ 及 $I$ 上给定全序相伴的
{\it 有序 {\v C}ech 复形}。
\end{definition}

\noindent
此复形有时也称为交替 {\v C}ech 复形，因为有序 {\v C}ech 复形
与交替 {\v C}ech 复形之间存在明显的比较映射。具体地，考虑映射
$$
c :
\check{\mathcal{C}}_{ord}^\bullet(\mathcal{U}, \mathcal{F})
\longrightarrow
\check{\mathcal{C}}^\bullet(\mathcal{U}, \mathcal{F})
$$
，定义为
$$
c(s)_{i_0\ldots i_p} =
\left\{
\begin{matrix}
0 &
\text{若} &
i_n = i_m \text{ 对某些 } n \not = m \text{ 成立}\\
\text{sgn}(\sigma) s_{i_{\sigma(0)}\ldots i_{\sigma(p)}} &
\text{若} &
i_{\sigma(0)} < i_{\sigma(1)} < \ldots < i_{\sigma(p)}
\end{matrix}
\right.
$$
其中 $\sigma$ 表示 $\{0, \ldots, p\}$ 的一个置换，
$\text{sgn}(\sigma)$ 表示其符号。文献中常借映射 $c$
将交替与有序 {\v C}ech 复形等同起来。具体地，有如下简单引理。

\begin{lemma}
\label{cohomology-lemma-ordered-alternating}
设 $X$ 为拓扑空间，$\mathcal{U} : U = \bigcup_{i \in I} U_i$
为开覆盖，并设 $I$ 带有全序。映射 $c$ 是复形态射。
事实上，它诱导复形同构
$$
c : \check{\mathcal{C}}_{ord}^\bullet(\mathcal{U}, \mathcal{F})
\to \check{\mathcal{C}}_{alt}^\bullet(\mathcal{U}, \mathcal{F})
$$
\end{lemma}

\begin{proof}
略。
\end{proof}

\noindent
另有映射
$$
\pi :
\check{\mathcal{C}}^\bullet(\mathcal{U}, \mathcal{F})
\longrightarrow
\check{\mathcal{C}}_{ord}^\bullet(\mathcal{U}, \mathcal{F})
$$
，其定义为
$$
\pi(s)_{i_0\ldots i_p} = s_{i_0\ldots i_p}
$$
，其中 $i_0 < i_1 < \ldots < i_p$。

\begin{lemma}
\label{cohomology-lemma-project-to-ordered}
设 $X$ 为拓扑空间，$\mathcal{U} : U = \bigcup_{i \in I} U_i$
为开覆盖，并设 $I$ 带有全序。映射
$\pi : \check{\mathcal{C}}^\bullet(\mathcal{U}, \mathcal{F})
\to \check{\mathcal{C}}_{ord}^\bullet(\mathcal{U}, \mathcal{F})$
是复形态射。它诱导复形同构
$$
\pi : \check{\mathcal{C}}_{alt}^\bullet(\mathcal{U}, \mathcal{F})
\to \check{\mathcal{C}}_{ord}^\bullet(\mathcal{U}, \mathcal{F})
$$
，且该同构是态射 $c$ 的左逆。
\end{lemma}

\begin{proof}
略。
\end{proof}

\begin{remark}
\label{cohomology-remark-compared-ordered-complexes}
这意味着，若指标集 $I$ 上有两个全序 $<_1$ 与 $<_2$，则得到复形同构
$\tau = \pi_2 \circ c_1 :
\check{\mathcal{C}}_{ord\text{-}1}(\mathcal{U}, \mathcal{F}) \to
\check{\mathcal{C}}_{ord\text{-}2}(\mathcal{U}, \mathcal{F})$.
显然
$$
\tau(s)_{i_0 \ldots i_p} =
\text{符号}(\sigma) s_{i_{\sigma(0)} \ldots i_{\sigma(p)}}
$$
，其中 $i_0 <_1 i_1 <_1 \ldots <_1 i_p$，且
$i_{\sigma(0)} <_2 i_{\sigma(1)} <_2 \ldots <_2 i_{\sigma(p)}$。
正是在此意义下，有序 {\v C}ech 复形与所选全序无关。
\end{remark}

\begin{lemma}
\label{cohomology-lemma-alternating-usual}
设 $X$ 为拓扑空间，$\mathcal{U} : U = \bigcup_{i \in I} U_i$
为开覆盖，并设 $I$ 带有全序。映射 $c \circ \pi$
与 $\check{\mathcal{C}}^\bullet(\mathcal{U}, \mathcal{F})$ 上的恒等映射链同伦。
特别地，包含映射
$\check{\mathcal{C}}_{alt}^\bullet(\mathcal{U}, \mathcal{F}) \to
\check{\mathcal{C}}^\bullet(\mathcal{U}, \mathcal{F})$
是链同伦等价。
\end{lemma}

\begin{proof}
对任意多重指标 $(i_0, \ldots, i_p) \in I^{p + 1}$，存在唯一置换
$\sigma : \{0, \ldots, p\} \to \{0, \ldots, p\}$，使得
$$
i_{\sigma(0)} \leq i_{\sigma(1)} \leq \ldots \leq i_{\sigma(p)}
\quad
\text{且}
\quad
\sigma(j) < \sigma(j + 1)
\quad
\text{若}
\quad
i_{\sigma(j)} = i_{\sigma(j + 1)}.
$$
将此置换记为 $\sigma = \sigma^{i_0 \ldots i_p}$。

\medskip\noindent
对任意置换 $\sigma : \{0, \ldots, p\} \to \{0, \ldots, p\}$
及任意 $0 \leq a \leq p$，以 $\sigma_a$ 表示 $\{0, \ldots, p\}$
的唯一置换，使得当 $0 \leq j < a$ 时
$\sigma_a(j) = \sigma(j)$，并且
$\sigma_a(a) < \sigma_a(a + 1) < \ldots < \sigma_a(p)$。
例如，若 $p = 3$ 且 $\sigma,\tau$ 由下式给出：
$$
\begin{matrix}
\text{id} & 0 & 1 & 2 & 3 \\
\sigma & 3 & 2 & 1 & 0
\end{matrix}
\quad \text{以及} \quad
\begin{matrix}
\text{id} & 0 & 1 & 2 & 3 \\
\tau & 3 & 0 & 2 & 1
\end{matrix}
$$
则
$$
\begin{matrix}
\text{id} & 0 & 1 & 2 & 3 \\
\sigma_0 & 0 & 1 & 2 & 3 \\
\sigma_1 & 3 & 0 & 1 & 2 \\
\sigma_2 & 3 & 2 & 0 & 1 \\
\sigma_3 & 3 & 2 & 1 & 0 \\
\end{matrix}
\quad \text{以及} \quad
\begin{matrix}
\text{id} & 0 & 1 & 2 & 3 \\
\tau_0 & 0 & 1 & 2 & 3 \\
\tau_1 & 3 & 0 & 1 & 2 \\
\tau_2 & 3 & 0 & 1 & 2 \\
\tau_3 & 3 & 0 & 2 & 1 \\
\end{matrix}
$$
显然总有 $\sigma_0 = \text{id}$ 与 $\sigma_p = \sigma$。

\medskip\noindent
引入上述记号后，对
$s \in \check{\mathcal{C}}^{p + 1}(\mathcal{U}, \mathcal{F})$，
定义元素 $h(s) \in \check{\mathcal{C}}^p(\mathcal{U}, \mathcal{F})$，
其分量为
\begin{equation}
\label{cohomology-equation-first-homotopy}
h(s)_{i_0\ldots i_p} =
\sum\nolimits_{0 \leq a \leq p}
(-1)^a \text{符号}(\sigma_a)
s_{i_{\sigma(0)} \ldots i_{\sigma(a)} i_{\sigma_a(a)} \ldots i_{\sigma_a(p)}}
\end{equation}
其中 $\sigma = \sigma^{i_0 \ldots i_p}$。指标 $i_{\sigma(a)}$ 在
$i_{\sigma(0)} \ldots i_{\sigma(a)} i_{\sigma_a(a)} \ldots i_{\sigma_a(p)}$
中出现两次：一次出现在前一组 $a + 1$ 个指标中，另一次出现在后一组
$p - a + 1$ 个指标中。这是因为按 $\sigma_a$ 的定义，对某个
$j \geq a$ 有 $\sigma_a(j) = \sigma(a)$。因此该和有意义，因为每个元素
$s_{i_{\sigma(0)} \ldots i_{\sigma(a)} i_{\sigma_a(a)} \ldots i_{\sigma_a(p)}}$
都定义在开集 $U_{i_0 \ldots i_p}$ 上。还要注意，当 $a = 0$ 时得到
$s_{i_0 \ldots i_p}$，而当 $a = p$ 时得到
$(-1)^p \text{符号}(\sigma) s_{i_{\sigma(0)} \ldots i_{\sigma(p)}}$.

\medskip\noindent
我们断言
$$
(dh + hd)(s)_{i_0 \ldots i_p} =
s_{i_0 \ldots i_p} -
\text{符号}(\sigma) s_{i_{\sigma(0)} \ldots i_{\sigma(p)}}
$$
，其中 $\sigma = \sigma^{i_0 \ldots i_p}$。此断言的验证从略。
（stacks-project 的 scripts 子目录中有一个名为 first-homotopy.gp
的 PARI/gp 脚本，可用来检验此断言的有限多个实例。
我们编写该脚本是为了确认符号正确。）
记
$$
\kappa :
\check{\mathcal{C}}^\bullet(\mathcal{U}, \mathcal{F})
\longrightarrow
\check{\mathcal{C}}^\bullet(\mathcal{U}, \mathcal{F})
$$
为按下式定义的算子：
$$
\kappa(s)_{i_0 \ldots i_p} =
\text{符号}(\sigma^{i_0 \ldots i_p}) s_{i_{\sigma(0)} \ldots i_{\sigma(p)}}.
$$
上述断言蕴含 $\kappa$ 是复形态射，且 $\kappa$ 与 {\v C}ech
复形的恒等映射链同伦。这并不立即蕴含引理，因为算子 $\kappa$
的像并非交替子复形。具体地，$\kappa$ 的像是``半交替''复形
$\check{\mathcal{C}}_{semi\text{-}alt}^p(\mathcal{U}, \mathcal{F})$
，其中 $s$ 是此复形的 $p$-上链，当且仅当
$$
s_{i_0 \ldots i_p} = \text{符号}(\sigma) s_{i_{\sigma(0)} \ldots i_{\sigma(p)}}
$$
对任意 $(i_0, \ldots, i_p) \in I^{p + 1}$ 成立，其中
$\sigma = \sigma^{i_0 \ldots i_p}$。
再引入 {\v C}ech 复形的另一变体，即定义为下式的半有序
{\v C}ech 复形：
$$
\check{\mathcal{C}}_{semi\text{-}ord}^p(\mathcal{U}, \mathcal{F})
=
\prod\nolimits_{i_0 \leq i_1 \leq \ldots \leq i_p}
\mathcal{F}(U_{i_0 \ldots i_p})
$$
容易看出，方程 (\ref{cohomology-equation-d-cech}) 同样定义了微分，因而得到一个复形。
还显然有（类似于引理 \ref{cohomology-lemma-project-to-ordered}）投影映射
$$
\check{\mathcal{C}}_{semi\text{-}alt}^\bullet(\mathcal{U}, \mathcal{F})
\longrightarrow
\check{\mathcal{C}}_{semi\text{-}ord}^\bullet(\mathcal{U}, \mathcal{F})
$$
是复形同构。

\medskip\noindent
因此，只要证明明显的包含映射
$$
\check{\mathcal{C}}_{ord}^p(\mathcal{U}, \mathcal{F})
\longrightarrow
\check{\mathcal{C}}_{semi\text{-}ord}^p(\mathcal{U}, \mathcal{F})
$$
是链同伦等价，就得到引理。为此使用链同伦
\begin{equation}
\label{cohomology-equation-second-homotopy}
h(s)_{i_0 \ldots i_p} =
\left\{
\begin{matrix}
0 & \text{若} & i_0 < i_1 < \ldots < i_p \\
(-1)^a s_{i_0 \ldots i_{a - 1} i_a i_a i_{a + 1} \ldots i_p}
& \text{若} & i_0 < i_1 < \ldots < i_{a - 1} < i_a = i_{a + 1}
\end{matrix}
\right.
\end{equation}
我们断言
$$
(dh + hd)(s)_{i_0 \ldots i_p} =
\left\{
\begin{matrix}
0 & \text{若} & i_0 < i_1 < \ldots < i_p \\
s_{i_0 \ldots i_p}
& \text{否则} &
\end{matrix}
\right.
$$
验证从略。（stacks-project 的 scripts 子目录中有一个名为
second-homotopy.gp 的 PARI/gp 脚本，可用来检验此断言的有限多个实例。
我们编写该脚本是为了确认符号正确。）
该断言清楚表明，投影与自然包含的复合
$$
\check{\mathcal{C}}_{semi\text{-}ord}^\bullet(\mathcal{U}, \mathcal{F})
\longrightarrow
\check{\mathcal{C}}_{ord}^\bullet(\mathcal{U}, \mathcal{F})
\longrightarrow
\check{\mathcal{C}}_{semi\text{-}ord}^\bullet(\mathcal{U}, \mathcal{F})
$$
按所需与恒等映射链同伦。
\end{proof}

\begin{lemma}
\label{cohomology-lemma-alternating-cech-trivial}
设 $X$ 为拓扑空间，$\mathcal{F}$ 为 $X$ 上的阿贝尔预层，
$\mathcal{U} : U = \bigcup_{i \in I} U_i$ 为开覆盖。
若对某个 $i \in I$ 有 $U_i = U$，则扩展的交替 {\v C}ech 复形
$$
\mathcal{F}(U) \to \check{\mathcal{C}}_{alt}^\bullet(\mathcal{U}, \mathcal{F})
$$
是通过把 $\mathcal{F}(U)$ 置于次数 $-1$ 并以
$\mathcal{F}(U)$ 到 $\check{\mathcal{C}}^0(\mathcal{U}, \mathcal{F})$
的典范映射为微分所得；它链同伦等价于 $0$。类似地，对 $I$
上的任意全序，扩展的有序 {\v C}ech 复形
$$
\mathcal{F}(U) \to
\check{\mathcal{C}}_{ord}^\bullet(\mathcal{U}, \mathcal{F})
$$
链同伦等价于 $0$。
\end{lemma}

\begin{proof}[第一个证明]
结合引理 \ref{cohomology-lemma-cech-trivial} 与 \ref{cohomology-lemma-alternating-usual}。
\end{proof}

\begin{proof}[第二个证明]
由于交替与有序 {\v C}ech 复形同构，只需对有序复形证明。
采用标准记号：扩展的有序 {\v C}ech 复形中次数为 $p$ 的上链 $s$
形如 $s = (s_{i_0 \ldots i_p})$，其中
$s_{i_0 \ldots i_p} \in \mathcal{F}(U_{i_0 \ldots i_p})$，且
$i_0 < \ldots < i_p$。在此记号下，
$$
d(x)_{i_0 \ldots i_{p + 1}} =
\sum\nolimits_j (-1)^j x_{i_0 \ldots \hat i_j \ldots i_p}
$$
固定满足 $U = U_i$ 的指标 $i \in I$。作为链同伦，使用映射
$$
h : \text{次数为 }p + 1\text{ 的上链} \to \text{次数为 }p\text{ 的上链}
$$
，定义为
$$
h(s)_{i_0 \ldots i_p} = 0 \text{，若 } i \in \{i_0, \ldots, i_p\}
\text{；且}
h(s)_{i_0 \ldots i_p} = 
(-1)^j s_{i_0 \ldots i_j i i_{j + 1} \ldots i_p} \text{，否则}
$$
在第二种情形中，$j$ 是使得 $i_j < i < i_{j + 1}$ 的唯一指标；
又因 $U = U_i$，有等式
$$
\mathcal{F}(U_{i_0 \ldots i_p}) =
\mathcal{F}(U_{i_0 \ldots i_j i i_{j + 1} \ldots i_p})
$$
故可将
$(-1)^j s_{i_0 \ldots i_j i i_{j + 1} \ldots i_p}$
视为 $\mathcal{F}(U_{i_0 \ldots i_p})$ 的元素。
下面通过计算证明 $d h + h d$ 等于恒等映射的负值，从而完成证明。
为此固定次数为 $p$ 的上链 $s$，并取 $I$ 中的元素
$i_0 < \ldots < i_p$。

\medskip\noindent
情形 I：$i \in \{i_0, \ldots, i_p\}$。设 $i = i_t$。
于是 $h(d(s))_{i_0 \ldots i_p} = 0$。另一方面，
$$
d(h(s))_{i_0 \ldots i_p} =
\sum (-1)^j h(s)_{i_0 \ldots \hat i_j \ldots i_p} =
(-1)^t h(s)_{i_0 \ldots \hat i \ldots i_p} =
(-1)^t (-1)^{t - 1} s_{i_0 \ldots i_p}
$$
因此按所需有
$(dh + hd)(s)_{i_0 \ldots i_p} = -s_{i_0 \ldots i_p}$。

\medskip\noindent
情形 II：$i \not \in \{i_0, \ldots, i_p\}$。取 $j$ 使得
$i_j < i < i_{j + 1}$。于是
\begin{align*}
h(d(s))_{i_0 \ldots i_p}
& =
(-1)^j d(s)_{i_0 \ldots i_j i i_{j + 1} \ldots i_p} \\
& =
\sum\nolimits_{j' \leq j} (-1)^{j + j'}
s_{i_0 \ldots \hat i_{j'} \ldots i_j i i_{j + 1} \ldots i_p} -
s_{i_0 \ldots i_p} \\
&
+ \sum\nolimits_{j' > j} (-1)^{j + j' + 1}
s_{i_0 \ldots i_j i i_{j + 1} \ldots \hat i_{j'} \ldots i_p}
\end{align*}
另一方面，
\begin{align*}
d(h(s))_{i_0 \ldots i_p}
& =
\sum\nolimits_{j'} (-1)^{j'} h(s)_{i_0 \ldots \hat i_{j'} \ldots i_p} \\
& =
\sum\nolimits_{j' \leq j} (-1)^{j' + j - 1}
s_{i_0 \ldots \hat i_{j'} \ldots i_j i i_{j + 1} \ldots i_p} \\
& +
\sum\nolimits_{j' > j} (-1)^{j' + j}
s_{i_0 \ldots i_j i i_{j + 1} \ldots \hat i_{j'} \ldots i_p}
\end{align*}
两式相加得到
$(dh + hd)(s)_{i_0 \ldots i_p} = - s_{i_0 \ldots i_p}$
，正是所需结论。
\end{proof}





\section{{\v C}ech 复形的另一种看法}
\label{cohomology-section-locally-finite-cech}

\noindent
本节讨论建立 {\v C}ech 复形与上同调之间关系的另一种方法。

\begin{lemma}
\label{cohomology-lemma-covering-resolution}
设 $X$ 为环空间，$\mathcal{U} : X = \bigcup_{i \in I} U_i$
为 $X$ 的开覆盖，$\mathcal{F}$ 为 $\mathcal{O}_X$-模。
以 $\mathcal{F}_{i_0 \ldots i_p}$ 表示 $\mathcal{F}$
在 $U_{i_0 \ldots i_p}$ 上的限制。存在 $\mathcal{O}_X$-模复形
${\mathfrak C}^\bullet(\mathcal{U}, \mathcal{F})$
，满足
$$
{\mathfrak C}^p(\mathcal{U}, \mathcal{F}) =
\prod\nolimits_{i_0 \ldots i_p}
(j_{i_0 \ldots i_p})_* \mathcal{F}_{i_0 \ldots i_p}
$$
，且其微分
$d : {\mathfrak C}^p(\mathcal{U}, \mathcal{F})
\to {\mathfrak C}^{p + 1}(\mathcal{U}, \mathcal{F})$
如方程 (\ref{cohomology-equation-d-cech}) 所示。此外，存在典范映射
$$
\mathcal{F} \to {\mathfrak C}^\bullet(\mathcal{U}, \mathcal{F})
$$
，它是拟同构；换言之，
${\mathfrak C}^\bullet(\mathcal{U}, \mathcal{F})$ 是 $\mathcal{F}$ 的消解。
\end{lemma}

\begin{proof}
只需验证
$$
0 \to \mathcal{F} \to \mathfrak{C}^0(\mathcal{U}, \mathcal{F}) \to
\mathfrak{C}^1(\mathcal{U}, \mathcal{F}) \to  \ldots
$$
在茎上正合。取 $x \in X$，并选取满足
$x \in U_{i_{\text{fix}}}$ 的 $i_{\text{fix}} \in I$。定义
$$
h : \mathfrak{C}^p(\mathcal{U}, \mathcal{F})_x
\to \mathfrak{C}^{p - 1}(\mathcal{U}, \mathcal{F})_x
$$
如下：若 $s \in \mathfrak{C}^p(\mathcal{U}, \mathcal{F})_x$，取其在
$x$ 的某个邻域 $V$ 上定义的代表元
$$
\widetilde{s} \in
\mathfrak{C}^p(\mathcal{U}, \mathcal{F})(V) =
\prod\nolimits_{i_0 \ldots i_p}
\mathcal{F}(V \cap U_{i_0} \cap \ldots \cap U_{i_p})
$$
，并置
$$
h(s)_{i_0 \ldots i_{p - 1}} =
\widetilde{s}_{i_{\text{fix}} i_0 \ldots i_{p - 1}, x}.
$$
同一公式（取 $p = 0$）给出映射
$\mathfrak{C}^{0}(\mathcal{U},\mathcal{F})_x \to \mathcal{F}_x$。
形式计算如下：
\begin{align*}
(dh + hd)(s)_{i_0 \ldots i_p}
& =
\sum\nolimits_{j = 0}^p
(-1)^j
h(s)_{i_0 \ldots \hat i_j \ldots i_p}
+
d(s)_{i_{\text{fix}} i_0 \ldots i_p}\\
& =
\sum\nolimits_{j = 0}^p
(-1)^j
s_{i_{\text{fix}} i_0 \ldots \hat i_j \ldots i_p}
+
s_{i_0 \ldots i_p}
+
\sum\nolimits_{j = 0}^p
(-1)^{j + 1}
s_{i_{\text{fix}} i_0 \ldots \hat i_j \ldots i_p} \\
& =
s_{i_0 \ldots i_p}
\end{align*}
这表明 $h$ 给出扩展复形
$$
0 \to \mathcal{F}_x \to \mathfrak{C}^0(\mathcal{U}, \mathcal{F})_x
\to \mathfrak{C}^1(\mathcal{U}, \mathcal{F})_x \to \ldots
$$
的恒等映射与零映射之间的链同伦，结论得证。
\end{proof}

\noindent
利用此引理，很容易重新证明引理
\ref{cohomology-lemma-cech-spectral-sequence} 的 {\v C}ech 到上同调谱序列。
具体地，设 $X,\mathcal{U},\mathcal{F}$ 如引理
\ref{cohomology-lemma-covering-resolution}，并设
$\mathcal{F} \to \mathcal{I}^\bullet$ 为内射消解。于是可考虑双复形
$$
A^{\bullet, \bullet} =
\Gamma(X, {\mathfrak C}^\bullet(\mathcal{U}, \mathcal{I}^\bullet)).
$$
按构造有
$$
A^{p, q} = \prod\nolimits_{i_0 \ldots i_p} \mathcal{I}^q(U_{i_0 \ldots i_p})
$$
考虑《同调代数》第 \ref{homology-section-double-complex} 节中
与此双复形相伴的两个谱序列，特别见《同调代数》引理
\ref{homology-lemma-ss-double-complex}。
对谱序列 $({}'E_r, {}'d_r)_{r \geq 0}$，有
${}'E_2^{p, q} = \check{H}^p(\mathcal{U}, \underline{H}^q(\mathcal{F}))$
，因为取积正合（《同调代数》引理
\ref{homology-lemma-product-abelian-groups-exact}）。
对谱序列 $({}''E_r, {}''d_r)_{r \geq 0}$，若 $p > 0$，则
${}''E_2^{p, q} = 0$，且 ${}''E_2^{0, q} = H^q(X, \mathcal{F})$。
事实上，对固定的 $q$，层复形
${\mathfrak C}^\bullet(\mathcal{U}, \mathcal{I}^q)$
是内射层 $\mathcal{I}^q$ 的内射层消解（引理
\ref{cohomology-lemma-covering-resolution}；另由引理
\ref{cohomology-lemma-cohomology-of-open}、\ref{cohomology-lemma-pushforward-injective-flat}
及《同调代数》引理 \ref{homology-lemma-product-injectives}）。
因此
$\Gamma(X, {\mathfrak C}^\bullet(\mathcal{U}, \mathcal{I}^q))$
的上同调在正次数为零，在次数 $0$ 等于
$\Gamma(X, \mathcal{I}^q)$。再对下一个微分取上同调，
便得到关于谱序列 $({}''E_r, {}''d_r)_{r \geq 0}$ 的断言。
两个谱序列都收敛到 $A^{\bullet, \bullet}$ 相伴总复形的上同调，故得结论。

\begin{definition}
\label{cohomology-definition-covering-locally-finite}
设 $X$ 为拓扑空间。若对每个 $x \in X$ 都存在 $x$ 的开邻域 $W$，
使得 $\{i \in I \mid W \cap U_i \not = \emptyset\}$ 有限，
则称开覆盖 $X = \bigcup_{i \in I} U_i$ {\it 局部有限}。
\end{definition}

\begin{remark}
\label{cohomology-remark-locally-finite-sections}
设 $X = \bigcup_{i \in I} U_i$ 为局部有限开覆盖，
并以 $j_i : U_i \to X$ 表示包含映射。设对每个 $i$，
都给定 $U_i$ 上的阿贝尔层 $\mathcal{F}_i$。考虑阿贝尔层
$\mathcal{G} = \bigoplus_{i \in I} (j_i)_*\mathcal{F}_i$。
则对开集 $V \subset X$，事实上有
$$
\Gamma(V, \mathcal{G}) = \prod\nolimits_{i \in I} \mathcal{F}_i(V \cap U_i).
$$
换言之，
$$
\bigoplus\nolimits_{i \in I} (j_i)_*\mathcal{F}_i =
\prod\nolimits_{i \in I} (j_i)_*\mathcal{F}_i
$$
起初这似乎有些奇怪，但要注意，一族层的直和是其直观预期对象的层化。
相关讨论见《模层》第 \ref{modules-section-kernels} 节。
因此在此情形下，引理 \ref{cohomology-lemma-covering-resolution} 中复形的各项为
$$
{\mathfrak C}^p(\mathcal{U}, \mathcal{F}) =
\bigoplus\nolimits_{i_0 \ldots i_p}
(j_{i_0 \ldots i_p})_* \mathcal{F}_{i_0 \ldots i_p}
$$
，这有时很有用。
\end{remark}











\section{复形的 {\v C}ech 上同调}
\label{cohomology-section-cech-cohomology-of-complexes}

\noindent
一般而言，对 $X$ 上的阿贝尔群层 ${\mathcal F}$ 与 ${\mathcal G}$，
存在杯积映射
$$
H^i(X, {\mathcal F}) \times H^j(X, {\mathcal G})
\longrightarrow
H^{i + j}(X, {\mathcal F} \otimes_{\mathbf Z} {\mathcal G}).
$$
本节用 {\v C}ech 上循环和一个显式杯积公式定义此映射。
如果担心上同调未必等于 {\v C}ech 上同调，则可使用超覆盖，
同时仍采用上循环记号。这样还有一个优点：同一方法可定义任意拓扑斯上
超上同调的杯积（此处插入未来参考文献）。

\medskip\noindent
设 ${\mathcal F}^\bullet$ 为 $X$ 上阿贝尔群预层的下有界复形。
我们常可用 {\v C}ech 上循环计算 $H^n(X, {\mathcal F}^\bullet)$。
具体地，设 ${\mathcal U} : X = \bigcup_{i \in I} U_i$
为 $X$ 的开覆盖。由于 {\v C}ech 复形
$\check{\mathcal{C}}^\bullet(\mathcal{U}, \mathcal{F})$
（定义 \ref{cohomology-definition-cech-complex}）关于预层 $\mathcal{F}$ 具有函子性，
可得到双复形
$\check{\mathcal{C}}^\bullet(\mathcal{U}, \mathcal{F}^\bullet)$.
与 $\check{\mathcal{C}}^\bullet({\mathcal U}, {\mathcal F}^\bullet)$
相伴的总复形，其次数 $n$ 项为
$$
\text{Tot}^n(\check{\mathcal{C}}^\bullet({\mathcal U}, {\mathcal F}^\bullet))
=
\bigoplus\nolimits_{p + q = n}
\prod\nolimits_{i_0\ldots i_p} {\mathcal F}^q(U_{i_0\ldots i_p})
$$
；见《同调代数》定义
\ref{homology-definition-associated-simple-complex}。
以 $\alpha = \{\alpha_{i_0\ldots i_p}\}$ 表示 $\text{Tot}^n$
中的典型元素，其中
$\alpha_{i_0 \ldots i_p} \in \mathcal{F}^q(U_{i_0\ldots i_p})$。
换言之，$\alpha_{i_0\ldots i_p}$ 的 $\mathcal{F}$-次数为
$q = n - p$。此记号要求我们始终留意 $\alpha$ 所在的次数。
用公式
$\deg_{\mathcal F}(\alpha_{i_0\ldots i_p}) = q$.
表示这一情形。按照《同调代数》定义
\ref{homology-definition-associated-simple-complex} 中的约定，
次数为 $n$ 的元素 $\alpha$ 的微分为
$$
d(\alpha)_{i_0\ldots i_{p + 1}}
=
\sum\nolimits_{j = 0}^{p + 1}
(-1)^j \alpha_{i_0 \ldots \hat i_j \ldots i_{p + 1}} + 
(-1)^{p + 1}d_{{\mathcal F}}(\alpha_{i_0 \ldots i_{p + 1}})
$$
，其中 $d_\mathcal{F}$ 表示复形 $\mathcal{F}^\bullet$ 的微分。
表达式 $\alpha_{i_0 \ldots \hat i_j \ldots i_{p + 1}}$ 表示把
$\alpha_{i_0 \ldots \hat i_j \ldots i_{p + 1}}
\in {\mathcal F}(U_{i_0\ldots\hat i_j\ldots i_{p + 1}})$
限制到 $U_{i_0 \ldots i_{p + 1}}$ 所得的元素。

\medskip\noindent
构造
$\text{Tot}(\check{\mathcal{C}}^\bullet({\mathcal U}, {\mathcal F}^\bullet))$
关于 ${\mathcal F}^\bullet$ 具有函子性。此外，存在复形的函子性变换
\begin{equation}
\label{cohomology-equation-global-sections-to-cech}
\Gamma(X, {\mathcal F}^\bullet)
\longrightarrow
\text{Tot}(\check{\mathcal{C}}^\bullet({\mathcal U}, {\mathcal F}^\bullet))
\end{equation}
，定义如下：截面 $s\in \Gamma(X, {\mathcal F}^n)$
映为元素 $\alpha = \{\alpha_{i_0\ldots i_p}\}$，其中
$\alpha_{i_0} = s|_{U_{i_0}}$，且当 $p > 0$ 时
$\alpha_{i_0\ldots i_p} = 0$。

\medskip\noindent
加细。设 ${\mathcal V} = \{ V_j \}_{j\in J}$ 是
${\mathcal U}$ 的一个加细。这意味着存在映射 $t: J \to I$，
使得对所有 $j\in J$ 都有 $V_j \subset U_{t(j)}$。由此得到函子性变换
\begin{equation}
\label{cohomology-equation-transformation}
T_t :
\text{Tot}(\check{\mathcal{C}}^\bullet({\mathcal U}, {\mathcal F}^\bullet))
\longrightarrow
\text{Tot}(\check{\mathcal{C}}^\bullet({\mathcal V}, {\mathcal F}^\bullet))
\end{equation}
，其定义规则为
$$
T_t(\alpha)_{j_0\ldots j_p}
=
\alpha_{t(j_0)\ldots t(j_p)}|_{V_{j_0\ldots j_p}}.
$$
给定如上的两个映射 $t, t' : J \to I$，上面构造的映射
$T_t$ 与 $T_{t'}$ 链同伦。此链同伦由下式给出：
$$
h(\alpha)_{j_0\ldots j_p}
=
\sum\nolimits_{a = 0}^{p}
(-1)^a
\alpha_{t(j_0)\ldots t(j_a) t'(j_a) \ldots t'(j_p)}
$$
，其中 $\alpha$ 为次数 $n$ 的元素。下面的计算表明此式确实成立；
这里仍设 $\alpha$ 为次数 $n$ 的元素（因而 $d(\alpha)$
的次数为 $n + 1$，而 $h(\alpha)$ 的次数为 $n - 1$）：
\begin{align*}
(
d(h(\alpha)) + h(d(\alpha))
)_{j_0\ldots j_p}
= &
\sum\nolimits_{k = 0}^p
(-1)^k
h(\alpha)_{j_0 \ldots \hat j_k \ldots j_p}
+ \\
&
(-1)^p
d_{\mathcal F}(h(\alpha)_{j_0 \ldots j_p})
+ \\
&
\sum\nolimits_{a = 0}^p
(-1)^a
d(\alpha)_{t(j_0) \ldots t(j_a) t'(j_a) \ldots t'(j_p)}
\\
= &
\sum\nolimits_{k = 0}^p
\sum\nolimits_{a = 0}^{k - 1}
(-1)^{k + a}
\alpha_{t(j_0)\ldots t(j_a)t'(j_a)\ldots \hat{t'(j_k)}\ldots t'(j_p)}
+ \\
&
\sum\nolimits_{k = 0}^p
\sum\nolimits_{a = k + 1}^p
(-1)^{k + a - 1}
\alpha_{t(j_0)\ldots \hat{t(j_k)}\ldots t(j_a)t'(j_a)\ldots t'(j_p)}
+ \\
&
\sum\nolimits_{a = 0}^p
(-1)^{p + a}
d_{\mathcal F}(\alpha_{t(j_0)\ldots t(j_a) t'(j_a) \ldots t'(j_p)})
+ \\
&
\sum\nolimits_{a = 0}^p
\sum\nolimits_{k = 0}^a
(-1)^{a + k}
\alpha_{t(j_0)\ldots\hat{t(j_k)}\ldots t(j_a)t'(j_a)\ldots t'(j_p)}
+ \\
&
\sum\nolimits_{a = 0}^p
\sum\nolimits_{k = a}^p
(-1)^{a + k + 1}
\alpha_{t(j_0) \ldots t(j_a) t'(j_a) \ldots \hat{t'(j_k)} \ldots t'(j_p)}
+ \\
&
\sum\nolimits_{a = 0}^p
(-1)^{a + p + 1}
d_{\mathcal F}(\alpha_{t(j_0)\ldots t(j_a) t'(j_a) \ldots t'(j_p)})
\\
= &
\alpha_{t'(j_0)\ldots t'(j_p)} +
(-1)^{2p + 1}\alpha_{t(j_0)\ldots t(j_p)}
\\
= &
T_{t'}(\alpha)_{j_0\ldots j_p} - T_t(\alpha)_{j_0\ldots j_p}
\end{align*}
我们把各项相消的验证留给读者。（注意，第一、第二以及第四、第五个
和式中同时含有 $k$ 与 $a$ 的各项会相消；例外是 $a = k$ 的项，
它们只出现在第四和第五个和式中，且彼此相消，仅留下所需的两项。）
由此可知，诱导映射
$$
H^n(T_t) :
H^n(
\text{Tot}(\check{\mathcal{C}}^\bullet({\mathcal U}, {\mathcal F}^\bullet))
)
\to
H^n(
\text{Tot}(\check{\mathcal{C}}^\bullet({\mathcal V}, {\mathcal F}^\bullet))
)
$$
与 $t$ 的选取无关。我们把 {\it Cech 超上同调}定义为
经由映射 $H^\bullet(T_t)$ 对所有加细取 Cech 上同调群的余极限。

\medskip\noindent
在（遍历 $X$ 的所有开覆盖的）余极限中，下列引理给出一个
从 Cech 超上同调到上同调的映射；它常常是同构，而若使用超覆盖，
则它总是同构。

\begin{lemma}
\label{cohomology-lemma-cech-complex-complex}
设 $(X, \mathcal{O}_X)$ 是环空间，且
$\mathcal{U} : X = \bigcup_{i \in I} U_i$ 是开覆盖。
对 $\mathcal{O}_X$-模的有下界复形 $\mathcal{F}^\bullet$，
存在典范映射
$$
\text{Tot}(\check{\mathcal{C}}^\bullet(\mathcal{U}, \mathcal{F}^\bullet))
\longrightarrow
R\Gamma(X, \mathcal{F}^\bullet)
$$
，它关于 $\mathcal{F}^\bullet$ 具有函子性，并与
(\ref{cohomology-equation-global-sections-to-cech}) 和
(\ref{cohomology-equation-transformation}) 相容。存在谱序列
$(E_r, d_r)_{r \geq 0}$，满足
$$
E_2^{p, q} =
H^p(\text{Tot}(\check{\mathcal{C}}^\bullet(\mathcal{U},
\underline{H}^q(\mathcal{F}^\bullet)))
$$
且收敛到 $H^{p + q}(X, \mathcal{F}^\bullet)$。
\end{lemma}

\begin{proof}
设 ${\mathcal I}^\bullet$ 是内射对象的有下界复形。
对 $\mathcal{I}^\bullet$ 而言，
映射 (\ref{cohomology-equation-global-sections-to-cech}) 为
$\Gamma(X, {\mathcal I}^\bullet) \to
\text{Tot}(\check{\mathcal{C}}^\bullet({\mathcal U}, {\mathcal I}^\bullet))$.
把《同调代数》引理
\ref{homology-lemma-double-complex-gives-resolution}
应用于双复形
$\check{\mathcal{C}}^\bullet({\mathcal U}, {\mathcal I}^\bullet)$，
并使用引理 \ref{cohomology-lemma-injective-trivial-cech}，可知这是
阿贝尔群复形的拟同构。设
${\mathcal F}^\bullet \to {\mathcal I}^\bullet$ 是从
${\mathcal F}^\bullet$ 到内射对象的有下界复形的拟同构。
由于 $R\Gamma(X, {\mathcal F}^\bullet)$ 由复形
$\Gamma(X, {\mathcal I}^\bullet)$ 表示，利用映射
$$
\text{Tot}(\check{\mathcal{C}}^\bullet({\mathcal U}, {\mathcal F}^\bullet))
\longrightarrow
\text{Tot}(\check{\mathcal{C}}^\bullet({\mathcal U}, {\mathcal I}^\bullet)).
$$
便得到引理中的映射。我们略去函子性与相容性的验证。
为构造引理中的谱序列，选取 Cartan--Eilenberg 消解
$\mathcal{F}^\bullet \to \mathcal{I}^{\bullet, \bullet}$，见
《导出范畴》引理 \ref{derived-lemma-cartan-eilenberg}。此时
$\mathcal{F}^\bullet \to \text{Tot}(\mathcal{I}^{\bullet, \bullet})$
是内射消解，因而
$$
\text{Tot}(\check{\mathcal{C}}^\bullet({\mathcal U},
\text{Tot}({\mathcal I}^{\bullet, \bullet})))
$$
如上所述，此复形计算 $R\Gamma(X, \mathcal{F}^\bullet)$。
由《同调代数》注记 \ref{homology-remark-triple-complex}，
可把它看成三重复形
$\check{\mathcal{C}}^\bullet({\mathcal U},
{\mathcal I}^{\bullet, \bullet})$ 所伴随的总复形；因此再次使用同一注记，
还可把它看成双复形 $A^{\bullet, \bullet}$ 所伴随的总复形，其各项为
$$
A^{n, m} =
\bigoplus\nolimits_{p + q = n}
\check{\mathcal{C}}^p({\mathcal U}, \mathcal{I}^{q, m})
$$
由于 $\mathcal{I}^{q, \bullet}$ 是 $\mathcal{F}^q$ 的内射消解，
可应用 $A^{\bullet, \bullet}$ 所伴随的第一谱序列
（《同调代数》引理 \ref{homology-lemma-ss-double-complex}），
得到满足下式的谱序列：
$$
E_1^{n, m} =
\bigoplus\nolimits_{p + q = n}
\check{\mathcal{C}}^p(\mathcal{U}, \underline{H}^m(\mathcal{F}^q))
$$
这是复形
$\text{Tot}(\check{\mathcal{C}}^\bullet(\mathcal{U},
\underline{H}^m(\mathcal{F}^\bullet))$ 的第 $n$ 项。因此得到
引理中所述的 $E_2$ 项。收敛性由《同调代数》引理
\ref{homology-lemma-first-quadrant-ss} 给出。
\end{proof}

\begin{lemma}
\label{cohomology-lemma-cech-complex-complex-computes}
设 $(X, \mathcal{O}_X)$ 是环空间，
$\mathcal{U} : X = \bigcup_{i \in I} U_i$ 是开覆盖，且
$\mathcal{F}^\bullet$ 是 $\mathcal{O}_X$-模的有下界复形。
若对所有 $i > 0$ 及所有 $p, i_0, \ldots, i_p, q$ 都有
$H^i(U_{i_0 \ldots i_p}, \mathcal{F}^q) = 0$，则引理
\ref{cohomology-lemma-cech-complex-complex} 中的映射
$
\text{Tot}(\check{\mathcal{C}}^\bullet(\mathcal{U}, \mathcal{F}^\bullet))
\to
R\Gamma(X, \mathcal{F}^\bullet)
$
是同构。
\end{lemma}

\begin{proof}
这立即由引理 \ref{cohomology-lemma-cech-complex-complex} 的谱序列得到。
\end{proof}

\begin{remark}
\label{cohomology-remark-shift-complex-cech-complex}
设 $(X, \mathcal{O}_X)$ 是环空间，
$\mathcal{U} : X = \bigcup_{i \in I} U_i$ 是开覆盖，
$\mathcal{F}^\bullet$ 是 $\mathcal{O}_X$-模的有下界复形，
而 $b$ 是整数。我们断言在导出范畴中有交换图
$$
\xymatrix{
\text{Tot}(\check{\mathcal{C}}^\bullet(\mathcal{U}, \mathcal{F}^\bullet))[b]
\ar[r] \ar[d]_\gamma &
R\Gamma(X, \mathcal{F}^\bullet)[b] \ar[d] \\
\text{Tot}(\check{\mathcal{C}}^\bullet(\mathcal{U}, \mathcal{F}^\bullet[b]))
\ar[r] &
R\Gamma(X, \mathcal{F}^\bullet[b])
}
$$
，其中映射 $\gamma$ 是《同调代数》注记
\ref{homology-remark-shift-double-complex} 中构造的复形映射。
这是有意义的，因为双复形
$\check{\mathcal{C}}^\bullet(\mathcal{U}, \mathcal{F}^\bullet[b])$
显然就是《同调代数》注记
\ref{homology-remark-shift-double-complex} 中引入的双复形
$\check{\mathcal{C}}^\bullet(\mathcal{U}, \mathcal{F}^\bullet)[0, b]$
。为验证该图交换，可像引理
\ref{cohomology-lemma-cech-complex-complex} 的证明那样选取内射消解
$\mathcal{F}^\bullet \to \mathcal{I}^\bullet$。作图追踪可知，
只需验证把 $\mathcal{F}^\bullet$ 替换为
$\mathcal{I}^\bullet$ 后该图交换。于是考虑扩充图
$$
\xymatrix{
\Gamma(X, \mathcal{I}^\bullet)[b] \ar[r] \ar[d] &
\text{Tot}(\check{\mathcal{C}}^\bullet(\mathcal{U}, \mathcal{I}^\bullet))[b]
\ar[r] \ar[d]_\gamma &
R\Gamma(X, \mathcal{I}^\bullet)[b] \ar[d] \\
\Gamma(X, \mathcal{I}^\bullet[b]) \ar[r] &
\text{Tot}(\check{\mathcal{C}}^\bullet(\mathcal{U}, \mathcal{I}^\bullet[b]))
\ar[r] &
R\Gamma(X, \mathcal{I}^\bullet[b])
}
$$
，其中左侧水平箭头为
(\ref{cohomology-equation-global-sections-to-cech})。在此情形下，
水平箭头在导出范畴中均为同构（见引理
\ref{cohomology-lemma-cech-complex-complex} 的证明），故只需证明左方块交换。
这确实成立，因为映射 $\gamma$ 在直和项
$\check{\mathcal{C}}^0(\mathcal{U}, \mathcal{I}^{q + b})$
上所用的符号为 $1$；参见《同调代数》注记
\ref{homology-remark-shift-double-complex} 中的公式。
\end{remark}

\noindent
设 $X$ 是拓扑空间，$\mathcal{U} : X = \bigcup_{i \in I} U_i$
是开覆盖，而 $\mathcal{F}^\bullet$ 是阿贝尔群预层的有下界复形。
考虑映射
$\tau :
\text{Tot}(\check{\mathcal{C}}^\bullet({\mathcal U}, {\mathcal F}^\bullet))
\to
\text{Tot}(\check{\mathcal{C}}^\bullet({\mathcal U}, {\mathcal F}^\bullet))$
，其定义为
$$
\tau(\alpha)_{i_0 \ldots i_p} = (-1)^{p(p + 1)/2} \alpha_{i_p \ldots i_0}.
$$
于是对次数 $n$ 的元素 $\alpha$ 有
\begin{align*}
& d(\tau(\alpha))_{i_0 \ldots i_{p + 1}} \\
& =
\sum\nolimits_{j = 0}^{p + 1}
(-1)^j
\tau(\alpha)_{i_0 \ldots \hat i_j \ldots i_{p + 1}}
+
(-1)^{p + 1}
d_{\mathcal F}(\tau(\alpha)_{i_0 \ldots i_{p + 1}})
\\
& =
\sum\nolimits_{j = 0}^{p + 1}
(-1)^{j + \frac{p(p + 1)}{2}}
\alpha_{i_{p + 1} \ldots \hat i_j \ldots i_0}
+
(-1)^{p + 1 + \frac{(p + 1)(p + 2)}{2}}
d_{\mathcal F}(\alpha_{i_{p + 1} \ldots i_0})
\end{align*}
另一方面有
\begin{align*}
& \tau(d(\alpha))_{i_0\ldots i_{p + 1}} \\
& =
(-1)^{\frac{(p + 1)(p + 2)}{2}} d(\alpha)_{i_{p + 1} \ldots i_0}
\\
& =
(-1)^{\frac{(p + 1)(p + 2)}{2}}
\left(
\sum\nolimits_{j = 0}^{p + 1}
(-1)^j
\alpha_{i_{p + 1}\ldots \hat i_{p + 1 - j} \ldots i_0}
+
(-1)^{p + 1}
d_{\mathcal F}(\alpha_{i_{p + 1}\ldots i_0})
\right)
\end{align*}
由于
$p(p + 1)/2 \equiv (p + 1)(p + 2)/2 + p + 1 \bmod 2$，
可得 $d(\tau(\alpha)) = \tau(d(\alpha))$。换言之，
$\tau$ 是复形
$\text{Tot}(\check{\mathcal{C}}^\bullet({\mathcal U}, {\mathcal F}^\bullet))$
的自同态。注意，下图
$$
\begin{matrix}
\Gamma(X, {\mathcal F}^\bullet) &
\longrightarrow &
\text{Tot}(\check{\mathcal{C}}^\bullet({\mathcal U}, {\mathcal F}^\bullet)) \\
\downarrow \text{id} & & \downarrow \tau \\
\Gamma(X, {\mathcal F}^\bullet) &
\longrightarrow &
\text{Tot}(\check{\mathcal{C}}^\bullet({\mathcal U}, {\mathcal F}^\bullet))
\end{matrix}
$$
交换。此外，$\tau$ 显然与加细相容。这表明 $\tau$ 在
Cech 上同调上按恒等映射作用（即在余极限中如此——前提是
Cech 超上同调与超上同调一致；使用超覆盖时总是如此）。
我们断言，在总 Cech 复形
$\text{Tot}(\check{\mathcal{C}}^\bullet({\mathcal U}, {\mathcal F}^\bullet))$
上，$\tau$ 实际上与恒等映射链同伦。为证明这一点，取链同伦
$$
h(\alpha)_{i_0 \ldots i_p}
=
\sum\nolimits_{a = 0}^p
\epsilon_p(a)
\alpha_{i_0 \ldots i_a i_p \ldots i_a}
\quad\text{其中}\quad
\epsilon_p(a) = (-1)^{\frac{(p - a)(p - a - 1)}{2} + p}
$$
，其中 $\alpha$ 的次数为 $n$。与往常一样，我们省略
$|_{U_{i_0 \ldots i_p}}$。下面的计算表明此式成立；
这里仍设 $\alpha$ 为次数 $n$ 的元素：
\begin{align*}
(d(h(\alpha)) + h(d(\alpha)))_{i_0 \ldots i_p}
= &
\sum\nolimits_{k = 0}^p
(-1)^k
h(\alpha)_{i_0 \ldots \hat i_k \ldots i_p}
+ \\
&
(-1)^p
d_{\mathcal F}(h(\alpha)_{i_0 \ldots i_p})
+ \\
&
\sum\nolimits_{a = 0}^p
\epsilon_p(a)
d(\alpha)_{i_0 \ldots i_a i_p \ldots i_a}
\\
= &
\sum\nolimits_{k = 0}^p
\sum\nolimits_{a = 0}^{k - 1}
(-1)^k \epsilon_{p - 1}(a)
\alpha_{i_0 \ldots i_a i_p \ldots \hat{i_k} \ldots i_a}
+ \\
&
\sum\nolimits_{k = 0}^p
\sum\nolimits_{a = k + 1}^p
(-1)^k \epsilon_{p - 1}(a - 1)
\alpha_{i_0 \ldots \hat{i_k} \ldots i_a i_p \ldots i_a}
+ \\
&
\sum\nolimits_{a = 0}^p
(-1)^p \epsilon_p(a)
d_{\mathcal F}(\alpha_{i_0 \ldots i_a i_p \ldots i_a})
+ \\
&
\sum\nolimits_{a = 0}^p
\sum\nolimits_{k = 0}^a
\epsilon_p(a) (-1)^k
\alpha_{i_0 \ldots \hat{i_k} \ldots i_a i_p \ldots i_a}
+ \\
&
\sum\nolimits_{a = 0}^p
\sum\nolimits_{k = a}^p
\epsilon_p(a) (-1)^{p + a + 1 - k}
\alpha_{i_0 \ldots i_a i_p \ldots \hat{i_k} \ldots i_a}
+ \\
&
\sum\nolimits_{a = 0}^p
\epsilon_p(a) (-1)^{p + 1}
d_{\mathcal F}(\alpha_{i_0 \ldots i_a i_p \ldots i_a})
\\
= &
\epsilon_p(0) \alpha_{i_p \ldots i_0} +
\epsilon_p(p) (-1)^{p + 1} \alpha_{i_0 \ldots i_p} \\
= &
(-1)^{\frac{p(p + 1)}{2}}\alpha_{i_p \ldots i_0}
- \alpha_{i_0 \ldots i_p}
\end{align*}
各项相消是因为
$$
(-1)^k \epsilon_{p - 1}(a) + \epsilon_p(a)(-1)^{p + a + 1 - k} = 0
\quad\text{且}\quad
(-1)^k\epsilon_{p - 1}(a - 1) + \epsilon_p(a) (-1)^k = 0
$$
我们把相消的验证留给读者。

\medskip\noindent
设有阿贝尔层的两个有下界复形
${\mathcal F}^\bullet$ 与 ${\mathcal G}^\bullet$。我们把复形
$\text{Tot}({\mathcal F}^\bullet\otimes_{\mathbf Z} {\mathcal G}^\bullet)$
定义为各项为
$\bigoplus_{p + q = n} {\mathcal F}^p \otimes {\mathcal G}^q$
且微分遵循下列规则的复形：
\begin{equation}
\label{cohomology-equation-differential-tensor-product-complexes}
d(\alpha \otimes \beta) =
d(\alpha)\otimes \beta + (-1)^{\deg(\alpha)} \alpha \otimes d(\beta)
\end{equation}
，其中 $\alpha$ 与 $\beta$ 为齐次元素；见《同调代数》定义
\ref{homology-definition-associated-simple-complex}。

\medskip\noindent
设 $M^\bullet$ 与 $N^\bullet$ 是阿贝尔群的两个有下界复形。
若 $m$（相应地 $n$）是 $M^\bullet$（相应地 $N^\bullet$）
的上循环，则 $m \otimes n$ 显然是
$\text{Tot}(M^\bullet\otimes N^\bullet)$ 的上循环。因此有杯积
$$
H^i(M^\bullet) \times H^j(N^\bullet)
\longrightarrow
H^{i + j}(Tot(M^\bullet\otimes N^\bullet)).
$$
《代数详论》第 \ref{more-algebra-section-products-tor} 节也讨论了这一点。

\medskip\noindent
因此，复形超上同调中杯积的构造归结为构造复形映射
\begin{equation}
\label{cohomology-equation-needs-signs}
\text{Tot}\left(
\text{Tot}(\check{\mathcal{C}}^\bullet({\mathcal U}, {\mathcal F}^\bullet))
\otimes_{\mathbf Z}
\text{Tot}(\check{\mathcal{C}}^\bullet({\mathcal U}, {\mathcal G}^\bullet))
\right)
\longrightarrow
\text{Tot}(
\check{\mathcal{C}}^\bullet({\mathcal U},
\text{Tot}({\mathcal F}^\bullet\otimes {\mathcal G}^\bullet)
))
\end{equation}
此映射记为 $\cup$，并由下列规则给出：
$$
(\alpha \cup \beta)_{i_0 \ldots i_p}
=
\sum\nolimits_{r = 0}^p
\epsilon(n, m, p, r)
\alpha_{i_0 \ldots i_r} \otimes \beta_{i_r \ldots i_p}.
$$
，其中 $\alpha$ 的次数为 $n$，$\beta$ 的次数为 $m$，并且
$$
\epsilon(n, m, p, r) = (-1)^{(p + r)n + rp + r}.
$$
注意 $\epsilon(n, m, p, n) = 1$。因此，若
$\mathcal{F}^\bullet = \mathcal{F}[0]$ 是只由置于次数 $0$ 的
一个阿贝尔层 $\mathcal{F}$ 构成的复形，则 $\cup$ 的公式中没有符号
（因为此时除非 $r = n$，否则
$\alpha_{i_0 \ldots i_r} = 0$）。关于为何必须出现符号以及如何计算它，
见 Deligne 撰写的 \cite[Exposee XVII]{SGA4}。为验证
(\ref{cohomology-equation-needs-signs}) 是复形映射，必须证明
$$
d(\alpha \cup \beta) =
d(\alpha) \cup \beta +
(-1)^{\deg(\alpha)} \alpha \cup d(\beta)
$$
；这是由
$\text{Tot}(
\text{Tot}(\check{\mathcal{C}}^\bullet({\mathcal U}, {\mathcal F}^\bullet))
\otimes_{\mathbf Z}
\text{Tot}(\check{\mathcal{C}}^\bullet({\mathcal U}, {\mathcal G}^\bullet))
)$
上的微分之定义所要求的，该定义见《同调代数》定义
\ref{homology-definition-associated-simple-complex}。先计算
\begin{align*}
d(\alpha \cup \beta)_{i_0 \ldots i_{p + 1}}
= &
\sum\nolimits_{j = 0}^{p + 1}
(-1)^j
(\alpha \cup \beta)_{i_0 \ldots \hat i_j \ldots i_{p + 1}}
+
(-1)^{p + 1}
d_{{\mathcal F} \otimes {\mathcal G}}
((\alpha \cup \beta)_{i_0 \ldots i_{p + 1}})
\\
= &
\sum\nolimits_{j = 0}^{p + 1}
\sum\nolimits_{r = 0}^{j - 1}
(-1)^j \epsilon(n, m, p, r)
\alpha_{i_0 \ldots i_r} \otimes \beta_{i_r \ldots \hat i_j \ldots i_{p + 1}}
+ \\
&
\sum\nolimits_{j = 0}^{p + 1}
\sum\nolimits_{r = j + 1}^{p + 1}
(-1)^j \epsilon(n, m, p, r - 1)
\alpha_{i_0 \ldots \hat i_j \ldots i_r} \otimes \beta_{i_r \ldots i_{p + 1}}
+ \\
&
\sum\nolimits_{r = 0}^{p + 1}
(-1)^{p + 1} \epsilon(n, m, p + 1, r)
d_{{\mathcal F} \otimes {\mathcal G}}
(\alpha_{i_0 \ldots i_r} \otimes \beta_{i_r \ldots i_{p + 1}})
\end{align*}
并注意最后一项中的各和式项等于
$$
(-1)^{p + 1} \epsilon(n, m, p + 1, r)
\left(
d_{\mathcal F}(\alpha_{i_0 \ldots i_r}) \otimes 
\beta_{i_r \ldots i_{p + 1}} +
(-1)^{n - r}
\alpha_{i_0 \ldots i_r} \otimes d_{\mathcal G}(\beta_{i_r \ldots i_{p + 1}})
\right).
$$
，因为
$\deg_\mathcal{F}(\alpha_{i_0 \ldots i_r}) = n - r$。
另一方面，
\begin{align*}
(d(\alpha) \cup \beta)_{i_0\ldots i_{p + 1}}
= &
\sum\nolimits_{r = 0}^{p + 1}
\epsilon(n + 1, m, p + 1, r)
d(\alpha)_{i_0\ldots i_r} \otimes \beta_{i_r\ldots i_{p + 1}}
\\
= &
\sum\nolimits_{r = 0}^{p + 1}
\sum\nolimits_{j = 0}^{r}
\epsilon(n + 1, m, p + 1, r) (-1)^j
\alpha_{i_0\ldots\hat{i_j}\ldots i_r} \otimes \beta_{i_r\ldots i_{p + 1}}
+ \\
&
\sum\nolimits_{r = 0}^{p + 1}
\epsilon(n + 1, m, p + 1, r) (-1)^r
d_{\mathcal F}(\alpha_{i_0 \ldots i_r}) \otimes \beta_{i_r\ldots i_{p + 1}}
\end{align*}
且
\begin{align*}
(\alpha \cup d(\beta))_{i_0\ldots i_{p + 1}}
= &
\sum\nolimits_{r = 0}^{p + 1}
\epsilon(n, m + 1, p + 1, r)
\alpha_{i_0 \ldots i_r} \otimes d(\beta)_{i_r \ldots i_{p + 1}}
\\
= &
\sum\nolimits_{r = 0}^{p + 1}
\sum\nolimits_{j = r}^{p + 1}
\epsilon(n, m + 1, p + 1, r) (-1)^{j - r}
\alpha_{i_0 \ldots i_r} \otimes \beta_{i_r \ldots \hat{i_j}\ldots i_{p + 1}}
+ \\
&
\sum\nolimits_{r = 0}^{p + 1}
\epsilon(n, m + 1, p + 1, r) (-1)^{p + 1 - r}
\alpha_{i_0 \ldots i_r} \otimes d_{\mathcal G}(\beta_{i_r \ldots i_{p + 1}})
\end{align*}
若下列等式成立，则所需等式成立：
\begin{align*}
(-1)^{p + 1} \epsilon(n, m, p + 1, r)
& =
\epsilon(n + 1, m, p + 1, r) (-1)^r \\
(-1)^{p + 1} \epsilon(n, m, p + 1, r) (-1)^{n - r}
& =
(-1)^n \epsilon(n, m + 1, p + 1, r) (-1)^{p + 1 - r} \\
\epsilon(n + 1, m, p + 1, r) (-1)^r
& =
(-1)^{1 + n} \epsilon(n, m + 1, p + 1, r - 1) \\
(-1)^j \epsilon(n, m, p, r)
& =
(-1)^n \epsilon(n, m + 1, p + 1, r) (-1)^{j - r} \\
(-1)^j \epsilon(n, m, p, r - 1)
& =
\epsilon(n + 1, m, p + 1, r) (-1)^j
\end{align*}
（第三个等式是使 $d(\alpha) \cup \beta$ 与
$(-1)^n \alpha \cup d(\beta)$ 中 $r = j$ 的各项彼此相消所必需的。）
我们把验证留给读者。（也可检查 Stacks Project 的
$scripts$ 子目录中的脚本 $signs.gp$。）

\medskip\noindent
杯积的结合律。设 ${\mathcal F}^\bullet$、
${\mathcal G}^\bullet$ 与 ${\mathcal H}^\bullet$ 是 $X$ 上
阿贝尔群的有下界复形。不涉及任何符号的显然映射
$$
\text{Tot}(
\text{Tot}({\mathcal F}^\bullet \otimes_{\mathbf Z} {\mathcal G}^\bullet)
\otimes_{\mathbf Z} {\mathcal H}^\bullet
)
\longrightarrow
\text{Tot}(
{\mathcal F}^\bullet \otimes_{\mathbf Z}
\text{Tot}({\mathcal G}^\bullet \otimes_{\mathbf Z} {\mathcal H}^\bullet)
).
$$
是复形的同构。换言之，三重复形
${\mathcal F}^\bullet \otimes_{\mathbf Z} {\mathcal G}^\bullet
\otimes_{\mathbf Z} {\mathcal H}^\bullet$ 给出定义良好的总复形，
其微分满足
$$
d(\alpha \otimes \beta \otimes \gamma) =
d(\alpha) \otimes \beta \otimes \gamma +
(-1)^{\deg(\alpha)} \alpha \otimes d(\beta) \otimes \gamma +
(-1)^{\deg(\alpha) + \deg(\beta)} \alpha \otimes \beta \otimes d(\gamma)
$$
，其中各元素为齐次元素。利用此映射容易验证
$$
(\alpha \cup \beta) \cup \gamma = \alpha \cup ( \beta \cup \gamma)
$$
。事实上，若 $\alpha$ 的次数为 $a$，$\beta$ 的次数为 $b$，
而 $\gamma$ 的次数为 $c$，则
\begin{align*}
((\alpha \cup \beta) \cup \gamma)_{i_0 \ldots i_p}
= &
\sum\nolimits_{r = 0}^p
\epsilon(a + b, c, p, r)
(\alpha \cup \beta)_{i_0 \ldots i_r} \otimes \gamma_{i_r \ldots i_p}
\\
= &
\sum\nolimits_{r = 0}^p
\sum\nolimits_{s = 0}^r
\epsilon(a + b, c, p, r) \epsilon(a, b, r, s)
\alpha_{i_0 \ldots i_s} \otimes
\beta_{i_s \ldots i_r} \otimes
\gamma_{i_r \ldots i_p}
\end{align*}
并且
\begin{align*}
(\alpha \cup (\beta \cup \gamma)_{i_0\ldots i_p}
= &
\sum\nolimits_{s = 0}^p
\epsilon(a, b + c, p, s)
\alpha_{i_0 \ldots i_s} \otimes (\beta \cup \gamma)_{i_s \ldots i_p}
\\
= &
\sum\nolimits_{s = 0}^p
\sum\nolimits_{r = s}^p
\epsilon(a, b + c, p, s) \epsilon(b, c, p - s, r - s)
\alpha_{i_0 \ldots i_s} \otimes \beta_{i_s \ldots i_r} \otimes
\gamma_{i_r \ldots i_p}
\end{align*}
一个直接的模 $2$ 计算表明这些符号相符。
（也可检查 Stacks Project 的 $scripts$ 子目录中的脚本 $signs.gp$。）

\medskip\noindent
最后说明杯积为何至少在上同调层面保持分次交换结构。为此使用上面引入的
算子 $\tau$。设 ${\mathcal F}^\bullet$ 是阿贝尔群的有下界复形，
并假设给定分次交换乘法
$$
\wedge^\bullet :
\text{Tot}({\mathcal F}^\bullet\otimes {\mathcal F}^\bullet)
\longrightarrow
{\mathcal F}^\bullet.
$$
这意味着：若 $s$ 是 ${\mathcal F}^a$ 的局部截面，而 $t$ 是
${\mathcal F}^b$ 的局部截面，则 $s \wedge t$ 是
${\mathcal F}^{a + b}$ 的局部截面。分次交换是指
$s \wedge t = (-1)^{ab} t \wedge s$。由于 $\wedge$ 是复形映射，故有
$d(s\wedge t) = d(s) \wedge t + (-1)^a s \wedge d(t)$.
复合
$$
\text{Tot}(
\text{Tot}(\check{\mathcal{C}}^\bullet({\mathcal U}, {\mathcal F}^\bullet))
\otimes
\text{Tot}(\check{\mathcal{C}}^\bullet({\mathcal U}, {\mathcal F}^\bullet))
)
\to
\text{Tot}(
\check{\mathcal{C}}^\bullet({\mathcal U},
\text{Tot}({\mathcal F}^\bullet\otimes_{\mathbf Z}{\mathcal F}^\bullet))
)
\to
\text{Tot}(\check{\mathcal{C}}^\bullet({\mathcal U}, {\mathcal F}^\bullet))
$$
在上同调上诱导杯积
$$
H^n(
\text{Tot}(\check{\mathcal{C}}^\bullet({\mathcal U}, {\mathcal F}^\bullet))
)
\times
H^m(
\text{Tot}(\check{\mathcal{C}}^\bullet({\mathcal U}, {\mathcal F}^\bullet))
)
\longrightarrow
H^{n + m}(
\text{Tot}(\check{\mathcal{C}}^\bullet({\mathcal U}, {\mathcal F}^\bullet))
)
$$
，从而在余极限中也诱导 Cech 上同调上的乘法，并因而（必要时使用
超覆盖）诱导 ${\mathcal F}^\bullet$ 的上同调上的乘法。
我们断言这个（上同调上的）乘法也分次交换。为证明这一点，先取
$\text{Tot}(\check{\mathcal{C}}^\bullet({\mathcal U}, {\mathcal F}^\bullet))$
中次数为 $n$ 的元素 $\alpha$，以及
$\text{Tot}(\check{\mathcal{C}}^\bullet({\mathcal U}, {\mathcal F}^\bullet))$
中次数为 $m$ 的元素 $\beta$，并计算
\begin{align*}
\wedge^\bullet(\alpha \cup \beta)_{i_0 \ldots i_p}
= &
\sum\nolimits_{r = 0}^p
\epsilon(n, m, p, r)
\alpha_{i_0 \ldots i_r} \wedge \beta_{i_r \ldots i_p} \\
= &
\sum\nolimits_{r = 0}^p
\epsilon(n, m, p, r)
(-1)^{\deg(\alpha_{i_0 \ldots i_r})\deg(\beta_{i_r \ldots i_p})}
\beta_{i_r \ldots i_p} \wedge \alpha_{i_0 \ldots i_r}
\end{align*}
，这里使用了 $\wedge$ 的分次交换性。另一方面有
\begin{align*}
\tau(\wedge^\bullet(\tau(\beta) \cup \tau(\alpha)))_{i_0 \ldots i_p}
= &
\chi(p)
\sum\nolimits_{r = 0}^p
\epsilon(m, n, p, r)
\tau(\beta)_{i_p \ldots i_{p - r}} \wedge \tau(\alpha)_{i_{p - r} \ldots i_0}
\\
= &
\chi(p)
\sum\nolimits_{r = 0}^p
\epsilon(m, n, p, r) \chi(r) \chi(p - r)
\beta_{i_{p - r} \ldots i_p} \wedge \alpha_{i_0 \ldots i_{p - r}}
\\
= &
\chi(p)
\sum\nolimits_{r = 0}^p
\epsilon(m, n, p, p - r) \chi(r) \chi(p - r)
\beta_{i_r \ldots i_p} \wedge \alpha_{i_0 \ldots i_r}
\end{align*}
，其中 $\chi(t) = (-1)^{\frac{t(t + 1)}{2}}$。前面已经证明
$\tau$ 在上同调上按恒等映射作用，因此只需验证
$$
\epsilon(n, m, p, r)
(-1)^{(n - r)(m - (p - r))}
=
(-1)^{nm} \chi(p)\epsilon(m, n, p, p - r) \chi(r) \chi(p - r)
$$
一个直接的模 $2$ 计算表明这些符号相符。
（也可检查 Stacks Project 的 $scripts$ 子目录中的脚本 $signs.gp$。）

\medskip\noindent
最后研究杯积与连接映射的相容性。设
$$
0
\to
{\mathcal F}_1^\bullet
\to
{\mathcal F}_2^\bullet
\to
{\mathcal F}_3^\bullet
\to
0
\quad\text{且}\quad
0
\leftarrow
{\mathcal G}_1^\bullet
\leftarrow
{\mathcal G}_2^\bullet
\leftarrow
{\mathcal G}_3^\bullet
\leftarrow
0
$$
是 $X$ 上阿贝尔层的有下界复形的短正合列。再设
${\mathcal H}^\bullet$ 是阿贝尔层的另一个有下界复形，并假设有复形映射
$$
\gamma_i :
\text{Tot}({\mathcal F}_i^\bullet \otimes_{\mathbf Z} {\mathcal G}_i^\bullet)
\longrightarrow
{\mathcal H}^\bullet
$$
，它们与这些复形之间的映射相容，即下列各图
$$
\xymatrix{
\text{Tot}({\mathcal F}_1^\bullet \otimes_{\mathbf Z} {\mathcal G}_1^\bullet)
\ar[d]_{\gamma_1}
&
\text{Tot}({\mathcal F}_1^\bullet \otimes_{\mathbf Z} {\mathcal G}_2^\bullet)
\ar[l] \ar[d]
\\
\mathcal{H}^\bullet &
\text{Tot}({\mathcal F}_2^\bullet \otimes_{\mathbf Z} {\mathcal G}_2^\bullet)
\ar[l]_-{\gamma_2}
}
$$
且
$$
\xymatrix{
\text{Tot}({\mathcal F}_2^\bullet \otimes_{\mathbf Z} {\mathcal G}_2^\bullet)
\ar[d]_{\gamma_2}
&
\text{Tot}({\mathcal F}_2^\bullet \otimes_{\mathbf Z} {\mathcal G}_3^\bullet)
\ar[l] \ar[d]
\\
\mathcal{H}^\bullet &
\text{Tot}({\mathcal F}_3^\bullet \otimes_{\mathbf Z} {\mathcal G}_3^\bullet)
\ar[l]_-{\gamma_3}
}
$$
交换。

\begin{lemma}
\label{cohomology-lemma-compute-sign-cup-product-boundaries}
在上述情形中，假设对各层 $\mathcal{F}_i^p$ 与
$\mathcal{G}_j^q$，Cech 上同调与上同调一致。设
$a_3 \in H^n(X, \mathcal{F}_3^\bullet)$ 且
$b_1 \in H^m(X, \mathcal{G}_1^\bullet)$。则有
$$
\gamma_1( \partial a_3 \cup b_1) =
(-1)^{n + 1} \gamma_3( a_3 \cup \partial b_1)
$$
在 $H^{n + m + 1}(X, \mathcal{H}^\bullet)$ 中成立，其中
$\partial$ 表示由上述复形短正合列所诱导的上同调连接映射。
\end{lemma}

\begin{proof}
采用下列约定与记号。把 ${\mathcal F}_1^p$ 看成
${\mathcal F}_2^p$ 的子层，并把 ${\mathcal G}_3^q$ 看成
${\mathcal G}_2^q$ 的子层。因此，若 $s$ 是
${\mathcal F}_1^p$ 的局部截面，也用 $s$ 表示
${\mathcal F}_2^p$ 的相应截面；对 ${\mathcal G}_3^q$ 的局部截面
也作同样约定。此外，若 $s$ 是 ${\mathcal F}_2^p$ 的局部截面，
则用 $\bar s$ 表示它在 ${\mathcal F}_3^p$ 中的像。
对映射 ${\mathcal G}_2^q \to {\mathcal G}^q_1$ 也作同样约定。
特别地，若 $s$ 是 ${\mathcal F}_2^p$ 的局部截面且
$\bar s = 0$，则 $s$ 是 ${\mathcal F}_1^p$ 的局部截面。
上述图的交换性蕴含：对 ${\mathcal F}_2^p$ 的局部截面 $s$
和 ${\mathcal G}_3^q$ 的局部截面 $t$，作为
${\mathcal H}^{p + q}$ 的截面有
$\gamma_2(s \otimes t) = \gamma_3(\bar s \otimes t)$。

\medskip\noindent
设 ${\mathcal U} : X =  \bigcup_{i \in I} U_i$ 是 $X$ 的开覆盖。
假设 $\alpha_3$（相应地 $\beta_1$）是
$\text{Tot}(
\check{\mathcal{C}}^\bullet({\mathcal U}, {\mathcal F}_3^\bullet))$
（相应地
$\text{Tot}(
\check{\mathcal{C}}^\bullet({\mathcal U}, {\mathcal G}_1^\bullet))$）
中次数为 $n$（相应地 $m$）的上循环，代表
$a_3$（相应地 $b_1$）。必要时加细 $\mathcal{U}$ 后，可以找到
$\text{Tot}(
\check{\mathcal{C}}^\bullet({\mathcal U}, {\mathcal F}_2^\bullet))$
中次数为 $n$ 的上链 $\alpha_2$，以及
$\text{Tot}(
\check{\mathcal{C}}^\bullet({\mathcal U}, {\mathcal G}_2^\bullet))$
中次数为 $m$ 的上链 $\beta_2$，它们分别映到
$\alpha_3$ 与 $\beta_1$。于是
$$
\overline{d(\alpha_2)} = d(\bar \alpha_2) = 0
\quad\text{且}\quad
\overline{d(\beta_2)} = d(\bar \beta_2) = 0.
$$
这意味着 $\alpha_1 = d(\alpha_2)$ 是
$\text{Tot}(\check{\mathcal{C}}^\bullet({\mathcal U}, {\mathcal F}_1^\bullet))$
中次数为 $n + 1$ 的上循环，代表 $\partial a_3$。类似地，
$\beta_3 = d(\beta_2)$ 是
$\text{Tot}(\check{\mathcal{C}}^\bullet({\mathcal U}, {\mathcal G}_3^\bullet))$
中次数为 $m + 1$ 的上循环，代表 $\partial b_1$。
因此可作如下计算：
\begin{align*}
d(\gamma_2(\alpha_2 \cup \beta_2))
& =
\gamma_2(d(\alpha_2 \cup \beta_2))
\\
& =
\gamma_2(d(\alpha_2) \cup \beta_2 + (-1)^n \alpha_2 \cup d(\beta_2) )
\\
& =
\gamma_2( \alpha_1 \cup \beta_2)  + (-1)^n \gamma_2( \alpha_2 \cup \beta_3)
\\
& =
\gamma_1(\alpha_1 \cup \beta_1) + (-1)^n \gamma_3(\alpha_3 \cup \beta_3)
\end{align*}
这甚至表明符号确如引理所述为 $(-1)^{n + 1}$
\footnote{该符号取决于与 $D(X)$ 中三角所对应的上同调长正合列的
符号约定。Stacks Project 的约定是：（a）特异三角对应于逐项分裂的
正合列；（b）长正合列中的连接映射由蛇引理中的映射给出，不另加符号。
见《导出范畴》第 \ref{derived-section-homotopy-triangulated} 节。}。
\end{proof}

\begin{lemma}
\label{cohomology-lemma-boundary-derivation-over-cup-product}
设 $X$ 是拓扑空间。设 $\mathcal{O}' \to \mathcal{O}$ 是环层的满射，
其核 $\mathcal{I} \subset \mathcal{O}'$ 的平方为零。则
$M = H^1(X, \mathcal{I})$ 是 $R = H^0(X, \mathcal{O})$-模，
并且由短正合列
$$
0 \to \mathcal{I} \to \mathcal{O}' \to \mathcal{O} \to 0
$$
所诱导的连接映射 $\partial : R \to M$ 是导子
（《代数》定义 \ref{algebra-definition-derivation}）。
\end{lemma}

\begin{proof}
由假设 $\mathcal{I} \cdot \mathcal{I} = 0$，映射
$\mathcal{O}' \to \SheafHom(\mathcal{I}, \mathcal{I})$
通过 $\mathcal{O}$ 分解。因此 $\mathcal{I}$ 是
$\mathcal{O}$-模层，这定义了 $M$ 上的 $R$-模结构。
连接映射是可加的，故只需证明 Leibniz 法则。设 $f \in R$。
选取开覆盖 $\mathcal{U} : X = \bigcup U_i$，使得存在
$f_i \in \mathcal{O}'(U_i)$ 提升
$f|_{U_i} \in \mathcal{O}(U_i)$。注意
$f_i - f_j$ 是 $\mathcal{I}(U_i \cap U_j)$ 的元素。
于是 $\partial(f)$ 对应于上循环 $\alpha$ 的 Cech 上同调类，
其中 $\alpha_{i_0i_1} = f_{i_0} - f_{i_1}$。
（注意，由引理 \ref{cohomology-lemma-cech-h1}，关于 $\mathcal{U}$ 的第一
Cech 上同调群是 $M$ 的子模。）
接着取另一个元素 $g \in R$，并假设（必要时加细开覆盖后）
$g_i \in \mathcal{O}'(U_i)$ 提升
$g|_{U_i} \in \mathcal{O}(U_i)$。于是
$\partial(g)$ 由上循环 $\beta$ 给出，其中
$\beta_{i_0i_1} = g_{i_0} - g_{i_1}$。由于
$f_ig_i \in \mathcal{O}'(U_i)$ 提升 $fg|_{U_i}$，
可知 $\partial(fg)$ 由上循环 $\gamma$ 给出，其中
$$
\gamma_{i_0i_1} = f_{i_0}g_{i_0} - f_{i_1}g_{i_1} =
(f_{i_0} - f_{i_1})g_{i_0} + f_{i_1}(g_{i_0} - g_{i_1}) =
\alpha_{i_0i_1}g + f\beta_{i_0i_1}
$$
，这里使用了我们对 $\mathcal{I}$ 上 $\mathcal{O}$-模结构的定义。
这就证明了 Leibniz 法则，证明完毕。
\end{proof}









\section{平坦消解}
\label{cohomology-section-flat}

\noindent
本节材料可参考 \cite{Spaltenstein}。设
$(X, \mathcal{O}_X)$ 是环空间。由《模层》引理
\ref{modules-lemma-module-quotient-flat}，任意 $\mathcal{O}_X$-模
都是某个平坦 $\mathcal{O}_X$-模的商。由《导出范畴》引理
\ref{derived-lemma-subcategory-left-resolution}，任意
$\mathcal{O}_X$-模的有上界复形都具有由平坦
$\mathcal{O}_X$-模的有上界复形给出的左消解。然而，对无界复形而言，
事实表明平坦消解还不够好。

\begin{lemma}
\label{cohomology-lemma-derived-tor-exact}
设 $(X, \mathcal{O}_X)$ 是环空间，而
$\mathcal{G}^\bullet$ 是 $\mathcal{O}_X$-模复形。函子
$$
K(\textit{Mod}(\mathcal{O}_X))
\longrightarrow
K(\textit{Mod}(\mathcal{O}_X)),
\quad
\mathcal{F}^\bullet \longmapsto
\text{Tot}(\mathcal{G}^\bullet \otimes_{\mathcal{O}_X} \mathcal{F}^\bullet)
$$
以及
$$
K(\textit{Mod}(\mathcal{O}_X))
\longrightarrow
K(\textit{Mod}(\mathcal{O}_X)),
\quad
\mathcal{F}^\bullet \longmapsto
\text{Tot}(\mathcal{F}^\bullet \otimes_{\mathcal{O}_X} \mathcal{G}^\bullet)
$$
是三角范畴之间的正合函子。
\end{lemma}

\begin{proof}
这由《导出范畴》注记
\ref{derived-remark-double-complex-as-tensor-product-of} 得到。
\end{proof}

\begin{definition}
\label{cohomology-definition-K-flat}
设 $(X, \mathcal{O}_X)$ 是环空间。称 $\mathcal{O}_X$-模复形
$\mathcal{K}^\bullet$ 为 {\it K-平坦的}，如果对任意
$\mathcal{O}_X$-模的无环复形 $\mathcal{F}^\bullet$，复形
$$
\text{Tot}(\mathcal{F}^\bullet \otimes_{\mathcal{O}_X} \mathcal{K}^\bullet)
$$
都是无环的。
\end{definition}

\begin{lemma}
\label{cohomology-lemma-K-flat-quasi-isomorphism}
设 $(X, \mathcal{O}_X)$ 是环空间，而
$\mathcal{K}^\bullet$ 是 K-平坦复形。则函子
$$
K(\textit{Mod}(\mathcal{O}_X))
\longrightarrow
K(\textit{Mod}(\mathcal{O}_X)), \quad
\mathcal{F}^\bullet
\longmapsto
\text{Tot}(\mathcal{F}^\bullet \otimes_{\mathcal{O}_X} \mathcal{K}^\bullet)
$$
将拟同构映为拟同构。
\end{lemma}

\begin{proof}
这由引理 \ref{cohomology-lemma-derived-tor-exact} 以及拟同构可由其锥无环刻画这一事实得到。
\end{proof}

\begin{lemma}
\label{cohomology-lemma-check-K-flat-stalks}
设 $(X, \mathcal{O}_X)$ 是环空间，且 $\mathcal{K}^\bullet$
是 $\mathcal{O}_X$-模复形。则 $\mathcal{K}^\bullet$ 是
K-平坦的，当且仅当对所有 $x \in X$，
$\mathcal{O}_{X, x}$-模复形 $\mathcal{K}_x^\bullet$ 都是
K-平坦的（《代数详论》定义 \ref{more-algebra-definition-K-flat}）。
\end{lemma}

\begin{proof}
若对所有 $x \in X$，$\mathcal{K}_x^\bullet$ 都是 K-平坦的，
则 $\mathcal{K}^\bullet$ 是 K-平坦的；这是因为 $\otimes$ 与直和
都和取茎交换，并且正合性可在茎上检验，见《模层》引理
\ref{modules-lemma-abelian}。反过来，假设
$\mathcal{K}^\bullet$ 是 K-平坦的。取 $x \in X$，并设
$M^\bullet$ 是 $\mathcal{O}_{X, x}$-模的无环复形。于是
$i_{x, *}M^\bullet$ 是 $\mathcal{O}_X$-模的无环复形，因而
$\text{Tot}(i_{x, *}M^\bullet \otimes_{\mathcal{O}_X}
\mathcal{K}^\bullet)$ 无环。在 $x$ 处取茎可知
$\text{Tot}(M^\bullet \otimes_{\mathcal{O}_{X, x}} \mathcal{K}_x^\bullet)$
无环。
\end{proof}

\begin{lemma}
\label{cohomology-lemma-tensor-product-K-flat}
设 $(X, \mathcal{O}_X)$ 是环空间。若
$\mathcal{K}^\bullet$、$\mathcal{L}^\bullet$ 是
$\mathcal{O}_X$-模的 K-平坦复形，则
$\text{Tot}(\mathcal{K}^\bullet \otimes_{\mathcal{O}_X} \mathcal{L}^\bullet)$
是 $\mathcal{O}_X$-模的 K-平坦复形。
\end{lemma}

\begin{proof}
这由同构
$$
\text{Tot}(\mathcal{M}^\bullet \otimes_{\mathcal{O}_X}
\text{Tot}(\mathcal{K}^\bullet \otimes_{\mathcal{O}_X} \mathcal{L}^\bullet))
=
\text{Tot}(\text{Tot}(\mathcal{M}^\bullet \otimes_{\mathcal{O}_X}
\mathcal{K}^\bullet) \otimes_{\mathcal{O}_X} \mathcal{L}^\bullet)
$$
和定义得到。
\end{proof}

\begin{lemma}
\label{cohomology-lemma-K-flat-two-out-of-three}
设 $(X, \mathcal{O}_X)$ 是环空间。设
$(\mathcal{K}_1^\bullet, \mathcal{K}_2^\bullet,
\mathcal{K}_3^\bullet)$ 是
$K(\textit{Mod}(\mathcal{O}_X))$ 中的特异三角。
若三个 $\mathcal{K}_i^\bullet$ 中有两个是 K-平坦的，
则第三个也是 K-平坦的。
\end{lemma}

\begin{proof}
这由引理 \ref{cohomology-lemma-derived-tor-exact} 以及下述事实得到：
在 $K(\textit{Mod}(\mathcal{O}_X))$ 的特异三角中，
若三者中有两个无环，则第三个也无环。
\end{proof}

\begin{lemma}
\label{cohomology-lemma-K-flat-two-out-of-three-ses}
设 $(X, \mathcal{O}_X)$ 是环空间。设
$0 \to \mathcal{K}_1^\bullet \to \mathcal{K}_2^\bullet \to
\mathcal{K}_3^\bullet \to 0$ 是复形的短正合列，并且
$\mathcal{K}_3^\bullet$ 的各项都是平坦
$\mathcal{O}_X$-模。若三个 $\mathcal{K}_i^\bullet$ 中有两个
是 K-平坦的，则第三个也是 K-平坦的。
\end{lemma}

\begin{proof}
由《模层》引理 \ref{modules-lemma-flat-tor-zero}，对每个复形
$\mathcal{L}^\bullet$ 都得到复形的短正合列
$$
0 \to
\text{Tot}(\mathcal{L}^\bullet \otimes_{\mathcal{O}_X} \mathcal{K}_1^\bullet)
\to
\text{Tot}(\mathcal{L}^\bullet \otimes_{\mathcal{O}_X} \mathcal{K}_1^\bullet)
\to
\text{Tot}(\mathcal{L}^\bullet \otimes_{\mathcal{O}_X} \mathcal{K}_1^\bullet)
\to 0
$$
。因此该引理由上同调层长正合列和 K-平坦复形的定义得到。
\end{proof}

\begin{lemma}
\label{cohomology-lemma-pullback-K-flat}
设 $f : (X, \mathcal{O}_X) \to (Y, \mathcal{O}_Y)$ 是环空间的态射。
$\mathcal{O}_Y$-模的 K-平坦复形的拉回是
$\mathcal{O}_X$-模的 K-平坦复形。
\end{lemma}

\begin{proof}
可在茎上检验，见引理 \ref{cohomology-lemma-check-K-flat-stalks}。
因此这由《层论》引理 \ref{sheaves-lemma-stalk-pullback-modules}
和《代数详论》引理 \ref{more-algebra-lemma-base-change-K-flat} 得到。
\end{proof}

\begin{lemma}
\label{cohomology-lemma-bounded-flat-K-flat}
设 $(X, \mathcal{O}_X)$ 是环空间。平坦
$\mathcal{O}_X$-模的有上界复形是 K-平坦的。
\end{lemma}

\begin{proof}
可在茎上检验，见引理 \ref{cohomology-lemma-check-K-flat-stalks}。
故该引理由《模层》引理 \ref{modules-lemma-flat-stalks-flat}
和《代数详论》引理
\ref{more-algebra-lemma-derived-tor-quasi-isomorphism} 得到。
\end{proof}

\noindent
在下列引理中，复形系统的余极限是指逐项余极限。

\begin{lemma}
\label{cohomology-lemma-colimit-K-flat}
设 $(X, \mathcal{O}_X)$ 是环空间，且
$\mathcal{K}_1^\bullet \to \mathcal{K}_2^\bullet \to \ldots$
是 K-平坦复形系统。则
$\colim_i \mathcal{K}_i^\bullet$ 是 K-平坦的。
\end{lemma}

\begin{proof}
由于取的是逐项余极限，显然
$$
\colim_i \text{Tot}(
\mathcal{F}^\bullet \otimes_{\mathcal{O}_X} \mathcal{K}_i^\bullet)
=
\text{Tot}(\mathcal{F}^\bullet \otimes_{\mathcal{O}_X}
\colim_i \mathcal{K}_i^\bullet)
$$
因此该引理由滤过余极限为正合这一事实得到。
\end{proof}

\begin{lemma}
\label{cohomology-lemma-resolution-by-direct-sums-extensions-by-zero}
设 $(X, \mathcal{O}_X)$ 是环空间。对任意
$\mathcal{O}_X$-模复形 $\mathcal{G}^\bullet$，存在
$\mathcal{O}_X$-模复形的交换图
$$
\xymatrix{
\mathcal{K}_1^\bullet \ar[d] \ar[r] &
\mathcal{K}_2^\bullet \ar[d] \ar[r] & \ldots \\
\tau_{\leq 1}\mathcal{G}^\bullet \ar[r] &
\tau_{\leq 2}\mathcal{G}^\bullet \ar[r] & \ldots
}
$$
，满足下列性质：（1）竖直箭头是拟同构且逐项满射；
（2）每个 $\mathcal{K}_n^\bullet$ 都是有上界复形，其各项为形如
$j_{U!}\mathcal{O}_U$ 的 $\mathcal{O}_X$-模的直和；
（3）映射
$\mathcal{K}_n^\bullet \to \mathcal{K}_{n + 1}^\bullet$
是逐项分裂单射，其余核为形如 $j_{U!}\mathcal{O}_U$ 的
$\mathcal{O}_X$-模的直和。此外，映射
$\colim \mathcal{K}_n^\bullet \to \mathcal{G}^\bullet$ 是拟同构。
\end{lemma}

\begin{proof}
该图以及性质（1）、（2）、（3）的存在性立即由《模层》引理
\ref{modules-lemma-module-quotient-flat} 和《导出范畴》引理
\ref{derived-lemma-special-direct-system} 得到。由于滤过余极限正合，
诱导映射
$\colim \mathcal{K}_n^\bullet \to \mathcal{G}^\bullet$ 是拟同构。
\end{proof}

\begin{lemma}
\label{cohomology-lemma-K-flat-resolution}
设 $(X, \mathcal{O}_X)$ 是环空间。对任意复形
$\mathcal{G}^\bullet$，存在各项为平坦 $\mathcal{O}_X$-模的
$K$-平坦复形 $\mathcal{K}^\bullet$，以及逐项满射的拟同构
$\mathcal{K}^\bullet \to \mathcal{G}^\bullet$。
\end{lemma}

\begin{proof}
选取引理 \ref{cohomology-lemma-resolution-by-direct-sums-extensions-by-zero}
中的图。每个复形 $\mathcal{K}_n^\bullet$ 都是平坦模的有上界复形，
见《模层》引理 \ref{modules-lemma-j-shriek-flat}。故由引理
\ref{cohomology-lemma-bounded-flat-K-flat}，$\mathcal{K}_n^\bullet$ 是
K-平坦的；再由引理 \ref{cohomology-lemma-colimit-K-flat}，
$\colim \mathcal{K}_n^\bullet$ 是 K-平坦的。按构造，诱导映射
$\colim \mathcal{K}_n^\bullet \to \mathcal{G}^\bullet$
是拟同构且逐项满射。引理
\ref{cohomology-lemma-resolution-by-direct-sums-extensions-by-zero} 的性质（3）
表明 $\colim \mathcal{K}_n^m$ 是平坦模的直和，因而平坦；
这就证明了最后一个断言。
\end{proof}

\begin{lemma}
\label{cohomology-lemma-derived-tor-quasi-isomorphism-other-side}
设 $(X, \mathcal{O}_X)$ 是环空间。设
$\alpha : \mathcal{P}^\bullet \to \mathcal{Q}^\bullet$ 是
$\mathcal{O}_X$-模的 K-平坦复形之间的拟同构。对任意
$\mathcal{O}_X$-模复形 $\mathcal{F}^\bullet$，诱导映射
$$
\text{Tot}(\text{id}_{\mathcal{F}^\bullet} \otimes \alpha) :
\text{Tot}(\mathcal{F}^\bullet \otimes_{\mathcal{O}_X} \mathcal{P}^\bullet)
\longrightarrow
\text{Tot}(\mathcal{F}^\bullet \otimes_{\mathcal{O}_X} \mathcal{Q}^\bullet)
$$
是拟同构。
\end{lemma}

\begin{proof}
选取拟同构 $\mathcal{K}^\bullet \to \mathcal{F}^\bullet$，
其中 $\mathcal{K}^\bullet$ 是 K-平坦复形，见引理
\ref{cohomology-lemma-K-flat-resolution}。考虑交换图
$$
\xymatrix{
\text{Tot}(\mathcal{K}^\bullet
\otimes_{\mathcal{O}_X} \mathcal{P}^\bullet) \ar[r] \ar[d] &
\text{Tot}(\mathcal{K}^\bullet
\otimes_{\mathcal{O}_X} \mathcal{Q}^\bullet) \ar[d] \\
\text{Tot}(\mathcal{F}^\bullet
\otimes_{\mathcal{O}_X} \mathcal{P}^\bullet) \ar[r] &
\text{Tot}(\mathcal{F}^\bullet
\otimes_{\mathcal{O}_X} \mathcal{Q}^\bullet)
}
$$
由引理 \ref{cohomology-lemma-K-flat-quasi-isomorphism}，竖直箭头和上方水平箭头
都是拟同构，故结论成立。
\end{proof}

\noindent
设 $(X, \mathcal{O}_X)$ 是环空间，且
$\mathcal{F}^\bullet$ 是 $D(\mathcal{O}_X)$ 的对象。
选取 K-平坦消解
$\mathcal{K}^\bullet \to \mathcal{F}^\bullet$，见引理
\ref{cohomology-lemma-K-flat-resolution}。由引理
\ref{cohomology-lemma-derived-tor-exact} 得到三角范畴的正合函子
$$
K(\mathcal{O}_X)
\longrightarrow
K(\mathcal{O}_X),
\quad
\mathcal{G}^\bullet
\longmapsto
\text{Tot}(\mathcal{G}^\bullet \otimes_{\mathcal{O}_X} \mathcal{K}^\bullet)
$$
由引理 \ref{cohomology-lemma-K-flat-quasi-isomorphism}，此函子诱导函子
$D(\mathcal{O}_X) \to D(\mathcal{O}_X)$；这是因为
$D(\mathcal{O}_X)$ 是 $K(\mathcal{O}_X)$ 关于拟同构的局部化。
由于 $\mathcal{F}^\bullet$ 的 $K$-平坦消解所成的范畴是余滤过的，
再结合引理 \ref{cohomology-lemma-derived-tor-quasi-isomorphism-other-side}，
所得函子（在同构意义下）与 $\mathcal{K}^\bullet$ 的选取无关。

\begin{definition}
\label{cohomology-definition-derived-tor}
设 $(X, \mathcal{O}_X)$ 是环空间，而
$\mathcal{F}^\bullet$ 是 $D(\mathcal{O}_X)$ 的对象。{\it 导出张量积}
$$
- \otimes_{\mathcal{O}_X}^{\mathbf{L}} \mathcal{F}^\bullet :
D(\mathcal{O}_X)
\longrightarrow
D(\mathcal{O}_X)
$$
是上面所述的三角范畴正合函子。
\end{definition}

\noindent
由上述显式构造，显然有典范同构
$$
\mathcal{F}^\bullet \otimes_{\mathcal{O}_X}^{\mathbf{L}} \mathcal{G}^\bullet
\cong
\mathcal{G}^\bullet \otimes_{\mathcal{O}_X}^{\mathbf{L}} \mathcal{F}^\bullet
$$
，其中 $\mathcal{G}^\bullet$ 与 $\mathcal{F}^\bullet$ 属于
$D(\mathcal{O}_X)$。这里使用《代数详论》第
\ref{more-algebra-section-sign-rules} 节给出的符号规则。因此写
$\mathcal{F}^\bullet \otimes_{\mathcal{O}_X}^{\mathbf{L}}
\mathcal{G}^\bullet$ 时，我们通常不区分使用哪一个变量来定义导出张量积。

\begin{definition}
\label{cohomology-definition-tor}
设 $(X, \mathcal{O}_X)$ 是环空间，而
$\mathcal{F}$、$\mathcal{G}$ 是 $\mathcal{O}_X$-模。
$\mathcal{F}$ 与 $\mathcal{G}$ 的 {\it Tor} 由公式
$$
\text{Tor}_p^{\mathcal{O}_X}(\mathcal{F}, \mathcal{G}) =
H^{-p}(\mathcal{F} \otimes_{\mathcal{O}_X}^\mathbf{L} \mathcal{G})
$$
定义，其中导出张量积如上定义。
\end{definition}

\noindent
此定义蕴含：对 $\mathcal{O}_X$-模的任意短正合列
$0 \to \mathcal{F}_1 \to \mathcal{F}_2 \to
\mathcal{F}_3 \to 0$，都有上同调长正合列
$$
\xymatrix{
\mathcal{F}_1 \otimes_{\mathcal{O}_X} \mathcal{G} \ar[r] &
\mathcal{F}_2 \otimes_{\mathcal{O}_X} \mathcal{G} \ar[r] &
\mathcal{F}_3 \otimes_{\mathcal{O}_X} \mathcal{G} \ar[r] & 0 \\
\text{Tor}_1^{\mathcal{O}_X}(\mathcal{F}_1, \mathcal{G}) \ar[r] &
\text{Tor}_1^{\mathcal{O}_X}(\mathcal{F}_2, \mathcal{G}) \ar[r] &
\text{Tor}_1^{\mathcal{O}_X}(\mathcal{F}_3, \mathcal{G}) \ar[ull]
}
$$
，其中 $\mathcal{G}$ 是任意 $\mathcal{O}_X$-模。称其为与该情形
相伴的 $\text{Tor}$ 长正合列。

\begin{lemma}
\label{cohomology-lemma-flat-tor-zero}
\begin{slogan}
Tor 衡量偏离平坦性的程度。
\end{slogan}
设 $(X, \mathcal{O}_X)$ 是环空间，而
$\mathcal{F}$ 是 $\mathcal{O}_X$-模。下列条件等价：
\begin{enumerate}
\item $\mathcal{F}$ 是平坦 $\mathcal{O}_X$-模；
\item $\text{Tor}_1^{\mathcal{O}_X}(\mathcal{F}, \mathcal{G}) = 0$
对每个 $\mathcal{O}_X$-模 $\mathcal{G}$ 都成立。
\end{enumerate}
\end{lemma}

\begin{proof}
若 $\mathcal{F}$ 平坦，则
$\mathcal{F} \otimes_{\mathcal{O}_X} -$ 是正合函子，且卫星函子消失。
反过来，假设（2）成立。若
$\mathcal{G} \to \mathcal{H}$ 是余核为 $\mathcal{Q}$ 的单射，
则 $\text{Tor}$ 长正合列表明映射
$\mathcal{F} \otimes_{\mathcal{O}_X} \mathcal{G} \to
\mathcal{F} \otimes_{\mathcal{O}_X} \mathcal{H}$
的核是
$\text{Tor}_1^{\mathcal{O}_X}(\mathcal{F}, \mathcal{Q})$
的商；后者依假设为零。因此 $\mathcal{F}$ 平坦。
\end{proof}

\begin{lemma}
\label{cohomology-lemma-factor-through-K-flat}
设 $(X, \mathcal{O}_X)$ 是环空间。设
$a : \mathcal{K}^\bullet \to \mathcal{L}^\bullet$ 是
$\mathcal{O}_X$-模复形的映射。若 $\mathcal{K}^\bullet$
是 K-平坦的，则存在复形 $\mathcal{N}^\bullet$ 和复形映射
$b : \mathcal{K}^\bullet \to \mathcal{N}^\bullet$、
$c : \mathcal{N}^\bullet \to \mathcal{L}^\bullet$，使得
\begin{enumerate}
\item $\mathcal{N}^\bullet$ 是 K-平坦的；
\item $c$ 是拟同构；
\item $a$ 与 $c \circ b$ 链同伦。
\end{enumerate}
若 $\mathcal{K}^\bullet$ 的各项平坦，则可选取
$\mathcal{N}^\bullet$、$b$、$c$，使得
$\mathcal{N}^\bullet$ 的各项也平坦。
\end{lemma}

\begin{proof}
我们将使用同伦范畴 $K(\textit{Mod}(\mathcal{O}_X))$ 为三角范畴这一事实，
见《导出范畴》命题
\ref{derived-proposition-homotopy-category-triangulated}。
选取特异三角
$\mathcal{K}^\bullet \to \mathcal{L}^\bullet \to
\mathcal{C}^\bullet \to \mathcal{K}^\bullet[1]$.
选取拟同构 $\mathcal{M}^\bullet \to \mathcal{C}^\bullet$，
其中 $\mathcal{M}^\bullet$ 为各项平坦的 K-平坦复形，见引理
\ref{cohomology-lemma-K-flat-resolution}。由三角范畴的公理，可把复合
$\mathcal{M}^\bullet \to \mathcal{C}^\bullet \to \mathcal{K}^\bullet[1]$
嵌入特异三角
$\mathcal{K}^\bullet \to \mathcal{N}^\bullet \to
\mathcal{M}^\bullet \to \mathcal{K}^\bullet[1]$.
由引理 \ref{cohomology-lemma-K-flat-two-out-of-three}，
$\mathcal{N}^\bullet$ 是 K-平坦的。再次使用三角范畴的公理，
可选取映射 $\mathcal{N}^\bullet \to \mathcal{L}^\bullet$，
把它嵌入下列特异三角的态射：
$$
\xymatrix{
\mathcal{K}^\bullet \ar[r] \ar[d] &
\mathcal{N}^\bullet \ar[r] \ar[d] &
\mathcal{M}^\bullet \ar[r] \ar[d] &
\mathcal{K}^\bullet[1] \ar[d] \\
\mathcal{K}^\bullet \ar[r] &
\mathcal{L}^\bullet \ar[r] &
\mathcal{C}^\bullet \ar[r] &
\mathcal{K}^\bullet[1]
}
$$
由于三个箭头中有两个是拟同构，由这些特异三角所伴随的上同调长正合列，
第三个箭头 $\mathcal{N}^\bullet \to \mathcal{L}^\bullet$ 也是拟同构
（也可考察此图在 $D(\mathcal{O}_X)$ 中的像，并使用《导出范畴》引理
\ref{derived-lemma-third-isomorphism-triangle}）。这就证明了
（1）、（2）、（3）。为证明最后一个断言，可以选取
$\mathcal{N}^\bullet$ 使得
$\mathcal{N}^n \cong \mathcal{M}^n \oplus \mathcal{K}^n$，见《导出范畴》引理
\ref{derived-lemma-improve-distinguished-triangle-homotopy}。
因此，若 $\mathcal{K}^\bullet$ 的各项平坦，就得到所需的平坦性。
\end{proof}











\section{导出拉回}
\label{cohomology-section-derived-pullback}

\noindent
设 $f : (X, \mathcal{O}_X) \to (Y, \mathcal{O}_Y)$
是环空间的态射。可以用 K-平坦消解定义导出拉回函子
$$
Lf^* : D(\mathcal{O}_Y) \to D(\mathcal{O}_X)
$$
。具体地，对任意 $\mathcal{O}_Y$-模复形
$\mathcal{G}^\bullet$，可选取 K-平坦消解
$\mathcal{K}^\bullet \to \mathcal{G}^\bullet$，并令
$Lf^*\mathcal{G}^\bullet = f^*\mathcal{K}^\bullet$。
利用引理 \ref{cohomology-lemma-pullback-K-flat}、
\ref{cohomology-lemma-K-flat-resolution} 和
\ref{cohomology-lemma-derived-tor-quasi-isomorphism-other-side} 可知此定义良好。
不过，为把所有细节完整处理，使用一些一般理论或许更为方便。

\begin{lemma}
\label{cohomology-lemma-derived-base-change}
上述构造与选择无关，并定义三角范畴的正合函子
$Lf^* : D(\mathcal{O}_Y) \to D(\mathcal{O}_X)$.
\end{lemma}

\begin{proof}
为此使用《导出范畴》第 \ref{derived-section-derived-functors} 节
发展的一般理论。令
$\mathcal{D} = K(\mathcal{O}_Y)$ 及
$\mathcal{D}' = D(\mathcal{O}_X)$。用
$F : \mathcal{D} \to \mathcal{D}'$ 表示由规则
$F(\mathcal{G}^\bullet) = f^*\mathcal{G}^\bullet$ 定义的
三角范畴正合函子。令 $S$ 为
$\mathcal{D} = K(\mathcal{O}_Y)$ 中所有拟同构所成的集合。
这给出《导出范畴》情形
\ref{derived-situation-derived-functor} 中的情形，因而可应用
《导出范畴》定义
\ref{derived-definition-right-derived-functor-defined}。
我们断言 $LF$ 处处有定义。取
$\mathcal{P} \subset \Ob(\mathcal{D})$ 为所有 $K$-平坦复形所成的族，
这由《导出范畴》引理
\ref{derived-lemma-find-existence-computes} 得到：
（1）由引理 \ref{cohomology-lemma-K-flat-resolution} 得到；为验证（2），
须证明对 $\mathcal{O}_Y$-模的 K-平坦复形之间的拟同构
$\mathcal{K}_1^\bullet  \to \mathcal{K}_2^\bullet$，映射
$f^*\mathcal{K}_1^\bullet  \to f^*\mathcal{K}_2^\bullet$
也是拟同构。把该映射写成
$$
f^{-1}\mathcal{K}_1^\bullet \otimes_{f^{-1}\mathcal{O}_Y} \mathcal{O}_X
\longrightarrow
f^{-1}\mathcal{K}_2^\bullet \otimes_{f^{-1}\mathcal{O}_Y} \mathcal{O}_X
$$
。函子 $f^{-1}$ 正合，故映射
$f^{-1}\mathcal{K}_1^\bullet \to
f^{-1}\mathcal{K}_2^\bullet$ 是拟同构。把引理
\ref{cohomology-lemma-pullback-K-flat} 应用于态射
$(X, f^{-1}\mathcal{O}_Y) \to (Y, \mathcal{O}_Y)$，可知复形
$f^{-1}\mathcal{K}_1^\bullet$ 与
$f^{-1}\mathcal{K}_2^\bullet$ 是
$f^{-1}\mathcal{O}_Y$-模的 K-平坦复形。因此引理
\ref{cohomology-lemma-derived-tor-quasi-isomorphism-other-side}
保证所显示的映射是拟同构。于是得到导出函子
$$
LF :
D(\mathcal{O}_Y) = S^{-1}\mathcal{D}
\longrightarrow
\mathcal{D}' = D(\mathcal{O}_X)
$$
；见《导出范畴》公式 (\ref{derived-equation-everywhere})。最后，
《导出范畴》引理 \ref{derived-lemma-find-existence-computes} 还保证
$LF(\mathcal{K}^\bullet) = F(\mathcal{K}^\bullet) = f^*\mathcal{K}^\bullet$
在 $\mathcal{K}^\bullet$ 为 K-平坦时成立；即
$Lf^* = LF$ 确实按上述方式计算。
\end{proof}

\begin{lemma}
\label{cohomology-lemma-derived-pullback-composition}
设 $f : X \to Y$ 与 $g : Y \to Z$ 是环空间的态射。
则作为函子 $D(\mathcal{O}_Z) \to D(\mathcal{O}_X)$，有
$Lf^* \circ Lg^* = L(g \circ f)^*$。
\end{lemma}

\begin{proof}
设 $E$ 是 $D(\mathcal{O}_Z)$ 的对象。按构造，为计算
$Lg^*E$，选取在 $Z$ 上表示 $E$ 的 K-平坦复形
$\mathcal{K}^\bullet$，并令
$Lg^*E = g^*\mathcal{K}^\bullet$。由引理
\ref{cohomology-lemma-pullback-K-flat}，$g^*\mathcal{K}^\bullet$ 在
$Y$ 上 K-平坦。于是 $Lf^*Lg^*E$ 由
$f^*g^*\mathcal{K}^\bullet = (g \circ f)^*\mathcal{K}^\bullet$
给出，而后者也表示 $L(g \circ f)^*E$。
\end{proof}

\begin{lemma}
\label{cohomology-lemma-pullback-tensor-product}
设 $f : (X, \mathcal{O}_X) \to (Y, \mathcal{O}_Y)$
是环空间的态射。存在典范的双函子性同构
$$
Lf^*(
\mathcal{F}^\bullet \otimes_{\mathcal{O}_Y}^{\mathbf{L}} \mathcal{G}^\bullet
) =
Lf^*\mathcal{F}^\bullet 
\otimes_{\mathcal{O}_X}^{\mathbf{L}}
Lf^*\mathcal{G}^\bullet 
$$
，其中
$\mathcal{F}^\bullet, \mathcal{G}^\bullet \in
\Ob(D(\mathcal{O}_Y))$。
\end{lemma}

\begin{proof}
可以假设 $\mathcal{F}^\bullet$ 与 $\mathcal{G}^\bullet$
都是 K-平坦复形。此时
$\mathcal{F}^\bullet \otimes_{\mathcal{O}_Y}^{\mathbf{L}} \mathcal{G}^\bullet$
就是双复形
$\mathcal{F}^\bullet \otimes_{\mathcal{O}_Y}
\mathcal{G}^\bullet$ 所伴随的总复形。由引理
\ref{cohomology-lemma-tensor-product-K-flat}，
$\text{Tot}(\mathcal{F}^\bullet \otimes_{\mathcal{O}_Y} \mathcal{G}^\bullet)$
也是 K-平坦的。因此引理中的同构来自同构
$$
\text{Tot}(f^*\mathcal{F}^\bullet \otimes_{\mathcal{O}_X}
f^*\mathcal{G}^\bullet)
\longrightarrow
f^*\text{Tot}(\mathcal{F}^\bullet \otimes_{\mathcal{O}_Y} \mathcal{G}^\bullet)
$$
，其各分量为《模层》引理
\ref{modules-lemma-tensor-product-pullback} 中的同构
$f^*\mathcal{F}^p \otimes_{\mathcal{O}_X} f^*\mathcal{G}^q \to
 f^*(\mathcal{F}^p \otimes_{\mathcal{O}_Y} \mathcal{G}^q)$。
\end{proof}

\begin{lemma}
\label{cohomology-lemma-variant-derived-pullback}
设 $f : (X, \mathcal{O}_X) \to (Y, \mathcal{O}_Y)$
是环空间的态射。存在典范的双函子性同构
$$
\mathcal{F}^\bullet
\otimes_{\mathcal{O}_X}^{\mathbf{L}}
Lf^*\mathcal{G}^\bullet
=
\mathcal{F}^\bullet 
\otimes_{f^{-1}\mathcal{O}_Y}^{\mathbf{L}}
f^{-1}\mathcal{G}^\bullet 
$$
，其中 $\mathcal{F}^\bullet$ 属于 $D(\mathcal{O}_X)$，
而 $\mathcal{G}^\bullet$ 属于 $D(\mathcal{O}_Y)$。
\end{lemma}

\begin{proof}
设 $\mathcal{F}$ 是 $\mathcal{O}_X$-模，而
$\mathcal{G}$ 是 $\mathcal{O}_Y$-模。则
$\mathcal{F} \otimes_{\mathcal{O}_X} f^*\mathcal{G} =
\mathcal{F} \otimes_{f^{-1}\mathcal{O}_Y} f^{-1}\mathcal{G}$
，这是因为
$f^*\mathcal{G} =
\mathcal{O}_X \otimes_{f^{-1}\mathcal{O}_Y} f^{-1}\mathcal{G}$.
引理由此及各定义得到。
\end{proof}

\begin{lemma}
\label{cohomology-lemma-tensor-pull-compatibility}
设 $f : (X, \mathcal{O}_X) \to (Y, \mathcal{O}_Y)$ 是环空间的态射。
设 $\mathcal{K}^\bullet$ 与 $\mathcal{M}^\bullet$ 是
$\mathcal{O}_Y$-模复形。下图
$$
\xymatrix{
Lf^*(\mathcal{K}^\bullet
\otimes_{\mathcal{O}_Y}^\mathbf{L}
\mathcal{M}^\bullet) \ar[r] \ar[d] &
Lf^*\text{Tot}(\mathcal{K}^\bullet
\otimes_{\mathcal{O}_Y}
\mathcal{M}^\bullet) \ar[d] \\
Lf^*\mathcal{K}^\bullet \otimes_{\mathcal{O}_X}^\mathbf{L}
Lf^*\mathcal{M}^\bullet \ar[d] &
f^*\text{Tot}(\mathcal{K}^\bullet
\otimes_{\mathcal{O}_Y}
\mathcal{M}^\bullet) \ar[d] \\
f^*\mathcal{K}^\bullet \otimes_{\mathcal{O}_X}^\mathbf{L}
f^*\mathcal{M}^\bullet \ar[r] &
\text{Tot}(f^*\mathcal{K}^\bullet \otimes_{\mathcal{O}_X}
f^*\mathcal{M}^\bullet)
}
$$
交换。
\end{lemma}

\begin{proof}
我们将使用引理 \ref{cohomology-lemma-pullback-K-flat} 中所述的 K-平坦消解的存在性。
若选取这样的消解
$\mathcal{P}^\bullet \to \mathcal{K}^\bullet$ 与
$\mathcal{Q}^\bullet \to \mathcal{M}^\bullet$，则可见
$$
\xymatrix{
Lf^*\text{Tot}(\mathcal{P}^\bullet
\otimes_{\mathcal{O}_Y}
\mathcal{Q}^\bullet) \ar[r] \ar[d] &
Lf^*\text{Tot}(\mathcal{K}^\bullet
\otimes_{\mathcal{O}_Y}
\mathcal{M}^\bullet) \ar[d] \\
f^*\text{Tot}(\mathcal{P}^\bullet
\otimes_{\mathcal{O}_Y}
\mathcal{Q}^\bullet) \ar[d] \ar[r] &
f^*\text{Tot}(\mathcal{K}^\bullet
\otimes_{\mathcal{O}_Y}
\mathcal{M}^\bullet) \ar[d] \\
\text{Tot}(f^*\mathcal{P}^\bullet \otimes_{\mathcal{O}_X}
f^*\mathcal{Q}^\bullet) \ar[r] &
\text{Tot}(f^*\mathcal{K}^\bullet \otimes_{\mathcal{O}_X}
f^*\mathcal{M}^\bullet)
}
$$
交换。不过，按我们对 $\mathcal{P}^\bullet$ 与
$\mathcal{Q}^\bullet$ 的选取以及引理
\ref{cohomology-lemma-tensor-product-K-flat}，此图的左边正是原图的左边。
\end{proof}









\section{无界复形的上同调}
\label{cohomology-section-unbounded}

\noindent
设 $(X, \mathcal{O}_X)$ 是环空间。范畴
$\textit{Mod}(\mathcal{O}_X)$ 是 Grothendieck 阿贝尔范畴：
它具有所有余极限，滤过余极限正合，并且具有生成元，即
$$
\bigoplus\nolimits_{U \subset X\text{开}} j_{U!}\mathcal{O}_U,
$$
；见《模层》第 \ref{modules-section-kernels} 节及引理
\ref{modules-lemma-j-shriek-flat}、
\ref{modules-lemma-module-quotient-flat}。由《内射对象》定理
\ref{injectives-theorem-K-injective-embedding-grothendieck}
，对 $\mathcal{O}_X$-模的每个复形 $\mathcal{F}^\bullet$，
都存在到 K-内射 $\mathcal{O}_X$-模复形的单射拟同构
$\mathcal{F}^\bullet \to \mathcal{I}^\bullet$，其中
$\mathcal{I}^\bullet$ 的各项均为内射 $\mathcal{O}_X$-模；
此外，可使此嵌入关于复形 $\mathcal{F}^\bullet$ 具有函子性。
由《导出范畴》引理
\ref{derived-lemma-enough-K-injectives-implies} 可得：
\begin{enumerate}
\item 任意到三角范畴 $\mathcal{D}$ 的正合函子
$F : K(\textit{Mod}(\mathcal{O}_X)) \to \mathcal{D}$
都具有右导出函子 $RF : D(\mathcal{O}_X) \to \mathcal{D}$；
\item 对任意到阿贝尔范畴 $\mathcal{A}$ 的加性函子
$F : \textit{Mod}(\mathcal{O}_X) \to \mathcal{A}$
，考虑由 $F$ 诱导的正合函子
$F : K(\textit{Mod}(\mathcal{O}_X)) \to K(\mathcal{A})$，
可得到右导出函子
$RF : D(\mathcal{O}_X) \to D(\mathcal{A})$。
\end{enumerate}
按构造有
$RF(\mathcal{F}^\bullet) = F(\mathcal{I}^\bullet)$，其中
$\mathcal{F}^\bullet \to \mathcal{I}^\bullet$ 如上。

\medskip\noindent
下面给出上述构造的一些例子：
\begin{enumerate}
\item 函子 $\Gamma(X, -) : \textit{Mod}(\mathcal{O}_X) \to
\text{Mod}_{\Gamma(X, \mathcal{O}_X)}$ 给出
$$
R\Gamma(X, -) : D(\mathcal{O}_X) \to D(\Gamma(X, \mathcal{O}_X))
$$
我们将用记号 $H^i(X, K) = H^i(R\Gamma(X, K))$ 表示上同调。
\item 对开集 $U \subset X$，考虑函子
$\Gamma(U, -) : \textit{Mod}(\mathcal{O}_X) \to
\text{Mod}_{\Gamma(U, \mathcal{O}_X)}$。它给出
$$
R\Gamma(U, -) : D(\mathcal{O}_X) \to D(\Gamma(U, \mathcal{O}_X))
$$
我们将用记号 $H^i(U, K) = H^i(R\Gamma(U, K))$ 表示上同调。
\item 对环空间态射
$f : (X, \mathcal{O}_X) \to (Y, \mathcal{O}_Y)$，考虑函子
$f_* : \textit{Mod}(\mathcal{O}_X) \to \textit{Mod}(\mathcal{O}_Y)$
，它在无界导出范畴上给出总直像
$$
Rf_* : D(\mathcal{O}_X) \longrightarrow D(\mathcal{O}_Y)
$$
\end{enumerate}

\begin{lemma}
\label{cohomology-lemma-adjoint}
设 $f : (X, \mathcal{O}_X) \to (Y, \mathcal{O}_Y)$ 是环空间的态射。
上面定义的函子 $Rf_*$ 与引理
\ref{cohomology-lemma-derived-base-change} 中定义的函子 $Lf^*$ 互为伴随：
$$
\Hom_{D(\mathcal{O}_X)}(Lf^*\mathcal{G}^\bullet, \mathcal{F}^\bullet)
=
\Hom_{D(\mathcal{O}_Y)}(\mathcal{G}^\bullet, Rf_*\mathcal{F}^\bullet)
$$
，它关于
$\mathcal{F}^\bullet \in \Ob(D(\mathcal{O}_X))$ 与
$\mathcal{G}^\bullet \in \Ob(D(\mathcal{O}_Y))$ 具有双函子性。
\end{lemma}

\begin{proof}
这形式地由 $Rf_*$ 与 $Lf^*$ 的存在性得到，见《导出范畴》引理
\ref{derived-lemma-derived-adjoint-functors}。
\end{proof}

\begin{lemma}
\label{cohomology-lemma-derived-pushforward-composition}
设 $f : X \to Y$ 与 $g : Y \to Z$ 是环空间的态射。
则作为函子 $D(\mathcal{O}_X) \to D(\mathcal{O}_Z)$，有
$Rg_* \circ Rf_* = R(g \circ f)_*$。
\end{lemma}

\begin{proof}
由引理 \ref{cohomology-lemma-adjoint}，$Rg_* \circ Rf_*$ 伴随于
$Lf^* \circ Lg^*$。由引理
\ref{cohomology-lemma-derived-pullback-composition} 有
$Lf^* \circ Lg^* = L(g \circ f)^*$，故由伴随函子的唯一性，
$Rg_* \circ Rf_* = R(g \circ f)_*$。
\end{proof}

\begin{remark}
\label{cohomology-remark-base-change}
无界导出函子 $Lf^*$ 与 $Rf_*$ 的构造使我们能在完全一般的情形下
构造基变换映射。具体地，假设
$$
\xymatrix{
X' \ar[r]_{g'} \ar[d]_{f'} &
X \ar[d]^f \\
S' \ar[r]^g &
S
}
$$
是环空间的交换图。设 $K$ 是 $D(\mathcal{O}_X)$ 的对象。
则在 $D(\mathcal{O}_{S'})$ 中存在典范基变换映射
$$
Lg^*Rf_*K \longrightarrow R(f')_*L(g')^*K
$$
。具体地，此映射伴随于映射
$L(f')^*Lg^*Rf_*K \to L(g')^*K$。
由于 $L(f')^*Lg^* = L(g')^*Lf^*$，这等同于映射
$L(g')^*Lf^*Rf_*K \to L(g')^*K$
，可取它为伴随映射 $Lf^*Rf_*K \to K$ 在 $L(g')^*$ 下的像。
\end{remark}

\begin{remark}
\label{cohomology-remark-compose-base-change}
考虑环空间的交换图
$$
\xymatrix{
X' \ar[r]_k \ar[d]_{f'} & X \ar[d]^f \\
Y' \ar[r]^l \ar[d]_{g'} & Y \ar[d]^g \\
Z' \ar[r]^m & Z
}
$$
。则注记 \ref{cohomology-remark-base-change} 中对应于两个方块的基变换映射
复合后给出对应于外部矩形的基变换映射。更确切地，复合
\begin{align*}
Lm^* \circ R(g \circ f)_*
& =
Lm^* \circ Rg_* \circ Rf_* \\
& \to Rg'_* \circ Ll^* \circ Rf_* \\
& \to Rg'_* \circ Rf'_* \circ Lk^* \\
& = R(g' \circ f')_* \circ Lk^*
\end{align*}
是该矩形的基变换映射。我们略去验证。
\end{remark}

\begin{remark}
\label{cohomology-remark-compose-base-change-horizontal}
考虑环空间的交换图
$$
\xymatrix{
X'' \ar[r]_{g'} \ar[d]_{f''} & X' \ar[r]_g \ar[d]_{f'} & X \ar[d]^f \\
Y'' \ar[r]^{h'} & Y' \ar[r]^h & Y
}
$$
。则注记 \ref{cohomology-remark-base-change} 中对应于两个方块的基变换映射
复合后给出对应于外部矩形的基变换映射。更确切地，复合
\begin{align*}
L(h \circ h')^* \circ Rf_*
& =
L(h')^* \circ Lh_* \circ Rf_* \\
& \to L(h')^* \circ Rf'_* \circ Lg^* \\
& \to Rf''_* \circ L(g')^* \circ Lg^* \\
& = Rf''_* \circ L(g \circ g')^*
\end{align*}
是该矩形的基变换映射。我们略去验证。
\end{remark}

\begin{lemma}
\label{cohomology-lemma-adjoints-push-pull-compatibility}
设 $f : (X, \mathcal{O}_X) \to (Y, \mathcal{O}_Y)$ 是环空间的态射。
设 $\mathcal{K}^\bullet$ 是 $\mathcal{O}_X$-模复形。图
$$
\xymatrix{
Lf^*f_*\mathcal{K}^\bullet \ar[r] \ar[d] &
f^*f_*\mathcal{K}^\bullet \ar[d] \\
Lf^*Rf_*\mathcal{K}^\bullet \ar[r] &
\mathcal{K}^\bullet
}
$$
来自复形上的 $Lf^* \to f^*$、$f_* \to Rf_*$ 以及伴随
$Lf^* \circ Rf_* \to \text{id}$，并在 $D(\mathcal{O}_X)$ 中交换。
\end{lemma}

\begin{proof}
我们将使用 K-平坦消解与 K-内射消解的存在性，见引理
\ref{cohomology-lemma-pullback-K-flat} 及上面的讨论。选取拟同构
$\mathcal{K}^\bullet \to \mathcal{I}^\bullet$，其中
$\mathcal{I}^\bullet$ 作为 $\mathcal{O}_X$-模复形是 K-内射的。
选取拟同构
$\mathcal{Q}^\bullet \to f_*\mathcal{I}^\bullet$，其中
$\mathcal{Q}^\bullet$ 作为 $\mathcal{O}_Y$-模复形是 K-平坦的。
可以选取 $\mathcal{O}_Y$-模的 K-平坦复形
$\mathcal{P}^\bullet$ 以及复形态射图
$$
\xymatrix{
\mathcal{P}^\bullet \ar[r] \ar[d] &
f_*\mathcal{K}^\bullet \ar[d] \\
\mathcal{Q}^\bullet \ar[r] & f_*\mathcal{I}^\bullet
}
$$
，它在同伦意义下交换，且上方水平箭头为拟同构。事实上，
由于拟同构在复形的同伦范畴中构成乘法系统，可以先对某个复形
$\mathcal{P}^\bullet$ 选取这样的图，然后用 K-平坦复形替换
$\mathcal{P}^\bullet$。取拉回后得到复形态射图
$$
\xymatrix{
f^*\mathcal{P}^\bullet \ar[r] \ar[d] &
f^*f_*\mathcal{K}^\bullet \ar[d] \ar[r] &
\mathcal{K}^\bullet \ar[d] \\
f^*\mathcal{Q}^\bullet \ar[r] &
f^*f_*\mathcal{I}^\bullet \ar[r] &
\mathcal{I}^\bullet
}
$$
，它在同伦意义下交换。外部矩形证明了引理中的断言。
\end{proof}

\begin{remark}
\label{cohomology-remark-cup-product}
设 $f : (X, \mathcal{O}_X) \to (Y, \mathcal{O}_Y)$ 是环空间的态射。
$Lf^*$ 与 $Rf_*$ 的伴随性使我们能构造相对杯积
$$
Rf_*K \otimes_{\mathcal{O}_Y}^\mathbf{L} Rf_*L
\longrightarrow
Rf_*(K \otimes_{\mathcal{O}_X}^\mathbf{L} L)
$$
；它对所有 $D(\mathcal{O}_X)$ 中的 $K, L$ 都是
$D(\mathcal{O}_Y)$ 中的映射。具体地，此映射伴随于映射
$Lf^*(Rf_*K \otimes_{\mathcal{O}_Y}^\mathbf{L} Rf_*L) \to
 K \otimes_{\mathcal{O}_X}^\mathbf{L} L$；可以把它取为同构
$Lf^*(Rf_*K \otimes_{\mathcal{O}_Y}^\mathbf{L} Rf_*L) =
Lf^*Rf_*K \otimes_{\mathcal{O}_X}^\mathbf{L} Lf^*Rf_*L$
（引理 \ref{cohomology-lemma-pullback-tensor-product}）与映射
$Lf^*Rf_*K \otimes_{\mathcal{O}_X}^\mathbf{L} Lf^*Rf_*L
\to K \otimes_{\mathcal{O}_X}^\mathbf{L} L$
的复合，其中后一个映射来自余单位
$Lf^* \circ Rf_* \to \text{id}$。
\end{remark}






\section{滤过复形的上同调}
\label{cohomology-section-cohomology-filtered-object}

\noindent
研究代数簇的上同调时，层的滤过复形经常自然出现；例如 de Rham
复形带有 Hodge 滤过。本节使用非常一般的《内射对象》引理
\ref{injectives-lemma-K-injective-embedding-filtration}
构造上同调上的谱序列，并把它们与先前构造的谱序列联系起来。

\begin{lemma}
\label{cohomology-lemma-spectral-sequence-filtered-object}
设 $(X, \mathcal{O}_X)$ 是环空间，而
$\mathcal{F}^\bullet$ 是 $\mathcal{O}_X$-模的滤过复形。
存在双分次 $\Gamma(X, \mathcal{O}_X)$-模的典范谱序列
$(E_r, \text{d}_r)_{r \geq 1}$，其中 $d_r$ 的双次数为
$(r, -r + 1)$，且
$$
E_1^{p, q} = H^{p + q}(X, \text{gr}^p\mathcal{F}^\bullet)
$$
若对每个 $n$ 都有
$$
H^n(X, F^p\mathcal{F}^\bullet) = 0\text{，当 }p \gg 0
\quad\text{且}\quad
H^n(X, F^p\mathcal{F}^\bullet) = H^n(X, \mathcal{F}^\bullet)\text{，当 }p \ll 0
$$
，则该谱序列有界并收敛到
$H^*(X, \mathcal{F}^\bullet)$。
\end{lemma}

\begin{proof}
（当该复形为具有有限滤过的模的有下界复形时，证明见下面的注记。）
选取《内射对象》引理
\ref{injectives-lemma-K-injective-embedding-filtration}
中的滤过复形映射
$j : \mathcal{F}^\bullet \to \mathcal{J}^\bullet$。
该谱序列是《同调代数》第
\ref{homology-section-filtered-complex} 节中与滤过复形
$$
\Gamma(X, \mathcal{J}^\bullet)
\quad\text{其中}\quad
F^p\Gamma(X, \mathcal{J}^\bullet) = \Gamma(X, F^p\mathcal{J}^\bullet)
$$
相伴的谱序列。由于上同调通过在 K-内射代表元上取值来计算，
可知 $E_1$ 页如引理所述。在所给条件下，收敛性与有界性由
《同调代数》引理 \ref{homology-lemma-ss-converges-trivial} 得到。
\end{proof}

\begin{remark}
\label{cohomology-remark-spectral-sequence-filtered-object}
设 $(X, \mathcal{O}_X)$ 是环空间，而
$\mathcal{F}^\bullet$ 是 $\mathcal{O}_X$-模的滤过复形。
若 $\mathcal{F}^\bullet$ 有下界，且对每个 $n$，
$\mathcal{F}^n$ 上的滤过有限，则可构造引理
\ref{cohomology-lemma-spectral-sequence-filtered-object} 中的谱序列，而无需使用《内射对象》引理
\ref{injectives-lemma-K-injective-embedding-filtration}。
事实上，由《导出范畴》引理
\ref{derived-lemma-right-resolution-by-filtered-injectives}，
存在滤过复形的滤过拟同构
$i : \mathcal{F}^\bullet \to \mathcal{I}^\bullet$，
其中 $\mathcal{I}^\bullet$ 有下界，对所有 $n$，
$\mathcal{I}^n$ 上的滤过有限，并且每个
$\text{gr}^p\mathcal{I}^n$ 都是内射 $\mathcal{O}_X$-模。
然后取与
$$
\Gamma(X, \mathcal{I}^\bullet)
\quad\text{其中}\quad
F^p\Gamma(X, \mathcal{I}^\bullet) = \Gamma(X, F^p\mathcal{I}^\bullet)
$$
相伴的谱序列。由于上同调可通过在内射对象的有下界复形上取值来计算，
可知 $E_1$ 页如引理所述。在所给条件下，收敛性与有界性由
《同调代数》引理 \ref{homology-lemma-biregular-ss-converges} 得到。
事实上，这是《导出范畴》引理
\ref{derived-lemma-ss-filtered-derived} 中谱序列的特例。
\end{remark}

\begin{example}
\label{cohomology-example-spectral-sequence}
设 $(X, \mathcal{O}_X)$ 是环空间，而
$\mathcal{F}^\bullet$ 是 $\mathcal{O}_X$-模复形。令
$F^p\mathcal{F}^\bullet =
\tau_{\leq -p}\mathcal{F}^\bullet$，即可应用引理
\ref{cohomology-lemma-spectral-sequence-filtered-object}。
（若 $\mathcal{F}^\bullet$ 有下界，可使用注记
\ref{cohomology-remark-spectral-sequence-filtered-object}。）于是得到谱序列
$$
E_1^{p, q} = H^{p + q}(X, H^{-p}(\mathcal{F}^\bullet)[p]) =
H^{2p + q}(X, H^{-p}(\mathcal{F}^\bullet))
$$
作重新编号 $p = -j$、$q = i + 2j$ 后可知，对任意
$K \in D(\mathcal{O}_X)$，存在双分次模的谱序列
$(E'_r, d'_r)_{r \geq 2}$，其中 $d'_r$ 的双次数为
$(r, -r + 1)$，并且
$$
(E'_2)^{i, j} = H^i(X, H^j(K))
$$
例如，若 $K$ 有下界，则此谱序列有界，并收敛到
$H^{i + j}(X, K)$。在有下界的情形下，此谱序列是《导出范畴》引理
\ref{derived-lemma-two-ss-complex-functor} 的第二谱序列
（用 Cartan--Eilenberg 消解构造）的一个例子。
\end{example}

\begin{example}
\label{cohomology-example-spectral-sequence-bis}
设 $(X, \mathcal{O}_X)$ 是环空间，而
$\mathcal{F}^\bullet$ 是 $\mathcal{O}_X$-模复形。令
$F^p\mathcal{F}^\bullet =
\sigma_{\geq p}\mathcal{F}^\bullet$，即可应用引理
\ref{cohomology-lemma-spectral-sequence-filtered-object}。于是得到谱序列
$$
E_1^{p, q} = H^{p + q}(X, \mathcal{F}^p[-p]) = H^q(X, \mathcal{F}^p)
$$
若 $\mathcal{F}^\bullet$ 有下界，则
\begin{enumerate}
\item 可以用注记 \ref{cohomology-remark-spectral-sequence-filtered-object}
构造此谱序列；
\item 此谱序列有界，并收敛到
$H^{i + j}(X, \mathcal{F}^\bullet)$；
\item 此谱序列等于《导出范畴》引理
\ref{derived-lemma-two-ss-complex-functor} 的第一谱序列
（用 Cartan--Eilenberg 消解构造）。
\end{enumerate}
\end{example}

\begin{lemma}
\label{cohomology-lemma-relative-spectral-sequence-filtered-object}
设 $f : (X, \mathcal{O}_X) \to (Y, \mathcal{O}_Y)$ 是环空间的态射。
设 $\mathcal{F}^\bullet$ 是 $\mathcal{O}_X$-模的滤过复形。
存在双分次 $\mathcal{O}_Y$-模的典范谱序列
$(E_r, \text{d}_r)_{r \geq 1}$，其中 $d_r$ 的双次数为
$(r, -r + 1)$，且
$$
E_1^{p, q} = R^{p + q}f_*\text{gr}^p\mathcal{F}^\bullet
$$
若对每个 $n$ 都有
$$
R^nf_*F^p\mathcal{F}^\bullet = 0 \text{，当 }p \gg 0
\quad\text{且}\quad
R^nf_*F^p\mathcal{F}^\bullet = R^nf_*\mathcal{F}^\bullet \text{，当 }p \ll 0
$$
，则该谱序列有界并收敛到
$Rf_*\mathcal{F}^\bullet$。
\end{lemma}

\begin{proof}
证明与引理 \ref{cohomology-lemma-spectral-sequence-filtered-object} 的证明完全相同。
\end{proof}





\section{Godement 消解}
\label{cohomology-section-godement}

\noindent
参考文献为 \cite{Godement}。

\medskip\noindent
设 $(X, \mathcal{O}_X)$ 是环空间。以 $X_{disc}$ 表示与 $X$
具有相同点集的离散拓扑空间，以 $f : X_{disc} \to X$
表示显然的连续映射。令
$\mathcal{O}_{X_{disc}} = f^{-1}\mathcal{O}_X$。
则 $f : (X_{disc}, \mathcal{O}_{X_{disc}}) \to (X, \mathcal{O}_X)$
是环空间的平坦态射。可将《单纯方法》第
\ref{simplicial-section-standard} 节中材料的
{\it 对偶}应用于模层上的伴随函子对 $f^*, f_*$。由此得到增广余单纯对象
$$
\xymatrix{
\text{id} \ar[r] &
f_*f^* \ar@<1ex>[r] \ar@<-1ex>[r] &
f_*f^*f_*f^* \ar@<0ex>[l] \ar@<-2ex>[r] \ar@<0ex>[r] \ar@<2ex>[r] &
f_*f^*f_*f^*f_*f^* \ar@<1ex>[l] \ar@<-1ex>[l]
}
$$
，它位于从 $\textit{Mod}(\mathcal{O}_X)$ 到自身的函子范畴中；
参见《单纯方法》引理 \ref{simplicial-lemma-standard-simplicial}。
此外，增广
$$
\xymatrix{
f^* \ar[r] &
f^*f_*f^* \ar@<1ex>[r] \ar@<-1ex>[r] &
f^*f_*f^*f_*f^* \ar@<0ex>[l] \ar@<-2ex>[r] \ar@<0ex>[r] \ar@<2ex>[r] &
f^*f_*f^*f_*f^*f_*f^* \ar@<1ex>[l] \ar@<-1ex>[l]
}
$$
是同伦等价；参见《单纯方法》引理
\ref{simplicial-lemma-standard-simplicial-homotopy}。

\begin{lemma}
\label{cohomology-lemma-godement-resolution}
设 $(X, \mathcal{O}_X)$ 是环空间。对每个
$\mathcal{O}_X$-模层 $\mathcal{F}$，存在消解
$$
0 \to
\mathcal{F} \to
f_*f^*\mathcal{F} \to
f_*f^*f_*f^*\mathcal{F} \to
f_*f^*f_*f^*f_*f^*\mathcal{F} \to \ldots
$$
，它关于 $\mathcal{F}$ 具有函子性，并且每一项
$f_*f^* \ldots f_*f^*\mathcal{F}$ 都是松弛
$\mathcal{O}_X$-模；此外，对所有 $x \in X$，映射
$$
\mathcal{F}_x[0] \to \Big(
(f_*f^*\mathcal{F})_x \to
(f_*f^*f_*f^*\mathcal{F})_x \to
(f_*f^*f_*f^*f_*f^*\mathcal{F})_x \to \ldots
\Big)
$$
是 $\mathcal{O}_{X, x}$-模复形范畴中的链同伦等价。
\end{lemma}

\begin{proof}
复形 $f_*f^*\mathcal{F} \to  f_*f^*f_*f^*\mathcal{F} \to
f_*f^*f_*f^*f_*f^*\mathcal{F} \to \ldots$ 是与上述余单纯对象相伴的复形，
其各项为
$f_*f^*\mathcal{F}, f_*f^*f_*f^*\mathcal{F},
f_*f^*f_*f^*f_*f^*\mathcal{F}, \ldots$；参见《单纯方法》第
\ref{simplicial-section-dold-kan-cosimplicial} 节。
增广给出所示的映射 $\mathcal{F} \to f_*f^*\mathcal{F}$。
对 $X_{disc}$ 上任意阿贝尔层 $\mathcal{H}$，推前
$f_*\mathcal{H}$ 都是松弛的，因为 $X_{disc}$ 是离散空间，
而松弛层的推前仍是松弛层。因此该复形的各项都是松弛
$\mathcal{O}_X$-模。

\medskip\noindent
若 $x \in X_{disc} = X$ 是一点，则对任意
$\mathcal{O}_X$-模 $\mathcal{G}$ 都有
$(f^*\mathcal{G})_x = \mathcal{G}_x$。因此 $f^*$ 是正合函子，并且
$\mathcal{O}_X$-模复形
$\mathcal{G}_1 \to \mathcal{G}_2 \to \mathcal{G}_3$
正合，当且仅当
$f^*\mathcal{G}_1 \to f^*\mathcal{G}_2 \to f^*\mathcal{G}_3$
正合（参见《模》引理 \ref{modules-lemma-abelian}）。
本节引言中提到的结果表明，应用 $f^*$ 后，常值余单纯对象
$f^*\mathcal{F}$ 到各项为
$f_*f^*\mathcal{F}, f_*f^*f_*f^*\mathcal{F},
f_*f^*f_*f^*f_*f^*\mathcal{F}, \ldots$
的余单纯对象之间得到一个同伦等价。由《单纯方法》引理
\ref{simplicial-lemma-homotopy-equivalence-s-Q} 可知
$$
f^*\mathcal{F}[0] \to \Big(
f^*f_*f^*\mathcal{F} \to
f^*f_*f^*f_*f^*\mathcal{F} \to
f^*f_*f^*f_*f^*f_*f^*\mathcal{F} \to \ldots
\Big)
$$
是链同伦等价。这立即推出引理其余两个断言。
\end{proof}

\begin{lemma}
\label{cohomology-lemma-godement-resolution-bounded-below}
设 $(X, \mathcal{O}_X)$ 是环空间，而
$\mathcal{F}^\bullet$ 是 $\mathcal{O}_X$-模的有下界复形。
存在拟同构
$\mathcal{F}^\bullet \to \mathcal{G}^\bullet$
，其中 $\mathcal{G}^\bullet$ 是松弛 $\mathcal{O}_X$-模的有下界复形；
并且对所有 $x \in X$，映射
$\mathcal{F}^\bullet_x \to \mathcal{G}^\bullet_x$
是 $\mathcal{O}_{X, x}$-模复形范畴中的链同伦等价。
\end{lemma}

\begin{proof}
令 $\mathcal{A}$ 为 $\mathcal{O}_X$-模复形的范畴，
并令 $\mathcal{B}$ 为 $\mathcal{O}_X$-模复形的范畴。
于是可将上述讨论应用于 $\mathcal{A}$ 与 $\mathcal{B}$
之间的伴随函子 $f^*$ 和 $f_*$。完全仿照引理
\ref{cohomology-lemma-godement-resolution} 的证明，可得消解
$$
0 \to
\mathcal{F}^\bullet \to
f_*f^*\mathcal{F}^\bullet \to
f_*f^*f_*f^*\mathcal{F}^\bullet \to
f_*f^*f_*f^*f_*f^*\mathcal{F}^\bullet \to \ldots
$$
，它位于阿贝尔范畴 $\mathcal{A}$ 中，并且每个
$f_*f^*\ldots f_*f^*\mathcal{F}^\bullet$ 的每一项都是松弛
$\mathcal{O}_X$-模；此外，对所有 $x \in X$，映射
$$
\mathcal{F}^\bullet_x[0] \to \Big(
(f_*f^*\mathcal{F}^\bullet)_x \to
(f_*f^*f_*f^*\mathcal{F}^\bullet)_x \to
(f_*f^*f_*f^*f_*f^*\mathcal{F}^\bullet)_x \to \ldots
\Big)
$$
是 $\mathcal{O}_{X, x}$-模的复形之复形范畴中的链同伦等价。
由于复形的复形就是双复形，可以考虑所诱导的映射
$$
\mathcal{F}^\bullet \to
\mathcal{G}^\bullet =
\text{Tot}(
f_*f^*\mathcal{F}^\bullet \to
f_*f^*f_*f^*\mathcal{F}^\bullet \to
f_*f^*f_*f^*f_*f^*\mathcal{F}^\bullet \to \ldots
)
$$
。由于复形 $\mathcal{F}^\bullet$ 有下界，
$\mathcal{G}^\bullet$ 也有下界。事实上，
$\mathcal{G}^\bullet$ 的每一项都是复形
$f_*f^*\ldots f_*f^*\mathcal{F}^\bullet$ 各项的有限直和，
因而是松弛的。引理的最后一个断言现由《同调代数》引理
\ref{homology-lemma-homotopy-complex-complexes} 推出。
这尤其表明
$\mathcal{F}^\bullet \to \mathcal{G}^\bullet$
是拟同构，故证明完成。
\end{proof}














\section{杯积}
\label{cohomology-section-cup-product}

\noindent
设 $(X, \mathcal{O}_X)$ 是环空间，而 $K, M$ 是
$D(\mathcal{O}_X)$ 的对象。令 $A = \Gamma(X, \mathcal{O}_X)$。
在此情形中，（整体）杯积是 $D(A)$ 中的映射
$$
\mu :
R\Gamma(X, K) \otimes_A^\mathbf{L} R\Gamma(X, M)
\longrightarrow
R\Gamma(X, K \otimes_{\mathcal{O}_X}^\mathbf{L} M)
$$
。通过 $D(pt, A) = D(A)$，按照注记
\ref{cohomology-remark-cup-product}，把它定义为环空间态射
$f : (X, \mathcal{O}_X) \to (pt, A)$ 的相对杯积。
此映射特别给出配对
$$
\cup :
H^i(X, K) \times H^j(X, M)
\longrightarrow
H^{i + j}(X, K \otimes_{\mathcal{O}_X}^\mathbf{L} M)
$$
具体地，给定 $\xi \in H^i(X, K) = H^i(R\Gamma(X, K))$ 与
$\eta \in H^j(X, M) = H^j(R\Gamma(X, M))$，可以先将二者“张量”，
得到
$H^{i + j}(R\Gamma(X, K) \otimes_A^\mathbf{L} R\Gamma(X, M))$
中的元素 $\xi \otimes \eta$；参见《更多代数》第
\ref{more-algebra-section-products-tor} 节。然后应用 $\mu$，
得到 $H^{i + j}(X, K \otimes_{\mathcal{O}_X}^\mathbf{L} M)$
中的所需元素
$\xi \cup \eta = \mu(\xi \otimes \eta)$。

\medskip\noindent
还可以用另一种方式理解 $\xi$ 与 $\eta$ 的杯积。具体地，可以写成
$$
H^i(R\Gamma(X, K)) = \Hom_{D(\mathcal{O}_X)}(\mathcal{O}_X[-i], K)
\quad\text{且}\quad
H^j(R\Gamma(X, M)) = \Hom_{D(\mathcal{O}_X)}(\mathcal{O}_X[-j], M)
$$
，因为 $\Hom(\mathcal{O}_X, -) = \Gamma(X, -)$。
因此，$\xi$ 与 $\eta$ 分别与映射
$$
\tilde \xi : \mathcal{O}_X[-i] \to K
\quad\text{及}\quad
\tilde \eta : \mathcal{O}_X[-j] \to M
$$
是“同一”回事。结合导出张量积的函子性，可得
$$
\mathcal{O}_X[-i - j] =
\mathcal{O}_X[-i] \otimes_{\mathcal{O}_X}^\mathbf{L} \mathcal{O}_X[-j]
\xrightarrow{\tilde \xi \otimes \tilde \eta}
K \otimes_{\mathcal{O}_X}^\mathbf{L} M
$$
；由与上面相同的理由，这给出
$H^{i + j}(X, K \otimes_{\mathcal{O}_X}^\mathbf{L} M)$ 的一个元素。

\begin{lemma}
\label{cohomology-lemma-second-cup-equals-first}
此构造给出杯积。
\end{lemma}

\begin{proof}
取上述 $f : (X, \mathcal{O}_X) \to (pt, A)$，则
$Rf_*(-) = R\Gamma(X, -)$，且映射 $\mu$ 与下列映射伴随：
$$
Lf^*(Rf_*K \otimes_A^\mathbf{L} Rf_*M) =
Lf^*Rf_*K \otimes_{\mathcal{O}_X}^\mathbf{L} Lf^*Rf_*M
\xrightarrow{\epsilon_K \otimes \epsilon_M}
K \otimes_{\mathcal{O}_X}^\mathbf{L} M
$$
，其中 $\epsilon$ 是 $Lf^*$ 与 $Rf_*$ 之间伴随的余单位。
若把 $\xi$ 与 $\eta$ 看作映射
$\xi : A[-i] \to R\Gamma(X, K)$ 和
$\eta : A[-j] \to R\Gamma(X, M)$，则张量
$\xi \otimes \eta$ 对应于映射\footnote{这里隐藏着一个符号。
具体地，该等式由复合
$$
A[-i - j] \to (A \otimes_A^\mathbf{L} A)[-i - j] \to
A[-i] \otimes_A^\mathbf{L} A[-j]
$$
定义；第二步使用《更多代数》第
(\ref{more-algebra-item-shift-tensor}) 项的等同，该等同原则上带有符号。
不过依照我们的约定，此处符号为 $+1$；即使它不是 $+1$ 也无妨，
因为在等同
$\mathcal{O}_X[-i - j] =
\mathcal{O}_X[-i] \otimes_{\mathcal{O}_X}^\mathbf{L} \mathcal{O}_X[-j]$
中使用了相同的符号。}
$$
A[-i - j] = A[-i] \otimes_A^\mathbf{L} A[-j]
\xrightarrow{\xi \otimes \eta}
R\Gamma(X, K) \otimes_A^\mathbf{L} R\Gamma(X, M)
$$
按照定义，杯积 $\xi \cup \eta$ 是映射
$A[-i - j] \to R\Gamma(X, K \otimes_{\mathcal{O}_X}^\mathbf{L} M)$
，它与下列映射伴随：
$$
(\epsilon_K \otimes \epsilon_M) \circ Lf^*(\xi \otimes \eta) =
(\epsilon_K \circ Lf^*\xi) \otimes (\epsilon_M \circ Lf^*\eta)
$$
然而，容易看出
$\epsilon_K \circ Lf^*\xi = \tilde \xi$，且
$\epsilon_M \circ Lf^*\eta = \tilde \eta$.
由此得到
$\widetilde{\xi \cup \eta} = \tilde \xi \otimes \tilde \eta$，
这正是所需的一致性。
\end{proof}

\begin{remark}
\label{cohomology-remark-cup-with-element-map-total-cohomology}
设 $(X, \mathcal{O}_X)$ 是环空间，而 $K, M$ 是
$D(\mathcal{O}_X)$ 的对象。令 $A = \Gamma(X, \mathcal{O}_X)$。
给定 $\xi \in H^i(X, K)$，可得相伴映射
$$
\xi = \text{“}\xi \cup -\text{”} :
R\Gamma(X, M)[-i]
\to
R\Gamma(X, K \otimes_{\mathcal{O}_X}^\mathbf{L} M)
$$
：先如引理 \ref{cohomology-lemma-second-cup-equals-first} 的证明中那样，
把 $\xi$ 表示为映射 $\xi : A[-i] \to R\Gamma(X, K)$，
再使用复合
$$
R\Gamma(X, M)[-i] = A[-i] \otimes_A^\mathbf{L} R\Gamma(X, M)
\xrightarrow{\xi \otimes 1}
R\Gamma(X, K) \otimes_A^\mathbf{L} R\Gamma(X, M)
\to
R\Gamma(X, K \otimes_{\mathcal{O}_X}^\mathbf{L} M)
$$
，其中第二个箭头是上述整体杯积 $\mu$。
由引理 \ref{cohomology-lemma-second-cup-equals-first} 及其证明可知，
在上同调上，这恢复了与 $\xi$ 作杯积。
\end{remark}

\noindent
下面陈述并证明相对杯积的一种自然相容性。设有环空间态射
$f : (X, \mathcal{O}_X) \to (Y, \mathcal{O}_Y)$，
并设 $\mathcal{K}^\bullet$ 与 $\mathcal{M}^\bullet$
为 $\mathcal{O}_X$-模复形。存在朴素杯积
$$
\text{Tot}(
f_*\mathcal{K}^\bullet
\otimes_{\mathcal{O}_Y}
f_*\mathcal{M}^\bullet)
\longrightarrow
f_*\text{Tot}(\mathcal{K}^\bullet
\otimes_{\mathcal{O}_X}
\mathcal{M}^\bullet)
$$
我们断言它与相对杯积相关。

\begin{lemma}
\label{cohomology-lemma-cup-compatible-with-naive}
在上述情形中，下图交换：
$$
\xymatrix{
f_*\mathcal{K}^\bullet
\otimes_{\mathcal{O}_Y}^\mathbf{L}
f_*\mathcal{M}^\bullet \ar[r] \ar[d]
&
Rf_*\mathcal{K}^\bullet
\otimes_{\mathcal{O}_Y}^\mathbf{L}
Rf_*\mathcal{M}^\bullet \ar[d]^{\text{注记 \ref{cohomology-remark-cup-product}}} \\
\text{Tot}(
f_*\mathcal{K}^\bullet
\otimes_{\mathcal{O}_Y}
f_*\mathcal{M}^\bullet) \ar[d]_{\text{朴素杯积}} &
Rf_*(\mathcal{K}^\bullet
\otimes_{\mathcal{O}_X}^\mathbf{L}
\mathcal{M}^\bullet) \ar[d] \\
f_*\text{Tot}(\mathcal{K}^\bullet
\otimes_{\mathcal{O}_X}
\mathcal{M}^\bullet) \ar[r] &
Rf_*\text{Tot}(\mathcal{K}^\bullet
\otimes_{\mathcal{O}_X}
\mathcal{M}^\bullet)
}
$$
\end{lemma}

\begin{proof}
由注记 \ref{cohomology-remark-cup-product} 中的构造可知，沿图顺时针绕行所得映射
$$
f_*\mathcal{K}^\bullet
\otimes_{\mathcal{O}_Y}^\mathbf{L}
f_*\mathcal{M}^\bullet 
\longrightarrow
Rf_*\text{Tot}(\mathcal{K}^\bullet
\otimes_{\mathcal{O}_X}
\mathcal{M}^\bullet)
$$
与下列映射伴随：
\begin{align*}
Lf^*(f_*\mathcal{K}^\bullet
\otimes_{\mathcal{O}_Y}^\mathbf{L}
f_*\mathcal{M}^\bullet)
& =
Lf^*f_*\mathcal{K}^\bullet
\otimes_{\mathcal{O}_Y}^\mathbf{L}
Lf^*f_*\mathcal{M}^\bullet \\
& \to
Lf^*Rf_*\mathcal{K}^\bullet
\otimes_{\mathcal{O}_Y}^\mathbf{L}
Lf^*Rf_*\mathcal{M}^\bullet \\
& \to
\mathcal{K}^\bullet
\otimes_{\mathcal{O}_Y}^\mathbf{L}
\mathcal{M}^\bullet \\
& \to
\text{Tot}(\mathcal{K}^\bullet
\otimes_{\mathcal{O}_X}
\mathcal{M}^\bullet)
\end{align*}
由引理 \ref{cohomology-lemma-adjoints-push-pull-compatibility}，这也等于
\begin{align*}
Lf^*(f_*\mathcal{K}^\bullet
\otimes_{\mathcal{O}_Y}^\mathbf{L}
f_*\mathcal{M}^\bullet)
& =
Lf^*f_*\mathcal{K}^\bullet
\otimes_{\mathcal{O}_Y}^\mathbf{L}
Lf^*f_*\mathcal{M}^\bullet \\
& \to
f^*f_*\mathcal{K}^\bullet
\otimes_{\mathcal{O}_Y}^\mathbf{L}
f^*f_*\mathcal{M}^\bullet \\
& \to
\mathcal{K}^\bullet
\otimes_{\mathcal{O}_Y}^\mathbf{L}
\mathcal{M}^\bullet \\
& \to
\text{Tot}(\mathcal{K}^\bullet
\otimes_{\mathcal{O}_X}
\mathcal{M}^\bullet)
\end{align*}
沿图逆时针绕行，得到与下列映射伴随的映射：
\begin{align*}
Lf^*(f_*\mathcal{K}^\bullet
\otimes_{\mathcal{O}_Y}^\mathbf{L}
f_*\mathcal{M}^\bullet)
& \to
Lf^*\text{Tot}(
f_*\mathcal{K}^\bullet
\otimes_{\mathcal{O}_Y}
f_*\mathcal{M}^\bullet) \\
& \to
Lf^*f_*\text{Tot}(\mathcal{K}^\bullet
\otimes_{\mathcal{O}_X}
\mathcal{M}^\bullet) \\
& \to
Lf^*Rf_*\text{Tot}(\mathcal{K}^\bullet
\otimes_{\mathcal{O}_X}
\mathcal{M}^\bullet) \\
& \to
\text{Tot}(\mathcal{K}^\bullet
\otimes_{\mathcal{O}_X}
\mathcal{M}^\bullet)
\end{align*}
由引理 \ref{cohomology-lemma-adjoints-push-pull-compatibility}，这也等于
\begin{align*}
Lf^*(f_*\mathcal{K}^\bullet
\otimes_{\mathcal{O}_Y}^\mathbf{L}
f_*\mathcal{M}^\bullet)
& \to
Lf^*\text{Tot}(
f_*\mathcal{K}^\bullet
\otimes_{\mathcal{O}_Y}
f_*\mathcal{M}^\bullet) \\
& \to
Lf^*f_*\text{Tot}(\mathcal{K}^\bullet
\otimes_{\mathcal{O}_X}
\mathcal{M}^\bullet) \\
& \to
f^*f_*\text{Tot}(\mathcal{K}^\bullet
\otimes_{\mathcal{O}_X}
\mathcal{M}^\bullet) \\
& \to
\text{Tot}(\mathcal{K}^\bullet
\otimes_{\mathcal{O}_X}
\mathcal{M}^\bullet)
\end{align*}
考察下图即可完成证明：
$$
\xymatrix{
Lf^*(f_*\mathcal{K}^\bullet
\otimes_{\mathcal{O}_Y}^\mathbf{L}
f_*\mathcal{M}^\bullet) \ar[d] \ar[rr] & &
Lf^*f_*\mathcal{K}^\bullet \otimes_{\mathcal{O}_X}^\mathbf{L}
Lf^*f_*\mathcal{M}^\bullet \ar[d] \\
Lf^*\text{Tot}(
f_*\mathcal{K}^\bullet
\otimes_{\mathcal{O}_Y}
f_*\mathcal{M}^\bullet) \ar[d]_{\text{朴素}} \ar[r] &
f^*\text{Tot}(
f_*\mathcal{K}^\bullet
\otimes_{\mathcal{O}_Y}
f_*\mathcal{M}^\bullet) \ar[ldd]^{\text{朴素}} \ar[dd] &
f^*f_*\mathcal{K}^\bullet \otimes_{\mathcal{O}_X}^\mathbf{L}
f^*f_*\mathcal{M}^\bullet \ar[dd] \ar[ldd] \\
Lf^*f_*\text{Tot}(\mathcal{K}^\bullet
\otimes_{\mathcal{O}_X}
\mathcal{M}^\bullet) \ar[d] \\
f^*f_*\text{Tot}(\mathcal{K}^\bullet \otimes_{\mathcal{O}_X}
\mathcal{M}^\bullet) \ar[rd] &
\text{Tot}(f^*f_*\mathcal{K}^\bullet \otimes_{\mathcal{O}_X}
f^*f_*\mathcal{M}^\bullet) \ar[d] &
\mathcal{K}^\bullet \otimes_{\mathcal{O}_X}^\mathbf{L}
\mathcal{M}^\bullet \ar[ld] \\
& \text{Tot}(\mathcal{K}^\bullet
\otimes_{\mathcal{O}_X}
\mathcal{M}^\bullet)
}
$$
此图中的所有多边形都交换。最上方的多边形由引理
\ref{cohomology-lemma-tensor-pull-compatibility} 交换。
含有两个朴素杯积的方块交换，因为
$Lf^* \to f^*$ 关于模复形具有函子性。
含有下列两个映射的方块也同样交换：
$\mathcal{A}^\bullet \otimes^\mathbf{L} \mathcal{B}^\bullet \to
\text{Tot}(\mathcal{A}^\bullet \otimes \mathcal{B}^\bullet)$.
最后，剩余方块在复形层次上交换，并且可视为朴素杯积的定义
（利用 $f^*$ 与 $f_*$ 的伴随性）。沿图的外边界绕行所得的
正是上面给出的两个映射，故证明完成。
\end{proof}

\noindent
设 $(X, \mathcal{O}_X)$ 是环空间，而 $\mathcal{K}^\bullet$ 与
$\mathcal{M}^\bullet$ 是 $\mathcal{O}_X$-模复形。
则有“朴素”杯积
$$
\mu' :
\text{Tot}(
\Gamma(X, \mathcal{K}^\bullet) \otimes_A \Gamma(X, \mathcal{M}^\bullet))
\longrightarrow
\Gamma(X, \text{Tot}(
\mathcal{K}^\bullet \otimes_{\mathcal{O}_X} \mathcal{M}^\bullet))
$$
将引理 \ref{cohomology-lemma-cup-compatible-with-naive}
应用于态射 $(X, \mathcal{O}_X) \to (pt, A)$，可知此朴素杯积与
本节第一段定义的杯积 $\mu$ 由下列交换图联系：
$$
\xymatrix{
\Gamma(X, \mathcal{K}^\bullet)
\otimes_A^\mathbf{L}
\Gamma(X, \mathcal{M}^\bullet) \ar[d] \ar[r] &
R\Gamma(X, \mathcal{K}^\bullet)
\otimes_A^\mathbf{L}
R\Gamma(X, \mathcal{M}^\bullet) \ar[d]^-\mu \\
\text{Tot}(\Gamma(X, \mathcal{K}^\bullet)
\otimes_A
\Gamma(X, \mathcal{M}^\bullet)) \ar[d]_-{\mu'} &
R\Gamma(X, \mathcal{K}^\bullet
\otimes_{\mathcal{O}_X}^\mathbf{L}
\mathcal{M}^\bullet) \ar[d] \\
\Gamma(X, \text{Tot}(\mathcal{K}^\bullet
\otimes_{\mathcal{O}_X}
\mathcal{M}^\bullet)) \ar[r] &
R\Gamma(X, \text{Tot}(\mathcal{K}^\bullet
\otimes_{\mathcal{O}_X}
\mathcal{M}^\bullet))
}
$$
该图位于 $D(A)$ 中。在上同调上得到交换图
$$
\xymatrix{
H^i(\Gamma(X, \mathcal{K}^\bullet)) \times
H^j(\Gamma(X, \mathcal{M}^\bullet)) \ar[d] \ar[r] &
H^{i + j}(X,
\text{Tot}(\mathcal{K}^\bullet \otimes_{\mathcal{O}_X} \mathcal{M}^\bullet)) \\
H^i(X, \mathcal{K}^\bullet) \times
H^j(X, \mathcal{M}^\bullet) \ar[r]^\cup &
H^{i + j}(X, \mathcal{K}^\bullet \otimes_{\mathcal{O}_X}^\mathbf{L}
\mathcal{M}^\bullet) \ar[u]
}
$$
，它把朴素杯积与真正的杯积联系起来。

\begin{lemma}
\label{cohomology-lemma-diagrams-commute}
设 $(X, \mathcal{O}_X)$ 是环空间。设
$\mathcal{K}^\bullet$ 与 $\mathcal{M}^\bullet$
是 $\mathcal{O}_X$-模的有下界复形，并设
$\mathcal{U} : X = \bigcup_{i \in I} U_i$ 是开覆盖。则
$$
\xymatrix{
\text{Tot}(\check{\mathcal{C}}^\bullet(\mathcal{U}, \mathcal{K}^\bullet))
\otimes_A^\mathbf{L}
\text{Tot}(\check{\mathcal{C}}^\bullet(\mathcal{U}, \mathcal{M}^\bullet))
\ar[d] \ar[r] &
R\Gamma(X, \mathcal{K}^\bullet)
\otimes_A^\mathbf{L}
R\Gamma(X, \mathcal{M}^\bullet) \ar[d]^\mu \\
\text{Tot}(
\text{Tot}(\check{\mathcal{C}}^\bullet(\mathcal{U}, \mathcal{K}^\bullet))
\otimes_A
\text{Tot}(\check{\mathcal{C}}^\bullet(\mathcal{U}, \mathcal{M}^\bullet)))
\ar[d]^{(\ref{cohomology-equation-needs-signs})} &
R\Gamma(X,
\mathcal{K}^\bullet \otimes_{\mathcal{O}_X}^\mathbf{L} \mathcal{M}^\bullet)
\ar[d] \\
\text{Tot}(
\check{\mathcal{C}}^\bullet({\mathcal U},
\text{Tot}(\mathcal{K}^\bullet \otimes_{\mathcal{O}_X} \mathcal{M}^\bullet)
)) \ar[r] &
R\Gamma(X,
\text{Tot}(\mathcal{K}^\bullet \otimes_{\mathcal{O}_X} \mathcal{M}^\bullet))
}
$$
在 $D(A)$ 中交换，其中水平箭头是引理
\ref{cohomology-lemma-cech-complex-complex} 中的箭头。
\end{lemma}

\begin{proof}
如引理 \ref{cohomology-lemma-godement-resolution-bounded-below} 中那样，
选取复形的拟同构
$a : \mathcal{K}^\bullet \to \mathcal{K}_1^\bullet$ 与
$b : \mathcal{M}^\bullet \to \mathcal{M}_1^\bullet$。
由于映射 $a$ 和 $b$ 在茎上都是链同伦等价，所诱导的映射
$$
\text{Tot}(\mathcal{K}^\bullet \otimes_{\mathcal{O}_X} \mathcal{M}^\bullet)
\to
\text{Tot}(\mathcal{K}_1^\bullet \otimes_{\mathcal{O}_X} \mathcal{M}_1^\bullet)
$$
在茎上也为链同伦等价（《更多代数》引理
\ref{more-algebra-lemma-derived-tor-homotopy}），因而是拟同构。因此两个图的目标
$$
R\Gamma(X,
\text{Tot}(\mathcal{K}^\bullet
\otimes_{\mathcal{O}_X} \mathcal{M}^\bullet)) =
R\Gamma(X,
\text{Tot}(\mathcal{K}_1^\bullet
\otimes_{\mathcal{O}_X} \mathcal{M}_1^\bullet))
$$
在 $D(A)$ 中相同。因此只需证明把 $\mathcal{K}$ 与 $\mathcal{M}$
分别替换为 $\mathcal{K}_1$ 与 $\mathcal{M}_1$ 后该图交换。
这就归结为下一段讨论的情形。

\medskip\noindent
假设 $\mathcal{K}^\bullet$ 与 $\mathcal{M}^\bullet$
是松弛 $\mathcal{O}_X$-模的有下界复形，并考察把杯积与
{\v C}ech 复形上的杯积 (\ref{cohomology-equation-needs-signs}) 联系起来的图。
此时可考察交换图
$$
\xymatrix{
\Gamma(X, \mathcal{K}^\bullet)
\otimes_A^\mathbf{L}
\Gamma(X, \mathcal{M}^\bullet) \ar[d] \ar[r] &
\text{Tot}(\check{\mathcal{C}}^\bullet(\mathcal{U}, \mathcal{K}^\bullet))
\otimes_A^\mathbf{L}
\text{Tot}(\check{\mathcal{C}}^\bullet(\mathcal{U}, \mathcal{M}^\bullet))
\ar[d] \\
\text{Tot}(\Gamma(X, \mathcal{K}^\bullet)
\otimes_A
\Gamma(X, \mathcal{M}^\bullet)) \ar[d] \ar[r] &
\text{Tot}(
\text{Tot}(\check{\mathcal{C}}^\bullet(\mathcal{U}, \mathcal{K}^\bullet))
\otimes_A
\text{Tot}(\check{\mathcal{C}}^\bullet(\mathcal{U}, \mathcal{M}^\bullet)))
\ar[d]^{(\ref{cohomology-equation-needs-signs})} \\
\Gamma(X, \text{Tot}(\mathcal{K}^\bullet
\otimes_{\mathcal{O}_X}
\mathcal{M}^\bullet)) \ar[r] &
\text{Tot}(
\check{\mathcal{C}}^\bullet({\mathcal U},
\text{Tot}(\mathcal{K}^\bullet \otimes_{\mathcal{O}_X} \mathcal{M}^\bullet)
))
}
$$
此图中的水平箭头是 $D(A)$ 中的同构，因为对
$\mathcal{K}^\bullet$ 这样的松弛模有下界复形，有
$$
\Gamma(X, \mathcal{K}^\bullet) =
\text{Tot}(\check{\mathcal{C}}^\bullet(\mathcal{U}, \mathcal{K}^\bullet)) =
R\Gamma(X, \mathcal{K}^\bullet)
$$
，此等式位于 $D(A)$ 中。它由引理 \ref{cohomology-lemma-flasque-acyclic}、
《导出范畴》引理 \ref{derived-lemma-leray-acyclicity} 以及引理 \ref{cohomology-lemma-cech-complex-complex-computes} 得到。
因此，引理之图中涉及 (\ref{cohomology-equation-needs-signs}) 的部分交换性，
由已经证明的引理 \ref{cohomology-lemma-cup-compatible-with-naive}
在 $f$ 为到一点的态射时的交换性得到（见引理
\ref{cohomology-lemma-cup-compatible-with-naive} 后的讨论）。
\end{proof}

\begin{lemma}
\label{cohomology-lemma-cup-product-associative}
设 $f : (X, \mathcal{O}_X) \to (Y, \mathcal{O}_Y)$
是环空间的态射。注记 \ref{cohomology-remark-cup-product} 的相对杯积具有结合性，
意即图
$$
\xymatrix{
Rf_*K \otimes_{\mathcal{O}_Y}^\mathbf{L}
Rf_*L \otimes_{\mathcal{O}_Y}^\mathbf{L}
Rf_*M \ar[r] \ar[d] &
Rf_*(K \otimes_{\mathcal{O}_X}^\mathbf{L} L)
\otimes_{\mathcal{O}_Y}^\mathbf{L} Rf_*M \ar[d] \\
Rf_*K \otimes_{\mathcal{O}_Y}^\mathbf{L}
Rf_*(L \otimes_{\mathcal{O}_X}^\mathbf{L} M) \ar[r] &
Rf_*(K \otimes_{\mathcal{O}_X}^\mathbf{L} 
L \otimes_{\mathcal{O}_X}^\mathbf{L} M)
}
$$
对 $D(\mathcal{O}_X)$ 中所有 $K, L, M$，都在
$D(\mathcal{O}_Y)$ 中交换。
\end{lemma}

\begin{proof}
沿任一侧绕行，都会得到与下列显然映射伴随的映射：
\begin{align*}
Lf^*(Rf_*K \otimes_{\mathcal{O}_Y}^\mathbf{L}
Rf_*L \otimes_{\mathcal{O}_Y}^\mathbf{L}
Rf_*M) & =
Lf^*(Rf_*K) \otimes_{\mathcal{O}_X}^\mathbf{L}
Lf^*(Rf_*L) \otimes_{\mathcal{O}_X}^\mathbf{L}
Lf^*(Rf_*M) \\
& \to
K \otimes_{\mathcal{O}_X}^\mathbf{L} 
L \otimes_{\mathcal{O}_X}^\mathbf{L} M
\end{align*}
该映射位于 $D(\mathcal{O}_X)$ 中。
\end{proof}

\begin{lemma}
\label{cohomology-lemma-cup-product-commutative}
设 $f : (X, \mathcal{O}_X) \to (Y, \mathcal{O}_Y)$
是环空间的态射。注记 \ref{cohomology-remark-cup-product} 的相对杯积具有交换性，
意即图
$$
\xymatrix{
Rf_*K \otimes_{\mathcal{O}_Y}^\mathbf{L} Rf_*L \ar[r] \ar[d]_\psi &
Rf_*(K \otimes_{\mathcal{O}_X}^\mathbf{L} L) \ar[d]^{Rf_*\psi} \\
Rf_*L \otimes_{\mathcal{O}_Y}^\mathbf{L} Rf_*K \ar[r] &
Rf_*(L \otimes_{\mathcal{O}_X}^\mathbf{L} K)
}
$$
对 $D(\mathcal{O}_X)$ 中所有 $K, L$，都在
$D(\mathcal{O}_Y)$ 中交换。这里 $\psi$ 是导出范畴上的交换约束
（引理 \ref{cohomology-lemma-symmetric-monoidal-derived}）。
\end{lemma}

\begin{proof}
从略。
\end{proof}

\begin{lemma}
\label{cohomology-lemma-compose-cup-product}
设 $f : (X, \mathcal{O}_X) \to (Y, \mathcal{O}_Y)$ 与
$g : (Y, \mathcal{O}_Y) \to (Z, \mathcal{O}_Z)$
是环空间的态射。注记 \ref{cohomology-remark-cup-product} 的相对杯积与复合相容，
意即图
$$
\xymatrix{
R(g \circ f)_*K \otimes_{\mathcal{O}_Z}^\mathbf{L} R(g \circ f)_*L
\ar@{=}[rr] \ar[d] & &
Rg_*Rf_*K \otimes_{\mathcal{O}_Z}^\mathbf{L} Rg_*Rf_*L \ar[d] \\
R(g \circ f)_*(K \otimes_{\mathcal{O}_X}^\mathbf{L} L) \ar@{=}[r] &
Rg_*Rf_*(K \otimes_{\mathcal{O}_X}^\mathbf{L} L) &
Rg_*(Rf_*K \otimes_{\mathcal{O}_Y}^\mathbf{L}  Rf_*L) \ar[l]
}
$$
对 $D(\mathcal{O}_X)$ 中所有 $K, L$，都在
$D(\mathcal{O}_Z)$ 中交换。
\end{lemma}

\begin{proof}
这是因为无论沿图以哪种方式绕行，都会得到与下列映射伴随的映射：
\begin{align*}
& L(g \circ f)^*\left(R(g \circ f)_*K
\otimes_{\mathcal{O}_Z}^\mathbf{L}
R(g \circ f)_*L\right) \\
& =
L(g \circ f)^*R(g \circ f)_*K
\otimes_{\mathcal{O}_X}^\mathbf{L}
L(g \circ f)^*R(g \circ f)_*L) \\
& \to
K \otimes_{\mathcal{O}_X}^\mathbf{L} L
\end{align*}
该映射位于 $D(\mathcal{O}_X)$ 中。为此要用到如下余单位复合：
$$
L(g \circ f)^*R(g \circ f)_* =
Lf^* Lg^* Rg_* Rf_*  \to
Lf^* Rf_* \to \text{id}
$$
它是 $L(g \circ f)^*$ 与 $R(g \circ f)_*$ 的余单位；参见《范畴》引理
\ref{categories-lemma-compose-counits}。
\end{proof}

\begin{lemma}
\label{cohomology-lemma-base-change-cup-product}
考察环空间的交换方块
$$
\xymatrix{
(X', \mathcal{O}_{X'}) \ar[r]_{g'} \ar[d]_{f'} &
(X, \mathcal{O}_X) \ar[d]^f \\
(Y', \mathcal{O}_{Y'}) \ar[r]^g & (Y, \mathcal{O}_Y) 
}
$$
。设 $K, L$ 属于 $D(\mathcal{O}_X)$。相对杯积与此方块相容，
意即图
$$
\xymatrix{
Lg^*(Rf_*K \otimes_{\mathcal{O}_Y}^\mathbf{L} Rf_*L)
\ar[r] \ar@{=}[d] &
Lg^*(Rf_*(K \otimes_{\mathcal{O}_X}^\mathbf{L} L)) \ar[d] \\
Lg^*Rf_*K \otimes_{\mathcal{O}_{Y'}}^\mathbf{L} Lg^*Rf_*L \ar[d] &
R(f')_*L(g')^*(K \otimes_{\mathcal{O}_X}^\mathbf{L} L) \ar@{=}[d] \\
R(f')_*(L(g')^*K \otimes_{\mathcal{O}_{Y'}} R(f')_*(L(g')^*L \ar[r] &
R(f')_*(L(g')^*K \otimes_{\mathcal{O}_{X'}}^\mathbf{L} L(g')^*L)
}
$$
在 $D(\mathcal{O}_{Y'})$ 中交换。水平箭头由相对杯积给出
（注记 \ref{cohomology-remark-cup-product}），竖直箭头由基变换映射
（注记 \ref{cohomology-remark-base-change}）与引理
\ref{cohomology-lemma-pullback-tensor-product} 给出。
\end{lemma}

\begin{proof}
从略。
\end{proof}














\section{K-内射复形的一些性质}
\label{cohomology-section-properties-K-injective}

\noindent
设 $(X, \mathcal{O}_X)$ 是环空间，$U \subset X$ 是开子集，并以
$j : (U, \mathcal{O}_U) \to (X, \mathcal{O}_X)$
表示相应的开浸入。拉回函子 $j^*$ 就是限制函子，因而正合。
所以在任意复形上，导出拉回 $Lj^*$ 只需限制该复形即可计算。
我们通常把相应函子简记为
$$
D(\mathcal{O}_X) \to D(\mathcal{O}_U),
\quad
E \mapsto j^*E = E|_U
$$
类似地，零延拓
$j_! : \textit{Mod}(\mathcal{O}_U) \to \textit{Mod}(\mathcal{O}_X)$
是正合函子（参见《层》第 \ref{sheaves-section-open-immersions} 节及
《模》引理 \ref{modules-lemma-j-shriek-exact}）。因此它诱导函子
$$
j_! : D(\mathcal{O}_U) \to D(\mathcal{O}_X),\quad
F \mapsto j_!F
$$
：只需对表示对象 $F$ 的任意复形逐项应用 $j_!$。

\begin{lemma}
\label{cohomology-lemma-restrict-K-injective-to-open}
设 $X$ 是环空间，而 $U \subset X$ 是开子空间。
$\mathcal{O}_X$-模的 K-内射复形限制到 $U$ 后，
是 $\mathcal{O}_U$-模的 K-内射复形。
\end{lemma}

\begin{proof}
这由《导出范畴》引理
\ref{derived-lemma-adjoint-preserve-K-injectives}
以及限制函子具有正合左伴随 $j_!$ 这一事实推出。
$j_!$ 的构造见《层》第 \ref{sheaves-section-open-immersions} 节，
其正合性见《模》引理 \ref{modules-lemma-j-shriek-exact}。
\end{proof}

\begin{lemma}
\label{cohomology-lemma-unbounded-cohomology-of-open}
设 $X$ 是环空间，而 $U \subset X$ 是开子空间。
对 $D(\mathcal{O}_X)$ 中的 $K$，有
$H^p(U, K) = H^p(U, K|_U)$.
\end{lemma}

\begin{proof}
取表示 $K$ 的 $\mathcal{O}_X$-模 K-内射复形
$\mathcal{I}^\bullet$。则
$$
H^q(U, K) = H^q(\Gamma(U, \mathcal{I}^\bullet)) =
H^q(\Gamma(U, \mathcal{I}^\bullet|_U))
$$
。这是由上同调的构造得到的。由引理
\ref{cohomology-lemma-restrict-K-injective-to-open}，
复形 $\mathcal{I}^\bullet|_U$ 是表示 $K|_U$ 的 K-内射复形，
故引理成立。
\end{proof}

\begin{lemma}
\label{cohomology-lemma-sheafification-cohomology}
设 $(X, \mathcal{O}_X)$ 是环空间，而 $K$ 是
$D(\mathcal{O}_X)$ 的对象。预层
$$
U \mapsto H^q(U, K) = H^q(U, K|_U)
$$
的层化是 $K$ 的第 $q$ 个上同调层 $H^q(K)$。
\end{lemma}

\begin{proof}
等式 $H^q(U, K) = H^q(U, K|_U)$ 由引理
\ref{cohomology-lemma-unbounded-cohomology-of-open} 成立。
选取表示 $K$ 的 K-内射复形 $\mathcal{I}^\bullet$。则
$$
H^q(U, K) =
\frac{\Ker(\mathcal{I}^q(U) \to \mathcal{I}^{q + 1}(U))}
{\Im(\mathcal{I}^{q - 1}(U) \to \mathcal{I}^q(U))}.
$$
此式由我们的上同调构造得到。又因为
$H^q(K) = \Ker(\mathcal{I}^q \to \mathcal{I}^{q + 1})/
\Im(\mathcal{I}^{q - 1} \to \mathcal{I}^q)$，结论显然。
\end{proof}

\begin{lemma}
\label{cohomology-lemma-restrict-direct-image-open}
设 $f : (X, \mathcal{O}_X) \to (Y, \mathcal{O}_Y)$ 是环空间的态射。
给定开子空间 $V \subset Y$，令 $U = f^{-1}(V)$，并以
$g : U \to V$ 表示所诱导的态射。则对 $D(\mathcal{O}_X)$ 中的 $E$，
有 $(Rf_*E)|_V = Rg_*(E|_U)$。
\end{lemma}

\begin{proof}
用 $\mathcal{O}_X$-模 K-内射复形
$\mathcal{I}^\bullet$ 表示 $E$。则
$Rf_*(E) = f_*\mathcal{I}^\bullet$，并且由引理
\ref{cohomology-lemma-restrict-K-injective-to-open}，
$Rg_*(E|_U) = g_*(\mathcal{I}^\bullet|_U)$。
对 $X$ 上任意层 $\mathcal{F}$，显然有
$(f_*\mathcal{F})|_V = g_*(\mathcal{F}|_U)$，故结论成立。
\end{proof}

\begin{lemma}
\label{cohomology-lemma-Leray-unbounded}
设 $f : X \to Y$ 是环空间的态射。作为函子
$D(\mathcal{O}_X) \to D(\Gamma(Y, \mathcal{O}_Y))$，有
$R\Gamma(Y, -) \circ Rf_* = R\Gamma(X, -)$。
更一般地，若 $V \subset Y$ 开且 $U = f^{-1}(V)$，则
$R\Gamma(U, -) = R\Gamma(V, -) \circ Rf_*$。
\end{lemma}

\begin{proof}
令 $Z$ 为底空间是单点且满足
$\Gamma(Z, \mathcal{O}_Z) = \Gamma(Y, \mathcal{O}_Y)$ 的环空间。
存在环空间的典范态射 $Y \to Z$，它在结构层的整体截面上诱导上述等同。
于是 $D(\mathcal{O}_Z) = D(\Gamma(Y, \mathcal{O}_Y))$。
因此，把引理 \ref{cohomology-lemma-derived-pushforward-composition}
应用于 $X \to Y \to Z$，即得
$R\Gamma(Y, -) \circ Rf_* = R\Gamma(X, -)$。

\medskip\noindent
将引理 \ref{cohomology-lemma-restrict-direct-image-open} 应用于第一个断言，
即得第二个（更一般的）断言。
\end{proof}

\begin{lemma}
\label{cohomology-lemma-unbounded-describe-higher-direct-images}
设 $f : (X, \mathcal{O}_X) \to (Y, \mathcal{O}_Y)$ 是环空间的态射，
而 $K$ 属于 $D(\mathcal{O}_X)$。则 $H^i(Rf_*K)$ 是与下列预层相伴的层：
$$
V \mapsto H^i(f^{-1}(V), K) = H^i(V, Rf_*K)
$$
\end{lemma}

\begin{proof}
对引理 \ref{cohomology-lemma-Leray-unbounded} 的第二个断言取上同调，
即得等式 $H^i(f^{-1}(V), K) = H^i(V, Rf_*K)$。
层化断言于是由引理 \ref{cohomology-lemma-sheafification-cohomology} 得到。
\end{proof}

\begin{lemma}
\label{cohomology-lemma-modules-abelian-unbounded}
设 $X$ 是环空间，$K$ 是 $D(\mathcal{O}_X)$ 的对象，并以
$K_{ab}$ 表示它在 $D(\underline{\mathbf{Z}}_X)$ 中的像。
\begin{enumerate}
\item 对任意开集 $U \subset X$，存在典范映射
$R\Gamma(U, K) \to R\Gamma(U, K_{ab})$
，它是 $D(\textit{Ab})$ 中的同构。
\item 设 $f : X \to Y$ 是环空间的态射。存在典范映射
$Rf_*K \to Rf_*(K_{ab})$，它是
$D(\underline{\mathbf{Z}}_Y)$ 中的同构。
\end{enumerate}
\end{lemma}

\begin{proof}
映射构造如下。选取表示 $K$ 的 K-内射复形
$\mathcal{I}^\bullet$，并选取拟同构
$\mathcal{I}^\bullet \to \mathcal{J}^\bullet$，其中
$\mathcal{J}^\bullet$ 是阿贝尔群的 K-内射复形。则 (1) 中的映射为
$\Gamma(U, \mathcal{I}^\bullet) \to \Gamma(U, \mathcal{J}^\bullet)$，
而 (2) 中的映射为
$f_*\mathcal{I}^\bullet \to f_*\mathcal{J}^\bullet$。
要证明这些映射是同构，只需证明它们在上同调群和上同调层上诱导同构。
由引理 \ref{cohomology-lemma-unbounded-cohomology-of-open} 与
\ref{cohomology-lemma-unbounded-describe-higher-direct-images}，只需证明映射
$$
H^0(X, K) \longrightarrow H^0(X, K_{ab})
$$
是同构。注意
$$
H^0(X, K) = \Hom_{D(\mathcal{O}_X)}(\mathcal{O}_X, K)
$$
，另一个群也有类似表达式。选取表示 $K$ 的任意
$\mathcal{O}_X$-模复形 $\mathcal{K}^\bullet$。
由导出范畴作为局部化的构造，有
$$
\Hom_{D(\mathcal{O}_X)}(\mathcal{O}_X, K) =
\colim_{s : \mathcal{F}^\bullet \to \mathcal{O}_X}
\Hom_{K(\mathcal{O}_X)}(\mathcal{F}^\bullet, \mathcal{K}^\bullet)
$$
，其中余极限取遍 $\mathcal{O}_X$-模复形的拟同构 $s$。类似地，有
$$
\Hom_{D(\underline{\mathbf{Z}}_X)}(\underline{\mathbf{Z}}_X, K) =
\colim_{s : \mathcal{G}^\bullet \to \underline{\mathbf{Z}}_X}
\Hom_{K(\underline{\mathbf{Z}}_X)}(\mathcal{G}^\bullet, \mathcal{K}^\bullet)
$$
接着注意，拟同构
$s : \mathcal{G}^\bullet \to \underline{\mathbf{Z}}_X$
中满足 $\mathcal{G}^\bullet$ 是平坦
$\underline{\mathbf{Z}}_X$-模的有上界复形者，在该系统中余终。
（这由《模》引理 \ref{modules-lemma-module-quotient-flat}
与《导出范畴》引理 \ref{derived-lemma-subcategory-left-resolution} 得到；
参见第 \ref{cohomology-section-flat} 节的讨论。）因此可以构造映射
$H^0(X, K) \longrightarrow H^0(X, K_{ab})$
的逆映射：把元素 $\xi \in H^0(X, K_{ab})$ 表示为一对
$$
(s : \mathcal{G}^\bullet \to \underline{\mathbf{Z}}_X,
a : \mathcal{G}^\bullet \to \mathcal{K}^\bullet)
$$
，其中 $\mathcal{G}^\bullet$ 是平坦
$\underline{\mathbf{Z}}_X$-模的有上界复形，并将其映到
$$
(\mathcal{G}^\bullet \otimes_{\underline{\mathbf{Z}}_X} \mathcal{O}_X
\to \mathcal{O}_X,
\mathcal{G}^\bullet  \otimes_{\underline{\mathbf{Z}}_X} \mathcal{O}_X
\to \mathcal{K}^\bullet)
$$
。这里唯一需要注意的是，由引理
\ref{cohomology-lemma-derived-tor-quasi-isomorphism-other-side} 与
\ref{cohomology-lemma-bounded-flat-K-flat}，第一个箭头是拟同构。
关于此构造确为逆映射的详细验证从略。
\end{proof}

\begin{lemma}
\label{cohomology-lemma-adjoint-lower-shriek-restrict}
设 $(X, \mathcal{O}_X)$ 是环空间，$U \subset X$ 是开子集，并以
$j : (U, \mathcal{O}_U) \to (X, \mathcal{O}_X)$
表示相应的开浸入。限制函子
$D(\mathcal{O}_X) \to D(\mathcal{O}_U)$ 是零延拓
$j_! : D(\mathcal{O}_U) \to D(\mathcal{O}_X)$ 的右伴随。
\end{lemma}

\begin{proof}
这由 $j_!$ 与 $j^*$ 互为伴随且均正合这一事实形式地推出
（因而 $Lj_! = j_!$ 与 $Rj^* = j^*$ 存在）；参见《导出范畴》引理
\ref{derived-lemma-derived-adjoint-functors}。
\end{proof}

\begin{lemma}
\label{cohomology-lemma-K-injective-flat}
设 $f : X \to Y$ 是环空间的平坦态射。
若 $\mathcal{I}^\bullet$ 是 $\mathcal{O}_X$-模的 K-内射复形，
则 $f_*\mathcal{I}^\bullet$ 作为 $\mathcal{O}_Y$-模复形是 K-内射的。
\end{lemma}

\begin{proof}
这是因为
$$
\Hom_{K(\mathcal{O}_Y)}(\mathcal{F}^\bullet, f_*\mathcal{I}^\bullet)
=
\Hom_{K(\mathcal{O}_X)}(f^*\mathcal{F}^\bullet, \mathcal{I}^\bullet)
$$
；此等式由《层》引理
\ref{sheaves-lemma-adjoint-pullback-pushforward-modules}
以及 $f$ 平坦从而 $f^*$ 正合这一事实得到。
\end{proof}





\section{无界 Mayer--Vietoris}
\label{cohomology-section-unbounded-mayer-vietoris}

\noindent
无界上同调也有 Mayer--Vietoris 序列。

\begin{lemma}
\label{cohomology-lemma-exact-sequence-lower-shriek}
设 $(X, \mathcal{O}_X)$ 是环空间，且
$X = U \cup V$ 是两个开子空间之并。对 $D(\mathcal{O}_X)$
的任意对象 $E$，有特异三角
$$
j_{U \cap V!}E|_{U \cap V} \to
j_{U!}E|_U \oplus j_{V!}E|_V \to E \to 
j_{U \cap V!}E|_{U \cap V}[1]
$$
，它位于 $D(\mathcal{O}_X)$ 中。
\end{lemma}

\begin{proof}
我们已在第 \ref{cohomology-section-properties-K-injective} 节看到，
限制函子与零延拓函子在任意复形上都只需直接应用相应函子即可计算。
设 $\mathcal{E}^\bullet$ 是表示 $E$ 的
$\mathcal{O}_X$-模复形。引理的特异三角，是与下列
$\mathcal{O}_X$-模复形短正合列相伴的特异三角
（见《导出范畴》第 \ref{derived-section-canonical-delta-functor} 节，
尤其是引理 \ref{derived-lemma-derived-canonical-delta-functor}）：
$$
0 \to j_{U \cap V!}\mathcal{E}^\bullet|_{U \cap V} \to
j_{U!}\mathcal{E}^\bullet|_U \oplus j_{V!}\mathcal{E}^\bullet|_V
\to \mathcal{E}^\bullet \to 0
$$
利用《层》引理 \ref{sheaves-lemma-j-shriek-modules}
在茎上检验即可看出此序列正合（计算从略）。
\end{proof}

\begin{lemma}
\label{cohomology-lemma-exact-sequence-j-star}
设 $(X, \mathcal{O}_X)$ 是环空间，且
$X = U \cup V$ 是两个开子空间之并。对 $D(\mathcal{O}_X)$
的任意对象 $E$，有特异三角
$$
E \to 
Rj_{U, *}E|_U \oplus Rj_{V, *}E|_V \to
Rj_{U \cap V, *}E|_{U \cap V} \to
E[1]
$$
，它位于 $D(\mathcal{O}_X)$ 中。
\end{lemma}

\begin{proof}
选取表示 $E$ 的 K-内射复形 $\mathcal{I}^\bullet$，
使其各项 $\mathcal{I}^n$ 都是
$\textit{Mod}(\mathcal{O}_X)$ 的内射对象；参见《内射对象》定理
\ref{injectives-theorem-K-injective-embedding-grothendieck}。
我们已经看到 $\mathcal{I}^\bullet|U$ 也是 K-内射复形
（引理 \ref{cohomology-lemma-restrict-K-injective-to-open}）。因此
$Rj_{U, *}E|_U$ 由 $j_{U, *}\mathcal{I}^\bullet|_U$ 表示；
$V$ 与 $U \cap V$ 的情形类似。于是引理的特异三角，是与下列
复形短正合列相伴的特异三角
（见《导出范畴》第 \ref{derived-section-canonical-delta-functor} 节，
尤其是引理 \ref{derived-lemma-derived-canonical-delta-functor}）：
$$
0 \to
\mathcal{I}^\bullet \to
j_{U, *}\mathcal{I}^\bullet|_U \oplus j_{V, *}\mathcal{I}^\bullet|_V \to
j_{U \cap V, *}\mathcal{I}^\bullet|_{U \cap V} \to
0.
$$
此序列正合，因为对任意开集 $W \subset X$ 及任意 $n$，序列
$$
0 \to
\mathcal{I}^n(W) \to
\mathcal{I}^n(W \cap U) \oplus \mathcal{I}^n(W \cap V) \to
\mathcal{I}^n(W \cap U \cap V) \to
0
$$
正合（见引理 \ref{cohomology-lemma-mayer-vietoris} 的证明）。
\end{proof}

\begin{lemma}
\label{cohomology-lemma-mayer-vietoris-hom}
设 $(X, \mathcal{O}_X)$ 是环空间，且 $X = U \cup V$
是 $X$ 的两个开子空间之并。对 $D(\mathcal{O}_X)$ 的对象
$E,F$，有 Mayer--Vietoris 序列
$$
\xymatrix{
& \ldots \ar[r] & \Ext^{-1}(E_{U \cap V}, F_{U \cap V}) \ar[lld] \\
\Hom(E, F) \ar[r] &
\Hom(E_U, F_U) \oplus
\Hom(E_V, F_V) \ar[r] &
\Hom(E_{U \cap V}, F_{U \cap V})
}
$$
，其中下标表示限制到相应开集，而 $\Hom$ 与 $\Ext$
均取自相应的导出范畴。
\end{lemma}

\begin{proof}
使用引理 \ref{cohomology-lemma-exact-sequence-lower-shriek} 的特异三角，
得到 $\Hom$ 的长正合列
（由《导出范畴》引理 \ref{derived-lemma-representable-homological}），
再使用
$$
\Hom_{D(\mathcal{O}_X)}(j_{U!}E|_U, F) =
\Hom_{D(\mathcal{O}_U)}(E|_U, F|_U)
$$
；此等式由引理 \ref{cohomology-lemma-adjoint-lower-shriek-restrict} 得到。
\end{proof}

\begin{lemma}
\label{cohomology-lemma-unbounded-mayer-vietoris}
设 $(X, \mathcal{O}_X)$ 是环空间，并假设
$X = U \cup V$ 是两个开子集之并。对 $D(\mathcal{O}_X)$
的对象 $E$，有特异三角
$$
R\Gamma(X, E) \to R\Gamma(U, E) \oplus R\Gamma(V, E) \to
R\Gamma(U \cap V, E) \to R\Gamma(X, E)[1]
$$
，从而特别有上同调长正合列
$$
\ldots \to
H^n(X, E) \to
H^n(U, E) \oplus H^0(V, E) \to
H^n(U \cap V, E) \to
H^{n + 1}(X, E) \to \ldots
$$
特异三角与长正合列的构造关于 $E$ 具有函子性。
\end{lemma}

\begin{proof}
选取表示 $E$ 的 K-内射复形 $\mathcal{I}^\bullet$。
可假设对所有 $n$，$\mathcal{I}^n$ 都是
$\textit{Mod}(\mathcal{O}_X)$ 的内射对象；参见《内射对象》定理
\ref{injectives-theorem-K-injective-embedding-grothendieck}。
于是 $R\Gamma(X, E)$ 由 $\Gamma(X, \mathcal{I}^\bullet)$ 计算。
由引理 \ref{cohomology-lemma-restrict-K-injective-to-open}，
$U$、$V$ 及 $U \cap V$ 的情形也类似。因此，引理的特异三角
是与下列复形短正合列相伴的特异三角
（见《导出范畴》第 \ref{derived-section-canonical-delta-functor} 节，
尤其是引理 \ref{derived-lemma-derived-canonical-delta-functor}）：
$$
0 \to
\mathcal{I}^\bullet(X) \to
\mathcal{I}^\bullet(U) \oplus \mathcal{I}^\bullet(V) \to
\mathcal{I}^\bullet(U \cap V) \to
0.
$$
我们已在引理 \ref{cohomology-lemma-mayer-vietoris} 的证明中看到这是短正合列。
最后一个断言由《内射对象》定理
\ref{injectives-theorem-K-injective-embedding-grothendieck}
中构造的函子性得到。
\end{proof}

\begin{lemma}
\label{cohomology-lemma-unbounded-relative-mayer-vietoris}
设 $f : X \to Y$ 是环空间的态射，并假设
$X = U \cup V$ 是两个开子集之并。以
$a = f|_U : U \to Y$、$b = f|_V : V \to Y$ 及
$c = f|_{U \cap V} : U \cap V \to Y$
表示相应的限制。对 $D(\mathcal{O}_X)$ 的每个对象 $E$，
存在特异三角
$$
Rf_*E \to
Ra_*(E|_U) \oplus Rb_*(E|_V) \to
Rc_*(E|_{U \cap V}) \to
Rf_*E[1]
$$
此三角关于 $E$ 具有函子性。
\end{lemma}

\begin{proof}
选取表示 $E$ 的 K-内射复形 $\mathcal{I}^\bullet$。
可假设对所有 $n$，$\mathcal{I}^n$ 都是
$\textit{Mod}(\mathcal{O}_X)$ 的内射对象；参见《内射对象》定理
\ref{injectives-theorem-K-injective-embedding-grothendieck}。
于是 $Rf_*E$ 由 $f_*\mathcal{I}^\bullet$ 计算。
由引理 \ref{cohomology-lemma-restrict-K-injective-to-open}，
$U$、$V$ 及 $U \cap V$ 的情形也类似。因此，引理的特异三角
是与下列复形短正合列相伴的特异三角
（见《导出范畴》第 \ref{derived-section-canonical-delta-functor} 节，
尤其是引理 \ref{derived-lemma-derived-canonical-delta-functor}）：
$$
0 \to
f_*\mathcal{I}^\bullet \to
a_*\mathcal{I}^\bullet|_U \oplus b_*\mathcal{I}^\bullet|_V \to
c_*\mathcal{I}^\bullet|_{U \cap V} \to
0.
$$
这是复形的短正合列；这由引理
\ref{cohomology-lemma-relative-mayer-vietoris} 以及下述事实得到：
对 $\textit{Mod}(\mathcal{O}_X)$ 的内射对象 $\mathcal{I}$，
$R^1f_*\mathcal{I} = 0$。
最后一个断言由《内射对象》定理
\ref{injectives-theorem-K-injective-embedding-grothendieck}
中构造的函子性得到。
\end{proof}

\begin{lemma}
\label{cohomology-lemma-pushforward-restriction}
设 $(X, \mathcal{O}_X)$ 是环空间。设 $j : U \to X$ 是开子空间，
而 $T \subset X$ 是包含于 $U$ 的闭子集。
\begin{enumerate}
\item 若 $E$ 是 $D(\mathcal{O}_X)$ 的对象，且其上同调层支撑在
$T$ 上，则 $E \to Rj_*(E|_U)$ 是同构。
\item 若 $F$ 是 $D(\mathcal{O}_U)$ 的对象，且其上同调层支撑在
$T$ 上，则 $j_!F \to Rj_*F$ 是同构。
\end{enumerate}
\end{lemma}

\begin{proof}
令 $V = X \setminus T$ 且 $W = U \cap V$。注意
$X = U \cup V$ 是 $X$ 的开覆盖，并以
$j_W : W \to V$ 表示开浸入。
设 $E$ 是 $D(\mathcal{O}_X)$ 的对象，且其上同调层支撑在 $T$ 上。
由引理 \ref{cohomology-lemma-restrict-direct-image-open}，有
$(Rj_*E|_U)|_V = Rj_{W, *}(E|_W) = 0$，因为根据假设 $E|_W = 0$。
另一方面，$Rj_*(E|_U)|_U = E|_U$。故 (1) 显然。
设 $F$ 是 $D(\mathcal{O}_U)$ 的对象，且其上同调层支撑在 $T$ 上。
由引理 \ref{cohomology-lemma-restrict-direct-image-open}，有
$(Rj_*F)|_V = Rj_{W, *}(F|_W) = 0$，因为根据假设 $F|_W = 0$。
还有 $(j_!F)|_V = j_{W!}(F|_W) = 0$
（第一个等式直接来自零延拓的定义）。由于
$(Rj_*F)|_U = F$ 且 $(j_!F)|_U = F$，可知 (2) 成立。
\end{proof}

\begin{lemma}
\label{cohomology-lemma-mayer-vietoris-cup}
设 $(X, \mathcal{O}_X)$ 是环空间，并令
$A = \Gamma(X, \mathcal{O}_X)$。假设 $X = U \cup V$
是两个开子集之并。对 $D(\mathcal{O}_X)$ 的对象 $K$ 与 $M$，
有特异三角的态射
$$
\xymatrix{
R\Gamma(X, K) \otimes_A^\mathbf{L} R\Gamma(X, M) \ar[r] \ar[d] &
R\Gamma(X, K \otimes_{\mathcal{O}_X}^\mathbf{L} M) \ar[d] \\
R\Gamma(X, K) \otimes_A^\mathbf{L}
(R\Gamma(U, M) \oplus R\Gamma(V, M)) \ar[r] \ar[d] &
R\Gamma(U, K \otimes_{\mathcal{O}_X}^\mathbf{L} M)
\oplus R\Gamma(V, K \otimes_{\mathcal{O}_X}^\mathbf{L} M)) \ar[d] \\
R\Gamma(X, K) \otimes_A^\mathbf{L} R\Gamma(U \cap V, M) \ar[r] \ar[d] &
R\Gamma(U \cap V, K \otimes_{\mathcal{O}_X}^\mathbf{L} M) \ar[d] \\
R\Gamma(X, K) \otimes_A^\mathbf{L} R\Gamma(X, M)[1] \ar[r] &
R\Gamma(X, K \otimes_{\mathcal{O}_X}^\mathbf{L} M)[1]
}
$$
，其中
\begin{enumerate}
\item 水平箭头由杯积给出；
\item 右侧是引理 \ref{cohomology-lemma-unbounded-mayer-vietoris}
对 $K \otimes_{\mathcal{O}_X}^\mathbf{L} M$ 所给出的特异三角；
\item 左侧是将正合函子
$R\Gamma(X, K) \otimes_A^\mathbf{L} - $
应用于引理 \ref{cohomology-lemma-unbounded-mayer-vietoris} 对 $M$
所给出的特异三角。
\end{enumerate}
\end{lemma}

\begin{proof}
选取表示 $R\Gamma(X, K)$ 的平坦 $A$-模 K-平坦复形
$T^\bullet$；参见《更多代数》引理
\ref{more-algebra-lemma-K-flat-resolution}。
以 $T^\bullet \otimes_A \mathcal{O}_X$ 表示 $T^\bullet$
沿环空间态射 $(X, \mathcal{O}_X) \to (pt, A)$ 的拉回。
在 $D(\mathcal{O}_X)$ 中有自然的伴随映射
$\epsilon : T^\bullet \otimes_A \mathcal{O}_X \to K$。
注意 $T^\bullet \otimes_A \mathcal{O}_X$
是逐项平坦的 $\mathcal{O}_X$-模 K-平坦复形；参见引理
\ref{cohomology-lemma-pullback-K-flat} 与《模》引理
\ref{modules-lemma-pullback-flat}。
由引理 \ref{cohomology-lemma-factor-through-K-flat}，可找到复形态射
$$
T^\bullet \otimes_A \mathcal{O}_X \longrightarrow \mathcal{K}^\bullet
$$
，它是表示 $\epsilon$ 的 $\mathcal{O}_X$-模复形态射，
且 $\mathcal{K}^\bullet$ 是逐项平坦的 K-平坦复形。
具体地，由 $D(\mathcal{O}_X)$ 的构造，可先用某个表示
$\epsilon$ 的 $\mathcal{O}_X$-模复形态射
$e : T^\bullet \otimes_A \mathcal{O}_X \to \mathcal{L}^\bullet$
表示 $\epsilon$，再把该引理应用于 $e$。
选取表示 $M$ 的 K-内射复形 $\mathcal{I}^\bullet$，
其各项都是内射 $\mathcal{O}_X$-模。最后选取拟同构
$$
\text{Tot}(\mathcal{K}^\bullet \otimes_\mathcal{O} \mathcal{I}^\bullet)
\longrightarrow
\mathcal{J}^\bullet
$$
，其目标为逐项是内射 $\mathcal{O}_X$-模的 K-内射复形。
注意，此箭头的源与目标都在 $D(\mathcal{O}_X)$ 中表示
$K \otimes_{\mathcal{O}_X}^\mathbf{L} M$。
此时，对任意开集 $W \subset X$，得到复形态射
$$
\text{Tot}(T^\bullet \otimes_A \mathcal{I}^\bullet(W))
\to
\text{Tot}(\mathcal{K}^\bullet(W) \otimes_A \mathcal{I}^\bullet(W))
\to
\mathcal{J}^\bullet(W)
$$
；它是 $A$-模复形的态射，其复合表示映射
$$
R\Gamma(X, K) \otimes_A^\mathbf{L} R\Gamma(W, M)
\longrightarrow
R\Gamma(W, K \otimes_{\mathcal{O}_X}^\mathbf{L} M)
$$
，该映射位于 $D(A)$ 中。显然，这些映射与限制映射相容。
现在可以考察下列 $A$-模复形的交换图：
$$
\xymatrix{
0 \ar[d] & 0 \ar[d] \\
\text{Tot}(T^\bullet \otimes_A \mathcal{I}^\bullet(X)) \ar[d] \ar[r] &
\mathcal{J}^\bullet(X) \ar[d] \\
\text{Tot}(T^\bullet \otimes_A
(\mathcal{I}^\bullet(U) \oplus \mathcal{I}^\bullet(V)) \ar[d] \ar[r] &
\mathcal{J}^\bullet(U) \oplus \mathcal{J}^\bullet(V) \ar[d] \\
\text{Tot}(T^\bullet \otimes_A \mathcal{I}^\bullet(U \cap V)) \ar[r] \ar[d] &
\mathcal{J}^\bullet(U \cap V) \ar[d] \\
0 & 0
}
$$
由引理 \ref{cohomology-lemma-mayer-vietoris} 的证明，各列都是
$A$-模复形的正合列（这里还用到：
由于 $T^\bullet$ 的各项是平坦 $A$-模，
$\text{Tot}(T^\bullet \otimes_A -)$
把 $A$-模复形的短正合列变为短正合列）。引理 \ref{cohomology-lemma-unbounded-mayer-vietoris} 的特异三角
正是与这些复形短正合列相伴的特异三角。因此，所需结论由
《导出范畴》第 \ref{derived-section-canonical-delta-functor} 节所讨论的
“取相伴特异三角”这一操作的函子性得到。
\end{proof}












\section{支撑在闭子集上的上同调，II}
\label{cohomology-section-cohomology-support-bis}

\noindent
继续第 \ref{cohomology-section-cohomology-support} 节开始的讨论。

\medskip\noindent
设 $(X, \mathcal{O}_X)$ 是环空间，而 $Z \subset X$ 是闭子集。
在此情形中，可以考察函子
$\textit{Mod}(\mathcal{O}_X) \to \textit{Mod}(\mathcal{O}_X(X))$
，它由 $\mathcal{F} \mapsto \Gamma_Z(X, \mathcal{F})$ 给出；
参见《模》定义 \ref{modules-definition-support} 与引理
\ref{modules-lemma-support-section-closed}。
使用 K-内射消解（见第 \ref{cohomology-section-unbounded} 节），得到右导出函子
$$
R\Gamma_Z(X, - ) : D(\mathcal{O}_X) \to D(\mathcal{O}_X(X))
$$
给定 $D(\mathcal{O}_X)$ 中的对象 $K$，以
$H^q_Z(X, K) = H^q(R\Gamma_Z(X, K))$
表示支撑在 $Z$ 上的上同调模。稍后将看到
（引理 \ref{cohomology-lemma-sections-support-abelian-unbounded}），
它与第 \ref{cohomology-section-cohomology-support} 节的构造一致。

\medskip\noindent
对 $\mathcal{O}_X$-模 $\mathcal{F}$，可以考察{\it 支撑在 $Z$
上的截面子层}，记为 $\mathcal{H}_Z(\mathcal{F})$，其定义为
$$
\mathcal{H}_Z(\mathcal{F})(U) =
\{s \in \mathcal{F}(U) \mid \text{Supp}(s) \subset U \cap Z\} =
\Gamma_{Z \cap U}(U, \mathcal{F}|_U)
$$
如《模》注记
\ref{modules-remark-sections-support-in-closed-modules} 所述，
可把 $\mathcal{H}_Z(\mathcal{F})$ 看作 $Z$ 上的
$\mathcal{O}_X|_Z$-模，并得到函子
$$
\textit{Mod}(\mathcal{O}_X) \longrightarrow \textit{Mod}(\mathcal{O}_X|_Z),
\quad
\mathcal{F} \longmapsto \mathcal{H}_Z(\mathcal{F})
\text{（看作 }\mathcal{O}_X|_Z\text{-模，定义在 }Z\text{ 上）}
$$
此函子左正合，但一般并不正合。与上面完全相同，可得右导出函子
$$
R\mathcal{H}_Z : D(\mathcal{O}_X) \longrightarrow D(\mathcal{O}_X|_Z)
$$
令 $\mathcal{H}^q_Z(K) = H^q(R\mathcal{H}_Z(K))$，于是
$\mathcal{H}^0_Z(\mathcal{F}) = \mathcal{H}_Z(\mathcal{F})$
对任意 $\mathcal{O}_X$-模层 $\mathcal{F}$ 都成立。

\begin{lemma}
\label{cohomology-lemma-cohomology-with-support-sheaf-on-support}
设 $(X, \mathcal{O}_X)$ 是环化空间。设 $i : Z \to X$ 是闭子集的嵌入。
\begin{enumerate}
\item $R\mathcal{H}_Z : D(\mathcal{O}_X) \to D(\mathcal{O}_X|_Z)$
是 $i_* : D(\mathcal{O}_X|_Z) \to D(\mathcal{O}_X)$ 的右伴随。
\item 对 $D(\mathcal{O}_X|_Z)$ 中的 $K$，有
$R\mathcal{H}_Z(i_*K) = K$。
\item 设 $\mathcal{G}$ 是 $Z$ 上的 $\mathcal{O}_X|_Z$-模层。则当
$p > 0$ 时，$\mathcal{H}^p_Z(i_*\mathcal{G}) = 0$。
\end{enumerate}
\end{lemma}

\begin{proof}
函子 $i_*$ 正合，故 $i_* = Ri_* = Li_*$。因此，本引理的 (1) 由
《模层》引理 \ref{modules-lemma-adjoint-section-with-support}
与
《导出范畴》引理 \ref{derived-lemma-derived-adjoint-functors}
推出。设 $K$ 如 (2) 中所述。可用一个由 $\mathcal{O}_X|_Z$-模组成的
K-内射复形 $\mathcal{I}^\bullet$ 表示 $K$。由引理 \ref{cohomology-lemma-K-injective-flat}，
表示 $i_*K$ 的复形 $i_*\mathcal{I}^\bullet$
是由 $\mathcal{O}_X$-模组成的 K-内射复形。因此，
$R\mathcal{H}_Z(i_*K)$ 由
$\mathcal{H}_Z(i_*\mathcal{I}^\bullet) = \mathcal{I}^\bullet$
计算，这就证明了 (2)。(3) 是 (2) 的特例。
\end{proof}

\noindent
设 $(X, \mathcal{O}_X)$ 是环化空间，且 $Z \subset X$ 是闭子集。
支撑包含于 $Z$ 的 $\mathcal{O}_X$-模所成的范畴，是所有
$\mathcal{O}_X$-模所成范畴的一个 Serre 子范畴；参见
《同调代数》定义 \ref{homology-definition-serre-subcategory}
以及
《模层》引理 \ref{modules-lemma-support-section-closed}。
记 $D_Z(\mathcal{O}_X)$ 为 $D(\mathcal{O}_X)$ 的严格满的饱和三角子范畴，
其对象为上同调层支撑在 $Z$ 上的复形；参见
《导出范畴》第 \ref{derived-section-triangulated-sub} 节。

\begin{lemma}
\label{cohomology-lemma-complexes-with-support-on-closed}
设 $(X, \mathcal{O}_X)$ 是环化空间。设 $i : Z \to X$ 是闭子集的嵌入。
\begin{enumerate}
\item 对 $D(\mathcal{O}_X|_Z)$ 中的 $K$，有
$i_*K \in D_Z(\mathcal{O}_X)$。
\item 函子 $i_* : D(\mathcal{O}_X|_Z) \to D_Z(\mathcal{O}_X)$
是等价，其拟逆为
$i^{-1}|_{D_Z(\mathcal{O}_X)} = R\mathcal{H}_Z|_{D_Z(\mathcal{O}_X)}$。
\item 函子
$i_* \circ R\mathcal{H}_Z : D(\mathcal{O}_X) \to D_Z(\mathcal{O}_X)$
是嵌入函子
$D_Z(\mathcal{O}_X) \to D(\mathcal{O}_X)$ 的右伴随。
\end{enumerate}
\end{lemma}

\begin{proof}
(1) 由定义立即得到。(3) 是 (2) 与引理 \ref{cohomology-lemma-cohomology-with-support-sheaf-on-support}
的形式推论。以下证明 (2)。

\medskip\noindent
将 $i$ 看作环化空间的态射
$i : (Z, \mathcal{O}_X|_Z) \to (X, \mathcal{O}_X)$。
回忆 $i^*$ 与 $i_*$ 构成一对伴随函子。由于 $i$ 是闭浸入，
$i_*$ 正合。又因 $i^{-1}\mathcal{O}_X = \mathcal{O}_X|_Z$
是 $(Z, \mathcal{O}_X|_Z)$ 的结构层，可见 $i^* = i^{-1}$ 正合，
且 $i^*i_* = i^{-1}i_*$ 同构于恒等函子。参见
《模层》引理 \ref{modules-lemma-exactness-pushforward-pullback} 和
\ref{modules-lemma-i-star-exact}。因此
$i_* : D(\mathcal{O}_X|_Z) \to D_Z(\mathcal{O}_X)$
是全忠实的，而 $i^{-1}$ 给出其左逆。另一方面，设 $K$ 是
$D_Z(\mathcal{O}_X)$ 的对象，并考虑伴随态射
$K \to i_*i^{-1}K$。
由 $i_*$ 与 $i^{-1}$ 的正合性，它在上同调层上诱导伴随态射
$H^n(K) \to i_*i^{-1}H^n(K)$。由于这些上同调层支撑在 $Z$ 上，
这些伴随态射都是同构，从而
$i_* : D(\mathcal{O}_X|_Z) \to D_Z(\mathcal{O}_X)$ 是等价。

\medskip\noindent
最后，只需证明：若 $K$ 是 $D_Z(\mathcal{O}_X)$ 的对象，则
$R\mathcal{H}_Z(K) = i^{-1}K$。为此，可使用刚刚证明的
$K = i_*i^{-1}K$。于是引理 \ref{cohomology-lemma-cohomology-with-support-sheaf-on-support}
给出所需结论。
\end{proof}

\begin{lemma}
\label{cohomology-lemma-sections-with-support-K-injective}
设 $(X, \mathcal{O}_X)$ 是环化空间。设 $i : Z \to X$ 是闭子集的嵌入。
若 $\mathcal{I}^\bullet$ 是由 $\mathcal{O}_X$-模组成的 K-内射复形，
则 $\mathcal{H}_Z(\mathcal{I}^\bullet)$
是由 $\mathcal{O}_X|_Z$-模组成的 K-内射复形。
\end{lemma}

\begin{proof}
由于 $i_* : \textit{Mod}(\mathcal{O}_X|_Z) \to
\textit{Mod}(\mathcal{O}_X)$ 正合，并且是 $\mathcal{H}_Z$ 的左伴随
（《模层》引理 \ref{modules-lemma-adjoint-section-with-support}），
结论由
《导出范畴》引理 \ref{derived-lemma-adjoint-preserve-K-injectives}
推出。
\end{proof}

\begin{lemma}
\label{cohomology-lemma-local-to-global-sections-with-support}
设 $(X, \mathcal{O}_X)$ 是环化空间。设 $i : Z \to X$ 是闭子集的嵌入。则
$R\Gamma(Z, - ) \circ R\mathcal{H}_Z = R\Gamma_Z(X, - )$
作为从 $D(\mathcal{O}_X)$ 到 $D(\mathcal{O}_X(X))$ 的函子成立。
\end{lemma}

\begin{proof}
这由利用 K-内射消解构造右导出函子、引理 \ref{cohomology-lemma-sections-with-support-K-injective}
以及等式 $\Gamma_Z(X, -) = \Gamma(Z, -) \circ \mathcal{H}_Z$ 得到。
\end{proof}

\begin{lemma}
\label{cohomology-lemma-triangle-sections-with-support}
设 $(X, \mathcal{O}_X)$ 是环化空间。设 $i : Z \to X$ 是闭子集的嵌入。
令 $U = X \setminus Z$。存在特异三角
$$
R\Gamma_Z(X, K) \to R\Gamma(X, K) \to R\Gamma(U, K) \to
R\Gamma_Z(X, K)[1]
$$
位于 $D(\mathcal{O}_X(X))$ 中，并关于 $D(\mathcal{O}_X)$ 中的
$K$ 具有函子性。
\end{lemma}

\begin{proof}
选择一个表示 $K$ 的 K-内射复形 $\mathcal{I}^\bullet$，其每一项都是
内射 $\mathcal{O}_X$-模；参见 第 \ref{cohomology-section-unbounded} 节。
回忆 $\mathcal{I}^\bullet|_U$ 是由 $\mathcal{O}_U$-模组成的
K-内射复形；参见引理 \ref{cohomology-lemma-restrict-K-injective-to-open}。
因此，特异三角中的每个导出函子，都可通过将相应原函子作用于
$\mathcal{I}^\bullet$ 来得到。故只需证明：对任意内射
$\mathcal{O}_X$-模 $\mathcal{I}$，都有短正合列
$$
0 \to \Gamma_Z(X, \mathcal{I}) \to \Gamma(X, \mathcal{I})
\to \Gamma(U, \mathcal{I}) \to 0
$$
这由引理 \ref{cohomology-lemma-injective-restriction-surjective} 和定义得到。
\end{proof}

\begin{lemma}
\label{cohomology-lemma-triangle-sections-with-support-sheaves}
设 $(X, \mathcal{O}_X)$ 是环化空间。设 $i : Z \to X$ 是闭子集的嵌入。
记 $j : U = X \setminus Z \to X$ 为其补集的嵌入。存在特异三角
$$
i_*R\mathcal{H}_Z(K) \to K \to Rj_*(K|_U) \to
i_*R\mathcal{H}_Z(K)[1]
$$
位于 $D(\mathcal{O}_X)$ 中，并关于 $D(\mathcal{O}_X)$ 中的
$K$ 具有函子性。
\end{lemma}

\begin{proof}
选择一个表示 $K$ 的 K-内射复形 $\mathcal{I}^\bullet$，其每一项都是
内射 $\mathcal{O}_X$-模；参见 第 \ref{cohomology-section-unbounded} 节。
回忆 $\mathcal{I}^\bullet|_U$ 是由 $\mathcal{O}_U$-模组成的
K-内射复形；参见引理 \ref{cohomology-lemma-restrict-K-injective-to-open}。
因此，特异三角中的每个导出函子，都可通过将相应原函子作用于
$\mathcal{I}^\bullet$ 来得到。故只需证明：对任意内射
$\mathcal{O}_X$-模 $\mathcal{I}$，都有短正合列
$$
0 \to i_*\mathcal{H}_Z(\mathcal{I}) \to \mathcal{I}
\to j_*(\mathcal{I}|_U) \to 0
$$
这由引理 \ref{cohomology-lemma-injective-restriction-surjective} 和定义得到。
\end{proof}

\begin{lemma}
\label{cohomology-lemma-sections-support-in-closed-disjoint-open}
设 $(X, \mathcal{O}_X)$ 是环化空间。设 $Z \subset X$ 是闭子集。
设 $j : U \to X$ 是满足 $U \cap Z = \emptyset$ 的开子集的嵌入。
则对 $D(\mathcal{O}_U)$ 中的所有 $K$，有
$R\mathcal{H}_Z(Rj_*K) = 0$。
\end{lemma}

\begin{proof}
选择一个由 $\mathcal{O}_U$-模组成且表示 $K$ 的 K-内射复形
$\mathcal{I}^\bullet$。于是 $j_*\mathcal{I}^\bullet$ 表示 $Rj_*K$。
由引理 \ref{cohomology-lemma-K-injective-flat}，复形
$j_*\mathcal{I}^\bullet$ 是由 $\mathcal{O}_X$-模组成的
K-内射复形。因此，$\mathcal{H}_Z(j_*\mathcal{I}^\bullet)$
表示 $R\mathcal{H}_Z(Rj_*K)$。故只需证明：对 $U$ 上任意阿贝尔层
$\mathcal{G}$，都有 $\mathcal{H}_Z(j_*\mathcal{G}) = 0$。
换言之，必须证明：若 $j_*\mathcal{G}$ 在某个开集 $W$ 上的截面
$s$ 支撑在 $W \cap Z$ 上，则 $s$ 为零。支撑条件意味着
$s|_{W \setminus W \cap Z} = 0$。由于
$j_*\mathcal{G}(W) =
\mathcal{G}(U \cap W) = j_*\mathcal{G}(W \setminus W \cap Z)$
，故 $s$ 确如所需为零。
\end{proof}

\begin{lemma}
\label{cohomology-lemma-sections-support-abelian-unbounded}
设 $(X, \mathcal{O}_X)$ 是环化空间。设 $Z \subset X$ 是闭子集。
设 $K$ 是 $D(\mathcal{O}_X)$ 的对象，并记 $K_{ab}$ 为其在
$D(\underline{\mathbf{Z}}_X)$ 中的像。
\begin{enumerate}
\item 存在典范态射 $R\Gamma_Z(X, K) \to R\Gamma_Z(X, K_{ab})$，
它是 $D(\textit{Ab})$ 中的同构。
\item 存在典范态射
$R\mathcal{H}_Z(K) \to R\mathcal{H}_Z(K_{ab})$
，它是 $D(\underline{\mathbf{Z}}_Z)$ 中的同构。
\end{enumerate}
\end{lemma}

\begin{proof}
证明 (1)。此态射构造如下。选择一个由 $\mathcal{O}_X$-模组成、表示
$K$ 的 K-内射复形 $\mathcal{I}^\bullet$。选择拟同构
$\mathcal{I}^\bullet \to \mathcal{J}^\bullet$，其中
$\mathcal{J}^\bullet$ 是由阿贝尔群组成的 K-内射复形。则 (1) 中的态射为
$$
\Gamma_Z(X, \mathcal{I}^\bullet) \to \Gamma_Z(X, \mathcal{J}^\bullet)
$$
；它由 $\Gamma_Z$ 是阿贝尔层上的函子这一事实确定。容易验证，
所得态射与引理 \ref{cohomology-lemma-modules-abelian-unbounded} 的典范态射
共同构成如下特异三角的态射：
$$
\xymatrix{
R\Gamma_Z(X, K) \ar[r] \ar[d] &
R\Gamma(X, K) \ar[r] \ar[d] &
R\Gamma(U, K) \ar[d] \\
R\Gamma_Z(X, K_{ab}) \ar[r] &
R\Gamma(X, K_{ab}) \ar[r] &
R\Gamma(U, K_{ab})
}
$$
这里使用引理 \ref{cohomology-lemma-triangle-sections-with-support}。
由所引引理，三个箭头中的两个是同构，故可由 《导出范畴》引理
\ref{derived-lemma-third-isomorphism-triangle}.
得出结论。

\medskip\noindent
省略 (2) 的证明。提示：将同一论证用于引理 \ref{cohomology-lemma-triangle-sections-with-support-sheaves}
的特异三角。
\end{proof}

\begin{remark}
\label{cohomology-remark-support-cup-product}
设 $(X, \mathcal{O}_X)$ 是环化空间。设 $i : Z \to X$ 是闭子集的嵌入。
给定 $D(\mathcal{O}_X)$ 中的 $K$ 与 $M$，存在典范态射
$$
K|_Z \otimes_{\mathcal{O}_X|_Z}^\mathbf{L} R\mathcal{H}_Z(M)
\longrightarrow
R\mathcal{H}_Z(K \otimes_{\mathcal{O}_X}^\mathbf{L} M)
$$
位于 $D(\mathcal{O}_X|_Z)$ 中。这里 $K|_Z = i^{-1}K$ 是 $K$ 在
$Z$ 上的限制，并看作 $D(\mathcal{O}_X|_Z)$ 的对象。由引理 \ref{cohomology-lemma-cohomology-with-support-sheaf-on-support}
中的 $i_*$ 与 $R\mathcal{H}_Z$ 的伴随性，为构造此态射，只需给出典范态射
$$
i_*\left(K|_Z \otimes_{\mathcal{O}_X|_Z}^\mathbf{L} R\mathcal{H}_Z(M)\right)
\longrightarrow
K \otimes_{\mathcal{O}_X}^\mathbf{L} M
$$
为构造此态射，选择一个表示 $M$、由 $\mathcal{O}_X$-模组成的
K-内射复形 $\mathcal{I}^\bullet$，以及一个表示 $K$、由
$\mathcal{O}_X$-模组成的 K-平坦复形 $\mathcal{K}^\bullet$。
注意，$\mathcal{K}^\bullet|_Z$ 是由 $\mathcal{O}_X|_Z$-模组成、
表示 $K|_Z$ 的 K-平坦复形；参见引理 \ref{cohomology-lemma-pullback-K-flat}。因此，需要构造复形态射
$$
i_*\text{Tot}\left(
\mathcal{K}^\bullet|_Z \otimes_{\mathcal{O}_X|_Z}
\mathcal{H}_Z(\mathcal{I}^\bullet)\right)
\longrightarrow
\text{Tot}(\mathcal{K}^\bullet \otimes_{\mathcal{O}_X} \mathcal{I}^\bullet)
$$
，它们由 $\mathcal{O}_X$-模组成。为此，只需给出态射
$$
i_*(\mathcal{K}^a|_Z \otimes_{\mathcal{O}_X|_Z}
\mathcal{H}_Z(\mathcal{I}^b))
\longrightarrow
\mathcal{K}^a \otimes_{\mathcal{O}_X} \mathcal{I}^b
$$
例如考察茎，可见此式左端等于
$\mathcal{K}^a \otimes_{\mathcal{O}_X} i_*\mathcal{H}_Z(\mathcal{I}^b)$
，再利用嵌入 $\mathcal{H}_Z(\mathcal{I}^b) \to \mathcal{I}^b$
即可得到所需态射。
\end{remark}

\begin{remark}
\label{cohomology-remark-support-cup-product-global}
采用注记 \ref{cohomology-remark-support-cup-product} 中的记号，可得到典范杯积
\begin{align*}
H^a(X, K) \times H^b_Z(X, M)
& =
H^a(X, K) \times H^b(Z, R\mathcal{H}_Z(M)) \\
& \to
H^a(Z, K|_Z) \times H^b(Z, R\mathcal{H}_Z(M)) \\
& \to
H^{a + b}(Z, K|_Z \otimes_{\mathcal{O}_X|_Z}^\mathbf{L} R\mathcal{H}_Z(M)) \\
& \to
H^{a + b}(Z, R\mathcal{H}_Z(K \otimes_{\mathcal{O}_X}^\mathbf{L} M)) \\
& =
H^{a + b}_Z(X, K \otimes_{\mathcal{O}_X}^\mathbf{L} M)
\end{align*}
这里的等号由引理 \ref{cohomology-lemma-local-to-global-sections-with-support}
给出；第一个箭头是在 $Z$ 上的限制，第二个箭头是杯积
（第 \ref{cohomology-section-cup-product} 节），第三个箭头则是注记 \ref{cohomology-remark-support-cup-product} 中的态射。
\end{remark}

\begin{lemma}
\label{cohomology-lemma-support-cup-product}
采用注记 \ref{cohomology-remark-support-cup-product} 中的记号，下图交换：
$$
\xymatrix{
H^i(X, K) \times H^j_Z(X, M) \ar[r] \ar[d] &
H^{i + j}_Z(X, K \otimes_{\mathcal{O}_X}^\mathbf{L} M) \ar[d] \\
H^i(X, K) \times H^j(X, M) \ar[r] &
H^{i + j}(X, K \otimes_{\mathcal{O}_X}^\mathbf{L} M)
}
$$
其中上方水平箭头是注记 \ref{cohomology-remark-support-cup-product-global}
中的杯积。
\end{lemma}

\begin{proof}
从略。
\end{proof}

\begin{remark}
\label{cohomology-remark-support-functorial}
设 $f : (X', \mathcal{O}_{X'}) \to (X, \mathcal{O}_X)$ 是环化空间的态射。
设 $Z \subset X$ 是闭子集，并令 $Z' = f^{-1}(Z)$。记
$f|_{Z'} : (Z', \mathcal{O}_{X'}|_{Z'}) \to (Z, \mathcal{O}_X|Z)$
为诱导的环化空间态射。对 $D(\mathcal{O}_X)$ 中任意 $K$，存在典范态射
$$
L(f|_{Z'})^*R\mathcal{H}_Z(K) \longrightarrow R\mathcal{H}_{Z'}(Lf^*K)
$$
位于 $D(\mathcal{O}_{X'}|_{Z'})$ 中。记 $i : Z \to X$ 与
$i' : Z' \to X'$ 为嵌入态射。将引理 \ref{cohomology-lemma-complexes-with-support-on-closed} 的 (2) 用于 $i'$，
则给出上述态射等价于给出态射
$$
i'_* L(f|_{Z'})^* R\mathcal{H}_Z(K)
\longrightarrow
i'_*R\mathcal{H}_{Z'}(Lf^*K)
$$
位于 $D_{Z'}(\mathcal{O}_{X'})$ 中。注记 \ref{cohomology-remark-base-change}
中的函子态射
$Lf^* \circ i_* \to i'_* \circ L(f|_{Z'})^*$
在此情形下是同构（可在茎上检验；这里使用 $i_*$ 与 $i'_*$ 的正合性、
$\mathcal{O}_{Z, z} = \mathcal{O}_{X, z}$，以及 $Z'$ 的类似等式）。
因此，只需构造下图上方的水平箭头
$$
\xymatrix{
Lf^* i_* R\mathcal{H}_Z(K) \ar[rr] \ar[rd] & &
i'_* R\mathcal{H}_{Z'}(Lf^*K) \ar[ld] \\
& Lf^*K
}
$$
复形 $Lf^* i_* R\mathcal{H}_Z(K)$ 支撑在 $Z'$ 上。东南方向的箭头来自
伴随态射 $i_*R\mathcal{H}_Z(K) \to K$
（引理 \ref{cohomology-lemma-cohomology-with-support-sheaf-on-support}）。
由引理 \ref{cohomology-lemma-complexes-with-support-on-closed} 的 (3)，伴随态射
$i'_* R\mathcal{H}_{Z'}(Lf^*K) \to Lf^*K$ 具有泛性质。因此，
东南方向的箭头唯一地经由西南方向的箭头分解，从而得到所需箭头。
\end{remark}

\begin{lemma}
\label{cohomology-lemma-support-functorial}
采用注记 \ref{cohomology-remark-support-functorial} 中的记号与假设，下图交换：
$$
\xymatrix{
H^p_Z(X, K) \ar[r] \ar[d] & H^p_{Z'}(X', Lf^*K) \ar[d] \\
H^p(X, K) \ar[r] & H^p(X', Lf^*K)
}
$$
其中上方水平箭头来自如下等同：
$H^p_Z(X, K) = H^p(Z, R\mathcal{H}_Z(K))$ 与
$H^p_{Z'}(X', Lf^*K) = H^p(Z', R\mathcal{H}_{Z'}(K'))$、拉回态射
$H^p(Z, R\mathcal{H}_Z(K)) \to H^p(Z', L(f|_{Z'})^*R\mathcal{H}_Z(K))$,
以及注记 \ref{cohomology-remark-support-functorial} 中构造的态射。
\end{lemma}

\begin{proof}
从略。提示：利用
$H^p(Z, R\mathcal{H}_Z(K)) = H^p(X, i_*R\mathcal{H}_Z(K))$
以及 $R\mathcal{H}_{Z'}(Lf^*K)$ 的类似等式，结论可由拉回态射的函子性
和定义注记 \ref{cohomology-remark-support-functorial} 中态射时所用的交换图得到。
\end{proof}






\section{逆系统与上同调，I}
\label{cohomology-section-inverse-systems}

\noindent
设 $A$ 是环，且 $I \subset A$ 是理想。我们证明若干关于
$A/I^n$-模层的逆系统的结果。

\begin{lemma}
\label{cohomology-lemma-ML-general}
设 $I$ 是环 $A$ 的理想。设 $X$ 是拓扑空间。设
$$
\ldots \to \mathcal{F}_3 \to \mathcal{F}_2 \to \mathcal{F}_1
$$
是 $X$ 上 $A$-模层的逆系统，并满足
$\mathcal{F}_n = \mathcal{F}_{n + 1}/I^n\mathcal{F}_{n + 1}$。
设 $p \geq 0$。假设
$$
\bigoplus\nolimits_{n \geq 0} H^{p + 1}(X, I^n\mathcal{F}_{n + 1})
$$
作为分次 $\bigoplus_{n \geq 0} I^n/I^{n + 1}$-模满足升链条件。
则逆系统 $M_n = H^p(X, \mathcal{F}_n)$ 满足 Mittag--Leffler
条件\footnote{事实上，存在 $c \geq 0$，使得对所有 $n \geq c$，
$\Im(M_n \to M_{n - c})$ 都是稳定像。}。
\end{lemma}

\begin{proof}
令 $N_n = H^{p + 1}(X, I^n\mathcal{F}_{n + 1})$，并令
$\delta_n : M_n \to N_n$ 为短正合列
$0 \to I^n\mathcal{F}_{n + 1} \to \mathcal{F}_{n + 1}
\to \mathcal{F}_n \to 0$
在上同调上诱导的边界态射。则
$\bigoplus \Im(\delta_n) \subset \bigoplus N_n$ 是分次子模。
事实上，若 $s \in M_n$ 且 $f \in I^m$，则有交换图
$$
\xymatrix{
0 \ar[r] &
I^n\mathcal{F}_{n + 1} \ar[d]_f \ar[r] &
\mathcal{F}_{n + 1} \ar[d]_f \ar[r] &
\mathcal{F}_n \ar[d]_f \ar[r] & 0 \\
0 \ar[r] &
I^{n + m}\mathcal{F}_{n + m + 1} \ar[r] &
\mathcal{F}_{n + m + 1} \ar[r] &
\mathcal{F}_{n + m} \ar[r] & 0
}
$$
中间的竖直态射由如下方式给出：先将 $\mathcal{F}_{n + 1}$ 的局部截面
提升为 $\mathcal{F}_{n + m + 1}$ 的截面，再乘以 $f$；其余竖直箭头亦然。
由此得到 $\delta_{n + m}(fs) = f \delta_n(s)$。根据假设，可找到
$s_j \in M_{n_j}$，$j = 1, \ldots, N$，使得
$\delta_{n_j}(s_j)$ 作为分次模生成
$\bigoplus \Im(\delta_n)$。设 $n > c = \max(n_j)$，并取 $s \in M_n$。
则可找到 $f_j \in I^{n - n_j}$，使得
$\delta_n(s) = \sum f_j \delta_{n_j}(s_j)$。因此
$\delta(s - \sum f_j s_j) = 0$；换言之，可找到 $s' \in M_{n + 1}$，
其在 $M_n$ 中的像为 $s - \sum f_js_j$。于是
$$
\Im(M_{n + 1} \to M_{n - c}) = \Im(M_n \to M_{n - c})
$$
这是因为各元素 $f_js_j$ 在 $M_{n - c}$ 中的像为零。引理得证。
\end{proof}

\begin{lemma}
\label{cohomology-lemma-ML-general-better}
设 $I$ 是环 $A$ 的理想。设 $X$ 是拓扑空间。设
$$
\ldots \to \mathcal{F}_3 \to \mathcal{F}_2 \to \mathcal{F}_1
$$
是 $X$ 上 $A$-模的逆系统，并满足
$\mathcal{F}_n = \mathcal{F}_{n + 1}/I^n\mathcal{F}_{n + 1}$。
设 $p \geq 0$。对给定的 $n$，定义
$$
N_n =
\bigcap\nolimits_{m \geq n}
\Im\left(
H^{p + 1}(X, I^n\mathcal{F}_{m + 1}) \to H^{p + 1}(X, I^n\mathcal{F}_{n + 1})
\right)
$$
若 $\bigoplus N_n$ 作为分次
$\bigoplus_{n \geq 0} I^n/I^{n + 1}$-模满足升链条件，则逆系统
$M_n = H^p(X, \mathcal{F}_n)$ 满足 Mittag--Leffler
条件\footnote{事实上，存在 $c \geq 0$，使得对所有 $n \geq c$，
$\Im(M_n \to M_{n - c})$ 都是稳定像。}。
\end{lemma}

\begin{proof}
证明与引理 \ref{cohomology-lemma-ML-general} 的证明完全相同。事实上，只要证明
$\bigoplus N_n$ 是
$\bigoplus H^{p + 1}(X, I^n\mathcal{F}_{n + 1})$
的分次 $\bigoplus_{n \geq 0} I^n/I^{n + 1}$-子模，并且边界态射
$\delta_n : M_n \to H^{p + 1}(X, I^n\mathcal{F}_{n + 1})$
的像包含于 $N_n$，结论便由其中的论证推出。

\medskip\noindent
设 $\xi \in N_n$ 且 $f \in I^k$。选择 $m \gg n + k$，并选择
$\xi' \in H^{p + 1}(X, I^n\mathcal{F}_{m + 1})$ 作为 $\xi$ 的提升。
考虑交换图
$$
\xymatrix{
0 \ar[r] &
I^n\mathcal{F}_{m + 1} \ar[d]_f \ar[r] &
\mathcal{F}_{m + 1} \ar[d]_f \ar[r] &
\mathcal{F}_n \ar[d]_f \ar[r] & 0 \\
0 \ar[r] &
I^{n + k}\mathcal{F}_{m + 1} \ar[r] &
\mathcal{F}_{m + 1} \ar[r] &
\mathcal{F}_{n + k} \ar[r] & 0
}
$$
其构造如引理 \ref{cohomology-lemma-ML-general} 的证明所示。它在上同调上诱导态射，
并可见 $f \xi' \in H^{p + 1}(X, I^{n + k}\mathcal{F}_{m + 1})$
映到 $f \xi$。由于这对所有 $m \gg n + k$ 都成立，故如所需，
$f\xi$ 属于 $N_{n + k}$。

\medskip\noindent
为证明边界态射 $\delta_n$ 的像包含于 $N_n$，考虑各图
$$
\xymatrix{
0 \ar[r] &
I^n\mathcal{F}_{m + 1} \ar[d] \ar[r] &
\mathcal{F}_{m + 1} \ar[d] \ar[r] &
\mathcal{F}_n \ar[d] \ar[r] & 0 \\
0 \ar[r] &
I^n\mathcal{F}_{n + 1} \ar[r] &
\mathcal{F}_{n + 1} \ar[r] &
\mathcal{F}_n \ar[r] & 0
}
$$
其中 $m \geq n$。考察诱导的上同调态射即可得出结论。
\end{proof}

\begin{lemma}
\label{cohomology-lemma-topology-I-adic-general}
设 $I$ 是环 $A$ 的理想。设 $X$ 是拓扑空间。设
$$
\ldots \to \mathcal{F}_3 \to \mathcal{F}_2 \to \mathcal{F}_1
$$
是 $X$ 上 $A$-模层的逆系统，并满足
$\mathcal{F}_n = \mathcal{F}_{n + 1}/I^n\mathcal{F}_{n + 1}$。
设 $p \geq 0$。假设
$$
\bigoplus\nolimits_{n \geq 0} H^p(X, I^n\mathcal{F}_{n + 1})
$$
作为分次 $\bigoplus_{n \geq 0} I^n/I^{n + 1}$-模满足升链条件。
则 $M = \lim H^p(X, \mathcal{F}_n)$ 上的极限拓扑就是 $I$-进拓扑。
\end{lemma}

\begin{proof}
对 $n \geq 1$，令 $F^n = \Ker(M \to H^p(X, \mathcal{F}_n))$，并令
$F^0 = M$。注意 $I F^n \subset F^{n + 1}$；特别地，
$I^n M \subset F^n$。因此，$I$-进拓扑比极限拓扑更细。反过来，
我们将证明：给定 $n$，存在 $m \geq n$，使得
$F^m \subset I^nM$\footnote{事实上，存在 $c \geq 0$，使得对所有
$n$，都有 $F^{n + c} \subset I^nM$。}。存在单射
$$
F^n/F^{n + 1} \longrightarrow H^p(X, \mathcal{F}_{n + 1})
$$
其像包含于态射
$H^p(X, I^n\mathcal{F}_{n + 1}) \to H^p(X, \mathcal{F}_{n + 1})$
的像中。记
$$
E_n \subset H^p(X, I^n\mathcal{F}_{n + 1})
$$
为 $F^n/F^{n + 1}$ 的逆像。则 $\bigoplus E_n$ 是
$\bigoplus H^p(X, I^n\mathcal{F}_{n + 1})$ 的分次
$\bigoplus I^n/I^{n + 1}$-子模，而
$\bigoplus E_n \to \bigoplus F^n/F^{n + 1}$ 是分次模同态；
细节从略。根据假设，$\bigoplus E_n$ 由有限多个齐次元素作为
$\bigoplus I^n/I^{n + 1}$-模生成。由于
$E_n \to F^n/F^{n + 1}$ 满射，可见
$\bigoplus F^n/F^{n + 1}$ 亦具有相同性质。因此，可找到 $r$、
$c_1, \ldots, c_r \geq 0$ 以及 $a_i \in F^{c_i}$，使其在
$\bigoplus F^n/F^{n + 1}$ 中的像生成该模。令 $c = \max(c_i)$。

\medskip\noindent
对 $n \geq c$，我们断言 $I F^n = F^{n + 1}$。由此断言可得
$F^{n + c} = I^nF^c \subset I^nM$，正合所需。为证明断言，设
$a \in F^{n + 1}$。$a$ 在 $F^{n + 1}/F^{n + 2}$ 中的像是各
$a_i$ 的线性组合。因此，对某些 $f_i \in I^{n + 1 - c_i}$，有
$a  - \sum f_i a_i \in F^{n + 2}$。由于 $n \geq c_i$ 时
$I^{n + 1 - c_i} = I \cdot I^{n - c_i}$，可写
$f_i = \sum g_{i, j} h_{i, j}$，其中 $g_{i, j} \in I$ 且
$h_{i, j}a_i \in F^n$。于是
$F^{n + 1} = F^{n + 2} + IF^n$.
简单归纳即得：对所有 $e > 0$，
$F^{n + 1} = F^{n + e} + IF^n$。因此，$IF^n$ 在
$F^{n + 1}$ 中稠密。选择 $I$ 的生成元
$k_1, \ldots, k_r$，并考虑连续态射
$$
u : (F^n)^{\oplus r} \longrightarrow F^{n + 1},\quad
(x_1, \ldots, x_r) \mapsto \sum k_i x_i
$$
（取极限拓扑）。由上可知，对所有 $m \geq n$，
$(F^m)^{\oplus r}$ 在 $u$ 下的像都在 $F^{m + 1}$ 中稠密。
由开映射引理（《更多代数》引理
\ref{more-algebra-lemma-open-mapping}），$u$ 是开映射，因而是满射。
故当 $n \geq c$ 时，$IF^n = F^{n + 1}$。证明完毕。
\end{proof}

\begin{lemma}
\label{cohomology-lemma-topology-I-adic-general-better}
设 $I$ 是环 $A$ 的理想。设 $X$ 是拓扑空间。设
$$
\ldots \to \mathcal{F}_3 \to \mathcal{F}_2 \to \mathcal{F}_1
$$
是 $X$ 上 $A$-模层的逆系统，并满足
$\mathcal{F}_n = \mathcal{F}_{n + 1}/I^n\mathcal{F}_{n + 1}$。
设 $p \geq 0$。对给定的 $n$，定义
$$
N_n =
\bigcap\nolimits_{m \geq n}
\Im\left(
H^p(X, I^n\mathcal{F}_{m + 1}) \to H^p(X, I^n\mathcal{F}_{n + 1})
\right)
$$
若 $\bigoplus N_n$ 作为分次
$\bigoplus_{n \geq 0} I^n/I^{n + 1}$-模满足升链条件，则
$M = \lim H^p(X, \mathcal{F}_n)$ 上的极限拓扑就是 $I$-进拓扑。
\end{lemma}

\begin{proof}
证明与引理 \ref{cohomology-lemma-topology-I-adic-general} 的证明完全相同。
事实上，只要证明 $\bigoplus N_n$ 是
$\bigoplus H^{p + 1}(X, I^n\mathcal{F}_{n + 1})$
的分次 $\bigoplus_{n \geq 0} I^n/I^{n + 1}$-子模，并且
$F^n/F^{n + 1} \subset H^p(X, \mathcal{F}_{n + 1})$
包含于 $N_n \to H^p(X, \mathcal{F}_{n + 1})$ 的像中，
结论便由其中的论证推出。在引理 \ref{cohomology-lemma-ML-general-better}
的证明中，我们已经看到了关于模结构的陈述。

\medskip\noindent
设 $t \in F^n$。选择元素
$s \in H^p(X, I^n\mathcal{F}_{n + 1})$，使其映到 $t$ 在
$H^p(X, \mathcal{F}_{n + 1})$ 中的像。我们必须证明 $s \in N_n$。
由于 $F^n$ 是态射 $M \to H^p(X, \mathcal{F}_n)$ 的核，故对所有
$m \geq n$，可将 $t$ 映到元素
$t_m \in H^p(X, \mathcal{F}_{m + 1})$，且它在
$H^p(X, \mathcal{F}_n)$ 中的像为零。考虑上同调序列
$$
H^{p - 1}(X, \mathcal{F}_n) \to
H^p(X, I^n\mathcal{F}_{m + 1}) \to
H^p(X, \mathcal{F}_{m + 1}) \to
H^p(X, \mathcal{F}_n)
$$
它由短正合列
$0 \to I^n\mathcal{F}_{m + 1} \to \mathcal{F}_{m + 1}
\to \mathcal{F}_n \to 0$
诱导。
可选择 $s_m \in H^p(X, I^n\mathcal{F}_{m + 1})$，使其映到 $t_m$。
将上述序列与 $m = n$ 时的序列比较，可见 $s_m$ 的像与 $s$
相差一个属于如下态射之像的元素：
$H^{p - 1}(X, \mathcal{F}_n) \to
H^p(X, I^n\mathcal{F}_{n + 1})$。
然而，此态射经由态射
$H^p(X, I^n\mathcal{F}_{m + 1}) \to H^p(X, I^n\mathcal{F}_{n + 1})$
分解，因而可见 $s$ 如所需属于其像。
\end{proof}





\section{逆系统与上同调，II}
\label{cohomology-section-inverse-systems-bis}

\noindent
本节在理想为主理想的情形下，继续
第 \ref{cohomology-section-inverse-systems} 节 中的讨论。

\begin{lemma}
\label{cohomology-lemma-equivalent-f-good}
设 $(X, \mathcal{O}_X)$ 是环化空间。设
$f \in \Gamma(X, \mathcal{O}_X)$。设
$$
\ldots \to \mathcal{F}_3 \to \mathcal{F}_2 \to \mathcal{F}_1
$$
是 $\mathcal{O}_X$-模的逆系统。考虑下列条件：
\begin{enumerate}
\item 对所有 $n \geq 1$，态射
$f : \mathcal{F}_{n + 1} \to \mathcal{F}_{n + 1}$
经由 $\mathcal{F}_{n + 1} \to \mathcal{F}_n$ 分解，并给出短正合列
$0 \to \mathcal{F}_n \to \mathcal{F}_{n + 1} \to \mathcal{F}_1 \to 0$,
\item 对所有 $n \geq 1$，态射
$f^n : \mathcal{F}_{n + 1} \to \mathcal{F}_{n + 1}$
经由 $\mathcal{F}_{n + 1} \to \mathcal{F}_1$ 分解，并给出短正合列
$0 \to \mathcal{F}_1 \to \mathcal{F}_{n + 1} \to \mathcal{F}_n \to 0$
\item 存在 $f$-可除的 $\mathcal{O}_X$-模 $\mathcal{G}$，使得
$\mathcal{F}_n = \mathcal{G}[f^n]$；并且
\item 存在无 $f$-挠的 $\mathcal{O}_X$-模 $\mathcal{F}$，使得
$\mathcal{F}_n = \mathcal{F}/f^n\mathcal{F}$.
\end{enumerate}
则 (4) $\Rightarrow$ (3) $\Leftrightarrow$ (2) $\Leftrightarrow$ (1)。
\end{lemma}

\begin{proof}
省略 (1) 与 (2) 等价的证明，也省略 (3) 推出 (1) 的证明。
给定 (1) 中的 $\mathcal{F}_n$，为证明 (3)，令
$\mathcal{G} = \colim \mathcal{F}_n$，其中各态射
$\mathcal{F}_1 \to \mathcal{F}_2 \to \mathcal{F}_3 \to \ldots$
如 (1) 中所述。态射 $f : \mathcal{G} \to \mathcal{G}$ 满射；
这是因为由 (1) 的短正合列，
$\mathcal{F}_{n + 1} \subset \mathcal{G}$ 的像是
$\mathcal{F}_n \subset \mathcal{G}$。因此，$\mathcal{G}$ 是
$f$-可除的 $\mathcal{O}_X$-模，并且
$\mathcal{F}_n = \mathcal{G}[f^n]$。

\medskip\noindent
设给定 (4) 中的 $\mathcal{F}$。态射
$\mathcal{F}/f^{n + 1}\mathcal{F} \to \mathcal{F}/f^n\mathcal{F}$
总是满射，其核是态射
$\mathcal{F}/f\mathcal{F} \to \mathcal{F}/f^{n + 1}\mathcal{F}$
的像；该态射由乘以 $f^n$ 诱导。为验证 (2)，只需说明态射
$f^n : \mathcal{F} \to \mathcal{F}/f^{n + 1}\mathcal{F}$
的核为 $f\mathcal{F}$。为此，只需证明：给定开集 $U \subset X$
上 $\mathcal{F}$ 的截面 $s,t$，若 $f^ns = f^{n + 1}t$，则
$s = ft$。这很清楚，因为 $\mathcal{F}$ 无 $f$-挠，故
$f : \mathcal{F} \to \mathcal{F}$ 是单射。
\end{proof}

\begin{lemma}
\label{cohomology-lemma-topology-I-adic-f}
设 $X$、$f$、$(\mathcal{F}_n)$ 满足引理 \ref{cohomology-lemma-equivalent-f-good} 的条件 (1)。设 $p \geq 0$，并令
$H^p = \lim H^p(X, \mathcal{F}_n)$.
则对所有 $c \geq 1$，$f^cH^p$ 是
$H^p \to H^p(X, \mathcal{F}_c)$ 的核。因此，$H^p$ 上的极限拓扑
就是 $f$-进拓扑。
\end{lemma}

\begin{proof}
设 $c \geq 1$。显然，$f^c H^p$ 在 $H^p(X, \mathcal{F}_c)$ 中映到零。
若 $\xi = (\xi_n) \in H^p$ 在极限拓扑中足够小，则 $\xi_c = 0$，
因而当 $n \geq c$ 时，$\xi_n$ 在
$H^p(X, \mathcal{F}_c)$ 中映到零。考虑短正合列的逆系统
$$
0 \to \mathcal{F}_{n - c} \xrightarrow{f^c} \mathcal{F}_n \to
\mathcal{F}_c \to 0
$$
以及相应的上同调长正合列的逆系统
$$
H^{p - 1}(X, \mathcal{F}_c) \to
H^p(X, \mathcal{F}_{n - c}) \to
H^p(X, \mathcal{F}_n) \to
H^p(X, \mathcal{F}_c)
$$
由于项 $H^{p - 1}(X, \mathcal{F}_c)$ 与 $n$ 无关，可以选择相容的元素序列
$\xi'_n \in H^p(X, \mathcal{F}_{n - c})$
作为 $\xi_n$ 的提升。令 $\xi' = (\xi'_n)$，便得到
$\xi = f^c \xi'$，正合所需。
\end{proof}

\begin{lemma}
\label{cohomology-lemma-limit-finite}
设 $A$ 是关于主理想 $(f)$ 完备的 Noether 环。设 $X$ 是拓扑空间。设
$$
\ldots \to \mathcal{F}_3 \to \mathcal{F}_2 \to \mathcal{F}_1
$$
是 $X$ 上 $A$-模层的逆系统。假设
\begin{enumerate}
\item $\Gamma(X, \mathcal{F}_1)$ 是有限 $A$-模；
\item $X$、$f$、$(\mathcal{F}_n)$ 满足引理 \ref{cohomology-lemma-equivalent-f-good} 的条件 (1)。
\end{enumerate}
则
$$
M = \lim \Gamma(X, \mathcal{F}_n)
$$
是有限 $A$-模，$f$ 是 $M$ 上的非零因子，而且 $M/fM$ 是
$M$ 在 $\Gamma(X, \mathcal{F}_1)$ 中的像。
\end{lemma}

\begin{proof}
由引理 \ref{cohomology-lemma-topology-I-adic-f}，有
$M/fM \subset H^0(X, \mathcal{F}_1)$。由 (1) 与 $A$ 的 Noether
性质，$M/fM$ 是有限 $A$-模。注意
$\bigcap f^nM = 0$，因为 $f^nM$ 在
$H^0(X, \mathcal{F}_n)$ 中映到零。由
《代数》引理 \ref{algebra-lemma-finite-over-complete-ring}
可得 $M$ 是有限 $A$-模。最后，设
$s = (s_n) \in M = \lim H^0(X, \mathcal{F}_n)$ 满足 $fs = 0$。
由引理 \ref{cohomology-lemma-equivalent-f-good} 的条件 (1)，
$s_{n + 1}$ 属于 $\mathcal{F}_{n + 1} \to \mathcal{F}_n$ 的核。
因此 $s_n = 0$。由于 $n$ 任意，故 $s = 0$。于是 $f$ 是 $M$
上的非零因子。
\end{proof}

\begin{lemma}
\label{cohomology-lemma-ML}
设 $A$ 是环，$f \in A$，且 $X$ 是拓扑空间。设
$$
\ldots \to \mathcal{F}_3 \to \mathcal{F}_2 \to \mathcal{F}_1
$$
是 $X$ 上 $A$-模层的逆系统。设 $p \geq 0$。假设
\begin{enumerate}
\item 或者 $H^{p + 1}(X, \mathcal{F}_1)$ 是有限长度的 $A$-模，
或者 $A$ 是 Noether 环且 $H^{p + 1}(X, \mathcal{F}_1)$ 是有限 $A$-模；
\item $X$、$f$、$(\mathcal{F}_n)$ 满足引理 \ref{cohomology-lemma-equivalent-f-good} 的条件 (1)。
\end{enumerate}
则逆系统 $M_n = H^p(X, \mathcal{F}_n)$ 满足 Mittag--Leffler 条件。
\end{lemma}

\begin{proof}
令 $I = (f)$。我们将使用引理 \ref{cohomology-lemma-ML-general} 的判别法。
注意，对所有 $n \geq 0$，
$f^n : \mathcal{F}_1 \to I^n\mathcal{F}_{n + 1}$ 都是同构。
因此，只需证明
$$
\bigoplus\nolimits_{n \geq 1} H^{p + 1}(X, \mathcal{F}_1) \cdot f^{n + 1}
$$
作为分次 $S = \bigoplus_{n \geq 0} A/(f) \cdot f^n$-模满足升链条件。
若 $A$ 不是 Noether 环，则 $H^{p + 1}(X, \mathcal{F}_1)$ 长度有限，
结论成立。若 $A$ 是 Noether 环，则 $S$ 是 Noether 环；由所假设的
$H^{p + 1}(X, \mathcal{F}_1)$ 的有限性，该模是有限 $S$-模，
故结论亦成立。省略若干细节。
\end{proof}

\begin{lemma}
\label{cohomology-lemma-ML-better}
设 $A$ 是环，$f \in A$，且 $X$ 是拓扑空间。设
$$
\ldots \to \mathcal{F}_3 \to \mathcal{F}_2 \to \mathcal{F}_1
$$
是 $X$ 上 $A$-模层的逆系统。设 $p \geq 0$。假设
\begin{enumerate}
\item 或者存在 $m \geq 1$，使得
$H^{p + 1}(X, \mathcal{F}_m) \to H^{p + 1}(X, \mathcal{F}_1)$
的像是有限长度的 $A$-模；或者 $A$ 是 Noether 环，且各态射
$H^{p + 1}(X, \mathcal{F}_m) \to H^{p + 1}(X, \mathcal{F}_1)$
之像的交是有限 $A$-模；
\item $X$、$f$、$(\mathcal{F}_n)$ 满足引理 \ref{cohomology-lemma-equivalent-f-good} 的条件 (1)。
\end{enumerate}
则逆系统 $M_n = H^p(X, \mathcal{F}_n)$ 满足 Mittag--Leffler 条件。
\end{lemma}

\begin{proof}
令 $I = (f)$。我们将使用引理 \ref{cohomology-lemma-ML-general-better}
中涉及模 $N_n$ 的判别法。当 $m \geq n$ 时，有
$I^n\mathcal{F}_{m + 1} = \mathcal{F}_{m + 1 - n}$。因此
$$
N_n = \bigcap\nolimits_{m \geq 1} \Im\left(
H^{p + 1}(X, \mathcal{F}_m) \to H^{p + 1}(X, \mathcal{F}_1)
\right)
$$
与 $n$ 无关，并且
$\bigoplus N_n = \bigoplus N_1 \cdot f^{n + 1}$。
于是，可完全仿照引理 \ref{cohomology-lemma-ML} 的证明得出结论。
\end{proof}

\begin{remark}
\label{cohomology-remark-compare-ML}
设 $(X, \mathcal{O}_X)$ 是环化空间。设
$f \in \Gamma(X, \mathcal{O}_X)$。设 $\mathcal{F}$ 是
$\mathcal{O}_X$-模。若 $\mathcal{F}$ 无 $f$-挠，则对每个
$p \geq 0$，都有逆系统的短正合列
$$
0 \to \{H^p(X, \mathcal{F})/f^nH^p(X, \mathcal{F})\}
\to \{H^p(X, \mathcal{F}/f^n\mathcal{F})\}
\to \{H^{p + 1}(X, \mathcal{F})[f^n]\} \to 0
$$
由于第一个逆系统满足 Mittag--Leffler 条件（ML），可由此得到三点：
\begin{enumerate}
\item 存在短正合列
$$
0 \to \widehat{H^p(X, \mathcal{F})} \to
\lim H^p(X, \mathcal{F}/f^n\mathcal{F}) \to
T_f(H^{p + 1}(X, \mathcal{F})) \to 0
$$
其中 $\widehat{\ }$ 表示通常的 $f$-进完备化，而 $T_f( - )$
表示
《更多代数》例
\ref{more-algebra-example-spectral-sequence-principal}
中的 $f$-进 Tate 模。
\item 有
$R^1\lim H^p(X, \mathcal{F}/f^n\mathcal{F}) =
R^1\lim H^{p + 1}(X, \mathcal{F})[f^n]$.
\item 系统 $\{H^{p + 1}(X, \mathcal{F})[f^n]\}$ 满足 ML，当且仅当
$\{H^p(X, \mathcal{F}/f^n\mathcal{F})\}$ 满足 ML。
\end{enumerate}
参见 《同调代数》引理 \ref{homology-lemma-Mittag-Leffler} 以及
《更多代数》引理 \ref{more-algebra-lemma-six-term-Rlim} 和
\ref{more-algebra-lemma-Mittag-Leffler}。
\end{remark}










\section{导出极限}
\label{cohomology-section-derived-limits}

\noindent
设 $(X, \mathcal{O}_X)$ 是环化空间。由于三角范畴
$D(\mathcal{O}_X)$ 有积
（《内射对象》引理 \ref{injectives-lemma-derived-products}），
故 $D(\mathcal{O}_X)$ 有导出极限；参见
《导出范畴》定义 \ref{derived-definition-derived-limit}。
若 $(K_n)$ 是 $D(\mathcal{O}_X)$ 中的逆系统，则记
$R\lim K_n$ 为其导出极限。

\begin{lemma}
\label{cohomology-lemma-RGamma-commutes-with-Rlim}
设 $(X, \mathcal{O}_X)$ 是环化空间。对开集 $U \subset X$，函子
$R\Gamma(U, -)$ 与 $R\lim$ 交换。此外，对 $D(\mathcal{O}_X)$
中的任意逆系统 $(K_n)$ 及任意 $m \in \mathbf{Z}$，存在短正合列
$$
0 \to
R^1\lim H^{m - 1}(U, K_n) \to H^m(U, R\lim K_n) \to
\lim H^m(U, K_n) \to 0
$$
\end{lemma}

\begin{proof}
第一个陈述由
《内射对象》引理 \ref{injectives-lemma-RF-commutes-with-Rlim}
推出。然后可将
《更多代数》注记 \ref{more-algebra-remark-compare-derived-limit}
用于 $R\lim R\Gamma(U, K_n) = R\Gamma(U, R\lim K_n)$，
从而得到上述短正合列。
\end{proof}

\begin{lemma}
\label{cohomology-lemma-Rf-commutes-with-Rlim}
设 $f : (X, \mathcal{O}_X) \to (Y, \mathcal{O}_Y)$ 是环化空间的态射。
则 $Rf_*$ 与 $R\lim$ 交换；亦即，$Rf_*$ 与导出极限交换。
\end{lemma}

\begin{proof}
设 $(K_n)$ 是 $D(\mathcal{O}_X)$ 中的逆系统。考虑定义性的特异三角
$$
R\lim K_n \to \prod K_n \to \prod K_n
$$
它位于 $D(\mathcal{O}_X)$ 中。应用正合函子 $Rf_*$，得到特异三角
$$
Rf_*(R\lim K_n) \to Rf_*\left(\prod K_n\right) \to Rf_*\left(\prod K_n\right)
$$
它位于 $D(\mathcal{O}_Y)$ 中。因此，只需证明 $Rf_*$
与导出范畴中的积交换（这些积并非仅由复形的积给出；参见
《内射对象》引理 \ref{injectives-lemma-derived-products}）。
然而，由引理 \ref{cohomology-lemma-adjoint}，$Rf_*$ 是右伴随，故结论形式地成立
（参见 《范畴》引理 \ref{categories-lemma-adjoint-exact}）。
注意：不能直接应用
《范畴》引理 \ref{categories-lemma-adjoint-exact}
，因为 $R\lim K_n$ 并不是 $D(\mathcal{O}_X)$ 中的极限。
\end{proof}

\begin{remark}
\label{cohomology-remark-discuss-derived-limit}
设 $(X, \mathcal{O}_X)$ 是环化空间。设 $(K_n)$ 是
$D(\mathcal{O}_X)$ 中的逆系统。令 $K = R\lim K_n$。对每个 $n$ 与 $m$，
令 $\mathcal{H}^m_n = H^m(K_n)$ 为 $K_n$ 的第 $m$ 个上同调层，并类似地令
$\mathcal{H}^m = H^m(K)$。记 $\underline{\mathcal{H}}^m_n$ 为预层
$$
U \longmapsto \underline{\mathcal{H}}^m_n(U) = H^m(U, K_n)
$$
类似地，令 $\underline{\mathcal{H}}^m(U) = H^m(U, K)$。
由引理 \ref{cohomology-lemma-sheafification-cohomology} 可见，
$\mathcal{H}^m_n$ 是 $\underline{\mathcal{H}}^m_n$ 的层化，而
$\mathcal{H}^m$ 是 $\underline{\mathcal{H}}^m$ 的层化。有图
$$
\xymatrix{
K \ar@{=}[d] &
\underline{\mathcal{H}}^m \ar[d] \ar[r] & 
\mathcal{H}^m \ar[d] \\
R\lim K_n &
\lim \underline{\mathcal{H}}^m_n \ar[r] & 
\lim \mathcal{H}^m_n
}
$$
一般而言，$\lim \mathcal{H}^m_n$ 未必是
$\lim \underline{\mathcal{H}}^m_n$ 的层化。若 $U \subset X$ 是开集，
则存在短正合列
\begin{equation}
\label{cohomology-equation-ses-Rlim-over-U}
0 \to
R^1\lim \underline{\mathcal{H}}^{m - 1}_n(U) \to
\underline{\mathcal{H}}^m(U) \to
\lim \underline{\mathcal{H}}^m_n(U) \to 0
\end{equation}
这由引理 \ref{cohomology-lemma-RGamma-commutes-with-Rlim} 得到。
\end{remark}

\noindent
下述引理适用于概形上转移态射满射的拟凝聚模逆系统。

\begin{lemma}
\label{cohomology-lemma-inverse-limit-is-derived-limit}
设 $(X, \mathcal{O}_X)$ 是环化空间。设 $(\mathcal{F}_n)$ 是
$\mathcal{O}_X$-模的逆系统。设 $\mathcal{B}$ 是 $X$ 的一族开集。假设
\begin{enumerate}
\item $X$ 的每个开集都有一个覆盖，其成员属于 $\mathcal{B}$；
\item 当 $p > 0$ 且 $U \in \mathcal{B}$ 时，
$H^p(U, \mathcal{F}_n) = 0$；
\item 当 $U \in \mathcal{B}$ 时，逆系统 $\mathcal{F}_n(U)$ 的
$R^1\lim$ 为零。
\end{enumerate}
则 $R\lim \mathcal{F}_n = \lim \mathcal{F}_n$，并且当 $p > 0$
且 $U \in \mathcal{B}$ 时，有
$H^p(U, \lim \mathcal{F}_n) = 0$。
\end{lemma}

\begin{proof}
令 $K_n = \mathcal{F}_n$ 且 $K = R\lim \mathcal{F}_n$。采用注记 \ref{cohomology-remark-discuss-derived-limit} 的记号并利用假设 (2)，
可见对 $U \in \mathcal{B}$，当 $m \not = 0$ 时有
$\underline{\mathcal{H}}_n^m(U) = 0$，而
$\underline{\mathcal{H}}_n^0(U) = \mathcal{F}_n(U)$。
由 Equation (\ref{cohomology-equation-ses-Rlim-over-U}) 与假设 (3)，
可见当 $m \not = 0$ 时 $\underline{\mathcal{H}}^m(U) = 0$，
而当 $m = 0$ 时它等于 $\lim \mathcal{F}_n(U)$。利用 (1) 层化，
可得当 $m \not = 0$ 时 $\mathcal{H}^m = 0$，而当 $m = 0$ 时
它等于 $\lim \mathcal{F}_n$。因此 $K = \lim \mathcal{F}_n$。
由于当 $m > 0$ 时
$H^m(U, K) = \underline{\mathcal{H}}^m(U) = 0$（见上），
第二个断言也成立。
\end{proof}

\begin{lemma}
\label{cohomology-lemma-cohomology-derived-limit-injective}
设 $(X, \mathcal{O}_X)$ 是环化空间。设 $(K_n)$ 是
$D(\mathcal{O}_X)$ 中的逆系统。设 $x \in X$ 且 $m \in \mathbf{Z}$。
假设存在整数 $n(x)$ 以及 $x$ 的开邻域基本系 $\mathfrak{U}_x$，
使得对 $U \in \mathfrak{U}_x$，
\begin{enumerate}
\item $R^1\lim H^{m - 1}(U, K_n) = 0$；并且
\item 当 $n \geq n(x)$ 时，
$H^m(U, K_n) \to H^m(U, K_{n(x)})$ 是单射。
\end{enumerate}
则茎上的态射 $H^m(R\lim K_n)_x \to H^m(K_{n(x)})_x$ 是单射。
\end{lemma}

\begin{proof}
设 $\gamma$ 是 $H^m(R\lim K_n)_x$ 中的元素，且它在
$H^m(K_{n(x)})_x$ 中映到零。由于 $H^m(R\lim K_n)$ 是预层
$U \mapsto H^m(U, R\lim K_n)$ 的层化
（由引理 \ref{cohomology-lemma-sheafification-cohomology}），
可选择 $U \in \mathfrak{U}_x$ 以及映到 $\gamma$ 的元素
$\tilde \gamma \in H^m(U, R\lim K_n)$。于是
$\tilde\gamma$ 映到
$\tilde\gamma_{n(x)} \in H^m(U, K_{n(x)})$。
再次利用引理 \ref{cohomology-lemma-sheafification-cohomology}，即
$H^m(K_{n(x)})$ 是预层 $U \mapsto H^m(U, K_{n(x)})$ 的层化，
可在缩小 $U$ 后假设 $\tilde\gamma_{n(x)} = 0$。
对这个 $U$，考虑短正合列
$$
0 \to
R^1\lim H^{m - 1}(U, K_n) \to H^m(U, R\lim K_n) \to
\lim H^m(U, K_n) \to 0
$$
它来自引理 \ref{cohomology-lemma-RGamma-commutes-with-Rlim}。
由假设 (1)，左端群为零；由假设 (2)，右端群单射到
$H^m(U, K_{n(x)})$。因此 $\tilde\gamma = 0$，从而如所需，
$\gamma = 0$。
\end{proof}

\begin{lemma}
\label{cohomology-lemma-is-limit-per-point}
设 $(X, \mathcal{O}_X)$ 是环化空间。设 $E \in D(\mathcal{O}_X)$。
假设对每个 $x \in X$，都存在函数
$p(x, -) : \mathbf{Z} \to \mathbf{Z}$ 以及 $x$ 的开邻域基本系
$\mathfrak{U}_x$，使得
$$
H^p(U, H^{m - p}(E)) = 0 \text{，其中 }
U \in \mathfrak{U}_x \text{ 且 } p > p(x, m)
$$
则
《导出范畴》注记
\ref{derived-remark-map-into-derived-limit-truncations}
中的态射 $E \to R\lim \tau_{\geq -n} E$
是 $D(\mathcal{O}_X)$ 中的同构。
\end{lemma}

\begin{proof}
令 $K_n = \tau_{\geq -n}E$ 且 $K = R\lim K_n$。典范态射
$E \to K$ 来自各典范态射 $E \to K_n = \tau_{\geq -n}E$。
必须证明 $E \to K$ 在上同调层上诱导同构
$H^m(E) \to H^m(K)$。在余下的证明中固定 $m$。若 $n \geq -m$，
则态射 $E \to \tau_{\geq -n}E = K_n$ 诱导同构
$H^m(E) \to H^m(K_n)$。为完成证明，只需说明：对每个 $x \in X$，
都存在整数 $n(x) \geq -m$，使得态射
$H^m(K)_x \to H^m(K_{n(x)})_x$ 是单射。事实上，此时复合
$$
H^m(E)_x \to H^m(K)_x \to H^m(K_{n(x)})_x
$$
是双射，而第二个箭头是单射，故第一个箭头是双射。令
$$
n(x) = 1 + \max\{-m, p(x, m - 1) - m, -1 + p(x, m) - m, -2 + p(x, m + 1) - m\}.
$$
于是总有 $n(x) \geq -m$。断言：态射
$$
H^{m - 1}(U, K_{n + 1}) \to H^{m - 1}(U, K_n)
\quad\text{且}\quad
H^m(U, K_{n + 1}) \to H^m(U, K_n)
$$
当 $n \geq n(x)$ 且 $U \in \mathfrak{U}_x$ 时都是同构。
此断言蕴含引理 \ref{cohomology-lemma-cohomology-derived-limit-injective}
的条件 (1) 与 (2) 成立，从而蕴含所需的单射性。回忆
（《导出范畴》注记
\ref{derived-remark-truncation-distinguished-triangle}）
，存在特异三角
$$
H^{-n - 1}(E)[n + 1] \to
K_{n + 1} \to K_n \to H^{-n - 1}(E)[n + 2]
$$
考察相应的上同调长正合列可知，只要
$$
H^{m + n}(U, H^{-n - 1}(E)),\quad
H^{m + n + 1}(U, H^{-n - 1}(E)),\quad
H^{m + n + 2}(U, H^{-n - 1}(E))
$$
当 $n \geq n(x)$ 且 $U \in \mathfrak{U}_x$ 时均为零，断言便成立。
而这由 $n(x)$ 的选取以及引理中的假设推出。
\end{proof}

\begin{lemma}
\label{cohomology-lemma-is-limit-spaltenstein}
\begin{reference}
\cite[Proposition 3.13]{Spaltenstein}
\end{reference}
设 $(X, \mathcal{O}_X)$ 是环化空间。设 $E \in D(\mathcal{O}_X)$。
假设对每个 $x \in X$，都存在整数 $d_x \geq 0$ 以及 $x$
的开邻域基本系 $\mathfrak{U}_x$，使得
$$
H^p(U, H^q(E)) = 0 \text{，其中 }
U \in \mathfrak{U}_x,\ p > d_x, \text{ 且 }q < 0
$$
则
《导出范畴》注记
\ref{derived-remark-map-into-derived-limit-truncations}
中的态射 $E \to R\lim \tau_{\geq -n} E$
是 $D(\mathcal{O}_X)$ 中的同构。
\end{lemma}

\begin{proof}
将引理 \ref{cohomology-lemma-is-limit-per-point} 用于
$p(x, m) = d_x + \max(0, m)$ 即得。
\end{proof}

\begin{lemma}
\label{cohomology-lemma-is-limit}
设 $(X, \mathcal{O}_X)$ 是环化空间。设 $E \in D(\mathcal{O}_X)$。
假设存在函数 $p(-) : \mathbf{Z} \to \mathbf{Z}$ 以及 $X$
的一族开集 $\mathcal{B}$，使得
\begin{enumerate}
\item $X$ 的每个开集都有一个覆盖，其成员属于 $\mathcal{B}$；并且
\item 当 $p > p(m)$ 且 $U \in \mathcal{B}$ 时，
$H^p(U, H^{m - p}(E)) = 0$。
\end{enumerate}
则
《导出范畴》注记
\ref{derived-remark-map-into-derived-limit-truncations}
中的态射 $E \to R\lim \tau_{\geq -n} E$
是 $D(\mathcal{O}_X)$ 中的同构。
\end{lemma}

\begin{proof}
将引理 \ref{cohomology-lemma-is-limit-per-point} 用于
$p(x, m) = p(m)$ 以及
$\mathfrak{U}_x = \{U \in \mathcal{B} \mid x \in U\}$。
即可。
\end{proof}

\begin{lemma}
\label{cohomology-lemma-is-limit-dimension}
设 $(X, \mathcal{O}_X)$ 是环化空间。设 $E \in D(\mathcal{O}_X)$。
假设存在整数 $d \geq 0$ 以及 $X$ 的拓扑的基 $\mathcal{B}$，使得
$$
H^p(U, H^q(E)) = 0 \text{，其中 }
U \in \mathcal{B},\ p > d, \text{ 且 }q < 0
$$
则
《导出范畴》注记
\ref{derived-remark-map-into-derived-limit-truncations}
中的态射 $E \to R\lim \tau_{\geq -n} E$
是 $D(\mathcal{O}_X)$ 中的同构。
\end{lemma}

\begin{proof}
将引理 \ref{cohomology-lemma-is-limit-spaltenstein} 用于 $d_x = d$ 以及
$\mathfrak{U}_x = \{U \in \mathcal{B} \mid x \in U\}$ 即可。
\end{proof}

\noindent
上述引理可用于在某些情形下计算上同调。

\begin{lemma}
\label{cohomology-lemma-cohomology-over-U-trivial}
设 $(X, \mathcal{O}_X)$ 是环化空间。设 $K$ 是
$D(\mathcal{O}_X)$ 的对象。设 $\mathcal{B}$ 是 $X$ 的一族开集。假设
\begin{enumerate}
\item $X$ 的每个开集都有一个覆盖，其成员属于 $\mathcal{B}$；
\item 对所有 $p > 0$、$q \in \mathbf{Z}$ 以及
$U \in \mathcal{B}$，都有 $H^p(U, H^q(K)) = 0$。
\end{enumerate}
则当 $q \in \mathbf{Z}$ 且 $U \in \mathcal{B}$ 时，
$H^q(U, K) = H^0(U, H^q(K))$。
\end{lemma}

\begin{proof}
由引理 \ref{cohomology-lemma-is-limit-dimension}（取 $d = 0$），注意有
$K = R\lim \tau_{\geq -n} K$。设 $U \in \mathcal{B}$。由
Equation (\ref{cohomology-equation-ses-Rlim-over-U})，得到短正合列
$$
0 \to R^1\lim H^{q - 1}(U, \tau_{\geq -n}K) \to
H^q(U, K) \to \lim H^q(U, \tau_{\geq -n}K) \to 0
$$
条件 (2) 蕴含
$H^q(U, \tau_{\geq -n} K) = H^0(U, H^q(\tau_{\geq -n} K))$
对所有 $q$ 都成立；这里使用例 \ref{cohomology-example-spectral-sequence} 的谱序列。
由于 $\tau_{\geq -n}K$ 下有界，该谱序列收敛。若 $n > -q$，
则 $H^q(\tau_{\geq -n}K) = H^q(K)$。因此，上述短正合列左右两端
的系统最终为常值，其值分别为 $H^0(U, H^{q - 1}(K))$ 与
$H^0(U, H^q(K))$。引理得证。
\end{proof}

\noindent
下面是另一个可以描述导出极限的情形。

\begin{lemma}
\label{cohomology-lemma-derived-limit-suitable-system}
设 $(X, \mathcal{O}_X)$ 是环化空间。设 $(K_n)$ 是
$D(\mathcal{O}_X)$ 中对象的逆系统。设 $\mathcal{B}$ 是
$X$ 的一族开集。假设
\begin{enumerate}
\item $X$ 的每个开集都有一个覆盖，其成员属于 $\mathcal{B}$；
\item 对所有 $U \in \mathcal{B}$ 与所有 $q \in \mathbf{Z}$，均有
\begin{enumerate}
\item 当 $p > 0$ 时，$H^p(U, H^q(K_n)) = 0$；
\item 逆系统 $H^0(U, H^q(K_n))$ 的 $R^1\lim$ 为零。
\end{enumerate}
\end{enumerate}
则对 $q \in \mathbf{Z}$，有
$H^q(R\lim K_n) = \lim H^q(K_n)$。
\end{lemma}

\begin{proof}
令 $K = R\lim K_n$，并取 $U \in \mathcal{B}$。由引理 \ref{cohomology-lemma-cohomology-over-U-trivial} 与 (2)(a)，有
$H^q(U, K_n) = H^0(U, H^q(K_n))$。由引理 \ref{cohomology-lemma-RGamma-commutes-with-Rlim} 与 (2)(b)，有
$H^q(U, K) = \lim H^0(U, H^q(K_n))$。因此，$H^q(U, K)$ 是各层
$H^q(K_n)$ 在 $U$ 上截面之逆极限。由于 $\lim H^q(K_n)$ 是层，
再利用假设 (1)，可见作为预层 $U \mapsto H^q(U, K)$ 的层化的
$H^q(K)$ 等于 $\lim H^q(K_n)$。引理得证。
\end{proof}








\section{构造 K-内射消解}
\label{cohomology-section-K-injective}


\noindent
设 $(X, \mathcal{O}_X)$ 是环化空间。设 $\mathcal{F}^\bullet$
是由 $\mathcal{O}_X$-模组成的复形。范畴
$\textit{Mod}(\mathcal{O}_X)$ 有足够多的内射对象，故可利用
《导出范畴》引理 \ref{derived-lemma-special-inverse-system}
构造图
$$
\xymatrix{
\ldots \ar[r] &
\tau_{\geq -2}\mathcal{F}^\bullet \ar[r] \ar[d] &
\tau_{\geq -1}\mathcal{F}^\bullet \ar[d] \\
\ldots \ar[r] & \mathcal{I}_2^\bullet \ar[r] & \mathcal{I}_1^\bullet
}
$$
它位于 $\mathcal{O}_X$-模复形的范畴中，并满足
\begin{enumerate}
\item 各竖直箭头是拟同构；
\item $\mathcal{I}_n^\bullet$ 是由内射对象组成的下有界复形；
\item 各箭头 $\mathcal{I}_{n + 1}^\bullet \to \mathcal{I}_n^\bullet$
逐项为分裂满射。
\end{enumerate}
由于 $\mathcal{O}_X$-模的范畴有极限（它们在预层层次上计算），
可形成逐项极限
$\mathcal{I}^\bullet = \lim_n \mathcal{I}_n^\bullet$。由
《导出范畴》引理
\ref{derived-lemma-bounded-below-injectives-K-injective} 和
\ref{derived-lemma-limit-K-injectives}
，这是 K-内射复形。一般而言，典范态射
\begin{equation}
\label{cohomology-equation-into-candidate-K-injective}
\mathcal{F}^\bullet \to \mathcal{I}^\bullet
\end{equation}
未必是拟同构。下述引理给出使它成为拟同构的一些条件。

\begin{lemma}
\label{cohomology-lemma-K-injective}
在上述情形中，记
$\mathcal{H}^m = H^m(\mathcal{F}^\bullet)$ 为第 $m$ 个上同调层。
设 $\mathcal{B}$ 是 $X$ 的一族开子集，且 $d \in \mathbf{N}$。假设
\begin{enumerate}
\item $X$ 的每个开集都有一个覆盖，其成员属于 $\mathcal{B}$；
\item 对每个 $U \in \mathcal{B}$，当 $p > d$ 且 $q < 0$ 时，有
$H^p(U, \mathcal{H}^q) = 0$\footnote{只需对所有 $m$ 都存在
$p(m)$，使得当 $p > p(m)$ 时，
$H^p(U. \mathcal{H}^{m - p}) = 0$；参见引理 \ref{cohomology-lemma-is-limit}。}。
\end{enumerate}
则 (\ref{cohomology-equation-into-candidate-K-injective}) 是拟同构。
\end{lemma}

\begin{proof}
由 《导出范畴》引理
\ref{derived-lemma-difficulty-K-injectives}，只需证明态射
$\mathcal{F}^\bullet \to R\lim \tau_{\geq -n} \mathcal{F}^\bullet$
是同构。这正是引理 \ref{cohomology-lemma-is-limit-dimension}。
\end{proof}

\noindent
下面给出一个技术引理，讨论层复形系统之逆极限的上同调层。
从某种意义上说，这并不是应当尝试证明的结论，因为应取导出极限，
而非实际的逆极限。

\begin{lemma}
\label{cohomology-lemma-inverse-limit-complexes}
设 $(X, \mathcal{O}_X)$ 是环化空间。设 $(\mathcal{F}_n^\bullet)$
是 $\mathcal{O}_X$-模复形的逆系统。设 $m \in \mathbf{Z}$。
假设存在 $X$ 的一族开子集 $\mathcal{B}$ 与整数 $n_0$，使得
\begin{enumerate}
\item $X$ 的每个开集都有一个覆盖，其成员属于 $\mathcal{B}$；
\item 对每个 $U \in \mathcal{B}$，
\begin{enumerate}
\item 阿贝尔群系统
$\mathcal{F}_n^{m - 2}(U)$ 与 $\mathcal{F}_n^{m - 1}(U)$
的 $R^1\lim$ 为零（例如，它们满足 Mittag--Leffler 条件）；
\item 阿贝尔群系统 $H^{m - 1}(\mathcal{F}_n^\bullet(U))$
的 $R^1\lim$ 为零（例如，它满足 Mittag--Leffler 条件）；并且
\item 对所有 $n \geq n_0$，有
$H^m(\mathcal{F}_n^\bullet(U)) = H^m(\mathcal{F}_{n_0}^\bullet(U))$
。
\end{enumerate}
\end{enumerate}
则态射
$H^m(\mathcal{F}^\bullet) \to \lim H^m(\mathcal{F}_n^\bullet) \to
H^m(\mathcal{F}_{n_0}^\bullet)$
都是层的同构，其中
$\mathcal{F}^\bullet = \lim \mathcal{F}_n^\bullet$ 是逐项逆极限。
\end{lemma}

\begin{proof}
设 $U \in \mathcal{B}$。注意，
$H^m(\mathcal{F}^\bullet(U))$ 是下列复形从左数第三处的上同调：
$$
\lim_n \mathcal{F}_n^{m - 2}(U) \to
\lim_n \mathcal{F}_n^{m - 1}(U) \to
\lim_n \mathcal{F}_n^m(U) \to
\lim_n \mathcal{F}_n^{m + 1}(U)
$$
由假设 (2)(a) 与 (2)(b)，可应用
《更多代数》引理 \ref{more-algebra-lemma-apply-Mittag-Leffler-again}
得到
$$
H^m(\mathcal{F}^\bullet(U)) = \lim H^m(\mathcal{F}_n^\bullet(U))
$$
由假设 (2)(c)，得到
$$
H^m(\mathcal{F}^\bullet(U)) = H^m(\mathcal{F}_n^\bullet(U))
$$
对所有 $n \geq n_0$ 成立。由假设 (1)，可得预层
$U \mapsto H^m(\mathcal{F}^\bullet(U))$ 的层化，等于预层
$U \mapsto H^m(\mathcal{F}_n^\bullet(U))$ 的层化，
这对所有 $n \geq n_0$ 都成立。因此，层的逆系统
$H^m(\mathcal{F}_n^\bullet)$ 从 $n \geq n_0$ 起为常值，
其值为 $H^m(\mathcal{F}^\bullet)$。引理得证。
\end{proof}












\section{逆系统与上同调，III}
\label{cohomology-section-inverse-systems-tri}

\noindent
本节利用导出极限，继续 第 \ref{cohomology-section-inverse-systems-bis} 节
中的讨论。

\begin{lemma}
\label{cohomology-lemma-formal-functions-principal}
设 $(X, \mathcal{O}_X)$ 是环化空间。设
$A \to \Gamma(X, \mathcal{O}_X)$ 是环同态，且 $f \in A$。
设 $E$ 是 $D(\mathcal{O}_X)$ 的对象。记
$$
E_n = E \otimes_{\mathcal{O}_X} (\mathcal{O}_X \xrightarrow{f^n} \mathcal{O}_X)
$$
并令 $E^\wedge = R\lim E_n$。对 $p \in \mathbf{Z}$，存在典范交换图
$$
\xymatrix{
& 0 & 0 \\
0 \ar[r] &
\widehat{H^p(X, E)} \ar[r] \ar[u] &
\lim H^p(X, E_n) \ar[r] \ar[u] &
T_f(H^{p + 1}(X, E)) \ar[r] &
0 \\
0 \ar[r] &
H^0(H^p(X, E)^\wedge) \ar[r] \ar[u] &
H^p(X, E^\wedge) \ar[r] \ar[u] &
T_f(H^{p + 1}(X, E)) \ar[r] \ar@{=}[u] &
0 \\
&
R^1\lim H^p(X, E)[f^n] \ar[u] \ar[r]^\cong &
R^1\lim H^{p - 1}(X, E_n) \ar[u] \\
& 0 \ar[u] & 0 \ar[u]
}
$$
其各行各列均正合；其中
$\widehat{H^p(X, E)} = \lim H^p(X, E)/f^n H^p(X, E)$
是通常的 $f$-进完备化，$H^p(X, E)^\wedge$ 是导出 $f$-进完备化，
而 $T_f(H^{p + 1}(X, E))$ 是 $f$-进 Tate 模；参见
《更多代数》例
\ref{more-algebra-example-spectral-sequence-principal}。
最后，有 $H^p(X, E^\wedge) = H^p(R\Gamma(X, E)^\wedge)$。
\end{lemma}

\begin{proof}
由引理 \ref{cohomology-lemma-Rf-commutes-with-Rlim}，注意有
$R\Gamma(X, E^\wedge) = R\lim R\Gamma(X, E_n)$。另一方面，
$$
R\Gamma(X, E_n) =
R\Gamma(X, E) \otimes_A^\mathbf{L} (A \xrightarrow{f^n} A)
$$
（细节从略）。可见 $R\Gamma(X, E^\wedge)$ 是导出 $f$-进完备化
$R\Gamma(X, E)^\wedge$。上述图因而由
《更多代数》引理
\ref{more-algebra-lemma-compare-completion-derived-completion}
得到。
\end{proof}

\begin{lemma}
\label{cohomology-lemma-stabilizes}
设 $\mathcal{A}$ 是阿贝尔范畴。设 $f : M \to M$ 是
$\mathcal{A}$ 中的态射。若
$M[f^n] = \Ker(f^n : M \to M)$ 稳定，则逆系统
$$
(M \xrightarrow{f^n} M)
\quad\text{与}\quad
\Coker(f^n : M \to M)
$$
在 $D(\mathcal{A})$ 中 pro-同构。
\end{lemma}

\begin{proof}
显然存在从第一个逆系统到第二个逆系统的态射。设
$M[f^c] = M[f^{c + 1}] = M[f^{c + 2}] = \ldots$.
于是，可借助下列各图，在 $D(\mathcal{A})$ 中定义反方向的逆系统态射：
$$
\xymatrix{
M/M[f^c] \ar[r]_-{f^{n + c}} \ar[d]_{f^c} &
M \ar[d]^1 \\
M \ar[r]^{f^n} & M
}
$$
由于上方水平箭头是单射，上行的复形拟同构于
$\Coker(f^{n + c} : M \to M)$。省略若干细节。
\end{proof}

\begin{example}
\label{cohomology-example-formal-functions-principal}
设 $(X, \mathcal{O}_X)$ 是环化空间。设
$A \to \Gamma(X, \mathcal{O}_X)$ 是环同态，且 $f \in A$。
设 $\mathcal{F}$ 是 $\mathcal{O}_X$-模。假设存在 $c$，使得对所有
$n \geq c$，都有
$\mathcal{F}[f^c] = \mathcal{F}[f^n]$。我们将把引理 \ref{cohomology-lemma-formal-functions-principal} 用于
$E = \mathcal{F}$。由引理 \ref{cohomology-lemma-stabilizes} 可见，
逆系统 $(E_n)$ 与逆系统 $(\mathcal{F}/f^n\mathcal{F})$ pro-同构。
因此，对 $p \in \mathbf{Z}$，得到交换图
$$
\xymatrix{
& 0 & 0 \\
0 \ar[r] &
\widehat{H^p(X, \mathcal{F})} \ar[r] \ar[u] &
\lim H^p(X, \mathcal{F}/f^n\mathcal{F}) \ar[r] \ar[u] &
T_f(H^{p + 1}(X, \mathcal{F})) \ar[r] &
0 \\
0 \ar[r] &
H^0(H^p(X, \mathcal{F})^\wedge) \ar[r] \ar[u] &
H^p(R\Gamma(X, \mathcal{F})^\wedge) \ar[r] \ar[u] &
T_f(H^{p + 1}(X, \mathcal{F})) \ar[r] \ar@{=}[u] &
0 \\
&
R^1\lim H^p(X, \mathcal{F})[f^n] \ar[u] \ar[r]^\cong &
R^1\lim H^{p - 1}(X, \mathcal{F}/f^n\mathcal{F}) \ar[u] \\
& 0 \ar[u] & 0 \ar[u]
}
$$
其各行各列均正合，其中
$\widehat{H^p(X, \mathcal{F})} =
\lim H^p(X, \mathcal{F})/f^n H^p(X, \mathcal{F})$
是通常的 $f$-进完备化，而对 $D(A)$ 中的 $M$，
$M^\wedge$ 表示导出 $f$-进完备化。
\end{example}
















\section{无界复形的 {\v C}ech 上同调}
\label{cohomology-section-cech-cohomology-of-unbounded-complexes}

\noindent
对无界复形而言，第 \ref{cohomology-section-cech-cohomology-of-complexes} 节
的构造并不“正确”。问题在于，Stacks 项目对双复形作全复形化时使用直和，
而这里必须将其替换为积。本节不这样处理，而是假设覆盖有限，并使用交替
{\v C}ech 复形。

\medskip\noindent
设 $(X, \mathcal{O}_X)$ 是环化空间。设 ${\mathcal F}^\bullet$
是 $\mathcal{O}_X$-模预层的复形。设
${\mathcal U} : X = \bigcup_{i \in I} U_i$ 是 $X$ 的一个
{\bf 有限}开覆盖。由于交替 {\v C}ech 复形
$\check{\mathcal{C}}_{alt}^\bullet(\mathcal{U}, \mathcal{F})$
（第 \ref{cohomology-section-alternating-cech} 节）
关于预层 $\mathcal{F}$ 具有函子性，可得到双复形
$\check{\mathcal{C}}^\bullet_{alt}(\mathcal{U}, \mathcal{F}^\bullet)$。
本节使用相应的全复形。构造
$\text{Tot}(\check{\mathcal{C}}^\bullet_{alt}({\mathcal U},
{\mathcal F}^\bullet))$
关于 ${\mathcal F}^\bullet$ 具有函子性。此外，还有函子变换
\begin{equation}
\label{cohomology-equation-global-sections-to-alternating-cech}
\Gamma(X, {\mathcal F}^\bullet)
\longrightarrow
\text{Tot}(\check{\mathcal{C}}^\bullet_{alt}({\mathcal U},
{\mathcal F}^\bullet))
\end{equation}
，它是复形之间的态射，并由下述规则定义：截面
$s\in \Gamma(X, {\mathcal F}^n)$ 映到元素
$\alpha = \{\alpha_{i_0\ldots i_p}\}$，其中
$\alpha_{i_0} = s|_{U_{i_0}}$，且当 $p > 0$ 时
$\alpha_{i_0\ldots i_p} = 0$。

\begin{lemma}
\label{cohomology-lemma-alternating-cech-complex-complex}
设 $(X, \mathcal{O}_X)$ 是环化空间。设
$\mathcal{U} : X = \bigcup_{i \in I} U_i$ 是有限开覆盖。
对由 $\mathcal{O}_X$-模组成的复形 $\mathcal{F}^\bullet$，
存在典范态射
$$
\text{Tot}(\check{\mathcal{C}}^\bullet_{alt}(\mathcal{U}, \mathcal{F}^\bullet))
\longrightarrow
R\Gamma(X, \mathcal{F}^\bullet)
$$
它关于 $\mathcal{F}^\bullet$ 具有函子性，并且与
(\ref{cohomology-equation-global-sections-to-alternating-cech}) 相容。
\end{lemma}

\begin{proof}
设 ${\mathcal I}^\bullet$ 是 K-内射复形，其各项均为内射
$\mathcal{O}_X$-模。对 $\mathcal{I}^\bullet$，
态射 (\ref{cohomology-equation-global-sections-to-alternating-cech}) 为
$\Gamma(X, {\mathcal I}^\bullet) \to
\text{Tot}(\check{\mathcal{C}}^\bullet_{alt}({\mathcal U},
{\mathcal I}^\bullet))$.
这是阿贝尔群复形的拟同构；此结论由
《同调代数》引理 \ref{homology-lemma-double-complex-gives-resolution}
应用于双复形
$\check{\mathcal{C}}^\bullet_{alt}({\mathcal U},
{\mathcal I}^\bullet)$，并利用引理 \ref{cohomology-lemma-injective-trivial-cech} 与 \ref{cohomology-lemma-alternating-usual}
得到。设 ${\mathcal F}^\bullet \to {\mathcal I}^\bullet$
是从 ${\mathcal F}^\bullet$ 到一个各项均内射的 K-内射复形的拟同构
（《内射对象》Theorem
\ref{injectives-theorem-K-injective-embedding-grothendieck}）。
由于 $R\Gamma(X, {\mathcal F}^\bullet)$ 由复形
$\Gamma(X, {\mathcal I}^\bullet)$ 表示，利用
$$
\text{Tot}(\check{\mathcal{C}}^\bullet_{alt}({\mathcal U},
{\mathcal F}^\bullet))
\longrightarrow
\text{Tot}(\check{\mathcal{C}}^\bullet_{alt}({\mathcal U},
{\mathcal I}^\bullet)).
$$
即可得到引理中的态射。省略函子性与相容性的验证。
\end{proof}

\begin{lemma}
\label{cohomology-lemma-alternating-cech-complex-complex-ss}
设 $(X, \mathcal{O}_X)$ 是环化空间。设
$\mathcal{U} : X = \bigcup_{i \in I} U_i$ 是有限开覆盖。
设 $\mathcal{F}^\bullet$ 是由 $\mathcal{O}_X$-模组成的复形。
设 $\mathcal{B}$ 是 $X$ 的一族开子集。假设
\begin{enumerate}
\item $X$ 的每个开集都有一个覆盖，其成员属于 $\mathcal{B}$；
\item 对所有 $i_0, \ldots, i_p \in I$，有
$U_{i_0\ldots i_p} \in \mathcal{B}$；
\item 对每个 $U \in \mathcal{B}$ 与 $p > 0$，都有
\begin{enumerate}
\item $H^p(U, \mathcal{F}^q) = 0$；
\item $H^p(U, \Coker(\mathcal{F}^{q - 1} \to \mathcal{F}^q)) = 0$；并且
\item $H^p(U, H^q(\mathcal{F})) = 0$。
\end{enumerate}
\end{enumerate}
则态射
$$
\text{Tot}(\check{\mathcal{C}}^\bullet_{alt}(\mathcal{U}, \mathcal{F}^\bullet))
\longrightarrow
R\Gamma(X, \mathcal{F}^\bullet)
$$
即引理 \ref{cohomology-lemma-alternating-cech-complex-complex} 中的态射，
是 $D(\textit{Ab})$ 中的同构。
\end{lemma}

\begin{proof}
先假设 $\mathcal{F}^\bullet$ 下有界。在这种情形下，态射
$$
\text{Tot}(\check{\mathcal{C}}^\bullet_{alt}(\mathcal{U}, \mathcal{F}^\bullet))
\longrightarrow
\text{Tot}(\check{\mathcal{C}}^\bullet(\mathcal{U}, \mathcal{F}^\bullet))
$$
由引理 \ref{cohomology-lemma-alternating-usual} 是拟同构。事实上，双复形的态射
$\check{\mathcal{C}}^\bullet_{alt}(\mathcal{U}, \mathcal{F}^\bullet) \to
\check{\mathcal{C}}^\bullet(\mathcal{U}, \mathcal{F}^\bullet)$
在这些复形所对应的第二谱序列的第一页之间诱导同构
（由 《同调代数》引理 \ref{homology-lemma-ss-double-complex}），
而这些谱序列收敛
（《同调代数》引理 \ref{homology-lemma-first-quadrant-ss}）。
因此，这种情形下的结论由引理 \ref{cohomology-lemma-cech-complex-complex-computes} 与假设 (3)(a) 得到。

\medskip\noindent
一般而言，由假设 (3)(c)，可选择如引理 \ref{cohomology-lemma-K-injective}
中的消解
$\mathcal{F}^\bullet \to \mathcal{I}^\bullet = \lim \mathcal{I}_n^\bullet$
。于是，引理中的态射变为
$$
\lim_n 
\text{Tot}(\check{\mathcal{C}}^\bullet_{alt}(\mathcal{U},
\tau_{\geq -n}\mathcal{F}^\bullet))
\longrightarrow
\Gamma(X, \mathcal{I}^\bullet) =
\lim_n \Gamma(X, \mathcal{I}_n^\bullet)
$$
这里箭头位于导出范畴中，而右侧的等式在复形层次上成立。
注意，(3)(b) 表明 $\tau_{\geq -n}\mathcal{F}^\bullet$
是满足本引理假设的下有界复形。因此，由下有界复形的情形可知，
每个态射
$$
\text{Tot}(\check{\mathcal{C}}^\bullet_{alt}(\mathcal{U},
\tau_{\geq -n}\mathcal{F}^\bullet))
\longrightarrow
\Gamma(X, \mathcal{I}_n^\bullet)
$$
都是拟同构。左侧复形在给定次数上的上同调最终为常值
（因为交替 {\v C}ech 复形有限），故右侧亦然。因此，由
《同调代数》引理 \ref{homology-lemma-apply-Mittag-Leffler-again}
（或更强的 《更多代数》引理
\ref{more-algebra-lemma-apply-Mittag-Leffler-again})
，右侧极限的上同调就是这个常值，结论得证。
\end{proof}






\section{Hom 复形}
\label{cohomology-section-hom-complexes}

\noindent
设 $(X, \mathcal{O}_X)$ 是环化空间。设
$\mathcal{L}^\bullet$ 与 $\mathcal{M}^\bullet$ 是两个
$\mathcal{O}_X$-模复形。我们构造 $\mathcal{O}_X$-模复形
$\SheafHom^\bullet(\mathcal{L}^\bullet, \mathcal{M}^\bullet)$。
具体地，对每个 $n$，令
$$
\SheafHom^n(\mathcal{L}^\bullet, \mathcal{M}^\bullet) =
\prod\nolimits_{n = p + q}
\SheafHom_{\mathcal{O}_X}(\mathcal{L}^{-q}, \mathcal{M}^p)
$$
最好将 $\SheafHom^n$ 看作如下 $\mathcal{O}_X$-模层：
从 $\mathcal{L}^\bullet$ 到 $\mathcal{M}^\bullet$
（二者看作分次 $\mathcal{O}_X$-模）的所有次数为 $n$ 的齐次
$\mathcal{O}_X$-线性态射。采用此术语，微分由规则
$$
\text{d}(f) =
\text{d}_\mathcal{M} \circ f - (-1)^n f \circ \text{d}_\mathcal{L}
$$
定义，其中
$f \in \SheafHom^n_{\mathcal{O}_X}(\mathcal{L}^\bullet,
\mathcal{M}^\bullet)$。省略 $\text{d}^2 = 0$ 的验证。
此构造是
《微分分次代数》例 \ref{dga-example-category-complexes}。
的特例。由构造立即可得
\begin{equation}
\label{cohomology-equation-cohomology-hom-complex}
H^n(\Gamma(U, \SheafHom^\bullet(\mathcal{L}^\bullet, \mathcal{M}^\bullet))) =
\Hom_{K(\mathcal{O}_U)}(\mathcal{L}^\bullet, \mathcal{M}^\bullet[n])
\end{equation}
这对所有 $n \in \mathbf{Z}$ 与每个开集 $U \subset X$ 都成立。

\begin{lemma}
\label{cohomology-lemma-compose}
设 $(X, \mathcal{O}_X)$ 是环化空间。给定 $\mathcal{O}_X$-模复形
$\mathcal{K}^\bullet,\mathcal{L}^\bullet,\mathcal{M}^\bullet$，
存在同构
$$
\SheafHom^\bullet(\mathcal{K}^\bullet,
\SheafHom^\bullet(\mathcal{L}^\bullet, \mathcal{M}^\bullet))
=
\SheafHom^\bullet(\text{Tot}(\mathcal{K}^\bullet \otimes_{\mathcal{O}_X}
\mathcal{L}^\bullet), \mathcal{M}^\bullet)
$$
它是 $\mathcal{O}_X$-模复形的同构，并关于
$\mathcal{K}^\bullet,\mathcal{L}^\bullet,\mathcal{M}^\bullet$
具有函子性。
\end{lemma}

\begin{proof}
从略。提示：证明与
《更多代数》引理 \ref{more-algebra-lemma-compose}
完全相同。
\end{proof}

\begin{lemma}
\label{cohomology-lemma-composition}
设 $(X, \mathcal{O}_X)$ 是环化空间。给定 $\mathcal{O}_X$-模复形
$\mathcal{K}^\bullet,\mathcal{L}^\bullet,\mathcal{M}^\bullet$，
存在典范态射
$$
\text{Tot}\left(
\SheafHom^\bullet(\mathcal{L}^\bullet, \mathcal{M}^\bullet)
\otimes_{\mathcal{O}_X}
\SheafHom^\bullet(\mathcal{K}^\bullet, \mathcal{L}^\bullet)
\right)
\longrightarrow
\SheafHom^\bullet(\mathcal{K}^\bullet, \mathcal{M}^\bullet)
$$
，它是 $\mathcal{O}_X$-模复形之间的态射。
\end{lemma}

\begin{proof}
从略。提示：证明与
《更多代数》引理 \ref{more-algebra-lemma-composition}
完全相同。
\end{proof}

\begin{lemma}
\label{cohomology-lemma-diagonal-better}
设 $(X, \mathcal{O}_X)$ 是环化空间。给定 $\mathcal{O}_X$-模复形
$\mathcal{K}^\bullet,\mathcal{L}^\bullet,\mathcal{M}^\bullet$，
存在典范态射
$$
\text{Tot}\left(
\mathcal{K}^\bullet \otimes_{\mathcal{O}_X}
\SheafHom^\bullet(\mathcal{M}^\bullet, \mathcal{L}^\bullet)
\right)
\longrightarrow
\SheafHom^\bullet(\mathcal{M}^\bullet,
\text{Tot}(\mathcal{K}^\bullet \otimes_{\mathcal{O}_X} \mathcal{L}^\bullet))
$$
，它是 $\mathcal{O}_X$-模复形之间的态射，并关于三个复形均具有函子性。
\end{lemma}

\begin{proof}
从略。提示：证明与
《更多代数》引理 \ref{more-algebra-lemma-diagonal-better}
完全相同。
\end{proof}

\begin{lemma}
\label{cohomology-lemma-diagonal}
设 $(X, \mathcal{O}_X)$ 是环化空间。给定 $\mathcal{O}_X$-模复形
$\mathcal{K}^\bullet,\mathcal{L}^\bullet$，存在典范态射
$$
\mathcal{K}^\bullet
\longrightarrow
\SheafHom^\bullet(\mathcal{L}^\bullet,
\text{Tot}(\mathcal{K}^\bullet \otimes_{\mathcal{O}_X} \mathcal{L}^\bullet))
$$
，它是 $\mathcal{O}_X$-模复形之间的态射，并关于两个复形均具有函子性。
\end{lemma}

\begin{proof}
从略。提示：证明与
《更多代数》引理 \ref{more-algebra-lemma-diagonal}
完全相同。
\end{proof}

\begin{lemma}
\label{cohomology-lemma-evaluate-and-more}
设 $(X, \mathcal{O}_X)$ 是环化空间。给定 $\mathcal{O}_X$-模复形
$\mathcal{K}^\bullet,\mathcal{L}^\bullet,\mathcal{M}^\bullet$，
存在典范态射
$$
\text{Tot}(\SheafHom^\bullet(\mathcal{L}^\bullet,
\mathcal{M}^\bullet) \otimes_{\mathcal{O}_X} \mathcal{K}^\bullet)
\longrightarrow
\SheafHom^\bullet(\SheafHom^\bullet(\mathcal{K}^\bullet,
\mathcal{L}^\bullet), \mathcal{M}^\bullet)
$$
，它是 $\mathcal{O}_X$-模复形之间的态射，并关于三个复形均具有函子性。
\end{lemma}

\begin{proof}
从略。提示：证明与
《更多代数》引理 \ref{more-algebra-lemma-evaluate-and-more}
完全相同。
\end{proof}

\begin{lemma}
\label{cohomology-lemma-RHom-into-K-injective}
设 $(X, \mathcal{O}_X)$ 是环化空间。设 $L$ 与 $M$ 是
$D(\mathcal{O}_X)$ 的对象。设 $\mathcal{I}^\bullet$
是表示 $M$ 的 $\mathcal{O}_X$-模 K-内射复形。设
$\mathcal{L}^\bullet$ 是表示 $L$ 的 $\mathcal{O}_X$-模复形。则
$$
H^0(\Gamma(U, \SheafHom^\bullet(\mathcal{L}^\bullet, \mathcal{I}^\bullet))) =
\Hom_{D(\mathcal{O}_U)}(L|_U, M|_U)
$$
对所有开集 $U \subset X$ 成立。
\end{lemma}

\begin{proof}
有
\begin{align*}
H^0(\Gamma(U, \SheafHom^\bullet(\mathcal{L}^\bullet, \mathcal{I}^\bullet)))
& =
\Hom_{K(\mathcal{O}_U)}(\mathcal{L}^\bullet|_U, \mathcal{I}^\bullet|_U) \\
& =
\Hom_{D(\mathcal{O}_U)}(L|_U, M|_U)
\end{align*}
第一个等号是 (\ref{cohomology-equation-cohomology-hom-complex})。
第二个等号成立，是因为由引理 \ref{cohomology-lemma-restrict-K-injective-to-open}，
$\mathcal{I}^\bullet|_U$ 是 K-内射复形。
\end{proof}

\begin{lemma}
\label{cohomology-lemma-RHom-well-defined}
设 $(X, \mathcal{O}_X)$ 是环化空间。设
$(\mathcal{I}')^\bullet \to \mathcal{I}^\bullet$
是 $\mathcal{O}_X$-模 K-内射复形之间的拟同构。设
$(\mathcal{L}')^\bullet \to \mathcal{L}^\bullet$
是 $\mathcal{O}_X$-模复形之间的拟同构。则
$$
\SheafHom^\bullet(\mathcal{L}^\bullet, (\mathcal{I}')^\bullet)
\longrightarrow
\SheafHom^\bullet((\mathcal{L}')^\bullet, \mathcal{I}^\bullet)
$$
是拟同构。
\end{lemma}

\begin{proof}
令 $M$ 为由 $\mathcal{I}^\bullet$ 与
$(\mathcal{I}')^\bullet$ 表示的 $D(\mathcal{O}_X)$ 对象。
令 $L$ 为由 $\mathcal{L}^\bullet$ 与
$(\mathcal{L}')^\bullet$ 表示的 $D(\mathcal{O}_X)$ 对象。
由引理 \ref{cohomology-lemma-RHom-into-K-injective} 可见，各层
$$
H^0(\SheafHom^\bullet(\mathcal{L}^\bullet, (\mathcal{I}')^\bullet))
\quad\text{与}\quad
H^0(\SheafHom^\bullet((\mathcal{L}')^\bullet, \mathcal{I}^\bullet))
$$
都等于预层
$$
U \longmapsto \Hom_{D(\mathcal{O}_U)}(L|_U, M|_U)
$$
所伴随的层。因此，此态射是拟同构。
\end{proof}

\begin{lemma}
\label{cohomology-lemma-RHom-from-K-flat-into-K-injective}
设 $(X, \mathcal{O}_X)$ 是环化空间。设 $\mathcal{I}^\bullet$
是 $\mathcal{O}_X$-模 K-内射复形，并设 $\mathcal{L}^\bullet$
是 $\mathcal{O}_X$-模 K-平坦复形。则
$\SheafHom^\bullet(\mathcal{L}^\bullet, \mathcal{I}^\bullet)$
是 $\mathcal{O}_X$-模 K-内射复形。
\end{lemma}

\begin{proof}
事实上，若 $\mathcal{K}^\bullet$ 是无上同调的
$\mathcal{O}_X$-模复形，则
\begin{align*}
\Hom_{K(\mathcal{O}_X)}(\mathcal{K}^\bullet,
\SheafHom^\bullet(\mathcal{L}^\bullet, \mathcal{I}^\bullet))
& =
H^0(\Gamma(X,
\SheafHom^\bullet(\mathcal{K}^\bullet,
\SheafHom^\bullet(\mathcal{L}^\bullet, \mathcal{I}^\bullet)))) \\
& =
H^0(\Gamma(X, \SheafHom^\bullet(\text{Tot}(
\mathcal{K}^\bullet \otimes_{\mathcal{O}_X} \mathcal{L}^\bullet),
\mathcal{I}^\bullet))) \\
& =
\Hom_{K(\mathcal{O}_X)}(
\text{Tot}(\mathcal{K}^\bullet \otimes_{\mathcal{O}_X} \mathcal{L}^\bullet),
\mathcal{I}^\bullet) \\
& =
0
\end{align*}
第一个等号由 (\ref{cohomology-equation-cohomology-hom-complex}) 得到。
第二个等号由引理 \ref{cohomology-lemma-compose} 得到。
第三个等号由 (\ref{cohomology-equation-cohomology-hom-complex}) 得到。
最后一个等号成立，是因为
$\text{Tot}(\mathcal{K}^\bullet \otimes_{\mathcal{O}_X} \mathcal{L}^\bullet)$
无上同调（因为 $\mathcal{L}^\bullet$ 是 K-平坦的；参见定义 \ref{cohomology-definition-K-flat}），且 $\mathcal{I}^\bullet$
是 K-内射的。
\end{proof}








\section{导出范畴中的内 Hom}
\label{cohomology-section-internal-hom}

\noindent
设 $(X, \mathcal{O}_X)$ 为环化空间，并设 $L,M$ 为
$D(\mathcal{O}_X)$ 的对象。我们希望构造 $D(\mathcal{O}_X)$ 的对象
$R\SheafHom(L, M)$，使得对 $D(\mathcal{O}_X)$ 的任意对象 $K$，
都存在典范双射
\begin{equation}
\label{cohomology-equation-internal-hom}
\Hom_{D(\mathcal{O}_X)}(K, R\SheafHom(L, M))
=
\Hom_{D(\mathcal{O}_X)}(K \otimes_{\mathcal{O}_X}^\mathbf{L} L, M)
\end{equation}
注意，由 Yoneda 引理
（《范畴》引理 \ref{categories-lemma-yoneda}），这个公式在唯一同构
意义下定义了 $R\SheafHom(L, M)$。

\medskip\noindent
为构造这样的对象，选取表示 $M$ 的 K-内射复形
$\mathcal{I}^\bullet$，并任取表示 $L$ 的
$\mathcal{O}_X$-模复形 $\mathcal{L}^\bullet$。定义
$$
R\SheafHom(L, M) = \SheafHom^\bullet(\mathcal{L}^\bullet, \mathcal{I}^\bullet)
$$
其中右端是第 \ref{cohomology-section-hom-complexes} 节构造的
$\mathcal{O}_X$-模复形。由引理 \ref{cohomology-lemma-RHom-well-defined}，
这一定义是良定的。我们得到函子
$$
D(\mathcal{O}_X)^{opp} \times D(\mathcal{O}_X) \longrightarrow D(\mathcal{O}_X),
\quad
(K, L) \longmapsto R\SheafHom(K, L)
$$
在证明式 (\ref{cohomology-equation-internal-hom}) 之前，我们先计算
$R\SheafHom(K,L)$ 的上同调群。

\begin{lemma}
\label{cohomology-lemma-section-RHom-over-U}
设 $(X,\mathcal{O}_X)$ 为环化空间，并设 $L,M$ 为
$D(\mathcal{O}_X)$ 的对象。对每个开集 $U$，有
$$
H^0(U, R\SheafHom(L, M)) =
\Hom_{D(\mathcal{O}_U)}(L|_U, M|_U)
$$
特别地，
$H^0(X,R\SheafHom(L,M))=\Hom_{D(\mathcal{O}_X)}(L,M)$。
\end{lemma}

\begin{proof}
选取表示 $M$ 的 K-内射 $\mathcal{O}_X$-模复形
$\mathcal{I}^\bullet$，以及表示 $L$ 的 K-平坦复形
$\mathcal{L}^\bullet$。于是
$\SheafHom^\bullet(\mathcal{L}^\bullet, \mathcal{I}^\bullet)$
由引理 \ref{cohomology-lemma-RHom-from-K-flat-into-K-injective} 是 K-内射的。
因此，只需在 $U$ 上取截面即可计算其在 $U$ 上的上同调，而结论由引理 \ref{cohomology-lemma-RHom-into-K-injective} 得出。
\end{proof}

\begin{lemma}
\label{cohomology-lemma-internal-hom}
设 $(X,\mathcal{O}_X)$ 为环化空间，并设 $K,L,M$ 为
$D(\mathcal{O}_X)$ 的对象。按上述构造，存在典范同构
$$
R\SheafHom(K, R\SheafHom(L, M)) =
R\SheafHom(K \otimes_{\mathcal{O}_X}^\mathbf{L} L, M)
$$
它位于 $D(\mathcal{O}_X)$ 中，并且关于 $K,L,M$ 具有函子性；
取 $H^0(X,-)$ 即恢复式 (\ref{cohomology-equation-internal-hom})。
\end{lemma}

\begin{proof}
选取表示 $M$ 的 K-内射复形 $\mathcal{I}^\bullet$，以及表示 $L$ 的
K-平坦 $\mathcal{O}_X$-模复形 $\mathcal{L}^\bullet$。任取表示 $K$ 的
$\mathcal{O}_X$-模复形 $\mathcal{K}^\bullet$。于是有
$$
\SheafHom^\bullet(\mathcal{K}^\bullet,
\SheafHom^\bullet(\mathcal{L}^\bullet, \mathcal{I}^\bullet))
=
\SheafHom^\bullet(
\text{Tot}(\mathcal{K}^\bullet \otimes_{\mathcal{O}_X} \mathcal{L}^\bullet),
\mathcal{I}^\bullet)
$$
这是引理 \ref{cohomology-lemma-compose} 的结论。注意，左端表示
$R\SheafHom(K,R\SheafHom(L,M))$（使用引理
\ref{cohomology-lemma-RHom-from-K-flat-into-K-injective}），而右端表示
$R\SheafHom(K \otimes_{\mathcal{O}_X}^\mathbf{L} L, M)$。
这就证明了引理中的显示公式。取整体截面并使用引理 \ref{cohomology-lemma-section-RHom-over-U}，即得式
(\ref{cohomology-equation-internal-hom})。
\end{proof}

\begin{lemma}
\label{cohomology-lemma-restriction-RHom-to-U}
设 $(X,\mathcal{O}_X)$ 为环化空间，并设 $K,L$ 为
$D(\mathcal{O}_X)$ 的对象。$R\SheafHom(K,L)$ 的构造与限制到开集的
操作交换；也就是说，对每个开集 $U$，有
$R\SheafHom(K|_U, L|_U) = R\SheafHom(K, L)|_U$.
\end{lemma}

\begin{proof}
这由构造和引理 \ref{cohomology-lemma-restrict-K-injective-to-open} 立即可得。
\end{proof}

\begin{lemma}
\label{cohomology-lemma-RHom-triangulated}
设 $(X,\mathcal{O}_X)$ 为环化空间。双函子 $R\SheafHom(-,-)$
在两个变元中都把特异三角变为特异三角。
\end{lemma}

\begin{proof}
只需注意，赋值
$$
(\mathcal{L}^\bullet, \mathcal{M}^\bullet) \longmapsto
\SheafHom^\bullet(\mathcal{L}^\bullet, \mathcal{M}^\bullet)
$$
在任一变元中都把逐项分裂的复形短正合列变为逐项分裂的短正合列。
细节从略。
\end{proof}

\begin{lemma}
\label{cohomology-lemma-internal-hom-composition}
设 $(X,\mathcal{O}_X)$ 为环化空间。给定 $D(\mathcal{O}_X)$ 中的
$K,L,M$，存在典范态射
$$
R\SheafHom(L, M) \otimes_{\mathcal{O}_X}^\mathbf{L} R\SheafHom(K, L)
\longrightarrow R\SheafHom(K, M)
$$
它位于 $D(\mathcal{O}_X)$ 中，并且关于 $K,L,M$ 具有函子性。
\end{lemma}

\begin{proof}
选取表示 $M$ 的 K-内射复形 $\mathcal{I}^\bullet$、表示 $L$ 的
K-内射复形 $\mathcal{J}^\bullet$，并任取表示 $K$ 的
$\mathcal{O}_X$-模复形 $\mathcal{K}^\bullet$。由引理
\ref{cohomology-lemma-composition}，存在复形态射
$$
\text{Tot}\left(
\SheafHom^\bullet(\mathcal{J}^\bullet, \mathcal{I}^\bullet)
\otimes_{\mathcal{O}_X}
\SheafHom^\bullet(\mathcal{K}^\bullet, \mathcal{J}^\bullet)
\right)
\longrightarrow
\SheafHom^\bullet(\mathcal{K}^\bullet, \mathcal{I}^\bullet)
$$
这些 $\mathcal{O}_X$-模复形
$\SheafHom^\bullet(\mathcal{J}^\bullet, \mathcal{I}^\bullet)$、
$\SheafHom^\bullet(\mathcal{K}^\bullet, \mathcal{J}^\bullet)$ 与
$\SheafHom^\bullet(\mathcal{K}^\bullet, \mathcal{I}^\bullet)$
分别表示 $R\SheafHom(L,M)$、$R\SheafHom(K,L)$ 和
$R\SheafHom(K,M)$。若选取 K-平坦复形 $\mathcal{H}^\bullet$
以及拟同构
$\mathcal{H}^\bullet \to
\SheafHom^\bullet(\mathcal{K}^\bullet, \mathcal{J}^\bullet)$,
则存在态射
$$
\text{Tot}\left(
\SheafHom^\bullet(\mathcal{J}^\bullet, \mathcal{I}^\bullet)
\otimes_{\mathcal{O}_X} \mathcal{H}^\bullet
\right)
\longrightarrow
\text{Tot}\left(
\SheafHom^\bullet(\mathcal{J}^\bullet, \mathcal{I}^\bullet)
\otimes_{\mathcal{O}_X}
\SheafHom^\bullet(\mathcal{K}^\bullet, \mathcal{J}^\bullet)
\right)
$$
其源表示
$R\SheafHom(L, M) \otimes_{\mathcal{O}_X}^\mathbf{L} R\SheafHom(K, L)$。
将上述两个显示态射复合，即得所需态射。构造具有函子性的证明从略。
\end{proof}

\begin{lemma}
\label{cohomology-lemma-internal-hom-diagonal-better}
设 $(X,\mathcal{O}_X)$ 为环化空间。给定 $D(\mathcal{O}_X)$ 中的
$K,L,M$，存在典范态射
$$
K \otimes_{\mathcal{O}_X}^\mathbf{L} R\SheafHom(M, L)
\longrightarrow
R\SheafHom(M, K \otimes_{\mathcal{O}_X}^\mathbf{L} L)
$$
它位于 $D(\mathcal{O}_X)$ 中，并且关于 $K,L,M$ 具有函子性。
\end{lemma}

\begin{proof}
选取表示 $K$ 的 K-平坦复形 $\mathcal{K}^\bullet$、表示 $L$ 的
K-内射复形 $\mathcal{I}^\bullet$，并任取表示 $M$ 的
$\mathcal{O}_X$-模复形 $\mathcal{M}^\bullet$。选取拟同构
$\text{Tot}(\mathcal{K}^\bullet \otimes_{\mathcal{O}_X} \mathcal{I}^\bullet)
\to \mathcal{J}^\bullet$
其中 $\mathcal{J}^\bullet$ 是 K-内射的。然后使用态射
$$
\text{Tot}\left(
\mathcal{K}^\bullet \otimes_{\mathcal{O}_X}
\SheafHom^\bullet(\mathcal{M}^\bullet, \mathcal{I}^\bullet)
\right)
\to
\SheafHom^\bullet(\mathcal{M}^\bullet,
\text{Tot}(\mathcal{K}^\bullet \otimes_{\mathcal{O}_X} \mathcal{I}^\bullet))
\to
\SheafHom^\bullet(\mathcal{M}^\bullet, \mathcal{J}^\bullet)
$$
其中第一个态射是引理 \ref{cohomology-lemma-diagonal-better} 中的态射。
\end{proof}

\begin{lemma}
\label{cohomology-lemma-internal-hom-diagonal}
设 $(X,\mathcal{O}_X)$ 为环化空间。给定 $D(\mathcal{O}_X)$ 中的
$K,L$，存在典范态射
$$
K \longrightarrow R\SheafHom(L, K \otimes_{\mathcal{O}_X}^\mathbf{L} L)
$$
它位于 $D(\mathcal{O}_X)$ 中，并且关于 $K$ 与 $L$ 均具有函子性。
\end{lemma}

\begin{proof}
选取表示 $K$ 的 K-平坦复形 $\mathcal{K}^\bullet$，并任取表示 $L$ 的
$\mathcal{O}_X$-模复形 $\mathcal{L}^\bullet$。选取 K-内射复形
$\mathcal{J}^\bullet$ 以及拟同构
$\text{Tot}(\mathcal{K}^\bullet \otimes_{\mathcal{O}_X} \mathcal{L}^\bullet)
\to \mathcal{J}^\bullet$。然后使用
$$
\mathcal{K}^\bullet \to
\SheafHom^\bullet(\mathcal{L}^\bullet,
\text{Tot}(\mathcal{K}^\bullet \otimes_{\mathcal{O}_X} \mathcal{L}^\bullet))
\to
\SheafHom^\bullet(\mathcal{L}^\bullet, \mathcal{J}^\bullet)
$$
其中第一个态射来自引理 \ref{cohomology-lemma-diagonal}。
\end{proof}

\begin{lemma}
\label{cohomology-lemma-dual}
设 $(X,\mathcal{O}_X)$ 为环化空间，并设 $L$ 为
$D(\mathcal{O}_X)$ 的对象。置
$L^\vee=R\SheafHom(L,\mathcal{O}_X)$。对 $D(\mathcal{O}_X)$ 中的
$M$，存在典范态射
\begin{equation}
\label{cohomology-equation-eval}
M \otimes^\mathbf{L}_{\mathcal{O}_X} L^\vee
\longrightarrow
R\SheafHom(L, M)
\end{equation}
它诱导典范态射
$$
H^0(X, M \otimes^\mathbf{L}_{\mathcal{O}_X} L^\vee)
\longrightarrow
\Hom_{D(\mathcal{O}_X)}(L, M)
$$
并且关于 $D(\mathcal{O}_X)$ 中的 $M$ 具有函子性。
\end{lemma}

\begin{proof}
利用等同 $M=R\SheafHom(\mathcal{O}_X,M)$，态射
(\ref{cohomology-equation-eval}) 是引理
\ref{cohomology-lemma-internal-hom-composition} 的特殊情形。
\end{proof}

\begin{lemma}
\label{cohomology-lemma-internal-hom-evaluate}
设 $(X,\mathcal{O}_X)$ 为环化空间，并设 $K,L,M$ 为
$D(\mathcal{O}_X)$ 的对象。存在典范态射
$$
R\SheafHom(L, M) \otimes_{\mathcal{O}_X}^\mathbf{L} K
\longrightarrow
R\SheafHom(R\SheafHom(K, L), M)
$$
它位于 $D(\mathcal{O}_X)$ 中，并且关于 $K,L,M$ 具有函子性。
\end{lemma}

\begin{proof}
选取表示 $M$ 的 K-内射复形 $\mathcal{I}^\bullet$、表示 $L$ 的
K-内射复形 $\mathcal{J}^\bullet$，以及表示 $K$ 的 K-平坦复形
$\mathcal{K}^\bullet$。所需态射由引理 \ref{cohomology-lemma-evaluate-and-more}
中的态射
$$
\text{Tot}(\SheafHom^\bullet(\mathcal{J}^\bullet,
\mathcal{I}^\bullet) \otimes_{\mathcal{O}_X} \mathcal{K}^\bullet)
\longrightarrow
\SheafHom^\bullet(\SheafHom^\bullet(\mathcal{K}^\bullet,
\mathcal{J}^\bullet), \mathcal{I}^\bullet)
$$
定义。根据我们对复形的特定选取，左端表示
$R\SheafHom(L, M) \otimes_{\mathcal{O}_X}^\mathbf{L} K$，右端表示
$R\SheafHom(R\SheafHom(K,L),M)$。它关于
$D(\mathcal{O}_X)$ 中三个对象均具有函子性的证明从略。
\end{proof}

\begin{remark}
\label{cohomology-remark-tensor-internal-hom}
设 $(X,\mathcal{O}_X)$ 为环化空间。对 $D(\mathcal{O}_X)$ 中的
$K,K',M,M'$，存在典范态射
$$
R\SheafHom(K, K') \otimes_{\mathcal{O}_X}^\mathbf{L}
R\SheafHom(M, M')
\longrightarrow
R\SheafHom(K \otimes_{\mathcal{O}_X}^\mathbf{L} M,
K' \otimes_{\mathcal{O}_X}^\mathbf{L} M')
$$
事实上，由式 (\ref{cohomology-equation-internal-hom})，给出这个态射等价于给出态射
$$
R\SheafHom(K, K') \otimes_{\mathcal{O}_X}^\mathbf{L}
R\SheafHom(M, M') \otimes_{\mathcal{O}_X}^\mathbf{L}
K \otimes_{\mathcal{O}_X}^\mathbf{L} M
\longrightarrow
K' \otimes_{\mathcal{O}_X}^\mathbf{L} M'
$$
为此，可以先交换中间两个因子（符号规则见 《更多代数》第 \ref{more-algebra-section-sign-rules} 节），再使用态射
$$
R\SheafHom(K, K') \otimes_{\mathcal{O}_X}^\mathbf{L} K \to K'
\quad\text{且}\quad
R\SheafHom(M, M') \otimes_{\mathcal{O}_X}^\mathbf{L} M \to M'
$$
这里应用引理 \ref{cohomology-lemma-internal-hom-composition}，并使用等同
$K = R\SheafHom(\mathcal{O}_X, K)$；对 $K'$、$M$ 和 $M'$
也作类似处理。
\end{remark}

\begin{remark}
\label{cohomology-remark-projection-formula-for-internal-hom}
设 $f:X\to Y$ 为环化空间的态射，并设 $K,L$ 为
$D(\mathcal{O}_X)$ 的对象。我们断言存在典范态射
$$
Rf_*R\SheafHom(L, K) \longrightarrow R\SheafHom(Rf_*L, Rf_*K)
$$
事实上，由式 (\ref{cohomology-equation-internal-hom})，给出这个态射等价于给出态射
$Rf_*R\SheafHom(L, K) \otimes_{\mathcal{O}_Y}^\mathbf{L} Rf_*L \to Rf_*K$.
为此，可以使用复合
$$
Rf_*R\SheafHom(L, K) \otimes_{\mathcal{O}_Y}^\mathbf{L} Rf_*L \to
Rf_*(R\SheafHom(L, K) \otimes_{\mathcal{O}_X}^\mathbf{L} L) \to
Rf_*K
$$
其中第一个箭头是相对杯积（注记 \ref{cohomology-remark-cup-product}），
第二个箭头则是对来自引理
\ref{cohomology-lemma-internal-hom-composition} 的典范态射
$R\SheafHom(L, K) \otimes_{\mathcal{O}_X}^\mathbf{L} L \to K$
应用 $Rf_*$ 所得（其中一个位置取 $\mathcal{O}_X$）。
\end{remark}

\begin{remark}
\label{cohomology-remark-relative-cup-and-composition}
设 $h:X\to Y$ 为环化空间的态射，并设 $K,M$ 为
$D(\mathcal{O}_Y)$ 的对象。图表
$$
\xymatrix{
Rf_*R\SheafHom_{\mathcal{O}_X}(K, M)
\otimes_{\mathcal{O}_Y}^\mathbf{L} Rf_*K
\ar[r] \ar[d] &
Rf_*\left(R\SheafHom_{\mathcal{O}_X}(K, M)
\otimes_{\mathcal{O}_X}^\mathbf{L} K\right)
\ar[d] \\
R\SheafHom_{\mathcal{O}_Y}(Rf_*K, Rf_*M) \otimes_{\mathcal{O}_Y}^\mathbf{L}
Rf_*K \ar[r] &
Rf_*M
}
$$
是交换的。这里，左侧竖直箭头来自注记
\ref{cohomology-remark-projection-formula-for-internal-hom}，顶部水平箭头是注记 \ref{cohomology-remark-cup-product} 中的箭头。另两个箭头是引理
\ref{cohomology-lemma-internal-hom-composition} 中态射的实例（把其中一个项换成
$\mathcal{O}_X$ 或 $\mathcal{O}_Y$）。
\end{remark}

\begin{remark}
\label{cohomology-remark-prepare-fancy-base-change}
设 $h:X\to Y$ 为环化空间的态射，并设 $K,L$ 为
$D(\mathcal{O}_Y)$ 的对象。我们断言存在典范态射
$$
Lh^*R\SheafHom(K, L) \longrightarrow R\SheafHom(Lh^*K, Lh^*L)
$$
它位于 $D(\mathcal{O}_X)$ 中。事实上，由式
(\ref{cohomology-equation-internal-hom})（该式见引理 \ref{cohomology-lemma-internal-hom}），
给出这样的态射等价于给出态射
$$
Lh^*R\SheafHom(K, L) \otimes^\mathbf{L} Lh^*K \longrightarrow Lh^*L
$$
由引理 \ref{cohomology-lemma-pullback-tensor-product}，这个箭头的源为
$Lh^*(\SheafHom(K,L)\otimes^\mathbf{L}K)$，故只需构造典范态射
$$
R\SheafHom(K, L) \otimes^\mathbf{L} K \longrightarrow L.
$$
为此，取经式 (\ref{cohomology-equation-internal-hom}) 与恒等态射
$$
\text{id} :
R\SheafHom(K, L)
\longrightarrow
R\SheafHom(K, L)
$$
相对应的箭头。
\end{remark}

\begin{remark}
\label{cohomology-remark-fancy-base-change}
设
$$
\xymatrix{
X' \ar[r]_h \ar[d]_{f'} &
X \ar[d]^f \\
S' \ar[r]^g &
S
}
$$
是环化空间的交换图表。设 $K,L$ 为 $D(\mathcal{O}_X)$ 的对象。
我们断言存在典范基变换态射
$$
Lg^*Rf_*R\SheafHom(K, L)
\longrightarrow
R(f')_*R\SheafHom(Lh^*K, Lh^*L)
$$
它位于 $D(\mathcal{O}_{S'})$ 中。具体地，取下列复合的伴随态射：
\begin{align*}
L(f')^*Lg^*Rf_*R\SheafHom(K, L)
& =
Lh^*Lf^*Rf_*R\SheafHom(K, L) \\
& \to
Lh^*R\SheafHom(K, L) \\
& \to
R\SheafHom(Lh^*K, Lh^*L)
\end{align*}
其中第一个箭头使用伴随态射 $Lf^*Rf_*\to\text{id}$，第二个箭头是注记 \ref{cohomology-remark-prepare-fancy-base-change} 中构造的典范态射。
\end{remark}





\section{Ext 层}
\label{cohomology-section-ext}

\noindent
设 $(X,\mathcal{O}_X)$ 为环化空间，并设
$K,L\in D(\mathcal{O}_X)$。利用导出范畴中内 Hom 的构造，取
$D(\mathcal{O}_X)$ 的对象 $R\SheafHom(K,L)$ 的第 $n$ 个上同调层，
得到良定的 $\mathcal{O}_X$-模层
$$
\SheafExt^n(K, L) = H^n(R\SheafHom(K, L))
$$
有时将这个对象写成 $\SheafExt^n_{\mathcal{O}_X}(K,L)$。
由引理 \ref{cohomology-lemma-section-RHom-over-U}，这个 $\SheafExt^n$-层是下列
对应的层化：
$$
U \longmapsto \Ext^n_{D(\mathcal{O}_U)}(K|_U, L|_U)
$$
由例 \ref{cohomology-example-spectral-sequence}，总存在谱序列
$$
E_2^{p, q} = H^p(X, \SheafExt^q(K, L))
$$
它在适当条件下收敛到 $\Ext^{p + q}_{D(\mathcal{O}_X)}(K,L)$
（例如 $L$ 下有界而 $K$ 上有界时）。





\section{整体导出 Hom}
\label{cohomology-section-global-RHom}

\noindent
设 $(X,\mathcal{O}_X)$ 为环化空间，并设
$K,L\in D(\mathcal{O}_X)$。利用导出范畴中内 Hom 的构造，得到良定对象
$$
R\Hom_X(K, L) = R\Gamma(X, R\SheafHom(K, L))
$$
它位于 $D(\Gamma(X,\mathcal{O}_X))$ 中。有时将这个对象写成
$R\Hom_{\mathcal{O}_X}(K,L)$。由引理
\ref{cohomology-lemma-section-RHom-over-U}，有
$$
H^0(R\Hom_X(K, L)) = \Hom_{D(\mathcal{O}_X)}(K, L),
\quad
H^p(R\Hom_X(K, L)) = \Ext_{D(\mathcal{O}_X)}^p(K, L)
$$
若 $f:Y\to X$ 是环化空间的态射，则存在典范态射
$$
R\Hom_X(K, L) \longrightarrow R\Hom_Y(Lf^*K, Lf^*L)
$$
它位于 $D(\Gamma(X,\mathcal{O}_X))$ 中，是对注记
\ref{cohomology-remark-prepare-fancy-base-change} 中定义的态射取整体截面所得。








\section{复形的粘合}
\label{cohomology-section-glueing-complexes}

\noindent
复形是可以粘合的！更准确地说，在某些情形下，可以把局部给定的
导出范畴对象粘成整体对象。我们先证明一些简单情形，然后在拓扑空间
与开覆盖的框架下证明十分一般的 \cite[Theorem 3.2.4]{BBD}。

\begin{lemma}
\label{cohomology-lemma-glue}
设 $(X,\mathcal{O}_X)$ 为环化空间，并设 $X=U\cup V$ 是 $X$ 的两个
开子空间之并。假设给定
\begin{enumerate}
\item $D(\mathcal{O}_U)$ 的对象 $A$，
\item $D(\mathcal{O}_V)$ 的对象 $B$，以及
\item 同构 $c:A|_{U\cap V}\to B|_{U\cap V}$。
\end{enumerate}
则存在 $D(\mathcal{O}_X)$ 的对象 $F$ 以及同构
$f:F|_U\to A$、$g:F|_V\to B$，使得
$c=g|_{U\cap V}\circ f^{-1}|_{U\cap V}$。此外，若给定
\begin{enumerate}
\item $D(\mathcal{O}_X)$ 的对象 $E$，
\item $D(\mathcal{O}_U)$ 中的态射 $a:A\to E|_U$，
\item $D(\mathcal{O}_V)$ 中的态射 $b:B\to E|_V$，
\end{enumerate}
满足
$$
a|_{U \cap V}  = b|_{U \cap V} \circ c.
$$
则存在 $D(\mathcal{O}_X)$ 中的态射 $F\to E$，其在 $U$ 上的限制为
$a\circ f$，在 $V$ 上的限制为 $b\circ g$。
\end{lemma}

\begin{proof}
以 $j_U$、$j_V$、$j_{U\cap V}$ 表示相应的开浸入。选取特异三角
$$
F \to Rj_{U, *}A \oplus Rj_{V, *}B \to Rj_{U \cap V, *}(B|_{U \cap V})
\to F[1]
$$
其中态射 $Rj_{V,*}B\to Rj_{U\cap V,*}(B|_{U\cap V})$ 是显然的，
而 $Rj_{U,*}A\to Rj_{U\cap V,*}(B|_{U\cap V})$ 是
$Rj_{U,*}A\to Rj_{U\cap V,*}(A|_{U\cap V})$ 与
$Rj_{U\cap V,*}c$ 的复合。限制到 $U$，得到
$$
F|_U \to A \oplus (Rj_{V, *}B)|_U \to (Rj_{U \cap V, *}(B|_{U \cap V}))|_U
\to F|_U[1]
$$
记 $j:U\cap V\to U$。限制到开集与上同调的相容性表明，
$(Rj_{V,*}B)|_U$ 与 $(Rj_{U\cap V,*}(B|_{U\cap V}))|_U$
都典范同构于 $Rj_*(B|_{U\cap V})$。因此，上一个显示图表的第二个
箭头有截面，从而态射 $F|_U\to A$ 是同构。类似地，
$F|_V\to B$ 也是同构。态射 $F\to E$ 的存在性来自 $\Hom$ 的
Mayer--Vietoris 序列；见引理 \ref{cohomology-lemma-mayer-vietoris-hom}。
\end{proof}

\begin{lemma}
\label{cohomology-lemma-vanishing-and-glueing}
设 $f:(X,\mathcal{O}_X)\to(Y,\mathcal{O}_Y)$ 为环化空间的态射，
并设 $\mathcal{B}$ 为 $Y$ 的拓扑基。
\begin{enumerate}
\item 假设 $K$ 属于 $D(\mathcal{O}_X)$，并且对 $V\in\mathcal{B}$，
当 $i<0$ 时有 $H^i(f^{-1}(V),K)=0$。则 $Rf_*K$ 的负次数上同调层
为零；对所有开集 $V\subset Y$，当 $i<0$ 时有
$H^i(f^{-1}(V),K)=0$；并且对应
$V\mapsto H^0(f^{-1}V,K)$ 是 $Y$ 上的层。
\item 假设 $K,L$ 属于 $D(\mathcal{O}_X)$，并且对
$V\in\mathcal{B}$，当 $i<0$ 时有
$\Ext^i(K|_{f^{-1}V},L|_{f^{-1}V})=0$。则对所有开集
$V\subset Y$，当 $i<0$ 时有
$\Ext^i(K|_{f^{-1}V},L|_{f^{-1}V})=0$；并且对应
$V\mapsto\Hom(K|_{f^{-1}V},L|_{f^{-1}V})$ 是 $Y$ 上的层。
\end{enumerate}
\end{lemma}

\begin{proof}
引理 \ref{cohomology-lemma-unbounded-describe-higher-direct-images} 告诉我们，
$H^i(Rf_*K)$ 是预层
$V\mapsto H^i(f^{-1}(V),K)=H^i(V,Rf_*K)$ 所伴随的层。
(1) 的假设蕴含 $Rf_*K$ 的次数 $<0$ 的上同调层为零。因此，对任意
开集 $V\subset Y$，上同调群 $H^i(V,Rf_*K)$ 在 $i<0$ 时为零，
在 $i=0$ 时等于 $H^0(V,H^0(Rf_*K))$。这证明了 (1)。

\medskip\noindent
为证明 (2)，对复形 $R\SheafHom(K,L)$ 应用 (1)，并使用引理
\ref{cohomology-lemma-section-RHom-over-U} 作相应转换。
\end{proof}

\begin{situation}
\label{cohomology-situation-locally-given}
设 $(X,\mathcal{O}_X)$ 为环化空间。给定
\begin{enumerate}
\item $X$ 的一族开集 $\mathcal{B}$；
\item 对每个 $U\in\mathcal{B}$，$D(\mathcal{O}_U)$ 中的对象 $K_U$；
\item 对满足 $V\subset U$ 且 $V,U\in\mathcal{B}$ 的每一对开集，
$D(\mathcal{O}_V)$ 中的同构 $\rho^U_V:K_U|_V\to K_V$；
\end{enumerate}
并要求每当 $W \subset V \subset U$ 且 $U, V, W$ 属于
$\mathcal{B}$ 时，
都有 $\rho^U_W=\rho^V_W\circ\rho^U_V|_W$。
\end{situation}

\noindent
若不再加假设，我们无法由此证明任何结论。一个有趣的情形是
$\mathcal{B}$ 为 $X$ 的拓扑基。另一个情形是给定拓扑空间的态射
$f:X\to Y$，而 $\mathcal{B}$ 的元素是 $Y$ 的某个拓扑基中各元素的
逆像。

\medskip\noindent
在情形 \ref{cohomology-situation-locally-given} 中，一个{\it 解}是二元组
$(K,\rho_U)$，其中 $K$ 是 $D(\mathcal{O}_X)$ 的对象，而对
$U\in\mathcal{B}$，$\rho_U:K|_U\to K_U$ 是同构，并且对所有
$V\subset U$、$U,V\in\mathcal{B}$，都有
$\rho^U_V\circ\rho_U|_V=\rho_V$。在某些情形下，解是唯一的。

\begin{lemma}
\label{cohomology-lemma-uniqueness}
在情形 \ref{cohomology-situation-locally-given} 中，假设
\begin{enumerate}
\item $X = \bigcup_{U \in \mathcal{B}} U$，并且
对 $U,V\in\mathcal{B}$ 有
$U \cap V = \bigcup_{W \in \mathcal{B}, W \subset U \cap V} W$,
\item 对任意 $U\in\mathcal{B}$，当 $i<0$ 时有
$\Ext^i(K_U,K_U)=0$。
\end{enumerate}
若解 $(K,\rho_U)$ 存在，则它在唯一同构意义下唯一；此外，当
$i<0$ 时有 $\Ext^i(K,K)=0$。
\end{lemma}

\begin{proof}
设 $(K,\rho_U)$ 与 $(K',\rho'_U)$ 是两个解。令 $f:X\to Y$ 为
《拓扑》引理 \ref{topology-lemma-create-map-from-subcollection}
中构造的连续映射，并置 $\mathcal{O}_Y=f_*\mathcal{O}_X$。
于是 $K,K'$ 和 $\mathcal{B}$ 满足引理
\ref{cohomology-lemma-vanishing-and-glueing} (2) 的情形。因此，$K$ 的负次数
Ext 群为零，并且对应
$$
V \longmapsto \Hom(K|_{f^{-1}V}, K'|_{f^{-1}V})
$$
是 $Y$ 上的层。由于 $(K,\rho_U)$ 与 $(K',\rho'_U)$ 都是解，故
$$
(\rho'_U)^{-1} \circ \rho_U : K|_U \longrightarrow K'|_U
$$
这些定义在 $U=f^{-1}(f(U))$ 上的态射在交叠处一致。因此得到上述层的
唯一整体截面，它定义所需同构 $K\to K'$，并与所有已有结构相容。
\end{proof}

\begin{remark}
\label{cohomology-remark-uniqueness}
采用引理 \ref{cohomology-lemma-uniqueness} 的记号与假设。设
$U,V\in\mathcal{B}$，并令 $\mathcal{B}'$ 为 $\mathcal{B}$ 中包含于
$U\cap V$ 的元素之集。则
$$
(\{K_{U'}\}_{U' \in \mathcal{B}'},
\{\rho_{V'}^{U'}\}_{V' \subset U'\text{ 且 }U', V' \in \mathcal{B}'})
$$
是环化空间 $U\cap V$ 上满足引理 \ref{cohomology-lemma-uniqueness} 假设的系统。
而且，$(K_U|_{U\cap V},\rho^U_{U'})$ 与
$(K_V|_{U\cap V},\rho^V_{U'})$ 都是该系统的解。由此引理，得到唯一同构
$$
\rho_{U, V} : K_U|_{U \cap V} \longrightarrow K_V|_{U \cap V}
$$
使得对每个 $U'\subset U\cap V$、$U'\in\mathcal{B}$，图表
$$
\xymatrix{
K_U|_{U'} \ar[rr]_{\rho_{U, V}|_{U'}} \ar[rd]_{\rho^U_{U'}} & &
K_V|_{U'} \ar[ld]^{\rho^V_{U'}} \\
& K_{U'}
}
$$
交换。再取第三个元素 $W\in\mathcal{B}$。得到同构
$\rho_{U, W} : K_U|_{U \cap W} \to K_W|_{U \cap W}$ 和
$\rho_{V, W} : K_U|_{V \cap W} \to K_W|_{V \cap W}$，
它们满足与 $\rho_{U,V}$ 类似的性质。最后，有
$$
\rho_{U, W}|_{U \cap V \cap W} =
\rho_{V, W}|_{U \cap V \cap W} \circ \rho_{U, V}|_{U \cap V \cap W}
$$
这由引理中的唯一性成立，因为等式两端都是与态射
$\rho^U_{U''}$ 和 $\rho^W_{U''}$ 相容的唯一同构，其中
$U''\subset U\cap V\cap W$、$U''\in\mathcal{B}$。略去一些次要细节。
集合 $(K_U,\rho_{U,V})$ 是 $X$ 的开覆盖
$\mathcal{U}:X=\bigcup_{U\in\mathcal{B}}U$ 在导出范畴中的下降数据。
用这种语言来说，寻找解的存在性，就是寻找“下降数据的有效性”。
\end{remark}

\begin{lemma}
\label{cohomology-lemma-solution-in-finite-case}
在情形 \ref{cohomology-situation-locally-given} 中，假设
\begin{enumerate}
\item $X = U_1 \cup \ldots \cup U_n$，其中 $U_i \in \mathcal{B}$；
\item 对 $U,V\in\mathcal{B}$，有
$U \cap V = \bigcup_{W \in \mathcal{B}, W \subset U \cap V} W$,
\item 对任意 $U\in\mathcal{B}$，当 $i<0$ 时有
$\Ext^i(K_U,K_U)=0$。
\end{enumerate}
则解存在，并且在唯一同构意义下唯一。
\end{lemma}

\begin{proof}
唯一性已由引理 \ref{cohomology-lemma-uniqueness} 得到。我们对 $n$ 作归纳来证明
本引理。$n=1$ 的情形是立即的。

\medskip\noindent
设 $n=2$。考虑注记 \ref{cohomology-remark-uniqueness} 中构造的同构
$\rho_{U_1, U_2} : K_{U_1}|_{U_1 \cap U_2} \to K_{U_2}|_{U_1 \cap U_2}$
。由引理 \ref{cohomology-lemma-glue}，得到 $D(\mathcal{O}_X)$ 中的对象 $K$，
以及与 $\rho_{U_1,U_2}$ 相容的同构
$\rho_{U_1}:K|_{U_1}\to K_{U_1}$ 与
$\rho_{U_2}:K|_{U_2}\to K_{U_2}$。取 $U\in\mathcal{B}$。
我们将构造同构 $\rho_U:K|_U\to K_U$，并把验证
$(K,\rho_U)$ 是解留给读者。令 $\mathcal{B}'$ 为 $\mathcal{B}$ 中
包含于 $U\cap U_1$ 或包含于 $U\cap U_2$ 的元素之集。于是
$(K_U,\rho^U_{U'})$ 是环化空间 $U$ 上系统
$(\{K_{U'}\}_{U' \in \mathcal{B}'},
\{\rho_{V'}^{U'}\}_{V' \subset U'\text{ 且 }U', V' \in \mathcal{B}'})$
的一个解。我们断言 $(K|_U,\tau_{U'})$ 是另一个解，其中对
$U'\in\mathcal{B}'$，如下选取 $\tau_{U'}$：若 $U'\subset U_1$，
则取复合
$$
K|_{U'} \xrightarrow{\rho_{U_1}|_{U'}}
K_{U_1}|_{U'} \xrightarrow{\rho^{U_1}_{U'}}
K_{U'}
$$
若 $U'\subset U_2$，则取复合
$$
K|_{U'} \xrightarrow{\rho_{U_2}|_{U'}}
K_{U_2}|_{U'} \xrightarrow{\rho^{U_2}_{U'}}
K_{U'}.
$$
为验证这是一个解，使用注记 \ref{cohomology-remark-uniqueness} 中所述
$\rho_{U_1,U_2}$ 的性质，以及 $\rho_{U_1}$、$\rho_{U_2}$ 与
$\rho_{U_1,U_2}$ 的相容性。由引理 \ref{cohomology-lemma-uniqueness}，得到唯一
同构 $K|_{U'}\to K_{U'}$；对 $U'\in\mathcal{B}'$，它与态射
$\tau_{U'}$ 和 $\rho^U_{U'}$ 相容。

\medskip\noindent
设 $n>2$。考虑开子空间 $X'=U_1\cup\ldots\cup U_{n-1}$，并令
$\mathcal{B}'$ 为 $\mathcal{B}$ 中包含于 $X'$ 的元素之集。于是得到
环化空间 $X'$ 上的系统
$(\{K_U\}_{U \in \mathcal{B}'}, \{\rho_V^U\}_{U, V \in \mathcal{B}'})$
，可以对它应用归纳假设。我们得到解 $(K_{X'},\rho^{X'}_U)$。
然后考虑 $X$ 的开集族
$\mathcal{B}^*=\mathcal{B}\cup\{X'\}$，并得到系统
$(\{K_U\}_{U \in \mathcal{B}^*},
\{\rho_V^U\}_{V \subset U\text{ 且 }U, V \in \mathcal{B}^*})$。
把引理 \ref{cohomology-lemma-uniqueness} 应用于解 $K_{X'}$，可知这个新系统
也满足条件 (3)。对该系统有 $X=X'\cup U_n$，因而归结到上面已经
处理的 $n=2$ 情形。
\end{proof}

\begin{lemma}
\label{cohomology-lemma-glueing-increasing-union}
设 $X$ 为环化空间，$E$ 为良序集，并设
$$
X = \bigcup\nolimits_{\alpha \in E} W_\alpha
$$
是一个开覆盖，满足 $W_\alpha\subset W_{\alpha+1}$，并且当 $\alpha$
不是后继元时，
$W_\alpha=\bigcup_{\beta<\alpha}W_\beta$。设 $K_\alpha$ 为
$D(\mathcal{O}_{W_\alpha})$ 的对象，且在 $i<0$ 时有
$\Ext^i(K_\alpha,K_\alpha)=0$。假设对所有 $\beta<\alpha$ 给定
同构
$\rho_\beta^\alpha :  K_\alpha|_{W_\beta} \to K_\beta$，
它位于 $D(\mathcal{O}_{W_\beta})$ 中；并且当
$\gamma<\beta<\alpha$ 时有
$\rho_\gamma^\alpha = \rho_\gamma^\beta \circ \rho^\alpha_\beta|_{W_\gamma}$
。则存在 $D(\mathcal{O}_X)$ 的对象 $K$，以及对
$\alpha\in E$ 的同构 $K|_{W_\alpha}\to K_\alpha$，它们与同构
$\rho_\beta^\alpha$ 相容。
\end{lemma}

\begin{proof}
在本证明中，$\alpha,\beta,\gamma,\ldots$ 表示 $E$ 的元素。
在 $W_\alpha$ 上选取表示 $K_\alpha$ 的 K-内射复形
$I_\alpha^\bullet$。对 $\beta<\alpha$，以
$j_{\beta,\alpha}:W_\beta\to W_\alpha$ 表示包含态射。我们将用
超限递归，对所有 $\beta<\alpha$ 构造复形态射
$$
\tau_{\beta, \alpha} :
(j_{\beta, \alpha})_!I_\beta^\bullet
\longrightarrow
I_\alpha^\bullet
$$
它表示同构
$\rho^\alpha_\beta:K_\alpha|_{W_\beta}\to K_\beta$ 之逆的伴随态射。
此外，我们将使得当 $\gamma<\beta<\alpha$ 时有
$$
\tau_{\gamma, \alpha} = \tau_{\beta, \alpha} \circ
(j_{\beta, \alpha})_!\tau_{\gamma, \beta}
$$
作为复形态射的等式。具体地，假设已对所有
$\gamma<\beta<\alpha$ 给定能正确复合的 $\tau_{\gamma,\beta}$。
若 $\alpha=\alpha'+1$ 是后继元，则任取复形态射
$$
(j_{\alpha', \alpha})_!I_{\alpha'}^\bullet \to I_\alpha^\bullet
$$
它是同构
$\rho^\alpha_{\alpha'}:K_\alpha|_{W_{\alpha'}}\to K_{\alpha'}$
之逆的伴随态射（由于 $I_\alpha^\bullet$ 是 K-内射的，这样选取是
可能的），并且对任意 $\beta<\alpha'$，置
$$
\tau_{\beta, \alpha} = \tau_{\alpha', \alpha} \circ
(j_{\alpha', \alpha})_!\tau_{\beta, \alpha'}
$$
若 $\alpha$ 不是后继元，则考虑 $W_\alpha$ 上的复形
$$
C^\bullet = \colim_{\beta < \alpha} (j_{\beta, \alpha})_!I_\beta^\bullet
$$
（逐项余极限），其中该系统的转移态射由
$\beta'<\beta<\alpha$ 时的态射 $\tau_{\beta',\beta}$ 给出。
我们断言 $C^\bullet$ 表示 $K_\alpha$。事实上，对
$\beta<\alpha$，余投影
$(j_{\beta,\alpha})_!I_\beta^\bullet\to C^\bullet$ 的限制给出态射
$$
\sigma_\beta : I_\beta^\bullet \longrightarrow C^\bullet|_{W_\beta}
$$
它是拟同构：若 $x\in W_\beta$，则考察茎得到
$$
(C^\bullet)_x =
\colim_{\beta' < \alpha}
\left((j_{\beta', \alpha})_!I_{\beta'}^\bullet\right)_x =
\colim_{\beta \leq \beta' < \alpha} (I_{\beta'}^\bullet)_x
\longleftarrow
(I_\beta^\bullet)_x
$$
这是拟同构。这里使用了取茎与余极限交换、滤过余极限正合，以及当
$\beta\leq\beta'<\alpha$ 时态射
$(I_\beta^\bullet)_x\to(I_{\beta'}^\bullet)_x$ 是拟同构。
因此，$(C^\bullet,\sigma_\beta^{-1})$ 是系统
$(\{K_\beta\}_{\beta<\alpha},
\{\rho^\beta_{\beta'}\}_{\beta'<\beta<\alpha})$ 的解。
由于 $(K_\alpha,\rho^\alpha_\beta)$ 是另一个解，我们得到
$D(\mathcal{O}_{W_\alpha})$ 中与所有态射相容的唯一同构
$\sigma:K_\alpha\to C^\bullet$；见引理
\ref{cohomology-lemma-solution-in-finite-case}（正是在这里用到了负次数 Ext 群
为零）。选取表示 $\sigma$ 的复形态射
$\tau:C^\bullet\to I_\alpha^\bullet$，并置
$$
\tau_{\beta, \alpha} = \tau|_{W_\beta} \circ \sigma_\beta
$$
由此得到所需态射。最后，取 $K$ 为下列复形所表示的导出范畴对象：
$$
K^\bullet = \colim_{\alpha \in E} (W_\alpha \to X)_!I_\alpha^\bullet
$$
其中转移态射由我们对 $\beta<\alpha$ 精心构造的
$\tau_{\beta,\alpha}$ 给出。与上面完全相同的论证表明，对所有
$\alpha$，余投影的限制确定同构
$$
K|_{W_\alpha} \longrightarrow K_\alpha
$$
它与给定态射 $\rho^\alpha_\beta$ 相容。
\end{proof}

\noindent
利用超限归纳，我们可以证明一般情形下的结论。

\begin{theorem}[BBD 粘合引理]
\label{cohomology-theorem-glueing-bbd-general}
\begin{reference}
\cite[Theorem 3.2.4]{BBD} 在不加有界性假设时的特殊情形。
\end{reference}
在情形 \ref{cohomology-situation-locally-given} 中，假设
\begin{enumerate}
\item $X = \bigcup_{U \in \mathcal{B}} U$,
\item 对 $U,V\in\mathcal{B}$，有
$U \cap V = \bigcup_{W \in \mathcal{B}, W \subset U \cap V} W$,
\item 对任意 $U\in\mathcal{B}$，当 $i<0$ 时有
$\Ext^i(K_U,K_U)=0$。
\end{enumerate}
则存在 $D(\mathcal{O}_X)$ 的对象 $K$，以及对
$U\in\mathcal{B}$ 的 $D(\mathcal{O}_U)$ 中的同构
$\rho_U:K|_U\to K_U$，使得对所有满足
$V\subset U$、$U,V\in\mathcal{B}$ 的开集，都有
$\rho^U_V\circ\rho_U|_V=\rho_V$。二元组 $(K,\rho_U)$ 在唯一同构
意义下唯一。
\end{theorem}

\begin{proof}
上文把二元组 $(K,\rho_U)$ 称为解。唯一性来自引理
\ref{cohomology-lemma-uniqueness}。若 $X$ 有由 $\mathcal{B}$ 中元素组成的
有限覆盖（例如 $X$ 拟紧），则本定理是引理
\ref{cohomology-lemma-solution-in-finite-case} 的推论。在一般情形下，用超限归纳
和引理 \ref{cohomology-lemma-glueing-increasing-union} 作完全相同的论证。

\medskip\noindent
首先用超限递归，对任意序数 $\alpha$ 选取开集
$W_\alpha\subset X$。具体地，置 $W_0=\emptyset$。若
$\alpha=\beta+1$ 是后继序数，则当 $W_\beta=X$ 时置
$W_\alpha=X$；当 $W_\beta \not = X$ 时，取不包含于 $W_\beta$ 的
$U_\alpha\in\mathcal{B}$，并置
$W_\alpha=W_\beta\cup U_\alpha$。若 $\alpha$ 是极限序数，则置
$W_\alpha=\bigcup_{\beta<\alpha}W_\beta$。于是当 $\alpha$ 足够大时，
$W_\alpha=X$。注意，对每个 $\alpha$，开集 $W_\alpha$ 都是
$\mathcal{B}$ 中元素之并。因此，若
$\mathcal{B}_\alpha=\{U\in\mathcal{B},U\subset W_\alpha\}$，则
$$
S_\alpha = (\{K_U\}_{U \in \mathcal{B}_\alpha},
\{\rho_V^U\}_{V \subset U\text{ 且 }U, V \in \mathcal{B}_\alpha})
$$
是环化空间 $W_\alpha$ 上引理 \ref{cohomology-lemma-uniqueness} 所述的系统。

\medskip\noindent
我们用超限归纳证明每个 $\alpha$ 的系统 $S_\alpha$ 都有解。
由于 $\alpha$ 足够大时该系统就是定理中给定的系统，这将证明定理。

\medskip\noindent
设 $\alpha=\beta+1$ 为后继序数。（上面引理的证明已经处理过这个
情形，但为清楚起见，我们重复该论证。）此时
$W_\alpha=W_\beta\cup U_\alpha$，其中
$U_\alpha\in\mathcal{B}$。由归纳假设，系统 $S_\beta$ 有解
$(K_{W_\beta}, \{\rho^{W_\beta}_U\}_{U \in \mathcal{B}_\beta})$
。于是可以考虑 $W_\alpha$ 的开集族
$\mathcal{B}_\alpha^* = \mathcal{B}_\alpha \cup \{W_\beta\}$
，并得到系统
$(\{K_U\}_{U \in \mathcal{B}_\alpha^*},
\{\rho_V^U\}_{V \subset U\text{ 且 }U, V \in \mathcal{B}_\alpha^*})$。
把引理 \ref{cohomology-lemma-uniqueness} 应用于解 $K_{W_\beta}$，可知这个
新系统也满足条件 (3)。对该系统有
$W_\alpha=W_\beta\cup U_\alpha$，因而归结到引理
\ref{cohomology-lemma-solution-in-finite-case} 已处理的情形。

\medskip\noindent
设 $\alpha$ 为极限序数。回忆此时
$W_\alpha=\bigcup_{\beta<\alpha}W_\beta$。对 $\beta<\alpha$，令
$(K_{W_\beta}, \{\rho^{W_\beta}_U\}_{U \in \mathcal{B}_\beta})$
为 $S_\beta$ 的解。当 $\gamma<\beta<\alpha$ 时，限制
$K_{W_\beta}|_{W_\gamma}$ 连同对
$U\in\mathcal{B}_\gamma$ 的态射 $\rho^{W_\beta}_U$，构成
$S_\gamma$ 的解。由唯一性，得到唯一同构
$\rho_{W_\gamma}^{W_\beta} : K_{W_\beta}|_{W_\gamma} \to K_{W_\gamma}$
它与对 $U\in\mathcal{B}_\gamma$ 的态射 $\rho^{W_\beta}_U$ 和
$\rho^{W_\gamma}_U$ 相容。这些态射以正确方式复合；也就是说，当
$\delta<\gamma<\beta<\alpha$ 时，
$\rho_{W_\delta}^{W_\gamma}\circ
\rho_{W_\gamma}^{W_\beta}|_{W_\delta}
=\rho^{W_\delta}_{W_\beta}$。因此，可以应用引理
\ref{cohomology-lemma-glueing-increasing-union}（注意，把引理
\ref{cohomology-lemma-uniqueness} 应用于解 $K_{W_\beta}$，可知
$K_{W_\beta}$ 的负次数 Ext 群为零），得到 $K_{W_\alpha}$ 和同构
$$
\rho_{W_\beta}^{W_\alpha} :
K_{W_\alpha}|_{W_\beta}
\longrightarrow
K_{W_\beta}
$$
它们与 $\gamma<\beta<\alpha$ 时的态射
$\rho_{W_\gamma}^{W_\beta}$ 相容。

\medskip\noindent
为证明 $K_{W_\alpha}$ 是解，还需对 $U\in\mathcal{B}_\alpha$
构造满足某些相容性的同构
$\rho_U^{W_\alpha}:K_{W_\alpha}|_U\to K_U$。我们把
$\rho_U^{W_\alpha}$ 选为唯一满足如下条件的态射：对任意
$\beta<\alpha$ 以及任意满足 $V\subset U$ 的
$V\in\mathcal{B}_\beta$，图表
$$
\xymatrix{
K_{W_\alpha}|_V \ar[r]_{\rho_U^{W_\alpha}|_V}
\ar[d]_{\rho_{W_\beta}^{W_\alpha}|_V}
& K_U|_V \ar[d]^{\rho_U^V} \\
K_{W_\beta} \ar[r]^{\rho_V^{W_\beta}}
& K_V
}
$$
交换。这一选取有意义，因为
$$
(\{K_V\}_{V \subset U, V \in \mathcal{B}_\beta\text{ 对某个 }\beta < \alpha},
\{\rho_V^{V'}\}_{V \subset V'\text{ 且 }V, V' \subset U
\text{，且 }V, V' \in \mathcal{B}_\beta\text{ 对某个 }\beta < \alpha})
$$
是环化空间 $U$ 上引理 \ref{cohomology-lemma-uniqueness} 所述的系统，并且
$(K_U,\rho^U_V)$ 与
$(K_{W_\alpha}|_U,  \rho_V^{W_\beta}\circ \rho_{W_\beta}^{W_\alpha}|_V)$
都是该系统的解。这给出存在性和唯一性。略去这些态射满足所需相容性的
证明（这只是记号整理）。
\end{proof}



\section{严格完美复形}
\label{cohomology-section-strictly-perfect}

\noindent
模的严格完美复形将在后面用于定义伪凝聚复形和完美复形。其定义如下。

\begin{definition}
\label{cohomology-definition-strictly-perfect}
设 $(X,\mathcal{O}_X)$ 为环化空间，并设 $\mathcal{E}^\bullet$
为 $\mathcal{O}_X$-模复形。若仅有有限多个 $i$ 使
$\mathcal{E}^i$ 非零，并且对所有 $i$，$\mathcal{E}^i$ 都是某个
有限自由 $\mathcal{O}_X$-模的直和项，则称
$\mathcal{E}^\bullet$ 是{\it 严格完美的}。
\end{definition}

\noindent
警告：由于我们没有假设 $X$ 是局部环化空间，有限自由
$\mathcal{O}_X$-模的直和项未必是有限局部自由的。

\begin{lemma}
\label{cohomology-lemma-cone}
严格完美复形之间态射的锥是严格完美的。
\end{lemma}

\begin{proof}
这由定义立即可得。
\end{proof}

\begin{lemma}
\label{cohomology-lemma-tensor}
两个严格完美复形的张量积所对应的全复形是严格完美的。
\end{lemma}

\begin{proof}
从略。
\end{proof}

\begin{lemma}
\label{cohomology-lemma-strictly-perfect-pullback}
设 $f:(X,\mathcal{O}_X)\to(Y,\mathcal{O}_Y)$ 为环化空间的态射。
若 $\mathcal{F}^\bullet$ 是严格完美的 $\mathcal{O}_Y$-模复形，
则 $f^*\mathcal{F}^\bullet$ 是严格完美的
$\mathcal{O}_X$-模复形。
\end{lemma}

\begin{proof}
有限自由模的拉回仍是有限自由的。函子 $f^*$ 是加性函子，因而保持
直和项。引理得证。
\end{proof}

\begin{lemma}
\label{cohomology-lemma-local-lift-map}
设 $(X,\mathcal{O}_X)$ 为环化空间。给定 $\mathcal{O}_X$-模的实线图表
$$
\xymatrix{
\mathcal{E} \ar@{..>}[dr] \ar[r] & \mathcal{F} \\
& \mathcal{G} \ar[u]_p
}
$$
其中 $\mathcal{E}$ 是某个有限自由 $\mathcal{O}_X$-模的直和项，
并且 $p$ 满射，则在 $X$ 上局部存在使图表交换的虚线箭头。
\end{lemma}

\begin{proof}
可以假设对某个 $n$ 有 $\mathcal{E}=\mathcal{O}_X^{\oplus n}$。
此时，寻找虚线箭头等价于提升基元素在
$\Gamma(X,\mathcal{F})$ 中的像。由层满射的刻画（《层》第 \ref{sheaves-section-exactness-points} 节），这在局部是可行的。
\end{proof}

\begin{lemma}
\label{cohomology-lemma-local-homotopy}
设 $(X,\mathcal{O}_X)$ 为环化空间。
\begin{enumerate}
\item 设 $\alpha:\mathcal{E}^\bullet\to\mathcal{F}^\bullet$
为 $\mathcal{O}_X$-模复形的态射，其中 $\mathcal{E}^\bullet$
严格完美，而 $\mathcal{F}^\bullet$ 无上同调。则 $\alpha$
在 $X$ 上局部零同伦。
\item 设 $\alpha:\mathcal{E}^\bullet\to\mathcal{F}^\bullet$
为 $\mathcal{O}_X$-模复形的态射，其中 $\mathcal{E}^\bullet$
严格完美，且当 $i<a$ 时有 $\mathcal{E}^i=0$，当 $i\geq a$ 时有
$H^i(\mathcal{F}^\bullet)=0$。则 $\alpha$ 在 $X$ 上局部零同伦。
\end{enumerate}
\end{lemma}

\begin{proof}
第一条由第二条推出，故只证明 (2)。我们对复形
$\mathcal{E}^\bullet$ 的长度作归纳。若对某个有限自由
$\mathcal{O}_X$-模的直和项 $\mathcal{E}$ 及整数 $n\geq a$，有
$\mathcal{E}^\bullet\cong\mathcal{E}[-n]$，则结论来自引理
\ref{cohomology-lemma-local-lift-map} 以及如下事实：
$\mathcal{F}^{n - 1} \to \Ker(\mathcal{F}^n \to \mathcal{F}^{n + 1})$
由所假设的 $H^n(\mathcal{F}^\bullet)$ 为零而满射。若
$\mathcal{E}^i$ 仅在 $i\in[a,b]$ 时可能非零，则有复形的分裂正合列
$$
0 \to \mathcal{E}^b[-b] \to \mathcal{E}^\bullet \to
\sigma_{\leq b - 1}\mathcal{E}^\bullet \to 0
$$
它在 $K(\mathcal{O}_X)$ 中确定一个特异三角。因此，由三角范畴的
公理得到正合列
$$
\Hom_{K(\mathcal{O}_X)}(
\sigma_{\leq b - 1}\mathcal{E}^\bullet, \mathcal{F}^\bullet)
\to
\Hom_{K(\mathcal{O}_X)}(\mathcal{E}^\bullet, \mathcal{F}^\bullet)
\to
\Hom_{K(\mathcal{O}_X)}(\mathcal{E}^b[-b], \mathcal{F}^\bullet)
$$
。复合 $\mathcal{E}^b[-b]\to\mathcal{F}^\bullet$ 在局部零同伦，
故可以假设我们的态射来自上述显示正合列左端的一个元素。由归纳假设，
这个元素在局部为零。
\end{proof}

\begin{lemma}
\label{cohomology-lemma-lift-through-quasi-isomorphism}
设 $(X,\mathcal{O}_X)$ 为环化空间。给定
$\mathcal{O}_X$-模复形的实线图表
$$
\xymatrix{
\mathcal{E}^\bullet \ar@{..>}[dr] \ar[r]_\alpha & \mathcal{F}^\bullet \\
& \mathcal{G}^\bullet \ar[u]_f
}
$$
其中 $\mathcal{E}^\bullet$ 严格完美，当 $j<a$ 时
$\mathcal{E}^j=0$，而 $H^j(f)$ 在 $j>a$ 时为同构、在 $j=a$ 时
满射，则在 $X$ 上局部存在使图表在同伦意义下交换的虚线箭头。
\end{lemma}

\begin{proof}
对 $f$ 的假设蕴含锥 $C(f)^\bullet$ 的次数 $\geq a$ 的上同调层
为零。因此，由引理 \ref{cohomology-lemma-local-homotopy}，存在开覆盖
$X=\bigcup U_i$，使得复合
$\mathcal{E}^\bullet \to \mathcal{F}^\bullet \to C(f)^\bullet$
在 $U_i$ 上零同伦。由于
$$
\mathcal{G}^\bullet \to \mathcal{F}^\bullet \to C(f)^\bullet \to
\mathcal{G}^\bullet[1]
$$
限制后成为 $K(\mathcal{O}_{U_i})$ 中的特异三角，故可把
$\alpha|_{U_i}$ 在同伦意义下提升为所需态射
$\alpha_i:\mathcal{E}^\bullet|_{U_i}\to\mathcal{G}^\bullet|_{U_i}$。
\end{proof}

\begin{lemma}
\label{cohomology-lemma-local-actual}
设 $(X,\mathcal{O}_X)$ 为环化空间。设
$\mathcal{E}^\bullet$、$\mathcal{F}^\bullet$ 为
$\mathcal{O}_X$-模复形，并且 $\mathcal{E}^\bullet$ 严格完美。
\begin{enumerate}
\item 对任意元素
$\alpha \in \Hom_{D(\mathcal{O}_X)}(\mathcal{E}^\bullet, \mathcal{F}^\bullet)$
，存在开覆盖 $X=\bigcup U_i$，使得 $\alpha|_{U_i}$ 由复形态射
$\alpha_i:\mathcal{E}^\bullet|_{U_i}\to
\mathcal{F}^\bullet|_{U_i}$ 给出。
\item 给定复形态射
$\alpha:\mathcal{E}^\bullet\to\mathcal{F}^\bullet$，若它在群
$\Hom_{D(\mathcal{O}_X)}(\mathcal{E}^\bullet, \mathcal{F}^\bullet)$
中的像为零，则存在开覆盖 $X=\bigcup U_i$，使得
$\alpha|_{U_i}$ 零同伦。
\end{enumerate}
\end{lemma}

\begin{proof}
证明 (1)。由导出范畴的构造，可以找到拟同构
$f:\mathcal{F}^\bullet\to\mathcal{G}^\bullet$ 以及复形态射
$\beta:\mathcal{E}^\bullet\to\mathcal{G}^\bullet$，使得
$\alpha=f^{-1}\beta$。因此，结论来自引理
\ref{cohomology-lemma-lift-through-quasi-isomorphism}。(2) 的证明从略。
\end{proof}

\begin{lemma}
\label{cohomology-lemma-Rhom-strictly-perfect}
设 $(X,\mathcal{O}_X)$ 为环化空间。设
$\mathcal{E}^\bullet$、$\mathcal{F}^\bullet$ 为
$\mathcal{O}_X$-模复形，并且 $\mathcal{E}^\bullet$ 严格完美。
则内 Hom $R\SheafHom(\mathcal{E}^\bullet,\mathcal{F}^\bullet)$
由复形 $\mathcal{H}^\bullet$ 表示，其各项为
$$
\mathcal{H}^n =
\bigoplus\nolimits_{n = p + q}
\SheafHom_{\mathcal{O}_X}(\mathcal{E}^{-q}, \mathcal{F}^p)
$$
，微分如第 \ref{cohomology-section-hom-complexes} 节所述。
\end{lemma}

\begin{proof}
选取到 K-内射复形的拟同构
$\mathcal{F}^\bullet\to\mathcal{I}^\bullet$。令
$(\mathcal{H}')^\bullet$ 为各项如下的复形：
$$
(\mathcal{H}')^n =
\prod\nolimits_{n = p + q}
\SheafHom_{\mathcal{O}_X}(\mathcal{E}^{-q}, \mathcal{I}^p)
$$
由第 \ref{cohomology-section-internal-hom} 节的构造，它表示
$R\SheafHom(\mathcal{E}^\bullet,\mathcal{F}^\bullet)$。只需证明态射
$$
\mathcal{H}^\bullet \longrightarrow (\mathcal{H}')^\bullet
$$
是拟同构。给定开集 $U\subset X$，直接考察可得
$$
H^0(\mathcal{H}^\bullet(U)) =
\Hom_{K(\mathcal{O}_U)}(\mathcal{E}^\bullet|_U, \mathcal{I}^\bullet|_U)
\to
H^0((\mathcal{H}')^\bullet(U)) =
\Hom_{D(\mathcal{O}_U)}(\mathcal{E}^\bullet|_U, \mathcal{I}^\bullet|_U)
$$
由引理 \ref{cohomology-lemma-local-actual}，对应
$U\mapsto H^0(\mathcal{H}^\bullet(U))$ 的层化等于对应
$U\mapsto H^0((\mathcal{H}')^\bullet(U))$ 的层化。其他上同调层也可作
类似论证。因此 $\mathcal{H}^\bullet$ 与
$(\mathcal{H}')^\bullet$ 拟同构，引理得证。
\end{proof}

\begin{lemma}
\label{cohomology-lemma-Rhom-strictly-perfect-K-flat}
在引理 \ref{cohomology-lemma-Rhom-strictly-perfect} 的情形下，若
$\mathcal{F}^\bullet$ 是 K-平坦的，则
$\mathcal{H}^\bullet$ 也是 K-平坦的。
\end{lemma}

\begin{proof}
注意，$\mathcal{H}^\bullet$ 就是 Hom 复形
$\SheafHom^\bullet(\mathcal{E}^\bullet, \mathcal{F}^\bullet)$
，因为严格完美复形 $\mathcal{E}^\bullet$ 的有界性保证 Hom 复形
定义中的乘积变成直和。设 $\mathcal{K}^\bullet$ 为无上同调的
$\mathcal{O}_X$-模复形。考虑引理 \ref{cohomology-lemma-diagonal-better} 中的态射
$$
\gamma :
\text{Tot}(\mathcal{K}^\bullet \otimes
\SheafHom^\bullet(\mathcal{E}^\bullet, \mathcal{F}^\bullet))
\longrightarrow
\SheafHom^\bullet(\mathcal{E}^\bullet,
\text{Tot}(\mathcal{K}^\bullet \otimes \mathcal{F}^\bullet))
$$
。由于 $\mathcal{F}^\bullet$ 是 K-平坦的，复形
$\text{Tot}(\mathcal{K}^\bullet\otimes\mathcal{F}^\bullet)$ 无上同调；
因而由引理 \ref{cohomology-lemma-local-actual}（或者也可用引理
\ref{cohomology-lemma-Rhom-strictly-perfect}），$\gamma$ 的目标也无上同调。
所以，为证明本引理，只需证明 $\gamma$ 是复形同构。为此，可以对
复形 $\mathcal{E}^\bullet$ 的长度作归纳。若长度 $\leq1$，则对某个
$n\geq0$ 和 $k\in\mathbf{Z}$，$\mathcal{E}^\bullet$ 是
$\mathcal{O}_X^{\oplus n}[k]$ 的直和项，此时直接考察即可得到结论。
若长度 $>1$，则考虑逐项分裂的复形短正合列
$$
0 \to \sigma_{\geq a + 1} \mathcal{E}^\bullet \to
\mathcal{E}^\bullet \to
\sigma_{\leq a} \mathcal{E}^\bullet \to 0
$$
，其中 $a\in\mathbf{Z}$ 适当；见 《同调代数》第 \ref{homology-section-truncations} 节。于是 $\gamma$ 嵌入逐项
分裂的复形短正合列之间的态射。由归纳假设，$\gamma$ 对
$\sigma_{\geq a+1}\mathcal{E}^\bullet$ 和
$\sigma_{\leq a}\mathcal{E}^\bullet$ 都是同构，因而对
$\mathcal{E}^\bullet$ 的结论也成立。略去一些细节。
\end{proof}

\begin{lemma}
\label{cohomology-lemma-Rhom-complex-of-direct-summands-finite-free}
设 $(X,\mathcal{O}_X)$ 为环化空间。设
$\mathcal{E}^\bullet$、$\mathcal{F}^\bullet$ 为
$\mathcal{O}_X$-模复形，满足
\begin{enumerate}
\item 当 $n\ll0$ 时，$\mathcal{F}^n=0$；
\item 当 $n\gg0$ 时，$\mathcal{E}^n=0$；并且
\item $\mathcal{E}^n$ 同构于某个有限自由
$\mathcal{O}_X$-模的直和项。
\end{enumerate}
则内 Hom $R\SheafHom(\mathcal{E}^\bullet,\mathcal{F}^\bullet)$
由复形 $\mathcal{H}^\bullet$ 表示，其各项为
$$
\mathcal{H}^n =
\bigoplus\nolimits_{n = p + q}
\SheafHom_{\mathcal{O}_X}(\mathcal{E}^{-q}, \mathcal{F}^p)
$$
，微分如第 \ref{cohomology-section-internal-hom} 节所述。
\end{lemma}

\begin{proof}
选取拟同构 $\mathcal{F}^\bullet\to\mathcal{I}^\bullet$，其中
$\mathcal{I}^\bullet$ 是下有界的内射复形。注意
$\mathcal{I}^\bullet$ 是 K-内射的（《导出范畴》引理
\ref{derived-lemma-bounded-below-injectives-K-injective}）。因此，
第 \ref{cohomology-section-internal-hom} 节的构造表明
$R\SheafHom(\mathcal{E}^\bullet,\mathcal{F}^\bullet)$ 由复形
$(\mathcal{H}')^\bullet$ 表示，其各项为
$$
(\mathcal{H}')^n =
\prod\nolimits_{n = p + q}
\SheafHom_{\mathcal{O}_X}(\mathcal{E}^{-q}, \mathcal{I}^p) =
\bigoplus\nolimits_{n = p + q}
\SheafHom_{\mathcal{O}_X}(\mathcal{E}^{-q}, \mathcal{I}^p)
$$
（之所以相等，是因为仅有有限多个非零项）。注意
$\mathcal{H}^\bullet$ 是各项为
$\SheafHom_{\mathcal{O}_X}(\mathcal{E}^{-q}, \mathcal{F}^p)$
的双复形所对应的全复形；$(\mathcal{H}')^\bullet$ 也类似。
自然态射 $(\mathcal{H}')^\bullet\to\mathcal{H}^\bullet$
来自双复形的态射。因此，为证明这个态射是拟同构，可以使用双复形的
谱序列（《同调代数》引理 \ref{homology-lemma-first-quadrant-ss}）
$$
{}'E_1^{p, q} =
H^p(\SheafHom_{\mathcal{O}_X}(\mathcal{E}^{-q}, \mathcal{F}^\bullet))
$$
，它收敛到 $H^{p+q}(\mathcal{H}^\bullet)$；
$(\mathcal{H}')^\bullet$ 也类似。为完成本引理的证明，只需证明
$\mathcal{F}^\bullet\to\mathcal{I}^\bullet$ 诱导上同调层的同构
$$
H^p(\SheafHom_{\mathcal{O}_X}(\mathcal{E}, \mathcal{F}^\bullet))
\longrightarrow
H^p(\SheafHom_{\mathcal{O}_X}(\mathcal{E}, \mathcal{I}^\bullet))
$$
，其中 $\mathcal{E}$ 是某个有限自由 $\mathcal{O}_X$-模的直和项。
当 $\mathcal{E}$ 有限自由时这是显然的，故结论成立。
\end{proof}





\section{伪凝聚模}
\label{cohomology-section-pseudo-coherent}

\noindent
本节讨论伪凝聚复形。

\begin{definition}
\label{cohomology-definition-pseudo-coherent}
设 $(X,\mathcal{O}_X)$ 为环化空间，$\mathcal{E}^\bullet$ 为
$\mathcal{O}_X$-模复形，并设 $m\in\mathbf{Z}$。
\begin{enumerate}
\item 若存在开覆盖 $X=\bigcup U_i$，并且对每个 $i$ 存在复形态射
$\alpha_i : \mathcal{E}_i^\bullet \to \mathcal{E}^\bullet|_{U_i}$
，其中 $\mathcal{E}_i^\bullet$ 在 $U_i$ 上严格完美，并且
$H^j(\alpha_i)$ 在 $j>m$ 时为同构，而 $H^m(\alpha_i)$ 满射，
则称 $\mathcal{E}^\bullet$ 是{\it $m$-伪凝聚的}。
\item 若 $\mathcal{E}^\bullet$ 对所有 $m$ 都是 $m$-伪凝聚的，
则称它是{\it 伪凝聚的}。
\item 当且仅当 $D(\mathcal{O}_X)$ 的对象 $E$ 可由一个
$m$-伪凝聚（相应地，伪凝聚）的 $\mathcal{O}_X$-模复形表示时，
称该对象是{\it $m$-伪凝聚的}（相应地，{\it 伪凝聚的}）。
\end{enumerate}
\end{definition}

\noindent
若 $X$ 拟紧，则 $D(\mathcal{O}_X)$ 的 $m$-伪凝聚对象属于
$D^-(\mathcal{O}_X)$。但若 $X$ 不拟紧，则未必如此。

\begin{lemma}
\label{cohomology-lemma-pseudo-coherent-independent-representative}
设 $(X,\mathcal{O}_X)$ 为环化空间，并设 $E$ 为
$D(\mathcal{O}_X)$ 的对象。
\begin{enumerate}
\item 若存在开覆盖 $X=\bigcup U_i$、$U_i$ 上的严格完美复形
$\mathcal{E}_i^\bullet$，以及 $D(\mathcal{O}_{U_i})$ 中的态射
$\alpha_i:\mathcal{E}_i^\bullet\to E|_{U_i}$，使得
$H^j(\alpha_i)$ 在 $j>m$ 时为同构，而 $H^m(\alpha_i)$ 满射，
则 $E$ 是 $m$-伪凝聚的。
\item 若 $E$ 是 $m$-伪凝聚的，则表示 $E$ 的任意复形都是
$m$-伪凝聚的。
\end{enumerate}
\end{lemma}

\begin{proof}
任取表示 $E$ 的复形 $\mathcal{F}^\bullet$，并设
$X=\bigcup U_i$ 及
$\alpha_i:\mathcal{E}_i^\bullet\to E|_{U_i}$ 如 (1) 所述。
我们将证明复形 $\mathcal{F}^\bullet$ 是 $m$-伪凝聚的，这将同时证明
(1) 和 (2)。由引理 \ref{cohomology-lemma-local-actual}，细化开覆盖
$X=\bigcup U_i$ 后，可以用复形态射
$\alpha_i:\mathcal{E}_i^\bullet\to\mathcal{F}^\bullet|_{U_i}$
表示各态射 $\alpha_i$。由假设，$H^j(\alpha_i)$ 在 $j>m$ 时为同构，
而 $H^m(\alpha_i)$ 满射，故 $\mathcal{F}^\bullet$ 是
$m$-伪凝聚的。
\end{proof}

\begin{lemma}
\label{cohomology-lemma-pseudo-coherent-pullback}
设 $f:(X,\mathcal{O}_X)\to(Y,\mathcal{O}_Y)$ 为环化空间的态射，
并设 $E$ 为 $D(\mathcal{O}_Y)$ 的对象。若 $E$ 是
$m$-伪凝聚的，则 $Lf^*E$ 也是 $m$-伪凝聚的。
\end{lemma}

\begin{proof}
用 $\mathcal{O}_Y$-模复形 $\mathcal{E}^\bullet$ 表示 $E$，并按定义
\ref{cohomology-definition-pseudo-coherent} 选取开覆盖 $Y=\bigcup V_i$ 和态射
$\alpha_i : \mathcal{E}_i^\bullet \to \mathcal{E}^\bullet|_{V_i}$
。置 $U_i=f^{-1}(V_i)$。由引理
\ref{cohomology-lemma-pseudo-coherent-independent-representative}，只需证明
$Lf^*\mathcal{E}^\bullet|_{U_i}$ 是 $m$-伪凝聚的。选取特异三角
$$
\mathcal{E}_i^\bullet \to
\mathcal{E}^\bullet|_{V_i} \to
C \to
\mathcal{E}_i^\bullet[1]
$$
对 $\alpha_i$ 的假设恰好意味着对所有 $j\geq m$，上同调层
$H^j(C)$ 为零。以 $f_i:U_i\to V_i$ 表示 $f$ 的限制。注意
$Lf^*\mathcal{E}^\bullet|_{U_i}=
Lf_i^*(\mathcal{E}|_{V_i})$。应用 $Lf_i^*$，得到特异三角
$$
Lf_i^*\mathcal{E}_i^\bullet \to
Lf_i^*\mathcal{E}|_{V_i} \to
Lf_i^*C \to
Lf_i^*\mathcal{E}_i^\bullet[1]
$$
由 $Lf_i^*$ 作为左导出函子的构造可知，对 $j\geq m$，
$H^j(Lf_i^*C)=0$（使用 《导出范畴》引理
\ref{derived-lemma-negative-vanishing} 的对偶）。因此
$H^j(Lf_i^*\alpha_i)$ 在 $j>m$ 时为同构，而
$H^m(Lf^*\alpha_i)$ 满射。另一方面，
$Lf_i^*\mathcal{E}_i^\bullet = f_i^*\mathcal{E}_i^\bullet$.
由引理 \ref{cohomology-lemma-strictly-perfect-pullback}，它严格完美。结论得证。
\end{proof}

\begin{lemma}
\label{cohomology-lemma-cone-pseudo-coherent}
设 $(X,\mathcal{O}_X)$ 为环化空间，$m\in\mathbf{Z}$，并设
$(K,L,M,f,g,h)$ 为 $D(\mathcal{O}_X)$ 中的特异三角。
\begin{enumerate}
\item 若 $K$ 是 $(m+1)$-伪凝聚的，而 $L$ 是 $m$-伪凝聚的，
则 $M$ 是 $m$-伪凝聚的。
\item 若 $K$ 与 $M$ 都是 $m$-伪凝聚的，则 $L$ 是 $m$-伪凝聚的。
\item 若 $L$ 是 $(m+1)$-伪凝聚的，而 $M$ 是 $m$-伪凝聚的，
则 $K$ 是 $(m+1)$-伪凝聚的。
\end{enumerate}
\end{lemma}

\begin{proof}
证明 (1)。选取开覆盖 $X=\bigcup U_i$，以及
$D(\mathcal{O}_{U_i})$ 中的态射
$\alpha_i:\mathcal{K}_i^\bullet\to K|_{U_i}$，使
$\mathcal{K}_i^\bullet$ 严格完美，且 $H^j(\alpha_i)$ 在
$j>m+1$ 时为同构、在 $j=m+1$ 时满射。可以把
$\mathcal{K}_i^\bullet$ 换成
$\sigma_{\geq m+1}\mathcal{K}_i^\bullet$，故可以假设当 $j<m+1$ 时
$\mathcal{K}_i^j=0$。细化开覆盖后，可以选取
$D(\mathcal{O}_{U_i})$ 中的态射
$\beta_i:\mathcal{L}_i^\bullet\to L|_{U_i}$，使
$\mathcal{L}_i^\bullet$ 严格完美，且 $H^j(\beta)$ 在 $j>m$ 时
为同构、在 $j=m$ 时满射。由引理
\ref{cohomology-lemma-lift-through-quasi-isomorphism}，再次细化覆盖后，可以找到
复形态射 $\gamma_i:\mathcal{K}^\bullet\to\mathcal{L}^\bullet$，
使图表
$$
\xymatrix{
K|_{U_i} \ar[r] & L|_{U_i} \\
\mathcal{K}_i^\bullet \ar[u]^{\alpha_i} \ar[r]^{\gamma_i} &
\mathcal{L}_i^\bullet \ar[u]_{\beta_i}
}
$$
在 $D(\mathcal{O}_{U_i})$ 中交换（这要求用实际复形态射表示
$\alpha_i$、$\beta_i$ 以及 $K|_{U_i}\to L|_{U_i}$；略去一些细节）。
锥 $C(\gamma_i)^\bullet$ 严格完美（引理 \ref{cohomology-lemma-cone}）。
图表的交换性蕴含存在特异三角之间的态射
$$
(\mathcal{K}_i^\bullet, \mathcal{L}_i^\bullet, C(\gamma_i)^\bullet)
\longrightarrow
(K|_{U_i}, L|_{U_i}, M|_{U_i}).
$$
由所诱导的上同调长正合列之间的态射，以及 《同调代数》引理
\ref{homology-lemma-four-lemma} 和 \ref{homology-lemma-five-lemma}，
可知 $C(\gamma_i)^\bullet\to M|_{U_i}$ 在次数 $>m$ 的上同调上诱导
同构，在次数 $m$ 上诱导满射。因此，由引理
\ref{cohomology-lemma-pseudo-coherent-independent-representative}，$M$ 是
$m$-伪凝聚的。

\medskip\noindent
把特异三角旋转后，由 (1) 即得断言 (2) 和 (3)。
\end{proof}

\begin{lemma}
\label{cohomology-lemma-tensor-pseudo-coherent}
设 $(X,\mathcal{O}_X)$ 为环化空间，并设 $K,L$ 为
$D(\mathcal{O}_X)$ 的对象。
\begin{enumerate}
\item 若 $K$ 是 $n$-伪凝聚的，且当 $i>a$ 时有 $H^i(K)=0$；
而 $L$ 是 $m$-伪凝聚的，且当 $j>b$ 时有 $H^j(L)=0$，则
$K\otimes_{\mathcal{O}_X}^\mathbf{L}L$ 是 $t$-伪凝聚的，其中
$t=\max(m+a,n+b)$。
\item 若 $K$ 和 $L$ 都是伪凝聚的，则
$K\otimes_{\mathcal{O}_X}^\mathbf{L}L$ 是伪凝聚的。
\end{enumerate}
\end{lemma}

\begin{proof}
证明 (1)。用开覆盖的各成员替换 $X$ 后，可以假设存在严格完美复形
$\mathcal{K}^\bullet$ 和 $\mathcal{L}^\bullet$，以及态射
$\alpha:\mathcal{K}^\bullet\to K$ 和
$\beta:\mathcal{L}^\bullet\to L$，使得 $H^i(\alpha)$ 在 $i>n$
时为同构、在 $i=n$ 时满射，而 $H^i(\beta)$ 在 $i>m$ 时为同构、
在 $i=m$ 时满射。于是态射
$$
\alpha \otimes^\mathbf{L} \beta :
\text{Tot}(\mathcal{K}^\bullet \otimes_{\mathcal{O}_X} \mathcal{L}^\bullet)
\to K \otimes_{\mathcal{O}_X}^\mathbf{L} L
$$
在 $i>t$ 时的第 $i$ 个上同调层上诱导同构，在 $i=t$ 时诱导满射。
这由 Tor 的谱序列得出（细节从略）。

\medskip\noindent
证明 (2)。先用某个开覆盖的各成员替换 $X$，便可归结到 $K$ 和 $L$
上有界的情形。此时断言立即由 (1) 得出。
\end{proof}

\begin{lemma}
\label{cohomology-lemma-summands-pseudo-coherent}
设 $(X,\mathcal{O}_X)$ 为环化空间，$m\in\mathbf{Z}$。若
$K\oplus L$ 在 $D(\mathcal{O}_X)$ 中是 $m$-伪凝聚的
（相应地，伪凝聚的），则 $K$ 和 $L$ 也如此。
\end{lemma}

\begin{proof}
假设 $K\oplus L$ 是 $m$-伪凝聚的。用开覆盖的各成员替换 $X$ 后，
可以假设 $K\oplus L\in D^-(\mathcal{O}_X)$，从而
$L\in D^-(\mathcal{O}_X)$。注意，存在特异三角
$$
(K \oplus L, K \oplus L, L \oplus L[1]) =
(K, K, 0) \oplus (L, L, L \oplus L[1])
$$
；见 《导出范畴》引理
\ref{derived-lemma-direct-sum-triangles}。由引理
\ref{cohomology-lemma-cone-pseudo-coherent}，$L\oplus L[1]$ 是
$m$-伪凝聚的，故 $L[1]\oplus L[2]$ 也为 $m$-伪凝聚。归纳可知
$L[n]\oplus L[n+1]$ 是 $m$-伪凝聚的。由于 $L$ 上有界，当 $n$
足够大时，$L[n]$ 是 $m$-伪凝聚的。因此，利用特异三角
$$
(L[n], L[n] \oplus L[n - 1], L[n - 1])
$$
反向递推，便得 $L[n-1],L[n-2],\ldots,L$ 都是 $m$-伪凝聚的，
如所需。
\end{proof}

\begin{lemma}
\label{cohomology-lemma-complex-pseudo-coherent-modules}
设 $(X,\mathcal{O}_X)$ 为环化空间，$m\in\mathbf{Z}$。设
$\mathcal{F}^\bullet$ 为（局部）上有界的 $\mathcal{O}_X$-模复形，
并且对所有 $i$，$\mathcal{F}^i$ 都是 $(m-i)$-伪凝聚的。
则 $\mathcal{F}^\bullet$ 是 $m$-伪凝聚的。
\end{lemma}

\begin{proof}
从略。提示：使用引理 \ref{cohomology-lemma-cone-pseudo-coherent} 和截断，
如 《更多代数》引理
\ref{more-algebra-lemma-complex-pseudo-coherent-modules} 的证明所示。
\end{proof}

\begin{lemma}
\label{cohomology-lemma-cohomology-pseudo-coherent}
设 $(X,\mathcal{O}_X)$ 为环化空间，$m\in\mathbf{Z}$，并设
$E$ 为 $D(\mathcal{O}_X)$ 的对象。若 $E$（局部）上有界，并且对所有
$i$，$H^i(E)$ 都是 $(m-i)$-伪凝聚的，则 $E$ 是 $m$-伪凝聚的。
\end{lemma}

\begin{proof}
从略。提示：使用引理 \ref{cohomology-lemma-cone-pseudo-coherent} 和截断，
如 《更多代数》引理
\ref{more-algebra-lemma-cohomology-pseudo-coherent} 的证明所示。
\end{proof}

\begin{lemma}
\label{cohomology-lemma-finite-cohomology}
设 $(X,\mathcal{O}_X)$ 为环化空间，$K$ 为
$D(\mathcal{O}_X)$ 的对象，并设 $m\in\mathbf{Z}$。
\begin{enumerate}
\item 若 $K$ 是 $m$-伪凝聚的，且当 $i>m$ 时有 $H^i(K)=0$，
则 $H^m(K)$ 是有限型 $\mathcal{O}_X$-模。
\item 若 $K$ 是 $m$-伪凝聚的，且当 $i>m+1$ 时有 $H^i(K)=0$，
则 $H^{m+1}(K)$ 是有限表示的 $\mathcal{O}_X$-模。
\end{enumerate}
\end{lemma}

\begin{proof}
证明 (1)。可以在 $X$ 上局部工作。因此，可以假设存在严格完美复形
$\mathcal{E}^\bullet$ 以及态射
$\alpha:\mathcal{E}^\bullet\to K$，它在次数 $>m$ 的上同调上诱导
同构，在次数 $m$ 上诱导满射。只需对 $\mathcal{E}^\bullet$
证明结论。令 $n$ 为使 $\mathcal{E}^n\not=0$ 的最大整数。若
$n=m$，则 $H^m(\mathcal{E}^\bullet)$ 是 $\mathcal{E}^n$ 的商，
结论显然。若 $n>m$，则由于 $H^n(E^\bullet)=0$，
$\mathcal{E}^{n-1}\to\mathcal{E}^n$ 满射。由引理
\ref{cohomology-lemma-local-lift-map}，可以在局部找到该满射的截面，并写成
$\mathcal{E}^{n-1}=\mathcal{E}'\oplus\mathcal{E}^n$。因此，只需对
复形 $(\mathcal{E}')^\bullet$ 证明结论；它与
$\mathcal{E}^\bullet$ 相同，只是次数 $n-1$ 的项换成
$\mathcal{E}'$，次数 $n$ 的项换成 $0$。对 $n$ 作归纳即得结论。

\medskip\noindent
证明 (2)。可以在 $X$ 上局部工作。因此，可以假设存在严格完美复形
$\mathcal{E}^\bullet$ 以及态射
$\alpha:\mathcal{E}^\bullet\to K$，它在次数 $>m$ 的上同调上诱导
同构，在次数 $m$ 上诱导满射。如 (1) 的证明那样，可以归结到当
$i>m+1$ 时 $\mathcal{E}^i=0$ 的情形。于是
$H^{m + 1}(K) \cong H^{m + 1}(\mathcal{E}^\bullet) =
\Coker(\mathcal{E}^m \to \mathcal{E}^{m + 1})$
，这是有限表示的。
\end{proof}

\begin{lemma}
\label{cohomology-lemma-n-pseudo-module}
设 $(X,\mathcal{O}_X)$ 为环化空间，并设 $\mathcal{F}$ 为
$\mathcal{O}_X$-模层。
\begin{enumerate}
\item 把 $\mathcal{F}$ 视为 $D(\mathcal{O}_X)$ 的对象时，它是
$0$-伪凝聚的，当且仅当 $\mathcal{F}$ 是有限型
$\mathcal{O}_X$-模；并且
\item 把 $\mathcal{F}$ 视为 $D(\mathcal{O}_X)$ 的对象时，它是
$(-1)$-伪凝聚的，当且仅当 $\mathcal{F}$ 是有限表示的
$\mathcal{O}_X$-模。
\end{enumerate}
\end{lemma}

\begin{proof}
一个方向的蕴含使用引理 \ref{cohomology-lemma-finite-cohomology}，另一个方向
使用引理 \ref{cohomology-lemma-cohomology-pseudo-coherent}。
\end{proof}





\section{Tor 维数}
\label{cohomology-section-tor}

\noindent
本节将更仔细地考察由平坦模给出的消解。

\begin{definition}
\label{cohomology-definition-tor-amplitude}
令 $(X, \mathcal{O}_X)$ 为环化空间。
令 $E$ 为 $D(\mathcal{O}_X)$ 的对象。
令 $a, b \in \mathbf{Z}$ 且 $a \leq b$。
\begin{enumerate}
\item 称 $E$ 的 {\it Tor 振幅位于 $[a, b]$}，
如果对每个 $\mathcal{O}_X$-模 $\mathcal{F}$ 以及每个 $i \not \in [a, b]$，都有
$H^i(E \otimes_{\mathcal{O}_X}^\mathbf{L} \mathcal{F}) = 0$。
\item 如果 $E$ 的 Tor 振幅位于 $[a, b]$，其中 $a, b$ 可适当选取，
则称其具有{\it 有限 Tor 维数}。
\item 如果存在开覆盖 $X = \bigcup U_i$，使得
对每个 $i$，$E|_{U_i}$ 都具有有限 Tor 维数，
则称 $E$ {\it 局部具有有限 Tor 维数}。
\end{enumerate}
如果把 $\mathcal{O}_X$-模 $\mathcal{F}$ 的 $\mathcal{F}[0]$
视为 $D(\mathcal{O}_X)$ 的对象时，其 Tor 振幅位于 $[-d, 0]$，
则称该模的 {\it Tor 维数 $\leq d$}。
\end{definition}

\noindent
注意，如果定义中的 $E$ 具有有限 Tor 维数，
那么 $E$ 是 $D^b(\mathcal{O}_X)$ 的对象；
在上述定义中取 $\mathcal{F} = \mathcal{O}_X$ 即可看出这一点。

\begin{lemma}
\label{cohomology-lemma-last-one-flat}
令 $(X, \mathcal{O}_X)$ 为环化空间。
令 $\mathcal{E}^\bullet$ 为由平坦
$\mathcal{O}_X$-模组成的上有界复形，其 Tor 振幅位于 $[a, b]$。
则 $\Coker(d_{\mathcal{E}^\bullet}^{a - 1})$ 是平坦
$\mathcal{O}_X$-模。
\end{lemma}

\begin{proof}
由于 $\mathcal{E}^\bullet$ 是由平坦模组成的上有界复形，对任意
$\mathcal{O}_X$-模 $\mathcal{F}$，都有
$\mathcal{E}^\bullet \otimes_{\mathcal{O}_X} \mathcal{F} =
\mathcal{E}^\bullet \otimes_{\mathcal{O}_X}^{\mathbf{L}} \mathcal{F}$。
因此，对每个 $\mathcal{O}_X$-模 $\mathcal{F}$，序列
$$
\mathcal{E}^{a - 2} \otimes_{\mathcal{O}_X} \mathcal{F} \to
\mathcal{E}^{a - 1} \otimes_{\mathcal{O}_X} \mathcal{F} \to
\mathcal{E}^a \otimes_{\mathcal{O}_X} \mathcal{F}
$$
在中间正合。由于
$\mathcal{E}^{a - 2} \to \mathcal{E}^{a - 1} \to \mathcal{E}^a \to
\Coker(d^{a - 1}) \to 0$
是平坦消解，故
$\text{Tor}_1^{\mathcal{O}_X}(\Coker(d^{a - 1}), \mathcal{F}) = 0$
对所有 $\mathcal{O}_X$-模 $\mathcal{F}$ 成立。这意味着
$\Coker(d^{a - 1})$ 平坦，见引理 \ref{cohomology-lemma-flat-tor-zero}。
\end{proof}

\begin{lemma}
\label{cohomology-lemma-tor-amplitude}
令 $(X, \mathcal{O}_X)$ 为环化空间。
令 $E$ 为 $D(\mathcal{O}_X)$ 的对象。
令 $a, b \in \mathbf{Z}$ 且 $a \leq b$。下列条件等价：
\begin{enumerate}
\item $E$ 的 Tor 振幅位于 $[a, b]$。
\item $E$ 可由平坦 $\mathcal{O}_X$-模组成的复形
$\mathcal{E}^\bullet$ 表示，并且
当 $i \not \in [a, b]$ 时 $\mathcal{E}^i = 0$。
\end{enumerate}
\end{lemma}

\begin{proof}
如果 (2) 成立，则可计算
$E \otimes_{\mathcal{O}_X}^\mathbf{L} \mathcal{F} =
\mathcal{E}^\bullet \otimes_{\mathcal{O}_X} \mathcal{F}$，
显然 (1) 成立。

\medskip\noindent
假设 (1) 成立。根据第 \ref{cohomology-section-flat} 节，
可用由平坦 $\mathcal{O}_X$-模组成的上有界复形
$\mathcal{K}^\bullet$ 表示 $E$。
令 $n$ 为满足 $\mathcal{K}^n \not = 0$ 的最大整数。
如果 $n > b$，则由 $H^n(\mathcal{K}^\bullet) = 0$ 可知
$\mathcal{K}^{n - 1} \to \mathcal{K}^n$ 是满射。由于 $\mathcal{K}^n$ 平坦，
$\Ker(\mathcal{K}^{n - 1} \to \mathcal{K}^n)$ 也是平坦的
（模，引理 \ref{modules-lemma-flat-ses}）。
因此可将 $\mathcal{K}^\bullet$ 替换为
$\tau_{\leq n - 1}\mathcal{K}^\bullet$。于是对 $n$ 归纳后，
可约化到如下情形：$K^\bullet$ 是由平坦
$\mathcal{O}_X$-模组成的复形，并且当 $i > b$ 时 $\mathcal{K}^i = 0$。

\medskip\noindent
令 $\mathcal{E}^\bullet = \tau_{\geq a}\mathcal{K}^\bullet$。
除 $\mathcal{E}^a$ 的平坦性以外，其余都很清楚；
这一平坦性立即由引理 \ref{cohomology-lemma-last-one-flat}
和定义得到。
\end{proof}

\begin{lemma}
\label{cohomology-lemma-tor-amplitude-pullback}
令 $f : (X, \mathcal{O}_X) \to (Y, \mathcal{O}_Y)$ 为环化空间的态射。
令 $E$ 为 $D(\mathcal{O}_Y)$ 的对象。
如果 $E$ 的 Tor 振幅位于 $[a, b]$，那么 $Lf^*E$ 的 Tor 振幅也位于
$[a, b]$。
\end{lemma}

\begin{proof}
假设 $E$ 的 Tor 振幅位于 $[a, b]$。由引理 \ref{cohomology-lemma-tor-amplitude}，
可用由平坦 $\mathcal{O}$-模组成的复形
$\mathcal{E}^\bullet$ 表示 $E$，且
当 $i \not \in [a, b]$ 时 $\mathcal{E}^i = 0$。于是
$f^*\mathcal{E}^\bullet$ 表示 $Lf^*E$。
由模，引理 \ref{modules-lemma-base-change-flat}，
各模 $f^*\mathcal{E}^i$ 都是平坦的。
因此由引理 \ref{cohomology-lemma-tor-amplitude}，
$Lf^*E$ 的 Tor 振幅位于 $[a, b]$。
\end{proof}

\begin{lemma}
\label{cohomology-lemma-tor-amplitude-stalk}
令 $(X, \mathcal{O}_X)$ 为环化空间。
令 $E$ 为 $D(\mathcal{O}_X)$ 的对象。
令 $a, b \in \mathbf{Z}$ 且 $a \leq b$。下列条件等价：
\begin{enumerate}
\item $E$ 的 Tor 振幅位于 $[a, b]$。
\item 对每个 $x \in X$，$D(\mathcal{O}_{X, x})$ 的对象 $E_x$
的 Tor 振幅位于 $[a, b]$。
\end{enumerate}
\end{lemma}

\begin{proof}
在 $x$ 处取茎，等同于沿环化空间态射
$(x, \mathcal{O}_{X, x}) \to (X, \mathcal{O}_X)$ 拉回。
因此 (1) $\Rightarrow$ (2) 由引理 \ref{cohomology-lemma-tor-amplitude-pullback} 得出。
反之，注意取茎与张量积可交换
（模，引理 \ref{modules-lemma-stalk-tensor-product}）。因此
$$
(E \otimes_{\mathcal{O}_X}^\mathbf{L} \mathcal{F})_x =
E_x \otimes_{\mathcal{O}_{X, x}}^\mathbf{L} \mathcal{F}_x
$$
另一方面，取茎是正合的，所以
$$
H^i(E \otimes_{\mathcal{O}_X}^\mathbf{L} \mathcal{F})_x =
H^i((E \otimes_{\mathcal{O}_X}^\mathbf{L} \mathcal{F})_x) =
H^i(E_x \otimes_{\mathcal{O}_{X, x}}^\mathbf{L} \mathcal{F}_x)
$$
而要检验
$H^i(E \otimes_{\mathcal{O}_X}^\mathbf{L} \mathcal{F})$ 是否为零，
只需检验其所有茎是否为零
（模，引理 \ref{modules-lemma-abelian}）。故 (2) 推出 (1)。
\end{proof}

\begin{lemma}
\label{cohomology-lemma-cone-tor-amplitude}
令 $(X, \mathcal{O}_X)$ 为环化空间。
令 $(K, L, M, f, g, h)$ 为 $D(\mathcal{O}_X)$ 中的特异
三角。令 $a, b \in \mathbf{Z}$。
\begin{enumerate}
\item 如果 $K$ 的 Tor 振幅位于 $[a + 1, b + 1]$，并且
$L$ 的 Tor 振幅位于 $[a, b]$，那么 $M$ 的
Tor 振幅位于 $[a, b]$。
\item 如果 $K$ 和 $M$ 的 Tor 振幅位于 $[a, b]$，那么
$L$ 的 Tor 振幅位于 $[a, b]$。
\item 如果 $L$ 的 Tor 振幅位于 $[a + 1, b + 1]$，
并且 $M$ 的 Tor 振幅位于 $[a, b]$，那么
$K$ 的 Tor 振幅位于 $[a + 1, b + 1]$。
\end{enumerate}
\end{lemma}

\begin{proof}
略去证明。提示：这直接来自特异三角所伴随的上同调长正合列，
以及
$- \otimes_{\mathcal{O}_X}^{\mathbf{L}} \mathcal{F}$
保持特异三角这一事实。
最容易证明的是 (2)，其余结论由平移得到。
\end{proof}

\begin{lemma}
\label{cohomology-lemma-tensor-tor-amplitude}
令 $(X, \mathcal{O}_X)$ 为环化空间。令 $K, L$ 为
$D(\mathcal{O}_X)$ 的对象。如果 $K$ 的 Tor 振幅位于 $[a, b]$，且
$L$ 的 Tor 振幅位于 $[c, d]$，那么 $K \otimes_{\mathcal{O}_X}^\mathbf{L} L$
的 Tor 振幅位于 $[a + c, b + d]$。
\end{lemma}

\begin{proof}
略去证明。提示：使用 Tor 的谱序列。
\end{proof}

\begin{lemma}
\label{cohomology-lemma-summands-tor-amplitude}
令 $(X, \mathcal{O}_X)$ 为环化空间。令 $a, b \in \mathbf{Z}$。
对于 $D(\mathcal{O}_X)$ 的对象 $K$、$L$，如果 $K \oplus L$ 的 Tor
振幅位于 $[a, b]$，那么 $K$ 和 $L$ 也都如此。
\end{lemma}

\begin{proof}
由 Tor 函子的可加性立即可得。
\end{proof}







\section{完美复形}
\label{cohomology-section-perfect}

\noindent
本节讨论环化空间上完美复形的
性质。

\begin{definition}
\label{cohomology-definition-perfect}
令 $(X, \mathcal{O}_X)$ 为环化空间。
令 $\mathcal{E}^\bullet$ 为 $\mathcal{O}_X$-模复形。
如果存在开覆盖 $X = \bigcup U_i$，使得对每个 $i$
都存在复形态射
$\mathcal{E}_i^\bullet \to \mathcal{E}^\bullet|_{U_i}$，
它是拟同构，且 $\mathcal{E}_i^\bullet$
是 $\mathcal{O}_{U_i}$-模的严格完美复形，
则称 $\mathcal{E}^\bullet$ 是{\it 完美的}。
如果 $D(\mathcal{O}_X)$ 的对象 $E$ 可由
完美的 $\mathcal{O}_X$-模复形表示，则称其是{\it 完美的}。
\end{definition}

\noindent
如果 $X$ 拟紧，那么 $D(\mathcal{O}_X)$ 的完美对象
属于 $D^b(\mathcal{O}_X)$。但当
$X$ 非拟紧时，未必如此。

\begin{lemma}
\label{cohomology-lemma-perfect-independent-representative}
令 $(X, \mathcal{O}_X)$ 为环化空间。
令 $E$ 为 $D(\mathcal{O}_X)$ 的对象。
\begin{enumerate}
\item 如果存在开覆盖 $X = \bigcup U_i$，以及 $U_i$ 上的
严格完美复形 $\mathcal{E}_i^\bullet$，使得
$\mathcal{E}_i^\bullet$ 在 $D(\mathcal{O}_{U_i})$ 中表示 $E|_{U_i}$，
那么 $E$ 是完美的。
\item 如果 $E$ 是完美的，那么表示 $E$ 的任意复形都是完美的。
\end{enumerate}
\end{lemma}

\begin{proof}
与引理 \ref{cohomology-lemma-pseudo-coherent-independent-representative}
的证明相同。
\end{proof}

\begin{lemma}
\label{cohomology-lemma-perfect-on-locally-ringed}
令 $(X, \mathcal{O}_X)$ 为环化空间。令 $E$ 为
$D(\mathcal{O}_X)$ 的对象。假设所有茎 $\mathcal{O}_{X, x}$
都是局部环。则下列条件等价：
\begin{enumerate}
\item $E$ 是完美的；
\item 存在开覆盖 $X = \bigcup U_i$，使得
$E|_{U_i}$ 可由有限局部自由
$\mathcal{O}_{U_i}$-模组成的有限复形表示；以及
\item 存在开覆盖 $X = \bigcup U_i$，使得
$E|_{U_i}$ 可由有限自由
$\mathcal{O}_{U_i}$-模组成的有限复形表示。
\end{enumerate}
\end{lemma}

\begin{proof}
这由引理 \ref{cohomology-lemma-perfect-independent-representative}
以及如下事实得出：在 $X$ 上，有限自由模的每个直和项
都是有限局部自由的。见模，引理
\ref{modules-lemma-direct-summand-of-locally-free-is-locally-free}。
\end{proof}

\begin{lemma}
\label{cohomology-lemma-perfect-precise}
令 $(X, \mathcal{O}_X)$ 为环化空间。
令 $E$ 为 $D(\mathcal{O}_X)$ 的对象。
令 $a \leq b$ 为整数。如果 $E$ 的 Tor 振幅位于 $[a, b]$
并且是 $(a - 1)$-伪凝聚的，那么 $E$ 是完美的。
\end{lemma}

\begin{proof}
用开覆盖的成员替换 $X$ 后，可假设存在
严格完美复形 $\mathcal{E}^\bullet$ 和态射
$\alpha : \mathcal{E}^\bullet \to E$，使得 $H^i(\alpha)$
对 $i \geq a$ 是同构。我们可以并且确实将 $\mathcal{E}^\bullet$ 替换为
$\sigma_{\geq a - 1}\mathcal{E}^\bullet$。选取特异三角
$$
\mathcal{E}^\bullet \to E \to C \to \mathcal{E}^\bullet[1]
$$
由 $E$ 和 $\mathcal{E}^\bullet$ 的上同调层消失性
以及关于 $\alpha$ 的假设，得到 $C \cong \mathcal{K}[2 - a]$，其中
$\mathcal{K} = \Ker(\mathcal{E}^{a - 1} \to \mathcal{E}^a)$。
令 $\mathcal{F}$ 为 $\mathcal{O}_X$-模。
应用 $- \otimes_{\mathcal{O}_X}^\mathbf{L} \mathcal{F}$ 后，
$E$ 的 Tor 振幅位于 $[a, b]$ 这一假设
推出 $\mathcal{K} \otimes_{\mathcal{O}_X} \mathcal{F} \to
\mathcal{E}^{a - 1} \otimes_{\mathcal{O}_X} \mathcal{F}$ 的像等于
$\Ker(\mathcal{E}^{a - 1} \otimes_{\mathcal{O}_X} \mathcal{F}
\to \mathcal{E}^a \otimes_{\mathcal{O}_X} \mathcal{F})$。
由此 $\text{Tor}_1^{\mathcal{O}_X}(\mathcal{E}', \mathcal{F}) = 0$，
其中 $\mathcal{E}' = \Coker(\mathcal{E}^{a - 1} \to \mathcal{E}^a)$。
因此 $\mathcal{E}'$ 平坦（引理 \ref{cohomology-lemma-flat-tor-zero}）。
于是由模，引理 \ref{modules-lemma-flat-locally-finite-presentation}，
$\mathcal{E}'$ 局部为某个有限自由模的直和项。
故复形
$$
\mathcal{E}' \to \mathcal{E}^{a + 1} \to \ldots \to \mathcal{E}^b
$$
局部拟同构于 $E$，从而 $E$ 是完美的。
\end{proof}

\begin{lemma}
\label{cohomology-lemma-perfect}
令 $(X, \mathcal{O}_X)$ 为环化空间。
令 $E$ 为 $D(\mathcal{O}_X)$ 的对象。
下列条件等价：
\begin{enumerate}
\item $E$ 是完美的；以及
\item $E$ 是伪凝聚的，并且局部具有有限 Tor 维数。
\end{enumerate}
\end{lemma}

\begin{proof}
假设 (1)。按定义，这意味着存在开覆盖
$X = \bigcup U_i$，使得 $E|_{U_i}$ 可由
严格完美复形表示。因此，由引理 \ref{cohomology-lemma-pseudo-coherent-independent-representative}，$E$ 是伪凝聚的（即
对所有 $m$ 都是 $m$-伪凝聚的）。
此外，有限自由模的直和项是平坦的，所以
由引理 \ref{cohomology-lemma-tor-amplitude}，
$E|_{U_i}$ 具有有限 Tor 维数。因此 (2) 成立。

\medskip\noindent
假设 (2)。用开覆盖的成员替换 $X$ 后，
可假设存在整数 $a \leq b$，使得 $E$
的 Tor 振幅位于 $[a, b]$。由于 $E$ 对所有 $m$
都是 $m$-伪凝聚的，应用引理 \ref{cohomology-lemma-perfect-precise} 即得结论。
\end{proof}

\begin{lemma}
\label{cohomology-lemma-perfect-pullback}
令 $f : (X, \mathcal{O}_X) \to (Y, \mathcal{O}_Y)$ 为环化
空间的态射。令 $E$ 为 $D(\mathcal{O}_Y)$ 的对象。如果 $E$ 在
$D(\mathcal{O}_Y)$ 中完美，那么 $Lf^*E$ 在 $D(\mathcal{O}_X)$ 中完美。
\end{lemma}

\begin{proof}
这由引理 \ref{cohomology-lemma-perfect}、
\ref{cohomology-lemma-tor-amplitude-pullback} 和
\ref{cohomology-lemma-pseudo-coherent-pullback} 得出。
（另一种证明是照搬引理
\ref{cohomology-lemma-pseudo-coherent-pullback} 的证明。）
\end{proof}

\begin{lemma}
\label{cohomology-lemma-two-out-of-three-perfect}
令 $(X, \mathcal{O}_X)$ 为环化空间。令 $(K, L, M, f, g, h)$
为 $D(\mathcal{O}_X)$ 中的特异三角。如果 $K, L, M$
三者中有两个完美，那么第三个也完美。
\end{lemma}

\begin{proof}
第一种证明：合并使用引理 \ref{cohomology-lemma-perfect}、\ref{cohomology-lemma-cone-pseudo-coherent} 和
\ref{cohomology-lemma-cone-tor-amplitude}。
第二种证明（略述）：设 $K$ 和 $L$ 完美。用开覆盖的成员
替换 $X$ 后，可假设 $K$ 和 $L$
分别由严格完美复形 $\mathcal{K}^\bullet$
和 $\mathcal{L}^\bullet$ 表示。再次用开覆盖的成员
替换 $X$ 后，可假设态射 $K \to L$ 由
复形态射 $\alpha : \mathcal{K}^\bullet \to \mathcal{L}^\bullet$ 给出，
见引理 \ref{cohomology-lemma-local-actual}。
于是 $M$ 同构于 $\alpha$ 的锥，而由引理 \ref{cohomology-lemma-cone}，
该锥是严格完美的。
\end{proof}

\begin{lemma}
\label{cohomology-lemma-tensor-perfect}
令 $(X, \mathcal{O}_X)$ 为环化空间。
如果 $K, L$ 是 $D(\mathcal{O}_X)$ 的完美对象，那么
$K \otimes_{\mathcal{O}_X}^\mathbf{L} L$ 也是完美对象。
\end{lemma}

\begin{proof}
这由引理 \ref{cohomology-lemma-perfect}、\ref{cohomology-lemma-tensor-pseudo-coherent} 和
\ref{cohomology-lemma-tensor-tor-amplitude} 得出。
\end{proof}

\begin{lemma}
\label{cohomology-lemma-summands-perfect}
令 $(X, \mathcal{O}_X)$ 为环化空间。
如果 $K \oplus L$ 是 $D(\mathcal{O}_X)$ 的完美对象，那么
$K$ 和 $L$ 也都是完美对象。
\end{lemma}

\begin{proof}
这由引理 \ref{cohomology-lemma-perfect}、\ref{cohomology-lemma-summands-pseudo-coherent} 和
\ref{cohomology-lemma-summands-tor-amplitude} 得出。
\end{proof}

\begin{lemma}
\label{cohomology-lemma-pushforward-perfect}
令 $(X, \mathcal{O}_X)$ 为环化空间。令 $j : U \to X$ 为
开子空间。令 $E$ 为 $D(\mathcal{O}_U)$ 的完美对象，其上同调
层支撑在闭子集 $T \subset U$ 上，并且 $j(T)$
在 $X$ 中闭。那么 $Rj_*E$ 是 $D(\mathcal{O}_X)$ 的完美对象。
\end{lemma}

\begin{proof}
完美复形这一性质在 $X$ 上是局部的。因此只需检验
$Rj_*E$ 限制到 $U$ 和 $V = X \setminus j(T)$ 后都是完美的。
我们有 $Rj_*E|_U = E$，它是完美的。又有
$Rj_*E|_V = 0$，因为 $E|_{U \setminus T} = 0$。
\end{proof}

\begin{lemma}
\label{cohomology-lemma-perfect-max-open-coh-loc-free}
令 $(X, \mathcal{O}_X)$ 为环化空间。令 $D(\mathcal{O}_X)$ 中的 $E$
为完美对象。假设所有茎 $\mathcal{O}_{X, x}$ 都是局部环。
那么集合
$$
U =
\{x \in X \mid
H^i(E)_x\text{ 是有限自由 }
\mathcal{O}_{X, x}\text{-模，对所有 }i\in \mathbf{Z}\}
$$
在 $X$ 中开，并且是满足如下性质的极大开集 $U \subset X$：
对所有 $i \in \mathbf{Z}$，$H^i(E)|_U$ 都是有限局部自由的。
\end{lemma}

\begin{proof}
注意，如果 $V \subset X$ 是某个开集，使得对所有 $i \in \mathbf{Z}$，
$H^i(E)|_V$ 都是有限局部自由的，那么 $V \subset U$。
令 $x \in U$。我们将证明 $x$ 的某个开邻域包含于 $U$，
并且对所有 $i$，$H^i(E)$ 在该邻域上都是有限局部自由的。
这将完成证明。
在证明过程中，我们可以（有限多次）用 $x$ 的开邻域
替换 $X$。
因此可假设 $E$ 由严格完美复形
$\mathcal{E}^\bullet$ 表示。设当 $i \not \in [a, b]$ 时
$\mathcal{E}^i = 0$。我们将对 $b - a$ 归纳证明结论。模
$H^b(E) = \Coker(d^{b - 1} : \mathcal{E}^{b - 1} \to \mathcal{E}^b)$
是有限表示的。由于 $H^b(E)_x$ 有限自由，
由模，引理 \ref{modules-lemma-finite-presentation-stalk-free}，
$H^b(E)$ 在 $x$ 的某个开邻域上有限自由。
因此，用一个（可能更小的）开邻域替换 $X$ 后，
可假设有直和分解
$\mathcal{E}^b = \Im(d^{b - 1}) \oplus H^b(E)$，并且
$H^b(E)$ 有限自由，见引理 \ref{cohomology-lemma-local-lift-map}。
再次使用同一论证，可假设
$\mathcal{E}^{b - 1} = \Ker(d^{b - 1}) \oplus \Im(d^{b - 1})$。
复形
$\mathcal{E}^a \to \ldots \to \mathcal{E}^{b - 2} \to \Ker(d^{b - 1})$
是严格完美复形，表示完美
对象 $E'$，且当 $i \not = b$ 时 $H^i(E) = H^i(E')$。
因此由归纳假设得到结论。
\end{proof}







\section{对偶}
\label{cohomology-section-duals}

\noindent
本节刻画复形范畴和导出范畴中的
可对偶对象。
特别地，我们将看到，$D(\mathcal{O}_X)$ 的对象
有对偶，当且仅当它是完美的（这由例 \ref{cohomology-example-dual-derived} 和引理 \ref{cohomology-lemma-left-dual-derived} 得出）。

\begin{lemma}
\label{cohomology-lemma-symmetric-monoidal-cat-complexes}
令 $(X, \mathcal{O}_X)$ 为环化空间。以
$\mathcal{F}^\bullet \otimes \mathcal{G}^\bullet =
\text{Tot}(\mathcal{F}^\bullet \otimes_{\mathcal{O}_X} \mathcal{G}^\bullet)$
定义张量积后，$\mathcal{O}_X$-模复形的范畴
是对称幺半范畴（符号规则见
更多代数，第 \ref{more-algebra-section-sign-rules} 节）。
\end{lemma}

\begin{proof}
略去证明。提示：取复形 $\mathbf{1}$ 为单位对象，
它在次数 $0$ 处为 $\mathcal{O}_X$，在其余次数处为零，并带有显然的同构
$\text{Tot}(\mathbf{1} \otimes_{\mathcal{O}_X} \mathcal{G}^\bullet) =
\mathcal{G}^\bullet$ 和
$\text{Tot}(\mathcal{F}^\bullet \otimes_{\mathcal{O}_X} \mathbf{1}) =
\mathcal{F}^\bullet$。
为证明此引理，需要检验若干图的交换性，见
范畴，定义
\ref{categories-definition-monoidal-category} 和
\ref{categories-definition-symmetric-monoidal-category}。
每种情形下的验证都很直接。
\end{proof}

\begin{example}
\label{cohomology-example-dual}
令 $(X, \mathcal{O}_X)$ 为环化空间。令 $\mathcal{F}^\bullet$
为 $\mathcal{O}_X$-模的局部有界复形，使得每个
$\mathcal{F}^n$ 都局部为某个有限自由
$\mathcal{O}_X$-模的直和项。换言之，存在开覆盖
$X = \bigcup U_i$，使得 $\mathcal{F}^\bullet|_{U_i}$ 是严格
完美复形。考虑第 \ref{cohomology-section-hom-complexes} 节中的复形
$$
\mathcal{G}^\bullet = \SheafHom^\bullet(\mathcal{F}^\bullet, \mathcal{O}_X)
$$
令
$$
\eta :
\mathcal{O}_X
\to
\text{Tot}(\mathcal{F}^\bullet \otimes_{\mathcal{O}_X} \mathcal{G}^\bullet)
\quad\text{且}\quad
\epsilon :
\text{Tot}(\mathcal{G}^\bullet \otimes_{\mathcal{O}_X} \mathcal{F}^\bullet)
\to
\mathcal{O}_X
$$
分别为 $\eta = \sum \eta_n$ 和 $\epsilon = \sum \epsilon_n$，
其中 $\eta_n : \mathcal{O}_X \to
\mathcal{F}^n \otimes_{\mathcal{O}_X} \mathcal{G}^{-n}$，
并且
$\epsilon_n : \mathcal{G}^{-n} \otimes_{\mathcal{O}_X} \mathcal{F}^n
\to \mathcal{O}_X$ 如模，例 \ref{modules-example-dual} 中所定义。
那么按范畴，定义 \ref{categories-definition-dual}，
$\mathcal{G}^\bullet, \eta, \epsilon$
是 $\mathcal{F}^\bullet$ 的左对偶。
我们略去下列等式的验证：
$(1 \otimes \epsilon) \circ (\eta \otimes 1) = \text{id}_{\mathcal{F}^\bullet}$
以及
$(\epsilon \otimes 1) \circ (1 \otimes \eta) =
\text{id}_{\mathcal{G}^\bullet}$。请比较
更多代数，引理 \ref{more-algebra-lemma-left-dual-complex}。
\end{example}

\begin{lemma}
\label{cohomology-lemma-left-dual-complex}
令 $(X, \mathcal{O}_X)$ 为环化空间。令 $\mathcal{F}^\bullet$
为 $\mathcal{O}_X$-模复形。如果 $\mathcal{F}^\bullet$
在 $\mathcal{O}_X$-模复形的幺半范畴中
有左对偶
（范畴，定义 \ref{categories-definition-dual}），
那么 $\mathcal{F}^\bullet$ 是局部有界复形，其各项都
局部为有限自由 $\mathcal{O}_X$-模的直和项，
并且该左对偶就是例 \ref{cohomology-example-dual} 中的构造。
\end{lemma}

\begin{proof}
由左对偶的唯一性
（范畴，注记 \ref{categories-remark-left-dual-adjoint}），
只要证明 $\mathcal{F}^\bullet$ 具有所述性质，便得到最后一个断言。
令 $\mathcal{G}^\bullet, \eta, \epsilon$ 为左对偶。
写成 $\eta = \sum \eta_n$ 和 $\epsilon = \sum \epsilon_n$，
其中 $\eta_n : \mathcal{O}_X \to
\mathcal{F}^n \otimes_{\mathcal{O}_X} \mathcal{G}^{-n}$，
并且
$\epsilon_n : \mathcal{G}^{-n} \otimes_{\mathcal{O}_X} \mathcal{F}^n
\to \mathcal{O}_X$。由于
$(1 \otimes \epsilon) \circ (\eta \otimes 1) = \text{id}_{\mathcal{F}^\bullet}$
以及
$(\epsilon \otimes 1) \circ (1 \otimes \eta) = \text{id}_{\mathcal{G}^\bullet}$，
由范畴，定义 \ref{categories-definition-dual}，立即可见
$(1 \otimes \epsilon_n) \circ (\eta_n \otimes 1) = \text{id}_{\mathcal{F}^n}$
以及
$(\epsilon_n \otimes 1) \circ (1 \otimes \eta_n) =
\text{id}_{\mathcal{G}^{-n}}$。
因此由模，引理 \ref{modules-lemma-left-dual-module}，
$\mathcal{F}^n$ 局部为某个有限自由
$\mathcal{O}_X$-模的直和项。
由于和 $\eta = \sum \eta_n$ 局部有限，可知
$\mathcal{F}^\bullet$ 局部有界。
\end{proof}

\begin{lemma}
\label{cohomology-lemma-internal-hom-evaluate-isom}
令 $(X, \mathcal{O}_X)$ 为环化空间。令 $K, L, M \in D(\mathcal{O}_X)$。
如果 $K$ 完美，那么引理 \ref{cohomology-lemma-internal-hom-evaluate} 中的态射
$$
R\SheafHom(L, M) \otimes_{\mathcal{O}_X}^\mathbf{L} K
\longrightarrow
R\SheafHom(R\SheafHom(K, L), M)
$$
是同构。
\end{lemma}

\begin{proof}
由于该态射整体定义，并且等式的左、右两端的构造都与局部化
交换（见引理 \ref{cohomology-lemma-restriction-RHom-to-U}），为证明断言可在
$X$ 上局部工作。因此可假设 $K$ 由严格
完美复形 $\mathcal{E}^\bullet$ 表示。

\medskip\noindent
如果 $K_1 \to K_2 \to K_3$ 是 $D(\mathcal{O}_X)$ 中的特异三角，
那么得到特异三角
$$
R\SheafHom(L, M) \otimes_{\mathcal{O}_X}^\mathbf{L} K_1 \to
R\SheafHom(L, M) \otimes_{\mathcal{O}_X}^\mathbf{L} K_2 \to
R\SheafHom(L, M) \otimes_{\mathcal{O}_X}^\mathbf{L} K_3
$$
以及
$$
R\SheafHom(R\SheafHom(K_1, L), M) \to
R\SheafHom(R\SheafHom(K_2, L), M)
R\SheafHom(R\SheafHom(K_3, L), M)
$$
见第 \ref{cohomology-section-flat} 节和引理 \ref{cohomology-lemma-RHom-triangulated}。引理 \ref{cohomology-lemma-internal-hom-evaluate} 的箭头对 $K$ 具有函子性，
因而得到这两个特异三角之间的态射。
于是，如果结论对 $K_1$ 和 $K_3$ 成立，则由导出范畴，引理
\ref{derived-lemma-third-isomorphism-triangle}，
结论对 $K_2$ 也成立。

\medskip\noindent
结合上述观察与朴素截断的特异三角
$$
\sigma_{\geq n}\mathcal{E}^\bullet \to \mathcal{E}^\bullet \to
\sigma_{\leq n - 1}\mathcal{E}^\bullet
$$
可约化到 $K$ 是放在某个次数上的
有限自由 $\mathcal{O}_X$-模的直和项这一情形。
再利用该问题与直和的显然相容性
（类似于上文所述）以及移位，可进一步约化到
$K = \mathcal{O}_X^{\oplus n}$ 的情形，其中 $n$ 是某个整数。
此时结论显然。
\end{proof}

\begin{lemma}
\label{cohomology-lemma-dual-perfect-complex}
令 $(X, \mathcal{O}_X)$ 为环化空间。令 $K$ 为
$D(\mathcal{O}_X)$ 的完美对象。那么
$K^\vee = R\SheafHom(K, \mathcal{O}_X)$ 也是完美对象，
且 $(K^\vee)^\vee \cong K$。存在函子性同构
$$
M \otimes^\mathbf{L}_{\mathcal{O}_X} K^\vee = R\SheafHom(K, M)
$$
以及
$$
H^0(X, M \otimes^\mathbf{L}_{\mathcal{O}_X} K^\vee) =
\Hom_{D(\mathcal{O}_X)}(K, M)
$$
其中 $M$ 属于 $D(\mathcal{O}_X)$。
\end{lemma}

\begin{proof}
由引理 \ref{cohomology-lemma-internal-hom-evaluate}，存在典范态射
$$
K = R\SheafHom(\mathcal{O}_X, \mathcal{O}_X)
\otimes_{\mathcal{O}_X}^\mathbf{L} K \longrightarrow
R\SheafHom(R\SheafHom(K, \mathcal{O}_X), \mathcal{O}_X) =
(K^\vee)^\vee
$$
由引理 \ref{cohomology-lemma-internal-hom-evaluate-isom}，它是同构。
为检验其余断言，我们将不再特别说明地使用如下事实：
内 Hom 的构造与到开集的限制可交换
（引理 \ref{cohomology-lemma-restriction-RHom-to-U}）。
可在 $X$ 上局部检验 $K^\vee$ 是完美的。
由引理 \ref{cohomology-lemma-dual}，
为看出最后一个断言，只需检验态射
(\ref{cohomology-equation-eval})
$$
M \otimes^\mathbf{L}_{\mathcal{O}_X} K^\vee
\longrightarrow
R\SheafHom(K, M)
$$
是同构。这也可在 $X$ 上局部检验。
因此，只需在 $K$ 由严格完美复形表示时，
证明这两个断言。

\medskip\noindent
假设 $K$ 由严格完美复形
$\mathcal{E}^\bullet$ 表示。由引理 \ref{cohomology-lemma-Rhom-strictly-perfect} 可知，
$K^\vee$ 由如下复形表示：其次数 $n$ 的项为
$(\mathcal{E}^{-n})^\vee =
\SheafHom_{\mathcal{O}_X}(\mathcal{E}^{-n}, \mathcal{O}_X)$。
由于 $\mathcal{E}^{-n}$ 是某个有限自由
$\mathcal{O}_X$-模的直和项，$(\mathcal{E}^{-n})^\vee$ 亦然。
因此 $K^\vee$ 也由严格完美复形表示，
从而可见 $K^\vee$ 完美。
为看出 (\ref{cohomology-equation-eval}) 是同构，用复形
$\mathcal{F}^\bullet$ 表示 $M$。
由引理 \ref{cohomology-lemma-Rhom-strictly-perfect}，复形
$R\SheafHom(K, M)$ 由以下各项组成的复形表示：
$$
\bigoplus\nolimits_{n = p + q}
\SheafHom_{\mathcal{O}_X}(\mathcal{E}^{-q}, \mathcal{F}^p)
$$
另一方面，对象 $M \otimes^\mathbf{L}_{\mathcal{O}_X} K^\vee$
由以下各项组成的复形表示：
$$
\bigoplus\nolimits_{n = p + q}
\mathcal{F}^p \otimes_{\mathcal{O}_X} (\mathcal{E}^{-q})^\vee
$$
因此，(\ref{cohomology-equation-eval}) 是同构这一断言
约化为如下典范态射
$$
\mathcal{F}
\otimes_{\mathcal{O}_X}
\SheafHom_{\mathcal{O}_X}(\mathcal{E}, \mathcal{O}_X)
\longrightarrow
\SheafHom_{\mathcal{O}_X}(\mathcal{E}, \mathcal{F})
$$
在 $\mathcal{E}$ 是某个有限自由 $\mathcal{O}_X$-模的直和项、
且 $\mathcal{F}$ 是任意 $\mathcal{O}_X$-模时为同构这一断言。
这立即由 $\mathcal{E}$ 有限自由时的相应断言得出。
\end{proof}


\begin{lemma}
\label{cohomology-lemma-symmetric-monoidal-derived}
令 $(X, \mathcal{O}_X)$ 为环化空间。导出范畴
$D(\mathcal{O}_X)$ 以导出张量积为张量积，并带有通常的结合
约束和交换约束，从而是对称幺半范畴（符号规则见
更多代数，第 \ref{more-algebra-section-sign-rules} 节）。
\end{lemma}

\begin{proof}
略去证明。可与引理 \ref{cohomology-lemma-symmetric-monoidal-cat-complexes} 比较。
\end{proof}

\begin{example}
\label{cohomology-example-dual-derived}
令 $(X, \mathcal{O}_X)$ 为环化空间。令 $K$ 为
$D(\mathcal{O}_X)$ 的完美对象。按引理
\ref{cohomology-lemma-dual-perfect-complex}，令
$K^\vee = R\SheafHom(K, \mathcal{O}_X)$。
则态射
$$
K \otimes_{\mathcal{O}_X}^\mathbf{L} K^\vee \longrightarrow R\SheafHom(K, K)
$$
是同构（由该引理）。记
$$
\eta :
\mathcal{O}_X
\longrightarrow
K \otimes_{\mathcal{O}_X}^\mathbf{L} K^\vee
$$
为如下态射：它将 $1$ 送到在上述同构下
对应于 $\text{id}_K$ 的截面。
记
$$
\epsilon : 
K^\vee
\otimes_{\mathcal{O}_X}^\mathbf{L} K
\longrightarrow
\mathcal{O}_X
$$
为求值态射（例如可用引理 \ref{cohomology-lemma-internal-hom-composition} 构造）。那么
按范畴，定义 \ref{categories-definition-dual}，
$K^\vee, \eta, \epsilon$ 是 $K$ 的左对偶。
我们略去下列等式的验证：
$(1 \otimes \epsilon) \circ (\eta \otimes 1) = \text{id}_K$
以及
$(\epsilon \otimes 1) \circ (1 \otimes \eta) =
\text{id}_{K^\vee}$。
\end{example}

\begin{lemma}
\label{cohomology-lemma-left-dual-derived}
令 $(X, \mathcal{O}_X)$ 为环化空间。令 $M$ 为
$D(\mathcal{O}_X)$ 的对象。如果 $M$ 在幺半范畴
$D(\mathcal{O}_X)$ 中有左对偶（范畴，定义
\ref{categories-definition-dual}），那么 $M$ 完美，并且该左对偶
就是例 \ref{cohomology-example-dual-derived} 中的构造。
\end{lemma}

\begin{proof}
令 $x \in X$。只需找到 $x$ 的开邻域 $U$，
使得 $M$ 限制到 $U$ 后为完美复形。因此在证明过程中，
可以（有限多次）用 $x$ 的开邻域替换 $X$。
令 $N, \eta, \epsilon$ 为左对偶。

\medskip\noindent
我们将多次使用如下论证。任取
表示 $M$ 的 $\mathcal{O}_X$-模复形
$\mathcal{M}^\bullet$。选取表示 $N$ 的 K-平坦复形
$\mathcal{N}^\bullet$，其各项都是平坦
$\mathcal{O}_X$-模，见引理 \ref{cohomology-lemma-K-flat-resolution}。
考虑态射
$$
\eta : \mathcal{O}_X \to
\text{Tot}(\mathcal{M}^\bullet \otimes_{\mathcal{O}_X} \mathcal{N}^\bullet)
$$
缩小 $X$ 后，可找到整数 $N$，以及对
$i = 1, \ldots, N$ 的整数 $n_i \in \mathbf{Z}$ 和
$\mathcal{M}^{n_i}$、$\mathcal{N}^{-n_i}$ 的截面
$f_i$、$g_i$，使得
$$
\eta(1) = \sum\nolimits_i f_i \otimes g_i
$$
令 $\mathcal{K}^\bullet \subset \mathcal{M}^\bullet$ 为任意
包含各截面 $f_i$（$i = 1, \ldots, N$）的
$\mathcal{O}_X$-模子复形。
由于各模 $\mathcal{N}^n$ 平坦，有
$\text{Tot}(\mathcal{K}^\bullet \otimes_{\mathcal{O}_X} \mathcal{N}^\bullet)
\subset
\text{Tot}(\mathcal{M}^\bullet \otimes_{\mathcal{O}_X} \mathcal{N}^\bullet)$，
故 $\eta$ 经由
$$
\tilde \eta :
\mathcal{O}_X \to
\text{Tot}(\mathcal{K}^\bullet \otimes_{\mathcal{O}_X} \mathcal{N}^\bullet)
$$
分解。记由 $\mathcal{K}^\bullet$ 表示的
$D(\mathcal{O}_X)$ 对象为 $K$，得到交换图
$$
\xymatrix{
M \ar[rr]_-{\eta \otimes 1} \ar[rrd]_{\tilde \eta \otimes 1} & &
M \otimes^\mathbf{L} N \otimes^\mathbf{L} M
\ar[r]_-{1 \otimes \epsilon} &
M \\
& &
K \otimes^\mathbf{L} N \otimes^\mathbf{L} M
\ar[u] \ar[r]^-{1 \otimes \epsilon} &
K \ar[u]
}
$$
由于上行的复合是 $M$ 上的恒等态射，
可知 $M$ 是 $D(\mathcal{O}_X)$ 中 $K$ 的直和项。

\medskip\noindent
第一次应用上述论证时，可选择子复形
$\mathcal{K}^\bullet = \sigma_{\geq a} \tau_{\leq b}\mathcal{M}^\bullet$，
其中对 $i = 1, \ldots, N$ 有 $a < n_i < b$。因此
$M$ 是 $D(\mathcal{O}_X)$ 中某个有界复形的直和项，
从而可假设 $M$ 属于 $D^b(\mathcal{O}_X)$。（回忆上述过程
涉及缩小 $X$。）

\medskip\noindent
由于 $M$ 属于 $D^b(\mathcal{O}_X)$，可选择
$\mathcal{M}^\bullet$ 为由平坦模组成的上有界复形
（由模，引理 \ref{modules-lemma-module-quotient-flat} 和
导出范畴，引理 \ref{derived-lemma-subcategory-left-resolution}）。
于是，在上述论证中可选择
$\mathcal{K}^\bullet = \sigma_{\geq a}\mathcal{M}^\bullet$，
其中对 $i = 1, \ldots, N$ 有 $a < n_i$。
由此可假设 $M$ 是
$D(\mathcal{O}_X)$ 中某个由平坦模组成的有界复形的直和项。
特别地，$M$ 的 Tor 振幅有限。

\medskip\noindent
设 $M$ 的 Tor 振幅位于 $[a, b]$。假设 $M$ 是 $m$-伪凝聚的，
我们将证明（缩小 $X$ 后）可假设 $M$
是 $(m - 1)$-伪凝聚的。结合引理 \ref{cohomology-lemma-perfect-precise} 以及 $M$
无论如何都是 $(b + 1)$-伪凝聚的这一事实，这将完成证明。
缩小 $X$ 后，可假设存在严格完美
复形 $\mathcal{E}^\bullet$ 以及 $D(\mathcal{O}_X)$ 中的态射
$\alpha : \mathcal{E}^\bullet \to M$，使得 $H^i(\alpha)$ 在
$i > m$ 时为同构，在 $i = m$ 时为满射。可以并且确实假设
当 $i < m$ 时 $\mathcal{E}^i = 0$。选取特异三角
$$
\mathcal{E}^\bullet \to M \to L \to \mathcal{E}^\bullet[1]
$$
注意，当 $i \geq m$ 时 $H^i(L) = 0$。因此可用满足
当 $i \geq m$ 时 $\mathcal{L}^i = 0$ 的复形
$\mathcal{L}^\bullet$ 表示 $L$。态射
$L \to \mathcal{E}^\bullet[1]$ 由复形态射
$\mathcal{L}^\bullet \to \mathcal{E}^\bullet[1]$ 给出；
该态射在除次数 $m - 1$ 之外的所有次数上为零，而在该次数
处得到态射 $\mathcal{L}^{m - 1} \to \mathcal{E}^m$，见
导出范畴，引理 \ref{derived-lemma-negative-exts}。
于是 $M$ 由复形
$$
\mathcal{M}^\bullet :
\ldots \to
\mathcal{L}^{m - 2} \to
\mathcal{L}^{m - 1} \to
\mathcal{E}^m \to
\mathcal{E}^{m + 1} \to \ldots
$$
表示。把第二段中的讨论应用于此复形，得到
$\mathcal{M}^{n_i}$ 的截面 $f_i$，其中 $i = 1, \ldots, N$。
对 $n < m$，令 $\mathcal{K}^n \subset \mathcal{L}^n$
为由满足 $n_i = n$ 的截面
$f_i$ 以及满足 $n_i = n - 1$ 的 $d(f_i)$ 生成的
$\mathcal{O}_X$-子模。
对 $n \geq m$，令 $\mathcal{K}^n = \mathcal{E}^n$。
显然有特异三角的态射
$$
\xymatrix{
\mathcal{E}^\bullet \ar[r] &
\mathcal{M}^\bullet \ar[r] &
\mathcal{L}^\bullet \ar[r] &
\mathcal{E}^\bullet[1] \\
\mathcal{E}^\bullet \ar[r] \ar[u] &
\mathcal{K}^\bullet \ar[r] \ar[u] &
\sigma_{\leq m - 1}\mathcal{K}^\bullet \ar[r] \ar[u] &
\mathcal{E}^\bullet[1] \ar[u]
}
$$
其中所有态射都如上所述。
记与复形 $\mathcal{K}^\bullet$ 对应的
$D(\mathcal{O}_X)$ 对象为 $K$。
由证明第二段中的论证，得到 $D(\mathcal{O}_X)$ 中的
态射 $s : M \to K$，使得复合
$M \to K \to M$ 是 $M$ 上的恒等态射。我们不知道图
$$
\xymatrix{
\mathcal{E}^\bullet \ar[r] &
\mathcal{K}^\bullet \ar@{=}[r] &
K \\
\mathcal{E}^\bullet \ar[u]^{\text{id}} \ar[r]^i &
\mathcal{M}^\bullet \ar@{=}[r] &
M \ar[u]_s
}
$$
是否交换，但知道它与态射 $K \to M$ 复合后交换。
由引理 \ref{cohomology-lemma-local-actual}，缩小 $X$ 后，可假设
$s \circ i$ 由复形态射
$\sigma : \mathcal{E}^\bullet \to \mathcal{K}^\bullet$ 给出。
再次使用同一引理，可假设 $\sigma$ 与含入
$\mathcal{K}^\bullet \subset \mathcal{M}^\bullet$ 的复合
由某个同伦
$\{h^i : \mathcal{E}^i \to \mathcal{M}^{i - 1}\}$
同伦于零。因此，把 $\mathcal{K}^{m - 1}$ 替换为
$\mathcal{K}^{m - 1} + \Im(h^m)$ 后（注意替换后
$\mathcal{K}^{m - 1}$ 仍由有限多个整体截面生成），
可见 $\sigma$ 本身同伦于零！
这意味着有如下实线部分交换的图：
$$
\xymatrix{
\mathcal{E}^\bullet \ar[r] &
M \ar[r] &
\mathcal{L}^\bullet \ar[r] &
\mathcal{E}^\bullet[1] \\
\mathcal{E}^\bullet \ar[r] \ar[u] &
K \ar[r] \ar[u] &
\sigma_{\leq m - 1}\mathcal{K}^\bullet \ar[r] \ar[u] &
\mathcal{E}^\bullet[1] \ar[u] \\
\mathcal{E}^\bullet \ar[r] \ar[u] &
M \ar[r] \ar[u]^s &
\mathcal{L}^\bullet \ar[r] \ar@{..>}[u] &
\mathcal{E}^\bullet[1] \ar[u]
}
$$
由三角范畴的公理，得到可嵌入该图的虚线
箭头。
考察次数 $m - 1$ 的上同调层，得到
$$
\xymatrix{
H^{m - 1}(M) \ar[r] &
H^{m - 1}(\mathcal{L}^\bullet) \ar[r] &
H^m(\mathcal{E}^\bullet) \\
H^{m - 1}(K) \ar[r] \ar[u] &
H^{m - 1}(\sigma_{\leq m - 1}\mathcal{K}^\bullet) \ar[r] \ar[u] &
H^m(\mathcal{E}^\bullet) \ar[u] \\
H^{m - 1}(M) \ar[r] \ar[u] &
H^{m - 1}(\mathcal{L}^\bullet) \ar[r] \ar[u] &
H^m(\mathcal{E}^\bullet) \ar[u]
}
$$
由于左列和右列中的竖直复合都是恒等态射，
可知中间的竖直复合
$H^{m - 1}(\mathcal{L}^\bullet) \to
H^{m - 1}(\sigma_{\leq m - 1}\mathcal{K}^\bullet) \to
H^{m - 1}(\mathcal{L}^\bullet)$ 是满射！
特别地，$H^{m - 1}(\sigma_{\leq m - 1}\mathcal{K}^\bullet) \to
H^{m - 1}(\mathcal{L}^\bullet)$ 是满射。
利用上述三角态射所诱导的上同调层长正合列态射，
追图可见，这推出 $H^i(K) \to H^i(M)$ 在
$i \geq m$ 时为同构，在 $i = m - 1$ 时为满射。
按构造，可选择 $r \geq 0$ 和满射
$\mathcal{O}_X^{\oplus r} \to \mathcal{K}^{m - 1}$。于是复合
$$
(\mathcal{O}_X^{\oplus r} \to \mathcal{E}^m \to
\mathcal{E}^{m + 1} \to \ldots ) \longrightarrow
K \longrightarrow M
$$
在次数 $\geq m$ 的上同调层上诱导同构，
并在次数 $m - 1$ 上诱导满射，证明完成。
\end{proof}






\section{杂项}
\label{cohomology-section-misc}

\noindent
一些无法归入别处的结果。

\begin{lemma}
\label{cohomology-lemma-colim-and-lim-of-duals}
令 $(X, \mathcal{O}_X)$ 为环化空间。令
$(K_n)_{n \in \mathbf{N}}$ 为 $D(\mathcal{O}_X)$ 中完美对象组成的系统。
令 $K = \text{hocolim} K_n$ 为导出余极限
（导出范畴，定义 \ref{derived-definition-derived-colimit}）。
那么对 $D(\mathcal{O}_X)$ 的任意对象 $E$，都有
$$
R\SheafHom(K, E) = R\lim E \otimes^\mathbf{L}_{\mathcal{O}_X} K_n^\vee
$$
其中 $(K_n^\vee)$ 是对偶完美复形组成的逆系统。
\end{lemma}

\begin{proof}
由引理 \ref{cohomology-lemma-dual-perfect-complex}，有
$R\lim E \otimes^\mathbf{L}_{\mathcal{O}_X} K_n^\vee =
R\lim R\SheafHom(K_n, E)$，
它嵌入特异三角
$$
R\lim R\SheafHom(K_n, E) \to
\prod R\SheafHom(K_n, E) \to
\prod R\SheafHom(K_n, E)
$$
由于 $K$ 同样嵌入特异三角
$\bigoplus K_n \to \bigoplus K_n \to K$，只需证明
$\prod R\SheafHom(K_n, E) = R\SheafHom(\bigoplus K_n, E)$。
这是 (\ref{cohomology-equation-internal-hom}) 以及导出张量积
与直和可交换的形式推论。
\end{proof}

\begin{lemma}
\label{cohomology-lemma-ext-composition-is-cup}
令 $(X, \mathcal{O}_X)$ 为环化空间。令 $K$ 和 $E$ 为
$D(\mathcal{O}_X)$ 的对象，其中 $E$ 完美。图
$$
\xymatrix{
H^0(X, K \otimes_{\mathcal{O}_X}^\mathbf{L} E^\vee) \times H^0(X, E)
\ar[r] \ar[d] &
H^0(X, K \otimes_{\mathcal{O}_X}^\mathbf{L} E^\vee
\otimes_{\mathcal{O}_X}^\mathbf{L} E) \ar[d] \\
\Hom_X(E, K) \times H^0(X, E) \ar[r] &
H^0(X, K)
}
$$
交换，其中上方水平箭头是杯积；右侧竖直箭头使用
$\epsilon : E^\vee \otimes_{\mathcal{O}_X}^\mathbf{L} E \to \mathcal{O}_X$
（例 \ref{cohomology-example-dual-derived}）；左侧竖直箭头使用引理 \ref{cohomology-lemma-dual-perfect-complex}；下方水平
箭头是显然的箭头。
\end{lemma}

\begin{proof}
我们将简写 $\otimes = \otimes_{\mathcal{O}_X}^\mathbf{L}$
和 $\mathcal{O} = \mathcal{O}_X$。我们把 $E$、$K$
分别与 $R\SheafHom(\mathcal{O}, E)$、$R\SheafHom(\mathcal{O}, K)$
等同，并把 $E^\vee$ 与 $R\SheafHom(E, \mathcal{O})$ 等同。

\medskip\noindent
令 $\xi \in H^0(X, K \otimes E^\vee)$ 且 $\eta \in H^0(X, E)$。
记 $D(\mathcal{O})$ 中相应的态射为
$\tilde \xi : \mathcal{O} \to K \otimes E^\vee$ 和
$\tilde \eta : \mathcal{O} \to E$。由引理
\ref{cohomology-lemma-second-cup-equals-first}，杯积
$\xi \cup \eta$ 对应于
$\tilde \xi \otimes \tilde \eta : \mathcal{O} \to
K \otimes E^\vee \otimes E$。

\medskip\noindent
我们断言，由引理 \ref{cohomology-lemma-dual-perfect-complex} 与 $\xi$
对应的态射 $\xi' : E \to K$ 是复合
$$
E = \mathcal{O} \otimes E
\xrightarrow{\tilde \xi \otimes 1_E}
K \otimes E^\vee \otimes E
\xrightarrow{1_K \otimes \epsilon}
K
$$
引理 \ref{cohomology-lemma-dual-perfect-complex} 中的构造
使用求值态射 (\ref{cohomology-equation-eval})；后者又使用
$E$ 与 $R\SheafHom(\mathcal{O}, E)$ 的等同，以及引理 \ref{cohomology-lemma-internal-hom-composition} 中构造的复合
$\underline{\circ}$。
因此 $\xi'$ 是复合
\begin{align*}
E = \mathcal{O} \otimes
R\SheafHom(\mathcal{O}, E)
& \xrightarrow{\tilde \xi \otimes 1}
R\SheafHom(\mathcal{O}, K) \otimes
R\SheafHom(E, \mathcal{O}) \otimes
R\SheafHom(\mathcal{O}, E) \\
& \xrightarrow{\underline{\circ} \otimes 1}
R\SheafHom(E, K) \otimes R\SheafHom(\mathcal{O}, E) \\
& \xrightarrow{\underline{\circ}}
R\SheafHom(\mathcal{O}, K) = K
\end{align*}
由此以及引理 \ref{cohomology-lemma-internal-hom-composition} 中构造的复合
$\underline{\circ}$ 具有结合性
（在此插入未来引用），并且例 \ref{cohomology-example-dual-derived} 中
$\epsilon$ 定义为复合
$\underline{\circ} : E^\vee \otimes E \to \mathcal{O}$，
立即得到该断言。

\medskip\noindent
使用前两段的结果，可知引理的断言等价于
$(1_K \otimes \epsilon) \circ (\tilde \xi \otimes \tilde \eta)$
等于
$(1_K \otimes \epsilon) \circ (\tilde \xi \otimes 1_E)
\circ (1_\mathcal{O} \otimes \tilde \eta)$，
这是立即的。
\end{proof}

\begin{lemma}
\label{cohomology-lemma-pullback-internal-hom}
令 $h : X \to Y$ 为环化空间的态射。
令 $K, M$ 为 $D(\mathcal{O}_Y)$ 的对象。注记
\ref{cohomology-remark-prepare-fancy-base-change} 中的典范态射
$$
Lh^*R\SheafHom(K, M) \longrightarrow R\SheafHom(Lh^*K, Lh^*M)
$$
在下列情形中是同构：
\begin{enumerate}
\item $K$ 完美；
\item $h$ 平坦，$K$ 伪凝聚，且 $M$（局部）下有界；
\item $\mathcal{O}_X$ 作为 $h^{-1}\mathcal{O}_Y$-模具有有限 Tor 维数，
$K$ 伪凝聚，且 $M$（局部）下有界。
\end{enumerate}
\end{lemma}

\begin{proof}
证明 (1)。问题在 $Y$ 上是局部的，故可假设
$K$ 由严格完美复形
$\mathcal{E}^\bullet$ 表示，见第 \ref{cohomology-section-perfect} 节。
选取表示 $M$ 的 K-平坦复形 $\mathcal{F}^\bullet$。
应用引理 \ref{cohomology-lemma-Rhom-strictly-perfect} 可知，
$R\SheafHom(K, L)$ 由复形
$\mathcal{H}^\bullet =
\SheafHom^\bullet(\mathcal{E}^\bullet, \mathcal{F}^\bullet)$ 表示，
其项为 $\mathcal{H}^n = \bigoplus\nolimits_{n = p + q}
\SheafHom_{\mathcal{O}_X}(\mathcal{E}^{-q}, \mathcal{F}^p)$。
由第 \ref{cohomology-section-derived-pullback} 节中 $Lh^*$ 的构造，
可知 $Lh^*K$ 由严格完美复形
$h^*\mathcal{E}^\bullet$ 表示（引理
\ref{cohomology-lemma-strictly-perfect-pullback}）。
类似地，对象 $Lh^*M$ 由
复形 $h^*\mathcal{F}^\bullet$ 表示。
最后，由引理 \ref{cohomology-lemma-Rhom-strictly-perfect-K-flat}，
$\mathcal{H}^\bullet$ 是 K-平坦的，故对象
$Lh^*R\SheafHom(K, M)$ 由 $h^*\mathcal{H}^\bullet$ 表示。
因此为完成证明，只需说明
$h^*\mathcal{H}^\bullet =
\SheafHom^\bullet(h^*\mathcal{E}^\bullet, h^*\mathcal{F}^\bullet)$。
为此只需注意，只要 $\mathcal{E}$ 是有限自由
$\mathcal{O}_X$-模的直和项，就有
$h^*\SheafHom(\mathcal{E}, \mathcal{F}) =
\SheafHom(h^*\mathcal{E}, \mathcal{F})$。

\medskip\noindent
证明 (2)。由于 $h$ 平坦，可以在任意 $\mathcal{O}_Y$-模复形上
直接使用 $h^*$ 来计算 $Lh^*$。
特别地，对所有 $i \in \mathbf{Z}$，有 $H^i(Lh^*K) = h^*H^i(K)$。
设当 $i < a$ 时 $H^i(M) = 0$。令 $K' \to K$ 为
$D(\mathcal{O}_Y)$ 中的态射，它对所有 $i \geq b$ 诱导同构
$H^i(K') \to H^i(K)$。那么相应的态射
$$
R\SheafHom(K, M) \to R\SheafHom(K', M)
$$
以及
$$
R\SheafHom(Lh^*K, Lh^*M) \to R\SheafHom(Lh^*K', Lh^*M)
$$
在次数 $< a - b$ 的上同调层上都是同构（略去细节）。
因此，为证明引理陈述中的态射在次数 $< a - b$ 的
上同调层上诱导同构，只需在这些次数上对 $K'$ 证明结论。
此外，与 (1) 的证明相同，问题在 $Y$ 上是局部的。
因此可假设 $K$ 由严格完美复形表示，见
第 \ref{cohomology-section-pseudo-coherent} 节。这就约化到情形 (1)。

\medskip\noindent
证明 (3)。证明与 (2) 相同，只是使用
$Lh^*$ 的上同调维数有界来得到所需的消失性。
略去细节。
\end{proof}

\begin{lemma}
\label{cohomology-lemma-stalk-internal-hom}
令 $X$ 为环化空间。令 $K, M$ 为 $D(\mathcal{O}_X)$ 的对象。
令 $x \in X$。典范态射
$$
R\SheafHom(K, M)_x \longrightarrow
R\Hom_{\mathcal{O}_{X, x}}(K_x, M_x)
$$
在下列情形中是同构：
\begin{enumerate}
\item $K$ 完美；
\item $K$ 伪凝聚且 $M$（局部）下有界。
\end{enumerate}
\end{lemma}

\begin{proof}
令 $Y = \{x\}$ 为单点环化空间，其结构层
由 $\mathcal{O}_{X, x}$ 给出。然后把引理
\ref{cohomology-lemma-pullback-internal-hom}
应用于平坦含入态射 $Y \to X$。
\end{proof}







\section{导出范畴中的可逆对象}
\label{cohomology-section-invertible-D-or-R}

\noindent
我们刻画环化空间的导出范畴中的可逆对象
（包括结构层的茎是局部环和不是局部环的情形）。

\begin{lemma}
\label{cohomology-lemma-category-summands-finite-free}
令 $(X, \mathcal{O}_X)$ 为环化空间。
置 $R = \Gamma(X, \mathcal{O}_X)$。作为有限自由
$\mathcal{O}_X$-模直和项的 $\mathcal{O}_X$-模所成的范畴，
等价于有限投射 $R$-模所成的范畴。
\end{lemma}

\begin{proof}
注意，有限投射 $R$-模等同于
有限自由 $R$-模的直和项。
该等价由函子 $\mathcal{E} \mapsto
\Gamma(X, \mathcal{E})$ 给出。逆函子由
模，引理 \ref{modules-lemma-construct-quasi-coherent-sheaves}
中的构造给出。
\end{proof}

\begin{lemma}
\label{cohomology-lemma-invertible-derived}
令 $(X, \mathcal{O}_X)$ 为环化空间。令 $M$ 为
$D(\mathcal{O}_X)$ 的对象。下列条件等价：
\begin{enumerate}
\item $M$ 在 $D(\mathcal{O}_X)$ 中可逆，见
范畴，定义 \ref{categories-definition-invertible}；以及
\item 存在局部有限的乘积分解
$$
\mathcal{O}_X = \prod\nolimits_{n \in \mathbf{Z}} \mathcal{O}_n
$$
并且对每个 $n$，存在可逆 $\mathcal{O}_n$-模
$\mathcal{H}^n$（模，定义 \ref{modules-definition-invertible}），
使得在 $D(\mathcal{O}_X)$ 中 $M = \bigoplus \mathcal{H}^n[-n]$。
\end{enumerate}
如果 (1) 和 (2) 成立，那么 $M$ 是 $D(\mathcal{O}_X)$ 的完美对象。如果
对所有 $x \in X$，$\mathcal{O}_{X, x}$ 都是局部环，那么这些条件
还等价于
\begin{enumerate}
\item[(3)] 存在开覆盖 $X = \bigcup U_i$，
并且对每个 $i$ 存在整数 $n_i$，使得 $M|_{U_i}$
由放在次数 $n_i$ 上的可逆 $\mathcal{O}_{U_i}$-模表示。
\end{enumerate}
\end{lemma}

\begin{proof}
假设 (2)。考虑对象 $R\SheafHom(M, \mathcal{O}_X)$
和复合态射
$$
R\SheafHom(M, \mathcal{O}_X) \otimes_{\mathcal{O}_X}^\mathbf{L} M \to
\mathcal{O}_X
$$
为证明这是同构，可在局部工作。因此可假设
$\mathcal{O}_X = \prod_{a \leq n \leq b} \mathcal{O}_n$
且 $M = \bigoplus_{a \leq n \leq b} \mathcal{H}^n[-n]$。
于是只需证明
$$
R\SheafHom(\mathcal{H}^m, \mathcal{O}_X)
\otimes_{\mathcal{O}_X}^\mathbf{L} \mathcal{H}^n
$$
在 $n \not = m$ 时为零，在 $n = m$ 时等于 $\mathcal{O}_n$。
$n \not = m$ 的情形由如下事实得出：$\mathcal{O}_n$ 和
$\mathcal{O}_m$ 都是平坦 $\mathcal{O}_X$-代数，并且
$\mathcal{O}_n \otimes_{\mathcal{O}_X} \mathcal{O}_m = 0$。
利用可逆 $\mathcal{O}_X$-模的局部结构
（模，引理 \ref{modules-lemma-invertible}）并在局部工作，
$n = m$ 情形中的同构可直接得到；
略去细节。由于 $D(\mathcal{O}_X)$ 是对称幺半范畴，
可知 $M$ 可逆。

\medskip\noindent
假设 (1)。(2) 中的描述给出了
$\mathcal{O}_n$ 的候选，即
$\SheafHom_{\mathcal{O}_X}(H^n(M), H^n(M))$。
如果这是一个局部有限的环层族，
并且 $\mathcal{O}_X = \prod \mathcal{O}_n$，那么利用
$\mathcal{O}_X$ 中来自该乘积分解的幂等元，立即得到
直和分解 $M = \bigoplus H^n(M)[-n]$。
这说明，为证明 (2)，可在 $X$ 上局部工作。

\medskip\noindent
选取 $D(\mathcal{O}_X)$ 的对象 $N$
以及同构
$M \otimes_{\mathcal{O}_X}^\mathbf{L} N \cong \mathcal{O}_X$。
令 $x \in X$。
那么 $N$ 是幺半范畴 $D(\mathcal{O}_X)$ 中 $M$ 的左对偶，
由引理 \ref{cohomology-lemma-left-dual-derived} 可知 $M$ 完美。
由对称性可见 $N$ 完美。用 $x$ 的开邻域替换 $X$ 后，
可假设 $M$ 和 $N$ 分别由严格完美
复形 $\mathcal{E}^\bullet$ 和 $\mathcal{F}^\bullet$ 表示。
于是 $M \otimes_{\mathcal{O}_X}^\mathbf{L} N$ 由
$\text{Tot}(\mathcal{E}^\bullet \otimes_{\mathcal{O}_X} \mathcal{F}^\bullet)$ 表示。
再次缩小 $X$ 后，可假设互逆同构
$\mathcal{O}_X \to M \otimes_{\mathcal{O}_X}^\mathbf{L} N$ 和
$M \otimes_{\mathcal{O}_X}^\mathbf{L} N \to \mathcal{O}_X$
由复形态射
$$
\alpha : \mathcal{O}_X \to
\text{Tot}(\mathcal{E}^\bullet \otimes_{\mathcal{O}_X} \mathcal{F}^\bullet)
\quad\text{且}\quad
\beta :
\text{Tot}(\mathcal{E}^\bullet \otimes_{\mathcal{O}_X} \mathcal{F}^\bullet)
\to \mathcal{O}_X
$$
给出。见引理 \ref{cohomology-lemma-local-actual}。于是作为复形态射，
$\beta \circ \alpha = 1$；而作为
$D(\mathcal{O}_X)$ 中的态射，$\alpha \circ \beta = 1$。
缩小 $X$ 后，由同一引理，可假设复合
$\alpha \circ \beta$ 通过某个同伦 $\theta$ 同伦于 $1$，
其分量为
$$
\theta^n :
\text{Tot}^n(\mathcal{E}^\bullet \otimes_{\mathcal{O}_X} \mathcal{F}^\bullet)
\to
\text{Tot}^{n - 1}(
\mathcal{E}^\bullet \otimes_{\mathcal{O}_X} \mathcal{F}^\bullet)
$$
置 $R = \Gamma(X, \mathcal{O}_X)$。由引理 \ref{cohomology-lemma-category-summands-finite-free}，
得到
\begin{enumerate}
\item $M^\bullet = \Gamma(X, \mathcal{E}^\bullet)$ 是由
有限投射 $R$-模组成的有界复形；
\item $N^\bullet = \Gamma(X, \mathcal{F}^\bullet)$ 是由
有限投射 $R$-模组成的有界复形；
\item $\alpha$ 和 $\beta$ 对应于复形态射
$a : R \to \text{Tot}(M^\bullet \otimes_R N^\bullet)$ 和
$b : \text{Tot}(M^\bullet \otimes_R N^\bullet) \to R$；
\item $\theta^n$ 对应于态射
$h^n : \text{Tot}^n(M^\bullet \otimes_R N^\bullet) \to
\text{Tot}^{n - 1}(M^\bullet \otimes_R N^\bullet)$；以及
\item $b \circ a = 1$ 且 $b \circ a - 1 = dh + hd$。
\end{enumerate}
由此 $M^\bullet$ 和 $N^\bullet$ 定义
$D(R)$ 中互逆的对象。由
更多代数，引理 \ref{more-algebra-lemma-invertible-derived}，
得到乘积分解 $R = \prod_{a \leq n \leq b} R_n$
以及可逆 $R_n$-模 $H^n$，使得
$M^\bullet \cong \bigoplus_{a \leq n \leq b} H^n[-n]$。
由于每个 $H^n$ 都是投射 $R$-模，$D(R)$ 中的该同构
可提升为复形态射
$$
\bigoplus H^n[-n] \longrightarrow M^\bullet
$$
相应地，再次使用引理 \ref{cohomology-lemma-category-summands-finite-free}，
得到态射
$$
\bigoplus H^n \otimes_R \mathcal{O}_X[-n] \to \mathcal{E}^\bullet
$$
它在 $D(\mathcal{O}_X)$ 中为同构。置
$\mathcal{O}_n = R_n \otimes_R \mathcal{O}_X$，可知 (2) 成立。

\medskip\noindent
如果 $\mathcal{O}_X$ 的所有茎都是局部环，那么
(2) 与 (3) 的等价性很容易证明。略去细节。
\end{proof}







\section{紧对象}
\label{cohomology-section-compact}

\noindent
本节研究环化空间上的模的导出范畴中的紧对象。
回忆，紧对象定义于
导出范畴，定义 \ref{derived-definition-compact-object}。
在适当的环化空间上，完美对象是紧的。

\begin{lemma}
\label{cohomology-lemma-when-jshriek-compact}
令 $X$ 为环化空间。令 $j : U \to X$ 为
开集的含入。如果存在整数 $d$，使得
\begin{enumerate}
\item 对所有 $p > d$，$H^p(U, \mathcal{F}) = 0$；并且
\item 函子 $\mathcal{F} \mapsto H^p(U, \mathcal{F})$
与直和可交换，
\end{enumerate}
那么 $\mathcal{O}_X$-模 $j_!\mathcal{O}_U$ 是
$D(\mathcal{O}_X)$ 的紧对象。
\end{lemma}

\begin{proof}
假设 (1) 和 (2)。由层，引理
\ref{sheaves-lemma-j-shriek-modules}，
$\Hom(j_!\mathcal{O}_U, \mathcal{F}) = \mathcal{F}(U)$，
故对 $D(\mathcal{O}_X)$ 中的 $K$，有
$\Hom(j_!\mathcal{O}_U, K) = R\Gamma(U, K)$。
因此必须证明 $R\Gamma(U, -)$ 与直和可交换。
第一个假设意味着函子
$F = H^0(U, -)$ 具有有限上同调维数。此外，第二个
假设意味着内射模的任意直和都对 $F$ 无环。
令 $K_i$ 为 $D(\mathcal{O}_X)$ 的一族对象。
选取表示 $K_i$ 且各项内射的 K-内射代表
$I_i^\bullet$，见内射对象，定理
\ref{injectives-theorem-K-injective-embedding-grothendieck}。
由于可通过把 $F$ 应用于任意无环对象复形来计算 $RF$
（导出范畴，引理 \ref{derived-lemma-unbounded-right-derived}），
并且 $\bigoplus K_i$ 由 $\bigoplus I_i^\bullet$ 表示
（内射对象，引理 \ref{injectives-lemma-derived-products}），
可知 $R\Gamma(U, \bigoplus K_i)$ 由
$\bigoplus H^0(U, I_i^\bullet)$ 表示。因此
$R\Gamma(U, -)$ 如所需地与直和可交换。
\end{proof}

\begin{lemma}
\label{cohomology-lemma-perfect-is-compact}
令 $X$ 为环化空间。假设 $X$ 的底拓扑空间
具有下列性质：
\begin{enumerate}
\item $X$ 拟紧；
\item 存在由拟紧开子集组成的基；并且
\item 任意两个拟紧开集的交仍拟紧。
\end{enumerate}
令 $K$ 为 $D(\mathcal{O}_X)$ 的完美对象。那么
\begin{enumerate}
\item[(a)] $K$ 在如下意义下是 $D^+(\mathcal{O}_X)$ 的紧对象：
如果 $M = \bigoplus_{i \in I} M_i$ 下有界，那么
$\Hom(K, M) = \bigoplus_{i \in I} \Hom(K, M_i)$。
\item[(b)] 如果 $X$ 具有有限上同调维数，即存在
$d$，使得当 $i > d$ 时 $H^i(X, \mathcal{F}) = 0$，那么
$K$ 是 $D(\mathcal{O}_X)$ 的紧对象。
\end{enumerate}
\end{lemma}

\begin{proof}
令 $K^\vee$ 为 $K$ 的对偶，见引理 \ref{cohomology-lemma-dual-perfect-complex}。那么对
$D(\mathcal{O}_X)$ 中的 $M$，函子性地有
$$
\Hom_{D(\mathcal{O}_X)}(K, M) =
H^0(X, K^\vee \otimes_{\mathcal{O}_X}^\mathbf{L} M)
$$
由于 $K^\vee \otimes_{\mathcal{O}_X}^\mathbf{L} -$ 与
直和可交换，只需证明 $R\Gamma(X, -)$
与相应的直和可交换。

\medskip\noindent
证明 (b)。由于 $R\Gamma(X, K) = R\Hom(\mathcal{O}_X, K)$，
并且由引理 \ref{cohomology-lemma-quasi-separated-cohomology-colimit}，
$H^p(X, -)$ 与直和可交换，这是引理 \ref{cohomology-lemma-when-jshriek-compact} 的特例。

\medskip\noindent
证明 (a)。令 $\mathcal{I}_i$（$i \in I$）为一族内射
$\mathcal{O}_X$-模。由引理 \ref{cohomology-lemma-quasi-separated-cohomology-colimit}，
可见
$$
H^p(X, \bigoplus\nolimits_{i \in I} \mathcal{I}_i) =
\bigoplus\nolimits_{i \in I} H^p(X, \mathcal{I}_i) = 0
$$
对所有 $p$ 成立。现在若 $M = \bigoplus M_i$ 如 (a) 中所述，
则存在 $a \in \mathbf{Z}$，使得当 $n < a$ 时
$H^n(M_i) = 0$。因此可选取表示 $M_i$ 的内射
$\mathcal{O}_X$-模复形 $\mathcal{I}_i^\bullet$，
使得当 $n < a$ 时 $\mathcal{I}_i^n = 0$，见
导出范畴，引理 \ref{derived-lemma-injective-resolutions-exist}。
由内射对象，引理 \ref{injectives-lemma-derived-products}，
直和复形 $\bigoplus \mathcal{I}_i^\bullet$
表示 $M$。由勒雷无环性
（导出范畴，引理 \ref{derived-lemma-leray-acyclicity}），
可见
$$
R\Gamma(X, M) = \Gamma(X, \bigoplus \mathcal{I}_i^\bullet) =
\bigoplus \Gamma(X, \bigoplus \mathcal{I}_i^\bullet) =
\bigoplus R\Gamma(X, M_i)
$$
如所需。
\end{proof}



















\section{投影公式}
\label{cohomology-section-projection-formula}

\noindent
本节汇集投影公式的若干变体。
最基本的版本是引理 \ref{cohomology-lemma-projection-formula}。
在陈述并证明它之后，我们将讨论一个涉及
完美复形的更一般版本。

\begin{lemma}
\label{cohomology-lemma-injective-tensor-finite-locally-free}
令 $X$ 为环化空间。
令 $\mathcal{I}$ 为内射 $\mathcal{O}_X$-模。
令 $\mathcal{E}$ 为 $\mathcal{O}_X$-模。
假设 $\mathcal{E}$ 在 $X$ 上有限局部自由，见
模，定义 \ref{modules-definition-locally-free}。
那么 $\mathcal{E} \otimes_{\mathcal{O}_X} \mathcal{I}$
是内射 $\mathcal{O}_X$-模。
\end{lemma}

\begin{proof}
这是因为在引理的假设下，有
$$
\Hom_{\mathcal{O}_X}(\mathcal{F},
\mathcal{E} \otimes_{\mathcal{O}_X} \mathcal{I})
=
\Hom_{\mathcal{O}_X}(
\mathcal{F} \otimes_{\mathcal{O}_X} \mathcal{E}^\vee, \mathcal{I})
$$
其中
$\mathcal{E}^\vee = \SheafHom_{\mathcal{O}_X}(\mathcal{E}, \mathcal{O}_X)$
是 $\mathcal{E}$ 的对偶，它同样有限局部自由。由于与有限局部自由层
作张量积是正合函子，由同调，引理
\ref{homology-lemma-characterize-injectives} 即得结论。
\end{proof}

\begin{lemma}
\label{cohomology-lemma-projection-formula}
令 $f : X \to Y$ 为环化空间的态射。
令 $\mathcal{F}$ 为 $\mathcal{O}_X$-模。
令 $\mathcal{E}$ 为 $\mathcal{O}_Y$-模。
假设 $\mathcal{E}$ 在 $Y$ 上有限局部自由，见
模，定义 \ref{modules-definition-locally-free}。
那么对所有 $q \geq 0$，存在同构
$$
\mathcal{E} \otimes_{\mathcal{O}_Y} R^qf_*\mathcal{F}
\longrightarrow
R^qf_*(f^*\mathcal{E} \otimes_{\mathcal{O}_X} \mathcal{F})
$$
事实上，在 $D^{+}(Y)$ 中存在关于 $\mathcal{F}$
具有函子性的同构
$$
\mathcal{E} \otimes_{\mathcal{O}_Y} Rf_*\mathcal{F}
\longrightarrow
Rf_*(f^*\mathcal{E} \otimes_{\mathcal{O}_X} \mathcal{F})
$$
\end{lemma}

\begin{proof}
在 $X$ 上选取内射消解 $\mathcal{F} \to \mathcal{I}^\bullet$。
注意，$f^*\mathcal{E}$ 也有限局部自由，故得到消解
$$
f^*\mathcal{E} \otimes_{\mathcal{O}_X} \mathcal{F}
\longrightarrow
f^*\mathcal{E} \otimes_{\mathcal{O}_X} \mathcal{I}^\bullet
$$
由引理 \ref{cohomology-lemma-injective-tensor-finite-locally-free}，
这是内射消解。
应用 $f_*$，得到
$$
Rf_*(f^*\mathcal{E} \otimes_{\mathcal{O}_X} \mathcal{F})
=
f_*(f^*\mathcal{E} \otimes_{\mathcal{O}_X} \mathcal{I}^\bullet).
$$
因此，只要能证明
$f_*(f^*\mathcal{E} \otimes_{\mathcal{O}_X} \mathcal{F}) =
\mathcal{E} \otimes_{\mathcal{O}_Y} f_*(\mathcal{F})$
关于 $\mathcal{O}_X$-模 $\mathcal{F}$ 具有函子性，引理即得证。
当 $\mathcal{E} = \mathcal{O}_Y^{\oplus n}$ 时这是显然的；
一般情形通过在 $Y$ 上局部工作得到。略去细节。
\end{proof}

\noindent
令 $f : X \to Y$ 为环化空间的态射。
令 $E \in D(\mathcal{O}_X)$ 且 $K \in D(\mathcal{O}_Y)$。
不作任何进一步假设，存在态射
\begin{equation}
\label{cohomology-equation-projection-formula-map}
Rf_*E \otimes^\mathbf{L}_{\mathcal{O}_Y} K
\longrightarrow
Rf_*(E \otimes^\mathbf{L}_{\mathcal{O}_X} Lf^*K)
\end{equation}
也就是说，它伴随于典范态射
$$
Lf^*(Rf_*E \otimes^\mathbf{L}_{\mathcal{O}_Y} K) =
Lf^*Rf_*E \otimes^\mathbf{L}_{\mathcal{O}_X} Lf^*K
\longrightarrow
E \otimes^\mathbf{L}_{\mathcal{O}_X} Lf^*K
$$
该态射来自 $Lf^*Rf_*E \to E$ 以及引理 \ref{cohomology-lemma-pullback-tensor-product} 和 \ref{cohomology-lemma-adjoint}。
投影公式一个相当一般的版本如下。

\begin{lemma}
\label{cohomology-lemma-projection-formula-perfect}
令 $f : X \to Y$ 为环化空间的态射。
令 $E \in D(\mathcal{O}_X)$ 且 $K \in D(\mathcal{O}_Y)$。
如果 $K$ 完美，那么在 $D(\mathcal{O}_Y)$ 中
$$
Rf_*E \otimes^\mathbf{L}_{\mathcal{O}_Y} K =
Rf_*(E \otimes^\mathbf{L}_{\mathcal{O}_X} Lf^*K)
$$
\end{lemma}

\begin{proof}
为检验 (\ref{cohomology-equation-projection-formula-map}) 是同构，
可在 $Y$ 上局部工作，也就是说，要找到开覆盖
$\{V_j \to Y\}$，使得该态射限制到 $V_j$ 后为同构。
按完美对象的定义，这意味着可假设 $K$ 由
$\mathcal{O}_Y$-模的严格完美复形表示。
注意，完全一般地，断言对
$K = K_1 \oplus K_2$ 成立，当且仅当它对
$K_1$ 和 $K_2$ 都成立。因此可假设 $K$ 是由
有限自由 $\mathcal{O}_Y$-模组成的有限复形。
此时，涉及朴素截断的简单论证把断言约化到
$K$ 由有限自由 $\mathcal{O}_Y$-模表示的情形。
由于断言关于变量 $K$ 在有限直和项下不变，可知
只需对 $K = \mathcal{O}_Y[n]$ 证明它，
而在该情形中结论平凡。
\end{proof}

\noindent
以下是一种投影公式完全一般地成立的情形。

\begin{lemma}
\label{cohomology-lemma-projection-formula-closed-immersion}
令 $f : X \to Y$ 为环化空间的态射，并且 $f$
是到某个闭子集上的同胚。那么
(\ref{cohomology-equation-projection-formula-map}) 总是同构。
\end{lemma}

\begin{proof}
由于 $f$ 是到某个闭子集上的同胚，函子 $f_*$
是正合的（模，引理 \ref{modules-lemma-i-star-exact}）。
因此，可通过把 $f_*$ 应用于任意代表复形来计算 $Rf_*$。
选取表示 $K$ 的 K-平坦 $\mathcal{O}_Y$-模复形
$\mathcal{K}^\bullet$，并任取表示 $E$ 的
$\mathcal{O}_X$-模复形 $\mathcal{E}^\bullet$。那么
$f^*\mathcal{K}^\bullet$ 表示 $Lf^*K$，并且它是
K-平坦 $\mathcal{O}_X$-模复形
（引理 \ref{cohomology-lemma-pullback-K-flat}）。因此
(\ref{cohomology-equation-projection-formula-map}) 的右端由
$$
f_*\text{Tot}(\mathcal{E}^\bullet
\otimes_{\mathcal{O}_X} f^*\mathcal{K}^\bullet)
$$
表示。同理，左端由
$$
\text{Tot}(f_*\mathcal{E}^\bullet \otimes_{\mathcal{O}_Y} \mathcal{K}^\bullet)
$$
表示。由于 $f_*$ 与直和可交换
（模，引理 \ref{modules-lemma-i-star-right-adjoint}），
只需证明
$$
f_*(\mathcal{E} \otimes_{\mathcal{O}_X} f^*\mathcal{K}) =
f_*\mathcal{E} \otimes_{\mathcal{O}_Y} \mathcal{K}
$$
对任意 $\mathcal{O}_X$-模 $\mathcal{E}$ 和 $\mathcal{O}_Y$-模
$\mathcal{K}$ 成立。我们将在茎上检验这一点。
令 $y \in Y$。如果 $y \not \in f(X)$，那么两端的茎
都为零。如果 $y = f(x)$，那么要证明
$$
\mathcal{E}_x \otimes_{\mathcal{O}_{X, x}}
(\mathcal{O}_{X, x} \otimes_{\mathcal{O}_{Y, y}} \mathcal{F}_y) =
\mathcal{E}_x \otimes_{\mathcal{O}_{Y, y}} \mathcal{F}_y
$$
（使用层，引理 \ref{sheaves-lemma-stalks-closed-pushforward}
和引理 \ref{sheaves-lemma-stalk-pullback-modules}）。
该等式成立，故引理得证。
\end{proof}

\begin{remark}
\label{cohomology-remark-compatible-with-diagram}
态射 (\ref{cohomology-equation-projection-formula-map}) 与注记
\ref{cohomology-remark-base-change} 中的基变换态射具有如下相容性。
具体而言，假设
$$
\xymatrix{
X' \ar[r]_{g'} \ar[d]_{f'} &
X \ar[d]^f \\
Y' \ar[r]^g &
Y
}
$$
是环化空间的交换图。
令 $E \in D(\mathcal{O}_X)$ 且 $K \in D(\mathcal{O}_Y)$。
那么图
$$
\xymatrix{
Lg^*(Rf_*E \otimes^\mathbf{L}_{\mathcal{O}_Y} K) \ar[r]_p \ar[d]_t &
Lg^*Rf_*(E \otimes^\mathbf{L}_{\mathcal{O}_X} Lf^*K) \ar[d]_b \\
Lg^*Rf_*E \otimes^\mathbf{L}_{\mathcal{O}_{Y'}} Lg^*K \ar[d]_b &
Rf'_*L(g')^*(E \otimes^\mathbf{L}_{\mathcal{O}_X} Lf^*K) \ar[d]_t \\
Rf'_*L(g')^*E \otimes^\mathbf{L}_{\mathcal{O}_{Y'}} Lg^*K \ar[rd]_p &
Rf'_*(L(g')^*E \otimes^\mathbf{L}_{\mathcal{O}_{Y'}} L(g')^*Lf^*K) \ar[d]_c \\
& Rf'_*(L(g')^*E \otimes^\mathbf{L}_{\mathcal{O}_{Y'}} L(f')^*Lg^*K)
}
$$
交换。这里，标记为 $t$ 的箭头由应用引理 \ref{cohomology-lemma-pullback-tensor-product} 得到，标记为 $b$ 的箭头
由应用注记 \ref{cohomology-remark-base-change} 得到，标记为 $p$ 的箭头
由应用 (\ref{cohomology-equation-projection-formula-map}) 得到，而
$c$ 来自 $L(g')^* \circ Lf^* = L(f')^* \circ Lg^*$。
略去验证。
\end{remark}

















\section{Berthelot 与 Ogus 引入的算子}
\label{cohomology-section-eta}

\noindent
本节继续更多代数，第
\ref{more-algebra-section-eta} 节中开始的讨论。
我们强烈建议读者先阅读该节。

\begin{lemma}
\label{cohomology-lemma-invertible-ideal-sheaf}
令 $(X, \mathcal{O}_X)$ 为环化空间。令
$\mathcal{I} \subset \mathcal{O}_X$ 为理想层。
考虑下列两个条件：
\begin{enumerate}
\item 对每个 $x \in X$，存在 $x$ 的开邻域
$U \subset X$ 和 $f \in \mathcal{I}(U)$，使得
$\mathcal{I}|_U = \mathcal{O}_U \cdot f$，并且
$f : \mathcal{O}_U \to \mathcal{O}_U$ 是单射；以及
\item $\mathcal{I}$ 作为 $\mathcal{O}_X$-模可逆。
\end{enumerate}
那么 (1) 推出 (2)；如果结构层的所有茎
$\mathcal{O}_{X, x}$ 都是局部环，则逆命题也成立。
\end{lemma}

\begin{proof}
略去证明。提示：使用模，引理
\ref{modules-lemma-invertible-is-locally-free-rank-1}。
\end{proof}

\begin{situation}
\label{cohomology-situation-eta}
令 $(X, \mathcal{O}_X)$ 为环化空间。令
$\mathcal{I} \subset \mathcal{O}_X$ 为满足引理 \ref{cohomology-lemma-invertible-ideal-sheaf} 中条件 (1)
的理想层\footnote{本节的讨论可推广到只要求
$\mathcal{I}$ 是模，第 \ref{modules-section-invertible} 节所定义的
可逆 $\mathcal{O}_X$-模这一情形。}。
\end{situation}

\begin{lemma}
\label{cohomology-lemma-I-torsion-free}
在情形 \ref{cohomology-situation-eta} 中，
令 $\mathcal{F}$ 为 $\mathcal{O}_X$-模。
下列条件等价：
\begin{enumerate}
\item 由被 $\mathcal{I}$ 零化的截面组成的子层
$\mathcal{F}[\mathcal{I}] \subset \mathcal{F}$ 为零；
\item 对所有 $n \geq 1$，子层 $\mathcal{F}[\mathcal{I}^n]$ 为零；
\item 乘法态射
$\mathcal{I} \otimes_{\mathcal{O}_X} \mathcal{F} \to \mathcal{F}$
是单射；
\item 对每个满足
$\mathcal{I}|_U = \mathcal{O}_U \cdot f$
（其中 $f \in \mathcal{I}(U)$）的开集 $U \subset X$，
态射 $f : \mathcal{F}|_U \to \mathcal{F}|_U$ 都是单射；
\item 对每个 $x \in X$ 以及理想
$\mathcal{I}_x \subset \mathcal{O}_{X, x}$ 的生成元 $f$，
元素 $f$ 在茎 $\mathcal{F}_x$ 上是非零因子。
\end{enumerate}
\end{lemma}

\begin{proof}
略去证明。
\end{proof}

\noindent
在情形 \ref{cohomology-situation-eta} 中，
令 $\mathcal{F}$ 为 $\mathcal{O}_X$-模。
如果引理 \ref{cohomology-lemma-I-torsion-free} 的等价条件成立，
则称 $\mathcal{F}$ {\it $\mathcal{I}$-无挠}。
在这种情况下，对任意 $i \in \mathbf{Z}$，记
$$
\mathcal{I}^i\mathcal{F} =
\mathcal{I}^{\otimes i} \otimes_{\mathcal{O}_X} \mathcal{F}
$$
于是有含入
$$
\ldots \subset
\mathcal{I}^{i + 1}\mathcal{F} \subset
\mathcal{I}^i\mathcal{F} \subset
\mathcal{I}^{i - 1}\mathcal{F} \subset \ldots
$$
各模 $\mathcal{I}^i\mathcal{F}$ 作为 $\mathcal{O}_X$-模
局部同构于 $\mathcal{F}$，但整体上未必同构。

\medskip\noindent
令 $\mathcal{F}^\bullet$ 为由 $\mathcal{I}$-无挠
$\mathcal{O}_X$-模组成的复形，其微分为
$d^i : \mathcal{F}^i \to \mathcal{F}^{i + 1}$。
此时定义 $\eta_\mathcal{I}\mathcal{F}^\bullet$
为具有下列各项的复形：
\begin{align*}
(\eta_\mathcal{I}\mathcal{F})^i
& =
\Ker\left(
d^i, -1 :
\mathcal{I}^i\mathcal{F}^i \oplus \mathcal{I}^{i + 1}\mathcal{F}^{i + 1}
\to
\mathcal{I}^i\mathcal{F}^{i + 1}
\right) \\
& =
\Ker\left(d^i :
\mathcal{I}^i\mathcal{F}^i
\to
\mathcal{I}^i\mathcal{F}^{i + 1}/
\mathcal{I}^{i + 1}\mathcal{F}^{i + 1}
\right)
\end{align*}
其微分由 $d^i$ 诱导。换言之，
$(\eta_\mathcal{I}\mathcal{F})^i$ 的局部截面 $s$
等同于 $\mathcal{I}^i\mathcal{F}^i$ 的局部截面 $s$，
且它在 $\mathcal{I}^i\mathcal{F}^{i + 1}$ 中的像 $d^i(s)$
属于子层 $\mathcal{I}^{i + 1}\mathcal{F}^{i + 1}$。
注意，$\eta_\mathcal{I}\mathcal{F}^\bullet$
仍是由 $\mathcal{I}$-无挠模组成的复形。

\medskip\noindent
令 $a^\bullet : \mathcal{F}^\bullet \to \mathcal{G}^\bullet$ 为
由 $\mathcal{I}$-无挠 $\mathcal{O}_X$-模组成的复形之间的态射。
于是由态射
$\mathcal{I}^i\mathcal{F}^i \to \mathcal{I}^i\mathcal{G}^i$
诱导复形态射
$$
\eta_\mathcal{I} a^\bullet :
\eta_\mathcal{I}\mathcal{F}^\bullet
\longrightarrow
\eta_\mathcal{I}\mathcal{G}^\bullet
$$
读者可检验，这给出由
$\mathcal{I}$-无挠 $\mathcal{O}_X$-模组成的复形范畴上的
自函子。

\medskip\noindent
如果 $a^\bullet, b^\bullet : \mathcal{F}^\bullet \to \mathcal{G}^\bullet$
是由 $\mathcal{I}$-无挠 $\mathcal{O}_X$-模组成的
复形之间的两个态射，并且
$h = \{h^i : \mathcal{F}^i \to \mathcal{G}^{i - 1}\}$
是 $a^\bullet$ 与 $b^\bullet$ 之间的同伦，那么定义
$\eta_\mathcal{I}h$ 为态射族
$(\eta_\mathcal{I}h)^i : (\eta_\mathcal{I}\mathcal{F})^i \to
(\eta_\mathcal{I}\mathcal{G})^{i - 1}$，
它把 $x$ 送到 $h^i(x)$；这是有意义的，因为
$x$ 是 $\mathcal{I}^i\mathcal{F}^i$ 的局部截面
意味着 $h^i(x)$ 是 $\mathcal{I}^i\mathcal{G}^{i - 1}$ 的局部截面，
而后者当然包含于 $(\eta_\mathcal{I}\mathcal{G})^{i - 1}$。
读者可检验，$\eta_\mathcal{I}h$ 是
$\eta_\mathcal{I}a^\bullet$ 与 $\eta_\mathcal{I}b^\bullet$ 之间的同伦。
综上，在 $\mathcal{I}$-无挠 $\mathcal{O}_X$-模的
加性范畴的同伦范畴上，得到函子
$$
\eta_f :
K(\mathcal{I}\text{-无挠 }\mathcal{O}_X\text{-模})
\longrightarrow
K(\mathcal{I}\text{-无挠 }\mathcal{O}_X\text{-模})
$$
（导出范畴，第 \ref{derived-section-homotopy} 节）。
$\eta_\mathcal{I}$ 在任何意义下都不是三角范畴的正合函子；
可与更多代数，例
\ref{more-algebra-example-eta-not-distinguished} 比较。

\begin{lemma}
\label{cohomology-lemma-eta-stalks}
在情形 \ref{cohomology-situation-eta} 中，
令 $\mathcal{F}^\bullet$ 为由 $\mathcal{I}$-无挠
$\mathcal{O}_X$-模组成的复形。
对 $x \in X$，选取生成元 $f \in \mathcal{I}_x$。那么茎
$(\eta_\mathcal{I}\mathcal{F}^\bullet)_x$
典范同构于更多代数，第 \ref{more-algebra-section-eta} 节中
构造的复形 $\eta_f\mathcal{F}^\bullet_x$。
\end{lemma}

\begin{proof}
略去证明。
\end{proof}

\begin{lemma}
\label{cohomology-lemma-eta-first-property}
在情形 \ref{cohomology-situation-eta} 中，
令 $\mathcal{F}^\bullet$ 为由 $\mathcal{I}$-无挠
$\mathcal{O}_X$-模组成的复形。存在上同调层的典范同构
$$
\mathcal{I}^{\otimes i} \otimes_{\mathcal{O}_X}
\left(
H^i(\mathcal{F}^\bullet)/H^i(\mathcal{F}^\bullet)[\mathcal{I}]
\right)
\longrightarrow H^i(\eta_\mathcal{I}\mathcal{F}^\bullet)
$$
\end{lemma}

\begin{proof}
如下定义态射：
$$
\mathcal{I}^{\otimes i} \otimes_{\mathcal{O}_X} H^i(\mathcal{F}^\bullet)
\longrightarrow
H^i(\eta_\mathcal{I}\mathcal{F}^\bullet)
$$
令 $g$ 为 $\mathcal{I}^{\otimes i}$ 的局部截面，
令 $\overline{s}$ 为 $H^i(\mathcal{F}^\bullet)$ 的局部截面。
那么 $\overline{s}$（局部地）是
$\Ker(d^i : \mathcal{F}^i \to \mathcal{F}^{i + 1})$
的局部截面 $s$ 的类。
我们把 $g \otimes \overline{s}$ 送到
$(\eta_\mathcal{I}\mathcal{F})^i \subset \mathcal{I}^i\mathcal{F}$
的局部截面 $gs$。
当然，$gs$ 属于 $\eta_\mathcal{I}\mathcal{F}^\bullet$ 上
$d^i$ 的核，因而定义
$H^i(\eta_\mathcal{I}\mathcal{F}^\bullet)$ 的局部截面。
检验其适定性没有困难。
我们断言，该态射经由引理中给出的同构分解。
这可在茎上检验，因此通过引理 \ref{cohomology-lemma-eta-stalks}，
它转化为更多代数，引理
\ref{more-algebra-lemma-eta-first-property} 的结果。
\end{proof}

\begin{lemma}
\label{cohomology-lemma-eta-qis}
在情形 \ref{cohomology-situation-eta} 中，
令 $\mathcal{F}^\bullet \to \mathcal{G}^\bullet$ 为
由 $\mathcal{I}$-无挠 $\mathcal{O}_X$-模组成的复形之间的态射。
那么诱导态射
$\eta_\mathcal{I}\mathcal{F}^\bullet \to \eta_\mathcal{I}\mathcal{G}^\bullet$
也是拟同构。
\end{lemma}

\begin{proof}
这是因为引理 \ref{cohomology-lemma-eta-first-property} 的同构
与复形态射相容。
\end{proof}


\begin{lemma}
\label{cohomology-lemma-Leta}
在情形 \ref{cohomology-situation-eta} 中，存在加性
函子\footnote{注意，此函子并不正合，即不把
特异三角变为特异三角。}
$L\eta_\mathcal{I} : D(\mathcal{O}_X) \to D(\mathcal{O}_X)$，
使得如果 $D(\mathcal{O}_X)$ 中的 $M$ 由
$\mathcal{I}$-无挠 $\mathcal{O}_X$-模复形
$\mathcal{F}^\bullet$ 表示，那么
$L\eta_\mathcal{I}M = \eta_\mathcal{I}\mathcal{F}^\bullet$。
态射也同样如此。
\end{lemma}

\begin{proof}
记 $\mathcal{T} \subset \textit{Mod}(\mathcal{O}_X)$ 为
$\mathcal{I}$-无挠 $\mathcal{O}_X$-模组成的全子范畴。
相应地，有含入
$$
K(\mathcal{T})
\quad\subset\quad
K(\textit{Mod}(\mathcal{O}_X)) = K(\mathcal{O}_X)
$$
它把 $K(\mathcal{T})$ 视为 $K(\mathcal{O}_X)$ 的全三角子范畴。
令 $S \subset \text{Arrows}(K(\mathcal{T}))$ 为拟同构组成的集合。
我们将应用
导出范畴，引理 \ref{derived-lemma-localization-subcategory}
证明态射
$$
S^{-1}K(\mathcal{T}) \longrightarrow D(\mathcal{O}_X)
$$
是三角范畴的等价。该引理说明，只需证明：
给定 $\mathcal{O}_X$-模复形 $\mathcal{G}^\bullet$，
存在拟同构 $\mathcal{F}^\bullet \to \mathcal{G}^\bullet$，
其中 $\mathcal{F}^\bullet$ 是由 $\mathcal{I}$-无挠
$\mathcal{O}_X$-模组成的复形。
由引理 \ref{cohomology-lemma-K-flat-resolution}，可找到拟同构
$\mathcal{F}^\bullet \to \mathcal{G}^\bullet$，使得复形
$\mathcal{F}^\bullet$ 是 K-平坦的（我们不会用到这一点），并且
由平坦 $\mathcal{O}_X$-模 $\mathcal{F}^i$ 组成。
由引理 \ref{cohomology-lemma-I-torsion-free} 的第三个刻画，
平坦 $\mathcal{O}_X$-模是
$\mathcal{I}$-无挠 $\mathcal{O}_X$-模，证明完成。

\medskip\noindent
有了这些准备，就可以定义 $L\eta_f$。
具体地，由本节引理 \ref{cohomology-lemma-I-torsion-free} 后面的讨论，
我们已经有良定义函子
$$
K(\mathcal{T}) \xrightarrow{\eta_f} K(\mathcal{T}) \to
K(\mathcal{O}_X) \to D(\mathcal{O}_X)
$$
按照引理 \ref{cohomology-lemma-eta-qis}，它把拟同构
送到拟同构。因此，由范畴，引理
\ref{categories-lemma-properties-left-localization}，
该函子经由
$S^{-1}K(\mathcal{T}) = D(\mathcal{O}_X)$ 分解。
\end{proof}

\noindent
在情形 \ref{cohomology-situation-eta} 中，我们构造 Bockstein 算子。
首先注意有交换图
$$
\xymatrix{
0 \ar[r] &
\mathcal{I}^{i + 1} \ar[r] \ar[d] &
\mathcal{I}^i \ar[r] \ar[d] &
\mathcal{I}^i/\mathcal{I}^{i + 1}
\ar[r] \ar@{=}[d] &
0 \\
0 \ar[r] &
\mathcal{I}^{i + 1}/\mathcal{I}^{i + 2} \ar[r] &
\mathcal{I}^i/\mathcal{I}^{i + 2} \ar[r] &
\mathcal{I}^i/\mathcal{I}^{i + 1} \ar[r] &
0
}
$$
其两行都是 $\mathcal{O}_X$-模的短正合列。
令 $M$ 为 $D(\mathcal{O}_X)$ 的对象。
把上述图与 $M$ 作张量积，得到特异三角的态射
$$
\xymatrix{
M \otimes^\mathbf{L} \mathcal{I}^{i + 1} \ar[r] \ar[d] &
M \otimes^\mathbf{L} \mathcal{I}^i \ar[r] \ar[d] &
M \otimes^\mathbf{L} \mathcal{I}^i/\mathcal{I}^{i + 1} \ar[d]^{\text{id}} \\
M \otimes^\mathbf{L} \mathcal{I}^{i + 1}/\mathcal{I}^{i + 2} \ar[r] &
M \otimes^\mathbf{L} \mathcal{I}^i/\mathcal{I}^{i + 2} \ar[r] &
M \otimes^\mathbf{L} \mathcal{I}^i/\mathcal{I}^{i + 1}
}
$$
尤其是，与下方三角相伴的上同调层长正合列
确定{\it Bockstein 算子}
$$
\beta = \beta^i :
H^i(M \otimes^\mathbf{L} \mathcal{I}^i/\mathcal{I}^{i + 1})
\longrightarrow
H^{i + 1}(M \otimes^\mathbf{L} \mathcal{I}^{i + 1}/\mathcal{I}^{i + 2})
$$
其中 $i \in \mathbf{Z}$。为供后用，我们在此记录：
由上述交换图，Bockstein 算子有分解
\begin{equation}
\label{cohomology-equation-factorization-bockstein}
\vcenter{
\xymatrix{
H^i(M \otimes^\mathbf{L} \mathcal{I}^i/\mathcal{I}^{i + 1})
\ar[r]_\delta \ar[rd]_\beta &
H^{i + 1}(M \otimes^\mathbf{L} \mathcal{I}^{i + 1}) \ar[d] \\
&
H^{i + 1}(M \otimes^\mathbf{L} \mathcal{I}^{i + 1}/\mathcal{I}^{i + 2})
}
}
\end{equation}
其中 $\delta$ 是来自上述交换图中上方特异三角的边界算子。
得到复形
\begin{equation}
\label{cohomology-equation-complex-bocksteins}
H^\bullet(M/\mathcal{I}) =
\left[
\begin{matrix}
\ldots \\
\downarrow \\
H^{i - 1}(M \otimes^\mathbf{L} \mathcal{I}^{i - 1}/\mathcal{I}^i) \\
\downarrow \beta \\
H^i(M \otimes^\mathbf{L} \mathcal{I}^i/\mathcal{I}^{i + 1}) \\
\downarrow \beta \\
H^{i + 1}(M \otimes^\mathbf{L} \mathcal{I}^{i + 1}/\mathcal{I}^{i + 2}) \\
\downarrow \\
\ldots
\end{matrix}
\right]
\end{equation}
也就是说，$\beta \circ \beta = 0$。这可以在茎上检验，
此时可由更多代数，第 \ref{more-algebra-section-eta} 节中
证明的相应代数结果推出。
另一种证明：短正合列
$0 \to \mathcal{I}^{i + 1}/\mathcal{I}^{i + 2}
\to \mathcal{I}^i/\mathcal{I}^{i + 2}
\to \mathcal{I}^i/\mathcal{I}^{i + 1} \to 0$
在 $D(\mathcal{O}_X)$ 中定义态射
$b^i : \mathcal{I}^i/\mathcal{I}^{i + 1} \to
(\mathcal{I}^{i + 1}/\mathcal{I}^{i + 2})[1]$；
这些态射与 $M$ 作张量积并取上同调层后，
诱导上述态射 $\beta$。
然后利用导出范畴，引理
\ref{derived-lemma-cup-ext-1-zero} 中的判据，以及
模 $\mathcal{I}^i/\mathcal{I}^{i + 3}$
是 $\mathcal{I}^{i + 1}/\mathcal{I}^{i + 3}$ 被
$\mathcal{I}^i/\mathcal{I}^{i + 1}$ 的扩张这一事实，可证明复合
$b^{i + 1}[1] \circ b^i : \mathcal{I}^i/\mathcal{I}^{i + 1} \to
(\mathcal{I}^{i + 1}/\mathcal{I}^{i + 2})[1] \to
(\mathcal{I}^{i + 2}/\mathcal{I}^{i + 3})[2]$
在 $D(\mathcal{O}_X)$ 中为零。


\begin{lemma}
\label{cohomology-lemma-eta-second-property}
在情形 \ref{cohomology-situation-eta} 中，令 $M$ 为
$D(\mathcal{O}_X)$ 的对象。存在 $D(\mathcal{O}_X)$ 中的典范同构
$$
L\eta_\mathcal{I}M \otimes^\mathbf{L} \mathcal{O}_X/\mathcal{I}
\longrightarrow
H^\bullet(M/\mathcal{I})
$$
其中右端是复形
(\ref{cohomology-equation-complex-bocksteins})。
\end{lemma}

\begin{proof}
根据引理 \ref{cohomology-lemma-eta-qis} 中 $L\eta_\mathcal{I}$ 的构造，
可假设 $M$ 由 $\mathcal{I}$-无挠
$\mathcal{O}_X$-模复形 $\mathcal{F}^\bullet$ 表示。那么
$L\eta_\mathcal{I}M$ 由复形
$\eta_\mathcal{I}\mathcal{F}^\bullet$ 表示，而后者同样
由 $\mathcal{I}$-无挠 $\mathcal{O}_X$-模组成。
因此 $L\eta_\mathcal{I}M \otimes^\mathbf{L} \mathcal{O}_X/\mathcal{I}$
由复形
$\eta_\mathcal{I}\mathcal{F}^\bullet \otimes \mathcal{O}_X/\mathcal{I}$
表示。类似地，复形 $H^\bullet(M/\mathcal{I})$ 的各项为
$H^i(\mathcal{F}^\bullet \otimes \mathcal{I}^i/\mathcal{I}^{i + 1})$。

\medskip\noindent
令 $f$ 为 $\mathcal{I}$ 的局部生成元。
令 $s$ 为 $(\eta_\mathcal{I}\mathcal{F})^i$ 的局部截面。
那么可写成 $s = f^is'$，其中 $s'$ 是
$\mathcal{F}^i$ 的局部截面；类似地，
$d^i(s) = f^{i + 1}t$，其中 $t$ 是
$\mathcal{F}^{i + 1}$ 的局部截面。因此 $d^i$ 把 $f^is'$
映到 $\mathcal{F}^{i + 1} \otimes \mathcal{I}^i/\mathcal{I}^{i + 1}$
中的零。
故可把 $s$ 映到
$H^i(\mathcal{F}^\bullet \otimes \mathcal{I}^i/\mathcal{I}^{i + 1})$
中 $f^is'$ 的类。
该规则定义 $\mathcal{O}_X$-模态射
$$
(\eta_\mathcal{I}\mathcal{F})^i \otimes \mathcal{O}_X/\mathcal{I}
\longrightarrow
H^i(\mathcal{F}^\bullet \otimes \mathcal{I}^i/\mathcal{I}^{i + 1})
$$
计算表明，这些态射与微分相容（本质上因为
$\beta$ 把 $f^is'$ 的类送到 $f^{i + 1}t$ 的类），
从而得到一个表示引理陈述中箭头的复形态射。

\medskip\noindent
为完成证明，注意上一段给出的构造在茎上
与更多代数，引理
\ref{more-algebra-lemma-eta-second-property}
中构造的态射一致，结论随即得出。
\end{proof}

\begin{lemma}
\label{cohomology-lemma-eta-tensor-invertible}
在情形 \ref{cohomology-situation-eta} 中，
令 $\mathcal{F}^\bullet$ 为由
$\mathcal{I}$-无挠 $\mathcal{O}_X$-模组成的复形。
令 $\mathcal{L}$ 为可逆 $\mathcal{O}_X$-模。
那么 $\eta_\mathcal{I}(\mathcal{F}^\bullet \otimes \mathcal{L}) =
(\eta_\mathcal{I}\mathcal{F}^\bullet) \otimes \mathcal{L}$。
\end{lemma}

\begin{proof}
由构造立即可得。
\end{proof}

\begin{lemma}
\label{cohomology-lemma-eta-cohomology-locally-free}
在情形 \ref{cohomology-situation-eta} 中，令 $M$ 为 $D(\mathcal{O}_X)$ 的对象。
令 $x \in X$ 且 $\mathcal{O}_{X, x}$ 非零。如果 $H^i(M)_x$
是有限自由 $\mathcal{O}_{X, x}$-模，那么
$H^i(L\eta_\mathcal{I}M)_x$ 是同秩的
有限自由 $\mathcal{O}_{X, x}$-模。
\end{lemma}

\begin{proof}
具体地，设 $f \in \mathcal{O}_{X, x}$ 生成茎 $\mathcal{I}_x$。
那么 $f$ 是 $\mathcal{O}_{X, x}$ 中的非零因子，因而
$H^i(M)_x[f] = 0$。因此由引理 \ref{cohomology-lemma-eta-first-property}，
$H^i(L\eta_\mathcal{I}M)_x$ 同构于
$\mathcal{I}^i_x \otimes_{\mathcal{O}_{X, x}} H^i(M)_x$，
后者如所需地是同秩自由模。
\end{proof}
% Generated mechanically from the frozen sites-cohomology.tex target.
% Do not edit this generated file; edit the bound locale source instead.

% 好，从这里开始。
%

\setcounter{chapter}{20}
\chapter{站点上的上同调}



\phantomsection
\label{sites-cohomology-section-phantom}


\section{引言}
\label{sites-cohomology-section-introduction}

\noindent
本文讨论层上同调的若干主题。
我们考察站点上的层会发生什么，
不过许多时候只是把拓扑情形中的讨论、
构造和证明在这里重述一遍。
基本参考文献为 \cite{SGA4}、\cite{Godement} 和 \cite{Iversen}。




\section{层上同调}
\label{sites-cohomology-section-cohomology-sheaves}

\noindent
设 $\mathcal{C}$ 为站点，见《站点》定义 \ref{sites-definition-site}。
设 $\mathcal{F}$ 为 $\mathcal{C}$ 上的阿贝尔层。
我们知道 $\mathcal{C}$ 上的阿贝尔层范畴有足够多的内射对象，
见《内射对象》定理 \ref{injectives-theorem-sheaves-injectives}。
因此可以选取一个内射分解
$\mathcal{F}[0] \to \mathcal{I}^\bullet$.
对站点 $\mathcal{C}$ 的任意对象 $U$，定义
\begin{equation}
\label{sites-cohomology-equation-cohomology-object-site}
H^i(U, \mathcal{F}) = H^i(\Gamma(U, \mathcal{I}^\bullet))
\end{equation}
为{\it 阿贝尔层 $\mathcal{F}$ 在对象 $U$ 上的第 $i$ 个上同调群}。
换言之，它们是函子 $\mathcal{F} \mapsto \mathcal{F}(U)$ 的右导出函子。
函子族 $H^i(U, -)$ 构成万有 $\delta$-函子
$\textit{Ab}(\mathcal{C}) \to \textit{Ab}$.

\medskip\noindent
站点 $\mathcal{C}$ 有时没有终对象。在这种情形下，
我们把 $\mathcal{C}$ 上集合预层 $\mathcal{F}$ 的{\it 整体截面}定义为集合
\begin{equation}
\label{sites-cohomology-equation-global-sections}
\Gamma(\mathcal{C}, \mathcal{F}) =
\Mor_{\textit{PSh}(\mathcal{C})}(e, \mathcal{F})
\end{equation}
其中 $e$ 是 $\mathcal{C}$ 上预层范畴的终对象。
在这种情形下，给定 $\mathcal{C}$ 上的阿贝尔层 $\mathcal{F}$，
如下定义{\it $\mathcal{F}$ 在 $\mathcal{C}$ 上的第 $i$ 个上同调群}：
\begin{equation}
\label{sites-cohomology-equation-cohomology}
H^i(\mathcal{C}, \mathcal{F}) = H^i(\Gamma(\mathcal{C}, \mathcal{I}^\bullet))
\end{equation}
换言之，它是整体截面函子的第 $i$ 个右导出函子。
函子族 $H^i(\mathcal{C}, -)$ 构成万有 $\delta$-函子
$\textit{Ab}(\mathcal{C}) \to \textit{Ab}$.

\medskip\noindent
设 $f : \Sh(\mathcal{C}) \to \Sh(\mathcal{D})$ 为拓扑斯的态射，
见《站点》定义 \ref{sites-definition-topos}。
对上述 $\mathcal{F}[0] \to \mathcal{I}^\bullet$，定义
\begin{equation}
\label{sites-cohomology-equation-higher-direct-image}
R^if_*\mathcal{F} = H^i(f_*\mathcal{I}^\bullet)
\end{equation}
为{\it $\mathcal{F}$ 的第 $i$ 个高阶直接像}。
它们是 $f_*$ 的右导出函子。
函子族 $R^if_*$ 构成从
$\textit{Ab}(\mathcal{C}) \to \textit{Ab}(\mathcal{D})$ 的万有 $\delta$-函子。

\medskip\noindent
设 $(\mathcal{C}, \mathcal{O})$ 为环化站点，
见《站点上的模》定义 \ref{sites-modules-definition-ringed-site}。
设 $\mathcal{F}$ 为 $\mathcal{O}$-模。
我们知道 $\mathcal{O}$-模范畴有足够多的内射对象，
见《内射对象》定理 \ref{injectives-theorem-sheaves-modules-injectives}。
因此可以选取一个内射分解
$\mathcal{F}[0] \to \mathcal{I}^\bullet$.
对站点 $\mathcal{C}$ 的任意对象 $U$，定义
\begin{equation}
\label{sites-cohomology-equation-cohomology-object-site-modules}
H^i(U, \mathcal{F}) = H^i(\Gamma(U, \mathcal{I}^\bullet))
\end{equation}
为{\it $\mathcal{F}$ 在 $U$ 上的第 $i$ 个上同调群}。
函子族 $H^i(U, -)$ 构成万有 $\delta$-函子
$\textit{Mod}(\mathcal{O}) \to \text{Mod}_{\mathcal{O}(U)}$。类似地，
\begin{equation}
\label{sites-cohomology-equation-cohomology-modules}
H^i(\mathcal{C}, \mathcal{F}) = H^i(\Gamma(\mathcal{C}, \mathcal{I}^\bullet))
\end{equation}
它是{\it $\mathcal{F}$ 在 $\mathcal{C}$ 上的第 $i$ 个上同调群}。
函子族 $H^i(\mathcal{C}, -)$ 构成万有 $\delta$-函子
$\textit{Mod}(\mathcal{C}) \to \text{Mod}_{\Gamma(\mathcal{C}, \mathcal{O})}$.

\medskip\noindent
设 $f : (\Sh(\mathcal{C}), \mathcal{O}) \to (\Sh(\mathcal{D}), \mathcal{O}')$
为环化拓扑斯的态射，
见《站点上的模》定义 \ref{sites-modules-definition-ringed-topos}。
对上述 $\mathcal{F}[0] \to \mathcal{I}^\bullet$，定义
\begin{equation}
\label{sites-cohomology-equation-higher-direct-image-modules}
R^if_*\mathcal{F} = H^i(f_*\mathcal{I}^\bullet)
\end{equation}
为{\it $\mathcal{F}$ 的第 $i$ 个高阶直接像}。
它们是 $f_*$ 的右导出函子。
函子族 $R^if_*$ 构成从
$\textit{Mod}(\mathcal{O}) \to \textit{Mod}(\mathcal{O}')$ 的万有 $\delta$-函子。








\section{导出函子}
\label{sites-cohomology-section-derived-functors}

\noindent
我们简要说明一种利用分解函子处理右导出函子的方法。
设 $(\mathcal{C}, \mathcal{O})$ 为环化站点。本章记
$$
K(\mathcal{O}) = K(\textit{Mod}(\mathcal{O}))
\quad
\text{和}
\quad
D(\mathcal{O}) = D(\textit{Mod}(\mathcal{O}))
$$
；对于《导出范畴》定义 \ref{derived-definition-complexes-notation}
和定义 \ref{derived-definition-unbounded-derived-category} 中引入的三角范畴，
其有界版本也采用类似记号。
由《导出范畴》注记 \ref{derived-remark-big-abelian-category}，
存在分解函子
$$
j = j_{(\mathcal{C}, \mathcal{O})} :
K^{+}(\textit{Mod}(\mathcal{O}))
\longrightarrow
K^{+}(\mathcal{I})
$$
其中 $\mathcal{I}$ 是 $\textit{Mod}(\mathcal{O})$ 中由内射
$\mathcal{O}$-模组成的严格满加性子范畴。
对任意阿贝尔范畴 $\mathcal{B}$ 以及任意左正合函子
$F : \textit{Mod}(\mathcal{O}) \to \mathcal{B}$，
记 $RF$ 为《导出范畴》第 \ref{derived-section-right-derived-functor} 节中
用上述分解函子 $j$ 构造的右导出函子：
\begin{equation}
\label{sites-cohomology-equation-RF}
RF = F \circ j' : D^{+}(\mathcal{O}) \longrightarrow D^{+}(\mathcal{B})
\end{equation}
记号见《导出范畴》引理 \ref{derived-lemma-right-derived-functor}。
注意，视情形而定，可以把 $RF$ 看作定义在
$\textit{Mod}(\mathcal{O})$、$\text{Comp}^{+}(\textit{Mod}(\mathcal{O}))$ 或
$K^{+}(\mathcal{O})$ 上。按照《导出范畴》定义
\ref{derived-definition-higher-derived-functors}，得到第 $i$ 个右导出函子
\begin{equation}
\label{sites-cohomology-equation-RFi}
R^iF = H^i \circ RF : \textit{Mod}(\mathcal{O}) \longrightarrow \mathcal{B}
\end{equation}
于是 $R^0F = F$，且 $\{R^iF, \delta\}_{i \geq 0}$ 是万有
$\delta$-函子；见《导出范畴》引理
\ref{derived-lemma-higher-derived-functors}。

\medskip\noindent
下面给出该构造的两个特例。给定环 $R$，记
$K(R) = K(\text{Mod}_R)$ 和 $D(R) = D(\text{Mod}_R)$，
其有界版本也类似。对 $\mathcal{C}$ 的任意对象 $U$，有左正合函子
$
\Gamma(U, -) :
\textit{Mod}(\mathcal{O})
\longrightarrow
\text{Mod}_{\mathcal{O}(U)}
$
由上述讨论，它给出
$$
R\Gamma(U, -) :
D^{+}(\mathcal{O})
\longrightarrow
D^{+}(\mathcal{O}(U))
$$
注意，$H^i(U, -) = R^i\Gamma(U, -)$ 与上面的
(\ref{sites-cohomology-equation-cohomology-object-site-modules}) 相容。类似地，有
$$
R\Gamma(\mathcal{C}, -) :
D^{+}(\mathcal{O})
\longrightarrow
D^{+}(\Gamma(\mathcal{C}, \mathcal{O}))
$$
它与 (\ref{sites-cohomology-equation-cohomology-modules}) 相容。若
$f : (\Sh(\mathcal{C}), \mathcal{O}) \to (\Sh(\mathcal{D}), \mathcal{O}')$
是环化拓扑斯的态射，则得到左正合函子
$f_* : \textit{Mod}(\mathcal{O}) \to \textit{Mod}(\mathcal{O}')$
，它给出{\it 导出推前}
$$
Rf_* : D^{+}(\mathcal{O}) \to D^+(\mathcal{O}')
$$
按照 (\ref{sites-cohomology-equation-higher-direct-image-modules})，
$Rf_*\mathcal{F}^\bullet$ 的第 $i$ 个上同调层记作
$R^if_*\mathcal{F}^\bullet$，称为第 $i$ 个{\it 高阶直接像}。
上面所示的函子都是导出范畴之间的正合函子。







\section{第一上同调与挠子}
\label{sites-cohomology-section-h1-torsors}

\begin{definition}
\label{sites-cohomology-definition-torsor}
设 $\mathcal{C}$ 为站点，$\mathcal{G}$ 为 $\mathcal{C}$ 上的
（不必交换的）群层。{\it 伪挠子}，更确切地说，
{\it 伪 $\mathcal{G}$-挠子}，是 $\mathcal{C}$ 上的集合层
$\mathcal{F}$，配备作用
$\mathcal{G} \times \mathcal{F} \to \mathcal{F}$，使得
\begin{enumerate}
\item 只要 $\mathcal{F}(U)$ 非空，作用
$\mathcal{G}(U) \times \mathcal{F}(U) \to \mathcal{F}(U)$
就是单传递的。
\end{enumerate}
{\it 伪 $\mathcal{G}$-挠子的态射}
$\mathcal{F} \to \mathcal{F}'$，是与 $\mathcal{G}$-作用相容的集合层态射。
{\it 挠子}，更确切地说，{\it $\mathcal{G}$-挠子}，
是还满足下列条件的伪 $\mathcal{G}$-挠子：
\begin{enumerate}
\item[(2)] 对每个 $U \in \Ob(\mathcal{C})$，存在 $U$ 的覆盖
$\{U_i \to U\}_{i \in I}$，使得对所有 $i \in I$，
$\mathcal{F}(U_i)$ 均非空。
\end{enumerate}
{\it $\mathcal{G}$-挠子的态射}就是伪 $\mathcal{G}$-挠子的态射。
{\it 平凡 $\mathcal{G}$-挠子}是配备显然左 $\mathcal{G}$-作用的层
$\mathcal{G}$。
\end{definition}

\noindent
显然，挠子的态射自动是同构。

\begin{lemma}
\label{sites-cohomology-lemma-trivial-torsor}
设 $\mathcal{C}$ 为站点，$\mathcal{G}$ 为 $\mathcal{C}$ 上的
（不必交换的）群层。$\mathcal{G}$-挠子 $\mathcal{F}$ 平凡，
当且仅当
$\Gamma(\mathcal{C}, \mathcal{F}) \not = \emptyset$.
\end{lemma}

\begin{proof}
略。
\end{proof}

\begin{lemma}
\label{sites-cohomology-lemma-torsors-h1}
设 $\mathcal{C}$ 为站点，$\mathcal{H}$ 为 $\mathcal{C}$ 上的阿贝尔层。
$\mathcal{H}$-挠子的同构类集合与 $H^1(\mathcal{C}, \mathcal{H})$
之间存在典范双射。
\end{lemma}

\begin{proof}
设 $\mathcal{F}$ 为 $\mathcal{H}$-挠子。考虑 $\mathcal{F}$ 上的自由阿贝尔层
$\mathbf{Z}[\mathcal{F}]$。它是如下规则的层化：对
$U \in \Ob(\mathcal{C})$，赋予所有有限形式和
$\sum n_i[s_i]$ 的集合，其中 $n_i \in \mathbf{Z}$ 且
$s_i \in \mathcal{F}(U)$。有自然映射
$$
\sigma : \mathbf{Z}[\mathcal{F}] \longrightarrow \underline{\mathbf{Z}}
$$
，它把局部截面 $\sum n_i[s_i]$ 映到 $\sum n_i$。
$\sigma$ 的核由形如 $[s] - [s']$ 的截面生成。有典范映射
$a : \Ker(\sigma) \to \mathcal{H}$
，它把 $[s] - [s']$ 映到 $h$，其中 $h$ 是满足
$h \cdot s' = s$ 的 $\mathcal{H}$ 的局部截面。考虑推出图
$$
\xymatrix{
0 \ar[r] &
\Ker(\sigma) \ar[r] \ar[d]^a &
\mathbf{Z}[\mathcal{F}] \ar[r] \ar[d] &
\underline{\mathbf{Z}} \ar[r] \ar[d] &
0 \\
0 \ar[r] &
\mathcal{H} \ar[r] &
\mathcal{E} \ar[r] &
\underline{\mathbf{Z}} \ar[r] &
0
}
$$
这里 $\mathcal{E}$ 是由推出得到的扩张。对下方短正合列所对应的
上同调长正合列，把边界算子作用于
$1 \in H^0(\mathcal{C}, \underline{\mathbf{Z}})$，得到元素
$\xi = \xi_\mathcal{F} \in H^1(\mathcal{C}, \mathcal{H})$
。

\medskip\noindent
反过来，给定 $\xi \in H^1(\mathcal{C}, \mathcal{H})$，
可如下为 $\xi$ 配置一个挠子。选取 $\mathcal{H}$ 到内射阿贝尔层
$\mathcal{I}$ 中的嵌入 $\mathcal{H} \to \mathcal{I}$。
令 $\mathcal{Q} = \mathcal{I}/\mathcal{H}$，于是有短正合列
$$
\xymatrix{
0 \ar[r] &
\mathcal{H} \ar[r] &
\mathcal{I} \ar[r] &
\mathcal{Q} \ar[r] &
0
}
$$
元素 $\xi$ 是某个整体截面
$q \in H^0(\mathcal{C}, \mathcal{Q})$
的像，因为 $H^1(\mathcal{C}, \mathcal{I}) = 0$（见《导出范畴》引理
\ref{derived-lemma-higher-derived-functors}）。
令 $\mathcal{F} \subset \mathcal{I}$ 为由那些在层 $\mathcal{Q}$ 中
映到 $q$ 的截面组成的（集合）子层。容易验证，
$\mathcal{F}$ 是 $\mathcal{H}$-挠子。

\medskip\noindent
上述两个构造互逆的验证留给读者。
\end{proof}






\section{第一上同调与扩张}
\label{sites-cohomology-section-h1-extensions}

\begin{lemma}
\label{sites-cohomology-lemma-h1-extensions}
设 $(\mathcal{C}, \mathcal{O})$ 为环化站点，
$\mathcal{F}$ 为 $\mathcal{C}$ 上的 $\mathcal{O}$-模层。存在典范双射
$$
\Ext^1_{\textit{Mod}(\mathcal{O})}(\mathcal{O}, \mathcal{F})
\longrightarrow
H^1(\mathcal{C}, \mathcal{F})
$$
，它把扩张
$$
0 \to \mathcal{F} \to \mathcal{E} \to \mathcal{O} \to 0
$$
映到 $1 \in \Gamma(\mathcal{C}, \mathcal{O})$ 在
$H^1(\mathcal{C}, \mathcal{F})$ 中的像。
\end{lemma}

\begin{proof}
我们构造引理中映射的逆。设 $\xi \in H^1(\mathcal{C}, \mathcal{F})$。
选取单射 $\mathcal{F} \subset \mathcal{I}$，其中 $\mathcal{I}$ 在
$\textit{Mod}(\mathcal{O})$ 中内射。令
$\mathcal{Q} = \mathcal{I}/\mathcal{F}$。由上同调长正合列可知，
$\xi$ 是某个截面
$\tilde \xi \in \Gamma(\mathcal{C}, \mathcal{Q}) =
\Hom_\mathcal{O}(\mathcal{O}, \mathcal{Q})$.
的像。现在只需作拉回
$$
\xymatrix{
0 \ar[r] &
\mathcal{F} \ar[r] \ar@{=}[d] &
\mathcal{E} \ar[r] \ar[d] &
\mathcal{O} \ar[r] \ar[d]^{\tilde \xi} &
0 \\
0 \ar[r] &
\mathcal{F} \ar[r] &
\mathcal{I} \ar[r] &
\mathcal{Q} \ar[r] &
0
}
$$
；见《同调代数》第 \ref{homology-section-extensions} 节。
\end{proof}

\noindent
下面的引理将由更一般的引理
\ref{sites-cohomology-lemma-cohomology-modules-abelian-agree} 取代。

\begin{lemma}
\label{sites-cohomology-lemma-h1-mod-ab-agree}
设 $(\mathcal{C}, \mathcal{O})$ 为环化站点，
$\mathcal{F}$ 为 $\mathcal{C}$ 上的 $\mathcal{O}$-模层。
以 $\mathcal{F}_{ab}$ 表示其底阿贝尔群层。则有函子性同构
$$
H^1(\mathcal{C}, \mathcal{F}_{ab})
=
H^1(\mathcal{C}, \mathcal{F})
$$
，其中左端是在 $\textit{Ab}(\mathcal{C})$ 中计算的上同调，
右端是在 $\textit{Mod}(\mathcal{O})$ 中计算的上同调。
\end{lemma}

\begin{proof}
以 $\underline{\mathbf{Z}}$ 表示常值层 $\mathbf{Z}$。由于
$\textit{Ab}(\mathcal{C}) = \textit{Mod}(\underline{\mathbf{Z}})$
，可以两次应用引理 \ref{sites-cohomology-lemma-h1-extensions}，
因而只需证明
$$
\Ext^1_{\textit{Mod}(\mathcal{O})}(\mathcal{O}, \mathcal{F})
=
\Ext^1_{\textit{Mod}(\underline{\mathbf{Z}})}(
\underline{\mathbf{Z}}, \mathcal{F}_{ab}).
$$
设 $0 \to \mathcal{F} \to \mathcal{E} \to \mathcal{O} \to 0$
是 $\textit{Mod}(\mathcal{O})$ 中的扩张。利用显然的阿贝尔层态射
$1 : \underline{\mathbf{Z}} \to \mathcal{O}$
作拉回，得到如下扩张 $\mathcal{E}_{ab}$：
$$
\xymatrix{
0 \ar[r] &
\mathcal{F}_{ab} \ar[r] \ar@{=}[d] &
\mathcal{E}_{ab} \ar[r] \ar[d] &
\underline{\mathbf{Z}} \ar[r] \ar[d]^{1} &
0 \\
0 \ar[r] &
\mathcal{F} \ar[r] &
\mathcal{E} \ar[r] &
\mathcal{O} \ar[r] &
0
}
$$
反向构造稍微有趣一些。设
$0 \to \mathcal{F}_{ab} \to \mathcal{E}_{ab} \to \underline{\mathbf{Z}} \to 0$
是 $\textit{Mod}(\underline{\mathbf{Z}})$ 中的扩张。
由于 $\underline{\mathbf{Z}}$ 是平坦的 $\underline{\mathbf{Z}}$-模，
可知序列
$$
0 \to \mathcal{F}_{ab} \otimes_{\underline{\mathbf{Z}}} \mathcal{O}
\to \mathcal{E}_{ab} \otimes_{\underline{\mathbf{Z}}} \mathcal{O}
\to \underline{\mathbf{Z}} \otimes_{\underline{\mathbf{Z}}} \mathcal{O}
\to 0
$$
正合；见《站点上的模》引理 \ref{sites-modules-lemma-flat-tor-zero}。
当然，
$\underline{\mathbf{Z}} \otimes_{\underline{\mathbf{Z}}} \mathcal{O}
= \mathcal{O}$.
因此可以沿（$\mathcal{O}$-线性的）乘法映射
$\mu : \mathcal{F} \otimes_{\underline{\mathbf{Z}}} \mathcal{O}
\to \mathcal{F}$ 作推出，得到如下 $\mathcal{O}$ 由
$\mathcal{F}$ 的扩张：
$$
\xymatrix{
0 \ar[r] &
\mathcal{F}_{ab} \otimes_{\underline{\mathbf{Z}}} \mathcal{O}
\ar[r] \ar[d]^\mu &
\mathcal{E}_{ab} \otimes_{\underline{\mathbf{Z}}} \mathcal{O}
\ar[r] \ar[d] &
\mathcal{O} \ar[r] \ar@{=}[d] &
0 \\
0 \ar[r] &
\mathcal{F} \ar[r] &
\mathcal{E} \ar[r] &
\mathcal{O} \ar[r] &
0
}
$$
这就是所需扩张。这两个构造互逆的验证从略。
\end{proof}





\section{第一上同调与可逆层}
\label{sites-cohomology-section-invertible-sheaves}

\noindent
环化站点的 Picard 群定义于《站点上的模》第
\ref{sites-modules-section-invertible} 节。

\begin{lemma}
\label{sites-cohomology-lemma-h1-invertible}
设 $(\mathcal{C}, \mathcal{O})$ 为局部环化站点。存在阿贝尔群的典范同构
$$
H^1(\mathcal{C}, \mathcal{O}^*) = \Pic(\mathcal{O}).
$$
。
\end{lemma}

\begin{proof}
设 $\mathcal{L}$ 为可逆 $\mathcal{O}$-模。考虑由下列规则定义的预层
$\mathcal{L}^*$：
$$
U \longmapsto \{s \in \mathcal{L}(U)
\text{使得 } \mathcal{O}_U \xrightarrow{s \cdot -} \mathcal{L}_U
\text{ 为同构}\}
$$
这个预层满足层条件。此外，若 $f \in \mathcal{O}^*(U)$ 且
$s \in \mathcal{L}^*(U)$，则显然 $fs \in \mathcal{L}^*(U)$。
同理，若 $s, s' \in \mathcal{L}^*(U)$，则存在唯一的
$f \in \mathcal{O}^*(U)$ 使得 $fs = s'$。再者，由《站点上的模》引理
\ref{sites-modules-lemma-invertible-is-locally-free-rank-1}，
层 $\mathcal{L}^*$ 局部有截面。换言之，
$\mathcal{L}^*$ 是 $\mathcal{O}^*$-挠子。因此得到映射
$$
\begin{matrix}
\text{在 }(\mathcal{C}, \mathcal{O})\text{ 上的可逆层} \\
\text{的同构类集合}
\end{matrix}
\longrightarrow
\begin{matrix}
\mathcal{O}^*\text{-挠子} \\
\text{的同构类集合}
\end{matrix}
$$
略去验证这是阿贝尔群同态。由引理 \ref{sites-cohomology-lemma-torsors-h1}，
右端与 $H^1(\mathcal{C}, \mathcal{O}^*)$ 典范双射。
因此只需证明这个映射既单又满。

\medskip\noindent
单射性。若挠子 $\mathcal{L}^*$ 平凡，则由引理
\ref{sites-cohomology-lemma-trivial-torsor}，$\mathcal{L}^*$ 有整体截面。
这恰好意味着 $\mathcal{L} \cong \mathcal{O}$ 是
$\Pic(\mathcal{O})$ 中的单位元。

\medskip\noindent
满射性。设 $\mathcal{F}$ 为 $\mathcal{O}^*$-挠子。考虑集合预层
$$
\mathcal{L}_1 : U \longmapsto
(\mathcal{F}(U) \times \mathcal{O}(U))/\mathcal{O}^*(U)
$$
其中 $f \in \mathcal{O}^*(U)$ 在 $(s, g)$ 上的作用为
$(fs, f^{-1}g)$。规定
$(s, g) + (s', g') = (s, g + (s'/s)g')$，其中 $s'/s$ 是满足
$fs = s'$ 的 $\mathcal{O}^*$ 的局部截面 $f$；并对
$\mathcal{O}$ 的局部截面 $h$ 规定 $h(s, g) = (s, hg)$，
则 $\mathcal{L}_1$ 成为 $\mathcal{O}$-模预层。
其层化 $\mathcal{L} = \mathcal{L}_1^\#$ 是可逆
$\mathcal{O}$-模，且相关的 $\mathcal{O}^*$-挠子
$\mathcal{L}^*$ 与 $\mathcal{F}$ 同构；验证从略。
\end{proof}









\section{上同调的局部性}
\label{sites-cohomology-section-locality}

\noindent
下面的引理说明，定义层 $\mathcal{F}$ 在站点的一个对象上的上同调时没有歧义。

\begin{lemma}
\label{sites-cohomology-lemma-cohomology-of-open}
设 $(\mathcal{C}, \mathcal{O})$ 为环化站点，
$U$ 为 $\mathcal{C}$ 的对象。
\begin{enumerate}
\item 若 $\mathcal{I}$ 是内射 $\mathcal{O}$-模，
则 $\mathcal{I}|_U$ 是内射 $\mathcal{O}_U$-模。
\item 对任意 $\mathcal{O}$-模层 $\mathcal{F}$，有
$H^p(U, \mathcal{F}) = H^p(\mathcal{C}/U, \mathcal{F}|_U)$.
\end{enumerate}
\end{lemma}

\begin{proof}
回忆，限制到 $U$ 的函子 $j_U^{-1}$ 是零延拓函子 $j_{U!}$ 的右伴随，
见《站点上的模》第 \ref{sites-modules-section-localize} 节。
此外，$j_{U!}$ 正合。因此 (1) 由《同调代数》引理
\ref{homology-lemma-adjoint-preserve-injectives} 得到。

\medskip\noindent
按照定义，$H^p(U, \mathcal{F}) = H^p(\mathcal{I}^\bullet(U))$，
其中 $\mathcal{F} \to \mathcal{I}^\bullet$ 是
$\textit{Mod}(\mathcal{O})$ 中的内射分解。
由上可知，$\mathcal{F}|_U \to \mathcal{I}^\bullet|_U$
是 $\textit{Mod}(\mathcal{O}_U)$ 中的内射分解。
故 $H^p(U, \mathcal{F}|_U)$ 等于
$H^p(\mathcal{I}^\bullet|_U(U))$。
当然，对 $\mathcal{C}$ 上的任意层 $\mathcal{F}$，都有
$\mathcal{F}(U) = \mathcal{F}|_U(U)$。这就给出 (2) 中的等式。
\end{proof}

\noindent
下面的引理将用于考察改变一个部分宇宙时会发生什么，
也用于比较小 étale 站点与大 étale 站点的上同调。

\begin{lemma}
\label{sites-cohomology-lemma-cohomology-bigger-site}
设 $\mathcal{C}$ 和 $\mathcal{D}$ 为站点，
$u : \mathcal{C} \to \mathcal{D}$ 为函子。
假设 $u$ 满足《站点》引理 \ref{sites-lemma-bigger-site} 的假设。
设 $g : \Sh(\mathcal{C}) \to \Sh(\mathcal{D})$
为相关的拓扑斯态射。对 $\mathcal{D}$ 上的任意阿贝尔层
$\mathcal{F}$，有同构
$$
R\Gamma(\mathcal{C}, g^{-1}\mathcal{F}) = R\Gamma(\mathcal{D}, \mathcal{F}),
$$
特别地，
$H^p(\mathcal{C}, g^{-1}\mathcal{F}) = H^p(\mathcal{D}, \mathcal{F})$
；并且对任意 $U \in \Ob(\mathcal{C})$，有同构
$$
R\Gamma(U, g^{-1}\mathcal{F}) = R\Gamma(u(U), \mathcal{F}),
$$
特别地，
$H^p(U, g^{-1}\mathcal{F}) = H^p(u(U), \mathcal{F})$。
所有这些同构关于 $\mathcal{F}$ 都是函子性的。
\end{lemma}

\begin{proof}
由假设 (e)，显然
$\Gamma(\mathcal{C}, g^{-1}\mathcal{F}) = \Gamma(\mathcal{D}, \mathcal{F})$
，所以只需证明 $g^{-1}$ 把内射阿贝尔层变为内射阿贝尔层。
照例应用《同调代数》引理
\ref{homology-lemma-adjoint-preserve-injectives}。
采用《站点》引理 \ref{sites-lemma-bigger-site} 的记号，
$g^{-1}$ 的左伴随是 $g_! = f^{-1}$，而它是正合函子。
因此该引理确实适用。
\end{proof}

\noindent
设 $(\mathcal{C}, \mathcal{O})$ 为环化站点，
$\mathcal{F}$ 为 $\mathcal{O}$-模层。
设 $\varphi : U \to V$ 为 $\mathcal{O}$ 的态射。
则有典范的{\it 限制映射}
\begin{equation}
\label{sites-cohomology-equation-restriction-mapping}
H^n(V, \mathcal{F})
\longrightarrow
H^n(U, \mathcal{F}), \quad
\xi \longmapsto \xi|_U
\end{equation}
，它关于 $\mathcal{F}$ 是函子性的。事实上，任取内射分解
$\mathcal{F} \to \mathcal{I}^\bullet$。层 $\mathcal{I}^p$ 的限制映射
给出复形态射
$$
\Gamma(V, \mathcal{I}^\bullet)
\longrightarrow
\Gamma(U, \mathcal{I}^\bullet)
$$
左端是表示 $R\Gamma(V, \mathcal{F})$ 的复形，右端是表示
$R\Gamma(U, \mathcal{F})$ 的复形。应用函子 $H^n$，得到上同调群上的映射。
如上所示，我们用记号 $\xi \mapsto \xi|_U$ 表示这个映射。
因此规则 $U \mapsto H^n(U, \mathcal{F})$ 是 $\mathcal{O}$-模预层。
这个预层通常记作 $\underline{H}^n(\mathcal{F})$。
引理 \ref{sites-cohomology-lemma-include} 将给出它的另一种解释。

\medskip\noindent
下面的引理说明，转到一个覆盖后可以消去高阶上同调类。

\begin{lemma}
\label{sites-cohomology-lemma-kill-cohomology-class-on-covering}
设 $(\mathcal{C}, \mathcal{O})$ 为环化站点，
$\mathcal{F}$ 为 $\mathcal{O}$-模层，$U$ 为 $\mathcal{C}$ 的对象。
设 $n > 0$ 且 $\xi \in H^n(U, \mathcal{F})$。
则存在 $\mathcal{C}$ 的覆盖 $\{U_i \to U\}$，
使得对所有 $i \in I$ 都有 $\xi|_{U_i} = 0$。
\end{lemma}

\begin{proof}
设 $\mathcal{F} \to \mathcal{I}^\bullet$ 为内射分解。则
$$
H^n(U, \mathcal{F}) =
\frac{\Ker(\mathcal{I}^n(U) \to \mathcal{I}^{n + 1}(U))}
{\Im(\mathcal{I}^{n - 1}(U) \to \mathcal{I}^n(U))}.
$$
选取表示上述商中上同调类的元素
$\tilde \xi \in \mathcal{I}^n(U)$。由于 $\mathcal{I}^\bullet$
是 $\mathcal{F}$ 的内射分解且 $n > 0$，复形
$\mathcal{I}^\bullet$ 在次数 $n$ 正合。因此
$\Im(\mathcal{I}^{n - 1} \to \mathcal{I}^n) =
\Ker(\mathcal{I}^n \to \mathcal{I}^{n + 1})$ 作为层成立。
由于 $\tilde \xi$ 是核层在 $U$ 上的截面，
故存在站点的覆盖 $\{U_i \to U\}$，使得
$\tilde \xi|_{U_i}$ 是某个截面
$\xi_i \in \mathcal{I}^{n - 1}(U_i)$ 在 $d$ 下的像。
按照限制 $\xi|_{U_i}$ 的定义，它对应于
$\tilde \xi|_{U_i}$ 的类，结论随即得到。
\end{proof}

\begin{lemma}
\label{sites-cohomology-lemma-higher-direct-images}
设 $f : (\mathcal{C}, \mathcal{O}_\mathcal{C}) \to
(\mathcal{D}, \mathcal{O}_\mathcal{D})$ 为环化站点的态射，
对应于连续函子 $u : \mathcal{D} \to \mathcal{C}$。
对任意 $\mathcal{F} \in \Ob(\textit{Mod}(\mathcal{O}_\mathcal{C}))$，
层 $R^if_*\mathcal{F}$ 是预层
$$
V \longmapsto H^i(u(V), \mathcal{F})
$$
的相关层。
\end{lemma}

\begin{proof}
设 $\mathcal{F} \to \mathcal{I}^\bullet$ 为内射分解。
按照定义，$R^if_*\mathcal{F}$ 是复形
$$
f_*\mathcal{I}^0 \to f_*\mathcal{I}^1 \to f_*\mathcal{I}^2 \to \ldots
$$
的第 $i$ 个上同调层。按照 $\mathcal{O}_\mathcal{D}$-模上的
阿贝尔范畴结构之定义，这个上同调层是预层
$$
V
\longmapsto
\frac{\Ker(f_*\mathcal{I}^i(V) \to f_*\mathcal{I}^{i + 1}(V))}
{\Im(f_*\mathcal{I}^{i - 1}(V) \to f_*\mathcal{I}^i(V))}
$$
的相关层，而它显然等于
$$
\frac{\Ker(\mathcal{I}^i(u(V)) \to \mathcal{I}^{i + 1}(u(V)))}
{\Im(\mathcal{I}^{i - 1}(u(V)) \to \mathcal{I}^i(u(V)))}
$$
，后者等于 $H^i(u(V), \mathcal{F})$，证毕。
\end{proof}






\section{{\v C}ech 复形与 {\v C}ech 上同调}
\label{sites-cohomology-section-cech}

\noindent
设 $\mathcal{C}$ 为范畴，$\mathcal{U} = \{U_i \to U\}_{i \in I}$
为具有固定靶的态射族，见《站点》定义
\ref{sites-definition-family-morphisms-fixed-target}。
假设所有纤维积 $U_{i_0} \times_U \ldots \times_U U_{i_p}$
都存在于 $\mathcal{C}$ 中。设 $\mathcal{F}$ 为 $\mathcal{C}$ 上的阿贝尔预层。令
$$
\check{\mathcal{C}}^p(\mathcal{U}, \mathcal{F})
=
\prod\nolimits_{(i_0, \ldots, i_p) \in I^{p + 1}}
\mathcal{F}(U_{i_0} \times_U \ldots \times_U U_{i_p}).
$$
这是阿贝尔群。对
$s \in \check{\mathcal{C}}^p(\mathcal{U}, \mathcal{F})$，
以 $s_{i_0\ldots i_p}$ 表示它在因子
$\mathcal{F}(U_{i_0} \times_U \ldots \times_U U_{i_p})$ 中的值。定义
$$
d : \check{\mathcal{C}}^p(\mathcal{U}, \mathcal{F})
\longrightarrow
\check{\mathcal{C}}^{p + 1}(\mathcal{U}, \mathcal{F})
$$
如下：
\begin{equation}
\label{sites-cohomology-equation-d-cech}
d(s)_{i_0\ldots i_{p + 1}} =
\sum\nolimits_{j = 0}^{p + 1}
(-1)^j s_{i_0\ldots \hat i_j \ldots i_{p + 1}}
|_{U_{i_0} \times_U \ldots \times_U U_{i_{p + 1}}}
\end{equation}
其中限制取自投影映射
$$
U_{i_0} \times_U \ldots \times_U U_{i_{p + 1}} \longrightarrow
U_{i_0} \times_U \ldots \times_U \widehat{U_{i_j}} \times_U
\ldots \times_U U_{i_{p + 1}}.
$$
不难看出 $d \circ d = 0$。换言之，
$\check{\mathcal{C}}^\bullet(\mathcal{U}, \mathcal{F})$ 是复形。

\begin{definition}
\label{sites-cohomology-definition-cech-complex}
设 $\mathcal{C}$ 为范畴，$\mathcal{U} = \{U_i \to U\}_{i \in I}$
为具有固定靶的态射族，使得所有纤维积
$U_{i_0} \times_U \ldots \times_U U_{i_p}$ 都存在于 $\mathcal{C}$ 中。
设 $\mathcal{F}$ 为 $\mathcal{C}$ 上的阿贝尔预层。
复形 $\check{\mathcal{C}}^\bullet(\mathcal{U}, \mathcal{F})$
称为与 $\mathcal{F}$ 和族 $\mathcal{U}$ 相关的{\it {\v C}ech 复形}。
它的上同调群
$H^i(\check{\mathcal{C}}^\bullet(\mathcal{U}, \mathcal{F}))$
称为 $\mathcal{F}$ 关于 $\mathcal{U}$ 的{\it {\v C}ech 上同调群}，
记作 $\check H^i(\mathcal{U}, \mathcal{F})$。
\end{definition}

\noindent
注意，站点 $\mathcal{C}$ 的任意覆盖 $\{U_i \to U\}$
都是这个定义适用的具有固定靶的态射族。

\begin{lemma}
\label{sites-cohomology-lemma-cech-h0}
设 $\mathcal{C}$ 为站点，$\mathcal{F}$ 为 $\mathcal{C}$ 上的阿贝尔预层。
下列条件等价：
\begin{enumerate}
\item $\mathcal{F}$ 是 $\mathcal{C}$ 上的阿贝尔层；
\item 对站点 $\mathcal{C}$ 的每个覆盖
$\mathcal{U} = \{U_i \to U\}_{i \in I}$，自然映射
$$
\mathcal{F}(U) \to \check{H}^0(\mathcal{U}, \mathcal{F})
$$
（见《站点》第 \ref{sites-section-sheafification} 节）是双射。
\end{enumerate}
\end{lemma}

\begin{proof}
这是因为层条件恰好是：对 $\mathcal{C}$ 的每个覆盖，
$\mathcal{F}(U) \to \check{H}^0(\mathcal{U}, \mathcal{F})$
都是双射。
\end{proof}

\noindent
设 $\mathcal{C}$ 为范畴，$\mathcal{U} = \{U_i \to U\}_{i\in I}$
为 $\mathcal{C}$ 中具有固定靶的态射族，并假设所有纤维积
$U_{i_0} \times_U \ldots \times_U U_{i_p}$ 都存在于 $\mathcal{C}$ 中。
设 $\mathcal{V} = \{V_j \to V\}_{j\in J}$ 为另一个这样的族。
设 $f : U \to V$、$\alpha : I \to J$ 和
$f_i : U_i \to V_{\alpha(i)}$ 构成具有固定靶的态射族之间的态射，
见《站点》第 \ref{sites-section-refinements} 节。
在这种情形下，得到 {\v C}ech 复形的映射
\begin{equation}
\label{sites-cohomology-equation-map-cech-complexes}
\varphi : \check{\mathcal{C}}^\bullet(\mathcal{V}, \mathcal{F})
\longrightarrow
\check{\mathcal{C}}^\bullet(\mathcal{U}, \mathcal{F})
\end{equation}
，它在次数 $p$ 由下式给出：
$$
\varphi(s)_{i_0 \ldots i_p} =
(f_{i_0} \times \ldots \times f_{i_p})^*s_{\alpha(i_0) \ldots \alpha(i_p)}
$$


\section{{\v C}ech 上同调作为预层上的函子}
\label{sites-cohomology-section-cech-functor}

\noindent
警告：本节只讨论范畴上的阿贝尔预层。
在层和层范畴的语境中，以下结果完全不成立！

\medskip\noindent
设 $\mathcal{C}$ 为范畴，$\mathcal{U} = \{U_i \to U\}_{i \in I}$
为具有固定靶的态射族，使得所有纤维积
$U_{i_0} \times_U \ldots \times_U U_{i_p}$ 都存在于 $\mathcal{C}$ 中。
设 $\mathcal{F}$ 为 $\mathcal{C}$ 上的阿贝尔预层。构造
$$
\mathcal{F} \longmapsto \check{\mathcal{C}}^\bullet(\mathcal{U}, \mathcal{F})
$$
关于 $\mathcal{F}$ 是函子性的。事实上，它是函子
\begin{equation}
\label{sites-cohomology-equation-cech-functor}
\check{\mathcal{C}}^\bullet(\mathcal{U}, -) :
\textit{PAb}(\mathcal{C})
\longrightarrow
\text{Comp}^{+}(\textit{Ab})
\end{equation}
。记号见《导出范畴》定义 \ref{derived-definition-complexes-notation}。
回忆，阿贝尔范畴中下有界复形所成的范畴仍是阿贝尔范畴，
见《同调代数》引理 \ref{homology-lemma-cat-cochain-abelian}。

\begin{lemma}
\label{sites-cohomology-lemma-cech-exact-presheaves}
方程 (\ref{sites-cohomology-equation-cech-functor}) 给出的函子是正合函子
（见《同调代数》引理 \ref{homology-lemma-exact-functor}）。
\end{lemma}

\begin{proof}
对 $\mathcal{C}$ 的任意对象 $W$，函子
$\mathcal{F} \mapsto \mathcal{F}(W)$ 是从
$\textit{PAb}(\mathcal{C})$ 到 $\textit{Ab}$ 的加性正合函子。
复形的各项 $\check{\mathcal{C}}^p(\mathcal{U}, \mathcal{F})$
是这些正合函子的乘积，因而也是正合的。
此外，复形列正合当且仅当它在每个给定次数上的项所成之列正合。
引理得证。
\end{proof}

\begin{lemma}
\label{sites-cohomology-lemma-cech-cohomology-delta-functor-presheaves}
设 $\mathcal{C}$ 为范畴，
$\mathcal{U} = \{U_i \to U\}_{i \in I}$ 为具有固定靶的态射族，
使得所有纤维积 $U_{i_0} \times_U \ldots \times_U U_{i_p}$
都存在于 $\mathcal{C}$ 中。函子
$\mathcal{F} \mapsto \check{H}^n(\mathcal{U}, \mathcal{F})$
构成从阿贝尔范畴 $\textit{PAb}(\mathcal{C})$ 到
$\mathbf{Z}$-模范畴的 $\delta$-函子
（见《同调代数》定义 \ref{homology-definition-cohomological-delta-functor}）。
\end{lemma}

\begin{proof}
由引理 \ref{sites-cohomology-lemma-cech-exact-presheaves}，阿贝尔预层的短正合列
$0 \to \mathcal{F}_1 \to \mathcal{F}_2 \to \mathcal{F}_3 \to 0$
被变为 $\mathbf{Z}$-模复形的短正合列。因此可应用《同调代数》引理
\ref{homology-lemma-long-exact-sequence-cochain}，得到边界映射
$\delta_{\mathcal{F}_1 \to \mathcal{F}_2 \to \mathcal{F}_3} :
\check{H}^n(\mathcal{U}, \mathcal{F}_3) \to
\check{H}^{n + 1}(\mathcal{U}, \mathcal{F}_1)$
以及相应的长正合列。这些映射与预层短正合列之间的映射相容，
其验证从略。
\end{proof}

\begin{lemma}
\label{sites-cohomology-lemma-cech-map-into}
设 $\mathcal{C}$ 为范畴，$\mathcal{U} = \{U_i \to U\}_{i \in I}$
为具有固定靶的态射族，使得所有纤维积
$U_{i_0} \times_U \ldots \times_U U_{i_p}$ 都存在于 $\mathcal{C}$ 中。
考虑阿贝尔预层的链复形 $\mathbf{Z}_{\mathcal{U}, \bullet}$：
$$
\ldots
\to
\bigoplus_{i_0i_1i_2} \mathbf{Z}_{U_{i_0} \times_U U_{i_1} \times_U U_{i_2}}
\to
\bigoplus_{i_0i_1} \mathbf{Z}_{U_{i_0} \times_U U_{i_1}}
\to
\bigoplus_{i_0} \mathbf{Z}_{U_{i_0}}
\to 0 \to \ldots
$$
其中最后一个非零项置于次数 $0$，且映射
$$
\mathbf{Z}_{U_{i_0} \times_U \ldots \times_u U_{i_{p + 1}}}
\longrightarrow
\mathbf{Z}_{U_{i_0} \times_U
\ldots \widehat{U_{i_j}} \ldots \times_U U_{i_{p + 1}}}
$$
由典范映射的 $(-1)^j$ 倍给出。则存在同构
$$
\Hom_{\textit{PAb}(\mathcal{C})}(\mathbf{Z}_{\mathcal{U}, \bullet}, \mathcal{F})
=
\check{\mathcal{C}}^\bullet(\mathcal{U}, \mathcal{F})
$$
，它关于 $\mathcal{F} \in \Ob(\textit{PAb}(\mathcal{C}))$ 是函子性的。
\end{lemma}

\begin{proof}
这直接来自下列事实：
\begin{align*}
\Hom_{\textit{PAb}(\mathcal{C})}(
\bigoplus_{i_0 \ldots i_p}
\mathbf{Z}_{U_{i_0} \times_U \ldots \times_U U_{i_p}},
\mathcal{F})
& =
\prod_{i_0 \ldots i_p}
\Hom_{\textit{PAb}(\mathcal{C})}(
\mathbf{Z}_{U_{i_0} \times_U \ldots \times_U U_{i_p}},
\mathcal{F}) \\
& =
\prod_{i_0 \ldots i_p}
\mathcal{F}(U_{i_0} \times_U \ldots \times_U U_{i_p})
\end{align*}
见《站点上的模》引理 \ref{sites-modules-lemma-obvious-adjointness}。
\end{proof}

\begin{lemma}
\label{sites-cohomology-lemma-homology-complex}
设 $\mathcal{C}$ 为范畴，
$\mathcal{U} = \{f_i : U_i \to U\}_{i \in I}$ 为具有固定靶的态射族，
使得所有纤维积 $U_{i_0} \times_U \ldots \times_U U_{i_p}$
都存在于 $\mathcal{C}$ 中。上面引理 \ref{sites-cohomology-lemma-cech-map-into}
中的预层链复形 $\mathbf{Z}_{\mathcal{U}, \bullet}$ 在正次数正合；
也就是说，对 $i > 0$，同调预层
$H_i(\mathbf{Z}_{\mathcal{U}, \bullet})$ 为零。
\end{lemma}

\begin{proof}
设 $V$ 为 $\mathcal{C}$ 的对象。我们需要证明阿贝尔群链复形
$\mathbf{Z}_{\mathcal{U}, \bullet}(V)$ 在次数 $> 0$ 正合。
这个复形是
$$
\xymatrix{
\ldots \ar[d] \\
\bigoplus_{i_0i_1i_2}
\mathbf{Z}[
\Mor_\mathcal{C}(V, U_{i_0} \times_U U_{i_1} \times_U U_{i_2})
]
\ar[d] \\
\bigoplus_{i_0i_1}
\mathbf{Z}[
\Mor_\mathcal{C}(V, U_{i_0} \times_U U_{i_1})
]
\ar[d] \\
\bigoplus_{i_0}
\mathbf{Z}[
\Mor_\mathcal{C}(V, U_{i_0})
] \ar[d] \\
0
}
$$
对任意态射 $\varphi : V \to U$，记
$\Mor_\varphi(V, U_i) = \{\varphi_i : V \to U_i \mid
f_i \circ \varphi_i = \varphi\}$。对
$\Mor_\varphi(V, U_{i_0} \times_U \ldots \times_U U_{i_p})$
采用类似记号。注意，与纤维积
$U_{i_0} \times_U \ldots \times_U U_{i_p}$ 之间的各个投影映射复合，
保持这些态射集合。因此上面的复形等于复形
$$
\xymatrix{
\ldots \ar[d] \\
\bigoplus_\varphi
\bigoplus_{i_0i_1i_2}
\mathbf{Z}[
\Mor_\varphi(V, U_{i_0} \times_U U_{i_1} \times_U U_{i_2})
]
\ar[d] \\
\bigoplus_\varphi
\bigoplus_{i_0i_1}
\mathbf{Z}[
\Mor_\varphi(V, U_{i_0} \times_U U_{i_1})
]
\ar[d] \\
\bigoplus_\varphi
\bigoplus_{i_0}
\mathbf{Z}[
\Mor_\varphi(V, U_{i_0})
] \ar[d] \\
0
}
$$
其次，注意有
$$
\Mor_\varphi(V, U_{i_0} \times_U \ldots \times_U U_{i_p})
=
\Mor_\varphi(V, U_{i_0}) \times \ldots
\times \Mor_\varphi(V, U_{i_p})
$$
利用这一点以及 $\mathbf{Z}[A] \oplus \mathbf{Z}[B] =
\mathbf{Z}[A \amalg B]$，可见该复形变为
$$
\xymatrix{
\ldots \ar[d] \\
\bigoplus_\varphi
\mathbf{Z}\left[
\coprod_{i_0i_1i_2}
\Mor_\varphi(V, U_{i_0}) \times \Mor_\varphi(V, U_{i_1}) \times
\Mor_\varphi(V, U_{i_2})
\right]
\ar[d] \\
\bigoplus_\varphi
\mathbf{Z}\left[
\coprod_{i_0i_1}
\Mor_\varphi(V, U_{i_0}) \times \Mor_\varphi(V, U_{i_1})
\right]
\ar[d] \\
\bigoplus_\varphi
\mathbf{Z}\left[
\coprod_{i_0}
\Mor_\varphi(V, U_{i_0})
\right] \ar[d] \\
0
}
$$
最后，令 $S_\varphi = \coprod_{i \in I} \Mor_\varphi(V, U_i)$，
得到
$$
\bigoplus\nolimits_\varphi \left(\ldots \to
\mathbf{Z}[S_\varphi \times S_\varphi \times S_\varphi] \to
\mathbf{Z}[S_\varphi \times S_\varphi] \to
\mathbf{Z}[S_\varphi] \to 0 \to \ldots
\right)
$$
这样，任务得到了简化。只需证明，对任意非空集合 $S$，
自由阿贝尔群的（增广）复形
$$
\ldots \to
\mathbf{Z}[S \times S \times S] \to
\mathbf{Z}[S \times S] \to
\mathbf{Z}[S] \xrightarrow{\Sigma} \mathbf{Z} \to 0 \to \ldots
$$
在所有次数都正合。为此固定元素 $s \in S$，并采用同伦
$$
n_{(s_0, \ldots, s_p)} \longmapsto n_{(s, s_0, \ldots, s_p)}
$$
，其中记号的含义显然。
\end{proof}

\begin{lemma}
\label{sites-cohomology-lemma-complex-tensored-still-exact}
\begin{slogan}
整数预层的 {\v C}ech 复形是整数常值预层的平坦分解。
\end{slogan}
设 $\mathcal{C}$ 为范畴，
$\mathcal{U} = \{f_i : U_i \to U\}_{i \in I}$ 为具有固定靶的态射族，
使得所有纤维积 $U_{i_0} \times_U \ldots \times_U U_{i_p}$
都存在于 $\mathcal{C}$ 中。设 $\mathcal{O}$ 为 $\mathcal{C}$ 上的环预层。
链复形
$$
\mathbf{Z}_{\mathcal{U}, \bullet}
\otimes_{p, \mathbf{Z}}
\mathcal{O}
$$
在正次数正合。这里 $\mathbf{Z}_{\mathcal{U}, \bullet}$
是引理 \ref{sites-cohomology-lemma-cech-map-into} 的链复形，张量积取在
取值为 $\mathbf{Z}$ 的常值环预层上。
\end{lemma}

\begin{proof}
设 $V$ 为 $\mathcal{C}$ 的对象。在引理
\ref{sites-cohomology-lemma-homology-complex} 的证明中，我们已经看到，
作为复形，$\mathbf{Z}_{\mathcal{U}, \bullet}(V)$ 同构于若干复形的直和，
而这些复形都同伦等价于置于次数零的 $\mathbf{Z}$。因此，
$\mathbf{Z}_{\mathcal{U}, \bullet}(V) \otimes_\mathbf{Z} \mathcal{O}(V)$
作为复形也同构于若干复形的直和，而这些复形都同伦等价于
置于次数零的 $\mathcal{O}(V)$。
也可以应用《站点上的模》引理
\ref{sites-modules-lemma-flat-resolution-of-flat}：
它适用是因为各预层 $\mathbf{Z}_{\mathcal{U}, i}$ 都平坦，
且引理 \ref{sites-cohomology-lemma-homology-complex} 的证明还说明
$H_0(\mathbf{Z}_{\mathcal{U}, \bullet})$ 是平坦预层。
\end{proof}

\begin{lemma}
\label{sites-cohomology-lemma-cech-cohomology-derived-presheaves}
设 $\mathcal{C}$ 为范畴，
$\mathcal{U} = \{f_i : U_i \to U\}_{i \in I}$ 为具有固定靶的态射族，
使得所有纤维积 $U_{i_0} \times_U \ldots \times_U U_{i_p}$
都存在于 $\mathcal{C}$ 中。作为 $\delta$-函子，{\v C}ech 上同调函子
$\check{H}^p(\mathcal{U}, -)$ 与函子
$$
\check{H}^0(\mathcal{U}, -) :
\textit{PAb}(\mathcal{C})
\longrightarrow
\textit{Ab}.
$$
的右导出函子典范同构。
此外，存在函子性的拟同构
$$
\check{\mathcal{C}}^\bullet(\mathcal{U}, \mathcal{F})
\longrightarrow
R\check{H}^0(\mathcal{U}, \mathcal{F})
$$
，其中右端表示左正合函子 $\check{H}^0(\mathcal{U}, -)$ 的导出函子
$$
R\check{H}^0(\mathcal{U}, -) :
D^{+}(\textit{PAb}(\mathcal{C}))
\longrightarrow
D^{+}(\mathbf{Z})
$$
\end{lemma}

\begin{proof}
注意，阿贝尔预层范畴有足够多的内射对象，见《内射对象》命题
\ref{injectives-proposition-presheaves-injectives}。
还要注意，$\check{H}^0(\mathcal{U}, -)$ 是从阿贝尔预层范畴到
$\mathbf{Z}$-模范畴的左正合函子。因此导出函子和右导出函子都存在，
见《导出范畴》第 \ref{derived-section-right-derived-functor} 节。

\medskip\noindent
设 $\mathcal{I}$ 为内射阿贝尔预层。在这种情形下，函子
$\Hom_{\textit{PAb}(\mathcal{C})}(-, \mathcal{I})$
在 $\textit{PAb}(\mathcal{C})$ 上正合。由引理
\ref{sites-cohomology-lemma-cech-map-into}，有
$$
\Hom_{\textit{PAb}(\mathcal{C})}(
\mathbf{Z}_{\mathcal{U}, \bullet}, \mathcal{I})
=
\check{\mathcal{C}}^\bullet(\mathcal{U}, \mathcal{I}).
$$
由引理 \ref{sites-cohomology-lemma-homology-complex}，
$\mathbf{Z}_{\mathcal{U}, \bullet}$ 在正次数正合。
再由上述向 $\mathcal{I}$ 取 Hom 的正合性可知，对所有 $i > 0$，
$\check{H}^i(\mathcal{U}, \mathcal{I}) = 0$。
因此 $\delta$-函子 $(\check{H}^n, \delta)$
（见引理 \ref{sites-cohomology-lemma-cech-cohomology-delta-functor-presheaves}）
满足《同调代数》引理 \ref{homology-lemma-efface-implies-universal}
的假设，从而是万有 $\delta$-函子。

\medskip\noindent
由《导出范畴》引理 \ref{derived-lemma-higher-derived-functors}，
序列 $R^i\check{H}^0(\mathcal{U}, -)$ 也构成万有 $\delta$-函子。
根据万有 $\delta$-函子的唯一性（见《同调代数》引理
\ref{homology-lemma-uniqueness-universal-delta-functor}），得到
$R^i\check{H}^0(\mathcal{U}, -) = \check{H}^i(\mathcal{U}, -)$.
这对大多数应用已经足够；读者可以跳过证明的剩余部分。

\medskip\noindent
设 $\mathcal{F}$ 为 $\mathcal{C}$ 上的任意阿贝尔预层。
在范畴 $\textit{PAb}(\mathcal{C})$ 中选取内射分解
$\mathcal{F} \to \mathcal{I}^\bullet$。考虑双复形
$\check{\mathcal{C}}^\bullet(\mathcal{U}, \mathcal{I}^\bullet)$
，其项为 $\check{\mathcal{C}}^p(\mathcal{U}, \mathcal{I}^q)$。
再考虑与这个双复形相关的全复形
$\text{Tot}(\check{\mathcal{C}}^\bullet(\mathcal{U}, \mathcal{I}^\bullet))$
，见《同调代数》第 \ref{homology-section-double-complexes} 节。
有复形态射
$$
\check{\mathcal{C}}^\bullet(\mathcal{U}, \mathcal{F})
\longrightarrow
\text{Tot}(\check{\mathcal{C}}^\bullet(\mathcal{U}, \mathcal{I}^\bullet))
$$
，它来自映射
$\check{\mathcal{C}}^p(\mathcal{U}, \mathcal{F})
\to \check{\mathcal{C}}^p(\mathcal{U}, \mathcal{I}^0)$
；另有复形态射
$$
\check{H}^0(\mathcal{U}, \mathcal{I}^\bullet)
\longrightarrow
\text{Tot}(\check{\mathcal{C}}^\bullet(\mathcal{U}, \mathcal{I}^\bullet))
$$
，它来自映射
$\check{H}^0(\mathcal{U}, \mathcal{I}^q) \to
\check{\mathcal{C}}^0(\mathcal{U}, \mathcal{I}^q)$.
应用《同调代数》引理
\ref{homology-lemma-double-complex-gives-resolution}，
这两个映射都是拟同构。事实上，由于作为函子的 {\v C}ech 复形正合
（引理 \ref{sites-cohomology-lemma-cech-exact-presheaves}），双复形的各列在正次数正合；
又由于我们刚刚看到内射预层 $\mathcal{I}^q$ 的高阶
{\v C}ech 上同调群为零，双复形的各行也在正次数正合。
拟同构在 $D^{+}(\mathbf{Z})$ 中可逆，因此得到引理中最后显示的态射。
该态射的函子性验证从略。
\end{proof}





\section{{\v C}ech 上同调与上同调}
\label{sites-cohomology-section-cech-cohomology-cohomology}

\noindent
上同调与 {\v C}ech 上同调之间的关系来自这样一个事实：
内射阿贝尔层的高阶 {\v C}ech 上同调为零。
为说明这一点，先注意内射阿贝尔层也是内射阿贝尔预层，
然后应用前一节关于 {\v C}ech 上同调的结果。

\begin{lemma}
\label{sites-cohomology-lemma-injective-abelian-sheaf-injective-presheaf}
设 $\mathcal{C}$ 为站点。内射阿贝尔层作为范畴
$\textit{PAb}(\mathcal{C})$ 中的对象也是内射的。
\end{lemma}

\begin{proof}
对范畴 $\mathcal{A} = \textit{Ab}(\mathcal{C})$、
$\mathcal{B} = \textit{PAb}(\mathcal{C})$，以及包含函子和层化函子，
应用《同调代数》引理 \ref{homology-lemma-adjoint-preserve-injectives}。
该引理的所有假设均满足，见《站点上的模》第
\ref{sites-modules-section-abelian-sheaves} 节。
\end{proof}

\begin{lemma}
\label{sites-cohomology-lemma-injective-trivial-cech}
设 $\mathcal{C}$ 为站点，
$\mathcal{U} = \{U_i \to U\}_{i \in I}$ 为 $\mathcal{C}$ 的覆盖。
设 $\mathcal{I}$ 为内射阿贝尔层，即
$\textit{Ab}(\mathcal{C})$ 中的内射对象。则
$$
\check{H}^p(\mathcal{U}, \mathcal{I}) =
\left\{
\begin{matrix}
\mathcal{I}(U) & \text{若} & p = 0 \\
0 & \text{若} & p > 0
\end{matrix}
\right.
$$
\end{lemma}

\begin{proof}
由引理 \ref{sites-cohomology-lemma-injective-abelian-sheaf-injective-presheaf}，
$\mathcal{I}$ 是 $\textit{PAb}(\mathcal{C})$ 中的内射对象。
因此可应用引理 \ref{sites-cohomology-lemma-cech-cohomology-derived-presheaves}
（或其证明），得到高阶 {\v C}ech 上同调群的消失。
次数零的情形见引理 \ref{sites-cohomology-lemma-cech-h0}。
\end{proof}

\begin{lemma}
\label{sites-cohomology-lemma-cech-cohomology}
设 $\mathcal{C}$ 为站点，
$\mathcal{U} = \{U_i \to U\}_{i \in I}$ 为 $\mathcal{C}$ 的覆盖。
存在函子
$$
\check{\mathcal{C}}^\bullet(\mathcal{U}, -)
\longrightarrow
R\Gamma(U, -)
$$
之间的变换，其中函子的定义域和值域为
$\textit{Ab}(\mathcal{C}) \to D^{+}(\mathbf{Z})$。
特别地，当 $\mathcal{F}$ 取遍 $\textit{Ab}(\mathcal{C})$ 时，
它给出函子变换
$\check{H}^p(\mathcal{U}, \mathcal{F}) \to H^p(U, \mathcal{F})$。
\end{lemma}

\begin{proof}
设 $\mathcal{F}$ 为阿贝尔层。选取内射分解
$\mathcal{F} \to \mathcal{I}^\bullet$。考虑双复形
$\check{\mathcal{C}}^\bullet(\mathcal{U}, \mathcal{I}^\bullet)$
，其项为 $\check{\mathcal{C}}^p(\mathcal{U}, \mathcal{I}^q)$。
再考虑相关的全复形
$\text{Tot}(\check{\mathcal{C}}^\bullet(\mathcal{U}, \mathcal{I}^\bullet))$,
见《同调代数》定义 \ref{homology-definition-associated-simple-complex}。
有复形态射
$$
\alpha :
\Gamma(U, \mathcal{I}^\bullet)
\longrightarrow
\text{Tot}(\check{\mathcal{C}}^\bullet(\mathcal{U}, \mathcal{I}^\bullet))
$$
，它来自映射
$\mathcal{I}^q(U) \to \check{H}^0(\mathcal{U}, \mathcal{I}^q)$
；另有复形态射
$$
\beta :
\check{\mathcal{C}}^\bullet(\mathcal{U}, \mathcal{F})
\longrightarrow
\text{Tot}(\check{\mathcal{C}}^\bullet(\mathcal{U}, \mathcal{I}^\bullet))
$$
，它来自映射 $\mathcal{F} \to \mathcal{I}^0$。
应用《同调代数》引理
\ref{homology-lemma-double-complex-gives-resolution}，可知
$\alpha$ 是拟同构。事实上，引理 \ref{sites-cohomology-lemma-injective-trivial-cech}
说明双复形 $\check{\mathcal{C}}^\bullet(\mathcal{U}, \mathcal{I}^\bullet)$
的第 $q$ 行是 $\Gamma(U, \mathcal{I}^q)$ 的分解。
故 $\alpha$ 在 $D^{+}(\mathbf{Z})$ 中可逆，引理中的变换是
$\beta$ 后接 $\alpha$ 的逆所得的复合。其函子性验证从略。
\end{proof}

\begin{lemma}
\label{sites-cohomology-lemma-cech-h1}
设 $\mathcal{C}$ 为站点，$\mathcal{G}$ 为 $\mathcal{C}$ 上的阿贝尔层，
$\mathcal{U} = \{U_i \to U\}_{i \in I}$ 为 $\mathcal{C}$ 的覆盖。映射
$$
\check{H}^1(\mathcal{U}, \mathcal{G})
\longrightarrow
H^1(U, \mathcal{G})
$$
为单射；并且经由引理 \ref{sites-cohomology-lemma-torsors-h1} 的双射，
它把 $\check{H}^1(\mathcal{U}, \mathcal{G})$ 等同于如下集合：
那些在每个 $U_i$ 上限制为平凡挠子的
$\mathcal{G}|_U$-挠子的同构类集合。
\end{lemma}

\begin{proof}
为此构造一个逆映射。设 $\mathcal{F}$ 为 $\mathcal{C}/U$ 上的
$\mathcal{G}|_U$-挠子，其在 $\mathcal{C}/U_i$ 上的限制平凡。
由引理 \ref{sites-cohomology-lemma-trivial-torsor}，这意味着存在截面
$s_i \in \mathcal{F}(U_i)$。在 $U_{i_0} \times_U U_{i_1}$ 上，
存在 $\mathcal{G}$ 的唯一截面 $s_{i_0i_1}$，使得
$s_{i_0i_1} \cdot s_{i_0}|_{U_{i_0} \times_U U_{i_1}} =
s_{i_1}|_{U_{i_0} \times_U U_{i_1}}$。直接计算可知，
$s_{i_0i_1}$ 是 {\v C}ech 余圈，而且它的类有良好定义
（即不依赖于截面 $s_i$ 的选择）。逆映射把 $\mathcal{F}$ 的同构类
映到余圈 $(s_{i_0i_1})$ 的上同调类。略去验证该映射确实为逆。
\end{proof}

\begin{lemma}
\label{sites-cohomology-lemma-include}
设 $\mathcal{C}$ 为站点。考虑函子
$i : \textit{Ab}(\mathcal{C}) \to \textit{PAb}(\mathcal{C})$.
它是左正合函子，其右导出函子为
$$
R^pi(\mathcal{F}) = \underline{H}^p(\mathcal{F}) :
U \longmapsto H^p(U, \mathcal{F})
$$
；见第 \ref{sites-cohomology-section-locality} 节的讨论。
\end{lemma}

\begin{proof}
显然 $i$ 左正合。选取内射分解
$\mathcal{F} \to \mathcal{I}^\bullet$。按照定义，$R^pi$ 是复形
$\mathcal{I}^\bullet$ 的第 $p$ 个上同调{\it 预层}。
换言之，$R^pi(\mathcal{F})$ 在 $\mathcal{C}$ 的对象 $U$ 上的截面为
$$
\frac{\Ker(\mathcal{I}^n(U) \to \mathcal{I}^{n + 1}(U))}
{\Im(\mathcal{I}^{n - 1}(U) \to \mathcal{I}^n(U))}.
$$
，这正是 $H^p(U, \mathcal{F})$ 的定义。
\end{proof}

\begin{lemma}
\label{sites-cohomology-lemma-cech-spectral-sequence}
设 $\mathcal{C}$ 为站点，
$\mathcal{U} = \{U_i \to U\}_{i \in I}$ 为 $\mathcal{C}$ 的覆盖。
对任意阿贝尔层 $\mathcal{F}$，存在谱序列
$(E_r, d_r)_{r \geq 0}$，满足
$$
E_2^{p, q} = \check{H}^p(\mathcal{U}, \underline{H}^q(\mathcal{F}))
$$
并收敛到 $H^{p + q}(U, \mathcal{F})$。
这个谱序列关于 $\mathcal{F}$ 是函子性的。
\end{lemma}

\begin{proof}
这是下列函子的 Grothendieck 谱序列（见《导出范畴》引理
\ref{derived-lemma-grothendieck-spectral-sequence}）：
$$
i :  \textit{Ab}(\mathcal{C}) \to \textit{PAb}(\mathcal{C})
\quad\text{和}\quad
\check{H}^0(\mathcal{U}, - ) : \textit{PAb}(\mathcal{C})
\to \textit{Ab}.
$$
由引理 \ref{sites-cohomology-lemma-cech-h0}，有
$\check{H}^0(\mathcal{U}, i(\mathcal{F})) = \mathcal{F}(U)$。
由引理 \ref{sites-cohomology-lemma-injective-trivial-cech}，$i(\mathcal{I})$
是 {\v C}ech 无环的。又由引理
\ref{sites-cohomology-lemma-cech-cohomology-derived-presheaves}，作为
$\textit{PAb}(\mathcal{C})$ 上的函子，有
$\check{H}^p(\mathcal{U}, -) = R^p\check{H}^0(\mathcal{U}, -)$。
合并这些事实即得引理。
\end{proof}

\begin{lemma}
\label{sites-cohomology-lemma-cech-spectral-sequence-application}
设 $\mathcal{C}$ 为站点，
$\mathcal{U} = \{U_i \to U\}_{i \in I}$ 为覆盖，
$\mathcal{F} \in \Ob(\textit{Ab}(\mathcal{C}))$。
假设对所有 $i > 0$、所有 $p \geq 0$ 以及所有
$i_0, \ldots, i_p \in I$，都有
$H^i(U_{i_0} \times_U \ldots \times_U U_{i_p}, \mathcal{F}) = 0$。
则 $\check{H}^p(\mathcal{U}, \mathcal{F}) = H^p(U, \mathcal{F})$。
\end{lemma}

\begin{proof}
应用引理 \ref{sites-cohomology-lemma-cech-spectral-sequence} 的谱序列。
假设意味着，只要 $q \not = 0$，就有 $E_2^{p, q} = 0$。
因此谱序列在 $E_2$ 退化，结论随即得到。
\end{proof}

\begin{lemma}
\label{sites-cohomology-lemma-ses-cech-h1}
设 $\mathcal{C}$ 为站点。设
$$
0 \to \mathcal{F} \to \mathcal{G} \to \mathcal{H} \to 0
$$
为 $\mathcal{C}$ 上阿贝尔层的短正合列。设 $U$ 为 $\mathcal{C}$ 的对象。
若存在 $U$ 的覆盖所成的共尾系，使得对其中每个覆盖 $\mathcal{U}$，
都有 $\check{H}^1(\mathcal{U}, \mathcal{F}) = 0$，
则映射 $\mathcal{G}(U) \to \mathcal{H}(U)$ 为满射。
\end{lemma}

\begin{proof}
取元素 $s \in \mathcal{H}(U)$。选取覆盖
$\mathcal{U} = \{U_i \to U\}_{i \in I}$，使得
(a) $\check{H}^1(\mathcal{U}, \mathcal{F}) = 0$，且
(b) $s|_{U_i}$ 是某个截面 $s_i \in \mathcal{G}(U_i)$ 的像。
显然可以找到满足 (b) 的覆盖；再由引理的假设，
可以找到同时满足 (a) 和 (b) 的覆盖。考虑截面
$$
s_{i_0i_1} =
s_{i_1}|_{U_{i_0} \times_U U_{i_1}} - s_{i_0}|_{U_{i_0} \times_U U_{i_1}}.
$$
由于 $s_i$ 提升 $s$，可知
$s_{i_0i_1} \in \mathcal{F}(U_{i_0} \times_U U_{i_1})$.
由 $\check{H}^1(\mathcal{U}, \mathcal{F})$ 的消失，
可以找到截面 $t_i \in \mathcal{F}(U_i)$，使得
$$
s_{i_0i_1} =
t_{i_1}|_{U_{i_0} \times_U U_{i_1}} - t_{i_0}|_{U_{i_0} \times_U U_{i_1}}.
$$
于是截面 $s_i - t_i$ 显然满足层条件，并粘合为
$\mathcal{G}$ 在 $U$ 上的一个截面，它映到 $s$。证毕。
\end{proof}

\begin{lemma}
\label{sites-cohomology-lemma-cech-vanish-collection}
（《层上同调》引理 \ref{cohomology-lemma-cech-vanish} 的变式。）
设 $\mathcal{C}$ 为站点，$\text{Cov}_\mathcal{C}$ 为
$\mathcal{C}$ 的覆盖所成的集合（见《站点》定义
\ref{sites-definition-site}）。设
$\mathcal{B} \subset \Ob(\mathcal{C})$，且
$\text{Cov} \subset \text{Cov}_\mathcal{C}$
为子集。设 $\mathcal{F}$ 为 $\mathcal{C}$ 上的阿贝尔层。假设
\begin{enumerate}
\item 对每个 $\mathcal{U} \in \text{Cov}$，若
$\mathcal{U} = \{U_i \to U\}_{i \in I}$，则
$U, U_i \in \mathcal{B}$，并且每个
$U_{i_0} \times_U \ldots \times_U U_{i_p} \in \mathcal{B}$.
\item 对每个 $U \in \mathcal{B}$，$\text{Cov}$ 中出现的
$U$ 的覆盖构成 $U$ 的覆盖所成的共尾系。
\item 对每个 $\mathcal{U} \in \text{Cov}$ 和所有 $p > 0$，都有
$\check{H}^p(\mathcal{U}, \mathcal{F}) = 0$。
\end{enumerate}
则对所有 $p > 0$ 以及任意 $U \in \mathcal{B}$，都有
$H^p(U, \mathcal{F}) = 0$。
\end{lemma}

\begin{proof}
设 $\mathcal{F}$ 和 $\text{Cov}$ 如引理所述。我们把这一点表述为：
“对任意 $\mathcal{U} \in \text{Cov}$，$\mathcal{F}$ 的高阶
{\v C}ech 上同调消失”。选取 $\mathcal{F}$ 到内射阿贝尔层
$\mathcal{I}$ 中的嵌入 $\mathcal{F} \to \mathcal{I}$。
由引理 \ref{sites-cohomology-lemma-injective-trivial-cech}，对任意
$\mathcal{U} \in \text{Cov}$，$\mathcal{I}$ 的高阶
{\v C}ech 上同调消失。令 $\mathcal{Q} = \mathcal{I}/\mathcal{F}$，
于是有短正合列
$$
0 \to \mathcal{F} \to \mathcal{I} \to \mathcal{Q} \to 0.
$$
由引理 \ref{sites-cohomology-lemma-ses-cech-h1} 和假设 (2)，
这个序列给出正合列
$$
0 \to \mathcal{F}(U) \to \mathcal{I}(U) \to \mathcal{Q}(U) \to 0.
$$
，它对每个 $U \in \mathcal{B}$ 都成立。因此对任意
$\mathcal{U} \in \text{Cov}$，得到 {\v C}ech 复形的短正合列
$$
0 \to
\check{\mathcal{C}}^\bullet(\mathcal{U}, \mathcal{F}) \to
\check{\mathcal{C}}^\bullet(\mathcal{U}, \mathcal{I}) \to
\check{\mathcal{C}}^\bullet(\mathcal{U}, \mathcal{Q}) \to 0
$$
，因为由假设 (1)，{\v C}ech 复形的每一项都是
$\mathcal{B}$ 中各元素上的取值之乘积。特别地，对任意覆盖
$\mathcal{U} \in \text{Cov}$，都有 {\v C}ech 上同调群的长正合列。
这说明 $\mathcal{Q}$ 也是阿贝尔层，并且对所有
$\mathcal{U} \in \text{Cov}$，其高阶 {\v C}ech 上同调消失。

\medskip\noindent
接着考察上同调长正合列
$$
\xymatrix{
0 \ar[r] &
H^0(U, \mathcal{F}) \ar[r] &
H^0(U, \mathcal{I}) \ar[r] &
H^0(U, \mathcal{Q}) \ar[lld] \\
&
H^1(U, \mathcal{F}) \ar[r] &
H^1(U, \mathcal{I}) \ar[r] &
H^1(U, \mathcal{Q}) \ar[lld] \\
&
\ldots & \ldots & \ldots \\
}
$$
，其中 $U \in \mathcal{B}$ 任意。由于 $\mathcal{I}$ 内射，
对 $n > 0$ 有 $H^n(U, \mathcal{I}) = 0$（见《导出范畴》引理
\ref{derived-lemma-higher-derived-functors}）。由上可知，
$H^0(U, \mathcal{I}) \to H^0(U, \mathcal{Q})$ 为满射，
故 $H^1(U, \mathcal{F}) = 0$。
由于 $\mathcal{F}$ 是任意一个对所有 $\mathcal{U} \in \text{Cov}$
高阶 {\v C}ech 上同调消失的阿贝尔层，而 $\mathcal{Q}$
也是这种层，所以同样有 $H^1(U, \mathcal{Q}) = 0$。
由长正合列，这又推出 $H^2(U, \mathcal{F}) = 0$。
如此继续即可。
\end{proof}














\section{第二上同调与 gerbe}
\label{sites-cohomology-section-gerbes}

\noindent
设 $p : \mathcal{S} \to \mathcal{C}$ 是站点上的一个 gerbe，
其所有自同构群均为交换群。在这种情形下，自同构层的第一与第二上同调群
（Stacks，引理 \ref{stacks-lemma-gerbe-abelian-auts}）
控制对象的存在性。

\medskip\noindent
以下引理将由我们今后补充的、关于这一关系的更完整讨论所取代。

\begin{lemma}
\label{sites-cohomology-lemma-existence}
设 $\mathcal{C}$ 为站点。设 $p : \mathcal{S} \to \mathcal{C}$
为此站点上的一个 gerbe，且其自同构层均为阿贝尔层。
令 $\mathcal{G}$ 为 Stacks，引理
\ref{stacks-lemma-gerbe-abelian-auts} 中构造的阿贝尔群层。
设 $U$ 为 $\mathcal{C}$ 的一个对象，并满足
\begin{enumerate}
\item $U$ 在 $\mathcal{C}$ 中存在一个由覆盖 $\{U_i \to U\}$
组成的共尾系，使得 $H^1(U_i, \mathcal{G}) = 0$ 且
$H^1(U_i \times_U U_j, \mathcal{G}) = 0$
对所有 $i, j$ 成立；
\item $H^2(U, \mathcal{G}) = 0$.
\end{enumerate}
则存在一个位于 $U$ 上的 $\mathcal{S}$ 中对象。
\end{lemma}

\begin{proof}
由 Stacks，定义 \ref{stacks-definition-gerbe}，
存在覆盖 $\mathcal{U} = \{U_i \to U\}$，并且对每个 $i$，
存在位于 $U_i$ 上的 $\mathcal{S}$ 中对象 $x_i$。
记 $U_{ij} = U_i \times_U U_j$。由 (1)，加细该覆盖后，
可以假设 $H^1(U_i, \mathcal{G}) = 0$ 且
$H^1(U_{ij}, \mathcal{G}) = 0$.
考虑 $\mathcal{C}/U_{ij}$ 上的层
$$
\mathcal{F}_{ij} =
\mathit{Isom}(x_i|_{U_{ij}}, x_j|_{U_{ij}})
$$
。由于
$\mathcal{G}|_{U_{ij}} = \mathit{Aut}(x_i|_{U_{ij}})$，可见存在由预复合给出的作用
$$
\mathcal{G}|_{U_{ij}} \times \mathcal{F}_{ij} \to \mathcal{F}_{ij}
$$
。显然，$\mathcal{F}_{ij}$ 是一个伪
$\mathcal{G}|_{U_{ij}}$-挠子；事实上它是挠子，因为 gerbe 的任意两个对象
局部同构。由我们对覆盖的选取以及引理 \ref{sites-cohomology-lemma-torsors-h1}，
这些挠子均平凡（从而由引理 \ref{sites-cohomology-lemma-trivial-torsor}，它们均有整体截面）。
换言之，可以选取同构
$$
\varphi_{ij} : x_i|_{U_{ij}} \longrightarrow x_j|_{U_{ij}}
$$
。为了得到位于 $U$ 上的对象 $x$，我们将适当调整这些
$\varphi_{ij}$ 的选取以得到下降数据（由于
$p : \mathcal{S} \to \mathcal{C}$ 是叠，该下降数据必为有效）。
成为下降数据的障碍在于上闭链条件可能不成立。令
$U_{ijk} = U_i \times_U U_j \times_U U_k$。于是可以考虑
$$
g_{ijk} = \varphi_{ik}^{-1}|_{U_{ijk}} \circ \varphi_{jk}|_{U_{ijk}} \circ
\varphi_{ij}|_{U_{ijk}}
$$
，这是 $U_{ijk}$ 上 $x_i$ 的一个自同构。因此可以并且确实将
$g_{ijk}$ 视为 $U_{ijk}$ 上 $\mathcal{G}$ 的一个截面。
计算（从略）表明，$(g_{i_0i_1i_2})$ 是 $\mathcal{G}$ 关于覆盖
$\mathcal{U}$ 的 {\v C}ech 复形
${\check C}^\bullet(\mathcal{U}, \mathcal{G})$ 中的一个 $2$-上闭链。
由引理 \ref{sites-cohomology-lemma-cech-spectral-sequence} 的谱序列以及
$H^1(U_i, \mathcal{G}) = 0$ 对所有 $i$ 成立，可知
${\check H}^2(\mathcal{U}, \mathcal{G}) \to H^2(U, \mathcal{G})$
为单射。因此，由假设 $H^2(U, \mathcal{G}) = 0$，
$(g_{i_0i_1i_2})$ 是上边缘。
从而可以找到截面 $g_{ij} \in \mathcal{G}(U_{ij})$，使得
$g_{ik}^{-1}|_{U_{ijk}} g_{jk}|_{U_{ijk}}  g_{ij}|_{U_{ijk}} = g_{ijk}$
对所有 $i,j,k$ 成立。
将 $\varphi_{ij}$ 替换为 $\varphi_{ij}g_{ij}^{-1}$ 后，
可见 $\varphi_{ij}$ 给出了位于 $U_i$ 上的对象 $x_i$ 的下降数据，
证明完毕。
\end{proof}












\section{模的上同调}
\label{sites-cohomology-section-cohomology-modules}

\noindent
前面对阿贝尔层上同调所述的一切，也适用于模的上同调，因为二者一致。

\begin{lemma}
\label{sites-cohomology-lemma-injective-module-injective-presheaf}
设 $(\mathcal{C}, \mathcal{O})$ 为环化站点。
一个内射模层作为范畴 $\textit{PMod}(\mathcal{O})$ 中的对象仍是内射的。
\end{lemma}

\begin{proof}
将《同调代数》引理 \ref{homology-lemma-adjoint-preserve-injectives}
应用于范畴 $\mathcal{A} = \textit{Mod}(\mathcal{O})$、
$\mathcal{B} = \textit{PMod}(\mathcal{O})$、嵌入函子及层化函子。
该引理的所有假设均满足，见《站点上的模》第
\ref{sites-modules-section-sheafification-presheaves-modules} 节。
\end{proof}

\begin{lemma}
\label{sites-cohomology-lemma-include-modules}
设 $(\mathcal{C}, \mathcal{O})$ 为环化站点。考虑函子
$i : \textit{Mod}(\mathcal{C}) \to \textit{PMod}(\mathcal{C})$.
这是左正合函子，其右导出函子由下式给出：
$$
R^pi(\mathcal{F}) = \underline{H}^p(\mathcal{F}) :
U \longmapsto H^p(U, \mathcal{F})
$$
；见第 \ref{sites-cohomology-section-locality} 节中的讨论。
\end{lemma}

\begin{proof}
显然 $i$ 左正合。在 $\textit{Mod}(\mathcal{O})$ 中选取内射分解
$\mathcal{F} \to \mathcal{I}^\bullet$。
根据定义，$R^pi$ 是复形 $\mathcal{I}^\bullet$ 的第 $p$ 个上同调
{\it 预层}。换言之，$R^pi(\mathcal{F})$ 在 $\mathcal{C}$ 的对象
$U$ 上的截面由下式给出：
$$
\frac{\Ker(\mathcal{I}^n(U) \to \mathcal{I}^{n + 1}(U))}
{\Im(\mathcal{I}^{n - 1}(U) \to \mathcal{I}^n(U))}.
$$
这正是 $H^p(U, \mathcal{F})$ 的定义。
\end{proof}

\begin{lemma}
\label{sites-cohomology-lemma-injective-module-trivial-cech}
设 $(\mathcal{C}, \mathcal{O})$ 为环化站点，
$\mathcal{U} = \{U_i \to U\}_{i \in I}$ 为 $\mathcal{C}$ 中的覆盖。
设 $\mathcal{I}$ 为内射 $\mathcal{O}$-模，即
$\textit{Mod}(\mathcal{O})$ 的内射对象。则
$$
\check{H}^p(\mathcal{U}, \mathcal{I}) =
\left\{
\begin{matrix}
\mathcal{I}(U) & \text{若} & p = 0 \\
0 & \text{若} & p > 0
\end{matrix}
\right.
$$
\end{lemma}

\begin{proof}
引理 \ref{sites-cohomology-lemma-cech-map-into} 给出下列一串等式中的第一个等式：
\begin{align*}
\check{\mathcal{C}}^\bullet(\mathcal{U}, \mathcal{I})
& =
\Mor_{\textit{PAb}(\mathcal{C})}(
\mathbf{Z}_{\mathcal{U}, \bullet}, \mathcal{I}) \\
& =
\Mor_{\textit{PMod}(\mathbf{Z})}(
\mathbf{Z}_{\mathcal{U}, \bullet}, \mathcal{I}) \\
& =
\Mor_{\textit{PMod}(\mathcal{O})}(
\mathbf{Z}_{\mathcal{U}, \bullet} \otimes_{p, \mathbf{Z}} \mathcal{O},
\mathcal{I})
\end{align*}
第三个等式来自《站点上的模》引理
\ref{sites-modules-lemma-adjointness-tensor-restrict-presheaves}。
由引理 \ref{sites-cohomology-lemma-injective-module-injective-presheaf}，
$\mathcal{I}$ 是 $\textit{PMod}(\mathcal{O})$ 中的内射对象。
因此 $\Hom_{\textit{PMod}(\mathcal{O})}(-, \mathcal{I})$ 是正合函子。
由引理 \ref{sites-cohomology-lemma-complex-tensored-still-exact}，高阶 {\v C}ech
上同调群均消失。零阶情形见引理 \ref{sites-cohomology-lemma-cech-h0}。
\end{proof}

\begin{lemma}
\label{sites-cohomology-lemma-cohomology-modules-abelian-agree}
设 $\mathcal{C}$ 为站点，$\mathcal{O}$ 为 $\mathcal{C}$ 上的环层。
设 $\mathcal{F}$ 为 $\mathcal{O}$-模，并以 $\mathcal{F}_{ab}$
记其底层阿贝尔群层。则
$$
H^i(\mathcal{C}, \mathcal{F}_{ab})
=
H^i(\mathcal{C}, \mathcal{F})
$$
，且对 $\mathcal{C}$ 的任意对象 $U$，还有
$$
H^i(U, \mathcal{F}_{ab})
=
H^i(U, \mathcal{F}).
$$
这里左边是在 $\textit{Ab}(\mathcal{C})$ 中计算的上同调，
右边是在 $\textit{Mod}(\mathcal{O})$ 中计算的上同调。
\end{lemma}

\begin{proof}
由《导出范畴》引理 \ref{derived-lemma-higher-derived-functors}，
$\delta$-函子 $(\mathcal{F} \mapsto H^p(U, \mathcal{F}))_{p \geq 0}$
是万有的。函子
$\textit{Mod}(\mathcal{O}) \to \textit{Ab}(\mathcal{C})$,
$\mathcal{F} \mapsto \mathcal{F}_{ab}$ 正合。因此
$(\mathcal{F} \mapsto H^p(U, \mathcal{F}_{ab}))_{p \geq 0}$
也是 $\delta$-函子。若能证明
$(\mathcal{F} \mapsto H^p(U, \mathcal{F}_{ab}))_{p \geq 0}$
也万有，则由万有 $\delta$-函子的唯一性可得引理的第二个断言；
见《同调代数》引理
\ref{homology-lemma-uniqueness-universal-delta-functor}。
由于 $\textit{Mod}(\mathcal{O})$ 有充分内射对象，只需证明：
对 $\textit{Mod}(\mathcal{O})$ 中任意内射对象 $\mathcal{I}$，
$H^i(U, \mathcal{I}_{ab}) = 0$；见《同调代数》引理
\ref{homology-lemma-efface-implies-universal}。

\medskip\noindent
设 $\mathcal{I}$ 为 $\textit{Mod}(\mathcal{O})$ 的内射对象。
在引理 \ref{sites-cohomology-lemma-cech-vanish-collection} 中取
$\mathcal{F} = \mathcal{I}$、$\mathcal{B} = \mathcal{C}$ 且
$\text{Cov} = \text{Cov}_\mathcal{C}$。
由引理 \ref{sites-cohomology-lemma-injective-module-trivial-cech}，该引理的假设 (3) 成立。
因此对 $\mathcal{C}$ 的每个对象 $U$，都有
$H^i(U, \mathcal{I}_{ab}) = 0$。

\medskip\noindent
若 $\mathcal{C}$ 有终对象，则这也推出第一个等式。否则，
由《站点》引理 \ref{sites-lemma-topos-good-site}，环化拓扑斯
$(\Sh(\mathcal{C}), \mathcal{O})$ 等价于一个底层站点确有终对象的
环化拓扑斯。因此引理成立。
\end{proof}

\begin{lemma}
\label{sites-cohomology-lemma-cohomology-products}
设 $\mathcal{C}$ 为站点，$I$ 为集合。对 $i \in I$，设
$\mathcal{F}_i$ 为 $\mathcal{C}$ 上的阿贝尔层。设
$U \in \Ob(\mathcal{C})$。典范映射
$$
H^p(U, \prod\nolimits_{i \in I} \mathcal{F}_i)
\longrightarrow
\prod\nolimits_{i \in I} H^p(U, \mathcal{F}_i)
$$
在 $p = 0$ 时为同构，在 $p = 1$ 时为单射。
\end{lemma}

\begin{proof}
对于 $p = 0$，断言成立，因为层的乘积等于其底层预层的乘积；
见《站点》引理 \ref{sites-lemma-limit-sheaf}。
下面证明 $p = 1$ 的情形。令 $\mathcal{F} = \prod \mathcal{F}_i$。
设 $\xi \in H^1(U, \mathcal{F})$ 在
$\prod H^1(U, \mathcal{F}_i)$ 中映为零。由上同调的局部性，
见引理 \ref{sites-cohomology-lemma-kill-cohomology-class-on-covering}，存在覆盖
$\mathcal{U} = \{U_j \to U\}$，使得对所有 $j$ 都有
$\xi|_{U_j} = 0$。由引理 \ref{sites-cohomology-lemma-cech-h1}，这意味着
$\xi$ 来自某个元素
$\check \xi \in \check H^1(\mathcal{U}, \mathcal{F})$.
由于对所有 $i$，映射
$\check H^1(\mathcal{U}, \mathcal{F}_i) \to H^1(U, \mathcal{F}_i)$
均为单射（由引理 \ref{sites-cohomology-lemma-cech-h1}），并且 $\xi$ 在
$\prod H^1(U, \mathcal{F}_i)$ 中的像为零，故其像
$\check \xi_i$ 在 $\check H^1(\mathcal{U}, \mathcal{F}_i)$ 中为零。
另一方面，由于 $\mathcal{F} = \prod \mathcal{F}_i$，
$\check{\mathcal{C}}^\bullet(\mathcal{U}, \mathcal{F})$ 是诸复形
$\check{\mathcal{C}}^\bullet(\mathcal{U}, \mathcal{F}_i)$ 的乘积。
因此由《同调代数》引理
\ref{homology-lemma-product-abelian-groups-exact}，可得所需结论
$\check \xi = 0$。
\end{proof}

\begin{lemma}
\label{sites-cohomology-lemma-restriction-along-monomorphism-surjective}
设 $(\mathcal{C}, \mathcal{O})$ 为环化站点，
$a : U' \to U$ 为 $\mathcal{C}$ 中的单态射。则对任意内射
$\mathcal{O}$-模 $\mathcal{I}$，限制映射
$\mathcal{I}(U) \to \mathcal{I}(U')$ 为满射。
\end{lemma}

\begin{proof}
令 $j : \mathcal{C}/U \to \mathcal{C}$ 与
$j' : \mathcal{C}/U' \to \mathcal{C}$ 为局部化态射
（《站点上的模》第 \ref{sites-modules-section-localize} 节）。
由于 $j_!$ 是限制函子的左伴随，对任意 $\mathcal{O}$-模层
$\mathcal{F}$，有
$$
\Hom_\mathcal{O}(j_!\mathcal{O}_U, \mathcal{F})
=
\Hom_{\mathcal{O}_U}(\mathcal{O}_U, \mathcal{F}|_U)
=
\mathcal{F}(U)
$$
。类似地，层 $j'_!\mathcal{O}_{U'}$ 表示函子
$\mathcal{F} \mapsto \mathcal{F}(U')$。此外，下面构造一个典范的
$\mathcal{O}$-模态射
$$
j'_!\mathcal{O}_{U'} \longrightarrow j_!\mathcal{O}_U
$$
；由 Yoneda 引理（《范畴》引理 \ref{categories-lemma-yoneda}），
它对应于限制映射 $\mathcal{F}(U) \to \mathcal{F}(U')$。
只需证明所显示的模态射是单射；见《同调代数》引理
\ref{homology-lemma-characterize-injectives}。

\medskip\noindent
为构造此映射，只需在如下两个预层之间构造映射：它们分别把
$\mathcal{C}$ 的对象 $V$ 送到 $\mathcal{O}(V)$-模
$$
\bigoplus\nolimits_{\varphi' \in \Mor_\mathcal{C}(V, U')} \mathcal{O}(V)
\quad\text{与}\quad
\bigoplus\nolimits_{\varphi \in \Mor_\mathcal{C}(V, U)} \mathcal{O}(V)
$$
；见《站点上的模》引理
\ref{sites-modules-lemma-extension-by-zero}。
取这样一个映射：它通过 $\mathcal{O}(V)$ 上的恒等映射，
把对应于 $\varphi'$ 的直和项映到对应于
$\varphi = a \circ \varphi'$ 的直和项。由于 $a$ 是单态射，
此映射为单射。层化正合，故结论成立。
\end{proof}






\section{完全无环层}
\label{sites-cohomology-section-limp}

\noindent
设 $(\mathcal{C}, \mathcal{O})$ 为环化站点。
设 $K$ 为 $\mathcal{C}$ 上的集合预层（这里特意使用正体大写字母，
以便与阿贝尔层区分）。给定 $\mathcal{O}$-模层 $\mathcal{F}$，令
$$
\mathcal{F}(K) =
\Mor_{\textit{PSh}(\mathcal{C})}(K, \mathcal{F}) =
\Mor_{\Sh(\mathcal{C})}(K^\#, \mathcal{F})
$$
。函子 $\mathcal{F} \mapsto \mathcal{F}(K)$ 是左正合函子
$\textit{Mod}(\mathcal{O}) \to \textit{Ab}$，故有右导出函子。
我们以 $H^p(K, \mathcal{F})$ 记这些右导出函子，从而
$H^0(K, \mathcal{F}) = \mathcal{F}(K)$。

\medskip\noindent
有如下几点观察：
\begin{enumerate}
\item 由于 $\mathcal{F}(K) = \mathcal{F}(K^\#)$，有
$H^p(K, \mathcal{F}) = H^p(K^\#, \mathcal{F})$.
上述定义允许 $K$ 为预层，仅是记号上的便利。
\item 假设对 $\mathcal{C}$ 的某个对象 $U$，有
$K = h_U$ 或 $K = h_U^\#$。则
$H^p(K, \mathcal{F}) = H^p(U, \mathcal{F})$，因为
$\Mor_{\Sh(\mathcal{C})}(h_U^\#, \mathcal{F}) = \mathcal{F}(U)$；
见《站点》第 \ref{sites-section-representable-sheaves} 节。
\item 若 $\mathcal{O} = \mathbf{Z}$（常值层），则上同调群给出函子
$H^p(K, - ) : \textit{Ab}(\mathcal{C}) \to \textit{Ab}$
，因为在此情形下
$\textit{Mod}(\mathcal{O}) = \textit{Ab}(\mathcal{C})$。
\end{enumerate}
用这种语言可以重述一些已经证明的结果。

\begin{lemma}
\label{sites-cohomology-lemma-compute-cohomology-on-sheaf-sets}
设 $(\mathcal{C}, \mathcal{O})$ 为环化站点，
$K$ 为 $\mathcal{C}$ 上的集合预层。设 $\mathcal{F}$ 为
$\mathcal{O}$-模，并以 $\mathcal{F}_{ab}$ 记其底层阿贝尔群层。
则 $H^p(K, \mathcal{F}) = H^p(K, \mathcal{F}_{ab})$。
\end{lemma}

\begin{proof}
可以用 $K$ 的层化替换 $K$，并假设 $K$ 是层。
注意，$H^p(K, \mathcal{F})$ 与 $H^p(K, \mathcal{F}_{ab})$
都只依赖于拓扑斯，而不依赖于底层站点。因此由《站点》引理
\ref{sites-lemma-topos-good-site}，可以用一个``更大''站点替换
$\mathcal{C}$，使得对 $\mathcal{C}$ 的某个对象 $U$，有
$K = h_U$。此时结论由引理
\ref{sites-cohomology-lemma-cohomology-modules-abelian-agree} 得出。
\end{proof}

\begin{lemma}
\label{sites-cohomology-lemma-cech-to-cohomology-sheaf-sets}
设 $\mathcal{C}$ 为站点，$K' \to K$ 为 $\mathcal{C}$ 上集合预层的
态射，且其层化为满射。令
$K'_p = K' \times_K \ldots \times_K K'$（$p + 1$ 个因子）。
对每个阿贝尔层 $\mathcal{F}$，存在谱序列，其
$E_1^{p, q} = H^q(K'_p, \mathcal{F})$，并收敛到
$H^{p + q}(K, \mathcal{F})$。
\end{lemma}

\begin{proof}
由于层化正合，可知 $(K_p')^\#$ 等于
$(K')^\# \times_{K^\#} \ldots \times_{K^\#} (K')^\#$
（$p + 1$ 个因子）。因此可以用各自的层化替换 $K$ 与 $K'$，
并假设 $K \to K'$ 是层的满射。按照《站点》引理
\ref{sites-lemma-topos-good-site}，用一个``更大''站点替换
$\mathcal{C}$ 后，可以假设 $K,K'$ 是 $\mathcal{C}$ 的对象，且
$\mathcal{U} = \{K' \to K\}$ 是覆盖。于是得到引理
\ref{sites-cohomology-lemma-cech-spectral-sequence} 的 {\v C}ech 到上同调谱序列，
其 $E_1$ 页如引理陈述所示。
\end{proof}

\begin{lemma}
\label{sites-cohomology-lemma-cohomology-on-sheaf-sets}
设 $\mathcal{C}$ 为站点，$K$ 为 $\mathcal{C}$ 上的集合层。
考虑拓扑斯态射
$j : \Sh(\mathcal{C}/K) \to \Sh(\mathcal{C})$，见《站点》引理
\ref{sites-lemma-localize-topos-site}。则 $j^{-1}$ 保持内射对象，且
$H^p(K, \mathcal{F}) = H^p(\mathcal{C}/K, j^{-1}\mathcal{F})$
对 $\mathcal{C}$ 上的任意阿贝尔层 $\mathcal{F}$ 成立。
\end{lemma}

\begin{proof}
由《站点》引理 \ref{sites-lemma-localize-topos} 与
\ref{sites-lemma-localize-topos-site}，拓扑斯态射 $j$
等价于一个局部化。因此结论由引理
\ref{sites-cohomology-lemma-cohomology-of-open} 得出。
\end{proof}

\noindent
结合引理 \ref{sites-cohomology-lemma-compute-cohomology-on-sheaf-sets} 可知，
对环化站点上的模层而言，以下定义也是``正确的''定义。

\begin{definition}
\label{sites-cohomology-definition-limp}
设 $\mathcal{C}$ 为站点。若对每个集合层 $K$ 及所有 $p \geq 1$，
都有 $H^p(K, \mathcal{F}) = 0$，则称阿贝尔层 $\mathcal{F}$ 为
{\it 完全无环的}\footnote{尽管这一术语用于
\cite[Vbis, Proposition 1.3.10]{SGA4}，它很可能不是标准术语。
在 \cite[V, Definition 4.1]{SGA4} 中，这一性质称为 ``flasque''，
但我们不能采用该术语，因为它会与我们对拓扑空间上 flasque 层的定义冲突。
如果你有更好的建议，请发送电子邮件至
\href{mailto:stacks.project@gmail.com}{stacks.project@gmail.com}
。}
\end{definition}

\noindent
显然，完全无环是内蕴性质，即在拓扑斯等价下保持。
完全无环层在站点的所有对象上都有消失的高阶上同调，
但一般而言，完全无环的条件严格更强。
下面给出完全无环层的一个有时很有用的刻画。

\begin{lemma}
\label{sites-cohomology-lemma-characterize-limp}
设 $\mathcal{C}$ 为站点，$\mathcal{F}$ 为阿贝尔层。若
\begin{enumerate}
\item 对 $p > 0$ 及 $U \in \Ob(\mathcal{C})$，有
$H^p(U, \mathcal{F}) = 0$；
\item 对集合层的每个满射 $K' \to K$，增广 {\v C}ech 复形
$$
0 \to H^0(K, \mathcal{F}) \to H^0(K', \mathcal{F}) \to
H^0(K' \times_K K', \mathcal{F}) \to \ldots
$$
正合，
\end{enumerate}
则 $\mathcal{F}$ 完全无环（反之亦然）。
\end{lemma}

\begin{proof}
由假设 (1)，对所有 $p > 0$ 及 $\mathcal{C}$ 的所有对象 $U$，有
$H^p(h_U^\#, g^{-1}\mathcal{I}) = 0$。注意，若
$K = \coprod K_i$ 是 $\mathcal{C}$ 上集合层的余积，则
$H^p(K, g^{-1}\mathcal{I}) = \prod H^p(K_i, g^{-1}\mathcal{I})$。
对任意集合层 $K$，存在满射
$$
K' = \coprod h_{U_i}^\# \longrightarrow K
$$
；见《站点》引理 \ref{sites-lemma-sheaf-coequalizer-representable}。
因此可得：(*) 对每个集合层 $K$，存在集合层的满射
$K' \to K$，使得对 $p > 0$ 有 $H^p(K', \mathcal{F}) = 0$。
我们断言，(*) 与条件 (2) 推出 $\mathcal{F}$ 完全无环。
注意，条件 (*) 与 (2) 只依赖于 $\mathcal{F}$ 作为拓扑斯
$\Sh(\mathcal{C})$ 中的对象，而不依赖于底层站点。
（在证明的其余部分将不再使用性质 (1)。）

\medskip\noindent
我们对 $n \geq 0$ 归纳证明 (*) 与 (2) 推出下列归纳假设 $IH_n$：
对所有 $0 < p \leq n$ 及所有集合层 $K$，有
$H^p(K, \mathcal{F}) = 0$。显然 $IH_0$ 成立。假设 $IH_n$ 成立。
选取集合层 $K$，并选取满射 $K' \to K$，使得对所有 $p > 0$，
$H^p(K', \mathcal{F}) = 0$。有谱序列
$$
E_1^{p, q} = H^q(K'_p, \mathcal{F})
$$
收敛到 $H^{p + q}(K, \mathcal{F})$；见引理
\ref{sites-cohomology-lemma-cech-to-cohomology-sheaf-sets}。
由 $IH_n$，对 $0 < q \leq n$ 有 $E_1^{p, q} = 0$；由假设 (2)，
对 $p > 0$ 有 $E_2^{p, 0} = 0$。最后，由 $K'$ 的选取，
$H^q(K', \mathcal{F}) = 0$，故对 $q > 0$ 有
$E_1^{0, q} = 0$。因此 $p + q = n + 1$ 的所有项
$E_2^{p, q}$ 均为零，从而 $H^{n + 1}(K, \mathcal{F}) = 0$。
\end{proof}







\section{Leray 谱序列}
\label{sites-cohomology-section-leray}

\noindent
证明 Leray 谱序列存在性的关键是下列引理。

\begin{lemma}
\label{sites-cohomology-lemma-direct-image-injective-sheaf}
设 $f : (\Sh(\mathcal{C}), \mathcal{O}_\mathcal{C}) \to
(\Sh(\mathcal{D}), \mathcal{O}_\mathcal{D})$ 为环化拓扑斯的态射。
则对 $\textit{Mod}(\mathcal{O}_\mathcal{C})$ 中的任意内射对象
$\mathcal{I}$，推前 $f_*\mathcal{I}$ 完全无环。
\end{lemma}

\begin{proof}
设 $K$ 为 $\mathcal{D}$ 上的集合层。由《站点上的模》引理
\ref{sites-modules-lemma-morphism-ringed-topoi-comes-from-morphism-ringed-sites}
，可以用``更大''站点替换 $\mathcal{C},\mathcal{D}$，使得 $f$
来自一个由连续函子 $u : \mathcal{D} \to \mathcal{C}$ 诱导的
环化站点态射，并且对 $\mathcal{D}$ 的某个对象 $V$，有
$K = h_V$。

\medskip\noindent
因此，当 $f$ 由环化站点态射给出时，只需证明：对 $q > 0$
及 $\mathcal{D}$ 的所有对象 $V$，$H^q(V, f_*\mathcal{I})$ 为零。
设 $\mathcal{V} = \{V_j \to V\}$ 为 $\mathcal{D}$ 的任意覆盖。
由于 $u$ 连续，$\mathcal{U} = \{u(V_j) \to u(V)\}$
是 $\mathcal{C}$ 的覆盖。于是由 $f_*$ 的定义，有 {\v C}ech 复形的等式
$$
\check{\mathcal{C}}^\bullet(\mathcal{V}, f_*\mathcal{I})
=
\check{\mathcal{C}}^\bullet(\mathcal{U}, \mathcal{I})
$$
。由引理 \ref{sites-cohomology-lemma-injective-module-trivial-cech}，
此复形的正次数上同调为零。结论遂由引理
\ref{sites-cohomology-lemma-cech-vanish-collection} 得出。
\end{proof}

\noindent
对于平坦态射，函子 $f_*$ 保持内射模。特别地，函子
$f_* : \textit{Ab}(\mathcal{C}) \to \textit{Ab}(\mathcal{D})$
总把内射阿贝尔层变为内射阿贝尔层。

\begin{lemma}
\label{sites-cohomology-lemma-pushforward-injective-flat}
设 $f : (\Sh(\mathcal{C}), \mathcal{O}_\mathcal{C}) \to
(\Sh(\mathcal{D}), \mathcal{O}_\mathcal{D})$ 为环化拓扑斯态射。
若 $f$ 平坦，则对任意内射 $\mathcal{O}_\mathcal{C}$-模
$\mathcal{I}$，$f_*\mathcal{I}$ 是内射
$\mathcal{O}_\mathcal{D}$-模。
\end{lemma}

\begin{proof}
此时函子 $f^*$ 正合，见《站点上的模》引理
\ref{sites-modules-lemma-flat-pullback-exact}。
因此结论由《同调代数》引理
\ref{homology-lemma-adjoint-preserve-injectives} 得出。
\end{proof}

\begin{lemma}
\label{sites-cohomology-lemma-limp-acyclic}
设 $(\Sh(\mathcal{C}), \mathcal{O}_\mathcal{C})$ 为环化拓扑斯。
完全无环层对下列函子右无环：
\begin{enumerate}
\item 对 $\mathcal{C}$ 的任意对象 $U$，函子 $H^0(U, -)$；
\item 对任意集合预层 $K$，函子 $\mathcal{F} \mapsto \mathcal{F}(K)$；
\item 整体截面函子 $\Gamma(\mathcal{C}, -)$；
\item 对环化拓扑斯的任意态射
$f : (\Sh(\mathcal{C}), \mathcal{O}_\mathcal{C}) \to
(\Sh(\mathcal{D}), \mathcal{O}_\mathcal{D})$，函子 $f_*$。
\end{enumerate}
\end{lemma}

\begin{proof}
(2) 是完全无环层的定义。正如完全无环层定义后的讨论所指出，
(1) 是 (2) 的推论。(3) 是 (2) 中取 $K = e$ 为
$\Sh(\mathcal{C})$ 的终对象时的特殊情形。

\medskip\noindent
为证明 (4)，由《站点上的模》引理
\ref{sites-modules-lemma-morphism-ringed-topoi-comes-from-morphism-ringed-sites}
，可以假设 $f$ 由站点态射给出。在此情形下，由引理
\ref{sites-cohomology-lemma-higher-direct-images} 中对高阶直接像的描述，
完全无环层的 $R^if_*$ 在 $i > 0$ 时为零。
\end{proof}

\begin{remark}
\label{sites-cohomology-remark-before-Leray}
作为上述结果的推论，《导出范畴》引理
\ref{derived-lemma-compose-derived-functors} 适用于多种情形。
例如，给定环化拓扑斯态射
$f : (\Sh(\mathcal{C}), \mathcal{O}_\mathcal{C}) \to
(\Sh(\mathcal{D}), \mathcal{O}_\mathcal{D})$，有
$$
R\Gamma(\mathcal{D}, Rf_*\mathcal{F}) = R\Gamma(\mathcal{C}, \mathcal{F})
$$
对任意 $\mathcal{O}_\mathcal{C}$-模层 $\mathcal{F}$ 成立。
事实上，对内射 $\mathcal{O}_\mathcal{X}$-模 $\mathcal{I}$，
由引理 \ref{sites-cohomology-lemma-direct-image-injective-sheaf}，
$\mathcal{O}_\mathcal{D}$-模 $f_*\mathcal{I}$ 完全无环；
而由引理 \ref{sites-cohomology-lemma-limp-acyclic}，完全无环层对
$\Gamma(\mathcal{D}, -)$ 无环。
\end{remark}

\begin{lemma}[Leray 谱序列]
\label{sites-cohomology-lemma-Leray}
设 $f : (\Sh(\mathcal{C}), \mathcal{O}_\mathcal{C}) \to
(\Sh(\mathcal{D}), \mathcal{O}_\mathcal{D})$ 为环化拓扑斯态射。
设 $\mathcal{F}^\bullet$ 为 $\mathcal{O}_\mathcal{C}$-模的下有界复形。
存在谱序列
$$
E_2^{p, q} = H^p(\mathcal{D}, R^qf_*(\mathcal{F}^\bullet))
$$
收敛到 $H^{p + q}(\mathcal{C}, \mathcal{F}^\bullet)$。
\end{lemma}

\begin{proof}
这正是《导出范畴》引理
\ref{derived-lemma-grothendieck-spectral-sequence} 中由函子复合
$\Gamma(\mathcal{C}, -) = \Gamma(\mathcal{D}, -) \circ f_*$.
产生的 Grothendieck 谱序列。关于《导出范畴》引理
\ref{derived-lemma-grothendieck-spectral-sequence} 的假设为何满足，
见引理 \ref{sites-cohomology-lemma-direct-image-injective-sheaf} 与
\ref{sites-cohomology-lemma-limp-acyclic}。
\end{proof}

\begin{lemma}
\label{sites-cohomology-lemma-apply-Leray}
设 $f : (\Sh(\mathcal{C}), \mathcal{O}_\mathcal{C}) \to
(\Sh(\mathcal{D}), \mathcal{O}_\mathcal{D})$ 为环化拓扑斯态射，
$\mathcal{F}$ 为 $\mathcal{O}_\mathcal{C}$-模。
\begin{enumerate}
\item 若对 $q > 0$ 有 $R^qf_*\mathcal{F} = 0$，则对所有 $p$，
$H^p(\mathcal{C}, \mathcal{F}) = H^p(\mathcal{D}, f_*\mathcal{F})$。
\item 若对所有 $q$ 及 $p > 0$ 有
$H^p(\mathcal{D}, R^qf_*\mathcal{F}) = 0$，则对所有 $q$，
$H^q(\mathcal{C}, \mathcal{F}) = H^0(\mathcal{D}, R^qf_*\mathcal{F})$。
\end{enumerate}
\end{lemma}

\begin{proof}
这是迫使 Leray 谱序列收敛的两个简单条件。
也可以直接证明这些事实（不使用谱序列），这是层上同调的一个好练习。
\end{proof}

\begin{lemma}[相对 Leray 谱序列]
\label{sites-cohomology-lemma-relative-Leray}
设
$f : (\Sh(\mathcal{C}), \mathcal{O}_\mathcal{C}) \to
(\Sh(\mathcal{D}), \mathcal{O}_\mathcal{D})$
与
$g : (\Sh(\mathcal{D}), \mathcal{O}_\mathcal{D}) \to
(\Sh(\mathcal{E}), \mathcal{O}_\mathcal{E})$
为环化拓扑斯态射。设 $\mathcal{F}$ 为
$\mathcal{O}_\mathcal{C}$-模。存在谱序列，其
$$
E_2^{p, q} = R^pg_*(R^qf_*\mathcal{F})
$$
收敛到 $R^{p + q}(g \circ f)_*\mathcal{F}$。
此谱序列关于 $\mathcal{F}$ 具有函子性，并且对
$\mathcal{O}_\mathcal{C}$-模的下有界复形也有相应版本。
\end{lemma}

\begin{proof}
这是函子复合的 Grothendieck 谱序列；见《导出范畴》引理
\ref{derived-lemma-grothendieck-spectral-sequence} 以及引理
\ref{sites-cohomology-lemma-direct-image-injective-sheaf} 与
\ref{sites-cohomology-lemma-limp-acyclic}。
\end{proof}







\section{基变换映射}
\label{sites-cohomology-section-base-change-map}

\noindent
本节在若干情形下构造基变换映射；一般情形在注记
\ref{sites-cohomology-remark-base-change} 中处理。本节的讨论限制于沿环化站点平坦态射
作基变换的情形，从而避免使用导出拉回。
在陈述结果之前，先讨论导出范畴上的平坦拉回。假设
$g : (\Sh(\mathcal{C}), \mathcal{O}_\mathcal{C})
\to (\Sh(\mathcal{D}), \mathcal{O}_\mathcal{D})$
为环化拓扑斯的平坦态射。由《站点上的模》引理
\ref{sites-modules-lemma-flat-pullback-exact}，函子
$g^* : \textit{Mod}(\mathcal{O}_\mathcal{D}) \to
\textit{Mod}(\mathcal{O}_\mathcal{C})$ 正合。因此它有导出函子
$$
g^* : D(\mathcal{O}_\mathcal{D}) \to D(\mathcal{O}_\mathcal{C})
$$
；它只需拉回 $D(\mathcal{O}_\mathcal{D})$ 中给定对象的一个代表即可计算，
见《导出范畴》引理 \ref{derived-lemma-right-derived-exact-functor}。
它保持上有界与下有界子范畴。因此如上所示，我们以 $g^*$
而非 $Lg^*$ 记此函子。

\begin{lemma}
\label{sites-cohomology-lemma-base-change-map-flat-case}
设
$$
\xymatrix{
(\Sh(\mathcal{C}'), \mathcal{O}_{\mathcal{C}'})
\ar[r]_{g'} \ar[d]_{f'} &
(\Sh(\mathcal{C}), \mathcal{O}_\mathcal{C}) \ar[d]^f \\
(\Sh(\mathcal{D}'), \mathcal{O}_{\mathcal{D}'})
\ar[r]^g &
(\Sh(\mathcal{D}), \mathcal{O}_\mathcal{D})
}
$$
为环化拓扑斯的交换图。设 $\mathcal{F}^\bullet$ 为
$\mathcal{O}_\mathcal{C}$-模的下有界复形。假设 $g$ 与 $g'$ 均平坦。
则 $D^{+}(\mathcal{O}_{\mathcal{D}'})$ 中存在典范基变换映射
$$
g^*Rf_*\mathcal{F}^\bullet
\longrightarrow
R(f')_*(g')^*\mathcal{F}^\bullet
$$
。
\end{lemma}

\begin{proof}
选取内射分解 $\mathcal{F}^\bullet \to \mathcal{I}^\bullet$ 与
$(g')^*\mathcal{F}^\bullet \to \mathcal{J}^\bullet$。
由引理 \ref{sites-cohomology-lemma-pushforward-injective-flat}，
$(g')_*\mathcal{J}^\bullet$ 是表示
$R(g')_*(g')^*\mathcal{F}^\bullet$ 的内射对象复形。因此由
《导出范畴》引理 \ref{derived-lemma-morphisms-lift} 与
\ref{derived-lemma-morphisms-equal-up-to-homotopy}，下图中的箭头
$\beta$ 存在且在同伦意义下唯一：
$$
\xymatrix{
(g')_*(g')^*\mathcal{F}^\bullet \ar[r] &
(g')_*\mathcal{J}^\bullet \\
\mathcal{F}^\bullet \ar[u]^{\text{伴随}} \ar[r] &
\mathcal{I}^\bullet \ar[u]_\beta
}
$$
。推前到 $\mathcal{D}$，得到
$$
f_*\beta :
f_*\mathcal{I}^\bullet
\longrightarrow
f_*(g')_*\mathcal{J}^\bullet
=
g_*(f')_*\mathcal{J}^\bullet
$$
。由 $g^*$ 与 $g_*$ 的伴随性，得到复形态射
$g^*f_*\mathcal{I}^\bullet \to (f')_*\mathcal{J}^\bullet$。
注意，此态射在同伦意义下唯一，因为整个过程中唯一的选取是
态射 $\beta$ 的选取，并且一切都在复形层次上完成。
\end{proof}








\section{上同调与余极限}
\label{sites-cohomology-section-limits}

\noindent
设 $(\mathcal{C}, \mathcal{O})$ 为环化站点。设
$\mathcal{I} \to \textit{Mod}(\mathcal{O})$，
$i \mapsto \mathcal{F}_i$ 为指标范畴 $\mathcal{I}$ 上的图式，
见《范畴》第 \ref{categories-section-limits} 节。
对每个 $i$ 都有典范映射
$\mathcal{F}_i \to \colim_i \mathcal{F}_i$，它诱导上同调之间的映射。
因此，对每个 $p \geq 0$ 及 $\mathcal{C}$ 的每个对象 $U$，
得到典范映射
$$
\colim_i H^p(U, \mathcal{F}_i)
\longrightarrow
H^p(U, \colim_i \mathcal{F}_i)
$$
。一般而言，这些映射不是同构，即使在 $p = 0$ 时也不一定是。

\medskip\noindent
下列引理是《站点》引理
\ref{sites-lemma-directed-colimits-sections} 在上同调中的对应结论。

\begin{lemma}
\label{sites-cohomology-lemma-colim-works-over-collection}
设 $\mathcal{C}$ 为站点，$\text{Cov}_\mathcal{C}$ 为
$\mathcal{C}$ 的覆盖所成的集合（见《站点》定义
\ref{sites-definition-site}）。设
$\mathcal{B} \subset \Ob(\mathcal{C})$，且
$\text{Cov} \subset \text{Cov}_\mathcal{C}$
为子集。假设
\begin{enumerate}
\item 对每个 $\mathcal{U} \in \text{Cov}$，有
$\mathcal{U} = \{U_i \to U\}_{i \in I}$，其中 $I$ 有限，
$U,U_i \in \mathcal{B}$，并且每个
$U_{i_0} \times_U \ldots \times_U U_{i_p} \in \mathcal{B}$；
\item 对每个 $U \in \mathcal{B}$，$\text{Cov}$ 中出现的
$U$ 的覆盖组成 $U$ 的覆盖的一个共尾系。
\end{enumerate}
则映射
$$
\colim_i H^p(U, \mathcal{F}_i)
\longrightarrow
H^p(U, \colim_i \mathcal{F}_i)
$$
对每个 $p \geq 0$、每个 $U \in \mathcal{B}$ 以及每个滤过图式
$\mathcal{I} \to \textit{Ab}(\mathcal{C})$ 均为同构。
\end{lemma}

\begin{proof}
对 $p$ 归纳证明此引理。注意，在 (1) 中我们要求覆盖
$\mathcal{U} \in \text{Cov}$ 有限，故 $\mathcal{B}$ 的所有元素均拟紧。
因此 (2) 与 (1) 推出任意 $U \in \mathcal{B}$ 均满足《站点》引理
\ref{sites-lemma-directed-colimits-sections} (4) 的假设。
于是结论对 $p = 0$ 成立。现在假设引理对 $p$ 成立，并证明
$p + 1$ 的情形。

\medskip\noindent
选取滤过图式
$\mathcal{F} : \mathcal{I} \to \textit{Ab}(\mathcal{C})$,
$i \mapsto \mathcal{F}_i$.
由于 $\textit{Ab}(\mathcal{C})$ 有函子性的内射嵌入，见《内射对象》
定理 \ref{injectives-theorem-sheaves-injectives}，可以找到滤过图式的态射
$\mathcal{F} \to \mathcal{I}$，使得每个
$\mathcal{F}_i \to \mathcal{I}_i$ 都是从阿贝尔层到内射阿贝尔层的
单射。以 $\mathcal{Q}_i$ 记余核，从而有短正合列
$$
0 \to
\mathcal{F}_i \to
\mathcal{I}_i \to
\mathcal{Q}_i \to 0.
$$
层的余极限是预层层次上的余极限的层化，层化正合，并且阿贝尔群的
滤过余极限正合（见《代数》引理
\ref{algebra-lemma-directed-colimit-exact}），故序列
$$
0 \to
\colim_i \mathcal{F}_i \to
\colim_i \mathcal{I}_i \to
\colim_i \mathcal{Q}_i \to 0.
$$
也是短正合列。我们断言，对所有 $U \in \mathcal{B}$ 及所有
$q \geq 1$，有 $H^q(U, \colim_i \mathcal{I}_i) = 0$。
暂且承认此断言，考虑交换图
$$
\xymatrix{
\colim_i H^p(U, \mathcal{I}_i) \ar[d] \ar[r] &
\colim_i H^p(U, \mathcal{Q}_i) \ar[d] \ar[r] &
\colim_i H^{p + 1}(U, \mathcal{F}_i) \ar[d] \ar[r] &
0 \ar[d] \\
H^p(U, \colim_i \mathcal{I}_i) \ar[r] &
H^p(U, \colim_i \mathcal{Q}_i) \ar[r] &
H^{p + 1}(U, \colim_i \mathcal{F}_i) \ar[r] &
0
}
$$
右下角的零来自该断言，右上角的零来自层 $\mathcal{I}_i$ 的内射性。
应用《代数》引理 \ref{algebra-lemma-directed-colimit-exact} 可知上行正合。
因此由蛇引理可得 $p + 1$ 的结论。

\medskip\noindent
还需证明该断言。我们将使用引理
\ref{sites-cohomology-lemma-cech-vanish-collection}。由 $p = 0$ 的结论，
对 $\mathcal{U} \in \text{Cov}$ 有
$$
\check{\mathcal{C}}^\bullet(\mathcal{U}, \colim_i \mathcal{I}_i)
=
\colim_i \check{\mathcal{C}}^\bullet(\mathcal{U}, \mathcal{I}_i)
$$
，因为所有 $U_{j_0} \times_U \ldots \times_U U_{j_p}$
均属于 $\mathcal{B}$。由引理 \ref{sites-cohomology-lemma-injective-trivial-cech}，
{\v C}ech 复形余极限中的每个复形都在次数 $\geq 1$ 无环。
因此由《代数》引理 \ref{algebra-lemma-directed-colimit-exact}，
{\v C}ech 复形
$\check{\mathcal{C}}^\bullet(\mathcal{U}, \colim_i \mathcal{I}_i)$
也在次数 $\geq 1$ 无环。换言之，对所有 $p \geq 1$，有
$\check{H}^p(\mathcal{U}, \colim_i \mathcal{I}_i) = 0$。
故引理 \ref{sites-cohomology-lemma-cech-vanish-collection} 的假设均满足，断言得证。
\end{proof}

\begin{lemma}
\label{sites-cohomology-lemma-colim-global}
设 $\mathcal{C}$ 为站点，$S \subset \Ob(\Sh(\mathcal{C}))$
为子集。以 $*$ 记 $\Sh(\mathcal{C})$ 的终对象。假设
\begin{enumerate}
\item 对某个 $K \in S$，映射 $K \to *$ 为满射；
\item 给定层的满射 $\mathcal{F} \to K$，其中 $K \in S$，
存在 $K' \in S$ 及态射 $K' \to \mathcal{F}$，使得复合
$K' \to K$ 为满射；
\item 给定 $K,K' \in S$，存在满射 $K'' \to K \times K'$，
其中 $K'' \in S$；
\item 给定 $a,b : K \to K'$，其中 $K,K' \in S$，存在满射
$K'' \to \text{等化子}(a, b)$，其中 $K'' \in S$；
\item 每个 $K \in S$ 均拟紧（《站点》定义
\ref{sites-definition-quasi-compact-topos}）。
\end{enumerate}
则对所有 $p \geq 0$，映射
$$
\colim_\lambda H^p(\mathcal{C}, \mathcal{F}_\lambda)
\longrightarrow
H^p(\mathcal{C}, \colim_\lambda \mathcal{F}_\lambda)
$$
对每个滤过图式 $\Lambda \to \textit{Ab}(\mathcal{C})$，
$\lambda \mapsto \mathcal{F}_\lambda$ 均为同构。
\end{lemma}

\begin{proof}
对 $p$ 归纳证明。基本情形 $p = 0$ 由《站点》引理
\ref{sites-lemma-directed-colimits-global-sections} 的 (4) 得出。
下面核对其假设；不过我们建议读者略过这一部分。
假设 $\mathcal{F} \to *$ 为满射。按 (1) 选取 $K \in S$，
使得 $K \to *$ 为满射。于是 $\mathcal{F} \times K \to K$ 为满射。
按 (2) 选取 $K' \to \mathcal{F} \times K$，其中 $K' \in S$，
且 $K' \to K$ 为满射。于是有态射 $K' \to \mathcal{F}$，
并且 $K' \to *$ 为满射。因此《站点》引理
\ref{sites-lemma-directed-colimits-global-sections} 的假设 (4)(a) 满足。
由《站点》引理 \ref{sites-lemma-image-of-quasi-compact} 以及假设
(3)、(5)，对所有 $K \in S$，$K \times K$ 均拟紧。
因此《站点》引理 \ref{sites-lemma-directed-colimits-global-sections}
的假设 (4)(b) 满足。基本情形证毕。

\medskip\noindent
归纳步骤。假设对于 $p \leq p_0$，以及所有可找到满足 (1)--(5)
的集合 $S \subset \Ob(\Sh(\mathcal{C}))$ 的拓扑斯
$\Sh(\mathcal{C})$，结论对 $H^p$ 成立。完全按照引理
\ref{sites-cohomology-lemma-colim-works-over-collection} 的证明论证可知，
只需证明：给定滤过余极限
$\mathcal{I} = \colim \mathcal{I}_\lambda$，其中
$\mathcal{I}_\lambda$ 是内射阿贝尔层，则
$H^{p_0 + 1}(\mathcal{C}, \mathcal{I}) = 0$。
按 (1) 选取满射 $K \to *$，其中 $K \in S$。
以 $K^n$ 记 $K$ 的 $n$ 重自乘积。考虑引理
\ref{sites-cohomology-lemma-cech-to-cohomology-sheaf-sets} 的谱序列
$$
E_1^{p, q} = H^q(K^{p + 1}, \mathcal{I}) \Rightarrow
H^{p + q}(*, \mathcal{I}) = H^{p + q}(\mathcal{C}, \mathcal{I})
$$
。回忆，对 $\mathcal{C}$ 上的任意阿贝尔层 $\mathcal{F}$，有
$H^q(K^{p + 1}, \mathcal{F}) = H^q(\mathcal{C}/K^{p + 1}, j^{-1}\mathcal{F})$,
；见引理 \ref{sites-cohomology-lemma-cohomology-on-sheaf-sets}。由于 $j^{-1}$
与余极限交换，有
$j^{-1}\mathcal{I} = \colim j^{-1}\mathcal{I}_\lambda$
。由引理 \ref{sites-cohomology-lemma-cohomology-of-open}，限制
$j^{-1}\mathcal{I}_\lambda$ 是 $\mathcal{C}/K^{p + 1}$ 上的
内射阿贝尔层。下面将证明归纳假设适用于
$\mathcal{C}/K^{p + 1}$，从而
$H^q(K^{p + 1}, \mathcal{I}) = \colim H^q(K^{p + 1}, \mathcal{I}_\lambda) = 0$
对 $q < p_0 + 1$ 成立（因 $\mathcal{I}_\lambda$ 内射而消失）。由此
$$
H^{p_0 + 1}(\mathcal{C}, \mathcal{I}) =
H^{p_0 + 1}\left(\ldots \to
H^0(K^{p_0}, \mathcal{I}) \to
H^0(K^{p_0 + 1}, \mathcal{I}) \to
H^0(K^{p_0 + 2}, \mathcal{I}) \to \ldots\right)
$$
再次使用归纳假设，右边所示的复形是下列复形关于 $\Lambda$ 的余极限：
$$
\ldots \to
H^0(K^{p_0}, \mathcal{I}_\lambda) \to
H^0(K^{p_0 + 1}, \mathcal{I}_\lambda) \to
H^0(K^{p_0 + 2}, \mathcal{I}_\lambda) \to \ldots
$$
由于 $\mathcal{I}_\lambda$ 是内射阿贝尔层，这些复形正合
（例如可由谱序列得出）。阿贝尔群范畴中的滤过余极限正合，
故得到所需的消失性。

\medskip\noindent
还需证明，对所有 $n \geq 1$，归纳假设均适用于站点
$\mathcal{C}/K^n$。回忆
$\Sh(\mathcal{C}/K^n) = \Sh(\mathcal{C})/K^n$，见《站点》引理
\ref{sites-lemma-localize-topos-site}。因此可以在
$\Sh(\mathcal{C})/K^n$ 中工作。以
$S_n \subset \Ob(\Sh(\mathcal{C}/K^n)$ 记形如
$K' \to K^n$ 的对象所成的集合。逐一核对各项性质：
\begin{enumerate}
\item 由 (3) 及归纳，存在满射 $K' \to K^n$，其中 $K' \in S$。
于是 $(K' \to K^n) \to (K^n \to K^n)$ 是
$\Sh(\mathcal{C})/K^n$ 中的满射，而 $K^n \to K^n$ 是
$\Sh(\mathcal{C})/K^n$ 的终对象。因此 (1) 对 $S_n$ 成立；
\item $S_n$ 的性质 (2) 立即由 $S$ 的性质 (2) 得出；
\item 设 $a : K_1 \to K^n$ 与 $b : K_2 \to K^n$ 属于 $S_n$。
则 $(K_1 \to K^n) \times (K_2 \to K^n)$ 是
$\Sh(\mathcal{C})/K^n$ 的对象
$K_1 \times_{K^n} K_2 \to K^n$。子层
$K_1 \times_{K^n} K_2 \subset K_1 \times K_2$ 是
$a \circ \text{pr}_1$ 与 $b \circ \text{pr}_2$ 的等化子。
写成 $a = (a_1, \ldots, a_n)$ 与 $b = (b_1, \ldots, b_n)$。
选取满射 $K_3 \to K_1 \times K_2$，其中 $K_3 \in S$；
由 $\mathcal{C}$ 的假设 (3)，这是可行的。选取满射
$$
K_4
\longrightarrow
\text{等化子}(K_3 \to K_1 \times K_2 \xrightarrow{a_1, b_1} K)
$$
，其中 $K_4 \in S$。由 $\mathcal{C}$ 的假设 (4)，这是可行的。
选取满射
$$
K_5
\longrightarrow
\text{等化子}(K_4 \to K_1 \times K_2 \xrightarrow{a_2, b_2} K)
$$
，其中 $K_5 \in S$。这仍然可行。以此方式继续，直至得到满射
$$
K_{3 + n}
\longrightarrow
\text{等化子}(K_{2 + n} \to K_1 \times K_2 \xrightarrow{a_n, b_n} K)
$$
，其中 $K_{3 + n} \in S$。由构造，
$K_{3 + n} \to K_1 \times_{K^n} K_2$ 为满射。因此
$(K_{3 + n} \to K^n)$ 属于 $S_n$，并满射到乘积
$(K_1 \to K^n) \times (K_2 \to K^n)$。故 (3) 对 $S_n$ 成立；
\item $S_n$ 的性质 (4) 立即由 $S$ 的性质 (4) 得出；
\item $S_n$ 的性质 (5) 由 $S$ 的性质 (5) 得出。事实上，
$\Sh(\mathcal{C})/K^n$ 的对象 $\mathcal{F} \to K^n$
对应于 $\Sh(\mathcal{C}/K^n)$ 的拟紧对象，当且仅当
$\mathcal{F}$ 是 $\Sh(\mathcal{C})$ 的拟紧对象。
\end{enumerate}
引理证毕。
\end{proof}

\begin{remark}
\label{sites-cohomology-remark-colim-global}
设 $\mathcal{C}$ 为站点，$\mathcal{B} \subset \Ob(\mathcal{C})$
为子集。令 $S \subset \Ob(\Sh(\mathcal{C}))$ 为形如
$$
K = \coprod\nolimits_{i = 1, \ldots, n} h_{U_i}^\#
$$
的层 $K$ 所成的集合，其中 $U_1, \ldots, U_n \in \mathcal{B}$。
可以问：这个集合何时满足引理 \ref{sites-cohomology-lemma-colim-global} 的假设？
一个充分条件是
\begin{enumerate}
\item 对某个 $n \geq 0$ 及 $U_1, \ldots, U_n \in \mathcal{B}$，
映射 $\coprod_{i = 1, \ldots, n} h_{U_i}^\# \to *$ 为满射；
\item $U \in \mathcal{B}$ 的每个覆盖都能由形如
$\{U_i \to U\}_{i = 1, \ldots, n}$ 的覆盖加细，
其中 $U_i \in \mathcal{B}$；
\item 给定 $U,U' \in \mathcal{B}$，存在 $n \geq 0$、
$U_1, \ldots, U_n \in \mathcal{B}$ 以及态射
$U_i \to U$ 与 $U_i \to U'$，使得
$\coprod_{i = 1, \ldots, n} h_{U_i}^\# \to
h_U^\# \times h_{U'}^\#$ 为满射；
\item 给定 $\mathcal{C}$ 中的态射 $a,b : U \to U'$，其中
$U,U' \in \mathcal{B}$，存在 $U_1, \ldots, U_n \in \mathcal{B}$
以及等化 $a,b$ 的态射 $U_i \to U$，使得
$\coprod_{i = 1, \ldots, n} h_{U_i}^\# \to
\text{等化子}(h_a^\#, h_b^\# : h_U^\# \to h_{U'}^\#)$
为满射。
\end{enumerate}
省略详细核对，只指出上述 (2) 保证 $\mathcal{B}$ 的每个元素均拟紧，
从而由《站点》引理 \ref{sites-lemma-coproduct-quasi-compact}，
每个 $K \in S$ 也拟紧。
\end{remark}

\begin{lemma}
\label{sites-cohomology-lemma-higher-direct-image-colimit}
设 $f : \mathcal{C} \to \mathcal{D}$ 为对应于连续函子
$u : \mathcal{D} \to \mathcal{C}$ 的站点态射。
设 $(\mathcal{F}_i, \varphi_{ii'})$ 为 $\mathcal{C}$ 上的阿贝尔层系统，
并令 $\mathcal{F} = \colim \mathcal{F}_i$。设整数 $p \geq 0$。
以 $\mathcal{B}$ 记满足
$H^p(u(V), \mathcal{F}) = \colim H^p(u(V), \mathcal{F}_i)$.
的 $V \in \Ob(\mathcal{D})$ 所成的集合。若 $\mathcal{D}$ 的每个对象
都有一个由 $\mathcal{B}$ 中元素组成的覆盖，则
$R^pf_*\mathcal{F} = \colim R^pf_*\mathcal{F}_i$。
\end{lemma}

\begin{proof}
回忆，$R^pf_*\mathcal{F}$ 是预层 $\mathcal{G}$ 的层化，
其中 $\mathcal{G}$ 把 $V$ 送到 $H^p(u(V), \mathcal{F})$；
见引理 \ref{sites-cohomology-lemma-higher-direct-images}。类似地，
$R^pf_*\mathcal{F}_i$ 是把 $V$ 送到
$H^p(u(V), \mathcal{F}_i)$ 的预层 $\mathcal{G}_i$ 的层化。
回忆，层化是从层到预层的嵌入函子的左伴随，见《站点》第
\ref{sites-section-sheafification} 节。因此层化与余极限交换，
见《范畴》引理 \ref{categories-lemma-adjoint-exact}。
故只需证明预层态射（余极限在预层范畴中取）
$$
\colim \mathcal{G}_i \longrightarrow \mathcal{G}
$$
在层化上诱导同构。这由《站点》引理
\ref{sites-lemma-isomorphism-sheafifications} 及关于
$\mathcal{B}$ 的假设得出。
\end{proof}

\begin{lemma}
\label{sites-cohomology-lemma-colim-sites-injective}
设 $\mathcal{I}$ 为余滤过指标范畴，
$(\mathcal{C}_i, f_a)$ 为 $\mathcal{I}$ 上的站点逆系统，
如《站点》情形 \ref{sites-situation-inverse-limit-sites}。
按照《站点》引理 \ref{sites-lemma-colimit-sites} 与
\ref{sites-lemma-compute-pullback-to-limit}，令
$\mathcal{C} = \colim \mathcal{C}_i$。此外，假设给定
\begin{enumerate}
\item 对所有 $i \in \Ob(\mathcal{I})$，$\mathcal{C}_i$
上的阿贝尔层 $\mathcal{F}_i$；
\item 对 $a : j \to i$，$\mathcal{C}_j$ 上阿贝尔层的态射
$\varphi_a : f_a^{-1}\mathcal{F}_i \to \mathcal{F}_j$
\end{enumerate}
，并且每当 $c = a \circ b$ 时，
$\varphi_c = \varphi_b \circ f_b^{-1}\varphi_a$。
则存在系统的态射
$(\mathcal{F}_i, \varphi_a) \to (\mathcal{G}_i, \psi_a)$，
使得 $\mathcal{F}_i \to \mathcal{G}_i$ 为单射，且
$\mathcal{G}_i$ 是内射阿贝尔层。
\end{lemma}

\begin{proof}
对每个 $i$，选取单射 $\mathcal{F}_i \to \mathcal{A}_i$，
其中 $\mathcal{A}_i$ 是 $\mathcal{C}_i$ 上的内射阿贝尔层。
于是可以考虑映射族
$$
\gamma_i :
\mathcal{F}_i
\longrightarrow
\prod\nolimits_{b : k \to i} f_{b, *}\mathcal{A}_k = \mathcal{G}_i
$$
，其分量映射是态射
$f_b^{-1}\mathcal{F}_i \to \mathcal{F}_k \to \mathcal{A}_k$
的伴随态射。对 $\mathcal{I}$ 中的 $a : j \to i$，存在典范映射
$$
\psi_a : f_a^{-1}\mathcal{G}_i \to \mathcal{G}_j
$$
，其分量为典范映射
$f_b^{-1}f_{a \circ b, *}\mathcal{A}_k \to f_{b, *}\mathcal{A}_k$
，其中 $b : k \to j$。于是得到系统的单射
$(\gamma_i) : (\mathcal{F}_i, \varphi_a) \to (\mathcal{G}_i, \psi_a)$
。注意，$\mathcal{G}_i$ 是 $\mathcal{C}_i$ 上的内射阿贝尔群层；
见引理 \ref{sites-cohomology-lemma-pushforward-injective-flat} 与《同调代数》引理
\ref{homology-lemma-product-injectives}。构造完毕。
\end{proof}

\begin{lemma}
\label{sites-cohomology-lemma-colimit}
在引理 \ref{sites-cohomology-lemma-colim-sites-injective} 的情形下，令
$\mathcal{F} = \colim f_i^{-1}\mathcal{F}_i$。设
$i \in \Ob(\mathcal{I})$，$X_i \in \text{Ob}(\mathcal{C}_i)$。则
$$
\colim_{a : j \to i} H^p(u_a(X_i), \mathcal{F}_j) =
H^p(u_i(X_i), \mathcal{F})
$$
对所有 $p \geq 0$ 成立。
\end{lemma}

\begin{proof}
$p = 0$ 的情形是《站点》引理 \ref{sites-lemma-colimit}。

\medskip\noindent
按照引理 \ref{sites-cohomology-lemma-colim-sites-injective} 选取
$(\mathcal{F}_i, \varphi_a) \to (\mathcal{G}_i, \psi_a)$。
完全依照引理 \ref{sites-cohomology-lemma-colim-works-over-collection} 的证明论证可知，
只需证明对 $p > 0$ 有
$H^p(X, \colim f_i^{-1}\mathcal{G}_i) = 0$。

\medskip\noindent
令 $\mathcal{G} = \colim f_i^{-1}\mathcal{G}_i$。
为证明 $\mathcal{G}$ 在 $\mathcal{C}$ 每个对象上的上同调消失，
只需证明 $\mathcal{G}$ 关于 $\mathcal{C}$ 的任意覆盖
$\mathcal{U}$ 的 {\v C}ech 上同调为零
（引理 \ref{sites-cohomology-lemma-cech-vanish-collection}）。
覆盖 $\mathcal{U}$ 来自某个 $\mathcal{C}_i$ 的覆盖
$\mathcal{U}_i$。由 $p = 0$ 的情形，有
$$
\check{\mathcal{C}}^\bullet(\mathcal{U}, \mathcal{G}) =
\colim_{a : j \to i}
\check{\mathcal{C}}^\bullet(u_a(\mathcal{U}_i), \mathcal{G}_j)
$$
。由引理 \ref{sites-cohomology-lemma-injective-trivial-cech}，右边是无环复形的滤过余极限，
故在正次数无环。见《代数》引理
\ref{algebra-lemma-directed-colimit-exact}。
\end{proof}















\section{平坦分辨}
\label{sites-cohomology-section-flat}

\noindent
本节在环化站点和环化拓扑斯的情形下，重新论证
《上同调》第 \ref{cohomology-section-flat} 节中的论证。

\begin{lemma}
\label{sites-cohomology-lemma-derived-tor-exact}
设 $(\mathcal{C}, \mathcal{O})$ 为环化站点。
设 $\mathcal{G}^\bullet$ 为一个 $\mathcal{O}$-模复形。
函子
$$
K(\textit{Mod}(\mathcal{O}))
\longrightarrow
K(\textit{Mod}(\mathcal{O})),
\quad
\mathcal{F}^\bullet \longmapsto
\text{Tot}(\mathcal{G}^\bullet \otimes_\mathcal{O} \mathcal{F}^\bullet)
$$
以及
$$
K(\textit{Mod}(\mathcal{O}))
\longrightarrow
K(\textit{Mod}(\mathcal{O})),
\quad
\mathcal{F}^\bullet \longmapsto
\text{Tot}(\mathcal{F}^\bullet \otimes_\mathcal{O} \mathcal{G}^\bullet)
$$
都是三角范畴之间的正合函子。
\end{lemma}

\begin{proof}
这由《导出范畴》注记
\ref{derived-remark-double-complex-as-tensor-product-of}.
\end{proof}

\begin{definition}
\label{sites-cohomology-definition-K-flat}
设 $(\mathcal{C}, \mathcal{O})$ 为环化站点。
若对于每个 $\mathcal{O}$-模的无环复形 $\mathcal{F}^\bullet$，复形
$$
\text{Tot}(\mathcal{F}^\bullet \otimes_\mathcal{O} \mathcal{K}^\bullet)
$$
都是无环的，则称 $\mathcal{O}$-模复形 $\mathcal{K}^\bullet$
为 {\it K-平坦}。
\end{definition}

\begin{lemma}
\label{sites-cohomology-lemma-K-flat-quasi-isomorphism}
设 $(\mathcal{C}, \mathcal{O})$ 为环化站点，且 $\mathcal{K}^\bullet$
为 K-平坦复形。则函子
$$
K(\textit{Mod}(\mathcal{O}))
\longrightarrow
K(\textit{Mod}(\mathcal{O})), \quad
\mathcal{F}^\bullet
\longmapsto
\text{Tot}(\mathcal{F}^\bullet \otimes_\mathcal{O} \mathcal{K}^\bullet)
$$
把拟同构映为拟同构。
\end{lemma}

\begin{proof}
这由引理
\ref{sites-cohomology-lemma-derived-tor-exact}
以及拟同构可由其锥为无环复形来刻画这一事实得到。
\end{proof}

\begin{lemma}
\label{sites-cohomology-lemma-restriction-K-flat}
设 $(\mathcal{C}, \mathcal{O})$ 为环化站点，设 $U$ 为 $\mathcal{C}$ 的对象。
若 $\mathcal{K}^\bullet$ 是 K-平坦的 $\mathcal{O}$-模复形，则
$\mathcal{K}^\bullet|_U$ 是 K-平坦的 $\mathcal{O}_U$-模复形。
\end{lemma}

\begin{proof}
设 $\mathcal{G}^\bullet$ 为 $\mathcal{O}_U$-模的正合复形。
由于 $j_{U!}$ 是正合的
(《站点上的模》引理 \ref{sites-modules-lemma-extension-by-zero-exact})
且 $\mathcal{K}^\bullet$ 是 K-平坦的 $\mathcal{O}$-模复形，故复形
$$
j_{U!}(\text{Tot}(\mathcal{G}^\bullet \otimes_{\mathcal{O}_U}
\mathcal{K}^\bullet|_U)) =
\text{Tot}(j_{U!}\mathcal{G}^\bullet \otimes_\mathcal{O} \mathcal{K}^\bullet)
$$
是正合的。这里的等式来自
《站点上的模》引理 \ref{sites-modules-lemma-j-shriek-and-tensor}
以及 $j_{U!}$ 与直和交换（因为它是左伴随）。由
《站点上的模》引理 \ref{sites-modules-lemma-j-shriek-reflects-exactness} 可知 $j_{U!}$ 反映正合性，故结论成立。
\end{proof}

\begin{lemma}
\label{sites-cohomology-lemma-tensor-product-K-flat}
设 $(\mathcal{C}, \mathcal{O})$ 为环化站点。
若 $\mathcal{K}^\bullet$、$\mathcal{L}^\bullet$ 是 $\mathcal{O}$-模的
K-平坦复形，则
$\text{Tot}(\mathcal{K}^\bullet \otimes_\mathcal{O} \mathcal{L}^\bullet)$
是 K-平坦的 $\mathcal{O}$-模复形。
\end{lemma}

\begin{proof}
这由同构
$$
\text{Tot}(\mathcal{M}^\bullet \otimes_\mathcal{O}
\text{Tot}(\mathcal{K}^\bullet \otimes_\mathcal{O} \mathcal{L}^\bullet))
=
\text{Tot}(\text{Tot}(\mathcal{M}^\bullet \otimes_\mathcal{O}
\mathcal{K}^\bullet) \otimes_\mathcal{O} \mathcal{L}^\bullet)
$$
以及定义得到。
\end{proof}

\begin{lemma}
\label{sites-cohomology-lemma-K-flat-two-out-of-three}
设 $(\mathcal{C}, \mathcal{O})$ 为环化站点。
设 $(\mathcal{K}_1^\bullet, \mathcal{K}_2^\bullet, \mathcal{K}_3^\bullet)$
是 $K(\textit{Mod}(\mathcal{O}))$ 中的有区别三角。
若 $\mathcal{K}_i^\bullet$ 中任意两个是 K-平坦的，则第三个也是。
\end{lemma}

\begin{proof}
这由引理 \ref{sites-cohomology-lemma-derived-tor-exact}
以及以下事实得到：在
$K(\textit{Mod}(\mathcal{O}))$
中的有区别三角中，若任意两个复形无环，则第三个也无环。
\end{proof}

\begin{lemma}
\label{sites-cohomology-lemma-K-flat-two-out-of-three-ses}
设 $(\mathcal{C}, \mathcal{O})$ 为环化站点。设
$0 \to \mathcal{K}_1^\bullet \to \mathcal{K}_2^\bullet \to
\mathcal{K}_3^\bullet \to 0$ 是复形的短正合序列，且
$\mathcal{K}_3^\bullet$ 的各项是平坦的 $\mathcal{O}$-模。
若 $\mathcal{K}_i^\bullet$ 中任意两个是 K-平坦的，则第三个也是。
\end{lemma}

\begin{proof}
由《站点上的模》引理 \ref{sites-modules-lemma-flat-tor-zero}，
对每个复形 $\mathcal{L}^\bullet$，我们得到短正合序列
$$
0 \to
\text{Tot}(\mathcal{L}^\bullet \otimes_\mathcal{O} \mathcal{K}_1^\bullet) \to
\text{Tot}(\mathcal{L}^\bullet \otimes_\mathcal{O} \mathcal{K}_2^\bullet) \to
\text{Tot}(\mathcal{L}^\bullet \otimes_\mathcal{O} \mathcal{K}_3^\bullet) \to 0
$$
为复形的短正合序列。因此由上同调层的长正合列和 K-平坦复形的定义，
引理得证。
\end{proof}

\begin{lemma}
\label{sites-cohomology-lemma-bounded-flat-K-flat}
设 $(\mathcal{C}, \mathcal{O})$ 为环化站点。平坦 $\mathcal{O}$-模的
上有界复形是 K-平坦的。
\end{lemma}

\begin{proof}
设 $\mathcal{K}^\bullet$ 为平坦 $\mathcal{O}$-模的上有界复形，
设 $\mathcal{L}^\bullet$ 为 $\mathcal{O}$-模的无环复形。注意
$\mathcal{L}^\bullet = \colim_m \tau_{\leq m}\mathcal{L}^\bullet$
其中余极限逐项取。因此也有
$$
\text{Tot}(\mathcal{K}^\bullet \otimes_\mathcal{O} \mathcal{L}^\bullet)
=
\colim_m \text{Tot}(
\mathcal{K}^\bullet \otimes_\mathcal{O} \tau_{\leq m}\mathcal{L}^\bullet)
$$
逐项成立。因此要证明左边的复形无环，只需证明右边的每个复形无环。
由于 $\tau_{\leq m}\mathcal{L}^\bullet$ 无环，这归结为
$\mathcal{L}^\bullet$ 上有界的情形。
此时《同调》的引理
《同调》引理 \ref{homology-lemma-first-quadrant-ss} 的谱序列为
$$
{}'E_1^{p, q} = H^p(\mathcal{L}^\bullet \otimes_R \mathcal{K}^q)
$$
中的谱序列含有上式，而上式为零，因为 $\mathcal{K}^q$ 平坦且
$\mathcal{L}^\bullet$ 无环。故结论成立。
\end{proof}

\begin{lemma}
\label{sites-cohomology-lemma-colimit-K-flat}
设 $(\mathcal{C}, \mathcal{O})$ 为环化站点。
设 $\mathcal{K}_1^\bullet \to \mathcal{K}_2^\bullet \to \ldots$
为 K-平坦复形组成的系统。
则 $\colim_i \mathcal{K}_i^\bullet$ 是 K-平坦的。
\end{lemma}

\begin{proof}
由于余极限逐项取，显然有
$$
\colim_i \text{Tot}(
\mathcal{F}^\bullet \otimes_\mathcal{O} \mathcal{K}_i^\bullet)
=
\text{Tot}(\mathcal{F}^\bullet \otimes_\mathcal{O}
\colim_i \mathcal{K}_i^\bullet)
$$
因此由滤过余极限正合这一事实，引理得证。
\end{proof}

\begin{lemma}
\label{sites-cohomology-lemma-resolution-by-direct-sums-extensions-by-zero}
设 $(\mathcal{C}, \mathcal{O})$ 为环化站点。对任意 $\mathcal{O}$-模复形
$\mathcal{G}^\bullet$，存在 $\mathcal{O}$-模复形的交换图
$$
\xymatrix{
\mathcal{K}_1^\bullet \ar[d] \ar[r] &
\mathcal{K}_2^\bullet \ar[d] \ar[r] & \ldots \\
\tau_{\leq 1}\mathcal{G}^\bullet \ar[r] &
\tau_{\leq 2}\mathcal{G}^\bullet \ar[r] & \ldots
}
$$
具有如下性质：(1) 垂直箭头是拟同构且逐项满射；
(2) 每个 $\mathcal{K}_n^\bullet$ 是上有界复形，其各项是形如
$j_{U!}\mathcal{O}_U$ 的 $\mathcal{O}$-模的直和；
(3) 映射 $\mathcal{K}_n^\bullet \to \mathcal{K}_{n + 1}^\bullet$ 是逐项可分裂的单射，
余核是形如 $j_{U!}\mathcal{O}_U$ 的 $\mathcal{O}$-模的直和。此外，映射
$\colim \mathcal{K}_n^\bullet \to \mathcal{G}^\bullet$ 是拟同构。
\end{lemma}

\begin{proof}
该图及性质 (1)、(2)、(3) 立即由
《站点上的模》引理 \ref{sites-modules-lemma-module-quotient-flat}
以及《导出范畴》引理 \ref{derived-lemma-special-direct-system} 得到。
诱导映射 $\colim \mathcal{K}_n^\bullet \to \mathcal{G}^\bullet$
是拟同构，因为滤过余极限正合。
\end{proof}

\begin{lemma}
\label{sites-cohomology-lemma-K-flat-resolution}
设 $(\mathcal{C}, \mathcal{O})$ 为环化站点。对任意复形
$\mathcal{G}^\bullet$，存在一个 K-平坦复形 $\mathcal{K}^\bullet$，其各项为平坦的
$\mathcal{O}$-模，并存在逐项满射的拟同构
$\mathcal{K}^\bullet \to \mathcal{G}^\bullet$。
\end{lemma}

\begin{proof}
选取如引理 \ref{sites-cohomology-lemma-resolution-by-direct-sums-extensions-by-zero} 中的图。
每个复形 $\mathcal{K}_n^\bullet$ 都是平坦模的上有界复形，见
《站点上的模》引理 \ref{sites-modules-lemma-j-shriek-flat}。
故由引理 \ref{sites-cohomology-lemma-bounded-flat-K-flat}，$\mathcal{K}_n^\bullet$ 是 K-平坦的。
于是由引理 \ref{sites-cohomology-lemma-colimit-K-flat}，$\colim \mathcal{K}_n^\bullet$ 是 K-平坦的。
由构造，诱导映射 $\colim \mathcal{K}_n^\bullet \to \mathcal{G}^\bullet$
是拟同构且逐项满射。引理
\ref{sites-cohomology-lemma-resolution-by-direct-sums-extensions-by-zero} 的性质 (3) 表明，
$\colim \mathcal{K}_n^m$ 是平坦模的直和，因而平坦，这就证明了最后的断言。
\end{proof}

\begin{lemma}
\label{sites-cohomology-lemma-derived-tor-quasi-isomorphism-other-side}
设 $(\mathcal{C}, \mathcal{O})$ 为环化站点。设
$\alpha : \mathcal{P}^\bullet \to \mathcal{Q}^\bullet$ 为 $\mathcal{O}$-模的
K-平坦复形之间的拟同构。对每个 $\mathcal{O}$-模复形 $\mathcal{F}^\bullet$，
诱导映射
$$
\text{Tot}(\text{id}_{\mathcal{F}^\bullet} \otimes \alpha) :
\text{Tot}(\mathcal{F}^\bullet \otimes_\mathcal{O} \mathcal{P}^\bullet)
\longrightarrow
\text{Tot}(\mathcal{F}^\bullet \otimes_\mathcal{O} \mathcal{Q}^\bullet)
$$
是拟同构。
\end{lemma}

\begin{proof}
选取拟同构 $\mathcal{K}^\bullet \to \mathcal{F}^\bullet$，其中
$\mathcal{K}^\bullet$ 是 K-平坦复形，见
引理 \ref{sites-cohomology-lemma-K-flat-resolution}。
考虑交换图
$$
\xymatrix{
\text{Tot}(\mathcal{K}^\bullet
\otimes_\mathcal{O} \mathcal{P}^\bullet) \ar[r] \ar[d] &
\text{Tot}(\mathcal{K}^\bullet
\otimes_\mathcal{O} \mathcal{Q}^\bullet) \ar[d] \\
\text{Tot}(\mathcal{F}^\bullet
\otimes_\mathcal{O} \mathcal{P}^\bullet) \ar[r] &
\text{Tot}(\mathcal{F}^\bullet
\otimes_\mathcal{O} \mathcal{Q}^\bullet)
}
$$
由引理 \ref{sites-cohomology-lemma-K-flat-quasi-isomorphism}，垂直箭头和上方水平箭头
都是拟同构，故结论成立。
\end{proof}

\noindent
设 $(\mathcal{C}, \mathcal{O})$ 为环化站点。
设 $\mathcal{F}^\bullet$ 为 $D(\mathcal{O})$ 的对象。
选取 K-平坦分辨 $\mathcal{K}^\bullet \to \mathcal{F}^\bullet$，见引理
\ref{sites-cohomology-lemma-K-flat-resolution}。由引理
\ref{sites-cohomology-lemma-derived-tor-exact}，得到三角范畴之间的正合函子
$$
K(\mathcal{O})
\longrightarrow
K(\mathcal{O}),
\quad
\mathcal{G}^\bullet
\longmapsto
\text{Tot}(\mathcal{G}^\bullet \otimes_\mathcal{O} \mathcal{K}^\bullet)
$$
由引理 \ref{sites-cohomology-lemma-K-flat-quasi-isomorphism}，此函子诱导出函子
$D(\mathcal{O}) \to D(\mathcal{O})$，这是因为
$D(\mathcal{O})$ 是 $K(\mathcal{O})$ 在拟同构处的局部化。
由于 $\mathcal{F}^\bullet$ 的 K-平坦分辨范畴是余滤的，且有引理
\ref{sites-cohomology-lemma-derived-tor-quasi-isomorphism-other-side}，所得函子
（在同构意义下）不依赖于 $\mathcal{K}^\bullet$ 的选择。

\begin{definition}
\label{sites-cohomology-definition-derived-tor}
设 $(\mathcal{C}, \mathcal{O})$ 为环化站点。
设 $\mathcal{F}^\bullet$ 为 $D(\mathcal{O})$ 的对象。
上面描述的 {\it 导出张量积}
$$
- \otimes_\mathcal{O}^{\mathbf{L}} \mathcal{F}^\bullet :
D(\mathcal{O})
\longrightarrow
D(\mathcal{O})
$$
是上面描述的三角范畴之间的正合函子。
\end{definition}

\noindent
从我们的显式构造立即可见，对于 $D(\mathcal{O})$ 中的
$\mathcal{G}^\bullet$ 和 $\mathcal{F}^\bullet$，有规范同构
$$
\mathcal{F}^\bullet \otimes_\mathcal{O}^{\mathbf{L}} \mathcal{G}^\bullet
\cong
\mathcal{G}^\bullet \otimes_\mathcal{O}^{\mathbf{L}} \mathcal{F}^\bullet
$$
这里使用《代数进阶》第 \ref{more-algebra-section-sign-rules} 节给出的符号规则。
因此，当我们写
$\mathcal{F}^\bullet \otimes_\mathcal{O}^{\mathbf{L}} \mathcal{G}^\bullet$
时，通常不区分用哪个变量来定义导出张量积。

\begin{definition}
\label{sites-cohomology-definition-tor}
设 $(\mathcal{C}, \mathcal{O})$ 为环化站点。
设 $\mathcal{F}$、$\mathcal{G}$ 为 $\mathcal{O}$-模。
$\mathcal{F}$ 与 $\mathcal{G}$ 的 {\it Tor} 由公式
$$
\text{Tor}_p^\mathcal{O}(\mathcal{F}, \mathcal{G}) =
H^{-p}(\mathcal{F} \otimes_\mathcal{O}^\mathbf{L} \mathcal{G})
$$
定义，其中导出张量积如上所述。
\end{definition}

\noindent
这个定义意味着，对每个 $\mathcal{O}$-模的短正合序列
$0 \to \mathcal{F}_1 \to \mathcal{F}_2 \to \mathcal{F}_3 \to 0$
都有上同调长正合列
$$
\xymatrix{
\mathcal{F}_1 \otimes_\mathcal{O} \mathcal{G} \ar[r] &
\mathcal{F}_2 \otimes_\mathcal{O} \mathcal{G} \ar[r] &
\mathcal{F}_3 \otimes_\mathcal{O} \mathcal{G} \ar[r] & 0 \\
\text{Tor}_1^\mathcal{O}(\mathcal{F}_1, \mathcal{G}) \ar[r] &
\text{Tor}_1^\mathcal{O}(\mathcal{F}_2, \mathcal{G}) \ar[r] &
\text{Tor}_1^\mathcal{O}(\mathcal{F}_3, \mathcal{G}) \ar[ull]
}
$$
对每个 $\mathcal{O}$-模 $\mathcal{G}$ 成立。我们将其称为
该情形所关联的 $\text{Tor}$ 长正合列。

\begin{lemma}
\label{sites-cohomology-lemma-flat-tor-zero}
设 $(\mathcal{C}, \mathcal{O})$ 为环化站点。
设 $\mathcal{F}$ 为 $\mathcal{O}$-模。以下条件等价：
\begin{enumerate}
\item $\mathcal{F}$ 是平坦的 $\mathcal{O}$-模；
\item $\text{Tor}_1^\mathcal{O}(\mathcal{F}, \mathcal{G}) = 0$
对每个 $\mathcal{O}$-模 $\mathcal{G}$ 成立。
\end{enumerate}
\end{lemma}

\begin{proof}
若 $\mathcal{F}$ 平坦，则 $\mathcal{F} \otimes_\mathcal{O} -$
是正合函子且卫星函子消失。反之，设 (2) 成立。若
$\mathcal{G} \to \mathcal{H}$ 为余核为 $\mathcal{Q}$ 的单射，则
$\text{Tor}$ 长正合列表明下列映射的核
$\mathcal{F} \otimes_\mathcal{O} \mathcal{G} \to
\mathcal{F} \otimes_\mathcal{O} \mathcal{H}$
是下列对象的商：
$\text{Tor}_1^\mathcal{O}(\mathcal{F}, \mathcal{Q})$
按假设该对象为零。因此 $\mathcal{F}$ 平坦。
\end{proof}

\begin{lemma}
\label{sites-cohomology-lemma-K-flat-flat-acyclic}
设 $(\mathcal{C}, \mathcal{O})$ 为环化站点。设 $\mathcal{K}^\bullet$
为各项平坦的 K-平坦无环复形。则
$\mathcal{F} = \Ker(\mathcal{K}^n \to \mathcal{K}^{n + 1})$
是平坦的 $\mathcal{O}$-模。
\end{lemma}

\begin{proof}
注意
$$
\ldots \to \mathcal{K}^{n - 2} \to \mathcal{K}^{n - 1} \to
\mathcal{F} \to 0
$$
是我们的模 $\mathcal{F}$ 的平坦分辨。由于平坦模的上有界复形是 K-平坦的
（引理 \ref{sites-cohomology-lemma-bounded-flat-K-flat}），我们可以用此分辨计算
$\text{Tor}_i(\mathcal{F}, \mathcal{G})$，其中 $\mathcal{G}$ 为任意
$\mathcal{O}$-模。一方面
$\mathcal{K}^\bullet \otimes_\mathcal{O}^\mathbf{L} \mathcal{G}$
在 $D(\mathcal{O})$ 中为零，因为 $\mathcal{K}^\bullet$ 无环；另一方面它由
$\mathcal{K}^\bullet \otimes_\mathcal{O} \mathcal{G}$.
因此
$$
\mathcal{K}^{n - 3} \otimes_\mathcal{O} \mathcal{G} \to
\mathcal{K}^{n - 2} \otimes_\mathcal{O} \mathcal{G} \to
\mathcal{K}^{n - 1} \otimes_\mathcal{O} \mathcal{G}
$$
是正合的。所以 $\text{Tor}_1(\mathcal{F}, \mathcal{G}) = 0$，
由引理 \ref{sites-cohomology-lemma-flat-tor-zero} 得出结论。
\end{proof}

\begin{lemma}
\label{sites-cohomology-lemma-factor-through-K-flat}
设 $(\mathcal{C}, \mathcal{O})$ 为环化站点。
设 $a : \mathcal{K}^\bullet \to \mathcal{L}^\bullet$ 为 $\mathcal{O}$-模复形的映射。
若 $\mathcal{K}^\bullet$ K-平坦，则存在复形 $\mathcal{N}^\bullet$ 以及复形映射
$b : \mathcal{K}^\bullet \to \mathcal{N}^\bullet$
和 $c : \mathcal{N}^\bullet \to \mathcal{L}^\bullet$，使得
\begin{enumerate}
\item $\mathcal{N}^\bullet$ 是 K-平坦的；
\item $c$ 是拟同构；
\item $a$ 与 $c \circ b$ 同伦。
\end{enumerate}
若 $\mathcal{K}^\bullet$ 的各项平坦，则可以选取
$\mathcal{N}^\bullet$、$b$ 和 $c$
使 $\mathcal{N}^\bullet$ 也具有同样性质。
\end{lemma}

\begin{proof}
我们将使用同伦范畴 $K(\textit{Mod}(\mathcal{O}))$ 是三角范畴这一事实，见《导出范畴》命题
\ref{derived-proposition-homotopy-category-triangulated}.
选取有区别三角
$\mathcal{K}^\bullet \to \mathcal{L}^\bullet \to
\mathcal{C}^\bullet \to \mathcal{K}^\bullet[1]$.
选取拟同构 $\mathcal{M}^\bullet \to \mathcal{C}^\bullet$，其中
$\mathcal{M}^\bullet$ K-平坦且各项平坦，见引理
\ref{sites-cohomology-lemma-K-flat-resolution}。由三角范畴公理，我们可以将复合映射
$\mathcal{M}^\bullet \to \mathcal{C}^\bullet \to \mathcal{K}^\bullet[1]$
嵌入有区别三角
$\mathcal{K}^\bullet \to \mathcal{N}^\bullet \to
\mathcal{M}^\bullet \to \mathcal{K}^\bullet[1]$.
由引理 \ref{sites-cohomology-lemma-K-flat-two-out-of-three}，$\mathcal{N}^\bullet$ 是 K-平坦的。
再次使用三角范畴公理，可以选取映射 $\mathcal{N}^\bullet \to \mathcal{L}^\bullet$，
使其组成下列有区别三角的态射
$$
\xymatrix{
\mathcal{K}^\bullet \ar[r] \ar[d] &
\mathcal{N}^\bullet \ar[r] \ar[d] &
\mathcal{M}^\bullet \ar[r] \ar[d] &
\mathcal{K}^\bullet[1] \ar[d] \\
\mathcal{K}^\bullet \ar[r] &
\mathcal{L}^\bullet \ar[r] &
\mathcal{C}^\bullet \ar[r] &
\mathcal{K}^\bullet[1]
}
$$
由于三个箭头中有两个是拟同构，故由这些有区别三角所关联的上同调长正合列，
第三个箭头 $\mathcal{N}^\bullet \to \mathcal{L}^\bullet$ 也是拟同构
（也可以将此图映入 $D(\mathcal{O})$，并使用《导出范畴》引理
\ref{derived-lemma-third-isomorphism-triangle}）。这就完成了 (1)、(2)、(3) 的证明。
为证明最后的断言，可以选取 $\mathcal{N}^\bullet$ 使
$\mathcal{N}^n \cong \mathcal{M}^n \oplus \mathcal{K}^n$，见
《导出范畴》引理 \ref{derived-lemma-improve-distinguished-triangle-homotopy}。
若 $\mathcal{K}^\bullet$ 的各项平坦，则由此得到所需的平坦性。
\end{proof}












\section{导出拉回}
\label{sites-cohomology-section-derived-pullback}

\noindent
设
$f : (\Sh(\mathcal{C}), \mathcal{O}) \to
(\Sh(\mathcal{C}'), \mathcal{O}')$
为环化拓扑斯之间的态射。我们可以使用 K-平坦分辨定义导出拉回函子
$$
Lf^* : D(\mathcal{O}') \to D(\mathcal{O})
$$

\begin{lemma}
\label{sites-cohomology-lemma-pullback-K-flat}
设 $f : (\Sh(\mathcal{C}'), \mathcal{O}') \to (\Sh(\mathcal{C}), \mathcal{O})$
为环化拓扑斯之间的态射。设 $\mathcal{K}^\bullet$ 为各项平坦的
$\mathcal{O}$-模 K-平坦复形。则 $f^*\mathcal{K}^\bullet$ 是各项平坦的
$\mathcal{O}'$-模 K-平坦复形。
\end{lemma}

\begin{proof}
由《站点上的模》引理 \ref{sites-modules-lemma-pullback-flat}，
$f^*\mathcal{K}^n$ 的各项是平坦的 $\mathcal{O}'$-模。选取图
$$
\xymatrix{
\mathcal{K}_1^\bullet \ar[d] \ar[r] &
\mathcal{K}_2^\bullet \ar[d] \ar[r] & \ldots \\
\tau_{\leq 1}\mathcal{K}^\bullet \ar[r] &
\tau_{\leq 2}\mathcal{K}^\bullet \ar[r] & \ldots
}
$$
如引理 \ref{sites-cohomology-lemma-resolution-by-direct-sums-extensions-by-zero} 中所示。
下文使用该引理的所有性质而不再另行说明。每个 $\mathcal{K}_n^\bullet$ 是平坦模的
上有界复形，见《站点上的模》引理 \ref{sites-modules-lemma-j-shriek-flat}。
考虑定义 $\mathcal{M}^\bullet$ 的复形短正合序列
$$
0 \to \mathcal{M}^\bullet \to
\colim \mathcal{K}_n^\bullet \to
\mathcal{K}^\bullet \to 0
$$
由引理 \ref{sites-cohomology-lemma-bounded-flat-K-flat} 和 \ref{sites-cohomology-lemma-colimit-K-flat}，
复形 $\colim \mathcal{K}_n^\bullet$ 是 K-平坦的；由《站点上的模》引理
\ref{sites-modules-lemma-colimits-flat}，其各项平坦。由《站点上的模》引理
\ref{sites-modules-lemma-flat-ses}，$\mathcal{M}^\bullet$ 的各项平坦；由引理
\ref{sites-cohomology-lemma-K-flat-two-out-of-three-ses}，$\mathcal{M}^\bullet$ 是 K-平坦的；
再由上同调长正合列，$\mathcal{M}^\bullet$ 无环（因为第二个箭头是拟同构）。拉回
$f^*(\colim \mathcal{K}_n^\bullet) = \colim f^*\mathcal{K}_n^\bullet$
是各项平坦的 $\mathcal{O}'$-模的下有界复形的余极限，因而 K-平坦
（使用上面的同一组引理）。拉回上述短正合序列
$$
0 \to f^*\mathcal{M}^\bullet \to
f^*(\colim \mathcal{K}_n^\bullet) \to
f^*\mathcal{K}^\bullet \to 0
$$
由《站点上的模》引理 \ref{sites-modules-lemma-pullback-ses} 是复形的短正合序列。
因此由引理 \ref{sites-cohomology-lemma-K-flat-two-out-of-three-ses}，只需证明
$f^*\mathcal{M}^\bullet$ K-平坦。这归结为下一段讨论的情形。

\medskip\noindent
设 $\mathcal{K}^\bullet$ 除 K-平坦外还是无环的，且各项平坦。
由引理 \ref{sites-cohomology-lemma-K-flat-flat-acyclic}，$\tau_{\leq n}\mathcal{K}^\bullet$
的所有项都是平坦的 $\mathcal{O}$-模。选取如上的图，并在此图中使用上面证明的
所有性质。记 $\mathcal{M}_n^\bullet$ 为复形映射
$\mathcal{K}_n^\bullet \to \tau_{\leq n}\mathcal{K}^\bullet$，
从而有复形的短正合序列
$$
0 \to \mathcal{M}_n^\bullet \to \mathcal{K}_n^\bullet \to
\tau_{\leq n}\mathcal{K}^\bullet \to 0
$$
由《站点上的模》引理 \ref{sites-modules-lemma-flat-ses}，复形
$\mathcal{M}_n^\bullet$ 的各项平坦。因此在此情形下
$\mathcal{M} = \colim \mathcal{M}_n^\bullet$ 是平坦模的下有界复形的滤过余极限。
于是 $f^*\mathcal{M}^\bullet$ K-平坦（论证同上），结论成立。
\end{proof}

\begin{lemma}
\label{sites-cohomology-lemma-derived-base-change}
设 $f : (\Sh(\mathcal{C}), \mathcal{O}) \to (\Sh(\mathcal{C}'), \mathcal{O}')$
为环化拓扑斯之间的态射。存在三角范畴之间的正合函子
$$
Lf^* : D(\mathcal{O}') \longrightarrow D(\mathcal{O})
$$
使得
$Lf^*\mathcal{K}^\bullet = f^*\mathcal{K}^\bullet$ 对任意各项平坦的
K-平坦复形 $\mathcal{K}^\bullet$ 成立，特别地，对任意平坦
$\mathcal{O}'$-模的上有界复形也成立。
\end{lemma}

\begin{proof}
为此使用《导出范畴》第 \ref{derived-section-derived-functors} 节发展的
一般理论。置 $\mathcal{D} = K(\mathcal{O}')$、$\mathcal{D}' = D(\mathcal{O})$。
记 $F : \mathcal{D} \to \mathcal{D}'$ 为由规则
$F(\mathcal{G}^\bullet) = f^*\mathcal{G}^\bullet$。
令 $S$ 为下列范畴中的拟同构集合：
$\mathcal{D} = K(\mathcal{O}')$.
这给出《导出范畴》情形 \ref{derived-situation-derived-functor} 中的情形，
因而适用《导出范畴》定义 \ref{derived-definition-right-derived-functor-defined}。
我们断言 $LF$ 处处有定义。这由《导出范畴》引理
\ref{derived-lemma-find-existence-computes} 得到，其中
$\mathcal{P} \subset \Ob(\mathcal{D})$ 是各项平坦的 K-平坦复形
$\mathcal{K}^\bullet$ 所成的集合。具体地，(1) 由引理
\ref{sites-cohomology-lemma-K-flat-resolution} 得出；为证明 (2)，需证明对于拟同构
$\mathcal{K}_1^\bullet  \to \mathcal{K}_2^\bullet$ 之间，且两者属于
$\mathcal{P}$，映射 $f^*\mathcal{K}_1^\bullet \to f^*\mathcal{K}_2^\bullet$
是拟同构。为此将其写成
$$
f^{-1}\mathcal{K}_1^\bullet \otimes_{f^{-1}\mathcal{O}'} \mathcal{O}
\longrightarrow
f^{-1}\mathcal{K}_2^\bullet \otimes_{f^{-1}\mathcal{O}'} \mathcal{O}
$$
函子 $f^{-1}$ 正合，故映射
$f^{-1}\mathcal{K}_1^\bullet \to f^{-1}\mathcal{K}_2^\bullet$ 是拟同构。复形
$f^{-1}\mathcal{K}_1^\bullet$ 与 $f^{-1}\mathcal{K}_2^\bullet$
由引理 \ref{sites-cohomology-lemma-pullback-K-flat} 是 $f^{-1}\mathcal{O}'$-模的 K-平坦复形，
因为可以考虑环化拓扑斯态射
$(\Sh(\mathcal{C}), f^{-1}\mathcal{O}') \to
(\Sh(\mathcal{C}'), \mathcal{O}')$。因此引理
\ref{sites-cohomology-lemma-derived-tor-quasi-isomorphism-other-side}
保证显示的映射是拟同构。于是得到导出函子
$$
LF :
D(\mathcal{O}') = S^{-1}\mathcal{D}
\longrightarrow
\mathcal{D}' = D(\mathcal{O})
$$
见《导出范畴》方程 (\ref{derived-equation-everywhere})。
最后，《导出范畴》引理 \ref{derived-lemma-find-existence-computes} 还保证
$LF(\mathcal{K}^\bullet) = F(\mathcal{K}^\bullet) = f^*\mathcal{K}^\bullet$
当 $\mathcal{K}^\bullet$ 属于 $\mathcal{P}$ 时成立。
注意到平坦模的上有界复形由引理 \ref{sites-cohomology-lemma-bounded-flat-K-flat}
属于 $\mathcal{P}$，即可完成证明。
\end{proof}

\begin{lemma}
\label{sites-cohomology-lemma-derived-pullback-composition}
考虑环化拓扑斯态射
$f : (\Sh(\mathcal{C}), \mathcal{O}_\mathcal{C}) \to
(\Sh(\mathcal{D}), \mathcal{O}_\mathcal{D})$
以及
$g : (\Sh(\mathcal{D}), \mathcal{O}_\mathcal{D}) \to
(\Sh(\mathcal{E}), \mathcal{O}_\mathcal{E})$.
则作为函子 $D(\mathcal{O}_\mathcal{E}) \to D(\mathcal{O}_\mathcal{C})$，
$Lf^* \circ Lg^* = L(g \circ f)^*$。
\end{lemma}

\begin{proof}
设 $E$ 为 $D(\mathcal{O}_\mathcal{E})$ 的对象。可用各项平坦的 K-平坦复形
$\mathcal{K}^\bullet$ 表示 $E$，见引理 \ref{sites-cohomology-lemma-K-flat-resolution}。
按构造，$Lg^*E$ 由 $g^*\mathcal{K}^\bullet$ 计算，见引理
\ref{sites-cohomology-lemma-derived-base-change}。由引理 \ref{sites-cohomology-lemma-pullback-K-flat}，复形
$g^*\mathcal{K}^\bullet$ K-平坦且各项平坦。因此 $Lf^*Lg^*E$ 由
$f^*g^*\mathcal{K}^\bullet$ 表示。另一方面 $L(g \circ f)^*E$ 也由
$(g \circ f)^*\mathcal{K}^\bullet = f^*g^*\mathcal{K}^\bullet$ 表示，故结论成立。
\end{proof}

\begin{lemma}
\label{sites-cohomology-lemma-pullback-tensor-product}
设 $f : (\Sh(\mathcal{C}), \mathcal{O}) \to (\Sh(\mathcal{D}), \mathcal{O}')$
为环化拓扑斯之间的态射。存在规范的双函子同构
$$
Lf^*(
\mathcal{F}^\bullet \otimes_{\mathcal{O}'}^{\mathbf{L}} \mathcal{G}^\bullet
) =
Lf^*\mathcal{F}^\bullet 
\otimes_{\mathcal{O}}^{\mathbf{L}}
Lf^*\mathcal{G}^\bullet 
$$
对 $\mathcal{F}^\bullet, \mathcal{G}^\bullet \in \Ob(D(\mathcal{O}'))$ 成立。
\end{lemma}

\begin{proof}
由引理 \ref{sites-cohomology-lemma-derived-base-change} 中导出拉回的构造以及引理
\ref{sites-cohomology-lemma-K-flat-resolution} 中分辨的存在性，可以将 $\mathcal{F}^\bullet$ 和
$\mathcal{G}^\bullet$ 换成各项平坦且 K-平坦的 $\mathcal{O}'$-模复形。在此情形下，
$\mathcal{F}^\bullet \otimes_{\mathcal{O}'}^{\mathbf{L}} \mathcal{G}^\bullet$
就是双复形
$\mathcal{F}^\bullet \otimes_{\mathcal{O}'} \mathcal{G}^\bullet$.
复形
$\text{Tot}(\mathcal{F}^\bullet \otimes_{\mathcal{O}'} \mathcal{G}^\bullet)$
由引理 \ref{sites-cohomology-lemma-tensor-product-K-flat} 和《站点上的模》引理
\ref{sites-modules-lemma-tensor-flats} 可知 K-平坦且各项平坦。
因此本引理的同构来自同构
$$
\text{Tot}(f^*\mathcal{F}^\bullet \otimes_{\mathcal{O}}
f^*\mathcal{G}^\bullet)
\longrightarrow
f^*\text{Tot}(\mathcal{F}^\bullet \otimes_{\mathcal{O}'} \mathcal{G}^\bullet)
$$
其组成部分是《站点上的模》引理
$f^*\mathcal{F}^p \otimes_{\mathcal{O}} f^*\mathcal{G}^q \to
f^*(\mathcal{F}^p \otimes_{\mathcal{O}'} \mathcal{G}^q)$，即
\ref{sites-modules-lemma-tensor-product-pullback} 的同构
\end{proof}

\begin{lemma}
\label{sites-cohomology-lemma-variant-derived-pullback}
设 $f : (\Sh(\mathcal{C}), \mathcal{O}) \to (\Sh(\mathcal{C}'), \mathcal{O}')$
为环化拓扑斯之间的态射。存在规范的双函子同构
$$
\mathcal{F}^\bullet
\otimes_\mathcal{O}^{\mathbf{L}}
Lf^*\mathcal{G}^\bullet
=
\mathcal{F}^\bullet 
\otimes_{f^{-1}\mathcal{O}'}^{\mathbf{L}}
f^{-1}\mathcal{G}^\bullet 
$$
其中 $\mathcal{F}^\bullet$ 属于 $D(\mathcal{O})$，
$\mathcal{G}^\bullet$ 属于 $D(\mathcal{O}')$。
\end{lemma}

\begin{proof}
设 $\mathcal{F}$ 为 $\mathcal{O}$-模，设 $\mathcal{G}$ 为 $\mathcal{O}'$-模。则
$\mathcal{F} \otimes_{\mathcal{O}} f^*\mathcal{G} =
\mathcal{F} \otimes_{f^{-1}\mathcal{O}'} f^{-1}\mathcal{G}$
因为
$f^*\mathcal{G} =
\mathcal{O} \otimes_{f^{-1}\mathcal{O}'} f^{-1}\mathcal{G}$.
由此及定义即可得到引理。
\end{proof}











\begin{lemma}
\label{sites-cohomology-lemma-check-K-flat-stalks}
设 $(\mathcal{C}, \mathcal{O})$ 为环化站点。
设 $\mathcal{K}^\bullet$ 为 $\mathcal{O}$-模复形。
\begin{enumerate}
\item 若 $\mathcal{K}^\bullet$ K-平坦，则对站点 $\mathcal{C}$ 的每个点 $p$，
$\mathcal{O}_p$-模复形 $\mathcal{K}_p^\bullet$ 按《代数进阶》定义
\ref{more-algebra-definition-K-flat} 的意义是 K-平坦的；
\item 若 $\mathcal{C}$ 有足够多的点，则逆命题成立。
\end{enumerate}
\end{lemma}

\begin{proof}
证明 (2)。若 $\mathcal{C}$ 有足够多的点，且对 $\mathcal{C}$ 的所有点 $p$，
$\mathcal{K}_p^\bullet$ 都是 K-平坦的，则 $\mathcal{K}^\bullet$ 是 K-平坦的，
因为 $\otimes$ 和直和与取茎交换，并且可在茎上检验正合性，见
《站点上的模》引理 \ref{sites-modules-lemma-check-exactness-stalks}。

\medskip\noindent
证明 (1)。设 $\mathcal{K}^\bullet$ K-平坦。选取拟同构
$a : \mathcal{L}^\bullet \to \mathcal{K}^\bullet$，使 $\mathcal{L}^\bullet$ K-平坦且各项平坦，见
引理 \ref{sites-cohomology-lemma-K-flat-resolution}。$\mathcal{L}^\bullet$ 的任意拉回都是 K-平坦的，见
引理 \ref{sites-cohomology-lemma-pullback-K-flat}。特别地，茎
$\mathcal{L}_p^\bullet$ 是 $\mathcal{O}_p$-模的 K-平坦复形。
于是 $a$ 的锥 $C(a)$ 是 K-平坦的（引理 \ref{sites-cohomology-lemma-K-flat-two-out-of-three}）
无环 $\mathcal{O}$-模复形，只需证明 $C(a)$ 的茎是 K-平坦的
（由《代数进阶》引理 \ref{more-algebra-lemma-K-flat-two-out-of-three}）。
因此可设 $\mathcal{K}^\bullet$ K-平坦且无环。

\medskip\noindent
设 $\mathcal{K}^\bullet$ 无环且 K-平坦。在继续之前，按《站点》引理
\ref{sites-lemma-topos-good-site} 将站点 $\mathcal{C}$ 换成另一个站点，
以确保 $\mathcal{C}$ 有所有有限极限。这意味着 $p$ 的邻域范畴是滤过的
（《站点》引理 \ref{sites-lemma-neighbourhoods-directed}），并且定义层的茎的余极限
是滤过的。设 $M$ 为有限表示的 $\mathcal{O}_p$-模。只需证明
$\mathcal{K}^\bullet \otimes_{\mathcal{O}_p} M$ 无环，见《代数进阶》引理
\ref{more-algebra-lemma-universally-acyclic-K-flat}。
由于 $\mathcal{O}_p$ 是 $\mathcal{O}(U)$（其中 $U$ 遍历 $p$ 的邻域）的滤过余极限，
可找到 $p$ 的邻域 $(U,x)$ 以及有限表示的 $\mathcal{O}(U)$-模 $M'$，其到
$\mathcal{O}_p$ 的基变换为 $M$，见《代数》引理
\ref{algebra-lemma-colimit-category-fp-modules}。由引理
\ref{sites-cohomology-lemma-restriction-K-flat}，可将 $\mathcal{C},\mathcal{O},\mathcal{K}^\bullet$
替换为 $\mathcal{C}/U,\mathcal{O}_U,\mathcal{K}^\bullet|_U$。
于是可设存在 $\mathcal{O}$-模 $\mathcal{F}$ 使 $M \cong \mathcal{F}_p$。
由于 $\mathcal{K}^\bullet$ K-平坦且无环，复形
$\mathcal{K}^\bullet \otimes_\mathcal{O} \mathcal{F}$ 无环
（按定义它计算导出张量积）。取茎是正合函子，故
$\mathcal{K}^\bullet \otimes_{\mathcal{O}_p} M$ 无环，正是所需。
\end{proof}

\begin{lemma}
\label{sites-cohomology-lemma-pullback-K-flat-points}
设 $f : (\Sh(\mathcal{C}), \mathcal{O}) \to (\Sh(\mathcal{C}'), \mathcal{O}')$
为环化拓扑斯之间的态射。若 $\mathcal{C}$ 有足够多的点，则
$\mathcal{O}'$-模 K-平坦复形的拉回是 $\mathcal{O}$-模的 K-平坦复形。
\end{lemma}

\begin{proof}
这由引理 \ref{sites-cohomology-lemma-check-K-flat-stalks}、《站点上的模》引理
\ref{sites-modules-lemma-pullback-stalk} 以及《代数进阶》引理
\ref{more-algebra-lemma-base-change-K-flat} 得到。
\end{proof}

\begin{lemma}
\label{sites-cohomology-lemma-tensor-pull-compatibility}
设 $f : (\Sh(\mathcal{C}), \mathcal{O}_\mathcal{C}) \to
(\Sh(\mathcal{D}), \mathcal{O}_\mathcal{D})$ 为环化拓扑斯之间的态射。
设 $\mathcal{K}^\bullet$ 和 $\mathcal{M}^\bullet$ 为 $\mathcal{O}_\mathcal{D}$-模复形。
下图
$$
\xymatrix{
Lf^*(\mathcal{K}^\bullet
\otimes_{\mathcal{O}_\mathcal{D}}^\mathbf{L}
\mathcal{M}^\bullet) \ar[r] \ar[d] &
Lf^*\text{Tot}(\mathcal{K}^\bullet
\otimes_{\mathcal{O}_\mathcal{D}}
\mathcal{M}^\bullet) \ar[d] \\
Lf^*\mathcal{K}^\bullet \otimes_{\mathcal{O}_\mathcal{C}}^\mathbf{L}
Lf^*\mathcal{M}^\bullet \ar[d] &
f^*\text{Tot}(\mathcal{K}^\bullet
\otimes_{\mathcal{O}_\mathcal{D}}
\mathcal{M}^\bullet) \ar[d] \\
f^*\mathcal{K}^\bullet \otimes_{\mathcal{O}_\mathcal{C}}^\mathbf{L}
f^*\mathcal{M}^\bullet \ar[r] &
\text{Tot}(f^*\mathcal{K}^\bullet \otimes_{\mathcal{O}_\mathcal{C}}
f^*\mathcal{M}^\bullet)
}
$$
交换。
\end{lemma}

\begin{proof}
我们将使用各项平坦的 K-平坦分辨的存在性
（引理 \ref{sites-cohomology-lemma-K-flat-resolution}），使用导出拉回由此类复形计算这一事实
（引理 \ref{sites-cohomology-lemma-derived-base-change}），以及拉回保持这些性质这一事实
（引理 \ref{sites-cohomology-lemma-pullback-K-flat}）。若选取此类分辨
$\mathcal{P}^\bullet \to \mathcal{K}^\bullet$ 和
$\mathcal{Q}^\bullet \to \mathcal{M}^\bullet$，则有
$$
\xymatrix{
Lf^*\text{Tot}(\mathcal{P}^\bullet
\otimes_{\mathcal{O}_\mathcal{D}}
\mathcal{Q}^\bullet) \ar[r] \ar[d] &
Lf^*\text{Tot}(\mathcal{K}^\bullet
\otimes_{\mathcal{O}_\mathcal{D}}
\mathcal{M}^\bullet) \ar[d] \\
f^*\text{Tot}(\mathcal{P}^\bullet
\otimes_{\mathcal{O}_\mathcal{D}}
\mathcal{Q}^\bullet) \ar[d] \ar[r] &
f^*\text{Tot}(\mathcal{K}^\bullet
\otimes_{\mathcal{O}_\mathcal{D}}
\mathcal{M}^\bullet) \ar[d] \\
\text{Tot}(f^*\mathcal{P}^\bullet \otimes_{\mathcal{O}_\mathcal{C}}
f^*\mathcal{Q}^\bullet) \ar[r] &
\text{Tot}(f^*\mathcal{K}^\bullet \otimes_{\mathcal{O}_\mathcal{C}}
f^*\mathcal{M}^\bullet)
}
$$
交换。不过，由于我们选择了 $\mathcal{P}^\bullet$、$\mathcal{Q}^\bullet$
以及引理 \ref{sites-cohomology-lemma-tensor-product-K-flat}，现在图的左端正是该图的左端。
\end{proof}









\section{无界复形的上同调}
\label{sites-cohomology-section-unbounded}

\noindent
设 $(\mathcal{C}, \mathcal{O})$ 为环化站点。
范畴 $\textit{Mod}(\mathcal{O})$ 是 Grothendieck 阿贝尔范畴：它有所有余极限，
滤过余极限正合，并且有生成元，即
$$
\bigoplus\nolimits_{U \in \Ob(\mathcal{C})} j_{U!}\mathcal{O}_U,
$$
见《站点上的模》第 \ref{sites-modules-section-kernels} 节以及引理
\ref{sites-modules-lemma-j-shriek-flat} 和
\ref{sites-modules-lemma-module-quotient-flat}。
由《内射对象》定理
\ref{injectives-theorem-K-injective-embedding-grothendieck}
对每个 $\mathcal{O}$-模复形 $\mathcal{F}^\bullet$，存在到 K-内射
$\mathcal{O}$-模复形 $\mathcal{I}^\bullet$ 的内射拟同构
$\mathcal{F}^\bullet \to \mathcal{I}^\bullet$，且此嵌入可选为关于
$\mathcal{F}^\bullet$ 的函子性的。由
导出范畴，引理 \ref{derived-lemma-enough-K-injectives-implies}
可知
\begin{enumerate}
\item 任意到三角范畴 $\mathcal{D}$ 的正合函子
$F : K(\textit{Mod}(\mathcal{O})) \to \mathcal{D}$ 都有右导出函子
$RF : D(\mathcal{O}) \to \mathcal{D}$；
\item 对任意到阿贝尔范畴 $\mathcal{A}$ 的加性函子
$F : \textit{Mod}(\mathcal{O}) \to \mathcal{A}$，考虑由 $F$ 诱导的正合函子
$F : K(\textit{Mod}(\mathcal{O})) \to K(\mathcal{A})$，则得到右导出函子
$RF : D(\mathcal{O}) \to D(\mathcal{A})$。
\end{enumerate}
按构造，$RF(\mathcal{F}^\bullet) = F(\mathcal{I}^\bullet)$，
其中 $\mathcal{F}^\bullet \to \mathcal{I}^\bullet$ 如上所述。

\medskip\noindent
下面给出上述构造的一些例子：
\begin{enumerate}
\item 函子 $\Gamma(\mathcal{C}, -) : \textit{Mod}(\mathcal{O}) \to
\text{Mod}_{\Gamma(\mathcal{C}, \mathcal{O})}$ 导出
$$
R\Gamma(\mathcal{C}, -) :
D(\mathcal{O})
\longrightarrow
D(\Gamma(\mathcal{C}, \mathcal{O}))
$$
我们用记号
$H^i(\mathcal{C}, K) = H^i(R\Gamma(\mathcal{C}, K))$ 表示上同调。
\item 对 $\mathcal{C}$ 的对象 $U$，考虑函子
$\Gamma(U, -) : \textit{Mod}(\mathcal{O}) \to
\text{Mod}_{\Gamma(U, \mathcal{O})}$。它导出
$$
R\Gamma(U, -) : D(\mathcal{O}) \to D(\Gamma(U, \mathcal{O}))
$$
我们用记号
$H^i(U, K) = H^i(R\Gamma(U, K))$ 表示上同调。
\item 对环化拓扑斯态射
$f : (\Sh(\mathcal{C}), \mathcal{O}) \to (\Sh(\mathcal{D}), \mathcal{O}')$
考虑函子
$f_* : \textit{Mod}(\mathcal{O}) \to \textit{Mod}(\mathcal{O}')$
它给出全导出直接像
$$
Rf_* : D(\mathcal{O}) \longrightarrow D(\mathcal{O}')
$$
在无界导出范畴上。
\end{enumerate}

\begin{lemma}
\label{sites-cohomology-lemma-adjoint}
设 $f : (\Sh(\mathcal{C}), \mathcal{O}) \to (\Sh(\mathcal{D}), \mathcal{O}')$
为环化拓扑斯之间的态射。上面定义的函子 $Rf_*$ 与引理
\ref{sites-cohomology-lemma-derived-base-change} 中定义的函子 $Lf^*$ 互为伴随：
$$
\Hom_{D(\mathcal{O})}(Lf^*\mathcal{G}^\bullet, \mathcal{F}^\bullet)
=
\Hom_{D(\mathcal{O}')}(\mathcal{G}^\bullet, Rf_*\mathcal{F}^\bullet)
$$
关于 $\mathcal{F}^\bullet \in \Ob(D(\mathcal{O}))$ 和
$\mathcal{G}^\bullet \in \Ob(D(\mathcal{O}'))$ 双函子性地成立。
\end{lemma}

\begin{proof}
这形式地由 $Rf_*$ 和 $Lf^*$ 存在这一事实得到，见《导出范畴》引理
\ref{derived-lemma-derived-adjoint-functors}。
\end{proof}

\begin{lemma}
\label{sites-cohomology-lemma-derived-pushforward-composition}
设
$f : (\Sh(\mathcal{C}), \mathcal{O}_\mathcal{C}) \to
(\Sh(\mathcal{D}), \mathcal{O}_\mathcal{D})$
以及
$g : (\Sh(\mathcal{D}), \mathcal{O}_\mathcal{D}) \to
(\Sh(\mathcal{E}), \mathcal{O}_\mathcal{E})$
为环化拓扑斯态射。则作为函子
$D(\mathcal{O}_\mathcal{C}) \to D(\mathcal{O}_\mathcal{E})$，有
$Rg_* \circ Rf_* = R(g \circ f)_*$。
\end{lemma}

\begin{proof}
由引理 \ref{sites-cohomology-lemma-adjoint}，$Rg_* \circ Rf_*$ 是 $Lf^* \circ Lg^*$ 的伴随。
由引理 \ref{sites-cohomology-lemma-derived-pullback-composition}，有
$Lf^* \circ Lg^* = L(g \circ f)^*$，因此由伴随函子的唯一性，
$Rg_* \circ Rf_* = R(g \circ f)_*$。
\end{proof}

\begin{remark}
\label{sites-cohomology-remark-base-change}
无界导出函子 $Lf^*$ 和 $Rf_*$ 的构造使我们能够在完全一般性下构造基变换映射。
具体地，设
$$
\xymatrix{
(\Sh(\mathcal{C}'), \mathcal{O}_{\mathcal{C}'})
\ar[r]_{g'} \ar[d]_{f'} &
(\Sh(\mathcal{C}), \mathcal{O}_\mathcal{C}) \ar[d]^f \\
(\Sh(\mathcal{D}'), \mathcal{O}_{\mathcal{D}'})
\ar[r]^g &
(\Sh(\mathcal{D}), \mathcal{O}_\mathcal{D})
}
$$
是环化拓扑斯的交换图。设 $K$ 为 $D(\mathcal{O}_\mathcal{C})$ 的对象。
则存在规范的基变换映射
$$
Lg^*Rf_*K \longrightarrow R(f')_*L(g')^*K
$$
它位于 $D(\mathcal{O}_{\mathcal{D}'})$ 中。具体地，该映射是下列映射的伴随：
$L(f')^*Lg^*Rf_*K \to L(g')^*K$.
由于 $L(f')^* \circ Lg^* = L(g')^* \circ Lf^*$，可见这与映射
$L(g')^*Lf^*Rf_*K \to L(g')^*K$ 相同；我们可以将其取为伴随映射
$Lf^*Rf_*K \to K$ 在 $L(g')^*$ 下的像。
\end{remark}

\begin{remark}
\label{sites-cohomology-remark-compose-base-change}
考虑交换图
$$
\xymatrix{
(\Sh(\mathcal{B}'), \mathcal{O}_{\mathcal{B}'})
\ar[r]_k \ar[d]_{f'} &
(\Sh(\mathcal{B}), \mathcal{O}_\mathcal{B}) \ar[d]^f \\
(\Sh(\mathcal{C}'), \mathcal{O}_{\mathcal{C}'})
\ar[r]^l \ar[d]_{g'} &
(\Sh(\mathcal{C}), \mathcal{O}_\mathcal{C}) \ar[d]^g \\
(\Sh(\mathcal{D}'), \mathcal{O}_{\mathcal{D}'})
\ar[r]^m &
(\Sh(\mathcal{D}), \mathcal{O}_\mathcal{D}) \\
}
$$
它是环化拓扑斯的图。则《基变换》注记
\ref{sites-cohomology-remark-base-change} 中两个方形的基变换映射复合，给出外部矩形的基变换映射。
更精确地，复合
\begin{align*}
Lm^* \circ R(g \circ f)_*
& =
Lm^* \circ Rg_* \circ Rf_* \\
& \to Rg'_* \circ Ll^* \circ Rf_* \\
& \to Rg'_* \circ Rf'_* \circ Lk^* \\
& = R(g' \circ f')_* \circ Lk^*
\end{align*}
是该矩形的基变换映射。验证从略。
\end{remark}

\begin{remark}
\label{sites-cohomology-remark-compose-base-change-horizontal}
考虑交换图
$$
\xymatrix{
(\Sh(\mathcal{C}''), \mathcal{O}_{\mathcal{C}''})
\ar[r]_{g'} \ar[d]_{f''} &
(\Sh(\mathcal{C}'), \mathcal{O}_{\mathcal{C}'})
\ar[r]_g \ar[d]_{f'} &
(\Sh(\mathcal{C}), \mathcal{O}_\mathcal{C}) \ar[d]^f \\
(\Sh(\mathcal{D}''), \mathcal{O}_{\mathcal{D}''})
\ar[r]^{h'} &
(\Sh(\mathcal{D}'), \mathcal{O}_{\mathcal{D}'})
\ar[r]^h &
(\Sh(\mathcal{D}), \mathcal{O}_\mathcal{D})
}
$$
它是环化拓扑斯的图。则《基变换》注记
\ref{sites-cohomology-remark-base-change} 中两个方形的基变换映射复合，给出外部矩形的基变换映射。
更精确地，复合
\begin{align*}
L(h \circ h')^* \circ Rf_*
& =
L(h')^* \circ Lh^* \circ Rf_* \\
& \to L(h')^* \circ Rf'_* \circ Lg^* \\
& \to Rf''_* \circ L(g')^* \circ Lg^* \\
& = Rf''_* \circ L(g \circ g')^*
\end{align*}
是该矩形的基变换映射。验证从略。
\end{remark}



\begin{lemma}
\label{sites-cohomology-lemma-adjoints-push-pull-compatibility}
设 $f : (\Sh(\mathcal{C}), \mathcal{O}_\mathcal{C}) \to
(\Sh(\mathcal{D}), \mathcal{O}_\mathcal{D})$ 为环化拓扑斯之间的态射。
设 $\mathcal{K}^\bullet$ 为 $\mathcal{O}_\mathcal{C}$-模复形。下图
$$
\xymatrix{
Lf^*f_*\mathcal{K}^\bullet \ar[r] \ar[d] &
f^*f_*\mathcal{K}^\bullet \ar[d] \\
Lf^*Rf_*\mathcal{K}^\bullet \ar[r] &
\mathcal{K}^\bullet
}
$$
来自复形上的 $Lf^* \to f^*$、复形上的 $f_* \to Rf_*$ 以及伴随
$Lf^* \circ Rf_* \to \text{id}$，并且在 $D(\mathcal{O}_\mathcal{C})$ 中交换。
\end{lemma}

\begin{proof}
我们将使用 K-平坦分辨和 K-内射分辨的存在性，见引理
\ref{sites-cohomology-lemma-K-flat-resolution}、\ref{sites-cohomology-lemma-derived-base-change}、
\ref{sites-cohomology-lemma-pullback-K-flat} 以及上面的讨论。
选取拟同构
$\mathcal{K}^\bullet \to \mathcal{I}^\bullet$，其中 $\mathcal{I}^\bullet$
是 $\mathcal{O}_\mathcal{C}$-模的 K-内射复形。选取拟同构
$\mathcal{Q}^\bullet \to f_*\mathcal{I}^\bullet$
其中 $\mathcal{Q}^\bullet$ 是各项平坦的 $\mathcal{O}_\mathcal{D}$-模 K-平坦复形。
可选取各项平坦的 $\mathcal{O}_\mathcal{D}$-模 K-平坦复形
$\mathcal{P}^\bullet$ 以及复形态射图
$$
\xymatrix{
\mathcal{P}^\bullet \ar[r] \ar[d] &
f_*\mathcal{K}^\bullet \ar[d] \\
\mathcal{Q}^\bullet \ar[r] & f_*\mathcal{I}^\bullet
}
$$
在同伦意义下交换，且上方水平箭头是拟同构。具体地，首先可对某个复形
$\mathcal{P}^\bullet$ 选取这样的图，因为拟同构在复形的同伦范畴中形成乘法系统；
然后可用各项平坦的 K-平坦复形对 $\mathcal{P}^\bullet$ 作分辨。
取拉回得到复形态射图
$$
\xymatrix{
f^*\mathcal{P}^\bullet \ar[r] \ar[d] &
f^*f_*\mathcal{K}^\bullet \ar[d] \ar[r] &
\mathcal{K}^\bullet \ar[d] \\
f^*\mathcal{Q}^\bullet \ar[r] &
f^*f_*\mathcal{I}^\bullet \ar[r] &
\mathcal{I}^\bullet
}
$$
在同伦意义下交换。外部矩形证明了引理中的断言。
\end{proof}



\begin{remark}
\label{sites-cohomology-remark-cup-product}
设 $f : (\Sh(\mathcal{C}), \mathcal{O}_\mathcal{C}) \to
(\Sh(\mathcal{D}), \mathcal{O}_\mathcal{D})$ 为环化拓扑斯之间的态射。
$Lf^*$ 与 $Rf_*$ 的伴随性允许我们构造相对杯积
$$
Rf_*K \otimes_{\mathcal{O}_\mathcal{D}}^\mathbf{L} Rf_*L
\longrightarrow
Rf_*(K \otimes_{\mathcal{O}_\mathcal{C}}^\mathbf{L} L)
$$
它位于 $D(\mathcal{O}_\mathcal{D})$ 中，对所有
$K,L \in D(\mathcal{O}_\mathcal{C})$ 成立。具体地，该映射是下列映射的伴随：
$Lf^*(Rf_*K \otimes_{\mathcal{O}_\mathcal{D}}^\mathbf{L} Rf_*L) \to
K \otimes_{\mathcal{O}_\mathcal{C}}^\mathbf{L} L$，其中可取同构
$Lf^*(Rf_*K \otimes_{\mathcal{O}_\mathcal{D}}^\mathbf{L} Rf_*L) =
Lf^*Rf_*K \otimes_{\mathcal{O}_\mathcal{C}}^\mathbf{L} Lf^*Rf_*L$
(引理 \ref{sites-cohomology-lemma-pullback-tensor-product})
与映射
$Lf^*Rf_*K \otimes_{\mathcal{O}_\mathcal{C}}^\mathbf{L} Lf^*Rf_*L
\to K \otimes_{\mathcal{O}_\mathcal{C}}^\mathbf{L} L$
后者来自余单位 $Lf^* \circ Rf_* \to \text{id}$。
\end{remark}

\begin{lemma}
\label{sites-cohomology-lemma-torsion}
设 $\mathcal{C}$ 为站点。设 $\mathcal{A} \subset \textit{Ab}(\mathcal{C})$
为由挠阿贝尔层组成的 Serre 子范畴。
则函子 $D(\mathcal{A}) \to D_\mathcal{A}(\mathcal{C})$ 是等价。
\end{lemma}

\begin{proof}
关键观察是，内射阿贝尔层 $\mathcal{I}$ 是可除的。具体地，若
$s \in \mathcal{I}(U)$ 是局部截面，则将 $s$ 解释为映射
$s : j_{U!}\mathbf{Z} \to \mathcal{I}$，并将内射对象的定义性质应用于
层的单射 $n : j_{U!}\mathbf{Z} \to j_{U!}\mathbf{Z}$，
从而得到存在 $s' \in \mathcal{I}(U)$ 使 $ns' = s$。

\medskip\noindent
对层 $\mathcal{F}$，记其挠子层为 $\mathcal{F}_{tor}$。
我们断言，若 $\mathcal{I}^\bullet$ 是内射阿贝尔层复形且其上同调层为挠的，则
$$
\mathcal{I}^\bullet_{tor} \to \mathcal{I}^\bullet
$$
是拟同构。具体地，由于 $\mathbf{Q}$ 在 $\mathbf{Z}$ 上平坦，有
$$
H^p(\mathcal{I}^\bullet) \otimes_\mathbf{Z} \mathbf{Q} =
H^p(\mathcal{I}^\bullet \otimes_\mathbf{Z} \mathbf{Q})
$$
它为零，因为上同调层是挠的。由上面证明的可除性，
$\mathcal{I}^\bullet \to \mathcal{I}^\bullet \otimes_\mathbf{Z} \mathbf{Q}$
是满射且核为 $\mathcal{I}^\bullet_{tor}$。由所得短正合序列所关联的
上同调层长正合列，断言成立。

\medskip\noindent
为证明引理，构造右伴随
$T : D(\mathcal{C}) \to D(\mathcal{A})$。具体地，对 $D(\mathcal{C})$ 中的 $K$，
可用上同调层为内射的 K-内射复形 $\mathcal{I}^\bullet$ 表示，见《内射对象》定理
\ref{injectives-theorem-K-injective-embedding-grothendieck}.
令 $T(K) = \mathcal{I}^\bullet_{tor}$，换言之，$T$ 是取挠子的右导出函子。
函子 $T$ 是下列函子的右伴随：
$i : D(\mathcal{A}) \to D_\mathcal{A}(\mathcal{C})$.
这立即由如下观察得到：若 $\mathcal{F}^\bullet$ 是挠层复形，则
$$
\Hom_{K(\mathcal{A})}(\mathcal{F}^\bullet, I^\bullet_{tor}) =
\Hom_{K(\textit{Ab}(\mathcal{C}))}(\mathcal{F}^\bullet, I^\bullet)
$$
特别地，$\mathcal{I}^\bullet_{tor}$ 是 $\mathcal{A}$ 的 K-内射复形。
细节从略；如有疑问，也可由更一般的《导出范畴》引理
\ref{derived-lemma-derived-adjoint-functors} 得到。
上面的断言给出，对 $D(\mathcal{A})$ 中的 $L$ 有 $L = T(i(L))$，
且当 $K \in D_\mathcal{A}(\mathcal{C})$ 时有 $i(T(K)) = K$。
由《范畴》引理 \ref{categories-lemma-adjoint-fully-faithful} 得到结论。
\end{proof}




\section{K-内射复形的一些性质}
\label{sites-cohomology-section-properties-K-injective}

\noindent
设 $(\mathcal{C}, \mathcal{O})$ 为环化站点。设 $U$ 为 $\mathcal{C}$ 的对象。记
$j : (\Sh(\mathcal{C}/U), \mathcal{O}_U) \to (\Sh(\mathcal{C}), \mathcal{O})$
为相应的局部化态射。拉回函子 $j^*$ 是正合的，因为它就是限制函子。
因此，导出拉回 $Lj^*$ 对任意复形都只需通过限制该复形来计算。我们常将相应函子
$$
D(\mathcal{O}) \to D(\mathcal{O}_U),
\quad
E \mapsto j^*E = E|_U
$$
类似地，延零函子
$j_! : \textit{Mod}(\mathcal{O}_U) \to \textit{Mod}(\mathcal{O})$（见《站点上的模》定义
\ref{sites-modules-definition-localize-ringed-site}）是正合函子
（《站点上的模》引理 \ref{sites-modules-lemma-extension-by-zero-exact}）。
因此它诱导函子
$$
j_! : D(\mathcal{O}_U) \to D(\mathcal{O}), \quad
F \mapsto j_!F
$$
只需将 $j_!$ 作用于表示对象 $F$ 的任意复形即可。

\begin{lemma}
\label{sites-cohomology-lemma-restrict-K-injective-to-open}
设 $(\mathcal{C}, \mathcal{O})$ 为环化站点，设 $U$ 为 $\mathcal{C}$ 的对象。
$\mathcal{O}$-模 K-内射复形限制到 $\mathcal{C}/U$ 后，是
$\mathcal{O}_U$-模的 K-内射复形。
\end{lemma}

\begin{proof}
立即由《导出范畴》引理
\ref{derived-lemma-adjoint-preserve-K-injectives} 以及限制函子有正合左伴随
$j_!$ 这一事实得到。见上面的讨论。
\end{proof}

\begin{lemma}
\label{sites-cohomology-lemma-unbounded-cohomology-of-open}
设 $(\mathcal{C}, \mathcal{O})$ 为环化站点。设 $U \in \Ob(\mathcal{C})$。
对 $D(\mathcal{O})$ 中的 $K$，有
$H^p(U, K) = H^p(\mathcal{C}/U, K|_{\mathcal{C}/U})$.
\end{lemma}

\begin{proof}
设 $\mathcal{I}^\bullet$ 为表示 $K$ 的 $\mathcal{O}$-模 K-内射复形。则
$$
H^q(U, K) = H^q(\Gamma(U, \mathcal{I}^\bullet)) =
H^q(\Gamma(\mathcal{C}/U, \mathcal{I}^\bullet|_{\mathcal{C}/U}))
$$
由上同调的构造成立。由引理 \ref{sites-cohomology-lemma-restrict-K-injective-to-open}，复形
$\mathcal{I}^\bullet|_{\mathcal{C}/U}$ 是表示 $K|_{\mathcal{C}/U}$ 的 K-内射复形，
引理得证。
\end{proof}

\begin{lemma}
\label{sites-cohomology-lemma-sheafification-cohomology}
设 $(\mathcal{C}, \mathcal{O})$ 为环化站点。设 $K$ 为 $D(\mathcal{O})$ 的对象。
下列预层的层化
$$
U \mapsto H^q(U, K) = H^q(\mathcal{C}/U, K|_{\mathcal{C}/U})
$$
就是 $K$ 的第 $q$ 个上同调层 $H^q(K)$。
\end{lemma}

\begin{proof}
等式 $H^q(U, K) = H^q(\mathcal{C}/U, K|_{\mathcal{C}/U})$
由引理 \ref{sites-cohomology-lemma-unbounded-cohomology-of-open} 成立。选取表示 $K$ 的 K-内射复形
$\mathcal{I}^\bullet$。则
$$
H^q(U, K) =
\frac{\Ker(\mathcal{I}^q(U) \to \mathcal{I}^{q + 1}(U))}
{\Im(\mathcal{I}^{q - 1}(U) \to \mathcal{I}^q(U))}.
$$
由上同调的构造成立。由于
$H^q(K) = \Ker(\mathcal{I}^q \to \mathcal{I}^{q + 1})/
\Im(\mathcal{I}^{q - 1} \to \mathcal{I}^q)$，结论显然。
\end{proof}

\begin{lemma}
\label{sites-cohomology-lemma-restrict-direct-image-open}
设 $f : (\mathcal{C}, \mathcal{O}_\mathcal{C}) \to
(\mathcal{D}, \mathcal{O}_\mathcal{D})$ 为对应于连续函子 $u : \mathcal{D} \to \mathcal{C}$
的环化站点态射。给定 $V \in \mathcal{D}$，置 $U = u(V)$，并记
$g : (\mathcal{C}/U, \mathcal{O}_U) \to (\mathcal{D}/V, \mathcal{O}_V)$
为诱导的环化站点态射（《站点上的模》引理
\ref{sites-modules-lemma-localize-morphism-ringed-sites}）。
则对 $D(\mathcal{O}_\mathcal{C})$ 中的 $E$，有
$(Rf_*E)|_{\mathcal{D}/V} = Rg_*(E|_{\mathcal{C}/U})$。
\end{lemma}

\begin{proof}
用 $\mathcal{O}_\mathcal{C}$-模 K-内射复形 $\mathcal{I}^\bullet$ 表示 $E$。则
$Rf_*(E) = f_*\mathcal{I}^\bullet$
且由引理 \ref{sites-cohomology-lemma-restrict-K-injective-to-open}，
$Rg_*(E|_{\mathcal{C}/U}) = g_*(\mathcal{I}^\bullet|_{\mathcal{C}/U})$。
由于显然有
$(f_*\mathcal{F})|_{\mathcal{D}/V} = g_*(\mathcal{F}|_{\mathcal{C}/U})$
对 $\mathcal{C}$ 上任意层 $\mathcal{F}$ 成立
（见《站点上的模》引理 \ref{sites-modules-lemma-localize-morphism-ringed-sites}，
或更基本的《站点》引理 \ref{sites-lemma-localize-morphism}），故结论成立。
\end{proof}

\begin{lemma}
\label{sites-cohomology-lemma-Leray-unbounded}
设 $f : (\mathcal{C}, \mathcal{O}_\mathcal{C}) \to
(\mathcal{D}, \mathcal{O}_\mathcal{D})$ 为对应于连续函子 $u : \mathcal{D} \to \mathcal{C}$ 的环化站点态射
。
作为函子 $D(\mathcal{O}_\mathcal{C}) \to D(\Gamma(\mathcal{O}_\mathcal{D}))$，有
$R\Gamma(\mathcal{D}, -) \circ Rf_* = R\Gamma(\mathcal{C}, -)$。
更一般地，对 $V \in \mathcal{D}$ 且 $U = u(V)$，有
$R\Gamma(U, -) = R\Gamma(V, -) \circ Rf_*$。
\end{lemma}

\begin{proof}
考虑点拓扑斯 $pt$，其 $\mathcal{O}_{pt}$ 由环
$\Gamma(\mathcal{O}_\mathcal{D})$ 给出。存在规范态射
$(\mathcal{D}, \mathcal{O}_\mathcal{D}) \to (pt, \mathcal{O}_{pt})$
它是环化拓扑斯态射，并在结构层的全局截面上诱导出该识别。于是
$D(\mathcal{O}_{pt}) = D(\Gamma(\mathcal{O}_\mathcal{D}))$.
断言
$R\Gamma(\mathcal{D}, -) \circ Rf_* = R\Gamma(\mathcal{C}, -)$
由引理 \ref{sites-cohomology-lemma-derived-pushforward-composition} 应用于
$$
(\mathcal{C}, \mathcal{O}_\mathcal{C}) \to
(\mathcal{D}, \mathcal{O}_\mathcal{D}) \to (pt, \mathcal{O}_{pt})
$$
第二个（更一般的）断言在应用引理 \ref{sites-cohomology-lemma-restrict-direct-image-open} 后由第一个断言得到。
\end{proof}

\begin{lemma}
\label{sites-cohomology-lemma-unbounded-describe-higher-direct-images}
设 $f : (\mathcal{C}, \mathcal{O}_\mathcal{C}) \to
(\mathcal{D}, \mathcal{O}_\mathcal{D})$ 为对应于连续函子 $u : \mathcal{D} \to \mathcal{C}$
的环化站点态射。设 $K \in D(\mathcal{O}_\mathcal{C})$。则 $H^i(Rf_*K)$ 是与下列预层
关联的层：
$$
V \mapsto H^i(u(V), K) = H^i(V, Rf_*K)
$$
\end{lemma}

\begin{proof}
等式 $H^i(u(V), K) = H^i(V, Rf_*K)$ 由引理
\ref{sites-cohomology-lemma-Leray-unbounded} 的第二个断言取上同调得到。层化的断言
由引理 \ref{sites-cohomology-lemma-sheafification-cohomology} 得到。
\end{proof}

\begin{lemma}
\label{sites-cohomology-lemma-modules-abelian-unbounded}
设 $(\mathcal{C}, \mathcal{O}_\mathcal{C})$ 为环化站点。
设 $K$ 为 $D(\mathcal{O}_\mathcal{C})$ 中的对象，并记 $K_{ab}$ 为其在
$D(\underline{\mathbf{Z}}_\mathcal{C})$ 中的像。
\begin{enumerate}
\item 存在典范映射
$R\Gamma(\mathcal{C}, K) \to R\Gamma(\mathcal{C}, K_{ab})$
，它是 $D(\textit{Ab})$ 中的同构。
\item 对任意 $U \in \mathcal{C}$，存在典范映射
$R\Gamma(U, K) \to R\Gamma(U, K_{ab})$
，它是 $D(\textit{Ab})$ 中的同构。
\item 设 $f : (\mathcal{C}, \mathcal{O}_\mathcal{C}) \to
(\mathcal{D}, \mathcal{O}_\mathcal{D})$ 为环化站点态射。存在典范映射
$Rf_*K \to Rf_*(K_{ab})$，它是
$D(\underline{\mathbf{Z}}_\mathcal{D})$ 中的同构。
\end{enumerate}
\end{lemma}

\begin{proof}
映射构造如下。选取表示 $K$ 的 K-内射复形
$\mathcal{I}^\bullet$。选取拟同构
$\mathcal{I}^\bullet \to \mathcal{J}^\bullet$，其中 $\mathcal{J}^\bullet$
是阿贝尔群的 K-内射复形。于是 (1) 中的映射为
$\Gamma(\mathcal{C}, \mathcal{I}^\bullet) \to
\Gamma(\mathcal{C}, \mathcal{J}^\bullet)$
(2) 中的映射为
$\Gamma(U, \mathcal{I}^\bullet) \to \Gamma(U, \mathcal{J}^\bullet)$
而 (3) 中的映射为
$f_*\mathcal{I}^\bullet \to f_*\mathcal{J}^\bullet$.
要证明这些映射是同构，只需证明它们在上同调群和上同调层上诱导同构。
由引理 \ref{sites-cohomology-lemma-unbounded-cohomology-of-open} 和
\ref{sites-cohomology-lemma-unbounded-describe-higher-direct-images}，只需证明映射
$$
H^0(\mathcal{C}, K) \longrightarrow H^0(\mathcal{C}, K_{ab})
$$
是同构。注意
$$
H^0(\mathcal{C}, K) =
\Hom_{D(\mathcal{O}_\mathcal{C})}(\mathcal{O}_\mathcal{C}, K)
$$
，另一个群同样如此。任取表示 $K$ 的 $\mathcal{O}_\mathcal{C}$-模复形
$\mathcal{K}^\bullet$。由将导出范畴构造为局部化可知
$$
\Hom_{D(\mathcal{O}_\mathcal{C})}(\mathcal{O}_\mathcal{C}, K) =
\colim_{s : \mathcal{F}^\bullet \to \mathcal{O}_\mathcal{C}}
\Hom_{K(\mathcal{O}_\mathcal{C})}(\mathcal{F}^\bullet, \mathcal{K}^\bullet)
$$
其中余极限取遍 $\mathcal{O}_\mathcal{C}$-模复形的拟同构 $s$。类似地，有
$$
\Hom_{D(\underline{\mathbf{Z}}_\mathcal{C})}
(\underline{\mathbf{Z}}_\mathcal{C}, K) =
\colim_{s : \mathcal{G}^\bullet \to \underline{\mathbf{Z}}_\mathcal{C}}
\Hom_{K(\underline{\mathbf{Z}}_\mathcal{C})}
(\mathcal{G}^\bullet, \mathcal{K}^\bullet)
$$
接着注意到，拟同构
$s : \mathcal{G}^\bullet \to \underline{\mathbf{Z}}_\mathcal{C}$
其中 $\mathcal{G}^\bullet$ 为平坦的 $\underline{\mathbf{Z}}_\mathcal{C}$-模的上有界复形，
在该系统中共尾。（这由《站点上的模》引理
\ref{sites-modules-lemma-module-quotient-flat} 和《导出范畴》引理
\ref{derived-lemma-subcategory-left-resolution} 得到；见第 \ref{sites-cohomology-section-flat} 节的讨论。）
因此可以构造映射
$H^0(\mathcal{C}, K) \longrightarrow H^0(\mathcal{C}, K_{ab})$
的逆映射：将 $\xi \in H^0(\mathcal{C}, K_{ab})$ 表示为一对
$$
(s : \mathcal{G}^\bullet \to \underline{\mathbf{Z}}_\mathcal{C},
a : \mathcal{G}^\bullet \to \mathcal{K}^\bullet)
$$
其中 $\mathcal{G}^\bullet$ 是平坦的 $\underline{\mathbf{Z}}_\mathcal{C}$-模的上有界复形，
并将其送到
$$
(\mathcal{G}^\bullet \otimes_{\underline{\mathbf{Z}}_\mathcal{C}}
\mathcal{O}_\mathcal{C}
\to \mathcal{O}_\mathcal{C},
\mathcal{G}^\bullet  \otimes_{\underline{\mathbf{Z}}_\mathcal{C}}
\mathcal{O}_\mathcal{C}
\to \mathcal{K}^\bullet)
$$
这里只需注意，第一个箭头由引理
\ref{sites-cohomology-lemma-derived-tor-quasi-isomorphism-other-side} 和
\ref{sites-cohomology-lemma-bounded-flat-K-flat} 是拟同构。具体验证该构造确为逆映射的细节从略。
\end{proof}

\begin{lemma}
\label{sites-cohomology-lemma-adjoint-lower-shriek-restrict}
设 $(\mathcal{C}, \mathcal{O})$ 为环化站点。设 $U$ 为 $\mathcal{C}$ 的对象。记
$j : (\Sh(\mathcal{C}/U), \mathcal{O}_U) \to (\Sh(\mathcal{C}), \mathcal{O})$
为相应的局部化态射。限制函子 $D(\mathcal{O}) \to D(\mathcal{O}_U)$ 是
延零函子 $j_! : D(\mathcal{O}_U) \to D(\mathcal{O})$ 的右伴随。
\end{lemma}

\begin{proof}
需证明
$$
\Hom_{D(\mathcal{O})}(j_!E, F) = \Hom_{D(\mathcal{O}_U)}(E, F|_U)
$$
选取表示 $E$ 的 $\mathcal{O}_U$-模复形 $\mathcal{E}^\bullet$，并选取表示 $F$ 的
K-内射复形 $\mathcal{I}^\bullet$。由引理
\ref{sites-cohomology-lemma-restrict-K-injective-to-open}，$\mathcal{I}^\bullet|_U$ 也是 K-内射的。
因此上式变为
$$
\Hom_{D(\mathcal{O})}(j_!\mathcal{E}^\bullet, \mathcal{I}^\bullet) =
\Hom_{D(\mathcal{O}_U)}(\mathcal{E}^\bullet, \mathcal{I}^\bullet|_U)
$$
这成立是因为 $|_U$ 与 $j_!$ 互为伴随函子
（《站点上的模》引理 \ref{sites-modules-lemma-extension-by-zero}），以及
《导出范畴》引理 \ref{derived-lemma-K-injective}。
\end{proof}

\begin{lemma}
\label{sites-cohomology-lemma-j-shriek-and-tensor}
设 $(\mathcal{C}, \mathcal{O})$ 为环化站点。设 $U \in \Ob(\mathcal{C})$。
对 $L \in D(\mathcal{O}_U)$ 和 $K \in D(\mathcal{O})$，有
$j_!L \otimes_\mathcal{O}^\mathbf{L} K =
j_!(L \otimes_{\mathcal{O}_U}^\mathbf{L} K|_U)$.
\end{lemma}

\begin{proof}
用 $\mathcal{O}_U$-模复形表示 $L$，用 $\mathcal{O}$-模 K-平坦复形表示 $K$，
并应用《站点上的模》引理 \ref{sites-modules-lemma-j-shriek-and-tensor}。
细节从略。
\end{proof}

\begin{lemma}
\label{sites-cohomology-lemma-K-injective-flat}
设 $f : (\Sh(\mathcal{C}), \mathcal{O}_\mathcal{C}) \to
(\Sh(\mathcal{D}), \mathcal{O}_\mathcal{D})$ 为环化拓扑斯之间的平坦态射。
若 $\mathcal{I}^\bullet$ 是 $\mathcal{O}_\mathcal{C}$-模的 K-内射复形，则
$f_*\mathcal{I}^\bullet$ 是 $\mathcal{O}_\mathcal{D}$-模的 K-内射复形。
\end{lemma}

\begin{proof}
这是因为
$$
\Hom_{K(\mathcal{O}_\mathcal{D})}(\mathcal{F}^\bullet, f_*\mathcal{I}^\bullet)
=
\Hom_{K(\mathcal{O}_\mathcal{C})}(f^*\mathcal{F}^\bullet, \mathcal{I}^\bullet)
$$
由《站点上的模》引理
\ref{sites-modules-lemma-adjoint-pullback-pushforward-modules} 以及 $f$ 平坦从而
$f^*$ 正合这一事实可知。
\end{proof}

\begin{lemma}
\label{sites-cohomology-lemma-hom-K-injective}
设 $\mathcal{C}$ 为站点。设 $\mathcal{O} \to \mathcal{O}'$ 为层环映射。
若 $\mathcal{I}^\bullet$ 是 $\mathcal{O}$-模的 K-内射复形，则
$\SheafHom_\mathcal{O}(\mathcal{O}', \mathcal{I}^\bullet)$
是 $\mathcal{O}'$-模的 K-内射复形。
\end{lemma}

\begin{proof}
这是因为
$\Hom_{K(\mathcal{O}')}(\mathcal{G}^\bullet,
\Hom_\mathcal{O}(\mathcal{O}', \mathcal{I}^\bullet)) =
\Hom_{K(\mathcal{O})}(\mathcal{G}^\bullet, \mathcal{I}^\bullet)$
由《站点上的模》引理 \ref{sites-modules-lemma-adjoint-hom-restrict} 成立。
\end{proof}






\section{局部化与上同调}
\label{sites-cohomology-section-localization}

\noindent
设 $\mathcal{C}$ 为站点，设 $f : X \to Y$ 为 $\mathcal{C}$ 的态射。
则得到拓扑斯态射
$$
j_{X/Y} : \Sh(\mathcal{C}/X) \longrightarrow \Sh(\mathcal{C}/Y)
$$
见《站点》第 \ref{sites-section-localize} 节和
\ref{sites-section-more-localize} 节。对于这类拓扑斯态射，某些上同调问题更容易处理。
下面给出一个得到平凡类型基变换定理的例子。

\begin{lemma}
\label{sites-cohomology-lemma-localize-cartesian-square}
设 $\mathcal{C}$ 为站点。设
$$
\xymatrix{
X' \ar[d] \ar[r] & X \ar[d] \\
Y' \ar[r] & Y
}
$$
为 $\mathcal{C}$ 的笛卡尔图。则有
$j_{Y'/Y}^{-1} \circ Rj_{X/Y, *} = Rj_{X'/Y', *} \circ j_{X'/X}^{-1}$
作为函子 $D(\mathcal{C}/X) \to D(\mathcal{C}/Y')$ 成立。
\end{lemma}

\begin{proof}
设 $E \in D(\mathcal{C}/X)$。选取表示 $E$ 的
$\mathcal{C}/X$ 上阿贝尔层的 K-内射复形 $\mathcal{I}^\bullet$。
由引理 \ref{sites-cohomology-lemma-restrict-K-injective-to-open}，$j_{X'/X}^{-1}\mathcal{I}^\bullet$
也是 K-内射的。因此可用 $j_{X'/Y', *}j_{X'/X}^{-1}\mathcal{I}^\bullet$
计算 $Rj_{X'/Y'}(j_{X'/X}^{-1}E)$。于是由《站点》引理
\ref{sites-lemma-localize-cartesian-square} 得到等式。
\end{proof}

\noindent
若 $(\mathcal{C}, \mathcal{O})$ 为环化站点且 $f : X \to Y$ 为 $\mathcal{C}$ 的态射，
则 $j_{X/Y}$ 变为环化拓扑斯态射
$$
j_{X/Y} :
(\Sh(\mathcal{C}/X), \mathcal{O}_X)
\longrightarrow
(\Sh(\mathcal{C}/Y), \mathcal{O}_Y)
$$
见《站点上的模》引理 \ref{sites-modules-lemma-relocalize}。

\begin{lemma}
\label{sites-cohomology-lemma-localize-cartesian-square-modules}
设 $(\mathcal{C}, \mathcal{O})$ 为环化站点。设
$$
\xymatrix{
X' \ar[d] \ar[r] & X \ar[d] \\
Y' \ar[r] & Y
}
$$
为 $\mathcal{C}$ 的笛卡尔图。则有
$j_{Y'/Y}^* \circ Rj_{X/Y, *} = Rj_{X'/Y', *} \circ j_{X'/X}^*$
作为函子 $D(\mathcal{O}_X) \to D(\mathcal{O}_{Y'})$ 成立。
\end{lemma}

\begin{proof}
由于 $j_{Y'/Y}^{-1}\mathcal{O}_Y = \mathcal{O}_{Y'}$，有
$j_{Y'/Y}^* = Lj_{Y'/Y}^* = j_{Y'/Y}^{-1}$。类似地，
$j_{X'/X}^* = Lj_{X'/X}^* = j_{X'/X}^{-1}$。因此由引理
\ref{sites-cohomology-lemma-modules-abelian-unbounded}，只需在阿贝尔层的导出范畴上证明结果，
而这已由引理 \ref{sites-cohomology-lemma-localize-cartesian-square} 完成。
\end{proof}












\section{逆系统与上同调}
\label{sites-cohomology-section-inverse-systems}

\noindent
本节证明关于模层逆系统的一些结果。

\begin{lemma}
\label{sites-cohomology-lemma-ML-general}
设 $I$ 为环 $A$ 的理想，设 $\mathcal{C}$ 为站点。设
$$
\ldots \to \mathcal{F}_3 \to \mathcal{F}_2 \to \mathcal{F}_1
$$
为 $\mathcal{C}$ 上 $A$-模层的逆系统，满足
$\mathcal{F}_n = \mathcal{F}_{n + 1}/I^n\mathcal{F}_{n + 1}$。
设 $p \geq 0$。假设
$$
\bigoplus\nolimits_{n \geq 0} H^{p + 1}(\mathcal{C}, I^n\mathcal{F}_{n + 1})
$$
作为分次 $\bigoplus_{n \geq 0} I^n/I^{n + 1}$-模满足升链条件。
则逆系统 $M_n = H^p(\mathcal{C}, \mathcal{F}_n)$ 满足 Mittag-Leffler 条件
\footnote{事实上，存在 $c \geq 0$，使得对所有 $n \geq c$，
$\Im(M_n \to M_{n - c})$ 都是稳定像。}。
\end{lemma}

\begin{proof}
置 $N_n = H^{p + 1}(\mathcal{C}, I^n\mathcal{F}_{n + 1})$，并令
$\delta_n : M_n \to N_n$ 为由短正合序列
$0 \to I^n\mathcal{F}_{n + 1} \to \mathcal{F}_{n + 1} \to \mathcal{F}_n \to 0$.
诱导的上同调边界映射。则 $\bigoplus \Im(\delta_n) \subset \bigoplus N_n$
是分次子模。具体地，若 $s \in M_n$ 且 $f \in I^m$，则有交换图
$$
\xymatrix{
0 \ar[r] &
I^n\mathcal{F}_{n + 1} \ar[d]_f \ar[r] &
\mathcal{F}_{n + 1} \ar[d]_f \ar[r] &
\mathcal{F}_n \ar[d]_f \ar[r] & 0 \\
0 \ar[r] &
I^{n + m}\mathcal{F}_{n + m + 1} \ar[r] &
\mathcal{F}_{n + m + 1} \ar[r] &
\mathcal{F}_{n + m} \ar[r] & 0
}
$$
中间的垂直映射通过将 $\mathcal{F}_{n + 1}$ 的局部截面提升为
$\mathcal{F}_{n + m + 1}$ 的截面，再乘以 $f$ 给出；其他垂直箭头类似。
因此 $\delta_{n + m}(fs) = f \delta_n(s)$。由假设，可找到
$s_j \in M_{n_j}$（$j = 1, \ldots, N$），使 $\delta_{n_j}(s_j)$
作为分次模生成 $\bigoplus \Im(\delta_n)$。令 $n > c = \max(n_j)$。
对 $s \in M_n$，可找到 $f_j \in I^{n - n_j}$ 使
$\delta_n(s) = \sum f_j \delta_{n_j}(s_j)$。于是
$\delta(s - \sum f_j s_j) = 0$，即存在 $s' \in M_{n + 1}$
映到 $M_n$ 中的 $s - \sum f_js_j$。从而
$$
\Im(M_{n + 1} \to M_{n - c}) = \Im(M_n \to M_{n - c})
$$
这是因为元素 $f_js_j$ 映到 $M_{n - c}$ 中为零。引理得证。
\end{proof}

\begin{lemma}
\label{sites-cohomology-lemma-ML-general-better}
设 $I$ 为环 $A$ 的理想，设 $\mathcal{C}$ 为站点。设
$$
\ldots \to \mathcal{F}_3 \to \mathcal{F}_2 \to \mathcal{F}_1
$$
为 $\mathcal{C}$ 上 $A$-模的逆系统，满足
$\mathcal{F}_n = \mathcal{F}_{n + 1}/I^n\mathcal{F}_{n + 1}$。
设 $p \geq 0$。对给定的 $n$，定义
$$
N_n =
\bigcap\nolimits_{m \geq n}
\Im\left(
H^{p + 1}(\mathcal{C}, I^n\mathcal{F}_{m + 1}) \to
H^{p + 1}(\mathcal{C}, I^n\mathcal{F}_{n + 1})
\right)
$$
若 $\bigoplus N_n$ 作为分次
$\bigoplus_{n \geq 0} I^n/I^{n + 1}$-模满足升链条件，则逆系统
$M_n = H^p(\mathcal{C}, \mathcal{F}_n)$ 满足 Mittag-Leffler 条件
\footnote{事实上，存在 $c \geq 0$，使得对所有 $n \geq c$，
$\Im(M_n \to M_{n - c})$ 都是稳定像。}。
\end{lemma}

\begin{proof}
证明与引理 \ref{sites-cohomology-lemma-ML-general} 的证明完全相同。事实上，只要证明
$\bigoplus N_n$ 是
$\bigoplus H^{p + 1}(\mathcal{C}, I^n\mathcal{F}_{n + 1})$ 的分次
$\bigoplus_{n \geq 0} I^n/I^{n + 1}$-子模，且边界映射
$\delta_n : M_n \to H^{p + 1}(\mathcal{C}, I^n\mathcal{F}_{n + 1})$
的像包含于 $N_n$，上面的论证即给出结果。

\medskip\noindent
设 $\xi \in N_n$ 且 $f \in I^k$。选取 $m \gg n + k$。选取
$\xi' \in H^{p + 1}(\mathcal{C}, I^n\mathcal{F}_{m + 1})$ 提升 $\xi$。
考虑图
$$
\xymatrix{
0 \ar[r] &
I^n\mathcal{F}_{m + 1} \ar[d]_f \ar[r] &
\mathcal{F}_{m + 1} \ar[d]_f \ar[r] &
\mathcal{F}_n \ar[d]_f \ar[r] & 0 \\
0 \ar[r] &
I^{n + k}\mathcal{F}_{m + 1} \ar[r] &
\mathcal{F}_{m + 1} \ar[r] &
\mathcal{F}_{n + k} \ar[r] & 0
}
$$
它按引理 \ref{sites-cohomology-lemma-ML-general} 的证明构造。得到上同调上的诱导映射，并注意到
$f \xi' \in H^{p + 1}(\mathcal{C}, I^{n + k}\mathcal{F}_{m + 1})$
映到 $f \xi$。由于对所有 $m \gg n + k$ 均如此，故 $f\xi \in N_{n + k}$，如所需。

\medskip\noindent
为证明边界映射 $\delta_n$ 的像包含于 $N_n$，考虑图
$$
\xymatrix{
0 \ar[r] &
I^n\mathcal{F}_{m + 1} \ar[d] \ar[r] &
\mathcal{F}_{m + 1} \ar[d] \ar[r] &
\mathcal{F}_n \ar[d] \ar[r] & 0 \\
0 \ar[r] &
I^n\mathcal{F}_{n + 1} \ar[r] &
\mathcal{F}_{n + 1} \ar[r] &
\mathcal{F}_n \ar[r] & 0
}
$$
其中 $m \geq n$。考察上同调上的诱导映射即可得出结论。
\end{proof}

\begin{lemma}
\label{sites-cohomology-lemma-topology-I-adic-general}
设 $I$ 为环 $A$ 的理想，设 $\mathcal{C}$ 为站点。设
$$
\ldots \to \mathcal{F}_3 \to \mathcal{F}_2 \to \mathcal{F}_1
$$
为 $\mathcal{C}$ 上 $A$-模层的逆系统，满足
$\mathcal{F}_n = \mathcal{F}_{n + 1}/I^n\mathcal{F}_{n + 1}$.
设 $p \geq 0$。假设
$$
\bigoplus\nolimits_{n \geq 0} H^p(\mathcal{C}, I^n\mathcal{F}_{n + 1})
$$
作为分次
$\bigoplus_{n \geq 0} I^n/I^{n + 1}$-模满足升链条件。则 $M = \lim H^p(\mathcal{C}, \mathcal{F}_n)$
上的极限拓扑就是 $I$-进拓扑。
\end{lemma}

\begin{proof}
对 $n \geq 1$ 置 $F^n = \Ker(M \to H^p(\mathcal{C}, \mathcal{F}_n))$，并置
$F^0 = M$。
注意到 $I F^n \subset F^{n + 1}$。特别地，$I^n M \subset F^n$。
因此 $I$-进拓扑比极限拓扑精细。反过来，我们将证明：对给定的 $n$，
存在 $m \geq n$ 使得 $F^m \subset I^nM$\footnote{事实上，
存在 $c \geq 0$，使得对所有 $n$ 都有 $F^{n + c} \subset I^nM$。}。
我们有单射
$$
F^n/F^{n + 1} \longrightarrow H^p(\mathcal{C}, \mathcal{F}_{n + 1})
$$
其像包含于
$H^p(\mathcal{C}, I^n\mathcal{F}_{n + 1}) \to
H^p(\mathcal{C}, \mathcal{F}_{n + 1})$.
记
$$
E_n \subset H^p(\mathcal{C}, I^n\mathcal{F}_{n + 1})
$$
为 $F^n/F^{n + 1}$ 的逆像。于是 $\bigoplus E_n$ 是
$\bigoplus H^p(\mathcal{C}, I^n\mathcal{F}_{n + 1})$ 的分次
$\bigoplus I^n/I^{n + 1}$-子模，并且
$\bigoplus E_n \to \bigoplus F^n/F^{n + 1}$ 是分次模同态；细节略去。
由假设，$\bigoplus E_n$ 在 $\bigoplus I^n/I^{n + 1}$ 上由有限个齐次元素生成。
由于 $E_n \to F^n/F^{n + 1}$ 是满射，可知
$\bigoplus F^n/F^{n + 1}$ 也具有同样性质。
因此可以找到 $r$、$c_1, \ldots, c_r \geq 0$ 以及
$a_i \in F^{c_i}$，使它们在 $\bigoplus F^n/F^{n + 1}$ 中的像生成该模。
置 $c = \max(c_i)$。

\medskip\noindent
对 $n \geq c$，我们断言 $I F^n = F^{n + 1}$。该断言表明
$F^{n + c} = I^nF^c \subset I^nM$，正如所需。为证明该断言，
设 $a \in F^{n + 1}$。$a$ 在 $F^{n + 1}/F^{n + 2}$ 中的像是
我们的 $a_i$ 的线性组合。因此，对某些 $f_i \in I^{n + 1 - c_i}$，有
$a  - \sum f_i a_i \in F^{n + 2}$。由于 $n \geq c_i$，
$I^{n + 1 - c_i} = I \cdot I^{n - c_i}$，所以可以写成
$f_i = \sum g_{i, j} h_{i, j}$，其中 $g_{i, j} \in I$
且 $h_{i, j}a_i \in F^n$。于是
$F^{n + 1} = F^{n + 2} + IF^n$。
简单的归纳论证给出，对所有 $e > 0$ 都有 $F^{n + 1} = F^{n + e} + IF^n$。
由此可知 $IF^n$ 在 $F^{n + 1}$ 中稠密。
取 $I$ 的生成元 $k_1, \ldots, k_r$，并考虑连续映射
$$
u : (F^n)^{\oplus r} \longrightarrow F^{n + 1},\quad
(x_1, \ldots, x_r) \mapsto \sum k_i x_i
$$
（在极限拓扑下）。
由上述论证，对所有 $m \geq n$，$u$ 在 $(F^m)^{\oplus r}$ 上的像在
$F^{m + 1}$ 中稠密。由《更多代数》引理
\ref{more-algebra-lemma-open-mapping}（开映射引理），可知 $u$ 是开映射。
因此 $u$ 满射，从而当 $n \geq c$ 时 $IF^n = F^{n + 1}$。
证明完毕。
\end{proof}

\begin{lemma}
\label{sites-cohomology-lemma-topology-I-adic-general-better}
设 $I$ 为环 $A$ 的理想，设 $\mathcal{C}$ 为站点。设
$$
\ldots \to \mathcal{F}_3 \to \mathcal{F}_2 \to \mathcal{F}_1
$$
为 $\mathcal{C}$ 上 $A$-模层的逆系统，满足
$\mathcal{F}_n = \mathcal{F}_{n + 1}/I^n\mathcal{F}_{n + 1}$.
设 $p \geq 0$。对给定的 $n$ 定义
$$
N_n =
\bigcap\nolimits_{m \geq n}
\Im\left(
H^p(\mathcal{C}, I^n\mathcal{F}_{m + 1}) \to
H^p(\mathcal{C}, I^n\mathcal{F}_{n + 1})
\right)
$$
若 $\bigoplus N_n$ 作为分次
$\bigoplus_{n \geq 0} I^n/I^{n + 1}$-模满足升链条件，则
$M = \lim H^p(\mathcal{C}, \mathcal{F}_n)$ 上的极限拓扑就是 $I$-进拓扑。
\end{lemma}

\begin{proof}
证明与引理 \ref{sites-cohomology-lemma-topology-I-adic-general} 的证明完全相同。
事实上，只要证明
$\bigoplus N_n$ 是 $\bigoplus H^{p + 1}(\mathcal{C}, I^n\mathcal{F}_{n + 1})$
的分次 $\bigoplus_{n \geq 0} I^n/I^{n + 1}$-子模，并且
$F^n/F^{n + 1} \subset H^p(\mathcal{C}, \mathcal{F}_{n + 1})$
包含于 $N_n \to H^p(\mathcal{C}, \mathcal{F}_{n + 1})$ 的像中，
就可由那里给出的论证得到结果。在引理
\ref{sites-cohomology-lemma-ML-general-better} 的证明中，我们已经看到了关于模结构的断言。

\medskip\noindent
设 $t \in F^n$。选取元素
$s \in H^p(\mathcal{C}, I^n\mathcal{F}_{n + 1})$
使其映到 $t$ 在 $H^p(\mathcal{C}, \mathcal{F}_{n + 1})$ 中的像。
我们需要证明 $s \in N_n$。现在 $F^n$ 是映射
$M \to H^p(\mathcal{C}, \mathcal{F}_n)$ 的核，因此对所有 $m \geq n$，可将 $t$
映到某个 $t_m \in H^p(\mathcal{C}, \mathcal{F}_{m + 1})$，且该元素在
$H^p(\mathcal{C}, \mathcal{F}_n)$ 中映到零。考虑上同调序列
$$
H^{p - 1}(\mathcal{C}, \mathcal{F}_n) \to
H^p(\mathcal{C}, I^n\mathcal{F}_{m + 1}) \to
H^p(\mathcal{C}, \mathcal{F}_{m + 1}) \to
H^p(\mathcal{C}, \mathcal{F}_n)
$$
它来自短正合序列
$0 \to I^n\mathcal{F}_{m + 1} \to \mathcal{F}_{m + 1} \to \mathcal{F}_n \to 0$.
可以选取 $s_m \in H^p(\mathcal{C}, I^n\mathcal{F}_{m + 1})$
映到 $t_m$。将上述序列与 $m = n$ 时的序列比较，可知
$s_m$ 映到 $s$，相差一个如下映射的像中的元素：
$H^{p - 1}(\mathcal{C}, \mathcal{F}_n) \to
H^p(\mathcal{C}, I^n\mathcal{F}_{n + 1})$.
然而，该映射分解经过映射
$H^p(\mathcal{C}, I^n\mathcal{F}_{m + 1}) \to
H^p(\mathcal{C}, I^n\mathcal{F}_{n + 1})$
，因此可知 $s$ 属于所需的像。
\end{proof}












\section{导出极限与同伦极限}
\label{sites-cohomology-section-derived-limits}

\noindent
设 $\mathcal{C}$ 为站点。考虑范畴 $\mathcal{C} \times \mathbf{N}$，其中
$\Mor((U, n), (V, m)) = \emptyset$（若 $n > m$），而
$\Mor((U, n), (V, m)) = \Mor(U, V)$（其他情形）。通过规定覆盖为满足
$\{U_i \to U\}$ 是 $\mathcal{C}$ 的覆盖的族
$\{(U_i, n) \to (U, n)\}$，我们给它赋予站点结构。
读者立即验证，$\mathcal{C} \times \mathbf{N}$ 上的层
正好等同于 $\mathcal{C}$ 上层的逆系统。
特别地，$\textit{Ab}(\mathcal{C} \times \mathbf{N})$
就是 $\mathcal{C}$ 上阿贝尔层的逆系统范畴。
现在考虑函子
$$
\lim : \textit{Ab}(\mathcal{C} \times \mathbf{N}) \to \textit{Ab}(\mathcal{C})
$$
它将逆系统取为其极限。这不过就是 $g_*$，其中
$g : \Sh(\mathcal{C} \times \mathbf{N}) \to \Sh(\mathcal{C})$
是与连续且余连续函子 $\mathcal{C} \times \mathbf{N} \to \mathcal{C}$
相伴的拓扑斯态射。（注意，$g^{-1}$ 将 $\mathcal{C}$ 上的层赋给
相应的常值逆系统。）

\medskip\noindent
由上面说明的一般构造，我们得到导出函子
$$
R\lim = Rg_* : D(\mathcal{C} \times \mathbf{N}) \to D(\mathcal{C}).
$$
如所示，该函子通常记作 $R\lim$。

\medskip\noindent
另一方面，连续且余连续函子
$\mathcal{C} \to \mathcal{C} \times \mathbf{N}$、$U \mapsto (U, n)$
定义了拓扑斯态射
$i_n : \Sh(\mathcal{C}) \to \Sh(\mathcal{C} \times \mathbf{N})$。
当然，$i_n^{-1}$ 是取逆系统第 $n$ 项的函子。因此有函子变换
$i_{n + 1}^{-1} \to i_n^{-1}$。于是给定
$K \in D(\mathcal{C} \times \mathbf{N})$，我们得到
$K_n = i_n^{-1}K \in D(\mathcal{C})$ 以及映射 $K_{n + 1} \to K_n$。
在《导出范畴》定义 \ref{derived-definition-derived-limit} 中，
我们定义了同伦极限的概念
$$
R\lim K_n \in D(\mathcal{C})
$$
我们断言这两个概念一致（在有意义的范围内）。

\begin{lemma}
\label{sites-cohomology-lemma-derived-limit-is-ok}
设 $\mathcal{C}$ 为站点，设 $K$ 为
$D(\mathcal{C} \times \mathbf{N})$ 的对象。按上面的方式置
$K_n = i_n^{-1}K$。则
$$
R\lim K \cong R\lim K_n
$$
在 $D(\mathcal{C})$ 中成立。
\end{lemma}

\begin{proof}
为在 $D(\mathcal{C} \times \mathbf{N})$ 的对象 $K$ 上计算 $R\lim$，
我们选取一个 K-内射表示 $\mathcal{I}^\bullet$，其各项是
$\textit{Ab}(\mathcal{C} \times \mathbf{N})$ 中的内射对象；见《内射对象》定理
\ref{injectives-theorem-K-injective-embedding-grothendieck}.
我们可以并且确实把 $\mathcal{I}^\bullet$ 看作复形
$(\mathcal{I}_n^\bullet)$ 的逆系统，于是得到
$$
R\lim K = \lim \mathcal{I}_n^\bullet
$$
其中右端是逐项逆极限。

\medskip\noindent
设 $\mathcal{J} = (\mathcal{J}_n)$ 是
$\textit{Ab}(\mathcal{C} \times \mathbf{N})$ 中的内射对象。态射
$(U, n) \to (U, n + 1)$ 是 $\mathcal{C} \times \mathbf{N}$ 中的单态射，故
$\mathcal{J}(U, n + 1) \to \mathcal{J}(U, n)$ 是满射
（引理 \ref{sites-cohomology-lemma-restriction-along-monomorphism-surjective}）。
因此 $\mathcal{J}_{n + 1} \to \mathcal{J}_n$ 作为预层映射是满射。

\medskip\noindent
注意函子 $i_n^{-1}$ 有一个正合左伴随 $i_{n, !}$。
具体地，$i_{n, !}\mathcal{F}$ 是逆系统
$\ldots 0 \to 0 \to \mathcal{F} \to \ldots \to \mathcal{F}$.
因此，由《导出范畴》引理
\ref{derived-lemma-adjoint-preserve-K-injectives}，复形
$i_n^{-1}\mathcal{I}^\bullet = \mathcal{I}_n^\bullet$ 是 K-内射的。

\medskip\noindent
由于我们选取的 K-内射复形各项都是内射的，我们得到
$$
0 \to  \lim \mathcal{I}_n^\bullet \to \prod \mathcal{I}_n^\bullet
\to \prod \mathcal{I}_n^\bullet \to 0
$$
是阿贝尔层复形的短正合序列，因为它也是阿贝尔预层复形的短正合序列。
此外，中间和右端的乘积表示 $D(\mathcal{C})$ 中的乘积；见
《内射对象》引理 \ref{injectives-lemma-derived-products} 及其证明
（这里正是使用 $\mathcal{I}_n^\bullet$ 为 K-内射的地方）。
因此，按三角范畴中同伦极限的定义，$R\lim K$ 是逆系统 $(K_n)$ 的同伦极限。
\end{proof}

\begin{lemma}
\label{sites-cohomology-lemma-RGamma-commutes-with-Rlim}
设 $(\mathcal{C}, \mathcal{O})$ 为环化站点。函子
$R\Gamma(\mathcal{C}, -)$ 以及对 $U \in \Ob(\mathcal{C})$ 的
$R\Gamma(U, -)$ 都与 $R\lim$ 交换。此外，对
$$
0 \to
R^1\lim H^{m - 1}(U, K_n) \to H^m(U, R\lim K_n) \to
\lim H^m(U, K_n) \to 0
$$
对 $D(\mathcal{O})$ 中任意逆系统 $(K_n)$ 及 $m \in \mathbf{Z}$，都有短正合序列。
对 $H^m(\mathcal{C}, R\lim K_n)$ 也有类似结论。
\end{lemma}

\begin{proof}
第一个断言由《内射对象》引理
\ref{injectives-lemma-RF-commutes-with-Rlim} 得到。然后可将《更多代数》备注
\ref{more-algebra-remark-compare-derived-limit} 应用于
$R\lim R\Gamma(U, K_n) = R\Gamma(U, R\lim K_n)$，从而得到这些短正合序列。
\end{proof}

\begin{lemma}
\label{sites-cohomology-lemma-Rf-commutes-with-Rlim}
设 $f : (\Sh(\mathcal{C}), \mathcal{O}) \to (\Sh(\mathcal{C}'), \mathcal{O}')$
为环化拓扑斯态射。则 $Rf_*$ 与 $R\lim$ 交换，即
$Rf_*$ 与导出极限交换。
\end{lemma}

\begin{proof}
设 $(K_n)$ 为 $D(\mathcal{O})$ 中对象的逆系统。对 $n$ 归纳，
可以选取 $\mathcal{O}$-模的实际复形 $\mathcal{K}_n^\bullet$ 以及复形映射
$\mathcal{K}_{n + 1}^\bullet \to \mathcal{K}_n^\bullet$，它们在
$D(\mathcal{O})$ 中表示映射 $K_{n + 1} \to K_n$。换言之，存在
$D(\mathcal{C} \times \mathbf{N})$ 中的对象 $K$，其对应的逆系统就是给定的逆系统。
接下来考虑交换图
$$
\xymatrix{
\Sh(\mathcal{C} \times \mathbf{N}) \ar[r]_g \ar[d]_{f \times 1} &
\Sh(\mathcal{C}) \ar[d]_f \\
\Sh(\mathcal{C}' \times \mathbf{N}) \ar[r]^{g'} &
\Sh(\mathcal{C}')
}
$$
其中各箭头为拓扑斯态射。于是
$R\lim R(f \times 1)_*K = Rf_* R\lim K$。逐一展开定义并使用引理
\ref{sites-cohomology-lemma-derived-limit-is-ok}，得到
$R\lim (Rf_*K_n) = Rf_*(R\lim K_n)$。

\medskip\noindent
当 $\mathcal{C}$ 有足够多的点时的另一证明。考虑定义性的有区别三角
$$
R\lim K_n \to \prod K_n \to \prod K_n
$$
它位于 $D(\mathcal{O})$ 中。应用正合函子 $Rf_*$，得到有区别三角
$$
Rf_*(R\lim K_n) \to Rf_*\left(\prod K_n\right) \to Rf_*\left(\prod K_n\right)
$$
它位于 $D(\mathcal{O}')$ 中。因此只需证明 $Rf_*$ 与导出范畴中的乘积交换
即可（这些乘积不只是复形的乘积，见《内射对象》引理
\ref{injectives-lemma-derived-products}）。然而，由于按引理
\ref{sites-cohomology-lemma-adjoint}，$Rf_*$ 是右伴随，这一点形式上成立（见《范畴》引理
\ref{categories-lemma-adjoint-exact}）。
注意：不能直接应用《范畴》引理
\ref{categories-lemma-adjoint-exact}，因为 $R\lim K_n$ 不是
$D(\mathcal{O})$ 中的极限。
\end{proof}

\begin{remark}
\label{sites-cohomology-remark-discuss-derived-limit}
设 $(\mathcal{C}, \mathcal{O})$ 为环化站点，设 $(K_n)$ 为
$D(\mathcal{O})$ 中的逆系统。置 $K = R\lim K_n$。对每个 $n$ 和 $m$，
令 $\mathcal{H}^m_n = H^m(K_n)$ 为 $K_n$ 的第 $m$ 个上同调层，并类似地置
$\mathcal{H}^m = H^m(K)$。记 $\underline{\mathcal{H}}^m_n$ 为预层
$$
U \longmapsto \underline{\mathcal{H}}^m_n(U) = H^m(U, K_n)
$$
类似地置 $\underline{\mathcal{H}}^m(U) = H^m(U, K)$。
由引理 \ref{sites-cohomology-lemma-sheafification-cohomology} 可知，
$\mathcal{H}^m_n$ 是 $\underline{\mathcal{H}}^m_n$ 的层化，
$\mathcal{H}^m$ 是 $\underline{\mathcal{H}}^m$ 的层化。
有如下交换图
$$
\xymatrix{
K \ar@{=}[d] &
\underline{\mathcal{H}}^m \ar[d] \ar[r] & 
\mathcal{H}^m \ar[d] \\
R\lim K_n &
\lim \underline{\mathcal{H}}^m_n \ar[r] & 
\lim \mathcal{H}^m_n
}
$$
一般而言，$\lim \mathcal{H}^m_n$ 不一定是
$\lim \underline{\mathcal{H}}^m_n$ 的层化。
若 $U \in \mathcal{C}$，则有短正合序列
\begin{equation}
\label{sites-cohomology-equation-ses-Rlim-over-U}
0 \to
R^1\lim \underline{\mathcal{H}}^{m - 1}_n(U) \to
\underline{\mathcal{H}}^m(U) \to
\lim \underline{\mathcal{H}}^m_n(U) \to 0
\end{equation}
这是由引理 \ref{sites-cohomology-lemma-RGamma-commutes-with-Rlim} 得到的。
\end{remark}

\noindent
下面的引理适用于代数空间或代数叠上的拟凝聚模逆系统，且其过渡映射为满射。

\begin{lemma}
\label{sites-cohomology-lemma-inverse-limit-is-derived-limit}
设 $(\mathcal{C}, \mathcal{O})$ 为环化站点，设 $(\mathcal{F}_n)$ 为
$\mathcal{O}$-模的逆系统。设 $\mathcal{B} \subset \Ob(\mathcal{C})$ 为子集。假设
\begin{enumerate}
\item $\mathcal{C}$ 的每个对象都有一个覆盖，其成员属于 $\mathcal{B}$；
\item 对 $p > 0$ 及 $U \in \mathcal{B}$，有 $H^p(U, \mathcal{F}_n) = 0$；
\item 对 $U \in \mathcal{B}$，逆系统 $\mathcal{F}_n(U)$ 的 $R^1\lim$ 消失。
\end{enumerate}
则 $R\lim \mathcal{F}_n = \lim \mathcal{F}_n$，且对 $p > 0$ 及
$U \in \mathcal{B}$，有 $H^p(U, \lim \mathcal{F}_n) = 0$。
\end{lemma}

\begin{proof}
置 $K_n = \mathcal{F}_n$，并置 $K = R\lim \mathcal{F}_n$。使用备注
\ref{sites-cohomology-remark-discuss-derived-limit} 的记号以及假设 (2)，可知对 $U \in \mathcal{B}$，
当 $m \not = 0$ 时 $\underline{\mathcal{H}}_n^m(U) = 0$，而
$\underline{\mathcal{H}}_n^0(U) = \mathcal{F}_n(U)$。
由等式（\ref{sites-cohomology-equation-ses-Rlim-over-U}）及假设 (3)，可知当 $m \not = 0$ 时
$\underline{\mathcal{H}}^m(U) = 0$，而当 $m = 0$ 时它等于
$\lim \mathcal{F}_n(U)$。利用 (1) 进行层化，得到当 $m \not = 0$ 时
$\mathcal{H}^m = 0$，而当 $m = 0$ 时它等于 $\lim \mathcal{F}_n$。
因此 $K = \lim \mathcal{F}_n$。由于对 $m > 0$ 有
$H^m(U, K) = \underline{\mathcal{H}}^m(U) = 0$（见上文），
故第二个断言成立。
\end{proof}

\begin{lemma}
\label{sites-cohomology-lemma-cohomology-derived-limit-injective}
设 $(\mathcal{C}, \mathcal{O})$ 为环化站点，设 $(K_n)$ 为
$D(\mathcal{O})$ 中的逆系统。设 $V \in \Ob(\mathcal{C})$ 且
$m \in \mathbf{Z}$。假设存在整数 $n(V)$ 以及 $V$ 的覆盖的共终系统
$\text{Cov}_V$，使得对 $\{V_i \to V\} \in \text{Cov}_V$，有
\begin{enumerate}
\item $R^1\lim H^{m - 1}(V_i, K_n) = 0$；
\item 当 $n \geq n(V)$ 时，$H^m(V_i, K_n) \to H^m(V_i, K_{n(V)})$
是单射。
\end{enumerate}
则截面上的映射 $H^m(R\lim K_n)(V) \to H^m(K_{n(V)})(V)$ 是单射。
\end{lemma}

\begin{proof}
设 $\gamma \in H^m(R\lim K_n)(V)$ 映到
$H^m(K_{n(V)})(V)$ 中的零元。由于 $H^m(R\lim K_n)$ 是预层
$U \mapsto H^m(U, R\lim K_n)$ 的层化（由引理
\ref{sites-cohomology-lemma-sheafification-cohomology}），可以选取 $\{V_i \to V\} \in \text{Cov}_V$
以及元素 $\tilde\gamma_i \in H^m(V_i, R\lim K_n)$，使其映到
$\gamma|_{V_i}$。
于是 $\tilde\gamma_i$ 映到 $\tilde\gamma_{i, n(V)} \in H^m(V_i, K_{n(V)})$。
由于 $H^m(K_{n(V)})$ 是预层 $U \mapsto H^m(U, K_{n(V)})$ 的层化
（再次使用引理 \ref{sites-cohomology-lemma-sheafification-cohomology}），可知将
$\{V_i \to V\}$ 替换为其细化后，可以假设对所有 $i$ 都有
$\tilde\gamma_{i, n(V)} = 0$。
对该覆盖，考虑引理 \ref{sites-cohomology-lemma-RGamma-commutes-with-Rlim} 给出的短正合序列
$$
0 \to
R^1\lim H^{m - 1}(V_i, K_n) \to H^m(V_i, R\lim K_n) \to
\lim H^m(V_i, K_n) \to 0
$$
。由假设 (1)，左端的群为零；由假设 (2)，右端的群到
$H^m(V_i, K_{n(V)})$ 的映射为单射。因此 $\tilde\gamma_i = 0$，
从而 $\gamma = 0$，正如所需。
\end{proof}

\begin{lemma}
\label{sites-cohomology-lemma-is-limit-per-object}
设 $(\mathcal{C}, \mathcal{O})$ 为环化站点，设 $E \in D(\mathcal{O})$。
设 $\mathcal{B} \subset \Ob(\mathcal{C})$ 为子集。假设
\begin{enumerate}
\item $\mathcal{C}$ 的每个对象都有一个覆盖，其成员属于 $\mathcal{B}$；
\item 对每个 $V \in \mathcal{B}$，存在函数
$p(V, -) : \mathbf{Z} \to \mathbf{Z}$ 以及 $V$ 的覆盖的共终系统 $\text{Cov}_V$，使得
$$
H^p(V_i, H^{m - p}(E)) = 0
$$
对所有 $\{V_i \to V\} \in \text{Cov}_V$ 及满足 $p > p(V, m)$ 的整数 $p,m$ 都成立。
\end{enumerate}
则《导出范畴》备注
\ref{derived-remark-map-into-derived-limit-truncations} 中的映射
$E \to R\lim \tau_{\geq -n} E$ 在 $D(\mathcal{O})$ 中是同构。
\end{lemma}

\begin{proof}
置 $K_n = \tau_{\geq -n}E$，并置 $K = R\lim K_n$。
典范映射 $E \to K$ 来自典范映射 $E \to K_n = \tau_{\geq -n}E$。
我们需要证明 $E \to K$ 在上同调层上诱导同构
$H^m(E) \to H^m(K)$。在证明的其余部分固定 $m$。若 $n \geq -m$，则
$E \to \tau_{\geq -n}E = K_n$ 诱导同构 $H^m(E) \to H^m(K_n)$。
要完成证明，只需证明对每个 $V \in \mathcal{B}$，存在整数 $n(V) \geq -m$，使映射
 $H^m(K)(V) \to H^m(K_{n(V)})(V)$ 为单射。事实上，此时复合映射
$$
H^m(E)(V) \to H^m(K)(V) \to H^m(K_{n(V)})(V)
$$
是双射，而第二个箭头是单射，故第一个箭头是双射。由性质 (1) 可知
$H^m(E) \to H^m(K)$ 是同构。置
$$
n(V) = 1 + \max\{-m, p(V, m - 1) - m, -1 + p(V, m) - m, -2 + p(V, m + 1) - m\}.
$$
从而无论如何都有 $n(V) \geq -m$。断言：当
$$
H^{m - 1}(V_i, K_{n + 1}) \to H^{m - 1}(V_i, K_n)
\quad\text{且}\quad
H^m(V_i, K_{n + 1}) \to H^m(V_i, K_n)
$$
对 $n \geq n(V)$ 且 $\{V_i \to V\} \in \text{Cov}_V$ 时，上述映射都是同构。
该断言意味着引理 \ref{sites-cohomology-lemma-cohomology-derived-limit-injective}
中的条件 (1) 和 (2) 得到满足，因而推出所需的单射性。
回忆《导出范畴》备注
\ref{derived-remark-truncation-distinguished-triangle}，我们有有区别三角
$$
H^{-n - 1}(E)[n + 1] \to
K_{n + 1} \to K_n \to H^{-n - 1}(E)[n + 2]
$$
考察相应的上同调长正合序列可知，只要
$$
H^{m + n}(V_i, H^{-n - 1}(E)),\quad
H^{m + n + 1}(V_i, H^{-n - 1}(E)),\quad
H^{m + n + 2}(V_i, H^{-n - 1}(E))
$$
在 $n \geq n(V)$ 且 $\{V_i \to V\} \in \text{Cov}_V$ 时为零，断言就成立。
这由我们对 $n(V)$ 的选择以及引理中的假设推出。
\end{proof}

\begin{lemma}
\label{sites-cohomology-lemma-is-limit-spaltenstein}
设 $(\mathcal{C}, \mathcal{O})$ 为环化站点，设 $E \in D(\mathcal{O})$。
设 $\mathcal{B} \subset \Ob(\mathcal{C})$ 为子集。假设
\begin{enumerate}
\item $\mathcal{C}$ 的每个对象都有一个覆盖，其成员属于 $\mathcal{B}$；
\item 对每个 $V \in \mathcal{B}$，存在整数 $d_V \geq 0$ 以及 $V$ 的覆盖的共终系统
$\text{Cov}_V$，使得
$$
H^p(V_i, H^q(E)) = 0 \text{（对 }
\{V_i \to V\} \in \text{Cov}_V,\ p > d_V, \text{ 且 }q < 0\text{）}
$$
\end{enumerate}
则《导出范畴》备注
\ref{derived-remark-map-into-derived-limit-truncations} 中的映射
$E \to R\lim \tau_{\geq -n} E$ 在 $D(\mathcal{O})$ 中是同构。
\end{lemma}

\begin{proof}
这由引理 \ref{sites-cohomology-lemma-is-limit-per-object}（取
$p(V, m) = d_V + \max(0, m)$）推出。
\end{proof}

\begin{lemma}
\label{sites-cohomology-lemma-is-limit}
设 $(\mathcal{C}, \mathcal{O})$ 为环化站点，设 $E \in D(\mathcal{O})$。
假设存在函数 $p(-) : \mathbf{Z} \to \mathbf{Z}$ 及子集
$\mathcal{B} \subset \Ob(\mathcal{C})$，使得
\begin{enumerate}
\item $\mathcal{C}$ 的每个对象都有一个覆盖，其成员属于 $\mathcal{B}$；
\item 对 $p > p(m)$ 及 $V \in \mathcal{B}$，有
$H^p(V, H^{m - p}(E)) = 0$。
\end{enumerate}
则《导出范畴》备注
\ref{derived-remark-map-into-derived-limit-truncations} 中的映射
$E \to R\lim \tau_{\geq -n} E$ 在 $D(\mathcal{O})$ 中是同构。
\end{lemma}

\begin{proof}
应用引理 \ref{sites-cohomology-lemma-is-limit-per-object}，取 $p(V, m) = p(m)$，并令
$\text{Cov}_V$ 等于所有满足对每个 $i$ 都有 $V_i \in \mathcal{B}$ 的
覆盖 $\{V_i \to V\}$ 所成的集合。
\end{proof}

\begin{lemma}
\label{sites-cohomology-lemma-is-limit-dimension}
设 $(\mathcal{C}, \mathcal{O})$ 为环化站点，设 $E \in D(\mathcal{O})$。
假设存在整数 $d \geq 0$ 及子集 $\mathcal{B} \subset \Ob(\mathcal{C})$，使得
\begin{enumerate}
\item $\mathcal{C}$ 的每个对象都有一个覆盖，其成员属于 $\mathcal{B}$；
\item 对 $p > d$、$q < 0$ 及 $V \in \mathcal{B}$，有
$H^p(V, H^q(E)) = 0$。
\end{enumerate}
则《导出范畴》备注
\ref{derived-remark-map-into-derived-limit-truncations} 中的映射
$E \to R\lim \tau_{\geq -n} E$ 在 $D(\mathcal{O})$ 中是同构。
\end{lemma}

\begin{proof}
应用引理 \ref{sites-cohomology-lemma-is-limit-spaltenstein}，取 $d_V = d$，并令
$\text{Cov}_V$ 等于所有满足对每个 $i$ 都有 $V_i \in \mathcal{B}$ 的
覆盖 $\{V_i \to V\}$ 所成的集合。
\end{proof}

\noindent
上述引理可用于在某些情形计算上同调。

\begin{lemma}
\label{sites-cohomology-lemma-cohomology-over-U-trivial}
设 $(\mathcal{C}, \mathcal{O})$ 为环化站点，设 $K$
为 $D(\mathcal{O})$ 的对象。设 $\mathcal{B} \subset \Ob(\mathcal{C})$ 为子集。假设
\begin{enumerate}
\item $\mathcal{C}$ 的每个对象都有一个覆盖，其成员属于 $\mathcal{B}$；
\item 对所有 $p > 0$、$q \in \mathbf{Z}$ 及 $U \in \mathcal{B}$，有
$H^p(U, H^q(K)) = 0$。
\end{enumerate}
则对 $q \in \mathbf{Z}$ 及 $U \in \mathcal{B}$，有
$H^q(U, K) = H^0(U, H^q(K))$。
\end{lemma}

\begin{proof}
注意，由引理 \ref{sites-cohomology-lemma-is-limit-dimension}（取 $d = 0$），有
$K = R\lim \tau_{\geq -n} K$。令 $U \in \mathcal{B}$。由等式
（\ref{sites-cohomology-equation-ses-Rlim-over-U}）得到短正合序列
$$
0 \to R^1\lim H^{q - 1}(U, \tau_{\geq -n}K) \to
H^q(U, K) \to \lim H^q(U, \tau_{\geq -n}K) \to 0
$$
由《导出范畴》引理
\ref{derived-lemma-two-ss-complex-functor} 的谱序列，条件 (2) 蕴含对所有 $q$ 有
$H^q(U, \tau_{\geq -n} K) = H^0(U, H^q(\tau_{\geq -n} K))$。
该谱序列收敛，因为 $\tau_{\geq -n}K$ 下有界。若 $n > -q$，则
$H^q(\tau_{\geq -n}K) = H^q(K)$。因此，所显示短正合序列左端和右端的系统
最终分别恒定为 $H^0(U, H^{q - 1}(K))$ 和 $H^0(U, H^q(K))$，引理得证。
\end{proof}

\noindent
下面是另一种可以描述导出极限的情形。

\begin{lemma}
\label{sites-cohomology-lemma-derived-limit-suitable-system}
设 $(\mathcal{C}, \mathcal{O})$ 为环化站点。设 $(K_n)$
为 $D(\mathcal{O})$ 中对象的逆系统。设 $\mathcal{B} \subset \Ob(\mathcal{C})$ 为子集。假设
\begin{enumerate}
\item $\mathcal{C}$ 的每个对象都有一个覆盖，其成员属于 $\mathcal{B}$；
\item 对所有 $U \in \mathcal{B}$ 及 $q \in \mathbf{Z}$，都有
\begin{enumerate}
\item 对 $p > 0$，有 $H^p(U, H^q(K_n)) = 0$；
\item 逆系统 $H^0(U, H^q(K_n))$ 的 $R^1\lim$ 消失。
\end{enumerate}
\end{enumerate}
则对 $q \in \mathbf{Z}$，有 $H^q(R\lim K_n) = \lim H^q(K_n)$。
\end{lemma}

\begin{proof}
置 $K = R\lim K_n$。我们使用备注
\ref{sites-cohomology-remark-discuss-derived-limit} 中的记号。令 $U \in \mathcal{B}$。
由引理 \ref{sites-cohomology-lemma-cohomology-over-U-trivial} 及 (2)(a)，有
$H^q(U, K_n) = H^0(U, H^q(K_n))$。
由于函子 $R\Gamma(U, -)$ 与导出极限交换，有
$$
H^q(U, K) = H^q(R\lim R\Gamma(U, K_n)) = \lim H^0(U, H^q(K_n))
$$
其中最后一个等号由《更多代数》备注
\ref{more-algebra-remark-compare-derived-limit} 及假设 (2)(b) 得到。因此 $H^q(U, K)$
是层 $H^q(K_n)$ 在 $U$ 上截面的逆极限。由于 $\lim H^q(K_n)$ 是一个层，
利用假设 (1) 可知，$H^q(K)$（它是预层 $U \mapsto H^q(U, K)$ 的层化）
等于 $\lim H^q(K_n)$。引理得证。
\end{proof}







\section{构造 K-内射分辨}
\label{sites-cohomology-section-K-injective}

\noindent
设 $(\mathcal{C}, \mathcal{O})$ 为环化站点。设 $\mathcal{F}^\bullet$ 为
$\mathcal{O}$-模复形。范畴 $\textit{Mod}(\mathcal{O})$ 有足够多的内射对象，
因此可使用《导出范畴》引理 \ref{derived-lemma-special-inverse-system}
构造如下图
$$
\xymatrix{
\ldots \ar[r] &
\tau_{\geq -2}\mathcal{F}^\bullet \ar[r] \ar[d] &
\tau_{\geq -1}\mathcal{F}^\bullet \ar[d] \\
\ldots \ar[r] & \mathcal{I}_2^\bullet \ar[r] & \mathcal{I}_1^\bullet
}
$$
它位于 $\mathcal{O}$-模复形范畴中，并满足
\begin{enumerate}
\item 竖直箭头是拟同构；
\item $\mathcal{I}_n^\bullet$ 是下有界的内射对象复形；
\item 箭头 $\mathcal{I}_{n + 1}^\bullet \to \mathcal{I}_n^\bullet$
是逐项可裂满射。
\end{enumerate}
 $\mathcal{O}$-模范畴有极限（在预层层面计算），因此可以形成逐项极限
$\mathcal{I}^\bullet = \lim_n \mathcal{I}_n^\bullet$。由《导出范畴》引理
\ref{derived-lemma-bounded-below-injectives-K-injective} 以及
\ref{derived-lemma-limit-K-injectives}
这是 K-内射复形。一般而言，典范映射
\begin{equation}
\label{sites-cohomology-equation-into-candidate-K-injective}
\mathcal{F}^\bullet \to \mathcal{I}^\bullet
\end{equation}
不一定是拟同构。下面的引理给出使其成为拟同构的一些条件。

\begin{lemma}
\label{sites-cohomology-lemma-K-injective}
在上述情形下，记
$\mathcal{H}^m = H^m(\mathcal{F}^\bullet)$ 为第 $m$ 个上同调层。
设 $\mathcal{B} \subset \Ob(\mathcal{C})$ 为子集，设 $d \in \mathbf{N}$。假设
\begin{enumerate}
\item $\mathcal{C}$ 的每个对象都有一个覆盖，其成员属于 $\mathcal{B}$；
\item 对每个 $U \in \mathcal{B}$，当 $p > d$ 且 $q < 0$ 时有
$H^p(U, \mathcal{H}^q) = 0$\footnote{只要
$\forall m$、$\exists p(m)$，且当 $p > p(m)$ 时有
$H^p(U. \mathcal{H}^{m - p}) = 0$ 即可；见引理
\ref{sites-cohomology-lemma-is-limit}。}。
\end{enumerate}
则（\ref{sites-cohomology-equation-into-candidate-K-injective}）是拟同构。
\end{lemma}

\begin{proof}
由《导出范畴》引理 \ref{derived-lemma-difficulty-K-injectives}，只需证明映射
$\mathcal{F}^\bullet \to R\lim \tau_{\geq -n} \mathcal{F}^\bullet$
是同构。这由引理 \ref{sites-cohomology-lemma-is-limit-dimension} 得到。
\end{proof}

\noindent
下面是关于模复形逆系统逐项极限的上同调层的技术引理。我们应尽可能避免使用此引理，
而应改用导出逆极限的论证。

\begin{lemma}
\label{sites-cohomology-lemma-inverse-limit-complexes}
设 $(\mathcal{C}, \mathcal{O})$ 为环化站点。设
$(\mathcal{F}_n^\bullet)$ 为 $\mathcal{O}$-模复形的逆系统。设 $m \in \mathbf{Z}$。
给定 $\mathcal{B} \subset \Ob(\mathcal{C})$ 及整数 $n_0$，使得
\begin{enumerate}
\item $\mathcal{C}$ 的每个对象都有一个覆盖，其成员属于 $\mathcal{B}$；
\item 对每个 $U \in \mathcal{B}$，
\begin{enumerate}
\item 阿贝尔群系统
$\mathcal{F}_n^{m - 2}(U)$ 和 $\mathcal{F}_n^{m - 1}(U)$ 的 $R^1\lim$ 消失
（例如，它们具有 Mittag-Leffler 性质）；
\item 阿贝尔群系统 $H^{m - 1}(\mathcal{F}_n^\bullet(U))$ 的 $R^1\lim$ 消失
（例如，它具有 Mittag-Leffler 性质）；
\item 有
$H^m(\mathcal{F}_n^\bullet(U)) = H^m(\mathcal{F}_{n_0}^\bullet(U))$
对所有 $n \geq n_0$ 都成立。
\end{enumerate}
\end{enumerate}
则映射 $H^m(\mathcal{F}^\bullet) \to \lim H^m(\mathcal{F}_n^\bullet)
\to H^m(\mathcal{F}_{n_0}^\bullet)$ 是层的同构，其中
$\mathcal{F}^\bullet = \lim \mathcal{F}_n^\bullet$ 是逐项逆极限。
\end{lemma}

\begin{proof}
令 $U \in \mathcal{B}$。注意 $H^m(\mathcal{F}^\bullet(U))$ 是下列复形
的上同调：
$$
\lim_n \mathcal{F}_n^{m - 2}(U) \to
\lim_n \mathcal{F}_n^{m - 1}(U) \to
\lim_n \mathcal{F}_n^m(U) \to
\lim_n \mathcal{F}_n^{m + 1}(U)
$$
中从左起第三个位置。由假设 (2)(a) 和 (2)(b)，可应用《更多代数》引理
\ref{more-algebra-lemma-apply-Mittag-Leffler-again}，从而得到
$$
H^m(\mathcal{F}^\bullet(U)) = \lim H^m(\mathcal{F}_n^\bullet(U))
$$
由假设 (2)(c)，得到
$$
H^m(\mathcal{F}^\bullet(U)) = H^m(\mathcal{F}_n^\bullet(U))
$$
对所有 $n \geq n_0$ 成立。由假设 (1)，预层
$U \mapsto H^m(\mathcal{F}^\bullet(U))$ 的层化等于对所有 $n \geq n_0$ 的预层
$U \mapsto H^m(\mathcal{F}_n^\bullet(U))$ 的层化。因此层的逆系统
$H^m(\mathcal{F}_n^\bullet)$ 在 $n \geq n_0$ 时恒定，值为
$H^m(\mathcal{F}^\bullet)$，引理得证。
\end{proof}







\section{有界上同调维数}
\label{sites-cohomology-section-bounded}

\noindent
本节讨论函子 $Rf_*$ 何时具有有界上同调维数。
当考虑无界复形时，这是一个相当微妙的问题。

\begin{situation}
\label{sites-cohomology-situation-olsson-laszlo}
设 $\mathcal{C}$ 为站点，设 $\mathcal{O}$ 为 $\mathcal{C}$ 上的层环。
设 $\mathcal{A} \subset \textit{Mod}(\mathcal{O})$ 为弱 Serre 子范畴。
我们假设以下条件成立：存在子集 $\mathcal{B} \subset \Ob(\mathcal{C})$，使得
\begin{enumerate}
\item $\mathcal{C}$ 的每个对象都有一个覆盖，其成员属于 $\mathcal{B}$；
\item 对每个 $V \in \mathcal{B}$，存在整数 $d_V$ 及 $V$ 的覆盖的共终系统
$\text{Cov}_V$，使得
$$
H^p(V_i, \mathcal{F}) = 0 \text{（对 }
\{V_i \to V\} \in \text{Cov}_V,\ p > d_V, \text{ 且 }
\mathcal{F} \in \Ob(\mathcal{A})
\text{）}
$$
\end{enumerate}
\end{situation}

\begin{lemma}
\label{sites-cohomology-lemma-olsson-laszlo}
\begin{reference}
这是带有略作修改假设的 \cite[Proposition 2.1.4]{six-I}；它是站点情形中
\cite[Proposition 3.13]{Spaltenstein} 的类似结果。
\end{reference}
在情形 \ref{sites-cohomology-situation-olsson-laszlo} 中，对任意
$E \in D_\mathcal{A}(\mathcal{O})$，《导出范畴》备注
\ref{derived-remark-map-into-derived-limit-truncations} 中的映射
$E \to R\lim \tau_{\geq -n} E$ 在 $D(\mathcal{O})$ 中是同构。
\end{lemma}

\begin{proof}
立即由引理 \ref{sites-cohomology-lemma-is-limit-spaltenstein} 得到。
\end{proof}

\begin{lemma}
\label{sites-cohomology-lemma-olsson-laszlo-modified}
在情形 \ref{sites-cohomology-situation-olsson-laszlo} 中，设
$(K_n)$ 为 $D_\mathcal{A}^+(\mathcal{O})$ 中的逆系统。
假设对每个 $j$，$\mathcal{A}$ 中的逆系统 $(H^j(K_n))$
最终恒定为 $\mathcal{H}^j$。则对所有 $j$，有
$H^j(R\lim K_n) = \mathcal{H}^j$。
\end{lemma}

\begin{proof}
设 $V \in \mathcal{B}$。设 $\{V_i \to V\}$ 属于情形
\ref{sites-cohomology-situation-olsson-laszlo} 中的集合 $\text{Cov}_V$。
由于 $K_n$ 下有界，有谱序列
$$
E_2^{p, q} = H^p(V_i, H^q(K_n))
$$
收敛到 $H^{p + q}(V_i, K_n)$。见《导出范畴》引理
\ref{derived-lemma-two-ss-complex-functor}。由假设，当 $p > d_V$ 时
$E_2^{p, q} = 0$。取 $n_0$ 使得
$$
\begin{matrix}
\mathcal{H}^{j + 1} & = & H^{j + 1}(K_n), \\
\mathcal{H}^j & = & H^j(K_n), \\
\ldots, \\
\mathcal{H}^{j - d_V - 2} & = & H^{j - d_V - 2}(K_n)
\end{matrix}
$$
对所有 $n \geq n_0$ 成立。比较上述关于 $K_n$ 和 $K_{n_0}$ 的谱序列，
可知当 $n \geq n_0$ 时，上同调群 $H^{j - 1}(V_i, K_n)$ 和 $H^j(V_i, K_n)$
与 $n$ 无关。因此对充分大的 $n$（取决于 $V$），截面上的映射
$H^j(R\lim K_n)(V) \to H^j(K_n)(V)$ 是单射，见引理
\ref{sites-cohomology-lemma-cohomology-derived-limit-injective}。
由于 $\mathcal{C}$ 的每个对象都可由 $\mathcal{B}$ 中的元素覆盖，故映射
$H^j(R\lim K_n) \to \mathcal{H}^j$ 是单射。

\medskip\noindent
满射性以类似方式证明。具体地，取 $U \in \Ob(\mathcal{C})$ 及
$\gamma \in \mathcal{H}^j(U)$。我们希望在用覆盖的成员替换 $U$ 后，
将 $\gamma$ 提升为 $H^j(R\lim K_n)$ 的截面。因此由情形
\ref{sites-cohomology-situation-olsson-laszlo} 的性质 (1)，可设 $U = V \in \mathcal{B}$。
取 $n_0$ 使得
$$
\begin{matrix}
\mathcal{H}^{j + 1} & = & H^{j + 1}(K_n), \\
\mathcal{H}^j & = & H^j(K_n), \\
\ldots, \\
\mathcal{H}^{j - d_V - 2} & = & H^{j - d_V - 2}(K_n)
\end{matrix}
$$
对所有 $n \geq n_0$ 成立。取 $\text{Cov}_V$ 中的覆盖
$\{V_i \to V\}$，使得
$\gamma|_{V_i} \in \mathcal{H}^j(V_i) = H^j(K_{n_0})(V_i)$
可提升为 $\gamma_{n_0, i} \in H^j(V_i, K_{n_0})$。
这是可能的，因为由引理 \ref{sites-cohomology-lemma-sheafification-cohomology}，
$H^j(K_{n_0})$ 是预层 $U \mapsto H^j(U, K_{n_0})$ 的层化。
由证明第一段的讨论，$H^{j - 1}(V_i, K_n)$ 和 $H^j(V_i, K_n)$
在 $n \geq n_0$ 时与 $n$ 无关。因此由引理
\ref{sites-cohomology-lemma-RGamma-commutes-with-Rlim}，$\gamma_{n_0, i}$ 可提升为
$\gamma_i \in H^j(V_i, R\lim K_n)$。证明完毕。
\end{proof}

\begin{lemma}
\label{sites-cohomology-lemma-olsson-laszlo-map-version-one}
\begin{reference}
这是带有略作修改假设的 \cite[Lemma 2.1.10]{six-I} 的一个版本。
\end{reference}
设 $f : (\Sh(\mathcal{C}), \mathcal{O}) \to (\Sh(\mathcal{C}'), \mathcal{O}')$
为环化拓扑斯态射。设 $\mathcal{A} \subset \textit{Mod}(\mathcal{O})$
和 $\mathcal{A}' \subset \textit{Mod}(\mathcal{O}')$ 为弱 Serre 子范畴。
假设存在整数 $N$ 使得
\begin{enumerate}
\item $\mathcal{C}, \mathcal{O}, \mathcal{A}$ 满足情形
\ref{sites-cohomology-situation-olsson-laszlo} 的假设；
\item $\mathcal{C}', \mathcal{O}', \mathcal{A}'$ 满足情形
\ref{sites-cohomology-situation-olsson-laszlo} 的假设；
\item 对 $p \geq 0$ 及 $\mathcal{F} \in \Ob(\mathcal{A})$，有
$R^pf_*\mathcal{F} \in \Ob(\mathcal{A}')$；
\item 对 $p > N$ 及 $\mathcal{F} \in \Ob(\mathcal{A})$，有
$R^pf_*\mathcal{F} = 0$；
\end{enumerate}
则对 $D_\mathcal{A}(\mathcal{O})$ 中的 $K$，有
\begin{enumerate}
\item[(a)] $Rf_*K$ 属于 $D_{\mathcal{A}'}(\mathcal{O}')$；
\item[(b)] 当 $j \geq N - n$ 时，映射
$H^j(Rf_*K) \to H^j(Rf_*(\tau_{\geq -n}K))$ 是同构。
\end{enumerate}
\end{lemma}

\begin{proof}
由引理 \ref{sites-cohomology-lemma-olsson-laszlo}，有 $K = R\lim \tau_{\geq -n}K$。
由引理 \ref{sites-cohomology-lemma-Rf-commutes-with-Rlim}，有
$Rf_*K = R\lim Rf_*\tau_{\geq -n}K$。
复形 $Rf_*\tau_{\geq -n}K$ 下有界。谱序列
$$
E_2^{p, q} = R^pf_*H^q(\tau_{\geq -n}K)
$$
收敛到 $H^{p + q}(Rf_*\tau_{\geq -n}K)$
（《导出范畴》引理 \ref{derived-lemma-two-ss-complex-functor}）。由假设 (3)，
$Rf_*\tau_{\geq -n}K$ 属于 $D^+_{\mathcal{A}'}(\mathcal{O}')$，
见《同调代数》引理 \ref{homology-lemma-biregular-ss-converges}。
注意，当 $m \geq n$ 时映射
$$
Rf_*(\tau_{\geq -m}K) \longrightarrow Rf_*(\tau_{\geq -n}K)
$$
由上述谱序列在 $j \geq -n + N$ 的次数上诱导上同调层的同构。因此可应用引理
\ref{sites-cohomology-lemma-olsson-laszlo-modified} 得到结论。
\end{proof}

\noindent
有时我们需要上述引理的一个变体，其中的假设略有不同。

\begin{situation}
\label{sites-cohomology-situation-olsson-laszlo-prime}
设
$f : (\mathcal{C}, \mathcal{O}) \to (\mathcal{C}', \mathcal{O}')$
为环化站点态射。设 $u : \mathcal{C}' \to \mathcal{C}$
为相应的站点连续函子。设 $\mathcal{A} \subset \textit{Mod}(\mathcal{O})$
为弱 Serre 子范畴。我们假设以下条件成立：存在子集
$\mathcal{B}' \subset \Ob(\mathcal{C}')$，使得
\begin{enumerate}
\item $\mathcal{C}'$ 的每个对象都有一个覆盖，其成员属于 $\mathcal{B}'$；
\item 对每个 $V' \in \mathcal{B}'$，存在整数 $d_{V'}$ 及 $V'$ 的覆盖的共终系统
$\text{Cov}_{V'}$，使得
$$
H^p(u(V'_i), \mathcal{F}) = 0 \text{（对 }
\{V'_i \to V'\} \in \text{Cov}_{V'},\ p > d_{V'}, \text{ 且 }
\mathcal{F} \in \Ob(\mathcal{A})\text{）}
$$
\end{enumerate}
\end{situation}

\begin{lemma}
\label{sites-cohomology-lemma-olsson-laszlo-map-version-two}
\begin{reference}
这是带有略作修改假设的 \cite[Lemma 2.1.10]{six-I} 的一个版本。
\end{reference}
设 $f : (\mathcal{C}, \mathcal{O}) \to (\mathcal{C}', \mathcal{O}')$
为环化站点态射。此外假设存在整数 $N$，使得
\begin{enumerate}
\item $\mathcal{C}, \mathcal{O}, \mathcal{A}$ 满足情形
\ref{sites-cohomology-situation-olsson-laszlo} 的假设；
\item $f : (\mathcal{C}, \mathcal{O}) \to (\mathcal{C}', \mathcal{O}')$
和 $\mathcal{A}$ 满足情形 \ref{sites-cohomology-situation-olsson-laszlo-prime} 的假设；
\item 对 $p > N$ 及 $\mathcal{F} \in \Ob(\mathcal{A})$，有
$R^pf_*\mathcal{F} = 0$；
\end{enumerate}
则对 $D_\mathcal{A}(\mathcal{O})$ 中的 $K$，当 $j \geq N - n$ 时，映射
$H^j(Rf_*K) \to H^j(Rf_*(\tau_{\geq -n}K))$ 是同构。
\end{lemma}

\begin{proof}
设 $K$ 属于 $D_\mathcal{A}(\mathcal{O})$。
根据引理 \ref{sites-cohomology-lemma-olsson-laszlo}，有 $K = R\lim \tau_{\geq -n}K$。
根据引理 \ref{sites-cohomology-lemma-Rf-commutes-with-Rlim}，有
$Rf_*K = R\lim Rf_*(\tau_{\geq -n}K)$。
取 $V' \in \mathcal{B}'$，并取 $\{V'_i \to V'\}$ 为
$\text{Cov}_{V'}$ 的一个元素。则有
$$
H^j(V'_i, Rf_*K) = H^j(u(V'_i), K)
\quad\text{且}\quad
H^j(V'_i, Rf_*(\tau_{\geq -n}K)) = H^j(u(V'_i), \tau_{\geq -n}K)
$$
情形 \ref{sites-cohomology-situation-olsson-laszlo-prime} 中的假设意味着，当 $n$ 足够大时，
最后一个群与 $n$ 无关；所需的 $n$ 只依赖于 $j$ 和 $d_{V'}$。细节略去。
对 $j$ 和 $j - 1$ 应用此结论，并通过引理
\ref{sites-cohomology-lemma-RGamma-commutes-with-Rlim} 得到
$$
H^j(V'_i, Rf_*K) = \lim H^j(V'_i, Rf_*(\tau_{\geq -n} K))
$$
且当 $n$ 大于某个仅依赖于 $d_{V'}$ 和 $j$ 的常数时，右侧系统为恒定系统。
因此引理 \ref{sites-cohomology-lemma-cohomology-derived-limit-injective} 蕴含
$$
H^j(Rf_*K)(V') \longrightarrow
\left(\lim H^j(Rf_*(\tau_{\geq -n}K))\right)(V')
$$
为单射。由于 $V' \in \mathcal{B}'$ 中的对象覆盖 $\mathcal{C}'$ 的每个对象，
我们得到映射
$H^j(Rf_*K) \to \lim H^j(Rf_*(\tau_{\geq -n}K))$
为单射。谱序列
$$
E_2^{p, q} = R^pf_*H^q(\tau_{\geq -n}K)
$$
收敛到 $H^{p + q}(Rf_*(\tau_{\geq -n}K))$
（导出范畴，引理 \ref{derived-lemma-two-ss-complex-functor}），
结合假设 (3) 可知
$H^j(Rf_*(\tau_{\geq -n}K))$ 在 $n \geq N - j$ 时为恒定群。
因此，当 $j \geq N - n$ 时，
$H^j(Rf_*K) \to H^j(Rf_*(\tau_{\geq -n}K))$ 为单射。

\medskip\noindent
于是我们已经证明，将引理最后一行的 ``isomorphism'' 替换为 ``injective''
后命题成立。现在取满足 $j \geq N - n$ 的 $j$ 和 $n$，考虑有区别三角
$$
\tau_{\leq -n - 1}K \to K \to \tau_{\geq -n}K \to (\tau_{\leq -n - 1}K)[1]
$$
参见导出范畴，注记 \ref{derived-remark-truncation-distinguished-triangle}。
由于 $\tau_{\geq -n}\tau_{\leq -n -1}K = 0$，对
$\tau_{-n - 1}K$ 已证明的单射性给出
$$
0 = H^j(Rf_*(\tau_{\leq -n - 1}K)) =
H^{j + 1}(Rf_*(\tau_{\leq -n - 1}K)) =
H^{j + 2}(Rf_*(\tau_{\leq -n - 1}K)) = \ldots
$$
由与该有区别三角相伴的上同调长正合序列
$$
Rf_*(\tau_{\leq -n - 1}K) \to Rf_*K \to Rf_*(\tau_{\geq -n}K) \to
Rf_*(\tau_{\leq -n - 1}K)[1]
$$
推出 $H^j(Rf_*K) \to H^j(Rf_*(\tau_{\geq -n}K))$
为同构。
\end{proof}








\section{Mayer-Vietoris}
\label{sites-cohomology-section-mayer-vietoris}

\noindent
通常的 Mayer--Vietoris 陈述及证明请参见上同调，第
\ref{cohomology-section-mayer-vietoris} 节。

\medskip\noindent
设 $(\mathcal{C}, \mathcal{O})$ 为环化站点。考虑范畴
$\mathcal{C}$ 中的交换图
$$
\xymatrix{
E \ar[d] \ar[r] & Y \ar[d] \\
Z \ar[r] & X
}
$$
在此情形，给定 $D(\mathcal{O})$ 中的对象 $K$，我们得到一个看似有区别三角
的开头
$$
R\Gamma(X, K) \to
R\Gamma(Z, K) \oplus
R\Gamma(Y, K) \to
R\Gamma(E, K)
$$
下面的引理将这一点说得更精确。

\begin{lemma}
\label{sites-cohomology-lemma-c-square}
在上述情形，选取表示 $K$ 的 $\mathcal{O}$-模 K-内射复形
$\mathcal{I}^\bullet$。将四个箭头之一的典范映射乘以 $-1$，得到复形映射
$$
\mathcal{I}^\bullet(X) \xrightarrow{\alpha}
\mathcal{I}^\bullet(Z) \oplus
\mathcal{I}^\bullet(Y) \xrightarrow{\beta}
\mathcal{I}^\bullet(E)
$$
且 $\beta \circ \alpha = 0$。于是有典范映射
$$
c^K_{X, Z, Y, E} :
\mathcal{I}^\bullet(X)
\longrightarrow
C(\beta)^\bullet[-1]
$$
所谓典范，是指表示 $K$ 的 K-内射复形取不同选择时，在阿贝尔群的导出范畴中
得到同构的箭头。若 $c^K_{X, Z, Y, E}$ 为同构，则利用其逆映射得到典范的
有区别三角
$$
R\Gamma(X, K) \to
R\Gamma(Z, K) \oplus
R\Gamma(Y, K) \to
R\Gamma(E, K) \to
R\Gamma(X, K)[1]
$$
所有这些构造关于 $K$ 都是函子性的。
\end{lemma}

\begin{proof}
此引理是自明的。例如，若 $\mathcal{J}^\bullet$ 是表示 $K$ 的另一个
K-内射复形，则可选取拟同构
$\mathcal{I}^\bullet \to \mathcal{J}^\bullet$
它在所涉及的全部复形之间确定拟同构。细节略去。关于锥的构造以及它与
有区别三角的关系，参见导出范畴，第
\ref{derived-section-cones} 节和
\ref{derived-section-homotopy-triangulated} 节。
\end{proof}

\begin{lemma}
\label{sites-cohomology-lemma-two-out-of-three-blow-up-square}
在上述情形，设
$K_1 \to K_2 \to K_3 \to K_1[1]$ 为 $D(\mathcal{O})$ 中的有区别三角。
若对 $\{1, 2, 3\}$ 中的两个 $i$，$c^{K_i}_{X, Z, Y, E}$ 是拟同构，
则对第三个 $i$ 它也是拟同构。
\end{lemma}

\begin{proof}
旋转该三角后，不妨设 $c^{K_1}_{X, Z, Y, E}$ 和
$c^{K_2}_{X, Z, Y, E}$ 为拟同构。选取表示 $K_1 \to K_2$ 的
$\mathcal{O}$-模 K-内射复形之间的映射
$f : \mathcal{I}^\bullet_1 \to \mathcal{I}^\bullet_2$。
则 $K_3$ 由 K-内射复形 $C(f)^\bullet$ 表示，参见
导出范畴，引理 \ref{derived-lemma-triangle-K-injective}。
于是态射 $c^{K_3}_{X, Z, Y, E}$ 是同构，因为它是
$K(\textit{Ab})$ 中有区别三角映射的第三条边，而另外两条边是拟同构。
细节略去；使用导出范畴，引理
\ref{derived-lemma-third-isomorphism-triangle}。
\end{proof}

\noindent
下面给出何时上述构造确实产生有区别三角的判据。

\begin{lemma}
\label{sites-cohomology-lemma-square-triangle}
在上述情形，假设
\begin{enumerate}
\item $h_X^\# = h_Y^\# \amalg_{h_E^\#} h_Z^\#$，且
\item $h_E^\# \to h_Y^\#$ 是单射。
\end{enumerate}
则引理 \ref{sites-cohomology-lemma-c-square} 的构造产生有区别三角
$$
R\Gamma(X, K) \to
R\Gamma(Z, K) \oplus
R\Gamma(Y, K) \to
R\Gamma(E, K) \to R\Gamma(X, K)[1]
$$
且对 $D(\mathcal{C})$ 中的 $K$ 是函子性的。
\end{lemma}

\begin{proof}
我们可以用各项均为内射阿贝尔层的 K-内射复形表示 $K$，参见
第 \ref{sites-cohomology-section-unbounded} 节。因此只需证明：若 $\mathcal{I}$ 是内射阿贝尔
层，则
$$
0 \to \mathcal{I}(X) \to
\mathcal{I}(Z) \oplus \mathcal{I}(Y) \to
\mathcal{I}(E) \to 0
$$
是短正合序列。第一个箭头是单射，因为由条件 (1)，映射
$h_Y \amalg h_Z \to h_X$ 层化后成为满射；这意味着
$\{Y \to X, Z \to X\}$ 可以由 $X$ 的一个覆盖细化。
最后一个箭头是满射，因为 $\mathcal{I}(Y) \to \mathcal{I}(E)$ 是满射。
具体地，有
$\mathcal{I}(E) = \Hom(\mathbf{Z}_E^\#, \mathcal{I})$,
$\mathcal{I}(Y) = \Hom(\mathbf{Z}_Y^\#, \mathcal{I})$,
由 (2)，映射 $\mathbf{Z}_E^\# \to \mathbf{Z}_Y^\#$ 是单射，且
$\mathcal{I}$ 是内射阿贝尔层。请与《站点上的模》，第
\ref{sites-modules-section-free-abelian-sheaf} 节相比较。
最后，设 $s \in \mathcal{I}(Y)$ 且
$t \in \mathcal{F}(Z)$ 映射到 $\mathcal{I}(E)$ 中的同一元素。
则 $s$ 和 $t$ 定义映射
$$
s \amalg t : h_Y^\# \amalg h_Z^\# \longrightarrow \mathcal{I}
$$
该映射按假设可分解经过 $h_Y^\# \amalg_{h_E^\#} h_Z^\#$。
于是由假设 (1) 得到唯一映射
$h_X^\# \to \mathcal{I}$，它对应于 $\mathcal{I}(X)$ 中的一个元素，
其在 $Y$ 上限制为 $s$，在 $Z$ 上限制为 $t$。
\end{proof}

\begin{lemma}
\label{sites-cohomology-lemma-square-triangle-general}
设 $\mathcal{C}$ 为站点。考虑交换图
$$
\xymatrix{
\mathcal{D} \ar[r] \ar[d] & \mathcal{F} \ar[d] \\
\mathcal{E} \ar[r] & \mathcal{G}
}
$$
其由 $\mathcal{C}$ 上的集合预层组成，并假设
\begin{enumerate}
\item $\mathcal{G}^\# =
\mathcal{E}^\# \amalg_{\mathcal{D}^\#} \mathcal{F}^\#$，且
\item $\mathcal{D}^\# \to \mathcal{F}^\#$ 是单射。
\end{enumerate}
则存在典范有区别三角
$$
R\Gamma(\mathcal{G}, K) \to
R\Gamma(\mathcal{E}, K) \oplus
R\Gamma(\mathcal{F}, K) \to
R\Gamma(\mathcal{D}, K) \to
R\Gamma(\mathcal{G}, K)[1]
$$
它关于 $K \in D(\mathcal{C})$ 是函子性的，其中 $R\Gamma(\mathcal{G}, -)$
是第 \ref{sites-cohomology-section-limp} 节所讨论的上同调。
\end{lemma}

\begin{proof}
由于层化是正合的，且
$R\Gamma(\mathcal{G}, -) = R\Gamma(\mathcal{G}^\#, -)$，
不妨假设 $\mathcal{D}, \mathcal{E}, \mathcal{F}, \mathcal{G}$ 是集合层。
此外，上同调 $R\Gamma(\mathcal{G}, -)$ 只依赖于拓扑斯，而不依赖于底层站点。
因此由
站点，引理 \ref{sites-lemma-topos-good-site}
可以将 $\mathcal{C}$ 替换为具有次典范拓扑的一个较大站点，
使得 $\mathcal{G} = h_X$、$\mathcal{F} = h_Y$、$\mathcal{E} = h_Z$ 且
$\mathcal{D} = h_E$，其中 $X, Y, Z, E$ 是 $\mathcal{C}$ 的对象。
此时结论由引理 \ref{sites-cohomology-lemma-square-triangle} 得出。
\end{proof}







\section{比较两个拓扑}
\label{sites-cohomology-section-compare}

\noindent
设 $\mathcal{C}$ 为范畴。令
$\text{Cov}(\mathcal{C}) \supset \text{Cov}'(\mathcal{C})$
为两种赋予 $\mathcal{C}$ 站点结构的方式。记 $\tau$ 为对应于
$\text{Cov}(\mathcal{C})$ 的拓扑，记 $\tau'$ 为对应于
$\text{Cov}'(\mathcal{C})$ 的拓扑。于是 $\mathcal{C}$ 上的恒等函子定义站点态射
$$
\epsilon : \mathcal{C}_\tau \longrightarrow \mathcal{C}_{\tau'}
$$
其中 $\epsilon_*$ 是底层预层上的恒等函子，而 $\epsilon^{-1}$ 是
$\tau'$-层的 $\tau$-层化。参见站点，例
\ref{sites-example-finer-topology} 和
\ref{sites-example-finer-topology-bis}。
在上述情形中有以下结论：
\begin{enumerate}
\item $\epsilon_* : \Sh(\mathcal{C}_\tau) \to \Sh(\mathcal{C}_{\tau'})$
是完全忠实的，且 $\epsilon^{-1} \circ \epsilon_* = \text{id}$；
\item $\epsilon_* : \textit{Ab}(\mathcal{C}_\tau) \to
\textit{Ab}(\mathcal{C}_{\tau'})$
是完全忠实的，且 $\epsilon^{-1} \circ \epsilon_* = \text{id}$；
\item $R\epsilon_* : D(\mathcal{C}_\tau) \to D(\mathcal{C}_{\tau'})$
是完全忠实的，且 $\epsilon^{-1} \circ R\epsilon_* = \text{id}$；
\item 若 $\mathcal{O}$ 是 $\tau$-拓扑下的环层，则 $\mathcal{O}$ 也是
$\tau'$-拓扑下的层，且 $\epsilon$ 成为环化站点的平坦态射
$$
\epsilon :
(\mathcal{C}_\tau, \mathcal{O}_\tau)
\longrightarrow
(\mathcal{C}_{\tau'}, \mathcal{O}_{\tau'})
$$
\item $\epsilon_* : \textit{Mod}(\mathcal{O}_\tau) \to
\textit{Mod}(\mathcal{O}_{\tau'})$
是完全忠实的，且 $\epsilon^* \circ \epsilon_* = \text{id}$；
\item $R\epsilon_* : D(\mathcal{O}_\tau) \to D(\mathcal{O}_{\tau'})$
是完全忠实的，且 $\epsilon^* \circ R\epsilon_* = \text{id}$。
\end{enumerate}
说明如下。

\medskip\noindent
(1)。设 $\mathcal{F}$ 是 $\tau$-拓扑下的集合层。
那么 $\epsilon_*\mathcal{F}$ 只是将 $\mathcal{F}$ 看作 $\tau'$-拓扑下的层。
应用 $\epsilon^{-1}$ 意味着对 $\mathcal{F}$ 作 $\tau$-层化；由于
$\mathcal{F}$ 已经是 $\tau$-层，这一操作不产生任何改变。因此
$\epsilon^{-1}(\epsilon_*\mathcal{F})) = \mathcal{F}$.
完全忠实性由范畴，引理
\ref{categories-lemma-adjoint-fully-faithful}.

\medskip\noindent
(2)。这是 (1) 的推论，因为阿贝尔层的拉回和推前与对底层集合层
进行相同操作是一回事。

\medskip\noindent
(3)。设 $K$ 是 $D(\mathcal{C}_\tau)$ 中的对象。为计算
$R\epsilon_*K$，选取表示 $K$ 的 K-内射复形
$\mathcal{I}^\bullet$，并令 $R\epsilon_*K = \epsilon_*\mathcal{I}^\bullet$。
由于 $\epsilon^{-1} : D(\mathcal{C}_{\tau'}) \to D(\mathcal{C}_\tau)$
对对象 $L$ 的计算方式是对表示 $L$ 的任意阿贝尔层复形应用正合函子
$\epsilon^{-1}$，故 $\epsilon^{-1}R\epsilon_*K$ 由
$\epsilon^{-1}\epsilon_*\mathcal{I}^\bullet$ 表示。由
(1) 有
$\mathcal{I}^\bullet = \epsilon^{-1}\epsilon_*\mathcal{I}^\bullet$.
换言之，$\epsilon^{-1} \circ R\epsilon_* = \text{id}$，结论如前。

\medskip\noindent
(4)。注意 $\epsilon^{-1}\mathcal{O}_{\tau'} = \mathcal{O}_\tau$，参见
(1) 中的讨论。因此 $\epsilon$ 是环化站点的平坦态射，参见
《站点上的模》，定义 \ref{sites-modules-definition-flat-morphism}。
不仅如此，显然在 $\mathcal{O}_{\tau'}$-模上有
$\epsilon^* = \epsilon^{-1}$（作为模的拉回与底层阿贝尔层的拉回具有相同
的底层阿贝尔层）。

\medskip\noindent
(5)。这由 (2) 及 (4) 中的论述立即得到。

\medskip\noindent
(6)。这与 (3) 类似，细节略去。












\section{上同调下降的形式性}
\label{sites-cohomology-section-formal-cohomological-descent}

\noindent
本节只讨论环化拓扑斯的态射在多大程度上确定从下方导出范畴
到上方导出范畴的嵌入。下面是一个典型结果。

\begin{lemma}
\label{sites-cohomology-lemma-downstairs}
设 $f : (\Sh(\mathcal{C}), \mathcal{O}_\mathcal{C}) \to
(\Sh(\mathcal{D}), \mathcal{O}_\mathcal{D})$ 为环化拓扑斯的态射。
考虑 $D(\mathcal{O}_\mathcal{D})$ 的满子范畴
$D' \subset D(\mathcal{O}_\mathcal{D})$，其对象 $K$ 满足
$$
K \longrightarrow Rf_*Lf^*K
$$
是同构。则 $D'$ 是 $D(\mathcal{O}_\mathcal{D})$ 的饱和三角严格满子范畴，
且函子 $Lf^* : D' \to D(\mathcal{O}_\mathcal{C})$ 完全忠实。
\end{lemma}

\begin{proof}
关于本情形中“饱和”的定义，参见导出范畴，定义
\ref{derived-definition-saturated}。关于三角子范畴，参见
导出范畴，引理 \ref{derived-lemma-triangulated-subcategory}。
引理中的典范映射是伴随函子对 $(Lf^*, Rf_*)$ 的单位，参见引理
\ref{sites-cohomology-lemma-adjoint}。有了这一点，$D'$ 为饱和三角子范畴的证明略去；
它形式上来自 $Lf^*$ 和 $Rf_*$ 是三角范畴之间的正合函子这一事实。
最后一部分形式上来自 $Lf^*$ 与 $Rf_*$ 互为伴随；可与范畴，引理
\ref{categories-lemma-adjoint-fully-faithful} 比较。
\end{proof}

\begin{lemma}
\label{sites-cohomology-lemma-upstairs}
设 $f : (\Sh(\mathcal{C}), \mathcal{O}_\mathcal{C}) \to
(\Sh(\mathcal{D}), \mathcal{O}_\mathcal{D})$ 为环化拓扑斯的态射。
考虑 $D(\mathcal{O}_\mathcal{C})$ 的满子范畴
$D' \subset D(\mathcal{O}_\mathcal{C})$，其对象 $K$ 满足
$$
Lf^*Rf_*K \longrightarrow K
$$
是同构。则 $D'$ 是 $D(\mathcal{O}_\mathcal{C})$ 的饱和三角严格满子范畴，
且函子 $Rf_* : D' \to D(\mathcal{O}_\mathcal{D})$ 完全忠实。
\end{lemma}

\begin{proof}
关于本情形中“饱和”的定义，参见导出范畴，定义
\ref{derived-definition-saturated}。关于三角子范畴，参见
导出范畴，引理 \ref{derived-lemma-triangulated-subcategory}。
引理中的典范映射是伴随函子对 $(Lf^*, Rf_*)$ 的余单位，参见引理
\ref{sites-cohomology-lemma-adjoint}。有了这一点，$D'$ 为饱和三角子范畴的证明略去；
它形式上来自 $Lf^*$ 和 $Rf_*$ 是三角范畴之间的正合函子这一事实。
最后一部分形式上来自 $Lf^*$ 与 $Rf_*$ 互为伴随；可与范畴，引理
\ref{categories-lemma-adjoint-fully-faithful} 比较。
\end{proof}

\begin{lemma}
\label{sites-cohomology-lemma-bounded-in-image-upstairs}
设 $f : (\Sh(\mathcal{C}), \mathcal{O}_\mathcal{C}) \to
(\Sh(\mathcal{D}), \mathcal{O}_\mathcal{D})$ 为环化拓扑斯的态射。
设 $K$ 是 $D(\mathcal{O}_\mathcal{C})$ 中的对象。假设
\begin{enumerate}
\item $f$ 是平坦的；
\item $K$ 下有界；
\item $f^*Rf_*H^q(K) \to H^q(K)$ 是同构。
\end{enumerate}
则 $f^*Rf_*K \to K$ 是同构。
\end{lemma}

\begin{proof}
注意，$f^*Rf_*K \to K$ 是同构，当且仅当它在上同调层 $H^j$ 上是同构。
注意
$H^j(f^*Rf_*K) = f^*H^j(Rf_*K) = f^*H^j(Rf_*\tau_{\leq j}K) =
H^j(f^*Rf_*\tau_{\leq j}K)$.
因此不妨假设 $K$ 有界。于是性质 (3) 告诉我们，上同调层属于
引理 \ref{sites-cohomology-lemma-upstairs} 中的三角子范畴
$D' \subset D(\mathcal{O}_\mathcal{C})$。因此 $K$ 也属于该子范畴。
\end{proof}

\begin{lemma}
\label{sites-cohomology-lemma-bounded-in-image-downstairs}
设 $f : (\Sh(\mathcal{C}), \mathcal{O}_\mathcal{C}) \to
(\Sh(\mathcal{D}), \mathcal{O}_\mathcal{D})$ 为环化拓扑斯的态射。
设 $K$ 是 $D(\mathcal{O}_\mathcal{D})$ 中的对象。假设
\begin{enumerate}
\item $f$ 是平坦的；
\item $K$ 下有界；
\item $H^q(K) \to Rf_*f^*H^q(K)$ 是同构。
\end{enumerate}
则 $K \to Rf_*f^*K$ 是同构。
\end{lemma}

\begin{proof}
注意，$K \to Rf_*f^*K$ 是同构，当且仅当它在上同调层 $H^j$ 上是同构。
注意
$H^j(Rf_*f^*K) = H^j(Rf_*\tau_{\leq j}f^*K) = H^j(Rf_*f^*\tau_{\leq j}K)$.
因此不妨假设 $K$ 有界。于是性质 (3) 告诉我们，上同调层属于
引理 \ref{sites-cohomology-lemma-downstairs} 中的三角子范畴
$D' \subset D(\mathcal{O}_\mathcal{D})$。因此 $K$ 也属于该子范畴。
\end{proof}

\begin{lemma}
\label{sites-cohomology-lemma-equivalence-bounded}
设 $f : (\Sh(\mathcal{C}), \mathcal{O}) \to (\Sh(\mathcal{C}'), \mathcal{O}')$
为环化拓扑斯的态射。设 $\mathcal{A} \subset \textit{Mod}(\mathcal{O})$
及 $\mathcal{A}' \subset \textit{Mod}(\mathcal{O}')$ 为弱 Serre 子范畴。假设
\begin{enumerate}
\item $f$ 是平坦的；
\item $f^*$ 诱导范畴等价
$\mathcal{A}' \to \mathcal{A}$；
\item 对 $\mathcal{F}' \in \Ob(\mathcal{A}')$，映射
$\mathcal{F}' \to Rf_*f^*\mathcal{F}'$ 是同构。
\end{enumerate}
则
$f^* : D_{\mathcal{A}'}^+(\mathcal{O}') \to D_\mathcal{A}^+(\mathcal{O})$
是范畴等价，其拟逆为
$Rf_* : D_\mathcal{A}^+(\mathcal{O}) \to D_{\mathcal{A}'}^+(\mathcal{O}')$.
\end{lemma}

\begin{proof}
由假设 (2)、(3) 以及引理
\ref{sites-cohomology-lemma-bounded-in-image-downstairs} 和 \ref{sites-cohomology-lemma-downstairs}，可知
$f^* : D_{\mathcal{A}'}^+(\mathcal{O}') \to D_\mathcal{A}^+(\mathcal{O})$
是完全忠实的。对 $\mathcal{F} \in \Ob(\mathcal{A})$，可写成
$\mathcal{F} = f^*\mathcal{F}'$。于是
$Rf_*\mathcal{F} = Rf_* f^*\mathcal{F}' = \mathcal{F}'$.
特别地，当 $p > 0$ 时 $R^pf_*\mathcal{F} = 0$，且
$f_*\mathcal{F} \in \Ob(\mathcal{A}')$。因此对任意
$K \in D^+_\mathcal{A}(\mathcal{O})$，利用收敛到 $R^{p + q}f_*K$ 的谱序列
$E_2^{p, q} = R^pf_*H^q(K)$，可知
$Rf_*K \in D^+_{\mathcal{A}'}(\mathcal{O}')$。
当然，由引理
\ref{sites-cohomology-lemma-bounded-in-image-upstairs} 和 \ref{sites-cohomology-lemma-upstairs} 可知
$Rf_* : D_\mathcal{A}^+(\mathcal{O}) \to D_{\mathcal{A}'}^+(\mathcal{O}')$
是完全忠实的。由于 $f^*$ 与 $Rf_*$ 互为伴随，遂得到引理结论；
例如可用范畴，引理 \ref{categories-lemma-adjoint-fully-faithful}。
\end{proof}

\begin{lemma}
\label{sites-cohomology-lemma-equivalence-unbounded-one}
\begin{reference}
这与 \cite[Theorem 2.2.3]{six-I} 类似。
\end{reference}
设 $f : (\Sh(\mathcal{C}), \mathcal{O}) \to (\Sh(\mathcal{C}'), \mathcal{O}')$
为环化拓扑斯的态射。设 $\mathcal{A} \subset \textit{Mod}(\mathcal{O})$
及 $\mathcal{A}' \subset \textit{Mod}(\mathcal{O}')$ 为弱 Serre 子范畴。假设
\begin{enumerate}
\item $f$ 是平坦的；
\item $f^*$ 诱导范畴等价 $\mathcal{A}' \to \mathcal{A}$；
\item 对 $\mathcal{F}' \in \Ob(\mathcal{A}')$，映射
$\mathcal{F}' \to Rf_*f^*\mathcal{F}'$ 是同构；
\item $\mathcal{C}, \mathcal{O}, \mathcal{A}$ 满足情形
\ref{sites-cohomology-situation-olsson-laszlo} 的假设；
\item $\mathcal{C}', \mathcal{O}', \mathcal{A}'$ 满足情形
\ref{sites-cohomology-situation-olsson-laszlo} 的假设。
\end{enumerate}
则 $f^* : D_{\mathcal{A}'}(\mathcal{O}') \to D_\mathcal{A}(\mathcal{O})$
是范畴等价，其拟逆为
$Rf_* : D_\mathcal{A}(\mathcal{O}) \to D_{\mathcal{A}'}(\mathcal{O}')$.
\end{lemma}

\begin{proof}
由于 $f^*$ 是正合的，显然 $f^*$ 定义函子
$f^* : D_{\mathcal{A}'}(\mathcal{O}') \to D_\mathcal{A}(\mathcal{O})$
如引理所述，且该函子还与截断函子 $\tau_{\geq -n}$ 交换。
我们已经知道，在相应的下有界范畴上，$f^*$ 与 $Rf_*$ 是互为拟逆的
等价，参见引理 \ref{sites-cohomology-lemma-equivalence-bounded}。
由引理 \ref{sites-cohomology-lemma-olsson-laszlo-map-version-one}（取 $N = 0$），可知
$Rf_*$ 确实定义函子
$Rf_* : D_\mathcal{A}(\mathcal{O}) \to D_{\mathcal{A}'}(\mathcal{O}')$
且该函子还与截断函子 $\tau_{\geq -n}$ 交换。
因此对 $D_\mathcal{A}(\mathcal{O})$ 中的 $K$，映射
$f^*Rf_*K \to K$ 是同构，因为它在各截断上是同构。
类似地，对 $D_{\mathcal{A}'}(\mathcal{O}')$ 中的 $K'$，映射
$K' \to Rf_*f^*K'$ 是同构，因为它在各截断上是同构。
证明完毕。
\end{proof}

\begin{lemma}
\label{sites-cohomology-lemma-equivalence-unbounded-two}
\begin{reference}
这与 \cite[Theorem 2.2.3]{six-I} 类似。
\end{reference}
设 $f : (\mathcal{C}, \mathcal{O}) \to (\mathcal{C}', \mathcal{O}')$
为环化站点态射。设 $\mathcal{A} \subset \textit{Mod}(\mathcal{O})$
及 $\mathcal{A}' \subset \textit{Mod}(\mathcal{O}')$ 为弱 Serre 子范畴。假设
\begin{enumerate}
\item $f$ 是平坦的；
\item $f^*$ 诱导范畴等价 $\mathcal{A}' \to \mathcal{A}$；
\item 对 $\mathcal{F}' \in \Ob(\mathcal{A}')$，映射
$\mathcal{F}' \to Rf_*f^*\mathcal{F}'$ 是同构；
\item $\mathcal{C}, \mathcal{O}, \mathcal{A}$ 满足情形
\ref{sites-cohomology-situation-olsson-laszlo} 的假设；
\item $f : (\mathcal{C}, \mathcal{O}) \to (\mathcal{C}', \mathcal{O}')$
和 $\mathcal{A}$ 满足情形
\ref{sites-cohomology-situation-olsson-laszlo-prime} 的假设。
\end{enumerate}
则 $f^* : D_{\mathcal{A}'}(\mathcal{O}') \to D_\mathcal{A}(\mathcal{O})$
是范畴等价，其拟逆为
$Rf_* : D_\mathcal{A}(\mathcal{O}) \to D_{\mathcal{A}'}(\mathcal{O}')$.
\end{lemma}

\begin{proof}
本引理的证明与引理 \ref{sites-cohomology-lemma-equivalence-unbounded-one} 的证明完全相同，
区别仅在于将对引理 \ref{sites-cohomology-lemma-olsson-laszlo-map-version-one} 的引用
替换为对
Lemma \ref{sites-cohomology-lemma-olsson-laszlo-map-version-two}.
\end{proof}








\section{比较两个拓扑，II}
\label{sites-cohomology-section-compare-II}

\noindent
设 $\mathcal{C}$ 为范畴。令
$\text{Cov}(\mathcal{C}) \supset \text{Cov}'(\mathcal{C})$
为两种赋予 $\mathcal{C}$ 站点结构的方式。记 $\tau$ 为对应于
$\text{Cov}(\mathcal{C})$ 的拓扑，记 $\tau'$ 为对应于
$\text{Cov}'(\mathcal{C})$ 的拓扑。于是 $\mathcal{C}$ 上的恒等函子定义站点态射
$$
\epsilon : \mathcal{C}_\tau \longrightarrow \mathcal{C}_{\tau'}
$$
其中 $\epsilon_*$ 是底层预层上的恒等函子，而 $\epsilon^{-1}$ 是
$\tau'$-层的 $\tau$-层化（因而显然正合）。设 $\mathcal{O}$ 是
$\tau$-拓扑下的环层。则 $\mathcal{O}$ 也是 $\tau'$-拓扑下的层，
且 $\epsilon$ 成为环化站点态射
$$
\epsilon :
(\mathcal{C}_\tau, \mathcal{O}_\tau)
\longrightarrow
(\mathcal{C}_{\tau'}, \mathcal{O}_{\tau'})
$$
更多讨论参见第 \ref{sites-cohomology-section-compare} 节。

\begin{lemma}
\label{sites-cohomology-lemma-compare-topologies-derived-adequate-modules}
设 $\epsilon : (\mathcal{C}_\tau, \mathcal{O}_\tau) \to
(\mathcal{C}_{\tau'}, \mathcal{O}_{\tau'})$ 如上。令
$\mathcal{B} \subset \Ob(\mathcal{C})$ 为子集。令
$\mathcal{A} \subset \textit{PMod}(\mathcal{O})$ 为满子范畴。假设
\begin{enumerate}
\item $\mathcal{A}$ 的每个对象都是 $\tau$-拓扑下的层；
\item $\mathcal{A}$ 是 $\textit{Mod}(\mathcal{O}_\tau)$ 的弱 Serre 子范畴；
\item $\mathcal{C}$ 的每个对象都有一个 $\tau'$-覆盖，其成员属于
$\mathcal{B}$；
\item 对每个 $U \in \mathcal{B}$，有
$H^p_\tau(U, \mathcal{F}) = 0$（对所有 $\mathcal{F} \in \mathcal{A}$ 及 $p > 0$）。
\end{enumerate}
则 $\mathcal{A}$ 是
$\textit{Mod}(\mathcal{O}_{\tau'})$ 的弱 Serre 子范畴，且存在三角范畴等价
$D_\mathcal{A}(\mathcal{O}_\tau) = D_\mathcal{A}(\mathcal{O}_{\tau'})$
由 $\epsilon^*$ 和 $R\epsilon_*$ 给出。
\end{lemma}

\begin{proof}
由于 $\epsilon^{-1}\mathcal{O}_{\tau'} = \mathcal{O}_\tau$，
可知 $\epsilon$ 是环化站点的平坦态射；事实上，在模层上
$\epsilon^{-1} = \epsilon^*$。由性质 (1)，可将 $\mathcal{A}$ 的每个对象
视为 $\mathcal{O}_\tau$-模层，也视为 $\mathcal{O}_{\tau'}$-模层。
换言之，有完全忠实的包含函子
$$
\mathcal{A} \to \textit{Mod}(\mathcal{O}_\tau) \to
\textit{Mod}(\mathcal{O}_{\tau'})
$$
为避免混淆，记 $\mathcal{A}' \subset \textit{Mod}(\mathcal{O}_{\tau'})$
为 $\mathcal{A}$ 的像。显然，$\epsilon_* : \mathcal{A} \to \mathcal{A}'$
和 $\epsilon^* : \mathcal{A}' \to \mathcal{A}$ 互为拟逆等价（参见引理前的
讨论，并注意 $\mathcal{A}'$ 的对象是 $\tau$-拓扑下的层）。

\medskip\noindent
条件 (3) 和 (4) 蕴含，当 $p > 0$ 且 $\mathcal{F} \in \Ob(\mathcal{A})$ 时，
$R^p\epsilon_*\mathcal{F} = 0$。这是因为 $R^p\epsilon_*$ 是预层
$U \mapsto H^p_\tau(U, \mathcal{F})$ 所关联的层，参见引理
\ref{sites-cohomology-lemma-higher-direct-images}。因此，对 $\mathcal{A}$ 中的任意正合复形
（这等价于 $\textit{Mod}(\mathcal{O}_\tau)$ 中各项属于 $\mathcal{A}$ 的正合复形，
参见同调，引理 \ref{homology-lemma-characterize-weak-serre-subcategory}），
应用函子 $\epsilon_*$ 后仍为正合。

\medskip\noindent
考虑正合序列
$$
\mathcal{F}'_0 \to \mathcal{F}'_1 \to
\mathcal{F}'_2 \to \mathcal{F}'_3 \to
\mathcal{F}'_4
$$
它位于 $\textit{Mod}(\mathcal{O}_{\tau'})$ 中，且
$\mathcal{F}'_0, \mathcal{F}'_1, \mathcal{F}'_3, \mathcal{F}'_4$ 属于
$\mathcal{A}'$。应用正合函子 $\epsilon^*$ 得到正合序列
$$
\epsilon^*\mathcal{F}'_0 \to \epsilon^*\mathcal{F}'_1 \to
\epsilon^*\mathcal{F}'_2 \to \epsilon^*\mathcal{F}'_3 \to
\epsilon^*\mathcal{F}'_4
$$
它位于 $\textit{Mod}(\mathcal{O}_\tau)$ 中。由于 $\mathcal{A}$ 是
弱 Serre 子范畴，且
$\epsilon^*\mathcal{F}'_0, \epsilon^*\mathcal{F}'_1,
\epsilon^*\mathcal{F}'_3, \epsilon^*\mathcal{F}'_4$ 属于
$\mathcal{A}$。由同调（参见定义
\ref{homology-definition-serre-subcategory}），得到
$\epsilon^*\mathcal{F}_2 \in \mathcal{A}$。考虑序列映射
$$
\xymatrix{
\mathcal{F}'_0 \ar[r] \ar[d] &
\mathcal{F}'_1 \ar[r] \ar[d] &
\mathcal{F}'_2 \ar[r] \ar[d] &
\mathcal{F}'_3 \ar[r] \ar[d] &
\mathcal{F}'_4 \ar[d] \\
\epsilon_*\epsilon^*\mathcal{F}'_0 \ar[r] &
\epsilon_*\epsilon^*\mathcal{F}'_1 \ar[r] &
\epsilon_*\epsilon^*\mathcal{F}'_2 \ar[r] &
\epsilon_*\epsilon^*\mathcal{F}'_3 \ar[r] &
\epsilon_*\epsilon^*\mathcal{F}'_4
}
$$
由前一段讨论，下行正合。指标为 $0$、$1$、$3$、$4$ 的竖直箭头
由第一段讨论是同构。由 $5$ 引理（同调，引理
\ref{homology-lemma-five-lemma}）可得
$\mathcal{F}'_2 \cong \epsilon_*\epsilon^*\mathcal{F}'_2$，
从而 $\mathcal{F}'_2 \in \mathcal{A}'$。这样便知
$\mathcal{A}'$ 是 $\textit{Mod}(\mathcal{O}_{\tau'})$ 的弱 Serre 子范畴，
参见同调，定义 \ref{homology-definition-serre-subcategory}。

\medskip\noindent
此时讨论导出范畴
$D_\mathcal{A}(\mathcal{O}_\tau)$ 和
$D_{\mathcal{A}'}(\mathcal{O}_{\tau'})$ 才有意义，参见导出范畴，
第 \ref{derived-section-triangulated-sub} 节。为完成证明，我们说明引理
\ref{sites-cohomology-lemma-equivalence-unbounded-two} 的条件 (1)--(5) 得到满足。
上面已经验证了 (1)、(2)、(3)。注意，由于每个对象都有由
$\mathcal{B}$ 中对象组成的 $\tau'$-覆盖，当然每个对象也有由
$\mathcal{B}$ 中对象组成的 $\tau$-覆盖。因此引理
\ref{sites-cohomology-lemma-equivalence-unbounded-two} 的条件 (4) 得到满足。
同理，条件 (5) 也得到满足。
\end{proof}

\begin{lemma}
\label{sites-cohomology-lemma-descent-squares}
设 $\epsilon : (\mathcal{C}_\tau, \mathcal{O}_\tau) \to
(\mathcal{C}_{\tau'}, \mathcal{O}_{\tau'})$ 如上。令 $A$ 为集合，并对
$\alpha \in A$ 令
$$
\xymatrix{
E_\alpha \ar[d] \ar[r] & Y_\alpha \ar[d] \\
Z_\alpha \ar[r] & X_\alpha
}
$$
为范畴 $\mathcal{C}$ 中的交换图。假设
\begin{enumerate}
\item $\tau'$-层 $\mathcal{F}'$ 是 $\tau$-层，当且仅当对所有 $\alpha$ 有
$\mathcal{F}'(X_\alpha) =
\mathcal{F}'(Z_\alpha) \times_{\mathcal{F}'(E_\alpha)} 
\mathcal{F}'(Y_\alpha)$；
\item 对于 $D(\mathcal{O}_{\tau'})$ 中属于 $R\epsilon_*$ 本质像的 $K'$，
引理 \ref{sites-cohomology-lemma-c-square} 的映射
$c^{K'}_{X_\alpha, Z_\alpha, Y_\alpha, E_\alpha}$
对所有 $\alpha$ 都是同构。
\end{enumerate}
则对 $K' \in D^+(\mathcal{O}_{\tau'})$，它属于 $R\epsilon_*$ 的本质像，
当且仅当映射 $c^{K'}_{X_\alpha, Z_\alpha, Y_\alpha, E_\alpha}$
对所有 $\alpha$ 都是同构。
\end{lemma}

\begin{proof}
“仅当”方向由假设 (2) 蕴含。
另一方面，若 $K'$ 只有一个非零上同调层，
则“若”方向由假设 (1) 得出。
一般地，我们将用归纳法证明“若”方向。称
$D^+(\mathcal{O}_{\tau'})$ 中的对象 $K'$ 满足 (P)，如果映射
满足 (P)，如果映射 $c^{K'}_{X_\alpha, Z_\alpha, Y_\alpha, E_\alpha}$
对所有 $\alpha \in A$ 都是同构。

\medskip\noindent
具体地，设 $K'$ 是满足 (P) 的 $D^+(\mathcal{O}_{\tau'})$ 中对象。
选取有区别三角
$$
K' \to R\epsilon_*\epsilon^{-1}K' \to M' \to K'[1]
$$
它位于 $D^+(\mathcal{O}_{\tau'})$ 中，其中第一个箭头是伴随映射。
由 (2) 和引理 \ref{sites-cohomology-lemma-two-out-of-three-blow-up-square} 可知 $M'$ 满足 (P)。另一方面，应用

$\epsilon^{-1}$，并利用第 \ref{sites-cohomology-section-compare} 节中
$\epsilon^{-1}R\epsilon_* = \text{id}$，得到 $\epsilon^{-1}M' = 0$。
下一段将证明 $M' = 0$，从而完成证明。

\medskip\noindent
设 $K'$ 是满足 (P) 且 $\epsilon^{-1}K' = 0$ 的 $D^+(\mathcal{O}_{\tau'})$ 中对象。我们将证明 $K' = 0$。

具体地，给定 $n \in \mathbf{Z}$，若 $i < n$ 时 $H^i(K') = 0$，
我们将证明 $H^n(K') = 0$。对 $\alpha \in A$，由引理 \ref{sites-cohomology-lemma-c-square} 有有区别三角

$$
R\Gamma_{\tau'}(X_\alpha, K') \to
R\Gamma_{\tau'}(Z_\alpha, K') \oplus
R\Gamma_{\tau'}(Y_\alpha, K') \to
R\Gamma_{\tau'}(E_\alpha, K') \to
R\Gamma_{\tau'}(X_\alpha, K')[1]
$$
取次数 $n$ 的上同调，并利用已假设的 $K'$ 上同调层消失，得到正合序列


$$
0 \to
H^n_{\tau'}(X_\alpha, K') \to
H^n_{\tau'}(Z_\alpha, K') \oplus
H^n_{\tau'}(Y_\alpha, K') \to
H^n_{\tau'}(E_\alpha, K')
$$
这与正合序列
$$
0 \to
\Gamma(X_\alpha, H^n(K')) \to
\Gamma(Z_\alpha, H^n(K')) \oplus
\Gamma(Y_\alpha, H^n(K')) \to
\Gamma(E_\alpha, H^n(K'))
$$
由假设 (1)，可得 $H^n(K')$ 是 $\tau$-层。然而，由于
$\tau$-层化 $\epsilon^{-1}H^n(K') = H^n(\epsilon^{-1}K')$，且
$\epsilon^{-1}H^n(K') = H^n(\epsilon^{-1}K')$ 为 $0$，因为
$\epsilon^{-1}K' = 0$，故其为 $0$；于是按要求
$H^n(K') = 0$。
\end{proof}

\begin{lemma}
\label{sites-cohomology-lemma-descent-squares-helper}
设 $\epsilon : (\mathcal{C}_\tau, \mathcal{O}_\tau) \to
(\mathcal{C}_{\tau'}, \mathcal{O}_{\tau'})$ 如上。令
$$
\xymatrix{
E \ar[d] \ar[r] & Y \ar[d] \\
Z \ar[r] & X
}
$$
为范畴 $\mathcal{C}$ 中满足下列条件的交换图：
\begin{enumerate}
\item $h_X^\# = h_Y^\# \amalg_{h_E^\#} h_Z^\#$；
\item $h_E^\# \to h_Y^\#$ 是单射。
\end{enumerate}
其中上标井号表示 $\tau$-层化。则对属于 $R\epsilon_*$ 本质像的
$K' \in D(\mathcal{O}_{\tau'})$，引理 \ref{sites-cohomology-lemma-c-square} 的映射
$c^{K'}_{X, Z, Y, E}$（使用 $\tau'$-拓扑）是同构。
\end{lemma}


\begin{proof}
此辅助引理几乎立即由引理 \ref{sites-cohomology-lemma-square-triangle} 得出，
我们强烈建议读者跳过证明。设 $K' = R\epsilon_*K$。
选取表示 $K$ 的 $\mathcal{O}_\tau$-模 K-内射复形
$\mathcal{J}^\bullet$。则 $\epsilon_*\mathcal{J}^\bullet$ 是表示 $K'$ 的
$\mathcal{O}_{\tau'}$-模 K-内射复形，参见引理
\ref{sites-cohomology-lemma-K-injective-flat}。接着，
$$
0 \to
\mathcal{J}^\bullet(X) \xrightarrow{\alpha}
\mathcal{J}^\bullet(Z) \oplus
\mathcal{J}^\bullet(Y) \xrightarrow{\beta}
\mathcal{J}^\bullet(E)
\to 0
$$
是阿贝尔群复形的短正合序列，参见引理
\ref{sites-cohomology-lemma-square-triangle} 及其证明。由于这正是定义
$c^{K'}_{X, Z, Y, E}$ 时所用的阿贝尔群复形序列，结论成立。
\end{proof}







\section{比较上同调}
\label{sites-cohomology-section-compare-general}

\noindent
我们发展一些一般理论，以帮助比较不同拓扑下的上同调。给定第
\ref{sites-cohomology-section-compare} 节中的 $\mathcal{C}$、$\tau$、$\tau'$，以及
$\mathcal{C}$ 中的态射 $f : X \to Y$，得到拓扑斯态射的交换图
\begin{equation}
\label{sites-cohomology-equation-commutative-epsilon}
\vcenter{
\xymatrix{
\Sh(\mathcal{C}_\tau/X) \ar[r]_{f_\tau} \ar[d]_{\epsilon_X} &
\Sh(\mathcal{C}_\tau/Y) \ar[d]^{\epsilon_Y} \\
\Sh(\mathcal{C}_{\tau'}/X) \ar[r]^{f_{\tau'}} &
\Sh(\mathcal{C}_{\tau'}/X)
}
}
\end{equation}
这里的态射 $\epsilon_X$（分别为 $\epsilon_Y$）是第
\ref{sites-cohomology-section-compare} 节中为带有 $\tau$ 和 $\tau'$ 两个拓扑的范畴
$\mathcal{C}/X$ 定义的比较态射。态射 $f_\tau$ 和 $f_{\tau'}$
是“重新局部化”态射（站点，引理 \ref{sites-lemma-relocalize}）。
图的交换性是站点，引理 \ref{sites-lemma-localize-morphism} 的特例
（取 $\mathcal{C} = \mathcal{C}_\tau/Y$、$\mathcal{D} = \mathcal{C}_{\tau'}/Y$、
$u = \text{id}$、$U = X$ 及 $V = X$）。此外还得到
$\epsilon_{X, *} \circ f_\tau^{-1} = f_{\tau'}^{-1} \circ \epsilon_{Y, *}$
这既可由该引理得到，也可因其显然性得到。

\begin{situation}
\label{sites-cohomology-situation-compare}
设 $\mathcal{C}$、$\tau$、$\tau'$ 如第 \ref{sites-cohomology-section-compare} 节。
假设给定子集 $\mathcal{P} \subset \text{Arrows}(\mathcal{C})$，并且对
$\mathcal{C}$ 的每个对象 $X$，给定弱 Serre 子范畴
$\mathcal{A}'_X \subset \textit{Ab}(\mathcal{C}_{\tau'}/X)$。
作如下假设：
\begin{enumerate}
\item
\label{sites-cohomology-item-base-change-P}
若 $f : X \to Y$ 属于 $\mathcal{P}$，且 $Y' \to Y$ 为任意态射，
则 $X \times_Y Y'$ 存在，且 $X \times_Y Y' \to Y'$ 属于 $\mathcal{P}$；
\item
\label{sites-cohomology-item-restriction-A}
$f_{\tau'}^{-1}$ 将 $\mathcal{A}'_Y$ 映入 $\mathcal{A}'_X$，
对 $\mathcal{C}$ 的任意态射 $f : X \to Y$ 均如此；
\item
\label{sites-cohomology-item-A-sheaf}
若 $X \in \mathcal{C}$ 且 $\mathcal{F}' \in \mathcal{A}'_X$，则
$\mathcal{F}'$ 满足 $\tau$-覆盖的层条件，即
$\mathcal{F}' = \epsilon_{X, *}\epsilon_X^{-1}\mathcal{F}'$,
\item
\label{sites-cohomology-item-A-and-P}
若 $f : X \to Y$ 属于 $\mathcal{P}$ 且
$\mathcal{F}' \in \Ob(\mathcal{A}'_X)$，则当 $i \geq 0$ 时
$R^if_{\tau', *}\mathcal{F}' \in \Ob(\mathcal{A}'_Y)$；
\item
\label{sites-cohomology-item-refine-tau-by-P}
若 $\{U_i \to U\}_{i \in I}$ 是 $\tau$-覆盖，则存在
\begin{enumerate}
\item $\tau'$-覆盖 $\{V_j \to U\}_{j \in J}$；
\item $\tau$-覆盖 $\{f_j : W_j \to V_j\}$，其中每个覆盖只含有一个
$f_j \in \mathcal{P}$；
\item $\tau'$-覆盖 $\{W_{jk} \to W_j\}_{k \in K_j}$
\end{enumerate}
使得 $\{W_{jk} \to U\}_{j \in J, k \in K_j}$ 是
$\{U_i \to U\}_{i \in I}$ 的细化。
\end{enumerate}
\end{situation}

\begin{lemma}
\label{sites-cohomology-lemma-A}
在情形 \ref{sites-cohomology-situation-compare} 中，对 $X \in \mathcal{C}$，记
$\mathcal{A}_X$ 为 $\textit{Ab}(\mathcal{C}_\tau/X)$ 中形如
$\epsilon_X^{-1}\mathcal{F}'$ 的对象，其中 $\mathcal{F}' \in \mathcal{A}'_X$。
则
\begin{enumerate}
\item 对 $\mathcal{F} \in \textit{Ab}(\mathcal{C}_\tau/X)$，有
$\mathcal{F} \in \mathcal{A}_X \Leftrightarrow
\epsilon_{X, *}\mathcal{F} \in \mathcal{A}'_X$；
\item 对 $\mathcal{C}$ 的任意态射 $f : X \to Y$，$f_\tau^{-1}$ 将
$\mathcal{A}_Y$ 映入 $\mathcal{A}_X$。
\end{enumerate}
\end{lemma}

\begin{proof}
第 (1) 部分由 (\ref{sites-cohomology-item-A-sheaf}) 得出；第 (2) 部分由
(\ref{sites-cohomology-item-restriction-A}) 以及 (\ref{sites-cohomology-equation-commutative-epsilon}) 的交换性得出，
其中
$\epsilon_X^{-1} \circ f_{\tau'}^{-1} = f_\tau^{-1} \circ \epsilon_Y^{-1}$.
\end{proof}

\noindent
我们的下一个目标是证明引理 \ref{sites-cohomology-lemma-compare-cohomology-general} 和
\ref{sites-cohomology-lemma-cohomological-descent-general}。我们将使用下面的归纳假设，
通过归纳论证完成证明。

\medskip\noindent
$(V_n)$ 对 $X \in \mathcal{C}$ 及 $\mathcal{F} \in \mathcal{A}_X$，有
$R^i\epsilon_{X, *}\mathcal{F} = 0$（$1 \leq i \leq n$）。

\begin{lemma}
\label{sites-cohomology-lemma-V-implies-C-general}
在情形 \ref{sites-cohomology-situation-compare} 中假设 $(V_n)$ 成立。
若 $f : X \to Y$ 属于 $\mathcal{P}$ 且 $\mathcal{F} \in \mathcal{A}_X$，
则当 $i \leq n$ 时有
$R^if_{\tau', *}\epsilon_{X, *}\mathcal{F} =
\epsilon_{Y, *}R^if_{\tau, *}\mathcal{F}$。
\end{lemma}

\begin{proof}
下文不再特别说明，我们将使用交换图
(\ref{sites-cohomology-equation-commutative-epsilon})。特别地，有
$$
Rf_{\tau', *}R\epsilon_{X, *}\mathcal{F} =
R\epsilon_{Y, *}Rf_{\tau, *}\mathcal{F}
$$
假设 $(V_n)$ 告诉我们
$\epsilon_{X, *}\mathcal{F} \to R\epsilon_{X, *}\mathcal{F}$
在次数 $\leq n$ 上是同构。因此
$Rf_{\tau', *}\epsilon_{X, *}\mathcal{F} \to
Rf_{\tau', *}R\epsilon_{X, *}\mathcal{F}$ 在次数 $\leq n$ 上是同构。
于是
$$
R^if_{\tau', *}\epsilon_{X, *}\mathcal{F} \to
H^i(R\epsilon_{Y, *}Rf_{\tau, *}\mathcal{F})
$$
当 $i \leq n$ 时是同构。我们考察引理
\ref{sites-cohomology-lemma-relative-Leray} 中关于
$R\epsilon_{Y, *}Rf_{\tau, *}\mathcal{F}$ 的谱序列的第二页来证明本引理。
图示如下：
$$
\begin{matrix}
\ldots &
\ldots &
\ldots &
\ldots \\
\epsilon_{Y, *}R^2f_{\tau, *}\mathcal{F} &
R^1\epsilon_{Y, *}R^2f_{\tau, *}\mathcal{F} &
R^2\epsilon_{Y, *}R^2f_{\tau, *}\mathcal{F} &
\ldots \\
\epsilon_{Y, *}R^1f_{\tau, *}\mathcal{F} &
R^1\epsilon_{Y, *}R^1f_{\tau, *}\mathcal{F} &
R^2\epsilon_{Y, *}R^1f_{\tau, *}\mathcal{F} &
\ldots \\
\epsilon_{Y, *}f_{\tau, *}\mathcal{F} &
R^1\epsilon_{Y, *}f_{\tau, *}\mathcal{F} &
R^2\epsilon_{Y, *}f_{\tau, *}\mathcal{F} &
\ldots
\end{matrix}
$$
令 $(C_m)$ 表示假设：
$R^if_{\tau', *}\epsilon_{X, *}\mathcal{F} =
\epsilon_{Y, *}R^if_{\tau, *}\mathcal{F}$（$i \leq m$）。
注意 $(C_0)$ 成立。我们将证明当 $m < n$ 时
$(C_{m - 1}) \Rightarrow (C_m)$。具体地，若 $(C_{m - 1})$ 成立，则当
$n \geq p > 0$ 且 $q \leq m - 1$ 时有
\begin{align*}
R^p\epsilon_{Y, *}R^qf_{\tau, *}\mathcal{F}
& =
R^p\epsilon_{Y, *}
\epsilon_Y^{-1} \epsilon_{Y, *} R^qf_{\tau, *}\mathcal{F} \\
& =
R^p\epsilon_{Y, *}
\epsilon_Y^{-1}R^qf_{\tau', *}\epsilon_{X, *}\mathcal{F} = 0
\end{align*}
第一个等号来自 $\epsilon_Y^{-1}\epsilon_{Y, *} = \text{id}$，
第二个来自 $(C_{m - 1})$，最后一个来自 $(V_n)$，因为由
(\ref{sites-cohomology-item-A-and-P})，$\epsilon_Y^{-1}R^qf_{\tau', *}\epsilon_{X, *}\mathcal{F}$
属于 $\mathcal{A}_Y$。考察谱序列可知
$E_2^{0, m} = \epsilon_{Y, *}R^mf_{\tau, *}\mathcal{F}$
是满足 $p + q = m$ 的 $E_2^{p, q}$ 中唯一非零项。记住
$\text{d}_r^{p, q} : E_r^{p, q} \to E_r^{p + r, q - r + 1}$。
因此不存在 $r \geq 2$ 的非零微分
$\text{d}_r^{p, q}$ 从该位置出发或进入该位置。于是
$H^m(R\epsilon_{Y, *}Rf_{\tau, *}\mathcal{F}) =
\epsilon_{Y, *}R^mf_{\tau, *}\mathcal{F}$，由上面的讨论得到
$(C_m)$。

\medskip\noindent
最后，假设 $(C_{n - 1})$。同样的分析表明
$E_2^{0, n} = \epsilon_{Y, *}R^nf_{\tau, *}\mathcal{F}$
是满足 $p + q = n$ 的 $E_2^{p, q}$ 中唯一非零项。
仍然没有非零微分进入该位置，但可能有非零微分从该位置出发。
具体地，映射
$d_{n + 1}^{0, n} : \epsilon_{Y, *}R^nf_{\tau, *}\mathcal{F} \to
R^{n + 1}\epsilon_{Y, *}f_{\tau, *}\mathcal{F}$.
因此存在正合序列
$$
0 \to R^nf_{\tau', *}\epsilon_{X, *}\mathcal{F} \to 
\epsilon_{Y, *}R^nf_{\tau, *}\mathcal{F} \to
R^{n + 1}\epsilon_{Y, *}f_{\tau, *}\mathcal{F}
$$
由 (\ref{sites-cohomology-item-A-and-P}) 和 (\ref{sites-cohomology-item-A-sheaf})，层
$R^nf_{\tau', *}\epsilon_{X, *}\mathcal{F}$
满足 $\tau$-覆盖的层性质，$\epsilon_{Y, *}R^nf_{\tau, *}\mathcal{F}$
也满足（使用第 \ref{sites-cohomology-section-compare} 节中对 $\epsilon_*$ 的描述）。
然而，$\tau'$-层
$R^{n + 1}\epsilon_{Y, *}f_{\tau, *}\mathcal{F}$
的 $\tau$-层化为零（由上同调的局部性；使用引理
\ref{sites-cohomology-lemma-kill-cohomology-class-on-covering} 和
\ref{sites-cohomology-lemma-higher-direct-images}）。因此
$R^nf_{\tau', *}\epsilon_{X, *}\mathcal{F} \to
\epsilon_{Y, *}R^nf_{\tau, *}\mathcal{F}$
必为同构，证明完成。
\end{proof}

\noindent
若 $E'$（分别地，$E$）是
$D(\mathcal{C}_{\tau'}/X)$（分别地，$D(\mathcal{C}_\tau/X)$）中的对象，
则对 $\mathcal{C}/X$ 中的对象 $U$，记
$H^n_{\tau'}(U, E')$（分别地，$H^n_\tau(U, E)$）为
$E'$（分别地，$E$）在 $U$ 上的上同调。

\begin{lemma}
\label{sites-cohomology-lemma-V-implies-cohomology-general}
在情形 \ref{sites-cohomology-situation-compare} 中，若 $(V_n)$ 成立，则对
$X \in \mathcal{C}$ 及 $L \in D(\mathcal{C}_{\tau'}/X)$，若
$i < 0$ 时 $H^i(L) = 0$，且 $0 \leq i \leq n$ 时
$H^i(L) \in \mathcal{A}'_X$，则有
$H^n_{\tau'}(X, L) = H^n_\tau(X, \epsilon_X^{-1}L)$.
\end{lemma}

\begin{proof}
由引理 \ref{sites-cohomology-lemma-Leray-unbounded}，有
$H^n_\tau(X, \epsilon_X^{-1}L) =
H^n_{\tau'}(X, R\epsilon_{X, *}\epsilon_X^{-1}L)$.
存在谱序列
$$
E_2^{p, q} = R^p\epsilon_{X, *}\epsilon_X^{-1}H^q(L)
$$
收敛到 $H^{p + q}(R\epsilon_{X, *}\epsilon_X^{-1}L)$。
由 $(V_n)$，当 $0 < p \leq n$ 且 $0 \leq q \leq n$ 时
$E_2^{p, q}$ 消失。因此
$E_2^{0, q} = \epsilon_{X, *}\epsilon_X^{-1}H^q(L) = H^q(L)$
是满足 $p + q \leq n$ 的 $E_2^{p, q}$ 中唯一可能的非零项。
于是映射
$$
L \longrightarrow R\epsilon_{X, *}\epsilon_X^{-1}L
$$
在次数 $\leq n$ 上是同构（细节略去）。因此
$H^i_{\tau'}(X, L) = H^i_{\tau'}(X, R\epsilon_{X, *}\epsilon_X^{-1}L)$
对 $i \leq n$ 成立。引理得证。
\end{proof}

\begin{lemma}
\label{sites-cohomology-lemma-V-implies-cohomology-extra-general}
在情形 \ref{sites-cohomology-situation-compare} 中，若 $(V_n)$ 成立，则对
$X \in \mathcal{C}$ 及 $\mathcal{F} \in \mathcal{A}_X$，映射
$H^{n + 1}_{\tau'}(X, \epsilon_{X, *}\mathcal{F}) \to
H^{n + 1}_\tau(X, \mathcal{F})$
是单射，其像正是那些在 $X$ 的某个 $\tau'$-覆盖上变为零的类。
\end{lemma}

\begin{proof}
记住 $\epsilon_X^{-1}\epsilon_{X, *}\mathcal{F} = \mathcal{F}$，
故该映射由沿 $\epsilon_X$ 拉回上同调类给出。Leray 谱序列（引理
\ref{sites-cohomology-lemma-Leray}）
$$
E_2^{p, q} = H^p_{\tau'}(X, R^q\epsilon_{X, *}\mathcal{F})
\Rightarrow
H^{p + q}_\tau(X, \mathcal{F})
$$
结合已假设的消失，得到正合序列
$$
0 \to
H^{n + 1}_{\tau'}(X, \epsilon_{X, *}\mathcal{F}) \to
H^{n + 1}_\tau(X, \mathcal{F}) \to
H^0_{\tau'}(X, R^{n + 1}\epsilon_{X, *}\mathcal{F})
$$
这正是引理的另一种表述。
\end{proof}

\begin{lemma}
\label{sites-cohomology-lemma-make-class-zero-general}
在情形 \ref{sites-cohomology-situation-compare} 中，设 $f : X \to Y$ 属于 $\mathcal{P}$，
且 $\{X \to Y\}$ 是 $\tau$-覆盖。设 $\mathcal{F}' \in \mathcal{A}'_Y$。
若 $n \geq 0$ 且
$$
\theta \in
\text{Equalizer}\left(
\xymatrix{
H^{n + 1}_{\tau'}(X, \mathcal{F}')
\ar@<1ex>[r] \ar@<-1ex>[r] &
H^{n + 1}_{\tau'}(X \times_Y X, \mathcal{F}')
}
\right)
$$
则存在 $\tau'$-覆盖 $\{Y_i \to Y\}$，使得 $\theta$ 在
$H^{n + 1}_{\tau'}(Y_i \times_Y X, \mathcal{F}')$ 中的限制为零。
\end{lemma}

\begin{proof}
由 (\ref{sites-cohomology-item-base-change-P})，$X \times_Y X$ 存在。对
$Z \in \mathcal{C}/Y$，记 $\mathcal{F}'|_Z$ 为 $\mathcal{F}'$ 到
$\mathcal{C}_{\tau'}/Z$ 的限制。记住
$H^{n + 1}_{\tau'}(X, \mathcal{F}') =
H^{n + 1}(\mathcal{C}_{\tau'}/X, \mathcal{F}'|_X)$，参见引理
\ref{sites-cohomology-lemma-cohomology-of-open}。本引理断言 $\theta$ 的像
$\overline{\theta} \in H^0(Y, R^{n + 1}f_{\tau', *}\mathcal{F}'|_X)$
为零。考虑笛卡尔图
$$
\xymatrix{
X \times_Y X \ar[d]_{\text{pr}_1} \ar[r]_{\text{pr}_2} &
X \ar[d]^f \\
X \ar[r]^f & Y
}
$$
由平凡基变换（引理 \ref{sites-cohomology-lemma-localize-cartesian-square}），有
$$
f_{\tau'}^{-1}R^{n + 1}f_{\tau', *}(\mathcal{F}'|_X) =
R^{n + 1}\text{pr}_{1, \tau', *}(\mathcal{F}'|_{X \times_Y X})
$$
若 $\text{pr}_1^{-1}\theta = \text{pr}_2^{-1}\theta$，则
$f_{\tau'}^{-1}\overline{\theta}$ 作为
$f_{\tau'}^{-1}R^{n + 1}f_{\tau', *}(\mathcal{F}'|_X)$ 的截面为零，
因为显然 $\text{pr}_1^{-1}\theta$ 映到
$H^0(X, R^{n + 1}\text{pr}_{1, \tau', *}(\mathcal{F}'|_{X \times_Y X}))$
中的零元素。由 (\ref{sites-cohomology-item-restriction-A})，$\mathcal{F}'|_X \in \mathcal{A}'_X$。
因此由 (\ref{sites-cohomology-item-A-and-P})，
$\mathcal{G}' = R^{n + 1}f_{\tau', *}(\mathcal{F}'|_X)$
是 $\mathcal{A}'_Y$ 中的对象。由 (\ref{sites-cohomology-item-A-sheaf})，
$\mathcal{G}'$ 满足 $\tau$-覆盖的层性质。
由于 $\{X \to Y\}$ 是 $\tau$-覆盖，限制
$\mathcal{G}'(Y) \to \mathcal{G}'(X)$ 是单射。因此
$\overline{\theta} = 0$。
\end{proof}

\begin{lemma}
\label{sites-cohomology-lemma-induction-step-V-C-general}
在情形 \ref{sites-cohomology-situation-compare} 中，有 $(V_n) \Rightarrow (V_{n + 1})$。
\end{lemma}

\begin{proof}
设 $X$ 是 $\mathcal{C}$ 中的对象，$\mathcal{F}$ 是 $\mathcal{A}_X$ 中的对象。
对某个 $U/X$，设 $\xi \in H^{n + 1}_\tau(U, \mathcal{F})$。
我们需要证明：$\xi$ 在
某个 $\tau'$-覆盖的各成员上的限制为零。参见
引理 \ref{sites-cohomology-lemma-higher-direct-images}。
由此可知，我们可以将
$U$ 替换为某个 $\tau'$-覆盖的各成员。

\medskip\noindent
由上同调的局部性
（引理 \ref{sites-cohomology-lemma-kill-cohomology-class-on-covering}），
我们可以选取一个 $\tau$-覆盖 $\{U_i \to U\}$，
使得 $\xi$ 在 $U_i$ 上的限制为零。
按照 (\ref{sites-cohomology-item-refine-tau-by-P})，选取 $\{V_j \to V\}$、$\{f_j : W_j \to V_j\}$，以及
$\{W_{jk} \to W_j\}$。
将 $U$ 替换为 $V_j$
并将 $\mathcal{F}$ 替换为其在
$\mathcal{C}_\tau/V_j$ 上的限制；由
(\ref{sites-cohomology-item-base-change-P})，这是允许的。于是我们归结为
下一段所讨论的情形。

\medskip\noindent
这里 $f : X \to Y$ 是 $\mathcal{P}$ 中的元素，
使得 $\{X \to Y\}$ 是 $\tau$-覆盖，
$\mathcal{F}$ 是 $\mathcal{A}_Y$ 中的对象，并且
$\xi \in H^{n + 1}_\tau(Y, \mathcal{F})$ 满足：
存在一个 $\tau'$-覆盖 $\{X_i \to X\}_{i \in I}$，
使得对所有 $i \in I$，$\xi$ 在 $X_i$ 上的限制为零。
问题：证明 $\xi$ 在 $Y$ 的某个 $\tau'$-覆盖上的限制为零。

\medskip\noindent
由引理 \ref{sites-cohomology-lemma-V-implies-cohomology-extra-general}，存在一个
唯一的 $\tau'$-上同调类
$\theta \in H^{n + 1}_{\tau'}(X, \epsilon_{X, *}\mathcal{F})$
其像为 $\xi|_X$。
由于 $\xi|_X$ 经两个投影拉回到 $X \times_Y X$ 后得到相同的类，
由此（利用唯一性），$\theta$ 也具有同样性质。
由引理 \ref{sites-cohomology-lemma-make-class-zero-general}
可知，将 $Y$ 替换为某个 $\tau'$-覆盖的各成员后，
可以设 $\theta = 0$。
因此可以设 $\xi|_X$ 为零。

\medskip\noindent
设 $f : X \to Y$ 是 $\mathcal{P}$ 中的一个元素，
使得 $\{X \to Y\}$ 是 $\tau$-覆盖，
$\mathcal{F}$ 是 $\mathcal{A}_Y$ 中的对象，并且
$\xi \in H^{n + 1}_\tau(Y, \mathcal{F})$
在 $H^{n + 1}_\tau(X, \mathcal{F})$ 中映为零。
问题：证明 $\xi$ 在 $Y$ 的某个 $\tau'$-覆盖上的限制为零。

\medskip\noindent
这些假设说明 $\xi$ 在下述态射下的像为零：
$$
\mathcal{F} \longrightarrow Rf_{\tau, *}f_\tau^{-1}\mathcal{F}
$$
使用引理 \ref{sites-cohomology-lemma-Leray-unbounded}。
利用截断的有区别三角
（《导出范畴》，评注
\ref{derived-remark-truncation-distinguished-triangle}）可作一个简单论证，说明
$\xi$ 在下述态射下的像为零：
$$
\mathcal{F} \longrightarrow \tau_{\leq n}Rf_{\tau, *}f_\tau^{-1}\mathcal{F}
$$
我们将把它与态射 $\epsilon_{Y, *}\mathcal{F} \to K$ 比较，
其中
$$
K = \tau_{\leq n}Rf_{\tau', *}f_{\tau'}^{-1}\epsilon_{Y, *}\mathcal{F} =
\tau_{\leq n}Rf_{\tau', *}\epsilon_{X, *}f_{\tau}^{-1}\mathcal{F}
$$
等式
$\epsilon_{X, *} f_\tau^{-1} = f_{\tau'}^{-1} \epsilon_{Y, *}$
是 (\ref{sites-cohomology-equation-commutative-epsilon}) 的一个性质。考虑态射
$$
Rf_{\tau', *}\epsilon_{X, *}f_{\tau}^{-1}\mathcal{F} \longrightarrow
Rf_{\tau', *}R\epsilon_{X, *}f_{\tau}^{-1}\mathcal{F} =
R\epsilon_{Y, *}Rf_{\tau, *}f_\tau^{-1}\mathcal{F}
$$
该态射用于引理 \ref{sites-cohomology-lemma-V-implies-C-general} 的证明，
并由伴随诱导出态射
$$
\epsilon_Y^{-1} Rf_{\tau', *}\epsilon_{X, *}f_{\tau}^{-1}\mathcal{F} \to
Rf_{\tau, *}f_\tau^{-1}\mathcal{F}
$$
对其作截断，得到态射
$$
\epsilon_Y^{-1}K
\longrightarrow
\tau_{\leq n}Rf_{\tau, *}f_\tau^{-1}\mathcal{F}
$$
由引理 \ref{sites-cohomology-lemma-V-implies-C-general}，该态射是同构；
该引理适用，因为 $f_\tau^{-1}\mathcal{F}$ 属于 $\mathcal{A}_X$，
这是引理 \ref{sites-cohomology-lemma-A} 的结论。选取一个有区别三角
$$
\epsilon_{Y, *}\mathcal{F} \to K \to L \to \epsilon_{Y, *}\mathcal{F}[1]
$$
态射 $\mathcal{F} \to f_{\tau, *}f_\tau^{-1}\mathcal{F}$
由于 $\{X \to Y\}$ 是 $\tau$-覆盖，该态射是单射。因此
$\epsilon_{Y, *}\mathcal{F} \to
\epsilon_{Y, *}f_{\tau, *}f_\tau^{-1}\mathcal{F} =
f_{\tau', *}f_{\tau'}^{-1}\epsilon_{Y, *}\mathcal{F}$
也为单射。
因此 $L$ 的非零上同调层只可能位于次数 $0, \ldots, n$。
由于 $f_{\tau', *}f_{\tau'}^{-1}\epsilon_{Y, *}\mathcal{F}$
根据 (\ref{sites-cohomology-item-restriction-A}) 和
(\ref{sites-cohomology-item-A-and-P}) 属于 $\mathcal{A}'_Y$，故
$$
H^0(L) = \Coker(\epsilon_{Y, *}\mathcal{F} \to
f_{\tau', *}f_{\tau'}^{-1}\epsilon_{Y, *}\mathcal{F})
$$
属于弱 Serre 子范畴 $\mathcal{A}'_Y$。当 $1 \leq i \leq n$ 时，
我们看到 $H^i(L) = R^if_{\tau', *}f_{\tau'}^{-1}\epsilon_{Y, *}\mathcal{F}$
根据 (\ref{sites-cohomology-item-restriction-A}) 和
(\ref{sites-cohomology-item-A-and-P}) 属于 $\mathcal{A}'_Y$。
用 $\epsilon_Y$ 拉回上面的有区别三角，得到有区别三角
$$
\mathcal{F} \to \tau_{\leq n}Rf_{\tau, *}f_\tau^{-1}\mathcal{F}
\to \epsilon_Y^{-1}L \to \mathcal{F}[1]
$$
由于 $\xi$ 在中间项中的像为零，我们得到
$\xi$ 是某个元素的像：
$\xi' \in H^n_\tau(Y, \epsilon_Y^{-1}L)$.
由引理 \ref{sites-cohomology-lemma-V-implies-cohomology-general}，有
$$
H^n_{\tau'}(Y, L) = H^n_\tau(Y, \epsilon_Y^{-1}L),
$$
因此可以将 $\xi'$ 提升为 $H^n_{\tau'}(Y, L)$ 中的元素，
并取其边界映到
$H^{n + 1}_{\tau'}(Y, \epsilon_{Y, *}\mathcal{F})$
从而 $\xi$ 属于规范态射的像：
$H^{n + 1}_{\tau'}(Y, \epsilon_{Y, *}\mathcal{F}) \to
H^{n + 1}_\tau(Y, \mathcal{F})$.
由上同调的局部性，对于
$H^{n + 1}_{\tau'}(Y, \epsilon_{Y, *}\mathcal{F})$，参见
引理 \ref{sites-cohomology-lemma-kill-cohomology-class-on-covering}，
得证。
\end{proof}

\begin{lemma}
\label{sites-cohomology-lemma-V-C-all-n-general}
在情形 \ref{sites-cohomology-situation-compare} 中，有
$(V_n)$ 对所有 $n$ 都成立。此外：
\begin{enumerate}
\item 对 $X$ 是 $\mathcal{C}$ 中的对象，且
$K' \in D^+_{\mathcal{A}'_X}(\mathcal{C}_{\tau'}/X)$ 上的态射
$K' \to R\epsilon_{X, *}(\epsilon_X^{-1}K')$ 是同构。
\item 对 $f : X \to Y$ 属于 $\mathcal{P}$，且
$K' \in D^+_{\mathcal{A}'_X}(\mathcal{C}_{\tau'}/X)$，有
$Rf_{\tau', *}K' \in D^+_{\mathcal{A}'_X}(\mathcal{C}_{\tau'}/Y)$，且
$\epsilon_Y^{-1}(Rf_{\tau', *}K') = Rf_{\tau, *}(\epsilon_X^{-1}K')$.
\end{enumerate}
\end{lemma}

\begin{proof}
注意 $(V_0)$ 成立，因为这是空条件。
于是由
引理 \ref{sites-cohomology-lemma-induction-step-V-C-general} 得到所有 $n$ 的 $(V_n)$。

\medskip\noindent
证明 (1)。对象 $K = \epsilon_X^{-1}K'$ 的上同调层
$H^i(K) = \epsilon_X^{-1}H^i(K')$ 属于 $\mathcal{A}_X$。
因此有谱序列
$$
E_2^{p, q} = R^p\epsilon_{X, *} H^q(K) \Rightarrow
H^{p + q}(R\epsilon_{X, *}K)
$$
由 $(V_n)$ 对所有 $n$ 成立，该谱序列退化，从而
$$
H^n(R\epsilon_{X, *}K) = \epsilon_{X, *}H^n(K) =
\epsilon_{X, *}\epsilon_X^{-1}H^i(K') = H^i(K').
$$
再次利用 $H^i(K')$ 属于 $\mathcal{A}'_X$ 这一事实。
因此规范态射 $K' \to R\epsilon_{X, *}(\epsilon_X^{-1}K')$
是同构。

\medskip\noindent
证明 (2)。利用谱序列
$$
E_2^{p, q} = R^pf_{\tau', *}H^q(K') \Rightarrow R^{p + q}f_{\tau', *}K'
$$
$R^pf_{\tau', *}H^q(K')$ 属于
$\mathcal{A}'_Y$（由 (\ref{sites-cohomology-item-A-and-P})），
以及 $\mathcal{A}'_Y$ 是
$\textit{Ab}(\mathcal{C}_{\tau'}/Y)$ 的弱 Serre 子范畴这一事实，
再由同调，引理 \ref{homology-lemma-biregular-ss-converges}
可知
$Rf_{\tau', *}K' \in D^+_{\mathcal{A}'_X}(\mathcal{C}_{\tau'}/X)$.
为完成证明，需证明基变换态射
$$
\epsilon_Y^{-1}(Rf_{\tau', *}K')
\longrightarrow
Rf_{\tau, *}(\epsilon_X^{-1}K')
$$
它是同构。将上面的谱序列与
谱序列
$$
E_2^{p, q} = R^pf_{\tau, *}H^q(\epsilon_X^{-1}K')
\Rightarrow R^{p + q}f_{\tau, *}\epsilon_X^{-1}K'
$$
由此可将问题化归为 $K'$ 只有一个非零
上同调层 $\mathcal{F}' \in \mathcal{A}'_X$ 的情形；细节略去。
然后由引理 \ref{sites-cohomology-lemma-V-implies-C-general} 得到
$\epsilon_Y^{-1}R^if_{\tau', *}\mathcal{F}' =
R^if_{\tau, *}\epsilon_X^{-1}\mathcal{F}'$ 对所有 $i$ 均成立，
证明完成。
\end{proof}

\begin{lemma}
\label{sites-cohomology-lemma-cohomological-descent-general}
在情形 \ref{sites-cohomology-situation-compare} 中，对于任意 $X$，
$\mathcal{C}$ 中的范畴
$\mathcal{A}_X \subset \textit{Ab}(\mathcal{C}_\tau/X)$
是弱 Serre 子范畴，且函子
$$
R\epsilon_{X, *} :
D^+_{\mathcal{A}_X}(\mathcal{C}_\tau/X)
\longrightarrow
D^+_{\mathcal{A}'_X}(\mathcal{C}_{\tau'}/X)
$$
是等价；其拟逆由 $\epsilon_X^{-1}$ 给出。
\end{lemma}

\begin{proof}
我们需要检查同调引理
\ref{homology-lemma-characterize-weak-serre-subcategory} 中列出的
关于 $\mathcal{A}_X$ 的条件。若 $\varphi : \mathcal{F} \to \mathcal{G}$ 是
$\mathcal{A}_X$ 中的态射，则 $\epsilon_{X, *}\varphi$：
$\epsilon_{X, *}\mathcal{F} \to \epsilon_{X, *}\mathcal{G}$
是 $\mathcal{A}'_X$ 中的态射。因此 $\Ker(\epsilon_{X, *}\varphi)$ 与
$\Coker(\epsilon_{X, *}\varphi)$ 是 $\mathcal{A}'_X$ 中的对象，
因为 $\mathcal{A}'_X$ 是 $\textit{Ab}(\mathcal{C}_{\tau'}/X)$ 的弱 Serre 子范畴。
应用 $\epsilon_X^{-1}$，得到正合序列
$$
0 \to
\epsilon_X^{-1}\Ker(\epsilon_{X, *}\varphi) \to
\mathcal{F} \to \mathcal{G} \to
\epsilon_X^{-1}\Coker(\epsilon_{X, *}\varphi) \to 0
$$
由此可见 $\Ker(\varphi)$ 和 $\Coker(\varphi)$ 属于
$\mathcal{A}_X$。最后，设
$$
0 \to \mathcal{F}_1 \to \mathcal{F}_2 \to \mathcal{F}_3 \to 0
$$
这是 $\textit{Ab}(\mathcal{C}_\tau/X)$ 中的短正合序列，
且 $\mathcal{F}_1$ 与 $\mathcal{F}_3$ 属于 $\mathcal{A}_X$。
随后应用 $\epsilon_{X, *}$，得到正合序列
$$
0 \to 
\epsilon_{X, *}\mathcal{F}_1 \to
\epsilon_{X, *}\mathcal{F}_2 \to
\epsilon_{X, *}\mathcal{F}_3 \to
R^1\epsilon_{X, *}\mathcal{F}_1 = 0
$$
由引理 \ref{sites-cohomology-lemma-V-C-all-n-general}，该项消失。
因此 $\epsilon_{X, *}\mathcal{F}_2$ 属于 $\mathcal{A}'_X$
因为它是 $\textit{Ab}(\mathcal{C}_{\tau'}/X)$ 的弱 Serre 子范畴。
用 $\epsilon_X$ 拉回，得到
$\mathcal{F}_2$ 属于 $\mathcal{A}_X$。

\medskip\noindent
因此 $\mathcal{A}_X$ 是
$\textit{Ab}(\mathcal{C}_\tau/X)$ 的弱 Serre 子范畴，于是考虑
范畴 $D^+_{\mathcal{A}_X}(\mathcal{C}_\tau/X)$ 是有意义的。
注意 $\epsilon_X^{-1} : \mathcal{A}'_X \to \mathcal{A}_X$
是等价，并且
$\mathcal{F}' \to R\epsilon_{X, *}\epsilon_X^{-1}\mathcal{F}'$
对 $\mathcal{F}' \in \mathcal{A}'_X$，这是同构，因为我们
根据引理 \ref{sites-cohomology-lemma-V-C-all-n-general} 已知 $(V_n)$ 对所有 $n$ 成立。
因此由引理 \ref{sites-cohomology-lemma-equivalence-bounded} 得出结论。
\end{proof}

\begin{lemma}
\label{sites-cohomology-lemma-compare-cohomology-general}
在情形 \ref{sites-cohomology-situation-compare} 中，设 $X$ 是 $\mathcal{C}$ 中的对象。
\begin{enumerate}
\item 对 $\mathcal{F}' \in \mathcal{A}'_X$，有
$H^n_{\tau'}(X, \mathcal{F}') = H^n_\tau(X, \epsilon_X^{-1}\mathcal{F}')$,
\item 对 $K' \in D^+_{\mathcal{A}'_X}(\mathcal{C}_{\tau'}/X)$，
有 $H^n_{\tau'}(X, K') = H^n_\tau(X, \epsilon_X^{-1}K')$。
\end{enumerate}
\end{lemma}

\begin{proof}
由引理 \ref{sites-cohomology-lemma-V-C-all-n-general}
以及评注 \ref{sites-cohomology-remark-before-Leray} 得出。
\end{proof}
























\section{豪斯多夫且局部拟紧空间上的上同调}
\label{sites-cohomology-section-cohomology-LC}

\noindent
我们沿用“豪斯多夫且局部拟紧”的说法，
而不采用文献中常见的“局部紧”说法。
记 $\textit{LC}$ 为如下范畴：其对象是豪斯多夫且
局部拟紧的拓扑空间，
其态射是连续映射。


\begin{lemma}
\label{sites-cohomology-lemma-LC-basic}
范畴 $\textit{LC}$ 有纤维积和终对象，因而有任意有限极限。若态射 $X \to Z$ 与 $Y \to Z$
在 $\textit{LC}$ 中，若
$X$ 与 $Y$ 拟紧，则
$X \times_Z Y$ 拟紧。
\end{lemma}

\begin{proof}
终对象是单点空间。给定 $X \to Z$ 与
$Y \to Z$ 这两个 $\textit{LC}$ 中的态射，纤维积 $X \times_Z Y$ 是
$X \times Y$ 的子空间。因此 $X \times_Z Y$ 是豪斯多夫的，因为
$X \times Y$ 由
拓扑，第 \ref{topology-section-Hausdorff} 节知是豪斯多夫的。

\medskip\noindent
若 $X$ 与 $Y$ 拟紧，则由
拓扑，定理 \ref{topology-theorem-tychonov}，$X \times Y$ 拟紧。
由于 $X \times_Z Y$ 是 $X \times Y$ 的闭子集
（拓扑，引理 \ref{topology-lemma-fibre-product-closed}），
由此根据
拓扑，引理 \ref{topology-lemma-closed-in-quasi-compact}，$X \times_Z Y$ 拟紧。

\medskip\noindent
最后回到一般情形。若 $x \in X$ 且 $y \in Y$，
可取拟紧邻域 $x \in E \subset X$ 以及
$y \in F \subset Y$，由上述结果可知 $E \times_Z F$ 是
$(x, y)$ 的拟紧邻域。因此 $X \times_Z Y$
由
拓扑，引理 \ref{topology-lemma-locally-quasi-compact-Hausdorff}，是 $\textit{LC}$ 中的对象。
\end{proof}

\noindent
我们可以在 $\textit{LC}$ 上赋予比通常拓扑更强的拓扑。

\begin{definition}
\label{sites-cohomology-definition-covering-LC}
设 $\{f_i : X_i \to X\}$ 是一族目标固定的态射，
其目标位于范畴 $\textit{LC}$ 中。我们称此族为
{\it qc 覆盖}\footnote{这是非标准记号。
我们采用它是为了使读者联想到概形的 fpqc 覆盖。}
若对每个 $x \in X$，存在 $i_1, \ldots, i_n \in I$ 以及
拟紧子集 $E_j \subset X_{i_j}$，使得
$\bigcup f_{i_j}(E_j)$ 是 $x$ 的一个邻域。
\end{definition}

\noindent
注意，$X = \bigcup U_i$ 作为 $\textit{LC}$ 中对象的开覆盖
给出 qc 覆盖 $\{U_i \to X\}$，因为 $X$ 是局部拟紧的。
我们先给出必要的引理。

\begin{lemma}
\label{sites-cohomology-lemma-qc}
设 $X$ 是豪斯多夫且局部拟紧的空间，换言之，
是 $\textit{LC}$ 中的对象。
\begin{enumerate}
\item 若 $X' \to X$ 是 $\textit{LC}$ 中的同构，则
$\{X' \to X\}$ 是 qc 覆盖。
\item 若 $\{f_i : X_i \to X\}_{i\in I}$ 是 qc 覆盖，且对每个
$i$ 有 qc 覆盖 $\{g_{ij} : X_{ij} \to X_i\}_{j\in J_i}$，则
$\{X_{ij} \to X\}_{i \in I, j\in J_i}$ 是 qc 覆盖。
\item 若 $\{X_i \to X\}_{i\in I}$ 是 qc 覆盖，
且 $X' \to X$ 是 $\textit{LC}$ 中的态射，则
$\{X' \times_X X_i \to X'\}_{i\in I}$ 是 qc 覆盖。
\end{enumerate}
\end{lemma}

\begin{proof}
第 (1) 项由上面关于开覆盖是 qc 覆盖的评注立即得到。

\medskip\noindent
证明 (2)。设 $x \in X$。选取 $i_1, \ldots, i_n \in I$ 以及
$E_a \subset X_{i_a}$ 为拟紧集，使得 $\bigcup f_{i_a}(E_a)$
是 $x$ 的邻域。对每个 $e \in E_a$，可找到
有限集 $J_e \subset J_{i_a}$ 以及拟紧集
$F_{e, j} \subset X_{ij}$（$j \in J_e$），使得 $\bigcup g_{ij}(F_{e, j})$
是 $e$ 的邻域。由于 $E_a$ 拟紧，可找到
有限个点 $e_1, \ldots, e_{m_a}$，使得
$$
E_a \subset
\bigcup\nolimits_{k = 1, \ldots, m_a}
\bigcup\nolimits_{j \in J_{e_k}} g_{ij}(F_{e_k, j})
$$
于是得到
$$
\bigcup\nolimits_{a = 1, \ldots, n}
\bigcup\nolimits_{k = 1, \ldots, m_a}
\bigcup\nolimits_{j \in J_{e_k}} f_i(g_{ij}(F_{e_k, j}))
$$
是 $x$ 的邻域。

\medskip\noindent
证明 (3)。设 $x' \in X'$ 是一个点，设其像 $x \in X$。
选取 $i_1, \ldots, i_n \in I$ 以及拟紧子集
$E_j \subset X_{i_j}$，使得 $\bigcup f_{i_j}(E_j)$ 是一个
$x$ 的邻域。选取 $X'$ 中 $x'$ 的拟紧邻域 $F \subset X'$，
其映入 $x$ 的拟紧邻域 $\bigcup f_{i_j}(E_j)$。于是
$F \times_X E_j \subset X' \times_X X_{i_j}$ 是
拟紧子集，而 $F$ 是态射
$\coprod F \times_X E_j \to F$ 的像。因此基变换是
qc 覆盖，证明完成。
\end{proof}

\noindent
由于 $\textit{LC}$ 的所有对象都是豪斯多夫空间，$\textit{LC}$ 中的任一态射
$f : X \to Y$ 都是拓扑空间之间的分离连续映射。于是 $f$ 是拓扑空间中的 proper 映射
当且仅当它是泛闭的。参见
拓扑，第 \ref{topology-section-proper} 节的讨论。


\begin{lemma}
\label{sites-cohomology-lemma-proper-surjective-is-qc-covering}
设 $f : X \to Y$ 是 $\textit{LC}$ 中的态射。
若 $f$ proper 且满射，则 $\{f : X \to Y\}$
是 qc 覆盖。
\end{lemma}

\begin{proof}
设 $y \in Y$ 是一个点。对每个 $x \in X_y$，选取一个拟紧
邻域 $E_x \subset X$。选取开集 $x \in U_x \subset E_x$。
由于 $f$ proper，纤维 $X_y$ 拟紧，因此可找到
$x_1, \ldots, x_n \in X_y$ 使得
$X_y \subset U_{x_1} \cup \ldots \cup U_{x_n}$。
我们断言 $f(E_{x_1}) \cup \ldots \cup f(E_{x_n})$ 是 $y$ 的邻域。
事实上，由于 $f$ 是闭映射
（拓扑，定理 \ref{topology-theorem-characterize-proper}）
可知 $Z = f(X \setminus U_{x_1} \cup \ldots \cup U_{x_n})$
是 $Y$ 中不含 $y$ 的闭子集。由于 $f$ 满射，
可知 $Y \setminus Z$ 包含于
$f(E_{x_1}) \cup \ldots \cup f(E_{x_n})$，正如所需。
\end{proof}
\noindent
除集合论方面的一些问题外，引理 \ref{sites-cohomology-lemma-qc} 表明
$\textit{LC}$ 与 qc 覆盖的集合构成一个站点。
我们将把这个站点（适当修改以克服
集合论问题）记为
$\textit{LC}_{qc}$。

\begin{remark}[集合论问题]
\label{sites-cohomology-remark-set-theoretic-LC}
范畴 $\textit{LC}$ 是一个“大的”范畴，因为其对象构成
真类；覆盖同样构成真类。
我们定义拓扑空间 $X$ 的 {\it size} 为
其点集的基数。选取基数上的函数
$Bound$，例如见集合，
等式 (\ref{sets-equation-bound})。
最后，令 $S_0$ 是 $\textit{LC}$ 中对象的初始集合，
例如 $S_0 = \{(\mathbf{R}, \text{欧几里得拓扑})\}$。
完全类似于集合，引理 \ref{sets-lemma-construct-category}
可选取极限序数 $\alpha$，使得
$\textit{LC}_\alpha = \textit{LC} \cap V_\alpha$
包含 $S_0$，并且在 $\textit{LC}$ 中存在的所有可数极限与
余极限下保持不变。此外，若 $X \in \textit{LC}_\alpha$
且 $Y \in \textit{LC}$，并且
$\text{size}(Y) \leq Bound(\text{size}(X))$，则 $Y$ 同构于
$\textit{LC}_\alpha$ 中的某个对象。
接着应用集合，引理 \ref{sets-lemma-coverings-site}
在 $\textit{LC}_\alpha$ 上选取 qc 覆盖的集合 $\text{Cov}$，
使得 $\textit{LC}_\alpha$ 中的每个 qc 覆盖都
在组合意义下等价于该集合中的某个覆盖。
这样得到站点 $(\textit{LC}_\alpha, \text{Cov})$，
记为 $\textit{LC}_{qc}$。
\end{remark}

\noindent
$\textit{LC}_{qc}$ 站点（评注 \ref{sites-cohomology-remark-set-theoretic-LC}）上还有第二个拓扑。
具体地，对给定对象
$X$，考虑 $\textit{LC}_{qc}$ 的所有覆盖 $\{X_i \to X\}$，
其中 $X_i \to X$ 是开浸入。我们将此站点记为
$\textit{LC}_{Zar}$。恒等函子
$\textit{LC}_{Zar} \to \textit{LC}_{qc}$ 连续，并定义
一个站点态射
$$
\epsilon : \textit{LC}_{qc} \longrightarrow \textit{LC}_{Zar}
$$
参见第 \ref{sites-cohomology-section-compare} 节。
对于豪斯多夫且局部拟紧的拓扑空间 $X$，更准确地说，
对于 $X \in \Ob(\textit{LC}_{qc})$，记诱导的态射为
$$
\epsilon_X : \textit{LC}_{qc}/X \longrightarrow \textit{LC}_{Zar}/X
$$
（参见站点，引理 \ref{sites-lemma-localize-morphism}）。
令 $X_{Zar}$ 为对象是 $X$ 的开集的站点，参见
站点，例 \ref{sites-example-site-topological}。
存在站点态射
$$
\pi_X : \textit{LC}_{Zar}/X \longrightarrow X_{Zar}
$$
由连续函子
$X_{Zar} \to \textit{LC}_{Zar}/X$，$U \mapsto U$ 给出。
具体地，$X_{Zar}$ 有纤维积和终对象，而上述
函子与它们交换，因此站点，命题 \ref{sites-proposition-get-morphism} 适用。
我们常将 $\pi$ 看作拓扑斯态射
$$
\pi_X : \Sh(\textit{LC}_{Zar}/X) \longrightarrow \Sh(X)
$$
利用等式 $\Sh(X_{Zar}) = \Sh(X)$。

\begin{lemma}
\label{sites-cohomology-lemma-describe-pullback-pi}
设 $X$ 是 $\textit{LC}_{qc}$ 中的对象，设 $\mathcal{F}$ 是
$X$ 上的层。该规则
$$
\textit{LC}_{qc}/X \longrightarrow \textit{Sets},\quad
(f : Y \to X) \longmapsto \Gamma(Y, f^{-1}\mathcal{F})
$$
是一个层，因而当然也是 $\textit{LC}_{Zar}/X$ 上的层。
该层在 $\textit{LC}_{Zar}/X$ 上等于
$\pi_X^{-1}\mathcal{F}$，并在 $\textit{LC}_{qc}/X$ 上等于
$\epsilon_X^{-1}\pi_X^{-1}\mathcal{F}$。
\end{lemma}

\begin{proof}
记由引理中的公式定义的预层为 $\mathcal{G}$。
当然，公式中的拉回 $f^{-1}$ 表示拓扑空间上层的通常
拉回。由定义立即可知，$\mathcal{G}$ 是 Zar
拓扑下的层。

\medskip\noindent
设 $Y \to X$ 是 $\textit{LC}_{qc}$ 中的态射。设
$\mathcal{V} = \{g_i : Y_i \to Y\}_{i \in I}$ 是 qc 覆盖。
要证明 $\mathcal{G}$ 是 qc 拓扑下的层，
只需证明
$\mathcal{G}(Y) \to H^0(\mathcal{V}, \mathcal{G})$
是同构（参见站点，第 \ref{sites-section-sheafification} 节）。
首先注意该态射是单射，因为 qc 覆盖
是满射，并且可以在茎上检验截面的相等性；
（使用层，引理 \ref{sheaves-lemma-sheaf-subset-stalks} 以及
\ref{sheaves-lemma-stalk-pullback-presheaf}）。因此
$\mathcal{G}$ 是 $\textit{LC}_{qc}$ 上的分离预层，
所以只需证明任意元素
$(s_i) \in H^0(\mathcal{V}, \mathcal{G})$
在将 $\mathcal{V}$ 替换为某个加细后，
映到 $\mathcal{G}(Y)$ 的像中
（站点，定理 \ref{sites-theorem-plus}）。

\medskip\noindent
将 $Y_{i, Zar}$ 上的层与 $Y_i$ 上的层等同，我们看到
$\mathcal{G}|_{Y_{i, Zar}}$ 是连续映射 $g_i : Y_i \to Y$
下的 $f^{-1}\mathcal{F}$ 的拉回。因此可以选取一个开覆盖
$Y_i = \bigcup V_{ij}$，使得对每个 $j$，存在开集
$W_{ij} \subset Y$ 以及截面 $t_{ij} \in \mathcal{G}(W_{ij})$，
满足 $V_{ij}$ 映入 $W_{ij}$，并且
$s|_{V_{ij}}$ 是 $t_{ij}$ 的拉回。换言之，将覆盖
$\{Y_i \to Y\}$ 加细后，可以假设存在开集 $W_i \subset Y$，
使得 $Y_i \to Y$ 经由 $W_i$ 分解，并且在 $W_i$ 上有
$\mathcal{G}$ 的截面 $t_i$，其限制是给定的截面 $s_i$。
此外，若 $y \in Y$ 同时属于 $Y_i \to Y$ 和 $Y_j \to Y$ 的像，
则茎 $f^{-1}\mathcal{F}_y$ 中的像 $t_{i, y}$ 与 $t_{j, y}$ 相等
（因为 $s_i$ 与 $s_j$ 在 $Y_i \times_Y Y_j$ 上相等）。
因此，对 $y \in Y$，存在一个定义良好的元素 $t_y \in
f^{-1}\mathcal{F}_y$；当 $y$ 属于 $Y_i \to Y$ 的像时，
它与 $t_{i, y}$ 相等。我们将证明元素 $(t_y)$ 来自 $Y$ 上
$f^{-1}\mathcal{F}$ 的一个全局截面，这就完成引理的证明。

\medskip\noindent
只需证明该结论在 $Y$ 上局部成立；参见层，第
\ref{sheaves-section-sheafification} 节。令 $y_0 \in Y$。
选取 $i_1, \ldots, i_n \in I$ 以及拟紧子集
$E_j \subset Y_{i_j}$，使得 $\bigcup g_{i_j}(E_j)$ 是 $y_0$
的一个邻域。令 $V \subset Y$ 是 $y_0$ 的一个开邻域，且包含于
$\bigcup g_{i_j}(E_j)$ 以及 $W_{i_1} \cap \ldots \cap W_{i_n}$。
由于 $t_{i_1, y_0} = \ldots = t_{i_n, y_0}$，缩小 $V$ 后可假设
$f^{-1}\mathcal{F}$ 的截面 $t_{i_j}|_V$（$j = 1, \ldots, n$）
彼此相等。由于 $V \subset \bigcup g_{i_j}(E_j)$，可知
$(t_y)_{y \in V}$ 来自这个截面。

\medskip\noindent
还需证明，在 $\textit{LC}_{qc}$ 上，$\mathcal{G}$ 等于
$\epsilon_X^{-1}\pi_X^{-1}\mathcal{F}$；分别地，在
$\textit{LC}_{Zar}$ 上等于 $\pi_X^{-1}\mathcal{F}$。
在这两种情形中，拉回都由取下述预层来定义：
$$
(f : Y \to X)
\longmapsto
\colim_{f(Y) \subset U \subset X} \mathcal{F}(U)
$$
然后作层化。在 Zar 拓扑下层化恰好得到我们的层
$\mathcal{G}$；又因为 $\mathcal{G}$ 是 qc 层，故在 qc 拓扑下
同样成立。
\end{proof}

\noindent
令 $X \in \Ob(\textit{LC}_{Zar})$，令 $\mathcal{H}$
是 $\textit{LC}_{Zar}/X$ 上的阿贝尔层。
对于 $\textit{LC}_{Zar}/X$ 的对象 $U$，我们将用
$H^n_{Zar}(U, \mathcal{H})$ 表示 $\mathcal{H}$ 在 $U$ 上的上同调。

\begin{lemma}
\label{sites-cohomology-lemma-collect-true-things-Zar}
令 $X$ 是 $\textit{LC}_{Zar}$ 中的对象。则
\begin{enumerate}
\item 对于 $\mathcal{F} \in \textit{Ab}(X)$，有
$H^n_{Zar}(X, \pi_X^{-1}\mathcal{F}) = H^n(X, \mathcal{F})$,
\item $\pi_{X, *} : \textit{Ab}(\textit{LC}_{Zar}/X) \to \textit{Ab}(X)$
是正合的，
\item 伴随的单位 $\text{id} \to \pi_{X, *} \circ \pi_X^{-1}$
是同构，
\item 对于 $K \in D(X)$，规范映射
$K \to R\pi_{X, *} \pi_X^{-1}K$ 是同构。
\end{enumerate}
令 $f : X \to Y$ 是 $\textit{LC}_{Zar}$ 中的态射。则
\begin{enumerate}
\item[(5)] 存在交换图
$$
\xymatrix{
\Sh(\textit{LC}_{Zar}/X) \ar[r]_{f_{Zar}} \ar[d]_{\pi_X} &
\Sh(\textit{LC}_{Zar}/Y) \ar[d]^{\pi_Y} \\
\Sh(X_{Zar}) \ar[r]^f &
\Sh(Y_{Zar})
}
$$
其中对象是拓扑斯，
\item[(6)] 对于 $L \in D^+(Y)$，有
$H^n_{Zar}(X, \pi_Y^{-1}L) = H^n(X, f^{-1}L)$,
\item[(7)] 若 $f$ 是 proper 映射，则有
\begin{enumerate}
\item $\pi_Y^{-1} \circ f_* = f_{Zar, *} \circ \pi_X^{-1}$ 是从
$\Sh(X)$ 到 $\Sh(\textit{LC}_{Zar}/Y)$ 的函子，
\item $\pi_Y^{-1} \circ Rf_* = Rf_{Zar, *} \circ \pi_X^{-1}$ 是从
$D^+(X)$ 到 $D^+(\textit{LC}_{Zar}/Y)$ 的函子。
\end{enumerate}
\end{enumerate}
\end{lemma}

\begin{proof}
证明 (1)。等式 $H^n_{Zar}(X, \pi_X^{-1}\mathcal{F}) = H^n(X, \mathcal{F})$
是一个一般事实，来自如下显然观察：$\textit{LC}_{Zar}$ 中 $X$ 的覆盖
与 $X$ 的开覆盖是同一回事。若希望看到详细证明，
应将引理 \ref{sites-cohomology-lemma-cohomology-bigger-site} 应用于函子
$X_{Zar} \to \textit{LC}_{Zar}$。

\medskip\noindent
证明 (2)。这是因为存在拓扑斯态射
$\tau_X : \Sh(X_{Zar}) \to \Sh(\textit{LC}_{Zar})$，使得
$\pi_{X, *} = \tau_X^{-1}$；这一点由将站点，引理
\ref{sites-lemma-bigger-site} 应用于定义 $\pi_X$ 所用的函子
$X_{Zar} \to \textit{LC}_{Zar}/X$ 得到。

\medskip\noindent
证明 (3)。这是因为根据站点，引理
\ref{sites-lemma-bigger-site}，$\tau_X^{-1} \circ \pi_X^{-1}$
是恒等函子。也可以从引理
\ref{sites-cohomology-lemma-describe-pullback-pi} 中对 $\pi_X^{-1}$ 的显式描述推出。

\medskip\noindent
证明 (4)。将 (3) 应用于表示 $K$ 的阿贝尔层复形。

\medskip\noindent
证明 (5)。拓扑斯态射 $f_{Zar}$ 来自应用站点，引理
\ref{sites-lemma-relocalize}；在此情形中，根据站点，引理
\ref{sites-lemma-relocalize-given-fibre-products}，它来自连续函子
$Z/Y \mapsto Z \times_Y X/X$。该图交换只是因为相应的连续函子
正确复合（参见站点，引理
\ref{sites-lemma-composition-morphisms-sites}）。

\medskip\noindent
证明 (6)。对于 $\mathcal{G} \in \textit{Ab}(Y)$，由引理
\ref{sites-cohomology-lemma-cohomology-of-open}，有
$H^n_{Zar}(X, \pi_Y^{-1}\mathcal{G}) =
H^n_{Zar}(X, f_{Zar}^{-1}\pi_Y^{-1}\mathcal{G})$。
根据 (5) 中图的交换性，这等于
$H^n_{Zar}(X, \pi_X^{-1}f^{-1}\mathcal{G})$。
当 $L$ 是集中在次数 $0$ 的单个层时，由 (1) 得到结论。
一般情形则通过用有下界的阿贝尔层复形表示
$L$ 得到。

\medskip\noindent
证明 (7a)。令 $\mathcal{F}$ 是 $X$ 上的层。
令 $g : Z \to Y$ 是 $\textit{LC}_{Zar}/Y$ 的对象。考虑纤维积
$$
\xymatrix{
Z' \ar[r]_{f'} \ar[d]_{g'} & Z \ar[d]^g \\
X \ar[r]^f & Y
}
$$
则有
$$
(f_{Zar, *}\pi_X^{-1}\mathcal{F})(Z/Y) =
(\pi_X^{-1}\mathcal{F})(Z'/X)  =
\Gamma(Z', (g')^{-1}\mathcal{F})  =
\Gamma(Z, f'_*(g')^{-1}\mathcal{F})
$$
第二个等号来自引理 \ref{sites-cohomology-lemma-describe-pullback-pi}。
另一方面，
$$
(\pi_Y^{-1}f_*\mathcal{F})(Z/Y) = \Gamma(Z, g^{-1}f_*\mathcal{F})
$$
同样由引理 \ref{sites-cohomology-lemma-describe-pullback-pi} 得到。
因此，根据集合层的真基变换定理
（上同调，引理 \ref{cohomology-lemma-proper-base-change-sheaves-of-sets}），
可知这两个集合规范同构。该同构与限制映射相容，
并定义出同构
$\pi_Y^{-1}f_*\mathcal{F} = f_{Zar, *}\pi_X^{-1}\mathcal{F}$.
从而得到函子同构
$\pi_Y^{-1} \circ f_* = f_{Zar, *} \circ \pi_X^{-1}$.

\medskip\noindent
证明 (7b)。令 $K \in D^+(X)$。根据引理
\ref{sites-cohomology-lemma-unbounded-describe-higher-direct-images}，
$Rf_{Zar, *}\pi_X^{-1}K$ 的第 $n$ 个上同调层是与下述预层
相伴的层：
$$
(g : Z \to Y) \longmapsto H^n_{Zar}(Z', \pi_X^{-1}K)
$$
记号如上。注意
\begin{align*}
H^n_{Zar}(Z', \pi_X^{-1}K)
& =
H^n(Z', (g')^{-1}K) \\
& =
H^n(Z, Rf'_*(g')^{-1}K) \\
& =
H^n(Z, g^{-1}Rf_*K) \\
& =
H^n_{Zar}(Z, \pi_Y^{-1}Rf_*K)
\end{align*}
第一个等式是将 (6) 应用于 $K$ 和 $g' : Z' \to X$ 的结果。
第二个等式是对 $f' : Z' \to Z$ 应用 Leray
（上同调，引理 \ref{cohomology-lemma-before-Leray}）。
第三个等式是真基变换定理
（上同调，定理 \ref{cohomology-theorem-proper-base-change}）。
第四个等式是将 (6) 应用于 $g : Z \to Y$ 和 $Rf_*K$ 的结果。
因此 $Rf_{Zar, *}\pi_X^{-1}K$ 与 $\pi_Y^{-1}Rf_*K$ 具有相同的
上同调层。我们省略对规范基变换映射
$\pi_Y^{-1}Rf_*K \to Rf_{Zar, *}\pi_X^{-1}K$
诱导出该同构的验证。
\end{proof}

\noindent
在引理 \ref{sites-cohomology-lemma-describe-pullback-pi} 的情形中，
$\epsilon$ 与 $\pi$ 的复合以及等式
$\Sh(X) = \Sh(X_{Zar})$ 确定一个拓扑斯态射
$$
a_X : \Sh(\textit{LC}_{qc}/X) \longrightarrow \Sh(X)
$$

\begin{lemma}
\label{sites-cohomology-lemma-push-pull-LC}
令 $f : X \to Y$ 是 $\textit{LC}_{qc}$ 中的态射。
则存在拓扑斯的交换图
$$
\vcenter{
\xymatrix{
\Sh(\textit{LC}_{qc}/X) \ar[r]_{f_{qc}} \ar[d]_{\epsilon_X} &
\Sh(\textit{LC}_{qc}/Y) \ar[d]^{\epsilon_Y} \\
\Sh(\textit{LC}_{Zar}/X) \ar[r]^{f_{Zar}} &
\Sh(\textit{LC}_{Zar}/Y)
}
}
\quad\text{和}\quad
\vcenter{
\xymatrix{
\Sh(\textit{LC}_{qc}/X) \ar[r]_{f_{qc}} \ar[d]_{a_X} &
\Sh(\textit{LC}_{qc}/Y) \ar[d]^{a_Y} \\
\Sh(X) \ar[r]^f &
\Sh(Y)
}
}
$$
其中 $a_X = \pi_X \circ \epsilon_X$，$a_Y = \pi_X \circ \epsilon_X$。
若 $f$ 是 proper 映射，则
$a_Y^{-1} \circ f_* = f_{qc, *} \circ a_X^{-1}$。
\end{lemma}

\begin{proof}
拓扑斯态射 $f_{qc}$ 是站点，引理
\ref{sites-lemma-relocalize} 中的那个态射；在此情形中，
它来自连续函子 $Z/Y \mapsto Z \times_Y X/X$，参见站点，
引理 \ref{sites-lemma-relocalize-given-fibre-products}。
左图交换是因为相应的连续函子正确复合
（参见站点，引理 \ref{sites-lemma-composition-morphisms-sites}）。
右图交换是因为左图交换，并且因为引理
\ref{sites-cohomology-lemma-collect-true-things-Zar} 的 (5) 部分。

\medskip\noindent
证明最后的断言。读者可以重复引理
\ref{sites-cohomology-lemma-collect-true-things-Zar} 的 (7a) 部分的证明；
这里我们改为从该部分推出结论。由于 $\epsilon_{Y, *}$ 是底层预层
上的恒等函子，它反映同构。引理
\ref{sites-cohomology-lemma-describe-pullback-pi} 中的描述表明
$\epsilon_{Y, *} \circ a_Y^{-1} = \pi_Y^{-1}$，对 $X$ 也同样成立。
要证明规范映射
$a_Y^{-1}f_*\mathcal{F} \to f_{qc, *}a_X^{-1}\mathcal{F}$
是同构，只需证明
$$
\pi_Y^{-1}f_*\mathcal{F} =
\epsilon_{Y, *}a_Y^{-1}f_*\mathcal{F} \to
\epsilon_{Y, *}f_{qc, *}a_X^{-1}\mathcal{F} =
f_{Zar, *}\epsilon_{X, *}a_X^{-1}\mathcal{F} =
f_{Zar, *}\pi_X^{-1}\mathcal{F}
$$
是同构。这正是引理
\ref{sites-cohomology-lemma-collect-true-things-Zar} 的 (7a) 部分。
\end{proof}

\begin{lemma}
\label{sites-cohomology-lemma-compare-qc-zar}
考虑比较态射
$\epsilon : \textit{LC}_{qc} \to \textit{LC}_{Zar}$。
令 $\mathcal{P}$ 表示拓扑空间的 proper 映射类。
对 $\textit{LC}_{Zar}$ 中的 $X$，记
$\mathcal{A}'_X \subset \textit{Ab}(\textit{LC}_{Zar}/X)$
为由形如 $\pi_X^{-1}\mathcal{F}$ 的层组成的全子范畴，
其中 $\mathcal{F} \in \textit{Ab}(X)$。
则
(\ref{sites-cohomology-item-base-change-P}),
(\ref{sites-cohomology-item-restriction-A}),
(\ref{sites-cohomology-item-A-sheaf}),
(\ref{sites-cohomology-item-A-and-P}) 和
(\ref{sites-cohomology-item-refine-tau-by-P})
这五个条件（情形 \ref{sites-cohomology-situation-compare}）均成立。
\end{lemma}

\begin{proof}
首先，通过检查同调，引理
\ref{homology-lemma-characterize-weak-serre-subcategory} 的条件
(1)、(2)、(3)、(4)，证明
$\mathcal{A}'_X \subset \textit{Ab}(\textit{LC}_{Zar}/X)$
是弱 Serre 子范畴。(1)、(2)、(3) 立即成立，因为根据引理
\ref{sites-cohomology-lemma-collect-true-things-Zar} 的 (3) 部分，$\pi_X^{-1}$
是正合且全忠实的。若
$0 \to \pi_X^{-1}\mathcal{F} \to \mathcal{G} \to \pi_X^{-1}\mathcal{F}' \to 0$
是 $\textit{Ab}(\textit{LC}_{Zar}/X)$ 中的短正合序列，则根据引理
\ref{sites-cohomology-lemma-collect-true-things-Zar} 的 (2) 部分，
$0 \to \mathcal{F} \to \pi_{X, *}\mathcal{G} \to \mathcal{F}' \to 0$
正合。因此 $\mathcal{G} = \pi_X^{-1}\pi_{X, *}\mathcal{G}$
属于 $\mathcal{A}'_X$，最后一个条件也得到验证。

\medskip\noindent
性质 (\ref{sites-cohomology-item-base-change-P}) 由引理 \ref{sites-cohomology-lemma-LC-basic} 以及
proper 映射的基变换仍是 proper 映射这一事实成立（参见拓扑，
定理 \ref{topology-theorem-characterize-proper} 和引理
\ref{topology-lemma-base-change-separated}）。

\medskip\noindent
性质 (\ref{sites-cohomology-item-restriction-A}) 来自引理
\ref{sites-cohomology-lemma-collect-true-things-Zar} 中的交换图 (5)。

\medskip\noindent
性质 (\ref{sites-cohomology-item-A-sheaf}) 就是引理 \ref{sites-cohomology-lemma-describe-pullback-pi}。

\medskip\noindent
性质 (\ref{sites-cohomology-item-A-and-P}) 是引理
\ref{sites-cohomology-lemma-collect-true-things-Zar} 的 (7)(b) 部分。

\medskip\noindent
证明 (\ref{sites-cohomology-item-refine-tau-by-P})。设给定一个 qc 覆盖
$\{U_i \to U\}$。对 $u \in U$，选取 $i_1, \ldots, i_m \in I$
以及拟紧子集 $E_j \subset U_{i_j}$，使得
$\bigcup f_{i_j}(E_j)$ 是 $u$ 的邻域。
注意，$Y = \coprod_{j = 1, \ldots, m} E_j \to U$
作为豪斯多夫拟紧空间之间的连续映射是 proper 的
（拓扑，引理 \ref{topology-lemma-closed-map}）。
选取包含于 $\bigcup f_{i_j}(E_j)$ 的开邻域 $u \in V$。
则 $Y \times_U V \to V$ 是满射 proper 态射，因而根据引理
\ref{sites-cohomology-lemma-proper-surjective-is-qc-covering} 是 $qc$ 覆盖。
对每个 $u \in U$ 都能如此处理，因此
(\ref{sites-cohomology-item-refine-tau-by-P}) 成立。
\end{proof}

\begin{lemma}
\label{sites-cohomology-lemma-V-C-all-n}
记号如上。
\begin{enumerate}
\item 对于 $X \in \Ob(\textit{LC}_{qc})$ 以及 $X$ 上的阿贝尔层
$\mathcal{F}$，有 $\epsilon_{X, *}a_X^{-1}\mathcal{F} = \pi_X^{-1}\mathcal{F}$，
并且当 $i > 0$ 时 $R^i\epsilon_{X, *}(a_X^{-1}\mathcal{F}) = 0$。
\item 对于 $\textit{LC}_{qc}$ 中的 proper 态射 $f : X \to Y$
以及 $X$ 上的阿贝尔层 $\mathcal{F}$，对所有 $i$ 有
$a_Y^{-1}(R^if_*\mathcal{F}) = R^if_{qc, *}(a_X^{-1}\mathcal{F})$。
\item 对于 $X \in \Ob(\textit{LC}_{qc})$ 以及 $D^+(X)$ 中的 $K$，
映射 $\pi_X^{-1}K \to R\epsilon_{X, *}(a_X^{-1}K)$ 是同构。
\item 对于 $\textit{LC}_{qc}$ 中的 proper 态射 $f : X \to Y$
以及 $D^+(X)$ 中的 $K$，有
$a_Y^{-1}(Rf_*K) = Rf_{qc, *}(a_X^{-1}K)$。
\end{enumerate}
\end{lemma}

\begin{proof}
根据引理 \ref{sites-cohomology-lemma-compare-qc-zar}，第
\ref{sites-cohomology-section-compare-general} 节中的引理都适用于当前情形。
为了将这些结果翻译过来，注意引理 \ref{sites-cohomology-lemma-A} 中的范畴
$\mathcal{A}_X$ 是下述函子的本质像：
$a_X^{-1} : \textit{Ab}(X) \to \textit{Ab}(\textit{LC}_{qc}/X)$。

\medskip\noindent
第 (1) 部分等价于对所有 $n$ 的 $(V_n)$，而这由引理
\ref{sites-cohomology-lemma-V-C-all-n-general} 成立。

\medskip\noindent
第 (2) 部分由将 $\epsilon_Y^{-1}$ 应用于引理
\ref{sites-cohomology-lemma-V-implies-C-general} 的结论得到。

\medskip\noindent
第 (3) 部分由引理 \ref{sites-cohomology-lemma-V-C-all-n-general} 的 (1) 部分得到，
因为 $\pi_X^{-1}K$ 属于
$D^+_{\mathcal{A}'_X}(\textit{LC}_{Zar}/X)$，且
$a_X^{-1} = \epsilon_X^{-1} \circ a_X^{-1}$。

\medskip\noindent
第 (4) 部分出于同样理由由引理
\ref{sites-cohomology-lemma-V-C-all-n-general} 的 (2) 部分得到。
\end{proof}

\begin{lemma}
\label{sites-cohomology-lemma-cohomological-descent-LC}
令 $X$ 是 $\textit{LC}_{qc}$ 中的对象。对于 $K \in D^+(X)$，映射
$$
K \longrightarrow Ra_{X, *}a_X^{-1}K
$$
是同构，其中 $a_X : \Sh(\textit{LC}_{qc}/X) \to \Sh(X)$ 如上所述。
\end{lemma}

\begin{proof}
首先将断言约化到 $K$ 由单个阿贝尔层给出的情形。具体地，
用有下界的复形 $\mathcal{F}^\bullet$ 表示 $K$。由单个层的情形可知
$\mathcal{F}^n = a_{X, *} a_X^{-1} \mathcal{F}^n$
并且当 $q > 0$ 时，层 $R^qa_{X, *}a_X^{-1}\mathcal{F}^n$
为零。将 Leray 无环引理（导出范畴，引理
\ref{derived-lemma-leray-acyclicity}）应用于 $a_X^{-1}\mathcal{F}^\bullet$
和函子 $a_{X, *}$，即可得到结论。以下设 $K = \mathcal{F}$。

\medskip\noindent
根据引理 \ref{sites-cohomology-lemma-describe-pullback-pi}，有
$a_{X, *}a_X^{-1}\mathcal{F} = \mathcal{F}$。因此只需证明当
$q > 0$ 时 $R^qa_{X, *}a_X^{-1}\mathcal{F} = 0$。
为此可使用 $a_X = \epsilon_X \circ \pi_X$ 以及 Leray 谱序列，
引理 \ref{sites-cohomology-lemma-relative-Leray}。根据引理
\ref{sites-cohomology-lemma-V-C-all-n}，当 $i > 0$ 时
$R^i\epsilon_{X, *}(a_X^{-1}\mathcal{F}) = 0$，并且
$\epsilon_{X, *}a_X^{-1}\mathcal{F} = \pi_X^{-1}\mathcal{F}$。
根据引理 \ref{sites-cohomology-lemma-collect-true-things-Zar}，当 $j > 0$ 时
$R^j\pi_{X, *}(\pi_X^{-1}\mathcal{F}) = 0$。
证明完成。
\end{proof}

\begin{lemma}
\label{sites-cohomology-lemma-compare-cohomology-LC}
令 $X \in \Ob(\textit{LC}_{qc})$，并令
$a_X : \Sh(\textit{LC}_{qc}/X) \to \Sh(X)$ 如上所述：
\begin{enumerate}
\item 对于 $X$ 上的阿贝尔层 $\mathcal{F}$，有
$H^n(X, \mathcal{F}) = H^n_{qc}(X, a_X^{-1}\mathcal{F})$,
\item 对于 $K \in D^+(X)$，有 $H^n(X, K) = H^n_{qc}(X, a_X^{-1}K)$。
\end{enumerate}
例如，若 $A$ 是阿贝尔群，则
$H^n(X, \underline{A}) = H^n_{qc}(X, \underline{A})$.
\end{lemma}

\begin{proof}
由引理 \ref{sites-cohomology-lemma-cohomological-descent-LC} 和评注
\ref{sites-cohomology-remark-before-Leray} 得到。
\end{proof}











\section{Ext 的谱序列}
\label{sites-cohomology-section-spectral-sequence-ext}

\noindent
本节汇集考虑 Ext 函子时出现的各种谱序列。对于环化站点
$(\mathcal{C}, \mathcal{O})$ 上的任意一对模复形
$\mathcal{G}^\bullet, \mathcal{F}^\bullet$，我们记
$$
\Ext^n_\mathcal{O}(\mathcal{G}^\bullet, \mathcal{F}^\bullet)
=
\Hom_{D(\mathcal{O})}(\mathcal{G}^\bullet, \mathcal{F}^\bullet[n])
$$
按照导出范畴第 \ref{derived-section-ext} 节中的一般约定。

\begin{example}
\label{sites-cohomology-example-hom-complex-into-sheaf}
令 $(\mathcal{C}, \mathcal{O})$ 是环化站点。
令 $\mathcal{K}^\bullet$ 是有上界的 $\mathcal{O}$-模复形，
令 $\mathcal{F}$ 是 $\mathcal{O}$-模。则存在以 $E_2$ 页为
$$
E_2^{i, j} =
\Ext_\mathcal{O}^i(H^{-j}(\mathcal{K}^\bullet), \mathcal{F})
\Rightarrow
\Ext_\mathcal{O}^{i + j}(\mathcal{K}^\bullet, \mathcal{F})
$$
以及另一个以 $E_1$ 页为
$$
E_1^{i, j} =
\Ext_\mathcal{O}^j(\mathcal{K}^{-i}, \mathcal{F})
\Rightarrow
\Ext_\mathcal{O}^{i + j}(\mathcal{K}^\bullet, \mathcal{F}).
$$
为构造这些谱序列，选取内射分解
$\mathcal{F} \to \mathcal{I}^\bullet$，并考虑来自双复形
$\Hom_\mathcal{O}(\mathcal{K}^\bullet, \mathcal{I}^\bullet)$
的两个谱序列，参见同调第 \ref{homology-section-double-complex} 节。
\end{example}








\section{杯积}
\label{sites-cohomology-section-cup-product}

\noindent
令 $(\mathcal{C}, \mathcal{O})$ 是环化站点。令
$K, M$ 是 $D(\mathcal{O})$ 中的对象。置
$A = \Gamma(\mathcal{C}, \mathcal{O})$。此情形中的（全局）杯积是
映射
$$
R\Gamma(\mathcal{C}, K) \otimes_A^\mathbf{L}
R\Gamma(\mathcal{C}, M)
\longrightarrow
R\Gamma(\mathcal{C}, K \otimes_\mathcal{O}^\mathbf{L} M)
$$
在 $D(A)$ 中。我们将其定义为环化拓扑斯态射
$(\Sh(\mathcal{C}), \mathcal{O}) \to (\Sh(pt), A)$
的相对杯积，见评注 \ref{sites-cohomology-remark-cup-product}。

\medskip\noindent
下面表述并证明相对杯积的一个自然相容性。设有环化拓扑斯态射
$f : (\Sh(\mathcal{C}), \mathcal{O}_\mathcal{C}) \to
(\Sh(\mathcal{D}), \mathcal{O}_\mathcal{D})$。
令 $\mathcal{K}^\bullet$ 和 $\mathcal{M}^\bullet$
是 $\mathcal{O}_\mathcal{C}$-模复形。存在朴素杯积
$$
\text{Tot}(
f_*\mathcal{K}^\bullet
\otimes_{\mathcal{O}_\mathcal{D}}
f_*\mathcal{M}^\bullet)
\longrightarrow
f_*\text{Tot}(\mathcal{K}^\bullet
\otimes_{\mathcal{O}_\mathcal{C}}
\mathcal{M}^\bullet)
$$
我们断言这与相对杯积有关。

\begin{lemma}
\label{sites-cohomology-lemma-cup-compatible-with-naive}
在上述情形，下图交换：
$$
\xymatrix{
f_*\mathcal{K}^\bullet
\otimes_{\mathcal{O}_\mathcal{D}}^\mathbf{L}
f_*\mathcal{M}^\bullet \ar[r] \ar[d]
&
Rf_*\mathcal{K}^\bullet
\otimes_{\mathcal{O}_\mathcal{D}}^\mathbf{L}
Rf_*\mathcal{M}^\bullet \ar[d]^{\text{评注 \ref{sites-cohomology-remark-cup-product}}} \\
\text{Tot}(
f_*\mathcal{K}^\bullet
\otimes_{\mathcal{O}_\mathcal{D}}
f_*\mathcal{M}^\bullet) \ar[d]_{\text{朴素杯积}} &
Rf_*(\mathcal{K}^\bullet
\otimes_{\mathcal{O}_\mathcal{C}}^\mathbf{L}
\mathcal{M}^\bullet) \ar[d] \\
f_*\text{Tot}(\mathcal{K}^\bullet
\otimes_{\mathcal{O}_\mathcal{C}}
\mathcal{M}^\bullet) \ar[r] &
Rf_*\text{Tot}(\mathcal{K}^\bullet
\otimes_{\mathcal{O}_\mathcal{C}}
\mathcal{M}^\bullet)
}
$$
\end{lemma}

\begin{proof}
根据评注 \ref{sites-cohomology-remark-cup-product} 中的构造可知，沿图顺时针走一周，
映射
$$
f_*\mathcal{K}^\bullet
\otimes_{\mathcal{O}_\mathcal{D}}^\mathbf{L}
f_*\mathcal{M}^\bullet 
\longrightarrow
Rf_*\text{Tot}(\mathcal{K}^\bullet
\otimes_{\mathcal{O}_\mathcal{C}}
\mathcal{M}^\bullet)
$$
是下述映射的伴随：
\begin{align*}
Lf^*(f_*\mathcal{K}^\bullet
\otimes_{\mathcal{O}_\mathcal{D}}^\mathbf{L}
f_*\mathcal{M}^\bullet)
& =
Lf^*f_*\mathcal{K}^\bullet
\otimes_{\mathcal{O}_\mathcal{D}}^\mathbf{L}
Lf^*f_*\mathcal{M}^\bullet \\
& \to
Lf^*Rf_*\mathcal{K}^\bullet
\otimes_{\mathcal{O}_\mathcal{D}}^\mathbf{L}
Lf^*Rf_*\mathcal{M}^\bullet \\
& \to
\mathcal{K}^\bullet
\otimes_{\mathcal{O}_\mathcal{D}}^\mathbf{L}
\mathcal{M}^\bullet \\
& \to
\text{Tot}(\mathcal{K}^\bullet
\otimes_{\mathcal{O}_\mathcal{C}}
\mathcal{M}^\bullet)
\end{align*}
根据引理 \ref{sites-cohomology-lemma-adjoints-push-pull-compatibility}，这也等于
\begin{align*}
Lf^*(f_*\mathcal{K}^\bullet
\otimes_{\mathcal{O}_\mathcal{D}}^\mathbf{L}
f_*\mathcal{M}^\bullet)
& =
Lf^*f_*\mathcal{K}^\bullet
\otimes_{\mathcal{O}_\mathcal{D}}^\mathbf{L}
Lf^*f_*\mathcal{M}^\bullet \\
& \to
f^*f_*\mathcal{K}^\bullet
\otimes_{\mathcal{O}_\mathcal{D}}^\mathbf{L}
f^*f_*\mathcal{M}^\bullet \\
& \to
\mathcal{K}^\bullet
\otimes_{\mathcal{O}_\mathcal{D}}^\mathbf{L}
\mathcal{M}^\bullet \\
& \to
\text{Tot}(\mathcal{K}^\bullet
\otimes_{\mathcal{O}_\mathcal{C}}
\mathcal{M}^\bullet)
\end{align*}
沿图逆时针走一周，我们得到下述映射的伴随：
\begin{align*}
Lf^*(f_*\mathcal{K}^\bullet
\otimes_{\mathcal{O}_\mathcal{D}}^\mathbf{L}
f_*\mathcal{M}^\bullet)
& \to
Lf^*\text{Tot}(
f_*\mathcal{K}^\bullet
\otimes_{\mathcal{O}_\mathcal{D}}
f_*\mathcal{M}^\bullet) \\
& \to
Lf^*f_*\text{Tot}(\mathcal{K}^\bullet
\otimes_{\mathcal{O}_\mathcal{C}}
\mathcal{M}^\bullet) \\
& \to
Lf^*Rf_*\text{Tot}(\mathcal{K}^\bullet
\otimes_{\mathcal{O}_\mathcal{C}}
\mathcal{M}^\bullet) \\
& \to
\text{Tot}(\mathcal{K}^\bullet
\otimes_{\mathcal{O}_\mathcal{C}}
\mathcal{M}^\bullet)
\end{align*}
根据引理 \ref{sites-cohomology-lemma-adjoints-push-pull-compatibility}，这也等于
\begin{align*}
Lf^*(f_*\mathcal{K}^\bullet
\otimes_{\mathcal{O}_\mathcal{D}}^\mathbf{L}
f_*\mathcal{M}^\bullet)
& \to
Lf^*\text{Tot}(
f_*\mathcal{K}^\bullet
\otimes_{\mathcal{O}_\mathcal{D}}
f_*\mathcal{M}^\bullet) \\
& \to
Lf^*f_*\text{Tot}(\mathcal{K}^\bullet
\otimes_{\mathcal{O}_\mathcal{C}}
\mathcal{M}^\bullet) \\
& \to
f^*f_*\text{Tot}(\mathcal{K}^\bullet
\otimes_{\mathcal{O}_\mathcal{C}}
\mathcal{M}^\bullet) \\
& \to
\text{Tot}(\mathcal{K}^\bullet
\otimes_{\mathcal{O}_\mathcal{C}}
\mathcal{M}^\bullet)
\end{align*}
现在考察下图即可完成证明：
$$
\xymatrix{
Lf^*(f_*\mathcal{K}^\bullet
\otimes_{\mathcal{O}_\mathcal{D}}^\mathbf{L}
f_*\mathcal{M}^\bullet) \ar[d] \ar[rr] & &
Lf^*f_*\mathcal{K}^\bullet \otimes_{\mathcal{O}_\mathcal{C}}^\mathbf{L}
Lf^*f_*\mathcal{M}^\bullet \ar[d] \\
Lf^*\text{Tot}(
f_*\mathcal{K}^\bullet
\otimes_{\mathcal{O}_\mathcal{D}}
f_*\mathcal{M}^\bullet) \ar[d]_{\text{朴素}} \ar[r] &
f^*\text{Tot}(
f_*\mathcal{K}^\bullet
\otimes_{\mathcal{O}_\mathcal{D}}
f_*\mathcal{M}^\bullet) \ar[ldd]^{\text{朴素}} \ar[dd] &
f^*f_*\mathcal{K}^\bullet \otimes_{\mathcal{O}_\mathcal{C}}^\mathbf{L}
f^*f_*\mathcal{M}^\bullet \ar[dd] \ar[ldd] \\
Lf^*f_*\text{Tot}(\mathcal{K}^\bullet
\otimes_{\mathcal{O}_\mathcal{C}}
\mathcal{M}^\bullet) \ar[d] \\
f^*f_*\text{Tot}(\mathcal{K}^\bullet \otimes_{\mathcal{O}_\mathcal{C}}
\mathcal{M}^\bullet) \ar[rd] &
\text{Tot}(f^*f_*\mathcal{K}^\bullet \otimes_{\mathcal{O}_\mathcal{C}}
f^*f_*\mathcal{M}^\bullet) \ar[d] &
\mathcal{K}^\bullet \otimes_{\mathcal{O}_\mathcal{C}}^\mathbf{L}
\mathcal{M}^\bullet \ar[ld] \\
& \text{Tot}(\mathcal{K}^\bullet
\otimes_{\mathcal{O}_\mathcal{C}}
\mathcal{M}^\bullet)
}
$$
该图中的所有多边形都交换。顶部的多边形由引理
\ref{sites-cohomology-lemma-tensor-pull-compatibility} 交换。
含有两个朴素杯积的方形交换，是因为
$Lf^* \to f^*$ 对模复形函子性。
含有下述两个映射的方形同理：
$\mathcal{A}^\bullet \otimes^\mathbf{L} \mathcal{B}^\bullet \to
\text{Tot}(\mathcal{A}^\bullet \otimes \mathcal{B}^\bullet)$.
最后，剩余方形的交换性在复形层面成立，并可视为朴素杯积的定义
（利用 $f^*$ 与 $f_*$ 的伴随性）。证明完成，因为沿图外圈走一周
所得的正是上面给出的两个映射。
\end{proof}

\begin{lemma}
\label{sites-cohomology-lemma-cup-product-associative}
令 $f : (\Sh(\mathcal{C}), \mathcal{O}) \to (\Sh(\mathcal{C}'), \mathcal{O}')$
是环化拓扑斯态射。评注 \ref{sites-cohomology-remark-cup-product} 中的相对杯积按下述
意义结合：对 $D(\mathcal{O})$ 中任意 $K,L,M$，下图
$$
\xymatrix{
Rf_*K \otimes_{\mathcal{O}'}^\mathbf{L}
Rf_*L \otimes_{\mathcal{O}'}^\mathbf{L}
Rf_*M \ar[r] \ar[d] &
Rf_*(K \otimes_\mathcal{O}^\mathbf{L} L)
\otimes_{\mathcal{O}'}^\mathbf{L} Rf_*M \ar[d] \\
Rf_*K \otimes_{\mathcal{O}'}^\mathbf{L}
Rf_*(L \otimes_\mathcal{O}^\mathbf{L} M) \ar[r] &
Rf_*(K \otimes_\mathcal{O}^\mathbf{L} 
L \otimes_\mathcal{O}^\mathbf{L} M)
}
$$
在 $D(\mathcal{O}')$ 中交换。
\end{lemma}

\begin{proof}
沿任一侧走一周，我们得到下述显然映射的伴随：
\begin{align*}
Lf^*(Rf_*K \otimes_{\mathcal{O}'}^\mathbf{L}
Rf_*L \otimes_{\mathcal{O}'}^\mathbf{L}
Rf_*M) & =
Lf^*(Rf_*K) \otimes_\mathcal{O}^\mathbf{L}
Lf^*(Rf_*L) \otimes_\mathcal{O}^\mathbf{L}
Lf^*(Rf_*M) \\
& \to
K \otimes_\mathcal{O}^\mathbf{L} 
L \otimes_\mathcal{O}^\mathbf{L} M
\end{align*}
该映射位于 $D(\mathcal{O})$ 中。
\end{proof}

\begin{lemma}
\label{sites-cohomology-lemma-cup-product-commutative}
令 $f : (\Sh(\mathcal{C}), \mathcal{O}) \to (\Sh(\mathcal{C}'), \mathcal{O}')$
是环化拓扑斯态射。评注 \ref{sites-cohomology-remark-cup-product} 中的相对杯积按下述
意义交换：下图
$$
\xymatrix{
Rf_*K \otimes_{\mathcal{O}'}^\mathbf{L} Rf_*L \ar[r] \ar[d]_\psi &
Rf_*(K \otimes_\mathcal{O}^\mathbf{L} L) \ar[d]^{Rf_*\psi} \\
Rf_*L \otimes_{\mathcal{O}'}^\mathbf{L} Rf_*K \ar[r] &
Rf_*(L \otimes_\mathcal{O}^\mathbf{L} K)
}
$$
对 $D(\mathcal{O})$ 中所有 $K,L$，在 $D(\mathcal{O}')$ 中交换。
这里 $\psi$ 是导出范畴中的交换约束
（引理 \ref{sites-cohomology-lemma-symmetric-monoidal-derived}）。
\end{lemma}

\begin{proof}
略去证明。
\end{proof}

\begin{lemma}
\label{sites-cohomology-lemma-compose-cup-product}
令 $f : (\Sh(\mathcal{C}), \mathcal{O}) \to (\Sh(\mathcal{C}'), \mathcal{O}')$
以及 $f' : (\Sh(\mathcal{C}'), \mathcal{O}') \to
(\Sh(\mathcal{C}''), \mathcal{O}'')$ 是环化拓扑斯态射。
评注 \ref{sites-cohomology-remark-cup-product} 中的相对杯积与复合相容，
即下图
$$
\xymatrix{
R(f' \circ f)_*K \otimes_{\mathcal{O}''}^\mathbf{L} R(f' \circ f)_*L
\ar@{=}[rr] \ar[d] & &
Rf'_*Rf_*K \otimes_{\mathcal{O}''}^\mathbf{L} Rf'_*Rf_*L \ar[d] \\
R(f' \circ f)_*(K \otimes_\mathcal{O}^\mathbf{L} L) \ar@{=}[r] &
Rf'_*Rf_*(K \otimes_\mathcal{O}^\mathbf{L} L) &
Rf'_*(Rf_*K \otimes_{\mathcal{O}'}^\mathbf{L}  Rf_*L) \ar[l]
}
$$
对 $D(\mathcal{O})$ 中所有 $K,L$，在 $D(\mathcal{O}'')$ 中交换。
\end{lemma}

\begin{proof}
这是因为沿图的任一方向走一周，都得到下述映射的伴随：
\begin{align*}
& L(f' \circ f)^*\left(R(f' \circ f)_*K
\otimes_{\mathcal{O}''}^\mathbf{L}
R(f' \circ f)_*L\right) \\
& =
L(f' \circ f)^*R(f' \circ f)_*K
\otimes_\mathcal{O}^\mathbf{L}
L(f' \circ f)^*R(f' \circ f)_*L) \\
& \to
K \otimes_\mathcal{O}^\mathbf{L} L
\end{align*}
该映射位于 $D(\mathcal{O})$ 中。为此使用如下余单位的复合：
$$
L(f' \circ f)^*R(f' \circ f)_* =
Lf^* L(f')^* Rf'_* Rf_*  \to
Lf^* Rf_* \to \text{id}
$$
它是 $L(f' \circ f)^*$ 与 $R(f' \circ f)_*$ 的余单位。
参见范畴，引理 \ref{categories-lemma-compose-counits}。
\end{proof}

\begin{lemma}
\label{sites-cohomology-lemma-base-change-cup-product}
考虑交换方形
$$
\xymatrix{
(\Sh(\mathcal{C}'), \mathcal{O}_{\mathcal{C}'}) \ar[r]_{g'} \ar[d]_{f'} &
(\Sh(\mathcal{C}, \mathcal{O}_{\mathcal{C}}) \ar[d]^f \\
(\Sh(\mathcal{D}'), \mathcal{O}_{\mathcal{D}'}) \ar[r]^g &
(\Sh(\mathcal{D}), \mathcal{O}_{\mathcal{D}}) 
}
$$
它由环化拓扑斯组成。令 $K,L \in D(\mathcal{O}_{\mathcal{C}})$。
相对杯积与该方形相容，意味着下图
$$
\xymatrix{
Lg^*(Rf_*K \otimes_{\mathcal{O}_{\mathcal{D}}}^\mathbf{L} Rf_*L)
\ar[r] \ar@{=}[d] &
Lg^*(Rf_*(K \otimes_{\mathcal{O}_{\mathcal{C}}}^\mathbf{L} L)) \ar[d] \\
Lg^*Rf_*K \otimes_{\mathcal{O}_{\mathcal{D}'}}^\mathbf{L} Lg^*Rf_*L \ar[d] &
R(f')_*L(g')^*(K \otimes_{\mathcal{O}_{\mathcal{C}}}^\mathbf{L} L) \ar@{=}[d] \\
R(f')_*(L(g')^*K \otimes_{\mathcal{O}_{\mathcal{D}'}} R(f')_*(L(g')^*L \ar[r] &
R(f')_*(L(g')^*K \otimes_{\mathcal{O}_{\mathcal{C}'}}^\mathbf{L} L(g')^*L)
}
$$
在 $D(\mathcal{O}_{\mathcal{D}'})$ 中交换。
水平箭头由相对杯积（评注 \ref{sites-cohomology-remark-cup-product}）给出，
竖直箭头由基变换映射（评注 \ref{sites-cohomology-remark-base-change}）
和引理 \ref{sites-cohomology-lemma-pullback-tensor-product} 给出。
\end{lemma}

\begin{proof}
略去证明。
\end{proof}






\section{Hom 复形}
\label{sites-cohomology-section-hom-complexes}

\noindent
令 $(\mathcal{C}, \mathcal{O})$ 是环化站点。令
$\mathcal{L}^\bullet$ 和 $\mathcal{M}^\bullet$ 是两个
$\mathcal{O}$-模复形。我们构造一个 $\mathcal{O}$-模复形
$\SheafHom^\bullet(\mathcal{L}^\bullet, \mathcal{M}^\bullet)$。
具体地，对每个 $n$ 置
$$
\SheafHom^n(\mathcal{L}^\bullet, \mathcal{M}^\bullet) =
\prod\nolimits_{n = p + q}
\SheafHom_\mathcal{O}(\mathcal{L}^{-q}, \mathcal{M}^p)
$$
可以将 $\SheafHom^n$ 理解为如下 $\mathcal{O}$-模层：
所有从 $\mathcal{L}^\bullet$ 到 $\mathcal{M}^\bullet$ 的
$\mathcal{O}$-线性映射（视为分次 $\mathcal{O}$-模）中，
次数齐次为 $n$ 的映射。在此术语下，按下式定义微分：
$$
\text{d}(f) =
\text{d}_\mathcal{M} \circ f - (-1)^n f \circ \text{d}_\mathcal{L}
$$
其中
$f \in \SheafHom^n_\mathcal{O}(\mathcal{L}^\bullet, \mathcal{M}^\bullet)$。
我们省略 $\text{d}^2 = 0$ 的验证。
该构造是微分分次代数，例 \ref{dga-example-category-complexes} 的特例。
由构造立即得到
\begin{equation}
\label{sites-cohomology-equation-cohomology-hom-complex}
H^n(\Gamma(U, \SheafHom^\bullet(\mathcal{L}^\bullet, \mathcal{M}^\bullet))) =
\Hom_{K(\mathcal{O}_U)}(\mathcal{L}^\bullet|_U, \mathcal{M}^\bullet[n]|_U)
\end{equation}
对所有 $n \in \mathbf{Z}$ 以及每个 $U \in \Ob(\mathcal{C})$ 均成立。
类似地，对全局截面复形有
\begin{equation}
\label{sites-cohomology-equation-global-cohomology-hom-complex}
H^n(\Gamma(\mathcal{C},
\SheafHom^\bullet(\mathcal{L}^\bullet, \mathcal{M}^\bullet))) =
\Hom_{K(\mathcal{O})}(\mathcal{L}^\bullet, \mathcal{M}^\bullet[n])
\end{equation}
。

\begin{lemma}
\label{sites-cohomology-lemma-compose}
令 $(\mathcal{C}, \mathcal{O})$ 是环化站点。给定
$\mathcal{O}$-模复形 $\mathcal{K}^\bullet, \mathcal{L}^\bullet, \mathcal{M}^\bullet$，
存在同构
$$
\SheafHom^\bullet(\mathcal{K}^\bullet,
\SheafHom^\bullet(\mathcal{L}^\bullet, \mathcal{M}^\bullet))
=
\SheafHom^\bullet(\text{Tot}(\mathcal{K}^\bullet \otimes_\mathcal{O}
\mathcal{L}^\bullet), \mathcal{M}^\bullet)
$$
它是关于 $\mathcal{K}^\bullet, \mathcal{L}^\bullet, \mathcal{M}^\bullet$
函子的 $\mathcal{O}$-模复形同构。
\end{lemma}

\begin{proof}
略去证明。提示：证明方式与代数进阶，引理
\ref{more-algebra-lemma-compose} 完全相同。
\end{proof}

\begin{lemma}
\label{sites-cohomology-lemma-composition}
令 $(\mathcal{C}, \mathcal{O})$ 是环化站点。给定
$\mathcal{O}$-模复形 $\mathcal{K}^\bullet, \mathcal{L}^\bullet, \mathcal{M}^\bullet$，
存在规范态射
$$
\text{Tot}\left(
\SheafHom^\bullet(\mathcal{L}^\bullet, \mathcal{M}^\bullet)
\otimes_\mathcal{O}
\SheafHom^\bullet(\mathcal{K}^\bullet, \mathcal{L}^\bullet)
\right)
\longrightarrow
\SheafHom^\bullet(\mathcal{K}^\bullet, \mathcal{M}^\bullet)
$$
它是 $\mathcal{O}$-模复形的态射。
\end{lemma}

\begin{proof}
略去证明。提示：证明方式与代数进阶，引理
\ref{more-algebra-lemma-composition} 完全相同。
\end{proof}

\begin{lemma}
\label{sites-cohomology-lemma-diagonal-better}
令 $(\mathcal{C}, \mathcal{O})$ 是环化站点。给定
$\mathcal{O}$-模复形 $\mathcal{K}^\bullet, \mathcal{L}^\bullet, \mathcal{M}^\bullet$，
存在规范态射
$$
\text{Tot}\left(
\mathcal{K}^\bullet \otimes_\mathcal{O}
\SheafHom^\bullet(\mathcal{M}^\bullet, \mathcal{L}^\bullet)
\right)
\longrightarrow
\SheafHom^\bullet(\mathcal{M}^\bullet,
\text{Tot}(\mathcal{K}^\bullet \otimes_\mathcal{O} \mathcal{L}^\bullet))
$$
它是关于这三个复形函子的 $\mathcal{O}$-模复形态射。
\end{lemma}

\begin{proof}
略去证明。提示：证明方式与代数进阶，引理
\ref{more-algebra-lemma-diagonal-better} 完全相同。
\end{proof}

\begin{lemma}
\label{sites-cohomology-lemma-diagonal}
令 $(\mathcal{C}, \mathcal{O})$ 是环化站点。给定
$\mathcal{O}$-模复形 $\mathcal{K}^\bullet, \mathcal{L}^\bullet, \mathcal{M}^\bullet$，
存在规范态射
$$
\mathcal{K}^\bullet
\longrightarrow
\SheafHom^\bullet(\mathcal{L}^\bullet,
\text{Tot}(\mathcal{K}^\bullet \otimes_\mathcal{O} \mathcal{L}^\bullet))
$$
它是关于两个复形函子的 $\mathcal{O}$-模复形态射。
\end{lemma}

\begin{proof}
略去证明。提示：证明方式与代数进阶，引理
\ref{more-algebra-lemma-diagonal} 完全相同。
\end{proof}

\begin{lemma}
\label{sites-cohomology-lemma-evaluate-and-more}
令 $(\mathcal{C}, \mathcal{O})$ 是环化站点。给定
$\mathcal{O}$-模复形 $\mathcal{K}^\bullet, \mathcal{L}^\bullet, \mathcal{M}^\bullet$，
存在规范态射
$$
\text{Tot}(\SheafHom^\bullet(\mathcal{L}^\bullet,
\mathcal{M}^\bullet) \otimes_\mathcal{O} \mathcal{K}^\bullet)
\longrightarrow
\SheafHom^\bullet(\SheafHom^\bullet(\mathcal{K}^\bullet,
\mathcal{L}^\bullet), \mathcal{M}^\bullet)
$$
它是关于这三个复形函子的 $\mathcal{O}$-模复形态射。
\end{lemma}

\begin{proof}
略去证明。提示：证明方式与代数进阶，引理
\ref{more-algebra-lemma-evaluate-and-more} 完全相同。
\end{proof}

\begin{lemma}
\label{sites-cohomology-lemma-RHom-into-K-injective}
令 $(\mathcal{C}, \mathcal{O})$ 是环化站点。令 $L,M$
是 $D(\mathcal{O})$ 中的对象。令 $\mathcal{I}^\bullet$
是表示 $M$ 的 K-内射 $\mathcal{O}$-模复形。令
$\mathcal{L}^\bullet$ 是表示 $L$ 的 $\mathcal{O}$-模复形。
则
$$
H^0(\Gamma(U, \SheafHom^\bullet(\mathcal{L}^\bullet, \mathcal{I}^\bullet))) =
\Hom_{D(\mathcal{O}_U)}(L|_U, M|_U)
$$
对所有 $U \in \Ob(\mathcal{C})$ 成立。类似地，
$H^0(\Gamma(\mathcal{C},
\SheafHom^\bullet(\mathcal{L}^\bullet, \mathcal{I}^\bullet))) =
\Hom_{D(\mathcal{O})}(L, M)$.
\end{lemma}

\begin{proof}
我们有
\begin{align*}
H^0(\Gamma(U, \SheafHom^\bullet(\mathcal{L}^\bullet, \mathcal{I}^\bullet)))
& =
\Hom_{K(\mathcal{O}_U)}(\mathcal{L}^\bullet|_U, \mathcal{I}^\bullet|_U) \\
& =
\Hom_{D(\mathcal{O}_U)}(L|_U, M|_U)
\end{align*}
第一个等式由 (\ref{sites-cohomology-equation-cohomology-hom-complex}) 得到。
第二个等式成立，是因为根据引理
\ref{sites-cohomology-lemma-restrict-K-injective-to-open}，$\mathcal{I}^\bullet|_U$
是 K-内射的。最后一个等式的证明类似，只是使用
(\ref{sites-cohomology-equation-global-cohomology-hom-complex})。
\end{proof}

\begin{lemma}
\label{sites-cohomology-lemma-RHom-well-defined}
令 $(\mathcal{C}, \mathcal{O})$ 是环化站点。令
$(\mathcal{I}')^\bullet \to \mathcal{I}^\bullet$
是 K-内射 $\mathcal{O}$-模复形之间的拟同构。
令 $(\mathcal{L}')^\bullet \to \mathcal{L}^\bullet$
是 $\mathcal{O}$-模复形之间的拟同构。则
$$
\SheafHom^\bullet(\mathcal{L}^\bullet, (\mathcal{I}')^\bullet)
\longrightarrow
\SheafHom^\bullet((\mathcal{L}')^\bullet, \mathcal{I}^\bullet)
$$
是拟同构。
\end{lemma}

\begin{proof}
令 $M$ 是由 $\mathcal{I}^\bullet$ 和 $(\mathcal{I}')^\bullet$
表示的 $D(\mathcal{O})$ 对象。令 $L$ 是由
$\mathcal{L}^\bullet$ 和 $(\mathcal{L}')^\bullet$ 表示的
$D(\mathcal{O})$ 对象。根据引理 \ref{sites-cohomology-lemma-RHom-into-K-injective}，
层
$$
H^0(\SheafHom^\bullet(\mathcal{L}^\bullet, (\mathcal{I}')^\bullet))
\quad\text{和}\quad
H^0(\SheafHom^\bullet((\mathcal{L}')^\bullet, \mathcal{I}^\bullet))
$$
都等于与下述预层相伴的层：
$$
U \longmapsto \Hom_{D(\mathcal{O}_U)}(L|_U, M|_U)
$$
因此该映射是拟同构。
\end{proof}

\begin{lemma}
\label{sites-cohomology-lemma-RHom-from-K-flat-into-K-injective}
令 $(\mathcal{C}, \mathcal{O})$ 是环化站点。令 $\mathcal{I}^\bullet$
是 K-内射的 $\mathcal{O}$-模复形，令 $\mathcal{L}^\bullet$ 是
K-平坦的 $\mathcal{O}$-模复形。则
$\SheafHom^\bullet(\mathcal{L}^\bullet, \mathcal{I}^\bullet)$
是 K-内射的 $\mathcal{O}$-模复形。
\end{lemma}

\begin{proof}
具体地，若 $\mathcal{K}^\bullet$ 是无环的 $\mathcal{O}$-模复形，则
\begin{align*}
\Hom_{K(\mathcal{O})}(\mathcal{K}^\bullet,
\SheafHom^\bullet(\mathcal{L}^\bullet, \mathcal{I}^\bullet))
& =
H^0(\Gamma(\mathcal{C},
\SheafHom^\bullet(\mathcal{K}^\bullet,
\SheafHom^\bullet(\mathcal{L}^\bullet, \mathcal{I}^\bullet)))) \\
& =
H^0(\Gamma(\mathcal{C}, \SheafHom^\bullet(\text{Tot}(
\mathcal{K}^\bullet \otimes_\mathcal{O} \mathcal{L}^\bullet),
\mathcal{I}^\bullet))) \\
& =
\Hom_{K(\mathcal{O})}(
\text{Tot}(\mathcal{K}^\bullet \otimes_\mathcal{O} \mathcal{L}^\bullet),
\mathcal{I}^\bullet) \\
& =
0
\end{align*}
第一个等式由 (\ref{sites-cohomology-equation-global-cohomology-hom-complex}) 得到。
第二个等式由引理 \ref{sites-cohomology-lemma-compose} 得到。
第三个等式由 (\ref{sites-cohomology-equation-global-cohomology-hom-complex}) 得到。
最后一个等式是因为
$\text{Tot}(\mathcal{K}^\bullet \otimes_\mathcal{O} \mathcal{L}^\bullet)$
是无环的：$\mathcal{L}^\bullet$ 是 K-平坦的
（定义 \ref{sites-cohomology-definition-K-flat}），而 $\mathcal{I}^\bullet$ 是 K-内射的。
\end{proof}








\section{导出范畴中的内部 Hom}
\label{sites-cohomology-section-internal-hom}

\noindent
令 $(\mathcal{C}, \mathcal{O})$ 是环化站点。令 $L,M$ 是
$D(\mathcal{O})$ 中的对象。我们希望构造 $D(\mathcal{O})$ 中的对象
$R\SheafHom(L,M)$，使得对 $D(\mathcal{O})$ 中任意第三个对象 $K$，
存在规范双射
\begin{equation}
\label{sites-cohomology-equation-internal-hom}
\Hom_{D(\mathcal{O})}(K, R\SheafHom(L, M))
=
\Hom_{D(\mathcal{O})}(K \otimes_\mathcal{O}^\mathbf{L} L, M)
\end{equation}
注意，根据 Yoneda 引理（范畴，引理
\ref{categories-lemma-yoneda}），该公式将 $R\SheafHom(L,M)$
确定到唯一同构。

\medskip\noindent
为构造这样的对象，选取表示 $M$ 的 K-内射
$\mathcal{O}$-模复形 $\mathcal{I}^\bullet$，以及任意表示 $L$ 的
$\mathcal{O}$-模复形 $\mathcal{L}^\bullet$。置
$$
R\SheafHom(L, M) = \SheafHom^\bullet(\mathcal{L}^\bullet, \mathcal{I}^\bullet)
$$
其中右端是第 \ref{sites-cohomology-section-hom-complexes} 节构造的
$\mathcal{O}$-模复形。根据引理 \ref{sites-cohomology-lemma-RHom-well-defined}，
该定义良好。得到函子
$$
D(\mathcal{O})^{opp} \times D(\mathcal{O}) \longrightarrow D(\mathcal{O}),
\quad
(K, L) \longmapsto R\SheafHom(K, L)
$$
作为证明 (\ref{sites-cohomology-equation-internal-hom}) 的准备，我们计算
$R\SheafHom(K,L)$ 的上同调群。

\begin{lemma}
\label{sites-cohomology-lemma-section-RHom-over-U}
令 $(\mathcal{C}, \mathcal{O})$ 是环化站点。令 $L,M$ 是
$D(\mathcal{O})$ 中的对象。对 $\mathcal{C}$ 的每个对象 $U$，有
$$
H^0(U, R\SheafHom(L, M)) =
\Hom_{D(\mathcal{O}_U)}(L|_U, M|_U)
$$
并且有 $H^0(\mathcal{C}, R\SheafHom(L,M)) =
\Hom_{D(\mathcal{O})}(L, M)$.
\end{lemma}

\begin{proof}
选取表示 $M$ 的 K-内射 $\mathcal{O}$-模复形
$\mathcal{I}^\bullet$，以及表示 $L$ 的 K-平坦复形
$\mathcal{L}^\bullet$。则
$\SheafHom^\bullet(\mathcal{L}^\bullet, \mathcal{I}^\bullet)$
根据引理 \ref{sites-cohomology-lemma-RHom-from-K-flat-into-K-injective}，
它是 K-内射的。因此只需取 $U$ 上的截面即可计算 $U$ 上的上同调，
结论由引理 \ref{sites-cohomology-lemma-RHom-into-K-injective} 得到。
\end{proof}

\begin{lemma}
\label{sites-cohomology-lemma-internal-hom}
令 $(\mathcal{C}, \mathcal{O})$ 是环化站点。令 $K,L,M$ 是
$D(\mathcal{O})$ 中的对象。利用上述构造，存在规范同构
$$
R\SheafHom(K, R\SheafHom(L, M)) =
R\SheafHom(K \otimes_\mathcal{O}^\mathbf{L} L, M)
$$
它位于 $D(\mathcal{O})$ 中，并且关于 $K,L,M$ 函子性；取
$H^0(\mathcal{C},-)$ 后恢复 (\ref{sites-cohomology-equation-internal-hom})。
\end{lemma}

\begin{proof}
选取表示 $M$ 的 K-内射复形 $\mathcal{I}^\bullet$，以及表示 $L$
的 K-平坦 $\mathcal{O}$-模复形 $\mathcal{L}^\bullet$。
对任意 $\mathcal{O}$-模复形 $\mathcal{K}^\bullet$，有
$$
\SheafHom^\bullet(\mathcal{K}^\bullet,
\SheafHom^\bullet(\mathcal{L}^\bullet, \mathcal{I}^\bullet))
=
\SheafHom^\bullet(
\text{Tot}(\mathcal{K}^\bullet \otimes_\mathcal{O} \mathcal{L}^\bullet),
\mathcal{I}^\bullet)
$$
由引理 \ref{sites-cohomology-lemma-compose} 得到。注意左端表示
$R\SheafHom(K,R\SheafHom(L,M))$（使用引理
\ref{sites-cohomology-lemma-RHom-from-K-flat-into-K-injective}），右端表示
$R\SheafHom(K \otimes_\mathcal{O}^\mathbf{L} L,M)$。
这就证明了引理中的显示公式。取全局截面并使用引理
\ref{sites-cohomology-lemma-section-RHom-over-U}，得到
(\ref{sites-cohomology-equation-internal-hom})。
\end{proof}

\begin{lemma}
\label{sites-cohomology-lemma-restriction-RHom-to-U}
令 $(\mathcal{C}, \mathcal{O})$ 是环化站点。令 $K,L$ 是
$D(\mathcal{O})$ 中的对象。$R\SheafHom(K,L)$ 的构造
与限制相容，即对 $\mathcal{C}$ 的每个对象 $U$，有
$R\SheafHom(K|_U, L|_U) = R\SheafHom(K, L)|_U$.
\end{lemma}

\begin{proof}
这由构造和引理 \ref{sites-cohomology-lemma-restrict-K-injective-to-open} 立即得到。
\end{proof}

\begin{lemma}
\label{sites-cohomology-lemma-RHom-triangulated}
令 $(\mathcal{C}, \mathcal{O})$ 是环化站点。双函子
$R\SheafHom(-,-)$ 将两个变量中的有区别三角变为有区别三角。
\end{lemma}

\begin{proof}
这由下述对应的性质得到：
$$
(\mathcal{L}^\bullet, \mathcal{M}^\bullet) \longmapsto
\SheafHom^\bullet(\mathcal{L}^\bullet, \mathcal{M}^\bullet)
$$
它将任一变量中逐项分裂的复形短正合序列变为逐项分裂的短正合序列。
细节略去。
\end{proof}

\begin{lemma}
\label{sites-cohomology-lemma-internal-hom-evaluate}
令 $(\mathcal{C}, \mathcal{O})$ 是环化站点。令 $K,L,M$ 是
$D(\mathcal{O})$ 中的对象。存在规范态射
$$
R\SheafHom(L, M) \otimes_\mathcal{O}^\mathbf{L} K
\longrightarrow
R\SheafHom(R\SheafHom(K, L), M)
$$
它位于 $D(\mathcal{O})$ 中，并且关于 $K,L,M$ 函子性。
\end{lemma}

\begin{proof}
选取表示 $M$ 的 K-内射复形 $\mathcal{I}^\bullet$，
表示 $L$ 的 K-内射复形 $\mathcal{J}^\bullet$，以及表示 $K$ 的
K-平坦复形 $\mathcal{K}^\bullet$。该映射利用引理
$$
\text{Tot}(\SheafHom^\bullet(\mathcal{J}^\bullet,
\mathcal{I}^\bullet) \otimes_\mathcal{O} \mathcal{K}^\bullet)
\longrightarrow
\SheafHom^\bullet(\SheafHom^\bullet(\mathcal{K}^\bullet,
\mathcal{J}^\bullet), \mathcal{I}^\bullet)
$$
\ref{sites-cohomology-lemma-evaluate-and-more} 中的映射定义。由所选复形，
左端表示 $R\SheafHom(L,M) \otimes_\mathcal{O}^\mathbf{L}K$，
右端表示 $R\SheafHom(R\SheafHom(K,L),M)$。
关于 $D(\mathcal{O})$ 中三个对象的函子性证明略去。
\end{proof}

\begin{lemma}
\label{sites-cohomology-lemma-internal-hom-composition}
\begin{slogan}
在 $R\SheafHom$ 上复合。
\end{slogan}
令 $(\mathcal{C}, \mathcal{O})$ 是环化站点。给定
$D(\mathcal{O})$ 中的 $K,L,M$，存在规范态射
$$
R\SheafHom(L, M) \otimes_\mathcal{O}^\mathbf{L} R\SheafHom(K, L)
\longrightarrow R\SheafHom(K, M)
$$
它位于 $D(\mathcal{O})$ 中。
\end{lemma}

\begin{proof}
选取表示 $M$ 的 K-内射复形 $\mathcal{I}^\bullet$、表示 $L$ 的
K-内射复形 $\mathcal{J}^\bullet$，以及任意表示 $K$ 的
$\mathcal{O}$-模复形 $\mathcal{K}^\bullet$。根据引理
\ref{sites-cohomology-lemma-composition}，存在复形映射
$$
\text{Tot}\left(
\SheafHom^\bullet(\mathcal{J}^\bullet, \mathcal{I}^\bullet)
\otimes_\mathcal{O}
\SheafHom^\bullet(\mathcal{K}^\bullet, \mathcal{J}^\bullet)
\right)
\longrightarrow
\SheafHom^\bullet(\mathcal{K}^\bullet, \mathcal{I}^\bullet)
$$
下列 $\mathcal{O}$-模复形
$\SheafHom^\bullet(\mathcal{J}^\bullet, \mathcal{I}^\bullet)$,
$\SheafHom^\bullet(\mathcal{K}^\bullet, \mathcal{J}^\bullet)$，以及
$\SheafHom^\bullet(\mathcal{K}^\bullet, \mathcal{I}^\bullet)$
分别表示 $R\SheafHom(L,M)$、$R\SheafHom(K,L)$ 和
$R\SheafHom(K,M)$。若选取 K-平坦复形 $\mathcal{H}^\bullet$ 以及拟同构
$\mathcal{H}^\bullet \to
\SheafHom^\bullet(\mathcal{K}^\bullet, \mathcal{J}^\bullet)$,
则存在映射
$$
\text{Tot}\left(
\SheafHom^\bullet(\mathcal{J}^\bullet, \mathcal{I}^\bullet)
\otimes_\mathcal{O} \mathcal{H}^\bullet
\right)
\longrightarrow
\text{Tot}\left(
\SheafHom^\bullet(\mathcal{J}^\bullet, \mathcal{I}^\bullet)
\otimes_\mathcal{O}
\SheafHom^\bullet(\mathcal{K}^\bullet, \mathcal{J}^\bullet)
\right)
$$
其源表示
$R\SheafHom(L,M) \otimes_\mathcal{O}^\mathbf{L} R\SheafHom(K,L)$。
复合显示的两个箭头即得到所需映射。构造的函子性证明略去。
\end{proof}

\begin{lemma}
\label{sites-cohomology-lemma-internal-hom-diagonal-better}
令 $(\mathcal{C}, \mathcal{O})$ 是环化站点。给定
$D(\mathcal{O})$ 中的 $K,L,M$，存在规范态射
$$
K \otimes_\mathcal{O}^\mathbf{L} R\SheafHom(M, L)
\longrightarrow
R\SheafHom(M, K \otimes_\mathcal{O}^\mathbf{L} L)
$$
它位于 $D(\mathcal{O})$ 中，并且关于 $K,L,M$ 函子性。
\end{lemma}

\begin{proof}
选取表示 $K$ 的 K-平坦复形 $\mathcal{K}^\bullet$，
表示 $L$ 的 K-内射复形 $\mathcal{I}^\bullet$，以及任意表示 $M$ 的
$\mathcal{O}$-模复形 $\mathcal{M}^\bullet$。选取拟同构
$\text{Tot}(\mathcal{K}^\bullet \otimes_{\mathcal{O}_X} \mathcal{I}^\bullet)
\to \mathcal{J}^\bullet$
其中 $\mathcal{J}^\bullet$ 是 K-内射的。然后使用映射
$$
\text{Tot}\left(
\mathcal{K}^\bullet \otimes_\mathcal{O}
\SheafHom^\bullet(\mathcal{M}^\bullet, \mathcal{I}^\bullet)
\right)
\to
\SheafHom^\bullet(\mathcal{M}^\bullet,
\text{Tot}(\mathcal{K}^\bullet \otimes_\mathcal{O} \mathcal{I}^\bullet))
\to
\SheafHom^\bullet(\mathcal{M}^\bullet, \mathcal{J}^\bullet)
$$
其中第一个映射来自引理 \ref{sites-cohomology-lemma-diagonal-better}。
\end{proof}

\begin{lemma}
\label{sites-cohomology-lemma-internal-hom-diagonal}
令 $(\mathcal{C}, \mathcal{O})$ 是环化站点。
给定 $D(\mathcal{O})$ 中的 $K,L$，存在规范态射
$$
K \longrightarrow R\SheafHom(L, K \otimes_\mathcal{O}^\mathbf{L} L)
$$
它位于 $D(\mathcal{O})$ 中，并且关于 $K,L$ 函子性。
\end{lemma}

\begin{proof}
选取表示 $K$ 的 K-平坦复形 $\mathcal{K}^\bullet$，以及表示 $L$ 的
任意 $\mathcal{O}$-模复形 $\mathcal{L}^\bullet$。选取 K-内射复形
$\mathcal{J}^\bullet$ 和拟同构
$\text{Tot}(\mathcal{K}^\bullet \otimes_\mathcal{O} \mathcal{L}^\bullet)
\to \mathcal{J}^\bullet$。然后使用
$$
\mathcal{K}^\bullet \to
\SheafHom^\bullet(\mathcal{L}^\bullet,
\text{Tot}(\mathcal{K}^\bullet \otimes_\mathcal{O} \mathcal{L}^\bullet))
\to
\SheafHom^\bullet(\mathcal{L}^\bullet, \mathcal{J}^\bullet)
$$
其中第一个映射来自引理 \ref{sites-cohomology-lemma-diagonal}。
\end{proof}

\begin{lemma}
\label{sites-cohomology-lemma-dual}
令 $(\mathcal{C}, \mathcal{O})$ 是环化站点。令 $L$ 是
$D(\mathcal{O})$ 中的对象。置 $L^\vee = R\SheafHom(L,\mathcal{O})$。
对于 $D(\mathcal{O})$ 中的 $M$，存在规范映射
\begin{equation}
\label{sites-cohomology-equation-eval}
M \otimes^\mathbf{L}_\mathcal{O} L^\vee \longrightarrow R\SheafHom(L, M)
\end{equation}
它诱导规范映射
$$
H^0(\mathcal{C}, M \otimes_\mathcal{O}^\mathbf{L} L^\vee)
\longrightarrow
\Hom_{D(\mathcal{O})}(L, M)
$$
且关于 $D(\mathcal{O})$ 中的 $M$ 函子性。
\end{lemma}

\begin{proof}
映射 (\ref{sites-cohomology-equation-eval}) 是引理
\ref{sites-cohomology-lemma-internal-hom-composition} 的特例，其中使用了识别
$M = R\SheafHom(\mathcal{O},M)$。
\end{proof}

\begin{remark}
\label{sites-cohomology-remark-projection-formula-for-internal-hom}
令 $f : (\Sh(\mathcal{C}), \mathcal{O}_\mathcal{C}) \to
(\Sh(\mathcal{D}), \mathcal{O}_\mathcal{D})$ 是环化拓扑斯态射。
令 $K,L$ 是 $D(\mathcal{O}_\mathcal{C})$ 中的对象。断言存在规范映射
$$
Rf_*R\SheafHom(L, K) \longrightarrow R\SheafHom(Rf_*L, Rf_*K)
$$
具体地，根据 (\ref{sites-cohomology-equation-internal-hom})，这等价于映射
$Rf_*R\SheafHom(L, K) \otimes_{\mathcal{O}_\mathcal{D}}^\mathbf{L} Rf_*L
\to Rf_*K$.
为此可使用复合
$$
Rf_*R\SheafHom(L, K) \otimes_{\mathcal{O}_\mathcal{D}}^\mathbf{L} Rf_*L \to
Rf_*(R\SheafHom(L, K) \otimes_{\mathcal{O}_\mathcal{C}}^\mathbf{L} L) \to
Rf_*K
$$
其中第一个箭头是相对杯积（评注 \ref{sites-cohomology-remark-cup-product}），
第二个箭头是将 $Rf_*$ 应用于来自引理
\ref{sites-cohomology-lemma-internal-hom-composition} 的规范映射
$R\SheafHom(L,K) \otimes_{\mathcal{O}_\mathcal{C}}^\mathbf{L} L \to K$
（其中一个位置取 $\mathcal{O}_\mathcal{C}$）。
\end{remark}

\begin{remark}
\label{sites-cohomology-remark-prepare-fancy-base-change}
令 $h : (\Sh(\mathcal{C}), \mathcal{O}) \to (\Sh(\mathcal{C}'), \mathcal{O}')$
是环化拓扑斯态射。令 $K,L$ 是 $D(\mathcal{O}')$ 中的对象。
断言存在规范映射
$$
Lh^*R\SheafHom(K, L) \longrightarrow R\SheafHom(Lh^*K, Lh^*L)
$$
它位于 $D(\mathcal{O})$ 中。具体地，根据
(\ref{sites-cohomology-equation-internal-hom})（该式由引理 \ref{sites-cohomology-lemma-internal-hom} 证明），
这样的映射等价于映射
$$
Lh^*R\SheafHom(K, L) \otimes^\mathbf{L} Lh^*K \longrightarrow Lh^*L
$$
根据引理 \ref{sites-cohomology-lemma-pullback-tensor-product}，该箭头的源是
$Lh^*(\SheafHom(K,L) \otimes^\mathbf{L} K)$，
因此只需构造规范映射
$$
R\SheafHom(K, L) \otimes^\mathbf{L} K \longrightarrow L.
$$
为此取与下式对应的箭头：
$$
\text{id} :
R\SheafHom(K, L)
\longrightarrow
R\SheafHom(K, L)
$$
通过 (\ref{sites-cohomology-equation-internal-hom})。
\end{remark}

\begin{remark}
\label{sites-cohomology-remark-fancy-base-change}
设
$$
\xymatrix{
(\Sh(\mathcal{C}'), \mathcal{O}_{\mathcal{C}'})
\ar[r]_h \ar[d]_{f'} &
(\Sh(\mathcal{C}), \mathcal{O}_\mathcal{C}) \ar[d]^f \\
(\Sh(\mathcal{D}'), \mathcal{O}_{\mathcal{D}'})
\ar[r]^g &
(\Sh(\mathcal{D}), \mathcal{O}_\mathcal{D})
}
$$
是环化拓扑斯的交换图。令 $K,L$ 是
$D(\mathcal{O}_\mathcal{C})$ 中的对象。断言存在规范基变换映射
$$
Lg^*Rf_*R\SheafHom(K, L)
\longrightarrow
R(f')_*R\SheafHom(Lh^*K, Lh^*L)
$$
它位于 $D(\mathcal{O}_{\mathcal{D}'})$ 中。具体地，取下述复合的伴随映射：
\begin{align*}
L(f')^*Lg^*Rf_*R\SheafHom(K, L)
& =
Lh^*Lf^*Rf_*R\SheafHom(K, L) \\
& \to
Lh^*R\SheafHom(K, L) \\
& \to
R\SheafHom(Lh^*K, Lh^*L)
\end{align*}
其中第一个箭头使用伴随映射 $Lf^*Rf_* \to \text{id}$，
第二个箭头是评注 \ref{sites-cohomology-remark-prepare-fancy-base-change} 中构造的规范映射。
\end{remark}




\section{全局导出 Hom}
\label{sites-cohomology-section-global-RHom}

\noindent
令 $(\Sh(\mathcal{C}), \mathcal{O})$ 是环化拓扑斯。
令 $K,L \in D(\mathcal{O})$。利用导出范畴中内部 Hom 的构造，
得到定义良好的对象
$$
R\Hom_\mathcal{O}(K, L) = R\Gamma(\mathcal{C}, R\SheafHom(K, L))
$$
它位于 $D(\Gamma(\mathcal{C}, \mathcal{O}))$ 中。根据引理
\ref{sites-cohomology-lemma-section-RHom-over-U}，有
$$
H^0(R\Hom_\mathcal{O}(K, L)) = \Hom_{D(\mathcal{O})}(K, L)
$$
并且
$$
H^p(R\Hom_\mathcal{O}(K, L)) = \Ext_{D(\mathcal{O})}^p(K, L)
$$
若 $f : (\mathcal{C}', \mathcal{O}') \to (\mathcal{C}, \mathcal{O})$
是环化拓扑斯态射，则存在规范映射
$$
R\Hom_\mathcal{O}(K, L) \longrightarrow R\Hom_{\mathcal{O}'}(Lf^*K, Lf^*L)
$$
它位于 $D(\Gamma(\mathcal{O}))$ 中，是取评注
\ref{sites-cohomology-remark-prepare-fancy-base-change} 所定义映射的全局截面得到的。









\section{导出下 shriek}
\label{sites-cohomology-section-derived-lower-shriek}

\noindent
本节研究环化拓扑斯的态射 $g$：除 $Lg^*$ 和 $Rg_*$ 外，
还存在导出函子 $Lg_!$。

\begin{lemma}
\label{sites-cohomology-lemma-pullback-injective-pre-limp}
令 $u : \mathcal{C} \to \mathcal{D}$ 是站点的连续且余连续函子。
令 $g : \Sh(\mathcal{C}) \to \Sh(\mathcal{D})$ 是相应的拓扑斯态射。
令 $\mathcal{O}_\mathcal{D}$ 是环层，令 $\mathcal{I}$ 是内射的
$\mathcal{O}_\mathcal{D}$-模。则
$H^p(U, g^{-1}\mathcal{I}) = 0$ 对所有 $p > 0$ 以及
$U \in \Ob(\mathcal{C})$ 成立。
\end{lemma}

\begin{proof}
若能证明对 $\mathcal{C}$ 的任意覆盖
$\mathcal{U} = \{U_i \to U\}$，所有高阶 Čech 上同调群
$\check H^p(\mathcal{U},g^{-1}\mathcal{I})$ 都消失，
则由引理 \ref{sites-cohomology-lemma-cech-vanish-collection} 得到本引理的消失结论。
由于 $u$ 连续，$u(\mathcal{U}) = \{u(U_i) \to u(U)\}$
是 $\mathcal{D}$ 的覆盖，并且
$u(U_{i_0} \times_U \ldots \times_U U_{i_n}) =
u(U_{i_0}) \times_{u(U)} \ldots \times_{u(U)} u(U_{i_n})$.
因此有
$$
\check H^p(\mathcal{U}, g^{-1}\mathcal{I}) =
\check H^p(u(\mathcal{U}), \mathcal{I})
$$
因为根据站点，引理 \ref{sites-lemma-when-shriek}，$g^{-1} = u^p$。
由于 $\mathcal{I}$ 是内射的 $\mathcal{O}_\mathcal{D}$-模，这些 Čech
上同调群消失，参见引理 \ref{sites-cohomology-lemma-injective-module-trivial-cech}。
\end{proof}

\begin{lemma}
\label{sites-cohomology-lemma-existence-derived-lower-shriek}
令 $u : \mathcal{C} \to \mathcal{D}$ 是站点的连续且余连续函子。
令 $g : \Sh(\mathcal{C}) \to \Sh(\mathcal{D})$ 是相应的拓扑斯态射。
令 $\mathcal{O}_\mathcal{D}$ 是环层，并置
$\mathcal{O}_\mathcal{C} = g^{-1}\mathcal{O}_\mathcal{D}$。
函子 $g_! : \textit{Mod}(\mathcal{O}_\mathcal{C}) \to
\textit{Mod}(\mathcal{O}_\mathcal{D})$
（参见站点上的模，引理
\ref{sites-modules-lemma-lower-shriek-modules}）具有左导出函子
$$
Lg_! : D(\mathcal{O}_\mathcal{C}) \longrightarrow D(\mathcal{O}_\mathcal{D})
$$
它是 $g^*$ 的左伴随。此外，对 $U \in \Ob(\mathcal{C})$，有
$$
Lg_!(j_{U!}\mathcal{O}_U) =
g_!j_{U!}\mathcal{O}_U =
j_{u(U)!} \mathcal{O}_{u(U)}.
$$
其中 $j_{U!}$ 和 $j_{u(U)!}$ 是与局部化态射相关的零延拓：
$j_U : \mathcal{C}/U \to \mathcal{C}$ 以及
$j_{u(U)} : \mathcal{D}/u(U) \to \mathcal{D}$。
\end{lemma}

\begin{proof}
我们将使用导出范畴，命题
\ref{derived-proposition-left-derived-exists} 构造 $Lg_!$。
为此需验证该命题的假设 (1)、(2)、(3)、(4)、(5)。首先，由于
$g_!$ 是左伴随，它右正合且与所有余极限交换，故 (5) 成立。
(3)、(4) 成立，因为环化站点上的模范畴是 Grothendieck 阿贝尔范畴。
令 $\mathcal{P} \subset \Ob(\textit{Mod}(\mathcal{O}_\mathcal{C}))$
为由形如 $j_{U!}\mathcal{O}_U$ 的模作直和得到的
$\mathcal{O}_\mathcal{C}$-模集合。注意
$g_!j_{U!}\mathcal{O}_U = j_{u(U)!}\mathcal{O}_{u(U)}$，参见站点上的模，
引理 \ref{sites-modules-lemma-lower-shriek-modules} 的证明。
每个 $\mathcal{O}_\mathcal{C}$-模都是 $\mathcal{P}$ 中某对象的商，参见
站点上的模，引理 \ref{sites-modules-lemma-module-quotient-flat}。
因此 (1) 成立。最后需证明 (2)。
令 $\mathcal{K}^\bullet$ 是有上界的无环
$\mathcal{O}_\mathcal{C}$-模复形，且所有 $n$ 都有 $\mathcal{K}^n \in \mathcal{P}$。
需证明 $g_!\mathcal{K}^\bullet$ 正合。为此，只需证明对每个内射的
$\mathcal{O}_\mathcal{D}$-模 $\mathcal{I}$，
$$
\Hom_{D(\mathcal{O}_\mathcal{D})}(
g_!\mathcal{K}^\bullet, \mathcal{I}[n]) = 0
$$
对所有 $n \in \mathbf{Z}$ 均成立。由于 $\mathcal{I}$ 内射，有
\begin{align*}
\Hom_{D(\mathcal{O}_\mathcal{D})}(
g_!\mathcal{K}^\bullet, \mathcal{I}[n])
& =
\Hom_{K(\mathcal{O}_\mathcal{D})}(
g_!\mathcal{K}^\bullet, \mathcal{I}[n]) \\
& =
H^n(\Hom_{\mathcal{O}_\mathcal{D}}(
g_!\mathcal{K}^\bullet, \mathcal{I})) \\
& =
H^n(\Hom_{\mathcal{O}_\mathcal{C}}(
\mathcal{K}^\bullet, g^{-1}\mathcal{I}))
\end{align*}
最后一个等式来自 $g_!$ 与 $g^{-1}$ 的伴随性。

\medskip\noindent
若 $g^{-1}\mathcal{I}$ 是内射的 $\mathcal{O}_\mathcal{C}$-模，
该群的消失显然成立。但 $g_!$ 一般不正合，所以 $g^{-1}\mathcal{I}$
不一定是内射的 $\mathcal{O}_\mathcal{C}$-模。不过我们知道
$$
\Ext^p_{\mathcal{O}_\mathcal{C}}(
j_{U!}\mathcal{O}_U, g^{-1}\mathcal{I}) =
H^p(U, g^{-1}\mathcal{I}) = 0 \text{，其中 }p \geq 1
$$
这里第一个等式来自
$\Hom_{\mathcal{O}_\mathcal{C}}(j_{U!}\mathcal{O}_U, \mathcal{H}) =
\mathcal{H}(U)$ 并取导出函子；而根据引理
\ref{sites-cohomology-lemma-pullback-injective-pre-limp}，对 $p > 0$ 和
$U \in \Ob(\mathcal{C})$，$H^p(U,g^{-1}\mathcal{I})$ 消失。
由于每个 $\mathcal{K}^{-q}$ 都是形如 $j_{U!}\mathcal{O}_U$ 的模的直和，
可知
$$
\Ext^p_{\mathcal{O}_\mathcal{C}}(\mathcal{K}^{-q}, g^{-1}\mathcal{I}) = 0
\text{ 对所有 }p \geq 1\text{ 和 }q
$$
使用谱序列（参见例 \ref{sites-cohomology-example-hom-complex-into-sheaf}）
$$
E_1^{p, q} = \Ext^q_{\mathcal{O}_\mathcal{C}}(
\mathcal{K}^{-p}, g^{-1}\mathcal{I})
\Rightarrow
\Ext^{p + q}_{\mathcal{O}_\mathcal{C}}(
\mathcal{K}^\bullet, g^{-1}\mathcal{I}) = 0.
$$
注意该谱序列收敛到零，因为 $\mathcal{K}^\bullet$ 无环
（因此在导出范畴中为零，从而 Ext 群消失）。由上面证明的高阶 Ext
消失，$E_1$ 页上唯一可能非零的项是
$E_1^{p, 0} = \Hom_{\mathcal{O}_\mathcal{C}}(
\mathcal{K}^{-p}, g^{-1}\mathcal{I})$.
因此复形
$\Hom_{\mathcal{O}_\mathcal{C}}(
\mathcal{K}^\bullet, g^{-1}\mathcal{I})$
是无环的，正如所需。

\medskip\noindent
因此左导出函子 $Lg_!$ 存在。它是
$g^{-1} = g^* = Rg^* = Lg^*$ 的左伴随，即
\begin{equation}
\label{sites-cohomology-equation-to-prove}
\Hom_{D(\mathcal{O}_\mathcal{C})}(K, g^*L) =
\Hom_{D(\mathcal{O}_\mathcal{D})}(Lg_!K, L)
\end{equation}
由导出范畴，引理 \ref{derived-lemma-derived-adjoint-functors} 得到。
证明完成。
\end{proof}

\begin{remark}
\label{sites-cohomology-remark-when-derived-shriek-equal}
警告！令 $u : \mathcal{C} \to \mathcal{D}$、$g$、$\mathcal{O}_\mathcal{D}$
和 $\mathcal{O}_\mathcal{C}$ 如引理
\ref{sites-cohomology-lemma-existence-derived-lower-shriek} 中所述。
一般而言，下图并不交换：
$$
\xymatrix{
D(\mathcal{O}_\mathcal{C}) \ar[r]_{Lg_!} \ar[d]_{forget} &
D(\mathcal{O}_\mathcal{D}) \ar[d]^{forget} \\
D(\mathcal{C}) \ar[r]^{Lg^{Ab}_!} &
D(\mathcal{D})
}
$$
其中函子 $Lg_!^{Ab}$ 是引理
\ref{sites-cohomology-lemma-existence-derived-lower-shriek} 中构造的函子，
但在 $\mathcal{C}$ 和 $\mathcal{D}$ 上都使用常值层 $\mathbf{Z}$ 作为结构层。
一般甚至不成立 $g_! = g_!^{Ab}$（参见站点上的模，评注
\ref{sites-modules-remark-when-shriek-equal}）；但即使
$g_! = g_!^{Ab}$，上述现象也可能发生！具体地，
引理 \ref{sites-cohomology-lemma-existence-derived-lower-shriek} 证明中 $Lg_!$ 的构造表明，
当且仅当下述规范映射
$$
Lg^{Ab}_!j_{U!}\mathcal{O}_U \longrightarrow j_{u(U)!}\mathcal{O}_{u(U)}
$$
对 $\mathcal{C}$ 中所有对象 $U$ 在 $D(\mathcal{D})$ 中是同构时，
$Lg_!$ 才与 $Lg_!^{\textit{Ab}}$ 一致。一般我们只能说存在自然变换
$$
Lg_!^{Ab} \circ forget \longrightarrow forget \circ Lg_!
$$
\end{remark}

\begin{lemma}
\label{sites-cohomology-lemma-pullback-injective-limp}
令 $u : \mathcal{C} \to \mathcal{D}$ 是站点的连续且余连续函子。
令 $g : \Sh(\mathcal{C}) \to \Sh(\mathcal{D})$ 是相应的拓扑斯态射。
令 $\mathcal{O}_\mathcal{D}$ 是环层，令 $\mathcal{I}$ 是内射的
$\mathcal{O}_\mathcal{D}$-模。若
$g_!^{Sh} : \Sh(\mathcal{C}) \to \Sh(\mathcal{D})$
与纤维积交换\footnote{若 $\mathcal{C}$ 有有限连通极限且 $u$ 与之交换则成立，
参见站点，引理 \ref{sites-lemma-preserve-equalizers}。}，则
$g^{-1}\mathcal{I}$ 完全无环。
\end{lemma}

\begin{proof}
我们将使用引理 \ref{sites-cohomology-lemma-characterize-limp} 的判别准则。
条件 (1) 由引理 \ref{sites-cohomology-lemma-pullback-injective-pre-limp} 成立。
令 $K' \to K$ 是 $\mathcal{C}$ 上集合层的满射。
由于 $g_!^{Sh}$ 是左伴随，$g_!^{Sh}K' \to g_!^{Sh}K$ 满射。
注意
\begin{align*}
H^0(K' \times_K \ldots \times_K K', g^{-1}\mathcal{I}) 
& =
H^0(g_!^{Sh}(K' \times_K \ldots \times_K K'), \mathcal{I}) \\
& =
H^0(g_!^{Sh}K' \times_{g_!^{Sh}K} \ldots \times_{g_!^{Sh}K} g_!^{Sh}K',
\mathcal{I})
\end{align*}
由我们对 $g_!^{Sh}$ 的假设成立。由于 $\mathcal{I}$ 是内射模，
根据引理 \ref{sites-cohomology-lemma-direct-image-injective-sheaf}（应用于恒等态射），
它完全无环。因此可使用引理 \ref{sites-cohomology-lemma-characterize-limp} 的逆命题，
得到复形
$$
0 \to H^0(K, g^{-1}\mathcal{I}) \to H^0(K', g^{-1}\mathcal{I}) \to
H^0(K' \times_K K', g^{-1}\mathcal{I}) \to \ldots
$$
正合，证毕。
\end{proof}

\begin{lemma}
\label{sites-cohomology-lemma-pullback-same-cohomology}
令 $u : \mathcal{C} \to \mathcal{D}$ 是站点的连续且余连续函子。
令 $g : \Sh(\mathcal{C}) \to \Sh(\mathcal{D})$ 是相应的拓扑斯态射。
令 $U \in \Ob(\mathcal{C})$。
\begin{enumerate}
\item 对于 $D(\mathcal{D})$ 中的 $M$，有
$R\Gamma(U, g^{-1}M) = R\Gamma(u(U), M)$.
\item 若 $\mathcal{O}_\mathcal{D}$ 是环层且
$\mathcal{O}_\mathcal{C} = g^{-1}\mathcal{O}_\mathcal{D}$，则对于
$D(\mathcal{O}_\mathcal{D})$ 中的 $M$，有
$R\Gamma(U, g^*M) = R\Gamma(u(U), M)$.
\end{enumerate}
\end{lemma}

\begin{proof}
在有下界情形中，用有下界的内射复形表示 $K$，并使用引理
\ref{sites-cohomology-lemma-pullback-injective-pre-limp} 以及 Leray 无环引理，
即可得到 (1)、(2)。在无界情形中，先注意 (1) 是 (2) 的特例。
对于 (2)，可使用
$$
R\Gamma(U, g^*M) =
R\Hom_{\mathcal{O}_\mathcal{C}}(j_{U!}\mathcal{O}_U, g^*M) =
R\Hom_{\mathcal{O}_\mathcal{D}}(j_{u(U)!}\mathcal{O}_{u(U)}, M) =
R\Gamma(u(U), M)
$$
其中中间等式是引理
\ref{sites-cohomology-lemma-existence-derived-lower-shriek} 的结果。
\end{proof}

\begin{lemma}
\label{sites-cohomology-lemma-special-square-cocontinuous}
假设给定交换图
$$
\xymatrix{
(\Sh(\mathcal{C}'), \mathcal{O}_{\mathcal{C}'})
\ar[r]_{(g', (g')^\sharp)} \ar[d]_{(f', (f')^\sharp)} &
(\Sh(\mathcal{C}), \mathcal{O}_\mathcal{C}) \ar[d]^{(f, f^\sharp)} \\
(\Sh(\mathcal{D}'), \mathcal{O}_{\mathcal{D}'}) \ar[r]^{(g, g^\sharp)} &
(\Sh(\mathcal{D}), \mathcal{O}_\mathcal{D})
}
$$
它由环化拓扑斯组成。假设
\begin{enumerate}
\item $f$、$f'$、$g$ 和 $g'$ 分别对应于余连续函子
$u$、$u'$、$v$ 和 $v'$，如《站点》引理
\ref{sites-lemma-cocontinuous-morphism-topoi} 所述；
\item $v \circ u' = u \circ v'$,
\item $v$ 和 $v'$ 既连续又余连续；
\item 对 $\mathcal{D}'$ 的任意对象 $V'$，由 $v$ 给出的函子
${}^{u'}_{V'}\mathcal{I} \to {}^{\ \ \ u}_{v(V')}\mathcal{I}$
是共终的；
\item $g^{-1}\mathcal{O}_{\mathcal{D}} = \mathcal{O}_{\mathcal{D}'}$
且 $(g')^{-1}\mathcal{O}_{\mathcal{C}} = \mathcal{O}_{\mathcal{C}'}$；
\item $g'_! : \textit{Ab}(\mathcal{C}') \to \textit{Ab}(\mathcal{C})$
是正合的\footnote{若 $\mathcal{C}'$ 中存在纤维积和等化子，且 $v'$ 与它们交换，
则此条件成立，见《站点上的模》引理
\ref{sites-modules-lemma-exactness-lower-shriek}。}。
\end{enumerate}
则作为函子
$D(\mathcal{O}_\mathcal{C}) \to D(\mathcal{O}_{\mathcal{D}'})$.
有 $Rf'_* \circ (g')^* = g^* \circ Rf_*$。
\end{lemma}

\begin{proof}
由条件 (5)，有
$g^* = Lg^* = g^{-1}$ 且 $(g')^* = L(g')^* = (g')^{-1}$。
根据引理 \ref{sites-cohomology-lemma-modules-abelian-unbounded}，只需在阿贝尔层的
导出范畴 $D(\mathcal{C})$ 上证明结论。选取对象
$K \in D(\mathcal{C})$。令 $\mathcal{I}^\bullet$ 是 $\mathcal{C}$ 上
表示 $K$ 的阿贝尔层 K-内射复形。由《导出范畴》引理
\ref{derived-lemma-adjoint-preserve-K-injectives} 和假设 (6)，
$(g')^{-1}\mathcal{I}^\bullet$ 是 $\mathcal{C}'$ 上阿贝尔层的 K-内射复形。
由《站点上的模》引理
\ref{sites-modules-lemma-special-square-cocontinuous}，有
$f'_*(g')^{-1}\mathcal{I}^\bullet = g^{-1}f_*\mathcal{I}^\bullet$。
由于 $f_*\mathcal{I}^\bullet$ 表示 $Rf_*K$，而
$f'_*(g')^{-1}\mathcal{I}^\bullet$ 表示 $Rf'_*(g')^{-1}K$，结论成立。
\end{proof}

\begin{lemma}
\label{sites-cohomology-lemma-special-square-continuous}
考虑交换图
$$
\xymatrix{
(\Sh(\mathcal{C}'), \mathcal{O}_{\mathcal{C}'})
\ar[r]_{(g', (g')^\sharp)} \ar[d]_{(f', (f')^\sharp)} &
(\Sh(\mathcal{C}), \mathcal{O}_\mathcal{C}) \ar[d]^{(f, f^\sharp)} \\
(\Sh(\mathcal{D}'), \mathcal{O}_{\mathcal{D}'}) \ar[r]^{(g, g^\sharp)} &
(\Sh(\mathcal{D}), \mathcal{O}_\mathcal{D})
}
$$
它由环化拓扑斯组成，并假设有函子
$$
\xymatrix{
\mathcal{C}' \ar[r]_{v'} &
\mathcal{C} \\
\mathcal{D}' \ar[r]^v \ar[u]^{u'} &
\mathcal{D} \ar[u]_u
}
$$
使得（采用《站点》第 \ref{sites-section-morphism-sites} 节和第
\ref{sites-section-cocontinuous-morphism-topoi} 节的记号）有
\begin{enumerate}
\item $u$ 和 $u'$ 连续，并分别给出态射 $f$ 和 $f'$；
\item $v$ 和 $v'$ 余连续，并分别给出态射 $g$ 和 $g'$；
\item $u \circ v = v' \circ u'$,
\item $v$ 和 $v'$ 既连续又余连续；
\item $g^{-1}\mathcal{O}_{\mathcal{D}} = \mathcal{O}_{\mathcal{D}'}$
且 $(g')^{-1}\mathcal{O}_{\mathcal{C}} = \mathcal{O}_{\mathcal{C}'}$。
\end{enumerate}
则作为函子 $D^+(\mathcal{O}_\mathcal{C}) \to
D^+(\mathcal{O}_{\mathcal{D}'})$，有
$Rf'_* \circ (g')^* = g^* \circ Rf_*$。若此外
\begin{enumerate}
\item[(6)] $g'_! : \textit{Ab}(\mathcal{C}') \to \textit{Ab}(\mathcal{C})$
是正合的\footnote{若 $\mathcal{C}'$ 中存在纤维积和等化子，且 $v'$ 与它们交换，
则此条件成立，见《站点上的模》引理
\ref{sites-modules-lemma-exactness-lower-shriek}。}，
\end{enumerate}
则作为函子 $D(\mathcal{O}_\mathcal{C}) \to
D(\mathcal{O}_{\mathcal{D}'})$，有
$Rf'_* \circ (g')^* = g^* \circ Rf_*$。
\end{lemma}

\begin{proof}
由条件 (5)，有
$g^* = Lg^* = g^{-1}$ 且 $(g')^* = L(g')^* = (g')^{-1}$。
根据引理 \ref{sites-cohomology-lemma-modules-abelian-unbounded}，只需在阿贝尔层的
导出范畴 $D^+(\mathcal{C})$ 或 $D(\mathcal{C})$ 上证明结论。

\medskip\noindent
选取对象 $K \in D^+(\mathcal{C})$。令 $\mathcal{I}^\bullet$ 是
$\mathcal{C}$ 上表示 $K$ 的下有界内射阿贝尔层复形。由引理
\ref{sites-cohomology-lemma-pullback-injective-pre-limp}，对所有 $p > 0$、任意 $q$
以及任意 $U' \in \Ob(\mathcal{C}')$，都有
$H^p(U', (g')^{-1}\mathcal{I}^q) = 0$。回忆
$R^pf'_*(g')^{-1}\mathcal{I}^q$ 是与预层
$V' \mapsto H^p(u'(V'), (g')^{-1}\mathcal{I}^q)$ 关联的层，见引理
\ref{sites-cohomology-lemma-higher-direct-images}。因此 $(g')^{-1}\mathcal{I}^q$
对函子 $f'_*$ 右无环。由 Leray 无环性引理（《导出范畴》引理
\ref{derived-lemma-leray-acyclicity}），
$f'_*(g')^*\mathcal{I}^\bullet$ 表示 $Rf'_*(g')^{-1}K$。
由《站点上的模》引理
\ref{sites-modules-lemma-special-square-continuous}，有
$f'_*(g')^{-1}\mathcal{I}^\bullet = g^{-1}f_*\mathcal{I}^\bullet$。
由于 $g^{-1}f_*\mathcal{I}^\bullet$ 表示 $g^{-1}Rf_*K$，结论成立。

\medskip\noindent
选取对象 $K \in D(\mathcal{C})$。令 $\mathcal{I}^\bullet$ 是
$\mathcal{C}$ 上表示 $K$ 的阿贝尔层 K-内射复形。由《导出范畴》引理
\ref{derived-lemma-adjoint-preserve-K-injectives} 和假设 (6)，
$(g')^{-1}\mathcal{I}^\bullet$ 是 $\mathcal{C}'$ 上阿贝尔层的
K-内射复形。由《站点上的模》引理
\ref{sites-modules-lemma-special-square-continuous}，有
$f'_*(g')^{-1}\mathcal{I}^\bullet = g^{-1}f_*\mathcal{I}^\bullet$。
由于 $f_*\mathcal{I}^\bullet$ 表示 $Rf_*K$，而
$f'_*(g')^{-1}\mathcal{I}^\bullet$ 表示 $Rf'_*(g')^{-1}K$，结论成立。
\end{proof}








\section{纤维化范畴的导出下 shriek}
\label{sites-cohomology-section-derived-lower-shriek-fibred}

\noindent
本节详细处理第 \ref{sites-cohomology-section-derived-lower-shriek} 节所讨论情形的若干特殊情况。
我们将确保模上的下 shriek 与阿贝尔群层上的下 shriek 相等。
建议读者初次阅读时跳过本节。

\begin{situation}
\label{sites-cohomology-situation-fibred-category}
设 $(\mathcal{D}, \mathcal{O}_\mathcal{D})$ 为环化站点，且
$p : \mathcal{C} \to \mathcal{D}$ 为纤维化范畴。赋予 $\mathcal{C}$
从 $\mathcal{D}$ 继承的拓扑（《叠》第
\ref{stacks-section-topology} 节）。以
$\pi : \Sh(\mathcal{C}) \to \Sh(\mathcal{D})$ 表示与 $p$ 关联的
拓扑斯态射（《叠》引理
\ref{stacks-lemma-topology-inherited-functorial}）。置
$\mathcal{O}_\mathcal{C} = \pi^{-1}\mathcal{O}_\mathcal{D}$，
从而得到环化拓扑斯的态射
$$
\pi :
(\Sh(\mathcal{C}), \mathcal{O}_\mathcal{C})
\longrightarrow
(\Sh(\mathcal{D}), \mathcal{O}_\mathcal{D})
$$
\end{situation}

\begin{lemma}
\label{sites-cohomology-lemma-fibred-category-with-object}
假设和记号如情形 \ref{sites-cohomology-situation-fibred-category}。对
$U \in \Ob(\mathcal{C})$，考虑诱导的拓扑斯态射
$$
\pi_U : \Sh(\mathcal{C}/U) \longrightarrow \Sh(\mathcal{D}/p(U))
$$
则存在拓扑斯态射
$$
\sigma : \Sh(\mathcal{D}/p(U)) \to \Sh(\mathcal{C}/U)
$$
使得 $\pi_U \circ \sigma = \text{id}$ 且 $\sigma^{-1} = \pi_{U, *}$。
\end{lemma}

\begin{proof}
注意，$\pi_U$ 是 $\pi$ 在局部化上的限制，见《站点》引理
\ref{sites-lemma-localize-cocontinuous}。对 $\mathcal{D}/p(U)$ 的对象
$V \to p(U)$，以 $V \times_{p(U)} U \to U$ 表示 $\mathcal{C}$ 中
$\mathcal{D}$ 上的强笛卡儿态射；它因 $p$ 是纤维化范畴而存在。函子
$$
v : \mathcal{D}/p(U) \to \mathcal{C}/U,\quad
V/p(U) \mapsto V \times_{p(U)} U/U
$$
依照 $\mathcal{C}$ 上拓扑的定义是连续的。此外，由强笛卡儿态射的定义，
它是 $p$ 的右伴随。因此我们处于《站点》第
\ref{sites-section-cocontinuous-adjoint} 节所讨论的情形，并且层
$\pi_{U, *}\mathcal{F}$ 等于
$V \mapsto \mathcal{F}(V \times_{p(U)} U)$（尤其见《站点》引理
\ref{sites-lemma-have-functor-other-way-morphism}）。

\medskip\noindent
但这里还有更多结论。事实上，函子 $v$ 也余连续（因为 $\mathcal{C}$
的覆盖中的所有态射都是强笛卡儿的）。因此 $v$ 定义了引理所述的态射
$\sigma$。等式 $\sigma^{-1} = \pi_{U, *}$ 由定义立即得到。由于
$\pi_U^{-1}\mathcal{G}$ 由规则
$U'/U \mapsto \mathcal{G}(p(U')/p(U))$ 给出，遂有
$\sigma^{-1} \circ \pi_U^{-1} = \text{id}$，这证明了等式
$\pi_U \circ \sigma = \text{id}$。
\end{proof}

\begin{situation}
\label{sites-cohomology-situation-morphism-fibred-categories}
设 $(\mathcal{D}, \mathcal{O}_\mathcal{D})$ 为环化站点。设
$u : \mathcal{C}' \to \mathcal{C}$ 为 $\mathcal{D}$ 上纤维化范畴的
$1$-态射（《范畴》定义
\ref{categories-definition-fibred-categories-over-C}）。赋予
$\mathcal{C}$ 和 $\mathcal{C}'$ 各自的继承拓扑（《叠》定义
\ref{stacks-definition-topology-inherited}），并令
$\pi : \Sh(\mathcal{C}) \to \Sh(\mathcal{D})$,
$\pi' : \Sh(\mathcal{C}') \to \Sh(\mathcal{D})$，以及
$g : \Sh(\mathcal{C}') \to \Sh(\mathcal{C})$
为相应的拓扑斯态射（《叠》引理
\ref{stacks-lemma-topology-inherited-functorial}）。置
$\mathcal{O}_\mathcal{C} = \pi^{-1}\mathcal{O}_\mathcal{D}$ 且
$\mathcal{O}_{\mathcal{C}'} = (\pi')^{-1}\mathcal{O}_\mathcal{D}$。
注意 $g^{-1}\mathcal{O}_\mathcal{C} = \mathcal{O}_{\mathcal{C}'}$，故
$$
\xymatrix{
(\Sh(\mathcal{C}'), \mathcal{O}_{\mathcal{C}'}) \ar[rd]_{\pi'} \ar[rr]_g & &
(\Sh(\mathcal{C}), \mathcal{O}_\mathcal{C}) \ar[ld]^\pi \\
& (\Sh(\mathcal{D}), \mathcal{O}_\mathcal{D})
}
$$
是环化拓扑斯态射的交换图。
\end{situation}

\begin{lemma}
\label{sites-cohomology-lemma-morphism-fibred-categories-with-object}
假设和记号如情形 \ref{sites-cohomology-situation-morphism-fibred-categories}。
对 $U' \in \Ob(\mathcal{C}')$，置 $U = u(U')$ 且 $V = p'(U')$，
并考虑诱导的环化拓扑斯态射
$$
\xymatrix{
(\Sh(\mathcal{C}'/U'), \mathcal{O}_{U'}) \ar[rd]_{\pi'_{U'}} \ar[rr]_{g'} & &
(\Sh(\mathcal{C}), \mathcal{O}_U) \ar[ld]^{\pi_U} \\
& (\Sh(\mathcal{D}/V), \mathcal{O}_V)
}
$$
则存在拓扑斯态射
$$
\sigma' : \Sh(\mathcal{D}/V) \to \Sh(\mathcal{C}'/U'),
$$
使得置 $\sigma = g' \circ \sigma'$ 后，有
$\pi'_{U'} \circ \sigma' = \text{id}$、$\pi_U \circ \sigma = \text{id}$、
$(\sigma')^{-1} = \pi'_{U', *}$ 且 $\sigma^{-1} = \pi_{U, *}$。
\end{lemma}

\begin{proof}
令 $v' : \mathcal{D}/V \to \mathcal{C}'/U'$ 为引理
\ref{sites-cohomology-lemma-fibred-category-with-object} 的证明中从
$p' : \mathcal{C}' \to \mathcal{D}'$ 和对象 $U'$ 出发所构造的函子。
由于 $u$ 是 $\mathcal{D}$ 上纤维化范畴的 $1$-态射，它把强笛卡儿态射
变为强笛卡儿态射，故函子 $v = u \circ v'$ 正是引理
\ref{sites-cohomology-lemma-fibred-category-with-object} 的证明中关于
$p : \mathcal{C} \to \mathcal{D}$ 和 $U$ 的函子。因此本引理由该引理得到。
\end{proof}

\begin{lemma}
\label{sites-cohomology-lemma-properties-lower-shriek-fibred-category}
假设和记号如情形 \ref{sites-cohomology-situation-morphism-fibred-categories}。
\begin{enumerate}
\item 在模和阿贝尔层上，$g^* = g^{-1}$ 分别有左伴随
$g_! : \textit{Mod}(\mathcal{O}_{\mathcal{C}'}) \to
\textit{Mod}(\mathcal{O}_\mathcal{C})$ 和
$g_!^{\textit{Ab}} : \textit{Ab}(\mathcal{C}') \to \textit{Ab}(\mathcal{C})$
。
\item 图
$$
\xymatrix{
\textit{Mod}(\mathcal{O}_{\mathcal{C}'}) \ar[d] \ar[r]_{g_!} &
\textit{Mod}(\mathcal{O}_\mathcal{C}) \ar[d] \\
\textit{Ab}(\mathcal{C}') \ar[r]^{g_!^{\textit{Ab}}} &
\textit{Ab}(\mathcal{C})
}
$$
交换。
\item 在模和阿贝尔层的导出范畴上，$g^* = g^{-1}$ 分别有左伴随
$Lg_! : D(\mathcal{O}_{\mathcal{C}'}) \to D(\mathcal{O}_\mathcal{C})$
以及
$Lg_!^{\textit{Ab}} : D(\mathcal{C}') \to D(\mathcal{C})$
。
\item 图
$$
\xymatrix{
D(\mathcal{O}_{\mathcal{C}'}) \ar[d] \ar[r]_{Lg_!} &
D(\mathcal{O}_\mathcal{C}) \ar[d] \\
D(\mathcal{C}') \ar[r]^{Lg_!^{\textit{Ab}}} &
D(\mathcal{C})
}
$$
交换。
\end{enumerate}
\end{lemma}

\begin{proof}
函子 $u$ 连续且余连续，见《叠》引理
\ref{stacks-lemma-topology-inherited-functorial}。因此函子
$g_!$、$g_!^{\textit{Ab}}$、$Lg_!$ 和 $Lg_!^{\textit{Ab}}$ 的存在性见
《站点上的模》第 \ref{sites-modules-section-exactness-lower-shriek} 节、
第 \ref{sites-modules-section-lower-shriek-modules} 节以及第
\ref{sites-cohomology-section-derived-lower-shriek} 节。

\medskip\noindent
为证明 (2)，只需证明规范映射
$$
g_!^{\textit{Ab}}j_{U'!}\mathcal{O}_{U'} \to j_{u(U')!}\mathcal{O}_{u(U')}
$$
对 $\mathcal{C}'$ 的所有对象 $U'$ 均为同构，见《站点上的模》评注
\ref{sites-modules-remark-when-shriek-equal}。类似地，为证明 (4)，
只需证明规范映射
$$
Lg_!^{\textit{Ab}}j_{U'!}\mathcal{O}_{U'} \to j_{u(U')!}\mathcal{O}_{u(U')}
$$
对 $\mathcal{C}'$ 的所有对象 $U'$ 在 $D(\mathcal{C})$ 中均为同构，
见评注 \ref{sites-cohomology-remark-when-derived-shriek-equal}。这也会推出前一个公式，
因此下面证明这一点。

\medskip\noindent
我们将使用如下事实：对局部化态射 $j$，函子 $j_!$ 与
$j_!^{\textit{Ab}}$ 一致（见《站点上的模》评注
\ref{sites-modules-remark-localize-shriek-equal}），并且 $j_!$ 正合
（《站点上的模》引理
\ref{sites-modules-lemma-extension-by-zero-exact}）。采用引理
\ref{sites-cohomology-lemma-morphism-fibred-categories-with-object} 的记号。由于
$Lg_!^{\textit{Ab}} \circ j_{U'!} = j_{U!} \circ L(g')^{\textit{Ab}}_!$
（由《站点》引理 \ref{sites-lemma-localize-cocontinuous} 的交换性
以及伴随函子的唯一性），只需证明
$L(g')^{\textit{Ab}}_!\mathcal{O}_{U'} = \mathcal{O}_U$。利用引理
\ref{sites-cohomology-lemma-morphism-fibred-categories-with-object} 的结论，对
$D(\mathcal{C}/u(U'))$ 的任意对象 $E$，有如下等式链：
\begin{align*}
\Hom_{D(\mathcal{C}/U)}(L(g')_!^{\textit{Ab}}\mathcal{O}_{U'}, E)
& =
\Hom_{D(\mathcal{C}'/U')}(\mathcal{O}_{U'}, (g')^{-1}E) \\
& =
\Hom_{D(\mathcal{C}'/U')}((\pi'_{U'})^{-1}\mathcal{O}_V, (g')^{-1}E) \\
& =
\Hom_{D(\mathcal{D}/V)}(\mathcal{O}_V, R\pi'_{U', *}(g')^{-1}E) \\
& =
\Hom_{D(\mathcal{D}/V)}(\mathcal{O}_V, (\sigma')^{-1}(g')^{-1}E) \\
& =
\Hom_{D(\mathcal{D}/V)}(\mathcal{O}_V, \sigma^{-1}E) \\
& =
\Hom_{D(\mathcal{D}/V)}(\mathcal{O}_V, \pi_{U, *}E) \\
& =
\Hom_{D(\mathcal{C}/U)}(\pi_U^{-1}\mathcal{O}_V, E) \\
& =
\Hom_{D(\mathcal{C}/U)}(\mathcal{O}_U, E)
\end{align*}
由 Yoneda 引理即得结论。
\end{proof}

\begin{remark}
\label{sites-cohomology-remark-fibred-category}
假设和记号如情形 \ref{sites-cohomology-situation-fibred-category}。注意，置
$\mathcal{C}' = \mathcal{D}$，并令 $u$ 等于 $\mathcal{C}$ 的结构函子，
就得到情形 \ref{sites-cohomology-situation-morphism-fibred-categories} 中的情形。因此引理
\ref{sites-cohomology-lemma-properties-lower-shriek-fibred-category} 表明，我们有函子
$\pi_!$、$\pi_!^{\textit{Ab}}$、$L\pi_!$ 和 $L\pi_!^{\textit{Ab}}$，满足
$forget \circ \pi_! = \pi_!^{\textit{Ab}} \circ forget$ 且
$forget \circ L\pi_! = L\pi_!^{\textit{Ab}} \circ forget$.
\end{remark}

\begin{remark}
\label{sites-cohomology-remark-morphism-fibred-categories}
假设和记号如情形 \ref{sites-cohomology-situation-morphism-fibred-categories}。设
$\mathcal{F}$ 为 $\mathcal{C}$ 上的阿贝尔层，$\mathcal{F}'$ 为
$\mathcal{C}'$ 上的阿贝尔层，并设
$t : \mathcal{F}' \to g^{-1}\mathcal{F}$ 为映射。则得到规范映射
$$
L\pi'_!(\mathcal{F}') \longrightarrow L\pi_!(\mathcal{F})
$$
它由 $t$ 的伴随 $g_!\mathcal{F}' \to \mathcal{F}$、映射
$Lg_!(\mathcal{F}') \to g_!\mathcal{F}'$ 以及等式
$L\pi'_! = L\pi_! \circ Lg_!$ 给出。
\end{remark}

\begin{lemma}
\label{sites-cohomology-lemma-compute-pi-shriek}
假设和记号如情形 \ref{sites-cohomology-situation-fibred-category}。对
$\textit{Ab}(\mathcal{C})$ 中的 $\mathcal{F}$，层 $\pi_!\mathcal{F}$
是与预层
$$
V \longmapsto \colim_{\mathcal{C}_V^{opp}} \mathcal{F}|_{\mathcal{C}_V}
$$
关联的层，其限制映射如证明中所示。
\end{lemma}

\begin{proof}
以 $\mathcal{H}$ 表示引理中的规则。对 $\mathcal{D}$ 的态射
$h : V' \to V$，有纤维化范畴的拉回函子
$h^* : \mathcal{C}_V \to \mathcal{C}_{V'}$（《范畴》定义
\ref{categories-definition-pullback-functor-fibred-category}）。此外，对
$U \in \Ob(\mathcal{C}_V)$，存在覆盖 $h$ 的强笛卡儿态射
$h^*U \to U$。沿这些强笛卡儿态射作限制，定义了函子的变换
$$
\mathcal{F}|_{\mathcal{C}_V}
\longrightarrow
\mathcal{F}|_{\mathcal{C}_{V'}} \circ h^*.
$$
从而得到余极限之间的映射
$\mathcal{H}(V) \to \mathcal{H}(V')$，见《范畴》引理
\ref{categories-lemma-functorial-colimit}。

\medskip\noindent
为证明本引理，我们证明
$$
\Mor_{\textit{PSh}(\mathcal{D})}(\mathcal{H}, \mathcal{G}) =
\Mor_{\Sh(\mathcal{C})}(\mathcal{F}, \pi^{-1}\mathcal{G})
$$
对每个 $\mathcal{C}$ 上的层 $\mathcal{G}$ 成立。左端的一个元素是对
$\mathcal{C}$ 中所有 $U$ 的相容映射系
$\mathcal{F}(U) \to \mathcal{G}(p(U))$。根据我们对 $\mathcal{C}$ 上
拓扑的选择，$\pi^{-1}\mathcal{G}(U) = \mathcal{G}(p(U))$，故右端也描述
同一对象，结论成立。
\end{proof}





\section{范畴上的同调}
\label{sites-cohomology-section-homology}

\noindent
在点上的范畴这一情形，我们把左导出下 shriek 函子称为同调函子。

\begin{example}[点上的范畴]
\label{sites-cohomology-example-category-to-point}
设 $\mathcal{C}$ 为范畴。赋予 $\mathcal{C}$ 混沌拓扑（《站点》例
\ref{sites-example-indiscrete}）。于是 $\mathcal{C}$ 上预层与层相同。
令 $*$ 为仅有一个对象和一个态射的范畴，则函子
$p : \mathcal{C} \to *$ 余连续且连续。令
$\pi : \Sh(\mathcal{C}) \to \Sh(*)$ 为相应的拓扑斯态射。设 $B$ 为环。
赋予 $*$ 环层 $B$，并赋予 $\mathcal{C}$ 结构层
$\mathcal{O}_\mathcal{C} = \pi^{-1}B$，将其记为 $\underline{B}$。这样
$$
\pi : (\Sh(\mathcal{C}), \underline{B}) \to (\Sh(*), B)
$$
就是情形 \ref{sites-cohomology-situation-fibred-category} 的一个例子。由评注
\ref{sites-cohomology-remark-fibred-category}，无需区分模上的 $\pi_!$ 与阿贝尔层上的
$\pi_!$。由引理 \ref{sites-cohomology-lemma-compute-pi-shriek}，有
$\pi_!\mathcal{F} = \colim_{\mathcal{C}^{opp}} \mathcal{F}$.
因此 $L_n\pi_!$ 是取余极限函子的第 $n$ 个左导出函子。以下记
$$
H_n(\mathcal{C}, \mathcal{F}) = L_n\pi_!(\mathcal{F})
$$
并称其为 $\mathcal{F}$ 在 $\mathcal{C}$ 上的{\it 第 $n$ 个同调群}。
\end{example}

\begin{example}[计算同调]
\label{sites-cohomology-example-left-derived-colimits}
在例 \ref{sites-cohomology-example-category-to-point} 中，可如下计算函子
$H_n(\mathcal{C}, -)$。设
$\mathcal{F} \in \Ob(\textit{Ab}(\mathcal{C}))$。考虑链复形
$$
K_\bullet(\mathcal{F}) :
\ \ldots \to
\bigoplus\nolimits_{U_2 \to U_1 \to U_0} \mathcal{F}(U_0)
\to
\bigoplus\nolimits_{U_1 \to U_0} \mathcal{F}(U_0)
\to
\bigoplus\nolimits_{U_0} \mathcal{F}(U_0)
$$
其中过渡映射由下式给出：
$$
(U_2 \to U_1 \to U_0, s)
\longmapsto
(U_1 \to U_0, s) - (U_2 \to U_0, s) + (U_2 \to U_1, s|_{U_1})
$$
其他次数类似。由构造，
$$
H_0(\mathcal{C}, \mathcal{F}) =
\colim_{\mathcal{C}^{opp}} \mathcal{F} =
H_0(K_\bullet(\mathcal{F})),
$$
见《范畴》引理
\ref{categories-lemma-colimits-coproducts-coequalizers}。
$K_\bullet(\mathcal{F})$ 的构造关于 $\mathcal{F}$ 是函子性的，并把
$\textit{Ab}(\mathcal{C})$ 的短正合列变为复形的短正合列。因此函子列
$\mathcal{F} \mapsto H_n(K_\bullet(\mathcal{F}))$ 构成 $\delta$-函子，见
《同调》定义 \ref{homology-definition-cohomological-delta-functor} 和引理
\ref{homology-lemma-long-exact-sequence-cochain}。当
$\mathcal{F} = j_{U!}\mathbf{Z}_U$ 时，复形 $K_\bullet(\mathcal{F})$
是单纯集 $X_\bullet$ 上自由 $\mathbf{Z}$-模所关联的复形，其各项为
$$
X_n = \coprod\nolimits_{U_n \to \ldots \to U_1 \to U_0}
\Mor_\mathcal{C}(U_0, U)
$$
该单纯集同伦等价于单点集 $\{*\}$ 上的常值单纯集。事实上，映射
$X_\bullet \to \{*\}$ 是显然的，映射 $\{*\} \to X_n$ 把 $*$ 映到
$(U \to \ldots \to U, \text{id}_U)$，而映射
$$
h_{n, i} : X_n \longrightarrow X_n
$$
（《单纯形》引理 \ref{simplicial-lemma-relations-homotopy}）定义两个映射
$X_\bullet \to X_\bullet$ 之间的同伦，并由规则
$$
h_{n, i} :
(U_n \to \ldots \to U_0, f)
\longmapsto
(U_n \to \ldots \to U_i \to U \to \ldots \to U, \text{id})
$$
给出，其中 $i > 0$，且 $h_{n, 0} = \text{id}$。验证略去。
这表明 $K_\bullet(j_{U!}\mathbf{Z}_U)$ 在负次数的上同调为零
（由《单纯形》评注 \ref{simplicial-remark-homotopy-better} 的函子性
以及《单纯形》引理 \ref{simplicial-lemma-homotopy-s-N} 的结论）。
因此，例如由《同调》引理
\ref{homology-lemma-efface-implies-universal} 和《导出范畴》引理
\ref{derived-lemma-right-derived-delta-functor} 的对偶，
$K_\bullet(\mathcal{F})$ 计算 $H_0(\mathcal{C}, -)$ 的左导出函子
$H_n(\mathcal{C}, -)$。
\end{example}

\begin{example}
\label{sites-cohomology-example-morphism-categories}
设 $u : \mathcal{C}' \to \mathcal{C}$ 为函子。如例
\ref{sites-cohomology-example-category-to-point}，赋予 $\mathcal{C}'$ 和 $\mathcal{C}$
混沌拓扑。令 $*$ 为仅有一个对象和一个态射的范畴，则函子
$u$、$\mathcal{C}' \to *$ 以及 $\mathcal{C} \to *$ 都余连续且连续。令
$g : \Sh(\mathcal{C}') \to \Sh(\mathcal{C})$、
$\pi' : \Sh(\mathcal{C}') \to \Sh(*)$ 和
$\pi : \Sh(\mathcal{C}) \to \Sh(*)$
为相应的拓扑斯态射。设 $B$ 为环。赋予 $*$ 环层 $B$，并赋予
$\mathcal{C}'$、$\mathcal{C}$ 常值层 $\underline{B}$。这样
$$
\xymatrix{
(\Sh(\mathcal{C}'), \underline{B}) \ar[rd]_{\pi'} \ar[rr]_g & &
(\Sh(\mathcal{C}), \underline{B}) \ar[ld]^\pi \\
& (\Sh(*), B)
}
$$
就是情形 \ref{sites-cohomology-situation-morphism-fibred-categories} 的一个例子。因此引理
\ref{sites-cohomology-lemma-properties-lower-shriek-fibred-category} 可应用于 $g$，
故无需区分模上的 $g_!$ 与阿贝尔层上的 $g_!$。特别地，评注
\ref{sites-cohomology-remark-morphism-fibred-categories} 给出规范映射
$$
H_n(\mathcal{C}', \mathcal{F}')
\longrightarrow
H_n(\mathcal{C}, \mathcal{F})
$$
只要给定 $\textit{Ab}(\mathcal{C})$ 中的 $\mathcal{F}$、
$\textit{Ab}(\mathcal{C}')$ 中的 $\mathcal{F}'$ 以及映射
$t : \mathcal{F}' \to g^{-1}\mathcal{F}$。用例
\ref{sites-cohomology-example-left-derived-colimits} 给出的同调计算表示时，
这些映射来自复形映射
$$
K_\bullet(\mathcal{F}') \longrightarrow K_\bullet(\mathcal{F})
$$
它由规则
$$
(U'_n \to \ldots \to U'_0, s') \longmapsto
(u(U'_n) \to \ldots \to u(U'_0), t(s'))
$$
给出，记号含义显然。
\end{example}

\begin{remark}
\label{sites-cohomology-remark-map-evaluation-to-derived}
记号和假设如例 \ref{sites-cohomology-example-category-to-point}。设
$\mathcal{F}^\bullet$ 为 $\mathcal{C}$ 上阿贝尔层的有界复形。
对 $\mathcal{C}$ 的任意对象 $U$，存在规范映射
$$
\mathcal{F}^\bullet(U) \longrightarrow L\pi_!(\mathcal{F}^\bullet)
$$
它位于 $D(\textit{Ab})$ 中。若 $\mathcal{F}^\bullet$ 是
$\underline{B}$-模复形，则该映射位于 $D(B)$ 中。为证明这一点，注意
$L\pi_!(\mathcal{F}^\bullet)$ 可通过取拟同构
$\mathcal{P}^\bullet \to \mathcal{F}^\bullet$ 来计算，其中
$\mathcal{P}^\bullet$ 是投射对象的复形。然而，由于拓扑是混沌的，
$\mathcal{P}^\bullet(U) \to \mathcal{F}^\bullet(U)$ 是拟同构，
故可在 $D(\textit{Ab})$（相应地在 $D(B)$）中求逆。再与从把 $\pi_!$
计算为余极限而来的规范映射
$\mathcal{P}^\bullet(U) \to \pi_!(\mathcal{P}^\bullet)$ 复合，
即得所需箭头。
\end{remark}

\begin{lemma}
\label{sites-cohomology-lemma-initial-final}
记号和假设如例 \ref{sites-cohomology-example-category-to-point}。若 $\mathcal{C}$ 有始对象
或终对象，则在 $D(\textit{Ab})$（相应地在 $D(B)$）上有
$L\pi_! \circ \pi^{-1} = \text{id}$。
\end{lemma}

\begin{proof}
若 $\mathcal{C}$ 有始对象，则 $\pi_!$ 由在该对象上取值来计算，
断言显然。若 $\mathcal{C}$ 有终对象，则 $R\pi_*$ 由在该对象上取值来计算，
故在 $D(\textit{Ab})$（相应地在 $D(B)$）上有
$R\pi_* \circ \pi^{-1} \cong \text{id}$。这表明
$\pi^{-1} : D(\textit{Ab}) \to D(\mathcal{C})$,
相应地，$\pi^{-1} : D(B) \to D(\underline{B})$ 完全忠实，见《范畴》引理
\ref{categories-lemma-adjoint-fully-faithful}。同一引理于是推出
$L\pi_! \circ \pi^{-1} = \text{id}$，正如所需。
\end{proof}

\begin{lemma}
\label{sites-cohomology-lemma-change-of-rings}
记号和假设如例 \ref{sites-cohomology-example-category-to-point}。设 $B \to B'$ 为环同态。
考虑环化拓扑斯的交换图
$$
\xymatrix{
(\Sh(\mathcal{C}), \underline{B}) \ar[d]_\pi &
(\Sh(\mathcal{C}), \underline{B'}) \ar[d]^{\pi'} \ar[l]^h \\
(*, B) & (*, B') \ar[l]_f
}
$$
则 $L\pi_! \circ Lh^* = Lf^* \circ L\pi'_!$。
\end{lemma}

\begin{proof}
两个函子都是显然函子 $D(B') \to D(\underline{B})$ 的右伴随。
\end{proof}

\begin{lemma}
\label{sites-cohomology-lemma-compute-by-cosimplicial-resolution}
记号和假设如例 \ref{sites-cohomology-example-category-to-point}。设 $U_\bullet$ 为
$\mathcal{C}$ 中的余单纯对象，使得对每个
$U \in \Ob(\mathcal{C})$，单纯集 $\Mor_\mathcal{C}(U_\bullet, U)$
同伦等价于单点集上的常值单纯集。则
$$
L\pi_!(\mathcal{F}) = \mathcal{F}(U_\bullet)
$$
在 $D(\textit{Ab})$（相应地在 $D(B)$）中成立，并且关于
$\textit{Ab}(\mathcal{C})$（相应地 $\textit{Mod}(\underline{B})$）中的
$\mathcal{F}$ 是函子性的。
\end{lemma}

\begin{proof}
由引理 \ref{sites-cohomology-lemma-properties-lower-shriek-fibred-category}，模上的
$L\pi_!$ 与阿贝尔层上的 $L\pi_!$ 一致，故只需在 $\mathcal{F}$ 为
阿贝尔层时证明。对 $U \in \Ob(\mathcal{C})$，阿贝尔层
$j_{U!}\mathbf{Z}_U$ 是 $\textit{Ab}(\mathcal{C})$ 的投射对象，因为
$\Hom(j_{U!}\mathbf{Z}_U, \mathcal{F}) = \mathcal{F}(U)$，而在混沌拓扑下
取截面是正合函子。由《站点上的模》引理
\ref{sites-modules-lemma-module-quotient-flat}，每个阿贝尔层都是若干
$j_{U!}\mathbf{Z}_U$ 的直和的商。因此可选取分辨
$$
\ldots \to \mathcal{G}^{-1} \to \mathcal{G}^0 \to \mathcal{F} \to 0
$$
其各项均为上述形式的层的直和，并取
$L\pi_!(\mathcal{F}) = \pi_!(\mathcal{G}^\bullet)$，以此计算
$L\pi_!(\mathcal{F})$。考虑双复形
$A^{\bullet, \bullet} = \mathcal{G}^\bullet(U_\bullet)$。映射
$\mathcal{G}^0 \to \mathcal{F}$ 给出复形映射
$A^{0, \bullet} \to \mathcal{F}(U_\bullet)$。由于 $\pi_!$ 通过取
$\mathcal{C}^{opp}$ 上的余极限来计算（引理
\ref{sites-cohomology-lemma-compute-pi-shriek}），可知两个复合
$\mathcal{G}^m(U_1) \to \mathcal{G}^m(U_0) \to \pi_!\mathcal{G}^m$
相等。因此得到规范的复形映射
$$
\text{Tot}(A^{\bullet, \bullet})
\longrightarrow
\pi_!(\mathcal{G}^\bullet) = L\pi_!(\mathcal{F})
$$
为证明本引理，只需证明复形
$$
\ldots \to \mathcal{G}^m(U_1) \to \mathcal{G}^m(U_0) \to
\pi_!\mathcal{G}^m \to 0
$$
正合，见《同调》引理
\ref{homology-lemma-double-complex-gives-resolution}。由于层
$\mathcal{G}^m$ 是若干 $j_{U!}\mathbf{Z}_U$ 的直和，归结为
$\mathcal{G} = j_{U!}\mathbf{Z}_U$。复形
$j_{U!}\mathbf{Z}_U(U_\bullet)$ 是单纯集
$\Mor_\mathcal{C}(U_\bullet, U)$ 上自由 $\mathbf{Z}$-模所关联的
阿贝尔群复形；我们假设该单纯集同伦等价于单点集。因此
$$
j_{U!}\mathbf{Z}_U(U_\bullet) \to \mathbf{Z}
$$
是阿贝尔群的同伦等价，因而是拟同构（《单纯形》评注
\ref{simplicial-remark-homotopy-better} 和引理
\ref{simplicial-lemma-homotopy-s-N}）。又因引理
\ref{sites-cohomology-lemma-properties-lower-shriek-fibred-category} 的证明已表明
$\pi_!j_{U!}\mathbf{Z}_U = \mathbf{Z}$，证明完成。
\end{proof}

\begin{lemma}
\label{sites-cohomology-lemma-get-it-now}
记号和假设如例 \ref{sites-cohomology-example-morphism-categories}。若存在
$\mathcal{C}'$ 的余单纯对象 $U'_\bullet$，使得引理
\ref{sites-cohomology-lemma-compute-by-cosimplicial-resolution} 同时适用于
$\mathcal{C}'$ 中的 $U'_\bullet$ 和 $\mathcal{C}$ 中的
$u(U'_\bullet)$，则作为函子 $D(\mathcal{C}) \to D(\textit{Ab})$
（相应地 $D(\mathcal{C}, \underline{B}) \to D(B)$），有
$L\pi'_! \circ g^{-1} = L\pi_!$。
\end{lemma}

\begin{proof}
这立即由引理 \ref{sites-cohomology-lemma-compute-by-cosimplicial-resolution} 以及
$g^{-1}$ 由与 $u$ 预复合给出这一事实得到。
\end{proof}

\begin{lemma}
\label{sites-cohomology-lemma-product-categories}
设 $\mathcal{C}_i$（$i = 1, 2$）为范畴，
$u_i : \mathcal{C}_1 \times \mathcal{C}_2 \to \mathcal{C}_i$ 为投影函子。
设 $B$ 为环，并设
$g_i : (\Sh(\mathcal{C}_1 \times \mathcal{C}_2), \underline{B}) \to
(\Sh(\mathcal{C}_i), \underline{B})$ 为相应的环化拓扑斯态射，见例
\ref{sites-cohomology-example-morphism-categories}。对 $K_i \in D(\mathcal{C}_i, B)$，有
$$
L(\pi_1 \times \pi_2)_!(
g_1^{-1}K_1 \otimes_{\underline{B}}^\mathbf{L} g_2^{-1}K_2)
=
L\pi_{1, !}(K_1) \otimes_B^\mathbf{L} L\pi_{2, !}(K_2)
$$
该等式在 $D(B)$ 中成立，记号含义显然。
\end{lemma}

\begin{proof}
由于两端都与余极限交换，只需对
$K_1 = j_{U!}\underline{B}_U$ 和 $K_2 = j_{V!}\underline{B}_V$
证明，其中 $U \in \Ob(\mathcal{C}_1)$ 且
$V \in \Ob(\mathcal{C}_2)$。见引理
\ref{sites-cohomology-lemma-existence-derived-lower-shriek} 中 $L\pi_!$ 的构造。
在此情形，
$$
g_1^{-1}K_1 \otimes_{\underline{B}}^\mathbf{L} g_2^{-1}K_2 =
g_1^{-1}K_1 \otimes_{\underline{B}} g_2^{-1}K_2 =
j_{(U, V)!}\underline{B}_{(U, V)}
$$
验证略去。引理公式的左右两端都取值为 $B$，故结论成立，见引理
\ref{sites-cohomology-lemma-existence-derived-lower-shriek} 中 $L\pi_!$ 的构造。
\end{proof}

\begin{lemma}
\label{sites-cohomology-lemma-eilenberg-zilber}
记号和假设如例 \ref{sites-cohomology-example-category-to-point}。若存在
$\mathcal{C}$ 的余单纯对象 $U_\bullet$，使得引理
\ref{sites-cohomology-lemma-compute-by-cosimplicial-resolution} 适用，则
$$
L\pi_!(K_1 \otimes^\mathbf{L}_{\underline{B}} K_2) =
L\pi_!(K_1) \otimes^\mathbf{L}_B L\pi_!(K_2)
$$
对所有 $K_i \in D(\underline{B})$ 成立。
\end{lemma}

\begin{proof}
考虑范畴和函子的图
$$
\xymatrix{
& & \mathcal{C} \\
\mathcal{C} \ar[r]^-u &
\mathcal{C} \times \mathcal{C} \ar[rd]^{u_2} \ar[ru]_{u_1} \\
& & \mathcal{C}
}
$$
其中 $u$ 为对角函子，$u_i$ 为投影函子。这给出环化拓扑斯态射
$g$、$g_1$、$g_2$。对 $\mathcal{C}$ 的任意对象 $(U_1, U_2)$，有
$$
\Mor_{\mathcal{C} \times \mathcal{C}}(u(U_\bullet), (U_1, U_2)) =
\Mor_\mathcal{C}(U_\bullet, U_1) \times \Mor_\mathcal{C}(U_\bullet, U_2)
$$
由《单纯形》引理 \ref{simplicial-lemma-products-homotopy}，它同伦等价于一点。
因此引理 \ref{sites-cohomology-lemma-get-it-now} 给出：对
$D(\mathcal{C} \times \mathcal{C}, B)$ 中任意 $K$，有
$L\pi_!(g^{-1}K) = L(\pi \times \pi)_!(K)$。取
$K = g_1^{-1}K_1 \otimes_B^\mathbf{L} g_2^{-1}K_2$。由于
$g^{-1} = g^* = Lg^*$ 与导出张量积交换（引理
\ref{sites-cohomology-lemma-pullback-tensor-product}），有
$g^{-1}K = K_1 \otimes^\mathbf{L}_{\underline{B}} K_2$。
最后应用引理 \ref{sites-cohomology-lemma-product-categories}。
\end{proof}

\begin{remark}[单纯模]
\label{sites-cohomology-remark-simplicial-modules}
令 $\mathcal{C} = \Delta$，并令 $B$ 为任意环。这是例
\ref{sites-cohomology-example-category-to-point} 的一个特殊情形，其中引理
\ref{sites-cohomology-lemma-compute-by-cosimplicial-resolution} 的假设成立。事实上，
令 $U_\bullet$ 为由恒等函子给出的 $\Delta$ 的余单纯对象。为验证条件，
必须证明对 $[m] \in \Ob(\Delta)$，单纯集
$\Delta[m] : n \mapsto \Mor_\Delta([n], [m])$ 同伦等价于一点。
这在《单纯形》例 \ref{simplicial-example-simplex-contractible} 中说明。

\medskip\noindent
在此情形，范畴 $\textit{Mod}(\underline{B})$ 就是单纯 $B$-模范畴，
而函子 $L\pi_!$ 把单纯 $B$-模 $M_\bullet$ 映到其关联的 $B$-模复形
$s(M_\bullet)$。因此上述结论可重新解释为关于单纯模的结论。例如，引理
\ref{sites-cohomology-lemma-eilenberg-zilber} 的一个特殊情形如下：若
$M_\bullet$、$M'_\bullet$ 是平坦单纯 $B$-模，则复形
$s(M_\bullet \otimes_B M'_\bullet)$ 拟同构于双复形
$s(M_\bullet) \otimes_B s(M'_\bullet)$ 所关联的全复形。
（提示：利用平坦性把导出张量积化为通常张量积。）
这是 Eilenberg--Zilber 定理的一个特殊情形，见 \cite{Eilenberg-Zilber}。
\end{remark}

\begin{lemma}
\label{sites-cohomology-lemma-O-homology-qis}
设 $\mathcal{C}$ 为范畴（赋予混沌拓扑）。设
$\mathcal{O} \to \mathcal{O}'$ 为 $\mathcal{C}$ 上环层的映射。假设
\begin{enumerate}
\item 存在引理 \ref{sites-cohomology-lemma-compute-by-cosimplicial-resolution} 所述的
$\mathcal{C}$ 中余单纯对象 $U_\bullet$；
\item $L\pi_!\mathcal{O} \to L\pi_!\mathcal{O}'$ 是同构。
\end{enumerate}
对 $D(\mathcal{O})$ 中的 $K$，有
$$
L\pi_!(K) = L\pi_!(K \otimes_\mathcal{O}^\mathbf{L} \mathcal{O}')
$$
该等式在 $D(\textit{Ab})$ 中成立。
\end{lemma}

\begin{proof}
注意：在本证明中，$L\pi_!$ 表示阿贝尔层上 $\pi_!$ 的左导出函子。
由于 $L\pi_!$ 与余极限交换，只需对 $\mathcal{O}$-模的上有界复形
证明（与《导出范畴》命题
\ref{derived-proposition-left-derived-exists} 的论证比较，或直接始终限于
上有界复形）。每个这样的复形都拟同构于一个上有界复形，其各项是
$j_{U!}\mathcal{O}_U$（$U \in \Ob(\mathcal{C})$）的直和，见
《站点上的模》引理 \ref{sites-modules-lemma-module-quotient-flat}。
因此只需对 $j_{U!}\mathcal{O}_U$ 证明本引理。由假设，
$$
S_\bullet = \Mor_\mathcal{C}(U_\bullet, U)
$$
是同伦等价于单点集上的常值单纯集的单纯集。置
$P_n = \mathcal{O}(U_n)$ 且 $P'_n = \mathcal{O}'(U_n)$。注意，
单纯阿贝尔群
$$
X_\bullet : n \longmapsto \bigoplus\nolimits_{s \in S_n} P_n
$$
所关联的复形由引理 \ref{sites-cohomology-lemma-compute-by-cosimplicial-resolution} 计算
$L\pi_!(j_{U!}\mathcal{O}_U)$。由于 $j_{U!}\mathcal{O}_U$ 是平坦
$\mathcal{O}$-模，有
$j_{U!}\mathcal{O}_U \otimes^\mathbf{L}_\mathcal{O} \mathcal{O}' =
j_{U!}\mathcal{O}'_U$，而它的 $L\pi_!$ 由单纯阿贝尔群
$$
X'_\bullet : n \longmapsto \bigoplus\nolimits_{s \in S_n} P'_n
$$
所关联的复形计算。把单纯集 $T_\bullet$ 映到各项为
$\bigoplus_{t \in T_n} P_n$ 的单纯阿贝尔群的规则是函子，故
$X_\bullet \to P_\bullet$ 是单纯阿贝尔群的同伦等价。类似地，把单纯集
$T_\bullet$ 映到各项为 $\bigoplus_{t \in T_n} P'_n$ 的单纯阿贝尔群的
规则也是函子。因此 $X'_\bullet \to P'_\bullet$ 是单纯阿贝尔群的
同伦等价。由假设，$P_\bullet \to P'_\bullet$ 是拟同构（因为由引理
\ref{sites-cohomology-lemma-compute-by-cosimplicial-resolution}，$P_\bullet$ 计算
$L\pi_!\mathcal{O}$，相应地 $P'_\bullet$ 计算 $L\pi_!\mathcal{O}'$）。
所以 $X_\bullet$ 与 $X'_\bullet$ 拟同构，正如所需。
\end{proof}

\begin{remark}
\label{sites-cohomology-remark-O-homology-B-homology-general}
设 $\mathcal{C}$ 和 $B$ 如例 \ref{sites-cohomology-example-category-to-point}。假设存在
引理 \ref{sites-cohomology-lemma-compute-by-cosimplicial-resolution} 所述的余单纯对象。
设 $\mathcal{O} \to \underline{B}$ 为 $\mathcal{C}$ 上环层的映射，
它诱导同构 $L\pi_!\mathcal{O} \to L\pi_!\underline{B}$。在此情形，
得到三角范畴的正合函子
$$
L\pi_! : D(\mathcal{O}) \longrightarrow D(B)
$$
事实上，对 $D(\mathcal{O})$ 的任意对象 $K$，由引理
\ref{sites-cohomology-lemma-O-homology-qis} 有
$L\pi^{\textit{Ab}}_!(K) =
L\pi^{\textit{Ab}}_!(K \otimes_{\mathcal{O}}^\mathbf{L} \underline{B})$
。因此可把所显示的函子定义为
$- \otimes^\mathbf{L}_\mathcal{O} \underline{B}$ 与函子
$L\pi_! : D(\underline{B}) \to D(B)$ 的复合。换言之，通过引理给出的
$L\pi_!(K)$ 与
$L\pi_!(K \otimes_\mathcal{O}^\mathbf{L} \underline{B})$ 的规范函子性识别，
得到 $L\pi_!(K)$ 上的 $B$-模结构。
\end{remark}







\section{计算导出下 shriek}
\label{sites-cohomology-section-calculate}

\noindent
本节应用第 \ref{sites-cohomology-section-homology} 节的结论，计算情形
\ref{sites-cohomology-situation-fibred-category} 中的 $L\pi_!$ 和情形
\ref{sites-cohomology-situation-morphism-fibred-categories} 中的 $Lg_!$。

\begin{lemma}
\label{sites-cohomology-lemma-compute-left-derived-pi-shriek-pre}
假设和记号如情形 \ref{sites-cohomology-situation-fibred-category}。对
$\textit{PAb}(\mathcal{C})$ 中的 $\mathcal{F}$ 和 $n \geq 0$，
考虑 $\mathcal{D}$ 上的阿贝尔层 $L_n(\mathcal{F})$，它是与预层
$$
V \longmapsto H_n(\mathcal{C}_V, \mathcal{F}|_{\mathcal{C}_V})
$$
关联的层，其限制映射如证明中所示。则
$L_n(\mathcal{F}) = L_n(\mathcal{F}^\#)$。
\end{lemma}

\begin{proof}
对 $\mathcal{D}$ 的态射 $h : V' \to V$，有纤维化范畴的拉回函子
$h^* : \mathcal{C}_V \to \mathcal{C}_{V'}$（《范畴》定义
\ref{categories-definition-pullback-functor-fibred-category}）。此外，对
$U \in \Ob(\mathcal{C}_V)$，存在覆盖 $h$ 的强笛卡儿态射
$h^*U \to U$。沿这些强笛卡儿态射作限制，定义了函子的变换
$$
\mathcal{F}|_{\mathcal{C}_V}
\longrightarrow
\mathcal{F}|_{\mathcal{C}_{V'}} \circ h^*.
$$
由例 \ref{sites-cohomology-example-morphism-categories}，得到所需的限制映射
$$
H_n(\mathcal{C}_V, \mathcal{F}|_{\mathcal{C}_V})
\longrightarrow
H_n(\mathcal{C}_{V'}, \mathcal{F}|_{\mathcal{C}_{V'}})
$$
以 $L_{n, p}(\mathcal{F})$ 表示该预层，于是
$L_n(\mathcal{F}) = L_{n, p}(\mathcal{F})^\#$。规范映射
$\gamma : \mathcal{F} \to \mathcal{F}^+$（《站点》定理
\ref{sites-theorem-plus}）定义规范映射
$L_{n, p}(\mathcal{F}) \to L_{n, p}(\mathcal{F}^+)$。
我们必须证明该映射在层化后成为同构。

\medskip\noindent
使用例 \ref{sites-cohomology-example-left-derived-colimits} 给出的同调计算。以
$K_\bullet(\mathcal{F}|_{\mathcal{C}_V})$ 表示 $\mathcal{F}$ 到纤维范畴
$\mathcal{C}_V$ 的限制所关联的复形。由上述说明，得到复形预层
$K_\bullet(\mathcal{F})$：
$$
V \longmapsto K_\bullet(\mathcal{F}|_{\mathcal{C}_V})
$$
它的上同调预层就是预层 $L_{n, p}(\mathcal{F})$。因此只需证明
$$
K_\bullet(\mathcal{F}) \longrightarrow K_\bullet(\mathcal{F}^+)
$$
在层化后成为同构。

\medskip\noindent
单射性。设 $V$ 为 $\mathcal{D}$ 的对象，且
$\xi \in K_n(\mathcal{F})(V)$ 为映到
$K_n(\mathcal{F}^+)(V)$ 中零元的元素。必须证明存在覆盖
$\{V_j \to V\}$，使得 $\xi|_{V_j}$ 在
$K_n(\mathcal{F})(V_j)$ 中为零。写成
$$
\xi = \sum (U_{i, n + 1} \to \ldots \to U_{i, 0}, \sigma_i)
$$
其中 $\sigma_i \in \mathcal{F}(U_{i, 0})$。可以安排为
$\mathcal{C}_V$ 的每个态射列 $U_n \to \ldots \to U_0$ 至多出现一次。
由于复形 $K_\bullet$ 的定义中的和是直和，该元素映到
$K_\bullet(\mathcal{F}^+)(V)$ 中零元的唯一可能是所有 $\sigma_i$ 都映到
$\mathcal{F}^+(U_{i, 0})$ 中的零元。由 $\mathcal{F}^+$ 的构造，存在覆盖
$\{U_{i, 0, j} \to U_{i, 0}\}$，使得
$\sigma_i|_{U_{i, 0, j}}$ 为零。由我们对 $\mathcal{C}$ 上拓扑的构造，
$U_{i, 0, j} \to U_{i, 0}$ 可写成某些态射 $V_{i, j} \to V$ 的拉回
（《范畴》定义 \ref{categories-definition-pullback-functor-fibred-category}），
而且每个 $\{V_{i, j} \to V\}$ 都是覆盖。选取同时加细所有覆盖
$\{V_{i, j} \to V\}$ 的覆盖 $\{V_j \to V\}$。显然
$\xi|_{V_j} = 0$。

\medskip\noindent
满射性。证明略去。提示：仿照单射性的证明论证。
\end{proof}

\begin{lemma}
\label{sites-cohomology-lemma-compute-left-derived-pi-shriek}
假设和记号如情形 \ref{sites-cohomology-situation-fibred-category}。对
$\textit{Ab}(\mathcal{C})$ 中的 $\mathcal{F}$ 和 $n \geq 0$，
层 $L_n\pi_!(\mathcal{F})$ 等于引理
\ref{sites-cohomology-lemma-compute-left-derived-pi-shriek-pre} 中构造的层
$L_n(\mathcal{F})$。
\end{lemma}

\begin{proof}
考虑从 $\textit{PAb}(\mathcal{C}) \to \textit{Ab}(\mathcal{C})$ 的函子列
$\mathcal{F} \mapsto L_n(\mathcal{F})$。由于对每个
$V \in \Ob(\mathcal{D})$，函子列 $H_n(\mathcal{C}_V, - )$ 构成
$\delta$-函子，函子 $\mathcal{F} \mapsto L_n(\mathcal{F})$ 也如此。
我们的目标是证明它们构成万有 $\delta$-函子。为此构造一些使这些函子
消失的阿贝尔预层。

\medskip\noindent
对 $U' \in \Ob(\mathcal{C})$，考虑阿贝尔预层
$\mathcal{F}_{U'} = j_{U'!}^{\textit{PAb}}\mathbf{Z}_{U'}$
（《站点上的模》评注
\ref{sites-modules-remark-localize-presheaves}）。回忆
$$
\mathcal{F}_{U'}(U) =
\bigoplus\nolimits_{\Mor_\mathcal{C}(U, U')} \mathbf{Z}
$$
若 $U$ 位于 $\mathcal{D})$ 中的 $V = p(U)$ 之上，且 $U'$ 位于
$V' = p(U')$ 之上，则任意态射 $a : U \to U'$ 都唯一分解为
$U \to h^*U' \to U'$，其中 $h = p(a) : V \to V'$（见《范畴》定义
\ref{categories-definition-pullback-functor-fibred-category}）。因此
$$
\mathcal{F}_{U'}|_{\mathcal{C}_V}
=
\bigoplus\nolimits_{h \in \Mor_\mathcal{D}(V, V')}
j_{h^*U'!}\mathbf{Z}_{h^*U'}
$$
其中 $j_{h^*U'} : \Sh(\mathcal{C}_V/h^*U') \to \Sh(\mathcal{C}_V)$
是局部化态射。层 $j_{h^*U'!}\mathbf{Z}_{h^*U'}$ 的高阶同调群消失
（见例 \ref{sites-cohomology-example-left-derived-colimits}）。所以对所有 $n > 0$ 和所有
$U'$，有 $L_n(\mathcal{F}_{U'}) = 0$。由此，任意阿贝尔预层
$\mathcal{F}$ 都是某个阿贝尔预层 $\mathcal{G}$ 的商，且对所有
$n > 0$ 有 $L_n(\mathcal{G}) = 0$（《站点上的模》引理
\ref{sites-modules-lemma-module-quotient-flat}）。由于
$L_n(\mathcal{F}) = L_n(\mathcal{F}^\#)$，阿贝尔层也具有同一性质。
因此函子列 $L_n(-)$ 是 $\textit{Ab}(\mathcal{C})$ 上的万有
$\delta$-函子（《同调》引理
\ref{homology-lemma-efface-implies-universal}）。由引理
\ref{sites-cohomology-lemma-compute-pi-shriek}，在 $n = 0$ 时它与
$H^{-n}(L\pi_!(-))$ 一致；再由万有 $\delta$-函子的唯一性（《同调》引理
\ref{homology-lemma-uniqueness-universal-delta-functor}）以及《导出范畴》引理
\ref{derived-lemma-right-derived-delta-functor}，得到结论。
\end{proof}

\begin{lemma}
\label{sites-cohomology-lemma-compute-left-derived-g-shriek}
假设和记号如情形 \ref{sites-cohomology-situation-morphism-fibred-categories}。对
$\mathcal{C}'$ 上的阿贝尔层 $\mathcal{F}'$，层
$L_ng_!(\mathcal{F}')$ 是与预层
$$
U \longmapsto H_n(\mathcal{I}_U, \mathcal{F}'_U)
$$
关联的层。记号和限制映射见证明。
\end{lemma}

\begin{proof}
设 $p(U) = V$。范畴 $\mathcal{I}_U$ 的对象是偶对
$(U', \varphi)$，其中 $\varphi : U \to u(U')$ 是 $\mathcal{C}$ 的态射，
且 $p(\varphi) = \text{id}_V$；换言之，$\varphi$ 是纤维范畴
$\mathcal{C}_V$ 的态射。态射
$(U'_1, \varphi_1) \to (U'_2, \varphi_2)$ 由纤维范畴
$\mathcal{C}'_V$ 的态射 $a : U'_1 \to U'_2$ 给出，并满足
$\varphi_2 = u(a) \circ \varphi_1$。预层 $\mathcal{F}'_U$ 把
$(U', \varphi)$ 映到 $\mathcal{F}'(U')$。下面构造限制映射。

\medskip\noindent
选取 $u$ 的分解
$$
\xymatrix{
\mathcal{C}' \ar@<1ex>[r]^{u'} &
\mathcal{C}'' \ar[r]^{u''} \ar@<1ex>[l]^w & \mathcal{C}
}
$$
如《范畴》引理
\ref{categories-lemma-ameliorate-morphism-fibred-categories} 所述。
则 $g_! = g''_! \circ g'_!$，导出函子也同样成立。另一方面，函子
$g'_!$ 正合，见《站点上的模》引理
\ref{sites-modules-lemma-have-left-adjoint}。因此
$Lg_!(\mathcal{F}') = Lg''_!(\mathcal{F}'')$，其中
$\mathcal{F}'' = g'_!\mathcal{F}'$。注意
$\mathcal{F}'' = h^{-1}\mathcal{F}'$，其中
$h : \Sh(\mathcal{C}'') \to \Sh(\mathcal{C}')$ 是与 $w$ 关联的
拓扑斯态射，见《站点》引理 \ref{sites-lemma-have-left-adjoint}。
函子 $u''$ 使 $\mathcal{C}''$ 成为 $\mathcal{C}$ 上的纤维化范畴，
故引理 \ref{sites-cohomology-lemma-compute-left-derived-pi-shriek} 可用于计算
$L_ng''_!$。又因为《范畴》引理
\ref{categories-lemma-ameliorate-morphism-fibred-categories} 的证明中
$\mathcal{C}''$ 的构造表明纤维范畴 $\mathcal{C}''_U$ 等于
$\mathcal{I}_U$，结论成立。此外，
$h^{-1}\mathcal{F}'|_{\mathcal{C}''_U}$ 由上述规则给出（由《叠》引理
\ref{stacks-lemma-topology-inherited-functorial}，$w$ 连续且余连续，
故可应用《站点》引理 \ref{sites-lemma-when-shriek}）。
\end{proof}








\section{单纯模}
\label{sites-cohomology-section-simplicial-modules}

\noindent
设 $A_\bullet$ 是一个单纯环。回忆我们可以将 $A_\bullet$
看作 $\Delta$（赋予混沌拓扑）上的层，见
《单纯对象》第 \ref{simplicial-section-simplicial-presheaves} 节。
于是，$A_\bullet$ 上的单纯模 $M_\bullet$ 正是 $\Delta$ 上的
$A_\bullet$-模层。换言之，对每个 $n \geq 0$，都有一个
$A_n$-模 $M_n$；对每个映射 $\varphi : [n] \to [m]$，都有相应的映射
$$
M_\bullet(\varphi) : M_m \longrightarrow M_n
$$
它是 $A_\bullet(\varphi)$-线性的，并且这些映射按通常方式复合。

\medskip\noindent
设 $\mathcal{C}$ 是一个站点。$\mathcal{C}$ 上的{\it 单纯环层}
$\mathcal{A}_\bullet$，是 $\mathcal{C}$ 上环层范畴中的单纯对象。
在这种情况下，指派 $U \mapsto \mathcal{A}_\bullet(U)$ 是一个
单纯环层；事实上，这两个概念等价。对于单纯阿贝尔层、
单纯 Lie 代数层等，也有类似的讨论。

\medskip\noindent
不过，与上面的单纯环情形一样，还可用另一种方式理解单纯层。
具体地，考虑投影
$$
p : \Delta \times \mathcal{C} \longrightarrow \mathcal{C}
$$
它定义了一个纤维化范畴，其中强笛卡儿态射恰为形如
$([n], U) \to ([n], V)$ 的态射。我们在范畴
$\Delta \times \mathcal{C}$ 上赋予从 $\mathcal{C}$ 继承的拓扑
（见《栈》第 \ref{stacks-section-topology} 节）。
$\Delta \times \mathcal{C}$ 中覆盖族的简单描述
（《栈》引理 \ref{stacks-lemma-topology-inherited}）立即表明：
$\mathcal{C}$ 上的单纯环层与 $\Delta \times \mathcal{C}$ 上的环层是同一回事。

\medskip\noindent
类比单纯环上的单纯模，我们如下定义单纯环层上的单纯模。

\begin{definition}
\label{sites-cohomology-definition-simplicial-module}
设 $\mathcal{C}$ 是一个站点，$\mathcal{A}_\bullet$ 是
$\mathcal{C}$ 上的单纯环层。一个
{\it 单纯 $\mathcal{A}_\bullet$-模} $\mathcal{F}_\bullet$
（有时称为{\it 单纯 $\mathcal{A}_\bullet$-模层}）是
$\Delta \times \mathcal{C}$ 上与 $\mathcal{A}_\bullet$ 相联系的环层之上的模层。
\end{definition}

\noindent
由此得到单纯模范畴 $\textit{Mod}(\mathcal{A}_\bullet)$ 以及相应的导出范畴
$D(\mathcal{A}_\bullet)$。给定单纯环层的映射
$\mathcal{A}_\bullet \to \mathcal{B}_\bullet$，可得函子
$$
- \otimes^\mathbf{L}_{\mathcal{A}_\bullet} \mathcal{B}_\bullet :
D(\mathcal{A}_\bullet)
\longrightarrow
D(\mathcal{B}_\bullet)
$$
此外，前几节的内容给出函子
$$
L\pi_! : D(\mathcal{A}_\bullet) \longrightarrow D(\mathcal{C})
$$
给定单纯模 $\mathcal{F}_\bullet$，对象
$L\pi_!(\mathcal{F}_\bullet)$ 由相联系的链复形
$s(\mathcal{F}_\bullet)$ 表示
（《单纯对象》第 \ref{simplicial-section-complexes} 节）。
这由引理 \ref{sites-cohomology-lemma-compute-left-derived-pi-shriek} 和
\ref{sites-cohomology-lemma-compute-by-cosimplicial-resolution} 得出。

\begin{lemma}
\label{sites-cohomology-lemma-base-change-by-qis}
设 $\mathcal{C}$ 是一个站点，并设
$\mathcal{A}_\bullet \to \mathcal{B}_\bullet$ 是 $\mathcal{C}$ 上单纯环层的同态。
如果 $L\pi_!\mathcal{A}_\bullet \to L\pi_!\mathcal{B}_\bullet$ 是
$D(\mathcal{C})$ 中的同构，那么
$$
L\pi_!(K) =
L\pi_!(K \otimes^\mathbf{L}_{\mathcal{A}_\bullet} \mathcal{B}_\bullet)
$$
对 $D(\mathcal{A}_\bullet)$ 中每个 $K$ 都成立。
\end{lemma}

\begin{proof}
设 $([n], U)$ 是 $\Delta \times \mathcal{C}$ 的一个对象。由于 $L\pi_!$
与余极限交换，只需对 $\mathcal{O}$-模的上有界复形证明此结论
（可与《导出范畴》命题 \ref{derived-proposition-left-derived-exists}
的论证比较，或者径直限于上有界复形）。每个这样的复形都拟同构于
一个各项均为平坦模的上有界复形，见《站点上的模》
引理 \ref{sites-modules-lemma-module-quotient-flat}。
因此，只需对平坦 $\mathcal{A}_\bullet$-模 $\mathcal{F}$ 证明该引理。
此时导出张量积就是通常的张量积，并且也是一个层。于是由
引理 \ref{sites-cohomology-lemma-compute-left-derived-pi-shriek}，可按
引理 \ref{sites-cohomology-lemma-compute-left-derived-pi-shriek-pre} 的步骤计算等式两边的同调层。
故只需对 $\mathcal{F}$ 在各纤维范畴（即 $\Delta \times U$）上的限制证明结论。
在此情形下，结论由引理 \ref{sites-cohomology-lemma-O-homology-qis} 得出。
\end{proof}

\begin{remark}
\label{sites-cohomology-remark-homology-augmentation}
设 $\mathcal{C}$ 是一个站点。设
$\epsilon : \mathcal{A}_\bullet \to \mathcal{O}$ 是环层范畴中的一个增广
（《单纯对象》定义 \ref{simplicial-definition-augmentation}）。
假设 $\epsilon$ 诱导拟同构
$s(\mathcal{A}_\bullet) \to \mathcal{O}$。
在此情形下，我们得到三角范畴间的正合函子
$$
L\pi_! : D(\mathcal{A}_\bullet) \longrightarrow D(\mathcal{O})
$$
事实上，对 $D(\mathcal{A}_\bullet)$ 的任意对象 $K$，由
引理 \ref{sites-cohomology-lemma-base-change-by-qis} 有
$L\pi_!(K) = L\pi_!(K \otimes_{\mathcal{A}_\bullet}^\mathbf{L} \mathcal{O})$
。因此，可以把所显示的函子定义为
$- \otimes^\mathbf{L}_{\mathcal{A}_\bullet} \mathcal{O}$ 与
注记 \ref{sites-cohomology-remark-fibred-category} 中函子
$L\pi_! : D(\Delta \times \mathcal{C}, \pi^{-1}\mathcal{O}) \to
D(\mathcal{O})$ 的复合。换言之，由该引理给出的
$L\pi_!(K)$ 与
$L\pi_!(K \otimes_{\mathcal{A}_\bullet}^\mathbf{L} \mathcal{O})$
之间规范且函子性的等同，我们在 $L\pi_!(K)$ 上得到一个
$\mathcal{O}$-模结构。
\end{remark}






\section{范畴上的上同调}
\label{sites-cohomology-section-cohomology}

\noindent
在例 \ref{sites-cohomology-example-category-to-point} 的情形下，除导出函子
$L\pi_!$ 外，我们还有函子 $R\pi_*$。对 $\mathcal{C}$ 上的阿贝尔层
$\mathcal{F}$，有
$H_n(\mathcal{C}, \mathcal{F}) = H^{-n}(L\pi_!\mathcal{F})$
和 $H^n(\mathcal{C}, \mathcal{F}) = H^n(R\pi_*\mathcal{F})$。

\begin{example}[计算上同调]
\label{sites-cohomology-example-right-derived-limits}
在例 \ref{sites-cohomology-example-category-to-point} 中，可如下计算函子
$H^n(\mathcal{C}, -)$。设
$\mathcal{F} \in \Ob(\textit{Ab}(\mathcal{C}))$，考虑余链复形
$$
K^\bullet(\mathcal{F}) :
\prod\nolimits_{U_0} \mathcal{F}(U_0)
\to
\prod\nolimits_{U_0 \to U_1} \mathcal{F}(U_0)
\to
\prod\nolimits_{U_0 \to U_1 \to U_2} \mathcal{F}(U_0)
\to \ldots
$$
其中转移映射为
$$
(s_{U_0 \to U_1})
\longmapsto
((U_0 \to U_1 \to U_2) \mapsto s_{U_0 \to U_1} - s_{U_0 \to U_2}
+ s_{U_1 \to U_2}|_{U_0})
$$
，其他次数中的情形类似。由构造，
$$
H^0(\mathcal{C}, \mathcal{F}) =
\lim_{\mathcal{C}^{opp}} \mathcal{F} =
H^0(K^\bullet(\mathcal{F})),
$$
见《范畴》引理 \ref{categories-lemma-limits-products-equalizers}。
$K^\bullet(\mathcal{F})$ 的构造关于 $\mathcal{F}$ 是函子性的，
并把 $\textit{Ab}(\mathcal{C})$ 的短正合列变为复形的短正合列。
因此，函子列
$\mathcal{F} \mapsto H^n(K^\bullet(\mathcal{F}))$ 构成一个 $\delta$-函子，见
《同调》定义 \ref{homology-definition-cohomological-delta-functor} 和
引理 \ref{homology-lemma-long-exact-sequence-cochain}。
对 $\mathcal{C}$ 的对象 $U$，以
$p_U : \Sh(*) \to \Sh(\mathcal{C})$ 表示相应的点，其中
$p_U^{-1}$ 等于在 $U$ 处取值，见
《站点》例 \ref{sites-example-indiscrete-points}。
设 $A$ 是一个阿贝尔群，并令 $\mathcal{F} = p_{U, *}A$。
此时，复形 $K^\bullet(\mathcal{F})$ 的各项为
$\text{Map}(X_n, A)$，其中
$$
X_n = \coprod\nolimits_{U_0 \to \ldots \to U_{n - 1} \to U_n}
\Mor_\mathcal{C}(U, U_0)
$$
这个单纯集同伦等价于单点集 $\{*\}$ 上的常单纯集。事实上，
映射 $X_\bullet \to \{*\}$ 是显然的；映射
$\{*\} \to X_n$ 把 $*$ 映到
$(U \to \ldots \to U, \text{id}_U)$；而定义两个映射
$X_\bullet \to X_\bullet$ 之间同伦的映射
$$
h_{n, i} : X_n \longrightarrow X_n
$$
（《单纯对象》引理 \ref{simplicial-lemma-relations-homotopy}）
由下述规则给出：
$$
h_{n, i} :
(U_0 \to \ldots \to U_n, f)
\longmapsto
(U \to \ldots \to U \to U_i \to \ldots \to U_n, \text{id})
$$
这里 $i > 0$，且 $h_{n, 0} = \text{id}$。验证从略。
由于 $\text{Map}(-, A)$ 是反变函子，这表明
$K^\bullet(p_{U, *}A)$ 在正次数中的上同调为零
（这使用《单纯对象》注记 \ref{simplicial-remark-homotopy-better}
的函子性以及《单纯对象》引理
\ref{simplicial-lemma-homotopy-s-Q} 的结论）。
由此还可知：若 $\mathcal{F}$ 是形如
$p_{U, *}A$ 的层的乘积，则 $K^\bullet(\mathcal{F})$ 在正次数中无环。
由于 $\mathcal{C}$ 上每个阿贝尔层都嵌入这样的乘积，故例如由
《同调》引理 \ref{homology-lemma-efface-implies-universal} 和
《导出范畴》引理 \ref{derived-lemma-right-derived-delta-functor}，
可知 $K^\bullet(\mathcal{F})$ 计算 $H^0(\mathcal{C}, -)$ 的左导出函子
$H^n(\mathcal{C}, -)$。
\end{example}

\begin{example}[计算 Ext]
\label{sites-cohomology-example-computing-exts}
在例 \ref{sites-cohomology-example-category-to-point} 中，进一步假设给定了
$\mathcal{C}$ 上的环层 $\mathcal{O}$。设
$\mathcal{F}$、$\mathcal{G}$ 是 $\mathcal{O}$-模。
考虑复形 $K^\bullet(\mathcal{G}, \mathcal{F})$，其 $n$ 次项为
$$
\prod\nolimits_{U_0 \to U_1 \to \ldots \to U_n}
\Hom_{\mathcal{O}(U_n)}(\mathcal{G}(U_n), \mathcal{F}(U_0))
$$
，转移映射为
$$
(\varphi_{U_0 \to U_1})
\longmapsto
((U_0 \to U_1 \to U_2) \mapsto
\varphi_{U_0 \to U_1} \circ \rho^{U_2}_{U_1}
- \varphi_{U_0 \to U_2}
+ \rho^{U_1}_{U_0} \circ \varphi_{U_1 \to U_2}
$$
，其他次数中的情形类似。这里的 $\rho$ 表示限制映射。由构造，
$$
\Hom_\mathcal{O}(\mathcal{G}, \mathcal{F}) =
H^0(K^\bullet(\mathcal{G}, \mathcal{F}))
$$
对任意一对 $\mathcal{O}$-模 $\mathcal{F}, \mathcal{G}$ 都成立。指派
$(\mathcal{G}, \mathcal{F}) \mapsto K^\bullet(\mathcal{G}, \mathcal{F})$
是一个双函子；它把第一变元中的直和变为乘积，并与第二变元中的乘积交换。
我们断言
$$
\Ext^i_\mathcal{O}(\mathcal{G}, \mathcal{F}) =
H^i(K^\bullet(\mathcal{G}, \mathcal{F}))
$$
对 $i \geq 0$ 成立，只要下列条件之一成立：
\begin{enumerate}
\item 对每个 $U \in \Ob(\mathcal{C})$，
$\mathcal{G}(U)$ 都是投射 $\mathcal{O}(U)$-模；或
\item 对每个 $U \in \Ob(\mathcal{C})$，
$\mathcal{F}(U)$ 都是内射 $\mathcal{O}(U)$-模。
\end{enumerate}
在情形 (1) 中，函子 $K^\bullet(\mathcal{G}, -)$ 是从
$\mathcal{O}$-模范畴到阿贝尔群余链复形范畴的正合函子。
因此，仿照例 \ref{sites-cohomology-example-right-derived-limits} 的论证，
只需证明：当 $\mathcal{F}=p_{U, *}A$，其中 $A$ 是一个
$\mathcal{O}(U)$-模时，$K^\bullet(\mathcal{G}, \mathcal{F})$
在正次数中无环。选取短正合列
\begin{equation}
\label{sites-cohomology-equation-split}
0 \to \mathcal{G}' \to \bigoplus j_{U_i!}\mathcal{O}_{U_i} \to
\mathcal{G} \to 0
\end{equation}
，见《站点上的模》引理
\ref{sites-modules-lemma-module-quotient-flat}。
由于 (1) 对中间的层和右端的层成立，它也对 $\mathcal{G}'$ 成立；
而在 $\mathcal{C}$ 的任一对象上对
(\ref{sites-cohomology-equation-split}) 取值，得到模的分裂正合列。
于是，对任意 $\mathcal{F}$，特别是对
$\mathcal{F} = p_{U, *}A$，得到复形的短正合列
$$
0 \to
K^\bullet(\mathcal{G}, \mathcal{F}) \to
\prod K^\bullet(j_{U_i!}\mathcal{O}_{U_i}, \mathcal{F}) \to
K^\bullet(\mathcal{G}', \mathcal{F}) \to 0
$$
。在 $H^0$ 上得到
$$
0 \to \Hom(\mathcal{G}, p_{U, *}A) \to
\Hom(\prod j_{U_i!}\mathcal{O}_{U_i}, p_{U, *}A) \to
\Hom(\mathcal{G}', p_{U, *}A) \to 0
$$
，它是正合的，因为
$\Hom(\mathcal{H}, p_{U, *}A) = \Hom_{\mathcal{O}(U)}(\mathcal{H}(U), A)$
，并且 (\ref{sites-cohomology-equation-split}) 在 $U$ 上的截面列分裂正合。
因此，利用维数转移可知，只需证明对所有
$U, U' \in \Ob(\mathcal{C})$，
$K^\bullet(j_{U'!}\mathcal{O}_{U'}, p_{U, *}A)$ 在正次数中无环。
在此情形下，
$K^n(j_{U'!}\mathcal{O}_{U'}, p_{U, *}A)$ 等于
$$
\prod\nolimits_{U \to U_0 \to U_1 \to \ldots \to U_n \to U'} A
$$
换言之，$K^\bullet(j_{U'!}\mathcal{O}_{U'}, p_{U, *}A)$
是各项为 $\text{Map}(X_\bullet, A)$ 的复形，其中
$$
X_n = \coprod\nolimits_{U_0 \to \ldots \to U_{n - 1} \to U_n}
\Mor_\mathcal{C}(U, U_0) \times \Mor_\mathcal{C}(U_n, U')
$$
这个单纯集同伦等价于单点集 $\{*\}$ 上的常单纯集；其证明与
例 \ref{sites-cohomology-example-right-derived-limits} 中相应陈述的证明完全相同。
这就完成了该断言的证明。

\medskip\noindent
情形 (2) 中的论证类似（但与之对偶）。
\end{example}






\section{范畴上的模}
\label{sites-cohomology-section-modules-cohomology}

\noindent
本节内容将用于定义概形范畴上的群胚栈之拟凝聚模导出范畴的一种变体。
见《栈上的层》第 \ref{stacks-sheaves-section-QC} 节。

\medskip\noindent
设 $\mathcal{C}$ 是一个范畴。我们把 $\mathcal{C}$ 看作赋予混沌拓扑的站点。
如例 \ref{sites-cohomology-example-computing-exts}，设 $\mathcal{O}$ 是
$\mathcal{C}$ 上的环层。换言之，$\mathcal{O}$ 是范畴
$\mathcal{C}$ 上的环预层，见《范畴》定义
\ref{categories-definition-presheaf}。

\begin{definition}
\label{sites-cohomology-definition-cartesian}
在上述情形下，以 $\mathit{QC}(\mathcal{C}, \mathcal{O})$，或简记为
{\it $\mathit{QC}(\mathcal{O})$}，表示
$D(\mathcal{O}) = D(\mathcal{C}, \mathcal{O})$ 的如下全子范畴：
其对象 $K$ 满足，对 $\mathcal{C}$ 中每个 $U \to V$，规范映射
$$
R\Gamma(V, K) \otimes_{\mathcal{O}(V)}^\mathbf{L} \mathcal{O}(U)
\longrightarrow
R\Gamma(U, K)
$$
是 $D(\mathcal{O}(U))$ 中的同构。
\end{definition}

\begin{lemma}
\label{sites-cohomology-lemma-cartesian-good-sub}
在上述情形下，子范畴 $\mathit{QC}(\mathcal{O})$ 是
$D(\mathcal{O})$ 的严格全、饱和三角子范畴，并且在任意直和下保持不变。
\end{lemma}

\begin{proof}
设 $U$ 是 $\mathcal{C}$ 的一个对象。由于 $\mathcal{C}$ 上的拓扑是混沌拓扑，
函子 $\mathcal{F} \mapsto \mathcal{F}(U)$ 正合且与直和交换。
因此，要计算正合函子 $K \mapsto R\Gamma(U, K)$，可用任意
$\mathcal{O}$-模复形 $\mathcal{F}^\bullet$ 表示 $K$，再取
$\mathcal{F}^\bullet(U)$。故 $R\Gamma(U, -)$ 与直和交换，见
《内射对象》引理 \ref{injectives-lemma-derived-products}。
类似地，给定 $\mathcal{C}$ 的态射 $U \to V$，导出张量积函子
$- \otimes_{\mathcal{O}(V)}^\mathbf{L} \mathcal{O}(U) :
D(\mathcal{O}(V)) \to D(\mathcal{O}(U))$
正合且与直和交换。由这些观察可直接得到该引理；细节从略。
\end{proof}

\begin{lemma}
\label{sites-cohomology-lemma-QC-bounded-above}
在上述情形下，假设 $M$ 是 $\mathit{QC}(\mathcal{O})$ 的一个对象，
且 $b \in \mathbf{Z}$ 满足对所有 $i > b$ 都有 $H^i(M) = 0$。
那么，$H^b(M)$ 是 $(\mathcal{C}, \mathcal{O})$ 上的拟凝聚模，
这里采用《站点上的模》定义 \ref{sites-modules-definition-site-local} 的含义。
\end{lemma}

\begin{proof}
由《站点上的模》引理
\ref{sites-modules-lemma-chaotic-quasi-coherent}，只需证明对
$\mathcal{C}$ 的每个态射 $U \to V$，映射
$$
H^p(M)(V) \otimes_{\mathcal{O}(V)} \mathcal{O}(U) \to H^b(M)(U)
$$
是同构。已知映射
$$
R\Gamma(V, M) \otimes_{\mathcal{O}(V)}^\mathbf{L} \mathcal{O}(U)
\to R\Gamma(U, M)
$$
是同构。于是，例如由 Tor 谱序列可得结论。细节从略。
\end{proof}

\begin{lemma}
\label{sites-cohomology-lemma-QC-final-object}
在上述情形下，假设 $\mathcal{C}$ 有终对象 $X$。
令 $R = \mathcal{O}(X)$，并以
$f : (\mathcal{C}, \mathcal{O}) \to (pt, R)$ 表示显然的站点态射。
那么，由 $Lf^*$ 和 $Rf_*$ 给出 $\mathit{QC}(\mathcal{O}) = D(R)$。
\end{lemma}

\begin{proof}
从略。
\end{proof}

\begin{lemma}
\label{sites-cohomology-lemma-QC-hom-out-of}
在上述情形下，假设 $K$ 是 $\mathit{QC}(\mathcal{O})$ 的一个对象，
而 $M$ 是 $D(\mathcal{O})$ 中的任意对象。对 $\mathcal{C}$ 的每个对象
$U$，都有
$$
\Hom_{D(\mathcal{O}_U)}(K|_U, M|_U) =
R\Hom_{\mathcal{O}(U)}(R\Gamma(U, K), R\Gamma(U, M))
$$
\end{lemma}

\begin{proof}
可用 $\mathcal{C}/U$ 代替 $\mathcal{C}$。因而可设
$U = X$ 是 $\mathcal{C}$ 的终对象。由
引理 \ref{sites-cohomology-lemma-QC-final-object}，可知 $K = Lf^*P$，其中
$P = R\Gamma(U, K) = R\Gamma(X, K) = Rf_*K$。
由于 $Lf^*$ 是 $Rf_*(-) = R\Gamma(U, -)$ 的左伴随，结论随即成立。
\end{proof}

\noindent
设 $(\mathcal{C}, \mathcal{O})$ 如上。对 $\mathcal{O}$-模复形
$\mathcal{F}^\bullet$，把 $\mathcal{F}^\bullet$ 的{\it 大小}
$|\mathcal{F}^\bullet|$ 定义为
$$
|\mathcal{F}^\bullet| =
\left|
\coprod\nolimits_{i \in \mathbf{Z},\ U \in \Ob(\mathcal{C})} \mathcal{F}^i(U)
\right|
$$
对 $D(\mathcal{O})$ 的对象 $K$，把 $K$ 的{\it 大小} $|K|$ 定义为基数
$$
|K| = \min
\left\{
\left|
\mathcal{F}^\bullet
\right|
\text{ 其中 }\mathcal{F}^\bullet\text{ 表示 }K
\right\}
$$
由基数的性质，这个最小值存在。

\begin{lemma}
\label{sites-cohomology-lemma-build-from}
在上述情形下，存在一个基数 $\kappa$，具有如下性质：
给定 $\mathcal{O}$-模复形 $\mathcal{F}^\bullet$ 以及子集
$\Omega^i_U \subset \mathcal{F}^i(U)$，存在子复形
$\mathcal{H}^\bullet \subset \mathcal{F}^\bullet$，使得
$\Omega^i_U \subset \mathcal{H}^i(U)$，并且
$|\mathcal{H}^\bullet| \leq \max(\kappa, |\bigcup \Omega^i_U|)$。
\end{lemma}

\begin{proof}
把 $\mathcal{H}^i(U)$ 定义为 $\mathcal{F}^i(U)$ 的如下
$\mathcal{O}(U)$-子模：它由 $\Omega^i_V$ 在沿任意态射
$f : U \to V$ 限制所得的像以及 $\text{d}(\Omega^{i - 1}_U)$ 生成。
$\mathcal{H}^i(U)$ 的基数不超过下列各基数的最大值：
$\aleph_0$、$\mathcal{O}(U)$ 的基数、
$\text{Arrows}(\mathcal{C})$ 的基数以及
$|\bigcup \Omega^i_U|$。细节从略。
\end{proof}

\begin{lemma}
\label{sites-cohomology-lemma-qis-small}
在上述情形下，存在一个基数 $\kappa$，具有如下性质：
给定表示 $D(\mathcal{O})$ 的对象 $K$ 的 $\mathcal{O}$-模复形
$\mathcal{F}^\bullet$，存在子复形
$\mathcal{H}^\bullet \subset \mathcal{F}^\bullet$，使得
$\mathcal{H}^\bullet$ 表示 $K$，且
$|\mathcal{H}^\bullet| \leq \max(\kappa, |K|)$。
\end{lemma}

\begin{proof}
首先，对每个 $i$ 和 $U$，选取子集
$\Omega^i_U \subset \Ker(\text{d} :
\mathcal{F}^i(U) \to \mathcal{F}^{i + 1}(U))$
，使它双射到
$H^i(K)(U) = H^i(\mathcal{F}^\bullet(U))$ 上。
由于可用大小为 $|K|$ 的复形表示 $K$，故
$|\Omega^i_U| \leq |K|$。应用引理 \ref{sites-cohomology-lemma-build-from}，
得到包含 $\Omega^i_U$ 的子复形
$\mathcal{S}^\bullet \subset \mathcal{F}^\bullet$，其大小至多为
$\max(\kappa, |K|)$；因而
$H^i(\mathcal{S}^\bullet) \to H^i(\mathcal{F}^\bullet)$
是层的满射。

\medskip\noindent
我们将归纳构造子复形
$$
\mathcal{S}^\bullet =
\mathcal{S}_0^\bullet \subset
\mathcal{S}_1^\bullet \subset
\mathcal{S}_2^\bullet \subset \ldots \subset \mathcal{F}^\bullet
$$
，使其大小均不超过 $\max(\kappa, |K|)$，并且
$H^i(\mathcal{S}_n^\bullet) \to H^i(\mathcal{F}^\bullet)$ 的核
与 $H^i(\mathcal{S}_n^\bullet) \to
H^i(\mathcal{S}_{n + 1}^\bullet)$ 的核相同。
构造完成后，可取
$\mathcal{H}^\bullet = \bigcup \mathcal{S}_n^\bullet$ 作为所求。

\medskip\noindent
现由 $\mathcal{S}_n^\bullet$ 构造
$\mathcal{S}_{n + 1}^\bullet$。对每个 $U$ 和 $i$，取子集
$\Omega^{i - 1}_U \subset \mathcal{F}^{i - 1}(U)$，使得
$\text{d} : \mathcal{F}^{i - 1}(U) \to \mathcal{F}^i(U)$
把 $\Omega^{i - 1}_U$ 双射到
$$
\mathcal{S}_n^i(U) \cap
\text{Im}(\text{d} : \mathcal{F}^{i - 1}(U)\to \mathcal{F}^i(U))
$$
上。注意 $|\Omega^i_U| \leq |K|$，因为
$\mathcal{S}_n^i(U)$ 具有这个界。随后对下列子集应用
引理 \ref{sites-cohomology-lemma-build-from}：
$$
\mathcal{S}^i(U) \cup \Omega^i_U \subset \mathcal{F}^i(U)
$$
即可得到 $\mathcal{S}_{n + 1}^\bullet$；其余均显然。
\end{proof}

\begin{lemma}
\label{sites-cohomology-lemma-cartesian-colimit-kappa}
在上述情形下，存在一个具有如下性质的基数 $\kappa$：
\begin{enumerate}
\item 对 $\mathit{QC}(\mathcal{O})$ 的每个非零对象 $K$，
存在 $\mathit{QC}(\mathcal{O})$ 中的非零态射 $E \to K$，
使得 $|E| \leq \kappa$；
\item 对 $\mathit{QC}(\mathcal{O})$ 中每个满足
$|E| \leq \kappa$ 的态射
$\alpha : E \to \bigoplus_n K_n$，存在
$\mathit{QC}(\mathcal{O})$ 中的态射 $E_n \to K_n$，其中
$|E_n| \leq \kappa$，使得 $\alpha$ 经由
$\bigoplus E_n \to \bigoplus K_n$ 分解。
\end{enumerate}
\end{lemma}

\begin{proof}
取 $\kappa$ 为下列基数集合的一个上界：
\begin{enumerate}
\item 对所有 $U \in \Ob(\mathcal{C})$ 的
$|\coprod_V j_{U!}\mathcal{O}_U(V)|$；
\item 对 $\mathcal{C}$ 中所有态射 $U \to V$，在《代数进阶》
引理 \ref{more-algebra-lemma-enlarge} 中得到的基数
$\kappa(\mathcal{O}(V) \to \mathcal{O}(U))$；
\item 引理 \ref{sites-cohomology-lemma-qis-small} 中得到的基数。
\end{enumerate}
我们断言：对表示 $\mathit{QC}(\mathcal{O})$ 中一个对象的任意复形
$\mathcal{F}^\bullet$，以及任意满足
$|\mathcal{S}^\bullet| \leq \kappa$ 的子复形
$\mathcal{S}^\bullet \subset \mathcal{F}^\bullet$，
存在 $\mathcal{F}^\bullet$ 的子复形 $\mathcal{H}^\bullet$，
它包含 $\mathcal{S}^\bullet$，表示
$\mathit{QC}(\mathcal{O})$ 中的一个对象，并且
$|\mathcal{H}^\bullet| \leq \kappa$。
下面两段说明该断言蕴含此引理。

\medskip\noindent
如 (1) 所述，设 $K$ 是 $\mathit{QC}(\mathcal{O})$ 的非零对象。
设 $K$ 由 $\mathcal{O}$-模复形 $\mathcal{F}^\bullet$ 表示。
那么对某个 $i$，$H^i(\mathcal{F}^\bullet)$ 非零。因此，
存在 $\mathcal{C}$ 的对象 $U$ 以及满足 $d(s) = 0$ 的截面
$s \in \mathcal{F}^i(U)$，它确定
$H^i(\mathcal{F}^\bullet)$ 在 $U$ 上的一个非零截面。
于是，$s : j_{U!}\mathcal{O}_U[-i] \to \mathcal{F}^\bullet$
的像是满足 $|\mathcal{S}^\bullet| \leq \kappa$ 的子复形
$\mathcal{S}^\bullet \subset \mathcal{F}^\bullet$。
应用该断言，得到 $\mathit{QC}(\mathcal{O})$ 中的非零态射
$\mathcal{H}^\bullet \to \mathcal{F}^\bullet$，其中
$|\mathcal{H}^\bullet| \leq \kappa$。故 (1) 成立。

\medskip\noindent
设 $\alpha : E \to \bigoplus K_n$ 如 (2)。
任取表示 $K_n$ 的复形 $\mathcal{K}_n^\bullet$。
于是 $\bigoplus \mathcal{K}_n^\bullet$ 表示 $\bigoplus K_n$。
由导出范畴的构造，可用复形 $\mathcal{E}^\bullet$ 表示 $E$，
使 $\alpha$ 由复形态射
$a : \mathcal{E}^\bullet \to \bigoplus \mathcal{K}_n^\bullet$
表示。由引理 \ref{sites-cohomology-lemma-qis-small} 以及上面对 $\kappa$ 的选取，
可设 $|\mathcal{E}^\bullet| \leq \kappa$。
应用该断言，得到子复形
$\mathcal{E}_n^\bullet \subset \mathcal{K}_n^\bullet$；
它们表示 $\mathit{QC}(\mathcal{O})$ 中满足
$|E_n| \leq \kappa$ 的对象 $E_n$，并包含
$a_n : \mathcal{E}^\bullet \to \mathcal{K}_n^\bullet$ 的像，正合所需。

\medskip\noindent
下面证明该断言。设 $\mathcal{F}^\bullet$ 是表示
$\mathit{QC}(\mathcal{O})$ 中一个对象的复形，并设
$\mathcal{S}^\bullet \subset \mathcal{F}^\bullet$ 是大小不超过
$\kappa$ 的子复形。我们将归纳构造子复形
$$
\mathcal{S}^\bullet = \mathcal{S}_0^\bullet \subset
\mathcal{S}_1^\bullet \subset
\mathcal{S}_2^\bullet \subset \ldots \subset \mathcal{F}^\bullet
$$
，使其大小均不超过 $\kappa$，且对 $\mathcal{C}$ 的每个态射
$f : U \to V$ 和每个 $i \in \mathbf{Z}$，满足
\begin{enumerate}
\item 映射
$H^i(\mathcal{S}_n^\bullet(V)
\otimes_{\mathcal{O}(V)}^\mathbf{L} \mathcal{O}(U)) \to
H^i(\mathcal{S}_n^\bullet(U))$
的核在
$H^i(\mathcal{S}_{n + 1}^\bullet(V)
\otimes_{\mathcal{O}(V)}^\mathbf{L} \mathcal{O}(U))$ 中的像为零；
\item 映射
$H^i(\mathcal{S}_n^\bullet(U)) \to H^i(\mathcal{S}_{n + 1}^\bullet(U))$
的像包含于下述映射的像：
$H^i(\mathcal{S}_{n + 1}^\bullet(V)
\otimes_{\mathcal{O}(V)}^\mathbf{L} \mathcal{O}(U)) \to
H^i(\mathcal{S}_{n + 1}^\bullet(U))$。
\end{enumerate}
构造完成后，可令
$\mathcal{H}^\bullet = \bigcup \mathcal{S}_n^\bullet$。
事实上，由于导出张量积以及对环上模复形取上同调都与滤过余极限交换，
条件 (1) 和 (2) 合起来保证
$$
\mathcal{H}^\bullet(V) \otimes_{\mathcal{O}(V)}^\mathbf{L} \mathcal{O}(U)
\longrightarrow
\mathcal{H}^\bullet(U)
$$
对所有 $\mathcal{C}$ 中的 $f : U \to V$，在每个次数的上同调上都是同构，
因而是 $D(\mathcal{O}(U))$ 中的同构。
所以 $\mathcal{H}^\bullet$ 如所需地表示
$\mathit{QC}(\mathcal{O})$ 中的一个对象。

\medskip\noindent
现由 $\mathcal{S}_n$ 构造 $\mathcal{S}_{n + 1}$。
对 $\mathcal{C}$ 的每个态射 $f : U \to V$，考虑交换图
$$
\xymatrix{
\mathcal{S}_n^\bullet(V)
\ar[r] \ar[d] &
\mathcal{S}_n^\bullet(U) \ar[d] \\
\mathcal{F}^\bullet(V)
\ar[r] &
\mathcal{F}^\bullet(U)
}
$$
对于环映射 $\mathcal{O}(V) \to \mathcal{O}(U)$，这是
《代数进阶》引理 \ref{more-algebra-lemma-enlarge} 中那样的图；
也就是说，底行诱导同构
$$
\mathcal{F}^\bullet(V) \otimes_{\mathcal{O}(V)}^\mathbf{L} \mathcal{O}(U)
\longrightarrow
\mathcal{F}^\bullet(U)
$$
于 $D(\mathcal{O}(U))$ 中。因此，可以按照《代数进阶》
引理 \ref{more-algebra-lemma-enlarge} 选取子复形
$$
\mathcal{S}_n^\bullet(V) \subset M^\bullet_f \subset \mathcal{F}^\bullet(V)
\quad\text{且}\quad
\mathcal{S}_n^\bullet(U) \subset N^\bullet_f \subset \mathcal{F}^\bullet(U)
$$
，特别有 $|N^i_f|, |M^i_f| \leq \kappa$。接着，利用子集
$$
\mathcal{S}_n^i(U) \amalg \coprod\nolimits_{f : U \to V} N^i_f
\amalg \coprod\nolimits_{g : W \to U} M^i_g
\subset
\mathcal{F}^i(U)
$$
应用引理 \ref{sites-cohomology-lemma-build-from}，得到子复形
$$
\mathcal{S}_n^\bullet \subset
\mathcal{S}_{n + 1}^\bullet \subset
\mathcal{F}^\bullet
$$
，它包含这些子集，并且
$|\mathcal{S}_{n + 1}^\bullet| \leq \kappa$。
条件 (1) 和 (2) 成立，因为由《代数进阶》引理
\ref{more-algebra-lemma-enlarge} 中的构造，相应陈述对
$\mathcal{S}_n^\bullet(V) \subset M^\bullet_f$ 和
$\mathcal{S}_n^\bullet(U) \subset N^\bullet_f$ 成立。
证明完毕。
\end{proof}

\begin{proposition}
\label{sites-cohomology-proposition-cartesian-brown}
设 $\mathcal{C}$ 是一个范畴，并将其视为赋予混沌拓扑的站点。
设 $\mathcal{O}$ 是 $\mathcal{C}$ 上的环层。令
$\mathit{QC}(\mathcal{O})$ 如定义 \ref{sites-cohomology-definition-cartesian}，则
\begin{enumerate}
\item $\mathit{QC}(\mathcal{O})$ 是 $D(\mathcal{O})$ 的严格全、饱和三角子范畴，
并且在任意直和下保持不变；
\item 任何把直和变为乘积的反变上同调函子
$H : \mathit{QC}(\mathcal{O}) \to \textit{Ab}$
都是可表的；
\item 任何把直和变为直和的三角范畴间正合函子
$F : \mathit{QC}(\mathcal{O}) \to \mathcal{D}$
都有正合右伴随；
\item 包含函子 $\mathit{QC}(\mathcal{O}) \to D(\mathcal{O})$
有正合右伴随。
\end{enumerate}
\end{proposition}

\begin{proof}
(1) 是引理 \ref{sites-cohomology-lemma-cartesian-good-sub}。
(2) 由引理 \ref{sites-cohomology-lemma-cartesian-colimit-kappa} 和
《导出范畴》引理 \ref{derived-lemma-brown-bis} 得出。
(3) 由引理 \ref{sites-cohomology-lemma-cartesian-colimit-kappa} 和
《导出范畴》命题 \ref{derived-proposition-brown-bis} 得出。
(4) 是 (3) 的特例。
\end{proof}

\noindent
设 $u : \mathcal{C}' \to \mathcal{C}$ 是范畴间的函子。
若把 $\mathcal{C}$ 和 $\mathcal{C}'$ 视为赋予混沌拓扑的站点，
则 $u$ 是连续且余连续的函子。因此得到拓扑斯态射
$g : \Sh(\mathcal{C}') \to \Sh(\mathcal{C})$，见
《站点》引理 \ref{sites-lemma-cocontinuous-morphism-topoi}。
此外，假设给定 $\mathcal{C}$ 上的环层 $\mathcal{O}$、
$\mathcal{C}'$ 上的环层 $\mathcal{O}'$，以及映射
$g^\sharp : g^{-1}\mathcal{O} \to \mathcal{O}'$。
把相应的环化拓扑斯态射简记为
$g : (\Sh(\mathcal{C}'), \mathcal{O}') \to (\Sh(\mathcal{C}), \mathcal{O})$,
见《站点上的模》第 \ref{sites-modules-section-ringed-topoi} 节。

\begin{lemma}
\label{sites-cohomology-lemma-cartesian-functorial}
设
$g : (\Sh(\mathcal{C}'), \mathcal{O}') \to
(\Sh(\mathcal{C}), \mathcal{O})$ 如上。那么，函子
$Lg^* : D(\mathcal{O}) \to D(\mathcal{O}')$
把 $\mathit{QC}(\mathcal{O})$ 映入 $\mathit{QC}(\mathcal{O}')$。
\end{lemma}

\begin{proof}
设 $U' \in \Ob(\mathcal{C}')$，其在 $\mathcal{C}$ 中的像为
$U = u(U')$。以 $pt$ 表示只有一个对象和一个态射的范畴。
以 $(\Sh(pt), \mathcal{O}'(U'))$ 和
$(\Sh(pt), \mathcal{O}(U))$ 表示所指的环化拓扑斯。
当然，我们把这些环化拓扑斯上模的导出范畴分别等同于
$D(\mathcal{O}'(U'))$ 和 $D(\mathcal{O}(U))$。
于是有环化拓扑斯的交换图
$$
\xymatrix{
(\Sh(pt), \mathcal{O}'(U')) \ar[rr]_{U'} \ar[d] & &
(\Sh(\mathcal{C}'), \mathcal{O}') \ar[d]^g \\
(\Sh(pt), \mathcal{O}(U)) \ar[rr]^U & &
(\Sh(\mathcal{C}), \mathcal{O})
}
$$
沿下方水平态射拉回，把 $D(\mathcal{O})$ 中的 $K$ 送到
$R\Gamma(U, K)$。沿左侧竖直箭头拉回，把 $M$ 送到
$M \otimes_{\mathcal{O}(U)}^\mathbf{L} \mathcal{O}'(U')$。
沿图的两个方向复合所得结果相同
（引理 \ref{sites-cohomology-lemma-derived-pullback-composition}），故
$$
R\Gamma(U', Lg^*K) =
R\Gamma(U, K) \otimes_{\mathcal{O}(U)}^\mathbf{L} \mathcal{O}'(U')
$$
最后，设 $f' : U' \to V'$ 是 $\mathcal{C}'$ 中的态射，并以
$f = u(f') : U = u(U') \to V = u(V')$ 表示其在 $\mathcal{C}$ 中的像。
若 $K$ 属于 $\mathit{QC}(\mathcal{O})$，则
\begin{align*}
R\Gamma(V', Lg^*K) \otimes_{\mathcal{O}'(V')}^\mathbf{L} \mathcal{O}'(U')
& =
R\Gamma(V, K) \otimes_{\mathcal{O}(V)}^\mathbf{L} \mathcal{O}'(V')
\otimes_{\mathcal{O}'(V')}^\mathbf{L} \mathcal{O}'(U') \\
& =
R\Gamma(V, K) \otimes_{\mathcal{O}(V)}^\mathbf{L} \mathcal{O}'(U') \\
& =
R\Gamma(V, K) \otimes_{\mathcal{O}(V)}^\mathbf{L} \mathcal{O}(U)
\otimes_{\mathcal{O}(U)}^\mathbf{L} \mathcal{O}'(U') \\
& =
R\Gamma(U, K) \otimes_{\mathcal{O}(U)}^\mathbf{L} \mathcal{O}'(U') \\
& =
R\Gamma(U', Lg^*K)
\end{align*}
，正合所需。这里对 $U'$ 和 $V'$ 都使用了上面的观察。
\end{proof}

\begin{lemma}
\label{sites-cohomology-lemma-cartesian-flat-quasi-coherent}
设 $\mathcal{C}$ 是一个范畴，并将其视为赋予混沌拓扑的站点。
设 $\mathcal{O}$ 是 $\mathcal{C}$ 上的环层。
假设对 $\mathcal{C}$ 中所有 $U \to V$，限制映射
$\mathcal{O}(V) \to \mathcal{O}(U)$ 都是平坦环映射。
那么，$\mathit{QC}(\mathcal{O})$ 等于
$D(\mathcal{O})$ 中由上同调层均拟凝聚的复形组成的子范畴
$D_\QCoh(\mathcal{O}) \subset D(\mathcal{O})$。
\end{lemma}

\begin{proof}
回忆在我们的假设下，
$\QCoh(\mathcal{O}) \subset \textit{Mod}(\mathcal{O})$
是弱 Serre 子范畴，见《站点上的模》引理
\ref{sites-modules-lemma-chaotic-flat-quasi-coherent}。
因此，可取 $D(\mathcal{O})$ 的全子范畴
$$
D_\QCoh(\mathcal{O}) = D_{\QCoh(\mathcal{O})}(\textit{Mod}(\mathcal{O}))
$$
，见《导出范畴》第 \ref{derived-section-triangulated-sub} 节。
（严格说来，在该引理的证明中并不需要这一点。）

\medskip\noindent
设 $M$ 是 $\mathit{QC}(\mathcal{O})$ 的一个对象。
由于对 $\mathcal{C}$ 的每个态射 $U \to V$，限制映射
$\mathcal{O}(V) \to \mathcal{O}(U)$ 都是平坦的，故
\begin{align*}
H^i(M)(U)
& =
H^i(R\Gamma(U, M)) \\
& =
H^i(R\Gamma(V, M) \otimes_{\mathcal{O}(V)}^\mathbf{L} \mathcal{O}(U)) \\
& =
H^i(R\Gamma(V, M)) \otimes_{\mathcal{O}(V)} \mathcal{O}(U) \\
& =
H^i(M)(V) \otimes_{\mathcal{O}(V)} \mathcal{O}(U)
\end{align*}
，从而由《站点上的模》引理
\ref{sites-modules-lemma-chaotic-quasi-coherent}，
$H^i(M)$ 是拟凝聚的。上面的第一个和最后一个等式源于如下事实：
由于 $\mathcal{C}$ 上的拓扑是混沌拓扑，在 $\mathcal{C}$ 的一个对象上取截面是正合函子。

\medskip\noindent
反过来，若 $M$ 是 $D_\QCoh(\mathcal{O})$ 的一个对象，则由
《站点上的模》引理
\ref{sites-modules-lemma-chaotic-quasi-coherent}，
映射 $R\Gamma(V, M) \to R\Gamma(U, M)$ 诱导同构
$H^i(M)(U) \to H^i(M)(V) \otimes_{\mathcal{O}(V)} \mathcal{O}(U)$。
于是，由 $\mathcal{O}(V) \to \mathcal{O}(U)$ 的平坦性，
$R\Gamma(V, K) \otimes_{\mathcal{O}(V)}^\mathbf{L} \mathcal{O}(U)
\to R\Gamma(U, K)$ 是 $D(\mathcal{O}(U))$ 中的同构，
故 $M$ 属于 $\mathit{QC}(\mathcal{O})$。
\end{proof}

\begin{lemma}
\label{sites-cohomology-lemma-cartesion-plus-topology}
设 $\epsilon : (\mathcal{C}_\tau, \mathcal{O}_\tau) \to
(\mathcal{C}_{\tau'}, \mathcal{O}_{\tau'})$ 如
第 \ref{sites-cohomology-section-compare} 节。假设
\begin{enumerate}
\item $\tau'$ 是范畴 $\mathcal{C}$ 上的混沌拓扑；
\item 对所有 $U \in \Ob(\mathcal{C})$ 以及所有
$\mathcal{O}(U)$-模的 K-平坦复形 $M^\bullet$，映射
$$
M^\bullet \longrightarrow
R\Gamma((\mathcal{C}/U)_\tau,
(M^\bullet \otimes_{\mathcal{O}(U)} \mathcal{O}_U)^\#)
$$
是拟同构（解释见证明）。
\end{enumerate}
那么，$\epsilon^*$ 和 $R\epsilon_*$ 在
$\mathit{QC}(\mathcal{O})$ 与
$D(\mathcal{C}_\tau, \mathcal{O}_\tau)$ 的如下全子范畴之间给出互为拟逆的等价：
该全子范畴由满足 $R\epsilon_*K$ 属于
$\mathit{QC}(\mathcal{O})$ 的对象 $K$ 组成\footnote{这意味着，对
$\mathcal{C}$ 中所有 $U \to V$，映射
$R\Gamma(V, K) \otimes_{\mathcal{O}(V)}^\mathbf{L} \mathcal{O}(U)
\to R\Gamma(U, K)$ 是同构。}。
\end{lemma}

\begin{proof}
我们将使用第 \ref{sites-cohomology-section-compare} 节中的观察而不再特别说明。
由于 $R\epsilon_*$ 完全忠实且
$\epsilon^* \circ R\epsilon_* = \text{id}$，为证明该引理，
只需证明对 $\mathit{QC}(\mathcal{O})$ 中的 $M$ 有
$R\epsilon_*(\epsilon^*M) = M$。条件 (2) 恰好是得到此结论所需的条件。
具体地，选取一个各项平坦、表示 $M$ 的
$\mathcal{O}$-模 K-平坦复形 $\mathcal{M}^\bullet$。
于是，$\epsilon^*M$ 由 $\mathcal{M}^\bullet$ 的
$\tau$-层化 $(\mathcal{M}^\bullet)^\#$ 表示。
设 $U \in \Ob(\mathcal{C})$。由 Leray 得
$$
R\Gamma(U, R\epsilon_*(\epsilon^*M)) =
R\Gamma((\mathcal{C}/U)_\tau, (\mathcal{M}^\bullet)^\#|_{\mathcal{C}/U}) =
R\Gamma((\mathcal{C}/U)_\tau, (\mathcal{M}^\bullet|_{\mathcal{C}/U})^\#)
$$
最后一个等式成立是因为层化与限制到 $\mathcal{C}/U$ 交换。
照例，以 $\mathcal{O}_U$ 表示 $\mathcal{O}$ 到
$\mathcal{C}/U$ 的限制。考虑 $\mathcal{O}_U$-模复形的映射
$$
\mathcal{M}^\bullet(U) \otimes_{\mathcal{O}(U)} \mathcal{O}_U
\longrightarrow
\mathcal{M}^\bullet|_{\mathcal{C}/U}
$$
（在 $\tau'$-拓扑中）。由我们对 $\mathcal{M}^\bullet$ 的选取，
复形 $\mathcal{M}^\bullet(U)$ 是 $\mathcal{O}(U)$-模的 K-平坦复形；
见引理 \ref{sites-cohomology-lemma-pullback-K-flat}，并利用 $U$ 到
$\mathcal{C}$ 的包含定义环化拓扑斯态射
$(\Sh(pt), \mathcal{O}(U)) \to (\Sh(\mathcal{C}_{\tau'}), \mathcal{O})$.
由于 $M$ 属于 $\mathit{QC}(\mathcal{O})$，可知所显示的箭头是拟同构。
又因层化正合，层化后仍然如此。因此
$$
R\Gamma(U, R\epsilon_*(\epsilon^*M)) =
R\Gamma((\mathcal{C}/U)_\tau,
(M^\bullet \otimes_{\mathcal{O}(U)} \mathcal{O}_U)^\#)
$$
，而假设 (2) 表明它等于
$R\Gamma(U, M) = \mathcal{M}^\bullet(U)$，正合所需。
\end{proof}

\begin{lemma}
\label{sites-cohomology-lemma-QC-hom-out-of-plus-topology}
采用引理 \ref{sites-cohomology-lemma-cartesion-plus-topology} 的记号和假设。
假设 $K$ 是 $\mathit{QC}(\mathcal{O})$ 的一个对象，而
$M$ 是 $D(\mathcal{O}_\tau)$ 中的任意对象。
对 $\mathcal{C}$ 的每个对象 $U$，都有
$$
\Hom_{D((\mathcal{O}_U)_\tau)}(\epsilon^*K|_U, M|_U) =
R\Hom_{\mathcal{O}(U)}(R\Gamma(U, K), R\Gamma(U, M))
$$
\end{lemma}

\begin{proof}
由伴随性，有
$$
\Hom_{D((\mathcal{O}_U)_\tau)}(\epsilon^*K|_U, M|_U) =
\Hom_{D((\mathcal{O}_U)_{\tau'})}(K|_U, R\epsilon_*M|_U)
$$
故由引理 \ref{sites-cohomology-lemma-QC-hom-out-of} 以及下式得到结论：
$$
R\Gamma(U, M) = R\Gamma(U, R\epsilon_*M)
$$
，其中使用了 Leray。
\end{proof}







\section{严格完美复形}
\label{sites-cohomology-section-strictly-perfect}

\noindent
本节是《上同调》第 \ref{cohomology-section-strictly-perfect} 节的对应版本。

\begin{definition}
\label{sites-cohomology-definition-strictly-perfect}
设 $(\mathcal{C}, \mathcal{O})$ 是一个环化站点，
$\mathcal{E}^\bullet$ 是 $\mathcal{O}$-模复形。
如果除有限多个 $i$ 外 $\mathcal{E}^i$ 均为零，并且对每个 $i$，
$\mathcal{E}^i$ 都是某个有限自由 $\mathcal{O}$-模的直和项，
则称 $\mathcal{E}^\bullet$ 是{\it 严格完美的}。
\end{definition}

\noindent
设 $U$ 是 $\mathcal{C}$ 的一个对象。我们常说
“设 $\mathcal{E}^\bullet$ 是 $\mathcal{O}_U$-模的严格完美复形”，
其含义是 $\mathcal{E}^\bullet$ 是环化站点
$(\mathcal{C}/U, \mathcal{O}_U)$ 上模的严格完美复形，见
《站点上的模》定义 \ref{sites-modules-definition-localize-ringed-site}。

\begin{lemma}
\label{sites-cohomology-lemma-cone}
严格完美复形之间态射的锥是严格完美的。
\end{lemma}

\begin{proof}
由定义立即可得。
\end{proof}

\begin{lemma}
\label{sites-cohomology-lemma-tensor}
与两个严格完美复形的张量积相联系的全复形是严格完美的。
\end{lemma}

\begin{proof}
从略。
\end{proof}

\begin{lemma}
\label{sites-cohomology-lemma-strictly-perfect-pullback}
设 $(f, f^\sharp) : (\mathcal{C}, \mathcal{O}_\mathcal{C}) \to
(\mathcal{D}, \mathcal{O}_\mathcal{D})$
是环化拓扑斯的态射。如果 $\mathcal{F}^\bullet$ 是
$\mathcal{O}_\mathcal{D}$-模的严格完美复形，那么
$f^*\mathcal{F}^\bullet$ 是
$\mathcal{O}_\mathcal{C}$-模的严格完美复形。
\end{lemma}

\begin{proof}
我们已在《站点上的模》引理
\ref{sites-modules-lemma-global-pullback} 中看到，有限自由模的拉回仍为有限自由模。
函子 $f^*$ 是加性函子，因而保持直和项。引理得证。
\end{proof}

\begin{lemma}
\label{sites-cohomology-lemma-local-lift-map}
设 $(\mathcal{C}, \mathcal{O})$ 是一个环化站点，$U$ 是
$\mathcal{C}$ 的一个对象。给定 $\mathcal{O}_U$-模的实线图
$$
\xymatrix{
\mathcal{E} \ar@{..>}[dr] \ar[r] & \mathcal{F} \\
& \mathcal{G} \ar[u]_p
}
$$
，其中 $\mathcal{E}$ 是某个有限自由 $\mathcal{O}_U$-模的直和项，
且 $p$ 为满射，则存在覆盖 $\{U_i \to U\}$，使得在每个 $U_i$ 上
都存在令该图交换的虚线箭头。
\end{lemma}

\begin{proof}
可设对某个 $n$ 有
$\mathcal{E} = \mathcal{O}_U^{\oplus n}$。
此时，找到虚线箭头等价于提升 $\Gamma(U, \mathcal{F})$ 中各基元素的像。
由层的满射之刻画，这在局部可以做到
（《站点》第 \ref{sites-section-sheaves-injective} 节）。
\end{proof}

\begin{lemma}
\label{sites-cohomology-lemma-local-homotopy}
设 $(\mathcal{C}, \mathcal{O})$ 是一个环化站点，$U$ 是
$\mathcal{C}$ 的一个对象。
\begin{enumerate}
\item 设 $\alpha : \mathcal{E}^\bullet \to \mathcal{F}^\bullet$
是 $\mathcal{O}_U$-模复形的态射，其中 $\mathcal{E}^\bullet$
严格完美，而 $\mathcal{F}^\bullet$ 无环。那么，存在覆盖
$\{U_i \to U\}$，使得每个 $\alpha|_{U_i}$ 都同伦于零。
\item 设 $\alpha : \mathcal{E}^\bullet \to \mathcal{F}^\bullet$
是 $\mathcal{O}_U$-模复形的态射，其中 $\mathcal{E}^\bullet$
严格完美，对 $i < a$ 有 $\mathcal{E}^i = 0$，并且对
$i \geq a$ 有 $H^i(\mathcal{F}^\bullet) = 0$。
那么，存在覆盖 $\{U_i \to U\}$，使得每个
$\alpha|_{U_i}$ 都同伦于零。
\end{enumerate}
\end{lemma}

\begin{proof}
第一条由第二条得出，故只证明 (2)。
我们对复形 $\mathcal{E}^\bullet$ 的长度作归纳。
如果对某个有限自由 $\mathcal{O}$-模的直和项 $\mathcal{E}$
以及整数 $n \geq a$，有
$\mathcal{E}^\bullet \cong \mathcal{E}[-n]$，则结论由
引理 \ref{sites-cohomology-lemma-local-lift-map} 以及下述事实得出：
$\mathcal{F}^{n - 1} \to \Ker(\mathcal{F}^n \to \mathcal{F}^{n + 1})$
因假设 $H^n(\mathcal{F}^\bullet)$ 为零而是满射。
若除 $i \in [a,b]$ 外 $\mathcal{E}^i$ 均为零，则有复形的分裂正合列
$$
0 \to \mathcal{E}^b[-b] \to \mathcal{E}^\bullet \to
\sigma_{\leq b - 1}\mathcal{E}^\bullet \to 0
$$
，它在 $K(\mathcal{O}_U)$ 中确定一个有区别三角，因而由三角范畴公理得到正合列
$$
\Hom_{K(\mathcal{O}_U)}(
\sigma_{\leq b - 1}\mathcal{E}^\bullet, \mathcal{F}^\bullet)
\to
\Hom_{K(\mathcal{O}_U)}(\mathcal{E}^\bullet, \mathcal{F}^\bullet)
\to
\Hom_{K(\mathcal{O}_U)}(\mathcal{E}^b[-b], \mathcal{F}^\bullet)
$$
。由上面的论证，复合
$\mathcal{E}^b[-b] \to \mathcal{F}^\bullet$
在 $U$ 的某个覆盖的各成员上同伦于零。因此，可设我们的映射来自
上面所显示正合列左端的一个元素。由归纳假设，该元素在
$U$ 的某个覆盖的各成员上为零。
\end{proof}

\begin{lemma}
\label{sites-cohomology-lemma-lift-through-quasi-isomorphism}
设 $(\mathcal{C}, \mathcal{O})$ 是一个环化站点，$U$ 是
$\mathcal{C}$ 的一个对象。给定 $\mathcal{O}_U$-模复形的实线图
$$
\xymatrix{
\mathcal{E}^\bullet \ar@{..>}[dr] \ar[r]_\alpha & \mathcal{F}^\bullet \\
& \mathcal{G}^\bullet \ar[u]_f
}
$$
，其中 $\mathcal{E}^\bullet$ 严格完美，对 $j < a$ 有
$\mathcal{E}^j = 0$，并且 $H^j(f)$ 对 $j > a$ 是同构、
对 $j = a$ 是满射，则存在覆盖 $\{U_i \to U\}$，并且对每个 $i$，
在 $U_i$ 上都存在使该图交换至同伦的虚线箭头。
\end{lemma}

\begin{proof}
关于 $f$ 的假设蕴含锥 $C(f)^\bullet$ 在次数 $\geq a$ 中的上同调层为零。
因此，引理 \ref{sites-cohomology-lemma-local-homotopy} 保证存在覆盖
$\{U_i \to U\}$，使得复合
$\mathcal{E}^\bullet \to \mathcal{F}^\bullet \to C(f)^\bullet$
在 $U_i$ 上同伦于零。由于
$$
\mathcal{G}^\bullet \to \mathcal{F}^\bullet \to C(f)^\bullet \to
\mathcal{G}^\bullet[1]
$$
限制为 $K(\mathcal{O}_{U_i})$ 中的有区别三角，可知
$\alpha|_{U_i}$ 能够提升至同伦为映射
$\alpha_i : \mathcal{E}^\bullet|_{U_i} \to
\mathcal{G}^\bullet|_{U_i}$，正合所需。
\end{proof}

\begin{lemma}
\label{sites-cohomology-lemma-local-actual}
设 $(\mathcal{C}, \mathcal{O})$ 是一个环化站点，$U$ 是
$\mathcal{C}$ 的一个对象。设 $\mathcal{E}^\bullet$、
$\mathcal{F}^\bullet$ 是 $\mathcal{O}_U$-模复形，其中
$\mathcal{E}^\bullet$ 严格完美。
\begin{enumerate}
\item 对任意元素
$\alpha \in \Hom_{D(\mathcal{O}_U)}(\mathcal{E}^\bullet, \mathcal{F}^\bullet)$
，存在覆盖 $\{U_i \to U\}$，使得
$\alpha|_{U_i}$ 由复形态射
$\alpha_i : \mathcal{E}^\bullet|_{U_i} \to
\mathcal{F}^\bullet|_{U_i}$ 给出。
\item 给定复形态射
$\alpha : \mathcal{E}^\bullet \to \mathcal{F}^\bullet$
，若它在群
$\Hom_{D(\mathcal{O}_U)}(\mathcal{E}^\bullet, \mathcal{F}^\bullet)$
中的像为零，则存在覆盖 $\{U_i \to U\}$，使得
$\alpha|_{U_i}$ 同伦于零。
\end{enumerate}
\end{lemma}

\begin{proof}
证明 (1)。由导出范畴的构造，可以找到拟同构
$f : \mathcal{F}^\bullet \to \mathcal{G}^\bullet$ 和复形映射
$\beta : \mathcal{E}^\bullet \to \mathcal{G}^\bullet$，
使得 $\alpha = f^{-1}\beta$。于是结论由
引理 \ref{sites-cohomology-lemma-lift-through-quasi-isomorphism} 得出。
(2) 的证明从略。
\end{proof}

\begin{lemma}
\label{sites-cohomology-lemma-Rhom-strictly-perfect}
设 $(\mathcal{C}, \mathcal{O})$ 是一个环化站点。
设 $\mathcal{E}^\bullet$、$\mathcal{F}^\bullet$ 是
$\mathcal{O}$-模复形，其中 $\mathcal{E}^\bullet$ 严格完美。
那么，内部 Hom
$R\SheafHom(\mathcal{E}^\bullet, \mathcal{F}^\bullet)$
由复形 $\mathcal{H}^\bullet$ 表示，其各项为
$$
\mathcal{H}^n =
\bigoplus\nolimits_{n = p + q}
\SheafHom_\mathcal{O}(\mathcal{E}^{-q}, \mathcal{F}^p)
$$
，微分如第 \ref{sites-cohomology-section-internal-hom} 节所述。
\end{lemma}

\begin{proof}
选取到 K-内射复形的拟同构
$\mathcal{F}^\bullet \to \mathcal{I}^\bullet$。
设 $(\mathcal{H}')^\bullet$ 是各项如下的复形：
$$
(\mathcal{H}')^n =
\prod\nolimits_{n = p + q}
\SheafHom_\mathcal{O}(\mathcal{L}^{-q}, \mathcal{I}^p)
$$
由第 \ref{sites-cohomology-section-internal-hom} 节中的构造，它表示
$R\SheafHom(\mathcal{E}^\bullet, \mathcal{F}^\bullet)$。
只需证明映射
$$
\mathcal{H}^\bullet \longrightarrow (\mathcal{H}')^\bullet
$$
是拟同构。给定 $\mathcal{C}$ 的对象 $U$，直接检验可得
$$
H^0(\mathcal{H}^\bullet(U)) =
\Hom_{K(\mathcal{O}_U)}(\mathcal{E}^\bullet|_U, \mathcal{I}^\bullet|_U)
\to
H^0((\mathcal{H}')^\bullet(U)) =
\Hom_{D(\mathcal{O}_U)}(\mathcal{E}^\bullet|_U, \mathcal{I}^\bullet|_U)
$$
由引理 \ref{sites-cohomology-lemma-local-actual}，
$U \mapsto H^0(\mathcal{H}^\bullet(U))$ 的层化等于
$U \mapsto H^0((\mathcal{H}')^\bullet(U))$ 的层化。
对其他上同调层可作类似论证。因此，
$\mathcal{H}^\bullet$ 与 $(\mathcal{H}')^\bullet$ 拟同构，引理得证。
\end{proof}

\begin{lemma}
\label{sites-cohomology-lemma-Rhom-complex-of-direct-summands-finite-free}
设 $(\mathcal{C}, \mathcal{O})$ 是一个环化站点。
设 $\mathcal{E}^\bullet$、$\mathcal{F}^\bullet$ 是
$\mathcal{O}$-模复形，并且
\begin{enumerate}
\item 对 $n \ll 0$ 有 $\mathcal{F}^n = 0$；
\item 对 $n \gg 0$ 有 $\mathcal{E}^n = 0$；
\item $\mathcal{E}^n$ 同构于某个有限自由 $\mathcal{O}$-模的直和项。
\end{enumerate}
那么，内部 Hom
$R\SheafHom(\mathcal{E}^\bullet, \mathcal{F}^\bullet)$
由复形 $\mathcal{H}^\bullet$ 表示，其各项为
$$
\mathcal{H}^n =
\bigoplus\nolimits_{n = p + q}
\SheafHom_\mathcal{O}(\mathcal{E}^{-q}, \mathcal{F}^p)
$$
，微分如第 \ref{sites-cohomology-section-internal-hom} 节所述。
\end{lemma}

\begin{proof}
选取拟同构 $\mathcal{F}^\bullet \to \mathcal{I}^\bullet$，
其中 $\mathcal{I}^\bullet$ 是内射对象的下有界复形。
注意 $\mathcal{I}^\bullet$ 是 K-内射的
（《导出范畴》引理
\ref{derived-lemma-bounded-below-injectives-K-injective}）。
因此，第 \ref{sites-cohomology-section-internal-hom} 节中的构造表明，
$R\SheafHom(\mathcal{E}^\bullet, \mathcal{F}^\bullet)$
由复形 $(\mathcal{H}')^\bullet$ 表示，其各项为
$$
(\mathcal{H}')^n =
\prod\nolimits_{n = p + q}
\SheafHom_\mathcal{O}(\mathcal{E}^{-q}, \mathcal{I}^p) =
\bigoplus\nolimits_{n = p + q}
\SheafHom_\mathcal{O}(\mathcal{E}^{-q}, \mathcal{I}^p)
$$
（等号成立是因为只有有限多个非零项）。
注意 $\mathcal{H}^\bullet$ 是与各项为
$\SheafHom_\mathcal{O}(\mathcal{E}^{-q}, \mathcal{F}^p)$
的双复形相联系的全复形；$(\mathcal{H}')^\bullet$ 也类似。
自然映射 $(\mathcal{H}')^\bullet \to \mathcal{H}^\bullet$
来自双复形的一个映射。
因此，为证明该映射是拟同构，可以使用双复形的谱序列
（《同调》引理 \ref{homology-lemma-first-quadrant-ss}）
$$
{}'E_1^{p, q} =
H^p(\SheafHom_\mathcal{O}(\mathcal{E}^{-q}, \mathcal{F}^\bullet))
$$
，它收敛到 $H^{p + q}(\mathcal{H}^\bullet)$；
对 $(\mathcal{H}')^\bullet$ 也有类似谱序列。
为完成引理的证明，只需证明：每当 $\mathcal{E}$ 是某个有限自由
$\mathcal{O}$-模的直和项时，
$\mathcal{F}^\bullet \to \mathcal{I}^\bullet$ 在上同调层上诱导同构
$$
H^p(\SheafHom_\mathcal{O}(\mathcal{E}, \mathcal{F}^\bullet))
\longrightarrow
H^p(\SheafHom_\mathcal{O}(\mathcal{E}, \mathcal{I}^\bullet))
$$
。当 $\mathcal{E}$ 有限自由时这是显然的，故一般结论也随之成立。
\end{proof}









\section{伪凝聚模}
\label{sites-cohomology-section-pseudo-coherent}

\noindent
本节讨论伪凝聚复形。

\begin{definition}
\label{sites-cohomology-definition-pseudo-coherent}
设 $(\mathcal{C}, \mathcal{O})$ 是一个环化站点，
$\mathcal{E}^\bullet$ 是 $\mathcal{O}$-模复形，并设
$m \in \mathbf{Z}$。
\begin{enumerate}
\item 如果对 $\mathcal{C}$ 的每个对象 $U$，都存在覆盖
$\{U_i \to U\}$，并且对每个 $i$ 都存在复形态射
$\alpha_i : \mathcal{E}_i^\bullet \to \mathcal{E}^\bullet|_{U_i}$，
其中 $\mathcal{E}_i^\bullet$ 是 $\mathcal{O}_{U_i}$-模的严格完美复形，
$H^j(\alpha_i)$ 对 $j > m$ 是同构，而
$H^m(\alpha_i)$ 是满射，则称 $\mathcal{E}^\bullet$
是{\it $m$-伪凝聚的}。
\item 如果 $\mathcal{E}^\bullet$ 对所有 $m$ 都是 $m$-伪凝聚的，
则称它是{\it 伪凝聚的}。
\item 当且仅当 $D(\mathcal{O})$ 的对象 $E$ 能由一个
$m$-伪凝聚（相应地，伪凝聚）的 $\mathcal{O}$-模复形表示时，
称 $E$ 是{\it $m$-伪凝聚的}（相应地，{\it 伪凝聚的}）。
\end{enumerate}
\end{definition}

\noindent
如果 $\mathcal{C}$ 有拟紧的终对象 $X$
（例如，$X$ 的每个覆盖都可由有限覆盖加细），那么
$D(\mathcal{O})$ 的 $m$-伪凝聚对象属于
$D^-(\mathcal{O})$。但一般而言未必如此。

\begin{lemma}
\label{sites-cohomology-lemma-pseudo-coherent-independent-representative}
设 $(\mathcal{C}, \mathcal{O})$ 是一个环化站点，
$E$ 是 $D(\mathcal{O})$ 的一个对象。
\begin{enumerate}
\item 如果 $\mathcal{C}$ 有终对象 $X$，并且存在覆盖
$\{U_i \to X\}$、$\mathcal{O}_{U_i}$-模的严格完美复形
$\mathcal{E}_i^\bullet$，以及 $D(\mathcal{O}_{U_i})$ 中的映射
$\alpha_i : \mathcal{E}_i^\bullet \to E|_{U_i}$，使得
$H^j(\alpha_i)$ 对 $j > m$ 是同构，且
$H^m(\alpha_i)$ 是满射，那么 $E$ 是 $m$-伪凝聚的。
\item 如果 $E$ 是 $m$-伪凝聚的，那么表示 $E$ 的任意
$\mathcal{O}$-模复形都是 $m$-伪凝聚的。
\item 如果对 $\mathcal{C}$ 的每个对象 $U$，都存在覆盖
$\{U_i \to U\}$，使得 $E|_{U_i}$ 是 $m$-伪凝聚的，
那么 $E$ 是 $m$-伪凝聚的。
\end{enumerate}
\end{lemma}

\begin{proof}
设 $\mathcal{F}^\bullet$ 是表示 $E$ 的任意复形，并设
$X$、$\{U_i \to X\}$ 以及
$\alpha_i : \mathcal{E}_i \to E|_{U_i}$ 如 (1)。
我们将证明复形 $\mathcal{F}^\bullet$ 是 $m$-伪凝聚的；
这将证明 (1)，并在 $\mathcal{C}$ 有终对象时证明 (2)。
由引理 \ref{sites-cohomology-lemma-local-actual}，加细覆盖
$\{U_i \to X\}$ 后，可以用复形映射
$\alpha_i : \mathcal{E}_i^\bullet \to
\mathcal{F}^\bullet|_{U_i}$ 表示这些映射。
根据假设，$H^j(\alpha_i)$ 对 $j > m$ 是同构，而
$H^m(\alpha_i)$ 是满射，故 $\mathcal{F}^\bullet$ 是 $m$-伪凝聚的。

\medskip\noindent
证明 (2)。由上可知，对 $\mathcal{C}$ 的所有对象 $U$，
$\mathcal{F}^\bullet|_U$ 作为 $\mathcal{O}_U$-模复形是
$m$-伪凝聚的。由定义形式地推出
$\mathcal{F}^\bullet$ 是 $m$-伪凝聚的。

\medskip\noindent
证明 (3)。由定义以及《站点》定义
\ref{sites-definition-site} 的 (2) 得出。
\end{proof}

\begin{lemma}
\label{sites-cohomology-lemma-pseudo-coherent-pullback}
设 $(f, f^\sharp) : (\mathcal{C}, \mathcal{O}_\mathcal{C}) \to
(\mathcal{D}, \mathcal{O}_\mathcal{D})$
是环化站点的态射。设 $E$ 是
$D(\mathcal{O}_\mathcal{C})$ 的一个对象。如果 $E$ 是
$m$-伪凝聚的，那么 $Lf^*E$ 是 $m$-伪凝聚的。
\end{lemma}

\begin{proof}
设 $f$ 由函子 $u : \mathcal{D} \to \mathcal{C}$ 给出。
设 $U$ 是 $\mathcal{C}$ 的一个对象。由《站点》引理
\ref{sites-lemma-morphism-of-sites-covering}，
可以找到覆盖 $\{U_i \to U\}$，并且对每个 $i$，都有从
$U_i$ 到某个 $\mathcal{D}$ 中对象 $V_i$ 的像 $u(V_i)$ 的态射。
由引理 \ref{sites-cohomology-lemma-pseudo-coherent-independent-representative}，
只需证明 $Lf^*E|_{U_i}$ 是 $m$-伪凝聚的。
为此，只需证明 $Lf^*E|_{u(V_i)}$ 是 $m$-伪凝聚的，
因为 $Lf^*E|_{U_i}$ 是 $Lf^*E|_{u(V_i)}$ 到
$\mathcal{C}/U_i$ 的限制（使用《站点上的模》引理
\ref{sites-modules-lemma-relocalize}）。
由《站点上的模》引理
\ref{sites-modules-lemma-localize-morphism-ringed-sites} 中的交换图，
只需对环化站点态射
$(\mathcal{C}/u(V_i), \mathcal{O}_{u(V_i)}) \to
(\mathcal{D}/V_i, \mathcal{O}_{V_i})$ 证明该引理。
因此，可设 $\mathcal{D}$ 有终对象 $Y$，且
$X = u(Y)$ 是 $\mathcal{C}$ 的终对象。

\medskip\noindent
设 $\{V_i \to Y\}$ 是一个覆盖，使得对每个 $i$，都存在
$\mathcal{O}_{V_i}$-模的严格完美复形
$\mathcal{F}_i^\bullet$ 以及 $D(\mathcal{O}_{V_i})$ 中的态射
$\alpha_i : \mathcal{F}_i^\bullet \to E|_{V_i}$，并且
$H^j(\alpha_i)$ 对 $j > m$ 是同构，而
$H^m(\alpha_i)$ 是满射。仿照上面的论证，只需对
$(\mathcal{C}/u(V_i), \mathcal{O}_{u(V_i)}) \to
(\mathcal{D}/V_i, \mathcal{O}_{V_i})$ 证明结论。
因此，可设存在 $\mathcal{O}_\mathcal{D}$-模的严格完美复形
$\mathcal{F}^\bullet$ 以及 $D(\mathcal{O}_\mathcal{D})$ 中的态射
$\alpha : \mathcal{F}^\bullet \to E$，使得
$H^j(\alpha)$ 对 $j > m$ 是同构，而
$H^m(\alpha)$ 是满射。在此情形下，选取有区别三角
$$
\mathcal{F}^\bullet \to E \to C \to \mathcal{F}^\bullet[1]
$$
关于 $\alpha$ 的假设恰好意味着对所有 $j \geq m$，
上同调层 $H^j(C)$ 都为零。应用 $Lf^*$，得到有区别三角
$$
Lf^*\mathcal{F}^\bullet \to Lf^*E \to Lf^*C \to Lf^*\mathcal{F}^\bullet[1]
$$
由 $Lf^*$ 作为左导出函子的构造，可知对 $j \geq m$ 有
$H^j(Lf^*C) = 0$（使用《导出范畴》引理
\ref{derived-lemma-negative-vanishing} 的对偶）。
因此，$H^j(Lf^*\alpha)$ 对 $j > m$ 是同构，而
$H^m(Lf^*\alpha)$ 是满射。另一方面，由于
$\mathcal{F}^\bullet$ 是平坦
$\mathcal{O}_\mathcal{D}$-模的上有界复形，有
$Lf^*\mathcal{F}^\bullet = f^*\mathcal{F}^\bullet$。
应用引理 \ref{sites-cohomology-lemma-strictly-perfect-pullback} 即得结论。
\end{proof}

\begin{lemma}
\label{sites-cohomology-lemma-cone-pseudo-coherent}
设 $(\mathcal{C}, \mathcal{O})$ 是一个环化站点，
$m \in \mathbf{Z}$。设
$(K, L, M, f, g, h)$ 是 $D(\mathcal{O})$ 中的有区别三角。
\begin{enumerate}
\item 如果 $K$ 是 $(m + 1)$-伪凝聚的，而 $L$ 是 $m$-伪凝聚的，
那么 $M$ 是 $m$-伪凝聚的。
\item 如果 $K$ 和 $M$ 都是 $m$-伪凝聚的，那么 $L$ 是 $m$-伪凝聚的。
\item 如果 $L$ 是 $(m + 1)$-伪凝聚的，而 $M$ 是 $m$-伪凝聚的，
那么 $K$ 是 $(m + 1)$-伪凝聚的。
\end{enumerate}
\end{lemma}

\begin{proof}
证明 (1)。设 $U$ 是 $\mathcal{C}$ 的一个对象。选取覆盖
$\{U_i \to U\}$ 以及 $D(\mathcal{O}_{U_i})$ 中的映射
$\alpha_i : \mathcal{K}_i^\bullet \to K|_{U_i}$，其中
$\mathcal{K}_i^\bullet$ 严格完美，$H^j(\alpha_i)$ 对
$j > m + 1$ 是同构，并且对 $j = m + 1$ 是满射。
可用 $\sigma_{\geq m + 1}\mathcal{K}_i^\bullet$ 代替
$\mathcal{K}_i^\bullet$，从而可设对 $j < m + 1$ 有
$\mathcal{K}_i^j = 0$。加细覆盖后，可以选取
$D(\mathcal{O}_{U_i})$ 中的映射
$\beta_i : \mathcal{L}_i^\bullet \to L|_{U_i}$，
其中 $\mathcal{L}_i^\bullet$ 严格完美，并且
$H^j(\beta)$ 对 $j > m$ 是同构、对 $j = m$ 是满射。
由引理 \ref{sites-cohomology-lemma-lift-through-quasi-isomorphism}，
再次加细覆盖后，可找到复形映射
$\gamma_i : \mathcal{K}^\bullet \to \mathcal{L}^\bullet$，
使得图
$$
\xymatrix{
K|_{U_i} \ar[r] & L|_{U_i} \\
\mathcal{K}_i^\bullet \ar[u]^{\alpha_i} \ar[r]^{\gamma_i} &
\mathcal{L}_i^\bullet \ar[u]_{\beta_i}
}
$$
在 $D(\mathcal{O}_{U_i})$ 中交换（这要求用实际的复形映射表示
$\alpha_i$、$\beta_i$ 以及 $K|_{U_i} \to L|_{U_i}$；
一些细节从略）。锥 $C(\gamma_i)^\bullet$ 严格完美
（引理 \ref{sites-cohomology-lemma-cone}）。
该图的交换性蕴含存在有区别三角的态射
$$
(\mathcal{K}_i^\bullet, \mathcal{L}_i^\bullet, C(\gamma_i)^\bullet)
\longrightarrow
(K|_{U_i}, L|_{U_i}, M|_{U_i}).
$$
由上同调长正合列上诱导的映射以及《同调》引理
\ref{homology-lemma-four-lemma} 和
\ref{homology-lemma-five-lemma}，可知
$C(\gamma_i)^\bullet \to M|_{U_i}$ 在次数 $>m$ 的上同调上诱导同构，
在次数 $m$ 上诱导满射。因此，由引理
\ref{sites-cohomology-lemma-pseudo-coherent-independent-representative}，
$M$ 是 $m$-伪凝聚的。

\medskip\noindent
旋转有区别三角，即由 (1) 得到断言 (2) 和 (3)。
\end{proof}

\begin{lemma}
\label{sites-cohomology-lemma-tensor-pseudo-coherent}
设 $(\mathcal{C}, \mathcal{O})$ 是一个环化站点，
$K,L$ 是 $D(\mathcal{O})$ 的对象。
\begin{enumerate}
\item 如果 $K$ 是 $n$-伪凝聚的，且对 $i > a$ 有
$H^i(K) = 0$；而 $L$ 是 $m$-伪凝聚的，且对 $j > b$ 有
$H^j(L) = 0$，那么
$K \otimes_\mathcal{O}^\mathbf{L} L$ 是 $t$-伪凝聚的，其中
$t = \max(m + a, n + b)$。
\item 如果 $K$ 和 $L$ 都是伪凝聚的，那么
$K \otimes_\mathcal{O}^\mathbf{L} L$ 是伪凝聚的。
\end{enumerate}
\end{lemma}

\begin{proof}
证明 (1)。设 $U$ 是 $\mathcal{C}$ 的一个对象。
用某个覆盖的各成员代替 $U$，并用局部化
$\mathcal{C}/U$ 代替 $\mathcal{C}$ 后，可设存在严格完美复形
$\mathcal{K}^\bullet$、$\mathcal{L}^\bullet$ 以及映射
$\alpha : \mathcal{K}^\bullet \to K$ 和
$\beta : \mathcal{L}^\bullet \to L$，使得
$H^i(\alpha)$ 对 $i > n$ 是同构、对 $i = n$ 是满射，并且
$H^i(\beta)$ 对 $i > m$ 是同构、对 $i = m$ 是满射。
于是映射
$$
\alpha \otimes^\mathbf{L} \beta :
\text{Tot}(\mathcal{K}^\bullet \otimes_\mathcal{O} \mathcal{L}^\bullet)
\to K \otimes_\mathcal{O}^\mathbf{L} L
$$
在 $i > t$ 时的第 $i$ 个上同调层上诱导同构，并在
$i=t$ 时诱导满射。这由 Tor 谱序列得出（细节从略）。

\medskip\noindent
证明 (2)。设 $U$ 是 $\mathcal{C}$ 的一个对象。
先用某个覆盖的各成员代替 $U$，并用局部化
$\mathcal{C}/U$ 代替 $\mathcal{C}$，可归结到 $K$ 和 $L$
均上有界的情形。此时断言立即由 (1) 得出。
\end{proof}

\begin{lemma}
\label{sites-cohomology-lemma-summands-pseudo-coherent}
设 $(\mathcal{C}, \mathcal{O})$ 是一个环化站点，
$m \in \mathbf{Z}$。如果 $D(\mathcal{O})$ 中的
$K \oplus L$ 是 $m$-伪凝聚的（相应地，伪凝聚的），
那么 $K$ 和 $L$ 也都是如此。
\end{lemma}

\begin{proof}
假设 $K \oplus L$ 是 $m$-伪凝聚的。设 $U$ 是
$\mathcal{C}$ 的一个对象。用某个覆盖的各成员代替 $U$ 后，
可设 $K \oplus L \in D^-(\mathcal{O}_U)$，从而
$L \in D^-(\mathcal{O}_U)$。注意有有区别三角
$$
(K \oplus L, K \oplus L, L \oplus L[1]) =
(K, K, 0) \oplus (L, L, L \oplus L[1])
$$
，见《导出范畴》引理
\ref{derived-lemma-direct-sum-triangles}。
由引理 \ref{sites-cohomology-lemma-cone-pseudo-coherent}，
$L \oplus L[1]$ 是 $m$-伪凝聚的。因此
$L[1] \oplus L[2]$ 也是 $m$-伪凝聚的。
归纳可知 $L[n] \oplus L[n + 1]$ 是 $m$-伪凝聚的。
由于 $L$ 上有界，对充分大的 $n$，$L[n]$ 是 $m$-伪凝聚的。
于是，利用有区别三角
$$
(L[n], L[n] \oplus L[n - 1], L[n - 1])
$$
反向推导，可知 $L[n - 1], L[n - 2], \ldots, L$ 都是
$m$-伪凝聚的，正合所需。
\end{proof}

\begin{lemma}
\label{sites-cohomology-lemma-finite-cohomology}
设 $(\mathcal{C}, \mathcal{O})$ 是一个环化站点，
$K$ 是 $D(\mathcal{O})$ 的一个对象，并设
$m \in \mathbf{Z}$。
\begin{enumerate}
\item 如果 $K$ 是 $m$-伪凝聚的，且对 $i > m$ 有
$H^i(K) = 0$，那么 $H^m(K)$ 是有限型 $\mathcal{O}$-模。
\item 如果 $K$ 是 $m$-伪凝聚的，且对 $i > m + 1$ 有
$H^i(K) = 0$，那么 $H^{m + 1}(K)$ 是有限表示
$\mathcal{O}$-模。
\end{enumerate}
\end{lemma}

\begin{proof}
证明 (1)。设 $U$ 是 $\mathcal{C}$ 的一个对象。需要证明：
$H^m(K)$ 在 $U$ 的某个覆盖的各成员上可由有限多个截面生成
（见《站点上的模》定义
\ref{sites-modules-definition-site-local}）。
因此，在证明过程中，可以有限多次选取覆盖
$\{U_i \to U\}$，并用 $\mathcal{C}/U_i$ 代替
$\mathcal{C}$、用 $U_i$ 代替 $U$。特别地，由我们的定义，
可设存在严格完美复形 $\mathcal{E}^\bullet$ 和映射
$\alpha : \mathcal{E}^\bullet \to K$，它在次数 $>m$ 的上同调上诱导同构，
并在次数 $m$ 上诱导满射。只需对 $\mathcal{E}^\bullet$ 证明结论。
设 $n$ 是使 $\mathcal{E}^n \not = 0$ 的最大整数。
若 $n=m$，则 $H^m(\mathcal{E}^\bullet)$ 是
$\mathcal{E}^n$ 的商，结论显然。
若 $n>m$，则由于 $H^n(E^\bullet)=0$，
$\mathcal{E}^{n - 1} \to \mathcal{E}^n$ 是满射。
由引理 \ref{sites-cohomology-lemma-local-lift-map}，用某个覆盖的各成员代替
$U$ 后，可找到该满射的截面，并写成
$\mathcal{E}^{n - 1} = \mathcal{E}' \oplus \mathcal{E}^n$。
因而，只需对复形 $(\mathcal{E}')^\bullet$ 证明结论；
它与 $\mathcal{E}^\bullet$ 相同，但在次数 $n-1$ 处为
$\mathcal{E}'$，在次数 $n$ 处为 $0$。对 $n$ 归纳即得结论。

\medskip\noindent
证明 (2)。取 $\mathcal{C}$ 的一个对象 $U$。
如 (1) 的证明，可在 $U$ 上局部工作。
因而可设存在严格完美复形 $\mathcal{E}^\bullet$ 和映射
$\alpha : \mathcal{E}^\bullet \to K$，它在次数 $>m$ 的上同调上诱导同构，
并在次数 $m$ 上诱导满射。如 (1) 的证明，可归结到对
$i > m+1$ 有 $\mathcal{E}^i=0$ 的情形。于是
$H^{m + 1}(K) \cong H^{m + 1}(\mathcal{E}^\bullet) =
\Coker(\mathcal{E}^m \to \mathcal{E}^{m + 1})$
是有限表示的。
\end{proof}








\section{Tor 维数}
\label{sites-cohomology-section-tor}

\noindent
本节更细致地考察平坦模分辨。

\begin{definition}
\label{sites-cohomology-definition-tor-amplitude}
设 $(\mathcal{C}, \mathcal{O})$ 是一个环化站点，
$E$ 是 $D(\mathcal{O})$ 的一个对象。
设 $a,b \in \mathbf{Z}$ 且 $a \leq b$。
\begin{enumerate}
\item 如果对所有 $\mathcal{O}$-模 $\mathcal{F}$ 以及所有
$i \not \in [a,b]$，都有
$H^i(E \otimes_\mathcal{O}^\mathbf{L} \mathcal{F}) = 0$，
则称 $E$ 的 {\it Tor 振幅在 $[a,b]$ 中}。
\item 如果对某些 $a,b$，$E$ 的 Tor 振幅在 $[a,b]$ 中，
则称 $E$ 具有{\it 有限 Tor 维数}。
\item 如果对 $\mathcal{C}$ 的任意对象 $U$，都存在覆盖
$\{U_i \to U\}$，使得所有 $E|_{U_i}$ 都具有有限 Tor 维数，
则称 $E$ {\it 局部具有有限 Tor 维数}。
\end{enumerate}
如果把 $\mathcal{F}[0]$ 看作 $D(\mathcal{O})$ 的对象时，
其 Tor 振幅在 $[-d,0]$ 中，则称 $\mathcal{O}$-模
$\mathcal{F}$ 的 {\it Tor 维数 $\leq d$}。
\end{definition}

\noindent
注意，如果定义中的 $E$ 具有有限 Tor 维数，那么
$E$ 是 $D^b(\mathcal{O})$ 的对象；这可由在上述定义中取
$\mathcal{F} = \mathcal{O}$ 看出。

\begin{lemma}
\label{sites-cohomology-lemma-last-one-flat}
设 $(\mathcal{C}, \mathcal{O})$ 是一个环化站点。
设 $\mathcal{E}^\bullet$ 是平坦 $\mathcal{O}$-模的上有界复形，
其 Tor 振幅在 $[a,b]$ 中。那么
$\Coker(d_{\mathcal{E}^\bullet}^{a - 1})$ 是平坦
$\mathcal{O}$-模。
\end{lemma}

\begin{proof}
由于 $\mathcal{E}^\bullet$ 是平坦模的上有界复形，对任意
$\mathcal{O}$-模 $\mathcal{F}$ 都有
$\mathcal{E}^\bullet \otimes_\mathcal{O} \mathcal{F} =
\mathcal{E}^\bullet \otimes_\mathcal{O}^{\mathbf{L}} \mathcal{F}$
。因此，对每个 $\mathcal{O}$-模 $\mathcal{F}$，序列
$$
\mathcal{E}^{a - 2} \otimes_\mathcal{O} \mathcal{F} \to
\mathcal{E}^{a - 1} \otimes_\mathcal{O} \mathcal{F} \to
\mathcal{E}^a \otimes_\mathcal{O} \mathcal{F}
$$
在中间正合。由于
$\mathcal{E}^{a - 2} \to \mathcal{E}^{a - 1} \to \mathcal{E}^a \to
\Coker(d^{a - 1}) \to 0$
是平坦分辨，这蕴含对所有 $\mathcal{O}$-模 $\mathcal{F}$ 都有
$\text{Tor}_1^\mathcal{O}(\Coker(d^{a - 1}), \mathcal{F}) = 0$
。这意味着 $\Coker(d^{a - 1})$ 平坦，见
引理 \ref{sites-cohomology-lemma-flat-tor-zero}。
\end{proof}

\begin{lemma}
\label{sites-cohomology-lemma-tor-amplitude}
设 $(\mathcal{C}, \mathcal{O})$ 是一个环化站点，
$E$ 是 $D(\mathcal{O})$ 的一个对象。
设 $a,b \in \mathbf{Z}$ 且 $a \leq b$。下列条件等价：
\begin{enumerate}
\item $E$ 的 Tor 振幅在 $[a,b]$ 中。
\item $E$ 由平坦 $\mathcal{O}$-模复形
$\mathcal{E}^\bullet$ 表示，并且对
$i \not \in [a,b]$ 有 $\mathcal{E}^i=0$。
\end{enumerate}
\end{lemma}

\begin{proof}
如果 (2) 成立，则可按下式计算：
$E \otimes_\mathcal{O}^\mathbf{L} \mathcal{F} =
\mathcal{E}^\bullet \otimes_\mathcal{O} \mathcal{F}$
，显然 (1) 成立。

\medskip\noindent
假设 (1) 成立。可用平坦 $\mathcal{O}$-模的上有界复形
$\mathcal{K}^\bullet$ 表示 $E$，见第 \ref{sites-cohomology-section-flat} 节。
设 $n$ 是使 $\mathcal{K}^n \not = 0$ 的最大整数。
若 $n>b$，则因 $H^n(\mathcal{K}^\bullet)=0$，
$\mathcal{K}^{n - 1} \to \mathcal{K}^n$ 是满射。
由于 $\mathcal{K}^n$ 平坦，
$\Ker(\mathcal{K}^{n - 1} \to \mathcal{K}^n)$ 平坦
（《站点上的模》引理 \ref{sites-modules-lemma-flat-ses}）。
因此，可以用 $\tau_{\leq n - 1}\mathcal{K}^\bullet$
代替 $\mathcal{K}^\bullet$。于是，对 $n$ 归纳，可归结到
$K^\bullet$ 是平坦 $\mathcal{O}$-模复形，并且对 $i>b$ 有
$\mathcal{K}^i=0$ 的情形。

\medskip\noindent
令 $\mathcal{E}^\bullet = \tau_{\geq a}\mathcal{K}^\bullet$。
除 $\mathcal{E}^a$ 的平坦性外，一切均显然；而这一点立即由
引理 \ref{sites-cohomology-lemma-last-one-flat} 和定义得出。
\end{proof}

\begin{lemma}
\label{sites-cohomology-lemma-bounded-below-tor-amplitude}
设 $(\mathcal{C}, \mathcal{O})$ 是一个环化站点，
$E$ 是 $D(\mathcal{O})$ 的一个对象。设
$a \in \mathbf{Z}$。下列条件等价：
\begin{enumerate}
\item $E$ 的 Tor 振幅在 $[a,\infty]$ 中。
\item $E$ 可由平坦 $\mathcal{O}$-模的 K-平坦复形
$\mathcal{E}^\bullet$ 表示，并且对
$i \not \in [a,\infty]$ 有 $\mathcal{E}^i=0$。
\end{enumerate}
此外，可以选取 $\mathcal{E}^\bullet$，使得沿任意环化站点态射的拉回
都是各项平坦的 K-平坦复形。
\end{lemma}

\begin{proof}
(2) $\Rightarrow$ (1) 是立即的。假设 (1) 成立。
首先选取表示 $E$ 且各项平坦的 K-平坦复形
$\mathcal{K}^\bullet$，见引理 \ref{sites-cohomology-lemma-K-flat-resolution}。
对任意 $\mathcal{O}$-模 $\mathcal{M}$，下列复形在中间的上同调
$$
\mathcal{K}^{n - 1} \otimes_\mathcal{O} \mathcal{M} \to
\mathcal{K}^n \otimes_\mathcal{O} \mathcal{M} \to
\mathcal{K}^{n + 1} \otimes_\mathcal{O} \mathcal{M}
$$
计算 $H^n(E \otimes_\mathcal{O}^\mathbf{L} \mathcal{M})$。
当 $n<a$ 时它总为零。因此，对复形
$\ldots \to \mathcal{K}^{a - 1} \to \mathcal{K}^a \to
\mathcal{K}^{a + 1}$ 应用引理 \ref{sites-cohomology-lemma-last-one-flat}，
可知
$\mathcal{N} = \Coker(\mathcal{K}^{a - 1} \to \mathcal{K}^a)$
是平坦 $\mathcal{O}$-模。令
$$
\mathcal{E}^\bullet = \tau_{\geq a}\mathcal{K}^\bullet =
(\ldots \to 0 \to \mathcal{N} \to \mathcal{K}^{a + 1} \to \ldots )
$$
映射 $\mathcal{K}^\bullet \to \mathcal{E}^\bullet$ 的核
$\mathcal{L}^\bullet$ 是复形
$$
\mathcal{L}^\bullet = (\ldots \to \mathcal{K}^{a - 1} \to
\mathcal{I} \to 0 \to \ldots)
$$
，其中 $\mathcal{I} \subset \mathcal{K}^a$ 是
$\mathcal{K}^{a - 1} \to \mathcal{K}^a$ 的像。
由于有短正合列
$0 \to \mathcal{I} \to \mathcal{K}^a \to \mathcal{N} \to 0$
，可知 $\mathcal{I}$ 是平坦 $\mathcal{O}$-模。
因此，$\mathcal{L}^\bullet$ 是平坦模的上有界复形，
故由引理 \ref{sites-cohomology-lemma-bounded-flat-K-flat} 可知它 K-平坦。
随后由引理 \ref{sites-cohomology-lemma-K-flat-two-out-of-three-ses} 可知
$\mathcal{E}^\bullet$ K-平坦。

\medskip\noindent
证明最后一个断言。设
$f : (\mathcal{C}', \mathcal{O}') \to (\mathcal{C}, \mathcal{O})$
是环化站点态射。由引理 \ref{sites-cohomology-lemma-pullback-K-flat}，
复形 $f^*\mathcal{K}^\bullet$ K-平坦且各项平坦。
复形 $f^*\mathcal{L}^\bullet$ 是平坦 $\mathcal{O}'$-模的上有界复形，
故 K-平坦。我们有 $\mathcal{O}'$-模复形的短正合列
$$
0 \to f^*\mathcal{L}^\bullet \to f^*\mathcal{K}^\bullet \to
f^*\mathcal{E}^\bullet \to 0
$$
，因为平坦模的短正合列
$0 \to \mathcal{I} \to \mathcal{K}^a \to \mathcal{N} \to 0$
拉回后仍为短正合列。由引理
\ref{sites-cohomology-lemma-K-flat-two-out-of-three-ses}，
复形 $f^*\mathcal{E}^\bullet$ K-平坦，证明完毕。
\end{proof}

\begin{lemma}
\label{sites-cohomology-lemma-tor-amplitude-pullback}
设 $(f, f^\sharp) : (\mathcal{C}, \mathcal{O}_\mathcal{C}) \to
(\mathcal{D}, \mathcal{O}_\mathcal{D})$
是环化站点的态射。设 $E$ 是
$D(\mathcal{O}_\mathcal{D})$ 的一个对象。
如果 $E$ 的 Tor 振幅在 $[a,b]$ 中，那么
$Lf^*E$ 的 Tor 振幅在 $[a,b]$ 中。
\end{lemma}

\begin{proof}
假设 $E$ 的 Tor 振幅在 $[a,b]$ 中。由
引理 \ref{sites-cohomology-lemma-tor-amplitude}，可用平坦
$\mathcal{O}$-模的复形 $\mathcal{E}^\bullet$ 表示 $E$，
其中对 $i \not \in [a,b]$ 有 $\mathcal{E}^i=0$。
于是 $Lf^*E$ 由 $f^*\mathcal{E}^\bullet$ 表示。
由《站点上的模》引理 \ref{sites-modules-lemma-pullback-flat}，
模 $f^*\mathcal{E}^i$ 是平坦的。因此，由
引理 \ref{sites-cohomology-lemma-tor-amplitude} 可知
$Lf^*E$ 的 Tor 振幅在 $[a,b]$ 中。
\end{proof}

\begin{lemma}
\label{sites-cohomology-lemma-cone-tor-amplitude}
设 $(\mathcal{C}, \mathcal{O})$ 是一个环化站点。
设 $(K,L,M,f,g,h)$ 是 $D(\mathcal{O})$ 中的有区别三角，
并设 $a,b \in \mathbf{Z}$。
\begin{enumerate}
\item 如果 $K$ 的 Tor 振幅在 $[a+1,b+1]$ 中，而
$L$ 的 Tor 振幅在 $[a,b]$ 中，那么
$M$ 的 Tor 振幅在 $[a,b]$ 中。
\item 如果 $K$ 和 $M$ 的 Tor 振幅都在 $[a,b]$ 中，那么
$L$ 的 Tor 振幅在 $[a,b]$ 中。
\item 如果 $L$ 的 Tor 振幅在 $[a+1,b+1]$ 中，而
$M$ 的 Tor 振幅在 $[a,b]$ 中，那么
$K$ 的 Tor 振幅在 $[a+1,b+1]$ 中。
\end{enumerate}
\end{lemma}

\begin{proof}
从略。提示：这直接由与有区别三角相联系的上同调长正合列以及
$- \otimes_\mathcal{O}^{\mathbf{L}} \mathcal{F}$
保持有区别三角这一事实得出。最容易证明的是 (2)，其余各条由平移得出。
\end{proof}

\begin{lemma}
\label{sites-cohomology-lemma-tensor-tor-amplitude}
设 $(\mathcal{C}, \mathcal{O})$ 是一个环化站点，
$K,L$ 是 $D(\mathcal{O})$ 的对象。如果 $K$ 的 Tor 振幅在
$[a,b]$ 中，而 $L$ 的 Tor 振幅在 $[c,d]$ 中，那么
$K \otimes_\mathcal{O}^\mathbf{L} L$ 的 Tor 振幅在
$[a+c,b+d]$ 中。
\end{lemma}

\begin{proof}
从略。提示：使用 Tor 谱序列。
\end{proof}

\begin{lemma}
\label{sites-cohomology-lemma-summands-tor-amplitude}
设 $(\mathcal{C}, \mathcal{O})$ 是一个环化站点，
$a,b \in \mathbf{Z}$。对 $D(\mathcal{O})$ 的对象 $K,L$，
如果 $K \oplus L$ 的 Tor 振幅在 $[a,b]$ 中，那么
$K$ 和 $L$ 也都如此。
\end{lemma}

\begin{proof}
由 Tor 函子的加性立即可得。
\end{proof}

\begin{lemma}
\label{sites-cohomology-lemma-bounded}
设 $(\mathcal{C}, \mathcal{O})$ 是一个环化站点，
$\mathcal{I} \subset \mathcal{O}$ 是理想层，
$K$ 是 $D(\mathcal{O})$ 的一个对象。
\begin{enumerate}
\item 如果
$K \otimes_\mathcal{O}^\mathbf{L} \mathcal{O}/\mathcal{I}$
上有界，那么对所有 $n$，
$K \otimes_\mathcal{O}^\mathbf{L} \mathcal{O}/\mathcal{I}^n$
一致上有界。
\item 如果
$K \otimes_\mathcal{O}^\mathbf{L} \mathcal{O}/\mathcal{I}$
作为 $D(\mathcal{O}/\mathcal{I})$ 的对象，其 Tor 振幅在
$[a,b]$ 中，那么对所有 $n$，
$K \otimes_\mathcal{O}^\mathbf{L} \mathcal{O}/\mathcal{I}^n$
作为 $D(\mathcal{O}/\mathcal{I}^n)$ 的对象，其 Tor 振幅也在
$[a,b]$ 中。
\end{enumerate}
\end{lemma}

\begin{proof}
证明 (1)。假设
$K \otimes_\mathcal{O}^\mathbf{L} \mathcal{O}/\mathcal{I}$
上有界，例如
$H^i(K \otimes_\mathcal{O}^\mathbf{L} \mathcal{O}/\mathcal{I}) = 0$
对 $i>b$ 成立。注意有有区别三角
$$
K \otimes_\mathcal{O}^\mathbf{L}
\mathcal{I}^n/\mathcal{I}^{n + 1} \to
K \otimes_\mathcal{O}^\mathbf{L}
\mathcal{O}/\mathcal{I}^{n + 1} \to
K \otimes_\mathcal{O}^\mathbf{L}
\mathcal{O}/\mathcal{I}^n \to
K \otimes_\mathcal{O}^\mathbf{L}
\mathcal{I}^n/\mathcal{I}^{n + 1}[1]
$$
，并且
$$
K \otimes_\mathcal{O}^\mathbf{L}
\mathcal{I}^n/\mathcal{I}^{n + 1} =
\left(
K \otimes_\mathcal{O}^\mathbf{L}
\mathcal{O}/\mathcal{I}\right)
\otimes_{\mathcal{O}/\mathcal{I}}^\mathbf{L}
\mathcal{I}^n/\mathcal{I}^{n + 1}
$$
。归纳可知，对所有 $n$ 以及所有 $i>b$，
$H^i(K \otimes_\mathcal{O}^\mathbf{L} \mathcal{O}/\mathcal{I}^n) = 0$
。

\medskip\noindent
证明 (2)。假设
$K \otimes_\mathcal{O}^\mathbf{L} \mathcal{O}/\mathcal{I}$
作为 $D(\mathcal{O}/\mathcal{I})$ 的对象，其 Tor 振幅在
$[a,b]$ 中。设 $\mathcal{F}$ 是
$\mathcal{O}/\mathcal{I}^n$-模层。则有有限滤过
$$
0 \subset \mathcal{I}^{n - 1}\mathcal{F} \subset \ldots
\subset \mathcal{I}\mathcal{F} \subset \mathcal{F}
$$
，其逐次商是 $\mathcal{O}/\mathcal{I}$-模层。因此，为证明
$K \otimes_\mathcal{O}^\mathbf{L} \mathcal{O}/\mathcal{I}^n$
的 Tor 振幅在 $[a,b]$ 中，只需证明对所有
$\mathcal{O}/\mathcal{I}$-模 $\mathcal{G}$ 以及所有
$i \not \in [a,b]$，都有
$H^i(K \otimes_\mathcal{O}^\mathbf{L} \mathcal{O}/\mathcal{I}^n
\otimes_{\mathcal{O}/\mathcal{I}^n}^\mathbf{L} \mathcal{G})$
为零。由于
$$
\left(K \otimes_\mathcal{O}^\mathbf{L} \mathcal{O}/\mathcal{I}^n\right)
\otimes_{\mathcal{O}/\mathcal{I}^n}^\mathbf{L} \mathcal{G}
=
\left(K \otimes_\mathcal{O}^\mathbf{L} \mathcal{O}/\mathcal{I}\right)
\otimes_{\mathcal{O}/\mathcal{I}}^\mathbf{L} \mathcal{G}
$$
对每个 $\mathcal{O}/\mathcal{I}$-模层 $\mathcal{G}$ 都成立，结论随即得出。
\end{proof}

\begin{lemma}
\label{sites-cohomology-lemma-tor-amplitude-stalk}
设 $(\mathcal{C}, \mathcal{O})$ 是一个环化站点，
$E$ 是 $D(\mathcal{O})$ 的一个对象，并设
$a,b \in \mathbf{Z}$。
\begin{enumerate}
\item 如果 $E$ 的 Tor 振幅在 $[a,b]$ 中，那么对站点
$\mathcal{C}$ 的每个点 $p$，
$D(\mathcal{O}_p)$ 的对象 $E_p$ 的 Tor 振幅也在 $[a,b]$ 中。
\item 如果 $\mathcal{C}$ 有足够多的点，那么逆命题成立。
\end{enumerate}
\end{lemma}

\begin{proof}
证明 (1)。这是因为在 $p$ 处取茎等同于沿环化站点态射
$(p, \mathcal{O}_p) \to (\mathcal{C}, \mathcal{O})$
拉回，因而可应用引理 \ref{sites-cohomology-lemma-tor-amplitude-pullback}。

\medskip\noindent
证明 (2)。如果 $\mathcal{C}$ 有足够多的点，则可在各茎处检验
$H^i(E \otimes_\mathcal{O}^\mathbf{L} \mathcal{F})$
的消失性，见《站点上的模》引理
\ref{sites-modules-lemma-check-exactness-stalks}。
由于
$H^i(E \otimes_\mathcal{O}^\mathbf{L} \mathcal{F})_p =
H^i(E_p \otimes_{\mathcal{O}_p}^\mathbf{L} \mathcal{F}_p)$，结论得证。
\end{proof}










\section{完美复形}
\label{sites-cohomology-section-perfect}

\noindent
本节讨论环化站点上完美复形的性质。

\begin{definition}
\label{sites-cohomology-definition-perfect}
设 $(\mathcal{C}, \mathcal{O})$ 是一个环化站点，
$\mathcal{E}^\bullet$ 是 $\mathcal{O}$-模复形。
如果对 $\mathcal{C}$ 的每个对象 $U$，都存在覆盖
$\{U_i \to U\}$，使得对每个 $i$ 都存在复形态射
$\mathcal{E}_i^\bullet \to \mathcal{E}^\bullet|_{U_i}$
，它是拟同构，且 $\mathcal{E}_i^\bullet$ 严格完美，
则称 $\mathcal{E}^\bullet$ 是{\it 完美的}。
如果 $D(\mathcal{O})$ 的对象 $E$ 能由完美
$\mathcal{O}$-模复形表示，则称 $E$ 是{\it 完美的}。
\end{definition}

\noindent
如果 $\Sh(\mathcal{C})$ 拟紧
（《站点》第 \ref{sites-section-quasi-compact} 节），
那么 $D(\mathcal{O})$ 的完美对象属于
$D^b(\mathcal{O})$。否则未必如此。

\begin{lemma}
\label{sites-cohomology-lemma-perfect-independent-representative}
设 $(\mathcal{C}, \mathcal{O})$ 是一个环化站点，
$E$ 是 $D(\mathcal{O})$ 的一个对象。
\begin{enumerate}
\item 如果 $\mathcal{C}$ 有终对象 $X$，并且存在覆盖
$\{U_i \to X\}$、$\mathcal{O}_{U_i}$-模的严格完美复形
$\mathcal{E}_i^\bullet$，以及 $D(\mathcal{O}_{U_i})$ 中的同构
$\alpha_i : \mathcal{E}_i^\bullet \to E|_{U_i}$，
那么 $E$ 是完美的。
\item 如果 $E$ 是完美的，那么表示 $E$ 的任意复形都是完美的。
\end{enumerate}
\end{lemma}

\begin{proof}
与引理 \ref{sites-cohomology-lemma-pseudo-coherent-independent-representative}
的证明完全相同。
\end{proof}

\begin{lemma}
\label{sites-cohomology-lemma-perfect-precise}
设 $(\mathcal{C}, \mathcal{O})$ 是一个环化站点，
$E$ 是 $D(\mathcal{O})$ 的一个对象。
设整数 $a \leq b$。如果 $E$ 的 Tor 振幅在 $[a,b]$ 中，
并且 $E$ 是 $(a-1)$-伪凝聚的，那么 $E$ 是完美的。
\end{lemma}

\begin{proof}
设 $U$ 是 $\mathcal{C}$ 的一个对象。用某个覆盖的各成员代替
$U$，并用局部化 $\mathcal{C}/U$ 代替 $\mathcal{C}$ 后，
可设存在严格完美复形 $\mathcal{E}^\bullet$ 和映射
$\alpha : \mathcal{E}^\bullet \to E$，使得
$H^i(\alpha)$ 对 $i \geq a$ 是同构。
我们可以并确实用
$\sigma_{\geq a - 1}\mathcal{E}^\bullet$ 代替
$\mathcal{E}^\bullet$。选取有区别三角
$$
\mathcal{E}^\bullet \to E \to C \to \mathcal{E}^\bullet[1]
$$
。由 $E$ 和 $\mathcal{E}^\bullet$ 的上同调层之消失性以及
关于 $\alpha$ 的假设，得到
$C \cong \mathcal{K}[2-a]$，其中
$\mathcal{K} = \Ker(\mathcal{E}^{a-1} \to \mathcal{E}^a)$。
设 $\mathcal{F}$ 是一个 $\mathcal{O}$-模。
应用 $- \otimes_\mathcal{O}^\mathbf{L} \mathcal{F}$，
关于 $E$ 的 Tor 振幅在 $[a,b]$ 中的假设蕴含
$\mathcal{K} \otimes_\mathcal{O} \mathcal{F} \to
\mathcal{E}^{a - 1} \otimes_\mathcal{O} \mathcal{F}$ 的像为
$\Ker(\mathcal{E}^{a - 1} \otimes_\mathcal{O} \mathcal{F}
\to \mathcal{E}^a \otimes_\mathcal{O} \mathcal{F})$.
因此
$\text{Tor}_1^\mathcal{O}(\mathcal{E}', \mathcal{F})=0$，
其中
$\mathcal{E}' = \Coker(\mathcal{E}^{a-1}\to\mathcal{E}^a)$。
故 $\mathcal{E}'$ 平坦（引理 \ref{sites-cohomology-lemma-flat-tor-zero}）。
于是，由《站点上的模》引理
\ref{sites-modules-lemma-flat-locally-finite-presentation}，
存在覆盖 $\{U_i \to U\}$，使得
$\mathcal{E}'|_{U_i}$ 是某个有限自由模的直和项。
因此复形
$$
\mathcal{E}'|_{U_i} \to \mathcal{E}^{a + 1}|_{U_i} \to \ldots \to
\mathcal{E}^b|_{U_i}
$$
与 $E|_{U_i}$ 拟同构，所以 $E$ 是完美的。
\end{proof}

\begin{lemma}
\label{sites-cohomology-lemma-perfect}
设 $(\mathcal{C}, \mathcal{O})$ 是一个环化站点，
$E$ 是 $D(\mathcal{O})$ 的一个对象。下列条件等价：
\begin{enumerate}
\item $E$ 是完美的；
\item $E$ 是伪凝聚的，并且局部具有有限 Tor 维数。
\end{enumerate}
\end{lemma}

\begin{proof}
假设 (1)。设 $U$ 是 $\mathcal{C}$ 的一个对象。
由定义，存在覆盖 $\{U_i \to U\}$，使得
$E|_{U_i}$ 由严格完美复形表示。因此，由引理
\ref{sites-cohomology-lemma-pseudo-coherent-independent-representative}，
$E$ 是伪凝聚的（即对所有 $m$ 都是 $m$-伪凝聚的）。
此外，有限自由模的直和项是平坦的，故由引理
\ref{sites-cohomology-lemma-tor-amplitude}，$E|_{U_i}$ 具有有限 Tor 维数。
所以 (2) 成立。

\medskip\noindent
假设 (2)。设 $U$ 是 $\mathcal{C}$ 的一个对象。
用某个覆盖的各成员代替 $U$ 后，可设存在整数
$a \leq b$，使得 $E|_U$ 的 Tor 振幅在 $[a,b]$ 中。
由于 $E|_U$ 对所有 $m$ 都是 $m$-伪凝聚的，
应用引理 \ref{sites-cohomology-lemma-perfect-precise} 即得结论。
\end{proof}

\begin{lemma}
\label{sites-cohomology-lemma-perfect-pullback}
设 $(f, f^\sharp) : (\mathcal{C}, \mathcal{O}_\mathcal{C}) \to
(\mathcal{D}, \mathcal{O}_\mathcal{D})$
是环化站点的态射。设 $E$ 是
$D(\mathcal{O}_\mathcal{D})$ 的一个对象。
如果 $E$ 在 $D(\mathcal{O}_\mathcal{D})$ 中是完美的，
那么 $Lf^*E$ 在 $D(\mathcal{O}_\mathcal{C})$ 中是完美的。
\end{lemma}

\begin{proof}
这由引理 \ref{sites-cohomology-lemma-perfect}、
\ref{sites-cohomology-lemma-tor-amplitude-pullback} 和
\ref{sites-cohomology-lemma-pseudo-coherent-pullback} 得出。
\end{proof}

\begin{lemma}
\label{sites-cohomology-lemma-two-out-of-three-perfect}
设 $(\mathcal{C}, \mathcal{O})$ 是一个环化站点。
设 $(K,L,M,f,g,h)$ 是 $D(\mathcal{O})$ 中的有区别三角。
如果 $K,L,M$ 中有两个是完美的，那么第三个也是完美的。
\end{lemma}

\begin{proof}
第一种证明：合并引理 \ref{sites-cohomology-lemma-perfect}、
\ref{sites-cohomology-lemma-cone-pseudo-coherent} 和
\ref{sites-cohomology-lemma-cone-tor-amplitude}。

第二种证明（概要）：设 $K$ 和 $L$ 完美。设 $U$ 是
$\mathcal{C}$ 的一个对象。用某个覆盖的各成员代替 $U$ 后，
可设 $K|_U$ 和 $L|_U$ 分别由严格完美复形
$\mathcal{K}^\bullet$ 和 $\mathcal{L}^\bullet$ 表示。
再次用某个覆盖的各成员代替 $U$ 后，可设映射
$K|_U \to L|_U$ 由复形映射
$\alpha : \mathcal{K}^\bullet \to \mathcal{L}^\bullet$ 给出，
见引理 \ref{sites-cohomology-lemma-local-actual}。
于是 $M|_U$ 同构于 $\alpha$ 的锥；由引理
\ref{sites-cohomology-lemma-cone}，该锥严格完美。
\end{proof}

\begin{lemma}
\label{sites-cohomology-lemma-tensor-perfect}
设 $(\mathcal{C}, \mathcal{O})$ 是一个环化站点。
如果 $K,L$ 是 $D(\mathcal{O})$ 的完美对象，那么
$K \otimes_\mathcal{O}^\mathbf{L} L$ 也是完美对象。
\end{lemma}

\begin{proof}
由引理 \ref{sites-cohomology-lemma-perfect}、
\ref{sites-cohomology-lemma-tensor-pseudo-coherent} 和
\ref{sites-cohomology-lemma-tensor-tor-amplitude} 得出。
\end{proof}

\begin{lemma}
\label{sites-cohomology-lemma-summands-perfect}
设 $(\mathcal{C}, \mathcal{O})$ 是一个环化站点。
如果 $K \oplus L$ 是 $D(\mathcal{O})$ 的完美对象，
那么 $K$ 和 $L$ 也都是完美对象。
\end{lemma}

\begin{proof}
由引理 \ref{sites-cohomology-lemma-perfect}、
\ref{sites-cohomology-lemma-summands-pseudo-coherent} 和
\ref{sites-cohomology-lemma-summands-tor-amplitude} 得出。
\end{proof}














\section{对偶}
\label{sites-cohomology-section-duals}

\noindent
本节刻画复形范畴和导出范畴中的可对偶对象。
特别地，我们将看到，$D(\mathcal{O})$ 的对象有对偶，当且仅当它是完美的
（这由例 \ref{sites-cohomology-example-dual-derived} 和
引理 \ref{sites-cohomology-lemma-left-dual-derived} 得出）。

\begin{lemma}
\label{sites-cohomology-lemma-symmetric-monoidal-cat-complexes}
设 $(\mathcal{C}, \mathcal{O})$ 是一个环化站点。以
$\mathcal{F}^\bullet \otimes \mathcal{G}^\bullet =
\text{Tot}(\mathcal{F}^\bullet \otimes_\mathcal{O} \mathcal{G}^\bullet)$
定义张量积的 $\mathcal{O}$-模复形范畴是一个对称幺半范畴。
\end{lemma}

\begin{proof}
略。提示：取在次数 $0$ 上为 $\mathcal{O}$、在其余次数上为零的复形
作为单位对象 $\mathbf{1}$，并取显然的同构
$\text{Tot}(\mathbf{1} \otimes_\mathcal{O} \mathcal{G}^\bullet) =
\mathcal{G}^\bullet$ 以及
$\text{Tot}(\mathcal{F}^\bullet \otimes_\mathcal{O} \mathbf{1}) =
\mathcal{F}^\bullet$。
要证明该引理，需要验证若干图表的交换性；见《范畴》定义
\ref{categories-definition-monoidal-category} 和
\ref{categories-definition-symmetric-monoidal-category}。
各项验证都很直接。
\end{proof}

\begin{example}
\label{sites-cohomology-example-dual}
设 $(\mathcal{C}, \mathcal{O})$ 是一个环化站点。设 $\mathcal{F}^\bullet$
为一个 $\mathcal{O}$-模复形，并且对每个
$U \in \Ob(\mathcal{C})$，都存在一个覆盖 $\{U_i \to U\}$，
使得 $\mathcal{F}^\bullet|_{U_i}$ 严格完美。考虑复形
$$
\mathcal{G}^\bullet = \SheafHom^\bullet(\mathcal{F}^\bullet, \mathcal{O})
$$
，如第 \ref{sites-cohomology-section-hom-complexes} 节所述。令
$$
\eta :
\mathcal{O}
\to
\text{Tot}(\mathcal{F}^\bullet \otimes_\mathcal{O} \mathcal{G}^\bullet)
\quad\text{和}\quad
\epsilon :
\text{Tot}(\mathcal{G}^\bullet \otimes_\mathcal{O} \mathcal{F}^\bullet)
\to
\mathcal{O}
$$
分别为 $\eta = \sum \eta_n$ 和 $\epsilon = \sum \epsilon_n$，
其中 $\eta_n : \mathcal{O} \to
\mathcal{F}^n \otimes_\mathcal{O} \mathcal{G}^{-n}$
以及
$\epsilon_n : \mathcal{G}^{-n} \otimes_\mathcal{O} \mathcal{F}^n
\to \mathcal{O}$ 如《站点上的模》例 \ref{sites-modules-example-dual}
中所定义。于是按照《范畴》定义 \ref{categories-definition-dual}，
$\mathcal{G}^\bullet, \eta, \epsilon$ 是 $\mathcal{F}^\bullet$ 的左对偶。
我们略去下列等式的验证：
$(1 \otimes \epsilon) \circ (\eta \otimes 1) = \text{id}_{\mathcal{F}^\bullet}$
和
$(\epsilon \otimes 1) \circ (1 \otimes \eta) =
\text{id}_{\mathcal{G}^\bullet}$。请与《代数详论》引理
\ref{more-algebra-lemma-left-dual-complex} 比较。
\end{example}

\begin{lemma}
\label{sites-cohomology-lemma-left-dual-complex}
设 $(\mathcal{C}, \mathcal{O})$ 是一个环化站点。设 $\mathcal{F}^\bullet$
为一个 $\mathcal{O}$-模复形。如果 $\mathcal{F}^\bullet$
在 $\mathcal{O}$-模复形的幺半范畴中有左对偶
（《范畴》定义 \ref{categories-definition-dual}），
那么对 $\mathcal{C}$ 的每个对象 $U$，都存在一个覆盖
$\{U_i \to U\}$，使得 $\mathcal{F}^\bullet|_{U_i}$
严格完美，并且该左对偶正是例 \ref{sites-cohomology-example-dual} 中构造的左对偶。
\end{lemma}

\begin{proof}
由左对偶的唯一性（《范畴》注记
\ref{categories-remark-left-dual-adjoint}），只要证明
$\mathcal{F}^\bullet$ 具有所述性质，就可得到最后一个断言。
设 $\mathcal{G}^\bullet, \eta, \epsilon$ 为一个左对偶。写成
$\eta = \sum \eta_n$ 和 $\epsilon = \sum \epsilon_n$，其中
$\eta_n : \mathcal{O} \to
\mathcal{F}^n \otimes_\mathcal{O} \mathcal{G}^{-n}$
以及
$\epsilon_n : \mathcal{G}^{-n} \otimes_\mathcal{O} \mathcal{F}^n
\to \mathcal{O}$。由于
$(1 \otimes \epsilon) \circ (\eta \otimes 1) = \text{id}_{\mathcal{F}^\bullet}$
和
$(\epsilon \otimes 1) \circ (1 \otimes \eta) = \text{id}_{\mathcal{G}^\bullet}$
由《范畴》定义 \ref{categories-definition-dual}，立即可见
$(1 \otimes \epsilon_n) \circ (\eta_n \otimes 1) = \text{id}_{\mathcal{F}^n}$
以及
$(\epsilon_n \otimes 1) \circ (1 \otimes \eta_n) =
\text{id}_{\mathcal{G}^{-n}}$.
换言之，$\mathcal{G}^{-n}$ 是 $\mathcal{F}^n$ 的左对偶，
而且《站点上的模》引理 \ref{sites-modules-lemma-left-dual-module}
适用于每个 $\mathcal{F}^n$。设 $U$ 是 $\mathcal{C}$ 的一个对象。
存在覆盖 $\{U_i \to U\}$，使得对每个 $i$，
只有有限多个 $\eta_n|_{U_i}$ 非零。因此，用 $U_i$ 替换 $U$ 后，
可以假设只有有限多个 $\eta_n|_U$ 非零；由所引引理，
这蕴含只有有限多个 $\mathcal{F}^n|_U$ 非零。
再次应用该引理，便可找到一个覆盖 $\{U_i \to U\}$，
使得每个 $\mathcal{F}^n|_{U_i}$ 都是有限自由
$\mathcal{O}$-模的直和项，证明完毕。
\end{proof}

\begin{lemma}
\label{sites-cohomology-lemma-dual-perfect-complex}
设 $(\mathcal{C}, \mathcal{O})$ 是一个环化站点。
设 $K$ 是 $D(\mathcal{O})$ 的一个完美对象。
则 $K^\vee = R\SheafHom(K, \mathcal{O})$ 也是完美对象，且
$(K^\vee)^\vee \cong K$。存在函子性同构
$$
M \otimes^\mathbf{L}_\mathcal{O} K^\vee = R\SheafHom_\mathcal{O}(K, M)
$$
和
$$
H^0(\mathcal{C}, M \otimes^\mathbf{L}_\mathcal{O} K^\vee) =
\Hom_{D(\mathcal{O})}(K, M)
$$
，其中 $M$ 是 $D(\mathcal{O})$ 的对象。
\end{lemma}

\begin{proof}
下文将不再特别说明地使用内部 Hom 的形成与限制可交换这一事实
（引理 \ref{sites-cohomology-lemma-restriction-RHom-to-U}）。设 $U$
是 $\mathcal{C}$ 的任意对象。要验证 $K^\vee$ 是完美的，
只需证明存在覆盖 $\{U_i \to U\}$，使得对所有 $i$，
$K^\vee|_{U_i}$ 都是完美的。存在典范映射
$$
K = R\SheafHom(\mathcal{O}_X, \mathcal{O}_X)
\otimes_{\mathcal{O}_X}^\mathbf{L} K \longrightarrow
R\SheafHom(R\SheafHom(K, \mathcal{O}_X), \mathcal{O}_X) =
(K^\vee)^\vee
$$
；见引理 \ref{sites-cohomology-lemma-internal-hom-evaluate}。只需证明存在覆盖
$\{U_i \to U\}$，使得该映射在 $\mathcal{C}/U_i$
上的限制对所有 $i$ 都是同构。由引理 \ref{sites-cohomology-lemma-dual}，
要证明最后一个断言，只需验证映射 (\ref{sites-cohomology-equation-eval})
$$
M \otimes^\mathbf{L}_\mathcal{O} K^\vee
\longrightarrow
R\SheafHom(K, M)
$$
是同构。这同样是上述意义下的局部问题。因此，只需在 $K$
由严格完美复形表示时证明该引理。

\medskip\noindent
假设 $K$ 由严格完美复形 $\mathcal{E}^\bullet$ 表示。
由引理 \ref{sites-cohomology-lemma-Rhom-strictly-perfect}，$K^\vee$
由下述复形表示：其次数 $-n$ 上的项为
$(\mathcal{E}^n)^\vee =
\SheafHom_\mathcal{O}(\mathcal{E}^n, \mathcal{O})$
。由于 $\mathcal{E}^n$ 是有限自由 $\mathcal{O}$-模的直和项，
$(\mathcal{E}^n)^\vee$ 也如此。因此，$K^\vee$
也由一个严格完美复形表示，从而 $K^\vee$ 是完美的。
映射 $K \to (K^\vee)^\vee$ 是同构，因为它在相差一个符号的意义下
由评价映射 $\mathcal{E}^n \to ((\mathcal{E}^n)^\vee)^\vee$
给出，而这些映射都是同构。要证明 (\ref{sites-cohomology-equation-eval}) 是同构，
用一个 K-平坦复形 $\mathcal{F}^\bullet$ 表示 $M$。
由引理 \ref{sites-cohomology-lemma-Rhom-strictly-perfect}，复形
$R\SheafHom(K, M)$ 由各项为
$$
\bigoplus\nolimits_{n = p + q}
\SheafHom_\mathcal{O}(\mathcal{E}^{-q}, \mathcal{F}^p)
$$
的复形表示。另一方面，对象
$M \otimes^\mathbf{L}_\mathcal{O} K^\vee$ 由各项为
$$
\bigoplus\nolimits_{n = p + q}
\mathcal{F}^p \otimes_\mathcal{O} (\mathcal{E}^{-q})^\vee
$$
的复形表示。因此，(\ref{sites-cohomology-equation-eval}) 是同构这一断言归结为：
当 $\mathcal{E}$ 是有限自由 $\mathcal{O}$-模的直和项且
$\mathcal{F}$ 是任意 $\mathcal{O}$-模时，典范映射
$$
\mathcal{F}
\otimes_\mathcal{O}
\SheafHom_\mathcal{O}(\mathcal{E}, \mathcal{O})
\longrightarrow
\SheafHom_\mathcal{O}(\mathcal{E}, \mathcal{F})
$$
是同构。这立即由 $\mathcal{E}$ 为有限自由模时的相应断言得出。
\end{proof}

\begin{lemma}
\label{sites-cohomology-lemma-symmetric-monoidal-derived}
设 $(\mathcal{C}, \mathcal{O})$ 是一个环化站点。导出范畴
$D(\mathcal{O})$ 是一个对称幺半范畴，其张量积由导出张量积给出，
并带有通常的结合性约束和交换性约束（符号规则见《代数详论》第
\ref{more-algebra-section-sign-rules} 节）。
\end{lemma}

\begin{proof}
略。与引理 \ref{sites-cohomology-lemma-symmetric-monoidal-cat-complexes} 比较。
\end{proof}

\begin{example}
\label{sites-cohomology-example-dual-derived}
设 $(\mathcal{C}, \mathcal{O})$ 是一个环化站点。设 $K$ 是
$D(\mathcal{O})$ 的一个完美对象。如引理
\ref{sites-cohomology-lemma-dual-perfect-complex} 中那样，令
$K^\vee = R\SheafHom(K, \mathcal{O})$。则映射
$$
K \otimes_\mathcal{O}^\mathbf{L} K^\vee \longrightarrow R\SheafHom(K, K)
$$
是同构（由该引理）。记映射
$$
\eta :
\mathcal{O}
\longrightarrow
K \otimes_\mathcal{O}^\mathbf{L} K^\vee
$$
为在上述同构下将 $1$ 映到与 $\text{id}_K$ 对应的截面的映射。
记
$$
\epsilon : 
K^\vee
\otimes_\mathcal{O}^\mathbf{L} K
\longrightarrow
\mathcal{O}
$$
为评价映射（例如，可以用引理
\ref{sites-cohomology-lemma-internal-hom-composition} 构造它）。于是按照《范畴》定义
\ref{categories-definition-dual}，$K^\vee, \eta, \epsilon$
是 $K$ 的左对偶。我们略去下列等式的验证：
$(1 \otimes \epsilon) \circ (\eta \otimes 1) = \text{id}_K$
和
$(\epsilon \otimes 1) \circ (1 \otimes \eta) =
\text{id}_{K^\vee}$.
\end{example}

\begin{lemma}
\label{sites-cohomology-lemma-left-dual-derived}
设 $(\mathcal{C}, \mathcal{O})$ 是一个环化站点。设 $M$ 是
$D(\mathcal{O})$ 的一个对象。如果 $M$ 在幺半范畴
$D(\mathcal{O})$ 中有左对偶（《范畴》定义
\ref{categories-definition-dual}），那么 $M$ 是完美的，
并且该左对偶正是例 \ref{sites-cohomology-example-dual-derived} 中构造的左对偶。
\end{lemma}

\begin{proof}
设 $N, \eta, \epsilon$ 为一个左对偶。注意，对 $\mathcal{C}$
的任意对象 $U$，限制 $N|_U, \eta|_U, \epsilon|_U$
都是 $M|_U$ 的一个左对偶。

\medskip\noindent
设 $U$ 是 $\mathcal{C}$ 的一个对象。只需找到 $\mathcal{C}$
的一个覆盖 $\{U_i \to U\}_{i \in I}$，使得 $M|_{U_i}$
是 $D(\mathcal{O}_{U_i})$ 的完美对象。因此，可以分别用
$\mathcal{C}/U, \mathcal{O}_U, M|_U, N|_U, \eta|_U, \epsilon|_U$
替换 $\mathcal{C}, \mathcal{O}, M, N, \eta, \epsilon$，并假设
$\mathcal{C}$ 有终对象 $X$。此外，在证明过程中，可以（有限多次）
用 $X$ 的覆盖 $\{U_i \to X\}$ 的成员替换 $X$。

\medskip\noindent
我们将多次使用下述论证。任取一个表示 $M$ 的 $\mathcal{O}$-模复形
$\mathcal{M}^\bullet$。选取一个表示 $N$ 的 K-平坦复形
$\mathcal{N}^\bullet$，其各项均为平坦 $\mathcal{O}$-模；
见引理 \ref{sites-cohomology-lemma-K-flat-resolution}。考虑映射
$$
\eta : \mathcal{O} \to
\text{Tot}(\mathcal{M}^\bullet \otimes_\mathcal{O} \mathcal{N}^\bullet)
$$
用某个覆盖的成员替换 $X$ 后，可以找到整数 $N$，并且对
$i = 1, \ldots, N$，找到整数 $n_i \in \mathbf{Z}$ 以及
$\mathcal{M}^{n_i}$ 和 $\mathcal{N}^{-n_i}$ 的截面 $f_i$ 和 $g_i$，使得
$$
\eta(1) = \sum\nolimits_i f_i \otimes g_i
$$
设 $\mathcal{K}^\bullet \subset \mathcal{M}^\bullet$
是任意一个包含截面 $f_i$（$i = 1, \ldots, N$）的
$\mathcal{O}$-模子复形。由于
$\text{Tot}(\mathcal{K}^\bullet \otimes_\mathcal{O} \mathcal{N}^\bullet)
\subset
\text{Tot}(\mathcal{M}^\bullet \otimes_\mathcal{O} \mathcal{N}^\bullet)$
（这是由模 $\mathcal{N}^n$ 的平坦性得到的），可见 $\eta$ 分解经过
$$
\tilde \eta :
\mathcal{O} \to
\text{Tot}(\mathcal{K}^\bullet \otimes_\mathcal{O} \mathcal{N}^\bullet)
$$
记 $K$ 为 $D(\mathcal{O})$ 中由 $\mathcal{K}^\bullet$
表示的对象，便得到交换图
$$
\xymatrix{
M \ar[rr]_-{\eta \otimes 1} \ar[rrd]_{\tilde \eta \otimes 1} & &
M \otimes^\mathbf{L} N \otimes^\mathbf{L} M
\ar[r]_-{1 \otimes \epsilon} &
M \\
& &
K \otimes^\mathbf{L} N \otimes^\mathbf{L} M
\ar[u] \ar[r]^-{1 \otimes \epsilon} &
K \ar[u]
}
$$
由于上行的复合是 $M$ 上的恒等映射，故 $M$
在 $D(\mathcal{O})$ 中是 $K$ 的一个直和项。

\medskip\noindent
首次应用上述论证时，可以选取子复形
$\mathcal{K}^\bullet = \sigma_{\geq a} \tau_{\leq b}\mathcal{M}^\bullet$
，其中对 $i = 1, \ldots, N$ 有 $a < n_i < b$。于是 $M$
在 $D(\mathcal{O})$ 中是一个有界复形的直和项，从而可以假设
$M$ 属于 $D^b(\mathcal{O})$。（回忆上述过程涉及用一个覆盖的成员替换
$X$。）

\medskip\noindent
由于 $M$ 属于 $D^b(\mathcal{O})$，可以选取 $\mathcal{M}^\bullet$
为一个由平坦模组成的上有界复形（由《模》引理
\ref{modules-lemma-module-quotient-flat} 和《导出范畴》引理
\ref{derived-lemma-subcategory-left-resolution}）。在上述论证中，
可以选取 $\mathcal{K}^\bullet = \sigma_{\geq a}\mathcal{M}^\bullet$，
其中对 $i = 1, \ldots, N$ 有 $a < n_i$。因此，可以假设
$M$ 在 $D(\mathcal{O})$ 中是一个由平坦模组成的有界复形的直和项。
特别地，$M$ 具有有限 Tor 振幅。

\medskip\noindent
设 $M$ 的 Tor 振幅包含于 $[a, b]$。假设 $M$ 是 $m$-伪凝聚的；
我们将证明（用一个覆盖的成员替换 $X$ 后）可以假设 $M$
是 $(m - 1)$-伪凝聚的。结合引理 \ref{sites-cohomology-lemma-perfect-precise}
以及 $M$ 无论如何都是 $(b + 1)$-伪凝聚的这一事实，这就会完成证明。
用一个覆盖的成员替换 $X$ 后，可以假设存在严格完美复形
$\mathcal{E}^\bullet$ 和 $D(\mathcal{O})$ 中的映射
$\alpha : \mathcal{E}^\bullet \to M$，使得 $H^i(\alpha)$
在 $i > m$ 时是同构，在 $i = m$ 时是满射。可以并且确实假设
$\mathcal{E}^i = 0$ 对 $i < m$ 成立。选取一个有区别三角
$$
\mathcal{E}^\bullet \to M \to L \to \mathcal{E}^\bullet[1]
$$
。注意，对 $i \geq m$ 有 $H^i(L) = 0$。因此，可以用一个满足
$\mathcal{L}^i = 0$（$i \geq m$）的复形 $\mathcal{L}^\bullet$
表示 $L$。映射 $L \to \mathcal{E}^\bullet[1]$ 由复形映射
$\mathcal{L}^\bullet \to \mathcal{E}^\bullet[1]$
给出；它除次数 $m - 1$ 外处处为零，而在次数 $m - 1$ 上得到映射
$\mathcal{L}^{m - 1} \to \mathcal{E}^m$；见《导出范畴》引理
\ref{derived-lemma-negative-exts}。于是 $M$ 由复形
$$
\mathcal{M}^\bullet :
\ldots \to
\mathcal{L}^{m - 2} \to
\mathcal{L}^{m - 1} \to
\mathcal{E}^m \to
\mathcal{E}^{m + 1} \to \ldots
$$
表示。把第二段中的论证应用于此复形，得到 $\mathcal{M}^{n_i}$
的截面 $f_i$，其中 $i = 1, \ldots, N$。对 $n < m$，令
$\mathcal{K}^n \subset \mathcal{L}^n$ 为由满足 $n_i = n$
的截面 $f_i$ 以及满足 $n_i = n - 1$ 的 $d(f_i)$ 所生成的
$\mathcal{O}$-子模。对 $n \geq m$，令
$\mathcal{K}^n = \mathcal{E}^n$。显然，存在有区别三角之间的态射
$$
\xymatrix{
\mathcal{E}^\bullet \ar[r] &
\mathcal{M}^\bullet \ar[r] &
\mathcal{L}^\bullet \ar[r] &
\mathcal{E}^\bullet[1] \\
\mathcal{E}^\bullet \ar[r] \ar[u] &
\mathcal{K}^\bullet \ar[r] \ar[u] &
\sigma_{\leq m - 1}\mathcal{K}^\bullet \ar[r] \ar[u] &
\mathcal{E}^\bullet[1] \ar[u]
}
$$
，其中所有态射均如上所示。记 $K$ 为 $D(\mathcal{O})$
中与复形 $\mathcal{K}^\bullet$ 对应的对象。
由证明第二段中的论证，得到 $D(\mathcal{O})$ 中的态射
$s : M \to K$，使得复合 $M \to K \to M$ 是 $M$ 上的恒等映射。
我们不知道图表
$$
\xymatrix{
\mathcal{E}^\bullet \ar[r] &
\mathcal{K}^\bullet \ar@{=}[r] &
K \\
\mathcal{E}^\bullet \ar[u]^{\text{id}} \ar[r]^i &
\mathcal{M}^\bullet \ar@{=}[r] &
M \ar[u]_s
}
$$
是否交换，但知道它与映射 $K \to M$ 复合后交换。由引理
\ref{sites-cohomology-lemma-local-actual}，用一个覆盖的成员替换 $X$ 后，
可以假设 $s \circ i$ 由复形映射
$\sigma : \mathcal{E}^\bullet \to \mathcal{K}^\bullet$ 给出。
由同一个引理，可以假设 $\sigma$ 与包含映射
$\mathcal{K}^\bullet \subset \mathcal{M}^\bullet$ 的复合通过某个同伦
$\{h^i : \mathcal{E}^i \to \mathcal{M}^{i - 1}\}$ 同伦于零。
因此，用 $\mathcal{K}^{m - 1} + \Im(h^m)$ 替换
$\mathcal{K}^{m - 1}$ 后（注意这样做之后，
$\mathcal{K}^{m - 1}$ 仍由有限多个整体截面生成），
可见 $\sigma$ 本身同伦于零！这意味着存在下列实线部分交换的图表
$$
\xymatrix{
\mathcal{E}^\bullet \ar[r] &
M \ar[r] &
\mathcal{L}^\bullet \ar[r] &
\mathcal{E}^\bullet[1] \\
\mathcal{E}^\bullet \ar[r] \ar[u] &
K \ar[r] \ar[u] &
\sigma_{\leq m - 1}\mathcal{K}^\bullet \ar[r] \ar[u] &
\mathcal{E}^\bullet[1] \ar[u] \\
\mathcal{E}^\bullet \ar[r] \ar[u] &
M \ar[r] \ar[u]^s &
\mathcal{L}^\bullet \ar[r] \ar@{..>}[u] &
\mathcal{E}^\bullet[1] \ar[u]
}
$$
。由三角范畴的公理，得到一条使图表成立的虚线箭头。
考察次数 $m - 1$ 上的上同调层，得到
$$
\xymatrix{
H^{m - 1}(M) \ar[r] &
H^{m - 1}(\mathcal{L}^\bullet) \ar[r] &
H^m(\mathcal{E}^\bullet) \\
H^{m - 1}(K) \ar[r] \ar[u] &
H^{m - 1}(\sigma_{\leq m - 1}\mathcal{K}^\bullet) \ar[r] \ar[u] &
H^m(\mathcal{E}^\bullet) \ar[u] \\
H^{m - 1}(M) \ar[r] \ar[u] &
H^{m - 1}(\mathcal{L}^\bullet) \ar[r] \ar[u] &
H^m(\mathcal{E}^\bullet) \ar[u]
}
$$
。由于左列和右列中的竖直复合都是恒等映射，故中列的竖直复合
$H^{m - 1}(\mathcal{L}^\bullet) \to
H^{m - 1}(\sigma_{\leq m - 1}\mathcal{K}^\bullet) \to
H^{m - 1}(\mathcal{L}^\bullet)$ 是满射！特别地，
$H^{m - 1}(\sigma_{\leq m - 1}\mathcal{K}^\bullet) \to
H^{m - 1}(\mathcal{L}^\bullet)$ 是满射。
利用上述三角之间的态射所诱导的上同调层长正合序列之间的映射，
追图可见，这蕴含 $H^i(K) \to H^i(M)$ 在 $i \geq m$
时是同构，在 $i = m - 1$ 时是满射。
由构造，可以选取 $r \geq 0$ 和一个满射
$\mathcal{O}^{\oplus r} \to \mathcal{K}^{m - 1}$。于是复合
$$
(\mathcal{O}^{\oplus r} \to \mathcal{E}^m \to
\mathcal{E}^{m + 1} \to \ldots ) \longrightarrow
K \longrightarrow M
$$
在次数 $\geq m$ 的上同调层上诱导同构，在次数 $m - 1$
上诱导满射，证明完毕。
\end{proof}

\begin{lemma}
\label{sites-cohomology-lemma-colim-and-lim-of-duals}
\begin{slogan}
完美对象系的平凡对偶性。
\end{slogan}
设 $(\mathcal{C}, \mathcal{O})$ 是一个环化站点。设
$(K_n)_{n \in \mathbf{N}}$ 是 $D(\mathcal{O})$ 中的一个完美对象系。
设 $K = \text{hocolim} K_n$ 为导出余极限
（《导出范畴》定义 \ref{derived-definition-derived-colimit}）。
那么对 $D(\mathcal{O})$ 的任意对象 $E$，有
$$
R\SheafHom(K, E) = R\lim E \otimes^\mathbf{L}_\mathcal{O} K_n^\vee
$$
，其中 $(K_n^\vee)$ 是由完美复形的对偶组成的逆系统。
\end{lemma}

\begin{proof}
由引理 \ref{sites-cohomology-lemma-dual-perfect-complex}，有
$R\lim E \otimes^\mathbf{L}_\mathcal{O} K_n^\vee =
R\lim R\SheafHom(K_n, E)$
，它嵌入有区别三角
$$
R\lim R\SheafHom(K_n, E) \to
\prod R\SheafHom(K_n, E) \to
\prod R\SheafHom(K_n, E)
$$
。由于 $K$ 类似地嵌入有区别三角
$\bigoplus K_n \to \bigoplus K_n \to K$，只需证明
$\prod R\SheafHom(K_n, E) = R\SheafHom(\bigoplus K_n, E)$。
这是 (\ref{sites-cohomology-equation-internal-hom}) 以及导出张量积与直和可交换这一事实
的形式推论。
\end{proof}






\section{导出范畴中的可逆对象}
\label{sites-cohomology-section-invertible-D-or-R}

\noindent
我们刻画环化站点的导出范畴中的可逆对象（既包括局部环化拓扑斯的情形，
也包括一般情形）。

\begin{lemma}
\label{sites-cohomology-lemma-category-summands-finite-free}
设 $(\mathcal{C}, \mathcal{O})$ 是一个环化站点。令
$R = \Gamma(\mathcal{C}, \mathcal{O})$。作为有限自由
$\mathcal{O}$-模的直和项的那些 $\mathcal{O}$-模所成的范畴，
等价于有限投射 $R$-模所成的范畴。
\end{lemma}

\begin{proof}
注意，有限投射 $R$-模恰好就是有限自由 $R$-模的直和项。
该等价由函子 $\mathcal{E} \mapsto
\Gamma(\mathcal{C}, \mathcal{E})$ 给出。
逆函子由下述构造给出。考虑拓扑斯态射
$f : \Sh(\mathcal{C}) \to \Sh(\text{pt})$，其中 $f_*$
由取整体截面给出，而 $f^{-1}$ 把集合 $S$，即
$\Sh(\text{pt})$ 的对象，映到取值为 $S$ 的常值层。
利用恒等映射 $R \to f_*\mathcal{O}$，得到环化拓扑斯的态射
$(f, f^\sharp) : (\Sh(\mathcal{C}), \mathcal{O}) \to (\Sh(\text{pt}), R)$
。于是逆函子由 $f^*$ 给出。
\end{proof}

\begin{lemma}
\label{sites-cohomology-lemma-invertible-derived}
设 $(\mathcal{C}, \mathcal{O})$ 是一个环化站点。设 $M$ 是
$D(\mathcal{O})$ 的一个对象。下列条件等价：
\begin{enumerate}
\item $M$ 在 $D(\mathcal{O})$ 中可逆；见《范畴》定义
\ref{categories-definition-invertible}；
\item 存在一个局部有限的\footnote{这意味着，对 $\mathcal{C}$
的每个对象 $U$，都存在一个覆盖 $\{U_i \to U\}$，使得对每个 $i$，
层 $\mathcal{O}_n|_{U_i}$ 只对有限多个 $n$ 非零。}直积分解
$$
\mathcal{O} = \prod\nolimits_{n \in \mathbf{Z}} \mathcal{O}_n
$$
，并且对每个 $n$，存在可逆 $\mathcal{O}_n$-模
$\mathcal{H}^n$
（《站点上的模》定义 \ref{sites-modules-definition-invertible-sheaf}），
而且在 $D(\mathcal{O})$ 中有 $M = \bigoplus \mathcal{H}^n[-n]$。
\end{enumerate}
如果 (1) 和 (2) 成立，那么 $M$ 是 $D(\mathcal{O})$ 的完美对象。
如果 $(\mathcal{C}, \mathcal{O})$ 是一个局部环化站点，
这些条件还等价于：
\begin{enumerate}
\item[(3)] 对 $\mathcal{C}$ 的每个对象 $U$，存在一个覆盖
$\{U_i \to U\}$，并且对每个 $i$ 存在整数 $n_i$，使得
$M|_{U_i}$ 由置于次数 $n_i$ 上的可逆
$\mathcal{O}_{U_i}$-模表示。
\end{enumerate}
\end{lemma}

\begin{proof}
假设 (2) 成立。考虑对象 $R\SheafHom(M, \mathcal{O})$ 和复合映射
$$
R\SheafHom(M, \mathcal{O}) \otimes_\mathcal{O}^\mathbf{L} M \to \mathcal{O}
$$
。要证明它是同构，可以局部地工作。因此，可以假设
$\mathcal{O} = \prod_{a \leq n \leq b} \mathcal{O}_n$ 且
$M = \bigoplus_{a \leq n \leq b} \mathcal{H}^n[-n]$。
只需证明
$$
R\SheafHom(\mathcal{H}^m, \mathcal{O})
\otimes_\mathcal{O}^\mathbf{L} \mathcal{H}^n
$$
在 $n \not = m$ 时为零，在 $n = m$ 时等于 $\mathcal{O}_n$。
$n \not = m$ 的情形来自如下事实：$\mathcal{O}_n$ 和
$\mathcal{O}_m$ 都是平坦 $\mathcal{O}$-代数，并且
$\mathcal{O}_n \otimes_\mathcal{O} \mathcal{O}_m = 0$。
利用可逆 $\mathcal{O}$-模的局部结构（《站点上的模》引理
\ref{sites-modules-lemma-invertible}）并局部地工作，
$n = m$ 情形中的同构可直接得出；我们略去细节。
由于 $D(\mathcal{O})$ 是对称幺半范畴，故 $M$ 可逆。

\medskip\noindent
假设 (1) 成立。(2) 中的描述给出 $\mathcal{O}_n$ 的一个候选，即
$\SheafHom_\mathcal{O}(H^n(M), H^n(M))$。
如果这是一个局部有限的环层族，并且
$\mathcal{O} = \prod \mathcal{O}_n$，那么利用来自该直积分解的
$\mathcal{O}$ 中的幂等元，立即得到直和分解
$M = \bigoplus H^n(M)[-n]$。这表明，为证明 (2)，可以在如下意义下
局部地工作。设 $U$ 是 $\mathcal{C}$ 的一个对象。必须证明存在
$U$ 的一个覆盖 $\{U_i \to U\}$，使得取上述 $\mathcal{O}_n$ 时，
限制到 $\mathcal{C}/U_i$ 后，上述各断言以及 (2) 中的断言都成立。
因此，可以假设 $\mathcal{C}$ 有终对象 $X$；在 (2) 的证明过程中，
可以有限多次用 $X$ 的某个覆盖的成员替换 $X$。

\medskip\noindent
选取 $D(\mathcal{O})$ 的一个对象 $N$ 以及一个同构
$M \otimes_\mathcal{O}^\mathbf{L} N \cong \mathcal{O}$。
于是 $N$ 是幺半范畴 $D(\mathcal{O})$ 中 $M$ 的左对偶；
由引理 \ref{sites-cohomology-lemma-left-dual-derived}，$M$ 是完美的。
由对称性，$N$ 也是完美的。用一个覆盖的成员替换 $X$ 后，
可以假设 $M$ 和 $N$ 分别由严格完美复形
$\mathcal{E}^\bullet$ 和 $\mathcal{F}^\bullet$ 表示。
于是 $M \otimes_\mathcal{O}^\mathbf{L} N$ 由
$\text{Tot}(\mathcal{E}^\bullet \otimes_\mathcal{O} \mathcal{F}^\bullet)$
表示。用 $X$ 的一个覆盖的成员替换 $X$ 后，
可以假设互逆同构
$\mathcal{O} \to M \otimes_\mathcal{O}^\mathbf{L} N$ 和
$M \otimes_\mathcal{O}^\mathbf{L} N \to \mathcal{O}$ 由复形映射
$$
\alpha : \mathcal{O} \to
\text{Tot}(\mathcal{E}^\bullet \otimes_\mathcal{O} \mathcal{F}^\bullet)
\quad\text{和}\quad
\beta :
\text{Tot}(\mathcal{E}^\bullet \otimes_\mathcal{O} \mathcal{F}^\bullet)
\to \mathcal{O}
$$
给出；见引理 \ref{sites-cohomology-lemma-local-actual}。于是，作为复形映射有
$\beta \circ \alpha = 1$，而作为 $D(\mathcal{O})$ 中的态射有
$\alpha \circ \beta = 1$。用 $X$ 的一个覆盖的成员替换 $X$ 后，
可以假设复合 $\alpha \circ \beta$ 通过某个同伦 $\theta$
同伦于 $1$，其分量为
$$
\theta^n :
\text{Tot}^n(\mathcal{E}^\bullet \otimes_\mathcal{O} \mathcal{F}^\bullet)
\to
\text{Tot}^{n - 1}(
\mathcal{E}^\bullet \otimes_\mathcal{O} \mathcal{F}^\bullet)
$$
；这是由先前同一个引理得到的。令
$R = \Gamma(\mathcal{C}, \mathcal{O})$。由引理
\ref{sites-cohomology-lemma-category-summands-finite-free}，得到：
\begin{enumerate}
\item $M^\bullet = \Gamma(X, \mathcal{E}^\bullet)$
是有限投射 $R$-模的有界复形；
\item $N^\bullet = \Gamma(X, \mathcal{F}^\bullet)$
是有限投射 $R$-模的有界复形；
\item $\alpha$ 和 $\beta$ 分别对应于复形映射
$a : R \to \text{Tot}(M^\bullet \otimes_R N^\bullet)$ 和
$b : \text{Tot}(M^\bullet \otimes_R N^\bullet) \to R$；
\item $\theta^n$ 对应于映射
$h^n : \text{Tot}^n(M^\bullet \otimes_R N^\bullet) \to
\text{Tot}^{n - 1}(M^\bullet \otimes_R N^\bullet)$；以及
\item $b \circ a = 1$ 且 $b \circ a - 1 = dh + hd$。
\end{enumerate}
由此，$M^\bullet$ 和 $N^\bullet$ 定义 $D(R)$ 中互逆的对象。
由《代数详论》引理 \ref{more-algebra-lemma-invertible-derived}，
得到直积分解 $R = \prod_{a \leq n \leq b} R_n$ 以及可逆
$R_n$-模 $H^n$，使得
$M^\bullet \cong \bigoplus_{a \leq n \leq b} H^n[-n]$。
$D(R)$ 中的这个同构可以提升为复形态射
$$
\bigoplus H^n[-n] \longrightarrow M^\bullet
$$
，因为每个 $H^n$ 作为 $R$-模都是投射的。相应地，
再次使用引理 \ref{sites-cohomology-lemma-category-summands-finite-free}，得到态射
$$
\bigoplus H^n \otimes_R \mathcal{O}[-n] \to \mathcal{E}^\bullet
$$
，它在 $D(\mathcal{O})$ 中是同构。这里 $M \otimes_R \mathcal{O}$
表示引理 \ref{sites-cohomology-lemma-category-summands-finite-free} 的证明中构造的
从有限投射 $R$-模到 $\mathcal{O}$-模的函子。
令 $\mathcal{O}_n = R_n \otimes_R \mathcal{O}$，便得出 (2) 成立。

\medskip\noindent
如果 $(\mathcal{C}, \mathcal{O})$ 是一个局部环化站点，
那么给定对象 $U$ 和有限直积分解
$\mathcal{O}|_U = \prod_{a \leq n \leq b} \mathcal{O}_n|_U$，
可以找到覆盖 $\{U_i \to U\}$，使得对每个 $i$，至多有一个 $n$
满足 $\mathcal{O}_n|_{U_i}$ 非零。这很容易由《站点上的模》引理
\ref{sites-modules-lemma-locally-ringed} 的 (2) 以及《站点上的模》定义
\ref{sites-modules-definition-locally-ringed} 中给出的局部环化站点定义
得出。由此易见 (2) $\Rightarrow$ (3)。不对环化站点作任何假设时，
(3) $\Rightarrow$ (2) 也成立。我们略去细节。
\end{proof}










\section{投影公式}
\label{sites-cohomology-section-projection-formula}

\noindent
设 $f : (\Sh(\mathcal{C}), \mathcal{O}_\mathcal{C}) \to
(\Sh(\mathcal{D}), \mathcal{O}_\mathcal{D})$ 是环化拓扑斯的一个态射。
设 $E \in D(\mathcal{O}_\mathcal{C})$ 且
$K \in D(\mathcal{O}_\mathcal{D})$。不作任何进一步假设，存在映射
\begin{equation}
\label{sites-cohomology-equation-projection-formula-map}
Rf_*E \otimes^\mathbf{L}_{\mathcal{O}_\mathcal{D}} K
\longrightarrow
Rf_*(E \otimes^\mathbf{L}_{\mathcal{O}_\mathcal{C}} Lf^*K)
\end{equation}
具体而言，它是典范映射
$$
Lf^*(Rf_*E \otimes^\mathbf{L}_{\mathcal{O}_\mathcal{D}} K) =
Lf^*Rf_*E \otimes^\mathbf{L}_{\mathcal{O}_\mathcal{C}} Lf^*K
\longrightarrow
E \otimes^\mathbf{L}_{\mathcal{O}_\mathcal{C}} Lf^*K
$$
的伴随映射；该典范映射来自映射 $Lf^*Rf_*E \to E$ 以及引理
\ref{sites-cohomology-lemma-pullback-tensor-product} 和 \ref{sites-cohomology-lemma-adjoint}。
下面是一个颇为一般的投影公式版本。

\begin{lemma}
\label{sites-cohomology-lemma-projection-formula}
设 $f : (\Sh(\mathcal{C}), \mathcal{O}_\mathcal{C}) \to
(\Sh(\mathcal{D}), \mathcal{O}_\mathcal{D})$ 是环化拓扑斯的一个态射。
设 $E \in D(\mathcal{O}_\mathcal{C})$ 且
$K \in D(\mathcal{O}_\mathcal{D})$。如果 $K$ 是完美的，那么
$$
Rf_*E \otimes^\mathbf{L}_{\mathcal{O}_\mathcal{D}} K =
Rf_*(E \otimes^\mathbf{L}_{\mathcal{O}_\mathcal{C}} Lf^*K)
$$
在 $D(\mathcal{O}_\mathcal{D})$ 中成立。
\end{lemma}

\begin{proof}
要验证 (\ref{sites-cohomology-equation-projection-formula-map}) 是同构，可以在
$\mathcal{D}$ 上局部地工作。也就是说，对 $\mathcal{D}$ 的任意对象
$V$，必须找到一个覆盖 $\{V_j \to V\}$，使得该映射在 $V_j$
上的限制是同构。由完美对象的定义，这意味着可以假设 $K$
由一个严格完美的 $\mathcal{O}_\mathcal{D}$-模复形表示。
注意，完全一般地，该断言对 $K = K_1 \oplus K_2$ 成立，当且仅当
它对 $K_1$ 和 $K_2$ 都成立。因此，可以假设 $K$
是由有限自由 $\mathcal{O}_\mathcal{D}$-模组成的有限复形。
此时，一个涉及朴素截断的简单论证把该断言归约到 $K$
由有限自由 $\mathcal{O}_\mathcal{D}$-模表示的情形。
由于该断言在变量 $K$ 的有限直和项下不变，故只需对
$K = \mathcal{O}_\mathcal{D}[n]$ 证明它；在这种情形下它是平凡的。
\end{proof}

\begin{remark}
\label{sites-cohomology-remark-compatible-with-diagram}
映射 (\ref{sites-cohomology-equation-projection-formula-map}) 在下述意义下与注记
\ref{sites-cohomology-remark-base-change} 的基变换映射相容。具体而言，假设
$$
\xymatrix{
(\Sh(\mathcal{C}'), \mathcal{O}_{\mathcal{C}'})
\ar[r]_{g'} \ar[d]_{f'} &
(\Sh(\mathcal{C}), \mathcal{O}_\mathcal{C}) \ar[d]^f \\
(\Sh(\mathcal{D}'), \mathcal{O}_{\mathcal{D}'})
\ar[r]^g &
(\Sh(\mathcal{D}), \mathcal{O}_\mathcal{D})
}
$$
是环化拓扑斯的交换图。设 $E \in D(\mathcal{O}_\mathcal{C})$
且 $K \in D(\mathcal{O}_\mathcal{D})$。则图表
$$
\xymatrix{
Lg^*(Rf_*E \otimes^\mathbf{L}_{\mathcal{O}_\mathcal{D}} K)
\ar[r]_p \ar[d]_t &
Lg^*Rf_*(E \otimes^\mathbf{L}_{\mathcal{O}_\mathcal{C}} Lf^*K)
\ar[d]_b \\
Lg^*Rf_*E \otimes^\mathbf{L}_{\mathcal{O}_{\mathcal{D}'}}
Lg^*K \ar[d]_b &
Rf'_*L(g')^*(E \otimes^\mathbf{L}_{\mathcal{O}_\mathcal{C}}
Lf^*K) \ar[d]_t \\
Rf'_*L(g')^*E \otimes^\mathbf{L}_{\mathcal{O}_{\mathcal{D}'}}
Lg^*K \ar[rd]_p &
Rf'_*(L(g')^*E \otimes^\mathbf{L}_{\mathcal{O}_{\mathcal{D}'}}
L(g')^*Lf^*K) \ar[d]_c \\
& Rf'_*(L(g')^*E \otimes^\mathbf{L}_{\mathcal{O}_{\mathcal{D}'}}
L(f')^*Lg^*K)
}
$$
交换。这里，标记为 $t$ 的箭头由应用引理
\ref{sites-cohomology-lemma-pullback-tensor-product} 得到，标记为 $b$ 的箭头由应用注记
\ref{sites-cohomology-remark-base-change} 得到，标记为 $p$ 的箭头由应用
(\ref{sites-cohomology-equation-projection-formula-map}) 得到，而 $c$ 来自
$L(g')^* \circ Lf^* = L(f')^* \circ Lg^*$。我们略去验证。
\end{remark}






\section{弱可缩对象}
\label{sites-cohomology-section-w-contractible}

\noindent
站点的对象 $U$ 称为{\it 弱可缩的}，如果集合层的每个满射
$\mathcal{F} \to \mathcal{G}$ 都诱导满射
$\mathcal{F}(U) \to \mathcal{G}(U)$；见《站点》定义
\ref{sites-definition-w-contractible}。

\begin{lemma}
\label{sites-cohomology-lemma-w-contractible}
设 $\mathcal{C}$ 是一个站点。设 $U$ 是 $\mathcal{C}$
的一个弱可缩对象。则：
\begin{enumerate}
\item 函子 $\mathcal{F} \mapsto \mathcal{F}(U)$ 是正合函子
$\textit{Ab}(\mathcal{C}) \to \textit{Ab}$；
\item $H^p(U, \mathcal{F}) = 0$
对每个阿贝尔层 $\mathcal{F}$ 以及所有 $p \geq 1$ 成立；以及
\item 对任意群层 $\mathcal{G}$，每个 $\mathcal{G}$-挠子
在 $U$ 上都有截面。
\end{enumerate}
\end{lemma}

\begin{proof}
第一个断言立即由定义得出（另见《同调》第
\ref{homology-section-functors} 节）。由《导出范畴》引理
\ref{derived-lemma-right-derived-exact-functor}，高阶导出函子消失。
设 $\mathcal{F}$ 是一个 $\mathcal{G}$-挠子。则
$\mathcal{F} \to *$ 是层的满射。因此，(3) 也由定义得出。
\end{proof}

\noindent
在此列出具有足够多弱可缩对象的一些推论会很方便。

\begin{proposition}
\label{sites-cohomology-proposition-enough-weakly-contractibles}
设 $\mathcal{C}$ 是一个站点。设 $\mathcal{B} \subset \Ob(\mathcal{C})$，
使得每个 $U \in \mathcal{B}$ 都是弱可缩的，并且 $\mathcal{C}$
的每个对象都有一个由 $\mathcal{B}$ 的元素组成的覆盖。
设 $\mathcal{O}$ 是 $\mathcal{C}$ 上的一个环层。则：
\begin{enumerate}
\item $\mathcal{O}$-模复形
$\mathcal{F}_1 \to \mathcal{F}_2 \to \mathcal{F}_3$ 正合，
当且仅当对所有 $U \in \mathcal{B}$，
$\mathcal{F}_1(U) \to \mathcal{F}_2(U) \to \mathcal{F}_3(U)$
都正合。
\item $D(\mathcal{O})$ 的每个对象 $K$ 都是其典范截断的导出极限：
$K = R\lim \tau_{\geq -n} K$。
\item 给定逆系统
$\ldots \to \mathcal{F}_3 \to \mathcal{F}_2 \to \mathcal{F}_1$
，若其过渡映射均为满射，则投影
$\lim \mathcal{F}_n \to \mathcal{F}_1$ 是满射。
\item $\textit{Mod}(\mathcal{O})$ 中的乘积是正合的。
\item $D(\mathcal{O})$ 中的乘积可以通过对任意表示复形取乘积来计算。
\item 如果 $(\mathcal{F}_n)$ 是一个 $\mathcal{O}$-模逆系统，
那么对所有 $p > 1$ 有 $R^p\lim \mathcal{F}_n = 0$，并且
$$
R^1\lim \mathcal{F}_n  =
\Coker(\prod \mathcal{F}_n \to \prod \mathcal{F}_n)
$$
，其中映射为 $(x_n) \mapsto (x_n - f(x_{n + 1}))$。
\item 如果 $(K_n)$ 是 $D(\mathcal{O})$ 的对象的一个逆系统，
那么存在短正合列
$$
0 \to R^1\lim H^{p - 1}(K_n) \to H^p(R\lim K_n) \to
\lim H^p(K_n) \to 0
$$
\end{enumerate}
\end{proposition}

\begin{proof}
证明 (1)。如果该序列正合，那么由引理 \ref{sites-cohomology-lemma-w-contractible}，
在 $\mathcal{C}$ 的任意弱可缩元素上取值都会得到正合序列。
反之，假设对所有 $U \in \mathcal{B}$，
$\mathcal{F}_1(U) \to \mathcal{F}_2(U) \to \mathcal{F}_3(U)$
都正合。设 $V$ 是 $\mathcal{C}$ 的一个对象，并设
$s \in \mathcal{F}_2(V)$ 是 $\mathcal{F}_2 \to \mathcal{F}_3$
的核中的元素。由假设，存在一个满足 $U_i \in \mathcal{B}$
的覆盖 $\{U_i \to V\}$。于是 $s|_{U_i}$ 可以提升为截面
$s_i \in \mathcal{F}_1(U_i)$。因此，$s$ 是像层
$\Im(\mathcal{F}_1 \to \mathcal{F}_2)$ 的一个截面。
换言之，序列
$\mathcal{F}_1 \to \mathcal{F}_2 \to \mathcal{F}_3$
正合。

\medskip\noindent
证明 (2)。这由引理 \ref{sites-cohomology-lemma-is-limit-dimension} 中 $d = 0$
的情形成立。

\medskip\noindent
证明 (3)。设 $(\mathcal{F}_n)$ 是如 (2) 中的系统，并令
$\mathcal{F} = \lim \mathcal{F}_n$。如果 $U \in \mathcal{B}$，那么
$\mathcal{F}(U) = \lim \mathcal{F}_n(U)$
满射到 $\mathcal{F}_1(U)$，因为所有过渡映射
$\mathcal{F}_{n + 1}(U) \to \mathcal{F}_n(U)$ 都是满射。
由《站点》定义 \ref{sites-definition-sheaves-injective-surjective}
以及每个对象都有一个由 $\mathcal{B}$ 的元素组成的覆盖这一假设，
$\mathcal{F} \to \mathcal{F}_1$ 是满射。

\medskip\noindent
证明 (4)。设
$\mathcal{F}_{i, 1} \to \mathcal{F}_{i, 2} \to \mathcal{F}_{i, 3}$
是一个 $\mathcal{O}$-模正合列族。要证明
$\prod \mathcal{F}_{i, 1} \to \prod \mathcal{F}_{i, 2} \to
\prod \mathcal{F}_{i, 3}$ 正合。使用 (1) 的判据。
设 $U \in \mathcal{B}$。则
$$
(\prod \mathcal{F}_{i, 1})(U) \to
(\prod \mathcal{F}_{i, 2})(U) \to
(\prod \mathcal{F}_{i, 3})(U)
$$
与
$$
\prod \mathcal{F}_{i, 1}(U) \to
\prod \mathcal{F}_{i, 2}(U) \to
\prod \mathcal{F}_{i, 3}(U)
$$
相同。每个序列
$\mathcal{F}_{i, 1}(U) \to \mathcal{F}_{i, 2}(U) \to \mathcal{F}_{i, 3}(U)$
均由 (1) 正合。因此，由《同调》引理
\ref{homology-lemma-product-abelian-groups-exact}，所显示的序列正合。
再次由 (1) 得出结论。

\medskip\noindent
证明 (5)。这由 (4) 以及《导出范畴》引理
\ref{derived-lemma-products} 的稍加推广版本得出。

\medskip\noindent
证明 (6) 和 (7)。关于导出极限、同伦极限及其关系的讨论，
参见第 \ref{sites-cohomology-section-derived-limits} 节。由《导出范畴》定义
\ref{derived-definition-derived-limit}，有有区别三角
$$
R\lim K_n \to \prod K_n \to \prod K_n \to R\lim K_n[1]
$$
。取上同调层的长正合序列，得到
$$
H^{p - 1}(\prod K_n) \to H^{p - 1}(\prod K_n) \to
H^p(R\lim K_n) \to H^p(\prod K_n) \to H^p(\prod K_n)
$$
。由于乘积由 (4) 正合，这变为
$$
\prod H^{p - 1}(K_n) \to \prod H^{p - 1}(K_n) \to
H^p(R\lim K_n) \to \prod H^p(K_n) \to \prod H^p(K_n)
$$
。先把它应用于 $K_n = \mathcal{F}_n[0]$ 的情形，
其中 $(\mathcal{F}_n)$ 如 (6) 中所述，便得出 (6) 成立。
再把它应用于 (7) 中的 $(K_n)$，便得出 (7) 成立。
\end{proof}






\section{紧对象}
\label{sites-cohomology-section-compact}

\noindent
本节研究环化站点上的模的导出范畴中的紧对象。回忆紧对象定义于
《导出范畴》定义 \ref{derived-definition-compact-object}。

\begin{lemma}
\label{sites-cohomology-lemma-compact-in-terms-of-generators}
设 $\mathcal{A}$ 是一个 Grothendieck 阿贝尔范畴。设
$S \subset \Ob(\mathcal{A})$ 是一个对象集，满足：
\begin{enumerate}
\item $\mathcal{A}$ 的任意对象都是 $S$ 中元素的某个直和的商；以及
\item 对任意 $E \in S$，函子 $\Hom_\mathcal{A}(E, -)$ 与直和可交换。
\end{enumerate}
那么，$D(\mathcal{A})$ 的每个紧对象在 $D(\mathcal{A})$ 中都是
一个有限复形的直和项，而该复形的每一项都是 $S$ 中元素的有限直和。
\end{lemma}

\begin{proof}
假设 $K \in D(\mathcal{A})$ 是一个紧对象。用复形 $K^\bullet$
表示 $K$，并考虑映射
$$
K^\bullet
\longrightarrow
\bigoplus\nolimits_{n \geq 0} \tau_{\geq n} K^\bullet
$$
，其中使用了典范截断；见《同调》第
\ref{homology-section-truncations} 节。这是有意义的，因为右边的直和
在每个次数上都是有限的。由假设，该映射分解经过一个有限直和。
因此，至少对一个 $n$，映射 $K \to \tau_{\geq n} K$ 为零，
即 $K$ 属于 $D^{-}(R)$。

\medskip\noindent
可以用一个上有界复形 $K^\bullet$ 表示 $K$，其每一项都是 $S$
中对象的直和；见《导出范畴》引理
\ref{derived-lemma-subcategory-left-resolution}。注意有
$$
K^\bullet = \bigcup\nolimits_{n \leq 0} \sigma_{\geq n}K^\bullet
$$
，其中使用了朴素截断；见《同调》第
\ref{homology-section-truncations} 节。因此，由《导出范畴》引理
\ref{derived-lemma-colim-hocolim} 和
\ref{derived-lemma-commutes-with-countable-sums}，可见在 $D(R)$ 中，
$1 : K^\bullet \to K^\bullet$ 分解经过
$\sigma_{\geq n}K^\bullet \to K^\bullet$。因此，
$1 : K^\bullet \to K^\bullet$ 在 $D(\mathcal{A})$ 中分解为
$$
K^\bullet \xrightarrow{\varphi} L^\bullet \xrightarrow{\psi} K^\bullet
$$
，其中复形 $L^\bullet$ 有界，且各项都是 $S$ 中元素的直和。
设 $L^i$ 对 $i \not \in [a, b]$ 为零。令 $c$ 为满足如下条件的
最大整数 $\leq b + 1$：对 $i < c$，$L^i$ 是 $S$
中元素的有限直和。断言：如果 $c < b + 1$，那么可以修改
$L^\bullet$ 以增大 $c$。对此断言作归纳，即可表明 $1_K$ 有分解
$$
K \xrightarrow{\varphi} L \xrightarrow{\psi} K
$$
，其中在 $D(\mathcal{A})$ 中，$L$ 可以由一个有限复形表示，
且其各项都是 $S$ 中元素的有限直和。注意
$e = \varphi \circ \psi \in \text{End}_{D(\mathcal{A})}(L)$
是一个幂等元。由《导出范畴》引理
\ref{derived-lemma-projectors-have-images-triangulated}，有
$L = \Ker(e) \oplus \Ker(1 - e)$。映射
$\varphi : K \to L$ 在 $D(R)$ 中诱导与 $\Ker(1 - e)$
的同构，结论得证。

\medskip\noindent
证明该断言。写成 $L^c = \bigoplus_{\lambda \in \Lambda} E_\lambda$。
由于 $L^{c - 1}$ 是 $S$ 中元素的有限直和，由假设 (2)，
可以找到有限子集 $\Lambda' \subset \Lambda$，使得
$L^{c - 1} \to L^c$ 分解经过
$\bigoplus_{\lambda \in \Lambda'} E_\lambda \subset L^c$。
考虑复形映射
$$
\pi :
L^\bullet
\longrightarrow
(\bigoplus\nolimits_{\lambda \in \Lambda \setminus \Lambda'} E_\lambda)[-i]
$$
，它在次数 $i$ 上由到与 $\Lambda \setminus \Lambda'$ 对应的因子的
投影给出。由关于 $K$ 的假设可见，必要时用一个更大的有限子集替换
$\Lambda'$ 后，可以假设 $D(\mathcal{A})$ 中
$\pi \circ \varphi = 0$。设
$(L')^\bullet \subset L^\bullet$ 为 $\pi$ 的核。
由于 $\pi$ 是满射，得到一个复形的短正合列，它在
$D(\mathcal{A})$ 中给出一个有区别三角（见《导出范畴》引理
\ref{derived-lemma-derived-canonical-delta-functor}）。
由于 $\Hom_{D(\mathcal{A})}(K, -)$ 是同调函子（见《导出范畴》引理
\ref{derived-lemma-representable-homological}）且
$\pi \circ \varphi = 0$，可以找到 $D(\mathcal{A})$ 中的态射
$\varphi' : K^\bullet \to (L')^\bullet$，使得它与
$(L')^\bullet \to L^\bullet$ 的复合等于 $\varphi$。
令 $\psi'$ 等于 $\psi$ 与 $(L')^\bullet \to L^\bullet$ 的复合，
便得到一个新的分解。由于 $(L')^\bullet$ 除次数 $c$ 外与
$L^\bullet$ 相同，并且
$(L')^c = \bigoplus_{\lambda \in \Lambda'} E_\lambda$，断言得证。
\end{proof}

\begin{lemma}
\label{sites-cohomology-lemma-compact-objects-if-enough-qc}
设 $(\mathcal{C}, \mathcal{O})$ 是一个环化站点。假设
$\mathcal{C}$ 的每个对象都有一个由拟紧对象组成的覆盖。
那么 $D(\mathcal{O})$ 的每个紧对象在 $D(\mathcal{O})$ 中都是
某个有限复形的直和项；该复形的各项都是形如
$j_!\mathcal{O}_U$ 的 $\mathcal{O}$-模的有限直和，
其中 $U$ 是 $\mathcal{C}$ 的一个拟紧对象。
\end{lemma}

\begin{proof}
应用引理 \ref{sites-cohomology-lemma-compact-in-terms-of-generators}，其中
$S \subset \Ob(\textit{Mod}(\mathcal{O}))$ 是所有形如
$j_!\mathcal{O}_U$ 的模所成的集合，且
$U \in \Ob(\mathcal{C})$ 拟紧。由《站点上的模》引理
\ref{sites-modules-lemma-module-quotient-flat} 以及每个 $U$
都可以被拟紧对象覆盖这一假设，假设 (1) 成立。假设 (2) 来自
$$
\Hom_\mathcal{O}(j_!\mathcal{O}_U, \mathcal{F}) = \mathcal{F}(U)
$$
，它由《站点》引理 \ref{sites-lemma-directed-colimits-sections}
与直和可交换。
\end{proof}

\noindent
在上述引理的情形中，模 $j_!\mathcal{O}_U$ 不一定是
$D(\mathcal{O})$ 的紧对象（即使 $U$ 是 $\mathcal{C}$
的拟紧对象也是如此）。下面两个引理处理这一问题。


\begin{lemma}
\label{sites-cohomology-lemma-when-jshriek-lower-compact}
设 $(\mathcal{C}, \mathcal{O})$ 是一个环化站点。设 $U$ 是
$\mathcal{C}$ 的一个对象。假设函子
$\mathcal{F} \mapsto H^p(U, \mathcal{F})$ 与直和可交换。
那么 $\mathcal{O}$-模 $j_!\mathcal{O}_U$ 在下述意义下是
$D^+(\mathcal{O})$ 的紧对象：如果 $D(\mathcal{O})$ 中的
$M = \bigoplus_{i \in I} M_i$ 下有界，那么
$\Hom(j_{U!}\mathcal{O}_U, M) =
\bigoplus_{i \in I} \Hom(j_{U!}\mathcal{O}_U, M_i)$。
\end{lemma}

\begin{proof}
由《站点上的模》公式
(\ref{sites-modules-equation-map-lower-shriek-OU-into-module})，
$\Hom(j_{U!}\mathcal{O}_U, -)$ 与函子
$\mathcal{F} \mapsto \mathcal{F}(U)$ 相同，因此只需证明
$H^p(U, M) = \bigoplus H^p(U, M_i)$。
设 $\mathcal{I}_i$（$i \in I$）是一个内射
$\mathcal{O}$-模族。由假设，有
$$
H^p(U, \bigoplus\nolimits_{i \in I} \mathcal{I}_i) =
\bigoplus\nolimits_{i \in I} H^p(U, \mathcal{I}_i) = 0
$$
对所有 $p$ 成立。由于 $M = \bigoplus M_i$ 下有界，
可见存在 $a \in \mathbf{Z}$，使得对 $n < a$ 有
$H^n(M_i) = 0$。因此，可以选取表示 $M_i$ 的内射
$\mathcal{O}$-模复形 $\mathcal{I}_i^\bullet$，使得对 $n < a$ 有
$\mathcal{I}_i^n = 0$；见《导出范畴》引理
\ref{derived-lemma-injective-resolutions-exist}。
由《内射对象》引理 \ref{injectives-lemma-derived-products}，
直和复形 $\bigoplus \mathcal{I}_i^\bullet$ 表示 $M$。
由 Leray 非循环性（《导出范畴》引理
\ref{derived-lemma-leray-acyclicity}），有
$$
R\Gamma(U, M) = \Gamma(U, \bigoplus \mathcal{I}_i^\bullet) =
\bigoplus \Gamma(U, \bigoplus \mathcal{I}_i^\bullet) =
\bigoplus R\Gamma(U, M_i)
$$
，正如所需。
\end{proof}

\begin{lemma}
\label{sites-cohomology-lemma-when-jshriek-lower-compact-worked-out}
设 $(\mathcal{C}, \mathcal{O})$ 是一个环化站点，
其覆盖集为 $\text{Cov}_\mathcal{C}$。设
$\mathcal{B} \subset \Ob(\mathcal{C})$ 且
$\text{Cov} \subset \text{Cov}_\mathcal{C}$ 为子集。假设：
\begin{enumerate}
\item 对每个 $\mathcal{U} \in \text{Cov}$，有
$\mathcal{U} = \{U_i \to U\}_{i \in I}$，其中 $I$ 有限，
$U, U_i \in \mathcal{B}$，并且每个
$U_{i_0} \times_U \ldots \times_U U_{i_p} \in \mathcal{B}$；
\item 对每个 $U \in \mathcal{B}$，$\text{Cov}$ 中出现的
$U$ 的覆盖构成 $U$ 的覆盖的一个共尾系统。
\end{enumerate}
那么，对 $U \in \mathcal{B}$，对象 $j_{U!}\mathcal{O}_U$
在下述意义下是 $D^+(\mathcal{O})$ 的紧对象：如果
$D(\mathcal{O})$ 中的 $M = \bigoplus_{i \in I} M_i$ 下有界，那么
$\Hom(j_{U!}\mathcal{O}_U, M) =
\bigoplus_{i \in I} \Hom(j_{U!}\mathcal{O}_U, M_i)$。
\end{lemma}

\begin{proof}
这由引理 \ref{sites-cohomology-lemma-when-jshriek-lower-compact} 和引理
\ref{sites-cohomology-lemma-colim-works-over-collection} 得出。
\end{proof}

\begin{lemma}
\label{sites-cohomology-lemma-when-jshriek-compact}
设 $(\mathcal{C}, \mathcal{O})$ 是一个环化站点。设 $U$ 是
$\mathcal{C}$ 的一个对象。如果存在整数 $d$ 使得：
\begin{enumerate}
\item 对所有 $p > d$，有 $H^p(U, \mathcal{F}) = 0$；以及
\item 函子 $\mathcal{F} \mapsto H^p(U, \mathcal{F})$ 与直和可交换，
\end{enumerate}
那么 $\mathcal{O}$-模 $j_!\mathcal{O}_U$ 是
$D(\mathcal{O})$ 的紧对象。
\end{lemma}

\begin{proof}
假设 (1) 和 (2)。回忆，对 $D(\mathcal{O})$ 中的 $K$，有
$\Hom(j_!\mathcal{O}_U, K) = R\Gamma(U, K)$。因此必须证明
$R\Gamma(U, -)$ 与直和可交换。第一个假设意味着函子
$F = H^0(U, -)$ 具有有限上同调维数。此外，第二个假设蕴含
内射模的任意直和对 $F$ 都是非循环的。设 $K_i$ 是
$D(\mathcal{O})$ 的一个对象族。选取表示 $K_i$ 且各项内射的
K-内射代表元 $I_i^\bullet$；见《内射对象》定理
\ref{injectives-theorem-K-injective-embedding-grothendieck}。
由于可以把 $F$ 应用于任意由非循环对象组成的复形来计算 $RF$
（《导出范畴》引理 \ref{derived-lemma-unbounded-right-derived}），
并且 $\bigoplus K_i$ 由 $\bigoplus I_i^\bullet$ 表示
（《内射对象》引理 \ref{injectives-lemma-derived-products}），
故 $R\Gamma(U, \bigoplus K_i)$ 由
$\bigoplus H^0(U, I_i^\bullet)$ 表示。因此，$R\Gamma(U, -)$
如所需地与直和可交换。
\end{proof}

\begin{lemma}
\label{sites-cohomology-lemma-quasi-compact-weakly-contractible-compact}
设 $(\mathcal{C}, \mathcal{O})$ 是一个环化站点。设 $U$
是 $\mathcal{C}$ 的一个拟紧且弱可缩的对象。那么
$j_!\mathcal{O}_U$ 是 $D(\mathcal{O})$ 的紧对象。
\end{lemma}

\begin{proof}
把引理 \ref{sites-cohomology-lemma-when-jshriek-compact} 和
\ref{sites-cohomology-lemma-w-contractible} 与《站点上的模》引理
\ref{sites-modules-lemma-sections-over-quasi-compact} 结合即可。
\end{proof}

\begin{lemma}
\label{sites-cohomology-lemma-perfect-is-compact}
设 $(\mathcal{C}, \mathcal{O})$ 是一个环化站点。假设
$\mathcal{C}$ 具有下列性质：
\begin{enumerate}
\item $\mathcal{C}$ 有一个拟紧终对象 $X$；
\item $\mathcal{C}$ 的每个拟紧对象都有一个由有限覆盖组成的
共尾覆盖系，而且这些覆盖均由拟紧对象组成；
\item 对于 $U$ 和 $U_i$ 均拟紧的有限覆盖
$\{U_i \to U\}_{i \in I}$，纤维积 $U_i \times_U U_j$ 均拟紧。
\end{enumerate}
设 $K$ 是 $D(\mathcal{O})$ 的一个完美对象。则：
\begin{enumerate}
\item[(a)] $K$ 在下述意义下是 $D^+(\mathcal{O})$ 的紧对象：
如果 $M = \bigoplus_{i \in I} M_i$ 下有界，那么
$\Hom(K, M) = \bigoplus_{i \in I} \Hom(K, M_i)$。
\item[(b)] 如果 $(\mathcal{C}, \mathcal{O})$ 具有有限上同调维数，
也就是说，存在 $d$，使得对任意 $\mathcal{O}$-模 $\mathcal{F}$
以及 $i > d$ 都有 $H^i(X, \mathcal{F}) = 0$，那么
$K$ 是 $D(\mathcal{O})$ 的紧对象。
\end{enumerate}
\end{lemma}

\begin{proof}
设 $K^\vee$ 为 $K$ 的对偶；见引理
\ref{sites-cohomology-lemma-dual-perfect-complex}。则有
$$
\Hom_{D(\mathcal{O})}(K, M) =
H^0(X, K^\vee \otimes_\mathcal{O}^\mathbf{L} M)
$$
，并且这关于 $D(\mathcal{O})$ 中的 $M$ 是函子性的。
由于 $K^\vee \otimes_\mathcal{O}^\mathbf{L} -$ 与直和可交换，
只需证明 $R\Gamma(X, -)$ 与有关的直和可交换。

\medskip\noindent
证明 (a)。按上述方式重新表述后，这是引理
\ref{sites-cohomology-lemma-when-jshriek-lower-compact-worked-out} 中 $U = X$
的一个特殊情形。

\medskip\noindent
证明 (b)。由于 $R\Gamma(X, K) = R\Hom(\mathcal{O}, K)$，
并且由引理 \ref{sites-cohomology-lemma-colim-works-over-collection}，
$H^p(X, -)$ 与直和可交换，这是引理
\ref{sites-cohomology-lemma-when-jshriek-compact} 的一个特殊情形。
\end{proof}






\section{上同调层局部常值的复形}
\label{sites-cohomology-section-locally-constant}

\noindent
局部常值层在《站点上的模》第
\ref{sites-modules-section-locally-constant} 节中引入。
设 $\mathcal{C}$ 是一个站点，$\Lambda$ 是一个环。
用 $D(\mathcal{C}, \Lambda)$ 表示 $\mathcal{C}$ 上
$\underline{\Lambda}$-模的阿贝尔范畴的导出范畴。

\begin{lemma}
\label{sites-cohomology-lemma-locally-constant}
设 $\mathcal{C}$ 是一个具有终对象 $X$ 的站点。
设 $\Lambda$ 是一个诺特环。设 $K \in D^b(\mathcal{C}, \Lambda)$，
且 $H^i(K)$ 是有限型 $\Lambda$-模的局部常值层。
那么存在覆盖 $\{U_i \to X\}$，使得每个 $K|_{U_i}$
都由一个有限型 $\Lambda$-模的局部常值层复形表示。
\end{lemma}

\begin{proof}
选取 $a \leq b$，使得对 $i \not \in [a, b]$ 有 $H^i(K) = 0$。
对 $b - a$ 作归纳，证明存在覆盖 $\{U_i \to X\}$，使得
$K|_{U_i}$ 可以由复形 $\underline{M^\bullet}_{U_i}$ 表示，
其中 $M^p$ 是有限型 $\Lambda$-模，且对
$p \not \in [a, b]$ 有 $M^p = 0$。当 $b = a$ 时，这是显然的。
一般地，可以用一个覆盖的成员替换 $X$，并假设 $H^b(K)$
为常值层，记为 $H^b(K) = \underline{M}$。由《站点上的模》引理
\ref{sites-modules-lemma-locally-constant-finite-type}，
模 $M$ 是有限 $\Lambda$-模。选取由 $M$ 的生成元
$x_1, \ldots, x_r$ 给出的满射 $\Lambda^{\oplus r} \to M$。

\medskip\noindent
对引理 \ref{sites-cohomology-lemma-kill-cohomology-class-on-covering} 稍作推广
（细节略去），可得一个覆盖 $\{U_i \to X\}$，使得
$x_i \in H^0(X, H^b(K))$ 可以提升为 $H^b(U_i, K)$ 的一个元素。
因此，用 $U_i$ 替换 $X$ 后，便得到存在映射
$\underline{\Lambda^{\oplus r}}[-b] \to K$
，它在次数 $b$ 的上同调层上诱导满射。选取一个有区别三角
$$
\underline{\Lambda^{\oplus r}}[-b] \to K \to L \to
\underline{\Lambda^{\oplus r}}[-b + 1]
$$
。此时，$L$ 的上同调层只在区间 $[a, b - 1]$ 中非零，
并在区间 $[a, b - 2]$ 中与 $K$ 的上同调层相同；在次数
$b - 1$ 上有短正合列
$$
0 \to H^{b - 1}(K) \to H^{b - 1}(L) \to
\underline{\Ker(\Lambda^{\oplus r} \to M)} \to 0
$$
。由《站点上的模》引理
\ref{sites-modules-lemma-kernel-finite-locally-constant}，
$H^{b - 1}(L)$ 是有限型的局部常值层。由归纳假设，得到
$D(\mathcal{C}, \underline{\Lambda})$ 中的同构
$\underline{M^\bullet} \to L$，其中 $M^p$
是有限 $\Lambda$-模，且对 $p \not \in [a, b - 1]$ 有 $M^p = 0$。
映射 $L \to \Lambda^{\oplus r}[-b + 1]$ 给出映射
$\underline{M^{b - 1}} \to \underline{\Lambda^{\oplus r}}$，
它局部为常值（《站点上的模》引理
\ref{sites-modules-lemma-morphism-locally-constant}）。
因此，可以假设它由映射 $M^{b - 1} \to \Lambda^{\oplus r}$ 给出。
该有区别三角表明复合
$M^{b - 2} \to M^{b - 1} \to \Lambda^{\oplus r}$ 为零，
而三角范畴的公理产生同构
$$
\underline{M^a \to \ldots \to M^{b - 1} \to \Lambda^{\oplus r}}
\longrightarrow K
$$
，它位于 $D(\mathcal{C}, \Lambda)$ 中。
\end{proof}

\noindent
设 $\mathcal{C}$ 是一个站点，$\Lambda$ 是一个环。利用态射
$\Sh(\mathcal{C}) \to \Sh(pt)$，可见存在函子
$D(\Lambda) \to D(\mathcal{C}, \Lambda)$，
$K \mapsto \underline{K}$。

\begin{lemma}
\label{sites-cohomology-lemma-map-out-of-locally-constant}
设 $\mathcal{C}$ 是一个具有终对象 $X$ 的站点。设 $\Lambda$ 是一个环。
给定：
\begin{enumerate}
\item $D(\Lambda)$ 的一个完美对象 $K$；
\item 一个表示 $K$ 的、由有限投射 $\Lambda$-模组成的有限复形
$K^\bullet$；
\item 一个 $\Lambda$-模层复形 $\mathcal{L}^\bullet$；以及
\item $D(\mathcal{C}, \Lambda)$ 中的一个映射
$\varphi : \underline{K} \to \mathcal{L}^\bullet$。
\end{enumerate}
那么存在覆盖 $\{U_i \to X\}$ 和复形映射
$\alpha_i : \underline{K}^\bullet|_{U_i} \to \mathcal{L}^\bullet|_{U_i}$
，后者表示 $\varphi|_{U_i}$。
\end{lemma}

\begin{proof}
立即由引理 \ref{sites-cohomology-lemma-local-actual} 得出。
\end{proof}

\begin{lemma}
\label{sites-cohomology-lemma-locally-constant-map}
设 $\mathcal{C}$ 是一个具有终对象 $X$ 的站点。
设 $\Lambda$ 是一个环。设 $K, L$ 是 $D(\Lambda)$ 的对象，
且 $K$ 完美。设 $\varphi : \underline{K} \to \underline{L}$
是 $D(\mathcal{C}, \Lambda)$ 中的映射。存在覆盖
$\{U_i \to X\}$，使得对某个 $D(\Lambda)$ 中的映射
$\alpha_i : K \to L$，有
$\varphi|_{U_i} = \underline{\alpha_i}$。
\end{lemma}

\begin{proof}
这由引理 \ref{sites-cohomology-lemma-map-out-of-locally-constant} 和《站点上的模》引理
\ref{sites-modules-lemma-morphism-locally-constant} 得出。
\end{proof}

\begin{lemma}
\label{sites-cohomology-lemma-locally-constant-tensor-product}
设 $\mathcal{C}$ 是一个站点。设 $\Lambda$ 是一个诺特环。
设 $K, L \in D^-(\mathcal{C}, \Lambda)$。如果 $K$ 和 $L$
的上同调层都是有限型 $\Lambda$-模的局部常值层，那么
$K \otimes_\Lambda^\mathbf{L} L$
的上同调层也是有限型 $\Lambda$-模的局部常值层。
\end{lemma}

\begin{proof}
把这作为引理 \ref{sites-cohomology-lemma-locally-constant} 的一个应用来证明。
注意，$H^i(K \otimes_\Lambda^\mathbf{L} L)$ 与
$H^i(\tau_{\geq i - 1}K \otimes_\Lambda^\mathbf{L} \tau_{\geq i - 1}L)$
相同。因此，可以假设 $K$ 和 $L$ 有界。由引理
\ref{sites-cohomology-lemma-locally-constant}，可以假设 $K$ 和 $L$
由有限型 $\Lambda$-模的局部常值层复形表示。
然后，可以用有限自由 $\Lambda$-模的上有界复形替换这些复形。
在这种情形下，结论是显然的。
\end{proof}

\begin{lemma}
\label{sites-cohomology-lemma-locally-constant-bounded}
设 $\mathcal{C}$ 是一个站点。设 $\Lambda$ 是一个诺特环。
设 $I \subset \Lambda$ 是一个理想。设
$K \in D^-(\mathcal{C}, \Lambda)$。如果
$K \otimes_\Lambda^\mathbf{L} \underline{\Lambda/I}$ 的上同调层
都是有限型 $\Lambda/I$-模的局部常值层，那么
$K \otimes_\Lambda^\mathbf{L} \underline{\Lambda/I^n}$
的上同调层对所有 $n \geq 1$ 都是有限型
$\Lambda/I^n$-模的局部常值层。
\end{lemma}

\begin{proof}
回忆有限型 $\Lambda$-模的局部常值层构成所有
$\underline{\Lambda}$-模所成范畴的一个弱 Serre 子范畴；见《站点上的模》
引理 \ref{sites-modules-lemma-kernel-finite-locally-constant}。
因此，由上同调层均为有限型 $\Lambda$-模的局部常值层的那些复形
组成的 $D(\mathcal{C}, \Lambda)$ 的子范畴，构成
$D(\mathcal{C}, \Lambda)$ 的一个严格满、饱和三角子范畴；
见《导出范畴》引理
\ref{derived-lemma-cohomology-in-serre-subcategory}。
接着，考虑有区别三角
$$
K \otimes_\Lambda^\mathbf{L} \underline{I^n/I^{n + 1}} \to
K \otimes_\Lambda^\mathbf{L} \underline{\Lambda/I^{n + 1}} \to
K \otimes_\Lambda^\mathbf{L} \underline{\Lambda/I^n} \to
K \otimes_\Lambda^\mathbf{L} \underline{I^n/I^{n + 1}}[1]
$$
以及同构
$$
K \otimes_\Lambda^\mathbf{L} \underline{I^n/I^{n + 1}}
=
\left(K \otimes_\Lambda^\mathbf{L} \underline{\Lambda/I}\right)
\otimes_{\Lambda/I}^\mathbf{L} \underline{I^n/I^{n + 1}}
$$
结合引理 \ref{sites-cohomology-lemma-locally-constant-tensor-product}，便得到结论。
\end{proof}
% Generated mechanically from the frozen dga.tex target.
% Do not edit this generated file; edit the bound locale source instead.

% OK, start here.
%

\setcounter{chapter}{21}
\chapter{微分分次代数}



\phantomsection
\label{dga-section-phantom}


\section{引言}
\label{dga-section-introduction}

\noindent
本章讨论微分分次代数、模、范畴等。一个基本参考文献是
\cite{Keller-Deriving}，综述论文可见 \cite{Keller-survey}。

\medskip\noindent
由于 Stacks project 不在意叙述篇幅，我们先在微分分次模范畴的框架中
发展这些材料。此后，再在微分分次范畴上的微分分次模框架中重新完成
这些构造。



\section{约定}
\label{dga-section-conventions}

\noindent
本章继续采用如下约定：{\it 环}指有 $1$ 的交换环。若 $R$ 是环，
则一个 {\it $R$-代数 $A$} 是一个 $R$-模 $A$，连同一个 $R$-双线性映射
$A \times A \to A$（乘法），使得该乘法结合且有单位元。换言之，
它们是含幺结合 $R$-代数，并且结构映射 $R \to A$ 的像包含在 $A$ 的中心中。

\medskip\noindent
{\bf 符号规则。} 本章将研究通常配备微分的分次代数和分次模。
底层复形上的符号规则总是采用（或兼容于）《更多代数》第
\ref{more-algebra-section-sign-rules} 节中引入的规则。
这有时会使乘法结构以意外方式发生扭曲，特别是在考虑左模，
或者左模与右模之间的关系时。





\section{微分分次代数}
\label{dga-section-dga}


\noindent
这里只给出定义。

\begin{definition}
\label{dga-definition-dga}
设 $R$ 是交换环。一个 {\it $R$ 上的微分分次代数}是以下两者之一：
\begin{enumerate}
\item 一个 $R$-模链复形 $A_\bullet$，配备 $R$-双线性映射
$A_n \times A_m \to A_{n + m}$，$(a, b) \mapsto ab$，使得
$$
\text{d}_{n + m}(ab) = \text{d}_n(a)b + (-1)^n a\text{d}_m(b)
$$
并且 $\bigoplus A_n$ 成为含幺结合 $R$-代数；或者
\item 一个 $R$-模上链复形 $A^\bullet$，配备 $R$-双线性映射
$A^n \times A^m \to A^{n + m}$，$(a, b) \mapsto ab$，使得
$$
\text{d}^{n + m}(ab) = \text{d}^n(a)b + (-1)^n a\text{d}^m(b)
$$
并且 $\bigoplus A^n$ 成为含幺结合 $R$-代数。
\end{enumerate}
\end{definition}

\noindent
我们常简写为 $A = \bigoplus A_n$ 或 $A = \bigoplus A^n$，并把它看作
配备一个 $\mathbf{Z}$-分次及一个平方为零且满足上述 Leibniz 规则的
$R$-线性算子 $\text{d}$ 的含幺结合 $R$-代数。在这种情形下，
我们常说“设 $(A, \text{d})$ 是一个微分分次代数”。

\medskip\noindent
微分分次 $R$-代数 $A$ 上联系微分与乘法的 Leibniz 规则，恰好意味着
乘法映射给出一个上链复形映射
$$
\text{Tot}(A^\bullet \otimes_R A^\bullet) \to A^\bullet
$$
这里 $A^\bullet$ 表示 $A$ 的底层上链复形。

\begin{definition}
\label{dga-definition-homomorphism-dga}
微分分次代数的一个 {\it 同态}
$f : (A, \text{d}) \to (B, \text{d})$，是一个与分次及 $\text{d}$ 相容的
代数映射 $f : A \to B$。
\end{definition}

\begin{definition}
\label{dga-definition-cdga}
若对次数为 $n$ 的 $a$ 和次数为 $m$ 的 $b$ 有
$ab = (-1)^{nm}ba$，则称微分分次代数 $(A, \text{d})$ 是
{\it 交换的}。若此外对 $\deg(a)$ 为奇数的 $a$ 有 $a^2 = 0$，
则称 $A$ 是 {\it 严格交换的}。
\end{definition}

\noindent
下述定义一般都有意义，但或许只有在对交换微分分次代数取张量积时
才是“正确的”。

\begin{definition}
\label{dga-definition-tensor-product}
设 $R$ 是环。设 $(A, \text{d})$、$(B, \text{d})$ 是 $R$ 上的微分分次代数。
$A$ 与 $B$ 的 {\it 张量积微分分次代数}是代数 $A \otimes_R B$，
其乘法定义为
$$
(a \otimes b)(a' \otimes b') = (-1)^{\deg(a')\deg(b)} aa' \otimes bb'
$$
并配备由下述规则定义的微分 $\text{d}$：
$\text{d}(a \otimes b) = \text{d}(a) \otimes b + (-1)^m a \otimes \text{d}(b)$
其中 $m = \deg(a)$。
\end{definition}

\begin{lemma}
\label{dga-lemma-total-complex-tensor-product}
设 $R$ 是环。设 $(A, \text{d})$、$(B, \text{d})$ 是 $R$ 上的微分分次代数。
用 $A^\bullet$、$B^\bullet$ 表示底层上链复形。作为 $R$-模上链复形，我们有
$$
(A \otimes_R B)^\bullet = \text{Tot}(A^\bullet \otimes_R B^\bullet).
$$
\end{lemma}

\begin{proof}
回忆全复形在 $A^p \otimes_R B^q$ 上的微分由
$\text{d}_1^{p, q} + (-1)^p \text{d}_2^{p, q}$ 给出。
这恰好与定义 \ref{dga-definition-tensor-product} 中 $A \otimes_R B$
上的微分规则相同。
\end{proof}






\section{微分分次模}
\label{dga-section-modules}

\noindent
本章默认使用右模；左模将在第 \ref{dga-section-left-modules} 节中讨论。

\begin{definition}
\label{dga-definition-dgm}
设 $R$ 是环。设 $(A, \text{d})$ 是 $R$ 上的微分分次代数。
$A$ 上的一个（右）{\it 微分分次模} $M$ 是一个右 $A$-模 $M$，
配备分次 $M = \bigoplus M^n$ 和微分 $\text{d}$，使得
$M^n A^m \subset M^{n + m}$、$\text{d}(M^n) \subset M^{n + 1}$，并且
$$
\text{d}(ma) = \text{d}(m)a + (-1)^n m\text{d}(a)
$$
对 $a \in A$ 和 $m \in M^n$ 成立。微分分次模的一个 {\it 同态}
$f : M \to N$，是一个与分次及微分相容的 $A$-模映射。
（右）微分分次 $A$-模的范畴记为 $\text{Mod}_{(A, \text{d})}$。
\end{definition}

\noindent
注意，可以把 $M$ 看作（右）$R$-模的上链复形 $M^\bullet$。
事实上，对 $r \in R$ 有 $\text{d}(r) = 0$，且 $r$ 映到 $A$ 的一个
零次元素，因而 $\text{d}(mr) = \text{d}(m)r$。

\medskip\noindent
微分分次 $R$-代数 $A$ 上的微分分次 $R$-模 $M$ 中，联系微分与乘法的
Leibniz 规则恰好意味着乘法映射给出一个上链复形映射
$$
\text{Tot}(M^\bullet \otimes_R A^\bullet) \to M^\bullet
$$
这里 $A^\bullet$ 和 $M^\bullet$ 分别表示 $A$ 和 $M$ 的底层上链复形。

\begin{lemma}
\label{dga-lemma-dgm-abelian}
设 $(A, d)$ 是微分分次代数。范畴 $\text{Mod}_{(A, \text{d})}$
是阿贝尔范畴，并且具有任意极限和余极限。
\end{lemma}

\begin{proof}
取核或余核与取底层 $A$-模可交换。直和和余极限也是如此。
换言之，$\text{Mod}_{(A, \text{d})}$ 中的这些运算与到 $A$-模范畴的
遗忘函子可交换。乘积和极限则并非如此。事实上，若 $N_i$（$i \in I$）
是一族微分分次 $A$-模，则 $\text{Mod}_{(A, \text{d})}$ 中的乘积
$\prod N_i$ 由 $(\prod N_i)^n = \prod N_i^n$ 以及
$\prod N_i = \bigoplus_n (\prod N_i)^n$ 给出。因此，乘积确实与到
分次 $A$-模范畴的遗忘函子可交换。具有乘积和等化子的范畴具有极限；
见《范畴》引理 \ref{categories-lemma-limits-products-equalizers}。
\end{proof}

\noindent
因此，若 $(A, \text{d})$ 是 $R$ 上的微分分次代数，则有阿贝尔范畴间的
正合函子
$$
\text{Mod}_{(A, \text{d})} \longrightarrow \text{Comp}(R)
$$
。对于微分分次模 $M$，其上同调群 $H^n(M)$ 定义为相应 $R$-模复形的
上同调。因此，微分分次模的短正合列 $0 \to K \to L \to M \to 0$
产生长正合列
\begin{equation}
\label{dga-equation-les}
H^n(K) \to H^n(L) \to H^n(M) \to H^{n + 1}(K)
\end{equation}
；它由上同调模组成，见《同调代数》引理
\ref{homology-lemma-long-exact-sequence-cochain}。

\medskip\noindent
此外，从现在起，我们沿用模复形的全部术语。例如，若对所有
$k \in \mathbf{Z}$ 都有 $H^k(M) = 0$，则称微分分次 $A$-模 $M$
是 {\it 无环的}。若微分分次 $A$-模的同态 $M \to N$ 对所有
$k \in \mathbf{Z}$ 都诱导同构 $H^k(M) \to H^k(N)$，则称它为
{\it 拟同构}；其余术语依此类推。

\begin{definition}
\label{dga-definition-shift}
设 $(A, \text{d})$ 是微分分次代数。设 $M$ 是微分分次模，
其底层 $R$-模复形为 $M^\bullet$。对任意 $k \in \mathbf{Z}$，
如下定义 {\it $k$-移位模} $M[k]$：
\begin{enumerate}
\item $M[k]$ 的底层 $R$-模复形是 $M^\bullet[k]$，即
$M[k]^n = M^{n + k}$ 且 $\text{d}_{M[k]} = (-1)^k\text{d}_M$；以及
\item 作为 $A$-模，其乘法
$$
(M[k])^n \times A^m \longrightarrow (M[k])^{n + m}
$$
等于给定的乘法 $M^{n + k} \times A^m \to M^{n + k + m}$。
\end{enumerate}
对于微分分次 $A$-模的态射 $f : M \to N$，令
$f[k] : M[k] \to N[k]$ 为在底层 $A$-模上等于 $f$ 的映射。
这定义了函子
$[k] : \text{Mod}_{(A, \text{d})} \to \text{Mod}_{(A, \text{d})}$.
\end{definition}

\noindent
来验证这个选择满足 Leibniz 规则。设 $x \in M[k]^n = M^{n + k}$、
$a \in A^m$，并以 $\cdot_{M[k]}$ 表示 $M[k]$ 中的乘法，则
\begin{align*}
\text{d}_{M[k]}(x \cdot_{M[k]} a)
& =
(-1)^k \text{d}_M(xa) \\
& =
(-1)^k \text{d}_M(x) a + (-1)^{k + n + k} x \text{d}(a) \\
& =
\text{d}_{M[k]}(x) a + (-1)^n x \text{d}(a) \\
& =
\text{d}_{M[k]}(x) \cdot_{M[k]} a + (-1)^n x \cdot_{M[k]} \text{d}(a)
\end{align*}
这正是所需结论，因为作为 $M[k]$ 的齐次元素，$x$ 的次数为 $n$。
还注意到，在这些选择下，可把乘法映射看作复形映射
$$
\text{Tot}(M^\bullet[k] \otimes _R A^\bullet) \to
\text{Tot}(M^\bullet \otimes _R A^\bullet)[k] \to
M^\bullet[k]
$$
其中第一个箭头是《更多代数》第 \ref{more-algebra-section-sign-rules} 节
(\ref{more-algebra-item-shift-tensor}) 中的箭头；在本情形下它不涉及符号。
（事实上，也可由这一观察推出 Leibniz 规则成立。）

\medskip\noindent
《同调代数》第 \ref{homology-section-homotopy-shift} 节中的说明适用。
特别地，我们将在不引入符号的情况下认同所有移位 $M[k]$ 的上同调群。

\medskip\noindent
至此，我们已有足够结构来讨论 {\it 三角}；见《导出范畴》定义
\ref{derived-definition-triangle}。下一步将发展充分的理论，以陈述并证明
微分分次模的同伦范畴是三角范畴。我们先定义同伦范畴。






\section{同伦范畴}
\label{dga-section-homotopy}

\noindent
我们的同伦考虑 $A$-模结构和分次，但（当然）不要求与微分相容。

\begin{definition}
\label{dga-definition-homotopy}
设 $(A, \text{d})$ 是微分分次代数。设 $f, g : M \to N$ 是微分分次
$A$-模的同态。$f$ 与 $g$ 之间的一个 {\it 同伦}是一个 $A$-模映射
$h : M \to N$，使得
\begin{enumerate}
\item 对所有 $n$ 有 $h(M^n) \subset N^{n - 1}$；并且
\item 对所有 $x \in M$ 有
$f(x) - g(x) = \text{d}_N(h(x)) + h(\text{d}_M(x))$。
\end{enumerate}
若存在这样的同伦，则称 $f$ 与 $g$ {\it 同伦}。
\end{definition}

\noindent
因此，$h$ 与 $A$-模结构及分次相容，但不与微分相容。若 $f = g$，
且 $h$ 是定义中的同伦，则 $h$ 在 $\text{Mod}_{(A, \text{d})}$ 中
定义一个态射 $h : M \to N[-1]$。

\begin{lemma}
\label{dga-lemma-compose-homotopy}
设 $(A, \text{d})$ 是微分分次代数。设 $f, g : L \to M$ 是微分分次
$A$-模的同态，并另给同态 $a : K \to L$ 和 $c : M \to N$。
若 $A$-模映射 $h : L \to M$ 定义了 $f$ 与 $g$ 之间的同伦，
则 $c \circ h \circ a$ 定义了 $c \circ f \circ a$ 与
$c \circ g \circ a$ 之间的同伦。
\end{lemma}

\begin{proof}
立即由《同调代数》引理 \ref{homology-lemma-compose-homotopy-cochain} 得出。
\end{proof}

\noindent
这个引理使我们可以如下定义同伦范畴。

\begin{definition}
\label{dga-definition-complexes-notation}
设 $(A, \text{d})$ 是微分分次代数。{\it 同伦范畴}记为
$K(\text{Mod}_{(A, \text{d})})$；其对象是
$\text{Mod}_{(A, \text{d})}$ 的对象，其态射是微分分次 $A$-模同态的
同伦类。
\end{definition}

\noindent
记号 $K(\text{Mod}_{(A, \text{d})})$ 并非标准记号，但至少与
Stacks project 其他部分中 $K(-)$ 的用法一致。

\begin{lemma}
\label{dga-lemma-homotopy-direct-sums}
设 $(A, \text{d})$ 是微分分次代数。同伦范畴
$K(\text{Mod}_{(A, \text{d})})$ 具有直和和乘积。
\end{lemma}

\begin{proof}
略。提示：直接使用引理 \ref{dga-lemma-dgm-abelian} 中的直和和乘积。
这是可行的，因为我们已经看到这些函子与到分次 $A$-模范畴的遗忘函子
可交换，并且 $\prod$ 在阿贝尔群族所成的范畴上是正合函子。
\end{proof}







\section{锥}
\label{dga-section-cones}

\noindent
我们为微分分次模范畴引入锥。

\begin{definition}
\label{dga-definition-cone}
设 $(A, \text{d})$ 是微分分次代数。设 $f : K \to L$ 是微分分次
$A$-模的同态。$f$ 的 {\it 锥}是微分分次 $A$-模 $C(f)$；
它由 $C(f) = L \oplus K$ 给出，其分次为
$C(f)^n = L^n \oplus K^{n + 1}$，微分为
$$
d_{C(f)} =
\left(
\begin{matrix}
\text{d}_L & f \\
0 & -\text{d}_K
\end{matrix}
\right)
$$
它配备由显然映射 $L \to C(f)$ 和 $C(f) \to K$ 所诱导的典范复形态射
$i : L \to C(f)$ 和 $p : C(f) \to K[1]$。
\end{definition}

\noindent
锥三角的构造在下述意义下是函子的。

\begin{lemma}
\label{dga-lemma-functorial-cone}
设 $(A, \text{d})$ 是微分分次代数。假设
$$
\xymatrix{
K_1 \ar[r]_{f_1} \ar[d]_a & L_1 \ar[d]^b \\
K_2 \ar[r]^{f_2} & L_2
}
$$
是一个在相差同伦意义下交换的微分分次 $A$-模同态图表。
那么存在态射 $c : C(f_1) \to C(f_2)$，它产生三角态射
$$
(a, b, c) : (K_1, L_1, C(f_1), f_1, i_1, p_1) \to
(K_1, L_1, C(f_1), f_2, i_2, p_2)
$$
，后者位于 $K(\text{Mod}_{(A, \text{d})})$ 中。
\end{lemma}

\begin{proof}
设 $h : K_1 \to L_2$ 是 $f_2 \circ a$ 与 $b \circ f_1$ 之间的同伦。
用矩阵
$$
c =
\left(
\begin{matrix}
b & h \\
0 & a
\end{matrix}
\right) :
L_1 \oplus K_1 \to L_2 \oplus K_2
$$
定义 $c$。矩阵计算表明 $c$ 是微分分次模的态射。
显然 $c \circ i_1 = i_2 \circ b$，也容易验证
$p_2 \circ c = a \circ p_1$。
\end{proof}











\section{容许短正合列}
\label{dga-section-admissible}

\noindent
容许短正合列是逐项分裂正合列在微分分次模框架中的对应物。

\begin{definition}
\label{dga-definition-admissible-ses}
设 $(A, \text{d})$ 是微分分次代数。
\begin{enumerate}
\item 若存在分次 $A$-模映射 $L \to K$，它是 $K \to L$ 的左逆，
则称微分分次 $A$-模的同态 $K \to L$ 为 {\it 容许单态射}。
\item 若存在分次 $A$-模映射 $M \to L$，它是 $L \to M$ 的右逆，
则称微分分次 $A$-模的同态 $L \to M$ 为 {\it 容许满态射}。
\item 若微分分次 $A$-模的短正合列 $0 \to K \to L \to M \to 0$
作为分次 $A$-模列分裂，则称它为 {\it 容许短正合列}。
\end{enumerate}
\end{definition}

\noindent
因此，这些分裂与除微分外的所有数据相容。给定一个容许短正合列，
我们得到一个三角；这正是要求分裂与 $A$-模结构相容的原因。

\begin{lemma}
\label{dga-lemma-admissible-ses}
设 $(A, \text{d})$ 是微分分次代数。设
$0 \to K \to L \to M \to 0$ 是微分分次 $A$-模的容许短正合列。
设 $s : M \to L$ 和 $\pi : L \to K$ 是满足
$\Ker(\pi) = \Im(s)$ 的分裂。那么得到态射
$$
\delta = \pi \circ \text{d}_L \circ s : M \to K[1]
$$
，它属于 $\text{Mod}_{(A, \text{d})}$，并诱导上同调长正合列
(\ref{dga-equation-les}) 中的边界映射。
\end{lemma}

\begin{proof}
由构造，映射 $\pi \circ \text{d}_L \circ s$ 与 $A$-模结构及分次相容。
由《同调代数》引理
\ref{homology-lemma-ses-termwise-split-cochain}，它也与微分相容。
设 $R$ 是使 $A$ 成为其上微分分次代数的环。映射的等式是关于
$R$-模的陈述。因此，结论由《同调代数》引理
\ref{homology-lemma-ses-termwise-split-cochain} 和
\ref{homology-lemma-ses-termwise-split-long-cochain} 得出。
\end{proof}

\begin{lemma}
\label{dga-lemma-make-commute-map}
设 $(A, \text{d})$ 是微分分次代数。设
$$
\xymatrix{
K \ar[r]_f \ar[d]_a & L \ar[d]^b \\
M \ar[r]^g & N
}
$$
是一个相差同伦交换的微分分次 $A$-模同态图表。
\begin{enumerate}
\item 若 $f$ 是容许单态射，则 $b$ 同伦于一个使图表交换的同态。
\item 若 $g$ 是容许满态射，则 $a$ 同伦于一个使图表交换的态射。
\end{enumerate}
\end{lemma}

\begin{proof}
设 $h : K \to N$ 是 $bf$ 与 $ga$ 之间的同伦，即
$bf - ga = \text{d}h + h\text{d}$。假设分次 $A$-模映射
$\pi : L \to K$ 是 $f$ 的左逆。取
$b' = b - \text{d}h\pi - h\pi \text{d}$。
假设分次 $A$-模映射 $s : N \to M$ 是 $g$ 的右逆。
取 $a' = a + \text{d}sh + sh\text{d}$。计算略。
\end{proof}

\begin{lemma}
\label{dga-lemma-make-injective}
设 $(A, \text{d})$ 是微分分次代数。设 $\alpha : K \to L$ 是微分分次
$A$-模的同态。存在分解
$$
\xymatrix{
K \ar[r]^{\tilde \alpha} \ar@/_1pc/[rr]_\alpha &
\tilde L \ar[r]^\pi & L
}
$$
，它位于 $\text{Mod}_{(A, \text{d})}$ 中，并且
\begin{enumerate}
\item $\tilde \alpha$ 是容许单态射（见定义
\ref{dga-definition-admissible-ses}）；
\item 存在态射 $s : L \to \tilde L$，使得
$\pi \circ s = \text{id}_L$，并且 $s \circ \pi$ 同伦于
$\text{id}_{\tilde L}$。
\end{enumerate}
\end{lemma}

\begin{proof}
证明与《导出范畴》引理 \ref{derived-lemma-make-injective} 的证明相同。
具体而言，置 $\tilde L = L \oplus C(1_K)$，并使用锥构造的初等性质。
\end{proof}

\begin{lemma}
\label{dga-lemma-sequence-maps-split}
设 $(A, \text{d})$ 是微分分次代数。设
$L_1 \to L_2 \to \ldots \to L_n$ 是微分分次 $A$-模的可复合同态列。
存在交换图表
$$
\xymatrix{
L_1 \ar[r] &
L_2 \ar[r] &
\ldots \ar[r] &
L_n \\
M_1 \ar[r] \ar[u] &
M_2 \ar[r] \ar[u] &
\ldots \ar[r] &
M_n \ar[u]
}
$$
，它位于 $\text{Mod}_{(A, \text{d})}$ 中，使得每个
$M_i \to M_{i + 1}$ 都是容许单态射，且每个 $M_i \to L_i$
都是同伦等价。
\end{lemma}

\begin{proof}
$n = 1$ 的情形没有内容。引理 \ref{dga-lemma-make-injective} 是 $n = 2$
的情形。假设除 $M_n$ 外图表均已构造。对复合
$M_{n - 1} \to L_{n - 1} \to L_n$ 应用引理
\ref{dga-lemma-make-injective}，所得分解
$M_{n - 1} \to M_n \to L_n$ 即满足要求。
\end{proof}



\begin{lemma}
\label{dga-lemma-nilpotent}
设 $(A, \text{d})$ 是微分分次代数。设
$0 \to K_i \to L_i \to M_i \to 0$（$i = 1, 2, 3$）是微分分次
$A$-模的容许短正合列。设 $b : L_1 \to L_2$ 和
$b' : L_2 \to L_3$ 是微分分次模的同态，使得
$$
\vcenter{
\xymatrix{
K_1 \ar[d]_0 \ar[r] &
L_1 \ar[r] \ar[d]_b &
M_1 \ar[d]_0 \\
K_2 \ar[r] & L_2 \ar[r] & M_2
}
}
\quad\text{且}\quad
\vcenter{
\xymatrix{
K_2 \ar[d]^0 \ar[r] &
L_2 \ar[r] \ar[d]^{b'} &
M_2 \ar[d]^0 \\
K_3 \ar[r] & L_3 \ar[r] & M_3
}
}
$$
相差同伦交换。那么 $b' \circ b$ 同伦于 $0$。
\end{lemma}

\begin{proof}
由引理 \ref{dga-lemma-make-commute-map}，可把 $b$ 和 $b'$ 替换为与之同伦的
映射，使得左图的右方块交换，而右图的左方块交换。换言之，
$\Im(b) \subset \Im(K_2 \to L_2)$ 且
$\Ker(b') \supset \Im(K_2 \to L_2)$。于是，作为模映射，
$b' \circ b = 0$。
\end{proof}
















\section{有区别三角}
\label{dga-section-distinguished}

\noindent
下述引理产生我们所用的有区别三角。

\begin{lemma}
\label{dga-lemma-triangle-independent-splittings}
设 $(A, \text{d})$ 是微分分次代数。设
$0 \to K \to L \to M \to 0$ 是微分分次 $A$-模的容许短正合列。
三角
\begin{equation}
\label{dga-equation-triangle-associated-to-admissible-ses}
K \to L \to M \xrightarrow{\delta} K[1]
\end{equation}
中的 $\delta$ 如引理 \ref{dga-lemma-admissible-ses} 所述。该三角在
$K(\text{Mod}_{(A, \text{d})})$ 中相差典范同构，不依赖于引理
\ref{dga-lemma-admissible-ses} 中所作的选择。
\end{lemma}

\begin{proof}
具体地，设 $(s', \pi')$ 是引理 \ref{dga-lemma-admissible-ses} 中那种分裂的
第二个选择。我们断言 $\delta$ 与 $\delta'$ 同伦。事实上，对某些唯一的
次数为 $-1$ 的 $A$-模同态 $h : M \to K$ 和 $g : M \to K$，写成
$s' = s + \alpha \circ h$ 和 $\pi' = \pi + g \circ \beta$。
则 $g = -h$，且 $g$ 是 $\delta$ 与 $\delta'$ 之间的同伦。
相关计算见《同调代数》引理
\ref{homology-lemma-ses-termwise-split-homotopy-cochain} 的证明。
\end{proof}

\begin{definition}
\label{dga-definition-distinguished-triangle}
设 $(A, \text{d})$ 是微分分次代数。
\begin{enumerate}
\item 若 $0 \to K \to L \to M \to 0$ 是微分分次 $A$-模的容许短正合列，
则 {\it 与 $0 \to K \to L \to M \to 0$ 关联的三角}是
$K(\text{Mod}_{(A, \text{d})})$ 中的三角
(\ref{dga-equation-triangle-associated-to-admissible-ses})。
\item 若 $K(\text{Mod}_{(A, \text{d})})$ 中的三角同构于某个与微分分次
$A$-模的容许短正合列关联的三角，则称它为 {\it 有区别三角}。
\end{enumerate}
\end{definition}









\section{锥与有区别三角}
\label{dga-section-cones-and-triangles}

\noindent
设 $(A, \text{d})$ 是微分分次代数。设 $f : K \to L$ 是微分分次
$A$-模的同态。那么 $(K, L, C(f), f, i, p)$ 构成三角
$$
K \to L \to C(f) \to K[1]
$$
，它位于 $\text{Mod}_{(A, \text{d})}$ 中，因而也位于
$K(\text{Mod}_{(A, \text{d})})$ 中。一般而言，锥{\bf 并不}是有区别三角，
但两者只相差一个符号或一次旋转（任君选择）。下面给出两个精确陈述。

\begin{lemma}
\label{dga-lemma-rotate-cone}
设 $(A, \text{d})$ 是微分分次代数。设 $f : K \to L$ 是微分分次模的同态。
三角 $(L, C(f), K[1], i, p, f[1])$ 是与容许短正合列
$$
0 \to L \to C(f) \to K[1] \to 0
$$
关联的三角；该短正合列来自 $f$ 的锥的定义。
\end{lemma}

\begin{proof}
立即由定义得出。
\end{proof}

\begin{lemma}
\label{dga-lemma-rotate-triangle}
设 $(A, \text{d})$ 是微分分次代数。设 $\alpha : K \to L$ 和
$\beta : L \to M$ 定义容许短正合列
$$
0 \to K \to L \to M \to 0
$$
，它由微分分次 $A$-模组成。设 $(K, L, M, \alpha, \beta, \delta)$
为关联三角。那么三角
$$
(M[-1], K, L, \delta[-1], \alpha, \beta)
\quad\text{且}\quad
(M[-1], K, C(\delta[-1]), \delta[-1], i, p)
$$
同构。
\end{lemma}

\begin{proof}
选取分裂后，写 $L = K \oplus M$，并把 $\alpha$ 和 $\beta$ 分别认同为
自然包含映射和投影映射。由 $\delta$ 的构造，有
$$
d_B =
\left(
\begin{matrix}
d_K & \delta \\
0 & d_M
\end{matrix}
\right)
$$
另一方面，$\delta[-1] : M[-1] \to K$ 的锥由
$C(\delta[-1]) = K \oplus M$ 给出，其微分与上面的矩阵完全相同！
引理由此得证。
\end{proof}

\begin{lemma}
\label{dga-lemma-third-isomorphism}
设 $(A, \text{d})$ 是微分分次代数。设 $f_1 : K_1 \to L_1$ 和
$f_2 : K_2 \to L_2$ 是微分分次 $A$-模的同态。设
$$
(a, b, c) :
(K_1, L_1, C(f_1), f_1, i_1, p_1)
\longrightarrow
(K_1, L_1, C(f_1), f_2, i_2, p_2)
$$
为 $K(\text{Mod}_{(A, \text{d})})$ 中的任意三角态射。
若 $a$ 和 $b$ 是同伦等价，则 $c$ 也是同伦等价。
\end{lemma}

\begin{proof}
设 $a^{-1} : K_2 \to K_1$ 是微分分次 $A$-模的同态，且它在
$K(\text{Mod}_{(A, \text{d})})$ 中是 $a$ 的逆。设
$b^{-1} : L_2 \to L_1$ 是微分分次 $A$-模的同态，且它在该范畴中是
$b$ 的逆。对 $f_1 \circ a^{-1} = b^{-1} \circ f_2$ 应用引理
\ref{dga-lemma-functorial-cone}，得到态射 $c' : C(f_2) \to C(f_1)$。
若能证明 $c \circ c'$ 和 $c' \circ c$ 在
$K(\text{Mod}_{(A, \text{d})})$ 中都是同构，结论即得。因此，只需证明：
给定 $K(\text{Mod}_{(A, \text{d})})$ 中形如
$(1, 1, c) : (K, L, C(f), f, i, p)$ 的三角态射，态射 $c$ 是该范畴中的
同构。由假设，下图中的两个方块
$$
\xymatrix{
L \ar[r] \ar[d]_1 &
C(f) \ar[r] \ar[d]_c &
K[1] \ar[d]_1 \\
L \ar[r] &
C(f) \ar[r] &
K[1]
}
$$
相差同伦交换。由 $C(f)$ 的构造，两行都是容许短正合列。
于是，由引理 \ref{dga-lemma-nilpotent}，在
$K(\text{Mod}_{(A, \text{d})})$ 中有 $(c - 1)^2 = 0$。
因此，$c$ 是该范畴中的同构，其逆为 $2 - c$。
\end{proof}

\noindent
下述引理表明，同伦范畴中由锥给出的三角族与有区别三角族相差同构而相同，
至少在不计符号时如此！

\begin{lemma}
\label{dga-lemma-the-same-up-to-isomorphisms}
设 $(A, \text{d})$ 是微分分次代数。
\begin{enumerate}
\item 给定微分分次 $A$-模的容许短正合列
$0 \to K \xrightarrow{\alpha} L \to M \to 0$
，存在同伦等价 $C(\alpha) \to M$，使得图表
$$
\xymatrix{
K \ar[r] \ar[d] & L \ar[d] \ar[r] &
C(\alpha) \ar[r]_{-p} \ar[d] & K[1] \ar[d] \\
K \ar[r]^\alpha & L \ar[r]^\beta &
M \ar[r]^\delta & K[1]
}
$$
在 $K(\text{Mod}_{(A, \text{d})})$ 中定义一个三角同构。
\item 给定复形态射 $f : K \to L$，存在三角同构
$$
\xymatrix{
K \ar[r] \ar[d] & \tilde L \ar[d] \ar[r] &
M \ar[r]_{\delta} \ar[d] & K[1] \ar[d] \\
K \ar[r] & L \ar[r] &
C(f) \ar[r]^{-p} & K[1]
}
$$
，其中上方三角是与容许短正合列 $K \to \tilde L \to M$ 关联的三角。
\end{enumerate}
\end{lemma}

\begin{proof}
证明 (1)。有 $C(\alpha) = L \oplus K$；我们简单地把
$C(\alpha) \to M$ 定义为先投影到 $L$，再与 $\beta$ 复合。
由于复合 $K^{n + 1} \to L^{n + 1} \to M^{n + 1}$ 为零，这定义了
微分分次模的态射。选择满足 $\Ker(\pi) = \Im(s)$ 的分裂
$s : M \to L$ 和 $\pi : L \to K$，并照例置
$\delta = \pi \circ \text{d}_L \circ s$。取由 $(s , -\delta)$ 给出的
$M \to C(\alpha)$ 作为同伦逆。它与微分相容，因为 $\delta^n$
可刻画为唯一满足下式的映射 $M^n \to K^{n + 1}$：
$\text{d} \circ s^n - s^{n + 1} \circ \text{d} = \alpha \circ \delta^n$,
见《同调代数》引理 \ref{homology-lemma-ses-termwise-split-cochain} 的证明。
复合 $M \to C(f) \to M$ 是恒等映射。复合
$C(f) \to M \to C(f)$ 等于态射
$$
\left(
\begin{matrix}
s \circ \beta & 0 \\
-\delta \circ \beta & 0
\end{matrix}
\right)
$$
为看出它同伦于恒等映射，使用由矩阵
$$
\left(
\begin{matrix}
0 & 0 \\
\pi & 0
\end{matrix}
\right) :
C(\alpha) = L \oplus K
\to
L \oplus K = C(\alpha)
$$
给出的同伦 $h : C(\alpha) \to C(\alpha)$。容易验证
$$
\left(
\begin{matrix}
1 & 0 \\
0 & 1
\end{matrix}
\right)
-
\left(
\begin{matrix}
s \\
-\delta
\end{matrix}
\right)
\left(
\begin{matrix}
\beta & 0
\end{matrix}
\right)
=
\left(
\begin{matrix}
\text{d} & \alpha \\
0 & -\text{d}
\end{matrix}
\right)
\left(
\begin{matrix}
0 & 0 \\
\pi & 0
\end{matrix}
\right)
+
\left(
\begin{matrix}
0 & 0 \\
\pi & 0
\end{matrix}
\right)
\left(
\begin{matrix}
\text{d} & \alpha \\
0 & -\text{d}
\end{matrix}
\right)
$$
为完成 (1) 的证明，还须证明态射 $-p : C(\alpha) \to K[1]$
（见定义 \ref{dga-definition-cone}）与复合 $C(\alpha) \to M \to K[1]$
相差同伦相等。这由上文显然可见。事实上，可使用同伦逆
$(s, -\delta) : M \to C(\alpha)$，转而验证两个映射
$M \to K[1]$ 相等。注意 $p \circ (s, -\delta) = -\delta$，正合所需。

\medskip\noindent
证明 (2)。令 $\tilde f : K \to \tilde L$、$s : L \to \tilde L$
和 $\pi : L \to L$ 如引理 \ref{dga-lemma-make-injective} 中所述。
由引理 \ref{dga-lemma-functorial-cone} 和 \ref{dga-lemma-third-isomorphism}，
三角 $(K, L, C(f), i, p)$ 与
$(K, \tilde L, C(\tilde f), \tilde i, \tilde p)$ 同构。
注意，可以复合三角的同构。因此，可以用 $\tilde L$ 替换 $L$，
并用 $\tilde f$ 替换 $f$。换言之，可以假设 $f$ 是容许单态射。
此时结论由 (1) 得出。
\end{proof}







\section{同伦范畴是三角范畴}
\label{dga-section-homotopy-triangulated}

\noindent
我们先证明它是预三角范畴。

\begin{lemma}
\label{dga-lemma-homotopy-category-pre-triangulated}
设 $(A, \text{d})$ 是微分分次代数。同伦范畴
$K(\text{Mod}_{(A, \text{d})})$ 连同其自然平移函子和有区别三角，
是一个预三角范畴。
\end{lemma}

\begin{proof}
证明 TR1。由定义，任何同构于有区别三角的三角都是有区别三角。
此外，任意三角 $(K, K, 0, 1, 0, 0)$ 都是有区别三角，因为
$0 \to K \to K \to 0 \to 0$ 是容许短正合列。最后，给定微分分次
$A$-模的任意同态 $f : K \to L$，由引理
\ref{dga-lemma-the-same-up-to-isomorphisms}，三角
$(K, L, C(f), f, i, -p)$ 是有区别三角。

\medskip\noindent
证明 TR2。设 $(X, Y, Z, f, g, h)$ 是一个三角，并假设
$(Y, Z, X[1], g, h, -f[1])$ 是有区别三角。则存在容许短正合列
$0 \to K \to L \to M \to 0$，使其关联三角
$(K, L, M, \alpha, \beta, \delta)$ 同构于
$(Y, Z, X[1], g, h, -f[1])$。反向旋转可见，
$(X, Y, Z, f, g, h)$ 同构于
$(M[-1], K, L, -\delta[-1], \alpha, \beta)$。
由引理 \ref{dga-lemma-rotate-triangle}，三角
$(M[-1], K, L, \delta[-1], \alpha, \beta)$ 同构于
$(M[-1], K, C(\delta[-1]), \delta[-1], i, p)$。
把前一三角同构预复合以 $Y$ 上的 $-1$，可得
$(X, Y, Z, f, g, h)$ 同构于
$(M[-1], K, C(\delta[-1]), \delta[-1], i, -p)$。
因此，由引理 \ref{dga-lemma-the-same-up-to-isomorphisms}，它是有区别三角。
反过来，假设 $(X, Y, Z, f, g, h)$ 是有区别三角。
由引理 \ref{dga-lemma-the-same-up-to-isomorphisms}，这意味着它同构于某个
$\text{Mod}_{(A, \text{d})}$ 中的态射 $f$ 所给出的三角
$(K, L, C(f), f, i, -p)$。于是旋转后的三角
$(Y, Z, X[1], g, h, -f[1])$ 同构于
$(L, C(f), K[1], i, -p, -f[1])$，后者又同构于三角
$(L, C(f), K[1], i, p, f[1])$。由引理 \ref{dga-lemma-rotate-cone}，
这个三角是有区别三角。因此，$(Y, Z, X[1], g, h, -f[1])$
如所需是有区别三角。

\medskip\noindent
证明 TR3。设 $(X, Y, Z, f, g, h)$ 和 $(X', Y', Z', f', g', h')$
是 $K(\mathcal{A})$ 的有区别三角，并设 $a : X \to X'$ 和
$b : Y \to Y'$ 是满足 $f' \circ a = b \circ f$ 的态射。
由引理 \ref{dga-lemma-the-same-up-to-isomorphisms}，可以假设
$(X, Y, Z, f, g, h) = (X, Y, C(f), f, i, -p)$，并且
$(X', Y', Z', f', g', h') = (X', Y', C(f'), f', i', -p')$。
此时，只需对 $f, f', a, b$ 给出的交换图表应用引理
\ref{dga-lemma-functorial-cone}。
\end{proof}

\noindent
在一般地证明 TR4 之前，先证明一个特殊情形。

\begin{lemma}
\label{dga-lemma-two-split-injections}
设 $(A, \text{d})$ 是微分分次代数。假设 $\alpha : K \to L$ 和
$\beta : L \to M$ 是微分分次 $A$-模的容许单态射。那么存在有区别三角
$(K, L, Q_1, \alpha, p_1, d_1)$、
$(K, M, Q_2, \beta \circ \alpha, p_2, d_2)$ 和
$(L, M, Q_3, \beta, p_3, d_3)$，使得 TR4 对它们成立。
\end{lemma}

\begin{proof}
设 $\pi_1 : L \to K$ 和 $\pi_3 : M \to L$ 是分次 $A$-模同态，
分别为 $\alpha$ 和 $\beta$ 的左逆。于是 $K \to M$ 也是容许单态射，
其左逆为 $\pi_2 = \pi_1 \circ \pi_3$。用 $Q_1$、$Q_2$、$Q_3$
分别表示 $K \to L$、$K \to M$、$L \to M$ 的余核。
则作为分次 $A$-模，有认同
$Q_1 = \Ker(\pi_1)$、$Q_3 = \Ker(\pi_3)$ 和
$Q_2 = \Ker(\pi_2)$。又有 $L = K \oplus Q_1$ 和
$M = L \oplus Q_3$，故 $M = K \oplus Q_1 \oplus Q_3$。
注意，$\pi_2 = \pi_1 \circ \pi_3$ 在 $Q_1$ 和 $Q_3$ 上均为零。
因此 $Q_2 = Q_1 \oplus Q_3$。考虑交换图表
$$
\begin{matrix}
0 & \to & K & \to & L & \to & Q_1 & \to & 0 \\
  &     & \downarrow &     & \downarrow &     & \downarrow  & \\
0 & \to & K & \to & M & \to & Q_2 & \to & 0 \\
  &     & \downarrow &     & \downarrow &     & \downarrow  & \\
0 & \to & L & \to & M & \to & Q_3 & \to & 0
\end{matrix}
$$
该图表的各行都是容许短正合列，因而按定义确定有区别三角。此外，
上图中的向下箭头与所选分裂相容，因而定义三角态射
$$
(K \to L \to Q_1 \to K[1])
\longrightarrow
(K \to M \to Q_2 \to K[1])
$$
以及
$$
(K \to M \to Q_2 \to K[1])
\longrightarrow
(L \to M \to Q_3 \to L[1]).
$$
注意，把图表中最下方列的分裂 $Q_3 \to M$ 与 $M \to Q_2$ 复合，
就得到分裂列 $0 \to Q_1 \to Q_2 \to Q_3 \to 0$ 的一个分裂。
由此容易看出，相应有区别三角
$$
(Q_1 \to Q_2 \to Q_3 \to Q_1[1])
$$
中的态射 $\delta : Q_3 \to Q_1[1]$ 等于复合
$Q_3 \to L[1] \to Q_1[1]$。因此，我们得到公理 TR4 结论中的结构。
\end{proof}

\noindent
下面给出最终结论。

\begin{proposition}
\label{dga-proposition-homotopy-category-triangulated}
设 $(A, \text{d})$ 是微分分次代数。微分分次 $A$-模的同伦范畴
$K(\text{Mod}_{(A, \text{d})})$ 连同其自然平移函子和有区别三角，
是一个三角范畴。
\end{proposition}

\begin{proof}
已知 $K(\text{Mod}_{(A, \text{d})})$ 是预三角范畴。因此，只需证明 TR4，
而证明它可使用《导出范畴》引理
\ref{derived-lemma-easier-axiom-four}。设 $K \to L$ 和 $L \to M$
是 $K(\text{Mod}_{(A, \text{d})})$ 中的可复合态射。由引理
\ref{dga-lemma-sequence-maps-split}，可以假设 $K \to L$ 和 $L \to M$
都是容许单态射。此时结论由引理 \ref{dga-lemma-two-split-injections} 得出。
\end{proof}











\section{左模}
\label{dga-section-left-modules}

\noindent
迄今所述的一切在左微分分次模的框架中都有对应版本，
但必须留意某些符号规则。

\medskip\noindent
设 $(A, \text{d})$ 是微分分次 $R$-代数。完全类似于右模，
把一个 {\it 左微分分次 $A$-模} $M$ 定义为一个左 $A$-模 $M$，
它配备分次 $M = \bigoplus M^n$ 和微分 $\text{d}$，使得
$A^n M^m \subset M^{n + m}$、$\text{d}(M^n) \subset M^{n + 1}$，并且
$$
\text{d}(am) = \text{d}(a) m + (-1)^{\deg(a)}a \text{d}(m)
$$
对齐次元素 $a \in A$ 和 $m \in M$ 成立。如前，这一 Leibniz 规则
恰好意味着乘法定义复形映射
$$
\text{Tot}(A^\bullet \otimes_R M^\bullet) \to M^\bullet
$$
这里 $A^\bullet$ 和 $M^\bullet$ 表示 $A$ 和 $M$ 的底层 $R$-模复形。

\begin{definition}
\label{dga-definition-opposite-dga}
设 $R$ 是环。设 $(A, \text{d})$ 是 $R$ 上的微分分次代数。
{\it 反微分分次代数}是 $R$ 上的微分分次代数
$(A^{opp}, \text{d})$；作为分次 $R$-模，有 $A^{opp} = A$，
微分满足 $\text{d} = \text{d}$，而乘法由
$$
a \cdot_{opp} b = (-1)^{\deg(a)\deg(b)} b a
$$
给出，其中 $a, b \in A$ 为齐次元素。
\end{definition}

\noindent
该定义成立，因为
\begin{align*}
\text{d}(a \cdot_{opp} b)
& =
(-1)^{\deg(a)\deg(b)} \text{d}(b a) \\
& =
(-1)^{\deg(a)\deg(b)} \text{d}(b) a +
(-1)^{\deg(a)\deg(b) + \deg(b)}b\text{d}(a) \\
& =
(-1)^{\deg(a)}a \cdot_{opp} \text{d}(b) + \text{d}(a) \cdot_{opp} b
\end{align*}
，正合所需。用底层 $R$-模复形表述，这意味着图表
$$
\xymatrix{
\text{Tot}(A^\bullet \otimes_R A^\bullet)
\ar[rrr]_-{A^{opp}\text{ 的乘法}}
\ar[d]_{\text{交换性约束}} & & &
A^\bullet \ar[d]^{\text{id}} \\
\text{Tot}(A^\bullet \otimes_R A^\bullet)
\ar[rrr]^-{A\text{ 的乘法}} & & &
A^\bullet
}
$$
交换。这里，$R$-模复形的对称幺半范畴上的交换性约束见《更多代数》
第 \ref{more-algebra-section-sign-rules} 节。

\medskip\noindent
设 $(A, \text{d})$ 是 $R$ 上的微分分次代数。设 $M$ 是左微分分次
$A$-模。把 $M$ 视为右 $A^{opp}$-模时记为 $M^{opp}$，其乘法
$\cdot_{opp}$ 由规则
$$
m \cdot_{opp} a = (-1)^{\deg(a)\deg(m)} a m
$$
定义，其中 $a$ 和 $m$ 为齐次元素。它与微分相容，因为也可以用图表
$$
\xymatrix{
\text{Tot}(M^\bullet \otimes_R A^\bullet)
\ar[rrr]_-{M^{opp}\text{ 上的乘法}}
\ar[d]_{\text{交换性约束}} & & &
M^\bullet \ar[d]^{\text{id}} \\
\text{Tot}(A^\bullet \otimes_R M^\bullet)
\ar[rrr]^-{M\text{ 上的乘法}} & & &
M^\bullet
}
$$
来定义 $M^{opp}$ 上的乘法 $\cdot_{opp}$。为看出它是结合乘法，
对齐次元素 $a, b \in A$ 和 $m \in M$ 计算如下：
\begin{align*}
m \cdot_{opp} (a \cdot_{opp} b)
& =
(-1)^{\deg(a)\deg(b)} m \cdot_{opp} (ba) \\
& =
(-1)^{\deg(a)\deg(b) + \deg(ab)\deg(m)} bam \\
& =
(-1)^{\deg(a)\deg(b) + \deg(ab)\deg(m) + \deg(b)\deg(am)}
(am) \cdot_{opp} b \\
& =
(-1)^{\deg(a)\deg(b) + \deg(ab)\deg(m) + \deg(b)\deg(am) + \deg(a)\deg(m)}
(m \cdot_{opp} a) \cdot_{opp} b \\
& =
(m \cdot_{opp} a) \cdot_{opp} b
\end{align*}
当然，也可利用 $R$-模复形的对称幺半范畴上结合性约束与交换性约束之间的
相容性证明这一点。

\begin{lemma}
\label{dga-lemma-left-right}
设 $(A, \text{d})$ 是微分分次 $R$-代数。从左微分分次 $A$-模范畴到
右微分分次 $A^{opp}$-模范畴的函子 $M \mapsto M^{opp}$ 是等价。
\end{lemma}

\begin{proof}
略。
\end{proof}

\noindent
接着讨论移位。设 $(A, \text{d})$ 是微分分次代数。设 $M$ 是左微分分次
$A$-模，其底层 $R$-模复形记为 $M^\bullet$。对任意
$k \in \mathbf{Z}$，如下定义 {\it $k$-移位模} $M[k]$：
\begin{enumerate}
\item $M[k]$ 的底层 $R$-模复形是 $M^\bullet[k]$；
\item 作为 $A$-模，其乘法
$$
A^n \times (M[k])^m \longrightarrow (M[k])^{n + m}
$$
等于给定乘法 $A^n \times M^{m + k} \to M^{n + m + k}$ 的
$(-1)^{nk}$ 倍。
\end{enumerate}
来验证该选择满足 Leibniz 规则。设 $a \in A^n$、
$x \in M[k]^m = M^{m + k}$，并以 $\cdot_{M[k]}$ 表示 $M[k]$ 中的乘法，
则
\begin{align*}
\text{d}_{M[k]}(a \cdot_{M[k]} x)
& =
(-1)^{k + nk} \text{d}_M(ax) \\
& =
(-1)^{k + nk} \text{d}(a) x + (-1)^{k + nk + n} a \text{d}_M(x) \\
& =
\text{d}(a) \cdot_{M[k]} x + (-1)^{nk + n} a \text{d}_{M[k]}(x) \\
& =
\text{d}(a) \cdot_{M[k]} x + (-1)^n a \cdot_{M[k]} \text{d}_{M[k]}(x)
\end{align*}
这正是所需结论，因为作为 $A$ 的齐次元素，$a$ 的次数为 $n$。
还注意到，在这些选择下，可把乘法映射看作复形映射
$$
\text{Tot}(A^\bullet \otimes_R M^\bullet[k]) \to
\text{Tot}(A^\bullet \otimes_R M^\bullet)[k] \to
M^\bullet[k]
$$
其中第一个箭头是《更多代数》第 \ref{more-algebra-section-sign-rules} 节
(\ref{more-algebra-item-shift-tensor}) 中的箭头；在本情形下它恰好涉及
上面所选的符号。（事实上，也可由这一观察推出 Leibniz 规则成立。）

\medskip\noindent
由上述规则，有典范认同
$$
(M[k])^{opp} = M^{opp}[k]
$$
，它是右微分分次 $A^{opp}$-模之间不引入符号的认同。换言之，
引理 \ref{dga-lemma-left-right} 的等价与移位函子相容。

\medskip\noindent
上述选择要求作出如下定义。

\begin{definition}
\label{dga-definition-shift-graded-module}
设 $R$ 是环。设 $A$ 是 $\mathbf{Z}$-分次 $R$-代数。
\begin{enumerate}
\item 给定右分次 $A$-模 $M$，把 {\it 第 $k$ 移位 $A$-模} $M[k]$
定义为与之相同的右 $A$-模，但分次改为 $(M[k])^n = M^{n + k}$。
\item 给定左分次 $A$-模 $M$，把 {\it 第 $k$ 移位 $A$-模} $M[k]$
定义为分次满足 $(M[k])^n = M^{n + k}$ 的模，其乘法
$A^n \times (M[k])^m \to (M[k])^{n + m}$
等于给定乘法 $A^n \times M^{m + k} \to M^{n + m + k}$ 的
$(-1)^{nk}$ 倍。
\end{enumerate}
\end{definition}

\noindent
设 $(A, \text{d})$ 是微分分次代数。设 $f, g : M \to N$ 是左微分分次
$A$-模的同态。$f$ 与 $g$ 之间的一个 {\it 同伦}是分次 $A$-模映射
$h : M \to N[-1]$（注意移位！），使得
$$
f(x) - g(x) = \text{d}_N(h(x)) + h(\text{d}_M(x))
$$
对所有 $x \in M$ 成立。若存在这样的同伦，则称 $f$ 与 $g$
{\it 同伦}。因此，$h$ 与 $A$-模结构（$N$ 上采用移位后的结构）及分次
（$N$ 上采用移位后的分次）相容，但不与微分相容。若 $f = g$ 且 $h$
是同伦，则 $h$ 定义左微分分次 $A$-模的态射 $h : M \to N[-1]$。

\medskip\noindent
由上述规则可知，$f, g : M \to N$ 同伦，当且仅当所诱导的态射
$f^{opp}, g^{opp} : M^{opp} \to N^{opp}$ 作为右微分分次
$A^{opp}$-模同态同伦（使用同一个同伦）。

\medskip\noindent
同伦范畴、锥、容许短正合列和有区别三角，全都按与右微分分次模完全相同的
方式定义（在底层 $R$-模复形上，一切都与 $R$-模复形的构造一致）。
这样也得到命题 \ref{dga-proposition-homotopy-category-triangulated}
关于左模的对应版本；或者也可通过研究反代数上的右模推出它。



\section{张量积}
\label{dga-section-tensor-product}

\noindent
设 $R$ 是环。设 $A$ 是 $R$-代数（见第 \ref{dga-section-conventions} 节）。
给定右 $A$-模 $M$ 和左 $A$-模 $N$，有 {\it 张量积}
$$
M \otimes_A N
$$
该张量积是 $R$-模。作为 $R$-模，$M \otimes_A N$ 由符号
$x \otimes y$（$x \in M$、$y \in N$）生成，并满足关系
$$
\begin{matrix}
(x_1 + x_2) \otimes y - x_1 \otimes y - x_2 \otimes y, \\
x \otimes (y_1 + y_2) - x \otimes y_1 - x \otimes y_2, \\
xa \otimes y - x \otimes ay
\end{matrix}
$$
，其中 $a \in A$、$x, x_1, x_2 \in M$ 且 $y, y_1, y_2 \in N$。
下面列出张量积的一些性质。

\medskip\noindent
张量积关于每个变元都是右正合的；事实上，它与直和及任意余极限可交换。

\medskip\noindent
张量积 $M \otimes_A N$ 承载普遍的 $A$-双线性映射
$M \times N \to M \otimes_A N$，$(x, y) \mapsto x \otimes y$。用公式表示为
$$
\text{双线性}_A(M \times N, Q) = \Hom_R(M \otimes_A N, Q)
$$
，对任意 $R$-模 $Q$ 成立。

\medskip\noindent
若 $A$ 是 $\mathbf{Z}$-分次代数，而 $M$、$N$ 是分次 $A$-模，
则 $M \otimes_A N$ 是分次 $R$-模。此时，$M \otimes_A N$ 的第 $n$
分次部分 $(M \otimes_A N)^n$ 等于
$$
\Coker\left(
\bigoplus\nolimits_{r + t + s = n}
M^r \otimes_R A^t \otimes_R N^s \to
\bigoplus\nolimits_{p + q = n} M^p \otimes_R N^q
\right)
$$
其中映射把 $x \otimes a \otimes y$ 送到
$x \otimes ay - xa \otimes y$；这里 $x \in M^r$、$y \in N^s$、
$a \in A^t$ 且 $r + s + t = n$。在这种情形下，映射
$M \times N \to M \otimes_A N$ 是 $A$-双线性的、与分次相容的，
并在下述意义下是普遍的：
$$
\text{分次双线性}_A(M \times N, Q) =
\Hom_{\text{分次 }R\text{-模}}(M \otimes_A N, Q)
$$
，对任意分次 $R$-模 $Q$ 成立；这里采用分次双线性映射的显然概念。

\medskip\noindent
若 $(A, \text{d})$ 是微分分次代数，而 $M$ 和 $N$ 分别是左、右微分分次
$A$-模，则 $M \otimes_A N$ 是微分分次 $R$-模，其微分为
$$
\text{d}(x \otimes y) =
\text{d}(x) \otimes y + (-1)^{\deg(x)}x \otimes \text{d}(y)
$$
，其中 $x \in M$ 和 $y \in N$ 为齐次元素。在这种情形下，映射
$M \times N \to M \otimes_A N$ 是 $A$-双线性的，并与分次及微分相容，
且在下述意义下是普遍的：
$$
\text{微分分次双线性}_A(M \times N, Q) =
\Hom_{\text{Comp}(R)}(M \otimes_A N, Q)
$$
，对任意微分分次 $R$-模 $Q$ 成立；这里采用微分分次双线性映射的显然概念。






\section{Hom 复形与微分分次模}
\label{dga-section-hom-complexes}

\noindent
我们强烈建议读者跳过本节。

\medskip\noindent
设 $R$ 是环，$M^\bullet$ 是 $R$-模复形。考虑 $R$-模复形
$$
E^\bullet = \Hom^\bullet(M^\bullet, M^\bullet)
$$
，它在《更多代数》第 \ref{more-algebra-section-hom-complexes} 节中引入。
由《更多代数》引理 \ref{more-algebra-lemma-composition}，有典范复合法则
$$
\text{Tot}(E^\bullet \otimes_R E^\bullet) \to E^\bullet
$$
，它是复形映射。因此，$E^\bullet$ 连同这一乘法是微分分次 $R$-代数，
记为 $(E, \text{d})$。此外，把 $M^\bullet$ 看作
$\Hom^\bullet(R, M^\bullet)$，可见复合定义乘法
$$
\text{Tot}(E^\bullet \otimes_R M^\bullet) \to M^\bullet
$$
，它使 $M^\bullet$ 成为一个{\bf 左}微分分次 $E$-模，记为 $M$。

\begin{lemma}
\label{dga-lemma-left-module-structure}
在上述情形中，设 $A$ 是微分分次 $R$-代数。在 $M$ 上给定左 $A$-模结构，
等价于给定微分分次 $R$-代数同态 $A \to E$。
\end{lemma}

\begin{proof}
证明略。注意，这一对应不涉及符号。
\end{proof}

\noindent
继续上述讨论，并假设另给一个 $R$-模复形 $N^\bullet$。考虑《更多代数》
第 \ref{more-algebra-section-hom-complexes} 节中引入的 $R$-模复形
$\Hom^\bullet(M^\bullet, N^\bullet)$。如上可见，复合
$$
\text{Tot}(\Hom^\bullet(M^\bullet, N^\bullet) \otimes_R E^\bullet)
\to \Hom^\bullet(M^\bullet, N^\bullet)
$$
定义一个乘法，使 $\Hom^\bullet(M^\bullet, N^\bullet)$ 成为一个
{\bf 右}微分分次 $E$-模。应用引理 \ref{dga-lemma-left-module-structure}，
可得：给定左微分分次 $A$-模 $M$ 和 $R$-模复形 $N^\bullet$，
存在一个典范的右微分分次 $A$-模，其底层 $R$-模复形为
$\Hom^\bullet(M^\bullet, N^\bullet)$，且乘法
$$
\Hom^n(M^\bullet, N^\bullet) \times A^m \longrightarrow
\Hom^{n + m}(M^\bullet, N^\bullet)
$$
把满足 $f_{p, q} \in \Hom(M^{-q}, N^p)$ 的
$f = (f_{p, q})_{p + q = n}$ 及 $a \in A^m$ 送到元素
$f \cdot a = (f_{p, q} \circ a)$，其中 $f_{p, q} \circ a$ 是映射
$$
M^{-q - m} \xrightarrow{a} M^{-q} \xrightarrow{f_{p, q}} N^p, \quad
x \longmapsto f_{p, q}(ax)
$$
；这里不引入符号。用记号 $\Hom(M, N^\bullet)$ 表示这个右微分分次
$A$-模。

\begin{lemma}
\label{dga-lemma-characterize-hom}
设 $R$ 是环。设 $(A, \text{d})$ 是微分分次 $R$-代数。
设 $M'$ 是右微分分次 $A$-模，$M$ 是左微分分次 $A$-模，
$N^\bullet$ 是 $R$-模复形。那么有
$$
\Hom_{\text{Mod}_{(A, d)}}(M', \Hom(M, N^\bullet)) =
\Hom_{\text{Comp}(R)}(M' \otimes_A M, N^\bullet)
$$
，其中如第 \ref{dga-section-tensor-product} 节那样，把
$M \otimes_A M$ 看作 $R$-模复形。
\end{lemma}

\begin{proof}
来说明等式两边都对应于与微分相容的分次 $A$-双线性映射
$$
M' \times M \longrightarrow N^\bullet
$$
。第 \ref{dga-section-tensor-product} 节已经说明右边确实如此。
给定左边的一个元素 $g$，《更多代数》引理
\ref{more-algebra-lemma-compose} 中的等式确定一个 $R$-模复形映射
$g' : \text{Tot}(M' \otimes_R M) \to N^\bullet$。
换言之，得到一个与微分相容的分次 $R$-双线性映射
$g'' : M' \times M \to N^\bullet$。$g$ 的 $A$-线性立即转化为
$g''$ 的 $A$-双线性。
\end{proof}

\noindent
设 $R$、$M^\bullet$、$E^\bullet$、$E$ 和 $M$ 如上。不过，现在假设给定
一个微分分次 $R$-代数 $A$ 以及 $M$ 上的一个{\bf 右}微分分次
$A$-模结构。于是可考虑映射
$$
\text{Tot}(A^\bullet \otimes_R M^\bullet)
\xrightarrow{\psi}
\text{Tot}(A^\bullet \otimes_R M^\bullet)
\to
M^\bullet
$$
其中第一个箭头是 $R$-模复形的微分分次范畴上的交换性约束。
这对应于 $R$-模复形映射
$$
\tau : A^\bullet \longrightarrow E^\bullet
$$
。回忆 $E^n = \prod_{p + q = n} \Hom_R(M^{-q}, M^p)$，并对
$a \in A^n$ 写 $\tau(a) = (\tau_{p, q}(a))_{p + q = n}$。于是
$$
\tau_{p, q}(a) : M^{-q} \longrightarrow M^p,\quad
x \longmapsto (-1)^{\deg(a)\deg(x)}x a = (-1)^{-nq}xa
$$
正如读者根据第 \ref{dga-section-left-modules} 节的讨论所应预料的，
它与 $A$ 上的乘法不相容。事实上，有
$$
\tau(a a') = (-1)^{\deg(a)\deg(a')}\tau(a') \tau(a)
$$
由此得到下述引理。

\begin{lemma}
\label{dga-lemma-right-module-structure}
在上述情形中，设 $A$ 是微分分次 $R$-代数。在 $M$ 上给定右 $A$-模结构，
等价于给定微分分次 $R$-代数同态 $\tau : A \to E^{opp}$。
\end{lemma}

\begin{proof}
见上面的讨论，并注意，由乘法映射
$M^n \times A^m \to M^{n + m}$ 构造 $\tau$ 时使用了符号。
\end{proof}

\noindent
设 $R$、$M^\bullet$、$E^\bullet$、$E$、$A$ 和 $M$ 如上，并如引理中那样
在 $M$ 上给定右微分分次 $A$-模结构。在这种情形下，存在一个典范的
左微分分次 $A$-模，其底层 $R$-模复形为
$\Hom^\bullet(M^\bullet, N^\bullet)$。具体而言，乘法可取为
\begin{align*}
\text{Tot}(A^\bullet \otimes_R \Hom^\bullet(M^\bullet, N^\bullet))
& \xrightarrow{\psi}
\text{Tot}(\Hom^\bullet(M^\bullet, N^\bullet) \otimes_R A^\bullet) \\
& \xrightarrow{\tau}
\text{Tot}(\Hom^\bullet(M^\bullet, N^\bullet) \otimes_R
\Hom^\bullet(M^\bullet, M^\bullet)) \\
& \to
\text{Tot}(\Hom^\bullet(M^\bullet, N^\bullet)
\end{align*}
第一个箭头使用 $R$-模复形范畴上的交换性约束，第二个箭头如上所述，
第三个箭头是 Hom 复形的复合法则。每个映射都是复形映射，因而所得映射
也是复形映射。事实上，这一构造使
$\Hom^\bullet(M^\bullet, N^\bullet)$ 成为左微分分次 $A$-模
（乘法的结合性可用对称幺半结构证明，也可使用下面的公式直接计算）。
把该乘法明确写为
$$
A^n \times \Hom^m(M^\bullet, N^\bullet) \longrightarrow
\Hom^{n + m}(M^\bullet, N^\bullet)
$$
。它把 $a \in A^n$ 以及满足 $f_{p, q} \in \Hom(M^{-q}, N^p)$ 的
$f = (f_{p, q})_{p + q = m}$ 送到元素 $a \cdot f$；后者的分量为
$$
(-1)^{nm}f_{p, q} \circ \tau_{-q, q + n}(a) =
(-1)^{nm - n(q + n)}f_{p, q} \circ a =
(-1)^{np + n} f_{p, q} \circ a
$$
，它位于 $\Hom_R(M^{-q - n}, N^p)$ 中；其中 $f_{p, q} \circ a$ 是映射
$$
M^{-q - n} \xrightarrow{a} M^{-q} \xrightarrow{f_{p, q}} N^p,\quad
x \longmapsto f_{p, q}(xa)
$$
。这里确实引入了符号 $(-1)^{np + n}$。用记号
$\Hom(M, N^\bullet)$ 表示这个左微分分次 $A$-模。

\begin{lemma}
\label{dga-lemma-characterize-hom-other-side}
设 $R$ 是环。设 $(A, \text{d})$ 是微分分次 $R$-代数。
设 $M$ 是右微分分次 $A$-模，$M'$ 是左微分分次 $A$-模，
$N^\bullet$ 是 $R$-模复形。那么有
$$
\Hom_{\text{左微分分次 }A\text{-模}}(M', \Hom(M, N^\bullet)) =
\Hom_{\text{Comp}(R)}(M \otimes_A M', N^\bullet)
$$
，其中如第 \ref{dga-section-tensor-product} 节那样，把
$M \otimes_A M'$ 看作 $R$-模复形。
\end{lemma}

\begin{proof}
来说明等式两边都对应于与微分相容的分次 $A$-双线性映射
$$
M \times M' \longrightarrow N^\bullet
$$
。第 \ref{dga-section-tensor-product} 节已经说明右边确实如此。
给定左边的一个元素 $g$，《更多代数》引理
\ref{more-algebra-lemma-compose} 中的等式确定复形映射
$g' : \text{Tot}(M' \otimes_R M) \to N^\bullet$。
预复合以交换性约束，得到
$$
\text{Tot}(M \otimes_R M') \xrightarrow{\psi}
\text{Tot}(M' \otimes_R M) \xrightarrow{g'}
N^\bullet
$$
，它对应于与微分相容的分次 $R$-双线性映射
$g'' : M \times M' \to N^\bullet$。$g$ 的 $A$-线性立即转化为
$g''$ 的 $A$-双线性。具体地，设 $x \in M^e$、
$x' \in (M')^{e'}$、$a \in A^n$。一方面有
\begin{align*}
g''(x, ax')
& =
(-1)^{e(n + e')} g'(ax' \otimes x) \\
& =
(-1)^{e(n + e')} g(ax')(x) \\
& =
(-1)^{e(n + e')} (a \cdot g(x'))(x) \\
& =
(-1)^{e(n + e') + n(n + e + e') + n} g(x')(xa)
\end{align*}
另一方面有
$$
g''(xa, x') = (-1)^{(e + n)e'} g'(x' \otimes xa) =
(-1)^{(e + n)e'} g(x')(xa) 
$$
对指数作一个简单的模 $2$ 计算，便知两者相同。
\end{proof}

\begin{remark}
\label{dga-remark-evaluation-map-left}
设 $R$ 是环，$A$ 是微分分次 $R$-代数，$M$ 是左微分分次 $A$-模，
$N^\bullet$ 是 $R$-模复形。上述构造先产生右微分分次 $A$-模
$\Hom(M, N^\bullet)$，再产生左微分分次 $A$-模
$\Hom(\Hom(M, N^\bullet), N^\bullet)$。我们断言存在求值映射
$$
ev : M \longrightarrow \Hom(\Hom(M, N^\bullet), N^\bullet)
$$
，它位于左微分分次 $A$-模范畴中。由引理
\ref{dga-lemma-characterize-hom}，为定义它，只需构造与分次及微分相容的
$A$-双线性配对
$$
\Hom(M, N^\bullet) \times M \longrightarrow N^\bullet
$$
。为此取
$$
(f, x) \longmapsto f(x)
$$
它与分次、微分及 $A$-双线性的相容性留给读者验证。$ev$ 在底层
$R$-模复形上的映射是《更多代数》条目
(\ref{more-algebra-item-evaluation})。
\end{remark}

\begin{remark}
\label{dga-remark-evaluation-map-right}
设 $R$ 是环，$A$ 是微分分次 $R$-代数，$M$ 是右微分分次 $A$-模，
$N^\bullet$ 是 $R$-模复形。上述构造先产生左微分分次 $A$-模
$\Hom(M, N^\bullet)$，再产生右微分分次 $A$-模
$\Hom(\Hom(M, N^\bullet), N^\bullet)$。我们断言存在求值映射
$$
ev : M \longrightarrow \Hom(\Hom(M, N^\bullet), N^\bullet)
$$
，它位于右微分分次 $A$-模范畴中。由引理
\ref{dga-lemma-characterize-hom}，为定义它，只需构造与分次及微分相容的
$A$-双线性配对
$$
M \times \Hom(M, N^\bullet) \longrightarrow N^\bullet
$$
。为此取
$$
(x, f) \longmapsto (-1)^{\deg(x)\deg(f)}f(x)
$$
它与分次、微分及 $A$-双线性的相容性留给读者验证。$ev$ 在底层
$R$-模复形上的映射是《更多代数》条目
(\ref{more-algebra-item-evaluation})。
\end{remark}

\begin{remark}
\label{dga-remark-shift-dual}
设 $R$ 是环，$A$ 是微分分次 $R$-代数，$M^\bullet$ 和 $N^\bullet$
是 $R$-模复形。设 $k \in \mathbf{Z}$，考虑 $R$-模复形同构
$$
\Hom^\bullet(M^\bullet, N^\bullet)[-k]
\longrightarrow
\Hom^\bullet(M^\bullet[k], N^\bullet)
$$
，它在《更多代数》条目 (\ref{more-algebra-item-shift-hom}) 中定义。
若 $M^\bullet$ 具有左（分别地，右）微分分次 $A$-模结构，
则这是右（分别地，左）微分分次 $A$-模的映射（模结构采用本节的定义）。
验证略。提醒读者：左分次 $A$-模的移位上的 $A$-模结构用一个符号定义；
见定义 \ref{dga-definition-shift-graded-module}。
\end{remark}






\section{代数上的投射模}
\label{dga-section-projectives-over-algebras}

\noindent
本节讨论代数上的投射模，类似于《代数》第
\ref{algebra-section-projective} 节。本节或许应移到别处。

\medskip\noindent
设 $R$ 是环，$A$ 是 $R$-代数；约定见第 \ref{dga-section-conventions} 节。
显然 $A$ 是投射右 $A$-模，因为对任意右 $A$-模 $M$ 都有
$\Hom_A(A, M) = M$（因而 $\Hom_A(A, -)$ 正合）。反之，设 $P$ 是投射
右 $A$-模。通过选择 $P$ 在 $A$ 上的一组生成元
$\{p_i\}_{i \in I}$，可以选取满射 $\bigoplus_{i \in I} A \to P$。
由于 $P$ 投射，该满射有左逆，故 $P$ 同构于某个自由模的直和项，
与交换情形完全相同（《代数》引理
\ref{algebra-lemma-characterize-projective}）。

\medskip\noindent
我们得到：
\begin{enumerate}
\item $A$-模范畴有足够多投射对象；
\item $A$ 是投射 $A$-模；
\item 每个 $A$-模都是若干个 $A$ 的直和的商；
\item 每个投射 $A$-模都是若干个 $A$ 的直和的一个直和项。
\end{enumerate}



\section{分次代数上的投射模}
\label{dga-section-projectives-over-graded-algebras}

\noindent
本节讨论分次代数上的投射分次模，类似于《代数》第
\ref{algebra-section-projective} 节。

\medskip\noindent
设 $R$ 是环。设 $A$ 是 $R$ 上的 $\mathbf{Z}$-分次代数。
约定见第 \ref{dga-section-conventions} 节。用 $\text{Mod}_A$ 表示分次右
$A$-模范畴。对整数 $k$，用 $A[k]$ 表示 $A$ 的移位。
对于分次右 $A$-模，有
$$
\Hom_{\text{Mod}_A}(A[k], M) = M^{-k}
$$
由于函子 $M \mapsto M^{-k}$ 在 $\text{Mod}_A$ 上正合，
$A[k]$ 是 $\text{Mod}_A$ 的投射对象。反之，假设 $P$ 是
$\text{Mod}_A$ 的投射对象。选择 $P$ 作为 $A$-模的一组齐次生成元，
可以找到满射
$$
\bigoplus\nolimits_{i \in I} A[k_i] \longrightarrow P
$$
因此，$\text{Mod}_A$ 的投射对象是若干个移位 $A[k]$ 的直和的一个直和项。

\medskip\noindent
我们得到：
\begin{enumerate}
\item 分次 $A$-模范畴有足够多投射对象；
\item 对每个 $k \in \mathbf{Z}$，$A[k]$ 都是投射 $A$-模；
\item 每个分次 $A$-模都是若干个模 $A[k]$（$k$ 可变）的直和的商；
\item 每个投射 $A$-模都是若干个模 $A[k]$（$k$ 可变）的直和的一个直和项。
\end{enumerate}




\section{投射模与微分分次代数}
\label{dga-section-projective-over-differential-graded}

\noindent
若 $(A, \text{d})$ 是微分分次代数，而 $P$ 是
$\text{Mod}_{(A, \text{d})}$ 的对象，则称 {\it $P$ 作为分次 $A$-模是投射的}，
有时也称 {\it $P$ 是分次投射的}；这意指 $P$ 是第
\ref{dga-section-projectives-over-graded-algebras} 节中分次 $A$-模的阿贝尔范畴
$\text{Mod}_A$ 的投射对象。

\begin{lemma}
\label{dga-lemma-target-graded-projective}
设 $(A, \text{d})$ 是微分分次代数。设 $M \to P$ 是微分分次 $A$-模的
满同态。若 $P$ 作为分次 $A$-模是投射的，则 $M \to P$ 是容许满态射。
\end{lemma}

\begin{proof}
这立即由定义得出。
\end{proof}

\begin{lemma}
\label{dga-lemma-hom-from-shift-free}
设 $(A, d)$ 是微分分次代数。那么，对任意微分分次 $A$-模 $M$，有
$$
\Hom_{\text{Mod}_{(A, \text{d})}}(A[k], M) =
\Ker(\text{d} : M^{-k} \to M^{-k + 1})
$$
并且
$$
\Hom_{K(\text{Mod}_{(A, \text{d})})}(A[k], M) = H^{-k}(M)
$$
。
\end{lemma}

\begin{proof}
立即由定义得出。
\end{proof}







\section{代数上的内射模}
\label{dga-section-modules-noncommutative}

\noindent
本节讨论代数上的内射模，类似于《更多代数》第
\ref{more-algebra-section-injectives-modules} 节。本节或许应移到别处。

\medskip\noindent
设 $R$ 是环，$A$ 是 $R$-代数；约定见第 \ref{dga-section-conventions} 节。
对右 $A$-模 $M$，置
$$
M^\vee = \Hom_\mathbf{Z}(M, \mathbf{Q}/\mathbf{Z})
$$
并借助乘法 $(a f)(x) = f(xa)$ 把它看作左 $A$-模。事实上，
$((ab)f)(x) = f(xab) = (bf)(xa) = (a(bf))(x)$。反之，若 $M$ 是左
$A$-模，则 $M^\vee$ 是右 $A$-模。由于 $\mathbf{Q}/\mathbf{Z}$
是内射阿贝尔群（《更多代数》引理
\ref{more-algebra-lemma-injective-abelian}），函子 $M \mapsto M^\vee$
是正合的（《更多代数》引理 \ref{more-algebra-lemma-vee-exact}）。
此外，对所有模 $M$，求值映射 $M \to (M^\vee)^\vee$ 是单射
（《更多代数》引理 \ref{more-algebra-lemma-ev-injective}）。

\medskip\noindent
我们断言 $A^\vee$ 是内射右 $A$-模。事实上，给定右 $A$-模 $N$，有
$$
\Hom_A(N, A^\vee) =
\Hom_A(N, \Hom_\mathbf{Z}(A, \mathbf{Q}/\mathbf{Z})) = N^\vee
$$
；结论由函子 $N \mapsto N^\vee$ 的正合性得出。第二个等式成立，是因为
$$
\Hom_\mathbf{Z}(N, \Hom_\mathbf{Z}(A, \mathbf{Q}/\mathbf{Z})) =
\Hom_\mathbf{Z}(N \otimes_\mathbf{Z} A, \mathbf{Q}/\mathbf{Z})
$$
，见《代数》引理 \ref{algebra-lemma-hom-from-tensor-product}。
在这个模中，$A$-线性恰好挑出双线性映射
$\varphi : N \times A \to \mathbf{Q}/\mathbf{Z}$，它们在
$x \otimes a$ 和 $xa \otimes 1$ 上取相同值；换言之，它们来自
$N^\vee$ 的元素。

\medskip\noindent
最后，对每个右 $A$-模 $M$，可选择满射
$\bigoplus_{i \in I} A \to M^\vee$，从而得到单射
$M \to (M^\vee)^\vee \to \prod_{i \in I} A^\vee$。

\medskip\noindent
我们得到：
\begin{enumerate}
\item $A$-模范畴有足够多内射对象；
\item $A^\vee$ 是内射 $A$-模；以及
\item 每个 $A$-模都单射到若干个 $A^\vee$ 的乘积中。
\end{enumerate}





\section{分次代数上的内射模}
\label{dga-section-modules-noncommutative-graded}

\noindent
本节讨论分次代数上的内射分次模，类似于《更多代数》第
\ref{more-algebra-section-injectives-modules} 节。

\medskip\noindent
设 $R$ 是环。设 $A$ 是 $R$ 上的 $\mathbf{Z}$-分次代数。
约定见第 \ref{dga-section-conventions} 节。若 $M$ 是分次 $R$-模，置
$$
M^\vee =
\bigoplus\nolimits_{n \in \mathbf{Z}}
\Hom_\mathbf{Z}(M^{-n}, \mathbf{Q}/\mathbf{Z}) =
\bigoplus\nolimits_{n \in \mathbf{Z}} (M^{-n})^\vee
$$
，把它视为分次 $R$-模（$R$ 在齐次部分上的作用不带符号）。
若 $M$ 具有左分次 $A$-模结构，则如下定义 $M^\vee$ 上的右分次
$A$-模结构：令 $a \in A^m$ 按
$$
(M^{-n})^\vee \to (M^{-n - m})^\vee, \quad
f \mapsto f \circ a
$$
作用，如第 \ref{dga-section-hom-complexes} 节所述。若 $M$ 具有右分次
$A$-模结构，则如下定义 $M^\vee$ 上的左分次 $A$-模结构：
令 $a \in A^n$ 按
$$
(M^{-m})^\vee \to (M^{-m - n})^\vee, \quad
f \mapsto (-1)^{nm}f \circ a
$$
作用，如第 \ref{dga-section-hom-complexes} 节所述（这个符号是被迫的，
因为研究微分分次模时我们希望使用同一个公式；若只关心分次模，
则可在此省略该符号）。在（左或右）分次 $A$-模范畴上，函子
$M \mapsto M^\vee$ 正合（在各分次部分上验证）。此外，有单射的求值映射
$$
ev : M \longrightarrow (M^\vee)^\vee, \quad
ev^n = (-1)^n \text{ 倍的求值映射 }M^n \to ((M^n)^\vee)^\vee
$$
，它是分次 $R$-模的映射；见《更多代数》条目
(\ref{more-algebra-item-evaluation})。若 $M$ 是左（分别地，右）$A$-模，
则这个求值映射是左（分别地，右）$A$-模同态；见评注
\ref{dga-remark-evaluation-map-left} 和 \ref{dga-remark-evaluation-map-right}。
最后，给定 $k \in \mathbf{Z}$，有典范同构
$$
M^\vee[-k] \longrightarrow (M[k])^\vee
$$
；它是使用一个符号的分次 $R$-模同构。若 $M$ 是左（分别地，右）
$A$-模，则它是右（分别地，左）$A$-模的同构。见评注
\ref{dga-remark-shift-dual}。

\medskip\noindent
我们断言 $A^\vee$ 是分次右 $A$-模范畴 $\text{Mod}_A$ 的内射对象。
事实上，给定分次右 $A$-模 $N$，有
$$
\Hom_{\text{Mod}_A}(N, A^\vee) =
\Hom_{\text{Comp}(\mathbf{Z})}(N \otimes_A A, \mathbf{Q}/\mathbf{Z})) =
(N^0)^\vee
$$
；这是引理 \ref{dga-lemma-characterize-hom} 应用于所有微分均为零的情形。
结论由函子 $N \mapsto (N^0)^\vee = (N^\vee)^0$ 的正合性得出。

\medskip\noindent
最后，对每个分次右 $A$-模 $M$，可以选择分次左 $A$-模的满射
$$
\bigoplus\nolimits_{i \in I} A[k_i] \to M^\vee
$$
，其中 $A[k_i]$ 表示 $A$ 由 $k_i \in \mathbf{Z}$ 给出的移位。
为此，选择 $M^\vee$ 的齐次生成元。这样得到单射
$$
M \to (M^\vee)^\vee \to \prod A[k_i]^\vee = \prod A^\vee[-k_i]
$$
注意，上式中的乘积是分次模范畴中的乘积（换言之，先在每个次数上取乘积，
再对各部分取直和）。

\medskip\noindent
我们得到：
\begin{enumerate}
\item 分次 $A$-模范畴有足够多内射对象；
\item 对每个 $k \in \mathbf{Z}$，模 $A^\vee[k]$ 都是内射的；以及
\item 每个 $A$-模都单射到分次模范畴中若干个移位 $A^\vee[k]$ 的乘积中。
\end{enumerate}





\section{内射模与微分分次代数}
\label{dga-section-modules-noncommutative-differential-graded}

\noindent
若 $(A, \text{d})$ 是微分分次代数，而 $I$ 是
$\text{Mod}_{(A, \text{d})}$ 的对象，则称 {\it $I$ 作为分次 $A$-模是内射的}，
有时也称 {\it $I$ 是分次内射的}；这意指 $I$ 是分次 $A$-模的阿贝尔范畴
$\text{Mod}_A$ 的内射对象。

\begin{lemma}
\label{dga-lemma-source-graded-injective}
设 $(A, \text{d})$ 是微分分次代数。设 $I \to M$ 是微分分次 $A$-模的
单同态。若 $I$ 是分次内射的，则 $I \to M$ 是容许单态射。
\end{lemma}

\begin{proof}
这立即由定义得出。
\end{proof}

\noindent
设 $(A, \text{d})$ 是微分分次代数。若 $M$ 是左（分别地，右）微分分次
$A$-模，则
$$
M^\vee = \Hom^\bullet(M^\bullet, \mathbf{Q}/\mathbf{Z})
$$
连同第 \ref{dga-section-modules-noncommutative-graded} 节中构造的 $A$-模结构，
由第 \ref{dga-section-hom-complexes} 节的讨论，是右（分别地，左）微分分次
$A$-模。由评注 \ref{dga-remark-evaluation-map-left} 和
\ref{dga-remark-evaluation-map-right}，第
\ref{dga-section-modules-noncommutative-graded} 节的求值映射
$$
M \longrightarrow (M^\vee)^\vee
$$
是左（分别地，右）微分分次 $A$-模的同态。

\begin{lemma}
\label{dga-lemma-map-into-dual}
设 $(A, \text{d})$ 是微分分次代数。若 $M$ 是左微分分次 $A$-模，
$N$ 是右微分分次 $A$-模，则
\begin{align*}
\Hom_{\text{Mod}_{(A, \text{d})}}(N, M^\vee)
& =
\Hom_{\text{Comp}(\mathbf{Z})}(N \otimes_A M, \mathbf{Q}/\mathbf{Z}) \\
& =
\text{微分分次双线性}_A(N \times M, \mathbf{Q}/\mathbf{Z})
\end{align*}
\end{lemma}

\begin{proof}
第一个等式是引理 \ref{dga-lemma-characterize-hom}，第二个等式已在第
\ref{dga-section-tensor-product} 节中证明。
\end{proof}

\begin{lemma}
\label{dga-lemma-hom-into-shift-dual-free}
设 $(A, \text{d})$ 是微分分次代数。那么有
$$
\Hom_{\text{Mod}_{(A, \text{d})}}(M, A^\vee[k]) =
\Ker(\text{d} : (M^\vee)^k \to (M^\vee)^{k + 1})
$$
并且
$$
\Hom_{K(\text{Mod}_{(A, \text{d})})}(M, A^\vee[k]) = H^k(M^\vee)
$$
；它们是关于微分分次 $A$-模 $M$ 的函子等式。
\end{lemma}

\begin{proof}
这由上面的讨论显然可见。
\end{proof}















\section{P-分辨}
\label{dga-section-P-resolutions}

\noindent
本节是《导出范畴》第 \ref{derived-section-unbounded} 节的对应版本。

\medskip\noindent
设 $(A, \text{d})$ 是微分分次代数。设 $P$ 是微分分次 $A$-模。
若存在滤过
$$
0 = F_{-1}P \subset F_0P \subset F_1P \subset \ldots \subset P
$$
，它由微分分次子模组成，并且
\begin{enumerate}
\item $P = \bigcup F_pP$；
\item 包含映射 $F_iP \to F_{i + 1}P$ 是容许单态射；
\item 作为微分分次 $A$-模，商 $F_{i + 1}P/F_iP$ 同构于若干个
$A[k]$ 的直和，
\end{enumerate}
则称 $P$ {\it 具有性质 (P)}。事实上，条件 (2) 是条件 (3) 的推论；
见引理 \ref{dga-lemma-target-graded-projective}。此外，读者可以验证，
作为分次 $A$-模，$P$ 同构于若干个 $A$ 的移位的直和。

\begin{lemma}
\label{dga-lemma-property-P-sequence}
设 $(A, \text{d})$ 是微分分次代数。设 $P$ 是微分分次 $A$-模。
若 $F_\bullet$ 是性质 (P) 中的滤过，则得到容许短正合列
$$
0 \to
\bigoplus\nolimits F_iP \to
\bigoplus\nolimits F_iP \to P \to 0
$$
，它由微分分次 $A$-模组成。
\end{lemma}

\begin{proof}
第二个映射是各包含映射的直和。第一个映射在源的直和项 $F_iP$ 上，
等于恒等映射 $F_iP \to F_iP$ 与包含映射
$F_iP \to F_{i + 1}P$ 的负映射之和。选择分次 $A$-模同态
$s_i : F_{i + 1}P \to F_iP$，使它们是包含映射的左逆。
复合后，对所有 $j > i$ 得到映射 $s_{j, i} : F_jP \to F_iP$。
于是，第一个箭头的一个左逆把 $x \in F_jP$ 映到
$\bigoplus F_iP$ 中的
$(s_{j, 0}(x), s_{j, 1}(x), \ldots, s_{j, j - 1}(x), 0, \ldots)$。
\end{proof}

\noindent
下述引理表明，具有性质 (P) 的微分分次模是 K-内射模的对偶概念
（即在某种意义下它们是 K-投射的）。见《导出范畴》定义
\ref{derived-definition-K-injective}。

\begin{lemma}
\label{dga-lemma-property-P-K-projective}
设 $(A, \text{d})$ 是微分分次代数。设 $P$ 是具有性质 (P) 的微分分次
$A$-模。那么
$$
\Hom_{K(\text{Mod}_{(A, \text{d})})}(P, N) = 0
$$
对所有无环微分分次 $A$-模 $N$ 成立。
\end{lemma}

\begin{proof}
将使用 $K(\text{Mod}_{(A, \text{d})})$ 是三角范畴这一事实
（命题 \ref{dga-proposition-homotopy-category-triangulated}）。
设 $F_\bullet$ 是 $P$ 上性质 (P) 中的滤过。引理
\ref{dga-lemma-property-P-sequence} 的短正合列产生一个有区别三角。
因此，由《导出范畴》引理
\ref{derived-lemma-representable-homological}，只需证明
$$
\Hom_{K(\text{Mod}_{(A, \text{d})})}(F_iP, N) = 0
$$
对所有无环微分分次 $A$-模 $N$ 和所有 $i$ 成立。每个微分分次模
$F_iP$ 都具有由容许单态射给出的有限滤过，其分次部分是若干个移位
$A[k]$ 的直和。因此，只需证明
$$
\Hom_{K(\text{Mod}_{(A, \text{d})})}(A[k], N) = 0
$$
对所有无环微分分次 $A$-模 $N$ 和所有 $k$ 成立。这由引理
\ref{dga-lemma-hom-from-shift-free} 得出。
\end{proof}

\begin{lemma}
\label{dga-lemma-good-quotient}
设 $(A, \text{d})$ 是微分分次代数。设 $M$ 是微分分次 $A$-模。
存在具有下列性质的微分分次 $A$-模同态 $P \to M$：
\begin{enumerate}
\item $P \to M$ 是满射；
\item $\Ker(\text{d}_P) \to \Ker(\text{d}_M)$ 是满射；以及
\item $P$ 位于容许短正合列 $0 \to P' \to P \to P'' \to 0$ 中，
其中 $P'$、$P''$ 是若干个 $A$ 的移位的直和。
\end{enumerate}
\end{lemma}

\begin{proof}
令 $P_k$ 为自由 $A$-模，其生成元 $x,y$ 的次数分别为 $k$ 和 $k+1$。
通过置 $\text{d}(x) = y$、$\text{d}(y) = 0$，在 $P_k$ 上定义微分分次
$A$-模结构。对每个元素 $m \in M^k$，都有同态 $P_k \to M$，
它把 $x$ 映到 $m$，把 $y$ 映到 $\text{d}(m)$。因此，若干个 $P_k$
的直和到 $M$ 有一个满射。这显然产生具有性质 (1) 和 (3) 的
$P \to M$。为得到性质 (2)，注意：若
$m \in \Ker(\text{d}_M)$ 的次数为 $k$，则有把 $1$ 映到 $m$ 的映射
$A[k] \to M$。因此，加入若干个 $A$ 的移位的直和即可满足 (2)。
\end{proof}

\begin{lemma}
\label{dga-lemma-resolve}
设 $(A, \text{d})$ 是微分分次代数。设 $M$ 是微分分次 $A$-模。
存在微分分次 $A$-模同态 $P \to M$，使得
\begin{enumerate}
\item $P \to M$ 是拟同构；并且
\item $P$ 具有性质 (P)。
\end{enumerate}
\end{lemma}

\begin{proof}
置 $M=M_0$。归纳选取短正合列
$$
0 \to M_{i + 1} \to P_i \to M_i \to 0
$$
，其中映射 $P_i \to M_i$ 按引理 \ref{dga-lemma-good-quotient} 选取。
这给出一个“分辨”
$$
\ldots \to P_2 \xrightarrow{f_2} P_1 \xrightarrow{f_1} P_0 \to M \to 0
$$
然后置
$$
P = \bigoplus\nolimits_{i \geq 0} P_i
$$
，把它视为 $A$-模，其分次由
$P^n = \bigoplus_{a + b = n} P_{-a}^b$ 给出，微分则如双复形关联全复形的
构造那样，由
$$
\text{d}_P(x) = f_{-a}(x) + (-1)^a \text{d}_{P_{-a}}(x)
$$
给出，其中 $x \in P_{-a}^b$。在这些约定下，$P$ 确实是微分分次
$A$-模。回忆每个 $P_i$ 都有两步滤过
$0 \to P_i' \to P_i \to P_i'' \to 0$，置
$$
F_{2i}P = \bigoplus\nolimits_{i \geq j \geq 0} P_j
\subset
\bigoplus\nolimits_{i \geq 0} P_i = P
$$
并把 $P'_{i + 1}$ 加到 $F_{2i}P$ 中得到 $F_{2i + 1}$。
这些都是微分分次子模，且相继的商是若干个 $A$ 的移位的直和。
由引理 \ref{dga-lemma-target-graded-projective}，包含映射
$F_iP \to F_{i + 1}P$ 是容许单态射。最后，还须证明由增广
$P_0 \to M$ 给出的映射 $P \to M$ 是拟同构。这由《同调代数》引理
\ref{homology-lemma-good-resolution-gives-qis} 得出。
\end{proof}





\section{I-分辨}
\label{dga-section-I-resolutions}

\noindent
本节是 P-分辨一节的对偶。

\medskip\noindent
设 $(A, \text{d})$ 是微分分次代数。设 $I$ 是微分分次 $A$-模。
若存在滤过
$$
I = F_0I \supset F_1I \supset F_2I \supset \ldots \supset 0
$$
，它由微分分次子模组成，并且
\begin{enumerate}
\item $I = \lim I/F_pI$；
\item 映射 $I/F_{i + 1}I \to I/F_iI$ 是容许满态射；
\item 作为微分分次 $A$-模，商 $F_iI/F_{i + 1}I$ 同构于第
\ref{dga-section-modules-noncommutative-differential-graded} 节构造的若干个
$A^\vee[k]$ 的直积，
\end{enumerate}
则称 $I$ {\it 具有性质 (I)}。事实上，条件 (2) 是条件 (3) 的推论；见引理
\ref{dga-lemma-source-graded-injective}。读者可以验证，作为分次模，$I$ 同构于
若干个 $A^\vee[k]$ 的直积。

\begin{lemma}
\label{dga-lemma-property-I-sequence}
设 $(A, \text{d})$ 是微分分次代数。设 $I$ 是微分分次 $A$-模。
若 $F_\bullet$ 是性质 (I) 中的滤过，则得到容许短正合列
$$
0 \to I \to
\prod\nolimits I/F_iI \to
\prod\nolimits I/F_iI \to 0
$$
，它由微分分次 $A$-模组成。
\end{lemma}

\begin{proof}
略。提示：这是引理 \ref{dga-lemma-property-P-sequence} 的对偶。
\end{proof}

\noindent
下述引理表明，具有性质 (I) 的微分分次模是 K-内射模的对应概念。见
《导出范畴》定义 \ref{derived-definition-K-injective}。

\begin{lemma}
\label{dga-lemma-property-I-K-injective}
设 $(A, \text{d})$ 是微分分次代数。设 $I$ 是具有性质 (I) 的微分分次
$A$-模。那么
$$
\Hom_{K(\text{Mod}_{(A, \text{d})})}(N, I) = 0
$$
对所有无环微分分次 $A$-模 $N$ 成立。
\end{lemma}

\begin{proof}
将使用 $K(\text{Mod}_{(A, \text{d})})$ 是三角范畴这一事实
（命题 \ref{dga-proposition-homotopy-category-triangulated}）。设 $F_\bullet$
是 $I$ 上性质 (I) 中的滤过。引理 \ref{dga-lemma-property-I-sequence}
的短正合列产生一个有区别三角。因此，由《导出范畴》引理
\ref{derived-lemma-representable-homological}，只需证明
$$
\Hom_{K(\text{Mod}_{(A, \text{d})})}(N, I/F_iI) = 0
$$
对所有无环微分分次 $A$-模 $N$ 和所有 $i$ 成立。每个微分分次模
$I/F_iI$ 都具有由容许单态射给出的有限滤过，其分次部分是若干个
$A^\vee[k]$ 的直积。因此，只需证明
$$
\Hom_{K(\text{Mod}_{(A, \text{d})})}(N, A^\vee[k]) = 0
$$
对所有无环微分分次 $A$-模 $N$ 和所有 $k$ 成立。这由引理
\ref{dga-lemma-hom-into-shift-dual-free} 以及 $(-)^\vee$ 是正合函子这一事实得出。
\end{proof}

\begin{lemma}
\label{dga-lemma-good-sub}
设 $(A, \text{d})$ 是微分分次代数。设 $M$ 是微分分次 $A$-模。
存在具有下列性质的微分分次 $A$-模同态 $M \to I$：
\begin{enumerate}
\item $M \to I$ 是单射；
\item $\Coker(\text{d}_M) \to \Coker(\text{d}_I)$ 是单射；以及
\item $I$ 位于容许短正合列 $0 \to I' \to I \to I'' \to 0$ 中，
其中 $I'$、$I''$ 是若干个 $A^\vee$ 的移位的直积。
\end{enumerate}
\end{lemma}

\begin{proof}
将使用第 \ref{dga-section-modules-noncommutative-differential-graded} 节构造的
函子 $N \mapsto N^\vee$（从左微分分次模到右微分分次模，以及从右微分
分次模到左微分分次模）及其全部性质。对每个 $k \in \mathbf{Z}$，令
$Q_k$ 为自由左 $A$-模，其生成元 $x,y$ 的次数分别为 $k$ 和 $k+1$。
通过置 $\text{d}(x) = y$、$\text{d}(y) = 0$，在 $Q_k$ 上定义左微分分次
$A$-模结构。完全仿照引理 \ref{dga-lemma-good-quotient} 的证明，可得左微分分次
$A$-模的满射
$$
\bigoplus\nolimits_{i \in I} Q_{k_i} \longrightarrow M^\vee
$$
。于是可以考虑单射
$$
M \to (M^\vee)^\vee \to (\bigoplus\nolimits_{i \in I} Q_{k_i})^\vee =
\prod\nolimits_{i \in I} I_{k_i}
$$
，其中 $I_k = Q_{-k}^\vee$ 是“对偶”右微分分次 $A$-模。此外，短正合列
$0 \to A[-k - 1] \to Q_k \to A[-k] \to 0$ 产生短正合列
$0 \to A^\vee[k] \to I_k \to A^\vee[k + 1] \to 0$.

\medskip\noindent
上一段的结果给出具有性质 (1) 和 (3) 的 $M \to I$。为得到性质 (2)，
设 $\overline{m} \in \Coker(\text{d}_M)$ 是次数为 $k$ 的非零元素。
选取映射 $\lambda : M^k \to \mathbf{Q}/\mathbf{Z}$，使它在
$\Im(M^{k - 1} \to M^k)$ 上为零，而在 $m$ 上不为零。由引理
\ref{dga-lemma-hom-into-shift-dual-free}，这对应于在 $m$ 上不为零的微分分次
$A$-模同态 $M \to A^\vee[k]$。因此，加入若干个 $A^\vee$ 的移位的副本
的直积，即可满足 (2)。
\end{proof}

\begin{lemma}
\label{dga-lemma-right-resolution}
设 $(A, \text{d})$ 是微分分次代数。设 $M$ 是微分分次 $A$-模。
存在微分分次 $A$-模同态 $M \to I$，使得
\begin{enumerate}
\item $M \to I$ 是拟同构；并且
\item $I$ 具有性质 (I)。
\end{enumerate}
\end{lemma}

\begin{proof}
置 $M=M_0$。归纳选取短正合列
$$
0 \to M_i \to I_i \to M_{i + 1} \to 0
$$
，其中映射 $M_i \to I_i$ 按引理 \ref{dga-lemma-good-sub} 选取。这给出一个“分辨”
$$
0 \to M \to I_0 \xrightarrow{f_0} I_1 \xrightarrow{f_1} I_1 \to \ldots
$$
以 $I$ 表示微分分次 $A$-模，其分次部分为
$$
I^n = \prod\nolimits_{i \geq 0} I^{n - i}_i
$$
，微分定义为
$$
\text{d}_I(x) = f_i(x) + (-1)^i \text{d}_{I_i}(x)
$$
，其中 $x \in I_i^{n-i}$。在这些约定下，$I$ 确实是微分分次 $A$-模。
回忆每个 $I_i$ 都有两步滤过 $0 \to I_i' \to I_i \to I_i'' \to 0$，置
$$
F_{2i}I^n = \prod\nolimits_{j \geq i} I^{n - j}_j
\subset
\prod\nolimits_{i \geq 0} I^{n - i}_i = I^n
$$
，并把因子 $I'_{i+1}$ 加到 $F_{2i}I$ 中得到 $F_{2i+1}I$。这些都是微分分次
子模，且相继的商是若干个 $A^\vee$ 的移位的直积。由引理
\ref{dga-lemma-source-graded-injective}，包含映射 $F_{i+1}I \to F_iI$ 是容许
单态射。最后，还须证明由增广 $M \to I_0$ 给出的映射 $M \to I$
是拟同构。这由《同调代数》引理
\ref{homology-lemma-good-right-resolution-gives-qis} 得出。
\end{proof}






\section{导出范畴}
\label{dga-section-derived}

\noindent
回忆，无环微分分次模和微分分次模之间的拟同构这两个概念均有定义
（见第 \ref{dga-section-modules} 节）。

\begin{lemma}
\label{dga-lemma-acyclic}
设 $(A, \text{d})$ 是微分分次代数。由无环模组成的
$K(\text{Mod}_{(A, \text{d})})$ 的满子范畴 $\text{Ac}$，是
$K(\text{Mod}_{(A, \text{d})})$ 的严格满且饱和的三角子范畴。相应的
$K(\text{Mod}_{(A, \text{d})})$ 的饱和乘法系（见《导出范畴》引理
\ref{derived-lemma-operations}）是拟同构类 $\text{Qis}$。特别地，局部化函子
$$
Q : K(\text{Mod}_{(A, \text{d})}) \to
\text{Qis}^{-1}K(\text{Mod}_{(A, \text{d})})
$$
的核是 $\text{Ac}$。此外，函子 $H^0$ 通过 $Q$ 分解。
\end{lemma}

\begin{proof}
由同调长正合列 (\ref{dga-equation-les})，知道 $H^0$ 是同调函子。
$H^0$ 的核是无环对象的子范畴，而在所有 $H^i$ 上诱导同构的箭头正是
拟同构。因此，本引理是《导出范畴》引理
\ref{derived-lemma-acyclic-general} 的特例。

\medskip\noindent
集合论说明。《导出范畴》命题
\ref{derived-proposition-construct-localization} 中的局部化构造，假定给定的
三角范畴是“小”的，即其底层对象族构成一个集合。令 $V_\alpha$ 是包含
$(A, \text{d})$ 的部分宇宙（如《集合》第
\ref{sets-section-sets-hierarchy} 节），并且 $\alpha$ 的共尾数大于
$\aleph_0$（见《集合》命题
\ref{sets-proposition-exist-ordinals-large-cofinality}）。于是可以考虑由
$V_\alpha$ 中的微分分次 $A$-模组成的范畴
$\text{Mod}_{(A, \text{d}), \alpha}$。直接核查可见，命题
\ref{dga-proposition-homotopy-category-triangulated} 的证明中使用的全部构造
都能在 $\text{Mod}_{(A, \text{d}), \alpha}$ 内进行（因为至多只需取微分
分次模的有限直和）。因此得到三角范畴
$\text{Qis}_\alpha^{-1}K(\text{Mod}_{(A, \text{d}), \alpha})$。
下文将看到，若 $\beta>\alpha$，则过渡函子
$$
\text{Qis}_\alpha^{-1}K(\text{Mod}_{(A, \text{d}), \alpha})
\longrightarrow
\text{Qis}_\beta^{-1}K(\text{Mod}_{(A, \text{d}), \beta})
$$
是全忠实的，因为商范畴中的态射集由同伦范畴中从 P-分辨出发的映射计算
（引理 \ref{dga-lemma-resolve} 的证明中构造 P-分辨时，要取可数直和以及由
给定模的子集作指标的直和）。因此，读者应把引理中的范畴视为这些子范畴
的并。
\end{proof}

\noindent
考虑到前一引理证明末尾的集合论说明，按如下方式定义导出范畴。

\begin{definition}
\label{dga-definition-unbounded-derived-category}
设 $(A, \text{d})$ 是微分分次代数。令 $\text{Ac}$ 和 $\text{Qis}$ 如引理
\ref{dga-lemma-acyclic} 所述。$(A, \text{d})$ 的{\it 导出范畴}是三角范畴
$$
D(A, \text{d}) =
K(\text{Mod}_{(A, \text{d})})/\text{Ac} =
\text{Qis}^{-1}K(\text{Mod}_{(A, \text{d})}).
$$
以 $H^0 : D(A, \text{d}) \to \text{Mod}_R$ 表示唯一的函子，使得它与
商函子的复合正好给出上文定义的函子 $H^0$。
\end{definition}

\noindent
下面是此前所称的、计算导出范畴中态射集的引理。

\begin{lemma}
\label{dga-lemma-hom-derived}
设 $(A, \text{d})$ 是微分分次代数。设 $M$ 和 $N$ 是微分分次 $A$-模。
\begin{enumerate}
\item 设 $P \to M$ 是引理 \ref{dga-lemma-resolve} 中的 P-分辨。那么
$$
\Hom_{D(A, \text{d})}(M, N) =
\Hom_{K(\text{Mod}_{(A, \text{d})})}(P, N)
$$
\item 设 $N \to I$ 是引理 \ref{dga-lemma-right-resolution} 中的 I-分辨。那么
$$
\Hom_{D(A, \text{d})}(M, N) =
\Hom_{K(\text{Mod}_{(A, \text{d})})}(M, I)
$$
\end{enumerate}
\end{lemma}

\begin{proof}
设 $P \to M$ 如 (1) 中所述。由于 $P \to M$ 是拟同构，由导出范畴的定义有
$$
\Hom_{D(A, \text{d})}(P, N) = \Hom_{D(A, \text{d})}(M, N)
$$
。$D(A, \text{d})$ 中的态射 $f:P\to N$ 等于 $s^{-1}f'$，其中
$f':P\to N'$ 是态射，而 $s:N\to N'$ 是拟同构。选择有区别三角
$$
N \to N' \to Q \to N[1]
$$
。由于 $s$ 是拟同构，$Q$ 是无环的。因此，由引理
\ref{dga-lemma-property-P-K-projective}，对所有 $k$ 都有
$\Hom_{K(\text{Mod}_{(A, \text{d})})}(P, Q[k]) = 0$。由于
$\Hom_{K(\text{Mod}_{(A, \text{d})})}(P, -)$
是上同调的，故 $f':P\to N'$ 可唯一提升为态射 $f:P\to N$。这就完成了证明。

\medskip\noindent
(2) 的证明是 (1) 的对偶：使用引理 \ref{dga-lemma-property-I-K-injective}
代替引理 \ref{dga-lemma-property-P-K-projective} 即可。
\end{proof}

\begin{lemma}
\label{dga-lemma-derived-products}
设 $(A, \text{d})$ 是微分分次代数。那么
\begin{enumerate}
\item $D(A, \text{d})$ 既有直和，也有直积；
\item 直和可由微分分次模的直和得到；
\item 直积可由微分分次模的直积得到。
\end{enumerate}
\end{lemma}

\begin{proof}
将使用如下事实：$\text{Mod}_{(A, \text{d})}$ 是具有任意直和与直积的
阿贝尔范畴，并且它们在 $K(\text{Mod}_{(A, \text{d})})$ 中产生直和与
直积。见引理 \ref{dga-lemma-dgm-abelian} 和
\ref{dga-lemma-homotopy-direct-sums}。

\medskip\noindent
设 $M_j$ 是一族微分分次 $A$-模。考虑分次直和
$M=\bigoplus M_j$；它是具有显然的……的微分分次 $A$-模。对微分分次
$A$-模 $N$，选取拟同构 $N\to I$，其中 $I$ 是具有性质 (I) 的微分分次
$A$-模；见引理 \ref{dga-lemma-right-resolution}。利用引理
\ref{dga-lemma-hom-derived}，有
\begin{align*}
\Hom_{D(A, \text{d})}(M, N)
& =
\Hom_{K(A, \text{d})}(M, I) \\
& =
\prod \Hom_{K(A, \text{d})}(M_j, I) \\
& =
\prod \Hom_{D(A, \text{d})}(M_j, N)
\end{align*}
由此得到 $D(A, \text{d})$ 中的直和，正如引理第 (2) 项所述。

\medskip\noindent
设 $M_j$ 是一族微分分次 $A$-模。考虑微分分次 $A$-模的直积
$M=\prod M_j$。对微分分次 $A$-模 $N$，选取拟同构 $P\to N$，
其中 $P$ 是具有性质 (P) 的微分分次 $A$-模；见引理
\ref{dga-lemma-resolve}。利用引理 \ref{dga-lemma-hom-derived}，有
\begin{align*}
\Hom_{D(A, \text{d})}(N, M)
& =
\Hom_{K(A, \text{d})}(P, M) \\
& =
\prod \Hom_{K(A, \text{d})}(P, M_j) \\
& =
\prod \Hom_{D(A, \text{d})}(N, M_j)
\end{align*}
由此得到 $D(A, \text{d})$ 中的直和，正如引理第 (3) 项所述。
\end{proof}

\begin{remark}
\label{dga-remark-P-resolution}
设 $R$ 是环，$(A, \text{d})$ 是微分分次 $R$-代数。利用 P-分辨，有时可将
关于 $D(A, \text{d})$ 中一般对象的命题约化为关于 $A[k]$ 的命题。
具体地，设 $T$ 是 $D(A, \text{d})$ 中对象的一个性质，并假设
\begin{enumerate}
\item 若 $K_i$（$i\in I$）是 $D(A, \text{d})$ 中的一族对象，且对所有
$i\in I$ 都有 $T(K_i)$，则 $T(\bigoplus K_i)$ 成立；
\item 若 $K\to L\to M\to K[1]$ 是 $D(A, \text{d})$ 中的有区别三角，
且 $T$ 对其中两个对象成立，则它对第三个对象也成立；以及
\item 对所有 $k\in\mathbf{Z}$，$T(A[k])$ 成立。
\end{enumerate}
那么 $T$ 对 $D(A, \text{d})$ 的所有对象成立。这由引理
\ref{dga-lemma-property-P-sequence} 和 \ref{dga-lemma-resolve} 即得。
\end{remark}







\section{典范 delta-函子}
\label{dga-section-canonical-delta-functor}

\noindent
设 $(A, \text{d})$ 是微分分次代数。考虑函子
$\text{Mod}_{(A, \text{d})} \to K(\text{Mod}_{(A, \text{d})})$。
一般而言，这个函子{\bf 不是} $\delta$-函子。然而，函子
$\text{Mod}_{(A, \text{d})} \to D(A, \text{d})$ 的确是一个
$\delta$-函子。为说明这一点，必须定义与阿贝尔范畴
$\text{Mod}_{(A, \text{d})}$ 中的短正合列
$$
0 \to K \xrightarrow{a} L \xrightarrow{b} M \to 0
$$
相联系的态射 $\delta$。考虑态射 $a$ 的锥 $C(a)$。有
$C(a)=L\oplus K$，并把 $q:C(a)\to M$ 定义为先投影到 $L$，再复合 $b$。
于是得到微分分次 $A$-模同态
$$
q : C(a) \longrightarrow M.
$$
显然 $q\circ i=b$，其中 $i$ 如定义 \ref{dga-definition-cone} 所述。注意，由于
$a$ 是单射，$q$ 的核可与 $\text{id}_K$ 的锥等同，而后者是无环的。
故 $q$ 是拟同构。由引理 \ref{dga-lemma-the-same-up-to-isomorphisms}，三角
$$
(K, L, C(a), a, i, -p)
$$
是 $K(\text{Mod}_{(A, \text{d})})$ 中的有区别三角。由于局部化函子
$K(\text{Mod}_{(A, \text{d})}) \to D(A, \text{d})$ 是正合的，
$(K,L,C(a),a,i,-p)$ 也是 $D(A, \text{d})$ 中的有区别三角。由于 $q$
是拟同构，它在 $D(A, \text{d})$ 中是同构。因此推出
$$
(K, L, M, a, b, -p \circ q^{-1})
$$
是 $D(A, \text{d})$ 中的有区别三角。这引出下述引理。

\begin{lemma}
\label{dga-lemma-derived-canonical-delta-functor}
设 $(A, \text{d})$ 是微分分次代数。上面定义的函子
$\text{Mod}_{(A, \text{d})} \to D(A, \text{d})$ 自然具有
$\delta$-函子结构，其中
$$
\delta_{K \to L \to M} = - p \circ q^{-1}
$$
，而 $p$ 和 $q$ 如上文所述。
\end{lemma}

\begin{proof}
已经看到，对复形的每个短正合列，这个选择都会给出有区别三角。
还须证明该构造的函子性；见《导出范畴》定义
\ref{derived-definition-delta-functor}。稍作推导即可由引理
\ref{dga-lemma-functorial-cone} 得出。可与《导出范畴》引理
\ref{derived-lemma-derived-canonical-delta-functor} 比较。
\end{proof}

\begin{lemma}
\label{dga-lemma-homotopy-colimit}
设 $(A, \text{d})$ 是微分分次代数，$M_n$ 是一个微分分次模系统。那么，
$D(A, \text{d})$ 中的导出余极限 $\text{hocolim} M_n$ 由微分分次模
$\colim M_n$ 表示。
\end{lemma}

\begin{proof}
置 $M=\colim M_n$。由《导出范畴》引理
\ref{derived-lemma-compute-colimit}（应用于其底层阿贝尔群复形），有微分分次
模的正合列
$$
0 \to \bigoplus M_n \to \bigoplus M_n \to M \to 0
$$
。由引理 \ref{dga-lemma-derived-products}，这些直和也是 $D(\mathcal{A})$ 中的
直和。因此，结论由《导出范畴》定义
\ref{derived-definition-derived-colimit} 中导出余极限的定义，以及复形的短正合列
给出有区别三角这一事实（引理
\ref{dga-lemma-derived-canonical-delta-functor}）得出。
\end{proof}








\section{线性范畴}
\label{dga-section-linear}

\noindent
这里只给出定义。

\begin{definition}
\label{dga-definition-linear-category}
设 $R$ 是环。{\it $R$-线性范畴 $\mathcal{A}$} 是这样的范畴：每个态射集
都带有 $R$-模结构，并且对 $x,y,z\in\Ob(\mathcal{A})$，复合律
$$
\Hom_\mathcal{A}(y, z) \times \Hom_\mathcal{A}(x, y)
\longrightarrow
\Hom_\mathcal{A}(x, z)
$$
是 $R$-双线性的。
\end{definition}

\noindent
因此，复合确定 $R$-模的 $R$-线性映射
$$
\Hom_\mathcal{A}(y, z) \otimes_R \Hom_\mathcal{A}(x, y)
\longrightarrow
\Hom_\mathcal{A}(x, z)
$$
。注意，这里并不假设 $R$-线性范畴是加性范畴。

\begin{definition}
\label{dga-definition-functor-linear-categories}
设 $R$ 是环。{\it $R$-线性范畴之间的函子}，也称为
{\it $R$-线性函子}，是函子 $F:\mathcal{A}\to\mathcal{B}$，使得对
$\mathcal{A}$ 的所有对象 $x,y$，映射
$F : \Hom_\mathcal{A}(x, y) \to \Hom_\mathcal{B}(F(x), F(y))$
是 $R$-模同态。
\end{definition}







\section{分次范畴}
\label{dga-section-graded}

\noindent
这里只给出一些定义。

\begin{definition}
\label{dga-definition-graded-category}
设 $R$ 是环。{\it $R$ 上的分次范畴 $\mathcal{A}$} 是这样的范畴：
每个态射集都带有分次 $R$-模结构，并且对
$x,y,z\in\Ob(\mathcal{A})$，复合是 $R$-双线性的，并诱导分次 $R$-模同态
$$
\Hom_\mathcal{A}(y, z) \otimes_R \Hom_\mathcal{A}(x, y)
\longrightarrow
\Hom_\mathcal{A}(x, z)
$$
（即保持次数的同态）。
\end{definition}

\noindent
在这种情形下，以 $\Hom_\mathcal{A}^i(x,y)$ 表示分次对象
$\Hom_\mathcal{A}(x,y)$ 的次数 $i$ 部分，从而
$$
\Hom_\mathcal{A}(x, y) =
\bigoplus\nolimits_{i \in \mathbf{Z}} \Hom_\mathcal{A}^i(x, y)
$$
是它按分次部分的直和分解。

\begin{definition}
\label{dga-definition-functor-graded-categories}
设 $R$ 是环。{\it $R$ 上分次范畴之间的函子}，也称为
{\it 分次函子}，是函子 $F:\mathcal{A}\to\mathcal{B}$，使得对
$\mathcal{A}$ 的所有对象 $x,y$，映射
$F : \Hom_\mathcal{A}(x, y) \to \Hom_\mathcal{A}(F(x), F(y))$
是分次 $R$-模同态。
\end{definition}

\noindent
给定一个分次范畴，往往关心其对应的、态射次数为 $0$ 的“通常”范畴。
正式定义如下。

\begin{definition}
\label{dga-definition-H0-of-graded-category}
设 $R$ 是环，$\mathcal{A}$ 是 $R$ 上的分次范畴。令
{\it $\mathcal{A}^0$} 为与 $\mathcal{A}$ 具有相同对象的范畴，并置
$$
\Hom_{\mathcal{A}^0}(x, y) = \Hom^0_\mathcal{A}(x, y)
$$
，即取 $\mathcal{A}$ 的态射分次模的次数 $0$ 部分。
\end{definition}

\begin{definition}
\label{dga-definition-graded-direct-sum}
设 $R$ 是环，$\mathcal{A}$ 是 $R$ 上的分次范畴。若 $\mathcal{A}$ 中的直和
$(x,y,z,i,j,p,q)$（记号同《同调代数》评注
\ref{homology-remark-direct-sum}）中的 $i,j,p,q$ 都是次数 $0$ 的齐次态射，
则称它为{\it 分次直和}。
\end{definition}

\begin{example}[分次对象的分次范畴]
\label{dga-example-graded-category-graded-objects}
设 $\mathcal{B}$ 是加性范畴。回忆，在《同调代数》定义
\ref{homology-definition-graded} 中已经定义了 $\mathcal{B}$ 的分次对象范畴
$\text{Gr}(\mathcal{B})$。本例将构造 $R=\mathbf{Z}$ 上的分次范畴
$\text{Gr}^{gr}(\mathcal{B})$，使其对应范畴
$\text{Gr}^{gr}(\mathcal{B})^0$ 恢复 $\text{Gr}(\mathcal{B})$。
取 $\mathcal{B}$ 的分次对象作为 $\text{Gr}^{gr}(\mathcal{B})$ 的对象。
给定 $\mathcal{B}$ 的分次对象 $A=(A^i)$ 和 $B=(B^i)$，置
$$
\Hom_{\text{Gr}^{gr}(\mathcal{B})}(A, B) =
\bigoplus\nolimits_{n \in \mathbf{Z}} \Hom^n(A, B)
$$
，其中次数 $n$ 的分次部分是从 $A$ 到 $B$ 的次数 $n$ 齐次映射所组成的
阿贝尔群。具体地，
$$
\Hom^n(A, B) = \prod\nolimits_{p + q = n} \Hom_\mathcal{B}(A^{-q}, B^p)
$$
（注意指标的反转，并注意这里是直积而不是直和）。换言之，从 $A$ 到 $B$
的次数 $n$ 态射 $f$ 可视为系统 $f=(f_{p,q})$，其中
$p,q\in\mathbf{Z}$、$p+q=n$，且 $f_{p,q}:A^{-q}\to B^p$ 是
$\mathcal{B}$ 的态射。给定 $\mathcal{B}$ 的分次对象 $A,B,C$，
$\text{Gr}^{gr}(\mathcal{B})$ 中态射的复合通过映射
$$
\Hom^m(B, C) \times \Hom^n(A, B) \longrightarrow \Hom^{n + m}(A, C)
$$
定义，即取分次对象齐次映射的普通复合 $(g,f)\mapsto g\circ f$。按分量写为
$$
(g \circ f)_{p, r} = g_{p, q} \circ f_{-q, r}
$$
，其中 $q$ 满足 $p+q=m$ 和 $-q+r=n$。
\end{example}

\begin{example}[分次模的分次范畴]
\label{dga-example-gm-gr-cat}
设 $A$ 是环 $R$ 上的 $\mathbf{Z}$-分次代数。构造 $R$ 上的分次范畴
$\text{Mod}^{gr}_A$，使其对应范畴 $(\text{Mod}^{gr}_A)^0$ 是分次
$A$-模范畴。取右分次 $A$-模作为 $\text{Mod}^{gr}_A$ 的对象（见第
\ref{dga-section-projectives-over-algebras} 节）。给定分次 $A$-模 $L$ 和 $M$，置
$$
\Hom_{\text{Mod}^{gr}_A}(L, M) =
\bigoplus\nolimits_{n \in \mathbf{Z}} \Hom^n(L, M)
$$
，其中 $\Hom^n(L,M)$ 是次数 $n$ 齐次的右 $A$-模映射 $L\to M$ 的集合，
即对所有 $i\in\mathbf{Z}$ 都有 $f(L^i)\subset M^{i+n}$。按分量写，
$$
\Hom^n(L, M) \subset \prod\nolimits_{p + q = n} \Hom_R(L^{-q}, M^p)
$$
（注意指标的反转）是由满足下式的 $f=(f_{p,q})$ 组成的子集：
$$
f_{p, q}(m a) = f_{p - i, q + i}(m)a
$$
，其中 $a\in A^i$、$m\in L^{-q-i}$。对分次 $A$-模 $K,L,M$，通过映射
$$
\Hom^m(L, M) \times \Hom^n(K, L) \longrightarrow \Hom^{n + m}(K, M)
$$
定义 $\text{Mod}^{gr}_A$ 中的复合，即取右 $A$-模映射的普通复合
$(g,f)\mapsto g\circ f$。
\end{example}

\begin{remark}
\label{dga-remark-graded-shift-functors}
设 $R$ 是环，$\mathcal{D}$ 是 $R$-线性范畴，并带有由
$n\in\mathbf{Z}$ 作指标的一族 $R$-线性函子
$[n]:\mathcal{D}\to\mathcal{D}$、$x\mapsto x[n]$，使得
$[n]\circ[m]=[n+m]$ 且 $[0]=\text{id}_\mathcal{D}$（作为函子相等）。
由此可以构造 $R$ 上与 $\mathcal{D}$ 具有相同对象的分次范畴
$\mathcal{D}^{gr}$，置
$$
\Hom_{\mathcal{D}^{gr}}(x, y) =
\bigoplus\nolimits_{n \in \mathbf{Z}} \Hom_\mathcal{D}(x, y[n])
$$
，其中 $x,y$ 是 $\mathcal{D}$ 的对象。注意
$(\mathcal{D}^{gr})^0=\mathcal{D}$（见定义
\ref{dga-definition-H0-of-graded-category}）。此外，分次范畴
$\mathcal{D}^{gr}$ 继承 $R$-线性分次函子 $[n]$，满足
$[n]\circ[m]=[n+m]$ 和 $[0]=\text{id}_{\mathcal{D}^{gr}}$，并具有性质
$$
\Hom_{\mathcal{D}^{gr}}(x, y[n]) = \Hom_{\mathcal{D}^{gr}}(x, y)[n]
$$
；这是与态射复合相容的分次 $R$-模同构。

\medskip\noindent
反过来，设给定 $R$ 上的分次范畴 $\mathcal{A}$，它带有一族 $R$-线性
分次函子 $[n]$，满足 $[n]\circ[m]=[n+m]$ 和
$[0]=\text{id}_\mathcal{A}$，并且还带有同构
$$
\Hom_\mathcal{A}(x, y[n]) = \Hom_\mathcal{A}(x, y)[n]
$$
；这是与态射复合相容的分次 $R$-模同构。读者容易证明
$\mathcal{A}=(\mathcal{A}^0)^{gr}$。

\medskip\noindent
下面是上述关系 $\mathcal{D}\leftrightarrow\mathcal{A}$ 的两个例子：
\begin{enumerate}
\item 设 $\mathcal{B}$ 是加性范畴。若
$\mathcal{D}=\text{Gr}(\mathcal{B})$，则
$\mathcal{A}=\text{Gr}^{gr}(\mathcal{B})$，如例
\ref{dga-example-graded-category-graded-objects} 所述。
\item 若 $A$ 是分次环，且 $\mathcal{D}=\text{Mod}_A$ 是分次右 $A$-模
范畴，则 $\mathcal{A}=\text{Mod}^{gr}_A$；见例
\ref{dga-example-gm-gr-cat}。
\end{enumerate}
\end{remark}






\section{微分分次范畴}
\label{dga-section-dga-categories}

\noindent
注意，若 $R$ 是环，则 $R$ 是自身之上的微分分次代数（当然取
$R=R^0$）。此时，微分分次 $R$-模与 $R$-模复形是一回事。特别地，
给定两个微分分次 $R$-模 $M$ 和 $N$，以 $M\otimes_R N$ 表示如下微分
分次 $R$-模：先取与 $M,N$ 对应的 $R$-模复形的张量积所得到的双复形，
再取其关联全复形。

\begin{definition}
\label{dga-definition-dga-category}
设 $R$ 是环。{\it $R$ 上的微分分次范畴 $\mathcal{A}$} 是这样的范畴：
每个态射集都带有微分分次 $R$-模结构，并且对
$x,y,z\in\Ob(\mathcal{A})$，复合是 $R$-双线性的，并诱导微分分次
$R$-模同态
$$
\Hom_\mathcal{A}(y, z) \otimes_R \Hom_\mathcal{A}(x, y)
\longrightarrow
\Hom_\mathcal{A}(x, z)
$$
。
\end{definition}

\noindent
定义中的最后一个条件表示：若
$f\in\Hom_\mathcal{A}^n(x,y)$ 和 $g\in\Hom_\mathcal{A}^m(y,z)$
分别是次数 $n$ 和 $m$ 的齐次态射，则在
$\Hom_\mathcal{A}^{n+m+1}(x,z)$ 中有
$$
\text{d}(g \circ f) = \text{d}(g) \circ f + (-1)^mg \circ \text{d}(f)
$$
。这由双复形全复形上的微分符号法则得出；见《同调代数》定义
\ref{homology-definition-associated-simple-complex}。

\begin{definition}
\label{dga-definition-functor-dga-categories}
设 $R$ 是环。{\it $R$ 上微分分次范畴之间的函子}是函子
$F:\mathcal{A}\to\mathcal{B}$，使得对 $\mathcal{A}$ 的所有对象 $x,y$，
映射
$F : \Hom_\mathcal{A}(x, y) \to \Hom_\mathcal{A}(F(x), F(y))$
是微分分次 $R$-模同态。
\end{definition}

\noindent
给定一个微分分次范畴，往往关心其对应的复形范畴和同伦范畴。正式定义如下。

\begin{definition}
\label{dga-definition-homotopy-category-of-dga-category}
设 $R$ 是环，$\mathcal{A}$ 是 $R$ 上的微分分次范畴。定义
\begin{enumerate}
\item {\it $\mathcal{A}$ 的复形范畴}\footnote{这可能不是标准术语。}
为范畴 $\text{Comp}(\mathcal{A})$；它与 $\mathcal{A}$ 具有相同对象，并且
$$
\Hom_{\text{Comp}(\mathcal{A})}(x, y) =
\Ker(d : \Hom^0_\mathcal{A}(x, y) \to \Hom^1_\mathcal{A}(x, y))
$$
\item {\it $\mathcal{A}$ 的同伦范畴}为范畴 $K(\mathcal{A})$；它与
$\mathcal{A}$ 具有相同对象，并且
$$
\Hom_{K(\mathcal{A})}(x, y) = H^0(\Hom_\mathcal{A}(x, y))
$$
\end{enumerate}
\end{definition}

\noindent
这里对符号 $K(\mathcal{A})$ 的用法并不标准，但至少与 Stacks 项目其他章中
$K(-)$ 的用法相容。

\begin{definition}
\label{dga-definition-dg-direct-sum}
设 $R$ 是环，$\mathcal{A}$ 是 $R$ 上的微分分次范畴。若 $\mathcal{A}$
中的直和 $(x,y,z,i,j,p,q)$（记号同《同调代数》评注
\ref{homology-remark-direct-sum}）中的 $i,j,p,q$ 都是次数 $0$ 的齐次闭态射，
即 $\text{d}(i)=0$ 等，则称它为{\it 微分分次直和}。
\end{definition}

\begin{lemma}
\label{dga-lemma-functorial}
设 $R$ 是环。$R$ 上微分分次范畴之间的函子
$F:\mathcal{A}\to\mathcal{B}$ 诱导函子
$\text{Comp}(\mathcal{A}) \to \text{Comp}(\mathcal{B})$
和 $K(\mathcal{A})\to K(\mathcal{B})$。
\end{lemma}

\begin{proof}
略。
\end{proof}

\begin{example}[复形的微分分次范畴]
\label{dga-example-category-complexes}
设 $\mathcal{B}$ 是加性范畴。构造 $R=\mathbf{Z}$ 上的微分分次范畴
$\text{Comp}^{dg}(\mathcal{B})$，使其对应的复形范畴是
$\text{Comp}(\mathcal{B})$，对应的同伦范畴是 $K(\mathcal{B})$。
取 $\mathcal{B}$ 的复形作为 $\text{Comp}^{dg}(\mathcal{B})$ 的对象。
给定 $\mathcal{B}$ 的复形 $A^\bullet$ 和 $B^\bullet$，有时也以
$A^\bullet$ 和 $B^\bullet$ 表示 $\mathcal{B}$ 中相应的分次对象
（即忘掉微分）。采用这一记号上的滥用，置
$$
\Hom_{\text{Comp}^{dg}(\mathcal{B})}(A^\bullet, B^\bullet) =
\Hom_{\text{Gr}^{gr}(\mathcal{B})}(A^\bullet, B^\bullet) =
\bigoplus\nolimits_{n \in \mathbf{Z}} \Hom^n(A, B)
$$
，把它视为分次 $\mathbf{Z}$-模，记号和定义同例
\ref{dga-example-graded-category-graded-objects}。换言之，第 $n$ 个分次部分是
分次对象之间次数 $n$ 的齐次态射所组成的阿贝尔群
$$
\Hom^n(A^\bullet, B^\bullet) =
\prod\nolimits_{p + q = n} \Hom_\mathcal{B}(A^{-q}, B^p)
$$
。注意指标的反转，并注意这里是直积而不是直和。对次数为 $n$ 的元素
$f\in\Hom^n(A^\bullet,B^\bullet)$，置
$$
\text{d}(f) = \text{d}_B \circ f - (-1)^n f \circ \text{d}_A
$$
。符号与《代数进阶》第 \ref{more-algebra-section-sign-rules} 节完全相同。
为使此式有意义，把 $\text{d}_B$ 和 $\text{d}_A$ 看作 $\mathcal{B}$ 的
分次对象之间次数 $1$ 的齐次映射，并在右侧使用范畴
$\text{Gr}^{gr}(\mathcal{B})$ 中的复合。按分量写，若
$f=(f_{p,q})$ 且 $f_{p,q}:A^{-q}\to B^p$，则
\begin{equation}
\label{dga-equation-differential-hom-complex}
\text{d}(f_{p, q}) =
\text{d}_B \circ f_{p, q} - (-1)^{p + q} f_{p, q} \circ \text{d}_A 
\end{equation}
注意，此表达式的第一项属于 $\Hom_\mathcal{B}(A^{-q},B^{p+1})$，
第二项属于 $\Hom_\mathcal{B}(A^{-q-1},B^p)$。读者可以验证：
\begin{enumerate}
\item $\text{d}$ 的平方为零；
\item $\Hom^n(A^\bullet,B^\bullet)$ 中的元素 $f$ 满足
$\text{d}(f)=0$，当且仅当 $\mathcal{B}$ 的分次对象之间的态射
$f:A^\bullet\to B^\bullet[n]$ 实际上是复形映射；
\item 特别地，$\text{Comp}^{dg}(\mathcal{B})$ 的复形范畴等于
$\text{Comp}(\mathcal{B})$；
\item (2) 中由 $f$ 定义的复形态射与零同伦，当且仅当存在
$g\in\Hom^{n-1}(A^\bullet,B^\bullet)$ 使得 $f=\text{d}(g)$；
\item 特别地，得到典范同构
$$
\Hom_{K(\mathcal{B})}(A^\bullet, B^\bullet)
\longrightarrow
H^0(\Hom_{\text{Comp}^{dg}(\mathcal{B})}(A^\bullet, B^\bullet))
$$
，并且 $\text{Comp}^{dg}(\mathcal{B})$ 的同伦范畴等于 $K(\mathcal{B})$。
\end{enumerate}
给定复形 $A^\bullet,B^\bullet,C^\bullet$，通过分次范畴
$\text{Gr}^{gr}(\mathcal{B})$ 中的复合 $(g,f)\mapsto g\circ f$ 定义复合
$$
\Hom^m(B^\bullet, C^\bullet) \times \Hom^n(A^\bullet, B^\bullet)
\longrightarrow
\Hom^{n + m}(A^\bullet, C^\bullet)
$$
；见例 \ref{dga-example-graded-category-graded-objects}。这定义微分分次模的映射
$$
\Hom_{\text{Comp}^{dg}(\mathcal{B})}(B^\bullet, C^\bullet)
\otimes_R
\Hom_{\text{Comp}^{dg}(\mathcal{B})}(A^\bullet, B^\bullet)
\longrightarrow
\Hom_{\text{Comp}^{dg}(\mathcal{B})}(A^\bullet, C^\bullet)
$$
，正如定义 \ref{dga-definition-dga-category} 所要求的，因为
\begin{align*}
\text{d}(g \circ f) & =
\text{d}_C \circ g \circ f - (-1)^{n + m} g \circ f \circ \text{d}_A \\
& =
\left(\text{d}_C \circ g - (-1)^m g \circ \text{d}_B\right) \circ f +
(-1)^m g \circ \left(\text{d}_B \circ f - (-1)^n f \circ \text{d}_A\right) \\
& =
\text{d}(g) \circ f + (-1)^m g \circ \text{d}(f)
\end{align*}
，这正是所需等式。
\end{example}

\begin{lemma}
\label{dga-lemma-additive-functor-induces-dga-functor}
设 $F:\mathcal{B}\to\mathcal{B}'$ 是加性范畴之间的加性函子。那么 $F$
诱导例 \ref{dga-example-category-complexes} 中微分分次范畴之间的函子
$$
F : \text{Comp}^{dg}(\mathcal{B}) \to \text{Comp}^{dg}(\mathcal{B}')
$$
，并在复形范畴和同伦范畴上诱导通常的函子。
\end{lemma}

\begin{proof}
略。
\end{proof}

\begin{example}[微分分次模的微分分次范畴]
\label{dga-example-dgm-dg-cat}
设 $(A, \text{d})$ 是环 $R$ 上的微分分次代数。构造 $R$ 上的微分分次
范畴 $\text{Mod}^{dg}_{(A, \text{d})}$，使其复形范畴是
$\text{Mod}_{(A, \text{d})}$，同伦范畴是
$K(\text{Mod}_{(A, \text{d})})$。取微分分次 $A$-模作为
$\text{Mod}^{dg}_{(A, \text{d})}$ 的对象。给定微分分次 $A$-模 $L,M$，置
$$
\Hom_{\text{Mod}^{dg}_{(A, \text{d})}}(L, M) =
\Hom_{\text{Mod}^{gr}_A}(L, M) = \bigoplus \Hom^n(L, M)
$$
，把它视为分次 $R$-模，其中右侧如例 \ref{dga-example-gm-gr-cat} 所定义。
换言之，第 $n$ 个分次部分 $\Hom^n(L,M)$ 是次数 $n$ 齐次的右 $A$-模映射
所组成的 $R$-模。对元素 $f\in\Hom^n(L,M)$，置
$$
\text{d}(f) = \text{d}_M \circ f - (-1)^n f \circ \text{d}_L
$$
。为使此式有意义，把 $\text{d}_M$ 和 $\text{d}_L$ 看作分次 $R$-模映射，
并采用分次 $R$-模映射的复合。显然，作为分次 $R$-模映射，
$\text{d}(f)$ 是次数 $n+1$ 的齐次映射；它还是 $A$-线性的，因为
\begin{align*}
\text{d}(f)(xa)
& =
\text{d}_M(f(x) a) - (-1)^n f (\text{d}_L(xa)) \\
& =
\text{d}_M(f(x)) a + (-1)^{\deg(x) + n} f(x) \text{d}(a) 
- (-1)^n f(\text{d}_L(x)) a - (-1)^{n + \deg(x)} f(x) \text{d}(a) \\
& = \text{d}(f)(x) a
\end{align*}
，这正是所需等式（注意，若 $\Hom$ 上微分的定义中没有这个符号，计算就
不能成立）。按分量表示 $f$ 的微分时，也有与例
\ref{dga-example-category-complexes} 类似的公式。读者可以仿照例
\ref{dga-example-category-complexes} 验证：
\begin{enumerate}
\item $\text{d}$ 的平方为零；
\item $\Hom^n(L,M)$ 中的元素 $f$ 满足 $\text{d}(f)=0$，当且仅当
$f:L\to M[n]$ 是微分分次 $A$-模同态；
\item 特别地，$\text{Mod}^{dg}_{(A, \text{d})}$ 的复形范畴是
$\text{Mod}_{(A, \text{d})}$；
\item (2) 中由 $f$ 定义的同态与零同伦，当且仅当存在
$g\in\Hom^{n-1}(L,M)$ 使得 $f=\text{d}(g)$；
\item 特别地，得到典范同构
$$
\Hom_{K(\text{Mod}_{(A, \text{d})})}(L, M)
\longrightarrow
H^0(\Hom_{\text{Mod}^{dg}_{(A, \text{d})}}(L, M))
$$
，并且 $\text{Mod}^{dg}_{(A, \text{d})}$ 的同伦范畴是
$K(\text{Mod}_{(A, \text{d})})$。
\end{enumerate}
给定微分分次 $A$-模 $K,L,M$，通过齐次右 $A$-模映射的复合
$(g,f)\mapsto g\circ f$ 定义复合
$$
\Hom^m(L, M) \times \Hom^n(K, L) \longrightarrow \Hom^{n + m}(K, M)
$$
。这定义微分分次模的映射
$$
\Hom_{\text{Mod}^{dg}_{(A, \text{d})}}(L, M) \otimes_R
\Hom_{\text{Mod}^{dg}_{(A, \text{d})}}(K, L) \longrightarrow
\Hom_{\text{Mod}^{dg}_{(A, \text{d})}}(K, M)
$$
，正如定义 \ref{dga-definition-dga-category} 所要求的，因为
\begin{align*}
\text{d}(g \circ f) & =
\text{d}_M \circ g \circ f - (-1)^{n + m} g \circ f \circ \text{d}_K \\
& =
\left(\text{d}_M \circ g - (-1)^m g \circ \text{d}_L\right) \circ f +
(-1)^m g \circ \left(\text{d}_L \circ f - (-1)^n f \circ \text{d}_K\right) \\
& =
\text{d}(g) \circ f + (-1)^m g \circ \text{d}(f)
\end{align*}
，这正是所需等式。
\end{example}

\begin{lemma}
\label{dga-lemma-homomorphism-induces-dga-functor}
设 $\varphi:(A,\text{d})\to(E,\text{d})$ 是微分分次代数同态。那么
$\varphi$ 诱导例 \ref{dga-example-dgm-dg-cat} 中微分分次范畴之间的函子
$$
F :
\text{Mod}^{dg}_{(E, \text{d})}
\longrightarrow
\text{Mod}^{dg}_{(A, \text{d})}
$$
，并在微分分次模范畴和同伦范畴上诱导显然的限制函子。
\end{lemma}

\begin{proof}
略。
\end{proof}

\begin{lemma}
\label{dga-lemma-construction}
设 $R$ 是环，$\mathcal{A}$ 是 $R$ 上的微分分次范畴，$x$ 是
$\mathcal{A}$ 的对象。令
$$
(E, \text{d}) = \Hom_\mathcal{A}(x, x)
$$
为 $x$ 的自同态微分分次 $R$-代数。令 $E$ 通过 $\mathcal{A}$ 中的复合作用
在 $\Hom_\mathcal{A}(x,y)$ 上，得到微分分次范畴之间的函子
$$
\mathcal{A} \longrightarrow \text{Mod}^{dg}_{(E, \text{d})},\quad
y \longmapsto \Hom_\mathcal{A}(x, y)
$$
。应用引理 \ref{dga-lemma-functorial}，此函子诱导函子
$$
\text{Comp}(\mathcal{A}) \to \text{Mod}_{(A, \text{d})}
\quad\text{且}\quad
K(\mathcal{A}) \to K(\text{Mod}_{(A, \text{d})})
$$
。
\end{lemma}

\begin{proof}
本引理不证自明。
\end{proof}









\section{得到三角范畴}
\label{dga-section-review}

\noindent
本节讨论《导出范畴》命题
\ref{derived-proposition-homotopy-category-triangulated} 和命题
\ref{dga-proposition-homotopy-category-triangulated} 的证明论证所适用的最一般框架。

\medskip\noindent
设 $R$ 是环，$\mathcal{A}$ 是 $R$ 上的微分分次范畴。为使论证成立，
对 $\mathcal{A}$ 施加下列公理：
\begin{enumerate}
\item[(A)] $\mathcal{A}$ 有零对象以及两个对象的微分分次直和
（如定义 \ref{dga-definition-dg-direct-sum}）。
\item[(B)] 存在微分分次范畴之间的函子
$[n]:\mathcal{A}\longrightarrow\mathcal{A}$，使得
$[0]=\text{id}_\mathcal{A}$ 且 $[n+m]=[n]\circ[m]$，并且给定与复合相容的
微分分次 $R$-模同构
$$
\Hom_\mathcal{A}(x, y[n]) = \Hom_\mathcal{A}(x, y)[n]
$$
。
\end{enumerate}

\noindent
对上述微分分次范畴 $\mathcal{A}$，规定：
\begin{enumerate}
\item 若在底层分次范畴中存在同构 $y\cong x\oplus z$，使得
$x\to z$ 和 $y\to z$ 是（余）投影，则称 $\text{Comp}(\mathcal{A})$
中的态射列 $x\to y\to z$ 为{\it 容许短正合列}。
\item 若 $\text{Comp}(\mathcal{A})$ 中的态射 $x\to y$ 可延拓为容许短正合列
$x\to y\to z$，则称它为{\it 容许单态射}。
\item 若 $\text{Comp}(\mathcal{A})$ 中的态射 $y\to z$ 可延拓为容许短正合列
$x\to y\to z$，则称它为{\it 容许满态射}。
\end{enumerate}
下一个引理说明，在公理 (A) 和 (B) 成立时，容许短正合列给出一个三角。

\begin{lemma}
\label{dga-lemma-get-triangle}
设 $\mathcal{A}$ 是满足公理 (A)、(B) 的微分分次范畴。给定容许短正合列
$x\to y\to z$，得到（见证明）$\text{Comp}(\mathcal{A})$ 中的三角
$$
x \to y \to z \to x[1]
$$
，其性质为：$z[-1]\to x\to y\to z\to x[1]$ 中任意两个相邻箭头的复合
在 $K(\mathcal{A})$ 中为零。
\end{lemma}

\begin{proof}
选择图表
$$
\xymatrix{
x \ar[rr]_1 \ar[rd]_a & & x \\
& y \ar[ru]_\pi \ar[rd]^b & \\
z \ar[rr]^1 \ar[ru]^s & & z
}
$$
，它给出容许短正合列定义中的分次对象同构 $y\cong x\oplus z$。在此情形下，
下列等式成立：
\begin{enumerate}
\item $1=\pi a$，从而 $\text{d}(\pi)a=0$；
\item $1=bs$，从而 $b\text{d}(s)=0$；
\item $1=a\pi+sb$，从而 $a\text{d}(\pi)+\text{d}(s)b=0$；
\item $\pi s=0$，从而 $\text{d}(\pi)s+\pi\text{d}(s)=0$；
\item $\text{d}(s)=a\pi\text{d}(s)$，因为
$\text{d}(s)=(a\pi+sb)\text{d}(s)$ 且 $b\text{d}(s)=0$；
\item $\text{d}(\pi)=\text{d}(\pi)sb$，因为
$\text{d}(\pi)=\text{d}(\pi)(a\pi+sb)$ 且 $\text{d}(\pi)a=0$；
\item $\text{d}(\pi\text{d}(s))=0$，因为将它后复合单态射 $a$，得到
$\text{d}(a\pi\text{d}(s))=\text{d}(\text{d}(s))=0$；以及
\item $\text{d}(\text{d}(\pi)s)=0$，因为由 (4)，它是
$\text{d}(\pi\text{d}(s))$ 的负值，而后者由 (7) 为 $0$。
\end{enumerate}
这里反复使用了 $\text{d}(a)=0$、$\text{d}(b)=0$ 和 $\text{d}(1)=0$。
由 (7)，
$$
\delta = \pi \text{d}(s) = - \text{d}(\pi) s : z \to x[1]
$$
是 $\text{Comp}(\mathcal{A})$ 中的态射。由 (5)，复合
$a\delta=a\pi\text{d}(s)=\text{d}(s)$ 与零同伦。由 (6)，复合
$\delta b=-\text{d}(\pi)sb=\text{d}(-\pi)$ 与零同伦。
\end{proof}

\noindent
除公理 (A)、(B) 外，还需要一个关于锥存在性的公理。将全部条件正式写成如下情形。

\begin{situation}
\label{dga-situation-ABC}
这里 $R$ 是环，$\mathcal{A}$ 是 $R$ 上满足公理 (A)、(B) 以及下列公理的
微分分次范畴：
\begin{enumerate}
\item[(C)] 给定次数为 $0$ 且满足 $\text{d}(f)=0$ 的箭头 $f:x\to y$，
$\text{Comp}(\mathcal{A})$ 中存在容许短正合列 $y\to c(f)\to x[1]$，
使引理 \ref{dga-lemma-get-triangle} 中的映射 $x[1]\to y[1]$ 等于 $f[1]$。
\end{enumerate}
\end{situation}

\noindent
称 $c(f)$ 为态射 $f$ 的一个{\it 锥}。若 (A)、(B)、(C) 成立，
则锥在一种较弱的意义下具有函子性。

\begin{lemma}
\label{dga-lemma-cone}
\begin{slogan}
同伦范畴是三角范畴。本引理证明三角范畴公理的一部分。
\end{slogan}
在情形 \ref{dga-situation-ABC} 中，设图表
$$
\xymatrix{
x_1 \ar[r]_{f_1} \ar[d]_a & y_1 \ar[d]^b \\
x_2 \ar[r]^{f_2} & y_2
}
$$
在 $\text{Comp}(\mathcal{A})$ 中在同伦意义下交换。则存在态射
$c:c(f_1)\to c(f_2)$，它在 $K(\mathcal{A})$ 中给出三角态射
$$
(a, b, c) : (x_1, y_1, c(f_1)) \to (x_1, y_1, c(f_1))
$$
。
\end{lemma}

\begin{proof}
假设意味着存在次数为 $-1$ 的态射 $h:x_1\to y_2$，使得
$\text{d}(h)=bf_1-f_2a$。选择与态射 $y_i\to c(f_i)\to x_i[1]$ 相容的
分次对象同构 $c(f_i)=y_i\oplus x_i[1]$。以
$a_i:y_i\to c(f_i)$、$b_i:c(f_i)\to x_i[1]$、
$s_i:x_i[1]\to c(f_i)$ 和 $\pi_i:c(f_i)\to y_i$ 表示这些给定态射。
回忆，$x_i[1]\to y_i[1]$ 由 $\pi_i\text{d}(s_i)$ 给出。由公理 (C)，这意味着
$$
f_i = \pi_i \text{d}(s_i) = - \text{d}(\pi_i) s_i
$$
（利用移位函子 $[1]$ 将 $\Hom(x_i,y_i)$ 与
$\Hom(x_i[1],y_i[1])$ 识别）。置
$c=a_2b\pi_1+s_2ab_1+a_2hb$。利用引理 \ref{dga-lemma-get-triangle}
证明中得到的等式，有
\begin{align*}
\text{d}(c)
& =
a_2 b \text{d}(\pi_1) + \text{d}(s_2) a b_1 + a_2 \text{d}(h) b_1 \\
& =
- a_2 b f_1 b_1 + a_2 f_2 a b_1 + a_2 (b f_1 - f_2 a) b_1 \\
& = 0
\end{align*}
（这里特别使用了
$\text{d}(\pi_1) = \text{d}(\pi_1) s_1 b_1 = f_1 b_1$ 以及
$\text{d}(s_2)=a_2\pi_2\text{d}(s_2)=a_2f_2$）。因此，$c$ 是
$\mathcal{A}$ 中次数为 $0$ 的态射 $c:c(f_1)\to c(f_2)$，并与给定态射
$y_i\to c(f_i)\to x_i[1]$ 相容。
\end{proof}

\noindent
在情形 \ref{dga-situation-ABC} 中，若存在容许短正合列 $x'\to y'\to z'$，
使得 $K(\mathcal{A})$ 中的三角 $(x,y,z,f,g,h)$ 作为三角同构于引理
\ref{dga-lemma-get-triangle} 构造的三角
$(x',y',z',x'\to y',y'\to z',\delta)$，则称它为{\it 有区别三角}。
下文将证明
$$
\boxed{
K(\mathcal{A})\text{ 是三角范畴}
}
$$
这个结果虽没有想象中那么一般，却适用于 Stacks 项目迄今所讨论情形的若干
自然推广。下面给出一些例子：
\begin{enumerate}
\item 设 $(X,\mathcal{O}_X)$ 是环化空间，$(A,d)$ 是微分分次
$\mathcal{O}_X$-代数层，$\mathcal{A}$ 是微分分次 $A$-模的微分分次范畴。
那么 $K(\mathcal{A})$ 是三角范畴。
\item 设 $(\mathcal{C},\mathcal{O})$ 是环化位点，$(A,d)$ 是微分分次
$\mathcal{O}$-代数层，$\mathcal{A}$ 是微分分次 $A$-模的微分分次范畴。
那么 $K(\mathcal{A})$ 是三角范畴。见《微分分次层》命题
\ref{sdga-proposition-homotopy-category-triangulated}。
\item 另外两个风格不同的例子见《例》第
\ref{examples-section-nongraded-differential-graded} 节。
\end{enumerate}

\noindent
下面这个简单引理是该构造的关键。

\begin{lemma}
\label{dga-lemma-id-cone-null}
在情形 \ref{dga-situation-ABC} 中，给定 $\mathcal{A}$ 的任意对象 $x$ 以及
恒等态射 $1_x:x\to x$ 的锥 $C(1_x)$，$C(1_x)$ 上的恒等态射与零同伦。
\end{lemma}

\begin{proof}
考虑公理 (C) 给出的容许短正合列
$$
\xymatrix{
x \ar@<0.5ex>[r]^a  &
C(1_x) \ar@<0.5ex>[l]^{\pi} \ar@<0.5ex>[r]^b &
x[1]\ar@<0.5ex>[l]^s
}
$$
。由引理 \ref{dga-lemma-get-triangle}，在移位下识别 Hom 集，有
$1_x=\pi d(s)=-d(\pi)s$，其中把 $s$ 看作
$\Hom_{\mathcal{A}}^{-1}(x,C(1_x))$ 中的态射。因此，利用引理
\ref{dga-lemma-get-triangle} 的公式 (5)，有 $a=a\pi d(s)=d(s)$；利用引理
\ref{dga-lemma-get-triangle} 的公式 (6)，有 $b=-d(\pi)sb=-d(\pi)$。故
$$
1_{C(1_x)} = a\pi + sb = d(s)\pi - sd(\pi) = d(s\pi)
$$
，因为 $s$ 的次数为 $-1$。
\end{proof}

\noindent
上述引理的更一般版本将见于引理 \ref{dga-lemma-cone-homotopy}。下述引理是
引理 \ref{dga-lemma-make-commute-map} 的对应版本。

\begin{lemma}
\label{dga-lemma-homo-change}
在情形 \ref{dga-situation-ABC} 中，给定 $\text{Comp}(\mathcal{A})$ 中在同伦
意义下交换的图表
$$
\xymatrix{x\ar[r]^f\ar[d]_a & y\ar[d]^b\\
z\ar[r]^g & w}
$$
。那么
\begin{enumerate}
\item 若 $f$ 是容许单态射，则 $b$ 同伦于某个使图表交换的态射 $b'$。
\item 若 $g$ 是容许满态射，则 $a$ 同伦于某个使图表交换的态射 $a'$。
\end{enumerate}
\end{lemma}

\begin{proof}
为证明 (1)，注意假设蕴含：存在 $\Hom_{\mathcal{A}}(x,w)$ 中次数为 $-1$
的某个 $h$，使得 $bf-ga=d(h)$。由于 $f$ 是容许单态射，在范畴
$\mathcal{A}$ 中存在次数为 $0$ 的态射 $\pi:y\to x$。令
$b'=b-d(h\pi)$。于是
\begin{align*}
b'f = bf - d(h\pi)f
= &
bf - d(h\pi f) \quad (\text{因为 }d(f) = 0) \\
= &
bf-d(h) \\
= &
ga
\end{align*}
，这正是所需等式。(2) 的证明略。
\end{proof}

\noindent
下述引理是引理 \ref{dga-lemma-make-injective} 的对应版本。

\begin{lemma}
\label{dga-lemma-factor}
在情形 \ref{dga-situation-ABC} 中，设 $\alpha:x\to y$ 是
$\text{Comp}(\mathcal{A})$ 中的态射。那么在 $\text{Comp}(\mathcal{A})$
中存在分解
$$
\xymatrix{
x \ar[r]^{\tilde{\alpha}}  &
\tilde{y} \ar@<0.5ex>[r]^{\pi} &
y\ar@<0.5ex>[l]^s
}
$$
，使得
\begin{enumerate}
\item $\tilde{\alpha}$ 是容许单态射，且
$\pi\tilde{\alpha}=\alpha$；
\item $\text{Comp}(\mathcal{A})$ 中存在态射 $s:y\to\tilde{y}$，
使得 $\pi s=1_y$，且 $s\pi$ 与 $1_{\tilde{y}}$ 同伦。
\end{enumerate}
\end{lemma}

\begin{proof}
由公理 (A)，可令 $\tilde y$ 为 $y$ 与 $C(1_x)$ 的微分分次直和；即存在图表
$$
\xymatrix@C=3pc{
y \ar@<0.5ex>[r]^s  &
y\oplus C(1_x) \ar@<0.5ex>[l]^{\pi} \ar@<0.5ex>[r]^{p} &
C(1_x)\ar@<0.5ex>[l]^t
}
$$
，其中所有态射次数均为零，且都在 $\text{Comp}(\mathcal{A})$ 中。
令 $\tilde y=y\oplus C(1_x)$。于是 $1_{\tilde y}=s\pi+tp$。现考虑图表
$$
\xymatrix{
x \ar[r]^{\tilde{\alpha}}  &
\tilde{y} \ar@<0.5ex>[r]^{\pi} &
y\ar@<0.5ex>[l]^s
}
$$
，其中 $\tilde\alpha$ 由态射 $x\xrightarrow{\alpha}y$ 以及下述容许短正合列
中的自然态射 $x\to C(1_x)$ 诱导：
$$
\xymatrix{
x \ar@<0.5ex>[r]  &
C(1_x) \ar@<0.5ex>[l] \ar@<0.5ex>[r] &
x[1]\ar@<0.5ex>[l]
}
$$
因此，该图表中次数为 $0$ 的态射 $C(1_x)\to x$ 与零态射 $y\to x$
共同诱导次数为 $0$ 的态射 $\beta:\tilde y\to x$。$\tilde\alpha$ 是容许
单态射，因为它位于容许短正合列
$$
\xymatrix{
x\ar[r]^{\tilde{\alpha}} &
\tilde{y} \ar[r] &
x[1]
}
$$
中。此外，由 $\tilde\alpha$ 的构造，$\pi\tilde\alpha=\alpha$；由第一个
图表，$\pi s=1_y$。只须证明 $s\pi$ 与 $1_{\tilde y}$ 同伦。把 $1_x$
写成某个次数为 $-1$ 的映射的微分 $d(h)$。于是，最后一个断言由下式得出：
\begin{align*}
1_{\tilde{y}} - s\pi
= &
tp \\
= &
t(dh)p\quad\text{（由引理 \ref{dga-lemma-id-cone-null}）} \\
= &
d(thp)
\end{align*}
，因为 $dt=dp=0$，且 $t$ 的次数为零。
\end{proof}

\noindent
下述引理是引理 \ref{dga-lemma-sequence-maps-split} 的对应版本。

\begin{lemma}
\label{dga-lemma-analogue-sequence-maps-split}
在情形 \ref{dga-situation-ABC} 中，设
$x_1\to x_2\to\ldots\to x_n$ 是 $\text{Comp}(\mathcal{A})$ 中一列可复合
态射。那么 $\text{Comp}(\mathcal{A})$ 中存在交换图表
$$
\xymatrix{x_1\ar[r] & x_2\ar[r] & \ldots\ar[r] & x_n\\
y_1\ar[r]\ar[u] & y_2\ar[r]\ar[u] & \ldots\ar[r] & y_n\ar[u]}
$$
，使每个 $y_i\to y_{i+1}$ 都是容许单态射，而每个 $y_i\to x_i$ 都是同伦等价。
\end{lemma}

\begin{proof}
$n=1$ 的情形是平凡的：只须取 $y_1=x_1$，而 $x_1$ 上的恒等态射当然是
同伦等价。$n=2$ 的情形由引理 \ref{dga-lemma-factor} 给出。假设已经把图表构造
到 $x_{n-1}$。对复合 $y_{n-1}\to x_{n-1}\to x_n$ 应用引理
\ref{dga-lemma-factor}，得到 $y_n$。于是 $y_{n-1}\to y_n$ 是容许单态射，
$y_n\to x_n$ 是同伦等价。
\end{proof}

\noindent
下述引理是引理 \ref{dga-lemma-nilpotent} 的对应版本。

\begin{lemma}
\label{dga-lemma-triseq}
在情形 \ref{dga-situation-ABC} 中，设 $x_i\to y_i\to z_i$ 是
$\mathcal{A}$ 中的态射（$i=1,2,3$），并且 $x_2\to y_2\to z_2$ 是容许
短正合列。设 $b:y_1\to y_2$ 和 $b':y_2\to y_3$ 是
$\text{Comp}(\mathcal{A})$ 中的态射，使得
$$
\vcenter{
\xymatrix{
x_1 \ar[d]_0 \ar[r] &
y_1 \ar[r] \ar[d]_b &
z_1 \ar[d]_0 \\
x_2 \ar[r] & y_2 \ar[r] & z_2
}
}
\quad\text{且}\quad
\vcenter{
\xymatrix{
x_2 \ar[d]^0 \ar[r] &
y_2 \ar[r] \ar[d]^{b'} &
z_2 \ar[d]^0 \\
x_3 \ar[r] & y_3 \ar[r] & z_3
}
}
$$
中的两个图表均在同伦意义下交换。那么 $b'\circ b$ 与 $0$ 同伦。
\end{lemma}

\begin{proof}
由引理 \ref{dga-lemma-homo-change}，可分别用与 $b,b'$ 同伦的映射
$\tilde b,\tilde b'$ 替代它们，使左图的右方块交换，右图的左方块交换。
设对 $\mathcal{A}$ 中次数为 $-1$ 的态射 $h,h'$，有
$b=\tilde b+d(h)$、$b'=\tilde b'+d(h')$。于是
$$
b'b = \tilde{b}'\tilde{b} + d(\tilde{b}'h + h'\tilde{b} + h'd(h))
$$
，因为 $d(\tilde b)=d(\tilde b')=0$。也就是说，$b'b$ 与
$\tilde b'\tilde b$ 同伦。现证明 $\tilde b'\tilde b=0$。由于
$x_2\xrightarrow{f}y_2\xrightarrow{g}z_2$ 是容许短正合列，存在次数为 $0$
的态射 $\pi:y_2\to x_2$ 和 $s:z_2\to y_2$，使得
$\text{id}_{y_2}=f\pi+sg$。因此
$$
\tilde{b}'\tilde{b} = \tilde{b}'(f\pi + sg)\tilde{b} = 0
$$
，因为两个交换方块分别蕴含 $g\tilde b=0$ 和 $\tilde b'f=0$。
\end{proof}

\noindent
下述引理是引理 \ref{dga-lemma-triangle-independent-splittings} 的对应版本。

\begin{lemma}
\label{dga-lemma-analogue-triangle-independent-splittings}
在情形 \ref{dga-situation-ABC} 中，设
$0\to x\to y\to z\to 0$ 是 $\text{Comp}(\mathcal{A})$ 中的容许短正合列。
以引理 \ref{dga-lemma-get-triangle} 中定义的 $\delta:z\to x[1]$ 构成三角
$$
\xymatrix{x\ar[r] & y\ar[r] & z\ar[r]^{\delta} & x[1]}
$$
。在 $K(\mathcal{A})$ 中，除典范同构外，它不依赖于引理
\ref{dga-lemma-get-triangle} 中所作的选择。
\end{lemma}

\begin{proof}
假设 $\delta$ 由分裂
$$
\xymatrix{
x \ar@<0.5ex>[r]^{a} &
y \ar@<0.5ex>[r]^b\ar@<0.5ex>[l]^{\pi} &
z \ar@<0.5ex>[l]^s
}
$$
定义，而 $\delta'$ 由以 $\pi',s'$ 代替 $\pi,s$ 的分裂定义。那么
$$
s'-s = (a\pi + sb)(s'-s) = a\pi s'
$$
，因为 $bs'=bs=1_z$ 且 $\pi s=0$。类似地，
$$
\pi' - \pi = (\pi' - \pi)(a\pi + sb) = \pi'sb
$$
由于按引理 \ref{dga-lemma-get-triangle} 的构造，$\delta=\pi d(s)$ 且
$\delta'=\pi'd(s')$，可计算
$$
\delta' = \pi'd(s') = (\pi + \pi'sb)d(s + a\pi s') = \delta + d(\pi s')
$$
；这里使用了 $\pi a=1_x$、$ba=0$，以及由引理
\ref{dga-lemma-get-triangle} 的公式 (5) 得出的
$\pi'sbd(s')=\pi'sba\pi d(s')=0$。
\end{proof}

\noindent
下述引理是引理 \ref{dga-lemma-rotate-cone} 的对应版本。

\begin{lemma}
\label{dga-lemma-restate-axiom-c}
在情形 \ref{dga-situation-ABC} 中，设 $f:x\to y$ 是
$\text{Comp}(\mathcal{A})$ 中的态射。三角 $(y,c(f),x[1],i,p,f[1])$
是与容许短正合列
$$
\xymatrix{y\ar[r] & c(f) \ar[r] & x[1]}
$$
相联系的三角，其中锥 $c(f)$ 如引理 \ref{dga-lemma-get-triangle} 中所定义。
\end{lemma}

\begin{proof}
这由公理 (C) 得出。
\end{proof}

\noindent
下述引理是引理 \ref{dga-lemma-rotate-triangle} 的对应版本。

\begin{lemma}
\label{dga-lemma-cone-rotate-isom}
在情形 \ref{dga-situation-ABC} 中，设 $\alpha:x\to y$ 和 $\beta:y\to z$
给出 $\text{Comp}(\mathcal{A})$ 中的容许短正合列
$$
\xymatrix{
x \ar[r] &
y\ar[r] &
z
}
$$
。令 $(x,y,z,\alpha,\beta,\delta)$ 为 $K(\mathcal{A})$ 中相联系的三角。
那么三角
$$
(z[-1], x, y, \delta[-1], \alpha, \beta)
\quad\text{且}\quad
(z[-1], x, c(\delta[-1]), \delta[-1], i, p)
$$
同构。
\end{lemma}

\begin{proof}
有如下形式的图表
$$
\xymatrix{
z[-1]\ar[r]^{\delta[-1]}\ar[d]^1 &
x\ar@<0.5ex>[r]^{\alpha}\ar[d]^1 &
y\ar@<0.5ex>[r]^{\beta}\ar@{.>}[d]\ar@<0.5ex>[l]^{\tilde{\alpha}} &
z\ar[d]^1\ar@<0.5ex>[l]^{\tilde\beta} \\
z[-1] \ar[r]^{\delta[-1]} &
x\ar@<0.5ex>[r]^i &
c(\delta[-1]) \ar@<0.5ex>[r]^p\ar@<0.5ex>[l]^{\tilde i} &
z\ar@<0.5ex>[l]^{\tilde p}
}
$$
，其中 $\alpha,\beta,i,p$ 的分裂分别由
$\tilde\alpha,\tilde\beta,\tilde i,\tilde p$ 给出。把态射
$y\to c(\delta[-1])$ 定义为 $i\tilde\alpha+\tilde p\beta$，把态射
$c(\delta[-1])\to y$ 定义为 $\alpha\tilde i+\tilde\beta p$。先验证它们
确实是 $\text{Comp}(\mathcal{A})$ 中的态射。由引理
\ref{dga-lemma-get-triangle} 中的恒等式，有关系
$\delta[-1]=\tilde\alpha d(\tilde\beta)=-d(\tilde\alpha)\tilde\beta$
以及 $\delta[-1]=\tilde i d(\tilde p)$。于是
\begin{align*}
d(\tilde{\alpha})
& =
d(\tilde{\alpha})\tilde{\beta}\beta \\
& =
-\delta[-1]\beta
\end{align*}
，其中第一个等号使用了引理 \ref{dga-lemma-get-triangle} 的公式 (6)，第二个
等号使用了前述关系。类似地，得到 $d(\tilde p)=i\delta[-1]$。因此
\begin{align*}
d(i\tilde{\alpha} + \tilde{p}\beta)
& =
d(i)\tilde{\alpha} + id(\tilde{\alpha}) +
d(\tilde{p})\beta + \tilde{p}d(\beta) \\
& =
id(\tilde{\alpha}) + d(\tilde{p})\beta \\
& =
-i\delta[-1]\beta + i\delta[-1]\beta \\
& =
0
\end{align*}
，故 $i\tilde\alpha+\tilde p\beta$ 确实是 $\text{Comp}(\mathcal{A})$
中的态射。类似计算表明，$\alpha\tilde i+\tilde\beta p$ 也是
$\text{Comp}(\mathcal{A})$ 中的态射。显然这些态射组成上述交换图表。计算得
\begin{align*}
(i\tilde{\alpha} + \tilde{p}\beta)(\alpha \tilde{i} + \tilde{\beta} p)
& =
i\tilde{\alpha}\alpha\tilde{i} + i\tilde{\alpha}\tilde{\beta}p
+ \tilde{p}\beta\alpha\tilde{i} + \tilde{p}\beta\tilde{\beta}p \\
& =
i\tilde{i} + \tilde{p}p \\
& =
1_{c(\delta[-1])}
\end{align*}
，其中自由使用了引理 \ref{dga-lemma-get-triangle} 中的恒等式。类似地，计算得
$(\alpha \tilde{i} + \tilde{\beta} p)(i\tilde{\alpha} + \tilde{p}\beta) = 1_y$,
，故 $y\cong c(\delta[-1])$。因此，所讨论的两个三角同构。
\end{proof}

\noindent
下述引理是引理 \ref{dga-lemma-third-isomorphism} 的对应版本。

\begin{lemma}
\label{dga-lemma-analogue-third-isomorphism}
在情形 \ref{dga-situation-ABC} 中，设 $f_1:x_1\to y_1$ 和
$f_2:x_2\to y_2$ 是 $\text{Comp}(\mathcal{A})$ 中的态射。设
$$
(a,b,c): (x_1,y_1,c(f_1), f_1, i_1, p_1) \to (x_2,y_2, c(f_2), f_2, i_1, p_1)
$$
是 $K(\mathcal{A})$ 中任意三角态射。若 $a$ 和 $b$ 是同伦等价，则 $c$
也是同伦等价。
\end{lemma}

\begin{proof}
由于 $a,b$ 是同伦等价，它们在 $K(\mathcal{A})$ 中可逆。以
$a^{-1},b^{-1}$ 表示它们在 $K(\mathcal{A})$ 中的逆，得到交换图表
$$
\xymatrix{
x_2\ar[d]^{a^{-1}}\ar[r]^{f_2} &
y_2\ar[d]^{b^{-1}}\ar[r]^{i_2} &
c(f_2)\ar[d]^{c'} \\
x_1\ar[r]^{f_1} &
y_1 \ar[r]^{i_1} &
c(f_1)
}
$$
，其中映射 $c'$ 通过把引理 \ref{dga-lemma-cone} 应用于上图左侧交换方块而定义。
由于图表在 $K(\mathcal{A})$ 中交换，由引理 \ref{dga-lemma-triseq}，只需证明：
给定 $K(\mathcal{A})$ 中的三角态射
$(1,1,c): (x,y,c(f),f,i,p)\to (x,y,c(f),f,i,p)$
，映射 $c$ 是 $K(\mathcal{A})$ 中的同构。在 $K(\mathcal{A})$ 中有交换图表
$$
\vcenter{
\xymatrix{
y\ar[d]^{1}\ar[r] &
c(f)\ar[d]^{c}\ar[r] &
x[1]\ar[d]^{1} \\
y\ar[r] &
c(f) \ar[r] &
x[1]
}
}
\quad\Rightarrow\quad
\vcenter{
\xymatrix{
y\ar[d]^{0}\ar[r] &
c(f)\ar[d]^{c-1}\ar[r] &
x[1]\ar[d]^{0} \\
y\ar[r] &
c(f) \ar[r] &
x[1]
}
}
$$
由于各行都是容许短正合列，由引理 \ref{dga-lemma-triseq} 得到恒等式
$(c-1)^2=0$。由此推出 $2-c$ 是 $c$ 在 $K(\mathcal{A})$ 中的逆，
因而 $c$ 是同构。
\end{proof}

\noindent
下述引理是引理 \ref{dga-lemma-the-same-up-to-isomorphisms} 的对应版本。

\begin{lemma}
\label{dga-lemma-cone-homotopy}
在情形 \ref{dga-situation-ABC} 中：
\begin{enumerate}
\item 给定容许短正合列 $x\xrightarrow{\alpha}y\xrightarrow{\beta}z$。
那么存在同伦等价 $e:C(\alpha)\to z$，使图表
\begin{equation}
\label{dga-equation-cone-isom-triangle}
\vcenter{
\xymatrix{
x\ar[r]^{\alpha}\ar[d] &
y\ar[r]^{b}\ar[d] &
C(\alpha)\ar[r]^{-c}\ar@{.>}[d]^{e} &
x[1]\ar[d] \\
x\ar[r]^{\alpha} &
y\ar[r]^{\beta} &
z\ar[r]^{\delta} & x[1]
}
}
\end{equation}
在 $K(\mathcal{A})$ 中定义三角同构。这里
$y\xrightarrow{b}C(\alpha)\xrightarrow{c}x[1]$ 是按公理 (C) 给出的容许
短正合列。
\item 给定 $\text{Comp}(\mathcal{A})$ 中的态射 $\alpha:x\to y$，令
$x\xrightarrow{\tilde\alpha}\tilde y\to y$ 为引理 \ref{dga-lemma-factor}
中给出的分解，其中容许单态射 $x\xrightarrow{\tilde\alpha}y$ 延拓为容许短正合列
$$
\xymatrix{
x \ar[r]^{\tilde{\alpha}} &
\tilde{y} \ar[r] & z
}
$$
。那么存在三角同构
$$
\xymatrix{
x \ar[r]^{\tilde{\alpha}} \ar[d] &
\tilde{y} \ar[r] \ar[d] &
z \ar[r]^{\delta} \ar@{.>}[d]^{e} &
x[1] \ar[d] \\
x \ar[r]^{\alpha} &
y \ar[r] &
C(\alpha) \ar[r]^{-c} &
x[1]
}
$$
，其中上方三角是与列 $x\xrightarrow{\tilde\alpha}\tilde y\to z$ 相联系的三角。
\end{enumerate}
\end{lemma}

\begin{proof}
对 (1)，考虑下述更完整的图表，其中\emph{不}改变 $c$ 的符号：
$$
\xymatrix{
x\ar@<0.5ex>[r]^{\alpha} \ar[d] &
y\ar@<0.5ex>[l]^{\pi} \ar@<0.5ex>[r]^{b}\ar[d] &
C(\alpha)\ar@<0.5ex>[l]^{p} \ar@<0.5ex>[r]^{c}\ar@{.>}@<0.5ex>[d]^{e} &
x[1]\ar@<0.5ex>[l]^{\sigma} \ar[d]\ar@<0.5ex>[r]^{\alpha} &
y[1]\ar@<0.5ex>[l]^{\pi} \\
x\ar@<0.5ex>[r]^{\alpha} &
y\ar@<0.5ex>[r]^{\beta} \ar@<0.5ex>[l]^{\pi} &
z\ar[r]^{\delta}\ar@<0.5ex>[l]^{s} \ar@{.>}@<0.5ex>[u]^{f} &
x[1]
}
$$
其中容许短正合列 $x\xrightarrow{\alpha}y\xrightarrow{\beta}z$ 带有分裂
$\pi,s$，容许短正合列
$y\xrightarrow{b}C(\alpha)\xrightarrow{c}x[1]$ 带有分裂 $p,\sigma$。
注意（在移位下识别 Hom 集），由引理 \ref{dga-lemma-get-triangle} 中的构造，
$$
\alpha = pd(\sigma) = -d(p)\sigma,\quad
\delta = \pi d(s) = -d(\pi)s
$$
。

\medskip\noindent
定义 $e=\beta p$ 和 $f=bs-\sigma\delta$。先验证它们是
$\text{Comp}(\mathcal{A})$ 中的态射。为证明 $d(e)=\beta d(p)$ 为零，
只需证明 $\beta d(p)b$ 和 $\beta d(p)\sigma$ 均为零，而
$$
\beta d(p)b = \beta d(pb) = \beta d(1_y) = 0,\quad
\beta d(p)\sigma = -\beta\alpha = 0
$$
。类似地，为验证 $d(f)=bd(s)-d(\sigma)\delta$ 为零，只需验证它与 $p$
及 $c$ 的后复合都为零，而
\begin{align*}
pbd(s) - pd(\sigma)\delta
= &
d(s)-\alpha\delta = d(s)-\alpha\pi d(s) = 0 \\
cbd(s)-cd(\sigma)\delta
= &
-cd(\sigma)\delta = -d(c\sigma)\delta = 0
\end{align*}
。图表 \ref{dga-equation-cone-isom-triangle} 左侧两个方块的交换性直接由定义得出。
在证明右侧方块在同伦意义下交换之前，先验证 $e$ 是同伦等价。显然，
$$
ef=\beta p (bs-\sigma\delta)=\beta s=1_z
$$
。为验证 $fe$ 与 $1_{C(\alpha)}$ 同伦，先注意
$$
b\alpha = bpd(\alpha) = d(\sigma),\quad
\alpha c = -d(p)\sigma c = -d(p),\quad
d(\pi)p = d(\pi)s\beta p = -\delta\beta p
$$
。利用这些恒等式，计算得
\begin{align*}
1_{C(\alpha)} = &
bp + \sigma c
\quad (\text{由 }y \xrightarrow{b} C(\alpha) \xrightarrow{c} x[1]) \\
= &
b(\alpha\pi + s\beta)p + \sigma(\pi\alpha)c
\quad (\text{由 }x \xrightarrow{\alpha} y \xrightarrow{\beta} z) \\
= &
d(\sigma)\pi p + bs\beta p - \sigma\pi d(p)
\quad (\text{由上述前两个恒等式}) \\
= &
d(\sigma)\pi p + bs\beta p - \sigma\delta\beta p
+ \sigma\delta\beta p - \sigma\pi d(p) \\
= &
(bs - \sigma\delta)\beta p + d(\sigma)\pi p
- \sigma d(\pi)p - \sigma\pi d(p)\quad
(\text{由上述第三个恒等式}) \\
= &
fe + d(\sigma \pi p)
\end{align*}
，因为 $\sigma\in\Hom^{-1}(x,C(\alpha))$（参见引理
\ref{dga-lemma-id-cone-null} 的证明）。因此 $e,f$ 互为同伦逆。最后，为验证
图表 \ref{dga-equation-cone-isom-triangle} 的右侧方块在同伦意义下交换，
只须验证 $-cf=\delta$。这由下式得出：
$$
-cf = -c(bs-\sigma\delta) = c\sigma\delta = \delta
$$
，因为 $cb=0$。

\medskip\noindent
对 (2)，考虑引理 \ref{dga-lemma-factor} 中给出的分解
$x\xrightarrow{\tilde\alpha}\tilde y\to y$，故第二个态射是同伦等价。
由引理 \ref{dga-lemma-cone} 和 \ref{dga-lemma-analogue-third-isomorphism}，
下列两个三角之间存在三角同构：
$$
x \xrightarrow{\alpha} y \to C(\alpha) \to x[1]
\quad\text{且}\quad
x \xrightarrow{\tilde{\alpha}} \tilde{y} \to C(\tilde{\alpha}) \to x[1]
$$
。由于可以复合三角同构，通过以 $\tilde\alpha$ 替代 $\alpha$、以
$\tilde y$ 替代 $y$、以 $C(\tilde\alpha)$ 替代 $C(\alpha)$，可假设
$\alpha$ 是容许单态射。在这种情形下，结论由 (1) 得出。
\end{proof}

\noindent
下述引理是引理 \ref{dga-lemma-homotopy-category-pre-triangulated} 的对应版本。

\begin{lemma}
\label{dga-lemma-analogue-homotopy-category-pre-triangulated}
在情形 \ref{dga-situation-ABC} 中，同伦范畴 $K(\mathcal{A})$ 连同其自然平移
函子和有区别三角，是预三角范畴。
\end{lemma}

\begin{proof}
逐一验证 TR1、TR2 和 TR3。

\medskip\noindent
TR1 的证明。由定义，每个与有区别三角同构的三角都是有区别的。由于
$$
\xymatrix{x\ar[r]^{1_x} & x\ar[r] & 0}
$$
是容许短正合列，$(x,x,0,1_x,0,0)$ 是有区别三角。此外，给定
$\text{Comp}(\mathcal{A})$ 中的态射 $\alpha:x\to y$，由引理
\ref{dga-lemma-cone-homotopy}，三角 $(x,y,c(\alpha),\alpha,i,-p)$ 是有区别的。

\medskip\noindent
TR2 的证明。设 $(x,y,z,\alpha,\beta,\gamma)$ 是三角，并假设
$(y,z,x[1],\beta,\gamma,-\alpha[1])$ 是有区别的。则存在容许短正合列
$0\to x'\to y'\to z'\to0$，使相联系的三角
$(x',y',z',\alpha',\beta',\gamma')$ 同构于
$(y,z,x[1],\beta,\gamma,-\alpha[1])$。旋转后，得出
$(x,y,z,\alpha,\beta,\gamma)$ 同构于
$(z'[-1],x',y',\gamma'[-1],\alpha',\beta')$。由引理
\ref{dga-lemma-cone-rotate-isom}，后者同构于
$(z'[-1],x',c(\gamma'[-1]),\gamma'[-1],i,p)$。按下图所示改变符号并复合
这两个同构：
$$
\xymatrix@C=3pc{
x\ar[r]^{\alpha}\ar[d] &
y\ar[r]^{\beta}\ar[d] &
z\ar[r]^{\gamma}\ar[d] &
x[1]\ar[d] \\
z'[-1]\ar[r]^{-\gamma'[-1]}\ar[d]_{-1_{z'[-1]}} &
x \ar[r]^{\alpha'}\ar@{=}[d] &
y' \ar[r]^{\beta'} \ar[d] &
z'\ar[d]^{-1_{z'}} \\
z'[-1]\ar[r]^{\gamma'[-1]} &
x \ar[r]^{\alpha'} &
c(\gamma'[-1]) \ar[r]^{-p} &
z'
}
$$
由引理 \ref{dga-lemma-cone-homotopy} (2)，推出
$(x,y,z,\alpha,\beta,\gamma)$ 是有区别的。反过来，假设
$(x,y,z,\alpha,\beta,\gamma)$ 是有区别的。由引理
\ref{dga-lemma-cone-homotopy} (1)，它同构于形如
$(x',y',c(\alpha'),\alpha',i,-p)$ 的三角，其中
$\alpha':x'\to y'$ 是 $\text{Comp}(\mathcal{A})$ 中的某个态射。
旋转后的三角 $(y,z,x[1],\beta,\gamma,-\alpha[1])$ 同构于
$(y',c(\alpha'),x'[1],i,-p,-\alpha[1])$，而后者同构于
$(y',c(\alpha'),x'[1],i,p,\alpha[1])$。由引理
\ref{dga-lemma-restate-axiom-c}，这个三角是有区别的，从而
$(y,z,x[1],\beta,\gamma,-\alpha[1])$ 是有区别的。

\medskip\noindent
TR3 的证明。设 $(x,y,z,\alpha,\beta,\gamma)$ 和
$(x',y',z',\alpha',\beta',\gamma')$ 是 $\text{Comp}(\mathcal{A})$
中的有区别三角，且 $f:x\to x'$、$g:y\to y'$ 是满足
$\alpha'\circ f=g\circ\alpha$ 的态射。由引理
\ref{dga-lemma-cone-homotopy}，可假设
$(x,y,z,\alpha,\beta,\gamma)= (x,y,c(\alpha),\alpha, i, -p)$
，且
$(x',y',z',\alpha',\beta',\gamma')=(x',y',c(\alpha'),\alpha',i',-p')$。
现应用引理 \ref{dga-lemma-cone} 即得结论。
\end{proof}

\noindent
下述引理是引理 \ref{dga-lemma-two-split-injections} 的对应版本。

\begin{lemma}
\label{dga-lemma-dgc-analogue-tr4}
在情形 \ref{dga-situation-ABC} 中，给定 $\mathcal{A}$ 中的容许单态射
$x\xrightarrow{\alpha}y$、$y\xrightarrow{\beta}z$，存在有区别三角
$(x,y,q_1,\alpha,p_1,\delta_1)$、$(x,z,q_2,\beta\alpha,p_2,\delta_2)$
和 $(y,z,q_3,\beta,p_3,\delta_3)$，使 TR4 成立。
\end{lemma}

\begin{proof}
给定容许单态射 $x\xrightarrow{\alpha}y$ 和 $y\xrightarrow{\beta}z$，
通过将它们延拓为容许短正合列，可以得到有区别三角
$$
\xymatrix{
x\ar@<0.5ex>[r]^{\alpha} &
y\ar@<0.5ex>[l]^{\pi_1} \ar@<0.5ex>[r]^{p_1} &
q_1 \ar[r]^{\delta_1} \ar@<0.5ex>[l]^{s_1} &
x[1]
}
$$
$$
\xymatrix{
x\ar@<0.5ex>[r]^{\beta\alpha} &
z\ar@<0.5ex>[l]^{\pi_1\pi_3} \ar@<0.5ex>[r]^{p_2} &
q_2 \ar[r]^{\delta_2} \ar@<0.5ex>[l]^{s_2} &
x[1]
}
$$
$$
\xymatrix{
y\ar@<0.5ex>[r]^{\beta} &
z\ar@<0.5ex>[l]^{\pi_3} \ar@<0.5ex>[r]^{p_3} &
q_3 \ar[r]^{\delta_3} \ar@<0.5ex>[l]^{s_3} &
x[1]
}
$$
。在这些图表中，映射 $\delta_i$ 类似于引理 \ref{dga-lemma-get-triangle}
中定义的映射，定义为 $\delta_i=\pi_i d(s_i)$。它们构成下述实线交换图表
$$
\xymatrix@C=5pc@R=3pc{
x\ar@<0.5ex>[r]^{\alpha} \ar@<0.5ex>[dr]^{\beta\alpha} &
y\ar@<0.5ex>[d]^{\beta} \ar@<0.5ex>[l]^{\pi_1} \ar@<0.5ex>[r]^{p_1} &
q_1 \ar[r]^{\delta_1} \ar@<0.5ex>[l]^{s_1} \ar@{.>}[dd]^{p_2\beta s_1} &
x[1] \\
 &
z \ar@<0.5ex>[u]^{\pi_3}\ar@<0.5ex>[d]^{p_3}
\ar@<0.5ex>[dr]^{p_2} \ar@<0.5ex>[ul]^{\pi_1\pi_3} & & \\
 &
q_3\ar@<0.5ex>[u]^{s_3} \ar[d]^{\delta_3} &
q_2 \ar@{.>}[l]^{p_3s_2} \ar@<0.5ex>[ul]^{s_2} \ar[dr]^{\delta_2} \\
 &
y[1] & & x[1]}
$$
，其中虚线箭头按图中所示定义。显然，它们的复合
$p_3s_2p_2\beta s_1=0$，因为 $s_2p_2=0$。断言两者都是
$\text{Comp}(\mathcal{A})$ 中的态射。可利用引理
\ref{dga-lemma-get-triangle} 中的等式验证这一点：
$$
d(p_2\beta s_1) = p_2\beta d(s_1) = p_2\beta\alpha\pi_1 d(s_1) = 0
$$
，因为 $p_2\beta\alpha=0$；并且
$$
d(p_3s_2) = p_3d(s_2) = p_3\beta\alpha\pi_1\pi_3 d(s_2) = 0
$$
，因为 $p_3\beta=0$。为验证 $q_1\to q_2\to q_3$ 是容许短正合列，
还须证明在底层分次范畴中 $q_2=q_1\oplus q_3$，且上述两个态射分别是
余投影和投影。为此，注意在底层分次范畴 $\mathcal{C}$ 中有
$$
y = x\oplus q_1,\quad
z = y\oplus q_3 = x\oplus q_1\oplus q_3
$$
，其中 $\pi_1\pi_3$ 给出到第一个因子的投影态射
$x\oplus q_1\oplus q_3\to z$。由 $\mathcal{A}$ 的公理 (A)，
$\mathcal{C}$ 是加性范畴。因此可应用《同调代数》引理
\ref{homology-lemma-additive-cat-biproduct-kernel}，得出
$$
\Ker(\pi_1\pi_3) = q_1\oplus q_3
$$
在 $\mathcal{C}$ 中成立。再把《同调代数》引理
\ref{homology-lemma-additive-cat-biproduct-kernel} 应用于
$z=x\oplus q_2$，得到 $\Ker(\pi_1\pi_3)=q_2$。故在 $\mathcal{C}$ 中
$q_2\cong q_1\oplus q_3$。显然，上述虚线态射给出余投影和投影。

\medskip\noindent
最后，还须验证容许短正合列 $q_1\to q_2\to q_3$ 诱导的态射
$\delta:q_3\to q_1[1]$ 与 $p_1\delta_3$ 一致。由引理
\ref{dga-lemma-get-triangle} 中的构造，态射 $\delta$ 由下式给出：
\begin{align*}
p_1\pi_3s_2d(p_2s_3)
= &
p_1\pi_3s_2p_2d(s_3) \\
= &
p_1\pi_3(1-\beta\alpha\pi_1\pi_3)d(s_3) \\
= &
p_1\pi_3d(s_3)\quad (\text{因为 }\pi_3\beta = 0) \\
= &
p_1\delta_3
\end{align*}
，这正是所需等式。证明完成。
\end{proof}

\noindent
汇集以上结果，终于得到命题
\ref{dga-proposition-homotopy-category-triangulated} 的对应版本。

\begin{proposition}
\label{dga-proposition-ABC-homotopy-category-triangulated}
在情形 \ref{dga-situation-ABC} 中，同伦范畴 $K(\mathcal{A})$ 连同其自然平移
函子和有区别三角，是三角范畴。
\end{proposition}

\begin{proof}
由引理 \ref{dga-lemma-analogue-homotopy-category-pre-triangulated}，
$K(\mathcal{A})$ 是预三角范畴。把引理
\ref{dga-lemma-analogue-sequence-maps-split}、\ref{dga-lemma-dgc-analogue-tr4}
与《导出范畴》引理 \ref{derived-lemma-easier-axiom-four} 结合，得出
$K(\mathcal{A})$ 是三角范畴。
\end{proof}

\begin{lemma}
\label{dga-lemma-functor-between-ABC}
设 $R$ 是环，$F:\mathcal{A}\to\mathcal{B}$ 是满足公理 (A)、(B)、(C)
的 $R$ 上微分分次范畴之间的函子，并且 $F(x[1])=F(x)[1]$。
那么 $F$ 诱导三角范畴之间的正合函子
$K(\mathcal{A})\to K(\mathcal{B})$。
\end{lemma}

\begin{proof}
确实，若 $x\to y\to z$ 是 $\text{Comp}(\mathcal{A})$ 中的容许短正合列，
则 $F(x)\to F(y)\to F(z)$ 是 $\text{Comp}(\mathcal{B})$ 中的容许短正合列。
此外，引理 \ref{dga-lemma-get-triangle} 中构造的“边界”态射
$\delta=\pi\text{d}(s):z\to x[1]$ 产生态射
$F(\delta):F(z)\to F(x[1])=F(x)[1]$；它等于容许短正合列
$F(x)\to F(y)\to F(z)$ 的边界映射 $F(\pi)\text{d}(F(s))$。
\end{proof}





\section{双模}
\label{dga-section-bimodules}

\noindent
继续第 \ref{dga-section-tensor-product} 节开始的讨论。

\begin{definition}
\label{dga-definition-bimodule}
双模。设 $R$ 是环。
\begin{enumerate}
\item 设 $A,B$ 是 $R$-代数。{\it $(A,B)$-双模}是带有 $R$-双线性映射
$$
A \times M \to M, (a, x) \mapsto ax
\quad\text{且}\quad
M \times B \to M, (x, b) \mapsto xb
$$
的 $R$-模 $M$，并满足
\begin{enumerate}
\item $a'(ax)=(a'a)x$ 且 $(xb)b'=x(bb')$；
\item $a(xb)=(ax)b$；以及
\item $1x=x=x1$。
\end{enumerate}
\item 设 $A,B$ 是 $\mathbf{Z}$-分次 $R$-代数。{\it 分次 $(A,B)$-双模}
是带有分次 $M=\bigoplus M^n$ 的 $(A,B)$-双模 $M$，使得
$A^nM^m\subset M^{n+m}$ 且 $M^nB^m\subset M^{n+m}$。
\item 设 $A,B$ 是微分分次 $R$-代数。{\it 微分分次 $(A,B)$-双模}
是带有次数为 $1$ 的齐次微分 $\text{d}:M\to M$ 的分次 $(A,B)$-双模，
使得对齐次元素 $a\in A$、$x\in M$、$b\in B$，有
$\text{d}(ax)=\text{d}(a)x+(-1)^{\deg(a)}a\text{d}(x)$ 且
$\text{d}(xb)=\text{d}(x)b+(-1)^{\deg(x)}x\text{d}(b)$。
\end{enumerate}
\end{definition}

\noindent
注意，微分分次 $(A,B)$-双模 $M$ 与下述数据是一回事：$M$ 是右微分分次
$B$-模，同时也是左微分分次 $A$-模，两个分次和微分彼此一致，并且
$A$-模结构与 $B$-模结构交换。精确表述如下。

\begin{lemma}
\label{dga-lemma-what-makes-a-bimodule-dg}
设 $R$ 是环，$(A,\text{d})$ 和 $(B,\text{d})$ 是 $R$ 上的微分分次代数，
$M$ 是右微分分次 $B$-模。$M$ 上与给定微分分次 $B$-模结构相容的
$(A,B)$-双模结构，与微分分次 $R$-代数同态
$$
A
\longrightarrow
\Hom_{\text{Mod}^{dg}_{(B, \text{d})}}(M, M)
$$
之间有一一对应。
\end{lemma}

\begin{proof}
设 $\mu:A\times M\to M$ 在 $M$ 的底层 $R$-模复形 $M^\bullet$ 上定义
左微分分次 $A$-模结构。由引理 \ref{dga-lemma-left-module-structure}，结构
$\mu$ 对应于微分分次 $R$-代数的映射
$\gamma:A\to\Hom^\bullet(M^\bullet,M^\bullet)$。本引理的断言恰是：
$\mu$ 与 $B$-作用交换，当且仅当 $\gamma$ 的像落在
$$
\Hom_{\text{Mod}^{dg}_{(B, \text{d})}}(M, M) \subset
\Hom^\bullet(M^\bullet, M^\bullet)
$$
中。详细计算略。
\end{proof}

\noindent
设 $M$ 是微分分次 $(A,B)$-双模。回忆第 \ref{dga-section-left-modules} 节：
左微分分次 $A$-模结构对应于右微分分次 $A^{opp}$-模结构。由于 $A$-模
结构与 $B$-模结构交换，这赋予 $M$ 微分分次
$A^{opp}\otimes_R B$-模结构：
$$
x \cdot (a \otimes b) = (-1)^{\deg(a)\deg(x)} axb
$$
。反过来，若有微分分次 $A^{opp}\otimes_R B$-模 $M$，则可利用上式得到
微分分次 $(A,B)$-双模。

\begin{lemma}
\label{dga-lemma-bimodule-over-tensor}
设 $R$ 是环，$(A,\text{d})$ 和 $(B,\text{d})$ 是 $R$ 上的微分分次代数。
上述构造定义范畴等价
$$
\begin{matrix}
\text{微分分次}\\
(A, B)\text{-双模}
\end{matrix}
\longleftrightarrow
\begin{matrix}
\text{右微分分次}\\
A^{opp} \otimes_R B\text{-模}
\end{matrix}
$$
\end{lemma}

\begin{proof}
由上述讨论立即得出。
\end{proof}

\noindent
设 $R$ 是环，$(A,\text{d})$ 和 $(B,\text{d})$ 是微分分次 $R$-代数。
设 $P$ 是微分分次 $(A,B)$-双模。若存在滤过
$$
0 = F_{-1}P \subset F_0P \subset F_1P \subset \ldots \subset P
$$
，它由微分分次 $(A,B)$-双模组成，并且
\begin{enumerate}
\item $P=\bigcup F_pP$；
\item 包含映射 $F_iP\to F_{i+1}P$ 作为分次 $(A,B)$-双模映射是分裂的；
\item 作为微分分次 $(A,B)$-双模，商 $F_{i+1}P/F_iP$ 同构于若干个
$(A\otimes_R B)[k]$ 的直和，
\end{enumerate}
则称 $P$ {\it 具有性质 (P)}。

\begin{lemma}
\label{dga-lemma-bimodule-resolve}
设 $R$ 是环，$(A,\text{d})$ 和 $(B,\text{d})$ 是微分分次 $R$-代数，
$M$ 是微分分次 $(A,B)$-双模。存在微分分次 $(A,B)$-双模同态 $P\to M$，
它是拟同构，并且 $P$ 具有上面定义的性质 (P)。
\end{lemma}

\begin{proof}
由引理 \ref{dga-lemma-bimodule-over-tensor} 和 \ref{dga-lemma-resolve} 立即得出。
\end{proof}

\begin{lemma}
\label{dga-lemma-bimodule-property-P-sequence}
设 $R$ 是环，$(A,\text{d})$ 和 $(B,\text{d})$ 是微分分次 $R$-代数。
设 $P$ 是具有性质 (P) 的微分分次 $(A,B)$-双模，其相应滤过为
$F_\bullet$。那么得到短正合列
$$
0 \to
\bigoplus\nolimits F_iP \to
\bigoplus\nolimits F_iP \to P \to 0
$$
，它由微分分次 $(A,B)$-双模组成，并且作为分次 $(A,B)$-双模列是分裂的。
\end{lemma}

\begin{proof}
由引理 \ref{dga-lemma-bimodule-over-tensor} 和
\ref{dga-lemma-property-P-sequence} 立即得出。
\end{proof}















\section{双模与张量积}
\label{dga-section-bimodules-tensor}

\noindent
设 $R$ 是环，$A,B$ 是 $R$-代数，$M$ 是右 $A$-模，$N$ 是
$(A,B)$-双模。那么 $M\otimes_A N$ 是右 $B$-模。

\medskip\noindent
若在上一段的情形中，$A,B$ 是 $\mathbf{Z}$-分次代数，$M$ 是分次
$A$-模，$N$ 是分次 $(A,B)$-双模，则 $M\otimes_A N$ 是右分次 $B$-模。
该构造关于 $M$ 具有函子性，并定义例 \ref{dga-example-gm-gr-cat} 中分次范畴
之间的函子
$$
- \otimes_A N :
\text{Mod}^{gr}_A
\longrightarrow
\text{Mod}^{gr}_B
$$
。确实，若 $M,M'$ 是分次 $A$-模，$f:M\to M'$ 是次数 $n$ 齐次的
$A$-模同态，则
$f\otimes\text{id}_N:M\otimes_A N\to M'\otimes_A N$ 是次数 $n$
齐次的 $B$-模同态。

\medskip\noindent
若在上一段的情形中，$(A,\text{d})$ 和 $(B,\text{d})$ 是微分分次代数，
$M$ 是微分分次 $A$-模，$N$ 是微分分次 $(A,B)$-双模，则
$M\otimes_A N$ 是右微分分次 $B$-模。

\begin{lemma}
\label{dga-lemma-tensor}
设 $R$ 是环，$(A,\text{d})$ 和 $(B,\text{d})$ 是 $R$ 上的微分分次代数，
$N$ 是微分分次 $(A,B)$-双模。那么 $M\mapsto M\otimes_A N$ 定义微分分次
范畴之间的函子
$$
- \otimes_A N :
\text{Mod}^{dg}_{(A, \text{d})}
\longrightarrow
\text{Mod}^{dg}_{(B, \text{d})}
$$
。应用引理 \ref{dga-lemma-functorial}，此函子诱导函子
$$
\text{Mod}_{(A, \text{d})} \to \text{Mod}_{(B, \text{d})}
\quad\text{以及}\quad
K(\text{Mod}_{(A, \text{d})}) \to K(\text{Mod}_{(B, \text{d})})
$$
。
\end{lemma}

\begin{proof}
上面已经看到该构造如何定义底层分次范畴之间的函子。因此，只需证明该构造
与微分相容。设 $M,M'$ 是微分分次 $A$-模，$f:M\to M'$ 是次数 $n$
齐次的 $A$-模同态。那么
$$
\text{d}(f) = \text{d}_{M'} \circ f - (-1)^n f \circ \text{d}_M
$$
。另一方面，
$$
\text{d}(f \otimes \text{id}_N) =
\text{d}_{M' \otimes_A N} \circ
(f \otimes \text{id}_N)
- (-1)^n
(f \otimes \text{id}_N) \circ \text{d}_{M \otimes_A N}
$$
。把它应用于元素 $x\otimes y$，其中 $x\in M$、$y\in N$ 是齐次元素，得到
\begin{align*}
\text{d}(f \otimes \text{id}_N)(x \otimes y)
= &
\text{d}_{M'}(f(x)) \otimes y + (-1)^{n + \deg(x)}f(x) \otimes \text{d}_N(y) \\
& - (-1)^n f(\text{d}_M(x)) \otimes y
- (-1)^{n + \deg(x)}f(x) \otimes \text{d}_N(y) \\
= &
\text{d}(f) (x \otimes y)
\end{align*}
故 $\text{d}(f)\otimes\text{id}_N=
\text{d}(f\otimes\text{id}_N)$。证明完成。
\end{proof}

\begin{remark}
\label{dga-remark-shift-tensor-no-sign}
设 $R$ 是环，$(A,\text{d})$ 和 $(B,\text{d})$ 是 $R$ 上的微分分次代数，
$N$ 是微分分次 $(A,B)$-双模，$M$ 是右微分分次 $A$-模。那么对每个
$k\in\mathbf{Z}$，有右微分分次 $B$-模同构
$$
(M \otimes_A N)[k] \longrightarrow
M[k] \otimes_A N
$$
，其定义不引入符号；见《代数进阶》第
\ref{more-algebra-section-sign-rules} 节。
\end{remark}

\noindent
若有环 $R$、$R$-代数 $A,B,C$、右 $A$-模 $M$、$(A,B)$-双模 $N$
以及 $(B,C)$-双模 $N'$，则 $N\otimes_B N'$ 是 $(A,C)$-双模，并且
$$
(M \otimes_A N) \otimes_B N' = M \otimes_A (N \otimes_B N')
$$
。这个等式在分次情形和微分分次情形中仍成立。符号法则见《代数进阶》第
\ref{more-algebra-section-sign-rules} 节。






\section{双模与内部 Hom}
\label{dga-section-bimodules-hom}

\noindent
设 $R$ 是环。若 $A$ 是 $R$-代数（见第 \ref{dga-section-conventions} 节中的约定），
而 $M,M'$ 是右 $A$-模，则像通常一样定义
$$
\Hom_A(M, M') = \{f : M \to M' \mid f \text{ 是 }A\text{-线性的}\}
$$
。

\medskip\noindent
设 $R$ 是环，$A,B$ 是 $R$-代数，$N$ 是 $(A,B)$-双模，$N'$ 是右
$B$-模。在这种情形下，通过前复合把
$$
\Hom_B(N, N')
$$
看作右 $A$-模。

\medskip\noindent
设 $R$ 是环，$A,B$ 是 $\mathbf{Z}$-分次 $R$-代数，$N$ 是分次
$(A,B)$-双模，$N'$ 是右分次 $B$-模。在这种情形下，通过前复合把例
\ref{dga-example-gm-gr-cat} 中定义的分次 $R$-模
$$
\Hom_{\text{Mod}^{gr}_B}(N, N')
$$
看作右分次 $A$-模。该构造关于 $N'$ 具有函子性，并定义例
\ref{dga-example-gm-gr-cat} 中分次范畴之间的函子
$$
\Hom_{\text{Mod}^{gr}_B}(N, -) :
\text{Mod}^{gr}_B
\longrightarrow
\text{Mod}^{gr}_A
$$
。确实，若 $N_1,N_2$ 是分次 $B$-模，而 $f:N_1\to N_2$ 是次数 $n$
齐次的 $B$-模同态，则诱导映射
$\Hom_{\text{Mod}^{gr}_B}(N,N_1)\to
\Hom_{\text{Mod}^{gr}_B}(N,N_2)$ 是次数 $n$ 齐次的 $A$-模同态。

\medskip\noindent
设 $R$ 是环，$A,B$ 是微分 $\mathbf{Z}$-分次 $R$-代数，$N$ 是微分分次
$(A,B)$-双模，$N'$ 是右微分分次 $B$-模。在这种情形下，通过如分次情形
中的前复合，把例 \ref{dga-example-dgm-dg-cat} 中定义的微分分次 $R$-模
$$
\Hom_{\text{Mod}^{dg}_{(B, \text{d})}}(N, N')
$$
看作右微分分次 $A$-模。这与微分相容，因为乘法是复合
$$
\Hom_{\text{Mod}^{dg}_B}(N, N') \otimes_R A \to
\Hom_{\text{Mod}^{dg}_B}(N, N') \otimes_R
\Hom_{\text{Mod}^{dg}_B}(N, N) \to
\Hom_{\text{Mod}^{dg}_B}(N, N')
$$
。第一个箭头使用引理 \ref{dga-lemma-what-makes-a-bimodule-dg} 中的映射，
第二个箭头是微分分次范畴 $\text{Mod}^{dg}_{(B, \text{d})}$ 中的复合。

\begin{lemma}
\label{dga-lemma-hom}
设 $R$ 是环，$(A,\text{d})$ 和 $(B,\text{d})$ 是 $R$ 上的微分分次代数，
$N$ 是微分分次 $(A,B)$-双模。上述构造定义微分分次范畴之间的函子
$$
\Hom_{\text{Mod}^{dg}_{(B, \text{d})}}(N, -) :
\text{Mod}^{dg}_{(B, \text{d})}
\longrightarrow
\text{Mod}^{dg}_{(A, \text{d})}
$$
。应用引理 \ref{dga-lemma-functorial}，此函子诱导函子
$$
\text{Mod}_{(B, \text{d})} \to \text{Mod}_{(A, \text{d})}
\quad\text{以及}\quad
K(\text{Mod}_{(B, \text{d})}) \to K(\text{Mod}_{(A, \text{d})})
$$
。
\end{lemma}

\begin{proof}
上面已经看到该构造如何定义底层分次范畴之间的函子。因此，只需证明该构造
与微分相容。设 $N_1,N_2$ 是微分分次 $B$-模。记
$$
H_{12} = \Hom_{\text{Mod}^{dg}_{(B, \text{d})}}(N_1, N_2),\quad
H_1 = \Hom_{\text{Mod}^{dg}_{(B, \text{d})}}(N, N_1),\quad
H_2 = \Hom_{\text{Mod}^{dg}_{(B, \text{d})}}(N, N_2)
$$
。考虑微分分次范畴 $\text{Mod}^{dg}_{(B, \text{d})}$ 中的复合
$$
c : H_{12} \otimes_R H_1 \longrightarrow H_2
$$
。设 $f:N_1\to N_2$ 是次数 $n$ 齐次的 $B$-模同态；换言之，
$f\in H_{12}^n$。引理中的函子把 $f$ 映到
$c_f:H_1\to H_2$、$g\mapsto c(f,g)$。$\text{d}(f)$ 也类似。另一方面，
$$
\Hom_{\text{Mod}^{dg}_{(A, \text{d})}}(H_1, H_2)
$$
上的微分把 $c_f$ 映到
$\text{d}_{H_2}\circ c_f-(-1)^n c_f\circ\text{d}_{H_1}$。由于 $c$ 是
$R$-模复形之间的态射，有
$\text{d}c(f,g)=c(\text{d}f,g)+(-1)^n c(f,\text{d}g)$。因此
\begin{align*}
(\text{d}c_f)(g)
& =
\text{d}c(f,g) - (-1)^n c(f, \text{d}g) \\
& =
c(\text{d}f, g) + (-1)^n c(f, \text{d}g)  - (-1)^n c(f, \text{d}g) \\
& =
c(\text{d}f, g) = c_{\text{d}f}(g)
\end{align*}
。证明完成。
\end{proof}

\begin{remark}
\label{dga-remark-shift-hom-no-sign}
设 $R$ 是环，$(A,\text{d})$ 和 $(B,\text{d})$ 是 $R$ 上的微分分次代数，
$N$ 是微分分次 $(A,B)$-双模，$N'$ 是右微分分次 $B$-模。那么对每个
$k\in\mathbf{Z}$，有右微分分次 $A$-模同构
$$
\Hom_{\text{Mod}^{gr}_B}(N, N')[k]
\longrightarrow
\Hom_{\text{Mod}^{gr}_B}(N, N'[k])
$$
，其定义不引入符号；见《代数进阶》第
\ref{more-algebra-section-sign-rules} 节。
\end{remark}

\begin{lemma}
\label{dga-lemma-tensor-hom-adjunction}
设 $R$ 是环，$A,B$ 是 $R$-代数，$M$ 是右 $A$-模，$N$ 是
$(A,B)$-双模，$N'$ 是右 $B$-模。那么有 $R$-模的典范同构
$$
\Hom_B(M \otimes_A N, N') = \Hom_A(M, \Hom_B(N, N'))
$$
。若 $A,B,M,N,N'$ 带有彼此相容的分次，则有分次 $R$-模的典范同构
$$
\Hom_{\text{Mod}_B^{gr}}(M \otimes_A N, N') =
\Hom_{\text{Mod}_A^{gr}}(M, \Hom_{\text{Mod}_B^{gr}}(N, N'))
$$
。若 $A,B,M,N,N'$ 带有彼此相容的微分分次结构，则有 $R$-模复形的典范同构
$$
\Hom_{\text{Mod}^{dg}_{(B, \text{d})}}(M \otimes_A N, N') =
\Hom_{\text{Mod}^{dg}_{(A, \text{d})}}(M,
\Hom_{\text{Mod}^{dg}_{(B, \text{d})}}(N, N'))
$$
。
\end{lemma}

\begin{proof}
略。提示：在非分次情形下，把两边都解释为右侧 $B$-线性的 $A$-双线性映射
$\psi:M\times N\to N'$。在（微分）分次情形下，使用《代数进阶》引理
\ref{more-algebra-lemma-compose} 的同构，并验证它与模结构相容。或者，
使用引理 \ref{dga-lemma-characterize-hom} 的同构，并证明它与 $B$-模结构相容。
\end{proof}







\section{导出 Hom}
\label{dga-section-restriction}

\noindent
本节与《代数进阶》第 \ref{more-algebra-section-RHom} 节对应。

\medskip\noindent
设 $R$ 是环，$(A,\text{d})$ 和 $(B,\text{d})$ 是 $R$ 上的微分分次代数，
$N$ 是微分分次 $(A,B)$-双模。考虑第 \ref{dga-section-bimodules-hom} 节的函子
\begin{equation}
\label{dga-equation-restriction}
\Hom_{\text{Mod}^{dg}_{(B, \text{d})}}(N, -) :
\text{Mod}_{(B, \text{d})}
\longrightarrow
\text{Mod}_{(A, \text{d})}
\end{equation}
。

\begin{lemma}
\label{dga-lemma-restriction-homotopy}
函子 (\ref{dga-equation-restriction}) 定义三角范畴之间的正合函子
$K(\text{Mod}_{(B, \text{d})})\to K(\text{Mod}_{(A, \text{d})})$。
\end{lemma}

\begin{proof}
通过引理 \ref{dga-lemma-hom} 和评注 \ref{dga-remark-shift-hom-no-sign}，
这由引理 \ref{dga-lemma-functor-between-ABC} 的一般原理得出。
\end{proof}

\noindent
回忆，有三角范畴之间的正合函子
$$
\Hom_{\text{Mod}^{dg}_{(B, \text{d})}}(N, -) :
K(\text{Mod}_{(B, \text{d})}) \to K(\text{Mod}_{(A, \text{d})})
$$
；见引理 \ref{dga-lemma-restriction-homotopy}。考虑图表
$$
\xymatrix{
K(\text{Mod}_{(B, \text{d})}) \ar[d] \ar[rr]_{\text{见上文}} \ar[rrd]_F & &
K(\text{Mod}_{(A, \text{d})}) \ar[d] \\
D(B, \text{d}) \ar@{..>}[rr] & &
D(A, \text{d})
}
$$
希望把虚线箭头构造为复合 $F$ 的{\it 右导出函子}。
（{\it 警告}：在多数有意义的情形中，图表并不交换。）具体地，在《导出范畴》
第 \ref{derived-section-derived-functors} 节的一般框架中，希望计算 $F$
关于 $K(\text{Mod}_{(A, \text{d})})$ 中拟同构乘法系的右导出函子。

\begin{lemma}
\label{dga-lemma-derived-restriction}
在上述情形中，$F$ 的右导出函子存在。记为
$R\Hom(N,-):D(B,\text{d})\to D(A,\text{d})$。
\end{lemma}

\begin{proof}
使用《导出范畴》引理 \ref{derived-lemma-find-existence-computes} 来证明。
取具有性质 (I) 的对象作为对象族 $\mathcal{I}$。性质 (1) 已在引理
\ref{dga-lemma-right-resolution} 中证明。性质 (2) 成立，因为若
$s:I\to I'$ 是具有性质 (I) 的模之间的拟同构，则由引理
\ref{dga-lemma-hom-derived}，$s$ 是同伦等价。
\end{proof}

\begin{lemma}
\label{dga-lemma-functoriality-derived-restriction}
设 $R$ 是环，$(A,\text{d})$ 和 $(B,\text{d})$ 是微分分次 $R$-代数，
$f:N\to N'$ 是微分分次 $(A,B)$-双模同态。那么 $f$ 诱导函子态射
$$
- \circ f : R\Hom(N', -) \longrightarrow R\Hom(N, -)
$$
。若 $f$ 是拟同构，则 $f\circ-$ 是函子同构。
\end{lemma}

\begin{proof}
记 $\mathcal{B}=\text{Mod}^{dg}_{(B, \text{d})}$ 为微分分次 $B$-模的
微分分次范畴；见例 \ref{dga-example-dgm-dg-cat}。设 $I$ 是具有性质 (I) 的
微分分次 $B$-模。那么
$f\circ-:\Hom_\mathcal{B}(N',I)\to\Hom_\mathcal{B}(N,I)$ 是微分分次
$A$-模之间的映射，并且关于 $I$ 具有函子性。由于函子
$R\Hom(N',-)$ 和 $R\Hom(N,-)$ 通过把 $\Hom_\mathcal{B}$ 应用于具有
性质 (I) 的对象来计算（引理 \ref{dga-lemma-derived-restriction}），便得到所示的
函子变换。

\medskip\noindent
假设 $f$ 是拟同构。设 $F_\bullet$ 是 $I$ 上的给定滤过。由于
$I=\lim I/F_pI$，有
$\Hom_\mathcal{B}(N', I) = \lim \Hom_\mathcal{B}(N', I/F_pI)$，并且
$\Hom_\mathcal{B}(N,I)=\lim\Hom_\mathcal{B}(N,I/F_pI)$。由于系统
$I/F_pI$ 的过渡映射作为分次模映射是分裂的，系统
$\Hom_\mathcal{B}(N',I/F_pI)$ 和 $\Hom_\mathcal{B}(N,I/F_pI)$ 的过渡
映射都是满射。因此，把 $\Hom_\mathcal{B}(N',I)$（分别地，
$\Hom_\mathcal{B}(N,I)$）看作阿贝尔群复形，它计算复形系统
$\Hom_\mathcal{B}(N',I/F_pI)$（分别地，
$\Hom_\mathcal{B}(N,I/F_pI)$）的 $R\lim$。见《代数进阶》引理
\ref{more-algebra-lemma-compute-Rlim}。因此只需证明每个映射
$$
\Hom_\mathcal{B}(N', I/F_pI) \to \Hom_\mathcal{B}(N, I/F_pI)
$$
都是拟同构。由于满射 $I/F_{p+1}I\to I/F_pI$ 作为分次 $B$-模映射是
分裂的，故
$$
0 \to \Hom_\mathcal{B}(N', F_pI/F_{p + 1}I) \to
\Hom_\mathcal{B}(N', I/F_{p + 1}I) \to
\Hom_\mathcal{B}(N', I/F_pI) \to 0
$$
是微分分次 $A$-模的短正合列。对 $N$ 有类似的列，并且 $f$ 诱导短正合列
之间的映射。因此，对 $p$ 归纳（从 $p=0$ 开始，此时 $I/F_0I=0$），
得出只需证明映射
$\Hom_\mathcal{B}(N', F_pI/F_{p + 1}I) \to \Hom_\mathcal{B}(N, F_pI/F_{p + 1}I)$
是拟同构。由于 $F_pI/F_{p+1}I$ 是若干个 $A^\vee$ 的移位的直积，只需证明
$\Hom_\mathcal{B}(N', B^\vee[k]) \to \Hom_\mathcal{B}(N, B^\vee[k])$
是拟同构。由引理 \ref{dga-lemma-hom-into-shift-dual-free}，只需证明
$(N')^\vee\to N^\vee$ 是拟同构。这成立，因为 $f$ 是拟同构，且
$(\ )^\vee$ 是正合函子。
\end{proof}

\begin{lemma}
\label{dga-lemma-derived-restriction-exts}
设 $(A,\text{d})$ 和 $(B,\text{d})$ 是环 $R$ 上的微分分次代数，$N$
是微分分次 $(A,B)$-双模。那么对每个 $n\in\mathbf{Z}$，有 $R$-模同构
$$
H^n(R\Hom(N, M)) = \Ext^n_{D(B, \text{d})}(N, M)
$$
，它关于 $M$ 具有函子性。关于引理
\ref{dga-lemma-functoriality-derived-restriction} 中所述的运算，它也关于 $N$
具有函子性。
\end{lemma}

\begin{proof}
在引理 \ref{dga-lemma-derived-restriction} 的证明中已经看到
$$
R\Hom(N, M) = 
\Hom_{\text{Mod}^{dg}_{(B, \text{d})}}(N, I)
$$
作为微分分次 $A$-模成立，其中 $M\to I$ 是从 $M$ 到具有性质 (I) 的微分
分次 $B$-模的拟同构。因此，由引理 \ref{dga-lemma-hom-derived}，该复形具有
正确的上同调模。这些等同的函子性讨论从略。
\end{proof}

\begin{lemma}
\label{dga-lemma-compute-derived-restriction}
设 $R$ 是环，$(A,\text{d})$ 和 $(B,\text{d})$ 是微分分次 $R$-代数，
$N$ 是微分分次 $(A,B)$-双模。若
$\Hom_{D(B, \text{d})}(N, N') = \Hom_{K(\text{Mod}_{(B, \text{d})})}(N, N')$
对所有 $N'\in K(B,\text{d})$ 成立（例如，若 $N$ 作为微分分次 $B$-模
具有性质 (P)），则在 $D(B,\text{d})$ 中关于 $M$ 函子地有
$$
R\Hom(N, M) = \Hom_{\text{Mod}^{dg}_{(B, \text{d})}}(N, M)
$$
。
\end{lemma}

\begin{proof}
由构造（引理 \ref{dga-lemma-derived-restriction}），为求 $R\Hom(N,M)$，
选择拟同构 $M\to I$，其中 $I$ 是具有性质 (I) 的微分分次 $B$-模，并置
$R\Hom(N, M) = \Hom_{\text{Mod}^{dg}_{(B, \text{d})}}(N, I)$.
由假设，$M\to I$ 诱导的映射
$$
\Hom_{\text{Mod}^{dg}_{(B, \text{d})}}(N, M) \longrightarrow
\Hom_{\text{Mod}^{dg}_{(B, \text{d})}}(N, I)
$$
是拟同构；见例 \ref{dga-example-dgm-dg-cat} 中的讨论。这证明了本引理。
若 $N$ 作为 $B$-模具有性质 (P)，则由引理 \ref{dga-lemma-hom-derived}，
假设成立。
\end{proof}




\section{导出 Hom 的一种变体}
\label{dga-section-variant}

\noindent
设 $\mathcal{A}$ 是阿贝尔范畴。考虑 $\mathcal{A}$ 的复形所组成的微分分次
范畴 $\text{Comp}^{dg}(\mathcal{A})$；见例
\ref{dga-example-category-complexes}。设 $K^\bullet$ 是 $\mathcal{A}$ 的复形。置
$$
(E, \text{d}) = \Hom_{\text{Comp}^{dg}(\mathcal{A})}(K^\bullet, K^\bullet)
$$
，并考虑引理 \ref{dga-lemma-construction} 中微分分次范畴之间的函子
$$
\text{Comp}^{dg}(\mathcal{A}) \longrightarrow \text{Mod}^{dg}_{(E, \text{d})},
\quad
X^\bullet
\longmapsto
\Hom_{\text{Comp}^{dg}(\mathcal{A})}(K^\bullet, X^\bullet)
$$
。

\begin{lemma}
\label{dga-lemma-existence-of-derived}
在上述情形中，若
$\Hom(K^\bullet,-):K(\mathcal{A})\to D(\textit{Ab})$ 的右导出函子
$R\Hom(K^\bullet,-)$ 在 $D(\mathcal{A})$ 上处处有定义，则得到三角范畴之间
的典范正合函子
$$
R\Hom(K^\bullet, -) : D(\mathcal{A}) \longrightarrow D(E, \text{d})
$$
；取相应的阿贝尔群复形后，它化为通常的函子。
\end{lemma}

\begin{proof}
注意，由引理 \ref{dga-lemma-construction}，有相应函子
$K(\mathcal{A})\to K(\text{Mod}_{(E, \text{d})})$。断言它是三角范畴之间
的正合函子。确实，设 $f:A^\bullet\to B^\bullet$ 是 $\mathcal{A}$ 的
复形映射。计算表明
$$
\Hom_{\text{Comp}^{dg}(\mathcal{A})}(K^\bullet, C(f)^\bullet)
=
C\left(
\Hom_{\text{Comp}^{dg}(\mathcal{A})}(K^\bullet, A^\bullet) \to
\Hom_{\text{Comp}^{dg}(\mathcal{A})}(K^\bullet, B^\bullet)
\right)
$$
，其中右侧是本章前文在 $\text{Mod}_{(E, \text{d})}$ 中定义的锥。这表明
该函子与锥相容，因而与有区别三角相容。设 $X^\bullet$ 是
$K(\mathcal{A})$ 的对象。考虑拟同构 $s:X^\bullet\to Y^\bullet$ 的范畴。
已知函子
$(s : X^\bullet \to Y^\bullet) \mapsto \Hom_\mathcal{A}(K^\bullet, Y^\bullet)$
在 $D(\textit{Ab})$ 中看时本质上是常值的。但由于遗忘函子
$D(E,\text{d})\to D(\textit{Ab})$ 与取上同调相容，在
$D(E,\text{d})$ 中也有同样结论。这证明了本引理。
\end{proof}

\noindent
{\bf 警告：}虽然本引理按其表述成立，并可能按此形式发挥作用，但除非对所有
$n\in\mathbf{Z}$ 都有
$H^n(E) = \Ext^n_{D(\mathcal{A})}(K^\bullet, K^\bullet)$
，否则微分代数 $E$ 不是“正确”的微分代数。





\section{导出张量积}
\label{dga-section-base-change}

\noindent
本节与《代数进阶》第
\ref{more-algebra-section-derived-base-change} 节对应。

\medskip\noindent
设 $R$ 是环，$(A,\text{d})$ 和 $(B,\text{d})$ 是 $R$ 上的微分分次代数，
$N$ 是微分分次 $(A,B)$-双模。考虑第 \ref{dga-section-bimodules-tensor} 节
定义的函子
\begin{equation}
\label{dga-equation-bc}
\text{Mod}_{(A, \text{d})}
\longrightarrow
\text{Mod}_{(B, \text{d})},\quad
M \longmapsto M \otimes_A N
\end{equation}
。

\begin{lemma}
\label{dga-lemma-bc-homotopy}
函子 (\ref{dga-equation-bc}) 定义三角范畴之间的正合函子
$K(\text{Mod}_{(A, \text{d})})\to K(\text{Mod}_{(B, \text{d})})$。
\end{lemma}

\begin{proof}
通过引理 \ref{dga-lemma-tensor} 和评注 \ref{dga-remark-shift-tensor-no-sign}，
这由引理 \ref{dga-lemma-functor-between-ABC} 的一般原理得出。
\end{proof}

\noindent
此时可以考虑图表
$$
\xymatrix{
K(\text{Mod}_{(A, \text{d})}) \ar[d] \ar[rr]_{- \otimes_A N} \ar[rrd]_F & &
K(\text{Mod}_{(B, \text{d})}) \ar[d] \\
D(A, \text{d}) \ar@{..>}[rr] & &
D(B, \text{d})
}
$$
下文构造的虚线箭头将是复合 $F$ 的{\it 左导出函子}。
（{\it 警告}：该图表并不交换。）具体地，在《导出范畴》第
\ref{derived-section-derived-functors} 节的一般框架中，希望计算 $F$
关于 $K(\text{Mod}_{(A, \text{d})})$ 中拟同构乘法系的左导出函子。

\begin{lemma}
\label{dga-lemma-derived-bc}
在上述情形中，$F$ 的左导出函子存在。记为
$-\otimes_A^\mathbf{L}N:D(A,\text{d})\to D(B,\text{d})$。
\end{lemma}

\begin{proof}
使用《导出范畴》引理 \ref{derived-lemma-find-existence-computes} 来证明。
取具有性质 (P) 的对象作为对象族 $\mathcal{P}$。性质 (1) 已在引理
\ref{dga-lemma-resolve} 中证明。性质 (2) 成立，因为若 $s:P\to P'$ 是具有
性质 (P) 的模之间的拟同构，则由引理 \ref{dga-lemma-hom-derived}，$s$
是同伦等价。
\end{proof}

\begin{lemma}
\label{dga-lemma-functoriality-bc}
设 $R$ 是环，$(A,\text{d})$ 和 $(B,\text{d})$ 是微分分次 $R$-代数，
$f:N\to N'$ 是微分分次 $(A,B)$-双模同态。那么 $f$ 诱导函子态射
$$
1\otimes f :
- \otimes_A^\mathbf{L} N
\longrightarrow
- \otimes_A^\mathbf{L} N'
$$
。若 $f$ 是拟同构，则 $1\otimes f$ 是函子同构。
\end{lemma}

\begin{proof}
设 $M$ 是具有性质 (P) 的微分分次 $A$-模。那么
$1\otimes f:M\otimes_A N\to M\otimes_A N'$ 是微分分次 $B$-模之间的映射，
并且关于 $M$ 具有函子性。由于函子 $-\otimes_A^\mathbf{L}N$ 和
$-\otimes_A^\mathbf{L}N'$ 通过在具有性质 (P) 的对象上取张量来计算
（引理 \ref{dga-lemma-derived-bc}），便得到所示的函子变换。

\medskip\noindent
假设 $f$ 是拟同构。设 $F_\bullet$ 是 $M$ 上的给定滤过。注意
$M \otimes_A N = \colim F_i(M) \otimes_A N$，并且
$M \otimes_A N' = \colim F_i(M) \otimes_A N'$。因此只需证明
$F_n(M) \otimes_A N \to F_n(M) \otimes_A N'$
是拟同构（滤过余极限是正合的；见《代数》引理
\ref{algebra-lemma-directed-colimit-exact}）。由于包含映射
$F_n(M)\to F_{n+1}(M)$ 作为分次 $A$-模映射是分裂的，故
$$
0 \to F_n(M) \otimes_A N \to F_{n + 1}(M) \otimes_A N \to
F_{n + 1}(M)/F_n(M) \otimes_A N \to 0
$$
是微分分次 $B$-模的短正合列。对 $N'$ 有类似的列，并且 $f$ 诱导短正合列
之间的映射。因此，对 $n$ 归纳（从 $n=-1$ 开始，此时
$F_{-1}(M)=0$），得出只需证明映射
$F_{n+1}(M)/F_n(M)\otimes_A N\to F_{n+1}(M)/F_n(M)\otimes_A N'$
是拟同构。这成立，因为 $F_{n+1}(M)/F_n(M)$ 是若干个 $A$ 的移位的直和，
而对 $A[k]$ 结论成立，原因是 $f:N\to N'$ 是拟同构。
\end{proof}

\begin{lemma}
\label{dga-lemma-compute-bc}
设 $R$ 是环，$(A,\text{d})$ 和 $(B,\text{d})$ 是微分分次 $R$-代数。
设 $N$ 是微分分次 $(A,B)$-双模，并且作为左微分分次 $A$-模具有性质 (P)。
那么对所有微分分次 $A$-模 $M$，$M\otimes_A^\mathbf{L}N$ 由
$M\otimes_A N$ 计算。
\end{lemma}

\begin{proof}
设 $f:M\to M'$ 是微分分次 $A$-模之间的拟同构。断言
$f\otimes\text{id}:M\otimes_A N\to M'\otimes_A N$ 是拟同构。若断言成立，
则由引理 \ref{dga-lemma-derived-bc} 证明中导出张量积的构造，得到所需结论。
映射 $f\otimes\text{id}$ 的构造只依赖于 $N$ 上的左微分分次 $A$-模结构。
此外，
$M\otimes_A N=N\otimes_{A^{opp}}M=N\otimes_{A^{opp}}^\mathbf{L}M$，
因为 $N$ 作为微分分次 $A^{opp}$-模具有性质 (P)。因此，断言由引理
\ref{dga-lemma-functoriality-bc} 得出。
\end{proof}

\begin{lemma}
\label{dga-lemma-tensor-hom-adjoint}
设 $R$ 是环，$(A,\text{d})$ 和 $(B,\text{d})$ 是微分分次 $R$-代数，
$N$ 是微分分次 $(A,B)$-双模。那么引理 \ref{dga-lemma-derived-bc} 中的函子
$$
- \otimes_A^\mathbf{L} N : D(A, \text{d}) \longrightarrow D(B, \text{d})
$$
是引理 \ref{dga-lemma-derived-restriction} 中函子
$$
R\Hom(N, -) : D(B, \text{d}) \longrightarrow D(A, \text{d})
$$
的左伴随。
\end{lemma}

\begin{proof}
这由《导出范畴》引理
\ref{derived-lemma-pre-derived-adjoint-functors-general}，以及引理
\ref{dga-lemma-tensor-hom-adjunction} 中
$-\otimes_A N$ 与 $\Hom_{\text{Mod}^{dg}_{(B, \text{d})}}(N,-)$
互为伴随这一事实得出。
\end{proof}

\begin{example}
\label{dga-example-map-hom-tensor}
设 $R$ 是环，$(A,\text{d})\to(B,\text{d})$ 是微分分次 $R$-代数同态。
可以把 $B$ 看作微分分次 $(A,B)$-双模，并得到函子
$$
- \otimes_A B : D(A, \text{d}) \longrightarrow D(B, \text{d})
$$
。由引理 \ref{dga-lemma-tensor-hom-adjoint}，它的左伴随是函子
$R\Hom(B,-)$。对微分分次 $B$-模 $N$，以 $N_A$ 表示通过 $A\to B$
限制得到的微分分次 $A$-模。那么显然有关于 $B$-模 $N$ 函子的典范同构
$$
\Hom_{\text{Mod}^{dg}_{(B, \text{d})}}(B, N) \longrightarrow N_A,\quad
f \longmapsto f(1)
$$
。因此，$R\Hom(B,-)$ 是限制函子，并得到
$$
\Hom_{D(A, \text{d})}(M, N_A) =
\Hom_{D(B, \text{d})}(M \otimes^\mathbf{L}_A B, N)
$$
；它对 $M$ 与 $N$ 双函子地成立，正如交换环的情形。最后注意，限制也是张量函子，
因为 $N_A = N \otimes_B {}_BB_A = N \otimes_B^\mathbf{L} {}_BB_A$，
其中 ${}_BB_A$ 表示把 $B$ 看作微分分次 $(B,A)$-双模。
\end{example}

\begin{lemma}
\label{dga-lemma-tensor-with-compact-fully-faithful}
沿用引理 \ref{dga-lemma-tensor-hom-adjoint} 的记号与假设。假设
\begin{enumerate}
\item $N$ 在 $D(B,\text{d})$ 中定义一个紧对象；
\item 对每个 $k\in\mathbf{Z}$，映射
$H^k(A) \to \Hom_{D(B, \text{d})}(N, N[k])$ 都是同构。
\end{enumerate}
则函子 $-\otimes_A^\mathbf{L}N$ 是全忠实的。
\end{lemma}

\begin{proof}
由引理 \ref{dga-lemma-tensor-hom-adjoint}，我们的函子有由
$R\Hom(N,-)$ 给出的左伴随。由《范畴》引理
\ref{categories-lemma-adjoint-fully-faithful}，只需证明：对任意微分分次
$A$-模 $M$，映射
$$
M \longrightarrow R\Hom(N, M \otimes_A^\mathbf{L} N)
$$
在 $D(A,\text{d})$ 中是同构。为此，只需证明
$$
H^n(M) \longrightarrow
\text{Ext}^n_{D(B, \text{d})}(N, M \otimes_A^\mathbf{L} N)
$$
是同构，见引理 \ref{dga-lemma-derived-restriction-exts}。由于 $N$ 是紧对象，
右端与直和交换。因此由评注 \ref{dga-remark-P-resolution}，只需在
$M=A[k]$ 时证明此映射为同构。由评注
\ref{dga-remark-shift-tensor-no-sign}，
$(A[k]\otimes_A^\mathbf{L}N)=N[k]$；而对 $N$ 的假设 (2)
正是说结论在这些情形下成立。
\end{proof}

\begin{lemma}
\label{dga-lemma-base-change-K-flat}
设 $R\to R'$ 是环映射，$(A,\text{d})$ 是微分分次 $R$-代数，
而 $(A',\text{d})$ 是其基变换，即 $A'=A\otimes_RR'$。
若 $A$ 作为 $R$-模复形是 K-平坦的，则
\begin{enumerate}
\item $- \otimes_A^\mathbf{L} A' : D(A, \text{d}) \to D(A', \text{d})$
等于函子
$$
K(A, \text{d}) \longrightarrow K(A', \text{d}),\quad
M \longmapsto M \otimes_R R'
$$
的右导出函子；
\item 图表
$$
\xymatrix{
D(A, \text{d}) \ar[r]_{- \otimes_A^\mathbf{L} A'} \ar[d]_{\text{限制}} &
D(A', \text{d}) \ar[d]^{\text{限制}} \\
D(R) \ar[r]^{- \otimes_R^\mathbf{L} R'} & D(R')
}
$$
交换；
\item 若 $M$ 作为 $R$-模复形是 K-平坦的，则微分分次 $A'$-模
$M\otimes_RR'$ 表示 $M\otimes_A^\mathbf{L}A'$。
\end{enumerate}
\end{lemma}

\begin{proof}
对任意微分分次 $A$-模 $M$，都有典范映射
$$
c_M : M \otimes_R R' \longrightarrow M \otimes_A A'
$$
。设 $P$ 是具有性质 (P) 的微分分次 $A$-模。我们断言 $c_P$ 是同构，
并且 $P$ 作为 $R$-模复形是 K-平坦的。利用引理
\ref{dga-lemma-derived-bc} 和《代数进阶》第
\ref{more-algebra-section-derived-tensor-product} 节中导出张量积的定义，
这将通过形式论证证明引理所述全部结论。

\medskip\noindent
设 $F_\bullet$ 是表明 $P$ 具有性质 (P) 的滤过。注意 $c_A$ 是同构，
且依假设 $A$ 作为 $R$-模复形是 K-平坦的。因此，对 $A$ 的平移的直和
同样成立（处理直和时也可使用《代数进阶》引理
\ref{more-algebra-lemma-colimit-K-flat}）。于是，这对复形
$F_{p+1}P/F_pP$ 成立。由于短正合列
$$
0 \to F_pP \to F_{p + 1}P \to F_{p + 1}P/F_pP \to 0
$$
作为分次模列是分裂正合的，可归纳得出：对每个 $p$，$c_{F_pP}$ 都是同构，
且 $F_pP$ 作为 $R$-模复形是 K-平坦的（使用《代数进阶》引理
\ref{more-algebra-lemma-K-flat-two-out-of-three}）。最后，由
$P=\colim F_pP$ 可知 $c_P$ 是同构，且 $P$ 作为 $R$-模复形是
K-平坦的（使用《代数进阶》引理
\ref{more-algebra-lemma-colimit-K-flat}）。
\end{proof}

\begin{lemma}
\label{dga-lemma-RHom-is-tensor}
设 $R$ 是环，$(A,\text{d})$ 与 $(B,\text{d})$ 是微分分次
$R$-代数，$T$ 是微分分次 $(A,B)$-双模。假设
\begin{enumerate}
\item $T$ 在 $D(B,\text{d})$ 中定义一个紧对象；
\item $S = \Hom_{\text{Mod}^{dg}_{(B, \text{d})}}(T, B)$
在 $D(A,\text{d})$ 中表示 $R\Hom(T,B)$。
\end{enumerate}
则 $S$ 具有微分分次 $(B,A)$-双模结构，并有同构
$$
N \otimes_B^\mathbf{L} S \longrightarrow R\Hom(T, N)
$$
，它关于 $D(B,\text{d})$ 中的 $N$ 是函子的。
\end{lemma}

\begin{proof}
记 $\mathcal{B}=\text{Mod}^{dg}_{(B,\text{d})}$。$S$ 上的右
$A$-模结构来自映射 $A\to\Hom_\mathcal{B}(T,T)$ 以及例
\ref{dga-example-dgm-dg-cat} 中定义的复合
$\Hom_\mathcal{B}(T, B) \otimes \Hom_\mathcal{B}(T, T)
\to \Hom_\mathcal{B}(T, B)$。再次使用这一乘法，可得映射
$$
c_N :
N \otimes_B S = \Hom_\mathcal{B}(B, N) \otimes_B \Hom_\mathcal{B}(T, B)
\longrightarrow
\Hom_\mathcal{B}(T, N)
$$
，它关于 $N$ 是函子的。给定 $N$，可选取拟同构
$P\to N\to I$，其中 $P$（相应地 $I$）是具有性质 (P)（相应地 (I)）
的微分分次 $B$-模。于是利用 $c_N$，在分别表示
$S\otimes_B^\mathbf{L}N$ 与 $R\Hom(T,N)$ 的对象之间得到映射
$P\otimes_BS\to\Hom_\mathcal{B}(T,I)$。显然，这定义了引理所述的
函子变换 $c$。

\medskip\noindent
为证明 $c$ 是函子同构，可假设 $N$ 是具有性质 (P) 的微分分次
$B$-模。由于 $T$ 在 $D(B,\text{d})$ 中定义紧对象，且箭头两端都定义
三角范畴的正合函子，利用引理 \ref{dga-lemma-property-P-sequence}，
可约化到 $N$ 具有有限滤过、其分次片是若干 $B[k]$ 的直和的情形。
再次利用箭头两端都是三角范畴的正合函子以及 $T$ 的紧性，可进一步
约化到 $N=B[k]$ 的情形。假设 (2) 恰好断言此时 $c$ 是同构。
\end{proof}






\section{导出张量积的复合}
\label{dga-section-compose-tensor-functors}

\noindent
我们建议读者略过本节。

\medskip\noindent
设 $R$ 是环，$(A,\text{d})$、$(B,\text{d})$ 和 $(C,\text{d})$
是微分分次 $R$-代数。设 $N$ 是微分分次 $(A,B)$-双模，$N'$ 是微分分次
$(B,C)$-模。以 $N_B$ 表示把双模 $N$ 看作微分分次 $B$-模
（忘掉 $A$-结构）所得的模。存在典范映射
\begin{equation}
\label{dga-equation-plain-versus-derived}
N_B \otimes_B^\mathbf{L} N'
\longrightarrow
(N \otimes_B N')_C
\end{equation}
于 $D(C,\text{d})$ 中。这里，$(N\otimes_BN')_C$ 表示把
$(A,C)$-双模 $N\otimes_BN'$ 看作微分分次 $C$-模所得的模。
事实上，此映射来自如下事实：导出张量积总有到普通张量积的映射
（因为它是左导出函子）。

\begin{lemma}
\label{dga-lemma-compose-tensor-functors-general}
设 $R$ 是环，$(A,\text{d})$、$(B,\text{d})$ 和 $(C,\text{d})$
是微分分次 $R$-代数。设 $N$ 是微分分次 $(A,B)$-双模，$N'$ 是微分分次
$(B,C)$-模。假设 (\ref{dga-equation-plain-versus-derived}) 是同构。
则复合
$$
\xymatrix{
D(A, \text{d}) \ar[rr]^{- \otimes_A^\mathbf{L} N} & &
D(B, \text{d}) \ar[rr]^{- \otimes_B^\mathbf{L} N'} & &
D(C, \text{d})
}
$$
同构于 $-\otimes_A^\mathbf{L}N''$，其中 $N''=N\otimes_BN'$
被看作 $(A,C)$-双模。
\end{lemma}

\begin{proof}
我们定义函子变换
$$
(- \otimes_A^\mathbf{L} N) \otimes_B^\mathbf{L} N'
\longrightarrow
- \otimes_A^\mathbf{L} N''
$$
。为此，设 $M$ 是具有性质 (P) 的微分分次 $A$-模。按照引理
\ref{dga-lemma-derived-bc} 的证明中函子 $-\otimes_A^\mathbf{L}N''$
的构造，普通张量积 $M\otimes_AN''$ 在 $D(C,\text{d})$ 中表示
$M\otimes_A^\mathbf{L}N''$。于是可写成
$$
M \otimes_A N'' =
M \otimes_A (N \otimes_B N') =
(M \otimes_A N) \otimes_B N'
$$
。模 $M\otimes_AN$ 在 $D(B,\text{d})$ 中表示
$M\otimes_A^\mathbf{L}N$。选取拟同构 $Q\to M\otimes_AN$，
其中 $Q$ 是具有性质 (P) 的微分分次 $B$-模。则 $Q\otimes_BN'$
在 $D(C,\text{d})$ 中表示
$(M\otimes_A^\mathbf{L}N)\otimes_B^\mathbf{L}N'$。
因此，可由下式定义所需映射：
$$
(M \otimes_A^\mathbf{L} N) \otimes_B^\mathbf{L} N' =
Q \otimes_B N' \to
M \otimes_A N \otimes_B N' =
M \otimes_A^\mathbf{L} N''
$$
。此映射的构造关于 $M$ 是函子的，并与有区别三角及直和相容；细节从略。
考虑 $D(A,\text{d})$ 的对象 $M$ 的性质 $T$，其含义是此映射为同构。则
\begin{enumerate}
\item 若每个 $M_i$ 都满足 $T$，则 $\bigoplus M_i$ 满足 $T$；
\item 对一个有区别三角，若三个对象中的两个满足 $T$，则第三个也满足；
\item $A[k]$ 满足 $T$，因为此时所得映射是
(\ref{dga-equation-plain-versus-derived}) 的一个平移，而我们已经假设后者是同构。
\end{enumerate}
因此由评注 \ref{dga-remark-P-resolution}，性质 $T$ 恒成立，证明完毕。
\end{proof}

\noindent
设 $R$ 是环，$(A,\text{d})$、$(B,\text{d})$ 和 $(C,\text{d})$
是微分分次 $R$-代数。暂以 $(A\otimes_RB)_B$ 表示把微分分次代数
$A\otimes_RB$ 看作（右）微分分次 $B$-模所得的模；以
${}_B(B\otimes_RC)_C$ 表示把微分分次代数 $B\otimes_RC$
看作微分分次 $(B,C)$-双模所得的双模。于是存在典范映射
\begin{equation}
\label{dga-equation-plain-versus-derived-algebras}
(A \otimes_R B)_B \otimes_B^\mathbf{L} {}_B(B \otimes_R C)_C
\longrightarrow
(A \otimes_R B \otimes_R C)_C
\end{equation}
于 $D(C,\text{d})$ 中，其中 $(A\otimes_RB\otimes_RC)_C$
表示把微分分次 $R$-代数 $A\otimes_RB\otimes_RC$ 看作（右）
微分分次 $C$-模所得的模。事实上，此映射来自等同
$$
(A \otimes_R B)_B \otimes_B {}_B(B \otimes_R C)_C =
(A \otimes_R B \otimes_R C)_C
$$
以及导出张量积总有到普通张量积的映射这一事实
（因为它是左导出函子）。

\begin{lemma}
\label{dga-lemma-compose-tensor-functors-general-algebra}
设 $R$ 是环，$(A,\text{d})$、$(B,\text{d})$ 和 $(C,\text{d})$
是微分分次 $R$-代数。假设
(\ref{dga-equation-plain-versus-derived-algebras}) 是同构。设 $N$ 是微分分次
$(A,B)$-双模，$N'$ 是微分分次 $(B,C)$-双模。则复合
$$
\xymatrix{
D(A, \text{d}) \ar[rr]^{- \otimes_A^\mathbf{L} N} & &
D(B, \text{d}) \ar[rr]^{- \otimes_B^\mathbf{L} N'} & &
D(C, \text{d})
}
$$
同构于 $-\otimes_A^\mathbf{L}N''$，其中微分分次 $(A,C)$-双模
$N''$ 将在证明中描述。
\end{lemma}

\begin{proof}
由引理 \ref{dga-lemma-functoriality-bc}，可把 $N$ 与 $N'$ 替换为与之
拟同构的双模。因此可假设 $N$（相应地 $N'$）作为微分分次
$(A,B)$-双模（相应地 $(B,C)$-双模）具有性质 (P)，见引理
\ref{dga-lemma-bimodule-resolve}。我们断言，对 $(A,C)$-双模
$N''=N\otimes_BN'$，引理成立。为证明这一点，只需证明
$$
N_B \otimes_B^\mathbf{L} N' \longrightarrow (N \otimes_B N')_C
$$
在 $D(C,\text{d})$ 中是同构，见引理
\ref{dga-lemma-compose-tensor-functors-general}。

\medskip\noindent
设 $F_\bullet$ 是双模性质 (P) 中 $N$ 上的滤过。由引理
\ref{dga-lemma-bimodule-property-P-sequence}，存在微分分次
$(A,B)$-双模的短正合列
$$
0 \to
\bigoplus\nolimits F_iN \to
\bigoplus\nolimits F_iN \to N \to 0
$$
，它作为分次 $(A,B)$-双模列是分裂的。因而它更是微分分次 $B$-模的
容许短正合列，并给出 $D(B,\text{d})$ 中的有区别三角
$$
\bigoplus\nolimits F_iN_B \to
\bigoplus\nolimits F_iN_B \to N_B \to
\bigoplus\nolimits F_iN_B[1]
$$
。利用 $-\otimes_B^\mathbf{L}N'$ 是三角范畴的正合函子且与直和交换，
并利用 $-\otimes_BN'$ 把容许正合列变为容许正合列且与直和交换，
可将问题约化为证明：对每个 $p$，
$$
(F_pN)_B \otimes_B^\mathbf{L} N' \longrightarrow (F_pN)_B \otimes_B N'
$$
都是拟同构。对 $(A,B)$-双模的短正合列
$$
0 \to F_pN \to F_{p + 1}N \to F_{p + 1}N/F_pN \to 0
$$
重复此论证；这些列作为分次 $(A,B)$-双模列是分裂的，故可约化为对
$F_{p+1}N/F_pN$ 证明同一断言。由于这些模是
$(A\otimes_RB)_B$ 的平移的直和，又可约化为证明
$$
(A \otimes_R B)_B \otimes_B^\mathbf{L} N'
\longrightarrow
(A \otimes_R B)_B \otimes_B N'
$$
是拟同构。

\medskip\noindent
选取双模性质 (P) 中 $N'$ 上的滤过 $F_\bullet$。选取微分分次
$B$-模的拟同构 $P\to(A\otimes_RB)_B$，其中 $P$ 具有性质 (P)。
由引理 \ref{dga-lemma-derived-bc} 中的构造，$P\otimes_BN'$ 在
$D(C,\text{d})$ 中表示
$(A\otimes_RB)_B\otimes_B^\mathbf{L}N'$；因此我们必须证明
$P\otimes_BN'\to(A\otimes_RB)_B\otimes_BN'$ 是拟同构。
由于 $N'=\colim F_pN'$，只需证明
$P\otimes_BF_pN'\to(A\otimes_RB)_B\otimes_BF_pN'$
是拟同构。利用短正合列
$0 \to F_pN' \to F_{p + 1}N' \to F_{p + 1}N'/F_pN' \to 0$
作为分次 $(B,C)$-双模列是分裂的，可约化为证明：对每个 $p$，
$P \otimes_B F_{p + 1}N'/F_pN' \to
(A \otimes_R B)_B \otimes_B F_{p + 1}N'/F_pN'$
是拟同构。最后，利用 $F_{p+1}N'/F_pN'$ 是
${}_B(B\otimes_RC)_C$ 的平移的直和，可知只需证明
$$
P \otimes_B {}_B(B \otimes_R C)_C \to
(A \otimes_R B)_B \otimes_B {}_B(B \otimes_R C)_C
$$
是拟同构。由于 $P\to(A\otimes_RB)_B$ 是由满足性质 (P) 的模
给出的分辨，此微分分次 $C$-模映射在 $D(C,\text{d})$ 中表示态射
(\ref{dga-equation-plain-versus-derived-algebras})，证明完毕。
\end{proof}

\begin{lemma}
\label{dga-lemma-compose-tensor-functors}
设 $R$ 是环，$(A,\text{d})$、$(B,\text{d})$ 和 $(C,\text{d})$
是微分分次 $R$-代数。若 $C$ 作为 $R$-模复形是 K-平坦的，则
(\ref{dga-equation-plain-versus-derived-algebras}) 是同构，并且引理
\ref{dga-lemma-compose-tensor-functors-general-algebra} 的结论成立。
\end{lemma}

\begin{proof}
选取微分分次 $B$-模的拟同构 $P\to(A\otimes_RB)_B$，其中 $P$
具有性质 (P)。于是只需证明
$$
P \otimes_B (B \otimes_R C) \longrightarrow
(A \otimes_R B) \otimes_B (B \otimes_R C)
$$
是拟同构。等价地，我们考察
$$
P \otimes_R C \longrightarrow
A \otimes_R B \otimes_R C
$$
。若 $C$ 作为 $R$-模复形是 K-平坦的，则由《代数进阶》引理
\ref{more-algebra-lemma-K-flat-quasi-isomorphism}，这是拟同构。
\end{proof}





\section{导出张量积的一种变体}
\label{dga-section-variant-base-change}

\noindent
设 $(\mathcal{C},\mathcal{O})$ 是环化位点。于是有函子
$$
\text{Comp}(\mathcal{O}) \to K(\mathcal{O}) \to D(\mathcal{O})
$$
；如上所见，还有微分分次增强 $\text{Comp}^{dg}(\mathcal{O})$。
具体地说，这是例 \ref{dga-example-category-complexes} 中与阿贝尔范畴
$\textit{Mod}(\mathcal{O})$ 相关联的微分分次范畴。设 $K^\bullet$
是 $\mathcal{O}$-模复形，换言之，是 $\text{Comp}^{dg}(\mathcal{O})$
的对象。令
$$
(E, \text{d}) =
\Hom_{\text{Comp}^{dg}(\mathcal{O})}(K^\bullet, K^\bullet)
$$
。这是微分分次 $\mathbf{Z}$-代数。我们断言，在这种情形下有导出基变换
的类似物。

\begin{lemma}
\label{dga-lemma-tensor-with-complex}
在上述情形下，存在微分分次范畴的函子
$$
- \otimes_E K^\bullet :
\text{Mod}^{dg}_{(E, \text{d})}
\longrightarrow
\text{Comp}^{dg}(\mathcal{O})
$$
。此函子把 $E$ 映到 $K^\bullet$，并与直和交换。
\end{lemma}

\begin{proof}
设 $M$ 是微分分次 $E$-模。对 $\mathcal{C}$ 的每个对象 $U$，复形
$K^\bullet(U)$ 既是左微分分次 $E$-模，也是右 $\mathcal{O}(U)$-模。
两个作用交换，故得到一个双模。因此，由第 \ref{dga-section-tensor-product}
节与第 \ref{dga-section-bimodules} 节中的构造，可以形成张量积
$$
M \otimes_E K^\bullet(U)
$$
；它是微分分次 $\mathcal{O}(U)$-模，即 $\mathcal{O}(U)$-模复形。
此构造关于 $U$ 是函子的，故可层化而得到 $\mathcal{O}$-模复形，记为
$$
M \otimes_E K^\bullet
$$
。此外，对每个 $U$，由引理 \ref{dga-lemma-tensor}，此构造确定微分分次
范畴的函子
$\text{Mod}^{dg}_{(E, \text{d})} \to \text{Comp}^{dg}(\mathcal{O}(U))$
。因此，显然得到引理所述函子。
\end{proof}

\begin{lemma}
\label{dga-lemma-tensor-with-complex-homotopy}
引理 \ref{dga-lemma-tensor-with-complex} 的函子定义三角范畴的正合函子
$K(\text{Mod}_{(E, \text{d})}) \to K(\mathcal{O})$.
\end{lemma}

\begin{proof}
由引理 \ref{dga-lemma-functorial}，该函子诱导同伦范畴之间的函子。
我们必须证明 $-\otimes_EK^\bullet$ 把有区别三角变为有区别三角。
设 $0\to K\to L\to M\to0$ 是微分分次 $E$-模的容许短正合列。
设 $s:M\to L$ 是分次 $E$-模同态，并且是 $L\to M$ 的左逆。
则 $s$ 定义分次 $\mathcal{O}$-模的映射
$M\otimes_EK^\bullet\to L\otimes_EK^\bullet$
（即保持 $\mathcal{O}$-模结构与分次，但不保持微分），它是
$L\otimes_EK^\bullet\to M\otimes_EK^\bullet$ 的左逆。因此
$$
0 \to K \otimes_E K^\bullet \to L \otimes_E K^\bullet \to
M \otimes_E K^\bullet \to 0
$$
是逐项分裂的复形短正合列，也就是说，它在 $K(\mathcal{O})$ 中定义
一个有区别三角。
\end{proof}

\begin{lemma}
\label{dga-lemma-tensor-with-complex-derived}
引理 \ref{dga-lemma-tensor-with-complex-homotopy} 的函子
$K(\text{Mod}_{(E, \text{d})})\to K(\mathcal{O})$ 有定义在整个
$D(E,\text{d})$ 上的左导出版本，记为
$- \otimes_E^\mathbf{L} K^\bullet : D(E, \text{d}) \to D(\mathcal{O})$.
\end{lemma}

\begin{proof}
使用《导出范畴》引理 \ref{derived-lemma-find-existence-computes}
来证明。取具有性质 (P) 的对象为对象族 $\mathcal{P}$。
性质 (1) 已在引理 \ref{dga-lemma-resolve} 中证明。性质 (2) 成立，
因为若 $s:P\to P'$ 是具有性质 (P) 的模之间的拟同构，则由引理
\ref{dga-lemma-hom-derived}，$s$ 是同伦等价。
\end{proof}

\begin{lemma}
\label{dga-lemma-upgrade-tensor-with-complex-derived}
设 $R$ 是环，$\mathcal{C}$ 是位点，$\mathcal{O}$ 是交换
$R$-代数层，$K^\bullet$ 是 $\mathcal{O}$-模复形。引理
\ref{dga-lemma-tensor-with-complex-derived} 的函子具有如下性质：对
$D(E,\text{d})$ 中的任意 $M,N$，在 $D(R)$ 中有典范映射
$$
R\Hom(M, N)
\longrightarrow
R\Hom_\mathcal{O}(M \otimes_E^\mathbf{L} K^\bullet,
N \otimes_E^\mathbf{L} K^\bullet)
$$
；它在上同调模上给出由函子 $-\otimes_E^\mathbf{L}K^\bullet$ 诱导的映射
$$
\Ext^n_{D(E, \text{d})}(M, N) \to
\Ext^n_{D(\mathcal{O})}
(M \otimes_E^\mathbf{L} K^\bullet, N \otimes_E^\mathbf{L} K^\bullet)
$$
。
\end{lemma}

\begin{proof}
箭头右端是《位点上的上同调》第
\ref{sites-cohomology-section-global-RHom} 节引入的全局导出 Hom，
它具有正确的上同调模。对左端，把 $M$ 看作 $(R,A)$-双模，并使用第
\ref{dga-section-restriction} 节引入的导出 $\Hom$；它也具有正确的上同调模。
为证明引理，可假设 $M,N$ 都是具有性质 (P) 的微分分次 $E$-模；
由引理 \ref{dga-lemma-functoriality-derived-restriction}，这不改变箭头左端。
由引理 \ref{dga-lemma-compute-derived-restriction}，箭头左端于是变为
$\Hom_{\text{Mod}^{dg}_{(B, \text{d})}}(M, N)$.
在引理 \ref{dga-lemma-tensor-with-complex}、
\ref{dga-lemma-tensor-with-complex-homotopy} 与
\ref{dga-lemma-tensor-with-complex-derived}
中，我们已经构造微分分次范畴的函子
$$
- \otimes_E K^\bullet :
\text{Mod}^{dg}_{(E, \text{d})}
\longrightarrow
\text{Comp}^{dg}(\mathcal{O})
$$
，并已证明，在具有性质 (P) 的微分分次 $E$-模上取此函子的值即可计算
$-\otimes_E^\mathbf{L}K^\bullet$。因此得到 $R$-模复形的映射
$$
\Hom_{\text{Mod}^{dg}_{(B, \text{d})}}(M, N)
\longrightarrow
\Hom_{\text{Comp}^{dg}(\mathcal{O})}
(M \otimes_E K^\bullet, N \otimes_E K^\bullet)
$$
。对任意 $\mathcal{O}$-模复形 $\mathcal{F}^\bullet$、
$\mathcal{G}^\bullet$，存在典范映射
$$
\Hom_{\text{Comp}^{dg}(\mathcal{O})}
(\mathcal{F}^\bullet, \mathcal{G}^\bullet) =
\Gamma(\mathcal{C},
\SheafHom^\bullet(\mathcal{F}^\bullet, \mathcal{G}^\bullet))
\longrightarrow
R\Hom_\mathcal{O}(\mathcal{F}^\bullet, \mathcal{G}^\bullet).
$$
合并这些映射，即得引理所需映射。
\end{proof}

\begin{lemma}
\label{dga-lemma-tensor-with-complex-hom-adjoint}
设 $(\mathcal{C},\mathcal{O})$ 是环化位点，$K^\bullet$ 是
$\mathcal{O}$-模复形。则引理 \ref{dga-lemma-tensor-with-complex-derived}
的函子
$$
- \otimes_E^\mathbf{L} K^\bullet :
D(E, \text{d})
\longrightarrow
D(\mathcal{O})
$$
是引理 \ref{dga-lemma-existence-of-derived} 的函子
$$
R\Hom(K^\bullet, -) : D(\mathcal{O}) \longrightarrow D(E, \text{d})
$$
的左伴随。
\end{lemma}

\begin{proof}
此断言的含义是，存在关于 $M$ 与 $L^\bullet$ 双函子的等式
$$
\Hom_{D(E, \text{d})}(M, R\Hom(K^\bullet, L^\bullet)) =
\Hom_{D(\mathcal{O})}(M \otimes^\mathbf{L}_E K^\bullet, L^\bullet)
$$
。为说明这一点，可把 $M$ 替换为具有性质 (P) 的微分分次 $E$-模 $P$，
也可把 $L^\bullet$ 替换为 K-内射的 $\mathcal{O}$-模复形
$I^\bullet$。结合引理 \ref{dga-lemma-hom-derived} 与陈述中所引诸引理给出的
导出函子计算，上式化为
$$
\Hom_{K(\text{Mod}_{(E, \text{d})})}
(P, \Hom_\mathcal{B}(K^\bullet, I^\bullet)) =
\Hom_{K(\mathcal{O})}(P \otimes_E K^\bullet, I^\bullet)
$$
，其中 $\mathcal{B}=\text{Comp}^{dg}(\mathcal{O})$。存在从右端到左端、
关于 $P$ 与 $I^\bullet$ 函子的求值映射（细节从略）。选取性质 (P)
定义中 $P$ 上的滤过 $F_\bullet$。由引理
\ref{dga-lemma-property-P-sequence}，并利用等式两端在
$K(\text{Mod}_{(E,\text{d})})$ 上都是关于 $P$ 的同调函子，
可约化到把 $P$ 换成微分分次 $E$-模 $\bigoplus F_iP$ 的情形。
由于两端都把变量 $P$ 中的直和变为直积，可再约化到 $P$ 是某个
微分分次 $E$-模 $F_iP$ 的情形。每个 $F_iP$ 都有由容许单态射给出的
有限滤过，其分次片是分次投射 $E$-模，故可约化到 $P$ 是分次投射
$E$-模的情形。此时显然有
$$
\Hom_{\text{Mod}^{dg}_{(E, \text{d})}}
(P, \Hom_\mathcal{B}(K^\bullet, I^\bullet)) =
\Hom_{\text{Comp}^{dg}(\mathcal{O})}(P \otimes_E K^\bullet, I^\bullet)
$$
作为分次 $\mathbf{Z}$-模的等式（因为此断言可约化到 $P=E[k]$
这一显然的情形）。该同构与微分相容，故得结论。
\end{proof}

\begin{lemma}
\label{dga-lemma-fully-faithful-in-compact-case}
设 $(\mathcal{C},\mathcal{O})$ 是环化位点，$K^\bullet$ 是
$\mathcal{O}$-模复形。假设
\begin{enumerate}
\item $K^\bullet$ 表示 $D(\mathcal{O})$ 的一个紧对象；
\item $E = \Hom_{\text{Comp}^{dg}(\mathcal{O})}(K^\bullet, K^\bullet)$
计算 $K^\bullet$ 在 $D(\mathcal{O})$ 中的 Ext 群。
\end{enumerate}
则引理 \ref{dga-lemma-tensor-with-complex-derived} 的函子
$$
- \otimes_E^\mathbf{L} K^\bullet :
D(E, \text{d})
\longrightarrow
D(\mathcal{O})
$$
是全忠实的。
\end{lemma}

\begin{proof}
由引理 \ref{dga-lemma-tensor-with-complex-hom-adjoint}，我们的函子有由
$R\Hom(K^\bullet,-)$ 给出的左伴随。因此只需证明：对微分分次
$E$-模 $M$，映射
$$
H^0(M) \longrightarrow
\Hom_{D(\mathcal{O})}(K^\bullet, M \otimes_E^\mathbf{L} K^\bullet)
$$
是同构。可假设 $M=P$ 是具有性质 (P) 的微分分次 $E$-模。
由于 $K^\bullet$ 定义紧对象，利用引理
\ref{dga-lemma-property-P-sequence}，可约化到 $P$ 具有有限滤过且其分次片
是若干 $E[k]$ 的直和的情形。再次利用紧性，可约化到 $P=E[k]$
的情形。对 $K^\bullet$ 的假设正是说结论在这些情形下成立。
\end{proof}











\section{紧对象的刻画}
\label{dga-section-compact}

\noindent
加性范畴的紧对象定义于《导出范畴》定义
\ref{derived-definition-compact-object}。本节刻画微分分次代数的导出范畴
中的紧对象。

\begin{remark}
\label{dga-remark-source-graded-projective}
设 $(A,\text{d})$ 是微分分次代数。能否刻画满足下列条件的微分分次
$A$-模 $P$：对每个微分分次 $A$-模 $M$，都有
$$
\Hom_{K(A, \text{d})}(P, M) =  \Hom_{D(A, \text{d})}(P, M)
$$
？设 $\mathcal{D}\subset K(A,\text{d})$ 是以满足上述条件的对象 $P$
为对象的满子范畴。则 $\mathcal{D}$ 是 $K(A,\text{d})$ 的严格满饱和
三角子范畴。若 $P$ 作为分次 $A$-模是投射的，则要判定 $P$ 是否为
$\mathcal{D}$ 的对象，只需检验：每当 $M$ 非循环时，
$\Hom_{K(A,\text{d})}(P,M)=0$。然而，一般来说，仅假设 $P$
作为分次 $A$-模是投射的还不够。例如，取 $A=R=k[\epsilon]$，其中
$k$ 是域，$k[\epsilon]=k[x]/(x^2)$ 是对偶数环。令 $P$ 为如下对象：
对每个 $n\in\mathbf{Z}$ 都有 $P^n=R$，且微分由乘以 $\epsilon$
给出。于是 $\text{id}_P\in\Hom_{K(A,\text{d})}(P,P)$ 是非零元素，
但 $P$ 非循环。
\end{remark}

\begin{remark}
\label{dga-remark-graded-projective-is-compact}
设 $(A,\text{d})$ 是微分分次代数。若微分分次 $A$-模 $M$ 作为右
$A$-模由有限多个元素生成，就称 $M$ 是{\it 有限的}。若微分分次
$A$-模 $P$ 是有限分次投射的，可以问：$P$ 是否给出
$D(A,\text{d})$ 的紧对象？推测一般并非如此，但我们不知道反例。
然而，若 $P$ 还是评注 \ref{dga-remark-source-graded-projective} 的范畴
$\mathcal{D}$ 的对象，则答案是肯定的（这来自如下事实：
$D(A,\text{d})$ 中的直和由模的直和给出；细节从略）。
\end{remark}

\begin{lemma}
\label{dga-lemma-factor-through-nicer}
设 $(A,\text{d})$ 是微分分次代数，$E$ 是 $D(A,\text{d})$ 的紧对象。
设 $P$ 是微分分次 $A$-模，它具有有限滤过
$$
0 = F_{-1}P \subset F_0P \subset F_1P \subset \ldots \subset F_nP = P
$$
，其各项为微分分次子模，并且对某些集合 $J_i$ 与整数 $k_{i,j}$ 有
$$
F_{i + 1}P/F_iP \cong \bigoplus\nolimits_{j \in J_i} A[k_{i, j}]
$$
作为微分分次 $A$-模的同构。设 $E\to P$ 是 $D(A,\text{d})$ 的态射。
则存在微分分次子模 $P'\subset P$，使得对某些有限子集
$J'_i\subset J_i$，$F_{i+1}P\cap P'/(F_iP\cap P')$ 等于
$\bigoplus_{j\in J'_i}A[k_{i,j}]$，并且 $E\to P$ 通过 $P'$ 分解。
\end{lemma}

\begin{proof}
我们对 $-1\leq m\leq n$ 向下归纳，证明存在微分分次子模
$P'\subset P$ 使得
\begin{enumerate}
\item $F_mP\subset P'$；
\item 对 $i\geq m$，商 $F_{i+1}P\cap P'/(F_iP\cap P')$ 同构于
$\bigoplus_{j\in J'_i}A[k_{i,j}]$，其中 $J'_i\subset J_i$ 为有限子集；
\item $E\to P$ 通过 $P'$ 分解。
\end{enumerate}
初始情形是 $m=n$，此时可取 $P'=P$。

\medskip\noindent
归纳步骤。假设 $P'$ 对 $m$ 有效。对 $i\geq m$ 与 $j\in J'_i$，
取次数为 $k_{i,j}$ 的齐次元素
$x_{i,j}\in F_{i+1}P\cap P'$，使其在
$F_{i+1}P\cap P'/(F_iP\cap P')$ 中的像是对应于 $j\in J_i$
的直和项的生成元。诸 $x_{i,j}$ 作为 $A$-模生成 $P'/F_mP$。写成
$$
\text{d}(x_{i, j}) = \sum x_{i', j'} a_{i, j}^{i', j'} + y_{i, j}
$$
，其中 $y_{i,j}\in F_mP$ 且 $a_{i,j}^{i',j'}\in A$。存在有限子集
$J'_{m-1}\subset J_{m-1}$，使每个 $y_{i,j}$ 都映到
$F_mP/F_{m-1}P$ 的子模
$\bigoplus_{j\in J'_{m-1}}A[k_{m-1,j}]$ 中。令
$P''\subset F_mP$ 为该子模在映射 $F_mP\to F_mP/F_{m-1}P$
下的逆像。于是可见 $A$-子模
$$
P'' + \sum x_{i, j}A
$$
是所求类型的微分分次子模。此外，
$$
P'/(P'' + \sum x_{i, j}A) =
\bigoplus\nolimits_{j \in J_{m - 1} \setminus J'_{m - 1}} A[k_{m - 1, j}]
$$
。由于 $E$ 是紧的，给定映射 $E\to P'$ 与商映射的复合通过上面所示模
的一个有限直和子项分解。因此，扩大 $J'_{m-1}$ 后，可假设
$E\to P'$ 通过 $P''+\sum x_{i,j}A$ 分解，正合所需。
\end{proof}

\noindent
并非 $D(A,\text{d})$ 的每个紧对象都来自有限分次投射的微分分次
$A$-模；见《例》第 \ref{examples-section-interesting-compact} 节。

\begin{proposition}
\label{dga-proposition-compact}
设 $(A,\text{d})$ 是微分分次代数，$E$ 是 $D(A,\text{d})$ 的对象。
则下列条件等价：
\begin{enumerate}
\item $E$ 是紧对象；
\item $E$ 是 $D(A,\text{d})$ 的某个对象的直和项，该对象由微分分次模
$P$ 表示；$P$ 有由微分分次子模组成的有限滤过 $F_\bullet$，并且
$F_iP/F_{i-1}P$ 是 $A$ 的平移的有限直和。
\end{enumerate}
\end{proposition}

\begin{proof}
假设 $E$ 是紧的。由引理 \ref{dga-lemma-resolve}，可假设 $E$ 由具有性质
(P) 的微分分次 $A$-模 $P$ 表示。考虑由引理
\ref{dga-lemma-property-P-sequence} 的容许短正合列给出的有区别三角
$$
\bigoplus F_iP \to \bigoplus F_iP \to P
\xrightarrow{\delta} \bigoplus F_iP[1]
$$
。由于 $E$ 是紧的，对某些态射
$\delta_i:P\to F_iP[1]$，有
$\delta = \sum_{i = 1, \ldots, n} \delta_i$
。$\delta$ 与映射
$\bigoplus F_iP[1]\to\bigoplus F_iP[1]$ 的复合为零
（《导出范畴》引理 \ref{derived-lemma-composition-zero}），故
$\delta=0$（这是因为 $\bigoplus F_iP\to\bigoplus F_iP$ 在直和项
$F_iP$ 上是恒等映射与到 $F_{i-1}P$ 的包含映射之差）。因此，由
《导出范畴》引理 \ref{derived-lemma-split}，$E$ 上的恒等映射在
$D(A,\text{d})$ 中通过 $\bigoplus F_iP$ 分解。再次利用 $P$ 的紧性，
可知映射 $E\to\bigoplus F_iP$ 对某个 $n$ 通过
$\bigoplus_{i=1,\ldots,n}F_iP$ 分解。换言之，$E$ 上的恒等映射通过
$\bigoplus_{i=1,\ldots,n}F_iP$ 分解。由引理
\ref{dga-lemma-factor-through-nicer}，$E$ 的恒等映射可分解为
$E\to P\to E$，其中 $P$ 如引理陈述的第 (2) 项所述。于是已经证明
(1) 蕴含 (2)。

\medskip\noindent
假设 (2)。由《导出范畴》引理
\ref{derived-lemma-compact-objects-subcategory}，只需证明 $P$ 给出紧对象。
注意 $P$ 具有性质 (P)，故由引理 \ref{dga-lemma-hom-derived}，对任意微分
分次模 $M$ 都有
$$
\Hom_{D(A, \text{d})}(P, M) = \Hom_{K(A, \text{d})}(P, M)
$$
。由于 $D(A,\text{d})$ 中的直和由分次模的直和给出（引理
\ref{dga-lemma-derived-products}），问题约化为证明
$\Hom_{K(A,\text{d})}(P,M)$ 与直和交换。利用 $K(A,\text{d})$
是三角范畴、$\Hom$ 在第一变量中是上同调函子以及 $P$ 上的滤过，
可约化到 $P$ 是 $A$ 的平移的有限直和的情形，进而约化到显然的情形
$P=A[k]$。
\end{proof}

\begin{lemma}
\label{dga-lemma-compact-implies-bounded}
设 $(A,\text{d})$ 是微分分次代数。对 $D(A,\text{d})$ 的每个紧对象
$E$，存在整数 $a\leq b$，使得只要对 $i\in[a,b]$ 有
$H^i(M)=0$，就有 $\Hom_{D(A,\text{d})}(E,M)=0$。
\end{lemma}

\begin{proof}
注意，$D(A,\text{d})$ 中存在这样一对整数的对象所成的子范畴，
是 $D(A,\text{d})$ 的饱和严格满三角子范畴。因此由命题
\ref{dga-proposition-compact}，只需在 $E$ 由如下微分分次模 $P$ 表示时证明：
$P$ 有由微分分次子模组成的有限滤过 $F_\bullet$，且
$F_iP/F_{i-1}P$ 是 $A$ 的平移的有限直和。利用与三角的相容性，
只需对 $P=A$ 证明。此时
$\Hom_{D(A,\text{d})}(A,M)=H^0(M)$，并且取 $a=b=0$ 即得结论。
\end{proof}

\noindent
若 $(A,\text{d})$ 只是置于次数 $0$ 且微分为零的代数，或者更一般地，
只分布于有限多个次数，则确实可以得到对紧对象更精确的描述。

\begin{lemma}
\label{dga-lemma-compact}
设 $(A,\text{d})$ 是微分分次代数，并假设当 $|n|\gg0$ 时
$A^n=0$。设 $E$ 是 $D(A,\text{d})$ 的对象。下列条件等价：
\begin{enumerate}
\item $E$ 是紧对象；
\item $E$ 可由微分分次 $A$-模 $P$ 表示，其中 $P$ 作为分次 $A$-模
是有限投射的，并且对每个微分分次 $A$-模 $M$ 都满足
$\Hom_{K(A, \text{d})}(P, M) = \Hom_{D(A, \text{d})}(P, M)$
。
\end{enumerate}
\end{lemma}

\begin{proof}
设 $\mathcal{D}\subset K(A,\text{d})$ 是评注
\ref{dga-remark-source-graded-projective} 中讨论的三角子范畴。设 $P$
是 $\mathcal{D}$ 的对象，并且作为分次 $A$-模是有限投射的。则由评注
\ref{dga-remark-graded-projective-is-compact}，$P$ 表示
$D(A,\text{d})$ 的紧对象。

\medskip\noindent
为证明逆命题，设 $E$ 是 $D(A,\text{d})$ 的紧对象。固定引理
\ref{dga-lemma-compact-implies-bounded} 中的 $a\leq b$。必要时减小 $a$
并增大 $b$，还可假设当 $i\notin[a,b]$ 时 $H^i(E)=0$
（这由命题 \ref{dga-proposition-compact} 及对 $A$ 的假设得出）。此外，
固定整数 $c>0$，使得当 $|n|\geq c$ 时 $A^n=0$。

\medskip\noindent
由命题 \ref{dga-proposition-compact}，可见在 $D(A,\text{d})$ 中，$E$
是某个微分分次 $A$-模 $P$ 的直和项；$P$ 有由微分分次子模组成的有限
滤过 $F_\bullet$，且 $F_iP/F_{i-1}P$ 是 $A$ 的平移的有限直和。
特别地，$P$ 具有性质 (P)，并且由引理 \ref{dga-lemma-hom-derived}，
对任意微分分次模 $M$ 有
$\Hom_{D(A, \text{d})}(P, M) = \Hom_{K(A, \text{d})}(P, M)$
。换言之，$P$ 是评注 \ref{dga-remark-source-graded-projective} 中讨论的
三角子范畴 $\mathcal{D}\subset K(A,\text{d})$ 的对象。注意 $P$
作为分次 $A$-模是有限自由的。

\medskip\noindent
选取 $n>0$ 使 $b+4c-n<a$。以微分分次 $A$-模的自同态
$\varphi:P\to P$ 表示到 $E$ 上的投影算子。考虑 $K(A,\text{d})$
中的有区别三角
$$
P \xrightarrow{1 - \varphi} P \to C \to P[1]
$$
，其中 $C$ 是第一个箭头的锥。则 $C$ 是 $\mathcal{D}$ 的对象，
在 $D(A,\text{d})$ 中有 $C\cong E\oplus E[1]$，并且 $C$
是有限分次自由 $A$-模。接着考虑 $K(A,\text{d})$ 中的有区别三角
$$
C[1] \to C \to C' \to C[2]
$$
，其中 $C'$ 是某个态射 $C[1]\to C$ 的锥，而该态射表示
$D(A,\text{d})$ 中的复合
$$
C[1] \cong E[1] \oplus E[2] \to E[1] \to E \oplus E[1] \cong C
$$
。于是可见 $C'$ 表示 $E\oplus E[2]$。依此继续，可找到微分分次
$A$-模 $P$，它是 $\mathcal{D}$ 的对象，作为分次 $A$-模是有限自由的，
并表示 $E\oplus E[n]$。

\medskip\noindent
选取 $P$ 作为 $A$-模的齐次元素基 $x_i$（$i\in I$）。令
$d_i=\deg(x_i)$。令 $P_1$ 为由满足 $d_i\leq a-c-1$ 的
$x_i$ 与 $\text{d}(x_i)$ 生成的 $P$ 的 $A$-子模；令 $P_2$
为由满足 $d_i\geq b-n+c$ 的 $x_i$ 与 $\text{d}(x_i)$ 生成的
$P$ 的 $A$-子模。注意
\begin{enumerate}
\item $P_1$ 与 $P_2$ 是 $P$ 的微分分次子模；
\item 当 $t\geq a$ 时，$P_1^t=0$；
\item 当 $t\leq a-2c$ 时，$P_1^t=P^t$；
\item 当 $t\leq b-n$ 时，$P_2^t=0$；
\item 当 $t\geq b-n+2c$ 时，$P_2^t=P^t$。
\end{enumerate}
由于按 $n$ 的选取有 $b-n+2c\geq a-2c$，得到微分分次 $A$-模的
短正合列
$$
0 \to P_1 \cap P_2 \to P_1 \oplus P_2 \xrightarrow{\pi} P \to 0
$$
。由于 $P$ 作为分次 $A$-模是投射的，这是容许短正合列
（引理 \ref{dga-lemma-target-graded-projective}）。因此在
$K(A,\text{d})$ 中得到边界映射
$\delta:P\to(P_1\cap P_2)[1]$，见引理 \ref{dga-lemma-admissible-ses}。
由于 $P=E\oplus E[n]$，且 $P_1\cap P_2$ 集中于次数 $(b-n,a)$，
可知
$\Hom_{D(A,\text{d})}(E\oplus E[n],(P_1\cap P_2)[1])$ 为零。
又因 $P$ 是 $\mathcal{D}$ 的对象，故 $\delta$ 作为
$K(A,\text{d})$ 中的态射为零。由《导出范畴》引理
\ref{derived-lemma-split}，可找到映射 $s:P\to P_1\oplus P_2$，使得
$\pi \circ s = \text{id}_P + \text{d}h + h\text{d}$
，其中 $h:P\to P$ 的次数为 $-1$。由于 $P_1\oplus P_2\to P$
是满射，且 $P$ 作为分次 $A$-模是投射的，可以选取 $h$ 的齐次提升
$\tilde h:P\to P_1\oplus P_2$。把 $s$ 改为
$s+\text{d}\tilde h+\tilde h\text{d}$，即可得到
$\pi\circ s=\text{id}_P$。这意味着得到直和分解
$P=s^{-1}(P_1)\oplus s^{-1}(P_2)$。由于 $s^{-1}(P_2)$ 在次数
$\geq b-n+2c$ 上等于 $P$，可见 $s^{-1}(P_2)\to P\to E$
是拟同构，即 $D(A,\text{d})$ 中的同构。证明完毕。
\end{proof}








\section{导出范畴的等价}
\label{dga-section-equivalence}

\noindent
设 $R$ 是环，$(A,\text{d})$ 与 $(B,\text{d})$ 是微分分次 $R$-代数。
一个自然出现的问题是：$D(A,\text{d})$ 与 $D(B,\text{d})$ 作为
$R$-线性三角范畴等价意味着什么？这是相当微妙的问题，而且事实将表明，
它并不总是正确的问题。尽管如此，本节仍汇集若干保证这种等价成立的条件。

\medskip\noindent
我们强烈建议读者参阅关于此主题的开创性论文 \cite{Rickard}。

\begin{lemma}
\label{dga-lemma-qis-equivalence}
设 $R$ 是环，$(A,\text{d})\to(B,\text{d})$ 是微分分次 $R$-代数的
同态，并在上同调代数上诱导同构。则
$$
- \otimes_A^\mathbf{L} B : D(A, \text{d}) \to D(B, \text{d})
$$
给出三角范畴的 $R$-线性等价，其拟逆为限制函子 $N\mapsto N_A$。
\end{lemma}

\begin{proof}
由引理 \ref{dga-lemma-tensor-with-compact-fully-faithful}，函子
$M\mapsto M\otimes_A^\mathbf{L}B$ 是全忠实的。由引理
\ref{dga-lemma-tensor-hom-adjoint}，函子
$N\mapsto R\Hom(B,N)=N_A$ 是右伴随，见例
\ref{dga-example-map-hom-tensor}。显然 $R\Hom(B,-)$ 的核为零。
结论遂由《导出范畴》引理
\ref{derived-lemma-fully-faithful-adjoint-kernel-zero} 得出。
\end{proof}

\noindent
分析上述证明可见，下列推广无需额外工作即可得到。

\begin{lemma}
\label{dga-lemma-tilting-equivalence}
设 $R$ 是环，$(A,\text{d})$ 与 $(B,\text{d})$ 是微分分次 $R$-代数，
$N$ 是微分分次 $(A,B)$-双模。假设
\begin{enumerate}
\item $N$ 在 $D(B,\text{d})$ 中定义紧对象；
\item 若 $N'\in D(B,\text{d})$，且对每个 $n\in\mathbf{Z}$ 有
$\Hom_{D(B,\text{d})}(N,N'[n])=0$，则 $N'=0$；
\item 对每个 $k\in\mathbf{Z}$，映射
$H^k(A)\to\Hom_{D(B,\text{d})}(N,N[k])$ 都是同构。
\end{enumerate}
则
$$
- \otimes_A^\mathbf{L} N : D(A, \text{d}) \to D(B, \text{d})
$$
给出三角范畴的 $R$-线性等价。
\end{lemma}

\begin{proof}
由引理 \ref{dga-lemma-tensor-with-compact-fully-faithful}，函子
$M\mapsto M\otimes_A^\mathbf{L}N$ 是全忠实的。由引理
\ref{dga-lemma-tensor-hom-adjoint}，函子 $N'\mapsto R\Hom(N,N')$
是右伴随。由假设 (3)，$R\Hom(N,-)$ 的核为零。因此结论由
《导出范畴》引理 \ref{derived-lemma-fully-faithful-adjoint-kernel-zero}
得出。
\end{proof}

\begin{remark}
\label{dga-remark-tilting-equivalence}
在引理 \ref{dga-lemma-tilting-equivalence} 中，可把条件 (2) 替换为
$N$ 是 $D_{compact}(B,d)$ 的经典生成元这一条件，见《导出范畴》命题
\ref{derived-proposition-generator-versus-classical-generator}。
此外，若已知 $R\Hom(N,B)$ 是 $D(A,\text{d})$ 的紧对象，则只需检验
$N$ 是 $D_{compact}(B,\text{d})$ 的弱生成元。证明从略；如果 Stacks
project 今后用到此结论，我们会在这里补上证明。
\end{remark}

\noindent
有时把下述引理中的 $B$-模 $P$ 称为“$(A,B)$-倾斜复形”。

\begin{lemma}
\label{dga-lemma-rickard}
设 $R$ 是环，$(A,\text{d})$ 与 $(B,\text{d})$ 是微分分次 $R$-代数。
假设 $A=H^0(A)$。下列条件等价：
\begin{enumerate}
\item $D(A,\text{d})$ 与 $D(B,\text{d})$ 作为 $R$-线性三角范畴等价；
\item 存在 $D(B,\text{d})$ 的对象 $P$，使得
\begin{enumerate}
\item $P$ 是 $D(B,\text{d})$ 的紧对象；
\item 若 $N\in D(B,\text{d})$，且对每个 $i\in\mathbf{Z}$ 有
$\Hom_{D(B,\text{d})}(P,N[i])=0$，则 $N=0$；
\item 当 $i\neq0$ 时，$\Hom_{D(B,\text{d})}(P,P[i])=0$；当 $i=0$
时，该群等于 $A$。
\end{enumerate}
\end{enumerate}
由 (2) 构造的等价 $D(A,\text{d})\to D(B,\text{d})$ 把 $A$ 映到 $P$。
\end{lemma}

\begin{proof}
设 $F:D(A,\text{d})\to D(B,\text{d})$ 是等价。则 $F$ 把紧对象映到
紧对象。因此 $P=F(A)$ 是紧的，即 (2)(a) 成立。条件 (2)(b) 与
(2)(c) 由 $F$ 是等价立即得出。

\medskip\noindent
设 $P$ 是 (2) 中的对象，以具有性质 (P) 的微分分次模表示 $P$。令
$$
(E, \text{d}) = \Hom_{\text{Mod}^{dg}_{(B, \text{d})}}(P, P)
$$
。由引理 \ref{dga-lemma-hom-derived} 与假设 (2)(c)，有 $H^0(E)=A$，
且当 $k\neq0$ 时 $H^k(E)=0$。把 $P$ 看作 $(E,B)$-双模，并使用引理
\ref{dga-lemma-tilting-equivalence} 与假设 (2)(b)，得到等价
$$
D(E, \text{d}) \to D(B, \text{d})
$$
，它把 $E$ 映到 $P$。令 $E'\subset E$ 为如下微分分次 $R$-子代数：
$$
(E')^i = \left\{
\begin{matrix}
E^i & \text{若 }i < 0 \\
\Ker(E^0 \to E^1) & \text{若 }i = 0 \\
0 & \text{若 }i > 0
\end{matrix}
\right.
$$
于是有微分分次代数的拟同构
$(A,\text{d})\leftarrow(E',\text{d})\rightarrow(E,\text{d})$。
因此由引理 \ref{dga-lemma-qis-equivalence} 得到等价
$$
D(A, \text{d}) \leftarrow D(E', \text{d}) \rightarrow D(E, \text{d})
\rightarrow D(B, \text{d})
$$
。
\end{proof}

\begin{remark}
\label{dga-remark-lift-equivalence-to-dga}
设 $R$ 是环，$(A,\text{d})$ 与 $(B,\text{d})$ 是微分分次 $R$-代数。
设给定三角范畴的 $R$-线性等价
$$
F : D(A, \text{d}) \longrightarrow D(B, \text{d})
$$
。令 $N=F(A)$。则 $N$ 是微分分次 $B$-模。由于 $F$ 是等价，且
$A$ 是 $D(A,\text{d})$ 的紧对象，可知 $N$ 是 $D(B,\text{d})$
的紧对象。由于 $A$ 生成 $D(A,\text{d})$ 且 $F$ 是等价，可见 $N$
生成 $D(B,\text{d})$。最后，
$H^k(A)=\Hom_{D(A,\text{d})}(A,A[k])$，而由 $F$ 是等价可知，
对每个 $k$，$F$ 诱导同构
$H^k(A)=\Hom_{D(B,\text{d})}(N,N[k])$。要得出存在由引理
\ref{dga-lemma-tilting-equivalence} 的构造产生的等价
$D(A,\text{d})\to D(B,\text{d})$，所需的只是在 $N$ 上给出一个
与微分相容、并与给定右 $B$-模结构交换的左 $A$-模结构。事实上，
只需在把 $N$ 替换为与之拟同构的微分分次 $B$-模后完成此事。
在某些情形下可以构造这一模结构。例如，若假设 $F$ 可提升为微分分次函子
$$
F^{dg} :
\text{Mod}^{dg}_{(A, \text{d})}
\longrightarrow
\text{Mod}^{dg}_{(B, \text{d})}
$$
（记号见例 \ref{dga-example-dgm-dg-cat}），它位于相应微分分次范畴之间，
则上述要求成立。另一种情形将在下面的命题中讨论。
\end{remark}

\begin{proposition}
\label{dga-proposition-rickard}
设 $R$ 是环，$(A,\text{d})$ 与 $(B,\text{d})$ 是微分分次 $R$-代数，
$F:D(A,\text{d})\to D(B,\text{d})$ 是三角范畴的 $R$-线性等价。
假设
\begin{enumerate}
\item $A=H^0(A)$；
\item $B$ 作为 $R$-模复形是 K-平坦的。
\end{enumerate}
则存在引理 \ref{dga-lemma-tilting-equivalence} 所述的 $(A,B)$-双模 $N$。
\end{proposition}

\begin{proof}
如上面的评注 \ref{dga-remark-lift-equivalence-to-dga}，在
$D(B,\text{d})$ 中令 $N=F(A)$。可假设 $N$ 是具有性质 (P) 的微分分次
$B$-模。令
$$
(E, \text{d}) = \Hom_{\text{Mod}^{dg}_{(B, \text{d})}}(N, N)
$$
。由引理 \ref{dga-lemma-hom-derived}，有 $H^0(E)=A$，且当 $k\neq0$
时 $H^k(E)=0$。此外，由评注 \ref{dga-remark-lift-equivalence-to-dga}
中的讨论及引理 \ref{dga-lemma-tilting-equivalence}，可见 $N$ 作为
$(E,B)$-双模诱导等价
$-\otimes_E^\mathbf{L}N:D(E,\text{d})\to D(B,\text{d})$。
令 $E'\subset E$ 为如下微分分次 $R$-子代数：
$$
(E')^i = \left\{
\begin{matrix}
E^i & \text{若 }i < 0 \\
\Ker(E^0 \to E^1) & \text{若 }i = 0 \\
0 & \text{若 }i > 0
\end{matrix}
\right.
$$
于是有微分分次代数的拟同构
$(A,\text{d})\leftarrow(E',\text{d})\rightarrow(E,\text{d})$。
因此由引理 \ref{dga-lemma-qis-equivalence} 得到等价
$$
D(A, \text{d}) \leftarrow D(E', \text{d}) \rightarrow D(E, \text{d})
\rightarrow D(B, \text{d})
$$
。注意，左侧箭头的拟逆 $D(A,\text{d})\to D(E',\text{d})$ 由
$M\mapsto M\otimes_A^\mathbf{L}A$ 给出，其中把 $A$ 看作
$(A,E')$-双模，见例 \ref{dga-example-map-hom-tensor}。另一方面，函子
$D(E',\text{d})\to D(B,\text{d})$ 由
$M\mapsto M\otimes_{E'}^\mathbf{L}N$ 给出，其中 $N$ 如上。
由引理 \ref{dga-lemma-compose-tensor-functors} 即得结论。
\end{proof}

\begin{remark}
\label{dga-remark-rickard}
设 $A,B,F,N$ 如命题 \ref{dga-proposition-rickard}。尚不清楚 $F$ 与函子
$G(-)=-\otimes_A^\mathbf{L}N$ 是否同构。按构造，在
$D(B,\text{d})$ 中有同构 $N=G(A)\to F(A)$。当 $M$ 是分次投射的
微分分次 $A$-模（例如 $A$ 的平移的直和）时，可以直接把它扩张为
函子同构 $G(M)\to F(M)$。进而，当 $M$ 是此类模之间某个映射的锥时，
可得 $G(M)\cong F(M)$。一般情形下是否还有更多结论，我们并不知道。
\end{remark}

\begin{lemma}
\label{dga-lemma-rickard-rings}
设 $R$ 是环，$A,B$ 是 $R$-代数。下列条件等价：
\begin{enumerate}
\item 存在三角范畴的 $R$-线性等价 $D(A)\to D(B)$；
\item 存在 $D(B)$ 的对象 $P$，使得
\begin{enumerate}
\item $P$ 可由有限投射 $B$-模组成的有限复形表示；
\item 若 $K\in D(B)$，且对每个 $i\in\mathbf{Z}$ 有
$\Ext^i_B(P,K)=0$，则 $K=0$；
\item 当 $i\neq0$ 时，$\Ext^i_B(P,P)=0$；当 $i=0$ 时，该群等于 $A$。
\end{enumerate}
\end{enumerate}
此外，若 $B$ 作为 $R$-模是平坦的，则上述条件还等价于
\begin{enumerate}
\item[(3)] 存在 $(A,B)$-双模 $N$，使得
$-\otimes_A^\mathbf{L}N:D(A)\to D(B)$ 是等价。
\end{enumerate}
\end{lemma}

\begin{proof}
(1) 与 (2) 的等价，是引理 \ref{dga-lemma-rickard} 结合引理
\ref{dga-lemma-compact} 对 $D(B)$ 的紧对象之刻画所得的特例
（略去一个小细节）。当 $B$ 是 $R$-平坦的时，与 (3) 的等价由命题
\ref{dga-proposition-rickard} 得出。
\end{proof}

\begin{remark}
\label{dga-remark-centers}
设 $R$ 是环，$A,B$ 是 $R$-代数。若 $D(A)$ 与 $D(B)$ 作为
$R$-线性三角范畴等价，则 $A$ 与 $B$ 的中心作为 $R$-代数同构。
特别地，若 $A,B$ 交换，则 $A\cong B$。这一颇为棘手的证明见
\cite[Proposition 9.2]{Rickard} 或 \cite[Proposition 6.3.2]{KZ}。
另一种途径或许是使用 Hochschild 上同调（见下面的评注）。
\end{remark}

\begin{remark}
\label{dga-remark-hochschild-cohomology}
设 $R$ 是环，$(A,\text{d})$ 与 $(B,\text{d})$ 是导出等价的微分分次
$R$-代数，即存在三角范畴的 $R$-线性等价
$D(A,\text{d})\to D(B,\text{d})$。我们想证明 $(A,\text{d})$ 与
$(B,\text{d})$ 的某些不变量相同。在许多情形下，可以更好地控制局面。
例如，可能存在形如
$$
- \otimes_A \Omega : D(A, \text{d}) \longrightarrow D(B, \text{d})
$$
的等价，其中 $\Omega$ 是某个微分分次 $(A,B)$-双模
（命题 \ref{dga-proposition-rickard} 的情形即如此；若等价来自几何构造，
也常常如此）。若该函子的拟逆又由
$$
- \otimes_B^\mathbf{L} \Omega' : D(B, \text{d}) \longrightarrow D(A, \text{d})
$$
给出，其中 $\Omega'$ 是微分分次 $(B,A)$-双模
（同前，这样的模 $\Omega'$ 在实际中常常存在），则可考虑双模导出范畴
之间的函子
$$
D(A^{opp} \otimes_R A, \text{d})
\longrightarrow
D(B^{opp} \otimes_R B, \text{d}),\quad
M \longmapsto \Omega' \otimes^\mathbf{L}_A M \otimes_A^\mathbf{L} \Omega
$$
（使用引理 \ref{dga-lemma-bimodule-over-tensor} 把双模化为右模）。注意，
此函子把 $(A,A)$-双模 $A$ 映到 $(B,B)$-双模 $B$。在适当条件下
（例如 $A,B,\Omega$ 在 $R$ 上的平坦性等），此函子也将是等价。
若确实如此，则得到 Hochschild 上同调群的同构
$$
HH^i(A, \text{d}) =
\Hom_{D(A^{opp} \otimes_R A, \text{d})}(A, A[i])
\longrightarrow
\Hom_{D(B^{opp} \otimes_R B, \text{d})}(B, B[i]) =
HH^i(B, \text{d}).
$$
例如，若 $A=H^0(A)$，则 $HH^0(A,\text{d})$ 等于 $A$ 的中心，
这就给出了评注 \ref{dga-remark-centers} 中所述结果的概念性证明。
如果今后用到本评注，我们将在此给出精确陈述与详细证明。
\end{remark}



\section{微分分次代数的分辨}
\label{dga-section-resolution-dgas}

\noindent
设 $R$ 是环。按我们的约定，集合 $S$ 上的自由 $R$-代数
$R\langle S\rangle$ 是如下代数：它以表达式
$$
s_1 s_2 \ldots s_n
$$
为 $R$-基，其中 $n\geq0$，而 $s_1,\ldots,s_n\in S$ 是 $S$
中元素的一个序列。乘法由连接给出：
$$
(s_1 s_2 \ldots s_n) \cdot (s'_1 s'_2 \ldots s'_m) =
s_1 \ldots s_n s'_1 \ldots s'_m
$$
。此代数由如下性质刻画：对每个 $R$-代数 $A$，映射
$$
\Mor_{R\text{-代数}}(R\langle S \rangle, A) \to
\text{Map}(S, A),\quad
\varphi \longmapsto (s \mapsto \varphi(s))
$$
是双射。

\medskip\noindent
在分次 $R$-代数范畴中，集合 $S$ 应带有分次，我们把它看作映射
$\deg:S\to\mathbf{Z}$。于是 $R\langle S\rangle$ 带有分次，使单项式
的次数为
$$
\deg(s_1 s_2 \ldots s_n) = \deg(s_1) + \ldots + \deg(s_n)
$$
。在此设定下，对每个分次 $R$-代数 $A$，映射
$$
\Mor_{\text{分次 }R\text{-代数}}(R\langle S \rangle, A) \to
\text{Map}_{\text{分次集合}}(S, A),\quad
\varphi \longmapsto (s \mapsto \varphi(s))
$$
是双射。

\medskip\noindent
若 $A$ 是分次 $R$-代数，$S$ 是分次集合，也可类似地构造
$A\langle S\rangle$。$A\langle S\rangle$ 的元素是形如
$$
a_0 s_1 a_1 s_2 \ldots a_{n - 1} s_n a_n
$$
的元素之和，其中 $a_i\in A$，并模去使这些表达式关于
$(a_0,\ldots,a_n)$ 为 $R$-多重线性的关系。因此，对 $S$ 中元素的
每个序列 $s_1,\ldots,s_n$，都有包含
$$
A \otimes_R \ldots \otimes_R A \subset A\langle S \rangle
$$
，而该代数是这些项的直和。由此定义，读者可证明映射
$$
\Mor_{\text{分次 }R\text{-代数}}(A\langle S \rangle, B) \to
\Mor_{\text{分次 }R\text{-代数}}(A, B) \times
\text{Map}_{\text{分次集合}}(S, B),
$$
把 $\varphi$ 映到 $(\varphi|_A,(s\mapsto\varphi(s)))$；
对每个分次 $R$-代数 $A$，该映射是双射。注意，若 $A$ 是自由分次
$R$-代数，则 $A\langle S\rangle$ 也是自由分次 $R$-代数。

\medskip\noindent
设 $A$ 是微分分次 $R$-代数，$S$ 是分次集合。再假设对每个
$s\in S$，给定齐次元素 $f_s\in A$，满足
$\deg(f_s)=\deg(s)+1$ 且 $\text{d}f_s=0$。则在
$A\langle S\rangle$ 上存在唯一的微分分次代数结构，使得
$\text{d}(s)=f_s$。例如，给定 $a,b,c\in A$ 与 $s,t\in S$，定义
\begin{align*}
\text{d}(asbtc)
& =
\text{d}(a)sbtc + (-1)^{\deg(a)}a f_s b t c +
(-1)^{\deg(a) + \deg(s)} as\text{d}(b)tc \\
& + (-1)^{\deg(a) + \deg(s) + \deg(b)} asb f_t c +
(-1)^{\deg(a) + \deg(s) + \deg(b) + \deg(t)} asbt\text{d}(c)
\end{align*}
。细节从略。

\begin{lemma}
\label{dga-lemma-K-flat-resolution}
设 $R$ 是环，$(B,\text{d})$ 是微分分次 $R$-代数。存在微分分次
$R$-代数的拟同构 $(A,\text{d})\to(B,\text{d})$，并具有下列性质：
\begin{enumerate}
\item $A$ 作为 $R$-模复形是 K-平坦的；
\item $A$ 是自由分次 $R$-代数。
\end{enumerate}
\end{lemma}

\begin{proof}
首先断言，可找到具有 (1)、(2) 且在上同调上诱导满射的
$(A_0,\text{d})\to(B,\text{d})$。具体地，取分次集合 $S$；
对每个 $s\in S$，取次数为 $\deg(s)$ 的齐次元素
$b_s\in\Ker(d:B\to B)$，使其在 $H^*(B)$ 中的类
$\overline b_s$ 作为 $R$-模生成 $H^*(B)$。可令
$A_0=R\langle S\rangle$，微分为零，并以 $s\mapsto b_s$
定义 $A_0\to B$。

\medskip\noindent
给定在上同调上诱导满射的 $A_0\to B$，归纳构造序列
$$
A_0 \to A_1 \to A_2 \to \ldots B
$$
。给定 $A_n\to B$，令 $S_n$ 为分次集合；对每个 $s\in S_n$，
令 $a_s\in\Ker(\text{d}:A_n\to A_n)$ 为次数
$\deg(s)+1$ 的齐次元素，其在 $H^*(A_n)$ 中的类
$\overline a_s$ 映到 $H^*(B)$ 中的零。取 $S_n$ 足够大，使诸
$\overline a_s$ 作为 $R$-模生成
$\Ker(H^*(A_n)\to H^*(B))$。然后令
$$
A_{n + 1} = A_n\langle S_n \rangle
$$
，其微分由 $\text{d}(s)=a_s$ 给出，见上面的讨论。于是每个
$(A_n,\text{d})$ 都满足 (1)、(2)；细节从略。映射
$H^*(A_n)\to H^*(B)$ 是满射，因为 $n=0$ 时已经如此。

\medskip\noindent
显然 $A=\colim A_n$ 是自由分次 $R$-代数。由《代数进阶》引理
\ref{more-algebra-lemma-colimit-K-flat}，它是 K-平坦的。映射
$H^*(A)\to H^*(B)$ 是同构：它是满射，也是单射。事实上，
$H^*(A)$ 的每个元素都来自某个 $H^*(A_n)$ 中的元素；若它在
$H^*(B)$ 中消失，则它在 $H^*(A_{n+1})$ 中消失，因而也在
$H^*(A)$ 中消失。
\end{proof}

\noindent
作为应用，我们证明引理
\ref{dga-lemma-compose-tensor-functors-general-algebra} 的“正确”版本。

\begin{lemma}
\label{dga-lemma-compose-tensor-functors-tor}
设 $R$ 是环，$(A,\text{d})$、$(B,\text{d})$ 和 $(C,\text{d})$
是微分分次 $R$-代数。假设 $A\otimes_RC$ 在 $D(R)$ 中表示
$A\otimes_R^\mathbf{L}C$。设 $N$ 是微分分次 $(A,B)$-双模，
$N'$ 是微分分次 $(B,C)$-双模。则复合
$$
\xymatrix{
D(A, \text{d}) \ar[rr]^{- \otimes_A^\mathbf{L} N} & &
D(B, \text{d}) \ar[rr]^{- \otimes_B^\mathbf{L} N'} & &
D(C, \text{d})
}
$$
同构于 $-\otimes_A^\mathbf{L}N''$，其中 $N''$ 是某个微分分次
$(A,C)$-双模。
\end{lemma}

\begin{proof}
利用引理 \ref{dga-lemma-K-flat-resolution}，选取拟同构
$(B',\text{d})\to(B,\text{d})$，其中 $B'$ 作为 $R$-模复形是
K-平坦的。由引理 \ref{dga-lemma-qis-equivalence}，函子
$-\otimes_{B'}^\mathbf{L}B:D(B',\text{d})\to D(B,\text{d})$
是等价，其拟逆由限制给出。注意，限制典范同构于函子
$-\otimes_B^\mathbf{L}B:D(B,\text{d})\to D(B',\text{d})$，
其中把 $B$ 看作 $(B,B')$-双模。因此，只需对下列复合证明引理：
$$
D(A) \to D(B) \to D(B'),\quad
D(B') \to D(B) \to D(C),\quad
D(A) \to D(B') \to D(C).
$$
第一个复合由引理 \ref{dga-lemma-compose-tensor-functors} 处理，因为 $B'$
作为 $R$-模复形是 K-平坦的。第二个复合成立，因为
$B \otimes_B^\mathbf{L} N' = N' = B \otimes_B N'$
，故可应用引理 \ref{dga-lemma-compose-tensor-functors-general}。于是约化到
$B$ 作为 $R$-模复形是 K-平坦的情形。

\medskip\noindent
假设 $B$ 作为 $R$-模复形是 K-平坦的。只需证明
(\ref{dga-equation-plain-versus-derived-algebras}) 是同构，见引理
\ref{dga-lemma-compose-tensor-functors-general-algebra}。选取拟同构
$L\to A$，其中 $L$ 是具有性质 (P) 的微分分次 $R$-模。显然
$P=L\otimes_RB$ 作为微分分次 $B$-模具有性质 (P)。因此必须证明
$P\to A\otimes_RB$ 诱导拟同构
$$
P \otimes_B (B \otimes_R C)
\longrightarrow
(A \otimes_R B) \otimes_B (B \otimes_R C)
$$
。可将其改写为
$$
P \otimes_R B \otimes_R C \longrightarrow A \otimes_R B \otimes_R C
$$
。由于 $B$ 作为 $R$-模复形是 K-平坦的，由《代数进阶》引理
\ref{more-algebra-lemma-K-flat-quasi-isomorphism}，只需证明
$$
P \otimes_R C \to A \otimes_R C
$$
是拟同构，而这恰是我们的假设。
\end{proof}

\noindent
下列引理其实不属于本节，但似乎没有更自然的安置位置。

\begin{lemma}
\label{dga-lemma-countable}
设 $(A,\text{d})$ 是微分分次代数，且对每个 $i$，$H^i(A)$ 都可数。
设 $M$ 是 $D(A,\text{d})$ 的对象。下列条件等价：
\begin{enumerate}
\item $M=\text{hocolim},E_n$，其中 $E_n$ 是 $D(A,\text{d})$ 的紧对象；
\item 对每个 $i$，$H^i(M)$ 都可数。
\end{enumerate}
\end{lemma}

\begin{proof}
假设 (1) 成立。由《导出范畴》引理
\ref{derived-lemma-cohomology-of-hocolim}，有
$H^i(M)=\colim H^i(E_n)$。因此只需证明对每个 $n$，$H^i(E_n)$
可数。由命题 \ref{dga-proposition-compact}，可见 $E_n$ 在
$D(A,\text{d})$ 中同构于某个微分分次模 $P$ 的直和项；$P$ 有由微分
分次子模组成的有限滤过 $F_\bullet$，且 $F_jP/F_{j-1}P$ 是 $A$
的平移的有限直和。依假设，群 $H^i(F_jP/F_{j-1}P)$ 可数。
对滤过长度归纳并使用上同调长正合列，即得 (2)。更有趣的是反向蕴含。

\medskip\noindent
我们断言，存在可数微分分次子代数 $A'\subset A$，使包含映射
$A'\to A$ 在上同调上定义同构。为构造 $A'$，选取可数微分分次子代数
$$
A_1 \subset A_2 \subset A_3 \subset \ldots
$$
，使得：(a) $H^i(A_1)\to H^i(A)$ 是满射；(b) 对 $n>1$，映射
$H^i(A_{n-1})\to H^i(A_n)$ 的核与映射
$H^i(A_{n-1})\to H^i(A)$ 的核相同。为构造 $A_1$，取生成 $A$
的上同调（作为环或作为分次阿贝尔群）的任意可数上链集
$S\subset A$，并令 $A_1$ 为由 $S$ 生成的 $A$ 的微分分次子代数。
给定 $A_{n-1}$ 后，为构造 $A_n$，对每个在 $H^i(A)$ 中映到零的上链
$a\in A_{n-1}^i$，选取 $s_a\in A^{i-1}$ 使
$\text{d}(s_a)=a$，并令 $A_n$ 为由 $A_{n-1}$ 与诸元素 $s_a$
生成的 $A$ 的微分分次子代数。最后取 $A'=\bigcup A_n$。

\medskip\noindent
由引理 \ref{dga-lemma-qis-equivalence}，限制映射
$D(A,\text{d})\to D(A',\text{d})$，$M\mapsto M_{A'}$，是等价。
由于 $M$ 与 $M_{A'}$ 的上同调群相同，只需对 $(A',\text{d})$
证明蕴含 (2) $\Rightarrow$ (1)。

\medskip\noindent
假设 $A$ 可数。用与上面完全相同类型的论证可见，对
$D(A,\text{d})$ 中的 $M$，下列两者等价：对每个 $i$，$H^i(M)$
可数；$M$ 可由可数微分分次模表示。因此，为证明蕴含
(2) $\Rightarrow$ (1)，可约化到下一段所述情形。

\medskip\noindent
假设 $A$ 可数，且 $M$ 是可数微分分次 $A$-模。我们断言，存在微分分次
$A$-模同态 $P\to M$，使得
\begin{enumerate}
\item $P\to M$ 是拟同构；
\item $P$ 具有性质 (P)；
\item $P$ 可数。
\end{enumerate}
查看引理 \ref{dga-lemma-resolve} 中构造 P-分辨的证明可见，只需说明在可数
微分分次模的设定下仍能证明引理 \ref{dga-lemma-good-quotient}。
这由该证明立即得出。

\medskip\noindent
假设 $A$ 可数，且 $M$ 是具有性质 (P) 的可数微分分次模。选取由微分分次
子模组成的滤过
$$
0 = F_{-1}P \subset F_0P \subset F_1P \subset \ldots \subset P
$$
，使得
\begin{enumerate}
\item $P=\bigcup F_pP$；
\item $F_iP\to F_{i+1}P$ 是容许单态射；
\item 对某些集合 $J_i$ 与整数 $k_j$，有微分分次模的同构
$F_iP/F_{i - 1}P \to \bigoplus_{j \in J_i} A[k_j]$
。
\end{enumerate}
当然，对每个 $i$，$J_i$ 都可数。对每个 $i$ 与 $j\in J_i$，
选取次数为 $k_j$ 的 $x_{i,j}\in F_iP$，使其在
$F_iP/F_{i-1}P$ 中的像生成对应于 $j$ 的直和项。

\medskip\noindent
断言：给定 $n$ 以及有限子集 $S_i\subset J_i$（$i=1,\ldots,n$），
存在有限子集 $S_i\subset T_i\subset J_i$（$i=1,\ldots,n$），使得
$P'=\bigoplus_{i\leq n}\bigoplus_{j\in T_i}Ax_{i,j}$ 是 $P$
的微分分次子模。这已在引理 \ref{dga-lemma-factor-through-nicer} 的证明中
说明，也容易直接证明：诸元素 $x_{i,j}$ 作为右 $A$-模自由生成 $P$。
$P$ 的结构表明
$$
\text{d}(x_{i, j}) = \sum\nolimits_{i' < i} x_{i', j'}a_{i', j'}
$$
，其中求和当然是有限的。因此，给定 $S_0,\ldots,S_n$，可先选取
$S_0\subset S'_0,\ldots,S_{n-1}\subset S'_{n-1}$，使得对每个
$j\in S_n$ 都有
$\text{d}(x_{n, j}) \in \bigoplus_{i' < n, j' \in S'_{i'}} x_{i', j'}A$
。然后对 $n$ 归纳，可以选取
$S'_0\subset T_0,\ldots,S'_{n-1}\subset T_{n-1}$，保证
$\bigoplus_{i'<n,j'\in T_{i'}}x_{i',j'}A$ 是微分分次 $A$-子模。
令 $T_n=S_n$，便得到所需的
$P'=\bigoplus_{i\leq n,j\in T_i}x_{i,j}A$。

\medskip\noindent
由此断言显然可知，$P=\bigcup P'_n$ 是上述 $P'_n$ 的可数递增并。
按构造，每个 $P'_n$ 都是具有性质 (P) 的微分分次模，其滤过有限，
且逐次商是 $A$ 的平移的有限直和。因此 $P'_n$ 定义
$D(A,\text{d})$ 的紧对象，例如见命题 \ref{dga-proposition-compact}。
由引理 \ref{dga-lemma-homotopy-colimit}，在 $D(A,\text{d})$ 中有
$P=\text{hocolim},P'_n$，故蕴含 (2) $\Rightarrow$ (1) 证明完毕。
\end{proof}
% Generated mechanically from the frozen dpa.tex target.
% Do not edit this generated file; edit the bound locale source instead.

% OK, start here.
%

\setcounter{chapter}{22}
\chapter{除幂代数}



\phantomsection
\label{dpa-section-phantom}


\section{引言}
\label{dpa-section-introduction}

\noindent
本章讨论除幂代数及其应用。一本参考书是 \cite{Berthelot}。





\section{除幂}
\label{dpa-section-divided-powers}

\noindent
本节汇集除幂环的一些结果。
我们采用约定 $0! = 1$（空乘积应当等于 $1$）。

\begin{definition}
\label{dpa-definition-divided-powers}
设 $A$ 是环，$I$ 是 $A$ 的理想。若一族映射
$\gamma_n : I \to I$（$n > 0$）对所有 $n \geq 0$、$m > 0$、
$x, y \in I$ 和 $a \in A$ 满足
\begin{enumerate}
\item $\gamma_1(x) = x$，并约定 $\gamma_0(x) = 1$；
\item $\gamma_n(x)\gamma_m(x) = \frac{(n + m)!}{n! m!} \gamma_{n + m}(x)$,
\item $\gamma_n(ax) = a^n \gamma_n(x)$,
\item $\gamma_n(x + y) = \sum_{i = 0, \ldots, n} \gamma_i(x)\gamma_{n - i}(y)$,
\item $\gamma_n(\gamma_m(x)) = \frac{(nm)!}{n! (m!)^n} \gamma_{nm}(x)$.
\end{enumerate}
则称它为 $I$ 上的{\it 除幂结构}。
\end{definition}

\noindent
注意，定义中出现的有理数 $\frac{(n + m)!}{n! m!}$
和 $\frac{(nm)!}{n! (m!)^n}$ 实际上都是整数；
前者是从 $n + m$ 个对象中选取 $n$ 个的方式数，
后者是把 $nm$ 个对象分成 $n$ 个各含 $m$ 个对象的组的方式数。
下面关于该定义的若干说明表明，$\gamma_n(x)$ 在 $I$ 中代替了 $x^n/n!$。

\begin{lemma}
\label{dpa-lemma-silly}
设 $A$ 是环，$I$ 是 $A$ 的理想。
\begin{enumerate}
\item 若 $\gamma$ 是 $I$ 上的除幂结构\footnote{这里以及下文中，
$\gamma$ 是从 $I$ 到 $I$ 的映射序列
$\gamma_1, \gamma_2, \gamma_3, \ldots$ 的简写。}，则对
$n \geq 1$、$x \in I$ 有 $n! \gamma_n(x) = x^n$。
\end{enumerate}
假设 $A$ 作为 $\mathbf{Z}$-模无挠。
\begin{enumerate}
\item[(2)] $I$ 上的除幂结构若存在，则唯一。
\item[(3)] 若 $\gamma_n : I \to I$ 是映射，则
$$
\gamma\text{ 是除幂结构}
\Leftrightarrow
n! \gamma_n(x) = x^n\ \forall x \in I, n \geq 1.
$$
\item[(4)] 理想 $I$ 具有除幂结构，当且仅当存在 $I$ 作为理想的一组生成元
$x_i$，使得对所有 $n \geq 1$ 都有 $x_i^n \in (n!)I$。
\end{enumerate}
\end{lemma}

\begin{proof}
\emph{(1) 的证明。}
若 $\gamma$ 是除幂结构，则条件 (2)（把其中的 $n,m$ 换成
$1,n-1$）给出
$n \gamma_n(x) = \gamma_1(x)\gamma_{n - 1}(x)$。
因此，由归纳法和条件 (1) 得到 $n! \gamma_n(x) = x^n$。

\medskip\noindent
假设 $A$ 作为 $\mathbf{Z}$-模无挠。
\emph{(2) 的证明。} 这由 (1) 立即得到。

\medskip\noindent
\emph{(3) 的证明。}
假设对所有 $x \in I$ 和 $n \geq 1$ 都有
$n! \gamma_n(x) = x^n$。由于
$A \subset A \otimes_{\mathbf{Z}} \mathbf{Q}$，只须在
$A$ 是 $\mathbf{Q}$-代数时证明定义
\ref{dpa-definition-divided-powers} 的公理 (1)--(5)。
此时 $\gamma_n(x) = x^n/n!$，直接验证即可得到 (1)--(5)；
例如，(4) 对应于二项式公式
$$
(x + y)^n = \sum_{i = 0, \ldots, n} \frac{n!}{i!(n - i)!} x^iy^{n - i}
$$
读者不妨完成这些验证，以确认各系数无误。

\medskip\noindent
\emph{(4) 的证明。}
假设 $x_i$ 是 $I$ 作为理想的一组生成元，并且对所有 $n \geq 1$
都有 $x_i^n \in (n!)I$。我们断言，对所有 $x \in I$ 都有
$x^n \in (n!)I$。若该断言成立，则可令
$\gamma_n(x) = x^n/n!$；由 (3)，这给出一个除幂结构。
为证明该断言，注意它对 $x = ax_i$ 成立。因此，该断言对
$I$ 作为阿贝尔群的一组生成元成立。对由这些生成元组成的表达式长度作归纳，
只须证明：若断言对 $x$ 和 $y$ 成立，则它对 $x+y$ 也成立。
这立即由二项式定理得到。
\end{proof}

\begin{example}
\label{dpa-example-ideal-generated-by-p}
设 $p$ 是素数。设 $A$ 是环，并且每个不能被 $p$ 整除的整数 $n$
在 $A$ 中都可逆；也就是说，$A$ 是 $\mathbf{Z}_{(p)}$-代数。
于是 $I = pA$ 具有典范除幂结构。具体地，对给定的
$x = pa \in I$，令
$$
\gamma_n(x) = \frac{p^n}{n!} a^n
$$
读者可以立即验证，对 $n \geq 1$ 有
$p^n/n! \in p\mathbf{Z}_{(p)}$（例如，这可由下述事实推出：
$n!$ 的素因子分解中 $p$ 的指数为
$\left\lfloor n/p \right\rfloor + \left\lfloor n/p^2 \right\rfloor
+ \left\lfloor n/p^3 \right\rfloor + \ldots$）。
因此该定义有意义，并给出一列映射 $\gamma_n : I \to I$。
直接验证即可看出定义 \ref{dpa-definition-divided-powers} 的条件 (1)--(5)
均成立。也可以先注意到该定义对
$A_0 = \mathbf{Z}_{(p)}$ 成立，再由引理 \ref{dpa-lemma-gamma-extends}
得到结论。
\end{example}

\noindent
注意，对 $A$ 的任意理想 $I$ 及其上的任意除幂结构 $\gamma$，
都有 $\gamma_n\left(0\right) = 0$。（这由定义
\ref{dpa-definition-divided-powers} 的公理 (3) 对 $a=0$ 的应用得到。）

\begin{lemma}
\label{dpa-lemma-check-on-generators}
设 $A$ 是环，$I$ 是 $A$ 的理想，并设
$\gamma_n : I \to I$（$n \geq 1$）是一列映射。假设
\begin{enumerate}
\item[(a)] 定义 \ref{dpa-definition-divided-powers} 的 (1)、(3) 和 (4)
对所有 $x,y\in I$ 成立；
\item[(b)] 性质 (2) 和 (5) 对 $I$ 作为理想的某组生成元中的 $x$ 成立。
\end{enumerate}
则 $\gamma$ 是 $I$ 上的除幂结构。
\end{lemma}

\begin{proof}
本证明中的编号 (1)、(2)、(3)、(4)、(5) 指定义
\ref{dpa-definition-divided-powers} 所列的条件。
应用 (3) 可知，若 (2) 和 (5) 对 $x$ 成立，则它们对所有
$a\in A$ 的 $ax$ 也成立。因此，(b) 蕴含 (2) 和 (5)
对 $I$ 作为阿贝尔群的一组生成元成立。
对由这些生成元组成的表达式长度作归纳，只须证明：若
(2) 和 (5) 对 $x,y\in I$ 成立，则它们对 $x+y$ 也成立。

\medskip\noindent
\emph{$x+y$ 的性质 (2)。} 由 (4)，
$$
\gamma_n(x + y)\gamma_m(x + y) =
\sum\nolimits_{i + j = n,\ k + l = m}
\gamma_i(x)\gamma_k(x)\gamma_j(y)\gamma_l(y)
$$
应用 $x$ 和 $y$ 的性质 (2)，上式等于
$$
\sum \frac{(i + k)!}{i!k!}\frac{(j + l)!}{j!l!}
\gamma_{i + k}(x)\gamma_{j + l}(y)
$$
将它与展开式
$$
\gamma_{n + m}(x + y) = \sum \gamma_a(x)\gamma_b(y)
$$
比较可知，只须证明：给定 $a+b=n+m$，有
$$
\sum\nolimits_{i + k = a,\ j + l = b,\ i + j = n,\ k + l = m}
\frac{(i + k)!}{i!k!}\frac{(j + l)!}{j!l!}
=
\frac{(n + m)!}{n!m!}.
$$
不直接证明此式，而是注意到它对多项式环 $\mathbf{Q}[x,y]$
中的理想 $I=(x,y)$ 成立，因为对 $f\in I$，
$\gamma_n(f)=f^n/n!$ 定义了 $I$ 上的除幂结构。
因此，上述有理数恒等式成立。

\medskip\noindent
\emph{$x+y$ 的性质 (5)。}
假设 (1)--(4) 成立，并且 (5) 对 $x$ 和 $y$ 成立。
我们再次把证明归结为一个有理数恒等式。具体地，应用 (4) 可写成
$\gamma_n(\gamma_m(x + y)) = \gamma_n(\sum \gamma_i(x)\gamma_j(y))$.
再次应用 (4)，可把 $\gamma_n(\gamma_m(x+y))$
写成若干项之和，其中每一项都是形如
$\gamma_k(\gamma_i(x)\gamma_j(y))$ 的因子的乘积。
若 $i>0$，则
\begin{align*}
\gamma_k(\gamma_i(x)\gamma_j(y)) & =
\gamma_j(y)^k\gamma_k(\gamma_i(x)) \\
& = \frac{(ki)!}{k!(i!)^k} \gamma_j(y)^k \gamma_{ki}(x) \\
& =
\frac{(ki)!}{k!(i!)^k} \frac{(kj)!}{(j!)^k} \gamma_{ki}(x) \gamma_{kj}(y)
\end{align*}
其中第一个等号使用 (3)，第二个等号使用 $x$ 的 (5)，第三个等号
恰好使用 $k$ 次 (2)。应用 $y$ 的 (5) 可知，该等式在 $i=0$
时同样成立。继续如此应用除 (5) 外的所有公理，可写成
$$
\gamma_n(\gamma_m(x + y)) =
\sum\nolimits_{i + j = nm} c_{ij}\gamma_i(x)\gamma_j(y)
$$
其中 $c_{ij}\in\mathbf{Z}$ 是某些普适常数。该等式在多项式环
$\mathbf{Q}[x,y]$ 中成立这一事实再次蕴含所有系数 $c_{ij}$
均等于 $(nm)!/n!(m!)^n$，正如所需。
\end{proof}

\begin{lemma}
\label{dpa-lemma-two-ideals}
设 $A$ 是环，并有两个理想 $I,J\subset A$。
设 $\gamma$ 是 $I$ 上的除幂结构，$\delta$ 是 $J$ 上的除幂结构。则
\begin{enumerate}
\item $\gamma$ 与 $\delta$ 在 $IJ$ 上一致；
\item 若 $\gamma$ 与 $\delta$ 在 $I\cap J$ 上一致，则它们是
$I+J$ 上唯一除幂结构 $\epsilon$ 的限制。
\end{enumerate}
\end{lemma}

\begin{proof}
设 $x\in I$ 且 $y\in J$。则
$$
\gamma_n(xy) = y^n\gamma_n(x) = n! \delta_n(y) \gamma_n(x) =
\delta_n(y) x^n = \delta_n(xy).
$$
因此，$\gamma$ 与 $\delta$ 在 $IJ$ 的一组（加法）生成元上相同。
由定义 \ref{dpa-definition-divided-powers} 的性质 (4)，它们在整个
$IJ$ 上相同。

\medskip\noindent
假设 $\gamma$ 与 $\delta$ 在 $I\cap J$ 上一致。
给定 $z\in I+J$，将其写成 $z=x+y$，其中 $x\in I$ 且 $y\in J$。
对所有 $n\geq 1$，令
$$
\epsilon_n(z) = \sum \gamma_i(x)\delta_{n - i}(y)
$$
为说明它是良定义的，设 $z=x'+y'$ 是另一种表示，其中
$x'\in I$ 且 $y'\in J$。于是
$w=x-x'=y'-y\in I\cap J$。因此
\begin{align*}
\sum\nolimits_{i + j = n} \gamma_i(x)\delta_j(y)
& =
\sum\nolimits_{i + j = n} \gamma_i(x' + w)\delta_j(y) \\
& =
\sum\nolimits_{i' + l + j = n} \gamma_{i'}(x')\gamma_l(w)\delta_j(y) \\
& =
\sum\nolimits_{i' + l + j = n} \gamma_{i'}(x')\delta_l(w)\delta_j(y) \\
& =
\sum\nolimits_{i' + j' = n} \gamma_{i'}(x')\delta_{j'}(y + w) \\
& =
\sum\nolimits_{i' + j' = n} \gamma_{i'}(x')\delta_{j'}(y')
\end{align*}
正如所需。因此，我们对所有 $n\geq 1$ 定义了映射
$\epsilon_n:I+J\to I+J$；容易看出
$\epsilon_n\mid_I=\gamma_n$ 且 $\epsilon_n\mid_J=\delta_n$。
下面验证映射族 $\epsilon_n$ 满足定义
\ref{dpa-definition-divided-powers} 的条件 (1)--(5)。
性质 (1) 和 (3) 是显然的。为验证 (4)，设
$z=x+y$、$z'=x'+y'$，其中 $x,x'\in I$ 且 $y,y'\in J$，计算
\begin{align*}
\epsilon_n(z + z') & =
\sum\nolimits_{a + b = n} \gamma_a(x + x')\delta_b(y + y') \\
& =
\sum\nolimits_{i + i' + j + j' = n}
\gamma_i(x) \gamma_{i'}(x')\delta_j(y)\delta_{j'}(y') \\
& =
\sum\nolimits_{k = 0, \ldots, n}
\sum\nolimits_{i+j=k} \gamma_i(x)\delta_j(y)
\sum\nolimits_{i'+j'=n-k} \gamma_{i'}(x')\delta_{j'}(y') \\
& =
\sum\nolimits_{k = 0, \ldots, n}\epsilon_k(z)\epsilon_{n-k}(z')
\end{align*}
正如所需。现在只须对 $I$ 或 $J$ 中的元素证明 (2) 和 (5)，
见引理 \ref{dpa-lemma-check-on-generators}。
由于 $\gamma$ 和 $\delta$ 都是除幂结构，这一点是显然的。

\medskip\noindent
由此证明了 $I+J$ 上存在除幂结构 $\epsilon$，其在 $I$ 和 $J$
上的限制分别为 $\gamma$ 和 $\delta$；其唯一性也很清楚。
\end{proof}

\begin{lemma}
\label{dpa-lemma-nil}
设 $p$ 是素数，$A$ 是环，$I\subset A$ 是理想，并设
$\gamma$ 是 $I$ 上的除幂结构。假设 $p$ 在 $A/I$ 中幂零。
则 $I$ 局部幂零，当且仅当 $p$ 在 $A$ 中幂零。
\end{lemma}

\begin{proof}
若在 $A$ 中 $p^N=0$，则对 $x\in I$ 有
$x^{pN}=(pN)!\gamma_{pN}(x)=0$，因为 $(pN)!$ 可被 $p^N$ 整除。
反之，假设 $I$ 局部幂零。我们还假设 $p$ 在 $A/I$ 中幂零，
故对某个 $r$ 有 $p^r\in I$；于是 $p^r$ 幂零，从而 $p$ 幂零。
\end{proof}








\section{除幂环}
\label{dpa-section-divided-power-rings}

\noindent
除幂环构成一个范畴。定义如下。

\begin{definition}
\label{dpa-definition-divided-power-ring}
{\it 除幂环}是三元组 $(A,I,\gamma)$，其中 $A$ 是环，
$I\subset A$ 是理想，而
$\gamma=(\gamma_n)_{n\geq 1}$ 是 $I$ 上的除幂结构。
{\it 除幂环同态}
$\varphi:(A,I,\gamma)\to(B,J,\delta)$ 是环同态
$\varphi:A\to B$，满足 $\varphi(I)\subset J$，并且对所有
$x\in I$ 和 $n\geq 1$ 都有
$\delta_n(\varphi(x))=\varphi(\gamma_n(x))$。
\end{definition}

\noindent
有时我们说“设 $(B,J,\delta)$ 是 $(A,I,\gamma)$ 上的除幂代数”，
意思是 $(B,J,\delta)$ 是一个除幂环，并配备了一个除幂环同态
$(A,I,\gamma)\to(B,J,\delta)$。

\begin{lemma}
\label{dpa-lemma-limits}
除幂环范畴具有所有极限，并且这些极限与环范畴中的极限一致。
\end{lemma}

\begin{proof}
空极限是零环（这有些奇怪，但我们需要如此）。
一族除幂环 $(A_t,I_t,\gamma_t)$（$t\in T$）的积为
$(\prod A_t,\prod I_t,\gamma)$，其中
$\gamma_n((x_t))=(\gamma_{t,n}(x_t))$。
两个同态
$\alpha,\beta:(A,I,\gamma)\to(B,J,\delta)$ 的等化子就是
$C=\{a\in A\mid\alpha(a)=\beta(a)\}$，配备理想 $C\cap I$
及其诱导的除幂。由此所有极限都存在，见
《范畴》引理 \ref{categories-lemma-limits-products-equalizers}。
\end{proof}

\noindent
下述引理以除幂代数为例说明一种很一般的范畴论现象。

\begin{lemma}
\label{dpa-lemma-a-version-of-brown}
设 $\mathcal{C}$ 是除幂环范畴，并设
$F:\mathcal{C}\to\textit{Sets}$ 是函子。假设
\begin{enumerate}
\item 存在基数 $\kappa$，使得对每个 $f\in F(A,I,\gamma)$，
都存在 $\mathcal{C}$ 中的态射
$(A',I',\gamma')\to(A,I,\gamma)$，其中 $|A'|\leq\kappa$，
并且 $f$ 是某个 $f'\in F(A',I',\gamma')$ 的像；
\item $F$ 与极限交换。
\end{enumerate}
则 $F$ 可表；也就是说，存在 $\mathcal{C}$ 的对象 $(B,J,\delta)$，
使得
$$
F(A, I, \gamma) = \Hom_\mathcal{C}((B, J, \delta), (A, I, \gamma))
$$
关于 $(A,I,\gamma)$ 函子性地成立。
\end{lemma}

\begin{proof}
这是《范畴》引理 \ref{categories-lemma-a-version-of-brown}
的一个特殊情形。
\end{proof}

\begin{lemma}
\label{dpa-lemma-colimits}
除幂环范畴具有所有余极限。
\end{lemma}

\begin{proof}
空余极限是带除幂理想 $(0)$ 的 $\mathbf{Z}$。
下面讨论一般的余极限。设 $\mathcal{C}$ 是范畴，并设
$c\mapsto(A_c,I_c,\gamma_c)$ 是一个图。考虑函子
$$
F(B, J, \delta) = \lim_{c \in \mathcal{C}}
Hom((A_c, I_c, \gamma_c), (B, J, \delta))
$$
注意，任意
$f=(f_c)_{c\in C}\in F(B,J,\delta)$ 都具有如下性质：
所有像 $f_c(A_c)$ 生成 $B$ 的一个基数受 $\kappa$ 控制的子环
$B'$，而所有像 $f_c(I_c)$ 在 $B'$ 中生成一个除幂子理想 $J'$。
于是 $f$ 分解为 $F(B')$ 中的某个 $f'$，再复合包含
$B'\to B$。此外，$F$ 与极限交换。因此可以应用引理
\ref{dpa-lemma-a-version-of-brown}，得到 $F$ 可表，结论成立。
\end{proof}

\begin{remark}
\label{dpa-remark-pushout}
设
$$
\xymatrix{
(B, J, \delta) \ar[r] & (B'', J'', \delta'') \\
(A, I, \gamma) \ar[r] \ar[u] & (B', J', \delta') \ar[u]
}
$$
是除幂环范畴中的一个推出（推出由引理 \ref{dpa-lemma-colimits} 存在）。则
\begin{enumerate}
\item $B''/J'' = B/J \otimes_{A/I} B'/J'$,
\item 映射 $\varphi : B \otimes_A B' \to B''$ 是满射；
\item 由 $\varphi$ 诱导的映射
$J \otimes_A B' + B \otimes_A J' \to J''$ 是满射。
\end{enumerate}
为证明 (1)，考虑映射
$(B'', J'', \delta'') \to (C, (0), \emptyset)$，并应用推出的定义。
为证明 (2)，考虑像 $\tilde B = \varphi(B \otimes_A B')$ 以及理想
$\tilde J = \varphi(J \otimes_A B' + B \otimes_A J')$。
由于 $\delta''$ 与 $\delta$、$\delta'$ 相容，显然对所有 $n \geq 1$ 都有
$\delta''_n(\tilde J) \subset \tilde J$。
因此，由推出的泛性质必有 $\tilde B = B''$ 且 $\tilde J = J''$。
\end{remark}

\begin{remark}
\label{dpa-remark-forgetful}
忘却函子 $(A, I, \gamma) \mapsto A$ 不与余极限交换。例如，设
$$
\xymatrix{
(B, J, \delta) \ar[r] & (B'', J'', \delta'') \\
(A, I, \gamma) \ar[r] \ar[u] & (B', J', \delta') \ar[u]
}
$$
是评注 \ref{dpa-remark-pushout} 中所讨论的除幂环范畴中的推出。
那么一般而言，映射 $B \otimes_A B' \to B''$ 不是同构。
一个具体例子如下：
$(A, I, \gamma) = (\mathbf{Z}, (0), \emptyset)$,
$(B, J, \delta) = (\mathbf{Z}/4\mathbf{Z}, 2\mathbf{Z}/4\mathbf{Z}, \delta)$,
以及
$(B', J', \delta') =
(\mathbf{Z}/4\mathbf{Z}, 2\mathbf{Z}/4\mathbf{Z}, \delta')$
其中 $\delta_2(2) = 2$，而 $\delta'_2(2) = 0$。
更准确地说，利用引理 \ref{dpa-lemma-need-only-gamma-p}，分别令
$\delta$、$\delta'$ 为 $J$、$J'$ 上唯一的除幂结构，使得
$\delta_2 : J \to J$、$\delta'_2 : J' \to J'$ 分别是映射
$0 \mapsto 0, 2 \mapsto 2$ 和 $0 \mapsto 0, 2 \mapsto 0$。
于是 $(B'', J'', \delta'') = (\mathbf{F}_2, (0), \emptyset)$，
这与张量积并不一致。
\end{remark}


\section{除幂结构的延拓}
\label{dpa-section-extend}

\noindent
定义如下。

\begin{definition}
\label{dpa-definition-extends}
给定除幂环 $(A, I, \gamma)$ 和环同态 $A \to B$，若 $IB$ 上存在
除幂结构 $\bar \gamma$，使得
$(A, I, \gamma) \to (B, IB, \bar\gamma)$ 是除幂环同态，
则称 $\gamma$ 可{\it 延拓}到 $B$。
\end{definition}

\begin{lemma}
\label{dpa-lemma-gamma-extends}
设 $(A, I, \gamma)$ 是除幂环，$A \to B$ 是环同态。
若 $\gamma$ 可延拓到 $B$，则其延拓唯一。
假设下列条件中至少有一个成立：
\begin{enumerate}
\item $IB = 0$,
\item $I$ 是主理想；
\item $A \to B$ 是平坦的。
\end{enumerate}
则 $\gamma$ 可延拓到 $B$。
\end{lemma}

\begin{proof}
$IB$ 的任意元素都可写成有限和
$\sum\nolimits_{i=1}^t b_ix_i$，其中 $b_i \in B$、$x_i \in I$。
若 $\gamma$ 延拓为 $IB$ 上的 $\bar\gamma$，则
$\bar\gamma_n(x_i) = \gamma_n(x_i)$。
因此，定义 \ref{dpa-definition-divided-powers} 中的条件 (3) 和 (4) 蕴含
$$
\bar\gamma_n(\sum\nolimits_{i=1}^t b_ix_i) =
\sum\nolimits_{n_1 + \ldots + n_t = n}
\prod\nolimits_{i = 1}^t b_i^{n_i}\gamma_{n_i}(x_i)
$$
由此可见，$\bar\gamma$ 若存在便是唯一的。

\medskip\noindent
若 $IB = 0$，令 $\bar\gamma_n(0) = 0$ 即可。若 $I = (x)$，
则定义 $\bar\gamma_n(bx) = b^n\gamma_n(x)$。这是良定义的：
若 $b'x = bx$，即 $(b - b')x = 0$，则
\begin{align*}
b^n\gamma_n(x) - (b')^n\gamma_n(x)
& =
(b^n - (b')^n)\gamma_n(x) \\
& =
(b^{n - 1} + \ldots + (b')^{n - 1})(b - b')\gamma_n(x) = 0
\end{align*}
因为 $\gamma_n(x)$ 可被 $x$ 整除（这是由于
$\gamma_n(I) \subset I$），从而被 $b-b'$ 零化。
下面证明定义 \ref{dpa-definition-divided-powers} 的条件 (1)--(5)。
(1)、(2)、(3)、(5) 由构造显然成立。
对于 (4)，设 $y,z\in IB$，写成 $y=bx$、$z=cx$。
则 $y+z=(b+c)x$，因而
\begin{align*}
\bar\gamma_n(y + z)
& =
(b + c)^n\gamma_n(x) \\
& =
\sum \frac{n!}{i!(n - i)!}b^ic^{n -i}\gamma_n(x) \\
& =
\sum b^ic^{n - i}\gamma_i(x)\gamma_{n - i}(x) \\
& =
\sum \bar\gamma_i(y)\bar\gamma_{n -i}(z)
\end{align*}
正是所需结论。

\medskip\noindent
假设 $A \to B$ 是平坦的。设 $b_1, \ldots, b_r \in B$ 且
$x_1, \ldots, x_r \in I$。若 $\bar\gamma_n$ 存在，则
$$
\bar\gamma_n(\sum b_ix_i) =
\sum b_1^{e_1} \ldots b_r^{e_r} \gamma_{e_1}(x_1) \ldots \gamma_{e_r}(x_r)
$$
其中求和取遍 $e_1 + \ldots + e_r = n$。
再设有 $c_1, \ldots, c_s \in B$ 和 $a_{ij} \in A$，使得
$b_i = \sum a_{ij}c_j$。令 $y_j = \sum a_{ij}x_i$，我们断言
$$
\sum b_1^{e_1} \ldots b_r^{e_r} \gamma_{e_1}(x_1) \ldots \gamma_{e_r}(x_r) =
\sum c_1^{d_1} \ldots c_s^{d_s} \gamma_{d_1}(y_1) \ldots \gamma_{d_s}(y_s)
$$
在 $B$ 中成立，其中右端求和取遍 $d_1 + \ldots + d_s = n$。
事实上，利用除幂结构的公理，可将两端都展开成若干项之和；
其系数属于 $\mathbf{Z}[a_{ij}]$，各项形如
$c_1^{d_1} \ldots c_s^{d_s}\gamma_{e_1}(x_1) \ldots \gamma_{e_r}(x_r)$。
为说明系数相同，注意当
$\gamma$ 是 $(x_1, \ldots, x_r)$ 上唯一的除幂结构时，
该结论在 $\mathbf{Q}[x_1, \ldots, x_r, c_1, \ldots, c_s, a_{ij}]$
中成立。由 Lazard 定理（《代数》定理 \ref{algebra-theorem-lazard}），
可将 $B$ 写成有限自由 $A$-模的有向余极限。
特别地，若 $z\in IB$ 既写成 $z=\sum x_ib_i$，又写成
$z=\sum x'_{i'}b'_{i'}$，则可找到 $c_1,\ldots,c_s\in B$ 以及
$a_{ij},a'_{i'j}\in A$，使得 $b_i=\sum a_{ij}c_j$、
$b'_{i'}=\sum a'_{i'j}c_j$，并且
$y_j = \sum x_ia_{ij} = \sum x'_{i'}a'_{i'j}$\footnote{也可以不借助
《代数》定理 \ref{algebra-theorem-lazard} 来证明这一点。事实上，若
$z = \sum x_ib_i$ 且 $z = \sum x'_{i'}b'_{i'}$，则
$\sum x_ib_i - \sum x'_{i'}b'_{i'} = 0$ 是 $A$-模 $B$ 中的一个关系。
因此，《代数》引理 \ref{algebra-lemma-flat-eq}（以 $x_i$ 和
$x'_{i'}$ 代替其中的 $f_i$，并以 $b_i$ 和 $b'_{i'}$ 代替其中的
$x_i$）给出所需的 $c_1, \ldots, c_s \in B$ 以及
$a_{ij}, a'_{i'j} \in A$。}。
因此，上述过程给出 $IB$ 上一个良定义的映射 $\bar\gamma_n$。
由构造，$\bar\gamma$ 满足条件 (1)、(3)、(4)。此外，对 $x\in I$ 有
$\bar\gamma_n(x)=\gamma_n(x)$。故由引理
\ref{dpa-lemma-check-on-generators}，$\bar\gamma$ 是 $IB$ 上的除幂结构。
\end{proof}

\begin{lemma}
\label{dpa-lemma-kernel}
设 $(A, I, \gamma)$ 是除幂环。
\begin{enumerate}
\item 若 $\varphi : (A, I, \gamma) \to (B, J, \delta)$ 是除幂环同态，
则对所有 $n\geq 1$，$\Ker(\varphi) \cap I$ 在 $\gamma_n$ 下不变。
\item 设 $\mathfrak a \subset A$ 是理想，并令
$I' = I \cap \mathfrak a$。下列条件等价：
\begin{enumerate}
\item 对所有 $n > 0$，$I'$ 在 $\gamma_n$ 下不变；
\item $\gamma$ 可延拓到 $A/\mathfrak a$；
\item 存在 $I'$ 作为理想的一组生成元 $x_i$，使得对所有 $n>0$，
$\gamma_n(x_i) \in I'$。
\end{enumerate}
\end{enumerate}
\end{lemma}

\begin{proof}
证明 (1)。这是显然的。假设 (2)(a)。对 $x\in I$，定义
$\bar\gamma_n(x \bmod I') = \gamma_n(x) \bmod I'$。
这是良定义的，因为由定义 \ref{dpa-definition-divided-powers} (4) 以及
假设中的 $\gamma_j(y)\in I'$，对 $y\in I'$ 有
$\gamma_n(x+y)=\gamma_n(x)\bmod I'$。由于 $\gamma$ 是除幂结构，
$\bar\gamma$ 显然也是除幂结构。因此 (2)(b) 成立。
另一方面，由 (1)，(2)(b) 蕴含 (2)(a)。显然 (2)(a) 蕴含 (2)(c)。
现在假设 (2)(c)。注意当 $x=ax_i$ 时，
$\gamma_n(x)=a^n\gamma_n(x_i)\in I'$。
因此，对 $I'$ 作为阿贝尔群的一组生成元，均有
$\gamma_n(x)\in I'$。对由这些生成元组成的表达式长度作归纳，
只须证明：若 $\forall n : \gamma_n(x), \gamma_n(y) \in I'$，则
$\forall n : \gamma_n(x+y)\in I'$。这立即由除幂结构的第四条公理得到。
\end{proof}

\begin{lemma}
\label{dpa-lemma-sub-dp-ideal}
设 $(A, I, \gamma)$ 是除幂环，$E\subset I$ 是子集。
则在 $\gamma$ 下不变且包含所有 $f\in E$ 的最小理想 $J\subset I$，
就是由 $\gamma_n(f)$（$n\geq 1$、$f\in E$）生成的理想 $J$。
\end{lemma}

\begin{proof}
立即由引理 \ref{dpa-lemma-kernel} 得到。
\end{proof}

\begin{lemma}
\label{dpa-lemma-extend-to-completion}
设 $(A, I, \gamma)$ 是除幂环，$p$ 是素数。若 $p$ 在 $A/I$ 中幂零，则
\begin{enumerate}
\item $p$-进完备化 $A^\wedge = \lim_e A/p^eA$ 满射到 $A/I$；
\item 该映射的核是 $I$ 的 $p$-进完备化 $I^\wedge$；
\item 每个 $\gamma_n$ 关于 $p$-进拓扑连续，并延拓为
$\gamma_n^\wedge : I^\wedge \to I^\wedge$，从而在 $I^\wedge$ 上定义
一个除幂结构。
\end{enumerate}
若进一步假设 $A$ 是 $\mathbf{Z}_{(p)}$-代数，则
\begin{enumerate}
\item[(4)] 当 $e$ 足够大时，理想 $p^eA \subset I$ 在除幂结构
$\gamma$ 下不变，并且
$$
(A^\wedge, I^\wedge, \gamma^\wedge) = \lim_e (A/p^eA, I/p^eA, \bar\gamma)
$$
在除幂环范畴中成立。
\end{enumerate}
\end{lemma}

\begin{proof}
取整数 $t\geq 1$，使得 $p^tA/I=0$，即 $p^tA\subset I$。
映射 $A^\wedge\to A/I$ 是复合
$A^\wedge\to A/p^tA\to A/I$，它是满射（例如由
《代数》引理 \ref{algebra-lemma-completion-generalities}）。
对 $e\geq t$，由
$p^eI \subset p^eA \cap I \subset p^{e - t}I$ 可见，
复合 $A^\wedge\to A/I$ 的核是 $I$ 的 $p$-进完备化。
映射 $\gamma_n$ 连续，因为
$$
\gamma_n(x + p^ey) =
\sum\nolimits_{i + j = n} p^{je}\gamma_i(x)\gamma_j(y) =
\gamma_n(x) \bmod p^eI
$$
由除幂结构的公理即得。显然，映射 $\gamma_n^\wedge$ 从
$\gamma_n$ 继承除幂结构的各条公理。最后证明末一断言，取 $e>t$。
此时 $p^eA\subset I$ 且 $\gamma_1(p^eA)\subset p^eA$；对 $n>1$，有
$$
\gamma_n(p^ea) = p^n \gamma_n(p^{e - 1}a) = \frac{p^n}{n!} p^{n(e - 1)}a^n
\in p^e A
$$
因为 $p^n/n! \in \mathbf{Z}_{(p)}$，并且 $n\geq 2$、$e\geq 2$，
故 $n(e-1)\geq e$。
这证明 $\gamma$ 可延拓到 $A/p^eA$，参见引理 \ref{dpa-lemma-kernel}。
关于极限的断言由引理 \ref{dpa-lemma-limits} 的证明中对极限的构造显然成立。
\end{proof}




\section{除幂多项式代数}
\label{dpa-section-divided-power-polynomial-ring}

\noindent
一个非常有用的例子是{\it 除幂多项式代数}。
设 $A$ 是环，$t\geq 1$。用
$A\langle x_1, \ldots, x_t \rangle$ 表示如下 $A$-代数：
作为 $A$-模，令
$$
A\langle x_1, \ldots, x_t \rangle =
\bigoplus\nolimits_{n_1, \ldots, n_t \geq 0} A x_1^{[n_1]} \ldots x_t^{[n_t]}
$$
其乘法由下式给出：
$$
x_i^{[n]}x_i^{[m]} = \frac{(n + m)!}{n!m!}x_i^{[n + m]}.
$$
另令 $x_i=x_i^{[1]}$。注意
$1=x_1^{[0]}\ldots x_t^{[0]}$。类似的构造给出具有无限多个变量的
除幂多项式代数。存在典范 $A$-代数同态
$A\langle x_1, \ldots, x_t \rangle \to A$，它在 $n>0$ 时将
$x_i^{[n]}$ 映为零。该映射的核记为
$A\langle x_1, \ldots, x_t \rangle_{+}$。

\begin{lemma}
\label{dpa-lemma-divided-power-polynomial-algebra}
设 $(A,I,\gamma)$ 是除幂环。下列理想上存在唯一的除幂结构 $\delta$：
$$
J = IA\langle x_1, \ldots, x_t \rangle + A\langle x_1, \ldots, x_t \rangle_{+}
$$
使得
\begin{enumerate}
\item $\delta_n(x_i) = x_i^{[n]}$；
\item $(A, I, \gamma) \to (A\langle x_1, \ldots, x_t \rangle, J, \delta)$
是除幂环同态。
\end{enumerate}
此外，$(A\langle x_1, \ldots, x_t \rangle, J, \delta)$ 具有如下
泛性质：一个除幂环同态
$\varphi : (A\langle x_1, \ldots, x_t \rangle, J, \delta) \to
(C, K, \epsilon)$
与一个除幂环同态 $A\to C$ 以及元素
$k_1, \ldots, k_t \in K$ 等价。
\end{lemma}

\begin{proof}
我们在单变量除幂多项式代数的情形证明该引理。一般情形可完全类似地证明；
也可以注意到，$A\langle x_1, \ldots, x_t\rangle$ 同构于逐个添加
除幂变量 $x_1,\ldots,x_t$ 所得到的环。

\medskip\noindent
设 $A\langle x \rangle_{+}$ 是由
$x, x^{[2]}, x^{[3]}, \ldots$ 生成的理想。
注意 $J = IA\langle x \rangle + A\langle x \rangle_{+}$，并且
$$
IA\langle x \rangle \cap A\langle x \rangle_{+} =
IA\langle x \rangle \cdot A\langle x \rangle_{+}
$$
故由引理 \ref{dpa-lemma-two-ideals}，只须证明理想
$IA\langle x \rangle$ 和 $A\langle x \rangle_{+}$ 上存在除幂结构。
由于 $A\to A\langle x\rangle$ 平坦，前者的存在性由引理
\ref{dpa-lemma-gamma-extends} 得到。对于后者，注意若 $A$ 无挠，
则可应用引理 \ref{dpa-lemma-silly} (4) 得到 $\delta$ 的存在性。
事实上，选取元素 $x^{[m]}$ 作为生成元，可见
$(x^{[m]})^n = \frac{(nm)!}{(m!)^n} x^{[nm]}$
且 $n!$ 整除整数 $\frac{(nm)!}{(m!)^n}$。
一般而言，将 $A$ 写成某个无挠环 $R$ 的商 $R/\mathfrak a$
（例如可取 $\mathbf{Z}$ 上的多项式环）。映射
$R\langle x \rangle \to A\langle x \rangle$ 的核是
$\bigoplus \mathfrak a x^{[m]}$。应用引理 \ref{dpa-lemma-kernel}
的判据 (2)(c)，可见 $R\langle x \rangle_{+}$ 上的除幂结构
按所需方式延拓到 $A\langle x \rangle$。

\medskip\noindent
证明泛性质。给定除幂环同态 $\varphi:A\to C$ 以及
$k_1,\ldots,k_t\in K$，考虑
$$
A\langle x_1, \ldots, x_t \rangle \to C,\quad
x_1^{[n_1]} \ldots x_t^{[n_t]} \longmapsto
\epsilon_{n_1}(k_1) \ldots \epsilon_{n_t}(k_t)
$$
并在系数上使用 $\varphi$。只须检验这是 $A$-代数同态（细节从略）。
反向构造是显然的。
\end{proof}

\begin{remark}
\label{dpa-remark-divided-power-polynomial-algebra}
设 $(A,I,\gamma)$ 是除幂环。引理
\ref{dpa-lemma-divided-power-polynomial-algebra} 有一个无限多个变量的版本。
首先注意，若 $s<t$，则存在典范映射
$$
A\langle x_1, \ldots, x_s \rangle \to A\langle x_1, \ldots, x_t\rangle
$$
因此，对任意集合 $W$，令
$$
A\langle x_w: w \in W\rangle =
\colim_{E \subset W} A\langle x_e:e \in E\rangle
$$
（余极限取遍 $W$ 的有限子集 $E$），转移映射如上。
由余极限的定义可见，$A\langle x_w:w\in W\rangle$ 的泛映射性质
与引理 \ref{dpa-lemma-divided-power-polynomial-algebra} 中所述的映射性质
完全类似。
\end{remark}

\noindent
下述引理可见于 \cite{BO}。

\begin{lemma}
\label{dpa-lemma-need-only-gamma-p}
设 $p$ 是素数。设 $A$ 是环，并且每个不能被 $p$ 整除的整数 $n$
在 $A$ 中都可逆；也就是说，$A$ 是 $\mathbf{Z}_{(p)}$-代数。
设 $I\subset A$ 是理想。$I$ 上的两个除幂结构
$\gamma,\gamma'$ 相等，当且仅当 $\gamma_p=\gamma'_p$。
此外，给定映射 $\delta:I\to I$，满足
\begin{enumerate}
\item 对所有 $x\in I$，$p!\delta(x)=x^p$\footnote{此条件可由另外两个
条件推出，参见评注 \ref{dpa-remark-ryo-suzuki}。}；
\item 对所有 $a\in A$、$x\in I$，$\delta(ax)=a^p\delta(x)$；
\item
$\delta(x + y) =
\delta(x) +
\sum\nolimits_{i + j = p, i,j \geq 1} \frac{1}{i!j!} x^i y^j +
\delta(y)$ 对所有 $x,y\in I$ 成立，
\end{enumerate}
则 $I$ 上存在唯一的除幂结构 $\gamma$，使得 $\gamma_p=\delta$。
\end{lemma}

\begin{proof}
若 $n$ 不能被 $p$ 整除，则
$\gamma_n(x)=c x\gamma_{n-1}(x)$，其中 $c$ 是
$\mathbf{Z}_{(p)}$ 中的单位。此外，
$$
\gamma_{pm}(x) = c \gamma_m(\gamma_p(x))
$$
其中 $c$ 是 $\mathbf{Z}_{(p)}$ 中的单位。因此第一条断言显然。
对于第二条断言，可反向利用这些等式定义所有 $\gamma_n$。
更准确地说，若
$n=a_0+a_1p+\ldots+a_ep^e$，其中 $a_i\in\{0,\ldots,p-1\}$，则令
$$
\gamma_n(x) = c_n x^{a_0} \delta(x)^{a_1} \ldots \delta^e(x)^{a_e}
$$
其中 $c_n\in\mathbf{Z}_{(p)}$ 定义为
$$
c_n =
{(p!)^{a_1 + a_2(1 + p) + \ldots + a_e(1 + \ldots + p^{e - 1})}}/{n!}.
$$
现在须证明除幂结构的公理 (1)--(5)，参见定义
\ref{dpa-definition-divided-powers}。公理 (1) 和 (3) 是立即可得的；
公理 (2) 和 (5) 可直接计算验证，细节从略。设 $x,y\in I$。
我们断言存在环同态
$$
\varphi : \mathbf{Z}_{(p)}\langle u, v \rangle \longrightarrow A
$$
将 $u^{[n]}$ 映为 $\gamma_n(x)$，将 $v^{[n]}$ 映为 $\gamma_n(y)$。
由 $\mathbf{Z}_{(p)}\langle u,v\rangle$ 的构造，这意味着须检验
$$
\gamma_n(x)\gamma_m(x) = \frac{(n + m)!}{n!m!} \gamma_{n + m}(x)
$$
在 $A$ 中成立，并对 $y$ 也成立。由于 $\gamma$ 满足 (2)，这确实成立。
令 $\epsilon$ 表示 $\mathbf{Z}_{(p)}\langle u,v\rangle$ 的理想
$\mathbf{Z}_{(p)}\langle u,v\rangle_{+}$ 上的除幂结构。
接着断言：对 $f\in\mathbf{Z}_{(p)}\langle u,v\rangle_{+}$ 和所有 $n$，
$\varphi(\epsilon_n(f))=\gamma_n(\varphi(f))$。
当 $n=0,1,\ldots,p-1$ 时这是显然的。当 $n=p$ 时，只须对理想
$\mathbf{Z}_{(p)}\langle u,v\rangle_{+}$ 的一组生成元证明，
因为 $\epsilon_p$ 和 $\gamma_p=\delta$ 都满足该引理的条件 (1)、(3)。
因此，只须证明
$\gamma_p(\gamma_n(x))=\frac{(pn)!}{p!(n!)^p}\gamma_{pn}(x)$，
并对 $y$ 也证明同样的等式；这由 $\gamma$ 满足 (5) 得到。
现在，若任意整数 $n=a_0+a_1p+\ldots+a_ep^e$ 按上述方式写成
$p$-进展开，则
$$
\epsilon_n(f) =
c_n f^{a_0} \gamma_p(f)^{a_1} \ldots \gamma_p^e(f)^{a_e}
$$
因为 $\epsilon$ 是除幂结构。因此，对所有 $n$ 都有
$\varphi(\epsilon_n(f))=\gamma_n(\varphi(f))$。
将其应用于 $f=u+v$，由 $\epsilon$ 是除幂结构可知
$\gamma$ 满足公理 (4)。
\end{proof}

\begin{remark}
\label{dpa-remark-ryo-suzuki}
设 $A\supset I$ 如引理 \ref{dpa-lemma-need-only-gamma-p} 中所述，
并设映射 $\delta:I\to I$ 满足该引理的 (2)、(3)。则 (1) 也成立。
首先，对 $n$ 作归纳，证明当 $n\geq1$ 时
$\delta(nx)=n\delta(x)+\frac{n^p-n}{p!}x^p$。$n=1$ 的情形成立。
假设它对某个 $n$ 成立，则
$\delta((n + 1)x) = \delta(nx) + \delta(x) +
(\sum_{i = 1}^{p - 1} \frac{n^i}{i!(p - i)!})x^p
= (n + 1)\delta(x) + \left( \frac{n^p - n}{p!} + \frac{(n + 1)^p}{p!}
- \frac{n^p + 1}{p!} \right)x^p =
(n + 1)\delta(x) + \frac{(n + 1)^p - (n + 1)}{p!}x^p$,
故它对 $n+1$ 也成立。另一方面，$\delta(nx)=n^p\delta(x)$，所以
$(n^p - n)\delta(x) = \frac{n^p - n}{p!}x^p$. 
但对任意 $p$，都存在 $n$ 使得 $p^2 \not \mid n^p-n$。
事实上，
$(\mathbf{Z}/p^2\mathbf{Z})^\times \cong \mathbf{Z}/p(p-1)\mathbf{Z}$.
因此存在 $n$，使得 $n^{p-1}\not\equiv1\mod p^2$。
于是得到 $p!\delta(x)=x^p$。
\end{remark}









\section{Tate 分辨}
\label{dpa-section-tate}

\noindent
本节简要讨论 \cite{Tate-homology} 和 \cite{AH} 中构造的分辨；
这些分辨将除幂结构与微分分次代数结合起来。
本节对微分分次代数使用{\it 同调记号}。
我们的微分分次代数位于非负同调次数。因此，微分分次代数
$(A,\text{d})$ 由如下链复形给出：
$$
\ldots \to A_2 \to A_1 \to A_0 \to 0 \to \ldots
$$
并配备一个乘法。

\medskip\noindent
设 $R$ 是环（照例假设交换）。本节常考虑负次数分量为零的分次
$R$-代数 $A=\bigoplus_{d\geq0}A_d$。令
$A_+=\bigoplus_{d>0}A_d$，并记
$A_{even}=\bigoplus_{d\geq0}A_{2d}$、
$A_{odd}=\bigoplus_{d\geq0}A_{2d+1}$。
回顾：若对齐次元素 $x,y$ 有
$xy=(-1)^{\deg(x)\deg(y)}yx$，则称 $A$ 分次交换。
若进一步对奇次数齐次元素 $x$ 有 $x^2=0$，则称 $A$ 严格分次交换。
最后，为理解下面的定义，请记住：若 $A$ 是 $\mathbf{Q}$-代数，
则 $\gamma_n(x)=x^n/n!$。

\begin{definition}
\label{dpa-definition-divided-powers-graded}
设 $R$ 是环，$A=\bigoplus_{d\geq0}A_d$ 是严格分次交换的分次
$R$-代数。若对所有 $n>0$ 定义的一族映射
$\gamma_n:A_{even,+}\to A_{even,+}$ 满足
\begin{enumerate}
\item 若 $x\in A_{2d}$，则 $\gamma_n(x)\in A_{2nd}$；
\item 对任意 $x$，$\gamma_1(x)=x$，并令 $\gamma_0(x)=1$；
\item $\gamma_n(x)\gamma_m(x) = \frac{(n + m)!}{n! m!} \gamma_{n + m}(x)$,
\item 对所有 $x\in A_{even}$、$y\in A_{even,+}$，
$\gamma_n(xy)=x^n\gamma_n(y)$；
\item 若 $x,y\in A_{odd}$ 齐次且 $n>1$，则 $\gamma_n(xy)=0$；
\item 若 $x,y\in A_{even,+}$，则
$\gamma_n(x+y)=\sum_{i=0,\ldots,n}\gamma_i(x)\gamma_{n-i}(y)$；
\item $\gamma_n(\gamma_m(x)) =
\frac{(nm)!}{n! (m!)^n} \gamma_{nm}(x)$ 对 $x\in A_{even,+}$ 成立。
\end{enumerate}
则称其为 $A$ 上的{\it 除幂结构}。
\end{definition}

\noindent
注意，条件 (2)、(3)、(4)、(6)、(7) 蕴含 $\gamma$ 是交换环
$A_{even}$ 的理想 $A_{even,+}$ 上的“通常”除幂结构，参见第
\ref{dpa-section-divided-powers}、
\ref{dpa-section-divided-power-rings}、
\ref{dpa-section-extend} 和
\ref{dpa-section-divided-power-polynomial-ring} 节。
特别地，对所有 $x\in A_{even,+}$ 有 $n!\gamma_n(x)=x^n$。
条件 (1) 表明 $\gamma$ 与分次相容，条件 (5) 表明当 $n>1$ 时，
$\gamma_n$ 在奇次数齐次元素的乘积上为零。但请注意，仍有可能
$$
\gamma_2(z_1 z_2 + z_3 z_4) = z_1z_2z_3z_4
$$
非零，其中 $z_1,z_2,z_3,z_4$ 是奇次数齐次元素。

\begin{example}[添加奇次变量]
\label{dpa-example-adjoining-odd}
设 $R$ 是环。设 $(A,\gamma)$ 是配备上述除幂结构的严格分次交换
分次 $R$-代数，并设 $d>0$ 是奇数。在此情形可向 $A$ 添加一个
次数为 $d$ 的变量 $T$。具体地，令
$$
A\langle T \rangle = A \oplus AT
$$
其分次由 $A\langle T\rangle_m=A_m\oplus A_{m-d}T$ 给出。
我们断言 $A\langle T\rangle$ 上存在唯一的除幂结构，
与 $A$ 上给定的除幂结构相容。具体地，对
$x\in A_{even,+}$、$y\in A_{odd}$，令
$$
\gamma_n(x + yT) = \gamma_n(x) + \gamma_{n - 1}(x)yT
$$
\end{example}

\begin{example}[添加偶次变量]
\label{dpa-example-adjoining-even}
设 $R$ 是环。设 $(A,\gamma)$ 是配备上述除幂结构的严格分次交换
分次 $R$-代数，并设 $d>0$ 是偶数。在此情形可向 $A$ 添加一个
次数为 $d$ 的变量 $T$。具体地，令
$$
A\langle T \rangle = A \oplus AT \oplus AT^{(2)} \oplus AT^{(3)} \oplus \ldots
$$
其乘法由
$$
T^{(n)} T^{(m)} = \frac{(n + m)!}{n!m!} T^{(n + m)}
$$
给出，分次由
$$
A\langle T \rangle_m =
A_m \oplus A_{m - d}T \oplus A_{m - 2d}T^{(2)} \oplus \ldots
$$
给出。我们断言 $A\langle T\rangle$ 上存在唯一的除幂结构，
与 $A$ 上给定的除幂结构相容，并满足
$\gamma_n(T^{(i)})=T^{(ni)}$。为定义该除幂结构，先在
$x_i\in A_{even}$ 时令
$$
\gamma_n\left(\sum\nolimits_{i > 0} x_i T^{(i)}\right) =
\sum \prod\nolimits_{n = \sum e_i} x_i^{e_i} T^{(ie_i)}
$$
若 $x_0\in A_{even,+}$，则令
$$
\gamma_n\left(\sum\nolimits_{i \geq 0} x_i T^{(i)}\right) =
\sum\nolimits_{a + b = n}
\gamma_a(x_0)\gamma_b\left(\sum\nolimits_{i > 0} x_iT^{(i)}\right)
$$
其中 $\gamma_b$ 如上定义。
\end{example}

\begin{remark}
\label{dpa-remark-adjoining-set-of-variables}
也可以在给定次数 $d>0$ 中添加一族（可能无限多个）外代数生成元
或除幂生成元，而不只是像例 \ref{dpa-example-adjoining-odd} 和
\ref{dpa-example-adjoining-even} 那样添加一个。具体地，沿用评注
\ref{dpa-remark-divided-power-polynomial-algebra}：对如上的 $(A,\gamma)$
及集合 $J$，令 $A\langle T_j:j\in J\rangle$ 为代数
$A\langle T_j:j\in S\rangle$ 在 $J$ 的所有有限子集 $S$ 上的
有向余极限。显然，该代数具有唯一的除幂结构，并与 $A$ 上给定的
结构以及每个生成元 $T_j$ 上的结构相容。
\end{remark}

\noindent
现在把除幂结构的定义与微分联系起来。为理解下面的定义，注意若
$A$ 是 $\mathbf{Q}$-代数且 $x\in A_{even,+}$，则
$\text{d}(x^n/n!)=\text{d}(x)x^{n-1}/(n-1)!$。

\begin{definition}
\label{dpa-definition-divided-powers-dga}
设 $R$ 是环，$A=\bigoplus_{d\geq0}A_d$ 是严格分次交换的微分分次
$R$-代数。若对所有 $x\in A_{even,+}$ 有
$\text{d}(\gamma_n(x))=\text{d}(x)\gamma_{n-1}(x)$，
则称 $A$ 上的除幂结构 $\gamma$ {\it 与微分分次结构相容}。
\end{definition}

\noindent
警告：设 $(A,\text{d},\gamma)$ 如定义
\ref{dpa-definition-divided-powers-dga}。即使 $x$ 是边界，
$\gamma_n(x)$ 也未必是边界。因此，一般而言 $\gamma$ 不在同调代数
$H(A)$ 上诱导除幂结构。有些论文的作者会加入另一条相容性条件，
以保证确实诱导这样的结构；这里我们选择不这样做。

\begin{lemma}
\label{dpa-lemma-dpdga-good}
设 $(A,\text{d},\gamma)$ 和 $(B,\text{d},\gamma)$ 如定义
\ref{dpa-definition-divided-powers-dga}，并设 $f:A\to B$ 是与除幂结构
相容的微分分次代数态射。假设
\begin{enumerate}
\item 对 $k>0$，$H_k(A)=0$；
\item $f$ 是满射。
\end{enumerate}
则 $\gamma$ 在分次 $R$-代数 $H(B)$ 上诱导一个除幂结构。
\end{lemma}

\begin{proof}
设 $x$ 和 $x'$ 是次数同为 $2d$ 的齐次元素，并在 $H(B)$ 中定义
同一个上同调类。写成 $x'-x=\text{d}(w)$。选取 $x$ 的提升
$y\in A_{2d}$ 和 $w$ 的提升 $z\in A_{2d+1}$。
则 $y'=y+\text{d}(z)$ 是 $x'$ 的提升。因此
$$
\gamma_n(y') = \sum \gamma_i(y) \gamma_{n - i}(\text{d}(z))
= \gamma_n(y) +
\sum\nolimits_{i < n} \gamma_i(y) \gamma_{n - i}(\text{d}(z))
$$
由于 $A$ 在正次数无环，并且对所有 $j$ 有
$\text{d}(\gamma_j(\text{d}(z)))=0$，可将其写成
$$
\gamma_n(y') = \gamma_n(y) +
\sum\nolimits_{i < n} \gamma_i(y) \text{d}(z_i)
$$
其中 $z_i$ 是 $A$ 中的某些元素。此外，对 $0<i<n$ 有
$$
\text{d}(\gamma_i(y) z_i) =
\text{d}(\gamma_i(y))z_i + \gamma_i(y)\text{d}(z_i) =
\text{d}(y) \gamma_{i - 1}(y) z_i + \gamma_i(y)\text{d}(z_i)
$$
由于 $\text{d}(y)$ 在 $B$ 中的像为零，第一项在 $B$ 中映为零。
因此，$\gamma_n(x')$ 和 $\gamma_n(x)$ 在 $H(B)$ 中映为同一元素。
由此，对所有 $d>0$、$n>0$ 得到良定义的映射
$\gamma_n:H_{2d}(B)\to H_{2nd}(B)$。关于这些映射在 $H(B)$ 上
定义除幂结构的验证从略。
\end{proof}

\begin{lemma}
\label{dpa-lemma-base-change-div}
设 $(A,\text{d},\gamma)$ 如定义
\ref{dpa-definition-divided-powers-dga}，并设 $R\to R'$ 是环同态。
则 $\text{d}$ 和 $\gamma$ 在 $A'=A\otimes_RR'$ 上诱导同类结构，
使得 $(A',\text{d},\gamma)$ 也如定义
\ref{dpa-definition-divided-powers-dga} 所述。
\end{lemma}

\begin{proof}
注意 $A'_{even}=A_{even}\otimes_RR'$ 且
$A'_{even,+}=A_{even,+}\otimes_RR'$。因此要证明的是除幂结构
$\gamma$ 可延拓到 $A'_{even}$（术语见定义 \ref{dpa-definition-extends}）。
一旦证明 $\gamma$ 可延拓，便容易看出该延拓具有全部所需性质。

\medskip\noindent
选取一个多项式 $R$-代数 $P$（其生成元集合任意）以及一个满的
$R$-代数同态 $P\to R'$。环同态
$A_{even}\to A_{even}\otimes_RP$ 是平坦的，故由引理
\ref{dpa-lemma-gamma-extends}，除幂结构 $\gamma$ 唯一地延拓到
$A_{even}\otimes_RP$。令 $J=\Ker(P\to R')$。
为证明 $\gamma$ 可延拓到 $A\otimes_RR'$，只须证明
$I'=\Ker(A_{even,+}\otimes_RP\to A_{even,+}\otimes_RR')$
由满足 $\gamma_n(z)\in I'$（对所有 $n>0$）的元素 $z$ 生成。
这是显然的，因为 $I'$ 由形如 $x\otimes f$ 的元素生成，其中
$x\in A_{even,+}$ 且 $f\in\Ker(P\to R')$。
\end{proof}

\begin{lemma}
\label{dpa-lemma-extend-differential}
设 $(A,\text{d},\gamma)$ 如定义
\ref{dpa-definition-divided-powers-dga}，并设 $d\geq1$ 是整数。
设 $A\langle T\rangle$ 是例 \ref{dpa-example-adjoining-odd} 或
\ref{dpa-example-adjoining-even} 中构造的、以 $T$ 为变量且
$\deg(T)=d$ 的分次除幂多项式代数。设 $f\in A_{d-1}$ 且
$\text{d}(f)=0$。则 $A\langle T\rangle$ 上存在唯一的微分
$\text{d}$，使得 $\text{d}(T)=f$，并且 $\text{d}$ 与
$A\langle T\rangle$ 上的除幂结构相容。
\end{lemma}

\begin{proof}
可由直接计算证明，细节从略。
\end{proof}

\noindent
在引理 \ref{dpa-lemma-compute-cohomology-adjoin-variable} 中，
我们将在若干特殊情形计算 $A\langle T\rangle$ 的上同调。
下面给出 Tate 的构造及 Avramov 和 Halperin 对它的推广。

\begin{lemma}
\label{dpa-lemma-tate-resolution}
设 $R\to S$ 是交换环同态。存在分解
$$
R \to A \to S
$$
具有如下性质：
\begin{enumerate}
\item $(A,\text{d},\gamma)$ 如定义
\ref{dpa-definition-divided-powers-dga} 所述；
\item 若给 $S$ 配备零微分，则 $A\to S$ 是拟同构；
\item $A_0=R[x_j:j\in J]\to S$ 是任意一个从多项式环到 $S$ 的满射；
\item $A$ 是 $R$ 上的分次除幂多项式代数。
\end{enumerate}
最后一个条件是指：从 $A_0$ 出发，像例 \ref{dpa-example-adjoining-odd}
或 \ref{dpa-example-adjoining-even} 那样，在每个正次数逐次添加一族变量
$T$，从而构造出 $A$。此外，若 $R$ 是诺特环且 $R\to S$ 为有限型，
则可以选取 $A$，使其在每个次数只有有限多个生成元。
\end{lemma}

\begin{proof}
我们写出 $R$ 为诺特环且 $R\to S$ 为有限型时的构造。
没有这些假设时，证明相同，只不过每个次数须使用一族
（可能无限多个）生成元。

\medskip\noindent
构造的起点：令 $A(0)=R[x_1,\ldots,x_n]$ 为通常的多项式环，
并设 $A(0)\to S$ 是满射。取分次
$A(0)_0=A(0)$，且当 $d\not=0$ 时 $A(0)_d=0$。
因此 $\text{d}=0$，并且当 $n>0$ 时 $\gamma_n$ 也为零。

\medskip\noindent
选取给定映射 $A(0)=R[x_1,\ldots,x_n]\to S$ 的核的一组生成元
$f_1,\ldots,f_m\in R[x_1,\ldots,x_n]$。
应用例 \ref{dpa-example-adjoining-odd} 共 $m$ 次，得到
$$
A(1) = A(0)\langle T_1, \ldots, T_m\rangle
$$
作为分次除幂多项式代数，其中 $\deg(T_i)=1$。令
$\text{d}(T_i)=f_i$。由于 $A(1)$ 是 $A(0)$ 上的除幂多项式代数，
且 $\text{d}(f_i)=0$，由引理 \ref{dpa-lemma-extend-differential}，
这唯一地延拓为 $A(1)$ 上的微分。

\medskip\noindent
归纳假设：设已给出分解
$$
R \to A(0) \to A(1) \to \ldots \to A(m) \to S
$$
其中 $A(0)$、$A(1)$ 如上；对 $2\leq m'\leq m$，每个
$R\to A(m')\to S$ 都满足引理陈述中的性质 (1)、(4)，
而将 (2) 换成条件：当 $m'>i\geq0$ 时，
$H_i(A(m'))\to H_i(S)$ 是同构。基础情形是 $m=1$。

\medskip\noindent
归纳步骤：假设已有归纳假设中的 $R\to A(m)\to S$。
考虑群 $H_m(A(m))$，它是 $H_0(A(m))=S$ 上的模。
事实上，它是 $A(m)_m$ 的一个子商，而 $A(m)_m$ 是
$A(m)_0=R[x_1,\ldots,x_n]$ 上的有限型模。因此可选取有限多个元素
$$
e_1, \ldots, e_t \in \Ker(\text{d} : A(m)_m \to A(m)_{m - 1})
$$
其像生成该模。应用例 \ref{dpa-example-adjoining-odd} 或
\ref{dpa-example-adjoining-even} 共 $t$ 次，得到
$$
A(m + 1) = A(m)\langle T_1, \ldots, T_t\rangle
$$
作为分次除幂代数，其中 $\deg(T_i)=m+1$。令
$\text{d}(T_i)=e_i$。由于 $A(m+1)$ 是 $A(m)$ 上的除幂多项式代数，
且 $\text{d}(e_i)=0$，这唯一地延拓为 $A(m+1)$ 上与除幂结构相容的
微分。由于只在次数 $m+1$ 及更高次数添加了材料，故当 $i<m$ 时
$H_i(A(m+1))=H_i(A(m))$。此外，由构造显然有
$H_m(A(m+1))=0$。

\medskip\noindent
为完成证明，注意我们已经证明存在映射序列
$$
R \to A(0) \to A(1) \to \ldots \to A(m) \to A(m + 1) \to \ldots \to S
$$
令 $A=\colim A(m)$ 即完成证明。
\end{proof}

\begin{lemma}
\label{dpa-lemma-tate-resoluton-pseudo-coherent-ring-map}
设 $R\to S$ 是伪凝聚环同态（《代数详论》定义
\ref{more-algebra-definition-pseudo-coherent-perfect}）。则引理
\ref{dpa-lemma-tate-resolution} 成立，并且可使 $S$ 的分辨 $A$
在每个次数只有有限多个生成元。
\end{lemma}

\begin{proof}
证明与引理 \ref{dpa-lemma-tate-resolution} 完全相同。唯一额外之处是：
给定 $A(m)\to S$，须证明 $H_m=H_m(A(m))$ 是有限
$R[x_1,\ldots,x_m]$-模（从而下一步只须添加有限多个变量）。
考虑复形
$$
\ldots \to A(m)_{m - 1} \to A(m)_m \to A(m)_{m - 1} \to
\ldots \to A(m)_0 \to S \to 0
$$
由于 $S$ 是伪凝聚 $R[x_1,\ldots,x_n]$-模，且 $A(m)_i$ 是有限自由
$R[x_1,\ldots,x_n]$-模，故该复形是伪凝聚复形，参见
《代数详论》引理 \ref{more-algebra-lemma-complex-pseudo-coherent-modules}。
由于该复形在同调次数 $>m$ 正合，由《代数详论》引理
\ref{more-algebra-lemma-finite-cohomology}，$H_m$ 是有限 $R$-模。
\end{proof}

\begin{lemma}
\label{dpa-lemma-uniqueness-tate-resolution}
设 $R$ 是交换环。假设 $(A,\text{d},\gamma)$ 和
$(B,\text{d},\gamma)$ 如定义 \ref{dpa-definition-divided-powers-dga}。
设 $\overline{\varphi}:H_0(A)\to H_0(B)$ 是 $R$-代数同态。假设
\begin{enumerate}
\item $A$ 是 $R$ 上的分次除幂多项式代数；
\item 对 $k>0$，$H_k(B)=0$。
\end{enumerate}
则存在与除幂结构相容的微分分次 $R$-代数态射
$\varphi:A\to B$，它提升 $\overline{\varphi}$。
\end{lemma}

\begin{proof}
该假设意味着：从 $R$ 出发，依次添加次数为零的某族多项式生成元、
正奇次数的外代数生成元和正偶次数的除幂生成元，便得到 $A$。
因此 $A$ 具有滤过 $R\subset A(0)\subset A(1)\subset\ldots$，
其中 $A(m+1)$ 由 $A(m)$ 添加次数为 $m+1$ 的适当类型生成元而得
（我们将它们统称为“除幂生成元”）。特别地，$A(0)\to H_0(A)$
是从 $R$ 上通常的多项式代数到 $H_0(A)$ 的满射。因此可将
$\overline{\varphi}$ 提升为 $R$-代数同态
$\varphi(0):A(0)\to B_0$。

\medskip\noindent
将 $A(1)$ 写成 $A(0)\langle T_j:j\in J\rangle$，其中 $J$ 是一族
次数为 $1$ 的除幂变量 $T_j$ 的指标集。令
$f_j=\varphi(0)(\text{d}(T_j))\in B_0$。由于 $\text{d}T_j$
在 $H_0(A)$ 中映为零，$f_j$ 在 $H_0(B)$ 中也映为零。
因此可找到 $b_j\in B_1$，使得 $\text{d}(b_j)=f_j$。
由引理 \ref{dpa-lemma-divided-power-polynomial-algebra} 中除幂多项式代数的
泛性质，可找到 $\varphi(0)$ 的提升 $\varphi(1):A(1)\to B$，
将 $T_j$ 映为 $f_j$。

\medskip\noindent
对某个 $m\geq1$ 构造出 $\varphi(m)$ 后，可完全以相同方式构造
$\varphi(m+1):A(m+1)\to B$。细节从略。
\end{proof}

\begin{lemma}
\label{dpa-lemma-divided-powers-on-tor}
设 $R$ 是交换环，$S$、$T$ 是交换 $R$-代数。则
$$
\operatorname{Tor}_*^R(S, T).
$$
上具有典范的、配备除幂结构的严格分次交换 $R$-代数结构。
\end{lemma}

\begin{proof}
选取如上的分解 $R\to A\to S$。由于 $A\to S$ 是拟同构，且
$A_d$ 是自由 $R$-模，微分分次代数 $B=A\otimes_RT$ 计算引理中
所示的 Tor 群。选取满射 $R[y_j:j\in J]\to T$。
于是 $B$ 是微分分次代数 $A[y_j:j\in J]$ 的商；后者的同调集中在
次数 $0$（等于 $S[y_j:j\in J]$）。由引理
\ref{dpa-lemma-base-change-div}，微分分次代数 $B$ 和
$A[y_j:j\in J]$ 都具有与微分相容的除幂结构。因此，由引理
\ref{dpa-lemma-dpdga-good}，在 $H(B)$ 上得到所需的除幂结构。

\medskip\noindent
以此方式构造的除幂代数结构与 $A$ 的选择无关。事实上，若 $A'$ 是
另一种选择，则引理 \ref{dpa-lemma-uniqueness-tate-resolution} 蕴含存在
保持全部结构以及到 $S$ 的增广的态射 $A\to A'$。
于是诱导态射
$B=A\otimes_RT\to A'\otimes_RT'=B'$ 也保持全部结构并且是拟同构。
因此，诱导的 Tor 代数同构与乘积和除幂相容。
\end{proof}





\section{在完全交中的应用}
\label{dpa-section-application-ci}

\noindent
设 $R$ 是环，$(A,\text{d},\gamma)$ 如定义
\ref{dpa-definition-divided-powers-dga}。次数为 $2$ 的{\it 导子}是具有
下列性质的 $R$-线性映射 $\theta:A\to A$：
\begin{enumerate}
\item $\theta(A_d) \subset A_{d - 2}$,
\item $\theta(xy) = \theta(x)y + x\theta(y)$,
\item $\theta$ 与 $\text{d}$ 交换；
\item $\theta(\gamma_n(x)) = \theta(x) \gamma_{n - 1}(x)$
对所有 $d$ 及所有 $x\in A_{2d}$ 成立。
\end{enumerate}
下述引理构造一个导子。

\begin{lemma}
\label{dpa-lemma-get-derivation}
设 $R$ 是环，$(A,\text{d},\gamma)$ 如定义
\ref{dpa-definition-divided-powers-dga}。设 $R'\to R$ 是环的满射，
其核平方为零且由一个元素 $f$ 生成。若 $A$ 是 $R$ 上的分次
除幂多项式代数，并且每个次数只有有限多个变量，则得到导子
$\theta:A/IA\to A/IA$，其中 $I$ 是 $f$ 在 $R$ 中的零化子。
\end{lemma}

\begin{proof}
由于 $A$ 是除幂多项式代数，可找到 $R'$ 上的除幂多项式代数
$A'$，使得 $A=A'\otimes_{R'}R$。此外，可以把 $\text{d}$ 提升为
$A'$ 上的 $R$-线性算子 $\text{d}$，使得
\begin{enumerate}
\item  $\text{d}(xy) = \text{d}(x)y + (-1)^{\deg(x)}x \text{d}(y)$
对 $A'$ 中的齐次元素 $x,y$ 成立；
\item  $\text{d}(\gamma_n(x)) = \text{d}(x) \gamma_{n - 1}(x)$ for
$x \in A'_{even, +}$。
\end{enumerate}
细节从略（提示：每次处理一个变量）。不过，$\text{d}^2$ 在 $A'$
上未必为零。显然，$\text{d}^2$ 将 $A'$ 映入
$fA'\cong A/IA$。因此 $\text{d}^2$ 零化 $fA'$，并分解为映射
$A\to A/IA$。由于 $\text{d}^2$ 是 $R$-线性的，得到所需映射
$\theta:A/IA\to A/IA$。导子各项性质的验证是立即可得的。
\end{proof}

\begin{lemma}
\label{dpa-lemma-compute-theta}
假设和记号如引理 \ref{dpa-lemma-get-derivation}。假设 $S=H_0(A)$ 同构于
$R[x_1, \ldots, x_n]/(f_1, \ldots, f_m)$
，其中 $n,m$ 为某些整数且 $f_j\in R[x_1,\ldots,x_n]$。
再假设给定关系
$$
\sum r_j f_j = 0
$$
其中 $r_j\in R[x_1,\ldots,x_n]$。选取 $r_j,f_j$ 的提升
$r'_j,f'_j\in R'[x_1,\ldots,x_n]$。将
$\sum r'_jf'_j$ 写成 $gf$，其中 $g\in R/I[x_1,\ldots,x_n]$。
若 $H_1(A)=0$ 且每个 $r_j$ 的所有系数都属于 $I$，则存在元素
$\xi\in H_2(A/IA)$，使得在 $S/IS$ 中 $\theta(\xi)=g$。
\end{lemma}

\begin{proof}
设 $A(0)\subset A(1)\subset A(2)\subset\ldots$ 是 $A$ 的滤过，
使得 $A(m)$ 由 $A(m-1)$ 添加次数为 $m$ 的除幂变量而得。
则 $A(0)$ 是 $R$ 上的多项式代数，并配备一个 $R$-代数满射
$A(0)\to S$。因此可以选取提升到 $S$ 的增广的映射
$$
\varphi : R[x_1, \ldots, x_n] \to A(0)
$$
接着，对某些次数为 $1$ 的除幂变量 $T_i$，有
$A(1)=A(0)\langle T_1,\ldots,T_t\rangle$。由于 $H_0(A)=S$，
可选取 $\xi_j\in\sum A(0)T_i$，使得
$\text{d}(\xi_j)=\varphi(f_j)$。于是
$$
\text{d}\left(\sum \varphi(r_j) \xi_j\right) =
\sum  \varphi(r_j) \varphi(f_j) =  \sum \varphi(r_jf_j) = 0
$$
由于 $H_1(A)=0$，可选取 $\xi\in A_2$，使得
$\text{d}(\xi)=\sum\varphi(r_j)\xi_j$。若 $r_j$ 的系数属于 $I$，
则 $\varphi(r_j)$ 也如此。此时 $\text{d}(\xi)$ 在 $A_1/IA_1$
中消失，因而 $\xi$ 在 $H_2(A/IA)$ 中定义一个类。

\medskip\noindent
引理 \ref{dpa-lemma-get-derivation} 的证明中对 $\theta$ 的构造，是逐次将
$A(i)$ 提升到 $A'(i)$ 并提升微分 $\text{d}$。将 $\varphi$ 提升为
$\varphi':R'[x_1,\ldots,x_n]\to A'(0)$。接着，
$A'(1)=A'(0)\langle T_1,\ldots,T_t\rangle$。此外，可将 $\xi_j$
提升为 $\xi'_j\in\sum A'(0)T_i$。于是，对某个 $a_j\in A'(0)$，
$\text{d}(\xi'_j)=\varphi'(f'_j)+fa_j$。
考虑 $\xi$ 的一个提升 $\xi'\in A'_2$。则
$$
\text{d}(\xi') = \sum \varphi'(r'_j)\xi'_j + \sum fb_iT_i
$$
其中 $b_i\in A(0)$。再次应用 $\text{d}$，得到
$$
\theta(\xi) = \sum \varphi'(r'_j)\varphi'(f'_j) +
\sum f \varphi'(r'_j) a_j + \sum fb_i \text{d}(T_i)
$$
第一项给出所需结果。第二项为零，因为 $r_j$ 的系数属于 $I$，
从而被 $f$ 零化。第三项在 $H_0$ 中映为零，因为
$\text{d}(T_i)$ 映为零。
\end{proof}

\noindent
下述引理的证明方法似乎源自 Gulliksen。

\begin{lemma}
\label{dpa-lemma-not-finite-pd}
设 $R'\to R$ 是诺特环的满射，其核平方为零且由一个元素 $f$ 生成。令
$S = R[x_1, \ldots, x_n]/(f_1, \ldots, f_m)$.
设 $\sum r_jf_j=0$ 是 $R[x_1,\ldots,x_n]$ 中的一个关系。假设
\begin{enumerate}
\item 每个 $r_j$ 的系数都属于 $f$ 在 $R$ 中的零化子 $I$；
\item 对某些提升 $r'_j,f'_j\in R'[x_1,\ldots,x_n]$，有
$\sum r'_jf'_j=gf$，其中 $g$ 在 $S/IS$ 中不幂零。
\end{enumerate}
则 $S$ 在 $R$ 上不具有有限 Tor 维数（即 $S$ 不是完美 $R$-代数）。
\end{lemma}

\begin{proof}
选取引理 \ref{dpa-lemma-tate-resolution} 中的 Tate 分辨
$R\to A\to S$。设 $\xi\in H_2(A/IA)$ 和
$\theta:A/IA\to A/IA$ 是引理 \ref{dpa-lemma-get-derivation} 和
\ref{dpa-lemma-compute-theta} 中得到的元素和导子。注意
$$
\theta^n(\gamma_n(\xi)) = g^n
$$
在 $H_0(A/IA)=S/IS$ 中成立。因此，若 $g$ 在 $S/IS$ 中不幂零，
则对所有 $n>0$，$\xi^n$ 在 $H_{2n}(A/IA)$ 中非零。由于
$H_{2n}(A/IA)=\text{Tor}^R_{2n}(S,R/I)$，结论成立。
\end{proof}

\noindent
下述结果可见于 \cite{Rodicio}。

\begin{lemma}
\label{dpa-lemma-injective}
设 $(A,\mathfrak m)$ 是诺特局部环，$I\subset J\subset A$ 是真理想。
若 $A/J$ 在 $A/I$ 上具有有限 Tor 维数，则
$I/\mathfrak mI\to J/\mathfrak mJ$ 是单射。
\end{lemma}

\begin{proof}
设 $f\in I$ 在 $I/\mathfrak mI$ 中的像非零，而在
$J/\mathfrak mJ$ 中映为零。可以选取理想 $I'$，满足
$\mathfrak mI\subset I'\subset I$，使得 $I/I'$ 由 $f$ 的像生成。
令 $R=A/I$、$R'=A/I'$。设 $J=(a_1,\ldots,a_m)$，其中
$a_j\in A$。则对某些 $b_j\in\mathfrak m$ 有
$f=\sum b_ja_j$。分别令 $r_j,f_j\in R$ 和
$r'_j,f'_j\in R'$ 为 $b_j,a_j$ 的像。于是处于引理
\ref{dpa-lemma-not-finite-pd} 的情形（其中该引理的理想 $I$ 等于
$\mathfrak m_R$），引理得证。
\end{proof}

\begin{lemma}
\label{dpa-lemma-regular-sequence}
设 $(A,\mathfrak m)$ 是诺特局部环，$I\subset J\subset A$ 是真理想。
假设
\begin{enumerate}
\item $A/J$ 在 $A/I$ 上具有有限 Tor 维数；
\item $J$ 由正则序列生成。
\end{enumerate}
则 $I$ 由正则序列生成，且 $J/I$ 由正则序列生成。
\end{lemma}

\begin{proof}
由引理 \ref{dpa-lemma-injective}，映射
$I/\mathfrak mI\to J/\mathfrak mJ$ 是单射。因此可找到 $s\leq r$
以及 $J$ 的一个极小生成元系 $f_1,\ldots,f_r$，使得
$f_1,\ldots,f_s$ 属于 $I$ 并构成 $I$ 的极小生成元系。
由《代数详论》引理
\ref{more-algebra-lemma-independence-of-generators} 和
\ref{more-algebra-lemma-noetherian-finite-all-equivalent}.
，$J$ 的任意极小生成元系都是正则序列，故引理成立。
\end{proof}

\begin{lemma}
\label{dpa-lemma-perfect-map-ci}
设 $R\to S$ 是诺特局部环之间的局部环同态。
设 $I\subset R$、$J\subset S$ 是理想且 $IS\subset J$。
若 $R\to S$ 平坦且 $S/\mathfrak m_RS$ 正则，则下列条件等价：
\begin{enumerate}
\item $J$ 由正则序列生成，且 $S/J$ 作为 $R/I$-模具有有限 Tor 维数；
\item $J$ 由正则序列生成，且只有有限多个 $p$ 使
$\text{Tor}^{R/I}_p(S/J,R/\mathfrak m_R)$ 非零；
\item $I$ 由正则序列生成，且 $J/IS$ 在 $S/IS$ 中由正则序列生成。
\end{enumerate}
\end{lemma}

\begin{proof}
若 (3) 成立，则 $J$ 由正则序列生成，例如参见《代数详论》引理
\ref{more-algebra-lemma-join-koszul-regular-sequences} 和
\ref{more-algebra-lemma-noetherian-finite-all-equivalent}.
此外，若 (3) 成立，则 $S/J=(S/I)/(J/I)$ 在 $S/IS$ 上具有有限
投射维数，因为 Koszul 复形是 $S/J$ 在 $S/IS$ 上的有限自由分辨。
由于 $R/I\to S/IS$ 平坦，由《代数详论》引理
\ref{more-algebra-lemma-flat-push-tor-amplitude}，$S/J$ 在 $R/I$
上具有有限 Tor 维数。因此 (3) 蕴含 (1)。

\medskip\noindent
(1) $\Rightarrow$ (2) 是平凡的。假设 (2)。由《代数详论》引理
\ref{more-algebra-lemma-perfect-over-regular-local-ring}，$S/J$ 在
$S/IS$ 上具有有限 Tor 维数。因此可应用引理
\ref{dpa-lemma-regular-sequence}，得到 $IS$ 和 $J/IS$ 都由正则序列生成。
设 $f_1,\ldots,f_r\in I$ 是 $I$ 的极小生成元系。由于 $R\to S$
平坦，$f_1,\ldots,f_r$ 在 $S$ 中构成 $IS$ 的极小生成元系。
因此，由《代数详论》引理
\ref{more-algebra-lemma-independence-of-generators} 和
\ref{more-algebra-lemma-noetherian-finite-all-equivalent}.
，元素序列 $f_1,\ldots,f_r\in R$ 的像在 $S$ 中构成正则序列。
故由《代数》引理 \ref{algebra-lemma-flat-increases-depth}，
$f_1,\ldots,f_r$ 在 $R$ 中是正则序列。
\end{proof}





\section{局部完全交环}
\label{dpa-section-lci}

\noindent
设 $(A,\mathfrak m)$ 是诺特完备局部环。由 Cohen 结构定理
（参见《代数》定理 \ref{algebra-theorem-cohen-structure-theorem}），
可将 $A$ 写成某个正则诺特完备局部环 $R$ 的商。
若存在某个满射 $R\to A$，其中 $R$ 是正则局部环且其核由正则序列
生成，则称 $A$ 是{\it 完全交}。下述引理表明该概念与满射的选择无关。

\begin{lemma}
\label{dpa-lemma-ci-well-defined}
设 $(A,\mathfrak m)$ 是诺特完备局部环。下列条件等价：
\begin{enumerate}
\item 对每个局部环满射 $R\to A$，只要 $R$ 是正则局部环，
$R\to A$ 的核便由正则序列生成；
\item 对某个局部环满射 $R\to A$，$R$ 是正则局部环，
且 $R\to A$ 的核由正则序列生成。
\end{enumerate}
\end{lemma}

\begin{proof}
设 $k$ 是 $A$ 的剩余域。若 $k$ 的特征为 $p>0$，则用 $\Lambda$
表示一个剩余域为 $k$ 的 Cohen 环（《代数》定义
\ref{algebra-definition-cohen-ring} 和引理
\ref{algebra-lemma-cohen-rings-exist}）。若 $k$ 的特征为 $0$，
令 $\Lambda=k$。回顾：对任意 $n$，
$\Lambda[[x_1,\ldots,x_n]]$ 在 $\mathfrak m$-进拓扑下分别在
$\mathbf{Z}$ 或 $\mathbf{Q}$ 上形式光滑，参见《代数详论》引理
\ref{more-algebra-lemma-power-series-ring-over-Cohen-fs}。
固定 Cohen 结构定理（《代数》定理
\ref{algebra-theorem-cohen-structure-theorem}）中的一个满射
$\Lambda[[x_1,\ldots,x_n]]\to A$。

\medskip\noindent
设 $R\to A$ 是从正则局部环 $R$ 出发的满射，并设
$f_1,\ldots,f_r$ 是 $\Ker(R\to A)$ 的极小生成元序列。
以下将反复使用而不再说明：诺特局部环中的理想由正则序列生成，
当且仅当其任意极小生成元集都是正则序列。由《代数》引理
\ref{algebra-lemma-flat-increases-depth} 和
\ref{algebra-lemma-completion-flat}，$f_1,\ldots,f_r$ 在 $R$ 中为
正则序列，当且仅当它们在完备化 $R^\wedge$ 中为正则序列。此外，
$$
R^\wedge/(f_1, \ldots, f_r)R^\wedge =
(R/(f_1, \ldots, f_n))^\wedge = A^\wedge = A
$$
因为 $A$ 是 $\mathfrak m_A$-进完备的（第一个等式由《代数》引理
\ref{algebra-lemma-completion-tensor} 得到）。最后，由于 $R$ 正则，
$R^\wedge$ 也正则（《代数详论》引理
\ref{more-algebra-lemma-completion-regular}）。故可假设 $R$ 完备。

\medskip\noindent
若 $R$ 完备，则可选取映射
$\Lambda[[x_1,\ldots,x_n]]\to R$，提升给定映射
$\Lambda[[x_1,\ldots,x_n]]\to A$，参见《代数详论》引理
\ref{more-algebra-lemma-lift-continuous}。再添加一些变量
$y_1,\ldots,y_m$，将其映到 $R\to A$ 的核的生成元，便可假设
$\Lambda[[x_1,\ldots,x_n,y_1,\ldots,y_m]]\to R$ 是满射
（略去若干细节）。于是可考虑交换图
$$
\xymatrix{
\Lambda[[x_1, \ldots, x_n, y_1, \ldots, y_m]] \ar[r] \ar[d] & R \ar[d] \\
\Lambda[[x_1, \ldots, x_n]] \ar[r] & A
}
$$
由《代数》引理 \ref{algebra-lemma-ci-well-defined}，关于 $R\to A$
的条件等价于关于固定映射
$\Lambda[[x_1,\ldots,x_n]]\to A$ 的条件。引理得证。
\end{proof}

\noindent
下面两个引理是对上述定义的合理性检验。

\begin{lemma}
\label{dpa-lemma-quotient-regular-ring-by-regular-sequence}
设 $R$ 是正则环，$\mathfrak p\subset R$ 是素理想，且
$f_1,\ldots,f_r\in\mathfrak p$ 是正则序列。则
$$
A = (R/(f_1, \ldots, f_r))_\mathfrak p =
R_\mathfrak p/(f_1, \ldots, f_r)R_\mathfrak p
$$
的完备化是上述意义下的完全交。
\end{lemma}

\begin{proof}
由诺特环 $R_\mathfrak p$ 上有限模的完备化是正合的
（《代数》引理 \ref{algebra-lemma-completion-tensor}），$A$ 的完备化为
$A^\wedge = R_\mathfrak p^\wedge/(f_1, \ldots, f_r)R_\mathfrak p^\wedge$
。由《代数》引理 \ref{algebra-lemma-completion-flat} 和
\ref{algebra-lemma-flat-increases-depth}，序列 $f_1,\ldots,f_r$ 在
$R_\mathfrak p$ 中的像是正则序列。此外，由《代数详论》引理
\ref{more-algebra-lemma-completion-regular}，$R_\mathfrak p^\wedge$
是正则局部环。因此结论由完备局部环的完全交定义成立。
\end{proof}

\noindent
下述引理是《代数》引理 \ref{algebra-lemma-lci} 的对应版本。

\begin{lemma}
\label{dpa-lemma-quotient-regular-ring}
设 $R$ 是正则环，$\mathfrak p\subset R$ 是素理想，
$I\subset\mathfrak p$ 是理想。令
$A=(R/I)_\mathfrak p=R_\mathfrak p/I_\mathfrak p$。下列条件等价：
\begin{enumerate}
\item $A$ 的完备化是上述意义下的完全交；
\item $I_\mathfrak p\subset R_\mathfrak p$ 由正则序列生成；
\item 模 $(I/I^2)_\mathfrak p$ 可由
$\dim(R_\mathfrak p)-\dim(A)$ 个元素生成；
\item 此处再添加内容。
\end{enumerate}
\end{lemma}

\begin{proof}
可以并确实将 $R$ 替换为它在 $\mathfrak p$ 处的局部化。
于是 $\mathfrak p=\mathfrak m$ 是 $R$ 的极大理想，且 $A=R/I$。
设 $f_1,\ldots,f_r\in I$ 是极小生成元序列。由诺特环
$R_\mathfrak p$ 上有限模的完备化是正合的（《代数》引理
\ref{algebra-lemma-completion-tensor}），$A$ 的完备化为
$A^\wedge = R^\wedge/(f_1, \ldots, f_r)R^\wedge$
。

\medskip\noindent
若 (1) 成立，则由假设，序列 $f_1,\ldots,f_r$ 在 $R^\wedge$ 中的像
是正则序列。故由《代数》引理 \ref{algebra-lemma-completion-flat} 和
\ref{algebra-lemma-flat-increases-depth}，它在 $R$ 中也是正则序列。
因此 (1) 蕴含 (2)。

\medskip\noindent
假设 (3) 成立。令 $c=\dim(R)-\dim(A)$，并取
$f_1,\ldots,f_c\in I$，使其像生成 $I/I^2$。由 Nakayama 引理
（《代数》引理 \ref{algebra-lemma-NAK}），
$I=(f_1,\ldots,f_c)$。由于 $R$ 正则，因而是 Cohen--Macaulay 的
（《代数》命题 \ref{algebra-proposition-CM-module}），再由《代数》命题
\ref{algebra-proposition-CM-module}，$f_1,\ldots,f_c$ 是正则序列。
因此 (3) 蕴含 (2)。最后，由引理
\ref{dpa-lemma-quotient-regular-ring-by-regular-sequence}，(2) 蕴含 (1)。
\end{proof}

\noindent
下述结果由 Avramov 得到，参见 \cite{Avramov}。

\begin{proposition}
\label{dpa-proposition-avramov}
设 $A\to B$ 是诺特局部环之间的平坦局部同态。下列条件等价：
\begin{enumerate}
\item $B^\wedge$ 是完全交；
\item $A^\wedge$ 和 $(B/\mathfrak m_AB)^\wedge$ 都是完全交。
\end{enumerate}
\end{proposition}

\begin{proof}
考虑交换图
$$
\xymatrix{
B \ar[r] & B^\wedge \\
A \ar[u] \ar[r] & A^\wedge \ar[u]
}
$$
由于水平映射忠实平坦（《代数》引理
\ref{algebra-lemma-completion-faithfully-flat}），右侧竖直箭头平坦
（例如由《代数》引理
\ref{algebra-lemma-criterion-flatness-fibre-Noetherian}）。此外，
$(B/\mathfrak m_A B)^\wedge = B^\wedge/\mathfrak m_{A^\wedge} B^\wedge$
由《代数》引理 \ref{algebra-lemma-completion-tensor} 成立。
故可假设 $A$、$B$ 是完备诺特局部环。

\medskip\noindent
假设 $A$、$B$ 是完备诺特局部环。选取《代数详论》引理
\ref{more-algebra-lemma-embed-map-Noetherian-complete-local-rings}
中的交换图
$$
\xymatrix{
S \ar[r] & B \\
R \ar[u] \ar[r] & A \ar[u]
}
$$
令 $I=\Ker(R\to A)$、$J=\Ker(S\to B)$。注意，由于
$R/I=A\to B=S/J$ 平坦，映射
$J/IS \otimes_R R/\mathfrak m_R \to J/J \cap \mathfrak m_R S$
是同构。因此，$J/IS$ 的一个极小生成元系映到
$\Ker(S/\mathfrak m_RS\to B/\mathfrak m_AB)$ 的一个极小生成元系。
最后，$S/\mathfrak m_RS$ 是正则局部环。

\medskip\noindent
假设 (1) 成立，即 $J$ 由正则序列生成。由于
$A=R/I\to B=S/J$ 平坦，引理 \ref{dpa-lemma-perfect-map-ci} 适用，
从而 $I$ 和 $J/IS$ 都由正则序列生成。由《代数》引理
\ref{algebra-lemma-dimension-base-fibre-equals-total}，有
$\dim(B)=\dim(A)+\dim(B/\mathfrak m_AB)$ 且
$\dim(S/IS)=\dim(A)+\dim(S/\mathfrak m_RS)$。
因此，$J/IS$ 由
$$
\dim(S/IS) - \dim(S/J) = \dim(S/\mathfrak m_R S) - \dim(B/\mathfrak m_A B)
$$
个元素生成（《代数》引理 \ref{algebra-lemma-one-equation}）。
因此，$\Ker(S/\mathfrak m_RS\to B/\mathfrak m_AB)$ 也由同样数目的
元素生成（见上）。故该核由正则序列生成，例如参见引理
\ref{dpa-lemma-quotient-regular-ring}。由此得到 (2)。

\medskip\noindent
若 (2) 成立，则 $I$ 和 $J/J\cap\mathfrak m_RS$ 都由正则序列生成。
提升这些生成元（见上），利用 $R/I\to S/IS$ 的平坦性以及
Grothendieck 引理（《代数》引理
\ref{algebra-lemma-grothendieck-regular-sequence}），可知 $J/IS$
在 $S/IS$ 中由正则序列生成。于是引理
\ref{dpa-lemma-perfect-map-ci} 表明 $J$ 由正则序列生成，故 (1) 成立。
\end{proof}

\begin{definition}
\label{dpa-definition-lci}
设 $A$ 是诺特环。
\begin{enumerate}
\item 若 $A$ 是局部环，并且其完备化是上述意义下的完全交，
则称 $A$ 是{\it 完全交}。
\item 一般而言，若 $A$ 的所有局部环都是完全交，则称 $A$
是{\it 局部完全交}。
\end{enumerate}
\end{definition}

\noindent
下面将检验这与《代数》定义 \ref{algebra-definition-lci-field} 和
\ref{algebra-definition-lci-local-ring} 中引入的术语不冲突。
但首先说明该定义“合理”：若 $A$ 是诺特局部完全交，则 $A$
是局部完全交，即它的所有局部环都是完全交。

\begin{lemma}
\label{dpa-lemma-ci-good}
设 $(A,\mathfrak m)$ 是诺特局部环，$\mathfrak p\subset A$ 是素理想。
若 $A$ 是完全交，则 $A_\mathfrak p$ 也是完全交。
\end{lemma}

\begin{proof}
选取 $A^\wedge$ 中位于 $\mathfrak p$ 上方的素理想 $\mathfrak q$
（这是可能的，因为由《代数》引理
\ref{algebra-lemma-completion-faithfully-flat}，$A\to A^\wedge$
忠实平坦）。于是 $A_\mathfrak p\to(A^\wedge)_\mathfrak q$
是平坦局部环同态。因此，由命题 \ref{dpa-proposition-avramov}，
$A_\mathfrak p$ 是完全交，当且仅当 $(A^\wedge)_\mathfrak q$
是完全交。故只须在 $A$ 完备时证明该引理（这是证明的关键步骤）。

\medskip\noindent
假设 $A$ 完备。由定义，可将 $A$ 写成
$R/(f_1,\ldots,f_r)$，其中 $f_1,\ldots,f_r$ 是正则局部环 $R$
中的正则序列。设 $\mathfrak q\subset R$ 是对应于 $\mathfrak p$
的素理想。注意 $f_1,\ldots,f_r\in\mathfrak q$，且
$A_\mathfrak p=R_\mathfrak q/(f_1,\ldots,f_r)R_\mathfrak q$。
故由引理 \ref{dpa-lemma-quotient-regular-ring-by-regular-sequence}，
$A_\mathfrak p$ 是完全交。
\end{proof}

\begin{lemma}
\label{dpa-lemma-check-lci-at-maximal-ideals}
设 $A$ 是诺特环。则 $A$ 是局部完全交，当且仅当对 $A$ 的每个极大理想
$\mathfrak m$，$A_\mathfrak m$ 都是完全交。
\end{lemma}

\begin{proof}
立即由引理 \ref{dpa-lemma-ci-good} 和定义得到。
\end{proof}

\begin{lemma}
\label{dpa-lemma-check-lci-agrees}
设 $S$ 是域 $k$ 上的有限型代数。
\begin{enumerate}
\item 对素理想 $\mathfrak q\subset S$，局部环 $S_\mathfrak q$
是《代数》定义 \ref{algebra-definition-lci-local-ring} 意义下的完全交，
当且仅当它是定义 \ref{dpa-definition-lci} 意义下的完全交；
\item $S$ 是《代数》定义 \ref{algebra-definition-lci-field} 意义下的
局部完全交，当且仅当它是定义 \ref{dpa-definition-lci} 意义下的
局部完全交。
\end{enumerate}
\end{lemma}

\begin{proof}
证明 (1)。设 $k[x_1,\ldots,x_n]\to S$ 是满射，
$\mathfrak p\subset k[x_1,\ldots,x_n]$ 是对应于 $\mathfrak q$
的素理想，并设 $I\subset k[x_1,\ldots,x_n]$ 是该满射的核。
注意 $k[x_1,\ldots,x_n]_\mathfrak p\to S_\mathfrak q$ 是满射，
其核为 $I_\mathfrak p$。由《代数》命题
\ref{algebra-proposition-finite-gl-dim-polynomial-ring}，
$k[x_1,\ldots,x_n]$ 是正则环。因此，将引理
\ref{dpa-lemma-quotient-regular-ring} 与《代数》引理
\ref{algebra-lemma-lci-local} 结合，便得到 (1) 中两个概念的等价性。

\medskip\noindent
证明 (1) 后，(2) 中的等价性由定义和《代数》引理
\ref{algebra-lemma-lci-global} 得到。
\end{proof}

\begin{lemma}
\label{dpa-lemma-avramov}
设 $A\to B$ 是诺特局部环之间的平坦局部同态。下列条件等价：
\begin{enumerate}
\item $B$ 是完全交；
\item $A$ 和 $B/\mathfrak m_AB$ 都是完全交。
\end{enumerate}
\end{lemma}

\begin{proof}
既然定义已经成立，这只是（非平凡的）命题
\ref{dpa-proposition-avramov} 的平凡改述。
\end{proof}









\section{局部完全交同态}
\label{dpa-section-lci-homomorphisms}

\noindent
设 $A\to B$ 是诺特完备局部环之间的局部同态。Cohen 结构定理的一个
推论是，可以找到诺特完备局部环的交换图
$$
\xymatrix{
S \ar[r] & B \\
& A \ar[lu] \ar[u]
}
$$
，其中 $S\to B$ 为满射，$A\to S$ 平坦，且
$S/\mathfrak m_AS$ 是正则局部环。这由《代数详论》引理
\ref{more-algebra-lemma-embed-map-Noetherian-complete-local-rings} 得到。
暂且称 $A\to S\to B$ 是 $A\to B$ 的{\it 良好分解}，若 $S$
是诺特局部环，$A\to S\to B$ 是局部环同态，$S\to B$ 为满射，
$A\to S$ 平坦，且 $S/\mathfrak m_AS$ 正则。若存在某个良好分解
$A\to S\to B$，使得 $S\to B$ 的核由正则序列生成，则称
$A\to B$ 是{\it 完全交同态}。下述引理表明该概念与图的选择无关。

\begin{lemma}
\label{dpa-lemma-ci-map-well-defined}
设 $A\to B$ 是诺特完备局部环之间的局部同态。下列条件等价：
\begin{enumerate}
\item 对某个良好分解 $A\to S\to B$，$S\to B$ 的核由正则序列生成；
\item 对每个良好分解 $A\to S\to B$，$S\to B$ 的核都由正则序列生成。
\end{enumerate}
\end{lemma}

\begin{proof}
设 $A\to S\to B$ 是良好分解。由于 $B$ 完备，得到分解
$A\to S^\wedge\to B$，其中 $S^\wedge$ 是 $S$ 的完备化。
注意这也是良好分解：环同态 $S\to S^\wedge$ 平坦
（《代数》引理 \ref{algebra-lemma-completion-flat}），故
$A\to S^\wedge$ 平坦。由于 $S/\mathfrak m_AS$ 正则，环
$S^\wedge/\mathfrak m_AS^\wedge=(S/\mathfrak m_AS)^\wedge$ 也正则
（《代数详论》引理 \ref{more-algebra-lemma-completion-regular}）。
设 $f_1,\ldots,f_r$ 是 $\Ker(S\to B)$ 的极小生成元序列。
以下将反复使用而不再说明：诺特局部环中的理想由正则序列生成，
当且仅当其任意极小生成元集都是正则序列。由《代数》引理
\ref{algebra-lemma-flat-increases-depth}，$f_1,\ldots,f_r$ 在 $S$
中为正则序列，当且仅当它们在完备化 $S^\wedge$ 中为正则序列。
此外，
$$
S^\wedge/(f_1, \ldots, f_r)R^\wedge =
(S/(f_1, \ldots, f_n))^\wedge = B^\wedge = B
$$
因为 $B$ 是 $\mathfrak m_B$-进完备的（第一个等式由《代数》引理
\ref{algebra-lemma-completion-tensor} 得到）。因此，$S\to B$ 的核
由正则序列生成，当且仅当 $S^\wedge\to B$ 的核由正则序列生成。
故只须考虑 $S$ 完备的良好分解。

\medskip\noindent
假设有两个分解 $A\to S\to B$ 和 $A\to S'\to B$，其中 $S$、$S'$
完备。由《代数详论》引理
\ref{more-algebra-lemma-dominate-two-surjections}，环
$S\times_BS'$ 是诺特完备局部环。因此，利用《代数详论》引理
\ref{more-algebra-lemma-embed-map-Noetherian-complete-local-rings}，
可选取良好分解 $A\to S''\to S\times_BS'$，其中 $S''$ 完备。
故只须证明：若 $A\to S'\to S\to B$ 是可比较的良好分解，则
$\Ker(S\to B)$ 由正则序列生成，当且仅当 $\Ker(S'\to B)$
由正则序列生成。

\medskip\noindent
设 $A\to S'\to S\to B$ 是可比较的良好分解。首先，由于
$S'/\mathfrak m_RS'\to S/\mathfrak m_RS$ 是正则局部环的满射，
其核由正则序列
$\overline{x}_1, \ldots, \overline{x}_c \in
\mathfrak m_{S'}/\mathfrak m_R S'$
生成；该序列可扩充为正则局部环 $S'/\mathfrak m_RS'$ 的一个
正则参数系，参见《代数》引理
\ref{algebra-lemma-regular-quotient-regular}。令 $I=\Ker(S'\to S)$。
由 $S$ 在 $R$ 上平坦，有
$$
I/\mathfrak m_R I =
\Ker(S'/\mathfrak m_R S' \to S/\mathfrak m_R S) =
(\overline{x}_1, \ldots, \overline{x}_c).
$$
选取提升 $x_1,\ldots,x_c\in I$。由于 $S'$ 在 $R$ 上平坦，
这些提升构成生成 $I$ 的正则序列，参见《代数》引理
\ref{algebra-lemma-grothendieck-regular-sequence}。

\medskip\noindent
因此，若 $\Ker(S\to B)$ 也由正则序列生成，则
$\Ker(S'\to B)$ 同样如此，参见《代数详论》引理
\ref{more-algebra-lemma-join-koszul-regular-sequences} 和
\ref{more-algebra-lemma-noetherian-finite-all-equivalent}。

\medskip\noindent
反之，假设 $J=\Ker(S'\to B)$ 由正则序列生成。由于 $I$ 的生成元
$x_1,\ldots,x_c$ 在 $\mathfrak m_{S'}/\mathfrak m_{S'}^2$
中映为线性无关的元素，映射
$I/\mathfrak m_{S'}I\to J/\mathfrak m_{S'}J$ 是单射。
因此 $J$ 有一个极小生成元系
$x_1,\ldots,x_c,y_1,\ldots,y_d$。该序列是正则序列，从而
$y_1,\ldots,y_d$ 在 $S$ 中的像构成生成 $\Ker(S\to B)$ 的正则序列。
引理得证。
\end{proof}

\noindent
在下面的命题中，请注意：若 $B$ 在 $A$ 上具有有限 Tor 维数，
特别是若 $B$ 在 $A$ 上平坦，则关于 Tor 消失的条件成立。

\begin{proposition}
\label{dpa-proposition-avramov-map}
设 $A\to B$ 是诺特局部环之间的局部同态。下列条件等价：
\begin{enumerate}
\item $B$ 是完全交，且只有有限多个 $p$ 使
$\text{Tor}^A_p(B,A/\mathfrak m_A)$ 非零；
\item $A$ 是完全交，且 $A^\wedge\to B^\wedge$ 是上述意义下的
完全交同态。
\end{enumerate}
\end{proposition}

\begin{proof}
设 $F_\bullet\to A/\mathfrak m_A$ 是由有限自由 $A$-模组成的分辨。
注意 $\text{Tor}^A_p(B,A/\mathfrak m_A)$ 是复形
$F_\bullet\otimes_AB$ 的第 $p$ 个同调。令完备化
$F_\bullet^\wedge=F_\bullet\otimes_AA^\wedge$。则
$F_\bullet^\wedge$ 是由有限自由 $A^\wedge$-模组成的
$A^\wedge/\mathfrak m_{A^\wedge}$ 的分辨（因为 $A\to A^\wedge$
平坦且有限模上的完备化正合，参见《代数》引理
\ref{algebra-lemma-completion-tensor} 和
\ref{algebra-lemma-completion-flat}）。因此
$$
F_\bullet^\wedge \otimes_{A^\wedge} B^\wedge =
F_\bullet \otimes_A B \otimes_B B^\wedge
$$
由 $B\to B^\wedge$ 的平坦性，得到
$$
\text{Tor}^{A^\wedge}_p(B^\wedge, A^\wedge/\mathfrak m_{A^\wedge}) =
\text{Tor}^A_p(B, A/\mathfrak m_A) \otimes_B B^\wedge
$$
由此可见，局部环同态 $A\to B$ 的条件 (1) 等价于局部环同态
$A^\wedge\to B^\wedge$ 的相同条件。因此可以假设 $A$、$B$
是完备诺特局部环（其余条件本来就是用完备化表述的）。

\medskip\noindent
假设 $A$、$B$ 是完备诺特局部环。选取《代数详论》引理
\ref{more-algebra-lemma-embed-map-Noetherian-complete-local-rings}
中的交换图
$$
\xymatrix{
S \ar[r] & B \\
R \ar[u] \ar[r] & A \ar[u]
}
$$
令 $I=\Ker(R\to A)$、$J=\Ker(S\to B)$。
命题现由引理 \ref{dpa-lemma-perfect-map-ci} 得到。
\end{proof}

\begin{remark}
\label{dpa-remark-no-good-ci-map}
对于一般诺特环之间的同态，似乎很难定义一个良好的“局部完全交同态”
概念。原因是：对诺特局部环 $A$，$A\to A^\wedge$ 的纤维不是局部
完全交环。因此，若 $A\to B$ 是诺特局部环之间的局部同态，且完备化
之间的映射 $A^\wedge\to B^\wedge$ 是上述意义下的完全交同态，
那么一般而言，$(A_\mathfrak p)^\wedge\to(B_\mathfrak q)^\wedge$
{\bf 不是}上述意义下的完全交同态。只研究优秀诺特环可以解决这个问题。
更一般地，可以研究形式纤维为完全交的诺特环，参见
\cite{Rodicio-ci}。我们将在《对偶化复形》第
\ref{dualizing-section-formal-fibres} 节发展这一理论。
\end{remark}

\noindent
作为本节的结尾，我们把上述概念与《代数详论》第
\ref{dpa-section-lci} 节中引入的概念作比较。

\begin{lemma}
\label{dpa-lemma-well-defined-if-you-can-find-good-factorization}
考虑诺特局部环的交换图
$$
\xymatrix{
S \ar[r] & B \\
& A \ar[lu] \ar[u]
}
$$
，其中 $S\to B$ 为满射，$A\to S$ 平坦，且
$S/\mathfrak m_AS$ 是正则局部环。下列条件等价：
\begin{enumerate}
\item $\Ker(S\to B)$ 由正则序列生成；
\item $A^\wedge\to B^\wedge$ 是上述定义的完全交同态。
\end{enumerate}
\end{lemma}

\begin{proof}
从略。提示：证明与引理 \ref{dpa-lemma-ci-map-well-defined} 的证明第一段
中的论证相同。
\end{proof}

\begin{lemma}
\label{dpa-lemma-finite-type-lci-map}
设 $A$ 是诺特环，$A\to B$ 是有限型环同态。下列条件等价：
\begin{enumerate}
\item $A\to B$ 是《代数详论》定义
\ref{more-algebra-definition-local-complete-intersection} 意义下的
局部完全交；
\item 对每个素理想 $\mathfrak q\subset B$，令
$\mathfrak p=A\cap\mathfrak q$，则环同态
$(A_\mathfrak p)^\wedge\to(B_\mathfrak q)^\wedge$ 是上述意义下的
完全交同态。
\end{enumerate}
\end{lemma}

\begin{proof}
选取满射 $R=A[x_1,\ldots,x_n]\to B$。注意 $A\to R$ 平坦且纤维正则。
设 $I$ 是 $R\to B$ 的核。假设 (2)。由引理
\ref{dpa-lemma-well-defined-if-you-can-find-good-factorization} 和《代数》引理
\ref{algebra-lemma-regular-sequence-in-neighbourhood}，$I$ 局部由正则序列
生成。换言之，(1) 成立。反之，假设 (1)。将 $R$ 和 $B$ 局部化后，
可假设 $I$ 由 Koszul 正则序列生成。由《代数详论》引理
\ref{more-algebra-lemma-noetherian-finite-all-equivalent}，$I$ 局部由
正则序列生成。故由引理
\ref{dpa-lemma-well-defined-if-you-can-find-good-factorization}，(2) 成立。
略去若干细节。
\end{proof}

\begin{lemma}
\label{dpa-lemma-avramov-map-finite-type}
设 $A$ 是诺特环，$A\to B$ 是有限型环同态，并且
$\Spec(B)\to\Spec(A)$ 的像包含 $\Spec(A)$ 的所有闭点。
下列条件等价：
\begin{enumerate}
\item $B$ 是完全交，且 $A\to B$ 具有有限 Tor 维数；
\item $A$ 是完全交，且 $A\to B$ 是《代数详论》定义
\ref{more-algebra-definition-local-complete-intersection} 意义下的
局部完全交。
\end{enumerate}
\end{lemma}

\begin{proof}
这是命题 \ref{dpa-proposition-avramov-map} 经由引理
\ref{dpa-lemma-finite-type-lci-map} 得到的改述。细节从略。
\end{proof}








\section{光滑环同态与对角}
\label{dpa-section-smooth-diagonal-perfect}

\noindent
本节利用上述材料，将光滑环同态刻画为对角完美的环同态。

\begin{lemma}
\label{dpa-lemma-local-perfect-diagonal}
设 $A\to B$ 是诺特局部环之间的局部同态，并且 $B$ 在 $A$ 上平坦且
本质有限型。若
$$
B \otimes_A B \longrightarrow B
$$
是完美环同态，即 $B$ 在 $B\otimes_AB$ 上具有有限 Tor 维数，
则 $B$ 是某个光滑 $A$-代数的局部化。
\end{lemma}

\begin{proof}
由于 $B$ 在 $A$ 上本质有限型，$B\otimes_AB$ 也如此；特别地，
$B\otimes_AB$ 是诺特环。因此，$B\otimes_AB$ 的商 $B$ 在
$B\otimes_AB$ 上伪凝聚（《代数详论》引理
\ref{more-algebra-lemma-Noetherian-pseudo-coherent}），这说明环同态的
完美性（《代数详论》定义
\ref{more-algebra-definition-pseudo-coherent-perfect}）与有限 Tor 维数
条件一致。

\medskip\noindent
可将 $B$ 写成 $R/K$，其中 $R$ 是 $A[x_1,\ldots,x_n]$ 在某个素理想处的
局部化，$K\subset R$ 是理想。记 $\mathfrak m\subset R\otimes_AR$
为极大理想，它是 $B$ 的极大理想在满射
$R\otimes_AR\to B\otimes_AB\to B$ 下的逆像。于是有满射
$$
(R \otimes_A R)_\mathfrak m \to (B \otimes_A B)_\mathfrak m \to B
$$
，从而得到引理 \ref{dpa-lemma-injective} 中的理想
$I\subset J\subset(R\otimes_AR)_\mathfrak m$。因此
$I/\mathfrak mI\to J/\mathfrak mJ$ 是单射。

\medskip\noindent
将 $K$ 写成 $(f_1,\ldots,f_r)$，其中 $r$ 极小。可以并确实假设
$f_i\in R$ 是 $A[x_1,\ldots,x_n]$ 中某个也记作 $f_i$ 的元素之像。
注意 $I$ 由 $f_1\otimes1,\ldots,f_r\otimes1$ 和
$1\otimes f_1,\ldots,1\otimes f_r$ 生成。我们断言这是 $I$ 的一个
极小生成元集。事实上，若 $\kappa$ 是 $R$、$B$、
$(R\otimes_AR)_\mathfrak m$ 和 $(B\otimes_AB)_\mathfrak m$ 的共同剩余域，
则有映射
$R \otimes_A R \to R \otimes_A \kappa \oplus \kappa \otimes_A R$
，它经由 $(R\otimes_AR)_\mathfrak m$ 分解。由于 $B$ 在 $A$ 上平坦，
并且有短正合列 $0\to K\to R\to B\to0$，故
$K\otimes_A\kappa\subset R\otimes_A\kappa$，参见《代数》引理
\ref{algebra-lemma-flat-tor-zero}。因此，将映射
$(R \otimes_A R)_\mathfrak m \to R \otimes_A \kappa \oplus \kappa \otimes_A R$
限制到 $I$，得到映射
$$
I \to K \otimes_A \kappa \oplus \kappa \otimes_A K \to
K \otimes_B \kappa \oplus \kappa \otimes_B K.
$$
元素
$f_1 \otimes 1, \ldots, f_r \otimes 1, 1 \otimes f_1, \ldots, 1 \otimes f_r$
映为该映射目标的一组基，因为由 Nakayama 引理（《代数》引理
\ref{algebra-lemma-NAK}），$f_1,\ldots,f_r$ 映为
$K\otimes_B\kappa$ 的一组基。这证明了断言。

\medskip\noindent
理想 $J$ 由 $f_1\otimes1,\ldots,f_r\otimes1$ 以及元素
$x_1\otimes1-1\otimes x_1,\ldots,x_n\otimes1-1\otimes x_n$
生成（对证明而言，只须知道这些元素属于理想 $J$）。现在可写成
$$
f_i \otimes 1 - 1 \otimes f_i =
\sum g_{ij} (x_j \otimes 1 - 1 \otimes x_j)
$$
其中 $g_{ij}\in(R\otimes_AR)_\mathfrak m$。这是关于
$A[x_1,\ldots,x_n]$ 中元素的一般事实，证明从略。记
$a_{ij}\in\kappa$ 为 $g_{ij}$ 的像。另一项计算表明，$a_{ij}$ 是
$\partial f_i/\partial x_j$ 在 $\kappa$ 中的像。
$I/\mathfrak mI\to J/\mathfrak mJ$ 的单射性与上述评注迫使矩阵
$(a_{ij})$ 具有极大秩 $r$。令
$$
C = A[x_1, \ldots, x_n]/(f_1, \ldots, f_r)
$$
并考虑朴素余切复形
$\NL_{C/A}\cong(C^{\oplus r}\to C^{\oplus n})$，其中映射由偏导数矩阵
给出。因此，$\NL_{C/A}\otimes_AB$ 同构于放在次数 $0$ 的秩
$n-r$ 自由 $B$-模。故对某个在 $B$ 中映为单位的 $g\in C$，
$C_g$ 在 $A$ 上光滑，参见《代数》引理
\ref{algebra-lemma-smooth-at-point}。证明完成。
\end{proof}

\begin{lemma}
\label{dpa-lemma-perfect-diagonal}
设 $A\to B$ 是诺特环之间的平坦有限型环同态。若
$$
B \otimes_A B \longrightarrow B
$$
是完美环同态，即 $B$ 在 $B\otimes_AB$ 上具有有限 Tor 维数，
则 $B$ 是光滑 $A$-代数。
\end{lemma}

\begin{proof}
这由引理 \ref{dpa-lemma-local-perfect-diagonal} 以及光滑环同态的一般事实
得到，参见《代数》引理 \ref{algebra-lemma-smooth-at-point} 和
\ref{algebra-lemma-locally-smooth}。或者，读者可以略微修改引理
\ref{dpa-lemma-local-perfect-diagonal} 的证明来证明本引理。
\end{proof}






\section{余法模的自由性}
\label{dpa-section-freeness-conormal}

\noindent
Tate 分辨及其上的导子可用于证明本节结果的（更强）版本，
参见 \cite{Iyengar}。另外两篇更初等的参考文献是
\cite{Vasconcelos} 和 \cite{Ferrand-lci}。

\begin{lemma}
\label{dpa-lemma-free-summand-in-ideal-finite-proj-dim}
\begin{reference}
\cite{Vasconcelos}
\end{reference}
设 $R$ 是诺特局部环，$I\subset R$ 是在 $R$ 上具有有限投射维数的理想。
若 $F\subset I/I^2$ 是同构于 $R/I$ 的直和项，则存在非零因子
$x\in I$，使得 $x$ 在 $I/I^2$ 中的像生成 $F$。
\end{lemma}

\begin{proof}
由假设，可选取有限自由分辨
$$
0 \to R^{\oplus n_e} \to R^{\oplus n_{e-1}} \to \ldots \to
R^{\oplus n_1} \to R \to R/I \to 0
$$
于是 $\varphi_1:R^{\oplus n_1}\to R$ 的秩为 $1$，并由《代数》命题
\ref{algebra-proposition-what-exact} 可知 $I$ 包含一个非零因子 $y$。
设 $\mathfrak p_1,\ldots,\mathfrak p_n$ 是 $R$ 的伴随素理想，
参见《代数》引理 \ref{algebra-lemma-finite-ass}。选取理想
$I^2\subset J\subset I$，使得 $J/I^2=F$。由于 $y^2\in J$ 且
$y^2\notin\mathfrak p_i$，对所有 $i$ 有
$J\not\subset\mathfrak p_i$，参见《代数》引理
\ref{algebra-lemma-ass-zero-divisors}。由 Nakayama 引理（《代数》引理
\ref{algebra-lemma-NAK}），$J\not\subset\mathfrak mJ+I^2$。
由《代数》引理 \ref{algebra-lemma-silly}，可选取 $x\in J$，使得
$x\notin\mathfrak mJ+I^2$，且对 $i=1,\ldots,n$，
$x\notin\mathfrak p_i$。于是 $x$ 是非零因子，并且 $x$ 在 $I/I^2$
中的像由 Nakayama 引理生成直和项 $J/I^2\cong R/I$。
\end{proof}

\begin{lemma}
\label{dpa-lemma-vasconcelos}
\begin{reference}
\cite[定理 1.1]{Vasconcelos} 的局部版本
\end{reference}
设 $R$ 是诺特局部环，$I\subset R$ 是在 $R$ 上具有有限投射维数的理想。
若 $F\subset I/I^2$ 是秩为 $r$ 的自由直和项，则存在正则序列
$x_1,\ldots,x_r\in I$，使得
$x_1\bmod I^2,\ldots,x_r\bmod I^2$ 生成 $F$。
\end{lemma}

\begin{proof}
若 $r=0$，则无须证明。假设 $r>0$。由引理
\ref{dpa-lemma-free-summand-in-ideal-finite-proj-dim}，可选取 $x\in I$，
使得 $x$ 是非零因子，且 $x\bmod I^2$ 生成 $F$ 中一个同构于
$R/I$ 的直和项。考虑环 $R'=R/(x)$ 和理想 $I'=I/(x)$。
显然 $R'/I'=R/I$。短正合列
$$
0 \to R/I \xrightarrow{x} I/xI \to I' \to 0
$$
分裂，因为映射 $I/xI\to I/I^2$ 将 $xR/xI$ 映为一个直和项。
现在 $I/xI=I\otimes_R^\mathbf{L}R'$ 在 $R'$ 上具有有限投射维数，
参见《代数详论》引理 \ref{more-algebra-lemma-perfect-module} 和
\ref{more-algebra-lemma-pull-perfect}。故直和项 $I'$ 在 $R'$ 上具有
有限投射维数。另一方面，有短正合列
$0\to xR/xI\to I/I^2\to I'/(I')^2\to0$，从而
$I'/(I')^2$ 具有秩为 $r-1$ 的自由直和项
$F'=F/(R/I\cdot x)$。对 $r$ 作归纳，可选取正则序列
$x'_2,\ldots,x'_r\in I'$，使其同余类自由生成 $F'$。
令 $x_1=x$，并任取 $R$ 中提升 $x'_1,\ldots,x'_r$ 的元素
$x_2,\ldots,x_r$，即可得到引理结论。
\end{proof}

\begin{proposition}
\label{dpa-proposition-regular-ideal}
\begin{reference}
\cite[推论 1]{Vasconcelos} 的一个变体。另见
\cite{Iyengar} 和 \cite{Ferrand-lci}。
\end{reference}
设 $R$ 是诺特环，$I\subset R$ 是具有有限投射维数的理想，并且
$I/I^2$ 在 $R/I$ 上有限局部自由。则 $I$ 是正则理想
（《代数详论》定义 \ref{more-algebra-definition-regular-ideal}）。
\end{proposition}

\begin{proof}
由《代数》引理 \ref{algebra-lemma-regular-sequence-in-neighbourhood}，
只须证明对每个 $\mathfrak p\supset I$，
$I_\mathfrak p\subset R_\mathfrak p$ 由正则序列生成。
故可假设 $R$ 是局部环。若 $I/I^2$ 的秩为 $r$，则由引理
\ref{dpa-lemma-vasconcelos}，存在生成 $I/I^2$ 的正则序列
$x_1,\ldots,x_r\in I$。由 Nakayama 引理（《代数》引理
\ref{algebra-lemma-NAK}），$I$ 由 $x_1,\ldots,x_r$ 生成。
\end{proof}

\noindent
对任意环的局部完全交同态 $A\to B$，朴素余切复形 $\NL_{B/A}$
是 Tor 振幅位于 $[-1,0]$ 的完美复形，参见《代数详论》引理
\ref{more-algebra-lemma-lci-NL}。利用上述结果，可以证明这在某些情形
刻画局部完全交同态。

\begin{lemma}
\label{dpa-lemma-perfect-NL-lci}
设 $A\to B$ 是诺特环之间的完美环同态（《代数详论》定义
\ref{more-algebra-definition-pseudo-coherent-perfect}）。下列条件等价：
\begin{enumerate}
\item $\NL_{B/A}$ 的 Tor 振幅位于 $[-1,0]$；
\item $\NL_{B/A}$ 是 $D(B)$ 中 Tor 振幅位于 $[-1,0]$ 的完美对象；
\item $A\to B$ 是局部完全交（《代数详论》定义
\ref{more-algebra-definition-local-complete-intersection}）。
\end{enumerate}
\end{lemma}

\begin{proof}
将 $B$ 写成 $A[x_1,\ldots,x_n]/I$。则 $\NL_{B/A}$ 由 $B$-模复形
$$
I/I^2 \longrightarrow \bigoplus B \text{d}x_i
$$
表示，其中 $I/I^2$ 放在次数 $-1$。由于次数 $0$ 项有限自由，
该复形的 Tor 振幅位于 $[-1,0]$，当且仅当 $I/I^2$ 是平坦 $B$-模，
参见《代数详论》引理 \ref{more-algebra-lemma-last-one-flat}。
由于 $I/I^2$ 是有限 $B$-模且 $B$ 诺特，这又等价于 $I/I^2$
是有限局部自由 $B$-模（《代数》引理
\ref{algebra-lemma-finite-projective}）。因此 (1) 与 (2) 的等价性显然。
此外，应用命题 \ref{dpa-proposition-regular-ideal}，并注意正则理想既是
Koszul 正则理想也是拟正则理想（参见《代数详论》第
\ref{more-algebra-section-ideals} 节），也得到 (1) 与 (3) 的等价性。
\end{proof}

\begin{lemma}
\label{dpa-lemma-flat-fp-NL-lci}
设 $A\to B$ 是有限表示的平坦环同态。下列条件等价：
\begin{enumerate}
\item $\NL_{B/A}$ 的 Tor 振幅位于 $[-1,0]$；
\item $\NL_{B/A}$ 是 $D(B)$ 中 Tor 振幅位于 $[-1,0]$ 的完美对象；
\item $A\to B$ 是 syntomic 同态（《代数》定义
\ref{algebra-definition-lci}）；
\item $A\to B$ 是局部完全交（《代数详论》定义
\ref{more-algebra-definition-local-complete-intersection}）。
\end{enumerate}
\end{lemma}

\begin{proof}
(3) 与 (4) 的等价性是《代数详论》引理
\ref{more-algebra-lemma-syntomic-lci}。

\medskip\noindent
若 $A\to B$ 是 syntomic 同态，则可找到余笛卡尔图
$$
\xymatrix{
B_0 \ar[r] & B \\
A_0 \ar[r] \ar[u] & A \ar[u]
}
$$
，使得 $A_0\to B_0$ 是 syntomic 同态且 $A_0$ 诺特，参见《代数》引理
\ref{algebra-lemma-limit-module-finite-presentation} 和
\ref{algebra-lemma-colimit-lci}。由引理 \ref{dpa-lemma-perfect-NL-lci}，
$\NL_{B_0/A_0}$ 是 Tor 振幅位于 $[-1,0]$ 的完美对象。
由《代数详论》引理 \ref{more-algebra-lemma-base-change-NL-flat}，
$\NL_{B/A}=\NL_{B_0/A_0}\otimes_{B_0}^\mathbf{L}B$ 也具有同样性质
（另见《代数详论》引理 \ref{more-algebra-lemma-pull-tor-amplitude} 和
\ref{more-algebra-lemma-pull-perfect}）。这证明 (3) 蕴含 (2)。

\medskip\noindent
假设 (1)。由《代数详论》引理
\ref{more-algebra-lemma-base-change-NL-flat}，对每个以域 $k$ 为目标的
环同态 $A\to k$，$\NL_{B\otimes_Ak/k}$ 的 Tor 振幅位于
$[-1,0]$（参见《代数详论》引理
\ref{more-algebra-lemma-pull-tor-amplitude}）。故由引理
\ref{dpa-lemma-perfect-NL-lci}，$k\to B\otimes_Ak$ 是局部完全交同态。
于是按定义 $A\to B$ 是 syntomic 同态。这证明 (1) 蕴含 (3)，
并完成证明。
\end{proof}





\section{Koszul 复形与 Tate 分辨}
\label{dpa-section-koszul-vs-tate}

\noindent
本节把《代数详论》引理
\ref{more-algebra-lemma-sequence-Koszul-complexes} 的结论“提升”到
带有与微分分次结构相容的除幂的微分分次代数范畴中（请注意，本节把
Koszul 复形表示为链复形，而上述引文中使用的是上链复形）。

\medskip\noindent
设 $R$ 是环。设 $I \subset R$ 是由 $f_1, \ldots, f_r \in R$ 生成的
理想。对 $n \geq 1$，记
$$
K_n = K_{n, \bullet} = R\langle \xi_1, \ldots, \xi_r\rangle
$$
为微分分次 Koszul 代数，其中 $\xi_i$ 的次数为 $1$，且
$\text{d}(\xi_i) = f_i^n$。其上存在唯一的除幂结构（如定义
\ref{dpa-definition-divided-powers-dga} 中所述），见例
\ref{dpa-example-adjoining-odd}。对 $m > n$，转移映射
$K_m \to K_n$ 是与除幂相容的微分分次代数映射，它把 $\xi_i$
送到 $f_i^{m - n}\xi_i$。

\begin{lemma}
\label{dpa-lemma-lift-tate-to-koszul}
在上述情形中，若 $R$ 是诺特环，则对每个 $n$，存在 $N \geq n$
以及映射
$$
K_N \to A \to R/(f_1^N, \ldots, f_r^N)\quad\text{以及}\quad A \to K_n
$$
满足下列性质：
\begin{enumerate}
\item $(A, \text{d}, \gamma)$ 如定义
\ref{dpa-definition-divided-powers-dga} 所述；
\item $A \to R/(f_1^N, \ldots, f_r^N)$ 是拟同构；
\item 复合 $K_N \to A \to R/(f_1^N, \ldots, f_r^N)$ 是典范映射；
\item 复合 $K_N \to A \to K_n$ 是转移映射；
\item $A_0 = R \to R/(f_1^N, \ldots, f_r^N)$ 是典范满射；
\item $A$ 是 $R$ 上的分次除幂多项式代数，且每个次数中只有有限多个
生成元；以及
\item $A \to K_n$ 是与除幂相容的微分分次 $R$-代数同态，并在
$0$ 次同调上诱导典范映射
$R/(f_1^N, \ldots, f_r^N) \to R/(f_1^n, \ldots, f_r^n)$。
\end{enumerate}
条件 (4) 意味着，$A$ 是从 $A_0$ 出发，依照例
\ref{dpa-example-adjoining-odd} 或 \ref{dpa-example-adjoining-even}，
在每个次数 $> 0$ 中逐次添加有限集变量 $T$ 构造而成的。
\end{lemma}

\begin{proof}
固定 $n$。若 $r = 0$，则可直接取 $A = R$。假设 $r > 0$。由
《代数详论》引理 \ref{more-algebra-lemma-sequence-Koszul-complexes}
（将其改写为链复形的语言），可选取
$$
n_{r} > n_{r - 1} > \ldots > n_1 > n_0 = n
$$
使得 Koszul 代数上的转移映射 $K_{n_{i + 1}} \to K_{n_i}$（见上）
在次数 $> 0$ 的同调上诱导零映射。我们将取 $N = n_r$ 来证明本引理。

\medskip\noindent
我们完全按照引理 \ref{dpa-lemma-tate-resolution} 的陈述与证明来构造
$A$。于是将有
$$
A = \colim A(m),\quad\text{且}\quad
A(0) \to A(1) \to A(2) \to \ldots \to R/(f_1^N, \ldots, f_r^N)
$$
这立即给出性质 (1)、(2)、(5) 和 (6)。为完成证明，我们将构造
$R$-代数映射 $K_N \to A \to K_n$。具体构造如下：
\begin{enumerate}
\item 同构 $A(1) \to K_N = K_{n_r}$；
\item 映射 $A(2) \to K_{n_{r - 1}}$；
\item $\ldots$
\item 映射 $A(r) \to K_{n_1}$；
\item 映射 $A(r + 1) \to K_{n_0} = K_n$；以及
\item 映射 $A \to K_n$。
\end{enumerate}
在每一步中，所构造的映射都位于以相容方式赋予除幂的微分分次代数之间；
每个映射都经由 Koszul 代数之间的转移映射与前一个映射相容，并在
$0$ 次同调上诱导显然的典范映射。

\medskip\noindent
回忆 $A(0) = R$。当 $m = 1$ 时，引理
\ref{dpa-lemma-tate-resolution} 的证明选取
$A(1) = R\langle T_1, \ldots, T_r\rangle$，其中 $T_i$ 的次数为
$1$，且 $\text{d}(T_i) = f_i^N$。也就是说，$f_i^N$ 是
$A(0) \to R/(f_1^N, \ldots, f_r^N)$ 的核的生成元。因此，对
$A(1) \to K_N = K_{n_r}$ 使用映射
$$
\varphi_1 : A(1) \longrightarrow K_{n_r},\quad T_i \longmapsto \xi_i
$$
它是同构。

\medskip\noindent
当 $m = 2$ 时，引理 \ref{dpa-lemma-tate-resolution} 的证明中的构造选取
生成元 $e_1, \ldots, e_t \in \Ker(\text{d} : A(1)_1 \to A(1)_0)$。
接着取
$A(2) = A(1)\langle T_1, \ldots, T_t\rangle$
作为除幂多项式代数，其中 $T_i$ 的次数为 $2$，且
$\text{d}(T_i) = e_i$。由于 $\varphi_1(e_i)$ 是 $K_{n_r}$ 中的
闭链，由上述对 $n_r$ 与 $n_{r - 1}$ 的选取可知，它在
$K_{n_{r - 1}}$ 中的像是边界。因此可以构造交换图
$$
\xymatrix{
A(1) \ar[d] \ar[r]_{\varphi_1} & K_{n_r} \ar[d] \\
A(2) \ar[r]^{\varphi_2} & K_{n_{r - 1}}
}
$$
方法是把 $T_i$ 送到次数 $2$ 中边界等于 $\varphi_1(e_i)$ 之像的
一个元素。由除幂多项式代数的泛性质，映射 $\varphi_2$ 存在，并与
微分和除幂相容。

\medskip\noindent
对 $m = 2, \ldots, r$，完全以相同方式构造代数 $A(m)$ 与映射
$\varphi_m : A(m) \to K_{n_{r + 1 - m}}$。

\medskip\noindent
给定映射 $A(r) \to K_{n_1}$，可见复合
$H_r(A(r)) \to H_r(K_{n_1}) \to H_r(K_{n_0}) \subset (K_{n_0})_r$
为零，故可将其延拓为 $A(r + 1) \to K_{n_0} = K_n$，只需把
$A(r + 1)$ 的新多项式生成元送到零。

\medskip\noindent
构造出 $A(r + 1) \to K_{n_0} = K_n$ 后，只需把新的多项式生成元
送到零，即可按唯一可能的方式延拓到
$A(r + 2), A(r + 3), \ldots$。证明完毕。
\end{proof}

\begin{remark}
\label{dpa-remark-pro-system-koszul}
在上述情形中，若 $R$ 是诺特环，则可归纳地选取整数列
$1 = n_0 < n_1 < n_2 < \ldots $，使得对
$i = 1, 2, 3, \ldots$，存在引理
\ref{dpa-lemma-lift-tate-to-koszul} 中的映射
$K_{n_i} \to A_i \to R/(f_1^{n_i}, \ldots, f_r^{n_i})$
与 $A_i \to K_{n_{i - 1}}$。把复合
$A_{i + 1} \to K_{n_i} \to A_i$ 记作 $A_{i + 1} \to A_i$。
则图
$$
\xymatrix{
K_{n_1} \ar[d] &
K_{n_2} \ar[d] \ar[l] &
K_{n_3} \ar[d] \ar[l] &
\ldots \ar[l] \\
A_1 \ar[d] &
A_2 \ar[l] \ar[d] &
A_3 \ar[l] \ar[d] &
\ldots \ar[l] \\
K_1 &
K_{n_1} \ar[l] &
K_{n_2} \ar[l] &
\ldots \ar[l]
}
$$
交换。由此可见，在带有相容除幂的微分分次 $R$-代数范畴中，逆系统
$(K_n)$ 与 $(A_n)$ 是 pro-同构的。
\end{remark}

\begin{lemma}
\label{dpa-lemma-compute-cohomology-adjoin-variable}
设 $(A, \text{d}, \gamma)$、$d \geq 1$、$f \in A_{d - 1}$ 以及
$A\langle T \rangle$ 如引理 \ref{dpa-lemma-extend-differential} 中所述。
\begin{enumerate}
\item 若 $d = 1$，则有长正合列
$$
\ldots \to H_0(A) \xrightarrow{f} H_0(A) \to H_0(A\langle T \rangle) \to 0
$$
\item 当 $d = 2$ 时，存在有界谱序列
$(E_1)_{i, j} = H_{j - i}(A) \cdot T^{[i]}$
收敛到 $H_{i + j}(A\langle T \rangle)$。微分
$(d_1)_{i, j} : H_{j - i}(A) \cdot T^{[i]} \to
H_{j - i + 1}(A) \cdot T^{[i - 1]}$
把 $\xi \cdot T^{[i]}$ 送到 $f \xi \cdot T^{[i - 1]}$ 的类。
\item 根据需要在此补充其他次数的情形。
\end{enumerate}
\end{lemma}

\begin{proof}
当 $d = 1$ 时，有复形的短正合列
$$
0 \to A \to A\langle T \rangle \to A \cdot T \to 0
$$
由此容易得到结论 (1)。当 $d = 2$ 时，把 $A\langle T \rangle$
视为带有子复形
$$
F^pA\langle T \rangle = \bigoplus\nolimits_{i \leq p} A \cdot T^{[i]}
$$
的滤过链复形。应用《同调代数》第
\ref{homology-section-filtered-complex} 节的谱序列（将其改写为链复形）
即得 (2)。
\end{proof}

\noindent
稍后将用到下述引理。

\begin{lemma}
\label{dpa-lemma-construct-some-maps}
在上述情形中，对所有 $n \geq t \geq 1$，存在 $N > n$ 以及映射
$$
K_t \longrightarrow K_n \otimes_R K_t
$$
位于左微分分次 $K_N$-模的导出范畴中，使它与乘法映射的复合
（无论按哪个方向复合）都是转移映射。
\end{lemma}

\begin{proof}
先对 $r = 1$ 证明。令 $f = f_1$。写
$K_t = R\langle x \rangle$、
$K_n = R\langle y \rangle$ 与 $K_N = R\langle z \rangle$，其中
$x$、$y$、$z$ 的次数均为 $1$，且
$\text{d}(x) = f^t$、$\text{d}(y) = f^n$、$\text{d}(z) = f^N$。
我们断言，对所有 $N > t$，存在拟同构
$$
B_{N, t} = R\langle x, z, u \rangle
\longrightarrow
K_t,\quad
x \mapsto x,\quad
z \mapsto f^{N - t}x,\quad
u \mapsto 0
$$
这里左端表示变量 $x,z$ 的次数为 $1$、变量 $u$ 的次数为 $2$ 的
除幂多项式代数，并且
$\text{d}(x) = f^t$、$\text{d}(z) = f^N$、且
$\text{d}(u) = z - f^{N - t}x$。为证明该断言，注意
$H_*(R\langle x, z\rangle)$ 的下列三个子模相同：
\begin{enumerate}
\item $H_*(R\langle x, z\rangle) \to H_*(K_t)$ 的核；
\item 映射
$z - f^{N - t}x : H_*(R\langle x, z\rangle) \to H_*(R\langle x, z\rangle)$
的像；以及
\item 映射
$z - f^{N - t}x : H_*(R\langle x, z\rangle) \to H_*(R\langle x, z\rangle)$
的核。
\end{enumerate}
这一点可由直接计算证明\footnote{提示：令
$z' = z - f^{N - t}x$，则
$R\langle x, z\rangle = R\langle x, z'\rangle$ 且
$\text{d}(z') = 0$；此外，映射
$R\langle x, z'\rangle \to K_t$ 正是把 $z'$ 杀掉的映射。}，
计算从略。再应用引理
\ref{dpa-lemma-compute-cohomology-adjoin-variable} 的第 (2) 部分，
即可得出该断言。

\medskip\noindent
通过把 $z$ 送到 $z$ 的微分分次 $R$-代数同态
$K_N \to B_{N, t}$，可把 $B_{N, t} \to K_t$ 看作左微分分次
$K_N$-模的拟同构。为定义引理陈述中的箭头，我们使用同态
$$
B_{N, t} = R\langle x, z, u \rangle \to K_n \otimes_R K_t,\quad
x \mapsto 1 \otimes x,\quad
z \mapsto f^{N - n}y \otimes 1,\quad
u \mapsto - f^{N - n - t}y \otimes x
$$
只要假设 $N \geq n + t$，这就有意义。一个简单明快的计算表明，
它与乘法映射的（前复合或后复合）是转移映射。

\medskip\noindent
当 $r > 1$ 时，把每个 Koszul 代数写成单变量 Koszul 代数的张量积，
并应用上述构造。换言之，写
$$
K_t = R\langle x_1, \ldots, x_r\rangle =
R\langle x_1\rangle \otimes_R \ldots \otimes_R R\langle x_r\rangle
$$
其中 $x_i$ 的次数为 $1$，且 $\text{d}(x_i) = f_i^t$。在
$r > 1$ 的情形下，使用
$$
B_{N, t} =
R\langle x_1, z_1, u_1 \rangle
\otimes_R \ldots \otimes_R
R\langle x_r, z_r, u_r \rangle
$$
其中 $x_i,z_i$ 的次数为 $1$，$u_i$ 的次数为 $2$，并且
$\text{d}(x_i) = f_i^t$、$\text{d}(z_i) = f_i^N$、且
$\text{d}(u_i) = z_i - f_i^{N - t}x_i$。张量积映射
$B_{N, t} \to K_t$ 是拟同构，因为它是自由 $R$-模的上有界复形之间
诸拟同构的张量积。最后，将映射
$$
B_{N, t}
\to
K_n \otimes_R K_t =
R\langle y_1, \ldots, y_r\rangle \otimes_R
R\langle x_1, \ldots, x_r\rangle
$$
定义为 $r = 1$ 情形中所构造映射的张量积；也可直接用规则
$x_i \mapsto 1 \otimes x_i$、
$z_i \mapsto f_i^{N - n}y_i \otimes 1$ 与
$u_i \mapsto - f_i^{N - n - t}y_i \otimes x_i$ 来定义，只要
$N \geq n + t$，这些规则就有意义。细节从略。
\end{proof}
% Generated mechanically from the frozen sdga.tex target.
% Do not edit this generated file; edit the bound locale source instead.

% OK, start here.
%

\setcounter{chapter}{23}
\chapter{微分分次层}



\phantomsection
\label{sdga-section-phantom}


\section{引言}
\label{sdga-section-introduction}

\noindent
本章继续《微分分次代数》第 \ref{dga-section-introduction} 节开始的
讨论。一篇综述论文是 \cite{Keller-survey}。




\section{约定}
\label{sdga-section-conventions}

\noindent
本章继续采用如下约定：{\it 环}是指有单位元 $1$ 的交换环。若 $R$
是环，则{\it $R$-代数 $A$}是一个 $R$-模 $A$，连同一个
$R$-双线性映射 $A \times A \to A$（乘法），使得乘法结合且具有
单位元。换言之，这是有单位元的结合 $R$-代数，并且结构映射
$R \to A$ 的像位于 $A$ 的中心中。








\section{分次代数层}
\label{sdga-section-ga}

\noindent
请跳过本节。

\begin{definition}
\label{sdga-definition-ga}
设 $(\mathcal{C}, \mathcal{O})$ 是环化位点。$(\mathcal{C},\mathcal O)$
上的{\it 分次 $\mathcal{O}$-代数层}或{\it 分次代数层}，由以
$n \in \mathbf{Z}$ 为指标的一族 $\mathcal{O}$-模
$\mathcal{A}^n$ 以及 $\mathcal{O}$-双线性映射
$$
\mathcal{A}^n \times \mathcal{A}^m \to \mathcal{A}^{n + m},\quad
(a, b) \longmapsto ab
$$
给出；这些映射称为乘法映射，并满足下列性质：
\begin{enumerate}
\item 乘法是结合的；以及
\item $\mathcal{A}^0$ 有一个全局截面 $1$，它是乘法的双边单位元。
\end{enumerate}
常把这样的结构记作 $\mathcal{A}$。一个{\it 分次
$\mathcal{O}$-代数同态} $f : \mathcal{A} \to \mathcal{B}$，
是与乘法映射相容的一族 $\mathcal{O}$-模映射
$f^n : \mathcal{A}^n \to \mathcal{B}^n$
。
\end{definition}

\noindent
给定一个分次 $\mathcal{O}$-代数 $\mathcal{A}$ 以及一个对象
$U \in \Ob(\mathcal{C})$，使用记号
$$
\mathcal{A}(U) =
\Gamma(U, \mathcal{A}) =
\bigoplus\nolimits_{n \in \mathbf{Z}} \mathcal{A}^n(U)
$$
这是一个分次 $\mathcal{O}(U)$-代数。

\begin{remark}
\label{sdga-remark-functoriality-ga}
设 $(f, f^\sharp) : (\Sh(\mathcal{C}), \mathcal{O}_\mathcal{C})
\to (\Sh(\mathcal{D}), \mathcal{O}_\mathcal{D})$
是环化拓扑斯的态射。则有：
\begin{enumerate}
\item 设 $\mathcal{A}$ 是分次 $\mathcal{O}_\mathcal{C}$-代数。
$\mathcal{A}$ 的乘法映射诱导乘法映射
$f_*\mathcal{A}^n \times f_*\mathcal{A}^m \to f_*\mathcal{A}^{n + m}$
，并且经由 $f^\sharp$，可把它们看作
$\mathcal{O}_\mathcal{D}$-双线性映射。把由此得到的分次
$\mathcal{O}_\mathcal{D}$-代数记作 $f_*\mathcal{A}$。
\item 设 $\mathcal{B}$ 是分次 $\mathcal{O}_\mathcal{D}$-代数。
$\mathcal{B}$ 的乘法映射诱导乘法映射
$f^*\mathcal{B}^n \times f^*\mathcal{B}^m \to f^*\mathcal{B}^{n + m}$
，并且利用 $f^\sharp$，可把它们看作
$\mathcal{O}_\mathcal{C}$-双线性映射。把由此得到的分次
$\mathcal{O}_\mathcal{C}$-代数记作 $f^*\mathcal{B}$。
\item 分次 $\mathcal{O}_\mathcal{C}$-代数同态
$f^*\mathcal{B} \to \mathcal{A}$ 的集合，与分次
$\mathcal{O}_\mathcal{C}$-代数同态
$\mathcal{B} \to f_*\mathcal{A}$ 的集合一一对应。
\end{enumerate}
第 (3) 部分立即由模层上的 $f^*$ 与 $f_*$ 之间通常的伴随得到。
\end{remark}




\section{分次模层}
\label{sdga-section-graded-modules}

\noindent
请跳过本节。

\begin{definition}
\label{sdga-definition-gm}
设 $(\mathcal{C}, \mathcal{O})$ 是环化位点，$\mathcal{A}$ 是
$(\mathcal{C}, \mathcal{O})$ 上的分次代数层。一个（右）
{\it 分次 $\mathcal{A}$-模}，或 $\mathcal{A}$ 上的（右）
{\it 分次模}，由以 $n \in \mathbf{Z}$ 为指标的一族
$\mathcal{O}$-模 $\mathcal{M}^n$ 以及 $\mathcal{O}$-双线性映射
$$
\mathcal{M}^n \times \mathcal{A}^m \to \mathcal{M}^{n + m},\quad
(x, a) \longmapsto xa
$$
给出；这些映射称为乘法映射，并满足下列性质：
\begin{enumerate}
\item 乘法满足 $(xa)a' = x(aa')$；
\item 对每个 $n$，$\mathcal{A}^0$ 的单位截面 $1$
在 $\mathcal{M}^n$ 上的作用都是恒等作用。
\end{enumerate}
我们常说“设 $\mathcal{M}$ 是一个分次 $\mathcal{A}$-模”来表示上述情形。
一个{\it 分次 $\mathcal{A}$-模同态}
$f : \mathcal{M} \to \mathcal{N}$，是与乘法映射相容的一族
$\mathcal{O}$-模映射 $f^n : \mathcal{M}^n \to \mathcal{N}^n$。
（右）分次 $\mathcal{A}$-模构成的范畴记作
$\textit{Mod}(\mathcal{A})$。
\end{definition}

\noindent
可以完全类似地定义{\it 左分次模}，但本章默认所说的模都是右模。

\medskip\noindent
给定一个分次 $\mathcal{A}$-模 $\mathcal{M}$ 和一个对象
$U \in \Ob(\mathcal{C})$，我们使用记号
$$
\mathcal{M}(U) =
\Gamma(U, \mathcal{M}) =
\bigoplus\nolimits_{n \in \mathbf{Z}} \mathcal{M}^n(U)
$$
这是一个（右）分次 $\mathcal{A}(U)$-模。

\begin{lemma}
\label{sdga-lemma-gm-abelian}
设 $(\mathcal{C}, \mathcal{O})$ 是环化位点，$\mathcal{A}$ 是分次
$\mathcal{O}$-代数。范畴 $\textit{Mod}(\mathcal{A})$ 是阿贝尔范畴，
并具有下列性质：
\begin{enumerate}
\item $\textit{Mod}(\mathcal{A})$ 有任意直和；
\item $\textit{Mod}(\mathcal{A})$ 有任意余极限；
\item $\textit{Mod}(\mathcal{A})$ 中的滤过余极限是正合的；
\item $\textit{Mod}(\mathcal{A})$ 有任意积；
\item $\textit{Mod}(\mathcal{A})$ 有任意极限。
\end{enumerate}
函子
$$
\textit{Mod}(\mathcal{A}) \longrightarrow \textit{Mod}(\mathcal{O}),\quad
\mathcal{M} \longmapsto \mathcal{M}^n
$$
把一个分次 $\mathcal{A}$-模映到其第 $n$ 项；它与所有极限和余极限交换。
\end{lemma}

\noindent
该引理说明，我们可以逐项取极限和余极限。它还说明（或者，如果你愿意，
蕴含）遗忘函子
$$
\textit{Mod}(\mathcal{A})
\longrightarrow
\text{分次 }\mathcal{O}\text{-模}
$$
与所有极限和余极限交换。

\begin{proof}
把引理陈述中的函子记作
$\text{gr}^n : \textit{Mod}(\mathcal{A}) \to \textit{Mod}(\mathcal{O})$
。考虑分次 $\mathcal{A}$-模同态
$f : \mathcal{M} \to \mathcal{N}$。把 $f$ 看成分次
$\mathcal{O}$-模映射时，它的核与余核还带有定义
\ref{sdga-definition-gm} 中的乘法映射。因此，它们也是
$\textit{Mod}(\mathcal{A})$ 中的核与余核。由此
$\textit{Mod}(\mathcal{A})$ 是阿贝尔范畴，并且取核与余核和
$\text{gr}^n$ 交换。

\medskip\noindent
为了证明极限和余极限存在，只需证明积和直和存在；参见《范畴》引理
\ref{categories-lemma-limits-products-equalizers} 和
\ref{categories-lemma-colimits-coproducts-coequalizers}。
同样的引理还表明，只要 $\text{gr}^n$ 与直和及积交换，
就可推出它与极限及余极限交换。

\medskip\noindent
设 $\mathcal{M}_t$（$t \in T$）是一族分次 $\mathcal{A}$-模。
可以考虑这样一个分次 $\mathcal{A}$-模：其次数 $n$ 项是
$\bigoplus_{t \in T} \mathcal{M}_t^n$（乘法映射是显然的）。
读者容易验证，这是 $\textit{Mod}(\mathcal{A})$ 中的直和。
积的情形类似。

\medskip\noindent
注意，对每个 $n$，$\text{gr}^n$ 都是正合函子；而
$\textit{Mod}(\mathcal{A})$ 中的复形
$\mathcal{M}_1 \to \mathcal{M}_2 \to \mathcal{M}_3$ 正合，当且仅当
$\text{gr}^n\mathcal{M}_1 \to \text{gr}^n\mathcal{M}_2 \to
\text{gr}^n\mathcal{M}_3$ 对每个 $n$ 都在
$\textit{Mod}(\mathcal{O})$ 中正合。因此，由于
$\textit{Mod}(\mathcal{O})$ 中的滤过余极限是正合的，我们得到 (3)；
它是格罗滕迪克阿贝尔范畴，参见《位点上的上同调》第
\ref{sites-cohomology-section-unbounded} 节。
\end{proof}






\section{分次模层的分次范畴}
\label{sdga-section-gm-gr-cat}

\noindent
请跳过本节。本节是《微分分次代数》例
\ref{dga-example-gm-gr-cat} 的类似物。关于分次范畴的约定，
请参见《微分分次代数》第 \ref{dga-section-graded} 节。

\medskip\noindent
设 $(\mathcal{C}, \mathcal{O})$ 是环化位点，$\mathcal{A}$ 是
$(\mathcal{C}, \mathcal{O})$ 上的分次代数层。我们将构造
$R = \Gamma(\mathcal{C}, \mathcal{O})$
上的分次范畴 $\textit{Mod}^{gr}(\mathcal{A})$，使其对应范畴
$(\textit{Mod}^{gr}(\mathcal{A}))^0$ 是分次 $\mathcal{A}$-模范畴。
取右分次 $\mathcal{A}$-模作为 $\textit{Mod}^{gr}(\mathcal{A})$
的对象（见第 \ref{sdga-section-graded-modules} 节）。给定分次
$\mathcal{A}$-模 $\mathcal{L}$ 和 $\mathcal{M}$，置
$$
\Hom_{\textit{Mod}^{gr}(\mathcal{A})}(\mathcal{L}, \mathcal{M}) =
\bigoplus\nolimits_{n \in \mathbf{Z}} \Hom^n(\mathcal{L}, \mathcal{M})
$$
，其中 $\Hom^n(\mathcal{L}, \mathcal{M})$ 是次数 $n$ 齐次的右
$\mathcal{A}$-模映射 $f : \mathcal{L} \to \mathcal{M}$ 所成的集合。
更确切地说，$f$ 由对 $i \in \mathbf{Z}$ 给出的一族映射
$f : \mathcal{L}^i \to \mathcal{M}^{i + n}$ 组成，并与乘法映射相容。
用分量表示，
$$
\Hom^n(\mathcal{L}, \mathcal{M})
\subset
\prod\nolimits_{p + q = n}
\Hom_\mathcal{O}(\mathcal{L}^{-q}, \mathcal{M}^p)
$$
（注意指标的反转）是由满足
$$
f_{p, q}(m a) = f_{p - i, q + i}(m)a
$$
的 $f = (f_{p, q})$ 组成的子集；这里 $a$ 是 $\mathcal{A}^i$
的局部截面，$m$ 是 $\mathcal{L}^{-q-i}$ 的局部截面。
对分次 $\mathcal{A}$-模 $\mathcal{K}$、$\mathcal{L}$、
$\mathcal{M}$，通过映射
$$
\Hom^m(\mathcal{L}, \mathcal{M}) \times
\Hom^n(\mathcal{K}, \mathcal{L}) \longrightarrow
\Hom^{n + m}(\mathcal{K}, \mathcal{M})
$$
定义 $\textit{Mod}^{gr}(\mathcal{A})$ 中的复合；它就是右
$\mathcal{A}$-模映射的通常复合：$(g,f)\mapsto g\circ f$。





\section{分次模层的张量积}
\label{sdga-section-tensor-product}

\noindent
请跳过本节。本节是《微分分次代数》第
\ref{dga-section-tensor-product} 节部分内容的类似物。

\medskip\noindent
设 $(\mathcal{C}, \mathcal{O})$ 是环化位点，$\mathcal{A}$ 是
$(\mathcal{C},\mathcal{O})$ 上的分次代数层。设 $\mathcal{M}$ 是右分次
$\mathcal{A}$-模，$\mathcal{N}$ 是左分次 $\mathcal{A}$-模。
定义{\it 张量积}
$\mathcal{M} \otimes_\mathcal{A} \mathcal{N}$
为这样的分次 $\mathcal{O}$-模：它的次数 $n$ 项是
$$
(\mathcal{M} \otimes_\mathcal{A} \mathcal{N})^n =
\Coker\left(
\bigoplus\nolimits_{r + s + t = n} \mathcal{M}^r \otimes_\mathcal{O}
\mathcal{A}^s \otimes_\mathcal{O} \mathcal{N}^t
\longrightarrow
\bigoplus\nolimits_{p + q = n} \mathcal{M}^p \otimes_\mathcal{O} \mathcal{N}^q
\right)
$$
，其中该映射把
$\mathcal{M}^r\otimes_\mathcal{O}\mathcal{A}^s
\otimes_\mathcal{O}\mathcal{N}^t$ 的局部截面
$x\otimes a\otimes y$ 映到 $xa\otimes y-x\otimes ay$。
按此定义，
$(\mathcal{M} \otimes_\mathcal{A} \mathcal{N})^n$
是预层
$U \mapsto (\mathcal{M}(U) \otimes_{\mathcal{A}(U)} \mathcal{N}(U))^n$
的层化；这里分次模的张量积按《微分分次代数》第
\ref{dga-section-tensor-product} 节中的定义。

\medskip\noindent
固定左分次 $\mathcal{A}$-模 $\mathcal{N}$，便得到函子
$$
- \otimes_\mathcal{A} \mathcal{N} :
\textit{Mod}(\mathcal{A})
\longrightarrow
\text{Gr}(\textit{Mod}(\mathcal{O})) =
\text{分次 }\mathcal{O}\text{-模}
$$
关于记号 $\text{Gr}(-)$，请参见《同调代数》定义
\ref{homology-definition-graded}。分次 $\mathcal{O}$-模的分次范畴记作
$\text{Gr}^{gr}(\textit{Mod}(\mathcal{O}))$；见《微分分次代数》例
\ref{dga-example-graded-category-graded-objects}。
上述函子可以提升为分次范畴之间的函子
$$
- \otimes_\mathcal{A} \mathcal{N} :
\textit{Mod}^{gr}(\mathcal{A})
\longrightarrow
\text{Gr}^{gr}(\textit{Mod}(\mathcal{O}))
$$
：把从 $\mathcal{M}$ 到 $\mathcal{M}'$ 的次数 $n$ 同态映到
从 $\mathcal{M}\otimes_\mathcal{A}\mathcal{N}$ 到
$\mathcal{M}'\otimes_\mathcal{A}\mathcal{N}$ 的诱导次数 $n$ 映射。






\section{分次模层的内 Hom}
\label{sdga-section-internal-hom-graded}

\noindent
我们强烈建议读者跳过本节。

\medskip\noindent
我们需要第 \ref{sdga-section-gm-gr-cat} 节中构造的层化版本。
设 $(\mathcal{C}, \mathcal{O})$、$\mathcal{A}$、
$\mathcal{M}$、$\mathcal{L}$ 如第 \ref{sdga-section-gm-gr-cat} 节所述。定义
$$
\SheafHom^{gr}_\mathcal{A}(\mathcal{M}, \mathcal{L})
$$
为这样的分次 $\mathcal{O}$-模：其次数 $n$ 项
$$
\SheafHom^n_\mathcal{A}(\mathcal{M}, \mathcal{L})
\subset
\prod\nolimits_{p + q = n}
\SheafHom_\mathcal{O}(\mathcal{L}^{-q}, \mathcal{M}^p)
$$
是由满足
$$
f_{p, q}(m a) = f_{p - i, q + i}(m)a
$$
的局部截面 $f=(f_{p,q})$ 组成的子层；这里 $a$ 是
$\mathcal{A}^i$ 的局部截面，$m$ 是 $\mathcal{L}^{-q-i}$
的局部截面。如同第 \ref{sdga-section-gm-gr-cat} 节，有复合映射
$$
\SheafHom^{gr}_\mathcal{A}(\mathcal{L}, \mathcal{M}) \otimes_\mathcal{O}
\SheafHom^{gr}_\mathcal{A}(\mathcal{K}, \mathcal{L})
\longrightarrow
\SheafHom^{gr}_\mathcal{A}(\mathcal{K}, \mathcal{M})
$$
，其中左端是第 \ref{sdga-section-tensor-product} 节所定义的分次
$\mathcal{O}$-模张量积。该映射由复合映射
$$
\SheafHom^m_\mathcal{A}(\mathcal{L}, \mathcal{M}) \otimes_\mathcal{O}
\SheafHom^n_\mathcal{A}(\mathcal{K}, \mathcal{L}) \longrightarrow
\SheafHom^{n + m}_\mathcal{A}(\mathcal{K}, \mathcal{M})
$$
给出；后者（局部地）按通常复合定义。

\medskip\noindent
按这些定义，有
$$
\Hom_{\textit{Mod}^{gr}(\mathcal{A})}(\mathcal{L}, \mathcal{M}) =
\Gamma(\mathcal{C}, \SheafHom^{gr}_\mathcal{A}(\mathcal{L}, \mathcal{M}))
$$
作为分次 $R$-模的等式，并且它与复合相容。







\section{分次双模层与张量--Hom 伴随}
\label{sdga-section-graded-bimodules}

\noindent
请跳过本节。

\begin{definition}
\label{sdga-definition-bimodule}
设 $(\mathcal{C},\mathcal{O})$ 是环化位点，$\mathcal{A}$ 和
$\mathcal{B}$ 是 $(\mathcal{C},\mathcal{O})$ 上的分次代数层。
一个{\it 分次 $(\mathcal{A},\mathcal{B})$-双模}，由以
$n\in\mathbf{Z}$ 为指标的一族 $\mathcal{O}$-模
$\mathcal{M}^n$ 以及 $\mathcal{O}$-双线性映射
$$
\mathcal{M}^n \times \mathcal{B}^m \to \mathcal{M}^{n + m},\quad
(x, b) \longmapsto xb
$$
和
$$
\mathcal{A}^n \times \mathcal{M}^m \to \mathcal{M}^{n + m},\quad
(a, x) \longmapsto ax
$$
给出；这些映射称为乘法映射，并满足下列性质：
\begin{enumerate}
\item 乘法满足 $a(a'x)=(aa')x$ 和 $(xb)b'=x(bb')$；
\item $(ax)b = a(xb)$；
\item $\mathcal{A}^0$ 的单位截面 $1$ 通过乘法作恒等作用；
\item $\mathcal{B}^0$ 的单位截面 $1$ 通过乘法作恒等作用。
\end{enumerate}
我们常把这样的结构记作 $\mathcal{M}$。一个{\it 分次
$(\mathcal{A},\mathcal{B})$-双模同态}
$f:\mathcal{M}\to\mathcal{N}$，是与乘法映射相容的一族
$\mathcal{O}$-模映射 $f^n:\mathcal{M}^n\to\mathcal{N}^n$。
\end{definition}

\noindent
给定分次 $(\mathcal{A},\mathcal{B})$-双模 $\mathcal{M}$ 和对象
$U\in\Ob(\mathcal{C})$，我们使用记号
$$
\mathcal{M}(U) =
\Gamma(U, \mathcal{M}) =
\bigoplus\nolimits_{n \in \mathbf{Z}} \mathcal{M}^n(U)
$$
这是分次 $(\mathcal{A}(U),\mathcal{B}(U))$-双模。

\medskip\noindent
设 $(\mathcal{C},\mathcal{O})$ 是环化位点，$\mathcal{A}$ 和
$\mathcal{B}$ 是 $(\mathcal{C},\mathcal{O})$ 上的分次代数层。
设 $\mathcal{M}$ 是右分次 $\mathcal{A}$-模，$\mathcal{N}$ 是分次
$(\mathcal{A},\mathcal{B})$-双模。在这种情形下，第
\ref{sdga-section-tensor-product} 节定义的分次张量积
$$
\mathcal{M} \otimes_\mathcal{A} \mathcal{N}
$$
是右分次 $\mathcal{B}$-模，其乘法映射是显然的。该构造定义一个函子
以及一个分次范畴之间的函子
$$
\otimes_\mathcal{A} \mathcal{N} :
\textit{Mod}(\mathcal{A})
\longrightarrow
\textit{Mod}(\mathcal{B})
\quad\text{且}\quad
\otimes_\mathcal{A} \mathcal{N} :
\textit{Mod}^{gr}(\mathcal{A})
\longrightarrow
\textit{Mod}^{gr}(\mathcal{B})
$$
：把从 $\mathcal{M}$ 到 $\mathcal{M}'$ 的次数 $n$ 同态映到
从 $\mathcal{M}\otimes_\mathcal{A}\mathcal{N}$ 到
$\mathcal{M}'\otimes_\mathcal{A}\mathcal{N}$ 的诱导次数 $n$ 映射。

\medskip\noindent
设 $(\mathcal{C},\mathcal{O})$ 是环化位点，$\mathcal{A}$ 和
$\mathcal{B}$ 是 $(\mathcal{C},\mathcal{O})$ 上的分次代数层。
设 $\mathcal{N}$ 是分次 $(\mathcal{A},\mathcal{B})$-双模，
$\mathcal{L}$ 是右分次 $\mathcal{B}$-模。在这种情形下，第
\ref{sdga-section-internal-hom-graded} 节定义的分次内 Hom
$$
\SheafHom_\mathcal{B}^{gr}(\mathcal{N}, \mathcal{L})
$$
是右分次 $\mathcal{A}$-模，其乘法映射为\footnote{我们在这里约定，
此式不涉及任何符号。}
$$
\SheafHom^n_\mathcal{B}(\mathcal{N}, \mathcal{L})
\times \mathcal{A}^m
\longrightarrow
\SheafHom^{n + m}_\mathcal{B}(\mathcal{N}, \mathcal{L})
$$
：它把 $U$ 上
$\SheafHom^n_\mathcal{B}(\mathcal{N},\mathcal{L})$ 的截面
$f=(f_{p,q})$ 和 $U$ 上 $\mathcal{A}^m$ 的截面 $a$，映到
$U$ 上 $\SheafHom^{n+m}_\mathcal{B}(\mathcal{N},\mathcal{L})$
的截面 $fa$；后者定义为映射族
$$
\mathcal{N}^{-q - m}|_U \xrightarrow{a \cdot -}
\mathcal{N}^{-q}|_U \xrightarrow{f_{p, q}}
\mathcal{M}^p|_U
$$
。略去其良定义性的验证。该构造定义一个函子以及一个分次范畴之间的函子
$$
\SheafHom_\mathcal{B}^{gr}(\mathcal{N}, -) :
\textit{Mod}(\mathcal{B})
\longrightarrow
\textit{Mod}(\mathcal{A})
\quad\text{且}\quad
\SheafHom_\mathcal{B}^{gr}(\mathcal{N}, -) :
\textit{Mod}^{gr}(\mathcal{B})
\longrightarrow
\textit{Mod}^{gr}(\mathcal{A})
$$
：把从 $\mathcal{L}$ 到 $\mathcal{L}'$ 的次数 $n$ 同态映到
从 $\SheafHom_\mathcal{B}^{gr}(\mathcal{N},\mathcal{L})$ 到
$\SheafHom_\mathcal{B}^{gr}(\mathcal{N},\mathcal{L}')$
的诱导次数 $n$ 映射。

\begin{lemma}
\label{sdga-lemma-tensor-hom-adjunction-gr}
设 $(\mathcal{C},\mathcal{O})$ 是环化位点，$\mathcal{A}$ 和
$\mathcal{B}$ 是 $(\mathcal{C},\mathcal{O})$ 上的分次代数层。
设 $\mathcal{M}$ 是右分次 $\mathcal{A}$-模，$\mathcal{N}$ 是分次
$(\mathcal{A},\mathcal{B})$-双模，$\mathcal{L}$ 是右分次
$\mathcal{B}$-模。采用上述约定，有
$$
\Hom_{\textit{Mod}^{gr}(\mathcal{B})}(
\mathcal{M} \otimes_\mathcal{A} \mathcal{N}, \mathcal{L}) =
\Hom_{\textit{Mod}^{gr}(\mathcal{A})}(
\mathcal{M}, \SheafHom_\mathcal{B}^{gr}(\mathcal{N}, \mathcal{L}))
$$
以及
$$
\SheafHom_\mathcal{B}^{gr}(
\mathcal{M} \otimes_\mathcal{A} \mathcal{N}, \mathcal{L}) =
\SheafHom_\mathcal{A}^{gr}(
\mathcal{M}, \SheafHom_\mathcal{B}^{gr}(\mathcal{N}, \mathcal{L}))
$$
；这些等式关于 $\mathcal{M}$、$\mathcal{N}$、$\mathcal{L}$
都具有函子性。
\end{lemma}

\begin{proof}
略。提示：把两端都解释为 $\mathcal{A}$-双线性的分次映射
$\psi : \mathcal{M} \times \mathcal{N} \to \mathcal{L}$
即可；这些映射在右侧关于 $\mathcal{B}$ 是线性的。
\end{proof}

\noindent
设 $(\mathcal{C},\mathcal{O})$ 是环化位点，$\mathcal{A}$ 和
$\mathcal{B}$ 是 $(\mathcal{C},\mathcal{O})$ 上的分次代数层。
作为上述情形的特例，设给定分次 $\mathcal{O}$-代数同态
$\varphi:\mathcal{A}\to\mathcal{B}$。于是得到一个函子和一个
分次范畴之间的函子
$$
\otimes_{\mathcal{A}, \varphi} \mathcal{B} :
\textit{Mod}(\mathcal{A})
\longrightarrow
\textit{Mod}(\mathcal{B})
\quad\text{且}\quad
\otimes_{\mathcal{A}, \varphi} \mathcal{B} :
\textit{Mod}^{gr}(\mathcal{A})
\longrightarrow
\textit{Mod}^{gr}(\mathcal{B})
$$
。另一方面，有限制标量函子
$$
res_\varphi :
\textit{Mod}(\mathcal{B})
\longrightarrow
\textit{Mod}(\mathcal{A})
\quad\text{且}\quad
res_\varphi :
\textit{Mod}^{gr}(\mathcal{B})
\longrightarrow
\textit{Mod}^{gr}(\mathcal{A})
$$
。由上一引理可以证明这些函子彼此伴随（如限制标量与基变换的通常情形）。
具体地，把 $\mathcal{B}$ 视为分次 $(\mathcal{A},\mathcal{B})$-双模时
记作 ${}_\mathcal{A}\mathcal{B}_\mathcal{B}$。则对任意右分次
$\mathcal{B}$-模 $\mathcal{L}$，有
$$
\SheafHom_\mathcal{B}^{gr}({}_\mathcal{A}\mathcal{B}_\mathcal{B}, \mathcal{L})
= res_\varphi(\mathcal{L})
$$
作为右分次 $\mathcal{A}$-模的等式。因此，引理
\ref{sdga-lemma-tensor-hom-adjunction-gr} 给出函子性同构
$$
\Hom_{\textit{Mod}^{gr}(\mathcal{B})}(
\mathcal{M} \otimes_{\mathcal{A}, \varphi} \mathcal{B}, \mathcal{L}) =
\Hom_{\textit{Mod}^{gr}(\mathcal{A})}(
\mathcal{M}, res_\varphi(\mathcal{L}))
$$
。若从上下文可知 $\varphi$，通常在此公式中省略对它的依赖。
以同样方式得到分次 $\mathcal{O}$-模的等式
$$
\SheafHom^{gr}_\mathcal{B}(
\mathcal{M} \otimes_\mathcal{A} \mathcal{B}, \mathcal{L}) =
\SheafHom_\mathcal{A}^{gr}(\mathcal{M}, \mathcal{L})
$$
。




\section{分次模层的拉回与直像}
\label{sdga-section-functoriality-graded}

\noindent
我们建议读者跳过本节。

\medskip\noindent
设 $(f, f^\sharp) : (\Sh(\mathcal{C}), \mathcal{O}_\mathcal{C})
\to (\Sh(\mathcal{D}), \mathcal{O}_\mathcal{D})$
是环化拓扑斯的态射。设 $\mathcal{A}$ 是分次
$\mathcal{O}_\mathcal{C}$-代数，$\mathcal{B}$ 是分次
$\mathcal{O}_\mathcal{D}$-代数。设给定分次
$f^{-1}\mathcal{O}_\mathcal{D}$-代数同态
$$
\varphi : f^{-1}\mathcal{B} \to \mathcal{A}
$$
。由限制标量与扩张标量的伴随，这等价于一个分次
$\mathcal{O}_\mathcal{C}$-代数同态
$\varphi:f^*\mathcal{B}\to\mathcal{A}$；也等价地，可以把
$\varphi$ 看作分次 $\mathcal{O}_\mathcal{D}$-代数同态
$$
\varphi : \mathcal{B} \to f_*\mathcal{A}
$$
。参见评注 \ref{sdga-remark-functoriality-ga}。

\medskip\noindent
先定义函子
$$
f_* :
\textit{Mod}(\mathcal{A})
\longrightarrow
\textit{Mod}(\mathcal{B})
$$
。给定分次 $\mathcal{A}$-模 $\mathcal{M}$，定义
$f_*\mathcal{M}$ 为这样的分次 $\mathcal{B}$-模：其次数 $n$ 项是
$f_*\mathcal{M}^n$，乘法取为
$$
f_*\mathcal{M}^n \times \mathcal{B}^m
\xrightarrow{(\text{id}, \varphi^m)}
f_*\mathcal{M}^n \times f_*\mathcal{A}^m
\xrightarrow{f_*\mu_{n, m}}
f_*\mathcal{M}^{n + m}
$$
，其中 $\mu_{n,m}:\mathcal{M}^n\times\mathcal{A}^m
\to\mathcal{M}^{n+m}$ 是 $\mathcal{M}$ 作为 $\mathcal{A}$-模的
乘法映射。这里使用了 $f_*$ 与积交换这一事实。该构造关于
$\mathcal{M}$ 显然具有函子性，因而得到所需函子。

\medskip\noindent
再定义函子
$$
f^* :
\textit{Mod}(\mathcal{B})
\longrightarrow
\textit{Mod}(\mathcal{A})
$$
为下列函子的复合：
$$
\textit{Mod}(\mathcal{B})
\xrightarrow{f^{-1}}
\textit{Mod}(f^{-1}\mathcal{B})
\xrightarrow{ - \otimes_{f^{-1}\mathcal{B}} \mathcal{A}}
\textit{Mod}(\mathcal{A})
$$
首先，给定分次 $\mathcal{B}$-模 $\mathcal{N}$，定义
$f^{-1}\mathcal{N}$ 为这样的分次 $f^{-1}\mathcal{B}$-模：
其次数 $n$ 项是 $f^{-1}\mathcal{N}^n$，乘法取为
$$
f^{-1}\nu_{n, m} :
f^{-1}\mathcal{N}^n \times f^{-1}\mathcal{B}^m
\longrightarrow
f^{-1}\mathcal{N}^{n + m}
$$
，其中 $\nu_{n,m}:\mathcal{N}^n\times\mathcal{B}^m
\to\mathcal{N}^{n+m}$ 是 $\mathcal{N}$ 作为 $\mathcal{B}$-模的
乘法映射。这里使用了 $f^{-1}$ 与积交换这一事实。该构造关于
$\mathcal{N}$ 显然具有函子性，从而得到函子 $f^{-1}$。
然后利用第 \ref{sdga-section-graded-bimodules} 节关于张量积的讨论，
定义函子
$$
- \otimes_{f^{-1}\mathcal{B}} \mathcal{A} :
\textit{Mod}(f^{-1}\mathcal{B})
\longrightarrow
\textit{Mod}(\mathcal{A})
$$
。最后，如上所预示，置
$$
f^*\mathcal{N} =
f^{-1}\mathcal{N} \otimes_{f^{-1}\mathcal{B}, \varphi} \mathcal{A}
$$
。

\medskip\noindent
函子 $f_*$ 和 $f^*$ 很容易提升为分次范畴之间的函子
$$
f_* :
\textit{Mod}^{gr}(\mathcal{A})
\longrightarrow
\textit{Mod}^{gr}(\mathcal{B})
\quad\text{且}\quad
f^* :
\textit{Mod}^{gr}(\mathcal{B})
\longrightarrow
\textit{Mod}^{gr}(\mathcal{A})
$$
；它们在底层对象上的作用不变，而在次数 $n$ 的齐次态射上由上述构造的
函子性定义。

\begin{lemma}
\label{sdga-lemma-adjunction-push-pull-gr}
在上述情形下，有
$$
\Hom_{\textit{Mod}^{gr}(\mathcal{B})}(
\mathcal{N}, f_*\mathcal{M}) =
\Hom_{\textit{Mod}^{gr}(\mathcal{A})}(
f^*\mathcal{N}, \mathcal{M})
$$
\end{lemma}

\begin{proof}
略。提示：先利用阿贝尔群层上 $f^{-1}$ 与 $f_*$ 的伴随，证明
$f^{-1}$ 与 $f_*$ 作为 $\textit{Mod}(\mathcal{B})$ 和
$\textit{Mod}(f^{-1}\mathcal{B})$ 之间的函子彼此伴随。
再使用第 \ref{sdga-section-graded-bimodules} 节给出的基变换与限制标量的伴随。
\end{proof}





\section{局部化与分次模层}
\label{sdga-section-localize-graded}

\noindent
我们建议读者跳过本节。

\medskip\noindent
设 $(\mathcal{C},\mathcal{O})$ 是环化位点，$U\in\Ob(\mathcal{C})$。
记
$$
j :
(\Sh(\mathcal{C}/U), \mathcal{O}_U)
\longrightarrow
(\Sh(\mathcal{C}), \mathcal{O})
$$
为相应的局部化态射（《位点上的模》第
\ref{sites-modules-section-localize} 节）。下面将使用如下事实：
对 $\mathcal{O}_U$-模 $\mathcal{M}_i$（$i=1,2$）和
$\mathcal{O}$-模 $\mathcal{A}$，有典范映射
$$
j_! :
\Hom_{\mathcal{O}_U}(
\mathcal{M}_1 \otimes_{\mathcal{O}_U} \mathcal{A}|_U, \mathcal{M}_2)
\longrightarrow
\Hom_\mathcal{O}(
j_!\mathcal{M}_1 \otimes_\mathcal{O} \mathcal{A}, j_!\mathcal{M}_2)
$$
。事实上，由《位点上的模》引理
\ref{sites-modules-lemma-j-shriek-and-tensor}，有
$j_!(\mathcal{M}_1 \otimes_{\mathcal{O}_U} \mathcal{A}|_U) =
j_!\mathcal{M}_1 \otimes_\mathcal{O} \mathcal{A}$。

\medskip\noindent
设 $\mathcal{A}$ 是分次 $\mathcal{O}$-代数。用
$\mathcal{A}_U$ 表示 $\mathcal{A}$ 在 $\mathcal{C}/U$ 上的限制；
换言之，$\mathcal{A}_U=j^*\mathcal{A}=j^{-1}\mathcal{A}$。
第 \ref{sdga-section-functoriality-graded} 节已经构造伴随函子
$$
j_* :
\textit{Mod}^{gr}(\mathcal{A}_U)
\longrightarrow
\textit{Mod}^{gr}(\mathcal{A})
\quad\text{且}\quad
j^* :
\textit{Mod}^{gr}(\mathcal{A})
\longrightarrow
\textit{Mod}^{gr}(\mathcal{A}_U)
$$
，其中 $j^*$ 是 $j_*$ 的左伴随。我们断言，此外还有正合函子
$$
j_! :
\textit{Mod}^{gr}(\mathcal{A}_U)
\longrightarrow
\textit{Mod}^{gr}(\mathcal{A})
$$
，它是 $j^*$ 的左伴随。具体地，给定分次 $\mathcal{A}_U$-模
$\mathcal{M}$，定义 $j_!\mathcal{M}$ 为这样的分次 $\mathcal{A}$-模：
其次数 $n$ 项为 $j_!\mathcal{M}^n$，乘法映射取为
$$
j_!\mu_{n, m} :
j_!\mathcal{M}^n \times \mathcal{A}^m \to
j_!\mathcal{M}^{n + m}
$$
，其中 $\mu_{m,n}:\mathcal{M}^n\times\mathcal{A}^m
\to\mathcal{M}^{n+m}$ 是给定的乘法映射。给定分次
$\mathcal{A}_U$-模之间次数 $n$ 的齐次映射
$f:\mathcal{M}\to\mathcal{M}'$，得到次数 $n$ 的齐次映射
$j_!f:j_!\mathcal{M}\to j_!\mathcal{M}'$。这样便得到所需函子。

\begin{lemma}
\label{sdga-lemma-extension-by-zero-graded}
在上述情形下，有
$$
\Hom_{\textit{Mod}^{gr}(\mathcal{A})}(
j_!\mathcal{M}, \mathcal{N}) =
\Hom_{\textit{Mod}^{gr}(\mathcal{A}_U)}(
\mathcal{M}, j^*\mathcal{N})
$$
\end{lemma}

\begin{proof}
由《位点上的模》第 \ref{sites-modules-section-localize} 节的讨论，
$\mathcal{O}$-模上的函子 $j_!$ 与 $j^*$ 彼此伴随。因此，若只看
$\mathcal{O}$-模结构，我们知道
$$
\Hom_{\text{Gr}^{gr}(\textit{Mod}(\mathcal{O}))}(
j_!\mathcal{M}, \mathcal{N}) =
\Hom_{\text{Gr}^{gr}(\textit{Mod}(\mathcal{O}_U))}(
\mathcal{M}, j^*\mathcal{N})
$$
（回忆，$\text{Gr}^{gr}(\textit{Mod}(\mathcal{O}))$ 表示分次
$\mathcal{O}$-模的分次范畴）。还须检验这些等同把左端的
$\mathcal{A}$-模映射对应到右端的 $\mathcal{A}_U$-模映射。
为此，设 $\mathcal{O}_U$-线性映射
$f^n : \mathcal{M}^n \to j^*\mathcal{N}^{n + d}$
对应于 $\mathcal{O}$-线性映射
$g^n : j_!\mathcal{M}^n \to \mathcal{N}^{n + d}$
。只需证明下图交换
$$
\xymatrix{
\mathcal{M}^n \otimes_{\mathcal{O}_U} \mathcal{A}_U^m
\ar[r]_{f^n \otimes 1} \ar[d] &
j^*\mathcal{N}^{n + d} \otimes_{\mathcal{O}_U} \mathcal{A}_U^m \ar[d] \\
\mathcal{M}^{n + m} \ar[r]^{f^{n + m}} &
j^*\mathcal{N}^{n + m + d}
}
$$
当且仅当下图交换：
$$
\xymatrix{
j_!\mathcal{M}^n \otimes_\mathcal{O} \mathcal{A}^m
\ar[r]_{g^n \otimes 1} \ar[d] &
\mathcal{N}^{n + d} \otimes_\mathcal{O} \mathcal{A}_U^m \ar[d] \\
j_!\mathcal{M}^{n + m} \ar[r]^{g^{n + m}} &
\mathcal{N}^{n + m + d}
}
$$
不过，我们知道
\begin{align*}
\Hom_{\mathcal{O}_U}(\mathcal{M}^n \otimes_{\mathcal{O}_U} \mathcal{A}_U^m,
j^*\mathcal{N}^{n + d + m})
& =
\Hom_\mathcal{O}(j_!(\mathcal{M}^n \otimes_{\mathcal{O}_U} \mathcal{A}_U^m),
\mathcal{N}^{n + d + m}) \\
& =
\Hom_\mathcal{O}(j_!\mathcal{M}^n \otimes_\mathcal{O} \mathcal{A}^m,
\mathcal{N}^{n + d + m})
\end{align*}
，这由上面已经使用的《位点上的模》引理
\ref{sites-modules-lemma-j-shriek-and-tensor} 得到。略去如下验证：
第一个群中第一幅图交换的障碍，映到最后一个群中第二幅图交换的障碍。
\end{proof}

\begin{lemma}
\label{sdga-lemma-tensor-with-extension-by-zero}
在上述情形下，设 $\mathcal{M}$ 是右分次 $\mathcal{A}_U$-模，
$\mathcal{N}$ 是左分次 $\mathcal{A}$-模。则
$$
j_!\mathcal{M} \otimes_\mathcal{A} \mathcal{N} =
j_!(\mathcal{M} \otimes_{\mathcal{A}_U} \mathcal{N}|_U)
$$
作为分次 $\mathcal{O}$-模的等式，并且关于 $\mathcal{M}$ 和
$\mathcal{N}$ 具有函子性。
\end{lemma}

\begin{proof}
回忆，$j_!\mathcal{M}\otimes_\mathcal{A}\mathcal{N}$ 的次数 $n$
分量是下列典范映射的余核：
$$
\bigoplus\nolimits_{r + s + t = n}
j_!\mathcal{M}^r \otimes_\mathcal{O}
\mathcal{A}^s \otimes_\mathcal{O}
\mathcal{N}^t
\longrightarrow
\bigoplus\nolimits_{p + q = n}
j_!\mathcal{M}^p \otimes_\mathcal{O} \mathcal{N}^q
$$
。参见第 \ref{sdga-section-tensor-product} 节。由《位点上的模》引理
\ref{sites-modules-lemma-j-shriek-and-tensor}，这等于下列映射的余核：
$$
\bigoplus\nolimits_{r + s + t = n}
j_!(\mathcal{M}^r \otimes_{\mathcal{O}_U}
\mathcal{A}^s|_U \otimes_{\mathcal{O}_U}
\mathcal{N}^t|_U)
\longrightarrow
\bigoplus\nolimits_{p + q = n}
j_!(\mathcal{M}^p \otimes_{\mathcal{O}_U} \mathcal{N}^q|_U)
$$
，结论即得。另一种证明是重做上述引理证明中的米田论证。
\end{proof}








\section{分次模层上的平移函子}
\label{sdga-section-shift}

\noindent
我们强烈建议读者跳过本节。分次代数上的分次模层正是《微分分次代数》
评注 \ref{dga-remark-graded-shift-functors} 所讨论现象的一个例子。

\medskip\noindent
设 $(\mathcal{C},\mathcal{O})$ 是环化位点，$\mathcal{A}$ 是
$(\mathcal{C},\mathcal{O})$ 上的分次代数层，$\mathcal{M}$ 是分次
$\mathcal{A}$-模，$k\in\mathbf{Z}$。定义 $\mathcal{M}$ 的
{\it 第 $k$ 次平移}，记作 $\mathcal{M}[k]$，为这样的分次
$\mathcal{A}$-模：其第 $n$ 个分次部分为
$$
(\mathcal{M}[k])^n = \mathcal{M}^{n + k}
$$
，即 $\mathcal{M}$ 的第 $n+k$ 个分次部分。它的乘法映射
$$
(\mathcal{M}[k])^n \times \mathcal{A}^m
\longrightarrow
(\mathcal{M}[k])^{n + m}
$$
直接取为 $\mathcal{M}$ 的乘法映射
$$
\mathcal{M}^{n + k} \times \mathcal{A}^m
\longrightarrow
\mathcal{M}^{n + m + k}
$$
。显然，这定义了函子 $[k]$，并且 $[k+l]=[k]\circ[l]$，还有
$$
\Hom_{\textit{Mod}^{gr}(\mathcal{A})}(\mathcal{L}, \mathcal{M}[k]) =
\Hom_{\textit{Mod}^{gr}(\mathcal{A})}(\mathcal{L}, \mathcal{M})[k]
$$
（不引入符号），并且关于 $\mathcal{M}$ 和 $\mathcal{L}$ 具有函子性。
因此，正如《微分分次代数》评注
\ref{dga-remark-graded-shift-functors} 所讨论的那样，分次
$\mathcal{A}$-模的分次范畴确实可以从分次 $\mathcal{A}$-模的通常范畴
和平移函子恢复出来。

\begin{lemma}
\label{sdga-lemma-gm-grothendieck-abelian}
设 $(\mathcal{C},\mathcal{O})$ 是环化位点，$\mathcal{A}$ 是分次
$\mathcal{O}$-代数。范畴 $\textit{Mod}(\mathcal{A})$
是格罗滕迪克阿贝尔范畴。
\end{lemma}

\begin{proof}
由引理 \ref{sdga-lemma-gm-abelian} 以及格罗滕迪克阿贝尔范畴的定义
（《单射对象》定义 \ref{injectives-definition-grothendieck-conditions}），
只须证明 $\textit{Mod}(\mathcal{A})$ 有生成元。我们断言
$$
\mathcal{G} = \bigoplus\nolimits_{k, U} j_{U!}\mathcal{A}_U[k]
$$
是一个生成元，其中直和遍历 $\mathcal{C}$ 的所有对象 $U$ 和
$k\in\mathbf{Z}$。事实上，给定分次 $\mathcal{A}$-模 $\mathcal{M}$，
若从 $\mathcal{G}$ 到 $\mathcal{M}$ 没有非零映射，则对所有
$k$ 和 $U$ 都有
$$
\Hom_{\textit{Mod}(\mathcal{A})}(j_{U!}\mathcal{A}_U[k], \mathcal{M}) =
\Hom_{\textit{Mod}(\mathcal{A}_U)}(\mathcal{A}_U[k], \mathcal{M}|_U) =
\Gamma(U, \mathcal{M}^{-k})
$$
等于零。因此 $\mathcal{M}$ 为零。
\end{proof}









\section{微分分次代数层}
\label{sdga-section-dga}

\noindent
本节是《微分分次代数》第 \ref{dga-section-dga} 节的类似物。

\begin{definition}
\label{sdga-definition-dga}
设 $(\mathcal{C},\mathcal{O})$ 是环化位点。$(\mathcal{C},\mathcal{O})$
上的{\it 微分分次 $\mathcal{O}$-代数层}或{\it 微分分次代数层}，
是一个 $\mathcal{O}$-模上链复形 $\mathcal{A}^\bullet$，并带有
$\mathcal{O}$-双线性映射
$$
\mathcal{A}^n \times \mathcal{A}^m \to \mathcal{A}^{n + m},\quad
(a, b) \longmapsto ab
$$
；这些映射称为乘法映射，并满足下列性质：
\begin{enumerate}
\item 乘法是结合的；
\item $\mathcal{A}^0$ 有一个全局截面 $1$，它是乘法的双边单位元；
\item 对 $U\in\Ob(\mathcal{C})$、$a\in\mathcal{A}^n(U)$ 和
$b\in\mathcal{A}^m(U)$，有
$$
\text{d}^{n + m}(ab) = \text{d}^n(a)b + (-1)^n a\text{d}^m(b)
$$
\end{enumerate}
我们常把这样的结构记作 $(\mathcal{A},\text{d})$。从
$(\mathcal{A},\text{d})$ 到 $(\mathcal{B},\text{d})$ 的
{\it 微分分次 $\mathcal{O}$-代数同态}，是与乘法映射相容的
$\mathcal{O}$-模复形映射
$f:\mathcal{A}^\bullet\to\mathcal{B}^\bullet$。
\end{definition}

\noindent
给定微分分次 $\mathcal{O}$-代数 $(\mathcal{A},\text{d})$ 和对象
$U\in\Ob(\mathcal{C})$，我们使用记号
$$
\mathcal{A}(U) =
\Gamma(U, \mathcal{A}) =
\bigoplus\nolimits_{n \in \mathbf{Z}} \mathcal{A}^n(U)
$$
这是微分分次 $\mathcal{O}(U)$-代数。

\medskip\noindent
我们将尽可能把微分分次 $\mathcal{O}$-代数
$(\mathcal{A},\text{d})$ 视为一个分次 $\mathcal{O}$-代数
$\mathcal{A}$，连同满足定义中莱布尼茨法则的次数 $1$ 算子
$\text{d}:\mathcal{A}\to\mathcal{A}$（这里把 $\mathcal{A}$
视为分次 $\mathcal{O}$-模）。

\begin{remark}
\label{sdga-remark-functoriality-dga}
设 $(f,f^\sharp):(\Sh(\mathcal{C}),\mathcal{O}_\mathcal{C})
\to (\Sh(\mathcal{D}), \mathcal{O}_\mathcal{D})$
是环化拓扑斯的态射。
\begin{enumerate}
\item 设 $(\mathcal{A},\text{d})$ 是微分分次
$\mathcal{O}_\mathcal{C}$-代数。其直像是微分分次
$\mathcal{O}_\mathcal{D}$-代数 $(f_*\mathcal{A},\text{d})$，
其中 $f_*\mathcal{A}$ 如评注 \ref{sdga-remark-functoriality-ga} 所述，
而映射 $f_*\mathcal{A}^n\to f_*\mathcal{A}^{n+1}$ 上的
$\text{d}=f_*\text{d}$。略去莱布尼茨法则成立的验证。
\item 设 $\mathcal{B}$ 是微分分次 $\mathcal{O}_\mathcal{D}$-代数。
其拉回是微分分次 $\mathcal{O}_\mathcal{C}$-代数
$(f^*\mathcal{B},\text{d})$，其中 $f^*\mathcal{B}$ 如评注
\ref{sdga-remark-functoriality-ga} 所述，而映射
$f^*\mathcal{B}^n\to f^*\mathcal{B}^{n+1}$ 上的
$\text{d}=f^*\text{d}$。略去莱布尼茨法则成立的验证。
\item 微分分次 $\mathcal{O}_\mathcal{C}$-代数同态
$f^*\mathcal{B}\to\mathcal{A}$ 的集合，与微分分次
$\mathcal{O}_\mathcal{D}$-代数同态
$\mathcal{B}\to f_*\mathcal{A}$ 的集合一一对应。
\end{enumerate}
第 (3) 条立即由模层上的 $f^*$ 与 $f_*$ 之间通常的伴随得到。
\end{remark}
















\section{微分分次模层}
\label{sdga-section-modules}

\noindent
本节是《微分分次代数》第 \ref{dga-section-modules} 节的类似物。

\begin{definition}
\label{sdga-definition-dgm}
设 $(\mathcal{C},\mathcal{O})$ 是环化位点，
$(\mathcal{A},\text{d})$ 是 $(\mathcal{C},\mathcal{O})$ 上的
微分分次代数层。一个（右）{\it 微分分次 $\mathcal{A}$-模}，
或 $\mathcal{A}$ 上的（右）{\it 微分分次模}，是一个上链复形
$\mathcal{M}^\bullet$，并带有 $\mathcal{O}$-双线性映射
$$
\mathcal{M}^n \times \mathcal{A}^m \to \mathcal{M}^{n + m},\quad
(x, a) \longmapsto xa
$$
；这些映射称为乘法映射，并满足下列性质：
\begin{enumerate}
\item 乘法满足 $(xa)a'=x(aa')$；
\item 对每个 $n$，$\mathcal{A}^0$ 的单位截面 $1$ 在
$\mathcal{M}^n$ 上作恒等作用；
\item 对 $U\in\Ob(\mathcal{C})$、$x\in\mathcal{M}^n(U)$ 和
$a\in\mathcal{A}^m(U)$，有
$$
\text{d}^{n + m}(xa) = \text{d}^n(x)a + (-1)^n x\text{d}^m(a)
$$
\end{enumerate}
我们常说“设 $\mathcal{M}$ 是微分分次 $\mathcal{A}$-模”来表示上述情形。
从 $\mathcal{M}$ 到 $\mathcal{N}$ 的
{\it 微分分次 $\mathcal{A}$-模同态}，是与乘法映射相容的
$\mathcal{O}$-模复形映射
$f:\mathcal{M}^\bullet\to\mathcal{N}^\bullet$。
（右）微分分次 $\mathcal{A}$-模构成的范畴记作
$\textit{Mod}(\mathcal{A},\text{d})$。
\end{definition}

\noindent
可以完全类似地定义{\it 左微分分次模}，但本章默认所说的模都是右模。

\medskip\noindent
给定微分分次 $\mathcal{A}$-模 $\mathcal{M}$ 和对象
$U\in\Ob(\mathcal{C})$，我们使用记号
$$
\mathcal{M}(U) =
\Gamma(U, \mathcal{M}) =
\bigoplus\nolimits_{n \in \mathbf{Z}} \mathcal{M}^n(U)
$$
这是一个（右）微分分次 $\mathcal{A}(U)$-模。

\begin{lemma}
\label{sdga-lemma-dgm-abelian}
设 $(\mathcal{C},\mathcal{O})$ 是环化位点，
$(\mathcal{A},\text{d})$ 是微分分次 $\mathcal{O}$-代数。
范畴 $\textit{Mod}(\mathcal{A},\text{d})$ 是阿贝尔范畴，并具有下列性质：
\begin{enumerate}
\item $\textit{Mod}(\mathcal{A},\text{d})$ 有任意直和；
\item $\textit{Mod}(\mathcal{A},\text{d})$ 有任意余极限；
\item $\textit{Mod}(\mathcal{A},\text{d})$ 中的滤过余极限是正合的；
\item $\textit{Mod}(\mathcal{A},\text{d})$ 有任意积；
\item $\textit{Mod}(\mathcal{A},\text{d})$ 有任意极限。
\end{enumerate}
遗忘函子
$$
\textit{Mod}(\mathcal{A}, \text{d})
\longrightarrow
\textit{Mod}(\mathcal{A})
$$
把微分分次 $\mathcal{A}$-模映到其底层分次模；它与所有极限和余极限交换。
\end{lemma}

\begin{proof}
把引理陈述中的函子记作
$F : \textit{Mod}(\mathcal{A}, \text{d}) \to \textit{Mod}(\mathcal{A})$
。注意，范畴 $\textit{Mod}(\mathcal{A})$ 具有性质 (1)--(5)；
见引理 \ref{sdga-lemma-gm-abelian}。

\medskip\noindent
考虑分次 $\mathcal{A}$-模同态 $f:\mathcal{M}\to\mathcal{N}$。
把 $f$ 看成分次 $\mathcal{A}$-模映射时，其核与余核还带有定义
\ref{sdga-definition-dgm} 中的微分。因此，它们也是
$\textit{Mod}(\mathcal{A},\text{d})$ 中的核与余核。由此
$\textit{Mod}(\mathcal{A},\text{d})$ 是阿贝尔范畴，并且取核与余核和
$F$ 交换。

\medskip\noindent
为了证明极限和余极限存在，只须证明积和直和存在；参见《范畴》引理
\ref{categories-lemma-limits-products-equalizers} 和
\ref{categories-lemma-colimits-coproducts-coequalizers}。
同样的引理表明，只要 $F$ 与直和及积交换，就可推出它与极限及余极限交换。

\medskip\noindent
设 $\mathcal{M}_t$（$t\in T$）是一族微分分次 $\mathcal{A}$-模。
把直和 $\bigoplus\mathcal{M}_t$ 看成分次 $\mathcal{A}$-模。
由于分次模的直和是逐项作出的，显然
$\bigoplus\mathcal{M}_t$ 带有微分。读者容易验证，这是
$\textit{Mod}(\mathcal{A},\text{d})$ 中的直和。积的情形类似。

\medskip\noindent
注意 $F$ 是正合函子，并且 $\textit{Mod}(\mathcal{A},\text{d})$
中的复形 $\mathcal{M}_1\to\mathcal{M}_2\to\mathcal{M}_3$ 正合，当且仅当
$F(\mathcal{M}_1) \to F(\mathcal{M}_2) \to F(\mathcal{M}_3)$
在 $\textit{Mod}(\mathcal{A})$ 中正合。由于
$\textit{Mod}(\mathcal{A})$ 中的滤过余极限正合，由此得到 (3)。
\end{proof}

\noindent
结合引理 \ref{sdga-lemma-dgm-abelian} 和 \ref{sdga-lemma-gm-abelian}，
得到阿贝尔范畴之间的正合忠实函子
$$
\textit{Mod}(\mathcal{A}, \text{d}) \longrightarrow \text{Comp}(\mathcal{O})
$$
。对微分分次 $\mathcal{A}$-模 $\mathcal{M}$，定义其上同调
$\mathcal{O}$-模 $H^i(\mathcal{M})$ 为与 $\mathcal{M}$ 对应的
$\mathcal{O}$-模复形的上同调。因此，微分分次 $\mathcal{A}$-模的短正合列
$0 \to \mathcal{K} \to \mathcal{L} \to \mathcal{M} \to 0$
给出上同调模的长正合列
\begin{equation}
\label{sdga-equation-les}
H^n(\mathcal{K}) \to H^n(\mathcal{L}) \to H^n(\mathcal{M}) \to
H^{n + 1}(\mathcal{K})
\end{equation}
；见《同调代数》引理 \ref{homology-lemma-long-exact-sequence-cochain}。

\medskip\noindent
此外，从现在起沿用模复形的全部术语。例如，若对所有
$k\in\mathbf{Z}$ 都有 $H^k(\mathcal{M})=0$，则称微分分次
$\mathcal{A}$-模 $\mathcal{M}$ 是{\it 无环的}。若微分分次
$\mathcal{A}$-模同态 $\mathcal{M}\to\mathcal{N}$ 对所有
$k\in\mathbf{Z}$ 都诱导同构
$H^k(\mathcal{M})\to H^k(\mathcal{N})$，则称它为{\it 拟同构}。
其余依此类推。


















\section{模的微分分次范畴}
\label{sdga-section-dgm-dg-cat}

\noindent
本节是《微分分次代数》例 \ref{dga-example-dgm-dg-cat} 的类似物。
关于微分分次范畴的约定，请参见《微分分次代数》第
\ref{dga-section-dga-categories} 节。

\medskip\noindent
设 $(\mathcal{C},\mathcal{O})$ 是环化位点，
$(\mathcal{A},\text{d})$ 是 $(\mathcal{C},\mathcal{O})$ 上的
微分分次代数层。我们将构造微分分次范畴
$$
\textit{Mod}^{dg}(\mathcal{A}, \text{d})
$$
，它在 $R=\Gamma(\mathcal{C},\mathcal{O})$ 上，且其对应的复形范畴
是微分分次 $\mathcal{A}$-模范畴：
$$
\textit{Mod}(\mathcal{A}, \text{d}) =
\text{Comp}(\textit{Mod}^{dg}(\mathcal{A}, \text{d}))
$$
。取右微分分次 $\mathcal{A}$-模作为
$\textit{Mod}^{dg}(\mathcal{A},\text{d})$ 的对象；见第
\ref{sdga-section-modules} 节。给定微分分次 $\mathcal{A}$-模
$\mathcal{L}$ 和 $\mathcal{M}$，置
$$
\Hom_{\textit{Mod}^{dg}(\mathcal{A}, \text{d})}(\mathcal{L}, \mathcal{M}) =
\Hom_{\textit{Mod}^{gr}(\mathcal{A})}(\mathcal{L}, \mathcal{M}) =
\bigoplus\nolimits_{n \in \mathbf{Z}} \Hom^n(\mathcal{L}, \mathcal{M})
$$
，把它作为分次 $R$-模；见第 \ref{sdga-section-gm-gr-cat} 节。
换言之，第 $n$ 个分次部分 $\Hom^n(\mathcal{L},\mathcal{M})$
是次数 $n$ 齐次的右 $\mathcal{A}$-模映射所成的 $R$-模。
对元素 $f\in\Hom^n(\mathcal{L},\mathcal{M})$，置
$$
\text{d}(f) =
\text{d}_\mathcal{M} \circ f - (-1)^n f \circ \text{d}_\mathcal{L}
$$
。为使此式有意义，把 $\text{d}_\mathcal{M}$ 和
$\text{d}_\mathcal{L}$ 看成分次 $\mathcal{O}$-模映射，并使用分次
$\mathcal{O}$-模映射的复合。显然，$\text{d}(f)$ 作为分次
$\mathcal{O}$-模映射是次数 $n+1$ 齐次的，并且关于
$\mathcal{A}$ 是线性的，因为对 $\mathcal{M}$ 和 $\mathcal{A}$
的齐次局部截面 $x$ 和 $a$，有
\begin{align*}
\text{d}(f)(xa)
& =
\text{d}_\mathcal{M}(f(x) a) - (-1)^n f (\text{d}_\mathcal{L}(xa)) \\
& =
\text{d}_\mathcal{M}(f(x)) a + (-1)^{\deg(x) + n} f(x) \text{d}(a) 
- (-1)^n f(\text{d}_\mathcal{L}(x)) a - (-1)^{n + \deg(x)} f(x) \text{d}(a) \\
& = \text{d}(f)(x) a
\end{align*}
，如所需（注意，若 $\Hom$ 上的微分定义中没有该符号，这一计算便不成立）。

\medskip\noindent
对微分分次 $\mathcal{A}$-模 $\mathcal{K}$、$\mathcal{L}$、
$\mathcal{M}$，第 \ref{sdga-section-gm-gr-cat} 节已经通过层映射的通常复合
定义了复合
$$
\Hom^m(\mathcal{L}, \mathcal{M}) \times
\Hom^n(\mathcal{K}, \mathcal{L}) \longrightarrow
\Hom^{n + m}(\mathcal{K}, \mathcal{M})
$$
。这定义微分分次模之间的映射
$$
\Hom_{\textit{Mod}^{dg}(\mathcal{A}, \text{d})}(\mathcal{L}, \mathcal{M})
\otimes_R
\Hom_{\textit{Mod}^{dg}(\mathcal{A}, \text{d})}(\mathcal{K}, \mathcal{L})
\longrightarrow
\Hom_{\textit{Mod}^{dg}(\mathcal{A}, \text{d})}(\mathcal{K}, \mathcal{M})
$$
，如《微分分次代数》定义 \ref{dga-definition-dga-category} 所要求，
因为
\begin{align*}
\text{d}(g \circ f) & =
\text{d}_\mathcal{M} \circ g \circ f
- (-1)^{n + m} g \circ f \circ \text{d}_\mathcal{K} \\
& =
\left(\text{d}_\mathcal{M} \circ g - (-1)^m g \circ \text{d}_L\right) \circ f
+ (-1)^m g \circ \left(\text{d}_\mathcal{L} \circ f
- (-1)^n f \circ \text{d}_\mathcal{K}\right) \\
& =
\text{d}(g) \circ f + (-1)^m g \circ \text{d}(f)
\end{align*}
，这里 $f$ 的次数为 $n$，$g$ 的次数为 $m$，正如所需。






\section{微分分次模层的张量积}
\label{sdga-section-tensor-product-dg}

\noindent
本节是《微分分次代数》第
\ref{dga-section-tensor-product} 节部分内容的类似物。

\medskip\noindent
设 $(\mathcal{C},\mathcal{O})$ 是环化位点，
$(\mathcal{A},\text{d})$ 是 $(\mathcal{C},\mathcal{O})$ 上的
微分分次代数层。设 $\mathcal{M}$ 是右微分分次 $\mathcal{A}$-模，
$\mathcal{N}$ 是左微分分次 $\mathcal{A}$-模。在这种情形下，
如下定义{\it 张量积} $\mathcal{M}\otimes_\mathcal{A}\mathcal{N}$。
作为分次 $\mathcal{O}$-模，它由第 \ref{sdga-section-tensor-product} 节的
构造给出，并带有微分
$$
\text{d}_{\mathcal{M} \otimes_\mathcal{A} \mathcal{N}} :
(\mathcal{M} \otimes_\mathcal{A} \mathcal{N})^n
\longrightarrow
(\mathcal{M} \otimes_\mathcal{A} \mathcal{N})^{n + 1}
$$
，其定义规则为
$$
\text{d}_{\mathcal{M} \otimes_\mathcal{A} \mathcal{N}}(x \otimes y) =
\text{d}_\mathcal{M}(x) \otimes y +
(-1)^{\deg(x)}x \otimes \text{d}_\mathcal{N}(y)
$$
，其中 $x$ 和 $y$ 分别是 $\mathcal{M}$ 和 $\mathcal{N}$ 的齐次局部截面。
为了看出其良定义性，必须证明
$\text{d}_{\mathcal{M}\otimes_\mathcal{A}\mathcal{N}}$ 把形如
$xa\otimes y-x\otimes ay$ 的元素映到零；这里 $x$、$a$、$y$
分别是 $\mathcal{M}$、$\mathcal{A}$、$\mathcal{N}$ 的齐次局部截面。
计算得
\begin{align*}
&
\text{d}_{\mathcal{M} \otimes_\mathcal{A} \mathcal{N}}(
xa \otimes y - x \otimes ay) \\
& =
\text{d}_\mathcal{M}(xa) \otimes y + (-1)^{\deg(x) + \deg(a)}
xa \otimes \text{d}_\mathcal{N}(y)
-\text{d}_\mathcal{M}(x) \otimes ay - (-1)^{\deg(x)}
x \otimes \text{d}_\mathcal{N}(ay) \\
& =
\text{d}_\mathcal{M}(x)a \otimes y + (-1)^{\deg(x)}x\text{d}(a) \otimes y
+ (-1)^{\deg(x) + \deg(a)} xa \otimes \text{d}_\mathcal{N}(y) \\
&
-\text{d}_\mathcal{M}(x) \otimes ay - (-1)^{\deg(x)}
x \otimes \text{d}(a)y - (-1)^{\deg(x) +
\deg(a)} x\otimes a\text{d}_\mathcal{N}(y)
\end{align*}
。注意下列元素
$$
\text{d}_\mathcal{M}(x)a \otimes y - \text{d}_\mathcal{M}(x) \otimes ay,\quad
x\text{d}(a) \otimes y - x \otimes \text{d}(a)y,\quad
\text{且}\quad
xa \otimes \text{d}_\mathcal{N}(y) - x\otimes a\text{d}_\mathcal{N}(y)
$$
在 $\mathcal{M}\otimes_\mathcal{A}\mathcal{N}$ 中都映到零，故得结论。
略去下式的验证：
$\text{d}_{\mathcal{M} \otimes_\mathcal{A} \mathcal{N}} \circ
\text{d}_{\mathcal{M} \otimes_\mathcal{A} \mathcal{N}} = 0$.

\medskip\noindent
固定左微分分次 $\mathcal{A}$-模 $\mathcal{N}$，便得到函子
$$
- \otimes_\mathcal{A} \mathcal{N} :
\textit{Mod}(\mathcal{A}, \text{d})
\longrightarrow
\text{Comp}(\mathcal{O})
$$
，其右端是 $\mathcal{O}$-模复形范畴。它可以提升为微分分次范畴之间的函子
$$
- \otimes_\mathcal{A} \mathcal{N} :
\textit{Mod}^{dg}(\mathcal{A}, \text{d})
\longrightarrow
\text{Comp}^{dg}(\mathcal{O})
$$
。在底层分次对象上，把次数 $n$ 的同态
$f:\mathcal{M}\to\mathcal{M}'$ 映到次数 $n$ 的映射
$f \otimes \text{id}_\mathcal{N} :
\mathcal{M} \otimes_\mathcal{A} \mathcal{N} \to
\mathcal{M}' \otimes_\mathcal{A} \mathcal{N}$，因为这正是第
\ref{sdga-section-tensor-product} 节中的做法。为证明该构造成立，
必须验证映射
$$
\Hom_{\textit{Mod}^{dg}(\mathcal{A}, \text{d})}(\mathcal{M}, \mathcal{M}')
\longrightarrow
\Hom_{\text{Comp}^{dg}(\mathcal{O})}(
\mathcal{M} \otimes_\mathcal{A} \mathcal{N},
\mathcal{M}' \otimes_\mathcal{A} \mathcal{N})
$$
与微分相容。对上述 $f$，须证明
$$
(\text{d}_{\mathcal{M}'} \circ f - (-1)^n f \circ \text{d}_\mathcal{M})
\otimes \text{id}_\mathcal{N}
$$
等于
$$
\text{d}_{\mathcal{M}' \otimes_\mathcal{A} \mathcal{N}}
\circ (f \otimes \text{id}_\mathcal{N})
- (-1)^n (f \otimes \text{id}_\mathcal{N}) \circ
\text{d}_{\mathcal{M} \otimes_\mathcal{A} \mathcal{N}}
$$
。计算这些算子在形如 $x\otimes y$ 的局部截面上的作用，其中
$x$ 和 $y$ 分别是 $\mathcal{M}$ 和 $\mathcal{N}$ 的齐次局部截面。
第一个算子给出
$$
(\text{d}_{\mathcal{M}'}(f(x)) - (-1)^n f(\text{d}_\mathcal{M}(x))) \otimes y
$$
，第二个算子给出
\begin{align*}
&\text{d}_{\mathcal{M}' \otimes_\mathcal{A} \mathcal{N}}(f(x) \otimes y)
- (-1)^n (f \otimes \text{id}_\mathcal{N})(
\text{d}_{\mathcal{M} \otimes_\mathcal{A} \mathcal{N}}(x \otimes y) \\
& =
\text{d}_{\mathcal{M}'}(f(x)) \otimes y +
(-1)^{\deg(x) + n}f(x) \otimes \text{d}_\mathcal{N}(y) \\
&
-(-1)^n f(\text{d}_\mathcal{M}(x)) \otimes y
-(-1)^n (-1)^{\deg(x)}f(x) \otimes \text{d}_\mathcal{N}(y)
\end{align*}
，这确实是同一个局部截面。










\section{微分分次模层的内 Hom}
\label{sdga-section-internal-hom-dg}

\noindent
我们需要第 \ref{sdga-section-dgm-dg-cat} 节中构造的层化版本。
设 $(\mathcal{C},\mathcal{O})$、$\mathcal{A}$、
$\mathcal{M}$、$\mathcal{L}$ 如第 \ref{sdga-section-dgm-dg-cat} 节所述。定义
$$
\SheafHom^{dg}_\mathcal{A}(\mathcal{M}, \mathcal{L}) =
\SheafHom^{gr}_\mathcal{A}(\mathcal{M}, \mathcal{L}) =
\bigoplus\nolimits_{n \in \mathbf{Z}}
\SheafHom^n_\mathcal{A}(\mathcal{M}, \mathcal{L})
$$
为分次 $\mathcal{O}$-模；见第
\ref{sdga-section-internal-hom-graded} 节。换言之，$U$ 上第 $n$ 个
分次部分 $\SheafHom^n_\mathcal{A}(\mathcal{L},\mathcal{M})$
的截面 $f$，是次数 $n$ 齐次的右 $\mathcal{A}_U$-模映射
$\mathcal{L}|_U\to\mathcal{M}|_U$。对这样的 $f$，置
$$
\text{d}(f) =
\text{d}_\mathcal{M}|_U \circ f - (-1)^n f \circ \text{d}_\mathcal{L}|_U
$$
。为使此式有意义，把 $\text{d}_\mathcal{M}|_U$ 和
$\text{d}_\mathcal{L}|_U$ 看成分次 $\mathcal{O}_U$-模映射，
并使用分次 $\mathcal{O}_U$-模映射的复合。显然，
$\text{d}(f)$ 作为分次 $\mathcal{O}_U$-模映射是次数 $n+1$
齐次的。使用与第 \ref{sdga-section-dgm-dg-cat} 节完全相同的计算，
可见 $\text{d}(f)$ 关于 $\mathcal{A}_U$ 是线性的。

\medskip\noindent
如同第 \ref{sdga-section-dgm-dg-cat} 节，有复合映射
$$
\SheafHom^{dg}_\mathcal{A}(\mathcal{L}, \mathcal{M}) \otimes_\mathcal{O}
\SheafHom^{dg}_\mathcal{A}(\mathcal{K}, \mathcal{L})
\longrightarrow
\SheafHom^{dg}_\mathcal{A}(\mathcal{K}, \mathcal{M})
$$
，其中左端是第 \ref{sdga-section-tensor-product-dg} 节定义的微分分次
$\mathcal{O}$-模张量积。该映射由复合映射
$$
\SheafHom^m(\mathcal{L}, \mathcal{M}) \otimes_\mathcal{O}
\SheafHom^n(\mathcal{K}, \mathcal{L}) \longrightarrow
\SheafHom^{n + m}(\mathcal{K}, \mathcal{M})
$$
给出；后者（局部地）按通常复合定义。在局部截面上使用与第
\ref{sdga-section-dgm-dg-cat} 节完全相同的计算，可见该复合映射是
微分分次 $\mathcal{O}$-模同态。

\medskip\noindent
按这些定义，有
$$
\Hom_{\textit{Mod}^{dg}(\mathcal{A})}(\mathcal{L}, \mathcal{M}) =
\Gamma(\mathcal{C}, \SheafHom^{dg}_\mathcal{A}(\mathcal{L}, \mathcal{M}))
$$
作为分次 $R$-模的等式，并且它与复合相容。









\section{微分分次双模层与张量--Hom 伴随}
\label{sdga-section-dg-bimodules}

\noindent
本节是《微分分次代数》第
\ref{dga-section-tensor-product} 节部分内容的类似物。

\begin{definition}
\label{sdga-definition-dg-bimodule}
设 $(\mathcal{C},\mathcal{O})$ 是环化位点，$\mathcal{A}$ 和
$\mathcal{B}$ 是 $(\mathcal{C},\mathcal{O})$ 上的微分分次代数层。
一个{\it 微分分次 $(\mathcal{A},\mathcal{B})$-双模}，由
$\mathcal{O}$-模复形 $\mathcal{M}^\bullet$ 以及
$\mathcal{O}$-双线性映射
$$
\mathcal{M}^n \times \mathcal{B}^m \to \mathcal{M}^{n + m},\quad
(x, b) \longmapsto xb
$$
和
$$
\mathcal{A}^n \times \mathcal{M}^m \to \mathcal{M}^{n + m},\quad
(a, x) \longmapsto ax
$$
给出；这些映射称为乘法映射，并满足下列性质：
\begin{enumerate}
\item 乘法满足 $a(a'x)=(aa')x$ 和 $(xb)b'=x(bb')$；
\item $(ax)b=a(xb)$；
\item $\text{d}(ax) = \text{d}(a) x + (-1)^{\deg(a)}a \text{d}(x)$ 和
$\text{d}(xb) = \text{d}(x) b + (-1)^{\deg(x)}x \text{d}(b)$；
\item $\mathcal{A}^0$ 的单位截面 $1$ 通过乘法作恒等作用；
\item $\mathcal{B}^0$ 的单位截面 $1$ 通过乘法作恒等作用。
\end{enumerate}
我们常把这样的结构记作 $\mathcal{M}$，有时写成
${}_\mathcal{A}\mathcal{M}_\mathcal{B}$。一个
{\it 微分分次 $(\mathcal{A},\mathcal{B})$-双模同态}
$f:\mathcal{M}\to\mathcal{N}$，是与乘法映射相容的
$\mathcal{O}$-模复形映射
$f:\mathcal{M}^\bullet\to\mathcal{N}^\bullet$。
\end{definition}

\noindent
给定微分分次 $(\mathcal{A},\mathcal{B})$-双模 $\mathcal{M}$
和对象 $U\in\Ob(\mathcal{C})$，我们使用记号
$$
\mathcal{M}(U) =
\Gamma(U, \mathcal{M}) =
\bigoplus\nolimits_{n \in \mathbf{Z}} \mathcal{M}^n(U)
$$
这是微分分次 $(\mathcal{A}(U),\mathcal{B}(U))$-双模。

\medskip\noindent
注意，微分分次 $(\mathcal{A},\mathcal{B})$-双模 $\mathcal{M}$
等价于一个右微分分次 $\mathcal{B}$-模，同时也是左微分分次
$\mathcal{A}$-模，且两者的分次和微分一致，并且
$\mathcal{A}$-模结构与 $\mathcal{B}$-模结构交换。精确陈述如下。

\begin{lemma}
\label{sdga-lemma-what-makes-a-bimodule-dg}
设 $(\mathcal{C},\mathcal{O})$ 是环化位点，$\mathcal{A}$ 和
$\mathcal{B}$ 是 $(\mathcal{C},\mathcal{O})$ 上的微分分次代数层，
$\mathcal{N}$ 是右微分分次 $\mathcal{B}$-模。$\mathcal{N}$ 上与给定
微分分次 $\mathcal{B}$-模结构相容的
$(\mathcal{A},\mathcal{B})$-双模结构，与微分分次
$\mathcal{O}$-代数同态
$$
\mathcal{A}
\longrightarrow
\SheafHom^{dg}_\mathcal{B}(\mathcal{N}, \mathcal{N})
$$
一一对应。
\end{lemma}

\begin{proof}
略。
\end{proof}

\noindent
设 $(\mathcal{C},\mathcal{O})$ 是环化位点，$\mathcal{A}$ 和
$\mathcal{B}$ 是 $(\mathcal{C},\mathcal{O})$ 上的微分分次代数层。
设 $\mathcal{M}$ 是右微分分次 $\mathcal{A}$-模，$\mathcal{N}$ 是
微分分次 $(\mathcal{A},\mathcal{B})$-双模。在这种情形下，第
\ref{sdga-section-tensor-product-dg} 节定义的微分分次张量积
$$
\mathcal{M} \otimes_\mathcal{A} \mathcal{N}
$$
是右微分分次 $\mathcal{B}$-模，其乘法映射如第
\ref{sdga-section-graded-bimodules} 节所述。该构造定义一个函子和一个
分次范畴之间的函子
$$
\otimes_\mathcal{A} \mathcal{N} :
\textit{Mod}(\mathcal{A}, \text{d})
\longrightarrow
\textit{Mod}(\mathcal{B}, \text{d})
\quad\text{且}\quad
\otimes_\mathcal{A} \mathcal{N} :
\textit{Mod}^{dg}(\mathcal{A}, \text{d})
\longrightarrow
\textit{Mod}^{dg}(\mathcal{B}, \text{d})
$$
：把从 $\mathcal{M}$ 到 $\mathcal{M}'$ 的次数 $n$ 同态映到
从 $\mathcal{M}\otimes_\mathcal{A}\mathcal{N}$ 到
$\mathcal{M}'\otimes_\mathcal{A}\mathcal{N}$ 的诱导次数 $n$ 映射。

\medskip\noindent
设 $(\mathcal{C},\mathcal{O})$ 是环化位点，$\mathcal{A}$ 和
$\mathcal{B}$ 是 $(\mathcal{C},\mathcal{O})$ 上的微分分次代数层。
设 $\mathcal{N}$ 是微分分次 $(\mathcal{A},\mathcal{B})$-双模，
$\mathcal{L}$ 是右微分分次 $\mathcal{B}$-模。在这种情形下，第
\ref{sdga-section-internal-hom-dg} 节定义的微分分次内 Hom
$$
\SheafHom_\mathcal{B}^{dg}(\mathcal{N}, \mathcal{L})
$$
是右微分分次 $\mathcal{A}$-模，其中右分次 $\mathcal{A}$-模结构
按第 \ref{sdga-section-graded-bimodules} 节定义。也可以用下列复合来定义乘法：
$$
\SheafHom_\mathcal{B}^{dg}(\mathcal{N}, \mathcal{L})
\otimes_\mathcal{O} \mathcal{A}
\to
\SheafHom_\mathcal{B}^{dg}(\mathcal{N}, \mathcal{L})
\otimes_\mathcal{O} \SheafHom_\mathcal{B}^{dg}(\mathcal{N}, \mathcal{N})
\to
\SheafHom_\mathcal{B}^{dg}(\mathcal{N}, \mathcal{L})
$$
其中第一个箭头来自引理 \ref{sdga-lemma-what-makes-a-bimodule-dg}，
第二个箭头是第 \ref{sdga-section-internal-hom-dg} 节的复合。
由于这两个箭头都与微分相容，所得对象确为微分分次
$\mathcal{A}$-模。该构造定义一个函子和一个微分分次范畴之间的函子
$$
\SheafHom_\mathcal{B}^{dg}(\mathcal{N}, -) :
\textit{Mod}(\mathcal{B}, \text{d})
\longrightarrow
\textit{Mod}(\mathcal{A})
\quad\text{且}\quad
\SheafHom_\mathcal{B}^{dg}(\mathcal{N}, -) :
\textit{Mod}^{dg}(\mathcal{B}, \text{d})
\longrightarrow
\textit{Mod}^{dg}(\mathcal{A}, \text{d})
$$
：把从 $\mathcal{L}$ 到 $\mathcal{L}'$ 的次数 $n$ 同态映到
从 $\SheafHom_\mathcal{B}^{dg}(\mathcal{N},\mathcal{L})$ 到
$\SheafHom_\mathcal{B}^{dg}(\mathcal{N},\mathcal{L}')$
的诱导次数 $n$ 映射。

\begin{lemma}
\label{sdga-lemma-tensor-hom-adjunction-dg}
设 $(\mathcal{C},\mathcal{O})$ 是环化位点，$\mathcal{A}$ 和
$\mathcal{B}$ 是 $(\mathcal{C},\mathcal{O})$ 上的微分分次代数层。
设 $\mathcal{M}$ 是右微分分次 $\mathcal{A}$-模，$\mathcal{N}$ 是
微分分次 $(\mathcal{A},\mathcal{B})$-双模，$\mathcal{L}$ 是右微分分次
$\mathcal{B}$-模。采用上述约定，有
$$
\Hom_{\textit{Mod}^{dg}(\mathcal{B}, \text{d})}(
\mathcal{M} \otimes_\mathcal{A} \mathcal{N}, \mathcal{L}) =
\Hom_{\textit{Mod}^{dg}(\mathcal{A}, \text{d})}(
\mathcal{M}, \SheafHom_\mathcal{B}^{dg}(\mathcal{N}, \mathcal{L}))
$$
以及
$$
\SheafHom_\mathcal{B}^{dg}(
\mathcal{M} \otimes_\mathcal{A} \mathcal{N}, \mathcal{L}) =
\SheafHom_\mathcal{A}^{dg}(
\mathcal{M}, \SheafHom_\mathcal{B}^{dg}(\mathcal{N}, \mathcal{L}))
$$
；这些等式关于 $\mathcal{M}$、$\mathcal{N}$、$\mathcal{L}$
都具有函子性。
\end{lemma}

\begin{proof}
略。提示：在分次层面，引理 \ref{sdga-lemma-tensor-hom-adjunction-gr}
已经证明了这一点。因此只须检验这些同构与微分相容；
这可通过局部截面层面的计算完成。
\end{proof}

\noindent
设 $(\mathcal{C},\mathcal{O})$ 是环化位点，$\mathcal{A}$ 和
$\mathcal{B}$ 是 $(\mathcal{C},\mathcal{O})$ 上的微分分次代数层。
作为上述情形的特例，设给定微分分次 $\mathcal{O}$-代数同态
$\varphi:\mathcal{A}\to\mathcal{B}$。于是得到一个函子和一个
微分分次范畴之间的函子
$$
\otimes_{\mathcal{A}, \varphi} \mathcal{B} :
\textit{Mod}(\mathcal{A}, \text{d})
\longrightarrow
\textit{Mod}(\mathcal{B}, \text{d})
\quad\text{且}\quad
\otimes_{\mathcal{A}, \varphi} \mathcal{B} :
\textit{Mod}^{dg}(\mathcal{A}, \text{d})
\longrightarrow
\textit{Mod}^{dg}(\mathcal{B}, \text{d})
$$
。另一方面，有限制标量函子
$$
res_\varphi :
\textit{Mod}(\mathcal{B}, \text{d})
\longrightarrow
\textit{Mod}(\mathcal{A}, \text{d})
\quad\text{且}\quad
res_\varphi :
\textit{Mod}^{dg}(\mathcal{B}, \text{d})
\longrightarrow
\textit{Mod}^{dg}(\mathcal{A}, \text{d})
$$
。由上一引理可以证明这些函子彼此伴随（如限制标量与基变换的通常情形）。
具体地，把 $\mathcal{B}$ 视为微分分次
$(\mathcal{A},\mathcal{B})$-双模时，记作
${}_\mathcal{A}\mathcal{B}_\mathcal{B}$。于是对任意右微分分次
$\mathcal{B}$-模 $\mathcal{L}$，有
$$
\SheafHom_\mathcal{B}^{dg}({}_\mathcal{A}\mathcal{B}_\mathcal{B}, \mathcal{L})
= res_\varphi(\mathcal{L})
$$
作为右微分分次 $\mathcal{A}$-模的等式。因此，引理
\ref{sdga-lemma-tensor-hom-adjunction-gr} 给出函子性同构
$$
\Hom_{\textit{Mod}^{dg}(\mathcal{B}, \text{d})}(
\mathcal{M} \otimes_{\mathcal{A}, \varphi} \mathcal{B}, \mathcal{L}) =
\Hom_{\textit{Mod}^{dg}(\mathcal{A}, \text{d})}(
\mathcal{M}, res_\varphi(\mathcal{L}))
$$
。若从上下文可知 $\varphi$，通常在此公式中省略对它的依赖。
以同样方式得到分次 $\mathcal{O}$-模的等式
$$
\SheafHom^{dg}_\mathcal{B}(
\mathcal{M} \otimes_\mathcal{A} \mathcal{B}, \mathcal{L}) =
\SheafHom_\mathcal{A}^{dg}(\mathcal{M}, \mathcal{L})
$$
。












\section{微分分次模层的拉回与直像}
\label{sdga-section-functoriality-dg}

\noindent
设 $(f,f^\sharp):(\Sh(\mathcal{C}),\mathcal{O}_\mathcal{C})
\to (\Sh(\mathcal{D}), \mathcal{O}_\mathcal{D})$
是环化拓扑斯的态射。设 $\mathcal{A}$ 是微分分次
$\mathcal{O}_\mathcal{C}$-代数，$\mathcal{B}$ 是微分分次
$\mathcal{O}_\mathcal{D}$-代数。设给定微分分次
$f^{-1}\mathcal{O}_\mathcal{D}$-代数同态
$$
\varphi : f^{-1}\mathcal{B} \to \mathcal{A}
$$
。由限制标量与扩张标量的伴随，这等价于微分分次
$\mathcal{O}_\mathcal{C}$-代数同态
$\varphi:f^*\mathcal{B}\to\mathcal{A}$；也等价地，可把
$\varphi$ 看作微分分次 $\mathcal{O}_\mathcal{D}$-代数同态
$$
\varphi : \mathcal{B} \to f_*\mathcal{A}
$$
。参见评注 \ref{sdga-remark-functoriality-dga}。

\medskip\noindent
先定义函子
$$
f_* :
\textit{Mod}(\mathcal{A}, \text{d})
\longrightarrow
\textit{Mod}(\mathcal{B}, \text{d})
$$
。给定微分分次 $\mathcal{A}$-模 $\mathcal{M}$，定义
$f_*\mathcal{M}$ 为第 \ref{sdga-section-functoriality-graded} 节构造的
分次 $\mathcal{B}$-模，其微分由映射
$f_*d:f_*\mathcal{M}^n\to f_*\mathcal{M}^{n+1}$ 给出。
该构造关于 $\mathcal{M}$ 显然具有函子性，从而得到所需函子。

\medskip\noindent
再定义函子
$$
f^* :
\textit{Mod}(\mathcal{B}, \text{d})
\longrightarrow
\textit{Mod}(\mathcal{A}, \text{d})
$$
。给定微分分次 $\mathcal{B}$-模 $\mathcal{N}$，定义
$f^*\mathcal{N}$ 为第 \ref{sdga-section-functoriality-graded} 节构造的
分次 $\mathcal{A}$-模。回忆
$$
f^*\mathcal{N} = f^{-1}\mathcal{N} \otimes_{f^{-1}\mathcal{B}} \mathcal{A}
$$
。由于 $f^{-1}\mathcal{N}$ 带有微分
$f^{-1}\text{d}:f^{-1}\mathcal{N}^n\to f^{-1}\mathcal{N}^{n+1}$，
可把该张量积视为第 \ref{sdga-section-dg-bimodules} 节所讨论张量积的例子，
因而得到微分。该构造关于 $\mathcal{N}$ 显然具有函子性，
从而得到函子 $f^*$。

\medskip\noindent
函子 $f_*$ 和 $f^*$ 很容易提升为微分分次范畴之间的函子
$$
f_* :
\textit{Mod}^{dg}(\mathcal{A}, \text{d})
\longrightarrow
\textit{Mod}^{dg}(\mathcal{B}, \text{d})
\quad\text{且}\quad
f^* :
\textit{Mod}^{dg}(\mathcal{B}, \text{d})
\longrightarrow
\textit{Mod}^{dg}(\mathcal{A}, \text{d})
$$
；它们在底层对象上的作用不变，而在次数 $n$ 的齐次态射上由上述构造的
函子性定义。

\begin{lemma}
\label{sdga-lemma-adjunction-push-pull-dg}
在上述情形下，有
$$
\Hom_{\textit{Mod}^{dg}(\mathcal{B}, \text{d})}(
\mathcal{N}, f_*\mathcal{M}) =
\Hom_{\textit{Mod}^{dg}(\mathcal{A}, \text{d})}(
f^*\mathcal{N}, \mathcal{M})
$$
\end{lemma}

\begin{proof}
略。提示：由引理 \ref{sdga-lemma-adjunction-push-pull-gr}，
该结论对底层分次范畴成立。一次计算表明这些同构与微分相容。
\end{proof}















\section{局部化与微分分次模层}
\label{sdga-section-localize-dg}

\noindent
设 $(\mathcal{C},\mathcal{O})$ 是环化位点，$U\in\Ob(\mathcal{C})$。
记
$$
j :
(\Sh(\mathcal{C}/U), \mathcal{O}_U)
\longrightarrow
(\Sh(\mathcal{C}), \mathcal{O})
$$
为相应的局部化态射（《位点上的模》第
\ref{sites-modules-section-localize} 节）。下面将使用如下事实：
对 $\mathcal{O}_U$-模 $\mathcal{M}_i$（$i=1,2$）和
$\mathcal{O}$-模 $\mathcal{A}$，有典范映射
$$
j_! :
\Hom_{\mathcal{O}_U}(
\mathcal{M}_1 \otimes_{\mathcal{O}_U} \mathcal{A}|_U, \mathcal{M}_2)
\longrightarrow
\Hom_\mathcal{O}(
j_!\mathcal{M}_1 \otimes_\mathcal{O} \mathcal{A}, j_!\mathcal{M}_2)
$$
。事实上，由《位点上的模》引理
\ref{sites-modules-lemma-j-shriek-and-tensor}，有
$j_!(\mathcal{M}_1 \otimes_{\mathcal{O}_U} \mathcal{A}|_U) =
j_!\mathcal{M}_1 \otimes_\mathcal{O} \mathcal{A}$。

\medskip\noindent
设 $\mathcal{A}$ 是微分分次 $\mathcal{O}$-代数。用
$\mathcal{A}_U$ 表示 $\mathcal{A}$ 在 $\mathcal{C}/U$ 上的限制；
换言之，$\mathcal{A}_U=j^*\mathcal{A}=j^{-1}\mathcal{A}$。
第 \ref{sdga-section-functoriality-dg} 节已经构造伴随函子
$$
j_* :
\textit{Mod}^{dg}(\mathcal{A}_U, \text{d})
\longrightarrow
\textit{Mod}^{dg}(\mathcal{A}, \text{d})
\quad\text{且}\quad
j^* :
\textit{Mod}^{dg}(\mathcal{A}, \text{d})
\longrightarrow
\textit{Mod}^{dg}(\mathcal{A}_U, \text{d})
$$
，其中 $j^*$ 是 $j_*$ 的左伴随。我们断言，此外还有正合函子
$$
j_! :
\textit{Mod}^{dg}(\mathcal{A}_U, \text{d})
\longrightarrow
\textit{Mod}^{dg}(\mathcal{A}, \text{d})
$$
，它是 $j_*$ 的右伴随。具体地，给定微分分次
$\mathcal{A}_U$-模 $\mathcal{M}$，定义 $j_!\mathcal{M}$ 为第
\ref{sdga-section-localize-graded} 节构造的分次 $\mathcal{A}$-模，
其微分为
$j_!\text{d} : j_!\mathcal{M}^n \to j_!\mathcal{M}^{n + 1}$.
给定微分分次 $\mathcal{A}_U$-模之间次数 $n$ 的齐次映射
$f:\mathcal{M}\to\mathcal{M}'$，得到微分分次 $\mathcal{A}$-模之间
次数 $n$ 的齐次映射
$j_!f:j_!\mathcal{M}\to j_!\mathcal{M}'$。略去该构造与微分相容的
直接验证。这样便得到所需函子。

\begin{lemma}
\label{sdga-lemma-extension-by-zero-dg}
在上述情形下，有
$$
\Hom_{\textit{Mod}^{dg}(\mathcal{A}, \text{d})}(
j_!\mathcal{M}, \mathcal{N}) =
\Hom_{\textit{Mod}^{dg}(\mathcal{A}_U, \text{d})}(
\mathcal{M}, j^*\mathcal{N})
$$
\end{lemma}

\begin{proof}
略。提示：引理 \ref{sdga-lemma-extension-by-zero-graded} 已经说明该结论在
分次层面成立。因此只须检验所得同构与微分相容。
\end{proof}

\begin{lemma}
\label{sdga-lemma-tensor-with-extension-by-zero-dg}
在上述情形下，设 $\mathcal{M}$ 是右微分分次 $\mathcal{A}_U$-模，
$\mathcal{N}$ 是左微分分次 $\mathcal{A}$-模。则
$$
j_!\mathcal{M} \otimes_\mathcal{A} \mathcal{N} =
j_!(\mathcal{M} \otimes_{\mathcal{A}_U} \mathcal{N}|_U)
$$
作为 $\mathcal{O}$-模复形的等式，并且关于 $\mathcal{M}$ 和
$\mathcal{N}$ 具有函子性。
\end{lemma}

\begin{proof}
作为分次模，该结论由引理
\ref{sdga-lemma-tensor-with-extension-by-zero} 得到。
略去此同构与微分相容的验证。
\end{proof}








\section{微分分次模层上的平移函子}
\label{sdga-section-shift-dg}

\noindent
设 $(\mathcal{C},\mathcal{O})$ 是环化位点，$\mathcal{A}$ 是
$(\mathcal{C},\mathcal{O})$ 上的微分分次代数层，$\mathcal{M}$
是微分分次 $\mathcal{A}$-模，$k\in\mathbf{Z}$。如下定义
$\mathcal{M}$ 的{\it 第 $k$ 次平移}，记作 $\mathcal{M}[k]$：
\begin{enumerate}
\item 作为分次 $\mathcal{A}$-模，$\mathcal{M}[k]$ 如第
\ref{sdga-section-shift} 节所定义；
\item 微分
$d_{\mathcal{M}[k]} : (\mathcal{M}[k])^n \to (\mathcal{M}[k])^{n + 1}$
定义为
$(-1)^k\text{d}_\mathcal{M} : \mathcal{M}^{n + k} \to \mathcal{M}^{n + k + 1}$.
\end{enumerate}
对次数 $n$ 齐次的 $\mathcal{A}$-模同态
$f:\mathcal{L}\to\mathcal{M}$，令
$f[k]:\mathcal{L}[k]\to\mathcal{M}[k]$ 由与 $f$ 相同的分量映射给出。
则 $f[k]$ 是次数 $n$ 齐次的 $\mathcal{A}$-模映射。这给出映射
$$
\Hom_{\textit{Mod}^{dg}(\mathcal{A}, \text{d})}(\mathcal{L}, \mathcal{M})
\longrightarrow
\Hom_{\textit{Mod}^{dg}(\mathcal{A}, \text{d})}
(\mathcal{L}[k], \mathcal{M}[k])
$$
，并与微分相容（这是因为 $\mathcal{L}$ 与 $\mathcal{M}$ 的微分符号
改变了相同的量）。这些选择与《微分分次代数》定义
\ref{dga-definition-shift} 中的选择相容。显然，这定义微分分次范畴之间的函子
$$
[k] :
\textit{Mod}^{dg}(\mathcal{A}, \text{d})
\longrightarrow
\textit{Mod}^{dg}(\mathcal{A}, \text{d})
$$
，并且 $[k+l]=[k]\circ[l]$。

\medskip\noindent
我们断言，在第 \ref{sdga-section-shift} 节对底层分次模定义的同构
$$
\Hom_{\textit{Mod}^{dg}(\mathcal{A}, \text{d})}(\mathcal{L}, \mathcal{M}[k])
=
\Hom_{\textit{Mod}^{dg}(\mathcal{A}, \text{d})}(\mathcal{L}, \mathcal{M})[k]
$$
与微分相容。为说明这一点，设
$f:\mathcal{L}\to\mathcal{M}[k]$ 是次数 $n$ 齐次的右
$\mathcal{A}$-模映射；它是左端的次数 $n$ 元素。用
$f':\mathcal{L}\to\mathcal{M}$ 表示次数 $n+k$ 齐次且与 $f$
具有{\bf 相同}分量映射的 $\mathcal{A}$-模映射。按我们的约定，
它就是右端对应的次数 $n$ 元素。按左端微分的定义，得到
$$
\text{d}_{LHS}(f) =
\text{d}_{\mathcal{M}[k]} \circ f - (-1)^n f \circ \text{d}_\mathcal{L} =
(-1)^k\text{d}_\mathcal{M} \circ f - (-1)^n f \circ \text{d}_\mathcal{L}
$$
；而对右端微分，得到
$$
\text{d}_{RHS}(f') =
(-1)^k\left(
\text{d}_\mathcal{M} \circ f' - (-1)^{n + k} f' \circ \text{d}_\mathcal{L}
\right) =
(-1)^k\text{d}_\mathcal{M} \circ f' - (-1)^n f' \circ \text{d}_\mathcal{L}
$$
。这些映射具有相同的分量映射，证明完成。










\section{同伦范畴}
\label{sdga-section-homotopy}

\noindent
本节对应《微分分次代数》第 \ref{dga-section-homotopy} 节。

\begin{definition}
\label{sdga-definition-homotopy}
设 $(\mathcal{C}, \mathcal{O})$ 是环化位点，
$\mathcal{A}$ 是 $(\mathcal{C}, \mathcal{O})$ 上的微分分次代数层。设
$f, g : \mathcal{M} \to \mathcal{N}$
是微分分次 $\mathcal{A}$-模同态。{\it $f$ 与 $g$ 之间的同伦}
是次数为 $-1$ 的齐次分次 $\mathcal{A}$-模映射
$h : \mathcal{M} \to \mathcal{N}$，且满足
$$
f - g = \text{d}_\mathcal{N} \circ h + h \circ \text{d}_\mathcal{M}
$$
若存在这样的同伦，则称 $f$ 与 $g$ {\it 同伦}。
\end{definition}

\noindent
在该定义的情形下，若有映射
$a : \mathcal{K} \to \mathcal{M}$ 和
$c : \mathcal{N} \to \mathcal{L}$，则
$$
\begin{matrix}
h\text{ 是一个同伦} \\
\text{，介于 }f\text{ 与 }g\text{ 之间}
\end{matrix}
\quad
\Rightarrow
\quad
\begin{matrix}
c \circ h \circ a\text{ 是}\\
c \circ f \circ a\text{ 与 }c\circ g \circ a\text{ 之间的同伦}
\end{matrix}
$$
因此可以定义 $\textit{Mod}(\mathcal{A}, \text{d})$
中态射同伦类的复合。

\begin{definition}
\label{sdga-definition-complexes-notation}
设 $(\mathcal{C}, \mathcal{O})$ 是环化位点，
$\mathcal{A}$ 是 $(\mathcal{C}, \mathcal{O})$ 上的微分分次代数层。
{\it 同伦范畴}，记作 $K(\textit{Mod}(\mathcal{A}, \text{d}))$，
是如下范畴：其对象为 $\textit{Mod}(\mathcal{A}, \text{d})$
的对象，其态射为微分分次 $\mathcal{A}$-模同态的同伦类。
\end{definition}

\noindent
记号 $K(\textit{Mod}(\mathcal{A}, \text{d}))$ 并非标准记号，
但至少与 Stacks 项目其他地方对 $K(-)$ 的用法一致。

\medskip\noindent
在《微分分次代数》定义
\ref{dga-definition-homotopy-category-of-dga-category} 中，
我们定义了微分分次范畴 $\mathcal{S}$ 的复形范畴
$\text{Comp}(\mathcal{S})$ 和同伦范畴 $K(\mathcal{S})$。
将此定义应用于微分分次范畴
$\textit{Mod}^{dg}(\mathcal{A}, \text{d})$，得到
$$
\textit{Mod}(\mathcal{A}, \text{d}) =
\text{Comp}(\textit{Mod}^{dg}(\mathcal{A}, \text{d}))
$$
（参见第 \ref{sdga-section-dgm-dg-cat} 节的讨论），并得到
$$
K(\textit{Mod}(\mathcal{A}, \text{d})) =
K(\textit{Mod}^{dg}(\mathcal{A}, \text{d}))
$$
要验证后一等式，注意当 $h$ 如定义 \ref{sdga-definition-homotopy} 中所述时，
有等式
$$
\text{d}_{\Hom_{\textit{Mod}^{dg}(\mathcal{A}, \text{d})}
(\mathcal{M}, \mathcal{N})}(h) =
\text{d}_\mathcal{N} \circ h + h \circ \text{d}_\mathcal{M}
$$
细节从略。

\begin{lemma}
\label{sdga-lemma-homotopy-direct-sums}
设 $(\mathcal{C}, \mathcal{O})$ 是环化位点，
$\mathcal{A}$ 是 $(\mathcal{C}, \mathcal{O})$ 上的微分分次代数层。
同伦范畴 $K(\textit{Mod}(\mathcal{A}, \text{d}))$ 有直和与直积。
\end{lemma}

\begin{proof}
从略。提示：只需使用引理 \ref{sdga-lemma-dgm-abelian} 中的直和与直积。
这是可行的，因为我们已经看到这些函子与到分次 $\mathcal{A}$-模范畴
的遗忘函子交换，而且在阿贝尔群族的范畴上，
$\prod$ 与 $\bigoplus$ 都是正合函子。
\end{proof}










\section{锥与三角}
\label{sdga-section-conclude-triangulated}

\noindent
本节利用《微分分次代数》第 \ref{dga-section-review} 节的内容，
证明微分分次 $\mathcal{A}$-模范畴的同伦范畴是三角范畴。

\begin{lemma}
\label{sdga-lemma-axioms-AB}
设 $(\mathcal{C}, \mathcal{O})$ 是环化位点，
$\mathcal{A}$ 是 $(\mathcal{C}, \mathcal{O})$ 上的微分分次代数层。
微分分次范畴 $\textit{Mod}^{dg}(\mathcal{A}, \text{d})$
满足《微分分次代数》第 \ref{dga-section-review} 节的公理 (A) 与 (B)。
\end{lemma}

\begin{proof}
设给定微分分次 $\mathcal{A}$-模 $\mathcal{M}$ 与 $\mathcal{N}$。
以显然的方式定义微分分次 $\mathcal{A}$-模
$\mathcal{M} \oplus \mathcal{N}$。此时余投射
$i : \mathcal{M} \to \mathcal{M} \oplus \mathcal{N}$、
$j : \mathcal{N} \to \mathcal{M} \oplus \mathcal{N}$ 以及投射
$p : \mathcal{M} \oplus \mathcal{N} \to \mathcal{N}$、
$q : \mathcal{M} \oplus \mathcal{N} \to \mathcal{M}$
都是微分分次 $\mathcal{A}$-模态射。因此 $i,j,p,q$
均为次数 $0$ 的齐次闭态射，即 $\text{d}(i)=0$，等等。
故此直和是《微分分次代数》定义
\ref{dga-definition-dg-direct-sum} 意义下的微分分次和。
这证明了公理 (A)。

\medskip\noindent
公理 (B) 已在第 \ref{sdga-section-shift-dg} 节中证明。
\end{proof}

\noindent
设 $(\mathcal{C}, \mathcal{O})$ 是环化位点，
$\mathcal{A}$ 是 $(\mathcal{C}, \mathcal{O})$ 上的微分分次代数层。
回忆，$\textit{Mod}(\mathcal{A}, \text{d})$ 中的序列
$$
0 \to \mathcal{K} \to \mathcal{L} \to \mathcal{N} \to 0
$$
若在 $\textit{Mod}(\mathcal{A})$ 中分裂，换言之，若作为分次
$\mathcal{A}$-模序列分裂，则称为{\it 容许短正合序列}
（《微分分次代数》第 \ref{dga-section-review} 节）。
以 $s : \mathcal{N} \to \mathcal{L}$ 和
$\pi : \mathcal{L} \to \mathcal{K}$ 表示分次
$\mathcal{A}$-模意义下的分裂映射。结合引理 \ref{sdga-lemma-axioms-AB}
与《微分分次代数》引理 \ref{dga-lemma-get-triangle}，得到三角
$$
\mathcal{K} \to \mathcal{L} \to \mathcal{N} \to \mathcal{K}[1]
$$
其中《微分分次代数》引理 \ref{dga-lemma-get-triangle} 的证明中
的箭头 $\mathcal{N} \to \mathcal{K}[1]$ 构造为
$$
\delta = \pi \circ
\text{d}_{\Hom_{\textit{Mod}^{dg}(\mathcal{A}, \text{d})}
(\mathcal{L}, \mathcal{M})}(s) =
\pi \circ \text{d}_\mathcal{L} \circ s -
\pi \circ s \circ \text{d}_\mathcal{N} =
\pi \circ \text{d}_\mathcal{L} \circ s
$$
请原谅这个可怕的记号。无论如何，在我们的情形下，
《微分分次代数》引理 \ref{dga-lemma-get-triangle} 所构造的边界映射
$\delta$，在底层 $\mathcal{O}$-模复形上与 Stacks 项目各处用于
逐项分裂的复形短正合序列的通常边界映射一致；参见《导出范畴》
定义 \ref{derived-definition-split-ses}。

\begin{definition}
\label{sdga-definition-cone}
设 $(\mathcal{C}, \mathcal{O})$ 是环化位点，
$\mathcal{A}$ 是 $(\mathcal{C}, \mathcal{O})$ 上的微分分次代数层。
设 $f : \mathcal{K} \to \mathcal{L}$ 是微分分次
$\mathcal{A}$-模同态。$f$ 的{\it 锥}是按如下方式定义的微分分次
$\mathcal{A}$-模 $C(f)$：
\begin{enumerate}
\item 其底层 $\mathcal{O}$-模复形是相应的 $\mathcal{A}$-模复形映射
$f : \mathcal{K}^\bullet \to \mathcal{L}^\bullet$ 的锥，即
$C(f)^n = \mathcal{L}^n \oplus \mathcal{K}^{n + 1}$，微分为
$$
d_{C(f)} =
\left(
\begin{matrix}
\text{d}_\mathcal{L} & f \\
0 & -\text{d}_\mathcal{K}
\end{matrix}
\right)
$$
\item 乘法映射
$$
C(f)^n \times \mathcal{A}^m \to C(f)^{n + m}
$$
是乘法映射
$\mathcal{L}^n \times \mathcal{A}^m \to \mathcal{L}^{n + m}$ 与
$\mathcal{K}^{n + 1} \times \mathcal{A}^m \to \mathcal{K}^{n + 1 + m}$.
的直和。
\end{enumerate}
显然的映射诱导出典范的微分分次 $\mathcal{A}$-模同态
$i : \mathcal{L} \to C(f)$ 与 $p : C(f) \to \mathcal{K}[1]$。
\end{definition}

\noindent
注意，在该定义的情形下，序列
$$
0 \to \mathcal{L} \to C(f) \to \mathcal{K}[1] \to 0
$$
是容许短正合序列。

\begin{lemma}
\label{sdga-lemma-axiom-C}
设 $(\mathcal{C}, \mathcal{O})$ 是环化位点，
$\mathcal{A}$ 是 $(\mathcal{C}, \mathcal{O})$ 上的微分分次代数层。
微分分次范畴 $\textit{Mod}^{dg}(\mathcal{A}, \text{d})$
满足《微分分次代数》情形 \ref{dga-situation-ABC} 中表述的公理 (C)。
\end{lemma}

\begin{proof}
设 $f : \mathcal{K} \to \mathcal{L}$ 是微分分次
$\mathcal{A}$-模同态。由上文有容许短正合序列
$$
0 \to \mathcal{L} \to C(f) \to \mathcal{K}[1] \to 0
$$
为完成证明，只需说明与之相伴的边界映射
$$
\delta : \mathcal{K}[1] \to \mathcal{L}[1]
$$
（参见上文的讨论）等于 $f[1]$。对截面
$s : \mathcal{K}[1] \to C(f)$，在次数 $n$ 上取嵌入
$\mathcal{K}^{n + 1} \to C(f)^n$。映射 $\pi$ 在次数 $n$
上由投射 $C(f)^n \to \mathcal{L}^n$ 给出。因此最后需要计算
$$
\delta = \pi \circ \text{d}_{C(f)} \circ s
$$
（参见上文的讨论）。用矩阵记号，此式等于
$$
\left(
\begin{matrix}
1 & 0
\end{matrix}
\right)
\left(
\begin{matrix}
\text{d}_\mathcal{L} & f \\
0 & -\text{d}_\mathcal{K}
\end{matrix}
\right)
\left(
\begin{matrix}
0 \\
1
\end{matrix}
\right) = f,
$$
正如所需。
\end{proof}

\noindent
至此可知，《微分分次代数》第 \ref{dga-section-review} 节中
证明的所有引理对微分分次范畴
$\textit{Mod}^{dg}(\mathcal{A}, \text{d})$ 均成立。
特别地，有如下命题。

\begin{proposition}
\label{sdga-proposition-homotopy-category-triangulated}
设 $(\mathcal{C}, \mathcal{O})$ 是环化位点，
$\mathcal{A}$ 是 $(\mathcal{C}, \mathcal{O})$ 上的微分分次代数层。
同伦范畴 $K(\textit{Mod}(\mathcal{A}, \text{d}))$ 是三角范畴，其中
\begin{enumerate}
\item 平移函子是第 \ref{sdga-section-shift-dg} 节中构造的函子；
\item 杰出三角是 $K(\textit{Mod}(\mathcal{A}, \text{d}))$
中作为三角同构于形如
$$
\mathcal{K} \to \mathcal{L} \to \mathcal{N}
\xrightarrow{\delta} \mathcal{K}[1],\quad\quad
\delta = \pi \circ \text{d}_\mathcal{L} \circ s
$$
的三角；该三角由上文 $\textit{Mod}(\mathcal{A}, \text{d})$
中的容许短正合序列
$0 \to \mathcal{K} \to \mathcal{L} \to \mathcal{N} \to 0$ 构造。
\end{enumerate}
\end{proposition}

\begin{proof}
回忆 $K(\textit{Mod}(\mathcal{A}, \text{d})) =
K(\textit{Mod}^{dg}(\mathcal{A}, \text{d}))$，参见第
\ref{sdga-section-homotopy} 节。于是命题由引理 \ref{sdga-lemma-axioms-AB}、
\ref{sdga-lemma-axiom-C} 以及《微分分次代数》命题
\ref{dga-proposition-ABC-homotopy-category-triangulated} 得出。
\end{proof}

\begin{remark}
\label{sdga-remark-cone-identity}
设 $(\mathcal{C}, \mathcal{O})$ 是环化位点，
$\mathcal{A}$ 是 $(\mathcal{C}, \mathcal{O})$ 上的微分分次代数层。
把恒等映射 $\mathcal{A}\to\mathcal{A}$ 视为微分分次
$\mathcal{A}$-模映射，并令 $C=C(\text{id}_\mathcal{A})$ 为其锥。则
$$
\Hom_{\textit{Mod}(\mathcal{A}, \text{d})}(C, \mathcal{M}) =
\{(x, y) \in
\Gamma(\mathcal{C}, \mathcal{M}^0) \times
\Gamma(\mathcal{C}, \mathcal{M}^{-1}) \mid
x = \text{d}(y)\}
$$
其中从左到右的映射把 $f$ 送到偶对 $(x,y)$：
$x$ 是 $C^{-1}=\mathcal{A}^{-1}\oplus\mathcal{A}^0$
的整体截面 $(0,1)$ 的像，而 $y$ 是
$C^0=\mathcal{A}^0\oplus\mathcal{A}^1$
的整体截面 $(1,0)$ 的像。
\end{remark}

\begin{lemma}
\label{sdga-lemma-dgm-grothendieck-abelian}
设 $(\mathcal{C}, \mathcal{O})$ 是环化位点，
$(\mathcal{A}, \text{d})$ 是微分分次 $\mathcal{O}$-代数。
范畴 $\textit{Mod}(\mathcal{A}, \text{d})$ 是格罗滕迪克阿贝尔范畴。
\end{lemma}

\begin{proof}
由引理 \ref{sdga-lemma-dgm-abelian} 以及格罗滕迪克阿贝尔范畴的定义
（《内射对象》定义 \ref{injectives-definition-grothendieck-conditions}），
只需证明 $\textit{Mod}(\mathcal{A}, \text{d})$ 有生成元。
对 $\mathcal{C}$ 的每个对象 $U$，以 $C_U$ 表示如评注
\ref{sdga-remark-cone-identity} 中恒等映射
$\mathcal{A}_U\to\mathcal{A}_U$ 的锥。我们断言
$$
\mathcal{G} = \bigoplus\nolimits_{k, U} j_{U!}C_U[k]
$$
是生成元，其中直和遍历 $\mathcal{C}$ 的所有对象 $U$
与所有 $k\in\mathbf{Z}$。事实上，给定微分分次
$\mathcal{A}$-模 $\mathcal{M}$，若从 $\mathcal{G}$
到 $\mathcal{M}$ 没有非零映射，则对所有 $k$ 与 $U$ 有
\begin{align*}
\Hom_{\textit{Mod}(\mathcal{A})}(j_{U!}C_U[k], \mathcal{M}) \\
& =
\Hom_{\textit{Mod}(\mathcal{A}_U)}(C_U[k], \mathcal{M}|_U) \\
& =
\{(x, y) \in \mathcal{M}^{-k}(U) \times \mathcal{M}^{-k - 1}(U) \mid
x = \text{d}(y)\}
\end{align*}
等于零。因此 $\mathcal{M}$ 为零。
\end{proof}











\section{平坦消解}
\label{sdga-section-P-resolutions}

\noindent
本节对应《微分分次代数》第 \ref{dga-section-P-resolutions} 节。

\medskip\noindent
设 $(\mathcal{C}, \mathcal{O})$ 是环化位点，
$\mathcal{A}$ 是 $(\mathcal{C}, \mathcal{O})$ 上的微分分次代数层。
若右微分分次 $\mathcal{A}$-模 $\mathcal{P}$ 满足
\begin{enumerate}
\item 函子
$\mathcal{N} \mapsto \mathcal{P} \otimes_\mathcal{A} \mathcal{N}$
在分次左 $\mathcal{A}$-模范畴上正合；
\item 若 $\mathcal{N}$ 是无环微分分次左 $\mathcal{A}$-模，则
$\mathcal{P} \otimes_\mathcal{A} \mathcal{N}$ 无环；
\item 对任意环化拓扑斯态射
$(f, f^\sharp) : (\Sh(\mathcal{C}'), \mathcal{O}')
\to (\Sh(\mathcal{C}), \mathcal{O})$、任意微分分次
$\mathcal{O}'$-代数 $\mathcal{A}'$，以及任意微分分次
$f^{-1}\mathcal{O}_\mathcal{D}$-代数映射
$\varphi : f^{-1}\mathcal{A} \to \mathcal{A}'$，
则把拉回 $f^*\mathcal{P}$（第 \ref{sdga-section-functoriality-dg} 节）
视为微分分次 $\mathcal{A}'$-模时，它满足性质 (1) 与 (2)。
\end{enumerate}
则称 $\mathcal{P}$ {\it 良好}。第一个条件表示 $\mathcal{P}$
作为右分次 $\mathcal{A}$-模是平坦的；第二个条件表示
$\mathcal{P}$ 在 Spaltenstein 的意义下是 K-平坦的
（参见《位点上的上同调》第
\ref{sites-cohomology-section-flat} 节）；第三个条件表示
任意基变换后这些性质仍成立。

\medskip\noindent
也许令人意外的是，良好模有很多。

\begin{lemma}
\label{sdga-lemma-supply-good}
设 $(\mathcal{C}, \mathcal{O})$ 是环化位点，
$\mathcal{A}$ 是 $(\mathcal{C}, \mathcal{O})$ 上的微分分次代数层，
$U \in \Ob(\mathcal{C})$。则 $j_!\mathcal{A}_U$
是良好微分分次 $\mathcal{A}$-模。
\end{lemma}

\begin{proof}
设 $\mathcal{N}$ 是分次左 $\mathcal{A}$-模。
由引理 \ref{sdga-lemma-tensor-with-extension-by-zero}，作为分次模有
$$
j_!\mathcal{A}_U \otimes_\mathcal{A} \mathcal{N} =
j_!(\mathcal{A}_U \otimes_{\mathcal{A}_U} \mathcal{N}|_U) =
j_!(\mathcal{N}_U)
$$
由于限制到 $U$ 与 $j_!$ 都正合，这证明了条件 (1)。
利用引理 \ref{sdga-lemma-tensor-with-extension-by-zero-dg}，
同一论证也证明条件 (2)。

\medskip\noindent
考虑环化拓扑斯态射
$(f, f^\sharp) : (\Sh(\mathcal{C}'), \mathcal{O}')
\to (\Sh(\mathcal{C}), \mathcal{O})$、微分分次
$\mathcal{O}'$-代数 $\mathcal{A}'$，以及微分分次
$f^{-1}\mathcal{O}$-代数映射
$\varphi : f^{-1}\mathcal{A} \to \mathcal{A}'$。
需要证明
$$
f^*j_!\mathcal{A}_U =
f^{-1}j_!\mathcal{A}_U \otimes_{f^{-1}\mathcal{A}} \mathcal{A}'
$$
对带有微分分次 $\mathcal{O}'$-代数层 $\mathcal{A}'$ 的环化拓扑斯
$(\Sh(\mathcal{C}'), \mathcal{O}')$ 满足 (1) 与 (2)。
为证明这一点，可以把 $(\Sh(\mathcal{C}), \mathcal{O})$
和 $(\Sh(\mathcal{C}'), \mathcal{O}')$ 替换为等价的环化拓扑斯。
因此由《位点上的模》引理
\ref{sites-modules-lemma-morphism-ringed-topoi-comes-from-morphism-ringed-sites}
可设 $f$ 来自连续函子 $u:\mathcal{C}\to\mathcal{C}'$
给出的位点态射 $f:\mathcal{C}\to\mathcal{C}'$。
此时置 $U'=u(U)$，并以
$j' : \Sh(\mathcal{C}'/U') \to \Sh(\mathcal{C}')$
表示相应的局部化态射。得到环化拓扑斯态射的交换方块
$$
\xymatrix{
(\Sh(\mathcal{C}'/U'), \mathcal{O}'_{U'})
\ar[rr]_{(j', (j')^\sharp)} \ar[d]_{(f', (f')^\sharp)} & &
(\Sh(\mathcal{C}'), \mathcal{O}')
\ar[d]^{(f, f^\sharp)} \\
(\Sh(\mathcal{C}/U), \mathcal{O}_U)
\ar[rr]^{(j, j^\sharp)} & &
(\Sh(\mathcal{C}), \mathcal{O}).
}
$$
且有 $f'_*(j')^{-1}=j^{-1}f_*$。参见《位点上的模》引理
\ref{sites-modules-lemma-localize-morphism-ringed-sites}。
由伴随的唯一性得到 $f^{-1}j_!=j'_!(f')^{-1}$。因此
\begin{align*}
f^*j_!\mathcal{A}_U
& =
f^{-1}j_!\mathcal{A}_U \otimes_{f^{-1}\mathcal{A}} \mathcal{A}' \\
& =
j'_!(f')^{-1}\mathcal{A}_U \otimes_{f^{-1}\mathcal{A}} \mathcal{A}' \\
& =
j'_!\left(
(f')^{-1}\mathcal{A}_U \otimes_{f^{-1}\mathcal{A}|_{U'}}
\mathcal{A}'|_{U'}\right) \\
& =
j'_!\mathcal{A}'_{U'}
\end{align*}
第一个等式是 $j_!\mathcal{A}_U$ 拉回为
$\mathcal{A}'$ 上微分分次模的定义。第二个等式来自
$f^{-1}j_!=j'_!(f')^{-1}$。第三个等式来自把引理
\ref{sdga-lemma-tensor-with-extension-by-zero-dg} 应用于环化位点
$(\mathcal{C}',f^{-1}\mathcal{O})$：其微分分次代数层为
$f^{-1}\mathcal{A}$，所用微分分次模分别是
$\mathcal{C}'/U'$ 上的 $(f')^{-1}\mathcal{A}_U$
与 $\mathcal{C}'$ 上的 $\mathcal{A}'$。第四个等式成立是因为
当然有 $(f')^{-1}\mathcal{A}_U=f^{-1}\mathcal{A}|_{U'}$。
由此可见，拉回仍是同类的模，而我们已在上文为它证明条件 (1) 与 (2)。
\end{proof}

\begin{lemma}
\label{sdga-lemma-good-admissible-ses}
设 $(\mathcal{C}, \mathcal{O})$ 是环化位点，
$\mathcal{A}$ 是 $(\mathcal{C}, \mathcal{O})$ 上的微分分次代数层。设
$0 \to \mathcal{P} \to \mathcal{P}' \to \mathcal{P}'' \to 0$
是微分分次 $\mathcal{A}$-模的容许短正合序列。
若这三个模中任意两个良好，则第三个也良好。
\end{lemma}

\begin{proof}
对条件 (1)，结论是立即的，因为该序列在分次层面是直和。
对条件 (2)，注意对任意微分分次左 $\mathcal{A}$-模，序列
$$
0 \to 
\mathcal{P} \otimes_\mathcal{A} \mathcal{N} \to
\mathcal{P}' \otimes_\mathcal{A} \mathcal{N} \to
\mathcal{P}'' \otimes_\mathcal{A} \mathcal{N} \to 0
$$
是微分分次 $\mathcal{O}$-模的容许短正合序列
（遗忘微分后，张量积正是在分次模范畴中取得的）。
因此，若三者中任意两个作为 $\mathcal{O}$-模复形正合，
则第三个也正合。最后，同一论证表明，给定环化拓扑斯态射
$(f, f^\sharp) : (\Sh(\mathcal{C}'), \mathcal{O}')
\to (\Sh(\mathcal{C}), \mathcal{O})$、微分分次
$\mathcal{O}'$-代数 $\mathcal{A}'$，以及微分分次
$f^{-1}\mathcal{O}$-代数映射
$\varphi:f^{-1}\mathcal{A}\to\mathcal{A}'$，序列
$$
0 \to f^*\mathcal{P} \to f^*\mathcal{P}' \to f^*\mathcal{P}'' \to 0
$$
是微分分次 $\mathcal{A}'$-模的容许短正合序列，
而上述论证在这里同样适用。
\end{proof}

\begin{lemma}
\label{sdga-lemma-good-direct-sum}
设 $(\mathcal{C}, \mathcal{O})$ 是环化位点，
$\mathcal{A}$ 是 $(\mathcal{C}, \mathcal{O})$ 上的微分分次代数层。
良好微分分次 $\mathcal{A}$-模的任意直和是良好的。
良好微分分次 $\mathcal{A}$-模的滤过余极限也是良好的。
\end{lemma}

\begin{proof}
从略。提示：直和与滤过余极限均同张量积及拉回交换。
\end{proof}

\begin{lemma}
\label{sdga-lemma-good-quotient}
设 $(\mathcal{C}, \mathcal{O})$ 是环化位点，
$\mathcal{A}$ 是 $(\mathcal{C}, \mathcal{O})$ 上的微分分次代数层，
$\mathcal{M}$ 是微分分次 $\mathcal{A}$-模。存在微分分次
$\mathcal{A}$-模同态 $\mathcal{P}\to\mathcal{M}$，满足
\begin{enumerate}
\item $\mathcal{P} \to \mathcal{M}$ 是满射；
\item $\Ker(\text{d}_\mathcal{P}) \to \Ker(\text{d}_\mathcal{M})$
是满射；
\item $\mathcal{P}$ 良好。
\end{enumerate}
\end{lemma}

\begin{proof}
考虑三元组 $(U,k,x)$，其中 $U$ 是 $\mathcal{C}$ 的对象，
$k\in\mathbf{Z}$，而 $x$ 是 $U$ 上满足
$\text{d}_\mathcal{M}(x)=0$ 的 $\mathcal{M}^k$ 截面。
于是得到唯一的微分分次 $\mathcal{A}_U$-模态射
$\varphi_x : \mathcal{A}_U[-k] \to \mathcal{M}|_U$
把 $1$ 映到 $x$。它伴随于态射
$\psi_x : j_{U!}\mathcal{A}_U[-k] \to \mathcal{M}$.
注意，$1\in\mathcal{A}_U(U)$ 对应于次数为 $k$ 的截面
$1\in j_{U!}\mathcal{A}_U[-k](U)$；其微分为零，且
$\psi_x$ 把它映到 $x$。因此考虑映射
$$
\bigoplus\nolimits_{(U, k, x)} j_{U!}\mathcal{A}_U[-k]
\longrightarrow
\mathcal{M}
$$
便得到条件 (2) 与 (3)。事实上，对象
$j_{U!}\mathcal{A}_U[-k]$ 是良好的（引理 \ref{sdga-lemma-supply-good}），
而良好对象的任意直和仍良好（引理 \ref{sdga-lemma-good-direct-sum}）。

\medskip\noindent
接着考虑三元组 $(U,k,x)$，其中 $U$ 是 $\mathcal{C}$ 的对象，
$k\in\mathbf{Z}$，而 $x$ 是 $\mathcal{M}^k$ 的截面
（不必被微分消去）。此时可以如评注 \ref{sdga-remark-cone-identity}
中那样考虑恒等映射 $\mathcal{A}_U\to\mathcal{A}_U$ 的锥 $C_U$。
元素 $x$ 决定映射
$\varphi_x:C_U[-k-1]\to\mathcal{A}_U$，参见评注
\ref{sdga-remark-cone-identity}。由于有容许短正合序列
$$
0 \to \mathcal{A}_U \to C_U \to \mathcal{A}_U[1] \to 0
$$
由引理 \ref{sdga-lemma-good-admissible-ses} 和已经用过的引理
\ref{sdga-lemma-supply-good}，可知 $j_{U!}C_U$ 是良好模。
同上，对伴随于 $\varphi_x$ 的映射
$\psi_x:j_{U!}C_U\to\mathcal{M}$ 取直和所得的映射
$$
\bigoplus\nolimits_{(U, k, x)} j_{U!}C_U \longrightarrow \mathcal{M}
$$
是满射。再与第一段构造的映射取直和，即得结论。
\end{proof}

\begin{remark}
\label{sdga-remark-sheaf-graded-sets}
设 $(\mathcal{C},\mathcal{O})$ 是环化位点。
$\mathcal{C}$ 上的{\it 分次集合层}是配有集合层映射
$\deg:\mathcal{S}\to\underline{\mathbf{Z}}$ 的集合层
$\mathcal{S}$。以 $\mathcal{O}[\mathcal{S}]$ 表示分次
$\mathcal{O}$-模，即由分次集合层 $\mathcal{S}$
生成的自由 $\mathcal{O}$-模。更确切地，
$\mathcal{O}[\mathcal{S}]$ 的第 $n$ 个分次部分是规则
$$
U \longmapsto
\bigoplus\nolimits_{s \in \mathcal{S}(U),\ \deg(s) = n} s \cdot \mathcal{O}(U)
$$
的层化。赋予零微分后，也可以把它看作微分分次
$\mathcal{O}$-模。设 $\mathcal{A}$ 是分次 $\mathcal{O}$-代数层。
类似地定义 $\mathcal{A}[\mathcal{S}]$ 为分次
$\mathcal{A}$-模，其第 $n$ 个分次部分是规则
$$
U \longmapsto
\bigoplus\nolimits_{s \in \mathcal{S}(U)} s \cdot \mathcal{A}^{n - \deg(s)}(U)
$$
的层化。若 $\mathcal{A}$ 是微分分次 $\mathcal{O}$-代数，
则对所有 $s\in\mathcal{S}(U)$ 置 $\text{d}(s)=0$ 并层化，
从而把它变成微分分次 $\mathcal{O}$-模。
\end{remark}

\begin{lemma}
\label{sdga-lemma-free-graded-module-good}
设 $(\mathcal{C},\mathcal{O})$ 是环化位点，
$\mathcal{A}$ 是微分分次 $\mathcal{A}$-代数，
$\mathcal{S}$ 是 $\mathcal{C}$ 上的分次集合层。
则由 $\mathcal{S}$ 生成的自由分次模
$\mathcal{A}[\mathcal{S}]$，赋予评注
\ref{sdga-remark-sheaf-graded-sets} 中的微分后，
是良好微分分次 $\mathcal{A}$-模。
\end{lemma}

\begin{proof}
设 $\mathcal{N}$ 是分次左 $\mathcal{A}$-模。则有
$$
\mathcal{A}[\mathcal{S}] \otimes_\mathcal{A} \mathcal{N} =
\mathcal{O}[\mathcal{S}] \otimes_\mathcal{O} \mathcal{N} =
\mathcal{N}[\mathcal{S}]
$$
其中 $\mathcal{N}[\mathcal{S}$ 是分次 $\mathcal{O}$-模，
其次数 $n$ 部分是与预层
$$
U \longmapsto
\bigoplus\nolimits_{s \in \mathcal{S}(U)} s \cdot \mathcal{N}^{n - \deg(s)}(U)
$$
相伴的层。显然，$\mathcal{N}\to\mathcal{N}[\mathcal{S}]$
是正合函子，因此 $\mathcal{A}[\mathcal{S}$
作为分次 $\mathcal{A}$-模是平坦的。接着，设
$\mathcal{N}$ 是微分分次左 $\mathcal{A}$-模。则作为分次
$\mathcal{O}$-模有
$$
H^*(\mathcal{A}[\mathcal{S}] \otimes_\mathcal{A} \mathcal{N}) =
H^*(\mathcal{O}[\mathcal{S}] \otimes_\mathcal{O} \mathcal{N})
$$
由其在 $\mathcal{O}$ 上的平坦性，这等于
$$
H^*(\mathcal{N})[\mathcal{S}]
$$
。因此，若 $\mathcal{N}$ 无环，则
$\mathcal{A}[\mathcal{S}]\otimes_\mathcal{A}\mathcal{N}$ 无环。

\medskip\noindent
最后，考虑环化拓扑斯态射
$(f, f^\sharp) : (\Sh(\mathcal{C}'), \mathcal{O}')
\to (\Sh(\mathcal{C}), \mathcal{O})$
、微分分次 $\mathcal{O}'$-代数 $\mathcal{A}'$，以及微分分次
$f^{-1}\mathcal{O}$-代数映射
$\varphi:f^{-1}\mathcal{A}\to\mathcal{A}'$。不难看出
$$
f^*\mathcal{A}[\mathcal{S}] = \mathcal{A}'[f^{-1}\mathcal{S}]
$$
这就完成了该模良好的证明。
\end{proof}

\begin{lemma}
\label{sdga-lemma-resolve}
设 $(\mathcal{C},\mathcal{O})$ 是环化位点，
$\mathcal{A}$ 是 $(\mathcal{C},\mathcal{O})$ 上的微分分次代数层，
$\mathcal{M}$ 是微分分次 $\mathcal{A}$-模。存在微分分次
$\mathcal{A}$-模同态 $\mathcal{P}\to\mathcal{M}$，满足
\begin{enumerate}
\item $\mathcal{P}\to\mathcal{M}$ 是拟同构；
\item $\mathcal{P}$ 良好。
\end{enumerate}
\end{lemma}

\begin{proof}[第一种证明]
令 $\mathcal{S}_0$ 为分次集合层
（评注 \ref{sdga-remark-sheaf-graded-sets}），其次数 $n$ 部分为
$\Ker(\text{d}_\mathcal{M}^n)$。考虑微分分次模同态
$$
\mathcal{P}_0 = \mathcal{A}[\mathcal{S}_0] \longrightarrow \mathcal{M}
$$
其中左端如评注 \ref{sdga-remark-sheaf-graded-sets} 所述，而该映射把
$\mathcal{S}_0$ 的局部截面 $s$ 映到
$\mathcal{M}^{\deg(s)}$ 的相应局部截面
（它属于微分的核，故这确实是微分分次模映射）。
按构造，诱导的上同调层映射
$H^n(\mathcal{P}_0)\to H^n(\mathcal{M})$ 是满射。
我们将归纳构造映射
$$
\mathcal{P}_0 \to \mathcal{P}_1 \to \mathcal{P}_2 \to \ldots \to \mathcal{M}
$$
显然，对所有 $i$，$H^*(\mathcal{P}_i)\to H^*(\mathcal{M})$
都将是满射。给定 $\mathcal{P}_i\to\mathcal{M}$，
以 $\mathcal{S}_{i+1}$ 表示如下分次集合层：其次数 $n$ 部分为
$$
\Ker(\text{d}_{\mathcal{P}_i}^{n + 1})
\times_{\mathcal{M}^{n + 1}, \text{d}}
\mathcal{M}^n
$$
置
$$
\mathcal{P}_{i + 1} = \mathcal{P}_i \oplus \mathcal{A}[\mathcal{S}_{i + 1}]
$$
作为分次 $\mathcal{A}$-模，其微分及到 $\mathcal{M}$ 的映射定义如下：
\begin{enumerate}
\item 对 $\mathcal{P}_i$ 的局部截面，使用 $\mathcal{P}_i$
上的微分以及给定的到 $\mathcal{M}$ 的映射；
\item 对 $\mathcal{S}_{i+1}$ 的局部截面 $s=(p,m)$，
置 $\text{d}(s)$ 等于 $p$（视为 $\mathcal{P}_i$
的次数 $\deg(s)+1$ 截面），并把 $s$ 映到 $\mathcal{M}$ 中的 $m$；
\item 唯一延拓微分，使其满足莱布尼茨法则。
\end{enumerate}
这是有意义的，因为 $\text{d}(m)$ 是 $p$ 的像，且
$\text{d}(p)=0$。最后置
$\mathcal{P}=\colim\mathcal{P}_i$，并赋予诱导的到
$\mathcal{M}$ 的映射。

\medskip\noindent
映射 $\mathcal{P}\to\mathcal{M}$ 是拟同构：
$H^n(\mathcal{P})=\colim H^n(\mathcal{P}_i)$，而对每个 $i$，
映射 $H^n(\mathcal{P}_i)\to H^n(\mathcal{M})$ 是满射；
按构造，其核被映射
$H^n(\mathcal{P}_i)\to H^n(\mathcal{P}_{i+1})$ 消去。
每个 $\mathcal{P}_i$ 都良好，因为由引理
\ref{sdga-lemma-free-graded-module-good}，$\mathcal{P}_0$ 良好，
而每个 $\mathcal{P}_{i+1}$ 都是容许短正合序列
$0 \to \mathcal{P}_i \to \mathcal{P}_{i + 1} \to
\mathcal{A}[\mathcal{S}_{i + 1}] \to 0$
的中间项；其两端由归纳假设均良好。因此由引理
\ref{sdga-lemma-good-admissible-ses}，$\mathcal{P}_{i+1}$ 良好。
最后，由引理 \ref{sdga-lemma-good-direct-sum} 得到 $\mathcal{P}$ 良好。
\end{proof}

\begin{proof}[第二种证明]
我们强烈建议读者在阅读本证明前，先阅读《微分分次代数》引理
\ref{dga-lemma-resolve} 的证明。置 $\mathcal{M}=\mathcal{M}_0$。
归纳选取短正合序列
$$
0 \to \mathcal{M}_{i + 1} \to \mathcal{P}_i \to \mathcal{M}_i \to 0
$$
其中映射 $\mathcal{P}_i\to\mathcal{M}_i$ 按引理
\ref{sdga-lemma-good-quotient} 选取。这给出一个``消解''
$$
\ldots \to \mathcal{P}_2 \xrightarrow{f_2}
\mathcal{P}_1 \xrightarrow{f_1} \mathcal{P}_0 \to \mathcal{M} \to 0
$$
然后令 $\mathcal{P}$ 为如下定义的微分分次 $\mathcal{A}$-模：
\begin{enumerate}
\item 作为分次 $\mathcal{A}$-模，置
$\mathcal{P}=\bigoplus_{a\leq 0}\mathcal{P}_{-a}[-a]$，即其次数
$n$ 部分为
$\mathcal{P}^n=\bigoplus\nolimits_{a+b=n}\mathcal{P}_{-a}^b$；
\item $\mathcal{P}$ 上的微分按与双复形相伴的全复形的构造给出：
$$
\text{d}_\mathcal{P}(x) = f_{-a}(x) + (-1)^a \text{d}_{\mathcal{P}_{-a}}(x)
$$
其中 $x$ 是 $\mathcal{P}_{-a}^b$ 的局部截面。
\end{enumerate}
在这些约定下，$\mathcal{P}$ 确为微分分次 $\mathcal{A}$-模；
细节从略。有微分分次 $\mathcal{A}$-模映射
$\mathcal{P}\to\mathcal{M}$：当 $a<0$ 时，它在直和项
$\mathcal{P}_{-a}[-a]$ 上为零；当 $a=0$ 时，它是给定映射
$\mathcal{P}_0\to\mathcal{M}$。注意
$$
\mathcal{P} = \colim_i F_i\mathcal{P}
$$
其中 $F_i\mathcal{P}\subset\mathcal{P}$ 是微分分次
$\mathcal{A}$-子模，其底层分次 $\mathcal{A}$-模为
$$
F_i\mathcal{P} =
\bigoplus\nolimits_{i \geq -a \geq 0} \mathcal{P}_{-a}[-a]
$$
立即可见，映射
$$
0 \to F_1\mathcal{P} \to F_2\mathcal{P} \to F_3\mathcal{P} \to
\ldots \to \mathcal{P}
$$
都是容许单态射，且有容许短正合序列
$$
0 \to F_i\mathcal{P} \to
F_{i + 1}\mathcal{P} \to
\mathcal{P}_{i + 1}[i + 1] \to 0
$$
由归纳法及引理 \ref{sdga-lemma-good-admissible-ses}，
$F_i\mathcal{P}$ 是良好微分分次 $\mathcal{A}$-模。
由于 $\mathcal{P}=\colim F_i\mathcal{P}$，由引理
\ref{sdga-lemma-good-direct-sum} 可知 $\mathcal{P}$ 良好。

\medskip\noindent
最后还需证明 $\mathcal{P}\to\mathcal{M}$ 是拟同构。
若 $\mathcal{C}$ 有足够多的点，则逐茎检验后，该结论由
《同调》引理 \ref{homology-lemma-good-resolution-gives-qis}
这一初等结果得出。一般情形可如下论证
（这个证明实在太长了；也可以像初等证明那样使用局部截面，
但那种论证也相当长）。由于滤过余极限在阿贝尔群层范畴上正合，有
$$
H^d(\mathcal{P}) = \colim H^d(F_i\mathcal{P})
$$
我们断言，对每个 $i\geq0$ 和 $d\in\mathbf{Z}$：
(a) 有短正合序列
$$
0 \to H^d(\mathcal{M}_{i + 1}[i]) \to
H^d(F_i\mathcal{P}) \to H^d(\mathcal{M}) \to 0
$$
其中第二个箭头来自
$F_i\mathcal{P}\to\mathcal{P}\to\mathcal{M}$；且
(b) 复合
$$
H^d(\mathcal{M}_{i + 1}[i]) \to
H^d(F_i\mathcal{P}) \to
H^d(F_{i + 1}\mathcal{P})
$$
为零。显然，该断言足以完成证明。

\medskip\noindent
证明该断言。对任意 $i\geq0$，有映射
$\mathcal{M}_{i+1}[i]\to F_i\mathcal{P}$，它来自
$\mathcal{M}_{i+1}$ 作为 $f_i$ 的核到 $\mathcal{P}_i$ 的包含。
考虑定义 $C_i$ 的 $\mathcal{O}$-模复形短正合序列
$$
0 \to \mathcal{M}_{i + 1}[i] \to F_i\mathcal{P} \to C_i \to 0
$$
注意 $C_0=\mathcal{M}_0=\mathcal{M}$。还要注意，
$C_i$ 是与双复形 $C_i^{\bullet,\bullet}$ 相伴的全复形；
该双复形的各列为
$$
\mathcal{M}_i = \mathcal{P}_i/\mathcal{M}_{i + 1},
\mathcal{P}_{i - 1}, \ldots, \mathcal{P}_0
$$
并分别处于次数 $-i,-i+1,\ldots,0$。有双复形映射
$C_i^{\bullet,\bullet}\to C_{i-1}^{\bullet,\bullet}$：
在次数 $-i$ 的列上为 $0$，在次数 $-i+1$ 上为满射
$\mathcal{P}_{i-1}\to\mathcal{M}_{i-1}$，在其余各列上为恒等映射。
因此有复形映射
$$
C_i \longrightarrow C_{i - 1}
$$
这些映射是满拟同构，因为其核是如下双复形的全复形：
次数 $-i,-i+1$ 的两列都是 $\mathcal{M}_i$，
而这两列之间是恒等映射。利用由此得到的等同
$H^d(C_i) = H^d(C_{i - 1} = \ldots = H^d(\mathcal{M})$
，已经可知由上述复形短正合序列得到长正合序列
$$
H^d(\mathcal{M}_{i + 1}[i]) \to
H^d(F_i\mathcal{P}) \to H^d(\mathcal{M}) \to
H^{d + 1}(\mathcal{M}_{i + 1}[i])
$$
。另一方面，还有交换图
$$
\xymatrix{
\mathcal{M}_{i + 2}[i + 1] \ar[r]_a & T_{i + 1} \ar[r] &
F_{i + 1}\mathcal{P} \ar[r] & C_{i + 1} \ar[d] \\
& \mathcal{M}_{i + 1}[i] \ar[r] \ar[u]^b &
F_i\mathcal{P} \ar[u] \ar[r] &
C_i
}
$$
其中 $T_{i+1}$ 是如下双复形的全复形：各列
$\mathcal{P}_{i+1},\mathcal{M}_{i+1}$ 分别置于次数
$-i-1,-i$。换言之，$T_{i+1}$ 是映射
$\mathcal{P}_{i+1}\to\mathcal{M}_{i+1}$ 的锥的一个平移。
于是 $a$ 是拟同构，而映射 $a^{-1}\circ b$ 是
$D(\mathcal{O})$ 中与短正合序列
$$
0 \to \mathcal{M}_{i + 2} \to \mathcal{P}_{i + 1} \to \mathcal{M}_{i + 1} \to 0
$$
相伴的杰出三角的第三个映射的一个平移。映射
$H^d(\mathcal{P}_{i+1})\to H^d(\mathcal{M}_{i+1})$
是满射，因为选取映射时要求
$\Ker(\text{d}_{\mathcal{P}_{i + 1}}) \to
\Ker(\text{d}_{\mathcal{M}_{i+1}})$ 为满射。因此
$a^{-1}\circ b$ 在上同调层上为零。这证明了断言的 (b)。
由于 $T_{i+1}$ 是复形满射
$F_{i+1}\mathcal{P}\to C_i$ 的核，得到上同调长正合序列的映射
$$
\xymatrix{
H^d(T_{i + 1}) \ar[r] &
H^d(F_{i + 1}\mathcal{P}) \ar[r] &
H^d(\mathcal{M}) \ar[r] &
H^{d + 1}(T_{i + 1}) \\
H^d(\mathcal{M}_{i + 1}[i]) \ar[r] \ar[u] &
H^d(F_i\mathcal{P}) \ar[r] \ar[u] &
H^d(\mathcal{M}) \ar[r] \ar[u] &
H^{d + 1}(\mathcal{M}_{i + 1}[i]) \ar[u]
}
$$
由上述讨论，外侧两个竖直映射为零。因此映射
$H^d(F_{i+1}\mathcal{P})\to H^d(\mathcal{M})$ 是满射，
断言的 (a) 随之成立。更确切地说，由此得到 $i>0$
时的断言；$i=0$ 的情形留给读者
（事实上，为得到本引理，只需对所有 $i\gg0$ 证明该断言）。
\end{proof}

\begin{lemma}
\label{sdga-lemma-acyclic-good}
设 $(\mathcal{C},\mathcal{O})$ 是环化位点，
$\mathcal{A}$ 是 $(\mathcal{C},\mathcal{O})$ 上的微分分次代数层。
设 $\mathcal{P}$ 是良好的无环右微分分次 $\mathcal{A}$-模。
\begin{enumerate}
\item 对任意微分分次左 $\mathcal{A}$-模 $\mathcal{N}$，
张量积 $\mathcal{P}\otimes_\mathcal{A}\mathcal{N}$ 无环；
\item 对任意环化拓扑斯态射
$(f, f^\sharp) : (\Sh(\mathcal{C}'), \mathcal{O}')
\to (\Sh(\mathcal{C}), \mathcal{O})$、任意微分分次
$\mathcal{O}'$-代数 $\mathcal{A}'$，以及任意微分分次
$f^{-1}\mathcal{O}$-代数映射
$\varphi:f^{-1}\mathcal{A}\to\mathcal{A}'$，
拉回 $f^*\mathcal{P}$ 无环且良好。
\end{enumerate}
\end{lemma}

\begin{proof}
证明 (1)。由引理 \ref{sdga-lemma-resolve}，可以选取良好左微分分次
$\mathcal{Q}$ 以及拟同构 $\mathcal{Q}\to\mathcal{N}$。于是
$\mathcal{P} \otimes_\mathcal{A} \mathcal{Q}$
无环，因为 $\mathcal{Q}$ 良好。令 $\mathcal{N}'$
为映射 $\mathcal{Q}\to\mathcal{N}$ 的锥。则
$\mathcal{P}\otimes_\mathcal{A}\mathcal{N}'$ 无环，因为
$\mathcal{P}$ 良好，而 $\mathcal{N}'$ 作为拟同构的锥是无环的。
按我们的三角结构构造，在
$K(\textit{Mod}(\mathcal{A},\text{d}))$ 中有杰出三角
$$
\mathcal{Q} \to \mathcal{N} \to \mathcal{N}' \to \mathcal{Q}[1]
$$
。由于 $\mathcal{P}\otimes_\mathcal{A}-$ 把杰出三角映到杰出三角，
在 $K(\textit{Mod}(\mathcal{O}))$ 中得到杰出三角
$$
\mathcal{P} \otimes_\mathcal{A} \mathcal{Q} \to
\mathcal{P} \otimes_\mathcal{A} \mathcal{N} \to
\mathcal{P} \otimes_\mathcal{A} \mathcal{N}' \to
\mathcal{P} \otimes_\mathcal{A} \mathcal{Q}[1]
$$
从而得出结论。

\medskip\noindent
证明 (2)。由良好模的定义，$f^*\mathcal{P}$ 良好。回忆
$f^*\mathcal{P} = f^{-1}\mathcal{P} \otimes_{f^{-1}\mathcal{A}} \mathcal{A}'$.
此时 $f^{-1}\mathcal{P}$ 是良好的无环微分分次
$f^{-1}\mathcal{A}$-模（因为 $f^{-1}$ 正合）。
因此由 (1) 可知 $f^*\mathcal{P}$ 无环。
\end{proof}






\section{分次模的微分分次包}
\label{sdga-section-dg-hull}

\noindent
分次模 $\mathcal{N}$ 的微分分次包，是应用下列引理中的函子 $G$ 所得的对象。

\begin{lemma}
\label{sdga-lemma-dg-hull}
设 $(\mathcal{C},\mathcal{O})$ 是环化位点，
$\mathcal{A}$ 是 $(\mathcal{C},\mathcal{O})$ 上的微分分次代数层。
遗忘函子
$F : \textit{Mod}(\mathcal{A}, \text{d}) \to \textit{Mod}(\mathcal{A})$
有左伴随
$G : \textit{Mod}(\mathcal{A})\to\textit{Mod}(\mathcal{A},\text{d})$。
\end{lemma}

\begin{proof}
为证明 $G$ 存在，可以使用伴随函子定理；参见《范畴》定理
\ref{categories-theorem-adjoint-functor}
（注意我们交换了 $F$ 与 $G$ 的角色）。由引理
\ref{sdga-lemma-dgm-abelian}，$F$ 满足正合性条件。
集合论条件可如下验证：设给定分次 $\mathcal{A}$-模
$\mathcal{N}$。对任意映射
$$
\varphi : \mathcal{N} \longrightarrow F(\mathcal{M})
$$
可以考虑满足
$\Im(\varphi)\subset F(\mathcal{M}')$ 的最小微分分次
$\mathcal{A}$-子模 $\mathcal{M}'\subset\mathcal{M}$。
显然，$\mathcal{M}'$ 是分次 $\mathcal{A}$-模映射
$$
\mathcal{N} \oplus
\mathcal{N}[-1] \otimes_\mathcal{O} \mathcal{A}
\longrightarrow
\mathcal{M}
$$
的像，该映射定义为
$$
(n, \sum n_i \otimes a_i) \longmapsto
\varphi(n) + \sum \text{d}(\varphi(n_i)) a_i
$$
，因为容易看出此映射的像是 $\mathcal{M}$ 的微分分次子模。
因此，这些 $\mathcal{M}'$ 的可能同构类数目有界，结论得证。
\end{proof}

\noindent
设 $(\mathcal{C},\mathcal{O})$ 是环化位点，
$\mathcal{A}$ 是 $(\mathcal{C},\mathcal{O})$ 上的微分分次代数层，
$\mathcal{M}$ 是微分分次 $\mathcal{A}$-模，并设在
$\textit{Mod}(\mathcal{A})$ 中有短正合序列
$$
0 \to \mathcal{N} \to F(\mathcal{M}) \to \mathcal{N}' \to 0
$$
。于是得到典范的分次 $\mathcal{A}$-模同态
$$
\overline{\text{d}} : \mathcal{N} \to \mathcal{N}'[1]
$$
，定义如下：给定 $\mathcal{N}$ 的局部截面 $x$，
把 $x$ 视为 $\mathcal{M}$ 的局部截面，并以
$\overline{\text{d}}(x)$ 表示
$\text{d}_\mathcal{M}(x)$ 在 $\mathcal{N}'$ 中的像。

\begin{lemma}
\label{sdga-lemma-dg-hull-acyclic}
引理 \ref{sdga-lemma-dg-hull} 中的函子 $F,G$ 具有如下性质。
给定分次 $\mathcal{A}$-模 $\mathcal{N}$，有
\begin{enumerate}
\item 余单位 $\mathcal{N}\to F(G(\mathcal{N}))$ 是单射；
\item 映射
$\overline{\text{d}}:\mathcal{N}\to
\Coker(\mathcal{N}\to F(G(\mathcal{N})))[1]$ 是同构；
\item $G(\mathcal{N})$ 是无环微分分次 $\mathcal{A}$-模。
\end{enumerate}
\end{lemma}

\begin{proof}
注意性质 (3) 是性质 (1) 与 (2) 的推论。事实上，若 $s$ 是
$F(G(\mathcal{N}))$ 的非零局部截面且 $\text{d}(s)=0$，
则 $s$ 不可能属于
$\mathcal{N}\to F(G(\mathcal{N}))$ 的像。因此，$s$
在余核中的像 $\overline{s}$ 可写为
$\overline{\text{d}}(s')$，其中 $s'$ 是 $\mathcal{N}$
的某个局部截面。于是 $s=\text{d}(s')$，因为差
$s-\text{d}(s')$ 仍属于 $\text{d}$ 的核，并且包含在余单位的像中。

\medskip\noindent
暂时以 $\mathcal{A}_{gr}$、$\mathcal{A}_{dg}$ 分别表示把层
$\mathcal{A}$ 视为其自身上的（右）分次模与（右）微分分次模。
本引理最重要的情形是理解 $G(\mathcal{A}_{gr})$。
当然，$G(\mathcal{A}_{gr})$ 是
$\textit{Mod}(\mathcal{A},\text{d})$ 中表示函子
$$
\mathcal{M} \longmapsto
\Hom_{\textit{Mod}(\mathcal{A})}(\mathcal{A}_{gr}, F(\mathcal{M})) =
\Gamma(\mathcal{C}, \mathcal{M})
$$
的对象。由评注 \ref{sdga-remark-cone-identity} 可知，此函子由
$C[-1]$ 表示，其中 $C$ 是 $\mathcal{A}_{dg}$ 恒等映射的锥。
在 $\textit{Mod}(\mathcal{A},\text{d})$ 中有短正合序列
$$
0 \to \mathcal{A}_{dg}[-1] \to C[-1] \to \mathcal{A}_{dg} \to 0
$$
，它在 $\textit{Mod}(\mathcal{A})$ 中由余单位
$\mathcal{A}_{gr}\to F(C[-1])$ 分裂。因此
$G(\mathcal{A}_{gr})$ 满足性质 (1) 与 (2)。

\medskip\noindent
设 $U$ 是 $\mathcal{C}$ 的对象。以
$j_U:\mathcal{C}/U\to\mathcal{C}$ 表示局部化态射，
以 $\mathcal{A}_U$ 表示 $\mathcal{A}$ 到 $U$ 的限制。
用记号 $\mathcal{A}_{U,gr}$ 表示把 $\mathcal{A}_U$
视为分次 $\mathcal{A}_U$-模。以
$F_U:\textit{Mod}(\mathcal{A}_U,\text{d})\to
\textit{Mod}(\mathcal{A}_U)$ 表示遗忘函子，以 $G_U$
表示其伴随。此时有交换图
$$
\vcenter{
\xymatrix{
\textit{Mod}(\mathcal{A}, \text{d}) \ar[d]_{j_U^*} \ar[r]_F &
\textit{Mod}(\mathcal{A}) \ar[d]^{j_U^*} \\
\textit{Mod}(\mathcal{A}_U, \text{d}) \ar[r]^{F_U} &
\textit{Mod}(\mathcal{A}_U)
}
}
\quad\text{且}\quad
\vcenter{
\xymatrix{
\textit{Mod}(\mathcal{A}_U, \text{d}) \ar[r]_{F_U} \ar[d]_{j_{U!}} &
\textit{Mod}(\mathcal{A}_U) \ar[d]^{j_{U!}} \\
\textit{Mod}(\mathcal{A}, \text{d}) \ar[r]^F &
\textit{Mod}(\mathcal{A})
}
}
$$
，这来自第 \ref{sdga-section-functoriality-graded}、
\ref{sdga-section-functoriality-dg}、\ref{sdga-section-localize-graded}
及 \ref{sdga-section-localize-dg} 节中对 $j_U^*$ 与 $j_{U!}$ 的构造。
由伴随的唯一性得到 $j_{U!}\circ G_U=G\circ j_{U!}$。
由于 $j_{U!}$ 是正合函子，证明前一部分已经得到余单位
$\mathcal{A}_{U,gr}\to F_U(G_U(\mathcal{A}_{U,gr}))$
的性质 (1) 与 (2)，故这些性质也适用于余单位
$j_{U!}\mathcal{A}_{U, gr} \to F(G(j_{U!}\mathcal{A}_{U, gr})) =
j_{U!}F_U(G_U(\mathcal{A}_{U, gr}))$.

\medskip\noindent
在引理 \ref{sdga-lemma-gm-grothendieck-abelian} 的证明中，
我们已经看到 $\textit{Mod}(\mathcal{A})$ 的任意对象都是若干
$j_{U!}\mathcal{A}_{U,gr}$ 副本之直和的商。
由于 $G$ 是左伴随，它与直和交换。因此，只要性质 (1) 与 (2)
对一族对象成立，就对它们的直和成立。由此，
$\textit{Mod}(\mathcal{A})$ 的每个对象 $\mathcal{N}$
都可置于正合序列
$$
\mathcal{N}_1 \to \mathcal{N}_0 \to \mathcal{N} \to 0
$$
中，使得 (1) 与 (2) 对 $\mathcal{N}_1$ 和 $\mathcal{N}_0$
成立。利用 $G$ 的右正合性，从中推出 $\mathcal{N}$
的 (1) 与 (2)，留给读者。
\end{proof}









\section{K-内射微分分次模}
\label{sdga-section-K-injective}

\noindent
本节是在微分分次代数层上的微分分次模层这一背景下，
与《内射对象》第 \ref{injectives-section-K-injective} 节对应的内容。

\begin{lemma}
\label{sdga-lemma-characterize-injectives}
设 $(\mathcal{C},\mathcal{O})$ 是环化位点，$\mathcal{A}$
是 $(\mathcal{C},\mathcal{O})$ 上的分次代数层。
存在集合 $T$，并且对每个 $t\in T$ 存在分次
$\mathcal{A}$-模的单射 $\mathcal{N}_t\to\mathcal{N}'_t$，
使得 $\textit{Mod}(\mathcal{A})$ 的对象 $\mathcal{I}$
是内射对象，当且仅当对每个实线图
$$
\xymatrix{
\mathcal{N}_t \ar[r] \ar[d] & \mathcal{I} \\
\mathcal{N}'_t \ar@{..>}[ru]
}
$$
在 $\textit{Mod}(\mathcal{A})$ 中存在使图交换的虚线箭头。
\end{lemma}

\begin{proof}
这在任意格罗滕迪克阿贝尔范畴中成立，参见《内射对象》引理
\ref{injectives-lemma-characterize-injective}。
由引理 \ref{sdga-lemma-gm-grothendieck-abelian}，
范畴 $\textit{Mod}(\mathcal{A})$ 是格罗滕迪克阿贝尔范畴。
\end{proof}

\begin{definition}
\label{sdga-definition-graded-injective}
设 $(\mathcal{C},\mathcal{O})$ 是环化位点，
$(\mathcal{A},\text{d})$ 是 $(\mathcal{C},\mathcal{O})$
上的微分分次代数层。若把 $\mathcal{M}$ 视为分次
$\mathcal{A}$-模时，它是分次 $\mathcal{A}$-模范畴
$\textit{Mod}(\mathcal{A})$ 的内射对象，则称微分分次
$\mathcal{A}$-模 $\mathcal{I}$ {\it 分次内射}
\footnote{这可能不是标准术语。}。
\end{definition}

\begin{remark}
\label{sdga-remark-why-graded-injective}
设 $(\mathcal{C},\mathcal{O})$ 是环化位点，
$(\mathcal{A},\text{d})$ 是 $(\mathcal{C},\mathcal{O})$
上的微分分次代数层。设 $\mathcal{I}$ 是分次内射微分分次
$\mathcal{A}$-模，并设
$$
0 \to \mathcal{M}_1 \to \mathcal{M}_2 \to \mathcal{M}_3 \to 0
$$
是微分分次 $\mathcal{A}$-模的短正合序列。
由于 $\mathcal{I}$ 分次内射，得到
$\Gamma(\mathcal{C},\mathcal{O})$-模复形的短正合序列
$$
0 \to
\Hom_{\textit{Mod}^{dg}(\mathcal{A}, \text{d})}(\mathcal{M}_3, \mathcal{I})
\to
\Hom_{\textit{Mod}^{dg}(\mathcal{A}, \text{d})}(\mathcal{M}_2, \mathcal{I})
\to
\Hom_{\textit{Mod}^{dg}(\mathcal{A}, \text{d})}(\mathcal{M}_1, \mathcal{I})
\to 0
$$
。取上同调得到同伦范畴中同态群的长正合序列
$$
\xymatrix{
\Hom_{K(\textit{Mod}(\mathcal{A}, \text{d}))}(\mathcal{M}_3, \mathcal{I})
\ar[d] &
\Hom_{K(\textit{Mod}(\mathcal{A}, \text{d}))}(\mathcal{M}_3, \mathcal{I})[1]
\ar[d] \\
\Hom_{K(\textit{Mod}(\mathcal{A}, \text{d}))}(\mathcal{M}_2, \mathcal{I})
\ar[d] &
\Hom_{K(\textit{Mod}(\mathcal{A}, \text{d}))}(\mathcal{M}_2, \mathcal{I})[1]
\ar[d] \\
\Hom_{K(\textit{Mod}(\mathcal{A}, \text{d}))}(\mathcal{M}_1, \mathcal{I})
\ar[ruu]
&
\Hom_{K(\textit{Mod}(\mathcal{A}, \text{d}))}(\mathcal{M}_1, \mathcal{I})[1]
}
$$
。关键在于，即使没有假设原短正合序列容许，我们仍得到这一序列
（所以一般而言，该短正合序列并不在同伦范畴中定义杰出三角）。
\end{remark}

\begin{lemma}
\label{sdga-lemma-product-graded-injective}
设 $(\mathcal{C},\mathcal{O})$ 是环化位点，
$(\mathcal{A},\text{d})$ 是 $(\mathcal{C},\mathcal{O})$
上的微分分次代数层。设 $T$ 是集合，并且对每个 $t\in T$，
$\mathcal{I}_t$ 是分次内射微分分次 $\mathcal{A}$-模。
则 $\prod\mathcal{I}_t$ 是分次内射微分分次 $\mathcal{A}$-模。
\end{lemma}

\begin{proof}
这是因为内射对象的直积仍内射，参见《同调》引理
\ref{homology-lemma-product-injectives}；还因为通过遗忘函子，
$\textit{Mod}(\mathcal{A},\text{d})$ 中的直积与
$\textit{Mod}(\mathcal{A})$ 中的直积相容。
\end{proof}

\begin{lemma}
\label{sdga-lemma-characterize-graded-injectives-in-dg}
设 $(\mathcal{C},\mathcal{O})$ 是环化位点，
$(\mathcal{A},\text{d})$ 是 $(\mathcal{C},\mathcal{O})$
上的微分分次代数层。存在集合 $T$，并且对每个 $t\in T$
存在无环微分分次 $\mathcal{A}$-模的单射
$\mathcal{M}_t\to\mathcal{M}'_t$，使得对
$\textit{Mod}(\mathcal{A},\text{d})$ 的对象 $\mathcal{I}$，
下列条件等价：
\begin{enumerate}
\item $\mathcal{I}$ 分次内射；
\item 对每个实线图
$$
\xymatrix{
\mathcal{M}_t \ar[r] \ar[d] & \mathcal{I} \\
\mathcal{M}'_t \ar@{..>}[ru]
}
$$
在 $\textit{Mod}(\mathcal{A},\text{d})$ 中存在使图交换的虚线箭头。
\end{enumerate}
\end{lemma}

\begin{proof}
取引理 \ref{sdga-lemma-characterize-injectives} 中的
$T$ 与 $\mathcal{N}_t\to\mathcal{N}'_t$。
以 $F:\textit{Mod}(\mathcal{A},\text{d})\to
\textit{Mod}(\mathcal{A})$ 表示遗忘函子，
以 $G$ 表示引理 \ref{sdga-lemma-dg-hull} 中 $F$ 的左伴随函子。置
$$
\mathcal{M}_t = G(\mathcal{N}_t) \to
G(\mathcal{N}'_t) = \mathcal{M}'_t
$$
由引理 \ref{sdga-lemma-dg-hull-acyclic}，这是无环微分分次
$\mathcal{A}$-模的单射。由于 $G$ 是 $F$ 的左伴随，
下图中存在虚线箭头
$$
\xymatrix{
\mathcal{M}_t \ar[r] \ar[d] & \mathcal{I} \\
\mathcal{M}'_t \ar@{..>}[ru]
}
$$
当且仅当下图中存在虚线箭头：
$$
\xymatrix{
\mathcal{N}_t \ar[r] \ar[d] & F(\mathcal{I}) \\
\mathcal{N}'_t \ar@{..>}[ru]
}
$$
因此，结论由箭头族
$\mathcal{N}_t\to\mathcal{N}_t'$ 的选取方式得出。
\end{proof}


\begin{lemma}
\label{sdga-lemma-small-acyclics}
设 $(\mathcal{C},\mathcal{O})$ 是环化位点，
$(\mathcal{A},\text{d})$ 是 $(\mathcal{C},\mathcal{O})$
上的微分分次代数层。存在集合 $S$，并且对每个 $s$
存在无环微分分次 $\mathcal{A}$-模 $\mathcal{M}_s$，
使得对每个非零无环微分分次 $\mathcal{A}$-模 $\mathcal{M}$，
存在 $s\in S$ 以及 $\textit{Mod}(\mathcal{A},\text{d})$
中的单射 $\mathcal{M}_s\to\mathcal{M}$。
\end{lemma}

\begin{proof}
开始之前回忆，我们的约定保证位点 $\mathcal{C}$
的对象和态射构成集合，并有一个覆盖集合
$\text{Cov}(\mathcal{C})$。若 $\mathcal{F}$
是微分分次 $\mathcal{A}$-模，定义 $|\mathcal{F}|$ 为
$$
\coprod\nolimits_{(U, n)} \mathcal{F}^n(U)
$$
的基数之和，其中 $U$ 遍历 $\mathcal{C}$ 的对象，
$n\in\mathbf{Z}$。选取无限基数 $\kappa$，使其大于
$|\Ob(\mathcal{C})|$、$|\text{Arrows}(\mathcal{C})|$、
$|\text{Cov}(\mathcal{C})|$、所有
$\{U_i\to U\}_{i\in I}\in\text{Cov}(\mathcal{C})$
所对应的 $\sup|I|$，以及 $|\mathcal{A}|$。

\medskip\noindent
设 $\mathcal{F}\subset\mathcal{M}$ 是微分分次
$\mathcal{A}$-模的包含。设给定集合 $K$，并且对每个 $k\in K$
给定三元组 $(U_k,n_k,x_k)$，其中 $U_k$ 是 $\mathcal{C}$
的对象，$n_k$ 是整数，而
$x_k\in\mathcal{M}^{n_k}(U_k)$ 是截面。可以考虑包含
$\mathcal{F}$ 以及所有 $k\in K$ 对应截面 $x_k$ 的最小微分分次
$\mathcal{A}$-子模 $\mathcal{F}'\subset\mathcal{M}$。可以把
$$
(\mathcal{F}')^n(U) \subset \mathcal{M}^n(U)
$$
描述为由满足如下条件的元素 $x\in\mathcal{M}^n(U)$ 组成的集合：
存在 $\{f_i:U_i\to U\}_{i\in I}\in\text{Cov}(\mathcal{C})$，
使得对每个 $i\in I$，存在有限集 $T_i$ 及态射
$g_t:U_i\to U_{k_t}$，且
$$
f_i^*x = y_i +
\sum\nolimits_{t \in T_i} a_{it}g_t^*x_{k_t} + b_{it}g_t^*\text{d}(x_{k_t})
$$
，其中 $y_i\in\mathcal{F}^n(U)$，并且
$a_{it}\in\mathcal{A}^{n-n_{k_t}}(U_i)$、
$b_{it}\in\mathcal{A}^{n-n_{k_t}-1}(U_i)$。
（细节从略；提示：这些截面显然属于 $\mathcal{F}'$，
反过来则证明这条规则定义微分分次 $\mathcal{A}$-子模。）
由此描述可知
$|\mathcal{F}'| \leq \max(|\mathcal{F}|, |K|, \kappa)$.

\medskip\noindent
设 $\mathcal{M}$ 是非零无环微分分次 $\mathcal{A}$-模。
则可找到整数 $n$ 以及 $\mathcal{M}^n$ 在
$\mathcal{C}$ 某个对象 $U$ 上的非零截面 $x$。令
$$
\mathcal{F}_0 \subset \mathcal{M}
$$
为包含 $x$ 的最小微分分次 $\mathcal{A}$-子模。
由上一段，$|\mathcal{F}_0|\leq\kappa$。归纳地，给定
$\mathcal{F}_0,\ldots,\mathcal{F}_n$ 后，如下定义
$\mathcal{F}_{n+1}$。考虑三元组集合
$$
L = \{(U, n, x)\}
\{U_i \to U\}_{i \in I}, (x_i)_{i \in I})\}
$$
，其中 $U$ 是 $\mathcal{C}$ 的对象，$n\in\mathbf{Z}$，
$x\in\mathcal{F}_n(U)$ 且 $\text{d}(x)=0$。
由于 $\mathcal{M}$ 无环，对每个三元组
$l=(U_l,n_l,x_l)\in L$，可以选取
$\{(U_{l, i} \to U_l\}_{i \in I_l} \in \text{Cov}(\mathcal{C})$，以及
$x_{l, i} \in \mathcal{M}^{n_l - 1}(U_{l, i})$，使得
$\text{d}(x_{l, i}) = x|_{U_{l, i}}$。然后置
$$
K = \{(U_{l, i}, n_l - 1, x_{l, i}) \mid l \in L, i \in I_l\}
$$
并令 $\mathcal{F}_{n+1}$ 为 $\mathcal{M}$ 中包含
$\mathcal{F}_n$ 及所有截面 $x_{l,i}$ 的最小微分分次
$\mathcal{A}$-子模。由于
$|K|\leq\max(\kappa,|\mathcal{F}_n|)$，归纳可得
$|\mathcal{F}_{n+1}|\leq\kappa$。

\medskip\noindent
按构造，包含 $\mathcal{F}_n\to\mathcal{F}_{n+1}$
在上同调层上诱导零映射。因此
$\mathcal{F}=\bigcup\mathcal{F}_n$ 是非零无环子模，且
$|\mathcal{F}|\leq\kappa$。由于满足 $|\mathcal{F}|$
有界的微分分次 $\mathcal{A}$-模 $\mathcal{F}$
的同构类只构成一个集合，结论得证。
\end{proof}

\begin{definition}
\label{sdga-definition-K-injective}
设 $(\mathcal{C},\mathcal{O})$ 是环化位点，
$(\mathcal{A},\text{d})$ 是 $(\mathcal{C},\mathcal{O})$
上的微分分次代数层。若对每个无环微分分次 $\mathcal{M}$ 都有
$$
\Hom_{K(\textit{Mod}(\mathcal{A}, \text{d}))}(\mathcal{M}, \mathcal{I}) = 0
$$
，则称微分分次 $\mathcal{A}$-模 $\mathcal{I}$ {\it K-内射}。
\end{definition}

\noindent
请注意这与《导出范畴》定义
\ref{derived-definition-K-injective} 的相似之处。

\begin{lemma}
\label{sdga-lemma-product-K-injective}
设 $(\mathcal{C},\mathcal{O})$ 是环化位点，
$(\mathcal{A},\text{d})$ 是 $(\mathcal{C},\mathcal{O})$
上的微分分次代数层。设 $T$ 是集合，并且对每个 $t\in T$，
$\mathcal{I}_t$ 是 K-内射微分分次 $\mathcal{A}$-模。
则 $\prod\mathcal{I}_t$ 是 K-内射微分分次 $\mathcal{A}$-模。
\end{lemma}

\begin{proof}
设 $\mathcal{K}$ 是无环微分分次 $\mathcal{A}$-模。则有
$$
\Hom_{\textit{Mod}^{dg}(\mathcal{A}, \text{d})}(\mathcal{K},
\prod\nolimits_{t \in T} \mathcal{I}_t)
=
\prod\nolimits_{t \in T}
\Hom_{\textit{Mod}^{dg}(\mathcal{A}, \text{d})}(\mathcal{K}, \mathcal{I}_t)
$$
，因为在 $\textit{Mod}(\mathcal{A},\text{d})$ 中取直积，
与到分次 $\mathcal{A}$-模范畴的遗忘函子交换。
由于取直积在阿贝尔群范畴上是正合函子，结论得证。
\end{proof}

\begin{lemma}
\label{sdga-lemma-first-property-dg-injective}
设 $(\mathcal{C},\mathcal{O})$ 是环化位点，
$(\mathcal{A},\text{d})$ 是 $(\mathcal{C},\mathcal{O})$
上的微分分次代数层。设 $\mathcal{I}$ 是
$\textit{Mod}(\mathcal{A},\text{d})$ 中既 K-内射又分次内射的对象。
对 $\textit{Mod}(\mathcal{A},\text{d})$ 中任意实线图
$$
\xymatrix{
\mathcal{M} \ar[r]_a \ar[d]_b & \mathcal{I} \\
\mathcal{M}' \ar@{..>}[ru]
}
$$
，若 $b$ 是单射且 $\mathcal{M}$ 无环，则存在使图交换的虚线箭头。
\end{lemma}

\begin{proof}
由于 $\mathcal{M}$ 无环且 $\mathcal{I}$ K-内射，
存在次数为 $-1$ 的分次 $\mathcal{A}$-模映射
$h:\mathcal{M}\to\mathcal{I}$，使得 $a=\text{d}(h)$。
由于 $\mathcal{I}$ 分次内射且 $b$ 是单射，
存在次数为 $-1$ 的分次 $\mathcal{A}$-模映射
$h':\mathcal{M}'\to\mathcal{I}$，使得 $h=h'\circ b$。
于是可取 $a'=\text{d}(h')$ 为虚线箭头。
\end{proof}

\begin{lemma}
\label{sdga-lemma-second-property-dg-injective}
设 $(\mathcal{C},\mathcal{O})$ 是环化位点，
$(\mathcal{A},\text{d})$ 是 $(\mathcal{C},\mathcal{O})$
上的微分分次代数层。设 $\mathcal{I}$ 是
$\textit{Mod}(\mathcal{A},\text{d})$ 中既 K-内射又分次内射的对象。
对 $\textit{Mod}(\mathcal{A},\text{d})$ 中任意实线图
$$
\xymatrix{
\mathcal{M} \ar[r]_a \ar[d]_b & \mathcal{I} \\
\mathcal{M}' \ar@{..>}[ru]
}
$$
，若 $b$ 是拟同构，则存在使图在同伦意义下交换的虚线箭头。
\end{lemma}

\begin{proof}
把 $\mathcal{M}'$ 替换为 $\mathcal{M}'$ 与
$\mathcal{M}$ 恒等映射之锥（后者无环）的直和后，
可设 $b$ 也是单射。此时 $b$ 的余核 $\mathcal{Q}$ 无环。因此
$$
\Hom_{K(\textit{Mod}(\mathcal{A}, \text{d}))}(\mathcal{Q}, \mathcal{I}) =
\Hom_{K(\textit{Mod}(\mathcal{A}, \text{d}))}(\mathcal{Q}, \mathcal{I})[1] = 0
$$
，因为 $\mathcal{I}$ K-内射。又因 $\mathcal{I}$ 分次内射，
由评注 \ref{sdga-remark-why-graded-injective} 可知
$$
\Hom_{K(\textit{Mod}(\mathcal{A}, \text{d}))}(\mathcal{M}', \mathcal{I})
\longrightarrow
\Hom_{K(\textit{Mod}(\mathcal{A}, \text{d}))}(\mathcal{M}, \mathcal{I})
$$
是双射，证明完成。
\end{proof}

\begin{lemma}
\label{sdga-lemma-better-set-of-monos}
设 $(\mathcal{C},\mathcal{O})$ 是环化位点，
$(\mathcal{A},\text{d})$ 是 $(\mathcal{C},\mathcal{O})$
上的微分分次代数层。存在集合 $R$，并且对每个 $r\in R$
存在无环微分分次 $\mathcal{A}$-模的单射
$\mathcal{M}_r\to\mathcal{M}'_r$，使得对
$\textit{Mod}(\mathcal{A},\text{d})$ 的对象 $\mathcal{I}$，
下列条件等价：
\begin{enumerate}
\item $\mathcal{I}$ 既 K-内射又分次内射；
\item 对每个实线图
$$
\xymatrix{
\mathcal{M}_r \ar[r] \ar[d] & \mathcal{I} \\
\mathcal{M}'_r \ar@{..>}[ru]
}
$$
在 $\textit{Mod}(\mathcal{A},\text{d})$ 中存在使图交换的虚线箭头。
\end{enumerate}
\end{lemma}

\begin{proof}
取引理 \ref{sdga-lemma-characterize-graded-injectives-in-dg} 中的
$T$ 与 $\mathcal{M}_t\to\mathcal{M}'_t$，并取引理
\ref{sdga-lemma-small-acyclics} 中的 $S$ 与 $\mathcal{M}_s$。
选取无环微分分次 $\mathcal{A}$-模的单射
$\mathcal{M}_s\to\mathcal{M}'_s$，使其零同伦。
这是可能的，因为可取 $\mathcal{M}'_s$ 为恒等映射的锥；
在此情形下，甚至 $\mathcal{M}'_s$ 上的恒等映射也零同伦，
参见《微分分次代数》引理 \ref{dga-lemma-id-cone-null}；
由第 \ref{sdga-section-conclude-triangulated} 节的讨论，该引理适用。
我们断言，$R=T\coprod S$ 连同给定映射满足要求。

\medskip\noindent
蕴含 (1) $\Rightarrow$ (2) 由引理
\ref{sdga-lemma-first-property-dg-injective} 得出。

\medskip\noindent
假设 (2)。首先，由引理
\ref{sdga-lemma-characterize-graded-injectives-in-dg} 可知
$\mathcal{I}$ 分次内射。其次，设 $\mathcal{M}$
是无环微分分次 $\mathcal{A}$-模。需要证明
$$
\Hom_{K(\textit{Mod}(\mathcal{A}, \text{d}))}(\mathcal{M}, \mathcal{I}) = 0
$$
其证明与《内射对象》引理
\ref{injectives-lemma-characterize-K-injective} 的证明完全相同。

\medskip\noindent
我们将按序数 $\alpha$ 归纳构造无环微分分次子模
$\mathcal{K}_\alpha\subset\mathcal{M}$。当 $\alpha=0$ 时，
置 $\mathcal{K}_0=0$。当 $\alpha>0$ 时，按如下方式进行：
\begin{enumerate}
\item 若 $\alpha=\beta+1$ 且 $\mathcal{K}_\beta=\mathcal{M}$，
则取 $\mathcal{K}_\alpha=\mathcal{K}_\beta$。
\item 若 $\alpha=\beta+1$ 且 $\mathcal{K}_\beta\ne\mathcal{M}$，
则 $\mathcal{M}/\mathcal{K}_\beta$ 是非零无环微分分次
$\mathcal{A}$-模。选取微分分次 $\mathcal{A}$-子模
$\mathcal{N}_\alpha\subset\mathcal{M}/\mathcal{K}_\beta$，
使其同构于某个 $s\in S$ 所对应的 $\mathcal{M}_s$；
参见引理 \ref{sdga-lemma-small-acyclics}。最后令
$\mathcal{K}_\alpha\subset\mathcal{M}$ 为 $\mathcal{N}_\alpha$ 的逆像。
\item 若 $\alpha$ 是极限序数，则置
$\mathcal{K}_\beta = \colim \mathcal{K}_\alpha$.
\end{enumerate}
显然，对某个充分大的序数 $\alpha$，有
$\mathcal{M}=\mathcal{K}_\alpha$。我们将对 $\alpha$
作超限归纳，证明
$$
\Hom_{K(\textit{Mod}(\mathcal{A}, \text{d}))}(\mathcal{K}_\alpha, \mathcal{I})
$$
为零。当 $\alpha=0$ 时成立，因为 $\mathcal{K}_0=0$。
假设对 $\beta$ 成立，且 $\alpha=\beta+1$。
在上述 (1) 情形中结论显然。在 (2) 情形中有短正合序列
$$
0 \to \mathcal{K}_\beta \to \mathcal{K}_\alpha \to \mathcal{N}_\alpha \to 0
$$
。由评注 \ref{sdga-remark-why-graded-injective} 以及已经证明的
$\mathcal{I}$ 分次内射性，得到正合序列
$$
\Hom_{K(\textit{Mod}(\mathcal{A}, \text{d}))}(\mathcal{K}_\beta, \mathcal{I})
\to
\Hom_{K(\textit{Mod}(\mathcal{A}, \text{d}))}(\mathcal{K}_\alpha, \mathcal{I})
\to
\Hom_{K(\textit{Mod}(\mathcal{A}, \text{d}))}(\mathcal{N}_\alpha, \mathcal{I})
$$
。左项由归纳假设为零。右项由假设 (2) 为零：任意映射
$\mathcal{M}_s\to\mathcal{I}$ 都经 $\mathcal{M}'_s$ 分解，
因而零同伦。所以中间群也为零。最后，设 $\alpha$ 是极限序数。
由于作为分次 $\mathcal{A}$-模还有
$\mathcal{K}_\alpha = \colim \mathcal{K}_\alpha$
，可见
$$
\Hom_{\textit{Mod}^{dg}(\mathcal{A}, \text{d})}
(\mathcal{K}_\alpha, \mathcal{I})
= \lim_{\beta < \alpha}
\Hom_{\textit{Mod}^{dg}(\mathcal{A}, \text{d})}
(\mathcal{K}_\beta, \mathcal{I})
$$
是阿贝尔群复形的等式。这些复形的上同调群计算
$K(\textit{Mod}(\mathcal{A},\text{d}))$ 中各平移之间的态射。
由于 $\mathcal{I}$ 分次内射，由评注
\ref{sdga-remark-why-graded-injective}，此复形系统中的转移映射为满射。
此外，对极限序数 $\beta\leq\alpha$，极限与该处的值相等。
故可应用《同调》引理 \ref{homology-lemma-ML-over-ordinals} 得出结论。
\end{proof}

\begin{lemma}
\label{sdga-lemma-functor-set-of-monos}
设 $(\mathcal{C},\mathcal{O})$ 是环化位点，
$(\mathcal{A},\text{d})$ 是 $(\mathcal{C},\mathcal{O})$
上的微分分次代数层。设 $R$ 是集合，并且对每个 $r\in R$
给定无环微分分次 $\mathcal{A}$-模的单射
$\mathcal{M}_r\to\mathcal{M}'_r$。存在函子
$M:\textit{Mod}(\mathcal{A},\text{d})\to
\textit{Mod}(\mathcal{A},\text{d})$ 与自然变换
$j:\text{id}\to M$，使得
\begin{enumerate}
\item $j_\mathcal{M}:\mathcal{M}\to M(\mathcal{M})$
是单射且为拟同构；
\item 对每个实线图
$$
\xymatrix{
\mathcal{M}_r \ar[r] \ar[d] & \mathcal{M} \ar[d]^{j_\mathcal{M}} \\
\mathcal{M}'_r \ar@{..>}[r] & M(\mathcal{M})
}
$$
在 $\textit{Mod}(\mathcal{A},\text{d})$ 中存在使图交换的虚线箭头。
\end{enumerate}
\end{lemma}

\begin{proof}
把 $M(\mathcal{M})$ 定义为下图中的推出：
$$
\xymatrix{
\bigoplus_{(r, \varphi)} \mathcal{M}_r \ar[r] \ar[d] &
\mathcal{M} \ar[d] \\
\bigoplus_{(r, \varphi)} \mathcal{M}'_r \ar[r] &
M(\mathcal{M})
}
$$
其中直和遍历满足 $r\in R$ 以及
$\varphi\in
\Hom_{\textit{Mod}(\mathcal{A}, \text{d})}(\mathcal{M}_r, \mathcal{M})$.
的所有偶对 $(r,\varphi)$。由于单射的推出仍为单射，
$\mathcal{M}\to M(\mathcal{M})$ 是单射。由于左侧竖直箭头
的余核无环，$\mathcal{M}\to M(\mathcal{M})$
与之同构的余核也无环，所以该映射是拟同构。
性质 (2) 由构造成立。此过程确实可构成函子的验证从略。
\end{proof}

\begin{theorem}
\label{sdga-theorem-qis-into-dg-injective}
设 $(\mathcal{C},\mathcal{O})$ 是环化位点，
$(\mathcal{A},\text{d})$ 是 $(\mathcal{C},\mathcal{O})$
上的微分分次代数层。对每个微分分次 $\mathcal{A}$-模
$\mathcal{M}$，存在拟同构 $\mathcal{M}\to\mathcal{I}$，
其中 $\mathcal{I}$ 是既分次内射又 K-内射的微分分次
$\mathcal{A}$-模。此外，该构造关于 $\mathcal{M}$ 具有函子性。
\end{theorem}

\begin{proof}
取引理 \ref{sdga-lemma-better-set-of-monos} 中找到的集合 $R$
以及 $\textit{Mod}(\mathcal{A},\text{d})$ 的态射
$\mathcal{M}_r\to\mathcal{M}'_r$。取引理
\ref{sdga-lemma-functor-set-of-monos} 利用这些数据构造的
$M$ 与自然变换 $\text{id}\to M$。使用超限递归，
如下定义函子序列 $M_\alpha$ 以及当 $\alpha<\beta$
时的自然变换 $M_\beta\to M_\alpha$：
\begin{enumerate}
\item $M_0=\text{id}$；
\item $M_{\alpha+1}=M\circ M_\alpha$；当 $\beta<\alpha+1$ 时，
自然变换 $M_\beta\to M_{\alpha+1}$ 来自已经构造的
$M_\beta\to M_\alpha$，以及由 $\text{id}\to M$ 所得的映射
$M_\alpha\to M\circ M_\alpha$；
\item 若 $\alpha$ 是极限序数，则
$M_\alpha=\colim_{\beta<\alpha}M_\beta$，并以余投射作为
当 $\alpha<\beta$ 时的变换 $M_\beta\to M_\alpha$。
\end{enumerate}
注意，对每个微分分次 $\mathcal{A}$-模，映射
$\mathcal{M} \to M_\beta(\mathcal{M}) \to M_\alpha(\mathcal{M})$
都是单的拟同构（因为滤过余极限正合）。

\medskip\noindent
回忆 $\textit{Mod}(\mathcal{A},\text{d})$ 是格罗滕迪克阿贝尔范畴。
因此，由《内射对象》命题
\ref{injectives-proposition-objects-are-small}
（应用于所有 $r\in R$ 对应的 $\mathcal{M}_r$ 的直和），
存在极限序数 $\alpha$，使每个 $r\in R$ 对应的
$\mathcal{M}_r$ 关于单射都是 $\alpha$-小的。
我们断言 $\mathcal{M}\to M_\alpha(\mathcal{M})$
正是所需的函子性嵌入，它把 $\mathcal{M}$
嵌入一个分次内射且 K-内射的模中。

\medskip\noindent
事实上，任意映射
$\mathcal{M}_r\to M_\alpha(\mathcal{M})$ 都经某个
$\beta<\alpha$ 对应的 $M_\beta(\mathcal{M})$ 分解。
而由 $M$ 的构造，这意味着映射
$\mathcal{M}_r\to M_{\beta+1}(\mathcal{M})
=M(M_\beta(\mathcal{M}))$ 经 $\mathcal{M}'_r$ 分解。
由于
$M_\beta(\mathcal{M})\subset M_{\beta+1}(\mathcal{M})
\subset M_\alpha(\mathcal{M})$，便得到所需的到
$M_\alpha(\mathcal{M})$ 的分解。最后由引理
\ref{sdga-lemma-better-set-of-monos} 中对
$R$ 及 $\mathcal{M}_r\to\mathcal{M}'_r$ 的选取得出结论。
\end{proof}










\section{导出范畴}
\label{sdga-section-derived}

\noindent
本节对应《微分分次代数》第 \ref{dga-section-derived} 节。

\medskip\noindent
设 $(\mathcal{C},\mathcal{O})$ 是环化位点，
$(\mathcal{A},\text{d})$ 是 $(\mathcal{C},\mathcal{O})$
上的微分分次代数层。我们将通过逆转
$K(\textit{Mod}(\mathcal{A},\text{d}))$ 中的拟同构来构造
导出范畴 $D(\mathcal{A},\text{d})$。

\begin{lemma}
\label{sdga-lemma-cohomology-homological}
设 $(\mathcal{C},\mathcal{O})$ 是环化位点，
$(\mathcal{A},\text{d})$ 是 $(\mathcal{C},\mathcal{O})$
上的微分分次代数层。第 \ref{sdga-section-modules} 节的函子
$H^0 : \textit{Mod}(\mathcal{A}, \text{d}) \to \textit{Mod}(\mathcal{O})$
经函子
$$
H^0 : K(\textit{Mod}(\mathcal{A}, \text{d})) \to \textit{Mod}(\mathcal{O})
$$
分解；后一函子在《导出范畴》定义
\ref{derived-definition-homological} 的意义下是同调函子。
\end{lemma}

\begin{proof}
由定义立即得到交换图
$$
\xymatrix{
\textit{Mod}(\mathcal{A}, \text{d}) \ar[r] \ar[d] &
K(\textit{Mod}(\mathcal{A}, \text{d})) \ar[d] \\
\text{Comp}(\mathcal{O}) \ar[r] &
K(\textit{Mod}(\mathcal{O}))
}
$$
由于 $H^0(\mathcal{M})$ 定义为 $\mathcal{M}$ 的底层
$\mathcal{O}$-模复形的第零上同调层，本引理由
$\mathcal{O}$-模复形的情形得出；后者是《导出范畴》引理
\ref{derived-lemma-cohomology-homological} 的特例。
\end{proof}

\begin{lemma}
\label{sdga-lemma-acyclics}
设 $(\mathcal{C},\mathcal{O})$ 是环化位点，
$(\mathcal{A},\text{d})$ 是 $(\mathcal{C},\mathcal{O})$
上的微分分次代数层。同伦范畴
$K(\textit{Mod}(\mathcal{A},\text{d}))$ 中由无环模组成的
全子范畴 $\text{Ac}$，是
$K(\textit{Mod}(\mathcal{A},\text{d}))$
的严格全饱和三角子范畴。
\end{lemma}

\begin{proof}
当然，$K(\textit{Mod}(\mathcal{A},\text{d}))$ 的对象
$\mathcal{M}$ 属于 $\text{Ac}$，当且仅当对所有 $i$，
$H^i(\mathcal{M})=H^0(\mathcal{M}[i])$ 均为零。
本引理由此、引理 \ref{sdga-lemma-cohomology-homological}
及《导出范畴》引理 \ref{derived-lemma-homological-functor-kernel}
得出。另见《导出范畴》定义 \ref{derived-definition-saturated}、
\ref{derived-definition-triangulated-subcategory}
及引理 \ref{derived-lemma-triangulated-subcategory}。
\end{proof}

\begin{lemma}
\label{sdga-lemma-qis}
设 $(\mathcal{C},\mathcal{O})$ 是环化位点，
$(\mathcal{A},\text{d})$ 是 $(\mathcal{C},\mathcal{O})$
上的微分分次代数层。考虑由拟同构组成的子类
$\text{Qis} \subset \text{Arrows}(K(\textit{Mod}(\mathcal{A}, \text{d})))$
。这是一个饱和乘法系，并与
$K(\textit{Mod}(\mathcal{A},\text{d}))$ 上的三角结构相容。
\end{lemma}

\begin{proof}
注意，若 $f,g:\mathcal{M}\to\mathcal{N}$ 是
$\textit{Mod}(\mathcal{A},\text{d})$ 中彼此同伦的态射，
则 $f$ 是拟同构当且仅当 $g$ 是拟同构。事实上，由引理
\ref{sdga-lemma-cohomology-homological}，映射
$H^i(f)=H^0(f[i])$ 与 $H^i(g)=H^0(g[i])$ 相同。
因此，说同伦范畴 $K(\textit{Mod}(\mathcal{A},\text{d}))$
的某个态射是拟同构并无歧义。关于``乘法系''、``饱和''
及``与三角结构相容''的定义，参见《导出范畴》定义
\ref{derived-definition-localization}，以及《范畴》定义
\ref{categories-definition-multiplicative-system}、
\ref{categories-definition-saturated-multiplicative-system}。

\medskip\noindent
为实际证明本引理，考虑三角范畴间正合函子的复合
$$
K(\textit{Mod}(\mathcal{A}, \text{d}))
\longrightarrow
K(\textit{Mod}(\mathcal{O}))
\longrightarrow
D(\mathcal{O})
$$
。注意，$K(\textit{Mod}(\mathcal{A},\text{d}))$ 的态射
$f:\mathcal{M}\to\mathcal{N}$ 属于 $\text{Qis}$，
当且仅当它在 $D(\mathcal{O})$ 中映为同构。
故本引理由《导出范畴》引理
\ref{derived-lemma-triangle-functor-localize} 得出。
\end{proof}

\noindent
在引理 \ref{sdga-lemma-qis} 的情形下，可以应用《导出范畴》命题
\ref{derived-proposition-construct-localization}，
得到三角范畴间的正合函子
$$
Q :
K(\textit{Mod}(\mathcal{A}, \text{d}))
\longrightarrow
\text{Qis}^{-1}K(\textit{Mod}(\mathcal{A}, \text{d}))
$$
。然而，由于 $\textit{Mod}(\mathcal{A},\text{d})$ 是``大''
范畴，即其对象构成真类，给定 $\mathcal{M}$ 与 $\mathcal{N}$ 后，
$\text{Qis}^{-1}K(\textit{Mod}(\mathcal{A},\text{d}))$
的构造是否产生态射的一个{\it 集合}
$$
\Mor_{\text{Qis}^{-1}K(\textit{Mod}(\mathcal{A}, \text{d}))}
(\mathcal{M}, \mathcal{N})
$$
，并非立即清楚。不过，借助我们对 K-内射复形的构造，
答案是肯定的。事实上，由定理
\ref{sdga-theorem-qis-into-dg-injective}，可选取拟同构
$s_0:\mathcal{N}\to\mathcal{I}$，其中 $\mathcal{I}$
是既分次内射又 K-内射的微分分次 $\mathcal{A}$-模。
回忆，上述集合的元素是偶对
$(f : \mathcal{M} \to \mathcal{N}', s : \mathcal{N} \to \mathcal{N}')$
的等价类，其中 $f$ 是
$K(\textit{Mod}(\mathcal{A},\text{d}))$ 的任意态射，
$s$ 是拟同构；参见《范畴》第
\ref{categories-section-localization} 节对左分式演算的描述。
由引理 \ref{sdga-lemma-second-property-dg-injective}，可选取下图中的虚线箭头
$$
\xymatrix{
\mathcal{M} \ar[rd]^f & &
\mathcal{N} \ar[ld]_s \ar[rd]^{s_0} \\
& \mathcal{N}' \ar@{..>}[rr]^{s'} & & \mathcal{I}
}
$$
使图在同伦范畴中交换。因此，偶对 $(f,s)$ 等价于
$(s'\circ f,s_0)$，从而这些等价类构成集合。

\begin{definition}
\label{sdga-definition-derived-category}
设 $(\mathcal{C},\mathcal{O})$ 是环化位点，
$(\mathcal{A},\text{d})$ 是 $(\mathcal{C},\mathcal{O})$
上的微分分次代数层。令 $\text{Qis}$ 如引理
\ref{sdga-lemma-qis} 中所述。{\it $(\mathcal{A},\text{d})$ 的导出范畴}
是上文详细讨论的三角范畴
$$
D(\mathcal{A}, \text{d}) =
\text{Qis}^{-1}K(\textit{Mod}(\mathcal{A}, \text{d}))
$$
\end{definition}

\noindent
下面证明关于这一构造的一些事实。

\begin{lemma}
\label{sdga-lemma-kernel-localization}
在定义 \ref{sdga-definition-derived-category} 中，局部化函子
$Q : K(\textit{Mod}(\mathcal{A}, \text{d})) \to D(\mathcal{A}, \text{d})$
的核是引理 \ref{sdga-lemma-acyclics} 中的范畴 $\text{Ac}$。
\end{lemma}

\begin{proof}
这立即来自《导出范畴》引理
\ref{derived-lemma-kernel-localization}，以及如下事实：
$0\to\mathcal{M}$ 是拟同构，当且仅当 $\mathcal{M}$ 无环。
\end{proof}

\begin{lemma}
\label{sdga-lemma-H0-over-D}
在定义 \ref{sdga-definition-derived-category} 中，函子
$H^0 : K(\textit{Mod}(\mathcal{A}, \text{d})) \to
\textit{Mod}(\mathcal{O})$ 经同调函子
$H^0 : D(\mathcal{A}, \text{d}) \to \textit{Mod}(\mathcal{O})$.
分解。
\end{lemma}

\begin{proof}
立即来自《导出范畴》引理
\ref{derived-lemma-universal-property-localization}。
\end{proof}

\noindent
下面给出此前承诺的计算导出范畴中态射集的引理。

\begin{lemma}
\label{sdga-lemma-hom-derived}
设 $(\mathcal{C},\mathcal{O})$ 是环化位点，
$(\mathcal{A},\text{d})$ 是 $(\mathcal{C},\mathcal{O})$
上的微分分次代数层。设 $\mathcal{M}$ 与 $\mathcal{N}$
是微分分次 $\mathcal{A}$-模，并设
$\mathcal{N}\to\mathcal{I}$ 是拟同构，其中 $\mathcal{I}$
是既分次内射又 K-内射的微分分次 $\mathcal{A}$-模。则
$$
\Hom_{D(\mathcal{A}, \text{d})}(\mathcal{M}, \mathcal{N}) =
\Hom_{K(\textit{Mod}(\mathcal{A}, \text{d}))}(\mathcal{M}, \mathcal{I})
$$
\end{lemma}

\begin{proof}
由于 $\mathcal{N}\to\mathcal{I}$ 是拟同构，有
$$
\Hom_{D(\mathcal{A}, \text{d})}(\mathcal{M}, \mathcal{N}) =
\Hom_{D(\mathcal{A}, \text{d})}(\mathcal{M}, \mathcal{I})
$$
。在定义 \ref{sdga-definition-derived-category} 之前的讨论中，
利用引理 \ref{sdga-lemma-second-property-dg-injective} 已经得到：
$D(\mathcal{A},\text{d})$ 中任意态射
$\mathcal{M}\to\mathcal{I}$ 都可由
$K(\textit{Mod}(\mathcal{A},\text{d}))$ 中的态射
$f:\mathcal{M}\to\mathcal{I}$ 表示。现在，若
$f,f':\mathcal{M}\to\mathcal{I}$ 是
$K(\textit{Mod}(\mathcal{A},\text{d}))$ 中的两个态射，
则它们在 $D(\mathcal{A},\text{d})$ 中定义同一态射，
当且仅当存在 $K(\textit{Mod}(\mathcal{A},\text{d}))$
中的拟同构 $g:\mathcal{I}\to\mathcal{K}$，使
$g\circ f=g\circ f'$；参见《范畴》引理
\ref{categories-lemma-equality-morphisms-left-localization}。
然而，由引理 \ref{sdga-lemma-second-property-dg-injective}，
存在映射 $h:\mathcal{K}\to\mathcal{I}$，使得在
$K(\textit{Mod}(\mathcal{A},\text{d}))$ 中
$h\circ g=\text{id}_\mathcal{I}$。故
$g\circ f=g\circ f'$ 蕴含 $f=f'$，证明完成。
\end{proof}

\begin{lemma}
\label{sdga-lemma-derived-products}
设 $(\mathcal{C},\mathcal{O})$ 是环化位点，
$(\mathcal{A},\text{d})$ 是 $(\mathcal{C},\mathcal{O})$
上的微分分次代数层。则
\begin{enumerate}
\item $D(\mathcal{A},\text{d})$ 既有直和也有直积；
\item 直和通过取微分分次 $\mathcal{A}$-模的直和得到；
\item 直积通过取 K-内射微分分次模的直积得到。
\end{enumerate}
\end{lemma}

\begin{proof}
我们将使用如下事实：
$\textit{Mod}(\mathcal{A},\text{d})$ 是具有任意直和与直积的
阿贝尔范畴，而且这些运算在
$K(\textit{Mod}(\mathcal{A},\text{d}))$ 中给出直和与直积。
参见引理 \ref{sdga-lemma-dgm-abelian} 与
\ref{sdga-lemma-homotopy-direct-sums}。

\medskip\noindent
设 $\mathcal{M}_j$ 是一族微分分次 $\mathcal{A}$-模。
把直和 $\mathcal{M}=\bigoplus\mathcal{M}_j$
视为微分分次 $\mathcal{A}$-模。对微分分次
$\mathcal{A}$-模 $\mathcal{N}$，选取拟同构
$\mathcal{N}\to\mathcal{I}$，其中 $\mathcal{I}$
作为微分分次 $\mathcal{A}$-模既分次内射又 K-内射；
参见定理 \ref{sdga-theorem-qis-into-dg-injective}。
利用引理 \ref{sdga-lemma-hom-derived}，有
\begin{align*}
\Hom_{D(\mathcal{A}, \text{d})}(\mathcal{M}, \mathcal{N})
& =
\Hom_{K(\mathcal{A}, \text{d})}(\mathcal{M}, \mathcal{I}) \\
& =
\prod \Hom_{K(\mathcal{A}, \text{d})}(\mathcal{M}_j, \mathcal{I}) \\
& =
\prod \Hom_{D(\mathcal{A}, \text{d})}(\mathcal{M}_j, \mathcal{I})
\end{align*}
由此得到 $D(A,\text{d})$ 中按本引理 (2) 所述的直和。

\medskip\noindent
设 $\mathcal{M}_j$ 是一族微分分次 $\mathcal{A}$-模。
对每个 $j$，选取拟同构
$\mathcal{M}\to\mathcal{I}_j$，其中 $\mathcal{I}_j$
作为微分分次 $\mathcal{A}$-模既分次内射又 K-内射。
考虑微分分次 $\mathcal{A}$-模的直积
$\mathcal{I}=\prod\mathcal{I}_j$。由引理
\ref{sdga-lemma-product-K-injective} 与
\ref{sdga-lemma-product-graded-injective}，$\mathcal{I}$
作为微分分次 $\mathcal{A}$-模既分次内射又 K-内射。
对微分分次 $\mathcal{A}$-模 $\mathcal{N}$，
利用引理 \ref{sdga-lemma-hom-derived} 有
\begin{align*}
\Hom_{D(\mathcal{A}, \text{d})}(\mathcal{N}, \mathcal{I})
& =
\Hom_{K(\mathcal{A}, \text{d})}(\mathcal{N}, \mathcal{I}) \\
& =
\prod \Hom_{K(\mathcal{A}, \text{d})}(\mathcal{N}, \mathcal{I}_j) \\
& =
\prod \Hom_{D(\mathcal{A}, \text{d})}(\mathcal{N}, \mathcal{M}_j)
\end{align*}
由此得到 $D(\mathcal{A},\text{d})$ 中按本引理 (3) 所述的直积。
\end{proof}







\section{典范 $\delta$-函子}
\label{sdga-section-canonical-delta-functor}

\noindent
设 $(\mathcal{C},\mathcal{O})$ 是环化位点，
$(\mathcal{A},\text{d})$ 是 $(\mathcal{C},\mathcal{O})$
上的微分分次代数层。考虑函子
$\textit{Mod}(\mathcal{A}, \text{d}) \to
K(\textit{Mod}(\mathcal{A}, \text{d}))$.
一般而言，该函子{\bf 不是} $\delta$-函子。
然而，函子
$\textit{Mod}(\mathcal{A},\text{d})\to D(A,\text{d})$
却是 $\delta$-函子。为说明这一点，必须定义与阿贝尔范畴
$\textit{Mod}(\mathcal{A},\text{d})$ 中短正合序列
$$
0 \to \mathcal{K} \xrightarrow{a} \mathcal{L} \xrightarrow{b} \mathcal{M} \to 0
$$
相伴的态射 $\delta$。考虑态射 $a$ 的锥 $C(a)$，
连同其典范态射 $i:\mathcal{L}\to C(a)$ 与
$p:C(a)\to\mathcal{K}[1]$；参见定义 \ref{sdga-definition-cone}。
存在微分分次 $\mathcal{A}$-模同态
$$
q : C(a) \longrightarrow \mathcal{M}
$$
，它由《微分分次代数》引理 \ref{dga-lemma-cone}
应用于图
$$
\xymatrix{
\mathcal{K} \ar[r]_a \ar[d] &
\mathcal{L} \ar[d]^b \\
0 \ar[r] &
\mathcal{M}
}
$$
得到（由第 \ref{sdga-section-conclude-triangulated} 节的讨论，
我们可以使用该引理）。映射 $q$ 是拟同构；例如，这是因为
在 $\mathcal{O}$-模复形的态射范畴中该结论成立，参见
《导出范畴》第 \ref{derived-section-canonical-delta-functor} 节的讨论。
根据《微分分次代数》引理 \ref{dga-lemma-cone-homotopy}
（由第 \ref{sdga-section-conclude-triangulated} 节的讨论，
我们可以使用该引理），三角
$$
(\mathcal{K}, \mathcal{L}, C(a), a, i, -p)
$$
是 $K(\textit{Mod}(\mathcal{A},\text{d}))$ 中的杰出三角。
由于局部化函子
$K(\textit{Mod}(\mathcal{A},\text{d}))\to
D(\mathcal{A},\text{d})$ 正合，
$(\mathcal{K},\mathcal{L},C(a),a,i,-p)$ 也是
$D(\mathcal{A},\text{d})$ 中的杰出三角。
由于 $q$ 是拟同构，它在 $D(\mathcal{A},\text{d})$ 中是同构。
因此推出
$$
(\mathcal{K}, \mathcal{L}, \mathcal{M}, a, b, -p \circ q^{-1})
$$
是 $D(\mathcal{A},\text{d})$ 中的杰出三角。这提示了下列引理。

\begin{lemma}
\label{sdga-lemma-derived-canonical-delta-functor}
设 $(\mathcal{C},\mathcal{O})$ 是环化位点，
$(\mathcal{A},\text{d})$ 是 $(\mathcal{C},\mathcal{O})$
上的微分分次代数层。局部化函子
$\textit{Mod}(\mathcal{A}, \text{d}) \to D(\mathcal{A}, \text{d})$
具有 $\delta$-函子的自然结构，其中
$$
\delta_{\mathcal{K} \to \mathcal{L} \to \mathcal{M}} = - p \circ q^{-1}
$$
，$p$ 与 $q$ 如上所述。
\end{lemma}

\begin{proof}
我们已经看到，对任意复形短正合序列，这一选择给出杰出三角。
还需证明此构造的函子性；参见《导出范畴》定义
\ref{derived-definition-delta-functor}。
稍作推演，这可由《微分分次代数》引理 \ref{dga-lemma-cone}
得出（由第 \ref{sdga-section-conclude-triangulated} 节的讨论，
我们可以使用该引理）。可比较《导出范畴》引理
\ref{derived-lemma-derived-canonical-delta-functor}。
\end{proof}

\begin{lemma}
\label{sdga-lemma-homotopy-colimit}
设 $(\mathcal{C},\mathcal{O})$ 是环化位点，
$(\mathcal{A},\text{d})$ 是 $(\mathcal{C},\mathcal{O})$
上的微分分次代数层。设 $\mathcal{M}_n$
是微分分次 $\mathcal{A}$-模系统。则
$D(\mathcal{A},\text{d})$ 中的导出余极限
$\text{hocolim}\mathcal{M}_n$ 由微分分次模
$\colim\mathcal{M}_n$ 表示。
\end{lemma}

\begin{proof}
置 $\mathcal{M}=\colim\mathcal{M}_n$。由《导出范畴》引理
\ref{derived-lemma-compute-colimit}
（应用于底层 $\mathcal{O}$-模复形），有微分分次
$\mathcal{A}$-模的正合序列
$$
0 \to \bigoplus \mathcal{M}_n \to \bigoplus \mathcal{M}_n \to \mathcal{M} \to 0
$$
。由引理 \ref{sdga-lemma-derived-products}，这些直和也是
$D(\mathcal{A},\text{d})$ 中的直和。因此，结论来自
《导出范畴》定义 \ref{derived-definition-derived-colimit}
中的导出余极限定义，以及复形短正合序列给出杰出三角这一事实
（引理 \ref{sdga-lemma-derived-canonical-delta-functor}）。
\end{proof}






\section{导出拉回}
\label{sdga-section-derived-pullback}

\noindent
设 $(f, f^\sharp) : (\Sh(\mathcal{C}), \mathcal{O}_\mathcal{C})
\to (\Sh(\mathcal{D}), \mathcal{O}_\mathcal{D})$
是环化拓扑斯态射。设 $\mathcal{A}$ 是微分分次
$\mathcal{O}_\mathcal{C}$-代数，$\mathcal{B}$ 是微分分次
$\mathcal{O}_\mathcal{D}$-代数。设给定微分分次
$f^{-1}\mathcal{O}_\mathcal{D}$-代数映射
$$
\varphi : f^{-1}\mathcal{B} \to \mathcal{A}
$$
。由标量限制与标量扩张的伴随性，这等价于微分分次
$\mathcal{O}_\mathcal{C}$-代数映射
$\varphi:f^*\mathcal{B}\to\mathcal{A}$；等价地，
$\varphi$ 可视为微分分次 $\mathcal{O}_\mathcal{D}$-代数映射
$$
\varphi : \mathcal{B} \to f_*\mathcal{A}
$$
。参见评注 \ref{sdga-remark-functoriality-dga}。

\medskip\noindent
除上述数据外，再设 $\mathcal{A}'$ 是另一个微分分次
$\mathcal{O}_\mathcal{C}$-代数，$\mathcal{N}$ 是微分分次
$(\mathcal{A},\mathcal{A}')$-双模。在此情形下可考虑函子
$$
\textit{Mod}(\mathcal{B}, \text{d})
\longrightarrow
\textit{Mod}(\mathcal{A}', \text{d}),\quad
\mathcal{M} \longmapsto f^*\mathcal{M} \otimes_{\mathcal{A}} \mathcal{N}
$$
注意，由第 \ref{sdga-section-functoriality-dg} 与
\ref{sdga-section-dg-bimodules} 节的讨论，它延拓为微分分次范畴间的函子
$$
\textit{Mod}^{dg}(\mathcal{B}, \text{d})
\longrightarrow
\textit{Mod}^{dg}(\mathcal{A}', \text{d}),\quad
\mathcal{M} \longmapsto f^*\mathcal{M} \otimes_{\mathcal{A}} \mathcal{N}
$$
。形式地随之得到三角范畴间的正合函子
\begin{equation}
\label{sdga-equation-pullback}
K(\textit{Mod}(\mathcal{B}, \text{d}))
\longrightarrow
K(\textit{Mod}(\mathcal{A}', \text{d})),\quad
\mathcal{M} \longmapsto f^*\mathcal{M} \otimes_{\mathcal{A}} \mathcal{N}
\end{equation}
。

\begin{lemma}
\label{sdga-lemma-derived-tensor-product}
在上述情形下，函子 (\ref{sdga-equation-pullback}) 与局部化函子
$K(\textit{Mod}(\mathcal{A}', \text{d})) \to D(\mathcal{A}', \text{d})$
的复合有左导出延拓
$D(\mathcal{B},\text{d})\to D(\mathcal{A}',\text{d})$；
它在良好右微分分次 $\mathcal{B}$-模 $\mathcal{P}$ 上的值为
$f^*\mathcal{P}\otimes_\mathcal{A}\mathcal{N}$。
\end{lemma}

\begin{proof}
回忆，对任意（右）微分分次 $\mathcal{B}$-模 $\mathcal{M}$，
存在拟同构 $\mathcal{P}\to\mathcal{M}$，其中 $\mathcal{P}$
是良好微分分次 $\mathcal{B}$-模；参见引理 \ref{sdga-lemma-resolve}。
因此，由《导出范畴》引理
\ref{derived-lemma-find-existence-computes}，只需证明：
给定良好微分分次 $\mathcal{B}$-模间的拟同构
$\mathcal{P}\to\mathcal{P}'$，诱导映射
$$
f^*\mathcal{P} \otimes_\mathcal{A} \mathcal{N}
\longrightarrow
f^*\mathcal{P}' \otimes_\mathcal{A} \mathcal{N}
$$
是拟同构。由引理 \ref{sdga-lemma-good-admissible-ses}，
$\mathcal{P}\to\mathcal{P}'$ 的锥 $\mathcal{P}''$
是良好微分分次 $\mathcal{A}$-模。由于在
$K(\textit{Mod}(\mathcal{B},\text{d}))$ 中有杰出三角
$$
\mathcal{P} \to \mathcal{P}' \to \mathcal{P}'' \to \mathcal{P}[1]
$$
，在 $K(\textit{Mod}(\mathcal{A}',\text{d}))$ 中得到杰出三角
$$
f^*\mathcal{P} \otimes_\mathcal{A} \mathcal{N} \to
f^*\mathcal{P}' \otimes_\mathcal{A} \mathcal{N} \to
f^*\mathcal{P}'' \otimes_\mathcal{A} \mathcal{N} \to
f^*\mathcal{P}[1] \otimes_\mathcal{A} \mathcal{N}
$$
。由引理 \ref{sdga-lemma-acyclic-good}，微分分次模
$f^*\mathcal{P}'' \otimes_\mathcal{A} \mathcal{N}$
无环，证明完成。
\end{proof}

\begin{definition}
\label{sdga-definition-pullback}
导出张量积与导出拉回。
\begin{enumerate}
\item 设 $(\mathcal{C},\mathcal{O})$ 是环化位点，
$\mathcal{A},\mathcal{B}$ 是微分分次 $\mathcal{O}$-代数，
$\mathcal{N}$ 是微分分次 $(\mathcal{A},\mathcal{B})$-双模。
引理 \ref{sdga-lemma-derived-tensor-product} 中构造的函子
$D(\mathcal{A},\text{d})\to D(\mathcal{B},\text{d})$
称为{\it 导出张量积}，记作
$-\otimes_\mathcal{A}^\mathbf{L}\mathcal{N}$。
\item 设 $(f, f^\sharp) : (\Sh(\mathcal{C}), \mathcal{O}_\mathcal{C})
\to (\Sh(\mathcal{D}), \mathcal{O}_\mathcal{D})$
是环化拓扑斯态射，$\mathcal{A}$ 是微分分次
$\mathcal{O}_\mathcal{C}$-代数，$\mathcal{B}$ 是微分分次
$\mathcal{O}_\mathcal{D}$-代数，并设
$\varphi:\mathcal{B}\to f_*\mathcal{A}$ 是微分分次
$\mathcal{O}_\mathcal{D}$-代数同态。引理
\ref{sdga-lemma-derived-tensor-product} 中构造的函子
$D(\mathcal{B},\text{d})\to D(\mathcal{A},\text{d})$
称为{\it 导出拉回}，记作 $Lf^*$。
\end{enumerate}
\end{definition}

\noindent
有了这些术语，可以表述一些显然的相容性。

\begin{lemma}
\label{sdga-lemma-compose-pullback-tensor}
引理 \ref{sdga-lemma-derived-tensor-product} 中的函子
$D(\mathcal{B},\text{d})\to D(\mathcal{A}',\text{d})$ 等于
$\mathcal{M} \mapsto
Lf^*\mathcal{M} \otimes_\mathcal{A}^\mathbf{L} \mathcal{N}$.
\end{lemma}

\begin{proof}
这是立即的，因为可用良好微分分次模表示对象来计算这些函子，
并且当 $\mathcal{P}$ 是良好微分分次 $\mathcal{B}$-模时，
$f^*\mathcal{P}$ 是良好微分分次 $\mathcal{A}$-模。
\end{proof}

\begin{lemma}
\label{sdga-lemma-compose-pullback}
设 $(f, f^\sharp) : (\Sh(\mathcal{C}), \mathcal{O})
\to (\Sh(\mathcal{C}'), \mathcal{O}')$ 与
$(g, g^\sharp) : (\Sh(\mathcal{C}'), \mathcal{O}')
\to (\Sh(\mathcal{C}''), \mathcal{O}'')$
是环化拓扑斯态射。设 $\mathcal{A},\mathcal{A}',\mathcal{A}''$
分别是微分分次 $\mathcal{O}$-代数、
$\mathcal{O}'$-代数、$\mathcal{O}''$-代数。设
$\varphi : \mathcal{A}' \to f_*\mathcal{A}$ 与
$\varphi' : \mathcal{A}'' \to g_*\mathcal{A}'$
分别是微分分次 $\mathcal{O}'$-代数同态与
$\mathcal{O}''$-代数同态。则
$L(g \circ f)^* = Lf^* \circ Lg^* :
D(\mathcal{A}'', \text{d}) \to D(\mathcal{A}, \text{d})$.
\end{lemma}

\begin{proof}
这是立即的，因为可用良好微分分次模表示对象来计算这些函子，
并且当 $\mathcal{P}$ 是良好微分分次 $\mathcal{A}$-模时，
$f^*\mathcal{P}$ 是良好微分分次 $\mathcal{A}'$-模。
\end{proof}

\noindent
设 $(\mathcal{C},\mathcal{O})$ 是环化位点，
$\mathcal{A},\mathcal{B}$ 是微分分次 $\mathcal{O}$-代数。
设 $\mathcal{N}\to\mathcal{N}'$ 是微分分次
$(\mathcal{A},\mathcal{B})$-双模同态。于是得到典范映射
$$
t :
\mathcal{M} \otimes_\mathcal{A}^\mathbf{L} \mathcal{N}
\longrightarrow
\mathcal{M} \otimes_\mathcal{A}^\mathbf{L} \mathcal{N}'
$$
，它关于 $D(\mathcal{A},\text{d})$ 中的 $\mathcal{M}$ 具有函子性，
并定义三角范畴间正合函子
$D(\mathcal{A},\text{d})\to D(\mathcal{B},\text{d})$
之间的自然变换。$t$ 在良好微分分次
$\mathcal{A}$-模 $\mathcal{P}$ 上的值是显然的映射
$$
\mathcal{P} \otimes_\mathcal{A}^\mathbf{L} \mathcal{N} =
\mathcal{P} \otimes_\mathcal{A} \mathcal{N}
\longrightarrow
\mathcal{P} \otimes_\mathcal{A} \mathcal{N}' =
\mathcal{P} \otimes_\mathcal{A}^\mathbf{L} \mathcal{N}'
$$

\begin{lemma}
\label{sdga-lemma-tensor-symmetry}
在上述情形下，若 $\mathcal{N}\to\mathcal{N}'$
在上同调层上是同构，则 $t$ 是函子之间的同构
$(- \otimes_\mathcal{A}^\mathbf{L} \mathcal{N}) \to
(- \otimes_\mathcal{A}^\mathbf{L} \mathcal{N}')$.
\end{lemma}

\begin{proof}
只需证明，对任意良好微分分次 $\mathcal{A}$-模
$\mathcal{P}$，映射
$\mathcal{P} \otimes_\mathcal{A} \mathcal{N} \to
\mathcal{P} \otimes_\mathcal{A} \mathcal{N}'$
在上同调层上是同构。为此，令 $\mathcal{N}''$
为把 $\mathcal{N}\to\mathcal{N}'$ 视为微分分次左
$\mathcal{A}$-模映射时的锥；参见定义 \ref{sdga-definition-cone}。
（当然，$\mathcal{N}''$ 也是双模，但我们不需要这一点。）
由张量积构造的函子性（它是微分分次范畴间的函子），
$\mathcal{P}\otimes_\mathcal{A}\mathcal{N}''$
作为 $\mathcal{O}$-模复形，是映射
$\mathcal{P} \otimes_\mathcal{A} \mathcal{N} \to
\mathcal{P} \otimes_\mathcal{A} \mathcal{N}'$.
的锥。因此只需证明
$\mathcal{P} \otimes_\mathcal{A} \mathcal{N}''$
无环。这是因为 $\mathcal{P}$ 良好，而 $\mathcal{N}''$
作为拟同构的锥是无环的。
\end{proof}

\begin{lemma}
\label{sdga-lemma-good-on-other-side}
设 $(\mathcal{C},\mathcal{O})$ 是环化位点，
$\mathcal{A},\mathcal{B}$ 是微分分次 $\mathcal{O}$-代数，
$\mathcal{N}$ 是微分分次 $(\mathcal{A},\mathcal{B})$-双模。
若 $\mathcal{N}$ 作为微分分次左 $\mathcal{A}$-模良好，
则对所有微分分次 $\mathcal{A}$-模 $\mathcal{M}$，有
$\mathcal{M}\otimes_\mathcal{A}^\mathbf{L}\mathcal{N}
=\mathcal{M}\otimes_\mathcal{A}\mathcal{N}$。
\end{lemma}

\begin{proof}
设 $\mathcal{P}\to\mathcal{M}$ 是拟同构，其中 $\mathcal{P}$
是良好（右）微分分次 $\mathcal{A}$-模。为证明本引理，
需要说明
$\mathcal{P} \otimes_\mathcal{A} \mathcal{N} \to
\mathcal{M} \otimes_\mathcal{A} \mathcal{N}$
是拟同构。映射 $\mathcal{P}\to\mathcal{M}$ 的锥 $C$
是无环右微分分次 $\mathcal{A}$-模。由于假设
$\mathcal{N}$ 作为微分分次左 $\mathcal{A}$-模良好，
$C\otimes_\mathcal{A}\mathcal{N}$ 无环。而作为
$\mathcal{O}$-模复形，$C\otimes_\mathcal{A}\mathcal{N}$
是映射
$\mathcal{P}\otimes_\mathcal{A}\mathcal{N}\to
\mathcal{M}\otimes_\mathcal{A}\mathcal{N}$ 的锥，故结论成立。
\end{proof}

\begin{lemma}
\label{sdga-lemma-compose-tensor}
设 $(\mathcal{C},\mathcal{O})$ 是环化位点，
$\mathcal{A},\mathcal{A}',\mathcal{A}''$ 是微分分次
$\mathcal{O}$-代数。设 $\mathcal{N}$ 与 $\mathcal{N}'$
分别是微分分次 $(\mathcal{A},\mathcal{A}')$-双模和
$(\mathcal{A}',\mathcal{A}'')$-双模。假设典范映射
$$
\mathcal{N} \otimes_{\mathcal{A}'}^\mathbf{L} \mathcal{N}'
\longrightarrow
\mathcal{N} \otimes_{\mathcal{A}'} \mathcal{N}'
$$
在 $D(\mathcal{A}'',\text{d})$ 中是拟同构。则
$$
(\mathcal{M}
\otimes_\mathcal{A}^\mathbf{L} \mathcal{N})
\otimes_{\mathcal{A}'}^\mathbf{L} \mathcal{N}'
=
\mathcal{M}
\otimes_\mathcal{A}^\mathbf{L}
(\mathcal{N} \otimes_{\mathcal{A}'} \mathcal{N}')
$$
是函子
$D(\mathcal{A},\text{d})\to D(\mathcal{A}'',\text{d})$ 的等式。
\end{lemma}

\begin{proof}
选取良好微分分次 $\mathcal{A}$-模 $\mathcal{P}$
以及拟同构 $\mathcal{P}\to\mathcal{M}$；参见引理
\ref{sdga-lemma-resolve}。则
$$
\mathcal{M}
\otimes_\mathcal{A}^\mathbf{L}
(\mathcal{N} \otimes_{\mathcal{A}'} \mathcal{N}') =
\mathcal{P} \otimes_\mathcal{A} \mathcal{N}
\otimes_{\mathcal{A}'} \mathcal{N}'
$$
并且
$$
(\mathcal{M}
\otimes_\mathcal{A}^\mathbf{L} \mathcal{N})
\otimes_{\mathcal{A}'}^\mathbf{L} \mathcal{N}' =
(\mathcal{P} \otimes_\mathcal{A} \mathcal{N})
\otimes_{\mathcal{A}'}^\mathbf{L} \mathcal{N}'
$$
因此需要证明典范映射
$$
(\mathcal{P} \otimes_\mathcal{A} \mathcal{N})
\otimes_{\mathcal{A}'}^\mathbf{L} \mathcal{N}'
\longrightarrow
\mathcal{P} \otimes_\mathcal{A} \mathcal{N}
\otimes_{\mathcal{A}'} \mathcal{N}'
$$
是拟同构。选取拟同构 $\mathcal{Q}\to\mathcal{N}'$，
其中 $\mathcal{Q}$ 是良好左微分分次 $\mathcal{A}'$-模
（引理 \ref{sdga-lemma-resolve}）。由引理
\ref{sdga-lemma-good-on-other-side}，上面的映射作为
$\mathcal{O}$-模导出范畴中的映射是
$$
\mathcal{P} \otimes_\mathcal{A} \mathcal{N}
\otimes_{\mathcal{A}'} \mathcal{Q}
\longrightarrow
\mathcal{P} \otimes_\mathcal{A} \mathcal{N}
\otimes_{\mathcal{A}'} \mathcal{N}'
$$
。依假设，
$\mathcal{N}\otimes_{\mathcal{A}'}\mathcal{Q}\to
\mathcal{N}\otimes_{\mathcal{A}'}\mathcal{N}'$ 是拟同构；
又因 $\mathcal{P}$ 是良好微分分次 $\mathcal{A}$-模，
由引理 \ref{sdga-lemma-tensor-symmetry}，此映射是拟同构
（左端与右端分别计算
$\mathcal{P}\otimes_\mathcal{A}^\mathbf{L}
(\mathcal{N}\otimes_{\mathcal{A}'}\mathcal{Q})$ 和
$\mathcal{P}\otimes_\mathcal{A}^\mathbf{L}
(\mathcal{N}\otimes_{\mathcal{A}'}\mathcal{N}')$；
也可以直接重复该引理证明中的论证）。
\end{proof}






\section{导出直像}
\label{sdga-section-derived-pushforward}

\noindent
足够多 K-内射对象的存在性保证：可以对同伦范畴上的任意正合函子
取右导出函子。

\begin{lemma}
\label{sdga-lemma-right-derived}
设 $(\mathcal{C},\mathcal{O})$ 是环化位点，
$(\mathcal{A},\text{d})$ 是 $(\mathcal{C},\mathcal{O})$
上的微分分次代数层。则三角范畴间的任意正合函子
$$
T : K(\textit{Mod}(\mathcal{A}, \text{d})) \longrightarrow \mathcal{D}
$$
都有右导出延拓
$RT:D(\mathcal{A},\text{d})\to\mathcal{D}$；
它在既分次内射又 K-内射的微分分次
$\mathcal{A}$-模 $\mathcal{I}$ 上的值为 $T(\mathcal{I})$。
\end{lemma}

\begin{proof}
由定理 \ref{sdga-theorem-qis-into-dg-injective}，对任意（右）微分分次
$\mathcal{A}$-模 $\mathcal{M}$，存在拟同构
$\mathcal{M}\to\mathcal{I}$，其中 $\mathcal{I}$
是既分次内射又 K-内射的微分分次 $\mathcal{A}$-模。
因此，由《导出范畴》引理
\ref{derived-lemma-find-existence-computes}，只需证明：
若 $\mathcal{I}\to\mathcal{I}'$ 是两个既分次内射又 K-内射的
微分分次 $\mathcal{A}$-模之间的拟同构，则
$T(\mathcal{I})\to T(\mathcal{I}')$ 是同构。这确实成立，
因为由引理 \ref{sdga-lemma-hom-derived} 可知
$\mathcal{I}\to\mathcal{I}'$ 在
$K(\textit{Mod}(\mathcal{A},\text{d}))$ 中是同构
（也可由引理 \ref{sdga-lemma-second-property-dg-injective} 推出）。
\end{proof}

\noindent
我们已经见过若干可应用此结果的函子。下面给出两个例子。

\begin{definition}
\label{sdga-definition-pushforward}
导出内 Hom 与导出直像。
\begin{enumerate}
\item 设 $(\mathcal{C},\mathcal{O})$ 是环化位点，
$\mathcal{A},\mathcal{B}$ 是微分分次 $\mathcal{O}$-代数，
$\mathcal{N}$ 是微分分次 $(\mathcal{A},\mathcal{B})$-双模。
内 Hom 函子 $\SheafHom_\mathcal{B}^{dg}(\mathcal{N},-)$
的右导出延拓
$$
R\SheafHom_\mathcal{B}(\mathcal{N}, -) :
D(\mathcal{B}, \text{d})
\longrightarrow
D(\mathcal{A}, \text{d})
$$
称为{\it 导出内 Hom}。
\item 设 $(f, f^\sharp) : (\Sh(\mathcal{C}), \mathcal{O}_\mathcal{C})
\to (\Sh(\mathcal{D}), \mathcal{O}_\mathcal{D})$
是环化拓扑斯态射，$\mathcal{A}$ 是微分分次
$\mathcal{O}_\mathcal{C}$-代数，$\mathcal{B}$ 是微分分次
$\mathcal{O}_\mathcal{D}$-代数，并设
$\varphi:\mathcal{B}\to f_*\mathcal{A}$ 是微分分次
$\mathcal{O}_\mathcal{D}$-代数同态。直像 $f_*$ 的右导出延拓
$$
Rf_* :
D(\mathcal{A}, \text{d})
\longrightarrow
D(\mathcal{B}, \text{d})
$$
称为{\it 导出直像}。
\end{enumerate}
\end{definition}

\noindent
事实证明，
$Rf_* : D(\mathcal{A}, \text{d}) \to D(\mathcal{B}, \text{d})$
与底层 $\mathcal{O}$-模复形上的导出直像一致；参见引理
\ref{sdga-lemma-pushforward-agrees}。

\medskip\noindent
这些函子分别是导出拉回与导出张量积的伴随。

\begin{lemma}
\label{sdga-lemma-derived-adjoint-tensor-hom}
设 $(\mathcal{C},\mathcal{O})$ 是环化位点，
$\mathcal{A},\mathcal{B}$ 是微分分次 $\mathcal{O}$-代数，
$\mathcal{N}$ 是微分分次 $(\mathcal{A},\mathcal{B})$-双模。则
$$
R\SheafHom_\mathcal{B}(\mathcal{N}, -) :
D(\mathcal{B}, \text{d})
\longrightarrow
D(\mathcal{A}, \text{d})
$$
是
$$
- \otimes_\mathcal{A}^\mathbf{L} \mathcal{N} :
D(\mathcal{A}, \text{d})
\longrightarrow
D(\mathcal{B}, \text{d})
$$
的右伴随。
\end{lemma}

\begin{proof}
这由《导出范畴》引理
\ref{derived-lemma-pre-derived-adjoint-functors-general}
和引理 \ref{sdga-lemma-tensor-hom-adjunction-dg} 得出。
\end{proof}

\begin{lemma}
\label{sdga-lemma-derived-adjoint-push-pull}
设 $(f, f^\sharp) : (\Sh(\mathcal{C}), \mathcal{O}_\mathcal{C})
\to (\Sh(\mathcal{D}), \mathcal{O}_\mathcal{D})$
是环化拓扑斯态射，$\mathcal{A}$ 是微分分次
$\mathcal{O}_\mathcal{C}$-代数，$\mathcal{B}$ 是微分分次
$\mathcal{O}_\mathcal{D}$-代数，并设
$\varphi:\mathcal{B}\to f_*\mathcal{A}$ 是微分分次
$\mathcal{O}_\mathcal{D}$-代数同态。则
$$
Rf_* : 
D(\mathcal{A}, \text{d})
\longrightarrow
D(\mathcal{B}, \text{d})
$$
是
$$
Lf^* :
D(\mathcal{B}, \text{d})
\longrightarrow
D(\mathcal{A}, \text{d})
$$
的右伴随。
\end{lemma}

\begin{proof}
这由《导出范畴》引理
\ref{derived-lemma-pre-derived-adjoint-functors-general}
和引理 \ref{sdga-lemma-adjunction-push-pull-dg} 得出。
\end{proof}

\noindent
下面讨论第 \ref{sdga-section-derived-pullback} 节所考虑的情形。

\medskip\noindent
设 $(f, f^\sharp) : (\Sh(\mathcal{C}), \mathcal{O}_\mathcal{C})
\to (\Sh(\mathcal{D}), \mathcal{O}_\mathcal{D})$
是环化拓扑斯态射。设 $\mathcal{A}$ 是微分分次
$\mathcal{O}_\mathcal{C}$-代数，$\mathcal{B}$ 是微分分次
$\mathcal{O}_\mathcal{D}$-代数。设给定微分分次
$f^{-1}\mathcal{O}_\mathcal{D}$-代数映射
$$
\varphi : f^{-1}\mathcal{B} \to \mathcal{A}
$$
。由标量限制与标量扩张的伴随性，这等价于微分分次
$\mathcal{O}_\mathcal{C}$-代数映射
$\varphi:f^*\mathcal{B}\to\mathcal{A}$；等价地，
$\varphi$ 可视为微分分次 $\mathcal{O}_\mathcal{D}$-代数映射
$$
\varphi : \mathcal{B} \to f_*\mathcal{A}
$$
。参见评注 \ref{sdga-remark-functoriality-dga}。

\medskip\noindent
除上述数据外，再设 $\mathcal{A}'$ 是另一个微分分次
$\mathcal{O}_\mathcal{C}$-代数，$\mathcal{N}$ 是微分分次
$(\mathcal{A},\mathcal{A}')$-双模。在此情形下可考虑函子
$$
\textit{Mod}(\mathcal{A}', \text{d})
\longrightarrow
\textit{Mod}(\mathcal{B}, \text{d}),\quad
\mathcal{M} \longmapsto
f_*\SheafHom^{dg}_{\mathcal{A}'}(\mathcal{N}, \mathcal{M})
$$
注意，由第 \ref{sdga-section-functoriality-dg} 与
\ref{sdga-section-dg-bimodules} 节的讨论，它延拓为微分分次范畴间的函子
$$
\textit{Mod}^{dg}(\mathcal{A}', \text{d})
\longrightarrow
\textit{Mod}^{dg}(\mathcal{B}, \text{d}),\quad
\mathcal{M} \longmapsto
f_*\SheafHom^{dg}_{\mathcal{A}'}(\mathcal{N}, \mathcal{M})
$$
。形式地随之得到三角范畴间的正合函子
\begin{equation}
\label{sdga-equation-pushforward}
K(\textit{Mod}(\mathcal{A}', \text{d}))
\longrightarrow
K(\textit{Mod}(\mathcal{B}, \text{d})),\quad
\mathcal{M} \longmapsto
f_*\SheafHom^{dg}_{\mathcal{A}'}(\mathcal{N}, \mathcal{M})
\end{equation}
。

\begin{lemma}
\label{sdga-lemma-compose-pushforward-hom}
在上述情形下，以
$RT:D(\mathcal{A}',\text{d})\to D(\mathcal{B},\text{d})$
表示 (\ref{sdga-equation-pushforward}) 的右导出延拓。则
$$
RT(\mathcal{M}) = Rf_* R\SheafHom(\mathcal{N}, \mathcal{M})
$$
关于 $\mathcal{M}$ 函子性地成立。
\end{lemma}

\begin{proof}
由引理 \ref{sdga-lemma-tensor-hom-adjunction-dg} 与
\ref{sdga-lemma-adjunction-push-pull-dg}，函子
(\ref{sdga-equation-pushforward}) 是函子
(\ref{sdga-equation-pullback}) 的右伴随。由《导出范畴》引理
\ref{derived-lemma-pre-derived-adjoint-functors-general}，
函子 $RT$ 是引理 \ref{sdga-lemma-derived-tensor-product}
中函子的右伴随；由引理 \ref{sdga-lemma-compose-pullback-tensor}，
后者等于
$Lf^*(-)\otimes_\mathcal{A}^\mathbf{L}\mathcal{N}$。
由引理 \ref{sdga-lemma-derived-adjoint-tensor-hom} 与
\ref{sdga-lemma-derived-adjoint-push-pull}，函子
$Lf^*(-)\otimes_\mathcal{A}^\mathbf{L}\mathcal{N}$ 的右伴随为
$Rf_* R\SheafHom(\mathcal{N}, -)$
。故由伴随的唯一性得出结论。
\end{proof}

\begin{lemma}
\label{sdga-lemma-compose-pushforward}
设 $(f, f^\sharp) : (\Sh(\mathcal{C}), \mathcal{O})
\to (\Sh(\mathcal{C}'), \mathcal{O}')$ 与
$(g, g^\sharp) : (\Sh(\mathcal{C}'), \mathcal{O}')
\to (\Sh(\mathcal{C}''), \mathcal{O}'')$
是环化拓扑斯态射。设 $\mathcal{A},\mathcal{A}',\mathcal{A}''$
分别是微分分次 $\mathcal{O}$-代数、
$\mathcal{O}'$-代数、$\mathcal{O}''$-代数，并设
$\varphi : \mathcal{A}' \to f_*\mathcal{A}$ 与
$\varphi' : \mathcal{A}'' \to g_*\mathcal{A}'$
分别是微分分次 $\mathcal{O}'$-代数同态与
$\mathcal{O}''$-代数同态。则
$R(g \circ f)_* = Rg_* \circ Rf_* :
D(\mathcal{A}, \text{d}) \to D(\mathcal{A}'', \text{d})$.
\end{lemma}

\begin{proof}
由引理 \ref{sdga-lemma-compose-pullback}、
\ref{sdga-lemma-derived-adjoint-push-pull} 及伴随的唯一性得出。
\end{proof}

\begin{lemma}
\label{sdga-lemma-compose-hom}
设 $(\mathcal{C},\mathcal{O})$ 是环化位点，
$\mathcal{A},\mathcal{A}',\mathcal{A}''$ 是微分分次
$\mathcal{O}$-代数。设 $\mathcal{N}$ 与 $\mathcal{N}'$
分别是微分分次 $(\mathcal{A},\mathcal{A}')$-双模和
$(\mathcal{A}',\mathcal{A}'')$-双模。假设典范映射
$$
\mathcal{N} \otimes_{\mathcal{A}'}^\mathbf{L} \mathcal{N}'
\longrightarrow
\mathcal{N} \otimes_{\mathcal{A}'} \mathcal{N}'
$$
在 $D(\mathcal{A}'',\text{d})$ 中是拟同构。则
$$
R\SheafHom_{\mathcal{A}''}
(\mathcal{N} \otimes_{\mathcal{A}'} \mathcal{N}', -)
=
R\SheafHom_{\mathcal{A}'}(\mathcal{N},
R\SheafHom_{\mathcal{A}''}(\mathcal{N}', -))
$$
是函子
$D(\mathcal{A}'',\text{d})\to D(\mathcal{A},\text{d})$ 的等式。
\end{lemma}

\begin{proof}
由引理 \ref{sdga-lemma-compose-tensor}、
\ref{sdga-lemma-derived-adjoint-tensor-hom} 及伴随的唯一性得出。
\end{proof}

\begin{lemma}
\label{sdga-lemma-pushforward-agrees}
设 $(f, f^\sharp) : (\Sh(\mathcal{C}), \mathcal{O}_\mathcal{C})
\to (\Sh(\mathcal{D}), \mathcal{O}_\mathcal{D})$
是环化拓扑斯态射。设 $\mathcal{A}$ 是微分分次
$\mathcal{O}_\mathcal{C}$-代数，$\mathcal{B}$ 是微分分次
$\mathcal{O}_\mathcal{D}$-代数，并设
$\varphi:\mathcal{B}\to f_*\mathcal{A}$ 是微分分次
$\mathcal{O}_\mathcal{D}$-代数同态。下图交换：
$$
\xymatrix{
D(\mathcal{A}, \text{d}) \ar[d]_{Rf_*} \ar[rr]_{forget} & &
D(\mathcal{O}_\mathcal{C}) \ar[d]^{Rf_*} \\
D(\mathcal{B}, \text{d}) \ar[rr]^{forget} & &
D(\mathcal{O}_\mathcal{D})
}
$$
\end{lemma}

\begin{proof}
除去若干范畴的等同，本引理由引理
\ref{sdga-lemma-compose-pushforward} 立即得出。

\medskip\noindent
把 $\mathcal{O}_\mathcal{C}$ 置于次数 $0$ 并赋予零微分，
便可将 $(\mathcal{O}_\mathcal{C},0)$ 视为微分分次
$\mathcal{O}_\mathcal{C}$-代数。显然有
$$
\textit{Mod}(\mathcal{O}_\mathcal{C}, 0) =
\text{Comp}(\mathcal{O}_\mathcal{C})
\quad\text{且}\quad
D(\mathcal{O}_\mathcal{C}, 0) = D(\mathcal{O}_\mathcal{C})
$$
在此等同下，遗忘函子
$\textit{Mod}(\mathcal{A}, \text{d}) \to
\text{Comp}(\mathcal{O}_\mathcal{C})$
是第 \ref{sdga-section-functoriality-dg} 节定义的``直像''
$\text{id}_{\mathcal{C},*}$；它对应于环化拓扑斯的恒等态射
$\text{id}_\mathcal{C} : (\mathcal{C}, \mathcal{O}_\mathcal{C}) \to
(\mathcal{C}, \mathcal{O}_\mathcal{C})$，以及微分分次
$\mathcal{O}_\mathcal{C}$-代数映射
$(\mathcal{O}_\mathcal{C},0)\to(\mathcal{A},\text{d})$。
由于 $\text{id}_{\mathcal{C},*}$ 正合，立即可知
$$
R\text{id}_{\mathcal{C}, *} = forget :
D(\mathcal{A}, \text{d}) \longrightarrow
D(\mathcal{O}_\mathcal{C}, 0) = D(\mathcal{O}_\mathcal{C})
$$
完全相同的论证表明
$$
R\text{id}_{\mathcal{D}, *} = forget :
D(\mathcal{B}, \text{d}) \longrightarrow
D(\mathcal{O}_\mathcal{D}, 0) = D(\mathcal{O}_\mathcal{D})
$$
此外，《位点上的上同调》第
\ref{sites-cohomology-section-unbounded} 节中
$Rf_* : D(\mathcal{O}_\mathcal{C}) \to D(\mathcal{O}_\mathcal{D})$
的构造，与定义 \ref{sdga-definition-pushforward} 中
$Rf_* : D(\mathcal{O}_\mathcal{C}, 0) \to D(\mathcal{O}_\mathcal{D}, 0)$
的构造一致，因为两个函子都定义为底层模复形之直像的右导出延拓。
由引理 \ref{sdga-lemma-compose-pushforward}，
$Rf_*\circ R\text{id}_{\mathcal{C},*}$ 与
$R\text{id}_{\mathcal{D},*}\circ Rf_*$ 都是
$f_*\circ forget=forget\circ f_*$ 的导出函子，
故由伴随的唯一性二者相等。
\end{proof}

\begin{lemma}
\label{sdga-lemma-cohomology-ext}
设 $(\mathcal{C},\mathcal{O})$ 是环化位点，
$\mathcal{A}$ 是微分分次 $\mathcal{O}$-代数，
$\mathcal{M}$ 是微分分次 $\mathcal{A}$-模，且
$n\in\mathbf{Z}$。则
$$
H^n(\mathcal{C}, \mathcal{M}) =
\Hom_{D(\mathcal{A}, \text{d})}(\mathcal{A}, \mathcal{M}[n])
$$
，其中左端是把 $\mathcal{M}$ 视为 $\mathcal{O}$-模复形所得的上同调。
\end{lemma}

\begin{proof}
为证明该公式，注意
$$
R\Gamma(\mathcal{C}, \mathcal{M}) = \Gamma(\mathcal{C}, \mathcal{I})
$$
，其中 $\mathcal{M}\to\mathcal{I}$ 是到既分次内射又 K-内射的
微分分次 $\mathcal{A}$-模 $\mathcal{I}$ 的拟同构
（合用引理 \ref{sdga-lemma-right-derived} 与
\ref{sdga-lemma-pushforward-agrees}）。由引理
\ref{sdga-lemma-hom-derived}，有
$$
\Hom_{D(\mathcal{A}, \text{d})}(\mathcal{A}, \mathcal{M}[n]) =
\Hom_{K(\textit{Mod}(\mathcal{A}, \text{d}))}(\mathcal{M}, \mathcal{I}[n]) =
H^0(\Gamma(\mathcal{C}, \mathcal{I}[n])) =
H^n(\Gamma(\mathcal{C}, \mathcal{I}))
$$
合并这两个结果即得所需等式。
\end{proof}








\section{导出范畴的等价}
\label{sdga-section-equivalence}

\noindent
本节是《微分分次代数》第 \ref{dga-section-equivalence} 节的类似物。

\begin{lemma}
\label{sdga-lemma-qis-equivalence}
设 $(\mathcal{C},\mathcal{O})$ 是环化位点。若
$\varphi:\mathcal{A}\to\mathcal{B}$ 是微分分次
$\mathcal{O}$-代数的同态，并在上同调层上诱导同构，则
$$
D(\mathcal{A}, \text{d}) \longrightarrow D(\mathcal{B}, \text{d}), \quad
\mathcal{M}
\longmapsto
\mathcal{M} \otimes_\mathcal{A}^\mathbf{L} \mathcal{B}
$$
是范畴等价。
\end{lemma}

\begin{proof}
回忆限制函子
$$
\textit{Mod}^{dg}(\mathcal{B}, \text{d}) \to
\textit{Mod}^{dg}(\mathcal{A}, \text{d}),\quad
\mathcal{N} \mapsto res_\varphi \mathcal{N}
$$
是函子
$$
\textit{Mod}^{dg}(\mathcal{A}, \text{d}) \to
\textit{Mod}^{dg}(\mathcal{B}, \text{d}),\quad
\mathcal{M} \mapsto \mathcal{M} \otimes_\mathcal{A} \mathcal{B}
$$
的右伴随；见第 \ref{sdga-section-dg-bimodules} 节。由于限制把拟同构
映为拟同构，故它显然有左导出延拓（仍由限制给出）。由
《导出范畴》引理
\ref{derived-lemma-pre-derived-adjoint-functors-general}，此函子是
$- \otimes_\mathcal{A}^\mathbf{L} \mathcal{B}$ 的右伴随。伴随态射
$$
\mathcal{M} \to
res_\varphi(\mathcal{M} \otimes_\mathcal{A}^\mathbf{L} \mathcal{B})
$$
在 $D(\mathcal{A}, \text{d})$ 中是同构，因为按照假设，
$\mathcal{A}\to\mathcal{B}$ 是（左）微分分次
$\mathcal{A}$-模的拟同构。特别地，本引理中的函子全忠实；见
《范畴》引理 \ref{categories-lemma-adjoint-fully-faithful}。
显然，限制函子
$D(\mathcal{B}, \text{d}) \to D(\mathcal{A}, \text{d})$
的核为零。于是结论由《导出范畴》引理
\ref{derived-lemma-fully-faithful-adjoint-kernel-zero} 得出。
\end{proof}













\section{微分分次代数的分辨}
\label{sdga-section-resolution-dgas}

\noindent
本节是《微分分次代数》第 \ref{dga-section-resolution-dgas} 节的类似物。

\medskip\noindent
设 $(\mathcal{C},\mathcal{O})$ 是环化位点。与评注
\ref{sdga-remark-sheaf-graded-sets} 一样，考虑 $\mathcal{C}$ 上的分次集合层
$\mathcal{S}$。把 $r$ 重自积
$\mathcal{S} \times \ldots \times \mathcal{S}$
看作按如下规则分次的集合层：
$\deg(s_1 \cdot \ldots \cdot s_r) = \sum \deg(s_i)$。
这里，给定局部截面 $s_i\in\mathcal{S}(U)$，$i=1,\ldots,r$，
用 $s_1\cdot\ldots\cdot s_r$ 表示
$\mathcal{S} \times \ldots \times \mathcal{S}$
在 $U$ 上的相应截面。记 $\mathcal{O}\langle\mathcal{S}\rangle$
为 $\mathcal{S}$ 上的自由分次 $\mathcal{O}$-代数。更准确地，置
$$
\mathcal{O}\langle \mathcal{S} \rangle =
\mathcal{O} \oplus
\bigoplus\nolimits_{r \geq 1}
\mathcal{O}[\mathcal{S} \times \ldots \times \mathcal{S}]
$$
，记号同评注 \ref{sdga-remark-sheaf-graded-sets}。按串接定义乘法：
$$
(s_1 \cdot \ldots \cdot s_r)  (s'_1 \cdot \ldots \cdot s'_{r'}) =
s_1 \cdot \ldots s_r \cdot s'_1 \cdot \ldots \cdot s'_{r'}
$$
，便使其成为分次 $\mathcal{O}$-代数层。可以令
$\mathcal{S}$ 的所有局部截面 $s$ 满足 $\text{d}(s)=0$，
再依 Leibniz 法则唯一延拓，从而在
$\mathcal{O}\langle\mathcal{S}\rangle$ 上赋予微分；不过考虑其他微分
同样重要。

\medskip\noindent
事实上，设给定如下系统：
\begin{enumerate}
\item 对 $i=0,1,2,\ldots$，给定分次集合层 $\mathcal{S}_i$；
\item 对 $i=0,1,2,\ldots$，给定态射
$$
\delta_{i + 1} : \mathcal{S}_{i + 1}
\longrightarrow
\mathcal{A}_i =
\mathcal{O}\langle \mathcal{S}_0 \amalg \ldots \amalg \mathcal{S}_i\rangle
$$
，它是次数为 $1$ 的分次集合层态射，且其像包含于靶上归纳定义的微分之核中。
\end{enumerate}
更准确地，先置
$\mathcal{A}_0 = \mathcal{O}\langle \mathcal{S}_0 \rangle$
，并在其上赋予满足 Leibniz 法则且使 $\mathcal{S}$ 的任意局部截面
$s$ 满足 $\text{d}(s)=0$ 的唯一微分。归纳地，设
$\mathcal{A}_i$ 上已给定微分 $\text{d}$。将它延拓为
$\mathcal{A}_{i+1}$ 上满足 Leibniz 法则并使
$$
\text{d}(s) = \delta(s)
$$
的唯一微分，其中若 $s$ 属于
$\mathcal{S}_0\amalg\ldots\amalg\mathcal{S}_{i+1}$ 的直和项
$\mathcal{S}_j$，则 $\delta(s)=\delta_j(s)$。此定义恰因
$\delta(s)$ 属于归纳定义的微分之核而有意义。

\begin{lemma}
\label{sdga-lemma-special-good}
在上述情形中，微分分次 $\mathcal{O}$-代数
$$
\mathcal{A} = \colim \mathcal{A}_i
$$
具有如下性质：对环化拓扑斯的任意态射
$(f, f^\sharp) : (\Sh(\mathcal{C}'), \mathcal{O}')
\to (\Sh(\mathcal{C}), \mathcal{O})$
，拉回 $f^*\mathcal{A}$ 作为分次 $\mathcal{O}'$-模是平坦的，
并且作为 $\mathcal{O}'$-模复形是 K-平坦的。
\end{lemma}

\begin{proof}
注意 $f^*\mathcal{A}=\colim f^*\mathcal{A}_i$，且
$$
f^*\mathcal{A}_i = \mathcal{O}'\langle
f^{-1}\mathcal{S}_0 \amalg \ldots \amalg f^{-1}\mathcal{S}_i\rangle
$$
，其微分由上述归纳过程以 $f^{-1}\delta_{i+1}$ 给出。因此，只需证明
$\mathcal{A}$ 作为分次 $\mathcal{O}$-模平坦，并且作为
$\mathcal{O}$-模复形 K-平坦。为此，只需证明每个
$\mathcal{A}_i$ 作为分次 $\mathcal{O}$-模平坦，并且作为
$\mathcal{O}$-模复形 K-平坦；可比较引理
\ref{sdga-lemma-good-direct-sum}。


\medskip\noindent
对 $i\geq1$，记
$\mathcal{S}=\mathcal{S}_0\amalg\ldots\amalg\mathcal{S}_i$，
于是作为分次 $\mathcal{O}$-代数，有
$\mathcal{A}_i=\mathcal{O}\langle\mathcal{S}\rangle$。我们要用
微分分次 $\mathcal{O}$-子模构造该代数的一个滤过。

\medskip\noindent
置 $W=\mathbf{Z}_{\geq0}^{i+1}$，并赋予字典序。也就是说，
对 $W$ 中的 $w=(w_0,\ldots w_i)$ 与
$w'=(w'_0,\ldots,w'_i)$，规定
$$
w > w' \Leftrightarrow
\exists j,\ 0 \leq j \leq i :
w_i = w'_i,\ w_{i - 1} = w'_{i - 1},\ \ldots ,
\ w_{j + 1} = w'_{j + 1},\ w_j > w'_j
$$
，依此类推。设给定
$s = s_1 \cdot \ldots \cdot s_r$，它是
$\mathcal{S}\times\ldots\times\mathcal{S}$ 在 $U$ 上的截面。若对每个
$a$，存在唯一的 $0\leq j_a\leq i$ 使
$s_a\in\mathcal{S}_{j_a}(U)$，便称{\it $s$ 的权已定义}。此时定义其权为
$$
w(s) = (w_0(s), \ldots, w_i(s)) \in W,\quad
w_j(s) = |\{a \mid j_a = j\}|
$$
$\mathcal{S}\times\ldots\times\mathcal{S}$ 的任意截面之权在局部均有定义。
读者容易验证，我们得到不交并分解
$$
\mathcal{S} \times \ldots \times \mathcal{S} =
\coprod\nolimits_{w \in W} \left(
\mathcal{S} \times \ldots \times \mathcal{S}\right)_w
$$
，其中各子层由给定权的截面组成。当然，对给定的 $r$，只会出现满足
$\sum_{0\leq j\leq i}w_j=r$ 的 $w\in W$。相应地得到分解
$$
\mathcal{A}_i = \mathcal{O} \oplus
\bigoplus\nolimits_{r \geq 1}
\bigoplus\nolimits_{w \in W}
\mathcal{O}[\left(\mathcal{S} \times \ldots \times \mathcal{S}\right)_w]
$$
证明的其余部分依赖如下显然观察：给定 $r$、$w$ 以及
$\left(\mathcal{S}\times\ldots\times\mathcal{S}\right)_w$ 的局部截面
$s=s_1\cdot\ldots\cdot s_r$，有
$$
\text{d}(s) \text{ 是如下对象的局部截面：} \mathcal{O} \oplus
\bigoplus\nolimits_{r' \geq 1}
\bigoplus\nolimits_{w' \in W,\ w' < w}
\mathcal{O}[\left(\mathcal{S} \times \ldots \times \mathcal{S}\right)_{w'}]
$$
这是因为在下列各表达式中
$$
(-1)^{\deg(s_1) + \ldots + \deg(s_{a - 1})}
s_1 \cdot \ldots s_{a - 1} \cdot \delta(s_a) \cdot
s_{a + 1} \cdot \ldots \cdot s_r
$$
（它们之和给出 $\text{d}(s)$），元素 $\delta(s_a)$ 在局部是诸元素
$s'_1\cdot\ldots\cdot s'_{r'}$ 的 $\mathcal{O}$-线性组合，其中
$s'_{a'}$ 属于 $\mathcal{S}_{j'_a}$，对某个
$0\leq j'_{a'}<j_a$；这里 $j_a$ 由 $s_a$ 是
$\mathcal{S}_{j_a}$ 的截面所确定。

\medskip\noindent
这意味着如下事实。对 $w\in W$，置
$$
F_w \mathcal{A}_i = \mathcal{O} \oplus \bigoplus\nolimits_{r \geq 1}
\bigoplus\nolimits_{w' \in W,\ w' \leq w}
\mathcal{O}[\left(\mathcal{S} \times \ldots \times \mathcal{S}\right)_{w'}]
$$
由上述观察，这是一个微分分次 $\mathcal{O}$-子模。得到微分分次
$\mathcal{A}$-模的容许短正合列
$$
0 \to \colim_{w' < w} F_{w'}\mathcal{A}_i \to
F_w\mathcal{A}_i \to
\bigoplus\nolimits_{r \geq 1}
\mathcal{O}[\left(\mathcal{S} \times \ldots \times \mathcal{S}\right)_w]
\to 0
$$
，其中右端的微分为零。

\medskip\noindent
现在对有序集 $W$ 作超限归纳以完成证明。微分分次复形
$F_0\mathcal{A}_0$ 是直和项 $\mathcal{O}$，因而 K-平坦且分次平坦。
对 $w\in W$，若结论对所有 $w'<w$ 的
$F_{w'}\mathcal{A}_i$ 成立，则由引理
\ref{sdga-lemma-good-direct-sum}、\ref{sdga-lemma-good-admissible-ses} 与
\ref{sdga-lemma-free-graded-module-good}，结论对 $w$ 也成立。最后，
$\mathcal{A}_i=\colim_{w\in W}F_w\mathcal{A}_i$，故得结论。
\end{proof}

\begin{lemma}
\label{sdga-lemma-good-dga}
设 $(\mathcal{C},\mathcal{O})$ 是环化位点，
$(\mathcal{B},\text{d})$ 是微分分次 $\mathcal{O}$-代数。存在微分分次
$\mathcal{O}$-代数的拟同构
$(\mathcal{A}, \text{d}) \to (\mathcal{B}, \text{d})$，使得
$\mathcal{A}$ 分次平坦，并且作为 $\mathcal{O}$-模复形
K-平坦；而且经环化拓扑斯的任意态射拉回后，同样的结论仍成立。
\end{lemma}

\begin{proof}
证明与引理 \ref{sdga-lemma-resolve} 的第一个证明完全相同，只是现在使用
自由分次代数而非自由分次模。

\medskip\noindent
我们像引理 \ref{sdga-lemma-special-good} 那样，通过构造
$$
\mathcal{A}_0 \to \mathcal{A}_1 \to \mathcal{A}_2 \to \ldots \to \mathcal{B}
$$
来构造 $\mathcal{A}=\colim\mathcal{A}_i$。令 $\mathcal{S}_0$ 为
次数 $n$ 部分是 $\Ker(\text{d}_\mathcal{B}^n)$ 的分次集合层
（评注 \ref{sdga-remark-sheaf-graded-sets}）。考虑微分分次模的同态
$$
\mathcal{A}_0 = \mathcal{O}\langle \mathcal{S}_0 \rangle
\longrightarrow
\mathcal{B}
$$
，该映射把 $\mathcal{S}_0$ 的局部截面 $s$ 映到
$\mathcal{A}^{\deg(s)}$ 中相应的局部截面（它属于微分之核，故我们的映射
确实是微分分次代数的映射）。按构造，所诱导的上同调层映射
$H^n(\mathcal{A}_0)\to H^n(\mathcal{B})$ 是满射，因而对所有 $i$
仍有同样结论。

\medskip\noindent
构造的归纳步骤如下。给定 $\mathcal{A}_i\to\mathcal{B}$，记
$\mathcal{S}_{i+1}$ 为次数 $n$ 部分是
$$
\Ker(\text{d}_{\mathcal{A}_i}^{n + 1})
\times_{\mathcal{B}^{n + 1}, \text{d}}
\mathcal{B}^n
$$
的分次集合层。它带有典范映射
$$
\delta_{i + 1} : \mathcal{S}_{i + 1} \longrightarrow
\mathcal{A}_i
$$
，按构造其像包含于 $\text{d}_{\mathcal{A}_i}$ 之核。因此
$\mathcal{A}_{i + 1}=\mathcal{O}\langle
\mathcal{S}_0\amalg\ldots\mathcal{S}_{i + 1}\rangle$ 带有延拓
$\mathcal{A}_i$ 上微分的微分；见本节开头的讨论。从
$\mathcal{A}_{i+1}$ 到 $\mathcal{B}$ 的映射是唯一的分次代数映射，
它在 $\mathcal{A}_i$ 上限制为给定映射，并把
$\mathcal{S}_{i+1}$ 的局部截面 $s=(a,b)$ 映到
$\mathcal{B}$ 中的 $b$。这与微分相容，恰因为 $\text{d}(b)$ 是
$a$ 在 $\mathcal{B}$ 中的像。

\medskip\noindent
映射 $\mathcal{A}\to\mathcal{B}$ 是拟同构：有
$H^n(\mathcal{A})=\colim H^n(\mathcal{A}_i)$，并且对每个 $i$，映射
$H^n(\mathcal{A}_i)\to H^n(\mathcal{B})$ 是满射，而按构造其核被映射
$H^n(\mathcal{A}_i)\to H^n(\mathcal{A}_{i+1})$ 零化。最后，
$\mathcal{A}$ 的平坦性条件已在引理 \ref{sdga-lemma-special-good} 中证明。
\end{proof}






\section{杂项}
\label{sdga-section-misc}

\noindent
设 $(f,f^\sharp):(\Sh(\mathcal{C}),\mathcal{O})
\to(\Sh(\mathcal{C}'),\mathcal{O}')$ 是环化拓扑斯的态射，且
$\mathcal{A}$ 是微分分次 $\mathcal{O}$-代数层。利用复合\footnote{更准确地，
应写成
$F(\mathcal{A}) \otimes_\mathcal{O}^\mathbf{L} F(\mathcal{A})
\to F(\mathcal{A} \otimes_\mathcal{O} \mathcal{A}) \to F(\mathcal{A})$
，其中 $F$ 表示到 $\mathcal{O}$-模复形的遗忘函子。还要注意，
$\mathcal{A}\otimes_\mathcal{O}\mathcal{A}$ 表示第
\ref{sdga-section-tensor-product-dg} 节的张量积，故
$F(\mathcal{A} \otimes_\mathcal{O} \mathcal{A}) =
\text{Tot}(F(\mathcal{A}) \otimes_\mathcal{O} F(\mathcal{A}))$.
该序列的第一个箭头是从两个 $\mathcal{O}$-模复形的导出张量积到
$\mathcal{O}$-模复形通常张量积的典范映射。}
$$
\mathcal{A} \otimes_\mathcal{O}^\mathbf{L} \mathcal{A}
\longrightarrow
\mathcal{A} \otimes_\mathcal{O} \mathcal{A}
\longrightarrow
\mathcal{A}
$$
以及相对杯积（见《站点上的上同调》评注
\ref{sites-cohomology-remark-cup-product} 与第
\ref{sites-cohomology-section-cup-product} 节），得到乘法\footnote{这里及下文，
$Rf_*:D(\mathcal{O})\to D(\mathcal{O}')$ 是《站点上的上同调》第
\ref{sites-cohomology-section-unbounded} 节及以后研究的导出函子。}
$$
\mu :
Rf_*\mathcal{A} \otimes_{\mathcal{O}'}^\mathbf{L} Rf_*\mathcal{A}
\longrightarrow
Rf_*\mathcal{A}
$$
，它位于 $D(\mathcal{O}')$ 中。此乘法具有结合性，即图
$$
\xymatrix{
Rf_*\mathcal{A} \otimes_{\mathcal{O}'}^\mathbf{L} Rf_*\mathcal{A}
\otimes_{\mathcal{O}'}^\mathbf{L} Rf_*\mathcal{A}
\ar[rr]_-{\mu \otimes 1} \ar[d]_{1 \otimes \mu} & &
Rf_*\mathcal{A} \otimes_{\mathcal{O}'}^\mathbf{L} Rf_*\mathcal{A}
\ar[d]^\mu \\
Rf_*\mathcal{A} \otimes_{\mathcal{O}'}^\mathbf{L} Rf_*\mathcal{A}
\ar[rr]^-\mu & &
Rf_*\mathcal{A}
}
$$
在 $D(\mathcal{O}')$ 中交换；这由《站点上的上同调》引理
\ref{sites-cohomology-lemma-cup-product-associative} 得出。
完全类似地，给定右微分分次 $\mathcal{A}$-模 $\mathcal{M}$，得到乘法
$$
\mu_\mathcal{M} :
Rf_*\mathcal{M} \otimes_{\mathcal{O}'}^\mathbf{L} Rf_*\mathcal{A}
\longrightarrow
Rf_*\mathcal{M}
$$
，它位于 $D(\mathcal{O}')$ 中。此乘法与上述 $\mu$ 相容，即图
$$
\xymatrix{
Rf_*\mathcal{M} \otimes_{\mathcal{O}'}^\mathbf{L} Rf_*\mathcal{A}
\otimes_{\mathcal{O}'}^\mathbf{L} Rf_*\mathcal{A}
\ar[rr]_-{\mu_\mathcal{M} \otimes 1} \ar[d]_{1 \otimes \mu} & &
Rf_*\mathcal{M} \otimes_{\mathcal{O}'}^\mathbf{L} Rf_*\mathcal{A}
\ar[d]^{\mu_\mathcal{M}} \\
Rf_*\mathcal{M} \otimes_{\mathcal{O}'}^\mathbf{L} Rf_*\mathcal{A}
\ar[rr]^-{\mu_\mathcal{M}} & &
Rf_*\mathcal{M}
}
$$
在 $D(\mathcal{O}')$ 中交换；这仍由《站点上的上同调》引理
\ref{sites-cohomology-lemma-cup-product-associative} 得出。

\medskip\noindent
上述构造的一个特殊例子是取 $f$ 为到单点拓扑斯 $\Sh(pt)$ 的态射。
此时 $\mu$ 正是杯积映射
$$
R\Gamma(\mathcal{C}, \mathcal{A})
\otimes_{\Gamma(\mathcal{C}, \mathcal{O})}^\mathbf{L}
R\Gamma(\mathcal{C}, \mathcal{A})
\longrightarrow
R\Gamma(\mathcal{C}, \mathcal{A}),
\quad \eta \otimes \theta \mapsto \eta \cup \theta
$$
，而类似地，$\mu_\mathcal{M}$ 是杯积映射
$$
R\Gamma(\mathcal{C}, \mathcal{M})
\otimes_{\Gamma(\mathcal{C}, \mathcal{O})}^\mathbf{L}
R\Gamma(\mathcal{C}, \mathcal{A})
\longrightarrow
R\Gamma(\mathcal{C}, \mathcal{M}),
\quad \eta \otimes \theta \mapsto \eta \cup \theta
$$
一般地，借助《站点上的上同调》评注
\ref{sites-cohomology-remark-before-Leray} 中的同一
$$
R\Gamma(\mathcal{C}, \mathcal{A}) =
R\Gamma(\mathcal{C}', Rf_*\mathcal{A})
\quad\text{且}\quad
R\Gamma(\mathcal{C}, \mathcal{M}) =
R\Gamma(\mathcal{C}', Rf_*\mathcal{M})
$$
，映射 $\mu_\mathcal{M}$ 在上同调上诱导杯积。为此，应用
《站点上的上同调》引理
\ref{sites-cohomology-lemma-compose-cup-product}，其中拓扑斯的第二个态射
取为上述从 $\Sh(\mathcal{C}')$ 到单点拓扑斯的态射。

\medskip\noindent
若 $\mathcal{M}_1\to\mathcal{M}_2$ 是右微分分次
$\mathcal{A}$-模的同态，则图
$$
\xymatrix{
Rf_*\mathcal{M}_1 \otimes_{\mathcal{O}'}^\mathbf{L} Rf_*\mathcal{A}
\ar[rr]_-{\mu_{\mathcal{M}_1}} \ar[d] & &
Rf_*\mathcal{M}_1 \ar[d] \\
Rf_*\mathcal{M}_2 \otimes_{\mathcal{O}'}^\mathbf{L} Rf_*\mathcal{A}
\ar[rr]^-{\mu_{\mathcal{M}_2}} & &
Rf_*\mathcal{M}_2
}
$$
在 $D(\mathcal{O}')$ 中交换；这源于相对杯积的函子性。设有右微分分次
$\mathcal{A}$-模的短正合列
$$
0 \to \mathcal{M}_1 \xrightarrow{a} \mathcal{M}_2 \to \mathcal{M}_3 \to 0
$$
。我们断言图
$$
\xymatrix{
Rf_*\mathcal{M}_3 \otimes_{\mathcal{O}'}^\mathbf{L} Rf_*\mathcal{A}
\ar[rr]_-{\mu_{\mathcal{M}_3}} \ar[d]_{Rf_*\delta \otimes \text{id}} & & 
Rf_*\mathcal{M}_3 \ar[d]^{Rf_*\delta} \\
Rf_*\mathcal{M}_1[1] \otimes_{\mathcal{O}'}^\mathbf{L} Rf_*\mathcal{A}
\ar[rr]^-{\mu_{\mathcal{M}_1[1]}} & &
Rf_*\mathcal{M}_1[1]
}
$$
在 $D(\mathcal{O}')$ 中交换，其中
$\delta:\mathcal{M}_3\to\mathcal{M}_1[1]$ 是给定短正合列所给出的
$D(\mathcal{O})$ 中的态射（见《导出范畴》第
\ref{derived-section-canonical-delta-functor} 节）。若我们的序列作为分次右
$\mathcal{A}$-模列分裂，则这一点很清楚，因为此时 $\delta$ 可由右
$\mathcal{A}$-模映射表示，从而可应用上述讨论。一般情形下，使用
$a$ 的锥和图
$$
\xymatrix{
\mathcal{M}_1 \ar[r]_a \ar[d] &
\mathcal{M}_2 \ar[r]_i \ar[d] &
C(a) \ar[r]_{-p} \ar[d]^q &
\mathcal{M}_1[1] \ar[d] \\
\mathcal{M}_1 \ar[r] &
\mathcal{M}_2 \ar[r] &
\mathcal{M}_3 \ar[r]^\delta &
\mathcal{M}_1[1]
}
$$
。由《导出范畴》引理
\ref{derived-lemma-derived-canonical-delta-functor} 中 $\delta$ 的定义，
右边的方块在 $D(\mathcal{O})$ 中交换。现在，锥 $C(a)$ 具有右微分分次
$\mathcal{A}$-模结构，使得 $i$、$p$、$q$ 是右微分分次
$\mathcal{A}$-模的同态；见定义 \ref{sdga-definition-cone}。因此由上可知，
对应图对态射 $q$ 与 $-p$ 交换。由于 $q$ 在 $D(\mathcal{O})$ 中是同构，
遂如所需得出对 $\delta$ 也有同样结论。

\medskip\noindent
在上述情形中，给定右微分分次 $\mathcal{A}$-模 $\mathcal{M}$，取
$$
\xi \in H^n(\mathcal{C}, \mathcal{M})
$$
。换言之，把 $\mathcal{M}$ 视为 $\mathcal{O}$-模复形时，$\xi$ 是其
次数 $n$ 上同调类。由引理 \ref{sdga-lemma-cohomology-ext}，可以构造右微分分次
$\mathcal{A}$-模的映射
$$
x : \mathcal{A} \rightarrow \mathcal{M}'[n]
\quad\text{且}\quad
s : \mathcal{M} \to \mathcal{M}'
$$
，其中 $s$ 是拟同构，并且 $\xi$ 是
$1\in H^0(\mathcal{C},\mathcal{A})$ 在导出范畴
$D(\mathcal{A},\text{d})$ 中经态射 $s[n]^{-1}\circ x$ 的像，
因而在导出范畴 $D(\mathcal{O})$ 中亦然。因此，$D(\mathcal{O})$
中的相应映射
$$
\xi' = (s[n])^{-1} \circ x : \mathcal{A} \longrightarrow \mathcal{M}[n]
$$
由以下两个性质唯一刻画：
\begin{enumerate}
\item $\xi'$ 可提升为 $D(\mathcal{A},\text{d})$ 中的态射；
\item 在
$H^0(\mathcal{C}, \mathcal{M}[n]) = H^n(\mathcal{C}, \mathcal{M})$
中有 $\xi=\xi'(1)$。
\end{enumerate}
由上面讨论的 $x$ 与 $s$ 同相对杯积的相容性可知，对环化拓扑斯的每个
态射\footnote{例如恒等态射。}
$(f, f^\sharp) : (\Sh(\mathcal{C}), \mathcal{O})
\to (\Sh(\mathcal{C}'), \mathcal{O}')$，$\xi'$ 的导出直像
$$
Rf_*\xi' : Rf_*\mathcal{A}  \longrightarrow Rf_*\mathcal{M}[n]
$$
与上述构造的映射 $\mu$ 和 $\mu_{\mathcal{M}[n]}$ 相容，即图
$$
\xymatrix{
Rf_*\mathcal{A} \otimes_{\mathcal{O}'}^\mathbf{L} Rf_*\mathcal{A}
\ar[rr]_-\mu \ar[d]_{Rf_*\xi' \otimes \text{id}}  & &
Rf_*\mathcal{A} \ar[d]^{Rf_*\xi'} \\
Rf_*\mathcal{M}[n] \otimes_{\mathcal{O}'}^\mathbf{L} Rf_*\mathcal{A}
\ar[rr]^-{\mu_{\mathcal{M}[n]}} & &
Rf_*\mathcal{M}[n]
}
$$
在 $D(\mathcal{O}')$ 中交换。把此相容性用于到单点拓扑斯的映射，
特别得到图
$$
\xymatrix{
R\Gamma(\mathcal{C}, \mathcal{A})
\otimes_{\Gamma(\mathcal{C}, \mathcal{O})}^\mathbf{L}
R\Gamma(\mathcal{C}, \mathcal{A})
\ar[d]_{\xi' \otimes \text{id}} \ar[r] &
R\Gamma(\mathcal{C}, \mathcal{A}) \ar[d]^{\xi'} \\
R\Gamma(\mathcal{C}, \mathcal{M}[n])
\otimes_{\Gamma(\mathcal{C}, \mathcal{O})}^\mathbf{L}
R\Gamma(\mathcal{C}, \mathcal{A})
\ar[r] &
R\Gamma(\mathcal{C}, \mathcal{M}[n])
}
$$
交换。结合 $\xi'(1)=\xi$，这说明所诱导的上同调映射
$$
\xi' : R\Gamma(\mathcal{C}, \mathcal{A}) \to
R\Gamma(\mathcal{C}, \mathcal{M}[n]), \quad \eta \mapsto \xi \cup \eta
$$
如所示由以 $\xi$ 作左杯乘给出。














\section{范畴上的微分分次模}
\label{sdga-section-modules-cohomology}

\noindent
本节续接《站点上的上同调》第
\ref{sites-cohomology-section-modules-cohomology} 节。

\medskip\noindent
设 $\mathcal{C}$ 是范畴。把 $\mathcal{C}$ 视为带混沌拓扑的位点。设
$\mathcal{O}$ 是 $\mathcal{C}$ 上的环层，
$(\mathcal{A},\text{d})$ 是微分分次 $\mathcal{O}$-代数层。换言之，
$\mathcal{O}$ 是范畴 $\mathcal{C}$ 上的环预层，而
$(\mathcal{A},\text{d})$ 是 $\mathcal{C}$ 上的微分分次
$\mathcal{O}$-代数预层；见《范畴》定义
\ref{categories-definition-presheaf}。

\begin{definition}
\label{sdga-definition-cartesian}
在上述情形中，以
{\it $\mathit{QC}(\mathcal{A}, \text{d})$}
表示 $D(\mathcal{A},\text{d})$ 的全子范畴，其对象为满足如下条件的
$M$：对 $\mathcal{C}$ 中所有 $U\to V$，典范映射
$$
R\Gamma(V, M) \otimes_{\mathcal{A}(V)}^\mathbf{L} \mathcal{A}(U)
\longrightarrow
R\Gamma(U, M)
$$
在 $D(\mathcal{A}(U),\text{d})$ 中是同构。
\end{definition}

\begin{lemma}
\label{sdga-lemma-cartesian-good-sub}
在上述情形中，子范畴 $\mathit{QC}(\mathcal{A},\text{d})$ 是
$D(\mathcal{A},\text{d})$ 的严格满的饱和三角子范畴，并在任意直和下保持。
\end{lemma}

\begin{proof}
设 $U$ 是 $\mathcal{C}$ 的对象。由于 $\mathcal{C}$ 上的拓扑是混沌的，
函子 $\mathcal{F}\mapsto\mathcal{F}(U)$ 正合且与直和交换。因此，
正合函子 $M\mapsto R\Gamma(U,M)$ 可如下计算：以任意微分分次
$\mathcal{A}$-模 $\mathcal{M}$ 表示 $K$，再取 $\mathcal{M}(U)$。
故 $R\Gamma(U,-)$ 与直和交换；见引理
\ref{sdga-lemma-derived-products}。类似地，给定 $\mathcal{C}$ 的态射
$U\to V$，导出张量积函子
$- \otimes_{\mathcal{O}(A)}^\mathbf{L} \mathcal{A}(U) :
D(\mathcal{A}(V)) \to D(\mathcal{A}(U))$
正合且与直和交换。本引理由这些观察直接得出；略去细节。
\end{proof}

\begin{remark}
\label{sdga-remark-cartesian-brown}
仍按上述记号，设 $\mathcal{C}$ 是视为带混沌拓扑位点的范畴，
$\mathcal{O}$ 是 $\mathcal{C}$ 上的环层，而
$(\mathcal{A},\text{d})$ 是微分分次 $\mathcal{O}$-代数层。则
《站点上的上同调》命题
\ref{sites-cohomology-proposition-cartesian-brown} 的类似结论对
$\mathit{QC}(\mathcal{A},\text{d})$ 成立，证明几乎完全相同：
\begin{enumerate}
\item 任何把直和变为积的反变上同调函子
$H : \mathit{QC}(\mathcal{A}, \text{d}) \to \textit{Ab}$
都是可表的；
\item 三角范畴之间任何把直和变为直和的正合函子
$F : \mathit{QC}(\mathcal{A}, \text{d}) \to \mathcal{D}$
都有正合右伴随；
\item 含入函子
$\mathit{QC}(\mathcal{A}, \text{d}) \to D(\mathcal{A}, \text{d})$
有正合右伴随。
\end{enumerate}
若今后需要此结论，我们将在这里精确表述并证明它。
\end{remark}

\noindent
设 $u:\mathcal{C}'\to\mathcal{C}$ 是范畴之间的函子。若把
$\mathcal{C}$ 与 $\mathcal{C}'$ 视为带混沌拓扑的位点，则 $u$ 是连续且
余连续的函子。因此得到拓扑斯的态射
$g:\Sh(\mathcal{C}')\to\Sh(\mathcal{C})$；见《位点》引理
\ref{sites-lemma-cocontinuous-morphism-topoi}。再设给定
$\mathcal{C}$ 上的环层 $\mathcal{O}$、$\mathcal{C}'$ 上的环层
$\mathcal{O}'$，以及映射
$g^\sharp : g^{-1}\mathcal{O} \to \mathcal{O}'$.
把相应的环化拓扑斯态射简记为
$g : (\Sh(\mathcal{C}'), \mathcal{O}') \to (\Sh(\mathcal{C}), \mathcal{O})$,
；见《位点上的模》第 \ref{sites-modules-section-ringed-topoi} 节。
最后，设 $(\mathcal{A},\text{d})$ 是微分分次 $\mathcal{O}$-代数层，
$(\mathcal{A}',\text{d})$ 是微分分次 $\mathcal{O}'$-代数层，并且给定
微分分次 $\mathcal{O}'$-代数的映射
$\varphi:g^*\mathcal{A}\to\mathcal{A}'$（见第
\ref{sdga-section-functoriality-dg} 节）。

\begin{lemma}
\label{sdga-lemma-cartesian-functorial}
设 $g:(\Sh(\mathcal{C}'),\mathcal{O}')
\to(\Sh(\mathcal{C}),\mathcal{O})$ 与
$\varphi:g^*\mathcal{A}\to\mathcal{A}'$ 如上。则函子
$Lg^* : D(\mathcal{A}, \text{d}) \to D(\mathcal{A}', \text{d})$
把 $\mathit{QC}(\mathcal{A},\text{d})$ 映入
$\mathit{QC}(\mathcal{A}',\text{d})$。
\end{lemma}

\begin{proof}
取 $U'\in\Ob(\mathcal{C}')$，其在 $\mathcal{C}$ 中的像记为
$U=u(U')$。以 $pt$ 表示仅有一个对象和一个态射的范畴。把所示环化拓扑斯
$(\Sh(pt),\mathcal{O}'(U'))$ 与 $(\Sh(pt),\mathcal{O}(U))$
分别赋予微分分次代数 $\mathcal{A}'(U)$ 与 $\mathcal{A}(U)$。
当然，把其上的微分分次模导出范畴分别等同于
$D(\mathcal{A}'(U'),\text{d})$ 与
$D(\mathcal{A}(U),\text{d})$。于是有环化拓扑斯的交换图
$$
\xymatrix{
(\Sh(pt), \mathcal{O}'(U')) \ar[rr]_{U'} \ar[d] & &
(\Sh(\mathcal{C}'), \mathcal{O}') \ar[d]^g \\
(\Sh(pt), \mathcal{O}(U)) \ar[rr]^U & &
(\Sh(\mathcal{C}), \mathcal{O})
}
$$
，各顶点均赋有相应的微分分次代数。沿下方水平态射拉回，把
$D(\mathcal{A},\text{d})$ 中的 $M$ 映为视作
$D(\mathcal{A}(U),\text{d})$ 中对象的 $R\Gamma(U,K)$。
沿左边竖直箭头拉回，把 $M$ 映为
$M \otimes_{\mathcal{A}(U)}^\mathbf{L} \mathcal{A}'(U')$.
沿此图任一方向复合都给出相同结果（引理
\ref{sdga-lemma-compose-pullback}），故得到
$$
R\Gamma(U', Lg^*K) =
R\Gamma(U, K) \otimes_{\mathcal{A}(U)}^\mathbf{L} \mathcal{A}'(U')
$$
最后，设 $f':U'\to V'$ 是 $\mathcal{C}'$ 中的态射，并把它在
$\mathcal{C}$ 中的像记为
$f=u(f'):U=u(U')\to V=u(V')$。若
$K\in\mathit{QC}(\mathcal{A},\text{d})$，则
\begin{align*}
R\Gamma(V', Lg^*K) \otimes_{\mathcal{A}'(V')}^\mathbf{L} \mathcal{A}'(U')
& =
R\Gamma(V, K) \otimes_{\mathcal{A}(V)}^\mathbf{L} \mathcal{A}'(V')
\otimes_{\mathcal{A}'(V')}^\mathbf{L} \mathcal{A}'(U') \\
& =
R\Gamma(V, K) \otimes_{\mathcal{A}(V)}^\mathbf{L} \mathcal{A}'(U') \\
& =
R\Gamma(V, K) \otimes_{\mathcal{A}(V)}^\mathbf{L} \mathcal{A}(U)
\otimes_{\mathcal{A}(U)}^\mathbf{L} \mathcal{A}'(U') \\
& =
R\Gamma(U, K) \otimes_{\mathcal{A}(U)}^\mathbf{L} \mathcal{A}'(U') \\
& =
R\Gamma(U', Lg^*K)
\end{align*}
，正如所需。这里对 $U'$ 与 $V'$ 都使用了上述观察。
\end{proof}






\section{范畴上的微分分次模（续）}
\label{sdga-section-modules-cohomology-bis}

\noindent
我们进一步给出第 \ref{sdga-section-modules-cohomology} 节所引入的拟凝聚模概念的
若干结果。

\begin{lemma}
\label{sdga-lemma-cartesian-cofinal-subcategory}
设 $\mathcal{C}$、$\mathcal{O}$、$\mathcal{A}$ 如第
\ref{sdga-section-modules-cohomology} 节。设
$\mathcal{C}'\subset\mathcal{C}$ 是具有如下性质的全子范畴：对每个
$U\in\Ob(\mathcal{C})$，由箭头 $U\to U'$ 构成的范畴
$U/\mathcal{C}'$ 是余滤过的。以 $\mathcal{O}'$、$\mathcal{A}'$
表示 $\mathcal{O}$、$\mathcal{A}$ 到 $\mathcal{C}'$ 的限制。则限制诱导等价
$\mathit{QC}(\mathcal{A}, \text{d}) \to \mathit{QC}(\mathcal{A}', \text{d})$.
\end{lemma}

\begin{proof}
我们构造该函子的拟逆。设
$M'$ 是 $\mathit{QC}(\mathcal{A}',\text{d})$ 的对象。可以用良好微分分次模
$\mathcal{M}'$ 表示 $M'$；见引理 \ref{sdga-lemma-resolve}。于是对每个
$U'\in\Ob(\mathcal{C}')$，微分分次
$\mathcal{A}'(U')$-模 $\mathcal{M}'(U)$ 是 K-平坦且分次平坦的；
对 $\mathcal{C}'$ 的每个态射 $U'_1\to U'_2$，映射
$$
\mathcal{M}'(U'_2) \otimes_{\mathcal{A}'(U'_2)} \mathcal{A}'(U'_1)
\longrightarrow
\mathcal{M}'(U'_1)
$$
是拟同构（因为源表示导出张量积）。考虑按如下规则定义的微分分次
$\mathcal{A}$-模 $\mathcal{M}$：
$$
\mathcal{M}(U) =
\colim_{U \to U' \in U/\mathcal{C}'}
\mathcal{M}'(U') \otimes_{\mathcal{A}'(U')} \mathcal{A}(U)
$$
。由本引理中的假设，这是复形的滤过余极限。由于
$M'\in\mathit{QC}(\mathcal{A}',\text{d})$，该系统中的所有转移映射都是
拟同构。又因滤过余极限正合，可知对任意态射 $U\to U'$，其中
$U'\in\Ob(\mathcal{C}')$，$D(\mathcal{A}(U),\text{d})$ 中的
$\mathcal{M}(U)$ 同构于
$\mathcal{M}'(U')\otimes_{\mathcal{A}'(U')}\mathcal{A}(U)$。

\medskip\noindent
我们断言 $\mathcal{M}\in\mathit{QC}(\mathcal{A},\text{d})$：事实上，
给定 $\mathcal{C}$ 中的 $U\to V$，选择映射 $V\to V'$，其中
$V'\in\Ob(\mathcal{C}')$。由上可知，映射
$\mathcal{M}(V)\to\mathcal{M}(U)$ 等同于映射
$$
\mathcal{M}'(V') \otimes_{\mathcal{A}'(V')} \mathcal{A}(V)
\longrightarrow
\mathcal{M}'(V') \otimes_{\mathcal{A}'(V')} \mathcal{A}(U)
$$
。由于 $\mathcal{M'}(V')$ 作为微分分次
$\mathcal{A}'(V')$-模是 K-平坦的，断言得证。

\medskip\noindent
自然映射 $\mathcal{M}|_{\mathcal{C}'}\to\mathcal{M}'$ 在
$D(\mathcal{A}',d)$ 中是同构，这立即由上得出。

\medskip\noindent
反之，若有 $\mathit{QC}(\mathcal{A},\text{d})$ 的对象 $E$，则以良好
微分分次模 $\mathcal{E}$ 表示它。置
$\mathcal{M}'=\mathcal{E}|_{\mathcal{C}'}$（这仍是良好微分分次模），
可知存在映射
$$
\mathcal{E} \to \mathcal{M}
$$
，它在 $\mathcal{C}$ 的 $U$ 上由映射
$$
\mathcal{E}(U)
\longrightarrow
\colim_{U \to U' \in U/\mathcal{C}'}
\mathcal{E}(U') \otimes_{\mathcal{A}'(U')} \mathcal{A}(U)
$$
给出；同理这是拟同构。因此限制与上述构造是互为拟逆的函子，正如所需。
\end{proof}

\begin{lemma}
\label{sdga-lemma-cartesian-qis-equivalence}
设 $\mathcal{C}$、$\mathcal{O}$ 如第
\ref{sdga-section-modules-cohomology} 节。设
$\varphi:\mathcal{A}\to\mathcal{B}$ 是微分分次
$\mathcal{O}$-代数的同态，并在上同调层上诱导同构，则引理
\ref{sdga-lemma-qis-equivalence} 的等价
$D(\mathcal{A}, \text{d}) \to D(\mathcal{B}, \text{d})$
诱导等价
$\mathit{QC}(\mathcal{A}, \text{d}) \to
\mathit{QC}(\mathcal{B}, \text{d})$.
\end{lemma}

\begin{proof}
只需证明如下事实：给定 $\mathcal{C}$ 的态射 $U\to V$ 与
$D(\mathcal{A},\text{d})$ 中的 $M$，下列条件等价：
\begin{enumerate}
\item $R\Gamma(V, M) \otimes_{\mathcal{A}(V)}^\mathbf{L} \mathcal{A}(U)
\to \Gamma(U, M)$ 在 $D(\mathcal{A}(U),\text{d})$ 中是同构；
\item $R\Gamma(V, M \otimes_\mathcal{A}^\mathbf{L} \mathcal{B})
\otimes_{\mathcal{B}(V)}^\mathbf{L} \mathcal{B}(U)
\to \Gamma(U, M \otimes_\mathcal{A}^\mathbf{L} \mathcal{B})$
在 $D(\mathcal{B}(U),\text{d})$ 中是同构。
\end{enumerate}
由于 $\mathcal{C}$ 上的拓扑是混沌的，这只归结为
$\mathcal{A}(U)\to\mathcal{B}(U)$ 与
$\mathcal{A}(V)\to\mathcal{B}(V)$ 是拟同构这一事实。略去细节。
\end{proof}






\section{微分分次代数的逆系统}
\label{sdga-section-qc-inverse-system}

\noindent
本节考虑第 \ref{sdga-section-modules-cohomology} 节所讨论情形的如下特殊情况：
\begin{enumerate}
\item $\mathcal{C}$ 是范畴 $\mathbf{N}$，当且仅当 $i\leq j$ 时，
存在唯一态射 $i\to j$；
\item $\mathcal{O}$ 是取值为给定环 $R$ 的常值环（预）层。
\end{enumerate}
在此设定下，微分分次 $\mathcal{O}$-代数层 $\mathcal{A}$ 等同于
微分分次 $R$-代数的逆系统 $(A_n)$。微分分次
$\mathcal{A}$-模层 $\mathcal{M}$ 等同于逆系统 $(M_n)$，其中
$M_n$ 是微分分次 $A_n$-模，而转移映射
$M_{n+1}\to M_n$ 是 $A_{n+1}$-模映射。本节在记号中省略微分，
并以 $\mathit{QC}(\mathcal{A})$ 表示定义
\ref{sdga-definition-cartesian} 中记为
$\mathit{QC}(\mathcal{A},\text{d})$ 的范畴。

\medskip\noindent
设 $\mathcal{B}=(B_n)$ 是微分分次 $R$-代数的另一个逆系统。给定
pro-对象的态射 $\varphi:(A_n)\to(B_n)$，我们将构造从
$\mathit{QC}(\mathcal{A})$ 到 $\mathit{QC}(\mathcal{B})$ 的正合函子。
具体地，按照《范畴》例
\ref{categories-example-pro-morphism-inverse-systems}，态射 $\varphi$
由整数序列 $\ldots\geq m(3)\geq m(2)\geq m(1)$ 及交换图
$$
\xymatrix{
\ldots \ar[r] &
A_{m(3)} \ar[d]^{\varphi_3} \ar[r] &
A_{m(2)} \ar[d]^{\varphi_2} \ar[r] &
A_{m(1)} \ar[d]^{\varphi_1} \\
\ldots \ar[r] &
B_3 \ar[r] &
B_2 \ar[r] &
B_1
}
$$
给出，其中各项是微分分次 $R$-代数。于是，给定表示
$\mathit{QC}(\mathcal{A})$ 中对象的良好微分分次
$\mathcal{A}$-模层 $\mathcal{M}=(M_n)$，可以置
$$
N_n = M_{m(n)} \otimes_{A_{m(n)}} B_n
$$
。此逆系统确定 $\mathit{QC}(\mathcal{B})$ 的对象，因为
$A_{m(n)}$-模 $M_{m(n)}$ 是 K-平坦的；略去细节。所得函子不依赖其构造中
所作选择这一事实也留给读者证明。

\begin{lemma}
\label{sdga-lemma-pro-isomorphic}
在上述情形中，设 $\mathcal{A}=(A_n)$ 和 $\mathcal{B}=(B_n)$ 是微分分次
$R$-代数的逆系统。若 $\varphi:(A_n)\to(B_n)$ 是 pro-对象的同构，则上面构造的
函子 $\mathit{QC}(\mathcal{A})\to\mathit{QC}(\mathcal{B})$ 是等价。
\end{lemma}

\begin{proof}
令 $\psi:(B_n)\to(A_n)$ 为与 $\varphi$ 互逆的 pro-对象态射。根据
《范畴》，例 \ref{categories-example-pro-morphism-inverse-systems} 中的讨论，
可设 $\varphi$ 由如上的一族映射给出，而 $\psi$ 由
$n(1)<n(2)<\ldots$ 以及如下交换图给出：
$$
\xymatrix{
\ldots \ar[r] &
B_{n(3)} \ar[d]^{\psi_3} \ar[r] &
B_{n(2)} \ar[d]^{\psi_2} \ar[r] &
B_{n(1)} \ar[d]^{\psi_1} \\
\ldots \ar[r] &
A_3 \ar[r] &
A_2 \ar[r] &
A_1
}
$$
其中各项是微分分次 $R$-代数。由于 $\varphi\circ\psi=\text{id}$，在可能增大
函数 $n(\cdot)$ 和 $m(\cdot)$ 的取值后，可以假设
$B_{n(m(i))}\to A_{m(i)}\to B_i$ 是转移映射。因此，函子的复合
$$
\mathit{QC}(\mathcal{B}) \to
\mathit{QC}(\mathcal{A}) \to
\mathit{QC}(\mathcal{B})
$$
把良好微分分次 $\mathcal{B}$-模层 $\mathcal{N}=(N_n)$ 送到取值为
$$
N'_i = N_{n(m(i))} \otimes_{B_{n(m(i))}} B_i
$$
的逆系统 $\mathcal{N}'=(N'_i)$；它与 $\mathcal{N}$ 典范拟同构，恰是因为
$\mathcal{N}$ 是 $\mathit{QC}(\mathcal{B})$ 的对象，且对所有 $j$，$N_j$
都是 K-平坦微分分次模。反向复合也有同样结论，故证毕。
\end{proof}

\noindent
令 $\mathcal{C}=\mathbf{N}$，令 $\mathcal{O}$ 为取值于环 $R$ 的常值层，并令
$\mathcal{A}$ 由微分分次 $R$-代数的逆系统 $(A_n)$ 给出。设有两个左微分分次
$\mathcal{A}$-模 $\mathcal{N}$ 和 $\mathcal{N}'$，它们分别由逆系统
$(N_n)$ 和 $(N'_n)$ 给出。因此每个 $N_n$ 和 $N'_n$ 都是左微分分次
$A_n$-模。暂称 $(N_n)$ 和 $(N'_n)$ {\it 在导出范畴中 pro-同构}，如果存在
整数序列
$$
1 = n_0 < n_1 < n_2 < n_3 < \ldots
$$
以及映射
$$
N_{n_{2i}} \to N'_{n_{2i - 1}}
\quad\text{于}\quad
D(A_{n_{2i}}^{opp}, \text{d})
$$
和
$$
N'_{n_{2i + 1}} \to N'_{n_{2i}}
\quad\text{于}\quad
D(A_{n_{2i + 1}}^{opp}, \text{d})
$$
使得复合 $N_{n_{2i}}\to N_{n_{2i-2}}$ 和
$N'_{n_{2i+1}}\to N'_{2i-1}$ 分别由相应系统的转移映射给出。

\begin{lemma}
\label{sdga-lemma-pro-isomorphic-systems-tensor}
若 $(N_n)$ 和 $(N'_n)$ 按上述定义在导出范畴中 pro-同构，则对
$D(\mathbf{N},\mathcal{A})$ 的每个对象 $(M_n)$，有
$$
R\lim (M_n \otimes_{A_n}^\mathbf{L} N_n) =
R\lim (M_n \otimes_{A_n}^\mathbf{L} N'_n)
$$
于 $D(R)$ 中。
\end{lemma}

\begin{proof}
该假设蕴涵 $D(R)$ 中的逆系统
$(M_n\otimes_{A_n}^\mathbf{L}N_n)$ 与逆系统
$(M_n\otimes_{A_n}^\mathbf{L}N'_n)$ pro-同构（通常意义下）。因此结论由
《代数详论》，引理 \ref{more-algebra-lemma-pro-isomorphism-Rlim} 得出。
\end{proof}

\begin{lemma}
\label{sdga-lemma-different-tensors}
\begin{reference}
这是 \cite[引理 3.5.4]{BS} 的一个变体。
\end{reference}
设 $R$ 为环，$f_1,\ldots,f_r\in R$。令 $K_n$ 为
$f_1^n,\ldots,f_r^n$ 上的 Koszul 复形，并将其视为微分分次 $R$-代数。
设 $(M_n)$ 是 $D(\mathbf{N},(K_n))$ 的对象。则对任意 $t\geq1$，有
$$
R\lim (M_n \otimes_R^\mathbf{L} K_t) =
R\lim (M_n \otimes_{K_n}^\mathbf{L} K_t)
$$
于 $D(R)$ 中。
\end{lemma}

\begin{proof}
固定 $t\geq1$。对 $n\geq t$，用 ${}_nK_t$ 表示微分分次 $R$-代数 $K_t$，
将其视为左微分分次 $K_n$-模。注意
$$
M_n \otimes_R^\mathbf{L} K_t =
M_n \otimes_{K_n}^\mathbf{L} (K_n \otimes_R^\mathbf{L} K_t) =
M_n \otimes_{K_n}^\mathbf{L} (K_n \otimes_R K_t)
$$
。因此由引理 \ref{sdga-lemma-pro-isomorphic-systems-tensor}，只需证明
$({}_nK_t)$ 和 $(K_n\otimes_RK_t)$ 在导出范畴中 pro-同构。乘法映射
$$
K_n \otimes_R K_t \longrightarrow {}_nK_t
$$
是左微分分次 $K_n$-模的映射。因此，为完成证明，只需证明对所有
$n\geq1$，存在 $N>n$ 以及映射
$$
{}_NK_t \longrightarrow {}_NK_n \otimes_R K_t
$$
于 $D(K_N^{opp},\text{d})$ 中，使它与乘法映射的复合（沿任一方向）是转移
映射。这已在《除幂代数》，引理 \ref{dpa-lemma-construct-some-maps} 中通过
一个显式构造完成。
\end{proof}

\begin{proposition}
\label{sdga-proposition-derived-complete}
设 $R$ 为诺特环，$I\subset R$ 为理想。以下三个范畴典范等价：
\begin{enumerate}
\item 令 $\mathcal{A}$ 为 $\mathbf{N}$ 上与 $R$-代数逆系统
$A_n=R/I^n$ 对应的 $R$-代数层。范畴 $\mathit{QC}(\mathcal{A})$。
\item 选取 $I$ 的生成元 $f_1,\ldots,f_r$。令 $\mathcal{B}$ 为
$\mathbf{N}$ 上与 $f_1^n,\ldots,f_r^n$ 上的 Koszul 代数逆系统对应的
微分分次 $R$-代数层。范畴 $\mathit{QC}(\mathcal{B})$。
\item 导出完备对象所成的全子范畴
$D_{comp}(R,I)\subset D(R)$；见《代数详论》，
定义 \ref{more-algebra-definition-derived-complete} 及其后的文字。
\end{enumerate}
\end{proposition}

\begin{proof}
考虑环化拓扑斯的显然态射
$f : (\Sh(\mathbf{N}), \mathcal{A}) \to (\Sh(pt), R)$
，并考虑伴随函子 $Lf^*$ 和 $Rf_*$。前者限制为函子
$$
F : D_{comp}(R, I) \longrightarrow \mathit{QC}(\mathcal{A})
$$
，它把由 K-平坦复形 $K^\bullet$ 表示的 $D_{comp}(R,I)$ 中对象 $K$ 送到
$\mathit{QC}(\mathcal{A})$ 的对象 $(K^\bullet\otimes_RR/I^n)$。后者限制为函子
$$
G : \mathit{QC}(\mathcal{A}) \longrightarrow D_{comp}(R, I)
$$
，它把 $\mathit{QC}(\mathcal{A})$ 的对象 $(M_n^\bullet)$ 送到
$R\lim M_n^\bullet$。例如，由《代数详论》，
引理 \ref{more-algebra-lemma-naive-derived-completion}，所得对象是导出完备的。
此外，由《代数详论》，
命题 \ref{more-algebra-proposition-noetherian-naive-completion-is-completion}，
有 $G\circ F=\text{id}$。因此，为看出 $F$ 和 $G$ 是互为拟逆的等价，只需看出
$G$ 的核为零（见《导出范畴》，
引理 \ref{derived-lemma-fully-faithful-adjoint-kernel-zero}）。
然而，直接证明这一点似乎并不容易！

\medskip\noindent
本段将证明 $\mathit{QC}(\mathcal{A})$ 和 $\mathit{QC}(\mathcal{B})$ 等价。
写 $\mathcal{B}=(B_n)$，其中 $B_n$ 是 Koszul 复形，并将它视为集中在次数
$-r,-r+1,\ldots,0$ 的上链复形。由《除幂代数》，
注 \ref{dpa-remark-pro-system-koszul}（但把链复形改成上链复形），可以找到
$1<n_1<n_2<\ldots$ 以及微分分次 $R$-代数的映射
$B_{n_i} \to E_i \to R/(f_1^{n_i}, \ldots, f_r^{n_i})$
和 $E_i\to B_{n_{i-1}}$，使得
$$
\xymatrix{
B_{n_1} \ar[d] &
B_{n_2} \ar[d] \ar[l] &
B_{n_3} \ar[d] \ar[l] &
\ldots \ar[l] \\
E_1 \ar[d] &
E_2 \ar[l] \ar[d] &
E_3 \ar[l] \ar[d] &
\ldots \ar[l] \\
B_1 &
B_{n_1} \ar[l] &
B_{n_2} \ar[l] &
\ldots \ar[l]
}
$$
是微分分次 $R$-代数的交换图，且
$E_i\to R/(f_1^{n_i},\ldots,f_r^{n_i})$ 是拟同构。由此得到：
\begin{enumerate}
\item $\mathit{QC}(\mathcal{B})$ 与 $\mathit{QC}((E_i))$ 之间有一个等价；
\item $\mathit{QC}((E_i))$ 与
$\mathit{QC}((R/(f_1^{n_i},\ldots,f_r^{n_i})))$ 之间有一个等价；
\item $\mathit{QC}((R/(f_1^{n_i},\ldots,f_r^{n_i})))$ 与
$\mathit{QC}(\mathcal{A})$ 之间有一个等价。
\end{enumerate}
事实上，对 (1)，可把引理 \ref{sdga-lemma-pro-isomorphic} 应用于上图；该图表明
$(E_i)$ 与 $(B_n)$ pro-同构。对 (2)，可把引理
\ref{sdga-lemma-cartesian-qis-equivalence} 应用于拟同构的逆系统
$E_i\to R/(f_1^{n_i},\ldots,f_r^{n_i})$。对 (3)，可应用引理
\ref{sdga-lemma-pro-isomorphic} 以及逆系统 $(R/I^n)$ 与
$(R/(f_1^{n_i},\ldots,f_r^{n_i})$ pro-同构这一初等事实。

\medskip\noindent
与证明的第一段完全一样，可以定义伴随函子\footnote{可以证明，通过上述等价，
这些函子与先前定义的 $F$ 和 $G$ 相容。}
$$
F' : D_{comp}(R, I) \longrightarrow \mathit{QC}(\mathcal{B})
\quad\text{及}\quad
G' : \mathit{QC}(\mathcal{B}) \longrightarrow D_{comp}(R, I).
$$
前者把由 K-平坦复形 $K^\bullet$ 表示的 $D_{comp}(R,I)$ 中对象 $K$ 送到
$\mathit{QC}(\mathcal{B})$ 的对象 $(K^\bullet\otimes_RB_n)$。后者把
$\mathit{QC}(\mathcal{B})$ 的对象 $(M_n)$ 送到 $R\lim M_n$。与上面一样论证，
只需证明 $G'$ 的核为零。因此，令 $\mathcal{M}=(M_n)$ 为 $\mathcal{B}$ 上的
良好微分分次模层，它表示 $\mathit{QC}(\mathcal{B})$ 中属于 $G'$ 之核的一个
对象。则
$$
0 = R\lim M_n \Rightarrow
0 = (R\lim M_n) \otimes_R^\mathbf{L} B_t =
R\lim (M_n \otimes_R^\mathbf{L} B_t)
$$
。由引理 \ref{sdga-lemma-different-tensors}，有
$R\lim (M_n \otimes_R^\mathbf{L} B_t) =
R\lim (M_n \otimes_{B_n}^\mathbf{L} B_t)$.
由于 $(M_n)$ 是 $\mathit{QC}(\mathcal{B})$ 的对象，可见逆系统
$M_n\otimes_{B_n}^\mathbf{L}B_t$ 最终为常值 $M_t$。所以如所需，$M_t=0$。
\end{proof}

\begin{remark}
\label{sdga-remark-proregular}
设 $R$ 为环，$f_1,\ldots,f_r\in R$ 为生成理想 $I$ 的元素序列。令 $K_n$ 为
$f_1^n,\ldots,f_r^n$ 上的 Koszul 复形，并将它视为微分分次 $R$-代数。
若对所有 $n$，都存在 $m>n$，使得 $K_m\to K_n$ 在除次数 $0$ 以外的
上同调上诱导零映射，则称 $f_1,\ldots,f_r$ 为{\it 弱 pro-正则序列}。
若是如此，则即使 $R$ 不是诺特环，命题 \ref{sdga-proposition-derived-complete}
证明中的论证仍然成立。特别地可见，只要 $I$ 可由弱 pro-正则序列生成，
$\mathit{QC}(\{R/I^n\})$ 作为 $R$-线性三角范畴就与导出完备对象所成的范畴
$D_{comp}(R,I)$ 等价。如有需要，我们将在此精确陈述并证明之。
\end{remark}
% Generated mechanically from the frozen hypercovering.tex target.
% Do not edit this generated file; edit the bound locale source instead.

% OK, start here.
%

\setcounter{chapter}{24}
\chapter{超覆盖}



\phantomsection
\label{hypercovering-section-phantom}


\section{引言}
\label{hypercovering-section-introduction}

\noindent
设 $\mathcal{C}$ 为位点，见《位点》定义 \ref{sites-definition-site}。
设 $X$ 为 $\mathcal{C}$ 的对象。
给定 $\mathcal{C}$ 上的阿贝尔层 $\mathcal{F}$，
我们希望计算其上同调群
$$
H^i(X, \mathcal{F}).
$$
按照一般定义（《位点上的上同调》第
\ref{sites-cohomology-section-cohomology-sheaves} 节），
可通过选取内射分辨
$
0 \to \mathcal{F} \to \mathcal{I}^0 \to \mathcal{I}^1 \to \ldots
$
并置
$$
H^i(X, \mathcal{F})
=
H^i(
\Gamma(X, \mathcal{I}^0) \to
\Gamma(X, \mathcal{I}^1) \to
\Gamma(X, \mathcal{I}^2)\to \ldots)
$$
本章的目标是证明，在 $\mathcal{C}$ 有纤维积的情形下，
也可以不选取内射分辨而计算这些上同调群。
为此我们将使用超覆盖。

\medskip\noindent
位点中的超覆盖是覆盖的推广；见
\cite[Expos\'e V, Sec. 7]{SGA4}。给定对象 $X$ 的超覆盖 $K$，
存在一个从 {\v C}ech 到上同调的谱序列，
它用该层在 $K$ 的各分量 $K_n$ 上的上同调来表示阿贝尔层
$\mathcal{F}$ 在 $X$ 上的上同调。事实表明，超覆盖总是足够多，
因而对所有超覆盖取余极限时，谱序列退化，且 $\mathcal{F}$
在 $X$ 上的上同调由 {\v C}ech 上同调群的余极限计算。

\medskip\noindent
还可考虑一种更一般的构造，即具有上同调下降的单纯增广；见
\cite[Expos\'e Vbis]{SGA4}。关于上同调下降，一份很好的手稿是
Brian Conrad 的文章，见
\url{https://math.stanford.edu/~conrad/papers/hypercover.pdf}。
我们将在关于单纯空间的章节中回到这些问题，并在那里证明，例如，
``局部紧''拓扑空间的真超覆盖具有上同调下降（《单纯空间》第
\ref{spaces-simplicial-section-proper-hypercovering} 节）。
我们的处理方法是把这一陈述归结为本章构造的从 {\v C}ech 到上同调的谱序列。






























\section{半可表对象}
\label{hypercovering-section-semi-representable}

\noindent
首先作如下定义。
字母 ``SR'' 表示半可表（Semi-Representable）。

\begin{definition}
\label{hypercovering-definition-SR}
设 $\mathcal{C}$ 为范畴。用 $\text{SR}(\mathcal{C})$
表示如下定义的{\it 半可表对象}范畴：
\begin{enumerate}
\item 对象是对象族 $\{U_i\}_{i \in I}$；
\item 态射 $\{U_i\}_{i \in I} \to \{V_j\}_{j \in J}$ 由映射
$\alpha : I \to J$ 以及对每个 $i \in I$ 给定的 $\mathcal{C}$ 中态射
$f_i : U_i \to V_{\alpha(i)}$ 构成。
\end{enumerate}
设 $X \in \Ob(\mathcal{C})$ 为 $\mathcal{C}$ 的对象。
{\it $X$ 上的半可表对象}范畴是范畴
$\text{SR}(\mathcal{C}, X) = \text{SR}(\mathcal{C}/X)$.
\end{definition}

\noindent
该定义实质上等价于
\cite[Expos\'e V, Subsection 7.3.0]{SGA4}。注意，这是一个``大''范畴。
稍后我们将``限制''构造 $X$ 的超覆盖所需指标集 $I$ 的大小，
届时便可重新定义 $\text{SR}(\mathcal{C}, X)$，使之成为一个范畴。
现明确写出 $\text{SR}(\mathcal{C}, X)$ 的对象与态射：
\begin{enumerate}
\item 对象是态射族 $\{U_i \to X\}_{i \in I}$；
\item 态射 $\{U_i \to X\}_{i \in I} \to
\{V_j \to X\}_{j \in J}$ 由映射 $\alpha : I \to J$
以及对每个 $i \in I$ 给定的 $X$ 上态射
$f_i : U_i \to V_{\alpha(i)}$ 构成。
\end{enumerate}
存在遗忘函子
$\text{SR}(\mathcal{C}, X) \to \text{SR}(\mathcal{C})$.

\begin{definition}
\label{hypercovering-definition-SR-F}
设 $\mathcal{C}$ 为范畴。
用 $F$ 表示{\it 把预层赋给半可表对象的函子}。以公式表示为
\begin{eqnarray*}
F : \text{SR}(\mathcal{C}) & \longrightarrow & \textit{PSh}(\mathcal{C}) \\
\{U_i\}_{i \in I} & \longmapsto & \amalg_{i\in I} h_{U_i}
\end{eqnarray*}
其中 $h_U$ 表示与对象 $U$ 相伴的可表预层。
\end{definition}

\noindent
给定态射 $U \to X$，得到可表预层的态射 $h_U \to h_X$。
因此，我们常把 $F$ 在 $\text{SR}(\mathcal{C}, X)$ 上的限制视为取值于
$h_X$ 上集合预层范畴、即 $\textit{PSh}(\mathcal{C})/h_X$ 的函子。
相应图示为：
$$
\xymatrix{
\text{SR}(\mathcal{C}, X) \ar[r]_F \ar[d] &
\textit{PSh}(\mathcal{C})/h_X \ar[d] \\
\text{SR}(\mathcal{C}) \ar[r]^F &
\textit{PSh}(\mathcal{C})
}
$$
下面讨论半可表对象范畴中极限的存在性。

\begin{lemma}
\label{hypercovering-lemma-coprod-prod-SR}
设 $\mathcal{C}$ 为范畴。
\begin{enumerate}
\item 范畴 $\text{SR}(\mathcal{C})$ 有余积，且 $F$ 与余积交换；
\item 函子 $F : \text{SR}(\mathcal{C}) \to \textit{PSh}(\mathcal{C})$
与极限交换；
\item 若 $\mathcal{C}$ 有纤维积，则 $\text{SR}(\mathcal{C})$ 有纤维积；
\item 若 $\mathcal{C}$ 有二元积，则 $\text{SR}(\mathcal{C})$ 有二元积；
\item 若 $\mathcal{C}$ 有等化子，则 $\text{SR}(\mathcal{C})$ 也有等化子；
\item 若 $\mathcal{C}$ 有终对象，则 $\text{SR}(\mathcal{C})$ 也有终对象。
\end{enumerate}
设 $X \in \Ob(\mathcal{C})$。
\begin{enumerate}
\item 范畴 $\text{SR}(\mathcal{C}, X)$ 有余积，且 $F$ 与余积交换；
\item 若 $\mathcal{C}$ 有纤维积，则 $\text{SR}(\mathcal{C}, X)$
有有限极限，且
$F : \text{SR}(\mathcal{C}, X) \to \textit{PSh}(\mathcal{C})/h_X$
与这些极限交换。
\end{enumerate}
\end{lemma}

\begin{proof}
先证明关于 $\text{SR}(\mathcal{C})$ 的各项结论。
证明 (1)。$\{U_i\}_{i \in I}$ 与 $\{V_j\}_{j \in J}$ 的余积是
$\{U_i\}_{i \in I} \amalg \{V_j\}_{j \in J}$；换言之，它是以
$I \amalg J$ 为指标集的对象族，对元素 $k \in I \amalg J$，若
$k = i \in I$ 则给出 $U_i$，若 $k = j \in J$ 则给出 $V_j$。
对象族之族的余积也类似。显然 $F$ 与这些余积交换。

\medskip\noindent
证明 (2)。对 $\Ob(\mathcal{C})$ 中的 $U$，考虑
$\text{SR}(\mathcal{C})$ 的对象 $\{U\}$。显然
$\Mor_{\text{SR}(\mathcal{C})}(\{U\}, K)) = F(K)(U)$
对 $K \in \Ob(\text{SR}(\mathcal{C}))$ 成立。由于预层的极限在截面层面计算
（《位点》第 \ref{sites-section-limits-colimits-PSh} 节），
故 $F$ 与极限交换。

\medskip\noindent
证明 (3)。设给定态射
$(\alpha, f_i) : \{U_i\}_{i \in I} \to \{V_j\}_{j \in J}$
以及态射
$(\beta, g_k) : \{W_k\}_{k \in K} \to \{V_j\}_{j \in J}$.
这些态射的纤维积由下式给出：
$$
\{ U_i \times_{f_i, V_j, g_k} W_k\}_{(i, j, k) \in I \times J \times K
\text{使得 } j = \alpha(i) = \beta(k)}
$$
若 $\mathcal{C}$ 有纤维积，则这些纤维积存在。

\medskip\noindent
证明 (4)。$\{U_i\}_{i \in I}$ 与 $\{V_j\}_{j \in J}$ 的积为
$\{U_i \times V_j\}_{i \in I, j \in J}$。若 $\mathcal{C}$ 有积，
则这些积存在。

\medskip\noindent
证明 (5)。两个映射
$(\alpha, f_i), (\alpha', f'_i) : \{U_i\}_{i \in I} \to \{V_j\}_{j \in J}$
的等化子为
$$
\{
\text{Eq}(f_i, f'_i : U_i \to V_{\alpha(i)})
\}_{i \in I,\ \alpha(i) = \alpha'(i)}
$$
若 $\mathcal{C}$ 有等化子，则这些等化子存在。

\medskip\noindent
证明 (6)。若 $X$ 是 $\mathcal{C}$ 的终对象，则 $\{X\}$ 是
$\text{SR}(\mathcal{C})$ 的终对象。

\medskip\noindent
下面证明关于 $\text{SR}(\mathcal{C}, X)$ 的陈述。
把上述结果应用于范畴 $\mathcal{C}/X$，并利用
$\text{SR}(\mathcal{C}/X) = \text{SR}(\mathcal{C}, X)$ 以及
$\textit{PSh}(\mathcal{C}/X) = \textit{PSh}(\mathcal{C})/h_X$
（把《位点》引理 \ref{sites-lemma-essential-image-j-shriek}
应用于赋予混沌拓扑的 $\mathcal{C}$），即可得到这些陈述。
不过我们也直接论证如下。显然，
$\{U_i \to X\}_{i \in I}$ 与 $\{V_j \to X\}_{j \in J}$ 的余积是
$\{U_i \to X\}_{i \in I} \amalg \{V_j \to X\}_{j \in J}$；
终点为 $X$ 的态射族之族的余积也类似。对象 $\{X \to X\}$ 是
$\text{SR}(\mathcal{C}, X)$ 的终对象。设给定态射
$(\alpha, f_i) : \{U_i \to X\}_{i \in I} \to \{V_j \to X\}_{j \in J}$
以及态射
$(\beta, g_k) : \{W_k \to X\}_{k \in K} \to \{V_j \to X\}_{j \in J}$.
这些态射的纤维积由下式给出：
$$
\{ U_i \times_{f_i, V_j, g_k} W_k \to X \}_{(i, j, k) \in I \times J \times K
\text{使得 } j = \alpha(i) = \beta(k)}
$$
由假设 $\mathcal{C}$ 有纤维积，这些纤维积存在。
因此 $\text{SR}(\mathcal{C}, X)$ 有有限极限；
见《范畴》引理 \ref{categories-lemma-finite-limits-exist}。
略去此情形下关于函子 $F$ 的陈述的验证。
\end{proof}





\section{超覆盖}
\label{hypercovering-section-hypercoverings}

\noindent
若再假设我们的范畴是位点，便可作如下定义。

\begin{definition}
\label{hypercovering-definition-covering-SR}
设 $\mathcal{C}$ 为位点。设
$f = (\alpha, f_i) : \{U_i\}_{i \in I} \to \{V_j\}_{j \in J}$
为范畴 $\text{SR}(\mathcal{C})$ 中的态射。若对每个 $j \in J$，
态射族 $\{U_i \to V_j\}_{i \in I, \alpha(i) = j}$
都是位点 $\mathcal{C}$ 的覆盖，则称 $f$ 为{\it 覆盖}。
设 $X$ 为 $\mathcal{C}$ 的对象。若 $\text{SR}(\mathcal{C}, X)$
中的态射 $K \to L$ 在 $\text{SR}(\mathcal{C})$ 中的像是覆盖，
则称它为{\it 覆盖}。
\end{definition}

\begin{lemma}
\label{hypercovering-lemma-covering-permanence}
设 $\mathcal{C}$ 为位点。
\begin{enumerate}
\item $\text{SR}(\mathcal{C})$ 中覆盖的复合仍为覆盖。
\item 若 $K \to L$ 是 $\text{SR}(\mathcal{C})$ 中的覆盖，
而 $L' \to L$ 是态射，则 $L' \times_L K$ 存在，且
$L' \times_L K \to L'$ 是覆盖。
\item 若 $\mathcal{C}$ 有二元积，并且 $A \to B$ 与 $K \to L$
是 $\text{SR}(\mathcal{C})$ 中的覆盖，则
$A \times K \to B \times L$ 是覆盖。
\end{enumerate}
设 $X \in \Ob(\mathcal{C})$。则 (1) 与 (2) 对
$\text{SR}(\mathcal{C}, X)$ 成立；若 $\mathcal{C}$ 有纤维积，
则 (3) 也成立。
\end{lemma}

\begin{proof}
(1) 由位点的公理立即得到。(2) 由引理
\ref{hypercovering-lemma-coprod-prod-SR} 的证明中 $\text{SR}(\mathcal{C})$
内纤维积的构造，以及 $\mathcal{C}$ 的覆盖中各态射须为可表态射的要求得到。
把 $A \times K \to B \times L$ 看作复合
$A \times K \to B \times K \to B \times L$，它因而是覆盖的基变换的复合，
故得 (3)。最后一个陈述由
$\text{SR}(\mathcal{C}, X) = \text{SR}(\mathcal{C}/X)$ 得出。
\end{proof}

\noindent
由引理 \ref{hypercovering-lemma-coprod-prod-SR} 与《单纯方法》引理
\ref{simplicial-lemma-existence-cosk}，若 $\mathcal{C}$ 有纤维积，
则 $\text{SR}(\mathcal{C}, X)$ 的截断单纯对象的余骨架存在。
因此下述定义是有意义的。

\begin{definition}
\label{hypercovering-definition-hypercovering}
设 $\mathcal{C}$ 为有纤维积的位点，且 $X \in \Ob(\mathcal{C})$。
{\it $X$ 的超覆盖}是 $\text{SR}(\mathcal{C}, X)$ 的单纯对象 $K$，满足：
\begin{enumerate}
\item 对象 $K_0$ 是位点 $\mathcal{C}$ 中 $X$ 的覆盖；
\item 对每个 $n \geq 0$，典范态射
$$
K_{n + 1} \longrightarrow (\text{cosk}_n \text{sk}_n K)_{n + 1}
$$
是上述意义下的覆盖。
\end{enumerate}
\end{definition}

\noindent
条件 (1) 是有意义的，因为 $\text{SR}(\mathcal{C}, X)$ 的每个对象
毕竟都是终点为 $X$ 的态射族。它也可以表述为：从 $K_0$ 到
$\text{SR}(\mathcal{C}, X)$ 的终对象的态射是覆盖。

\begin{example}[{\v C}ech 超覆盖]
\label{hypercovering-example-cech}
设 $\mathcal{C}$ 为有纤维积的位点，且
$\{U_i \to X\}_{i \in I}$ 是 $\mathcal{C}$ 的覆盖。置
$K_0 = \{U_i \to X\}_{i \in I}$。则 $K_0$ 是
$\text{SR}(\mathcal{C}, X)$ 的 $0$-截断单纯对象，故可构造
$$
K = \text{cosk}_0 K_0.
$$
显然 $K$ 满足定义 \ref{hypercovering-definition-hypercovering} 的条件 (1)。
由《单纯方法》引理 \ref{simplicial-lemma-cosk-up}，所有态射
$K_{n + 1} \to (\text{cosk}_n \text{sk}_n K)_{n + 1}$ 都是同构，
故它也满足条件 (2)。注意，各项 $K_n$ 正是通常的
$$
K_n = \{
U_{i_0} \times_X U_{i_1} \times_X \ldots \times_X U_{i_n} \to X
\}_{(i_0, i_1, \ldots, i_n) \in I^{n + 1}}
$$
这种形式的 $X$ 的超覆盖称为 $X$ 的{\it {\v C}ech 超覆盖}。
\end{example}

\begin{example}[由位点的单纯对象给出的超覆盖]
\label{hypercovering-example-hypercovering-in-C}
设 $\mathcal{C}$ 为有纤维积的位点，$X \in \Ob(\mathcal{C})$，
并设 $U$ 为 $\mathcal{C}$ 的单纯对象。仍记 $U_n = U([n])$。
最后，假设给定增广
$$
a : U \to X
$$
在此情形下，可考虑 $\text{SR}(\mathcal{C}, X)$ 的单纯对象 $K$，
其各项为 $K_n = \{U_n \to X\}$。当且仅当下列三个条件成立时，
$K$ 是定义 \ref{hypercovering-definition-hypercovering} 意义下的 $X$ 的超覆盖
\footnote{由于 $\mathcal{C}$ 有纤维积，范畴 $\mathcal{C}/X$
有全部有限极限。因此，由《单纯方法》引理
\ref{simplicial-lemma-existence-cosk}，所需余骨架存在。}：
\begin{enumerate}
\item $\{U_0 \to X\}$ 是 $\mathcal{C}$ 的覆盖；
\item $\{U_1 \to U_0 \times_X U_0\}$ 是 $\mathcal{C}$ 的覆盖；
\item $\{U_{n + 1} \to (\text{cosk}_n\text{sk}_n U)_{n + 1}\}$
对 $n \geq 1$ 是 $\mathcal{C}$ 的覆盖。
\end{enumerate}
略去直接的验证。
\end{example}

\begin{example}[与一个覆盖相伴的 {\v C}ech 超覆盖]
\label{hypercovering-example-cech-cover}
设 $\mathcal{C}$ 为有纤维积的位点。设 $U \to X$ 是
$\mathcal{C}$ 的态射，使得 $\{U \to X\}$ 是 $\mathcal{C}$ 的覆盖
\footnote{$\mathcal{C}$ 中具有此性质的态射有时称为``覆盖态射''。}。
考虑 $\text{SR}(\mathcal{C}, X)$ 的单纯对象 $K$，其各项为
$$
K_n = \{U \times_X U \times_X \ldots \times_X U \to X\}
\quad (n + 1 \text{ 个因子})
$$
则 $K$ 是 $X$ 的超覆盖。该例同时是例 \ref{hypercovering-example-cech}
与例 \ref{hypercovering-example-hypercovering-in-C} 的特殊情形。
\end{example}

\begin{lemma}
\label{hypercovering-lemma-hypercoverings-set}
设 $\mathcal{C}$ 为有纤维积的位点，且 $X \in \Ob(\mathcal{C})$。
$X$ 的所有超覆盖所成的类是集合。
\end{lemma}

\begin{proof}
由于 $\mathcal{C}$ 是位点，$X$ 的所有覆盖所成的类是集合。
因此，可取的 $K_0$ 所成的类是集合。假设已经证明可取的
$K_0, \ldots, K_n$ 所成的类是集合。于是只需证明，给定
$K_0, \ldots, K_n$ 后，可取的 $K_{n + 1}$ 所成的类是集合。
这显然成立，因为 $K_{n + 1}$ 必须从
$(\text{cosk}_n \text{sk}_n K)_{n + 1}$ 的所有可能覆盖中选取。
\end{proof}

\begin{remark}
\label{hypercovering-remark-hypercoverings-really-set}
该引理不只是说，存在由超覆盖的选择组成且为集合的共尾系统；
它断言超覆盖本身确实组成一个集合。
\end{remark}

\noindent
$\mathcal{C}$ 上的预层范畴有有限（余）极限。因此，
集合预层上存在函子 $\text{cosk}_n$。

\begin{lemma}
\label{hypercovering-lemma-hypercovering-F}
设 $\mathcal{C}$ 为有纤维积的位点，$X \in \Ob(\mathcal{C})$，
并设 $K$ 为 $X$ 的超覆盖。考虑 $\textit{PSh}(\mathcal{C})$
的单纯对象 $F(K)$，并赋予它到常值单纯预层 $h_X$ 的增广。
\begin{enumerate}
\item 预层态射 $F(K)_0 \to h_X$ 在层化后成为满射。
\item 态射
$$
(d^1_0, d^1_1) :
F(K)_1
\longrightarrow
F(K)_0 \times_{h_X} F(K)_0
$$
在层化后成为满射。
\item 对每个 $n \geq 1$，态射
$$
F(K)_{n + 1} \longrightarrow (\text{cosk}_n \text{sk}_n F(K))_{n + 1}
$$
在层化后成为满射。
\end{enumerate}
\end{lemma}

\begin{proof}
我们将使用如下事实：若 $\{U_i \to U\}_{i \in I}$
是位点 $\mathcal{C}$ 的覆盖，则态射
$$
\amalg_{i \in I} h_{U_i} \to h_U
$$
在层化后成为满射；见《位点》引理
\ref{sites-lemma-covering-surjective-after-sheafification}。
故第一个断言立即成立。

\medskip\noindent
对第二个断言，由《单纯方法》例 \ref{simplicial-example-cosk0}，
单纯对象 $\text{cosk}_0 \text{sk}_0 K$ 的各项为
$K_0 \times \ldots \times K_0$。因此，由超覆盖的定义，
$(d^1_0, d^1_1) : K_1 \to K_0 \times K_0$ 是覆盖。
结合上述断言以及 $F$ 把积变为 $h_X$ 上的纤维积这一事实，便得到 (2)。

\medskip\noindent
对第三个断言，我们声称
$\text{cosk}_n \text{sk}_n F(K) =
F(\text{cosk}_n \text{sk}_n K)$
对 $n \geq 1$ 成立。为证明这一点，暂用 $F'$ 表示函子
$\text{SR}(\mathcal{C}, X) \to \textit{PSh}(\mathcal{C})/h_X$.
由引理 \ref{hypercovering-lemma-coprod-prod-SR}，函子 $F'$ 与有限极限交换。
根据《单纯方法》第 \ref{simplicial-section-skeleton} 节中
对 $\text{cosk}_n$ 函子的描述，有
$\text{cosk}_n \text{sk}_n F'(K) =
F'(\text{cosk}_n \text{sk}_n K)$.
回忆《单纯方法》引理 \ref{simplicial-lemma-existence-cosk}
中描述 $(\text{cosk}_n U)_m$ 时所用的范畴是
$(\Delta/[m])^{opp}_{\leq n}$。一个有趣的练习是证明，只要
$n \geq 1$，范畴 $(\Delta/[m])_{\leq n}$ 就是连通的（见《范畴》定义
\ref{categories-definition-category-connected}）。因此，由《范畴》引理
\ref{categories-lemma-connected-limit-over-X}，有
$\text{cosk}_n \text{sk}_n F'(K) =
\text{cosk}_n \text{sk}_n F(K)$，从而该断言成立。
由此得到性质 (2)，因为此时可见，(2) 中的态射是把函子 $F$
应用于定义 \ref{hypercovering-definition-covering-SR} 意义下的覆盖所得；
结论遂由本证明开头提到的第一个事实得出。
\end{proof}



\section{无环性}
\label{hypercovering-section-acyclicity}

\noindent
设 $\mathcal{C}$ 为位点。对集合预层 $\mathcal{F}$，用
$\mathbf{Z}_\mathcal{F}$ 表示由下述规则定义的阿贝尔群预层：
$$
\mathbf{Z}_\mathcal{F}(U) = \text{由 }\mathcal{F}(U)\text{ 生成的自由阿贝尔群}.
$$
有时称它为{\it $\mathcal{F}$ 上的自由阿贝尔预层}。
当然，构造 $\mathcal{F} \mapsto \mathbf{Z}_\mathcal{F}$ 是函子，
且是遗忘函子
$\textit{PAb}(\mathcal{C}) \to \textit{PSh}(\mathcal{C})$.
的左伴随。当然，层化 $\mathbf{Z}_\mathcal{F}^\#$ 是阿贝尔群层，
而函子 $\mathcal{F} \mapsto \mathbf{Z}_\mathcal{F}^\#$ 也为左伴随。
有时称 $\mathbf{Z}_\mathcal{F}^\#$ 为
{\it $\mathcal{F}$ 上的自由阿贝尔层}。

\medskip\noindent
对位点 $\mathcal{C}$ 的对象 $X$，用 $\mathbf{Z}_X$ 表示
$h_X$ 上的自由阿贝尔预层，并用 $\mathbf{Z}_X^\#$ 表示其层化。

\begin{definition}
\label{hypercovering-definition-homology}
设 $\mathcal{C}$ 为位点，$K$ 为 $\textit{PSh}(\mathcal{C})$
的单纯对象。由上可得 $\textit{Ab}(\mathcal{C})$ 的单纯对象
$\mathbf{Z}_K^\#$。可取其相伴的阿贝尔预层复形
$s(\mathbf{Z}_K^\#)$；见《单纯方法》第
\ref{simplicial-section-complexes} 节。$K$ 的{\it 同调}是
阿贝尔层复形 $s(\mathbf{Z}_K^\#)$ 的同调。
\end{definition}

\noindent
换言之，$K$ 的{\it 第 $i$ 个同调 $H_i(K)$}是阿贝尔群层
$H_i(K) = H_i(s(\mathbf{Z}_K^\#))$。本节研究 $K$ 是
$\mathcal{C}$ 的对象 $X$ 的超覆盖时的同调。

\begin{lemma}
\label{hypercovering-lemma-compare-cosk0}
设 $\mathcal{C}$ 为位点，$\mathcal{F} \to \mathcal{G}$
为集合预层的态射。用 $K$ 表示 $\textit{PSh}(\mathcal{C})$
的单纯对象，其第 $n$ 项是 $\mathcal{F}$ 在 $\mathcal{G}$ 上的
$(n + 1)$ 重纤维积；见《单纯方法》例
\ref{simplicial-example-fibre-products-simplicial-object}。
若 $\mathcal{F} \to \mathcal{G}$ 在层化后为满射，则
$$
H_i(K) =
\left\{
\begin{matrix}
0 & \text{若} & i > 0\\
\mathbf{Z}_\mathcal{G}^\# & \text{若} & i = 0
\end{matrix}
\right.
$$
次数 $0$ 上的同构由态射
$H_0(K) \to \mathbf{Z}_\mathcal{G}^\#$ 给出；该态射来自映射
$(\mathbf{Z}_K^\#)_0 =
\mathbf{Z}_\mathcal{F}^\# \to \mathbf{Z}_\mathcal{G}^\#$。
\end{lemma}

\begin{proof}
令 $\mathcal{G}' \subset \mathcal{G}$ 为态射
$\mathcal{F} \to \mathcal{G}$ 的像。取 $U \in \Ob(\mathcal{C})$，置
$A = \mathcal{F}(U)$ 和 $B = \mathcal{G}'(U)$。
则单纯集 $K(U)$ 等于第 $n$ 个单形由下式给出的单纯集：
$$
A \times_B A \times_B \ldots \times_B A\ (n + 1 \text{ 个因子)}.
$$
由《单纯方法》引理 \ref{simplicial-lemma-cosk-minus-one-equivalence}，
态射 $K(U) \to B$ 是平凡 Kan 纤维化，因而是同伦等价
（《单纯方法》引理 \ref{simplicial-lemma-trivial-kan-homotopy}）。
对它应用``生成自由阿贝尔群''函子，得到
$$
\mathbf{Z}_K(U) \longrightarrow \mathbf{Z}_B
$$
是同伦等价。注意，$s(\mathbf{Z}_B)$ 是复形
$$
\ldots \to
\bigoplus\nolimits_{b \in B}\mathbf{Z} \xrightarrow{0}
\bigoplus\nolimits_{b \in B}\mathbf{Z} \xrightarrow{1}
\bigoplus\nolimits_{b \in B}\mathbf{Z} \xrightarrow{0}
\bigoplus\nolimits_{b \in B}\mathbf{Z} \to 0
$$
；见《单纯方法》引理 \ref{simplicial-lemma-homology-eilenberg-maclane}。
因此，对 $i > 0$ 有
$H_i(s(\mathbf{Z}_K(U))) = 0$，且
$H_0(s(\mathbf{Z}_K(U))) = \bigoplus_{b \in B}\mathbf{Z}
= \bigoplus_{s \in \mathcal{G}'(U)} \mathbf{Z}$.
这些等同与限制映射相容。

\medskip\noindent
由此推出 $H_i(s(\mathbf{Z}_K)) = 0$ 对 $i > 0$ 成立，且
$H_0(s(\mathbf{Z}_K)) = \mathbf{Z}_{\mathcal{G}'}$；这里在
$\textit{PAb}(\mathcal{C})$ 中计算同调群。由于层化是正合函子，
便得到引理的结论。确切地说，正合性蕴含
$H_0(s(\mathbf{Z}_K))^\# = H_0(s(\mathbf{Z}_K^\#))$，
其他指标也同样。
\end{proof}

\begin{lemma}
\label{hypercovering-lemma-acyclicity}
设 $\mathcal{C}$ 为位点，$f : L \to K$ 为
$\textit{PSh}(\mathcal{C})$ 的单纯对象之间的态射，
且 $n \geq 0$ 为整数。假设：
\begin{enumerate}
\item 对 $i < n$，态射 $L_i \to K_i$ 是同构；
\item 态射 $L_n \to K_n$ 在层化后为满射；
\item 典范映射 $L \to \text{cosk}_n \text{sk}_n L$ 是同构；
\item 典范映射 $K \to \text{cosk}_n \text{sk}_n K$ 是同构。
\end{enumerate}
则 $H_i(f) : H_i(L) \to H_i(K)$ 是同构。
\end{lemma}

\begin{proof}
该证明与上面引理 \ref{hypercovering-lemma-compare-cosk0} 的证明完全相同。
确切地说，先令 $K_n' \subset K_n$ 为映射 $L_n \to K_n$
之像所成的子预层。假设 (2) 意味着 $K_n'$ 的层化等于 $K_n$ 的层化。
此外，由于对所有 $i < n$ 都有 $L_i = K_i$，取
$U_0 = L_0 = K_0, \ldots, U_{n - 1} = L_{n - 1} = K_{n - 1}, U_n = K'_n$.
便得到 $n$-截断单纯预层 $U$。记单纯预层
$K' = \text{cosk}_n U$。由于可将 $K'_m$ 构造为有限极限，
而层化又正合，故 $(K'_m)^\# = K_m$。换言之，
$(K')^\# = K^\#$。再次由层化的正合性，推出
$H_i(K) = H_i(K')$。因此，只需对态射 $L \to K'$ 证明该引理；
也就是说，可以假设 $L_n \to K_n$ 是{\it 预层}的满射！

\medskip\noindent
在此情形下，对 $\mathcal{C}$ 的任意对象 $U$，单纯集的态射
$$
L(U) \longrightarrow K(U)
$$
满足《单纯方法》引理 \ref{simplicial-lemma-section} 的全部假设。
故它是平凡 Kan 纤维化，尤其是同伦等价
（《单纯方法》引理 \ref{simplicial-lemma-trivial-kan-homotopy}）。因此
$$
\mathbf{Z}_L(U) \longrightarrow \mathbf{Z}_K(U)
$$
也是同伦等价。这对所有 $U$ 都成立，结论随之得出。
\end{proof}

\begin{lemma}
\label{hypercovering-lemma-acyclic-hypercover-sheaves}
设 $\mathcal{C}$ 为位点，$K$ 为单纯预层，$\mathcal{G}$ 为预层，
并设 $K \to \mathcal{G}$ 为 $K$ 到 $\mathcal{G}$ 的增广。假设：
\begin{enumerate}
\item 预层态射 $K_0 \to \mathcal{G}$ 在层化后成为满射；
\item 态射
$$
(d^1_0, d^1_1) :
K_1
\longrightarrow
K_0 \times_\mathcal{G} K_0
$$
在层化后成为满射；
\item 对每个 $n \geq 1$，态射
$$
K_{n + 1} \longrightarrow (\text{cosk}_n \text{sk}_n K)_{n + 1}
$$
在层化后成为满射。
\end{enumerate}
则对 $i > 0$ 有 $H_i(K) = 0$，且
$H_0(K) = \mathbf{Z}_\mathcal{G}^\#$。
\end{lemma}

\begin{proof}
对 $n \geq 1$，记 $K^n = \text{cosk}_n \text{sk}_n K$。
把 $K^0$ 定义为这样的单纯对象：$(K^0)_n$ 等于 $(n + 1)$ 重纤维积
$K_0 \times_\mathcal{G} \ldots \times_\mathcal{G} K_0$；
见《单纯方法》例 \ref{simplicial-example-fibre-products-simplicial-object}。
存在态射
$$
K \longrightarrow \ldots \to K^n \to K^{n - 1} \to \ldots \to K^1 \to K^0.
$$
态射 $K \to K^i$ 以及 $j \geq i \geq 1$ 时的 $K^j \to K^i$
来自各函子 $\text{cosk}_n$ 的泛性质。态射 $K^1 \to K^0$
是《单纯方法》评注 \ref{simplicial-remark-augmentation} 中的典范态射。
还回忆，$K^0 \to \text{cosk}_1 \text{sk}_1 K^0$ 是同构；
见《单纯方法》引理 \ref{simplicial-lemma-cosk-minus-one}。

\medskip\noindent
由引理 \ref{hypercovering-lemma-compare-cosk0}，对 $i > 0$ 有
$H_i(K^0) = 0$，且 $H_0(K^0) = \mathbf{Z}_\mathcal{G}^\#$。

\medskip\noindent
取 $n \geq 1$，考虑态射 $K^n \to K^{n - 1}$。
它在次数 $< n$ 的各项上是同构。注意，
$K^n \to \text{cosk}_n \text{sk}_n K^n$ 与
$K^{n - 1} \to \text{cosk}_n \text{sk}_n K^{n - 1}$ 都是同构；又有
$(K^n)_n = K_n$ 及
$(K^{n - 1})_n = (\text{cosk}_{n - 1} \text{sk}_{n - 1} K)_n$。
因此，由假设，$(K^n)_n \to (K^{n - 1})_n$ 是层化后成为满射的预层态射。
由引理 \ref{hypercovering-lemma-acyclicity}，得到
$H_i(K^n) = H_i(K^{n - 1})$。结合上文即证明该引理。
\end{proof}

\begin{lemma}
\label{hypercovering-lemma-hypercovering-acyclic}
设 $\mathcal{C}$ 为有纤维积的位点，$X$ 为 $\mathcal{C}$ 的对象，
并设 $K$ 为 $X$ 的超覆盖。单纯预层 $F(K)$ 的同调在次数 $> 0$
时为 $0$，在次数 $0$ 时等于 $\mathbf{Z}_X^\#$。
\end{lemma}

\begin{proof}
合用引理 \ref{hypercovering-lemma-acyclic-hypercover-sheaves}
与 \ref{hypercovering-lemma-hypercovering-F}。
\end{proof}












\section{{\v C}ech 上同调与超覆盖}
\label{hypercovering-section-hyper-cech}

\noindent
设 $\mathcal{C}$ 为位点。考虑位点 $\mathcal{C}$ 上的阿贝尔群预层
$\mathcal{F}$。它定义函子
\begin{eqnarray*}
\mathcal{F} : \text{SR}(\mathcal{C})^{opp}
& \longrightarrow &
\textit{Ab} \\
\{U_i\}_{i \in I} &
\longmapsto &
\prod\nolimits_{i \in I} \mathcal{F}(U_i)
\end{eqnarray*}
于是，$\text{SR}(\mathcal{C})$ 的单纯对象 $K$ 被变为
$\textit{Ab}$ 的余单纯对象 $\mathcal{F}(K)$。与 $\mathcal{F}(K)$
相伴的上链复形 $s(\mathcal{F}(K))$（《单纯方法》第
\ref{simplicial-section-dold-kan-cosimplicial} 节）称为
$\mathcal{F}$ 关于单纯对象 $K$ 的 {\v C}ech 复形。置
$$
\check{H}^i(K, \mathcal{F})
=
H^i(s(\mathcal{F}(K))).
$$
并称之为 $\mathcal{F}$ 关于 $K$ 的第 $i$ 个 {\v C}ech 上同调群。
本节证明《层上同调》第 \ref{cohomology-section-cech}、
\ref{cohomology-section-cech-functor} 与
\ref{cohomology-section-cech-cohomology-cohomology} 节中
关于开覆盖的 {\v C}ech 上同调若干结果的类似物。

\begin{lemma}
\label{hypercovering-lemma-h0-cech}
设 $\mathcal{C}$ 为有纤维积的位点，$X$ 为 $\mathcal{C}$ 的对象，
$K$ 为 $X$ 的超覆盖，且 $\mathcal{F}$ 为 $\mathcal{C}$ 上的阿贝尔群层。
则 $\check{H}^0(K, \mathcal{F}) = \mathcal{F}(X)$。
\end{lemma}

\begin{proof}
有
$$
\check{H}^0(K, \mathcal{F})
=
\Ker(\mathcal{F}(K_0) \longrightarrow \mathcal{F}(K_1))
$$
写 $K_0 = \{U_i \to X\}$。它是位点 $\mathcal{C}$ 中的覆盖。
此外，$K_1 \to K_0 \times K_0$ 是
$\text{SR}(\mathcal{C}, X)$ 中的覆盖。因此可写
$K_1 = \amalg_{i_0, i_1 \in I} \{V_{i_0i_1j} \to X\}$，
使态射 $K_1 \to K_0 \times K_0$ 由位点 $\mathcal{C}$ 的覆盖
$\{V_{i_0i_1j} \to U_{i_0} \times_X U_{i_1}\}$ 给出。
从而还可把
$$
\check{H}^0(K, \mathcal{F})
=
\Ker(
\prod\nolimits_i \mathcal{F}(U_i)
\longrightarrow
\prod\nolimits_{i_0i_1 j} \mathcal{F}(V_{i_0i_1j})
)
$$
按显然的映射作此等同。$\mathcal{F}$ 的层性质蕴含
$\check{H}^0(K, \mathcal{F}) = H^0(X, \mathcal{F})$。
\end{proof}

\noindent
事实上，该性质当然刻画了 $\mathcal{C}$ 上所有阿贝尔预层中的阿贝尔层。
在此情形下，《层上同调》引理 \ref{hypercovering-lemma-injective-trivial-cech}
的类似物如下。

\begin{lemma}
\label{hypercovering-lemma-injective-trivial-cech}
设 $\mathcal{C}$ 是具有纤维积的位点，$X$ 是 $\mathcal{C}$ 的对象，
$K$ 是 $X$ 的超覆盖，$\mathcal{I}$ 是 $\mathcal{C}$ 上的内射阿贝尔群层。则
$$
\check{H}^p(K, \mathcal{I}) =
\left\{
\begin{matrix}
\mathcal{I}(X) & \text{若} & p = 0 \\
0 & \text{若} & p > 0
\end{matrix}
\right.
$$
\end{lemma}

\begin{proof}
注意，对 $\text{SR}(\mathcal{C}, X)$ 的任意对象
$Z = \{U_i \to X\}$ 以及 $\mathcal{C}$ 上的任意阿贝尔层
$\mathcal{F}$，都有
\begin{eqnarray*}
\mathcal{F}(Z)
& = &
\prod \mathcal{F}(U_i) \\
& = &
\prod \Mor_{\textit{PSh}(\mathcal{C})}(h_{U_i}, \mathcal{F})\\
& = &
\Mor_{\textit{PSh}(\mathcal{C})}(F(Z), \mathcal{F})\\
& = &
\Mor_{\textit{PAb}(\mathcal{C})}(\mathbf{Z}_{F(Z)}, \mathcal{F}) \\
& = &
\Mor_{\textit{Ab}(\mathcal{C})}(\mathbf{Z}_{F(Z)}^\#, \mathcal{F})
\end{eqnarray*}
由此可见，对 $\text{SR}(\mathcal{C}, X)$ 的任意单纯对象 $K$，都有
\begin{equation}
\label{hypercovering-equation-identify-cech}
s(\mathcal{F}(K))
=
\Hom_{\textit{Ab}(\mathcal{C})}(s(\mathbf{Z}_{F(K)}^\#), \mathcal{F})
\end{equation}
记号见定义 \ref{hypercovering-definition-homology}。
若 $K$ 是超覆盖，则层复形 $s(\mathbf{Z}_{F(K)}^\#)$
拟同构于 $\mathbf{Z}_X^\#$，见引理 \ref{hypercovering-lemma-hypercovering-acyclic}。
因此，若 $\mathcal{I}$ 是内射阿贝尔层而 $K$ 是超覆盖，
则复形 $s(\mathcal{I}(K))$ 除了次数 $0$ 外是无环的。换言之，
$$
\check{H}^i(K, \mathcal{I}) = 0
$$
对 $i > 0$ 成立。结合引理 \ref{hypercovering-lemma-h0-cech} 即得结论。
\end{proof}

\noindent
下面给出《位点上的上同调》引理
\ref{sites-cohomology-lemma-cech-spectral-sequence} 的类似物。
设 $\mathcal{C}$ 是位点，$\mathcal{F}$ 是 $\mathcal{C}$ 上的阿贝尔群层。
回忆，$\underline{H}^i(\mathcal{F})$ 表示 $\mathcal{C}$ 上由规则
$\underline{H}^i(\mathcal{F}) : U \longmapsto H^i(U, \mathcal{F})$
定义的阿贝尔群预层。依照本节引言中的做法，把它扩张到
$\text{SR}(\mathcal{C})$ 上。

\begin{lemma}
\label{hypercovering-lemma-cech-spectral-sequence}
设 $\mathcal{C}$ 是具有纤维积的位点，$X$ 是 $\mathcal{C}$ 的对象，
$K$ 是 $X$ 的超覆盖，$\mathcal{F}$ 是 $\mathcal{C}$ 上的阿贝尔群层。
在 $D^{+}(\textit{Ab})$ 中存在关于 $\mathcal{F}$ 函子性的态射
$$
s(\mathcal{F}(K))
\longrightarrow
R\Gamma(X, \mathcal{F})
$$
它诱导函子 $\textit{Ab}(\mathcal{C}) \to \textit{Ab}$ 之间的自然变换
$$
\check{H}^i(K, -) \longrightarrow H^i(X, -)
$$
此外，存在谱序列 $(E_r, d_r)_{r \geq 0}$，满足
$$
E_2^{p, q} = \check{H}^p(K, \underline{H}^q(\mathcal{F}))
$$
并收敛到 $H^{p + q}(X, \mathcal{F})$。
此谱序列关于 $\mathcal{F}$ 和超覆盖 $K$ 都是函子性的。
\end{lemma}

\begin{proof}
可以采用《上同调》一章相应引理中的同一方法证明。这里改用双复形论证。

\medskip\noindent
在 $\mathcal{C}$ 上阿贝尔层的范畴中选取内射分辨
$\mathcal{F} \to \mathcal{I}^\bullet$。考虑各项为
$$
A^{p, q} = \mathcal{I}^q(K_p)
$$
的双复形 $A^{\bullet, \bullet}$，其中微分
$d_1^{p, q} : A^{p, q} \to A^{p + 1, q}$ 来自与余单纯阿贝尔群
$\mathcal{I}^p(K)$ 相伴的复形 $s(\mathcal{I}^q(K))$ 上的微分，
而微分 $d_2^{p, q} : A^{p, q} \to A^{p, q + 1}$ 来自微分
$\mathcal{I}^q \to \mathcal{I}^{q + 1}$。
用 $\text{Tot}(A^{\bullet, \bullet})$ 表示与双复形
$A^{\bullet, \bullet}$ 相伴的全复形，见《同调代数》第
\ref{homology-section-double-complexes} 节。我们将使用与此双复形相伴的
两个谱序列 $({}'E_r, {}'d_r)$ 和 $({}''E_r, {}''d_r)$，见
《同调代数》第 \ref{homology-section-double-complex} 节。

\medskip\noindent
由引理 \ref{hypercovering-lemma-injective-trivial-cech}，复形
$s(\mathcal{I}^q(K))$ 在正次数无环，且其 $H^0$ 等于
$\mathcal{I}^q(X)$。因此，由《同调代数》引理
\ref{homology-lemma-double-complex-gives-resolution}，自然映射
$$
\mathcal{I}^\bullet(X) \longrightarrow \text{Tot}(A^{\bullet, \bullet})
$$
是阿贝尔群复形的拟同构。特别地，
$H^n(\text{Tot}(A^{\bullet, \bullet})) = H^n(X, \mathcal{F})$。

\medskip\noindent
引理中的态射 $s(\mathcal{F}(K)) \longrightarrow R\Gamma(X, \mathcal{F})$
是态射 $s(\mathcal{F}(K)) \to \text{Tot}(A^{\bullet, \bullet})$
与上述拟同构之逆的复合。这是可行的，因为 $\mathcal{I}^\bullet(X)$
是 $R\Gamma(X, \mathcal{F})$ 的一个代表元。

\medskip\noindent
考虑谱序列 $({}'E_r, {}'d_r)_{r \geq 0}$。由《同调代数》引理
\ref{homology-lemma-ss-double-complex}，有
$$
{}'E_2^{p, q} = H^p_I(H^q_{II}(A^{\bullet, \bullet}))
$$
换言之，先对 $d_2$ 取上同调，得到群
${}'E_1^{p, q} = \underline{H}^q(\mathcal{F})(K_p)$。
因此，由微分 ${}'d_1$ 的描述，确有
${}'E_2^{p, q} = \check{H}^p(K, \underline{H}^q(\mathcal{F}))$。
由以上所述和《同调代数》引理
\ref{homology-lemma-first-quadrant-ss}，可见它按要求收敛到
$H^n(X, \mathcal{F})$。

\medskip\noindent
关于上述构造对阿贝尔层 $\mathcal{F}$ 和超覆盖 $K$ 的函子性，
其证明从略。
\end{proof}









\section{Verdier 意义下的超覆盖}
\label{hypercovering-section-hypercoverings-verdier}

\noindent
敏锐的读者会注意到，为了得到对象 $X$ 的超覆盖所对应的从
{\v C}ech 上同调到上同调的谱序列，我们所需的只是引理
\ref{hypercovering-lemma-hypercovering-F} 的结论。因此下面的定义是合理的。

\begin{definition}
\label{hypercovering-definition-hypercovering-variant}
设 $\mathcal{C}$ 是位点，并假设 $\mathcal{C}$ 具有等化子和纤维积。
设 $\mathcal{G}$ 是集合预层。$\mathcal{G}$ 的一个{\it 超覆盖}，是
$\text{SR}(\mathcal{C})$ 的单纯对象 $K$ 连同一个增广
$F(K) \to \mathcal{G}$，使得
\begin{enumerate}
\item $F(K_0) \to \mathcal{G}$ 在层化后成为满射；
\item $F(K_1) \to F(K_0) \times_\mathcal{G} F(K_0)$
在层化后成为满射；并且
\item $F(K_{n + 1}) \longrightarrow F((\text{cosk}_n \text{sk}_n K)_{n + 1})$
对 $n \geq 1$ 在层化后成为满射。
\end{enumerate}
若 $\text{SR}(\mathcal{C})$ 的单纯对象 $K$ 是
$\textit{PSh}(\mathcal{C})$ 的终对象 $*$ 的超覆盖，则称 $K$ 为一个
{\it 超覆盖}。
\end{definition}

\noindent
假设 $\mathcal{C}$ 具有纤维积和等化子，便保证
$\text{SR}(\mathcal{C})$ 具有纤维积和等化子，并且 $F$ 与它们交换
（引理 \ref{hypercovering-lemma-coprod-prod-SR}）；这足以定义所用的余骨架函子
（见《单纯方法》评注 \ref{simplicial-remark-existence-cosk} 和
《范畴论》引理 \ref{categories-lemma-fibre-products-equalizers-exist}）。
对一般的 $\mathcal{C}$，可把条件 (3) 替换为：
$F(K_{n + 1}) \longrightarrow ((\text{cosk}_n \text{sk}_n F(K))_{n + 1})$
对 $n \geq 1$ 在层化后成为满射；本节的结果仍然成立。

\medskip\noindent
设 $\mathcal{F}$ 是 $\mathcal{C}$ 上的阿贝尔层。上一节定义了
$\mathcal{F}$ 关于 $\text{SR}(\mathcal{C})$ 的单纯对象 $K$ 的
{\v C}ech 复形。现在，对给定预层 $\mathcal{G}$，置
$$
H^0(\mathcal{G}, \mathcal{F}) =
\Mor_{\textit{PSh}(\mathcal{C})}(\mathcal{G}, \mathcal{F}) =
\Mor_{\Sh(\mathcal{C})}(\mathcal{G}^\#, \mathcal{F}) =
H^0(\mathcal{G}^\#, \mathcal{F})
$$
记号同《位点上的上同调》第 \ref{sites-cohomology-section-limp} 节。
这是左正合函子，其高阶导出函子（《位点上的上同调》第
\ref{sites-cohomology-section-limp} 节作了简要研究）记为
$H^i(\mathcal{G}, \mathcal{F})$。我们将证明，给定 $\mathcal{G}$ 的
超覆盖 $K$，存在一个从 {\v C}ech 上同调到上同调的谱序列，
收敛到上同调 $H^i(\mathcal{G}, \mathcal{F})$。注意，若
$\mathcal{G} = *$，则 $H^i(*, \mathcal{F}) = H^i(\mathcal{C}, \mathcal{F})$，
这恢复了 $\mathcal{F}$ 在位点 $\mathcal{C}$ 上的上同调。

\begin{lemma}
\label{hypercovering-lemma-h0-cech-variant}
设 $\mathcal{C}$ 是具有等化子和纤维积的位点，
$\mathcal{G}$ 是 $\mathcal{C}$ 上的预层，$K$ 是 $\mathcal{G}$ 的超覆盖，
$\mathcal{F}$ 是 $\mathcal{C}$ 上的阿贝尔群层。则
$\check{H}^0(K, \mathcal{F}) = H^0(\mathcal{G}, \mathcal{F})$。
\end{lemma}

\begin{proof}
这由 $H^0(\mathcal{G}, \mathcal{F})$ 的定义以及下图在层化后成为余等化子图的事实得出：
$$
\xymatrix{
F(K_1) \ar@<1ex>[r] \ar@<-1ex>[r] &
F(K_0) \ar[r] & \mathcal{G}
}
$$
\end{proof}

\begin{lemma}
\label{hypercovering-lemma-injective-trivial-cech-variant}
设 $\mathcal{C}$ 是具有等化子和纤维积的位点，
$\mathcal{G}$ 是 $\mathcal{C}$ 上的预层，$K$ 是 $\mathcal{G}$ 的超覆盖，
$\mathcal{I}$ 是 $\mathcal{C}$ 上的内射阿贝尔群层。则
$$
\check{H}^p(K, \mathcal{I}) =
\left\{
\begin{matrix}
H^0(\mathcal{G}, \mathcal{I}) & \text{若} & p = 0 \\
0 & \text{若} & p > 0
\end{matrix}
\right.
$$
\end{lemma}

\begin{proof}
由 (\ref{hypercovering-equation-identify-cech})，有
$$
s(\mathcal{F}(K))
=
\Hom_{\textit{Ab}(\mathcal{C})}(s(\mathbf{Z}_{F(K)}^\#), \mathcal{F})
$$
复形 $s(\mathbf{Z}_{F(K)}^\#)$ 拟同构于
$\mathbf{Z}_\mathcal{G}^\#$，见引理
\ref{hypercovering-lemma-acyclic-hypercover-sheaves}。因此，若 $\mathcal{I}$ 是内射阿贝尔层，
则复形 $s(\mathcal{I}(K))$ 除了次数 $0$ 外是无环的。换言之，
对 $i > 0$ 有 $\check{H}^i(K, \mathcal{I}) = 0$。
结合引理 \ref{hypercovering-lemma-h0-cech-variant} 即得结论。
\end{proof}

\begin{lemma}
\label{hypercovering-lemma-cech-spectral-sequence-variant}
设 $\mathcal{C}$ 是具有等化子和纤维积的位点，
$\mathcal{G}$ 是 $\mathcal{C}$ 上的预层，$K$ 是 $\mathcal{G}$ 的超覆盖，
$\mathcal{F}$ 是 $\mathcal{C}$ 上的阿贝尔群层。
在 $D^{+}(\textit{Ab})$ 中存在关于 $\mathcal{F}$ 函子性的态射
$$
s(\mathcal{F}(K)) \longrightarrow R\Gamma(\mathcal{G}, \mathcal{F})
$$
它诱导函子 $\textit{Ab}(\mathcal{C}) \to \textit{Ab}$ 之间的自然变换
$$
\check{H}^i(K, -) \longrightarrow H^i(\mathcal{G}, -)
$$
此外，存在谱序列 $(E_r, d_r)_{r \geq 0}$，满足
$$
E_2^{p, q} = \check{H}^p(K, \underline{H}^q(\mathcal{F}))
$$
并收敛到 $H^{p + q}(\mathcal{G}, \mathcal{F})$。
此谱序列关于 $\mathcal{F}$ 和超覆盖 $K$ 都是函子性的。
\end{lemma}

\begin{proof}
在 $\mathcal{C}$ 上阿贝尔层的范畴中选取内射分辨
$\mathcal{F} \to \mathcal{I}^\bullet$。考虑各项为
$$
A^{p, q} = \mathcal{I}^q(K_p)
$$
的双复形 $A^{\bullet, \bullet}$，其中微分
$d_1^{p, q} : A^{p, q} \to A^{p + 1, q}$ 来自微分
$\mathcal{I}^p \to \mathcal{I}^{p + 1}$，而微分
$d_2^{p, q} : A^{p, q} \to A^{p, q + 1}$ 来自与余单纯阿贝尔群
$\mathcal{I}^p(K)$ 相伴的复形 $s(\mathcal{I}^p(K))$ 上的微分，
如上所述。我们将使用与此双复形相伴的两个谱序列
$({}'E_r, {}'d_r)$ 和 $({}''E_r, {}''d_r)$，见《同调代数》第
\ref{homology-section-double-complex} 节。

\medskip\noindent
由引理 \ref{hypercovering-lemma-injective-trivial-cech-variant}，复形
$s(\mathcal{I}^p(K))$ 在正次数无环，且其 $H^0$ 等于
$H^0(\mathcal{G}, \mathcal{I}^p)$。因此，由《同调代数》引理
\ref{homology-lemma-double-complex-gives-resolution} 及其证明，
谱序列 $({}'E_r, {}'d_r)$ 退化，并且自然映射
$$
H^0(\mathcal{G}, \mathcal{I}^\bullet) \longrightarrow
\text{Tot}(A^{\bullet, \bullet})
$$
是阿贝尔群复形的拟同构。引理中的态射
$s(\mathcal{F}(K)) \longrightarrow R\Gamma(\mathcal{G}, \mathcal{F})$
是自然映射 $s(\mathcal{F}(K)) \to \text{Tot}(A^{\bullet, \bullet})$
与上述拟同构之逆的复合。这是可行的，因为
$H^0(\mathcal{G}, \mathcal{I}^\bullet)$ 是
$R\Gamma(\mathcal{G}, \mathcal{F})$ 的一个代表元。

\medskip\noindent
考虑谱序列 $({}''E_r, {}''d_r)_{r \geq 0}$。由《同调代数》引理
\ref{homology-lemma-ss-double-complex}，有
$$
{}''E_2^{p, q} = H^p_{II}(H^q_I(A^{\bullet, \bullet}))
$$
换言之，先对 $d_1$ 取上同调，得到群
${}''E_1^{p, q} = \underline{H}^p(\mathcal{F})(K_q)$。
因此，由微分 ${}''d_1$ 的描述，确有
${}''E_2^{p, q} = \check{H}^p(K, \underline{H}^q(\mathcal{F}))$。
由于此谱序列收敛到 $\text{Tot}(A^{\bullet, \bullet})$ 的上同调，
证明完成。
\end{proof}

\begin{lemma}
\label{hypercovering-lemma-cech-spectral-sequence-verdier}
设 $\mathcal{C}$ 是具有等化子和纤维积的位点，$K$ 是超覆盖，
$\mathcal{F}$ 是阿贝尔层。存在谱序列 $(E_r, d_r)_{r \geq 0}$，满足
$$
E_2^{p, q} = \check{H}^p(K, \underline{H}^q(\mathcal{F}))
$$
并收敛到整体上同调群 $H^{p + q}(\mathcal{F})$。
\end{lemma}

\begin{proof}
这是引理 \ref{hypercovering-lemma-cech-spectral-sequence-variant} 的特殊情形。
\end{proof}







\section{超覆盖的覆盖细化}
\label{hypercovering-section-covering}

\noindent
下面给出若干构造超覆盖的方法。注意，由于范畴
$\text{SR}(\mathcal{C}, X)$ 具有纤维积，它的单纯对象范畴也具有纤维积，
见《单纯方法》引理 \ref{simplicial-lemma-fibre-product}。

\begin{lemma}
\label{hypercovering-lemma-funny-gamma}
设 $\mathcal{C}$ 是具有纤维积的位点，$X$ 是 $\mathcal{C}$ 的对象，
$K,L,M$ 是 $\text{SR}(\mathcal{C}, X)$ 的单纯对象，并设
$a : K \to L$、$b : M \to L$ 是态射。假设
\begin{enumerate}
\item $K$ 是 $X$ 的超覆盖；
\item 态射 $M_0 \to L_0$ 是覆盖；并且
\item 对每个 $n \geq 0$，下图中的箭头 $\gamma$ 是覆盖：
$$
\xymatrix{
M_{n + 1} \ar[dd] \ar[rr] \ar[rd]^\gamma &
&
(\text{cosk}_n \text{sk}_n M)_{n + 1} \ar[dd] \\
&
L_{n + 1}
\times_{(\text{cosk}_n \text{sk}_n L)_{n + 1}}
(\text{cosk}_n \text{sk}_n M)_{n + 1}
\ar[ld] \ar[ru]
& \\
L_{n + 1} \ar[rr] & & (\text{cosk}_n \text{sk}_n L)_{n + 1}
}
$$
\end{enumerate}
则纤维积 $K \times_L M$ 是 $X$ 的超覆盖。
\end{lemma}

\begin{proof}
由 (2)，态射 $(K \times_L M)_0 = K_0 \times_{L_0} M_0 \to K_0$
是覆盖的基变换，因而是覆盖，见引理 \ref{hypercovering-lemma-covering-permanence}。
由 (1)，$K_0 \to \{X \to X\}$ 是覆盖。因此，由引理
\ref{hypercovering-lemma-covering-permanence}，
$(K \times_L M)_0 \to \{X \to X\}$ 是覆盖。故
$K \times_L M$ 满足定义 \ref{hypercovering-definition-hypercovering} 的第一个条件。

\medskip\noindent
还需验证，对所有 $n \geq 0$，态射
$$
K_{n + 1} \times_{L_{n + 1}} M_{n + 1} = (K \times_L M)_{n + 1}
\longrightarrow
(\text{cosk}_n \text{sk}_n (K \times_L M))_{n + 1}
$$
是覆盖。简记
$A = (\text{cosk}_n \text{sk}_n K)_{n + 1}$,
$B = (\text{cosk}_n \text{sk}_n L)_{n + 1}$，以及
$C = (\text{cosk}_n \text{sk}_n M)_{n + 1}$.
函子 $\text{cosk}_n \text{sk}_n$ 与纤维积交换，见《单纯方法》引理
\ref{simplicial-lemma-cosk-fibre-product}。因此，上式右端等于
$A \times_B C$。考虑交换图
$$
\xymatrix{
K_{n + 1} \times_{L_{n + 1}} M_{n + 1} \ar[r] \ar[d] &
M_{n + 1} \ar[d] \ar[rd]_\gamma \ar[rrd] &
& \\
K_{n + 1} \ar[r] \ar[rd] &
L_{n + 1} \ar[rrd] &
L_{n + 1} \times_B C \ar[l] \ar[r] &
C \ar[d] \\
&
A \ar[rr] &
&
B
}
$$
此图表明
$$
K_{n + 1} \times_{L_{n + 1}} M_{n + 1}
=
(K_{n + 1} \times_B C)
\times_{(L_{n + 1} \times_B C), \gamma}
M_{n + 1}
$$
现在，$K_{n + 1} \times_B C \to A \times_B C$ 是覆盖
$K_{n + 1} \to A$ 沿态射 $A \times_B C \to A$ 的基变换，
故为覆盖。由假设 (3)，态射 $\gamma$ 是覆盖。因此态射
$$
(K_{n + 1} \times_B C)
\times_{(L_{n + 1} \times_B C), \gamma}
M_{n + 1}
\longrightarrow
K_{n + 1} \times_B C
$$
作为覆盖的基变换也是覆盖。覆盖的复合仍是覆盖，故引理成立。
\end{proof}

\begin{lemma}
\label{hypercovering-lemma-product-hypercoverings}
设 $\mathcal{C}$ 是具有纤维积的位点，$X$ 是 $\mathcal{C}$ 的对象。
若 $K,L$ 是 $X$ 的超覆盖，则 $K \times L$ 是 $X$ 的超覆盖。
\end{lemma}

\begin{proof}
既可直接验证，也可应用上述引理 \ref{hypercovering-lemma-funny-gamma}，并验证
$L \to \{X \to X\}$ 具有性质 (3)。
\end{proof}


\noindent
设 $\mathcal{C}$ 是具有纤维积的位点，$X$ 是 $\mathcal{C}$ 的对象。
由于范畴 $\text{SR}(\mathcal{C}, X)$ 具有余积和有限极限，
对某些单纯集 $U$（例如只有有限多个非退化单形者）以及
$\text{SR}(\mathcal{C}, X)$ 的任意单纯对象 $K$，可以讨论对象
$U \times K$ 和 $\Hom(U, K)$。见《单纯方法》第
\ref{simplicial-section-product-with-simplicial-sets} 节和第
\ref{simplicial-section-hom-from-simplicial-sets} 节。

\begin{lemma}
\label{hypercovering-lemma-covering}
设 $\mathcal{C}$ 是具有纤维积的位点，$X$ 是 $\mathcal{C}$ 的对象，
$K$ 是 $X$ 的超覆盖，$k \geq 0$ 是整数。设
$u : Z \to K_k$ 是 $\text{SR}(\mathcal{C}, X)$ 中的覆盖。
则存在超覆盖的态射 $f: L \to K$，使得 $L_k \to K_k$ 经由 $u$ 分解。
\end{lemma}

\begin{proof}
记 $Y = K_k$。设 $C[k]$ 是《单纯方法》例
\ref{simplicial-example-simplex-cosimplicial-set} 中定义的余单纯集。
我们将使用《单纯方法》引理
\ref{simplicial-lemma-morphism-into-product} 给出的
$\Hom(C[k], Y)$ 和 $\Hom(C[k], Z)$ 的描述。
存在对应于 $\text{id} : K_k = Y \to Y$ 的典范态射
$K \to \Hom(C[k], Y)$。考虑态射
$\Hom(C[k], Z) \to \Hom(C[k], Y)$；其次数 $n$ 项上的态射为
$$
\prod\nolimits_{\alpha : [k] \to [n]} Z
\longrightarrow
\prod\nolimits_{\alpha : [k] \to [n]} Y
$$
即在每个因子上使用给定态射 $Z \to Y$。置
$$
L = K \times_{\Hom(C[k], Y)} \Hom(C[k], Z).
$$
态射 $L_k \to K_k$ 位于交换图
$$
\xymatrix{
L_k \ar[r] \ar[d] &
\prod_{\alpha : [k] \to [k]} Z \ar[r]^-{\text{pr}_{\text{id}_{[k]}}} \ar[d] &
Z \ar[d] \\
K_k \ar[r] &
\prod_{\alpha : [k] \to [k]} Y \ar[r]^-{\text{pr}_{\text{id}_{[k]}}} &
Y
}
$$
中。由于底部两个箭头的复合是恒等态射，便得到所需分解。

\medskip\noindent
还需证明 $L$ 是 $X$ 的超覆盖。为此应用引理
\ref{hypercovering-lemma-funny-gamma}。条件 (1) 由假设成立。对于 (2)，态射
$$
\Hom(C[k], Z)_0 \to \Hom(C[k], Y)_0
$$
是覆盖，因为态射 $[k] \to [0]$ 仅有一个，故它同构于 $Z \to Y$。

\medskip\noindent
考虑 $n = 0$ 时的条件 (3)。由于
$(\text{cosk}_0 T)_1 = T \times T$
（《单纯方法》例 \ref{simplicial-example-cosk0}）且
$\Hom(C[k], Z)_1 = \prod_{\alpha : [k] \to [1]} Z$，得到图
$$
\xymatrix{
\prod\nolimits_{\alpha : [k] \to [1]} Z \ar[r] \ar[d] &
Z \times Z \ar[d] \\
\prod\nolimits_{\alpha : [k] \to [1]} Y \ar[r] &
Y \times Y
}
$$
其中水平箭头对应于向两个非满射 $\alpha$ 所对应因子的投影。
因此箭头 $\gamma$ 是态射
$$
\prod\nolimits_{\alpha : [k] \to [1]} Z
\longrightarrow
\prod\nolimits_{\alpha : [k] \to [1]\text{ 非满}} Z
\times
\prod\nolimits_{\alpha : [k] \to [1]\text{ 满}} Y
$$
它是覆盖的乘积，故由引理 \ref{hypercovering-lemma-covering-permanence} 仍为覆盖。

\medskip\noindent
考虑 $n > 0$ 时的条件 (3)。我们断言，存在有限集的单射
$\tau : S' \to S$，使得对 $\text{SR}(\mathcal{C}, X)$ 的任意对象
$T$，态射
\begin{equation}
\label{hypercovering-equation-map}
\Hom(C[k], T)_{n + 1} \to
(\text{cosk}_n \text{sk}_n \Hom(C[k], T))_{n + 1}
\end{equation}
关于 $T$ 函子性地同构于投影
$\prod_{s \in S} T \to \prod_{s' \in S'} T$。若此言成立，
则仿照上一段可见，箭头 $\gamma$ 是态射
$$
\prod\nolimits_{s \in S} Z
\longrightarrow
\prod\nolimits_{s \in S'} Z
\times
\prod\nolimits_{s \not\in \tau(S')} Y
$$
它是覆盖的乘积，故由引理 \ref{hypercovering-lemma-covering-permanence} 仍为覆盖。
由构造，有
$\Hom(C[k], T)_{n + 1} = \prod_{\alpha : [k] \to [n + 1]} T$
（见《单纯方法》引理 \ref{simplicial-lemma-morphism-into-product}）。
相应地取 $S = \text{Map}([k], [n + 1])$。另一方面，
《单纯方法》引理 \ref{simplicial-lemma-formula-limit} 把
$(\text{cosk}_n \text{sk}_n \Hom(C[k], T))_{n + 1}$
的点描述为 $\Hom(C[k], T)_n$ 的点列
$(f_0, \ldots, f_{n + 1})$，满足
$d^n_{j - 1} f_i = d^n_i f_j$（$0 \leq i < j \leq n + 1$）。
可写 $f_i = (f_{i, \alpha})$，其中 $f_{i, \alpha}$ 是 $T$ 的点，
$\alpha \in \text{Map}([k], [n])$。这些条件化为
$$
f_{i, \delta^n_{j - 1} \circ \beta} = f_{j, \delta_i^n \circ \beta}
$$
对任意 $0 \leq i < j \leq n + 1$ 和 $\beta : [k] \to [n - 1]$
成立。因此
$$
S' = \{0, \ldots, n + 1\} \times \text{Map}([k], [n]) / \sim
$$
其中等价关系由等价式
$$
(i, \delta^n_{j - 1} \circ \beta) \sim (j, \delta_i^n \circ \beta)
$$
生成，其中 $0 \leq i < j \leq n + 1$ 且
$\beta : [k] \to [n - 1]$。一个略去的计算表明，态射
(\ref{hypercovering-equation-map}) 对应于映射 $S' \to S$，它把 $(i, \alpha)$
送到 $\delta^{n + 1}_i \circ \alpha \in S$。（由《单纯方法》引理
\ref{simplicial-lemma-relations-face-degeneracy} 的 (1)，此映射确实良定义，
这也许能使读者安心。）为完成证明，只需说明：若
$\alpha, \alpha' : [k] \to [n]$、$0 \leq i < j \leq n + 1$，且
$$
\delta^{n + 1}_i \circ \alpha = \delta^{n + 1}_j \circ \alpha'
$$
则存在某个 $\beta : [k] \to [n - 1]$，使得
$\alpha = \delta^n_{j - 1} \circ \beta$ 且
$\alpha' = \delta_i^n \circ \beta$。这很容易验证，故略去。
\end{proof}

\begin{lemma}
\label{hypercovering-lemma-covering-sheaf}
设 $\mathcal{C}$ 是具有纤维积的位点，$X$ 是 $\mathcal{C}$ 的对象，
$K$ 是 $X$ 的超覆盖，$n \geq 0$ 是整数。设
$u : \mathcal{F} \to F(K_n)$ 是层化后成为满射的预层态射。
则存在超覆盖的态射 $f: L \to K$，使得
$F(f_n) : F(L_n) \to F(K_n)$ 经由 $u$ 分解。
\end{lemma}

\begin{proof}
写成 $K_n = \{U_i \to X\}_{i \in I}$。于是映射 $u$ 是集合预层态射
$u : \mathcal{F} \to \amalg h_{u_i}$。关于 $u$ 的假设意味着：
对每个 $i \in I$，存在位点 $\mathcal{C}$ 的覆盖
$\{U_{ij} \to U_i\}_{j \in I_i}$ 和预层态射
$t_{ij} : h_{U_{ij}} \to \mathcal{F}$，使得 $u \circ t_{ij}$
是由态射 $U_{ij} \to U_i$ 得到的映射
$h_{U_{ij}} \to h_{U_i}$。置 $J = \amalg_{i \in I} I_i$，并令
$\alpha : J \to I$ 为显然的映射。对 $j \in J$，记
$V_j = U_{\alpha(j)j}$，并置 $Z = \{V_j \to X\}_{j \in J}$。
最后，考虑由 $\alpha : J \to I$ 和上述态射
$V_j = U_{\alpha(j)j} \to U_{\alpha(j)}$ 给出的态射
$u' : Z \to K_n$。显然，这是范畴 $\text{SR}(\mathcal{C}, X)$ 中的覆盖；
由构造，$F(u') : F(Z) \to F(K_n)$ 经由 $u$ 分解。
故结论由上述引理 \ref{hypercovering-lemma-covering} 得出。
\end{proof}


\section{添加单形}
\label{hypercovering-section-adding-simplices}

\noindent
本节证明后面将用到的一些技术性引理。设 $\mathcal{C}$ 是具有纤维积的位点，
$X$ 是 $\mathcal{C}$ 的对象。如上述第 \ref{hypercovering-section-covering} 节所指出，
对某些单纯集 $U$ 和 $\text{SR}(\mathcal{C}, X)$ 的任意单纯对象 $K$，
对象 $U \times K$ 和 $\Hom(U, K)$ 都已定义。见《单纯方法》第
\ref{simplicial-section-product-with-simplicial-sets} 节和第
\ref{simplicial-section-hom-from-simplicial-sets} 节。

\begin{lemma}
\label{hypercovering-lemma-one-more-simplex}
设 $\mathcal{C}$ 是具有纤维积的位点，$X$ 是 $\mathcal{C}$ 的对象，
$K$ 是 $X$ 的超覆盖。设 $U \subset V$ 是单纯集，且对所有 $n$，
$U_n,V_n$ 都是有限非空集。假设 $U$ 只有有限多个非退化单形。
设 $n \geq 0$ 且 $x \in V_n$、$x \not \in U_n$，并满足
\begin{enumerate}
\item 对 $i < n$ 有 $V_i = U_i$；
\item $V_n = U_n \cup \{x\}$；
\item 对 $j > n$，任意 $z \in V_j$、$z \not \in U_j$ 都是退化的。
\end{enumerate}
则 $\text{SR}(\mathcal{C}, X)$ 中的态射
$$
\Hom(V, K)_0
\longrightarrow
\Hom(U, K)_0
$$
是覆盖。
\end{lemma}

\begin{proof}
若 $n = 0$，则容易推出 $V = U \amalg \Delta[0]$（见下文）。
此时 $\Hom(V, K)_0 = \Hom(U, K)_0 \times K_0$，结论由引理
\ref{hypercovering-lemma-covering-permanence} 得出。

\medskip\noindent
设 $a : \Delta[n] \to V$ 是《单纯方法》引理
\ref{simplicial-lemma-simplex-map} 中与 $x$ 相伴的态射。
记 $\partial \Delta[n] = i_{(n-1)!} \text{sk}_{n - 1} \Delta[n]$
为 $\Delta[n]$ 的 $(n - 1)$-骨架。设
$b : \partial \Delta[n] \to U$ 是 $a$ 在 $\Delta[n]$ 的
$(n - 1)$-骨架上的限制。由《单纯方法》引理
\ref{simplicial-lemma-glue-simplex}，有
$V = U \amalg_{\partial \Delta[n]} \Delta[n]$。由《单纯方法》引理
\ref{simplicial-lemma-hom-from-coprod}，图
$$
\xymatrix{
\Hom(V, K)_0 \ar[r] \ar[d] &
\Hom(U, K)_0 \ar[d] \\
\Hom(\Delta[n], K)_0 \ar[r] &
\Hom(\partial \Delta[n], K)_0
}
$$
是纤维积方块。因此只需证明底部水平箭头是覆盖。由《单纯方法》引理
\ref{simplicial-lemma-cosk-shriek}，此箭头等同于
$$
K_n \to (\text{cosk}_{n - 1} \text{sk}_{n - 1} K)_n
$$
故由超覆盖的定义，它是覆盖。
\end{proof}

\begin{lemma}
\label{hypercovering-lemma-add-simplices}
设 $\mathcal{C}$ 是具有纤维积的位点，$X$ 是 $\mathcal{C}$ 的对象，
$K$ 是 $X$ 的超覆盖。设 $U \subset V$ 是单纯集，且对所有 $n$，
$U_n,V_n$ 都是有限非空集。假设 $U$ 和 $V$ 都只有有限多个非退化单形。
则 $\text{SR}(\mathcal{C}, X)$ 中的态射
$$
\Hom(V, K)_0
\longrightarrow
\Hom(U, K)_0
$$
是覆盖。
\end{lemma}

\begin{proof}
由上述引理 \ref{hypercovering-lemma-one-more-simplex}，只需证明关于引理中这类单纯集
包含 $U \subset V$ 的一个简单引理。这恰是《单纯方法》引理
\ref{simplicial-lemma-add-simplices} 的结论。
\end{proof}

\begin{lemma}
\label{hypercovering-lemma-degeneracy-maps-coverings}
设 $\mathcal{C}$ 是具有纤维积的位点，$X$ 是 $\mathcal{C}$ 的对象，
$K$ 是 $X$ 的超覆盖。则
\begin{enumerate}
\item 对每个 $n \geq 0$，$K_n$ 是 $X$ 的覆盖；
\item 对所有 $n \geq 1$ 和 $0 \leq i \leq n$，
$d^n_i : K_n \to K_{n - 1}$ 是覆盖。
\end{enumerate}
\end{lemma}

\begin{proof}
回忆，由定义 \ref{hypercovering-definition-hypercovering}，$K_0$ 是 $X$ 的覆盖；
这等价于说 $K_0 \to \{X \to X\}$ 是定义
\ref{hypercovering-definition-covering-SR} 意义下的覆盖。因此 (1) 由 (2) 得出，
因为 (2) 将证明复合
$K_n \to K_{n - 1} \to \ldots \to K_0 \to \{X \to X\}$
由引理 \ref{hypercovering-lemma-covering-permanence} 是覆盖。

\medskip\noindent
证明 (2)。由《单纯方法》引理
\ref{simplicial-lemma-exists-hom-from-simplicial-set-finite}，有
$\Mor(\Delta[n], K)_0 = K_n$。因此，把引理
\ref{hypercovering-lemma-add-simplices} 应用于 $n + 1$ 个不同的包含
$\Delta[n - 1] \to \Delta[n]$，便得 (2)。
\end{proof}

\begin{remark}
\label{hypercovering-remark-P-covering}
引理 \ref{hypercovering-lemma-add-simplices} 和
\ref{hypercovering-lemma-degeneracy-maps-coverings} 有如下有用的特殊情形。
设范畴 $\mathcal{C}$ 具有纤维积。设
$P \subset \text{Arrows}(\mathcal{C})$ 在基变换和复合下稳定，
并包含所有同构。所谓一个 {\it $P$-超覆盖}，是从
$\mathcal{C}$ 的单纯对象出发的增广 $a : U \to X$，使得
\begin{enumerate}
\item $U_0 \to X$ 属于 $P$；
\item $U_1 \to U_0 \times_X U_0$ 属于 $P$；
\item $U_{n + 1} \to (\text{cosk}_n\text{sk}_n U)_{n + 1}$
对 $n \geq 1$ 属于 $P$。
\end{enumerate}
范畴 $\mathcal{C}/X$ 具有所有有限极限，故上述表述中所用的余骨架存在
（见《范畴论》引理 \ref{categories-lemma-finite-limits-exist}）。
我们断言，态射 $U_n \to X$ 和 $d^n_i : U_n \to U_{n - 1}$ 都属于
$P$。把 $\mathcal{C}$ 视为一个位点，其覆盖为满足 $f \in P$ 的
$\{f : V \to U\}$，再取由 $K_n = \{U_n \to X\}$ 给出的 $K$，
便可由上述引理得到此结论。
\end{remark}


\section{同伦}
\label{hypercovering-section-homotopies}

\noindent
设 $\mathcal{C}$ 是具有纤维积的位点，$X$ 是 $\mathcal{C}$ 的对象，
$L$ 是 $\text{SR}(\mathcal{C}, X)$ 的单纯对象。由《单纯方法》引理
\ref{simplicial-lemma-exists-hom-from-simplicial-set-finite}，
范畴 $\text{Simp}(\text{SR}(\mathcal{C}, X))$ 中存在对象
$\Hom(\Delta[1], L)$，它表示函子
$$
T
\longmapsto
\Mor_{\text{Simp}(\text{SR}(\mathcal{C}, X))}(\Delta[1] \times T, L)
$$
由 $e_i : \Delta[0] \to \Delta[1]$ 和等同
$\Hom(\Delta[0], L) = L$，得到典范态射
$$
\Hom(\Delta[1], L) \to L \times L
$$

\begin{lemma}
\label{hypercovering-lemma-hom-hypercovering}
设 $\mathcal{C}$ 是具有纤维积的位点，$X$ 是 $\mathcal{C}$ 的对象，
$L$ 是 $\text{SR}(\mathcal{C}, X)$ 的单纯对象。设 $n \geq 0$。
考虑由上述态射得到的交换图
\begin{equation}
\label{hypercovering-equation-diagram}
\xymatrix{
\Hom(\Delta[1], L)_{n + 1} \ar[r] \ar[d] &
(\text{cosk}_n \text{sk}_n \Hom(\Delta[1], L))_{n + 1} \ar[d] \\
(L \times L)_{n + 1} \ar[r] &
(\text{cosk}_n \text{sk}_n (L \times L))_{n + 1}
}
\end{equation}
此图各项可作如下等同，其中
$\partial \Delta[n + 1] = i_{n!}\text{sk}_n \Delta[n + 1]$
是 $(n + 1)$-单形的 $n$-骨架：
\begin{eqnarray*}
\Hom(\Delta[1], L)_{n + 1}
& = &
\Hom(\Delta[1] \times \Delta[n + 1], L)_0 \\
(\text{cosk}_n \text{sk}_n \Hom(\Delta[1], L))_{n + 1}
& = &
\Hom(\Delta[1] \times \partial \Delta[n + 1], L)_0 \\
(L \times L)_{n + 1}
& = &
\Hom(
(\Delta[n + 1] \amalg \Delta[n + 1], L)_0 \\
(\text{cosk}_n \text{sk}_n (L \times L))_{n + 1}
& = &
\Hom(
\partial \Delta[n + 1]
\amalg
\partial \Delta[n + 1], L)_0
\end{eqnarray*}
而 $\text{SR}(\mathcal{C}, X)$ 中这些对象之间的态射来自单纯集交换图
\begin{equation}
\label{hypercovering-equation-dual-diagram}
\xymatrix{
\Delta[1] \times \Delta[n + 1] &
\Delta[1] \times \partial\Delta[n + 1] \ar[l] \\
\Delta[n + 1] \amalg \Delta[n + 1] \ar[u] &
\partial\Delta[n + 1] \amalg \partial\Delta[n + 1]
\ar[l] \ar[u]
}
\end{equation}
此外，图 (\ref{hypercovering-equation-diagram}) 中底部箭头与右侧箭头的纤维积等于
$$
\Hom(U, L)_0
$$
其中 $U \subset \Delta[1] \times \Delta[n + 1]$ 是最小单纯子集，
使得 $\Delta[n + 1] \amalg \Delta[n + 1]$ 和
$\Delta[1] \times \partial\Delta[n + 1]$ 都映入其中。
\end{lemma}

\begin{proof}
第一和第三个等式是《单纯方法》引理
\ref{simplicial-lemma-exists-hom-from-simplicial-set-finite}。
第二和第四个等式由该引理与《单纯方法》引理
\ref{simplicial-lemma-cosk-shriek} 合并得出。最后一个断言由如下事实得出：
$U$ 是图 (\ref{hypercovering-equation-dual-diagram}) 的底部和右侧箭头的推出；
再应用《单纯方法》引理 \ref{simplicial-lemma-hom-from-coprod}。
要看出 $U$ 等于此推出，只需看出
$\Delta[n + 1] \amalg \Delta[n + 1]$ 与
$\Delta[1] \times \partial\Delta[n + 1]$ 在
$\Delta[1] \times \Delta[n + 1]$ 中的交等于
$\partial\Delta[n + 1] \amalg \partial\Delta[n + 1]$。
这留给读者验证。
\end{proof}

\begin{lemma}
\label{hypercovering-lemma-homotopy}
设 $\mathcal{C}$ 是具有纤维积的位点，$X$ 是 $\mathcal{C}$ 的对象，
$K,L$ 是 $X$ 的超覆盖，$a,b : K \to L$ 是超覆盖的态射。
则存在超覆盖的态射 $c : K' \to K$，使得
$a \circ c$ 同伦于 $b \circ c$。
\end{lemma}

\begin{proof}
考虑交换图
$$
\xymatrix{
K' \ar@{=}[r]^-{def} \ar[rd]_c &
K \times_{(L \times L)} \Hom(\Delta[1], L)
\ar[r] \ar[d] & \Hom(\Delta[1], L) \ar[d] \\
& K \ar[r]^{(a, b)} & L \times L
}
$$
由 $\Hom(\Delta[1], L)$ 的函子性质，水平态射的复合对应于一个态射
$K' \times \Delta[1] \to L$；它定义 $c \circ a$ 与 $c \circ b$
之间的同伦。因此，只要证明 $K'$ 是 $X$ 的超覆盖，就可得到引理。
为此，把引理 \ref{hypercovering-lemma-funny-gamma} 应用于态射对
$K \to L \times L$ 和 $\Hom(\Delta[1], L) \to L \times L$。
引理 \ref{hypercovering-lemma-funny-gamma} 的条件 (1) 成立。引理
\ref{hypercovering-lemma-funny-gamma} 的条件 (2) 也成立，
因为 $\Hom(\Delta[1], L)_0 = L_1$；又因假设 $L$ 是超覆盖，态射
$(d^1_0, d^1_1) : L_1 \to L_0 \times L_0$ 是
$\text{SR}(\mathcal{C}, X)$ 中的覆盖。为证明引理
\ref{hypercovering-lemma-funny-gamma} 的条件 (3)，应用上述引理
\ref{hypercovering-lemma-hom-hypercovering}。根据该引理，引理
\ref{hypercovering-lemma-funny-gamma} 的条件 (3) 中的态射
$\gamma$ 是态射
$$
\Hom(\Delta[1] \times \Delta[n + 1], L)_0
\longrightarrow
\Hom(U, L)_0
$$
其中 $U \subset \Delta[1] \times \Delta[n + 1]$。由引理
\ref{hypercovering-lemma-add-simplices}，这是覆盖，故断言得证。
\end{proof}

\begin{remark}
\label{hypercovering-remark-contractible-category}
注意，证明的关键在于应用引理 \ref{hypercovering-lemma-add-simplices}。
这个引理完全是一般性的，并不在意单纯集的确切形状
（只要它们仅有有限多个非退化单形）。因此完全有理由期待如下结果：
给定任意态射 $a : K \times \partial \Delta[k] \to L$，其中 $K,L$
是超覆盖，则存在超覆盖的态射 $c : K' \to K$ 和态射
$g : K' \times \Delta[k] \to L$，使得
$g|_{K' \times \partial \Delta[k]} =
a \circ (c \times \text{id}_{\partial \Delta[k]})$.
换言之，超覆盖的范畴在适当意义下是可缩的。
\end{remark}

















\section{上同调与超覆盖}
\label{hypercovering-section-cohomology}

\noindent
设 $\mathcal{C}$ 是具有纤维积的位点，$X$ 是 $\mathcal{C}$ 的对象，
$\mathcal{F}$ 是 $\mathcal{C}$ 上的阿贝尔群层，$K,L$ 是 $X$ 的超覆盖。
若 $a,b : K \to L$ 是同伦映射，则
$\mathcal{F}(a), \mathcal{F}(b) : \mathcal{F}(K) \to \mathcal{F}(L)$
也是同伦映射，见《单纯方法》引理
\ref{simplicial-lemma-functorial-homotopy}。因此，它们在相伴上链复形的
上同调群上产生相同作用，见《单纯方法》引理
\ref{simplicial-lemma-homotopy-s-Q}。我们将用这一点定义遍历所有超覆盖的余极限。

\medskip\noindent
暂记 $\text{HC}(\mathcal{C}, X)$ 为如下范畴：其对象是 $X$ 的超覆盖，
其态射是 $X$ 的超覆盖之间的映射的同伦类。我们已经看到这是一个范畴，
而且不是“大”范畴，见引理 \ref{hypercovering-lemma-hypercoverings-set}。
$\text{HC}(\mathcal{C}, X)$ 的反范畴将作为下图的指标范畴；术语见
《范畴论》第 \ref{categories-section-limits} 节。考虑图
$$
\check{H}^i(-, \mathcal{F}) :
\text{HC}(\mathcal{C}, X)^{opp}
\longrightarrow
\textit{Ab}.
$$
由引理 \ref{hypercovering-lemma-product-hypercoverings}、\ref{hypercovering-lemma-homotopy}
以及上述关于同伦的说明，此图是有向的，见《范畴论》定义
\ref{categories-definition-directed}。因此余极限
$$
\check{H}^i_{\text{HC}}(X, \mathcal{F}) =
\colim_{K \in \text{HC}(\mathcal{C}, X)} \check{H}^i(K, \mathcal{F})
$$
有特别简单的描述（见所引位置）。

\begin{theorem}
\label{hypercovering-theorem-cohomology-hypercoverings}
设 $\mathcal{C}$ 是具有纤维积的位点，$X$ 是 $\mathcal{C}$ 的对象，
$i \geq 0$。函子
\begin{eqnarray*}
\textit{Ab}(\mathcal{C}) & \longrightarrow & \textit{Ab} \\
\mathcal{F} & \longmapsto & H^i(X, \mathcal{F}) \\
\mathcal{F} & \longmapsto & \check{H}^i_{\text{HC}}(X, \mathcal{F})
\end{eqnarray*}
典范同构。
\end{theorem}

\begin{proof}[使用谱序列的证明]
设对某个 $p \geq 0$ 有 $\xi \in H^p(X, \mathcal{F})$。
我们证明：对 $X$ 的某个超覆盖 $K$，$\xi$ 位于引理
\ref{hypercovering-lemma-cech-spectral-sequence} 的映射
$\check{H}^p(X, \mathcal{F}) \to H^p(X, \mathcal{F})$ 的像中。

\medskip\noindent
若 $p = 0$，这由引理 \ref{hypercovering-lemma-h0-cech} 成立。若 $p = 1$，
选取例 \ref{hypercovering-example-cech} 中 $X$ 的 {\v C}ech 超覆盖 $K$，
从位点 $\mathcal{C}$ 中满足 $\xi|_{U_i} = 0$ 的覆盖
$K_0 = \{U_i \to X\}$ 开始；见《位点上的上同调》引理
\ref{sites-cohomology-lemma-kill-cohomology-class-on-covering}。
此时立即由引理 \ref{hypercovering-lemma-cech-spectral-sequence} 的谱序列得出，
$\xi$ 来自 $\check{H}^1(K, \mathcal{F})$ 的一个元素。
一般地，选取 $X$ 的任意超覆盖 $K$，使得 $\xi$ 在
$\underline{H}^p(\mathcal{F})(K_0)$ 中的像为零（再次使用例
\ref{hypercovering-example-cech} 和《位点上的上同调》引理
\ref{sites-cohomology-lemma-kill-cohomology-class-on-covering}）。
由引理 \ref{hypercovering-lemma-cech-spectral-sequence} 的谱序列，$\xi$ 来自
$\check{H}^p(K, \mathcal{F})$ 的一个元素的障碍，是一列元素
$\xi_1, \ldots, \xi_{p - 1}$，其中
$\xi_q \in \check{H}^{p - q}(K, \underline{H}^q(\mathcal{F}))$
（更精确地说，是 $\xi_q$ 在这些群的某些子商中的像）。

\medskip\noindent
可以归纳地用细化替换超覆盖 $K$，使障碍
$\xi_1, \ldots, \xi_{p - 1}$ 的限制为零（而不仅是其在子商中的像为零，
故此处没有微妙之处）。事实上，假设已经达到
$\xi_{q + 1}, \ldots, \xi_{p - 1}$ 都为零的情形。注意，
$\xi_q \in \check{H}^{p - q}(K, \underline{H}^q(\mathcal{F}))$
是某个元素
$$
\tilde \xi_q \in
\underline{H}^q(\mathcal{F})(K_{p - q}) =
\prod H^q(U_i, \mathcal{F})
$$
的类，其中 $K_{p - q} = \{U_i \to X\}_{i \in I}$。设
$\xi_{q, i}$ 为 $\tilde \xi_q$ 在 $H^q(U_i, \mathcal{F})$ 中的分量。
由于 $q \geq 1$，可再次应用《位点上的上同调》引理
\ref{sites-cohomology-lemma-kill-cohomology-class-on-covering}，
选取位点的覆盖 $\{U_{i, j} \to U_i\}$，使每个限制
$\xi_{q, i}|_{U_{i, j}} = 0$。考虑范畴
$\text{SR}(\mathcal{C}, X)$ 的对象 $Z = \{U_{i, j} \to X\}$ 及其显然态射
$u : Z \to K_{p - q}$。显然 $u$ 是覆盖，见定义
\ref{hypercovering-definition-covering-SR}。由引理 \ref{hypercovering-lemma-covering}，存在 $X$
的超覆盖的态射 $L \to K$，使 $L_{p - q} \to K_{p - q}$ 经由 $u$
分解。于是显然，$\xi_q$ 在
$\underline{H}^q(\mathcal{F})(L_{p - q})$ 中的像为零。
由于引理 \ref{hypercovering-lemma-cech-spectral-sequence} 的谱序列是函子性的，
用 $L$ 替换 $K$ 后便达到 $\xi_q, \ldots, \xi_{p - 1}$ 均为零的情形。
如此继续，最终得到一个使它们全为零的超覆盖，因而 $\xi$ 位于映射
$\check{H}^p(X, \mathcal{F}) \to H^p(X, \mathcal{F})$ 的像中。

\medskip\noindent
设 $K$ 是 $X$ 的超覆盖，$\xi \in \check{H}^p(K, \mathcal{F})$，
且 $\xi$ 在引理 \ref{hypercovering-lemma-cech-spectral-sequence} 的映射
$\check{H}^p(X, \mathcal{F}) \to H^p(X, \mathcal{F})$ 下的像为零。
为完成定理证明，需说明存在超覆盖的态射 $L \to K$，使 $\xi$
在 $\check{H}^p(L, \mathcal{F})$ 中的限制为零。由引理
\ref{hypercovering-lemma-cech-spectral-sequence} 的谱序列，$\xi$ 在
$H^p(X, \mathcal{F})$ 中的像消失，意味着存在元素
$\xi_1, \ldots, \xi_{p - 2}$，其中
$\xi_q \in \check{H}^{p - 1 - q}(K, \underline{H}^q(\mathcal{F}))$
（更精确地说，是它们在某些子商中的像），使得谱序列中的像
$d_{q + 1}^{p - 1 - q, q}\xi_q$ 之和为 $\xi$。因此，完全按照上述机制，
可以找到超覆盖的态射 $L \to K$，使元素
$\xi_q$（$q = 1, \ldots, p - 2$）在
$\check{H}^{p - 1 - q}(L, \underline{H}^q(\mathcal{F}))$ 中的限制为零。
于是，由于态射 $L \to K$ 依照引理
\ref{hypercovering-lemma-cech-spectral-sequence} 诱导谱序列的态射，可得 $\xi$ 为零。
\end{proof}

\begin{proof}[不使用谱序列的证明]
我们已经看到 $i = 0$ 时的结论，见引理 \ref{hypercovering-lemma-h0-cech}。
已知函子 $H^i(X, -)$ 构成泛 $\delta$-函子，见《导出范畴》引理
\ref{derived-lemma-higher-derived-functors}。为证明定理，只需说明函子列
$\check{H}^i_{HC}(X, -)$ 构成 $\delta$-函子。事实上，已知
{\v C}ech 上同调在内射层上为零（引理
\ref{hypercovering-lemma-injective-trivial-cech}），于是可以应用《同调代数》引理
\ref{homology-lemma-efface-implies-universal}。

\medskip\noindent
设
$$
0 \to \mathcal{F} \to \mathcal{G} \to \mathcal{H} \to 0
$$
是 $\mathcal{C}$ 上阿贝尔层的短正合列。设
$\xi \in \check{H}^p_{HC}(X, \mathcal{H})$。选取 $X$ 的超覆盖 $K$
以及在上同调中代表 $\xi$ 的元素 $\sigma \in \mathcal{H}(K_p)$。
存在相应的复形正合列
$$
0 \to s(\mathcal{F}(K)) \to s(\mathcal{G}(K)) \to s(\mathcal{H}(K))
$$
但不能保证右端也有一个零；这正是妨碍我们直接应用蛇引理定义
$\delta(\xi)$ 的唯一因素。回忆，若 $K_p = \{U_i \to X\}$，则
$$
\mathcal{H}(K_p) = \prod \mathcal{H}(U_i)
$$
写 $\sigma =\prod \sigma_i$，其中 $\sigma_i \in \mathcal{H}(U_i)$。
由于 $\mathcal{G} \to \mathcal{H}$ 是层的满射，存在覆盖
$\{U_{i, j} \to U_i\}$，使 $\sigma_i|_{U_{i, j}}$ 是某个元素
$\tau_{i, j} \in \mathcal{G}(U_{i, j})$ 的像。考虑范畴
$\text{SR}(\mathcal{C}, X)$ 的对象 $Z = \{U_{i, j} \to X\}$ 及其显然态射
$u : Z \to K_p$。显然 $u$ 是覆盖，见定义
\ref{hypercovering-definition-covering-SR}。由引理 \ref{hypercovering-lemma-covering}，存在 $X$
的超覆盖的态射 $L \to K$，使 $L_p \to K_p$ 经由 $u$ 分解。
因此用 $L$ 替换 $K$ 后，可以假设 $\sigma$ 是某个元素
$\tau \in \mathcal{G}(K_p)$ 的像。注意 $d(\sigma) = 0$，
但不一定有 $d(\tau) = 0$。因而
$d(\tau) \in \mathcal{F}(K_{p + 1})$ 是闭链。在此情形下，定义
$\delta(\xi)$ 为闭链 $d(\tau)$ 在
$\check{H}^{p + 1}_{HC}(X, \mathcal{F})$ 中的类。

\medskip\noindent
此时有几件事需要验证：(a) $\delta(\xi)$ 不依赖于 $\tau$ 的选择；
(b) $\delta(\xi)$ 不依赖于使 $\sigma$ 可提升的超覆盖 $L \to K$
的选择；(c) $\delta(\xi)$ 不依赖于起初为表示 $\xi$ 所选的超覆盖和
$\sigma$。我们略去 (a)、(b)、(c) 的验证；其中对超覆盖选择的独立性，
实质上归结为引理 \ref{hypercovering-lemma-product-hypercoverings} 和
\ref{hypercovering-lemma-homotopy}。我们还略去验证 $\delta$ 关于 $\mathcal{C}$ 上
阿贝尔层短正合列之间的态射具有函子性。

\medskip\noindent
最后，必须验证在这个 $\delta$ 的定义下，上述阿贝尔层短正合列给出
{\v C}ech 上同调群的长正合列。先证明：若
$\delta(\xi) = 0$（$\xi$ 如上），则 $\xi$ 是某个元素
$\xi' \in \check{H}^p_{HC}(X, \mathcal{G})$ 的像。事实上，若
$\delta(\xi) = 0$，则沿用上述记号，$d(\tau)$ 的类在
$\check{H}^{p + 1}_{HC}(X, \mathcal{F})$ 中为零。因此存在超覆盖的态射
$L \to K$，使 $d(\tau)$ 限制所得的 $\mathcal{F}(L_{p + 1})$ 中的元素，
等于某个 $\upsilon \in \mathcal{F}(L_p)$ 的 $d(\upsilon)$。
这蕴含 $\tau|_{L_p} + \upsilon$ 构成闭链，并确定一个类
$\xi' \in \check{H}^p(L, \mathcal{G})$，它按要求映到 $\xi$。

\medskip\noindent
我们略去以下断言的证明：若
$\xi' \in \check{H}^{p + 1}_{HC}(X, \mathcal{F})$ 在
$\check{H}^{p + 1}_{HC}(X, \mathcal{G})$ 中映为零，则对某个
$\xi \in \check{H}^p_{HC}(X, \mathcal{H})$ 有
$\xi' = \delta(\xi)$。
\end{proof}

\noindent
下面用一个巧妙手法推出定理 \ref{hypercovering-theorem-cohomology-hypercoverings}
的 Verdier 情形。

\begin{proposition}
\label{hypercovering-proposition-cohomology-hypercoverings}
设 $\mathcal{C}$ 是具有纤维积和二元积的位点，
$\mathcal{F}$ 是 $\mathcal{C}$ 上的阿贝尔层，$i \geq 0$。则
\begin{enumerate}
\item 对每个 $\xi \in H^i(\mathcal{F})$，存在超覆盖 $K$，使 $\xi$
属于典范映射 $\check{H}^i(K, \mathcal{F}) \to H^i(\mathcal{F})$ 的像；
\item 若 $K,L$ 是超覆盖，且
$\xi_K \in \check{H}^i(K, \mathcal{F})$、
$\xi_L \in \check{H}^i(L, \mathcal{F})$ 映到
$H^i(\mathcal{F})$ 的同一元素，则存在超覆盖 $M$ 以及态射
$M \to K$ 和 $M \to L$，使 $\xi_K$ 与 $\xi_L$ 映到
$\check{H}^i(M, \mathcal{F})$ 的同一元素。
\end{enumerate}
换言之，忽略集合论问题，$\mathcal{F}$ 在 $\mathcal{C}$ 上的上同调群，
是 $\mathcal{F}$ 遍历所有超覆盖的 {\v C}ech 上同调群的余极限。
\end{proposition}

\begin{proof}
此结果是定理 \ref{hypercovering-theorem-cohomology-hypercoverings} 的直接推论。
事实上，可以人为地用一个稍大的位点 $\mathcal{C}'$ 替换
$\mathcal{C}$，使得 (I) $\mathcal{C}'$ 有终对象 $X$；(II)
$\mathcal{C}$ 中的超覆盖与 $\mathcal{C}'$ 中 $X$ 的超覆盖大致相同。
但由于问题本身的性质，需要处理相当多的记账细节。

\medskip\noindent
若 $\mathcal{C}$ 中具有固定靶的态射族 $\{U_i \to U\}$ 所对应的映射
$\coprod_{i \in I} h_{U_i} \to h_U$ 在层化后成为满射，则称它为
{\it 弱覆盖}。如下构造新位点 $\mathcal{C}'$：
\begin{enumerate}
\item 作为范畴，置
$\Ob(\mathcal{C}') = \Ob(\mathcal{C}) \amalg \{X\}$，
并从 $\mathcal{C}'$ 的每个对象向 $X$ 添加唯一态射；
\item 当 $\mathcal{C}$ 中存在纤维积和二元积时，
$\mathcal{C}'$ 具有相应的纤维积；
\item $\mathcal{C}'$ 的覆盖包括 $\mathcal{C}$ 的弱覆盖，以及满足如下条件的
$\{U_i \to X\}_{i \in I}$：或者某个 $i$ 有 $U_i = X$；或者对所有
$i$ 都有 $U_i \not = X$，且 $\mathcal{C}$ 上预层的映射
$\coprod h_{U_i} \to *$ 在 $\mathcal{C}$ 上层化后成为满射；
\item 应用《集合论》引理 \ref{sets-lemma-coverings-site} 限制覆盖，
从而得到位点 $\mathcal{C}'$。
\end{enumerate}
于是 $\Sh(\mathcal{C}') = \Sh(\mathcal{C})$，因为包含函子
$\mathcal{C} \to \mathcal{C}'$ 是特殊余连续函子
（见《位点与层》定义 \ref{sites-definition-special-cocontinuous-functor}）。
直接的验证从略。

\medskip\noindent
选取 $\mathcal{C}'$ 的覆盖 $\{U_i \to X\}$，使对所有 $i$，$U_i$ 都是
$\mathcal{C}$ 的对象（这是可能的，因为
$\mathcal{C} \to \mathcal{C}'$ 特殊余连续）。于是
$K_0 = \{U_i \to X\}$ 是上述位点 $\mathcal{C}'$ 中的覆盖。
把 $K_0$ 看成 $\text{SR}(\mathcal{C}', X)$ 的对象，并置
$K_{init} = \text{cosk}_0(K_0)$。于是 $K_{init}$ 是 $X$ 的超覆盖，
见例 \ref{hypercovering-example-cech}。注意，每个 $K_{init, n}$ 都具有
$\{W_j \to X\}$ 的形状，其中 $W_j \in \Ob(\mathcal{C})$。

\medskip\noindent
证明 (1)。选取
$\xi \in H^i(\mathcal{F}) = H^i(X, \mathcal{F}')$，其中
$\mathcal{F}'$ 是 $\mathcal{C}'$ 上与 $\mathcal{C}$ 上的
$\mathcal{F}$ 对应的阿贝尔层。由定理
\ref{hypercovering-theorem-cohomology-hypercoverings}，存在 $\mathcal{C}'$ 中 $X$
的超覆盖的态射 $K' \to K_{init}$，使 $\xi$ 来自
$\check{H}^i(K', \mathcal{F})$ 的一个元素。写
$K'_n = \{U_{n, j} \to X\}$。由于 $K'_n$ 映到
$K_{init, n}$，可见 $U_{n, j}$ 是 $\mathcal{C}$ 的对象。因此置
$K_n = \{U_{n, j}\}$，便可定义 $\text{SR}(\mathcal{C})$ 的单纯对象
$K$。由于 $\mathcal{C}'$ 中由 $\mathcal{C}$ 的态射族构成的覆盖，
是弱覆盖，可见 $K$ 是定义 \ref{hypercovering-definition-hypercovering-variant}
意义下的超覆盖。最后，由于 $\mathcal{F}'$ 是 $\mathcal{C}'$ 上
限制到 $\mathcal{C}$ 等于 $\mathcal{F}$ 的唯一层，{\v C}ech 复形
$s(\mathcal{F}(K))$ 与 $s(\mathcal{F}'(K'))$ 相同，于是 (1) 成立。
（略去与到上同调群的映射之相容性的验证。）

\medskip\noindent
证明 (2)。设 $K,L$ 是 $\mathcal{C}$ 中的超覆盖。令 $K',L'$ 为通过函子
$\text{SR}(\mathcal{C}) \to \text{SR}(\mathcal{C}', X)$,
$\{U_i\} \mapsto \{U_i \to X\}$ 由 $K,L$ 得到的
$\text{SR}(\mathcal{C}', X)$ 的单纯对象。与前面一样，{\v C}ech 复形相等，
于是得到 $\xi_{K'}$ 和 $\xi_{L'}$，它们映到 $\mathcal{C}'$ 上
$\mathcal{F}'$ 的同一上同调类。由于集合论问题，必要时扩大
$\mathcal{C}'$ 中覆盖的选择后，可以假设 $K',L'$ 是
$\mathcal{C}'$ 中 $X$ 的超覆盖；这由定义
\ref{hypercovering-definition-hypercovering-variant} 中的超覆盖定义，以及
$\mathcal{C}$ 中的弱覆盖给出 $\mathcal{C}'$ 中的覆盖这一事实成立。
由定理 \ref{hypercovering-theorem-cohomology-hypercoverings}，存在
$\mathcal{C}'$ 中 $X$ 的超覆盖 $M'$ 以及态射
$M' \to K'$、$M' \to L'$ 和 $M' \to K_{init}$，使
$\xi_{K'}$ 与 $\xi_{L'}$ 限制到
$\check{H}^i(M', \mathcal{F})$ 的同一元素。像上面一样展开此陈述，
便得 (2)。
\end{proof}





\section{空间的超覆盖}
\label{hypercovering-section-hypercoverings-spaces}

\noindent
即使对拓扑空间，上述理论也颇有趣。在这种情形下，可以具体写出超覆盖，
并看清结果实际表达的内容。

\medskip\noindent
设 $X$ 是拓扑空间。考虑《位点与层》例
\ref{sites-example-site-topological} 中的位点 $X_{Zar}$。
回忆，$X_{Zar}$ 的对象就是 $X$ 的开集，而 $X_{Zar}$ 的态射就是包含。
那么，位点 $X_{Zar}$ 上 $X$ 的超覆盖是什么？

\medskip\noindent
先展开定义 \ref{hypercovering-definition-SR}。$\text{SR}(X_{Zar}, X)$ 的一个对象，
就是一个集合 $I$ 以及对每个 $i \in I$ 给定的开集 $U_i \subset X$。
由于态射 $U_i \to X$ 不会混淆，把它记为 $\{U_i\}_{i \in I}$。
两个这样的对象之间的态射
$\{U_i\}_{i \in I} \to \{V_j\}_{j \in J}$，由集合映射
$\alpha : I \to J$ 给出，并满足对所有 $i \in I$ 有
$U_i \subset V_{\alpha(i)}$。这样的态射何时是覆盖？当且仅当对每个
$j \in J$ 都有
$V_j = \bigcup_{i\in I, \ \alpha(i) = j} U_i$（并且这是位点
$X_{Zar}$ 中的覆盖）。

\medskip\noindent
由以上内容得到位点 $X_{Zar}$ 中超覆盖的如下描述。
$X_{Zar}$ 中 $X$ 的一个超覆盖由以下数据给出：
\begin{enumerate}
\item 一个单纯集 $I$（见《单纯方法》第
\ref{simplicial-section-simplicial-set} 节）；
\item 对每个 $n \geq 0$ 和每个 $i \in I_n$，一个开集 $U_i \subset X$。
\end{enumerate}
用 $(I, \{U_i\})$ 表示这样一组数据。为使它成为 $X$ 的超覆盖，
要求满足以下性质：
\begin{itemize}
\item 对 $i \in I_n$ 和 $0 \leq a \leq n$，有
$U_i \subset U_{d^n_a(i)}$；
\item 对 $i \in I_n$ 和 $0 \leq a \leq n$，有
$U_i = U_{s^n_a(i)}$；
\item 有
\begin{equation}
\label{hypercovering-equation-covering-X}
X = \bigcup\nolimits_{i \in I_0} U_i,
\end{equation}
\item 对每个 $i_0, i_1 \in I_0$，有
\begin{equation}
\label{hypercovering-equation-covering-two}
U_{i_0} \cap U_{i_1} =
\bigcup\nolimits_{i \in I_1, \ d^1_0(i) = i_0, \ d^1_1(i) = i_1} U_i,
\end{equation}
\item 对每个 $n \geq 1$ 以及每个满足
$d^n_{b - 1}(i_a) = d^n_a(i_b)$（$0\leq a < b\leq n + 1$）的
$(i_0, \ldots, i_{n + 1}) \in (I_n)^{n + 2}$，有
\begin{equation}
\label{hypercovering-equation-covering-general}
U_{i_0} \cap \ldots \cap U_{i_{n + 1}} =
\bigcup\nolimits_{i \in I_{n + 1},
\ d^{n + 1}_a(i) = i_a, \ a = 0, \ldots, n + 1} U_i,
\end{equation}
\item 开覆盖 (\ref{hypercovering-equation-covering-X})、
(\ref{hypercovering-equation-covering-two}) 和 (\ref{hypercovering-equation-covering-general})
都属于 $\text{Cov}(X_{Zar})$（这是一个集合论条件，用来限制覆盖的指标集大小）。
\end{itemize}
例如，条件 (\ref{hypercovering-equation-covering-X}) 和
(\ref{hypercovering-equation-covering-two}) 应当可从《空间上的层》一章中认出，
而条件 (\ref{hypercovering-equation-covering-general}) 是其自然推广。

\begin{remark}
\label{hypercovering-remark-not-covering-set}
这个描述的一个特点是：若某个多重交
$U_{i_0} \cap \ldots \cap U_{i_{n + 1}}$ 为空，则右侧的覆盖可以是空覆盖。
因此，映射 $I_{n + 1} \to (\text{cosk}_n\text{sk}_n I)_{n + 1}$
不一定自动为满射。这意味着 $I$ 的几何实现可能是一个有趣的
（非可缩）空间。

\medskip\noindent
事实上，令 $I'_n \subset I_n$ 为由满足 $U_i \not = \emptyset$ 的单形
$i \in I_n$ 组成的子集。容易看出 $I' \subset I$ 是单纯子集，
且 $(I', \{U_i\})$ 是超覆盖。因此总能把一个超覆盖细化为所有开集
$U_i$ 都非空的超覆盖。
\end{remark}

\begin{remark}
\label{hypercovering-remark-repackage-into-simplicial-space}
再用另一种方式重新组织这些信息。设 $(I, \{U_i\})$ 是拓扑空间
$X$ 的超覆盖。由这些数据，置
$$
U_n = \coprod\nolimits_{i \in I_n} U_i,
$$
可构造单纯拓扑空间 $U_\bullet$；对给定的
$\varphi : [n] \to [m]$，令态射 $U(\varphi) : U_n \to U_m$
为来自各个 $i \in I_n$ 所对应包含 $U_i \subset U_{\varphi(i)}$ 的态射。
这个单纯拓扑空间带有增广 $\epsilon : U_\bullet \to X$。
凭借此态射，单纯空间 $U_\bullet$ 成为 $X$ 的超覆盖；沿此超覆盖，
存在 \cite[Expos\'e Vbis]{SGA4} 意义下的上同调下降。换言之，
$H^n(U_\bullet, \epsilon^*\mathcal{F}) = H^n(X, \mathcal{F})$。
（此处以后插入关于单纯空间上的上同调，以及用这种语言表述的上同调下降的参考。）
设 $\mathcal{F}$ 是 $X$ 上的阿贝尔层。此时，引理
\ref{hypercovering-lemma-cech-spectral-sequence} 的谱序列成为具有如下 $E_1$-项的谱序列：
$$
E_1^{p, q} = H^q(U_p, \epsilon_q^*\mathcal{F})
\Rightarrow
H^{p + q}(U_\bullet, \epsilon^*\mathcal{F}) = H^{p + q}(X, \mathcal{F})
$$
它把 $\epsilon^*\mathcal{F}$ 的全上同调与 $\mathcal{F}$ 在
$U_\bullet$ 各片上的上同调群作比较。（此处以后插入关于此谱序列的参考。）
\end{remark}

\noindent
在拓扑学中，我们常希望找到 $X$ 的超覆盖，使所有 $U_i$ 都来自
$X$ 的一个给定拓扑基，而且所有覆盖 (\ref{hypercovering-equation-covering-two})
和 (\ref{hypercovering-equation-covering-general}) 都来自一个给定的共终覆盖族。
下面给出两个示例性引理。

\begin{lemma}
\label{hypercovering-lemma-basis-hypercovering}
设 $X$ 是拓扑空间，$\mathcal{B}$ 是 $X$ 的拓扑基。
则存在 $X$ 的超覆盖 $(I, \{U_i\})$，使每个 $U_i$ 都属于
$\mathcal{B}$。
\end{lemma}

\begin{proof}
设 $n \geq 0$。所谓 $X$ 的一个{\it $n$-截断超覆盖}，是由一个
$n$-截断单纯集 $I$，以及对每个 $i \in I_a$、$0 \leq a \leq n$
给定的 $X$ 的开集 $U_i$ 所组成，并要求定义超覆盖的条件凡有意义时都成立。
换言之，只有当其中出现的所有单形都是满足 $a \leq n$ 的
$a$-单形时，才要求相应的包含关系和覆盖条件。若能证明：给定所有
$U_i \in \mathcal{B}$ 的 $n$-截断超覆盖 $(I, \{U_i\})$，可以在不添加
任何 $a \leq n$ 的 $a$-单形的情况下，把它扩张为
$(n + 1)$-截断超覆盖，那么引理即得证。具体如下。先考虑由
$I' = \text{sk}_{n + 1}(\text{cosk}_n I)$.
定义的 $(n + 1)$-截断单纯集 $I'$。回忆，
$$
I'_{n + 1} =
\left\{
\begin{matrix}
(i_0, \ldots, i_{n + 1}) \in (I_n)^{n + 2} \text{，使得}\\
d^n_{b - 1}(i_a) = d^n_a(i_b) \text{ 对所有 }0\leq a < b\leq n + 1
\end{matrix}
\right\}
$$
若 $i' \in I'_{n + 1}$ 是退化的，比如 $i' = s^n_a(i)$，则置
$U_{i'} = U_i$（无论如何，第二个条件已迫使我们这样做）。
此时还置 $J_{i'} = \{i'\}$。若 $i' \in I'_{n + 1}$ 是非退化的，
比如 $i' = (i_0, \ldots, i_{n + 1})$，则选取集合 $J_{i'}$
和开覆盖
\begin{equation}
\label{hypercovering-equation-choose-covering}
U_{i_0} \cap \ldots \cap U_{i_{n + 1}} =
\bigcup\nolimits_{i \in J_{i'}} U_i,
\end{equation}
其中对 $i \in J_{i'}$ 有 $U_i \in \mathcal{B}$。置
$$
I_{n + 1} = \coprod\nolimits_{i' \in I'_{n + 1}} J_{i'}
$$
存在典范映射 $\pi : I_{n + 1} \to I'_{n + 1}$；由构造，它在
$I'_{n + 1}$ 的退化单形集合上是双射。对 $i \in I_{n + 1}$，定义
$d^{n + 1}_a(i) = d^{n + 1}_a(\pi(i))$。对 $i \in I_n$，把
$s^n_a(i) \in I_{n + 1}$ 定义为位于退化单形
$s^n_a(i) \in I'_{n + 1}$ 上方的唯一单形。略去验证这定义了
$X$ 的一个 $(n + 1)$-截断超覆盖。
\end{proof}

\begin{lemma}
\label{hypercovering-lemma-quasi-separated-quasi-compact-hypercovering}
设 $X$ 是拓扑空间，$\mathcal{B}$ 是 $X$ 的拓扑基。假设
\begin{enumerate}
\item $X$ 是拟紧的；
\item 每个 $U \in \mathcal{B}$ 都是拟紧开集；
\item $X$ 中任意两个拟紧开集的交仍拟紧。
\end{enumerate}
则存在 $X$ 的超覆盖 $(I, \{U_i\})$，具有以下性质：
\begin{enumerate}
\item 每个 $U_i$ 都属于基 $\mathcal{B}$；
\item 每个 $I_n$ 都是有限集；特别地，
\item 覆盖 (\ref{hypercovering-equation-covering-X})、
(\ref{hypercovering-equation-covering-two}) 和 (\ref{hypercovering-equation-covering-general})
都是有限的。
\end{enumerate}
\end{lemma}

\begin{proof}
在引理 \ref{hypercovering-lemma-basis-hypercovering} 证明的构造中，若在
(\ref{hypercovering-equation-choose-covering}) 处选取由 $\mathcal{B}$ 中元素组成的
有限覆盖，便直接得到结论。细节从略。
\end{proof}







\section{构造超覆盖}
\label{hypercovering-section-hypercovering-sites}

\noindent
设 $\mathcal{C}$ 是位点。本节按如下方式看待
$\text{SR}(\mathcal{C})$ 的单纯对象。照例置 $K_n = K([n])$，
并用 $K(\varphi) : K_n \to K_m$ 表示与
$\varphi : [m] \to [n]$ 相伴的态射。可以写
$K_n = \{U_{n, i}\}_{i \in I_n}$。对
$\varphi : [m] \to [n]$，态射 $K(\varphi) : K_n \to K_m$
由映射 $\alpha(\varphi) : I_n \to I_m$ 以及对 $i \in I_n$ 的态射
$f_{\varphi, i} : U_{n, i} \to U_{m, \alpha(\varphi)(i)}$ 给出。
$K$ 是 $\text{SR}(\mathcal{C})$ 的单纯对象这一事实蕴含
$(I_n, \alpha(\varphi))$ 是单纯集，而且当
$\psi : [l] \to [m]$ 时，
$f_{\psi, \alpha(\varphi)(i)} \circ f_{\varphi, i} =
f_{\varphi \circ \psi, i}$。

\begin{lemma}
\label{hypercovering-lemma-split}
设 $\mathcal{C}$ 是位点，$K$ 是 $\text{SR}(\mathcal{C})$ 的
$r$-截断单纯对象。以下条件等价：
\begin{enumerate}
\item $K$ 是分裂的（《单纯方法》定义
\ref{simplicial-definition-split}）；
\item $f_{\varphi, i} : U_{n, i} \to U_{m, \alpha(\varphi)(i)}$
对 $r \geq n \geq 0$、满射 $\varphi : [m] \to [n]$ 和
$i \in I_n$ 是同构；
\item $f_{\sigma^n_j, i} : U_{n, i} \to U_{n + 1, \alpha(\sigma^n_j)(i)}$
对 $0 \leq j \leq n < r$ 和 $i \in I_n$ 是同构。
\end{enumerate}
对单纯对象也有同样结论，只需在 (2) 和 (3) 中置 $r = \infty$。
\end{lemma}

\begin{proof}
单纯集的分裂是唯一的，并由各次数 $n$ 中的非退化指标
$N(I_n)$ 给出，见《单纯方法》引理
\ref{simplicial-lemma-splitting-simplicial-sets}。
$\text{SR}(\mathcal{C})$ 的两个对象
$\{U_i\}_{i \in I}$ 和 $\{U_j\}_{j \in J}$ 的余积，按显然记号由
$\{U_l\}_{l \in I \amalg J}$ 给出。因此 $K$ 的分裂必由
$N(K_n) = \{U_i\}_{i \in N(I_n)}$ 给出。展开定义便得到 (1) 与 (2)
的等价性。任意满射 $\varphi : [m] \to [n]$ 都是态射
$\sigma^k_j$（$k = n, n + 1, \ldots, m - 1$）的复合，故 (2) 与 (3)
等价。
\end{proof}

\begin{lemma}
\label{hypercovering-lemma-hypercovering-object}
设 $\mathcal{C}$ 是具有纤维积的位点，
$\mathcal{B} \subset \Ob(\mathcal{C})$ 是子集。假设
\begin{enumerate}
\item $\mathcal{C}$ 的任意对象 $U$ 都有一个满足
$U_j \in \mathcal{B}$ 的覆盖 $\{U_j \to U\}_{j \in J}$；
\item 若 $\{U_j \to U\}_{j \in J}$ 是满足
$U_j \in \mathcal{B}$ 的覆盖，且 $\{U' \to U\}$ 是满足
$U' \in \mathcal{B}$ 的态射，则
$\{U_j \to U\}_{j \in J} \amalg \{U' \to U\}$ 是覆盖。
\end{enumerate}
则对 $\mathcal{C}$ 中的任意 $X$，存在 $X$ 的超覆盖 $K$，使得
$K_n = \{U_{n, i}\}_{i \in I_n}$，并且对所有 $i \in I_n$ 都有
$U_{n, i} \in \mathcal{B}$。
\end{lemma}

\begin{proof}
引理 \ref{hypercovering-lemma-basis-hypercovering} 的证明是本证明的预备练习；
建议读者先阅读那个证明。

\medskip\noindent
先用位点 $\mathcal{C}/X$ 替换 $\mathcal{C}$。这样做以后，可以假设
$X$ 是 $\mathcal{C}$ 的终对象，并且 $\mathcal{C}$ 具有所有有限极限
（《范畴论》引理 \ref{categories-lemma-finite-limits-exist}）。

\medskip\noindent
设 $n \geq 0$。所谓 $X$ 的一个
{\it $n$-截断 $\mathcal{B}$-超覆盖}，是
$\text{SR}(\mathcal{C})$ 的一个 $n$-截断单纯对象 $K$，满足：
对 $i \in I_a$、$0 \leq a \leq n$ 有
$U_{a, i} \in \mathcal{B}$；$K_0$ 是 $X$ 的覆盖；且
$K_{a + 1} \to (\text{cosk}_a \text{sk}_a K)_{a + 1}$
对 $a = 0, \ldots, n - 1$ 是定义
\ref{hypercovering-definition-covering-SR} 意义下的覆盖。

\medskip\noindent
由假设，$X$ 有一个满足 $U_i \in \mathcal{B}$ 的覆盖
$\{U_{0, i} \to X\}_{i \in I_0}$，故得到 $X$ 的一个
$0$-截断 $\mathcal{B}$-超覆盖。注意，$X$ 的任意
$0$-截断 $\mathcal{B}$-超覆盖都是分裂的，见引理 \ref{hypercovering-lemma-split}。

\medskip\noindent
若能证明：对 $n \geq 0$，给定 $X$ 的一个分裂
$n$-截断 $\mathcal{B}$-超覆盖 $K$，可把它扩张为 $X$ 的一个分裂
$(n + 1)$-截断 $\mathcal{B}$-超覆盖，则引理成立。

\medskip\noindent
构造该扩张。考虑 $\text{SR}(\mathcal{C})$ 的 $(n + 1)$-截断单纯对象
$K' = \text{sk}_{n + 1}(\text{cosk}_n K)$。写成
$$
K'_{n + 1} = \{U'_{n + 1, i}\}_{i \in I'_{n + 1}}
$$
由于 $K = \text{sk}_n K'$，对 $0 \leq a \leq n$ 有
$K_a = K'_a$。对每个 $i' \in I'_{n + 1}$，选取覆盖
\begin{equation}
\label{hypercovering-equation-choose-covering-B}
\{g_{n + 1, j} : U_{n + 1, j} \to U'_{n + 1, i'}\}_{j \in J_{i'}}
\end{equation}
使得对 $j \in J_{i'}$ 有 $U_{n + 1, j} \in \mathcal{B}$。
由引理中关于 $\mathcal{B}$ 的假设，这是可能的。对
$0 \leq m \leq n$，用 $N_m \subset I_m$ 表示非退化指标的子集。置
$$
I_{n + 1} =
\coprod\nolimits_{\varphi : [n + 1] \to [m]\text{ 满射，}0\leq m \leq n}
N_m \amalg
\coprod\nolimits_{i' \in I'_{n + 1}} J_{i'}
$$
对 $j \in I_{n + 1}$，置
$$
U_{n + 1, j} =
\left\{
\begin{matrix}
U_{m, i} & \text{若} &
j = (\varphi, i) & \text{其中} & \varphi : [n + 1] \to [m], i \in N_m \\
U_{n + 1, j} & \text{若} &
j \in J_{i'} & \text{其中} & i' \in I'_{n + 1}
\end{matrix}
\right.
$$
记号是显然的。置 $K_{n + 1} = \{U_{n + 1, j}\}_{j \in I_{n + 1}}$。
由构造，对所有 $j \in I_{n + 1}$，$U_{n + 1, j}$ 都属于
$\mathcal{B}$。定义相容映射
$$
I_{n + 1} \to I'_{n + 1}
\quad\text{和}\quad
K_{n + 1} \to K'_{n + 1}
$$
具体地，第一个映射由
$(\varphi, i) \mapsto \alpha'(\varphi)(i)$ 和
$(j \in J_{i'}) \mapsto i'$ 给出；第二个映射使用态射
$$
f'_{\varphi, i} : U_{m, i} \to U'_{n + 1, \alpha'(\varphi)(i)}
\quad\text{和}\quad
g_{n + 1, j} : U_{n + 1, j} \to U'_{n + 1, i'}
$$
我们断言，态射
$$
K_{n + 1} \to K'_{n + 1} =
(\text{cosk}_n \text{sk}_n K')_{n + 1} =
(\text{cosk}_n K)_{n + 1}
$$
是定义 \ref{hypercovering-definition-covering-SR} 意义下的覆盖。事实上，若
$i' \in I'_{n + 1}$，则或者 $i'$ 非退化，此时 $i'$ 在
$I_{n + 1}$ 中的逆像等于 $J_{i'}$，由选择
(\ref{hypercovering-equation-choose-covering-B}) 得到 $U'_{n + 1, i'}$ 的一个覆盖；
或者 $i'$ 退化，此时对唯一的偶对 $(\varphi, i)$，$i'$ 在
$I_{n + 1}$ 中的逆像为 $J_{i'} \amalg \{(\varphi, i)\}$，
由选择 (\ref{hypercovering-equation-choose-covering-B}) 和引理的假设 (2) 得到覆盖。

\medskip\noindent
为完成证明，还需定义与态射
$\varphi : [m] \to [n + 1]$、$0 \leq m \leq n$ 对应的态射
$K(\varphi) : K_{n + 1} \to K_m$，以及与态射
$\varphi : [n + 1] \to [m]$、$0 \leq m \leq n$ 对应的态射
$K(\varphi) : K_m \to K_{n + 1}$，并使它们满足适当的复合关系。
对第一类，使用复合
$$
K_{n + 1} \to K'_{n + 1} \xrightarrow{K'(\varphi)} K'_m = K_m
$$
来定义 $K(\varphi) : K_{n + 1} \to K_m$。对第二类，设给定
$\varphi : [n + 1] \to [m]$、$0 \leq m \leq n$。
如下定义相应态射 $K(\varphi) : K_m \to K_{n + 1}$：
\begin{enumerate}
\item 对 $i \in I_m$，存在唯一满射 $\psi : [m] \to [m_0]$
以及唯一的非退化 $i_0 \in I_{m_0}$，使得
$\alpha(\psi)(i_0) = i$\footnote{例如，若 $i$ 非退化，则
$m = m_0$ 且 $\psi = \text{id}_{[m]}$。}；
\item 置 $\varphi_0 = \psi_0 \circ \varphi : [n + 1] \to [m_0]$，
并把 $i \in I_m$ 映到 $(\varphi_0, i_0) \in I_{n + 1}$；换言之，
$\alpha(\varphi)(i) = (\varphi_0, i_0)$；
\item 态射
$f_{\varphi, i} : U_{m, i} \to U_{n + 1, \alpha(\varphi)(i)} = U_{m_0, i_0}$
是同构 $f_{\psi, i_0} : U_{m_0, i_0} \to U_{m, i}$ 的逆
（见引理 \ref{hypercovering-lemma-split}）。
\end{enumerate}
略去直接但繁琐的验证：这定义了 $X$ 的一个分裂
$(n + 1)$-截断 $\mathcal{B}$-超覆盖，扩张给定的 $n$-截断者。
事实上，除需验证对所有 $0 \leq a,b \leq n + 1$ 的
$\varphi : [a] \to [b]$，态射 $K(\varphi)$ 的复合正确之外，
其余一切均由以上构造显然成立。
\end{proof}

\begin{lemma}
\label{hypercovering-lemma-hypercovering-site}
设 $\mathcal{C}$ 是具有等化子和纤维积的位点，
$\mathcal{B} \subset \Ob(\mathcal{C})$ 是子集。假设
$\mathcal{C}$ 的任意对象都有一个各成员属于 $\mathcal{B}$ 的覆盖。
则存在超覆盖 $K$，使得 $K_n = \{U_i\}_{i \in I_n}$，
并且对所有 $i \in I_n$ 都有 $U_i \in \mathcal{B}$。
\end{lemma}

\begin{proof}
此证明与引理 \ref{hypercovering-lemma-hypercovering-object} 的证明几乎相同。
这里只说明不同之处。

\medskip\noindent
设 $n \geq 1$。所谓一个{\it $n$-截断 $\mathcal{B}$-超覆盖}，
是 $\text{SR}(\mathcal{C})$ 的一个 $n$-截断单纯对象 $K$，满足：
对 $i \in I_a$、$0 \leq a \leq n$ 有
$U_{a, i} \in \mathcal{B}$；并且
\begin{enumerate}
\item $F(K_0)^\# \to *$ 是满射；
\item $F(K_1)^\# \to F(K_0)^\# \times F(K_0)^\#$ 是满射；
\item $F(K_{a + 1})^\# \to F((\text{cosk}_a \text{sk}_a K)_{a + 1})^\#$
对 $a = 1, \ldots, n - 1$ 是满射。
\end{enumerate}
先显式构造一个分裂 $1$-截断 $\mathcal{B}$-超覆盖。

\medskip\noindent
取 $I_0 = \mathcal{B}$ 和 $K_0 = \{U\}_{U \in \mathcal{B}}$。
由关于 $\mathcal{B}$ 的假设，(1) 成立。置
$$
\Omega =
\{(U, V, W, a, b) \mid U, V, W \in \mathcal{B}, a : U \to V, b : U \to W\}
$$
再置 $I_1 = I_0 \amalg \Omega$。对 $i \in I_1$，若 $i \in I_0$，
置 $U_{1, i} = U_{0, i}$；若
$i = (U, V, W, a, b) \in \Omega$，置 $U_{1, i} = U$。
映射 $K(\sigma^0_0) : K_0 \to K_1$ 对应于包含
$\alpha(\sigma^0_0) : I_0 \to I_1$ 和对象上的恒等态射
$f_{\sigma^0_0, i} : U_{0, i} \to U_{1, i}$。映射
$K(\delta^1_0), K(\delta^1_1) : K_1 \to K_0$ 对应于两个映射
$I_1 \to I_0$：它们在 $I_0 \subset I_1$ 上为恒等，并分别把
$(U, V, W, a, b) \in \Omega \subset I_1$ 映到 $V$ 和 $W$。
相应态射 $f_{\delta^1_0, i}, f_{\delta^1_1, i} : U_{1, i} \to U_{0, i}$，
在 $i \in I_0$ 时为恒等；在
$i = (U, V, W, a, b) \in \Omega$ 时分别为 $a,b$。
(2) 成立的理由如下：$F(K_0)^\# \times F(K_0)^\#$ 在
$\mathcal{C}$ 的对象 $U$ 上的任意截面，在用某个覆盖的各成员替换
$U$ 后，都来自映射 $U \to F(K_0) \times F(K_0)$。
这又意味着存在 $V,W \in \mathcal{B}$ 和两个态射
$U \to V$、$U \to W$。进一步用某个覆盖的各成员替换 $U$ 后，
可以按要求假设 $U \in \mathcal{B}$。

\medskip\noindent
若能证明：对 $n \geq 1$，给定一个分裂
$n$-截断 $\mathcal{B}$-超覆盖 $K$，可以把它扩张为一个分裂
$(n + 1)$-截断 $\mathcal{B}$-超覆盖，则引理成立。这里的论证与引理
\ref{hypercovering-lemma-hypercovering-object} 的证明完全相同。除以下说明外，
精确细节从略。第一，在本引理的证明中不需要假设 (2)，因为无需态射
$K_{n + 1} \to (\text{cosk}_n K)_{n + 1}$ 本身是覆盖；只需它在相伴集合层上
诱导满射，而这由《位点与层》引理
\ref{sites-lemma-covering-surjective-after-sheafification} 得出。
第二，$\mathcal{C}$ 具有纤维积和等化子的假设保证
$\text{SR}(\mathcal{C})$ 具有纤维积和等化子，并且 $F$ 与它们交换
（引理 \ref{hypercovering-lemma-coprod-prod-SR}）。这足以保证所用的余骨架函子存在
（见《单纯方法》评注 \ref{simplicial-remark-existence-cosk} 和
《范畴论》引理 \ref{categories-lemma-fibre-products-equalizers-exist}）。
\end{proof}

\begin{lemma}
\label{hypercovering-lemma-hypercovering-morphism-sites}
设 $f : \mathcal{C} \to \mathcal{D}$ 是由函子
$u : \mathcal{D} \to \mathcal{C}$ 给出的位点态射。假设
$\mathcal{D}$ 和 $\mathcal{C}$ 具有等化子和纤维积，且 $u$ 与它们交换。
若 $\text{SR}(\mathcal{D})$ 的单纯对象 $K$ 是超覆盖，
则 $u(K)$ 是超覆盖。
\end{lemma}

\begin{proof}
如本节引言中那样写 $K_n = \{U_{n, i}\}_{i \in I_n}$，则 $u(K)$
是 $\text{SR}(\mathcal{C})$ 的对象，其中
$u(K_n) = \{u(U_i)\}_{i \in I_n}$。由《位点与层》引理
\ref{sites-lemma-pullback-representable-sheaf}，对
$U \in \Ob(\mathcal{D})$ 有 $f^{-1}h_U^\# = h_{u(U)}^\#$。
这意味着对所有 $n$，
$f^{-1}F(K_n)^\# = F(u(K_n))^\#$。验证定义
\ref{hypercovering-definition-hypercovering-variant} 中 $u(K)$ 成为超覆盖的条件
(1)、(2)、(3)。由于 $f^{-1}$ 是正合函子，可得
$$
F(u(K_0))^\# = f^{-1}F(K_0)^\# \to f^{-1}* = *
$$
作为满射的拉回是满射，故得到 (1)。类似地，
$$
F(u(K_1))^\# = f^{-1}F(K_1)^\# \to
f^{-1} (F(K_0) \times F(K_0))^\# = F(u(K_0))^\# \times F(u(K_0))^\#
$$
作为拉回是满射，故得到 (2)。对于条件 (3)，要用相同方法得出结论，
只需有
$$
F((\text{cosk}_n \text{sk}_n u(K))_{n + 1})^\# =
f^{-1}F((\text{cosk}_n \text{sk}_n K)_{n + 1})^\#
$$
以上说明 $f^{-1}F(-) = F(u(-))$。因此只需证明，对 $n \geq 1$，
$u$ 与定义 $(\text{cosk}_n \text{sk}_n K)_{n + 1}$ 时所用的极限交换。
由《单纯方法》评注 \ref{simplicial-remark-existence-cosk}，
这些极限是有限连通极限；由假设，$u$ 与它们交换。
\end{proof}

\begin{lemma}
\label{hypercovering-lemma-hypercovering-continuous-functor}
设 $\mathcal{C},\mathcal{D}$ 是位点，
$u : \mathcal{D} \to \mathcal{C}$ 是连续函子。假设
$\mathcal{D}$ 和 $\mathcal{C}$ 具有纤维积，且 $u$ 与纤维积交换。
设 $Y \in \mathcal{D}$，并设
$K \in \text{SR}(\mathcal{D}, Y)$ 是 $Y$ 的超覆盖。
则 $u(K)$ 是 $u(Y)$ 的超覆盖。
\end{lemma}

\begin{proof}
这比引理 \ref{hypercovering-lemma-hypercovering-morphism-sites} 的证明容易，
因为“是某个对象的超覆盖”这一概念更强，见定义
\ref{hypercovering-definition-hypercovering} 和 \ref{hypercovering-definition-covering-SR}。
事实上，由位点态射的定义，$u$ 把覆盖送到覆盖。只需验证，对
$n \geq 1$，$u$ 与定义
$(\text{cosk}_n \text{sk}_n K)_{n + 1}$ 时所用的极限交换。
这是显然的，因为诱导函子
$\mathcal{D}/Y \to \mathcal{C}/X$ 与所有有限极限交换
（而由《范畴论》引理 \ref{categories-lemma-finite-limits-exist}，
源和靶都具有所有有限极限）。
\end{proof}

\begin{lemma}
\label{hypercovering-lemma-w-contractible}
设 $\mathcal{C}$ 是位点，$\mathcal{B} \subset \Ob(\mathcal{C})$
是子集。假设
\begin{enumerate}
\item $\mathcal{C}$ 具有纤维积；
\item 对所有 $X \in \Ob(\mathcal{C})$，存在满足
$U_i \in \mathcal{B}$ 的有限覆盖 $\{U_i \to X\}_{i \in I}$；
\item 若 $\{U_i \to X\}_{i \in I}$ 是满足
$U_i \in \mathcal{B}$ 的有限覆盖，且 $U \to X$ 是满足
$U \in \mathcal{B}$ 的态射，则
$\{U_i \to X\}_{i \in I} \amalg \{U \to X\}$ 是覆盖。
\end{enumerate}
则对每个 $X$，存在 $X$ 的超覆盖 $K$，使每个
$K_n = \{U_{n, i} \to X\}_{i \in I_n}$ 中 $I_n$ 有限，且
$U_{n, i} \in \mathcal{B}$。
\end{lemma}

\begin{proof}
此引理是引理
\ref{hypercovering-lemma-quasi-separated-quasi-compact-hypercovering} 的位点版本。
证明时完全依照引理 \ref{hypercovering-lemma-hypercovering-object} 的证明，
并注意以下两点：
\begin{enumerate}
\item[(a)] 选取初始覆盖 $\{U_{0, i} \to X\}_{i \in I_0}$，
使 $U_{0, i} \in \mathcal{B}$ 且指标集 $I_0$ 有限；
\item[(b)] 选取覆盖 (\ref{hypercovering-equation-choose-covering-B}) 时，
取 $J_{i'}$ 为有限集。
\end{enumerate}
读者容易看出，经过这些修改，最终对所有 $n$ 都得到有限指标集 $I_n$。
\end{proof}

\begin{remark}
\label{hypercovering-remark-taking-disjoint-unions}
设 $\mathcal{C}$ 是位点，$K,L$ 是 $\text{SR}(\mathcal{C})$ 的对象。
写 $K = \{U_i\}_{i \in I}$、$L = \{V_j\}_{j \in J}$。
假设 $U = \coprod_{i \in I} U_i$ 和
$V = \coprod_{j \in J} V_j$ 存在。于是得到
$$
\Mor_{\text{SR}(\mathcal{C})}(K, L) \longrightarrow \Mor_\mathcal{C}(U, V)
$$
如下。给定由 $\alpha : I \to J$ 和
$f_i : U_i \to V_{\alpha(i)}$ 给出的 $f : K \to L$，得到函子的变换
$$
\Mor_\mathcal{C}(V, -) =
\prod\nolimits_{j \in J} \Mor_\mathcal{C}(V_j, -)
\to
\prod\nolimits_{i \in I} \Mor_\mathcal{C}(U_i, -) =
\Mor_\mathcal{C}(U, -)
$$
它把 $(g_j)_{j \in J}$ 送到
$(g_{\alpha(i)} \circ f_i)_{i \in I}$。因此 Yoneda 引理给出相应映射
$U \to V$。当然，$U \to V$ 通过态射 $f_i$ 把直和项 $U_i$
映入直和项 $V_{\alpha(i)}$。
\end{remark}

\begin{remark}
\label{hypercovering-remark-take-unions-hypercovering}
设 $\mathcal{C}$ 是具有纤维积和等化子的位点，$K$ 是超覆盖。
写 $K_n = \{U_{n, i}\}_{i \in I_n}$。假设
\begin{enumerate}
\item[(a)] $U_n = \coprod_{i \in I_n} U_{n, i}$ 存在；
\item[(b)] $\coprod_{i \in I_n} h_{U_{n, i}} \to h_{U_n}$
在层化后诱导同构。
\end{enumerate}
于是得到 $\text{SR}(\mathcal{C})$ 的另一个单纯对象 $L$，其中
$L_n = \{U_n\}$，见评注 \ref{hypercovering-remark-taking-disjoint-unions}。
我们断言 $L$ 是超覆盖。为此验证定义
\ref{hypercovering-definition-hypercovering-variant} 的条件 (1)、(2)、(3)。
条件 (1) 由 (b) 和 $K$ 的条件 (1) 得出；条件 (2) 完全类似。
条件 (3) 由下式得出：
\begin{align*}
F((\text{cosk}_n \text{sk}_n L)_{n + 1})^\#
& =
((\text{cosk}_n \text{sk}_n F(L)^\#)_{n + 1}) \\
& =
((\text{cosk}_n \text{sk}_n F(K)^\#)_{n + 1}) \\
& =
F((\text{cosk}_n \text{sk}_n K)_{n + 1})^\#
\end{align*}
其中 $n \geq 1$；因此完全如 (1)、(2) 那样，$K$ 的相应条件蕴含
$L$ 的相应条件。注意，$F$ 与连通极限交换且层化正合，这证明第一和
最后一个等式；中间的等式由 (b) 中
$F(K)^\# = F(L)^\#$ 得出。
\end{remark}

\begin{remark}
\label{hypercovering-remark-take-unions-hypercovering-X}
设 $\mathcal{C}$ 是位点，$X \in \Ob(\mathcal{C})$。假设
$\mathcal{C}$ 具有纤维积，并设 $K$ 是 $X$ 的超覆盖。
写 $K_n = \{U_{n, i}\}_{i \in I_n}$。假设
\begin{enumerate}
\item[(a)] $U_n = \coprod_{i \in I_n} U_{n, i}$ 存在；
\item[(b)] 给定 $\text{SR}(\mathcal{C})$ 中的态射
$(\alpha, f_i) : \{U_i\}_{i \in I} \to \{V_j\}_{j \in J}$ 和
$(\beta, g_k) : \{W_k\}_{k \in K} \to \{V_j\}_{j \in J}$
，且
$U = \coprod U_i$、$V = \coprod V_j$ 和 $W = \coprod W_j$
存在，则 $U \times_V W =
\coprod_{(i, j, k), \alpha(i) = j = \beta(k)} U_i \times_{V_j} W_k$；
\item[(c)] 若 $(\alpha, f_i) : \{U_i\}_{i \in I} \to \{V_j\}_{j \in J}$
是定义 \ref{hypercovering-definition-covering-SR} 意义下的覆盖，且
$U = \coprod U_i$ 和 $V = \coprod V_j$ 存在，则评注
\ref{hypercovering-remark-taking-disjoint-unions} 中相应的态射 $U \to V$
是 $\mathcal{C}$ 的覆盖。
\end{enumerate}
于是得到 $\text{SR}(\mathcal{C})$ 的另一个单纯对象 $L$，其中
$L_n = \{U_n\}$，见评注 \ref{hypercovering-remark-taking-disjoint-unions}。
我们断言 $L$ 是 $X$ 的超覆盖。为此验证定义
\ref{hypercovering-definition-hypercovering} 的条件 (1)、(2)。条件 (1) 由 (c)
和 $K$ 的条件 (1) 得出，因为 $K$ 的条件 (1) 表明
$K_0 = \{U_{0, i}\}_{i \in I_0}$ 是定义
\ref{hypercovering-definition-covering-SR} 意义下 $\{X\}$ 的覆盖。
条件 (2) 成立，是因为 $\mathcal{C}/X$ 具有所有有限极限，故
$\text{SR}(\mathcal{C}/X)$ 也具有所有有限极限，而条件 (b) 表明
“取不交并”的构造与这些有限极限交换。因此态射
$$
L_{n + 1} \longrightarrow (\text{cosk}_n \text{sk}_n L)_{n + 1}
$$
是覆盖，因为它是把“取不交并”函子应用于态射
$$
K_{n + 1} \longrightarrow (\text{cosk}_n \text{sk}_n K)_{n + 1}
$$
所得的结果；由 $K$ 的条件 (2)，后者被假设为定义
\ref{hypercovering-definition-covering-SR} 意义下的覆盖。这一说法有意义，
因为性质 (b) 特别保证：若从 $X$ 上一张可取不交并的半可表对象有限图出发，
则该图在 $\text{SR}(\mathcal{C}/X)$ 中的极限仍是 $X$ 上可取不交并的
半可表对象。
\end{remark}
% Generated mechanically from the frozen schemes.tex target.
% Do not edit this generated file; edit the bound locale source instead.

% OK, start here.
%

\setcounter{chapter}{25}
\chapter{概形}



\phantomsection
\label{schemes-section-phantom}


\section{引言}
\label{schemes-section-introduction}

\noindent
本文定义概形。
一部基本参考文献是 \cite{EGA}。









\section{局部环层空间}
\label{schemes-section-locally-ringed-spaces}

\noindent
回顾《层》第 \ref{sheaves-section-ringed-spaces} 节中定义的环层空间。
简言之，环层空间是一个二元组 $(X, \mathcal{O}_X)$，其中 $X$ 是拓扑空间，
$\mathcal{O}_X$ 是环层。环层空间之间的态射
$f : (X, \mathcal{O}_X) \to (Y, \mathcal{O}_Y)$ 由一个连续映射
$f : X \to Y$ 和一个环层的 $f$-映射
$f^\sharp : \mathcal{O}_Y \to \mathcal{O}_X$ 给出。也可以把 $f^\sharp$
看作映射 $\mathcal{O}_Y \to f_*\mathcal{O}_X$；参见《层》的定义
\ref{sheaves-definition-f-map} 和引理 \ref{sheaves-lemma-f-map}。

\medskip\noindent
一个值得牢记的几何例子是 $\mathcal{C}^\infty$-流形及
$\mathcal{C}^\infty$-流形之间的态射。
具体地，若 $M$ 是 $\mathcal{C}^\infty$-流形，则光滑函数层
$\mathcal{C}^\infty_M$ 是 $M$ 上的环层。流形间的任意映射
$f : M \to N$ 是光滑的，当且仅当对每个局部截面 $h$——这里指
$\mathcal{C}^\infty_N$ 的局部截面——复合 $h \circ f$ 都是
$\mathcal{C}^\infty_M$ 的局部截面。
因此，光滑映射 $f$ 自然给出环层空间之间的态射
$$
f : (M , \mathcal{C}^\infty_M) \longrightarrow (N, \mathcal{C}^\infty_N)
$$
；参见《层》的例 \ref{sheaves-example-continuous-map-ringed}。
考察茎上的情形很有启发性。设 $m \in M$，且其像为 $f(m) = n \in N$。
回顾茎 $\mathcal{C}^\infty_{M, m}$ 是 $m$ 处光滑函数芽组成的环；
参见《层》的例 \ref{sheaves-example-germs-functions}。
$(M, m)$ 上的函数芽代数是一个局部环，其极大理想由在 $m$ 处消失的函数组成。
$\mathcal{C}^\infty_{N, n}$ 的情形相同。茎上的映射
$f^\sharp : \mathcal{C}^\infty_{N, n} \to \mathcal{C}^\infty_{M, m}$
把极大理想映入极大理想；这直接源于 $f(m) = n$。

\medskip\noindent
代数几何研究概形。概形上的环层并不由其底空间的某种内蕴性质决定。
给环 $R$ 的谱（见《代数》第 \ref{algebra-section-spectrum-ring} 节）
配上由 $R$ 构造出的环层（见下文），便得到我们的基本构件。
稍后将看到，$\mathcal{O}$ 在 $\Spec(R)$ 上的各个茎正是 $R$ 在相应素理想处的局部环。
在这一背景下引入局部环层空间有两个原因：
(1) 一般而言，没有一种机制能由谱之间的连续映射自动给出相应环之间的映射，
因此我们把 $f^\sharp$ 作为附加数据；
(2) 若在环层空间范畴中考虑这些谱之间的态射，则其茎上的映射未必是局部同态，
而几何直观要求它们应当如此。为此作如下定义。

\begin{definition}
\label{schemes-definition-locally-ringed-space}
局部环层空间。
\begin{enumerate}
\item {\it 局部环层空间 $(X, \mathcal{O}_X)$} 是一个二元组，其中
$X$ 是拓扑空间，$\mathcal{O}_X$ 是环层，并且它的每个茎都是局部环。
\item 给定局部环层空间 $(X, \mathcal{O}_X)$，称
$\mathcal{O}_{X, x}$ 为 {\it $X$ 在 $x$ 处的局部环}。
以 $\mathfrak{m}_{X, x}$（或简记为 $\mathfrak{m}_x$）表示
$\mathcal{O}_{X, x}$ 的极大理想。此外，{\it $X$ 在 $x$ 处的剩余域} 是
$\kappa(x) = \mathcal{O}_{X, x}/\mathfrak{m}_x$。
\item {\it 局部环层空间之间的态射}
$(f, f^\sharp) : (X, \mathcal{O}_X) \to (Y, \mathcal{O}_Y)$
是满足下述条件的环层空间态射：对每个 $x \in X$，诱导的环同态
$\mathcal{O}_{Y, f(x)} \to \mathcal{O}_{X, x}$ 都是局部环同态。
\end{enumerate}
\end{definition}

\noindent
讨论局部环层空间时，我们通常在记号中略去环层 $\mathcal{O}_X$，
而只说“局部环层空间 $X$”。我们还滥用记号，以 $X$ 同时表示其底拓扑空间。
相应的环层 $\mathcal{O}_X$ 称为 {\it $X$ 的结构层}。
此外，局部环 $\mathcal{O}_{X, x}$ 的极大理想通常记作
$\mathfrak{m}_{X, x}$，或简记作 $\mathfrak{m}_x$。
我们会说“设 $f : X \to Y$ 是局部环层空间之间的态射”，
从而也略去结构层。此时又滥用记号，以 $f : X\to Y$ 同时表示底拓扑空间之间的连续映射。
与 $f$ 对应的 $f$-映射通常记作 $f^\sharp$。
$f$ 是局部环层空间态射这一条件，可以表述为：对每个 $x\in X$，茎上的映射
$$
f^\sharp_x : \mathcal{O}_{Y, f(x)} \longrightarrow \mathcal{O}_{X, x}
$$
把极大理想 $\mathfrak m_{Y, f(x)}$ 映入 $\mathfrak m_{X, x}$。

\medskip\noindent
下面利用这些记号约定说明，局部环层空间及其态射构成一个范畴。
为此只须验证，两个局部环层空间态射的复合仍是局部环层空间态射。
设 $f : X \to Y$ 和 $g : Y \to Z$ 是局部环层空间态射。
$f$ 与 $g$ 的复合定义于《层》的定义
\ref{sheaves-definition-composition-maps-ringed-spaces}。
取 $x \in X$。由《层》的引理 \ref{sheaves-lemma-compose-f-maps-stalks}，复合映射
$$
\mathcal{O}_{Z, g(f(x))}
\xrightarrow{g^\sharp}
\mathcal{O}_{Y, f(x)}
\xrightarrow{f^\sharp}
\mathcal{O}_{X, x}
$$
就是态射 $g \circ f$ 所诱导的茎上映射。由于局部环同态的复合仍是局部环同态，
结论随即成立。

\medskip\noindent
这个定义有一个令人满意的性质：函子
$$
\textit{局部环化空间}
\longrightarrow
\textit{环化空间}
$$
反映同构（事实上还不止于此）。下面给出一个较为具体的表述。

\begin{lemma}
\label{schemes-lemma-isomorphism-locally-ringed}
\begin{slogan}
局部环层空间之间的环层空间同构也是局部环层空间同构。
\end{slogan}
设 $X$、$Y$ 为局部环层空间。若 $f : X \to Y$ 是环层空间同构，
则 $f$ 是局部环层空间同构。
\end{lemma}

\begin{proof}
这直接归结为代数中的相应事实：若 $A$、$B$ 是局部环，
则任意环同构 $A \to B$ 都是局部环同构。
\end{proof}













\section{局部环层空间的开浸入}
\label{schemes-section-open-immersion}

\begin{definition}
\label{schemes-definition-immersion-locally-ringed-spaces}
设 $f : X \to Y$ 是局部环层空间态射。称 $f$ 为 {\it 开浸入}，如果
$f$ 将 $X$ 同胚地映到 $Y$ 的一个开子集上，并且映射
$f^{-1}\mathcal{O}_Y \to \mathcal{O}_X$ 是同构。
\end{definition}

\noindent
下述构造与《层》的定义 \ref{sheaves-definition-restriction} (3) 平行。

\begin{example}
\label{schemes-example-open-subspace}
设 $X$ 为局部环层空间，$U \subset X$ 为开子集，并令
$\mathcal{O}_U = \mathcal{O}_X|_U$ 为 $\mathcal{O}_X$ 在 $U$ 上的限制。
对 $u \in U$，茎 $\mathcal{O}_{U, u}$ 等于茎 $\mathcal{O}_{X, u}$，因而是局部环。
所以 $(U, \mathcal{O}_U)$ 是局部环层空间，且态射
$j : (U, \mathcal{O}_U) \to (X, \mathcal{O}_X)$ 是开浸入。
\end{example}

\begin{definition}
\label{schemes-definition-open-subspace}
设 $X$ 为局部环层空间，$U \subset X$ 为开子集。
上面例 \ref{schemes-example-open-subspace} 中的局部环层空间
$(U, \mathcal{O}_U)$ 称为 {\it $X$ 的与 $U$ 相伴的开子空间}。
\end{definition}

\begin{lemma}
\label{schemes-lemma-open-immersion}
设 $f : X \to Y$ 是局部环层空间的开浸入，并令
$j : V = f(X) \to Y$ 为 $Y$ 的与 $f$ 的像相伴的开子空间。
存在唯一的局部环层空间同构 $f' : X \cong V$，使得 $f = j \circ f'$。
\end{lemma}

\begin{proof}
令 $f'$ 为 $X$ 与 $V$ 之间由 $f$ 诱导的同胚。作为拓扑空间之间的映射，
$f = j \circ f'$。由于存在层同构
$f^\sharp : f^{-1}(\mathcal{O}_Y) \to \mathcal{O}_X$，故有环同构
$f^\sharp : \Gamma(U, f^{-1}(\mathcal{O}_Y)) \to \Gamma(U, \mathcal{O}_X)$，
其中 $U \subset X$ 为任意开子集
。由于 $\mathcal{O}_V = j^{-1}\mathcal{O}_Y$ 且 $f^{-1} = f'^{-1} j^{-1}$
（《层》的引理 \ref{sheaves-lemma-pullback-composition}），可得
$f^{-1}\mathcal{O}_Y = f'^{-1}\mathcal{O}_V$，从而映射
$\Gamma(U, f'^{-1}(\mathcal{O}_V)) \to \Gamma(U, f^{-1}(\mathcal{O}_Y))$
对每个开集 $U \subset X$ 都是同构。将这些同构复合，便对每个开子集
$U \subset X$ 得到环同构
$$
\Gamma(U, f'^{-1}(\mathcal{O}_V)) \to \Gamma(U, \mathcal{O}_X)
$$
，因而得到层同构 $f^{-1}(\mathcal{O}_V) \to \mathcal{O}_X$。
换言之，我们得到同构 $f'^{\sharp} : f'^{-1}(\mathcal{O}_V) \to \mathcal{O}_X$，
从而得到局部环层空间同构
$(f', f'^{\sharp}) : (X, \mathcal{O}_X) \to (V, \mathcal{O}_V)$
（使用引理 \ref{schemes-lemma-isomorphism-locally-ringed}）。按构造，
作为局部环层空间态射也有 $f = j \circ f'$。

\medskip\noindent
再设有态射 $f'' : (X, \mathcal{O}_X) \to (V, \mathcal{O}_V)$，满足
$f = j \circ f''$。对任意 $x \in X$，有 $j(f'(x)) = j(f''(x))$，由此推出
$f'(x) = f''(x)$，因为 $j$ 是包含映射。因此 $f'$ 与 $f''$ 作为拓扑空间态射相同。
在结构层上，对每个开子集 $U \subset X$，有交换图
$$
\xymatrix @R=5em{
\Gamma(U, f^{-1}(\mathcal{O}_Y)) \ar[d]_\cong\ar[r]^\cong &
\Gamma(U, \mathcal{O}_X) \\
\Gamma(U, f'^{-1}(\mathcal{O}_V)) \ar@/^/[ru]^{f'^\sharp}
\ar@/_/[ru]_{f''^\sharp} &
}
$$
由此可见，$f'^\sharp$ 与 $f''^\sharp$ 定义同一个层态射。
\end{proof}

\noindent
从现在起，我们不再区分开子集与其相伴的开子空间。

\begin{lemma}
\label{schemes-lemma-restrict-map-to-opens}
设 $f : X \to Y$ 是局部环层空间态射，$U \subset X$、$V \subset Y$ 为开子集，
并设 $f(U) \subset V$。存在唯一的局部环层空间态射 $f|_U : U \to V$，
使下图成为局部环层空间的交换方块：
$$
\xymatrix{
U \ar[d]_{f|_U} \ar[r] & X \ar[d]^f \\
V \ar[r] & Y
}
$$
\end{lemma}

\begin{proof}
略。
\end{proof}

\noindent
下文将不加说明地使用由上引理得到的如下事实。给定任意局部环层空间态射
$f : Y \to X$，以及开子集 $U \subset X$，并设它满足 $f(Y) \subset U$，
存在唯一的局部环层空间态射 $Y \to U$，使复合 $Y \to U \to X$ 等于 $f$。
事实上，我们甚至会滥用记号直接写成 $f : Y \to U$，因为这很少引起混淆。









\section{局部环层空间的闭浸入}
\label{schemes-section-closed-immersion}

\noindent
沿用《模》的定义 \ref{modules-definition-closed-immersion} 中引入的约定。

\begin{definition}
\label{schemes-definition-closed-immersion-locally-ringed-spaces}
设 $i : Z \to X$ 是局部环层空间态射。若满足下列条件，则称 $i$ 为
{\it 闭浸入}：
\begin{enumerate}
\item 映射 $i$ 将 $Z$ 同胚地映到 $X$ 的一个闭子集上。
\item 映射 $\mathcal{O}_X \to i_*\mathcal{O}_Z$ 是满射；以
$\mathcal{I}$ 表示其核。
\item $\mathcal{O}_X$-模 $\mathcal{I}$ 局部地由截面生成。
\end{enumerate}
\end{definition}

\begin{lemma}
\label{schemes-lemma-closed-local-target}
设 $f : Z \to X$ 是局部环层空间态射。要使 $f$ 为闭浸入，只需存在开覆盖
$X = \bigcup U_i$，使每个 $f : f^{-1}U_i \to U_i$ 都是闭浸入。
\end{lemma}

\begin{proof}
略。
\end{proof}

\begin{example}
\label{schemes-example-closed-subspace}
设 $X$ 是局部环层空间，$\mathcal{I} \subset \mathcal{O}_X$ 是理想层，
并且作为 $\mathcal{O}_X$-模层局部地由截面生成。令 $Z$ 为环层
$\mathcal{O}_X/\mathcal{I}$ 的支撑。由《模》的引理
\ref{modules-lemma-support-sheaf-rings-closed}，它是 $X$ 的闭子集。
以 $i : Z \to X$ 表示包含映射。由《模》的引理
\ref{modules-lemma-i-star-exact}，存在唯一的环层 $\mathcal{O}_Z$ 定义在 $Z$ 上，
满足 $i_*\mathcal{O}_Z = \mathcal{O}_X/\mathcal{I}$。
对任意 $z \in Z$，茎 $\mathcal{O}_{Z, z}$ 等于局部环的非零商
$\mathcal{O}_{X, i(z)}/\mathcal{I}_{i(z)}$，因而仍是局部环。
所以 $i : (Z, \mathcal{O}_Z) \to (X, \mathcal{O}_X)$ 是局部环层空间的闭浸入。
\end{example}

\begin{definition}
\label{schemes-definition-closed-subspace}
设 $X$ 是局部环层空间，$\mathcal{I}$ 是 $X$ 上局部地由截面生成的理想层。
上面例 \ref{schemes-example-closed-subspace} 中的局部环层空间
$(Z, \mathcal{O}_Z)$ 称为 {\it $X$ 的与理想层 $\mathcal{I}$ 相伴的闭子空间}。
\end{definition}

\begin{lemma}
\label{schemes-lemma-closed-immersion}
设 $f : X \to Y$ 是局部环层空间的闭浸入，$\mathcal{I}$ 是映射
$\mathcal{O}_Y \to f_*\mathcal{O}_X$ 的核，并令 $i : Z \to Y$ 为
$Y$ 的与 $\mathcal{I}$ 相伴的闭子空间。存在唯一的局部环层空间同构
$f' : X \cong Z$，使得 $f = i \circ f'$。
\end{lemma}

\begin{proof}
略。
\end{proof}

\begin{lemma}
\label{schemes-lemma-characterize-closed-subspace}
设 $X$、$Y$ 为局部环层空间，$\mathcal{I} \subset \mathcal{O}_X$ 是局部地由截面生成的理想层，
且 $i : Z \to X$ 是相伴的闭子空间。态射 $f : Y \to X$ 经 $Z$ 分解，
当且仅当映射 $f^*\mathcal{I} \to f^*\mathcal{O}_X = \mathcal{O}_Y$ 为零。
在这种情形下，态射 $g : Y \to Z$ 若满足 $f = i \circ g$，则它是唯一的。
\end{lemma}

\begin{proof}
显然，若 $f$ 分解为 $Y \to Z \to X$，则映射
$f^*\mathcal{I} \to \mathcal{O}_Y$ 为零。反过来，设
$f^*\mathcal{I} \to \mathcal{O}_Y$ 为零。任取 $y \in Y$，考虑环同态
$f^\sharp_y : \mathcal{O}_{X, f(y)} \to \mathcal{O}_{Y, y}$。
由于复合映射
$\mathcal{I}_{f(y)} \to \mathcal{O}_{X, f(y)} \to \mathcal{O}_{Y, y}$
按假设为零，且 $f^\sharp_y(1) = 1$，可知
$1 \not \in \mathcal{I}_{f(y)}$，即
$\mathcal{I}_{f(y)} \not = \mathcal{O}_{X, f(y)}$。因此
$f(Y) \subset Z = \text{Supp}(\mathcal{O}_X/\mathcal{I})$。
于是 $f = i \circ g$，其中 $g : Y \to Z$ 连续。考虑映射
$f^\sharp : \mathcal{O}_X \to f_*\mathcal{O}_Y$。
假设 $f^*\mathcal{I} \to \mathcal{O}_Y$ 为零，意味着复合
$\mathcal{I} \to \mathcal{O}_X \to f_*\mathcal{O}_Y$ 为零；
这是由 $f_*$ 与 $f^*$ 的伴随性得到的。
换言之，我们得到环层态射
$\overline{f^\sharp} : \mathcal{O}_X/\mathcal{I} \to f_*\mathcal{O}_Y$。
注意 $f_*\mathcal{O}_Y = i_*g_*\mathcal{O}_Y$，且
$\mathcal{O}_X/\mathcal{I} = i_*\mathcal{O}_Z$。
由《层》的引理 \ref{sheaves-lemma-equivalence-categories-closed-structures}，
得到唯一的环层态射 $g^\sharp : \mathcal{O}_Z \to g_*\mathcal{O}_Y$，
其沿 $i$ 的正像为 $\overline{f^\sharp}$。我们略去如下验证：
$(g, g^\sharp)$ 定义局部环层空间态射，并且作为局部环层空间态射有
$f = i \circ g$。$(g, g^\sharp)$ 的唯一性已在上面指出。
\end{proof}

\begin{lemma}
\label{schemes-lemma-restrict-map-to-closed}
设 $f : X \to Y$ 是局部环层空间态射，
$\mathcal{I} \subset \mathcal{O}_Y$ 是局部地由截面生成的理想层，
且 $i : Z \to Y$ 是与 $\mathcal{I}$ 相伴的闭子空间。
令 $\mathcal{J}$ 为映射
$f^*\mathcal{I} \to f^*\mathcal{O}_Y = \mathcal{O}_X$ 的像。
则这个理想局部地由截面生成。再令 $i' : Z' \to X$ 为相伴的 $X$ 的闭子空间。
存在唯一的局部环层空间态射 $f' : Z' \to Z$，使下图成为局部环层空间的交换方块：
$$
\xymatrix{
Z' \ar[d]_{f'} \ar[r]_{i'} & X \ar[d]^f \\
Z \ar[r]^{i} & Y
}
$$
此外，这个图在局部环层空间范畴中是纤维方块。
\end{lemma}

\begin{proof}
由《模》的引理 \ref{modules-lemma-pullback-locally-generated}，
理想 $\mathcal{J}$ 局部地由截面生成。引理的其余部分由上面的引理
\ref{schemes-lemma-characterize-closed-subspace} 对“态射经闭子空间分解”的刻画得到。
\end{proof}














\section{仿射概形}
\label{schemes-section-affine-schemes}

\noindent
设 $R$ 为环。考虑拓扑空间 $\Spec(R)$，它与 $R$ 相伴；参见《代数》第
\ref{algebra-section-spectrum-ring} 节。我们将在这个空间上赋予一个环层
$\mathcal{O}_{\Spec(R)}$，所得二元组
$(\Spec(R), \mathcal{O}_{\Spec(R)})$ 将是一个仿射概形。

\medskip\noindent
回顾 $\Spec(R)$ 有一个由开集 $D(f)$（$f \in R$）组成的基；这些开集称为
标准开集，参见《代数》的定义 \ref{algebra-definition-Zariski-topology}。
此外，两个标准开集的交仍为标准开集：
$D(f) \cap D(g) = D(fg)$，其中 $f, g\in R$。

\begin{lemma}
\label{schemes-lemma-standard-open}
设 $R$ 为环，$f \in R$。
\begin{enumerate}
\item 若 $g\in R$ 且 $D(g) \subset D(f)$，则
\begin{enumerate}
\item $f$ 在 $R_g$ 中可逆；
\item $g^e = af$，其中 $e \geq 1$ 且 $a \in R$；
\item 存在典范环同态 $R_f \to R_g$；
\item 存在典范的 $R_f$-模同态 $M_f \to M_g$，适用于任意 $R$-模 $M$。
\end{enumerate}
\item $D(f)$ 的任意开覆盖都可加细为形如
$D(f) = \bigcup_{i = 1}^n D(g_i)$ 的有限开覆盖。
\item 若 $g_1, \ldots, g_n \in R$，则 $D(f) \subset \bigcup D(g_i)$
当且仅当 $g_1, \ldots, g_n$ 在 $R_f$ 中生成单位理想。
\end{enumerate}
\end{lemma}

\begin{proof}
回顾 $D(g) = \Spec(R_g)$（见《代数》的引理
\ref{algebra-lemma-standard-open}）。(a) 成立，是因为 $f$ 在 $R_g$ 中的像
不属于任何素理想，因而可逆；参见《代数》的引理
\ref{algebra-lemma-Zariski-topology}。把 $f$ 在 $R_g$ 中的逆元写成
$a/g^d$。这意味着 $g^d - af$ 被 $g$ 的某个幂零化，从而得到 (b)。
对于 (c)，由 (a) 及局部化的泛性质，映射 $R_f \to R_g$ 存在；
也可以把 $b/f^n$ 映到 $a^nb/g^{ne}$ 来定义它。利用等式
$M_f = M \otimes_R R_f$ 可以得到模之间的映射 $M_f \to M_g$；
也可以把 $x/f^n$ 映到 $a^nx/g^{ne}$ 来定义它。

\medskip\noindent
回顾 $D(f)$ 是拟紧的，参见《代数》的引理
\ref{algebra-lemma-qc-open}。由于标准开集构成这个拓扑的一组基，
第二个断言立即成立。

\medskip\noindent
第三个断言直接来自《代数》的引理
\ref{algebra-lemma-Zariski-topology}。
\end{proof}

\noindent
在《层》第 \ref{sheaves-section-bases} 节中，我们定义了基上的层，
并说明这一概念与空间上的层本质上等价；参见《层》的引理
\ref{sheaves-lemma-extend-off-basis} 和
\ref{sheaves-lemma-extend-off-basis-structures}。此外，《层》的引理
\ref{sheaves-lemma-cofinal-systems-coverings-standard-case} 表明，
对每个标准开集，只须在一个开覆盖的共尾系上检验层条件。
由上面的引理，只须对标准开集构成的有限覆盖检验这一条件。

\begin{definition}
\label{schemes-definition-standard-covering}
设 $R$ 为环。
\begin{enumerate}
\item $\Spec(R)$ 的一个 {\it 标准开覆盖} 是形如
$\Spec(R) = \bigcup_{i = 1}^n D(f_i)$ 的覆盖，其中
$f_1, \ldots, f_n \in R$。
\item 设 $D(f) \subset \Spec(R)$ 是标准开集。$D(f)$ 的一个
{\it 标准开覆盖} 是形如 $D(f) = \bigcup_{i = 1}^n D(g_i)$ 的覆盖，
其中 $g_1, \ldots, g_n \in R$。
\end{enumerate}
\end{definition}

\noindent
设 $R$ 为环，$M$ 为 $R$-模。我们将在标准开集构成的基上定义预层
$\widetilde M$。设 $U \subset \Spec(R)$ 是标准开集。若 $f, g \in R$
满足 $D(f) = D(g)$，则由上面的引理 \ref{schemes-lemma-standard-open}，
存在互为逆映射的典范映射 $M_f \to M_g$ 和 $M_g \to M_f$。
因此可以任选 $f$，使其满足 $U = D(f)$，并定义
$$
\widetilde M(U) = M_f.
$$
注意，若 $D(g) \subset D(f)$，则由上面的引理
\ref{schemes-lemma-standard-open}，存在典范映射
$$
\widetilde M(D(f)) = M_f \longrightarrow M_g = \widetilde M(D(g)).
$$
显然，这在标准开集构成的基上定义了一个阿贝尔群预层。若 $M = R$，
则 $\widetilde R$ 是该基上的环预层。

\medskip\noindent
下面计算 $\widetilde M$ 在点 $x \in \Spec(R)$ 处的茎。设 $x$ 对应于素理想
$\mathfrak p \subset R$。由茎的定义可得
$$
\widetilde M_x = \colim_{f\in R, f\not\in \mathfrak p} M_f
$$
这里，集合 $\{f \in R, f \not \in \mathfrak p\}$ 由规则
$f \geq f' \Leftrightarrow D(f) \subset D(f')$ 赋予预序。
若 $f_1, f_2 \in R \setminus \mathfrak p$，则在这个预序中
$f_1f_2 \geq f_1$。因此，由《代数》的引理
\ref{algebra-lemma-localization-colimit} 可得
$$
\widetilde M_x = M_{\mathfrak p}.
$$

\medskip\noindent
接下来检验标准开覆盖的层条件。若
$D(f) = \bigcup_{i = 1}^n D(g_i)$，则这个覆盖的层条件等价于下列序列的正合性：
$$
0 \to M_f \to \bigoplus M_{g_i} \to \bigoplus M_{g_ig_j}.
$$
注意 $D(g_i) = D(fg_i)$，因此可将这个序列改写为
$$
0 \to M_f \to \bigoplus M_{fg_i} \to \bigoplus M_{fg_ig_j}.
$$
此外，由上面的引理 \ref{schemes-lemma-standard-open}，
$g_1, \ldots, g_n$ 在 $R_f$ 中生成单位理想。因此可把《代数》的引理
\ref{algebra-lemma-cover-module} 应用于模 $M_f$、环 $R_f$ 以及元素
$g_1, \ldots, g_n$。由此可知该序列正合。根据前面的说明，
$\widetilde M$ 是标准开集构成的基上的一个层。

\medskip\noindent
由《层》第 \ref{sheaves-section-bases} 节的内容可知，存在唯一的环层
$\mathcal{O}_{\Spec(R)}$，它在标准开集上与 $\widetilde R$ 一致。
注意，由上面对茎的计算，这个环层的所有茎都是局部环。

\medskip\noindent
类似地，对任意 $R$-模 $M$，存在唯一的 $\mathcal{O}_{\Spec(R)}$-模层
$\mathcal{F}$，它在标准开集上与 $\widetilde M$ 一致；参见《层》的引理
\ref{sheaves-lemma-extend-off-basis-module}。

\begin{definition}
\label{schemes-definition-structure-sheaf}
设 $R$ 为环。
\begin{enumerate}
\item {\it 结构层 $\mathcal{O}_{\Spec(R)}$，即 $R$ 的谱的结构层}，是唯一的环层
$\mathcal{O}_{\Spec(R)}$，它在标准开集构成的基上与 $\widetilde R$ 一致。
\item 局部环层空间 $(\Spec(R), \mathcal{O}_{\Spec(R)})$ 称为 $R$ 的
{\it 谱}，记作 $\Spec(R)$。
\item 这个 $\mathcal{O}_{\Spec(R)}$-模层把 $\widetilde M$ 扩张到
$\Spec(R)$ 的全部开集上；它称为一个 $\mathcal{O}_{\Spec(R)}$-模层，
且与 $M$ 相伴。这个层仍记作 $\widetilde M$。
\end{enumerate}
\end{definition}

\noindent
下面汇总至此所得的结果。

\begin{lemma}
\label{schemes-lemma-spec-sheaves}
设 $R$ 为环，$M$ 为 $R$-模，并设 $\widetilde M$ 是一个
$\mathcal{O}_{\Spec(R)}$-模层，它与 $M$ 相伴。
\begin{enumerate}
\item 有 $\Gamma(\Spec(R), \mathcal{O}_{\Spec(R)}) = R$。
\item 有 $\Gamma(\Spec(R), \widetilde M) = M$，且这是作为 $R$-模的等式。
\item 对每个 $f \in R$，有
$\Gamma(D(f), \mathcal{O}_{\Spec(R)}) = R_f$。
\item 对每个 $f\in R$，有 $\Gamma(D(f), \widetilde M) = M_f$，
且这是作为 $R_f$-模的等式。
\item 只要 $D(g) \subset D(f)$，$\mathcal{O}_{\Spec(R)}$ 和
$\widetilde M$ 上的限制映射就是引理 \ref{schemes-lemma-standard-open} 中的映射
$R_f \to R_g$ 和 $M_f \to M_g$。
\item 设 $\mathfrak p$ 是 $R$ 的素理想，$x \in \Spec(R)$ 是相应的点。
则有 $\mathcal{O}_{\Spec(R), x} = R_{\mathfrak p}$。
\item 设 $\mathfrak p$ 是 $R$ 的素理想，$x \in \Spec(R)$ 是相应的点。
有 $\widetilde M_x = M_{\mathfrak p}$，且这是作为 $R_{\mathfrak p}$-模的等式。
\end{enumerate}
此外，所有这些识别都关于 $R$-模 $M$ 具有函子性。特别地，函子
$M \mapsto \widetilde M$ 是从 $R$-模范畴到
$\mathcal{O}_{\Spec(R)}$-模范畴的正合函子。
\end{lemma}

\begin{proof}
断言 (1)--(7) 由上面的讨论显然成立。函子 $M \mapsto \widetilde M$
的正合性来自以下两个事实：函子 $M \mapsto M_{\mathfrak p}$ 是正合的，
而短正合列的正合性可以在茎上检验；参见《模》的引理
\ref{modules-lemma-abelian}。
\end{proof}

\begin{definition}
\label{schemes-definition-affine-scheme}
{\it 仿射概形} 是这样的局部环层空间：它作为局部环层空间同构于
$\Spec(R)$，其中 $R$ 是某个环。{\it 仿射概形之间的态射} 是局部环层空间范畴中的态射。
\end{definition}

\noindent
仿射概形在所有局部环层空间中起着特殊作用；下一节将讨论这一点。


















\section{仿射概形范畴}
\label{schemes-section-category-affine-schemes}

\noindent
注意，若 $Y$ 是仿射概形，则存在典范的 $1-1$ 双射，
把它的点对应到 $\Gamma(Y, \mathcal{O}_Y)$ 中的素理想。

\begin{lemma}
\label{schemes-lemma-morphism-into-affine-where-point-goes}
设 $X$ 为局部环层空间，$Y$ 为仿射概形，并设
$f \in \Mor(X, Y)$ 是局部环层空间态射。给定点 $x \in X$，考虑环同态
$$
\Gamma(Y, \mathcal{O}_Y) \xrightarrow{f^\sharp}
\Gamma(X, \mathcal{O}_X) \to \mathcal{O}_{X, x}
$$
以 $\mathfrak p \subset \Gamma(Y, \mathcal{O}_Y)$ 表示 $\mathfrak m_x$ 的原像。
设 $y \in Y$ 是相应的点，则 $f(x) = y$。
\end{lemma}

\begin{proof}
考虑交换图
$$
\xymatrix{
\Gamma(X, \mathcal{O}_X) \ar[r] &
\mathcal{O}_{X, x} \\
\Gamma(Y, \mathcal{O}_Y) \ar[r] \ar[u] &
\mathcal{O}_{Y, f(x)} \ar[u]
}
$$
（参见《层》的定义 \ref{sheaves-definition-f-map} 下方关于 $f$-映射的讨论）。
由于右侧竖直箭头是局部同态，可知 $\mathfrak m_{f(x)}$ 是
$\mathfrak m_x$ 的原像。结论随即成立。
\end{proof}

\begin{lemma}
\label{schemes-lemma-f-open}
设 $X$ 为局部环层空间，$f \in \Gamma(X, \mathcal{O}_X)$。集合
$$
D(f) = \{x \in X \mid \text{像 }f \not\in \mathfrak m_x\}
$$
是开集。此外，$f|_{D(f)}$ 有逆元。
\end{lemma}

\begin{proof}
这是《模》的引理 \ref{modules-lemma-s-open} 的特殊情形；这里也给出直接证明。
设 $U \subset X$ 和 $V \subset X$ 是两个开子集，且 $f|_U$ 有逆元 $g$，
$f|_V$ 有逆元 $h$。显然 $g|_{U\cap V} = h|_{U\cap V}$。
因此，只须说明 $f$ 在任意 $x \in D(f)$ 的某个开邻域上可逆。
这是显然的，因为 $f \not \in \mathfrak m_x$ 意味着
$f \in \mathcal{O}_{X, x}$ 有逆元 $g \in \mathcal{O}_{X, x}$；
这又意味着存在某个开邻域 $x \in U \subset X$，使得
$g \in \mathcal{O}_X(U)$ 且 $g\cdot f|_U = 1$。
\end{proof}

\begin{lemma}
\label{schemes-lemma-f-open-affine}
在上面的引理 \ref{schemes-lemma-f-open} 中，若 $X$ 是仿射概形，
则开集 $D(f)$ 与前面定义的标准开集 $D(f)$ 相同
（见《代数》的定义 \ref{algebra-definition-spectrum-ring}）。
\end{lemma}

\begin{proof}
略。
\end{proof}

\begin{lemma}
\label{schemes-lemma-morphism-into-affine}
\begin{reference}
这一事实可参见 \cite[II, Err 1, Prop. 1.8.1]{EGA}，其中将其归于 J. Tate。
\end{reference}
设 $X$ 为局部环层空间，$Y$ 为仿射概形。映射
$$
\Mor(X, Y)
\longrightarrow
\Hom(\Gamma(Y, \mathcal{O}_Y), \Gamma(X, \mathcal{O}_X))
$$
把 $f$ 映到全局截面上的 $f^\sharp$；这个映射是双射。
\end{lemma}

\begin{proof}
由于 $Y$ 是仿射概形，有
$(Y, \mathcal{O}_Y) \cong (\Spec(R), \mathcal{O}_{\Spec(R)})$
，其中 $R$ 是某个环。在证明中，我们将使用关于 $Y$ 及其结构层的一些事实；它们直接来自已知的环谱性质，
例如引理 \ref{schemes-lemma-spec-sheaves}。

\medskip\noindent
受上述引理启发，我们构造逆映射。设
$\psi_Y : \Gamma(Y, \mathcal{O}_Y) \to \Gamma(X, \mathcal{O}_X)$
是环同态。首先定义相应的空间映射
$$
\Psi : X \longrightarrow Y
$$
，其规则由引理 \ref{schemes-lemma-morphism-into-affine-where-point-goes} 给出。
换言之，给定 $x \in X$，定义 $\Psi(x)$ 为 $Y$ 中与下述素理想相应的点：
该素理想位于 $\Gamma(Y, \mathcal{O}_Y)$ 中，并且是 $\mathfrak m_x$
在如下复合映射下的原像：
$
\Gamma(Y, \mathcal{O}_Y) \xrightarrow{\psi_Y}
\Gamma(X, \mathcal{O}_X) \to
\mathcal{O}_{X, x}
$.

\medskip\noindent
我们断言映射 $\Psi : X \to Y$ 连续。标准开集 $D(g)$（其中
$g \in \Gamma(Y, \mathcal{O}_Y)$）构成 $Y$ 的拓扑的一组基。
因此只须证明 $\Psi^{-1}(D(g))$ 是开集。按 $\Psi$ 的构造，原像
$\Psi^{-1}(D(g))$ 恰为集合 $D(\psi_Y(g)) \subset X$；
由引理 \ref{schemes-lemma-f-open}，它是开集。因此 $\Psi$ 连续。

\medskip\noindent
接着构造一个层的 $\Psi$-映射，它从 $\mathcal{O}_Y$ 到 $\mathcal{O}_X$。
由《层》的引理 \ref{sheaves-lemma-f-map-basis-below-structures}，
只须定义与限制映射相容的环同态
$\psi_{D(g)} : \Gamma(D(g), \mathcal{O}_Y) \to
\Gamma(\Psi^{-1}(D(g)), \mathcal{O}_X)$。
存在典范同构
$\Gamma(D(g), \mathcal{O}_Y) = \Gamma(Y, \mathcal{O}_Y)_g$，
因为 $Y$ 是仿射概形。
由于 $\psi_Y(g)$ 在 $D(\psi_Y(g))$ 上可逆，存在典范映射
$$
\Gamma(Y, \mathcal{O}_Y)_g
\longrightarrow
\Gamma(\Psi^{-1}(D(g)), \mathcal{O}_X)
=
\Gamma(D(\psi_Y(g)), \mathcal{O}_X)
$$
；由局部化的泛性质，它扩张映射 $\psi_Y$。注意，这里只能取这个典范映射！
将它与典范识别
$\Gamma(D(g), \mathcal{O}_Y) = \Gamma(Y, \mathcal{O}_Y)_g$ 结合，
定义为 $\psi_{D(g)}$。这与局部化相容，因为仿射概形上的限制映射同样由
局部化的泛性质定义；参见引理 \ref{schemes-lemma-spec-sheaves} 和
\ref{schemes-lemma-standard-open}。

\medskip\noindent
这样就定义了环层空间态射
$(\Psi, \psi) : (X, \mathcal{O}_X) \to (Y, \mathcal{O}_Y)$
，它在全局截面上给回 $\psi_Y$。要说明它是局部环层空间态射，
必须证明所诱导的局部环同态
$$
\psi_x : \mathcal{O}_{Y, \Psi(x)} \longrightarrow \mathcal{O}_{X, x}
$$
都是局部的。这立即来自引理
\ref{schemes-lemma-morphism-into-affine-where-point-goes} 的证明中的交换图以及
$\Psi$ 的定义。

\medskip\noindent
最后要说明，构造 $(\Psi, \psi) \mapsto \psi_Y$ 与构造
$\psi_Y \mapsto (\Psi, \psi)$ 互为逆构造。显然，
$\psi_Y \mapsto (\Psi, \psi) \mapsto \psi_Y$。因此只须证明：给定
$\psi_Y$，至多有一个二元组 $(\Psi, \psi)$ 产生它。
$\Psi$ 的唯一性已在引理 \ref{schemes-lemma-morphism-into-affine-where-point-goes}
中证明；给定 $\Psi$ 的唯一性后，映射 $\psi$ 的唯一性也已在上面的证明过程中指出。
\end{proof}

\begin{lemma}
\label{schemes-lemma-category-affine-schemes}
仿射概形范畴与环范畴的反范畴等价。这个等价由如下函子给出：
它把仿射概形映到其结构层的全局截面。
\end{lemma}

\begin{proof}
这由定义 \ref{schemes-definition-affine-scheme} 和引理
\ref{schemes-lemma-morphism-into-affine} 立即得到。
\end{proof}

\begin{lemma}
\label{schemes-lemma-standard-open-affine}
设 $Y$ 为仿射概形，$f \in \Gamma(Y, \mathcal{O}_Y)$。
开子空间 $D(f)$ 是仿射概形。
\end{lemma}

\begin{proof}
可以假设 $Y = \Spec(R)$ 且 $f \in R$。考虑仿射概形态射
$\phi : U = \Spec(R_f) \to \Spec(R) = Y$，它由环同态
$R \to R_f$ 诱导。由《代数》的引理
\ref{algebra-lemma-standard-open}，它同胚地映到 $D(f)$ 上。
另一方面，映射 $\phi^{-1}\mathcal{O}_Y \to \mathcal{O}_U$ 在茎上是同构，
因而本身是同构。所以 $\phi$ 是开浸入。由引理 \ref{schemes-lemma-open-immersion}，
$D(f)$ 同构于 $U$。
\end{proof}

\begin{lemma}
\label{schemes-lemma-fibre-product-affine-schemes}
仿射概形范畴有有限积和纤维积。换言之，它有有限极限。
此外，仿射概形范畴中的积和纤维积与局部环层空间范畴中的相同。
用公式表示，在局部环层空间范畴中有
$$
\Spec(R) \times \Spec(S) =
\Spec(R \otimes_{\mathbf{Z}} S)
$$
；而给定环同态 $R \to A$、$R \to B$，有
$$
\Spec(A) \times_{\Spec(R)} \Spec(B)
=
\Spec(A \otimes_R B).
$$
\end{lemma}

\begin{proof}
这只是引理 \ref{schemes-lemma-morphism-into-affine} 的一个应用。
首先，由该引理，仿射概形 $\Spec(\mathbf{Z})$ 是局部环层空间范畴中的终对象。
因此第一个展示公式由第二个推出。为证明第二个公式，注意对任意局部环层空间 $X$，有
\begin{eqnarray*}
\Mor(X, \Spec(A \otimes_R B))
& = &
\Hom(A \otimes_R B, \mathcal{O}_X(X)) \\
& = &
\Hom(A, \mathcal{O}_X(X))
\times_{\Hom(R, \mathcal{O}_X(X))}
\Hom(B, \mathcal{O}_X(X)) \\
& = &
\Mor(X, \Spec(A))
\times_{\Mor(X, \Spec(R))}
\Mor(X, \Spec(B))
\end{eqnarray*}
这就证明了该公式。有关定义参见《范畴》第
\ref{categories-section-fibre-products} 节。
\end{proof}

\begin{lemma}
\label{schemes-lemma-disjoint-union-affines}
设 $X$ 为局部环层空间。假设 $X = U \amalg V$，其中 $U$ 和 $V$ 是开集，
且 $U$、$V$ 都是仿射概形，则 $X$ 是仿射概形。
\end{lemma}

\begin{proof}
令 $R = \Gamma(X, \mathcal{O}_X)$。由层性质，
$R = \mathcal{O}_X(U) \times \mathcal{O}_X(V)$。
由引理 \ref{schemes-lemma-morphism-into-affine}，存在典范局部环层空间态射
$X \to \Spec(R)$。由《代数》的引理 \ref{algebra-lemma-spec-product}，
作为拓扑空间有
$\Spec(\mathcal{O}_X(U)) \amalg \Spec(\mathcal{O}_X(V)) =
\Spec(R)$
，其中的映射来自环同态 $R \to \mathcal{O}_X(U)$ 和
$R \to \mathcal{O}_X(V)$。这当然也意味着 $\Spec(R)$ 是局部环层空间范畴中的余积。
按假设，态射 $X \to \Spec(R)$ 诱导 $\Spec(\mathcal{O}_X(U))$ 与 $U$ 的同构，
对 $V$ 也同样如此。因此 $X \to \Spec(R)$ 是同构。
\end{proof}















\section{仿射概形上的拟凝聚层}
\label{schemes-section-quasi-coherent-affine}

\noindent
回顾，我们已在《模》的定义 \ref{modules-definition-quasi-coherent}
中定义了拟凝聚层这一抽象概念。本节将证明，仿射概形 $\Spec(R)$
上的任意拟凝聚层都对应于层 $\widetilde M$；这里它与某个 $R$-模 $M$
相关联。

\begin{lemma}
\label{schemes-lemma-compare-constructions}
设 $(X, \mathcal{O}_X) = (\Spec(R), \mathcal{O}_{\Spec(R)})$
为仿射概形。设 $M$ 为 $R$-模。层
$\widetilde M$（定义 \ref{schemes-definition-structure-sheaf}）由这个 $R$-模
$M$ 所关联；它与层 $\mathcal{F}_M$ 之间存在典范同构，后者由 $R$-模
$M$ 所关联
（《模》的定义 \ref{modules-definition-sheaf-associated}）。
该同构关于 $M$ 具有函子性。特别地，各层 $\widetilde M$ 都是拟凝聚的。
此外，它们由如下映射性质刻画：
$$
\Hom_{\mathcal{O}_X}(\widetilde M, \mathcal{F})
=
\Hom_R(M, \Gamma(X, \mathcal{F}))
$$
其中可取任意一个 $\mathcal{O}_X$-模层 $\mathcal{F}$。这里，映射
$\alpha : \widetilde M \to \mathcal{F}$ 对应于它在整体截面上诱导的映射。
\end{lemma}

\begin{proof}
由《模》的引理 \ref{modules-lemma-construct-quasi-coherent-sheaves}，
存在一个态射 $\mathcal{F}_M \to \widetilde M$，它对应于映射
$M \to \Gamma(X, \widetilde M) = M$。设 $x \in X$ 对应于素理想
$\mathfrak p \subset R$。在茎上诱导的映射是
$\mathcal{O}_{X, x} \otimes_R M \to M_{\mathfrak p}$；它们都是同构，
因为 $R_{\mathfrak p} \otimes_R M = M_{\mathfrak p}$。
所以映射 $\mathcal{F}_M \to \widetilde M$ 是同构。所述映射性质则来自
各层 $\mathcal{F}_M$ 的映射性质。
\end{proof}

\begin{lemma}
\label{schemes-lemma-widetilde-constructions}
设 $(X, \mathcal{O}_X) = (\Spec(R), \mathcal{O}_{\Spec(R)})$
为仿射概形。存在如下典范同构：
\begin{enumerate}
\item
$
\widetilde{M \otimes_R N}
\cong
\widetilde M \otimes_{\mathcal{O}_X} \widetilde N
$,
参见《模》第 \ref{modules-section-tensor-product} 节。
\item
$
\widetilde{\text{T}^n(M)}
\cong
\text{T}^n(\widetilde M)
$,
$
\widetilde{\text{Sym}^n(M)}
\cong
\text{Sym}^n(\widetilde M)
$，以及
$
\widetilde{\wedge^n(M)}
\cong
\wedge^n(\widetilde M)
$，
参见《模》第 \ref{modules-section-symmetric-exterior} 节。
\item 若 $M$ 是有限表示的 $R$-模，则
$
\SheafHom_{\mathcal{O}_X}(\widetilde M, \widetilde N)
\cong
\widetilde{\Hom_R(M,  N)}
$，
参见《模》第 \ref{modules-section-internal-hom} 节。
\end{enumerate}
\end{lemma}

\begin{proof}[第一种证明]
由引理 \ref{schemes-lemma-compare-constructions} 和《模》的引理
\ref{modules-lemma-construct-quasi-coherent-sheaves} 可知，函子
$M \mapsto \widetilde M$ 可以看成 $\pi^*$，其中 $\pi$ 是某个环层空间态射。
再由《模》的引理 \ref{modules-lemma-tensor-product-pullback} 和
\ref{modules-lemma-pullback-tensor-algebra}，模的拉回与张量构造相交换。
例如，态射 $\pi : (X, \mathcal{O}_X) \to (\{*\}, R)$ 是平坦的，
因为 $\mathcal{O}_X$ 的各个茎都是 $R$ 的局部化（引理
\ref{schemes-lemma-spec-sheaves}），因而在 $R$ 上平坦。因此，当第一个模是
有限表示的时，由《模》的引理 \ref{modules-lemma-pullback-internal-hom}，
沿 $\pi$ 的拉回与内部 Hom 相交换。
\end{proof}

\begin{proof}[第二种证明]
(1) 的证明。由引理 \ref{schemes-lemma-compare-constructions}，要给出从
$\widetilde{M \otimes_R N}$ 到
$\widetilde M \otimes_{\mathcal{O}_X} \widetilde N$ 的映射，
只须给出整体截面上的映射
$M \otimes_R N \to
\Gamma(X, \widetilde M \otimes_{\mathcal{O}_X} \widetilde N)$
它由模层张量积的定义给出。要证明该映射是同构，只须检验它在各个茎上
都是同构。这来自 $\widetilde{M}$ 的茎的描述（见引理
\ref{schemes-lemma-spec-sheaves}，或《模》的引理
\ref{modules-lemma-construct-quasi-coherent-sheaves}）、张量积与局部化
相交换这一事实（《代数》的引理
\ref{algebra-lemma-tensor-product-localization}），以及《模》的引理
\ref{modules-lemma-stalk-tensor-product}。

\medskip\noindent
(2) 的证明。这与 (1) 的证明类似，只须使用《代数》的引理
\ref{algebra-lemma-tensor-algebra-localization} 和《模》的引理
\ref{modules-lemma-stalk-tensor-algebra}。

\medskip\noindent
(3) 的证明。由于构造 $M \mapsto \widetilde{M}$ 具有函子性，存在一个
$R$-线性映射
$\Hom_R(M, N) \to \Hom_{\mathcal{O}_X}(\widetilde{M}, \widetilde{N})$.
该映射的目标是
$\SheafHom_{\mathcal{O}_X}(\widetilde M, \widetilde N)$ 的整体截面。
所以由引理 \ref{schemes-lemma-compare-constructions}，得到一个
$\mathcal{O}_X$-模映射 $\widetilde{\Hom_R(M,  N)} \to
\SheafHom_{\mathcal{O}_X}(\widetilde M, \widetilde N)$。
比较各茎即可检验它是同构。若 $M$ 作为 $R$-模是有限表示的，则
$\widetilde M$ 作为 $\mathcal{O}_X$-模具有整体有限表示。因此结论来自
《代数》的引理 \ref{algebra-lemma-hom-from-finitely-presented} 和《模》的
引理 \ref{modules-lemma-stalk-internal-hom}。
\end{proof}

\begin{proof}[(1) 的第三种证明]
对任意 $\mathcal{O}_X$-模 $\mathcal{F}$，有如下关于 $M$、$N$ 和
$\mathcal{F}$ 具有函子性的同构：
\begin{align*}
\Hom_{\mathcal{O}_X}(\widetilde{M} \otimes _{\mathcal{O} _X} \widetilde{N},
\mathcal{F})
& =
\Hom_{\mathcal{O}_X}(\widetilde{M},
\SheafHom_{\mathcal{O} _X} (\widetilde{N}, \mathcal{F})) \\
& =
\Hom_R(M, \Gamma(X,
\SheafHom_{\mathcal{O}_X}(\widetilde{N}, \mathcal{F}))) \\
& =
\Hom_R(M, \Hom_{\mathcal{O}_X}(\widetilde{N}, \mathcal{F})) \\
& =
\Hom_R(M, \Hom_R(N, \Gamma(X,\mathcal{F}))) \\
& =
\Hom_R(M \otimes_R N, \Gamma(X, \mathcal{F})) \\
& =
\Hom_{\mathcal{O}_X}(\widetilde{M \otimes_R N}, \mathcal{F})
\end{align*}
第一个等式是《模》的引理 \ref{modules-lemma-internal-hom}。
第二个等式是 $\widetilde{M}$ 的泛性质，参见引理
\ref{schemes-lemma-compare-constructions}。第三个等式由 $\SheafHom$ 的定义成立。
第四个等式是 $\widetilde{N}$ 的泛性质。第五个等式是《代数》的引理
\ref{algebra-lemma-hom-from-tensor-product}。最后一个等式是
$\widetilde{M \otimes_R N}$ 的泛性质。由米田引理（《范畴》的引理
\ref{categories-lemma-yoneda}）即得 (1)。
\end{proof}

\begin{lemma}
\label{schemes-lemma-widetilde-pullback}
设
$(X, \mathcal{O}_X) = (\Spec(S), \mathcal{O}_{\Spec(S)})$,
$(Y, \mathcal{O}_Y) = (\Spec(R), \mathcal{O}_{\Spec(R)})$
为仿射概形。设 $\psi : (X, \mathcal{O}_X) \to (Y, \mathcal{O}_Y)$
为仿射概形的态射，对应于环同态 $\psi^\sharp : R \to S$
（见引理 \ref{schemes-lemma-category-affine-schemes}）。
\begin{enumerate}
\item 有 $\psi^* \widetilde M = \widetilde{S \otimes_R M}$，且该等式关于
$R$-模 $M$ 具有函子性。
\item 有 $\psi_* \widetilde N = \widetilde{N_R}$，且该等式关于
$S$-模 $N$ 具有函子性。
\end{enumerate}
\end{lemma}

\begin{proof}
第一个断言来自引理 \ref{schemes-lemma-compare-constructions} 中的认同以及《模》的
引理 \ref{modules-lemma-restrict-quasi-coherent}。第二个断言来自等式
$\psi^{-1}(D(f)) = D(\psi^\sharp(f))$，因而
$$
\psi_* \widetilde N(D(f)) = \widetilde N(D(\psi^\sharp(f))) =
N_{\psi^\sharp(f)} = (N_R)_f = \widetilde{N_R}(D(f))
$$
正是所需的等式。
\end{proof}

\noindent
上面的引理 \ref{schemes-lemma-widetilde-pullback} 特别说明：把层
$\widetilde M$ 限制到标准仿射开子空间 $D(f)$ 上，所得就是
$\widetilde{M_f}$。下文将直接使用这一事实，不再另行说明。

\begin{lemma}
\label{schemes-lemma-quasi-coherent-affine}
设 $(X, \mathcal{O}_X) = (\Spec(R), \mathcal{O}_{\Spec(R)})$
为仿射概形，$\mathcal{F}$ 为拟凝聚 $\mathcal{O}_X$-模。那么
$\mathcal{F}$ 同构于由 $R$-模 $\Gamma(X, \mathcal{F})$ 所关联的层。
\end{lemma}

\begin{proof}
设 $\mathcal{F}$ 为拟凝聚 $\mathcal{O}_X$-模。由于每个标准开集
$D(f)$ 都是拟紧的，可见 $X$ 是局部拟紧的；也就是说，每一点都有一个
由拟紧邻域组成的基本系，参见《拓扑》的定义
\ref{topology-definition-locally-quasi-compact}。因此，由《模》的引理
\ref{modules-lemma-quasi-coherent-module}，对每个素理想
$\mathfrak p \subset R$ 及其对应的点 $x \in X$，都存在开邻域
$x \in U \subset X$，
使得 $\mathcal{F}|_U$ 同构于由某个 $\mathcal{O}_X(U)$-模 $M$ 所关联的
拟凝聚层。换言之，我们得到一个由具有这一性质的 $U$ 组成的开覆盖。
例如，由引理 \ref{schemes-lemma-standard-open}，可把这个覆盖加细为标准开覆盖。
于是得到覆盖 $\Spec(R) = \bigcup D(f_i)$、$R_{f_i}$-模 $M_i$，以及同构
$\varphi_i : \mathcal{F}|_{D(f_i)} \to \mathcal{F}_{M_i}$，其中再次把
相应的 $R_{f_i}$-模记为 $M_i$。在交叠处得到同构
$$
\xymatrix{
\mathcal{F}_{M_i}|_{D(f_if_j)}
\ar[rr]^{\varphi_i^{-1}|_{D(f_if_j)}}
& &
\mathcal{F}|_{D(f_if_j)}
\ar[rr]^{\varphi_j|_{D(f_if_j)}}
& &
\mathcal{F}_{M_j}|_{D(f_if_j)}.
}
$$
把这些同构记作 $\psi_{ij}$。显然，在三重交叠处有余循环条件
$$
\psi_{jk}|_{D(f_if_jf_k)}
\circ
\psi_{ij}|_{D(f_if_jf_k)}
=
\psi_{ik}|_{D(f_if_jf_k)}
$$
。

\medskip\noindent
回顾，每个开子空间 $D(f_i)$、$D(f_if_j)$、$D(f_if_jf_k)$ 都是仿射概形。
因此，各层 $\mathcal{F}_{M_i}$ 都同构于层 $\widetilde M_i$；这是上面的
引理 \ref{schemes-lemma-compare-constructions} 的结论。特别地，
$\mathcal{F}_{M_i}(D(f_if_j)) = (M_i)_{f_j}$，等等。再由上面的引理
\ref{schemes-lemma-compare-constructions} 可知，$\psi_{ij}$ 对应于唯一的
$R_{f_if_j}$-模同构
$$
\psi_{ij} : (M_i)_{f_j} \longrightarrow (M_j)_{f_i}
$$
即 $\psi_{ij}$ 在 $D(f_if_j)$ 上的截面所诱导的映射。此外，这些同构满足
如下余循环条件：
$$
\xymatrix{
(M_i)_{f_jf_k}
\ar[rd]_{\psi_{ij}}
\ar[rr]^{\psi_{ik}}
& &
(M_k)_{f_if_j} \\
&
(M_j)_{f_if_k} \ar[ru]_{\psi_{jk}}
}
$$
对任意三元组 $i, j, k$，该图交换。

\medskip\noindent
现在，《代数》的引理 \ref{algebra-lemma-glue-modules} 表明，存在一个
$R$-模 $M$，使得 $M_i = M_{f_i}$，并与各态射 $\psi_{ij}$ 相容。
考虑 $\mathcal{F}_M = \widetilde M$。此时，要证明 $\widetilde M$ 同构于
我们最初的拟凝聚层 $\mathcal{F}$，只剩一个形式性步骤。确切地说，
层 $\mathcal{F}$ 与 $\widetilde M$ 给出的 $\mathcal{O}_X$-模层粘合数据
彼此同构；这里所取覆盖为 $X = \bigcup D(f_i)$。参见《层》第
\ref{sheaves-section-glueing-sheaves} 节，尤其是引理
\ref{sheaves-lemma-mapping-property-glue}。具体到当前情形，这归结为如下论证：
我们构造一个 $R$-模映射
$$
M \longrightarrow \Gamma(X, \mathcal{F}).
$$
具体而言，给定 $m \in M$，得到 $m_i = m/1 \in M_{f_i} = M_i$；
这是 $M$ 的构造所致。由 $M_i$ 的构造，这对应于一个截面
$s_i \in \mathcal{F}(U_i)$（即 $\varphi^{-1}_i(m_i)$）。我们断言
$s_i|_{D(f_if_j)} = s_j|_{D(f_if_j)}$。这成立是因为，由 $M$ 的构造，
有 $\psi_{ij}(m_i) = m_j$，并且各 $\psi_{ij}$ 正是按上述方式构造的。
由 $\mathcal{F}$ 的层条件，这组截面唯一地给出一个截面 $s$；它是
$\mathcal{F}$ 在 $X$ 上的截面。留给读者验证 $m \mapsto s$ 是
$R$-模映射。由引理
\ref{schemes-lemma-compare-constructions}，得到相应的 $\mathcal{O}_X$-模映射
$$
\widetilde M \longrightarrow \mathcal{F}.
$$
按构造，该映射化为同构 $\varphi_i^{-1}$；在每个 $D(f_i)$ 上皆如此，
因而该映射本身是同构。
\end{proof}

\begin{lemma}
\label{schemes-lemma-equivalence-quasi-coherent}
设 $(X, \mathcal{O}_X) = (\Spec(R), \mathcal{O}_{\Spec(R)})$
为仿射概形。函子 $M \mapsto \widetilde M$ 与
$\mathcal{F} \mapsto \Gamma(X, \mathcal{F})$ 给出如下范畴之间互为拟逆的等价：
$$
\xymatrix{
\QCoh(\mathcal{O}_X)
\ar@<1ex>[r]
&
\text{Mod}_R
\ar@<1ex>[l]
}
$$
即拟凝聚 $\mathcal{O}_X$-模范畴与 $R$-模范畴之间的等价。
\end{lemma}

\begin{proof}
见上面的引理 \ref{schemes-lemma-compare-constructions} 和
\ref{schemes-lemma-quasi-coherent-affine}。
\end{proof}

\noindent
从现在起，我们不再区分仿射概形上的拟凝聚层与形如 $\widetilde M$ 的层。

\begin{lemma}
\label{schemes-lemma-kernel-cokernel-quasi-coherent}
设 $X = \Spec(R)$ 为仿射概形。拟凝聚 $\mathcal{O}_X$-模之间映射的核与余核
都是拟凝聚的。
\end{lemma}

\begin{proof}
这来自函子 $\widetilde{\ }$ 的正合性，因为由引理
\ref{schemes-lemma-compare-constructions} 可知，任意映射
$\psi : \widetilde{M} \to \widetilde{N}$ 都来自一个 $R$-模映射
$\varphi : M \to N$。（所以有
$\Ker(\psi) = \widetilde{\Ker(\varphi)}$ 和
$\Coker(\psi) = \widetilde{\Coker(\varphi)}$。）
\end{proof}

\begin{lemma}
\label{schemes-lemma-colimit-quasi-coherent}
设 $X = \Spec(R)$ 为仿射概形。$X$ 上任意一族拟凝聚层的直和仍是拟凝聚的；
余极限也同样如此。
\end{lemma}

\begin{proof}
设 $\mathcal{F}_i$（$i \in I$）是 $X$ 上的一族拟凝聚层。由上面的引理
\ref{schemes-lemma-equivalence-quasi-coherent}，可写成
$\mathcal{F}_i = \widetilde{M_i}$，其中它来自某个 $R$-模 $M_i$。
令 $M = \bigoplus M_i$，并考虑层
$\widetilde{M}$。对每个标准开集 $D(f)$，有
$$
\widetilde{M}(D(f)) = M_f =
\left(\bigoplus M_i\right)_f =
\bigoplus M_{i, f}.
$$
因此可见，拟凝聚 $\mathcal{O}_X$-模 $\widetilde{M}$ 正是各层
$\mathcal{F}_i$ 的直和。一般余极限的情形可用类似论证处理。
\end{proof}

\begin{lemma}
\label{schemes-lemma-extension-quasi-coherent}
设 $(X, \mathcal{O}_X) = (\Spec(R), \mathcal{O}_{\Spec(R)})$
为仿射概形。设
$$
0 \to
\mathcal{F}_1 \to
\mathcal{F}_2 \to
\mathcal{F}_3 \to
0
$$
是 $\mathcal{O}_X$-模层的短正合列。若三项中有两项是拟凝聚的，
则第三项也是拟凝聚的。
\end{lemma}

\begin{proof}
若 $\mathcal{F}_1$ 和 $\mathcal{F}_2$ 都是拟凝聚的，则断言显然成立，
因为函子 $M \mapsto \widetilde M$ 是正合的；见引理
\ref{schemes-lemma-spec-sheaves}。$\mathcal{F}_2$ 和 $\mathcal{F}_3$ 都拟凝聚的情形
也类似。现在设 $\mathcal{F}_1 = \widetilde M_1$ 和
$\mathcal{F}_3 = \widetilde M_3$ 都是拟凝聚的。令
$M_2 = \Gamma(X, \mathcal{F}_2)$。我们断言，只须证明序列
$$
0 \to M_1 \to M_2 \to M_3 \to 0
$$
正合。事实上，若它正合，则利用引理 \ref{schemes-lemma-compare-constructions}
的映射性质可得到交换图
$$
\xymatrix{
0 \ar[r] &
\widetilde M_1 \ar[r] \ar[d] &
\widetilde M_2 \ar[r] \ar[d] &
\widetilde M_3 \ar[r] \ar[d] &
0 \\
0 \ar[r] &
\mathcal{F}_1 \ar[r] &
\mathcal{F}_2 \ar[r] &
\mathcal{F}_3 \ar[r] &
0
}
$$
于是由蛇引理即得结论。

\medskip\noindent
这里的“标准”论证应当先证明 $H^1(X, \mathcal{F}) = 0$，其中
$\mathcal{F}$ 是任意拟凝聚层。这其实并不十分困难，但最好留到后面再做。
这里我们改用一个小技巧。

\medskip\noindent
取 $m \in M_3 = \Gamma(X, \mathcal{F}_3)$。考虑如下集合：
$$
I = \{ f \in R \mid \text{元素 }fm\text{ 来自 }M_2\}.
$$
显然这是一个理想。只须证明 $1 \in I$。因此，只须证明对任意素理想
$\mathfrak p$，都存在 $f \in I$，且 $f \not\in \mathfrak p$。
设 $x \in X$ 为对应于 $\mathfrak p$ 的点。由于满射性可以在茎上检验，
存在一个开邻域 $U$，它是 $x$ 的邻域，并使得 $m|_U$ 来自某个局部截面
$s \in \mathcal{F}_2(U)$。事实上，可以假定 $U = D(f)$ 是标准开集，
即 $f \in R$ 且 $f \not \in \mathfrak p$。我们将证明，对某个
$N \gg 0$ 有 $f^N \in I$，这便完成证明。

\medskip\noindent
任取一点 $z \in V(f)$，设它对应于素理想 $\mathfrak q \subset R$。
还可以找到 $g \in R$，满足 $g \not \in \mathfrak q$，使得
$m|_{D(g)}$ 可提升为某个 $s' \in \mathcal{F}_2(D(g))$。考虑差
$s|_{D(fg)} - s'|_{D(fg)}$。这是一个元素 $m'$，它属于
$\mathcal{F}_1(D(fg)) = (M_1)_{fg}$。对某个整数
$n = n(z)$，元素 $f^n m'$ 来自某个 $m'_1 \in (M_1)_g$。可见
$f^n s$ 延拓为一个截面 $\sigma$，它是 $\mathcal{F}_2$ 在
$D(f) \cup D(g)$ 上的截面；这是因为它与 $f^n s' + m'_1$ 的限制
在 $D(f) \cap D(g) = D(fg)$ 上相同。此外，$\sigma$ 映到 $f^n m$ 在
$D(f) \cup D(g)$ 上的限制。

\medskip\noindent
由于 $V(f)$ 是拟紧的，存在有限多个元素 $g_1, \ldots, g_m \in R$，
使得 $V(f) \subset \bigcup D(g_j)$；还存在一个整数 $n > 0$ 以及截面
$\sigma_j \in \mathcal{F}_2(D(f) \cup D(g_j))$，使得
$\sigma_j|_{D(f)} = f^n s$，并且 $\sigma_j$ 映到截面
$f^nm|_{D(f) \cup D(g_j)}$；后者是 $\mathcal{F}_3$ 的截面。考虑这些差：
$$
\sigma_j|_{D(f) \cup D(g_jg_k)}
-
\sigma_k|_{D(f) \cup D(g_jg_k)}.
$$
它们对应于 $\mathcal{F}_1$ 在 $D(f) \cup D(g_jg_k)$ 上的截面，
并且在 $D(f)$ 上为零。特别地，它们在
$\mathcal{F}_1(D(g_jg_k)) = (M_1)_{g_jg_k}$ 中的像，在
$(M_1)_{g_jg_kf}$ 中为零。因此，$f$ 的某个高次幂把每一个这样的像都
零化。换言之，各元素 $f^N \sigma_j$ 对某个 $N \gg 0$ 满足层性质的
粘合条件，并给出一个截面 $\sigma $；它属于 $\mathcal{F}_2$，定义在
$\bigcup (D(f) \cup D(g_j)) = X$ 上，正如所需。
\end{proof}







\section{仿射概形的闭子空间}
\label{schemes-section-closed-in-affine}

\begin{example}
\label{schemes-example-closed-immersion-affines}
设 $R$ 为环，$I \subset R$ 为理想。考虑仿射概形的态射
$i : Z = \Spec(R/I) \to \Spec(R) = X$。由《代数》的引理
\ref{algebra-lemma-spec-closed}，它把 $Z$ 同胚地映到 $X$ 的一个闭子集上。
此外，若 $I \subset \mathfrak p \subset R$ 是对应于点 $x = i(z)$ 的
素理想，其中 $x \in X$、$z \in Z$，则在茎上得到映射
$$
\mathcal{O}_{X, x} = R_{\mathfrak p}
\longrightarrow
R_{\mathfrak p}/IR_{\mathfrak p} = \mathcal{O}_{Z, z}
$$
由此可见，$i$ 是局部环层空间的闭浸入；见定义
\ref{schemes-definition-closed-immersion-locally-ringed-spaces}。显然，它同构于由
拟凝聚理想层 $\widetilde I$ 所关联的闭子空间，如例
\ref{schemes-example-closed-subspace} 所述。
\end{example}

\begin{lemma}
\label{schemes-lemma-closed-immersion-affine-case}
\begin{slogan}
对仿射概形，闭浸入对应于理想。
\end{slogan}
设 $(X, \mathcal{O}_X) = (\Spec(R), \mathcal{O}_{\Spec(R)})$
为仿射概形，$i : Z \to X$ 为局部环层空间的任意闭浸入。则存在唯一理想
$I \subset R$，使得态射 $i : Z \to X$ 可与闭浸入
$\Spec(R/I) \to \Spec(R)$ 认同；后一个闭浸入是在上面的例
\ref{schemes-example-closed-immersion-affines} 中构造的。
\end{lemma}

\begin{proof}
这几乎是显然的！事实上，由引理 \ref{schemes-lemma-closed-immersion}，可以把
$Z \to X$ 与由理想层 $\mathcal{I} \subset \mathcal{O}_X$ 所关联的闭子空间
认同，如定义 \ref{schemes-definition-closed-subspace} 和例
\ref{schemes-example-closed-subspace} 所述。依照我们的约定，该理想层作为
$\mathcal{O}_X$-模层局部地由截面生成。因此，商层
$\mathcal{O}_X / \mathcal{I}$ 在 $X$ 上局部地是某个映射
$\bigoplus_{j \in J} \mathcal{O}_U \to \mathcal{O}_U$ 的余核。
所以按定义，$\mathcal{O}_X / \mathcal{I}$ 是拟凝聚的。由第
\ref{schemes-section-quasi-coherent-affine} 节的结果，它形如 $\widetilde S$，
是由某个 $R$-模 $S$ 得到的。此外，由于
$\mathcal{O}_X = \widetilde R \to \widetilde S$ 是满射，由引理
\ref{schemes-lemma-extension-quasi-coherent} 可见，$\mathcal{I}$ 也是拟凝聚的，
记作 $\mathcal{I} = \widetilde I$。当然，$I \subset R$ 且 $S = R/I$，
一切都已清楚。
\end{proof}













\section{概形}
\label{schemes-section-schemes}

\begin{definition}
\label{schemes-definition-scheme}
\begin{history}
在 \cite{EGA1} 中，我们现在所谓的概形称为“pre-sch\'ema”，而“sch\'ema”
一词留给 Stacks Project 中所谓的分离概形。在第二版
\cite{EGA1-second} 中，术语改成了如今的标准用法。不过，偶尔仍会见到
“prescheme”这一术语，例如 \cite{Murre-lectures}。
\end{history}
{\it 概形}是满足如下性质的局部环层空间：每一点都有一个本身为仿射概形的
开邻域。{\it 概形态射}是局部环层空间的态射。概形范畴记作 $\Sch$。
\end{definition}

\noindent
设 $X$ 为概形。我们将作如下（极轻微的）语言滥用：若开子空间
$U \subset X$ 中的 $U$ 是仿射概形，就称它为{\it 仿射开集}或
{\it 开仿射集}。
我们常写 $U = \Spec(R)$，以表示 $U$ 同构于 $\Spec(R)$，并且还表示
我们将暂时认同 $U$ 与 $\Spec(R)$。

\begin{lemma}
\label{schemes-lemma-open-subspace-scheme}
设 $X$ 为概形，$j : U \to X$ 为局部环层空间的开浸入。则 $U$ 是概形。
特别地，$X$ 的任意开子空间都是概形。
\end{lemma}

\begin{proof}
设 $U \subset X$，并取 $u \in U$。选取仿射开邻域
$u \in V \subset X$。由于 $V$ 的标准开集构成 $V$ 上拓扑的一组基，
存在 $f\in \mathcal{O}_V(V)$，使得 $u \in D(f) \subset U$。而由引理
\ref{schemes-lemma-standard-open-affine}，$D(f)$ 是仿射概形。这就证明了 $U$ 的
每一点都有一个仿射开邻域。
\end{proof}

\noindent
显然，该引理（或其证明）表明，任意概形 $X$ 的拓扑都有一组由仿射开集
构成的基（见《拓扑》第 \ref{topology-section-bases} 节）。

\begin{example}
\label{schemes-example-not-affine}
设 $k$ 为域。开子空间
$U = \Spec(k[x, y]) \setminus \{ (x, y)\}$ 位于仿射概形
$X =\Spec(k[x, y])$ 中，并给出了一个非仿射概形的例子。
它由以下两个仿射开集覆盖：
\begin{center}
$D(x) = \Spec(k[x, y, 1/x])$，\qquad
$D(y) = \Spec(k[x, y, 1/y])$。
\end{center}
二者之交为 $D(xy) = \Spec(k[x, y, 1/xy])$。由 $\mathcal{O}_U$ 的层性质，存在正合列
$$
0 \to
\Gamma(U, \mathcal{O}_U) \to
k[x, y, 1/x] \times k[x, y, 1/y] \to
k[x, y, 1/xy]
$$
由此可知，映射 $k[x, y] \to \Gamma(U, \mathcal{O}_U)$（它来自态射
$U \to X$）是同构。因此 $U$ 不可能是仿射的；否则由引理
\ref{schemes-lemma-category-affine-schemes} 将有 $U \cong X$。
\end{example}











\section{概形的浸入}
\label{schemes-section-immersions}

\noindent
在引理 \ref{schemes-lemma-open-subspace-scheme} 中，我们看到概形的任意开子空间
仍是概形。下面将证明，概形的闭子空间也有同样的性质。

\medskip\noindent
注意，拟凝聚 $\mathcal{O}_X$-模层的概念对任意环层空间 $X$ 都有定义，
特别是在 $X$ 为概形时。由第 \ref{schemes-section-quasi-coherent-affine} 节的结果
可知，这样的层在任意仿射开集 $U \subset X$ 上都形如 $\widetilde M$，
其中它来自某个 $\mathcal{O}_X(U)$-模 $M$。

\begin{lemma}
\label{schemes-lemma-closed-subspace-scheme}
设 $X$ 为概形，$i : Z \to X$ 为局部环层空间的闭浸入。
\begin{enumerate}
\item 局部环层空间 $Z$ 是概形；
\item 核 $\mathcal{I}$ 是拟凝聚理想层；这里所取映射为
$\mathcal{O}_X \to i_*\mathcal{O}_Z$；
\item 对任意仿射开集 $U = \Spec(R)$，若它属于 $X$，则态射
$i^{-1}(U) \to U$ 可与 $\Spec(R/I) \to \Spec(R)$ 认同，其中
$I \subset R$ 是某个理想；
\item 有 $\mathcal{I}|_U = \widetilde I$。
\end{enumerate}
特别地，任意局部由截面生成的理想层都是拟凝聚理想层（反之亦然），
而且 $X$ 的任意闭子空间都是概形。
\end{lemma}

\begin{proof}
设 $i : Z \to X$ 为闭浸入，并取一点 $z \in Z$。任取一个仿射开邻域
$i(z) \in U \subset X$，记 $U = \Spec(R)$。由引理
\ref{schemes-lemma-closed-immersion-affine-case} 可知，$i^{-1}(U) \to U$ 可与
仿射概形的态射 $\Spec(R/I) \to \Spec(R)$ 认同。首先，这说明
$z \in i^{-1}(U) \subset Z$ 是 $z$ 的仿射邻域，因而 $Z$ 是概形。
其次，这说明 $\mathcal{I}|_U$ 就是 $\widetilde I$。换言之，对每一点
$x \in i(Z)$，都存在一个开邻域，使得 $\mathcal{I}$ 在该邻域中拟凝聚。
注意，$\mathcal{I}|_{X \setminus i(Z)}
\cong \mathcal{O}_{X \setminus i(Z)}$。因此，该理想层在
$X \setminus i(Z)$ 上的限制也是拟凝聚的。由此可知，$\mathcal{I}$ 拟凝聚。
\end{proof}

\begin{definition}
\label{schemes-definition-immersion}
设 $X$ 为概形。
\begin{enumerate}
\item 若一个概形态射作为局部环层空间态射是开浸入，就称它为
{\it 开浸入}（见定义 \ref{schemes-definition-immersion-locally-ringed-spaces}）。
\item $X$ 的{\it 开子概形}是 $X$ 的开子空间，其含义见定义
\ref{schemes-definition-open-subspace}；$X$ 的开子概形由引理
\ref{schemes-lemma-open-subspace-scheme} 可知是概形。
\item 若一个概形态射作为局部环层空间态射是闭浸入，就称它为
{\it 闭浸入}（见定义
\ref{schemes-definition-closed-immersion-locally-ringed-spaces}）。
\item $X$ 的{\it 闭子概形}是 $X$ 的闭子空间，其含义见定义
\ref{schemes-definition-closed-subspace}；由引理 \ref{schemes-lemma-closed-subspace-scheme}，
闭子概形是概形。
\item 若概形态射 $f : X \to Y$ 可分解为 $j \circ i$，其中 $i$ 是闭浸入、
$j$ 是开浸入，就称它为{\it 浸入}或{\it 局部闭浸入}。
\end{enumerate}
\end{definition}

\noindent
由第 \ref{schemes-section-open-immersion} 节和第
\ref{schemes-section-closed-immersion} 节中的引理可知，概形的任意开浸入
（相应地，闭浸入）都同构于目标中某个开子概形（相应地，闭子概形）的
包含映射。

\medskip\noindent
我们对闭浸入的定义介于 Hartshorne 与 EGA 的定义之间。Hartshorne 把
闭浸入定义为概形态射 $f : X \to Y$：它把 $X$ 同胚地映到 $Y$ 的一个
闭子集上，并且 $f^\# : \mathcal{O}_Y \to f_*\mathcal{O}_X$ 是满射；见
\cite[Page 85]{H}。我们将在引理
\ref{schemes-lemma-characterize-closed-immersions} 中证明，这与我们的概念等价。
在 \cite{EGA} 中，Grothendieck 和 Dieudonn\'e 首先利用拟凝聚理想层，
通过例 \ref{schemes-example-closed-subspace} 的构造定义闭子概形；然后把闭浸入
定义为一个态射 $f : X \to Y$，它诱导出与某个闭子概形的同构。由引理
\ref{schemes-lemma-closed-subspace-scheme} 可知，这与我们的概念一致。

\medskip\noindent
从教学角度说，上面的定义简直是一场灾难。给学生讲授这些材料时，我们发现
常常更方便的做法是把闭浸入定义为概形的仿射态射 $f : X \to Y$，并要求
$f^\# : \mathcal{O}_Y \to f_*\mathcal{O}_X$ 是满射。事实上，仿射态射这一
概念（《态射》第 \ref{morphisms-section-affine} 节）十分自然且容易理解。

\medskip\noindent
关于闭浸入的更多信息，建议读者参阅《态射》第
\ref{morphisms-section-closed-immersions} 节和第
\ref{morphisms-section-closed-immersions-quasi-coherent} 节。

\medskip\noindent
本节末尾将讨论局部闭子概形与浸入。

\begin{remark}
\label{schemes-remark-not-reverse-open-closed}
若 $f : X \to Y$ 是概形的浸入，则一般不能把 $f$ 分解成先开浸入、
后闭浸入的复合。见《态射》的例 \ref{morphisms-example-thibaut}。
\end{remark}

\begin{lemma}
\label{schemes-lemma-immersion-when-closed}
设 $f : Y \to X$ 为概形的浸入。则 $f$ 是闭浸入，当且仅当
$f(Y) \subset X$ 是闭子集。
\end{lemma}

\begin{proof}
若 $f$ 是闭浸入，则按定义 $f(Y)$ 为闭集。反过来，设 $f(Y)$ 为闭集。
按定义，存在开子概形 $U \subset X$，使得 $f$ 是闭浸入
$i : Y \to U$ 与开浸入 $j : U \to X$ 的复合。设
$\mathcal{I} \subset \mathcal{O}_U$ 为与闭浸入 $i$ 关联的拟凝聚理想层。
注意
$\mathcal{I}|_{U \setminus i(Y)}
= \mathcal{O}_{U \setminus i(Y)}
= \mathcal{O}_{X \setminus i(Y)}|_{U \setminus i(Y)}$.
所以可以把它们粘合（见《层》第 \ref{sheaves-section-glueing-sheaves} 节）：
将 $\mathcal{I}$ 与 $\mathcal{O}_{X \setminus i(Y)}$ 粘合成理想层
$\mathcal{J} \subset \mathcal{O}_X$。由于 $X$ 的每一点都有一个邻域，
使得 $\mathcal{J}$ 在该邻域上拟凝聚，可见 $\mathcal{J}$ 是拟凝聚的
（特别地，它局部由截面生成）。按构造，$\mathcal{O}_X/\mathcal{J}$
的支集包含在 $U$ 中，并且限制到那里就等于
$\mathcal{O}_U/\mathcal{I}$。因此，与 $\mathcal{I}$ 和 $\mathcal{J}$
关联的闭子空间典范同构；见例 \ref{schemes-example-closed-subspace}。特别地，
$U$ 中与 $\mathcal{I}$ 关联的闭子空间同构于 $X$ 的一个闭子空间。
由于 $Y \to U$ 与由 $\mathcal{I}$ 关联的闭子空间认同（见引理
\ref{schemes-lemma-closed-immersion}），所以 $Y \to U \to X$ 是闭浸入。
\end{proof}

\noindent
设 $f : Y \to X$ 为浸入。令
$Z = \overline{f(Y)} \setminus f(Y)$；这是 $X$ 的闭子集。再令
$U = X \setminus Z$。该引理说明，$U$ 是 $X$ 的最大开子空间，使得
$f : Y \to X$ 可通过一个进入 $U$ 的闭浸入分解。我们把 $X$ 的一个
{\it 局部闭子概形}定义为二元组 $(Z, U)$，其中 $Z$ 是开子概形 $U$ 的
闭子概形，而后者又是 $X$ 的开子概形，并且还满足
$\overline{Z} \cup U = X$。通常只说“设 $Z$ 为 $X$ 的局部闭子概形”，
因为可以恢复 $U$，所依据的是态射 $Z \to X$。上面的论述表明，任意浸入
$f : Y \to X$ 唯一地分解为 $Y \to Z \to X$，其中 $Z$ 是 $X$ 的
局部闭子空间，而 $Y \to Z$ 是同构。

\medskip\noindent
这一做法的意义在于，$X$ 的所有局部闭子概形构成一个集合。可以在该集合上
定义一个偏序；出于显然的理由，我们称之为包含关系。具体地，若
$Z \to X$ 和 $Z' \to X$ 是 $X$ 的两个局部闭子概形，那么我们称
{\it $Z$ 包含于 $Z'$}，当且仅当态射 $Z \to X$ 通过 $Z'$ 分解。若确实如此，
则 $Z$ 当然与 $Z'$ 的唯一一个局部闭子概形认同，依此类推。

\medskip\noindent
关于浸入的更多信息，参见《态射》第 \ref{morphisms-section-immersions} 节。










\section{概形的 Zariski 拓扑}
\label{schemes-section-topology}

\noindent
关于适用于概形 Zariski 拓扑的一些基础拓扑材料，见《拓扑》第
\ref{topology-section-introduction} 节。

\begin{lemma}
\label{schemes-lemma-scheme-sober}
设 $X$ 为概形。$X$ 的任意不可约闭子集都有唯一泛点。换言之，
$X$ 是清醒拓扑空间；见《拓扑》的定义
\ref{topology-definition-generic-point}。
\end{lemma}

\begin{proof}
设 $Z \subset X$ 为不可约闭子集。对每个仿射开集 $U \subset X$，
$U = \Spec(R)$，都有 $Z \cap U = V(I)$，其中 $I \subset R$ 是唯一根理想。
注意，$Z \cap U$ 或为空，或不可约。在后一种情形（它至少对一个 $U$
发生），可见 $I = \mathfrak p$ 是素理想；它给出一个泛点 $\xi$，
属于 $Z \cap U$。于是 $Z = \overline{\{\xi\}}$；换言之，
$\xi$ 是 $Z$ 的泛点。若 $\xi'$ 是另一个泛点，则
$\xi' \in Z \cap U$，从而立即得到 $\xi' = \xi$。
\end{proof}

\begin{lemma}
\label{schemes-lemma-basis-affine-opens}
设 $X$ 为概形。$X$ 的所有仿射开集构成 $X$ 上拓扑的一组基。
\end{lemma}

\begin{proof}
这来自第 \ref{schemes-section-schemes} 节中关于开子概形的讨论。
\end{proof}

\begin{remark}
\label{schemes-remark-intersection-affine-opens}
一般而言，$X$ 中两个仿射开集的交不再是仿射开集。见例
\ref{schemes-example-affine-space-zero-doubled}。
\end{remark}

\begin{lemma}
\label{schemes-lemma-locally-quasi-compact}
任意概形的底拓扑空间都是局部拟紧的；见《拓扑》的定义
\ref{topology-definition-locally-quasi-compact}。
\end{lemma}

\begin{proof}
这来自上面的引理 \ref{schemes-lemma-basis-affine-opens}，以及环的谱是拟紧的这一
事实；见《代数》的引理 \ref{algebra-lemma-quasi-compact}。
\end{proof}

\begin{lemma}
\label{schemes-lemma-standard-open-two-affines}
设 $X$ 为概形。设 $U, V$ 为 $X$ 的仿射开集，且 $x \in U \cap V$。
存在一个仿射开邻域 $W$，它是 $x$ 的邻域，并使得 $W$ 同时是 $U$ 和
$V$ 的标准开集。
\end{lemma}

\begin{proof}
写成 $U = \Spec(A)$ 和 $V = \Spec(B)$。设 $x$ 对应于素理想
$\mathfrak p \subset A$ 以及素理想 $\mathfrak q \subset B$。可以选取
$f \in A$，$f \not \in \mathfrak p$，使得 $D(f) \subset U \cap V$。
注意，$D(f)$ 的任意标准开集都是 $\Spec(A) = U$ 的标准开集。因此可以
假定 $U \subset V$。换言之，现在可以把 $U$ 看成 $V$ 的一个仿射开集。
接着选取 $g \in B$，$g \not \in \mathfrak q$，使得 $D(g) \subset U$。
此时可见 $D(g) = D(g_A)$，其中 $g_A \in A$ 表示 $g$ 在映射
$B \to A$ 下的像。引理得证。
\end{proof}

\begin{lemma}
\label{schemes-lemma-good-subcover}
设 $X$ 为概形。设 $X = \bigcup_i U_i$ 为仿射开覆盖，并设
$V \subset X$ 为仿射开集。存在标准开覆盖
$V = \bigcup_{j = 1, \ldots, m} V_j$（见定义
\ref{schemes-definition-standard-covering}），使得每个 $V_j$ 都是某个 $U_i$
中的标准开集。
\end{lemma}

\begin{proof}
取 $v \in V$。则有 $v \in U_i$，其中 $i$ 是某个指标。由上面的引理
\ref{schemes-lemma-standard-open-two-affines}，存在开集
$v \in W_v \subset V \cap U_i$，使得 $W_v$ 同时是 $V$ 和 $U_i$
中的标准开集。由于 $V$ 拟紧，引理得证。
\end{proof}

\begin{lemma}
\label{schemes-lemma-sheaf-on-affines}
设 $X$ 为概形。设 $\mathcal{B}$ 为 $X$ 的所有仿射开集构成的集合。
设 $\mathcal{F}$ 为 $\mathcal{B}$ 上的集合预层，见《层》的定义
\ref{sheaves-definition-presheaf-basis}。以下条件等价：
\begin{enumerate}
\item $\mathcal{F}$ 是 $X$ 上某个层限制到 $\mathcal{B}$ 的结果；
\item $\mathcal{F}$ 是 $\mathcal{B}$ 上的层；
\item $\mathcal{F}(\emptyset)$ 是单点集，并且每当
$U = V \cup W$，其中 $U, V, W \in \mathcal{B}$ 且
$V, W \subset U$ 是标准开集（《代数》的定义
\ref{algebra-definition-Zariski-topology}），映射
$$
\mathcal{F}(U) \longrightarrow \mathcal{F}(V) \times \mathcal{F}(W)
$$
是单射，其像由所有二元组 $(s, t)$ 构成，并且这些二元组满足
$s|_{V \cap W} = t|_{V \cap W}$。
\end{enumerate}
\end{lemma}

\begin{proof}
(1) 与 (2) 的等价性就是《层》的引理
\ref{sheaves-lemma-restrict-basis-equivalence}。显然，(2) 蕴含 (3)。
所以只须证明 (3) 蕴含 (2)。由《层》的引理
\ref{sheaves-lemma-cofinal-systems-coverings-standard-case} 和引理
\ref{schemes-lemma-standard-open}，只须对标准开覆盖（定义
\ref{schemes-definition-standard-covering}）证明层条件；这里覆盖的是
$\mathcal{B}$ 中的元素。设
$U = U_1 \cup \ldots \cup U_n$ 是标准开覆盖，其中 $U \subset X$ 为
仿射开集。我们对 $n$ 归纳证明该覆盖满足层条件。若 $n = 0$，则 $U$
为空，层条件由假设成立。若 $n = 1$，则无事可证。若 $n = 2$，
这正是条件 (3)。若 $n > 2$，则写 $U_i = D(f_i)$，其中
$f_i \in A = \mathcal{O}_X(U)$。设 $s_i \in \mathcal{F}(U_i)$ 是截面，
且有 $s_i|_{U_i \cap U_j} = s_j|_{U_i \cap U_j}$，这对所有
$1 \leq i < j \leq n$ 都成立。由于 $U = U_1 \cup \ldots \cup U_n$，有
$1 = \sum_{i = 1, \ldots, n} a_i f_i$；该等式在 $A$ 中成立，其中某个
$a_i \in A$。见
《代数》的引理 \ref{algebra-lemma-Zariski-topology}。令
$g = \sum_{i = 1, \ldots, n - 1} a_if_i$。于是
$U = D(g) \cup D(f_n)$。注意，
$D(g) = D(gf_1) \cup \ldots \cup D(gf_{n - 1})$ 是标准开覆盖。
由归纳假设，存在唯一截面 $s' \in \mathcal{F}(D(g))$，它与
$s_i|_{D(gfi)}$ 对 $i = 1, \ldots, n - 1$ 的每个值都相符。我们断言
$s'$ 与 $s_n$ 在 $D(gf_n)$ 上的限制相同。这来自归纳假设以及覆盖
$D(gf_n) = D(gf_nf_1) \cup \ldots \cup D(gf_nf_{n - 1})$。因此存在唯一截面
$s \in \mathcal{F}(U)$，它在 $D(g)$ 上的限制是 $s'$，在 $D(f_n)$ 上的
限制是 $s_n$。我们略去如下验证：$s$ 的限制是 $s_i$，相应开集为
$D(f_i)$，其中 $i = 1, \ldots, n - 1$；也略去 $s$ 唯一性的验证。
\end{proof}

\begin{lemma}
\label{schemes-lemma-scheme-finite-discrete-affine}
设 $X$ 为概形，其底拓扑空间是有限离散集。则 $X$ 是仿射的。
\end{lemma}

\begin{proof}
写成 $X = \{x_1, \ldots, x_n\}$。则 $U_i = \{x_i\}$ 是 $x_i$ 的开邻域。
由引理 \ref{schemes-lemma-basis-affine-opens}，它是仿射的。因此 $X$ 是有限多个
仿射概形的不交并，进而由引理 \ref{schemes-lemma-disjoint-union-affines} 可知它是
仿射的。
\end{proof}

\begin{example}
\label{schemes-example-scheme-without-closed-points}
存在没有闭点的概形。事实上，设 $R$ 为局部整环，其谱形如
$(0) = \mathfrak p_0 \subset \mathfrak p_1 \subset \mathfrak p_2
\subset \ldots \subset \mathfrak m$。那么开子概形
$\Spec(R) \setminus \{\mathfrak m\}$ 没有闭点。为说明这样的环 $R$
确实存在，使用如下事实：给定任意全序群 $(\Gamma, \geq)$，存在赋值环
$A$，其值群为 $(\Gamma, \geq)$；见 \cite{Krull}。记号见《代数》第
\ref{algebra-section-valuation-rings} 节。我们取
$\Gamma = \mathbf{Z}x_1 \oplus \mathbf{Z}x_2 \oplus \mathbf{Z}x_3 \oplus
\ldots$，并定义 $\sum_i a_i x_i \geq 0$ 当且仅当第一个非零的 $a_i$
满足 $> 0$，或所有 $a_i = 0$。于是
$x_1 \geq x_2 \geq x_3 \geq  \ldots \geq 0$。各子集
$x_i + \Gamma_{\geq 0}$ 是 $(\Gamma, \geq)$ 的素理想；参见《代数》的
引理 \ref{algebra-lemma-ideals-valuation-ring} 上方的记号。它们连同
$\emptyset$ 与 $\Gamma_{\geq 0}$ 是全部素理想。因此，$A$ 是一个环的
例子，其谱具有上述结构；这来自《代数》的引理
\ref{algebra-lemma-ideals-valuation-ring}。
\end{example}



















\section{约化概形}
\label{schemes-section-reduced}

\begin{definition}
\label{schemes-definition-reduced}
设 $X$ 为概形。称 $X$ 是{\it 约化的}，如果每个局部环
$\mathcal{O}_{X, x}$ 都是约化的。
\end{definition}

\begin{lemma}
\label{schemes-lemma-reduced}
概形 $X$ 是约化的，当且仅当 $\mathcal{O}_X(U)$ 是约化环；这里
$U \subset X$ 是任意开集。
\end{lemma}

\begin{proof}
先设 $X$ 约化。取截面 $f \in \mathcal{O}_X(U)$，并设 $f^n = 0$。
则 $f$ 在 $\mathcal{O}_{U, u}$ 中的像为零；这对所有 $u \in U$ 都成立。
所以 $f$ 为零；见《层》的引理
\ref{sheaves-lemma-sheaf-subset-stalks}。反过来，设 $\mathcal{O}_X(U)$
约化；这对所有开集 $U$ 都成立。任取非零元素
$f \in \mathcal{O}_{X, x}$。任意代表 $(U, f \in \mathcal{O}(U))$，
只要它代表 $f$，就非零，因而不是幂零元。
所以 $f$ 在 $\mathcal{O}_{X, x}$ 中不是幂零元。
\end{proof}

\begin{lemma}
\label{schemes-lemma-affine-reduced}
仿射概形 $\Spec(R)$ 约化，当且仅当 $R$ 约化。
\end{lemma}

\begin{proof}
正向蕴含立即来自上面的引理 \ref{schemes-lemma-reduced}。反向蕴含则来自如下事实：
约化环的任意局部化仍约化，特别地，约化环的各局部环都约化。
\end{proof}

\begin{lemma}
\label{schemes-lemma-reduced-closed-subscheme}
设 $X$ 为概形，$T \subset X$ 为闭子集。存在唯一闭子概形
$Z \subset X$，满足：(a) $Z$ 的底拓扑空间等于 $T$；(b) $Z$ 约化。
\end{lemma}

\begin{proof}
设 $\mathcal{I} \subset \mathcal{O}_X$ 为由下式定义的子预层：
$$
\mathcal{I}(U) = \{f \in \mathcal{O}_X(U) \mid
f(t) = 0\text{ 对所有 }t \in T\cap U\}
$$
这里用 $f(t)$ 表示 $f$ 在剩余域 $\kappa(t)$ 中的像；该域是 $X$ 于
$t$ 处的剩余域。
由于该条件具有局部性，显然 $\mathcal{I}$ 是理想层。此外，设
$U = \Spec(R)$ 为仿射开集。可以写成 $T \cap U = V(I)$，其中
$I \subset R$ 是唯一根理想。给定素理想 $\mathfrak p \in V(I)$，
它对应于 $t \in T \cap U$；再给定元素 $f \in R$，则有
$f(t) = 0 \Leftrightarrow f \in \mathfrak p$。所以，
$\mathcal{I}(U) = \bigcap_{\mathfrak p \in V(I)} \mathfrak p
= I$；这是《代数》的引理 \ref{algebra-lemma-Zariski-topology} 的结论。
此外，对任意标准开集 $D(g) \subset \Spec(R) = U$，同样的论证给出
$\mathcal{I}(D(g)) = I_g$。所以 $\widetilde I$ 与 $\mathcal{I}|_U$
在一组开集基上作为理想相同，因而二者相等。因此，$\mathcal{I}$ 是拟凝聚
理想层。

\medskip\noindent
现在可以把 $Z$ 定义为由理想层 $\mathcal{I}$ 所关联的闭子空间。
对每个仿射开集 $U = \Spec(R)$，若它属于 $X$，则有
$Z \cap U = \Spec(R/I)$，其中 $I$ 是根理想，因而 $Z$ 约化
（由上面的引理 \ref{schemes-lemma-affine-reduced}）。按构造，$Z$ 的底闭子集就是
$T$。因此已经找到满足 (a)、(b) 的闭子概形。

\medskip\noindent
设 $Z' \subset X$ 是另一个满足 (a)、(b) 的闭子概形。对每个仿射开集
$U = \Spec(R)$，若它属于 $X$，则有
$Z' \cap U = \Spec(R/I')$，其中 $I' \subset R$ 是某个理想。由引理
\ref{schemes-lemma-affine-reduced}，环 $R/I'$ 约化，因而 $I'$ 是根理想。由于
$V(I') = T \cap U = V(I)$，得到 $I = I'$；这是《代数》的引理
\ref{algebra-lemma-Zariski-topology} 的结论。所以 $Z'$ 与 $Z$
由同一个理想层定义，因而相等。
\end{proof}

\begin{definition}
\label{schemes-definition-reduced-induced-scheme}
设 $X$ 为概形，$Z \subset X$ 为闭子集。$Z$ 上的一个{\it 概形结构}
由闭子概形 $Z'$ 给出；它属于 $X$，且其底集合等于 $Z$。我们常说“设
$(Z, \mathcal{O}_Z)$ 是 $Z$ 上的一个概形结构”来表达这一点。$Z$ 上的
{\it 约化诱导概形结构}是引理 \ref{schemes-lemma-reduced-closed-subscheme}
中构造的结构。{\it 约化 $X_{red}$，其对象为 $X$}，是 $X$ 本身上的
约化诱导概形结构。
\end{definition}

\noindent
我们说“设 $Z \subset X$ 为 $X$ 的一个不可约分支”时，通常使用约化诱导
概形结构，把 $Z$ 看成 $X$ 的约化闭子概形。

\begin{remark}
\label{schemes-remark-reduced-induced-locally-closed}
设 $X$ 为概形，$T \subset X$ 为局部闭子集。在这种情形下，我们有时也使用
“$T$ 上的约化诱导概形结构”这一说法。它指如下结构：按照定义
\ref{schemes-definition-reduced-induced-scheme}，把 $T$ 看成开子概形
$X \setminus \partial T$ 中的闭子集；后者是 $X$ 的开子概形。这里
$\partial T = \overline{T} \setminus T$ 是 $T$ 在 $X$ 的拓扑空间中的“边界”。
\end{remark}

\begin{lemma}
\label{schemes-lemma-map-into-reduction}
设 $X$ 为概形，$Z \subset X$ 为闭子概形，$Y$ 为约化概形。态射
$f : Y \to X$ 通过 $Z$ 分解，当且仅当 $f(Y) \subset Z$（在集合论意义下）。
特别地，任意态射 $Y \to X$ 都分解为 $Y \to X_{red} \to X$。
\end{lemma}

\begin{proof}
设 $f(Y) \subset Z$（在集合论意义下）。设
$\mathcal{I} \subset \mathcal{O}_X$ 为 $Z$ 的理想层。对任意仿射开集
$U \subset X$、$\Spec(B) = V \subset Y$，若 $f(V) \subset U$，并任取
$g \in \mathcal{I}(U)$，则其拉回
$b = f^\sharp(g) \in \Gamma(V, \mathcal{O}_Y) = B$ 在任意 $y \in V$
的剩余域中都映为零。换言之，
$b \in \bigcap_{\mathfrak p \subset B} \mathfrak p$。这推出 $b = 0$，
因为 $B$ 约化（引理 \ref{schemes-lemma-reduced} 以及《代数》的引理
\ref{algebra-lemma-Zariski-topology}）。因此，$f$ 通过 $Z$ 分解；
这是引理 \ref{schemes-lemma-characterize-closed-subspace} 的结论。
\end{proof}










\section{概形的点}
\label{schemes-section-points}

\noindent
给定概形 $X$，可以定义函子
$$
h_X : \Sch^{opp}
\longrightarrow
\textit{Sets}, \quad
T \longmapsto \Mor(T, X).
$$
见《范畴》的例 \ref{categories-example-hom-functor}。它称为 $X$ 的
{\it 点函子}。概形理论中颇有趣的一项工作，是描述 $X$ 的内部几何；
这里所用的是函子 $h_X$。本节将给出描述 $X$ 的点的一种简单方法。

\medskip\noindent
设 $X$ 为概形，$R$ 为局部环，其极大理想为 $\mathfrak m \subset R$。
设 $f : \Spec(R) \to X$ 为概形态射。令 $x \in X$ 为闭点
$\mathfrak m \in \Spec(R)$ 的像。于是得到局部环之间的局部同态
$$
f^\sharp :
\mathcal{O}_{X, x}
\longrightarrow
\mathcal{O}_{\Spec(R), \mathfrak m} = R.
$$

\begin{lemma}
\label{schemes-lemma-morphism-from-spec-local-ring}
设 $X$ 为概形，$R$ 为局部环。上述构造在态射 $\Spec(R) \to X$ 与二元组
$(x, \varphi)$ 之间给出一一对应；后二者由点 $x \in X$ 和局部环的局部同态
$\varphi : \mathcal{O}_{X, x} \to R$ 组成。
\end{lemma}

\begin{proof}
设 $A$ 为环。对任意环同态 $\psi : A \to R$，存在唯一素理想
$\mathfrak p \subset A$ 以及分解 $A \to A_{\mathfrak p} \to R$，
使得后一个映射是局部环的局部同态。事实上，
$\mathfrak p = \psi^{-1}(\mathfrak m)$。由引理
\ref{schemes-lemma-morphism-into-affine}，这就证明了当 $X$ 是仿射概形时的断言。

\medskip\noindent
现设 $X$ 为一般概形。任意 $x \in X$ 都包含在某个仿射开集
$U \subset X$ 中。由仿射情形可知，每个二元组 $(x, \varphi)$ 都由引理前的
构造产生。

\medskip\noindent
要完成证明，只须说明任意态射 $f : \Spec(R) \to X$ 的像都包含在每个
含有闭点像 $x$ 的仿射开集中；这里闭点属于 $\Spec(R)$。事实上，设
$x \in V \subset X$ 是任意含 $x$ 的开邻域。则
$f^{-1}(V) \subset \Spec(R)$ 是包含唯一闭点的开集，因而等于 $\Spec(R)$。
\end{proof}

\noindent
作为上面引理的特例，对每一点 $x$，若它属于概形 $X$，便得到典范态射
\begin{equation}
\label{schemes-equation-canonical-morphism}
\Spec(\mathcal{O}_{X, x}) \longrightarrow X
\end{equation}
它对应于 $X$ 在 $x$ 处局部环上的恒等映射。上面的引理也可以重述为：
对任意态射 $f : \Spec(R) \to X$，存在唯一点 $x \in X$，使得 $f$ 分解为
$\Spec(R) \to \Spec(\mathcal{O}_{X, x}) \to X$
其中第一个映射来自局部同态 $\mathcal{O}_{X, x} \to R$。

\medskip\noindent
若有概形态射 $f : X \to S$，且点 $x$ 映到点 $s \in S$，则得到交换图
$$
\xymatrix{
\Spec(\mathcal{O}_{X, x}) \ar[r] \ar[d] & X \ar[d] \\
\Spec(\mathcal{O}_{S, s}) \ar[r] & S
}
$$
其中左侧竖直映射对应于局部环同态
$f^\sharp_x : \mathcal{O}_{S, s} \to \mathcal{O}_{X, x}$。

\begin{lemma}
\label{schemes-lemma-specialize-points}
设 $X$ 为概形。设 $x, x' \in X$ 为 $X$ 的点。则 $x' \in X$ 是 $x$ 的泛化，
当且仅当 $x'$ 属于典范态射 $\Spec(\mathcal{O}_{X, x}) \to X$ 的像。
\end{lemma}

\begin{proof}
连续映射保持特化／泛化关系。由于 $\Spec(\mathcal{O}_{X, x})$ 的每个点
都是闭点的泛化，可见 $\Spec(\mathcal{O}_{X, x}) \to X$ 像中的每一点
都是 $x$ 的泛化。反过来，设 $x'$ 是 $x$ 的泛化。选取仿射开邻域
$U = \Spec(R)$，使它成为 $x$ 的邻域。则 $x' \in U$。设素理想
$\mathfrak p \subset R$ 与
$\mathfrak p' \subset R$ 分别对应于 $x$ 和 $x'$。由于 $x'$ 是 $x$ 的
泛化，可见 $\mathfrak p' \subset \mathfrak p$。这意味着 $\mathfrak p'$
属于态射
$\Spec(\mathcal{O}_{X, x}) = \Spec(R_{\mathfrak p})
\to \Spec(R) = U \subset X$ as desired.
\end{proof}

\noindent
现在讨论从域的谱出发的态射。设 $(R, \mathfrak m, \kappa)$ 为局部环，
其极大理想为 $\mathfrak m$，剩余域为 $\kappa$。设 $K$ 为域。按定义，
局部同态 $R \to K$ 分解为 $R \to \kappa \to K$；也就是说，它等同于态射
$\kappa \to K$。所以可见，态射
$$
\Spec(K) \longrightarrow X
$$
对应于二元组 $(x, \kappa(x) \to K)$。可以在从域的谱到 $X$ 的态射上
定义预序：若 $\Spec(K) \to X$ 通过 $\Spec(L) \to X$ 分解，
就说 $\Spec(K) \to X$ 支配 $\Spec(L) \to X$。这提示如下概念：暂且称
两个态射 $p : \Spec(K) \to X$ 与 $q : \Spec(L) \to X$ {\it 等价}，
如果存在第三个域 $\Omega$ 以及交换图
$$
\xymatrix{
\Spec(\Omega) \ar[r] \ar[d] &
\Spec(L) \ar[d]^q \\
\Spec(K) \ar[r]^p &
X
}
$$
当然，这立即说明三个概形 $\Spec(K)$、$\Spec(L)$ 和 $\Spec(\Omega)$ 的
唯一点都映到同一个 $x \in X$。因此，由上面的讨论，这样的图对应于一点
$x \in X$ 以及如下域的交换图：
$$
\xymatrix{
\Omega &
L \ar[l] \\
K \ar[u] &
\kappa(x) \ar[l] \ar[u]
}
$$
这定义了一个等价关系，因为给定任意一族域扩张 $K_i/\kappa$，都存在某个
域扩张 $\Omega/\kappa$，使得所有域扩张 $K_i$ 都包含在扩张 $\Omega$ 中。

\begin{lemma}
\label{schemes-lemma-characterize-points}
设 $X$ 为概形。$X$ 的点与从域的谱到 $X$ 的态射之等价类一一对应。
此外，每个等价类都包含一个最小元 $\Spec(\kappa(x)) \to X$，
它在唯一同构意义下唯一。
\end{lemma}

\begin{proof}
这来自上面的讨论。
\end{proof}

\noindent
当然，态射 $\Spec(\kappa(x)) \to X$ 通过典范态射
$\Spec(\mathcal{O}_{X, x}) \to X$ 分解。在当前语境下，引理
\ref{schemes-lemma-specialize-points} 的内容是：态射 $\Spec(\kappa(x')) \to X$
分解为
$\Spec(\kappa(x')) \to \Spec(\mathcal{O}_{X, x}) \to X$
只要 $x'$ 是 $x$ 的泛化。若有概形态射 $f : X \to S$，且点 $x$ 映到点
$s \in S$，则得到交换图
$$
\xymatrix{
\Spec(\kappa(x)) \ar[r] \ar[d] &
\Spec(\mathcal{O}_{X, x}) \ar[r] \ar[d] &
X \ar[d] \\
\Spec(\kappa(s)) \ar[r] &
\Spec(\mathcal{O}_{S, s}) \ar[r] &
S.
}
$$










\section{概形的粘合}
\label{schemes-section-glueing-schemes}

\noindent
设 $I$ 为集合。对每个 $i \in I$，设 $(X_i, \mathcal{O}_i)$ 为局部环层空间。
（事实上，下面的构造对环层空间同样适用。）对每一对 $i, j \in I$，设
$U_{ij} \subset X_i$ 为开子空间。对每一对 $i, j \in I$，设
$$
\varphi_{ij} : U_{ij} \to U_{ji}
$$
为局部环层空间的同构。为方便起见，假设 $U_{ii} = X_i$。对每个三元组
$i, j, k \in I$，假设：
\begin{enumerate}
\item 有
$\varphi_{ij}^{-1}(U_{ji} \cap U_{jk}) =  U_{ij} \cap U_{ik}$；
\item 下图
$$
\xymatrix{
U_{ij} \cap U_{ik} \ar[rr]_{\varphi_{ik}} \ar[rd]_{\varphi_{ij}} & &
U_{ki} \cap U_{kj} \\
& U_{ji} \cap U_{jk} \ar[ru]_{\varphi_{jk}}
}
$$
交换。
\end{enumerate}
（特别地，由此可得 $\varphi_{ii} = \text{id}_{X_i}$。）称一个族
$(I, (X_i)_{i\in I}, (U_{ij})_{i, j\in I}, (\varphi_{ij})_{i, j\in I})$
若满足上述条件，就是一组粘合数据。

\begin{lemma}
\label{schemes-lemma-glue}
\begin{slogan}
若有两个局部环层空间，并且第一个的一个子空间同构于另一个的一个子空间，
就可以把它们粘合成一个大的局部环层空间。
\end{slogan}
给定局部环层空间的任意粘合数据，存在局部环层空间 $X$、开子空间
$U_i \subset X$，以及局部环层空间的同构 $\varphi_i : X_i \to U_i$，使得
\begin{enumerate}
\item $X=\bigcup_{i\in I} U_i$,
\item $\varphi_i(U_{ij}) = U_i \cap U_j$, and
\item $\varphi_{ij} =
\varphi_j^{-1}|_{U_i \cap U_j} \circ \varphi_i|_{U_{ij}}$.
\end{enumerate}
局部环层空间 $X$ 由如下映射性质刻画：给定局部环层空间 $Y$，有
\begin{eqnarray*}
\Mor(X, Y) & = & \{ (f_i)_{i\in I} \mid
f_i : X_i \to Y, \ f_j \circ \varphi_{ij} = f_i|_{U_{ij}}\} \\
f & \mapsto & (f|_{U_i} \circ \varphi_i)_{i \in I} \\
\Mor(Y, X) & = &
\left\{
\begin{matrix}
\text{开覆盖 }Y = \bigcup\nolimits_{i \in I} V_i\text{ 且 }
(g_i : V_i \to X_i)_{i \in I}
\text{ 使得}\\
g_i^{-1}(U_{ij}) = V_i \cap V_j
\text{ 且 }
g_j|_{V_i \cap V_j} = \varphi_{ij} \circ g_i|_{V_i \cap V_j}
\end{matrix}
\right\} \\
g & \mapsto &
V_i = g^{-1}(U_i), \ g_i = \varphi_i^{-1} \circ g|_{V_i}
\end{eqnarray*}
\end{lemma}

\begin{proof}
分阶段构造 $X$。作为集合，取
$$
X = (\coprod X_i) / \sim.
$$
这里，给定 $x \in X_i$ 与 $x' \in X_j$，规定 $x \sim x'$ 当且仅当
$x \in U_{ij}$、$x' \in U_{ji}$ 且 $\varphi_{ij}(x) = x'$。这是等价关系：
若 $x \in X_i$、$x' \in X_j$、$x'' \in X_k$，并且 $x \sim x'$、
$x' \sim x''$，则 $x' \in U_{ji} \cap U_{jk}$；于是由粘合数据的条件 (1)，
还有 $x \in U_{ij} \cap U_{ik}$ 与 $x'' \in U_{ki} \cap U_{kj}$，再由条件
(2) 得到 $\varphi_{ik}(x) = x''$。（自反性与对称性来自假设
$U_{ii} = X_i$ 和 $\varphi_{ii} = \text{id}_{X_i}$。）把自然映射记作
$\varphi_i : X_i \to X$，并记 $U_i = \varphi_i(X_i) \subset X$。
注意，$\varphi_i : X_i \to U_i$ 是双射。

\medskip\noindent
$X$ 上的拓扑按如下规则定义：$U \subset X$ 为开集，当且仅当
$\varphi_i^{-1}(U)$ 对所有 $i$ 都是开集。留给读者验证，这确实定义了拓扑。
特别地，注意 $U_i$ 是开集，因为 $\varphi_j^{-1}(U_i)
= U_{ji}$ 是 $X_j$ 中的开集；这对所有 $j$ 都成立。此外，对任意开集
$W \subset X_i$，其像
$\varphi_i(W) \subset U_i$ 是开集，因为
$\varphi_j^{-1}(\varphi_i(W)) = \varphi_{ji}^{-1}(W \cap U_{ij})$。
因此，$\varphi_i : X_i \to U_i$ 是同胚。

\medskip\noindent
要得到局部环层空间，还须构造环层 $\mathcal{O}_X$。我们通过粘合各环层
$\mathcal{O}_{U_i} := \varphi_{i, *} \mathcal{O}_i$ 来完成。具体而言，
在交换图
$$
\xymatrix{
U_{ij} \ar[rr]_{\varphi_{ij}} \ar[rd]_{\varphi_i|_{U_{ij}}}
& &
U_{ji} \ar[ld]^{\varphi_j|_{U_{ji}}} \\
& U_i \cap U_j &
}
$$
中，上方箭头是环层空间的同构，因而得到唯一的环层同构
$$
\mathcal{O}_{U_i}|_{U_i \cap U_j}
\longrightarrow
\mathcal{O}_{U_j}|_{U_i \cap U_j}.
$$
它们满足《层》第 \ref{sheaves-section-glueing-sheaves} 节中的余循环条件。
由该节的结果，得到环层 $\mathcal{O}_X$；它定义在 $X$ 上，并使得
$\mathcal{O}_X|_{U_i}$ 与 $\mathcal{O}_{U_i}$ 同构，并与上面展示的粘合映射
相容。特别地，$(X, \mathcal{O}_X)$ 是局部环层空间，因为
$\mathcal{O}_X$ 的茎等于对应点处 $\mathcal{O}_i$ 的茎。

\medskip\noindent
映射性质的证明略去。
\end{proof}

\begin{lemma}
\label{schemes-lemma-glue-schemes}
\begin{slogan}
概形可以粘合成新的概形。
\end{slogan}
在上面的引理 \ref{schemes-lemma-glue} 中，假设所有 $X_i$ 都是概形。
则所得局部环层空间 $X$ 是概形。
\end{lemma}

\begin{proof}
这是显然的，因为每个 $U_i$ 都是概形，因而每个 $x \in X$ 都有仿射邻域。
\end{proof}

\noindent
通常把 $X_i$ 看成 $X$ 的开子空间；所借助的是同构 $\varphi_i$。
下面两个例子中都采用这一约定。

\begin{example}[零点加倍的仿射空间]
\label{schemes-example-affine-space-zero-doubled}
设 $k$ 为域，$n \geq 1$。令 $X_1 = \Spec(k[x_1, \ldots, x_n])$，
令 $X_2 = \Spec(k[y_1, \ldots, y_n])$。令 $0_1 \in X_1$ 为对应于极大理想
$(x_1, \ldots, x_n) \subset k[x_1, \ldots, x_n]$ 的点。令 $0_2 \in X_2$
为对应于极大理想 $(y_1, \ldots, y_n) \subset k[y_1, \ldots, y_n]$ 的点。
令 $U_{12} = X_1 \setminus \{0_1\}$，并令
$U_{21} = X_2 \setminus \{0_2\}$。设 $\varphi_{12} : U_{12} \to U_{21}$
为来自 $k$-代数同构
$k[y_1, \ldots, y_n] \to k[x_1, \ldots, x_n]$ 的同构；该同构把 $y_i$
映到 $x_i$（并诱导出 $X_1 \cong X_2$，把 $0_1$ 映到 $0_2$）。
设 $X$ 为由粘合数据
$(X_1, X_2, U_{12}, U_{21}, \varphi_{12},
\varphi_{21} = \varphi_{12}^{-1})$
所得的概形。按照例子上方引入的轻微记号滥用，把 $X_1, X_2 \subset X$
看成开子概形。存在态射 $f : X \to \Spec(k[t_1, \ldots, t_n])$，
它在 $X_1$ 上（相应地，在 $X_2$ 上）对应于 $k$-代数同态
$k[t_1, \ldots, t_n] \to k[x_1, \ldots, x_n]$
(resp.\ $k[t_1, \ldots, t_n] \to k[y_1, \ldots, y_n]$)
它把 $t_i$ 映到 $x_i$（相应地，把 $t_i$ 映到 $y_i$）。容易看出，
该态射把 $k[t_1, \ldots, t_n]$ 与 $\Gamma(X, \mathcal{O}_X)$ 认同。
由于 $f(0_1) = f(0_2)$，可见 $X$ 不是仿射的。

\medskip\noindent
注意，$X_1$ 与 $X_2$ 都是 $X$ 的仿射开集。但若 $n = 2$，则
$X_1 \cap X_2$ 是例 \ref{schemes-example-not-affine} 中描述的概形，因而不是仿射的。
所以一般而言，概形的仿射开集之交不再仿射。（更一般地，该事实对任意
$n > 1$ 都成立。）

\medskip\noindent
该例还有如下奇特性质。若 $n > 1$，则存在许多不可约闭子集 $T \subset X$
（例如取 $X_1$ 中任意非闭点的闭包）。但除非 $T = \{0_1\}$ 或
$T = \{0_2\}$，总有 $0_1 \in T \Leftrightarrow 0_2 \in T$。证明略去。
\end{example}

\begin{example}[射影直线]
\label{schemes-example-projective-line}
设 $k$ 为域。令 $X_1 = \Spec(k[x])$，令 $X_2 = \Spec(k[y])$。
令 $0 \in X_1$ 为对应于极大理想 $(x) \subset k[x]$ 的点。
令 $\infty \in X_2$ 为对应于极大理想 $(y) \subset k[y]$ 的点。
令 $U_{12} = X_1 \setminus \{0\} = D(x) = \Spec(k[x, 1/x])$，并令
$U_{21} = X_2 \setminus \{\infty\} = D(y) = \Spec(k[y, 1/y])$。
设 $\varphi_{12} : U_{12} \to U_{21}$ 为来自 $k$-代数同构
$k[y, 1/y] \to k[x, 1/x]$ 的同构；该同构把 $y$ 映到 $1/x$。
设 $\mathbf{P}^1_k$ 为由粘合数据
$(X_1, X_2, U_{12}, U_{21}, \varphi_{12},
\varphi_{21} = \varphi_{12}^{-1})$
所得的概形。按照例子上方引入的轻微记号滥用，把
$X_i \subset \mathbf{P}^1_k$ 看成开子概形。在此情形下，
$\Gamma(\mathbf{P}^1_k, \mathcal{O}) = k$，因为满足如下条件的多项式
$g(x)$ 只有常数多项式：它是 $x$ 中的多项式，$g(1/y)$ 也是 $y$ 中的
多项式。由于
$\mathbf{P}^1_k$ 是无限的，可见 $\mathbf{P}^1_k$ 不是仿射的。

\medskip\noindent
我们断言存在仿射开集 $U \subset \mathbf{P}^1_k$，它同时包含 $0$ 与
$\infty$。事实上，令 $U = \mathbf{P}^1_k \setminus \{1\}$，其中 $1$ 是
$X_1$ 中对应于极大理想 $(x - 1)$ 的点，也是 $X_2$ 中对应于极大理想
$(y - 1)$ 的点。容易看出
$s = 1/(x - 1) = y/(1 - y) \in \Gamma(U, \mathcal{O}_U)$。事实上，可以证明
$\Gamma(U, \mathcal{O}_U)$ 等于多项式环 $k[s]$，且相应态射
$U \to \Spec(k[s])$ 是概形同构。细节略去。
\end{example}




\section{一个可表性判据}
\label{schemes-section-representable}

\noindent
本节用函子的语言重新表述第 \ref{schemes-section-glueing-schemes} 节的粘合引理。
先回顾《范畴》第 \ref{categories-section-opposite} 节的一些材料。
回顾，给定概形 $X$，可以定义函子
$$
h_X : \Sch^{opp}
\longrightarrow
\textit{Sets}, \quad
T \longmapsto \Mor(T, X).
$$
它称为 $X$ 的{\it 点函子}。

\medskip\noindent
设 $F$ 为从概形范畴到集合范畴的反变函子。写成公式即
$$
F : \Sch^{opp}
\longrightarrow
\textit{Sets}.
$$
我们采用《位点》第 \ref{sites-section-presheaves} 节中的相同术语。
具体地，给定概形 $T$、元素 $\xi \in F(T)$ 以及态射 $f : T' \to T$，
用 $f^*\xi$ 表示元素 $F(f)(\xi)$，有时甚至采用记号 $\xi|_{T'}$。

\begin{definition}
\label{schemes-definition-representable-functor}
（见《范畴》的定义 \ref{categories-definition-representable-functor}。）
设 $F$ 为从概形范畴到集合范畴的反变函子（如上）。称 $F$
{\it 由概形表示}，或称它{\it 可表}，如果存在概形 $X$，使得
$h_X \cong F$。
\end{definition}

\noindent
设 $F$ 由概形 $X$ 表示，且 $s : h_X \to F$ 是同构。由《范畴》的米田引理
\ref{categories-lemma-yoneda}，二元组 $(X, s : h_X \to F)$ 若存在，
就在唯一同构意义下唯一。此外，米田引理说明，对上述任意反变函子 $F$
和任意概形 $Y$，都有双射
$$
\Mor_{\text{Fun}(\Sch^{opp}, \textit{Sets})} (h_Y, F)
\longrightarrow
F(Y), \quad
s \longmapsto s(\text{id}_Y).
$$
反向构造如下：给定任意 $\xi \in F(Y)$，函子变换
$s_\xi : h_Y \to F$ 把任意态射 $f : T \to Y$ 送到元素
$f^*\xi \in F(T)$。

\medskip\noindent
特别地，当 $F$ 可表时，存在概形 $X$ 和元素 $\xi \in F(X)$，使得相应态射
$h_X \to F$ 是同构。在这种情形下，也说{\it 二元组 $(X, \xi)$ 表示 $F$}。
元素 $\xi \in F(X)$ 常称为{\it “泛族”}；等到讨论代数栈时，这一名称的
理由会更加清楚（此处将来插入参考）。目前只须注意：若二元组 $(X, \xi)$
表示 $F$，则每个元素 $\xi' \in F(T)$（其中 $T$ 任意）都形如
$\xi' = f^*\xi$，其中 $f : T \to X$ 唯一。

\begin{example}
\label{schemes-example-global-sections}
考虑如下规则：把每个概形 $T$ 送到集合
$F(T) = \Gamma(T, \mathcal{O}_T)$。
对态射 $f : T' \to T$ 使用拉回映射
$f^\sharp : \Gamma(T, \mathcal{O}_T) \to \Gamma(T', \mathcal{O}_{T'})$，
即可把它变成反变函子。给定环 $R$ 和元素 $t \in R$，存在唯一环同态
$\mathbf{Z}[x] \to R$，把 $x$ 映到 $t$。因此，利用引理
\ref{schemes-lemma-morphism-into-affine}，可见
$$
\Mor(T, \Spec(\mathbf{Z}[x])) =
\Hom(\mathbf{Z}[x], \Gamma(T, \mathcal{O}_T)) =
\Gamma(T, \mathcal{O}_T).
$$
这确实给出同构 $h_{\Spec(\mathbf{Z}[x])} \to F$。“泛族” $\xi$ 是什么？
要得到它，必须把上面的认同应用于
$\text{id}_{\Spec(\mathbf{Z}[x])}$。显然，在这些认同下得到
$\xi = x \in \Gamma(\Spec(\mathbf{Z}[x]),
\mathcal{O}_{\Spec(\mathbf{Z}[x])}) = \mathbf{Z}[x]$
，正如预期。
\end{example}

\begin{definition}
\label{schemes-definition-representable-by-open-immersions}
设 $F$ 为概形范畴上取值于集合的反变函子。
\begin{enumerate}
\item 称 $F$ {\it 满足 Zariski 拓扑的层性质}，如果对每个概形 $T$、
每个开覆盖 $T = \bigcup_{i \in I} U_i$，
以及任意一族元素 $\xi_i \in F(U_i)$，只要
$\xi_i|_{U_i \cap U_j} =
\xi_j|_{U_i \cap U_j}$，就存在唯一元素
$\xi \in F(T)$，使得 $\xi_i = \xi|_{U_i}$ 在 $F(U_i)$ 中成立。
\item {\it 子函子 $H \subset F$} 是如下规则：对每个概形 $T$，赋予子集
$H(T) \subset F(T)$，使得映射 $F(f) : F(T) \to F(T')$ 把 $H(T)$
映入 $H(T')$；这对所有概形态射 $f : T' \to T$ 都成立。
\item 设 $H \subset F$ 为子函子。称 $H \subset F$ {\it 由开浸入表示}，
如果对所有二元组 $(T, \xi)$，其中 $T$ 是概形且 $\xi \in F(T)$，
都存在具有如下性质的开子概形 $U_\xi \subset T$：
\begin{itemize}
\item[(*)] 态射 $f : T' \to T$ 通过 $U_\xi$ 分解，当且仅当
$f^*\xi \in H(T')$。
\end{itemize}
\item 设 $I$ 为集合。对每个 $i \in I$，设 $H_i \subset F$ 为子函子。
称族 $(H_i)_{i \in I}$ {\it 覆盖 $F$}，当且仅当对每个
$\xi \in F(T)$，都存在开覆盖 $T = \bigcup U_i$，使得
$\xi|_{U_i} \in H_i(U_i)$。
\end{enumerate}
\end{definition}

\noindent
在条件 (4) 中，若 $H_i \subset F$ 对所有 $i$ 都由开浸入表示，
则要检验 $(H_i)_{i \in I}$ 覆盖 $F$，只须检验
$F(T) = \bigcup H_i(T)$；这里 $T$ 是域的谱。

\begin{lemma}
\label{schemes-lemma-glue-functors}
设 $F$ 为概形范畴上取值于集合范畴的反变函子。假设：
\begin{enumerate}
\item $F$ 满足 Zariski 拓扑的层性质；
\item 存在集合 $I$ 以及一族子函子 $F_i \subset F$，使得
\begin{enumerate}
\item 每个 $F_i$ 可表；
\item 每个 $F_i \subset F$ 由开浸入表示；
\item 族 $(F_i)_{i \in I}$ 覆盖 $F$。
\end{enumerate}
\end{enumerate}
则 $F$ 可表。
\end{lemma}

\begin{proof}
令 $X_i$ 为表示 $F_i$ 的概形，并令
$\xi_i \in F_i(X_i) \subset F(X_i)$ 为“泛族”。
由于 $F_j \subset F$ 由开浸入表示，
存在开子集 $U_{ij} \subset X_i$，使得
$T \to X_i$ 经由 $U_{ij}$ 分解，当且仅当
$\xi_i|_T \in F_j(T)$。特别地，
$\xi_i|_{U_{ij}} \in F_j(U_{ij})$，因而得到典范态射
$\varphi_{ij} : U_{ij} \to X_j$，满足
$\varphi_{ij}^*\xi_j = \xi_i|_{U_{ij}}$。按照 $U_{ji}$ 的定义，
这表明 $\varphi_{ij}$ 经由 $U_{ji}$ 分解。
由于 $(\varphi_{ij} \circ \varphi_{ji})^*\xi_j
=\varphi_{ji}^*(\varphi_{ij}^*\xi_j) =
\varphi_{ji}^*\xi_i = \xi_j$，我们断定
$\varphi_{ij} \circ \varphi_{ji} = \text{id}_{U_{ji}}$，
这是因为二元组 $(X_j, \xi_j)$ 表示 $F_j$。
特别地，映射 $\varphi_{ij} : U_{ij} \to U_{ji}$
都是概形同构。
接着须证明
$\varphi_{ij}^{-1}(U_{ji} \cap U_{jk}) = U_{ij} \cap U_{ik}$。
这是因为：(a) $U_{ji} \cap U_{jk}$ 是 $U_{ji}$ 中使
$\xi_j$ 限制为 $F_k$ 的一个元素的最大开子集；
(b) $U_{ij} \cap U_{ik}$ 是 $U_{ij}$ 中使
$\xi_i$ 限制为 $F_k$ 的一个元素的最大开子集；
(c) $\varphi_{ij}^*\xi_j = \xi_i$。此外，
第 \ref{schemes-section-glueing-schemes} 节中的余圈条件成立，因为
$\varphi_{jk}|_{U_{ji} \cap U_{jk}} \circ
\varphi_{ij}|_{U_{ij} \cap U_{ik}}$ 与
$\varphi_{ik}|_{U_{ij} \cap U_{ik}}$ 都把 $\xi_k$
拉回为元素 $\xi_i$。
因此可以应用引理 \ref{schemes-lemma-glue-schemes}，
得到概形 $X$、开覆盖
$X = \bigcup U_i$，以及同构 $\varphi_i : X_i \to U_i$，
它们具有引理 \ref{schemes-lemma-glue} 中的性质。
令 $\xi_i' = (\varphi_i^{-1})^* \xi_i$。
引理 \ref{schemes-lemma-glue} 的条件蕴含
$\xi_i'|_{U_i \cap U_j} = \xi_j'|_{U_i \cap U_j}$。
因此，由于 $F$ 满足 Zariski 拓扑的层条件，
可知存在元素
$\xi' \in F(X)$，使得
$\xi_i = \varphi_i^*\xi'|_{U_i}$ 对所有 $i$ 都成立。
由于 $\varphi_i$ 是同构，还可得到
$(U_i, \xi'|_{U_i})$ 表示函子 $F_i$。

\medskip\noindent
我们声称二元组 $(X, \xi')$ 表示函子 $F$。
为证此点，令 $T$ 为概形，并令 $\xi \in F(T)$。
我们将构造唯一态射 $g : T \to X$，使得
$g^*\xi' = \xi$。确切地说，由子函子
$F_i$ 覆盖 $F$ 这一条件，存在开覆盖 $T = \bigcup V_i$，
使得对每个 $i$，限制 $\xi|_{V_i} \in F_i(V_i)$。
此外，由于每个包含 $F_i \subset F$ 都由开浸入表示，
可假设每个 $V_i \subset T$ 都是具有此性质的最大开子集。
由于 $(U_i, \xi'|_{U_i})$ 表示函子 $F_i$，
得到唯一态射 $g_i : V_i \to U_i$，使得
$g_i^*\xi'|_{U_i} = \xi|_{V_i}$。在交集 $V_i \cap V_j$ 上，
态射 $g_i$ 与 $g_j$ 相同；例如，这是因为它们都把
$\xi'|_{U_i \cap U_j} \in F_i(U_i \cap U_j)$
拉回为同一个元素。因此态射 $g_i$ 粘合成唯一态射
$T \to X$，正合所需。
\end{proof}

\begin{remark}
\label{schemes-remark-representable-locally-ringed}
设函子 $F$ 定义在所有局部环层空间上，
并把引理 \ref{schemes-lemma-glue-functors} 的条件替换为以下条件：
\begin{enumerate}
\item $F$ 在局部环层空间范畴上满足层性质；
\item 存在集合 $I$ 以及一族子函子
$F_i \subset F$，使得
\begin{enumerate}
\item 每个 $F_i$ 都由概形表示；
\item 每个 $F_i \subset F$ 在局部环层空间范畴上都由开浸入表示；
\item 族 $(F_i)_{i \in I}$ 在局部环层空间范畴上的函子意义下覆盖 $F$。
\end{enumerate}
\end{enumerate}
我们把进一步写出细节留给读者。
最终结论是：函子 $F$ 在局部环层空间范畴中可表，
并且表示对象是一个概形。
\end{remark}



\section{概形纤维积的存在性}
\label{schemes-section-existence-fibre-products}

\noindent
一个非常基本的问题是：概形范畴中是否存在乘积与纤维积。
我们先抽象地证明乘积与纤维积存在，下一节再说明如何以一种
合理的方式理解概形的纤维积。

\begin{lemma}
\label{schemes-lemma-fibre-products}
概形范畴具有终对象、乘积和纤维积。
换言之，概形范畴具有有限极限；见《范畴》，
引理 \ref{categories-lemma-finite-limits-exist}。
\end{lemma}

\begin{proof}
请跳过这个证明。更重要的是学习如何使用下一节所说明的纤维积。

\medskip\noindent
由引理 \ref{schemes-lemma-morphism-into-affine}，
概形 $\Spec(\mathbf{Z})$ 是局部环层空间范畴中的终对象。
因此只须证明纤维积存在。

\medskip\noindent
令 $f : X \to S$ 与 $g : Y \to S$ 为概形态射。须证明函子
\begin{eqnarray*}
F : \Sch^{opp} & \longrightarrow & \textit{Sets} \\
T & \longmapsto &
\Mor(T, X) \times_{\Mor(T, S)} \Mor(T, Y)
\end{eqnarray*}
可表。我们声称引理 \ref{schemes-lemma-glue-functors}
适用于函子 $F$。证明此点便证得本引理。

\medskip\noindent
首先证明 $F$ 满足 Zariski 拓扑的层性质。具体而言，设 $T$ 为概形，
$T = \bigcup_{i \in I} U_i$ 为开覆盖，并且
$\xi_i \in F(U_i)$ 满足
$\xi_i|_{U_i \cap U_j} =  \xi_j|_{U_i \cap U_j}$；这对所有二元组
$i, j$ 都成立。按照定义，$\xi_i$ 对应于二元组 $(a_i, b_i)$，其中
$a_i : U_i \to X$ 与 $b_i : U_i \to Y$ 是概形态射，满足
$f \circ a_i = g \circ b_i$。粘合条件说明
$a_i|_{U_i \cap U_j} =  a_j|_{U_i \cap U_j}$
and
$b_i|_{U_i \cap U_j} =  b_j|_{U_i \cap U_j}$.
因此，粘合态射 $a_i$ 得到一个局部环层空间态射
（即概形态射）$a : T \to X$；同样得到 $b : T \to Y$
（例如见引理 \ref{schemes-lemma-glue} 的映射性质）。此外，在某个开覆盖的
各成员上，复合 $f \circ a$ 与 $g \circ b$ 相同。因此
$f \circ a = g \circ b$，二元组 $(a, b)$ 定义了 $F(T)$ 的一个元素，
它限制为二元组 $(a_i, b_i)$，在每个 $U_i$ 上皆如此。层条件得证。

\medskip\noindent
接着构造子函子族。选取由仿射开子集构成的开覆盖
$S = \bigcup\nolimits_{i \in I} U_i$。
对每个 $i \in I$，选取由仿射开子集构成的开覆盖
$f^{-1}(U_i) = \bigcup\nolimits_{j \in J_i} V_j$，以及
$g^{-1}(U_i) = \bigcup\nolimits_{k \in K_i} W_k$。
注意，$X = \bigcup_{i \in I} \bigcup_{j \in J_i} V_j$
是开覆盖；$Y$ 也同样。对任意 $i \in I$ 以及每一对
$(j, k) \in J_i \times K_i$，都有交换图
$$
\xymatrix{
    & W_k \ar[d] \ar[rd] &   \\
V_j \ar[rd] \ar[r] & U_i \ar[rd] & Y \ar[d] \\
    & X \ar[r]  & S
}
$$
其中所有斜箭头都是开浸入。对这样的三元组，得到函子
\begin{eqnarray*}
F_{i, j, k} : \Sch^{opp} & \longrightarrow & \textit{Sets} \\
T & \longmapsto &
\Mor(T, V_j) \times_{\Mor(T, U_i)} \Mor(T, W_k).
\end{eqnarray*}
存在显然的函子变换 $F_{i, j, k} \to F$
（来自上面的巨大交换图），它是单射，因此可把 $F_{i, j, k}$
视为 $F$ 的子函子。

\medskip\noindent
检验引理 \ref{schemes-lemma-glue-functors} 的条件 (2)(a)。
这直接由引理 \ref{schemes-lemma-fibre-product-affine-schemes} 得出。
（注意，这里用到了：仿射概形范畴中的纤维积也是整个局部环层空间
范畴中的纤维积。）

\medskip\noindent
检验引理 \ref{schemes-lemma-glue-functors} 的条件 (2)(b)。
令 $T$ 为概形，并令 $\xi \in F(T)$。换言之，
$\xi = (a, b)$，其中 $a : T \to X$ 与 $b : T \to Y$ 是概形态射，
满足 $f \circ a = g \circ b$。
令 $V_{i, j, k} = a^{-1}(V_j) \cap b^{-1}(W_k)$。对任意另一个态射
$h : T' \to T$，都有
$h^*\xi = (a \circ h, b \circ h)$。因此可知，
$h^*\xi \in F_{i, j, k}(T')$ 当且仅当
$a(h(T')) \subset V_j$ 且 $b(h(T')) \subset W_k$。
换言之，当且仅当 $h(T') \subset V_{i, j, k}$。这证明了条件 (2)(b)。

\medskip\noindent
检验引理 \ref{schemes-lemma-glue-functors} 的条件 (2)(c)。
令 $T$ 为概形，并令 $\xi = (a, b) \in F(T)$ 如上。
仍令 $V_{i, j, k} = a^{-1}(V_j) \cap b^{-1}(W_k)$。
条件 (2)(c) 正是说 $T = \bigcup V_{i, j, k}$，这显然成立。
因此引理得证，纤维积存在。
\end{proof}

\begin{remark}
\label{schemes-remark-fibre-product-schemes-locally-ringed}
利用注 \ref{schemes-remark-representable-locally-ringed}，
可以证明：概形态射的纤维积在局部环层空间范畴中存在，
并且是概形。
\end{remark}




\section{概形的纤维积}
\label{schemes-section-fibre-products}

\noindent
尽管已经证明了概形的纤维积存在，这里仍回顾其一般定义。

\begin{definition}
\label{schemes-definition-fibre-product}
给定概形态射 $f : X \to S$ 与 $g : Y \to S$，
{\it 纤维积}是概形 $X \times_S Y$ 连同投影态射
$p : X \times_S Y \to X$ 与 $q : X \times_S Y \to Y$，
它们处于下列交换图中：
$$
\xymatrix{
X \times_S Y \ar[r]_q \ar[d]_p & Y \ar[d]^g \\
X \ar[r]^f & S
}
$$
并且此图在所有这类图中具有泛性质；见《范畴》，
定义 \ref{categories-definition-fibre-products}。
\end{definition}

\noindent
换言之，给定任意由概形态射组成的实线交换图
$$
\xymatrix{
T \ar[rrrd] \ar@{-->}[rrd] \ar[rrdd]
&
&
\\
&
&
X \times_S Y \ar[d] \ar[r]
&
Y \ar[d]
\\
&
&
X \ar[r]
&
S
}
$$
存在唯一的虚线箭头使该图交换。我们将证明若干引理，
以说明应如何理解纤维积。

\begin{lemma}
\label{schemes-lemma-fibre-product-affines}
令 $f : X \to S$ 与 $g : Y \to S$ 为具有相同目标的概形态射。
若 $X, Y, S$ 都是仿射的，则 $X \times_S Y$ 是仿射的。
\end{lemma}

\begin{proof}
设 $X = \Spec(A)$、$Y = \Spec(B)$，且
$S = \Spec(R)$。由引理 \ref{schemes-lemma-fibre-product-affine-schemes}，
仿射概形 $\Spec(A \otimes_R B)$ 是局部环层空间范畴中的纤维积
$X \times_S Y$。因而它当然也是概形范畴中的纤维积。
\end{proof}

\begin{lemma}
\label{schemes-lemma-open-fibre-product}
令 $f : X \to S$ 与 $g : Y \to S$ 为具有相同目标的概形态射。
令 $X \times_S Y$、$p$、$q$ 为该纤维积。设
$U \subset S$，$V \subset X$、$W \subset Y$ 是开子概形，满足
$f(V) \subset U$ 且 $g(W) \subset U$。则典范态射
$V \times_U W \to X \times_S Y$ 是开浸入，它把
$V \times_U W$ 等同于 $p^{-1}(V) \cap q^{-1}(W)$。
\end{lemma}

\begin{proof}
令 $T$ 为概形。设 $a : T \to V$ 与 $b : T \to W$ 为态射，
并且有 $f \circ a = g \circ b$，这里把二者视为到 $U$ 的态射。
那么它们作为到 $S$ 的态射也相同。由纤维积的泛性质，
得到唯一态射 $T \to X \times_S Y$。当然，该态射的像包含在开集
$p^{-1}(V) \cap q^{-1}(W)$ 中。因此
$p^{-1}(V) \cap q^{-1}(W)$ 是 $V$ 与 $W$ 在 $U$ 上的纤维积。
结论由纤维积的唯一性得出；见《范畴》，
第 \ref{categories-section-fibre-products} 节。
\end{proof}

\noindent
特别地，这说明在该引理的情形下
$V \times_U W = V \times_S W$。此外，若 $U, V, W$ 都是仿射的，
则已知 $V \times_U W$ 是仿射的。当然，可以覆盖 $X \times_S Y$，
所用的是这样的仿射开集 $V \times_U W$。将此表述为一个引理。

\begin{lemma}
\label{schemes-lemma-affine-covering-fibre-product}
\begin{slogan}
纤维积的直接构造：概形纤维积的仿射开覆盖可以由各部分相容的
仿射开覆盖拼装而成。
\end{slogan}
令 $f : X \to S$ 与 $g : Y \to S$ 为具有相同目标的概形态射。
令 $S = \bigcup U_i$ 为 $S$ 的任意仿射开覆盖。对每个 $i \in I$，
令 $f^{-1}(U_i) = \bigcup_{j \in J_i} V_j$ 为 $f^{-1}(U_i)$ 的仿射开覆盖，
并令 $g^{-1}(U_i) = \bigcup_{k \in K_i} W_k$ 为 $g^{-1}(U_i)$ 的仿射开覆盖。
则
$$
X \times_S Y =
\bigcup\nolimits_{i \in I}
\bigcup\nolimits_{j \in J_i, \ k \in K_i}
V_j \times_{U_i} W_k
$$
是 $X \times_S Y$ 的仿射开覆盖。
\end{lemma}

\begin{proof}
见本引理上方的讨论。
\end{proof}

\noindent
换言之，我们本可以利用前一个引理，通过粘合仿射概形直接构造纤维积。
（当然，这恰好正是引理 \ref{schemes-lemma-fibre-products} 的证明中所做的。）
下面给出描述概形纤维积点集的一种方式。

\begin{lemma}
\label{schemes-lemma-points-fibre-product}
令 $f : X \to S$ 与 $g : Y \to S$ 为具有相同目标的概形态射。
点 $z$（视为 $X \times_S Y$ 的点）与下列四元组一一对应：
$$
(x, y, s, \mathfrak p)
$$
其中 $x \in X$、$y \in Y$、$s \in S$ 是满足
$f(x) = s$、$g(y) = s$ 的点，而 $\mathfrak p$ 是环
$\kappa(x) \otimes_{\kappa(s)} \kappa(y)$ 的素理想。
$z$ 的剩余域对应于素理想 $\mathfrak p$ 的剩余域。
\end{lemma}

\begin{proof}
令 $z$ 为 $X \times_S Y$ 的一个点，并构造上述四元组。
回忆可以把 $z$ 看成态射
$\Spec(\kappa(z)) \to X \times_S Y$；见引理
\ref{schemes-lemma-characterize-points}。此态射对应于态射
$a : \Spec(\kappa(z)) \to X$ 与
$b : \Spec(\kappa(z)) \to Y$，它们满足
$f \circ a = g \circ b$。再次应用同一引理，得到点
$x \in X$、$y \in Y$，它们位于同一点 $s \in S$ 上方；还得到域映射
$\kappa(x) \to \kappa(z)$、$\kappa(y) \to \kappa(z)$，使得复合
$\kappa(s) \to \kappa(x) \to \kappa(z)$ 与
$\kappa(s) \to \kappa(y) \to \kappa(z)$ 相同。换言之，得到环映射
$\kappa(x) \otimes_{\kappa(s)} \kappa(y) \to \kappa(z)$。
令 $\mathfrak p$ 为此映射的核。

\medskip\noindent
反过来，给定四元组 $(x, y, s, \mathfrak p)$，得到实线交换图
$$
\xymatrix{
X \times_S Y
\ar@/_/[dddr] \ar@/^/[rrrd]
& & & \\
&
\Spec(\kappa(x) \otimes_{\kappa(s)} \kappa(y)/\mathfrak p)
\ar[r] \ar[d] \ar@{-->}[lu]
&
\Spec(\kappa(y)) \ar[d] \ar[r] &
Y \ar[dd] \\
&
\Spec(\kappa(x)) \ar[r] \ar[d] &
\Spec(\kappa(s)) \ar[rd] &
\\
&
X \ar[rr] &
&
S
}
$$
见第 \ref{schemes-section-points} 节中的讨论。因而得到虚线箭头。
相应的点 $z$（作为 $X \times_S Y$ 的点）是下列概形泛点的像：
$\Spec(\kappa(x) \otimes_{\kappa(s)} \kappa(y)/\mathfrak p)$。
我们略去验证这两个构造互逆的过程。
\end{proof}

\begin{lemma}
\label{schemes-lemma-fibre-product-immersion}
令 $f : X \to S$ 与 $g : Y \to S$ 为具有相同目标的概形态射。
\begin{enumerate}
\item 若 $f : X \to S$ 是闭浸入，则
$X \times_S Y \to Y$ 是闭浸入。此外，若 $X \to S$ 对应于拟凝聚理想层
$\mathcal{I} \subset \mathcal{O}_S$，则 $X \times_S Y \to Y$ 对应于理想层
$\Im(g^*\mathcal{I} \to \mathcal{O}_Y)$。
\item 若 $f : X \to S$ 是开浸入，则
$X \times_S Y \to Y$ 是开浸入。
\item 若 $f : X \to S$ 是浸入，则
$X \times_S Y \to Y$ 是浸入。
\end{enumerate}
\end{lemma}

\begin{proof}
假设 $X \to S$ 是对应于拟凝聚理想层
$\mathcal{I} \subset \mathcal{O}_S$ 的闭浸入。由引理
\ref{schemes-lemma-restrict-map-to-closed}，闭子空间 $Z \subset Y$，
它由理想层 $\Im(g^*\mathcal{I} \to \mathcal{O}_Y)$ 定义，
是局部环层空间范畴中的纤维积。
由引理 \ref{schemes-lemma-closed-subspace-scheme}，$Z$ 是概形。
因此 $Z = X \times_S Y$，第一项得证。例如，第二项由引理
\ref{schemes-lemma-open-fibre-product} 得出。第三项是前两项的结合。
\end{proof}

\begin{definition}
\label{schemes-definition-inverse-image-closed-subscheme}
令 $f : X \to Y$ 为概形态射。令 $Z \subset Y$ 为 $Y$ 的闭子概形。
所谓{\it 逆像 $f^{-1}(Z)$，即闭子概形 $Z$ 的逆像}，是闭子概形
$Z \times_Y X$（视为 $X$ 的闭子概形）。见上面的引理
\ref{schemes-lemma-fibre-product-immersion}。
\end{definition}

\noindent
有时我们也把这个术语用于局部闭子概形和开子概形。









\section{代数几何中的基变换}
\label{schemes-section-base-change}

\noindent
引入概形语言的一个动机是：它精确刻画了在某个给定域上定义一个簇
是什么意思。例如，簇 $X$ 在 $\mathbf{Q}$ 上，等价于
（《簇》，定义 \ref{varieties-definition-variety}）一个有限型、分离、
不可约且约化的态射 $X \to \Spec(\mathbf{Q})$\footnote{当然，代数几何学家
仍会争论是否应当要求 $X$ 在 $\mathbf{Q}$ 上几何不可约。}。
无论如何，更一般的思想是在给定的{\it 基底概形}上研究概形，
这个基底通常记为 $S$。我们说“令 $X$ 为 $S$ 上的概形”，
仅仅是指 $X$ 配备了一个态射 $X \to S$。在图中，我们会尽量把
{\it 结构态射} $X \to S$ 画成从 $X$ 向下指向 $S$ 的箭头。
我们往往更关心 $X$ 相对于 $S$ 的性质，而不是 $X$ 的内在几何。
例如，我们希望了解态射 $X \to S$ 的纤维、$X$ 在基变换后会发生什么，等等。

\medskip\noindent
下面引入一些惯用术语。当然，这些术语只是研究一个范畴中某个给定对象
之上对象的范畴这一思想的特例；见《范畴》，
例 \ref{categories-example-category-over-X}。

\begin{definition}
\label{schemes-definition-base-change}
令 $S$ 为概形。
\begin{enumerate}
\item 我们称 $X$ 为{\it $S$ 上的概形}，是指 $X$ 配备一个概形态射
$X \to S$。态射 $X \to S$ 有时称为{\it 结构态射}。
\item 若 $R$ 为环，我们说 $X$ 是{\it $R$ 上的概形}，
而不说 $X$ 是 $\Spec(R)$ 上的概形。
\item 所谓{\it 态射 $f : X \to Y$，即 $S$ 上概形的态射}，
是指一个概形态射，使得复合 $X \to Y \to S$
（由 $f$ 与 $Y$ 的结构态射组成）等于 $X$ 的结构态射。
\item 记 $\Mor_S(X, Y)$ 为所有从 $X$ 到 $Y$ 且位于 $S$ 上的态射所成的集合。
\item 令 $X$ 为 $S$ 上的概形，并令 $S' \to S$ 为概形态射。
$X$ 的{\it 基变换}是概形 $X_{S'} = S' \times_S X$，视为 $S'$ 上的概形。
\item 令 $f : X \to Y$ 为 $S$ 上的概形态射，并令 $S' \to S$ 为概形态射。
$f$ 的{\it 基变换}是诱导态射 $f' : X_{S'} \to Y_{S'}$
（即态射 $\text{id}_{S'} \times_{\text{id}_S} f$）。
\item 令 $R$ 为环。令 $X$ 为 $R$ 上的概形。令 $R \to R'$ 为环同态。
{\it 基变换} $X_{R'}$ 是概形 $\Spec(R') \times_{\Spec(R)} X$，
视为 $R'$ 上的概形。
\end{enumerate}
\end{definition}

\noindent
下面是一个典型结论。

\begin{lemma}
\label{schemes-lemma-base-change-immersion}
令 $S$ 为概形。令 $f : X \to Y$ 为 $S$ 上概形的浸入
（分别地，闭浸入、开浸入）。则 $f$ 的任意基变换仍为浸入
（分别地，闭浸入、开浸入）。
\end{lemma}

\begin{proof}
可以把 $f$ 沿态射 $S' \to S$ 的基变换看作下列交换图中
左上方的竖直箭头：
$$
\xymatrix{
X_{S'} \ar[r] \ar[d] & X \ar[d] \ar@/^4ex/[dd] \\
Y_{S'} \ar[r] \ar[d] & Y \ar[d] \\
S' \ar[r] & S
}
$$
此图蕴含 $X_{S'} \cong Y_{S'} \times_Y X$，
本引理由引理 \ref{schemes-lemma-fibre-product-immersion} 得出。
\end{proof}

\noindent
事实上，这类结论如此典型，以至于有一个专门术语来表述它。

\begin{definition}
\label{schemes-definition-preserved-by-base-change}
性质与基变换。
\begin{enumerate}
\item 令 $\mathcal{P}$ 为基底上概形的一种性质。我们称
$\mathcal{P}$ 在{\it 任意基变换下保持}，或简称
$\mathcal{P}$ 在{\it 基变换下保持}，如果每当 $X/S$ 具有
$\mathcal{P}$ 时，任意基变换 $X_{S'}/S'$ 都具有 $\mathcal{P}$。
\item 令 $\mathcal{P}$ 为基底上概形态射的一种性质。我们称
$\mathcal{P}$ 在{\it 任意基变换下保持}，或简称在{\it 基变换下保持}，
如果每当 $f : X \to Y$ 在 $S$ 上具有 $\mathcal{P}$ 时，任意基变换
$f' : X_{S'} \to Y_{S'}$ 在 $S'$ 上都具有 $\mathcal{P}$。
\end{enumerate}
\end{definition}

\noindent
现在可以说，“是闭浸入”这一性质在任意基变换下保持。

\begin{definition}
\label{schemes-definition-fibre}
令 $f : X \to S$ 为概形态射。令 $s \in S$ 为一点。
所谓{\it 概形论纤维 $X_s$，即 $f$ 在 $s$ 上方的纤维}，
或简称{\it $f$ 在 $s$ 上方的纤维}，是处于下列纤维积图中的概形：
$$
\xymatrix{
X_s = \Spec(\kappa(s)) \times_S X \ar[r] \ar[d] &
X \ar[d] \\
\Spec(\kappa(s)) \ar[r] &
S
}
$$
我们始终把纤维 $X_s$ 看作 $\kappa(s)$ 上的概形。
\end{definition}

\begin{lemma}
\label{schemes-lemma-fibre-topological}
令 $f : X \to S$ 为概形态射。考虑下列两图：
$$
\vcenter{
\xymatrix{
X_s \ar[r] \ar[d] & X \ar[d] \\
\Spec(\kappa(s)) \ar[r] & S
}
}
\quad\text{且}\quad
\vcenter{
\xymatrix{
\Spec(\mathcal{O}_{S, s}) \times_S X \ar[r] \ar[d] & X \ar[d] \\
\Spec(\mathcal{O}_{S, s}) \ar[r] & S
}
}
$$
在两种情形中，这些图都诱导拓扑空间的纤维方块；特别地，
上方的水平箭头是到其像的同胚。
\end{lemma}

\begin{proof}
选取仿射开集 $U \subset S$，使其包含 $s$。
下方的水平态射经由 $U$ 分解；例如见引理
\ref{schemes-lemma-morphism-from-spec-local-ring}。
因此可假设 $S$ 是仿射的。若 $X$ 也是仿射的，则结论由
《代数》，注 \ref{algebra-remark-fundamental-diagram} 得出。
一般情形可通过用仿射开集覆盖 $X$ 得出。
\end{proof}

\begin{lemma}
\label{schemes-lemma-local-ring-fibre}
令 $f : X \to S$ 为概形态射。令 $x \in X$ 的像为 $s \in S$。
把 $x$ 看作 $X_s$ 的一点（见引理 \ref{schemes-lemma-fibre-topological}），则有
$$
\mathcal{O}_{X_s, x} \cong
\mathcal{O}_{X, x}/\mathfrak m_s\mathcal{O}_{X, x} \cong
\mathcal{O}_{X, x} \otimes_{\mathcal{O}_{S, s}} \kappa(s)
$$
\end{lemma}

\begin{proof}
与引理 \ref{schemes-lemma-fibre-topological} 的证明相同，这归结为 $X$ 与 $S$
均仿射的情形。在仿射情形中，该陈述转化为《代数》，
注 \ref{algebra-remark-local-ring-fibre}。
\end{proof}







\section{拟紧态射}
\label{schemes-section-quasi-compact}

\noindent
若一个概形的底拓扑空间拟紧，则称该概形{\it 拟紧}。
相对版本的概念定义如下。

\begin{definition}
\label{schemes-definition-quasi-compact}
若一个概形态射的底拓扑空间映射拟紧，则称该态射{\it 拟紧}；
见《拓扑》，定义 \ref{topology-definition-quasi-compact}。
\end{definition}

\begin{lemma}
\label{schemes-lemma-quasi-compact-affine}
令 $f : X \to S$ 为概形态射。下列条件等价：
\begin{enumerate}
\item $f : X \to S$ 拟紧；
\item 每个仿射开集的逆像拟紧；
\item 存在某个仿射开覆盖 $S = \bigcup_{i \in I} U_i$，
使得 $f^{-1}(U_i)$ 对所有 $i$ 都拟紧。
\end{enumerate}
\end{lemma}

\begin{proof}
设给定 (3) 中的覆盖 $S = \bigcup_{i \in I} U_i$。
首先，令 $U \subset S$ 为任意仿射开集。对任意 $u \in U$，
可以找到指标 $i(u) \in I$，使得 $u \in U_{i(u)}$。
由于标准开集构成 $U_{i(u)}$ 的拓扑基，可以找到
$W_u \subset U \cap U_{i(u)}$，它是 $U_{i(u)}$ 中的标准开集。
由拟紧性，可以找到有限多个点 $u_1, \ldots, u_n \in U$，使得
$U = \bigcup_{j = 1}^n W_{u_j}$。对每个 $j$，把
$f^{-1}U_{i(u_j)} = \bigcup_{k \in K_j} V_{jk}$ 写成有限多个仿射开集的并。
由于 $W_{u_j} \subset U_{i(u_j)}$ 是标准开集，可知
$f^{-1}(W_{u_j}) \cap V_{jk}$ 是 $V_{jk}$ 的标准开集；见《代数》，
引理 \ref{algebra-lemma-spec-functorial}。因此
$f^{-1}(W_{u_j}) \cap V_{jk}$ 是仿射的，从而
$f^{-1}(W_{u_j})$ 是有限多个仿射概形的并。这证明任意仿射开集的
逆像都是有限多个仿射开集的并。

\medskip\noindent
接着，假设每个仿射开集的逆像都是有限多个仿射开集的并。
令 $K \subset S$ 为任意拟紧开集。由于 $S$ 具有由仿射开集组成的拓扑基，
可知 $K$ 是有限多个仿射开集的并。因此 $K$ 的逆像是有限多个仿射开集的并，
所以 $f$ 拟紧。

\medskip\noindent
最后，假设 $f$ 拟紧。此时，前一段的论证说明任意仿射开集的逆像
都是有限多个仿射开集的并。
\end{proof}

\begin{lemma}
\label{schemes-lemma-quasi-compact-preserved-base-change}
“拟紧”是基底上概形态射的一种性质，并且在任意基变换下保持。
\end{lemma}

\begin{proof}
略。
\end{proof}

\begin{lemma}
\label{schemes-lemma-composition-quasi-compact}
拟紧态射的复合拟紧。
\end{lemma}

\begin{proof}
这由定义及《拓扑》，
引理 \ref{topology-lemma-composition-quasi-compact} 得出。
\end{proof}

\begin{lemma}
\label{schemes-lemma-closed-immersion-quasi-compact}
闭浸入拟紧。
\end{lemma}

\begin{proof}
这由定义及《拓扑》，
引理 \ref{topology-lemma-closed-in-quasi-compact} 得出。
\end{proof}

\begin{example}
\label{schemes-example-open-immersion-not-quasi-compact}
一般而言，开浸入不拟紧。标准例子是开子空间
$U \subset X$，其中 $X = \Spec(k[x_1, x_2, x_3, \ldots])$，
$U$ 是 $X \setminus \{0\}$，而 $0$ 是 $X$ 中对应于极大理想
$(x_1, x_2, x_3, \ldots)$ 的点。
\end{example}

\begin{lemma}
\label{schemes-lemma-image-quasi-compact-closed}
令 $f : X \to S$ 为拟紧概形态射。下列条件等价：
\begin{enumerate}
\item $f(X) \subset S$ 是闭集；
\item $f(X) \subset S$ 在特化下稳定。
\end{enumerate}
\end{lemma}

\begin{proof}
由《拓扑》，引理 \ref{topology-lemma-open-closed-specialization}，
有 (1) $\Rightarrow$ (2)。假设 (2)。令 $U \subset S$ 为仿射开集。
只须证明 $f(X) \cap U$ 闭。由于 $U \cap f(X)$ 在 $U$ 中的特化下稳定，
问题归结为 $S$ 仿射的情形。因为 $f$ 拟紧，可推出
$X = f^{-1}(S)$ 拟紧，这是由于 $S$ 仿射。因此可以写成
$X = \bigcup_{i = 1}^n U_i$，其中 $U_i \subset X$ 是仿射开集。
设 $S = \Spec(R)$，并令 $U_i = \Spec(A_i)$，其中
$R$-代数为 $A_i$。于是
$f(X) = \Im(\Spec(A_1 \times \ldots \times A_n)
\to \Spec(R))$。因此本引理由《代数》，
引理 \ref{algebra-lemma-image-stable-specialization-closed} 得出。
\end{proof}

\begin{lemma}
\label{schemes-lemma-quasi-compact-closed}
令 $f : X \to S$ 为拟紧概形态射。则 $f$ 是闭映射，当且仅当特化可沿
$f$ 提升；见《拓扑》，定义 \ref{topology-definition-lift-specializations}。
\end{lemma}

\begin{proof}
由《拓扑》，引理 \ref{topology-lemma-closed-open-map-specialization}，
若 $f$ 是闭映射，则特化可沿 $f$ 提升。反过来，假设特化可沿 $f$ 提升。
令 $Z \subset X$ 为闭子集。赋予 $Z$ 约化诱导概形结构，便可把它看作概形；
见定义 \ref{schemes-definition-reduced-induced-scheme}。由于 $Z \subset X$ 闭，
$f$ 在 $Z$ 上的限制仍拟紧。此外，特化也可沿 $Z \to S$ 提升；
见《拓扑》，引理 \ref{topology-lemma-lift-specialization-composition}。
因此，只须证明 $f(X)$ 闭，前提是特化可沿 $f$ 提升。
特别地，$f(X)$ 在特化下稳定；见《拓扑》，
引理 \ref{topology-lemma-lift-specializations-images}。于是由引理
\ref{schemes-lemma-image-quasi-compact-closed}，$f(X)$ 闭。
\end{proof}










\section{万有闭性的赋值判据}
\label{schemes-section-valuative-criterion-universal-closedness}

\noindent
《拓扑》第 \ref{topology-section-proper} 节讨论了固有映射：
可以把它看作所有纤维都拟紧的拓扑空间闭映射，
也可以看作所有基变换都是闭映射的映射。下面给出代数几何中的对应概念。

\begin{definition}
\label{schemes-definition-universally-closed}
若概形态射 $f : X \to S$ 的每个基变换
$f' : X_{S'} \to S'$ 都是闭映射，则称它{\it 万有闭}。
\end{definition}

\noindent
事实上，形容词“万有”经常这样使用。换言之，给定态射的一种性质
$\mathcal{P}$，我们说“$X \to S$ 万有地具有 $\mathcal{P}$”，
当且仅当每个基变换 $X_{S'} \to S'$ 都具有 $\mathcal{P}$。

\medskip\noindent
关于万有闭态射性质的更详细讨论，请见《态射》，
第 \ref{morphisms-section-proper} 节。本节只讨论万有闭态射与满足
赋值判据存在性部分的态射之间的关系。

\begin{lemma}
\label{schemes-lemma-specializations-lift}
令 $f : X \to S$ 为概形态射。
\begin{enumerate}
\item 若 $f$ 万有闭，则特化可沿 $f$ 的任意基变换提升；
见《拓扑》，定义 \ref{topology-definition-lift-specializations}。
\item 若 $f$ 拟紧，且特化可沿 $f$ 的任意基变换提升，则 $f$ 万有闭。
\end{enumerate}
\end{lemma}

\begin{proof}
(1) 直接由《拓扑》，引理
\ref{topology-lemma-closed-open-map-specialization} 得出。
(2) 由引理 \ref{schemes-lemma-quasi-compact-closed} 与
\ref{schemes-lemma-quasi-compact-preserved-base-change} 得出。
\end{proof}

\begin{definition}
\label{schemes-definition-valuative-criterion}
令 $f : X \to S$ 为概形态射。称 $f${\it 满足赋值判据的存在性部分}，
如果给定任意实线交换图
$$
\xymatrix{
\Spec(K) \ar[r] \ar[d] & X \ar[d] \\
\Spec(A) \ar[r] \ar@{-->}[ru] & S
}
$$
其中 $A$ 是以 $K$ 为分式域的赋值环，虚线箭头总存在。称 $f$
{\it 满足赋值判据的唯一性部分}，如果对任意上述图，虚线箭头
至多有一个（当然不要求存在）。
\end{definition}

\noindent
{\it 赋值环}是其分式域中支配关系下极大的局部整环；见《代数》，
定义 \ref{algebra-definition-valuation-ring}。因此，赋值环的谱有唯一泛点
$\eta$ 和唯一闭点 $0$，当然还有特化 $\eta \leadsto 0$。
赋值环的重要性在于：任意概形中点的任何特化，都是某个赋值环之谱的
某个态射下 $\eta \leadsto 0$ 的像。精确结论如下。

\begin{lemma}
\label{schemes-lemma-points-specialize}
\begin{slogan}
特化由赋值环见证。
\end{slogan}
令 $S$ 为概形。令 $s' \leadsto s$ 为 $S$ 中点的特化。则：
\begin{enumerate}
\item 存在赋值环 $A$ 和态射 $f : \Spec(A) \to S$，使得泛点
$\eta$（属于 $\Spec(A)$）映到 $s'$，而特殊点映到 $s$；
\item 给定域扩张 $K/\kappa(s')$，可以安排使扩张
$\kappa(\eta)/\kappa(s')$（由 $f$ 诱导）同构于给定扩张。
\end{enumerate}
\end{lemma}

\begin{proof}
令 $s' \leadsto s$ 为 $S$ 中的特化，并令 $K/\kappa(s')$ 为域扩张。
由引理 \ref{schemes-lemma-specialize-points} 及引理
\ref{schemes-lemma-characterize-points} 后面的讨论，得到环映射
$\mathcal{O}_{S, s} \to \kappa(s') \to K$。
令 $A \subset K$ 为任意赋值环，其分式域为 $K$，并且它支配
$\mathcal{O}_{S, s} \to K$ 的像；见《代数》，引理
\ref{algebra-lemma-dominate}。环同态 $\mathcal{O}_{S, s} \to A$
诱导态射 $f : \Spec(A) \to S$；见引理
\ref{schemes-lemma-morphism-from-spec-local-ring}。按照构造，此态射具有所有所需性质。
\end{proof}

\begin{lemma}
\label{schemes-lemma-lift-specializations-valuative}
令 $f : X \to S$ 为概形态射。下列条件等价：
\begin{enumerate}
\item 特化可沿 $f$ 的任意基变换提升；
\item 态射 $f$ 满足赋值判据的存在性部分。
\end{enumerate}
\end{lemma}

\begin{proof}
假设 (1) 成立。给定定义 \ref{schemes-definition-valuative-criterion} 中的实线图。
为了找到虚线箭头，可以把 $X \to S$ 替换为
$X_{\Spec(A)} \to \Spec(A)$，因为该假设在基变换下保持。
因此可假设 $S = \Spec(A)$。令 $x' \in X$ 为
$\Spec(K) \to X$ 的像，于是 $\kappa(x') \subset K$；
见引理 \ref{schemes-lemma-characterize-points}。
由假设，存在特化 $x' \leadsto x$，它发生在 $X$ 中，使得 $x$ 映到
$S = \Spec(A)$ 的闭点。得到局部环同态
$A \to \mathcal{O}_{X, x}$ 以及环同态
$\mathcal{O}_{X, x} \to \kappa(x')$；见引理
\ref{schemes-lemma-specialize-points} 及引理 \ref{schemes-lemma-characterize-points}
后面的讨论。复合
$A \to \mathcal{O}_{X, x} \to \kappa(x') \to K$ 是给定的单射
$A \to K$。由于 $A \to \mathcal{O}_{X, x}$ 是局部同态，
$\mathcal{O}_{X, x} \to K$ 的像支配 $A$，因而等于 $A$；
见《代数》，定义 \ref{algebra-definition-valuation-ring}。
因此得到环同态 $\mathcal{O}_{X, x} \to A$，进而得到态射
$\Spec(A) \to X$（见引理 \ref{schemes-lemma-morphism-from-spec-local-ring}
及其后的讨论）。这证明了 (2)。

\medskip\noindent
反过来，假设 (2) 成立。考察下列交换图，立即可知赋值判据的存在性部分
对任意基变换 $X_{S'} \to S'$ 都成立；这里该基变换来自 $f$：
$$
\xymatrix{
\Spec(K) \ar[r] \ar[d] & X_{S'} \ar[r] \ar[d] & X \ar[d] \\
\Spec(A) \ar[r] \ar@{-->}[ru] \ar@{-->}[rru] & S' \ar[r] & S
}
$$
具体而言，由纤维积的定义，较为水平的虚线箭头会给出另一个虚线箭头。
因此显然只须证明特化可沿 $f$ 提升。
令 $s' \leadsto s$ 为 $S$ 中的特化，并令 $x' \in X$ 为位于 $s'$ 上方的点。
将引理 \ref{schemes-lemma-points-specialize} 应用于 $s' \leadsto s$ 以及域扩张
$K = \kappa(x')/\kappa(s')$，得到交换图
$$
\xymatrix{
\Spec(K) \ar[rr] \ar[d] & & X \ar[d] \\
\Spec(A) \ar[r] \ar@{-->}[rru] &
\Spec(\mathcal{O}_{S, s}) \ar[r] & S
}
$$
由条件 (2) 得到虚线箭头。闭点的像 $x$，这里闭点属于 $\Spec(A)$，
而该像位于 $X$ 中，便是所求解。也就是说，$x$ 是 $x'$ 的特化，
并映到 $s$。
\end{proof}

\begin{proposition}[万有闭性的赋值判据]
\label{schemes-proposition-characterize-universally-closed}
令 $f$ 为拟紧概形态射。则 $f$ 万有闭，当且仅当 $f$
满足赋值判据的存在性部分。
\end{proposition}

\begin{proof}
这是上面的引理 \ref{schemes-lemma-specializations-lift} 与
\ref{schemes-lemma-lift-specializations-valuative} 的形式推论。
\end{proof}

\begin{example}
\label{schemes-example-projective-line-universally-closed}
令 $k$ 为域。考虑结构态射
$p : \mathbf{P}^1_k \to \Spec(k)$，它来自 $k$ 上的射影直线；
见例 \ref{schemes-example-projective-line}。
利用上面的赋值判据证明 $p$ 万有闭。
按照构造，$\mathbf{P}^1_k$ 由两个仿射开集覆盖，因而 $p$ 拟紧。
给定交换图
$$
\xymatrix{
\Spec(K) \ar[r]_\xi \ar[d] & \mathbf{P}^1_k \ar[d] \\
\Spec(A) \ar[r]^\varphi & \Spec(k)
}
$$
其中 $A$ 是赋值环，$K$ 是其分式域。回忆 $\mathbf{P}^1_k$ 由
$\Spec(k[x])$ 与 $\Spec(k[y])$ 粘合而得，其中把 $D(x)$ 与 $D(y)$
经由 $x = y^{-1}$ 粘合（或更对称地写成 $xy = 1$）。
为了证明存在态射 $\Spec(A) \to \mathbf{P}^1_k$ 斜向嵌入上图，
由对称性可假设 $\xi$ 映入开集 $\Spec(k[x])$。这给出下列环的交换图：
$$
\xymatrix{
K & k[x] \ar[l]^{\xi^\sharp} \\
A \ar[u] & k \ar[u] \ar[l]_{\varphi^\sharp}
}
$$
由《代数》，引理 \ref{algebra-lemma-valuation-ring-x-or-x-inverse}，
可知或者 $\xi^\sharp(x) \in A$，或者 $\xi^\sharp(x)^{-1} \in A$。
在第一种情形，得到环同态
$$
k[x] \to A,
\ \lambda \mapsto \varphi^\sharp(\lambda),
\  x \mapsto \xi^\sharp(x)
$$
它嵌入上面的环图中，结论得证。在第二种情形，得到环同态
$$
k[y] \to A,
\ \lambda \mapsto \varphi^\sharp(\lambda),
\ y \mapsto \xi^\sharp(x)^{-1}.
$$
这给出态射
$\Spec(A) \to \Spec(k[y]) \to \mathbf{P}^1_k$，
它斜向嵌入本例开头的交换图中（验证略）。
\end{example}





























\section{分离公理}
\label{schemes-section-separation-axioms}

\noindent
拓扑空间 $X$ 是豪斯多夫的，当且仅当对角线
$\Delta \subset X \times X$ 是闭子集。代数几何中的类似做法是：
给定概形 $X$，它位于基底概形 $S$ 上，考虑对角态射
$$
\Delta_{X/S} : X \longrightarrow X \times_S X.
$$
这是满足
$\text{pr}_1 \circ \Delta_{X/S} = \text{id}_X$ 以及
$\text{pr}_2 \circ \Delta_{X/S} = \text{id}_X$ 的唯一概形态射
（它在任何具有纤维积的范畴中都存在）。

\begin{lemma}
\label{schemes-lemma-diagonal-affines-closed}
仿射概形之间态射的对角态射是闭浸入。
\end{lemma}

\begin{proof}
与态射 $\Spec(S) \to \Spec(R)$ 相关联的对角态射，是由环同态
$S \otimes_R S \to S$、$a \otimes b \mapsto ab$ 所对应的谱上态射。
此映射显然满射，故有 $S \cong S \otimes_R S/J$，其中
$J \subset S \otimes_R S$ 是某个理想。因此 $\Delta$ 是闭浸入；
见例 \ref{schemes-example-closed-immersion-affines}。
\end{proof}

\begin{lemma}
\label{schemes-lemma-diagonal-immersion}
\begin{slogan}
相对概形的对角态射是浸入。
\end{slogan}
令 $X$ 为 $S$ 上的概形。对角态射 $\Delta_{X/S}$ 是浸入。
\end{lemma}

\begin{proof}
回忆，若 $V \subset X$ 是仿射开集并映入仿射开集 $U \subset S$，
则 $V \times_U V$ 是 $X \times_S X$ 中的仿射开集；见引理
\ref{schemes-lemma-fibre-product-affines} 与 \ref{schemes-lemma-open-fibre-product}。
考虑开子概形 $W$（属于 $X \times_S X$），它是这些仿射开集
$V \times_U V$ 的并。由引理 \ref{schemes-lemma-closed-local-target}，
只须证明每个态射
$\Delta_{X/S}^{-1}(V \times_U V) \to V \times_U V$ 都是闭浸入。
由于 $V = \Delta_{X/S}^{-1}(V \times_U V)$，这正是在检验
$\Delta_{V/U}$ 是闭浸入，而这就是引理
\ref{schemes-lemma-diagonal-affines-closed}。
\end{proof}

\begin{definition}
\label{schemes-definition-separated}
令 $f : X \to S$ 为概形态射。
\begin{enumerate}
\item 称 $f${\it 分离}，如果对角态射 $\Delta_{X/S}$ 是闭浸入。
\item 称 $f${\it 拟分离}，如果对角态射 $\Delta_{X/S}$ 拟紧。
\item 称概形 $Y${\it 分离}，如果态射 $Y \to \Spec(\mathbf{Z})$ 分离。
\item 称概形 $Y${\it 拟分离}，如果态射 $Y \to \Spec(\mathbf{Z})$ 拟分离。
\end{enumerate}
\end{definition}

\noindent
由引理 \ref{schemes-lemma-diagonal-immersion} 与 \ref{schemes-lemma-immersion-when-closed}，
可知 $\Delta_{X/S}$ 是闭浸入，当且仅当
$\Delta_{X/S}(X) \subset X \times_S X$ 是闭子集。此外，由引理
\ref{schemes-lemma-closed-immersion-quasi-compact}，分离态射必拟分离。
引入拟分离态射的原因是：研究代数簇时（尤其在模空间、代数栈等问题中）
自然会出现非分离态射，但它们多数仍然拟分离。

\begin{example}
\label{schemes-example-not-quasi-separated}
下面给出一个非拟分离态射的例子。设
$X = X_1 \cup X_2 \to S = \Spec(k)$，其中
$X_1 = X_2 = \Spec(k[t_1, t_2, t_3, \ldots])$，
并沿 $\{0\} = \{(t_1, t_2, t_3, \ldots)\}$ 的补集粘合
（粘合法同例 \ref{schemes-example-affine-space-zero-doubled}）。
此时，仿射概形 $X_1 \times_S X_2$ 在 $\Delta_{X/S}$ 下的逆像是概形
$\Spec(k[t_1, t_2, t_3, \ldots]) \setminus \{0\}$，它不拟紧。
\end{example}

\begin{lemma}
\label{schemes-lemma-where-are-they-equal}
令 $X$、$Y$ 为 $S$ 上的概形。令 $a, b : X \to Y$ 为 $S$ 上的概形态射。
存在最大的局部闭子概形 $Z \subset X$，使得 $a|_Z = b|_Z$。
事实上，$Z$ 是 $(a, b)$ 的等化子。此外，若 $Y$ 在 $S$ 上分离，
则 $Z$ 是闭子概形。
\end{lemma}

\begin{proof}
由范畴论理由，$(a, b)$ 的等化子是下图中的纤维积 $Z$：
$$
\xymatrix{
Z = Y \times_{(Y \times_S Y)} X \ar[r] \ar[d] &
 X \ar[d]^{(a , b)} \\
Y \ar[r]^-{\Delta_{Y/S}} & Y \times_S Y
}
$$
因此本引理由引理 \ref{schemes-lemma-base-change-immersion}、
\ref{schemes-lemma-diagonal-immersion} 及定义 \ref{schemes-definition-separated} 得出。
\end{proof}

\begin{lemma}
\label{schemes-lemma-characterize-quasi-separated}
令 $f : X \to S$ 为概形态射。下列条件等价：
\begin{enumerate}
\item 态射 $f$ 拟分离。
\item 对每一对仿射开集 $U, V \subset X$，若它们映入 $S$ 的同一个
仿射开集，则交集 $U \cap V$ 是 $X$ 的有限多个仿射开集的并。
\item 存在仿射开覆盖 $S = \bigcup_{i \in I} U_i$，并且对每个 $i$
存在仿射开覆盖 $f^{-1}U_i = \bigcup_{j \in I_i} V_j$，使得对每个 $i$
以及每一对 $j, j' \in I_i$，交集 $V_j \cap V_{j'}$ 都是 $X$ 的
有限多个仿射开集的并。
\end{enumerate}
\end{lemma}

\begin{proof}
先证明 (3) 蕴含 (1)。由引理 \ref{schemes-lemma-affine-covering-fibre-product}，覆盖
$X \times_S X = \bigcup_i \bigcup_{j, j'} V_j \times_{U_i} V_{j'}$
是 $X \times_S X$ 的仿射开覆盖。此外，
$\Delta_{X/S}^{-1}(V_j \times_{U_i} V_{j'}) = V_j \cap V_{j'}$。
故该蕴含由引理 \ref{schemes-lemma-quasi-compact-affine} 得出。

\medskip\noindent
(1) $\Rightarrow$ (2) 来自如下事实：在 (2) 的假设下，纤维积
$U \times_S V$ 是 $X \times_S X$ 的仿射开集。
(2) $\Rightarrow$ (3) 是平凡的。
\end{proof}

\begin{lemma}
\label{schemes-lemma-characterize-separated}
令 $f : X \to S$ 为概形态射。
\begin{enumerate}
\item 若 $f$ 分离，则对每一对仿射开集 $(U, V)$（属于 $X$），
只要它们映入 $S$ 的同一个仿射开集，就有：
\begin{enumerate}
\item 交集 $U \cap V$ 是仿射的；
\item 环同态
$\mathcal{O}_X(U) \otimes_{\mathbf{Z}} \mathcal{O}_X(V)
\to \mathcal{O}_X(U \cap V)$
满射。
\end{enumerate}
\item 若任意一对点 $x_1, x_2 \in X$ 位于同一点 $s \in S$ 上方，
都分别包含在仿射开集 $x_1 \in U$、$x_2 \in V$ 中，且这些开集映入
$S$ 的同一个仿射开集并满足 (a)、(b)，则 $f$ 分离。
\end{enumerate}
\end{lemma}

\begin{proof}
假设 $f$ 分离。设 $(U, V)$ 是 (1) 中的一对开集。
令 $W = \Spec(R)$ 为 $S$ 的仿射开集，并同时包含 $f(U)$ 与 $f(V)$。
写成 $U = \Spec(A)$、$V = \Spec(B)$，其中用到 $R$-代数 $A$ 与 $B$。
由引理 \ref{schemes-lemma-open-fibre-product}，
$U \times_S V = U \times_W V = \Spec(A \otimes_R B)$
是 $X \times_S X$ 的仿射开集。因而由引理
\ref{schemes-lemma-closed-subspace-scheme} 可知，
$\Delta^{-1}(U \times_S V) \to U \times_S V$ 可等同于
$\Spec((A \otimes_R B)/J) \to \Spec(A \otimes_R B)$，
其中 $J \subset A \otimes_R B$ 是某个理想。
因此 $U \cap V = \Delta^{-1}(U \times_S V)$ 是仿射的。
断言 (1)(b) 成立，因为
$A \otimes_{\mathbf{Z}} B \to (A \otimes_R B)/J$ 满射。

\medskip\noindent
假设 (2) 中表述的假设成立。显然，仿射开集 $U \times_S V$，
对应于 (2) 中各对 $(U, V)$，构成 $X \times_S X$ 的仿射开覆盖
（例如见引理 \ref{schemes-lemma-affine-covering-fibre-product}）。
因此，只须证明每个态射
$U \cap V = \Delta_{X/S}^{-1}(U \times_S V) \to U \times_S V$
都是闭浸入；见引理 \ref{schemes-lemma-closed-local-target}。
由假设 (a)，有 $U \cap V = \Spec(C)$，其中 $C$ 为某个环。
选取仿射开集 $W = \Spec(R)$（属于 $S$），使 $U$ 与 $V$ 都映入其中，
并写成 $U = \Spec(A)$、$V = \Spec(B)$，则假设 (b) 是说复合
$$
A \otimes_{\mathbf{Z}} B \to A \otimes_R B \to C
$$
满射。因此 $A \otimes_R B \to C$ 满射，从而
$\Spec(C) \to \Spec(A \otimes_R B)$ 是闭浸入。
\end{proof}

\begin{example}
\label{schemes-example-projective-line-separated}
令 $k$ 为域。考虑结构态射
$p : \mathbf{P}^1_k \to \Spec(k)$，它来自 $k$ 上的射影直线；
见例 \ref{schemes-example-projective-line}。利用上面的引理证明 $p$ 分离。
按照构造，$\mathbf{P}^1_k$ 由两个仿射开集
$U = \Spec(k[x])$ 与 $V = \Spec(k[y])$ 覆盖，其交集为
$U \cap V = \Spec(k[x, y]/(xy - 1))$（采用显然的记号）。
因此，只须检验引理 \ref{schemes-lemma-characterize-separated} 的条件 (2)(a)、(2)(b)
对仿射开集对 $(U, U)$、$(U, V)$、$(V, U)$ 及 $(V, V)$ 成立。
对 $(U, U)$ 与 $(V, V)$ 而言这是平凡的。对 $(U, V)$ 而言，
这相当于证明 $U \cap V$ 仿射（确实如此），并且环同态
$$
k[x] \otimes_{\mathbf{Z}} k[y] \longrightarrow k[x, y]/(xy - 1)
$$
满射。这是显然的，因为右端的任意元素都可写成关于 $x$ 的多项式
与关于 $y$ 的多项式之和。
\end{example}

\begin{lemma}
\label{schemes-lemma-fibre-product-after-map}
令 $f : X \to T$ 与 $g : Y \to T$ 为具有相同目标的概形态射，
并令 $h : T \to S$ 为概形态射。则诱导态射
$i : X \times_T Y \to X \times_S Y$ 是浸入。若 $T \to S$ 分离，
则 $i$ 是闭浸入；若 $T \to S$ 拟分离，则 $i$ 拟紧。
\end{lemma}

\begin{proof}
由一般范畴论，下图
$$
\xymatrix{
X \times_T Y \ar[r] \ar[d] & X \times_S Y \ar[d] \\
T \ar[r]^{\Delta_{T/S}} \ar[r] & T \times_S T
}
$$
是纤维积图。本引理由引理 \ref{schemes-lemma-diagonal-immersion}、
\ref{schemes-lemma-fibre-product-immersion} 与
\ref{schemes-lemma-quasi-compact-preserved-base-change} 得出。
\end{proof}

\begin{lemma}
\label{schemes-lemma-semi-diagonal}
令 $g : X \to Y$ 为 $S$ 上的概形态射。态射
$i : X \to X \times_S Y$ 是浸入。若 $Y$ 在 $S$ 上分离，
则它是闭浸入；若 $Y$ 在 $S$ 上拟分离，则它拟紧。
\end{lemma}

\begin{proof}
这是把引理 \ref{schemes-lemma-fibre-product-after-map} 应用于态射
$X = X \times_Y Y \to X \times_S Y$ 所得的特殊情形。
\end{proof}

\begin{lemma}
\label{schemes-lemma-section-immersion}
令 $f : X \to S$ 为概形态射。令 $s : S \to X$ 为 $f$ 的截面
（用公式写为 $f \circ s = \text{id}_S$）。则 $s$ 是浸入。
若 $f$ 分离，则 $s$ 是闭浸入；若 $f$ 拟分离，则 $s$ 拟紧。
\end{lemma}

\begin{proof}
这是把引理 \ref{schemes-lemma-semi-diagonal} 应用于 $g =s$ 所得的特殊情形；
此时态射为 $i = s : S \to S \times_S X$。
\end{proof}

\begin{lemma}
\label{schemes-lemma-separated-permanence}
恒久性质。
\begin{enumerate}
\item 分离态射的复合分离。
\item 拟分离态射的复合拟分离。
\item 分离态射的基变换分离。
\item 拟分离态射的基变换拟分离。
\item 分离态射的（纤维）积分离。
\item 拟分离态射的（纤维）积拟分离。
\end{enumerate}
\end{lemma}

\begin{proof}
令 $X \to Y \to Z$ 为态射。假设 $X \to Y$ 与 $Y \to Z$ 分离。复合
$$
X \to X \times_Y X \to X \times_Z X
$$
是闭浸入，因为第一个态射依假设是闭浸入，第二个态射则由引理
\ref{schemes-lemma-fibre-product-after-map} 是闭浸入。同一论证也适用于“拟分离”
（使用相同的引理）。

\medskip\noindent
令 $f : X \to Y$ 为基底 $S$ 上的概形态射。令 $S' \to S$ 为概形态射。
令 $f' : X_{S'} \to Y_{S'}$ 为 $f$ 的基变换。则 $f'$ 的对角态射为
$$
\Delta_{f'} :
X_{S'} = S' \times_S X
\longrightarrow
X_{S'} \times_{Y_{S'}} X_{S'} = S' \times _S (X \times_Y X)
$$
容易看出，这是 $\Delta_f$ 的基变换。因此，(3)、(4) 由如下事实得出：
闭浸入与拟紧态射在任意基变换下保持（引理
\ref{schemes-lemma-fibre-product-immersion} 与
\ref{schemes-lemma-quasi-compact-preserved-base-change}）。

\medskip\noindent
若 $f : X \to Y$ 与 $g : U \to V$ 是基底 $S$ 上的概形态射，则
$f \times g$ 是 $X \times_S U \to X \times_S V$（$g$ 的基变换）与
$X \times_S V \to Y \times_S V$（$f$ 的基变换）的复合。
因此 (5)、(6) 由 (1)--(4) 得出。
\end{proof}

\begin{lemma}
\label{schemes-lemma-compose-after-separated}
\begin{slogan}
分离态射与拟分离态射满足消去律。
\end{slogan}
令 $f : X \to Y$ 与 $g : Y \to Z$ 为概形态射。
若 $g \circ f$ 分离，则 $f$ 也分离；
若 $g \circ f$ 拟分离，则 $f$ 也拟分离。
\end{lemma}

\begin{proof}
假设 $g \circ f$ 分离。考虑其对角态射的分解
$X \to X \times_Y X \to X \times_Z X$；这里所说的对角态射属于
$g \circ f$。由引理 \ref{schemes-lemma-fibre-product-after-map}，后一个态射是浸入。
依假设，$X$ 在 $X \times_Z X$ 中的像闭，因而它在
$X \times_Y X$ 中也闭。因此由引理 \ref{schemes-lemma-immersion-when-closed}，
$X \to X \times_Y X$ 是闭浸入。

\medskip\noindent
假设 $g \circ f$ 拟分离。令 $V \subset Y$ 为映入 $Z$ 的某个仿射开集
的仿射开集。令 $U_1, U_2 \subset X$ 为映入 $V$ 的仿射开集。于是
$U_1 \cap U_2$ 是有限多个仿射开集的并，因为 $U_1, U_2$ 映入
$Z$ 的同一个仿射开集。又因为可以覆盖 $Y$，所用仿射开集形如 $V$，
本引理由引理
\ref{schemes-lemma-characterize-quasi-separated} 得出。
\end{proof}

\begin{lemma}
\label{schemes-lemma-quasi-compact-permanence}
令 $f : X \to Y$ 与 $g : Y \to Z$ 为概形态射。
若 $g \circ f$ 拟紧且 $g$ 拟分离，则 $f$ 拟紧。
\end{lemma}

\begin{proof}
这是因为 $f$ 等于复合
$(1, f) : X \to X \times_Z Y \to Y$。第一个映射由引理
\ref{schemes-lemma-section-immersion} 拟紧，因为它是拟分离态射
$X \times_Z Y \to X$ 的截面（该态射是 $g$ 的基变换；
见引理 \ref{schemes-lemma-separated-permanence}）。第二个映射也拟紧，
因为它是 $g \circ f$ 的基变换；见引理
\ref{schemes-lemma-quasi-compact-preserved-base-change}。最后，拟紧态射的复合拟紧；
见引理 \ref{schemes-lemma-composition-quasi-compact}。
\end{proof}

\begin{lemma}
\label{schemes-lemma-affine-separated}
仿射概形分离。从仿射概形到另一个概形的态射分离。
\end{lemma}

\begin{proof}
令 $U = \Spec(A)$ 为仿射概形。由引理
\ref{schemes-lemma-diagonal-affines-closed}，$U \to \Spec(\mathbf{Z})$ 的对角闭。
因此由定义 \ref{schemes-definition-separated}，$U$ 分离。
若 $U \to X$ 是概形态射，则可把引理
\ref{schemes-lemma-compose-after-separated} 应用于态射
$U \to X \to \Spec(\mathbf{Z})$，从而得出 $U \to X$ 分离。
\end{proof}

\noindent
读者或许会问：为了断定一对仿射开集的交仍仿射，
是否真的必须只考虑其像包含在某个仿射开集中的那些对？往往并不需要！

\begin{lemma}
\label{schemes-lemma-curiosity}
令 $f : X \to S$ 为态射。假设 $f$ 分离且 $S$ 是分离概形。
设 $U \subset X$ 与 $V \subset X$ 为仿射开集。则 $U \cap V$ 仿射
（并且是 $U \times V$ 的闭子概形）。
\end{lemma}

\begin{proof}
在此情形，由引理 \ref{schemes-lemma-separated-permanence}，$X$ 分离。
因此 $U \cap V$ 仿射；这是把引理 \ref{schemes-lemma-characterize-separated}
应用于态射 $X \to \Spec(\mathbf{Z})$ 所得的结论。
\end{proof}

\noindent
另一方面，下例说明不能指望仿射概形的像包含在某个仿射开集中。

\begin{example}
\label{schemes-example-image-affine-projective}
考虑非仿射概形 $U = \Spec(k[x, y]) \setminus \{(x, y)\}$，
见例 \ref{schemes-example-not-affine}。另一方面，考虑概形
$$
\mathbf{GL}_{2, k} = \Spec(k[a, b, c, d, 1/ad - bc]).
$$
存在态射 $\mathbf{GL}_{2, k} \to U$，它对应于环同态
$x \mapsto a$、$y \mapsto b$。容易看出该态射满射，故其像不包含在
$U$ 的任何仿射开集中。事实上，仿射概形 $\mathbf{GL}_{2, k}$ 也满射到
$\mathbf{P}^1_k$，而 $\mathbf{P}^1_k$ 甚至不能浸入{\it 任何}仿射概形。
\end{example}

\begin{remark}
\label{schemes-remark-quasi-compact-and-quasi-separated}
令 $P$ 为以下 $4$ 种性质之一：“拟分离”、“分离”、
“拟紧且拟分离”或“拟紧且分离”。则下列陈述成立：
\begin{enumerate}
\item 任意仿射概形都具有 $P$。
\item 给定概形态射 $f : X \to Y$，若 $Y$ 具有 $P$，则
$f$ 具有 $P$ 当且仅当 $X$ 具有 $P$。
\item 若 $X \to Y$ 与 $Z \to Y$ 为概形态射，并且
$X$、$Y$、$Z$ 都具有 $P$，则 $X \times_Y Z$ 具有 $P$。
\end{enumerate}
陈述 (1) 显然。令 $f : X \to Y$ 如 (2)。若 $f$ 具有 $P$，则 $X$ 也具有，
这是把引理 \ref{schemes-lemma-separated-permanence} 与
\ref{schemes-lemma-composition-quasi-compact} 应用于
$X \to Y \to \Spec(\mathbf{Z})$ 所得。若 $X$ 具有 $P$，则 $f$ 也具有，
这是把引理 \ref{schemes-lemma-compose-after-separated} 与
\ref{schemes-lemma-quasi-compact-permanence} 应用于
$X \to Y \to \Spec(\mathbf{Z})$ 所得。
令 $X \to Y$ 与 $Z \to Y$ 如 (3)。投影 $X \times_Y Z \to Z$ 具有 $P$，
因为它是箭头 $X \to Y$ 的基变换，而该箭头由 (2) 具有 $P$；
见引理 \ref{schemes-lemma-separated-permanence} 与
\ref{schemes-lemma-quasi-compact-preserved-base-change}。因此由 (2)，
$X \times_Y Z$ 具有 $P$。
\end{remark}





\section{分离性的赋值判据}
\label{schemes-section-valuative-separatedness}

\begin{lemma}
\label{schemes-lemma-separated-implies-valuative}
令 $f : X \to S$ 为概形态射。若 $f$ 分离，则 $f$ 满足
赋值判据的唯一性部分。
\end{lemma}

\begin{proof}
给定定义 \ref{schemes-definition-valuative-criterion} 中的一个图。假设有两个态射
$a, b : \Spec(A) \to X$ 嵌入该图。令 $Z \subset \Spec(A)$ 为 $a$ 与 $b$
的等化子。由引理 \ref{schemes-lemma-where-are-they-equal}，它是 $\Spec(A)$ 的
闭子概形。依假设，它包含 $\Spec(A)$ 的泛点。由于 $A$ 是整环，
这蕴含 $Z = \Spec(A)$。故如所需，$a = b$。
\end{proof}

\begin{lemma}[分离性的赋值判据]
\label{schemes-lemma-valuative-criterion-separatedness}
\begin{reference}
\cite[II Proposition 7.2.3]{EGA}
\end{reference}
令 $f : X \to S$ 为态射。假设：
\begin{enumerate}
\item 态射 $f$ 拟分离；
\item 态射 $f$ 满足赋值判据的唯一性部分。
\end{enumerate}
则 $f$ 分离。
\end{lemma}

\begin{proof}
由假设 (1)、命题 \ref{schemes-proposition-characterize-universally-closed}，
以及引理 \ref{schemes-lemma-diagonal-immersion}、\ref{schemes-lemma-immersion-when-closed}，
可知只须证明态射
$\Delta_{X/S} : X \to X \times_S X$ 满足赋值判据的存在性部分。
给定实线交换图
$$
\xymatrix{
\Spec(K) \ar[r] \ar[d] & X \ar[d] \\
\Spec(A) \ar[r] \ar@{-->}[ru] & X \times_S X
}
$$
右下箭头对应于一对态射 $a, b : \Spec(A) \to X$，它们位于 $S$ 上。
由 (2) 可知 $a = b$。因此取 $a$ 为虚线箭头即可。
\end{proof}





\section{单态射}
\label{schemes-section-monomorphisms}

\begin{definition}
\label{schemes-definition-monomorphism}
若一个概形态射是概形范畴中的单态射，则称它为{\it 单态射}；
见《范畴》，定义 \ref{categories-definition-mono-epi}。
\end{definition}

\begin{lemma}
\label{schemes-lemma-monomorphism}
\begin{slogan}
概形态射是单态射，当且仅当其对角态射是同构。
\end{slogan}
令 $j : X \to Y$ 为概形态射。则 $j$ 是单态射，当且仅当对角态射
$\Delta_{X/Y} : X \to X \times_Y X$ 是同构。
\end{lemma}

\begin{proof}
这在任何具有纤维积的范畴中都成立。
\end{proof}

\begin{lemma}
\label{schemes-lemma-monomorphism-separated}
概形单态射分离。
\end{lemma}

\begin{proof}
这是因为同构是闭浸入，并应用上面的引理 \ref{schemes-lemma-monomorphism}。
\end{proof}

\begin{lemma}
\label{schemes-lemma-composition-monomorphism}
单态射的复合是单态射。
\end{lemma}

\begin{proof}
这在任何范畴中都成立。
\end{proof}

\begin{lemma}
\label{schemes-lemma-base-change-monomorphism}
单态射的基变换是单态射。
\end{lemma}

\begin{proof}
这在任何具有纤维积的范畴中都成立。
\end{proof}

\begin{lemma}
\label{schemes-lemma-injective-points}
令 $j : X \to Y$ 为概形态射。若 $j$ 在点上单射，则 $j$ 分离。
\end{lemma}

\begin{proof}
令 $z$ 为 $X \times_Y X$ 的点。则 $x = \text{pr}_1(z)$ 与
$\text{pr}_2(z)$ 相同，因为 $j$ 把这两点映到同一点 $y$（属于 $Y$）。
于是可选取仿射开邻域 $V \subset Y$，它包含 $y$；以及仿射开邻域
$U \subset X$，它包含 $x$，并使得 $j(U) \subset V$。于是
$z \in U \times_V U \subset X \times_Y X$。
因此 $X \times_Y X$ 是仿射开集 $U \times_V U$ 的并。由于
$\Delta_{X/Y}^{-1}(U \times_V U) = U$，并且
$U \to U \times_V U$ 是闭浸入，可断定 $\Delta_{X/Y}$ 是闭浸入
（见引理 \ref{schemes-lemma-diagonal-immersion} 证明中的论证）。
\end{proof}

\begin{lemma}
\label{schemes-lemma-injective-points-surjective-stalks}
令 $j : X \to Y$ 为概形态射。若
\begin{enumerate}
\item $j$ 在点上单射；
\item 对任意 $x \in X$，环同态
$j^\sharp_x : \mathcal{O}_{Y, j(x)} \to \mathcal{O}_{X, x}$
满射，
\end{enumerate}
则 $j$ 是单态射。
\end{lemma}

\begin{proof}
令 $a, b : Z \to X$ 为两个概形态射，满足
$j \circ a  = j \circ b$。由 (1)，作为底拓扑空间映射有 $a = b$。
对任意 $z \in Z$，有
$a^\sharp_z \circ j^\sharp_{a(z)} = b^\sharp_z \circ j^\sharp_{b(z)}$，
这是从 $\mathcal{O}_{Y, j(a(z))} \to \mathcal{O}_{Z, z}$ 的映射。
映射 $j^\sharp_x$ 的满射性迫使
$a^\sharp_z = b^\sharp_z$，$\forall z \in Z$。这蕴含
$a^\sharp = b^\sharp$，因此如所需，作为概形态射有 $a = b$。
\end{proof}

\begin{lemma}
\label{schemes-lemma-immersions-monomorphisms}
概形浸入是单态射。特别地，任意浸入都分离。
\end{lemma}

\begin{proof}
可以通过检验引理 \ref{schemes-lemma-injective-points-surjective-stalks}
的判据适用来证明。或许更优雅的是，利用引理
\ref{schemes-lemma-restrict-map-to-opens} 与
\ref{schemes-lemma-characterize-closed-subspace}：它们蕴含开浸入和闭浸入
都是单态射，因而任意浸入（它是这类态射的复合）都是单态射。
\end{proof}

\begin{lemma}
\label{schemes-lemma-subscheme-of-separated-scheme}
令 $f : X \to S$ 为分离态射。任意局部闭子概形
$Z \subset X$ 在 $S$ 上分离。
\end{lemma}

\begin{proof}
这由引理 \ref{schemes-lemma-immersions-monomorphisms} 以及分离态射的复合分离
这一事实得出（引理 \ref{schemes-lemma-separated-permanence}）。
\end{proof}

\begin{example}
\label{schemes-example-Q-over-Z}
态射 $\Spec(\mathbf{Q}) \to \Spec(\mathbf{Z})$ 是单态射，
因为 $\mathbf{Q} \otimes_{\mathbf{Z}} \mathbf{Q} = \mathbf{Q}$。
更一般地，对任意概形 $S$ 及任意点 $s \in S$，典范态射
$$
\Spec(\mathcal{O}_{S, s}) \longrightarrow S
$$
是单态射。
\end{example}

\begin{lemma}
\label{schemes-lemma-mono-towards-spec-field}
令 $k_1, \ldots, k_n$ 为域。对任意概形单态射
$X \to \Spec(k_1 \times \ldots \times k_n)$，存在子集
$I \subset \{1, \ldots, n\}$，使得
$X \cong \Spec(\prod_{i \in I} k_i)$；这里把它们视为
$\Spec(k_1 \times \ldots \times k_n)$ 上的概形。
更一般地，若 $X = \coprod_{i \in I} \Spec(k_i)$ 是域的谱的不交并，
且 $Y \to X$ 是单态射，则存在子集 $J \subset I$，使得
$Y = \coprod_{i \in J} \Spec(k_i)$。
\end{lemma}

\begin{proof}
先归结为 $n = 1$（或 $\# I = 1$）的情形；方法是取开闭子概形
$\Spec(k_i)$ 的逆像。此时 $X$ 只有一个点，因而是仿射的。
相应的代数问题如下：若 $k \to R$ 为代数同态，并且
$R \otimes_k R \cong R$，则 $R \cong k$ 或 $R = 0$。
这由维数理由成立。另见《代数》，
引理 \ref{algebra-lemma-epimorphism-over-field}。
\end{proof}










\section{拟凝聚模的函子性}
\label{schemes-section-quasi-coherent}

\noindent
令 $X$ 为概形。记 $\QCoh(\mathcal{O}_X)$ 为拟凝聚
$\mathcal{O}_X$-模所成的范畴，定义见《模》，
定义 \ref{modules-definition-quasi-coherent}。我们已在第
\ref{schemes-section-quasi-coherent-affine} 节看到，范畴 $\QCoh(\mathcal{O}_X)$
在 $X$ 仿射时具有许多良好性质。由于拟凝聚性在 $X$ 上是局部的，
这些性质也由任意概形 $X$ 上的拟凝聚层范畴继承。现列举如下。
\begin{enumerate}
\item $\mathcal{O}_X$-模层 $\mathcal{F}$ 拟凝聚，当且仅当
$\mathcal{F}$ 在每个仿射开集 $U = \Spec(R)$ 上的限制都形如
$\widetilde M$，其中它来自某个 $R$-模 $M$。
\item $\mathcal{O}_X$-模层 $\mathcal{F}$ 拟凝聚，当且仅当
$\mathcal{F}$ 在某个仿射开覆盖的每个成员上的限制拟凝聚。
\item 拟凝聚层的任意直和拟凝聚。
\item 拟凝聚层的任意余极限拟凝聚。
\item
\label{schemes-item-kernel}
拟凝聚层态射的核与余核拟凝聚。
\item 给定 $\mathcal{O}_X$-模的短正合列
$0 \to \mathcal{F}_1 \to \mathcal{F}_2 \to \mathcal{F}_3 \to 0$
，若其中两个拟凝聚，则第三个也拟凝聚。
\item 给定概形态射 $f : Y \to X$，拟凝聚 $\mathcal{O}_X$-模的拉回是
拟凝聚 $\mathcal{O}_Y$-模。见《模》，引理
\ref{modules-lemma-pullback-quasi-coherent}。
\item 给定两个拟凝聚 $\mathcal{O}_X$-模，其张量积拟凝聚；见《模》，
引理 \ref{modules-lemma-tensor-product-permanence}。
\item 给定拟凝聚 $\mathcal{O}_X$-模 $\mathcal{F}$，其上的张量代数、
对称代数与外代数 $\mathcal{F}$ 都拟凝聚；见《模》，引理
\ref{modules-lemma-whole-tensor-algebra-permanence}。
\item 给定两个拟凝聚 $\mathcal{O}_X$-模 $\mathcal{F}$、$\mathcal{G}$，
并且 $\mathcal{F}$ 是有限表示的，则内部 Hom
$\SheafHom_{\mathcal{O}_X}(\mathcal{F}, \mathcal{G})$
拟凝聚；见《模》，引理
\ref{modules-lemma-internal-hom-locally-kernel-direct-sum} 及上面的
(\ref{schemes-item-kernel})。
\end{enumerate}
另一方面，拟凝聚模的推前一般并不拟凝聚。下面给出它确实拟凝聚的一种情形。

\begin{lemma}
\label{schemes-lemma-push-forward-quasi-coherent}
令 $f : X \to S$ 为概形态射。若 $f$ 拟紧且拟分离，则
$f_*$ 把拟凝聚 $\mathcal{O}_X$-模变为拟凝聚 $\mathcal{O}_S$-模。
\end{lemma}

\begin{proof}
此问题在 $S$ 上是局部的，故可假设 $S$ 仿射。由于 $X$ 拟紧，可以写成
$X = \bigcup_{i = 1}^n U_i$，其中每个 $U_i$ 都是仿射开集。
由于 $f$ 拟分离，可以把
$U_i \cap U_j = \bigcup_{k = 1}^{n_{ij}} U_{ijk}$ 写成某些仿射开集
$U_{ijk}$ 的并；见引理 \ref{schemes-lemma-characterize-quasi-separated}。
记 $f_i : U_i \to S$ 与 $f_{ijk} : U_{ijk} \to S$ 为 $f$ 的限制。
对任意开集 $V$（属于 $S$）以及任意层 $\mathcal{F}$（定义在 $X$ 上），有
\begin{eqnarray*}
f_*\mathcal{F}(V) & = & \mathcal{F}(f^{-1}V) \\
& = &
\Ker\left(
\bigoplus\nolimits_i \mathcal{F}(f^{-1}V \cap U_i)
\to
\bigoplus\nolimits_{i, j, k} \mathcal{F}(f^{-1}V \cap U_{ijk})\right) \\
& = &
\Ker\left(
\bigoplus\nolimits_i f_{i, *}(\mathcal{F}|_{U_i})(V)
\to
\bigoplus\nolimits_{i, j, k} f_{ijk, *}(\mathcal{F}|_{U_{ijk}})(V)\right) \\
& = &
\Ker\left(
\bigoplus\nolimits_i f_{i, *}(\mathcal{F}|_{U_i})
\to
\bigoplus\nolimits_{i, j, k} f_{ijk, *}(\mathcal{F}|_{U_{ijk}})\right)(V)
\end{eqnarray*}
换言之，有层的正合列
$$
0 \to f_*\mathcal{F}
\to \bigoplus f_{i, *}\mathcal{F}_i
\to \bigoplus f_{ijk, *}\mathcal{F}_{ijk}
$$
其中 $\mathcal{F}_i, \mathcal{F}_{ijk}$ 表示 $\mathcal{F}$ 在相应开集上的限制。
若 $\mathcal{F}$ 是拟凝聚 $\mathcal{O}_X$-模，则 $\mathcal{F}_i$ 是拟凝聚
$\mathcal{O}_{U_i}$-模，而 $\mathcal{F}_{ijk}$ 是拟凝聚
$\mathcal{O}_{U_{ijk}}$-模。因此由引理 \ref{schemes-lemma-widetilde-pullback}，
正合列的第二、第三项都是拟凝聚 $\mathcal{O}_S$-模。故
$f_*\mathcal{F}$ 是拟凝聚 $\mathcal{O}_S$-模。
\end{proof}

\noindent
利用这一点，可以如下刻画概形的（闭）浸入。

\begin{lemma}
\label{schemes-lemma-characterize-closed-immersions}
令 $f : X \to Y$ 为概形态射。假设：
\begin{enumerate}
\item $f$ 诱导 $X$ 到 $Y$ 的某个闭子集上的同胚；
\item $f^\sharp : \mathcal{O}_Y \to f_*\mathcal{O}_X$
满射。
\end{enumerate}
则 $f$ 是概形的闭浸入。
\end{lemma}

\begin{proof}
假设 (1)、(2)。由 (1)，态射 $f$ 拟紧
（见《拓扑》，引理 \ref{topology-lemma-closed-in-quasi-compact}）。
条件 (1)、(2) 蕴含引理 \ref{schemes-lemma-injective-points-surjective-stalks}
的条件 (1)、(2)。因此 $f : X \to Y$ 是单态射。特别地，$f$ 分离；
见引理 \ref{schemes-lemma-monomorphism-separated}。于是可以应用上面的引理
\ref{schemes-lemma-push-forward-quasi-coherent}，得出 $f_*\mathcal{O}_X$ 是拟凝聚
$\mathcal{O}_Y$-模。因此，$\mathcal{O}_Y \to f_*\mathcal{O}_X$ 的核拟凝聚；
见引理 \ref{schemes-lemma-extension-quasi-coherent}。由于拟凝聚层局部由截面生成
（见《模》，定义 \ref{modules-definition-quasi-coherent}），这蕴含 $f$
是闭浸入；见定义 \ref{schemes-definition-closed-immersion-locally-ringed-spaces}。
\end{proof}

\noindent
可以用这个引理证明下面的引理。

\begin{lemma}
\label{schemes-lemma-composition-immersion}
概形浸入的复合是浸入；概形闭浸入的复合是闭浸入；
概形开浸入的复合是开浸入。
\end{lemma}

\begin{proof}
开浸入的情形显然，因为开子空间的开子空间仍是开子空间。

\medskip\noindent
设 $a : Z \to Y$ 与 $b : Y \to X$ 为概形闭浸入。我们验证
$c = b \circ a$ 也是闭浸入。假设蕴含 $a$ 与 $b$ 都是到闭子集上的同胚，
因而 $c = b \circ a$ 也是到闭子集上的同胚。此外，映射
$\mathcal{O}_X \to c_*\mathcal{O}_Z$ 满射，因为它分解为满射
$\mathcal{O}_X \to b_*\mathcal{O}_Y$ 与
$b_*\mathcal{O}_Y \to b_*a_*\mathcal{O}_Z$ 的复合（后者满射是因为
$b_*$ 正合；见《模》，引理 \ref{modules-lemma-i-star-exact}）。
因此由上面的引理 \ref{schemes-lemma-characterize-closed-immersions}，$c$ 是闭浸入。

\medskip\noindent
最后处理一般浸入。设 $a : Z \to Y$ 与 $b : Y \to X$ 为概形浸入。
这意味着存在开子概形 $V \subset Y$ 与 $U \subset X$，使得
$a(Z) \subset V$、$b(Y) \subset U$，且 $a : Z \to V$ 与
$b : Y \to U$ 都是闭浸入。由于 $Y$ 上的拓扑由 $U$ 上的拓扑诱导，
可以找到开集 $U' \subset U$，使得 $V = b^{-1}(U')$。于是
$Z \to V = b^{-1}(U') \to U'$ 是闭浸入的复合，因而是闭浸入。
这证明 $Z \to X$ 是浸入，结论得证。
\end{proof}
% Generated mechanically from the frozen constructions.tex target.
% Do not edit this generated file; edit the bound locale source instead.

% OK, start here.
%

\setcounter{chapter}{26}
\chapter{概形的构造}



\phantomsection
\label{constructions-section-phantom}


\section{引言}
\label{constructions-section-introduction}

\noindent
本章介绍由已知概形构造新概形的一些方法。
一部基本参考文献是 \cite{EGA}。





\section{相对粘合}
\label{constructions-section-relative-glueing}

\noindent
当我们试图构造概形 $X$（它位于 $S$ 上），并且已经知道如何构造 $X$ 在 $S$ 的
各仿射开集上的限制时，下面的引理正好适用。事实上，这个结果完全一般，
在（局部）环层空间的框架下也成立；不过我们的证明采用概形的语言。

\begin{lemma}
\label{constructions-lemma-relative-glueing}
设 $S$ 为概形，$\mathcal{B}$ 为 $S$ 的拓扑的一组基。
假设给定以下数据：
\begin{enumerate}
\item 对每个 $U \in \mathcal{B}$，给定概形 $f_U : X_U \to U$，
它位于 $U$ 上；
\item 对 $U, V \in \mathcal{B}$，若 $V \subset U$，给定态射
$\rho^U_V : X_V \to X_U$，它位于 $U$ 上。
\end{enumerate}
假设：
\begin{enumerate}
\item[(a)] 每个 $\rho^U_V$ 都诱导概形同构
$X_V \to f_U^{-1}(V)$，它位于 $V$ 上；
\item[(b)] 每当 $W, V, U \in \mathcal{B}$ 且
$W \subset V \subset U$ 时，都有 $\rho^U_W = \rho^U_V \circ \rho ^V_W$。
\end{enumerate}
则存在概形态射 $f : X \to S$ 以及同构
$i_U : f^{-1}(U) \to X_U$（$U \in \mathcal{B}$），使得对
$V, U \in \mathcal{B}$ 且 $V \subset U$，复合
$$
\xymatrix{
X_V \ar[r]^{i_V^{-1}} &
f^{-1}(V) \ar[rr]^{inclusion} & &
f^{-1}(U) \ar[r]^{i_U} &
X_U
}
$$
就是态射 $\rho^U_V$。此外，$X$ 在 $S$ 上唯一同构的意义下唯一。
\end{lemma}

\begin{proof}
证明将使用《概形》引理 \ref{schemes-lemma-glue-functors}。
首先定义从概形范畴到集合范畴的反变函子 $F$。具体地，对概形 $T$，令
$$
F(T) =
\left\{
\begin{matrix}
(g, \{h_U\}_{U \in \mathcal{B}}),
\ g : T \to S, \ h_U : g^{-1}(U) \to X_U, \\
f_U \circ h_U = g|_{g^{-1}(U)},
\ h_U|_{g^{-1}(V)} = \rho^U_V \circ h_V
\ \forall\ V, U \in \mathcal{B}, V \subset U
\end{matrix}
\right\}.
$$
限制映射 $F(T) \to F(T')$ 由态射 $T' \to T$ 给出，并且仅由复合得到。
对任意 $W \in \mathcal{B}$，考虑子函子 $F_W \subset F$；它由满足
条件的那些系统 $(g, \{h_U\})$ 组成；此条件是 $g(T) \subset W$。

\medskip\noindent
先证明 $F$ 满足 Zariski 拓扑的层性质。设 $T$ 为概形，
$T = \bigcup V_i$ 为开覆盖，并且元素 $\xi_i \in F(V_i)$ 满足
$\xi_i|_{V_i \cap V_j} = \xi_j|_{V_i \cap V_j}$。
记 $\xi_i = (g_i, \{h_{i, U}\})$。立即可见，态射 $g_i$ 粘合成唯一的
整体态射 $g : T \to S$。此外，显然有
$g^{-1}(U) = \bigcup g_i^{-1}(U)$。因此态射
$h_{i, U} : g_i^{-1}(U) \to X_U$ 粘合成唯一态射
$h_U : g^{-1}(U) \to X_U$。容易验证，系统 $(g, \{h_U\})$ 是 $F(T)$
的元素。因此 $F$ 满足 Zariski 拓扑的层性质。

\medskip\noindent
接着验证每个 $F_W$（$W \in \mathcal{B}$）都可表。也就是说，断言函子变换
$$
F_W \longrightarrow \Mor(-, X_W), \ (g, \{h_U\}) \longmapsto h_W
$$
是同构。为此，设 $T$ 为概形，$\alpha : T \to X_W$ 为态射，并令
$g = f_W \circ \alpha$。对任意 $U \in \mathcal{B}$，若 $U \subset W$，
可以把 $h_U : g^{-1}(U) \to X_U$ 定义为复合
$(\rho^W_U)^{-1} \circ \alpha|_{g^{-1}(U)}$。这是可行的，因为像
$\alpha(g^{-1}(U))$ 包含于 $f_W^{-1}(U)$，并且引理的条件 (a) 成立。
显然，有 $f_U \circ h_U = g|_{g^{-1}(U)}$；这里的 $U$ 如上。
此外，若还有 $V \in \mathcal{B}$ 且 $V \subset U \subset W$，
则由引理的性质 (b)，有
$\rho^U_V \circ h_V = h_U|_{g^{-1}(V)}$。还需要定义 $h_U$（对任意元素
$U \in \mathcal{B}$）。由于 $\mathcal{B}$ 是 $S$ 的拓扑的一组基，
可以找到开覆盖 $U \cap W = \bigcup U_i$，其中 $U_i \in \mathcal{B}$。
因为 $g$ 的像位于 $W$ 中，所以
$g^{-1}(U) = g^{-1}(U \cap W) = \bigcup g^{-1}(U_i)$.
考虑态射
$h_i = \rho^U_{U_i} \circ h_{U_i} : g^{-1}(U_i) \to X_U$。
利用引理的条件 (b)，不难证明
$h_i|_{g^{-1}(U_i) \cap g^{-1}(U_j)} = h_j|_{g^{-1}(U_i) \cap g^{-1}(U_j)}$.
因此这些态射粘合成所需的态射 $h_U : g^{-1}(U) \to X_U$。
系统 $(g, \{h_U\})$ 是 $F_W(T)$ 的元素，并且在上面显示的箭头下映到
$\alpha$；略去这一点的简单验证。

\medskip\noindent
接着验证每个 $F_W \subset F$ 都由开浸入可表。这由定义立即可见。

\medskip\noindent
最后需要验证族 $(F_W)_{W \in \mathcal{B}}$ 覆盖 $F$。
这由构造以及 $\mathcal{B}$ 是 $S$ 的拓扑的一组基立即可见。

\medskip\noindent
设 $X$ 为表示函子 $F$ 的概形，并设
$(f, \{i_U\}) \in F(X)$ 为一个“泛族”。由于每个 $F_W$ 都由 $X_W$
表示（通过上面显示的函子态射），可见 $i_W : f^{-1}(W) \to X_W$
正是所需的同构。引理得证。
\end{proof}

\begin{lemma}
\label{constructions-lemma-relative-glueing-sheaves}
设 $S$ 为概形，$\mathcal{B}$ 为 $S$ 的拓扑的一组基。
假设给定以下数据：
\begin{enumerate}
\item 对每个 $U \in \mathcal{B}$，给定概形 $f_U : X_U \to U$，
它位于 $U$ 上；
\item 对每个 $U \in \mathcal{B}$，给定拟凝聚层 $\mathcal{F}_U$，
它定义在 $X_U$ 上；
\item 对每一对 $U, V \in \mathcal{B}$，若 $V \subset U$，给定态射
$\rho^U_V : X_V \to X_U$；
\item 对每一对 $U, V \in \mathcal{B}$，若 $V \subset U$，给定态射
$\theta^U_V : (\rho^U_V)^*\mathcal{F}_U \to \mathcal{F}_V$。
\end{enumerate}
假设：
\begin{enumerate}
\item[(a)] 每个 $\rho^U_V$ 都诱导概形同构
$X_V \to f_U^{-1}(V)$，它位于 $V$ 上；
\item[(b)] 每个 $\theta^U_V$ 都是同构；
\item[(c)] 每当 $W, V, U \in \mathcal{B}$ 且
$W \subset V \subset U$ 时，都有 $\rho^U_W = \rho^U_V \circ \rho ^V_W$；
\item[(d)] 每当 $W, V, U \in \mathcal{B}$ 且
$W \subset V \subset U$ 时，都有
$\theta^U_W = \theta^V_W \circ (\rho^V_W)^*\theta^U_V$。
\end{enumerate}
则存在概形态射 $f : X \to S$、拟凝聚层 $\mathcal{F}$（定义在 $X$ 上），
以及同构 $i_U : f^{-1}(U) \to X_U$ 与
$\theta_U : i_U^*\mathcal{F}_U \to \mathcal{F}|_{f^{-1}(U)}$
（$U \in \mathcal{B}$），使得对 $V, U \in \mathcal{B}$ 且
$V \subset U$，复合
$$
\xymatrix{
X_V \ar[r]^{i_V^{-1}} &
f^{-1}(V) \ar[rr]^{inclusion} & &
f^{-1}(U) \ar[r]^{i_U} &
X_U
}
$$
就是态射 $\rho^U_V$，而复合
\begin{equation}
\label{constructions-equation-glue}
(\rho^U_V)^*\mathcal{F}_U
=
(i_V^{-1})^*((i_U^*\mathcal{F}_U)|_{f^{-1}(V)})
\xrightarrow{\theta_U|_{f^{-1}(V)}}
(i_V^{-1})^*(\mathcal{F}|_{f^{-1}(V)})
\xrightarrow{\theta_V^{-1}}
\mathcal{F}_V
\end{equation}
等于 $\theta^U_V$。此外，$(X, \mathcal{F})$ 在 $S$ 上唯一同构的意义下唯一。
\end{lemma}

\begin{proof}
由引理 \ref{constructions-lemma-relative-glueing}，得到概形 $X$（它位于 $S$ 上）以及同构 $i_U$。
令 $\mathcal{F}'_U = i_U^*\mathcal{F}_U$，其中 $U \in \mathcal{B}$。
这是拟凝聚 $\mathcal{O}_{f^{-1}(U)}$-模。映射
$$
\mathcal{F}'_U|_{f^{-1}(V)} =
i_U^*\mathcal{F}_U|_{f^{-1}(V)} =
i_V^*(\rho^U_V)^*\mathcal{F}_U \xrightarrow{i_V^*\theta^U_V}
i_V^*\mathcal{F}_V = \mathcal{F}'_V
$$
定义了同构
$(\theta')^U_V : \mathcal{F}'_U|_{f^{-1}(V)} \to \mathcal{F}'_V$
（其中 $V \subset U$ 是 $\mathcal{B}$ 的元素）。条件 (d) 恰好说明，
当 $W \subset V \subset U$ 是 $\mathcal{B}$ 的三个元素时，这些同构彼此相容。
由此可以得到良定义的同构
$$
\varphi_{12} :
\mathcal{F}'_{U_1}|_{f^{-1}(U_1 \cap U_2)}
\longrightarrow
\mathcal{F}'_{U_2}|_{f^{-1}(U_1 \cap U_2)}
$$
（其中 $U_1, U_2 \in \mathcal{B}$）：覆盖交集
$U_1 \cap U_2 = \bigcup V_j$，其中各 $V_j$ 都属于 $\mathcal{B}$，并取
$$
\varphi_{12}|_{V_j} =
\left((\theta')^{U_2}_{V_j}\right)^{-1}
\circ
(\theta')^{U_1}_{V_j}.
$$
略去这些映射确实粘合成 $\varphi_{12}$ 的验证，也略去层的粘合数据所需的
余循环条件的验证（见《层》第 \ref{sheaves-section-glueing-sheaves} 节）。
由《层》引理 \ref{sheaves-lemma-glue-sheaves}，得到 $\mathcal{F}$（定义在
$X$ 上）。等式 (\ref{constructions-equation-glue}) 的验证从略。
\end{proof}

\begin{remark}
\label{constructions-remark-relative-glueing-functorial}
引理 \ref{constructions-lemma-relative-glueing} 与
\ref{constructions-lemma-relative-glueing-sheaves} 中说明的构造具有函子性。
具体地，设给定引理 \ref{constructions-lemma-relative-glueing} 中的两组数据
$(f_U : X_U \to U, \rho^U_V)$ 与
$(g_U : Y_U \to U, \sigma^U_V)$。假设对每个 $U \in \mathcal{B}$，
给定态射 $h_U : X_U \to Y_U$，它位于 $U$ 上，且与限制态射
$\rho^U_V$ 和 $\sigma^U_V$ 相容。函子性意味着，这会产生概形态射
$h : X \to Y$（它位于 $S$ 上），它在各片上的限制就是态射 $h_U$；其中
$f : X \to S$ 由数据 $(f_U : X_U \to U, \rho^U_V)$ 得到，而
$g : Y \to S$ 由数据 $(g_U : Y_U \to U, \sigma^U_V)$ 得到。

\medskip\noindent
类似地，设给定两组数据
$(f_U : X_U \to U, \mathcal{F}_U, \rho^U_V, \theta^U_V)$ 与
$(g_U : Y_U \to U, \mathcal{G}_U, \sigma^U_V, \eta^U_V)$
，如引理 \ref{constructions-lemma-relative-glueing-sheaves} 所述。
假设对每个 $U \in \mathcal{B}$，给定态射
$h_U : X_U \to Y_U$，它位于 $U$ 上，且与限制态射 $\rho^U_V$ 和 $\sigma^U_V$ 相容；
还给定态射 $\tau_U : h_U^*\mathcal{G}_U \to \mathcal{F}_U$，它与映射
$\theta^U_V$ 和 $\eta^U_V$ 相容。函子性意味着，这些数据会产生
概形态射 $h : X \to Y$（它位于 $S$ 上），它在各片上的限制就是态射 $h_U$；
还产生一个态射 $h^*\mathcal{G} \to \mathcal{F}$，它在各片上的限制就是映射
$h_U$。这里，$(f : X \to S, \mathcal{F})$ 由数据
$(f_U : X_U \to U, \mathcal{F}_U, \rho^U_V, \theta^U_V)$ 得到，而
$(g : Y \to S, \mathcal{G})$ 由数据
$(g_U : Y_U \to U, \mathcal{G}_U, \sigma^U_V, \eta^U_V)$ 得到。

\medskip\noindent
略去这些验证，也不再给出相对粘合数据与相对对象之间“范畴等价”的适当表述。
\end{remark}















\section{通过粘合构造相对谱}
\label{constructions-section-spec-via-glueing}

\begin{situation}
\label{constructions-situation-relative-spec}
这里 $S$ 是概形，$\mathcal{A}$ 是拟凝聚 $\mathcal{O}_S$-代数。
这意味着 $\mathcal{A}$ 是一层 $\mathcal{O}_S$-代数，并且作为
$\mathcal{O}_S$-模是拟凝聚的。
\end{situation}

\noindent
本节概述如何构造概形态射
$$
\underline{\Spec}_S(\mathcal{A}) \longrightarrow S
$$
其方法是粘合各谱 $\Spec(\Gamma(U, \mathcal{A}))$，其中 $U$ 遍历 $S$ 的
仿射开集。先证明 $\mathcal{A}$ 在各仿射开集上的取值之谱组成一族适当的概形，
如引理 \ref{constructions-lemma-relative-glueing} 所需。

\begin{lemma}
\label{constructions-lemma-spec-inclusion}
在情形 \ref{constructions-situation-relative-spec} 中，设
$U \subset U' \subset S$ 为仿射开集。令 $A = \mathcal{A}(U)$，并令
$A' = \mathcal{A}(U')$。环映射 $A' \to A$ 诱导态射
$\Spec(A) \to \Spec(A')$，且图
$$
\xymatrix{
\Spec(A) \ar[r] \ar[d] &
\Spec(A') \ar[d] \\
U \ar[r] &
U'
}
$$
是笛卡尔图。
\end{lemma}

\begin{proof}
令 $R = \mathcal{O}_S(U)$，并令 $R' = \mathcal{O}_S(U')$。
注意，映射 $R \otimes_{R'} A' \to A$ 是同构，因为 $\mathcal{A}$ 拟凝聚
（例如见《概形》引理 \ref{schemes-lemma-widetilde-pullback}）。
结论来自《概形》引理 \ref{schemes-lemma-fibre-product-affine-schemes}
中对仿射概形纤维积的描述。
\end{proof}

\noindent
特别地，引理中的态射 $\Spec(A) \to \Spec(A')$ 是开浸入。

\begin{lemma}
\label{constructions-lemma-transitive-spec}
在情形 \ref{constructions-situation-relative-spec} 中，设
$U \subset U' \subset U'' \subset S$ 为仿射开集。令
$A = \mathcal{A}(U)$、$A' = \mathcal{A}(U')$ 以及
$A'' = \mathcal{A}(U'')$。引理 \ref{constructions-lemma-spec-inclusion} 中的态射
$\Spec(A) \to \Spec(A')$ 与
$\Spec(A') \to \Spec(A'')$ 的复合，正是引理
\ref{constructions-lemma-spec-inclusion} 中的态射 $\Spec(A) \to \Spec(A'')$。
\end{lemma}

\begin{proof}
这是因为映射 $A'' \to A$ 是 $A'' \to A'$ 与
$A' \to A$ 的复合（因为 $\mathcal{A}$ 是层）。
\end{proof}

\begin{lemma}
\label{constructions-lemma-glue-relative-spec}
在情形 \ref{constructions-situation-relative-spec} 中，存在概形态射
$$
\pi : \underline{\Spec}_S(\mathcal{A}) \longrightarrow S
$$
，它具有以下性质：
\begin{enumerate}
\item 对每个仿射开集 $U \subset S$，存在同构
$i_U : \pi^{-1}(U) \to \Spec(\mathcal{A}(U))$，它位于 $U$ 上；
\item 对仿射开集 $U \subset U' \subset S$，复合
$$
\xymatrix{
\Spec(\mathcal{A}(U)) \ar[r]^{i_U^{-1}} &
\pi^{-1}(U) \ar[rr]^{inclusion} & &
\pi^{-1}(U') \ar[r]^{i_{U'}} &
\Spec(\mathcal{A}(U'))
}
$$
就是上面引理 \ref{constructions-lemma-spec-inclusion} 中的开浸入。
\end{enumerate}
此外，$\underline{\Spec}_S(\mathcal{A})$ 在 $S$ 上唯一同构的意义下唯一。
\end{lemma}

\begin{proof}
这立即由引理 \ref{constructions-lemma-relative-glueing}、
\ref{constructions-lemma-spec-inclusion} 和
\ref{constructions-lemma-transitive-spec} 得出。
唯一性正是引理 \ref{constructions-lemma-relative-glueing} 最后一句的内容。
\end{proof}












\section{作为函子的相对谱}
\label{constructions-section-spec}

\noindent
沿用情形 \ref{constructions-situation-relative-spec} 的设定，即 $S$ 是概形，
$\mathcal{A}$ 是一层拟凝聚 $\mathcal{O}_S$-代数。

\medskip\noindent
对任意 $f : T \to S$，拉回 $f^*\mathcal{A}$ 是一层拟凝聚
$\mathcal{O}_T$-代数。我们要考虑二元组 $(f : T \to S, \varphi)$，其中
$f$ 是概形态射，而 $\varphi : f^*\mathcal{A} \to \mathcal{O}_T$ 是
$\mathcal{O}_T$-代数态射。注意，这等价于给定 $f^{-1}\mathcal{O}_S$-代数同态
$\varphi : f^{-1}\mathcal{A} \to \mathcal{O}_T$；见《层》引理
\ref{sheaves-lemma-adjointness-tensor-restrict}。这也等价于给定
$\mathcal{O}_S$-代数映射 $\varphi : \mathcal{A} \to f_*\mathcal{O}_T$；
见《层》引理 \ref{sheaves-lemma-adjoint-push-pull-modules}。
下文将不加说明地使用对 $\varphi$ 的这三种理解方式。

\medskip\noindent
给定这样的二元组 $(f : T \to S, \varphi)$ 以及态射 $a : T' \to T$，
得到另一个二元组 $(f' = f \circ a, \varphi' = a^*\varphi)$；称它为
$(f, \varphi)$ 的拉回。描述 $\varphi' = a^*\varphi$ 的一种方式，是把它写成复合
$\mathcal{A} \to f_*\mathcal{O}_T \to f'_*\mathcal{O}_{T'}$
；其中第二个映射是 $f_*a^\sharp$，这里
$a^\sharp : \mathcal{O}_T \to a_*\mathcal{O}_{T'}$。
由此定义了函子
\begin{eqnarray}
\label{constructions-equation-spec}
F : \Sch^{opp} & \longrightarrow & \textit{Sets} \\
T & \longmapsto & F(T) = \{\text{对 }(f, \varphi) \text{ 如上}\}
\nonumber
\end{eqnarray}

\begin{lemma}
\label{constructions-lemma-spec-base-change}
在情形 \ref{constructions-situation-relative-spec} 中，令 $F$ 为上面与
$(S, \mathcal{A})$ 对应的函子。设 $g : S' \to S$ 为概形态射。
令 $\mathcal{A}' = g^*\mathcal{A}$，并令 $F'$ 为上面与
$(S', \mathcal{A}')$ 对应的函子。则存在函子的典范同构
$$
F' \cong h_{S'} \times_{h_S} F
$$
。
\end{lemma}

\begin{proof}
二元组 $(f' : T \to S', \varphi' : (f')^*\mathcal{A}' \to \mathcal{O}_T)$
等价于二元组 $(f, \varphi : f^*\mathcal{A} \to \mathcal{O}_T)$ 连同 $f$ 的分解
$f = g \circ f'$。事实上，在这些记号下有
$(f')^* \mathcal{A}' = (f')^*g^*\mathcal{A} = f^*\mathcal{A}$.
引理得证。
\end{proof}

\begin{lemma}
\label{constructions-lemma-spec-affine}
在情形 \ref{constructions-situation-relative-spec} 中，令 $F$ 为上面与
$(S, \mathcal{A})$ 对应的函子。若 $S$ 仿射，则 $F$ 由仿射概形
$\Spec(\Gamma(S, \mathcal{A}))$ 表示。
\end{lemma}

\begin{proof}
写成 $S = \Spec(R)$，并令 $A = \Gamma(S, \mathcal{A})$。
于是 $A$ 是 $R$-代数，并且 $\mathcal{A} = \widetilde A$。
环映射 $R \to A$ 给出典范映射
$$
f_{univ} : \Spec(A)
\longrightarrow
S = \Spec(R).
$$
有
$f_{univ}^*\mathcal{A} =  \widetilde{A \otimes_R A}$
；见《概形》引理 \ref{schemes-lemma-widetilde-pullback}。因此存在典范映射
$$
\varphi_{univ} :
f_{univ}^*\mathcal{A} = \widetilde{A \otimes_R A}
\longrightarrow
\widetilde A = \mathcal{O}_{\Spec(A)}
$$
，它来自 $A$-模映射 $A \otimes_R A \to A$，
$a \otimes a' \mapsto aa'$。断言在此情形下，二元组
$(f_{univ}, \varphi_{univ})$ 表示 $F$。换言之，断言对任意概形 $T$，映射
$$
\Mor(T, \Spec(A)) \longrightarrow \{\text{对 } (f, \varphi)\},\quad
a \longmapsto (f_{univ} \circ a, a^*\varphi_{univ})
$$
是双射。

\medskip\noindent
下面构造逆映射。对任意二元组 $(f : T \to S, \varphi)$，得到诱导的环映射
$$
\xymatrix{
A = \Gamma(S, \mathcal{A}) \ar[r]^{f^*} &
\Gamma(T, f^*\mathcal{A}) \ar[r]^{\varphi} &
\Gamma(T, \mathcal{O}_T)
}
$$
它诱导概形态射 $T \to \Spec(A)$；见《概形》引理
\ref{schemes-lemma-morphism-into-affine}。

\medskip\noindent
略去这个映射是上面所显示映射之逆的验证。
\end{proof}

\begin{lemma}
\label{constructions-lemma-spec}
在情形 \ref{constructions-situation-relative-spec} 中，函子 $F$ 由某个概形表示。
\end{lemma}

\begin{proof}
将使用《概形》引理 \ref{schemes-lemma-glue-functors}。

\medskip\noindent
先验证 $F$ 满足 Zariski 拓扑的层性质。具体地，设 $T$ 为概形，
$T = \bigcup_{i \in I} U_i$ 为开覆盖，并且
$(f_i, \varphi_i) \in F(U_i)$ 满足
$(f_i, \varphi_i)|_{U_i \cap U_j} = (f_j, \varphi_j)|_{U_i \cap U_j}$。
这蕴含态射 $f_i : U_i \to S$ 粘合成概形态射 $f : T \to S$，且满足
$f|_{U_i} = f_i$；见《概形》第 \ref{schemes-section-glueing-schemes} 节。
因此 $f_i^*\mathcal{A} = f^*\mathcal{A}|_{U_i}$，而由假设，态射
$\varphi_i$ 在 $U_i \cap U_j$ 上相同。于是由《层》第
\ref{sheaves-section-glueing-sheaves} 节，它们粘合成
$\mathcal{O}_T$-代数态射 $f^*\mathcal{A} \to \mathcal{O}_T$。
这证明 $F$ 满足 Zariski 拓扑的层条件。

\medskip\noindent
设 $S = \bigcup_{i \in I} U_i$ 为仿射开覆盖。令 $F_i \subset F$ 为子函子，
它由那些二元组 $(f : T \to S, \varphi)$ 组成；这些二元组满足
$f(T) \subset U_i$。

\medskip\noindent
需要证明每个 $F_i$ 都可表。由引理 \ref{constructions-lemma-spec-base-change}，
$F_i$ 与以下函子等同：它对应于 $U_i$ 以及其上的拟凝聚
$\mathcal{O}_{U_i}$-代数 $\mathcal{A}|_{U_i}$。
因此结论由引理 \ref{constructions-lemma-spec-affine} 得出。

\medskip\noindent
接着证明 $F_i \subset F$ 由开浸入可表。设
$(f : T \to S, \varphi) \in F(T)$，并考虑 $V_i = f^{-1}(U_i)$。
由 $F_i$ 的定义，对给定的 $a : T' \to T$，有
$a^*(f, \varphi) \in F_i(T')$ 当且仅当 $a(T') \subset V_i$。
这正是所需的结论。

\medskip\noindent
最后需要证明族 $(F_i)_{i \in I}$ 覆盖 $F$。设
$(f : T \to S, \varphi) \in F(T)$，并考虑 $V_i = f^{-1}(U_i)$。
由于 $S = \bigcup_{i \in I} U_i$ 是 $S$ 的开覆盖，可见
$T = \bigcup_{i \in I} V_i$ 是 $T$ 的开覆盖。此外，
$(f, \varphi)|_{V_i} \in F_i(V_i)$。引理得证。
\end{proof}

\begin{lemma}
\label{constructions-lemma-glueing-gives-functor-spec}
在情形 \ref{constructions-situation-relative-spec} 中，引理
\ref{constructions-lemma-glue-relative-spec} 构造的概形
$\pi : \underline{\Spec}_S(\mathcal{A}) \to S$
与表示函子 $F$ 的概形，作为 $S$ 上的概形典范同构。
\end{lemma}

\begin{proof}
设 $X \to S$ 为表示函子 $F$ 的概形。考虑 $\mathcal{O}_S$-代数层
$\mathcal{R} = \pi_*\mathcal{O}_{\underline{\Spec}_S(\mathcal{A})}$.
按照 $\underline{\Spec}_S(\mathcal{A})$ 的构造，有同构
$\mathcal{A}(U) \to \mathcal{R}(U)$（对每个仿射开集 $U \subset S$）；
这由引理 \ref{constructions-lemma-glue-relative-spec} 的第 (1) 部分得出。
对开集 $U \subset U' \subset S$，这些同构与限制映射相容；
这由引理 \ref{constructions-lemma-glue-relative-spec} 的第 (2) 部分得出。
因此由《层》引理 \ref{sheaves-lemma-restrict-basis-equivalence-modules}，
这些同构来自一个 $\mathcal{O}_S$-代数同构
$\varphi : \mathcal{A} \to \mathcal{R}$。于是得到元素
$(\pi, \varphi) \in F(\underline{\Spec}_S(\mathcal{A}))$。
由于 $X$ 表示函子 $F$，得到相应的概形态射
$can : \underline{\Spec}_S(\mathcal{A}) \to X$，它位于 $S$ 上。

\medskip\noindent
任取仿射开集 $U \subset S$。令 $F_U \subset F$ 为 $F$ 的子函子，
它对应于二元组 $(f, \varphi)$；这些二元组定义在概形 $T$ 上，并满足
$f(T) \subset U$。显然，基变换 $X_U$ 表示 $F_U$。此外，由引理
\ref{constructions-lemma-spec-affine}，$F_U$ 由
$\Spec(\mathcal{A}(U)) = \pi^{-1}(U)$ 表示。换言之，
$X_U \cong \pi^{-1}(U)$。略去以下验证：这个等同确实由态射 $can$
到 $U$ 的基变换给出。
\end{proof}

\begin{definition}
\label{constructions-definition-relative-spec}
设 $S$ 为概形。设 $\mathcal{A}$ 为一层拟凝聚 $\mathcal{O}_S$-代数。
把 $\mathcal{A}$ 在 $S$ 上的{\it 相对谱}，或简称把 $\mathcal{A}$ 在
$S$ 上的{\it 谱}，定义为引理 \ref{constructions-lemma-glue-relative-spec} 中构造并表示
函子 $F$（\ref{constructions-equation-spec}）的概形；见引理
\ref{constructions-lemma-glueing-gives-functor-spec}。
把它记作 $\pi : \underline{\Spec}_S(\mathcal{A}) \to S$。
“泛族”是 $\mathcal{O}_S$-代数态射
$$
\mathcal{A}
\longrightarrow
\pi_*\mathcal{O}_{\underline{\Spec}_S(\mathcal{A})}
$$
\end{definition}

\noindent
下面的引理尤其说明，相对谱的构造与基变换可交换。

\begin{lemma}
\label{constructions-lemma-spec-properties}
设 $S$ 为概形。设 $\mathcal{A}$ 为一层拟凝聚
$\mathcal{O}_S$-代数。令
$\pi : \underline{\Spec}_S(\mathcal{A}) \to S$
为 $\mathcal{A}$ 在 $S$ 上的相对谱。
\begin{enumerate}
\item 对每个仿射开集 $U \subset S$，逆像 $\pi^{-1}(U)$ 是仿射的。
\item 对每个态射 $g : S' \to S$，都有
$S' \times_S \underline{\Spec}_S(\mathcal{A}) =
\underline{\Spec}_{S'}(g^*\mathcal{A})$.
\item
泛映射
$$
\mathcal{A}
\longrightarrow
\pi_*\mathcal{O}_{\underline{\Spec}_S(\mathcal{A})}
$$
是 $\mathcal{O}_S$-代数同构。
\end{enumerate}
\end{lemma}

\begin{proof}
第 (1) 部分来自相对谱的粘合描述；见引理 \ref{constructions-lemma-glue-relative-spec}。
第 (2) 部分立即由引理 \ref{constructions-lemma-spec-base-change} 得出。
第 (3) 部分在 $S$ 上是局部的，而当 $S$ 仿射时，例如由引理
\ref{constructions-lemma-spec-affine} 可知它显然成立。
\end{proof}

\begin{lemma}
\label{constructions-lemma-canonical-morphism}
设 $f : X \to S$ 为拟紧且拟分离的概形态射。由《概形》引理
\ref{schemes-lemma-push-forward-quasi-coherent}，层
$f_*\mathcal{O}_X$ 是一层拟凝聚 $\mathcal{O}_S$-代数。存在典范态射
$$
can : X \longrightarrow \underline{\Spec}_S(f_*\mathcal{O}_X)
$$
，它是 $S$ 上的概形态射。对任意仿射开集 $U \subset S$，限制
$can|_{f^{-1}(U)}$ 与典范态射
$$
f^{-1}(U) \longrightarrow \Spec(\Gamma(f^{-1}(U), \mathcal{O}_X))
$$
等同；后者来自《概形》引理 \ref{schemes-lemma-morphism-into-affine}。
\end{lemma}

\begin{proof}
按照把 $\underline{\Spec}$ 定义为表示函子 $F$ 的概形的方式，该态射来自典范映射
$\varphi : f^*f_*\mathcal{O}_X \to \mathcal{O}_X$（由推前与拉回的伴随性，
它对应于 $\text{id} : f_*\mathcal{O}_X \to f_*\mathcal{O}_X$）。
关于在 $f^{-1}(U)$ 上限制的断言，由仿射概形上相对谱的描述得出；
见引理 \ref{constructions-lemma-spec-affine}。
\end{proof}











\section{仿射 n 维空间}
\label{constructions-section-affine-n-space}

\noindent
作为相对谱的一项应用，我们如下定义基概形上的仿射 $n$ 维空间。
设基概形为 $S$。对任意整数 $n \geq 0$，可以考虑拟凝聚
$\mathcal{O}_S$-代数层 $\mathcal{O}_S[T_1, \ldots, T_n]$。
它是拟凝聚的，因为把它看作 $\mathcal{O}_S$-模层时，它就是由多重指标标号的
$\mathcal{O}_S$ 的诸副本之直和。

\begin{definition}
\label{constructions-definition-affine-n-space}
设 $S$ 为概形且 $n \geq 0$。概形
$$
\mathbf{A}^n_S =
\underline{\Spec}_S(\mathcal{O}_S[T_1, \ldots, T_n])
$$
把它看作 $S$ 上的概形，称为{\it 仿射 $n$ 维空间（在 $S$ 上）}。
若 $S = \Spec(R)$ 是仿射的，则也称它为{\it 仿射 $n$ 维空间（在 $R$ 上）}，
并记作 $\mathbf{A}^n_R$。
\end{definition}

\noindent
注意：
\par\noindent
$\mathbf{A}^n_R = \Spec(R[T_1, \ldots, T_n])$。
对概形的任意态射 $g : S' \to S$，有
\par\noindent
$g^*\mathcal{O}_S[T_1, \ldots, T_n] = \mathcal{O}_{S'}[T_1, \ldots, T_n]$
。
\par\noindent
因而 $\mathbf{A}^n_{S'} = S' \times_S \mathbf{A}^n_S$ 是相应的基变换。
所以也可用下式定义仿射 $n$ 维空间：
$$
\mathbf{A}^n_S = S \times_{\Spec(\mathbf{Z})} \mathbf{A}^n_{\mathbf{Z}}.
$$
此外，从 $S$-概形 $f : X \to S$ 到 $\mathbf{A}^n_S$ 的一个态射，
由一个 $\mathcal{O}_S$-代数同态给出：
\par\noindent
$\mathcal{O}_S[T_1, \ldots, T_n] \to f_*\mathcal{O}_X$。
这显然等价于给出各 $T_i$ 的像。换言之，从 $X$ 到 $\mathbf{A}^n_S$
的 $S$ 上态射，等价于给出 $n$ 个元素：
\par\noindent
$h_1, \ldots, h_n \in \Gamma(X, \mathcal{O}_X)$。

















\section{向量丛}
\label{constructions-section-vector-bundle}

\noindent
设 $S$ 为概形。
设 $\mathcal{E}$ 为一层拟凝聚 $\mathcal{O}_S$-模。
由《模》引理 \ref{modules-lemma-whole-tensor-algebra-permanence}，
对称代数 $\text{Sym}(\mathcal{E})$（即 $\mathcal{E}$ 在 $\mathcal{O}_S$
上的对称代数）是一层拟凝聚 $\mathcal{O}_S$-代数。
因此可以对它取相对谱。

\begin{definition}
\label{constructions-definition-vector-bundle}
设 $S$ 为概形。设 $\mathcal{E}$ 为拟凝聚
$\mathcal{O}_S$-模\footnote{读者或许会预期这里要求
$\mathcal{E}$ 是有限局部自由的。为了与 \cite[II, Definition 1.7.8]{EGA}
保持一致，我们不作此要求。}。
{\it 与 $\mathcal{E}$ 关联的向量丛}定义为
$$
\mathbf{V}(\mathcal{E}) = \underline{\Spec}_S(\text{Sym}(\mathcal{E})).
$$
\end{definition}

\noindent
与 $\mathcal{E}$ 关联的向量丛带有一些附加结构。具体而言，有分次
$$
\pi_*\mathcal{O}_{\mathbf{V}(\mathcal{E})} =
\bigoplus\nolimits_{n \geq 0} \text{Sym}^n(\mathcal{E}).
$$
它使 $\pi_*\mathcal{O}_{\mathbf{V}(\mathcal{E})}$ 成为分次
$\mathcal{O}_S$-代数。反过来，可以恢复 $\mathcal{E}$，方法是取它的次数
$1$ 部分。因此如下定义抽象向量丛。

\begin{definition}
\label{constructions-definition-abstract-vector-bundle}
设 $S$ 为概形。{\it 向量丛 $\pi : V \to S$（在 $S$ 上）}是一个
仿射概形态射，使得 $\pi_*\mathcal{O}_V$ 带有分次
$\mathcal{O}_S$-代数结构
$\pi_*\mathcal{O}_V = \bigoplus\nolimits_{n \geq 0} \mathcal{E}_n$
，满足 $\mathcal{E}_0 = \mathcal{O}_S$，并且映射
$$
\text{Sym}^n(\mathcal{E}_1) \longrightarrow \mathcal{E}_n
$$
对所有 $n \geq 0$ 都是同构。{\it $S$ 上向量丛的态射}是态射
$f : V \to V'$，使得诱导映射
$$
f^* : \pi'_*\mathcal{O}_{V'} \longrightarrow \pi_*\mathcal{O}_V
$$
与给定的分次相容。
\end{definition}

\noindent
概形 $S$ 上的向量丛有如下例子：仿射 $n$ 维空间
$\mathbf{A}^n_S$（在 $S$ 上）；见定义 \ref{constructions-definition-affine-n-space}。这是因为
$\mathcal{O}_S[T_1, \ldots, T_n] = \text{Sym}(\mathcal{O}_S^{\oplus n})$.

\begin{lemma}
\label{constructions-lemma-category-vector-bundles}
概形 $S$ 上的向量丛范畴与拟凝聚 $\mathcal{O}_S$-模范畴反等价。
\end{lemma}

\begin{proof}
略。提示：一个方向使用函子
$\underline{\Spec}_S(\text{Sym}^*_{\mathcal{O}_S}(-))$
，另一个方向使用函子
$(\pi : V \to S) \leadsto (\pi_*\mathcal{O}_V)_1$；这里下标表示取次数
$1$ 部分。
\end{proof}




\section{锥}
\label{constructions-section-cone}

\noindent
在代数几何中，锥对应于分次代数。按照我们的约定，分次环或分次代数
$A$ 带有由非负整数给出的分次
$A = \bigoplus_{d \geq 0} A_d$；见《代数》章节
\ref{algebra-section-graded}。

\begin{definition}
\label{constructions-definition-cone}
设 $S$ 为概形。设 $\mathcal{A}$ 为拟凝聚分次
$\mathcal{O}_S$-代数。假设 $\mathcal{O}_S \to \mathcal{A}_0$
是同构\footnote{通常还会假设 $\mathcal{A}$ 由 $\mathcal{A}_1$ 在
$\mathcal{O}_S$ 上生成。为了与 \cite[II, (8.3.1)]{EGA} 保持一致，
我们不作此假设。}。{\it 与 $\mathcal{A}$ 关联的锥}，也称
{\it 与 $\mathcal{A}$ 关联的仿射锥}，定义为
$$
C(\mathcal{A}) = \underline{\Spec}_S(\mathcal{A}).
$$
\end{definition}

\noindent
与一层分次 $\mathcal{O}_S$-代数关联的锥带有一些附加结构。具体而言，
得到分次
$$
\pi_*\mathcal{O}_{C(\mathcal{A})} =
\bigoplus\nolimits_{n \geq 0} \mathcal{A}_n
$$
因此可以如下定义抽象锥。

\begin{definition}
\label{constructions-definition-abstract-cone}
设 $S$ 为概形。{\it 锥 $\pi : C \to S$（在 $S$ 上）}是一个
仿射概形态射，使得 $\pi_*\mathcal{O}_C$ 带有分次 $\mathcal{O}_S$-代数结构
$\pi_*\mathcal{O}_C = \bigoplus\nolimits_{n \geq 0} \mathcal{A}_n$
，满足 $\mathcal{A}_0 = \mathcal{O}_S$。从 $\pi : C \to S$ 到
$\pi' : C' \to S$ 的{\it 锥态射}，是态射 $f : C \to C'$，
使得诱导映射
$$
f^* : \pi'_*\mathcal{O}_{C'} \longrightarrow \pi_*\mathcal{O}_C
$$
与给定的分次相容。
\end{definition}

\noindent
任意向量丛都是锥的一个例子。事实上，$S$ 上的向量丛范畴是
$S$ 上锥范畴的满子范畴。










\section{分次环的 Proj}
\label{constructions-section-proj}

\noindent
本节依照 \cite[II, Section 2]{EGA} 构造分次环的 Proj。

\medskip\noindent
设 $S$ 为分次环。考虑拓扑空间 $\text{Proj}(S)$，它与 $S$ 关联；
见《代数》章节 \ref{algebra-section-proj}。我们将在这个空间上赋予一层环
$\mathcal{O}_{\text{Proj}(S)}$，使所得二元组
$(\text{Proj}(S), \mathcal{O}_{\text{Proj}(S)})$ 成为概形。

\medskip\noindent
回顾 $\text{Proj}(S)$ 有一个由开集 $D_{+}(f)$ 构成的基，其中
$f \in S_d$ 且 $d \geq 1$；我们称这些开集为{\it 标准开集}，
见《代数》章节 \ref{algebra-section-proj}。即使不再明说，这一术语也总是蕴含
$f$ 是正次数齐次元。此外，两个标准开集之交仍是标准开集：
$D_{+}(f) \cap D_{+}(g) = D_{+}(fg)$，其中 $f, g \in S$ 是正次数齐次元。

\begin{lemma}
\label{constructions-lemma-standard-open}
设 $S$ 为分次环。设 $f \in S$ 为正次数齐次元。
\begin{enumerate}
\item 若 $g\in S$ 是正次数齐次元且
$D_{+}(g) \subset D_{+}(f)$，则
\begin{enumerate}
\item $f$ 在 $S_g$ 中可逆，且
$f^{\deg(g)}/g^{\deg(f)}$ 在 $S_{(g)}$ 中可逆；
\item 有 $g^e = af$，其中 $e \geq 1$，且 $a \in S$ 是齐次元；
\item 存在典范 $S$-代数映射 $S_f \to S_g$；
\item 存在典范 $S_0$-代数映射 $S_{(f)} \to S_{(g)}$，
它与映射 $S_f \to S_g$ 相容；
\item 映射 $S_{(f)} \to S_{(g)}$ 诱导同构
$$
(S_{(f)})_{g^{\deg(f)}/f^{\deg(g)}} \cong S_{(g)},
$$
\item 这些映射诱导拓扑空间的交换图
$$
\xymatrix{
D_{+}(g) \ar[d] &
\{\mathbf{Z}\text{-分次素理想，位于 }S_g\} \ar[l] \ar[r] \ar[d] &
\Spec(S_{(g)}) \ar[d] \\
D_{+}(f) &
\{\mathbf{Z}\text{-分次素理想，位于 }S_f\} \ar[l] \ar[r] &
\Spec(S_{(f)})
}
$$
，其中水平映射是同胚，竖直映射是开浸入；
\item 存在相互相容的典范 $S_f$-模映射和 $S_{(f)}$-模映射
$M_f \to M_g$ 与 $M_{(f)} \to M_{(g)}$，它们对任意分次 $S$-模
$M$ 都有定义；并且
\item 映射 $M_{(f)} \to M_{(g)}$ 诱导同构
$$
(M_{(f)})_{g^{\deg(f)}/f^{\deg(g)}} \cong M_{(g)}.
$$
\end{enumerate}
\item $D_{+}(f)$ 的任意开覆盖都可细化为形如
$D_{+}(f) = \bigcup_{i = 1}^n D_{+}(g_i)$ 的有限开覆盖。
\item 设 $g_1, \ldots, g_n \in S$ 为正次数齐次元。则
$D_{+}(f) \subset \bigcup D_{+}(g_i)$ 当且仅当
$g_1^{\deg(f)}/f^{\deg(g_1)}, \ldots, g_n^{\deg(f)}/f^{\deg(g_n)}$
在 $S_{(f)}$ 中生成单位理想。
\end{enumerate}
\end{lemma}

\begin{proof}
回顾 $D_{+}(g) = \Spec(S_{(g)})$，其等同由环映射
$S \to S_g \leftarrow S_{(g)}$ 给出；见《代数》引理
\ref{algebra-lemma-topology-proj}。因此
$f^{\deg(g)}/g^{\deg(f)}$ 是 $S_{(g)}$ 中不属于任何素理想的元素，
从而可逆；见《代数》引理 \ref{algebra-lemma-Zariski-topology}。
这证明了 (a)。把 $f$ 在 $S_g$ 中的逆元写成 $a/g^d$。
可将 $a$ 替换为它的次数 $d\deg(g) - \deg(f)$ 齐次部分。
这意味着 $g^d - af$ 被 $g$ 的某个幂零化，因而有 $g^e = af$，
其中某个 $a \in S$ 是次数为 $e\deg(g) - \deg(f)$ 的齐次元。
这证明了 (b)。关于 (c)，由局部化的泛性质以及 (a)，映射
$S_f \to S_g$ 存在；也可把 $b/f^n$ 映为 $a^nb/g^{ne}$ 来定义它。
这显然也诱导次数零元素子环之间的映射 $S_{(f)} \to S_{(g)}$。
类似地，可定义 $M_f \to M_g$ 和 $M_{(f)} \to M_{(g)}$，其定义是
把 $x/f^n$ 映为 $a^nx/g^{ne}$。由上述公式立刻可见，
$S_{(g)}$（相应地 $M_{(g)}$）是 $S_{(f)}$（相应地 $M_{(f)}$）
的主局部化。拓扑空间交换图中的映射对应于上述环映射。
由《代数》引理 \ref{algebra-lemma-topology-proj}，水平箭头是同胚。
竖直箭头是开浸入，因为左侧箭头是开子集的包含映射。

\medskip\noindent
开集 $D_{+}(f)$ 是拟紧的，因为它同胚于 $\Spec(S_{(f)})$；
见《代数》引理 \ref{algebra-lemma-quasi-compact}。
因此第二个断言直接由标准开集构成拓扑基这一事实得出。

\medskip\noindent
第三个断言直接由《代数》引理 \ref{algebra-lemma-Zariski-topology} 得出。
\end{proof}

\noindent
在《层》章节 \ref{sheaves-section-bases} 中，我们定义了基上的层，
并证明它与空间上的层本质上等价；见《层》引理
\ref{sheaves-lemma-extend-off-basis} 和
\ref{sheaves-lemma-extend-off-basis-structures}。此外，《层》引理
\ref{sheaves-lemma-cofinal-systems-coverings-standard-case} 证明，
对每个标准开集，只须在开覆盖的一个共尾系统上检验层条件。
由上面的引理，只须对由标准开集给出的有限覆盖作检验。

\begin{definition}
\label{constructions-definition-standard-covering}
设 $S$ 为分次环。假设 $D_{+}(f) \subset \text{Proj}(S)$ 是标准开集。
$D_{+}(f)$ 的{\it 标准开覆盖}是形如
$D_{+}(f) = \bigcup_{i = 1}^n D_{+}(g_i)$ 的覆盖，其中
$g_1, \ldots, g_n \in S$ 是正次数齐次元。
\end{definition}

\noindent
设 $S$ 为分次环。设 $M$ 为分次 $S$-模。我们将在标准开集构成的基上定义
预层 $\widetilde M$。假设 $U \subset \text{Proj}(S)$ 是标准开集。
若 $f, g \in S$ 是正次数齐次元且 $D_{+}(f) = D_{+}(g)$，
则由上面的引理 \ref{constructions-lemma-standard-open}，存在互为逆映射的典范映射
$M_{(f)} \to M_{(g)}$ 和 $M_{(g)} \to M_{(f)}$。因此可任选一个
$f$，使 $U = D_{+}(f)$，并定义
$$
\widetilde M(U) = M_{(f)}.
$$
注意，若 $D_{+}(g) \subset D_{+}(f)$，则由上面的引理
\ref{constructions-lemma-standard-open} 有典范映射
$$
\widetilde M(D_{+}(f)) = M_{(f)} \longrightarrow
M_{(g)} = \widetilde M(D_{+}(g)).
$$
显然，这在标准开集构成的基上定义了阿贝尔群预层。若 $M = S$，则
$\widetilde S$ 是该基上的环预层。对一般的 $M$，可见 $\widetilde M$
是该基上的 $\widetilde S$-模预层。

\medskip\noindent
来计算 $\widetilde M$ 在一点 $x \in \text{Proj}(S)$ 处的茎。
假设 $x$ 对应于齐次素理想 $\mathfrak p \subset S$。
由茎的定义可见
$$
\widetilde M_x
=
\colim_{f\in S_d, d > 0, f\not\in \mathfrak p} M_{(f)}
$$
。这里集合 $\{f \in S_d, d > 0, f \not \in \mathfrak p\}$ 按如下规则预序：
$f \geq f' \Leftrightarrow D_{+}(f) \subset D_{+}(f')$。
若 $f_1, f_2 \in S \setminus \mathfrak p$ 是正次数齐次元，则在此序中
$f_1f_2 \geq f_1$。在《代数》章节 \ref{algebra-section-proj} 中，
我们把 $M_{(\mathfrak p)}$ 定义为由分式 $x/f$ 构成的模，其中
$x, f$ 齐次，$\deg(x) = \deg(f)$，且 $f \not \in \mathfrak p$。
由于 $\mathfrak p \in \text{Proj}(S)$，至少存在一个正次数齐次元
$f_0 \in S$ 满足 $f_0 \not\in \mathfrak p$。于是
$x/f = f_0x/ff_0$，可见总能假设 $M_{(\mathfrak p)}$ 中元素的分母
具有正次数。由这些说明容易得出
$$
\widetilde M_x = M_{(\mathfrak p)}.
$$

\medskip\noindent
接下来检验标准开覆盖的层条件。若
$D_{+}(f) = \bigcup_{i = 1}^n D_{+}(g_i)$，则这个覆盖的层条件等价于
序列
$$
0 \to M_{(f)} \to \bigoplus M_{(g_i)} \to \bigoplus M_{(g_ig_j)}.
$$
的正合性。注意 $D_{+}(g_i) = D_{+}(fg_i)$，因此可将这个序列改写为
$$
0 \to M_{(f)} \to \bigoplus M_{(fg_i)} \to \bigoplus M_{(fg_ig_j)}.
$$
由引理 \ref{constructions-lemma-standard-open} 可见，
\par\noindent
$g_1^{\deg(f)}/f^{\deg(g_1)}, \ldots, g_n^{\deg(f)}/f^{\deg(g_n)}$
\par\noindent
在 $S_{(f)}$ 中生成单位理想，而模 $M_{(fg_i)}$、$M_{(fg_ig_j)}$
是 $S_{(f)}$-模 $M_{(f)}$ 关于这些元素及其乘积的主局部化。
因此可将《代数》引理 \ref{algebra-lemma-cover-module} 应用于模
$M_{(f)}$（它在 $S_{(f)}$ 上）以及元素
\par\noindent
$g_1^{\deg(f)}/f^{\deg(g_1)}, \ldots, g_n^{\deg(f)}/f^{\deg(g_n)}$.
\par\noindent
从而该序列正合。由上述说明可见，$\widetilde M$ 是标准开集构成的基上的层。

\medskip\noindent
于是由《层》章节 \ref{sheaves-section-bases} 的内容可知，存在唯一的环层
$\mathcal{O}_{\text{Proj}(S)}$，它在标准开集上与 $\widetilde S$ 一致。
注意，由上面对茎的计算以及《代数》引理
\ref{algebra-lemma-proj-prime}，这个环层的各个茎都是局部环。

\medskip\noindent
类似地，对任意分次 $S$-模 $M$，存在唯一的
$\mathcal{O}_{\text{Proj}(S)}$-模层 $\mathcal{F}$，它在标准开集上与
$\widetilde M$ 一致；见《层》引理
\ref{sheaves-lemma-extend-off-basis-module}。

\begin{definition}
\label{constructions-definition-structure-sheaf}
设 $S$ 为分次环。
\begin{enumerate}
\item 齐次谱的{\it 结构层 $\mathcal{O}_{\text{Proj}(S)}$（关于
$S$）}是唯一的环层
$\mathcal{O}_{\text{Proj}(S)}$
，它在标准开集构成的基上与 $\widetilde S$ 一致。
\item 局部环层空间
$(\text{Proj}(S), \mathcal{O}_{\text{Proj}(S)})$
称为 $S$ 的{\it 齐次谱}，记作 $\text{Proj}(S)$。
\item 设有一个 $\mathcal{O}_{\text{Proj}(S)}$-模层，它把
$\widetilde M$ 延拓到 $\text{Proj}(S)$ 的所有开集；称此层为一个
$\mathcal{O}_{\text{Proj}(S)}$-模层，即与 $M$ 关联的模层。
这个层也记作 $\widetilde M$。
\end{enumerate}
\end{definition}

\noindent
总结目前所得的结果。

\begin{lemma}
\label{constructions-lemma-proj-sheaves}
设 $S$ 为分次环。设 $M$ 为分次 $S$-模。设 $\widetilde M$ 为一个
$\mathcal{O}_{\text{Proj}(S)}$-模层，即与 $M$ 关联的层。
\begin{enumerate}
\item 对每个正次数齐次元 $f \in S$，有
$$
\Gamma(D_{+}(f), \mathcal{O}_{\text{Proj}(S)}) = S_{(f)}.
$$
\item 对每个正次数齐次元 $f\in S$，有
$\Gamma(D_{+}(f), \widetilde M) = M_{(f)}$，这是 $S_{(f)}$-模的等式。
\item 每当 $D_{+}(g) \subset D_{+}(f)$ 时，
$\mathcal{O}_{\text{Proj}(S)}$ 和 $\widetilde M$ 上的限制映射，
就是引理 \ref{constructions-lemma-standard-open} 中的映射
$S_{(f)} \to S_{(g)}$ 和 $M_{(f)} \to M_{(g)}$。
\item 设 $\mathfrak p$ 为 $S$ 的、不包含 $S_{+}$ 的齐次素理想，
并设 $x \in \text{Proj}(S)$ 为相应的点。则有
$\mathcal{O}_{\text{Proj}(S), x} = S_{(\mathfrak p)}$.
\item 设 $\mathfrak p$ 为 $S$ 的、不包含 $S_{+}$ 的齐次素理想，
并设 $x \in \text{Proj}(S)$ 为相应的点。则有
$(\widetilde M)_x = M_{(\mathfrak p)}$，这是 $S_{(\mathfrak p)}$-模的等式。
\item
\label{constructions-item-map}
存在典范环映射
\par\noindent
$
S_0 \longrightarrow \Gamma(\text{Proj}(S), \widetilde S)
$
\par\noindent
以及典范 $S_0$-模映射
\par\noindent
$
M_0 \longrightarrow \Gamma(\text{Proj}(S), \widetilde M)
$
\par\noindent
，它们与上面对标准开集上截面及茎的描述相容。
\end{enumerate}
此外，所有这些等同都关于分次 $S$-模 $M$ 具有函子性。特别地，函子
$M \mapsto \widetilde M$ 是从分次 $S$-模范畴到
$\mathcal{O}_{\text{Proj}(S)}$-模范畴的正合函子。
\end{lemma}

\begin{proof}
断言 (1)--(5) 由上述讨论即得。断言 (6) 成立，是因为存在与 (3) 中限制映射
相容的典范映射 $M_0 \to M_{(f)}$、$x \mapsto x/1$。
函子 $M \mapsto \widetilde M$ 的正合性来自如下事实：函子
$M \mapsto M_{(\mathfrak p)}$ 是正合的（见《代数》引理
\ref{algebra-lemma-proj-prime}），而短正合序列的正合性可在茎上检验；
见《模》引理 \ref{modules-lemma-abelian}。
\end{proof}

\begin{remark}
\label{constructions-remark-global-sections-not-isomorphism}
从 $M_0$ 到 $\widetilde M$ 的整体截面的映射一般远非同构。一个平凡例子是
取 $S = k[x, y, z]$，其中 $1 = \deg(x) = \deg(y) = \deg(z)$
（也可取任意多个变量），再取 $M = S/(x^{100}, y^{100}, z^{100})$。
容易看出 $\widetilde M = 0$，但 $M_0 = k$。
\end{remark}

\begin{lemma}
\label{constructions-lemma-standard-open-proj}
设 $S$ 为分次环。设 $f \in S$ 为正次数齐次元。假设
$D(g) \subset \Spec(S_{(f)})$ 是标准开集。则存在正次数齐次元
$h \in S$，使得在《代数》引理 \ref{algebra-lemma-topology-proj}
的同胚下，$D(g)$ 对应于 $D_{+}(h) \subset D_{+}(f)$。事实上，
可取 $h$ 使得 $g = h/f^n$，其中 $n$ 为某个整数。
\end{lemma}

\begin{proof}
写成 $g = h/f^n$，其中 $h$ 是某个正次数齐次元，且 $n \geq 1$。
若 $D_{+}(h)$ 不包含于 $D_{+}(f)$，则把 $h$ 替换为 $hf$，并把
$n$ 替换为 $n + 1$。此时 $h$ 具有所需形式，并且
$D_{+}(h) \subset D_{+}(f)$ 对应于 $D(g) \subset \Spec(S_{(f)})$。
\end{proof}

\begin{lemma}
\label{constructions-lemma-proj-scheme}
设 $S$ 为分次环。局部环层空间 $\text{Proj}(S)$ 是概形。
标准开集 $D_{+}(f)$ 是仿射开集。对任意分次 $S$-模 $M$，层
$\widetilde M$ 是拟凝聚 $\mathcal{O}_{\text{Proj}(S)}$-模层。
\end{lemma}

\begin{proof}
考虑标准开集 $D_{+}(f) \subset \text{Proj}(S)$。由引理
\ref{constructions-lemma-standard-open} 和 \ref{constructions-lemma-proj-sheaves}，有
$\Gamma(D_{+}(f), \mathcal{O}_{\text{Proj}(S)}) = S_{(f)}$，并有同胚
$\varphi : D_{+}(f) \to \Spec(S_{(f)})$。对任意标准开集
$D(g) \subset \Spec(S_{(f)})$，可按引理
\ref{constructions-lemma-standard-open-proj} 选取 $h \in S_{+}$。于是
$\varphi^{-1}(D(g)) = D_{+}(h)$，并由引理
\ref{constructions-lemma-proj-sheaves} 和 \ref{constructions-lemma-standard-open} 可得
$$
\Gamma(D_{+}(h), \mathcal{O}_{\text{Proj}(S)})
=
S_{(h)}
=
(S_{(f)})_{h^{\deg(f)}/f^{\deg(h)}}
=
(S_{(f)})_g
=
\Gamma(D(g), \mathcal{O}_{\Spec(S_{(f)})}).
$$
因此 $\mathcal{O}_{\text{Proj}(S)}$ 在 $D_{+}(f)$ 上的限制，在同胚
$\varphi$ 下恰好对应于层 $\mathcal{O}_{\Spec(S_{(f)})}$；
后者按《概形》章节 \ref{schemes-section-affine-schemes} 定义。
从而 $D_{+}(f)$ 是仿射概形，并同构于 $\Spec(S_{(f)})$，
同构由 $\varphi$ 给出；因此 $\text{Proj}(S)$ 是概形。

\medskip\noindent
完全相同的论证表明，$\widetilde M$ 是拟凝聚
$\mathcal{O}_{\text{Proj}(S)}$-模层。具体而言，上述论证给出
$$
\widetilde M|_{D_{+}(f)} \cong \varphi^*\left(\widetilde{M_{(f)}}\right)
$$
，这表明 $\widetilde M$ 是拟凝聚的。
\end{proof}

\begin{lemma}
\label{constructions-lemma-proj-separated}
设 $S$ 为分次环。概形 $\text{Proj}(S)$ 是分离的。
\end{lemma}

\begin{proof}
需要证明典范态射 $\text{Proj}(S) \to \Spec(\mathbf{Z})$ 是分离的。
使用《概形》引理 \ref{schemes-lemma-characterize-separated}。
因此只须对任意一对标准开集 $D_{+}(f)$ 和 $D_{+}(g)$ 证明：
$D_{+}(f) \cap D_{+}(g) = D_{+}(fg)$ 是仿射的（这是显然的），
并且环映射
$$
S_{(f)} \otimes_{\mathbf{Z}} S_{(g)} \longrightarrow S_{(fg)}
$$
是满射。任意元素 $s$（在 $S_{(fg)}$ 中）都形如 $s = h/(f^ng^m)$，
其中 $h \in S$ 是次数为 $n\deg(f) + m\deg(g)$ 的齐次元。
可以把 $h$ 乘以适当的单项式 $f^ig^j$，并假设
$n = n' \deg(g)$ 且 $m = m' \deg(f)$。于是可将 $s$ 改写为
$s = h/f^{(n' + m')\deg(g)} \cdot f^{m'\deg(g)}/g^{m'\deg(f)}$。
所以 $s$ 的确位于所示箭头的像中。
\end{proof}

\begin{lemma}
\label{constructions-lemma-proj-quasi-compact}
设 $S$ 为分次环。概形 $\text{Proj}(S)$ 拟紧，当且仅当存在有限多个齐次元
$f_1, \ldots, f_n \in S_{+}$，使得
$S_{+} \subset \sqrt{(f_1, \ldots, f_n)}$。在这种情形下，
$\text{Proj}(S) = D_+(f_1) \cup \ldots \cup D_+(f_n)$。
\end{lemma}

\begin{proof}
给定这样一族元素，由《代数》引理 \ref{algebra-lemma-topology-proj}，
标准仿射开集 $D_{+}(f_i)$ 覆盖 $\text{Proj}(S)$。反过来，若
$\text{Proj}(S)$ 拟紧，则可用有限多个标准开集 $D_{+}(f_i)$，
$i = 1, \ldots, n$ 覆盖它；由上述引理可见
$S_{+} \subset \sqrt{(f_1, \ldots, f_n)}$。
\end{proof}

\begin{lemma}
\label{constructions-lemma-structure-morphism-proj}
设 $S$ 为分次环。概形 $\text{Proj}(S)$ 有一个到仿射概形
$\Spec(S_0)$ 的典范态射，它与《代数》定义
\ref{algebra-definition-proj} 给出的拓扑空间映射一致。
\end{lemma}

\begin{proof}
上文已经看到，我们对 $\widetilde S$（相应地 $\widetilde M$）的构造给出
一层 $S_0$-代数（相应地 $S_0$-模）。因此由《概形》引理
\ref{schemes-lemma-morphism-into-affine} 得到一个态射。这个态射限制到
$D_{+}(f)$ 时，来自典范环映射 $S_0 \to S_{(f)}$。映射
$S \to S_f$、$S_{(f)} \to S_f$ 都是 $S_0$-代数映射；
见引理 \ref{constructions-lemma-standard-open}。因此，若齐次素理想
$\mathfrak p \subset S$ 对应于 $\mathbf{Z}$-分次素理想
$\mathfrak p' \subset S_f$ 以及（通常意义下的）素理想
$\mathfrak p'' \subset S_{(f)}$，则它们在 $S_0$ 中的逆像相同。
\end{proof}

\begin{lemma}
\label{constructions-lemma-proj-valuative-criterion}
设 $S$ 为分次环。若 $S$ 作为 $S_0$ 上的代数有限生成，则态射
$\text{Proj}(S) \to \Spec(S_0)$ 满足赋值判据的存在性部分与唯一性部分；
见《概形》定义 \ref{schemes-definition-valuative-criterion}。
\end{lemma}

\begin{proof}
唯一性部分来自 $\text{Proj}(S)$ 是分离的这一事实（引理
\ref{constructions-lemma-proj-separated} 以及《概形》引理
\ref{schemes-lemma-separated-implies-valuative}）。选取齐次元
$x_i \in S_{+}$，$i = 1, \ldots, n$；它们生成 $S$，作为 $S_0$-代数。
令 $d_i = \deg(x_i)$，并置 $d = \text{lcm}\{d_i\}$。假设给定交换图
$$
\xymatrix{
\Spec(K) \ar[r] \ar[d] & \text{Proj}(S) \ar[d] \\
\Spec(A) \ar[r] & \Spec(S_0)
}
$$
，如《概形》定义 \ref{schemes-definition-valuative-criterion} 所述。
记 $v : K^* \to \Gamma$ 为 $A$ 的赋值；见《代数》定义
\ref{algebra-definition-value-group}。可以选取齐次元 $f \in S_{+}$，
使 $\Spec(K)$ 映入 $D_{+}(f)$。于是得到环映射的交换图
$$
\xymatrix{
K & S_{(f)} \ar[l]^{\varphi} \\
A \ar[u] & S_0 \ar[l] \ar[u]
}
$$
重新编号后，可以假设 $\varphi(x_i^{\deg(f)}/f^{d_i})$ 对
$i = 1, \ldots, r$ 非零，而对 $i = r + 1, \ldots, n$ 为零。
由于开集 $D_{+}(x_i)$ 覆盖 $\text{Proj}(S)$，可见 $r \geq 1$。
取指标 $i_0 \in \{1, \ldots, r\}$，使
$\gamma_i = (d/d_i)v(\varphi(x_i^{\deg(f)}/f^{d_i}))$ 在 $\Gamma$ 中最小。
为方便起见，置 $x_0 = x_{i_0}$ 与 $d_0 = d_{i_0}$。环映射
$\varphi$ 分解经过映射 $\varphi' : S_{(fx_0)} \to K$，后者给出环映射
$S_{(x_0)} \to S_{(fx_0)} \to K$。代数 $S_{(x_0)}$ 在 $S_0$ 上由元素
$x_1^{e_1} \ldots x_n^{e_n}/x_0^{e_0}$ 生成，其中
$\sum e_i d_i = e_0 d_0$。若 $e_i > 0$ 对某个 $i > r$ 成立，则
$\varphi'(x_1^{e_1} \ldots x_n^{e_n}/x_0^{e_0}) = 0$。
若 $e_i = 0$ 对 $i > r$ 成立，则
\begin{align*}
d \deg(f) v(\varphi'(x_1^{e_1} \ldots x_r^{e_r}/x_0^{e_0}))
& =
d v(\varphi'(x_1^{e_1 \deg(f)} \ldots x_r^{e_r \deg(f)}/x_0^{e_0 \deg(f)})) \\
& =
d \sum e_i v(\varphi'(x_i^{\deg(f)}/f^{d_i}))
- e_0 v(\varphi'(x_0^{\deg(f)}/f^{d_0})) \\
& =
\sum e_i d_i \gamma_i - e_0 d_0 \gamma_0 \\
& \geq
\sum e_i d_i \gamma_0 - e_0 d_0 \gamma_0 = 0
\end{align*}
，因为 $\gamma_0$ 在诸 $\gamma_i$ 中最小。这说明 $S_{(x_0)}$ 映入
$A$，映射由 $\varphi'$ 给出。相应的概形态射
$\Spec(A) \to \Spec(S_{(x_0)}) = D_{+}(x_0)
\subset \text{Proj}(S)$ 给出嵌入本证明第一个交换图中的所需态射。
\end{proof}

\noindent
在引理 \ref{constructions-lemma-proj-valuative-criterion} 的证明中已经看到，在该引理的
假设下，态射 $\text{Proj}(S) \to \Spec(S_0)$ 也是拟紧的。因此由
《概形》命题 \ref{schemes-proposition-characterize-universally-closed}，
在该引理的情形中 $\text{Proj}(S) \to \Spec(S_0)$ 是万有闭的。
下面给出若干例子，说明若不对分次环 $S$ 作某种假设，这些结果不成立。

\begin{example}
\label{constructions-example-not-existence-valuative-big-proj}
设 $\mathbf{C}[X_1, X_2, X_3, \ldots]$ 为分次 $\mathbf{C}$-代数，
其中每个 $X_i$ 的次数都是 $1$。考虑环映射
$$
\mathbf{C}[X_1, X_2, X_3, \ldots]
\longrightarrow
\mathbf{C}[t^\alpha ; \alpha \in \mathbf{Q}_{\geq 0}]
$$
，它把 $X_i$ 映为 $t^{1/i}$。右端成为赋值环 $A$，这是在理想
$\mathfrak m = (t^\alpha ; \alpha > 0)$ 处局部化后得到的。
设 $K$ 为 $A$ 的分式域。上述映射给出态射
$\Spec(K) \to \text{Proj}(\mathbf{C}[X_1, X_2, X_3, \ldots])$，
它不能延拓为定义在整个 $\Spec(A)$ 上的态射。理由是，$\Spec(A)$ 的像
将包含于某个 $D_{+}(X_i)$ 中，但此时 $X_{i + 1}/X_i$ 应映为 $A$ 中的元素；
事实上并非如此，因为它映为 $t^{1/(i + 1) - 1/i}$。
\end{example}

\begin{example}
\label{constructions-example-not-existence-valuative-small-proj}
设 $R = \mathbf{C}[t]$，并令
$$
S = R[X_1, X_2, X_3, \ldots]/(X_i^2 - tX_{i + 1}).
$$
分次满足 $R = S_0$ 且 $\deg(X_i) = 2^{i - 1}$。注意，若
$\mathfrak p \in \text{Proj}(S)$，则 $t \not \in \mathfrak p$；
否则 $\mathfrak p$ 必须包含所有 $X_i$，这对齐次谱的元素是不允许的。
因此有 $D_{+}(X_i) = D_{+}(X_{i + 1})$，对所有 $i$ 均成立。所以
$\text{Proj}(S)$ 是拟紧的；事实上它是仿射的，因为它等于
$D_{+}(X_1)$。容易看出，$\text{Proj}(S) \to \Spec(R)$ 的像是
$D(t)$。因此态射 $\text{Proj}(S) \to \Spec(R)$ 不是闭态射。
所以赋值判据不能适用，因为它会蕴含该态射是闭的（见《概形》命题
\ref{schemes-proposition-characterize-universally-closed}）。
\end{example}

\begin{example}
\label{constructions-example-trivial-proj}
设 $A$ 为环。令 $S = A[T]$ 为分次 $A$-代数，其中 $T$ 的次数为 $1$。
则典范态射 $\text{Proj}(S) \to \Spec(A)$（见引理
\ref{constructions-lemma-structure-morphism-proj}）是同构。
\end{example}

\begin{example}
\label{constructions-example-open-subset-proj}
设 $X = \Spec(A)$ 为仿射概形，并设 $U \subset X$ 为开子概形。
给 $A[T]$ 分次，规定 $\deg T = 1$。定义 $S$ 为 $A[T]$ 的如下子环：
它由 $A$ 和所有 $fT^i$ 生成，其中 $i \ge 0$，而 $f \in A$ 满足
$D(f) \subset U$。我们断言 $S$ 是满足 $S_0 = A$ 的分次环，并且
$\text{Proj}(S) \cong U$；此同构把引理
\ref{constructions-lemma-structure-morphism-proj} 的典范态射
$\text{Proj}(S) \to \Spec(A)$ 与包含映射 $U \subset X$ 等同起来。

\medskip\noindent
假设 $\mathfrak p \in \text{Proj}(S)$ 满足每个 $fT \in S_1$ 都属于
$\mathfrak p$。则每个生成元 $fT^i$（其中 $i \ge 1$）都属于
$\mathfrak p$，因为 $(fT^i)^2 = (fT)(fT^{2i-1}) \in \mathfrak p$，
且 $\mathfrak p$ 是根理想。但这会给出 $\mathfrak p \supset S_+$，
不可能。因此 $\text{Proj}(S)$ 由标准仿射开子集
$\{D_+(fT)\}_{fT \in S_1}$ 覆盖。

\medskip\noindent
注意，若 $fT \in S_1$，则包含映射 $S \subset A[T]$ 诱导
$S[(fT)^{-1}]$ 与 $A[T, T^{-1}, f^{-1}]$ 的分次同构。因此标准开子集
$D_+(fT) \cong \Spec(S_{(fT)})$ 同构于
$\Spec(A[T, T^{-1}, f^{-1}]_0) = \Spec(A[f^{-1}])$。显然，此同构是
典范态射 $\text{Proj}(S) \to \Spec(A)$ 的限制。若还有 $gT \in S_1$，则
$S[(fT)^{-1}, (gT)^{-1}] \cong A[T, T^{-1}, f^{-1}, g^{-1}]$
是分次环同构，故
$D_+(fT) \cap D_+(gT) \cong \Spec(A[f^{-1}, g^{-1}])$。因此
$\text{Proj}(S)$ 是诸开子概形 $D_+(fT)$ 的并；在典范态射下，
它们同构于开子概形 $D(f) \subset X$，并且这些开子概形在
$\text{Proj}(S)$ 中的相交方式与它们在 $X$ 中的相交方式相同。
由此可知，典范态射给出 $\text{Proj}(S)$ 与所有
$D(f) \subset U$ 之并的同构，而这个并就是 $U$。
\end{example}





















\section{Proj 上的拟凝聚层}
\label{constructions-section-quasi-coherent-proj}

\noindent
设 $S$ 为分次环。设 $M$ 为分次 $S$-模。在引理
\ref{constructions-lemma-proj-sheaves} 中，我们已经构造了拟凝聚模层
$\widetilde{M}$（它位于 $\text{Proj}(S)$ 上）以及映射
\begin{equation}
\label{constructions-equation-map-global-sections}
M_0 \longrightarrow \Gamma(\text{Proj}(S), \widetilde{M})
\end{equation}
，它把次数 $0$ 部分（属于 $M$）映到 $\widetilde{M}$ 的整体截面。
考虑次数 $0$ 部分：它在第 $n$ 次扭转 $M(n)$ 中（原分次模为
$M$；见《代数》章节 \ref{algebra-section-graded}）等于 $M_n$。
因此得到映射
\begin{equation}
\label{constructions-equation-map-global-sections-degree-n}
M_n \longrightarrow \Gamma(\text{Proj}(S), \widetilde{M(n)}).
\end{equation}
我们希望对任意拟凝聚层 $\mathcal{F}$（在 $\text{Proj}(S)$ 上）都能执行此操作。
为此，将与结构层的第 $n$ 次扭转作张量积；见定义
\ref{constructions-definition-twist}。为联系这两个概念，将使用下面的引理。

\begin{lemma}
\label{constructions-lemma-widetilde-tensor}
设 $S$ 为分次环。设
$(X, \mathcal{O}_X) = (\text{Proj}(S), \mathcal{O}_{\text{Proj}(S)})$
为引理 \ref{constructions-lemma-proj-scheme} 中的概形。设 $f \in S_{+}$ 为齐次元。
设 $x \in X$ 为对应于齐次素理想 $\mathfrak p \subset S$ 的点。
设 $M$、$N$ 为分次 $S$-模。存在典范的
$\mathcal{O}_{\text{Proj}(S)}$-模映射
$$
\widetilde M \otimes_{\mathcal{O}_X} \widetilde N
\longrightarrow
\widetilde{M \otimes_S N}
$$
，它诱导典范映射
$
M_{(f)} \otimes_{S_{(f)}} N_{(f)}
\to
(M \otimes_S N)_{(f)}
$
（在 $D_{+}(f)$ 上的截面上），并诱导典范映射
$
M_{(\mathfrak p)} \otimes_{S_{(\mathfrak p)}} N_{(\mathfrak p)}
\to
(M \otimes_S N)_{(\mathfrak p)}
$
（在点 $x$ 的茎上）。此外，下图交换：
$$
\xymatrix{
M_0 \otimes_{S_0} N_0 \ar[r] \ar[d] &
(M \otimes_S N)_0 \ar[d] \\
\Gamma(X, \widetilde M \otimes_{\mathcal{O}_X} \widetilde N) \ar[r] &
\Gamma(X, \widetilde{M \otimes_S N})
}
$$
其中竖直映射由（\ref{constructions-equation-map-global-sections}）给出。
\end{lemma}

\begin{proof}
构造所示态射等价于构造 $\mathcal{O}_X$-双线性映射
$$
\widetilde M \times \widetilde N
\longrightarrow
\widetilde{M \otimes_S N}
$$
；见《模》章节 \ref{modules-section-tensor-product}。只须在开集
$D_{+}(f)$ 上的截面上定义它，并使之与限制映射相容。在 $D_{+}(f)$ 上，
使用 $S_{(f)}$-双线性映射
$M_{(f)} \times N_{(f)} \to (M \otimes_S N)_{(f)}$,
$(x/f^n, y/f^m) \mapsto (x \otimes y)/f^{n + m}$。细节从略。
\end{proof}

\begin{remark}
\label{constructions-remark-not-isomorphism}
一般而言，上面引理 \ref{constructions-lemma-widetilde-tensor} 中构造的映射不是同构。
下面给出一个例子。设 $k$ 为域。设 $S = k[x, y, z]$，其中 $k$ 的次数为
$0$，而 $\deg(x) = 1$、$\deg(y) = 2$、$\deg(z) = 3$。
设 $M = S(1)$ 且 $N = S(2)$；记号见《代数》章节
\ref{algebra-section-graded}。于是 $M \otimes_S N = S(3)$。注意
\begin{eqnarray*}
S_z
& = &
k[x, y, z, 1/z] \\
S_{(z)}
& = &
k[x^3/z, xy/z, y^3/z^2]
\cong
k[u, v, w]/(uw - v^3) \\
M_{(z)} & = & S_{(z)} \cdot x + S_{(z)} \cdot y^2/z \subset S_z \\
N_{(z)} & = & S_{(z)} \cdot y + S_{(z)} \cdot x^2 \subset S_z \\
S(3)_{(z)} & = & S_{(z)} \cdot z \subset S_z
\end{eqnarray*}
考虑极大理想 $\mathfrak m = (u, v, w) \subset S_{(z)}$。不难看出，
$M_{(z)}/\mathfrak mM_{(z)}$ 和 $N_{(z)}/\mathfrak mN_{(z)}$
的维数都是 $2$（在 $\kappa(\mathfrak m)$ 上）。但
$S(3)_{(z)}/\mathfrak mS(3)_{(z)}$ 的维数是 $1$。因此映射
$M_{(z)} \otimes N_{(z)} \to S(3)_{(z)}$ 不是同构。
\end{remark}









\section{Proj 上的可逆层}
\label{constructions-section-invertible-on-proj}

\noindent
回顾《代数》章节 \ref{algebra-section-graded} 中，对分次环上的分次模所关联的
扭转模 $M(n)$ 之构造。

\begin{definition}
\label{constructions-definition-twist}
设 $S$ 为分次环。设 $X = \text{Proj}(S)$。
\begin{enumerate}
\item 定义 $\mathcal{O}_X(n) = \widetilde{S(n)}$。它称为第 $n$ 个
{\it $\text{Proj}(S)$ 的结构层扭转}。
\item 对任意 $\mathcal{O}_X$-模层 $\mathcal{F}$，置
$\mathcal{F}(n) = \mathcal{F} \otimes_{\mathcal{O}_X} \mathcal{O}_X(n)$.
\end{enumerate}
\end{definition}

\noindent
将使用引理 \ref{constructions-lemma-widetilde-tensor} 构造一些典范映射。
由于 $S(n) \otimes_S S(m) = S(n + m)$，可见存在典范映射
\begin{equation}
\label{constructions-equation-multiply}
\mathcal{O}_X(n) \otimes_{\mathcal{O}_X} \mathcal{O}_X(m)
\longrightarrow
\mathcal{O}_X(n + m).
\end{equation}
这些映射一般不是同构；见注记 \ref{constructions-remark-not-isomorphism} 中的例子。
同一个例子表明，$\mathcal{O}_X(n)$ 一般{\it 不是} $X$ 上的可逆层。
与任意 $\mathcal{O}_X$-模 $\mathcal{F}$ 作张量积，得到映射
\begin{equation}
\label{constructions-equation-multiply-on-sheaf}
\mathcal{O}_X(n) \otimes_{\mathcal{O}_X} \mathcal{F}(m)
\longrightarrow
\mathcal{F}(n + m).
\end{equation}
整体截面上的映射（\ref{constructions-equation-map-global-sections-degree-n}）给出分次环映射
\begin{equation}
\label{constructions-equation-global-sections}
S \longrightarrow \bigoplus\nolimits_{n \geq 0} \Gamma(X, \mathcal{O}_X(n)).
\end{equation}
而对任意 $\mathcal{O}_X$-模 $\mathcal{F}$，映射
（\ref{constructions-equation-multiply-on-sheaf}）给出分次模结构
\begin{equation}
\label{constructions-equation-global-sections-module}
\bigoplus\nolimits_{n \geq 0} \Gamma(X, \mathcal{O}_X(n))
\times
\bigoplus\nolimits_{m \in \mathbf{Z}} \Gamma(X, \mathcal{F}(m))
\longrightarrow
\bigoplus\nolimits_{m \in \mathbf{Z}} \Gamma(X, \mathcal{F}(m))
\end{equation}
，并且通过（\ref{constructions-equation-global-sections}）也得到 $S$-模结构。
更一般地，给定任意分次 $S$-模 $M$，有
$M(n) = M \otimes_S S(n)$。因此得到映射
\begin{equation}
\label{constructions-equation-multiply-more-generally}
\widetilde M(n)
=
\widetilde M
\otimes_{\mathcal{O}_X}
\mathcal{O}_X(n)
\longrightarrow
\widetilde{M(n)}.
\end{equation}
在整体截面上，（\ref{constructions-equation-map-global-sections-degree-n}）定义分次
$S$-模映射
\begin{equation}
\label{constructions-equation-global-sections-more-generally}
M \longrightarrow
\bigoplus\nolimits_{n \in \mathbf{Z}} \Gamma(X, \widetilde{M(n)}).
\end{equation}
下面是一个基本上直接由定义得出的重要事实。

\begin{lemma}
\label{constructions-lemma-when-invertible}
设 $S$ 为分次环。置 $X = \text{Proj}(S)$。设 $f \in S$ 为次数
$d > 0$ 的齐次元。层 $\mathcal{O}_X(nd)|_{D_{+}(f)}$ 是可逆的；事实上，
对所有 $n \in \mathbf{Z}$ 它们都是平凡的（见《模》定义
\ref{modules-definition-invertible}）。映射（\ref{constructions-equation-multiply}）
限制到 $D_{+}(f)$ 后为
$$
\mathcal{O}_X(nd)|_{D_{+}(f)} \otimes_{\mathcal{O}_{D_{+}(f)}}
\mathcal{O}_X(m)|_{D_{+}(f)}
\longrightarrow
\mathcal{O}_X(nd + m)|_{D_{+}(f)},
$$
映射（\ref{constructions-equation-multiply-on-sheaf}）限制到 $D_+(f)$ 后为
$$
\mathcal{O}_X(nd)|_{D_{+}(f)} \otimes_{\mathcal{O}_{D_{+}(f)}}
\mathcal{F}(m)|_{D_{+}(f)}
\longrightarrow
\mathcal{F}(nd + m)|_{D_{+}(f)},
$$
而映射（\ref{constructions-equation-multiply-more-generally}）限制到 $D_{+}(f)$ 后为
$$
\widetilde M(nd)|_{D_{+}(f)}
=
\widetilde M|_{D_{+}(f)}
\otimes_{\mathcal{O}_{D_{+}(f)}}
\mathcal{O}_X(nd)|_{D_{+}(f)}
\longrightarrow
\widetilde{M(nd)}|_{D_{+}(f)}
$$
；对所有 $n, m \in \mathbf{Z}$，这些映射都是同构。
\end{lemma}

\begin{proof}
分次 $S$-模映射 $S \to S(nd)$ 与 $M \to M(nd)$ 由
$x \mapsto f^n x$ 给出；在使 $f$ 可逆后，它们成为同构。第一个映射表明
$S_{(f)} \cong S(nd)_{(f)}$，从而给出同构
$\mathcal{O}_{D_{+}(f)} \cong \mathcal{O}_X(nd)|_{D_{+}(f)}$。
第二个映射表明
$S(nd)_{(f)} \otimes_{S_{(f)}} M_{(f)} \to M(nd)_{(f)}$
是同构。映射（\ref{constructions-equation-multiply-on-sheaf}）的情形由映射
（\ref{constructions-equation-multiply}）的情形得出。
\end{proof}

\begin{lemma}
\label{constructions-lemma-apply-modules}
设 $S$ 为分次环。设 $M$ 为分次 $S$-模。置 $X = \text{Proj}(S)$。
假设 $X$ 由标准开集 $D_+(f)$ 覆盖，其中 $f \in S_1$；例如，当 $S$
由 $S_1$ 在 $S_0$ 上生成时即如此。则层 $\mathcal{O}_X(n)$ 是可逆的，
而映射（\ref{constructions-equation-multiply}）、（\ref{constructions-equation-multiply-on-sheaf}）与
（\ref{constructions-equation-multiply-more-generally}）都是同构。特别地，它们诱导同构
$$
\mathcal{O}_X(1)^{\otimes n} \cong
\mathcal{O}_X(n)
\quad
\text{且}
\quad
\widetilde{M} \otimes_{\mathcal{O}_X} \mathcal{O}_X(n) =
\widetilde{M}(n) \cong \widetilde{M(n)}
$$
因此（\ref{constructions-equation-map-global-sections-degree-n}）成为映射
\begin{equation}
\label{constructions-equation-map-global-sections-degree-n-simplified}
M_n \longrightarrow \Gamma(X, \widetilde{M}(n))
\end{equation}
，而（\ref{constructions-equation-global-sections-more-generally}）成为映射
\begin{equation}
\label{constructions-equation-global-sections-more-generally-simplified}
M \longrightarrow
\bigoplus\nolimits_{n \in \mathbf{Z}} \Gamma(X, \widetilde{M}(n)).
\end{equation}
\end{lemma}

\begin{proof}
在本引理的假设下，$X$ 由开子集 $D_{+}(f)$ 覆盖，其中 $f \in S_1$；
于是本引理由上面的引理 \ref{constructions-lemma-when-invertible} 得出。
\end{proof}

\begin{lemma}
\label{constructions-lemma-where-invertible}
设 $S$ 为分次环。置 $X = \text{Proj}(S)$。固定整数 $d \geq 1$。
$X$ 的下列开子集相等：
\begin{enumerate}
\item 最大开子集 $W = W_d \subset X$，满足每个
$\mathcal{O}_X(dn)|_W$ 都可逆，且所有乘法映射
$\mathcal{O}_X(nd)|_W \otimes_{\mathcal{O}_W} \mathcal{O}_X(md)|_W
\to \mathcal{O}_X(nd + md)|_W$
（见 \ref{constructions-equation-multiply}）都是同构。
\item 开子集 $D_{+}(fg)$ 的并，其中 $f, g \in S$ 齐次且
$\deg(f) = \deg(g) + d$。
\end{enumerate}
此外，所有映射
$\widetilde M(nd)|_W = \widetilde M|_W \otimes_{\mathcal{O}_W}
\mathcal{O}_X(nd)|_W \to \widetilde{M(nd)}|_W$
（见 \ref{constructions-equation-multiply-more-generally}）都是同构。
\end{lemma}

\begin{proof}
若 $x \in D_{+}(fg)$ 且 $\deg(f) = \deg(g) + d$，则在 $D_{+}(fg)$ 上，
层 $\mathcal{O}_X(dn)$ 由元素 $(f/g)^n = f^{2n}/(fg)^n$ 生成。
像引理 \ref{constructions-lemma-when-invertible} 的证明那样论证，可知 $x$ 属于 (1) 中
定义的开子集 $W$。

\medskip\noindent
反过来，假设 $\mathcal{O}_X(dn)$ 对所有 $n$ 都在一个开邻域 $V$
（它是点 $x \in X$ 的邻域）上秩一自由，并且所有乘法映射
$\mathcal{O}_X(nd)|_V \otimes_{\mathcal{O}_V} \mathcal{O}_X(md)|_V
\to \mathcal{O}_X(nd + md)|_V$ 都是同构。可选取齐次元 $h \in S_{+}$，使得
$x \in D_{+}(h) \subset V$。由结构层扭转的定义可知，存在元素
$s$（属于 $(S_h)_d$），使 $s^n$ 是 $(S_h)_{nd}$ 作为
$S_{(h)}$-模的一组基；此事对所有 $n \in \mathbf{Z}$ 成立。可以写成
$s = f/h^m$，其中 $m \geq 1$ 且 $f \in S_{d + m \deg(h)}$。
置 $g = h^m$，于是 $s = f/g$。由构造可知 $x \in D_{+}(g)$。
注意 $g^d \in (S_h)_{d\deg(g)}$。由假设，可把它写成
$s^{\deg(g)} = f^{\deg(g)}/g^{\deg(g)}$ 的倍数，即
$g^d = a/g^e \cdot f^{\deg(g)}/g^{\deg(g)}$。于是得到
$g^{d + e + \deg(g)} = a f^{\deg(g)}$，从而还有
$x \in D_{+}(f)$。所以 $x$ 属于 (2) 中定义的集合。

\medskip\noindent
生成截面 $s = f/g$ 的存在性（它位于仿射开集 $D_{+}(fg)$ 上），
连同其各次幂自由生成模层
$\mathcal{O}_X(nd)$；这容易蕴含乘法映射
$\widetilde M(nd)|_W = \widetilde M|_W \otimes_{\mathcal{O}_W}
\mathcal{O}_X(nd)|_W \to \widetilde{M(nd)}|_W$
（见 \ref{constructions-equation-multiply-more-generally}）是同构。可与引理
\ref{constructions-lemma-when-invertible} 的证明比较。
\end{proof}

\noindent
回顾《模》引理 \ref{modules-lemma-s-open}：给定可逆层 $\mathcal{L}$
（在局部环层空间 $X$ 上），并给定整体截面 $s$（属于 $\mathcal{L}$），则集合
$X_s = \{x \in X \mid s \not \in \mathfrak m_x\mathcal{L}_x\}$ 是开的。

\begin{lemma}
\label{constructions-lemma-principal-open}
设 $S$ 为分次环。置 $X = \text{Proj}(S)$。固定整数 $d \geq 1$。
设 $W = W_d \subset X$ 为引理 \ref{constructions-lemma-where-invertible} 中定义的
开子概形。设 $n \geq 1$ 且 $f \in S_{nd}$。用
$s \in \Gamma(W, \mathcal{O}_W(nd))$ 表示如下截面：它是 $f$ 经由
（\ref{constructions-equation-global-sections}）在 $W$ 上的限制所得的像。则
$$
W_s = D_{+}(f) \cap W.
$$
\end{lemma}

\begin{proof}
设 $D_{+}(ab) \subset W$ 为标准仿射开集，其中 $a, b \in S$ 齐次且
$\deg(a) = \deg(b) + d$。注意
$D_{+}(ab) \cap D_{+}(f) = D_{+}(abf)$。另一方面，$s$ 在
$D_{+}(ab)$ 上的限制对应于元素
$f/1 = b^nf/a^n (a/b)^n \in (S_{ab})_{nd}$。在引理
\ref{constructions-lemma-where-invertible} 的证明中已经看到，$(a/b)^n$ 是
$\mathcal{O}_W(nd)$ 在 $D_{+}(ab)$ 上的生成元。因此
$W_s \cap D_{+}(ab)$ 是与
$b^nf/a^n \in \mathcal{O}_X(D_{+}(ab))$ 关联的主开集。
本引理的结论由此显然。
\end{proof}

\noindent
下面的引理陈述若干性质；以后将用它们刻画带有丰沛可逆层的概形。

\begin{lemma}
\label{constructions-lemma-ample-on-proj}
设 $S$ 为分次环。设 $X = \text{Proj}(S)$。设 $Y \subset X$ 为拟紧开子概形。
用 $\mathcal{O}_Y(n)$ 表示 $\mathcal{O}_X(n)$ 在 $Y$ 上的限制。
存在整数 $d \geq 1$，使得
\begin{enumerate}
\item 子概形 $Y$ 包含于引理 \ref{constructions-lemma-where-invertible} 中定义的开集
$W_d$；
\item 层 $\mathcal{O}_Y(dn)$ 对所有 $n \in \mathbf{Z}$ 都可逆；
\item 所有映射
$\mathcal{O}_Y(nd) \otimes_{\mathcal{O}_Y} \mathcal{O}_Y(m)
\longrightarrow
\mathcal{O}_Y(nd + m)$
（来自等式（\ref{constructions-equation-multiply}））都是同构；
\item 所有映射
$\widetilde M(nd)|_Y = \widetilde M|_Y \otimes_{\mathcal{O}_Y}
\mathcal{O}_X(nd)|_Y \to \widetilde{M(nd)}|_Y$
（见 \ref{constructions-equation-multiply-more-generally}）都是同构；
\item 给定 $f \in S_{nd}$，用 $s \in \Gamma(Y, \mathcal{O}_Y(nd))$
表示 $f$ 经由（\ref{constructions-equation-global-sections}）在 $Y$ 上限制所得的像，
则 $D_{+}(f) \cap Y = Y_s$；
\item $Y$ 的拓扑有一个基，它由开集 $Y_s$ 构成，其中
$s \in \Gamma(Y, \mathcal{O}_Y(nd))$ 且 $n \geq 1$；并且
\item $Y$ 的拓扑有一个基，它由那些仿射开集 $Y_s \subset Y$ 构成，其中
$s \in \Gamma(Y, \mathcal{O}_Y(nd))$ 且 $n \geq 1$。
\end{enumerate}
\end{lemma}

\begin{proof}
由于 $Y$ 拟紧，存在有限多个齐次元 $f_i \in S_{+}$，
$i = 1, \ldots, n$，使标准开集 $D_{+}(f_i)$ 给出 $Y$ 的开覆盖。
令 $d_i = \deg(f_i)$，并置 $d = d_1 \ldots d_n$。注意
$D_{+}(f_i) = D_{+}(f_i^{d/d_i})$，因此由引理
\ref{constructions-lemma-where-invertible} 的刻画 (2)，或使用引理
\ref{constructions-lemma-when-invertible} 的 (1)，立刻可见 $Y \subset W_d$。
注意，由引理 \ref{constructions-lemma-where-invertible}，(1) 蕴含 (2)、(3) 与 (4)。
（注意 (3) 是 (4) 的特殊情形。）断言 (5) 来自引理
\ref{constructions-lemma-principal-open}。断言 (6) 和 (7) 成立，是因为开子集
$D_{+}(f)$ 构成 $X$ 的拓扑基，并且都是仿射的。
\end{proof}

\begin{lemma}
\label{constructions-lemma-comparison-proj-quasi-coherent}
设 $S$ 为分次环。置 $X = \text{Proj}(S)$。设 $\mathcal{F}$ 为拟凝聚
$\mathcal{O}_X$-模。利用（\ref{constructions-equation-global-sections-module}）和
（\ref{constructions-equation-global-sections}），把
$M = \bigoplus_{n \in \mathbf{Z}} \Gamma(X, \mathcal{F}(n))$ 视为分次
$S$-模。则存在典范 $\mathcal{O}_X$-模映射
$$
\widetilde{M} \longrightarrow \mathcal{F}
$$
，它关于 $\mathcal{F}$ 具有函子性，并且诱导映射
$M_0 \to \Gamma(X, \mathcal{F})$ 是恒等映射。
\end{lemma}

\begin{proof}
设 $f \in S$ 为次数 $d > 0$ 的齐次元。回顾由引理
\ref{constructions-lemma-proj-sheaves}，$\widetilde{M}|_{D_{+}(f)}$ 对应于
$S_{(f)}$-模 $M_{(f)}$。因此可以定义典范映射
$$
M_{(f)} \longrightarrow \Gamma(D_+(f), \mathcal{F}),\quad
m/f^n \longmapsto m|_{D_+(f)} \otimes f|_{D_+(f)}^{-n}
$$
；它有意义，因为 $f|_{D_+(f)}$ 是可逆层
$\mathcal{O}_X(d)|_{D_+(f)}$ 的平凡化截面；见引理
\ref{constructions-lemma-when-invertible} 及其证明。由于 $\widetilde{M}$ 拟凝聚，
这给出典范映射
$$
\widetilde{M}|_{D_+(f)} \longrightarrow \mathcal{F}|_{D_+(f)}
$$
；这里使用《概形》引理 \ref{schemes-lemma-compare-constructions}。
若证明所示映射可在交叠处粘合，就得到整体映射。这一点的证明从略。
最后一个断言的证明也从略。
\end{proof}










\section{Proj 的函子性}
\label{constructions-section-functoriality-proj}

\noindent
分次环映射 $\psi : A \to B$ 不总能给出相应射影齐次谱之间的态射。
原因是，逆像 $\psi^{-1}(\mathfrak q)$（其中
$\mathfrak q \subset B$ 是齐次素理想）可能包含无关素理想 $A_{+}$，即使
$\mathfrak q$ 不包含 $B_{+}$。正确的结果如下。

\begin{lemma}
\label{constructions-lemma-morphism-proj}
设 $A$、$B$ 为两个分次环。置 $X = \text{Proj}(A)$ 且
$Y = \text{Proj}(B)$。设 $\psi : A \to B$ 为分次环映射。置
$$
U(\psi)
=
\bigcup\nolimits_{f \in A_{+}\ \text{齐次}} D_{+}(\psi(f))
\subset Y.
$$
则存在典范概形态射
$$
r_\psi :
U(\psi)
\longrightarrow
X
$$
以及 $\mathbf{Z}$-分次 $\mathcal{O}_{U(\psi)}$-代数映射
$$
\theta = \theta_\psi :
r_\psi^*\left(
\bigoplus\nolimits_{d \in \mathbf{Z}} \mathcal{O}_X(d)
\right)
\longrightarrow
\bigoplus\nolimits_{d \in \mathbf{Z}} \mathcal{O}_{U(\psi)}(d).
$$
三元组 $(U(\psi), r_\psi, \theta)$ 由下列性质刻画：
\begin{enumerate}
\item 对每个 $d \geq 0$，图
$$
\xymatrix{
A_d \ar[d] \ar[rr]_{\psi} & &
B_d \ar[d] \\
\Gamma(X, \mathcal{O}_X(d)) \ar[r]^-\theta &
\Gamma(U(\psi), \mathcal{O}_Y(d)) &
\Gamma(Y, \mathcal{O}_Y(d)) \ar[l]
}
$$
交换。
\item 对任意齐次元 $f \in A_{+}$，有
$r_\psi^{-1}(D_{+}(f)) = D_{+}(\psi(f))$，并且 $r_\psi$ 在
$D_{+}(\psi(f))$ 上的限制，对应于环映射
$A_{(f)} \to B_{(\psi(f))}$；此映射由 $\psi$ 诱导。
\end{enumerate}
\end{lemma}

\begin{proof}
显然，条件 (2) 唯一确定概形态射与开子集 $U(\psi)$。选取
$f \in A_d$，其中 $d \geq 1$。注意，$\mathcal{O}_X(n)|_{D_{+}(f)}$
对应于 $A_{(f)}$-模 $(A_f)_n$，而
$\mathcal{O}_Y(n)|_{D_{+}(\psi(f))}$ 对应于 $B_{(\psi(f))}$-模
$(B_{\psi(f)})_n$。换言之，$\theta$ 限制到 $D_{+}(\psi(f))$ 时，
对应于 $\mathbf{Z}$-分次 $B_{(\psi(f))}$-代数映射
$$
A_f \otimes_{A_{(f)}} B_{(\psi(f))}
\longrightarrow
B_{\psi(f)}
$$
条件 (1) 确定 $A$ 的所有元素之像。由于 $f$ 是可逆元且映为
$\psi(f)$，可见 $1/f^m$ 映为 $1/\psi(f)^m$。由此容易看出，
$\theta$ 被唯一确定；具体由规则
$$
a/f^m \otimes b/\psi(f)^e \longmapsto \psi(a)b/\psi(f)^{m + e}.
$$
给出。为证明存在性，注意上述唯一性证明给出如下良定义规则：把态射 $r$
和映射 $\theta$ 限制到每个形如 $D_{+}(\psi(f)) \subset U(\psi)$ 的
标准开集，并映入 $D_{+}(f)$。称它们为 $r_f$ 与 $\theta_f$。
因此只须验证：若 $D_{+}(f) \subset D_{+}(g)$ 对某些齐次元
$f, g \in A_{+}$ 成立，则 $r_g$ 在 $D_{+}(\psi(f))$ 上的限制
与 $r_f$ 一致。这由上面对 $r$ 与 $\theta$ 给出的公式显然成立。
\end{proof}

\begin{lemma}
\label{constructions-lemma-morphism-proj-transitive}
设 $A$、$B$ 与 $C$ 为分次环。置 $X = \text{Proj}(A)$、
$Y = \text{Proj}(B)$ 与 $Z = \text{Proj}(C)$。设
$\varphi : A \to B$、$\psi : B \to C$ 为分次环映射。则有
$$
U(\psi \circ \varphi) = r_\psi^{-1}(U(\varphi))
\quad
\text{且}
\quad
r_{\psi \circ \varphi}
=
r_\varphi \circ r_\psi|_{U(\psi \circ \varphi)}.
$$
此外有
$$
\theta_\psi \circ r_\psi^*\theta_\varphi
=
\theta_{\psi \circ \varphi}
$$
，其中记号含义显然。
\end{lemma}

\begin{proof}
略。
\end{proof}

\begin{lemma}
\label{constructions-lemma-surjective-graded-rings-map-proj}
沿用上面引理 \ref{constructions-lemma-morphism-proj} 的假设与记号。假设
$A_d \to B_d$ 对所有 $d \gg 0$ 都是满射。则
\begin{enumerate}
\item $U(\psi) = Y$,
\item $r_\psi : Y \to X$ 是闭浸入；并且
\item 映射 $\theta : r_\psi^*\mathcal{O}_X(n) \to \mathcal{O}_Y(n)$
都是满射，但一般不是同构（即使 $A \to B$ 是满射）。
\end{enumerate}
\end{lemma}

\begin{proof}
第 (1) 部分来自 $U(\psi)$ 的定义以及如下事实：
$D_{+}(f) = D_{+}(f^n)$ 对任意 $n > 0$ 成立。对齐次元
$f \in A_{+}$，可见
$A_{(f)} \to B_{(\psi(f))}$ 是满射，因为 $B_{(\psi(f))}$ 的任意元素
都可表示为分式 $b/\psi(f)^n$，其中 $n$ 可任意大（这迫使
$b \in B$ 的次数很大）。这证明了 (2)。同一论证表明映射
$$
A_f \to B_{\psi(f)}
$$
是满射，从而证明 $\theta$ 的满射性。作为这个映射不是同构的例子，
考虑分次环 $A = k[x, y]$，其中 $k$ 为域，且 $\deg(x) = 1$、
$\deg(y) = 2$。置 $I = (x)$，从而 $B = k[y]$。注意在此情形下
$\mathcal{O}_Y(1) = 0$。但容易看出 $r_\psi^*\mathcal{O}_X(1)$ 非零。
（也有不那么愚蠢的例子。）
\end{proof}

\begin{lemma}
\label{constructions-lemma-eventual-iso-graded-rings-map-proj}
沿用上面引理 \ref{constructions-lemma-morphism-proj} 的假设与记号。假设
$A_d \to B_d$ 对所有 $d \gg 0$ 都是同构。则
\begin{enumerate}
\item $U(\psi) = Y$,
\item $r_\psi : Y \to X$ 是同构；并且
\item 映射 $\theta : r_\psi^*\mathcal{O}_X(n) \to \mathcal{O}_Y(n)$
都是同构。
\end{enumerate}
\end{lemma}

\begin{proof}
由引理 \ref{constructions-lemma-surjective-graded-rings-map-proj} 得到 (1)。设
$f \in A_{+}$ 为齐次元。关于 $\psi$ 的假设蕴含
$A_f \to B_f$ 是同构（细节从略）。因此显然 $r_\psi$ 与 $\theta$
限制到 $D_{+}(f)$ 上都是同构。本引理得证。
\end{proof}

\begin{lemma}
\label{constructions-lemma-surjective-graded-rings-generated-degree-1-map-proj}
沿用上面引理 \ref{constructions-lemma-morphism-proj} 的假设与记号。
假设 $A_d \to B_d$ 对于 $d \gg 0$ 是满射，并且 $A$ 由 $A_1$
在 $A_0$ 上生成。则
\begin{enumerate}
\item $U(\psi) = Y$,
\item $r_\psi : Y \to X$ 是闭浸入；并且
\item 各映射 $\theta : r_\psi^*\mathcal{O}_X(n) \to \mathcal{O}_Y(n)$
均为同构。
\end{enumerate}
\end{lemma}

\begin{proof}
由引理 \ref{constructions-lemma-eventual-iso-graded-rings-map-proj} 与
\ref{constructions-lemma-morphism-proj-transitive}
可知，我们可将 $B$ 替换为 $A \to B$ 的像，
而不改变 $X$ 或各层 $\mathcal{O}_X(n)$。
因此可假设 $A \to B$ 是满射。由引理
\ref{constructions-lemma-surjective-graded-rings-map-proj} 得到 (1)、(2)
以及 (3) 中映射的满射性。
由引理 \ref{constructions-lemma-apply-modules} 可知，
$\mathcal{O}_X(n)$ 与 $\mathcal{O}_Y(n)$
都是可逆层。因此 $\theta$ 是同构。
\end{proof}

\begin{lemma}
\label{constructions-lemma-base-change-map-proj}
\begin{reference}
\cite[II, Proposition 2.8.10]{EGA}
\end{reference}
沿用上面引理 \ref{constructions-lemma-morphism-proj} 的假设与记号。
假设存在环同态 $R \to A_0$ 与环同态
$R \to R'$，使得 $B = R' \otimes_R A$。则
\begin{enumerate}
\item $U(\psi) = Y$,
\item 图表
$$
\xymatrix{
Y = \text{Proj}(B) \ar[r]_{r_\psi} \ar[d] &
\text{Proj}(A) = X \ar[d] \\
\Spec(R') \ar[r] &
\Spec(R)
}
$$
是纤维积方块；并且
\item 各映射 $\theta : r_\psi^*\mathcal{O}_X(n) \to \mathcal{O}_Y(n)$
均为同构。
\end{enumerate}
\end{lemma}

\begin{proof}
只需考察标准开集 $D_{+}(f)$（其中 $f \in A_{+}$）上的情形，
即刻可得结论。
\end{proof}

\begin{lemma}
\label{constructions-lemma-localization-map-proj}
沿用上面引理 \ref{constructions-lemma-morphism-proj} 的假设与记号。
假设存在 $g \in A_0$，使得 $\psi$ 诱导同构
$A_g \to B$。则
$U(\psi) = Y$，$r_\psi : Y \to X$ 是开浸入，
并诱导 $Y$ 与
$D(g) \subset \Spec(A_0)$ 的逆像之间的同构。此外，映射 $\theta$
是同构。
\end{lemma}

\begin{proof}
这是上面引理 \ref{constructions-lemma-base-change-map-proj} 的特殊情形。
\end{proof}

\begin{lemma}
\label{constructions-lemma-d-uple}
设 $S$ 为分次环。设 $d \geq 1$。采用代数章第
\ref{algebra-section-graded} 节的记号，置 $S' = S^{(d)}$。置
$X = \text{Proj}(S)$ 与 $X' = \text{Proj}(S')$。存在概形的典范
同构 $i : X \to X'$，使得
\begin{enumerate}
\item 对任一分次 $S$-模 $M$，置 $M' = M^{(d)}$，
都有典范同构 $\widetilde{M} \to i^*\widetilde{M'}$；
\item 有典范同构
$\mathcal{O}_{X}(nd) \to i^*\mathcal{O}_{X'}(n)$
\end{enumerate}
且这些同构与引理 \ref{constructions-lemma-widetilde-tensor} 的乘法映射相容，
从而与下列映射相容：
(\ref{constructions-equation-multiply}),
(\ref{constructions-equation-multiply-on-sheaf}),
(\ref{constructions-equation-global-sections}),
(\ref{constructions-equation-global-sections-module}),
(\ref{constructions-equation-multiply-more-generally})，以及
(\ref{constructions-equation-global-sections-more-generally})（精确表述见证明。
\end{lemma}

\begin{proof}
单射环同态 $S' \to S$（按我们的约定，它不是分次环同态）
诱导映射 $j : \Spec(S) \to \Spec(S')$。
给定分次素理想 $\mathfrak p \subset S$，可见
$\mathfrak p' = j(\mathfrak p) = S' \cap \mathfrak p$
是 $S'$ 的分次素理想。
此外，若 $f \in S_+$ 齐次且 $f \not \in \mathfrak p$，则
$f^d \in S'_+$ 且 $f^d \not \in \mathfrak p'$。反之，若
$\mathfrak p' \subset S'$ 是不含某个齐次元
$f \in S'_+$ 的分次素理想，则
$\mathfrak p = \{g \in S \mid g^d \in \mathfrak p'\}$ 是
$S$ 的一个不含 $f$ 的分次素理想，并且它在 $j$ 下的像
是 $\mathfrak p'$。为说明 $\mathfrak p$ 是理想，注意若
$g, h \in \mathfrak p$，则由二项式公式有
$(g + h)^{2d} \in \mathfrak p'$，
进而 $g + h \in \mathfrak p'$，因为 $\mathfrak p'$ 是素理想。
由此可见，$j$ 诱导同胚 $i : X \to X'$。
此外，给定齐次元 $f \in S_+$，有
$S_{(f)} \cong S'_{(f^d)}$。由于这些同构与
引理 \ref{constructions-lemma-standard-open} 的限制映射相容，可知存在结构层的
同构 $i^\sharp : i^{-1}\mathcal{O}_{X'} \to \mathcal{O}_X$，
它位于 $X$ 与 $X'$ 上，因而 $i$ 是概形的同构。

\medskip\noindent
设 $M$ 为分次 $S$-模。给定齐次元 $f \in S_+$，有
$M_{(f)} \cong M'_{(f^d)}$，故以与上面完全相同的方式
得到 (1) 中的同构。(2) 中的同构是 (1) 在
$M = S(nd)$ 时的特殊情形，此时 $M' = S'(n)$。设 $M$ 与 $N$
为分次 $S$-模。于是有
$$
M' \otimes_{S'} N' =
(M \otimes_S N)^{(d)} =
(M \otimes_S N)'
$$
这可由定义直接验证。于是，与引理
\ref{constructions-lemma-widetilde-tensor} 的乘法映射相容，正是指下图交换：
$$
\xymatrix{
\widetilde M \otimes_{\mathcal{O}_X} \widetilde N
\ar[d]_{(1) \otimes (1)} \ar[r] &
\widetilde{M \otimes_S N} \ar[d]^{(1)} \\
i^*(\widetilde{M'} \otimes_{\mathcal{O}_{X'}} \widetilde{N'}) \ar[r] &
i^*(\widetilde{M' \otimes_{S'} N'})
}
$$
考察开集 $D_+(f) = D_+(f^d)$ 上这些映射的构造即可看出这一点：
上方水平箭头由映射
$M_{(f)} \times N_{(f)} \to (M \otimes_S N)_{(f)}$
给出，而下方水平箭头由映射
$M'_{(f^d)} \times N'_{(f^d)} \to (M' \otimes_{S'} N')_{(f^d)}$
给出。由于这些映射经由诸如
$M_{(f)} = M'_{(f^d)}$ 等识别彼此相同，便得到所需的相容性。
其余相容性的证明从略。
\end{proof}











\section{到 Proj 的态射}
\label{constructions-section-morphisms-proj}

\noindent
设 $S$ 为分次环。
设 $X = \text{Proj}(S)$ 为 $S$ 的齐次谱。
设 $d \geq 1$ 为整数。
考虑开子概形
\begin{equation}
\label{constructions-equation-Ud}
U_d = \bigcup\nolimits_{f  \in S_d} D_{+}(f)
\quad\subset\quad
X = \text{Proj}(S)
\end{equation}
注意 $d | d' \Rightarrow U_d \subset U_{d'}$，且
$X = \bigcup_d U_d$。$X$ 与 $U_d$ 都不必拟紧，见代数章引理
\ref{algebra-lemma-topology-proj}。
记 $\mathcal{O}_{U_d}(n) = \mathcal{O}_X(n)|_{U_d}$。
由引理 \ref{constructions-lemma-when-invertible} 可知，
$\mathcal{O}_{U_d}(nd)$（$n \in \mathbf{Z}$）
是可逆 $\mathcal{O}_{U_d}$-模，并且所有乘法映射
$\mathcal{O}_{U_d}(nd) \otimes_{\mathcal{O}_{U_d}} \mathcal{O}_{U_d}(m)
\to \mathcal{O}_{U_d}(nd + m)$，即
(\ref{constructions-equation-multiply}) 中的映射，都是同构。特别地，有
$\mathcal{O}_{U_d}(nd) \cong \mathcal{O}_{U_d}(d)^{\otimes n}$.
整体截面上的分次环映射 (\ref{constructions-equation-global-sections}) 与到 $U_d$ 的限制
合在一起给出分次环同态
\begin{equation}
\label{constructions-equation-psi-d}
\psi^d : S^{(d)} \longrightarrow \Gamma_*(U_d, \mathcal{O}_{U_d}(d)).
\end{equation}
关于记号 $S^{(d)}$，见代数章第 \ref{algebra-section-graded} 节。
关于记号 $\Gamma_*$，见模章定义 \ref{modules-definition-gamma-star}。
此外，由于 $U_d$ 被诸开集 $D_{+}(f)$（$f \in S_d$）覆盖，
可知 $\mathcal{O}_{U_d}(d)$ 由下列映射像中的截面整体生成：
$\psi^d_1 : S^{(d)}_1 = S_d \to \Gamma(U_d, \mathcal{O}_{U_d}(d))$，
见模章定义 \ref{modules-definition-globally-generated}。

\medskip\noindent
设 $Y$ 为概形，并设 $\varphi : Y \to X$ 为概形态射。
假设像 $\varphi(Y)$ 包含于开子概形 $U_d$ 中，
后者是 $X$ 的开子概形。
由模章定义 \ref{modules-definition-gamma-star} 后的讨论，
得到分次环同态
$$
\Gamma_*(U_d, \mathcal{O}_{U_d}(d))
\longrightarrow
\Gamma_*(Y, \varphi^*\mathcal{O}_X(d)).
$$
它与 $\psi^d$ 的复合给出分次环同态
\begin{equation}
\label{constructions-equation-psi-phi-d}
\psi_\varphi^d :
S^{(d)}
\longrightarrow
\Gamma_*(Y, \varphi^*\mathcal{O}_X(d))
\end{equation}
并且可逆层 $\varphi^*\mathcal{O}_X(d)$ 由下列映射像中的截面整体生成：
\par\noindent
$(S^{(d)})_1 = S_d \to \Gamma(Y, \varphi^*\mathcal{O}_X(d))$.

\begin{lemma}
\label{constructions-lemma-converse-construction}
设 $S$ 为分次环，并置 $X = \text{Proj}(S)$。
设 $d \geq 1$，且 $U_d \subset X$ 如上。
设 $Y$ 为概形。
设 $\mathcal{L}$ 为 $Y$ 上的可逆层。
设 $\psi : S^{(d)} \to \Gamma_*(Y, \mathcal{L})$ 为分次环同态，
并且 $\mathcal{L}$ 由下列映射像中的截面生成：
$\psi|_{S_d} : S_d \to \Gamma(Y, \mathcal{L})$.
则存在态射 $\varphi : Y \to X$，满足 $\varphi(Y) \subset U_d$，
并存在同构 $\alpha : \varphi^*\mathcal{O}_{U_d}(d) \to \mathcal{L}$，
使得 $\psi_\varphi^d$ 与 $\psi$ 经由 $\alpha$ 一致：
$$
\xymatrix{
\Gamma_*(Y, \mathcal{L}) &
\Gamma_*(Y, \varphi^*\mathcal{O}_{U_d}(d)) \ar[l]^-\alpha &
\Gamma_*(U_d, \mathcal{O}_{U_d}(d)) \ar[l]^-{\varphi^*} \\
S^{(d)} \ar[u]^\psi & &
S^{(d)} \ar[u]^{\psi^d} \ar[ul]^{\psi^d_\varphi} \ar[ll]_{\text{id}}
}
$$
该图交换。此外，二元组 $(\varphi, \alpha)$ 是唯一的。
\end{lemma}

\begin{proof}
取 $f \in S_d$。记 $s = \psi(f) \in \Gamma(Y, \mathcal{L})$。
在开集 $Y_s$ 上，$s$ 不消失；乘以 $s$ 诱导同构
$\mathcal{O}_{Y_s} \to \mathcal{L}|_{Y_s}$，
见模章引理 \ref{modules-lemma-s-open}。我们将此映射的逆记为
$x \mapsto x/s$，对 $\mathcal{L}$ 的幂也采用类似记号。利用这一点，
定义环同态 $\psi_{(f)} : S_{(f)} \to \Gamma(Y_s, \mathcal{O}_Y)$，
它将分式 $a/f^n$ 映到 $\psi(a)/s^n$。
由概形章引理 \ref{schemes-lemma-morphism-into-affine}，
这对应于态射
$\varphi_f : Y_s \to \Spec(S_{(f)}) = D_{+}(f)$。
再引入同构
$\alpha_f : \varphi_f^*\mathcal{O}_{D_{+}(f)}(d) \to \mathcal{L}|_{Y_s}$，
它将平凡化截面 $f$ 在 $D_{+}(f)$ 上的拉回映到
平凡化截面 $s$ 在 $Y_s$ 上的限制。
按此选择，引理中的图在将 $Y$ 替换为 $Y_s$、
将 $\varphi$ 替换为 $\varphi_f$、将 $\alpha$ 替换为 $\alpha_f$ 后交换；
验证从略。

\medskip\noindent
设 $f' \in S_d$ 是第二个元素，并记
$s' = \psi(f') \in \Gamma(Y, \mathcal{L})$。于是
$Y_s \cap Y_{s'} = Y_{ss'}$，且类似地，
$D_{+}(f) \cap D_{+}(f') = D_{+}(ff')$。
在引理 \ref{constructions-lemma-ample-on-proj} 中已经看到，
$D_{+}(f') \cap D_{+}(f)$ 正是 $D_{+}(f)$ 中
$\mathcal{O}_X(d)$ 的由 $f'$ 所定义的截面不消失的点集。因此
$\varphi_f^{-1}(D_{+}(f') \cap D_{+}(f)) = Y_s \cap Y_{s'}
= \varphi_{f'}^{-1}(D_{+}(f') \cap D_{+}(f))$.
在 $D_{+}(f) \cap D_{+}(f')$ 上，分式 $f/f'$ 是结构层的可逆截面，
其逆为 $f'/f$。注意
$\psi_{(f')}(f/f') = \psi(f)/s' = s/s'$，且
$\psi_{(f)}(f'/f) = \psi(f')/s = s'/s$。我们断言，存在唯一的环同态
$S_{(ff')} \to \Gamma(Y_{ss'}, \mathcal{O}_Y)$，
它使下图交换：
$$
\xymatrix{
\Gamma(Y_s, \mathcal{O}_Y) \ar[r] &
\Gamma(Y_{ss'}, \mathcal{O}_Y) &
\Gamma(Y_{s, '} \mathcal{O}_Y) \ar[l]\\
S_{(f)} \ar[r] \ar[u]^{\psi_{(f)}} &
S_{(ff')} \ar[u] &
S_{(f')} \ar[l] \ar[u]^{\psi_{(f')}}
}
$$
该同态存在，因为可以使用规则
$x/(ff')^n \mapsto \psi(x)/(ss')^n$；由上面的公式可知这一规则“成立”。
由于 $\text{Proj}(S)$ 是分离的，唯一性由引理
\ref{constructions-lemma-proj-separated} 及其证明得出。这表明态射
$\varphi_f$ 与 $\varphi_{f'}$ 在 $Y_s \cap Y_{s'}$ 上一致。
$\alpha_f$ 与 $\alpha_{f'}$ 的限制在 $Y_s \cap Y_{s'}$ 上一致，
因为正则函数 $s/s'$ 与 $\psi_{(f')}(f/f')$ 一致。这证明态射
$\psi_f$ 可黏合为从 $Y$ 到 $U_d \subset X$ 的整体态射，且
各映射 $\alpha_f$ 黏合为满足引理条件的同构。

\medskip\noindent
还需证明二元组 $(\varphi, \alpha)$ 唯一。
设 $(\varphi', \alpha')$ 是第二个这样的二元组。
设 $f \in S_d$。由引理中的图交换可知，$D_{+}(f)$ 在 $\varphi$ 与
$\varphi'$ 下的逆像都等于 $Y_{\psi(f)}$。由于诸开集 $D_{+}(f)$ 构成
$X$ 的拓扑的一组基，而 $X$ 是清醒拓扑空间
（见概形章引理 \ref{schemes-lemma-scheme-sober}），
这意味着映射 $\varphi$ 与 $\varphi'$ 在底拓扑空间上相同。用
$s = \psi(f)$ 将可逆层 $\mathcal{L}$ 在 $Y_{\psi(f)}$ 上平凡化。
由这些图的交换性，有
$\alpha^{\otimes n}(\psi^d_{\varphi}(x)) =
\psi(x) = (\alpha')^{\otimes n}(\psi^d_{\varphi'}(x))$
对所有 $x \in S_{nd}$ 成立。由 $\psi^d_{\varphi}$ 与
$\psi^d_{\varphi'}$ 的构造，有
$\psi^d_{\varphi}(x) = \varphi^\sharp(x/f^n) \psi^d_{\varphi}(f^n)$
在 $Y_{\psi(f)}$ 上成立，对 $\psi^d_{\varphi'}$ 也类似。由引理中的图交换可得
$\varphi^\sharp(x/f^n) = (\varphi')^\sharp(x/f^n)$.
这证明 $\varphi$ 与 $\varphi'$ 诱导从 $Y_{\psi(f)}$ 到仿射概形
$D_{+}(f) = \Spec(S_{(f)})$ 的同一个态射。
因此 $\varphi$ 与 $\varphi'$ 作为态射相同。
最后，还需证明，在给定 $\varphi$ 后，引理中图的交换性唯一确定
$\alpha$。验证从略。
\end{proof}

\noindent
我们继续引理前的讨论。
设 $S$ 为分次环。
设 $Y$ 为概形。我们考虑{\it 三元组}
$(d, \mathcal{L}, \psi)$，其中
\begin{enumerate}
\item $d \geq 1$ 是整数；
\item $\mathcal{L}$ 是可逆 $\mathcal{O}_Y$-模；并且
\item $\psi : S^{(d)} \to \Gamma_*(Y, \mathcal{L})$ 是分次环同态，
且 $\mathcal{L}$ 由整体截面 $\psi(f)$（其中 $f \in S_d$）生成。
\end{enumerate}
给定态射 $h : Y' \to Y$ 与三元组
$(d, \mathcal{L}, \psi)$（它位于 $Y$ 上），可将它拉回为三元组
$(d, h^*\mathcal{L}, h^* \circ \psi)$。
给定两个三元组 $(d, \mathcal{L}, \psi)$ 与
$(d, \mathcal{L}', \psi')$，它们具有相同的整数 $d$；若存在同构
$\beta : \mathcal{L} \to \mathcal{L}'$，使得作为分次环同态有
$\beta \circ \psi = \psi'$，其中这些同态为
$S^{(d)} \to \Gamma_*(Y, \mathcal{L}')$，则称它们{\it 严格等价}。

\medskip\noindent
对每个整数 $d \geq 1$，定义
\begin{eqnarray*}
F_d : \Sch^{opp} & \longrightarrow & \textit{Sets}, \\
Y & \longmapsto &
\{\text{三元组的严格等价类： }
(d, \mathcal{L}, \psi)
\text{ 如上}\}
\end{eqnarray*}
其拉回按上面所述定义。

\begin{lemma}
\label{constructions-lemma-proj-functor-strict}
设 $S$ 为分次环。
设 $X = \text{Proj}(S)$。
开子概形 $U_d \subset X$（\ref{constructions-equation-Ud}）表示函子
$F_d$，且上面定义的三元组 $(d, \mathcal{O}_{U_d}(d), \psi^d)$
是泛族（见概形章第 \ref{schemes-section-representable} 节）。
\end{lemma}

\begin{proof}
这是引理 \ref{constructions-lemma-converse-construction} 的重新表述。
\end{proof}

\begin{lemma}
\label{constructions-lemma-apply}
设 $S$ 为分次环，它作为 $S_0$-代数由 $S_1$ 中的元素生成。
在这种情形下，概形 $X = \text{Proj}(S)$ 表示如下函子：它将概形
$Y$ 映到二元组 $(\mathcal{L}, \psi)$ 的集合，其中
\begin{enumerate}
\item $\mathcal{L}$ 是可逆 $\mathcal{O}_Y$-模；并且
\item $\psi : S \to \Gamma_*(Y, \mathcal{L})$ 是分次环同态，
且 $\mathcal{L}$ 由整体截面 $\psi(f)$（其中 $f \in S_1$）生成，
\end{enumerate}
并按上述严格等价关系取等价类。
\end{lemma}

\begin{proof}
在本引理的假设下有 $X = U_1$，而本引理正是上面引理
\ref{constructions-lemma-proj-functor-strict} 的重新表述。
\end{proof}

\noindent
本节最后讨论一般分次环的 $\text{Proj}(S)$ 所对应的函子，其中该环记为 $S$。
建议读者跳过本节余下部分。

\medskip\noindent
固定任意分次环 $S$。设 $T$ 为概形。
称两个三元组 $(d, \mathcal{L}, \psi)$ 与
$(d', \mathcal{L}', \psi')$（它们在 $T$ 上，并允许整数
$d$、$d'$ 不同）{\it 等价}，如果存在同构
$\beta : \mathcal{L}^{\otimes d'} \to (\mathcal{L}')^{\otimes d}$，
它是 $T$ 上可逆层的同构，
使得 $\beta \circ \psi|_{S^{(dd')}}$ 与 $\psi'|_{S^{(dd')}}$ 作为分次环同态
$S^{(dd')} \to \Gamma_*(Y, (\mathcal{L}')^{\otimes dd'})$ 一致。

\begin{lemma}
\label{constructions-lemma-equivalent}
设 $S$ 为分次环。置 $X = \text{Proj}(S)$。设 $T$ 为概形。
设 $(d, \mathcal{L}, \psi)$ 与 $(d', \mathcal{L}', \psi')$
为 $T$ 上的两个三元组。下列条件等价：
\begin{enumerate}
\item 设 $n = \text{lcm}(d, d')$。写成 $n = ad = a'd'$。存在同构
$\beta : \mathcal{L}^{\otimes a} \to (\mathcal{L}')^{\otimes a'}$
使得 $\beta \circ \psi|_{S^{(n)}}$ 与 $\psi'|_{S^{(n)}}$ 作为分次环同态
$S^{(n)} \to \Gamma_*(Y, (\mathcal{L}')^{\otimes n})$ 一致。
\item 三元组 $(d, \mathcal{L}, \psi)$ 与 $(d', \mathcal{L}', \psi')$
等价。
\item 对某个正整数 $n = ad = a'd'$，存在同构
$\beta : \mathcal{L}^{\otimes a} \to (\mathcal{L}')^{\otimes a'}$
使得 $\beta \circ \psi|_{S^{(n)}}$ 与 $\psi'|_{S^{(n)}}$ 作为分次环同态
$S^{(n)} \to \Gamma_*(Y, (\mathcal{L}')^{\otimes n})$ 一致。
\item 态射 $\varphi : T \to X$ 与 $\varphi' : T \to X$
（它们分别与 $(d, \mathcal{L}, \psi)$ 及 $(d', \mathcal{L}', \psi')$ 相伴）相等。
\end{enumerate}
\end{lemma}

\begin{proof}
显然，限制到可被更多整数整除的次数并取可逆层的幂，可得
(1) 蕴含 (2)，且 (2) 蕴含 (3)。
由引理 \ref{constructions-lemma-converse-construction} 中的唯一性陈述，
(3) 也蕴含 (4)。
因此只需证明 (4) 蕴含 (1)。假设 (4)，即 $\varphi = \varphi'$。
注意，这意味着可写成
$\mathcal{L} = \varphi^*\mathcal{O}_X(d)$，且
$\mathcal{L}' = \varphi^*\mathcal{O}_X(d')$.
此外，经由这些识别，分次环同态 $\psi$ 与 $\psi'$ 对应于典范分次环同态
$$
S \longrightarrow \bigoplus\nolimits_{n \geq 0} \Gamma(X, \mathcal{O}_X(n))
$$
到 $S^{(d)}$ 与 $S^{(d')}$ 的限制，再与 $\varphi$ 的拉回复合
（再次使用引理 \ref{constructions-lemma-converse-construction}）。因此取 $\beta$ 为同构
$$
(\varphi^*\mathcal{O}_X(d))^{\otimes a} =
\varphi^*\mathcal{O}_X(n) =
(\varphi^*\mathcal{O}_X(d'))^{\otimes a'}
$$
即可。
\end{proof}

\noindent
设 $S$ 为分次环。
设 $X = \text{Proj}(S)$。
在开子概形 $U_d \subset X = \text{Proj}(S)$
（\ref{constructions-equation-Ud}）上，有三元组
$(d, \mathcal{O}_{U_d}(d), \psi^d)$。显然，若 $d | d'$，则三元组
$(d, \mathcal{O}_{U_d}(d), \psi^d)$ 与
$(d', \mathcal{O}_{U_{d'}}(d'), \psi^{d'})$ 限制到开集 $U_d$
（它是 $U_{d'}$ 的子集）后等价。
结合引理 \ref{constructions-lemma-converse-construction}，这表明态射 $Y \to X$
大致对应于 $Y$ 上三元组的等价类。这并不完全正确，因为若 $Y$
不是拟紧的，可能不存在一个处处适用的单一三元组。
因此，定义相应函子时必须稍加谨慎。

\medskip\noindent
以下是一种可行的处理方式。假设 $d' = ad$。
考虑函子变换 $F_d \to F_{d'}$，它将三元组
$(d, \mathcal{L}, \psi)$（位于 $T$ 上）映为三元组
$(d', \mathcal{L}^{\otimes a}, \psi|_{S^{(d')}})$。
引理 \ref{constructions-lemma-equivalent} 的一个推论是：变换 $F_d \to F_{d'}$ 是单射！
对拟紧概形 $T$，定义
$$
F(T) = \bigcup\nolimits_{d \in \mathbf{N}} F_d(T)
$$
其中转移映射如上所述。显然，这在拟紧概形范畴上定义了一个取值于集合的
反变函子。对一般概形 $T$，定义
$$
F(T)
=
\lim_{V \subset T\text{ 拟紧开}} F(V).
$$
换言之，一个元素 $\xi$（属于 $F(T)$）对应于一族相容的选择
$\xi_V \in F(V)$，其中 $V$ 遍历 $T$ 的拟紧开集。
拉回映射 $F(T) \to F(T')$（对应于概形态射 $T' \to T$）的定义从略。
由此定义了函子
\begin{eqnarray*}
F : \Sch^{opp} & \longrightarrow & \textit{Sets}
\end{eqnarray*}

\begin{lemma}
\label{constructions-lemma-proj-functor}
设 $S$ 为分次环。
设 $X = \text{Proj}(S)$。
上面定义的函子 $F$ 由概形 $X$ 表示。
\end{lemma}

\begin{proof}
上面已看到，函子 $F_d$ 对应于开子概形 $U_d \subset X$。此外，
上面定义的函子变换 $F_d \to F_{d'}$（若 $d | d'$）
对应于包含态射 $U_d \to U_{d'}$（见上面的讨论）。
因此，要证明 $F$ 由 $X$ 表示，只需证明：态射 $T \to X$ 在概形
$T$ 拟紧时其像落在某个 $U_d$ 中；并且对一般概形 $T$ 有
$$
\Mor(T, X)
=
\lim_{V \subset T\text{ 拟紧开}} \Mor(V, X).
$$
这些验证从略。
\end{proof}







\section{射影空间}
\label{constructions-section-projective-space}

\noindent
射影空间是代数几何所研究的基本对象之一。本节只给出它作为多项式环之
$\text{Proj}$ 的构造。稍后我们将发现它的许多优美性质。

\begin{lemma}
\label{constructions-lemma-projective-space}
设 $S = \mathbf{Z}[T_0, \ldots, T_n]$，其中 $\deg(T_i) = 1$。
概形
$$
\mathbf{P}^n_{\mathbf{Z}} = \text{Proj}(S)
$$
表示如下函子：它将概形 $Y$ 映到二元组
$(\mathcal{L}, (s_0, \ldots, s_n))$ 的集合，其中
\begin{enumerate}
\item $\mathcal{L}$ 是可逆 $\mathcal{O}_Y$-模；并且
\item $s_0, \ldots, s_n$ 是 $\mathcal{L}$ 的整体截面，且生成
$\mathcal{L}$，
\end{enumerate}
并按下列等价关系取等价类：
$(\mathcal{L}, (s_0, \ldots, s_n)) \sim
(\mathcal{N}, (t_0, \ldots, t_n))$ $\Leftrightarrow$ 存在同构
$\beta : \mathcal{L} \to \mathcal{N}$，使得
$\beta(s_i) = t_i$ 对 $i = 0, \ldots, n$ 成立。
\end{lemma}

\begin{proof}
这是上面引理 \ref{constructions-lemma-apply} 的特殊情形。
事实上，对任意分次环 $A$ 有
\begin{eqnarray*}
\Mor_{\text{分次环}}(\mathbf{Z}[T_0, \ldots, T_n], A)
& = &
A_1 \times \ldots \times A_1 \\
\psi & \mapsto & (\psi(T_0), \ldots, \psi(T_n))
\end{eqnarray*}
且次数 $1$ 部分在 $\Gamma_*(Y, \mathcal{L})$ 中正是
$\Gamma(Y, \mathcal{L})$。
\end{proof}

\begin{definition}
\label{constructions-definition-projective-space}
概形
$\mathbf{P}^n_{\mathbf{Z}} = \text{Proj}(\mathbf{Z}[T_0, \ldots, T_n])$
称为{\it 射影 $n$ 维空间（在 $\mathbf{Z}$ 上）}。
它的基变换 $\mathbf{P}^n_S$（到概形 $S$）称为
{\it 射影 $n$ 维空间（在 $S$ 上）}。若 $R$ 为环，则到
$\Spec(R)$ 的基变换记为 $\mathbf{P}^n_R$，并称为
{\it 射影 $n$ 维空间（在 $R$ 上）}。
\end{definition}

\noindent
给定概形 $Y$（它在 $S$ 上）以及引理 \ref{constructions-lemma-projective-space} 中的二元组
$(\mathcal{L}, (s_0, \ldots, s_n))$，其诱导的到 $\mathbf{P}^n_S$ 的态射记为
$$
\varphi_{(\mathcal{L}, (s_0, \ldots, s_n))} :
Y \longrightarrow \mathbf{P}^n_S
$$
这是有意义的，因为该二元组定义了到 $\mathbf{P}^n_{\mathbf{Z}}$ 的态射，
而我们已有到 $S$ 的结构态射；二者合在一起给出到
$\mathbf{P}^n_S = \mathbf{P}^n_{\mathbf{Z}} \times S$ 的态射。
借助引理 \ref{constructions-lemma-relative-proj-base-change}，
可将 $\mathbf{P}^n_S$ 识别为分次 $\mathcal{O}_S$-代数
$\mathcal{O}_S[T_0, \ldots, T_n]$ 的相对 Proj，并可将
$\varphi_{(\mathcal{L}, (s_0, \ldots, s_n))}$ 刻画为唯一的 $S$-态射，满足
$$
\mathcal{L} =
\varphi_{(\mathcal{L}, (s_0, \ldots, s_n))}^*\mathcal{O}_{\mathbf{P}^n_S}(1)
\quad
\text{且}
\quad
s_i = \varphi_{(\mathcal{L}, (s_0, \ldots, s_n))}^*T_i
$$
这里把 $T_i$ 看作 $\mathcal{O}_{\mathbf{P}^n_S}(1)$ 的整体截面；
所用映射为 $\psi$，见引理 \ref{constructions-lemma-glue-relative-proj-twists}；细节从略。

\begin{lemma}
\label{constructions-lemma-standard-covering-projective-space}
射影 $n$ 维空间（在 $\mathbf{Z}$ 上）由 $n + 1$ 个标准开集覆盖：
$$
\mathbf{P}^n_{\mathbf{Z}} =
\bigcup\nolimits_{i = 0, \ldots, n} D_{+}(T_i)
$$
其中每个 $D_{+}(T_i)$ 都同构于 $\mathbf{A}^n_{\mathbf{Z}}$，即
仿射 $n$ 维空间（在 $\mathbf{Z}$ 上）。
\end{lemma}

\begin{proof}
这是因为
$\mathbf{Z}[T_0, \ldots, T_n]_{+} = (T_0, \ldots, T_n)$，并且
并且
$$
\Spec
\left(
\mathbf{Z}
\left[\frac{T_0}{T_i}, \ldots, \frac{T_n}{T_i}
\right]
\right)
\cong
\mathbf{A}^n_{\mathbf{Z}}
$$
以显然的方式成立。
\end{proof}

\begin{lemma}
\label{constructions-lemma-projective-space-separated}
设 $S$ 为概形。结构态射 $\mathbf{P}^n_S \to S$ 具有下列性质：
\begin{enumerate}
\item 分离；
\item 拟紧；
\item 满足赋值判据的存在性部分与唯一性部分；并且
\item 万有闭。
\end{enumerate}
\end{lemma}

\begin{proof}
这些性质在基变换下都稳定（对后两个性质这是显然的；对前两个性质，见概形章引理
\ref{schemes-lemma-separated-permanence} 与
\ref{schemes-lemma-quasi-compact-preserved-base-change}).
因此只需对态射
$\mathbf{P}^n_{\mathbf{Z}} \to \Spec(\mathbf{Z})$.
证明这些性质。分离性是引理 \ref{constructions-lemma-proj-separated}。拟紧性由引理
\ref{constructions-lemma-standard-covering-projective-space} 得出。
赋值判据的存在性与唯一性由引理 \ref{constructions-lemma-proj-valuative-criterion} 得出。
万有闭性由以上结果与概形章命题
\ref{schemes-proposition-characterize-universally-closed} 得出。
\end{proof}

\begin{remark}
\label{constructions-remark-missing-finite-type}
上述性质列表中还缺少什么？当然还缺少有限型性质。这里不列出它，
是因为到目前为止我们还没有定义有限型的概念。（另一个缺少的性质是
“光滑性”。相信读者还能想到许多其他性质。）
\end{remark}

\begin{lemma}[Segre 嵌入]
\label{constructions-lemma-segre-embedding}
设 $S$ 为概形。存在闭浸入
$$
\mathbf{P}^n_S \times_S \mathbf{P}^m_S
\longrightarrow
\mathbf{P}^{nm + n + m}_S
$$
称为{\it Segre 嵌入}。
\end{lemma}

\begin{proof}
只需在 $S = \Spec(\mathbf{Z})$ 时证明。于是略去下标 $S$，
在绝对情形下工作。写成
$\mathbf{P}^n = \text{Proj}(\mathbf{Z}[X_0, \ldots, X_n])$，
$\mathbf{P}^m = \text{Proj}(\mathbf{Z}[Y_0, \ldots, Y_m])$,
以及
$\mathbf{P}^{nm + n + m} =
\text{Proj}(\mathbf{Z}[Z_0, \ldots, Z_{nm + n + m}])$.
为了映到 $\mathbf{P}^{nm + n + m}$，必须在左边写下一个可逆层
$\mathcal{L}$，以及生成它的 $(n + 1)(m + 1)$ 个截面 $s_i$。
见引理 \ref{constructions-lemma-projective-space}。
所取的可逆层为
$$
\mathcal{L} =
\text{pr}_1^*\mathcal{O}_{\mathbf{P}^n}(1)
\otimes
\text{pr}_2^*\mathcal{O}_{\mathbf{P}^m}(1)
$$
所取的截面为
$$
s_0 = X_0Y_0, \ s_1 = X_1Y_0, \ldots, \ s_n = X_nY_0,
\ s_{n + 1} = X_0Y_1, \ldots, \ s_{nm + n + m} = X_nY_m.
$$
这些截面生成 $\mathcal{L}$，因为截面 $X_i$ 生成
$\mathcal{O}_{\mathbf{P}^n}(1)$，而截面 $Y_j$ 生成
$\mathcal{O}_{\mathbf{P}^m}(1)$。诱导态射 $\varphi$ 满足
$$
\varphi^{-1}(D_{+}(Z_{i + (n + 1)j})) = D_{+}(X_i) \times D_{+}(Y_j).
$$
因此它是仿射态射。在情形 $(i, j) = (0, 0)$ 下，相应的环同态为
$$
\mathbf{Z}[Z_1/Z_0, \ldots, Z_{nm + n + m}/Z_0]
\longrightarrow
\mathbf{Z}[X_1/X_0, \ldots, X_n/X_0, Y_1/Y_0, \ldots, Y_n/Y_0]
$$
它将 $Z_i/Z_0$ 映到元素 $X_i/X_0$（对 $i \leq n$），并将元素
$Z_{(n + 1)j}/Z_0$ 映到元素 $Y_j/Y_0$。因此它是满射。
对其他仿射开子集也可作类似论证。因此态射 $\varphi$ 是闭浸入
（见概形章引理 \ref{schemes-lemma-closed-local-target} 与例
\ref{schemes-example-closed-immersion-affines}）。
\end{proof}

\noindent
下面两个引理是稍后更一般结果的特殊情形，但现在在这里直接证明它们或许更为合适。
第一个是除子章引理 \ref{divisors-lemma-closed-subscheme-proj} 的特殊情形。

\begin{lemma}
\label{constructions-lemma-closed-in-projective-space}
设 $R$ 为环。设 $Z \subset \mathbf{P}^n_R$ 为闭子概形。令
$$
I_d = \Ker\left(
R[T_0, \ldots, T_n]_d
\longrightarrow
\Gamma(Z, \mathcal{O}_{\mathbf{P}^n_R}(d)|_Z)\right)
$$
则 $I = \bigoplus I_d \subset R[T_0, \ldots, T_n]$ 是分次理想，并且
$Z = \text{Proj}(R[T_0, \ldots, T_n]/I)$。
\end{lemma}

\begin{proof}
显然 $I$ 是分次理想。
置 $Z' = \text{Proj}(R[T_0, \ldots, T_n]/I)$。
由引理 \ref{constructions-lemma-surjective-graded-rings-generated-degree-1-map-proj} 可知，
$Z'$ 是 $\mathbf{P}^n_R$ 的闭子概形。
要验证等式 $Z = Z'$，只需在标准仿射开集 $D_{+}(T_i)$ 上验证。
通过重编号齐次坐标，可假设 $i = 0$。设 $Z \cap D_{+}(T_0)$ 与
$Z' \cap D_{+}(T_0)$ 分别由理想 $J$ 与 $J'$ 截出；这些理想位于
$R[T_1/T_0, \ldots, T_n/T_0]$ 中。则 $J'$ 是由元素 $F/T_0^{\deg(F)}$ 生成的理想，
其中 $F \in I$ 为齐次元。
假设 $F \in I$ 的次数为 $d$。由于 $F$ 作为
$\mathcal{O}_{\mathbf{P}^n_R}(d)$ 限制到 $Z$ 上的截面消失，可知
$F/T_0^d$ 是 $J$ 的元素。因此 $J' \subset J$。

\medskip\noindent
反之，假设 $f \in J$。若 $f$ 的总次数为
$d$（视为 $T_1/T_0, \ldots, T_n/T_0$ 的多项式），则可写成
$f = F/T_0^d$，其中 $F \in R[T_0, \ldots, T_n]_d$。
取 $i \in \{1, \ldots, n\}$。则 $Z \cap D_{+}(T_i)$ 由某个理想
$J_i \subset R[T_0/T_i, \ldots, T_n/T_i]$ 截出。此外，
$$
J \cdot
R\left[
\frac{T_1}{T_0}, \ldots, \frac{T_n}{T_0},
\frac{T_0}{T_i}, \ldots, \frac{T_n}{T_i}
\right]
=
J_i \cdot
R\left[
\frac{T_1}{T_0}, \ldots, \frac{T_n}{T_0},
\frac{T_0}{T_i}, \ldots, \frac{T_n}{T_i}
\right]
$$
左边是 $J$ 关于元素 $T_i/T_0$ 的局部化，右边是 $J_i$ 关于元素
$T_0/T_i$ 的局部化。因此，
$T_0^{d_i}F/T_i^{d + d_i}$ 是 $J_i$ 的元素，其中 $d_i$ 充分大。这证明
$T_0^{\max(d_i)}F$ 是 $I$ 的元素，因为它在每个标准仿射开集
$D_{+}(T_i)$ 上的限制都在闭子概形 $Z \cap D_{+}(T_i)$ 上消失。
因此 $f \in J'$，从而得到所需的 $J \subset J'$。
\end{proof}

\noindent
下面的引理是性质章引理 \ref{properties-lemma-ample-quasi-coherent} 或
\ref{properties-lemma-proj-quasi-coherent} 这一更一般结果的特殊情形。

\begin{lemma}
\label{constructions-lemma-quasi-coherent-projective-space}
设 $R$ 为环。
设 $\mathcal{F}$ 为 $\mathbf{P}^n_R$ 上的拟凝聚层。
对 $d \geq 0$，置
$$
M_d
=
\Gamma(\mathbf{P}^n_R,
\mathcal{F} \otimes_{\mathcal{O}_{\mathbf{P}^n_R}}
\mathcal{O}_{\mathbf{P}^n_R}(d))
=
\Gamma(\mathbf{P}^n_R, \mathcal{F}(d))
$$
则 $M = \bigoplus_{d \geq 0} M_d$ 是分次 $R[T_0, \ldots, R_n]$-模，
并且存在典范同构 $\mathcal{F} = \widetilde{M}$。
\end{lemma}

\begin{proof}
乘法映射
$$
R[T_0, \ldots, R_n]_e \times M_d \longrightarrow M_{d + e}
$$
来自自然同构
$$
\mathcal{O}_{\mathbf{P}^n_R}(e)
\otimes_{\mathcal{O}_{\mathbf{P}^n_R}}
\mathcal{F}(d)
\longrightarrow
\mathcal{F}(e + d)
$$
见方程 (\ref{constructions-equation-global-sections-module})。来构造映射
$c : \widetilde{M} \to \mathcal{F}$。在每个标准仿射开集
$U_i = D_{+}(T_i)$ 上，可见 $\Gamma(U_i, \widetilde{M}) = (M[1/T_i])_0$，
其中下标 ${}_0$ 表示次数 $0$ 部分。其中的元素可写成
$m/T_i^d$，其中 $m \in M_d$。由于 $T_i$ 是
$\mathcal{O}(1)$ 的一个生成元（在 $U_i$ 上），总可以写成
$m|_{U_i} = m_i \otimes T_i^d$，其中
$m_i \in \Gamma(U_i, \mathcal{F})$ 是唯一截面。因此自然的猜测是
$c(m/T_i^d) = m_i$。一个略去的小论证表明，这会给出良定义映射
$c : \widetilde{M} \to \mathcal{F}$，只要能证明
$$
(T_i/T_j)^d m_i|_{U_i \cap U_j} = m_j|_{U_i \cap U_j}
$$
在 $M[1/T_iT_j]$ 中成立。
而这是显然的，因为在交叠上，生成元 $T_i$ 与 $T_j$
（它们属于 $\mathcal{O}(1)$）相差可逆函数 $T_i/T_j$。

\medskip\noindent
$c$ 的单射性。可在仿射开集 $U_i$ 上检验单射性。设
$i \in \{0, \ldots, n\}$，并设 $s$ 为元素
$s = m/T_i^d \in \Gamma(U_i, \widetilde{M})$，满足
$c(m/T_i^d) = 0$。由上面对 $c$ 的描述，这意味着
$m_i = 0$，故 $m|_{U_i} = 0$。因此 $T_i^em = 0$ 在 $M$ 中成立，
其中 $e$ 为某个整数。于是如所需，$s = m/T_i^d = T_i^e/T_i^{e + d} = 0$。

\medskip\noindent
$c$ 的满射性。可在仿射开集 $U_i$ 上检验满射性。通过重编号，
只需在 $U_0$ 上检验。设 $s \in \mathcal{F}(U_0)$。
写成 $\mathcal{F}|_{U_i} = \widetilde{N_i}$；这里所用的是某个
$R[T_0/T_i, \ldots, T_0/T_i]$-模 $N_i$。这是可能的，因为
$\mathcal{F}$ 拟凝聚。于是 $s$ 对应于元素 $x \in N_0$。此时有
$$
(N_i)_{T_j/T_i} \cong (N_j)_{T_i/T_j}
$$
（下标表示“关于该元的主局部化”），它们是下列环上的模：
$$
R\left[
\frac{T_0}{T_i}, \ldots, \frac{T_n}{T_i},
\frac{T_0}{T_j}, \ldots, \frac{T_n}{T_j}
\right].
$$
这意味着，对某个充分大的整数 $d$，存在元素
$s_i \in N_i$（$i = 1, \ldots, n$），使得
$$
s = (T_i/T_0)^d s_i
$$
在 $U_0 \cap U_i$ 上成立。接下来考察差
$$
t_{ij} = s_i - (T_j/T_i)^d s_j
$$
在 $U_i \cap U_j$ 上（$0 < i < j$）。由对 $s_i$ 的选择可知
$t_{ij}|_{U_0 \cap U_i \cap U_j} = 0$。因此存在充分大的整数
$e$，使得 $(T_0/T_i)^et_{ij} = 0$。置 $s_i' = (T_0/T_i)^es_i$，
并置 $s_0' = s$。于是有
$$
s_a' = (T_b/T_a)^{e + d} s_b'
$$
在 $U_a \cap U_b$ 上对所有 $a, b$ 成立。这恰好是诸元素 $s'_a$
黏合为整体截面 $m \in \Gamma(\mathbf{P}^n_R, \mathcal{F}(e + d))$ 的条件。
此外，由构造有 $c(m/T_0^{e + d}) = s$。因此 $c$ 是满射，证毕。
\end{proof}

\begin{lemma}
\label{constructions-lemma-globally-generated-omega-twist-1}
设 $X$ 为概形。设 $\mathcal{L}$ 为可逆层，并设
$s_0, \ldots, s_n$ 为生成 $\mathcal{L}$ 的整体截面。设
$\mathcal{F}$ 为诱导映射 $\mathcal{O}_X^{\oplus n + 1} \to \mathcal{L}$ 的核。
则 $\mathcal{F} \otimes \mathcal{L}$ 整体生成。
\end{lemma}

\begin{proof}
事实上，当 $X$ 是任意局部环层空间时，结论仍成立。
层 $\mathcal{F}$ 是有限局部自由 $\mathcal{O}_X$-模，秩为 $n$。诸元素
$$
s_{ij} = (0, \ldots, 0, s_j, 0, \ldots, 0, -s_i, 0, \ldots, 0)
\in \Gamma(X, \mathcal{L}^{\oplus n + 1})
$$
中，$s_j$ 位于第 $i$ 个位置，$-s_i$ 位于第 $j$ 个位置；它们在
$\mathcal{L}^{\otimes 2}$ 中映为零。因此
$s_{ij} \in \Gamma(X, \mathcal{F} \otimes_{\mathcal{O}_X} \mathcal{L})$.
局部计算表明，这些截面生成 $\mathcal{F} \otimes \mathcal{L}$。

\medskip\noindent
另一证明。考虑态射 $\varphi : X \to \mathbf{P}^n_\mathbf{Z}$，它与二元组
$(\mathcal{L}, (s_0, \ldots, s_n))$ 相伴。由于 $\mathcal{O}(1)$ 的拉回是
$\mathcal{L}$，且 $T_i$ 的拉回是 $s_i$，只需在
$\mathbf{P}^n_\mathbf{Z}$ 的情形证明本引理。在这种情形下，层
$\mathcal{F}$ 对应于分次 $S = \mathbf{Z}[T_0, \ldots, T_n]$-模
$M$，它适合短正合列
$$
0 \to M \to S^{\oplus n + 1} \to S(1) \to 0
$$
其中第二个映射由 $T_0, \ldots, T_n$ 给出。在这种情形下，
上述陈述转化为：诸元素
$$
T_{ij} = (0, \ldots, 0, T_j, 0, \ldots, 0, -T_i, 0, \ldots, 0)
\in M(1)_0
$$
生成分次模 $M(1)$（作为 $S$-模）。细节从略。
\end{proof}








\section{可逆层与到 Proj 的态射}
\label{constructions-section-invertible-proj}

\noindent
设 $T$ 为概形，并设 $\mathcal{L}$ 为 $T$ 上的可逆层。对截面
$s \in \Gamma(T, \mathcal{L})$，用 $T_s$ 表示 $s$ 不消失的点所成开子集。
见模章引理 \ref{modules-lemma-s-open}。可将下面的引理看作引理
\ref{constructions-lemma-apply} 的轻微推广。它也是引理 \ref{constructions-lemma-morphism-proj} 的推广。

\begin{lemma}
\label{constructions-lemma-invertible-map-into-proj}
设 $A$ 为分次环。
置 $X = \text{Proj}(A)$。
设 $T$ 为概形。
设 $\mathcal{L}$ 为可逆 $\mathcal{O}_T$-模。
设 $\psi : A \to \Gamma_*(T, \mathcal{L})$ 为分次环同态。置
$$
U(\psi) = \bigcup\nolimits_{f \in A_{+}\text{ 齐次}} T_{\psi(f)}
$$
态射 $\psi$ 诱导概形的典范态射
$$
r_{\mathcal{L}, \psi} :
U(\psi) \longrightarrow X
$$
以及 $\mathbf{Z}$-分次 $\mathcal{O}_T$-代数的映射
$$
\theta :
r_{\mathcal{L}, \psi}^*\left(
\bigoplus\nolimits_{d \in \mathbf{Z}} \mathcal{O}_X(d)
\right)
\longrightarrow
\bigoplus\nolimits_{d \in \mathbf{Z}} \mathcal{L}^{\otimes d}|_{U(\psi)}.
$$
三元组 $(U(\psi), r_{\mathcal{L}, \psi}, \theta)$ 由下列性质刻画：
\begin{enumerate}
\item 对齐次元 $f \in A_{+}$，有
$r_{\mathcal{L}, \psi}^{-1}(D_{+}(f)) = T_{\psi(f)}$.
\item 对每个 $d \geq 0$，图
$$
\xymatrix{
A_d \ar[d]_{(\ref{constructions-equation-global-sections})} \ar[r]_{\psi} &
\Gamma(T, \mathcal{L}^{\otimes d}) \ar[d]^{restrict} \\
\Gamma(X, \mathcal{O}_X(d)) \ar[r]^{\theta} &
\Gamma(U(\psi), \mathcal{L}^{\otimes d})
}
$$
交换。
\end{enumerate}
此外，对任意 $d \geq 1$ 与任意开子概形 $V \subset T$，若
$\psi(A_d)$ 中的截面生成 $\mathcal{L}^{\otimes d}|_V$，则态射
$r_{\mathcal{L}, \psi}|_V$ 与态射 $\varphi : V \to \text{Proj}(A)$ 一致，
且映射 $\theta|_V$ 与映射
$\alpha : \varphi^*\mathcal{O}_X(d) \to \mathcal{L}^{\otimes d}|_V$ 一致；
这里 $(\varphi, \alpha)$ 是引理 \ref{constructions-lemma-converse-construction} 中与
$\psi|_{A^{(d)}} : A^{(d)} \to \Gamma_*(V, \mathcal{L}^{\otimes d})$.
相伴的二元组。
\end{lemma}

\begin{proof}
假设有两个满足 (1) 与 (2) 的三元组 $(U, r : U \to X, \theta)$
和 $(U', r' : U' \to X, \theta')$。性质 (1) 蕴含
$U = U' = U(\psi)$，并且 $r = r'$ 作为底拓扑空间的映射相等，
因为诸开集 $D_{+}(f)$ 构成 $X$ 的拓扑的一组基，而且 $X$ 是清醒拓扑空间
（见代数章第 \ref{algebra-section-proj} 节与概形章引理
\ref{schemes-lemma-scheme-sober}）。
设 $f \in A_{+}$ 为齐次元。注意
$\Gamma(D_{+}(f), \bigoplus_{n \in \mathbf{Z}} \mathcal{O}_X(n)) = A_f$
作为 $\mathbf{Z}$-分次代数成立。考虑两个 $\mathbf{Z}$-分次环同态
$$
\theta, \theta' :
A_f
\longrightarrow
\Gamma(T_{\psi(f)}, \bigoplus \mathcal{L}^{\otimes n}).
$$
已知乘以 $f$（相应地，乘以 $\psi(f)$）在左边（相应地，右边）
是同构。还由 (2) 知道，
$\theta(x/1) = \theta'(x/1) = \psi(x)|_{T_{\psi(f)}}$
对所有 $x \in A$ 都成立。
因此容易推出所需的 $\theta = \theta'$。考察次数 $0$ 部分可得
$r^\sharp = (r')^\sharp$，即 $r = r'$ 作为概形态射相等。
这证明了唯一性。

\medskip\noindent
现在证明存在性。由刚证明的唯一性，只需局部构造二元组
$(r, \theta)$，其所在概形为 $T$。因此可假设 $T = \Spec(R)$ 仿射，
$\mathcal{L} = \mathcal{O}_T$，并且对某个齐次元 $f \in A_{+}$，
$\psi(f)$ 生成 $\mathcal{O}_T = \mathcal{O}_T^{\otimes \deg(f)}$。
换言之，$\psi(f) = u \in R^*$ 是单位。在这种情形下，映射
$\psi$ 是分次环同态
$$
A \longrightarrow R[x] = \Gamma_*(T, \mathcal{O}_T)
$$
它将 $f$ 映到 $ux^{\deg(f)}$。显然，这唯一地延拓为
$\mathbf{Z}$-分次环同态 $\theta : A_f \to R[x, x^{-1}]$，
方式是将 $1/f$ 映到 $u^{-1}x^{-\deg(f)}$。此映射的次数零部分给出
环同态 $A_{(f)} \to R$，后者给出态射
$r : T = \Spec(R) \to \Spec(A_{(f)}) = D_{+}(f) \subset X$。
因此在这一特殊情形中构造出了 $(r, \theta)$。

\medskip\noindent
最后证明引理的最后一个陈述。根据引理 \ref{constructions-lemma-converse-construction}，
那里构造的态射是使其陈述中所示图交换的唯一态射。
本引理中图的交换性蕴含其限制到 $V$ 与 $A^{(d)}$ 后仍交换。
结论随即得出。
\end{proof}

\begin{remark}
\label{constructions-remark-not-in-invertible-locus}
假设与上面引理 \ref{constructions-lemma-invertible-map-into-proj} 相同。
态射 $r_{\mathcal{L}, \psi}$ 的像不必包含于层
$\mathcal{O}_X(1)$ 可逆的轨迹中。
下面给出一个例子。
设 $k$ 为域。
设 $S = k[A, B, C]$ 按 $\deg(A) = 1$、$\deg(B) = 2$、$\deg(C) = 3$ 分次。
置 $X = \text{Proj}(S)$。
设 $T = \mathbf{P}^2_k = \text{Proj}(k[X_0, X_1, X_2])$。
回忆 $\mathcal{L} = \mathcal{O}_T(1)$ 可逆，且
$\mathcal{O}_T(n) = \mathcal{L}^{\otimes n}$。
考虑这些映射的复合 $\psi$：
$$
S \to k[X_0, X_1, X_2] \to \Gamma_*(T, \mathcal{L}).
$$
这里第一个映射为 $A \mapsto X_0$、$B \mapsto X_1^2$、
$C \mapsto X_2^3$，第二个映射为 (\ref{constructions-equation-global-sections})。
由本引理，这对应于态射
$r_{\mathcal{L}, \psi} : T \to X = \text{Proj}(S)$
并且容易看出它是满射。另一方面，在注记 \ref{constructions-remark-not-isomorphism} 中，
我们已经说明层 $\mathcal{O}_X(1)$ 并非在 $X$ 的所有点处都可逆。
\end{remark}










\section{通过黏合构造相对 Proj}
\label{constructions-section-relative-proj-via-glueing}

\begin{situation}
\label{constructions-situation-relative-proj}
这里 $S$ 为概形，$\mathcal{A}$ 为拟凝聚分次 $\mathcal{O}_S$-代数。
\end{situation}

\noindent
本节概述如何构造概形态射
$$
\underline{\text{Proj}}_S(\mathcal{A}) \longrightarrow S
$$
其方法是黏合齐次谱 $\text{Proj}(\Gamma(U, \mathcal{A}))$，其中
$U$ 遍历 $S$ 的仿射开集。首先说明，$\mathcal{A}$ 在诸仿射开集上的取值之
齐次谱构成一族合适的概形，如引理 \ref{constructions-lemma-relative-glueing} 所述。

\begin{lemma}
\label{constructions-lemma-proj-inclusion}
在情形 \ref{constructions-situation-relative-proj} 中。
假设 $U \subset U' \subset S$ 为仿射开集。
设 $A = \mathcal{A}(U)$ 且 $A' = \mathcal{A}(U')$。
分次环同态 $A' \to A$ 诱导态射
$r : \text{Proj}(A) \to \text{Proj}(A')$，且图
$$
\xymatrix{
\text{Proj}(A) \ar[r] \ar[d] &
\text{Proj}(A') \ar[d] \\
U \ar[r] &
U'
}
$$
是笛卡尔的。此外，存在与乘法映射相容的典范同构
\par\noindent
$\theta : r^*\mathcal{O}_{\text{Proj}(A')}(n) \to
\mathcal{O}_{\text{Proj}(A)}(n)$。
\end{lemma}

\begin{proof}
设 $R = \mathcal{O}_S(U)$ 且 $R' = \mathcal{O}_S(U')$。
注意，映射 $R \otimes_{R'} A' \to A$ 是同构，因为 $\mathcal{A}$ 拟凝聚
（例如见概形章引理 \ref{schemes-lemma-widetilde-pullback}）。
因此本引理由引理 \ref{constructions-lemma-base-change-map-proj} 得出。
\end{proof}

\noindent
特别地，引理中的态射 $\text{Proj}(A) \to \text{Proj}(A')$ 是开浸入。

\begin{lemma}
\label{constructions-lemma-transitive-proj}
在情形 \ref{constructions-situation-relative-proj} 中。
假设 $U \subset U' \subset U'' \subset S$ 为仿射开集。
设 $A = \mathcal{A}(U)$、$A' = \mathcal{A}(U')$ 且 $A'' = \mathcal{A}(U'')$。
引理 \ref{constructions-lemma-proj-inclusion} 中态射
$r : \text{Proj}(A) \to \text{Proj}(A')$ 与
$r' : \text{Proj}(A') \to \text{Proj}(A'')$ 的复合，等于同一引理中的态射
$r'' : \text{Proj}(A) \to \text{Proj}(A'')$，即引理
\ref{constructions-lemma-proj-inclusion} 的相应态射。对同构 $\theta$ 也有类似陈述。
\end{lemma}

\begin{proof}
这由引理 \ref{constructions-lemma-morphism-proj-transitive} 得出，因为映射
$A'' \to A$ 是 $A'' \to A'$ 与 $A' \to A$ 的复合。
\end{proof}

\begin{lemma}
\label{constructions-lemma-glue-relative-proj}
在情形 \ref{constructions-situation-relative-proj} 中。存在概形态射
$$
\pi : \underline{\text{Proj}}_S(\mathcal{A}) \longrightarrow S
$$
具有下列性质：
\begin{enumerate}
\item 对每个仿射开集 $U \subset S$，存在同构
$i_U : \pi^{-1}(U) \to \text{Proj}(A)$，其中 $A = \mathcal{A}(U)$；并且
\item 对仿射开集 $U \subset U' \subset S$，复合
$$
\xymatrix{
\text{Proj}(A) \ar[r]^{i_U^{-1}} &
\pi^{-1}(U) \ar[rr]^{inclusion} & &
\pi^{-1}(U') \ar[r]^{i_{U'}} &
\text{Proj}(A')
}
$$
（其中 $A = \mathcal{A}(U)$、$A' = \mathcal{A}(U')$）
是上面引理 \ref{constructions-lemma-proj-inclusion} 的开浸入。
\end{enumerate}
\end{lemma}

\begin{proof}
这立即由引理 \ref{constructions-lemma-relative-glueing}、
\ref{constructions-lemma-proj-inclusion} 与
\ref{constructions-lemma-transitive-proj}.
得出。
\end{proof}

\begin{lemma}
\label{constructions-lemma-glue-relative-proj-twists}
在情形 \ref{constructions-situation-relative-proj} 中。
引理 \ref{constructions-lemma-glue-relative-proj} 的态射
$\pi : \underline{\text{Proj}}_S(\mathcal{A}) \to S$ 带有下列附加结构。
存在拟凝聚 $\mathbf{Z}$-分次
$\mathcal{O}_{\underline{\text{Proj}}_S(\mathcal{A})}$-代数层
$\bigoplus\nolimits_{n \in \mathbf{Z}}
\mathcal{O}_{\underline{\text{Proj}}_S(\mathcal{A})}(n)$,
以及分次 $\mathcal{O}_S$-代数的态射
$$
\psi :
\mathcal{A}
\longrightarrow
\bigoplus\nolimits_{n \geq 0}
\pi_*\left(\mathcal{O}_{\underline{\text{Proj}}_S(\mathcal{A})}(n)\right)
$$
它们由下列性质唯一确定：对每个仿射开集 $U \subset S$，若
$A = \mathcal{A}(U)$，则存在同构
$$
\theta_U :
i_U^*\left(
\bigoplus\nolimits_{n \in \mathbf{Z}} \mathcal{O}_{\text{Proj}(A)}(n)
\right)
\longrightarrow
\left(
\bigoplus\nolimits_{n \in \mathbf{Z}}
\mathcal{O}_{\underline{\text{Proj}}_S(\mathcal{A})}(n)
\right)|_{\pi^{-1}(U)}
$$
（它是 $\mathbf{Z}$-分次 $\mathcal{O}_{\pi^{-1}(U)}$-代数的同构），使得图
$$
\xymatrix{
A_n
\ar[rr]_\psi
\ar[dr]_-{(\ref{constructions-equation-global-sections})}
& &
\Gamma(\pi^{-1}(U),
\mathcal{O}_{\underline{\text{Proj}}_S(\mathcal{A})}(n)) \\
&
\Gamma(\text{Proj}(A),
\mathcal{O}_{\text{Proj}(A)}(n))
\ar[ru]_-{\theta_U}
&
}
$$
交换。
\end{lemma}

\begin{proof}
我们将使用引理 \ref{constructions-lemma-relative-glueing-sheaves}，把各 $\mathbf{Z}$-分次代数层
$\bigoplus_{n \in \mathbf{Z}} \mathcal{O}_{\text{Proj}(A)}(n)$
黏合起来；这里 $A = \mathcal{A}(U)$，$U \subset S$ 为仿射开集，
黏合所得概形为 $\underline{\text{Proj}}_S(\mathcal{A})$。
所需数据已在引理 \ref{constructions-lemma-proj-inclusion} 中构造，并且引理
\ref{constructions-lemma-relative-glueing-sheaves} 的条件 (d) 已在引理
\ref{constructions-lemma-transitive-proj} 中检验。因此得到 $\mathbf{Z}$-分次
$\mathcal{O}_{\underline{\text{Proj}}_S(\mathcal{A})}$-代数层
$\bigoplus_{n \in \mathbf{Z}}
\mathcal{O}_{\underline{\text{Proj}}_S(\mathcal{A})}(n)$
以及同构 $\theta_U$，它们对所有仿射开集 $U \subset S$ 与所有
$n \in \mathbf{Z}$ 都存在。对每个仿射开集 $U \subset S$，若
$A = \mathcal{A}(U)$，则有映射
$A \to \Gamma(\text{Proj}(A),
\bigoplus_{n \geq 0} \mathcal{O}_{\text{Proj}(A)}(n))$.
因此由相对黏合的函子性，映射 $\psi$ 存在，见注记
\ref{constructions-remark-relative-glueing-functorial}。引理中的图由构造可知交换。
这刻画了 $\mathbf{Z}$-分次 $\mathcal{O}_{\underline{\text{Proj}}_S(\mathcal{A})}$-代数层
$\bigoplus \mathcal{O}_{\underline{\text{Proj}}_S(\mathcal{A})}(n)$
因为引理 \ref{constructions-lemma-morphism-proj} 的证明表明，要求这些图交换会唯一确定
各映射 $\theta_U$。略去一些细节。
\end{proof}















\section{作为函子的相对 Proj}
\label{constructions-section-relative-proj}

\noindent
我们处于情形 \ref{constructions-situation-relative-proj}。因此 $S$ 为概形，且
$\mathcal{A} = \bigoplus_{d \geq 0} \mathcal{A}_d$ 为拟凝聚分次
$\mathcal{O}_S$-代数。本节通过构造一个由相对齐次谱表示的函子，
将 $\text{Proj}$ 的构造相对化。由此将构造概形态射
$$
\underline{\text{Proj}}_S(\mathcal{A}) \longrightarrow S
$$
它在 $S$ 的仿射开集上看起来像分次环的齐次谱。讨论仿照第
\ref{constructions-section-spec} 节对相对谱的讨论。较容易的方法是黏合形如
$\text{Proj}(A)$ 的概形，其中 $A = \Gamma(U, \mathcal{A})$，
$U \subset S$ 为仿射开集；该方法见第
\ref{constructions-section-relative-proj-via-glueing} 节，并将证明本节结果与之同构。

\medskip\noindent
暂时固定整数 $d \geq 1$。
仿照代数章第 \ref{algebra-section-graded} 节的记号，记
$\mathcal{A}^{(d)} = \bigoplus_{n \geq 0} \mathcal{A}_{nd}$。
设 $T$ 为概形。考虑{\it 四元组 $(d, f : T \to S, \mathcal{L}, \psi)$
（在 $T$ 上）}，其中
\begin{enumerate}
\item $d$ 是上面固定的整数；
\item $f : T \to S$ 是概形态射；
\item $\mathcal{L}$ 是可逆 $\mathcal{O}_T$-模；并且
\item
$\psi : f^*\mathcal{A}^{(d)} \to \bigoplus_{n \geq 0}\mathcal{L}^{\otimes n}$
是分次 $\mathcal{O}_T$-代数同态，且
$f^*\mathcal{A}_d \to \mathcal{L}$ 是满射。
\end{enumerate}
给定态射 $h : T' \to T$ 与四元组 $(d, f, \mathcal{L}, \psi)$（在 $T$ 上），
可将它拉回为四元组 $(d, f \circ h, h^*\mathcal{L}, h^*\psi)$（在 $T'$ 上）。
给定两个四元组 $(d, f, \mathcal{L}, \psi)$ 与
$(d, f', \mathcal{L}', \psi')$（在 $T$ 上且整数 $d$ 相同），若
$f = f'$，并存在同构 $\beta : \mathcal{L} \to \mathcal{L}'$，使得
$\beta \circ \psi = \psi'$ 作为分次 $\mathcal{O}_T$-代数同态
$f^*\mathcal{A}^{(d)} \to \bigoplus_{n \geq 0} (\mathcal{L}')^{\otimes n}$ 成立，
则称它们{\it 严格等价}。

\medskip\noindent
对每个整数 $d \geq 1$，定义
\begin{eqnarray*}
F_d : \Sch^{opp} & \longrightarrow & \textit{Sets}, \\
T & \longmapsto &
\{\text{下列对象的严格等价类： }
(d, f : T \to S, \mathcal{L}, \psi)
\text{ 如上}\}
\end{eqnarray*}
其拉回按上面所述定义。

\begin{lemma}
\label{constructions-lemma-proj-base-change}
在情形 \ref{constructions-situation-relative-proj} 中。设 $d \geq 1$。
设 $F_d$ 为上面与 $(S, \mathcal{A})$ 相伴的函子。
设 $g : S' \to S$ 为概形态射。置 $\mathcal{A}' = g^*\mathcal{A}$。
设 $F_d'$ 为上面与 $(S', \mathcal{A}')$ 相伴的函子。
则存在函子的典范同构
$$
F'_d \cong h_{S'} \times_{h_S} F_d
$$
\end{lemma}

\begin{proof}
四元组
$(d, f' : T \to S', \mathcal{L}',
\psi' : (f')^*(\mathcal{A}')^{(d)} \to
\bigoplus_{n \geq 0} (\mathcal{L}')^{\otimes n})$
等同于一个四元组
$(d, f, \mathcal{L},
\psi : f^*\mathcal{A}^{(d)} \to
\bigoplus_{n \geq 0} \mathcal{L}^{\otimes n})$
连同 $f$ 的分解 $f = g \circ f'$。具体而言，该对应由
$f = g \circ f'$、$\mathcal{L} = \mathcal{L}'$ 与 $\psi = \psi'$ 给出，
其中使用识别
$(f')^*(\mathcal{A}')^{(d)} = (f')^*g^*(\mathcal{A}^{(d)}) =
f^*\mathcal{A}^{(d)}$。引理得证。
\end{proof}

\begin{lemma}
\label{constructions-lemma-relative-proj-affine}
在情形 \ref{constructions-situation-relative-proj} 中。设 $F_d$ 为上面与
$(d, S, \mathcal{A})$ 相伴的函子。若 $S$ 仿射，则 $F_d$ 由开子概形
$U_d$（\ref{constructions-equation-Ud}）表示；该开子概形位于概形
$\text{Proj}(\Gamma(S, \mathcal{A}))$ 中。
\end{lemma}

\begin{proof}
写成 $S = \Spec(R)$ 且 $A = \Gamma(S, \mathcal{A})$。
则 $A$ 是分次 $R$-代数，且 $\mathcal{A} = \widetilde A$。
为证明本引理，必须把函子 $F_d$ 与第 \ref{constructions-section-morphisms-proj} 节定义的
三元组函子 $F_d^{triples}$ 识别起来。

\medskip\noindent
设 $(d, f : T \to S, \mathcal{L}, \psi)$ 为四元组。
可将 $\psi$ 看作 $\mathcal{O}_S$-模同态
$\mathcal{A}^{(d)} \to \bigoplus_{n \geq 0} f_*\mathcal{L}^{\otimes n}$.
由于 $\mathcal{A}^{(d)}$ 拟凝聚，这等同于 $R$-线性的分次环同态
$A^{(d)} \to \Gamma(S, \bigoplus_{n \geq 0} f_*\mathcal{L}^{\otimes n})$.
显然，$\Gamma(S, \bigoplus_{n \geq 0} f_*\mathcal{L}^{\otimes n}) =
\Gamma_*(T, \mathcal{L})$。因此可将三元组 $(d, \mathcal{L}, \psi)$
与该四元组相伴。

\medskip\noindent
反之，设 $(d, \mathcal{L}, \psi)$ 为三元组。
复合 $R \to A_0 \to \Gamma(T, \mathcal{O}_T)$ 确定态射
$f : T \to S = \Spec(R)$，见概形章引理
\ref{schemes-lemma-morphism-into-affine}。按此选择 $f$，映射
$A^{(d)} \to \Gamma(S, \bigoplus_{n \geq 0} f_*\mathcal{L}^{\otimes n})$
是 $R$-线性的，因而对应于 $\psi$；它可用于四元组
$(d, f : T \to S, \mathcal{L}, \psi)$。
由此建立函子同构 $F_d = F_d^{triples}$ 的验证从略。
\end{proof}

\begin{lemma}
\label{constructions-lemma-relative-proj-d}
在情形 \ref{constructions-situation-relative-proj} 中。函子 $F_d$ 可由概形表示。
\end{lemma}

\begin{proof}
我们将使用概形章引理 \ref{schemes-lemma-glue-functors}。

\medskip\noindent
首先检验 $F_d$ 对 Zariski 拓扑满足层性质。设 $T$ 为概形，
$T = \bigcup_{i \in I} U_i$ 为开覆盖，并且
$(d, f_i, \mathcal{L}_i, \psi_i) \in F_d(U_i)$，使得
$(d, f_i, \mathcal{L}_i, \psi_i)|_{U_i \cap U_j}$ 与
$(d, f_j, \mathcal{L}_j, \psi_j)|_{U_i \cap U_j}$ 严格等价。
这蕴含态射 $f_i : U_i \to S$ 黏合为概形态射
$f : T \to S$，满足 $f|_{I_i} = f_i$，见概形章第
\ref{schemes-section-glueing-schemes} 节。因此
$f_i^*\mathcal{A}^{(d)} = f^*\mathcal{A}^{(d)}|_{U_i}$。
它还蕴含存在同构
$\beta_{ij} : \mathcal{L}_i|_{U_i \cap U_j} \to \mathcal{L}_j|_{U_i \cap U_j}$
使得 $\beta_{ij} \circ \psi_i = \psi_j$ 在 $U_i \cap U_j$ 上成立。
注意，同构 $\beta_{ij}$ 由这一要求唯一确定，因为映射
$f_i^*\mathcal{A}_d \to \mathcal{L}_i$ 是满射。特别地，可见
$\beta_{jk} \circ \beta_{ij} = \beta_{ik}$
在 $U_i \cap U_j \cap U_k$ 上成立。因此由层章第
\ref{sheaves-section-glueing-sheaves} 节，可逆层 $\mathcal{L}_i$
黏合为可逆 $\mathcal{O}_T$-模 $\mathcal{L}$，而态射 $\psi_i$
黏合为 $\mathcal{O}_T$-代数的态射
$\psi : f^*\mathcal{A}^{(d)} \to \bigoplus_{n \geq 0} \mathcal{L}^{\otimes n}$.
这证明 $F_d$ 对 Zariski 拓扑满足层条件。

\medskip\noindent
设 $S = \bigcup_{i \in I} U_i$ 为仿射开覆盖。
设 $F_{d, i} \subset F_d$ 为由二元组
$(f : T \to S, \varphi)$ 所组成的子函子，其中 $f(T) \subset U_i$。

\medskip\noindent
必须证明每个 $F_{d, i}$ 都可表示。由引理 \ref{constructions-lemma-proj-base-change}，
$F_{d, i}$ 与下述函子识别：它与 $U_i$ 以及其上的拟凝聚分次
$\mathcal{O}_{U_i}$-代数 $\mathcal{A}|_{U_i}$ 相伴。
因此结论由引理 \ref{constructions-lemma-relative-proj-affine} 得出。

\medskip\noindent
接着证明 $F_{d, i} \subset F_d$ 可由开浸入表示。
设 $(f : T \to S, \varphi) \in F_d(T)$。考虑 $V_i = f^{-1}(U_i)$。
由 $F_{d, i}$ 的定义，对任意 $a : T' \to T$，有
$a^*(f, \varphi) \in F_{d, i}(T')$ 当且仅当 $a(T') \subset V_i$。
这正是所需结论。

\medskip\noindent
最后必须证明族 $(F_{d, i})_{i \in I}$ 覆盖 $F_d$。设
$(f : T \to S, \varphi) \in F_d(T)$。考虑 $V_i = f^{-1}(U_i)$。
由于 $S = \bigcup_{i \in I} U_i$ 是 $S$ 的开覆盖，可见
$T = \bigcup_{i \in I} V_i$ 是 $T$ 的开覆盖。此外，
$(f, \varphi)|_{V_i} \in F_{d, i}(V_i)$。引理证毕。
\end{proof}

\noindent
现在可以在当前相对情形下重做第 \ref{constructions-section-morphisms-proj} 节末尾的内容，
并定义一个由 $\underline{\text{Proj}}_S(\mathcal{A})$ 表示的函子。
为此，引入两个四元组
$(d, f : T \to S, \mathcal{L}, \psi)$ 与
$(d', f' : T \to S, \mathcal{L}', \psi')$ 之间的等价概念，其中整数
$d, d'$ 可以取不同值。具体而言，若 $f = f'$，并存在同构
$\beta : \mathcal{L}^{\otimes d'} \to (\mathcal{L}')^{\otimes d}$，使得
$\beta \circ \psi|_{f^*\mathcal{A}^{(dd')}} = \psi'|_{f^*\mathcal{A}^{(dd')}}$.
则称它们{\it 等价}。下面的引理蕴含这确实定义了等价关系。
（这并非完全平凡。）

\begin{lemma}
\label{constructions-lemma-equivalent-relative}
在情形 \ref{constructions-situation-relative-proj} 中。设 $T$ 为概形。
设 $(d, f, \mathcal{L}, \psi)$、$(d', f', \mathcal{L}', \psi')$
为 $T$ 上的两个四元组。下列条件等价：
\begin{enumerate}
\item 设 $m = \text{lcm}(d, d')$。写成 $m = ad = a'd'$。
有 $f = f'$，并且存在同构
$\beta : \mathcal{L}^{\otimes a} \to (\mathcal{L}')^{\otimes a'}$
使得 $\beta \circ \psi|_{f^*\mathcal{A}^{(m)}}$ 与
$\psi'|_{f^*\mathcal{A}^{(m)}}$ 作为分次环同态
$f^*\mathcal{A}^{(m)} \to \bigoplus_{n \geq 0} (\mathcal{L}')^{\otimes mn}$.
一致。
\item 四元组 $(d, f, \mathcal{L}, \psi)$ 与
$(d', f', \mathcal{L}', \psi')$ 等价。
\item 有 $f = f'$，并且对某个正整数 $m = ad = a'd'$，存在同构
$\beta : \mathcal{L}^{\otimes a} \to (\mathcal{L}')^{\otimes a'}$
使得 $\beta \circ \psi|_{f^*\mathcal{A}^{(m)}}$ 与
$\psi'|_{f^*\mathcal{A}^{(m)}}$ 作为分次环同态
$f^*\mathcal{A}^{(m)} \to \bigoplus_{n \geq 0} (\mathcal{L}')^{\otimes mn}$.
一致。
\end{enumerate}
\end{lemma}

\begin{proof}
显然，通过限制到可被更多整数整除的次数并取可逆层的幂，
(1) 蕴含 (2)，且 (2) 蕴含 (3)。假设 (3) 对某个整数
$m = ad = a'd'$ 成立。置
$m_0 = \text{lcm}(d, d')$，并写成 $m_0 = a_0d = a'_0d'$。
给定同构
$\beta : \mathcal{L}^{\otimes a} \to (\mathcal{L}')^{\otimes a'}$
满足 (3) 所述性质。要寻找同构
$\beta_0 : \mathcal{L}^{\otimes a_0} \to (\mathcal{L}')^{\otimes a'_0}$
也满足该性质。由于按假设映射
$\psi : f^*\mathcal{A}_d \to \mathcal{L}$ 与
$\psi' : (f')^*\mathcal{A}_{d'} \to \mathcal{L}'$ 是满射，
下列映射也同样是满射：
$\psi : f^*\mathcal{A}_{m_0} \to \mathcal{L}^{\otimes a_0}$
以及 $\psi' : (f')^*\mathcal{A}_{m_0} \to (\mathcal{L}')^{\otimes a_0}$。
因此，若 $\beta_0$ 存在，则由条件 $\beta_0 \circ \psi = \psi'$ 唯一确定。
这意味着可以在 $T$ 上局部工作。因此可假设
$f = f' : T \to S$ 的像落在某个仿射开集中，换言之可假设 $S$ 仿射。
此时结论由三元组的相应结果（见引理 \ref{constructions-lemma-equivalent}）以及仿射底情形下
三元组与四元组之间的对应（见引理 \ref{constructions-lemma-relative-proj-affine} 的证明）得出。
\end{proof}

\noindent
假设 $d' = ad$。考虑函子变换 $F_d \to F_{d'}$，它将四元组
$(d, f, \mathcal{L}, \psi)$（在 $T$ 上）映为四元组
$(d', f, \mathcal{L}^{\otimes a}, \psi|_{f^*\mathcal{A}^{(d')}})$.
引理 \ref{constructions-lemma-equivalent-relative} 的一个推论是：变换
$F_d \to F_{d'}$ 是单射！对拟紧概形 $T$，定义
$$
F(T) = \bigcup\nolimits_{d \in \mathbf{N}} F_d(T)
$$
其中转移映射如上所述。显然，这在拟紧概形范畴上定义了取值于集合的反变函子。
对一般概形 $T$，定义
$$
F(T)
=
\lim_{V \subset T\text{ 拟紧开}} F(V).
$$
换言之，一个元素 $\xi$（属于 $F(T)$）对应于一族相容的选择
$\xi_V \in F(V)$，其中 $V$ 遍历 $T$ 的拟紧开集。
拉回映射 $F(T) \to F(T')$（对应概形态射 $T' \to T$）的定义从略。
由此定义函子
\begin{equation}
\label{constructions-equation-proj}
F : \Sch^{opp} \longrightarrow \textit{Sets}
\end{equation}

\begin{lemma}
\label{constructions-lemma-relative-proj}
在情形 \ref{constructions-situation-relative-proj} 中。上面的函子 $F$ 可由概形表示。
\end{lemma}

\begin{proof}
设 $U_d \to S$ 为表示上面函子 $F_d$ 的概形。设 $\mathcal{L}_d$ 与
$\psi^d : \pi_d^*\mathcal{A}^{(d)} \to
\bigoplus_{n \geq 0} \mathcal{L}_d^{\otimes n}$ 为泛对象。
若 $d | d'$，则可考虑四元组
$(d', \pi_d, \mathcal{L}_d^{\otimes d'/d}, \psi^d|_{\mathcal{A}^{(d')}})$
它确定典范态射 $U_d \to U_{d'}$（在 $S$ 上）。
由构造，此态射对应于上面定义的函子变换 $F_d \to F_{d'}$。

\medskip\noindent
对每个仿射开集 $\Spec(R) = V \subset S$，置
$A = \Gamma(V, \mathcal{A})$，则基变换 $U_{d, V}$ 与
$\text{Proj}(A)$ 的相应开子概形有典范识别，见引理
\ref{constructions-lemma-relative-proj-affine}。此外，上面构造的态射
$U_{d, V} \to U_{d', V}$ 对应于 $\text{Proj}(A)$ 中开集的包含。
因此 $U_d \to U_{d'}$ 是开浸入。

\medskip\noindent
于是可构造 $X$：将概形 $U_d$ 沿开浸入 $U_d \to U_{d'}$ 黏合。
技术上，方便的做法是选取序列 $d_1 | d_2 | d_3 | \ldots$，
使每个正整数整除某个 $d_i$，然后利用上述开浸入直接取
$X = \bigcup U_{d_i}$。接着容易证明 $X$ 表示函子 $F$。
\end{proof}

\begin{lemma}
\label{constructions-lemma-glueing-gives-functor-proj}
在情形 \ref{constructions-situation-relative-proj} 中。引理 \ref{constructions-lemma-glue-relative-proj}
构造的概形 $\pi : \underline{\text{Proj}}_S(\mathcal{A}) \to S$，
与表示函子 $F$ 的概形，作为 $S$ 上的概形典范同构。
\end{lemma}

\begin{proof}
设 $X$ 为表示函子 $F$ 的概形。注意，$X$ 是 $S$ 上的概形，
因为函子 $F$ 带有自然变换 $F \to h_S$。
写成 $Y = \underline{\text{Proj}}_S(\mathcal{A})$。
必须证明 $X \cong Y$ 作为 $S$-概形成立。给出两个论证。

\medskip\noindent
第一个论证使用引理 \ref{constructions-lemma-relative-proj} 的证明中将 $X$ 构造为
诸概形 $U_d$（它们表示 $F_d$）之并的做法。在 $S$ 的每个仿射开集上，
可将 $X$ 与 $\mathcal{A}$ 在该开集上的截面之齐次谱识别，因为对开集
$U_d$ 已经如此。此外，这些识别与进一步限制到更小的仿射开集相容。
另一方面，$Y$ 正是通过黏合这些齐次谱构造的。
因此可黏合这些同构，得到 $X$ 与
$\underline{\text{Proj}}_S(\mathcal{A})$ 之间所需的同构。细节从略。

\medskip\noindent
下面是第二个论证。引理 \ref{constructions-lemma-glue-relative-proj-twists} 表明存在分次代数的态射
$$
\psi : \pi^*\mathcal{A}
\longrightarrow
\bigoplus\nolimits_{n \geq 0} \mathcal{O}_Y(n)
$$
它位于 $Y$ 上，并且在 $S$ 的仿射开集上的截面层面与
(\ref{constructions-equation-global-sections}) 一致。因此对每个 $y \in Y$，存在
开邻域 $V \subset Y$（它是 $y$ 的邻域）与整数 $d \geq 1$，使得当 $d | n$ 时，层
$\mathcal{O}_Y(n)|_V$ 可逆，且乘法映射
$\mathcal{O}_Y(n)|_V \otimes_{\mathcal{O}_V} \mathcal{O}_Y(m)|_V
\to \mathcal{O}_Y(n + m)|_V$ 都是同构。因此，将 $\psi$ 限制到层
$\pi^*\mathcal{A}^{(d)}|_V$，得到 $F_d(V)$ 的一个元素。由于诸开集
$V$ 覆盖 $Y$，可见“$\psi$”给出 $F(Y)$ 的一个元素，
从而给出典范态射 $Y \to X$（在 $S$ 上）。由于此构造完全典范，
为验证它是同构，可以在 $S$ 上局部工作。因此归约到 $S$ 仿射的情形，
此时结论显然。
\end{proof}

\begin{definition}
\label{constructions-definition-relative-proj}
设 $S$ 为概形。设 $\mathcal{A}$ 为拟凝聚分次 $\mathcal{O}_S$-代数层。
{\it $\mathcal{A}$ 在 $S$ 上的相对齐次谱}，或
{\it $\mathcal{A}$ 在 $S$ 上的齐次谱}，或
{\it $\mathcal{A}$ 在 $S$ 上的相对 Proj}，是引理
\ref{constructions-lemma-glue-relative-proj} 中构造且表示函子 $F$
（\ref{constructions-equation-proj}）的概形，见引理
\ref{constructions-lemma-glueing-gives-functor-proj}。将它记为
$\pi : \underline{\text{Proj}}_S(\mathcal{A}) \to S$。
\end{definition}

\noindent
相对 Proj 带有拟凝聚 $\mathbf{Z}$-分次代数层
$\bigoplus_{n \in \mathbf{Z}}
\mathcal{O}_{\underline{\text{Proj}}_S(\mathcal{A})}(n)$
（结构层的扭转），以及分次代数的“泛”同态
$$
\psi_{univ} :
\mathcal{A}
\longrightarrow
\pi_*\left(
\bigoplus\nolimits_{n \geq 0}
\mathcal{O}_{\underline{\text{Proj}}_S(\mathcal{A})}(n)
\right)
$$
见引理 \ref{constructions-lemma-glue-relative-proj-twists}。也可将其视为同态
$$
\psi_{univ} :
\pi^*\mathcal{A}
\longrightarrow
\bigoplus\nolimits_{n \geq 0}
\mathcal{O}_{\underline{\text{Proj}}_S(\mathcal{A})}(n)
$$
。下面的引理表述了此对象的泛性质。

\begin{lemma}
\label{constructions-lemma-tie-up-psi}
在情形 \ref{constructions-situation-relative-proj} 中。设
$(f : T \to S, d, \mathcal{L}, \psi)$ 为四元组。设
$r_{d, \mathcal{L}, \psi} : T \to \underline{\text{Proj}}_S(\mathcal{A})$
为相伴的 $S$-态射。存在 $\mathbf{Z}$-分次 $\mathcal{O}_T$-代数的同构
$$
\theta :
r_{d, \mathcal{L}, \psi}^*\left(
\bigoplus\nolimits_{n \in \mathbf{Z}}
\mathcal{O}_{\underline{\text{Proj}}_S(\mathcal{A})}(nd)
\right)
\longrightarrow
\bigoplus\nolimits_{n \in \mathbf{Z}} \mathcal{L}^{\otimes n}
$$
使得下图交换：
$$
\xymatrix{
\mathcal{A}^{(d)} \ar[rr]_-{\psi}
 \ar[rd]_-{\psi_{univ}} & &
f_*\left(
\bigoplus\nolimits_{n \in \mathbf{Z}}
\mathcal{L}^{\otimes n}
\right) \\
 &
\pi_*\left(
\bigoplus\nolimits_{n \geq 0}
\mathcal{O}_{\underline{\text{Proj}}_S(\mathcal{A})}(nd)
\right) \ar[ru]_\theta
}
$$
此图的交换性唯一确定 $\theta$。
\end{lemma}

\begin{proof}
注意，四元组 $(f : T \to S, d, \mathcal{L}, \psi)$ 定义了
$F_d(T)$ 的一个元素。设
$U_d \subset \underline{\text{Proj}}_S(\mathcal{A})$
为下述轨迹：层 $\mathcal{O}_{\underline{\text{Proj}}_S(\mathcal{A})}(d)$
可逆，并由下列映射的像生成：
$\psi_{univ} : \pi^*\mathcal{A}_d \to
\mathcal{O}_{\underline{\text{Proj}}_S(\mathcal{A})}(d)$.
回忆 $U_d$ 表示函子 $F_d$，见引理 \ref{constructions-lemma-relative-proj} 的证明。
因此，只需证明四元组
$(U_d \to S, d, \mathcal{O}_{U_d}(d), \psi_{univ}|_{\mathcal{A}^{(d)}})$
是泛族，即 $F_d(U_d)$ 中的表示对象。可以在限制到 $S$ 的仿射开集后证明，
因为 (a) 函子 $F_d$ 的形成与基变换交换（见引理
\ref{constructions-lemma-proj-base-change}），且 (b) 二元组
$(\bigoplus_{n \in \mathbf{Z}}
\mathcal{O}_{\underline{\text{Proj}}_S(\mathcal{A})}(n),
\psi_{univ})$
是在 $S$ 的仿射开集上黏合构造的（见引理
\ref{constructions-lemma-glue-relative-proj-twists}）。因此可假设 $S$ 仿射。
此时，四元组函子 $F_d$ 与三元组函子 $F_d$ 一致
（见引理 \ref{constructions-lemma-relative-proj-affine} 的证明），而且引理
\ref{constructions-lemma-proj-functor-strict} 表明
$(d, \mathcal{O}_{U_d}(d), \psi^d)$ 是 $U_d$ 上的泛三元组。
逆向追溯引理 \ref{constructions-lemma-relative-proj-affine} 证明中的各项识别，可知
$(U_d \to S, d, \mathcal{O}_{U_d}(d), \psi_{univ}|_{\mathcal{A}^{(d)}})$
正是所需的泛四元组。
\end{proof}

\begin{lemma}
\label{constructions-lemma-relative-proj-separated}
设 $S$ 为概形，$\mathcal{A}$ 为拟凝聚分次 $\mathcal{O}_S$-代数层。态射
$\pi : \underline{\text{Proj}}_S(\mathcal{A}) \to S$
是分离的。
\end{lemma}

\begin{proof}
要证明态射分离，可以在底上局部工作，见概形章第
\ref{schemes-section-separation-axioms} 节。由构造，
$\underline{\text{Proj}}_S(\mathcal{A})$ 在任意仿射开集 $U \subset S$ 上
都同构于 $\text{Proj}(A)$，其中 $A = \mathcal{A}(U)$。由引理
\ref{constructions-lemma-proj-separated} 可见 $\text{Proj}(A)$ 分离。
因此 $\text{Proj}(A) \to U$ 分离（见概形章引理
\ref{schemes-lemma-compose-after-separated}），如所需。
\end{proof}

\begin{lemma}
\label{constructions-lemma-relative-proj-base-change}
设 $S$ 为概形，$\mathcal{A}$ 为拟凝聚分次 $\mathcal{O}_S$-代数层。
设 $g : S' \to S$ 为任意概形态射。则存在典范同构
$$
r :
\underline{\text{Proj}}_{S'}(g^*\mathcal{A})
\longrightarrow
S' \times_S \underline{\text{Proj}}_S(\mathcal{A})
$$
以及相应的同构
$$
\theta :
r^*\text{pr}_2^*\left(\bigoplus\nolimits_{d \in \mathbf{Z}}
\mathcal{O}_{\underline{\text{Proj}}_S(\mathcal{A})}(d)\right)
\longrightarrow
\bigoplus\nolimits_{d \in \mathbf{Z}}
\mathcal{O}_{\underline{\text{Proj}}_{S'}(g^*\mathcal{A})}(d)
$$
，它是 $\mathbf{Z}$-分次
$\mathcal{O}_{\underline{\text{Proj}}_{S'}(g^*\mathcal{A})}$-代数的同构。
\end{lemma}

\begin{proof}
这由引理 \ref{constructions-lemma-proj-base-change} 以及引理 \ref{constructions-lemma-relative-proj} 中
将 $\underline{\text{Proj}}_S(\mathcal{A})$ 构造为诸概形 $U_d$
（它们表示函子 $F_d$）之并的做法得出。用通过黏合构造相对 Proj 的语言，此同构由引理
\ref{constructions-lemma-base-change-map-proj} 中构造的同构给出，后者也提供同构
$\theta$。略去一些细节。
\end{proof}

\begin{lemma}
\label{constructions-lemma-apply-relative}
设 $S$ 为概形。
设 $\mathcal{A}$ 为拟凝聚分次 $\mathcal{O}_S$-代数层，它作为
$\mathcal{A}_0$-代数由 $\mathcal{A}_1$ 生成。在这种情形下，概形
$X = \underline{\text{Proj}}_S(\mathcal{A})$ 表示函子 $F_1$；该函子将
概形 $f : T \to S$（在 $S$ 上）映到二元组 $(\mathcal{L}, \psi)$ 的集合，其中
\begin{enumerate}
\item $\mathcal{L}$ 是可逆 $\mathcal{O}_T$-模；并且
\item $\psi : f^*\mathcal{A} \to \bigoplus_{n \geq 0} \mathcal{L}^{\otimes n}$
是分次 $\mathcal{O}_T$-代数同态，且
$f^*\mathcal{A}_1 \to \mathcal{L}$ 是满射，
\end{enumerate}
并按上述严格等价关系取等价类。此外，在这种情形下，所有拟凝聚层
$\mathcal{O}_{\underline{\text{Proj}}(\mathcal{A})}(n)$
都是可逆 $\mathcal{O}_{\underline{\text{Proj}}(\mathcal{A})}$-模，
并且乘法映射诱导同构
$
\mathcal{O}_{\underline{\text{Proj}}(\mathcal{A})}(n)
\otimes_{\mathcal{O}_{\underline{\text{Proj}}(\mathcal{A})}}
\mathcal{O}_{\underline{\text{Proj}}(\mathcal{A})}(m) =
\mathcal{O}_{\underline{\text{Proj}}(\mathcal{A})}(n + m)$.
\end{lemma}

\begin{proof}
在本引理的假设下，诸层
$\mathcal{O}_{\underline{\text{Proj}}(\mathcal{A})}(n)$
可逆，且乘法映射是同构；这在 $S$ 的仿射开集上由引理
\ref{constructions-lemma-relative-proj} 与引理 \ref{constructions-lemma-apply} 得出。
因此 $X$ 的确表示函子 $F_1$，见引理 \ref{constructions-lemma-relative-proj} 的证明。
\end{proof}












\section{相对 Proj 上的拟凝聚层}
\label{constructions-section-quasi-coherent-relative-proj}

\noindent
我们简要讨论如何在相对情形下处理分次模。

\medskip\noindent
我们置身于情形 \ref{constructions-situation-relative-proj}。因此 $S$ 是一个概形，而
$\mathcal{A}$ 是一个拟凝聚分次 $\mathcal{O}_S$-代数。设
$\mathcal{M} = \bigoplus_{n \in \mathbf{Z}} \mathcal{M}_n$
是一个分次 $\mathcal{A}$-模，并且作为 $\mathcal{O}_S$-模是拟凝聚的。
我们将描述 $\underline{\text{Proj}}_S(\mathcal{A})$ 上与之相伴的拟凝聚模层。
首先描述这个层在映入相对 Proj 的概形 $T$ 上的值。

\medskip\noindent
设 $T$ 是一个概形，并设 $(d, f : T \to S, \mathcal{L}, \psi)$
为第 \ref{constructions-section-relative-proj} 节所述的 $T$ 上四元组。我们如下定义一个拟凝聚层
$\widetilde{\mathcal{M}}_T$，它是一个 $\mathcal{O}_T$-模层：
\begin{equation}
\label{constructions-equation-widetilde-M}
\widetilde{\mathcal{M}}_T =
\left(
f^*\mathcal{M}^{(d)}
\otimes_{f^*\mathcal{A}^{(d)}}
\left(\bigoplus\nolimits_{n \in \mathbf{Z}} \mathcal{L}^{\otimes n}\right)
\right)_0
\end{equation}
因此 $\widetilde{\mathcal{M}}_T$ 是 $0$ 次部分；它取自两个分次
$f^*\mathcal{A}^{(d)}$-模 $\mathcal{M}^{(d)}$ 与
$\bigoplus\nolimits_{n \in \mathbf{Z}} \mathcal{L}^{\otimes n}$
的张量积。注意，尽管记号中略去了这种依赖，层
$\widetilde{\mathcal{M}}_T$ 仍依赖于该四元组。这个构造具有如下良好性质：
给定任意态射 $g : T' \to T$，都有
$\widetilde{\mathcal{M}}_{T'} = g^*\widetilde{\mathcal{M}}_T$，
其中 $\widetilde{\mathcal{M}}_{T'}$ 表示与拉回四元组
$(d, f \circ g, g^*\mathcal{L}, g^*\psi)$ 相伴的拟凝聚层。

\medskip\noindent
由于 (\ref{constructions-equation-widetilde-M}) 中的所有层都是拟凝聚的，我们可以在仿射开集
$\Spec(C) = V \subset T$ 上具体写出这个构造；该开集映入仿射开集
$\Spec(R) = U \subset S$。具体而言，设 $\mathcal{A}|_U$ 对应于分次 $R$-代数 $A$，
$\mathcal{M}|_U$ 对应于分次 $A$-模 $M$，而 $\mathcal{L}|_V$
对应于可逆 $C$-模 $L$。映射 $\psi$ 诱导一个分次 $R$-代数映射
$\gamma : A^{(d)} \to \bigoplus_{n \geq 0} L^{\otimes n}$。
（这里 $L$ 的张量幂均取在 $C$ 上。）于是 $(\widetilde{\mathcal{M}}_T)|_V$
就是与下列 $C$-模相伴的拟凝聚层：
$$
N_{R, C, A, M, \gamma} =
\left(
M^{(d)} \otimes_{A^{(d)}, \gamma}
\left(\bigoplus\nolimits_{n \in \mathbf{Z}} L^{\otimes n}\right)
\right)_0
$$
根据假设，我们甚至可以用若干仿射开集覆盖 $T$；这些开集记为 $V$，并且存在某个
$a \in A_d$，使得 $\gamma(a) \in L$ 是一个 $C$-基，即模 $L$ 的一组基。
在这种情况下，$N_{R, C, A, M, \gamma}$ 的任意元素都是纯张量之和
$\sum m_i \otimes \gamma(a)^{-n_i}$，其中 $m_i \in M_{n_id}$。
事实上，我们可以把每个 $m_i$ 乘以 $a$ 的某个适当正幂，然后合并各项，
由此看出 $N_{R, C, A, M, \gamma}$ 的每个元素都可写成
$m \otimes \gamma(a)^{-n}$，其中 $m \in M_{nd}$ 且 $n \gg 0$。
换言之，在此情形下有
$$
N_{R, C, A, M, \gamma} = M_{(a)} \otimes_{A_{(a)}} C
$$
其中映射 $A_{(a)} \to C$ 是
$x/a^n \mapsto \gamma(x)/\gamma(a)^n$。换言之，这是
$\widetilde{M}$ 在 $D_{+}(a) \subset \text{Proj}(A)$ 上的值拉回到
$\Spec(C)$ 后所得的模；这里使用的是态射
$\Spec(C) \to D_{+}(a)$，它来自 $\gamma$。

\begin{lemma}
\label{constructions-lemma-relative-proj-modules}
在情形 \ref{constructions-situation-relative-proj} 中，给定任意拟凝聚分次
$\mathcal{A}$-模层 $\mathcal{M}$（它位于 $S$ 上），都存在一个典范相伴的
$\mathcal{O}_{\underline{\text{Proj}}_S(\mathcal{A})}$-模层
$\widetilde{\mathcal{M}}$，它具有如下性质：
\begin{enumerate}
\item 给定概形 $T$ 以及四元组
$(T \to S, d, \mathcal{L}, \psi)$（它位于 $T$ 上），若该四元组对应于态射
$h : T \to \underline{\text{Proj}}_S(\mathcal{A})$，则存在典范同构
$\widetilde{\mathcal{M}}_T = h^*\widetilde{\mathcal{M}}$，
其中 $\widetilde{\mathcal{M}}_T$ 由 (\ref{constructions-equation-widetilde-M}) 定义。
\item (1) 中的同构与拉回相容。
\item 存在典范映射
$$
\pi^*\mathcal{M}_0 \longrightarrow \widetilde{\mathcal{M}}.
$$
\item 构造 $\mathcal{M} \mapsto \widetilde{\mathcal{M}}$ 关于 $\mathcal{M}$ 是函子性的。
\item 构造 $\mathcal{M} \mapsto \widetilde{\mathcal{M}}$ 是正合的。
\item 存在如下典范映射
$$
\widetilde{\mathcal{M}}
\otimes_{\mathcal{O}_{\underline{\text{Proj}}_S(\mathcal{A})}}
\widetilde{\mathcal{N}}
\longrightarrow
\widetilde{\mathcal{M} \otimes_\mathcal{A} \mathcal{N}}
$$
它们与引理 \ref{constructions-lemma-widetilde-tensor} 中的映射相同。
\item 存在典范映射
$$
\pi^*\mathcal{M}
\longrightarrow
\bigoplus\nolimits_{n \in \mathbf{Z}}
\widetilde{\mathcal{M}(n)}
$$
它们推广了 (\ref{constructions-equation-global-sections-more-generally})。
\item $\widetilde{\mathcal{M}}$ 的形成与基变换可交换。
\end{enumerate}
\end{lemma}

\begin{proof}
从略。我们应当把这个引理拆成若干部分，并分别证明各部分。
\end{proof}












\section{相对 Proj 的函子性}
\label{constructions-section-functoriality-relative-proj}

\noindent
本节是相对 Proj 情形下第 \ref{constructions-section-functoriality-proj} 节的类比。
设 $S$ 是一个概形。分次 $\mathcal{O}_S$-代数映射
$\psi : \mathcal{A} \to \mathcal{B}$ 并不总能诱导相伴相对 Proj 之间的态射。
正确的结果如下。

\begin{lemma}
\label{constructions-lemma-morphism-relative-proj}
设 $S$ 是一个概形。设 $\mathcal{A}$、$\mathcal{B}$ 是两个分次拟凝聚
$\mathcal{O}_S$-代数。置
$p : X = \underline{\text{Proj}}_S(\mathcal{A}) \to S$ 及
$q : Y = \underline{\text{Proj}}_S(\mathcal{B}) \to S$。设
$\psi : \mathcal{A} \to \mathcal{B}$ 是分次 $\mathcal{O}_S$-代数同态。
存在典范开集 $U(\psi) \subset Y$ 以及典范概形态射
$$
r_\psi :
U(\psi)
\longrightarrow
X
$$
它在 $S$ 上；此外还存在一个 $\mathbf{Z}$-分次
$\mathcal{O}_{U(\psi)}$-代数映射
$$
\theta = \theta_\psi :
r_\psi^*\left(
\bigoplus\nolimits_{d \in \mathbf{Z}} \mathcal{O}_X(d)
\right)
\longrightarrow
\bigoplus\nolimits_{d \in \mathbf{Z}} \mathcal{O}_{U(\psi)}(d).
$$
三元组 $(U(\psi), r_\psi, \theta)$ 由如下性质刻画：对任意仿射开集
$W \subset S$，三元组
$$
(U(\psi) \cap p^{-1}W,\quad
r_\psi|_{U(\psi) \cap p^{-1}W} : U(\psi) \cap p^{-1}W \to q^{-1}W,\quad
\theta|_{U(\psi) \cap p^{-1}W})
$$
等于与 $\psi : \mathcal{A}(W) \to \mathcal{B}(W)$ 相伴的三元组，
即引理 \ref{constructions-lemma-morphism-proj} 中的三元组；这里采用恒等
$p^{-1}W = \text{Proj}(\mathcal{A}(W))$ 与
$q^{-1}W = \text{Proj}(\mathcal{B}(W))$，它们来自第
\ref{constructions-section-relative-proj-via-glueing} 节。
\end{lemma}

\begin{proof}
把这些局部三元组粘合起来，便立即得到本引理。
\end{proof}

\begin{lemma}
\label{constructions-lemma-morphism-relative-proj-transitive}
设 $S$ 是一个概形。设 $\mathcal{A}$、$\mathcal{B}$ 和 $\mathcal{C}$ 是
拟凝聚分次 $\mathcal{O}_S$-代数。置
$X = \underline{\text{Proj}}_S(\mathcal{A})$、
$Y = \underline{\text{Proj}}_S(\mathcal{B})$ 及
$Z = \underline{\text{Proj}}_S(\mathcal{C})$。设
$\varphi : \mathcal{A} \to \mathcal{B}$、
$\psi : \mathcal{B} \to \mathcal{C}$ 是分次 $\mathcal{O}_S$-代数映射。
则有
$$
U(\psi \circ \varphi) = r_\varphi^{-1}(U(\psi))
\quad
\text{且}
\quad
r_{\psi \circ \varphi}
=
r_\varphi \circ r_\psi|_{U(\psi \circ \varphi)}.
$$
此外还有
$$
\theta_\psi \circ r_\psi^*\theta_\varphi
=
\theta_{\psi \circ \varphi}
$$
其中记号含义显然。
\end{lemma}

\begin{proof}
从略。
\end{proof}

\begin{lemma}
\label{constructions-lemma-surjective-graded-rings-map-relative-proj}
沿用上述引理 \ref{constructions-lemma-morphism-relative-proj} 的假设与记号。假设
$\mathcal{A}_d \to \mathcal{B}_d$ 对 $d \gg 0$ 是满射。则
\begin{enumerate}
\item $U(\psi) = Y$,
\item $r_\psi : Y \to X$ 是闭浸入，并且
\item 映射 $\theta : r_\psi^*\mathcal{O}_X(n) \to \mathcal{O}_Y(n)$
是满射，但一般不是同构（即使 $\mathcal{A} \to \mathcal{B}$ 是满射）。
\end{enumerate}
\end{lemma}

\begin{proof}
由引理 \ref{constructions-lemma-morphism-relative-proj} 与引理
\ref{constructions-lemma-surjective-graded-rings-map-proj} 合并即得。
\end{proof}

\begin{lemma}
\label{constructions-lemma-eventual-iso-graded-rings-map-relative-proj}
沿用上述引理 \ref{constructions-lemma-morphism-relative-proj} 的假设与记号。假设
$\mathcal{A}_d \to \mathcal{B}_d$ 对所有 $d \gg 0$ 都是同构。则
\begin{enumerate}
\item $U(\psi) = Y$,
\item $r_\psi : Y \to X$ 是同构，并且
\item 映射 $\theta : r_\psi^*\mathcal{O}_X(n) \to \mathcal{O}_Y(n)$
是同构。
\end{enumerate}
\end{lemma}

\begin{proof}
由引理 \ref{constructions-lemma-morphism-relative-proj} 与引理
\ref{constructions-lemma-eventual-iso-graded-rings-map-proj} 合并即得。
\end{proof}

\begin{lemma}
\label{constructions-lemma-surjective-generated-degree-1-map-relative-proj}
沿用上述引理 \ref{constructions-lemma-morphism-relative-proj} 的假设与记号。假设
$\mathcal{A}_d \to \mathcal{B}_d$ 对 $d \gg 0$ 是满射，并且
$\mathcal{A}$ 由 $\mathcal{A}_1$ 在 $\mathcal{A}_0$ 上生成。则
\begin{enumerate}
\item $U(\psi) = Y$,
\item $r_\psi : Y \to X$ 是闭浸入，并且
\item 映射 $\theta : r_\psi^*\mathcal{O}_X(n) \to \mathcal{O}_Y(n)$
是同构。
\end{enumerate}
\end{lemma}

\begin{proof}
由引理 \ref{constructions-lemma-morphism-relative-proj} 与引理
\ref{constructions-lemma-surjective-graded-rings-generated-degree-1-map-proj} 合并即得。
\end{proof}

















\section{可逆层与映入相对 Proj 的态射}
\label{constructions-section-invertible-relative-proj}

\noindent
下述引理似乎会在某处用到。情形如下：
\begin{enumerate}
\item 设 $S$ 是一个概形。
\item 设 $\mathcal{A}$ 是一个拟凝聚分次 $\mathcal{O}_S$-代数。
\item 记 $\pi : \underline{\text{Proj}}_S(\mathcal{A}) \to S$ 为 $S$ 上的相对齐次谱。
\item 设 $f : X \to S$ 是一个概形态射。
\item 设 $\mathcal{L}$ 是一个可逆 $\mathcal{O}_X$-模。
\item 设 $\psi : f^*\mathcal{A} \to
\bigoplus_{d \geq 0} \mathcal{L}^{\otimes d}$
是一个分次 $\mathcal{O}_X$-代数同态。
\end{enumerate}
给定这些数据，置
$$
U(\psi) = \bigcup\nolimits_{(U, V, a)} U_{\psi(a)}
$$
其中 $(U, V, a)$ 满足：
\begin{enumerate}
\item $V \subset S$ 是仿射开集，
\item $U = f^{-1}(V)$，并且
\item $a \in \mathcal{A}(V)_{+}$ 是齐次的。
\end{enumerate}
此时 $\psi(a) \in \Gamma(U, \mathcal{L}^{\otimes \deg(a)})$，而
$U_{\psi(a)}$ 是对应的开集（见《模》中的引理 \ref{modules-lemma-s-open}）。

\begin{lemma}
\label{constructions-lemma-invertible-map-into-relative-proj}
沿用上述假设与记号。态射 $\psi$ 诱导如下 $S$ 上的典范概形态射：
$$
r_{\mathcal{L}, \psi} :
U(\psi) \longrightarrow \underline{\text{Proj}}_S(\mathcal{A})
$$
并且还诱导如下分次 $\mathcal{O}_{U(\psi)}$-代数映射：
$$
\theta :
r_{\mathcal{L}, \psi}^*\left(
\bigoplus\nolimits_{d \geq 0}
\mathcal{O}_{\underline{\text{Proj}}_S(\mathcal{A})}(d)
\right)
\longrightarrow
\bigoplus\nolimits_{d \geq 0} \mathcal{L}^{\otimes d}|_{U(\psi)}
$$
它们由下列性质刻画：
\begin{enumerate}
\item 对每个开集 $V \subset S$ 以及每个 $d \geq 0$，图
$$
\xymatrix{
\mathcal{A}_d(V) \ar[d]_{\psi} \ar[r]_{\psi} &
\Gamma(f^{-1}(V), \mathcal{L}^{\otimes d}) \ar[d]^{restrict} \\
\Gamma(\pi^{-1}(V),
\mathcal{O}_{\underline{\text{Proj}}_S(\mathcal{A})}(d)) \ar[r]^{\theta} &
\Gamma(f^{-1}(V) \cap U(\psi), \mathcal{L}^{\otimes d})
}
$$
是交换的。
\item 对任意 $d \geq 1$ 以及任意开子概形 $W \subset X$，若
$\psi|_W : f^*\mathcal{A}_d|_W \to \mathcal{L}^{\otimes d}|_W$ 是满射，
则态射 $r_{\mathcal{L}, \psi}$ 的限制与态射
$W \to \underline{\text{Proj}}_S(\mathcal{A})$ 一致；后者由相对齐次谱的构造存在，
见定义 \ref{constructions-definition-relative-proj}。
\item 对任意仿射开集 $V \subset S$，限制
$$
(U(\psi) \cap f^{-1}(V), r_{\mathcal{L}, \psi}|_{U(\psi) \cap f^{-1}(V)},
\theta|_{U(\psi) \cap f^{-1}(V)})
$$
经由 $i_V$（见引理 \ref{constructions-lemma-glue-relative-proj}）与三元组
$(U(\psi'), r_{\mathcal{L}, \psi'}, \theta')$ 一致；这是引理
\ref{constructions-lemma-invertible-map-into-proj} 中与映射
$\psi' : A = \mathcal{A}(V) \to \Gamma_*(f^{-1}(V), \mathcal{L}|_{f^{-1}(V)})$
相伴的三元组，而该映射由 $\psi$ 诱导。
\end{enumerate}
\end{lemma}

\begin{proof}
利用刻画 (3) 构造态射 $r_{\mathcal{L}, \psi}$ 与 $\theta$；这个构造是在
$S$ 上局部进行的。利用引理 \ref{constructions-lemma-invertible-map-into-proj} 的唯一性说明这些构造能够粘合。
细节从略。
\end{proof}







\section{可逆层扭曲与相对 Proj}
\label{constructions-section-twisting-and-proj}

\noindent
设 $S$ 是一个概形。设
$\mathcal{A} = \bigoplus_{d \geq 0} \mathcal{A}_d$ 是一个拟凝聚分次
$\mathcal{O}_S$-代数。设 $\mathcal{L}$ 是 $S$ 上的一个可逆层。
在这种情形下，我们得到另一个拟凝聚分次 $\mathcal{O}_S$-代数，即
$$
\mathcal{B}
=
\bigoplus\nolimits_{d \geq 0}
\mathcal{A}_d \otimes_{\mathcal{O}_S} \mathcal{L}^{\otimes d}
$$
事实表明，$\mathcal{A}$ 与 $\mathcal{B}$ 具有同构的相对齐次谱。

\begin{lemma}
\label{constructions-lemma-twisting-and-proj}
沿用上述记号 $S$、$\mathcal{A}$、$\mathcal{L}$ 和 $\mathcal{B}$。
存在典范同构
$$
\xymatrix{
P = \underline{\text{Proj}}_S(\mathcal{A})
\ar[rr]_g \ar[rd]_\pi & &
\underline{\text{Proj}}_S(\mathcal{B}) = P'
\ar[ld]^{\pi'} \\
& S &
}
$$
它具有如下性质：
\begin{enumerate}
\item 存在同构
$\theta_n : g^*\mathcal{O}_{P'}(n)
\to
\mathcal{O}_P(n) \otimes \pi^*\mathcal{L}^{\otimes n}$
它们拼合为如下 $\mathbf{Z}$-分次代数同构：
$$
\theta :
g^*\left(
\bigoplus\nolimits_{n \in \mathbf{Z}} \mathcal{O}_{P'}(n)
\right)
\longrightarrow
\bigoplus\nolimits_{n \in \mathbf{Z}} \mathcal{O}_P(n)
\otimes \pi^*\mathcal{L}^{\otimes n}
$$
\item 对每个开集 $V \subset S$，图
$$
\xymatrix{
\mathcal{A}_n(V) \otimes \mathcal{L}^{\otimes n}(V)
\ar[r]_{multiply} \ar[d]^{\psi \otimes \pi^*}
&
\mathcal{B}_n(V) \ar[dd]^\psi \\
\Gamma(\pi^{-1}V, \mathcal{O}_P(n)) \otimes
\Gamma(\pi^{-1}V, \pi^*\mathcal{L}^{\otimes n})
\ar[d]^{multiply} \\
\Gamma(\pi^{-1}V, \mathcal{O}_P(n) \otimes \pi^*\mathcal{L}^{\otimes n})
&
\Gamma(\pi'^{-1}V, \mathcal{O}_{P'}(n)) \ar[l]_-{\theta_n}
}
$$
都是交换的。
\item 视需要在此补充更多内容。
\end{enumerate}
\end{lemma}

\begin{proof}
当 $\mathcal{L} \cong \mathcal{O}_S$ 时，这就是恒等映射。一般地，选取
$S$ 的一个开覆盖，使得 $\mathcal{L}$ 在各片上平凡化，再粘合相应的映射。
细节从略。
\end{proof}
















\section{射影丛}
\label{constructions-section-projective-bundle}

\noindent
设 $S$ 是一个概形。设 $\mathcal{E}$ 是一个拟凝聚 $\mathcal{O}_S$-模层。
由《模》中的引理 \ref{modules-lemma-whole-tensor-algebra-permanence}，
对称代数 $\text{Sym}(\mathcal{E})$，即 $\mathcal{E}$ 在 $\mathcal{O}_S$
上的对称代数，是一个拟凝聚 $\mathcal{O}_S$-代数层。注意，它由 $1$ 次部分在
$\mathcal{O}_S$ 上生成。因此可以对它应用上一节的构造，特别是引理
\ref{constructions-lemma-relative-proj} 和 \ref{constructions-lemma-apply-relative}。

\begin{definition}
\label{constructions-definition-projective-bundle}
设 $S$ 是一个概形。设 $\mathcal{E}$ 是一个拟凝聚
$\mathcal{O}_S$-模\footnote{读者或许预期这里要求 $\mathcal{E}$ 是有限局部自由的。
为了与 \cite[II, Definition 4.1.1]{EGA} 保持一致，我们不作此要求。}。
我们记
$$
\pi :
\mathbf{P}(\mathcal{E}) = \underline{\text{Proj}}_S(\text{Sym}(\mathcal{E}))
\longrightarrow
S
$$
并称它为{\it 与 $\mathcal{E}$ 相伴的射影丛}。记号
$\mathcal{O}_{\mathbf{P}(\mathcal{E})}(n)$ 表示引理
\ref{constructions-lemma-apply-relative} 中的可逆 $\mathcal{O}_{\mathbf{P}(\mathcal{E})}$-模，
称为{\it 结构层的第 $n$ 次扭曲}。
\end{definition}

\noindent
由引理 \ref{constructions-lemma-glue-relative-proj-twists}，存在典范 $\mathcal{O}_S$-模同态
$$
\text{Sym}^n(\mathcal{E})
\longrightarrow
\pi_*\mathcal{O}_{\mathbf{P}(\mathcal{E})}(n)
\quad\text{等价地}\quad
\pi^*\text{Sym}^n(\mathcal{E})
\longrightarrow
\mathcal{O}_{\mathbf{P}(\mathcal{E})}(n)
$$
对所有 $n \geq 0$ 都成立。特别地，当 $n = 1$ 时有
$$
\mathcal{E}
\longrightarrow
\pi_*\mathcal{O}_{\mathbf{P}(\mathcal{E})}(1)
\quad\text{等价地}\quad
\pi^*\mathcal{E}
\longrightarrow
\mathcal{O}_{\mathbf{P}(\mathcal{E})}(1)
$$
并且由引理 \ref{constructions-lemma-apply-relative}，映射
$\pi^*\mathcal{E} \to \mathcal{O}_{\mathbf{P}(\mathcal{E})}(1)$ 是满射。
这是记忆我们如何规范化 $\mathbf{P}(\mathcal{E})$ 的构造的一种好方法。

\medskip\noindent
警告：在有些文献中，概形 $\mathbf{P}(\mathcal{E})$ 仅当 $\mathcal{E}$
在 $S$ 上有限局部自由时才有定义。此外，有时 $\mathbf{P}(\mathcal{E})$
实际上被定义为我们记作 $\mathbf{P}(\mathcal{E}^\vee)$ 的概形，其中
$\mathcal{E}^\vee$ 是 $\mathcal{E}$ 的对偶（并且这种定义也只在
$\mathcal{E}$ 有限局部自由时采用）。

\medskip\noindent
设 $S$、$\mathcal{E}$、$\mathbf{P}(\mathcal{E}) \to S$ 如定义
\ref{constructions-definition-projective-bundle} 所述。设 $f : T \to S$ 是 $S$ 上的一个概形。
设 $\psi : f^*\mathcal{E} \to \mathcal{L}$ 是满射，其中 $\mathcal{L}$
是一个可逆 $\mathcal{O}_T$-模。所诱导的分次 $\mathcal{O}_T$-代数映射
$$
f^*\text{Sym}(\mathcal{E}) = \text{Sym}(f^*\mathcal{E}) \to
\text{Sym}(\mathcal{L}) = \bigoplus\nolimits_{n \geq 0} \mathcal{L}^{\otimes n}
$$
对应于态射
$$
\varphi_{\mathcal{L}, \psi} : T \longrightarrow \mathbf{P}(\mathcal{E})
$$
它位于 $S$ 上。这来自我们把相对 Proj 构造为表示第
\ref{constructions-section-relative-proj} 节中函子 $F$ 的概形。另一方面，给定态射
$\varphi : T \to \mathbf{P}(\mathcal{E})$（它位于 $S$ 上），可以置
$\mathcal{L} = \varphi^*\mathcal{O}_{\mathbf{P}(\mathcal{E})}(1)$，并令
$\psi : f^*\mathcal{E} \to \mathcal{L}$ 等于沿 $\varphi$ 拉回典范满射
$\pi^*\mathcal{E} \to \mathcal{O}_{\mathbf{P}(\mathcal{E})}(1)$ 所得的映射。
由引理 \ref{constructions-lemma-apply-relative}，这些构造给出如下两个集合之间的互逆双射：
偶对 $(\mathcal{L}, \psi)$ 的同构类集合，以及态射
$\varphi : T \to \mathbf{P}(\mathcal{E})$（它位于 $S$ 上）的集合。因此
$\mathbf{P}(\mathcal{E})$ 表示这样的函子：它给 $f : T \to S$ 配上所有
$\mathcal{O}_T$-模商，其中被取商的模是 $f^*\mathcal{E}$，并且这些商局部自由秩为
$1$。

\begin{example}[向量空间的射影空间]
\label{constructions-example-projective-space}
设 $k$ 是一个域。设 $V$ 是一个 $k$-向量空间。相应的{\it 射影空间}是 $k$-概形
$$
\mathbf{P}(V) = \text{Proj}(\text{Sym}(V))
$$
其中 $\text{Sym}(V)$ 是 $V$ 在 $k$ 上的对称代数。显然有
$\mathbf{P}(V) \cong \mathbf{P}^n_k$，只要 $\dim(V) = n + 1$，因为此时
$V$ 的对称代数同构于一个含 $n + 1$ 个变量的多项式环。若把 $V$ 看作
$\Spec(k)$ 上的拟凝聚模，则 $\mathbf{P}(V)$ 是 $\Spec(k)$ 上相应的射影空间丛。
由上述讨论，一个 $k$-值点 $p$（属于 $\mathbf{P}(V)$）对应于一个 $k$-向量空间满射
$V \to L_p$，其中 $\dim(L_p) = 1$。更一般地，设 $X$ 是 $k$ 上的概形，
设 $\mathcal{L}$ 是一个可逆 $\mathcal{O}_X$-模，并设
$\psi : V \to \Gamma(X, \mathcal{L})$ 是一个 $k$-线性映射，使得
$\mathcal{L}$ 作为 $\mathcal{O}_X$-模由 $\psi$ 像中的截面生成。
于是上述讨论给出典范态射
$$
\varphi_{\mathcal{L}, \psi} : X \longrightarrow \mathbf{P}(V)
$$
它是 $k$ 上的概形态射，并且存在同构
$\theta : \varphi_{\mathcal{L}, \psi}^*\mathcal{O}_{\mathbf{P}(V)}(1)
\to \mathcal{L}$，同时 $\psi$ 与下列复合一致：
$$
V \to
\Gamma(\mathbf{P}(V), \mathcal{O}_{\mathbf{P}(V)}(1))
\to
\Gamma(X, \varphi_{\mathcal{L}, \psi}^*\mathcal{O}_{\mathbf{P}(V)}(1))
\to
\Gamma(X, \mathcal{L})
$$
见引理 \ref{constructions-lemma-invertible-map-into-proj}。若
$V \subset \Gamma(X, \mathcal{L})$ 是一个子空间，则把上述构造的态射简记为
$\varphi_{\mathcal{L}, V}$。若 $\dim(V) = n + 1$，并选定一组基
$v_0, \ldots, v_n$（它是 $V$ 的基），则图
$$
\xymatrix{
X \ar@{=}[d] \ar[rr]_{\varphi_{\mathcal{L}, \psi}} & &
\mathbf{P}(V) \ar[d]^{\cong} \\
X \ar[rr]^{\varphi_{(\mathcal{L}, (s_0, \ldots, s_n))}} & &
\mathbf{P}^n_k
}
$$
是交换的，其中 $s_i = \psi(v_i) \in \Gamma(X, \mathcal{L})$，
$\varphi_{(\mathcal{L}, (s_0, \ldots, s_n))}$ 如第
\ref{constructions-section-projective-space} 节所述，而右侧竖直箭头对应于同构
$k[T_0, \ldots, T_n] \to \text{Sym}(V)$；该同构把 $T_i$ 映到 $v_i$。
\end{example}

\begin{example}
\label{constructions-example-projective-bundle}
映射 $\text{Sym}^n(\mathcal{E}) \to
\pi_*(\mathcal{O}_{\mathbf{P}(\mathcal{E})}(n))$
在 $\mathcal{E}$ 局部自由时是同构，但一般不必是同构。事实上，我们将给出一个该映射在 $n = 1$ 时
不是单射的例子。置 $S = \Spec(A)$，其中
$$
A = k[u, v, s_1, s_2, t_1, t_2]/I
$$
这里 $k$ 是一个域，并且
$$
I = (-us_1 + vt_1 + ut_2, vs_1 + us_2 - vt_2, vs_2, ut_1).
$$
记 $\overline{u}$ 为 $u$ 在 $A$ 中的类；其他变量也采用类似记号。
设 $M = (Ax \oplus Ay)/A(\overline{u}x + \overline{v}y)$，于是
$$
\text{Sym}(M) = A[x, y]/(\overline{u}x + \overline{v}y)
= k[x, y, u, v, s_1, s_2, t_1, t_2]/J
$$
其中
$$
J = (-us_1 + vt_1 + ut_2, vs_1 + us_2 - vt_2, vs_2, ut_1, ux + vy).
$$
在这种情形下，与拟凝聚层 $\mathcal{E} = \widetilde{M}$（它位于
$S = \Spec(A)$ 上）相伴的射影丛是概形
$$
P =
\text{Proj}(\text{Sym}(M)).
$$
注意，这个概形具有仿射开覆盖
$P = D_{+}(x) \cup D_{+}(y)$.
考虑元素 $m \in M$，它是元素 $us_1x + vt_2y$ 的像。注意
$$
x(us_1x + vt_2y) = (s_1x + s_2y)(ux + vy) \bmod I
$$
并且
$$
y(us_1x + vt_2y) = (t_1x + t_2y)(ux + vy) \bmod I.
$$
第一个等式蕴含 $m$ 作为 $\mathcal{O}_P(1)$ 在 $D_{+}(x)$ 上的截面映到零，
第二个等式则蕴含它作为 $\mathcal{O}_P(1)$ 在 $D_{+}(y)$ 上的截面映到零。
这说明 $m$ 在 $\Gamma(P, \mathcal{O}_P(1))$ 中映到零。另一方面，我们断言
$m \not = 0$；因此 $m$ 给出 $\mathcal{E}$ 的一个非零整体截面映到
$\Gamma(P, \mathcal{O}_P(1))$ 中零的例子。反设 $m = 0$，以求矛盾。
此时存在元素 $f \in k[u, v, s_1, s_2, t_1, t_2]$，使得
$$
us_1x + vt_2y = f(ux + vy) \bmod I
$$
由于 $I$ 由次数为 $2$ 的齐次多项式生成，我们可以把 $f$ 分解为各齐次分量，
并取其次数为 1 的分量。换言之，可以假设
$$
f = au + bv + \alpha_1s_1 + \alpha_2s_2 + \beta_1t_1 + \beta_2t_2
$$
其中 $a, b, \alpha_1, \alpha_2, \beta_1, \beta_2 \in k$。所得条件为
$$
\begin{matrix}
us_1 - u(au + bv + \alpha_1s_1 + \alpha_2s_2 + \beta_1t_1 + \beta_2t_2)
\in I \\
vt_2 - v(au + bv + \alpha_1s_1 + \alpha_2s_2 + \beta_1t_1 + \beta_2t_2)
\in I
\end{matrix}
$$
生成元中没有项 $u^2, uv, v^2$；这里说的是 $I$ 的生成元。因此可见 $a = b = 0$。
于是得到关系
$$
\begin{matrix}
us_1 - u(\alpha_1s_1 + \alpha_2s_2 + \beta_1t_1 + \beta_2t_2)
\in I \\
vt_2 - v(\alpha_1s_1 + \alpha_2s_2 + \beta_1t_1 + \beta_2t_2)
\in I
\end{matrix}
$$
可以用 $I$ 的第一个生成元把每次出现的 $us_1$ 替换为 $vt_1 + ut_2$，
用 $I$ 的第二个生成元把每次出现的 $vs_1$ 替换为 $-us_2 + vt_2$，
用第三个生成元消去 $vs_2$，再用第三个生成元消去 $ut_1$。
于是得到关系
$$
\begin{matrix}
(1 - \alpha_1)vt_1 + (1 - \alpha_1)ut_2 - \alpha_2us_2 - \beta_2ut_2 = 0 \\
(1 - \alpha_1)vt_2 + \alpha_1us_2 - \beta_1vt_1 - \beta_2vt_2 = 0
\end{matrix}
$$
这蕴含 $\alpha_1$ 必须同时等于 $0$ 和 $1$，这正是所需的矛盾。
\end{example}


\begin{lemma}
\label{constructions-lemma-projective-bundle-separated}
设 $S$ 是一个概形。结构态射 $\mathbf{P}(\mathcal{E}) \to S$
（它来自 $S$ 上的一个射影丛）是分离的。
\end{lemma}

\begin{proof}
由引理 \ref{constructions-lemma-relative-proj-separated} 立即可得。
\end{proof}

\begin{lemma}
\label{constructions-lemma-projective-space-bundle}
设 $S$ 是一个概形。设 $n \geq 0$。则 $\mathbf{P}^n_S$ 是 $S$ 上的一个射影丛。
\end{lemma}

\begin{proof}
注意
$$
\mathbf{P}^n_{\mathbf{Z}} =
\text{Proj}(\mathbf{Z}[T_0, \ldots, T_n]) =
\underline{\text{Proj}}_{\Spec(\mathbf{Z})}
\left(\widetilde{\mathbf{Z}[T_0, \ldots, T_n]}\right)
$$
其中环 $\mathbf{Z}[T_0, \ldots, T_n]$ 的分次由 $\deg(T_i) = 1$ 给出，
而 $\mathbf{Z}$ 的元素位于 $0$ 次。回忆 $\mathbf{P}^n_S$ 定义为
$\mathbf{P}^n_{\mathbf{Z}} \times_{\Spec(\mathbf{Z})} S$。
此外，相对齐次谱的形成与基变换可交换，见引理
\ref{constructions-lemma-relative-proj-base-change}。对任意概形
$g : S \to \Spec(\mathbf{Z})$，有
$g^*\mathcal{O}_{\Spec(\mathbf{Z})}[T_0, \ldots, T_n]
= \mathcal{O}_S[T_0, \ldots, T_n]$.
综合以上结果可见
$$
\mathbf{P}^n_S = \underline{\text{Proj}}_S(\mathcal{O}_S[T_0, \ldots, T_n]).
$$
最后注意
$\mathcal{O}_S[T_0, \ldots, T_n] = \text{Sym}(\mathcal{O}_S^{\oplus n + 1})$。
因此 $\mathbf{P}^n_S$ 是 $S$ 上的一个射影丛。
\end{proof}



\section{格拉斯曼概形}
\label{constructions-section-grassmannian}

\noindent
本节引入标准的格拉斯曼函子，并说明它们可由概形表示。选取整数 $k$、$n$，
满足 $0 < k < n$。我们将构造函子
\begin{equation}
\label{constructions-equation-gkn}
G(k, n) : \Sch \longrightarrow \textit{Sets}
\end{equation}
粗略地说，它参数化 $k$ 维子空间，这些子空间位于 $n$ 维空间中。
不过，出于技术原因，参数化 $(n - k)$ 维商空间更为方便，我们将采用这种做法。

\medskip\noindent
更准确地说，$G(k, n)$ 给概形 $S$ 配上集合 $G(k, n)(S)$，其元素是如下满射的同构类：
$$
q : \mathcal{O}_S^{\oplus n} \longrightarrow \mathcal{Q}
$$
其中 $\mathcal{Q}$ 是一个有限局部自由 $\mathcal{O}_S$-模，其秩为 $n - k$。
这确实是一个集合，例如可由《模》中的引理
\ref{modules-lemma-set-isomorphism-classes-finite-type-modules} 得知；也可以注意到，
满射 $q$ 的同构类由 $q$ 的核决定（而给定一个层，其子层构成一个集合）。
给定概形态射 $f : T \to S$，令
$G(k, n)(f) : G(k, n)(S) \to G(k, n)(T)$ 把
$q : \mathcal{O}_S^{\oplus n} \longrightarrow \mathcal{Q}$ 的同构类映到
$f^*q : \mathcal{O}_T^{\oplus n} \longrightarrow f^*\mathcal{Q}$ 的同构类。
这是有意义的，因为 (1) $f^*\mathcal{O}_S = \mathcal{O}_T$；
(2) $f^*$ 是可加的；(3) $f^*$ 保持局部自由模（《模》，引理
\ref{modules-lemma-pullback-locally-free}）；(4) $f^*$ 是右正合的（《模》，引理
\ref{modules-lemma-exactness-pushforward-pullback}）。

\begin{lemma}
\label{constructions-lemma-gkn-representable}
设 $0 < k < n$。则 (\ref{constructions-equation-gkn}) 中的函子 $G(k, n)$ 可由概形表示。
\end{lemma}

\begin{proof}
置 $F = G(k, n)$。为证明本引理，我们使用《概形》中的引理
\ref{schemes-lemma-glue-functors} 所给判据。$F$ 对 Zariski 拓扑满足层性质，
这是因为层可以粘合；见《层》第 \ref{sheaves-section-glueing-sheaves} 节
（略去一些细节）。

\medskip\noindent
子函子族 $F_i$。设 $I$ 是 $\{1, \ldots, n\}$ 的所有基数为 $n - k$
的子集所成的集合。给定概形 $S$ 以及 $j \in \{1, \ldots, n\}$，
记 $e_j$ 为整体截面
$$
e_j = (0, \ldots, 0, 1, 0, \ldots, 0)\quad(1\text{ 于 }j\text{次位置})
$$
它属于 $\mathcal{O}_S^{\oplus n}$。显然，这些截面自由生成
$\mathcal{O}_S^{\oplus n}$。类似地，对 $j \in \{1, \ldots, n - k\}$，
记 $f_j$ 为 $\mathcal{O}_S^{\oplus n - k}$ 的整体截面；除第 $j$ 个分量取
$1$ 外，其余分量均为零。对 $i \in I$，令
$$
s_i : \mathcal{O}_S^{\oplus n - k} \longrightarrow \mathcal{O}_S^{\oplus n}
$$
它是余投射 $\mathcal{O}_S \to \mathcal{O}_S^{\oplus n}$ 的直和；
这些余投射对应于 $I$ 的元素。更准确地，若
$i = \{i_1, \ldots, i_{n - k}\}$ 且 $i_1 < i_2 < \ldots < i_{n - k}$，
则 $s_i$ 把 $f_j$ 映到 $e_{i_j}$，其中 $j \in \{1, \ldots, n - k\}$。
用此记号，置
$$
F_i(S) = \{q : \mathcal{O}_S^{\oplus n} \to \mathcal{Q} \in F(S) \mid
q \circ s_i \text{ 为满射}\}
\subset F(S)
$$
给定概形态射 $f : T \to S$，拉回 $f^*s_i$ 是 $T$ 上相应的映射。
由于 $f^*$ 是右正合的（《模》，引理
\ref{modules-lemma-exactness-pushforward-pullback}），可知 $F_i$ 是 $F$ 的子函子。

\medskip\noindent
$F_i$ 的可表示性。为证明这一点，重新编号后可以假设
$i = \{1, \ldots, n - k\}$。这意味着 $s_i$ 是前 $n - k$ 个直和分量的嵌入。
注意，若 $q \circ s_i$ 是满射，则 $q \circ s_i$ 是同秩有限局部自由模之间的满射，
因而是同构（《模》，引理 \ref{modules-lemma-map-finite-locally-free}）。
因此，若 $q : \mathcal{O}_S^{\oplus n} \to \mathcal{Q}$ 是 $F_i(S)$ 的元素，
则可以用 $q \circ s_i$ 把 $\mathcal{Q}$ 与 $\mathcal{O}_S^{\oplus n - k}$ 等同。
这样便得到
$$
q : \mathcal{O}_S^{\oplus n} \longrightarrow \mathcal{O}_S^{\oplus n - k}
$$
它把 $e_j$ 映到 $f_j$（记号同上），其中 $j = 1, \ldots, n - k$。
为了完全确定 $q$，必须指定
$q(e_{n - k + 1}), \ldots, q(e_n)$ 在
$\Gamma(S, \mathcal{O}_S^{\oplus n - k})$ 中的像。因此 $F_i$ 同构于函子
$$
S \longmapsto
\prod\nolimits_{j = n - k + 1, \ldots, n}
\Gamma(S,  \mathcal{O}_S^{\oplus n - k})
$$
这个函子同构于 $k(n - k)$ 重自积；该自积取自函子
$S \mapsto \Gamma(S, \mathcal{O}_S)$。由《概形》中的例
\ref{schemes-example-global-sections}，
后一个函子由 $\mathbf{A}^1_\mathbf{Z}$ 表示。因此 $F_i$ 由
$\mathbf{A}^{k(n - k)}_\mathbf{Z}$ 表示，因为在概形范畴中，积就是在
$\Spec(\mathbf{Z})$ 上的纤维积。

\medskip\noindent
嵌入 $F_i \subset F$ 可由开浸入表示。设 $S$ 是一个概形，并设
$q : \mathcal{O}_S^{\oplus n} \to \mathcal{Q}$ 是 $F(S)$ 的元素。
由《模》中的引理 \ref{modules-lemma-finite-type-surjective-on-stalk}，
集合 $U_i = \{s \in S \mid (q \circ s_i)_s\text{ 满射}\}$
在 $S$ 中是开的。由于 $\mathcal{O}_{S, s}$ 是局部环，而 $\mathcal{Q}_s$
是有限 $\mathcal{O}_{S, s}$-模，由 Nakayama 引理（《代数》，引理
\ref{algebra-lemma-NAK}）可得
$$
s \in U_i \Leftrightarrow
\left(
\text{映射 }
\kappa(s)^{\oplus n - k} \to \mathcal{Q}_s/\mathfrak m_s\mathcal{Q}_s
\text{ 由下列对象诱导： }
(q \circ s_i)_s
\text{ 为满射}
\right)
$$
设 $f : T \to S$ 是一个概形态射，并设 $t \in T$ 是映到 $s \in S$ 的点。有
$(f^*\mathcal{Q})_t =
\mathcal{Q}_s \otimes_{\mathcal{O}_{S, s}} \mathcal{O}_{T, t}$
（《层》，引理 \ref{sheaves-lemma-stalk-pullback-modules}），其余类似。因此映射
$$
\kappa(t)^{\oplus n - k} \to (f^*\mathcal{Q})_t/\mathfrak m_t(f^*\mathcal{Q})_t
$$
由 $(f^*q \circ f^*s_i)_t$ 诱导，它是上述映射
$\kappa(s)^{\oplus n - k} \to \mathcal{Q}_s/\mathfrak m_s\mathcal{Q}_s$
沿域扩张 $\kappa(t)/\kappa(s)$ 的基变换。因此 $s \in U_i$ 当且仅当 $t$
属于 $f^*q$ 所对应的开集。特别地，$T \to S$ 通过 $U_i$ 分解，当且仅当
$f^*q \in F_i(T)$，这正是所需的结论。

\medskip\noindent
族 $F_i$（$i \in I$）覆盖 $F$。设
$q : \mathcal{O}_S^{\oplus n} \to \mathcal{Q}$ 是 $F(S)$ 的元素。
我们必须说明，对每个点 $s$（它属于 $S$），都存在 $i \in I$，使得 $s_i$
在 $s$ 的某个邻域上是满射。因此必须说明，下列复合之一
$$
\kappa(s)^{\oplus n - k} \xrightarrow{s_i}
\kappa(s)^{\oplus n} \rightarrow
\mathcal{Q}_s/\mathfrak m_s\mathcal{Q}_s
$$
是满射（见上一段）。由于 $\mathcal{Q}_s/\mathfrak m_s\mathcal{Q}_s$
是维数为 $n - k$ 的向量空间，这由向量空间理论可得。
\end{proof}

\begin{definition}
\label{constructions-definition-grassmannian}
设 $0 < k < n$。概形 $\mathbf{G}(k, n)$ 表示函子 $G(k, n)$，它
称为{\it $\mathbf{Z}$ 上的格拉斯曼概形}。其基变换
$\mathbf{G}(k, n)_S$（取到概形 $S$）称为{\it $S$ 上的格拉斯曼概形}。若 $R$ 是一个环，
则到 $\Spec(R)$ 的基变换记为 $\mathbf{G}(k, n)_R$，称为
{\it $R$ 上的格拉斯曼概形}。
\end{definition}

\noindent
这个定义是有意义的，因为我们已在引理 \ref{constructions-lemma-gkn-representable} 中说明这些函子确实可表示。

\begin{lemma}
\label{constructions-lemma-projective-space-grassmannian}
设 $n \geq 1$。存在典范同构
$\mathbf{G}(n, n + 1) = \mathbf{P}^n_\mathbf{Z}$。
\end{lemma}

\begin{proof}
由引理 \ref{constructions-lemma-projective-space}，概形 $\mathbf{P}^n_\mathbf{Z}$ 表示如下函子：
它给概形 $S$ 配上偶对 $(\mathcal{L}, (s_0, \ldots, s_n))$ 的同构类集合；
这样的偶对由可逆模 $\mathcal{L}$ 以及一个 $(n + 1)$ 元整体截面组构成，
该截面组生成 $\mathcal{L}$。给定这样的偶对，得到一个商
$$
\mathcal{O}_S^{\oplus n + 1} \longrightarrow \mathcal{L},\quad
(h_0, \ldots, h_n) \longmapsto \sum h_i s_i.
$$
反过来，给定元素 $q : \mathcal{O}_S^{\oplus n + 1} \to \mathcal{Q}$，
它属于 $G(n, n + 1)(S)$，便得到这样的偶对，即
$(\mathcal{Q}, (q(e_1), \ldots, q(e_{n + 1})))$。这里 $e_i$（
$i = 1, \ldots, n + 1$）是自由模 $\mathcal{O}_S^{\oplus n + 1}$ 的标准生成截面。
这些构造定义函子之间互逆变换的验证从略。
\end{proof}
% Generated mechanically from the frozen properties.tex target.
% Do not edit this generated file; edit the bound locale source instead.

% OK, start here.
%

\setcounter{chapter}{27}
\chapter{概形的性质}



\phantomsection
\label{properties-section-phantom}


\section{引言}
\label{properties-section-introduction}

\noindent
本章介绍概形的一些绝对性质。
基础参考文献是 \cite{EGA}。




\section{可构造集}
\label{properties-section-constructible}

\noindent
可构造集与局部可构造集在《拓扑》第
\ref{topology-section-constructible} 节中引入。
概形的局部可构造子集可刻画如下。

\begin{lemma}
\label{properties-lemma-locally-constructible}
令 $X$ 为概形。
子集 $E$ 在 $X$ 中局部可构造，当且仅当在 $X$ 上，
$E \cap U$ 在 $U$ 中可构造，其中 $U$ 取遍 $X$ 的仿射开集。
\end{lemma}

\begin{proof}
假设 $E$ 局部可构造。则存在开覆盖
$X = \bigcup U_i$，使得 $E \cap U_i$ 在 $U_i$ 中可构造，而这对每个 $i$ 成立。
令 $V \subset X$ 为任意仿射开集。可以找到有限仿射开覆盖
$V = V_1 \cup \ldots \cup V_m$，使得对每个 $j$，都有
$V_j \subset U_i$，其中 $i = i(j)$。由《拓扑》引理
\ref{topology-lemma-open-immersion-constructible-inverse-image}，
每个 $E \cap V_j$ 在 $V_j$ 中可构造。由于包含映射
$V_j \to V$ 拟紧（见《概形》引理
\ref{schemes-lemma-quasi-compact-affine}），由《拓扑》引理
\ref{topology-lemma-collate-constructible} 可知，$E \cap V$ 在 $V$ 中可构造。
逆向蕴含是立即的。
\end{proof}

\begin{lemma}
\label{properties-lemma-generic-point-in-constructible}
令 $X$ 为概形，$E \subset X$ 为局部可构造子集。
令 $\xi \in X$ 为 $X$ 的一个不可约分支的泛点。
\begin{enumerate}
\item 若 $\xi \in E$，则 $\xi$ 的某个开邻域包含于 $E$。
\item 若 $\xi \not \in E$，则 $\xi$ 的某个开邻域与 $E$ 不相交。
\end{enumerate}
\end{lemma}

\begin{proof}
局部可构造子集的补集仍局部可构造，故只需证明 (2)。可以假设 $X$
仿射，因而 $E$ 可构造（引理 \ref{properties-lemma-locally-constructible}）。
在这种情形下，$X$ 是谱空间
（《代数》引理 \ref{algebra-lemma-spec-spectral}）。于是，
$\xi \not \in E$ 蕴含 $\xi \not \in \overline{E}$；这是由《拓扑》引理
\ref{topology-lemma-constructible-stable-specialization-closed}
以及如下事实得出的：$X$ 中没有异于 $\xi$ 且特化到 $\xi$ 的点。
\end{proof}

\begin{lemma}
\label{properties-lemma-quasi-separated-quasi-compact-open-retrocompact}
令 $X$ 为拟分离概形。$X$ 的任意两个拟紧开集之交是 $X$ 的拟紧开集。
$X$ 的每个拟紧开集都在 $X$ 中逆拟紧。
\end{lemma}

\begin{proof}
若 $U$ 与 $V$ 是拟紧开集，则
$U \cap V = \Delta^{-1}(U \times V)$，其中 $\Delta : X \to X \times X$
是对角映射。由于 $X$ 拟分离，$\Delta$ 拟紧。可知 $U \cap V$ 拟紧，
因为 $U \times V$ 拟紧（略去细节；使用《概形》引理
\ref{schemes-lemma-affine-covering-fibre-product}
可知 $U \times V$ 是有限多个仿射概形之并）。
其余断言由第一个断言及《拓扑》引理
\ref{topology-lemma-topology-quasi-separated-scheme} 得出。
\end{proof}

\begin{lemma}
\label{properties-lemma-quasi-compact-quasi-separated-spectral}
令 $X$ 为拟紧且拟分离的概形。
则 $X$ 的底拓扑空间是谱空间。
\end{lemma}

\begin{proof}
根据《拓扑》定义 \ref{topology-definition-spectral-space}，
需要检验 $X$ 清醒、拟紧，具有由拟紧开集组成的拓扑基，并且任意两个拟紧开集
之交仍拟紧。这由《概形》引理 \ref{schemes-lemma-scheme-sober}、
\ref{schemes-lemma-basis-affine-opens}
以及上面的引理
\ref{properties-lemma-quasi-separated-quasi-compact-open-retrocompact}
得出。
\end{proof}

\begin{lemma}
\label{properties-lemma-constructible-quasi-compact-quasi-separated}
令 $X$ 为拟紧且拟分离的概形。
$X$ 的任意局部可构造子集都是可构造的。
\end{lemma}

\begin{proof}
由于 $X$ 拟紧，可以选取有限仿射开覆盖
$X = V_1 \cup \ldots \cup V_m$。由于 $X$ 拟分离，根据引理
\ref{properties-lemma-quasi-separated-quasi-compact-open-retrocompact}，每个 $V_i$
都在 $X$ 中逆拟紧。因此，由《拓扑》引理
\ref{topology-lemma-collate-constructible} 可知，$E \subset X$ 在 $X$
中可构造，当且仅当各个 $E \cap V_j$ 在 $V_j$ 中可构造。
于是由引理 \ref{properties-lemma-locally-constructible} 得到结论。
\end{proof}

\begin{lemma}
\label{properties-lemma-retrocompact}
令 $X$ 为概形。子集 $E$ 在 $X$ 中逆拟紧，当且仅当在 $X$ 上，
$E \cap U$ 拟紧，其中 $U$ 取遍 $X$ 的仿射开集。
\end{lemma}

\begin{proof}
这立即来自如下事实：$X$ 的每个拟紧开集都是有限多个仿射开集之并。
\end{proof}

\begin{lemma}
\label{properties-lemma-stratification-locally-finite-constructible}
一个分割 $X = \coprod_{i \in I} X_i$，其中 $X$ 是概形，若各部分均逆拟紧，
则该分割局部有限，当且仅当各部分均局部可构造。
\end{lemma}

\begin{proof}
关于逆拟紧、分割和局部有限的定义，见《拓扑》的定义
\ref{topology-definition-quasi-compact}、
\ref{topology-definition-paritition} 和
\ref{topology-definition-locally-finite}。

\medskip\noindent
若该分割局部有限，且 $U \subset X$ 是仿射开集，则
$U = \coprod_{i \in I} U \cap X_i$ 是有限分割（更准确地说，除有限多个部分外，
其余部分均为空）。因此 $U \cap X_i$ 拟紧，且它在 $U$ 中的补集作为有限多个
逆拟紧部分之并，仍逆拟紧。于是由《拓扑》引理
\ref{topology-lemma-locally-closed-constructible-image}，$U \cap X_i$
可构造。再由引理 \ref{properties-lemma-locally-constructible}，$X_i$ 局部可构造。

\medskip\noindent
假设各部分均局部可构造。则对任意仿射开集 $U \subset X$，得到由可构造子集
组成的覆盖 $U = \coprod X_i \cap U$。可构造拓扑是拟紧的，见《拓扑》引理
\ref{topology-lemma-constructible-hausdorff-quasi-compact}，
故此覆盖具有有限加细，即该分割局部有限。
\end{proof}




\section{整概形、不可约概形与约化概形}
\label{properties-section-integral}

\begin{definition}
\label{properties-definition-integral}
令 $X$ 为概形。若 $X$ 非空，并且对每个非空仿射开集
$\Spec(R) = U \subset X$，环 $R$ 都是整环，则称它是{it 整的}。
\end{definition}

\begin{lemma}
\label{properties-lemma-characterize-reduced}
令 $X$ 为概形。
下列条件等价。
\begin{enumerate}
\item 概形 $X$ 是约化的，见《概形》定义
\ref{schemes-definition-reduced}。
\item 存在仿射开覆盖 $X = \bigcup U_i$，使得每个
$\Gamma(U_i, \mathcal{O}_X)$ 都约化。
\item 对每个仿射开集 $U \subset X$，环
$\mathcal{O}_X(U)$ 都约化。
\item 对每个开集 $U \subset X$，环 $\mathcal{O}_X(U)$ 都约化。
\end{enumerate}
\end{lemma}

\begin{proof}
见《概形》引理 \ref{schemes-lemma-reduced} 和
\ref{schemes-lemma-affine-reduced}。
\end{proof}

\begin{lemma}
\label{properties-lemma-characterize-irreducible}
令 $X$ 为概形。
下列条件等价。
\begin{enumerate}
\item 概形 $X$ 不可约。
\item 存在仿射开覆盖 $X = \bigcup_{i \in I} U_i$，使得 $I$ 非空，
$U_i$ 对所有 $i \in I$ 都不可约，并且
$U_i \cap U_j \not = \emptyset$ 对所有 $i, j \in I$ 都成立。
\item 概形 $X$ 非空，并且每个非空仿射开集 $U \subset X$ 都不可约。
\end{enumerate}
\end{lemma}

\begin{proof}
假设 (1)。由《概形》引理 \ref{schemes-lemma-scheme-sober}，可知
$X$ 有唯一泛点 $\eta$。于是
$X = \overline{\{\eta\}}$。故 $\eta$ 属于每个非空仿射开集
$U \subset X$。这蕴含 $\eta \in U$ 稠密，因而 $U$ 不可约。
它还蕴含任意两个非空仿射开集相交。
因此 (1) 同时蕴含 (2) 与 (3)。

\medskip\noindent
假设 (2)。设 $X = Z_1 \cup Z_2$ 是两个闭子集之并。
对每个 $i$，要么 $U_i \subset Z_1$，要么 $U_i \subset Z_2$。
选取某个 $i \in I$，并假设 $U_i \subset Z_1$（必要时将
$Z_1$、$Z_2$ 重新编号）。对任意 $j \in I$，开子集
$U_i \cap U_j$ 在 $U_j$ 中稠密，并包含于闭子集 $Z_1 \cap U_j$。
由此可知 $U_j \subset Z_1$。故如所需，$X = Z_1$。

\medskip\noindent
假设 (3)。选取仿射开覆盖 $X = \bigcup_{i \in I} U_i$。
可以假设每个 $U_i$ 都非空。
由于 $X$ 非空，可知 $I$ 非空。
由假设，每个 $U_i$ 都不可约。
假设 $U_i \cap U_j = \emptyset$ 对某对 $i, j \in I$ 成立。
则开集 $U_i \amalg  U_j = U_i \cup U_j$ 是仿射的，见《概形》引理
\ref{schemes-lemma-disjoint-union-affines}。
故由假设它不可约，这与上式矛盾。因此 (3) 蕴含 (2)。引理得证。
\end{proof}

\begin{lemma}
\label{properties-lemma-characterize-integral}
概形 $X$ 是整的，当且仅当它既约化又不可约。
\end{lemma}

\begin{proof}
若 $X$ 不可约，则每个仿射开集 $\Spec(R) = U \subset X$
都不可约。若 $X$ 约化，则由上面的引理 \ref{properties-lemma-characterize-reduced}，
$R$ 约化。因此 $R$ 约化且 $(0)$ 是素理想，即 $R$ 是整环。

\medskip\noindent
若 $X$ 是整的，则对每个非空仿射开集
$\Spec(R) = U \subset X$，环 $R$ 都约化，
因而由引理 \ref{properties-lemma-characterize-reduced}，$X$ 约化。
此外，每个非空仿射开集都不可约。
因此 $X$ 不可约，见引理 \ref{properties-lemma-characterize-irreducible}。
\end{proof}

\noindent
在《例》第
\ref{examples-section-connected-locally-integral-not-integral} 节中，
我们构造了一个连通仿射概形，它的所有局部环都是整环，但它并不是整概形。










\section{由环的性质定义的概形类型}
\label{properties-section-properties-rings}

\noindent
本节研究环的哪些性质能够用来定义概形的局部性质。

\begin{definition}
\label{properties-definition-property-local}
令 $P$ 为环的一种性质。
若下列条件成立，则称 $P$ 是{it 局部的}：
\begin{enumerate}
\item 对任意环 $R$ 及任意 $f \in R$，有
$P(R) \Rightarrow P(R_f)$。
\item 对任意环 $R$ 及满足下列条件的 $f_i \in R$：
$(f_1, \ldots, f_n) = R$，有
$\forall i, P(R_{f_i}) \Rightarrow P(R)$。
\end{enumerate}
\end{definition}

\begin{definition}
\label{properties-definition-locally-P}
令 $P$ 为环的一种性质。令 $X$ 为概形。
称 $X${\it 局部具有 $P$ 性质}，如果对任意 $x \in X$，存在仿射开邻域
$U$，它是 $x$ 在 $X$ 中的邻域，并且 $\mathcal{O}_X(U)$ 具有性质 $P$。
\end{definition}

\noindent
只有当该性质是局部性质时，这才是一个合适的概念。
即使 $P$ 是局部性质，我们也不会自动使用这个定义，说某个概形
“局部具有 $P$ 性质”，除非我们还在别处明确给出该定义。

\begin{lemma}
\label{properties-lemma-locally-P}
令 $X$ 为概形。令 $P$ 为环的一种局部性质。
下列条件等价：
\begin{enumerate}
\item 概形 $X$ 局部具有 $P$ 性质。
\item 对每个仿射开集 $U \subset X$，性质
$P(\mathcal{O}_X(U))$ 成立。
\item 存在仿射开覆盖 $X = \bigcup U_i$，使得每个
$\mathcal{O}_X(U_i)$ 都满足 $P$。
\item 存在开覆盖 $X = \bigcup X_j$，使得每个开子概形 $X_j$
都局部具有 $P$ 性质。
\end{enumerate}
此外，若 $X$ 局部具有 $P$ 性质，则每个开子概形也局部具有 $P$ 性质。
\end{lemma}

\begin{proof}
显然 (1) $\Leftrightarrow$ (3)，且 (2) $\Rightarrow$ (1)。
若 (3) $\Rightarrow$ (2)，则引理的最后一个陈述成立，并容易推出 (4)
也等价于 (1)。因此只需证明 (3) $\Rightarrow$ (2)。

\medskip\noindent
令 $X = \bigcup U_i$ 为仿射开覆盖，记
$U_i = \Spec(R_i)$。假设 $P(R_i)$。
令 $\Spec(R) = U \subset X$ 为任意仿射开集。
由《概形》引理 \ref{schemes-lemma-good-subcover}，存在 $U = \Spec(R)$
的一个标准覆盖，由标准开集 $D(f_j)$ 组成，使得每个环 $R_{f_j}$
都是某个环 $R_i$ 的主局部化。由定义
\ref{properties-definition-property-local} (1)，得到 $P(R_{f_j})$。
进而由定义 \ref{properties-definition-property-local} (2)，得到 $P(R)$。
\end{proof}

\noindent
下面给出一个应用示例。

\begin{lemma}
\label{properties-lemma-reduced-is-locally-reduced}
令 $X$ 为概形。则 $X$ 约化，当且仅当 $X$ 在定义
\ref{properties-definition-locally-P} 的意义下“局部约化”。
\end{lemma}

\begin{proof}
这由引理 \ref{properties-lemma-characterize-reduced} 立即可见。
\end{proof}

\begin{lemma}
\label{properties-lemma-properties-local}
环 $R$ 的下列性质是局部性质。
\begin{enumerate}
\item （Cohen--Macaulay。）
环 $R$ 是 Noether 环且为 CM，见《代数》定义
\ref{algebra-definition-ring-CM}。
\item （正则。）
环 $R$ 是 Noether 环且正则，见《代数》定义
\ref{algebra-definition-regular}。
\item （绝对 Noether。）
环 $R$ 在 $Z$ 上有限型。
\item 可按需补充更多性质。\footnote{但这里只列出那些我们尚未在别处分别处理的性质。}
\end{enumerate}
\end{lemma}

\begin{proof}
略。
\end{proof}















\section{Noether 概形}
\label{properties-section-noetherian}

\noindent
回顾：若环 $R$ 满足理想的升链条件，则称它是{it Noether 环}。
等价地，$R$ 的每个理想都有限生成。

\begin{definition}
\label{properties-definition-noetherian}
令 $X$ 为概形。
\begin{enumerate}
\item 称 $X$ 是{it 局部 Noether 的}，如果每个 $x \in X$ 都有一个仿射开邻域
$\Spec(R) = U \subset X$，使得环 $R$ 是 Noether 环。
\item 若 $X$ 局部 Noether 且拟紧，则称 $X$ 是{it Noether 的}。
\end{enumerate}
\end{definition}

\noindent
下面是刻画局部 Noether 概形的标准结果。

\begin{lemma}
\label{properties-lemma-locally-Noetherian}
令 $X$ 为概形。下列条件等价：
\begin{enumerate}
\item 概形 $X$ 局部 Noether。
\item 对每个仿射开集 $U \subset X$，环 $\mathcal{O}_X(U)$
都是 Noether 环。
\item 存在仿射开覆盖 $X = \bigcup U_i$，使得每个
$\mathcal{O}_X(U_i)$ 都是 Noether 环。
\item 存在开覆盖 $X = \bigcup X_j$，使得每个开子概形 $X_j$
都局部 Noether。
\end{enumerate}
此外，若 $X$ 局部 Noether，则每个开子概形也局部 Noether。
\end{lemma}

\begin{proof}
只需证明 Noether 性是环的一种局部性质，见引理 \ref{properties-lemma-locally-P}。
Noether 环的任意局部化仍是 Noether 环，见《代数》引理
\ref{algebra-lemma-Noetherian-permanence}。
由《代数》引理 \ref{algebra-lemma-cover}，可知定义
\ref{properties-definition-property-local} 中的第二个性质成立。
\end{proof}

\begin{lemma}
\label{properties-lemma-immersion-into-noetherian}
若浸入 $Z \to X$ 的目标 $X$ 局部 Noether，则该浸入拟紧。
\end{lemma}

\begin{proof}
闭浸入显然拟紧。
拟紧态射的复合仍拟紧，见《拓扑》引理
\ref{topology-lemma-composition-quasi-compact}。
因此只需证明到局部 Noether 概形中的开浸入拟紧。
利用《概形》引理 \ref{schemes-lemma-quasi-compact-affine}，
可化约到 $X$ 仿射的情形。
Noether 环的谱的任意开子集都拟紧（例如，将《代数》引理
\ref{algebra-lemma-Noetherian-topology} 与《拓扑》引理
\ref{topology-lemma-Noetherian}、
\ref{topology-lemma-Noetherian-quasi-compact} 结合即可）。
\end{proof}

\begin{lemma}
\label{properties-lemma-locally-Noetherian-quasi-separated}
局部 Noether 概形是拟分离的。
\end{lemma}

\begin{proof}
由《概形》引理 \ref{schemes-lemma-characterize-quasi-separated}，
需要证明两个仿射开集之交 $U \cap V$ 在 $X$ 中拟紧。
例如，考虑开浸入 $U \cap V \to U$，结论就由上面的引理
\ref{properties-lemma-immersion-into-noetherian} 得出。
（其实这只不过是因为 Noether 环的谱的任意开集都拟紧。）
\end{proof}

\begin{lemma}
\label{properties-lemma-Noetherian-topology}
一个（局部）Noether 概形的底拓扑空间是（局部）Noether 的，
见《拓扑》定义 \ref{topology-definition-noetherian}。
\end{lemma}

\begin{proof}
这是因为 Noether 概形是有限多个 Noether 环的谱之并，并可应用
《代数》引理 \ref{algebra-lemma-Noetherian-topology} 与
《拓扑》引理 \ref{topology-lemma-finite-union-Noetherian}。
\end{proof}

\begin{lemma}
\label{properties-lemma-locally-closed-in-Noetherian}
一个（局部）Noether 概形的任意局部闭子概形都是（局部）Noether 的。
\end{lemma}

\begin{proof}
略。提示：Noether 环的任意商环与任意局部化都是 Noether 环。
对于 Noether 情形，再次使用如下事实：Noether 空间的任意子集在诱导拓扑下
都是 Noether 空间。
\end{proof}

\begin{lemma}
\label{properties-lemma-Noetherian-irreducible-components}
Noether 概形只有有限多个不可约分支。
\end{lemma}

\begin{proof}
Noether 概形的底拓扑空间是 Noether 的
（引理 \ref{properties-lemma-Noetherian-topology}），而 Noether 拓扑空间只有有限多个
不可约分支（《拓扑》引理 \ref{topology-lemma-Noetherian}），故结论成立。
\end{proof}

\begin{lemma}
\label{properties-lemma-morphism-Noetherian-schemes-quasi-compact}
概形态射 $f : X \to Y$ 若具有 Noether 源概形 $X$，则它拟紧。
\end{lemma}

\begin{proof}
使用引理 \ref{properties-lemma-Noetherian-topology}，并使用 Noether 拓扑空间的任意子集
都是拟紧的这一事实（见《拓扑》引理
\ref{topology-lemma-Noetherian} 和
\ref{topology-lemma-Noetherian-quasi-compact}）。
\end{proof}

\noindent
下面是一个有趣的引理。
它说明每个局部 Noether 概形都有大量闭点（每个闭子集中至少有一个）。

\begin{lemma}
\label{properties-lemma-locally-Noetherian-closed-point}
任意非空局部 Noether 概形都有闭点。
局部 Noether 概形的任意非空闭子集都有闭点。
等价地，局部 Noether 概形的任意一点都特化到一个闭点。
\end{lemma}

\begin{proof}
第二个断言由第一个断言得出（使用《概形》引理
\ref{schemes-lemma-reduced-closed-subscheme} 与引理
\ref{properties-lemma-locally-closed-in-Noetherian}）。
考虑任意非空仿射开集 $U \subset X$。
令 $x \in U$ 为闭点。若 $x$ 是 $X$ 的闭点，则已经完成。
否则，令 $X_0 \subset X$ 为 $\overline{\{x\}}$ 上的既约诱导闭子概形结构。
于是由《概形》引理 \ref{schemes-lemma-closed-subspace-scheme}，
$U_0 = U \cap X_0$ 是 $X_0$ 的仿射开集，并且
$U_0 = \{x\}$。令 $y \in X_0$、$y \not = x$ 为 $x$ 的一个特化。
考虑局部环 $R = \mathcal{O}_{X_0, y}$。
由于引理 \ref{properties-lemma-locally-closed-in-Noetherian} 表明 $X_0$ 是 Noether 的，
这是一个 Noether 局部环。记 $V \subset \Spec(R)$ 为 $U_0$ 在
$\Spec(R)$ 中通过典范态射 $\Spec(R) \to X_0$ 得到的逆像（见《概形》第
\ref{schemes-section-points} 节）。
按构造，$V$ 是单点集，其唯一点对应于 $x$（使用《概形》引理
\ref{schemes-lemma-specialize-points}）。
由《代数》引理
\ref{algebra-lemma-Noetherian-local-domain-dim-2-infinite-opens}，可知
$\dim(R) = 1$。
换言之，$y$ 是 $x$ 的直接特化（见《拓扑》定义
\ref{topology-definition-dimension-function}）。
再换言之，任意满足 $y \not = x$ 且 $x \leadsto y$ 的点，
都是 $x$ 的直接特化。显然这些点中的每一个都是闭点，如所需。
\end{proof}

\begin{lemma}
\label{properties-lemma-locally-Noetherian-specialization-dvr}
令 $X$ 为局部 Noether 概形。
令 $x' \leadsto x$ 为 $X$ 中两点之间的特化。则
\begin{enumerate}
\item 存在离散赋值环 $R$ 与态射
$f : \Spec(R) \to X$，使得泛点 $\eta$（它属于 $\Spec(R)$）
映到 $x'$，而特殊点映到 $x$；并且
\item 在 $x \not = x'$ 的前提下，给定有限生成域扩张
$K/\kappa(x')$，可以作出安排，使得扩张
$\kappa(\eta)/\kappa(x')$（它由 $f$ 诱导）与给定扩张同构。
\end{enumerate}
\end{lemma}

\begin{proof}
首先假设 $x' \leadsto x$ 是 $X$ 中的一个特化，且
$x \not = x'$；令 $K/\kappa(x')$ 为有限生成域扩张。由《概形》引理
\ref{schemes-lemma-specialize-points} 以及《概形》引理
\ref{schemes-lemma-characterize-points} 后面的讨论，得到环同态
$\mathcal{O}_{X, x} \to \kappa(x') \to K$。
由于 $x \not = x'$，$\mathcal{O}_{X, x}$ 在 $K$ 中的像不是域（略去细节）。
令 $R \subset K$ 为任意分式域为 $K$ 且支配
$\mathcal{O}_{X, x} \to K$ 的像的离散赋值环，见《代数》引理
\ref{algebra-lemma-exists-dvr}。
环同态 $\mathcal{O}_{X, x} \to R$ 诱导态射
$f : \Spec(R) \to X$，见《概形》引理
\ref{schemes-lemma-morphism-from-spec-local-ring}。
按构造，该态射具有所需的全部性质。
若 $x = x'$，则可以取 $R = \kappa(x)[t]_{(t)}$。
\end{proof}

\begin{lemma}
\label{properties-lemma-thin-infinite-sequence}
令 $S$ 为 Noether 概形。令 $T \subset S$ 为无限子集。
则存在无限子集 $T' \subset T$，使得 $T'$ 的各点之间不存在非平凡特化。
\end{lemma}

\begin{proof}
令 $T_0 \subset T$ 为由 $t \in T$ 组成的集合，其中这些点不特化到 $T$ 中另一点。
若 $T_0$ 无限，则取 $T' = T_0$ 即可。因此可以且确实假设 $T_0$ 有限。
归纳地，对 $i > 0$，考虑集合 $T_i \subset T$，它由满足下列条件的
$t \in T$ 组成：
\begin{enumerate}
\item $t \not \in T_{i - 1} \cup T_{i - 2} \cup \ldots \cup T_0$；
\item 存在非平凡特化 $t \leadsto t'$，且 $t' \in T_{i - 1}$；并且
\item 对任意非平凡特化 $t \leadsto t'$，若 $t' \in T$，则有
$t' \in T_{i - 1} \cup T_{i - 2} \cup \ldots \cup T_0$。
\end{enumerate}
同样，若 $T_i$ 无限，则取 $T' = T_i$ 即可。
令 $d$ 为各局部环 $\mathcal{O}_{S, t}$（其中 $t \in T_0$）维数的最大值；
那么 $d$ 是整数，因为 $T_0$ 有限，并且由《代数》命题
\ref{algebra-proposition-dimension} 可知这些局部环的维数有限。
于是，$T_i = \emptyset$ 对 $i > d$ 成立。
事实上，若 $t \in T_i$，则可以找到非平凡特化序列
$t = t_i \leadsto t_{i - 1} \leadsto \ldots \leadsto t_0$，
其中 $t_0 \in T_0$。由于各点
$t = t_i, t_{i - 1}, \ldots, t_0$ 均属于
$\Spec(\mathcal{O}_{S, t_0})$
（《概形》引理 \ref{schemes-lemma-specialize-points}），可知 $i \leq d$。
因此 $\bigcup T_i = T_d \cup \ldots \cup T_0$ 是 $T$ 的有限子集。

\medskip\noindent
假设 $t \in T$ 不属于 $\bigcup T_i$。则必存在特化
$t \leadsto t'$，其中 $t' \in T$ 且 $t' \not \in \bigcup T_i$。
（事实上，若 $t$ 的每个特化都属于有限集 $T_d \cup \ldots \cup T_0$，
则存在最大的 $i$，使得有某个特化 $t \leadsto t'$ 满足 $t' \in T_i$；
此时按构造 $t \in T_{i + 1}$。）
于是得到无限序列
$$
t \leadsto t' \leadsto t'' \leadsto \ldots
$$
，它由 $T \setminus \bigcup T_i$ 中各点之间的非平凡特化组成。
这是不可能的，因为由引理 \ref{properties-lemma-locally-Noetherian-quasi-separated}，
$S$ 的底拓扑空间是 Noether 的。
\end{proof}

\begin{lemma}
\label{properties-lemma-maximal-points}
令 $S$ 为 Noether 概形。令 $T \subset S$ 为子集。令
$T_0 \subset T$ 为由如下 $t \in T$ 组成的集合：不存在满足
非平凡特化 $t' \leadsto t$，其中 $t' \in T'$。则：(a) $T_0$ 的各点之间
不存在特化；(b) $T$ 的每一点都是 $T_0$ 中某一点的特化；(c) $T$ 与
$T_0$ 的闭包相同。
\end{lemma}

\begin{proof}
回顾：$\dim(\mathcal{O}_{S, s}) < \infty$ 对任意 $s \in S$ 成立，
见《代数》命题 \ref{algebra-proposition-dimension}。令 $t \in T$。
若 $t' \leadsto t$，则由维数理论，
$\dim(\mathcal{O}_{S, t'}) \leq \dim(\mathcal{O}_{S, t})$，
且等号成立当且仅当 $t' = t$。因此，若选取 $t' \leadsto t$，使得
$\dim(\mathcal{O}_{T, t'})$ 最小，则 $t' \in T_0$。
换言之，每个 $t \in T$ 都是 $T_0$ 中某元素的特化。
\end{proof}

\begin{lemma}
\label{properties-lemma-countable-dense-subset}
令 $S$ 为 Noether 概形。令 $T \subset S$ 为无限稠密子集。
则存在可数子集 $E \subset T$，它在 $S$ 中稠密。
\end{lemma}

\begin{proof}
令 $T'$ 为由如下 $s \in S$ 组成的集合：$\overline{\{s\}} \cap T$
含有一个可数子集，其闭包是 $\overline{\{s\}}$。
由于有限集可数，有 $T \subset T'$。
对 $s \in T'$，选取这样的可数子集
$E_s \subset \overline{\{s\}} \cap T$。
令 $E' = \{s_1, s_2, s_3, \ldots\} \subset T'$ 为可数子集。
则 $E'$ 在 $S$ 中的闭包，就是可数子集 $\bigcup_n E_{s_n}$ 的闭包，
而该子集包含于 $T$。
由此可知，若 $Z$ 是 $E'$ 的闭包的一个不可约分支，则 $Z$ 的泛点属于 $T'$。

\medskip\noindent
记 $T'_0 \subset T'$ 为由如下 $t \in T'$ 组成的子集：不存在满足
非平凡特化 $t' \leadsto t$，其中 $t' \in T'$；这正如引理
\ref{properties-lemma-maximal-points}，下文将直接使用该引理的结果。
若 $T'_0$ 无限，则选取可数子集 $E' \subset T'_0$。
由第一段的论证，$E'$ 的闭包的各个不可约分支的泛点都属于 $T'$。
然而，由于其中一点特化到 $E' \subset T'_0$ 中无限多个互不相同的元素，
这产生矛盾。因此 $T'_0$ 有限，设
$T'_0 = \{s_1, \ldots, s_m\}$。于是，$S$ 是 $T$ 的闭包，且包含于
$\{s_1, \ldots, s_m\}$ 的闭包；后者又包含于可数子集
$E_{s_1} \cup \ldots \cup E_{s_m} \subset T$ 的闭包，如所需。
\end{proof}








\section{Jacobson 概形}
\label{properties-section-jacobson}

\noindent
回顾：若闭点在每个闭子集中都稠密，则称一个空间是{it Jacobson 的}，
见《拓扑》第 \ref{topology-section-space-jacobson} 节。

\begin{definition}
\label{properties-definition-jacobson}
称概形 $S$ 是{it Jacobson 概形}，如果它的底拓扑空间是 Jacobson 的。
\end{definition}

\noindent
回顾：若环 $R$ 的每个根理想都是极大理想之交，则称 $R$ 是 Jacobson 环，
见《代数》定义 \ref{algebra-definition-ring-jacobson}。

\begin{lemma}
\label{properties-lemma-affine-jacobson}
仿射概形 $\Spec(R)$ 是 Jacobson 概形，当且仅当环 $R$ 是 Jacobson 环。
\end{lemma}

\begin{proof}
这就是《代数》引理 \ref{algebra-lemma-jacobson}。
\end{proof}

\noindent
下面是刻画 Jacobson 概形的标准结果。
直观地说，它断言 Jacobson $\Leftrightarrow$ 局部 Jacobson。

\begin{lemma}
\label{properties-lemma-locally-jacobson}
令 $X$ 为概形。下列条件等价：
\begin{enumerate}
\item 概形 $X$ 是 Jacobson 概形。
\item 概形 $X$ 在定义 \ref{properties-definition-locally-P} 的意义下“局部 Jacobson”。
\item 对每个仿射开集 $U \subset X$，环 $\mathcal{O}_X(U)$ 是 Jacobson 环。
\item 存在仿射开覆盖 $X = \bigcup U_i$，使得每个 $\mathcal{O}_X(U_i)$
都是 Jacobson 环。
\item 存在开覆盖 $X = \bigcup X_j$，使得每个开子概形 $X_j$
都是 Jacobson 概形。
\end{enumerate}
此外，若 $X$ 是 Jacobson 概形，则每个开子概形也都是 Jacobson 概形。
\end{lemma}

\begin{proof}
引理的最后一个断言由《拓扑》引理
\ref{topology-lemma-jacobson-inherited} 得出。
(5) 与 (1) 的等价性就是《拓扑》引理
\ref{topology-lemma-jacobson-local}。
因此，使用引理 \ref{properties-lemma-affine-jacobson}，可知
(1) $\Leftrightarrow$ (2)。
要完成引理的证明，只需证明“Jacobson”是环的一种局部性质，见引理
\ref{properties-lemma-locally-P}。
Jacobson 环关于一个元素的任意局部化仍为 Jacobson 环，见《代数》引理
\ref{algebra-lemma-Jacobson-invert-element}。
设 $R$ 为环，$f_1, \ldots, f_n \in R$ 生成单位理想，并且每个
$R_{f_i}$ 都是 Jacobson 环。于是
$\Spec(R) = \bigcup D(f_i)$ 是若干开子集之并，这些开子集全都是 Jacobson 的；
因此再次由《拓扑》引理 \ref{topology-lemma-jacobson-local}，
$\Spec(R)$ 是 Jacobson 概形。
这证明了定义 \ref{properties-definition-property-local} 的第二个性质。
\end{proof}

\noindent
代数几何中常用的许多概形都是 Jacobson 概形，见《态射》引理
\ref{morphisms-lemma-ubiquity-Jacobson-schemes}。
这里指出以下有趣的情形。

\begin{lemma}
\label{properties-lemma-complement-closed-point-Jacobson}
以下是 Noether Jacobson 概形的例子。
\begin{enumerate}
\item 若 $(R, \mathfrak m)$ 是 Noether 局部环，则穿孔谱
$\Spec(R) \setminus \{\mathfrak m\}$ 是 Jacobson 概形。
\item 若 $R$ 是 Jacobson 根为 $\text{rad}(R)$ 的 Noether 环，
则 $\Spec(R) \setminus V(\text{rad}(R))$ 是 Jacobson 概形。
\item 若 $(R, I)$ 是 Zariski 对（《更多代数》定义
\ref{more-algebra-definition-zariski-pair}），且 $R$ 是 Noether 环，
则 $\Spec(R) \setminus V(I)$ 是 Jacobson 概形。
\end{enumerate}
\end{lemma}

\begin{proof}
证明 (3)。注意，$\Spec(R) - V(I)$ 由仿射开集
$\Spec(R_f)$（其中 $f \in I$）覆盖。由《更多代数》引理
\ref{more-algebra-lemma-noetherian-zariski-jacobson-complement}，
环 $R_f$ 是 Jacobson 环。
因此由引理 \ref{properties-lemma-locally-jacobson}，
$\Spec(R) \setminus V(I)$ 是 Jacobson 概形。
(1) 与 (2) 是 (3) 的特殊情形。

\medskip\noindent
下面直接证明情形 (1)。
由于 $\Spec(R)$ 是 Noether 概形，
$S$ 是 Noether 概形（引理 \ref{properties-lemma-locally-closed-in-Noetherian}）。
因此 $S$ 是清醒的 Noether 拓扑空间（使用《概形》引理
\ref{schemes-lemma-scheme-sober}）。
假设 $S$ 不是 Jacobson 概形，以导出矛盾。
由《拓扑》引理
\ref{topology-lemma-non-jacobson-Noetherian-characterize}，
存在某个非闭点 $\xi \in S$，使得 $\{\xi\}$ 局部闭。
这对应于素理想 $\mathfrak p \subset R$，并且：(1) 存在素理想
$\mathfrak q$，满足 $\mathfrak p \subset \mathfrak q \subset \mathfrak m$，
两个包含关系都是严格的；(2) $\{\mathfrak p\}$ 在
$\Spec(R/\mathfrak p)$ 中是开的。这与《代数》引理
\ref{algebra-lemma-Noetherian-local-domain-dim-2-infinite-opens} 矛盾。
\end{proof}








\section{正规概形}
\label{properties-section-normal}

\noindent
回顾：若环 $R$ 的所有局部环都是正规整环，则称该环正规，
见《代数》定义 \ref{algebra-definition-ring-normal}。
正规整环是指在其分式域中整闭的整环，见《代数》定义
\ref{algebra-definition-domain-normal}。
因此可以如下定义正规概形。

\begin{definition}
\label{properties-definition-normal}
概形 $X$ 正规，当且仅当对所有 $x \in X$，局部环
$\mathcal{O}_{X, x}$ 都是正规整环。
\end{definition}

\noindent
这似乎是 EGA 所采用的定义，见 \cite[0, 4.1.4]{EGA}。
假设 $X = \Spec(A)$，且 $A$ 约化。此时，说 $X$ 正规并不等价于说
$A$ 在其总分式环中整闭。然而，若 $A$ 是 Noether 环，则二者等价
（见《代数》引理 \ref{algebra-lemma-characterize-reduced-ring-normal}）。

\begin{lemma}
\label{properties-lemma-locally-normal}
令 $X$ 为概形。下列条件等价：
\begin{enumerate}
\item 概形 $X$ 正规。
\item 对每个仿射开集 $U \subset X$，环 $\mathcal{O}_X(U)$ 正规。
\item 存在仿射开覆盖 $X = \bigcup U_i$，使得每个 $\mathcal{O}_X(U_i)$
都正规。
\item 存在开覆盖 $X = \bigcup X_j$，使得每个开子概形 $X_j$ 都正规。
\end{enumerate}
此外，若 $X$ 正规，则每个开子概形也正规。
\end{lemma}

\begin{proof}
这由定义立即可见。
\end{proof}

\begin{lemma}
\label{properties-lemma-normal-reduced}
正规概形是约化的。
\end{lemma}

\begin{proof}
这立即来自定义。
\end{proof}

\begin{lemma}
\label{properties-lemma-integral-normal}
令 $X$ 为整概形。
则 $X$ 正规，当且仅当对每个非空仿射开集
$U \subset X$，环 $\mathcal{O}_X(U)$ 都是正规整环。
\end{lemma}

\begin{proof}
这由《代数》引理 \ref{algebra-lemma-normality-is-local} 得出。
\end{proof}

\begin{lemma}
\label{properties-lemma-normal-locally-finite-nr-irreducibles}
令 $X$ 为概形，并假设任意拟紧开集都只有有限多个不可约分支。
下列条件等价：
\begin{enumerate}
\item $X$ 正规；
\item $X$ 是若干正规整概形的不交并。
\end{enumerate}
\end{lemma}

\begin{proof}
由定义立即可见，(2) 蕴含 (1)。
令 $X$ 为正规概形，并假设每个拟紧开集都只有有限多个不可约分支。
若 $X$ 仿射，则由《代数》引理
\ref{algebra-lemma-characterize-reduced-ring-normal}，$X$ 满足 (2)。
对于一般的 $X$，令 $X = \bigcup X_i$ 为仿射开覆盖。
注意，每个 $X_i$ 也只有有限多个不可约分支，并且该引理对每个 $X_i$ 成立。
令 $T \subset X$ 为一个不可约分支。由仿射情形，每个交
$T \cap X_i$ 在 $X_i$ 中是开的，并且是正规整概形。
因此 $T \subset X$ 是开子概形，并且是正规整概形。
这证明 $X$ 是其不可约分支的不交并，而这些分支都是正规整概形。
\end{proof}

\begin{lemma}
\label{properties-lemma-normal-Noetherian}
令 $X$ 为 Noether 概形。
下列条件等价：
\begin{enumerate}
\item $X$ 正规；
\item $X$ 是有限多个正规整概形的不交并。
\end{enumerate}
\end{lemma}

\begin{proof}
这是引理 \ref{properties-lemma-normal-locally-finite-nr-irreducibles} 的特殊情形，
因为 Noether 概形具有 Noether 底拓扑空间（引理
\ref{properties-lemma-Noetherian-topology} 以及《拓扑》引理
\ref{topology-lemma-Noetherian}）。
\end{proof}

\begin{lemma}
\label{properties-lemma-normal-locally-Noetherian}
令 $X$ 为局部 Noether 概形。
下列条件等价：
\begin{enumerate}
\item $X$ 正规；
\item $X$ 是若干正规整概形的不交并。
\end{enumerate}
\end{lemma}

\begin{proof}
略。提示：由引理 \ref{properties-lemma-normal-Noetherian}，这纯粹是一个拓扑事实。
\end{proof}

\begin{remark}
\label{properties-remark-normal-connected-irreducible}
令 $X$ 为正规概形。若 $X$ 局部 Noether，则可知
$X$ 是整概形，当且仅当 $X$ 连通，见引理
\ref{properties-lemma-normal-locally-Noetherian}。
但是，存在连通仿射概形 $X$，使得 $\mathcal{O}_{X, x}$ 对所有
$x \in X$ 都是整环，而 $X$ 不可约，见《例》第
\ref{examples-section-connected-locally-integral-not-integral} 节。
这个例子甚至还是正规概形（略去证明），务请注意！
\end{remark}

\begin{lemma}
\label{properties-lemma-normal-integral-sections}
\begin{slogan}
正规概形上的函数环是正规环。
\end{slogan}
令 $X$ 为正规整概形。
则 $\Gamma(X, \mathcal{O}_X)$ 是正规整环。
\end{lemma}

\begin{proof}
置 $R = \Gamma(X, \mathcal{O}_X)$。
显然 $R$ 是整环。
假设 $f = a/b$ 是其分式域中的元素，并且在 $R$ 上整。设有
$f^d + \sum_{i = 0, \ldots, d - 1} a_i f^i = 0$，其中
$a_i \in R$。令 $U \subset X$ 为非空仿射开集。
由于 $b \in R$ 非零，并且 $X$ 是整概形，可知
$b|_U \in \mathcal{O}_X(U)$ 也非零。
因此 $a/b$ 是 $\mathcal{O}_X(U)$ 的分式域中的元素，并且在
$\mathcal{O}_X(U)$ 上整（因为可以使用同一个多项式
$f^d + \sum_{i = 0, \ldots, d - 1} a_i|_U f^i = 0$，在 $U$ 上亦然）。
由于 $\mathcal{O}_X(U)$ 是正规整环
（引理 \ref{properties-lemma-integral-normal}），可知
$f_U = (a|_U)/(b|_U) \in \mathcal{O}_X(U)$。显然，
$f_U|_V = f_V$ 只要 $V \subset U \subset X$ 是非空仿射开集。
因此各局部截面 $f_U$ 粘合成一个元素
$g \in R = \Gamma(X, \mathcal{O}_X)$。于是 $bg$ 与 $a$
限制为 $\mathcal{O}_X(U)$ 的同一元素，对所有上述 $U$ 都如此，故
$bg = a$；换言之，$g$ 映到 $f$，而这是在 $R$ 的分式域中。
\end{proof}


\section{Cohen-Macaulay 概形}
\label{properties-section-Cohen-Macaulay}

\noindent
回顾《代数》定义 \ref{algebra-definition-local-ring-CM}：若 Noether 局部环
$(R, \mathfrak m)$ 满足
$\text{depth}_{\mathfrak m}(R) = \dim(R)$，则称其为 Cohen-Macaulay 环。
又回顾：若 Noether 环 $R$ 的每个局部环 $R_{\mathfrak p}$ 都是
Cohen-Macaulay 环，则称 $R$ 为 Cohen-Macaulay 环，见《代数》定义
\ref{algebra-definition-ring-CM}。

\begin{definition}
\label{properties-definition-Cohen-Macaulay}
令 $X$ 为概形。称 $X$ 是{\it Cohen-Macaulay 概形}，如果对每个
$x \in X$，都存在仿射开邻域 $U \subset X$ 包含 $x$，使得环
$\mathcal{O}_X(U)$ 是 Noether 环且为 Cohen-Macaulay 环。
\end{definition}

\begin{lemma}
\label{properties-lemma-characterize-Cohen-Macaulay}
令 $X$ 为概形。下列条件等价：
\begin{enumerate}
\item $X$ 是 Cohen-Macaulay 概形；
\item $X$ 局部 Noether，且它的所有局部环都是 Cohen-Macaulay 环；
以及
\item $X$ 局部 Noether，且对任意闭点 $x \in X$，局部环
$\mathcal{O}_{X, x}$ 都是 Cohen-Macaulay 环。
\end{enumerate}
\end{lemma}

\begin{proof}
《代数》引理 \ref{algebra-lemma-localize-CM} 说明，Cohen-Macaulay
局部环的局部化仍是 Cohen-Macaulay 环。将此与引理
\ref{properties-lemma-locally-Noetherian}、局部 Noether 概形上闭点的存在性
（引理 \ref{properties-lemma-locally-Noetherian-closed-point}）以及各定义结合，
即得本引理。
\end{proof}

\begin{lemma}
\label{properties-lemma-locally-Cohen-Macaulay}
令 $X$ 为概形。下列条件等价：
\begin{enumerate}
\item 概形 $X$ 是 Cohen-Macaulay 概形。
\item 对每个仿射开集 $U \subset X$，环 $\mathcal{O}_X(U)$
都是 Noether 环且为 Cohen-Macaulay 环。
\item 存在仿射开覆盖 $X = \bigcup U_i$，使得每个
$\mathcal{O}_X(U_i)$ 都是 Noether 环且为 Cohen-Macaulay 环。
\item 存在开覆盖 $X = \bigcup X_j$，使得每个开子概形 $X_j$
都是 Cohen-Macaulay 概形。
\end{enumerate}
此外，若 $X$ 是 Cohen-Macaulay 概形，则每个开子概形也都是
Cohen-Macaulay 概形。
\end{lemma}

\begin{proof}
结合引理 \ref{properties-lemma-locally-Noetherian}
和 \ref{properties-lemma-characterize-Cohen-Macaulay}。
\end{proof}

\noindent
关于 Cohen-Macaulay 概形和深度的更多信息，见《概形上同调》第
\ref{coherent-section-depth} 节。








\section{正则概形}
\label{properties-section-regular}

\noindent
回顾《代数》定义 \ref{algebra-definition-regular-local}：若 Noether
局部环 $(R, \mathfrak m)$ 的理想 $\mathfrak m$ 可由
$\dim(R)$ 个元素生成，则称该局部环{\it 正则}。
又回顾：若 Noether 环 $R$ 的每个局部环 $R_{\mathfrak p}$ 都正则，
则称 $R$ 为{\it 正则环}，见《代数》定义
\ref{algebra-definition-regular}。

\begin{definition}
\label{properties-definition-regular}
令 $X$ 为概形。称 $X$ 是{\it 正则概形}，或称{\it 非奇异概形}，
如果对每个 $x \in X$，都存在仿射开邻域 $U \subset X$ 包含 $x$，
使得环 $\mathcal{O}_X(U)$ 是 Noether 正则环。
\end{definition}

\begin{lemma}
\label{properties-lemma-characterize-regular}
令 $X$ 为概形。下列条件等价：
\begin{enumerate}
\item $X$ 正则；
\item $X$ 局部 Noether，且它的所有局部环都正则；
以及
\item $X$ 局部 Noether，且对任意闭点 $x \in X$，局部环
$\mathcal{O}_{X, x}$ 都正则。
\end{enumerate}
\end{lemma}

\begin{proof}
由《代数》定义 \ref{algebra-definition-regular} 之前的讨论可知，
正则局部环的局部化仍正则。将此与引理
\ref{properties-lemma-locally-Noetherian}、局部 Noether 概形上闭点的存在性
（引理 \ref{properties-lemma-locally-Noetherian-closed-point}）以及各定义结合，
即得本引理。
\end{proof}

\begin{lemma}
\label{properties-lemma-locally-regular}
令 $X$ 为概形。下列条件等价：
\begin{enumerate}
\item 概形 $X$ 正则。
\item 对每个仿射开集 $U \subset X$，环 $\mathcal{O}_X(U)$
都是 Noether 正则环。
\item 存在仿射开覆盖 $X = \bigcup U_i$，使得每个
$\mathcal{O}_X(U_i)$ 都是 Noether 正则环。
\item 存在开覆盖 $X = \bigcup X_j$，使得每个开子概形 $X_j$
都正则。
\end{enumerate}
此外，若 $X$ 正则，则每个开子概形也正则。
\end{lemma}

\begin{proof}
结合引理 \ref{properties-lemma-locally-Noetherian}
和 \ref{properties-lemma-characterize-regular}。
\end{proof}

\begin{lemma}
\label{properties-lemma-regular-normal}
正则概形是正规的。
\end{lemma}

\begin{proof}
见《代数》引理 \ref{algebra-lemma-regular-normal}。
\end{proof}





\section{维数}
\label{properties-section-dimension}

\noindent
概形的维数就是其底拓扑空间的维数。

\begin{definition}
\label{properties-definition-dimension}
令 $X$ 为概形。
\begin{enumerate}
\item $X$ 的{\it 维数}就是 $X$ 作为拓扑空间的维数，见《拓扑》定义
\ref{topology-definition-Krull}。
\item 对 $x \in X$，按照《拓扑》定义
\ref{topology-definition-Krull}，以 $\dim_x(X)$ 表示 $X$ 的底拓扑空间
在 $x$ 处的维数。称 $\dim_x(X)$ 为{\it $X$ 在 $x$ 处的维数}。
\end{enumerate}
\end{definition}

\noindent
由于概形的底拓扑空间是清醒空间
（《概形》引理 \ref{schemes-lemma-scheme-sober}），我们可以把 $X$
的维数算作下列不可约闭子集链的长度 $n$ 的上确界：
$$
T_0 \subset T_1 \subset \ldots \subset T_n
$$
其中各集合都是 $X$ 的不可约闭子集；也可以把它算作下列点的特化链
的长度 $n$ 的上确界：
$$
\xi_n \leadsto \xi_{n - 1} \leadsto \ldots \leadsto \xi_0
$$
其中各点属于 $X$。

\begin{lemma}
\label{properties-lemma-dimension}
令 $X$ 为概形。下列三个量相等：
\begin{enumerate}
\item $X$ 的维数；
\item $X$ 的各局部环维数的上确界；
\item $\dim_x(X)$ 在 $x \in X$ 时的上确界。
\end{enumerate}
\end{lemma}

\begin{proof}
注意，给定 $X$ 中点的特化链
$$
\xi_n \leadsto \xi_{n - 1} \leadsto \ldots \leadsto \xi_0
$$
由《概形》引理 \ref{schemes-lemma-specialize-points}，所有点
$\xi_i$ 都对应于 $X$ 在 $\xi_0$ 处的局部环的素理想。
因此，$X$ 的维数就是其各局部环维数的上确界。特别地，
$\dim_x(X) \geq \dim(\mathcal{O}_{X, x})$，因为 $\dim_x(X)$
是 $x$ 的各开邻域维数的最小值。因此
$\sup_{x \in X} \dim_x(X) \geq \dim(X)$。另一方面，显然
$\sup_{x \in X} \dim_x(X) \leq \dim(X)$，因为
$\dim(U) \leq \dim(X)$ 对 $X$ 的任意开子集都成立。
\end{proof}

\begin{lemma}
\label{properties-lemma-codimension-local-ring}
令 $X$ 为概形，$Y \subset X$ 为不可约闭子集，并令 $\xi \in Y$
为其泛点。则
$$
\text{codim}(Y, X) = \dim(\mathcal{O}_{X, \xi})
$$
其中余维按《拓扑》定义 \ref{topology-definition-codimension} 定义。
\end{lemma}

\begin{proof}
由《拓扑》引理 \ref{topology-lemma-codimension-at-generic-point}，
可以把 $X$ 替换为 $\xi$ 的一个仿射开邻域。在这种情形下，结论容易由
《代数》引理 \ref{algebra-lemma-irreducible-components-containing-x} 得出。
\end{proof}

\begin{lemma}
\label{properties-lemma-generic-point}
令 $X$ 为概形，$x \in X$。则 $x$ 是 $X$ 的某个不可约分支的泛点，
当且仅当 $\dim(\mathcal{O}_{X, x}) = 0$。
\end{lemma}

\begin{proof}
例如，这可由引理 \ref{properties-lemma-codimension-local-ring} 得出。
\end{proof}

\begin{lemma}
\label{properties-lemma-locally-Noetherian-dimension-0}
维数为 $0$ 的局部 Noether 概形是不交并；其每个分支都是一个
Artin 局部环的谱。
\end{lemma}

\begin{proof}
维数为 $0$ 的 Noether 环是有限多个 Artin 局部环的直积，见《代数》命题
\ref{algebra-proposition-dimension-zero-ring}。因此，局部 Noether 概形
$X$ 的维数为 $0$ 时，其任意仿射开集都有离散的底拓扑空间。这说明 $X$
上的拓扑是离散的。本引理容易由这些观察得出。
\end{proof}

\begin{lemma}
\label{properties-lemma-dimension-zero}
\begin{reference}
Ofer Gabber 于 2016 年 6 月 4 日发来的电子邮件
\end{reference}
令 $X$ 为零维概形。下列条件等价：
\begin{enumerate}
\item $X$ 拟分离；
\item $X$ 分离；
\item $X$ 是豪斯多夫空间；
\item 每个仿射开集都是闭集。
\end{enumerate}
在这种情形下，$X$ 的连通分支都是点，并且 $X$ 的每个拟紧开集都仿射。
特别地，若 $X$ 拟紧，则 $X$ 仿射。
\end{lemma}

\begin{proof}
由于 $X$ 的维数为零，对任意仿射开集 $U \subset X$，空间 $U$
都是预有限空间，并且满足我们将在下文自由使用的许多其他性质；见《代数》引理
\ref{algebra-lemma-ring-with-only-minimal-primes}。选择仿射开覆盖
$X = \bigcup U_i$。

\medskip\noindent
若 (4) 成立，则 $U_i \cap U_j$ 是 $U_i$ 的闭子集，因而拟紧；所以由
《概形》引理 \ref{schemes-lemma-characterize-quasi-separated}，$X$
拟分离，即 (1) 成立。

\medskip\noindent
若 (1) 成立，则 $U_i \cap U_j$ 是 $U_i$ 的拟紧开集，因而在 $U_i$
中闭。于是 $U_i \cap U_j \to U_i$ 是像为闭集的开浸入，故它是闭浸入。
特别地，$U_i \cap U_j$ 仿射，并且
$\mathcal{O}(U_i) \to \mathcal{O}_X(U_i \cap U_j)$ 满射。因此由《概形》
引理 \ref{schemes-lemma-characterize-separated}，$X$ 分离，即 (2) 成立。

\medskip\noindent
假设 (2) 成立，并令 $x, y \in X$。设 $x \in U_i$。若 $y \in U_i$，
则可以找到 $x$ 与 $y$ 的不交开邻域，因为 $U_i$ 是豪斯多夫空间。
设 $y \not \in U_i$ 且 $y \in U_j$。则 $y \not \in U_i \cap U_j$；
该交是 $U_j$ 的仿射开集，因而在 $U_j$ 中闭。因此可以找到 $y$ 的一个
不与 $U_i$ 相交的开邻域，从而断定 $X$ 是豪斯多夫空间，即 (3) 成立。

\medskip\noindent
假设 (3) 成立。令 $U \subset X$ 为仿射开集。由《拓扑》引理
\ref{topology-lemma-quasi-compact-in-Hausdorff}，$U$ 在 $X$ 中闭。
这证明 (4) 成立。

\medskip\noindent
假设 $X$ 满足等价条件 (1) -- (4)。我们证明本引理的最后几个断言。
设 $x, y \in X$ 且 $x \not = y$。由于 $y$ 不特化到 $x$，可以选择仿射开集
$U \subset X$，使得 $x \in U$ 且 $y \not \in U$。于是
$X = U \amalg (X \setminus U)$ 是开闭子集的一个分解；这说明 $x$ 与 $y$
不属于 $X$ 的同一个连通分支。其次，假设 $U \subset X$ 是拟紧开集。
将 $U = U_1 \cup \ldots \cup U_n$ 写成仿射开集之并。我们对 $n$ 归纳，
证明 $U$ 仿射。这立即将问题归结到 $n = 2$ 的情形。此时
$U = (U_1 \setminus U_2) \amalg (U_1 \cap U_2) \amalg (U_2 \setminus U_1)$，
而上述论证说明每一部分都仿射。
\end{proof}

\begin{lemma}
\label{properties-lemma-isolated-dim-0-locally-noetherian}
令 $x$ 为局部 Noether 概形 $X$ 的一点。则 $\dim_x(X) = 0$，
当且仅当 $x$ 是 $X$ 的孤立点。
\end{lemma}

\begin{proof}
若 $x$ 是孤立点，则 $\{x\}$ 是 $X$ 的开子集
（《拓扑》定义 \ref{topology-definition-isolated-point}），因而由定义
$\dim_x(X) = \dim(\{x\}) = 0$。反之，若 $\dim_x(X) = 0$，则存在
开邻域 $U \subset X$ 包含 $x$，使得 $\dim(U) = 0$
（《拓扑》定义 \ref{topology-definition-Krull}）。由引理
\ref{properties-lemma-locally-Noetherian-dimension-0}，$U$ 上的拓扑是离散的，故
$x$ 是孤立点。
\end{proof}






\section{链状概形}
\label{properties-section-catenary}

\noindent
回顾：若对拓扑空间 $X$ 的每一对不可约闭子集 $T \subset T'$，都存在
不可约闭子集的极大链
$$
T = T_0 \subset T_1 \subset \ldots \subset T_e = T'
$$
并且每条这样的链都有相同长度，则称该空间是{\it 链状的}。见《拓扑》定义
\ref{topology-definition-catenary}。

\begin{definition}
\label{properties-definition-catenary}
令 $S$ 为概形。若 $S$ 的底拓扑空间是链状的，则称 $S$ 是
{\it 链状概形}。
\end{definition}

\noindent
回顾：若对环 $A$ 的任意一对素理想
$\mathfrak p \subset \mathfrak q$，都存在素理想的极大链
$$
\mathfrak p =
\mathfrak p_0 \subset \ldots \subset \mathfrak p_e
= \mathfrak q
$$
并且所有这样的链都有相同长度，则称该环为{\it 链环}。见《代数》定义
\ref{algebra-definition-catenary}。

\begin{lemma}
\label{properties-lemma-catenary-local}
令 $S$ 为概形。下列条件等价：
\begin{enumerate}
\item $S$ 是链状概形；
\item 存在 $S$ 的开覆盖，其每个成员都是链状概形；
\item 对每个仿射开集 $\Spec(R) = U \subset S$，环 $R$ 都是链环；以及
\item 存在仿射开覆盖 $S = \bigcup U_i$，使得每个 $U_i$ 都是一个链环的谱。
\end{enumerate}
此外，在这种情形下，$S$ 的任意局部闭子概形也都是链状概形。
\end{lemma}

\begin{proof}
结合《拓扑》引理 \ref{topology-lemma-catenary} 与《代数》引理
\ref{algebra-lemma-catenary}。
\end{proof}

\begin{lemma}
\label{properties-lemma-catenary-dimension-function}
令 $S$ 为局部 Noether 概形。下列条件等价：
\begin{enumerate}
\item $S$ 是链状概形；以及
\item 在 Zariski 拓扑下局部地，$S$ 上存在维数函数
（见《拓扑》定义 \ref{topology-definition-dimension-function}）。
\end{enumerate}
\end{lemma}

\begin{proof}
这由《拓扑》引理
\ref{topology-lemma-catenary}、
\ref{topology-lemma-dimension-function-catenary} 和
\ref{topology-lemma-locally-dimension-function}，
《概形》引理 \ref{schemes-lemma-scheme-sober}，以及最后的引理
\ref{properties-lemma-Noetherian-topology} 得出。
\end{proof}

\noindent
可以证明，一个概形是链状概形，当且仅当它的各局部环都是链环。

\begin{lemma}
\label{properties-lemma-catenary-local-rings-catenary}
令 $X$ 为概形。下列条件等价：
\begin{enumerate}
\item $X$ 是链状概形；以及
\item 对任意 $x \in X$，局部环 $\mathcal{O}_{X, x}$ 都是链环。
\end{enumerate}
\end{lemma}

\begin{proof}
假设 $X$ 是链状概形。令 $x \in X$。由引理
\ref{properties-lemma-catenary-local}，可以把 $X$ 替换为 $x$ 的一个仿射开邻域；
此时 $\Gamma(X, \mathcal{O}_X)$ 是链环。由《代数》引理
\ref{algebra-lemma-localization-catenary}，链环的任意局部化仍为链环。
所以 $\mathcal{O}_{X, x}$ 是链环。

\medskip\noindent
反过来，假设 $X$ 的所有局部环都是链环。令 $Y \subset Y'$ 为 $X$
的不可约闭子集之间的包含关系，并令 $\xi \in Y$ 为泛点。令
$\mathfrak p \subset \mathcal{O}_{X, \xi}$ 为对应于 $Y'$ 的泛点的素理想；
见《概形》引理 \ref{schemes-lemma-specialize-points}。由同一引理，
$X$ 中介于 $Y$ 与 $Y'$ 之间的不可约闭子集，对应于素理想
$\mathfrak q \subset \mathcal{O}_{X, \xi}$，这些素理想满足
$\mathfrak p \subset \mathfrak q \subset \mathfrak m_{\xi}$。因为
$\mathcal{O}_{X, \xi}$ 是链环，可见这些素理想的所有极大链都是有限的，
并且具有相同长度。
\end{proof}






\section{Serre 条件}
\label{properties-section-Rk}

\noindent
下面给出两个经常有用的技术概念。另见《概形上同调》第
\ref{coherent-section-depth} 节。

\begin{definition}
\label{properties-definition-Rk}
令 $X$ 为局部 Noether 概形，并令 $k \geq 0$。
\begin{enumerate}
\item 称 $X$ {\it 在余维 $k$ 内正则}，或称 $X$ 具有性质
{\it $(R_k)$}，如果对每个 $x \in X$ 都有
$$
\dim(\mathcal{O}_{X, x}) \leq k
\Rightarrow
\mathcal{O}_{X, x}\text{ 正则}
$$
\item 称 $X$ 具有性质 {\it $(S_k)$}，如果对每个 $x \in X$ 都有
如下不等式：\par\noindent
$\text{depth}(\mathcal{O}_{X, x}) \geq \min(k, \dim(\mathcal{O}_{X, x}))$。
\end{enumerate}
\end{definition}

\noindent
“在余维 $k$ 内正则”这一说法是有意义的，因为我们已在第
\ref{properties-section-catenary} 节看到：若 $Y \subset X$ 是泛点为 $x$ 的不可约
闭子集，则 $\dim(\mathcal{O}_{X, x}) = \text{codim}(Y, X)$。例如，条件
$(R_0)$ 意味着：对每个泛点 $\eta \in X$，其中该点是 $X$ 的一个不可约分支的泛点，
局部环 $\mathcal{O}_{X, \eta}$ 都是域。但是，对一般 Noether 概形，$X$
的正则轨迹可能表现得很糟糕，所以必须谨慎。

\begin{lemma}
\label{properties-lemma-scheme-regular-iff-all-Rk}
令 $X$ 为局部 Noether 概形。则 $X$ 正则，当且仅当 $X$ 具有性质
$(R_k)$，对所有 $k \geq 0$ 都如此。
\end{lemma}

\begin{proof}
这由引理 \ref{properties-lemma-characterize-regular} 及各定义得出。
\end{proof}

\begin{lemma}
\label{properties-lemma-scheme-CM-iff-all-Sk}
令 $X$ 为局部 Noether 概形。则 $X$ 是 Cohen-Macaulay 概形，当且仅当
$X$ 具有性质 $(S_k)$，对所有 $k \geq 0$ 都如此。
\end{lemma}

\begin{proof}
由引理 \ref{properties-lemma-characterize-Cohen-Macaulay}，只须考察局部环。
因此，本引理成立，因为 Noether 局部环是 Cohen-Macaulay 环，当且仅当
其深度等于其维数。
\end{proof}

\begin{lemma}
\label{properties-lemma-criterion-reduced}
令 $X$ 为局部 Noether 概形。则 $X$ 约化，当且仅当 $X$ 具有性质
$(S_1)$ 和 $(R_0)$。
\end{lemma}

\begin{proof}
这就是《代数》引理 \ref{algebra-lemma-criterion-reduced}。
\end{proof}

\begin{lemma}
\label{properties-lemma-criterion-normal}
令 $X$ 为局部 Noether 概形。则 $X$ 正规，当且仅当 $X$ 具有性质
$(S_2)$ 和 $(R_1)$。
\end{lemma}

\begin{proof}
这就是《代数》引理 \ref{algebra-lemma-criterion-normal}。
\end{proof}

\begin{lemma}
\label{properties-lemma-normal-dimension-1-regular}
令 $X$ 为正规局部 Noether 概形，且维数 $\leq 1$。则 $X$ 正则。
\end{lemma}

\begin{proof}
这由引理 \ref{properties-lemma-criterion-normal} 及各定义得出。
\end{proof}

\begin{lemma}
\label{properties-lemma-normal-dimension-2-Cohen-Macaulay}
令 $X$ 为正规局部 Noether 概形，且维数 $\leq 2$。则 $X$ 是
Cohen-Macaulay 概形。
\end{lemma}

\begin{proof}
这由引理 \ref{properties-lemma-criterion-normal} 及各定义得出。
\end{proof}








\section{Japanese 概形与 Nagata 概形}
\label{properties-section-nagata}

\noindent
本节所考虑的概念在 EGA 中并没有得到醒目的定义。EGA 的
\cite[IV Corollary 5.11.4]{EGA} 提到并定义了“泛 Japanese 概形”；其
\cite[IV Remark 10.4.14 (ii)]{EGA} 提到“Japanese 概形”，但没有给出定义。
文献中有少数地方使用下面定义的 Nagata 概形
（例如见 \cite[Definition 8.2.30]{Liu} 和 \cite[Page 142]{Greco}）。

\medskip\noindent
简要回顾：若整环 $R$ 在其分式域的任意有限扩张中的整闭包都是 $R$ 上的
有限模，则称 $R$ 是{\it Japanese 整环}。若环 $R$ 对任意有限型环同态
$R \to S$，只要 $S$ 是整环，$S$ 就是 Japanese 整环，则称该环是
{\it 泛 Japanese 环}。若环 $R$ 是 Noether 环，并且
$R/\mathfrak p$ 对每个素理想 $\mathfrak p$ 都是 Japanese 整环，则称 $R$ 为
{\it Nagata 环}。

\begin{definition}
\label{properties-definition-nagata}
令 $X$ 为概形。
\begin{enumerate}
\item 假设 $X$ 是整概形。称 $X$ 是{\it Japanese 概形}，如果对每个
$x \in X$，都存在仿射开邻域 $x \in U \subset X$，使得环
$\mathcal{O}_X(U)$ 是 Japanese 整环
（见《代数》定义 \ref{algebra-definition-N}）。
\item 称 $X$ 是{\it 泛 Japanese 概形}，如果对每个 $x \in X$，都存在
仿射开邻域 $x \in U \subset X$，使得环 $\mathcal{O}_X(U)$ 是泛
Japanese 环（见《代数》定义 \ref{algebra-definition-nagata}）。
\item 称 $X$ 是{\it Nagata 概形}，如果对每个 $x \in X$，都存在仿射
开邻域 $x \in U \subset X$，使得环 $\mathcal{O}_X(U)$ 是 Nagata 环
（见《代数》定义 \ref{algebra-definition-nagata}）。
\end{enumerate}
\end{definition}

\noindent
成为 Nagata 概形等价于局部 Noether 且为泛 Japanese 概形；见引理
\ref{properties-lemma-nagata-universally-Japanese}。

\begin{remark}
\label{properties-remark-non-integral-Japanese}
在 \cite{Hoobler-finite} 中，称一个（局部 Noether）概形 $X$ 是
Japanese 概形，如果对每个 $x \in X$ 以及每个伴随素理想
$\mathfrak p$（它是 $\mathcal{O}_{X, x}$ 的伴随素理想），环
$\mathcal{O}_{X, x}/\mathfrak p$ 都是
Japanese 整环。我们不使用这个定义，因为存在一维 Noether 整环，其局部环
都是优秀环（特别地，都是 Japanese 整环），但它的正规化不是有限的。见
\cite[Example 1]{Hochster-loci}、\cite{Heinzer-Levy} 或
\cite[Expos\'e XIX]{Traveaux}。另一方面，可以用如下方式避开这个问题：
称概形 $X$ 是 Japanese 概形，如果对每个仿射开集
$\Spec(A) \subset X$，环 $A/\mathfrak p$ 对每个伴随素理想
$\mathfrak p$（属于 $A$）都是 Japanese 整环。
\end{remark}

\begin{lemma}
\label{properties-lemma-nagata-locally-Noetherian}
Nagata 概形是局部 Noether 的。
\end{lemma}

\begin{proof}
这是因为 Nagata 环按定义是 Noether 环。
\end{proof}

\begin{lemma}
\label{properties-lemma-locally-Japanese}
令 $X$ 为整概形。下列条件等价：
\begin{enumerate}
\item 概形 $X$ 是 Japanese 概形。
\item 对每个仿射开集 $U \subset X$，整环 $\mathcal{O}_X(U)$ 都是
Japanese 整环。
\item 存在仿射开覆盖 $X = \bigcup U_i$，使得每个
$\mathcal{O}_X(U_i)$ 都是 Japanese 整环。
\item 存在开覆盖 $X = \bigcup X_j$，使得每个开子概形 $X_j$ 都是
Japanese 概形。
\end{enumerate}
此外，若 $X$ 是 Japanese 概形，则每个开子概形也是 Japanese 概形。
\end{lemma}

\begin{proof}
这由引理 \ref{properties-lemma-locally-P} 以及《代数》引理
\ref{algebra-lemma-localize-N} 和
\ref{algebra-lemma-Japanese-local} 得出。
\end{proof}

\begin{lemma}
\label{properties-lemma-locally-universally-Japanese}
令 $X$ 为概形。下列条件等价：
\begin{enumerate}
\item 概形 $X$ 是泛 Japanese 概形。
\item 对每个仿射开集 $U \subset X$，环 $\mathcal{O}_X(U)$ 都是
泛 Japanese 环。
\item 存在仿射开覆盖 $X = \bigcup U_i$，使得每个
$\mathcal{O}_X(U_i)$ 都是泛 Japanese 环。
\item 存在开覆盖 $X = \bigcup X_j$，使得每个开子概形 $X_j$ 都是
泛 Japanese 概形。
\end{enumerate}
此外，若 $X$ 是泛 Japanese 概形，则每个开子概形也是泛 Japanese 概形。
\end{lemma}

\begin{proof}
这由引理 \ref{properties-lemma-locally-P} 以及《代数》引理
\ref{algebra-lemma-universally-japanese} 和
\ref{algebra-lemma-nagata-local} 得出。
\end{proof}

\begin{lemma}
\label{properties-lemma-locally-nagata}
令 $X$ 为概形。下列条件等价：
\begin{enumerate}
\item 概形 $X$ 是 Nagata 概形。
\item 对每个仿射开集 $U \subset X$，环 $\mathcal{O}_X(U)$ 都是 Nagata 环。
\item 存在仿射开覆盖 $X = \bigcup U_i$，使得每个
$\mathcal{O}_X(U_i)$ 都是 Nagata 环。
\item 存在开覆盖 $X = \bigcup X_j$，使得每个开子概形 $X_j$ 都是
Nagata 概形。
\end{enumerate}
此外，若 $X$ 是 Nagata 概形，则每个开子概形也是 Nagata 概形。
\end{lemma}

\begin{proof}
这由引理 \ref{properties-lemma-locally-P} 以及《代数》引理
\ref{algebra-lemma-nagata-localize} 和
\ref{algebra-lemma-nagata-local} 得出。
\end{proof}

\begin{lemma}
\label{properties-lemma-characterize-nagata}
令 $X$ 为局部 Noether 概形。则 $X$ 是 Nagata 概形，当且仅当每个整的
闭子概形 $Z \subset X$ 都是 Japanese 概形。
\end{lemma}

\begin{proof}
假设 $X$ 是 Nagata 概形。令 $Z \subset X$ 为整的闭子概形。令 $z \in Z$。
取仿射开集 $\Spec(A) = U \subset X$ 包含 $z$，使得 $A$ 是 Nagata 环。
则有 $Z \cap U \cong \Spec(A/\mathfrak p)$，其中 $\mathfrak p$ 为某个素理想；
见《概形》引理
\ref{schemes-lemma-closed-subspace-scheme}
（以及定义 \ref{properties-definition-integral}）。由《代数》定义
\ref{algebra-definition-nagata} 可知 $A/\mathfrak p$ 是 Japanese 整环。
因此，按定义，$Z$ 是 Japanese 概形。

\medskip\noindent
假设 $X$ 的每个整的闭子概形都是 Japanese 概形。任取仿射开集
$\Spec(A) = U \subset X$。因为 $X$ 局部 Noether，由引理
\ref{properties-lemma-locally-Noetherian} 可知 $A$ 是 Noether 环。令
$\mathfrak p \subset A$ 为素理想。我们须证明 $A/\mathfrak p$ 是
Japanese 整环。令 $T \subset U$ 为闭子集
$V(\mathfrak p) \subset \Spec(A)$，并令 $\overline{T} \subset X$ 为其
闭包。由于 $\overline{T}$ 是不可约子集的闭包，故它不可约。因此，由
$\overline{T}$ 定义的约化闭子概形是整的闭子概形（仍称为
$\overline{T}$）；见《概形》引理
\ref{schemes-lemma-reduced-closed-subscheme}。换言之，
$\Spec(A/\mathfrak p)$ 是 $X$ 的一个整的闭子概形的仿射开集。依假设，
这个子概形是 Japanese 概形；再由引理 \ref{properties-lemma-locally-Japanese}，
可知 $A/\mathfrak p$ 是 Japanese 整环。
\end{proof}

\begin{lemma}
\label{properties-lemma-nagata-universally-Japanese}
令 $X$ 为概形。下列条件等价：
\begin{enumerate}
\item $X$ 是 Nagata 概形；以及
\item $X$ 局部 Noether 且为泛 Japanese 概形。
\end{enumerate}
\end{lemma}

\begin{proof}
这就是《代数》命题
\ref{algebra-proposition-nagata-universally-japanese}。
\end{proof}

\noindent
这一讨论将在《态射》第 \ref{morphisms-section-nagata} 节继续。




\section{G-概形}
\label{properties-section-G-scheme}

\noindent
我们指出：环 $R$ 是 G-环，当且仅当 $R$ 是 Noether 环，并且对每个素理想
$\mathfrak p$（属于 $R$），环同态
$R_\mathfrak p \to (R_\mathfrak p)^\wedge$ 都是正则的。

\begin{definition}
\label{properties-definition-G-scheme}
令 $X$ 为概形。称 $X$ 是{\it G-概形}\footnote{这可能不是标准术语。若
$G$ 是有限群或群概形，则 $G$-概形有时指带有 $G$ 作用的概形。}，如果对每个
$x \in X$，都存在仿射开邻域 $x \in U \subset X$，使得环
$\mathcal{O}_X(U)$ 是 G-环（见《代数续篇》定义
\ref{more-algebra-definition-G-ring}）。
\end{definition}

\begin{lemma}
\label{properties-lemma-characterize-G-scheme}
令 $X$ 为概形。下列条件等价：
\begin{enumerate}
\item $X$ 是 G-概形；
\item $X$ 局部 Noether，并且对所有 $x \in X$，环同态
$\mathcal{O}_{X, x} \to \mathcal{O}_{X, x}^\wedge$ 都正则；以及
\item $X$ 局部 Noether，并且对任意闭点 $x \in X$，环同态
$\mathcal{O}_{X, x} \to \mathcal{O}_{X, x}^\wedge$ 都正则。
\end{enumerate}
\end{lemma}

\begin{proof}
显然，(1) 蕴含 (2)，而 (2) 蕴含 (3)。假设 (3)。任取 $x \in X$。
由引理 \ref{properties-lemma-locally-Noetherian-closed-point}，存在特化
$x \leadsto x'$，其中 $x'$ 在 $X$ 中闭。于是，由《代数续篇》引理
\ref{more-algebra-lemma-check-G-ring-maximal-ideals}，
$\mathcal{O}_{X, x'}$ 是 G-环。由于 $\mathcal{O}_{X, x}$ 是
$\mathcal{O}_{X, x'}$ 在一个素理想处的局部化
（《概形》引理 \ref{schemes-lemma-specialize-points}），可知
$\mathcal{O}_{X, x} \to \mathcal{O}_{X, x}^\wedge$ 是正则环同态。
由于 $x \in X$ 任意，可得：对任意仿射开集 $U \subset X$，Noether 环
（引理 \ref{properties-lemma-locally-Noetherian}）$\mathcal{O}_X(U)$ 都是 G-环。
这证明 (1) 成立。
\end{proof}

\begin{lemma}
\label{properties-lemma-locally-G-scheme}
令 $X$ 为概形。下列条件等价：
\begin{enumerate}
\item 概形 $X$ 是 G-概形。
\item 对每个仿射开集 $U \subset X$，环 $\mathcal{O}_X(U)$ 都是 G-环。
\item 存在仿射开覆盖 $X = \bigcup U_i$，使得每个
$\mathcal{O}_X(U_i)$ 都是 G-环。
\item 存在开覆盖 $X = \bigcup X_j$，使得每个开子概形 $X_j$ 都是
G-概形。
\end{enumerate}
此外，若 $X$ 是 G-概形，则每个开子概形也都是 G-概形。
\end{lemma}

\begin{proof}
结合引理 \ref{properties-lemma-locally-Noetherian}
和 \ref{properties-lemma-characterize-G-scheme}。
\end{proof}



\section{奇异轨迹}
\label{properties-section-singular-locus}

\noindent
定义如下。

\begin{definition}
\label{properties-definition-singular-locus}
令 $X$ 为局部 Noether 概形。{\it 正则轨迹} $\text{Reg}(X)$，即 $X$
的所有 $x \in X$ 中满足 $\mathcal{O}_{X, x}$ 为正则局部环的点所成的集合。
{\it 奇异轨迹} $\text{Sing}(X)$ 是补集
$X \setminus \text{Reg}(X)$，即所有 $x \in X$ 中满足
$\mathcal{O}_{X, x}$ 不是正则局部环的点所成的集合。
\end{definition}

\noindent
局部 Noether 概形的正则轨迹在泛化下稳定；见《代数》定义
\ref{algebra-definition-regular} 之前的讨论。然而，对一般局部 Noether
概形，正则轨迹未必是开的。读者可以在《代数续篇》第
\ref{more-algebra-section-singular-locus} 节找到它成为开集的一些判据。
我们将在《态射》第 \ref{morphisms-section-singular-locus} 节进一步讨论。






\section{局部不可约性}
\label{properties-section-unibranch}

\noindent
回顾：我们已在《代数续篇》第 \ref{more-algebra-section-unibranch} 节
引入了（几何）单枝局部环的概念。

\begin{definition}
\label{properties-definition-unibranch}
\begin{reference}
\cite[Chapter IV (6.15.1)]{EGA4}
\end{reference}
令 $X$ 为概形，并令 $x \in X$。称 $X$ {\it 在 $x$ 处单枝}，如果局部环
$\mathcal{O}_{X, x}$ 单枝。称 $X$ {\it 在 $x$ 处几何单枝}，如果局部环
$\mathcal{O}_{X, x}$ 几何单枝。称 $X$ {\it 单枝}，如果 $X$ 在其所有点处
都单枝。称 $X$ {\it 几何单枝}，如果 $X$ 在其所有点处都几何单枝。
\end{definition}

\noindent
确实可能发生如下情形：局部环 $A$ 在《代数续篇》定义
\ref{more-algebra-definition-unibranch} 的意义下几何单枝，但概形
$\Spec(A)$ 在定义 \ref{properties-definition-unibranch} 的意义下并非几何单枝。
例如，当 $A$ 是一条具有普通二重点奇点（结点）的不可约平面曲线之锥
在顶点处的局部环时，就会发生这种情形。

\begin{lemma}
\label{properties-lemma-normal-geometrically-unibranch}
正规概形是几何单枝的。
\end{lemma}

\begin{proof}
这由各定义得出。具体地说，概形正规是指其各局部环都是正规整环；
由《代数续篇》定义 \ref{more-algebra-definition-unibranch} 立即可知，
正规局部整环几何单枝。
\end{proof}

\begin{lemma}
\label{properties-lemma-geometrically-unibranch}
\begin{reference}
与 \cite[Proposition 2.3]{Etale-coverings} 比较
\end{reference}
令 $X$ 为 Noether 概形。下列条件等价：
\begin{enumerate}
\item $X$ 几何单枝（定义 \ref{properties-definition-unibranch}）；
\item 对每个点 $x \in X$，若它不是 $X$ 的某个不可约分支之泛点，则
严格 Hensel 化 $\mathcal{O}_{X, x}^{sh}$ 的穿孔谱都连通。
\end{enumerate}
\end{lemma}

\begin{proof}
《代数续篇》引理 \ref{more-algebra-lemma-geometrically-unibranch}
说明，(1) 蕴含 (2) 中的穿孔谱不可约，因而特别地连通。

\medskip\noindent
假设 (2)。令 $x \in X$。我们须证明 $\mathcal{O}_{X, x}$ 几何单枝。
对 $\dim(\mathcal{O}_{X, x})$ 归纳，可以假设结论对 $x$ 的每个非平凡
泛化都成立。可以把 $X$ 替换为 $\Spec(\mathcal{O}_{X, x})$。换言之，
可以假设 $X = \Spec(A)$，其中 $A$ 是局部环，并且 $A_\mathfrak p$
对每个非极大素理想 $\mathfrak p \subset A$ 都几何单枝。

\medskip\noindent
令 $A^{sh}$ 为 $A$ 的严格 Hensel 化。若
$\mathfrak q \subset A^{sh}$ 是位于 $\mathfrak p \subset A$ 上方的素理想，
则 $A_\mathfrak p \to A^{sh}_\mathfrak q$ 是 étale 代数的滤过余极限。
因此，$A_\mathfrak p$ 与 $A^{sh}_\mathfrak q$ 的严格 Hensel 化同构。
由此及《代数续篇》引理
\ref{more-algebra-lemma-geometrically-unibranch} 可知：
$A^{sh}_\mathfrak q$ 对每个非极大素理想 $\mathfrak q$（属于 $A^{sh}$）
都有唯一极小素理想。

\medskip\noindent
令 $\mathfrak q_1, \ldots, \mathfrak q_r$ 为 $A^{sh}$ 的极小素理想。
我们须证明 $r = 1$。由上文可知，
$V(\mathfrak q_1) \cap V(\mathfrak q_j) = \{\mathfrak m^{sh}\}$ 对
$j = 2, \ldots, r$ 成立。
因此，$V(\mathfrak q_1) \setminus \{\mathfrak m^{sh}\}$ 是 $A^{sh}$ 的
穿孔谱的开闭子集；这与该穿孔谱连通的假设矛盾，除非 $r = 1$。
\end{proof}

\begin{definition}
\label{properties-definition-number-of-branches}
令 $X$ 为概形，并令 $x \in X$。{\it $X$ 在 $x$ 处的分枝数}，是局部环
$\mathcal{O}_{X, x}$ 的分枝数，其定义见《代数续篇》定义
\ref{more-algebra-definition-number-of-branches}。{\it $X$ 在 $x$ 处的
几何分枝数}，是局部环 $\mathcal{O}_{X, x}$ 的几何分枝数，其定义见
《代数续篇》定义 \ref{more-algebra-definition-number-of-branches}。
\end{definition}

\noindent
我们常想把它与完备局部环的分枝作比较，但一般而言，这种比较并不直接；
关于这个问题的一些信息见《代数续篇》第
\ref{more-algebra-section-branches-completion} 节。

\begin{lemma}
\label{properties-lemma-number-of-branches-irreducible-components}
令 $X$ 为概形且 $x \in X$。令 $X_i$（$i \in I$）为 $X$ 的所有经过
$x$ 的不可约分支。则 $X$ 在 $x$ 处的（几何）分枝数，等于对
$i \in I$ 求和所得的 $X_i$ 在 $x$ 处的（几何）分枝数之和。
\end{lemma}

\begin{proof}
把 $X_i$ 看作 $X$ 的整的闭子概形；见《概形》定义
\ref{schemes-definition-reduced-induced-scheme} 和引理
\ref{properties-lemma-characterize-integral}。注意，$X_i$ 在 $x$ 处
的（几何）分枝数至少为 $1$，对所有 $i$ 都如此（本质上由定义）。回顾：
$X_i$ 按 $1$ 对 $1$ 的方式与极小素理想
$\mathfrak p_i \subset \mathcal{O}_{X, x}$ 对应；见《代数》引理
\ref{algebra-lemma-irreducible-components-containing-x}。因此，若 $I$
无限，则 $\mathcal{O}_{X, x}$ 有无限多个极小素理想，从而
$\mathcal{O}_{X, x}^h$ 与 $\mathcal{O}_{X, x}^{sh}$ 都有无限多个极小
素理想（结合《代数》引理
\ref{algebra-lemma-injective-minimal-primes-in-image}、
\ref{algebra-lemma-minimal-prime-image-minimal-prime} 以及同态
$\mathcal{O}_{X, x} \to  \mathcal{O}_{X, x}^h \to \mathcal{O}_{X, x}^{sh}$
的单射性）。在这种情形下，$X$ 在 $x$ 处的（几何）分枝数按定义为
$\infty$，求和所得结果也是如此。因此可假设 $I$ 有限。
令 $A'$ 为 $\mathcal{O}_{X, x}$ 在总分式环 $Q$ 中的整闭包，而该环
是 $(\mathcal{O}_{X, x})_{red}$ 的总分式环。令 $A'_i$ 为
$\mathcal{O}_{X, x}/\mathfrak p_i$ 在总分式环 $Q_i$ 中的整闭包，
而该环是 $\mathcal{O}_{X, x}/\mathfrak p_i$ 的总分式环。由《代数》引理
\ref{algebra-lemma-total-ring-fractions-no-embedded-points}，
$Q = \prod_{i \in I} Q_i$。因此 $A' = \prod A'_i$。本引理的等式随即由
《代数续篇》引理 \ref{more-algebra-lemma-number-of-branches-1} 得出；
该引理用 $A'$ 的极大理想表出（几何）分枝数。
\end{proof}

\begin{lemma}
\label{properties-lemma-number-of-branches-1}
令 $X$ 为概形，并令 $x \in X$。
\begin{enumerate}
\item $X$ 在 $x$ 处的分枝数为 $1$，当且仅当 $X$ 在 $x$ 处单枝。
\item $X$ 在 $x$ 处的几何分枝数为 $1$，当且仅当 $X$ 在 $x$ 处几何单枝。
\end{enumerate}
\end{lemma}

\begin{proof}
本引理由各定义以及环的相应结论立即得出；见《代数续篇》引理
\ref{more-algebra-lemma-number-of-branches-1}。
\end{proof}




\section{有限型与有限表现模的刻画}
\label{properties-section-characterizing-finite-type-presentation}

\noindent
令 $X$ 为概形。
令 $\mathcal{F}$ 为准凝聚 $\mathcal{O}_X$-模。
下述引理表明，$\mathcal{F}$ 是有限型的
（见《模》定义 \ref{modules-definition-finite-type}）
当且仅当 $\mathcal{F}$ 在每个仿射开子集 $\Spec(A) = U \subset X$ 上
都具有形式 $\widetilde M$，即来自某个有限 $A$-模 $M$。
类似地，$\mathcal{F}$ 是有限表现的
（见《模》定义 \ref{modules-definition-finite-presentation}）
当且仅当 $\mathcal{F}$ 在每个仿射开子集 $\Spec(A) = U \subset X$ 上
都具有形式 $\widetilde M$，其中它来自某个有限表现 $A$-模 $M$。

\begin{lemma}
\label{properties-lemma-finite-type-module}
令 $X = \Spec(R)$ 为仿射概形。
准凝聚 $\mathcal{O}_X$-模层
$\widetilde M$ 是有限型 $\mathcal{O}_X$-模，当且仅当
$M$ 是有限 $R$-模。
\end{lemma}

\begin{proof}
假设 $\widetilde M$ 是有限型 $\mathcal{O}_X$-模。
这意味着存在 $X$ 的一个开覆盖，使得 $\widetilde M$ 在该覆盖各成员
上的限制由有限多个截面整体生成。因此还存在一个标准开覆盖
$X = \bigcup_{i = 1, \ldots, n} D(f_i)$，使得 $\widetilde M|_{D(f_i)}$
由有限多个截面生成。于是 $M_{f_i}$ 对每个 $i$ 都有限生成。
故结论由《代数》引理 \ref{algebra-lemma-cover} 得出。
\end{proof}

\begin{lemma}
\label{properties-lemma-finite-presentation-module}
令 $X = \Spec(R)$ 为仿射概形。准凝聚 $\mathcal{O}_X$-模层
$\widetilde M$ 是有限表现 $\mathcal{O}_X$-模，当且仅当
$M$ 是有限表现 $R$-模。
\end{lemma}

\begin{proof}
假设 $\widetilde M$ 是有限表现 $\mathcal{O}_X$-模。
由引理 \ref{properties-lemma-finite-type-module} 可知 $M$ 是有限 $R$-模。
选取一个满射 $R^n \to M$，记其核为 $K$。由《概形》引理
\ref{schemes-lemma-spec-sheaves}，存在短正合列
$$
0 \to \widetilde{K} \to
\bigoplus \mathcal{O}_X^{\oplus n} \to
\widetilde{M} \to 0
$$
由《模》引理
\ref{modules-lemma-kernel-surjection-finite-free-onto-finite-presentation}
可知 $\widetilde{K}$ 是有限型 $\mathcal{O}_X$-模。
再次应用引理 \ref{properties-lemma-finite-type-module}，
可知 $K$ 是有限 $R$-模。因此 $M$ 是有限表现 $R$-模。
\end{proof}

\section{主开集上的截面}
\label{properties-section-principal-opens}

\noindent
下面是这类结果的一个典型例子。在后面的引理中，我们将采用一种更朴素但
也更直接的证明方法。

\begin{lemma}
\label{properties-lemma-invert-f-sections}
\begin{slogan}
准凝聚层的截面在无穷远处只有亚纯奇点。
\end{slogan}
令 $X$ 为概形。令 $f \in \Gamma(X, \mathcal{O}_X)$。
把 $X_f \subset X$ 记为 $f$ 可逆之处的开集，见《概形》引理
\ref{schemes-lemma-f-open}。
若 $X$ 拟紧且拟分离，则典范映射
$$
\Gamma(X, \mathcal{O}_X)_f \longrightarrow \Gamma(X_f, \mathcal{O}_X)
$$
是同构。此外，若 $\mathcal{F}$ 是准凝聚 $\mathcal{O}_X$-模层，则映射
$$
\Gamma(X, \mathcal{F})_f \longrightarrow \Gamma(X_f, \mathcal{F})
$$
是同构。
\end{lemma}

\begin{proof}
记 $R = \Gamma(X, \mathcal{O}_X)$。考虑概形的典范态射
$$
\varphi : X \longrightarrow \Spec(R)
$$
，见《概形》引理
\ref{schemes-lemma-morphism-into-affine}。
右端标准开集 $D(f)$ 的逆像就是左端的 $X_f$。
此外，因为假设 $X$ 拟紧且拟分离，态射 $\varphi$ 也拟紧且拟分离；
见《概形》引理 \ref{schemes-lemma-quasi-compact-affine} 和
\ref{schemes-lemma-compose-after-separated}。因此由《概形》引理
\ref{schemes-lemma-push-forward-quasi-coherent} 可知
$\varphi_*\mathcal{F}$ 是准凝聚的。
于是 $\varphi_*\mathcal{F} = \widetilde M$，其中
$M = \Gamma(X, \mathcal{F})$ 是一个 $R$-模。因此
$$
\Gamma(X_f, \mathcal{F}) =
\Gamma(D(f), \varphi_*\mathcal{F}) =
\Gamma(D(f), \widetilde M) = M_f
$$
，这恰是本引理的内容。令 $\mathcal{F} = \mathcal{O}_X$，
即可得到本引理中第一个陈列同构。
\end{proof}

\noindent
回顾：给定概形 $X$、可逆层 $\mathcal{L}$（在 $X$ 上），以及
$\mathcal{O}_X$-模层 $\mathcal{F}$，可得到分次环
$\Gamma_*(X, \mathcal{L}) =
\bigoplus\nolimits_{n \geq 0} \Gamma(X, \mathcal{L}^{\otimes n})$
以及分次 $\Gamma_*(X, \mathcal{L})$-模
$\Gamma_*(X, \mathcal{L}, \mathcal{F}) =
\bigoplus\nolimits_{n \in \mathbf{Z}}
\Gamma(X, \mathcal{F} \otimes_{\mathcal{O}_X} \mathcal{L}^{\otimes n})$；
见《模》定义 \ref{modules-definition-gamma-star}。
若还给定一个截面 $s \in \Gamma(X, \mathcal{L})$，则得到映射
\begin{equation}
\label{properties-equation-module-invert-s}
\Gamma_*(X, \mathcal{L}, \mathcal{F})_{(s)}
\longrightarrow
\Gamma(X_s, \mathcal{F}|_{X_s})
\end{equation}
，它把 $t/s^n$（其中
$t \in \Gamma(X, \mathcal{F} \otimes_{\mathcal{O}_X} \mathcal{L}^{\otimes n})$）
映到 $t|_{X_s} \otimes s|_{X_s}^{-n}$。这是有意义的，因为按定义
$X_s \subset X$ 是 $s$ 有逆的开集；见《模》引理
\ref{modules-lemma-s-open}。

\begin{lemma}
\label{properties-lemma-invert-s-sections}
\begin{reference}
\cite[Section 6.8]{EGA1-second}
\end{reference}
令 $X$ 为概形。令 $\mathcal{L}$ 为 $X$ 上的可逆层。
令 $s \in \Gamma(X, \mathcal{L})$。令 $\mathcal{F}$ 为准凝聚
$\mathcal{O}_X$-模。
\begin{enumerate}
\item 若 $X$ 拟紧，则 (\ref{properties-equation-module-invert-s}) 是单射；
\item 若 $X$ 拟紧且拟分离，则
(\ref{properties-equation-module-invert-s}) 是同构。
\end{enumerate}
特别地，若 $X$ 拟紧且拟分离，则典范映射
$$
\Gamma_*(X, \mathcal{L})_{(s)}
\longrightarrow
\Gamma(X_s, \mathcal{O}_X),\quad
a/s^n \longmapsto a \otimes s^{-n}
$$
是同构。
\end{lemma}

\begin{proof}
假设 $X$ 拟紧。选取一个有限仿射开覆盖
$X = U_1 \cup \ldots \cup U_m$，其中 $U_j$ 仿射且
$\mathcal{L}|_{U_j} \cong \mathcal{O}_{U_j}$。经此同构，
像 $s|_{U_j}$ 对应于某个
$f_j \in \Gamma(U_j, \mathcal{O}_{U_j})$。于是
$X_s \cap U_j = D(f_j)$。

\medskip\noindent
(1) 的证明。令 $t/s^n$ 为
(\ref{properties-equation-module-invert-s}) 的核中元素。则 $t|_{X_s} = 0$。
故 $(t|_{U_j})|_{D(f_j)} = 0$。由引理
\ref{properties-lemma-invert-f-sections} 可知
$f_j^{e_j} t|_{U_j} = 0$，其中某个 $e_j \geq 0$。
令 $e = \max(e_j)$。于是 $t \otimes s^e$ 在 $U_j$ 上的限制对所有
$j$ 均为零，
故它为零。由于 $t/s^n$ 等于
$t \otimes s^e/s^{n + e}$（作为
$\Gamma_*(X, \mathcal{L}, \mathcal{F})_{(s)}$ 中的元素），
故如所需有 $t/s^n = 0$。

\medskip\noindent
(2) 的证明。假设 $X$ 拟紧且拟分离。
则 $U_j \cap U_{j'}$ 对所有指标对 $j, j'$ 均拟紧；见《概形》引理
\ref{schemes-lemma-characterize-quasi-separated}。
由 (1) 可知 (\ref{properties-equation-module-invert-s}) 是单射。
令 $t' \in \Gamma(X_s, \mathcal{F}|_{X_s})$。对每个 $j$，存在整数
$e_j \geq 0$ 以及
$t'_j \in \Gamma(U_j, \mathcal{F}|_{U_j})$，使得
$t'|_{D(f_j)}$ 经由引理 \ref{properties-lemma-invert-f-sections} 的同构对应于
$t'_j/f_j^{e_j}$。
令 $e = \max(e_j)$，并置
$$
t_j = f_j^{e - e_j} t'_j \otimes q_j^e \in
\Gamma(U_j,
(\mathcal{F} \otimes_{\mathcal{O}_X} \mathcal{L}^{\otimes e})|_{U_j})
$$
，其中 $q_j \in \Gamma(U_j, \mathcal{L}|_{U_j})$ 是来自同构
$\mathcal{L}|_{U_j} \cong \mathcal{O}_{U_j}$ 的平凡化截面。特别地，
$s|_{U_j} = f_j q_j$。利用这一点计算可见，
$t_j|_{U_j \cap U_{j'}}$ 与 $t_{j'}|_{U_j \cap U_{j'}}$
映到 $\mathcal{F}$ 在 $U_j \cap U_{j'} \cap X_s$ 上的同一截面。
由 $U_j \cap U_{j'}$ 的拟紧性及 (1)，存在整数
$e' \geq 0$，使得
$$
t_j|_{U_j \cap U_{j'}} \otimes s^{e'}|_{U_j \cap U_{j'}} =
t_{j'}|_{U_j \cap U_{j'}} \otimes s^{e'}|_{U_j \cap U_{j'}}
$$
作为 $\mathcal{F} \otimes \mathcal{L}^{\otimes e + e'}$ 在
$U_j \cap U_{j'}$ 上的截面相等。
可选择同一个 $e'$ 对所有指标对 $j, j'$ 都适用。于是层条件给出一个截面
$t \in \Gamma(X, \mathcal{F} \otimes \mathcal{L}^{\otimes e + e'})$，
其在 $U_j$ 上的限制为 $t_j \otimes s^{e'}|_{U_j}$。
简单计算表明，$t/s^{e + e'}$ 映到 $t'$，如所需。
\end{proof}

\noindent
令 $X$ 为概形。令 $\mathcal{L}$ 为可逆 $\mathcal{O}_X$-模。
令 $\mathcal{F}$ 与 $\mathcal{G}$ 为准凝聚 $\mathcal{O}_X$-模。
考虑分次 $\Gamma_*(X, \mathcal{L})$-模
$$
M = \bigoplus\nolimits_{n \in \mathbf{Z}} \Hom_{\mathcal{O}_X}(\mathcal{F},
\mathcal{G} \otimes_{\mathcal{O}_X} \mathcal{L}^{\otimes n})
$$
接着，令 $s \in \Gamma(X, \mathcal{L})$ 为截面。于是有典范映射
\begin{equation}
\label{properties-equation-hom-invert-s}
M_{(s)} \longrightarrow
\Hom_{\mathcal{O}_{X_s}}(\mathcal{F}|_{X_s}, \mathcal{G}|_{X_s})
\end{equation}
，它把 $\alpha/s^n$ 映到映射
$\alpha|_{X_s} \otimes s|_{X_s}^{-n}$。
下述引理与引理 \ref{properties-lemma-extend-finite-presentation} 合在一起，大致说明：
若 $X$ 拟紧且拟分离，则有限表现 $\mathcal{O}_{X_s}$-模范畴就是
有限表现 $\mathcal{O}_X$-模范畴把映射的乘法系
$s^n: \mathcal{F} \to
\mathcal{F} \otimes_{\mathcal{O}_X} \mathcal{L}^{\otimes n}$ 中的映射都
变为可逆之后所得的范畴。

\begin{lemma}
\label{properties-lemma-section-maps-backwards}
令 $X$ 为概形。令 $\mathcal{L}$ 为可逆 $\mathcal{O}_X$-模。
令 $s \in \Gamma(X, \mathcal{L})$ 为截面。
令 $\mathcal{F}$、$\mathcal{G}$ 为准凝聚 $\mathcal{O}_X$-模。
\begin{enumerate}
\item 若 $X$ 拟紧且 $\mathcal{F}$ 是有限型的，则
(\ref{properties-equation-hom-invert-s}) 是单射；
\item 若 $X$ 拟紧且拟分离，并且 $\mathcal{F}$ 是有限表现的，则
(\ref{properties-equation-hom-invert-s}) 是双射。
\end{enumerate}
\end{lemma}

\begin{proof}
先证明 $X = \Spec(A)$ 为仿射且
$\mathcal{L} = \mathcal{O}_X$ 的情形。此时 $s$ 对应于某个
$f \in A$。设
$\mathcal{F} = \widetilde{M}$ 且 $\mathcal{G} = \widetilde{N}$，
其中它们是某些 $A$-模 $M$ 和 $N$。于是本引理
（通过引理 \ref{properties-lemma-finite-type-module} 与
\ref{properties-lemma-finite-presentation-module}）转化为下列代数陈述：
\begin{enumerate}
\item 若 $M$ 是有限 $A$-模，且 $\varphi : M \to N$ 是
$A$-模映射，使得诱导映射 $M_f \to N_f$ 为零，
则 $f^n\varphi = 0$ 对某个 $n$ 成立。
\item 若 $M$ 是有限表现 $A$-模，则
$\Hom_A(M, N)_f = \Hom_{A_f}(M_f, N_f)$。
\end{enumerate}
第二个陈述是《代数》引理
\ref{algebra-lemma-hom-from-finitely-presented}；第一个陈述的证明从略。

\medskip\noindent
接着证明一般 $X$ 的 (1)。
假设 $X$ 拟紧，并选取有限仿射开覆盖
$X = U_1 \cup \ldots \cup U_m$，其中 $U_j$ 仿射且
$\mathcal{L}|_{U_j} \cong \mathcal{O}_{U_j}$。经此同构，
像 $s|_{U_j}$ 对应于某个
$f_j \in \Gamma(U_j, \mathcal{O}_{U_j})$。于是
$X_s \cap U_j = D(f_j)$。
令 $\alpha/s^n$ 为
(\ref{properties-equation-hom-invert-s}) 的核中元素。则 $\alpha|_{X_s} = 0$。
故 $(\alpha|_{U_j})|_{D(f_j)} = 0$。由上面处理的仿射情形，
$f_j^{e_j} \alpha|_{U_j} = 0$ 对某个 $e_j \geq 0$ 成立。
令 $e = \max(e_j)$。于是 $\alpha \otimes s^e$ 在 $U_j$ 上的限制对所有
$j$ 均为零，故它为零。由于 $\alpha/s^n$ 等于
$\alpha \otimes s^e/s^{n + e}$（作为 $M_{(s)}$ 中的元素），
故如所需有 $\alpha/s^n = 0$。

\medskip\noindent
(2) 的证明。由于 $\mathcal{F}$ 是有限表现的，层
$\SheafHom_{\mathcal{O}_X}(\mathcal{F}, \mathcal{G})$ 是准凝聚的；
见《概形》第 \ref{schemes-section-quasi-coherent} 节。此外，显然对所有
$n$ 都有
$$
\SheafHom_{\mathcal{O}_X}(\mathcal{F},
\mathcal{G} \otimes_{\mathcal{O}_X} \mathcal{L}^{\otimes n}) =
\SheafHom_{\mathcal{O}_X}(\mathcal{F}, \mathcal{G})
\otimes_{\mathcal{O}_X} \mathcal{L}^{\otimes n}
$$
因此在这种情形下，结论由引理 \ref{properties-lemma-invert-s-sections} 对
$\SheafHom_{\mathcal{O}_X}(\mathcal{F}, \mathcal{G})$ 的应用得出。
\end{proof}

\section{拟仿射概形}
\label{properties-section-quasi-affine}

\begin{definition}
\label{properties-definition-quasi-affine}
若概形 $X$ 拟紧并且同构于某个仿射概形的开子概形，则称它是
{\it 拟仿射的}。
\end{definition}

\begin{lemma}
\label{properties-lemma-quasi-coherent-quasi-affine}
令 $A$ 为环，并令 $U \subset \Spec(A)$ 为拟紧开子概形。
对于准凝聚层 $\mathcal{F}$（在 $U$ 上），典范映射
$$
\widetilde{H^0(U, \mathcal{F})}|_U \to \mathcal{F}
$$
是同构。
\end{lemma}

\begin{proof}
以 $j : U \to \Spec(A)$ 表示包含态射。则
$H^0(U, \mathcal{F}) = H^0(\Spec(A), j_*\mathcal{F})$，并且
$j_*\mathcal{F}$ 由《概形》引理
\ref{schemes-lemma-push-forward-quasi-coherent} 是准凝聚的。
因此，由《概形》引理 \ref{schemes-lemma-equivalence-quasi-coherent}，
$j_*\mathcal{F} = \widetilde{H^0(U, \mathcal{F})}$。
再限制回 $U$ 即得本引理。
\end{proof}

\begin{lemma}
\label{properties-lemma-invert-f-affine}
令 $X$ 为概形。令 $f \in \Gamma(X, \mathcal{O}_X)$。
假设 $X$ 拟紧且拟分离，并假设 $X_f$ 是仿射的。则典范态射
$$
j : X \longrightarrow \Spec(\Gamma(X, \mathcal{O}_X))
$$
（来自《概形》引理 \ref{schemes-lemma-morphism-into-affine}）
诱导一个同构，把 $X_f = j^{-1}(D(f))$ 同构到标准仿射开集
$D(f) \subset \Spec(\Gamma(X, \mathcal{O}_X))$ 上。
\end{lemma}

\begin{proof}
这是显然的，因为由上面的引理 \ref{properties-lemma-invert-f-sections}，
$j$ 诱导环同构
$\Gamma(X, \mathcal{O}_X)_f \to \mathcal{O}_X(X_f)$。
\end{proof}

\begin{lemma}
\label{properties-lemma-quasi-affine}
令 $X$ 为概形。则 $X$ 是拟仿射的，当且仅当来自《概形》引理
\ref{schemes-lemma-morphism-into-affine} 的典范态射
$$
X \longrightarrow \Spec(\Gamma(X, \mathcal{O}_X))
$$
是拟紧开浸入。
\end{lemma}

\begin{proof}
若所陈列的态射是拟紧开浸入，则 $X$ 同构于
$\Spec(\Gamma(X, \mathcal{O}_X))$ 的一个拟紧开子概形，显然
$X$ 是拟仿射的。

\medskip\noindent
假设 $X$ 是拟仿射的，即 $X \subset \Spec(R)$ 是拟紧开集。
特别地，这蕴涵 $X$ 是分离的；见《概形》引理
\ref{schemes-lemma-subscheme-of-separated-scheme}。
令 $A = \Gamma(X, \mathcal{O}_X)$。
考虑环映射 $R \to A$；它来自
$R = \Gamma(\Spec(R), \mathcal{O}_{\Spec(R)})$
以及层 $\mathcal{O}_{\Spec(R)}$ 的限制映射。
由《概形》引理 \ref{schemes-lemma-morphism-into-affine}，得到包含态射的分解：
$$
X \longrightarrow
\Spec(A) \longrightarrow
\Spec(R)
$$
令 $x \in X$。选取 $r \in R$，使得
$x \in D(r)$ 且 $D(r) \subset X$。以 $f \in A$ 表示 $r$ 在 $A$ 中的像。
上面引理 \ref{properties-lemma-invert-f-sections} 的开集 $X_f$ 等于
$D(r) \subset X$，故由该引理的结论有 $A_f \cong R_r$。
因此 $D(r) \to \Spec(A)$ 是到标准仿射开集 $D(f)$（属于
$\Spec(A)$）上的同构。由于 $X$ 可由这样的仿射开集 $D(f)$ 覆盖，结论得证。
\end{proof}

\begin{lemma}
\label{properties-lemma-cartesian-diagram-quasi-affine}
令 $U \to V$ 为拟仿射概形的开浸入。则
$$
\xymatrix{
U \ar[d] \ar[rr]_-j & & \Spec(\Gamma(U, \mathcal{O}_U)) \ar[d] \\
U \ar[r] & V \ar[r]^-{j'} & \Spec(\Gamma(V, \mathcal{O}_V))
}
$$
是笛卡儿图。
\end{lemma}

\begin{proof}
由《概形》引理 \ref{schemes-lemma-morphism-into-affine}，该图交换。
记 $A = \Gamma(U, \mathcal{O}_U)$ 且
$B = \Gamma(V, \mathcal{O}_V)$。令 $g \in B$ 使得 $V_g$ 仿射且包含于
$U$。这意味着，若 $f$ 是 $g$ 在 $A$ 中的像，则 $U_f = V_g$。
由引理 \ref{properties-lemma-invert-f-affine} 可知，$j'$ 诱导一个同构，把
$V_g$ 同构到标准开集 $D(g)$（属于 $\Spec(B)$）上。因此
$V_g \times_{\Spec(B)} \Spec(A) \to \Spec(A)$ 是到
$D(f) \subset \Spec(A)$ 上的同构。再次由引理
\ref{properties-lemma-invert-f-affine}，$j$ 把 $U_f$ 同构地映到 $D(f)$。
因此可见 $U_f = U_f \times_{\Spec(B)} \Spec(A)$。由于由引理
\ref{properties-lemma-quasi-affine}，可覆盖 $U$，所用开集为上述 $V_g = U_f$，故
$U \to U \times_{\Spec(B)} \Spec(A)$ 是同构。
\end{proof}

\begin{lemma}
\label{properties-lemma-quasi-affine-presentation}
令 $X$ 为拟仿射概形。存在整数 $n \geq 0$、仿射概形 $T$ 以及态射
$T \to X$，使得对每个态射 $X' \to X$，若 $X'$ 仿射，则纤维积
$X' \times_X T$ 同构于 $\mathbf{A}^n_{X'}$，并且这是在 $X'$ 上的同构。
\end{lemma}

\begin{proof}
由定义，存在环 $A$，使得 $X$ 同构于拟紧开子概形
$U \subset \Spec(A)$。回顾：标准开集 $D(f) \subset \Spec(A)$ 构成拓扑的
一组基；见《代数》第 \ref{algebra-section-spectrum-ring} 节。因为 $U$
拟紧，可选取 $f_1, \ldots, f_n \in A$，使得
$U = D(f_1) \cup \ldots \cup D(f_n)$。因此不妨设
$X = \Spec(A) \setminus V(I)$，其中 $I = (f_1, \ldots, f_n)$。置
$$
T = \Spec(A[t, x_1, \ldots, x_n]/(f_1 x_1 + \ldots + f_n x_n - 1))
$$
结构态射 $T \to \Spec(A)$ 通过开集 $X$ 分解，从而给出态射
$T \to X$。若 $X' = \Spec(A')$，且态射 $X' \to X$ 对应于环映射
$A \to A'$，则 $f'_1, \ldots, f'_n \in A'$（即
$f_1, \ldots, f_n$ 的像）生成 $A'$ 中的单位理想。
设 $1 = f'_1 a'_1 + \ldots + f'_n a'_n$。
基变换 $X' \times_X T$ 是下列环的谱：
$A'[t, x_1, \ldots, x_n]/(f'_1 x_1 + \ldots + f'_n x_n - 1)$。
我们断言，下列 $A'$-代数同态
$$
\varphi :
A'[y_1, \ldots, y_n]
\longrightarrow
A'[t, x_1, \ldots, x_n, x_{n + 1}]/(f'_1 x_1 + \ldots + f'_n x_n - 1)
$$
把 $y_i$ 送到 $a'_i t + x_i$，并且它是同构。该断言即完成本引理的证明。
$\varphi$ 的逆由下列 $A'$-代数同态给出：
$$
\psi :
A'[t, x_1, \ldots, x_n, x_{n + 1}]/(f'_1 x_1 + \ldots + f'_n x_n - 1)
\longrightarrow
A'[y_1, \ldots, y_n]
$$
它把 $t$ 送到 $-1 + f'_1 y_1 + \ldots + f'_n y_n$，并把 $x_i$ 送到
$y_i + a'_i - a'_i(f'_1 y_1 + \ldots + f'_n y_n)$，其中
$i = 1, \ldots, n$。这是有意义的，因为 $\sum f'_ix_i$ 被映到
$$
\begin{matrix}
\sum f'_i(y_i + a'_i - a'_i(\sum f'_j y_j)) =
(\sum f'_iy_i) + 1 - (\sum f'_j y_j) = 1
\end{matrix}
$$
为看出这两个映射互逆，作如下计算：
$$
\begin{matrix}
\varphi(\psi(t) = \varphi(-1 +  \sum f'_i y_i) =
-1 + \sum f'_i (a'_i t + x_i) = t \\
\varphi(\psi(x_i)) = \varphi(y_i + a'_i - a'_i(\sum f'_j y_j)) =
a'_i t + x_i + a'_i - a'_i(\sum f'_ja'_jt + f'_jx_j) = x_i \\
\psi(\varphi(y_i)) = \psi(a'_i t + x_i) =
a'_i(-1 + \sum f'_j y_j) + y_i + a'_i - a'_i(\sum f'_j y_j) = y_i
\end{matrix}
$$
证明完毕。
\end{proof}

\section{平坦模}
\label{properties-section-flat}

\noindent
在任意环层空间 $(X, \mathcal{O}_X)$ 上，我们知道一个
$\mathcal{O}_X$-模（在一点处）平坦意味着什么；见《模》定义
\ref{modules-definition-flat}
（定义 \ref{modules-definition-flat-at-point}）。
对于仿射概形上的准凝聚层，这与代数章中定义的概念一致。

\begin{lemma}
\label{properties-lemma-flat-module}
\begin{slogan}
模的平坦性与层的平坦性相同。
\end{slogan}
令 $X = \Spec(R)$ 为仿射概形。
令 $\mathcal{F} = \widetilde{M}$，它来自某个 $R$-模 $M$。
准凝聚层 $\mathcal{F}$ 是平坦 $\mathcal{O}_X$-模，当且仅当
$M$ 是平坦 $R$-模。
\end{lemma}

\begin{proof}
$\mathcal{F}$ 的平坦性可在茎上检验；见《模》引理
\ref{modules-lemma-flat-stalks-flat}。
环上模的情形也是如此；见《代数》引理
\ref{algebra-lemma-flat-localization}。
又因为 $\mathcal{F}_x = M_{\mathfrak p}$（当 $x$ 对应于
$\mathfrak p$ 时），本引理成立。
\end{proof}

\section{局部自由模}
\label{properties-section-finite-locally-free}

\noindent
在任意环层空间上，我们知道一个 $\mathcal{O}_X$-模是（有限）局部自由的
意味着什么。在仿射概形上，这与代数章中定义的概念一致。

\begin{lemma}
\label{properties-lemma-locally-free-module}
令 $X = \Spec(R)$ 为仿射概形。
令 $\mathcal{F} = \widetilde{M}$，它来自某个 $R$-模 $M$。
准凝聚层 $\mathcal{F}$ 是（有限）局部自由 $\mathcal{O}_X$-模，当且仅当
$M$ 是（有限）局部自由 $R$-模。
\end{lemma}

\begin{proof}
这由定义得出；见《模》定义
\ref{modules-definition-locally-free}
以及《代数》定义 \ref{algebra-definition-locally-free}。
\end{proof}

\noindent
有限局部自由模可以用许多不同方式刻画。

\begin{lemma}
\label{properties-lemma-finite-locally-free}
令 $X$ 为概形。
令 $\mathcal{F}$ 为准凝聚 $\mathcal{O}_X$-模。
下列条件等价：
\begin{enumerate}
\item $\mathcal{F}$ 是有限表现的平坦 $\mathcal{O}_X$-模；
\item $\mathcal{F}$ 是有限表现 $\mathcal{O}_X$-模，并且对所有
$x \in X$，茎 $\mathcal{F}_x$ 都是自由 $\mathcal{O}_{X, x}$-模；
\item $\mathcal{F}$ 是有限型局部自由 $\mathcal{O}_X$-模；
\item $\mathcal{F}$ 是有限局部自由 $\mathcal{O}_X$-模；
\item $\mathcal{F}$ 是有限型 $\mathcal{O}_X$-模，对每个 $x \in X$，
茎 $\mathcal{F}_x$ 都是自由 $\mathcal{O}_{X, x}$-模，并且函数
$$
\rho_\mathcal{F} : X \to \mathbf{Z}, \quad
x \longmapsto
\dim_{\kappa(x)} \mathcal{F}_x \otimes_{\mathcal{O}_{X, x}} \kappa(x)
$$
在 $X$ 的 Zariski 拓扑中局部常值。
\end{enumerate}
\end{lemma}

\begin{proof}
本引理立即归结为仿射情形。在这种情形下，它是《代数》引理
\ref{algebra-lemma-finite-projective} 的重述。该转译使用引理
\ref{properties-lemma-finite-type-module}、
\ref{properties-lemma-finite-presentation-module}、
\ref{properties-lemma-flat-module} 和
\ref{properties-lemma-locally-free-module}。
\end{proof}

\begin{lemma}
\label{properties-lemma-finite-locally-free-reduced}
令 $X$ 为约化概形。令 $\mathcal{F}$ 为准凝聚
$\mathcal{O}_X$-模。则引理 \ref{properties-lemma-finite-locally-free} 的等价条件
还等价于
\begin{enumerate}
\item[(6)] $\mathcal{F}$ 是有限型 $\mathcal{O}_X$-模，并且函数
$$
\rho_\mathcal{F} : X \to \mathbf{Z}, \quad
x \longmapsto
\dim_{\kappa(x)} \mathcal{F}_x \otimes_{\mathcal{O}_{X, x}} \kappa(x)
$$
在 $X$ 的 Zariski 拓扑中局部常值。
\end{enumerate}
\end{lemma}

\begin{proof}
本引理立即归结为仿射情形。在这种情形下，它是《代数》引理
\ref{algebra-lemma-finite-projective-reduced} 的重述。
\end{proof}

\section{局部投射模}
\label{properties-section-locally-projective}

\noindent
代数章中已有工作的一个推论是，可以如下定义局部投射模。

\begin{definition}
\label{properties-definition-locally-projective}
令 $X$ 为概形。令 $\mathcal{F}$ 为准凝聚
$\mathcal{O}_X$-模。称 $\mathcal{F}$ 是{\it 局部投射的}，如果对每个
仿射开集 $U \subset X$，$\mathcal{O}_X(U)$-模
$\mathcal{F}(U)$ 都是投射的。
\end{definition}

\begin{lemma}
\label{properties-lemma-locally-projective}
令 $X$ 为概形。
令 $\mathcal{F}$ 为准凝聚 $\mathcal{O}_X$-模。
下列条件等价：
\begin{enumerate}
\item $\mathcal{F}$ 局部投射；
\item 存在仿射开覆盖 $X = \bigcup U_i$，使得
$\mathcal{O}_X(U_i)$-模 $\mathcal{F}(U_i)$ 对每个 $i$ 都是投射的。
\end{enumerate}
特别地，若 $X = \Spec(A)$ 且 $\mathcal{F} = \widetilde{M}$，
则 $\mathcal{F}$ 局部投射，当且仅当 $M$ 是投射 $A$-模。
\end{lemma}

\begin{proof}
首先注意，若 $M$ 是投射 $A$-模且 $A \to B$ 是环映射，则
$M \otimes_A B$ 是投射 $B$-模；见《代数》引理
\ref{algebra-lemma-ascend-properties-modules}。
因此，若 $U$ 是仿射开集，且 $\mathcal{F}(U)$ 是投射
$\mathcal{O}_X(U)$-模，则标准开集 $D(f)$ 是仿射开集，并且
$\mathcal{F}(D(f))$ 是投射 $\mathcal{O}_X(D(f))$-模，这对所有
$f \in \mathcal{O}_X(U)$ 都成立。
假设 (2) 成立。令 $U \subset X$ 为任意仿射开集。
可以找到由有限多个标准开集组成的开覆盖
$U = \bigcup_{j = 1, \ldots, m} D(f_j)$；其中这些标准开集
$D(f_j)$ 满足：对每个 $j$，开集 $D(f_j)$ 都是某个 $U_i$ 的标准开集；
见《概形》引理
\ref{schemes-lemma-standard-open-two-affines}。
因此，若置 $A = \mathcal{O}_X(U)$，且 $M$ 是一个 $A$-模，使
$\mathcal{F}|_U$ 对应于 $M$，则 $M_{f_j}$ 是投射
$A_{f_j}$-模。由此，$A \to B = \prod A_{f_j}$ 是忠实平坦环映射，
且 $M \otimes_A B$ 是投射 $B$-模。因此由《代数》定理
\ref{algebra-theorem-ffdescent-projectivity}，$M$ 是投射的。
\end{proof}

\begin{lemma}
\label{properties-lemma-locally-projective-pullback}
令 $f : X \to Y$ 为概形态射。
令 $\mathcal{G}$ 为准凝聚 $\mathcal{O}_Y$-模。
若 $\mathcal{G}$ 在 $Y$ 上局部投射，则 $f^*\mathcal{G}$
在 $X$ 上局部投射。
\end{lemma}

\begin{proof}
这由《代数》引理 \ref{algebra-lemma-ascend-properties-modules}
和引理 \ref{properties-lemma-locally-projective} 得出。
\end{proof}

\section{准凝聚层的延拓}
\label{properties-section-extending-quasi-coherent-sheaves}

\noindent
有时，能够证明开子概形上给定的准凝聚层可延拓到整个概形，是很有用的。

\begin{lemma}
\label{properties-lemma-extend-trivial}
令 $j : U \to X$ 为概形的拟紧开浸入。
\begin{enumerate}
\item $U$ 上任意准凝聚层都可延拓为 $X$ 上的准凝聚层。
\item 令 $\mathcal{F}$ 为 $X$ 上的准凝聚层。令
$\mathcal{G} \subset \mathcal{F}|_U$ 为准凝聚子层。则存在
准凝聚子层 $\mathcal{H}$（属于 $\mathcal{F}$），使得
$\mathcal{H}|_U = \mathcal{G}$，这里的等式是在
$\mathcal{F}|_U$ 的子层意义下成立。
\item 令 $\mathcal{F}$ 为 $X$ 上的准凝聚层。令 $\mathcal{G}$ 为
$U$ 上的准凝聚层。令
$\varphi : \mathcal{G} \to \mathcal{F}|_U$ 为
$\mathcal{O}_U$-模态射。则存在准凝聚层 $\mathcal{H}$（为
$\mathcal{O}_X$-模）和映射 $\psi : \mathcal{H} \to \mathcal{F}$，使得
$\mathcal{H}|_U = \mathcal{G}$ 且 $\psi|_U = \varphi$。
\end{enumerate}
\end{lemma}

\begin{proof}
浸入是分离的（见《概形》引理
\ref{schemes-lemma-immersions-monomorphisms}），并且 $j$ 依假设拟紧。
因此，对任意准凝聚层 $\mathcal{G}$（在 $U$ 上），层
$j_*\mathcal{G}$ 是到
$X$ 上的一个延拓。见《概形》引理
\ref{schemes-lemma-push-forward-quasi-coherent} 和《层》第
\ref{sheaves-section-open-immersions} 节。

\medskip\noindent
假设 $\mathcal{F}$、$\mathcal{G}$ 如 (2) 中。则
$j_*\mathcal{G}$ 是 $X$ 上的准凝聚层（见上文）。
它是 $j_*j^*\mathcal{F}$ 的子层。因此核
$$
\mathcal{H} =
\Ker(\mathcal{F} \oplus j_* \mathcal{G}
\longrightarrow j_*j^*\mathcal{F})
$$
也是准凝聚的；见《概形》第 \ref{schemes-section-quasi-coherent} 节。
形式地检验可知 $\mathcal{H} \subset \mathcal{F}$ 且
$\mathcal{H}|_U = \mathcal{G}$（再次使用《层》第
\ref{sheaves-section-open-immersions} 节的材料）。

\medskip\noindent
(3) 的证明方式与 (2) 相同。只需取
$\mathcal{H} = \Ker(\mathcal{F} \oplus j_* \mathcal{G}
\to j_*j^*\mathcal{F})$，连同它到 $\mathcal{F}$ 的显然映射，
以及它与 $\mathcal{G}$ 的显然等同（在 $U$ 上）。
\end{proof}

\begin{lemma}
\label{properties-lemma-extend}
令 $X$ 为拟紧且拟分离概形。令 $U \subset X$ 为拟紧开集。
令 $\mathcal{F}$ 为准凝聚 $\mathcal{O}_X$-模。
令 $\mathcal{G} \subset \mathcal{F}|_U$ 为有限型准凝聚
$\mathcal{O}_U$-子模。则存在有限型准凝聚子模
$\mathcal{G}' \subset \mathcal{F}$，使得
$\mathcal{G}'|_U = \mathcal{G}$。
\end{lemma}

\begin{proof}
令 $n$ 为仿射开集 $U_i \subset X$ 的最小数目，其中
$i = 1, \ldots , n$ 且 $X = U \cup \bigcup U_i$。
（这里使用 $X$ 拟紧。）假设可证明 $n = 1$ 的情形。于是可依次把
$\mathcal{G}$ 延拓为 $\mathcal{G}_1$（定义在 $U \cup U_1$ 上），
再延拓为 $\mathcal{G}_2$（定义在 $U \cup U_1 \cup U_2$ 上），
再延拓为 $\mathcal{G}_3$（定义在 $U \cup U_1 \cup U_2 \cup U_3$ 上），
依此类推。因此归结为 $n = 1$ 的情形。

\medskip\noindent
所以可假设 $X = U \cup V$，其中 $V$ 仿射。
因为 $X$ 拟分离，并且 $U$、$V$ 是拟紧开集，所以
$U \cap V$ 是拟紧开集。只需对系统
$(V, U \cap V, \mathcal{F}|_V, \mathcal{G}|_{U \cap V})$
证明本引理，因为可把所得的层 $\mathcal{G}'$（在 $V$ 上）与给定的
层 $\mathcal{G}$（在 $U$ 上）沿着它们在 $U \cap V$ 上的公共值粘合。
因此归结为 $X$ 仿射的情形。

\medskip\noindent
假设 $X = \Spec(R)$。写成 $\mathcal{F} = \widetilde M$，其中
它来自某个 $R$-模 $M$。由上面的引理 \ref{properties-lemma-extend-trivial}，可找到
准凝聚子层 $\mathcal{H} \subset \mathcal{F}$，它限制为
$\mathcal{G}$（在 $U$ 上）。写成 $\mathcal{H} = \widetilde N$，其中它来自某个
$R$-模 $N$。对每个 $u \in U$，存在 $f \in R$，使得
$u \in D(f) \subset U$ 且 $N_f$ 有限生成；见引理
\ref{properties-lemma-finite-type-module}。
由于 $U$ 拟紧，可用有限多个 $D(f_i)$ 覆盖它，使得
$N_{f_i}$ 由有限多个元素生成，设这些元素为
$x_{i, 1}/f_i^N, \ldots, x_{i, r_i}/f_i^N$。
令 $N' \subset N$ 为由元素 $x_{i, j}$ 生成的子模。则子层
$\mathcal{G}' = \widetilde{N'} \subset \mathcal{H} \subset \mathcal{F}$
满足要求。
\end{proof}

\begin{lemma}
\label{properties-lemma-quasi-coherent-colimit-finite-type}
令 $X$ 为拟紧且拟分离概形。
任意准凝聚 $\mathcal{O}_X$-模层都是其有限型准凝聚
$\mathcal{O}_X$-子模的有向余极限。
\end{lemma}

\begin{proof}
该余极限是有向的，因为若 $\mathcal{G}_1$、$\mathcal{G}_2$
是有限型准凝聚子层，则
$\mathcal{G}_1 \oplus \mathcal{G}_2 \to \mathcal{F}$ 的像是有限型
准凝聚子模。
令 $U \subset X$ 为任意仿射开集，并令
$s \in \Gamma(U, \mathcal{F})$ 为任意截面。
令 $\mathcal{G} \subset \mathcal{F}|_U$ 为由 $s$ 生成的子层。
则显然 $\mathcal{G}$ 是准凝聚的，并且作为
$\mathcal{O}_U$-模是有限型的。由引理 \ref{properties-lemma-extend} 可知，
$\mathcal{G}$ 是某个有限型准凝聚子层
$\mathcal{G}' \subset \mathcal{F}$ 的限制。由于 $X$ 的拓扑具有由仿射
开集组成的一组基，可知 $\mathcal{F}$ 的每个局部截面都局部地包含在
某个有限型准凝聚子模中。故结论成立。
\end{proof}

\begin{lemma}
\label{properties-lemma-extend-finite-presentation}
令 $X$ 为拟紧且拟分离概形。
令 $\mathcal{F}$ 为准凝聚 $\mathcal{O}_X$-模。
令 $U \subset X$ 为拟紧开集。
令 $\mathcal{G}$ 为有限表现 $\mathcal{O}_U$-模。
令 $\varphi : \mathcal{G} \to \mathcal{F}|_U$ 为
$\mathcal{O}_U$-模态射。
则存在有限表现 $\mathcal{O}_X$-模 $\mathcal{G}'$，以及
$\mathcal{O}_X$-模态射 $\varphi' : \mathcal{G}' \to \mathcal{F}$，
使得 $\mathcal{G}'|_U = \mathcal{G}$ 且
$\varphi'|_U = \varphi$。
\end{lemma}

\begin{proof}
证明的开头重复引理 \ref{properties-lemma-extend} 证明的开头。不过我们仍将它仔细写出。

\medskip\noindent
令 $n$ 为仿射开集 $U_i \subset X$ 的最小数目，其中
$i = 1, \ldots , n$ 且 $X = U \cup \bigcup U_i$。
（这里使用 $X$ 拟紧。）假设可证明 $n = 1$ 的情形。于是可依次把对
$(\mathcal{G}, \varphi)$ 延拓为对
$(\mathcal{G}_1, \varphi_1)$（定义在 $U \cup U_1$ 上），
再延拓为对 $(\mathcal{G}_2, \varphi_2)$（定义在
$U \cup U_1 \cup U_2$ 上），
再延拓为对 $(\mathcal{G}_3, \varphi_3)$（定义在
$U \cup U_1 \cup U_2 \cup U_3$ 上），依此类推。
因此归结为 $n = 1$ 的情形。

\medskip\noindent
所以可假设 $X = U \cup V$，其中 $V$ 仿射。
因为 $X$ 拟分离且 $U$ 拟紧，所以 $U \cap V \subset V$ 拟紧。
假设对系统
$(V, U \cap V, \mathcal{F}|_V, \mathcal{G}|_{U \cap V}, \varphi|_{U \cap V})$
证明了本引理，从而产生
$(\mathcal{G}', \varphi')$（在 $V$ 上）。那么可把
$\mathcal{G}'$（在 $V$ 上）与给定的层 $\mathcal{G}$（在 $U$ 上）
沿着它们在 $U \cap V$ 上的公共值粘合；
类似地，可把映射 $\varphi'$ 与映射 $\varphi$ 沿着它们在
$U \cap V$ 上的公共值粘合。因此归结为 $X$ 仿射的情形。

\medskip\noindent
假设 $X = \Spec(R)$。
由上面的引理 \ref{properties-lemma-extend-trivial}，可找到准凝聚层
$\mathcal{H}$，以及映射
$\psi : \mathcal{H} \to \mathcal{F}$（在 $X$ 上）；它限制为
$\mathcal{G}$ 和 $\varphi$（在 $U$ 上）。
由引理 \ref{properties-lemma-extend}，可找到有限型准凝聚
$\mathcal{O}_X$-子模 $\mathcal{H}' \subset \mathcal{H}$，使得
$\mathcal{H}'|_U = \mathcal{G}$。因此，替换 $\mathcal{H}$，所用的是
$\mathcal{H}'$；并替换 $\psi$，所用的是 $\psi$ 在 $\mathcal{H}'$ 上的限制。此后，
可假设 $\mathcal{H}$ 是有限型的。
由引理 \ref{properties-lemma-finite-presentation-module} 可知
$\mathcal{H} = \widetilde{N}$，其中 $N$ 是有限生成 $R$-模。
因此存在如下准凝聚 $\mathcal{O}_X$-模短正合列中的满射：
$$
0 \to \mathcal{K} \to \mathcal{O}_X^{\oplus n} \to \mathcal{H} \to 0
$$
其中 $\mathcal{K}$ 定义为该核。
由于 $\mathcal{G}$ 是有限表现的，且由《模》引理
\ref{modules-lemma-kernel-surjection-finite-free-onto-finite-presentation}
有 $\mathcal{H}|_U = \mathcal{G}$，故限制
$\mathcal{K}|_U$ 是有限型 $\mathcal{O}_U$-模。因此再次由引理
\ref{properties-lemma-extend}，存在有限型准凝聚 $\mathcal{O}_X$-子模
$\mathcal{K}' \subset \mathcal{K}$，使得
$\mathcal{K}'|_U = \mathcal{K}|_U$。本引理所提问题的解是置
$$
\mathcal{G}' = \mathcal{O}_X^{\oplus n}/\mathcal{K}'
$$
它显然是有限表现的，并且限制后给出 $\mathcal{G}$（在 $U$ 上），
其中 $\varphi'$ 等于复合
$$
\mathcal{G}' = \mathcal{O}_X^{\oplus n}/\mathcal{K}'
\to \mathcal{O}_X^{\oplus n}/\mathcal{K} = \mathcal{H} \xrightarrow{\psi}
\mathcal{F}.
$$
本引理证毕。
\end{proof}

\begin{lemma}
\label{properties-lemma-lift-finite-presentation}
令 $X$ 为拟紧且拟分离概形。令 $U \subset X$ 为拟紧开集。
令 $\mathcal{G}$ 为 $\mathcal{O}_U$-模。
\begin{enumerate}
\item 若 $\mathcal{G}$ 是有限型准凝聚模，则存在有限型准凝聚
$\mathcal{O}_X$-模 $\mathcal{G}'$，使得
$\mathcal{G}'|_U = \mathcal{G}$。
\item 若 $\mathcal{G}$ 是有限表现的，则存在有限表现
$\mathcal{O}_X$-模 $\mathcal{G}'$，使得
$\mathcal{G}'|_U = \mathcal{G}$。
\end{enumerate}
\end{lemma}

\begin{proof}
(2) 是引理 \ref{properties-lemma-extend-finite-presentation} 中
$\mathcal{F} = 0$ 的特殊情形。对于 (1)，先由引理
\ref{properties-lemma-extend-trivial} 把
$\mathcal{G} = \mathcal{F}|_U$ 写为某个准凝聚
$\mathcal{O}_X$-模的限制，然后对 $\mathcal{G} = \mathcal{F}|_U$
应用引理 \ref{properties-lemma-extend}。
\end{proof}

\noindent
下述引理说明，拟紧且拟分离概形上的每个准凝聚层都是有限表现
$\mathcal{O}$-模的滤过余极限。事实上，在下一个引理中，我们将其改述为
（也许更熟悉的）有向集上的有向余极限。

\begin{lemma}
\label{properties-lemma-directed-colimit-diagram-finite-presentation}
\begin{slogan}
拟紧且拟分离概形上的准凝聚模是有限表现模的滤过余极限。
\end{slogan}
令 $X$ 为概形。假设 $X$ 拟紧且拟分离。
令 $\mathcal{F}$ 为准凝聚 $\mathcal{O}_X$-模。
则存在
\begin{enumerate}
\item 一个滤过指标范畴 $\mathcal{I}$（见《范畴》定义
\ref{categories-definition-directed}）；
\item 一个图 $\mathcal{I} \to \textit{Mod}(\mathcal{O}_X)$（见《范畴》第
\ref{categories-section-limits} 节），
$i \mapsto \mathcal{F}_i$；
\item $\mathcal{O}_X$-模态射
$\varphi_i : \mathcal{F}_i \to \mathcal{F}$；
\end{enumerate}
使得每个 $\mathcal{F}_i$ 都是有限表现的，并且态射 $\varphi_i$
诱导同构
$$
\colim_i \mathcal{F}_i
=
\mathcal{F}.
$$
\end{lemma}

\begin{proof}
选取集合 $I$，并对每个 $i \in I$ 选取一个有限表现
$\mathcal{O}_X$-模，以及一个 $\mathcal{O}_X$-模同态
$\varphi_i : \mathcal{F}_i \to \mathcal{F}$，使其具有如下性质：
对任意 $\psi : \mathcal{G} \to \mathcal{F}$，若 $\mathcal{G}$
有限表现，则存在 $i \in I$，以及同构
$\alpha : \mathcal{F}_i \to \mathcal{G}$，使得
$\varphi_i = \psi \circ \alpha$。由《模》引理
\ref{modules-lemma-set-isomorphism-classes-finite-type-modules}
显然存在这样的集合（也见其证明）。
以 $\mathcal{I}$ 表示如下范畴：
$\Ob(\mathcal{I}) = I$，并且给定 $i, i' \in I$，置
$$
\Mor_\mathcal{I}(i, i') =
\{\alpha : \mathcal{F}_i \to \mathcal{F}_{i'} \mid
\alpha \circ \varphi_{i'} = \varphi_i
\}.
$$
我们断言 $\mathcal{I}$ 是滤过范畴，并且
$\mathcal{F} = \colim_i \mathcal{F}_i$。

\medskip\noindent
令 $i, i' \in I$。可以考虑态射
$$
\mathcal{F}_i \oplus \mathcal{F}_{i'} \longrightarrow \mathcal{F}
$$
，它是 $\varphi_i$ 与 $\varphi_{i'}$ 的直和。
因为有限表现 $\mathcal{O}_X$-模的直和仍有限表现，故存在某个
$i'' \in I$，使得
$\varphi_{i''} : \mathcal{F}_{i''} \to \mathcal{F}$
同构于上面陈列的到 $\mathcal{F}$ 的箭头。
由于有交换图
$$
\xymatrix{
\mathcal{F}_i \ar[r] \ar[d] & \mathcal{F} \ar@{=}[d] \\
\mathcal{F}_i \oplus \mathcal{F}_{i'} \ar[r] & \mathcal{F}
}
\quad
\text{且}
\quad
\xymatrix{
\mathcal{F}_{i'} \ar[r] \ar[d] & \mathcal{F} \ar@{=}[d] \\
\mathcal{F}_i \oplus \mathcal{F}_{i'} \ar[r] & \mathcal{F}
}
$$
可知有态射 $i \to i''$ 和 $i' \to i''$，它们属于 $\mathcal{I}$。
接着，假设给定 $i, i' \in I$ 以及态射
$\alpha, \beta : i \to i'$（它们对应于 $\mathcal{O}_X$-模映射
$\alpha, \beta : \mathcal{F}_i \to \mathcal{F}_{i'}$）。
此时考虑余等化子
$$
\mathcal{G} =
\Coker(
\mathcal{F}_i \xrightarrow{\alpha - \beta} \mathcal{F}_{i'}
)
$$
注意，$\mathcal{G}$ 是有限表现 $\mathcal{O}_X$-模。
因为根据范畴 $\mathcal{I}$ 中态射的定义，
$\varphi_{i'} \circ \alpha = \varphi_{i'} \circ \beta$，
所以得到诱导映射 $\psi : \mathcal{G} \to \mathcal{F}$。
因此，对 $(\mathcal{G}, \psi)$ 再次同构于对
$(\mathcal{F}_{i''}, \varphi_{i''})$，其中某个 $i''$。
从而存在态射 $i' \to i''$（在 $\mathcal{I}$ 中），它等化
$\alpha$ 与 $\beta$。故已证明范畴 $\mathcal{I}$ 是滤过的。

\medskip\noindent
还须证明该图的余极限是 $\mathcal{F}$。根据余极限的定义以及我们对范畴
$\mathcal{I}$ 的定义，有典范映射
$$
\varphi :
\colim_i \mathcal{F}_i
\longrightarrow
\mathcal{F}.
$$
选取 $x \in X$。我们来证明 $\varphi_x$ 是同构。回顾
$$
(\colim_i \mathcal{F}_i)_x
=
\colim_i \mathcal{F}_{i, x},
$$
见《层》第 \ref{sheaves-section-limits-sheaves} 节。
先证明映射 $\varphi_x$ 是单射。
假设 $s \in \mathcal{F}_{i, x}$ 为一个元素，并且 $s$ 映到
$\mathcal{F}_x$ 中的零。则存在拟紧开集 $U$，使得 $s$ 来自
$s \in \mathcal{F}_i(U)$，并且
$\varphi_i(s) = 0$ 在 $\mathcal{F}(U)$ 中成立。
由引理 \ref{properties-lemma-extend}，可找到有限型准凝聚子层
$\mathcal{K} \subset \Ker(\varphi_i)$，它限制为准凝聚
$\mathcal{O}_U$-子模（属于 $\mathcal{F}_i$），而该子模由 $s$ 生成：
$\mathcal{K}|_U = \mathcal{O}_U\cdot s \subset \mathcal{F}_i|_U$。
显然，$\mathcal{F}_i/\mathcal{K}$ 是有限表现的，并且映射
$\varphi_i$ 通过商映射
$\mathcal{F}_i \to \mathcal{F}_i/\mathcal{K}$ 分解。因此可找到
$i' \in I$ 以及态射
$\alpha : \mathcal{F}_i \to \mathcal{F}_{i'}$（在 $\mathcal{I}$ 中），
它可与商映射
$\mathcal{F}_i \to \mathcal{F}_i/\mathcal{K}$ 等同。
于是截面 $s$ 映到 $\mathcal{F}_{i'}(U)$ 中的零，特别地，它在
$(\colim_i \mathcal{F}_i)_x =
\colim_i \mathcal{F}_{i, x}$ 中也为零。单射性得证。
最后证明映射 $\varphi_x$ 是满射。选取 $s \in \mathcal{F}_x$。
选取拟紧开邻域 $U \subset X$（它包含 $x$），使得 $s$ 对应于截面
$s \in \mathcal{F}(U)$。考虑映射
$s : \mathcal{O}_U \to \mathcal{F}$（乘以 $s$）。
由引理 \ref{properties-lemma-extend-finite-presentation}，存在有限表现
$\mathcal{O}_X$-模 $\mathcal{G}$，以及 $\mathcal{O}_X$-模映射
$\mathcal{G} \to \mathcal{F}$，使得 $\mathcal{G}|_U \to \mathcal{F}|_U$
可等同于
$s : \mathcal{O}_U \to \mathcal{F}$。
再次根据 $\mathcal{I}$ 的定义，存在 $i \in I$，使得
$\mathcal{G} \to \mathcal{F}$ 同构于
$\varphi_i : \mathcal{F}_i \to \mathcal{F}$。显然存在截面
$s' \in \mathcal{F}_i(U)$，它映到 $s \in \mathcal{F}(U)$。
这证明了满射性，也完成了本引理的证明。
\end{proof}

\begin{lemma}
\label{properties-lemma-directed-colimit-finite-presentation}
令 $X$ 为概形。假设 $X$ 拟紧且拟分离。
令 $\mathcal{F}$ 为准凝聚 $\mathcal{O}_X$-模。
则存在
\begin{enumerate}
\item 一个有向集 $I$（见《范畴》定义
\ref{categories-definition-directed-set}）；
\item 一个系统 $(\mathcal{F}_i, \varphi_{ii'})$，它定义在 $I$ 上并属于
$\textit{Mod}(\mathcal{O}_X)$（见《范畴》定义
\ref{categories-definition-system-over-poset}）；
\item $\mathcal{O}_X$-模态射
$\varphi_i : \mathcal{F}_i \to \mathcal{F}$；
\end{enumerate}
使得每个 $\mathcal{F}_i$ 都是有限表现的，并且态射 $\varphi_i$
诱导同构
$$
\colim_i \mathcal{F}_i
=
\mathcal{F}.
$$
\end{lemma}

\begin{proof}
这是引理 \ref{properties-lemma-directed-colimit-diagram-finite-presentation} 与
《范畴》引理 \ref{categories-lemma-directed-category-system} 的直接推论
（再结合 $\mathcal{O}_X$-模层范畴中余极限存在这一事实；见《层》第
\ref{sheaves-section-limits-sheaves} 节）。
\end{proof}

\begin{lemma}
\label{properties-lemma-finite-directed-colimit-surjective-maps}
令 $X$ 为概形。假设 $X$ 拟紧且拟分离。
令 $\mathcal{F}$ 为有限型准凝聚 $\mathcal{O}_X$-模。
则可写成 $\mathcal{F} = \colim \mathcal{F}_i$，其中
$\mathcal{F}_i$ 是有限表现的，并且所有过渡映射
$\mathcal{F}_i \to \mathcal{F}_{i'}$ 都是满射。
\end{lemma}

\begin{proof}
把 $\mathcal{F} = \colim \mathcal{G}_i$ 写成有限表现
$\mathcal{O}_X$-模的滤过余极限
（引理 \ref{properties-lemma-directed-colimit-finite-presentation}）。
我们断言 $\mathcal{G}_i \to \mathcal{F}$ 对某个 $i$ 是满射。
事实上，选取有限仿射开覆盖 $X = U_1 \cup \ldots \cup U_m$。
选取截面 $s_{jl} \in \mathcal{F}(U_j)$ 来生成
$\mathcal{F}|_{U_j}$；见引理 \ref{properties-lemma-finite-type-module}。
由《层》引理 \ref{sheaves-lemma-directed-colimits-sections}，
$s_{jl}$ 属于 $\mathcal{G}_i \to \mathcal{F}$ 的像，只要 $i$ 足够大。
因此，$\mathcal{G}_i \to \mathcal{F}$ 是满射，只要 $i$ 足够大。
选取这样的 $i$，并令 $\mathcal{K} \subset \mathcal{G}_i$ 为映射
$\mathcal{G}_i \to \mathcal{F}$ 的核。把
$\mathcal{K} = \colim \mathcal{K}_a$ 写成其有限型准凝聚子模的滤过余极限
（引理 \ref{properties-lemma-quasi-coherent-colimit-finite-type}）。于是
$\mathcal{F} = \colim \mathcal{G}_i/\mathcal{K}_a$ 是本引理所提问题的解。
\end{proof}

\begin{lemma}
\label{properties-lemma-application-directed-colimit}
令 $X$ 为拟紧且拟分离概形。
令 $\mathcal{F}$ 为有限型准凝聚 $\mathcal{O}_X$-模。
令 $U \subset X$ 为拟紧开集，并使 $\mathcal{F}|_U$ 是有限表现的。
则存在 $\mathcal{O}_X$-模映射
$\varphi : \mathcal{G} \to \mathcal{F}$，满足
(a) $\mathcal{G}$ 有限表现，
(b) $\varphi$ 为满射，并且
(c) $\varphi|_U$ 为同构。
\end{lemma}

\begin{proof}
把 $\mathcal{F} = \colim \mathcal{F}_i$ 写成有向余极限，其中每个
$\mathcal{F}_i$ 都有限表现；见引理
\ref{properties-lemma-directed-colimit-finite-presentation}。
选取有限仿射开覆盖 $X = \bigcup V_j$，并选取有限多个截面
$s_{jl} \in \mathcal{F}(V_j)$ 生成 $\mathcal{F}|_{V_j}$；见引理
\ref{properties-lemma-finite-type-module}。
由《层》引理 \ref{sheaves-lemma-directed-colimits-sections}，
$s_{jl}$ 属于 $\mathcal{F}_i \to \mathcal{F}$ 的像，只要 $i$ 足够大。
因此，$\mathcal{F}_i \to \mathcal{F}$ 是满射，只要 $i$ 足够大。
选取这样的 $i$，并令 $\mathcal{K} \subset \mathcal{F}_i$ 为映射
$\mathcal{F}_i \to \mathcal{F}$ 的核。由于 $\mathcal{F}_U$ 是有限表现的，
可知 $\mathcal{K}|_U$ 是有限型的；见《模》引理
\ref{modules-lemma-kernel-surjection-finite-free-onto-finite-presentation}。
因此可找到有限型准凝聚子模
$\mathcal{K}' \subset \mathcal{K}$，满足
$\mathcal{K}'|_U = \mathcal{K}|_U$；见引理 \ref{properties-lemma-extend}。于是
$\mathcal{G} = \mathcal{F}_i/\mathcal{K}'$ 连同给定映射
$\mathcal{G} \to \mathcal{F}$ 即为所求。
\end{proof}

\noindent
令 $X$ 为概形。在下述引理中，我们使用{\it 有限表现准凝聚
$\mathcal{O}_X$-代数 $\mathcal{A}$}这一概念。它意味着，对每个仿射开集
$\Spec(R) \subset X$，都有 $\mathcal{A} = \widetilde{A}$，其中 $A$
是一个（交换）$R$-代数，并且作为 $R$-代数有限表现。

\begin{lemma}
\label{properties-lemma-algebra-directed-colimit-finite-presentation}
令 $X$ 为概形。假设 $X$ 拟紧且拟分离。
令 $\mathcal{A}$ 为准凝聚 $\mathcal{O}_X$-代数。
则存在
\begin{enumerate}
\item 一个有向集 $I$（见《范畴》定义
\ref{categories-definition-directed-set}）；
\item 一个系统 $(\mathcal{A}_i, \varphi_{ii'})$，它定义在 $I$ 上并属于
$\mathcal{O}_X$-代数范畴；
\item $\mathcal{O}_X$-代数态射
$\varphi_i : \mathcal{A}_i \to \mathcal{A}$；
\end{enumerate}
使得每个 $\mathcal{A}_i$ 都是有限表现准凝聚
$\mathcal{O}_X$-代数，并且态射 $\varphi_i$ 诱导同构
$$
\colim_i \mathcal{A}_i
=
\mathcal{A}.
$$
\end{lemma}

\begin{proof}
先如引理 \ref{properties-lemma-directed-colimit-finite-presentation}，
把 $\mathcal{A} = \colim_i \mathcal{F}_i$ 写成有限表现准凝聚层的有向余极限。
对每个 $i$，令 $\mathcal{B}_i = \text{Sym}(\mathcal{F}_i)$ 为
$\mathcal{F}_i$ 在 $\mathcal{O}_X$ 上的对称代数。记
$\mathcal{I}_i = \Ker(\mathcal{B}_i \to \mathcal{A})$。写成
$\mathcal{I}_i = \colim_j \mathcal{F}_{i, j}$，其中
$\mathcal{F}_{i, j}$ 是 $\mathcal{I}_i$ 的有限型准凝聚子模；见引理
\ref{properties-lemma-quasi-coherent-colimit-finite-type}。
令 $\mathcal{I}_{i, j} \subset \mathcal{I}_i$ 等于
$\mathcal{B}_i$-理想，该理想由 $\mathcal{F}_{i, j}$ 生成。置
$\mathcal{A}_{i, j} = \mathcal{B}_i/\mathcal{I}_{i, j}$。
则 $\mathcal{A}_{i, j}$ 是有限表现准凝聚
$\mathcal{O}_X$-代数。定义 $(i, j) \leq (i', j')$ 当且仅当
$i \leq i'$，并且映射 $\mathcal{B}_i \to \mathcal{B}_{i'}$
把理想 $\mathcal{I}_{i, j}$ 映入理想 $\mathcal{I}_{i', j'}$。
于是显然 $\mathcal{A} = \colim_{i, j} \mathcal{A}_{i, j}$。
\end{proof}

\noindent
令 $X$ 为概形。在下述引理中，我们使用{\it 有限型准凝聚
$\mathcal{O}_X$-代数 $\mathcal{A}$}这一概念。它意味着，对每个仿射开集
$\Spec(R) \subset X$，都有 $\mathcal{A} = \widetilde{A}$，其中 $A$
是一个（交换）$R$-代数，并且作为 $R$-代数有限型。

\begin{lemma}
\label{properties-lemma-algebra-directed-colimit-finite-type}
令 $X$ 为概形。假设 $X$ 拟紧且拟分离。
令 $\mathcal{A}$ 为准凝聚 $\mathcal{O}_X$-代数。
则 $\mathcal{A}$ 是其有限型准凝聚 $\mathcal{O}_X$-子代数的有向余极限。
\end{lemma}

\begin{proof}
若 $\mathcal{A}_1, \mathcal{A}_2 \subset \mathcal{A}$ 是有限型准凝聚
$\mathcal{O}_X$-子代数，则
$\mathcal{A}_1 \otimes_{\mathcal{O}_X} \mathcal{A}_2 \to \mathcal{A}$
的像也是有限型准凝聚 $\mathcal{O}_X$-子代数（略去若干细节），并且包含
$\mathcal{A}_1$ 和 $\mathcal{A}_2$。由此可见该系统是有向的。
为证明 $\mathcal{A}$ 是该系统的余极限，如引理
\ref{properties-lemma-algebra-directed-colimit-finite-presentation}，
把 $\mathcal{A} = \colim_i \mathcal{A}_i$ 写成有限表现准凝聚
$\mathcal{O}_X$-代数的有向余极限。则像
$\mathcal{A}'_i = \Im(\mathcal{A}_i \to \mathcal{A})$
是 $\mathcal{A}$ 的有限型准凝聚子代数。由于
$\mathcal{A}$ 是这些子代数的余极限，结论得证。
\end{proof}
\noindent
令 $X$ 为概形。在下述引理中，我们使用{\it 有限（分别地，整的）准凝聚
$\mathcal{O}_X$-代数 $\mathcal{A}$}这一概念。它意味着，对每个仿射开集
$\Spec(R) \subset X$，都有 $\mathcal{A} = \widetilde{A}$，其中 $A$
是一个（交换）$R$-代数，并且作为 $R$-代数是有限的（分别地，是整的）。

\begin{lemma}
\label{properties-lemma-finite-algebra-directed-colimit-finite-finitely-presented}
令 $X$ 为概形。假设 $X$ 拟紧且拟分离。
令 $\mathcal{A}$ 为有限准凝聚 $\mathcal{O}_X$-代数。
则 $\mathcal{A} = \colim \mathcal{A}_i$ 是有限且有限表现的准凝聚
$\mathcal{O}_X$-代数的有向余极限，并且所有过渡映射
$\mathcal{A}_{i'} \to \mathcal{A}_i$ 都是满射。
\end{lemma}

\begin{proof}
由引理 \ref{properties-lemma-finite-directed-colimit-surjective-maps}，存在有限表现的
$\mathcal{O}_X$-模 $\mathcal{F}$ 以及满射
$\mathcal{F} \to \mathcal{A}$。利用代数结构，得到满射
$$
\text{Sym}^*_{\mathcal{O}_X}(\mathcal{F}) \longrightarrow \mathcal{A}
$$
以 $\mathcal{J}$ 表示其核。把 $\mathcal{J} = \colim \mathcal{E}_i$
写成有限型 $\mathcal{O}_X$-子模 $\mathcal{E}_i$ 的滤过余极限
（引理 \ref{properties-lemma-quasi-coherent-colimit-finite-type}）。置
$$
\mathcal{A}_i = \text{Sym}^*_{\mathcal{O}_X}(\mathcal{F})/(\mathcal{E}_i)
$$
其中 $(\mathcal{E}_i)$ 表示由映射
$\mathcal{E}_i \to \text{Sym}^*_{\mathcal{O}_X}(\mathcal{F})$ 的像生成的理想层。
于是每个 $\mathcal{A}_i$ 都是有限表现的 $\mathcal{O}_X$-代数，
过渡映射都是满射，并且 $\mathcal{A} = \colim \mathcal{A}_i$。
为完成证明，还须说明 $\mathcal{A}_i$ 是有限
$\mathcal{O}_X$-代数，只要 $i$ 足够大。为此，选取仿射开覆盖
$X = U_1 \cup \ldots \cup U_m$。选取生成元
$f_{j, 1}, \ldots, f_{j, N_j} \in \Gamma(U_i, \mathcal{F})$。
因为 $\mathcal{A}(U_j)$ 是有限 $\mathcal{O}_X(U_j)$-代数，所以对每个
$k$，都存在首一多项式 $P_{j, k} \in \mathcal{O}(U_j)[T]$，使得
$P_{j, k}(f_{j, k})$ 在 $\mathcal{A}(U_j)$ 中为零。由于依构造有
$\mathcal{A} = \colim \mathcal{A}_i$，故有
$P_{j, k}(f_{j, k}) = 0$ 在 $\mathcal{A}_i(U_j)$ 中成立，只要 $i$ 足够大。对这样的
$i$，代数 $\mathcal{A}_i$ 是有限的。
\end{proof}

\begin{lemma}
\label{properties-lemma-integral-algebra-directed-colimit-finite}
令 $X$ 为概形。假设 $X$ 拟紧且拟分离。
令 $\mathcal{A}$ 为整的准凝聚 $\mathcal{O}_X$-代数。则
\begin{enumerate}
\item $\mathcal{A}$ 是其有限准凝聚 $\mathcal{O}_X$-子代数的有向余极限；
\item $\mathcal{A}$ 是有限且有限表现的准凝聚 $\mathcal{O}_X$-代数的直余极限。
\end{enumerate}
\end{lemma}

\begin{proof}
由引理 \ref{properties-lemma-algebra-directed-colimit-finite-type}，有
$\mathcal{A} = \colim \mathcal{A}_i$，其中
$\mathcal{A}_i \subset \mathcal{A}$ 遍历有限型准凝聚
$\mathcal{O}_X$-代数。任意有限型准凝聚 $\mathcal{O}_X$-子代数（属于
$\mathcal{A}$）都是有限的（把《代数》引理
\ref{algebra-lemma-characterize-finite-in-terms-of-integral}
应用于 $\mathcal{A}_i(U) \subset \mathcal{A}(U)$，其中 $U$ 是 $X$
中的仿射开集）。这证明了 (1)。

\medskip\noindent
为证明 (2)，利用引理 \ref{properties-lemma-directed-colimit-finite-presentation}，
把 $\mathcal{A} = \colim \mathcal{F}_i$ 写成有限表现
$\mathcal{O}_X$-模的余极限。对每个 $i$，令 $\mathcal{J}_i$ 为映射
$$
\text{Sym}^*_{\mathcal{O}_X}(\mathcal{F}_i) \longrightarrow \mathcal{A}
$$
的核。若 $i' \geq i$，则有诱导映射
$\mathcal{J}_i \to \mathcal{J}_{i'}$，并且
$\mathcal{A} =
\colim \text{Sym}^*_{\mathcal{O}_X}(\mathcal{F}_i)/\mathcal{J}_i$。
此外，准凝聚 $\mathcal{O}_X$-代数
$\text{Sym}^*_{\mathcal{O}_X}(\mathcal{F}_i)/\mathcal{J}_i$ 是有限的
（见上文）。把 $\mathcal{J}_i = \colim \mathcal{E}_{ik}$ 写成有限表现
$\mathcal{O}_X$-模的余极限。给定 $i' \geq i$ 与 $k$，存在 $k'$，使得有映射
$\mathcal{E}_{ik} \to \mathcal{E}_{i'k'}$，并使图
$$
\xymatrix{
\mathcal{J}_i \ar[r] & \mathcal{J}_{i'} \\
\mathcal{E}_{ik} \ar[u] \ar[r] & \mathcal{E}_{i'k'} \ar[u]
}
$$
交换。这由《模》引理
\ref{modules-lemma-finite-presentation-quasi-compact-colimit} 得出。
它诱导映射
$$
\mathcal{A}_{ik} =
\text{Sym}^*_{\mathcal{O}_X}(\mathcal{F}_i)/(\mathcal{E}_{ik})
\longrightarrow
\text{Sym}^*_{\mathcal{O}_X}(\mathcal{F}_{i'})/(\mathcal{E}_{i'k'}) =
\mathcal{A}_{i'k'}
$$
其中 $(\mathcal{E}_{ik})$ 表示由 $\mathcal{E}_{ik}$ 生成的理想。
准凝聚 $\mathcal{O}_X$-代数 $\mathcal{A}_{ki}$ 当 $k$ 足够大时
既有限表现又有限（见引理
\ref{properties-lemma-finite-algebra-directed-colimit-finite-finitely-presented} 的证明）。
最后，有
$$
\colim \mathcal{A}_{ik} = \colim \mathcal{A}_i = \mathcal{A}
$$
事实上，第一个等式已在引理
\ref{properties-lemma-finite-algebra-directed-colimit-finite-finitely-presented}
的证明中说明；第二个等式成立，是因为 $\mathcal{A}$ 是模
$\mathcal{F}_i$ 的余极限。
\end{proof}





\section{Gabber 的结果}
\label{properties-section-gabber}

\noindent
本节证明 Gabber 的一个结果：对每个概形，都存在一个基数 $\kappa$，
使得每个准凝聚模 $\mathcal{F}$ 都是它的 $\kappa$-生成准凝聚子层之并。
由此可知，概形上的准凝聚层范畴是具有极限和充分多内射对象的
Grothendieck 阿贝尔范畴\footnote{Akhil Mathew 在一篇
\href{https://amathew.wordpress.com/2011/07/30/quasi-coherent-sheaves-presentable-categories-and-a-result-of-gabber/}{博客文章}
中给出了很好的解释。}。

\begin{definition}
\label{properties-definition-kappa-generated}
令 $(X, \mathcal{O}_X)$ 为环层空间，令 $\kappa$ 为无限基数。
称 $\mathcal{O}_X$-模层 $\mathcal{F}$ 是{\it $\kappa$-生成的}，
如果存在开覆盖 $X = \bigcup U_i$，使得 $\mathcal{F}|_{U_i}$ 由子集
$R_i \subset \mathcal{F}(U_i)$ 生成，并且该子集的基数至多为
$\kappa$。
\end{definition}

\noindent
注意，至多 $\kappa$ 个 $\kappa$-生成模的直和仍是
$\kappa$-生成的，因为 $\kappa \otimes \kappa = \kappa$；见《集合》第
\ref{sets-section-cardinals} 节。
特别地，两个 $\kappa$-生成模的直和具有此性质。
此外，$\kappa$-生成层的商层仍是 $\kappa$-生成的。
（但对子模未必如此。）

\begin{lemma}
\label{properties-lemma-set-of-iso-classes}
令 $(X, \mathcal{O}_X)$ 为环层空间，令 $\kappa$ 为基数。
则存在集合 $T$ 以及族 $(\mathcal{F}_t)_{t \in T}$，其成员是
$\kappa$-生成 $\mathcal{O}_X$-模，使得每个 $\kappa$-生成
$\mathcal{O}_X$-模都同构于某个 $\mathcal{F}_t$。
\end{lemma}

\begin{proof}
$X$ 的覆盖组成一个集合（只要不允许重复）。
假设 $X = \bigcup U_i$ 是一个覆盖，并且 $\mathcal{F}_i$
是 $\mathcal{O}_{U_i}$-模。则具有性质
$\mathcal{O}_X$-模 $\mathcal{F}$ 的同构类组成一个集合，其中
$\mathcal{F}|_{U_i} \cong \mathcal{F}_i$，因为粘合映射组成一个集合。
这样，问题化为证明：商映射
$\oplus_{k \in \kappa} \mathcal{O}_X \to \mathcal{F}$ 的（同构类）
对任意环层空间 $X$ 都组成一个集合。这是显然的。
\end{proof}

\noindent
下面就是本节标题所指的结果。

\begin{lemma}
\label{properties-lemma-colimit-kappa}
令 $X$ 为概形。则存在一个基数 $\kappa$，使得每个准凝聚模
$\mathcal{F}$ 都是它的 $\kappa$-生成准凝聚子模的有向余极限。
\end{lemma}

\begin{proof}
选取仿射开覆盖 $X = \bigcup_{i \in I} U_i$。对每一对
$i, j$，选取仿射开覆盖
$U_i \cap U_j = \bigcup_{k \in I_{ij}} U_{ijk}$。
写成 $U_i = \Spec(A_i)$ 和 $U_{ijk} = \Spec(A_{ijk})$。
令 $\kappa$ 为任意无限基数，并且它 $\geq$ 各集合
$I$、$I_{ij}$ 的基数。

\medskip\noindent
令 $\mathcal{F}$ 为准凝聚层。置 $M_i = \mathcal{F}(U_i)$ 和
$M_{ijk} = \mathcal{F}(U_{ijk})$。注意
$$
M_i \otimes_{A_i} A_{ijk} = M_{ijk} = M_j \otimes_{A_j} A_{ijk}.
$$
见《概形》引理 \ref{schemes-lemma-widetilde-pullback}。
利用选择公理，选取映射
$$
(i, j, k, m) \mapsto S(i, j, k, m)
$$
它对每个 $i, j \in I$、$k \in I_{ij}$ 和 $m \in M_i$，
指定有限子集 $S(i, j, k, m) \subset M_j$，使得
$$
m \otimes 1 = \sum\nolimits_{m' \in S(i, j, k, m)} m' \otimes a_{m'}
$$
在 $M_{ijk}$ 中成立，其中某些 $a_{m'} \in A_{ijk}$。
此外，约定 $S(i, i, k, m) = \{m\}$ 对上面所有
$i, j = i, k, m$ 都成立。固定这样的映射。

\medskip\noindent
给定子集族 $\mathcal{S} = (S_i)_{i \in I}$，其中
$S_i \subset M_i$ 的基数至多为 $\kappa$，置
$\mathcal{S}' = (S'_i)$，其中
$$
S'_j = \bigcup\nolimits_{(i, k, m)\text{ 使得 }m \in S_i}
S(i, j, k, m)
$$
注意 $S_i \subset S'_i$。又注意，$S'_i$ 的基数至多为
$\kappa$，因为它是以一个基数至多为 $\kappa$ 的集合为指标的有限集之并。
置 $\mathcal{S}^{(0)} = \mathcal{S}$、
$\mathcal{S}^{(1)} = \mathcal{S}'$，并归纳地置
$\mathcal{S}^{(n + 1)} = (\mathcal{S}^{(n)})'$。再置
$\mathcal{S}^{(\infty)} = \bigcup_{n \geq 0} \mathcal{S}^{(n)}$。
写 $\mathcal{S}^{(\infty)} = (S^{(\infty)}_i)$，则对任意元素
$m \in S^{(\infty)}_i$，$m$ 在 $M_{ijk}$ 中的像都可写成有限和
$\sum m' \otimes a_{m'}$，其中 $m' \in S_j^{(\infty)}$。
由此可见，若置
$$
N_i = A_i\text{-子模，属于 }M_i\text{ 由下列对象生成： }S^{(\infty)}_i
$$
则有
$$
N_i \otimes_{A_i} A_{ijk} = N_j \otimes_{A_j} A_{ijk}.
$$
把两边视为 $M_{ijk}$ 的子模。因此，存在准凝聚子层
$\mathcal{G} \subset \mathcal{F}$，使得 $\mathcal{G}(U_i) = N_i$。
此外，依构造，层 $\mathcal{G}$ 是 $\kappa$-生成的。

\medskip\noindent
令 $\{\mathcal{G}_t\}_{t \in T}$ 为所有 $\kappa$-生成准凝聚子层组成的集合。
若 $t, t' \in T$，则 $\mathcal{G}_t + \mathcal{G}_{t'}$
也是 $\kappa$-生成准凝聚子层，因为它是映射
$\mathcal{G}_t \oplus \mathcal{G}_{t'} \to \mathcal{F}$ 的像。
因此，该系统（按包含关系排序）是有向的。
上述论证表明，$\mathcal{F}$ 在 $U_i$ 上的每个截面都属于某个
$\mathcal{G}_t$（因为可以从 $\mathcal{S}$ 开始，并使给定截面属于
$S_i$）。所以
$\colim_t \mathcal{G}_t \to \mathcal{F}$ 既单又满，正是所需。
\end{proof}

\begin{proposition}
\label{properties-proposition-coherator}
令 $X$ 为概形。
\begin{enumerate}
\item 范畴 $\QCoh(\mathcal{O}_X)$ 是 Grothendieck 阿贝尔范畴。
因此，$\QCoh(\mathcal{O}_X)$ 具有充分多内射对象和所有极限。
\item 包含函子
$\QCoh(\mathcal{O}_X) \to \textit{Mod}(\mathcal{O}_X)$
具有右伴随\footnote{这个函子有时称为{\it 凝聚化函子}。}
$$
Q : \textit{Mod}(\mathcal{O}_X) \longrightarrow \QCoh(\mathcal{O}_X)
$$
使得对每个准凝聚层 $\mathcal{F}$，伴随映射
$Q(\mathcal{F}) \to \mathcal{F}$ 都是同构。
\end{enumerate}
\end{proposition}

\begin{proof}
(1) 意味着 $\QCoh(\mathcal{O}_X)$ 满足：(a) 具有所有余极限，
(b) 滤过余极限正合，(c) 具有生成对象；见《内射对象》第
\ref{injectives-section-grothendieck-conditions} 节。
由《概形》第 \ref{schemes-section-quasi-coherent} 节，
$\QCoh(\mathcal{O}_X)$ 中的余极限存在，并且与
$\textit{Mod}(\mathcal{O}_X)$ 中的余极限一致。由《模》引理
\ref{modules-lemma-limits-colimits}，滤过余极限正合。
因此 (a) 和 (b) 成立。为构造生成对象 $U$，选取引理
\ref{properties-lemma-colimit-kappa} 中的基数 $\kappa$。选取引理
\ref{properties-lemma-set-of-iso-classes} 中的族
$(\mathcal{F}_t)_{t \in T}$，它由 $\kappa$-生成准凝聚层组成。置
$U = \bigoplus_{t \in T} \mathcal{F}_t$。由于
$\QCoh(\mathcal{O}_X)$ 的每个对象都是 $\kappa$-生成准凝聚模的滤过余极限，
即同构于 $\mathcal{F}_t$ 的对象的滤过余极限，显然 $U$ 是生成对象。
极限和内射对象的断言在任意 Grothendieck 阿贝尔范畴中都成立；见《内射对象》
定理 \ref{injectives-theorem-injective-embedding-grothendieck} 和引理
\ref{injectives-lemma-grothendieck-products}。

\medskip\noindent
证明 (2)。为构造 $Q$，使用下述一般步骤。
给定对象 $\mathcal{F}$，它属于 $\textit{Mod}(\mathcal{O}_X)$，考虑函子
$$
\QCoh(\mathcal{O}_X)^{opp} \longrightarrow \textit{Sets},\quad
\mathcal{G} \longmapsto \Hom_X(\mathcal{G}, \mathcal{F})
$$
该函子把余极限变为极限，因而可表；见《内射对象》引理
\ref{injectives-lemma-grothendieck-brown}。
因此存在准凝聚层 $Q(\mathcal{F})$ 以及函子性同构
$\Hom_X(\mathcal{G}, \mathcal{F}) = \Hom_X(\mathcal{G}, Q(\mathcal{F}))$，
其中 $\mathcal{G}$ 属于 $\QCoh(\mathcal{O}_X)$。由 Yoneda 引理
（《范畴》引理 \ref{categories-lemma-yoneda}），构造
$\mathcal{F} \leadsto Q(\mathcal{F})$ 关于 $\mathcal{F}$ 是函子性的。
依构造，$Q$ 是包含函子的右伴随。
映射 $Q(\mathcal{F}) \to \mathcal{F}$ 是同构，只要 $\mathcal{F}$ 准凝聚；
这是包含函子
$\QCoh(\mathcal{O}_X) \to \textit{Mod}(\mathcal{O}_X)$
为全忠实函子这一事实的形式推论。
\end{proof}





\section{支集包含于闭子集的截面}
\label{properties-section-sections-with-support-in-closed}

\noindent
给定任意拓扑空间 $X$、闭子集 $Z \subset X$ 和阿贝尔层
$\mathcal{F}$，可以取由支集包含于 $Z$ 的截面组成的子层。
若 $X$ 为概形，$Z$ 为闭子概形，而 $\mathcal{F}$ 为准凝聚模，
则还有一种变体：取概形论意义下支集位于 $Z$ 上的截面。
不过，在概形情形中必须谨慎，因为所得 $\mathcal{O}_X$-模未必准凝聚。

\begin{lemma}
\label{properties-lemma-quasi-coherent-finite-type-ideals}
令 $X$ 为拟紧且拟分离概形。
令 $U \subset X$ 为开子概形。下列条件等价：
\begin{enumerate}
\item $U$ 在 $X$ 中逆拟紧；
\item $U$ 拟紧；
\item $U$ 是有限多个仿射开集之并；
\item 存在有限型准凝聚理想层
$\mathcal{I} \subset \mathcal{O}_X$，使得 $X \setminus U = V(\mathcal{I})$
（作为集合）。
\end{enumerate}
\end{lemma}

\begin{proof}
(1)、(2)、(3) 的等价性由引理
\ref{properties-lemma-quasi-separated-quasi-compact-open-retrocompact} 得出。
假设 (1)、(2)、(3) 成立。令 $T = X \setminus U$。由《概形》引理
\ref{schemes-lemma-reduced-closed-subscheme}，存在唯一准凝聚理想层
$\mathcal{J}$，它在 $T$ 上切出约化诱导闭子概形结构。
注意 $\mathcal{J}|_U = \mathcal{O}_U$，而后者是有限型
$\mathcal{O}_U$-模。由引理 \ref{properties-lemma-extend}，存在有限型准凝聚子层
$\mathcal{I} \subset \mathcal{J}$，并且它满足
$\mathcal{I}|_U = \mathcal{J}|_U$。于是
$X \setminus U = V(\mathcal{I})$，从而得到 (4)。反之，若
$\mathcal{I}$ 如 (4)，并且 $W = \Spec(R) \subset X$ 为仿射开集，则
$\mathcal{I}|_W = \widetilde{I}$，其中某个有限生成理想
$I \subset R$；见引理 \ref{properties-lemma-finite-type-module}。
因此 $U \cap W = \Spec(R) \setminus V(I)$ 拟紧；见《代数》引理
\ref{algebra-lemma-qc-open}。于是 $U \subset X$ 由引理
\ref{properties-lemma-retrocompact} 可知是逆拟紧的。
\end{proof}

\begin{lemma}
\label{properties-lemma-sections-annihilated-by-ideal}
令 $X$ 为概形。
令 $\mathcal{I} \subset \mathcal{O}_X$ 为准凝聚理想层。
令 $\mathcal{F}$ 为准凝聚 $\mathcal{O}_X$-模。
考虑 $\mathcal{O}_X$-模层 $\mathcal{F}'$，它对每个开集
$U \subset X$ 赋予
$$
\mathcal{F}'(U)
=
\{s \in \mathcal{F}(U) \mid
\mathcal{I}s = 0\}
$$
假设 $\mathcal{I}$ 有限型。则
\begin{enumerate}
\item $\mathcal{F}'$ 是准凝聚 $\mathcal{O}_X$-模层；
\item 对任意仿射开集 $U \subset X$，有
$\mathcal{F}'(U) = \{s \in \mathcal{F}(U) \mid \mathcal{I}(U)s = 0\}$；
\item $\mathcal{F}'_x = \{s \in \mathcal{F}_x \mid \mathcal{I}_x s = 0\}$。
\end{enumerate}
\end{lemma}

\begin{proof}
显然，定义 $\mathcal{F}'$ 的规则给出 $\mathcal{F}$ 的子层
（层条件容易检验）。因此可以在 $X$ 上局部地验证其他断言。
换言之，可以假设
$X = \Spec(A)$、$\mathcal{F} = \widetilde{M}$ 且
$\mathcal{I} = \widetilde{I}$。在此情形下，显然
$\mathcal{F}'(U) = \{x \in M \mid Ix = 0\} =: M'$，因为
$\widetilde{I}$ 由它的整体截面 $I$ 生成；这证明了 (2)。
为证明 $\mathcal{F}'$ 准凝聚，只须说明对每个 $f \in A$ 都有
$\{x \in M_f \mid I_f x = 0\} = (M')_f$。
写 $I = (g_1, \ldots, g_t)$；这是可行的，因为 $\mathcal{I}$ 有限型；
见引理 \ref{properties-lemma-finite-type-module}。
若 $x = y/f^n$ 且 $I_fx = 0$，则这意味着对每个 $i$ 都存在
$m \geq 0$，使得 $f^mg_ix = 0$。
可以选取一个 $m$，它对所有 $i$ 都适用（这里正用到了 $I$ 有限生成）。
于是 $f^mx \in M'$，并且
$x/f^n = f^mx/f^{n + m}$ 在 $(M')_f$ 中成立，正是所需。
(3) 的证明类似，故略去。
\end{proof}

\begin{definition}
\label{properties-definition-subsheaf-sections-annihilated-by-ideal}
令 $X$ 为概形。
令 $\mathcal{I} \subset \mathcal{O}_X$ 为有限型准凝聚理想层。
令 $\mathcal{F}$ 为准凝聚 $\mathcal{O}_X$-模。
把上面引理 \ref{properties-lemma-sections-annihilated-by-ideal} 定义的子层
$\mathcal{F}' \subset \mathcal{F}$ 称为{\it 被 $\mathcal{I}$ 零化的截面子层}。
\end{definition}

\begin{lemma}
\label{properties-lemma-push-sections-annihilated-by-ideal}
令 $f : X \to Y$ 为概形的拟紧拟分离态射。
令 $\mathcal{I} \subset \mathcal{O}_Y$ 为有限型准凝聚理想层。
令 $\mathcal{F}$ 为准凝聚 $\mathcal{O}_X$-模。令
$\mathcal{F}' \subset \mathcal{F}$ 为被
$f^{-1}\mathcal{I}\mathcal{O}_X$ 零化的截面子层。
则 $f_*\mathcal{F}' \subset f_*\mathcal{F}$ 是被
$\mathcal{I}$ 零化的截面子层。
\end{lemma}

\begin{proof}
略。（提示：$f$ 拟紧且拟分离这一假设蕴含
$f_*\mathcal{F}$ 准凝聚，因而引理
\ref{properties-lemma-sections-annihilated-by-ideal} 可应用于
$\mathcal{I}$ 和 $f_*\mathcal{F}$。）
\end{proof}

\noindent
对于拓扑空间上的阿贝尔层，我们已在《模》评注
\ref{modules-remark-sections-support-in-closed} 中讨论过
支集包含于闭子集的截面子层。对于准凝聚模，这个子模未必准凝聚；
但若该闭子集的补集逆拟紧，则它准凝聚。

\begin{lemma}
\label{properties-lemma-sections-supported-on-closed-subset}
令 $X$ 为概形。令 $Z \subset X$ 为闭子集。
令 $\mathcal{F}$ 为准凝聚 $\mathcal{O}_X$-模。
考虑 $\mathcal{O}_X$-模层 $\mathcal{F}'$，它对每个开集
$U \subset X$ 赋予
$$
\mathcal{F}'(U)
=
\{s \in \mathcal{F}(U) \mid
\text{下列对象的支集： }s\text{ 包含于 }Z \cap U\}
$$
若 $X \setminus Z$ 是 $X$ 的逆拟紧开集，则
\begin{enumerate}
\item 对仿射开集 $U \subset X$，存在有限生成理想
$I \subset \mathcal{O}_X(U)$，使得 $Z \cap U = V(I)$；
\item 对 (1) 中的 $U$ 和 $I$，有
$\mathcal{F}'(U) = \{x \in \mathcal{F}(U) \mid
I^nx = 0 \text{ 对某个 } n\}$；
\item $\mathcal{F}'$ 是准凝聚 $\mathcal{O}_X$-模层。
\end{enumerate}
\end{lemma}

\begin{proof}
(1) 是《代数》引理 \ref{algebra-lemma-qc-open}。
令 $U = \Spec(A)$，并令 $I$ 如 (1)。
则 $\mathcal{F}|_U$ 是与某个 $A$-模 $M$ 相伴的准凝聚层。我们有
$$
\mathcal{F}'(U) = \{x \in M \mid x = 0\text{ 于 }M_\mathfrak p
\text{ 对所有 }\mathfrak p \not \in Z\}.
$$
这是由《模》定义 \ref{modules-definition-support} 得出的。因此
$x \in \mathcal{F}'(U)$ 当且仅当
$V(\text{Ann}(x)) \subset V(I)$；见《代数》引理
\ref{algebra-lemma-support-element}。由于 $I$ 有限生成，这等价于
$I^n x = 0$ 对某个 $n$ 成立。这证明了 (2)。

\medskip\noindent
证明 (3)。注意，对给定开集 $U \subset X$，有正合序列
$$
0 \to \mathcal{F}'(U) \to \mathcal{F}(U) \to \mathcal{F}(U \setminus Z)
$$
若以 $j : X \setminus Z \to X$ 表示包含态射，则注意
$\mathcal{F}(U \setminus Z)$ 是模
$j_*(\mathcal{F}|_{X \setminus Z})$ 在 $U$ 上的截面。因此有正合序列
$$
0 \to \mathcal{F}' \to \mathcal{F} \to j_*(\mathcal{F}|_{X \setminus Z})
$$
限制 $\mathcal{F}|_{X \setminus Z}$ 是准凝聚的。
所以 $j_*(\mathcal{F}|_{X \setminus Z})$ 准凝聚；这是由《概形》引理
\ref{schemes-lemma-push-forward-quasi-coherent}
及关于 $j$ 拟紧的假设（任意开浸入都分离）得出的。
于是 $\mathcal{F}'$ 作为准凝聚模间映射的核是准凝聚的；见《概形》第
\ref{schemes-section-quasi-coherent} 节。
\end{proof}

\begin{definition}
\label{properties-definition-subsheaf-sections-supported-on-closed}
令 $X$ 为概形。
令 $T \subset X$ 为闭子集，并使其补集在 $X$ 中逆拟紧。
令 $\mathcal{F}$ 为准凝聚 $\mathcal{O}_X$-模。
把引理 \ref{properties-lemma-sections-supported-on-closed-subset} 定义的准凝聚子层
$\mathcal{F}' \subset \mathcal{F}$ 称为{\it 支集位于 $T$ 上的截面子层}。
\end{definition}

\begin{lemma}
\label{properties-lemma-push-sections-supported-on-closed-subset}
令 $f : X \to Y$ 为概形的拟紧拟分离态射。
令 $Z \subset Y$ 为闭子集，使得 $Y \setminus Z$ 在 $Y$ 中逆拟紧。
令 $\mathcal{F}$ 为准凝聚 $\mathcal{O}_X$-模。令
$\mathcal{F}' \subset \mathcal{F}$ 为支集位于 $f^{-1}Z$ 上的截面子层。
则 $f_*\mathcal{F}' \subset f_*\mathcal{F}$ 是支集位于
$Z$ 上的截面子层。
\end{lemma}

\begin{proof}
略。（提示：先证明 $X \setminus f^{-1}Z$ 在 $X$ 中逆拟紧，因为
$Y \setminus Z$ 在 $Y$ 中逆拟紧。因此引理
\ref{properties-lemma-sections-supported-on-closed-subset} 可应用于
$f^{-1}Z$ 和 $\mathcal{F}$。由于 $f$ 拟紧且拟分离，可知
$f_*\mathcal{F}$ 准凝聚。因此引理
\ref{properties-lemma-sections-supported-on-closed-subset} 可应用于
$Z$ 和 $f_*\mathcal{F}$。最后，直接对照这些层。）
\end{proof}





\section{准凝聚层的截面}
\label{properties-section-sections-quasi-coherent}

\noindent
下面计算仿射谱的拟紧开集上准凝聚层的截面。

\begin{lemma}
\label{properties-lemma-sections-over-quasi-compact-open-in-affine}
令 $A$ 为环。
令 $I \subset A$ 为有限生成理想。
令 $M$ 为 $A$-模。
则存在典范映射
$$
\colim_n \Hom_A(I^n, M)
\longrightarrow
\Gamma(\Spec(A) \setminus V(I), \widetilde{M}).
$$
这个映射总是单射。
若对所有 $x \in M$ 都有 $Ix = 0 \Rightarrow x = 0$，
则该映射是同构。一般地，置
$M_n = \{x \in M \mid I^nx = 0\}$，则存在同构
$$
\colim_n \Hom_A(I^n, M/M_n)
\longrightarrow
\Gamma(\Spec(A) \setminus V(I), \widetilde{M}).
$$
\end{lemma}

\begin{proof}
由于 $I^{n + 1} \subset I^n$ 且 $M_n \subset M_{n + 1}$，
可以通过与这些映射复合，得到典范 $A$-模映射
$$
\Hom_A(I^n, M)
\longrightarrow
\Hom_A(I^{n + 1}, M)
$$
以及
$$
\Hom_A(I^n, M/M_n)
\longrightarrow
\Hom_A(I^{n + 1}, M/M_{n + 1})
$$
并把它们用作这些系统中的过渡映射。给定
$A$-模映射 $\varphi : I^n \to M$，得到层映射
$\widetilde{\varphi} : \widetilde{I^n} \to \widetilde{M}$，
可以把它限制到开集 $\Spec(A) \setminus V(I)$。
由于 $\widetilde{I^n}$ 限制到该开集后给出结构层，便得到
$\Gamma(\Spec(A) \setminus V(I), \widetilde{M})$ 的一个元素。
略去它与系统 $\Hom_A(I^n, M)$ 中过渡映射相容的验证。
这给出第一条箭头。为得到第二条箭头，注意
$\widetilde{M}$ 和 $\widetilde{M/M_n}$ 在开集
$\Spec(A) \setminus V(I)$ 上一致，因为层 $\widetilde{M_n}$
显然支集位于 $V(I)$ 上。因此可以使用与前面相同的机制。

\medskip\noindent
下面用代数语言具体说明如何定义这条箭头。
设 $I = (f_1, \ldots, f_t)$。则
$\Spec(A) \setminus V(I) = \bigcup_{i = 1, \ldots, t} D(f_i)$。
因此
$$
0 \to
\Gamma(\Spec(A) \setminus V(I), \widetilde{M}) \to
\bigoplus\nolimits_i M_{f_i} \to
\bigoplus\nolimits_{i, j} M_{f_if_j}
$$
正合。假设 $\varphi : I^n \to M$ 是 $A$-模映射。
考虑元素向量 $\varphi(f_i^n)/f_i^n \in M_{f_i}$。
容易看出，这个向量在上面正合序列的第二个直和中映为零。
因而得到 $\Gamma(\Spec(A) \setminus V(I), \widetilde{M})$ 的一个元素。
略去此描述与上面描述一致的验证。

\medskip\noindent
利用这一描述来证明第一条箭头是单射。
事实上，若 $\varphi$ 映为零，则对每个 $i$，元素
$\varphi(f_i^n)/f_i^n$ 在 $M_{f_i}$ 中为零。换言之，对每个
$i$，都有 $f_i^m\varphi(f_i^n) = 0$ 对某个 $m \geq 0$ 成立。
可以选取一个 $m$，它对所有 $i$ 都适用。
于是 $\varphi(f_i^{n + m}) = 0$ 对所有 $i$ 都成立。
容易看出，这意味着
$\varphi|_{I^{t(n + m - 1) + 1}} = 0$；换言之，$\varphi$
在余极限的第 $t(n + m - 1) + 1$ 项中映为零。
单射性由此得出。

\medskip\noindent
注意，每个 $M_n = 0$，只要 $Ix = 0 \Rightarrow x = 0$
对 $x \in M$ 成立。因此，为完成本引理的证明，只须说明第二条箭头是同构。

\medskip\noindent
尝试构造本引理第二个映射的逆。令
$s \in \Gamma(\Spec(A) \setminus V(I), \widetilde{M})$。
它对应于上面正合序列第一个直和中的向量
$x_i/f_i^n$，其中 $x_i \in M$。
因此，对每个 $i, j$，都存在 $m \geq 0$，使得
$f_i^m f_j^m (f_j^n x_i - f_i^n x_j) = 0$ 在 $M$ 中成立。
可以选取一个 $m$，它对所有有序对 $i, j$ 都适用。
把 $x_i$ 替换为 $f_i^mx_i$，并把 $n$ 替换为 $n + m$ 后，
可知得到 $f_j^nx_i = f_i^nx_j$，它在 $M$ 中对所有 $i, j$ 都成立。
引入
$$
K_n = \{x \in M \mid f_1^nx = \ldots = f_t^nx = 0\}
$$
我们断言存在 $A$-模映射
$$
\varphi :
I^{t(n - 1) + 1}
\longrightarrow
M/K_n
$$
它把单项式
$f_1^{e_1} \ldots f_t^{e_t}$（其中 $\sum e_i = t(n - 1) + 1$）
映到模 $K_n$ 的类，该类由表达式
$f_1^{e_1} \ldots f_i^{e_i - n} \ldots f_t^{e_t}x_i$
给出，其中选取 $i$ 使得 $e_i \geq n$
（注意，至少存在一个这样的 $i$）。
为说明这一点，假设
$$
\sum\nolimits_{E = (e_1, \ldots, e_t), |E| = t(n - 1) + 1}
a_E f_1^{e_1} \ldots f_t^{e_t} = 0
$$
是这些单项式间的关系，其系数 $a_E$ 属于 $A$。
则会把它映到
$$
z =
\sum\nolimits_{E = (e_1, \ldots, e_t), |E| = t(n - 1) + 1}
a_E f_1^{e_1} \ldots f_{i(E)}^{e_{i(E)} - n} \ldots f_t^{e_t}x_{i(E)}
$$
其中对每个多重指标 $E$，选定一个特定的 $i(E)$，使得
$e_{i(E)} \geq n$。
注意，若把它乘以 $f_j^n$，其中 $j$ 任意，便得到零，因为利用上面的关系
$f_j^nx_i = f_i^nx_j$ 可得
\begin{align*}
f_j^nz & = \sum\nolimits_{E = (e_1, \ldots, e_t), |E| = t(n - 1) + 1}
a_E f_1^{e_1} \ldots f_j^{e_j + n}
\ldots f_{i(E)}^{e_{i(E)} - n} \ldots f_t^{e_t}x_{i(E)} \\
& =
\sum\nolimits_{E = (e_1, \ldots, e_t), |E| = t(n - 1) + 1}
a_E f_1^{e_1} \ldots f_t^{e_t}x_j
= 0.
\end{align*}
因此 $z \in K_n$，从而每个关系在 $M/K_n$ 中都映为零。
这证明了断言。

\medskip\noindent
注意 $K_n \subset M_{t(n - 1) + 1}$。因此映射 $\varphi$
特别给出 $A$-模映射
$I^{t(n - 1) + 1} \to M/M_{t(n - 1) + 1}$。
这证明了本引理的第二条箭头是满射。略去单射性的证明。
\end{proof}

\begin{example}
\label{properties-example-does-not-work-in-general}
下面给出两个例子，说明引理
\ref{properties-lemma-sections-over-quasi-compact-open-in-affine}
的第一条陈列映射并非同构。

\medskip\noindent
令 $k$ 为域。考虑环
$$
A = k[x, y, z_1, z_2, \ldots]/(x^nz_n).
$$
置 $I = (x)$ 并令 $M = A$。则元素 $y/x$ 定义了 $\Spec(A)$
的结构层在 $D(x) = \Spec(A) \setminus V(I)$ 上的一个截面。
我们断言，$y/x$ 不在典范映射
$\colim \Hom_A(I^n, A) \to A_x = \mathcal{O}(D(x))$ 的像中。
事实上，若在像中，则它会来自同态
$\varphi : I^n \to A$，其中某个 $n$。置 $a = \varphi(x^n)$。
于是会有 $x^m(xa - x^ny) = 0$，其中某个 $m > 0$。
这意味着 $x^{m + 1}a = x^{m + n}y$，继而意味着
$\varphi(x^{n + m + 1}) = x^{m + n}y$。这导致矛盾，因为它将蕴含
$$
0 = \varphi(0) = \varphi(z_{n + m + 1} x^{n + m + 1}) =
x^{m + n}y z_{n + m + 1}
$$
而这在环 $A$ 中不成立。

\medskip\noindent
令 $k$ 为域。考虑环
$$
A = k[f, g, x, y, \{a_n, b_n\}_{n \geq 1}]/
(fy - gx, \{a_nf^n + b_ng^n\}_{n \geq 1}).
$$
置 $I = (f, g)$ 并令 $M = A$。则 $x/f \in A_f$ 和
$y/g \in A_g$ 在 $A_{fg}$ 中映到同一个元素。因此，它们定义了截面
$s$；这是 $\Spec(A)$ 的结构层在
$D(f) \cup D(g) = \Spec(A) \setminus V(I)$ 上的截面。
但是，不存在 $n \geq 0$，使得 $s$ 来自如引理
\ref{properties-lemma-sections-over-quasi-compact-open-in-affine}
第一条陈列箭头之源中的 $A$-模映射
$\varphi : I^n \to A$。事实上，给定这样的模映射，置
$x_n = \varphi(f^n)$ 和 $y_n = \varphi(g^n)$。
则有 $f^mx_n = f^{n + m - 1}x$ 以及
$g^my_n = g^{n + m - 1}y$，其中某个 $m \geq 0$
（见本引理的证明）。但此时会有
$0 = \varphi(0) =
\varphi(a_{n + m}f^{n + m} + b_{n + m}g^{n + m}) =
a_{n + m}f^{n + m - 1}x + b_{n + m}g^{n + m - 1}y$，
而这在环 $A$ 中不成立。
\end{example}

\noindent
在 Noether 情形中，我们将改进下述引理；见《概形上同调》引理
\ref{coherent-lemma-homs-over-open}。

\begin{lemma}
\label{properties-lemma-sections-over-quasi-compact-open}
令 $X$ 为拟紧概形。
令 $\mathcal{I} \subset \mathcal{O}_X$ 为有限型准凝聚理想层。
令 $Z \subset X$ 为由 $\mathcal{I}$ 定义的闭子概形，并置
$U = X \setminus Z$。
令 $\mathcal{F}$ 为准凝聚 $\mathcal{O}_X$-模。
典范映射
$$
\colim_n \Hom_{\mathcal{O}_X}(\mathcal{I}^n,
\mathcal{F})
\longrightarrow
\Gamma(U, \mathcal{F})
$$
是单射。进一步假设 $X$ 拟分离。
令 $\mathcal{F}_n \subset \mathcal{F}$ 为被
$\mathcal{I}^n$ 零化的截面子层。典范映射
$$
\colim_n \Hom_{\mathcal{O}_X}(\mathcal{I}^n,
\mathcal{F}/\mathcal{F}_n)
\longrightarrow
\Gamma(U, \mathcal{F})
$$
是同构。
\end{lemma}

\begin{proof}
令 $\Spec(A) = W \subset X$ 为仿射开集。写
$\mathcal{F}|_W = \widetilde{M}$，其中某个 $A$-模 $M$，并写
$\mathcal{I}|_W = \widetilde{I}$，其中某个有限型理想
$I \subset A$。把本引理的第一条陈列映射限制到 $W$，得到引理
\ref{properties-lemma-sections-over-quasi-compact-open-in-affine}
的第一条陈列映射。由于可以用有限多个仿射开集覆盖 $X$，
这证明了本引理的第一条陈列映射是单射。

\medskip\noindent
有 $\mathcal{F}_n|_W = \widetilde{M_n}$，其中
$M_n \subset M$ 如引理
\ref{properties-lemma-sections-over-quasi-compact-open-in-affine}
中定义（略去细节）。该引理保证存在双射
$$
\colim_n  \Hom_{\mathcal{O}_W}(
\mathcal{I}^n|_W, (\mathcal{F}/\mathcal{F}_n)|_W)
\longrightarrow
\Gamma(U \cap W, \mathcal{F})
$$
这对任意这样的仿射开集 $W$ 成立。

\medskip\noindent
为证明本引理的第二条陈列箭头是双射，选取有限仿射开覆盖
$X = \bigcup_{j = 1, \ldots, m} W_j$。单射性立即由上文和覆盖的有限性得出。
若 $X$ 拟分离，则对每一对 $j, j'$，选取有限仿射开覆盖
$$
W_j \cap W_{j'} = \bigcup\nolimits_{k = 1, \ldots, m_{jj'}} W_{jj'k}.
$$
令 $s \in \Gamma(U, \mathcal{F})$。如上所见，对每个 $j$，存在
$n_j$ 以及映射
$\varphi_j : \mathcal{I}^{n_j}|_{W_j} \to
(\mathcal{F}/\mathcal{F}_{n_j})|_{W_j}$，
它对应于 $s|_{U \cap W_j}$。
同理，对每个三元组 $(j, j', k)$，存在整数 $n_{jj'k}$，
使得 $\varphi_j$ 与 $\varphi_{j'}$ 作为映射
$\mathcal{I}^{n_{jj'k}} \to \mathcal{F}/\mathcal{F}_{n_{jj'k}}$
限制到 $W_{jj'k}$ 上时一致。令
$n = \max\{n_j, n_{jj'k}\}$，则 $\varphi_j$ 可粘合为映射
$\mathcal{I}^n \to \mathcal{F}/\mathcal{F}_n$，该映射定义在 $X$ 上。
这证明了该映射的满射性。
\end{proof}





\section{丰沛可逆层}
\label{properties-section-ample}

\noindent
回顾《模》引理 \ref{modules-lemma-s-open}：给定可逆层
$\mathcal{L}$，它在局部环层空间 $X$ 上；并给定整体截面
$s$，它属于 $\mathcal{L}$，则集合
$X_s = \{x \in X \mid s \not \in \mathfrak m_x\mathcal{L}_x\}$
是开的。一般地，有
$X_s \cap X_{s'} = X_{ss'}$，其中 $ss'$ 表示截面
$s \otimes s' \in \Gamma(X, \mathcal{L} \otimes \mathcal{L}')$。

\begin{definition}
\label{properties-definition-ample}
\begin{reference}
\cite[II Definition 4.5.3]{EGA}
\end{reference}
令 $X$ 为概形。
令 $\mathcal{L}$ 为可逆 $\mathcal{O}_X$-模。称 $\mathcal{L}$ {\it 丰沛}，若
\begin{enumerate}
\item $X$ 拟紧；
\item 对每个 $x \in X$，都存在 $n \geq 1$ 和
$s \in \Gamma(X, \mathcal{L}^{\otimes n})$，使得
$x \in X_s$ 且 $X_s$ 仿射，
\end{enumerate}
\end{definition}

\begin{lemma}
\label{properties-lemma-ample-power-ample}
\begin{reference}
\cite[II Proposition 4.5.6(i)]{EGA}
\end{reference}
令 $X$ 为概形。令 $\mathcal{L}$ 为可逆 $\mathcal{O}_X$-模。
令 $n \geq 1$。则 $\mathcal{L}$ 丰沛，当且仅当
$\mathcal{L}^{\otimes n}$ 丰沛。
\end{lemma}

\begin{proof}
这由 $X_{s^n} = X_s$ 这一事实得出。
\end{proof}

\begin{lemma}
\label{properties-lemma-ample-on-closed}
令 $X$ 为概形。
令 $\mathcal{L}$ 为丰沛可逆 $\mathcal{O}_X$-模。
对任意闭子概形 $Z \subset X$，$\mathcal{L}$ 到 $Z$ 的限制丰沛。
\end{lemma}

\begin{proof}
这是显然的，因为拟紧空间的闭子集拟紧，而仿射概形的闭子概形仿射
（见《概形》引理 \ref{schemes-lemma-closed-immersion-affine-case}）。
\end{proof}

\begin{lemma}
\label{properties-lemma-affine-cap-s-open}
令 $X$ 为概形。令 $\mathcal{L}$ 为可逆 $\mathcal{O}_X$-模。
令 $s \in \Gamma(X, \mathcal{L})$。对任意仿射集
$U \subset X$，交 $U \cap X_s$ 是仿射的。
\end{lemma}

\begin{proof}
这转化为下述代数问题。
令 $R$ 为环。令 $N$ 为可逆 $R$-模
（即秩为 1 的局部自由模）。令 $s \in N$ 为元素。
则 $V = \{\mathfrak p \mid s \not \in \mathfrak p N\}$ 是
$\Spec(R)$ 的仿射开子集。

\medskip\noindent
令 $A = \bigoplus_{n \geq 0} A_n$ 为 $N$ 的对称代数
（它是交换的），并把 $s$ 看作 $A_1$ 的元素。
置 $B = A/(s - 1)A$。这是一个 $R$-代数，其构造与任意基变换
$R \to R'$ 交换。因此，
$B' = B \otimes_R R'$ 是零环，只要 $s$ 在
$N' = N \otimes_R R'$ 中映为零。
由此，若 $x \in \Spec(R) \setminus V$，则
$B \otimes_R \kappa(x) = 0$。
我们断定 $\Spec(B) \to \Spec(R)$ 通过 $V$ 分解，因为在
$x \not \in V$ 上的纤维为空。另一方面，若
$\Spec(R') \subset V$ 为仿射开集，则 $s$ 映到 $N'$ 的一个基元素，
并且可知 $B' = R'[s]/(s - 1) \cong R'$。
因此 $\Spec(B) \to V$ 是同构，而 $V$ 的确仿射。
\end{proof}

\begin{lemma}
\label{properties-lemma-ample-tensor-globally-generated}
\begin{reference}
\cite[II Proposition 4.5.6(ii)]{EGA}
\end{reference}
令 $X$ 为概形。令 $\mathcal{L}$ 和 $\mathcal{M}$
为可逆 $\mathcal{O}_X$-模。若
\begin{enumerate}
\item $\mathcal{L}$ 丰沛；
\item 开集 $X_t$，其中
$t \in \Gamma(X, \mathcal{M}^{\otimes m})$ 且 $m > 0$，覆盖 $X$，
\end{enumerate}
则 $\mathcal{L} \otimes \mathcal{M}$ 丰沛。
\end{lemma}

\begin{proof}
检验定义 \ref{properties-definition-ample} 的条件。
由于 $\mathcal{L}$ 丰沛，可知 $X$ 拟紧。
令 $x \in X$。选取 $n \geq 1$、$m \geq 1$、
$s \in \Gamma(X, \mathcal{L}^{\otimes n})$ 和
$t \in \Gamma(X, \mathcal{M}^{\otimes m})$，使得
$x \in X_s$、$x \in X_t$ 且 $X_s$ 仿射。
则
$s^mt^n \in \Gamma(X, (\mathcal{L} \otimes \mathcal{M})^{\otimes nm})$、
$x \in X_{s^mt^n}$，并且由引理
\ref{properties-lemma-affine-cap-s-open}，$X_{s^mt^n}$ 仿射。
\end{proof}

\begin{lemma}
\label{properties-lemma-affine-s-opens}
令 $X$ 为概形。令 $\mathcal{L}$ 为可逆 $\mathcal{O}_X$-模。
假设开集 $X_s$，其中
$s \in \Gamma(X, \mathcal{L}^{\otimes n})$ 且 $n \geq 1$，
构成 $X$ 的拓扑基。
则在这些开集之中，仿射的开集 $X_s$ 构成 $X$ 的拓扑基。
\end{lemma}

\begin{proof}
令 $x \in X$。选取仿射开邻域
$\Spec(R) = U \subset X$，它含有 $x$。由假设，存在
$n \geq 1$ 和 $s \in \Gamma(X, \mathcal{L}^{\otimes n})$，
使得 $X_s \subset U$。由上面的引理
\ref{properties-lemma-affine-cap-s-open}，交 $X_s = U \cap X_s$ 仿射。
由于 $U$ 可任意缩小，结论得证。
\end{proof}

\begin{lemma}
\label{properties-lemma-affine-s-opens-cover-quasi-separated}
令 $X$ 为概形，$\mathcal{L}$ 为可逆 $\mathcal{O}_X$-模。
假设对每个点 $x$（属于 $X$），都存在 $n \geq 1$ 和
$s \in \Gamma(X, \mathcal{L}^{\otimes n})$，使得
$x \in X_s$ 且 $X_s$ 仿射。则 $X$ 分离。
\end{lemma}

\begin{proof}
先证明 $X$ 拟分离。由假设，可以找到 $X$ 的覆盖，该覆盖由形如
$X_s$ 的仿射开集组成。由引理
\ref{properties-lemma-affine-cap-s-open}，任意两个这样的集合之交仿射，
所以《概形》引理
\ref{schemes-lemma-characterize-quasi-separated} 蕴含 $X$ 拟分离。

\medskip\noindent
为证明 $X$ 分离，可以使用赋值判据，即《概形》引理
\ref{schemes-lemma-valuative-criterion-separatedness}。
因此，令 $A$ 为分式域为 $K$ 的赋值环，并考虑两个态射
$f, g : \Spec(A) \to X$，使得两个复合
$\Spec(K) \to \Spec(A) \to X$ 相同。
由于 $A$ 是局部环，存在 $p, q \ge 1$、
$s \in \Gamma(X,
\mathcal{L}^{\otimes p})$ 和
$t \in \Gamma(X,
\mathcal{L}^{\otimes q})$，使得
$X_s$ 与 $X_t$ 仿射，
$f(\Spec A) \subseteq X_s$ 且 $g(\Spec A)
\subseteq X_t$。
现在把 $s$ 替换为 $s^q$，把 $t$ 替换为 $t^p$，
并把 $\mathcal{L}$ 替换为 $\mathcal{L}^{\otimes pq}$。
这是无害的，因为 $X_s = X_{s^q}$ 且 $X_t = X_{t^p}$，
而现在 $s$ 和 $t$ 都是同一个层 $\mathcal{L}$ 的截面。

\medskip\noindent
准凝聚模 $f^*\mathcal{L}$ 对应于 $A$-模 $M$，而
$g^*\mathcal{L}$ 对应于 $A$-模 $N$；这是由仿射概形上准凝聚模的分类
（《概形》引理 \ref{schemes-lemma-quasi-coherent-affine}）得出的。
$A$-模 $M$ 与 $N$ 都是秩为 $1$ 的局部自由模
（引理 \ref{properties-lemma-locally-free-module}），而由于 $A$ 局部，它们自由
（《代数》引理 \ref{algebra-lemma-K0-local}）。因此，可以把
$M$ 和 $N$ 等同为 $A$-子模，它们分别属于 $M \otimes_A K$ 和 $N
\otimes_A K$。等式
$f|_{\Spec(K)} = g|_{\Spec(K)}$ 确定一个同构
$\phi \colon M \otimes_A K \to N
\otimes_A K$。

\medskip\noindent
令 $x \in M$ 和 $y \in N$ 分别为把 $s$ 沿 $f$ 与 $g$ 拉回所得的元素。
它们满足 $\phi(x \otimes 1) = y \otimes 1$。
$f$ 的像包含于 $X_s$，故 $x \not\in \mathfrak{m}_A M$，
即 $x$ 生成 $M$。因此，$\phi$ 确定 $M$ 到 $N$ 中由
$y$ 生成的子模的同构。对称地使用 $t$，$\phi^{-1}$
确定 $N$ 到 $M$ 的一个子模的同构。因此，$\phi$ 限制为
$M$ 与 $N$ 之间的同构。由于 $x$ 生成 $M$，它的像 $y$ 生成
$N$，从而 $y \not\in \mathfrak{m}_A N$。所以
$g(\Spec(A)) \subseteq X_s$。因为 $X_s$ 仿射，由《概形》引理
\ref{schemes-lemma-affine-separated}，它分离，因而得出 $f = g$。
\end{proof}

\begin{lemma}
\label{properties-lemma-ample-separated}
令 $X$ 为概形。若 $X$ 上存在丰沛可逆层，则 $X$ 分离。
\end{lemma}

\begin{proof}
立即由引理 \ref{properties-lemma-affine-s-opens-cover-quasi-separated} 和定义
\ref{properties-definition-ample} 得出。
\end{proof}

\begin{lemma}
\label{properties-lemma-map-into-proj}
令 $X$ 为概形。
令 $\mathcal{L}$ 为可逆 $\mathcal{O}_X$-模。
把 $S = \Gamma_*(X, \mathcal{L})$ 视为分次环。
若 $X$ 的每个点都包含于某个开子概形 $X_s$，其中某个齐次
$s \in S_{+}$，则存在到 $S$ 的齐次谱（见《构造》第
\ref{constructions-section-proj} 节）的典范概形态射
$$
f : X \longrightarrow Y = \text{Proj}(S),
$$
它具有下列性质：
\begin{enumerate}
\item 有 $f^{-1}(D_{+}(s)) = X_s$，对任意齐次
$s \in S_{+}$；
\item 存在与乘法映射相容的 $\mathcal{O}_X$-模映射
$f^*\mathcal{O}_Y(n) \to \mathcal{L}^{\otimes n}$；见《构造》方程
(\ref{constructions-equation-multiply})；
\item 复合
$S_n \to \Gamma(Y, \mathcal{O}_Y(n)) \to \Gamma(X, \mathcal{L}^{\otimes n})$
是恒等映射；
\item 对每个 $x \in X$，都存在整数 $d \geq 1$ 以及开邻域
$U \subset X$，它包含 $x$，使得
$f^*\mathcal{O}_Y(dn)|_U \to \mathcal{L}^{\otimes dn}|_U$
都是同构，对所有 $n \in \mathbf{Z}$。
\end{enumerate}
\end{lemma}

\begin{proof}
以 $\psi : S \to \Gamma_*(X, \mathcal{L})$ 表示恒等映射。
将使用《构造》引理
\ref{constructions-lemma-invertible-map-into-proj} 的三元组
$(U(\psi), r_{\mathcal{L}, \psi}, \theta)$。
由假设，开子概形 $U(\psi)$ 的等于 $X$。因此
$r_{\mathcal{L}, \psi} : U(\psi) \to Y$ 在整个 $X$ 上有定义。
置 $f = r_{\mathcal{L}, \psi}$。
(2) 中的映射是 $\theta$ 的分量。
(3) 由上引引理中的条件 (2) 得出。
(1) 由 (3) 与上引引理中的条件 (1) 结合得出。
(4) 由《构造》引理
\ref{constructions-lemma-invertible-map-into-proj} 的最后一个陈述得出，
因为其中提到的映射 $\alpha$ 是同构。
\end{proof}

\begin{lemma}
\label{properties-lemma-map-into-proj-quasi-compact}
令 $X$ 为概形。令 $\mathcal{L}$ 为可逆 $\mathcal{O}_X$-模。
置 $S = \Gamma_*(X, \mathcal{L})$。
假设 (a) $X$ 的每个点都包含于某个开子概形 $X_s$，其中某个齐次
$s \in S_{+}$，并且 (b) $X$ 拟紧。则上面引理
\ref{properties-lemma-map-into-proj} 的典范概形态射
$f : X \longrightarrow \text{Proj}(S)$ 拟紧且像稠密。
\end{lemma}

\begin{proof}
为证明 $f$ 拟紧，只须说明 $f^{-1}(D_{+}(s))$ 拟紧，对任意齐次
$s \in S_{+}$。把
$X = \bigcup_{i = 1, \ldots, n} X_i$ 写成有限多个仿射开集之并。
由引理 \ref{properties-lemma-affine-cap-s-open}，每个交
$X_s \cap X_i$ 都仿射。因此
$X_s = \bigcup_{i = 1, \ldots, n} X_s \cap X_i$ 拟紧。
假设 $f$ 的像不稠密以导出矛盾。由于开集
$D_+(s)$（其中 $s \in S_+$ 齐次）构成
$\text{Proj}(S)$ 的拓扑基，可以找到这样的 $s$，使得
$D_+(s) \not = \emptyset$ 且 $f(X) \cap D_+(s) = \emptyset$。
由引理 \ref{properties-lemma-map-into-proj}，这意味着 $X_s = \emptyset$。
由引理 \ref{properties-lemma-invert-s-sections}，这意味着某个幂 $s^n$
是 $\mathcal{L}^{\otimes n\deg(s)}$ 的零截面。
继而这意味着 $D_+(s) = \emptyset$，即所需矛盾。
\end{proof}

\begin{lemma}
\label{properties-lemma-ample-immersion-into-proj}
令 $X$ 为概形。令 $\mathcal{L}$ 为可逆 $\mathcal{O}_X$-模。
置 $S = \Gamma_*(X, \mathcal{L})$。
假设 $\mathcal{L}$ 丰沛。则引理 \ref{properties-lemma-map-into-proj} 的典范概形态射
$f : X \longrightarrow \text{Proj}(S)$ 是像稠密的开浸入。
\end{lemma}

\begin{proof}
由引理 \ref{properties-lemma-affine-s-opens-cover-quasi-separated}，可知
$X$ 拟分离。选取有限多个齐次
$s_1, \ldots, s_n \in S_{+}$，使得 $X_{s_i}$ 仿射并且
$X = \bigcup X_{s_i}$。
设 $s_i$ 的次数为 $d_i$。$D_{+}(s_i)$ 在 $f$ 下的逆像是
$X_{s_i}$；见引理 \ref{properties-lemma-map-into-proj}。
由引理 \ref{properties-lemma-invert-s-sections}，环映射
$$
(S^{(d_i)})_{(s_i)} = \Gamma(D_{+}(s_i), \mathcal{O}_{\text{Proj}(S)})
\longrightarrow
\Gamma(X_{s_i}, \mathcal{O}_X)
$$
是同构。因此 $f$ 诱导同构
$X_{s_i} \to D_{+}(s_i)$。从而 $f$ 是 $X$ 到开子概形
$\bigcup_{i = 1, \ldots, n} D_{+}(s_i)$ 上的同构；该开子概形属于
$\text{Proj}(S)$。
像由引理 \ref{properties-lemma-map-into-proj-quasi-compact} 可知稠密。
\end{proof}

\begin{lemma}
\label{properties-lemma-open-in-proj-ample}
令 $X$ 为概形。
令 $S$ 为分次环。假设 $X$ 拟紧，并且存在开浸入
$$
j : X \longrightarrow Y = \text{Proj}(S).
$$
则 $j^*\mathcal{O}_Y(d)$ 是丰沛可逆层，其中某个 $d > 0$。
\end{lemma}

\begin{proof}
这是《构造》引理 \ref{constructions-lemma-ample-on-proj}。
\end{proof}

\begin{proposition}
\label{properties-proposition-characterize-ample}
令 $X$ 为拟紧概形。
令 $\mathcal{L}$ 为 $X$ 上的可逆层。
置 $S = \Gamma_*(X, \mathcal{L})$。
下列条件等价：
\begin{enumerate}
\item
\label{properties-item-ample}
$\mathcal{L}$ 丰沛；
\item
\label{properties-item-immersion}
开集 $X_s$（其中 $s \in S_{+}$ 齐次）覆盖 $X$，且相伴态射
$X \to \text{Proj}(S)$ 是开浸入；
\item
\label{properties-item-s-basis}
开集 $X_s$（其中 $s \in S_{+}$ 齐次）构成 $X$ 的拓扑基；
\item
\label{properties-item-s-affine-basis}
仿射的开集 $X_s$（其中 $s \in S_{+}$ 齐次）构成 $X$ 的拓扑基；
\item
\label{properties-item-qc-gg}
对每个准凝聚层 $\mathcal{F}$（在 $X$ 上），典范映射
$$
\Gamma(X, \mathcal{F} \otimes_{\mathcal{O}_X} \mathcal{L}^{\otimes n})
\otimes_{\mathbf{Z}} \mathcal{L}^{\otimes -n}
\longrightarrow
\mathcal{F}
$$
之像的和（其中 $n \geq 1$）等于 $\mathcal{F}$；
\item
\label{properties-item-qc-i-gg}
与 (\ref{properties-item-qc-gg}) 相同的性质，其中 $\mathcal{F}$
遍历所有准凝聚理想层；
\item
\label{properties-item-c-gg}
$X$ 拟分离，并且对每个有限型准凝聚层
$\mathcal{F}$（在 $X$ 上），都存在整数 $n_0$，使得
$\mathcal{F} \otimes_{\mathcal{O}_X} \mathcal{L}^{\otimes n}$
对所有 $n \geq n_0$ 都全局生成；
\item
\label{properties-item-c-q}
$X$ 拟分离，并且对每个有限型准凝聚层
$\mathcal{F}$（在 $X$ 上），都存在整数 $n > 0$、$k \geq 0$，使得
$\mathcal{F}$ 是 $k$ 个 $\mathcal{L}^{\otimes - n}$ 的直和的商；
\item
\label{properties-item-c-i-q}
与 (\ref{properties-item-c-q}) 相同，其中 $\mathcal{F}$ 遍历
$X$ 上的所有有限型理想层。
\end{enumerate}
\end{proposition}

\begin{proof}
引理 \ref{properties-lemma-ample-immersion-into-proj} 给出
(\ref{properties-item-ample}) $\Rightarrow$ (\ref{properties-item-immersion})。
引理 \ref{properties-lemma-ample-power-ample} 与
\ref{properties-lemma-open-in-proj-ample} 给出蕴含
(\ref{properties-item-ample}) $\Leftarrow$ (\ref{properties-item-immersion})。
蕴含
(\ref{properties-item-immersion}) $\Rightarrow$
(\ref{properties-item-s-affine-basis}) $\Rightarrow$ (\ref{properties-item-s-basis})
由《构造》第 \ref{constructions-section-proj} 节是显然的。
引理 \ref{properties-lemma-affine-s-opens} 给出
(\ref{properties-item-s-basis}) $\Rightarrow$ (\ref{properties-item-ample})。
因此，前四个条件全都等价。

\medskip\noindent
假设等价条件 (1)--(4) 成立。
特别地，注意 $X$ 分离（因为它是分离概形
$\text{Proj}(S)$ 的开子概形）。
令 $\mathcal{F}$ 为 $X$ 上的准凝聚层。
选取齐次 $s \in S_{+}$，使得 $X_s$ 仿射。
我们断言，任意截面 $m \in \Gamma(X_s, \mathcal{F})$
都属于上面 (\ref{properties-item-qc-gg}) 所陈列的某个映射之像。
由于这些仿射集 $X_s$ 覆盖 $X$，这将蕴含
(\ref{properties-item-qc-gg})。
事实上，由引理 \ref{properties-lemma-invert-s-sections}，可以把
$m$ 写成 $m' \otimes s^{-n}$ 的像，其中某个
$n \geq 1$ 和某个
$m' \in \Gamma(X, \mathcal{F} \otimes \mathcal{L}^{\otimes n})$。
这证明了断言。

\medskip\noindent
显然，(\ref{properties-item-qc-gg}) $\Rightarrow$ (\ref{properties-item-qc-i-gg})。
假设 (\ref{properties-item-qc-i-gg})，并证明 $\mathcal{L}$ 丰沛。
选取 $x \in X$。令 $U \subset X$ 为包含 $x$ 的仿射开集。
置 $Z = X \setminus U$。可把 $Z$ 看作约化闭子概形；见《概形》第
\ref{schemes-section-reduced} 节。
令 $\mathcal{I} \subset \mathcal{O}_X$ 为与闭子概形 $Z$
对应的准凝聚理想层。由假设 (\ref{properties-item-qc-i-gg})，存在
$n \geq 1$ 以及截面
$s \in \Gamma(X, \mathcal{I} \otimes \mathcal{L}^{\otimes n})$，
使得 $s$ 在 $x$ 处不消失（更准确地说，
$s \not \in \mathfrak m_x \mathcal{I}_x \otimes \mathcal{L}_x^{\otimes n}$）。
可把 $s$ 看作 $\mathcal{L}^{\otimes n}$ 的截面。
由于它显然沿 $Z$ 消失，可知 $X_s \subset U$。
因此 $X_s$ 仿射；见引理 \ref{properties-lemma-affine-cap-s-open}。
这证明了 $\mathcal{L}$ 丰沛。至此已经证明 (1)--(6) 等价。

\medskip\noindent
假设等价条件 (1)--(6) 成立。下文将不再另行说明地使用如下事实：
两个全局生成模层的张量积仍全局生成（见《模》引理
\ref{modules-lemma-tensor-product-globally-generated}）。
由 (1)，可找到元素 $s_i \in S_{d_i}$，其中 $d_i \geq 1$，使得
$X = \bigcup_{i = 1, \ldots, n} X_{s_i}$。
置 $d = d_1\ldots d_n$。由此，$\mathcal{L}^{\otimes d}$ 由
$$
s_1^{d/d_1}, \ldots, s_n^{d/d_n}.
$$
全局生成。这意味着，若 $\mathcal{L}^{\otimes j}$ 全局生成，
则 $\mathcal{L}^{\otimes j + dn}$ 也全局生成，对所有 $n \geq 0$。
固定 $j \in \{0, \ldots, d - 1\}$。对任意点 $x \in X$，存在
$n \geq 1$ 以及整体截面 $s$，它属于 $\mathcal{L}^{j + dn}$，
它在 $x$ 处不消失；这是把 (\ref{properties-item-qc-gg}) 应用于
$\mathcal{F} = \mathcal{L}^{\otimes j}$ 和丰沛可逆层
$\mathcal{L}^{\otimes d}$ 得出的。由于 $X$ 拟紧，可以找到有限列整数
$n_i$ 以及整体截面 $s_i$，它们属于
$\mathcal{L}^{\otimes j + dn_i}$，并且在 $X$ 的每一点并非全都消失。
由于 $\mathcal{L}^{\otimes d}$ 全局生成，这意味着
$\mathcal{L}^{\otimes j + dn}$ 全局生成，其中
$n = \max\{n_i\}$。因为对模 $d$ 的每个同余类都证明了这一点，
所以存在 $n_0 = n_0(\mathcal{L})$，使得
$\mathcal{L}^{\otimes n}$ 对所有 $n \geq n_0$ 都全局生成。
至此可知，若 $\mathcal{F}$ 全局生成，则
$\mathcal{F} \otimes \mathcal{L}^{\otimes n}$ 对所有
$n \geq n_0$ 也全局生成。

\medskip\noindent
继续假设等价条件 (1)--(6) 成立。
令 $\mathcal{F}$ 为有限型准凝聚 $\mathcal{O}_X$-模层。
以 $\mathcal{F}_n \subset \mathcal{F}$ 表示
(\ref{properties-item-qc-gg}) 中典范映射的像。依构造，
$\mathcal{F}_n \otimes \mathcal{L}^{\otimes n}$ 全局生成。
由 (\ref{properties-item-qc-gg})，$\mathcal{F}$ 是子层
$\mathcal{F}_n$（$n \geq 1$）之和。由《模》引理
\ref{modules-lemma-finite-type-quasi-compact-colimit}，可知
$\mathcal{F} = \sum_{n = 1, \ldots, N} \mathcal{F}_n$，其中某个
$N \geq 1$。因此，
$\mathcal{F} \otimes \mathcal{L}^{\otimes n}$ 就全局生成，只要
$n \geq N + n_0(\mathcal{L})$，其中 $n_0(\mathcal{L})$
如上。由此断定 (1)--(6) 蕴含 (\ref{properties-item-c-gg})。

\medskip\noindent
假设 (\ref{properties-item-c-gg})。令 $\mathcal{F}$ 为有限型准凝聚
$\mathcal{O}_X$-模层。由 (\ref{properties-item-c-gg})，存在整数
$n \geq 1$，使得典范映射
$$
\Gamma(X, \mathcal{F} \otimes_{\mathcal{O}_X} \mathcal{L}^{\otimes n})
\otimes_{\mathbf{Z}} \mathcal{L}^{\otimes -n}
\longrightarrow
\mathcal{F}
$$
满射。令 $I$ 为
$\Gamma(X, \mathcal{F} \otimes_{\mathcal{O}_X} \mathcal{L}^{\otimes n})$
的有限子集组成的集合，并按包含关系偏序。则 $I$ 是有向偏序集。
对 $i = \{s_1, \ldots, s_{r(i)}\}$，令
$\mathcal{F}_i \subset \mathcal{F}$ 为映射
$$
\bigoplus\nolimits_{j = 1, \ldots, r(i)} \mathcal{L}^{\otimes -n}
\longrightarrow
\mathcal{F}
$$
的像；在该映射中，乘以 $s_j$ 这一作用施加在第 $j$ 个因子上。
上面的满射性蕴含
$\mathcal{F} = \colim_{i \in I} \mathcal{F}_i$。
因此可应用《模》引理
\ref{modules-lemma-finite-type-quasi-compact-colimit}，并得出
$\mathcal{F} = \mathcal{F}_i$，其中某个 $i$。
所以已经证明 (\ref{properties-item-c-q})。换言之，
(\ref{properties-item-c-gg}) $\Rightarrow$ (\ref{properties-item-c-q})。

\medskip\noindent
蕴含 (\ref{properties-item-c-q}) $\Rightarrow$ (\ref{properties-item-c-i-q}) 是平凡的。

\medskip\noindent
最后，假设 (\ref{properties-item-c-i-q})。
令 $\mathcal{I} \subset \mathcal{O}_X$ 为准凝聚理想层。
由引理 \ref{properties-lemma-quasi-coherent-colimit-finite-type}
（这里使用了 $X$ 拟分离这一条件），可知
$\mathcal{I} = \colim_\alpha I_\alpha$，其中每个
$I_\alpha$ 都是有限型准凝聚层。由于依假设，每个 $I_\alpha$
都是 $\mathcal{L}$ 的负张量幂的商，所以对 $\mathcal{I}$
也得到同样的结论（当然，不再有幂的有限性或有界性）。
因此断定 (\ref{properties-item-c-i-q}) 蕴含 (\ref{properties-item-qc-i-gg})。
命题证毕。
\end{proof}

\begin{lemma}
\label{properties-lemma-ample-on-locally-closed}
设 $X$ 为概形。设 $\mathcal{L}$ 为丰沛可逆
$\mathcal{O}_X$-模。设 $i : X' \to X$ 为概形态射。
假设下列条件中至少有一个成立：
\begin{enumerate}
\item $i$ 是拟紧浸入；
\item $X'$ 拟紧，且 $i$ 是浸入；
\item $i$ 拟紧，并在 $X'$ 与 $i(X')$ 之间诱导同胚；
\item $X'$ 拟紧，并且 $i$ 在 $X'$ 与 $i(X')$ 之间诱导同胚。
\end{enumerate}
则 $i^*\mathcal{L}$ 在 $X'$ 上丰沛。
\end{lemma}

\begin{proof}
注意，在情形 (1) 与 (3) 中，概形 $X'$ 拟紧，
因为依定义 \ref{properties-definition-ample}，$X$ 拟紧。
所以只需证明 (2) 与 (4)。由于 (2) 是 (4) 的特殊情形，
只需证明 (4)。

\medskip\noindent
假设条件 (4) 成立。对于 $s \in \Gamma(X, \mathcal{L}^{\otimes d})$，
以 $s' = i^*s$ 表示 $s$ 在 $X'$ 上的拉回。注意，
$s'$ 是 $(i^*\mathcal{L})^{\otimes d}$ 的截面。由命题
\ref{properties-proposition-characterize-ample}，开集 $X_s$（其中
$s \in \Gamma(X, \mathcal{L}^{\otimes d})$）构成 $X$ 的拓扑基。
由《模》备注 \ref{modules-remark-pullback-s-open}，
$X'_{s'} = i^{-1}(X_s)$；又因为 $X' \to i(X')$ 是同胚，
所以开集 $X'_{s'}$ 构成 $X'$ 的拓扑基。因此，由命题
\ref{properties-proposition-characterize-ample}，$i^*\mathcal{L}$ 丰沛。
\end{proof}

\begin{lemma}
\label{properties-lemma-ample-on-product}
设 $S$ 为拟分离概形。设 $X$、$Y$ 为 $S$ 上的概形。
设 $\mathcal{L}$ 为丰沛可逆 $\mathcal{O}_X$-模，
并设 $\mathcal{N}$ 为丰沛可逆 $\mathcal{O}_Y$-模。
则 $\mathcal{M} = \text{pr}_1^*\mathcal{L}
\otimes_{\mathcal{O}_{X \times_S Y}} \text{pr}_2^*\mathcal{N}$
是 $X \times_S Y$ 上的丰沛可逆层。
\end{lemma}

\begin{proof}
态射 $i : X \times_S Y \to X \times Y$ 是拟紧浸入；见《概形》引理
\ref{schemes-lemma-fibre-product-after-map}。另一方面，$\mathcal{M}$
是沿 $i$ 拉回 $X \times Y$ 上相应可逆模层所得。
由引理 \ref{properties-lemma-ample-on-locally-closed}，只需对 $X \times Y$
证明该引理。我们检验定义 \ref{properties-definition-ample} 中
$\mathcal{M}$ 在 $X \times Y$ 上满足 (1) 与 (2)。

\medskip\noindent
由于 $X$ 和 $Y$ 都拟紧，$X \times Y$ 也拟紧。
令 $z \in X \times Y$ 为一点。令 $x \in X$ 和 $y \in Y$
为其投影。选取 $n > 0$ 以及
$s \in \Gamma(X, \mathcal{L}^{\otimes n})$，
使得 $X_s$ 是 $x$ 的仿射开邻域。选取 $m > 0$ 以及
$t \in \Gamma(Y, \mathcal{N}^{\otimes m})$，
使得 $Y_t$ 是 $y$ 的仿射开邻域。于是
$r = \text{pr}_1^*s \otimes \text{pr}_2^*t$ 是 $\mathcal{M}$ 的截面，且
$(X \times Y)_r = X_s \times Y_t$。这是 $z$ 的仿射开邻域，
证明完成。
\end{proof}







\section{仿射概形与拟仿射概形}
\label{properties-section-affine-quasi-affine}

\begin{lemma}
\label{properties-lemma-quasi-affine-O-ample}
设 $X$ 为概形。
则 $X$ 拟仿射当且仅当 $\mathcal{O}_X$ 丰沛。
\end{lemma}

\begin{proof}
假设 $X$ 拟仿射。置 $A = \Gamma(X, \mathcal{O}_X)$。
考虑由引理 \ref{properties-lemma-quasi-affine} 给出的如下开浸入：
$$
j : X \longrightarrow \Spec(A)
$$
注意，$\Spec(A) = \text{Proj}(A[T])$；见《构造》例
\ref{constructions-example-trivial-proj}。
所以可应用引理 \ref{properties-lemma-open-in-proj-ample}，断定
$\mathcal{O}_X$ 丰沛。

\medskip\noindent
假设 $\mathcal{O}_X$ 丰沛。注意，作为分次环，
$\Gamma_*(X, \mathcal{O}_X) \cong A[T]$。
所以，由引理 \ref{properties-lemma-ample-immersion-into-proj} 与
\ref{properties-lemma-quasi-affine}，并结合上面已见的等式
$\Spec(A) = \text{Proj}(A[T])$ 对任意环 $A$ 都成立这一事实，
即可得到结论。
\end{proof}

\begin{lemma}
\label{properties-lemma-quasi-affine-locally-closed}
设 $X$ 为拟仿射概形。对任意拟紧浸入
$i : X' \to X$，概形 $X'$ 都拟仿射。
\end{lemma}

\begin{proof}
这也可不使用丰沛可逆层的材料而直接证明；我们建议读者在餐巾纸上
自行完成。由于 $X$ 拟仿射，由引理
\ref{properties-lemma-quasi-affine-O-ample}，$\mathcal{O}_X$ 丰沛。
于是，由引理 \ref{properties-lemma-ample-on-locally-closed}，
$\mathcal{O}_{X'}$ 丰沛。再由引理
\ref{properties-lemma-quasi-affine-O-ample}，$X'$ 拟仿射。
\end{proof}

\begin{lemma}
\label{properties-lemma-characterize-affine}
设 $X$ 为概形。假设存在有限多个元素
$f_1, \ldots, f_n \in \Gamma(X, \mathcal{O}_X)$，使得
\begin{enumerate}
\item 每个 $X_{f_i}$ 都是 $X$ 的仿射开集；并且
\item $f_1, \ldots, f_n$ 在
$\Gamma(X, \mathcal{O}_X)$ 中生成的理想等于单位理想。
\end{enumerate}
则 $X$ 仿射。
\end{lemma}

\begin{proof}
假设有引理所述的 $f_1, \ldots, f_n$。
可写成 $1 = \sum g_i f_i$，其中某些
$g_j \in \Gamma(X, \mathcal{O}_X)$，所以显然
$X = \bigcup X_{f_i}$。
（$f_i$ 不可能在某点全都消失。）由于每个 $X_{f_i}$
都拟紧（因为它仿射），可知 $X$ 拟紧。所以，由上面的引理
\ref{properties-lemma-quasi-affine-O-ample}，$X$ 拟仿射。
考虑开浸入
$$
j : X \to \Spec(\Gamma(X, \mathcal{O}_X)),
$$
见引理 \ref{properties-lemma-quasi-affine}。右端标准开集
$D(f_i)$ 的逆像等于左端的 $X_{f_i}$，并且态射 $j$ 诱导同构
$X_{f_i} \cong D(f_i)$；见引理
\ref{properties-lemma-invert-f-affine}。由于 $f_i$ 生成单位理想，可知
$\Spec(\Gamma(X, \mathcal{O}_X))
= \bigcup_{i = 1, \ldots, n} D(f_i)$。所以 $j$ 是同构。
\end{proof}







\section{准凝聚层与丰沛可逆层}
\label{properties-section-ample-quasi-coherent}

\noindent
本节的主题是：在有丰沛可逆层时，每个准凝聚层都来自某个分次模。

\begin{situation}
\label{properties-situation-ample}
设 $X$ 为概形。
设 $\mathcal{L}$ 为 $X$ 上的丰沛可逆层。
置 $S = \Gamma_*(X, \mathcal{L})$，视为分次环。
置 $Y = \text{Proj}(S)$。
设 $f : X \to Y$ 为引理 \ref{properties-lemma-map-into-proj} 的典范态射。
它配备一个 $\mathbf{Z}$-分次 $\mathcal{O}_X$-代数映射
$\bigoplus f^*\mathcal{O}_Y(n) \to \bigoplus \mathcal{L}^{\otimes n}$。
\end{situation}

\noindent
下一个引理其实是再下一个引理的特殊情形，不过先明确指出它成立
似乎更妥当。

\begin{lemma}
\label{properties-lemma-ample-gcd-is-one}
在情形 \ref{properties-situation-ample} 中，典范态射 $f : X \to Y$
把 $X$ 映入开子概形 $W = W_1 \subset Y$；在其上，
$\mathcal{O}_Y(1)$ 可逆，并且所有乘法映射
$\mathcal{O}_Y(n) \otimes_{\mathcal{O}_Y} \mathcal{O}_Y(m) \to
\mathcal{O}_Y(n + m)$
都是同构（见《构造》引理
\ref{constructions-lemma-where-invertible}）。此外，所有映射
$f^*\mathcal{O}_Y(n) \to \mathcal{L}^{\otimes n}$ 都是同构。
\end{lemma}

\begin{proof}
由命题 \ref{properties-proposition-characterize-ample}，存在整数 $n_0$，使得
$\mathcal{L}^{\otimes n}$ 对所有 $n \geq n_0$ 都全局生成。
令 $x \in X$ 为一点。由上文，可找到
$a \in S_{n_0}$ 与 $b \in S_{n_0 + 1}$，使得
$a$ 和 $b$ 在 $x$ 处都不消失。因此
$f(x) \in D_{+}(a) \cap D_{+}(b) = D_{+}(ab)$。由《构造》引理
\ref{constructions-lemma-where-invertible}，可知
$f(x) \in W_1$，如所需。构造映射 $f$ 时曾使用《构造》引理
\ref{constructions-lemma-invertible-map-into-proj}；由该引理，映射
$f^*\mathcal{O}_Y(n_0) \to \mathcal{L}^{\otimes n_0}$ 与
$f^*\mathcal{O}_Y(n_0 + 1) \to \mathcal{L}^{\otimes n_0 + 1}$
在 $x$ 的某个邻域上是同构。结合代数结构的相容性以及 $f$ 映入 $W$
这一事实，可断定所有映射
$f^*\mathcal{O}_Y(n) \to \mathcal{L}^{\otimes n}$
在 $x$ 的某个邻域上都是同构。结论得证。
\end{proof}

\noindent
回顾《模》定义 \ref{modules-definition-gamma-star}：给定局部环层空间
$X$、可逆层 $\mathcal{L}$ 以及 $\mathcal{O}_X$-模 $\mathcal{F}$，
我们有分次 $\Gamma_*(X, \mathcal{L})$-模
$$
\Gamma_*(X, \mathcal{L}, \mathcal{F}) =
\bigoplus\nolimits_{n \in \mathbf{Z}}
\Gamma(X, \mathcal{F} \otimes_{\mathcal{O}_X} \mathcal{L}^{\otimes n}).
$$
下一个引理说明，在情形 \ref{properties-situation-ample} 中，可以从这个分次模
恢复准凝聚 $\mathcal{O}_X$-模 $\mathcal{F}$。还可参看《构造》引理
\ref{constructions-lemma-quasi-coherent-projective-space}，其中我们在
特殊情形 $X = \mathbf{P}^n_R$ 下证明了这个引理。

\begin{lemma}
\label{properties-lemma-ample-quasi-coherent}
在情形 \ref{properties-situation-ample} 中，设 $\mathcal{F}$ 为 $X$ 上的准凝聚层。
置 $M = \Gamma_*(X, \mathcal{L}, \mathcal{F})$，视为分次 $S$-模。
存在同构
$$
f^*\widetilde{M} \longrightarrow \mathcal{F}
$$
且关于 $\mathcal{F}$ 函子性成立，并使得
$M_0 \to \Gamma(\text{Proj}(S), \widetilde{M}) \to \Gamma(X, \mathcal{F})$
为恒等映射。
\end{lemma}

\begin{proof}
令 $s \in S_{+}$ 为齐次元素，使得 $X_s$ 是 $X$ 中的仿射开集。
回顾，$\widetilde{M}|_{D_{+}(s)}$ 对应于
$S_{(s)}$-模 $M_{(s)}$；见《构造》引理
\ref{constructions-lemma-proj-sheaves}。又回顾，
$f^{-1}(D_{+}(s)) = X_s$。由于 $X$ 带有丰沛可逆层，
它拟紧且拟分离；见第 \ref{properties-section-ample} 节。
由引理 \ref{properties-lemma-invert-s-sections}，存在典范同构
$M_{(s)} = \Gamma_*(X, \mathcal{L}, \mathcal{F})_{(s)} \to
\Gamma(X_s, \mathcal{F})$。由于 $\mathcal{F}$ 准凝聚，这给出典范同构
$$
f^*\widetilde{M}|_{X_s} \to \mathcal{F}|_{X_s}
$$
由于 $\mathcal{L}$ 在 $X$ 上丰沛，可知 $X$ 被形如 $X_s$ 的仿射开集
覆盖。因此，只需证明所显示的映射可在交叠上粘合。证明从略。
\end{proof}

\begin{remark}
\label{properties-remark-neurotic}
采用引理 \ref{properties-lemma-ample-quasi-coherent} 的假设与记号。
把该引理所显示的映射记为 $\theta_\mathcal{F}$。注意，引理
\ref{properties-lemma-ample-gcd-is-one} 的同构
$f^*\mathcal{O}_Y(n) \to \mathcal{L}^{\otimes n}$ 正是
$\theta_{\mathcal{L}^{\otimes n}}$。考虑乘法映射
$$
\widetilde{M} \otimes_{\mathcal{O}_Y} \mathcal{O}_Y(n)
\longrightarrow
\widetilde{M(n)}
$$
见《构造》方程
(\ref{constructions-equation-multiply-more-generally})。
把它拉回到 $X$，并考虑
$$
\xymatrix{
f^*\widetilde{M} \otimes_{\mathcal{O}_X} f^*\mathcal{O}_Y(n)
\ar[r]
\ar[d]_{\theta_\mathcal{F} \otimes \theta_{\mathcal{L}^{\otimes n}}}
&
f^*\widetilde{M(n)}
\ar[d]^{\theta_{\mathcal{F} \otimes \mathcal{L}^{\otimes n}}}
\\
\mathcal{F} \otimes \mathcal{L}^{\otimes n} \ar[r]^{\text{id}} &
\mathcal{F} \otimes \mathcal{L}^{\otimes n}
}
$$
这里使用了显然的等同
$M(n) = \Gamma_*(X, \mathcal{L}, \mathcal{F} \otimes \mathcal{L}^{\otimes n})$。
该图交换。证明从略。
\end{remark}

\noindent
下面的引理应当可以从引理 \ref{properties-lemma-ample-quasi-coherent}
推出（反之亦然），但直接重复证明似乎更简单。

\begin{lemma}
\label{properties-lemma-proj-quasi-coherent}
设 $S$ 为分次环，并且 $X = \text{Proj}(S)$ 拟紧。
设 $\mathcal{F}$ 为准凝聚 $\mathcal{O}_X$-模。置
$M = \bigoplus_{n \in \mathbf{Z}} \Gamma(X, \mathcal{F}(n))$，
视为分次 $S$-模；见《构造》第
\ref{constructions-section-invertible-on-proj} 节。《构造》引理
\ref{constructions-lemma-comparison-proj-quasi-coherent} 的映射
$$
\widetilde{M} \longrightarrow \mathcal{F}
$$
是同构。如果 $X$ 被标准开集 $D_+(f)$ 覆盖，其中 $f$ 的次数为 $1$，
则所诱导的映射
$M_n \to \Gamma(X, \mathcal{F}(n))$ 是恒等映射。
\end{lemma}

\begin{proof}
由于 $X$ 拟紧，可以找到正次数的齐次元素
$f_1, \ldots, f_n \in S$，使得
$X = D_+(f_1) \cup \ldots \cup D_+(f_n)$。令 $d$ 为
$f_1, \ldots, f_n$ 各次数的最小公倍数。用 $f_i$ 的某次幂替换它后，
可假设每个 $f_i$ 的次数都是 $d$。于是可知
$\mathcal{L} = \mathcal{O}_X(d)$ 可逆，乘法映射
$\mathcal{O}_X(ad) \otimes \mathcal{O}_X(bd) \to \mathcal{O}_X((a + b)d)$
都是同构，并且每个 $f_i$ 决定一个整体截面 $s_i$，它属于 $\mathcal{L}$，
使得 $X_{s_i} = D_+(f_i)$；见《构造》引理
\ref{constructions-lemma-where-invertible} 与
\ref{constructions-lemma-principal-open}。因此
$\Gamma(X, \mathcal{F}(ad)) =
\Gamma(X, \mathcal{F} \otimes \mathcal{L}^{\otimes a})$。
回顾，$\widetilde{M}|_{D_{+}(f_i)}$ 对应于
$S_{(f_i)}$-模 $M_{(f_i)}$；见《构造》引理
\ref{constructions-lemma-proj-sheaves}。由于 $f_i$ 的次数是 $d$，
$M_{(f_i)}$ 的同构类只依赖于 $M$ 中次数可被 $d$ 整除的齐次分量。
更准确地说，$M_{(f_i)}$ 的同构类只依赖于分次
$\Gamma_*(X, \mathcal{L})$-模
$\Gamma_*(X, \mathcal{L}, \mathcal{F})$，以及像 $s_i$，即 $f_i$ 在
$\Gamma_*(X, \mathcal{L})$ 中的像。概形 $X$ 依假设拟紧，
并且由《构造》引理
\ref{constructions-lemma-proj-separated} 分离。由引理
\ref{properties-lemma-invert-s-sections}，存在典范同构
$$
M_{(f_i)} = \Gamma_*(X, \mathcal{L}, \mathcal{F})_{(s_i)} \to
\Gamma(X_{s_i}, \mathcal{F}).
$$
于是，《构造》引理
\ref{constructions-lemma-comparison-proj-quasi-coherent}
中该映射的构造表明，它在 $D_+(f_i)$ 上是同构；
由于这些开集覆盖 $X$，它就是同构。最后一个陈述的证明从略。
\end{proof}




\section{寻找适当的仿射开集}
\label{properties-section-finding-affine-opens}

\noindent
本节汇集在一般程度各异的情形下有关仿射开集存在性的一些结果。

\begin{lemma}
\label{properties-lemma-maximal-points-affine}
设 $X$ 为拟分离概形。设 $Z_1, \ldots, Z_n$ 为 $X$
的两两不同的不可约分支；见《拓扑》第
\ref{topology-section-irreducible-components} 节。设
$\eta_i \in Z_i$ 为它们的泛点；见《概形》引理
\ref{schemes-lemma-scheme-sober}。存在仿射开邻域
$\eta_i \in U_i$，使得 $U_i \cap U_j = \emptyset$
对所有 $i \not = j$ 都成立。特别地，
$U = U_1 \cup \ldots \cup U_n$ 是包含所有点
$\eta_1, \ldots, \eta_n$ 的仿射开集。
\end{lemma}

\begin{proof}
令 $V_i$ 为任意仿射开集，它包含 $\eta_i$，并且与闭集
$Z_1 \cup \ldots \hat Z_i \ldots \cup Z_n$ 不交。
由于 $X$ 拟分离，对每个 $i$，并集
$W_i = \bigcup_{j, j \not = i} V_i \cap V_j$ 是 $V_i$
中不包含 $\eta_i$ 的拟紧开集。由《代数》引理
\ref{algebra-lemma-standard-open-containing-maximal-point}，
可找到开邻域 $U_i \subset V_i$，它包含 $\eta_i$ 且与 $W_i$ 不交。
最后，$U$ 仿射，因为它是环
$R_1 \times \ldots \times R_n$ 的谱，其中
$R_i = \mathcal{O}_X(U_i)$；见《概形》引理
\ref{schemes-lemma-disjoint-union-affines}。
\end{proof}

\begin{remark}
\label{properties-remark-maximal-points-affine}
若 $X$ 不拟分离，则上面的引理
\ref{properties-lemma-maximal-points-affine} 不成立。下面给出一例。取
$R = \mathbf{Q}[x, y_1, y_2, \ldots]/((x-i)y_i)$。
考虑极小素理想 $\mathfrak p = (y_1, y_2, \ldots)$，它是 $R$ 的理想。
把 $\Spec(R)$ 的两个副本沿着（非拟紧）开集
$\Spec(R) \setminus V(\mathfrak p)$ 粘合，
得到概形 $X$（按《概形》例
\ref{schemes-example-affine-space-zero-doubled} 的方式粘合）。
那么，$X$ 中对应于 $\mathfrak p$ 的两个极大点不包含在同一个
仿射开集中。理由是，$\Spec(R)$ 中任何包含 $\mathfrak p$ 的开集
都包含无穷多条“直线” $x = i$、$y_j = 0$，其中
$j \not = i$，其参数为 $y_i$。细节从略。
\end{remark}

\noindent
尽管有上面的例子，对于“大多数”有限族不可约闭子集，至少在
$X$ 拟紧时，仍可应用上面的引理
\ref{properties-lemma-maximal-points-affine}。这是因为 $X$ 含有稠密的分离开集。

\begin{lemma}
\label{properties-lemma-quasi-compact-dense-open-separated}
设 $X$ 为拟紧概形。存在分离的稠密开集 $V \subset X$。
\end{lemma}

\begin{proof}
设 $X = \bigcup_{i = 1, \ldots, n} U_i$ 是 $n$ 个仿射开子概形之并。
我们对 $n$ 归纳证明。$n = 1$ 时显然。设
$V' \subset \bigcup_{i = 1, \ldots, n - 1} U_i$ 为分离的稠密开子概形；
归纳假设保证它存在。考虑
$$
V = V' \amalg (U_n \setminus \overline{V'}).
$$
显然，$V$ 是 $X$ 的分离稠密开子概形。
\end{proof}

\noindent
事实表明，即使 $X$ 既拟分离又拟紧，也未必存在分离、拟紧的
稠密开集；见《例》引理
\ref{examples-lemma-no-dense-separated-quasi-compact-open-in-qcqs}。
下面稍微加强上面的引理 \ref{properties-lemma-maximal-points-affine}。

\begin{lemma}
\label{properties-lemma-point-and-maximal-points-affine}
设 $X$ 为拟分离概形。设 $Z_1, \ldots, Z_n$ 为 $X$
的两两不同的不可约分支。设 $\eta_i \in Z_i$ 为其泛点。
任取 $x \in X$。存在仿射开集 $U \subset X$，包含
$x$ 以及所有 $\eta_i$。
\end{lemma}

\begin{proof}
假设 $x \in Z_1 \cap \ldots \cap Z_r$，且
$x \not \in Z_{r + 1}, \ldots, Z_n$。于是可选取仿射开集
$W \subset X$，使得 $x \in W$，并且
$W \cap Z_i = \emptyset$ 对 $i = r + 1, \ldots, n$ 成立。
注意，显然 $\eta_i \in W$ 对 $i = 1, \ldots, r$ 成立。
由引理 \ref{properties-lemma-maximal-points-affine}，可选取两两不交的仿射开集
$U_i \subset X$，使得 $\eta_i \in U_i$ 对
$i = r + 1, \ldots, n$ 成立。由于 $X$ 拟分离，开集
$W \cap U_i$ 拟紧，而且不包含 $\eta_i$，其中
$i = r + 1, \ldots, n$。所以，由《代数》引理
\ref{algebra-lemma-standard-open-containing-maximal-point}，可缩小 $U_i$，
使得 $W \cap U_i = \emptyset$ 对
$i = r + 1, \ldots, n$ 成立。
于是并集
$U = W \cup \bigcup_{i = r + 1, \ldots, n} U_i$ 是不交并，因而
（由《概形》引理 \ref{schemes-lemma-disjoint-union-affines}）
是所需的仿射开集。
\end{proof}

\begin{lemma}
\label{properties-lemma-ample-finite-set-in-affine}
设 $X$ 为概形。假设下列条件之一成立：
\begin{enumerate}
\item 概形 $X$ 拟仿射；
\item 概形 $X$ 同构于某个仿射概形的局部闭子概形；
\item $X$ 上存在丰沛可逆层；
\item 概形 $X$ 同构于 $\text{Proj}(S)$ 的局部闭子概形，
其中 $S$ 为某个分次环。
\end{enumerate}
则对任意有限子集 $E \subset X$，存在仿射开集
$U \subset X$，使得 $E \subset U$。
\end{lemma}

\begin{proof}
由《性质》定义 \ref{properties-definition-quasi-affine}，拟仿射概形是某个仿射概形的
拟紧开子概形。任意仿射概形 $\Spec(R)$ 都同构于
$\text{Proj}(R[X])$，其中 $R[X]$ 以 $\deg(X) = 1$ 分次。
由命题 \ref{properties-proposition-characterize-ample}，若 $X$ 有丰沛可逆层，
则 $X$ 同构于 $\text{Proj}(S)$ 的开子概形，其中 $S$ 为某个分次环。
所以，只需在情形 (4) 中证明引理。
（我们建议读者自行直接证明情形 (2)。）

\medskip\noindent
因此，假设 $X \subset \text{Proj}(S)$ 是局部闭子概形，其中 $S$
是某个分次环。令 $T = \overline{X} \setminus X$。
回顾，标准开集 $D_{+}(f)$ 构成 $\text{Proj}(S)$ 的拓扑基。
由于 $E$ 有限，可选取有限多个齐次元素 $f_i \in S_{+}$，使得
$$
E \subset
D_{+}(f_1) \cup \ldots \cup D_{+}(f_n) \subset
\text{Proj}(S) \setminus T
$$
假设 $E = \{\mathfrak p_1, \ldots, \mathfrak p_m\}$
是 $\text{Proj}(S)$ 的子集。
考虑理想 $I = (f_1, \ldots, f_n) \subset S$。
由于 $I \not \subset \mathfrak p_j$ 对所有
$j = 1, \ldots, m$ 都成立，由《代数》引理
\ref{algebra-lemma-graded-silly} 可知，存在齐次元素
$f \in I$，并且 $f \not \in \mathfrak p_j$ 对所有
$j = 1, \ldots, m$ 都成立。于是 $E \subset D_{+}(f) \subset
D_{+}(f_1) \cup \ldots \cup D_{+}(f_n)$。由于 $D_{+}(f)$ 与
$T$ 不交，可知 $X \cap D_{+}(f)$ 是仿射概形
$D_{+}(f)$ 的闭子概形，因而是 $X$ 的仿射开集，如所需。
\end{proof}

\begin{lemma}
\label{properties-lemma-ample-finite-set-in-principal-affine}
设 $X$ 为概形。设 $\mathcal{L}$ 为 $X$ 上的丰沛可逆层。设
$$
E \subset W \subset X
$$
其中 $E$ 有限，且 $W$ 是 $X$ 中的开集。则存在 $n > 0$
以及截面 $s \in \Gamma(X, \mathcal{L}^{\otimes n})$，使得
$X_s$ 仿射，并且 $E \subset X_s \subset W$。
\end{lemma}

\begin{proof}
读者可修改引理 \ref{properties-lemma-ample-finite-set-in-affine} 的证明来证明本引理；
我们改从该引理推出它。由引理
\ref{properties-lemma-ample-finite-set-in-affine}，可选取仿射开集
$U \subset W$，使得 $E \subset U$。考虑分次环
$S = \Gamma_*(X, \mathcal{L}) =
\bigoplus_{n \geq 0} \Gamma(X, \mathcal{L}^{\otimes n})$。
对每个 $x \in E$，令 $\mathfrak p_x \subset S$ 为由在 $x$ 处消失的
截面组成的分次理想。显然，$\mathfrak p_x$ 是素理想；又因为
$\mathcal{L}$ 的某个幂全局生成，显然
$S_{+} \not \subset \mathfrak p_x$。令 $I \subset S$ 为由在
$X \setminus U$ 的所有点处消失的截面组成的分次理想。
由于集合 $X_s$ 构成拓扑基，可知
$I \not \subset \mathfrak p_x$ 对所有 $x \in E$ 都成立。
由（分次）素理想回避（《代数》引理
\ref{algebra-lemma-graded-silly}），可找到齐次 $s \in I$，
使得 $s \not \in \mathfrak p_x$ 对所有 $x \in E$ 都成立。
于是 $E \subset X_s \subset U$，并且由引理
\ref{properties-lemma-affine-cap-s-open}，$X_s$ 仿射。
\end{proof}

\begin{lemma}
\label{properties-lemma-quasi-affine-invertible-nonvanishing-section}
设 $X$ 为拟仿射概形。设 $\mathcal{L}$ 为可逆
$\mathcal{O}_X$-模。设 $E \subset W \subset X$，其中 $E$ 有限且
$W$ 开。则存在 $s \in \Gamma(X, \mathcal{L})$，使得
$X_s$ 仿射，并且 $E \subset X_s \subset W$。
\end{lemma}

\begin{proof}
本引理的证明与《代数》引理 \ref{algebra-lemma-silly} 的证明有许多
共同之处。设 $E = \{x_1, \ldots, x_n\}$。若
$E = W = \emptyset$，则 $s = 0$ 可用。若 $W \not = \emptyset$，
则必要时添加一点后，可假设 $E \not = \emptyset$。所以可假设
$n \geq 1$。我们对 $n$ 归纳证明本引理。

\medskip\noindent
基础情形：$n = 1$。用 $W$ 中 $x_1$ 的一个仿射开邻域替换 $W$ 后，
可假设 $W$ 仿射。结合引理 \ref{properties-lemma-quasi-affine-O-ample} 与命题
\ref{properties-proposition-characterize-ample}，可知每个准凝聚
$\mathcal{O}_X$-模都全局生成。因此，存在整体截面 $s$，
它属于 $\mathcal{L}$，且在 $x_1$ 处不消失。另一方面，令 $Z \subset X$ 为
$X \setminus W$ 上的约化诱导闭子概形。把全局生成性应用于
准凝聚理想层 $\mathcal{I}$，即 $Z$ 的理想层，可找到整体截面 $f$，
它属于 $\mathcal{I}$，
它在 $x_1$ 处不消失。于是 $s' = fs$ 是 $\mathcal{L}$ 的整体截面，
它在 $x_1$ 处不消失，并且 $X_{s'} \subset W$。由引理
\ref{properties-lemma-affine-cap-s-open}，$X_{s'}$ 仿射。

\medskip\noindent
归纳步骤：$n > 1$。若存在特化关系 $x_i \leadsto x_j$，其中
$i \not = j$，则只需对
$\{x_1, \ldots, x_n\} \setminus \{x_i\}$ 证明引理，归纳即得结论。
因此可假设各 $x_i$ 之间没有特化关系。由引理
\ref{properties-lemma-ample-finite-set-in-affine} 或引理
\ref{properties-lemma-ample-finite-set-in-principal-affine}，可假设 $W$ 仿射。
由归纳，可找到整体截面 $s$，它属于 $\mathcal{L}$，并使得
$X_s \subset W$ 仿射，并包含 $x_1, \ldots, x_{n - 1}$。
若 $x_n \in X_s$，则已经完成。假设 $s$ 在 $x_n$ 处为零。
由 $n = 1$ 的情形，可找到整体截面 $s'$，它属于 $\mathcal{L}$，并使得
$\{x_n\} \subset X_{s'} \subset
W \setminus \overline{\{x_1, \ldots, x_{n - 1}\}}$。
这里用到了 $x_n$ 不是 $x_1, \ldots, x_{n - 1}$ 的特化。
于是 $s + s'$ 是 $\mathcal{L}$ 的整体截面，它在
$x_1, \ldots, x_n$ 处都不消失，并且
$X_{s + s'} \subset W$；如前即可得出结论。
\end{proof}

\begin{lemma}
\label{properties-lemma-ring-affine-open-injective-local-ring}
设 $X$ 为概形，且 $x \in X$ 为一点。在下列每种情形中，都存在
仿射开邻域 $U \subset X$，它包含 $x$，并使得典范映射
$\mathcal{O}_X(U) \to \mathcal{O}_{X, x}$ 为单射：
\begin{enumerate}
\item $X$ 整；
\item $X$ 局部诺特；
\item $X$ 约化且只有有限多个不可约分支。
\end{enumerate}
\end{lemma}

\begin{proof}
翻译成代数后，这由《代数》引理
\ref{algebra-lemma-subring-of-local-ring} 得出。
\end{proof}

\begin{lemma}
\label{properties-lemma-standard-open-two-affines}
设 $U$、$V$ 为仿射概形，并设 $W \to U$ 与 $W \to V$
为开浸入。对任意 $w \in W$，存在仿射开邻域
$W' \subset W$，它包含 $w$，并使得 $W'$ 映至 $U$ 和 $V$ 中的标准开集。
\end{lemma}

\begin{proof}
令 $X$ 为把 $U$ 与 $V$ 沿 $W$ 粘合所得的概形，并应用《概形》引理
\ref{schemes-lemma-standard-open-two-affines}。
\end{proof}
% Generated mechanically from the frozen morphisms.tex target.
% Do not edit this generated file; edit the bound locale source instead.

% OK, start here.
%

\setcounter{chapter}{28}
\chapter{概形的态射}



\phantomsection
\label{morphisms-section-phantom}


\section{引言}
\label{morphisms-section-introduction}

\noindent
本章介绍若干类概形态射。基本参考文献为 \cite{EGA}。




















\section{闭浸入}
\label{morphisms-section-closed-immersions}

\noindent
本节阐明此前所得的若干概形闭浸入结果。回顾，概形态射
$i : Z \to X$ 按定义是闭浸入，如果：(a) $i$ 诱导到 $X$
的某个闭子集上的同胚；(b)
$i^\sharp : \mathcal{O}_X \to i_*\mathcal{O}_Z$ 满射；并且 (c)
$i^\sharp$ 的核局部由截面生成；见《概形》定义
\ref{schemes-definition-immersion} 与
\ref{schemes-definition-closed-immersion-locally-ringed-spaces}。
事实表明，在已知 $Z$ 与 $X$ 都是概形时，有许多不同方式可以刻画
闭浸入。

\begin{lemma}
\label{morphisms-lemma-closed-immersion}
设 $i : Z \to X$ 为概形态射。下列条件等价：
\begin{enumerate}
\item 态射 $i$ 是闭浸入。
\item 对每个仿射开集 $\Spec(R) = U \subset X$，都存在理想
$I \subset R$，使得 $i^{-1}(U) = \Spec(R/I)$ 作为
$U = \Spec(R)$ 上的概形成立。
\item 存在仿射开覆盖 $X = \bigcup_{j \in J} U_j$，其中
$U_j = \Spec(R_j)$，并且对每个 $j \in J$ 都存在理想
$I_j \subset R_j$，使得 $i^{-1}(U_j) = \Spec(R_j/I_j)$ 作为
$U_j = \Spec(R_j)$ 上的概形成立。
\item 态射 $i$ 诱导 $Z$ 到 $X$ 的某个闭子集上的同胚，并且
$i^\sharp : \mathcal{O}_X \to i_*\mathcal{O}_Z$ 满射。
\item 态射 $i$ 诱导 $Z$ 到 $X$ 的某个闭子集上的同胚，映射
$i^\sharp : \mathcal{O}_X \to i_*\mathcal{O}_Z$ 满射，并且核
$\Ker(i^\sharp)\subset \mathcal{O}_X$ 是准凝聚理想层。
\item 态射 $i$ 诱导 $Z$ 到 $X$ 的某个闭子集上的同胚，映射
$i^\sharp : \mathcal{O}_X \to i_*\mathcal{O}_Z$ 满射，并且核
$\Ker(i^\sharp)\subset \mathcal{O}_X$ 是局部由截面生成的理想层。
\end{enumerate}
\end{lemma}

\begin{proof}
条件 (6) 就是闭浸入的定义；见《概形》定义
\ref{schemes-definition-closed-immersion-locally-ringed-spaces}
与 \ref{schemes-definition-immersion}。所以
(6) $\Leftrightarrow$ (1)。由《概形》引理
\ref{schemes-lemma-closed-subspace-scheme}，有
(1) $\Rightarrow$ (2)。显然，(2) $\Rightarrow$ (3)。

\medskip\noindent
假设 (3)。每个态射
$\Spec(R_j/I_j) \to \Spec(R_j)$ 都是闭浸入；见《概形》例
\ref{schemes-example-closed-immersion-affines}。因此，
$i^{-1}(U_j) \to U_j$ 是到其像上的同胚，并且
$i^\sharp|_{U_j}$ 满射。所以 $i$ 是到其像上的同胚，并且
$i^\sharp$ 满射，因为这些性质可以局部检验。由此断定
(3) $\Rightarrow$ (4)。

\medskip\noindent
蕴含 (4) $\Rightarrow$ (1) 是《概形》引理
\ref{schemes-lemma-characterize-closed-immersions}。
蕴含 (5) $\Rightarrow$ (6) 显然。蕴含
(6) $\Rightarrow$ (5) 则由《概形》引理
\ref{schemes-lemma-closed-subspace-scheme} 得出。
\end{proof}



\begin{lemma}
\label{morphisms-lemma-closed-immersion-ideals}
设 $X$ 为概形。设 $i : Z \to X$ 与 $i' : Z' \to X$ 为闭浸入，
并考虑理想层
$\mathcal{I} = \Ker(i^\sharp)$ 与
$\mathcal{I}' = \Ker((i')^\sharp)$，它们都是
$\mathcal{O}_X$ 的理想层。
\begin{enumerate}
\item 态射 $i : Z \to X$ 可分解为 $Z \to Z' \to X$，
其中相应态射为某个 $a : Z \to Z'$，当且仅当
$\mathcal{I}' \subset \mathcal{I}$。若这成立，则 $a$ 是闭浸入。
\item $Z \cong Z'$ 作为 $X$ 上的同构成立，当且仅当
$\mathcal{I} = \mathcal{I}'$。
\end{enumerate}
\end{lemma}

\begin{proof}
这由《概形》第 \ref{schemes-section-closed-immersion} 节中关于闭子空间的
讨论得出，特别是《概形》引理
\ref{schemes-lemma-closed-immersion} 与
\ref{schemes-lemma-characterize-closed-subspace}。
也可以直接由上面引理 \ref{morphisms-lemma-closed-immersion} 中的刻画 (3) 得出。
\end{proof}

\begin{lemma}
\label{morphisms-lemma-closed-immersion-bijection-ideals}
设 $X$ 为概形。设 $\mathcal{I} \subset \mathcal{O}_X$ 为理想层。
下列条件等价：
\begin{enumerate}
\item $\mathcal{I}$ 作为 $\mathcal{O}_X$-模层局部由截面生成；
\item $\mathcal{I}$ 作为 $\mathcal{O}_X$-模层准凝聚；并且
\item 存在概形的闭浸入 $i : Z \to X$，其相应理想层
$\Ker(i^\sharp)$ 等于 $\mathcal{I}$。
\end{enumerate}
\end{lemma}

\begin{proof}
(1) 与 (2) 的等价性直接来自《概形》引理
\ref{schemes-lemma-closed-subspace-scheme}。若 (1) 成立，则由《概形》定义
\ref{schemes-definition-closed-subspace} 与例
\ref{schemes-example-closed-subspace}，存在闭子空间
$i : Z \to X$，使得 $\mathcal{I} = \Ker(i^\sharp)$。
由《概形》引理 \ref{schemes-lemma-closed-subspace-scheme}，
这是概形的闭浸入，故 (3) 成立。反过来，若 (3) 成立，
则由《概形》引理 \ref{schemes-lemma-closed-subspace-scheme}，
(2) 成立（该引理适用，因为概形的闭浸入尤其是局部环层空间的
闭浸入）。
\end{proof}

\begin{lemma}
\label{morphisms-lemma-base-change-closed-immersion}
闭浸入的基变换仍是闭浸入。
\end{lemma}

\begin{proof}
见《概形》引理 \ref{schemes-lemma-base-change-immersion}。
\end{proof}

\begin{lemma}
\label{morphisms-lemma-composition-closed-immersion}
闭浸入的复合仍是闭浸入。
\end{lemma}

\begin{proof}
我们已在《概形》引理
\ref{schemes-lemma-composition-immersion} 中见过这一点，
不过这里再给出一个证明。确切地说，这由引理
\ref{morphisms-lemma-closed-immersion} 中闭浸入的刻画 (3) 得出。
因为若 $I \subset R$ 是理想，且
$\overline{J} \subset R/I$ 是理想，则有 $\overline{J} = J/I$，
其中 $J \subset R$ 是某个包含 $I$ 的理想，并且
$(R/I)/\overline{J} = R/J$。
\end{proof}

\begin{lemma}
\label{morphisms-lemma-closed-immersion-quasi-compact}
闭浸入拟紧。
\end{lemma}

\begin{proof}
本引理与《概形》引理
\ref{schemes-lemma-closed-immersion-quasi-compact} 重复。
\end{proof}

\begin{lemma}
\label{morphisms-lemma-closed-immersion-separated}
闭浸入分离。
\end{lemma}

\begin{proof}
本引理是《概形》引理
\ref{schemes-lemma-immersions-monomorphisms} 的特殊情形。
\end{proof}




\section{浸入}
\label{morphisms-section-immersions}

\noindent
本节汇集若干关于浸入的事实。

\begin{lemma}
\label{morphisms-lemma-immersion-permanence}
设 $Z \to Y \to X$ 为概形态射。
\begin{enumerate}
\item 若 $Z \to X$ 是浸入，则 $Z \to Y$ 是浸入。
\item 若 $Z \to X$ 是拟紧浸入且 $Y \to X$ 拟分离，
则 $Z \to Y$ 是拟紧浸入。
\item 若 $Z \to X$ 是闭浸入且 $Y \to X$ 分离，
则 $Z \to Y$ 是闭浸入。
\end{enumerate}
\end{lemma}

\begin{proof}
在每种情形中，证明都归结为考察交换图
$$
\xymatrix{
Z \ar[r] \ar[rd] & Y \times_X Z \ar[r] \ar[d] & Z \ar[d] \\
& Y \ar[r] & X
}
$$
其中上方两条水平箭头的复合为恒等态射。我们证明 (1)。
第一条水平箭头是 $Y \times_X Z \to Z$ 的截面，因而由《概形》引理
\ref{schemes-lemma-section-immersion}，它是浸入。箭头
$Y \times_X Z \to Y$ 是 $Z \to X$ 的基变换，故为浸入
（《概形》引理 \ref{schemes-lemma-base-change-immersion}）。
最后，浸入的复合仍是浸入
（《概形》引理 \ref{schemes-lemma-composition-immersion}）。
这证明了 (1)。另两个结论以完全相同的方式证明。
\end{proof}

\begin{lemma}
\label{morphisms-lemma-factor-quasi-compact-immersion}
设 $h : Z \to X$ 为浸入。若 $h$ 拟紧，则可将
$h = i \circ j$ 分解，其中 $j : Z \to \overline{Z}$ 为开浸入，
而 $i : \overline{Z} \to X$ 为闭浸入。
\end{lemma}

\begin{proof}
注意，$h$ 拟紧且拟分离（见《概形》引理
\ref{schemes-lemma-immersions-monomorphisms}）。因此，由《概形》引理
\ref{schemes-lemma-push-forward-quasi-coherent}，
$h_*\mathcal{O}_Z$ 是准凝聚 $\mathcal{O}_X$-模层。由此可知
$\mathcal{I} = \Ker(\mathcal{O}_X \to h_*\mathcal{O}_Z)$
是准凝聚理想层；见《概形》第
\ref{schemes-section-quasi-coherent} 节。设
$\overline{Z} \subset X$ 为与 $\mathcal{I}$ 对应的闭子概形；
见引理 \ref{morphisms-lemma-closed-immersion-bijection-ideals}。
由《概形》引理 \ref{schemes-lemma-characterize-closed-subspace}，
态射 $h$ 可分解为 $h = i \circ j$，其中
$i : \overline{Z} \to X$ 为包含态射。为说明 $j$ 是开浸入，
选取开子概形 $U \subset X$，使得 $h$ 诱导 $Z$ 到 $U$ 的闭浸入。
此时显然，$\mathcal{I}|_U$ 正是与闭浸入 $Z \to U$ 对应的理想层。
因此，$Z = \overline{Z} \cap U$。
\end{proof}

\begin{lemma}
\label{morphisms-lemma-factor-reduced-immersion}
设 $h : Z \to X$ 为浸入。若 $Z$ 约化，则可将
$h = i \circ j$ 分解，其中 $j : Z \to \overline{Z}$ 为开浸入，
而 $i : \overline{Z} \to X$ 为闭浸入。
\end{lemma}

\begin{proof}
设 $\overline{Z} \subset X$ 为 $h(Z)$ 的闭包，并赋予约化诱导闭子概形
结构；见《概形》定义
\ref{schemes-definition-reduced-induced-scheme}。
由《概形》引理 \ref{schemes-lemma-map-into-reduction}，
态射 $h$ 可分解为 $h = i \circ j$，其中
$i : \overline{Z} \to X$ 为包含态射，而
$j : Z \to \overline{Z}$。由浸入的定义可知，存在开子概形
$U \subset X$，使得 $h$ 可经到 $U$ 的闭浸入分解。因此，
$\overline{Z} \cap U$ 与 $h(Z)$ 都是 $U$ 的约化闭子概形，
并且具有相同的底闭集。于是由《概形》引理
\ref{schemes-lemma-reduced-closed-subscheme} 中的唯一性可知
$h(Z) \cong \overline{Z} \cap U$。所以 $j$ 诱导 $Z$ 到
$\overline{Z} \cap U$ 的同构。换言之，$j$ 是开浸入。
\end{proof}



\begin{example}
\label{morphisms-example-thibaut}
下面给出一个不能表示为先开浸入后闭浸入之复合的浸入例子。
设 $k$ 为域。设 $X = \Spec(k[x_1, x_2, x_3, \ldots])$。
设 $U = \bigcup_{n = 1}^{\infty} D(x_n)$。则
$U \to X$ 是开浸入。考虑理想
$$
I_n =
(x_1^n, x_2^n, \ldots, x_{n - 1}^n, x_n - 1, x_{n + 1}, x_{n + 2}, \ldots)
\subset
k[x_1, x_2, x_3, \ldots][1/x_n].
$$
注意，$I_n k[x_1, x_2, x_3, \ldots][1/x_nx_m] = (1)$
对任意 $m \not = n$ 都成立。因此，准凝聚理想
$\widetilde I_n$ 在 $D(x_n)$ 上给出，并且它们在
$D(x_nx_m)$ 上相容；确切地说，
$\widetilde I_n|_{D(x_nx_m)} = \mathcal{O}_{D(x_n x_m)}$
在 $n \not = m$ 时成立。
故这些理想粘合为准凝聚理想层
$\mathcal{I} \subset \mathcal{O}_U$。设 $Z \subset U$ 为与
$\mathcal{I}$ 对应的闭子概形。于是 $Z \to X$ 是浸入。

\medskip\noindent
我们断言，不能将 $Z \to X$ 分解为
$Z \to \overline{Z} \to X$，其中 $\overline{Z} \to X$ 为闭浸入，
而 $Z \to \overline{Z}$ 为开浸入。事实上，$\overline{Z}$ 必须由
某个理想 $I \subset k[x_1, x_2, x_3, \ldots]$ 定义，使得
$I_n = I k[x_1, x_2, x_3, \ldots][1/x_n]$。但是，唯一的元素
$f \in k[x_1, x_2, x_3, \ldots]$，其最终属于所有 $I_n$，
是 $0$！因此，$I$ 不存在。
\end{example}

\begin{lemma}
\label{morphisms-lemma-check-immersion}
设 $f : Y \to X$ 为概形态射。若对所有 $y \in Y$，都存在开子概形
$f(y) \in U \subset X$，使得
$f|_{f^{-1}(U)} : f^{-1}(U) \to U$ 是浸入，则 $f$ 是浸入。
\end{lemma}

\begin{proof}
这一陈述直接来自《概形》第 \ref{schemes-section-immersions} 节末尾
关于闭子概形的讨论，不过我们也给出详细证明。设
$Z \subset X$ 为 $f(Y)$ 的闭包。由于取闭包与限制到开集可交换，
由假设可知 $f(Y) \subset Z$ 是开集。故
$Z' = Z \setminus f(Y)$ 是闭集。因此，
$X' = X \setminus Z'$ 是 $X$ 的开子概形，并且 $f$ 可分解为
$f : Y \to X'$，随后接包含态射。若 $y \in Y$ 且
$U \subset X$ 如引理陈述中所述，则
$U' = X' \cap U$ 是 $f'(y)$ 的开邻域，使得
$(f')^{-1}(U') \to U'$ 是像为闭集的浸入
（引理 \ref{morphisms-lemma-immersion-permanence}）。因此它是闭浸入；
见《概形》引理 \ref{schemes-lemma-immersion-when-closed}。
由于闭浸入性在目标上是局部的
（例如由引理 \ref{morphisms-lemma-closed-immersion}），可断定
$f'$ 如所要求的是闭浸入。
\end{proof}






\section{闭浸入与准凝聚层}
\label{morphisms-section-closed-immersions-quasi-coherent}

\noindent
下列引理终于对概形上的准凝聚层完成了《模》引理
\ref{modules-lemma-i-star-exact} 对阿贝尔群层所完成的工作。
另见《模》第 \ref{modules-section-closed-immersion} 节中的讨论。

\begin{lemma}
\label{morphisms-lemma-i-star-equivalence}
设 $i : Z \to X$ 为概形的闭浸入。设
$\mathcal{I} \subset \mathcal{O}_X$ 为切出 $Z$ 的准凝聚理想层。
函子
$$
i_* :
\QCoh(\mathcal{O}_Z)
\longrightarrow
\QCoh(\mathcal{O}_X)
$$
正合且全忠实，其本质像由准凝聚 $\mathcal{O}_X$-模
$\mathcal{G}$ 中满足 $\mathcal{I}\mathcal{G} = 0$ 的那些模组成。
\end{lemma}

\begin{proof}
闭浸入拟紧且分离；见引理
\ref{morphisms-lemma-closed-immersion-quasi-compact} 与
\ref{morphisms-lemma-closed-immersion-separated}。因此《概形》引理
\ref{schemes-lemma-push-forward-quasi-coherent} 适用，并且
$Z$ 上准凝聚层的正像确实是 $X$ 上的准凝聚层。

\medskip\noindent
由《模》引理 \ref{modules-lemma-i-star-equivalence}，函子
$i_*$ 全忠实。

\medskip\noindent
现在来描述函子 $i_*$ 的本质像。我们有
$\mathcal{I}(i_*\mathcal{F}) = 0$ 对任意准凝聚
$\mathcal{O}_Z$-模成立；
例如可由《模》引理 \ref{modules-lemma-i-star-equivalence} 得出。
接着，设 $\mathcal{G}$ 为任意准凝聚 $\mathcal{O}_X$-模，且满足
$\mathcal{I}\mathcal{G} = 0$。只需证明典范映射
$$
\mathcal{G} \longrightarrow i_* i^*\mathcal{G}
$$
是同构\footnote{这已在更一般的情形中于《模》引理
\ref{modules-lemma-i-star-equivalence} 的证明中得到。}。
在概形与准凝聚模的情形中，在 $X$ 上仿射局部地工作，并使用引理
\ref{morphisms-lemma-closed-immersion} 与《概形》引理
\ref{schemes-lemma-widetilde-pullback}，即可归结为证明下列代数陈述：
给定环 $R$、理想 $I$ 以及 $R$-模 $N$，并假设 $IN = 0$，典范映射
$$
N \longrightarrow N \otimes_R R/I,\quad
n \longmapsto n \otimes 1
$$
是 $R$-模同构。这个简单代数事实的证明从略。
\end{proof}

\noindent
设 $i : Z \to X$ 为闭浸入。由于上引理，我们常滥用记号，用
$\mathcal{F}$ 表示层 $i_*\mathcal{F}$，并将其视为 $X$ 上的层。

\begin{lemma}
\label{morphisms-lemma-largest-quasi-coherent-subsheaf}
设 $X$ 为概形。设 $\mathcal{F}$ 为准凝聚
$\mathcal{O}_X$-模。设 $\mathcal{G} \subset \mathcal{F}$ 为
$\mathcal{O}_X$-子模。存在唯一的准凝聚
$\mathcal{O}_X$-子模 $\mathcal{G}' \subset \mathcal{G}$，具有下列性质：
对每个准凝聚 $\mathcal{O}_X$-模 $\mathcal{H}$，映射
$$
\Hom_{\mathcal{O}_X}(\mathcal{H}, \mathcal{G}')
\longrightarrow
\Hom_{\mathcal{O}_X}(\mathcal{H}, \mathcal{G})
$$
是双射。特别地，$\mathcal{G}'$ 是最大的准凝聚
$\mathcal{O}_X$-子模，它作为 $\mathcal{F}$ 的子模包含于
$\mathcal{G}$。
\end{lemma}

\begin{proof}
设 $\mathcal{G}_a$（$a \in A$）为所有准凝聚
$\mathcal{O}_X$-子模中包含于 $\mathcal{G}$ 的那些所成的集合。则映射
$$
\bigoplus\nolimits_{a \in A} \mathcal{G}_a \longrightarrow \mathcal{F}
$$
的像 $\mathcal{G}'$ 是准凝聚的，因为 $X$ 上准凝聚层之间映射的像
是准凝聚的，且准凝聚层的直和仍准凝聚；见《概形》第
\ref{schemes-section-quasi-coherent} 节。模 $\mathcal{G}'$ 包含于
$\mathcal{G}$。因此，这就是最大的准凝聚 $\mathcal{O}_X$-模，
并包含于 $\mathcal{G}$。

\medskip\noindent
为证明该公式，设 $\mathcal{H}$ 为准凝聚 $\mathcal{O}_X$-模，并设
$\alpha : \mathcal{H} \to \mathcal{G}$ 为 $\mathcal{O}_X$-模映射。
复合 $\mathcal{H} \to \mathcal{G} \to \mathcal{F}$ 的像是准凝聚的，
因为它是准凝聚层之间映射的像。因此它包含于 $\mathcal{G}'$。
故 $\alpha$ 如所要求地经 $\mathcal{G}'$ 分解。
\end{proof}

\begin{lemma}
\label{morphisms-lemma-i-upper-shriek}
设 $i : Z \to X$ 为概形的闭浸入。存在函子
\footnote{这很可能是非标准记号。}
$i^! : \QCoh(\mathcal{O}_X) \to \QCoh(\mathcal{O}_Z)$，
它是 $i_*$ 的右伴随。（比较《模》引理
\ref{modules-lemma-i-star-right-adjoint}。）
\end{lemma}

\begin{proof}
给定准凝聚 $\mathcal{O}_X$-模 $\mathcal{G}$，考察子层
$\mathcal{H}_Z(\mathcal{G})$；它是 $\mathcal{G}$ 中被
$\mathcal{I}$ 零化的局部截面所组成的子层。由引理
\ref{morphisms-lemma-largest-quasi-coherent-subsheaf}，存在典范的最大准凝聚
$\mathcal{O}_X$-子模 $\mathcal{H}_Z(\mathcal{G})'$。由构造，对任意
准凝聚 $\mathcal{O}_Z$-模 $\mathcal{F}$，都有
$$
\Hom_{\mathcal{O}_X}(i_*\mathcal{F}, \mathcal{H}_Z(\mathcal{G})')
=
\Hom_{\mathcal{O}_X}(i_*\mathcal{F}, \mathcal{G})
$$
。因此可以令 $i^!\mathcal{G} = i^*(\mathcal{H}_Z(\mathcal{G})')$。
细节从略。
\end{proof}

\noindent
利用准凝聚理想层与闭子概形之间的 $1$-to-$1$ 对应
（见引理 \ref{morphisms-lemma-closed-immersion-bijection-ideals}），
可以定义闭子概形的概形论交与概形论并。

\begin{definition}
\label{morphisms-definition-scheme-theoretic-intersection-union}
设 $X$ 为概形。设 $Z, Y \subset X$ 为分别对应于准凝聚理想层
$\mathcal{I}, \mathcal{J} \subset \mathcal{O}_X$ 的闭子概形。
$Z$ 与 $Y$ 的{\it 概形论交}是 $X$ 中由
$\mathcal{I} + \mathcal{J}$ 切出的闭子概形。
$Z$ 与 $Y$ 的{\it 概形论并}是 $X$ 中由
$\mathcal{I} \cap \mathcal{J}$ 切出的闭子概形。
\end{definition}

\noindent

\begin{lemma}
\label{morphisms-lemma-scheme-theoretic-intersection}
设 $X$ 为概形。设 $Z, Y \subset X$ 为闭子概形。设
$Z \cap Y$ 为 $Z$ 与 $Y$ 的概形论交。则
$Z \cap Y \to Z$ 与 $Z \cap Y \to Y$ 都是闭浸入，并且
$$
\xymatrix{
Z \cap Y \ar[r] \ar[d] & Z \ar[d] \\
Y \ar[r] & X
}
$$
是概形的笛卡尔图，即 $Z \cap Y = Z \times_X Y$。
\end{lemma}

\begin{proof}
态射 $Z \cap Y \to Z$ 与 $Z \cap Y \to Y$ 由引理
\ref{morphisms-lemma-closed-immersion-ideals} 可知都是闭浸入。
设 $U = \Spec(A)$ 为 $X$ 的仿射开集，并设 $Z \cap U$ 与
$Y \cap U$ 分别对应于理想 $I \subset A$ 与 $J \subset A$。
则 $Z \cap Y \cap U$ 对应于 $I + J \subset A$。由于
$A/I \otimes_A A/J = A/(I + J)$，由《概形》第
\ref{schemes-section-fibre-products} 节中对概形纤维积的描述可知，
该图是笛卡尔的。
\end{proof}

\begin{lemma}
\label{morphisms-lemma-scheme-theoretic-union}
设 $S$ 为概形。设 $X, Y \subset S$ 为闭子概形。设
$X \cup Y$ 为 $X$ 与 $Y$ 的概形论并。设
$X \cap Y$ 为 $X$ 与 $Y$ 的概形论交。则
$X \to X \cup Y$ 与 $Y \to X \cup Y$ 都是闭浸入，存在短正合列
$$
0 \to \mathcal{O}_{X \cup Y} \to \mathcal{O}_X \times \mathcal{O}_Y
\to \mathcal{O}_{X \cap Y} \to 0
$$
作为 $\mathcal{O}_S$-模的短正合列，并且图
$$
\xymatrix{
X \cap Y \ar[r] \ar[d] & X \ar[d] \\
Y \ar[r] & X \cup Y
}
$$
在概形范畴中是余笛卡尔的，即
$X \cup Y = X \amalg_{X \cap Y} Y$。
\end{lemma}

\begin{proof}
态射 $X \to X \cup Y$ 与 $Y \to X \cup Y$ 由引理
\ref{morphisms-lemma-closed-immersion-ideals} 可知都是闭浸入。
在短正合列中，我们利用引理 \ref{morphisms-lemma-i-star-equivalence} 的等价，
把 $S$ 的闭子概形上的准凝聚模看作 $S$ 上的准凝聚模。
对于该列的第一个映射，我们使用典范映射
$\mathcal{O}_{X \cup Y} \to \mathcal{O}_X$ 与
$\mathcal{O}_{X \cup Y} \to \mathcal{O}_Y$；
对于第二个映射，我们使用典范映射
$\mathcal{O}_X \to \mathcal{O}_{X \cap Y}$ 以及典范映射
$\mathcal{O}_Y \to \mathcal{O}_{X \cap Y}$ 的负映射。
于是可仿射局部地检验正合性。设 $U = \Spec(A)$ 为 $S$ 的仿射开集，
并设 $X \cap U$ 与 $Y \cap U$ 分别对应于理想
$I \subset A$ 与 $J \subset A$。则
$(X \cup Y) \cap U$ 对应于 $I \cap J \subset A$，而
$X \cap Y \cap U$ 对应于 $I + J \subset A$。
因此，正合性来自下列序列的正合性：
$$
0 \to A/I \cap J \to A/I \times A/J \to A/(I + J) \to 0
$$
为证明该图是余笛卡尔的，设给定概形 $T$ 以及概形态射
$f : X \to T$、$g : Y \to T$，它们作为态射
$X \cap Y \to T$ 相同。目标：证明存在唯一态射
$h : X \cup Y \to T$ 与 $f$ 和 $g$ 相同。为构造 $h$，可以在
$X \cup Y$ 上仿射局部地工作；见《概形》第
\ref{schemes-section-glueing-schemes} 节。若 $s \in X$ 且
$s \not \in Y$，则 $X \to X \cup Y$ 在 $s$ 的某邻域上是同构，
此时显然可构造 $h$。对 $s \in Y$ 且 $s \not \in X$ 也同样。
若 $s \in X \cap Y$，可选取仿射开集
$V = \Spec(B) \subset T$，它包含 $f(s) = g(s)$。
随后可选取仿射开集 $U = \Spec(A) \subset S$，它包含 $s$，
并使得 $f(X \cap U)$ 与 $g(Y \cap U)$ 都包含于 $V$。态射
$f|_{X \cap U}$ 与 $g|_{Y \cap V}$ 到 $V$，并对应于环映射
$$
B \to A/I
\quad\text{且}\quad
B \to A/J
$$
，它们作为到 $A/(I + J)$ 的映射相同。由上面显示的短正合列，
这些环同态存在唯一的提升，得到所需的环映射
$B \to A/I \cap J$。
\end{proof}









\section{模的支集}
\label{morphisms-section-support}

\noindent
本节汇集概形上准凝聚模支集的若干初等结果。回顾，模层的支集已在
《模》第 \ref{modules-section-support} 节中定义。另一方面，模的支集
已在《代数》第 \ref{algebra-section-support} 节中定义。二者相符。

\begin{lemma}
\label{morphisms-lemma-support-affine-open}
设 $X$ 为概形。设 $\mathcal{F}$ 为 $X$ 上的准凝聚层。
设 $\Spec(A) = U \subset X$ 为仿射开集，并令
$M = \Gamma(U, \mathcal{F})$。设 $x \in U$，并设
$\mathfrak p \subset A$ 为对应素理想。下列条件等价：
\begin{enumerate}
\item $\mathfrak p$ 属于 $M$ 的支集；以及
\item $x$ 属于 $\mathcal{F}$ 的支集。
\end{enumerate}
\end{lemma}

\begin{proof}
这来自等式 $\mathcal{F}_x = M_{\mathfrak p}$；见《概形》引理
\ref{schemes-lemma-spec-sheaves} 以及相应定义。
\end{proof}

\begin{lemma}
\label{morphisms-lemma-support-closed-specialization}
设 $X$ 为概形。设 $\mathcal{F}$ 为 $X$ 上的准凝聚层。
$\mathcal{F}$ 的支集在特殊化下封闭。
\end{lemma}

\begin{proof}
若 $x' \leadsto x$ 是特殊化且 $\mathcal{F}_x = 0$，则
$\mathcal{F}_{x'}$ 为零，因为 $\mathcal{F}_{x'}$ 是模
$\mathcal{F}_x$ 的局部化。因此，$\text{Supp}(\mathcal{F})$ 的补集
在一般化下封闭。
\end{proof}

\noindent
对有限型准凝聚模而言，支集是闭的，可以在纤维上检验，并与基变换
可交换。

\begin{lemma}
\label{morphisms-lemma-support-finite-type}
设 $\mathcal{F}$ 为概形 $X$ 上的有限型准凝聚模。则：
\begin{enumerate}
\item $\mathcal{F}$ 的支集是闭的。
\item 对 $x \in X$，有
$$
x \in \text{Supp}(\mathcal{F})
\Leftrightarrow
\mathcal{F}_x \not = 0
\Leftrightarrow
\mathcal{F}_x \otimes_{\mathcal{O}_{X, x}} \kappa(x) \not = 0.
$$
\item 对任意概形态射 $f : Y \to X$，拉回
$f^*\mathcal{F}$ 也为有限型，并且
$\text{Supp}(f^*\mathcal{F}) = f^{-1}(\text{Supp}(\mathcal{F}))$。
\end{enumerate}
\end{lemma}

\begin{proof}
部分 (1) 是《模》引理
\ref{modules-lemma-support-finite-type-closed} 的重新表述。
也可结合引理 \ref{morphisms-lemma-support-affine-open}、《性质》引理
\ref{properties-lemma-finite-type-module} 以及《代数》引理
\ref{algebra-lemma-support-closed} 看出这一点。(2) 中第一个等价是
支集的定义，第二个等价来自 Nakayama 引理；见《代数》引理
\ref{algebra-lemma-NAK}。设 $f : Y \to X$ 为概形态射。注意，
$f^*\mathcal{F}$ 由《模》引理
\ref{modules-lemma-pullback-finite-type} 可知是有限型。
对于最后一个断言，设 $y \in Y$ 的像为 $x \in X$。回顾
$$
(f^*\mathcal{F})_y =
\mathcal{F}_x \otimes_{\mathcal{O}_{X, x}} \mathcal{O}_{Y, y},
$$
见《层》引理 \ref{sheaves-lemma-stalk-pullback-modules}。
因此，$(f^*\mathcal{F})_y \otimes \kappa(y)$ 非零，当且仅当
$\mathcal{F}_x \otimes \kappa(x)$ 非零。由 (2)，这推出
$x \in \text{Supp}(\mathcal{F})$ 当且仅当
$y \in \text{Supp}(f^*\mathcal{F})$，这正是断言 (3) 的内容。
\end{proof}

\begin{lemma}
\label{morphisms-lemma-scheme-theoretic-support}
设 $\mathcal{F}$ 为概形 $X$ 上的有限型准凝聚模。存在最小闭子概形
$i : Z \to X$，使得存在准凝聚 $\mathcal{O}_Z$-模
$\mathcal{G}$，满足 $i_*\mathcal{G} \cong \mathcal{F}$。此外：
\begin{enumerate}
\item 若 $\Spec(A) \subset X$ 是任意仿射开集，并且
$\mathcal{F}|_{\Spec(A)} = \widetilde{M}$，则
$Z \cap \Spec(A) = \Spec(A/I)$，其中 $I = \text{Ann}_A(M)$。
\item 准凝聚层 $\mathcal{G}$ 在唯一同构意义下唯一。
\item 准凝聚层 $\mathcal{G}$ 是有限型的。
\item $\mathcal{G}$ 与 $\mathcal{F}$ 的支集都是 $Z$。
\end{enumerate}
\end{lemma}

\begin{proof}
假设 $i' : Z' \to X$ 是满足引理中仿射开集上描述的闭子概形。
那么由引理 \ref{morphisms-lemma-i-star-equivalence} 可知，有
$\mathcal{F} \cong i'_*\mathcal{G}'$，其中 $\mathcal{G}'$ 是
$Z'$ 上某个唯一的准凝聚层。此外，显然
$Z'$ 是具有此性质的最小闭子概形（仍由同一引理）。
最后，使用《性质》引理
\ref{properties-lemma-finite-type-module} 与《代数》引理
\ref{algebra-lemma-finite-over-subring} 可知 $\mathcal{G}'$ 是有限型的。
由《代数》引理 \ref{algebra-lemma-support-closed}，有
$\text{Supp}(\mathcal{G}') = Z$。因此，为证明本引理，只需说明
(1) 中的刻画确实定义了闭子概形。为此，又只需证明给定规则产生
准凝聚理想层；见引理
\ref{morphisms-lemma-closed-immersion-bijection-ideals}。
这归结为下列代数事实：若 $A$ 为环，$f \in A$，且 $M$ 为有限
$A$-模，则 $\text{Ann}_A(M)_f = \text{Ann}_{A_f}(M_f)$。
证明从略。
\end{proof}

\begin{definition}
\label{morphisms-definition-scheme-theoretic-support}
设 $X$ 为概形。设 $\mathcal{F}$ 为有限型准凝聚
$\mathcal{O}_X$-模。$\mathcal{F}$ 的{\it 概形论支集}是引理
\ref{morphisms-lemma-scheme-theoretic-support} 中构造的闭子概形
$Z \subset X$。
\end{definition}

\noindent
在此情形中，我们常把 $\mathcal{F}$ 看作 $Z$ 上的有限型准凝聚层
（通过引理 \ref{morphisms-lemma-i-star-equivalence} 的范畴等价）。










\section{概形论像}
\label{morphisms-section-scheme-theoretic-image}

\noindent
注意：本节中的一些材料处于极其一般的层次，其表现可能不同于预期。

\begin{lemma}
\label{morphisms-lemma-scheme-theoretic-image}
设 $f : X \to Y$ 为概形态射。存在闭子概形 $Z \subset Y$，
使得 $f$ 经 $Z$ 分解，并且对任意其他闭子概形 $Z' \subset Y$，
若 $f$ 经 $Z'$ 分解，则有 $Z \subset Z'$。
\end{lemma}

\begin{proof}
设 $\mathcal{I} = \Ker(\mathcal{O}_Y \to f_*\mathcal{O}_X)$。
若 $\mathcal{I}$ 准凝聚，就取 $Z$ 为由 $\mathcal{I}$ 决定的闭子概形；
见引理 \ref{morphisms-lemma-closed-immersion-bijection-ideals}。
由引理 \ref{morphisms-lemma-closed-immersion-ideals}，这样确实可行。
一般而言，同一引理要求证明存在最大准凝聚理想层
$\mathcal{I}'$，它包含于 $\mathcal{I}$。这由引理
\ref{morphisms-lemma-largest-quasi-coherent-subsheaf} 得出。
\end{proof}

\begin{definition}
\label{morphisms-definition-scheme-theoretic-image}
设 $f : X \to Y$ 为概形态射。$f$ 的{\it 概形论像}是最小闭子概形
$Z \subset Y$，且 $f$ 经它分解；见上面的引理
\ref{morphisms-lemma-scheme-theoretic-image}。
\end{definition}

\noindent
对概形态射 $f : X \to Y$ 及其概形论像 $Z$，我们常以
$f : X \to Z$ 表示 $f$ 经其概形论像的分解。若态射 $f$ 不拟紧，
则一般而言：
\begin{enumerate}
\item 集合论包含 $\overline{f(X)} \subset Z$ 不是等式，即
$f(X) \subset Z$ 不是稠密子集；并且
\item 概形论像的构造不与到 $Y$ 的开子概形的限制可交换。
\end{enumerate}
在《例》第 \ref{examples-section-scheme-theoretic-image} 节中，
读者可找到同时表现这两个现象的例子。这些现象甚至可对浸入出现；
见《例》第 \ref{examples-section-strange-immersion} 节。
不过，下一个引理表明，态射拟紧时两个灾难都会避免。

\begin{lemma}
\label{morphisms-lemma-quasi-compact-scheme-theoretic-image}
设 $f : X \to Y$ 为概形态射。设 $Z \subset Y$ 为 $f$ 的概形论像。
若 $f$ 拟紧，则：
\begin{enumerate}
\item 理想层
$\mathcal{I} = \Ker(\mathcal{O}_Y \to f_*\mathcal{O}_X)$
是准凝聚的；
\item 概形论像 $Z$ 是由 $\mathcal{I}$ 决定的闭子概形；
\item 对任意开集 $U \subset Y$，态射
$f|_{f^{-1}(U)} : f^{-1}(U) \to U$ 的概形论像等于
$Z \cap U$；并且
\item 像 $f(X) \subset Z$ 是 $Z$ 的稠密子集，换言之，态射
$X \to Z$ 是支配的（见定义 \ref{morphisms-definition-dominant}）。
\end{enumerate}
\end{lemma}

\begin{proof}
部分 (4) 由部分 (3) 得出。为证明 (3)，只需证明 (1)，因为
$\mathcal{I}$ 的形成与限制到 $Y$ 的开子概形可交换。
而若 (1) 成立，则我们已在引理
\ref{morphisms-lemma-scheme-theoretic-image} 的证明中证明 (2)。
因此只需证明 $\mathcal{I}$ 准凝聚。由于准凝聚性是局部性质，
可以假设 $Y$ 仿射。由于 $f$ 拟紧，可取有限仿射开覆盖
$X = \bigcup_{i = 1, \ldots, n} U_i$。以 $f'$ 表示复合
$$
X' = \coprod U_i \longrightarrow X \longrightarrow Y.
$$
于是 $f_*\mathcal{O}_X$ 是 $f'_*\mathcal{O}_{X'}$ 的子层，因而
$\mathcal{I} = \Ker(\mathcal{O}_Y \to f'_*\mathcal{O}_{X'})$。
由《概形》引理 \ref{schemes-lemma-push-forward-quasi-coherent}，
层 $f'_*\mathcal{O}_{X'}$ 在 $Y$ 上准凝聚。因此得证。
\end{proof}

\begin{example}
\label{morphisms-example-scheme-theoretic-image}
若 $A \to B$ 是核为 $I$ 的环映射，则
$\Spec(B) \to \Spec(A)$ 的概形论像是闭子概形
$\Spec(A/I)$，它位于 $\Spec(A)$ 中。这由引理
\ref{morphisms-lemma-quasi-compact-scheme-theoretic-image} 得出。
\end{example}

\noindent
若态射拟紧，则概形论像只添加像中各点的特殊化点。

\begin{lemma}
\label{morphisms-lemma-reach-points-scheme-theoretic-image}
设 $f : X \to Y$ 为拟紧态射。设 $Z$ 为 $f$ 的概形论像。
设 $z \in Z$\footnote{由引理
\ref{morphisms-lemma-quasi-compact-scheme-theoretic-image}，从集合论看，
$Z$ 与 $f(X)$ 在 $Y$ 中的闭包相同。}。存在赋值环 $A$，其分式域为
$K$，并且存在交换图
$$
\xymatrix{
\Spec(K) \ar[rr] \ar[d] & & X \ar[d] \ar[ld] \\
\Spec(A) \ar[r] & Z \ar[r] & Y
}
$$
使得 $\Spec(A)$ 的闭点映到 $z$。特别地，$Z$ 中任意一点都是
$f(X)$ 中某一点的特殊化。
\end{lemma}

\begin{proof}
设 $z \in \Spec(R) = V \subset Y$ 为 $z$ 的仿射开邻域。
由引理 \ref{morphisms-lemma-quasi-compact-scheme-theoretic-image}，
交 $Z \cap V$ 是 $f^{-1}(V) \to V$ 的概形论像。因此可将 $Y$
替换为 $V$，并假设 $Y = \Spec(R)$ 仿射。在此情形中，$X$ 拟紧，
因为 $f$ 拟紧。设 $X = U_1 \cup \ldots \cup U_n$ 为有限仿射开覆盖。
写成 $U_i = \Spec(A_i)$。令
$I = \Ker(R \to A_1 \times \ldots \times A_n)$。
再次由引理 \ref{morphisms-lemma-quasi-compact-scheme-theoretic-image} 可知，
$Z$ 对应于闭子概形 $\Spec(R/I)$，它位于 $Y$ 中。若
$\mathfrak p \subset R$ 为与 $z$ 对应的素理想，则
$I_{\mathfrak p} \subset R_{\mathfrak p}$ 不是等式。因此
（由于局部化正合，见《代数》命题
\ref{algebra-proposition-localization-exact}）可知
$R_{\mathfrak p} \to
(A_1)_{\mathfrak p} \times \ldots \times (A_n)_{\mathfrak p}$
不是零映射。所以某个环 $(A_i)_{\mathfrak p}$ 非零。
因此存在 $i$ 以及素理想 $\mathfrak q_i \subset A_i$，它位于某个
素理想 $\mathfrak p_i \subset \mathfrak p$ 之上。由《代数》引理
\ref{algebra-lemma-dominate}，可选择赋值环
$A \subset K = \kappa(\mathfrak q_i)$，它支配局部环
$R_{\mathfrak p}/\mathfrak p_iR_{\mathfrak p} \subset \kappa(\mathfrak q_i)$。
这给出所需交换图。若干细节从略。
\end{proof}

\begin{lemma}
\label{morphisms-lemma-factor-factor}
设
$$
\xymatrix{
X_1 \ar[d] \ar[r]_{f_1} & Y_1 \ar[d] \\
X_2 \ar[r]^{f_2} & Y_2
}
$$
为概形的交换图。设 $Z_i \subset Y_i$（$i = 1, 2$）为 $f_i$
的概形论像。则态射 $Y_1 \to Y_2$ 诱导态射 $Z_1 \to Z_2$ 以及交换图
$$
\xymatrix{
X_1 \ar[r] \ar[d] & Z_1 \ar[d] \ar[r] & Y_1 \ar[d] \\
X_2 \ar[r] & Z_2 \ar[r] & Y_2
}
$$
\end{lemma}

\begin{proof}
$Z_2$ 在 $Y_1$ 中的概形论逆像是 $Y_1$ 的闭子概形，并且 $f_1$
经它分解。因此 $Z_1$ 包含于其中。引理得证。
\end{proof}

\begin{lemma}
\label{morphisms-lemma-scheme-theoretic-image-reduced}
设 $f : X \to Y$ 为概形态射。若 $X$ 约化，则 $f$ 的概形论像是
$\overline{f(X)}$ 上的约化诱导概形结构。
\end{lemma}

\begin{proof}
这是因为 $\overline{f(X)}$ 上的约化诱导概形结构显然是 $Y$ 的
最小闭子概形，并且 $f$ 经它分解；见《概形》引理
\ref{schemes-lemma-map-into-reduction}。
\end{proof}

\begin{lemma}
\label{morphisms-lemma-scheme-theoretic-image-of-partial-section}
设 $f : X \to Y$ 为概形的分离态射。设 $V \subset Y$ 为逆拟紧开集。
设 $s : V \to X$ 为满足 $f \circ s = \text{id}_V$ 的态射。
设 $Y'$ 为 $s$ 的概形论像。则 $Y' \to Y$ 在 $V$ 上为同构。
\end{lemma}

\begin{proof}
假设 $V$ 在 $Y$ 中逆拟紧（《拓扑》定义
\ref{topology-definition-quasi-compact}），这意味着
$V \to Y$ 是拟紧态射。由《概形》引理
\ref{schemes-lemma-quasi-compact-permanence}，态射
$s : V \to X$ 拟紧。因此，由引理
\ref{morphisms-lemma-quasi-compact-scheme-theoretic-image}，概形论像
$Y'$（它是 $s$ 的概形论像）的构造与到开集的限制可交换。特别地，
$Y' \cap f^{-1}(V)$ 是分离态射 $f^{-1}(V) \to V$ 的某个截面的
概形论像。由于分离态射的截面是闭浸入
（《概形》引理 \ref{schemes-lemma-section-immersion}），
可断定 $Y' \cap f^{-1}(V) \to V$ 如所要求地为同构。
\end{proof}






\section{概形论闭包与稠密性}
\label{morphisms-section-scheme-theoretic-closure}

\noindent
我们采用 \cite[IV, Definition 11.10.2]{EGA} 中的下列定义。

\begin{definition}
\label{morphisms-definition-scheme-theoretically-dense}
设 $X$ 为概形。设 $U \subset X$ 为开子概形。
\begin{enumerate}
\item 态射 $U \to X$ 的概形论像称为{\it $U$ 在 $X$ 中的概形论闭包}。
\item 称 $U$ 在 $X$ 中{\it 概形论稠密}，如果对每个开集
$V \subset X$，$U \cap V$ 在 $V$ 中的概形论闭包等于 $V$。
\end{enumerate}
\end{definition}

\noindent
采用此定义，$U$ 在 $X$ 中概形论稠密{\bf 并不}等价于 $U$ 的概形论
闭包为 $X$；见例
\ref{morphisms-example-scheme-theretically-dense-not-dense}。这略显不雅；
不过可参见下面的引理
\ref{morphisms-lemma-scheme-theoretically-dense-quasi-compact} 与
\ref{morphisms-lemma-reduced-scheme-theoretically-dense}。另一方面，采用此定义，
$U$ 在 $X$ 中概形论稠密，当且仅当对每个开集 $V \subset X$，环映射
$\mathcal{O}_X(V) \to \mathcal{O}_X(U \cap V)$ 都是单射；见下面的引理
\ref{morphisms-lemma-characterize-scheme-theoretically-dense}。特别地，概形论稠密
蕴含稠密，这是令人满意的。

\begin{example}
\label{morphisms-example-scheme-theretically-dense-not-dense}
下面给出一个概形论闭包为 $X$，但其底拓扑空间中并不稠密的例子。
设 $k$ 为域。令
$A = k[x, z_1, z_2, \ldots]/(x^n z_n)$，并令
$I = (z_1, z_2, \ldots) \subset A$。考虑仿射概形
$X = \Spec(A)$ 与开子概形 $U = X \setminus V(I)$。
由于 $A \to \prod_n A_{z_n}$ 单射，可知 $U$ 的概形论闭包是 $X$。
考虑态射 $X \to \Spec(k[x])$。此态射满射
（令所有 $z_n = 0$ 即可看出）。但此态射到 $U$ 的限制不满射，
因为它映到点 $x = 0$。因此，$U$ 不可能在 $X$ 中拓扑稠密。
\end{example}

\begin{lemma}
\label{morphisms-lemma-scheme-theoretically-dense-quasi-compact}
设 $X$ 为概形。设 $U \subset X$ 为开子概形。
若包含态射 $U \to X$ 拟紧，则 $U$ 在 $X$ 中概形论稠密，当且仅当
$U$ 在 $X$ 中的概形论闭包是 $X$。
\end{lemma}

\begin{proof}
这由引理 \ref{morphisms-lemma-quasi-compact-scheme-theoretic-image} 的部分 (3)
得出。
\end{proof}

\begin{example}
\label{morphisms-example-scheme-theoretic-closure}
设 $A$ 为环且 $X = \Spec(A)$。设 $f_1, \ldots, f_n \in A$，并设
$U = D(f_1) \cup \ldots \cup D(f_n)$。令
$I = \Ker(A \to \prod A_{f_i})$。则 $U$ 在 $X$ 中的概形论闭包是
闭子概形 $\Spec(A/I)$，它位于 $X$ 中。注意，$U \to X$ 拟紧。因此由引理
\ref{morphisms-lemma-scheme-theoretically-dense-quasi-compact} 可知，$U$ 在
$X$ 中概形论稠密，当且仅当 $I = 0$。
\end{example}

\begin{lemma}
\label{morphisms-lemma-characterize-scheme-theoretically-dense}
设 $j : U \to X$ 为概形的开浸入。则 $U$ 在 $X$ 中概形论稠密，
当且仅当 $\mathcal{O}_X \to j_*\mathcal{O}_U$ 单射。
\end{lemma}

\begin{proof}
若 $\mathcal{O}_X \to j_*\mathcal{O}_U$ 单射，则限制到任意开集
$V$ 后仍然如此，其中该开集是 $X$ 的开集。因此，$U \cap V$ 在 $V$
中的概形论闭包等于 $V$；
见引理 \ref{morphisms-lemma-scheme-theoretic-image} 的证明。反过来，假设
$U \cap V$ 的概形论闭包等于 $V$，并且对所有开集 $V$ 都如此。假设
$\mathcal{O}_X \to j_*\mathcal{O}_U$ 不单射。则可找到仿射开集
$\Spec(A) = V \subset X$ 与非零元素 $f \in A$，使得 $f$ 在
$\Gamma(V \cap U, \mathcal{O}_X)$ 中的像为零。在此情形中，
$V \cap U$ 在 $V$ 中的概形论闭包显然包含于
$\Spec(A/(f))$，矛盾。
\end{proof}

\begin{lemma}
\label{morphisms-lemma-intersection-scheme-theoretically-dense}
设 $X$ 为概形。若 $U$、$V$ 是 $X$ 的概形论稠密开子概形，
则 $U \cap V$ 也是。
\end{lemma}

\begin{proof}
设 $W \subset X$ 为任意开集。考虑映射
$\mathcal{O}_X(W) \to \mathcal{O}_X(W \cap V)
\to \mathcal{O}_X(W \cap V \cap U)$。
由引理 \ref{morphisms-lemma-characterize-scheme-theoretically-dense}，
两个映射都单射。因此复合单射。故再次由引理
\ref{morphisms-lemma-characterize-scheme-theoretically-dense}，
$U \cap V$ 在 $X$ 中概形论稠密。
\end{proof}

\begin{lemma}
\label{morphisms-lemma-quasi-compact-immersion}
设 $h : Z \to X$ 为浸入。假设 $h$ 拟紧，或 $Z$ 约化。
设 $\overline{Z} \subset X$ 为 $h$ 的概形论像。则态射
$Z \to \overline{Z}$ 是开浸入，它将 $Z$ 认同为 $\overline{Z}$ 的
概形论稠密开子概形。此外，$Z$ 在 $\overline{Z}$ 中拓扑稠密。
\end{lemma}

\begin{proof}
由引理 \ref{morphisms-lemma-factor-quasi-compact-immersion} 或引理
\ref{morphisms-lemma-factor-reduced-immersion}，可将 $Z \to X$ 分解为
$Z \to \overline{Z}_1 \to X$，其中 $Z \to \overline{Z}_1$ 为开浸入，
而 $\overline{Z}_1 \to X$ 为闭浸入。另一方面，设
$Z \to \overline{Z} \subset X$ 为 $Z \to X$ 的概形论像。
可断定 $\overline{Z} \subset \overline{Z}_1$。由于 $Z$ 是
$\overline{Z}_1$ 的开子概形，故 $Z$ 也是 $\overline{Z}$ 的
开子概形。在 $Z$ 约化的情形中，已知
$Z \subset \overline{Z}_1$ 拓扑稠密；这是由引理
\ref{morphisms-lemma-factor-reduced-immersion} 的证明中对
$\overline{Z}_1$ 的构造所得。
因此，$\overline{Z}_1$ 与 $\overline{Z}$ 具有相同的底拓扑空间。
于是 $\overline{Z} \subset \overline{Z}_1$ 是到约化概形的闭浸入，
并在底拓扑空间上诱导双射，故它是同构。
在 $Z \to X$ 拟紧的情形中，论证如下：$Z$ 在 $\overline{Z}$ 中
概形论稠密这一断言由引理
\ref{morphisms-lemma-quasi-compact-scheme-theoretic-image} 的部分 (3) 得出。
最后一个断言由引理
\ref{morphisms-lemma-quasi-compact-scheme-theoretic-image} 的部分 (4) 得出。
\end{proof}

\begin{lemma}
\label{morphisms-lemma-reduced-scheme-theoretically-dense}
设 $X$ 为约化概形，并设 $U \subset X$ 为开子概形。下列条件等价：
\begin{enumerate}
\item $U$ 在 $X$ 中拓扑稠密；
\item $U$ 在 $X$ 中的概形论闭包是 $X$；并且
\item $U$ 在 $X$ 中概形论稠密。
\end{enumerate}
\end{lemma}

\begin{proof}
这由引理 \ref{morphisms-lemma-quasi-compact-immersion} 以及下述事实得出：
若闭子概形 $Z$ 位于 $X$ 中，且其底拓扑空间等于 $X$，则作为概形
它必等于 $X$。
\end{proof}

\begin{lemma}
\label{morphisms-lemma-reduced-subscheme-closure}
设 $X$ 为概形，并设 $U \subset X$ 为约化开子概形。下列条件等价：
\begin{enumerate}
\item $U$ 在 $X$ 中的概形论闭包是 $X$；以及
\item $U$ 在 $X$ 中概形论稠密。
\end{enumerate}
若这些条件成立，则 $X$ 是约化概形。
\end{lemma}

\begin{proof}
这由引理 \ref{morphisms-lemma-quasi-compact-immersion} 以及下述事实得出：
由引理 \ref{morphisms-lemma-scheme-theoretic-image-reduced}，$U$ 在 $X$ 中的
概形论闭包约化。
\end{proof}


\begin{lemma}
\label{morphisms-lemma-equality-of-morphisms}
设 $S$ 为概形。设 $X$、$Y$ 为 $S$ 上的概形。
设 $f, g : X \to Y$ 为 $S$ 上的概形态射。设 $U \subset X$
为满足 $f|_U = g|_U$ 的开子概形。若 $U$ 在 $X$ 中的概形论闭包
是 $X$，且 $Y \to S$ 分离，则 $f = g$。
\end{lemma}

\begin{proof}
这由定义以及《概形》引理
\ref{schemes-lemma-where-are-they-equal} 得出。
\end{proof}




\section{支配态射}
\label{morphisms-section-dominant}

\noindent
如果习惯了簇的支配态射概念，可能会发现概形态射为支配态射的定义
与预期略有不同。

\begin{definition}
\label{morphisms-definition-dominant}
概形态射 $f : X \to S$ 称为{\it 支配的}，如果 $f$ 的像是 $S$
的稠密子集。
\end{definition}

\noindent
例如，若 $k$ 是无限域，且 $\lambda_1, \lambda_2, \ldots$ 是 $k$
中两两不同元素的可数族，则态射
$$
\coprod\nolimits_{i = 1, 2, \ldots } \Spec(k)
\longrightarrow
\Spec(k[x])
$$
是支配的，其中第 $i$ 个因子映到点 $x = \lambda_i$。

\begin{lemma}
\label{morphisms-lemma-generic-points-in-image-dominant}
设 $f : X \to S$ 为概形态射。若 $S$ 的每个不可约分支的每个泛点
都属于 $f$ 的像，则 $f$ 支配。
\end{lemma}

\begin{proof}
这是一个拓扑事实，直接来自以下两点：概形的底拓扑空间清醒，
见《概形》引理 \ref{schemes-lemma-scheme-sober}；并且 $S$ 的每个点
都包含在 $S$ 的某个不可约分支中，见《拓扑》引理
\ref{topology-lemma-irreducible}。
\end{proof}

\noindent
若态射拟紧，则只有在目标泛点属于像时态射才支配这一预期确实成立。

\begin{lemma}
\label{morphisms-lemma-quasi-compact-dominant}
\begin{slogan}
像包含泛点的态射是支配的
\end{slogan}
设 $f : X \to S$ 为概形的拟紧态射。则 $f$ 支配，当且仅当对每个
不可约分支 $Z \subset S$，$Z$ 的泛点属于 $f$ 的像。
\end{lemma}

\begin{proof}
设 $V \subset S$ 为仿射开集。由于 $f$ 拟紧，可以选择有限多个仿射
开集 $U_i \subset f^{-1}(V)$（$i = 1, \ldots, n$）覆盖
$f^{-1}(V)$。考虑仿射概形的态射
$$
f' :
\coprod\nolimits_{i = 1, \ldots, n} U_i
\longrightarrow
V.
$$
仿射概形的不交并仍仿射；见《概形》引理
\ref{schemes-lemma-disjoint-union-affines}。$V$ 的不可约分支的泛点，
恰好是 $S$ 中与 $V$ 相交的那些不可约分支的泛点。此外，$f$ 支配
当且仅当 $f'$ 支配，不论上面如何选择 $V, n, U_i$。
因此已将本引理归结为仿射概形态射的情形。仿射情形即《代数》引理
\ref{algebra-lemma-image-dense-generic-points}。
\end{proof}

\begin{lemma}
\label{morphisms-lemma-base-change-dominant}
设 $f : X \to S$ 为概形的拟紧支配态射。设 $g : S' \to S$
为概形态射，并设 $f' : X' \to S'$ 为 $f$ 沿 $g$ 的基变换。
若一般化可沿 $g$ 提升，则 $f'$ 支配。
\end{lemma}

\noindent
我们将看到，若 $g$ 开（引理 \ref{morphisms-lemma-open-generizing}）或平坦
（引理 \ref{morphisms-lemma-generalizations-lift-flat}），则一般化可沿 $g$ 提升。

\begin{proof}
注意，由《概形》引理
\ref{schemes-lemma-quasi-compact-preserved-base-change}，$f'$ 拟紧。
设 $\eta' \in S'$ 为 $S'$ 某个不可约分支的泛点。
若一般化可沿 $g$ 提升，则 $\eta = g(\eta')$ 是 $S$ 某个不可约分支
的泛点。由引理 \ref{morphisms-lemma-quasi-compact-dominant} 可知，$\eta$
属于 $f$ 的像。于是由《概形》引理
\ref{schemes-lemma-points-fibre-product}，$\eta'$ 属于 $f'$ 的像。
再由引理 \ref{morphisms-lemma-quasi-compact-dominant}，$f'$ 支配。
\end{proof}

\noindent
其中一部分对任意支配态射（不必拟紧）仍成立。

\begin{lemma}
\label{morphisms-lemma-open-base-change-dominant}
设 $f : X \to S$ 为概形的支配态射。设 $g : S' \to S$
为概形的开态射，并设 $f' : X' \to S'$ 为 $f$ 沿 $g$ 的基变换。
则 $f'$ 支配。
\end{lemma}

\begin{proof}
若 $f'$ 不支配，则 $f'(X')$ 包含于某个闭子集 $B$ 中，且该子集
不等于 $S'$。于是 $S'-B$ 是 $S'$ 中非空的开集。由于 $g$ 开，
$g(S'-B)$ 是 $S$ 中非空的开集。由于 $f(X)$ 在 $S$ 中稠密，
存在点 $x\in X$，使得 $f(x) \in g(S'-B)$。令
$W = \Spec(\kappa (f(x)))$，并用 $X_W$、$S_W$、$S'_W$ 表示这些概形
从 $S$ 到 $W$ 的拉回。则 $X_W$ 与 $(S'-B)_W$ 都非空，因而
$X_W \times_W (S'-B)_W$ 非空，因为 $W$ 是域的谱。
这与 $f'(X \times_S S')$ 包含于 $B \subset S'$ 矛盾。
所以 $f'$ 确实支配。
\end{proof}

\begin{lemma}
\label{morphisms-lemma-quasi-compact-generic-point-not-in-image}
设 $f : X \to S$ 为概形的拟紧态射。设 $\eta \in S$ 为 $S$
某个不可约分支的泛点。若 $\eta \not \in f(X)$，则存在开邻域
$V \subset S$，它是 $\eta$ 的邻域，并使得
$f^{-1}(V) = \emptyset$。
\end{lemma}

\begin{proof}
设 $Z \subset S$ 为 $f$ 的概形论像。需要证明
$\eta \not \in Z$。这由引理
\ref{morphisms-lemma-reach-points-scheme-theoretic-image} 得出，也可如下看出。
由引理 \ref{morphisms-lemma-quasi-compact-scheme-theoretic-image}，态射
$X \to Z$ 支配；由引理 \ref{morphisms-lemma-quasi-compact-dominant}，
这意味着 $Z$ 的所有不可约分支的所有泛点都属于 $X \to Z$ 的像。
由假设可知 $\eta \not \in Z$，因为若 $\eta$ 属于 $Z$，它将是
$Z$ 某个不可约分支的泛点。
\end{proof}

\noindent
还有一种情形，其中支配等同于所有不可约分支的泛点都属于像。

\begin{lemma}
\label{morphisms-lemma-dominant-finite-number-irreducible-components}
设 $f : X \to S$ 为概形态射。假设 $X$ 只有有限多个不可约分支。
则 $f$ 支配（当且）仅当对每个不可约分支 $Z \subset S$，$Z$ 的泛点
属于 $f$ 的像。若如此，则 $S$ 也只有有限多个不可约分支。
\end{lemma}

\begin{proof}
假设 $f$ 支配。设 $X = Z_1 \cup Z_2 \cup \ldots \cup Z_n$ 为 $X$
分解成不可约分支的分解。设 $\xi_i \in Z_i$ 为其泛点，因而
$Z_i = \overline{\{\xi_i\}}$。注意，$f(Z_i)$ 是 $S$ 的不可约子集。
因此
$$
S = \overline{f(X)} = \bigcup \overline{f(Z_i)} =
\bigcup \overline{\{f(\xi_i)\}}
$$
是有限多个不可约子集的并，而这些子集的泛点属于 $f$ 的像。
引理得证。
\end{proof}

\begin{lemma}
\label{morphisms-lemma-dominant-between-integral}
设 $f : X \to Y$ 为整概形之间的态射。下列条件等价：
\begin{enumerate}
\item $f$ 支配；
\item $f$ 将 $X$ 的泛点映到 $Y$ 的泛点；
\item 对某些非空仿射开集 $U \subset X$ 与 $V \subset Y$，若
$f(U) \subset V$，则环映射
$\mathcal{O}_Y(V) \to \mathcal{O}_X(U)$ 单射；
\item 对所有非空仿射开集 $U \subset X$ 与 $V \subset Y$，若
$f(U) \subset V$，则环映射
$\mathcal{O}_Y(V) \to \mathcal{O}_X(U)$ 单射；
\item 对某个 $x \in X$，其像为 $y = f(x) \in Y$，局部环映射
$\mathcal{O}_{Y, y} \to \mathcal{O}_{X, x}$ 单射；并且
\item 对所有 $x \in X$，其像为 $y = f(x) \in Y$，局部环映射
$\mathcal{O}_{Y, y} \to \mathcal{O}_{X, x}$ 单射。
\end{enumerate}
\end{lemma}

\begin{proof}
(1) 与 (2) 的等价性由引理
\ref{morphisms-lemma-dominant-finite-number-irreducible-components} 得出。
设 $U \subset X$ 与 $V \subset Y$ 为满足 $f(U) \subset V$ 的非空
仿射开集。回顾，环 $A = \mathcal{O}_X(U)$ 与
$B = \mathcal{O}_Y(V)$ 都是整环。态射 $f|_U : U \to V$ 对应于环映射
$\varphi : B \to A$。$X$ 与 $Y$ 的泛点分别对应于素理想
$(0) \subset A$ 与 $(0) \subset B$。因此，(2) 等价于条件
$(0) = \varphi^{-1}((0))$，即等价于 $\varphi$ 单射。
这样便看出 (2)、(3) 与 (4) 等价。类似地，给定 (5) 中的 $x$ 与
$y$，局部环 $\mathcal{O}_{X, x}$ 与 $\mathcal{O}_{Y, y}$ 都是整环，
而素理想 $(0) \subset \mathcal{O}_{X, x}$ 与
$(0) \subset \mathcal{O}_{Y, y}$ 对应于 $X$ 与 $Y$ 的泛点
（通过将 $x$ 处局部环的谱认同为特殊化到 $x$ 的点集；见《概形》引理
\ref{schemes-lemma-specialize-points}）。因此可以完全按照上面的方式
论证，得到 (2)、(5) 与 (6) 等价。
\end{proof}






\section{满射态射}
\label{morphisms-section-surjective}

\begin{definition}
\label{morphisms-definition-surjective}
若概形态射在底拓扑空间上满射，则称它{\it 满射}。
\end{definition}

\begin{lemma}
\label{morphisms-lemma-composition-surjective}
满射态射的复合满射。
\end{lemma}

\begin{proof}
从略。
\end{proof}

\begin{lemma}
\label{morphisms-lemma-when-point-maps-to-pair}
设 $X$ 与 $Y$ 为基概形 $S$ 上的概形。给定点 $x \in X$ 与
$y \in Y$，存在 $X \times_S Y$ 中的点，在投影下分别映到 $x$ 与
$y$，当且仅当 $x$ 与 $y$ 位于 $S$ 的同一点之上。
\end{lemma}

\begin{proof}
该条件显然必要，反向则来自《概形》引理
\ref{schemes-lemma-points-fibre-product} 的证明。
\end{proof}

\begin{lemma}
\label{morphisms-lemma-base-change-surjective}
满射态射的基变换满射。
\end{lemma}

\begin{proof}
设 $f: X \to Y$ 为基概形 $S$ 上的概形态射。
若 $S' \to S$ 为概形态射，设 $p: X_{S'} \to X$ 与
$q: Y_{S'} \to Y$ 为典范投影。交换方块
$$
\xymatrix{
X_{S'} \ar[d]_{f_{S'}} \ar[r]_p & X \ar[d]^{f} \\
Y_{S'} \ar[r]^{q} & Y.
}
$$
将 $X_{S'}$ 认同为 $X \to Y$ 与 $Y_{S'} \to Y$ 的纤维积。
设 $Z$ 为 $X$ 的底拓扑空间的子集。则
$q^{-1}(f(Z)) = f_{S'}(p^{-1}(Z))$，因为
$y' \in q^{-1}(f(Z))$ 当且仅当 $q(y') = f(x)$ 对某个
$x \in Z$ 成立；由引理 \ref{morphisms-lemma-when-point-maps-to-pair}，
这又当且仅当存在 $x' \in X_{S'}$，使得
$f_{S'}(x') = y'$ 且 $p(x') = x$。特别地，取 $Z = X$，可知若
$f$ 满射，则其基变换 $f_{S'}: X_{S'} \to Y_{S'}$ 也满射。
\end{proof}

\begin{example}
\label{morphisms-example-injective-not-preserved-base-change}
双射性在基变换下不稳定，因而单射性也不稳定。例如考虑双射
$\Spec(\mathbf{C}) \to \Spec(\mathbf{R})$。基变换
$\Spec(\mathbf{C} \otimes_{\mathbf{R}} \mathbf{C}) \to
\Spec(\mathbf{C})$
不单射，因为存在同构
$\mathbf{C} \otimes_{\mathbf{R}} \mathbf{C} \cong \mathbf{C} \times \mathbf{C}$
（该分解来自幂等元
$\frac{1 \otimes 1 + i \otimes i}{2}$），因而
$\Spec(\mathbf{C} \otimes_{\mathbf{R}} \mathbf{C})$ 有两个点。
\end{example}

\begin{lemma}
\label{morphisms-lemma-surjection-from-quasi-compact}
设
$$
\xymatrix{
X \ar[rr]_f \ar[rd]_p & &
Y \ar[dl]^q \\
& Z
}
$$
为概形态射的交换图。若 $f$ 满射且 $p$ 拟紧，则 $q$ 拟紧。
\end{lemma}

\begin{proof}
设 $W \subset Z$ 为拟紧开集。由假设，$p^{-1}(W)$ 拟紧。
因此由《拓扑》引理 \ref{topology-lemma-image-quasi-compact}，
逆像 $q^{-1}(W) = f(p^{-1}(W))$ 也拟紧。引理得证。
\end{proof}

\section{根态射与泛单射态射}
\label{morphisms-section-radicial}

\noindent
本节定义概形态射何时为{\it 根态射}，以及何时为
{\it 泛单射}，并证明这两个概念等价。之所以同时引入二者，
是因为对代数空间而言也有相应概念，而它们未必总是等价。

\begin{definition}
\label{morphisms-definition-universally-injective}
设 $f : X \to S$ 为一个态射。
\begin{enumerate}
\item 称 $f$ 为{\it 泛单射}，如果且仅如果对概形的任意态射
$S' \to S$，基变换 $f' : X_{S'} \to S'$ 在底层拓扑空间上为单射。
\item 称 $f$ 为{\it 根态射}，如果 $f$ 作为拓扑空间之间的映射
为单射，且对每个 $x \in X$，域扩张
$\kappa(x)/\kappa(f(x))$ 都纯不可分。
\end{enumerate}
\end{definition}

\begin{lemma}
\label{morphisms-lemma-universally-injective}
\begin{reference}
\cite[Proposition 3.7.1]{EGA1-second}
\end{reference}
设 $f : X \to S$ 为概形态射。下列条件等价：
\begin{enumerate}
\item 对每个域 $K$，诱导映射
$\Mor(\Spec(K), X) \to \Mor(\Spec(K), S)$
都是单射。
\item 态射 $f$ 泛单射。
\item 态射 $f$ 是根态射。
\item 对角态射 $\Delta_{X/S} : X \longrightarrow X \times_S X$
满射。
\end{enumerate}
\end{lemma}

\begin{proof}
设 $K$ 为一个域，$s : \Spec(K) \to S$ 为一个态射。
给出一个态射 $x : \Spec(K) \to X$ 使得 $f \circ x = s$，
等价于给出投影 $X_K = \Spec(K) \times_S X \to \Spec(K)$
的一个截面，而这又等价于给出一个点 $x \in X_K$，
其剩余域为 $K$。因此 (2) 蕴含 (1)。

\medskip\noindent
反之，设 (1) 成立。假设 $x, x' \in X_{S'}$ 映到同一点
$s' \in S'$。选取域的交换图
$$
\xymatrix{
K & \kappa(x) \ar[l] \\
\kappa(x') \ar[u] & \kappa(s') \ar[l] \ar[u]
}
$$
由《概形》引理 \ref{schemes-lemma-characterize-points}，
得到两个态射 $a, a' : \Spec(K) \to X_{S'}$。其中一个对应于
点 $x$ 及嵌入 $\kappa(x) \subset K$，另一个对应于点 $x'$ 及嵌入
$\kappa(x') \subset K$。此外，$f' \circ a = f' \circ a'$。
条件 (1) 于是推出 $a$ 与 $a'$ 同 $X_{S'} \to X$ 的复合相等。
由于 $X_{S'}$ 是 $S'$ 与 $X$ 在 $S$ 上的纤维积，故
$a = a'$，从而 $x = x'$。因此 (1) 蕴含 (2)。

\medskip\noindent
若有两个不同的点 $x, x' \in X$ 映到同一个点 $s$，
则 (2) 不成立。若对某个 $s = f(x)$、$x \in X$，域扩张
$\kappa(x)/\kappa(s)$ 不是纯不可分的，则可以找到域扩张
$K/\kappa(s)$，使得 $\kappa(x)$ 有两个 $\kappa(s)$-同态到
$K$。由《概形》引理
\ref{schemes-lemma-characterize-points}，这意味着映射
$\Mor(\Spec(K), X) \to \Mor(\Spec(K), S)$
不是单射，因而 (1) 不成立。因此，等价条件 (1) 与 (2)
蕴含 $f$ 是根态射，即蕴含 (3)。

\medskip\noindent
假设 (3)。由《概形》引理
\ref{schemes-lemma-characterize-points}，一个态射
$\Spec(K) \to X$ 由一对 $(x, \kappa(x) \to K)$ 给出。
性质 (3) 恰好说明，把这一对 $(x, \kappa(x) \to K)$
对应到 $(s, \kappa(s) \to \kappa(x) \to K)$ 的映射是单射。
换言之，(1) 成立。至此已知 (1)、(2)、(3) 彼此等价。

\medskip\noindent
最后证明 (4) 与 (1)、(2)、(3) 等价。$X \times_S X$ 的一个点
由四元组 $(x_1, x_2, s, \mathfrak p)$ 给出，其中
$x_1, x_2 \in X$，$f(x_1) = f(x_2) = s$，并且
$\mathfrak p \subset \kappa(x_1) \otimes_{\kappa(s)} \kappa(x_2)$
是一个素理想；见《概形》引理
\ref{schemes-lemma-points-fibre-product}。若 $f$ 泛单射，
则在泛单射的定义中取 $S'=X$。由于 $\Delta_{X/S}$ 是单射态射
$X \times_S X  \longrightarrow X$ 的一个截面，它必为满射。
反之，若 $\Delta_{X/S}$ 满射，则总有 $x_1 = x_2 = x$，
且恰有一个这样的素理想 $\mathfrak p$；这说明
$\kappa(s) \subset \kappa(x)$ 纯不可分。因此 $f$ 是根态射。
或者，若 $\Delta_{X/S}$ 满射，则对任意 $S' \to S$，
其基变换 $\Delta_{X_{S'}/S'}$ 都满射，这推出 $f$ 泛单射。
引理证毕。
\end{proof}

\begin{lemma}
\label{morphisms-lemma-universally-injective-separated}
泛单射态射是分离的。
\end{lemma}

\begin{proof}
将引理 \ref{morphisms-lemma-universally-injective} 与以下事实结合即可：
$X \to S$ 分离，当且仅当 $\Delta_{X/S}$ 的像在
$X \times_S X$ 中闭；见《概形》定义
\ref{schemes-definition-separated} 及其后的讨论。
\end{proof}

\begin{lemma}
\label{morphisms-lemma-base-change-universally-injective}
泛单射态射的基变换仍泛单射。
\end{lemma}

\begin{proof}
这是形式的。
\end{proof}

\begin{lemma}
\label{morphisms-lemma-composition-universally-injective}
根态射的复合仍为根态射，因而等价的泛单射条件在复合下也保持。
\end{lemma}

\begin{proof}
略。
\end{proof}










\section{仿射态射}
\label{morphisms-section-affine}

\begin{definition}
\label{morphisms-definition-affine}
若概形态射 $f : X \to S$ 满足：$S$ 的每个仿射开集的逆像
都是 $X$ 的仿射开集，则称它为{\it 仿射的}。
\end{definition}

\begin{lemma}
\label{morphisms-lemma-affine-separated}
仿射态射是分离且拟紧的。
\end{lemma}

\begin{proof}
设 $f : X \to S$ 仿射。由《概形》引理
\ref{schemes-lemma-quasi-compact-affine} 立即得到拟紧性。
我们用《概形》引理 \ref{schemes-lemma-characterize-separated}
证明 $f$ 分离。设 $x_1, x_2 \in X$ 为 $X$ 中映到同一点
$s \in S$ 的两个点。任取仿射开集 $W \subset S$，使其包含 $s$。
依假设，$f^{-1}(W)$ 仿射。在上述引理中取
$U = V = f^{-1}(W)$ 即可。
\end{proof}

\begin{lemma}
\label{morphisms-lemma-characterize-affine}
\begin{reference}
\cite[II, Corollary 1.3.2]{EGA}
\end{reference}
设 $f : X \to S$ 为概形态射。下列条件等价：
\begin{enumerate}
\item 态射 $f$ 仿射。
\item 存在仿射开覆盖 $S = \bigcup W_j$，使每个
$f^{-1}(W_j)$ 都仿射。
\item 存在一个准凝聚 $\mathcal{O}_S$-代数层
$\mathcal{A}$，以及同构
$X \cong \underline{\Spec}_S(\mathcal{A})$；它是 $S$ 上概形的同构。记号见《构造》
第 \ref{constructions-section-spec} 节。
\end{enumerate}
此外，在这种情形下，$X = \underline{\Spec}_S(f_*\mathcal{O}_X)$。
\end{lemma}

\begin{proof}
(1) 蕴含 (2) 是显然的。

\medskip\noindent
设 $S = \bigcup_{j \in J} W_j$ 是仿射开覆盖，且每个
$f^{-1}(W_j)$ 都仿射。由《概形》引理
\ref{schemes-lemma-quasi-compact-affine}，$f$ 拟紧。由《概形》引理
\ref{schemes-lemma-characterize-quasi-separated}，态射 $f$ 拟分离。
因此由《概形》引理 \ref{schemes-lemma-push-forward-quasi-coherent}，
层 $\mathcal{A} = f_*\mathcal{O}_X$ 是准凝聚
$\mathcal{O}_S$-代数层。于是有概形
$g : Y = \underline{\Spec}_S(\mathcal{A}) \to S$，它在 $S$ 上。
恒等映射
$\text{id} : \mathcal{A} = f_*\mathcal{O}_X \to f_*\mathcal{O}_X$
由相对谱的定义给出态射 $can : X \to Y$，它在 $S$ 上；见《构造》
引理 \ref{constructions-lemma-canonical-morphism}。由假设和刚引用的引理，
限制 $can|_{f^{-1}(W_j)} : f^{-1}(W_j) \to g^{-1}(W_j)$
是同构。因此 $can$ 是同构。我们已经证明 (2) 蕴含 (3)。

\medskip\noindent
假设 (3)。由《构造》引理
\ref{constructions-lemma-spec-properties}，每个仿射开集的逆像都仿射，
所以按定义该态射仿射。
\end{proof}

\begin{remark}
\label{morphisms-remark-direct-argument}
也可如下直接证明上述引理 \ref{morphisms-lemma-characterize-affine} 中
(2) 蕴含 (1)。设 $S = \bigcup W_j$ 是仿射开覆盖，且每个
$f^{-1}(W_j)$ 都仿射。先如上述证明那样论证
$\mathcal{A} = f_*\mathcal{O}_X$ 准凝聚。
设 $\Spec(R) = V \subset S$ 为仿射开集。需要证明
$f^{-1}(V)$ 仿射。令
$A = \mathcal{A}(V) = f_*\mathcal{O}_X(V) = \mathcal{O}_X(f^{-1}(V))$。
由《概形》引理 \ref{schemes-lemma-morphism-into-affine}，存在
典范态射 $\psi : f^{-1}(V) \to \Spec(A)$，它在
$\Spec(R) = V$ 上。
由《概形》引理 \ref{schemes-lemma-good-subcover}，存在整数
$n \geq 0$、标准开覆盖
$V = \bigcup_{i = 1, \ldots, n} D(h_i)$（$h_i \in R$），
以及映射 $a : \{1, \ldots, n\} \to J$，使得每个 $D(h_i)$
也是仿射概形 $W_{a(i)}$ 的标准开集。仿射概形态射下标准开集的逆像
仍为标准开集，见《代数》引理 \ref{algebra-lemma-spec-functorial}。
因此 $f^{-1}(D(h_i))$ 是 $f^{-1}(W_{a(i)})$ 的标准开集，
特别地，$f^{-1}(D(h_i))$ 仿射。由于 $\mathcal{A}$ 准凝聚，
$A_{h_i} = \mathcal{A}(D(h_i)) = \mathcal{O}_X(f^{-1}(D(h_i)))$，
所以 $f^{-1}(D(h_i))$ 是 $A_{h_i}$ 的谱。由此，态射 $\psi$
把开集 $f^{-1}(D(h_i))$ 与开集 $\Spec(A_{h_i})$ 同构起来，
后者是 $\Spec(A)$ 的开集。由于 $f^{-1}(V) = \bigcup f^{-1}(D(h_i))$
且 $\Spec(A) = \bigcup \Spec(A_{h_i})$，结论成立。
\end{remark}


\begin{lemma}
\label{morphisms-lemma-affine-equivalence-algebras}
设 $S$ 为概形。存在范畴的反等价
$$
\begin{matrix}
\text{Schemes affine} \\
\text{相对于 }S
\end{matrix}
\longleftrightarrow
\begin{matrix}
\text{quasi-coherent sheaves} \\
\text{属于 }\mathcal{O}_S\text{-algebras}
\end{matrix}
$$
它把 $f : X \to S$ 对应到层 $f_*\mathcal{O}_X$。
此外，这一等价与任意基变换相容。
\end{lemma}

\begin{proof}
从右到左的函子由 $\underline{\Spec}_S$ 给出。
由引理 \ref{morphisms-lemma-characterize-affine} 和《构造》引理
\ref{constructions-lemma-spec-properties} 的 (3)，这两个函子互逆。
最后一个断言即《构造》引理
\ref{constructions-lemma-spec-properties} 的 (2)。
\end{proof}

\begin{lemma}
\label{morphisms-lemma-quasi-coherence-scalar-restriction}
\begin{reference}
\cite[I, Proposition 9.6.1]{EGA}
\end{reference}
设 $S$ 为概形，$\mathcal{A}$ 为准凝聚 $\mathcal{O}_S$-代数。
一个 $\mathcal{A}$-模作为 $\mathcal{O}_S$-模准凝聚，当且仅当
它作为 $\mathcal{A}$-模准凝聚。
\end{lemma}

\begin{proof}
设 $\mathcal{F}$ 为一个 $\mathcal{A}$-模。若 $\mathcal{F}$ 作为
$\mathcal{A}$-模准凝聚，则对每个 $s \in S$，存在开邻域
$U$（它包含 $s$）以及 $\mathcal{A}|_U$-模的正合列
$$
\bigoplus\nolimits_J\mathcal{A}|_U\to
\bigoplus\nolimits_I\mathcal{A}|_U\to
\mathcal{F}|_U
$$
它同时也是 $\mathcal{O}_U$-模的正合列。因此，$\mathcal{F}|_U$
作为概形上两个准凝聚 $\mathcal{O}_U$-模之间态射的余核而准凝聚。
所以 $\mathcal{F}$ 作为 $\mathcal{O}_X$-模准凝聚。

\medskip\noindent
反之，假设 $\mathcal{F}$ 作为 $\mathcal{O}_X$-模准凝聚。
取仿射开集 $\Spec(R) = U \subset S$。有 $\mathcal{O}_U$-模同构
$\mathcal{A}|_U \cong \widetilde{A}$ 与
$\mathcal{F}|_U \cong \widetilde{M}$，其中 $R$-代数 $A$ 和
$R$-模 $M$ 适当选取。$\mathcal{A}$-模在 $\mathcal{F}$ 上的结构
转化为 $A$-模在 $M$ 上的结构，并与给定的 $R$-模结构相容
（细节略）。选取 $A$-模的正合列
$$
\bigoplus\nolimits_J A \to
\bigoplus\nolimits_I A \to
M \to 0
$$
由于函子 $\widetilde{\ }$ 正合，这给出 $\widetilde{A}$-模的正合列
$$
\bigoplus\nolimits_J \widetilde{A}\to
\bigoplus\nolimits_I \widetilde{A}\to
\widetilde{M} \to 0
$$
这说明 $\mathcal{F}$ 作为 $\mathcal{A}$-模准凝聚。
\end{proof}

\begin{lemma}
\label{morphisms-lemma-affine-equivalence-modules}
设 $f : X \to S$ 为概形的仿射态射，并令
$\mathcal{A} = f_*\mathcal{O}_X$。函子
$\mathcal{F} \mapsto f_*\mathcal{F}$ 诱导范畴等价
$$
\left\{
\begin{matrix}
\text{category of quasi-coherent}\\
\mathcal{O}_X\text{-模}
\end{matrix}
\right\}
\longrightarrow
\left\{
\begin{matrix}
\text{category of quasi-coherent}\\
\mathcal{A}\text{-模}
\end{matrix}
\right\}
$$
此外，一个 $\mathcal{A}$-模作为 $\mathcal{O}_S$-模准凝聚，
当且仅当它作为 $\mathcal{A}$-模准凝聚。
\end{lemma}

\begin{proof}
最后一个断言即引理 \ref{morphisms-lemma-quasi-coherence-scalar-restriction}。
由引理 \ref{morphisms-lemma-affine-separated} 和《概形》引理
\ref{schemes-lemma-push-forward-quasi-coherent}，正向像
$f_*\mathcal{F}$（其中 $\mathcal{O}_X$-模 $\mathcal{F}$ 准凝聚）
是准凝聚 $\mathcal{O}_S$-模。因此得到引理所述的函子。

\medskip\noindent
我们为这一函子构造拟逆 $h$。设 $\mathcal{G}$ 为准凝聚
$\mathcal{A}$-模。定义
$$
h(\mathcal{G}) = f^*\mathcal{G} \otimes_{f^*\mathcal{A}} \mathcal{O}_X =
f^*\mathcal{G} \otimes_{f^*f_*\mathcal{O}_X} \mathcal{O}_X
$$
说明如下：拉回 $f^*\mathcal{A} = f^*f_*\mathcal{O}_X$ 是
$\mathcal{O}_X$-代数，伴随映射
$f^*f_*\mathcal{O}_X \to \mathcal{O}_X$ 是代数同态，而拉回
$f^*\mathcal{G}$ 是 $f^*\mathcal{A}$-模。由于拉回和换环保持
拟凝聚性，$h(\mathcal{G})$ 准凝聚。存在函子性映射
$$
h(f_*\mathcal{F}) =
f^*f_*\mathcal{F} \otimes_{f^*f_*\mathcal{O}_X} \mathcal{O}_X \to \mathcal{F}
$$
它来自伴随映射 $f^*f_*\mathcal{F} \to \mathcal{F}$；还存在函子性映射
$$
\mathcal{G} \to f_*h(\mathcal{G}) \quad\text{adjoint to the map}\quad
f^*\mathcal{G} \to f^*\mathcal{G} \otimes_{f^*f_*\mathcal{O}_X} \mathcal{O}_X
$$
它把局部截面 $s$（它属于 $f^*\mathcal{F}$）送到 $s \otimes 1$。
要完成证明，只需证明这些映射对上述 $\mathcal{F}$ 和
$\mathcal{G}$ 都是同构。这可以在一个仿射覆盖的各成员上检验，
也就是在 $X$ 和 $S$ 都仿射时检验。

\medskip\noindent
使上述论证成立的关键代数观察如下。设 $R \to A$ 为环同态，
$N$ 为 $A$-模，则
$$
(N \otimes_R A) \otimes_{(A \otimes_R A)} A = N
$$
即等式左端是把 $h$ 作用于准凝聚 $\widetilde{A}$-模
$\widetilde{N}$ 所得的结果；该模在 $\Spec(R)$ 上。细节略。
\end{proof}

\begin{lemma}
\label{morphisms-lemma-composition-affine}
仿射态射的复合仍为仿射态射。
\end{lemma}

\begin{proof}
设 $f : X \to Y$ 与 $g : Y \to Z$ 为仿射态射。
设 $U \subset Z$ 为仿射开集。$g^{-1}(U)$ 由关于 $g$ 的假设仿射；
继而 $f^{-1}(g^{-1}(U))$ 由关于 $f$ 的假设仿射。
因此 $(g \circ f)^{-1}(U)$ 仿射。
\end{proof}

\begin{lemma}
\label{morphisms-lemma-base-change-affine}
仿射态射的基变换仍为仿射态射。
\end{lemma}

\begin{proof}
设 $f : X \to S$ 为仿射态射，$S' \to S$ 为任意态射。
以 $f' : X_{S'} = S' \times_S X \to S'$ 表示 $f$ 的基变换。
对每个 $s' \in S'$，都存在仿射开邻域 $s' \in V \subset S'$，
它映入某个仿射开集 $U \subset S$。由假设，$f^{-1}(U)$ 仿射。
由《概形》第 \ref{schemes-section-fibre-products} 节的内容，
$f^{-1}(U)_V = V \times_U f^{-1}(U)$ 仿射且等于
$(f')^{-1}(V)$。这证明 $S'$ 有一个仿射开覆盖，其在 $f'$ 下的逆像
均仿射。最后应用上述引理 \ref{morphisms-lemma-characterize-affine}。
\end{proof}

\begin{lemma}
\label{morphisms-lemma-closed-immersion-affine}
闭浸入是仿射态射。
\end{lemma}

\begin{proof}
这一事实最初见于《概形》引理
\ref{schemes-lemma-closed-immersion-affine-case}；完整陈述见《概形》引理
\ref{schemes-lemma-closed-subspace-scheme}。
\end{proof}

\begin{lemma}
\label{morphisms-lemma-affine-s-open}
设 $X$ 为概形，$\mathcal{L}$ 为可逆 $\mathcal{O}_X$-模，且
$s \in \Gamma(X, \mathcal{L})$。包含态射 $j : X_s \to X$ 仿射。
\end{lemma}

\begin{proof}
这由《性质》引理 \ref{properties-lemma-affine-cap-s-open} 及定义推出。
\end{proof}

\begin{lemma}
\label{morphisms-lemma-affine-permanence}
设 $g : X \to Y$ 为 $S$ 上的概形态射。
\begin{enumerate}
\item 若 $X$ 在 $S$ 上仿射，且 $\Delta : Y \to Y \times_S Y$ 仿射，
则 $g$ 仿射。
\item 若 $X$ 在 $S$ 上仿射，且 $Y$ 在 $S$ 上分离，则 $g$ 仿射。
\item 从仿射概形到具有仿射对角态射的概形的态射是仿射的。
\item 从仿射概形到分离概形的态射是仿射的。
\end{enumerate}
\end{lemma}

\begin{proof}
证明 (1)。基变换 $X \times_S Y \to Y$ 由引理
\ref{morphisms-lemma-base-change-affine} 仿射。态射
$(1, g) : X \to X \times_S Y$ 是 $Y \to Y \times_S Y$
沿态射 $X \times_S Y \to Y \times_S Y$ 的基变换。因此由引理
\ref{morphisms-lemma-base-change-affine}，它仿射。仿射态射的复合仿射
（见引理 \ref{morphisms-lemma-composition-affine}），故 (1) 成立。
由 (1) 得 (2)，因为闭浸入仿射（见引理
\ref{morphisms-lemma-closed-immersion-affine}），而 $Y/S$ 分离意味着
$\Delta$ 是闭浸入。(3) 与 (4) 分别是 (1) 与 (2) 的特殊情形。
\end{proof}

\begin{lemma}
\label{morphisms-lemma-morphism-affines-affine}
仿射概形之间的态射是仿射的。
\end{lemma}

\begin{proof}
在引理 \ref{morphisms-lemma-affine-permanence} 中取
$S = \Spec(\mathbf{Z})$ 即得。也可由引理
\ref{morphisms-lemma-characterize-affine} 中 (1) 与 (2) 的等价直接得到。
\end{proof}

\begin{lemma}
\label{morphisms-lemma-Artinian-affine}
设 $S$ 为概形，$A$ 为阿廷环。任意态射 $\Spec(A) \to S$ 都仿射。
\end{lemma}

\begin{proof}
略。
\end{proof}

\begin{lemma}
\label{morphisms-lemma-get-affine}
设 $j : Y \to X$ 为概形的浸入。假设存在开集 $U \subset X$，
其补集为 $Z = X \setminus U$，且满足
\begin{enumerate}
\item $U \to X$ 仿射；
\item $j^{-1}(U) \to U$ 仿射；
\item $j(Y) \cap Z$ 闭。
\end{enumerate}
则 $j$ 仿射。特别地，若 $X$ 仿射，则 $Y$ 也仿射。
\end{lemma}

\begin{proof}
由《概形》定义 \ref{schemes-definition-immersion}，存在开子概形
$W \subset X$，使 $j$ 分解为闭浸入 $i : Y \to W$，随后接包含态射
$W \to X$。由于闭浸入仿射（引理
\ref{morphisms-lemma-closed-immersion-affine}），对每个 $x \in W$，
$x$ 在 $X$ 中都有一个仿射开邻域，其在 $j$ 下的逆像仿射。
若 $x \in U$，则由假设 (2) 同样成立。最后，设 $x \in Z$
且 $x \not \in W$。于是 $x \not \in j(Y) \cap Z$。
由假设 (3)，可找到仿射开邻域 $V \subset X$，它包含 $x$
且与 $j(Y) \cap Z$ 不相交。此时
$j^{-1}(V) = j^{-1}(V \cap U)$，由假设 (1)、(2) 它仿射。
因此由引理 \ref{morphisms-lemma-characterize-affine}，$j$ 仿射。
\end{proof}

\section{丰沛可逆模族}
\label{morphisms-section-families-ample-invertible-modules}

\noindent
本节简述丰沛可逆模族的概念。

\begin{definition}
\label{morphisms-definition-family-ample-invertible-modules}
\begin{reference}
\cite[II Definition 2.2.4]{SGA6}
\end{reference}
设 $X$ 为概形，$\{\mathcal{L}_i\}_{i \in I}$ 为一族可逆
$\mathcal{O}_X$-模。若满足下列条件，则称
$\{\mathcal{L}_i\}_{i \in I}$ 为{\it $X$ 上的丰沛可逆模族}：
\begin{enumerate}
\item $X$ 拟紧；
\item 对每个 $x \in X$，存在 $i \in I$、$n \geq 1$ 以及
$s \in \Gamma(X, \mathcal{L}_i^{\otimes n})$，使得
$x \in X_s$ 且 $X_s$ 仿射。
\end{enumerate}
\end{definition}

\noindent
若 $\{\mathcal{L}_i\}_{i \in I}$ 是概形 $X$ 上的丰沛可逆模族，
则存在有限子集 $I' \subset I$，使得
$\{\mathcal{L}_i\}_{i \in I'}$ 仍是 $X$ 上的丰沛可逆模族
（这由拟紧性立即推出）。由下一个引理，具有丰沛可逆模族的概形
具有仿射对角态射，因而尤其拟分离。

\begin{lemma}
\label{morphisms-lemma-affine-diagonal}
设 $X$ 为概形，并且对每个点 $x \in X$，都存在可逆
$\mathcal{O}_X$-模 $\mathcal{L}$ 和全局截面
$s \in \Gamma(X, \mathcal{L})$，使得 $x \in X_s$ 且
$X_s$ 仿射。那么 $X$ 的对角态射仿射。
\end{lemma}

\begin{proof}
给定可逆 $\mathcal{O}_X$-模 $\mathcal{L}$、$\mathcal{M}$ 以及全局截面
$s \in \Gamma(X, \mathcal{L})$、
$t \in \Gamma(X, \mathcal{M})$，使得 $X_s$ 与 $X_t$ 仿射，
我们只需证明 $X_s \cap X_t$ 仿射。事实上，把引理
\ref{morphisms-lemma-characterize-affine} 应用于
$\Delta : X \to X \times X$，并利用
$\Delta^{-1}(X_s \times X_t) = X_s \cap X_t$，即可看出
$\Delta$ 仿射。而 $X_s \cap X_t$ 仿射则由《性质》引理
\ref{properties-lemma-affine-cap-s-open} 得到。
\end{proof}

\begin{remark}
\label{morphisms-remark-affine-s-opens-cover-family}
《性质》引理 \ref{properties-lemma-affine-s-opens-cover-quasi-separated}
表明，具有一个丰沛可逆模的概形是分离的。对于具有丰沛可逆模族的概形，
这一断言并不成立。具体地，取《概形》例
\ref{schemes-example-affine-space-zero-doubled} 中的 $X$，令
$n = 1$，即零点加倍的仿射直线。沿用该例的记号，但以 $x$
表示 $x_1$，以 $y$ 表示 $y_1$。对每个整数 $n$，都有一个可逆层
$\mathcal{L}_n$ 在 $X$ 上；它在 $X_1$ 与 $X_2$ 上平凡，且其转移函数
$U_{12} \to U_{21}$ 为 $f(x) \mapsto y^n f(y)$。
$\mathcal{L}_n$ 的全局截面是二元组
$(f(x), g(y)) \in k[x] \oplus
k[y]$，满足 $y^n f(y) = g(y)$。截面 $s = (1,
y)$ 属于 $\mathcal{L}_1$，截面 $t = (x, 1)$ 属于
$\mathcal{L}_{-1}$；二者确定一个仿射开覆盖，
因为 $X_s = X_1$ 且 $X_t = X_2$。因此 $X$ 具有丰沛可逆模族，
但它不分离。
\end{remark}

\section{拟仿射态射}
\label{morphisms-section-quasi-affine}

\noindent
回忆：若概形 $X$ 拟紧且同构于某个仿射概形的开子概形，
则称它为{\it 拟仿射的}；见《性质》定义
\ref{properties-definition-quasi-affine}。

\begin{definition}
\label{morphisms-definition-quasi-affine}
若概形态射 $f : X \to S$ 满足：$S$ 的每个仿射开集的逆像
都是拟仿射概形，则称它为{\it 拟仿射的}。
\end{definition}

\begin{lemma}
\label{morphisms-lemma-quasi-affine-separated}
拟仿射态射是分离且拟紧的。
\end{lemma}

\begin{proof}
设 $f : X \to S$ 拟仿射。由《概形》引理
\ref{schemes-lemma-quasi-compact-affine} 立即得到拟紧性。
设 $U \subset S$ 为仿射开集。若能证明 $f^{-1}(U)$ 是分离概形，
则 $f$ 分离（《概形》引理
\ref{schemes-lemma-characterize-separated} 表明分离性在基底上是局部的）。
依假设，$f^{-1}(U)$ 同构于某个仿射概形的开子概形。
仿射概形分离，故其每个开子概形也分离，如所需。
\end{proof}

\begin{lemma}
\label{morphisms-lemma-characterize-quasi-affine}
设 $f : X \to S$ 为概形态射。下列条件等价：
\begin{enumerate}
\item 态射 $f$ 拟仿射。
\item 存在仿射开覆盖 $S = \bigcup W_j$，使每个
$f^{-1}(W_j)$ 都拟仿射。
\item 存在准凝聚 $\mathcal{O}_S$-代数层 $\mathcal{A}$ 以及拟紧开浸入
$$
\xymatrix{
X \ar[rr] \ar[rd] & & \underline{\Spec}_S(\mathcal{A}) \ar[dl] \\
& S &
}
$$
在 $S$ 上。
\item 与 (3) 相同，但取 $\mathcal{A} = f_*\mathcal{O}_X$，
并令水平箭头为《构造》引理
\ref{constructions-lemma-canonical-morphism} 的典范态射。
\end{enumerate}
\end{lemma}

\begin{proof}
(1) 蕴含 (2) 以及 (4) 蕴含 (3) 都是显然的。

\medskip\noindent
设 $S = \bigcup_{j \in J} W_j$ 是仿射开覆盖，且每个
$f^{-1}(W_j)$ 都拟仿射。由《概形》引理
\ref{schemes-lemma-quasi-compact-affine}，$f$ 拟紧。由《概形》引理
\ref{schemes-lemma-characterize-quasi-separated}，态射 $f$ 拟分离。
因此由《概形》引理 \ref{schemes-lemma-push-forward-quasi-coherent}，
层 $\mathcal{A} = f_*\mathcal{O}_X$ 是准凝聚
$\mathcal{O}_X$-代数层。于是有概形
$g : Y = \underline{\Spec}_S(\mathcal{A}) \to S$，它在 $S$ 上。
恒等映射
$\text{id} : \mathcal{A} = f_*\mathcal{O}_X \to f_*\mathcal{O}_X$
由相对谱的定义给出态射 $can : X \to Y$，它在 $S$ 上；见《构造》
引理 \ref{constructions-lemma-canonical-morphism}。由假设、刚引用的引理
以及《性质》引理 \ref{properties-lemma-quasi-affine}，限制
$can|_{f^{-1}(W_j)} : f^{-1}(W_j) \to g^{-1}(W_j)$ 是拟紧开浸入。
因此 $can$ 是拟紧开浸入。我们已经证明 (2) 蕴含 (4)。

\medskip\noindent
假设 (3)。任取仿射开集 $U \subset S$。由《构造》引理
\ref{constructions-lemma-spec-properties}，$U$ 在相对谱中的逆像仿射。
因此 $f^{-1}(U)$ 拟仿射（注意拟紧性也已编码在 (3) 中）。
故 (3) 蕴含 (1)。
\end{proof}

\begin{lemma}
\label{morphisms-lemma-composition-quasi-affine}
拟仿射态射的复合仍为拟仿射态射。
\end{lemma}

\begin{proof}
设 $f : X \to Y$ 与 $g : Y \to Z$ 为拟仿射态射。
设 $U \subset Z$ 为仿射开集。$g^{-1}(U)$ 由关于 $g$ 的假设拟仿射。
设 $j : g^{-1}(U) \to V$ 是到仿射概形 $V$ 的拟紧开浸入。
由上述引理 \ref{morphisms-lemma-characterize-quasi-affine}，
$f^{-1}(g^{-1}(U))$ 是相对谱
$\underline{\Spec}_{g^{-1}(U)}(\mathcal{A})$ 的拟紧开子概形，其中有某个
准凝聚 $\mathcal{O}_{g^{-1}(U)}$-代数层 $\mathcal{A}$。
由《概形》引理 \ref{schemes-lemma-push-forward-quasi-coherent}，层
$\mathcal{A}' = j_*\mathcal{A}$ 是准凝聚 $\mathcal{O}_V$-代数层，
且满足 $j^*\mathcal{A}' = \mathcal{A}$。因此得到交换图
$$
\xymatrix{
f^{-1}(g^{-1}(U)) \ar[r] &
\underline{\Spec}_{g^{-1}(U)}(\mathcal{A})
\ar[r] \ar[d] &
\underline{\Spec}_V(\mathcal{A}') \ar[d] \\
& g^{-1}(U) \ar[r]^j & V
}
$$
其中方块为纤维方块；见《构造》引理
\ref{constructions-lemma-spec-properties}。注意右上角是仿射概形。
因此 $(g \circ f)^{-1}(U)$ 拟仿射。
\end{proof}

\begin{lemma}
\label{morphisms-lemma-base-change-quasi-affine}
拟仿射态射的基变换仍为拟仿射态射。
\end{lemma}

\begin{proof}
设 $f : X \to S$ 为拟仿射态射。由上述引理
\ref{morphisms-lemma-characterize-quasi-affine}，可以找到准凝聚
$\mathcal{O}_S$-代数层 $\mathcal{A}$ 和拟紧开浸入
$X \to \underline{\Spec}_S(\mathcal{A})$，后者在 $S$ 上。
设 $g : S' \to S$ 为任意态射。以
$f' : X_{S'} = S' \times_S X \to S'$ 表示 $f$ 的基变换。
拟紧开浸入的基变换仍为拟紧开浸入，故
$X_{S'} \to \underline{\Spec}_{S'}(g^*\mathcal{A})$
是拟紧开浸入（这里使用了《概形》引理
\ref{schemes-lemma-quasi-compact-preserved-base-change}、
\ref{schemes-lemma-base-change-immersion} 以及《构造》引理
\ref{constructions-lemma-spec-properties}）。再次由引理
\ref{morphisms-lemma-characterize-quasi-affine}，$X_{S'} \to S'$ 拟仿射。
\end{proof}

\begin{lemma}
\label{morphisms-lemma-quasi-compact-immersion-quasi-affine}
拟紧浸入是拟仿射态射。
\end{lemma}

\begin{proof}
设 $X \to S$ 为拟紧浸入。需要证明每个仿射开集的逆像拟仿射。
因此，可假设 $S$ 是仿射概形，并证明 $X$ 拟仿射。
由引理 \ref{morphisms-lemma-quasi-compact-immersion}，态射 $X \to S$ 分解为
$X \to Z \to S$，其中 $Z$ 是 $S$ 的闭子概形，且 $X \subset Z$
是拟紧开集。由于 $S$ 仿射，引理 \ref{morphisms-lemma-closed-immersion}
推出 $Z$ 仿射。结论得证。
\end{proof}

\begin{lemma}
\label{morphisms-lemma-affine-quasi-affine}
设 $S$ 为概形，$X$ 为仿射概形。态射 $f : X \to S$ 拟仿射，
当且仅当它拟紧。特别地，从仿射概形到拟分离概形的任意态射都拟仿射。
\end{lemma}

\begin{proof}
设 $V \subset S$ 为仿射开集。则 $f^{-1}(V)$ 是仿射概形 $X$ 的
开子概形，所以它拟仿射当且仅当它拟紧。这证明第一个断言。
若 $f : X \to S$ 中 $X$ 仿射而 $S$ 拟分离，则该态射的拟紧性
由《概形》引理 \ref{schemes-lemma-quasi-compact-permanence} 应用于
$X \to S \to \Spec(\mathbf{Z})$ 得到。
\end{proof}

\begin{lemma}
\label{morphisms-lemma-quasi-affine-permanence}
设 $g : X \to Y$ 为 $S$ 上的概形态射。若 $X$ 在 $S$ 上拟仿射，
且 $Y$ 在 $S$ 上拟分离，则 $g$ 拟仿射。特别地，从拟仿射概形
到拟分离概形的任意态射都拟仿射。
\end{lemma}

\begin{proof}
基变换 $X \times_S Y \to Y$ 由引理
\ref{morphisms-lemma-base-change-quasi-affine} 拟仿射。态射
$X \to X \times_S Y$ 是拟紧浸入，因为 $Y \to S$ 拟分离；见《概形》引理
\ref{schemes-lemma-section-immersion}。由引理
\ref{morphisms-lemma-quasi-compact-immersion-quasi-affine}，拟紧浸入拟仿射；
而拟仿射态射的复合拟仿射（见引理
\ref{morphisms-lemma-composition-quasi-affine}）。结论得证。
\end{proof}

\section{由环同态性质定义的态射类型}
\label{morphisms-section-properties-ring-maps}

\noindent
本节研究环同态的哪些性质可用来定义概形态射的局部性质。

\begin{definition}
\label{morphisms-definition-property-local}
设 $P$ 为环同态的一个性质。
\begin{enumerate}
\item 若下列条件成立，则称 $P$ 是{\it 局部的}：
\begin{enumerate}
\item 对任意环同态 $R \to A$ 以及任意 $f \in R$，有
$P(R \to A) \Rightarrow P(R_f \to A_f)$。
\item 对任意环 $R$、$A$，任意 $f \in R$、$a\in A$，以及任意环同态
$R_f \to A$，有 $P(R_f \to A) \Rightarrow P(R \to A_a)$。
\item 对任意环同态 $R \to A$，若 $a_i \in A$ 且
$(a_1, \ldots, a_n) = A$，则
$\forall i, P(R \to A_{a_i}) \Rightarrow P(R \to A)$。
\end{enumerate}
\item 称 $P$ 在{\it 基变换下稳定}，如果对任意环同态
$R \to A$、$R \to R'$ 都有
$P(R \to A) \Rightarrow P(R' \to R' \otimes_R A)$。
\item 称 $P$ 在{\it 复合下稳定}，如果对任意环同态
$A \to B$、$B \to C$ 都有
$P(A \to B) \wedge P(B \to C) \Rightarrow P(A \to C)$。
\end{enumerate}
\end{definition}

\begin{definition}
\label{morphisms-definition-locally-P}
设 $P$ 为环同态的一个性质，$f : X \to S$ 为概形态射。
称 $f$ {\it 局部为 $P$ 型}，如果对任意 $x \in X$，
存在仿射开邻域 $U$，它是 $x$ 在 $X$ 中的邻域，并映入某个仿射开集
$V \subset S$，而诱导环同态
$\mathcal{O}_S(V) \to \mathcal{O}_X(U)$ 具有性质 $P$。
\end{definition}

\noindent
除非性质 $P$ 是局部性质，否则这不是一个“好”定义。
即使 $P$ 是局部性质，我们也不会自动用这个定义称一个态射
“局部为 $P$ 型”，除非还在别处明确给出相应定义。

\begin{lemma}
\label{morphisms-lemma-locally-P}
设 $f : X \to S$ 为概形态射，$P$ 为环同态的一个性质。
设 $U$ 为 $X$ 的仿射开集，$V$ 为 $S$ 的仿射开集，且
$f(U) \subset V$。若 $f$ 局部为 $P$ 型且 $P$ 是局部的，
则 $P(\mathcal{O}_S(V) \to \mathcal{O}_X(U))$ 成立。
\end{lemma}

\begin{proof}
因为 $f$ 局部为 $P$ 型，所以对每个 $u \in U$，存在仿射开集
$U_u \subset X$，它映入仿射开集 $V_u \subset S$，且
$P(\mathcal{O}_S(V_u) \to \mathcal{O}_X(U_u))$ 成立。
选取开邻域 $U'_u \subset U \cap U_u$，使它包含 $u$ 且同时是
$U$ 与 $U_u$ 中的标准仿射开集；见《概形》引理
\ref{schemes-lemma-standard-open-two-affines}。由定义
\ref{morphisms-definition-property-local} (1)(b)，
$P(\mathcal{O}_S(V_u) \to \mathcal{O}_X(U'_u))$ 成立。
因此可以假设 $U_u \subset U$ 是标准仿射开集。
选取开邻域 $V'_u \subset V \cap V_u$，使它包含 $f(u)$ 且同时是
$V$ 与 $V_u$ 中的标准仿射开集；见《概形》引理
\ref{schemes-lemma-standard-open-two-affines}。于是
$U'_u = f^{-1}(V'_u) \cap U_u$ 是 $U_u$（因而也是 $U$）中的
标准仿射开集，并由定义 \ref{morphisms-definition-property-local} (1)(a) 有
$P(\mathcal{O}_S(V'_u) \to \mathcal{O}_X(U'_u))$。
因此可以假设 $U_u \subset U$ 与 $V_u \subset V$ 都是标准仿射开集。
再应用一次定义 \ref{morphisms-definition-property-local} (1)(b)，得到
$P(\mathcal{O}_S(V) \to \mathcal{O}_X(U_u))$。因为 $U$ 拟紧，
可选取有限多个点 $u_1, \ldots, u_n \in U$，使得
$$
U = U_{u_1} \cup \ldots \cup U_{u_n}.
$$
由定义 \ref{morphisms-definition-property-local} (1)(c)，得
$P(\mathcal{O}_S(V) \to \mathcal{O}_X(U))$。
\end{proof}

\begin{lemma}
\label{morphisms-lemma-locally-P-characterize}
设 $P$ 为环同态的局部性质，$f : X \to S$ 为概形态射。
下列条件等价：
\begin{enumerate}
\item 态射 $f$ 局部为 $P$ 型。
\item 对每一对仿射开集 $U \subset X$、$V \subset S$，若
$f(U) \subset V$，则 $P(\mathcal{O}_S(V) \to \mathcal{O}_X(U))$ 成立。
\item 存在开覆盖 $S = \bigcup_{j \in J} V_j$ 以及开覆盖
$f^{-1}(V_j) = \bigcup_{i \in I_j} U_i$，使每个态射
$U_i \to V_j$（$j\in J, i\in I_j$）都局部为 $P$ 型。
\item 存在仿射开覆盖 $S = \bigcup_{j \in J} V_j$ 以及仿射开覆盖
$f^{-1}(V_j) = \bigcup_{i \in I_j} U_i$，使得
$P(\mathcal{O}_S(V_j) \to \mathcal{O}_X(U_i))$ 对所有
$j\in J, i\in I_j$ 都成立。
\end{enumerate}
此外，若 $f$ 局部为 $P$ 型，则对任意开子概形
$U \subset X$、$V \subset S$，只要 $f(U) \subset V$，限制
$f|_U : U \to V$ 就局部为 $P$ 型。
\end{lemma}

\begin{proof}
这由上述引理 \ref{morphisms-lemma-locally-P} 推出。
\end{proof}

\begin{lemma}
\label{morphisms-lemma-composition-type-P}
设 $P$ 为环同态的一个性质，并假设 $P$ 是局部的且在复合下稳定。
局部为 $P$ 型的态射之复合仍局部为 $P$ 型。
\end{lemma}

\begin{proof}
设 $f : X \to Y$ 与 $g : Y \to Z$ 为局部 $P$ 型的态射。
设 $x \in X$。选取仿射开邻域 $W \subset Z$，使它包含
$g(f(x))$；再选取仿射开邻域 $V \subset g^{-1}(W)$，使它包含
$f(x)$；最后选取仿射开邻域 $U \subset f^{-1}(V)$，使它包含
$x$。由引理
\ref{morphisms-lemma-locally-P-characterize}，环同态
$\mathcal{O}_Z(W) \to \mathcal{O}_Y(V)$ 与
$\mathcal{O}_Y(V) \to \mathcal{O}_X(U)$ 满足 $P$。
故 $\mathcal{O}_Z(W) \to \mathcal{O}_X(U)$ 满足 $P$，
因为假设 $P$ 在复合下稳定。
\end{proof}

\begin{lemma}
\label{morphisms-lemma-base-change-type-P}
设 $P$ 为环同态的一个性质，并假设 $P$ 是局部的且在基变换下稳定。
局部为 $P$ 型的态射之基变换仍局部为 $P$ 型。
\end{lemma}

\begin{proof}
设 $f : X \to S$ 为局部 $P$ 型的态射，$S' \to S$ 为任意态射。
以 $f' : X_{S'} = S' \times_S X \to S'$ 表示 $f$ 的基变换。
对每个 $s' \in S'$，都存在仿射开邻域 $s' \in V' \subset S'$，
它映入某个仿射开集 $V \subset S$。由引理
\ref{morphisms-lemma-locally-P-characterize}，开集 $f^{-1}(V)$ 是仿射开集
$U_i$ 的并，并且环同态
$\mathcal{O}_S(V) \to \mathcal{O}_X(U_i)$ 都满足 $P$。
由《概形》第 \ref{schemes-section-fibre-products} 节的内容，
$f^{-1}(U)_{V'} = V' \times_V f^{-1}(V)$ 是仿射开集
$V' \times_V U_i$ 的并。由于
$\mathcal{O}_{X_{S'}}(V' \times_V U_i) =
\mathcal{O}_{S'}(V') \otimes_{\mathcal{O}_S(V)} \mathcal{O}_X(U_i)$，
环同态
$\mathcal{O}_{S'}(V') \to \mathcal{O}_{X_{S'}}(V' \times_V U_i)$
满足 $P$，因为假设 $P$ 在基变换下稳定。
\end{proof}

\begin{lemma}
\label{morphisms-lemma-properties-local}
环同态 $R \to A$ 的下列性质是局部的：
\begin{enumerate}
\item （局部环上的同构。）
对每个素理想 $\mathfrak q$（它属于 $A$），若它位于
$\mathfrak p \subset R$ 之上，则环同态 $R \to A$ 诱导同构
$R_{\mathfrak p} \to A_{\mathfrak q}$。
\item （开浸入。）
对每个素理想 $\mathfrak q$（它属于 $A$），存在 $f \in R$，满足
$\varphi(f) \not \in \mathfrak q$，且环同态 $\varphi : R \to A$
诱导同构 $R_f \to A_f$。
\item （既约纤维。）
对每个素理想 $\mathfrak p$（它属于 $R$），纤维环
$A \otimes_R \kappa(\mathfrak p)$ 是既约的。
\item （维数至多为 $n$ 的纤维。）
对每个素理想 $\mathfrak p$（它属于 $R$），纤维环
$A \otimes_R \kappa(\mathfrak p)$ 的 Krull 维数至多为 $n$。
\item （靶上局部诺特。）
环同态 $R \to A$ 具有性质：$A$ 是诺特环。
\item 按需在此添加更多性质\footnote{但仅添加那些尚未在别处单独处理的性质。}。
\end{enumerate}
\end{lemma}

\begin{proof}
略。
\end{proof}

\begin{lemma}
\label{morphisms-lemma-properties-base-change}
环同态的下列性质在基变换下稳定：
\begin{enumerate}
\item （局部环上的同构。）
对每个素理想 $\mathfrak q$（它属于 $A$），若它位于
$\mathfrak p \subset R$ 之上，则环同态 $R \to A$ 诱导同构
$R_{\mathfrak p} \to A_{\mathfrak q}$。
\item （开浸入。）
对每个素理想 $\mathfrak q$（它属于 $A$），存在 $f \in R$，满足
$\varphi(f) \not \in \mathfrak q$，且环同态 $\varphi : R \to A$
诱导同构 $R_f \to A_f$。
\item 按需在此添加更多性质\footnote{但仅添加那些尚未在别处单独处理的性质。}。
\end{enumerate}
\end{lemma}

\begin{proof}
略。
\end{proof}

\begin{lemma}
\label{morphisms-lemma-properties-composition}
环同态的下列性质在复合下稳定：
\begin{enumerate}
\item （局部环上的同构。）
对每个素理想 $\mathfrak q$（它属于 $A$），若它位于
$\mathfrak p \subset R$ 之上，则环同态 $R \to A$ 诱导同构
$R_{\mathfrak p} \to A_{\mathfrak q}$。
\item （开浸入。）
对每个素理想 $\mathfrak q$（它属于 $A$），存在 $f \in R$，满足
$\varphi(f) \not \in \mathfrak q$，且环同态 $\varphi : R \to A$
诱导同构 $R_f \to A_f$。
\item （靶上局部诺特。）
环同态 $R \to A$ 具有性质：$A$ 是诺特环。
\item 按需在此添加更多性质\footnote{但仅添加那些尚未在别处单独处理的性质。}。
\end{enumerate}
\end{lemma}

\begin{proof}
略。
\end{proof}

\section{有限型态射}
\label{morphisms-section-finite-type}

\noindent
回忆：若环同态 $R \to A$ 使得 $A$ 同构于
$R[x_1, \ldots, x_n]$ 的一个商（作为 $R$-代数），则称它为有限型；见《代数》定义
\ref{algebra-definition-finite-type}。

\begin{definition}
\label{morphisms-definition-finite-type}
设 $f : X \to S$ 为概形态射。
\begin{enumerate}
\item 称 $f$ {\it 在 $x \in X$ 处为有限型}，如果存在仿射开邻域
$\Spec(A) = U \subset X$，它包含 $x$，以及仿射开集
$\Spec(R) = V \subset S$，满足 $f(U) \subset V$，且诱导环同态
$R \to A$ 为有限型。
\item 若 $f$ 在 $X$ 的每一点处都为有限型，则称它{\it 局部为有限型}。
\item 若 $f$ 局部为有限型且拟紧，则称它为{\it 有限型}。
\end{enumerate}
\end{definition}

\begin{lemma}
\label{morphisms-lemma-locally-finite-type-characterize}
设 $f : X \to S$ 为概形态射。下列条件等价：
\begin{enumerate}
\item 态射 $f$ 局部为有限型。
\item 对所有仿射开集 $U \subset X$、$V \subset S$，若
$f(U) \subset V$，则环同态
$\mathcal{O}_S(V) \to \mathcal{O}_X(U)$ 为有限型。
\item 存在开覆盖 $S = \bigcup_{j \in J} V_j$ 以及开覆盖
$f^{-1}(V_j) = \bigcup_{i \in I_j} U_i$，使每个态射
$U_i \to V_j$（$j\in J, i\in I_j$）都局部为有限型。
\item 存在仿射开覆盖 $S = \bigcup_{j \in J} V_j$ 以及仿射开覆盖
$f^{-1}(V_j) = \bigcup_{i \in I_j} U_i$，使环同态
$\mathcal{O}_S(V_j) \to \mathcal{O}_X(U_i)$ 对所有
$j\in J, i\in I_j$ 都为有限型。
\end{enumerate}
此外，若 $f$ 局部为有限型，则对任意开子概形
$U \subset X$、$V \subset S$，只要 $f(U) \subset V$，限制
$f|_U : U \to V$ 就局部为有限型。
\end{lemma}

\begin{proof}
若证明“$R \to A$ 为有限型”是局部性质，结论即由上述引理
\ref{morphisms-lemma-locally-P} 得到。我们检验定义
\ref{morphisms-definition-property-local} 的条件 (a)、(b)、(c)。
由《代数》引理 \ref{algebra-lemma-base-change-finiteness}，有限型性
在基变换下稳定，故 (a) 成立。由《代数》引理
\ref{algebra-lemma-compose-finite-type}，有限型性在复合下稳定；
而对任意环 $R$，环同态 $R \to R_f$ 显然为有限型。
故 (b) 成立。最后，(c) 由《代数》引理
\ref{algebra-lemma-cover-upstairs} 得到。
\end{proof}

\begin{lemma}
\label{morphisms-lemma-composition-finite-type}
两个局部有限型态射的复合局部为有限型。两个有限型态射的复合也为有限型。
\end{lemma}

\begin{proof}
在引理 \ref{morphisms-lemma-locally-finite-type-characterize} 的证明中，
我们已经看到有限型性是环同态的局部性质。因此，引理的第一个断言
由引理 \ref{morphisms-lemma-composition-type-P} 以及有限型性在环同态复合下稳定
这一事实推出；见《代数》引理
\ref{algebra-lemma-compose-finite-type}。再结合上述结论以及拟紧态射的复合
仍拟紧这一事实（见《概形》引理
\ref{schemes-lemma-composition-quasi-compact}），可知有限型态射的复合
为有限型。
\end{proof}

\begin{lemma}
\label{morphisms-lemma-base-change-finite-type}
局部有限型态射的基变换局部为有限型。有限型态射的基变换也为有限型。
\end{lemma}

\begin{proof}
在引理 \ref{morphisms-lemma-locally-finite-type-characterize} 的证明中，
我们已经看到有限型性是环同态的局部性质。因此，引理的第一个断言
由引理 \ref{morphisms-lemma-base-change-type-P} 以及有限型性在环同态基变换下稳定
这一事实推出；见《代数》引理
\ref{algebra-lemma-base-change-finiteness}。再结合上述结论以及拟紧态射的基变换
仍拟紧这一事实（见《概形》引理
\ref{schemes-lemma-quasi-compact-preserved-base-change}），可知有限型态射的基变换
为有限型态射。
\end{proof}

\begin{lemma}
\label{morphisms-lemma-immersion-locally-finite-type}
闭浸入为有限型。浸入局部为有限型。
\end{lemma}

\begin{proof}
这是因为开浸入是局部同构，而闭浸入显然为有限型。
\end{proof}

\begin{lemma}
\label{morphisms-lemma-finite-type-noetherian}
设 $f : X \to S$ 为态射。若 $S$（局部）诺特且 $f$（局部）为有限型，
则 $X$（局部）诺特。
\end{lemma}

\begin{proof}
有限型诺特环代数仍为诺特环，故结论立即成立；见《代数》引理
\ref{algebra-lemma-Noetherian-permanence}。（也可使用：拟紧态射的靶若拟紧，
则源也拟紧。）
\end{proof}

\begin{lemma}
\label{morphisms-lemma-finite-type-Noetherian-quasi-separated}
设 $f : X \to S$ 局部为有限型，且 $S$ 局部诺特。则 $f$ 拟分离。
\end{lemma}

\begin{proof}
事实上，$X$ 拟分离；见《性质》引理
\ref{properties-lemma-locally-Noetherian-quasi-separated} 及上述引理
\ref{morphisms-lemma-finite-type-noetherian}。再应用《概形》引理
\ref{schemes-lemma-compose-after-separated}，得到 $f$ 拟分离。
\end{proof}

\begin{lemma}
\label{morphisms-lemma-permanence-finite-type}
设 $X \to Y$ 为基概形 $S$ 上的概形态射。若 $X$ 在 $S$ 上
局部为有限型，则 $X \to Y$ 局部为有限型。
\end{lemma}

\begin{proof}
通过引理 \ref{morphisms-lemma-locally-finite-type-characterize}，这转化为下列
代数事实：给定环同态 $A \to B \to C$，若 $A \to C$ 为有限型，
则 $B \to C$ 为有限型。（见《代数》引理
\ref{algebra-lemma-compose-finite-type}。）
\end{proof}

\section{有限型点与 Jacobson 概形}
\label{morphisms-section-points-finite-type}

\noindent
设 $S$ 为概形。有限型点 $s$（属于 $S$）是使态射
$\Spec(\kappa(s)) \to S$ 为有限型的点。研究它的原因是：
在某种意义和某些情形下，有限型点可以代替闭点。例如，它们总是足够多。
此外，一个概形是 Jacobson 概形，当且仅当所有有限型点都是闭点。

\begin{lemma}
\label{morphisms-lemma-point-finite-type}
设 $S$ 为概形，$k$ 为域，$f : \Spec(k) \to S$ 为态射。
下列条件等价：
\begin{enumerate}
\item 态射 $f$ 为有限型。
\item 态射 $f$ 局部为有限型。
\item 存在仿射开集 $U = \Spec(R)$（它属于 $S$），使 $f$ 对应于有限环同态
$R \to k$。
\item 存在仿射开集 $U = \Spec(R)$（它属于 $S$），使 $f$ 的像由
一个闭点 $u$ 组成，该点在 $U$ 中，且域扩张 $k/\kappa(u)$ 有限。
\end{enumerate}
\end{lemma}

\begin{proof}
由于 $\Spec(k)$ 仅含一点，故从它出发的任意态射都拟紧，
所以 (1) 与 (2) 的等价是显然的。

\medskip\noindent
设 $f$ 局部为有限型。任取仿射开集 $\Spec(R) = U \subset S$，
使 $f$ 的像包含在 $U$ 中且环同态 $R \to k$ 为有限型。
设 $\mathfrak p \subset R$ 为其核。于是 $R/\mathfrak p \subset k$
为有限型。由《代数》引理
\ref{algebra-lemma-field-finite-type-over-domain}，存在
$\overline{f} \in R/\mathfrak p$，使 $(R/\mathfrak p)_{\overline{f}}$
为域，且 $(R/\mathfrak p)_{\overline{f}} \to k$ 为有限域扩张。
若 $f \in R$ 是 $\overline{f}$ 的一个提升，则 $k$ 是有限
$R_f$-模。因此 (2) $\Rightarrow$ (3)。

\medskip\noindent
设 $\Spec(R) = U \subset S$ 为仿射开集，并且 $f$ 对应于有限环同态
$R \to k$。由引理 \ref{morphisms-lemma-locally-finite-type-characterize}，
$f$ 局部为有限型。因此 (3) $\Rightarrow$ (2)。

\medskip\noindent
设 $R \to k$ 有限。$R \to k$ 的像是一个域，并且由《代数》引理
\ref{algebra-lemma-integral-under-field}，$k$ 在该域上有限。
因此 $R \to k$ 的核是极大理想。故 (3) $\Rightarrow$ (4)。

\medskip\noindent
(4) $\Rightarrow$ (3) 是立即的。
\end{proof}

\begin{lemma}
\label{morphisms-lemma-artinian-finite-type}
设 $S$ 为概形，$A$ 为剩余域是 $\kappa$ 的阿廷局部环，
$f : \Spec(A) \to S$ 为概形态射。则 $f$ 为有限型，当且仅当复合
$\Spec(\kappa) \to \Spec(A) \to S$ 为有限型。
\end{lemma}

\begin{proof}
由于态射 $\Spec(\kappa) \to \Spec(A)$ 为有限型，若 $f$ 为有限型，
则复合 $\Spec(\kappa) \to S$ 也为有限型，这是清楚的
（见引理 \ref{morphisms-lemma-composition-finite-type}）。反之，注意
$\Spec(A) \to S$ 的像包含在某个仿射开集
$U = \Spec(B)$ 中，该开集属于 $S$，因为 $\Spec(A)$ 只有一个点。最后对
$B \to A$ 应用《代数》引理
\ref{algebra-lemma-essentially-of-finite-type-into-artinian-local}。
\end{proof}

\noindent
回忆：给定一点 $s$，它属于概形 $S$，则有典范态射
$\Spec(\kappa(s)) \to S$；见《概形》第
\ref{schemes-section-points} 节。

\begin{definition}
\label{morphisms-definition-finite-type-point}
设 $S$ 为概形。若一点 $s$（属于 $S$）的典范态射
$\Spec(\kappa(s)) \to S$ 为有限型，则称该点为{\it 有限型点}。
用 $S_{\text{ft-pts}}$ 表示 $S$ 的有限型点集。
\end{definition}

\noindent
有限型点集可以刻画如下。

\begin{lemma}
\label{morphisms-lemma-identify-finite-type-points}
设 $S$ 为概形。我们有
$$
S_{\text{ft-pts}} = \bigcup\nolimits_{U \subset S\text{ 开}} U_0
$$
其中 $U_0$ 是 $U$ 的闭点集。这里可以让 $U$ 遍历 $S$ 的所有开集，
也可以只遍历所有仿射开集。
\end{lemma}

\begin{proof}
这由引理 \ref{morphisms-lemma-point-finite-type} 立即得到。
\end{proof}

\begin{lemma}
\label{morphisms-lemma-finite-type-points-morphism}
设 $f : T \to S$ 为概形态射。若 $f$ 局部为有限型，则
$f(T_{\text{ft-pts}}) \subset S_{\text{ft-pts}}$。
\end{lemma}

\begin{proof}
若 $T$ 是域的谱，这就是引理 \ref{morphisms-lemma-point-finite-type}。
一般情形则由局部有限型态射的复合仍局部为有限型得到
（引理 \ref{morphisms-lemma-composition-finite-type}）。
\end{proof}

\begin{lemma}
\label{morphisms-lemma-finite-type-points-surjective-morphism}
设 $f : T \to S$ 为概形态射。若 $f$ 局部为有限型且满射，则
$f(T_{\text{ft-pts}}) = S_{\text{ft-pts}}$。
\end{lemma}

\begin{proof}
由引理 \ref{morphisms-lemma-finite-type-points-morphism}，有
$f(T_{\text{ft-pts}}) \subset S_{\text{ft-pts}}$。设 $s \in S$
为有限型点。由于 $f$ 满射，概形
$T_s = \Spec(\kappa(s)) \times_S T$ 非空；因此由引理
\ref{morphisms-lemma-identify-finite-type-points}，它有有限型点 $t \in T_s$。
现在 $T_s \to T$ 是有限型态射，因为它是 $s \to S$ 的基变换
（引理 \ref{morphisms-lemma-base-change-finite-type}）。因此由引理
\ref{morphisms-lemma-finite-type-points-morphism}，$t$ 在 $T$ 中的像是有限型点，
并且按构造映到 $s$。
\end{proof}

\begin{lemma}
\label{morphisms-lemma-enough-finite-type-points}
设 $S$ 为概形。对任意局部闭子集 $T \subset S$，有
$$
T \not = \emptyset
\Rightarrow
T \cap S_{\text{ft-pts}} \not = \emptyset.
$$
特别地，对任意闭子集 $T \subset S$，$T \cap S_{\text{ft-pts}}$
在 $T$ 中稠密。
\end{lemma}

\begin{proof}
注意 $T$ 带有概形结构（见《概形》引理
\ref{schemes-lemma-reduced-closed-subscheme}），使 $T \to S$ 是局部闭浸入。
任意局部闭浸入都局部为有限型，见引理
\ref{morphisms-lemma-immersion-locally-finite-type}。因此由引理
\ref{morphisms-lemma-finite-type-points-morphism}，有
$T_{\text{ft-pts}} \subset S_{\text{ft-pts}}$。最后，$T$ 的任意非空
仿射开集都至少有一个闭点；由引理
\ref{morphisms-lemma-identify-finite-type-points}，该点是 $T$ 的有限型点。
\end{proof}

\noindent
因此，《拓扑》第 \ref{topology-section-space-jacobson} 节的大部分内容
在把闭点集替换为有限型点集后仍成立。事实上，若 $S$ 是 Jacobson
概形，则有限型点正好恢复为闭点。

\begin{lemma}
\label{morphisms-lemma-jacobson-finite-type-points}
设 $S$ 为概形。下列条件等价：
\begin{enumerate}
\item 概形 $S$ 是 Jacobson 的；
\item $S_{\text{ft-pts}}$ 是 $S$ 的闭点集；
\item 对所有局部有限型态射 $T \to S$，闭点映到闭点；
\item 对所有局部有限型态射 $T \to S$，闭点 $t \in T$ 映到闭点
$s \in S$，且 $\kappa(s) \subset \kappa(t)$ 有限。
\end{enumerate}
\end{lemma}

\begin{proof}
显然有 (4) $\Rightarrow$ (3) $\Rightarrow$ (2)。引理
\ref{morphisms-lemma-enough-finite-type-points} 表明 (2) 蕴含 (1)。
故只需证明 (1) 蕴含 (4)。设 $T \to S$ 局部为有限型。
取闭点 $t \in T$，并令 $s \in S$ 为其像。选取仿射开邻域
$\Spec(R) = U \subset S$，它包含 $s$；再选取仿射开邻域
$\Spec(A) = V \subset T$，它包含 $t$，并使 $V$ 映入 $U$。诱导环同态
$R \to A$ 为有限型（见引理
\ref{morphisms-lemma-locally-finite-type-characterize}），且由《性质》引理
\ref{properties-lemma-locally-jacobson}，$R$ 是 Jacobson 环。
于是结论由《代数》命题
\ref{algebra-proposition-Jacobson-permanence} 得到。
\end{proof}

\begin{lemma}
\label{morphisms-lemma-Jacobson-universally-Jacobson}
设 $S$ 为 Jacobson 概形。任意在 $S$ 上局部为有限型的概形都是
Jacobson 概形。
\end{lemma}

\begin{proof}
这由《代数》命题 \ref{algebra-proposition-Jacobson-permanence}
直接得到（同时使用《性质》引理
\ref{properties-lemma-locally-jacobson} 和引理
\ref{morphisms-lemma-locally-finite-type-characterize}）。
\end{proof}

\begin{lemma}
\label{morphisms-lemma-ubiquity-Jacobson-schemes}
下列各类概形是 Jacobson 概形：
\begin{enumerate}
\item 任意在域上局部有限型的概形。
\item 任意在 $\mathbf{Z}$ 上局部有限型的概形。
\item 任意在具有无穷多个素理想的 $1$ 维诺特整环上局部有限型的概形。
\item 形如 $\Spec(R) \setminus \{\mathfrak m\}$ 的概形，其中
$(R, \mathfrak m)$ 是诺特局部环；以及任意在它上面局部有限型的概形。
\end{enumerate}
\end{lemma}

\begin{proof}
以下不再明言地使用引理 \ref{morphisms-lemma-Jacobson-universally-Jacobson}。
域的谱显然是 Jacobson 概形。$\mathbf{Z}$ 的谱是 Jacobson 概形，
见《代数》引理 \ref{algebra-lemma-pid-jacobson}。
(3) 见《代数》引理 \ref{algebra-lemma-noetherian-dim-1-Jacobson}。
(4) 见《性质》引理
\ref{properties-lemma-complement-closed-point-Jacobson}。
\end{proof}

\section{普遍链状概形}
\label{morphisms-section-universally-catenary}

\noindent
回顾：若对拓扑空间 $X$ 的每一对不可约闭子集 $T \subset T'$，
都存在一条不可约闭子集的极大链
$$
T = T_0 \subset T_1 \subset \ldots \subset T_e = T'
$$
且每条这样的链都具有相同长度，则称该空间为{\it 链状的}。
见《拓扑学》定义 \ref{topology-definition-catenary}。
回顾：若一个概形的底拓扑空间是链状的，则称该概形为链状概形。
见《性质》定义 \ref{properties-definition-catenary}。

\begin{definition}
\label{morphisms-definition-universally-catenary}
令 $S$ 为一个概形，并假设 $S$ 局部诺特。我们称 $S$
为{\it 普遍链状概形}，如果对每个局部有限型态射
$X \to S$，概形 $X$ 都是链状概形。
\end{definition}

\noindent
这是一个比链状性“更好”的概念，因为存在链状但并非普遍链状的诺特概形。
见《例》节 \ref{examples-section-non-catenary-Noetherian-local}。
许多概形都是普遍链状的，见下文引理 \ref{morphisms-lemma-ubiquity-uc}。

\medskip\noindent
回顾：称环 $A$ 为{\it 链环}，如果对任意一对素理想
$\mathfrak p \subset \mathfrak q$，都存在一条素理想的极大链
$$
\mathfrak p =
\mathfrak p_0 \subset \ldots \subset \mathfrak p_e
= \mathfrak q
$$
且所有这样的链都具有相同长度。
见《代数》定义 \ref{algebra-definition-catenary}。
我们已在《性质》节 \ref{properties-section-catenary}
中讨论了链状概形与链环之间的关系。回顾：称环 $A$
为{\it 普遍链环}，如果 $A$ 是诺特环，并且对每个有限型环同态
$A \to B$，环 $B$ 都是链环。见《代数》定义
\ref{algebra-definition-universally-catenary}。
代数几何中出现的许多重要环都具有这一性质。

\begin{lemma}
\label{morphisms-lemma-universally-catenary-local}
令 $S$ 为局部诺特概形。下列条件等价：
\begin{enumerate}
\item $S$ 是普遍链状概形；
\item 存在 $S$ 的一个开覆盖，其每个成员都是普遍链状概形；
\item 对每个仿射开集 $\Spec(R) = U \subset S$，环 $R$
都是普遍链环；以及
\item 存在仿射开覆盖 $S = \bigcup U_i$，使得每个 $U_i$
都是一个普遍链环的谱。
\end{enumerate}
此外，在这种情形下，任何在 $S$ 上局部有限型的概形也都是普遍链状概形。
\end{lemma}

\begin{proof}
由引理 \ref{morphisms-lemma-immersion-locally-finite-type}，开浸入是局部有限型的。
局部有限型态射的复合仍然局部有限型（引理
\ref{morphisms-lemma-composition-finite-type}）。因此显然，若 $S$ 是普遍链状概形，
则其任意开子概形以及任意在 $S$ 上局部有限型的概形也都是普遍链状概形。
这就证明了引理的最后一个断言，并证明了 (1) 蕴含 (2)。

\medskip\noindent
若 $\Spec(R)$ 是普遍链状概形，则每个概形 $\Spec(A)$，
其中 $A$ 是有限型 $R$-代数，都为链状概形。因此由《代数》引理
\ref{algebra-lemma-catenary}，所有这些环 $A$ 都是链环。
故 $R$ 是普遍链环。结合以上论述，我们得到 (1) 蕴含 (3)，
而 (2) 蕴含 (4)。当然，(3) 显然蕴含 (4)。

\medskip\noindent
为完成证明，我们来证明 (4) 蕴含 (1)。假设 (4)，并令
$X \to S$ 为局部有限型态射。可以找到仿射开覆盖
$X = \bigcup V_j$，使每个 $V_j \to S$ 的像都包含于某个 $U_i$ 中。
由引理 \ref{morphisms-lemma-locally-finite-type-characterize}，诱导环同态
$\mathcal{O}(U_i) \to \mathcal{O}(V_j)$ 是有限型的。
因此 $\mathcal{O}(V_j)$ 是链环。于是由《性质》引理
\ref{properties-lemma-catenary-local}，$X$ 是链状概形。
\end{proof}

\begin{lemma}
\label{morphisms-lemma-universally-catenary-local-rings-universally-catenary}
令 $S$ 为局部诺特概形。下列条件等价：
\begin{enumerate}
\item $S$ 是普遍链状概形；以及
\item 所有局部环 $\mathcal{O}_{S, s}$（这里所指的概形是 $S$）
都是普遍链环。
\end{enumerate}
\end{lemma}

\begin{proof}
假设 $S$ 的所有局部环都是普遍链环。令 $f : X \to S$ 为局部有限型态射。
我们知道，$X$ 是链状概形，当且仅当 $\mathcal{O}_{X, x}$
对所有 $x \in X$ 都是链环。若 $f(x) = s$，则
$\mathcal{O}_{X, x}$ 在 $\mathcal{O}_{S, s}$ 上本质有限型。
因此 $\mathcal{O}_{X, x}$ 是链环，因为我们假设
$\mathcal{O}_{S, s}$ 是普遍链环。

\medskip\noindent
反过来，假设 $S$ 是普遍链状概形。令 $s \in S$。
由引理 \ref{morphisms-lemma-universally-catenary-local}，可以把 $S$
替换为 $s$ 的一个仿射开邻域。设 $S = \Spec(R)$，并设 $s$
对应素理想 $\mathfrak p$。任意有限型 $R_{\mathfrak p}$-代数
$A'$ 都具有形式 $A_{\mathfrak p}$，其中某个有限型 $R$-代数记为 $A$。
由假设（如果愿意，也可再用引理
\ref{morphisms-lemma-universally-catenary-local}），环 $A$ 是链环，
从而 $A'$（$A$ 的一个局部化）也是链环。因此
$R_{\mathfrak p}$ 是普遍链环。
\end{proof}

\begin{lemma}
\label{morphisms-lemma-catenary-check-irreducible}
令 $S$ 为局部诺特概形。则 $S$ 是普遍链状概形，当且仅当
$S$ 的各不可约分支都是普遍链状概形。
\end{lemma}

\begin{proof}
略。仿射情形见《代数》引理
\ref{algebra-lemma-catenary-check-irreducible}。
\end{proof}

\begin{lemma}
\label{morphisms-lemma-ubiquity-uc}
下列各类概形都是普遍链状概形：
\begin{enumerate}
\item 任意在域上局部有限型的概形。
\item 任意在 Cohen--Macaulay 概形上局部有限型的概形。
\item 任意在 $\mathbf{Z}$ 上局部有限型的概形。
\item 任意在 $1$ 维诺特整环上局部有限型的概形。
\item 等等。
\end{enumerate}
\end{lemma}

\begin{proof}
这些结论都源于 Cohen--Macaulay 环是普遍链环这一事实，
见《代数》引理 \ref{algebra-lemma-CM-ring-catenary}。
还要使用引理 \ref{morphisms-lemma-universally-catenary-local} 的最后一个断言。
略去一些细节。
\end{proof}

\section{再论 Nagata 概形}
\label{morphisms-section-nagata}

\noindent
关于 Nagata 概形与泛 Japanese 概形的定义和基本性质，
见《性质》节 \ref{properties-section-nagata}。

\begin{lemma}
\label{morphisms-lemma-finite-type-nagata}
令 $f : X \to S$ 为一个态射。若 $S$ 是 Nagata 概形，且 $f$
局部有限型，则 $X$ 是 Nagata 概形。若 $S$ 是泛 Japanese 概形，
且 $f$ 局部有限型，则 $X$ 是泛 Japanese 概形。
\end{lemma}

\begin{proof}
对于“泛 Japanese”，这来自《代数》引理
\ref{algebra-lemma-universally-japanese}。
对于“Nagata”，这来自《代数》命题
\ref{algebra-proposition-nagata-universally-japanese}。
\end{proof}

\begin{lemma}
\label{morphisms-lemma-ubiquity-nagata}
下列各类概形都是 Nagata 概形：
\begin{enumerate}
\item 任意在域上局部有限型的概形。
\item 任意在诺特完备局部环上局部有限型的概形。
\item 任意在 $\mathbf{Z}$ 上局部有限型的概形。
\item 任意在特征为零的 Dedekind 环上局部有限型的概形。
\item 等等。
\end{enumerate}
\end{lemma}

\begin{proof}
由引理 \ref{morphisms-lemma-finite-type-nagata}，只需证明以上提到的环都是
Nagata 环。为此见《代数》命题
\ref{algebra-proposition-ubiquity-nagata}。
\end{proof}

\section{再论奇异轨迹}
\label{morphisms-section-singular-locus}

\noindent
我们要寻找一个判据，使得对任意在基底上局部有限型的概形，其正则轨迹都是开的。
定义如下。

\begin{definition}
\label{morphisms-definition-J}
令 $X$ 为局部诺特概形。称 $X$ 是{\it J-2 的}，如果对每个局部有限型态射
$Y \to X$，正则轨迹 $\text{Reg}(Y)$ 在 $Y$ 中都是开的。
\end{definition}

\noindent
这是诺特环的相应概念的类比，见《代数续篇》定义
\ref{more-algebra-definition-J}。

\begin{lemma}
\label{morphisms-lemma-J}
令 $X$ 为局部诺特概形。下列条件等价：
\begin{enumerate}
\item $X$ 是 J-2 的；
\item 存在 $X$ 的一个开覆盖，其每个成员都是 J-2 概形；
\item 对每个仿射开集 $\Spec(R) = U \subset X$，环 $R$ 都是 J-2 的；以及
\item 存在仿射开覆盖 $S = \bigcup U_i$，使得每个
$\mathcal{O}(U_i)$ 对所有 $i$ 都是 J-2 的。
\end{enumerate}
此外，在这种情形下，任何在 $X$ 上局部有限型的概形也都是 J-2 的。
\end{lemma}

\begin{proof}
由引理 \ref{morphisms-lemma-immersion-locally-finite-type}，开浸入是局部有限型的。
局部有限型态射的复合仍然局部有限型（引理
\ref{morphisms-lemma-composition-finite-type}）。因此显然，若 $X$ 是 J-2 的，
则任意开子概形以及任意在 $X$ 上局部有限型的概形也都是 J-2 的。
这证明了引理的最后一个断言。

\medskip\noindent
若 $\Spec(R)$ 是 J-2 的，则对每个有限型 $R$-代数 $A$，
概形 $\Spec(A)$ 的正则轨迹都是开的。因此按照定义，$R$ 是 J-2 的
（见《代数续篇》定义 \ref{more-algebra-definition-J}）。
结合以上论述，我们得到 (1) 蕴含 (3)，而 (2) 蕴含 (4)。
当然，(1) $\Rightarrow$ (2) 和 (3) $\Rightarrow$ (4) 都是显然的。

\medskip\noindent
为完成证明，我们来证明 (4) 蕴含 (1)。假设 (4)，并令
$Y \to X$ 为局部有限型态射。可以找到仿射开覆盖
$Y = \bigcup V_j$，使每个 $V_j \to X$ 的像都包含于某个 $U_i$ 中。
由引理 \ref{morphisms-lemma-locally-finite-type-characterize}，诱导环同态
$\mathcal{O}(U_i) \to \mathcal{O}(V_j)$ 是有限型的。因此
$V_j = \Spec(\mathcal{O}(V_j))$ 的正则轨迹是开的。由于
$\text{Reg}(Y) \cap V_j = \text{Reg}(V_j)$，可知
$\text{Reg}(Y)$ 是开的，这正是所求。
\end{proof}

\begin{lemma}
\label{morphisms-lemma-ubiquity-J-2}
下列各类概形都是 J-2 的：
\begin{enumerate}
\item 任意在域上局部有限型的概形。
\item 任意在诺特完备局部环上局部有限型的概形。
\item 任意在 $\mathbf{Z}$ 上局部有限型的概形。
\item 任意在维数为 $1$ 的诺特局部环上局部有限型的概形。
\item 任意在维数为 $1$ 的 Nagata 环上局部有限型的概形。
\item 任意在特征为零的 Dedekind 环上局部有限型的概形。
\item 等等。
\end{enumerate}
\end{lemma}

\begin{proof}
由引理 \ref{morphisms-lemma-J}，只需证明以上提到的环都是 J-2 的。
为此见《代数续篇》命题
\ref{more-algebra-proposition-ubiquity-J-2}。
\end{proof}

\section{优秀概形}
\label{morphisms-section-excellent}

\noindent
回顾：若一个环是 G-环并且是 J-2 的，则称它为拟优秀环。
若一个环是拟优秀环并且是普遍链环，则称它为优秀环。

\begin{definition}
\label{morphisms-definition-excellent}
令 $X$ 为一个概形。
\begin{enumerate}
\item 称 $X$ 是{\it 拟优秀概形}，如果对每个 $x \in X$，都存在仿射开邻域
$x \in U \subset X$，使环 $\mathcal{O}_X(U)$ 是拟优秀环
（见《代数续篇》定义 \ref{more-algebra-definition-excellent}）。
\item 称 $X$ 是{\it 优秀概形}，如果对每个 $x \in X$，都存在仿射开邻域
$x \in U \subset X$，使环 $\mathcal{O}_X(U)$ 是优秀环
（见《代数续篇》定义 \ref{more-algebra-definition-excellent}）。
\end{enumerate}
\end{definition}

\noindent
拟优秀概形是局部诺特的（G-环是诺特环）。

\begin{lemma}
\label{morphisms-lemma-characterize-quasi-excellent}
令 $X$ 为一个概形。下列条件等价：
\begin{enumerate}
\item $X$ 是拟优秀概形；以及
\item $X$ 是 G-概形并且是 J-2 的。
\end{enumerate}
\end{lemma}

\begin{proof}
将拟优秀概形的定义（定义 \ref{morphisms-definition-excellent}）、
拟优秀环的定义（《代数续篇》定义
\ref{more-algebra-definition-excellent}）、G-概形的定义
（《性质》定义 \ref{properties-definition-G-scheme}）以及引理
\ref{morphisms-lemma-J} 结合起来，即得 (1) $\Rightarrow$ (2)。
反过来，若 $X$ 是 G-概形并且是 J-2 的，则对任意 $x \in X$，
都可以找到仿射开邻域 $x \in U \subset X$，使
$\mathcal{O}_X(U)$ 是 G-环。由引理 \ref{morphisms-lemma-J}，环
$\mathcal{O}_X(U)$ 也是 J-2 的，因而是拟优秀环。
\end{proof}

\begin{lemma}
\label{morphisms-lemma-quasi-excellent-nagata}
拟优秀概形是 Nagata 概形。
\end{lemma}

\begin{proof}
见《代数续篇》引理
\ref{more-algebra-lemma-quasi-excellent-nagata}。
\end{proof}

\begin{lemma}
\label{morphisms-lemma-characterize-excellent}
令 $X$ 为一个概形。下列条件等价：
\begin{enumerate}
\item $X$ 是优秀概形；以及
\item $X$ 是拟优秀概形并且是普遍链状概形。
\end{enumerate}
\end{lemma}

\begin{proof}
假设 (1)。由于优秀环是拟优秀环，显然 $X$ 是拟优秀概形。
由于优秀环是普遍链环，由引理
\ref{morphisms-lemma-universally-catenary-local} 还可推出 $X$ 是普遍链状概形。
假设 (2)。令 $x \in X$。由于 $X$ 是拟优秀概形，存在仿射开邻域
$x \in U \subset X$，使 $\mathcal{O}_X(U)$ 是拟优秀环。
由于 $X$ 是普遍范畴的，由引理
\ref{morphisms-lemma-universally-catenary-local}，环 $\mathcal{O}_X(U)$
也是普遍链环，因而是优秀环。
\end{proof}

\begin{lemma}
\label{morphisms-lemma-locally-excellent}
令 $X$ 为一个概形。下列条件等价：
\begin{enumerate}
\item 概形 $X$ 是（拟）优秀概形。
\item 对每个仿射开集 $U \subset X$，环 $\mathcal{O}_X(U)$
都是（拟）优秀环。
\item 存在仿射开覆盖 $X = \bigcup U_i$，使每个
$\mathcal{O}_X(U_i)$ 都是（拟）优秀环。
\item 存在开覆盖 $X = \bigcup X_j$，使每个开子概形
$X_j$ 都是（拟）优秀概形。
\end{enumerate}
此外，若 $X$ 是（拟）优秀概形，则每个开子概形也都是（拟）优秀概形。
\end{lemma}

\begin{proof}
由引理 \ref{morphisms-lemma-characterize-quasi-excellent}，拟优秀等价于既是
G-概形又是 J-2 的。因此在拟优秀情形，该引理由引理
\ref{morphisms-lemma-J} 和《性质》引理
\ref{properties-lemma-locally-G-scheme} 得出。
在优秀情形，使用引理 \ref{morphisms-lemma-characterize-excellent}、
拟优秀情形以及引理 \ref{morphisms-lemma-universally-catenary-local}。
\end{proof}

\begin{lemma}
\label{morphisms-lemma-finite-type-excellent}
令 $f : X \to S$ 为一个态射。若 $S$ 是（拟）优秀概形且 $f$
局部有限型，则 $X$ 是（拟）优秀概形。
\end{lemma}

\begin{proof}
见《代数续篇》引理
\ref{more-algebra-lemma-finite-type-over-excellent}。
\end{proof}

\begin{lemma}
\label{morphisms-lemma-ubiquity-excellent}
下列各类概形都是优秀概形：
\begin{enumerate}
\item 任意在域上局部有限型的概形。
\item 任意在诺特完备局部环上局部有限型的概形。
\item 任意在 $\mathbf{Z}$ 上局部有限型的概形。
\item 任意在特征为零的 Dedekind 环上局部有限型的概形。
\end{enumerate}
\end{lemma}

\begin{proof}
由引理 \ref{morphisms-lemma-locally-excellent} 和
\ref{morphisms-lemma-finite-type-excellent}，只需证明以上提到的环都是优秀环。
为此见《代数续篇》命题
\ref{more-algebra-proposition-ubiquity-excellent}。
\end{proof}

\section{拟有限态射}
\label{morphisms-section-quasi-finite}

\noindent
对拟有限态射进行扎实处理，是后续许多发展的基础。这将导向 Zariski
主定理的各种版本、纤维维数的行为、\'etale 态射的下降等诸多内容。
阅读本节之前，最好先看一下《代数》节
\ref{algebra-section-quasi-finite} 中的代数结果。

\medskip\noindent
回顾：若有限型环同态 $R \to A$ 的纤维中由素理想 $\mathfrak q$
确定的点是孤立点，则称该环同态在 $\mathfrak q$ 处拟有限；见《代数》定义
\ref{algebra-definition-quasi-finite}。

\begin{definition}
\label{morphisms-definition-quasi-finite}
\begin{reference}
\cite[II Definition 6.2.3]{EGA}
\end{reference}
令 $f : X \to S$ 为概形态射。
\begin{enumerate}
\item 称 $f$ 在点 $x \in X$ 处{\it 拟有限}，如果存在仿射邻域
$\Spec(A) = U \subset X$（它是 $x$ 的邻域）以及仿射开集
$\Spec(R) = V \subset S$，使得 $f(U) \subset V$，环同态
$R \to A$ 是有限型的，并且 $R \to A$ 在 $A$ 中对应于 $x$
的素理想处拟有限（见上文）。
\item 称 $f$ {\it 局部拟有限}，如果 $f$ 在每一点 $x$
（它属于 $X$）处都拟有限。
\item 称 $f$ {\it 拟有限}，如果 $f$ 是有限型的，并且每一点 $x$
都是其所在纤维的孤立点。
\end{enumerate}
\end{definition}

\noindent
显然，局部拟有限态射是局部有限型的。下文将看到，对一个局部有限型态射
$f$，它在 $x$ 处拟有限，当且仅当 $x$ 是其所在纤维的孤立点。
此外，一个态射拟有限的点所成的集合是开的；我们将在关于 Zariski
主定理的节 \ref{morphisms-section-Zariski} 中看到这一点。

\begin{lemma}
\label{morphisms-lemma-algebraic-residue-field-extension-closed-point-fibre}
令 $f : X \to S$ 为概形态射。令 $x \in X$ 为一点，并置
$s = f(x)$。若 $\kappa(x)/\kappa(s)$ 是代数域扩张，则：
\begin{enumerate}
\item $x$ 是其所在纤维的闭点；以及
\item 若此外 $s$ 是 $S$ 的闭点，则 $x$ 是 $X$ 的闭点。
\end{enumerate}
\end{lemma}

\begin{proof}
第二个断言由第一个断言和初等拓扑立即得到。根据《概形》引理
\ref{schemes-lemma-fibre-topological}，为了证明第一个断言，
可以把 $X$ 替换为 $X_s$，把 $S$ 替换为 $\Spec(\kappa(s))$。
因此可以假设 $S = \Spec(k)$ 是一个域的谱。在这种情形下，令
$\Spec(A) = U \subset X$ 为任意包含 $x$ 的仿射开集。点 $x$
对应于素理想 $\mathfrak q \subset A$，且
$\kappa(\mathfrak q)/k$ 是代数域扩张。由《代数》引理
\ref{algebra-lemma-finite-residue-extension-closed}，可知
$\mathfrak q$ 是极大理想，即 $x \in U$ 是闭点。由于仿射开集构成
$X$ 的拓扑基，可知 $\{x\}$ 是闭集。
\end{proof}

\noindent
下一个引理是 Hilbert 零点定理的一个版本。

\begin{lemma}
\label{morphisms-lemma-closed-point-fibre-locally-finite-type}
令 $f : X \to S$ 为概形态射。令 $x \in X$ 为一点，并置
$s = f(x)$。假设 $f$ 局部有限型。则 $x$ 是其所在纤维的闭点，
当且仅当 $\kappa(x)/\kappa(s)$ 是有限域扩张。
\end{lemma}

\begin{proof}
若该扩张有限，则由上面的引理
\ref{morphisms-lemma-algebraic-residue-field-extension-closed-point-fibre}，
$x$ 是其所在纤维的闭点。反过来，假设 $x$ 是其所在纤维的闭点。
选取仿射开集 $\Spec(A) = U \subset X$ 和
$\Spec(R) = V \subset S$，使得 $f(U) \subset V$。
由引理 \ref{morphisms-lemma-locally-finite-type-characterize}，环同态
$R \to A$ 是有限型的。令 $\mathfrak q \subset A$ 和
$\mathfrak p \subset R$ 分别为对应于 $x$ 和 $s$ 的素理想。
考虑纤维环 $\overline{A} = A \otimes_R \kappa(\mathfrak p)$。
令 $\overline{\mathfrak q}$ 为 $\overline{A}$ 中对应于
$\mathfrak q$ 的素理想。$x$ 是其所在纤维闭点这一假设蕴含
$\overline{\mathfrak q}$ 是 $\overline{A}$ 的极大理想。由于
$\overline{A}$ 是域 $\kappa(\mathfrak p)$ 上的有限型代数，
由 Hilbert 零点定理（见《代数》定理
\ref{algebra-theorem-nullstellensatz}），可知
$\kappa(\overline{\mathfrak q})$ 是 $\kappa(\mathfrak p)$
的有限扩张。由于 $\kappa(s) = \kappa(\mathfrak p)$ 且
$\kappa(x) = \kappa(\mathfrak q) = \kappa(\overline{\mathfrak q})$，
结论成立。
\end{proof}

\begin{lemma}
\label{morphisms-lemma-base-change-closed-point-fibre-locally-finite-type}
令 $f : X \to S$ 为局部有限型的概形态射。令 $g : S' \to S$
为任意态射，并记其基变换为 $f' : X' \to S'$。若
$x' \in X'$ 映到的点 $x \in X$ 在 $X_{f(x)}$ 中是闭点，
则 $x'$ 在 $X'_{f'(x')}$ 中是闭点。
\end{lemma}

\begin{proof}
剩余域 $\kappa(x')$ 是
$\kappa(f'(x')) \otimes_{\kappa(f(x))} \kappa(x)$ 的一个商，
见《概形》引理 \ref{schemes-lemma-points-fibre-product}。
因此它是 $\kappa(f'(x'))$ 的有限扩张，因为由引理
\ref{morphisms-lemma-closed-point-fibre-locally-finite-type}，
$\kappa(x)$ 是 $\kappa(f(x))$ 的有限扩张。
再应用该引理一次，便可知 $x'$ 在其所在纤维中是闭点。
\end{proof}

\begin{lemma}
\label{morphisms-lemma-residue-field-quasi-finite}
令 $f : X \to S$ 为概形态射。令 $x \in X$ 为一点，并置
$s = f(x)$。若 $f$ 在 $x$ 处拟有限，则剩余域扩张
$\kappa(x)/\kappa(s)$ 是有限的。
\end{lemma}

\begin{proof}
这直接来自《代数》定义 \ref{algebra-definition-quasi-finite}。
\end{proof}

\begin{lemma}
\label{morphisms-lemma-quasi-finite-at-point-characterize}
令 $f : X \to S$ 为概形态射。令 $x \in X$ 为一点，并置
$s = f(x)$。令 $X_s$ 为 $f$ 在 $s$ 处的纤维。假设 $f$
局部有限型。下列条件等价：
\begin{enumerate}
\item 态射 $f$ 在 $x$ 处拟有限。
\item 点 $x$ 在 $X_s$ 中是孤立点。
\item 点 $x$ 在 $X_s$ 中是闭点，并且不存在点
$x' \in X_s$ 满足 $x' \not = x$ 且特化到 $x$。
\item 对任意一对仿射开集
$\Spec(A) = U \subset X$、$\Spec(R) = V \subset S$，若
$f(U) \subset V$，且 $x \in U$ 对应于 $\mathfrak q \subset A$，
则环同态 $R \to A$ 在 $\mathfrak q$ 处拟有限。
\end{enumerate}
\end{lemma}

\begin{proof}
假设 $f$ 在 $x$ 处拟有限。由假设，存在开集
$U \subset X$、$V \subset S$，使得 $f(U) \subset V$、$x \in U$，
并且 $x$ 是 $U_s$ 的孤立点。因此 $\{x\} \subset U_s$ 是开子集。
由于 $U_s = U \cap X_s \subset X_s$ 也是开的，可知
$\{x\} \subset X_s$ 也是开子集。所以 $x$ 是 $X_s$ 的孤立点。

\medskip\noindent
注意，由引理 \ref{morphisms-lemma-ubiquity-Jacobson-schemes}
（以及引理 \ref{morphisms-lemma-base-change-finite-type}），$X_s$
是 Jacobson 概形。若 $x$ 在 $X_s$ 中是孤立点，即
$\{x\} \subset X_s$ 是开的，则 $\{x\}$ 包含一个闭点
（由 Jacobson 性质），因此 $x$ 在 $X_s$ 中是闭点。
显然，不存在点 $x' \in X_s$，它不同于 $x$ 且特化到 $x$。

\medskip\noindent
假设 $x$ 在 $X_s$ 中是闭点，并且不存在点 $x' \in X_s$，
它不同于 $x$ 且特化到 $x$。考虑一对仿射开集
$\Spec(A) = U \subset X$、$\Spec(R) = V \subset S$，满足
$f(U) \subset V$ 且 $x \in U$。令 $\mathfrak q \subset A$
对应于 $x$，并令 $\mathfrak p \subset R$ 对应于 $s$。
由引理 \ref{morphisms-lemma-locally-finite-type-characterize}，环同态
$R \to A$ 是有限型的。考虑纤维环
$\overline{A} = A \otimes_R \kappa(\mathfrak p)$。
令 $\overline{\mathfrak q}$ 为 $\overline{A}$ 中对应于
$\mathfrak q$ 的素理想。由于 $\Spec(\overline{A})$ 是纤维
$X_s$ 的开子概形，可知 $\overline{q}$ 是 $\overline{A}$ 的极大理想，
并且 $\Spec(\overline{A})$ 中没有点特化到
$\overline{\mathfrak q}$。这蕴含
$\dim(\overline{A}_{\overline{q}}) = 0$。因此由《代数》定义
\ref{algebra-definition-quasi-finite}，可知 $R \to A$ 在
$\mathfrak q$ 处拟有限；按照定义，这就是说 $X \to S$ 在 $x$
处拟有限。

\medskip\noindent
至此，我们已经证明条件 (1)--(3) 等价。显然 (4) 蕴含 (1)。
(2) 也显然蕴含 (4)，因为若 $x$ 是 $X_s$ 的孤立点，则它在
$U_s$ 中也为孤立点，其中 $U$ 是任意包含它的开集。
\end{proof}

\begin{lemma}
\label{morphisms-lemma-finite-fibre}
令 $f : X \to S$ 为概形态射，并令 $s \in S$。假设：
\begin{enumerate}
\item $f$ 局部有限型；以及
\item $f^{-1}(\{s\})$ 是有限集。
\end{enumerate}
则 $X_s$ 是有限离散拓扑空间，并且 $f$ 在 $X$ 中位于 $s$
上方的每一点处都拟有限。
\end{lemma}

\begin{proof}
设 $T$ 为一个概形，它 (a) 在域 $k$ 上局部有限型，并且 (b) 只有有限多个点。
则引理 \ref{morphisms-lemma-ubiquity-Jacobson-schemes} 表明 $T$ 是 Jacobson 概形。
有限 Jacobson 空间是离散的，见《拓扑学》引理
\ref{topology-lemma-finite-jacobson}。把这一论述应用于纤维 $X_s$；
它在 $\Spec(\kappa(s))$ 上局部有限型，由此得到第一个断言。
最后，应用引理 \ref{morphisms-lemma-quasi-finite-at-point-characterize}
即得第二个断言。
\end{proof}

\begin{lemma}
\label{morphisms-lemma-locally-quasi-finite-fibres}
\begin{slogan}
具有离散纤维的有限型态射是拟有限的。
\end{slogan}
令 $f : X \to S$ 为概形态射。假设 $f$ 局部有限型。
则下列条件等价：
\begin{enumerate}
\item $f$ 局部拟有限；
\item 对每个 $s \in S$，纤维 $X_s$ 都是离散拓扑空间；以及
\item 对每个态射 $\Spec(k) \to S$（其中 $k$ 是域），基变换
$X_k$ 的底拓扑空间都是离散的。
\end{enumerate}
\end{lemma}

\begin{proof}
(3) 蕴含 (2) 是立即的。引理
\ref{morphisms-lemma-quasi-finite-at-point-characterize} 表明 (2) 与 (1) 等价。
假设 (2)，并令 $\Spec(k) \to S$ 如 (3) 所述。以 $s \in S$
记 $\Spec(k) \to S$ 的像。则 $X_k$ 是 $X_s$ 经
$\Spec(k) \to \Spec(\kappa(s))$ 的基变换。因此由引理
\ref{morphisms-lemma-base-change-closed-point-fibre-locally-finite-type}，
$X_k$ 的每一点都是闭点。由于 $X_k \to \Spec(k)$ 局部有限型
（由引理 \ref{morphisms-lemma-base-change-finite-type}），可以应用引理
\ref{morphisms-lemma-quasi-finite-at-point-characterize}，得到 $X_k$ 的每一点
都是孤立点，即 $X_k$ 的底拓扑空间是离散的。
\end{proof}

\begin{lemma}
\label{morphisms-lemma-quasi-finite-locally-quasi-compact}
令 $f : X \to S$ 为概形态射。则 $f$ 拟有限，当且仅当 $f$
局部拟有限且拟紧。
\end{lemma}

\begin{proof}
假设 $f$ 拟有限。由定义 \ref{morphisms-definition-finite-type}，它是拟紧的。
令 $x \in X$。由引理 \ref{morphisms-lemma-quasi-finite-at-point-characterize}，
可知 $f$ 在 $x$ 处拟有限。因此 $f$ 既拟紧又局部拟有限。

\medskip\noindent
假设 $f$ 拟紧且局部拟有限。则 $f$ 是有限型的。令 $x \in X$
为一点。由引理 \ref{morphisms-lemma-quasi-finite-at-point-characterize}，
可知 $x$ 是其所在纤维的孤立点。引理得证。
\end{proof}

\begin{lemma}
\label{morphisms-lemma-quasi-finite}
令 $f : X \to S$ 为概形态射。下列条件等价：
\begin{enumerate}
\item $f$ 拟有限；以及
\item $f$ 局部有限型、拟紧，并且具有有限纤维。
\end{enumerate}
\end{lemma}

\begin{proof}
假设 $f$ 拟有限。特别地，$f$ 局部有限型且拟紧（因为它是有限型的）。
令 $s \in S$。由于每个 $x \in X_s$ 都是 $X_s$ 的孤立点，可见
$X_s = \bigcup_{x \in X_s} \{x\}$ 是一个开覆盖。由于 $f$ 拟紧，
纤维 $X_s$ 是拟紧的。因此 $X_s$ 是有限的。

\medskip\noindent
反过来，假设 $f$ 局部有限型、拟紧且具有有限纤维。
由引理 \ref{morphisms-lemma-finite-fibre}，它局部拟有限。因此由引理
\ref{morphisms-lemma-quasi-finite-locally-quasi-compact}，它拟有限。
\end{proof}

\noindent
回顾：若环同态 $R \to A$ 是有限型的，并且在 $A$ 的{\it 所有}素理想处
都拟有限，则称该环同态拟有限；见《代数》定义
\ref{algebra-definition-quasi-finite}。

\begin{lemma}
\label{morphisms-lemma-locally-quasi-finite-characterize}
令 $f : X \to S$ 为概形态射。下列条件等价：
\begin{enumerate}
\item 态射 $f$ 局部拟有限。
\item 对每一对仿射开集 $U \subset X$、$V \subset S$，若
$f(U) \subset V$，则环同态
$\mathcal{O}_S(V) \to \mathcal{O}_X(U)$ 都拟有限。
\item 存在开覆盖 $S = \bigcup_{j \in J} V_j$，以及开覆盖
$f^{-1}(V_j) = \bigcup_{i \in I_j} U_i$，使每个态射
$U_i \to V_j$（$j\in J, i\in I_j$）都局部拟有限。
\item 存在仿射开覆盖 $S = \bigcup_{j \in J} V_j$，以及仿射开覆盖
$f^{-1}(V_j) = \bigcup_{i \in I_j} U_i$，使环同态
$\mathcal{O}_S(V_j) \to \mathcal{O}_X(U_i)$ 对所有
$j\in J, i\in I_j$ 都拟有限。
\end{enumerate}
此外，若 $f$ 局部拟有限，则对任何开子概形
$U \subset X$、$V \subset S$，若 $f(U) \subset V$，则限制
$f|_U : U \to V$ 都局部拟有限。
\end{lemma}

\begin{proof}
对环同态 $R \to A$，定义 $P(R \to A)$ 表示“$R \to A$ 拟有限”
（见引理之前的评注）。我们断言 $P$ 是环同态的局部性质。
来验证定义 \ref{morphisms-definition-property-local} 的条件 (a)、(b) 和 (c)。
在引理 \ref{morphisms-lemma-locally-finite-type-characterize} 的证明中，
我们已经看到性质“有限型”满足 (a)、(b) 和 (c)。注意，对于有限型环同态
$R \to A$，$R \to A$ 在 $\mathfrak q$ 处拟有限这一性质，只依赖于局部环
$A_{\mathfrak q}$ 作为 $R_{\mathfrak p}$-代数的结构，其中
$\mathfrak p = R \cap \mathfrak q$（通常的记号滥用）。由这些评注，
定义 \ref{morphisms-definition-property-local} 的 (a)、(b) 和 (c) 立即成立。
例如，设 $R \to A$ 为环同态，并且所有环同态
$R \to A_{a_i}$ 都拟有限，其中 $a_1, \ldots, a_n \in A$
生成单位理想。
可知 $R \to A$ 是有限型的。此外，对任意素理想
$\mathfrak q \subset A$，局部环 $A_{\mathfrak q}$ 作为 $R$-代数，
同构于局部环 $(A_{a_i})_{\mathfrak q_i}$，其中 $i$ 是某个指标，
且 $\mathfrak q_i \subset A_{a_i}$ 是某个素理想。因此
$R \to A$ 在 $\mathfrak q$ 处拟有限。

\medskip\noindent
可知，引理 \ref{morphisms-lemma-locally-P} 可应用于上一段中的 $P$。
因此只需证明：$f$ 局部拟有限，等价于 $f$ 局部为
$P$ 型。由于
$P(R \to A)$ 表示“$R \to A$ 拟有限”，也就是 $R \to A$
在 $A$ 的每个素理想处拟有限，这来自引理
\ref{morphisms-lemma-quasi-finite-at-point-characterize}。
\end{proof}

\begin{lemma}
\label{morphisms-lemma-composition-quasi-finite}
两个局部拟有限态射的复合仍然局部拟有限。对拟有限态射也有相同结论。
\end{lemma}

\begin{proof}
在引理 \ref{morphisms-lemma-locally-quasi-finite-characterize} 的证明中，
我们看到 $P = $“拟有限”是环同态的局部性质，并且概形态射局部拟有限，
当且仅当它局部为 $P$ 型（如定义 \ref{morphisms-definition-locally-P}）。
因此，引理的第一个断言来自引理 \ref{morphisms-lemma-composition-type-P}，
以及拟有限是一个在复合下稳定的环同态性质这一事实；见《代数》引理
\ref{algebra-lemma-quasi-finite-composition}。
结合上述结论、引理 \ref{morphisms-lemma-quasi-finite-locally-quasi-compact}，
以及拟紧态射的复合拟紧这一事实（见《概形》引理
\ref{schemes-lemma-composition-quasi-compact}），可知拟有限态射的复合
仍然拟有限。
\end{proof}

\noindent
稍后将看到（引理 \ref{morphisms-lemma-quasi-finite-points-open}），
下一个引理中的集合 $U$ 是开的。

\begin{lemma}
\label{morphisms-lemma-base-change-quasi-finite}
\begin{slogan}
（局部）拟有限态射在基变换下稳定。
\end{slogan}
令 $f : X \to S$ 为概形态射。令 $g : S' \to S$ 为概形态射。
以 $f' : X' \to S'$ 记 $f$ 经 $g$ 的基变换，
并以 $g' : X' \to X$ 记投影。假设 $X$ 在 $S$ 上局部有限型。
\begin{enumerate}
\item 令 $U \subset X$（相应地，$U' \subset X'$）为 $f$
（相应地，$f'$）拟有限的点所成的集合。则
$U' = U \times_S S' = (g')^{-1}(U)$。
\item 局部拟有限态射的基变换局部拟有限。
\item 拟有限态射的基变换拟有限。
\end{enumerate}
\end{lemma}

\begin{proof}
第一和第二个断言来自相应的代数结果，见《代数》引理
\ref{algebra-lemma-quasi-finite-base-change}（并结合引理
\ref{morphisms-lemma-base-change-finite-type} 所给出的 $f'$ 也局部有限型这一事实）。
结合上述结论、引理 \ref{morphisms-lemma-quasi-finite-locally-quasi-compact}，
以及拟紧态射的基变换拟紧这一事实（见《概形》引理
\ref{schemes-lemma-quasi-compact-preserved-base-change}），可知拟有限态射的
基变换拟有限。
\end{proof}

\begin{lemma}
\label{morphisms-lemma-quasi-finite-at-a-finite-number-of-points}
令 $f : X \to S$ 为有限型概形态射，并令 $s \in S$。
$X$ 中位于 $s$ 上方且 $f$ 在其处拟有限的点至多只有有限多个。
\end{lemma}

\begin{proof}
纤维 $X_s$ 是域上的有限型概形，因而是诺特概形
（引理 \ref{morphisms-lemma-finite-type-noetherian}）。所以 $X_s$
上的拓扑是诺特的（《性质》引理
\ref{properties-lemma-Noetherian-topology}），由初等拓扑可知它至多有
有限多个孤立点。因此本引理由引理
\ref{morphisms-lemma-quasi-finite-at-point-characterize} 得出。
\end{proof}

\begin{lemma}
\label{morphisms-lemma-monomorphism-loc-finite-type-loc-quasi-finite}
令 $f : X \to Y$ 为概形态射。若 $f$ 局部有限型且为单态射，
则 $f$ 是分离且局部拟有限的。
\end{lemma}

\begin{proof}
由《概形》引理 \ref{schemes-lemma-monomorphism-separated}，
单态射是分离的。单态射是单射，因此例如由引理
\ref{morphisms-lemma-quasi-finite-at-point-characterize}，可知 $f$
在每个 $x \in X$ 处拟有限。
\end{proof}

\begin{lemma}
\label{morphisms-lemma-immersion-locally-quasi-finite}
任意浸入都局部拟有限。
\end{lemma}

\begin{proof}
这是因为开浸入是局部同构，而闭浸入显然拟有限。
\end{proof}

\begin{lemma}
\label{morphisms-lemma-permanence-quasi-finite}
令 $X \to Y$ 为基概形 $S$ 上的概形态射。令 $x \in X$。
若 $X \to S$ 在 $x$ 处拟有限，则 $X \to Y$ 在 $x$ 处拟有限。
若 $X$ 在 $S$ 上局部拟有限，则 $X \to Y$ 局部拟有限。
\end{lemma}

\begin{proof}
通过引理 \ref{morphisms-lemma-locally-quasi-finite-characterize}，这转化为如下代数事实：
给定环同态 $A \to B \to C$，若 $A \to C$ 拟有限，则
$B \to C$ 拟有限。这来自《代数》引理
\ref{algebra-lemma-four-rings}，其中取 $R = A$、$S = S' = C$
和 $R' = B$。
\end{proof}

\begin{lemma}
\label{morphisms-lemma-quasi-finite-local-source}
令 $f : X \to Y$ 和 $g : Y \to S$ 为概形态射。若 $f$ 满射、
$g \circ f$ 局部拟有限且 $g$ 局部有限型，则 $g : Y \to S$
局部拟有限。
\end{lemma}

\begin{proof}
令 $x \in X$ 的像依次为 $y \in Y$ 和 $s \in S$。
由于 $g \circ f$ 局部拟有限，由引理
\ref{morphisms-lemma-residue-field-quasi-finite}，扩张
$\kappa(x)/\kappa(s)$ 是有限的。因此
$\kappa(y)/\kappa(s)$ 是有限的。由引理
\ref{morphisms-lemma-algebraic-residue-field-extension-closed-point-fibre}，
$y$ 是 $Y_s$ 的闭点。由于 $f$ 满射，可知 $Y$ 的每一点在其
$S$ 上的纤维中都是闭点。因此由引理
\ref{morphisms-lemma-quasi-finite-at-point-characterize}，可知 $g$
在每一点处拟有限。
\end{proof}

\section{有限表示态射}
\label{morphisms-section-finite-presentation}

\noindent
回顾：若环同态 $R \to A$ 使得 $A$ 同构于
$R[x_1, \ldots, x_n]/(f_1, \ldots, f_m)$，且这是一个 $R$-代数同构，
其中 $n, m$
为某些整数且 $f_j$ 为某些多项式，则称该环同态为有限表示的；
见《代数》定义 \ref{algebra-definition-finite-type}。

\begin{definition}
\label{morphisms-definition-finite-presentation}
令 $f : X \to S$ 为概形态射。
\begin{enumerate}
\item 称 $f$ 在 $x \in X$ 处{\it 有限表示}，如果存在仿射开邻域
$\Spec(A) = U \subset X$（它是 $x$ 的邻域）以及仿射开集
$\Spec(R) = V \subset S$，满足 $f(U) \subset V$，使得诱导环同态
$R \to A$ 是有限表示的。
\item 称 $f$ {\it 局部有限表示}，如果它在 $X$ 的每一点处都有限表示。
\item 称 $f$ {\it 有限表示}，如果它局部有限表示、拟紧且拟分离。
\end{enumerate}
\end{definition}

\noindent
注意，有限表示态射{\bf 不}只是局部有限表示的拟紧态射。
稍后我们将把局部有限表示态射刻画为如下态射：对于仿射概形的
任意有向系 $T_i$（它们都在 $S$ 上），都有
$$
\colim \Mor_S(T_i, X) = \Mor_S(\lim T_i, X)
$$
成立。见《极限》命题
\ref{limits-proposition-characterize-locally-finite-presentation}。
在《极限》节 \ref{limits-section-descending-relative} 中，我们证明：
若 $S = \lim_i S_i$ 是仿射概形的极限，则任意概形 $X$，若它在
$S$ 上有限表示，都能下降为概形 $X_i$；它在 $S_i$ 上，其中 $i$
为某个指标。

\begin{lemma}
\label{morphisms-lemma-locally-finite-presentation-characterize}
令 $f : X \to S$ 为概形态射。下列条件等价：
\begin{enumerate}
\item 态射 $f$ 局部有限表示。
\item 对每一对仿射开集 $U \subset X$、$V \subset S$，若
$f(U) \subset V$，则环同态
$\mathcal{O}_S(V) \to \mathcal{O}_X(U)$ 是有限表示的。
\item 存在开覆盖 $S = \bigcup_{j \in J} V_j$，以及开覆盖
$f^{-1}(V_j) = \bigcup_{i \in I_j} U_i$，使每个态射
$U_i \to V_j$（$j\in J, i\in I_j$）都局部有限表示。
\item 存在仿射开覆盖 $S = \bigcup_{j \in J} V_j$，以及仿射开覆盖
$f^{-1}(V_j) = \bigcup_{i \in I_j} U_i$，使环同态
$\mathcal{O}_S(V_j) \to \mathcal{O}_X(U_i)$ 对所有
$j\in J, i\in I_j$ 都是有限表示的。
\end{enumerate}
此外，若 $f$ 局部有限表示，则对任何开子概形 $U \subset X$、
$V \subset S$，若 $f(U) \subset V$，则限制 $f|_U : U \to V$
局部有限表示。
\end{lemma}

\begin{proof}
只要证明性质“$R \to A$ 是有限表示的”为局部性质，结论便来自引理
\ref{morphisms-lemma-locally-P-characterize}。来验证定义
\ref{morphisms-definition-property-local} 的条件 (a)、(b) 和 (c)。
由《代数》引理 \ref{algebra-lemma-base-change-finiteness}，
有限表示性在基变换下稳定，故 (a) 成立。由《代数》引理
\ref{algebra-lemma-compose-finite-type}，有限表示性在复合下稳定；
并且对任意环 $R$，环同态 $R \to R_f$ 显然是有限表示的。
故 (b) 成立。最后，由《代数》引理
\ref{algebra-lemma-cover-upstairs}，性质 (c) 成立。
\end{proof}

\begin{lemma}
\label{morphisms-lemma-composition-finite-presentation}
两个局部有限表示态射的复合仍然局部有限表示。
有限表示态射也有同样的结论。
\end{lemma}

\begin{proof}
在引理 \ref{morphisms-lemma-locally-finite-presentation-characterize} 的证明中，
我们看到有限表示性是环同态的局部性质。因此引理的第一个断言来自引理
\ref{morphisms-lemma-composition-type-P}，以及有限表示性是一个在复合下稳定的
环同态性质这一事实；见《代数》引理
\ref{algebra-lemma-compose-finite-type}。
结合上述结论以及拟紧拟分离态射的复合仍然拟紧拟分离这一事实
（见《概形》引理 \ref{schemes-lemma-composition-quasi-compact} 和
\ref{schemes-lemma-separated-permanence}），可知有限表示态射的复合
仍然有限表示。
\end{proof}

\begin{lemma}
\label{morphisms-lemma-base-change-finite-presentation}
局部有限表示态射的基变换仍然局部有限表示。
有限表示态射也有同样的结论。
\end{lemma}

\begin{proof}
在引理 \ref{morphisms-lemma-locally-finite-presentation-characterize} 的证明中，
我们看到有限表示性是环同态的局部性质。因此引理的第一个断言来自引理
\ref{morphisms-lemma-composition-type-P}，以及有限表示性是一个在基变换下稳定的
环同态性质这一事实；见《代数》引理
\ref{algebra-lemma-base-change-finiteness}。
结合上述结论以及拟紧拟分离态射的基变换仍然拟紧拟分离这一事实
（见《概形》引理
\ref{schemes-lemma-quasi-compact-preserved-base-change} 和
\ref{schemes-lemma-separated-permanence}），可知有限表示态射的基变换
仍然是有限表示态射。
\end{proof}

\begin{lemma}
\label{morphisms-lemma-open-immersion-locally-finite-presentation}
任意开浸入都局部有限表示。
\end{lemma}

\begin{proof}
这是因为开浸入是局部同构。
\end{proof}

\begin{lemma}
\label{morphisms-lemma-quasi-compact-open-immersion-finite-presentation}
开浸入有限表示，当且仅当它拟紧。
\end{lemma}

\begin{proof}
我们已经看到（引理
\ref{morphisms-lemma-open-immersion-locally-finite-presentation}），开浸入局部有限表示。
我们也已经看到（《概形》引理
\ref{schemes-lemma-immersions-monomorphisms}），浸入是分离的，因而拟分离。
结合这一点和定义 \ref{morphisms-definition-finite-presentation}，即得结论。
\end{proof}

\begin{lemma}
\label{morphisms-lemma-closed-immersion-finite-presentation}
\begin{slogan}
有限表示的闭浸入对应于有限型拟凝聚理想层。
\end{slogan}
闭浸入 $i : Z \to X$ 有限表示，当且仅当相伴的拟凝聚理想层
$\mathcal{I} = \Ker(\mathcal{O}_X \to i_*\mathcal{O}_Z)$
作为 $\mathcal{O}_X$-模是有限型的。
\end{lemma}

\begin{proof}
在任意仿射开集 $\Spec(R) \subset X$ 上，我们有
$i^{-1}(\Spec(R)) = \Spec(R/I)$ 和
$\mathcal{I} = \widetilde{I}$。此外，$\mathcal{I}$ 是有限型的，
当且仅当对每个这样的仿射开集，$I$ 都是有限 $R$-模
（见《性质》引理 \ref{properties-lemma-finite-type-module}）。
而 $R/I$ 在 $R$ 上有限表示，当且仅当 $I$ 是有限 $R$-模。
因此结论成立。
\end{proof}

\begin{lemma}
\label{morphisms-lemma-finite-presentation-finite-type}
局部有限表示态射局部有限型。有限表示态射有限型。
\end{lemma}

\begin{proof}
略。
\end{proof}

\begin{lemma}
\label{morphisms-lemma-noetherian-finite-type-finite-presentation}
\begin{slogan}
在局部诺特基底上，有限型就是有限表示。
\end{slogan}
令 $f : X \to S$ 为一个态射。
\begin{enumerate}
\item 若 $S$ 局部诺特且 $f$ 局部有限型，则 $f$ 局部有限表示。
\item 若 $S$ 局部诺特且 $f$ 有限型，则 $f$ 有限表示。
\end{enumerate}
\end{lemma}

\begin{proof}
第一个断言来自以下事实：诺特环上的有限型环是有限表示的；见《代数》引理
\ref{algebra-lemma-Noetherian-finite-type-is-finite-presentation}。
假设 $f$ 有限型且 $S$ 局部诺特。则由 (1)，$f$ 拟紧且局部有限表示。
所以只需证明 $f$ 拟分离。这来自引理
\ref{morphisms-lemma-finite-type-Noetherian-quasi-separated}
（以及引理 \ref{morphisms-lemma-finite-presentation-finite-type}）。
\end{proof}

\begin{lemma}
\label{morphisms-lemma-finite-presentation-quasi-compact-quasi-separated}
令 $S$ 为拟紧拟分离概形。若 $X$ 在 $S$ 上有限表示，则 $X$
拟紧且拟分离。
\end{lemma}

\begin{proof}
略。
\end{proof}

\begin{lemma}
\label{morphisms-lemma-finite-presentation-permanence}
令 $f : X \to Y$ 为 $S$ 上的概形态射。
\begin{enumerate}
\item 若 $X$ 在 $S$ 上局部有限表示，并且 $Y$ 在 $S$ 上局部有限型，
则 $f$ 局部有限表示。
\item 若 $X$ 在 $S$ 上有限表示，并且 $Y$ 拟分离且在 $S$ 上局部有限型，
则 $f$ 有限表示。
\end{enumerate}
\end{lemma}

\begin{proof}
证明 (1)。通过引理 \ref{morphisms-lemma-locally-finite-presentation-characterize}，
这转化为如下代数事实：给定环同态 $A \to B \to C$，若
$A \to C$ 是有限表示的，且 $A \to B$ 是有限型的，则
$B \to C$ 是有限表示的。见《代数》引理
\ref{algebra-lemma-compose-finite-type}。

\medskip\noindent
(2) 来自 (1) 以及《概形》引理
\ref{schemes-lemma-compose-after-separated} 和
\ref{schemes-lemma-quasi-compact-permanence}。
\end{proof}

\begin{lemma}
\label{morphisms-lemma-diagonal-morphism-finite-type}
令 $f : X \to Y$ 为概形态射，其对角态射为
$\Delta : X \to X \times_Y X$。若 $f$ 局部有限型，则
$\Delta$ 局部有限表示。若 $f$ 拟分离且局部有限型，则
$\Delta$ 有限表示。
\end{lemma}

\begin{proof}
注意，$\Delta$ 是 $X$ 上的概形态射（经由第二投影
$X \times_Y X \to X$）。假设 $f$ 局部有限型。
注意，$X$ 在 $X$ 上有限表示，而 $X \times_Y X$ 在 $X$ 上局部有限型
（由引理 \ref{morphisms-lemma-base-change-finite-type}）。因此第一个断言来自引理
\ref{morphisms-lemma-finite-presentation-permanence}。
第二个断言来自第一个断言、各定义，以及对角态射是单态射因而分离这一事实
（《概形》引理 \ref{schemes-lemma-monomorphism-separated}）。
\end{proof}

\section{可构造集}
\label{morphisms-section-constructible}

\noindent
我们已在《性质》节 \ref{properties-section-constructible}
中讨论了概形的可构造集与局部可构造集。本节证明一些关于（局部）可构造集
的像与逆像的结果。主要结果是 Chevalley 定理：局部可构造集在有限表示态射
下的像仍然局部可构造。

\begin{lemma}
\label{morphisms-lemma-inverse-image-constructible}
令 $f : X \to Y$ 为概形态射，并令 $E \subset Y$ 为子集。
若 $E$ 在 $Y$ 中（局部）可构造，则 $f^{-1}(E)$ 在 $X$
中（局部）可构造。
\end{lemma}

\begin{proof}
为了证明每个可构造子集的逆像都是可构造的，只需证明每个逆拟紧开集
$V$（它是 $Y$ 的开集）的逆像在 $X$ 中逆拟紧；见《拓扑学》引理
\ref{topology-lemma-inverse-images-constructibles}。$V$ 在 $Y$
中逆拟紧的含义正是开浸入 $V \to Y$ 拟紧。因此基变换
$f^{-1}(V) = X \times_Y V \to X$ 也是拟紧的；见《概形》引理
\ref{schemes-lemma-quasi-compact-preserved-base-change}。
所以 $f^{-1}(V)$ 在 $X$ 中逆拟紧。
假设 $E$ 在 $Y$ 中局部可构造。选取 $x \in X$。选取仿射邻域
$V$（它是 $f(x)$ 的邻域），并选取仿射邻域 $U \subset X$
（它是 $x$ 的邻域），使得
$f(U) \subset V$。于是把 $f|_U : U \to V$ 看成以 $V$ 为靶的态射。
由《性质》引理 \ref{properties-lemma-locally-constructible}，
$E \cap V$ 在 $V$ 中可构造。由可构造情形，
$(f|_U)^{-1}(E \cap V)$ 在 $U$ 中可构造。由于
$(f|_U)^{-1}(E \cap V) = f^{-1}(E) \cap U$，结论成立。
\end{proof}

\begin{lemma}
\label{morphisms-lemma-chevalley}
令 $f : X \to Y$ 为概形态射。假设：
\begin{enumerate}
\item $f$ 拟紧且局部有限表示；以及
\item $Y$ 拟紧且拟分离。
\end{enumerate}
则 $X$ 的每个可构造子集的像在 $Y$ 中可构造。
\end{lemma}

\begin{proof}
由《性质》引理
\ref{properties-lemma-constructible-quasi-compact-quasi-separated}，
只需在 $Y$ 为仿射概形时证明本引理。在这种情形下，$X$ 拟紧。
因此可以写成 $X = U_1 \cup \ldots \cup U_n$，其中每个 $U_i$
都是 $X$ 的仿射开集。若 $E \subset X$ 可构造，则每个
$E \cap U_i$ 也可构造；见《拓扑学》引理
\ref{topology-lemma-open-immersion-constructible-inverse-image}。
由于 $f(E) = \bigcup f(E \cap U_i)$，且可构造集的有限并仍可构造，
问题归结为 $X$ 为仿射概形的情形。在这种情形下，结论就是《代数》定理
\ref{algebra-theorem-chevalley}。
\end{proof}

\begin{theorem}[Chevalley 定理]
\label{morphisms-theorem-chevalley}
\begin{reference}
\cite[IV, Theorem 1.8.4]{EGA}
\end{reference}
令 $f : X \to Y$ 为概形态射。假设 $f$ 拟紧且局部有限表示。
则每个局部可构造子集的像仍然局部可构造。
\end{theorem}

\begin{proof}
令 $E \subset X$ 局部可构造。必须证明 $f(E)$ 也局部可构造。
我们将证明 $f(E) \cap V$ 是可构造的，其中 $V \subset Y$
为任意仿射开集。由此可归结到 $Y$ 为仿射概形的情形。在这种情形下，
$X$ 拟紧。因此可以写成 $X = U_1 \cup \ldots \cup U_n$，
其中每个 $U_i$ 都是 $X$ 的仿射开集。若 $E \subset X$
局部可构造，则每个 $E \cap U_i$ 都可构造；见《性质》引理
\ref{properties-lemma-locally-constructible}。由于
$f(E) = \bigcup f(E \cap U_i)$，且可构造集的有限并仍可构造，
问题归结为 $X$ 为仿射概形的情形。在这种情形下，结论就是《代数》定理
\ref{algebra-theorem-chevalley}。
\end{proof}

\begin{lemma}
\label{morphisms-lemma-constructible-containing-open}
令 $X$ 为概形，令 $x \in X$，并令 $E \subset X$ 为局部可构造子集。
若 $\{x' \mid x' \leadsto x\} \subset E$，则 $E$ 包含 $x$
的一个开邻域。
\end{lemma}

\begin{proof}
假设 $\{x' \mid x' \leadsto x\} \subset E$。可以假设 $X$ 为仿射概形。
在这种情形下，由《性质》引理
\ref{properties-lemma-locally-constructible}，$E$ 可构造。
特别地，补集 $E^c$ 也可构造。由《代数》引理
\ref{algebra-lemma-constructible-is-image}，可以找到仿射概形态射
$f : Y \to X$，使得 $E^c = f(Y)$。令 $Z \subset X$ 为 $f$
的概形论像。由引理 \ref{morphisms-lemma-reach-points-scheme-theoretic-image}
和假设 $\{x' \mid x' \leadsto x\} \subset E$，可知
$x \not \in Z$。因此 $X \setminus Z \subset E$ 是 $x$
的开邻域，并且包含于 $E$。
\end{proof}





\section{开态射}
\label{morphisms-section-open}

\begin{definition}
\label{morphisms-definition-open}
令 $f : X \to S$ 为态射。
\begin{enumerate}
\item 称 $f$ 为{\it 开态射}，如果它在底层拓扑空间上诱导的映射为开映射。
\item 称 $f$ 为{\it 泛开态射}，如果对任意概形态射 $S' \to S$，
基变换 $f' : X_{S'} \to S'$ 都是开态射。
\end{enumerate}
\end{definition}

\noindent
根据《拓扑学》引理
\ref{topology-lemma-closed-open-map-specialization}，
泛化可沿某些类型的拓扑空间开映射提升。事实上，泛化可沿任意概形的开态射提升
（见引理 \ref{morphisms-lemma-open-generizing}）。我们还将看到，泛化可沿概形的平坦态射
提升（引理 \ref{morphisms-lemma-generalizations-lift-flat}）。这有时反过来推出该态射是开态射。

\begin{lemma}
\label{morphisms-lemma-locally-finite-presentation-universally-open}
令 $f : X \to S$ 为态射。
\begin{enumerate}
\item 若 $f$ 局部有限表示，且泛化可沿 $f$ 提升，则 $f$ 是开态射。
\item 若 $f$ 局部有限表示，且泛化可沿 $f$ 的每个基变换提升，
则 $f$ 是泛开态射。
\end{enumerate}
\end{lemma}

\begin{proof}
只需证明第一个断言。问题可归结到 $X$ 和 $S$ 都是仿射概形的情形。
在这种情形下，结论来自《代数》引理
\ref{algebra-lemma-going-up-down-specialization} 和命题
\ref{algebra-proposition-fppf-open}。
\end{proof}

\noindent
关于有限表示平坦态射的情形，另见引理 \ref{morphisms-lemma-fppf-open}。

\begin{lemma}
\label{morphisms-lemma-composition-open}
（泛）开态射的复合仍为（泛）开态射。
\end{lemma}

\begin{proof}
略。
\end{proof}

\begin{lemma}
\label{morphisms-lemma-scheme-over-field-universally-open}
令 $k$ 为域，令 $X$ 为 $k$ 上的概形。结构态射
$X \to \Spec(k)$ 是泛开态射。
\end{lemma}

\begin{proof}
令 $S \to \Spec(k)$ 为态射。必须证明基变换
$X_S \to S$ 是开态射。这个问题在 $S$ 和 $X$ 上都是局部的，
故可假设 $S$ 和 $X$ 都是仿射概形。在这种情形下，结论就是《代数》引理
\ref{algebra-lemma-map-into-tensor-algebra-open}。
\end{proof}

\begin{lemma}
\label{morphisms-lemma-open-generizing}
\begin{reference}
由 \cite[IV, Corollary 1.10.4]{EGA} 中 (a) $\Rightarrow$ (b) 得出。
\end{reference}
令 $\varphi : X \to Y$ 为概形态射。若 $\varphi$ 是开态射，
则 $\varphi$ 是保泛化的（即泛化可沿 $\varphi$ 提升）。
若 $\varphi$ 是泛开态射，则 $\varphi$ 是泛保泛化的。
\end{lemma}

\begin{proof}
假设 $\varphi$ 是开态射。令 $y' \leadsto y$ 为 $Y$ 中点的一次特化。
令 $x \in X$ 且 $\varphi(x) = y$。选取仿射开集
$U \subset X$ 和 $V \subset Y$，使得 $\varphi(U) \subset V$ 且
$x \in U$。于是 $y' \in V$。因此可将 $X$ 替换为 $U$，将 $Y$
替换为 $V$，并假设 $X$、$Y$ 都是仿射概形。仿射情形就是《代数》引理
\ref{algebra-lemma-open-going-down}（结合《代数》引理
\ref{algebra-lemma-going-up-down-specialization}）。
\end{proof}

\begin{lemma}
\label{morphisms-lemma-descent-quasi-compact}
令 $f : X \to Y$ 为概形态射。令 $g : Y' \to Y$ 为开满射，
并假设基变换 $f' : X' \to Y'$ 拟紧。则 $f$ 拟紧。
\end{lemma}

\begin{proof}
令 $V \subset Y$ 为拟紧开集。由于 $g$ 开且满，可以找到拟紧开集
$W' \subset W$，使得 $g(W') = V$。依假设，$(f')^{-1}(W')$
拟紧。$(f')^{-1}(W')$ 在 $X$ 中的像等于 $f^{-1}(V)$；见引理
\ref{morphisms-lemma-when-point-maps-to-pair}。因此 $f^{-1}(V)$ 作为拟紧空间的像
是拟紧的；见《拓扑学》引理 \ref{topology-lemma-image-quasi-compact}。
所以 $f$ 拟紧。
\end{proof}





\section{商态射}
\label{morphisms-section-submersive}

\begin{definition}
\label{morphisms-definition-submersive}
令 $f : X \to Y$ 为概形态射。
\begin{enumerate}
\item 称 $f$ 为{\it 商态射}\footnote{这与微分流形中的浸没概念大不相同。}，
如果底层拓扑空间之间的连续映射是商映射；见《拓扑学》定义
\ref{topology-definition-submersive}。
\item 称 $f$ 为{\it 泛商态射}，如果对每个概形态射 $Y' \to Y$，
基变换 $Y' \times_Y X \to Y'$ 都是商态射。
\end{enumerate}
\end{definition}

\noindent
注意，商态射特别是满射。

\begin{lemma}
\label{morphisms-lemma-base-change-universally-submersive}
概形的泛商态射沿任意概形态射的基变换仍为泛商态射。
\end{lemma}

\begin{proof}
这由定义立即得到。
\end{proof}

\begin{lemma}
\label{morphisms-lemma-composition-universally-submersive}
两个概形的（泛）商态射的复合仍为（泛）商态射。
\end{lemma}

\begin{proof}
略。
\end{proof}


\section{平坦态射}
\label{morphisms-section-flat}

\noindent
平坦性是代数几何中最重要的技术工具之一。本节引入这一概念。除引理
\ref{morphisms-lemma-fppf-open} 外，我们有意只讨论直接的观察。非常重要的一类结果，
即平坦性判别准则，见《代数》各节
\ref{algebra-section-criteria-flatness}、
\ref{algebra-section-flatness-artinian}、
\ref{algebra-section-more-flatness-criteria}，以及《态射进阶》节
\ref{more-morphisms-section-criterion-flat-fibres}。另有一章专门讨论概形平坦态射
的进阶材料，即《平坦性进阶》节 \ref{flat-section-introduction}。

\medskip\noindent
回忆：模 $M$ 在环 $R$ 上称为{\it 平坦的}，如果函子
$-\otimes_R M : \text{Mod}_R \to \text{Mod}_R$ 正合。环同态
$R \to A$ 称为{\it 平坦的}，如果 $A$ 作为 $R$-模平坦。见《代数》定义
\ref{algebra-definition-flat}。

\begin{definition}
\label{morphisms-definition-flat}
令 $f : X \to S$ 为概形态射。令 $\mathcal{F}$ 为拟凝聚
$\mathcal{O}_X$-模层。
\begin{enumerate}
\item 称 $f$ 在点 $x \in X$ 处{\it 平坦}，如果局部环
$\mathcal{O}_{X, x}$ 在局部环 $\mathcal{O}_{S, f(x)}$ 上平坦。
\item 称 $\mathcal{F}$ {\it 在 $S$ 上于点 $x \in X$ 处平坦}，
如果茎 $\mathcal{F}_x$ 是平坦的 $\mathcal{O}_{S, f(x)}$-模。
\item 称 $f$ {\it 平坦}，如果 $f$ 在 $X$ 的每一点处平坦。
\item 称 $\mathcal{F}$ {\it 在 $S$ 上平坦}，如果 $\mathcal{F}$
在 $S$ 上于每一点 $x$ 处平坦，而这些点取自 $X$。
\end{enumerate}
\end{definition}

\noindent
因此，$f$ 平坦当且仅当结构层 $\mathcal{O}_X$ 在 $S$ 上平坦。

\begin{lemma}
\label{morphisms-lemma-flat-module-characterize}
令 $f : X \to S$ 为概形态射。令 $\mathcal{F}$ 为拟凝聚
$\mathcal{O}_X$-模层。下列条件等价：
\begin{enumerate}
\item 层 $\mathcal{F}$ 在 $S$ 上平坦。
\item 对每对仿射开集 $U \subset X$、$V \subset S$，若
$f(U) \subset V$，则 $\mathcal{O}_S(V)$-模 $\mathcal{F}(U)$ 平坦。
\item 存在开覆盖 $S = \bigcup_{j \in J} V_j$ 以及开覆盖
$f^{-1}(V_j) = \bigcup_{i \in I_j} U_i$，使得每个模层
$\mathcal{F}|_{U_i}$ 在 $V_j$ 上平坦，其中所有
$j\in J, i\in I_j$ 均如此。
\item 存在仿射开覆盖 $S = \bigcup_{j \in J} V_j$ 以及仿射开覆盖
$f^{-1}(V_j) = \bigcup_{i \in I_j} U_i$，使得
$\mathcal{F}(U_i)$ 是平坦的 $\mathcal{O}_S(V_j)$-模，其中所有
$j\in J, i\in I_j$ 均如此。
\end{enumerate}
此外，若 $\mathcal{F}$ 在 $S$ 上平坦，则对任意开子概形
$U \subset X$、$V \subset S$，若 $f(U) \subset V$，限制
$\mathcal{F}|_U$ 在 $V$ 上平坦。
\end{lemma}

\begin{proof}
令 $R \to A$ 为环同态，令 $M$ 为 $A$-模。若 $M$ 在 $R$ 上平坦，
则对所有素理想 $\mathfrak q$，模 $M_{\mathfrak q}$ 在
$R_{\mathfrak p}$ 上平坦，其中 $\mathfrak p$ 是 $R$ 中位于
$\mathfrak q$ 之下的素理想。反之，若 $M_{\mathfrak q}$ 在
$R_{\mathfrak p}$ 上平坦，且这对每个素理想 $\mathfrak q$
（它属于 $A$）均成立，
则 $M$ 在 $R$ 上平坦。见《代数》引理
\ref{algebra-lemma-flat-localization}。这一等价性容易推出本引理的各个断言。
\end{proof}

\begin{lemma}
\label{morphisms-lemma-flat-characterize}
令 $f : X \to S$ 为概形态射。下列条件等价：
\begin{enumerate}
\item 态射 $f$ 平坦。
\item 对每对仿射开集 $U \subset X$、$V \subset S$，若
$f(U) \subset V$，则环同态
$\mathcal{O}_S(V) \to \mathcal{O}_X(U)$ 平坦。
\item 存在开覆盖 $S = \bigcup_{j \in J} V_j$ 以及开覆盖
$f^{-1}(V_j) = \bigcup_{i \in I_j} U_i$，使得每个态射
$U_i \to V_j$ 都平坦，其中 $j\in J, i\in I_j$。
\item 存在仿射开覆盖 $S = \bigcup_{j \in J} V_j$ 以及仿射开覆盖
$f^{-1}(V_j) = \bigcup_{i \in I_j} U_i$，使得
$\mathcal{O}_S(V_j) \to \mathcal{O}_X(U_i)$ 平坦，其中所有
$j\in J, i\in I_j$ 均如此。
\end{enumerate}
此外，若 $f$ 平坦，则对任意开子概形 $U \subset X$、$V \subset S$，
若 $f(U) \subset V$，限制 $f|_U : U \to V$ 平坦。
\end{lemma}

\begin{proof}
这是上面引理 \ref{morphisms-lemma-flat-module-characterize} 的特殊情形。
\end{proof}

\begin{lemma}
\label{morphisms-lemma-pushforward-flat-affine}
令 $f : X \to Y$ 为基概形 $S$ 上的仿射概形态射。令
$\mathcal{F}$ 为拟凝聚 $\mathcal{O}_X$-模。则 $\mathcal{F}$ 在
$S$ 上平坦，当且仅当 $f_*\mathcal{F}$ 在 $S$ 上平坦。
\end{lemma}

\begin{proof}
由引理 \ref{morphisms-lemma-flat-module-characterize} 以及 $f$ 是仿射态射这一事实，
问题归结到仿射情形。设 $X \to Y \to S$ 对应于环同态
$C \leftarrow B \leftarrow A$。令 $N$ 为一个 $C$-模，它对应于
$\mathcal{F}$。回忆，$f_*\mathcal{F}$ 对应于把 $N$ 看成 $B$-模；见《概形》引理
\ref{schemes-lemma-widetilde-pullback}。因此结论显然。
\end{proof}

\begin{lemma}
\label{morphisms-lemma-composition-module-flat}
令 $X \to Y \to Z$ 为概形态射。令 $\mathcal{F}$ 为拟凝聚
$\mathcal{O}_X$-模。令 $x \in X$，其像为 $y$，位于 $Y$ 中。
若 $\mathcal{F}$ 在 $Y$ 上于点 $x$ 处平坦，并且 $Y$ 在 $Z$
上于点 $y$ 处平坦，则 $\mathcal{F}$ 在 $Z$ 上于点 $x$ 处平坦。
\end{lemma}

\begin{proof}
见《代数》引理 \ref{algebra-lemma-composition-flat}。
\end{proof}

\begin{lemma}
\label{morphisms-lemma-composition-flat}
平坦态射的复合仍平坦。
\end{lemma}

\begin{proof}
这是引理 \ref{morphisms-lemma-composition-module-flat} 的特殊情形。
\end{proof}

\begin{lemma}
\label{morphisms-lemma-base-change-module-flat}
令 $f : X \to S$ 为概形态射。令 $\mathcal{F}$ 为拟凝聚
$\mathcal{O}_X$-模层。令 $g : S' \to S$ 为概形态射。
记投影为 $g' : X' = X_{S'} \to X$。令 $x' \in X'$ 为一点，
其像为 $x = g'(x') \in X$。若 $\mathcal{F}$ 在 $S$ 上于点
$x$ 处平坦，则 $(g')^*\mathcal{F}$ 在 $S'$ 上于点 $x'$ 处平坦。
特别地，若 $\mathcal{F}$ 在 $S$ 上平坦，则
$(g')^*\mathcal{F}$ 在 $S'$ 上平坦。
\end{lemma}

\begin{proof}
见《代数》引理 \ref{algebra-lemma-flat-base-change}。
\end{proof}

\begin{lemma}
\label{morphisms-lemma-base-change-flat}
平坦态射的基变换仍平坦。
\end{lemma}

\begin{proof}
这是引理 \ref{morphisms-lemma-base-change-module-flat} 的特殊情形。
\end{proof}

\begin{lemma}
\label{morphisms-lemma-generalizations-lift-flat}
令 $f : X \to S$ 为概形的平坦态射。则泛化可沿 $f$ 提升；见《拓扑学》定义
\ref{topology-definition-lift-specializations}。
\end{lemma}

\begin{proof}
见《代数》节 \ref{algebra-section-going-up}。
\end{proof}

\begin{lemma}
\label{morphisms-lemma-fppf-open}
局部有限表示的平坦态射是泛开态射。
\end{lemma}

\begin{proof}
这来自引理 \ref{morphisms-lemma-generalizations-lift-flat} 以及上面的引理
\ref{morphisms-lemma-locally-finite-presentation-universally-open}。
也可以如下直接论证。

\medskip\noindent
令 $f : X \to S$ 平坦且局部有限表示。由引理
\ref{morphisms-lemma-base-change-flat} 和
\ref{morphisms-lemma-base-change-finite-presentation}，$f$ 的任意基变换都平坦且
局部有限表示。因此只需证明 $f$ 是开态射。为证明 $f$ 是开态射，
只需证明可以用仿射开集覆盖 $X$，即 $X = \bigcup U_i$，使得
$U_i \to S$ 是开态射。可以覆盖 $X$，所用仿射开集为 $U_i \subset X$，
使得每个 $U_i$ 都映入仿射开集 $V_i \subset S$，并且所诱导的环同态
$\mathcal{O}_S(V_i) \to \mathcal{O}_X(U_i)$ 平坦且有限表示（引理
\ref{morphisms-lemma-flat-characterize} 和
\ref{morphisms-lemma-locally-finite-presentation-characterize}）。于是由《代数》命题
\ref{algebra-proposition-fppf-open}，$U_i \to V_i$ 是开态射，证明完成。
\end{proof}

\begin{lemma}
\label{morphisms-lemma-pf-flat-module-open}
令 $f : X \to Y$ 为概形态射。令 $\mathcal{F}$ 为拟凝聚
$\mathcal{O}_X$-模。假设 $f$ 局部有限表示，$\mathcal{F}$ 有限型，
$X = \text{Supp}(\mathcal{F})$，并且 $\mathcal{F}$ 在 $Y$ 上平坦。
则 $f$ 是泛开态射。
\end{lemma}

\begin{proof}
由引理 \ref{morphisms-lemma-base-change-module-flat}、
\ref{morphisms-lemma-base-change-finite-presentation} 和
\ref{morphisms-lemma-support-finite-type}，这些假设在基变换下保持。
由引理 \ref{morphisms-lemma-locally-finite-presentation-universally-open}，
只需证明泛化可沿 $f$ 提升。这来自《代数》引理
\ref{algebra-lemma-going-down-flat-module}。
\end{proof}

\begin{lemma}
\label{morphisms-lemma-fpqc-quotient-topology}
\begin{reference}
\cite[Expose VIII, Corollaire 4.3]{SGA1} 和
\cite[IV, Corollaire 2.3.12]{EGA}
\end{reference}
令 $f : X \to Y$ 为拟紧、满且平坦的态射。子集 $T \subset Y$ 为开集
（相应地，闭集），当且仅当 $f^{-1}(T)$ 为开集（相应地，闭集）。换言之，
$f$ 是商态射。
\end{lemma}

\begin{proof}
这个问题在 $Y$ 上是局部的，故可假设 $Y$ 为仿射概形。在这种情形下，
$X$ 也拟紧，因为 $f$ 拟紧。把 $X = X_1 \cup \ldots \cup X_n$
写成仿射开集的有限并。于是
$f' : X' = X_1 \amalg \ldots \amalg X_n \to Y$ 是仿射概形之间的
满平坦态射。注意，对 $T \subset Y$，有
$(f')^{-1}(T) = f^{-1}(T) \cap X_1 \amalg \ldots \amalg f^{-1}(T) \cap X_n$。
因此 $f^{-1}(T)$ 为开集，当且仅当 $(f')^{-1}(T)$ 为开集。
所以可假设 $X$ 和 $Y$ 都是仿射概形。

\medskip\noindent
令 $f : \Spec(B) \to \Spec(A)$ 为仿射概形的满态射，对应于平坦环同态
$A \to B$。假设 $f^{-1}(T)$ 为闭集，比如
$f^{-1}(T) = V(J)$，其中 $J \subset B$ 为理想。于是
$T = f(f^{-1}(T)) = f(V(J))$ 是 $\Spec(B/J) \to \Spec(A)$ 的像
（这里用到了 $f$ 的满性）。另一方面，泛化可沿 $f$ 提升（引理
\ref{morphisms-lemma-generalizations-lift-flat}）。因此，由《拓扑学》引理
\ref{topology-lemma-lift-specializations-images}，
$Y \setminus T = f(X \setminus f^{-1}(T))$ 在泛化下稳定。所以 $T$
在特化下稳定（《拓扑学》引理
\ref{topology-lemma-open-closed-specialization}）。由《代数》引理
\ref{algebra-lemma-image-stable-specialization-closed}，$T$ 是闭集。
\end{proof}

\begin{lemma}
\label{morphisms-lemma-flat-permanence}
令 $h : X \to Y$ 为 $S$ 上的概形态射。令 $\mathcal{G}$ 为 $Y$
上的拟凝聚层。令 $x \in X$，且 $y = h(x) \in Y$。若 $h$ 在 $x$
处平坦，则
$$
\mathcal{G}\text{ flat over }S\text{ 在 }y
\Leftrightarrow
h^*\mathcal{G}\text{ flat over }S\text{ 在 }x.
$$
特别地：若 $h$ 满且平坦，则 $\mathcal{G}$ 在 $S$ 上平坦，当且仅当
$h^*\mathcal{G}$ 在 $S$ 上平坦。若 $h$ 满且平坦，并且 $X$ 在 $S$
上平坦，则 $Y$ 在 $S$ 上平坦。
\end{lemma}

\begin{proof}
可以应用《代数》引理
\ref{algebra-lemma-flatness-descends-more-general} 来证明。下面给出直接证明。
令 $s \in S$ 为 $y$ 的像。考虑局部环同态
$\mathcal{O}_{S, s} \to \mathcal{O}_{Y, y} \to \mathcal{O}_{X, x}$。
由假设，环同态 $\mathcal{O}_{Y, y} \to \mathcal{O}_{X, x}$ 忠实平坦；
见《代数》引理 \ref{algebra-lemma-local-flat-ff}。令
$N = \mathcal{G}_y$。注意，
$h^*\mathcal{G}_x = N \otimes_{\mathcal{O}_{Y, y}} \mathcal{O}_{X, x}$；
见《层》引理 \ref{sheaves-lemma-stalk-pullback-modules}。令
$M' \to M$ 为 $\mathcal{O}_{S, s}$-模的单射。由上述忠实平坦性，有
\begin{align*}
\Ker(
M' \otimes_{\mathcal{O}_{S, s}} N \to M \otimes_{\mathcal{O}_{S, s}} N)
\otimes_{\mathcal{O}_{Y, y}} \mathcal{O}_{X, x} \\
=
\Ker(
M' \otimes_{\mathcal{O}_{S, s}} N
\otimes_{\mathcal{O}_{Y, y}} \mathcal{O}_{X, x}
\to
M \otimes_{\mathcal{O}_{S, s}} N
\otimes_{\mathcal{O}_{Y, y}} \mathcal{O}_{X, x})
\end{align*}
所以本引理的等价性来自《代数》引理 \ref{algebra-lemma-flat}
中平坦性的第二种刻画。
\end{proof}

\begin{lemma}
\label{morphisms-lemma-flat-pullback-support}
令 $f : Y \to X$ 为概形态射。令 $\mathcal{F}$ 为有限型拟凝聚
$\mathcal{O}_X$-模，其概形论支集为 $Z \subset X$。若 $f$ 平坦，
则 $f^{-1}(Z)$ 是 $f^*\mathcal{F}$ 的概形论支集。
\end{lemma}

\begin{proof}
使用引理 \ref{morphisms-lemma-scheme-theoretic-support} 给出的仿射概形上概形论支集的刻画，
问题归结到《代数》引理
\ref{algebra-lemma-annihilator-flat-base-change}。
\end{proof}

\begin{lemma}
\label{morphisms-lemma-flat-morphism-scheme-theoretically-dense-open}
令 $f : X \to Y$ 为概形的平坦态射。令 $V \subset Y$ 为逆拟紧开集，
且概形论稠密。则 $f^{-1}V$ 在 $X$ 中概形论稠密。
\end{lemma}

\begin{proof}
我们将使用引理 \ref{morphisms-lemma-characterize-scheme-theoretically-dense} 的刻画。
必须证明，对任意开集 $U \subset X$，映射
$\mathcal{O}_X(U) \to \mathcal{O}_X(U \cap f^{-1}V)$ 都是单射。
只需在 $U$ 是映入某个仿射开集 $W \subset Y$ 的仿射开集时证明这一点。
设 $W = \Spec(A)$ 且 $U = \Spec(B)$。由《代数》引理
\ref{algebra-lemma-qc-open}，有
$V \cap W = D(f_1) \cup \ldots \cup D(f_n)$，其中某些
$f_i \in A$。
因此必须证明
$B \to B_{f_1} \times \ldots \times B_{f_n}$ 是单射。已知
$A \to A_{f_1} \times \ldots \times A_{f_n}$ 是单射，并且
$A \to B$ 平坦。由于 $B_{f_i} = A_{f_i} \otimes_A B$，结论成立。
\end{proof}

\begin{lemma}
\label{morphisms-lemma-flat-base-change-scheme-theoretic-image}
\begin{slogan}
在拟紧情形下，取概形论像与平坦基变换可交换
\end{slogan}
令 $f : X \to Y$ 为概形的平坦态射。令 $g : V \to Y$ 为概形的拟紧态射。
令 $Z \subset Y$ 为 $g$ 的概形论像，并令 $Z' \subset X$ 为基变换
$V \times_Y X \to X$ 的概形论像。则 $Z' = f^{-1}Z$。
\end{lemma}

\begin{proof}
回忆，$Z$ 由下式所截：
$\mathcal{I} = \Ker(\mathcal{O}_Y \to g_*\mathcal{O}_V)$，
而 $Z'$ 由下式所截：
$\mathcal{I}' = \Ker(\mathcal{O}_X \to
(V \times_Y X \to X)_*\mathcal{O}_{V \times_Y X})$；见引理
\ref{morphisms-lemma-quasi-compact-scheme-theoretic-image}。因此这个问题在 $X$ 和
$Y$ 上是局部的，可以假设 $X$ 和 $Y$ 都是仿射概形。注意，可以将
$V$ 替换为 $\coprod V_i$，其中 $V = V_1 \cup \ldots \cup V_n$
是有限仿射开覆盖。因此可假设 $g$ 为仿射态射。在这种情形下，
$(V \times_Y X \to X)_*\mathcal{O}_{V \times_Y X}$ 是
$g_*\mathcal{O}_V$ 沿 $f$ 的拉回。由于 $f$ 平坦，可得
$f^*\mathcal{I} = \mathcal{I}'$，本引理成立。
\end{proof}




\section{平坦闭浸入}
\label{morphisms-section-flat-closed-immersions}

\noindent
概形的连通分支不总是开集，但它们总有典范概形结构。本节说明这一点。

\begin{lemma}
\label{morphisms-lemma-characterize-flat-closed-immersions}
令 $X$ 为概形。把 $X$ 的每个闭子概形对应到底层闭子集的规则定义了双射
$$
\left\{
\begin{matrix}
\text{closed subschemes }Z \subset X \\
\text{使得 }Z \to X\text{ 平坦}
\end{matrix}
\right\}
\leftrightarrow
\left\{
\begin{matrix}
\text{closed subsets }Z \subset X \\
\text{closed under generalizations}
\end{matrix}
\right\}
$$
若 $Z \subset X$ 是这样的闭子概形，则每个概形态射
$g : Y \to X$，若在集合论意义下 $g(Y) \subset Z$，都在概形论意义下
唯一地通过 $Z$ 分解。
\end{lemma}

\begin{proof}
双射的仿射情形就是《代数》引理
\ref{algebra-lemma-pure-open-closed-specializations}。对一般概形 $X$，
用仿射开集覆盖 $X$ 并黏合即可得到双射。细节略去。关于最后一个断言，
注意投影 $Z \times_{X, g} Y \to Y$ 是平坦的（引理
\ref{morphisms-lemma-base-change-flat}）闭浸入，并且在底层拓扑空间上是双射，
所以由证明第一部分所建立的双射，它必为同构。
\end{proof}

\begin{lemma}
\label{morphisms-lemma-flat-closed-immersions-finite-presentation}
有限表示的平坦闭浸入是某个开闭子概形的开浸入。
\end{lemma}

\begin{proof}
仿射情形就是《代数》引理
\ref{algebra-lemma-finitely-generated-pure-ideal}。一般情形下，
用仿射开集覆盖 $X$ 即得本引理。细节略去。
\end{proof}

\noindent
注意，连通分支 $T$（属于概形 $X$）是在泛化下稳定的闭子集。
因此下述定义是有意义的。

\begin{definition}
\label{morphisms-definition-scheme-structure-connected-component}
令 $X$ 为概形。令 $T \subset X$ 为连通分支。$T$ 上的
{\it 典范概形结构}是 $T$ 上唯一的概形结构，使得闭浸入
$T \to X$ 平坦；见引理
\ref{morphisms-lemma-characterize-flat-closed-immersions}。
\end{definition}

\noindent
利用前一引理，可以判定每个有限平坦 $\mathcal{O}_X$-模何时有限局部自由。

\begin{lemma}
\label{morphisms-lemma-finite-flat-is-finite-locally-free}
令 $X$ 为概形。下列条件等价：
\begin{enumerate}
\item 每个有限平坦拟凝聚 $\mathcal{O}_X$-模都是有限局部自由的；以及
\item 每个在泛化下闭合的闭子集 $Z \subset X$ 都是开集。
\end{enumerate}
\end{lemma}

\begin{proof}
在仿射情形下，这就是《代数》引理
\ref{algebra-lemma-finite-flat-module-finitely-presented}。
概形情形不能直接从仿射情形推出，所以我们直接重复论证。

\medskip\noindent
假设 (1)。考虑闭浸入 $i : Z \to X$，并假设 $i$ 平坦。于是
$i_*\mathcal{O}_Z$ 拟凝聚且平坦，故由 (1) 有限局部自由。因此
$Z = \text{Supp}(i_*\mathcal{O}_Z)$ 也是开集，从而 (2) 成立。
所以由引理 \ref{morphisms-lemma-characterize-flat-closed-immersions}
中对平坦闭浸入的刻画，可得 (1) $\Rightarrow$ (2)。

\medskip\noindent
反之，假设 $X$ 满足 (2)。令 $\mathcal{F}$ 为有限平坦拟凝聚
$\mathcal{O}_X$-模。令 $Z = \text{Supp}(\mathcal{F})$；这个集合是
$\mathcal{F}$ 的支集，并且为闭集；见《模》引理
\ref{modules-lemma-support-finite-type-closed}。另一方面，若
$x \leadsto x'$ 是一次特化，则由《代数》引理
\ref{algebra-lemma-finite-flat-local}，模 $\mathcal{F}_{x'}$ 在
$\mathcal{O}_{X, x'}$ 上自由，并且
$$
\mathcal{F}_x =
\mathcal{F}_{x'} \otimes_{\mathcal{O}_{X, x'}} \mathcal{O}_{X, x}.
$$
因此
$x' \in \text{Supp}(\mathcal{F}) \Rightarrow x \in \text{Supp}(\mathcal{F})$，
换言之，该支集在泛化下闭合。由于 $X$ 满足 (2)，可知
$\mathcal{F}$ 的支集是开闭集。模
$\wedge^i(\mathcal{F})$，$i = 1, 2, 3, \ldots$，也都是有限平坦拟凝聚
$\mathcal{O}_X$-模；见《模》节
\ref{modules-section-symmetric-exterior}。注意
$\text{Supp}(\wedge^{i + 1}(\mathcal{F})) \subset
\text{Supp}(\wedge^i(\mathcal{F}))$。
因此存在开闭子集的分解
$$
X = U_0 \amalg U_1 \amalg U_2 \amalg \ldots
$$
使得 $\wedge^i(\mathcal{F})$ 的支集都是
$U_i \cup U_{i + 1} \cup \ldots$，对所有 $i$ 均如此。令 $x$ 为 $X$ 的一点，设
$x \in U_r$。注意
$\wedge^i(\mathcal{F})_x \otimes \kappa(x) =
\wedge^i(\mathcal{F}_x \otimes \kappa(x))$。
因此，$x \in U_r$ 推出 $\mathcal{F}_x \otimes \kappa(x)$
是维数为 $r$ 的向量空间。由 Nakayama 引理（见《代数》引理
\ref{algebra-lemma-NAK}），可以选取仿射开邻域
$U \subset U_r \subset X$（它是 $x$ 的邻域），以及截面
$s_1, \ldots, s_r \in \mathcal{F}(U)$，使得所诱导的映射
$$
\mathcal{O}_U^{\oplus r} \longrightarrow \mathcal{F}|_U, \quad
(f_1, \ldots, f_r) \longmapsto \sum f_i s_i
$$
满射。这意味着 $\wedge^r(\mathcal{F}|_U)$ 是有限平坦拟凝聚
$\mathcal{O}_U$-模，其支集是整个 $U$。由上面的论证，它由单个元素
$s_1 \wedge \ldots \wedge s_r$ 生成。因此
$\wedge^r(\mathcal{F}|_U) \cong \mathcal{O}_U/\mathcal{I}$，
其中 $\mathcal{I}$ 是某个拟凝聚理想层，使得
$\mathcal{O}_U/\mathcal{I}$ 在 $\mathcal{O}_U$ 上平坦，并且
$V(\mathcal{I}) = U$。应用引理
\ref{morphisms-lemma-characterize-flat-closed-immersions}，可得
$\mathcal{I} = 0$。所以 $s_1 \wedge \ldots \wedge s_r$ 是
$\wedge^r(\mathcal{F}|_U)$ 的一组基，进而上述映射既单又满。
这证明了 $\mathcal{F}$ 如所要求的是有限局部自由的。
\end{proof}






\section{一般平坦性}
\label{morphisms-section-generic-flatness}

\noindent
整基底上的有限型概形在基底的某个稠密开集上平坦。在《代数》节
\ref{algebra-section-generic-flatness} 中，我们证明了 Noether 情形、
有限表示态射的情形以及一般情形。本节只陈述并证明一般情形。不过，
命题 \ref{morphisms-proposition-generic-flatness-reduced} 将取代这一结果；
它表明只假设基底约化，结论仍然成立。

\begin{proposition}[一般平坦性]
\label{morphisms-proposition-generic-flatness}
令 $f : X \to S$ 为概形态射。令 $\mathcal{F}$ 为拟凝聚
$\mathcal{O}_X$-模层。假设：
\begin{enumerate}
\item $S$ 是整概形；
\item $f$ 有限型；并且
\item $\mathcal{F}$ 是有限型 $\mathcal{O}_X$-模。
\end{enumerate}
则存在稠密开子概形 $U \subset S$，使得 $X_U \to U$ 平坦且有限表示，
并且 $\mathcal{F}|_{X_U}$ 在 $U$ 上平坦且在
$\mathcal{O}_{X_U}$ 上有限表示。
\end{proposition}

\begin{proof}
由于 $S$ 是整概形，它不可约（见《性质》引理
\ref{properties-lemma-characterize-integral}），且任意非空开集都稠密。
因此可以用 $S$ 的一个仿射开集替换 $S$，并假设
$S = \Spec(A)$ 为仿射概形。由于 $S$ 是整概形，可知 $A$ 是整环。
由于 $f$ 有限型，它拟紧，所以 $X$ 拟紧。因此可以找到有限仿射开覆盖
$X = \bigcup_{i = 1, \ldots, n} X_i$。写成 $X_i = \Spec(B_i)$。
由引理 \ref{morphisms-lemma-locally-finite-type-characterize}，$B_i$ 是有限型
$A$-代数。此外，存在有限型 $B_i$-模 $M_i$，使得
$\mathcal{F}|_{X_i}$ 是对应于 $B_i$-模 $M_i$ 的拟凝聚层；见《性质》引理
\ref{properties-lemma-finite-type-module}。接着，对每对指标 $i, j$，
选取理想 $I_{ij} \subset B_i$，使得
$X_i \setminus X_i \cap X_j = V(I_{ij})$ 在
$X_i = \Spec(B_i)$ 中成立。令
$M_{ij} = B_i/I_{ij}$，并把它看成 $B_i$-模。于是
$V(I_{ij}) = \text{Supp}(M_{ij})$，且 $M_{ij}$ 是有限 $B_i$-模。

\medskip\noindent
现在对各对 $(A \to B_i, M_{ij})$ 以及各对
$(A \to B_i, M_i)$ 应用《代数》引理
\ref{algebra-lemma-generic-flatness}。因此得到非零元
$f_{ij}, f_i \in A$，使得：(a)
$A_{f_{ij}} \to B_{i, f_{ij}}$ 平坦且有限表示，
$M_{ij, f_{ij}}$ 在 $A_{f_{ij}}$ 上平坦且在 $B_{i, f_{ij}}$
上有限表示；并且 (b) $B_{i, f_i}$ 在 $A_f$ 上平坦且有限表示，
$M_{i, f_i}$ 在 $B_{i, f_i}$ 上平坦且有限表示。令
$f = (\prod f_i) (\prod f_{ij})$。我们断言取 $U = D(f)$ 即可。

\medskip\noindent
为证明这一断言，可以将 $A$ 替换为 $A_f$，也就是施行基变换
$U = \Spec(A_f) \to S$。基变换后，每个 $A \to B_i$ 都平坦且
有限表示，并且 $M_i$、$M_{ij}$ 在 $A$ 上平坦且在 $B_i$ 上
有限表示。这已经证明 $X \to S$ 拟紧、局部有限表示且平坦，并证明
$\mathcal{F}$ 在 $S$ 上平坦且在 $\mathcal{O}_X$ 上有限表示；见引理
\ref{morphisms-lemma-locally-finite-presentation-characterize} 和《性质》引理
\ref{properties-lemma-finite-presentation-module}。由于 $M_{ij}$ 在
$B_i$ 上有限表示，可知
$X_i \cap X_j = X_i \setminus \text{Supp}(M_{ij})$ 是 $X_i$
的拟紧开集；见《代数》引理
\ref{algebra-lemma-support-finite-presentation-constructible}。
所以由《概形》引理 \ref{schemes-lemma-characterize-quasi-separated}，
$X \to S$ 拟分离。这证明了本命题。
\end{proof}

\noindent
事实上，在任意约化基底上也有一般平坦性版本。陈述如下。

\begin{proposition}[一般平坦性，约化情形]
\label{morphisms-proposition-generic-flatness-reduced}
令 $f : X \to S$ 为概形态射。令 $\mathcal{F}$ 为拟凝聚
$\mathcal{O}_X$-模层。假设：
\begin{enumerate}
\item $S$ 约化；
\item $f$ 有限型；并且
\item $\mathcal{F}$ 是有限型 $\mathcal{O}_X$-模。
\end{enumerate}
则存在稠密开子概形 $U \subset S$，使得 $X_U \to U$ 平坦且有限表示，
并且 $\mathcal{F}|_{X_U}$ 在 $U$ 上平坦且在
$\mathcal{O}_{X_U}$ 上有限表示。
\end{proposition}

\begin{proof}
对急于得知结果的读者：本证明重复命题
\ref{morphisms-proposition-generic-flatness} 的证明，只是用《代数》引理
\ref{algebra-lemma-generic-flatness-reduced} 代替《代数》引理
\ref{algebra-lemma-generic-flatness}。

\medskip\noindent
由于平坦和有限表示都在基底上是局部性质（见引理
\ref{morphisms-lemma-flat-module-characterize} 和
\ref{morphisms-lemma-locally-finite-presentation-characterize}），可以在 $S$
上仿射局部地工作。因此可假设 $S = \Spec(A)$，其中 $A$ 为约化环
（见《性质》引理 \ref{properties-lemma-characterize-reduced}）。
由于 $f$ 有限型，它拟紧，所以 $X$ 拟紧。因此可以找到有限仿射开覆盖
$X = \bigcup_{i = 1, \ldots, n} X_i$。写成 $X_i = \Spec(B_i)$。
由引理 \ref{morphisms-lemma-locally-finite-type-characterize}，$B_i$ 是有限型
$A$-代数。此外，存在有限型 $B_i$-模 $M_i$，使得
$\mathcal{F}|_{X_i}$ 是对应于 $B_i$-模 $M_i$ 的拟凝聚层；见《性质》引理
\ref{properties-lemma-finite-type-module}。接着，对每对指标 $i, j$，
选取理想 $I_{ij} \subset B_i$，使得
$X_i \setminus X_i \cap X_j = V(I_{ij})$ 在
$X_i = \Spec(B_i)$ 中成立。令
$M_{ij} = B_i/I_{ij}$，并把它看成 $B_i$-模。于是
$V(I_{ij}) = \text{Supp}(M_{ij})$，且 $M_{ij}$ 是有限 $B_i$-模。

\medskip\noindent
现在对各对 $(A \to B_i, M_{ij})$ 以及各对
$(A \to B_i, M_i)$ 应用《代数》引理
\ref{algebra-lemma-generic-flatness-reduced}。因此得到稠密开集
$U(A \to B_i, M_{ij}) \subset S$ 以及稠密开集
$U(A \to B_i, M_i) \subset S$，记号如《代数》方程
(\ref{algebra-equation-good-locus})。由于稠密开集的有限交仍为稠密开集，
可知
$$
U =
\bigcap\nolimits_{i, j} U(A \to B_i, M_{ij})
\quad\cap\quad
\bigcap\nolimits_i U(A \to B_i, M_i)
$$
在 $S$ 中开且稠密。我们断言 $U$ 就是所需开集。

\medskip\noindent
取 $u \in U$。由轨迹 $U(A \to B_i, M_{ij})$ 和
$U(A \to B, M_i)$ 的定义，存在 $f_{ij}, f_i \in A$，使得：
(a) $u \in D(f_i)$ 且 $u \in D(f_{ij})$；
(b) $A_{f_{ij}} \to B_{i, f_{ij}}$ 平坦且有限表示，
$M_{ij, f_{ij}}$ 在 $A_{f_{ij}}$ 上平坦且在 $B_{i, f_{ij}}$
上有限表示；并且
(c) $B_{i, f_i}$ 在 $A_{f_i}$ 上平坦且有限表示，
$M_{i, f_i}$ 在 $B_{i, f_i}$ 上有限表示且在 $A_{f_i}$ 上平坦。
令 $f = (\prod f_i) (\prod f_{ij})$。现在只需证明
$X \to S$ 在 $D(f)$ 上平坦且有限表示，并且 $\mathcal{F}$ 限制到
$X_{D(f)}$ 后在 $D(f)$ 上平坦且在 $X_{D(f)}$ 的结构层上有限表示。

\medskip\noindent
因此可以将 $A$ 替换为 $A_f$，也就是施行基变换
$\Spec(A_f) \to S$。基变换后，每个 $A \to B_i$ 都平坦且有限表示，
并且 $M_i$、$M_{ij}$ 在 $A$ 上平坦且在 $B_i$ 上有限表示。
这已经证明 $X \to S$ 拟紧、局部有限表示且平坦，并证明
$\mathcal{F}$ 在 $S$ 上平坦且在 $\mathcal{O}_X$ 上有限表示；见引理
\ref{morphisms-lemma-locally-finite-presentation-characterize} 和《性质》引理
\ref{properties-lemma-finite-presentation-module}。由于 $M_{ij}$ 在
$B_i$ 上有限表示，可知
$X_i \cap X_j = X_i \setminus \text{Supp}(M_{ij})$ 是 $X_i$
的拟紧开集；见《代数》引理
\ref{algebra-lemma-support-finite-presentation-constructible}。
所以由《概形》引理 \ref{schemes-lemma-characterize-quasi-separated}，
$X \to S$ 拟分离。这证明了本命题。
\end{proof}









\begin{remark}
\label{morphisms-remark-flattening}
上述结果是概形态射的更精细平坦化技术的第一步。Raynaud 和 Gruson
的文章 \cite{GruRay} 包含许多这方面的精彩结果。
\end{remark}







\section{态射与纤维的维数}
\label{morphisms-section-dimension-fibres}

\noindent
令 $X$ 为拓扑空间，令 $x \in X$。回忆，$\dim_x(X)$ 已定义为
$x$ 的开邻域（在 $X$ 中）的维数的最小值。见《拓扑学》定义
\ref{topology-definition-Krull}。

\begin{lemma}
\label{morphisms-lemma-dimension-fibre-at-a-point}
令 $f : X \to S$ 为概形态射。令 $x \in X$，并置 $s = f(x)$。
假设 $f$ 局部有限型。则
$$
\dim_x(X_s) =
\dim(\mathcal{O}_{X_s, x}) + \text{trdeg}_{\kappa(s)}(\kappa(x)).
$$
\end{lemma}

\begin{proof}
问题立即归结到 $S = s$ 且 $X$ 为仿射概形的情形。在这种情形下，
结论来自《代数》引理
\ref{algebra-lemma-dimension-at-a-point-finite-type-field}。
\end{proof}

\begin{lemma}
\label{morphisms-lemma-dimension-fibre-at-a-point-additive}
令 $f : X \to Y$ 和 $g : Y \to S$ 为概形态射。令 $x \in X$，
并置 $y = f(x)$、$s = g(y)$。假设 $f$ 和 $g$ 局部有限型。则
$$
\dim_x(X_s) \leq \dim_x(X_y) + \dim_y(Y_s).
$$
此外，若 $\mathcal{O}_{X_s, x}$ 在 $\mathcal{O}_{Y_s, y}$ 上平坦，
则等号成立；例如，若 $\mathcal{O}_{X, x}$ 在 $\mathcal{O}_{Y, y}$
上平坦，就满足这一条件。
\end{lemma}

\begin{proof}
注意
$\text{trdeg}_{\kappa(s)}(\kappa(x)) =
\text{trdeg}_{\kappa(y)}(\kappa(x)) + \text{trdeg}_{\kappa(s)}(\kappa(y))$。
因此由引理 \ref{morphisms-lemma-dimension-fibre-at-a-point}，该断言等价于
$$
\dim(\mathcal{O}_{X_s, x})
\leq
\dim(\mathcal{O}_{X_y, x}) + \dim(\mathcal{O}_{Y_s, y}).
$$
对此见《代数》引理
\ref{algebra-lemma-dimension-base-fibre-total}。平坦情形见《代数》引理
\ref{algebra-lemma-dimension-base-fibre-equals-total}。
\end{proof}

\begin{lemma}
\label{morphisms-lemma-dimension-fibre-after-base-change}
令
$$
\xymatrix{
X' \ar[r]_{g'} \ar[d]_{f'} & X \ar[d]^f \\
S' \ar[r]^g & S
}
$$
为概形的纤维积图。假设 $f$ 局部有限型。设 $x' \in X'$、
$x = g'(x')$、$s' = f'(x')$，并且 $s = g(s') = f(x)$。则：
\begin{enumerate}
\item $\dim_x(X_s) = \dim_{x'}(X'_{s'})$；
\item 若 $F$ 是态射 $X'_{s'} \to X_s$ 在 $x$ 上的纤维，则
$$
\dim(\mathcal{O}_{F, x'}) =
\dim(\mathcal{O}_{X'_{s'}, x'}) - \dim(\mathcal{O}_{X_s, x}) =
\text{trdeg}_{\kappa(s)}(\kappa(x)) -
\text{trdeg}_{\kappa(s')}(\kappa(x'))
$$
有 $\dim(\mathcal{O}_{X'_{s'}, x'}) \geq \dim(\mathcal{O}_{X_s, x})$，
且 $\text{trdeg}_{\kappa(s)}(\kappa(x)) \geq
\text{trdeg}_{\kappa(s')}(\kappa(x'))$。
\item 给定 $s', s, x$，存在一种对 $x'$ 的选择，使得
$\dim(\mathcal{O}_{X'_{s'}, x'}) = \dim(\mathcal{O}_{X_s, x})$，并且
$\text{trdeg}_{\kappa(s)}(\kappa(x)) = \text{trdeg}_{\kappa(s')}(\kappa(x'))$。
\end{enumerate}
\end{lemma}

\begin{proof}
(1) 立即来自《代数》引理
\ref{algebra-lemma-dimension-at-a-point-preserved-field-extension}。
(2) 和 (3) 来自《代数》引理
\ref{algebra-lemma-inequalities-under-field-extension}。
\end{proof}

\noindent
下一个引理来自一个非平凡的代数结果，即 Zariski 主定理的代数版本。

\begin{lemma}
\label{morphisms-lemma-openness-bounded-dimension-fibres}
\begin{reference}
\cite[IV Theorem 13.1.3]{EGA}
\end{reference}
令 $f : X \to S$ 为概形态射。令 $n \geq 0$。假设 $f$ 局部有限型。
集合
$$
U_n = \{x \in X \mid \dim_x X_{f(x)} \leq n\}
$$
在 $X$ 中是开集。
\end{lemma}

\begin{proof}
这立即来自《代数》引理
\ref{algebra-lemma-dimension-fibres-bounded-open-upstairs}
\end{proof}

\begin{lemma}
\label{morphisms-lemma-morphism-finite-type-bounded-dimension}
令 $f : X \to Y$ 为有限型态射，并假设 $Y$ 拟紧。则 $f$ 的纤维维数有界。
\end{lemma}

\begin{proof}
由引理 \ref{morphisms-lemma-openness-bounded-dimension-fibres}，集合
$U_n \subset X$ 中各点的纤维维数均为 $\leq n$，且该集合开。由于 $f$ 有限型，
每一点都包含在某个 $U_n$ 中（因为域上有限型代数的维数有限）。
由于 $Y$ 拟紧且 $f$ 有限型，可知 $X$ 拟紧。因此
$X = U_n$ 对某个 $n$ 成立。
\end{proof}

\begin{lemma}
\label{morphisms-lemma-openness-bounded-dimension-fibres-finite-presentation}
令 $f : X \to S$ 为概形态射。令 $n \geq 0$。假设 $f$ 局部有限表示。
引理 \ref{morphisms-lemma-openness-bounded-dimension-fibres} 中的开集
$$
U_n = \{x \in X \mid \dim_x X_{f(x)} \leq n\}
$$
在 $X$ 中逆拟紧。（见《拓扑学》定义
\ref{topology-definition-quasi-compact}。）
\end{lemma}

\begin{proof}
拓扑空间 $X$ 有一个由如下仿射开集组成的拓扑基：$U \subset X$，
且所诱导的态射 $f|_U : U \to S$ 通过某个仿射开集
$V \subset S$ 分解。因此只需证明 $U \cap U_n$ 拟紧，其中 $U$
是这样的开集。注意，$U_n \cap U$ 与开集
$\{x \in U \mid \dim_x U_{f(x)} \leq n\}$ 相同。
这把问题归结到 $X$ 和 $S$ 都是仿射概形的情形。在这种情形下，
本引理来自《代数》引理
\ref{algebra-lemma-dimension-fibres-bounded-quasi-compact-open-upstairs}
（以及引理 \ref{morphisms-lemma-locally-finite-presentation-characterize}）。
\end{proof}

\begin{lemma}
\label{morphisms-lemma-dimension-fibre-specialization}
令 $f : X \to S$ 为概形态射。令 $x \leadsto x'$ 为 $X$ 中点的一次
非平凡特化，并假设二者位于同一点 $s \in S$ 上。假设 $f$ 局部有限型。
则：
\begin{enumerate}
\item $\dim_x(X_s) \leq \dim_{x'}(X_s)$；
\item $\dim(\mathcal{O}_{X_s, x}) < \dim(\mathcal{O}_{X_s, x'})$；并且
\item $\text{trdeg}_{\kappa(s)}(\kappa(x)) >
\text{trdeg}_{\kappa(s)}(\kappa(x'))$。
\end{enumerate}
\end{lemma}

\begin{proof}
(1) 来自如下事实：$X_s$ 的任何含 $x'$ 的开集也含 $x$。
(2) 来自如下事实：$\mathcal{O}_{X_s, x}$ 是
$\mathcal{O}_{X_s, x'}$ 在一个素理想处的局部化，因此
$\mathcal{O}_{X_s, x}$ 中的任何素理想链都是
$\mathcal{O}_{X_s, x'}$ 中一条严格更长的素理想链的一部分。
最后一个不等式来自《代数》引理
\ref{algebra-lemma-tr-deg-specialization}。
\end{proof}






\section{指定相对维数的态射}
\label{morphisms-section-relative-dimension}

\noindent
为了方便地谈论指定相对维数的态射，引入下述概念。

\begin{definition}
\label{morphisms-definition-relative-dimension-d}
令 $f : X \to S$ 为概形态射。假设 $f$ 局部有限型。
\begin{enumerate}
\item 称 $f$ 具有{\it 相对维数 $\leq d$}（在 $x$ 处），如果
$\dim_x(X_{f(x)}) \leq d$。
\item 称 $f$ 具有{\it 相对维数 $\leq d$}，如果
$\dim_x(X_{f(x)}) \leq d$ 对所有 $x \in X$ 成立。
\item 称 $f$ 具有{\it 相对维数 $d$}，如果所有非空纤维 $X_s$
都等维且维数为 $d$。
\end{enumerate}
\end{definition}

\noindent
这个概念并非特别良态，但在许多情形下很好用。

\begin{lemma}
\label{morphisms-lemma-base-change-relative-dimension-d}
令 $f : X \to S$ 为局部有限型概形态射。若 $f$ 相对维数为 $d$，
则 $f$ 的任意基变换也如此。相对维数 $\leq d$ 的情形同样成立。
\end{lemma}

\begin{proof}
这立即来自引理 \ref{morphisms-lemma-dimension-fibre-after-base-change}。
\end{proof}

\begin{lemma}
\label{morphisms-lemma-composition-relative-dimension-d}
令 $f : X \to Y$、$g : Y \to Z$ 局部有限型。若 $f$ 相对维数
$\leq d$，且 $g$ 相对维数 $\leq e$，则 $g \circ f$ 相对维数
$\leq d + e$。若
\begin{enumerate}
\item $f$ 相对维数为 $d$；
\item $g$ 相对维数为 $e$；并且
\item $f$ 平坦，
\end{enumerate}
则 $g \circ f$ 相对维数为 $d + e$。
\end{lemma}

\begin{proof}
这立即来自引理 \ref{morphisms-lemma-dimension-fibre-at-a-point-additive}。
\end{proof}

\noindent
一般而言，不能把一个态射分解成相对维数分别为指定数值的各部分。
不过，当该态射的纤维是 Cohen--Macaulay 的并且该态射平坦且有限表示时，
可以作这种分解。

\begin{lemma}
\label{morphisms-lemma-flat-finite-presentation-CM-fibres-relative-dimension}
\begin{slogan}
Cohen--Macaulay 态射分解为纯相对维数的开闭部分
\end{slogan}
令 $f : X \to S$ 为概形态射。假设：
\begin{enumerate}
\item $f$ 平坦；
\item $f$ 局部有限表示；并且
\item 对所有 $s \in S$，纤维 $X_s$ 都是 Cohen--Macaulay 的
（《性质》定义 \ref{properties-definition-Cohen-Macaulay}）。
\end{enumerate}
则存在开闭子概形 $X_d \subset X$，使得
$X = \coprod_{d \geq 0} X_d$，并且 $f|_{X_d} : X_d \to S$
相对维数为 $d$。
\end{lemma}

\begin{proof}
这立即来自《代数》引理
\ref{algebra-lemma-relative-dimension-CM}。
\end{proof}

\begin{lemma}
\label{morphisms-lemma-locally-quasi-finite-rel-dimension-0}
令 $f : X \to S$ 为概形态射。假设 $f$ 局部有限型。令 $x \in X$，
且 $s = f(x)$。则 $f$ 在 $x$ 处拟有限，当且仅当
$\dim_x(X_s) = 0$。特别地，$f$ 局部拟有限，当且仅当 $f$ 相对维数为
$0$。
\end{lemma}

\begin{proof}
第一种证明。若 $f$ 在 $x$ 处拟有限，则 $\kappa(x)$ 是
$\kappa(s)$ 的有限扩张（由引理
\ref{morphisms-lemma-residue-field-quasi-finite}），且 $x$ 在 $X_s$ 中孤立
（由引理 \ref{morphisms-lemma-quasi-finite-at-point-characterize}），故由引理
\ref{morphisms-lemma-dimension-fibre-at-a-point} 有 $\dim_x(X_s) = 0$。
反之，若 $\dim_x(X_s) = 0$，则由引理
\ref{morphisms-lemma-dimension-fibre-at-a-point} 可知 $\kappa(s) \subset \kappa(x)$
是代数扩张，并且 $X_s$ 中没有其他点特化到 $x$。因此由引理
\ref{morphisms-lemma-algebraic-residue-field-extension-closed-point-fibre}，$x$
在其纤维中闭；再由引理
\ref{morphisms-lemma-quasi-finite-at-point-characterize} (3)，可得 $f$ 在 $x$
处拟有限。

\medskip\noindent
第二种证明。纤维 $X_s$ 是域上的局部有限型概形，因而局部 Noether
（引理 \ref{morphisms-lemma-finite-type-noetherian}）。结论现在来自引理
\ref{morphisms-lemma-quasi-finite-at-point-characterize} 和《性质》引理
\ref{properties-lemma-isolated-dim-0-locally-noetherian}。
\end{proof}

\begin{lemma}
\label{morphisms-lemma-rel-dimension-dimension}
令 $f : X \to Y$ 为局部 Noether 概形之间的态射，并假设它平坦、
局部有限型且相对维数为 $d$。对每一点 $x$（它属于 $X$），若其像为
$y$（位于 $Y$ 中），则有 $\dim_x(X) = \dim_y(Y) + d$。
\end{lemma}

\begin{proof}
把 $X$ 和 $Y$ 分别缩小到 $x$ 和 $y$ 的开邻域后，由概形在一点处的
维数的定义（《性质》定义 \ref{properties-definition-dimension}），
可以假设 $\dim(X) = \dim_x(X)$ 且 $\dim(Y) = \dim_y(Y)$。
由引理 \ref{morphisms-lemma-noetherian-finite-type-finite-presentation} 和
\ref{morphisms-lemma-fppf-open}，态射 $f$ 是开态射。因此可以缩小 $Y$，使得
$f$ 满。还需证明 $\dim(X) = \dim(Y) + d$。

\medskip\noindent
令 $a$ 为 $X$ 中一点，其像为 $b$，位于 $Y$ 中。由《代数》引理
\ref{algebra-lemma-dimension-base-fibre-equals-total}，
$$
\dim(\mathcal{O}_{X,a}) = \dim(\mathcal{O}_{Y,b}) + \dim(\mathcal{O}_{X_b, a}).
$$
对所有点 $a$（属于 $X$）取上确界，可得
$\dim(X) = \dim(Y) + d$，正如所需；见《性质》引理
\ref{properties-lemma-dimension}。
\end{proof}









\section{Syntomic 态射}
\label{morphisms-section-syntomic}

\noindent
$A$ 是域 $k$ 上的代数；称它为 $k$ 上的{\it 全局完全交}，如果
$A \cong k[x_1, \ldots, x_n]/(f_1, \ldots, f_c)$ 且
$\dim(A) = n - c$。$A$ 是域 $k$ 上的代数；称它为{\it 局部完全交}，
如果 $\Spec(A)$ 可以由标准开集覆盖，且每个标准开集都是 $k$
上的全局完全交。见《代数》节 \ref{algebra-section-lci}。
回忆，环同态 $R \to A$ 称为 {\it syntomic}，如果它有限表示、
平坦且其纤维为局部完全交环；见《代数》定义
\ref{algebra-definition-lci}。

\begin{definition}
\label{morphisms-definition-syntomic}
令 $f : X \to S$ 为概形态射。
\begin{enumerate}
\item 称 $f$ {\it 在 $x \in X$ 处 syntomic}，如果存在仿射开邻域
$\Spec(A) = U \subset X$（它是 $x$ 的邻域）以及仿射开集
$\Spec(R) = V \subset S$，使得 $f(U) \subset V$，且所诱导的环同态
$R \to A$ 是 syntomic 的。
\item 称 $f$ {\it syntomic}，如果它在 $X$ 的每一点处 syntomic。
\item 若 $S = \Spec(k)$ 且 $f$ syntomic，则称 $X$ 是
{\it $k$ 上的局部完全交}。
\item 仿射概形态射 $f : X \to S$ 称为{\it 标准 syntomic}，
如果存在全局相对完全交
$R \to R[x_1, \ldots, x_n]/(f_1, \ldots, f_c)$（见《代数》定义
\ref{algebra-definition-relative-global-complete-intersection}），使得
$X \to S$ 同构于
$$
\Spec(R[x_1, \ldots, x_n]/(f_1, \ldots, f_c)) \to \Spec(R).
$$
\end{enumerate}
\end{definition}

\noindent
文献中有时把 syntomic 态射称为{\it 平坦局部完全交态射}。
这是一类方便的态射。例如，可以用它们定义 syntomic 拓扑；
该拓扑比光滑拓扑和 \'etale 拓扑更细，但具有许多相同的形式性质。

\medskip\noindent
全局相对完全交（我们曾用它定义标准 syntomic 环同态）特别是平坦的。
在《态射进阶》节 \ref{more-morphisms-section-lci} 中，我们将考虑局部形如
$X \to S$ 的态射
$$
\Spec(R[x_1, \ldots, x_n]/(f_1, \ldots, f_c)) \to \Spec(R).
$$
其中 $f_1, \ldots, f_r$ 是 $R[x_1, \ldots, x_n]$ 中某个
Koszul--正则序列。这样的态射将称为{\it 局部完全交态射}。
一旦给出这个定义，态射 syntomic 当且仅当它是平坦的局部完全交态射。

\medskip\noindent
注意，syntomic 态射的定义中没有分离性或拟紧性假设。因此，
syntomic 是源上的局部性质。精确结果如下。

\begin{lemma}
\label{morphisms-lemma-syntomic-characterize}
令 $f : X \to S$ 为概形态射。下列条件等价：
\begin{enumerate}
\item 态射 $f$ syntomic。
\item 对每对仿射开集 $U \subset X$、$V \subset S$，若
$f(U) \subset V$，则环同态
$\mathcal{O}_S(V) \to \mathcal{O}_X(U)$ syntomic。
\item 存在开覆盖 $S = \bigcup_{j \in J} V_j$ 以及开覆盖
$f^{-1}(V_j) = \bigcup_{i \in I_j} U_i$，使得每个态射
$U_i \to V_j$ 都 syntomic，其中 $j\in J, i\in I_j$。
\item 存在仿射开覆盖 $S = \bigcup_{j \in J} V_j$ 以及仿射开覆盖
$f^{-1}(V_j) = \bigcup_{i \in I_j} U_i$，使得环同态
$\mathcal{O}_S(V_j) \to \mathcal{O}_X(U_i)$ syntomic，其中所有
$j\in J, i\in I_j$ 均如此。
\end{enumerate}
此外，若 $f$ syntomic，则对任意开子概形 $U \subset X$、$V \subset S$，
若 $f(U) \subset V$，限制 $f|_U : U \to V$ syntomic。
\end{lemma}

\begin{proof}
若能证明性质“$R \to A$ 是 syntomic 的”是局部的，结论就来自引理
\ref{morphisms-lemma-locally-P}。我们核验定义
\ref{morphisms-definition-property-local} 的条件 (a)、(b) 和 (c)。
由《代数》引理 \ref{algebra-lemma-base-change-syntomic}，
syntomic 在基变换下稳定，故条件 (a) 成立。由《代数》引理
\ref{algebra-lemma-composition-syntomic}，syntomic 在复合下稳定；
并且显然地，对任意环 $R$，环同态 $R \to R_f$ syntomic。
所以条件 (b) 成立。最后，由《代数》引理
\ref{algebra-lemma-local-syntomic}，条件 (c) 成立。
\end{proof}

\begin{lemma}
\label{morphisms-lemma-composition-syntomic}
两个 syntomic 态射的复合仍 syntomic。
\end{lemma}

\begin{proof}
在引理 \ref{morphisms-lemma-syntomic-characterize} 的证明中，我们看到
syntomic 是环同态的局部性质。因此，本引理的第一个断言来自引理
\ref{morphisms-lemma-composition-type-P}，并结合 syntomic 是在复合下稳定的
环同态性质这一事实；见《代数》引理
\ref{algebra-lemma-composition-syntomic}。
\end{proof}

\begin{lemma}
\label{morphisms-lemma-base-change-syntomic}
syntomic 态射的基变换仍 syntomic。
\end{lemma}

\begin{proof}
在引理 \ref{morphisms-lemma-syntomic-characterize} 的证明中，我们看到
syntomic 是环同态的局部性质。因此，本引理来自引理
\ref{morphisms-lemma-composition-type-P}，并结合 syntomic 是在基变换下稳定的
环同态性质这一事实；见《代数》引理
\ref{algebra-lemma-base-change-syntomic}。
\end{proof}

\begin{lemma}
\label{morphisms-lemma-open-immersion-syntomic}
任意开浸入都 syntomic。
\end{lemma}

\begin{proof}
这是因为开浸入是局部同构。
\end{proof}

\begin{lemma}
\label{morphisms-lemma-syntomic-locally-finite-presentation}
syntomic 态射局部有限表示。
\end{lemma}

\begin{proof}
这是因为按照定义，syntomic 环同态有限表示。
\end{proof}

\begin{lemma}
\label{morphisms-lemma-syntomic-flat}
syntomic 态射平坦。
\end{lemma}

\begin{proof}
这是因为按照定义，syntomic 环同态平坦。
\end{proof}

\begin{lemma}
\label{morphisms-lemma-syntomic-open}
syntomic 态射是泛开态射。
\end{lemma}

\begin{proof}
结合引理 \ref{morphisms-lemma-syntomic-locally-finite-presentation}、
\ref{morphisms-lemma-syntomic-flat} 和
\ref{morphisms-lemma-fppf-open}。
\end{proof}

\noindent
令 $k$ 为域。令 $A$ 为本质有限型的局部 $k$-代数（在 $k$ 上）。
回忆，$A$ 称为 $k$ 上的{\it 完全交}，如果可写成
$A \cong R/(f_1, \ldots, f_c)$，其中 $R$ 是本质有限型的正则局部环
（在 $k$ 上），且 $f_1, \ldots, f_c$ 是 $R$ 中的正则序列；
见《代数》定义 \ref{algebra-definition-lci-local-ring}。

\begin{lemma}
\label{morphisms-lemma-local-complete-intersection}
令 $k$ 为域。令 $X$ 为 $k$ 上的局部有限型概形。下列条件等价：
\begin{enumerate}
\item $X$ 是 $k$ 上的局部完全交；
\item 对每个 $x \in X$，存在仿射开集
$U = \Spec(R) \subset X$，它是 $x$ 的邻域，并且
$R \cong k[x_1, \ldots, x_n]/(f_1, \ldots, f_c)$ 是 $k$
上的全局完全交；并且
\item 对每个 $x \in X$，局部环 $\mathcal{O}_{X, x}$ 是 $k$
上的完全交。
\end{enumerate}
\end{lemma}

\begin{proof}
相应的代数结果见《代数》引理
\ref{algebra-lemma-lci-at-prime} 和
\ref{algebra-lemma-lci-global}。
\end{proof}

\noindent
下一个引理表明，每个 syntomic 态射在局部都是标准 syntomic 的。
因此可以用标准 syntomic 态射作为 syntomic 态射的{\it 局部模型}。
此外，它还表明，有限表示平坦态射 syntomic，当且仅当其纤维为局部完全交概形。

\begin{lemma}
\label{morphisms-lemma-syntomic-locally-standard-syntomic}
令 $f : X  \to S$ 为概形态射。令 $x \in X$ 为一点，其像为
$s = f(x)$。令 $V \subset S$ 为 $s$ 的仿射开邻域。下列条件等价：
\begin{enumerate}
\item 态射 $f$ 在 $x$ 处 syntomic。
\item 存在仿射开集 $U \subset X$，使得 $x \in U$、
$f(U) \subset V$，并且 $f|_U : U \to V$ 标准 syntomic。
\item 态射 $f$ 在 $x$ 处有限表示，局部环同态
$\mathcal{O}_{S, s} \to \mathcal{O}_{X, x}$ 平坦，并且
$\mathcal{O}_{X, x}/\mathfrak m_s \mathcal{O}_{X, x}$ 是
$\kappa(s)$ 上的完全交（见《代数》定义
\ref{algebra-definition-lci-local-ring}）。
\end{enumerate}
\end{lemma}

\begin{proof}
这来自各定义以及《代数》引理
\ref{algebra-lemma-syntomic}。
\end{proof}

\begin{lemma}
\label{morphisms-lemma-syntomic-flat-fibres}
令 $f : X \to S$ 为概形态射。若 $f$ 平坦、局部有限表示，且所有纤维
$X_s$ 都是局部完全交，则 $f$ syntomic。
\end{lemma}

\begin{proof}
由引理 \ref{morphisms-lemma-local-complete-intersection} 和
\ref{morphisms-lemma-syntomic-locally-standard-syntomic}，以及局部环同构
$
\mathcal{O}_{X, x}/\mathfrak m_s \mathcal{O}_{X, x}
\cong
\mathcal{O}_{X_s, x}
$，结论显然。
\end{proof}

\begin{lemma}
\label{morphisms-lemma-set-points-where-fibres-lci}
令 $f : X \to S$ 为概形态射。假设 $f$ 局部有限型。集合
$$
T = \{x \in X \mid \mathcal{O}_{X_{f(x)}, x}
\text{ is a complete intersection over }\kappa(f(x))\}
$$
的形成与任意基变换可交换：对任意态射 $g : S' \to S$，考虑基变换
$f' : X' \to S'$（它来自 $f$）以及投影 $g' : X' \to X$。
则集合 $T'$ 对应于态射 $f'$，且等于 $T' = (g')^{-1}(T)$。
特别地，若假设 $f$ 平坦且局部有限表示，则 $f$ 为 syntomic 的点所成
开集也有同一性质。
\end{lemma}

\begin{proof}
令 $s' \in S'$ 为一点，并令 $s = g(s')$。则有
$$
X'_{s'} =
\Spec(\kappa(s')) \times_{\Spec(\kappa(s))} X_s
$$
换言之，基变换的纤维就是纤维的基变换。因此第一部分等价于《代数》引理
\ref{algebra-lemma-lci-field-change-local}。第二部分来自第一部分，
因为在这种情形下，依引理
\ref{morphisms-lemma-syntomic-locally-standard-syntomic}，$T$ 正是 $f$
为 syntomic 的点所成集合。
\end{proof}

\begin{lemma}
\label{morphisms-lemma-standard-syntomic-relative-dimension}
令 $R$ 为环。令
$R \to A = R[x_1, \ldots, x_n]/(f_1, \ldots, f_c)$ 为全局相对完全交。
置 $S = \Spec(R)$ 且 $X = \Spec(A)$。考虑态射
$f : X \to S$，它与环同态 $R \to A$ 相伴。函数
$x \mapsto \dim_x(X_{f(x)})$ 是取值 $n - c$ 的常值函数。
\end{lemma}

\begin{proof}
由《代数》定义
\ref{algebra-definition-relative-global-complete-intersection}，
$R \to A$ 为全局相对完全交意味着所有非零纤维环的维数都是
$n - c$。因此，对素理想 $\mathfrak p$（属于 $R$），纤维环
$\kappa(\mathfrak p)[x_1, \ldots, x_n]/(\overline{f}_1, \ldots, \overline{f}_c)$
或者为零，或者是维数为 $n - c$ 的全局完全交环。由《代数》定义
\ref{algebra-definition-lci-field} 后面的讨论，这推出它等维且维数为
$n - c$。本引理得证。
\end{proof}

\begin{lemma}
\label{morphisms-lemma-syntomic-relative-dimension}
令 $f : X \to S$ 为 syntomic 态射。函数
$x \mapsto \dim_x(X_{f(x)})$ 在 $X$ 上局部常值。
\end{lemma}

\begin{proof}
由引理 \ref{morphisms-lemma-syntomic-locally-standard-syntomic}，态射 $f$
局部看起来像仿射概形之间的标准 syntomic 态射。因此结论来自引理
\ref{morphisms-lemma-standard-syntomic-relative-dimension}。
\end{proof}

\noindent
引理 \ref{morphisms-lemma-syntomic-relative-dimension} 表明下述定义是有意义的。

\begin{definition}
\label{morphisms-definition-syntomic-relative-dimension}
令 $d \geq 0$ 为整数。概形态射 $f : X \to S$ 被称为
{\it 相对维数 $d$ 的 syntomic 态射}，是指 $f$ syntomic，且函数满足
$\dim_x(X_{f(x)}) = d$，其中该等式对所有 $x \in X$ 成立。
\end{definition}

\noindent
换言之，$f$ syntomic，且非空纤维等维且维数为 $d$。

\begin{lemma}
\label{morphisms-lemma-syntomic-permanence}
令
$$
\xymatrix{
X \ar[rr]_f \ar[rd]_p & &
Y \ar[dl]^q \\
& S
}
$$
为概形态射的交换图。假设：
\begin{enumerate}
\item $f$ 满且 syntomic；
\item $p$ syntomic；并且
\item $q$ 局部有限表示\footnote{事实上，若 $f$ 满、平坦且局部有限表示，
并且 $p$ syntomic，则 $q$ 和 $f$ 都 syntomic；见《下降》引理
\ref{descent-lemma-syntomic-permanence}。}。
\end{enumerate}
则 $q$ syntomic。
\end{lemma}

\begin{proof}
由引理 \ref{morphisms-lemma-flat-permanence}，可知 $q$ 平坦。因此只需证明
$Y \to S$ 的纤维都是局部完全交；见引理
\ref{morphisms-lemma-syntomic-flat-fibres}。令 $s \in S$。考虑态射
$X_s \to Y_s$。这是态射 $X \to Y$ 的基变换，因而满且 syntomic
（引理 \ref{morphisms-lemma-base-change-syntomic}）。出于同一理由，$X_s$
在 $\kappa(s)$ 上 syntomic。此外，$Y_s$ 在 $\kappa(s)$ 上局部有限型
（引理 \ref{morphisms-lemma-base-change-finite-type}）。这样，问题归结到 $S$
是一个域的谱的情形。

\medskip\noindent
假设 $S = \Spec(k)$。令 $y \in Y$。选取仿射开集
$\Spec(A) \subset Y$，作为 $y$ 的邻域。令 $\Spec(B) \subset X$ 为仿射开集，使得
$f(\Spec(B)) \subset \Spec(A)$，并且包含某点 $x \in X$，满足
$f(x) = y$。选取满射 $k[x_1, \ldots, x_n] \to A$，其核为 $I$。
选取满射 $A[y_1, \ldots, y_m] \to B$；它又诱导满射
$k[x_i, y_j] \to B$，其核为 $J$。令
$\mathfrak q \subset k[x_i, y_j]$ 为对应于 $y \in \Spec(B)$ 的素理想，
并令 $\mathfrak p \subset k[x_i]$ 为对应于 $x \in \Spec(A)$ 的素理想。
由于 $x$ 映到 $y$，有 $\mathfrak p = \mathfrak q \cap k[x_i]$。
考虑下列局部环交换图：
$$
\xymatrix{
\mathcal{O}_{X, x} \ar@{=}[r] &
B_{\mathfrak q} &
k[x_1, \ldots, x_n, y_1, \ldots, y_m]_{\mathfrak q} \ar[l] \\
\mathcal{O}_{Y, y} \ar@{=}[r] \ar[u] & A_{\mathfrak p} \ar[u] &
k[x_1, \ldots, x_n]_{\mathfrak p} \ar[l] \ar[u]
}
$$
我们断言《代数》引理
\ref{algebra-lemma-lci-permanence-initial} 的假设均满足。条件 (1) 和 (2)
显然。由于 $X \to Y$ 平坦，条件 (4) 成立。条件 (3) 成立，
因为环 $\mathcal{O}_{Y, y}$ 和
$\mathcal{O}_{X_y, x} = \mathcal{O}_{X, x}/\mathfrak m_y\mathcal{O}_{X, x}$
由 $f$ 和 $p$ syntomic 这一假设都是完全交环；见引理
\ref{morphisms-lemma-syntomic-locally-standard-syntomic}。《代数》引理
\ref{algebra-lemma-lci-permanence-initial} 的结论恰好是
$\mathcal{O}_{Y, y}$ 为完全交环！所以再次由引理
\ref{morphisms-lemma-syntomic-locally-standard-syntomic}，可知 $Y$ 在 $k$
上于 $y$ 处 syntomic，正如所需。
\end{proof}

\section{浸入的余法层}
\label{morphisms-section-conormal-sheaf}

\noindent
令 $i : Z \to X$ 为闭浸入。令
$\mathcal{I} \subset \mathcal{O}_X$ 为相应的拟凝聚理想层。考虑 $X$
上拟凝聚层的短正合列
$$
0 \to \mathcal{I}^2 \to \mathcal{I} \to \mathcal{I}/\mathcal{I}^2 \to 0
$$
。由于层 $\mathcal{I}/\mathcal{I}^2$ 被 $\mathcal{I}$ 零化，故由引理
\ref{morphisms-lemma-i-star-equivalence}，它对应于 $Z$ 上的一个层。这个拟凝聚
$\mathcal{O}_Z$-模称为 $Z$ 在 $X$ 中的{\it 余法层}；按照
第 \ref{morphisms-section-closed-immersions-quasi-coherent} 节所述的记号滥用，
常将其简记为 $\mathcal{I}/\mathcal{I}^2$。

\medskip\noindent
若 $i : Z \to X$ 是（局部闭）浸入，则把 $i$ 的余法层定义为闭浸入
$i : Z \to X \setminus \partial Z$ 的余法层，其中
$\partial Z = \overline{Z} \setminus Z$。它常记作
$\mathcal{I}/\mathcal{I}^2$，这里 $\mathcal{I}$ 是闭浸入
$i : Z \to X \setminus \partial Z$ 的理想层。

\begin{definition}
\label{morphisms-definition-conormal-sheaf}
令 $i : Z \to X$ 为浸入。称 $\mathcal{C}_{Z/X}$ 为 $Z$ 在 $X$
中的{\it 余法层}，也称为 $i$ 的{\it 余法层}；它是上述拟凝聚
$\mathcal{O}_Z$-模 $\mathcal{I}/\mathcal{I}^2$。
\end{definition}

\noindent
在 \cite[IV Definition 16.1.2]{EGA} 中，这个层记作
$\mathcal{N}_{Z/X}$。我们不采用这一约定，因为希望保留记号
$\mathcal{N}_{Z/X}$ 表示{\it 该浸入的法层}。当余法层有限表示时，
它定义为
$$
\mathcal{N}_{Z/X} =
\SheafHom_{\mathcal{O}_Z}(\mathcal{C}_{Z/X}, \mathcal{O}_Z) =
\SheafHom_{\mathcal{O}_Z}(\mathcal{I}/\mathcal{I}^2, \mathcal{O}_Z)
$$
（否则法层甚至未必拟凝聚）。稍后我们会回到法层（此处待插未来引用）。

\begin{lemma}
\label{morphisms-lemma-affine-conormal}
令 $i : Z \to X$ 为浸入。$i$ 的余法层具有下列性质：
\begin{enumerate}
\item 令 $U \subset X$ 为任意开子概形，使 $i$ 可分解为
$Z \xrightarrow{i'} U \to X$，其中 $i'$ 是闭浸入。令
$\mathcal{I} = \Ker((i')^\sharp) \subset \mathcal{O}_U$。则
$$
\mathcal{C}_{Z/X} = (i')^*\mathcal{I}\quad\text{且}\quad
i'_*\mathcal{C}_{Z/X} = \mathcal{I}/\mathcal{I}^2
$$
\item 对任意仿射开集 $\Spec(R) = U \subset X$，若
$Z \cap U = \Spec(R/I)$，则有典范同构
$\Gamma(Z \cap U, \mathcal{C}_{Z/X}) = I/I^2$。
\end{enumerate}
\end{lemma}

\begin{proof}
大体上由定义即得。注意，对环 $R$ 及其理想 $I$（属于 $R$），有
$I/I^2 = I \otimes_R R/I$。细节从略。
\end{proof}

\begin{lemma}
\label{morphisms-lemma-conormal-functorial}
令
$$
\xymatrix{
Z \ar[r]_i \ar[d]_f & X \ar[d]^g \\
Z' \ar[r]^{i'} & X'
}
$$
为概形范畴中的交换图。假设 $i$、$i'$ 均为浸入。存在典范的
$\mathcal{O}_Z$-模同态
$$
f^*\mathcal{C}_{Z'/X'}
\longrightarrow
\mathcal{C}_{Z/X}
$$
，它由下述性质刻画：对每一对仿射开集
$(\Spec(R) = U \subset X, \Spec(R') = U' \subset X')$，若
$g(U) \subset U'$，且
$Z \cap U = \Spec(R/I)$ 和 $Z' \cap U' = \Spec(R'/I')$，则诱导同态
$$
\Gamma(Z' \cap U', \mathcal{C}_{Z'/X'}) = I'/I'^2
\longrightarrow
I/I^2 = \Gamma(Z \cap U, \mathcal{C}_{Z/X})
$$
由环同态 $g^\sharp : R' \to R$ 诱导；该环同态满足
$g^\sharp(I') \subset I$。
\end{lemma}

\begin{proof}
令 $\partial Z' = \overline{Z'} \setminus Z'$，并令
$\partial Z = \overline{Z} \setminus Z$。它们分别是 $X'$ 和 $X$
的闭子集。把 $X'$ 替换为 $X' \setminus \partial Z'$，并把 $X$
替换为 $X \setminus \big(g^{-1}(\partial Z') \cup \partial Z\big)$，
便可假设 $i$ 和 $i'$ 都是闭浸入。

\medskip\noindent
$g \circ i$ 分解通过 $i'$ 这一事实蕴含：$g^*\mathcal{I}'$ 映入
$\mathcal{I}$；这里所用的典范同态为 $g^*\mathcal{I}' \to \mathcal{O}_X$。
见《概形》引理
\ref{schemes-lemma-characterize-closed-subspace} 和
\ref{schemes-lemma-restrict-map-to-closed}。因此得到拟凝聚层的诱导同态
$g^*(\mathcal{I}'/(\mathcal{I}')^2) \to \mathcal{I}/\mathcal{I}^2$。
再由 $i$ 拉回，得到
$i^*g^*(\mathcal{I}'/(\mathcal{I}')^2) \to i^*(\mathcal{I}/\mathcal{I}^2)$。
注意 $i^*(\mathcal{I}/\mathcal{I}^2) = \mathcal{C}_{Z/X}$。另一方面，
$i^*g^*(\mathcal{I}'/(\mathcal{I}')^2) =
f^*(i')^*(\mathcal{I}'/(\mathcal{I}')^2) = f^*\mathcal{C}_{Z'/X'}$。
这就给出了所需同态。

\medskip\noindent
展开定义即可验证该同态在局部上由给定同态
$I'/(I')^2 \to I/I^2$ 描述，细节从略。还要注意，对任意
$x \in i(Z)$，确实存在仿射开邻域 $U$、$U'$，使得
$f(U) \subset U'$，并使 $Z \cap U$ 和 $U' \cap Z'$ 都闭，且
$x \in U$。证明从略。因此，本引理的要求确实刻画了该同态
（也可以用来定义它）。
\end{proof}

\begin{lemma}
\label{morphisms-lemma-conormal-functorial-flat}
令
$$
\xymatrix{
Z \ar[r]_i \ar[d]_f & X \ar[d]^g \\
Z' \ar[r]^{i'} & X'
}
$$
为概形范畴中的纤维积图，其中 $i$、$i'$ 为浸入。则引理
\ref{morphisms-lemma-conormal-functorial} 的典范同态
$f^*\mathcal{C}_{Z'/X'} \to \mathcal{C}_{Z/X}$ 是满射。若 $g$ 平坦，
则它是同构。
\end{lemma}

\begin{proof}
令 $R' \to R$ 为环同态，$I' \subset R'$ 为理想。置 $I = I'R$。
则 $I'/(I')^2 \otimes_{R'} R \to I/I^2$ 是满射。若
$R' \to R$ 平坦，则 $I = I' \otimes_{R'} R$ 且
$I^2 = (I')^2 \otimes_{R'} R$，从而可见该同态为同构。
\end{proof}

\begin{lemma}
\label{morphisms-lemma-transitivity-conormal}
令 $Z \to Y \to X$ 为概形的浸入。则有典范正合列
$$
i^*\mathcal{C}_{Y/X} \to
\mathcal{C}_{Z/X} \to
\mathcal{C}_{Z/Y} \to 0
$$
，其中同态来自引理 \ref{morphisms-lemma-conormal-functorial}，而
$i : Z \to Y$ 是第一个态射。
\end{lemma}

\begin{proof}
通过引理 \ref{morphisms-lemma-conormal-functorial}，这转化为下述代数事实。
设 $C \to B \to A$ 为满环同态。令 $I = \Ker(B \to A)$、
$J = \Ker(C \to A)$，并令 $K = \Ker(C \to B)$。则有正合列
$$
K/K^2 \otimes_B A \to J/J^2 \to I/I^2 \to 0.
$$
这立即由 $I = J/K$ 得出。
\end{proof}










\section{态射的微分层}
\label{morphisms-section-sheaf-differentials}

\noindent
我们建议读者参阅交换代数章中的相应一节
（《代数》第 \ref{algebra-section-differentials} 节），以及模层章中的相应一节
（《模》第 \ref{modules-section-differentials} 节）。

\begin{definition}
\label{morphisms-definition-sheaf-differentials}
令 $f : X \to S$ 为概形态射。{\it 微分层 $\Omega_{X/S}$}，即 $X$ 在
$S$ 上的微分层，是把 $f$ 视为环层空间态射时的微分层
（《模》定义 \ref{modules-definition-differentials}），并带有其{\it 泛
$S$-导子}
$$
\text{d}_{X/S} : \mathcal{O}_X \longrightarrow \Omega_{X/S}.
$$
\end{definition}

\noindent
$\Omega_{X/S}$ 是拟凝聚 $\mathcal{O}_X$-模。例如，它同构于对角态射
$\Delta : X \to X \times_S X$ 的余法层（引理
\ref{morphisms-lemma-differentials-diagonal}）。我们已通过泛性质定义 $X$ 在 $S$
上的微分模，即把它定义为泛导子的取值对象。若对相对微分层另有任何
满足这一泛性质的构造，则由 Yoneda 引理，它与上述构造典范同构。
为方便起见，我们在此重述这一泛性质。

\begin{lemma}
\label{morphisms-lemma-universal-derivation-universal}
令 $f : X \to S$ 为概形态射。映射
$$
\Hom_{\mathcal{O}_X}(\Omega_{X/S}, \mathcal{F})
\longrightarrow
\text{Der}_S(\mathcal{O}_X, \mathcal{F}),\quad
\alpha \longmapsto \alpha \circ \text{d}_{X/S}
$$
是函子 $\textit{Mod}(\mathcal{O}_X) \to \textit{Sets}$ 的同构。
\end{lemma}

\begin{proof}
这只是定义的重述。
\end{proof}

\begin{lemma}
\label{morphisms-lemma-differentials-restrict-open}
令 $f : X \to S$ 为概形态射。令 $U \subset X$、$V \subset S$ 为开子概形，
且 $f(U) \subset V$。则存在唯一同构
$\Omega_{X/S}|_U = \Omega_{U/V}$（作为 $\mathcal{O}_U$-模），使得
$\text{d}_{X/S}|_U = \text{d}_{U/V}$。
\end{lemma}

\begin{proof}
利用典范恒等式
$f^{-1}\mathcal{O}_S|_U = (f|_U)^{-1}\mathcal{O}_V$，这就是《模》引理
\ref{modules-lemma-localize-differentials} 的特殊情形。
\end{proof}

\noindent
以下我们将采用这些典范恒等式，直接用 $\Omega_{U/S}$ 或
$\Omega_{U/V}$ 表示 $\Omega_{X/S}$ 在 $U$ 上的限制。

\begin{lemma}
\label{morphisms-lemma-affine-case-derivation}
令 $R \to A$ 为环同态。令 $\mathcal{F}$ 为 $\mathcal{O}_X$-模层，
定义在 $X = \Spec(A)$ 上。置 $S = \Spec(R)$。把 $S$-导子作用于全局截面，
由此得到 $\mathcal{F}$ 的 $S$-导子之集（作用于 $\mathcal{F}$）与
$R$-导子之集（定义在 $M = \Gamma(X, \mathcal{F})$ 上）之间的双射。
\end{lemma}

\begin{proof}
令 $D : A \to M$ 为 $R$-导子。我们须证明，存在唯一的 $S$-导子，
作用于 $\mathcal{F}$，并且它在全局截面上给出 $D$。令
$U = D(f) \subset X$ 为标准仿射开集。$\Gamma(U, \mathcal{O}_X)$
的任意元素形如 $a/f^n$，其中 $a \in A$ 且 $n \geq 0$。由 Leibniz 法则，
在 $\Gamma(U, \mathcal{F})$ 中有
$$
D(a)|_U = a/f^n D(f^n)|_U + f^n D(a/f^n)
$$
。由于 $f$ 在 $\Gamma(U, \mathcal{F})$ 上可逆，这完全确定了
$D(a/f^n) \in \Gamma(U, \mathcal{F})$ 的值，从而证明唯一性。
存在性只需定义
$$
D(a/f^n) := (1/f^n) D(a)|_U - a/f^{2n} D(f^n)|_U
$$
，再证明它具有一切所需性质（在 $X$ 的标准开集基上）即可。细节从略。
\end{proof}

\begin{lemma}
\label{morphisms-lemma-differentials-affine}
令 $f : X \to S$ 为概形态射。对任意一对仿射开集
$\Spec(A) = U \subset X$、$\Spec(R) = V \subset S$，若 $f(U) \subset V$，
则存在唯一同构
$$
\Gamma(U, \Omega_{X/S}) = \Omega_{A/R}.
$$
它与 $\text{d}_{X/S}$ 以及 $\text{d} : A \to \Omega_{A/R}$ 相容。
\end{lemma}

\begin{proof}
由引理 \ref{morphisms-lemma-differentials-restrict-open}，可将 $X$、$S$ 分别替换为
$U$、$V$。因而可假设 $X = \Spec(A)$ 且 $S = \Spec(R)$，并须在
$U = X$、$V = S$ 时证明本引理。考虑 $A$-模
$M = \Gamma(X, \Omega_{X/S})$ 及其 $R$-导子
$\text{d}_{X/S} : A \to M$。令 $N$ 为另一个 $A$-模，并以
$\widetilde{N}$ 表示拟凝聚 $\mathcal{O}_X$-模，它与 $N$ 相伴；见《概形》
第 \ref{schemes-section-quasi-coherent-affine} 节。与
$\text{d}_{X/S} : A \to M$ 预复合，得到箭头
$$
\alpha : \Hom_A(M, N) \longrightarrow \text{Der}_R(A, N)
$$
。利用引理 \ref{morphisms-lemma-universal-derivation-universal} 和
\ref{morphisms-lemma-affine-case-derivation}，得到恒等式
$$
\Hom_{\mathcal{O}_X}(\Omega_{X/S}, \widetilde{N}) =
\text{Der}_S(\mathcal{O}_X, \widetilde{N}) =
\text{Der}_R(A, N)
$$
。取全局截面给出箭头
$\Hom_{\mathcal{O}_X}(\Omega_{X/S}, \widetilde{N}) \to \Hom_R(M, N)$。
把这个箭头与上述恒等式合并，得到箭头
$$
\beta : \text{Der}_R(A, N) \longrightarrow \Hom_R(M, N)
$$
。考察全局截面可知，$\alpha$ 与 $\beta$ 互为逆。因此
$\text{d}_{X/S} : A \to M$ 与 $\text{d} : A \to \Omega_{A/R}$
满足同一泛性质；见《代数》引理 \ref{algebra-lemma-universal-omega}。
于是 Yoneda 引理（《范畴》引理 \ref{categories-lemma-yoneda}）蕴含存在唯一的
$A$-模同构 $M \cong \Omega_{A/R}$，并且它与导子相容。
\end{proof}

\begin{remark}
\label{morphisms-remark-differentials-glue}
上述引理给出了构造微分模的第二种方法。具体地，令
$f : X \to S$ 为概形态射。考虑所有仿射开集 $U \subset X$，它们映入
$S$ 的某个仿射开集。它们构成 $X$ 的拓扑的一个基。因此，只须定义
$\Gamma(U, \Omega_{X/S})$（对这样的 $U$）。对这样的开集，我们直接置
$\Gamma(U, \Omega_{X/S}) = \Omega_{A/R}$；这里 $A$、$R$ 如上述引理
\ref{morphisms-lemma-differentials-affine} 中所述。这一办法可行，但要构造限制映射
并验证该方案满足层条件，还需要略多一些代数准备。
\end{remark}

\noindent
下面的引理给出了定义微分层的又一种方法，并特别表明 $\Omega_{X/S}$
在 $X$、$S$ 为概形时拟凝聚。

\begin{lemma}
\label{morphisms-lemma-differentials-diagonal}
令 $f : X \to S$ 为概形态射。$\Omega_{X/S}$ 与对角态射
$\Delta_{X/S} : X \longrightarrow X \times_S X$ 的余法层之间存在典范同构。
\end{lemma}

\begin{proof}
我们先构造一对“全局”层及层的全局同态，再在后面描述某些仿射开集上的构造。

\medskip\noindent
回忆 $\Delta = \Delta_{X/S} : X \to X \times_S X$ 是浸入；见《概形》引理
\ref{schemes-lemma-diagonal-immersion}。令 $\mathcal{J}$ 为该浸入的理想层；
它定义在某个开子概形 $W$ 上，该开子概形属于 $X \times_S X$，使得
$\Delta(X) \subset W$ 闭。取《概形》引理
\ref{schemes-lemma-diagonal-immersion} 的证明中找到的那个开集。注意，环层
$\mathcal{O}_W/\mathcal{J}^2$ 的支集包含于 $\Delta(X)$。此外，它位于层的短正合列
$$
0 \to \mathcal{J}/\mathcal{J}^2
\to \mathcal{O}_W/\mathcal{J}^2
\to \Delta_*\mathcal{O}_X
\to 0.
$$
中。利用 $\Delta^{-1}$，可把它看成满的 $f^{-1}\mathcal{O}_S$-代数层同态，
其核为 $\Delta$ 的余法层（见定义 \ref{morphisms-definition-conormal-sheaf} 和引理
\ref{morphisms-lemma-affine-conormal}）：
$$
0 \to \mathcal{C}_{X/X \times_S X}
\to \Delta^{-1}(\mathcal{O}_W/\mathcal{J}^2)
\to \mathcal{O}_X
\to 0
$$
。这把我们置于《模》引理
\ref{modules-lemma-double-structure-gives-derivation} 的情形中。投影态射
$p_i : X \times_S X \to X$（$i = 1, 2$）诱导环层同态
$(p_i)^\sharp : (p_i)^{-1}\mathcal{O}_X \to \mathcal{O}_{X \times_S X}$。
将它限制到 $W$ 并模去 $\mathcal{J}^2$，得到
$(p_i)^{-1}\mathcal{O}_X \to \mathcal{O}_W/\mathcal{J}^2$。由于
$\Delta^{-1}p_i^{-1}\mathcal{O}_X = \mathcal{O}_X$，得到同态
$$
s_i : \mathcal{O}_X \to \Delta^{-1}(\mathcal{O}_W/\mathcal{J}^2).
$$
$s_1$ 与 $s_2$ 都是同态
$\Delta^{-1}(\mathcal{O}_W/\mathcal{J}^2) \to \mathcal{O}_X$ 的截面，
正如《模》引理 \ref{modules-lemma-double-structure-gives-derivation} 中那样。
因此得到 $S$-导子
$\text{d} = s_2 - s_1 : \mathcal{O}_X \to \mathcal{C}_{X/X \times_S X}$。
由微分模的泛性质，得到唯一的 $\mathcal{O}_X$-线性同态
$$
\Omega_{X/S} \longrightarrow \mathcal{C}_{X/X \times_S X},\quad
f\text{d}g \longmapsto fs_2(g) - fs_1(g)
$$
。为证明该同态是同构，我们在适当的仿射开集上将它具体写出。可以覆盖
$X$，所用仿射开集为 $\Spec(A) = U \subset X$，使其像包含于仿射开集
$\Spec(R) = V \subset S$ 中。根据《概形》引理
\ref{schemes-lemma-diagonal-immersion} 的证明，
$U \times_V U \subset X \times_S X$ 是包含于上述 $W$ 中的仿射开集。此外，
$U \times_V U = \Spec(A \otimes_R A)$。层 $\mathcal{J}$ 对应于理想
$J = \Ker(A \otimes_R A \to A)$。该短正合列对应于
$A \otimes_R A$-模的短正合列
$$
0 \to J/J^2 \to (A \otimes_R A)/J^2 \to A \to 0
$$
。截面 $s_i$ 对应于环同态
$$
A \longrightarrow (A \otimes_R A)/J^2,\quad
s_1 : a \mapsto a \otimes 1,\quad
s_2 : a \mapsto 1 \otimes a.
$$
由引理 \ref{morphisms-lemma-affine-conormal}，有
$\Gamma(U, \mathcal{C}_{X/X \times_S X}) = J/J^2$；由引理
\ref{morphisms-lemma-differentials-affine}，有
$\Gamma(U, \Omega_{X/S}) = \Omega_{A/R}$。上述同态就是
$a \text{d}b \mapsto a \otimes b - ab \otimes 1$，它在《代数》引理
\ref{algebra-lemma-differentials-diagonal} 中已证明为同构。
\end{proof}

\begin{lemma}
\label{morphisms-lemma-functoriality-differentials}
令
$$
\xymatrix{
X' \ar[d] \ar[r]_f & X \ar[d] \\
S' \ar[r] & S
}
$$
为概形的交换图。将典范同态
$\mathcal{O}_X \to f_*\mathcal{O}_{X'}$ 与同态
$f_*\text{d}_{X'/S'} : f_*\mathcal{O}_{X'} \to f_*\Omega_{X'/S'}$ 复合，
得到一个 $S$-导子。因此得到典范的 $\mathcal{O}_X$-模同态
$\Omega_{X/S} \to f_*\Omega_{X'/S'}$；再由 $f_*$ 与 $f^*$ 的伴随性，
得到典范的 $\mathcal{O}_{X'}$-模同态
$$
c_f : f^*\Omega_{X/S} \longrightarrow \Omega_{X'/S'}.
$$
它由下述性质唯一刻画：$f^*\text{d}_{X/S}(h)$ 映到
$\text{d}_{X'/S'}(f^* h)$，其中 $h$ 是 $\mathcal{O}_X$ 的任意局部截面。
\end{lemma}

\begin{proof}
这是《模》引理
\ref{modules-lemma-functoriality-differentials-ringed-spaces} 的特殊情形。
对于概形，也可利用余法层的函子性（见引理
\ref{morphisms-lemma-conormal-functorial}）以及引理
\ref{morphisms-lemma-differentials-diagonal} 来定义 $c_f$。还可以利用本引理最后一行的
刻画，把仿射片上定义的同态粘合起来（见《代数》方程
(\ref{algebra-equation-functorial-omega})）。
\end{proof}

\begin{lemma}
\label{morphisms-lemma-triangle-differentials}
令 $f : X \to Y$、$g : Y \to S$ 为概形态射。则有典范正合列
$$
f^*\Omega_{Y/S} \to \Omega_{X/S} \to \Omega_{X/Y} \to 0
$$
，其中同态来自引理 \ref{morphisms-lemma-functoriality-differentials} 的应用。
\end{lemma}

\begin{proof}
这是《代数》引理 \ref{algebra-lemma-exact-sequence-differentials} 的层化版本。
也可以使用环层空间态射的一般版本；见《模》引理
\ref{modules-lemma-triangle-differentials}。
\end{proof}

\begin{lemma}
\label{morphisms-lemma-base-change-differentials}
令 $X \to S$ 为概形态射。令 $g : S' \to S$ 为概形态射。
令 $X' = X_{S'}$ 为 $X$ 的基变换。以 $g' : X' \to X$ 表示投影。
则引理 \ref{morphisms-lemma-functoriality-differentials} 的同态
$$
(g')^*\Omega_{X/S} \to \Omega_{X'/S'}
$$
为同构。
\end{lemma}

\begin{proof}
这是《代数》引理 \ref{algebra-lemma-differentials-base-change} 的层化版本。
\end{proof}

\begin{lemma}
\label{morphisms-lemma-differential-product}
令 $f : X \to S$ 与 $g : Y \to S$ 为具有同一目标的概形态射。令
$p : X \times_S Y \to X$ 与 $q : X \times_S Y \to Y$ 为投影态射。
引理 \ref{morphisms-lemma-functoriality-differentials} 给出的同态
$$
p^*\Omega_{X/S} \oplus q^*\Omega_{Y/S}
\longrightarrow
\Omega_{X \times_S Y/S}
$$
是同构。
\end{lemma}

\begin{proof}
由引理 \ref{morphisms-lemma-base-change-differentials}，复合
$p^*\Omega_{X/S} \to \Omega_{X \times_S Y/S} \to \Omega_{X \times_S Y/Y}$
为同构；对 $q$ 亦然。此外，由引理
\ref{morphisms-lemma-triangle-differentials}，同态
$p^*\Omega_{X/S} \to \Omega_{X \times_S Y/S}$ 的余核是
$\Omega_{X \times_S Y/X}$。结论随即成立。
\end{proof}

\begin{lemma}
\label{morphisms-lemma-finite-type-differentials}
令 $f : X \to S$ 为概形态射。若 $f$ 局部有限型，则
$\Omega_{X/S}$ 是有限型 $\mathcal{O}_X$-模。
\end{lemma}

\begin{proof}
立即由《代数》引理
\ref{algebra-lemma-differentials-finitely-generated}、引理
\ref{morphisms-lemma-differentials-affine}、引理
\ref{morphisms-lemma-locally-finite-type-characterize} 以及《性质》引理
\ref{properties-lemma-finite-type-module} 得出。
\end{proof}

\begin{lemma}
\label{morphisms-lemma-finite-presentation-differentials}
令 $f : X \to S$ 为概形态射。若 $f$ 局部有限表示，则
$\Omega_{X/S}$ 是有限表示 $\mathcal{O}_X$-模。
\end{lemma}

\begin{proof}
立即由《代数》引理
\ref{algebra-lemma-differentials-finitely-presented}、引理
\ref{morphisms-lemma-differentials-affine}、引理
\ref{morphisms-lemma-locally-finite-presentation-characterize} 以及《性质》引理
\ref{properties-lemma-finite-presentation-module} 得出。
\end{proof}

\begin{lemma}
\label{morphisms-lemma-immersion-differentials}
若 $X \to S$ 是浸入，或更一般地是单态射，则 $\Omega_{X/S}$ 为零。
\end{lemma}

\begin{proof}
这是因为此时 $\Delta_{X/S}$ 是同构，因而其余法层平凡。所以由引理
\ref{morphisms-lemma-differentials-diagonal}，有 $\Omega_{X/S} = 0$。其代数版本是
《代数》引理 \ref{algebra-lemma-trivial-differential-surjective}。
\end{proof}

\begin{lemma}
\label{morphisms-lemma-differentials-relative-immersion}
令 $i : Z \to X$ 为 $S$ 上概形的浸入。存在典范正合列
$$
\mathcal{C}_{Z/X} \to i^*\Omega_{X/S} \to \Omega_{Z/S} \to 0
$$
，其中第一个箭头由 $\text{d}_{X/S}$ 诱导，第二个箭头来自引理
\ref{morphisms-lemma-functoriality-differentials}。
\end{lemma}

\begin{proof}
这是《代数》引理 \ref{algebra-lemma-differential-seq} 的层化版本。不过，
我们应当确认第一个箭头可以全局定义。因此在此说明“由
$\text{d}_{X/S}$ 诱导”的含义。具体地，可以假设 $i$ 是闭浸入
（通过缩小 $X$）。令 $\mathcal{I} \subset \mathcal{O}_X$ 为对应于
$Z \subset X$ 的理想层。于是
$\text{d}_{X/S} : \mathcal{I} \to \Omega_{X/S}$ 把子层
$\mathcal{I}^2 \subset \mathcal{I}$ 映入
$\mathcal{I}\Omega_{X/S}$。因此它诱导同态
$\mathcal{I}/\mathcal{I}^2 \to \Omega_{X/S}/\mathcal{I}\Omega_{X/S}$，
该同态是 $\mathcal{O}_X/\mathcal{I}$-线性的。由引理
\ref{morphisms-lemma-i-star-equivalence}，这对应于所需同态
$\mathcal{C}_{Z/X} \to i^*\Omega_{X/S}$。
\end{proof}

\begin{lemma}
\label{morphisms-lemma-differentials-relative-immersion-section}
令 $i : Z \to X$ 为 $S$ 上概形的浸入，并假设 $i$（局部）有左逆。
则引理 \ref{morphisms-lemma-differentials-relative-immersion} 的典范序列
$$
0 \to \mathcal{C}_{Z/X} \to i^*\Omega_{X/S} \to \Omega_{Z/S} \to 0
$$
（局部）分裂正合。特别地，若 $s : S \to X$ 是结构态射
$X \to S$ 的截面，则同态
$\mathcal{C}_{S/X} \to s^*\Omega_{X/S}$（由 $\text{d}_{X/S}$ 诱导）是同构。
\end{lemma}

\begin{proof}
由《代数》引理 \ref{algebra-lemma-differential-seq-split}，当 $X, Z, S$
均仿射时结论成立。一般情形下，由此得出该序列局部分裂正合。若
$g : X \to Z$ 是 $i$ 的左逆，则由《模》引理
\ref{modules-lemma-check-functoriality-differentials}，$i^*c_g$ 是同态
$i^*\Omega_{X/S} \to \Omega_{Z/S}$ 的右逆。于是由《同调》引理
\ref{homology-lemma-ses-split}，该序列分裂。最后，若 $s$ 是截面，
则它是浸入 $s : Z = S \to X$，并且这是 $S$ 上的浸入（见《概形》引理
\ref{schemes-lemma-section-immersion}），此时 $\Omega_{Z/S} = 0$。
\end{proof}

\begin{remark}
\label{morphisms-remark-differentials-diagonal}
令 $X \to S$ 为概形态射。由引理
\ref{morphisms-lemma-differential-product}，有
$$
\Omega_{X \times_S X/S} =
\text{pr}_1^*\Omega_{X/S} \oplus \text{pr}_2^*\Omega_{X/S}
$$
。另一方面，对角态射 $\Delta : X \to X \times_S X$ 是浸入，且局部有左逆。
因此由引理 \ref{morphisms-lemma-differentials-relative-immersion-section}，得到典范短正合列
$$
0 \to \mathcal{C}_{X/X \times_S X} \to \Omega_{X/S} \oplus \Omega_{X/S}
\to \Omega_{X/S} \to 0
$$
。注意右侧箭头是 $(1, 1)$，它确实是分裂满射。另一方面，由引理
\ref{morphisms-lemma-differentials-diagonal}，有恒等式
$\Omega_{X/S} = \mathcal{C}_{X/X \times_S X}$。由于在这一恒等式中选择了
$\text{d}_{X/S}(f) = s_2(f) - s_1(f)$，左侧箭头遂为同态
$(-1, 1)$\footnote{具体地，局部截面
$\text{d}_{X/S}(f) = 1 \otimes f - f \otimes 1$（属于对角态射 $\Delta$ 的理想层）经由
$\text{d}_{X \times_S X/X}$ 映到局部截面
$1 \otimes 1 \otimes 1 \otimes f - 1 \otimes f \otimes 1 \otimes 1
-1 \otimes 1 \otimes f \otimes 1 + f \otimes 1 \otimes 1 \otimes 1 =
\text{pr}_2^*\text{d}_{X/S}(f) - \text{pr}_1^*\text{d}_{X/S}(f)$。}。
\end{remark}

\begin{lemma}
\label{morphisms-lemma-two-immersions}
令
$$
\xymatrix{
Z \ar[r]_i \ar[rd]_j & X \ar[d] \\
& Y
}
$$
为概形的交换图，其中 $i$ 和 $j$ 为浸入。则有典范正合列
$$
\mathcal{C}_{Z/Y} \to
\mathcal{C}_{Z/X} \to
i^*\Omega_{X/Y} \to 0
$$
，其中第一个箭头来自引理
\ref{morphisms-lemma-conormal-functorial}，第二个箭头来自引理
\ref{morphisms-lemma-differentials-relative-immersion}。
\end{lemma}

\begin{proof}
其代数版本是《代数》引理 \ref{algebra-lemma-application-NL}。
\end{proof}










\section{有限阶微分算子}
\label{morphisms-section-differential-operators}

\noindent
我们建议读者参阅交换代数章中的相应一节
（《代数》第 \ref{algebra-section-differential-operators} 节），以及模层章中的
相应一节（《模》第 \ref{modules-section-differential-operators} 节）。

\begin{lemma}
\label{morphisms-lemma-affine-case-differential-operators}
令 $R \to A$ 为环同态。以 $f : X \to S$ 表示相应的仿射概形态射。
令 $\mathcal{F}$ 和 $\mathcal{G}$ 为 $\mathcal{O}_X$-模。若
$\mathcal{F}$ 拟凝聚，则把微分算子送到其在全局截面上的作用所得的映射
$$
\text{Diff}^k_{X/S}(\mathcal{F}, \mathcal{G}) \to
\text{Diff}^k_{A/R}(\Gamma(X, \mathcal{F}), \Gamma(X, \mathcal{G}))
$$
是双射。
\end{lemma}

\begin{proof}
把 $\mathcal{F} = \widetilde{M}$ 写成某个 $A$-模 $M$ 的形式。置
$N = \Gamma(X, \mathcal{G})$。令 $D : M \to N$ 为阶数 $k$ 的微分算子。
我们须证明，存在唯一的微分算子
$\mathcal{F} \to \mathcal{G}$，其阶数为 $k$，并且它在全局截面上给出 $D$。令
$U = D(f) \subset X$ 为标准仿射开集。则 $\mathcal{F}(U) = M_f$ 是局部化。
由《代数》引理 \ref{algebra-lemma-invert-system-differential-operators}，
微分算子 $D$ 唯一延拓为微分算子
$$
D_f : \mathcal{F}(U) = \widetilde{M}(U) = M_f \to N_f = \widetilde{N}(U)
$$
。唯一性表明这些同态 $D_f$ 粘合，给出层同态
$\widetilde{M} \to \widetilde{N}$；这一粘合发生在 $X$ 的所有标准开集构成的基上。因此，由《层》第
\ref{sheaves-section-bases} 节的内容，得到唯一的层同态
$\widetilde{D} : \widetilde{M} \to \widetilde{N}$，它与这些同态相合。
由于 $\widetilde{D}$ 在标准开集上由阶数 $k$ 的微分算子给出，故
$\widetilde{D}$ 是阶数 $k$ 的微分算子（略去一个小细节）。最后，与典范的
$\mathcal{O}_X$-模同态 $c : \widetilde{N} \to \mathcal{G}$ 后复合
（《概形》引理 \ref{schemes-lemma-compare-constructions}），得到
$c \circ \widetilde{D} : \mathcal{F} \to \mathcal{G}$；由《模》引理
\ref{modules-lemma-composition-differential-operators}，这是阶数 $k$
的微分算子。存在性得证。唯一性的证明从略。
\end{proof}

\begin{lemma}
\label{morphisms-lemma-base-change-differential-operators}
令 $a : X \to S$ 与 $b : Y \to S$ 为概形态射。令 $\mathcal{F}$ 和
$\mathcal{F}'$ 为拟凝聚 $\mathcal{O}_X$-模。令
$D : \mathcal{F} \to \mathcal{F}'$ 为阶数 $k$ 的微分算子，定义在 $X/S$ 上。
令 $\mathcal{G}$ 为拟凝聚 $\mathcal{O}_Y$-模。则存在唯一的微分算子
$$
D' :
\text{pr}_1^*\mathcal{F} \otimes_{\mathcal{O}_{X \times_S Y}}
\text{pr}_2^*\mathcal{G}
\longrightarrow
\text{pr}_1^*\mathcal{F}' \otimes_{\mathcal{O}_{X \times_S Y}}
\text{pr}_2^*\mathcal{G}
$$
，它的阶数为 $k$，定义在 $X \times_S Y / Y$ 上，并且
$
D'(s \otimes t) = D(s) \otimes t
$
，其中 $s$ 是 $\mathcal{F}$ 的局部截面，$t$ 是 $\mathcal{G}$ 的局部截面。
\end{lemma}

\begin{proof}
当 $X$、$Y$、$S$ 仿射时，利用引理
\ref{morphisms-lemma-affine-case-differential-operators}，结论由相应的代数结果得出；
见《代数》引理 \ref{algebra-lemma-base-change-differential-operators}。
一般情形下，利用仿射覆盖（例如《概形》引理
\ref{schemes-lemma-affine-covering-fibre-product} 中的覆盖）全局构造 $D'$。
细节从略。
\end{proof}

\begin{remark}
\label{morphisms-remark-base-change-differential-operators}
令 $a : X \to S$ 与 $b : Y \to S$ 为概形态射。以
$p : X \times_S Y \to X$ 和 $q : X \times_S Y \to Y$ 表示投影。
在本注中，给定 $\mathcal{O}_X$-模 $\mathcal{F}$ 和 $\mathcal{O}_Y$-模
$\mathcal{G}$，置
$$
\mathcal{F} \boxtimes \mathcal{G} =
p^*\mathcal{F} \otimes_{\mathcal{O}_{X \times_S Y}} q^*\mathcal{G}
$$
。以 $\mathcal{A}_{X/S}$ 表示加性范畴，其对象为拟凝聚
$\mathcal{O}_X$-模，态射为 $X/S$ 上的有限阶微分算子。
$\mathcal{A}_{Y/S}$ 和 $\mathcal{A}_{X \times_S Y/S}$ 同理。
引理 \ref{morphisms-lemma-base-change-differential-operators} 的构造确定了函子
$$
\boxtimes : 
\mathcal{A}_{X/S} \times \mathcal{A}_{Y/S} \longrightarrow
\mathcal{A}_{X \times_S Y/S},
\quad
(\mathcal{F}, \mathcal{G}) \longmapsto \mathcal{F} \boxtimes \mathcal{G}
$$
，它在态射上双线性。若 $X = \Spec(A)$、$Y = \Spec(B)$ 且
$S = \Spec(R)$，则通过拟凝聚层与模的恒等，此函子在对象上由
$(M, N) \mapsto M \otimes_R N$ 给出，并把态射
$(D, D') : (M, N) \to (M', N')$ 送到
$D \otimes D' : M \otimes_R N \to M' \otimes_R N'$。
\end{remark}










\section{光滑态射}
\label{morphisms-section-smooth}

\noindent
令 $f : X \to Y$ 为拓扑空间的连续映射。考虑下列条件：对每个
$x \in X$，存在开邻域 $x \in U \subset X$、$f(x) \in V \subset Y$，
以及整数 $d$，使得 $f(U) \subset V$，并得到交换图
$$
\xymatrix{
X \ar[d] &
U \ar[l] \ar[d] \ar[r]_-\pi &
V \times \mathbf{R}^d \ar[ld] \\
Y & V \ar[l]
}
$$
，其中 $\pi$ 是到某个开子集上的同胚。概形的光滑态射是概形范畴中
这类映射的类似物。见引理 \ref{morphisms-lemma-smooth-locally-standard-smooth} 和引理
\ref{morphisms-lemma-smooth-etale-over-affine-space}。

\medskip\noindent
与（或许的）预期相反，光滑环同态的概念并非只用微分模定义。回忆：
$R \to A$ 是{\it 光滑环同态}，是指 $A$ 在 $R$ 上有限表示，并且 $A$
在 $R$ 上的朴素余切复形拟同构于置于次数 $0$ 的投射模；见《代数》定义
\ref{algebra-definition-smooth}。

\begin{definition}
\label{morphisms-definition-smooth}
令 $f : X \to S$ 为概形态射。
\begin{enumerate}
\item 称 $f$ 在 $x \in X$ 处{\it 光滑}，若存在仿射开邻域
$\Spec(A) = U \subset X$（它是 $x$ 的邻域）以及仿射开集
$\Spec(R) = V \subset S$，满足 $f(U) \subset V$，且诱导环同态
$R \to A$ 光滑。
\item 称 $f$ {\it 光滑}，若它在 $X$ 的每一点处都光滑。
\item 若仿射概形态射 $f : X \to S$ 同构于
某个标准光滑环同态
\begin{center}
$R \to R[x_1, \ldots, x_n]/(f_1, \ldots, f_c)$
\end{center}
（见《代数》定义
\ref{algebra-definition-standard-smooth}）所给出的态射 $X \to S$，即
$$
\Spec(R[x_1, \ldots, x_n]/(f_1, \ldots, f_c)) \to \Spec(R).
$$
则称它{\it 标准光滑}。
\end{enumerate}
\end{definition}

\noindent
这个定义的一个良好特征是：态射光滑的点所成之集自动为开集。

\medskip\noindent
注意定义中没有分离性或拟紧性的假设。因此，光滑性在源上是局部性质。
精确结果如下。

\begin{lemma}
\label{morphisms-lemma-smooth-characterize}
令 $f : X \to S$ 为概形态射。下列条件等价：
\begin{enumerate}
\item 态射 $f$ 光滑。
\item 对每一对仿射开集 $U \subset X$、$V \subset S$，若
$f(U) \subset V$，则环同态
$\mathcal{O}_S(V) \to \mathcal{O}_X(U)$ 光滑。
\item 存在开覆盖 $S = \bigcup_{j \in J} V_j$，以及开覆盖
$f^{-1}(V_j) = \bigcup_{i \in I_j} U_i$，使得每个态射
$U_i \to V_j$（$j\in J, i\in I_j$）都光滑。
\item 存在仿射开覆盖 $S = \bigcup_{j \in J} V_j$，以及仿射开覆盖
$f^{-1}(V_j) = \bigcup_{i \in I_j} U_i$，使得环同态
$\mathcal{O}_S(V_j) \to \mathcal{O}_X(U_i)$ 对所有
$j\in J, i\in I_j$ 都光滑。
\end{enumerate}
此外，若 $f$ 光滑，则对任意开子概形 $U \subset X$、$V \subset S$，
若 $f(U) \subset V$，限制 $f|_U : U \to V$ 光滑。
\end{lemma}

\begin{proof}
若能证明“$R \to A$ 光滑”是局部性质，则结论由引理
\ref{morphisms-lemma-locally-P} 得出。我们验证定义
\ref{morphisms-definition-property-local} 的条件 (a)、(b)、(c)。由《代数》引理
\ref{algebra-lemma-base-change-smooth}，光滑性在基变换下稳定，故 (a) 成立。
由《代数》引理 \ref{algebra-lemma-compose-smooth}，光滑性在复合下稳定；
并且对任意环 $R$，环同态 $R \to R_f$（标准）光滑，故 (b) 成立。
最后，由《代数》引理 \ref{algebra-lemma-locally-smooth}，性质 (c) 成立。
\end{proof}

\noindent
下面的引理把光滑态射刻画为纤维光滑的平坦有限表示态射。注意，
域上的光滑概形会在《簇》第 \ref{varieties-section-smooth} 节中更详细地讨论。

\begin{lemma}
\label{morphisms-lemma-smooth-flat-smooth-fibres}
令 $f : X \to S$ 为概形态射。若 $f$ 平坦、局部有限表示，且所有纤维
$X_s$ 都光滑，则 $f$ 光滑。
\end{lemma}

\begin{proof}
由《代数》引理 \ref{algebra-lemma-flat-fibre-smooth} 得出。
\end{proof}

\begin{lemma}
\label{morphisms-lemma-composition-smooth}
两个光滑态射的复合光滑。
\end{lemma}

\begin{proof}
在引理 \ref{morphisms-lemma-smooth-characterize} 的证明中，我们已看到光滑性是
环同态的局部性质。因此，本引理的第一项断言由引理
\ref{morphisms-lemma-composition-type-P} 以及光滑环同态性质在复合下稳定这一事实得出；
见《代数》引理 \ref{algebra-lemma-compose-smooth}。
\end{proof}

\begin{lemma}
\label{morphisms-lemma-base-change-smooth}
光滑态射的基变换光滑。
\end{lemma}

\begin{proof}
在引理 \ref{morphisms-lemma-smooth-characterize} 的证明中，我们已看到光滑性是
环同态的局部性质。因此，本引理由引理
\ref{morphisms-lemma-composition-type-P} 以及光滑环同态性质在基变换下稳定这一事实得出；
见《代数》引理 \ref{algebra-lemma-base-change-smooth}。
\end{proof}

\begin{lemma}
\label{morphisms-lemma-open-immersion-smooth}
任意开浸入都光滑。
\end{lemma}

\begin{proof}
这是因为开浸入是局部同构。
\end{proof}

\begin{lemma}
\label{morphisms-lemma-smooth-syntomic}
光滑态射是 syntomic 态射。
\end{lemma}

\begin{proof}
见《代数》引理 \ref{algebra-lemma-smooth-syntomic}。
\end{proof}

\begin{lemma}
\label{morphisms-lemma-smooth-locally-finite-presentation}
光滑态射局部有限表示。
\end{lemma}

\begin{proof}
这是因为按定义，光滑环同态有限表示。
\end{proof}

\begin{lemma}
\label{morphisms-lemma-smooth-flat}
光滑态射平坦。
\end{lemma}

\begin{proof}
合并引理 \ref{morphisms-lemma-syntomic-flat} 与 \ref{morphisms-lemma-smooth-syntomic}。
\end{proof}

\begin{lemma}
\label{morphisms-lemma-smooth-open}
光滑态射泛开。
\end{lemma}

\begin{proof}
合并引理 \ref{morphisms-lemma-smooth-flat}、
\ref{morphisms-lemma-smooth-locally-finite-presentation} 和
\ref{morphisms-lemma-fppf-open}。或者，也可合并引理
\ref{morphisms-lemma-smooth-syntomic} 与
\ref{morphisms-lemma-syntomic-open}。
\end{proof}

\noindent
下面的引理说明任意光滑态射在局部上标准光滑。因此，可用标准光滑态射作为
光滑态射的{\it 局部模型}。

\begin{lemma}
\label{morphisms-lemma-smooth-locally-standard-smooth}
\begin{slogan}
光滑态射局部标准光滑。
\end{slogan}
令 $f : X  \to S$ 为概形态射。令 $x \in X$ 为一点。令
$V \subset S$ 为 $f(x)$ 的仿射开邻域。下列条件等价：
\begin{enumerate}
\item 态射 $f$ 在 $x$ 处光滑。
\item 存在仿射开集 $U \subset X$，满足 $x \in U$、$f(U) \subset V$，
并且诱导态射 $f|_U : U \to V$ 标准光滑。
\end{enumerate}
\end{lemma}

\begin{proof}
由定义以及《代数》引理 \ref{algebra-lemma-standard-smooth} 和
\ref{algebra-lemma-smooth-syntomic} 得出。
\end{proof}

\begin{lemma}
\label{morphisms-lemma-smooth-omega-finite-locally-free}
令 $f : X \to S$ 为概形态射。假设 $f$ 光滑。则微分模
$\Omega_{X/S}$（即 $X$ 在 $S$ 上的微分模）有限局部自由，并且
$$
\text{rank}_x(\Omega_{X/S}) = \dim_x(X_{f(x)})
$$
对每个 $x \in X$ 都成立。
\end{lemma}

\begin{proof}
该断言在 $X$ 和 $S$ 上是局部的。由上述引理
\ref{morphisms-lemma-smooth-locally-standard-smooth}，可以假设 $f$ 是仿射概形的
标准光滑态射。此时结论由《代数》引理
\ref{algebra-lemma-standard-smooth} 得出（以及全局相对完全交的定义；见《代数》
定义 \ref{algebra-definition-relative-global-complete-intersection}）。
\end{proof}

\noindent
引理 \ref{morphisms-lemma-smooth-omega-finite-locally-free} 表明下述定义是有意义的。

\begin{definition}
\label{morphisms-definition-smooth-relative-dimension}
令 $d \geq 0$ 为整数。称概形态射 $f : X \to S$
为{\it 相对维数 $d$ 的光滑态射}，若 $f$ 光滑，且
$\Omega_{X/S}$ 是常秩 $d$ 的有限局部自由模。
\end{definition}

\noindent
换言之，$f$ 光滑，且非空纤维等维、维数为 $d$。由下面的引理
\ref{morphisms-lemma-smooth-at-point}，这也等价于要求：(a) $f$ 局部有限表示，
(b) $f$ 平坦，(c) 所有非空纤维等维且维数为 $d$，以及 (d)
$\Omega_{X/S}$ 是秩 $d$ 的有限局部自由模。仅仅假设 $f$ 平坦、有限表示，
且 $\Omega_{X/S}$ 是秩 $d$ 的有限局部自由模，并不足够。反例为
$\Spec(\mathbf{F}_p[t]) \to \Spec(\mathbf{F}_p[t^p])$。

\medskip\noindent
下面给出一点处光滑性的微分判据。这一结果有许多变体，都可能在某处有用。
我们会按需在这里加入它们。

\begin{lemma}
\label{morphisms-lemma-smooth-at-point}
令 $f : X \to S$ 为概形态射。令 $x \in X$。置 $s = f(x)$。
假设 $f$ 局部有限表示。下列条件等价：
\begin{enumerate}
\item 态射 $f$ 在 $x$ 处光滑。
\item 局部环同态 $\mathcal{O}_{S, s} \to \mathcal{O}_{X, x}$ 平坦，且
$X_s \to \Spec(\kappa(s))$ 在 $x$ 处光滑。
\item 局部环同态 $\mathcal{O}_{S, s} \to \mathcal{O}_{X, x}$ 平坦，且
$\mathcal{O}_{X, x}$-模 $\Omega_{X/S, x}$ 可由至多
$\dim_x(X_{f(x)})$ 个元素生成。
\item 局部环同态 $\mathcal{O}_{S, s} \to \mathcal{O}_{X, x}$ 平坦，且
$\kappa(x)$-向量空间
$$
\Omega_{X_s/s, x} \otimes_{\mathcal{O}_{X_s, x}} \kappa(x) =
\Omega_{X/S, x} \otimes_{\mathcal{O}_{X, x}} \kappa(x)
$$
可由至多 $\dim_x(X_{f(x)})$ 个元素生成。
\item 存在仿射开集 $U \subset X$ 与 $V \subset S$，使得
$x \in U$、$f(U) \subset V$，并且诱导态射
$f|_U : U \to V$ 标准光滑。
\item 存在仿射开集 $\Spec(A) = U \subset X$ 和
$\Spec(R) = V \subset S$，其中 $x \in U$ 对应于
$\mathfrak q \subset A$，且 $f(U) \subset V$，并存在表示
$$
A = R[x_1, \ldots, x_n]/(f_1, \ldots, f_c)
$$
使得
$$
g =
\det
\left(
\begin{matrix}
\partial f_1/\partial x_1 &
\partial f_2/\partial x_1 &
\ldots &
\partial f_c/\partial x_1 \\
\partial f_1/\partial x_2 &
\partial f_2/\partial x_2 &
\ldots &
\partial f_c/\partial x_2 \\
\ldots & \ldots & \ldots & \ldots \\
\partial f_1/\partial x_c &
\partial f_2/\partial x_c &
\ldots &
\partial f_c/\partial x_c
\end{matrix}
\right)
$$
在 $A$ 中的像不属于 $\mathfrak q$。
\end{enumerate}
\end{lemma}

\begin{proof}
注意，若 $f$ 在 $x$ 处光滑，则由引理
\ref{morphisms-lemma-smooth-locally-standard-smooth} 可知 (5) 成立，而 (6) 是 (5)
略微减弱的版本。此外，$f$ 光滑蕴含局部环同态
$\mathcal{O}_{S, s} \to \mathcal{O}_{X, x}$ 平坦（见引理
\ref{morphisms-lemma-smooth-flat}），并且 $\Omega_{X/S}$ 有限局部自由，其秩等于
$\dim_x(X_s)$（见引理 \ref{morphisms-lemma-smooth-omega-finite-locally-free}）。
因此 (1) 蕴含 (3) 和 (4)。由引理 \ref{morphisms-lemma-base-change-smooth}，
还可知 (1) 蕴含 (2)。

\medskip\noindent
由引理 \ref{morphisms-lemma-base-change-differentials}，微分模
$\Omega_{X_s/s}$（即纤维 $X_s$ 在 $\kappa(s)$ 上的微分模）是微分模
$\Omega_{X/S}$（即 $X$ 在 $S$ 上的微分模）的拉回。这给出本引理第 (4) 项中的等式。由引理
\ref{morphisms-lemma-finite-type-differentials}，这些模有限型。因此，模
$\Omega_{X/S, x}$ 与 $\Omega_{X_s/s, x}$ 的最小生成元数相同，并且由
Nakayama 引理（《代数》引理 \ref{algebra-lemma-NAK}）等于这个
$\kappa(x)$-向量空间的维数。这尤其表明 (3) 与 (4) 等价。

\medskip\noindent
《代数》引理 \ref{algebra-lemma-flat-fibre-smooth} 表明 (2) 蕴含 (1)。
《代数》引理 \ref{algebra-lemma-characterize-smooth-over-field} 表明 (3)
和 (4) 蕴含 (2)。最后，(6) 蕴含 (5)，例如见《代数》例
\ref{algebra-example-make-standard-smooth}；而由《代数》引理
\ref{algebra-lemma-standard-smooth}，(5) 蕴含 (1)。
\end{proof}

\begin{lemma}
\label{morphisms-lemma-set-points-where-fibres-smooth}
令
$$
\xymatrix{
X' \ar[r]_{g'} \ar[d]_{f'} & X \ar[d]^f \\
S' \ar[r]^g & S
}
$$
为概形的笛卡尔图。令 $W \subset X$ 和 $W' \subset X'$ 分别为
$f$ 与 $f'$ 光滑的点所成的开子概形。若下列任一条件成立，则
$W' = (g')^{-1}(W)$：
\begin{enumerate}
\item $f$ 平坦且局部有限表示；或
\item $f$ 局部有限表示且 $g$ 平坦。
\end{enumerate}
\end{lemma}

\begin{proof}
先假设 $f$ 局部有限型。考虑集合
$$
T = \{x \in X \mid X_{f(x)}\text{ is smooth over }\kappa(f(x))\text{ 在 }x\}
$$
以及集合 $T' \subset X'$（对应于 $f'$）。我们断言
$T' = (g')^{-1}(T)$。事实上，令 $s' \in S'$ 为一点，并令
$s = g(s')$。则有
$$
X'_{s'} =
\Spec(\kappa(s')) \times_{\Spec(\kappa(s))} X_s
$$
。换言之，基变换的纤维是纤维的基变换。因此，该断言等价于《代数》引理
\ref{algebra-lemma-smooth-field-change-local}。

\medskip\noindent
这样，情形 (1) 成立，因为在情形 (1) 中，由引理
\ref{morphisms-lemma-smooth-at-point}，$T$ 正是 $f$ 光滑的点所成的（开）集。

\medskip\noindent
在情形 (2) 中，令 $x' \in W'$。则 $g'$ 在 $x'$ 处平坦（引理
\ref{morphisms-lemma-base-change-module-flat}），并且 $g \circ f'$ 在 $x'$ 处平坦
（引理 \ref{morphisms-lemma-composition-module-flat}）。因此，由引理
\ref{morphisms-lemma-flat-permanence}，$f$ 在 $x = g'(x')$ 处平坦。另一方面，
由于 $x' \in T'$（引理 \ref{morphisms-lemma-base-change-smooth}），可知
$x \in T$。所以，由引理 \ref{morphisms-lemma-smooth-at-point}，$f$ 在 $x$ 处光滑。
\end{proof}

\noindent
下面的引理确实用到了光滑环同态的朴素余切复形之 $H^{-1}$ 消失。

\begin{lemma}
\label{morphisms-lemma-triangle-differentials-smooth}
令 $f : X \to Y$、$g : Y \to S$ 为概形态射。假设 $f$ 光滑。则
$$
0 \to f^*\Omega_{Y/S} \to \Omega_{X/S} \to \Omega_{X/Y} \to 0
$$
（见引理 \ref{morphisms-lemma-triangle-differentials}）短正合。
\end{lemma}

\begin{proof}
本引理的代数版本如下：给定环同态 $A \to B \to C$，其中
$B \to C$ 光滑，则《代数》引理
\ref{algebra-lemma-exact-sequence-differentials} 的序列
$$
0 \to C \otimes_B \Omega_{B/A} \to \Omega_{C/A} \to \Omega_{C/B} \to 0
$$
正合。这就是《代数》引理
\ref{algebra-lemma-triangle-differentials-smooth}。
\end{proof}

\begin{lemma}
\label{morphisms-lemma-differentials-relative-immersion-smooth}
令 $i : Z \to X$ 为 $S$ 上概形的浸入。假设 $Z$ 在 $S$ 上光滑。
则引理 \ref{morphisms-lemma-differentials-relative-immersion} 的典范正合列
$$
0 \to \mathcal{C}_{Z/X} \to i^*\Omega_{X/S} \to \Omega_{Z/S} \to 0
$$
短正合。
\end{lemma}

\begin{proof}
本引理的代数版本如下：给定环同态 $A \to B \to C$，其中
$A \to C$ 光滑，且 $B \to C$ 满、核为 $J$，则《代数》引理
\ref{algebra-lemma-differential-seq} 的序列
$$
0 \to J/J^2 \to C \otimes_B \Omega_{B/A} \to \Omega_{C/A} \to 0
$$
正合。这就是《代数》引理
\ref{algebra-lemma-differential-seq-smooth}。
\end{proof}

\begin{lemma}
\label{morphisms-lemma-two-immersions-smooth}
令
$$
\xymatrix{
Z \ar[r]_i \ar[rd]_j & X \ar[d] \\
& Y
}
$$
为概形的交换图，其中 $i$ 与 $j$ 是浸入，并且 $X \to Y$ 光滑。
则引理 \ref{morphisms-lemma-two-immersions} 的典范正合列
$$
0 \to  \mathcal{C}_{Z/Y} \to \mathcal{C}_{Z/X} \to i^*\Omega_{X/Y} \to 0
$$
正合。
\end{lemma}

\begin{proof}
本引理的代数版本如下：给定环同态 $A \to B \to C$，其中
$A \to C$ 满，且 $A \to B$ 光滑，则《代数》引理
\ref{algebra-lemma-application-NL} 的序列
$$
0 \to I/I^2 \to J/J^2 \to C \otimes_B \Omega_{B/A} \to 0
$$
正合。这就是《代数》引理 \ref{algebra-lemma-application-NL-smooth}。
\end{proof}

\begin{lemma}
\label{morphisms-lemma-smooth-permanence}
令
$$
\xymatrix{
X \ar[rr]_f \ar[rd]_p & &
Y \ar[dl]^q \\
& S
}
$$
为概形态射的交换图。假设：
\begin{enumerate}
\item $f$ 满且光滑；
\item $p$ 光滑；并且
\item $q$ 局部有限表示\footnote{事实上，这由 (1) 和 (2) 蕴含；见《下降》
引理 \ref{descent-lemma-flat-finitely-presented-permanence}。此外，只需假设
$f$ 满、平坦且局部有限表示；见《下降》引理
\ref{descent-lemma-smooth-permanence}。}。
\end{enumerate}
则 $q$ 光滑。
\end{lemma}

\begin{proof}
由引理 \ref{morphisms-lemma-flat-permanence}，可知 $q$ 平坦。取点 $y \in Y$。
取一点 $x \in X$，它映到 $y$。设 $f$ 的相对维数为 $a$（在 $x$ 处），且
$p$ 的相对维数为 $b$（在 $x$ 处）。由引理
\ref{morphisms-lemma-smooth-omega-finite-locally-free}，这意味着
$\Omega_{X/S, x}$ 为秩 $b$ 的自由模，且 $\Omega_{X/Y, x}$ 为秩 $a$
的自由模。由引理 \ref{morphisms-lemma-triangle-differentials-smooth} 的短正合列，
这意味着 $(f^*\Omega_{Y/S})_x$ 为秩 $b - a$ 的自由模。由 Nakayama
引理，这蕴含 $\Omega_{Y/S, y}$ 可由 $b - a$ 个元素生成。此外，
由引理 \ref{morphisms-lemma-dimension-fibre-at-a-point-additive}，有
$\dim_y(Y_s) = b - a$。所以，由引理 \ref{morphisms-lemma-smooth-at-point} 第 (2) 项，
可得 $Y \to S$ 在 $y$ 处光滑。
\end{proof}

\noindent
在下面引理的情形中，$\sigma$ 的像在 $X$ 上局部由正则序列切出；见《除子》
引理 \ref{divisors-lemma-section-smooth-regular-immersion}。

\begin{lemma}
\label{morphisms-lemma-section-smooth-morphism}
令 $f : X \to S$ 为概形态射。令 $\sigma : S \to X$ 为 $f$ 的截面。
令 $s \in S$ 为一点，使得 $f$ 在 $x = \sigma(s)$ 处光滑。则存在
仿射开邻域 $\Spec(A) = U \subset S$（属于 $s$）和
$\Spec(B) = V \subset X$（属于 $x$），使得：
\begin{enumerate}
\item $f(V) \subset U$ 且 $\sigma(U) \subset V$；
\item 对于 $I = \Ker(\sigma^\# : B \to A)$，模 $I/I^2$ 是自由 $A$-模；并且
\item 有 $B^\wedge \cong A[[x_1, \ldots, x_d]]$（作为 $A$-代数），其中
$B^\wedge$ 表示 $B$ 关于 $I$ 的完备化。
\end{enumerate}
\end{lemma}

\begin{proof}
取仿射开集 $U \subset S$，它包含 $s$。取仿射开集
$V \subset f^{-1}(U)$，它包含 $x$。取仿射开集
$U' \subset \sigma^{-1}(V)$，它包含 $s$。注意，$V' = f^{-1}(U') \cap V$ 仿射，
因为它等于纤维积 $V' = U' \times_U V$。于是 $U'$ 与 $V'$ 满足 (1)。
写 $U' = \Spec(A')$ 且 $V' = \Spec(B')$。由《代数》引理
\ref{algebra-lemma-section-smooth}，模 $I'/(I')^2$ 是有限局部自由
$A'$-模。因此，把 $U'$ 替换为较小的仿射开集 $U'' \subset U'$，并把
$V'$ 替换为 $V'' = V' \cap f^{-1}(U'')$ 后，得到
$I''/(I'')^2$ 自由的情形，即 (2) 成立。在此情形下，由《代数》引理
\ref{algebra-lemma-section-smooth}，(3) 也成立。
\end{proof}

\noindent
概形 $X$ 在一点 $x$ 处的维数（《性质》定义
\ref{properties-definition-dimension}），就是把 $X$ 视为拓扑空间时在
$x$ 处的维数；见《拓扑》定义 \ref{topology-definition-Krull}。一般而言，
这不是局部环 $\mathcal{O}_{X,x}$ 的维数。

\begin{lemma}
\label{morphisms-lemma-smoothness-dimension}
令 $f : X \to Y$ 为局部 Noether 概形之间的光滑态射。对每个点
$x$（位于 $X$），若其像 $y$ 位于 $Y$，则
$$
\dim_x(X) = \dim_y(Y) + \dim_x(X_y),
$$
其中 $X_y$ 表示 $y$ 上的纤维。
\end{lemma}

\begin{proof}
把 $X$ 替换为 $x$ 的开邻域后，存在自然数 $d$，使得
$X \to Y$ 的所有纤维在每一点处的维数都是 $d$；见引理
\ref{morphisms-lemma-smooth-omega-finite-locally-free}。于是 $f$ 平坦（引理
\ref{morphisms-lemma-smooth-flat}）、局部有限型（引理
\ref{morphisms-lemma-smooth-locally-finite-presentation}），且相对维数为 $d$。
因此结论由引理 \ref{morphisms-lemma-rel-dimension-dimension} 得出。
\end{proof}

\section{非分歧态射}
\label{morphisms-section-unramified}

\noindent
在讨论（也许）更有趣的平展态射这一类之前，我们先简要讨论非分歧态射。
回顾：若环同态 $R \to A$ 是有限型的并且 $\Omega_{A/R} = 0$，
则称它是{\it 非分歧的}（这是 \cite{Henselian} 的定义）。
若环同态 $R \to A$ 是有限表示的并且 $\Omega_{A/R} = 0$，
则称它是 {\it G-非分歧的}（这是 \cite{EGA} 的定义）。见
代数，引理 \ref{algebra-definition-unramified}。

\begin{definition}
\label{morphisms-definition-unramified}
设 $f : X \to S$ 是概形的态射。
\begin{enumerate}
\item 称 $f$ {\it 在 $x \in X$ 处非分歧}，如果存在仿射开邻域
$\Spec(A) = U \subset X$ 包含 $x$，以及仿射开集
$\Spec(R) = V \subset S$，满足 $f(U) \subset V$，并且所诱导的
环同态 $R \to A$ 非分歧。
\item 称 $f$ {\it 在 $x \in X$ 处 G-非分歧}，如果存在仿射开邻域
$\Spec(A) = U \subset X$ 包含 $x$，以及仿射开集
$\Spec(R) = V \subset S$，满足 $f(U) \subset V$，并且所诱导的
环同态 $R \to A$ 是 G-非分歧的。
\item 若 $f$ 在 $X$ 的每一点处非分歧，则称其{\it 非分歧}。
\item 若 $f$ 在 $X$ 的每一点处 G-非分歧，则称其
{\it G-非分歧}。
\end{enumerate}
\end{definition}

\noindent
注意，G-非分歧态射是非分歧态射。因此，关于非分歧态射的任何结果
都会蕴含关于 G-非分歧态射的相应结果。此外，若 $S$ 局部 Noether，
则 G-非分歧态射与非分歧态射没有区别；见引理
\ref{morphisms-lemma-noetherian-unramified}。这个定义的一个可取之处是：
一个态射非分歧（相应地，G-非分歧）的点集自动是开集。

\begin{lemma}
\label{morphisms-lemma-unramified-omega-zero}
设 $f : X \to S$ 是概形的态射。那么
\begin{enumerate}
\item $f$ 非分歧，当且仅当 $f$ 局部为有限型且
$\Omega_{X/S} = 0$；并且
\item $f$ 是 G-非分歧的，当且仅当 $f$ 局部为有限表示且
$\Omega_{X/S} = 0$。
\end{enumerate}
\end{lemma}

\begin{proof}
根据定义，环同态 $R \to A$ 非分歧（相应地，G-非分歧），当且仅当
它是有限型的（相应地，有限表示的）并且 $\Omega_{A/R} = 0$。
因此，该引理直接由定义和引理 \ref{morphisms-lemma-differentials-affine} 得出。
\end{proof}

\noindent
注意，定义中没有分离性或拟紧性假设。因此，非分歧性在源上具有
局部性质。下面是精确的结果。

\begin{lemma}
\label{morphisms-lemma-unramified-characterize}
设 $f : X \to S$ 是概形的态射。下列条件等价：
\begin{enumerate}
\item 态射 $f$ 非分歧（相应地，G-非分歧）。
\item 对每个仿射开集 $U \subset X$、$V \subset S$，若
$f(U) \subset V$，则环同态
$\mathcal{O}_S(V) \to \mathcal{O}_X(U)$ 非分歧
（相应地，G-非分歧）。
\item 存在开覆盖 $S = \bigcup_{j \in J} V_j$ 以及开覆盖
$f^{-1}(V_j) = \bigcup_{i \in I_j} U_i$，使得每个态射
$U_i \to V_j$（$j\in J, i\in I_j$）都非分歧
（相应地，G-非分歧）。
\item 存在仿射开覆盖 $S = \bigcup_{j \in J} V_j$ 以及仿射开覆盖
$f^{-1}(V_j) = \bigcup_{i \in I_j} U_i$，使得环同态
$\mathcal{O}_S(V_j) \to \mathcal{O}_X(U_i)$ 对所有
$j\in J, i\in I_j$ 都非分歧（相应地，G-非分歧）。
\end{enumerate}
此外，若 $f$ 非分歧（相应地，G-非分歧），则对任意开子概形
$U \subset X$、$V \subset S$，只要 $f(U) \subset V$，其限制
$f|_U : U \to V$ 就非分歧（相应地，G-非分歧）。
\end{lemma}

\begin{proof}
只要证明性质“$R \to A$ 非分歧”是局部性质，结论便由引理
\ref{morphisms-lemma-locally-P} 得出。我们检验定义
\ref{morphisms-definition-property-local} 的条件 (a)、(b) 和 (c)。
这些性质已在代数，引理 \ref{algebra-lemma-unramified} 中证明。
\end{proof}

\begin{lemma}
\label{morphisms-lemma-composition-unramified}
两个非分歧态射的复合仍非分歧。对 G-非分歧态射也有同样的结论。
\end{lemma}

\begin{proof}
引理 \ref{morphisms-lemma-unramified-characterize} 的证明表明，
非分歧（相应地，G-非分歧）是环同态的局部性质。
因而第一项断言由引理 \ref{morphisms-lemma-composition-type-P}
以及如下事实得出：非分歧（相应地，G-非分歧）是一个在复合下
稳定的环同态性质；见代数，引理 \ref{algebra-lemma-unramified}。
\end{proof}

\begin{lemma}
\label{morphisms-lemma-base-change-unramified}
非分歧态射的基变换仍非分歧。对 G-非分歧态射也有同样的结论。
\end{lemma}

\begin{proof}
引理 \ref{morphisms-lemma-unramified-characterize} 的证明表明，
非分歧（相应地，G-非分歧）是环同态的局部性质。
因而该引理由引理 \ref{morphisms-lemma-base-change-type-P} 以及如下事实得出：
非分歧（相应地，G-非分歧）是一个在基变换下稳定的环同态性质；
见代数，引理 \ref{algebra-lemma-unramified}。
\end{proof}

\begin{lemma}
\label{morphisms-lemma-noetherian-unramified}
设 $f : X \to S$ 是概形的态射。假设 $S$ 局部 Noether。
那么 $f$ 非分歧，当且仅当 $f$ 是 G-非分歧的。
\end{lemma}

\begin{proof}
这由定义和引理
\ref{morphisms-lemma-noetherian-finite-type-finite-presentation} 得出。
\end{proof}


\begin{lemma}
\label{morphisms-lemma-open-immersion-unramified}
任何开浸入都是 G-非分歧的。
\end{lemma}

\begin{proof}
这是因为开浸入是局部同构。
\end{proof}

\begin{lemma}
\label{morphisms-lemma-closed-immersion-unramified}
闭浸入 $i : Z \to X$ 是非分歧的。它是 G-非分歧的，当且仅当
相应的拟凝聚理想层
$\mathcal{I} = \Ker(\mathcal{O}_X \to i_*\mathcal{O}_Z)$
（作为 $\mathcal{O}_X$-模）是有限型的。
\end{lemma}

\begin{proof}
这由引理 \ref{morphisms-lemma-closed-immersion-finite-presentation} 和
代数，引理 \ref{algebra-lemma-unramified} 得出。
\end{proof}

\begin{lemma}
\label{morphisms-lemma-unramified-locally-finite-type}
非分歧态射局部为有限型。G-非分歧态射局部为有限表示。
\end{lemma}

\begin{proof}
根据定义，非分歧环同态是有限型的。
根据定义，G-非分歧环同态是有限表示的。
\end{proof}

\begin{lemma}
\label{morphisms-lemma-unramified-quasi-finite}
设 $f : X \to S$ 是概形的态射。若 $f$ 在 $x$ 处非分歧，
则 $f$ 在 $x$ 处拟有限。特别地，非分歧态射局部拟有限。
\end{lemma}

\begin{proof}
见代数，引理 \ref{algebra-lemma-unramified-quasi-finite}。
\end{proof}

\begin{lemma}
\label{morphisms-lemma-unramified-over-field}
非分歧态射的纤维具有下列性质。
\begin{enumerate}
\item 设 $X$ 是域 $k$ 上的概形。结构态射
$X \to \Spec(k)$ 非分歧，当且仅当 $X$ 是 $k$ 的有限可分
域扩张的谱的不交并。
\item 若 $f : X \to S$ 是非分歧态射，则对每个 $s \in S$，
纤维 $X_s$ 都是 $\kappa(s)$ 的有限可分域扩张的谱的不交并。
\end{enumerate}
\end{lemma}

\begin{proof}
第 (2) 项由第 (1) 项和引理
\ref{morphisms-lemma-base-change-unramified} 得出。我们来证明第 (1) 项。
先使用代数，引理 \ref{algebra-lemma-characterize-unramified}。
该引理说明，若 $X$ 是 $k$ 的有限可分域扩张的谱的不交并，
则 $X \to \Spec(k)$ 非分歧。反过来，假设
$X \to \Spec(k)$ 非分歧。由代数，引理
\ref{algebra-lemma-unramified-at-prime}，对每个 $x \in X$，
剩余域扩张 $\kappa(x)/k$ 都是有限可分的。由于
$X \to \Spec(k)$ 局部拟有限（引理
\ref{morphisms-lemma-unramified-quasi-finite}），由引理
\ref{morphisms-lemma-quasi-finite-at-point-characterize} 可知，$X$ 的所有点
都是孤立闭点。因此 $X$ 是离散空间，特别地，它是其局部环的谱的
不交并。再次使用代数，引理
\ref{algebra-lemma-unramified-at-prime}，这些局部环都是域，证毕。
\end{proof}

\noindent
下面的引理将非分歧态射刻画为具有非分歧纤维的局部有限型态射。

\begin{lemma}
\label{morphisms-lemma-unramified-etale-fibres}
设 $f : X \to S$ 是概形的态射。
\begin{enumerate}
\item 若 $f$ 非分歧，则对任意 $x \in X$，域扩张
$\kappa(x)/\kappa(f(x))$ 是有限可分的。
\item 若 $f$ 局部为有限型，并且对每个 $s \in S$，纤维
$X_s$ 都是 $\kappa(s)$ 的有限可分域扩张的谱的不交并，
则 $f$ 非分歧。
\item 若 $f$ 局部为有限表示，并且对每个 $s \in S$，纤维
$X_s$ 都是 $\kappa(s)$ 的有限可分域扩张的谱的不交并，
则 $f$ 是 G-非分歧的。
\end{enumerate}
\end{lemma}

\begin{proof}
这由代数中的引理
\ref{algebra-lemma-unramified-at-prime} 和
\ref{algebra-lemma-characterize-unramified} 得出。
\end{proof}

\noindent
下面用对角态射给出非分歧态射的一种刻画。

\begin{lemma}
\label{morphisms-lemma-diagonal-unramified-morphism}
设 $f : X \to S$ 是一个态射。
\begin{enumerate}
\item 若 $f$ 非分歧，则对角态射
$\Delta : X \to X \times_S X$ 是开浸入。
\item 若 $f$ 局部为有限型且 $\Delta$ 是开浸入，则 $f$ 非分歧。
\item 若 $f$ 局部为有限表示且 $\Delta$ 是开浸入，则 $f$
是 G-非分歧的。
\end{enumerate}
\end{lemma}

\begin{proof}
第一项断言由代数，引理
\ref{algebra-lemma-diagonal-unramified} 得出。
第二项断言来自如下事实：$\Omega_{X/S}$ 是对角态射的余法层
（引理 \ref{morphisms-lemma-differentials-diagonal}），因而若
$\Delta$ 是开浸入，它显然为零。
\end{proof}

\begin{lemma}
\label{morphisms-lemma-unramified-at-point}
设 $f : X \to S$ 是概形的态射。设 $x \in X$。
令 $s = f(x)$。假设 $f$ 局部为有限型
（相应地，局部为有限表示）。下列条件等价：
\begin{enumerate}
\item 态射 $f$ 在 $x$ 处非分歧（相应地，G-非分歧）。
\item 纤维 $X_s$ 作为 $\kappa(s)$ 上的概形在 $x$ 处非分歧。
\item $\mathcal{O}_{X, x}$-模 $\Omega_{X/S, x}$ 为零。
\item $\mathcal{O}_{X_s, x}$-模 $\Omega_{X_s/s, x}$ 为零。
\item $\kappa(x)$-向量空间
$$
\Omega_{X_s/s, x} \otimes_{\mathcal{O}_{X_s, x}} \kappa(x) =
\Omega_{X/S, x} \otimes_{\mathcal{O}_{X, x}} \kappa(x)
$$
为零。
\item 有 $\mathfrak m_s\mathcal{O}_{X, x} = \mathfrak m_x$，
并且域扩张 $\kappa(x)/\kappa(s)$ 是有限可分的。
\end{enumerate}
\end{lemma}

\begin{proof}
注意，若 $f$ 在 $x$ 处非分歧，则由定义以及
第 \ref{morphisms-section-sheaf-differentials} 节关于微分模的结果可知，
$\Omega_{X/S} = 0$ 在 $x$ 的某个邻域中成立。因此 (1) 蕴含
(3)，也蕴含 (5) 中右侧向量空间的消失。它还蕴含 (2)，因为由
引理 \ref{morphisms-lemma-base-change-differentials}，微分模
$\Omega_{X_s/s}$（即纤维 $X_s$ 在 $\kappa(s)$ 上的微分模）
是微分模 $\Omega_{X/S}$（即 $X$ 在 $S$ 上的微分模）的拉回。
这个关于微分模的事实还蕴含
第 (4) 项中所展示的向量空间等式。由引理
\ref{morphisms-lemma-finite-type-differentials}，模
$\Omega_{X/S, x}$ 和 $\Omega_{X_s/s, x}$ 都是有限型的。
因此，由 Nakayama 引理（代数，引理
\ref{algebra-lemma-NAK}），模
$\Omega_{X/S, x}$ 和 $\Omega_{X_s/s, x}$ 为零，当且仅当
(4) 中相应的 $\kappa(x)$-向量空间为零。
特别地，这表明 (3)、(4) 和 (5) 等价。$\Omega_{X/S}$ 的支集
在 $X$ 中闭；见模，引理
\ref{modules-lemma-support-finite-type-closed}。假设 (3) 蕴含
$x$ 不在该支集中。因此 $\Omega_{X/S}$ 在 $x$ 的某个邻域中为零，
这蕴含 (1)。将 (1) 与 (3) 的等价性应用于 $X_s \to s$，
便得到 (2) 与 (4) 的等价性。至此，我们已经看到
(1) -- (5) 等价。

\medskip\noindent
或者，可以使用代数，引理 \ref{algebra-lemma-unramified}
更直接地看出 (1) -- (5) 的等价性。

\medskip\noindent
(1) 与 (6) 的等价性由引理
\ref{morphisms-lemma-unramified-etale-fibres} 得出。也可更直接地由
代数中的引理 \ref{algebra-lemma-unramified-at-prime} 和
\ref{algebra-lemma-characterize-unramified} 得出。
\end{proof}

\begin{lemma}
\label{morphisms-lemma-set-points-where-fibres-unramified}
设 $f : X \to S$ 是概形的态射。假设 $f$ 局部为有限型。
下列开集的形成
\begin{align*}
T
& =
\{x \in X \mid X_{f(x)}\text{ is unramified over }\kappa(f(x))\text{ 在 }x\} \\
& =
\{x \in X \mid X\text{ is unramified over }S\text{ 在 }x\}
\end{align*}
与任意基变换交换：对任意态射 $g : S' \to S$，考虑基变换
$f' : X' \to S'$（它是 $f$ 的基变换）以及投影
$g' : X' \to X$。则相应集合 $T'$（对于态射 $f'$）等于
$T' = (g')^{-1}(T)$。
若假设 $f$ 局部为有限表示，则对 $f$ 为 G-非分歧的点组成的
开集也有同样的结论。
\end{lemma}

\begin{proof}
设 $s' \in S'$ 是一点，并令 $s = g(s')$。于是
$$
X'_{s'} =
\Spec(\kappa(s')) \times_{\Spec(\kappa(s))} X_s
$$
换言之，基变换的纤维是纤维的基变换。特别地，
$$
\Omega_{X_s/s, x} \otimes_{\mathcal{O}_{X_s, x}} \kappa(x')
=
\Omega_{X'_{s'}/s', x'} \otimes_{\mathcal{O}_{X'_{s'}, x'}} \kappa(x')
$$
见引理 \ref{morphisms-lemma-base-change-differentials}。于是由引理
\ref{morphisms-lemma-unramified-at-point}，$x' \in T'$ 当且仅当
$x \in T$。第二项断言由第一项得出，因为在这种情况下，
根据引理 \ref{morphisms-lemma-unramified-at-point}，$T$ 正是 $f$
为 G-非分歧的点组成的（开）集。
\end{proof}

\begin{lemma}
\label{morphisms-lemma-unramified-permanence}
设 $f : X \to Y$ 是 $S$ 上概形的态射。
\begin{enumerate}
\item 若 $X$ 在 $S$ 上非分歧，则 $f$ 非分歧。
\item 若 $X$ 在 $S$ 上 G-非分歧，且 $Y$ 在 $S$ 上局部为有限型，
则 $f$ 是 G-非分歧的。
\end{enumerate}
\end{lemma}

\begin{proof}
假设 $X$ 在 $S$ 上非分歧。由引理
\ref{morphisms-lemma-permanence-finite-type} 可知 $f$ 局部为有限型。
根据假设有 $\Omega_{X/S} = 0$。因而由引理
\ref{morphisms-lemma-triangle-differentials}，$\Omega_{X/Y} = 0$。
所以 $f$ 非分歧。若 $X$ 在 $S$ 上 G-非分歧且 $Y$ 在 $S$
上局部为有限型，则由引理
\ref{morphisms-lemma-finite-presentation-permanence} 可知 $f$
局部为有限表示，从而 $f$ 是 G-非分歧的。
\end{proof}

\begin{lemma}
\label{morphisms-lemma-value-at-one-point}
设 $S$ 为概形。设 $X$、$Y$ 是 $S$ 上的概形。
设 $f, g : X \to Y$ 是 $S$ 上的态射。设 $x \in X$。假设
\begin{enumerate}
\item 结构态射 $Y \to S$ 非分歧；
\item 有 $f(x) = g(x)$（在 $Y$ 中），记 $y = f(x) = g(x)$；并且
\item 所诱导的映射 $f^\sharp, g^\sharp : \kappa(y) \to \kappa(x)$
相等。
\end{enumerate}
则 $x$ 在 $X$ 中有一个开邻域，使得 $f$ 与 $g$ 在其上相等。
\end{lemma}

\begin{proof}
考虑态射 $(f, g) : X \to Y \times_S Y$。由假设 (1) 和引理
\ref{morphisms-lemma-diagonal-unramified-morphism}，
$\Delta_{Y/S}(Y)$ 的逆像在 $X$ 中是开集。
而假设 (2) 和 (3) 蕴含 $x$ 位于这个开子集中。
\end{proof}

\section{平展态射}
\label{morphisms-section-etale}

\noindent
概形的 Zariski 拓扑是一种非常粗糙的拓扑。考察
$\mathbf{C}$ 上的簇时，这一点尤其明显。事实表明，把平展态射
规定为拓扑学中局部同构的类似物，会引入一种细得多的拓扑。
在 $\mathbf{C}$ 上的簇上，用有限系数计算上同调时，这种拓扑
给出“正确的”Betti 数。另一个可观察的现象是：若
$f : X \to Y$ 是 $\mathbf{C}$ 上簇的平展态射，而 $x$
是 $X$ 的闭点，则 $f$ 诱导完备局部环的同构
$\mathcal{O}^{\wedge}_{Y, f(x)} \to \mathcal{O}^{\wedge}_{X, x}$。

\medskip\noindent
本节开始研究这些问题。事实上，我们有意把讨论限制在最低限度，
因为更有趣的结果将在别处讨论。回顾：若环同态 $R \to A$
光滑且 $\Omega_{A/R} = 0$，则称其{\it 平展}；见
代数，定义 \ref{algebra-definition-etale}。

\begin{definition}
\label{morphisms-definition-etale}
设 $f : X \to S$ 是概形的态射。
\begin{enumerate}
\item 称 $f$ {\it 在 $x \in X$ 处平展}，如果存在仿射开邻域
$\Spec(A) = U \subset X$ 包含 $x$，以及仿射开集
$\Spec(R) = V \subset S$，满足 $f(U) \subset V$，并且所诱导的
环同态 $R \to A$ 平展。
\item 若 $f$ 在 $X$ 的每一点处平展，则称其{\it 平展}。
\item 仿射概形的态射 $f : X \to S$ 称为{\it 标准平展的}，
如果 $X \to S$ 同构于
$$
\Spec(R[x]_h/(g)) \to \Spec(R)
$$
其中 $R \to R[x]_h/(g)$ 是标准平展环同态；见
代数，定义 \ref{algebra-definition-standard-etale}；也就是说，
$g$ 是首一的，并且 $g'$ 在 $R[x]_h/(g)$ 中可逆。
\end{enumerate}
\end{definition}

\noindent
一个态射平展，当且仅当它是相对维数 $0$ 的光滑态射
（见定义 \ref{morphisms-definition-smooth-relative-dimension}）。
该定义的一个可取之处是：态射平展的点集自动是开集。

\medskip\noindent
注意，定义中没有分离性或拟紧性假设。因此，平展性在源上具有
局部性质。下面是精确的结果。

\begin{lemma}
\label{morphisms-lemma-etale-characterize}
设 $f : X \to S$ 是概形的态射。下列条件等价：
\begin{enumerate}
\item 态射 $f$ 平展。
\item 对每个仿射开集 $U \subset X$、$V \subset S$，若
$f(U) \subset V$，则环同态
$\mathcal{O}_S(V) \to \mathcal{O}_X(U)$ 平展。
\item 存在开覆盖 $S = \bigcup_{j \in J} V_j$ 以及开覆盖
$f^{-1}(V_j) = \bigcup_{i \in I_j} U_i$，使得每个态射
$U_i \to V_j$（$j\in J, i\in I_j$）都平展。
\item 存在仿射开覆盖 $S = \bigcup_{j \in J} V_j$ 以及仿射开覆盖
$f^{-1}(V_j) = \bigcup_{i \in I_j} U_i$，使得环同态
$\mathcal{O}_S(V_j) \to \mathcal{O}_X(U_i)$ 对所有
$j\in J, i\in I_j$ 都平展。
\end{enumerate}
此外，若 $f$ 平展，则对任意开子概形 $U \subset X$、
$V \subset S$，只要 $f(U) \subset V$，其限制
$f|_U : U \to V$ 就平展。
\end{lemma}

\begin{proof}
只要证明性质“$R \to A$ 平展”是局部性质，结论便由引理
\ref{morphisms-lemma-locally-P} 得出。我们检验定义
\ref{morphisms-definition-property-local} 的条件 (a)、(b) 和 (c)。
这些条件都由代数，引理 \ref{algebra-lemma-etale} 得出。
\end{proof}

\begin{lemma}
\label{morphisms-lemma-composition-etale}
两个平展态射的复合仍平展。
\end{lemma}

\begin{proof}
在引理 \ref{morphisms-lemma-etale-characterize} 的证明中，我们已经看到，
平展是环同态的局部性质。因此，该引理的第一项断言由引理
\ref{morphisms-lemma-composition-type-P} 以及如下事实得出：
平展是一个在复合下稳定的环同态性质；见
代数，引理 \ref{algebra-lemma-etale}。
\end{proof}

\begin{lemma}
\label{morphisms-lemma-base-change-etale}
平展态射的基变换仍平展。
\end{lemma}

\begin{proof}
在引理 \ref{morphisms-lemma-etale-characterize} 的证明中，我们已经看到，
平展是环同态的局部性质。因此，该引理由引理
\ref{morphisms-lemma-composition-type-P} 以及如下事实得出：
平展是一个在基变换下稳定的环同态性质；见
代数，引理 \ref{algebra-lemma-etale}。
\end{proof}

\begin{lemma}
\label{morphisms-lemma-etale-smooth-unramified}
设 $f : X \to S$ 是概形的态射。设 $x \in X$。那么 $f$
在 $x$ 处平展，当且仅当 $f$ 在 $x$ 处既光滑又非分歧。
\end{lemma}

\begin{proof}
这直接由定义得出。
\end{proof}

\begin{lemma}
\label{morphisms-lemma-etale-locally-quasi-finite}
平展态射局部拟有限。
\end{lemma}

\begin{proof}
由引理 \ref{morphisms-lemma-etale-smooth-unramified}，平展态射非分歧。
由引理 \ref{morphisms-lemma-unramified-quasi-finite}，非分歧态射局部拟有限。
\end{proof}

\begin{lemma}
\label{morphisms-lemma-etale-over-field}
\begin{slogan}
域上的平展概形和平展态射的纤维的描述。
\end{slogan}
平展态射的纤维具有下列性质。
\begin{enumerate}
\item 设 $X$ 是域 $k$ 上的概形。结构态射
$X \to \Spec(k)$ 平展，当且仅当 $X$ 是 $k$
的有限可分域扩张的谱的不交并。
\item 若 $f : X \to S$ 是平展态射，则对每个 $s \in S$，
纤维 $X_s$ 都是 $\kappa(s)$ 的有限可分域扩张的谱的不交并。
\end{enumerate}
\end{lemma}

\begin{proof}
可以从引理 \ref{morphisms-lemma-unramified-over-field} 推出此结论，
其中使用上面的引理 \ref{morphisms-lemma-etale-smooth-unramified}。
下面给出一个直接证明。

\medskip\noindent
我们将使用代数，引理 \ref{algebra-lemma-etale-over-field}。
因此显然，若 $X$ 是 $k$ 的有限可分域扩张的谱的不交并，
则 $X \to \Spec(k)$ 平展。反过来，假设
$X \to \Spec(k)$ 平展。那么对任意仿射开集 $U \subset X$，
$U$ 都是 $k$ 的有限可分域扩张的谱的有限不交并。
因而 $X$ 的所有点都是闭点（例如见引理
\ref{morphisms-lemma-algebraic-residue-field-extension-closed-point-fibre}）。
所以 $X$ 是离散空间，证毕。
\end{proof}

\noindent
下面的引理将平展态射刻画为具有“平展纤维”的平坦有限表示态射。

\begin{lemma}
\label{morphisms-lemma-etale-flat-etale-fibres}
设 $f : X \to S$ 是概形的态射。若 $f$ 平坦、局部为有限表示，
并且对每个 $s \in S$，纤维 $X_s$ 都是 $\kappa(s)$
的有限可分域扩张的谱的不交并，则 $f$ 平展。
\end{lemma}

\begin{proof}
可以从代数，引理 \ref{algebra-lemma-characterize-etale}
推出此结论。下面给出另一个证明。

\medskip\noindent
由引理 \ref{morphisms-lemma-etale-over-field}，纤维 $X_s$ 平展，
从而在 $s$ 上光滑。由引理
\ref{morphisms-lemma-smooth-flat-smooth-fibres} 可知 $X \to S$ 光滑。
由引理 \ref{morphisms-lemma-unramified-etale-fibres} 可知 $f$ 非分歧。
最后应用引理 \ref{morphisms-lemma-etale-smooth-unramified}。
\end{proof}

\begin{lemma}
\label{morphisms-lemma-open-immersion-etale}
任何开浸入都平展。
\end{lemma}

\begin{proof}
这是因为开浸入是局部同构。
\end{proof}

\begin{lemma}
\label{morphisms-lemma-etale-syntomic}
平展态射是 syntomic 态射。
\end{lemma}

\begin{proof}
见代数，引理 \ref{algebra-lemma-smooth-syntomic}，并使用如下事实：
平展态射恰好是相对维数 $0$ 的光滑态射。
\end{proof}

\begin{lemma}
\label{morphisms-lemma-etale-locally-finite-presentation}
平展态射局部为有限表示。
\end{lemma}

\begin{proof}
这是因为根据定义，平展环同态是有限表示的。
\end{proof}

\begin{lemma}
\label{morphisms-lemma-etale-flat}
平展态射是平坦态射。
\end{lemma}

\begin{proof}
结合引理 \ref{morphisms-lemma-syntomic-flat} 和
\ref{morphisms-lemma-etale-syntomic}。
\end{proof}

\begin{lemma}
\label{morphisms-lemma-etale-open}
平展态射是开态射。
\end{lemma}

\begin{proof}
结合引理 \ref{morphisms-lemma-etale-flat}、
\ref{morphisms-lemma-etale-locally-finite-presentation} 和
\ref{morphisms-lemma-fppf-open}。
\end{proof}

\noindent
下面的引理说明，局部地，每个平展态射都是标准平展态射。
要在完全一般的情形下证明这个结果，实际上颇为棘手。
棘手的部分隐藏在交换代数一章中。因此，标准平展态射是一般
平展态射的一个{\it 局部模型}。

\begin{lemma}
\label{morphisms-lemma-etale-locally-standard-etale}
设 $f : X  \to S$ 是概形的态射。设 $x \in X$ 是一点。
设 $V \subset S$ 是 $f(x)$ 的仿射开邻域。下列条件等价：
\begin{enumerate}
\item 态射 $f$ 在 $x$ 处平展。
\item 存在仿射开集 $U \subset X$，满足 $x \in U$ 和
$f(U) \subset V$，使得所诱导的态射 $f|_U : U \to V$
是标准平展的（见定义 \ref{morphisms-definition-etale}）。
\end{enumerate}
\end{lemma}

\begin{proof}
这由定义和代数，命题
\ref{algebra-proposition-etale-locally-standard} 得出。
\end{proof}

\noindent
下面给出一点处平展性的微分判据。这个结果有许多变体，
它们各自可能在某处有用。我们只在需要时把它们添加在这里。

\begin{lemma}
\label{morphisms-lemma-etale-at-point}
设 $f : X \to S$ 是概形的态射。设 $x \in X$。
令 $s = f(x)$。假设 $f$ 局部为有限表示。下列条件等价：
\begin{enumerate}
\item 态射 $f$ 在 $x$ 处平展。
\item 局部环同态 $\mathcal{O}_{S, s} \to \mathcal{O}_{X, x}$
平坦，并且 $X_s \to \Spec(\kappa(s))$ 在 $x$ 处平展。
\item 局部环同态 $\mathcal{O}_{S, s} \to \mathcal{O}_{X, x}$
平坦，并且 $X_s \to \Spec(\kappa(s))$ 在 $x$ 处非分歧。
\item 局部环同态 $\mathcal{O}_{S, s} \to \mathcal{O}_{X, x}$
平坦，并且 $\mathcal{O}_{X, x}$-模 $\Omega_{X/S, x}$ 为零。
\item 局部环同态 $\mathcal{O}_{S, s} \to \mathcal{O}_{X, x}$
平坦，并且 $\kappa(x)$-向量空间
$$
\Omega_{X_s/s, x} \otimes_{\mathcal{O}_{X_s, x}} \kappa(x) =
\Omega_{X/S, x} \otimes_{\mathcal{O}_{X, x}} \kappa(x)
$$
为零。
\item 局部环同态 $\mathcal{O}_{S, s} \to \mathcal{O}_{X, x}$
平坦，且有 $\mathfrak m_s\mathcal{O}_{X, x} = \mathfrak m_x$，
并且域扩张 $\kappa(x)/\kappa(s)$ 是有限可分的。
\item 存在仿射开集 $U \subset X$ 和 $V \subset S$，满足
$x \in U$、$f(U) \subset V$，且所诱导的态射
$f|_U : U \to V$ 是相对维数 $0$ 的标准光滑态射。
\item 存在仿射开集 $\Spec(A) = U \subset X$ 和
$\Spec(R) = V \subset S$，其中 $x \in U$ 对应于
$\mathfrak q \subset A$，且 $f(U) \subset V$，使得存在表示
$$
A = R[x_1, \ldots, x_n]/(f_1, \ldots, f_n)
$$
并且
$$
g =
\det
\left(
\begin{matrix}
\partial f_1/\partial x_1 &
\partial f_2/\partial x_1 &
\ldots &
\partial f_n/\partial x_1 \\
\partial f_1/\partial x_2 &
\partial f_2/\partial x_2 &
\ldots &
\partial f_n/\partial x_2 \\
\ldots & \ldots & \ldots & \ldots \\
\partial f_1/\partial x_n &
\partial f_2/\partial x_n &
\ldots &
\partial f_n/\partial x_n
\end{matrix}
\right)
$$
在 $A$ 中的像不属于 $\mathfrak q$。
\item 存在仿射开集 $U \subset X$ 和 $V \subset S$，满足
$x \in U$、$f(U) \subset V$，且所诱导的态射
$f|_U : U \to V$ 是标准平展的。
\item 存在仿射开集 $\Spec(A) = U \subset X$ 和
$\Spec(R) = V \subset S$，其中 $x \in U$ 对应于
$\mathfrak q \subset A$，且 $f(U) \subset V$，使得存在表示
$$
A = R[x]_Q/(P) = R[x, 1/Q]/(P)
$$
其中 $P, Q \in R[x]$，$P$ 是首一的，并且
$P' = \text{d}P/\text{d}x$ 在 $A$ 中的像不属于
$\mathfrak q$。
\end{enumerate}
\end{lemma}

\begin{proof}
使用引理 \ref{morphisms-lemma-etale-locally-standard-etale} 和定义，
可见 (1) 蕴含其余所有条件。对于条件 (2) -- (10) 中的每一个，
结合引理 \ref{morphisms-lemma-smooth-at-point} 和
\ref{morphisms-lemma-unramified-at-point}，通过证明 $f$ 在 $x$
处既光滑又非分歧，再应用引理
\ref{morphisms-lemma-etale-smooth-unramified}，可见 (1) 成立。
省略若干细节。
\end{proof}

\begin{lemma}
\label{morphisms-lemma-flat-unramified-etale}
一个态射在一点处平展，当且仅当它在该点处平坦且 G-非分歧。
一个态射平展，当且仅当它平坦且 G-非分歧。
\end{lemma}

\begin{proof}
这由引理 \ref{morphisms-lemma-etale-at-point} 和
\ref{morphisms-lemma-unramified-at-point} 得出。
\end{proof}

\begin{lemma}
\label{morphisms-lemma-set-points-where-fibres-etale}
设
$$
\xymatrix{
X' \ar[r]_{g'} \ar[d]_{f'} & X \ar[d]^f \\
S' \ar[r]^g & S
}
$$
是概形的笛卡尔图。设 $W \subset X$（相应地，
$W' \subset X'$）是 $f$（相应地，$f'$）平展的点组成的开子概形。
若下列任一条件成立，则 $W' = (g')^{-1}(W)$：
\begin{enumerate}
\item $f$ 平坦且局部为有限表示；或
\item $f$ 局部为有限表示且 $g$ 平坦。
\end{enumerate}
\end{lemma}

\begin{proof}
先假设 $f$ 局部为有限型。考虑集合
$$
T = \{x \in X \mid f\text{ 在下列点非分歧： }x\}
$$
以及相应集合 $T' \subset X'$（对于 $f'$）。由引理
\ref{morphisms-lemma-set-points-where-fibres-unramified}，
$T' = (g')^{-1}(T)$。

\medskip\noindent
因此情形 (1) 成立，因为在情形 (1) 中，由引理
\ref{morphisms-lemma-flat-unramified-etale}，$T$ 正是 $f$ 平展的点组成的
（开）集。

\medskip\noindent
在情形 (2) 中，设 $x' \in W'$。那么 $g'$ 在 $x'$ 处平坦
（引理 \ref{morphisms-lemma-base-change-module-flat}），并且
$g \circ f'$ 在 $x'$ 处平坦（引理
\ref{morphisms-lemma-composition-module-flat}）。由引理
\ref{morphisms-lemma-flat-permanence} 可知，$f$ 在 $x = g'(x')$ 处平坦。
另一方面，由于 $x' \in T'$（引理
\ref{morphisms-lemma-base-change-smooth}），可知 $x \in T$。因此由引理
\ref{morphisms-lemma-etale-at-point}，$f$ 在 $x$ 处平展。
\end{proof}

\begin{lemma}
\label{morphisms-lemma-etale-permanence-better}
\begin{reference}
例如见 \cite[Remark 18.30 (2).]{GWII}。
\end{reference}
设 $f : X \to Y$ 是 $S$ 上概形的态射。若 $X$ 在 $S$ 上平展，
且 $Y$ 在 $S$ 上非分歧，则 $f$ 平展。
\end{lemma}

\begin{proof}
考虑分解 $X \to X \times_S Y \to Y$，其中第一支箭头由
$\text{id}_X$ 和 $f$ 给出，第二支箭头是投影。我们断言两支箭头
都平展，因而由引理 \ref{morphisms-lemma-composition-etale}，$f$ 平展。
确实，该投影平展，因为它是 $X \to S$ 的基变换；见引理
\ref{morphisms-lemma-base-change-etale}。第一支箭头是对角态射
$Y \to Y \times_S Y$ 的基变换，因为方块
$$
\xymatrix{
X \ar[d] \ar[r] & X \times_S Y \ar[d] \\
Y \ar[r] & Y \times_S Y
}
$$
是笛卡尔的。由引理 \ref{morphisms-lemma-diagonal-unramified-morphism}，
对角态射 $Y \to Y \times_S Y$ 是开浸入。开浸入的基变换仍是
开浸入（《概形》，引理
\ref{schemes-lemma-base-change-immersion}），而开浸入平展
（引理 \ref{morphisms-lemma-open-immersion-etale}）。
\end{proof}

\begin{lemma}
\label{morphisms-lemma-etale-permanence}
\begin{slogan}
平展态射的消去律。
\end{slogan}
设 $f : X \to Y$ 是 $S$ 上概形的态射。若 $X$ 和 $Y$
都在 $S$ 上平展，则 $f$ 平展。
\end{lemma}

\begin{proof}
由引理 \ref{morphisms-lemma-etale-smooth-unramified}，$Y \to S$
非分歧，所以可以应用引理
\ref{morphisms-lemma-etale-permanence-better}。另见
代数，引理 \ref{algebra-lemma-map-between-etale}。
\end{proof}

\begin{lemma}
\label{morphisms-lemma-etale-permanence-two}
设
$$
\xymatrix{
X \ar[rr]_f \ar[rd]_p & &
Y \ar[dl]^q \\
& S
}
$$
是概形态射的交换图。假设
\begin{enumerate}
\item $f$ 满射且平展；
\item $p$ 平展；并且
\item $q$ 局部为有限表示\footnote{事实上，这由 (1) 和 (2)
蕴含；见下降，引理
\ref{descent-lemma-flat-finitely-presented-permanence}。
此外，只需假设 $f$ 满射、平坦且局部为有限表示；见
下降，引理 \ref{descent-lemma-smooth-permanence}。}。
\end{enumerate}
那么 $q$ 平展。
\end{lemma}

\begin{proof}
由引理 \ref{morphisms-lemma-smooth-permanence} 可知 $q$ 光滑。
因此只需证明 $q$ 的相对维数为 $0$。这由引理
\ref{morphisms-lemma-dimension-fibre-at-a-point-additive} 以及如下事实得出：
$f$ 和 $p$ 的相对维数都是 $0$。
\end{proof}

\noindent
光滑态射的最后一种刻画是：光滑态射 $f : X \to S$
局部地是一个平展态射与投影 $\mathbf{A}_S^d \to S$ 的复合。

\begin{lemma}
\label{morphisms-lemma-smooth-etale-over-affine-space}
\begin{slogan}
光滑概形在平展局部形如仿射空间。
\end{slogan}
设 $\varphi : X \to Y$ 是概形的态射。设 $x \in X$。
设 $V \subset Y$ 是 $\varphi(x)$ 的仿射开邻域。
若 $\varphi$ 在 $x$ 处光滑，则存在整数 $d \geq 0$
和仿射开集 $U \subset X$，满足 $x \in U$ 和
$\varphi(U) \subset V$，并且存在交换图
$$
\xymatrix{
X \ar[d] & U \ar[l] \ar[d] \ar[r]_-\pi & \mathbf{A}^d_V \ar[ld] \\
Y & V \ar[l]
}
$$
其中 $\pi$ 平展。
\end{lemma}

\begin{proof}
由引理 \ref{morphisms-lemma-smooth-locally-standard-smooth}，可以找到引理中
所述的仿射开集 $U$，使得 $\varphi|_U : U \to V$ 标准光滑。
写成 $U = \Spec(A)$ 和 $V = \Spec(R)$，从而可以写
$$
A = R[x_1, \ldots, x_n]/(f_1, \ldots, f_c)
$$
并且
$$
g =
\det
\left(
\begin{matrix}
\partial f_1/\partial x_1 &
\partial f_2/\partial x_1 &
\ldots &
\partial f_c/\partial x_1 \\
\partial f_1/\partial x_2 &
\partial f_2/\partial x_2 &
\ldots &
\partial f_c/\partial x_2 \\
\ldots & \ldots & \ldots & \ldots \\
\partial f_1/\partial x_c &
\partial f_2/\partial x_c &
\ldots &
\partial f_c/\partial x_c
\end{matrix}
\right)
$$
在 $A$ 中的像可逆。那么显然
$R[x_{c + 1}, \ldots, x_n] \to A$ 是相对维数 $0$
的标准光滑环同态。因此它是相对维数 $0$ 的光滑环同态。
换言之，环同态 $R[x_{c + 1}, \ldots, x_n] \to A$ 平展。
由于 $\mathbf{A}^{n - c}_V = \Spec(R[x_{c + 1}, \ldots, x_n])$，
取 $d = n - c$ 即得本引理。
\end{proof}

\section{相对丰沛层}
\label{morphisms-section-relatively-ample}

\noindent
设 $X$ 是概形，$\mathcal{L}$ 是 $X$ 上的可逆层。
那么 $\mathcal{L}$ 在 $X$ 上丰沛，如果 $X$ 拟紧，并且
$X$ 的每一点都包含在形如 $X_s$ 的仿射开集中，其中
$s \in \Gamma(X, \mathcal{L}^{\otimes n})$ 且 $n \geq 1$；见性质，定义
\ref{properties-definition-ample}。我们如下把它变成一个相对概念。

\begin{definition}
\label{morphisms-definition-relatively-ample}
\begin{reference}
\cite[II Definition 4.6.1]{EGA}
\end{reference}
设 $f : X \to S$ 是概形的态射。设 $\mathcal{L}$ 是可逆
$\mathcal{O}_X$-模。若 $\mathcal{L}$ {\it 相对丰沛}，
也称{\it $f$-相对丰沛}、{\it 在 $X/S$ 上丰沛}或
{\it $f$-丰沛}，其含义是：$f : X \to S$ 拟紧，并且对每个
仿射开集 $V \subset S$，$\mathcal{L}$ 在开子概形
$f^{-1}(V)$（它是 $X$ 的开子概形）上的限制是丰沛的。
\end{definition}

\noindent
注意，$X$ 上相对丰沛层的存在性并不迫使态射 $X \to S$
是有限型的。

\begin{lemma}
\label{morphisms-lemma-ample-power-ample}
设 $X \to S$ 是概形的态射。设 $\mathcal{L}$ 是可逆
$\mathcal{O}_X$-模。设 $n \geq 1$。那么
$\mathcal{L}$ 是 $f$-丰沛的，当且仅当
$\mathcal{L}^{\otimes n}$ 是 $f$-丰沛的。
\end{lemma}

\begin{proof}
这由性质，引理 \ref{properties-lemma-ample-power-ample} 得出。
\end{proof}

\begin{lemma}
\label{morphisms-lemma-relatively-ample-separated}
设 $f : X \to S$ 是概形的态射。若存在一个 $f$-丰沛可逆层，
则 $f$ 分离。
\end{lemma}

\begin{proof}
分离性在基上是局部的（例如见《概形》，引理
\ref{schemes-lemma-characterize-separated}；这也很容易直接由定义
得出）。因此可以假设 $S$ 仿射，并且 $X$ 有丰沛可逆层。
在这种情况下，结果由性质，引理
\ref{properties-lemma-ample-separated} 得出。
\end{proof}

\noindent
相对丰沛可逆层有许多种刻画方式，类似于性质，命题
\ref{properties-proposition-characterize-ample} 中的等价条件。
我们将在需要时把这些刻画添加在这里。

\begin{lemma}
\label{morphisms-lemma-characterize-relatively-ample}
\begin{reference}
\cite[II, Proposition 4.6.3]{EGA}
\end{reference}
设 $f : X \to S$ 是概形的拟紧态射。设 $\mathcal{L}$
是 $X$ 上的可逆层。下列条件等价：
\begin{enumerate}
\item 可逆层 $\mathcal{L}$ 是 $f$-丰沛的。
\item 存在开覆盖 $S = \bigcup V_i$，使得每个
$\mathcal{L}|_{f^{-1}(V_i)}$ 相对于
$f^{-1}(V_i) \to V_i$ 都是丰沛的。
\item 存在仿射开覆盖 $S = \bigcup V_i$，使得每个
$\mathcal{L}|_{f^{-1}(V_i)}$ 都丰沛。
\item 存在拟凝聚分次 $\mathcal{O}_S$-代数 $\mathcal{A}$
以及分次 $\mathcal{O}_X$-代数的映射
$\psi : f^*\mathcal{A} \to \bigoplus_{d \geq 0} \mathcal{L}^{\otimes d}$，
使得 $U(\psi) = X$，并且
$$
r_{\mathcal{L}, \psi} :
X
\longrightarrow
\underline{\text{Proj}}_S(\mathcal{A})
$$
是开浸入（记号见构造，引理
\ref{constructions-lemma-invertible-map-into-relative-proj}）。
\item 态射 $f$ 拟分离，并且上面的第 (4) 项在
$\mathcal{A} = f_*(\bigoplus_{d \geq 0} \mathcal{L}^{\otimes d})$
和 $\psi$ 为伴随映射时成立。
\item 与 (4) 相同，但只要求 $r_{\mathcal{L}, \psi}$ 是浸入。
\end{enumerate}
\end{lemma}

\begin{proof}
由定义立即可知 (1) 蕴含 (2)，而 (2) 蕴含 (3)。
显然 (5) 蕴含 (4)。

\medskip\noindent
假设 (3) 对仿射开覆盖 $S = \bigcup V_i$ 成立。
我们将证明 (5) 成立。由于每个 $f^{-1}(V_i)$ 都有丰沛可逆层，
可知 $f^{-1}(V_i)$ 分离（性质，引理
\ref{properties-lemma-ample-separated}）。因而 $f$ 分离。
由《概形》，引理 \ref{schemes-lemma-push-forward-quasi-coherent}，
可知 $\mathcal{A} = f_*(\bigoplus_{d \geq 0} \mathcal{L}^{\otimes d})$
是拟凝聚分次 $\mathcal{O}_S$-代数。用
$\psi : f^*\mathcal{A} \to \bigoplus_{d \geq 0} \mathcal{L}^{\otimes d}$
表示伴随映射。构造，第
\ref{constructions-section-invertible-relative-proj} 节中对开集
$U(\psi)$ 的描述，以及 $\mathcal{L}|_{f^{-1}(V_i)}$
的丰沛性定义表明 $U(\psi) = X$。此外，构造，引理
\ref{constructions-lemma-invertible-map-into-relative-proj} 第 (3) 项
表明，$r_{\mathcal{L}, \psi}$ 在 $f^{-1}(V_i)$ 上的限制，
与性质，引理 \ref{properties-lemma-map-into-proj} 中的态射相同；
根据性质，引理
\ref{properties-lemma-ample-immersion-into-proj}，该态射是开浸入。
所以 (5) 成立。

\medskip\noindent
下面证明 (4) 蕴含 (1)。假设 (4)。用
$\pi : \underline{\text{Proj}}_S(\mathcal{A}) \to S$
表示结构态射。选取仿射开集 $V \subset S$。由构造，定义
\ref{constructions-definition-relative-proj}，可知
$\pi^{-1}(V) \subset \underline{\text{Proj}}_S(\mathcal{A})$
等于 $\text{Proj}(A)$，其中 $A = \mathcal{A}(V)$ 视为分次环。
因此，$r_{\mathcal{L}, \psi}$ 把 $f^{-1}(V)$ 同构地映到
$\text{Proj}(A)$ 的一个拟紧开集中。此外，
$\mathcal{L}^{\otimes d}$ 同构于
$\mathcal{O}_{\text{Proj}(A)}(d)$ 的拉回，其中某个 $d \geq 1$。
（见构造，引理
\ref{constructions-lemma-invertible-map-into-relative-proj} 第 (3) 项，
以及构造，引理
\ref{constructions-lemma-invertible-map-into-proj} 的最后一项断言。）
由性质中的引理 \ref{properties-lemma-open-in-proj-ample} 和
\ref{properties-lemma-ample-power-ample}，这蕴含
$\mathcal{L}|_{f^{-1}(V)}$ 丰沛。

\medskip\noindent
假设 (6)。由上面 (1) - (5) 的等价性可知，在 $X/S$
上相对丰沛这一性质在 $S$ 上是局部的。因此可以假设 $S$
仿射，并且必须证明 $\mathcal{L}$ 在 $X$ 上丰沛。在这种情况下，
态射 $r_{\mathcal{L}, \psi}$ 与由映射
$r_{\mathcal{L}, \psi} : X \to \text{Proj}(A)$
所给出的态射相等；后者也记作
$\psi : A = \mathcal{A}(V) \to \Gamma_*(X, \mathcal{L})$。
（见上面的引用。）如上还可看出，$\mathcal{L}^{\otimes d}$
是层 $\mathcal{O}_{\text{Proj}(A)}(d)$ 的拉回，其中某个
$d \geq 1$。此外，由于 $X$ 拟紧，
可知 $X$ 被等同于某个拟紧开子概形
$Y \subset \text{Proj}(A)$ 的闭子概形。由构造，引理
\ref{constructions-lemma-ample-on-proj}（另见性质，引理
\ref{properties-lemma-open-in-proj-ample}），可知
$\mathcal{O}_Y(d')$ 是 $Y$ 上的丰沛可逆层，其中某个
$d' \geq 1$。
由于丰沛层在闭子概形上的限制仍丰沛，见性质，引理
\ref{properties-lemma-ample-on-closed}，可知
$\mathcal{O}_Y(d')$ 的拉回丰沛。结合这些结果和性质，引理
\ref{properties-lemma-ample-power-ample}，可得
$\mathcal{L}$ 如所需是丰沛的。
\end{proof}

\begin{lemma}
\label{morphisms-lemma-ample-over-affine}
\begin{reference}
\cite[II Corollary 4.6.6]{EGA}
\end{reference}
设 $f : X \to S$ 是概形的态射。设 $\mathcal{L}$ 是可逆
$\mathcal{O}_X$-模。假设 $S$ 仿射。那么 $\mathcal{L}$
是 $f$-相对丰沛的，当且仅当 $\mathcal{L}$ 在 $X$ 上丰沛。
\end{lemma}

\begin{proof}
这直接由引理 \ref{morphisms-lemma-characterize-relatively-ample} 和定义得出。
\end{proof}

\begin{lemma}
\label{morphisms-lemma-quasi-affine-O-ample}
\begin{reference}
\cite[II Proposition 5.1.6]{EGA}
\end{reference}
设 $f : X \to S$ 是概形的态射。那么 $f$ 拟仿射，当且仅当
$\mathcal{O}_X$ 是 $f$-相对丰沛的。
\end{lemma}

\begin{proof}
这由性质，引理 \ref{properties-lemma-quasi-affine-O-ample}
和定义得出。
\end{proof}

\begin{lemma}
\label{morphisms-lemma-pullback-ample-tensor-relatively-ample}
设 $f : X \to Y$ 是概形的态射，$\mathcal{M}$ 是可逆
$\mathcal{O}_Y$-模，$\mathcal{L}$ 是可逆
$\mathcal{O}_X$-模。
\begin{enumerate}
\item 若 $\mathcal{L}$ 是 $f$-丰沛的且 $\mathcal{M}$ 丰沛，
则 $\mathcal{L} \otimes f^*\mathcal{M}^{\otimes a}$ 在
$a \gg 0$ 时丰沛。
\item 若 $\mathcal{M}$ 丰沛且 $f$ 拟仿射，则
$f^*\mathcal{M}$ 丰沛。
\end{enumerate}
\end{lemma}

\begin{proof}
假设 $\mathcal{L}$ 是 $f$-丰沛的且 $\mathcal{M}$ 丰沛。
根据假设，$Y$ 和 $f$ 都拟紧（见定义
\ref{morphisms-definition-relatively-ample} 和性质，定义
\ref{properties-definition-ample}）。因此 $X$ 拟紧。
由性质，引理 \ref{properties-lemma-ample-separated}，概形
$Y$ 分离；由引理 \ref{morphisms-lemma-relatively-ample-separated}，
态射 $f$ 分离。因此，由《概形》，引理
\ref{schemes-lemma-separated-permanence}，$X$ 分离。
选取 $x \in X$。可以选择 $m \geq 1$ 和
$t \in \Gamma(Y, \mathcal{M}^{\otimes m})$，使得 $Y_t$
仿射且 $f(x) \in Y_t$。由于 $\mathcal{L}$ 在
$f^{-1}(Y_t) = X_{f^*t}$ 上的限制是丰沛可逆层，
可以选择 $n \geq 1$ 和
$s \in \Gamma(X_{f^*t}, \mathcal{L}^{\otimes n})$，使得
$x \in (X_{f^*t})_s$ 且 $(X_{f^*t})_s$ 仿射。
由性质，引理 \ref{properties-lemma-invert-s-sections} 第 (2) 项
（上述情形满足其假设），存在整数 $e \geq 1$ 和截面
$s' \in \Gamma(X, \mathcal{L}^{\otimes n} \otimes f^*\mathcal{M}^{\otimes em})$，
它限制为 $s(f^*t)^e$（在 $X_{f^*t}$ 上）。对任意 $b > 0$，
考虑截面 $s'' = s'(f^*t)^b$，它是
$\mathcal{L}^{\otimes n} \otimes f^*\mathcal{M}^{\otimes (e + b)m}$
的截面。那么
$X_{s''} = (X_{f^*t})_s$ 是 $X$ 中包含 $x$ 的仿射开集。
选取 $b$ 使得 $n$ 整除 $e + b$，可见
$\mathcal{L}^{\otimes n} \otimes f^*\mathcal{M}^{\otimes (e + b)m}$
是 $n$ 次幂，其底为
$\mathcal{L} \otimes f^*\mathcal{M}^{\otimes a}$，其中某个
$a$；而且可以得到任何
$a$，只要它能被 $m$ 整除且足够大。由于 $X$ 拟紧，有限多个这样的仿射开集
覆盖 $X$。由此可知，对于某个充分整除且足够大的 $a$，可逆层
$\mathcal{L} \otimes f^*\mathcal{M}^{\otimes a}$ 在 $X$ 上丰沛。
另一方面，可知 $\mathcal{M}^{\otimes c}$（因而它在 $X$ 上的拉回）
对所有 $c \gg 0$ 都全局生成；见性质，命题
\ref{properties-proposition-characterize-ample}。因此，
$\mathcal{L} \otimes f^*\mathcal{M}^{\otimes a + c}$ 丰沛
（性质，引理
\ref{properties-lemma-ample-tensor-globally-generated}），其中
$c \gg 0$，从而 (1) 得证。

\medskip\noindent
第 (2) 项由引理 \ref{morphisms-lemma-quasi-affine-O-ample}、性质，引理
\ref{properties-lemma-ample-power-ample} 和第 (1) 项得出。
\end{proof}

\begin{lemma}
\label{morphisms-lemma-ample-composition}
设 $g : Y \to S$ 和 $f : X \to Y$ 是概形的态射。
设 $\mathcal{M}$ 是可逆 $\mathcal{O}_Y$-模。
设 $\mathcal{L}$ 是可逆 $\mathcal{O}_X$-模。
若 $S$ 拟紧、$\mathcal{M}$ 是 $g$-丰沛的且
$\mathcal{L}$ 是 $f$-丰沛的，则
$\mathcal{L} \otimes f^*\mathcal{M}^{\otimes a}$
是 $g \circ f$-丰沛的，其中 $a \gg 0$。
\end{lemma}

\begin{proof}
设 $S = \bigcup_{i = 1, \ldots, n} V_i$ 是有限仿射开覆盖。
由引理 \ref{morphisms-lemma-characterize-relatively-ample}，只需证明
$\mathcal{L} \otimes f^*\mathcal{M}^{\otimes a}$ 在
$(g \circ f)^{-1}(V_i)$ 上对 $i = 1, \ldots, n$ 丰沛。
因此，该引理由引理
\ref{morphisms-lemma-pullback-ample-tensor-relatively-ample} 得出。
\end{proof}

\begin{lemma}
\label{morphisms-lemma-ample-base-change}
设 $f : X \to S$ 是概形的态射。设 $\mathcal{L}$ 是可逆
$\mathcal{O}_X$-模。设 $S' \to S$ 是概形的态射。
设 $f' : X' \to S'$ 是 $f$ 的基变换，并用 $\mathcal{L}'$
表示 $\mathcal{L}$ 在 $X'$ 上的拉回。若 $\mathcal{L}$
是 $f$-丰沛的，则 $\mathcal{L}'$ 是 $f'$-丰沛的。
\end{lemma}

\begin{proof}
由引理 \ref{morphisms-lemma-characterize-relatively-ample}，只需找到仿射开覆盖
$S' = \bigcup U'_i$，使得 $\mathcal{L}'$ 在
$(f')^{-1}(U_i')$ 上的限制对所有 $i$ 都是丰沛可逆层。
可以选择 $U'_i$ 映入某个仿射开集 $U_i \subset S$。
在这种情况下，态射
$(f')^{-1}(U'_i) \to f^{-1}(U_i)$ 是仿射的，因为它是仿射态射
$U'_i \to U_i$ 的基变换（引理 \ref{morphisms-lemma-base-change-affine}）。
因此，由引理
\ref{morphisms-lemma-pullback-ample-tensor-relatively-ample}，
$\mathcal{L}'|_{(f')^{-1}(U'_i)}$ 丰沛。
\end{proof}

\begin{lemma}
\label{morphisms-lemma-ample-permanence}
设 $g : Y \to S$ 和 $f : X \to Y$ 是概形的态射。
设 $\mathcal{L}$ 是可逆 $\mathcal{O}_X$-模。
若 $\mathcal{L}$ 是 $g \circ f$-丰沛的，并且 $f$
拟紧\footnote{若 $g$ 拟分离，则这由《概形》，引理
\ref{schemes-lemma-quasi-compact-permanence} 得出。}，
则 $\mathcal{L}$ 是 $f$-丰沛的。
\end{lemma}

\begin{proof}
假设 $f$ 拟紧且 $\mathcal{L}$ 是 $g \circ f$-丰沛的。
设 $U \subset S$ 是仿射开集，设 $V \subset Y$ 是满足
$g(V) \subset U$ 的仿射开集。那么根据假设，
$\mathcal{L}|_{(g \circ f)^{-1}(U)}$ 在
$(g \circ f)^{-1}(U)$ 上丰沛。由于
$f^{-1}(V) \subset (g \circ f)^{-1}(U)$，由性质，引理
\ref{properties-lemma-ample-on-locally-closed} 可知，
$\mathcal{L}|_{f^{-1}(V)}$ 在 $f^{-1}(V)$ 上丰沛。
确切地说，由《概形》，引理
\ref{schemes-lemma-quasi-compact-permanence}，
$f^{-1}(V) \to (g \circ f)^{-1}(U)$ 是拟紧开浸入，因为
$(g \circ f)^{-1}(U)$ 分离（性质，引理
\ref{properties-lemma-ample-separated}），并且
$f^{-1}(V)$ 拟紧（因为 $f$ 拟紧）。因此，由引理
\ref{morphisms-lemma-characterize-relatively-ample}，可得
$\mathcal{L}$ 是 $f$-丰沛的。
\end{proof}

\section{相对甚丰沛层}
\label{morphisms-section-very-ample}

\noindent
回顾：给定拟凝聚层 $\mathcal{E}$（在概形 $S$ 上），
与 $\mathcal{E}$ 相伴的{\it 射影丛}是态射
$\mathbf{P}(\mathcal{E}) \to S$，其中
$\mathbf{P}(\mathcal{E}) = \underline{\text{Proj}}_S(\text{Sym}(\mathcal{E}))$；
见构造，定义 \ref{constructions-definition-projective-bundle}。

\begin{definition}
\label{morphisms-definition-very-ample}
设 $f : X \to S$ 是概形的态射。设 $\mathcal{L}$ 是可逆
$\mathcal{O}_X$-模。若 $\mathcal{L}$ {\it 相对甚丰沛}，
更确切地称为{\it $f$-相对甚丰沛}，也称{\it 在 $X/S$ 上甚丰沛}
或{\it $f$-甚丰沛}，其含义是：存在拟凝聚
$\mathcal{O}_S$-模 $\mathcal{E}$ 和浸入
$i : X \to \mathbf{P}(\mathcal{E})$（在 $S$ 上），使得
$\mathcal{L} \cong i^*\mathcal{O}_{\mathbf{P}(\mathcal{E})}(1)$。
\end{definition}



\noindent
由于这个定义中没有拟紧性假设，一般而言，相对甚丰沛可逆层
未必是相对丰沛可逆层。

\begin{lemma}
\label{morphisms-lemma-ample-very-ample}
\begin{reference}
\cite[II, Proposition 4.6.2]{EGA}
\end{reference}
设 $f : X \to S$ 是概形的态射。设 $\mathcal{L}$ 是可逆
$\mathcal{O}_X$-模。若 $f$ 拟紧且 $\mathcal{L}$
是相对甚丰沛可逆层，则 $\mathcal{L}$ 是相对丰沛可逆层。
\end{lemma}

\begin{proof}
根据定义，存在拟凝聚 $\mathcal{O}_S$-模 $\mathcal{E}$
和浸入 $i : X \to \mathbf{P}(\mathcal{E})$（在 $S$ 上），使得
$\mathcal{L} \cong i^*\mathcal{O}_{\mathbf{P}(\mathcal{E})}(1)$。
令 $\mathcal{A} = \text{Sym}(\mathcal{E})$，于是根据定义，
$\mathbf{P}(\mathcal{E}) = \underline{\text{Proj}}_S(\mathcal{A})$。
分次 $\mathcal{O}_S$-代数 $\mathcal{A}$ 配备映射
$$
\psi :
\mathcal{A} \to
\bigoplus\nolimits_{n \geq 0}
\pi_*\mathcal{O}_{\mathbf{P}(\mathcal{E})}(n) \to
\bigoplus\nolimits_{n \geq 0}
f_*\mathcal{L}^{\otimes n}
$$
其中第二支箭头使用同一
$\mathcal{L} \cong i^*\mathcal{O}_{\mathbf{P}(\mathcal{E})}(1)$。
由 $f_*$ 和 $f^*$ 的伴随性，得到态射
$\psi : f^*\mathcal{A} \to \bigoplus_{n \geq 0}\mathcal{L}^{\otimes n}$。
我们省略验证：与这个映射相伴的态射
$r_{\mathcal{L}, \psi}$ 恰好就是浸入 $i$。
因此，结果由引理 \ref{morphisms-lemma-characterize-relatively-ample}
第 (6) 项得出。
\end{proof}

\noindent
为得到该引理的正确逆命题，我们要问：给定相对丰沛可逆层
$\mathcal{L}$，是否存在整数 $n \geq 1$，使得
$\mathcal{L}^{\otimes n}$ 相对甚丰沛？一般而言并非如此。
有若干因素会阻止这一结论成立：
\begin{enumerate}
\item 即使 $S$ 仿射，也可能不存在可行的有限整数 $n$，
因为 $X \to S$ 不是有限型的；见例
\ref{morphisms-example-not-finite-type-proj}。
\item 基不拟紧意味着，即使 $f$ 为有限型，该结果也可能不成立。
确切地说，给定域 $k$，存在概形 $X_d$，它在 $k$ 上有限型，
带有丰沛可逆层 $\mathcal{O}_{X_d}(1)$，使得
$\mathcal{O}_{X_d}(1)$ 的甚丰沛张量幂中最小的是第 $d$ 次幂。
见例 \ref{morphisms-example-not-bounded}。取 $f$ 为这些概形 $X_d$
映到 $\Spec(k)$ 的各个副本的不交并的态射，即得一个例子。
\end{enumerate}
逆命题的本版本见下面的引理
\ref{morphisms-lemma-finite-type-ample-very-ample}。证明它之前先做一些准备工作。

\begin{example}
\label{morphisms-example-very-ample}
设 $S$ 是概形。设 $\mathcal{A}$ 是拟凝聚分次
$\mathcal{O}_S$-代数，它由 $\mathcal{A}_1$ 生成
（在 $\mathcal{A}_0$ 上）。令
$X = \underline{\text{Proj}}_S(\mathcal{A})$。在这种情况下，
$\mathcal{O}_X(1)$ 可逆且在 $X/S$ 上甚丰沛。确切地说，
与分次 $\mathcal{O}_S$-代数映射
$$
\text{Sym}_{\mathcal{O}_X}^*(\mathcal{A}_1)
\longrightarrow
\mathcal{A}
$$
相伴的态射是闭浸入 $X \to \mathbf{P}(\mathcal{A}_1)$，
它把 $\mathcal{O}_{\mathbf{P}(\mathcal{A}_1)}(1)$ 拉回为
$\mathcal{O}_X(1)$；见构造，引理
\ref{constructions-lemma-surjective-generated-degree-1-map-relative-proj}。
\end{example}

\begin{example}
\label{morphisms-example-not-finite-type-proj}
设 $k$ 是域。考虑分次 $k$-代数
$$
A = k[U, V, Z_1, Z_2, Z_3, \ldots]/I
\quad
\text{其中}
\quad
I = (U^2 - Z_1^2, U^4 - Z_2^2, U^6 - Z_3^2, \ldots)
$$
其分次由 $\deg(U) = \deg(V) = \deg(Z_1) = 1$
和 $\deg(Z_d) = d$ 给出。注意，$X = \text{Proj}(A)$
由 $D_{+}(U)$ 和 $D_{+}(V)$ 覆盖。因此层
$\mathcal{O}_X(n)$ 全都可逆，并同构于
$\mathcal{O}_X(1)^{\otimes n}$。特别地，$\mathcal{O}_X(1)$
是丰沛的，并且是 $f$-丰沛的，其中态射为
$f : X \to \Spec(k)$。
我们断言，$\mathcal{O}_X(1)$ 的任何幂都不是 $f$-相对甚丰沛的。
确实，容易看出 $\Gamma(X, \mathcal{O}_X(n))$ 是第 $n$
个分次分量（代数 $A$ 的）。因此，若 $\mathcal{O}_X(n)$ 甚丰沛，
则 $X$ 会是 $k$ 上某个射影空间的闭子概形，从而在 $k$
上为有限型。另一方面，$D_{+}(V)$ 是
$k[t, t_1, t_2, \ldots]/(t^2 - t_1^2, t^4 - t_2^2, t^6 - t_3^2, \ldots)$
的谱，而后者在 $k$ 上不是有限型的。
\end{example}

\begin{example}
\label{morphisms-example-not-bounded}
设 $k$ 是无限域。设 $\lambda_1, \lambda_2, \lambda_3, \ldots$
是 $k^*$ 中两两不同的元素。（这并非严格必要；事实上，即使所有
$\lambda_i$ 都等于 $1$，这个例子也完全成立。）
考虑分次 $k$-代数
$$
A_d = k[U, V, Z]/I_d
\quad
\text{其中}
\quad
I_d = (Z^2 - \prod\nolimits_{i = 1}^{2d} (U - \lambda_i V)).
$$
其分次由 $\deg(U) = \deg(V) = 1$ 和 $\deg(Z) = d$ 给出。
那么 $X_d = \text{Proj}(A_d)$ 有丰沛可逆层
$\mathcal{O}_{X_d}(1)$。我们断言，若
$\mathcal{O}_{X_d}(n)$ 甚丰沛，则 $n \geq d$。
理由是 $Z$ 的次数为 $d$，从而
$\Gamma(X_d, \mathcal{O}_{X_d}(n)) =
k[U, V]_n$ 对 $n < d$ 成立。省略细节。
\end{example}

\begin{lemma}
\label{morphisms-lemma-relatively-very-ample-separated}
设 $f : X \to S$ 是概形的态射。设 $\mathcal{L}$ 是 $X$ 上的可逆层。
若 $\mathcal{L}$ 在 $X/S$ 上相对甚丰沛，则 $f$ 分离。
\end{lemma}

\begin{proof}
分离性在基底上是局部的（见概形，节
\ref{schemes-section-separation-axioms}）。浸入是分离的
（见概形，引理 \ref{schemes-lemma-immersions-monomorphisms}）。
因此引理成立，因为局部地，$X$ 可浸入一个分次环的齐次谱，
而该齐次谱是分离的；见构造，引理
\ref{constructions-lemma-proj-separated}。
\end{proof}

\begin{lemma}
\label{morphisms-lemma-relatively-very-ample}
设 $f : X \to S$ 是概形的态射。设 $\mathcal{L}$ 是 $X$ 上的可逆层。
假设 $f$ 拟紧。下列条件等价：
\begin{enumerate}
\item $\mathcal{L}$ 在 $X/S$ 上相对甚丰沛；
\item 存在开覆盖 $S = \bigcup V_j$，使得
$\mathcal{L}|_{f^{-1}(V_j)}$ 在 $f^{-1}(V_j)/V_j$ 上相对甚丰沛，
对所有 $j$ 均成立；
\item 存在拟凝聚分次 $\mathcal{O}_S$-代数层 $\mathcal{A}$，
它由次数 $1$ 的部分在 $\mathcal{O}_S$ 上生成，并且存在分次
$\mathcal{O}_X$-代数映射
$\psi : f^*\mathcal{A} \to \bigoplus_{n \geq 0} \mathcal{L}^{\otimes n}$，
使得 $f^*\mathcal{A}_1 \to \mathcal{L}$ 满射，且相伴态射
$r_{\mathcal{L}, \psi} : X \to \underline{\text{Proj}}_S(\mathcal{A})$
是浸入；
\item $f$ 拟分离，典范映射
$\psi : f^*f_*\mathcal{L} \to \mathcal{L}$ 满射，且相伴映射
$r_{\mathcal{L}, \psi} : X \to \mathbf{P}(f_*\mathcal{L})$ 是浸入。
\end{enumerate}
\end{lemma}

\begin{proof}
显然，(1) 蕴含 (2)。同样显然，(4) 蕴含 (1)；(4) 中的拟分离性假设
用于通过概形，引理 \ref{schemes-lemma-push-forward-quasi-coherent}
保证 $f_*\mathcal{L}$ 拟凝聚。

\medskip\noindent
假设 (2)。我们证明 (4)。令 $S = \bigcup V_j$ 是 (2) 中的开覆盖。
令 $X_j = f^{-1}(V_j)$，并令 $f_j : X_j \to V_j$ 为 $f$ 的限制。
由引理 \ref{morphisms-lemma-relatively-very-ample-separated} 可知 $f$ 分离
（因为分离性在基底上是局部的）。依假设，存在拟凝聚
$\mathcal{O}_{V_j}$-模 $\mathcal{E}_j$ 和浸入
$i_j : X_j \to \mathbf{P}(\mathcal{E}_j)$，满足
$\mathcal{L}|_{X_j} \cong i_j^*\mathcal{O}_{\mathbf{P}(\mathcal{E}_j)}(1)$。
态射 $i_j$ 对应于满射
$f_j^*\mathcal{E}_j \to \mathcal{L}|_{X_j}$；见构造，节
\ref{constructions-section-projective-bundle}。这个映射伴随于一个映射
$\mathcal{E}_j \to f_*\mathcal{L}|_{V_j}$，使得复合
$$
f_j^*\mathcal{E}_j \to (f^*f_*\mathcal{L})|_{X_j} \to \mathcal{L}|_{X_j}
$$
满射。于是
$\psi : f^*f_*\mathcal{L} \to \mathcal{L}$ 满射。令
$r_{\mathcal{L}, \psi} : X \to \mathbf{P}(f_*\mathcal{L})$
为相伴态射。还须证明 $r_{\mathcal{L}, \psi}$ 是浸入；我们建议读者自行证明。
$\mathcal{O}_{V_j}$-模映射 $\mathcal{E}_j \to f_*\mathcal{L}|_{V_j}$
诱导对称代数之间的同态，后者又定义态射
$$
\mathbf{P}(f_*\mathcal{L}|_{V_j}) \supset U_j \longrightarrow
\mathbf{P}(\mathcal{E}_j)
$$
其中 $U_j$ 是构造，引理
\ref{constructions-lemma-morphism-relative-proj} 中的开子概形。
$\psi$ 与 $\mathcal{E}_j \to f_*\mathcal{L}|_{V_j}$ 的相容性说明
$r_{\mathcal{L}, \psi}(X_j) \subset U_j$，并且存在分解
$$
\xymatrix{
X_j \ar[r]^-{r_{\mathcal{L}, \psi}} & U_j \ar[r] &
\mathbf{P}(\mathcal{E}_j)
}
$$
我们省略验证。这说明 $r_{\mathcal{L}, \psi}$ 是浸入。

\medskip\noindent
至此可见，(1)、(2) 与 (4) 等价。显然 (4) 蕴含 (3)。
假设 (3)。我们证明 (1)。令 $\mathcal{A}$ 是拟凝聚分次
$\mathcal{O}_S$-代数层，它由次数 $1$ 的部分在 $\mathcal{O}_S$ 上生成。
考虑分次 $\mathcal{O}_S$-代数映射
$\text{Sym}(\mathcal{A}_1) \to \mathcal{A}$。依假设，此映射满射，
因而诱导闭浸入
$$
\underline{\text{Proj}}_S(\mathcal{A})
\longrightarrow
\mathbf{P}(\mathcal{A}_1)
$$
它把 $\mathcal{O}(1)$ 拉回为 $\mathcal{O}(1)$；见构造，引理
\ref{constructions-lemma-surjective-generated-degree-1-map-relative-proj}。
因此，(3) 蕴含 (1) 是显然的。
\end{proof}

\begin{lemma}
\label{morphisms-lemma-very-ample-base-change}
设 $f : X \to S$ 是概形的态射。设 $\mathcal{L}$ 是可逆
$\mathcal{O}_X$-模。设 $S' \to S$ 是概形的态射。令
$f' : X' \to S'$ 为 $f$ 的基变换，并用 $\mathcal{L}'$ 表示
$\mathcal{L}$ 到 $X'$ 的拉回。若 $\mathcal{L}$ 是 $f$-甚丰沛的，
则 $\mathcal{L}'$ 是 $f'$-甚丰沛的。
\end{lemma}

\begin{proof}
由定义 \ref{morphisms-definition-very-ample}，存在拟凝聚
$\mathcal{O}_S$-模 $\mathcal{E}$ 和浸入
$i : X \to \mathbf{P}(\mathcal{E})$（在 $S$ 上），使得
$\mathcal{L} \cong i^*\mathcal{O}_{\mathbf{P}(\mathcal{E})}(1)$。
$\mathbf{P}(\mathcal{E})$ 到 $S'$ 的基变换是与模 $\mathcal{E}'$
相伴的射影丛，其中该模是 $\mathcal{E}$ 的拉回；而
$\mathcal{O}_{\mathbf{P}(\mathcal{E})}(1)$ 的拉回是
$\mathcal{O}_{\mathbf{P}(\mathcal{E}')}(1)$；见构造，引理
\ref{constructions-lemma-relative-proj-base-change}。最后，浸入的基变换仍是浸入
（概形，引理 \ref{schemes-lemma-base-change-immersion}）。
\end{proof}







\section{相对于有限型态射的丰沛层与甚丰沛层}
\label{morphisms-section-ample-finite-type}

\noindent
事实上，本节大部分内容涉及尚未定义的（拟）射影态射概念。

\begin{lemma}
\label{morphisms-lemma-very-ample-finite-type-over-affine}
设 $f : X \to S$ 是概形的态射。设 $\mathcal{L}$ 是 $X$ 上的可逆层。
假设：
\begin{enumerate}
\item 可逆层 $\mathcal{L}$ 在 $X/S$ 上甚丰沛；
\item 态射 $X \to S$ 为有限型；并且
\item $S$ 仿射。
\end{enumerate}
则存在 $n \geq 0$ 以及浸入
$i : X \to \mathbf{P}^n_S$（在 $S$ 上），使得
$\mathcal{L} \cong i^*\mathcal{O}_{\mathbf{P}^n_S}(1)$。
\end{lemma}

\begin{proof}
假设 (1)、(2) 和 (3)。条件 (3) 意味着
$S = \Spec(R)$，其中 $R$ 为某个环。根据定义，条件 (1) 意味着存在拟凝聚
$\mathcal{O}_S$-模 $\mathcal{E}$ 和浸入
$\alpha : X \to \mathbf{P}(\mathcal{E})$，使得
$\mathcal{L} = \alpha^*\mathcal{O}_{\mathbf{P}(\mathcal{E})}(1)$。
写成 $\mathcal{E} = \widetilde{M}$，其中某个 $R$-模为 $M$。于是
$$
\mathbf{P}(\mathcal{E}) = \text{Proj}(\text{Sym}_R(M)).
$$
由于 $\alpha$ 是浸入，而
$\text{Proj}(\text{Sym}_R(M))$ 的拓扑由标准开集
$D_{+}(f)$ 生成，其中 $f \in \text{Sym}_R^d(M)$、$d \geq 1$，
故对每个 $x \in X$，可找到
$f \in \text{Sym}_R^d(M)$、$d \geq 1$，使得
$\alpha(x) \in D_{+}(f)$，并使
$$
\alpha|_{\alpha^{-1}(D_{+}(f))} : \alpha^{-1}(D_{+}(f)) \to D_{+}(f)
$$
为闭浸入。条件 (2) 蕴含 $X$ 拟紧。因此可找到有限多个元素
$f_j \in \text{Sym}_R^{d_j}(M)$、$d_j \geq 1$，使得对每个
$f = f_j$，上面显示的映射为闭浸入，并且
$\alpha(X) \subset \bigcup D_{+}(f_j)$。写
$U_j = \alpha^{-1}(D_{+}(f_j))$。注意，$U_j$ 是仿射概形
$D_{+}(f_j)$ 的闭子概形，因而仿射。写 $U_j = \Spec(A_j)$。
条件 (2) 还蕴含 $A_j$ 在 $R$ 上为有限型；见引理
\ref{morphisms-lemma-locally-finite-type-characterize}。选择有限多个
$x_{j, k} \in A_j$，使它们生成 $A_j$（作为 $R$-代数）。
由于 $\alpha|_{U_j}$ 是闭浸入，可见 $x_{j, k}$ 是某个元素
$$
f_{j, k}/f_j^{e_{j, k}} \in \text{Sym}_R(M)_{(f_j)}
=
\Gamma(D_{+}(f_j), \mathcal{O}_{\text{Proj}(\text{Sym}_R(M))}).
$$
的像。最后，选择 $n \geq 1$ 以及元素 $y_0, \ldots, y_n \in M$，
使得每个多项式 $f_j, f_{j, k} \in \text{Sym}_R(M)$ 都是元素
$y_t$ 的多项式，其系数在 $R$ 中。考虑分次环映射
$$
\psi : R[Y_0, \ldots, Y_n] \longrightarrow \text{Sym}_R(M),
\quad Y_i \longmapsto y_i.
$$
用 $F_j$、$F_{j, k}$ 表示 $R[Y_0, \ldots, Y_n]$ 中满足
$\psi(F_j) = f_j$ 和 $\psi(F_{j, k}) = f_{j, k}$ 的元素。
由构造，引理 \ref{constructions-lemma-morphism-proj}，得到开子概形
$$
U(\psi) \subset \text{Proj}(\text{Sym}_R(M))
$$
和态射 $r_\psi : U(\psi) \to \mathbf{P}^n_R$。此态射满足
$r_\psi^{-1}(D_{+}(F_j)) = D_{+}(f_j)$，因而可见
$\alpha(X) \subset U(\psi)$。此外，显然
$$
i = r_\psi \circ \alpha : X \longrightarrow \mathbf{P}^n_R
$$
仍是浸入，因为依构造
$i^\sharp(F_{j, k}/F_j^{e_{j, k}}) = x_{j, k} \in
A_j = \Gamma(U_j, \mathcal{O}_X)$。此外，态射 $r_\psi$ 配备映射
$\theta : r_\psi^*\mathcal{O}_{\mathbf{P}^n_R}(1)
\to \mathcal{O}_{\text{Proj}(\text{Sym}_R(M))}(1)|_{U(\psi)}$，
在此情形中该映射为同构（$\theta$ 的构造见上引引理；略去一些细节）。
由于原映射 $\alpha$ 按假设满足
$\mathcal{L} = \alpha^*\mathcal{O}_{\text{Proj}(\text{Sym}_R(M))}(1)$，
结论成立。
\end{proof}

\begin{lemma}
\label{morphisms-lemma-quasi-affine-finite-type-over-S}
设 $\pi : X \to S$ 是概形的态射。假设 $X$ 拟仿射，并且
$\pi$ 局部为有限型。则存在 $n \geq 0$ 以及浸入
$i : X \to \mathbf{A}^n_S$（在 $S$ 上）。
\end{lemma}

\begin{proof}
令 $A = \Gamma(X, \mathcal{O}_X)$。依假设，$X$ 拟紧，并被视为
$\Spec(A)$ 的开子概形；见性质，引理
\ref{properties-lemma-quasi-affine}。此外，诸开集 $X_f$ 中满足
$f \in A$ 且 $X_f$ 仿射的那些，构成 $X$ 的拓扑的一组基；
见性质，引理 \ref{properties-lemma-quasi-affine} 的证明。
因此可找到有限多个 $f_j \in A$（$j = 1, \ldots, m$），使得
$X = \bigcup X_{f_j}$，并且有 $\pi(X_{f_j}) \subset V_j$，其中
$V_j \subset S$ 是某个仿射开集。由引理
\ref{morphisms-lemma-locally-finite-type-characterize}，环映射
$\mathcal{O}(V_j) \to \mathcal{O}(X_{f_j}) = A_{f_j}$ 为有限型。
因此可选择 $a_1, \ldots, a_N \in A$，使得元素
$a_1, \ldots, a_N, 1/f_j$ 生成 $A_{f_j}$（作为
$\mathcal{O}(V_j)$-代数），对每个 $j$ 均如此。取 $n = m + N$，令
$$
i : X \longrightarrow \mathbf{A}^n_S
$$
为由全局截面 $f_1, \ldots, f_m, a_1, \ldots, a_N$ 给出的态射；
这些是 $X$ 的结构层的全局截面。
令 $D(x_j) \subset \mathbf{A}^n_S$ 为第 $j$ 个坐标函数非零处的开子概形。
那么，对 $1 \leq j \leq m$，有
$i^{-1}(D(x_j)) = X_{f_j}$，并且诱导态射
$X_{f_j} \to D(x_j)$ 经由仿射开集
\begin{center}
$\Spec(\mathcal{O}(V_j)[x_1, \ldots, x_n, 1/x_j])$
\end{center}
分解；
该开集位于 $D(x_j)$ 中。
由于环映射
$\mathcal{O}(V_j)[x_1, \ldots, x_n, 1/x_j] \to A_{f_j}$
依构造满射，故可得 $i^{-1}(D(x_j)) \to D(x_j)$
如所需是浸入。
\end{proof}

\begin{lemma}
\label{morphisms-lemma-quasi-projective-finite-type-over-S}
设 $f : X \to S$ 是概形的态射。设 $\mathcal{L}$ 是 $X$ 上的可逆层。
假设：
\begin{enumerate}
\item 可逆层 $\mathcal{L}$ 在 $X$ 上丰沛；并且
\item 态射 $X \to S$ 局部为有限型。
\end{enumerate}
则存在 $d_0 \geq 1$，使得对每个 $d \geq d_0$，存在
$n \geq 0$ 以及浸入 $i : X \to \mathbf{P}^n_S$（在 $S$ 上），满足
$\mathcal{L}^{\otimes d} \cong i^*\mathcal{O}_{\mathbf{P}^n_S}(1)$。
\end{lemma}

\begin{proof}
令
$A = \Gamma_*(X, \mathcal{L}) =
\bigoplus_{d \geq 0} \Gamma(X, \mathcal{L}^{\otimes d})$。
由性质，命题 \ref{properties-proposition-characterize-ample}，
仿射开集 $X_a$ 中满足 $a \in A_{+}$ 为齐次元的那些，
构成 $X$ 的拓扑的一组基。
因此可找到有限多个这样的元素 $a_0, \ldots, a_n \in A_{+}$，使得：
\begin{enumerate}
\item $X = \bigcup_{i = 0, \ldots, n} X_{a_i}$；
\item 每个 $X_{a_i}$ 都仿射；并且
\item 每个 $X_{a_i}$ 都映入某个仿射开集 $V_i \subset S$。
\end{enumerate}
由引理 \ref{morphisms-lemma-locally-finite-type-characterize}，可见环映射
$\mathcal{O}_S(V_i) \to \mathcal{O}_X(X_{a_i})$ 为有限型。
因此可找到有限多个元素
$f_{ij} \in \mathcal{O}_X(X_{a_i})$（$j = 1, \ldots, n_i$），
它们生成 $\mathcal{O}_X(X_{a_i})$（作为
$\mathcal{O}_S(V_i)$-代数）。由性质，引理
\ref{properties-lemma-invert-s-sections}，可将每个
$f_{ij}$ 写为 $a_{ij}/a_i^{e_{ij}}$，其中某个
$a_{ij} \in A_{+}$ 为齐次元。令 $N$ 是所有元素
$a_i$、$a_{ij}$ 的次数的一个正公倍数。考虑元素
$$
a_i^{N/\deg(a_i)}, \ a_{ij}a_i^{(N/\deg(a_i)) - e_{ij}} \in A_N.
$$
依构造，它们生成可逆层 $\mathcal{L}^{\otimes N}$（在 $X$ 上）。
因此它们给出态射
$$
j : X \longrightarrow
\mathbf{P}_S^{m}
\quad
\text{其中 } m = n + \sum n_i
$$
（在 $S$ 上）；见构造，引理
\ref{constructions-lemma-projective-space} 及定义
\ref{constructions-definition-projective-space}。此外，
$j^*\mathcal{O}_{\mathbf{P}_S}(1) = \mathcal{L}^{\otimes N}$。
我们将齐次坐标命名为 $T_0, \ldots, T_n, T_{ij}$，而不是
$T_0, \ldots, T_m$。对 $i = 0, \ldots, n$，有
$j^{-1}(D_{+}(T_i)) = X_{a_i}$。此外，拉回元素
$T_{ij}/T_i$（通过 $j^\sharp$）时，得到元素
$f_{ij} \in \mathcal{O}_X(X_{a_i})$。
因此，态射 $j$ 在 $X_{a_i}$ 上的限制给出闭浸入：将
$X_{a_i}$ 浸入仿射开集
$D_{+}(T_i) \cap \mathbf{P}^m_{V_i}$；该开集位于
$\mathbf{P}^N_S$ 中。故可得态射 $j$ 为浸入。
这说明引理对某些 $d$ 和 $n$ 成立，这在几乎所有应用中已经足够。

\medskip\noindent
这证明：对某个 $d_2 \geq 1$（即上面的 $d_2 = N$）和某个
$m \geq 0$，存在由全局截面给出的浸入
$j : X \to \mathbf{P}^m_S$；这些截面为
$s'_0, \ldots, s'_m \in \Gamma(X, \mathcal{L}^{\otimes d_2})$。
由性质，命题 \ref{properties-proposition-characterize-ample}，
存在整数 $d_1$，使得 $\mathcal{L}^{\otimes d}$ 全局生成，
对所有 $d \geq d_1$ 都如此。令 $d_0 = d_1 + d_2$。
我们断言引理对这个 $d_0$ 成立。确实，给定整数
$d \geq d_0$，可选择
$s''_1, \ldots, s''_t
\in \Gamma(X, \mathcal{L}^{\otimes d - d_2})$，使其生成
$\mathcal{L}^{\otimes d - d_2}$（在 $X$ 上）。令 $k = (m + 1)t$，并用
$s_0, \ldots, s_k$ 表示截面的集合
$s'_\alpha s''_\beta$，其中 $\alpha = 0, \ldots, m$、
$\beta = 1, \ldots, t$。这些截面生成
$\mathcal{L}^{\otimes d}$（在 $X$ 上），因而定义态射
$$
i : X \longrightarrow \mathbf{P}^{k - 1}_S
$$
使得 $i^*\mathcal{O}_{\mathbf{P}^n_S}(1) \cong \mathcal{L}^{\otimes d}$。
为看出 $i$ 是浸入，注意 $i$ 是复合
$$
X \longrightarrow \mathbf{P}^m_S \times_S \mathbf{P}^{t - 1}_S
\longrightarrow \mathbf{P}^{k - 1}_S
$$
其中第一个态射是 $(j, j')$，$j'$ 由
$s''_1, \ldots, s''_t$ 给出，而第二个态射是 Segre 嵌入
（构造，引理 \ref{constructions-lemma-segre-embedding}）。
由于 $j$ 是浸入，$(j, j')$ 也是浸入（将引理
\ref{morphisms-lemma-immersion-permanence} 应用于
$X \to \mathbf{P}^m_S \times_S \mathbf{P}^{t - 1}_S
\to \mathbf{P}^m_S$）。因此 $i$ 是浸入的复合，从而仍为浸入
（概形，引理 \ref{schemes-lemma-composition-immersion}）。
\end{proof}

\begin{lemma}
\label{morphisms-lemma-finite-type-over-affine-ample-very-ample}
设 $f : X \to S$ 是概形的态射。设 $\mathcal{L}$ 是可逆
$\mathcal{O}_X$-模。假设 $S$ 仿射且 $f$ 为有限型。下列条件等价：
\begin{enumerate}
\item $\mathcal{L}$ 在 $X$ 上丰沛；
\item $\mathcal{L}$ 是 $f$-丰沛的；
\item $\mathcal{L}^{\otimes d}$ 是 $f$-甚丰沛的，对某个 $d \geq 1$ 成立；
\item $\mathcal{L}^{\otimes d}$ 是 $f$-甚丰沛的，对所有 $d \gg 1$ 成立；
\item 对某个 $d \geq 1$，存在 $n \geq 1$ 以及浸入
$i : X \to \mathbf{P}^n_S$，使得
$\mathcal{L}^{\otimes d} \cong i^*\mathcal{O}_{\mathbf{P}^n_S}(1)$；并且
\item 对所有 $d \gg 1$，存在 $n \geq 1$ 以及浸入
$i : X \to \mathbf{P}^n_S$，使得
$\mathcal{L}^{\otimes d} \cong i^*\mathcal{O}_{\mathbf{P}^n_S}(1)$。
\end{enumerate}
\end{lemma}

\begin{proof}
(1) 与 (2) 的等价性就是引理 \ref{morphisms-lemma-ample-over-affine}。
(2) $\Rightarrow$ (6) 是引理
\ref{morphisms-lemma-quasi-projective-finite-type-over-S}。显然，(6) 蕴含 (5)。
由于 $\mathbf{P}^n_S$ 是 $S$ 上的射影丛（见构造，引理
\ref{constructions-lemma-projective-space-bundle}），由相对甚丰沛层的定义，
可见 (5) 蕴含 (3)，(6) 蕴含 (4)。显然，(4) 蕴含 (3)。
最后还须证明 (3) 蕴含 (2)，这由引理
\ref{morphisms-lemma-ample-very-ample} 及引理
\ref{morphisms-lemma-ample-power-ample} 得出。
\end{proof}

\begin{lemma}
\label{morphisms-lemma-finite-type-ample-very-ample}
设 $f : X \to S$ 是概形的态射。设 $\mathcal{L}$ 是可逆
$\mathcal{O}_X$-模。假设 $S$ 拟紧且 $f$ 为有限型。下列条件等价：
\begin{enumerate}
\item $\mathcal{L}$ 是 $f$-丰沛的；
\item $\mathcal{L}^{\otimes d}$ 是 $f$-甚丰沛的，对某个 $d \geq 1$ 成立；
\item $\mathcal{L}^{\otimes d}$ 是 $f$-甚丰沛的，对所有 $d \gg 1$ 成立。
\end{enumerate}
\end{lemma}

\begin{proof}
显然，(3) 蕴含 (2)。引理 \ref{morphisms-lemma-ample-very-ample} 保证 (2)
蕴含 (1)，因为根据定义，有限型态射拟紧。假设 $\mathcal{L}$ 是
$f$-丰沛的。选取有限仿射开覆盖
$S = V_1 \cup \ldots \cup V_m$。写 $X_i = f^{-1}(V_i)$。
由上面的引理 \ref{morphisms-lemma-finite-type-over-affine-ample-very-ample} 可见，
存在 $d_0$，使得 $\mathcal{L}^{\otimes d}$ 在 $X_i/V_i$ 上相对甚丰沛，
对所有 $d \geq d_0$ 都如此。因此，由引理
\ref{morphisms-lemma-relatively-very-ample}，可得 (1) 蕴含 (3)。
\end{proof}

\noindent
下面两个引理给出有限型态射情形中，相对甚丰沛可逆层与
相对丰沛可逆层最常用也最有用的刻画。

\begin{lemma}
\label{morphisms-lemma-characterize-very-ample-on-finite-type}
设 $f : X \to S$ 是概形的态射。设 $\mathcal{L}$ 是 $X$ 上的可逆层。
假设 $f$ 为有限型。下列条件等价：
\begin{enumerate}
\item $\mathcal{L}$ 是 $f$-相对甚丰沛的；并且
\item 存在开覆盖 $S = \bigcup V_j$；对每个 $j$，存在整数 $n_j$
以及浸入
$$
i_j :
X_j = f^{-1}(V_j) = V_j \times_S X
\longrightarrow
\mathbf{P}^{n_j}_{V_j}
$$
（在 $V_j$ 上），使得
$\mathcal{L}|_{X_j} \cong i_j^*\mathcal{O}_{\mathbf{P}^{n_j}_{V_j}}(1)$。
\end{enumerate}
\end{lemma}

\begin{proof}
取 $S$ 的一个仿射开覆盖，并将引理
\ref{morphisms-lemma-very-ample-finite-type-over-affine} 应用于 $f$ 和
$\mathcal{L}$ 的每个限制，便可见 (1) 蕴含 (2)。由引理
\ref{morphisms-lemma-relatively-very-ample} 可见 (2) 蕴含 (1)。
\end{proof}

\begin{lemma}
\label{morphisms-lemma-characterize-ample-on-finite-type}
设 $f : X \to S$ 是概形的态射。设 $\mathcal{L}$ 是 $X$ 上的可逆层。
假设 $f$ 为有限型。下列条件等价：
\begin{enumerate}
\item $\mathcal{L}$ 是 $f$-相对丰沛的；并且
\item 存在开覆盖 $S = \bigcup V_j$；对每个 $j$，存在整数
$d_j \geq 1$、$n_j \geq 0$ 以及浸入
$$
i_j :
X_j = f^{-1}(V_j) = V_j \times_S X
\longrightarrow
\mathbf{P}^{n_j}_{V_j}
$$
（在 $V_j$ 上），使得
$\mathcal{L}^{\otimes d_j}|_{X_j} \cong
i_j^*\mathcal{O}_{\mathbf{P}^{n_j}_{V_j}}(1)$。
\end{enumerate}
\end{lemma}

\begin{proof}
取 $S$ 的一个仿射开覆盖，并将引理
\ref{morphisms-lemma-finite-type-over-affine-ample-very-ample} 应用于 $f$ 和
$\mathcal{L}$ 的每个限制，便可见 (1) 蕴含 (2)。由引理
\ref{morphisms-lemma-characterize-relatively-ample} 可见 (2) 蕴含 (1)。
\end{proof}

\begin{lemma}
\label{morphisms-lemma-invertible-add-enough-ample-very-ample}
设 $f : X \to S$ 是概形的态射。设 $\mathcal{N}$、$\mathcal{L}$
为可逆 $\mathcal{O}_X$-模。假设 $S$ 拟紧、$f$ 为有限型，
且 $\mathcal{L}$ 是 $f$-丰沛的。则
$\mathcal{N} \otimes_{\mathcal{O}_X} \mathcal{L}^{\otimes d}$
是 $f$-甚丰沛的，对所有 $d \gg 1$ 都如此。
\end{lemma}

\begin{proof}
由引理 \ref{morphisms-lemma-characterize-very-ample-on-finite-type}，可归约到
$S$ 仿射的情形。结合引理
\ref{morphisms-lemma-finite-type-over-affine-ample-very-ample} 与性质，命题
\ref{properties-proposition-characterize-ample}，可找到整数 $d_0$，使得
$\mathcal{N} \otimes \mathcal{L}^{\otimes d_0}$ 全局生成。选择它的全局截面
$s_0, \ldots, s_n$，即 $\mathcal{N} \otimes \mathcal{L}^{\otimes d_0}$
的生成元。这确定一个态射 $j : X \to \mathbf{P}^n_S$（在 $S$ 上）。
由引理 \ref{morphisms-lemma-finite-type-over-affine-ample-very-ample}，还可选取整数
$d_1$，使得对所有 $d \geq d_1$，存在截面
$t_{d, 0}, \ldots, t_{d, n(d)}$（属于 $\mathcal{L}^{\otimes d}$），
它们生成该层并定义浸入
$$
j_d = \varphi_{\mathcal{L}^{\otimes d}, t_{d, 0}, \ldots, t_{d, n(d)}} :
X
\longrightarrow
\mathbf{P}^{n(d)}_S
$$
（在 $S$ 上）。于是，对 $d \geq d_0 + d_1$，可考虑态射
$$
\varphi_{\mathcal{N} \otimes \mathcal{L}^{\otimes d},
s_j \otimes t_{d - d_0, i}} :
X
\longrightarrow
\mathbf{P}^{(n + 1)(n(d - d_0) + 1) - 1}_S
$$
此态射是浸入，因为它是复合
$$
X \to \mathbf{P}^n_S \times_S \mathbf{P}^{n(d - d_0)}_S
\to \mathbf{P}^{(n + 1)(n(d - d_0) + 1) - 1}_S
$$
其中第一个态射是 $(j, j_{d - d_0})$，第二个态射是 Segre 嵌入
（构造，引理 \ref{constructions-lemma-segre-embedding}）。
由于 $j$ 是浸入，$(j, j_{d - d_0})$ 也是浸入（应用引理
\ref{morphisms-lemma-immersion-permanence}）。这是浸入的复合，因而仍为浸入
（概形，引理 \ref{schemes-lemma-composition-immersion}）。
\end{proof}







\section{拟射影态射}
\label{morphisms-section-quasi-projective}

\noindent
上一节的讨论提示了下列定义。我们采用 \cite{EGA} 中的拟射影定义；
带字母 ``H'' 的版本是 \cite{H} 中的定义。

\begin{definition}
\label{morphisms-definition-quasi-projective}
\begin{reference}
\cite[II, Definition 5.3.1]{EGA} and \cite[page 103]{H}
\end{reference}
设 $f : X \to S$ 是概形的态射。
\begin{enumerate}
\item 称 $f$ {\it 拟射影}，如果 $f$ 为有限型，且存在一个
$f$-相对丰沛的可逆 $\mathcal{O}_X$-模。
\item 称 $f$ {\it H-拟射影}，如果存在拟紧浸入
$X \to \mathbf{P}^n_S$（在 $S$ 上），其中 $n$ 为某个整数。
\footnote{这与 Hartshorne 中的定义不完全相同。确切地说，
Hartshorne 中的定义（1997 年第八次修订印刷）要求 $f$
是开浸入后接 H-射影态射的复合（见定义
\ref{morphisms-definition-projective}），而这并不蕴含 $f$ 拟紧。反方向的蕴含见引理
\ref{morphisms-lemma-H-quasi-projective-open-H-projective}。}
\item 称 $f$ {\it 局部拟射影}，如果存在开覆盖
$S = \bigcup V_j$，使得每个 $f^{-1}(V_j) \to V_j$ 都拟射影。
\end{enumerate}
\end{definition}

\noindent
正如这个定义所表明的，拟射影性质在 $S$ 上不是局部的。
稍后我们将能进一步说明拟射影概形的范畴；见态射进阶，节
\ref{more-morphisms-section-quasi-projective}。

\begin{lemma}
\label{morphisms-lemma-base-change-quasi-projective}
拟射影态射的基变换仍为拟射影态射。
\end{lemma}

\begin{proof}
这由引理 \ref{morphisms-lemma-base-change-finite-type} 和
\ref{morphisms-lemma-ample-base-change} 得出。
\end{proof}

\begin{lemma}
\label{morphisms-lemma-composition-quasi-projective}
设 $f : X \to Y$ 和 $g : Y \to S$ 是概形的态射。若 $S$ 拟紧，
且 $f$ 和 $g$ 都拟射影，则 $g \circ f$ 拟射影。
\end{lemma}

\begin{proof}
这由引理 \ref{morphisms-lemma-composition-finite-type} 和
\ref{morphisms-lemma-ample-composition} 得出。
\end{proof}

\begin{lemma}
\label{morphisms-lemma-quasi-projective-properties}
设 $f : X \to S$ 是概形的态射。若 $f$ 拟射影、H-拟射影或局部拟射影，
则 $f$ 是有限型的分离态射。
\end{lemma}

\begin{proof}
省略。
\end{proof}

\begin{lemma}
\label{morphisms-lemma-H-quasi-projective-quasi-projective}
H-拟射影态射是拟射影态射。
\end{lemma}

\begin{proof}
省略。
\end{proof}

\begin{lemma}
\label{morphisms-lemma-characterize-locally-quasi-projective}
设 $f : X \to S$ 是概形的态射。下列条件等价：
\begin{enumerate}
\item 态射 $f$ 局部拟射影。
\item 存在开覆盖 $S = \bigcup V_j$，使得每个
$f^{-1}(V_j) \to V_j$ 都 H-拟射影。
\end{enumerate}
\end{lemma}

\begin{proof}
由引理 \ref{morphisms-lemma-H-quasi-projective-quasi-projective} 可见 (2) 蕴含 (1)。
假设 (1)。这个问题在 $S$ 上是局部的，故可假设 $S$ 仿射、
$X$ 在 $S$ 上为有限型，并且 $\mathcal{L}$ 是 $X/S$ 上的相对丰沛可逆层。
由引理 \ref{morphisms-lemma-finite-type-over-affine-ample-very-ample}，可假设
$\mathcal{L}$ 在 $X$ 上丰沛。由引理
\ref{morphisms-lemma-quasi-projective-finite-type-over-S}，可见存在 $X$
到 $S$ 上某个射影空间的浸入，即 $X$ 在 $S$ 上 H-拟射影，
如所需。
\end{proof}

\begin{lemma}
\label{morphisms-lemma-quasi-affine-finite-type-quasi-projective}
\begin{reference}
\cite[II, Proposition 5.3.4 (i)]{EGA}
\end{reference}
有限型拟仿射态射是拟射影态射。
\end{lemma}

\begin{proof}
这由引理 \ref{morphisms-lemma-quasi-affine-O-ample} 得出。
\end{proof}

\begin{lemma}
\label{morphisms-lemma-quasi-projective-permanence}
设 $g : Y \to S$ 和 $f : X \to Y$ 是概形的态射。若
$g \circ f$ 拟射影且 $f$ 拟紧\footnote{若 $g$ 拟分离，则这由概形，引理
\ref{schemes-lemma-quasi-compact-permanence} 得出。}，
则 $f$ 拟射影。
\end{lemma}

\begin{proof}
由引理 \ref{morphisms-lemma-permanence-finite-type}，注意到 $f$ 为有限型。
因而引理由引理 \ref{morphisms-lemma-ample-permanence} 及定义得出。
\end{proof}





\section{固有态射}
\label{morphisms-section-proper}

\noindent
固有态射的概念在代数几何中起着重要作用。固有态射的一个重要例子
将是结构态射 $\mathbf{P}^n_S \to S$，即射影 $n$-空间的结构态射；事实上，
正是这个启发性例子引出了该定义。

\begin{definition}
\label{morphisms-definition-proper}
设 $f : X \to S$ 是概形的态射。若 $f$ 分离、为有限型且万有闭，
则称 $f$ {\it 固有}。
\end{definition}

\noindent
从零点加倍的仿射直线到仿射直线的态射为有限型且万有闭，
因此上述定义中的分离条件是必要的。本节其余部分将证明固有态射
和万有闭态射的一些基本性质。

\begin{lemma}
\label{morphisms-lemma-universally-closed-local-on-the-base}
设 $f : X \to S$ 是概形的态射。下列条件等价：
\begin{enumerate}
\item 态射 $f$ 万有闭。
\item 存在开覆盖 $S = \bigcup V_j$，使得
$f^{-1}(V_j) \to V_j$ 对所有指标 $j$ 均万有闭。
\end{enumerate}
\end{lemma}

\begin{proof}
这由定义显然可知。
\end{proof}

\begin{lemma}
\label{morphisms-lemma-proper-local-on-the-base}
设 $f : X \to S$ 是概形的态射。下列条件等价：
\begin{enumerate}
\item 态射 $f$ 固有。
\item 存在开覆盖 $S = \bigcup V_j$，使得
$f^{-1}(V_j) \to V_j$ 对所有指标 $j$ 均固有。
\end{enumerate}
\end{lemma}

\begin{proof}
省略。
\end{proof}

\begin{lemma}
\label{morphisms-lemma-composition-proper}
固有态射的复合是固有态射。万有闭态射也同样如此。
\end{lemma}

\begin{proof}
闭态射的复合是闭态射。若
$X \to Y \to Z$ 是万有闭态射，且 $Z' \to Z$ 是任意态射，
则可见
$Z' \times_Z X = (Z' \times_Z Y) \times_Y X  \to Z' \times_Z Y$
是闭的，并且 $Z' \times_Z Y \to Z'$ 是闭的。
这就证明了万有闭态射的结论。我们已经看到，“分离”和“有限型”
在复合下保持（概形，引理
\ref{schemes-lemma-separated-permanence} 及引理
\ref{morphisms-lemma-composition-finite-type}）。因此，固有态射的结论也成立。
\end{proof}

\begin{lemma}
\label{morphisms-lemma-base-change-proper}
固有态射的基变换是固有态射。万有闭态射也同样如此。
\end{lemma}

\begin{proof}
对万有闭态射而言，这由定义成立。对分离态射而言，这成立
（概形，引理 \ref{schemes-lemma-separated-permanence}）。
对有限型态射而言，这成立（引理 \ref{morphisms-lemma-base-change-finite-type}）。
因此，对固有态射也成立。
\end{proof}

\begin{lemma}
\label{morphisms-lemma-closed-immersion-proper}
闭浸入是固有的，因而更是万有闭的。
\end{lemma}

\begin{proof}
闭浸入的基变换是闭浸入（概形，引理
\ref{schemes-lemma-base-change-immersion}），因而万有闭。
闭浸入分离（概形，引理
\ref{schemes-lemma-immersions-monomorphisms}）。
闭浸入为有限型（引理 \ref{morphisms-lemma-immersion-locally-finite-type}）。
因此闭浸入固有。
\end{proof}

\begin{lemma}
\label{morphisms-lemma-image-proper-scheme-closed}
设给定概形的交换图
$$
\xymatrix{
X \ar[rr] \ar[rd] & &
Y \ar[ld] \\
& S &
}
$$
其中 $Y$ 在 $S$ 上分离。
\begin{enumerate}
\item 若 $X \to S$ 万有闭，则态射 $X \to Y$ 万有闭。
\item 若 $X$ 在 $S$ 上固有，则态射 $X \to Y$ 固有。
\end{enumerate}
特别地，在两种情形中，$X$ 在 $Y$ 中的像都是闭的。
\end{lemma}

\begin{proof}
假设 $X \to S$ 万有闭（相应地，固有）。将该态射分解为
$X \to X \times_S Y \to Y$。第一个态射是闭浸入；见概形，引理
\ref{schemes-lemma-semi-diagonal}。因此，第一个态射固有
（引理 \ref{morphisms-lemma-closed-immersion-proper}）。投影
$X \times_S Y \to Y$ 是一个万有闭（相应地，固有）态射的基变换，
因而万有闭（相应地，固有）；见引理 \ref{morphisms-lemma-base-change-proper}。
于是，作为万有闭（相应地，固有）态射的复合，$X \to Y$
万有闭（相应地，固有）（引理 \ref{morphisms-lemma-composition-proper}）。
\end{proof}

\noindent
下列引理的证明归功于 Bjorn Poonen；见
\href{https://mathoverflow.net/questions/23337/is-a-universally-closed-morphism-of-schemes-quasi-compact/23528#23528}{此处}。

\begin{lemma}
\label{morphisms-lemma-universally-closed-quasi-compact}
\begin{reference}
Due to Bjorn Poonen.
\end{reference}
概形的万有闭态射是拟紧的。
\end{lemma}

\begin{proof}
设 $f : X \to S$ 是态射。假设 $f$ 不拟紧。目标是证明 $f$
不万有闭。由概形，引理 \ref{schemes-lemma-quasi-compact-affine}，
存在仿射开集 $V \subset S$，使得 $f^{-1}(V)$ 不拟紧。
为达到目标，只须证明 $f^{-1}(V) \to V$ 不万有闭；因此可假设
$S = \Spec(A)$，其中 $A$ 表示某个环。

\medskip\noindent
写 $X = \bigcup_{i \in I} X_i$，其中诸 $X_i$ 是 $X$ 的仿射开子概形。
令 $T = \Spec(A[y_i ; i \in I])$。令 $T_i = D(y_i) \subset T$。
令 $Z$ 为闭集
$(X \times_S T) - \bigcup_{i \in I} (X_i \times_S T_i)$。
只须证明像 $f_T(Z)$ 不是闭的；这里，$Z$ 在
$f_T : X \times_S T \to T$ 下取像。

\medskip\noindent
存在点 $s \in S$，使得不存在邻域 $U$（它是 $s$ 在 $S$ 中的邻域），
满足 $X_U$ 拟紧。否则可覆盖 $S$，只需有限多个这样的 $U$，而概形，引理
\ref{schemes-lemma-quasi-compact-affine} 将蕴含 $f$ 拟紧。
固定这样的 $s \in S$。

\medskip\noindent
首先验证 $f_T(Z_s) \ne T_s$。令 $t \in T$ 为位于 $s$ 上方的点，
满足 $\kappa(t) = \kappa(s)$，并且 $y_i = 1$
在 $\kappa(t)$ 中成立，对所有 $i$ 都如此。于是有 $t \in T_i$，
对所有 $i$ 都如此；而
$Z_s \to T_s$ 在 $t$ 上方的纤维同构于
$(X - \bigcup_{i \in I} X_i)_s$，后者为空。因此
$t \in T_s - f_T(Z_s)$。

\medskip\noindent
假设 $f_T(Z)$ 在 $T$ 中闭。则存在元素
$g \in A[y_i; i \in I]$，使得 $f_T(Z) \subset V(g)$，但
$t \not \in V(g)$。因此 $g$ 在 $\kappa(t)$ 中的像非零。
特别地，$g$ 的某个系数在 $\kappa(s)$ 中的像非零，因而这个系数
在某个邻域 $U$ 上可逆，其中该邻域包含 $s$。令 $J$ 为所有满足
$j \in I$ 且 $y_j$ 出现在 $g$ 中的指标所成的有限集。由于 $X_U$ 不拟紧，
可选择点 $x \in X - \bigcup_{j \in J} X_j$，它位于某个
$u \in U$ 上方。由于 $g$ 有一个在 $U$ 上可逆的系数，
可找到点 $t' \in T$，它位于 $u$ 上方，使得
$t' \not \in V(g)$，并且 $t' \in V(y_i)$ 对所有
$i \notin J$ 都成立。这是因为
$V(y_i; i \in I, i \not\in J) = \Spec(A[y_j; j\in J])$，
而这个概形中位于 $u$ 上方的点集与
$\Spec(\kappa(u)[y_j; j \in J])$ 双射。换言之，
$t' \notin T_i$ 对每个 $i \notin J$ 都成立。由概形，引理
\ref{schemes-lemma-points-fibre-product}，可找到
点 $z$，它属于 $X \times_S T$，并映到 $x \in X$ 和 $t' \in T$。
由于 $x \not \in X_j$ 对 $j \in J$ 成立，而
$t' \not \in T_i$ 对 $i \in I \setminus J$ 成立，故 $z \in Z$。
另一方面，$f_T(z) = t' \not \in V(g)$，这与
$f_T(Z) \subset V(g)$ 矛盾。因此，“$f_T(Z)$ 闭”的假设是错误的，
确实可得 $f_T$ 不闭，如所需。
\end{proof}

\noindent
下列引理说明，固有概形在基上有限型的分离概形中的像仍然固有。

\begin{lemma}
\label{morphisms-lemma-image-proper-is-proper}
设 $S$ 是概形。设 $f : X \to Y$ 是 $S$ 上概形的态射。
若 $X$ 在 $S$ 上万有闭且 $f$ 满射，则 $Y$ 在 $S$ 上万有闭。
特别地，若 $Y$ 在 $S$ 上还分离且局部为有限型，则 $Y$ 在 $S$ 上固有。
\end{lemma}

\begin{proof}
假设 $X$ 万有闭且 $f$ 满射。记 $p : X \to S$、$q : Y \to S$
为结构态射。设 $S' \to S$ 是概形的态射。基变换
$f' : X_{S'} \to Y_{S'}$ 是满射
（引理 \ref{morphisms-lemma-base-change-surjective}），而基变换
$p' : X_{S'} \to S'$ 是闭态射。若 $T \subset Y_{S'}$ 是闭集，
则 $(f')^{-1}(T) \subset X_{S'}$ 是闭集，因而
$p'((f')^{-1}(T)) = q'(T)$ 是闭集。所以 $q'$ 是闭态射。
这证明了第一个断言。于是由引理
\ref{morphisms-lemma-universally-closed-quasi-compact}，$Y \to S$ 拟紧；
因此，若另外 $Y \to S$ 局部为有限型且分离，则根据定义，
$Y \to S$ 固有。
\end{proof}

\begin{lemma}
\label{morphisms-lemma-scheme-theoretic-image-is-proper}
设给定概形的交换图
$$
\xymatrix{
X \ar[rr]_h \ar[rd]_f & & Y \ar[ld]^g \\
& S
}
$$
假设
\begin{enumerate}
\item $X \to S$ 是万有闭态射（例如固有态射），并且
\item $Y \to S$ 分离且局部为有限型。
\end{enumerate}
则概形论像 $Z \subset Y$（即 $h$ 的概形论像）在 $S$ 上固有，且
$X \to Z$ 是满射。
\end{lemma}

\begin{proof}
$h$ 的概形论像在第 \ref{morphisms-section-scheme-theoretic-image} 节中构造。
由于 $f$ 拟紧（引理 \ref{morphisms-lemma-universally-closed-quasi-compact}），
可知 $h$ 拟紧（概形，引理
\ref{schemes-lemma-quasi-compact-permanence}）。因此
$h(X) \subset Z$ 稠密（引理
\ref{morphisms-lemma-quasi-compact-scheme-theoretic-image}）。另一方面，
$h(X)$ 在 $Y$ 中闭（引理
\ref{morphisms-lemma-image-proper-scheme-closed}），故 $X \to Z$ 是满射。
于是 $Z \to S$ 固有（引理 \ref{morphisms-lemma-image-proper-is-proper}）。
\end{proof}

\noindent
分离概形在满射万有闭态射下的目标是分离的。

\begin{lemma}
\label{morphisms-lemma-image-universally-closed-separated}
设 $S$ 是概形。设 $f : X \to Y$ 是 $S$ 上概形的满射万有闭态射。
\begin{enumerate}
\item 若 $X$ 拟分离，则 $Y$ 拟分离。
\item 若 $X$ 分离，则 $Y$ 分离。
\item 若 $X$ 在 $S$ 上拟分离，则 $Y$ 在 $S$ 上拟分离。
\item 若 $X$ 在 $S$ 上分离，则 $Y$ 在 $S$ 上分离。
\end{enumerate}
\end{lemma}

\begin{proof}
取 $S = \Spec(\mathbf{Z})$ 时，(1) 和 (2) 是 (3) 和 (4) 的推论
（见概形，定义 \ref{schemes-definition-separated}）。考虑交换图
$$
\xymatrix{
X \ar[d] \ar[rr]_{\Delta_{X/S}} & & X \times_S X \ar[d] \\
Y \ar[rr]^{\Delta_{Y/S}} & & Y \times_S Y
}
$$
左侧竖直箭头是满射（即泛满射）。右侧竖直箭头是万有闭的，因为它是
万有闭态射 $X \times_S X \to X \times_S Y \to Y \times_S Y$
的复合。因此它也拟紧；见引理
\ref{morphisms-lemma-universally-closed-quasi-compact}。

\medskip\noindent
假设 $X$ 在 $S$ 上拟分离，即 $\Delta_{X/S}$ 拟紧。
若 $V \subset Y \times_S Y$ 是拟紧开集，则
$V \times_{Y \times_S Y} X \to \Delta_{Y/S}^{-1}(V)$ 是满射，
且由上面的说明，$V \times_{Y \times_S Y} X$ 拟紧。
于是 $\Delta_{Y/S}$ 拟紧，即 $Y$ 在 $S$ 上拟分离。

\medskip\noindent
假设 $X$ 在 $S$ 上分离，即 $\Delta_{X/S}$ 是闭浸入。
那么 $X \to Y \times_S Y$ 作为闭态射的复合是闭态射。
由于 $X \to Y$ 满射，可知 $\Delta_{Y/S}(Y)$ 在
$Y \times_S Y$ 中闭。因此，由概形，定义
\ref{schemes-definition-separated} 后面的讨论，$Y$ 在 $S$ 上分离。
\end{proof}








\section{赋值判据}
\label{morphisms-section-valuative-criteria}

\noindent
我们已经在概形，第
\ref{schemes-section-valuative-criterion-universal-closedness} 节和第
\ref{schemes-section-valuative-separatedness} 节中讨论了万有闭性和分离性的
赋值判据。本节将讨论若干推论和变体。在极限，第
\ref{limits-section-Noetherian-valuative-criterion} 节中，我们将证明，
对局部 Noether 概形及有限型态射，只须考虑离散赋值环。

\begin{lemma}[固有性的赋值判据]
\label{morphisms-lemma-characterize-proper}
\begin{reference}
\cite[II Theorem 7.3.8]{EGA}
\end{reference}
设 $S$ 是概形。设 $f : X \to Y$ 是 $S$ 上概形的态射。
假设 $f$ 为有限型且拟分离。则下列条件等价：
\begin{enumerate}
\item $f$ 固有；
\item $f$ 满足赋值判据
（概形，定义 \ref{schemes-definition-valuative-criterion}）；
\item 对任意交换实线图
$$
\xymatrix{
\Spec(K) \ar[r] \ar[d] & X \ar[d] \\
\Spec(A) \ar[r] \ar@{-->}[ru] & Y
}
$$
其中 $A$ 是以 $K$ 为分式域的赋值环，都存在唯一的虚线箭头使图交换。
\end{enumerate}
\end{lemma}

\begin{proof}
(3) 是 (2) 的改述。因此，该引理是概形，命题
\ref{schemes-proposition-characterize-universally-closed}、引理
\ref{schemes-lemma-separated-implies-valuative}、引理
\ref{schemes-lemma-valuative-criterion-separatedness} 以及诸定义的形式推论。
\end{proof}

\noindent
检验赋值判据时，通常不必考虑所有可能的图。我们把下一个引理中的
赋值判据称为“精细赋值判据”。

\begin{lemma}
\label{morphisms-lemma-refined-valuative-criterion-universally-closed}
设 $f : X \to S$ 和 $h : U \to X$ 是概形的态射。
假设 $f$ 和 $h$ 拟紧，且 $h(U)$ 在 $X$ 中稠密。若对任意交换实线图
$$
\xymatrix{
\Spec(K) \ar[r] \ar[d] & U \ar[r]^h & X \ar[d]^f \\
\Spec(A) \ar[rr] \ar@{-->}[rru] & & S
}
$$
其中 $A$ 是以 $K$ 为分式域的赋值环，都存在唯一的虚线箭头使图交换，
则 $f$ 万有闭。若此外 $f$ 拟分离，则 $f$ 分离。
\end{lemma}

\begin{proof}
为证明 $f$ 万有闭，我们验证 $f$ 的赋值判据的存在性部分；由概形，命题
\ref{schemes-proposition-characterize-universally-closed}，这已经足够。
为此，考虑交换图
$$
\xymatrix{
\Spec(K) \ar[r] \ar[d] & X \ar[d] \\
\Spec(A) \ar[r] & S
}
$$
其中 $A$ 是赋值环，而 $K$ 是 $A$ 的分式域。注意，由于赋值环和域
都是既约的，根据概形，引理 \ref{schemes-lemma-map-into-reduction}，
可分别用各自的既约概形替换 $U$、$X$ 和 $S$。在这种情形下，
$h(U)$ 稠密这一假设意味着 $h : U \to X$ 的概形论像就是 $X$；
见引理 \ref{morphisms-lemma-scheme-theoretic-image-reduced}。还可以把 $S$ 替换为
$\Spec(A) \to S$ 所经过的一个仿射开集。因此可假设
$S = \Spec(R)$。

\medskip\noindent
设 $\Spec(B) \subset X$ 是 $\Spec(K) \to X$ 所经过的仿射开集。
选择一个多项式代数 $P$（定义在 $B$ 上）及一个 $B$-代数满射 $P \to K$。
于是 $\Spec(P) \to X$ 平坦。因此，由引理
\ref{morphisms-lemma-flat-base-change-scheme-theoretic-image}，态射
$U \times_X \Spec(P) \to \Spec(P)$ 的概形论像是 $\Spec(P)$。
由引理 \ref{morphisms-lemma-reach-points-scheme-theoretic-image}，可找到交换图
$$
\xymatrix{
\Spec(K') \ar[r] \ar[d] & U \times_X \Spec(P) \ar[d] \\
\Spec(A') \ar[r] & \Spec(P)
}
$$
其中 $A'$ 是赋值环，$K'$ 是 $A'$ 的分式域，并且
$\Spec(A')$ 的闭点映到 $\Spec(K) \subset \Spec(P)$。换言之，
存在 $B$-代数同态
$\varphi : K \to A'/\mathfrak m_{A'}$。选择一个赋值环
$A'' \subset A'/\mathfrak m_{A'}$，它支配 $\varphi(A)$，并以
$K'' = A'/\mathfrak m_{A'}$ 为分式域
（代数，引理 \ref{algebra-lemma-dominate}）。令
$$
C = \{\lambda \in A' \mid \lambda \bmod \mathfrak m_{A'} \in A''\}.
$$
由代数，引理 \ref{algebra-lemma-stack-valuation-rings}，这是一个赋值环。
由于 $C$ 是 $R$-代数，并以 $K'$ 为分式域，得到交换图
$$
\xymatrix{
\Spec(K') \ar[r] \ar[d] & U \ar[r] & X \ar[d] \\
\Spec(C) \ar[rr] \ar@{-->}[rru] & & S
}
$$
它具有引理陈述中的形式。因此存在按图所示的虚线箭头。由引理中的
唯一性假设，复合 $\Spec(A') \to \Spec(C) \to X$ 与给定态射
$\Spec(A') \to \Spec(P) \to \Spec(B) \subset X$ 一致。
因此，把该态射限制到
$C/\mathfrak m_{A'} = A''$ 的谱上，就得到给定态射
$\Spec(K'') = \Spec(A'/\mathfrak m_{A'}) \to \Spec(K) \to X$。
令 $x \in X$ 是 $\Spec(A'') \to X$ 下闭点的像。所诱导的环同态
$\mathcal{O}_{X, x} \to A''$ 的像是一个局部子环，它包含在
$K \subset K''$ 中。由于 $A$ 在 $K$ 中关于支配关系是极大的，
并且 $A \subset A''$，故 $A = K \cap A''$。因此
$\mathcal{O}_{X, x} \to A''$ 通过 $A \subset A''$ 分解。
这样便得到所需箭头 $\Spec(A) \to X$。

\medskip\noindent
最后，假设 $f$ 拟分离。那么 $\Delta : X \to X \times_S X$ 拟紧。
给定实线图
$$
\xymatrix{
\Spec(K) \ar[r] \ar[d] & U \ar[r]^h & X \ar[d]^\Delta \\
\Spec(A) \ar[rr] \ar@{-->}[rru] & & X \times_S X
}
$$
其中 $A$ 是以 $K$ 为分式域的赋值环，存在唯一的虚线箭头使图交换。
事实上，下方的水平箭头等价于一对态射 $\Spec(A) \to X$；
它们都可在引理的图中充当虚线箭头。因此，所要求的唯一性表明
下方的水平箭头通过 $\Delta$ 分解。于是可把刚证明的结果应用于
$\Delta : X \to X \times_S X$ 和 $h : U \to X$，得到
$\Delta$ 万有闭。显然，这意味着 $f$ 分离。
\end{proof}

\begin{remark}
\label{morphisms-remark-check-val-on-open}
引理 \ref{morphisms-lemma-refined-valuative-criterion-universally-closed} 中关于虚线箭头
唯一性的假设是必要的（细节省略）。当然，若 $f$ 分离，则唯一性有保证
（概形，引理 \ref{schemes-lemma-separated-implies-valuative}）。
\end{remark}

\begin{lemma}
\label{morphisms-lemma-morphism-defined-local-ring}
设 $S$ 是概形。设 $X$、$Y$ 是 $S$ 上的概形。设 $s \in S$，
而 $x \in X$、$y \in Y$ 是位于 $s$ 上方的点。
\begin{enumerate}
\item 设 $f, g : X \to Y$ 是 $S$ 上的态射，满足
$f(x) = g(x) = y$，并且
$f^\sharp_x = g^\sharp_x : \mathcal{O}_{Y, y} \to \mathcal{O}_{X, x}$。
在下列各情形中，存在开邻域 $U \subset X$，使得
$f|_U = g|_U$：
\begin{enumerate}
\item $Y$ 在 $S$ 上局部为有限型；
\item $X$ 整；
\item $X$ 局部 Noether；或
\item $X$ 既约且只有有限多个不可约分支。
\end{enumerate}
\item 设 $\varphi : \mathcal{O}_{Y, y} \to \mathcal{O}_{X, x}$
是局部 $\mathcal{O}_{S, s}$-代数同态。在下列各情形中，存在
开邻域 $U \subset X$，它包含 $x$，以及态射 $f : U \to Y$，
它把 $x$ 映到 $y$，并满足 $f^\sharp_x = \varphi$：
\begin{enumerate}
\item $Y$ 在 $S$ 上局部为有限表示；
\item $Y$ 局部为有限型且 $X$ 整；
\item $Y$ 局部为有限型且 $X$ 局部 Noether；或
\item $Y$ 局部为有限型且 $X$ 既约并只有有限多个不可约分支。
\end{enumerate}
\end{enumerate}
\end{lemma}

\begin{proof}
证明 (1)。可以把 $X$、$Y$、$S$ 分别替换为 $x$、$y$、$s$
的适当仿射开邻域，并归结为下列代数问题：给定环 $R$ 和两个
$R$-代数同态 $\varphi, \psi : B \to A$，满足
\begin{enumerate}
\item $R \to B$ 为有限型，或 $A$ 是整环，或 $A$ 是 Noether 环，
或 $A$ 既约且只有有限多个极小素理想；
\item 两个同态 $B \to A_\mathfrak p$ 对某个素理想
$\mathfrak p \subset A$ 相同；
\end{enumerate}
要证明，$\varphi, \psi$ 定义相同的同态 $B \to A_g$，其中可适当选择
$g \in A$，且 $g \not \in \mathfrak p$。若 $R \to B$ 为有限型，
设 $t_1, \ldots, t_m \in B$ 是 $B$ 作为 $R$-代数的一组生成元。
对每个 $j$，可找到 $g_j \in A$、$g_j \not \in \mathfrak p$，
使得 $\varphi(t_j)$ 和 $\psi(t_j)$ 在 $A_{g_j}$ 中的像相同。
然后令 $g = \prod g_j$。在其他情形中（即 $A$ 是整环、Noether 环，
或既约且只有有限多个极小素理想），可找到 $g \in A$、
$g \not \in \mathfrak p$，使得 $A_g \subset A_\mathfrak p$。
见代数，引理 \ref{algebra-lemma-subring-of-local-ring}。
因此同态 $B \to A_g$ 相同，如所需。

\medskip\noindent
证明 (2)。为此，可以把 $X$、$Y$ 和 $S$ 替换为适当的仿射开集。
设 $X = \Spec(A)$、$Y = \Spec(B)$ 且 $S = \Spec(R)$。
设 $\mathfrak p \subset A$ 是对应于 $x$ 的素理想。
设 $\mathfrak q \subset B$ 是对应于 $y$ 的素理想。于是
$\varphi$ 是局部 $R$-代数同态
$\varphi : B_\mathfrak q \to A_\mathfrak p$。
若 $R \to B$ 是有限表示的环同态，则存在
$g \in A \setminus \mathfrak p$ 及 $R$-代数同态 $B \to A_g$，使图
$$
\xymatrix{
B_\mathfrak q \ar[r]_\varphi & A_\mathfrak p \\
B \ar[u] \ar[r] & A_g \ar[u]
}
$$
交换；见代数，引理 \ref{algebra-lemma-characterize-finite-presentation} 和
\ref{algebra-lemma-localization-colimit}。所诱导的态射
$\Spec(A_g) \to \Spec(B)$ 即满足要求。若 $B$ 在 $R$ 上为有限型，
设 $t_1, \ldots, t_m \in B$ 是 $B$ 作为 $R$-代数的一组生成元。
于是可选择 $g_j \in A$、$g_j \not \in \mathfrak p$，使得
$\varphi(t_j) \in \Im(A_{g_j} \to A_\mathfrak p)$。
因此，把 $A$ 替换为 $A[1/\prod g_j]$ 后，可假设 $B$ 映入
$A \to A_\mathfrak p$ 的像中。若能找到 $g \in A$、
$g \not \in \mathfrak p$，使得 $A_g \to A_\mathfrak p$ 是单射，
就会得到所需的 $R$-代数同态 $B \to A_g$。于是，再次应用代数，引理
\ref{algebra-lemma-subring-of-local-ring} 即完成证明。
\end{proof}

\begin{lemma}
\label{morphisms-lemma-extend-across}
设 $S$ 是概形。设 $X$、$Y$ 是 $S$ 上的概形。设 $x \in X$。
设 $U \subset X$ 是开集，并设 $f : U \to Y$ 是 $S$ 上的态射。
假设
\begin{enumerate}
\item $x$ 在 $U$ 的闭包中；
\item $X$ 既约且只有有限多个不可约分支，或 $X$ 是 Noether 概形；
\item $\mathcal{O}_{X, x}$ 是赋值环；
\item $Y \to S$ 固有。
\end{enumerate}
则存在开集 $U \subset U' \subset X$，它包含 $x$，以及一个
$S$-态射 $f' : U' \to Y$，它延拓 $f$。
\end{lemma}

\begin{proof}
把 $X$ 替换为一个开邻域是无妨的：该邻域包含 $x$，并位于 $X$ 中
（略去一个小细节）。由性质，引理
\ref{properties-lemma-ring-affine-open-injective-local-ring}，
可假设 $X$ 仿射，且
$\Gamma(X, \mathcal{O}_X) \subset \mathcal{O}_{X, x}$。
特别地，$X$ 是整概形，具有唯一的泛点 $\xi$；其剩余域是分式域
$K$，即赋值环 $\mathcal{O}_{X, x}$ 的分式域。由于 $x$ 位于 $U$ 的闭包中，
可知 $U$ 非空，因而 $U$ 包含 $\xi$。于是，由固有性的赋值判据
（引理 \ref{morphisms-lemma-characterize-proper}），存在态射
$t : \Spec(\mathcal{O}_{X, x}) \to Y$，使其嵌入交换图
$$
\xymatrix{
\Spec(K) \ar[d]_\xi \ar[r] & \Spec(\mathcal{O}_{X, x}) \ar[d]_t \\
U \ar[r]^f & Y
}
$$
这个图由 $S$ 上概形的态射组成。应用引理
\ref{morphisms-lemma-morphism-defined-local-ring}，其中 $y = t(x)$ 且
$\varphi = t^\sharp_x$，可得到开邻域 $V \subset X$，它包含 $x$，
以及态射 $g : V \to Y$；该态射定义在 $S$ 上，并把 $x$ 映到
$y$，并满足 $g^\sharp_x = t^\sharp_x$。由于 $Y \to S$ 分离，
等化子 $E$（对应于 $f|_{U \cap V}$ 与 $g|_{U \cap V}$）是
$U \cap V$ 的闭子概形；见概形，引理
\ref{schemes-lemma-where-are-they-equal}。根据构造，$f$ 和 $g$
决定相同的态射 $\Spec(K) \to Y$，故 $E$ 包含整概形
$U \cap V$ 的泛点。因此 $E = U \cap V$，从而 $f$ 和 $g$
粘合为态射 $U' = U \cup V \to Y$，如所需。
\end{proof}





\section{射影态射}
\label{morphisms-section-projective}

\noindent
我们采用 \cite{EGA} 中射影态射的定义。带字母“H”的定义版本来自
\cite{H}。由此得到的定义并不相同，但两者都有用。

\begin{definition}
\label{morphisms-definition-projective}
设 $f : X \to S$ 是概形的态射。
\begin{enumerate}
\item 称 $f$ {\it 射影}，如果 $X$ 作为 $S$-概形同构于某个射影丛
$\mathbf{P}(\mathcal{E})$ 的闭子概形，其中所用的有限型
$\mathcal{O}_S$-模 $\mathcal{E}$ 是拟凝聚的。
\item 称 $f$ {\it H-射影}，如果存在整数 $n$ 以及闭浸入
$X \to \mathbf{P}^n_S$，后者定义在 $S$ 上。
\item 称 $f$ {\it 局部射影}，如果存在开覆盖 $S = \bigcup U_i$，使得每个
$f^{-1}(U_i) \to U_i$ 都射影。
\end{enumerate}
\end{definition}

\noindent
正如预期，射影态射拟射影；见引理
\ref{morphisms-lemma-projective-quasi-projective}。反之，拟射影态射常常是开浸入
与射影态射的复合；见引理
\ref{morphisms-lemma-quasi-projective-open-projective}。关于拟射影基上射影态射的
性质概览，见态射进阶，第 \ref{more-morphisms-section-projective} 节。

\begin{example}
\label{morphisms-example-projective}
设 $S$ 是概形。设 $\mathcal{A}$ 是拟凝聚分次
$\mathcal{O}_S$-代数，并由 $\mathcal{A}_1$ 在 $\mathcal{A}_0$ 上生成。
再假设 $\mathcal{A}_1$ 在 $\mathcal{O}_S$ 上为有限型。令
$X = \underline{\text{Proj}}_S(\mathcal{A})$。此时 $X \to S$ 射影。
事实上，由分次 $\mathcal{O}_S$-代数同态
$$
\text{Sym}_{\mathcal{O}_X}^*(\mathcal{A}_1)
\longrightarrow
\mathcal{A}
$$
所关联的态射是闭浸入；见构造，引理
\ref{constructions-lemma-surjective-generated-degree-1-map-relative-proj}。
\end{example}
 
\begin{lemma}
\label{morphisms-lemma-H-projective}
H-射影态射是 H-拟射影态射。H-射影态射是射影态射。
\end{lemma}

\begin{proof}
第一个断言由定义立即成立。第二个断言成立，因为
$\mathbf{P}^n_S$ 是 $S$ 上的射影丛；见构造，引理
\ref{constructions-lemma-projective-space-bundle}。
\end{proof}

\begin{lemma}
\label{morphisms-lemma-characterize-locally-projective}
设 $f : X \to S$ 是概形的态射。下列条件等价：
\begin{enumerate}
\item 态射 $f$ 局部射影。
\item 存在开覆盖 $S = \bigcup U_i$，使得每个
$f^{-1}(U_i) \to U_i$ 都 H-射影。
\end{enumerate}
\end{lemma}

\begin{proof}
由引理 \ref{morphisms-lemma-H-projective} 可见 (2) 蕴含 (1)。假设 (1)。
对每个点 $s \in S$，可找到仿射开邻域
$\Spec(R) = U \subset S$，它包含 $s$，使得 $X_U$ 同构于
$\mathbf{P}(\mathcal{E})$ 的某个闭子概形，其中所用的 $\mathcal{O}_U$-模
$\mathcal{E}$ 是有限型且拟凝聚的。写 $\mathcal{E} = \widetilde{M}$，其中
$R$-模 $M$ 是有限型的（见性质，引理
\ref{properties-lemma-finite-type-module}）。选择
$x_0, \ldots, x_n \in M$ 作为 $M$ 的一组 $R$-模生成元。
考虑分次 $R$-代数满射
$$
R[X_0, \ldots, X_n] \longrightarrow \text{Sym}_R(M).
$$
根据构造，引理
\ref{constructions-lemma-surjective-graded-rings-map-proj}，相应的态射
$$
\mathbf{P}(\mathcal{E}) \to \mathbf{P}^n_R
$$
是闭浸入。因此可知 $f^{-1}(U)$ 同构于
$\mathbf{P}^n_U$ 的闭子概形（作为 $U$ 上概形）。换言之，(2) 成立。
\end{proof}

\begin{lemma}
\label{morphisms-lemma-locally-projective-proper}
局部射影态射是固有态射。
\end{lemma}

\begin{proof}
设 $f : X \to S$ 局部射影。为证明 $f$ 固有，可以在基上局部地处理；
见引理 \ref{morphisms-lemma-proper-local-on-the-base}。因此，由上面的引理
\ref{morphisms-lemma-characterize-locally-projective}，可假设存在闭浸入
$X \to \mathbf{P}^n_S$。由引理 \ref{morphisms-lemma-composition-proper} 和
\ref{morphisms-lemma-closed-immersion-proper}，只须证明
$\mathbf{P}^n_S \to S$ 固有。由于 $\mathbf{P}^n_S \to S$ 是
$\mathbf{P}^n_{\mathbf{Z}} \to \Spec(\mathbf{Z})$ 的基变换，
由引理 \ref{morphisms-lemma-base-change-proper}，只须证明
$\mathbf{P}^n_{\mathbf{Z}} \to \Spec(\mathbf{Z})$ 固有。
由构造，引理 \ref{constructions-lemma-proj-separated}，概形
$\mathbf{P}^n_{\mathbf{Z}}$ 分离。由构造，引理
\ref{constructions-lemma-proj-quasi-compact}，概形
$\mathbf{P}^n_{\mathbf{Z}}$ 拟紧。显然
$\mathbf{P}^n_{\mathbf{Z}} \to \Spec(\mathbf{Z})$ 局部为有限型，因为
$\mathbf{P}^n_{\mathbf{Z}}$ 被仿射开集 $D_{+}(X_i)$ 覆盖，每个开集都是
有限型 $\mathbf{Z}$-代数
$$
\mathbf{Z}[X_0/X_i, \ldots, X_n/X_i].
$$
的谱。最后，还须证明
$\mathbf{P}^n_{\mathbf{Z}} \to \Spec(\mathbf{Z})$ 万有闭。这由构造，引理
\ref{constructions-lemma-proj-valuative-criterion} 及赋值判据得出
（见概形，命题
\ref{schemes-proposition-characterize-universally-closed}）。
\end{proof}

\begin{lemma}
\label{morphisms-lemma-proper-ample-locally-projective}
设 $f : X \to S$ 是概形的固有态射。若存在 $f$-丰沛可逆层，
且它定义在 $X$ 上，
则 $f$ 局部射影。
\end{lemma}

\begin{proof}
若存在 $f$-丰沛可逆层，则在 $S$ 上局部地可找到浸入
$i : X \to \mathbf{P}^n_S$；见引理
\ref{morphisms-lemma-finite-type-over-affine-ample-very-ample}。由于 $X \to S$ 固有，
态射 $i$ 是闭浸入；见引理 \ref{morphisms-lemma-image-proper-scheme-closed}。
\end{proof}

\begin{lemma}
\label{morphisms-lemma-H-projective-composition}
H-射影态射的复合是 H-射影态射。
\end{lemma}

\begin{proof}
假设 $X \to Y$ 和 $Y \to Z$ 都 H-射影。于是存在闭浸入
$X \to \mathbf{P}^n_Y$（定义在 $Y$ 上），以及闭浸入
$Y \to \mathbf{P}^m_Z$（定义在 $Z$ 上）。考虑下图
$$
\xymatrix{
X \ar[r] \ar[d] &
\mathbf{P}^n_Y \ar[r] \ar[dl] &
\mathbf{P}^n_{\mathbf{P}^m_Z} \ar[dl] \ar@{=}[r] &
\mathbf{P}^n_Z \times_Z \mathbf{P}^m_Z \ar[r] &
\mathbf{P}^{nm + n + m}_Z \ar[ddllll] \\
Y \ar[r] \ar[d] & \mathbf{P}^m_Z \ar[dl] & \\
Z & &
}
$$
这里，最右侧的顶部水平箭头是 Segre 嵌入；见构造，引理
\ref{constructions-lemma-segre-embedding}。该图把 $X$ 识别为
$\mathbf{P}^{nm + n + m}_Z$ 的闭子概形，如所需。
\end{proof}

\begin{lemma}
\label{morphisms-lemma-H-projective-base-change}
H-射影态射的基变换是 H-射影态射。
\end{lemma}

\begin{proof}
这是因为概形上射影空间的基变换仍是射影空间，并且闭浸入的基变换
仍是闭浸入；见概形，引理 \ref{schemes-lemma-base-change-immersion}。
\end{proof}

\begin{lemma}
\label{morphisms-lemma-base-change-projective}
（局部）射影态射的基变换是（局部）射影态射。
\end{lemma}

\begin{proof}
这是因为概形上射影丛的基变换仍是射影丛，有限型
$\mathcal{O}$-模的拉回仍为有限型
（模，引理 \ref{modules-lemma-pullback-finite-type}），并且闭浸入的
基变换仍是闭浸入；见概形，引理
\ref{schemes-lemma-base-change-immersion}。省略一些细节。
\end{proof}

\begin{lemma}
\label{morphisms-lemma-projective-quasi-projective}
射影态射是拟射影态射。
\end{lemma}

\begin{proof}
设 $f : X \to S$ 是射影态射。选择闭浸入
$i : X \to \mathbf{P}(\mathcal{E})$，其中 $\mathcal{E}$ 是拟凝聚的
有限型 $\mathcal{O}_S$-模。于是
$\mathcal{L} = i^*\mathcal{O}_{\mathbf{P}(\mathcal{E})}(1)$ 是
$f$-甚丰沛的。由于 $f$ 固有（引理
\ref{morphisms-lemma-locally-projective-proper}），它拟紧。因此，引理
\ref{morphisms-lemma-ample-very-ample} 蕴含 $\mathcal{L}$ 是 $f$-丰沛的。
由于 $f$ 固有，它为有限型。这样就验证了拟射影的全部定义性质，
结论成立。
\end{proof}

\begin{lemma}
\label{morphisms-lemma-H-quasi-projective-open-H-projective}
设 $f : X \to S$ 是 H-拟射影态射。则 $f$ 分解为
$X \to X' \to S$，其中 $X \to X'$ 是开浸入，而
$X' \to S$ 是 H-射影态射。
\end{lemma}

\begin{proof}
根据定义，可把 $f$ 分解为拟紧浸入
$i : X \to \mathbf{P}^n_S$，再复合投影 $\mathbf{P}^n_S \to S$。
由引理 \ref{morphisms-lemma-quasi-compact-immersion}，存在闭子概形
$X' \subset \mathbf{P}^n_S$，使得 $i$ 通过开浸入
$X \to X'$ 分解。引理得证。
\end{proof}

\begin{lemma}
\label{morphisms-lemma-quasi-projective-open-projective}
设 $f : X \to S$ 是拟射影态射，且 $S$ 拟紧并拟分离。
则 $f$ 分解为 $X \to X' \to S$，其中 $X \to X'$ 是开浸入，
而 $X' \to S$ 是射影态射。
\end{lemma}

\begin{proof}
设 $\mathcal{L}$ 为 $f$-丰沛层。由于 $f$ 为有限型且 $S$ 拟紧，
$\mathcal{L}^{\otimes n}$ 为 $f$-甚丰沛，其中 $n > 0$；见引理
\ref{morphisms-lemma-finite-type-ample-very-ample}。把
$\mathcal{L}$ 替换为 $\mathcal{L}^{\otimes n}$。写
$\mathcal{F} = f_*\mathcal{L}$。这是拟凝聚 $\mathcal{O}_S$-模；见概形，
引理 \ref{schemes-lemma-push-forward-quasi-coherent}
（拟射影态射拟紧且分离；见引理
\ref{morphisms-lemma-quasi-projective-properties}）。由性质，引理
\ref{properties-lemma-directed-colimit-finite-presentation}，可找到有向集
$I$ 及其上的有限型拟凝聚 $\mathcal{O}_S$-模系统
$\mathcal{E}_i$，使得该系统以 $I$ 为指标且
$\mathcal{F} = \colim \mathcal{E}_i$。考虑复合
$\psi_i : f^*\mathcal{E}_i \to f^*\mathcal{F} \to \mathcal{L}$。
选择有限仿射开覆盖 $S = \bigcup_{j = 1, \ldots, m} V_j$。
对每个 $j$，可选择截面
$$
s_{j, 0}, \ldots, s_{j, n_j} \in
\Gamma(f^{-1}(V_j), \mathcal{L}) = f_*\mathcal{L}(V_j) = \mathcal{F}(V_j)
$$
它们生成 $\mathcal{L}$ 在 $f^{-1}V_j$ 上的限制，并定义浸入
$$
f^{-1}V_j \longrightarrow \mathbf{P}^{n_j}_{V_j},
$$
见引理 \ref{morphisms-lemma-very-ample-finite-type-over-affine}。选择 $i$，使得存在
截面 $e_{j, t} \in \mathcal{E}_i(V_j)$，它们映到 $s_{j, t}$，
所取像位于 $\mathcal{F}$ 中；这一点对所有 $j = 1, \ldots, m$ 和
$t = 1, \ldots, n_j$ 都成立。于是同态 $\psi_i$ 满射，因为截面
$f^*e_{j, t}$ 与截面 $s_{j, t}$ 有相同的像，而后者生成
$\mathcal{L}|_{f^{-1}V_j}$。因此得到 $S$ 上的态射
$$
r_{\mathcal{L}, \psi_i} : X \longrightarrow \mathbf{P}(\mathcal{E}_i)
$$
使得在 $V_j$ 上，上述浸入分解为
$$
f^{-1}V_j \to \mathbf{P}(\mathcal{E}_i)|_{V_j} \to \mathbf{P}^{n_j}_{V_j}
$$
由引理 \ref{morphisms-lemma-immersion-permanence}，可知
$r_{\mathcal{L}, \psi_i}|_{V_j}$ 是浸入。由于 $S = \bigcup V_j$，
可知 $r_{\mathcal{L}, \psi_i}$ 是浸入。注意，
$r_{\mathcal{L}, \psi_i}$ 拟紧，因为 $X \to S$ 拟紧，且
$\mathbf{P}(\mathcal{E}_i) \to S$ 分离
（见概形，引理 \ref{schemes-lemma-quasi-compact-permanence}）。
由引理 \ref{morphisms-lemma-quasi-compact-immersion}，存在闭子概形
$X' \subset \mathbf{P}(\mathcal{E}_i)$，使得 $i$ 通过开浸入
$X \to X'$ 分解。于是 $X' \to S$ 根据定义射影，结论成立。
\end{proof}

\begin{lemma}
\label{morphisms-lemma-projective-is-quasi-projective-proper}
设 $S$ 是拟紧且拟分离的概形。设 $f : X \to S$ 是概形的态射。则
\begin{enumerate}
\item $f$ 射影，当且仅当 $f$ 拟射影且固有；并且
\item $f$ H-射影，当且仅当 $f$ H-拟射影且固有。
\end{enumerate}
\end{lemma}

\begin{proof}
若 $f$ 射影，则由引理 \ref{morphisms-lemma-projective-quasi-projective}，$f$ 拟射影，
并由引理 \ref{morphisms-lemma-locally-projective-proper}，它固有。反之，若
$X \to S$ 拟射影且固有，则由引理
\ref{morphisms-lemma-quasi-projective-open-projective}，可选择开浸入
$X \to X'$，其中 $X' \to S$ 射影。由于 $X \to S$ 固有，
可知 $X$ 在 $X'$ 中闭（引理
\ref{morphisms-lemma-image-proper-scheme-closed}），即 $X \to X'$ 是（开且）闭浸入。
由于 $X'$ 同构于 $S$ 上某个射影丛的闭子概形
（定义 \ref{morphisms-definition-projective}），可知 $X$ 也如此，即
$X \to S$ 是射影态射。这证明了 (1)。对 (2) 的证明相同，只须使用引理
\ref{morphisms-lemma-H-projective} 和
\ref{morphisms-lemma-H-quasi-projective-open-H-projective}。
\end{proof}

\begin{lemma}
\label{morphisms-lemma-composition-projective}
设 $f : X \to Y$ 和 $g : Y \to S$ 是概形的态射。若 $S$ 拟紧且拟分离，
并且 $f$ 和 $g$ 都射影，则 $g \circ f$ 射影。
\end{lemma}

\begin{proof}
由引理 \ref{morphisms-lemma-projective-quasi-projective} 和
\ref{morphisms-lemma-locally-projective-proper}，可见 $f$ 和 $g$ 拟射影且固有。
由引理 \ref{morphisms-lemma-composition-proper} 和
\ref{morphisms-lemma-composition-quasi-projective}，可见 $g \circ f$ 固有且拟射影。
因此，由引理 \ref{morphisms-lemma-projective-is-quasi-projective-proper}，
$g \circ f$ 射影。
\end{proof}

\begin{lemma}
\label{morphisms-lemma-projective-permanence}
设 $g : Y \to S$ 和 $f : X \to Y$ 是概形的态射。若 $g \circ f$ 射影，
且 $g$ 分离，则 $f$ 射影。
\end{lemma}

\begin{proof}
选择闭浸入 $X \to \mathbf{P}(\mathcal{E})$，其中 $\mathcal{E}$ 是拟凝聚的
有限型 $\mathcal{O}_S$-模。于是得到态射
$X \to \mathbf{P}(\mathcal{E}) \times_S Y$。该态射是闭浸入，因为它是复合
$$
X \to X \times_S Y \to \mathbf{P}(\mathcal{E}) \times_S Y
$$
其中第一个态射由概形，引理 \ref{schemes-lemma-semi-diagonal}
（以及 $g$ 分离这一事实）是闭浸入，第二个态射则是闭浸入的基变换。
最后，纤维积 $\mathbf{P}(\mathcal{E}) \times_S Y$ 同构于
$\mathbf{P}(g^*\mathcal{E})$，而拉回保持拟凝聚有限型模。
\end{proof}

\begin{lemma}
\label{morphisms-lemma-projective-over-quasi-projective-is-H-projective}
设 $S$ 是具有丰沛可逆层的概形。则
\begin{enumerate}
\item 任意射影态射 $X \to S$ 都 H-射影；并且
\item 任意拟射影态射 $X \to S$ 都 H-拟射影。
\end{enumerate}
\end{lemma}

\begin{proof}
关于 $S$ 的假设蕴含 $S$ 拟紧且分离；见性质，定义
\ref{properties-definition-ample}、引理
\ref{properties-lemma-ample-immersion-into-proj} 以及构造，引理
\ref{constructions-lemma-proj-separated}。因此，引理
\ref{morphisms-lemma-quasi-projective-open-projective} 适用，并可见 (1) 蕴含 (2)。
设 $\mathcal{E}$ 是有限型拟凝聚 $\mathcal{O}_S$-模。根据射影态射的定义，
只须证明 $\mathbf{P}(\mathcal{E}) \to S$ H-射影。若 $\mathcal{E}$
由有限多个全局截面生成，则相应的满射
$\mathcal{O}_S^{\oplus n} \to \mathcal{E}$ 诱导闭浸入
$$
\mathbf{P}(\mathcal{E}) \longrightarrow
\mathbf{P}(\mathcal{O}_S^{\oplus n}) = \mathbf{P}^{n - 1}_S
$$
如所需。一般地，设 $\mathcal{L}$ 是 $S$ 上的丰沛可逆层。由性质，命题
\ref{properties-proposition-characterize-ample}，存在整数 $n$，使得
$\mathcal{E} \otimes_{\mathcal{O}_S} \mathcal{L}^{\otimes n}$
由有限多个全局截面生成。由于
$\mathbf{P}(\mathcal{E}) =
\mathbf{P}(\mathcal{E} \otimes_{\mathcal{O}_S} \mathcal{L}^{\otimes n})$；见
构造，引理 \ref{constructions-lemma-twisting-and-proj}，证明完成。
\end{proof}

\begin{lemma}
\label{morphisms-lemma-proper-ample-is-proj}
设 $f : X \to S$ 是万有闭态射。设 $\mathcal{L}$ 是 $f$-丰沛可逆
$\mathcal{O}_X$-模。则引理 \ref{morphisms-lemma-characterize-relatively-ample} 中的
典范态射
$$
r : X
\longrightarrow
\underline{\text{Proj}}_S
\left(
\bigoplus\nolimits_{d \geq 0} f_*\mathcal{L}^{\otimes d}
\right)
$$
是同构。
\end{lemma}

\begin{proof}
注意 $f$ 拟紧，因为存在 $f$-丰沛可逆模这一事实迫使 $f$ 拟紧。
由所引引理，态射 $r$ 是开浸入。另一方面，由引理
\ref{morphisms-lemma-image-proper-scheme-closed}，$r$ 的像是闭的
（由构造，引理 \ref{constructions-lemma-relative-proj-separated}，
$r$ 的目标在 $S$ 上分离）。最后，由性质，引理
\ref{properties-lemma-ample-immersion-into-proj}，$r$ 的像稠密
（这里还使用了引理 \ref{morphisms-lemma-characterize-relatively-ample} 的证明中
已经表明的事实：态射 $r$ 在 $S$ 的仿射开集上由性质，引理
\ref{properties-lemma-map-into-proj} 的典范态射给出）。因此可知
$r$ 是满射开浸入，即同构。
\end{proof}

\begin{lemma}
\label{morphisms-lemma-proper-ample-delete-affine}
设 $f : X \to S$ 是万有闭态射。设 $\mathcal{L}$ 是 $f$-丰沛可逆
$\mathcal{O}_X$-模。设 $s \in \Gamma(X, \mathcal{L})$。则
$X_s \to S$ 是仿射态射。
\end{lemma}

\begin{proof}
该问题在 $S$ 上是局部的（引理 \ref{morphisms-lemma-characterize-affine}），
故可假设 $S$ 仿射。由引理 \ref{morphisms-lemma-proper-ample-is-proj}，可写
$X = \text{Proj}(A)$，其中 $A$ 是分次环，而 $s$ 对应于
$f \in A_1$，并且 $X_s = D_+(f)$（性质，引理
\ref{properties-lemma-map-into-proj}）。根据 $\text{Proj}(A)$ 的构造，
这证明了引理；见构造，第 \ref{constructions-section-proj} 节。
\end{proof}

\section{整态射与有限态射}
\label{morphisms-section-integral}

\noindent
回顾：若环同态 $R \to A$ 的每个 $A$ 中元素都满足一个系数在 $R$ 中的
首一方程，则称该同态{\it 整}。再回顾：若对环同态 $R \to A$，
$A$ 作为 $R$-模有限，则称该同态{\it 有限}。见代数，定义
\ref{algebra-definition-integral-ring-map}。

\begin{definition}
\label{morphisms-definition-integral}
设 $f : X \to S$ 是概形的态射。
\begin{enumerate}
\item 称 $f$ {\it 整}，如果 $f$ 仿射，并且对每个仿射开集
$\Spec(R) = V \subset S$，其逆像为
$\Spec(A) = f^{-1}(V) \subset X$ 时，相应的环同态 $R \to A$ 整。
\item 称 $f$ {\it 有限}，如果 $f$ 仿射，并且对每个仿射开集
$\Spec(R) = V \subset S$，其逆像为
$\Spec(A) = f^{-1}(V) \subset X$ 时，相应的环同态 $R \to A$ 有限。
\end{enumerate}
\end{definition}

\noindent
显然，整态射和有限态射分离且拟紧。有限态射显然也是有限型态射。
本节的大多数引理都是完全标准的。不过，请注意本节末尾有趣的引理
\ref{morphisms-lemma-integral-universally-closed}。

\begin{lemma}
\label{morphisms-lemma-integral-local}
设 $f : X \to S$ 是概形的态射。下列条件等价：
\begin{enumerate}
\item 态射 $f$ 整。
\item 存在仿射开覆盖 $S = \bigcup U_i$，使得每个
$f^{-1}(U_i)$ 仿射，并且
$\mathcal{O}_S(U_i) \to \mathcal{O}_X(f^{-1}(U_i))$ 整。
\item 存在开覆盖 $S = \bigcup U_i$，使得每个
$f^{-1}(U_i) \to U_i$ 都整。
\end{enumerate}
此外，若 $f$ 整，则对每个开子概形 $U \subset S$，态射
$f : f^{-1}(U) \to U$ 都整。
\end{lemma}

\begin{proof}
见代数，引理 \ref{algebra-lemma-integral-local}。省略一些细节。
\end{proof}

\begin{lemma}
\label{morphisms-lemma-finite-local}
设 $f : X \to S$ 是概形的态射。下列条件等价：
\begin{enumerate}
\item 态射 $f$ 有限。
\item 存在仿射开覆盖 $S = \bigcup U_i$，使得每个
$f^{-1}(U_i)$ 仿射，并且
$\mathcal{O}_S(U_i) \to \mathcal{O}_X(f^{-1}(U_i))$ 有限。
\item 存在开覆盖 $S = \bigcup U_i$，使得每个
$f^{-1}(U_i) \to U_i$ 都有限。
\end{enumerate}
此外，若 $f$ 有限，则对每个开子概形 $U \subset S$，态射
$f : f^{-1}(U) \to U$ 都有限。
\end{lemma}

\begin{proof}
见代数，引理 \ref{algebra-lemma-integral-local}。省略一些细节。
\end{proof}

\begin{lemma}
\label{morphisms-lemma-finite-integral}
有限态射是整态射。局部为有限型的整态射是有限态射。
\end{lemma}

\begin{proof}
见代数，引理 \ref{algebra-lemma-finite-is-integral} 和引理
\ref{algebra-lemma-characterize-finite-in-terms-of-integral}。
\end{proof}

\begin{lemma}
\label{morphisms-lemma-composition-finite}
有限态射的复合是有限态射。整态射也同样如此。
\end{lemma}

\begin{proof}
见代数，引理 \ref{algebra-lemma-finite-transitive} 和
\ref{algebra-lemma-integral-transitive}。
\end{proof}

\begin{lemma}
\label{morphisms-lemma-base-change-finite}
有限态射的基变换是有限态射。整态射也同样如此。
\end{lemma}

\begin{proof}
见代数，引理 \ref{algebra-lemma-base-change-integral}。
\end{proof}

\begin{lemma}
\label{morphisms-lemma-integral-universally-closed}
\begin{slogan}
整 $=$ 仿射 $+$ 万有闭
\end{slogan}
设 $f : X \to S$ 是概形的态射。下列条件等价：
\begin{enumerate}
\item $f$ 整；
\item $f$ 仿射且万有闭。
\end{enumerate}
\end{lemma}

\begin{proof}
设 (1) 成立。整态射按定义是仿射的。整态射的基变换仍整，故为证明 (2)，
只须证明整态射闭。这由代数，引理
\ref{algebra-lemma-integral-going-up} 和
\ref{algebra-lemma-going-up-closed} 得出。

\medskip\noindent
设 (2) 成立。可假设 $f$ 是态射 $f : \Spec(A) \to \Spec(R)$，
它由环同态 $R \to A$ 给出。取 $a$ 为 $A$ 的元素。我们须证明
$a$ 在 $R$ 上整，即在核 $I$ 中存在首一多项式；这里的同态是
$R[x] \to A$，它把 $x$ 映到 $a$。考虑环
$B = A[x]/(ax -1)$，并令 $J$ 为复合
$R[x]\to A[x] \to B$ 的核。若 $f\in J$，则存在 $q\in A[x]$，使得
$f = (ax-1)q$ 在 $A[x]$ 中成立。因此，若
$f = \sum_i f_ix^i$ 且 $q = \sum_iq_ix^i$，则对所有 $i \geq 0$，有
$f_i = aq_{i-1} - q_i$。对 $n \geq \deg q + 1$，多项式
$$
\sum\nolimits_{i \geq 0} f_i x^{n - i} =
\sum\nolimits_{i \geq 0} (a q_{i - 1} - q_i) x^{n - i} =
(a - x) \sum\nolimits_{i \geq 0} q_i x^{n - i - 1}
$$
显然属于 $I$；若 $f_0 = 1$，该多项式还首一。因此问题归结为证明
$J$ 含有常数项为 $1$ 的多项式。为此，我们证明
$\Spec(R[x]/(J + (x))$ 为空。


\medskip\noindent
由于 $f$ 万有闭，基变换 $\Spec(A[x]) \to \Spec(R[x])$ 闭。因此闭子集
$\Spec(B) \subset \Spec(A[x])$ 的像是闭子集
$\Spec(R[x]/J) \subset \Spec(R[x])$；见例
\ref{morphisms-example-scheme-theoretic-image} 和引理
\ref{morphisms-lemma-quasi-compact-scheme-theoretic-image}。特别地，
$\Spec(B) \to \Spec(R[x]/J)$ 满射。考虑下图，其中每个方块都是拉回方块：
$$
\xymatrix{
\Spec(B) \ar@{->>}[r]^g &
\Spec(R[x]/J)  \ar[r] &
\Spec(R[x])\\
\emptyset \ar[u] \ar[r] &
\Spec(R[x]/(J + (x)))\ar[u] \ar[r] &
\Spec(R) \ar[u]^0
}
$$
左下角为空，因为它是 $R\otimes_{R[x]} B$ 的谱，其中同态
$R[x]\to B$ 把 $x$ 映到一个可逆元，而同态 $R[x]\to R$ 把 $x$ 映到
$0$。由于 $g$ 满射，这蕴含 $\Spec(R[x]/(J + (x)))$ 为空，正如所需。
\end{proof}

\begin{lemma}
\label{morphisms-lemma-integral-fibres}
设 $f : X \to S$ 是整态射。则 $X$ 的每个点在其纤维中都是闭点。
\end{lemma}

\begin{proof}
见代数，引理 \ref{algebra-lemma-integral-no-inclusion}。
\end{proof}

\begin{lemma}
\label{morphisms-lemma-integral-dimension}
设 $f : X \to Y$ 是整态射。则 $\dim(X) \leq \dim(Y)$。
若 $f$ 满射，则 $\dim(X) = \dim(Y)$。
\end{lemma}

\begin{proof}
由于 $X$ 和 $Y$ 的维数是仿射开覆盖中各成员维数的上确界，可假设
$Y$ 和 $X$ 仿射。不等式由代数，引理
\ref{algebra-lemma-integral-dim-up} 得出。等式继而由代数，引理
\ref{algebra-lemma-dimension-going-up} 和
\ref{algebra-lemma-integral-going-up} 得出。
\end{proof}

\begin{lemma}
\label{morphisms-lemma-finite-quasi-finite}
有限态射是拟有限态射。
\end{lemma}

\begin{proof}
这由代数，引理 \ref{algebra-lemma-quasi-finite} 和引理
\ref{morphisms-lemma-quasi-finite-locally-quasi-compact} 推出。另一种证明是：由引理
\ref{morphisms-lemma-integral-fibres}（以及有限态射整这一事实），纤维中的所有点
都是闭点；再用引理 \ref{morphisms-lemma-quasi-finite-at-point-characterize} (3)，
可见 $f$ 在 $x$ 处拟有限，对所有 $x \in X$ 均如此。
\end{proof}

\begin{lemma}
\label{morphisms-lemma-finite-proper}
设 $f : X \to S$ 是概形的态射。下列条件等价：
\begin{enumerate}
\item $f$ 有限；
\item $f$ 仿射且固有。
\end{enumerate}
\end{lemma}

\begin{proof}
这形式地由引理 \ref{morphisms-lemma-integral-universally-closed}、有限态射整且分离这一
事实、固有态射等价于有限型、分离且万有闭的态射这一事实，以及有限型整态射
有限这一事实（引理 \ref{morphisms-lemma-finite-integral}）得出。
\end{proof}

\begin{lemma}
\label{morphisms-lemma-closed-immersion-finite}
闭浸入有限（从而尤其整）。
\end{lemma}

\begin{proof}
这是因为闭浸入仿射（引理 \ref{morphisms-lemma-closed-immersion-affine}），
而满射环同态有限且整。
\end{proof}

\begin{lemma}
\label{morphisms-lemma-finite-union-finite}
设 $X_i \to Y$，$i = 1, \ldots, n$ 是概形的有限态射。则
$X_1 \amalg \ldots \amalg X_n \to Y$ 也有限。
\end{lemma}

\begin{proof}
这由如下代数事实得出：若 $R \to A_i$，$i = 1, \ldots, n$ 是有限环同态，
则 $R \to A_1 \times \ldots \times A_n$ 也有限。
\end{proof}

\begin{lemma}
\label{morphisms-lemma-finite-permanence}
设 $f : X \to Y$ 和 $g : Y \to Z$ 是态射。
\begin{enumerate}
\item 若 $g \circ f$ 有限且 $g$ 分离，则 $f$ 有限。
\item 若 $g \circ f$ 整且 $g$ 分离，则 $f$ 整。
\end{enumerate}
\end{lemma}

\begin{proof}
设 $g \circ f$ 有限（相应地，整），且 $g$ 分离。由引理
\ref{morphisms-lemma-base-change-finite}，基变换 $X \times_Z Y \to Y$
有限（相应地，整）。态射 $X \to X \times_Z Y$ 是闭浸入，
因为 $Y \to Z$ 分离；见概形，引理
\ref{schemes-lemma-section-immersion}。闭浸入有限（相应地，整）；见引理
\ref{morphisms-lemma-closed-immersion-finite}。有限（相应地，整）态射的复合仍有限
（相应地，整）；见引理 \ref{morphisms-lemma-composition-finite}。证毕。
\end{proof}

\begin{lemma}
\label{morphisms-lemma-finite-monomorphism-closed}
设 $f : X \to Y$ 是概形的态射。若 $f$ 有限且为单态射，则 $f$ 是闭浸入。
\end{lemma}

\begin{proof}
这归结为代数，引理
\ref{algebra-lemma-finite-epimorphism-surjective}。
\end{proof}

\begin{lemma}
\label{morphisms-lemma-finite-projective}
有限态射是射影态射。
\end{lemma}

\begin{proof}
设 $f : X \to S$ 是有限态射。则 $f_*\mathcal{O}_X$ 是拟凝聚
$\mathcal{O}_S$-模（引理 \ref{morphisms-lemma-affine-equivalence-algebras}），
并且是有限型的（由有限态射的定义和性质，引理
\ref{properties-lemma-finite-type-module}）。我们断言存在 $S$ 上的闭浸入
$$
\sigma :
X
\longrightarrow
\mathbf{P}(f_*\mathcal{O}_X) =
\underline{\text{Proj}}_S(\text{Sym}^*_{\mathcal{O}_S}(f_*\mathcal{O}_X))
$$
这就完成证明。具体地，令 $\sigma$ 为（通过构造，引理
\ref{constructions-lemma-apply-relative}）对应于满射
$$
f^*f_*\mathcal{O}_X \longrightarrow \mathcal{O}_X
$$
的态射；该满射来自伴随映射 $f^*f_* \to \text{id}$。那么 $\sigma$
是闭浸入。事实上，在 $S$ 上仿射局部地，可写
$X = \Spec(A)$ 和 $S = \Spec(R)$。由于 $X$ 在 $S$ 上有限，
可选取 $a_1, \ldots, a_n \in A$，使其生成 $A$ 作为 $R$-模。于是
$R$-代数同态 $R[T_0, T_1, \ldots, T_n] \to \text{Sym}_R^*(A)$，
它把 $T_0$ 映到 $1$，并把 $T_i$ 映到 $a_i$（$1 \leq i \leq n$），
是满射。因此，由构造，引理
\ref{constructions-lemma-surjective-graded-rings-map-proj}，
$\mathbf{P}(f_*\mathcal{O}_X)$ 是 $\mathbf{P}^n_S$ 的闭子概形。
所以只须证明诱导态射 $X \to \mathbf{P}^n_S$ 是闭浸入
（例如见引理 \ref{morphisms-lemma-closed-immersion-ideals}）。读者可检验
$X \to \mathbf{P}^n_S$ 的像包含在开集
$D_+(T_0) \cong \mathbf{A}^n_S$ 中，并且 $X \to \mathbf{A}^n_S$
对应于满射 $R$-代数同态
$R[x_1, \ldots, x_n] \to \text{Sym}_R^*(A)$，它把 $x_i$ 映到 $a_i$。
因此 $X \to \mathbf{P}^n_S$ 是浸入（作为闭浸入与开浸入的复合）。由于
$X$ 在 $S$ 上有限，该浸入的像闭（这里使用引理
\ref{morphisms-lemma-finite-proper}、引理
\ref{morphisms-lemma-image-proper-scheme-closed} 和构造，引理
\ref{constructions-lemma-projective-space-separated}）。最后由概形，引理
\ref{schemes-lemma-immersion-when-closed} 得出结论。
\end{proof}

\section{泛同胚}
\label{morphisms-section-universal-homeomorphisms}

\noindent
下面的定义其实是多余的，因为泛同胚恰好就是整、泛单射且满射的态射；
见引理 \ref{morphisms-lemma-universal-homeomorphism}。

\begin{definition}
\label{morphisms-definition-universal-homeomorphism}
若概形态射 $f : X \to Y$ 的基变换
$f' : Y' \times_Y X \to Y'$ 对每个态射 $Y' \to Y$ 都是同胚，
则称它为{\it 泛同胚}。
\end{definition}

\noindent
先陈述几个必备引理。

\begin{lemma}
\label{morphisms-lemma-base-change-universal-homeomorphism}
概形的泛同胚沿任意概形态射作基变换后仍为泛同胚。
\end{lemma}

\begin{proof}
这由定义立即得出。
\end{proof}

\begin{lemma}
\label{morphisms-lemma-composition-universal-homeomorphism}
两个概形泛同胚的复合仍为泛同胚。
\end{lemma}

\begin{proof}
略。
\end{proof}

\noindent
下面的简单引理是刻画泛同胚的关键。

\begin{lemma}
\label{morphisms-lemma-homeomorphism-affine}
设 $f : X \to Y$ 是概形的态射。若 $f$ 是到 $Y$ 的某个闭子集上的同胚，
则 $f$ 仿射。
\end{lemma}

\begin{proof}
取点 $y \in Y$。若 $y \not \in f(X)$，则存在 $y$ 的仿射邻域与
$f(X)$ 不交。若 $y \in f(X)$，令 $x \in X$ 为 $X$ 中唯一映到 $y$
的点。取仿射开邻域 $y \in V$。取 $U \subset X$ 为 $x$ 的仿射开邻域，
使其映入 $V$。$f(U) \subset V \cap f(X)$ 由关于 $f$ 的假设
在诱导拓扑中开。因此可选取 $h \in \Gamma(V, \mathcal{O}_Y)$，使得
$y \in D(h)$ 且 $D(h) \cap f(X) \subset f(U)$。以
$h' \in \Gamma(U, \mathcal{O}_X)$ 表示 $f^\sharp(h)$ 在 $U$ 上的限制。
于是可见 $D(h') \subset U$ 等于 $f^{-1}(D(h))$。换言之，$Y$ 的每个点
都有逆像为仿射概形的开邻域。因此 $f$ 仿射；见引理
\ref{morphisms-lemma-characterize-affine}。
\end{proof}

\begin{lemma}
\label{morphisms-lemma-universal-homeomorphism}
设 $f : X \to Y$ 是概形的态射。下列条件等价：
\begin{enumerate}
\item $f$ 是泛同胚；
\item $f$ 整、泛单射且满射。
\end{enumerate}
\end{lemma}

\begin{proof}
设 $f$ 是泛同胚。由引理 \ref{morphisms-lemma-homeomorphism-affine}，可见 $f$
仿射。由于 $f$ 显然万有闭，引理
\ref{morphisms-lemma-integral-universally-closed} 表明 $f$ 整。还显然，$f$
泛单射且满射。

\medskip\noindent
设 $f$ 整、泛单射且满射。由引理
\ref{morphisms-lemma-integral-universally-closed}，$f$ 万有闭。它还泛双射（见引理
\ref{morphisms-lemma-base-change-surjective}），故它是泛同胚。
\end{proof}

\begin{lemma}
\label{morphisms-lemma-reduction-universal-homeomorphism}
设 $X$ 是概形。典范闭浸入 $X_{red} \to X$（见概形，定义
\ref{schemes-definition-reduced-induced-scheme}）是泛同胚。
\end{lemma}

\begin{proof}
略。
\end{proof}

\begin{lemma}
\label{morphisms-lemma-check-closed-infinitesimally}
设 $f : X \to S$ 和 $S' \to S$ 是概形的态射。假设
\begin{enumerate}
\item $S' \to S$ 是闭浸入；
\item $S' \to S$ 在点上双射；
\item $X \times_S S' \to S'$ 是闭浸入；
\item $X \to S$ 是有限型态射，或者 $S' \to S$ 是有限表示态射。
\end{enumerate}
则 $f : X \to S$ 是闭浸入。
\end{lemma}

\begin{proof}
假设 (1) 和 (2) 蕴含 $S' \to S$ 是泛同胚（例如，因为
$S_{red} = S'_{red}$，并使用引理
\ref{morphisms-lemma-reduction-universal-homeomorphism}）。因此 (3) 蕴含
$X \to S$ 是到 $S$ 的某个闭子集上的同胚。于是由引理
\ref{morphisms-lemma-homeomorphism-affine}，$X \to S$ 仿射。取仿射开集
$U \subset S$，记作 $U = \Spec(A)$。则由 (1)，
$S' = \Spec(A/I)$；并且由 (2)，$I$ 是局部幂零理想。由于 $f$ 仿射，
可写 $f^{-1}(U) = \Spec(B)$。假设 (4) 告诉我们，$B$ 是有限型
$A$-代数（引理 \ref{morphisms-lemma-locally-finite-type-characterize}），或者
$I$ 有限生成（引理 \ref{morphisms-lemma-closed-immersion-finite-presentation}）。
假设 (3) 即 $A/I \to B/IB$ 满射。若 $A \to B$ 为有限型，则由代数，
引理 \ref{algebra-lemma-surjective-mod-locally-nilpotent}；若 $I$ 有限生成
从而幂零，则由代数，引理 \ref{algebra-lemma-NAK}；均可推出
$A \to B$ 满射。这意味着 $f$ 是闭浸入；见引理
\ref{morphisms-lemma-closed-immersion}。
\end{proof}

\begin{lemma}
\label{morphisms-lemma-permanence-universal-homeomorphism}
设 $f : X \to Z$ 是两个态射 $g : X \to Y$ 和 $h : Y \to Z$ 的复合。
若态射 $\{f, g, h\}$ 中任意两个是泛同胚，则第三个也是泛同胚。
\end{lemma}

\begin{proof}
若 $g$ 和 $h$ 都是泛同胚，则由引理
\ref{morphisms-lemma-composition-universal-homeomorphism}，$f$ 也是泛同胚。

\medskip\noindent
设 $f$ 和 $g$ 都是泛同胚。我们要证明 $h$ 也是。沿概形的任意态射
$\alpha : Z' \to Z$ 对该图作基变换，得到下图，其中所有方块均笛卡尔：
$$
\xymatrix{
X' \ar[r]^{g'} \ar[d] & Y' \ar[r]^{h'} \ar[d] & Z' \ar[d] \\
X \ar[r]^{g} & Y \ar[r]^{h} & Z.
}
$$
我们的假设蕴含复合 $f'= h' \circ g' : X' \to Z'$ 和
$g' : X' \to Y'$ 都是同胚，因此 $h'$ 也是同胚。这就证明了 $h$
是泛同胚。

\medskip\noindent
最后，设 $f$ 和 $h$ 是泛同胚。我们要证明 $g$ 是泛同胚。
取概形的任意态射 $\beta : Y' \to Y$。得到下图，其中所有方块均笛卡尔：
$$
\xymatrix{
X' \ar[r]^{g'} \ar[d] & Y' \ar[d]^{\gamma} & \\
X'' \ar[r]^{g''} \ar[d] & Y'' \ar[r]^{h''} \ar[d] & Y' \ar[d]^{h \circ \beta} \\
X \ar[r]^{g} & Y \ar[r]^{h} & Z.
}
$$
这里，态射 $\gamma : Y' \to Y''$ 由纤维积的泛性质以及两个态射
$id_{Y'} : Y' \to Y'$ 和 $\beta : Y' \to Y$ 定义。我们将证明 $g'$
是同胚。由于成为同胚这一性质满足三取二性质，可见 $g''$ 是同胚。
观察上方方块，只须证明 $\gamma$ 是泛同胚。由于 $h''$ 是同胚，引理
\ref{morphisms-lemma-homeomorphism-affine} 表明它是仿射态射，从而尤其分离
（引理 \ref{morphisms-lemma-affine-separated}）。由于 $h'' \circ \gamma$ 是恒等态射，
概形，引理 \ref{schemes-lemma-section-immersion} 表明 $\gamma$ 是闭浸入。
由于 $h''$ 双射，可见 $\gamma$ 是双射闭浸入，因而是泛同胚
（例如由引理 \ref{morphisms-lemma-universal-homeomorphism} 的刻画），如所需。
\end{proof}

\section{仿射概形的泛同胚}
\label{morphisms-section-universal-homeomorphisms-affine}

\noindent
本节刻画仿射概形的泛同胚。

\begin{lemma}
\label{morphisms-lemma-subalgebra-inherits}
设 $A \to B$ 是环同态，使得诱导的概形态射
$f : \Spec(B) \to \Spec(A)$ 是泛同胚；相应地，是诱导剩余域同构的
泛同胚；相应地，万有闭；相应地，万有闭且泛单射。则对任意
$A$-子代数 $B' \subset B$，$f' : \Spec(B') \to \Spec(A)$ 也具有
相应性质。
\end{lemma}

\begin{proof}
若 $f$ 万有闭，则由引理 \ref{morphisms-lemma-integral-universally-closed}，
$B$ 在 $A$ 上整。因此 $B'$ 在 $A$ 上整，且 $f'$ 万有闭
（由同一引理）。这证明了 $f$ 万有闭的情形。

\medskip\noindent
继续可见 $B$ 在 $B'$ 上整（代数，引理
\ref{algebra-lemma-integral-permanence}），这蕴含
$\Spec(B) \to \Spec(B')$ 满射（代数，引理
\ref{algebra-lemma-integral-overring-surjective}）。因此，若
$A \to B$ 诱导剩余域的纯不可分扩张，则 $A \to B'$ 也如此。
这证明了 $f$ 万有闭且泛单射的情形；见引理
\ref{morphisms-lemma-universally-injective}。

\medskip\noindent
$f$ 为泛同胚的情形由上述论述、引理
\ref{morphisms-lemma-universal-homeomorphism} 以及如下显然观察得出：若 $f$
满射，则 $f'$ 也满射。

\medskip\noindent
若 $A \to B$ 诱导剩余域同构，则 $A \to B'$ 也如此
（见第二段的论证）。这样便证明了剩余情形。
\end{proof}

\begin{lemma}
\label{morphisms-lemma-colimit-inherits}
设 $A$ 是环。设 $B = \colim B_\lambda$ 是 $A$-代数的滤过余极限。
若每个 $f_\lambda : \Spec(B_\lambda) \to \Spec(A)$ 都是泛同胚；
相应地，是诱导剩余域同构的泛同胚；相应地，万有闭；相应地，万有闭且泛单射，
则 $f : \Spec(B) \to \Spec(A)$ 也具有相应性质。
\end{lemma}

\begin{proof}
若 $f_\lambda$ 万有闭，则由引理
\ref{morphisms-lemma-integral-universally-closed}，$B_\lambda$ 在 $A$ 上整。
因此 $B$ 在 $A$ 上整，且 $f$ 万有闭（由同一引理）。
这证明了每个 $f_\lambda$ 都万有闭的情形。

\medskip\noindent
对素理想 $\mathfrak q \subset B$，若它位于 $\mathfrak p \subset A$ 上方，
以 $\mathfrak q_\lambda \subset B_\lambda$ 表示其逆像。则
$\kappa(\mathfrak q) = \colim \kappa(\mathfrak q_\lambda)$。因此，若
$A \to B_\lambda$ 诱导剩余域的纯不可分扩张，则 $A \to B$ 也如此。
这证明了 $f_\lambda$ 万有闭且泛单射的情形；见引理
\ref{morphisms-lemma-universally-injective}。

\medskip\noindent
$f$ 为泛同胚的情形由上述论述和引理
\ref{morphisms-lemma-universal-homeomorphism} 得出，并使用如下事实：$B$ 中的素理想
等同于所有 $B_\lambda$ 中素理想的相容序列。

\medskip\noindent
若 $A \to B_\lambda$ 诱导剩余域同构，则 $A \to B$ 也如此
（见第二段的论证）。这样便证明了剩余情形。
\end{proof}

\begin{lemma}
\label{morphisms-lemma-special-elements-and-localization}
设 $A \subset B$ 是环扩张，$S \subset A$ 是乘法子集。设 $n \geq 1$，
并取 $b_i \in B$，其中 $1 \leq i \leq n$。任意
$x \in S^{-1}B$，若满足
$$
x \not \in S^{-1}A\text{ 且 } b_i x^i \in S^{-1}A\text{ 对 }i = 1, \ldots, n
$$
则等于 $s^{-1}y$，其中 $s \in S$ 且 $y \in B$，并满足
$$
y \not \in A\text{ 且 } b_i y^i \in A\text{ 对 }i = 1, \ldots, n
$$
\end{lemma}

\begin{proof}
略。提示：通分。
\end{proof}

\begin{lemma}
\label{morphisms-lemma-nth-and-nplusone-implies-square-and-cube}
设 $A \subset B$ 是环扩张。若存在 $b \in B$、$b \not \in A$ 和整数
$n \geq 2$，使得 $b^n \in A$ 且 $b^{n + 1} \in A$，则存在
$b' \in B$、$b' \not \in A$，使得 $(b')^2 \in A$ 且 $(b')^3 \in A$。
\end{lemma}

\begin{proof}
取引理中的 $b$ 和 $n$。则 $b$ 的所有充分高次幂均属于 $A$。
事实上，$(b^n)^k(b^{n + 1})^i = b^{(k + i)n + i}$，这蕴含任意幂
$b^m$ 在 $m \geq n^2$ 时都属于 $A$。因此，若 $i \geq 1$ 是使得
$b^i \not \in A$ 的最大整数，则 $(b^i)^2 \in A$ 且 $(b^i)^3 \in A$。
\end{proof}

\begin{lemma}
\label{morphisms-lemma-square-and-cube}
设 $A \subset B$ 是环扩张，使得 $\Spec(B) \to \Spec(A)$ 是诱导剩余域
同构的泛同胚。若 $A \not = B$，则存在 $b \in B$、$b \not \in A$，
使得 $b^2 \in A$ 且 $b^3 \in A$。
\end{lemma}

\begin{proof}
回顾：$A \subset B$ 整（引理
\ref{morphisms-lemma-integral-universally-closed}）。由引理
\ref{morphisms-lemma-subalgebra-inherits}，可假设 $B$ 在 $A$ 上由单个元素生成。
因此 $B$ 在 $A$ 上有限（代数，引理
\ref{algebra-lemma-characterize-finite-in-terms-of-integral}）。故
$B/A$ 作为 $A$-模的支撑闭且非空（代数，引理
\ref{algebra-lemma-support-closed} 和
\ref{algebra-lemma-support-zero}）。取 $\mathfrak p \subset A$ 为该支撑的
一个极小素理想。把 $A \subset B$ 替换为
$A_\mathfrak p \subset B_\mathfrak p$ 后（引理
\ref{morphisms-lemma-special-elements-and-localization} 允许这样做），可假设
$(A, \mathfrak m)$ 是局部环，$B$ 在 $A$ 上有限，并且 $B/A$ 的支撑是
$\{\mathfrak m\}$，这里视其为 $A$-模。由于 $B/A$ 是有限模，可见
$I = \text{Ann}_A(B/A)$ 满足 $\mathfrak m = \sqrt{I}$（代数，引理
\ref{algebra-lemma-support-closed}）。令 $\mathfrak m' \subset B$ 为位于
$\mathfrak m$ 上方的唯一素理想。由于
$\Spec(B) \to \Spec(A)$ 是同胚，可得 $\mathfrak m' = \sqrt{IB}$。
对 $f \in \mathfrak m'$，取 $n \geq 1$，使得 $f^n \in IB$。
则还有 $f^{n + 1} \in IB$。$IB \subset A$ 由我们对 $I$ 的选取成立，故
$f^n, f^{n + 1} \in A$。使用引理
\ref{morphisms-lemma-nth-and-nplusone-implies-square-and-cube}，若
$\mathfrak m' \not \subset A$，便可断定引理成立。然而，若
$\mathfrak m' \subset A$，则 $\mathfrak m' = \mathfrak m$；再由假设中
剩余域也同构，可得 $A = B$。这一矛盾完成证明。
\end{proof}

\begin{lemma}
\label{morphisms-lemma-pth-power-and-multiple}
设 $A \subset B$ 是环扩张，使得 $\Spec(B) \to \Spec(A)$ 是泛同胚。
若 $A \not = B$，则或者存在 $b \in B$、$b \not \in A$，使得
$b^2 \in A$ 且 $b^3 \in A$；或者存在素数 $p$ 和 $b \in B$、
$b \not \in A$，使得 $pb \in A$ 且 $b^p \in A$。
\end{lemma}

\begin{proof}
论证几乎与引理 \ref{morphisms-lemma-square-and-cube} 的证明完全相同，但为确保它
确实成立，我们写出全部细节。

\medskip\noindent
回顾：$A \subset B$ 整（引理
\ref{morphisms-lemma-integral-universally-closed}）。由引理
\ref{morphisms-lemma-subalgebra-inherits}，可假设 $B$ 在 $A$ 上由单个元素生成。
因此 $B$ 在 $A$ 上有限（代数，引理
\ref{algebra-lemma-characterize-finite-in-terms-of-integral}）。故
$B/A$ 作为 $A$-模的支撑闭且非空（代数，引理
\ref{algebra-lemma-support-closed} 和
\ref{algebra-lemma-support-zero}）。取 $\mathfrak p \subset A$ 为该支撑的
一个极小素理想。把 $A \subset B$ 替换为
$A_\mathfrak p \subset B_\mathfrak p$ 后（引理
\ref{morphisms-lemma-special-elements-and-localization} 允许这样做），可假设
$(A, \mathfrak m)$ 是局部环，$B$ 在 $A$ 上有限，并且 $B/A$ 的支撑是
$\{\mathfrak m\}$，这里视其为 $A$-模。由于 $B/A$ 是有限模，可见
$I = \text{Ann}_A(B/A)$ 满足 $\mathfrak m = \sqrt{I}$（代数，引理
\ref{algebra-lemma-support-closed}）。令 $\mathfrak m' \subset B$ 为位于
$\mathfrak m$ 上方的唯一素理想。由于
$\Spec(B) \to \Spec(A)$ 是同胚，可得 $\mathfrak m' = \sqrt{IB}$。
对 $f \in \mathfrak m'$，取 $n \geq 1$，使得 $f^n \in IB$。
则还有 $f^{n + 1} \in IB$。$IB \subset A$ 由我们对 $I$ 的选取成立，故
$f^n, f^{n + 1} \in A$。使用引理
\ref{morphisms-lemma-nth-and-nplusone-implies-square-and-cube}，若
$\mathfrak m' \not \subset A$，便可断定引理成立。若
$\mathfrak m' \subset A$，则 $\mathfrak m' = \mathfrak m$。由于
$A \not = B$，可知剩余域的同态
$\kappa = A/\mathfrak m \to B/\mathfrak m' = \kappa'$ 不可能是同构。
由引理 \ref{morphisms-lemma-universally-injective}，可知 $\kappa$ 的特征是某个素数
$p$，且扩张 $\kappa'/\kappa$ 纯不可分。取 $b \in B$，使其在
$\kappa'$ 中的像不属于 $\kappa$，但其 $p$ 次幂属于 $\kappa$。于是
$b \not \in A$、$b^p \in A$ 且 $pb \in A$（因为
$pb \in \mathfrak m' = \mathfrak m \subset A$），如所需。
\end{proof}

\begin{proposition}
\label{morphisms-proposition-universal-homeomorphism-equal-residue-fields}
设 $A \subset B$ 是环扩张。下列条件等价：
\begin{enumerate}
\item $\Spec(B) \to \Spec(A)$ 是诱导剩余域同构的泛同胚；
\item 每个有限子集 $E \subset B$ 都包含在扩张
$$
A[b_1, \ldots, b_n] \subset B
$$
中，并且有 $b_i^2, b_i^3 \in A[b_1, \ldots, b_{i - 1}]$，
其中 $i = 1, \ldots, n$。
\end{enumerate}
\end{proposition}

\begin{proof}
设 (1) 成立。用超限递归，对每个序数 $\alpha$ 如下构造一个
$A$-子代数 $B_\alpha \subset B$。置 $B_0 = A$。若 $\alpha$ 是极限序数，
则置 $B_\alpha = \colim_{\beta < \alpha} B_\beta$。若
$\alpha = \beta + 1$，则或者 $B_\beta = B$，这时置
$B_\alpha = B_\beta$；或者 $B_\beta \not = B$，这时应用引理
\ref{morphisms-lemma-square-and-cube}，选择 $b_\alpha \in B$、
$b_\alpha \not \in B_\beta$，使得 $b_\alpha^2, b_\alpha^3 \in B_\beta$，
并置 $B_\alpha = B_\beta[b_\alpha] \subset B$。显然
$B = \colim B_\alpha$（事实上，$B = B_\alpha$ 对某个序数 $\alpha$
成立，考察基数即可看出）。我们将用超限归纳证明：(2) 对
$A \to B_\alpha$ 对每个序数 $\alpha$ 成立。对 $\alpha = 0$ 显然成立。
设该陈述对每个 $\beta < \alpha$ 成立，并令 $E \subset B_\alpha$ 为
有限子集。若 $\alpha$ 是极限序数，则
$B_\alpha = \bigcup_{\beta < \alpha} B_\beta$，可见
$E \subset B_\beta$ 对某个 $\beta < \alpha$ 成立，这证明了该情形。若
$\alpha = \beta + 1$，则 $B_\alpha = B_\beta[b_\alpha]$。因此每个
$e \in E$ 都可写成多项式 $e = \sum d_{e, i}b_\alpha^i$，其中
$d_{e, i} \in B_\beta$。令 $D \subset B_\beta$ 为集合
$D = \{d_{e, i}\} \cup \{b_\alpha^2, b_\alpha^3\}$。由归纳假设，存在
命题中所述的 $A$-子代数 $A[b_1, \ldots, b_n] \subset B_\beta$，
它包含 $D$。于是 $A[b_1, \ldots, b_n, b_\alpha] \subset B_\alpha$
是命题中所述的 $A$-子代数，位于 $B_\alpha$ 中，并且包含 $E$。

\medskip\noindent
设 (2) 成立。把 $B = \colim B_\lambda$ 写成其有限 $A$-子代数的余极限。
由引理 \ref{morphisms-lemma-colimit-inherits}，只须证明
$\Spec(B_\lambda) \to \Spec(A)$ 是诱导剩余域同构的泛同胚。
万有闭态射的复合万有闭；诱导剩余域同构的态射也同样如此。因此只须证明：
若 $A \subset B$，且 $B$ 由单个元素 $b$ 生成，并满足
$b^2, b^3 \in A$，则 (1) 成立。这样的扩张整，故
$\Spec(B) \to \Spec(A)$ 万有闭且满射（引理
\ref{morphisms-lemma-integral-universally-closed} 和代数，引理
\ref{algebra-lemma-integral-overring-surjective}）。注意
$(b^2)^3 = (b^3)^2$ 在 $A$ 中成立。对任意环同态
$\varphi : A \to K$ 到域 $K$，可见存在 $\lambda \in K$，使得
$\varphi(b^2) = \lambda^2$ 且 $\varphi(b^3) = \lambda^3$。事实上，若
置 $\lambda = 0$ 若 $\varphi(b^2) = 0$；否则置
$\lambda = \varphi(b^3)/\varphi(b^2)$。因此 $B \otimes_A K$ 是
$K[x]/(x^2 - \lambda^2, x^3 - \lambda^3)$ 的商。该环恰有一个素理想，
其剩余域是 $K$。这蕴含 $\Spec(B) \to \Spec(A)$ 双射并诱导剩余域同构。
结合万有闭性可知 (1) 成立；见引理
\ref{morphisms-lemma-universal-homeomorphism} 和
\ref{morphisms-lemma-universally-injective}。
\end{proof}

\begin{proposition}
\label{morphisms-proposition-universal-homeomorphism}
设 $A \subset B$ 是环扩张。下列条件等价：
\begin{enumerate}
\item $\Spec(B) \to \Spec(A)$ 是泛同胚；
\item 每个有限子集 $E \subset B$ 都包含在扩张
$$
A[b_1, \ldots, b_n] \subset B
$$
中，并且对 $i = 1, \ldots, n$，有
\begin{enumerate}
\item $b_i^2, b_i^3 \in A[b_1, \ldots, b_{i - 1}]$；或
\item 存在素数 $p$，使得
$pb_i, b_i^p \in A[b_1, \ldots, b_{i - 1}]$。
\end{enumerate}
\end{enumerate}
\end{proposition}

\begin{proof}
证明与命题
\ref{morphisms-proposition-universal-homeomorphism-equal-residue-fields} 的证明完全相同，
但须作如下改动：
\begin{enumerate}
\item 用引理 \ref{morphisms-lemma-pth-power-and-multiple} 代替引理
\ref{morphisms-lemma-square-and-cube}。这意味着对每个后继序数
$\alpha = \beta + 1$，或者有 $b_\alpha^2, b_\alpha^3 \in B_\beta$，
或者存在素数 $p$，并有 $pb_\alpha, b_\alpha^p \in B_\beta$。
\item 若 $\alpha$ 是后继序数，则取
$D = \{d_{e, i}\} \cup \{b_\alpha^2, b_\alpha^3\}$，或取
$D = \{d_{e, i}\} \cup \{pb_\alpha, b_\alpha^p\}$，具体取决于
$\alpha$ 落入哪个情形。
\item 在 (2) $\Rightarrow$ (1) 的证明中，还须考虑如下情形：$B$ 在
$A$ 上由单个元素 $b$ 生成，且有 $pb, b^p \in B$，其中 $p$ 是某个素数。
此时 $A \subset B$ 在谱上诱导泛同胚，例如由代数，
引理 \ref{algebra-lemma-p-ring-map}。
\end{enumerate}
这完成证明。
\end{proof}

\begin{lemma}
\label{morphisms-lemma-universal-homeo-iso-if-invert-p}
设 $p$ 是素数。设 $A \to B$ 是诱导同构
$A[1/p] \to B[1/p]$ 的环同态（例如，若 $p$ 在 $A$ 中幂零）。
下列条件等价：
\begin{enumerate}
\item $\Spec(B) \to \Spec(A)$ 是泛同胚；
\item $A \to B$ 的核是局部幂零理想，并且对每个 $b \in B$，存在
$p$-幂 $q$，使得 $qb$ 和 $b^q$ 都在 $A \to B$ 的像中。
\end{enumerate}
\end{lemma}

\begin{proof}
若 (2) 成立，则由代数，引理 \ref{algebra-lemma-p-ring-map}，(1) 成立。
设 (1) 成立。由代数，引理
\ref{algebra-lemma-image-dense-generic-points}，$A \to B$ 的核由幂零元组成。
因此可把 $A$ 替换为 $A \to B$ 的像，并假设 $A \subset B$。由代数，引理
\ref{algebra-lemma-help-with-powers}，集合
$$
B' = \{b \in B \mid p^nb, b^{p^n} \in A\text{ 对某个 }n \geq 0\}
$$
是一个 $A$-子代数，包含于 $B$（对乘法封闭是显然的）。我们须证明
$B' = B$。否则，由引理 \ref{morphisms-lemma-pth-power-and-multiple}，存在
$b \in B$、$b \not \in B'$，使得或者 $b^2, b^3 \in B'$，或者存在素数
$\ell$，使得 $\ell b, b^\ell \in B'$。我们将证明两种情形都导致矛盾，
从而证明引理。

\medskip\noindent
由于 $A[1/p] = B[1/p]$，可选择 $p$-幂 $q$，使得 $qb \in A$。

\medskip\noindent
若 $b^2, b^3 \in B'$，则还有 $b^q \in B'$。由 $B'$ 的定义，
$(b^q)^{q'} \in A$ 对某个 $p$-幂 $q'$ 成立。于是
$qq'b, b^{qq'} \in A$，从而 $b \in B'$，矛盾。

\medskip\noindent
现在设存在素数 $\ell$，使得 $\ell b, b^\ell \in B'$。若
$\ell \not = p$，则 $\ell b \in B'$ 和 $qb \in A \subset B'$ 蕴含
$b \in B'$，矛盾。因此 $\ell = p$ 且 $b^p \in B'$，于是完全如前
得到矛盾。
\end{proof}

\begin{lemma}
\label{morphisms-lemma-make-universal-homeo}
设 $A$ 是环。设 $x, y \in A$。
\begin{enumerate}
\item 若 $x^3 = y^2$ 在 $A$ 中成立，则
$A \to B = A[t]/(t^2 - x, t^3 - y)$ 诱导剩余域上的双射，并诱导谱上的
泛同胚。
\item 若存在素数 $p$，使得 $p^px = y^p$ 在 $A$ 中成立，则
$A \to B = A[t]/(t^p - x, pt - y)$ 诱导谱上的泛同胚。
\end{enumerate}
\end{lemma}

\begin{proof}
我们用引理 \ref{morphisms-lemma-universal-homeomorphism} 的判据来检验。在两种情形中，
环同态都整。因此只须证明：给定域 $k$ 和环同态
$\varphi : A \to k$，$k$-代数 $B \otimes_A k$ 有唯一的素理想；
在情形 (1)，其剩余域等于 $k$；在情形 (2)，其剩余域在 $k$ 上纯不可分。
见引理 \ref{morphisms-lemma-universally-injective}。

\medskip\noindent
在情形 (1)，置 $\lambda = 0$ 若 $\varphi(x) = 0$；否则置
$\lambda = \varphi(y)/\varphi(x)$。那么
$B = k[t]/(t^2 - \lambda^2, t^3 - \lambda^2)$。结论于是显然。

\medskip\noindent
在情形 (2)，若 $k$ 的特征是 $p$，则得到 $\varphi(y) = 0$，且
$B = k[t]/(t^p - \varphi(x))$；这是局部 Artin $k$-代数，其剩余域或者是
$k$，或者是次数为 $p$ 的 $k$ 的纯不可分扩张。若 $k$ 的特征不是 $p$，
则置 $\lambda = \varphi(y)/p$，可见
$B = k[t]/(t - \lambda) = k$，同样得到结论。
\end{proof}

\begin{lemma}
\label{morphisms-lemma-universal-homeo-limit}
设 $A \to B$ 是环同态。
\begin{enumerate}
\item 若 $A \to B$ 在谱上诱导泛同胚，则
$B = \colim B_i$ 是有限表示 $A$-代数 $B_i$ 的滤过余极限，其中
$A \to B_i$ 在谱上诱导泛同胚。
\item 若 $A \to B$ 诱导剩余域同构，并在谱上诱导泛同胚，则
$B = \colim B_i$ 是有限表示 $A$-代数 $B_i$ 的滤过余极限，其中
$A \to B_i$ 诱导剩余域同构，并在谱上诱导泛同胚。
\end{enumerate}
\end{lemma}

\begin{proof}
(1) 的证明。我们将使用代数，引理
\ref{algebra-lemma-when-colimit} 的判据。设 $A \to C$ 为有限表示，
并设 $\varphi : C \to B$ 是 $A$-代数同态。令
$B' = \varphi(C) \subset B$ 为其像。由引理
\ref{morphisms-lemma-subalgebra-inherits}，$A \to B'$ 在谱上诱导泛同胚。
由代数，引理 \ref{algebra-lemma-ring-colimit-fp}，可写
$B' = \colim_{i \in I} B_i$，其中 $A \to B_i$ 为有限表示，且转移同态满射。
由代数，引理 \ref{algebra-lemma-characterize-finite-presentation}，
可选择指标 $0 \in I$，以及分解 $C \to B_0 \to B'$，它分解同态
$C \to B'$。我们断言，$\Spec(B_i) \to \Spec(A)$ 对所有充分大的
$i$ 都是泛同胚。该断言完成 (1) 的证明。

\medskip\noindent
断言的证明。由引理 \ref{morphisms-lemma-reduction-universal-homeomorphism}，环同态
$A_{red} \to B'_{red}$ 在谱上诱导泛同胚。因此，由代数，引理
\ref{algebra-lemma-image-dense-generic-points}，
$A_{red} \subset B'_{red}$。置 $A' = \Im(A \to B')$，则有诱导双射的满射
$A \to A' \to A_{red}$，且
$\Spec(A_{red}) = \Spec(A') = \Spec(A)$。因此
$A' \subset B'$ 在谱上诱导泛同胚。由命题
\ref{morphisms-proposition-universal-homeomorphism} 以及 $B'$ 在 $A'$ 上有限型这一事实，
可找到 $n$ 和 $b'_1, \ldots, b'_n \in B'$，使得
$B' = A'[b'_1, \ldots, b'_n]$，并且对 $j = 1, \ldots, n$，有
\begin{enumerate}
\item $(b'_j)^2, (b'_j)^3 \in A'[b'_1, \ldots, b'_{j - 1}]$；或
\item 存在素数 $p$，使得
$pb'_j, (b'_j)^p \in A'[b'_1, \ldots, b'_{j - 1}]$。
\end{enumerate}
选择 $b_1, \ldots, b_n \in B_0$，它们提升
$b'_1, \ldots, b'_n$。对 $i \geq 0$，以 $b_{j, i}$ 表示 $b_j$ 在
$B_i$ 中的像。对充分大的 $i$，我们将对 $j = 1, \ldots, n$ 有
\begin{enumerate}
\item $b_{j, i}^2, b_{j, i}^3 \in A_i[b_{1, i}, \ldots, b_{j - 1, i}]$；或
\item 存在素数 $p$，使得
$pb_{j, i}, b_{j, i}^p \in A_i[b_{1, i}, \ldots, b_{j - 1, i}]$。
\end{enumerate}
这里 $A_i \subset B_i$ 是 $A \to B_i$ 的像。注意 $A \to A_i$ 是满射
环同态，其核为局部幂零理想。必要时再增大 $i$，可假设 $B_i$ 由
$b_1, \ldots, b_n$ 在 $A_i$ 上生成，换言之，
$B_i = A_i[b_1, \ldots, b_n]$。由代数，引理
\ref{algebra-lemma-p-ring-map} 和
\ref{algebra-lemma-2-3-ring-map}，可知
$A \to A_i \to A_i[b_1] \to \ldots \to A_i[b_1, \ldots, b_n] = B_i$
在谱上诱导泛同胚。这完成断言的证明。

\medskip\noindent
(2) 的证明完全相同。
\end{proof}

\section{绝对弱正规化与半正规化}
\label{morphisms-section-seminormalization}

\noindent
受上一节所证结果的启发，我们给出如下定义。

\begin{definition}
\label{morphisms-definition-seminormal-ring}
设 $A$ 是环。
\begin{enumerate}
\item 若 $A$ 具有如下性质：对所有 $x, y \in A$，只要
$x^3 = y^2$，就存在唯一的 $a \in A$，使得 $x = a^2$ 且 $y = a^3$，
则称它{\it 半正规}。
\item 若 $A$ 满足：(a) $A$ 半正规；并且 (b) 对任意素数 $p$ 和
$x, y \in A$，只要 $p^px = y^p$，就存在唯一的 $a \in A$，使得
$x = a^p$ 且 $y = pa$，则称它{\it 绝对弱正规}。
\end{enumerate}
\end{definition}

\noindent
一个有趣的观察（见 \cite{Costa}）是：在半正规环的定义中，只须
假设\footnote{设 $A$ 是环，并假设对所有 $x, y \in A$，只要
$x^3 = y^2$，都存在 $a \in A$，使得
$x = a^2$ 且 $y = a^3$。那么 $A$ 既约：若 $x^2 = 0$，则
$x^2 = x^3$，因而存在 $a$，使得 $x = a^3$ 且 $x = a^2$。于是
$x = a^3 = ax = a^4 = x^2 = 0$。最后，若
$a_1^2 = a_2^2$ 且 $a_1^3 = a_2^3$，其中 $a_1, a_2$ 属于既约环，则
$(a_1 - a_2)^3 = a_1^3 - 3a_1^2a_2 +3a_1a_2^2 - a_2^3 =
(1 - 3 + 3 - 1)a_1^3 = 0$，故 $a_1 = a_2$。}$a$ 的存在性。
绝对弱正规概形定义于 \cite[Appendix B]{rydh_descent}。

\begin{lemma}
\label{morphisms-lemma-seminormal-local-property}
成为半正规或绝对弱正规是环的局部性质；见性质，定义
\ref{properties-definition-property-local}。
\end{lemma}

\begin{proof}
设 $A$ 半正规，且 $f \in A$。取 $x', y' \in A_f$，并满足
$(x')^3 = (y')^2$。写成 $x' = x/f^{2n}$ 和 $y' = y/f^{3n}$，
其中 $n \geq 0$ 且 $x, y \in A$。把 $x, y$ 分别替换为
$f^{2m}x, f^{3m}y$，并把 $n$ 替换为 $n + m$ 后，可见
$x^3 = y^2$ 在 $A$ 中成立。于是找到唯一的 $a \in A$，使得
$x = a^2$ 且 $y = a^3$。置 $a' = a/f^n$，便得到
$x' = (a')^2$ 且 $y' = (a')^3$，如所需。$a'$ 的唯一性由 $a$ 的
唯一性得出。完全同样地，读者可证明：若 $A$ 绝对弱正规，则 $A_f$
绝对弱正规。

\medskip\noindent
设 $A$ 是环，且 $f_1, \ldots, f_n \in A$ 生成单位理想。假设
$A_{f_i}$ 对每个 $i$ 都半正规。取 $x, y \in A$，并满足
$x^3 = y^2$。对每个 $i$，可找到唯一的 $a_i \in A_{f_i}$，使得
$x = a_i^2$ 且 $y = a_i^3$ 在 $A_{f_i}$ 中成立。由唯一性和第一段的结果（它告诉我们
$A_{f_if_j}$ 半正规），可见 $a_i$ 与 $a_j$ 映到 $A_{f_if_j}$ 的同一元素。
由代数，引理 \ref{algebra-lemma-standard-covering}，可找到唯一的
$a \in A$，它映到 $a_i$ 在 $A_{f_i}$ 中，对所有 $i$ 均如此。同理可得
$x = a^2$ 且 $y = a^3$。显然这个 $a$ 唯一。因此 $A$ 半正规。
若假设 $A_{f_i}$ 绝对弱正规，则完全相同的论证表明 $A$ 绝对弱正规。
\end{proof}

\noindent
接下来定义半正规概形和绝对弱正规概形。

\begin{definition}
\label{morphisms-definition-seminormal}
设 $X$ 是概形。
\begin{enumerate}
\item 若 $X$ 具有如下性质：每个 $x \in X$ 都有仿射开邻域
$\Spec(R) = U \subset X$，且环 $R$ 半正规，则称它{\it 半正规}。
\item 若 $X$ 具有如下性质：每个 $x \in X$ 都有仿射开邻域
$\Spec(R) = U \subset X$，且环 $R$ 绝对弱正规，则称它
{\it 绝对弱正规}。
\end{enumerate}
\end{definition}

\noindent
下面是必备引理。

\begin{lemma}
\label{morphisms-lemma-locally-seminormal}
设 $X$ 是概形。下列条件等价：
\begin{enumerate}
\item 概形 $X$ 半正规。
\item 对每个仿射开集 $U \subset X$，环 $\mathcal{O}_X(U)$ 半正规。
\item 存在仿射开覆盖 $X = \bigcup U_i$，使得每个
$\mathcal{O}_X(U_i)$ 都半正规。
\item 存在开覆盖 $X = \bigcup X_j$，使得每个开子概形 $X_j$ 都半正规。
\end{enumerate}
此外，若 $X$ 半正规，则每个开子概形都半正规。把“半正规”替换为
“绝对弱正规”，上述陈述仍成立。
\end{lemma}

\begin{proof}
结合性质，引理 \ref{properties-lemma-locally-P} 和引理
\ref{morphisms-lemma-seminormal-local-property}。
\end{proof}

\begin{lemma}
\label{morphisms-lemma-seminormal-reduced}
半正规概形或环既约。从而，绝对弱正规概形或环也同样既约。
\end{lemma}

\begin{proof}
设 $A$ 是环。若 $a \in A$ 非零但 $a^2 = 0$，则 $a^2 = 0^2$ 且
$a^3 = 0^3$，从而 $A$ 不半正规。
\end{proof}

\begin{lemma}
\label{morphisms-lemma-seminormalization-ring}
设 $A$ 是环。
\begin{enumerate}
\item 在谱上诱导泛同胚的环同态 $A \to B$ 所成的范畴有终对象
$A \to A^{awn}$。
\item 给定 (1) 的范畴中的 $A \to B$，所得同态
$B \to A^{awn}$ 是同构，当且仅当 $B$ 绝对弱正规。
\item 在剩余域上诱导同构且在谱上诱导泛同胚的环同态 $A \to B$
所成的范畴有终对象 $A \to A^{sn}$。
\item 给定 (3) 的范畴中的 $A \to B$，所得同态
$B \to A^{sn}$ 是同构，当且仅当 $B$ 半正规。
\end{enumerate}
对任意环同态 $\varphi : A \to A'$，存在唯一的同态
$\varphi^{awn} : A^{awn} \to (A')^{awn}$ 和
$\varphi^{sn} : A^{sn} \to (A')^{sn}$ 与 $\varphi$ 相容。
\end{lemma}

\begin{proof}
我们证明 (1) 和 (2)，略去 (3)、(4) 以及最后陈述的证明。考虑如下形式的
$A$-代数所成的范畴：
$$
B = A[x_1, \ldots, x_n]/J
$$
其中 $J$ 是有限生成理想，使得 $A \to B$ 在谱上定义泛同胚。我们断言该范畴
有向（范畴，定义 \ref{categories-definition-directed}）。事实上，给定
$$
B = A[x_1, \ldots, x_n]/J
\quad\text{且}\quad
B' = A[x_1, \ldots, x_{n'}]/J'
$$
则可考虑
$$
B'' = A[x_1, \ldots, x_{n + n'}]/J''
$$
其中 $J''$ 由 $J$ 的元素以及元素
$f(x_{n + 1}, \ldots, x_{n + n'})$ 生成，其中 $f \in J'$。于是有
$A$-代数同态 $B \to B''$ 和 $B' \to B''$，它们诱导同构
$B \otimes_A B' \to B''$。由引理
\ref{morphisms-lemma-base-change-universal-homeomorphism} 和
\ref{morphisms-lemma-composition-universal-homeomorphism} 可知
$\Spec(B'') \to \Spec(A)$ 是泛同胚，故 $A \to B''$ 属于我们的范畴。
最后，给定该范畴中的 $\varphi, \varphi' : B \to B'$，其中 $B$ 如上显示，
考虑商 $B''$；它是 $B'$ 除以由 $\varphi(x_i) - \varphi'(x_i)$
（$i = 1, \ldots, n$）生成的理想所得的商。由于
$\Spec(B') = \Spec(B)$，可见
$\Spec(B'') \to \Spec(B')$ 是双射闭浸入，故为泛同胚。因此 $B''$
属于我们的范畴，并且 $\varphi, \varphi'$ 被 $B' \to B''$ 等化。
这完成断言的证明。置
$$
A^{awn} = \colim B
$$
其中余极限取遍刚描述的范畴。注意，由引理
\ref{morphisms-lemma-colimit-inherits}，$A \to A^{awn}$ 在谱上诱导泛同胚
（此处用到该范畴有向）。

\medskip\noindent
给定在谱上诱导泛同胚的有限表示环同态 $A \to B$，得到典范同态
$B \to A^{awn}$；这是由 $A^{awn}$ 的构造本身给出的。由于 (1) 中每个 $A \to B$ 都是 (1) 中
有限表示的 $A \to B$ 的滤过余极限（引理
\ref{morphisms-lemma-universal-homeo-limit}），可见 $A \to A^{awn}$ 是 (1) 的范畴中的
终对象。

\medskip\noindent
设 $x, y \in A^{awn}$ 是满足 $x^3 = y^2$ 的元素。由引理
\ref{morphisms-lemma-make-universal-homeo}，
$A^{awn} \to A^{awn}[t]/(t^2 - x, t^3 - y)$ 在谱上诱导泛同胚。因此
$A \to A^{awn}[t]/(t^2 - x, t^3 - y)$ 属于 (1) 的范畴，于是得到唯一的
$A$-代数同态
$A^{awn}[t]/(t^2 - x, t^3 - y) \to A^{awn}$。$a \in A^{awn}$ 为
$t$ 的像，因而是满足 $a^2 = x$ 且 $a^3 = y$ 的唯一元素，
这些等式在 $A^{awn}$ 中成立。完全同样地，给定素数 $p$ 和
$x, y \in A^{awn}$，若 $p^px = y^p$，则找到唯一的 $a \in A^{awn}$，
使得 $a^p = x$ 且 $pq = y$。因此按定义，$A^{awn}$ 绝对弱正规。

\medskip\noindent
最后，设 $A \to B$ 属于 (1) 的范畴，且 $B$ 绝对弱正规。由于
$A^{awn} \to B^{awn}$ 在谱上诱导泛同胚，并且 $A^{awn}$ 既约（引理
\ref{morphisms-lemma-seminormal-reduced}），可得 $A^{awn} \subset B^{awn}$
（见代数，引理 \ref{algebra-lemma-image-dense-generic-points}）。若此包含
不是等式，则引理 \ref{morphisms-lemma-pth-power-and-multiple} 蕴含存在元素
$b \in B^{awn}$、$b \not \in A^{awn}$，使得或者
$b^2, b^3 \in A^{awn}$，或者 $pb, b^p \in A^{awn}$ 对某个素数 $p$ 成立。
然而，由定义 \ref{morphisms-definition-seminormal-ring} 中的存在唯一性，这迫使
$b \in A^{awn}$，从而得到矛盾，完成证明。
\end{proof}

\begin{lemma}
\label{morphisms-lemma-seminormalization}
设 $X$ 是概形。
\begin{enumerate}
\item 泛同胚 $Y \to X$ 所成的范畴有始对象 $X^{awn} \to X$。
\item 给定 (1) 的范畴中的 $Y \to X$，所得态射
$X^{awn} \to Y$ 是同构，当且仅当 $Y$ 绝对弱正规。
\item 诱导剩余域同构的泛同胚 $Y \to X$ 所成的范畴有始对象
$X^{sn} \to X$。
\item 给定 (3) 的范畴中的 $Y \to X$，所得态射
$X^{sn} \to Y$ 是同构，当且仅当 $Y$ 半正规。
\end{enumerate}
对概形的任意态射 $h : X' \to X$，存在唯一的态射
$h^{awn} : (X')^{awn} \to X^{awn}$ 和
$h^{sn} : (X')^{sn} \to X^{sn}$ 与 $h$ 相容。
\end{lemma}

\begin{proof}
我们证明 (1) 和 (2)，略去 (3) 和 (4) 的证明。设 $h : X' \to X$
是概形的态射。若 (1) 对 $X$ 和 $X'$ 成立，则
$X' \times_X X^{awn} \to X'$ 是泛同胚，故存在态射
$(X')^{awn} \to X' \times_X X^{awn}$，它在 $X'$ 上唯一；这是由
$(X')^{awn} \to X'$ 的泛性质得到的。与投影
$X' \times_X X^{awn} \to X^{awn}$ 复合，得到 $h^{awn}$。若进一步 (2)
对 $X$ 和 $X'$ 成立，且 $h$ 是开浸入，则
$X' \times_X X^{awn}$ 绝对弱正规（引理
\ref{morphisms-lemma-locally-seminormal}），于是可得
$(X')^{awn} \to X' \times_X X^{awn}$ 是同构。

\medskip\noindent
回顾：任意泛同胚仿射；见引理 \ref{morphisms-lemma-homeomorphism-affine}。因此若
$X$ 仿射，则 (1) 和 (2) 立即由引理
\ref{morphisms-lemma-seminormalization-ring} 得出。设 $X$ 是概形，并令
$\mathcal{B}$ 为 $X$ 的仿射开集所成的集合。对每个
$U \in \mathcal{B}$，得到 $U^{awn} \to U$；对
$V \subset U$，其中 $V, U \in \mathcal{B}$，由上一段的论述得到典范同构
$\rho_{V, U} : V^{awn} \to V \times_U U^{awn}$。因此，由相对粘合
（构造，引理 \ref{constructions-lemma-relative-glueing}），得到态射
$X^{awn} \to X$；它限制为 $U^{awn}$，位于 $U$ 上，并与
$\rho_{V, U}$ 相容。接着，设 $Y \to X$ 是泛同胚。则
$U \times_X Y \to U$ 对 $U \in \mathcal{B}$ 是泛同胚，并得到唯一态射
$g_U : U^{awn} \to U \times_X Y$，它在 $U$ 上。这些 $g_U$ 与态射
$\rho_{V, U}$ 相容；细节略。因此存在唯一态射
$g : X^{awn} \to Y$，它在 $X$ 上，并与 $g_U$ 一致，在 $U$ 上；见构造，评注
\ref{constructions-remark-relative-glueing-functorial}。这对 $X$ 证明了 (1)。
(2) 因其在仿射局部成立而得出。
\end{proof}

\begin{definition}
\label{morphisms-definition-seminormalization}
设 $X$ 是概形。
\begin{enumerate}
\item 引理 \ref{morphisms-lemma-seminormalization} 中构造的态射
$X^{sn} \to X$ 称为 $X$ 的{\it 半正规化}。
\item 引理 \ref{morphisms-lemma-seminormalization} 中构造的态射
$X^{awn} \to X$ 称为 $X$ 的{\it 绝对弱正规化}。
\end{enumerate}
\end{definition}

\noindent
确实，$X^{sn}$ 是 $X$ 的半正规化，也是半正规概形，而绝对弱正规化
$X^{awn}$ 是绝对弱正规概形。此外，对概形的任意态射
$h : Y \to X$，得到典范交换图
$$
\xymatrix{
Y^{awn} \ar[d]^{h^{awn}} \ar[r] &
Y^{sn} \ar[d]^{h^{sn}} \ar[r] &
Y \ar[d]^h \\
X^{awn} \ar[r] &
X^{sn} \ar[r] &
X
}
$$
；其中箭头 $h^{sn}$ 和 $h^{awn}$ 是与 $h$ 相容的唯一箭头。

\begin{lemma}
\label{morphisms-lemma-characterize-seminormal}
设 $X$ 是概形。下列条件等价：
\begin{enumerate}
\item $X$ 半正规；
\item $X$ 等于其自身的半正规化，即态射 $X^{sn} \to X$ 是同构；
\item 若 $\pi : Y \to X$ 是诱导剩余域同构的泛同胚，且 $Y$ 既约，
则 $\pi$ 是同构。
\end{enumerate}
\end{lemma}

\begin{proof}
(1) 与 (2) 的等价性由引理 \ref{morphisms-lemma-seminormalization} 显然。若 (3) 成立，
则 $X^{sn} \to X$ 是同构，从而 (2) 成立。

\medskip\noindent
设 (2) 成立，并令 $\pi : Y \to X$ 是诱导剩余域同构的泛同胚，且
$Y$ 既约。由引理 \ref{morphisms-lemma-seminormalization}，存在分解
$X \to Y \to X$，它分解 $\text{id}_X$。于是 $X \to Y$ 是闭浸入
（由概形，引理 \ref{schemes-lemma-section-immersion}，以及 $\pi$ 分离这一
事实，例如由引理 \ref{morphisms-lemma-universally-injective-separated}）。由于
$X \to Y$ 还在点上双射，而 $Y$ 既约，它必为同构。这完成证明。
\end{proof}

\begin{lemma}
\label{morphisms-lemma-characterize-absolutely-weakly-normal}
设 $X$ 是概形。下列条件等价：
\begin{enumerate}
\item $X$ 绝对弱正规；
\item $X$ 等于其自身的绝对弱正规化，即态射 $X^{awn} \to X$ 是同构；
\item 若 $\pi : Y \to X$ 是泛同胚，且 $Y$ 既约，则 $\pi$ 是同构。
\end{enumerate}
\end{lemma}

\begin{proof}
证明与引理 \ref{morphisms-lemma-characterize-seminormal} 完全相同。
\end{proof}





\section{有限局部自由态射}
\label{morphisms-section-finite-locally-free}

\noindent
许多论文中的作者使用有限平坦态射，而他们实际所指的是有限局部自由态射。
原因是，当基底局部诺特时，这两个概念相同。但一般而言并非如此，见
《习题》练习 \ref{exercises-exercise-flat-not-projective}。

\begin{definition}
\label{morphisms-definition-finite-locally-free}
设 $f : X \to S$ 是概形的态射。称 $f$ {\it 有限局部自由}，
如果 $f$ 是仿射态射，并且 $f_*\mathcal{O}_X$ 是有限局部自由
$\mathcal{O}_S$-模。在这种情况下，称 $f$ 具有{\it 秩}或
{\it 次数} $d$，如果层 $f_*\mathcal{O}_X$ 是次数为 $d$
的有限局部自由层。
\end{definition}

\noindent
注意，若 $f : X \to S$ 有限局部自由，则 $S$ 是开闭子概形
$S_d$ 的不交并，使得 $f^{-1}(S_d) \to S_d$ 是次数为 $d$
的有限局部自由态射。

\begin{lemma}
\label{morphisms-lemma-finite-flat}
设 $f : X \to S$ 是概形的态射。下列条件等价：
\begin{enumerate}
\item $f$ 有限局部自由；
\item $f$ 有限、平坦且局部有限表示。
\end{enumerate}
若 $S$ 局部诺特，则它们还等价于
\begin{enumerate}
\item[(3)] $f$ 有限且平坦。
\end{enumerate}
\end{lemma}

\begin{proof}
设 $V \subset S$ 是仿射开集。在上述三种情形中，态射都是仿射的，
因而 $f^{-1}(V)$ 是仿射的。因此可写
$V = \Spec(R)$ 和 $f^{-1}(V) = \Spec(A)$，其中取某个
$R$-代数 $A$。假设 (1) 成立。这意味着 $S$ 可由仿射开集
$V = \Spec(R)$ 覆盖，使得 $A$ 作为 $R$-模是有限自由的。
由《代数》引理 \ref{algebra-lemma-finite-finite-type}，
$R \to A$ 是有限表示的。因此 (2) 成立。反之，假设 (2) 成立。
对每个仿射开集 $V = \Spec(R)$（它是 $S$ 的开集），环同态
$R \to A$ 是有限且有限表示的，并且 $A$ 作为 $R$-模是平坦的。
由《代数》引理
\ref{algebra-lemma-finite-finitely-presented-extension}，
$A$ 作为 $R$-模是有限表示的。因此《代数》引理
\ref{algebra-lemma-finite-projective} 推出 $A$ 有限局部自由，
故 (1) 成立。诺特情形来自以下事实：诺特环上的有限模是有限表示模；
见《代数》引理
\ref{algebra-lemma-Noetherian-finite-type-is-finite-presentation}。
\end{proof}

\begin{lemma}
\label{morphisms-lemma-composition-finite-locally-free}
有限局部自由态射的复合仍是有限局部自由态射。
\end{lemma}

\begin{proof}
略。
\end{proof}

\begin{lemma}
\label{morphisms-lemma-base-change-finite-locally-free}
有限局部自由态射的基变换仍是有限局部自由态射。
\end{lemma}

\begin{proof}
略。
\end{proof}

\begin{lemma}
\label{morphisms-lemma-finite-locally-free}
设 $f : X \to S$ 是概形的有限局部自由态射。存在由开闭子概形给出的
不交并分解 $S = \coprod_{d \geq 0} S_d$，使得令
$X_d = f^{-1}(S_d)$ 后，各限制 $f|_{X_d}$ 都是有限局部自由态射
$X_d \to S_d$，其次数为 $d$。
\end{lemma}

\begin{proof}
这是因为有限局部自由层在局部具有良定义的秩。细节略。
\end{proof}

\begin{lemma}
\label{morphisms-lemma-massage-finite}
设 $f : Y \to X$ 是有限态射，且 $X$ 仿射。存在图
$$
\xymatrix{
Z' \ar[rd] &
Y' \ar[l]^i \ar[d] \ar[r] &
Y \ar[d] \\
 & X' \ar[r] & X
}
$$
其中
\begin{enumerate}
\item $Y' \to Y$ 和 $X' \to X$ 都是满的有限局部自由态射；
\item $Y' = X' \times_X Y$；
\item $i : Y' \to Z'$ 是闭浸入；
\item $Z' \to X'$ 有限局部自由；并且
\item $Z' = \bigcup_{j = 1, \ldots, m} Z'_j$ 是闭子概形的
（集合论意义下的）有限并，其中每一个都同构地映到 $X'$。
\end{enumerate}
\end{lemma}

\begin{proof}
写成 $X = \Spec(A)$ 和 $Y = \Spec(B)$。另见《代数续篇》章节
\ref{more-algebra-section-descent-flatness-integral}。
设 $x_1, \ldots, x_n \in B$ 是 $B$ 在 $A$ 上的一组生成元。
对每个 $i$，可以选取首一多项式 $P_i(T) \in A[T]$，
使得 $P(x_i) = 0$ 在 $B$ 中成立。由《代数》引理
\ref{algebra-lemma-adjoin-roots}（应用 $n$ 次），存在有限局部自由环扩张
$A \subset A'$，使得每个 $P_i$ 完全分裂：
$$
P_i(T) = \prod\nolimits_{k = 1, \ldots, d_i} (T - \alpha_{ik})
$$
，其中 $\alpha_{ik} \in A'$。令
$$
C = A'[T_1, \ldots, T_n]/(P_1(T_1), \ldots, P_n(T_n))
$$
以及 $B' = A' \otimes_A B$。同态
$C \to B'$，$T_i \mapsto 1 \otimes x_i$ 是 $A'$-代数满同态。
令 $X' = \Spec(A')$、$Y' = \Spec(B')$ 和 $Z' = \Spec(C)$，
便知 (1)--(4) 成立。(5) 成立是因为，在集合论意义下，
$\Spec(C)$ 是由下列理想切出的闭子概形的并：
$$
(T_1 - \alpha_{1k_1}, T_2 - \alpha_{2k_2}, \ldots, T_n - \alpha_{nk_n})
$$
，其中任意 $1 \leq k_i \leq d_i$。
\end{proof}

\noindent
下一个引理以正确的一般性表述在下文引理
\ref{morphisms-lemma-image-nowhere-dense-quasi-finite} 中。

\begin{lemma}
\label{morphisms-lemma-image-nowhere-dense-finite}
设 $f : Y \to X$ 是概形的有限态射。设 $T \subset Y$ 是 $Y$ 的
闭无处稠密子集。则 $f(T) \subset X$ 是 $X$ 的闭无处稠密子集。
\end{lemma}

\begin{proof}
由引理 \ref{morphisms-lemma-finite-proper}，已知 $f(T) \subset X$ 是闭的。
设 $X = \bigcup X_i$ 是仿射开覆盖。由于 $T$ 在 $Y$ 中无处稠密，
$T \cap f^{-1}(X_i)$ 在 $f^{-1}(X_i)$ 中也无处稠密。因此，若能在仿射
情形证明该定理，便知 $f(T) \cap X_i$ 无处稠密。于是由《拓扑》引理
\ref{topology-lemma-nowhere-dense-local} 可知 $T$ 在 $X$ 中无处稠密。

\medskip\noindent
假设 $X$ 仿射。选取引理 \ref{morphisms-lemma-massage-finite} 中的图
$$
\xymatrix{
Z' \ar[rd] &
Y' \ar[l]^i \ar[d]^{f'} \ar[r]_a &
Y \ar[d]^f \\
 & X' \ar[r]^b & X
}
$$
。态射 $a$、$b$ 都是开的，因为它们有限局部自由
（引理 \ref{morphisms-lemma-finite-flat} 和 \ref{morphisms-lemma-fppf-open}）。
因此 $T' = a^{-1}(T)$ 无处稠密，见《拓扑》引理
\ref{topology-lemma-open-inverse-image-closed-nowhere-dense}。
态射 $b$ 是满且开的。因此，若能证明
$f'(T') = b^{-1}(f(T))$ 无处稠密，则由《拓扑》引理
\ref{topology-lemma-open-inverse-image-closed-nowhere-dense}，
$f(T)$ 无处稠密。由于 $i$ 是闭浸入，由《拓扑》引理
\ref{topology-lemma-closed-image-nowhere-dense}，
$i(T') \subset Z'$ 是闭且无处稠密的。这样便把问题归结为下一段
所讨论的情形。

\medskip\noindent
假设 $Y = \bigcup_{i = 1, \ldots, n} Y_i$ 是闭子集的有限并，
每个子集都同构地映到 $X$。考虑 $T_i = Y_i \cap T$。
若每个 $T_i$ 在相应的 $Y_i$ 中都无处稠密，则每个 $f(T_i)$ 在 $X$
中都无处稠密，因为 $Y_i \to X$ 是同构。因此
$f(T) = f(T_i)$ 是 $X$ 中无处稠密闭子集的有限并，于是结论成立；
见《拓扑》引理 \ref{topology-lemma-nowhere-dense}。
若不然，不妨说 $T_1$ 包含一个非空开集
$V \subset Y_1$。我们将说明这会导致矛盾。
考虑 $Y_2 \cap V \subset V$。它或者是一个真闭子集，或者等于 $V$。
在第一种情形，将 $V$ 替换为 $V \setminus V \cap Y_2$，于是
$V \subset T_1$ 在 $Y_1$ 中是开集且不与 $Y_2$ 相交。
在第二种情形，$V \subset Y_1 \cap Y_2$ 在 $Y_1$ 和 $Y_2$
中都是开集。依次对 $i = 3, \ldots, n$ 重复此过程。结果得到不交并分解
$$
\{1, \ldots, n\} = I_1 \amalg I_2, \quad 1 \in I_1
$$
，以及一个开集 $V$（它是 $Y_1$ 的开集且包含于 $T_1$ 中），使得
$V \subset Y_i$ 对 $i \in I_1$ 成立，而
$V \cap Y_i = \emptyset$ 对 $i \in I_2$ 成立。令 $U = f(V)$。
这是 $X$ 的开集，因为 $f|_{Y_1} : Y_1 \to X$ 是同构。于是
$$
f^{-1}(U) = V\ \amalg\ \bigcup\nolimits_{i \in I_2} (Y_i \cap f^{-1}(U))
$$
。由于 $\bigcup_{i \in I_2} Y_i$ 是闭的，这说明
$V \subset f^{-1}(U)$ 是开的，从而 $V \subset Y$ 是开集。
这与 $T$ 在 $Y$ 中无处稠密的假设矛盾，证毕。
\end{proof}

\section{有理映射}
\label{morphisms-section-rational-maps}

\noindent
设 $X$ 是概形。注意，若 $U$、$V$ 是 $X$ 中的稠密开集，
则 $U \cap V$ 也是稠密开集。

\begin{definition}
\label{morphisms-definition-rational-map}
设 $X$、$Y$ 是概形。
\begin{enumerate}
\item 设 $f : U \to Y$、$g : V \to Y$ 是概形态射，定义在
$U$、$V$ 上；这些子集属于 $X$。称 $f$ {\it 等价}于 $g$，如果
$f|_W = g|_W$，其中 $W \subset U \cap V$ 是 $X$ 中的稠密开集。
\item 从 $X$ 到 $Y$ 的{\it 有理映射}是由 (1) 中等价关系定义的等价类。
\item 若 $X$、$Y$ 是在基概形 $S$ 上的概形，则称从 $X$ 到 $Y$ 的有理映射是
一个{\it $S$-有理映射}，即从 $X$ 到 $Y$ 的有理映射，如果该等价类存在一个代表
$f : U \to Y$，且它是一个 $S$-态射。
\end{enumerate}
\end{definition}

\noindent
对于定义 (1) 中的两个态射 $f$、$g$，我们称它们定义了同一个有理映射，
而不说它们是等价的。在某些情形中，有理映射由泛点处局部环之间的映射决定。

\begin{lemma}
\label{morphisms-lemma-rational-map-finite-presentation}
设 $S$ 是概形。设 $X$ 和 $Y$ 是在 $S$ 上的概形。假设 $X$ 有限个
不可约分支，其泛点为 $x_1, \ldots, x_n$。设 $s_i \in S$ 是
$x_i$ 的像。考虑映射
$$
\left\{
\begin{matrix}
S\text{-rational maps} \\
\text{from }X\text{ to }Y
\end{matrix}
\right\}
\longrightarrow
\left\{
\begin{matrix}
(y_1, \varphi_1, \ldots, y_n, \varphi_n)\text{ 其中 }
y_i \in Y\text{ lies over }s_i\text{ 且}\\
\varphi_i : \mathcal{O}_{Y, y_i} \to \mathcal{O}_{X, x_i}
\text{ is a local }\mathcal{O}_{S, s_i}\text{-algebra map}
\end{matrix}
\right\}
$$
，它将 $f : U \to Y$ 映为 $2n$-元组，其中
$y_i = f(x_i)$ 且 $\varphi_i = f^\sharp_{x_i}$。则有
\begin{enumerate}
\item 若 $Y \to S$ 局部有限型，则该映射是单射；
\item 若 $Y \to S$ 局部有限表示，则该映射是双射；
\item 若 $Y \to S$ 局部有限型且 $X$ 既约，则该映射是双射。
\end{enumerate}
\end{lemma}

\begin{proof}
注意，$X$ 的任意稠密开集都包含点 $x_i$，所以该构造有意义。
为证明 (1) 或 (2)，可以将 $X$ 替换为任意稠密开集。因此，若
$Z_1, \ldots, Z_n$ 是 $X$ 的不可约分支，则可以将 $X$ 替换为
$X \setminus \bigcup_{i \not = j} Z_i \cap Z_j$。
完成这一步后，$X$ 是其不可约分支（视为开闭子概形）的不交并。
于是箭头的右端和左端都是关于不可约分支的乘积，问题归结为
$X$ 不可约的情形。

\medskip\noindent
假设 $X$ 不可约，其泛点 $x$ 位于 $s \in S$ 之上。
(1) 来自引理 \ref{morphisms-lemma-morphism-defined-local-ring} 的 (1)。
(2) 和 (3) 来自同一引理的 (2)。
\end{proof}

\begin{definition}
\label{morphisms-definition-rational-function}
设 $X$ 是概形。$X$ 上的{\it 有理函数}是从 $X$ 到
$\mathbf{A}^1_{\mathbf{Z}}$ 的有理映射。
\end{definition}

\noindent
关于仿射直线 $\mathbf{A}^1$ 的定义，见《构造》定义
\ref{constructions-definition-affine-n-space}。设 $X$ 是在 $S$ 上的概形。
对任意开集 $U \subset X$，$U \to \mathbf{A}^1_{\mathbf{Z}}$ 的态射
等同于 $U \to \mathbf{A}^1_S$ 的 $S$-态射。因此，有理函数也等同于
一个 $S$-有理映射，即从 $X$ 到 $\mathbf{A}^1_S$ 的映射。

\medskip\noindent
回忆，有典范同一化
$\Mor(T, \mathbf{A}^1_{\mathbf{Z}}) = \Gamma(T, \mathcal{O}_T)$，其中 $T$ 是任意概形；
见《概形》例 \ref{schemes-example-global-sections}。因此
$\mathbf{A}^1_{\mathbf{Z}}$ 是概形范畴中的环对象。更具体地说，态射
\begin{eqnarray*}
+ : \mathbf{A}^1_{\mathbf{Z}} \times \mathbf{A}^1_{\mathbf{Z}}
& \longrightarrow &
\mathbf{A}^1_{\mathbf{Z}} \\
(f, g) & \longmapsto & f + g \\
* : \mathbf{A}^1_{\mathbf{Z}} \times \mathbf{A}^1_{\mathbf{Z}}
& \longrightarrow &
\mathbf{A}^1_{\mathbf{Z}} \\
(f, g) & \longmapsto & fg
\end{eqnarray*}
满足环中加法和乘法的全部公理（总是取含幺 $1$ 的交换环）。
因此，取值于 $\mathbf{A}^1_{\mathbf{Z}}$ 的有理映射集合也自然成为一个环。

\begin{definition}
\label{morphisms-definition-ring-of-rational-functions}
设 $X$ 是概形。$X$ 上的{\it 有理函数环}是环 $R(X)$，
其元素是有理函数，并配备上述加法和乘法。
\end{definition}

\noindent
对于只有有限个不可约分支的概形，我们可以计算这个环。

\begin{lemma}
\label{morphisms-lemma-integral-scheme-rational-functions}
设 $X$ 是只有有限个不可约分支 $X_1, \ldots, X_n$ 的概形。
若 $\eta_i \in X_i$ 是泛点，则
$$
R(X) = \mathcal{O}_{X, \eta_1} \times \ldots \times \mathcal{O}_{X, \eta_n}
$$
若 $X$ 既约，则这等于 $\prod \kappa(\eta_i)$。
若 $X$ 整，则 $R(X) = \mathcal{O}_{X, \eta} = \kappa(\eta)$
是一个域。
\end{lemma}

\begin{proof}
设 $U \subset X$ 是稠密开子集。则
$U_i = (U \cap X_i) \setminus (\bigcup_{j \not = i} X_j)$
是非空开集，因为它包含 $\eta_i$，包含于 $X_i$，并且
$\bigcup U_i \subset U \subset X$ 是稠密的。
因此，引理中的同一化来自下面这一串等式：
\begin{align*}
R(X)
& =
\colim_{U \subset X\text{ open dense}} \Mor(U, \mathbf{A}^1_\mathbf{Z}) \\
& =
\colim_{U \subset X\text{ open dense}} \mathcal{O}_X(U) \\
& =
\colim_{\eta_i \in U_i \subset X\text{ 开}} \prod \mathcal{O}_X(U_i) \\
& =
\prod \colim_{\eta_i \in U_i \subset X\text{ 开}} \mathcal{O}_X(U_i) \\
& =
\prod \mathcal{O}_{X, \eta_i}
\end{align*}
其中第二个等式来自《概形》例
\ref{schemes-example-global-sections}。最后一个断言来自
《代数》引理 \ref{algebra-lemma-minimal-prime-reduced-ring}。
\end{proof}

\begin{definition}
\label{morphisms-definition-function-field}
设 $X$ 是整概形。$X$ 的{\it 函数域}，或{\it 有理函数域}，
是域 $R(X)$。
\end{definition}

\noindent
我们有时用 $k(X)$ 而不是 $R(X)$ 表示这个域。
利用函数域可以阐明整概形上的分离条件。注意，根据引理
\ref{morphisms-lemma-integral-scheme-rational-functions}，在整概形上，每个局部环
$\mathcal{O}_{X, x}$ 都可视为 $R(X)$ 的局部子环。

\begin{lemma}
\label{morphisms-lemma-distinct-local-rings}
设 $X$ 是整分离概形。设 $Z_1$、$Z_2$ 是 $X$ 中不同的不可约闭子集。
设 $\eta_i$ 是 $Z_i$ 的泛点。若 $Z_1 \not\subset Z_2$，则有
$\mathcal{O}_{X, \eta_1} \not \subset \mathcal{O}_{X, \eta_2}$
作为 $R(X)$ 的子环。特别地，若 $Z_1 = \{x\}$ 由一个闭点 $x$ 组成，则存在一个在 $x$
的某邻域上正则的函数，但它不属于 $\mathcal{O}_{X, \eta_{2}}$。
\end{lemma}

\begin{proof}
首先注意，在假设 $X$ 分离时，存在唯一的概形态射
$\Spec(\mathcal{O}_{X, \eta_2}) \to X$（在 $X$ 上），使得复合
$$
\Spec(R(X)) \longrightarrow
\Spec(\mathcal{O}_{X, \eta_2}) \longrightarrow X
$$
是典范映射 $\Spec(R(X)) \to X$。
确实，有典范映射
$can : \Spec(\mathcal{O}_{X, \eta_2}) \to X$，见《概形》方程
\ref{schemes-equation-canonical-morphism}。给定第二个态射 $a$ 到 $X$，
由假设，$a$ 与 $can$ 在
$\Spec(\mathcal{O}_{X, \eta_2})$ 的泛点处相同。
现在，$X$ 分离保证这两个映射相同的
$\Spec(\mathcal{O}_{X, \eta_2})$ 子集是闭的；见《概形》引理
\ref{schemes-lemma-where-are-they-equal}。因此 $a = can$ 在整个
$\Spec(\mathcal{O}_{X, \eta_2})$ 上成立。

\medskip\noindent
假设 $Z_1 \not \subset Z_2$，并反设
$\mathcal{O}_{X, \eta_{1}} \subset \mathcal{O}_{X, \eta_{2}}$
作为子环包含于 $R(X)$。
于是得到第二个态射
$$
\Spec(\mathcal{O}_{X, \eta_{2}}) \longrightarrow
\Spec(\mathcal{O}_{X, \eta_{1}}) \longrightarrow
X.
$$
根据上面的结论，这个复合必须等于 $can$。
这意味着 $\eta_2$ 特化到 $\eta_1$（见《概形》引理
\ref{schemes-lemma-specialize-points}）。但这与假设
$Z_1 \not \subset Z_2$ 矛盾。
\end{proof}

\begin{definition}
\label{morphisms-definition-domain-of-definition}
设 $\varphi$ 是概形 $X$ 与 $Y$ 之间的有理映射。若 $\varphi$ 在点
$x \in X$ 处有定义，即存在代表 $(U, f)$ 的 $\varphi$ 使得 $x \in U$，
则称其{\it 在该点处有定义}。
$\varphi$ 的{\it 定义域}是所有 $\varphi$ 有定义的点组成的集合。
\end{definition}

\noindent
由此定义，一般而言并不成立：$\varphi$ 存在一个代表，其定义域包含
整个定义域。

\begin{lemma}
\label{morphisms-lemma-rational-map-from-reduced-to-separated}
设 $X$、$Y$ 是概形。假设 $X$ 既约且 $Y$ 分离。
设 $\varphi$ 是从 $X$ 到 $Y$ 的有理映射，其定义域为
$U \subset X$。则存在唯一态射 $f : U \to Y$ 表示 $\varphi$。
若 $X$ 和 $Y$ 是在分离概形 $S$ 上的概形，且 $\varphi$ 是
$S$-有理映射，则 $f$ 是在 $S$ 上的态射。
\end{lemma}

\begin{proof}
设 $(V, g)$ 和 $(V', g')$ 是 $\varphi$ 的代表。则 $g, g'$ 在
稠密开子概形 $W \subset V \cap V'$ 上相同。另一方面，
$E$ 是 $g|_{V \cap V'}$ 与 $g'|_{V \cap V'}$ 的等化子，且是
$V \cap V'$ 的闭子概形（《概形》引理
\ref{schemes-lemma-where-are-they-equal}）。现在
$W \subset E$ 蕴含集合论意义下 $E = V \cap V'$。
由于 $V \cap V'$ 既约，得到概形论意义下
$E = V \cap V'$，即
$g|_{V \cap V'} = g'|_{V \cap V'}$。因此可以将代表
$g : V \to Y$ 的 $\varphi$ 粘合为态射
$f : U \to Y$，见《概形》引理 \ref{schemes-lemma-glue}。
最后一个断言的证明略去。
\end{proof}

\noindent
一般而言，复合有理映射没有意义。原因是第一个有理映射的某个代表的像
可能与第二个有理映射的定义域没有交集。然而，若假设概形不可约并考察
支配有理映射，则可以复合有理映射。

\begin{definition}
\label{morphisms-definition-dominant-rational}
设 $X$ 与 $Y$ 是不可约概形。从 $X$ 到 $Y$ 的有理映射称为
{\it 支配的}，如果任一代表 $f : U \to Y$ 都是概形的支配态射。
\end{definition}

\noindent
由引理 \ref{morphisms-lemma-dominant-finite-number-irreducible-components}，以下要求等价：
$\eta \in X$ 是泛点并映到泛点 $\xi$，它属于 $Y$。对任一代表，都有
$f(\eta) = \xi$，其中代表写为 $f : U \to Y$。我们可以复合支配有理映射
$\varphi$；它介于不可约概形 $X$ 与 $Y$ 之间；以及任意有理映射 $\psi$，它从 $Y$ 到 $Z$。
具体地，取代表
$f : U \to Y$，其中 $U \subset X$ 是稠密开集，并取代表
$g : V \to Z$，其中 $V \subset Y$ 是稠密开集。于是
$W = f^{-1}(V) \subset X$ 是非空开集（因为它包含 $X$ 的泛点），并令
$\psi \circ \varphi$ 为 $g \circ f|_W : W \to Z$ 的等价类。验证良定义性留给读者。

\medskip\noindent
这样得到一个范畴，其对象是不可约概形，态射是支配有理映射。给定基概形 $S$，
同样可以定义一个范畴，其对象是 $S$ 上的不可约概形，态射是支配的 $S$-有理映射。

\begin{definition}
\label{morphisms-definition-birational-schemes}
设 $X$ 与 $Y$ 是不可约概形。
\begin{enumerate}
\item 如果 $X$ 与 $Y$ 在不可约概形及支配有理映射所成的范畴中同构，
  则称 $X$ 与 $Y$ 是{\it 双有理的}。
\item 假设 $X$ 与 $Y$ 是概形，基概形为 $S$。称 $X$ 与 $Y$ 是
  {\it $S$-双有理的}，如果 $X$ 与 $Y$ 在 $S$ 上的不可约概形及支配的
  $S$-有理映射所成的范畴中同构。
\end{enumerate}
\end{definition}

\noindent
若 $X$ 与 $Y$ 是双有理的不可约概形，则从 $X$ 到 $Z$ 的有理映射集合
与从 $Y$ 到 $Z$ 的有理映射集合对所有概形 $Z$ 都存在双射（关于 $Z$ 具有函子性）。对“通常的”
不可约概形而言，这只是可能的一个定义。另一种定义是要求 $X$ 与 $Y$ 具有同构的有理函数环。
对于簇，这些条件是等价的，见引理 \ref{morphisms-lemma-criterion-birational-finite-presentation}。

\begin{lemma}
\label{morphisms-lemma-birational-integral}
设 $X$ 与 $Y$ 是不可约概形。
\begin{enumerate}
\item 概形 $X$ 与 $Y$ 双有理，当且仅当它们具有同构的非空开子概形。
\item 假设 $X$ 与 $Y$ 是基概形 $S$ 上的概形。则 $X$ 与 $Y$ 是
  $S$-双有理的，当且仅当存在非空开集 $U \subset X$ 与 $V \subset Y$，
  使它们在 $S$ 上同构。
\end{enumerate}
\end{lemma}

\begin{proof}
假设 $X$ 与 $Y$ 双有理。设 $f : U \to Y$ 与 $g : V \to X$ 分别定义从
$X$ 到 $Y$ 以及从 $Y$ 到 $X$ 的互逆支配有理映射。我们可以假设 $V$ 是仿射的，
并可将 $U$ 替换为 $f^{-1}(V)$ 的一个仿射开子集。由于 $g \circ f$ 作为支配有理映射
是恒等映射，复合 $U \to V \to X$ 在 $U$ 的一个稠密开集上是恒等映射。因此，
再将 $U$ 换成较小的仿射开集后，可假设 $U \to V \to X$ 是 $U$ 到 $X$ 的包含映射。
于是 $U \to V$ 是一个浸入（将《概形》引理 \ref{schemes-lemma-section-immersion}
应用于 $U \to g^{-1}(U) \to U$）。然而，交换 $U$ 与 $V$ 的角色并重复上述论证，
可见存在非空仿射开集 $V' \subset V$，使包含映射分解为 $V' \to U \to V$。
于是 $V' \to U$ 必然是开浸入。事实上，$V' \to f^{-1}(V') \to V'$ 是单态射
（《概形》引理 \ref{schemes-lemma-immersions-monomorphisms}），并且复合为恒等映射，
所以它们是同构。因而 $V'$ 同时同构于 $X$ 与 $Y$ 的某个开子概形。
在 $S$-有理映射的情形，完全相同的论证成立。
\end{proof}

\begin{remark}
\label{morphisms-remark-generalize-category}
下面给出不可约概形与支配有理映射所成范畴的一个推广。对概形 $X$，记 $X^0$ 为满足
$x \in X$ 且 $\dim(\mathcal{O}_{X, x}) = 0$ 的点的集合；换言之，$X^0$ 是
$X$ 的不可约分支的泛点集合。于是可以考虑如下范畴：
\begin{enumerate}
\item 对象是这样的概形 $X$，其中每个拟紧开集都有有限多个不可约分支；
\item 从 $X$ 到 $Y$ 的态射是有理映射 $f : U \to Y$；它从 $X$ 到 $Y$，
  并满足 $f(U^0) = Y^0$。
\end{enumerate}
如果 $U \subset X$ 是概形的稠密开集，则 $U^0 \subset X^0$ 不一定是等式；
但若 $X$ 是该范畴的对象，则情况确实如此。因此，该范畴中任意两个态射的复合
是良定义的，并且仍是该范畴中的态射。
\end{remark}

\begin{remark}
\label{morphisms-remark-pseudo-morphisms}
定义 \ref{morphisms-definition-rational-map} 有一个变体：只考虑态射 $U \to Y$，其定义在
概形论意义下稠密的开子概形 $U \subset X$ 上。利用引理
\ref{morphisms-lemma-intersection-scheme-theoretically-dense}，可知由此得到一个等价关系。
这些等价类称为从 $X$ 到 $Y$ 的{\it 伪态射}。若 $X$ 既约，则这两个概念相同。
\end{remark}

\section{双有理态射}
\label{morphisms-section-birational}

\noindent
你可能熟悉簇的双有理映射这一概念，它具有在目标的某个开子集上为同构的性质。
然而，一般而言，双有理态射可能在任何非空开集上都不是同构，见例
\ref{morphisms-example-birational-not-iso-over-open}。下面给出形式定义。

\begin{definition}
\label{morphisms-definition-birational}
\begin{reference}
\cite[(2.2.9)]{EGA1}
\end{reference}
设 $X$、$Y$ 是概形。假设 $X$ 与 $Y$ 都只有有限多个不可约分支。
若态射 $f : X \to Y$ 满足以下条件，则称其为{\it 双有理的}：
\begin{enumerate}
\item $f$ 在 $X$ 的不可约分支的泛点集合与 $Y$ 的不可约分支的泛点集合之间诱导双射；
\item 对泛点 $\eta \in X$（它属于 $X$ 的一个不可约分支），局部环映射
  $\mathcal{O}_{Y, f(\eta)} \to \mathcal{O}_{X, \eta}$ 是同构。
\end{enumerate}
\end{definition}

\noindent
下面将看到，双有理态射在泛点上的纤维是单点集。此外，在实际中遇到的大多数情形下，
不可约概形 $X$ 与 $Y$ 之间存在双有理态射，就意味着 $X$ 与 $Y$ 是双有理概形。

\begin{lemma}
\label{morphisms-lemma-birational-dominant}
设 $f : X \to Y$ 是概形态射，且只有有限多个不可约分支。若 $f$ 是双有理的，
则 $f$ 是支配的。
\end{lemma}

\begin{proof}
由引理 \ref{morphisms-lemma-generic-points-in-image-dominant} 及定义立即得到。
\end{proof}

\begin{lemma}
\label{morphisms-lemma-birational-generic-fibres}
设 $f : X \to Y$ 是双有理概形态射，且只有有限多个不可约分支。若 $y \in Y$ 是
一个不可约分支的泛点，则基变换
$X \times_Y \Spec(\mathcal{O}_{Y, y}) \to \Spec(\mathcal{O}_{Y, y})$
是同构。
\end{lemma}

\begin{proof}
不妨设 $Y = \Spec(B)$ 是仿射且不可约的。于是 $X$ 也不可约。若对任意非空
仿射开集 $U \subset X$ 结论成立，则对 $X$ 结论也成立（略去一个小论证）。因此不妨设
$X$ 也是仿射的，记作 $X = \Spec(A)$。
令 $y \in Y$ 对应于极小素理想 $\mathfrak q \subset B$。根据假设，$A$ 有唯一的极小素理想
$\mathfrak p$，它位于 $\mathfrak q$ 之上，并且
$B_\mathfrak q \to A_\mathfrak p$ 是同构。
于是 $A_\mathfrak q \to \kappa(\mathfrak p)$ 是满射，因而
$\mathfrak p A_\mathfrak q$ 是极大理想。另一方面，
$\mathfrak p A_\mathfrak q$ 是 $A_\mathfrak q$ 的唯一极小素理想。
因此 $\mathfrak p A_\mathfrak q$ 是 $A_\mathfrak q$ 的唯一素理想，且
$A_\mathfrak q = A_\mathfrak p$。由于 $A_\mathfrak q = A \otimes_B B_\mathfrak q$，
引理得证。
\end{proof}

\begin{example}
\label{morphisms-example-birational-not-iso-over-open}
下面给出两个双有理态射的例子，它们在目标的任何开集上都不是同构。

\medskip\noindent
第一个例子。设 $k$ 是无限域。设 $A = k[x]$。令
$B = k[x, \{y_{\alpha}\}_{\alpha \in k}]/
((x-\alpha)y_\alpha, y_\alpha y_\beta)$。
存在包含 $A \subset B$ 以及回缩 $B \to A$，后者通过令所有 $y_\alpha$ 都为零得到。
态射 $\Spec(A) \to \Spec(B)$ 与态射 $\Spec(B) \to \Spec(A)$ 都是双有理的，
但在任何开集上都不是同构。

\medskip\noindent
第二个例子。设 $A$ 是整环。设 $S \subset A$ 是不含 $0$ 的乘法闭子集。
令 $B = S^{-1}A$，则态射 $f : \Spec(B) \to \Spec(A)$ 是双有理的。
若存在开集 $U$，它属于 $\Spec(A)$，使得 $f^{-1}(U) \to U$ 是同构，
则存在 $a \in A$，使 $S$ 的每个元素都在主局部化 $A_a$ 中变为可逆元。
取 $A = \mathbf{Z}$ 且 $S$ 为奇整数集合，即得反例。
\end{example}

\begin{lemma}
\label{morphisms-lemma-birational-birational}
设 $f : X \to Y$ 是概形的双有理态射，且这些概形在基概形 $S$ 上并只有有限多个不可约分支。
假设下列条件之一成立：
\begin{enumerate}
\item $f$ 局部有限型且 $Y$ 既约；
\item $f$ 局部有限表示。
\end{enumerate}
则存在稠密开集 $U \subset X$ 与 $V \subset Y$，使得
$f(U) \subset V$ 且 $f|_U : U \to V$ 是同构。特别地，若 $X$ 与 $Y$ 不可约，
则 $X$ 与 $Y$ 是 $S$-双有理的。
\end{lemma}

\begin{proof}
立即可将情形约化为 $X$ 与 $Y$ 不可约的情形，此处略去。进一步缩小后，可假设 $X$ 与
$Y$ 是仿射的，即 $X = \Spec(A)$ 且 $Y = \Spec(B)$。
根据假设，$A$（分别地，$B$）有唯一的极小素理想 $\mathfrak p$（分别地，
$\mathfrak q$），素理想 $\mathfrak p$ 位于 $\mathfrak q$ 之上，并且
$B_\mathfrak q = A_\mathfrak p$。
由引理 \ref{morphisms-lemma-birational-generic-fibres}，有
$B_\mathfrak q = A_\mathfrak q = A_\mathfrak p$。

\medskip\noindent
假设 $B \to A$ 是有限型的，设 $A = B[x_1, \ldots, x_n]$。
存在 $b_i \in B$ 及 $g_i \in B \setminus \mathfrak q$，使得 $b_i/g_i$ 映到
$x_i$ 在 $A_\mathfrak q$ 中的像。因此，$b_i  - g_ix_i$ 在
$A_{g_i'}$ 中映为零，其中 $g_i' \in B \setminus \mathfrak q$。
令 $g = \prod g_i g'_i$，则 $B_g \to A_g$ 是满射。若进一步 $Y$ 既约，则映射
$B_g \to B_\mathfrak q$ 是单射，从而 $B_g \to A_g$ 也是单射。这证明了情形 (1)。

\medskip\noindent
情形 (2) 的证明。由上一段论证可假设 $B \to A$ 是满射。由于 $f$ 局部有限表示，
核 $J \subset B$ 是有限生成理想。设 $J = (b_1, \ldots, b_r)$。由于
$B_\mathfrak q = A_\mathfrak q$，存在 $g_i \in B \setminus \mathfrak q$ 使得
$g_i b_i = 0$。令 $g = \prod g_i$，则 $B_g \to A_g$ 是同构。
\end{proof}

\begin{lemma}
\label{morphisms-lemma-criterion-birational-finite-presentation}
设 $S$ 是概形。设 $X$、$Y$ 是不可约概形，并且在 $S$ 上局部有限表示。
设 $x \in X$、$y \in Y$ 是泛点。以下条件等价：
\begin{enumerate}
\item $X$ 与 $Y$ 是 $S$-双有理的；
\item $X$ 与 $Y$ 有非空开子概形，并且这些开子概形在 $S$ 上同构；
\item $x$ 与 $y$ 映到同一点 $s$（它属于 $S$），且
  $\mathcal{O}_{X, x}$ 与 $\mathcal{O}_{Y, y}$ 作为
  $\mathcal{O}_{S, s}$-代数同构。
\end{enumerate}
\end{lemma}

\begin{proof}
我们已经在引理 \ref{morphisms-lemma-birational-integral} 中证明了 (1) 与 (2) 的等价性。
(2) 蕴含 (3) 是直接的。为完成证明，假设 (3) 成立，并证明 (1)。
由引理 \ref{morphisms-lemma-rational-map-finite-presentation}，存在有理映射
$f : U \to Y$，它把 $x \in U$ 映到 $y$，并诱导给定的同构
$\mathcal{O}_{Y, y} \cong \mathcal{O}_{X, x}$。
因此 $f$ 是双有理态射，从而由引理
\ref{morphisms-lemma-birational-birational} 在非空开子集上诱导同构。
证明完毕。
\end{proof}

\begin{lemma}
\label{morphisms-lemma-common-open}
设 $S$ 是概形。设 $X$、$Y$ 是在 $S$ 上局部有限型的整概形。
设 $x \in X$、$y \in Y$ 是泛点。以下条件等价：
\begin{enumerate}
\item $X$ 与 $Y$ 是 $S$-双有理的；
\item $X$ 与 $Y$ 有非空开子概形，并且这些开子概形在 $S$ 上同构；
\item $x$ 与 $y$ 映到同一点 $s \in S$，且
  $\kappa(x) \cong \kappa(y)$ 作为 $\kappa(s)$-扩张同构。
\end{enumerate}
\end{lemma}

\begin{proof}
我们已经在引理 \ref{morphisms-lemma-birational-integral} 中证明了 (1) 与 (2) 的等价性。
(2) 蕴含 (3) 是直接的。为完成证明，假设 (3) 成立，并证明 (1)。
由代数引理 \ref{algebra-lemma-minimal-prime-reduced-ring}，有
$\mathcal{O}_{X, x} = \kappa(x)$ 且 $\mathcal{O}_{Y, y} = \kappa(y)$。
由引理 \ref{morphisms-lemma-rational-map-finite-presentation}，存在有理映射
$f : U \to Y$，它把 $x \in U$ 映到 $y$，并诱导给定的同构
$\mathcal{O}_{Y, y} \cong \mathcal{O}_{X, x}$。
因此 $f$ 是双有理态射，从而由引理
\ref{morphisms-lemma-birational-birational} 在非空开子集上诱导同构。
证明完毕。
\end{proof}
\section{概形态射的泛有限性}
\label{morphisms-section-generically-finite}

\noindent
本节刻画局部有限型且在某种意义下“泛有限”的概形之间的映射。

\begin{lemma}
\label{morphisms-lemma-generically-finite}
设 $X$、$Y$ 是概形。设 $f : X \to Y$ 是局部有限型的。
设 $\eta \in Y$ 是 $Y$ 的某个不可约分支的泛点。以下条件等价：
\begin{enumerate}
\item 集合 $f^{-1}(\{\eta\})$ 是有限的；
\item 存在仿射开集 $U_i \subset X$（$i = 1, \ldots, n$）以及
  $V \subset Y$，使 $f(U_i) \subset V$、$\eta \in V$ 且
  $f^{-1}(\{\eta\}) \subset \bigcup U_i$，并且每个
  $f|_{U_i} : U_i \to V$ 都是有限的。
\end{enumerate}
若 $f$ 拟分离，则这些条件还等价于
\begin{enumerate}
\item[(3)] 存在仿射开集 $V \subset Y$ 以及 $U \subset X$，使
  $f(U) \subset V$、$\eta \in V$ 且 $f^{-1}(\{\eta\}) \subset U$，
  并且 $f|_U : U \to V$ 是有限的。
\end{enumerate}
若 $f$ 拟紧且拟分离，则这些条件还等价于
\begin{enumerate}
\item[(4)] 存在仿射开集 $V \subset Y$，$\eta \in V$，使得
  $f^{-1}(V) \to V$ 是有限的。
\end{enumerate}
\end{lemma}

\begin{proof}
问题在基上是局部的。因此可用 $Y$ 的一个包含 $\eta$ 的仿射邻域替换，以下证明始终假设
$Y$ 是仿射的，写作 $Y = \Spec(R)$。

\medskip\noindent
显然 (2) 蕴含 (1)。假设
$f^{-1}(\{\eta\}) = \{\xi_1, \ldots, \xi_n\}$ 是有限集。
取仿射开集 $U_i \subset X$，其中 $\xi_i \in U_i$。由代数引理
\ref{algebra-lemma-generically-finite}，将 $Y$ 替换为 $Y$ 中的某个标准开集后，
每个态射 $U_i \to Y$ 都是有限的。换言之，(2) 成立。

\medskip\noindent
显然 (3) 蕴含 (1)。现在假设 $f$ 拟分离且 (1) 成立。写成
$f^{-1}(\{\eta\}) = \{\xi_1, \ldots, \xi_n\}$。由引理
\ref{morphisms-lemma-finite-fibre}，这些 $\xi_i$ 之间没有特化关系。
由于 $\xi_i$ 映到不可约分支的泛点 $\eta$，它属于 $Y$；不存在
非平凡特化 $\xi \leadsto \xi_i$，该特化发生在 $X$ 中（因为 $\xi$ 也会映到 $\eta$）。
因此每个 $\xi_i$ 都是 $X$ 的某个不可约分支的泛点。
由于 $Y$ 仿射且 $f$ 拟分离，$X$ 也是拟分离的。由性质引理
\ref{properties-lemma-maximal-points-affine}，可找到仿射开集 $U \subset X$，其中包含每个
$\xi_i$。由代数引理 \ref{algebra-lemma-generically-finite}，将 $Y$ 替换为
$Y$ 中的某个标准开集后，态射 $U \to Y$ 是有限的。换言之，(3) 成立。

\medskip\noindent
显然 (4) 在不对 $f$ 作任何额外假设时蕴含 (1)--(3)。现在假设 $f$ 拟紧且拟分离。
我们需要证明等价条件 (1)--(3) 蕴含 (4)。
设 $U$、$V$ 如 (3) 中所述，并以 $Y$ 替换为 $V$。由于 $f$ 拟紧且 $Y$ 拟紧（因为它仿射），
$X$ 拟紧。因此 $Z = X \setminus U$ 拟紧，故态射 $f|_Z : Z \to Y$ 拟紧。
由 $Z$ 的构造，$\eta \not \in f(Z)$。由引理
\ref{morphisms-lemma-quasi-compact-generic-point-not-in-image}，存在仿射开邻域 $V'$，包含 $\eta$ 且位于 $Y$ 中，
使 $f^{-1}(V') \cap Z = \emptyset$。于是 $f^{-1}(V') \subset U$，这意味着
$f^{-1}(V') \to V'$ 是有限的。
\end{proof}

\begin{example}
\label{morphisms-example-counter-generically-finite}
令 $A = \prod_{n \in \mathbf{N}} \mathbf{F}_2$。$A$ 的每个元素都是幂等元，
因此每个素理想都是极大理想，剩余域为 $\mathbf{F}_2$。于是
$X = \Spec(A)$ 的拓扑完全不连通且拟紧。投影映射 $A \to \mathbf{F}_2$ 定义
$\Spec(A)$ 的开点。由于 $X$ 拟紧，不可能 $X$ 的所有点都是开的。取 $x \in X$ 为
一个不是开点的闭点。于是可以构造概形 $Y$，把 $X$ 的两个副本沿
$X \setminus \{x\}$ 粘合；换言之，这是 $X$，其中 $x$ 被加倍，比较
《概形》例 \ref{schemes-example-affine-space-zero-doubled}。
态射 $f : Y \to X$ 是拟紧、有限型且具有有限纤维，但不是拟分离的。
$x \in X$ 是 $X$ 的某个不可约分支的泛点（因为 $X$ 完全不连通）。
然而，引理 \ref{morphisms-lemma-generically-finite} 的性质 (3) 与 (4) 不成立。
原因是，对任意开邻域 $x \in U \subset X$，逆像 $f^{-1}(U)$ 都不是仿射的，
因为 $f^{-1}(U)$ 上的函数不能分离位于 $x$ 上方的两个点（证明略；这是一个很好的练习）。
因此，在引理的 (3) 与 (4) 中，要求 $f$ 拟分离是必要的。
\end{example}

\begin{remark}
\label{morphisms-remark-quasi-finite-finite-over-dense-open}
引理 \ref{morphisms-lemma-generically-finite} 的一个替代说法是：拟有限态射在目标的某个稠密开集上有限。
这将在《更多态射》引理
\ref{more-morphisms-lemma-quasi-finite-finite-over-dense-open} 中证明。
\end{remark}
\begin{lemma}
\label{morphisms-lemma-quasi-finiteness-over-generic-point}
设 $X$、$Y$ 是概形。设 $f : X \to Y$ 是局部有限型的。
令 $X^0$、分别的 $Y^0$ 表示 $X$、$Y$ 的不可约分支的泛点集合。令
$\eta \in Y^0$。以下条件等价：
\begin{enumerate}
\item $f^{-1}(\{\eta\}) \subset X^0$；
\item $f$ 在所有位于 $\eta$ 上方的点处拟有限；
\item $f$ 在所有 $\xi \in X^0$（它们位于 $\eta$ 上方）处拟有限。
\end{enumerate}
\end{lemma}

\begin{proof}
条件 (1) 蕴含纤维 $X_\eta$ 的点之间没有特化关系。因此由引理
\ref{morphisms-lemma-quasi-finite-at-point-characterize}，(2) 成立。
蕴含 (2) $\Rightarrow$ (3) 是直接的。
由于 $\eta$ 是 $Y$ 的泛点，$X_\eta$ 的泛点也是 $X$ 的泛点。因此由 (3) 及引理
\ref{morphisms-lemma-quasi-finite-at-point-characterize}，$X_\eta$ 的泛点也都是闭点。
于是 $X_\eta$ 的所有点都是泛点，从而 (1) 成立。
\end{proof}

\begin{lemma}
\label{morphisms-lemma-finite-over-dense-open}
设 $X$、$Y$ 是概形。设 $f : X \to Y$ 是局部有限型的。
令 $X^0$、分别的 $Y^0$ 表示 $X$、$Y$ 的不可约分支的泛点集合。假设
\begin{enumerate}
\item $X^0$ 与 $Y^0$ 有限，且 $f^{-1}(Y^0) = X^0$；
\item $f$ 拟紧或 $f$ 分离。
\end{enumerate}
则存在稠密开集 $V \subset Y$，使得 $f^{-1}(V) \to V$ 是有限的。
\end{lemma}

\begin{proof}
由于 $Y$ 只有有限多个不可约分支，可以找到一个稠密开集，它是不交的不可约分支之并。
因此不妨设 $Y$ 是具有泛点 $\eta$ 的不可约仿射概形。由于 $\eta$ 上的纤维有限，因为 $X^0$ 有限。

\medskip\noindent
假设 $f$ 分离且 $Y$ 不可约仿射。取 $V \subset Y$ 与 $U \subset X$，如引理
\ref{morphisms-lemma-generically-finite} 的 (3) 所述。由于 $f|_U : U \to V$ 有限，依据引理
\ref{morphisms-lemma-image-proper-scheme-closed} 与 \ref{morphisms-lemma-finite-proper}，
$U \subset f^{-1}(V)$ 既闭又开。因此有 $f^{-1}(V) = U \amalg W$，其中 $W$ 是
$X$ 的开子概形。然而，由于 $U$ 包含 $X$ 的所有泛点，得到 $W = \emptyset$，
如所需。

\medskip\noindent
假设 $f$ 拟紧且 $Y$ 不可约仿射。则 $X$ 拟紧，因而存在稠密开子概形 $U \subset X$，它是分离的
（性质引理 \ref{properties-lemma-quasi-compact-dense-open-separated}）。
由于泛点集合 $X^0$ 有限，得到 $X^0 \subset U$。于是
$\eta \not \in f(X \setminus U)$。由于 $X \setminus U \to Y$ 拟紧，得到一个非空开集
$V \subset Y$，使 $f^{-1}(V) \subset U$，见引理 \ref{morphisms-lemma-quasi-compact-dominant}。
以 $X$ 替换为 $f^{-1}(V)$，以 $Y$ 替换为 $V$，便约化为上一段已经处理的分离情形。
\end{proof}

\begin{lemma}
\label{morphisms-lemma-birational-isomorphism-over-dense-open}
设 $X$、$Y$ 是概形。设 $f : X \to Y$ 是双有理态射，且这两个概形只有有限多个不可约分支。
假设
\begin{enumerate}
\item $f$ 拟紧或 $f$ 分离；
\item $f$ 局部有限型且 $Y$ 既约，或 $f$ 局部有限表示。
\end{enumerate}
则存在稠密开集 $V \subset Y$，使得 $f^{-1}(V) \to V$ 是同构。
\end{lemma}

\begin{proof}
由引理 \ref{morphisms-lemma-finite-over-dense-open}，不妨设 $f$ 是有限的。
由于 $Y$ 只有有限多个不可约分支，可以找到一个稠密开集，它是不交的不可约分支之并。
因此不妨设 $Y$ 不可约。由引理 \ref{morphisms-lemma-birational-birational}，存在非空开集
$U \subset X$，使 $f|_U : U \to Y$ 是开浸入。去掉闭子集
$f$ 有限时去掉闭子集 $f(X \setminus U)$，它来自 $Y$，后得到 $f$ 是同构。
\end{proof}
\begin{lemma}
\label{morphisms-lemma-finite-degree}
设 $X$、$Y$ 是整概形。设 $f : X \to Y$ 是局部有限型的。
假设 $f$ 支配。以下条件等价：
\begin{enumerate}
\item 扩张 $R(Y) \subset R(X)$ 的超越次数为 $0$；
\item 扩张 $R(Y) \subset R(X)$ 是有限的；
\item 存在非空仿射开集 $U \subset X$ 与 $V \subset Y$，使 $f(U) \subset V$，
  且 $f|_U : U \to V$ 是有限的；
\item $X$ 的泛点是 $X$ 中映到 $Y$ 的泛点的唯一点。
\end{enumerate}
若 $f$ 分离或 $f$ 拟紧，则这些条件还等价于
\begin{enumerate}
\item[(5)] 存在非空仿射开集 $V \subset Y$，使得
  $f^{-1}(V) \to V$ 是有限的。
\end{enumerate}
\end{lemma}

\begin{proof}
取任意仿射开集 $\Spec(A) = U \subset X$ 与 $\Spec(R) = V \subset Y$，使 $f(U) \subset V$。
按定义，$R$ 与 $A$ 是整环。环映射 $R \to A$ 是有限型的
（见引理 \ref{morphisms-lemma-locally-finite-type-characterize}）。由引理
\ref{morphisms-lemma-dominant-finite-number-irreducible-components}，$X$ 的泛点映到 $Y$ 的泛点，
因此 $R \to A$ 是单射。令 $K = R(Y)$ 为 $R$ 的分式域，令 $L = R(X)$ 为 $A$ 的分式域。
于是 $L/K$ 是有限生成域扩张，从而 (1) 等价于 (2)。

\medskip\noindent
假设 (2) 成立。设 $x_1, \ldots, x_n \in A$ 是 $A$ 在 $R$ 上的生成元。
根据假设，存在非零多项式 $P_i(X) \in R[X]$，使 $P_i(x_i) = 0$。
令 $f_i \in R$ 为 $P_i$ 的首项系数。于是
$R_{f_1 \ldots f_n} \to A_{f_1 \ldots f_n}$ 是有限的，即 (3) 成立。
注意 (3) 蕴含 (2)。因此 (1)、(2)、(3) 全部等价。

\medskip\noindent
令 $\eta$ 是 $X$ 的泛点，并令 $\eta' \in Y$ 是 $Y$ 的泛点。假设 (4)。则
$\dim_\eta(X_{\eta'}) = 0$，由引理 \ref{morphisms-lemma-dimension-fibre-at-a-point} 可见
$R(X) = \kappa(\eta)$ 的超越次数为 $0$，相对于 $R(Y) = \kappa(\eta')$。
换言之，(1) 成立。现在假设等价条件 (1)、(2)、(3) 成立。设 $x \in X$ 是映到 $\eta'$ 的点。
由于 $x$ 是 $\eta$ 的特化点，得到包含关系
$R(Y) \subset \mathcal{O}_{X, x} \subset R(X)$，这意味着 $\mathcal{O}_{X, x}$ 是域，见
代数引理 \ref{algebra-lemma-integral-over-field}。因此 $x = \eta$。于是 (1)--(4) 全部等价。

\medskip\noindent
显然 (5) 在不对 $f$ 作任何额外假设时蕴含 (3)。剩下要证明的是：若 $f$ 分离或拟紧，
则等价条件 (1)--(4) 蕴含 (5)。这由引理 \ref{morphisms-lemma-finite-over-dense-open} 得出。
\end{proof}
\begin{definition}
\label{morphisms-definition-degree}
令 $X$ 与 $Y$ 为整概形。
令 $f : X \to Y$ 为局部有限型且支配的态射。
假设 $[R(X) : R(Y)] < \infty$，或引理 \ref{morphisms-lemma-finite-degree} 的等价条件 (1)--(4) 中的任一条件。
则正整数
$$
\deg(X/Y) = [R(X) : R(Y)]
$$
称为 $X$ 在 $Y$ 上的\emph{次数}。
\end{definition}

\noindent
可以将这个概念推广到态射
$f : X \to Y$，如果 (a) $Y$ 是带泛点 $\eta$ 的整概形，(b) $f$ 是局部有限型的，并且 (c) $f^{-1}(\{\eta\})$ 是有限的。
在此情形，我们定义
$$
\deg(X/Y)
=
\sum\nolimits_{\xi \in X, \ f(\xi) = \eta}
\dim_{R(Y)} (\mathcal{O}_{X, \xi}).
$$
也就是说，由于 $R(Y) = \kappa(\eta) = \mathcal{O}_{Y, \eta}$
（引理 \ref{morphisms-lemma-integral-scheme-rational-functions}），上述维数由上面的引理
\ref{morphisms-lemma-generically-finite} 可知都是有限的。
不过，对大多数应用而言，上面给出的定义才是适用的定义。

\begin{lemma}
\label{morphisms-lemma-degree-composition}
令 $X$、$Y$、$Z$ 为整概形。
令 $f : X \to Y$ 与 $g : Y \to Z$ 为局部有限型的支配态射。假设 $[R(X) : R(Y)] < \infty$ 且
$[R(Y) : R(Z)] < \infty$。则
$$
\deg(X/Z) = \deg(X/Y) \deg(Y/Z).
$$
\end{lemma}

\begin{proof}
这是域的有限扩张塔中次数的可乘性所致，见
《域》，引理 \ref{fields-lemma-multiplicativity-degrees}。
\end{proof}

\begin{remark}
\label{morphisms-remark-definition-generically-finite}
令 $f : X \to Y$ 为局部有限型的概形态射。
对于定义\emph{泛有限}态射，至少有两种可采用的性质。这两种性质分别对应于希望该性质在源上局部成立还是在目标上局部成立：
\begin{enumerate}
\item（目标上的局部性质；Ravi Vakil 建议。）
假设 $Y$ 的每个拟紧开子集都有有限多个不可约分支
（例如 $Y$ 局部 Noether）。
要求每个泛点的逆像都是有限的，见引理 \ref{morphisms-lemma-generically-finite}。
\item（源上的局部性质。）要求存在稠密开集 $U \subset X$，使得 $U \to Y$ 是局部拟有限的。
\end{enumerate}
在情形 (1) 中，不借助关于 $Y$ 的辅助条件也可以表述该要求，但对一般概形而言这可能不是正确的概念。按 (2) 表述的性质并不蕴含泛点上的纤维有限；但是，如果 $f$ 是拟紧的且 $Y$ 满足 (1) 中的条件，则确实蕴含这一点。
\end{remark}

\begin{definition}
\label{morphisms-definition-modification}
令 $X$ 为整概形。$X$ 的一个\emph{修改}是一个双有理的真态射 $f : X' \to X$，其中 $X'$ 为整概形。
\end{definition}

\noindent
令 $f : X' \to X$ 是定义中的修改。由引理 \ref{morphisms-lemma-finite-degree}，存在非空的 $U \subset X$，使得 $f^{-1}(U) \to U$ 是有限的。由一般平坦性
（命题 \ref{morphisms-proposition-generic-flatness}），我们可以假设 $f^{-1}(U) \to U$ 是平坦且有限表示的。
因此 $f^{-1}(U) \to U$ 是有限局部自由的（引理 \ref{morphisms-lemma-finite-flat}）。
由于 $f$ 是双有理的，$X'$ 在 $X$ 上的次数为 $1$。
所以 $f^{-1}(U) \to U$ 是次数为 $1$ 的有限局部自由态射，换言之，它是一个同构。
因此，我们可以将\emph{修改}重新定义为整概形之间的真态射 $f : X' \to X$，使得 $f^{-1}(U) \to U$ 是同构，其中 $U \subset X$ 是某个非空开集。

\begin{definition}
\label{morphisms-definition-alteration}
\begin{reference}
\cite[Definition 2.20]{alterations}
\end{reference}
令 $X$ 为整概形。$X$ 的一个\emph{变更（alteration）}是一个真支配态射 $f : Y \to X$，其中 $Y$ 为整概形，并且 $f^{-1}(U) \to U$ 是有限的，其中 $U \subset X$ 是某个非空开集。
\end{definition}

\noindent
这是 \cite{alterations} 中给出的定义，区别在于这里不要求 $X$ 与 $Y$ 是 Noether 概形。按照上面的论证，我们看到，变更是整概形之间的真支配态射 $f : Y \to X$，它诱导函数域的有限扩张，即满足引理 \ref{morphisms-lemma-finite-degree} 的等价条件。

\section{维数公式}
\label{morphisms-section-dimension-formula}

\noindent
对于 Noether 概形之间的态射，我们可以进一步说明局部环的维数。下面是一个重要（且不难证明）的结果。回忆 $R(X)$ 表示整概形 $X$ 的函数域。

\begin{lemma}
\label{morphisms-lemma-dimension-formula}
设 $S$ 为概形。
设 $f : X \to S$ 为概形态射。
令 $x \in X$，并置 $s = f(x)$。
假设
\begin{enumerate}
\item $S$ 是局部 Noether 的；
\item $f$ 是局部有限型的；
\item $X$ 与 $S$ 是整概形；并且
\item $f$ 是支配的。
\end{enumerate}
我们有
\begin{equation}
\label{morphisms-equation-dimension-formula}
\dim(\mathcal{O}_{X, x})
\leq
\dim(\mathcal{O}_{S, s}) + \text{trdeg}_{R(S)}R(X)
- \text{trdeg}_{\kappa(s)} \kappa(x).
\end{equation}
此外，如果 $S$ 是普遍链状的，则等号成立。
\end{lemma}

\begin{proof}
相应的代数表述是《代数》中的引理 \ref{algebra-lemma-dimension-formula}。
\end{proof}

\begin{lemma}
\label{morphisms-lemma-dimension-formula-general}
设 $S$ 为概形。设 $f : X \to S$ 为概形态射。
令 $x \in X$，并置 $s = f(x)$。假设 $S$ 是局部 Noether 的且 $f$ 是局部有限型的，
我们有
\begin{equation}
\label{morphisms-equation-dimension-formula-general}
\dim(\mathcal{O}_{X, x})
\leq
\dim(\mathcal{O}_{S, s}) + E - \text{trdeg}_{\kappa(s)} \kappa(x).
\end{equation}
其中 $E$ 是所有 $\text{trdeg}_{\kappa(f(\xi))}(\kappa(\xi))$ 的最大值，
其中 $\xi$ 遍历 $X$ 中包含 $x$ 的不可约分支的泛点。
\end{lemma}

\begin{proof}
令 $X_1, \ldots, X_n$ 为 $X$ 的不可约分支中包含 $x$ 的那些，并赋予它们约化的诱导概形结构。
这些分支对应于 $\mathfrak q_i$（即 $\mathcal{O}_{X, x}$ 的极小素理想），因此它们只有有限多个
（《概形》引理 \ref{schemes-lemma-specialize-points} 与《代数》引理
\ref{algebra-lemma-Noetherian-irreducible-components}）。
于是
$\dim(\mathcal{O}_{X, x}) = \max \dim(\mathcal{O}_{X, x}/\mathfrak q_i)
= \max \dim(\mathcal{O}_{X_i, x})$。
在定义中出现的 $\xi$ 恰好是 $E$ 所取的那些泛点，即不可约分支的泛点 $\xi_i \in X_i$。令
$Z_i = \overline{\{f(\xi_i)\}} \subset S$，并赋予它约化的诱导概形结构。
$X_i \to X \to S$ 的复合经过 $Z_i$ 分解
（《概形》引理 \ref{schemes-lemma-map-into-reduction}）。因此可将维数公式（引理
\ref{morphisms-lemma-dimension-formula}）应用于
$\dim(\mathcal{O}_{X_i, x}) \leq \dim(\mathcal{O}_{Z_i, x}) +
\text{trdeg}_{\kappa(f(\xi))}(\kappa(\xi)) -
\text{trdeg}_{\kappa(s)} \kappa(x)$。
汇总以上各式即得本引理。
\end{proof}
\noindent
一个应用是：在赋有维数函数的普遍链状概形上有限型的任意概形上构造维数函数。关于维数函数的定义，见
《拓扑》定义 \ref{topology-definition-dimension-function}。

\begin{lemma}
\label{morphisms-lemma-dimension-function-propagates}
令 $S$ 为局部 Noether 且普遍链状的概形。
令 $\delta : S \to \mathbf{Z}$ 为维数函数。
令 $f : X \to S$ 为概形态射。
假设 $f$ 局部有限型。
则映射
\begin{align*}
\delta = \delta_{X/S} : X & \longrightarrow \mathbf{Z} \\
x & \longmapsto \delta(f(x)) + \text{trdeg}_{\kappa(f(x))} \kappa(x)
\end{align*}
是 $X$ 上的维数函数。
\end{lemma}

\begin{proof}
令 $f : X \to S$ 为局部有限型的。
令 $x \leadsto y$、$x \not = y$ 是 $X$ 中的一个特化。
我们需要证明 $\delta_{X/S}(x) > \delta_{X/S}(y)$，并且 $\delta_{X/S}(x) = \delta_{X/S}(y) + 1$（当 $y$ 是 $x$ 的直接特化时）。

\medskip\noindent
选取仿射开集 $V \subset S$，其包含 $y$ 的像；并选取仿射开集 $U \subset X$，它映入 $V$ 且包含 $y$。
显然可以将 $X$ 换成 $U$，将 $S$ 换成 $V$。因此不妨设 $X = \Spec(A)$、
$S = \Spec(R)$，且 $f$ 由环映射 $R \to A$ 给出。环 $R$ 是普遍链状的
（引理 \ref{morphisms-lemma-universally-catenary-local}），且环映射 $R \to A$ 是有限型的
（引理 \ref{morphisms-lemma-locally-finite-type-characterize}）。

\medskip\noindent
令 $\mathfrak q \subset A$ 为对应于点 $x$ 的素理想，令 $\mathfrak p \subset R$ 为对应于 $f(x)$ 的素理想。
将 $\delta'$ 定义为 $\delta$ 在 $S' = \Spec(R/\mathfrak p) \subset S$ 上的限制；它是维数函数。
环 $R/\mathfrak p$ 是普遍链状的。
将 $\delta_{X/S}$ 限制到 $X' = \Spec(A/\mathfrak q)$，显然等于函数 $\delta_{X'/S'}$，后者由维数函数 $\delta'$ 构造。
因此，除上述假设外，我们还可以假设 $R \subset A$ 是整环，换言之，$X$ 与 $S$ 是整概形，且 $x$ 是 $X$ 的泛点、$f(x)$ 是 $S$ 的泛点。

\medskip\noindent
注意 $\mathcal{O}_{X, x} = R(X)$，并且由于 $x \leadsto y$ 且 $x \not = y$，环 $\mathcal{O}_{X, y}$ 的谱至少有两个点
（《概形》引理 \ref{schemes-lemma-specialize-points}），因此 $\dim(\mathcal{O}_{X, y}) > 0$。
如果 $y$ 是 $x$ 的直接特化，则 $\Spec(\mathcal{O}_{X, y}) = \{x, y\}$ 且
$\dim(\mathcal{O}_{X, y}) = 1$。

\medskip\noindent
写 $s = f(x)$，$t = f(y)$。我们计算
\begin{align*}
\delta_{X/S}(x) - \delta_{X/S}(y)
& =
\delta(s) + \text{trdeg}_{\kappa(s)} \kappa(x)
- \delta(t) - \text{trdeg}_{\kappa(t)} \kappa(y) \\
& =
\delta(s) - \delta(t) +
\text{trdeg}_{R(S)} R(X) - \text{trdeg}_{\kappa(t)} \kappa(y) \\
& =
\delta(s) - \delta(t) + \dim(\mathcal{O}_{X, y})
- \dim(\mathcal{O}_{S, t})
\end{align*}
其中最后一步使用了（\ref{morphisms-equation-dimension-formula}）中的等式。
由于 $\delta$ 是概形 $S$ 上的维数函数，且 $s \in S$ 是泛点，差值
$\delta(s) - \delta(t)$ 由《拓扑》引理
\ref{topology-lemma-dimension-function-catenary} 可知等于 $\text{codim}(\overline{\{t\}}, S)$。
由《性质》引理
\ref{properties-lemma-codimension-local-ring}，这等于 $\dim(\mathcal{O}_{S, t})$。
因此我们得到
$$
\delta_{X/S}(x) - \delta_{X/S}(y) = \dim(\mathcal{O}_{X, y})
$$
结合上面对 $\dim(\mathcal{O}_{X, y})$ 的说明即得本引理。
\end{proof}
\noindent
维数公式的另一个应用是：维数在“变更”（稍后定义）下不变。

\begin{lemma}
\label{morphisms-lemma-alteration-dimension}
令 $f : X \to Y$ 为概形态射。假设
\begin{enumerate}
\item $Y$ 是局部 Noether 的；
\item $X$ 与 $Y$ 是整概形；
\item $f$ 是支配的；并且
\item $f$ 是局部有限型的。
\end{enumerate}
则
$$
\dim(X) \leq \dim(Y) + \text{trdeg}_{R(Y)} R(X).
$$
如果 $f$ 是闭的\footnote{例如若 $f$ 是真态射，见《概形》定义
\ref{morphisms-definition-proper}。}，则等号成立。
\end{lemma}

\begin{proof}
令 $f : X \to Y$ 如引理中所述。
令 $\xi_0 \leadsto \xi_1 \leadsto \ldots \leadsto \xi_e$ 是 $X$ 中的一列特化。置 $x = \xi_e$，
$y = f(x)$。注意 $e \leq \dim(\mathcal{O}_{X, x})$，因为给出的特化发生在
$\mathcal{O}_{X, x}$ 的谱中，见《概形》引理
\ref{schemes-lemma-specialize-points}。
由维数公式（引理 \ref{morphisms-lemma-dimension-formula}），我们看到
\begin{align*}
e & \leq \dim(\mathcal{O}_{X, x}) \\
& \leq \dim(\mathcal{O}_{Y, y}) +
\text{trdeg}_{R(Y)} R(X) - \text{trdeg}_{\kappa(y)} \kappa(x) \\
& \leq \dim(\mathcal{O}_{Y, y}) + \text{trdeg}_{R(Y)} R(X)
\end{align*}
因此如愿得到 $e \leq \dim(Y) + \text{trdeg}_{R(Y)} R(X)$。

\medskip\noindent
现在再假设 $f$ 也是闭的。
设 $\overline{\xi}_0 \leadsto \overline{\xi}_1 \leadsto \ldots
\leadsto \overline{\xi}_d$ 是 $Y$ 中的一列特化。
我们要证明 $\dim(X) \geq d + r$。
不妨设 $\overline{\xi}_0 = \eta$ 是 $Y$ 的泛点。
$X_\eta$ 是在 $\kappa(\eta) = R(Y)$ 上局部有限型的概形。由于 $f$ 支配，它非空。因此由引理
\ref{morphisms-lemma-ubiquity-Jacobson-schemes}，它是 Jacobson 概形。
于是由引理 \ref{morphisms-lemma-jacobson-finite-type-points}，可以找到闭点
$\xi_0 \in X_\eta$，并且扩张 $\kappa(\eta) \subset \kappa(\xi_0)$ 是有限扩张。
注意 $\mathcal{O}_{X, \xi_0} = \mathcal{O}_{X_\eta, \xi_0}$，因为 $\eta$ 是 $Y$ 的泛点。
因此，由维数公式引理 \ref{morphisms-lemma-dimension-formula} 可知 $\dim(\mathcal{O}_{X, \xi_0}) = r$；该引理应用于概形 $X_\eta$（它在普遍链状概形
$\Spec(\kappa(\eta))$ 上；见引理 \ref{morphisms-lemma-ubiquity-uc}）及点 $\xi_0$。
这意味着我们可以找到
$\xi_{-r} \leadsto \ldots \leadsto \xi_{-1} \leadsto \xi_0$ 于 $X$ 中。
另一方面，由于 $f$ 是闭的，特化沿着 $f$ 提升，见
《拓扑》引理 \ref{topology-lemma-closed-open-map-specialization}。
因此，由于 $\xi_0$ 位于
$\eta = \overline{\xi}_0$ 上方，我们可以找到特化
$\xi_0 \leadsto \xi_1 \leadsto \ldots \leadsto \xi_d$，它们分别位于
$\overline{\xi}_0 \leadsto \overline{\xi}_1 \leadsto \ldots
\leadsto \overline{\xi}_d$ 上方。换言之，我们有
$$
\xi_{-r} \leadsto \ldots \leadsto \xi_{-1} \leadsto \xi_0
\leadsto \xi_1 \leadsto \ldots \leadsto \xi_d
$$
这意味着 $\dim(X) \geq d + r$，如愿以偿。
\end{proof}

\begin{lemma}
\label{morphisms-lemma-alteration-dimension-general}
令 $f : X \to Y$ 为概形态射。假设 $Y$ 是局部 Noether 的且 $f$ 是局部有限型的。
则
$$
\dim(X) \leq \dim(Y) + E
$$
其中 $E$ 是所有 $\text{trdeg}_{\kappa(f(\xi))}(\kappa(\xi))$ 的上确界，
其中 $\xi$ 遍历 $X$ 的不可约分支的泛点。
\end{lemma}

\begin{proof}
这是引理 \ref{morphisms-lemma-dimension-formula-general} 与《性质》引理
\ref{properties-lemma-dimension} 的直接推论。
\end{proof}
\section{相对正规化}
\label{morphisms-section-normalization-X-in-Y}

\noindent
本节构造“一个概形在另一个概形中的正规化”。

\begin{lemma}
\label{morphisms-lemma-integral-closure}
设 $X$ 为概形。设 $\mathcal{A}$ 为一个准凝聚的 $\mathcal{O}_X$-代数层。按下列规则定义的子层
$\mathcal{A}' \subset \mathcal{A}$：
$$
U \longmapsto \{f \in \mathcal{A}(U) \mid
f_x \in \mathcal{A}_x \text{ integral over } \mathcal{O}_{X, x}
\text{ 对所有 }x \in U\}
$$
是准凝聚的 $\mathcal{O}_X$-代数层，茎 $\mathcal{A}'_x$ 是 $\mathcal{O}_{X, x}$ 在
$\mathcal{A}_x$ 中的整闭包，并且对任意仿射开集 $U \subset X$，环
$\mathcal{A}'(U) \subset \mathcal{A}(U)$ 是 $\mathcal{O}_X(U)$ 在
$\mathcal{A}(U)$ 中的整闭包。
\end{lemma}

\begin{proof}
由于这些条件具有局部性质，这是一个子层。
由《代数》引理 \ref{algebra-lemma-integral-closure-is-ring}，它是一个 $\mathcal{O}_X$-代数层。
令 $U \subset X$ 为仿射开集。设 $U = \Spec(R)$，并设 $\mathcal{A}$ 是与
$R$-代数 $A$ 关联的准凝聚层。根据
《代数》引理 \ref{algebra-lemma-integral-closure-stalks}，$\mathcal{A}'$ 在
$U$ 上的值由 $A'$（即 $R$ 在 $A$ 中的整闭包）给出。这证明了引理的最后一个断言。
为证明 $\mathcal{A}'$ 准凝聚，只需证明 $\mathcal{A}'(D(f)) = A'_f$。
这由整闭包与局部化可交换这一事实得出，见
《代数》引理 \ref{algebra-lemma-integral-closure-localize}。
同一事实也表明茎具有所述形式。
\end{proof}

\begin{definition}
\label{morphisms-definition-integral-closure}
设 $X$ 为概形。设 $\mathcal{A}$ 为准凝聚的 $\mathcal{O}_X$-代数层。$\mathcal{O}_X$
在 $\mathcal{A}$ 中的\emph{整闭包}是上面引理 \ref{morphisms-lemma-integral-closure} 中构造的
准凝聚 $\mathcal{O}_X$-子代数 $\mathcal{A}' \subset \mathcal{A}$。
\end{definition}

\noindent
在上述定义的情形下，我们可以考虑相对谱的态射
$$
\xymatrix{
Y = \underline{\Spec}_X(\mathcal{A}) \ar[rr] \ar[rd] & &
X' = \underline{\Spec}_X(\mathcal{A}') \ar[ld] \\
& X &
}
$$
见引理 \ref{morphisms-lemma-affine-equivalence-algebras}。
态射 $X' \to X$ 将是 $X$ 在概形 $Y$ 中的正规化。
下面给出稍微一般的情形。设我们有概形的拟紧且拟分离态射 $f : Y \to X$。
在此情形，$\mathcal{O}_X$-代数层 $f_*\mathcal{O}_Y$ 是准凝聚的，见
《概形》引理 \ref{schemes-lemma-push-forward-quasi-coherent}。
取 $\mathcal{O}' \subset f_*\mathcal{O}_Y$ 的整闭包，得到准凝聚的 $\mathcal{O}_X$-代数层，
其相对谱就是 $X$ 在 $Y$ 中的正规化。下面是形式定义。

\begin{definition}
\label{morphisms-definition-normalization-X-in-Y}
设 $f : Y \to X$ 为概形的拟紧且拟分离态射。
令 $\mathcal{O}'$ 为 $\mathcal{O}_X$ 在 $f_*\mathcal{O}_Y$ 中的整闭包。$X$ 在 $Y$ 中的\emph{正规化}是概形
\footnote{概形 $X'$ 不必是正规概形；例如若 $Y = X$ 且 $f = \text{id}_X$，则 $X' = X$。}
$$
\nu : X' = \underline{\Spec}_X(\mathcal{O}') \to X
$$
它位于 $X$ 上。它带有初始态射 $f$ 的自然分解
$$
Y \xrightarrow{f'} X' \xrightarrow{\nu} X
$$
。\end{definition}

\noindent
该分解是典范态射
$Y \to \underline{\Spec}_X(f_*\mathcal{O}_Y)$（见《构造》引理
\ref{constructions-lemma-canonical-morphism}）与由包含映射
$\mathcal{O}' \to f_*\mathcal{O}_Y$ 产生的相对谱态射的复合。我们可以如下刻画正规化。

\begin{lemma}
\label{morphisms-lemma-characterize-normalization}
设 $f : Y \to X$ 为概形的拟紧且拟分离态射。
分解 $f = \nu \circ f'$（其中 $\nu : X' \to X$ 是 $X$ 在 $Y$ 中的正规化）由以下两个性质刻画：
\begin{enumerate}
\item 态射 $\nu$ 是整的；并且
\item 对任意分解 $f = \pi \circ g$（其中 $\pi : Z \to X$ 是整的），存在交换图
$$
\xymatrix{
Y \ar[d]_{f'} \ar[r]_g & Z \ar[d]^\pi \\
X' \ar[ru]^h \ar[r]^\nu & X
}
$$
以及唯一的态射 $h : X' \to Z$。
\end{enumerate}
此外，态射 $f' : Y \to X'$ 是支配的，并且在 (2) 中态射 $h : X' \to Z$ 是 $Z$ 在 $Y$ 中的正规化。
\end{lemma}
\begin{proof}
令 $\mathcal{O}' \subset f_*\mathcal{O}_Y$ 为定义 \ref{morphisms-definition-normalization-X-in-Y} 中的
$\mathcal{O}_X$ 的整闭包。
按构造，态射 $\nu$ 是整的，这证明了 (1)。
假设给定分解 $f = \pi \circ g$，其中 $\pi : Z \to X$ 如 (2) 中是整的。
由定义 \ref{morphisms-definition-integral}，$\pi$ 是仿射的，因此 $Z$ 是某个准凝聚的
$\mathcal{O}_X$-代数层 $\mathcal{B}$ 的相对谱。
态射 $g : Y \to Z$ 对应于一个 $\mathcal{O}_X$-代数映射
$\chi : \mathcal{B} \to f_*\mathcal{O}_Y$。由于 $\mathcal{B}(U)$ 都整于 $\mathcal{O}_X(U)$，其中 $U \subset X$ 为任意仿射开集
（由定义 \ref{morphisms-definition-integral}），
由引理 \ref{morphisms-lemma-integral-closure} 可知
$\chi(\mathcal{B}) \subset \mathcal{O}'$。
由相对谱的函子性及引理 \ref{morphisms-lemma-affine-equivalence-algebras}，这给出唯一的态射
$h : X' \to Z$。交换图的验证略去。

\medskip\noindent
(1) 与 (2) 显然刻画分解 $f = \nu \circ f'$，因为它将该分解刻画为范畴中的始对象。

\medskip\noindent
由 (2) 的泛性质可见，$f'$ 不经过 $X'$ 的真闭子概形。因此 $f'$ 的概形论像是 $X'$。
由于 $f'$ 是拟紧的（由《概形》引理
\ref{schemes-lemma-quasi-compact-permanence} 以及 $\nu$ 作为仿射态射而分离这一事实），
我们看到 $f'(Y)$ 在 $X'$ 中稠密。因此 $f'$ 是支配的。

\medskip\noindent
注意，由《概形》引理 \ref{schemes-lemma-compose-after-separated} 与
\ref{schemes-lemma-quasi-compact-permanence}，$g$ 是拟紧且拟分离的。因此引理的最后一个断言是有意义的。
由引理 \ref{morphisms-lemma-finite-permanence}，(2) 中的态射 $h$ 是整的。
给定分解 $g = \pi' \circ g'$，其中 $\pi' : Z' \to Z$ 是整的，我们得到分解
$f = (\pi \circ \pi') \circ g'$，并得到态射 $h' : X' \to Z'$。
由唯一性，$\pi' \circ h' = h$。因此，刻画 (1)、(2) 适用于态射
$h : X' \to Z$，从而得到引理的最后一个断言。
\end{proof}

\begin{lemma}
\label{morphisms-lemma-functoriality-normalization}
令
$$
\xymatrix{
Y_2 \ar[d]_{f_2} \ar[r] & Y_1 \ar[d]^{f_1} \\
X_2 \ar[r] & X_1
}
$$
为概形态射的交换图。
假设 $f_1$、$f_2$ 拟紧且拟分离。
令 $f_i = \nu_i \circ f_i'$，$i = 1, 2$，为典范分解，其中
$\nu_i : X_i' \to X_i$ 是 $X_i$ 在 $Y_i$ 中的正规化。则存在唯一的
箭头 $X'_2 \to X'_1$，使其适合于一个交换图
$$
\xymatrix{
Y_2 \ar[d]_{f_2'} \ar[r] & Y_1 \ar[d]^{f_1'} \\
X_2' \ar[d]_{\nu_2} \ar[r] & X_1' \ar[d]^{\nu_1} \\
X_2 \ar[r] & X_1
}
$$
\end{lemma}

\begin{proof}
由引理 \ref{morphisms-lemma-characterize-normalization} (1) 与
\ref{morphisms-lemma-base-change-finite}，基变换 $X_2 \times_{X_1} X'_1 \to X_2$
是整的。注意 $f_2$ 经过这个态射分解。
因此，由引理 \ref{morphisms-lemma-characterize-normalization} (2)，我们得到唯一的态射
$X'_2 \to X_2 \times_{X_1} X'_1$。
这给出了适合交换图的箭头 $X'_2 \to X'_1$，唯一性也随之成立。
\end{proof}

\begin{lemma}
\label{morphisms-lemma-normalization-localization}
设 $f : Y \to X$ 为概形的拟紧且拟分离态射。
令 $U \subset X$ 为开子概形，并置 $V = f^{-1}(U)$。
则 $U$ 在 $V$ 中的正规化是 $U$ 在 $X$ 中的 $Y$ 的正规化的逆像。
\end{lemma}

\begin{proof}
由构造立即可见。
\end{proof}

\begin{lemma}
\label{morphisms-lemma-normalization-is-normalization}
设 $f : Y \to X$ 为概形的拟紧且拟分离态射。
令 $X'$ 为 $X$ 在 $Y$ 中的正规化。则 $X'$ 在 $Y$ 中的正规化就是 $X'$。
\end{lemma}

\begin{proof}
若 $Y \to X'' \to X'$ 是 $X'$ 在 $Y$ 中的正规化，则可将引理
\ref{morphisms-lemma-characterize-normalization} 应用于复合 $X'' \to X$，从而得到典范态射
$h : X' \to X''$（在 $X$ 上）。我们略去验证态射 $h$ 与 $X'' \to X'$ 互为逆态射（使用引理中分解的唯一性）。
\end{proof}

\begin{lemma}
\label{morphisms-lemma-normalization-in-reduced}
设 $f : Y \to X$ 为概形的拟紧且拟分离态射。
设 $X' \to X$ 为 $X$ 在 $Y$ 中的正规化。若 $Y$ 约化，则 $X'$ 也约化。
\end{lemma}

\begin{proof}
这是约化环的子环仍为约化环这一事实的直接结果。
细节略。
\end{proof}
\begin{lemma}
\label{morphisms-lemma-normalization-generic}
设 $f : Y \to X$ 为概形的拟紧且拟分离态射。
设 $X' \to X$ 为 $X$ 在 $Y$ 中的正规化。$X'$ 的每个不可约分支的泛点都是 $Y$ 的某个不可约分支的泛点的像。
\end{lemma}

\begin{proof}
由引理 \ref{morphisms-lemma-normalization-localization}，不妨设 $X = \Spec(A)$ 是仿射的。
选取有限仿射开覆盖 $Y = \bigcup \Spec(B_i)$。
于是 $X' = \Spec(A')$，态射 $\Spec(B_i) \to Y \to X'$ 共同给出单射的 $A$-代数映射
$A' \to \prod B_i$。
因此引理由《代数》引理 \ref{algebra-lemma-injective-minimal-primes-in-image} 得出。
\end{proof}

\begin{lemma}
\label{morphisms-lemma-normalization-in-disjoint-union}
设 $f : Y \to X$ 为概形的拟紧且拟分离态射。
假设 $Y = Y_1 \amalg Y_2$ 是两个概形的不交并。
记 $f_i = f|_{Y_i}$。令 $X_i'$ 为 $X$ 在 $Y_i$ 中的正规化。
则 $X_1' \amalg X_2'$ 是 $X$ 在 $Y$ 中的正规化。
\end{lemma}

\begin{proof}
用整闭包的语言，这对应于以下事实：
令 $A \to B$ 为环映射。假设 $B = B_1 \times B_2$。
令 $A_i'$ 为 $A$ 在 $B_i$ 中的整闭包。则
$A_1' \times A_2'$ 是 $A$ 在 $B$ 中的整闭包。
之所以成立，是因为元素 $(1, 0)$ 与 $(0, 1)$ 属于 $B$，是幂等元，因而整于 $A$。
所以整闭包 $A'$ 是 $A$ 在 $B$ 中的一个乘积；不难看出，其因子正是上面所述的整闭包 $A'_i$（细节略）。
\end{proof}

\begin{lemma}
\label{morphisms-lemma-normalization-in-universally-closed}
设 $f : X \to S$ 为概形的拟紧、拟分离且万有闭态射。
则 $f_*\mathcal{O}_X$ 整于 $\mathcal{O}_S$。换言之，$S$ 在 $X$ 中的正规化等于构造中的引理
\ref{constructions-lemma-canonical-morphism} 给出的分解
$$
X \longrightarrow \underline{\Spec}_S(f_*\mathcal{O}_X)
\longrightarrow S
$$
。
\end{lemma}

\begin{proof}
问题在 $S$ 上是局部的，因此不妨设 $S = \Spec(R)$ 是仿射的。
令 $h \in \Gamma(X, \mathcal{O}_X)$。我们需要证明 $h$ 满足一个系数在 $R$ 中的首一方程。
将 $h$ 看作如下交换图中的态射：
$$
\xymatrix{
X \ar[rr]_h \ar[rd]_f & & \mathbf{A}^1_S \ar[ld] \\
& S &
}
$$
令 $Z \subset \mathbf{A}^1_S$ 为 $h$ 的概形论像，见定义 \ref{morphisms-definition-scheme-theoretic-image}。
态射 $h$ 是拟紧的，因为 $f$ 拟紧且 $\mathbf{A}^1_S \to S$ 分离，见
《概形》引理 \ref{schemes-lemma-quasi-compact-permanence}。
由引理 \ref{morphisms-lemma-quasi-compact-scheme-theoretic-image}，态射 $X \to Z$ 是支配的。
由引理 \ref{morphisms-lemma-image-proper-scheme-closed}，态射 $X \to Z$ 是闭的。因此（按集合论意义）$h(X) = Z$。
于是可使用引理 \ref{morphisms-lemma-image-proper-is-proper}，断定 $Z \to S$ 是万有闭的（甚至是固有的）。
由于 $Z \subset \mathbf{A}^1_S$，可见 $Z \to S$ 是仿射且固有的，因而由引理
\ref{morphisms-lemma-integral-universally-closed} 是整的。
写 $\mathbf{A}^1_S = \Spec(R[T])$，我们得到理想 $I \subset R[T]$（它是 $Z$ 的理想）
包含一个首一多项式 $P(T) \in R[T]$。
于是 $P(h) = 0$，证毕。
\end{proof}

\begin{lemma}
\label{morphisms-lemma-normalization-in-integral}
设 $f : Y \to X$ 为整态射。
则 $X$ 在 $Y$ 中的正规化等于 $Y$。
\end{lemma}

\begin{proof}
由引理 \ref{morphisms-lemma-integral-universally-closed}，这是引理
\ref{morphisms-lemma-normalization-in-universally-closed} 的特殊情形。
\end{proof}
\begin{lemma}
\label{morphisms-lemma-normal-normalization}
设 $f : Y \to X$ 为概形的拟紧且拟分离态射。令 $X'$ 为 $X$ 在 $Y$ 中的正规化。假设
\begin{enumerate}
\item $Y$ 是正规概形；
\item $Y$ 的拟紧开集具有有限多个不可约分支。
\end{enumerate}
则 $X'$ 是整正规概形的不交并。此外，态射 $Y \to X'$ 是支配的，并且在不可约分支之间诱导双射。
\end{lemma}

\begin{proof}
令 $U \subset X$ 为仿射开集。考虑 $U'$，它是 $U$ 在 $X'$ 中的逆像。
置 $V = f^{-1}(U)$。由引理 \ref{morphisms-lemma-normalization-localization}，
$V \to U' \to U$ 是 $U$ 在 $V$ 中的正规化。设 $U = \Spec(A)$。
则 $V$ 拟紧，因而由假设它有有限多个不可约分支。所以
$V = \coprod_{i = 1, \ldots n} V_i$ 是正规整概形的有限不交并，依据
《性质》引理 \ref{properties-lemma-normal-locally-finite-nr-irreducibles}。
由引理 \ref{morphisms-lemma-normalization-in-disjoint-union}，可见
$U' = \coprod_{i = 1, \ldots, n} U_i'$，其中 $U'_i$ 是 $U$ 在 $V_i$ 中的正规化。
由《性质》引理 \ref{properties-lemma-normal-integral-sections}，
$B_i = \Gamma(V_i, \mathcal{O}_{V_i})$ 是正规整环。
注意 $U_i' = \Spec(A_i')$，其中 $A_i' \subset B_i$
是 $A$ 在 $B_i$ 中的整闭包，见引理 \ref{morphisms-lemma-integral-closure}。由
《代数》引理 \ref{algebra-lemma-integral-closure-in-normal}，可见 $A_i' \subset B_i$ 是正规整环。
因此 $U' = \coprod U_i'$ 是正规整概形的有限不交并，因而是正规概形。

\medskip\noindent
由于 $X'$ 有由概形 $U'$ 构成的开覆盖，由《性质》引理
\ref{properties-lemma-locally-normal} 可知 $X'$ 是正规的。
另一方面，每个 $U'$ 都是不可约概形的有限不交并，因此 $X'$ 的每个拟紧开集都有有限多个不可约分支
（拓扑论证略）。所以由《性质》引理
\ref{properties-lemma-normal-locally-finite-nr-irreducibles}，$X'$ 是正规整概形的不交并。
从上面对 $X'$ 的描述显然可知，$Y \to X'$ 是支配的，并且在不可约分支之间诱导双射 $V \to U'$，对每个仿射开集 $U \subset X$。
态射 $Y \to X'$ 的不可约分支双射由此及拓扑论证得到（略）。
\end{proof}
\begin{lemma}
\label{morphisms-lemma-nagata-normalization-finite-general}
设 $f : X \to S$ 为态射。假设
\begin{enumerate}
\item $S$ 是 Nagata 概形；
\item $f$ 拟紧且拟分离；
\item $X$ 的拟紧开集具有有限多个不可约分支；
\item 若 $x \in X$ 是不可约分支的泛点，则域扩张
$\kappa(x)/\kappa(f(x))$ 是有限生成的；并且
\item $X$ 是约化的。
\end{enumerate}
则正规化 $\nu : S' \to S$（即 $S$ 在 $X$ 中的正规化）是有限的。
\end{lemma}

\begin{proof}
由假设 (1)，立即可将情形约化为 $S = \Spec(R)$，其中 $R$ 是 Nagata 环。
我们需要证明整闭包 $A$（即 $R$ 在 $\Gamma(X, \mathcal{O}_X)$ 中的整闭包）是有限的 $R$-模。
由假设 (2)，$f$ 拟紧，因此可写成 $X = \bigcup_{i = 1, \ldots, n} U_i$，其中每个 $U_i$ 都是仿射的。
设 $U_i = \Spec(B_i)$。由假设 (5)，每个 $B_i$ 都是约化的，并且由假设 (3) 以及
《代数》引理 \ref{algebra-lemma-irreducible}，它有有限多个极小素理想
$\mathfrak q_{i1}, \ldots, \mathfrak q_{im_i}$。
我们有
$$
\Gamma(X, \mathcal{O}_X) \subset B_1 \times \ldots \times B_n
\subset
\prod\nolimits_{i = 1, \ldots, n}
\prod\nolimits_{j = 1, \ldots, m_i} (B_i)_{\mathfrak q_{ij}}
$$
第二个包含关系来自《代数》引理 \ref{algebra-lemma-reduced-ring-sub-product-fields}。
由《代数》引理 \ref{algebra-lemma-minimal-prime-reduced-ring}，有
$\kappa(\mathfrak q_{ij}) = (B_i)_{\mathfrak q_{ij}}$。
因此，整闭包 $A$（即 $R$ 在 $\Gamma(X, \mathcal{O}_X)$ 中的整闭包）
包含于整闭包 $A_{ij}$（即 $R$ 在 $\kappa(\mathfrak q_{ij})$ 中的整闭包）的乘积中。
由于 $R$ 是 Noether 环，只需证明 $A_{ij}$ 是有限的 $R$-模，对每个 $i, j$ 均如此。
令 $\mathfrak p_{ij} \subset R$ 为 $\mathfrak q_{ij}$ 的像。
由假设 (4)，$\kappa(\mathfrak q_{ij})/\kappa(\mathfrak p_{ij})$
是有限生成的域扩张，因此 $R \to \kappa(\mathfrak q_{ij})$ 本质上是有限型的。
于是由《代数》引理
\ref{algebra-lemma-nagata-in-reduced-finite-type-finite-integral-closure}，$R \to A_{ij}$ 是有限的。
\end{proof}

\begin{lemma}
\label{morphisms-lemma-nagata-normalization-finite}
设 $f : X \to S$ 为态射。假设
\begin{enumerate}
\item $S$ 是 Nagata 概形；
\item $f$ 是有限型的；
\item $X$ 是约化的。
\end{enumerate}
则正规化 $\nu : S' \to S$（即 $S$ 在 $X$ 中的正规化）是有限的。
\end{lemma}

\begin{proof}
这是引理 \ref{morphisms-lemma-nagata-normalization-finite-general} 的特殊情形。
具体而言，(2) 成立，因为按定义有限型态射 $f$ 是拟紧的，并且由引理
\ref{morphisms-lemma-finite-type-Noetherian-quasi-separated} 是拟分离的。
条件 (3) 成立，因为 $X$ 局部 Noether，见引理 \ref{morphisms-lemma-finite-type-noetherian}。
最后，条件 (4) 成立，因为有限型态射诱导有限生成的剩余域扩张。
\end{proof}
\begin{lemma}
\label{morphisms-lemma-relative-normalization-normal-codim-1}
设 $f : Y \to X$ 为概形的有限型态射，其中 $Y$ 约化且 $X$ 为 Nagata 概形。令 $X'$ 为 $X$ 在 $Y$ 中的正规化。
令 $x' \in X'$ 为满足下列条件的点：
\begin{enumerate}
\item $\dim(\mathcal{O}_{X', x'}) = 1$，并且
\item $Y \to X'$ 在 $x'$ 上的纤维为空。
\end{enumerate}
则 $\mathcal{O}_{X', x'}$ 是离散赋值环。
\end{lemma}

\begin{proof}
可以将 $X$ 替换为包含 $x'$ 的像的仿射邻域。因此不妨设 $X = \Spec(A)$，其中 $A$ 是 Nagata 环。
由引理 \ref{morphisms-lemma-nagata-normalization-finite}，态射 $X' \to X$ 是有限的。因此可写成
$X' = \Spec(A')$，其中 $A$-代数 $A'$ 是有限的。
由引理 \ref{morphisms-lemma-normalization-is-normalization}，将 $X$ 替换为 $X'$ 后，
我们约化到下一段所述的情形。

\medskip\noindent
情形 $X = X' = \Spec(A)$，其中 $A$ 为 Noether 环。
令 $\mathfrak p \subset A$ 为对应于点 $x'$ 的素理想。
选取 $g \in \mathfrak p$，使其不包含于 $A$ 的任何极小素理想中（使用素理想避免以及 $A$ 只有有限多个极小素理想这一事实，见
《代数》引理 \ref{algebra-lemma-silly} 与
\ref{algebra-lemma-Noetherian-irreducible-components}）。
置 $Z = f^{-1}V(g) \subset Y$；它是 $Y$ 的闭子概形。
由于 $f(Z)$ 按 $g$ 的选择不包含任何泛点，且不包含 $x'$，因为 $x'$ 不在 $f$ 的像中。
$f(Z)$ 的闭包是 $f(Z)$ 中点的特化点集合，这是引理 \ref{morphisms-lemma-reach-points-scheme-theoretic-image} 的结论。
因此 $f(Z)$ 的闭包不包含 $x'$，因为条件 $\dim(\mathcal{O}_{X', x'}) = 1$ 蕴含只有 $X = X'$ 的泛点特化到 $x'$。
换言之，将 $X$ 替换为包含 $x'$ 的一个仿射开邻域后，不妨设 $f^{-1}V(g) = \emptyset$。
于是 $g$ 映为 $Y$ 上的可逆整体函数，并得到分解
$$
A \to A_g \to \Gamma(Y, \mathcal{O}_Y)
$$
由于 $X = X'$，这意味着 $A$ 等于整闭包（即 $A$ 在 $A_g$ 中的整闭包）。由
《代数》引理 \ref{algebra-lemma-integral-closure-localize}，我们断定
$A_\mathfrak p$ 是 $A_\mathfrak p$ 的整闭包（位于 $A_\mathfrak p[1/g]$ 中）。
由 $g$ 的选择、$\dim(A_\mathfrak p) = 1$ 且 $A$ 约化，可知 $A_\mathfrak p[1/g]$
是域的有限乘积（包含于 $\mathfrak p$ 的极小素理想的剩余域之乘积）。因此由
《代数》引理 \ref{algebra-lemma-characterize-reduced-ring-normal}，$A_\mathfrak p$ 是正规的，证明完成。细节略。
\end{proof}

\section{正规化}
\label{morphisms-section-normalization}
\noindent
接下来讨论概形 $X$ 的正规化。只有当 $X$ 的不可约分支局部有限时，我们才定义并构造它。
令 $X$ 为一个概形，使每个拟紧开集都有有限多个不可约分支。令
$X^{(0)} \subset X$ 为 $X$ 的不可约分支的泛点集合。令
\begin{equation}
\label{morphisms-equation-generic-points}
f :
Y = \coprod\nolimits_{\eta \in X^{(0)}} \Spec(\kappa(\eta))
\longrightarrow
X
\end{equation}
为利用《概形》第 \ref{schemes-section-points} 节中的典范映射，将泛点包含到 $X$ 中的态射。
注意，由假设该态射是拟紧的，并且由于 $Y$ 是分离的而拟分离（见《概形》第
\ref{schemes-section-separation-axioms} 节）。

\begin{definition}
\label{morphisms-definition-normalization}
设 $X$ 为一个概形，使每个拟紧开集都有有限多个不可约分支。我们将 $X$ 的\emph{正规化}定义为态射
$$
\nu : X^\nu \longrightarrow X
$$
它是 $X$ 在上面构造的态射 $f : Y \to X$（\ref{morphisms-equation-generic-points}）中的正规化。
\end{definition}

\noindent
任何局部 Noether 概形都有局部有限的不可约分支集合，因此该定义适用。
通常只对约化概形定义正规化。按上述定义，$X$ 的正规化与约化 $X_{red}$（它来自 $X$）的正规化相同。

\begin{lemma}
\label{morphisms-lemma-normalization-reduced}
设 $X$ 为一个概形，使每个拟紧开集都有有限多个不可约分支。正规化态射
$\nu$ 经过约化 $X_{red}$ 分解，并且 $X^\nu \to X_{red}$ 是 $X_{red}$ 的正规化。
\end{lemma}

\begin{proof}
令 $f : Y \to X$ 为态射（\ref{morphisms-equation-generic-points}）。由《概形》引理
\ref{schemes-lemma-map-into-reduction}，我们得到分解
$Y \to X_{red} \to X$，它分解 $f$。
由引理 \ref{morphisms-lemma-characterize-normalization}，得到典范态射 $X^\nu \to X_{red}$，
并且 $X^\nu$ 是 $X_{red}$ 在 $Y$ 中的正规化。
引理得证，因为 $Y \to X_{red}$ 与为 $X_{red}$ 构造的态射（\ref{morphisms-equation-generic-points}）相同。
\end{proof}

\noindent
若 $X$ 是约化的，则 $X$ 的正规化与 $\mathcal{O}_X$ 在亚纯函数层 $\mathcal{K}_X$ 中的整闭包的相对谱相同
（见《除子》第 \ref{divisors-section-meromorphic-functions} 节）。
也就是说，在此情形 $\mathcal{K}_X = f_*\mathcal{O}_Y$，见
《除子》引理 \ref{divisors-lemma-reduced-finite-irreducible} 及其证明。下面对此作明确描述。

\begin{lemma}
\label{morphisms-lemma-description-normalization}
设 $X$ 为约化概形，使每个拟紧开集都有有限多个不可约分支。令 $\Spec(A) = U \subset X$
为仿射开集。则
\begin{enumerate}
\item $A$ 有有限多个极小素理想
$\mathfrak q_1, \ldots, \mathfrak q_t$；
\item 总分式环 $Q(A)$（属于 $A$）是
$Q(A/\mathfrak q_1) \times \ldots \times Q(A/\mathfrak q_t)$；
\item 整闭包 $A'$（即 $A$ 在 $Q(A)$ 中的整闭包）是各整环 $A/\mathfrak q_i$
在域 $Q(A/\mathfrak q_i)$ 中的整闭包的乘积；并且
\item $\nu^{-1}(U)$ 可识别为 $A'$ 的谱，其中
$\nu : X^\nu \to X$ 是正规化态射。
\end{enumerate}
\end{lemma}

\begin{proof}
极小素理想对应于不可约分支（《代数》引理 \ref{algebra-lemma-irreducible}），
因此由假设得到 (1)。有
$(0) = \mathfrak q_1 \cap \ldots \cap \mathfrak q_t$，因为 $A$ 约化
（《代数》引理 \ref{algebra-lemma-Zariski-topology}）。
于是由《代数》引理 \ref{algebra-lemma-total-ring-fractions-no-embedded-points} 与
\ref{algebra-lemma-minimal-prime-reduced-ring}，有
$Q(A) = \prod A_{\mathfrak q_i} = \prod \kappa(\mathfrak q_i)$。
这证明了 (2)。(3) 由《代数》引理
\ref{algebra-lemma-characterize-reduced-ring-normal} 或引理 \ref{morphisms-lemma-normalization-in-disjoint-union} 得出。
(4) 成立是因为显然 $f^{-1}(U) \to U$ 是态射
$$
\Spec\left(\prod \kappa(\mathfrak q_i)\right)
\longrightarrow
\Spec(A)
$$
其中 $f : Y \to X$ 是态射（\ref{morphisms-equation-generic-points}）。
\end{proof}
\begin{lemma}
\label{morphisms-lemma-stalk-normalization}
设 $X$ 为一个概形，使每个拟紧开集都有有限多个不可约分支。设 $\nu : X^\nu \to X$
为 $X$ 的正规化。令 $x \in X$。则下列对象作为 $\mathcal{O}_{X, x}$-代数典范同构：
\begin{enumerate}
\item 茎 $(\nu_*\mathcal{O}_{X^\nu})_x$；
\item $\mathcal{O}_{X, x}$ 在 $(\mathcal{O}_{X, x})_{red}$ 的总分式环中的整闭包；
\item $\mathcal{O}_{X, x}$ 在由 $\mathcal{O}_{X, x}$ 的极小素理想的剩余域组成的乘积中的整闭包
（而这些极小素理想只有有限多个）。
\end{enumerate}
\end{lemma}

\begin{proof}
将 $X$ 替换为 $x$ 的一个仿射开邻域后，不妨设 $X$ 有有限多个不可约分支，并且 $x$ 包含在每一个分支中。
于是茎 $(\nu_*\mathcal{O}_{X^\nu})_x$ 是 $A = \mathcal{O}_{X, x}$ 在积 $L$ 中的整闭包，
其中该积由 $A$ 的极小素理想的剩余域组成。
这由正规化的构造和引理 \ref{morphisms-lemma-integral-closure} 得出。
或者，可以使用引理 \ref{morphisms-lemma-description-normalization} 以及正规化与局部化可交换这一事实
（《代数》引理 \ref{algebra-lemma-integral-closure-localize}）。
由于 $A_{red}$ 有有限多个极小素理想
（因为它们恰好对应于 $X$ 中经过 $x$ 的不可约分支的泛点），
由《代数》引理 \ref{algebra-lemma-total-ring-fractions-no-embedded-points} 可知，$L$ 是 $A_{red}$ 的总分式环。
因此，我们的环也是 $A$ 在 $A_{red}$ 的总分式环中的整闭包。
\end{proof}

\begin{lemma}
\label{morphisms-lemma-normalization-normal}
设 $X$ 为一个概形，使每个拟紧开集都有有限多个不可约分支。
\begin{enumerate}
\item 正规化 $X^\nu$ 是整正规概形的不交并。
\item 态射 $\nu : X^\nu \to X$ 是整的、满射的，并在不可约分支之间诱导双射。
\item 对任意整态射 $\alpha : X' \to X$，如果对于拟紧开集 $U \subset X$，逆像 $\alpha^{-1}(U)$
有有限多个不可约分支，并且
$\alpha|_{\alpha^{-1}(U)} : \alpha^{-1}(U) \to U$ 是双有理的\footnote{这种笨拙的表述是必要的，因为我们只在源和目标都有有限多个不可约分支时定义了态射为双有理的含义。若
$X'_{red} \to X_{red}$ 满足该条件，则已经足够。}，
则存在分解 $X^\nu \to X' \to X$，并且 $X^\nu \to X'$ 是 $X'$ 的正规化。
\item 对任意态射 $Z \to X$，其中 $Z$ 是正规概形，并且 $Z$ 的每个不可约分支都支配 $X$ 的一个不可约分支，
存在唯一分解 $Z \to X^\nu \to X$。
\end{enumerate}
\end{lemma}
\begin{proof}
令 $f : Y \to X$ 如（\ref{morphisms-equation-generic-points}）中所述。
$X^\nu$ 是整正规概形的不交并，因为 $Y$ 是正规的，且 $Y$ 的每个仿射开集都有有限多个不可约分支，见
引理 \ref{morphisms-lemma-normal-normalization}。这证明了 (1)。
或者，也可由引理 \ref{morphisms-lemma-normalization-reduced} 与
\ref{morphisms-lemma-description-normalization} 推出 (1)。

\medskip\noindent
由引理 \ref{morphisms-lemma-characterize-normalization}，态射 $\nu$ 是整的。
由引理 \ref{morphisms-lemma-normal-normalization}，态射 $Y \to X^\nu$ 在不可约分支之间诱导双射；
由 $Y$ 的构造，这意味着 $X^\nu \to X$ 在不可约分支之间诱导双射。
按构造 $f : Y \to X$ 是支配的，因此 $\nu$ 也是支配的。由于整态射是闭的
（引理 \ref{morphisms-lemma-integral-universally-closed}），这意味着 $\nu$ 是满射。这证明了 (2)。

\medskip\noindent
假设 $\alpha : X' \to X$ 如 (3) 中所述。显然 $X'$ 满足定义正规化所需的假设。
令 $f' : Y' \to X'$ 为从 $X'$ 出发按（\ref{morphisms-equation-generic-points}）构造的态射。
由于 $\alpha$ 局部双有理，显然 $Y' = Y$ 且 $f = \alpha \circ f'$。
因此分解 $X^\nu \to X' \to X$ 存在，并且由引理
\ref{morphisms-lemma-characterize-normalization}，$X^\nu \to X'$ 是 $X'$ 的正规化。这证明了 (3)。

\medskip\noindent
令 $g : Z \to X$ 为态射，其定义域是正规概形，并且每个不可约分支都支配 $X$ 的一个不可约分支。
由引理 \ref{morphisms-lemma-normalization-reduced}，有 $X^\nu = X_{red}^\nu$，并且由
《概形》引理 \ref{schemes-lemma-map-into-reduction}，$Z \to X$ 经过 $X_{red}$ 分解。
因此可将 $X$ 替换为 $X_{red}$，并假设 $X$ 约化。此外，由于分解唯一，只需在 $Z$ 上局部构造它。
令 $W \subset Z$ 与 $U \subset X$ 为满足 $g(W) \subset U$ 的仿射开集。
写 $U = \Spec(A)$、$W = \Spec(B)$，其中 $g|_W$ 由 $\varphi : A \to B$ 给出。
我们将自由使用引理 \ref{morphisms-lemma-description-normalization} 的结果。
令 $\mathfrak p_1, \ldots, \mathfrak p_t$ 为 $A$ 的极小素理想。
由于 $Z$ 正规，$B$ 是正规环，特别是约化的。此外，由假设，对任意极小素理想
$\mathfrak q \subset B$，$\varphi^{-1}(\mathfrak q)$ 是 $A$ 的极小素理想。
因此，若 $x \in A$ 是非零因子，即 $x \not \in \bigcup \mathfrak p_i$，则 $\varphi(x)$ 在 $B$ 中是非零因子。
于是得到典范环映射 $Q(A) \to Q(B)$。
由于 $B$ 是正规的，它等于自己在 $Q(B)$ 中的整闭包（见
《代数》引理 \ref{algebra-lemma-normal-ring-integrally-closed}）。
因此，整闭包 $A' \subset Q(A)$（即 $A$ 的整闭包）映入 $B$，通过典范映射 $Q(A) \to Q(B)$。
由于 $\nu^{-1}(U) = \Spec(A')$，这给出典范分解
$W \to \nu^{-1}(U) \to U$，它分解 $\nu|_W$。
唯一性的验证略去。
\end{proof}
\begin{lemma}
\label{morphisms-lemma-normalization-in-terms-of-components}
设 $X$ 为一个概形，使每个拟紧开集都有有限多个不可约分支。令 $Z_i \subset X$，$i \in I$，
为赋予约化诱导结构的 $X$ 的不可约分支。令 $Z_i^\nu \to Z_i$ 为正规化。
则 $\coprod_{i \in I} Z_i^\nu \to X$ 是 $X$ 的正规化。
\end{lemma}

\begin{proof}
由引理 \ref{morphisms-lemma-normalization-reduced}，不妨设 $X$ 约化。
于是引理可由引理 \ref{morphisms-lemma-description-normalization} 中的局部描述得到，
或者由引理 \ref{morphisms-lemma-normalization-normal} 的 (3) 得到，因为
$\coprod Z_i \to X$ 是整的且局部双有理的（$X$ 约化且具有局部有限多个不可约分支）。
\end{proof}

\begin{lemma}
\label{morphisms-lemma-normalization-birational}
设 $X$ 为具有有限多个不可约分支的约化概形。
则正规化态射 $X^\nu \to X$ 是双有理的。
\end{lemma}

\begin{proof}
由引理 \ref{morphisms-lemma-normalization-normal}，正规化在不可约分支之间诱导双射。
令 $\eta \in X$ 为 $X$ 的一个不可约分支的泛点，并令 $\eta^\nu \in X^\nu$
为 $X^\nu$ 中对应不可约分支的泛点。
则 $\eta^\nu \mapsto \eta$；为了完成证明，需要证明
$\mathcal{O}_{X, \eta} \to \mathcal{O}_{X^\nu, \eta^\nu}$
是同构，见定义 \ref{morphisms-definition-birational}。
由于 $X$ 与 $X^\nu$ 都约化，可见两个局部环都等于其剩余域
（《代数》引理 \ref{algebra-lemma-minimal-prime-reduced-ring}）。
另一方面，由正规化作为
$X$ 在 $Y = \coprod \Spec(\kappa(\eta))$ 中的正规化的构造，我们有
$\kappa(\eta) \subset \kappa(\eta^\nu) \subset \kappa(\eta)$，证明完成。
\end{proof}

\begin{lemma}
\label{morphisms-lemma-finite-birational-over-normal}
整概形之间的有限（甚至整）双有理态射 $f : X \to Y$，
若 $Y$ 为正规概形，则是同构。
\end{lemma}

\begin{proof}
令 $V \subset Y$ 为仿射开集，并令其逆像 $U \subset X$ 也是仿射开集。
由于 $f$ 是整概形之间的双有理态射，环同态
$\mathcal{O}_Y(V) \to \mathcal{O}_X(U)$ 是整环的单射，
并且诱导分式域的同构。由于 $Y$ 正规，环 $\mathcal{O}_Y(V)$
在其分式域中整闭。由于 $f$ 是有限的（或整的），
$\mathcal{O}_X(U)$ 的每个元素都在 $\mathcal{O}_Y(V)$ 上整。
因此 $\mathcal{O}_Y(V) = \mathcal{O}_X(U)$。这就证明了 $f$ 是同构，证毕。
\end{proof}

\begin{lemma}
\label{morphisms-lemma-normalization-trivial}
设 $X$ 为每个局部都只有有限多个不可约分支的概形。
正规化态射 $\nu : X^\nu \to X$ 是同构，当且仅当 $X$ 正规。
\end{lemma}

\begin{proof}
若 $\nu$ 是同构，则由引理 \ref{morphisms-lemma-normalization-normal}，$X$ 正规。
反之，设 $X$ 正规。由引理 \ref{morphisms-lemma-normalization-localization} 以及
《性质》引理 \ref{properties-lemma-normal-locally-finite-nr-irreducibles}，
不妨设 $X$ 是整概形。由引理 \ref{morphisms-lemma-normalization-normal}，
态射 $\nu$ 是整的，且 $X^\nu$ 只有一个不可约分支
（所以它是整概形）。由引理 \ref{morphisms-lemma-normalization-birational} 和
\ref{morphisms-lemma-finite-birational-over-normal} 得出结论。
\end{proof}

\begin{lemma}
\label{morphisms-lemma-Japanese-normalization}
设 $X$ 为整的 Japanese 概形。
正规化 $\nu : X^\nu \to X$ 是有限态射。
\end{lemma}

\begin{proof}
这由定义（《性质》定义 \ref{properties-definition-nagata}）和
引理 \ref{morphisms-lemma-description-normalization} 得出。具体而言，在这种情况下，
该引理说明 $\nu^{-1}(\Spec(A))$ 是 $A$ 在其分式域中的整闭包的谱。
\end{proof}

\begin{lemma}
\label{morphisms-lemma-nagata-normalization}
设 $X$ 为 Nagata 概形。
正规化 $\nu : X^\nu \to X$ 是有限态射。
\end{lemma}

\begin{proof}
注意 Nagata 概形是局部诺特的，因而定义 \ref{morphisms-definition-normalization}
确实适用。该引理现在是引理
\ref{morphisms-lemma-nagata-normalization-finite-general} 的特例，
但也可以直接如下证明。
将 $X^\nu \to X$ 写成复合
$X^\nu \to X_{red} \to X$。由于 $X_{red} \to X$ 是闭浸入，
它是有限的。因此（由引理 \ref{morphisms-lemma-composition-finite}）
只需对约化的 Nagata 概形证明该引理。
令 $\Spec(A) = U \subset X$ 为仿射开集。由引理
\ref{morphisms-lemma-description-normalization}，有
$\nu^{-1}(U) = \Spec(\prod A_i')$，其中 $A_i'$ 是
$A/\mathfrak q_i$ 在其分式域中的整闭包。由于 $A$ 是 Nagata 环
（见《性质》引理 \ref{properties-lemma-locally-nagata}），
每个环扩张 $A/\mathfrak q_i \subset A'_i$ 都是有限的。
因此 $A \to \prod A'_i$ 是有限环映射，结论成立。
\end{proof}

\begin{lemma}
\label{morphisms-lemma-normalization-geom-unibranch-univ-homeo}
设 $X$ 为不可约、几何单分支概形。
正规化态射 $\nu : X^\nu \to X$ 是泛同胚。
\end{lemma}

\begin{proof}
我们需要证明 $\nu$ 是整的、泛单射且满射，见引理
\ref{morphisms-lemma-universal-homeomorphism}。由引理
\ref{morphisms-lemma-normalization-normal}，态射 $\nu$ 是整的。
令 $x \in X$，并置 $A = \mathcal{O}_{X, x}$。
由于 $X$ 不可约，$A$ 只有一个极小素理想 $\mathfrak p$，且
$A_{red} = A/\mathfrak p$。由引理 \ref{morphisms-lemma-stalk-normalization}，
茎 $A' = (\nu_*\mathcal{O}_{X^\nu})_x$ 是 $A$ 在 $A_{red}$ 的分式域中的整闭包。
由《代数》定义 \ref{more-algebra-definition-unibranch}，可知 $A'$
只有一个素理想 $\mathfrak m'$，它位于 $\mathfrak m_x \subset A$ 之上，并且
$\kappa(\mathfrak m')/\kappa(x)$ 是纯不可分扩张。因此由引理
\ref{morphisms-lemma-universally-injective}，$\nu$ 是双射（因而满射）且泛单射。
\end{proof}





\section{Weak normalization}
\label{morphisms-section-weak-normalization}

\noindent
我们只在概形局部具有有限多个不可约分支时定义其弱正规化；这类似于正规化的情形。

\begin{lemma}
\label{morphisms-lemma-relative-weak-normalization-algebra}
设环映射 $A \to B$ 诱导谱的支配态射
$\Spec(B) \to \Spec(A)$。存在一个 $A$-子代数
$B' \subset B$，满足
\begin{enumerate}
\item $\Spec(B') \to \Spec(A)$ 是泛同胚；
\item 给定分解 $A \to C \to B$，其中
$\Spec(C) \to \Spec(A)$ 是泛同胚，则 $C \to B$ 的像包含于 $B'$ 中。
\end{enumerate}
\end{lemma}

\begin{proof}
以下反复使用引理 \ref{morphisms-lemma-reduction-universal-homeomorphism}，不再另行说明。
考虑交换图
$$
\xymatrix{
B \ar[r] & B_{red} \\
A \ar[u] \ar[r] & A_{red} \ar[u]
}
$$
对于 (2) 中任意一个分解 $A \to C \to B$（其所分解的映射为 $A \to B$），可见
$A_{red} \to C_{red} \to B_{red}$ 是
$A_{red} \to B_{red}$ 的满足 (2) 的分解。
因此，如果引理对 $A_{red} \to B_{red}$ 成立并产生
$A_{red}$-子代数 $B'_{red} \subset B_{red}$，那么令 $B' \subset B$
等于 $B'_{red}$ 的逆像，即可为 $A \to B$ 解决该引理。
于是问题归约到下一段所讨论的情形。

\medskip\noindent
设 $A$ 与 $B$ 都约化。在此情形，由《代数》引理
\ref{algebra-lemma-image-dense-generic-points}，有 $A \subset B$。
令 $A \to C \to B$ 为 (2) 中的一个分解。
于是可以将命题 \ref{morphisms-proposition-universal-homeomorphism}
应用于 $A \subset C$，从而每个 $C$ 的元素都包含于一个扩张
$A[c_1, \ldots, c_n] \subset C$ 中，并且对 $i = 1, \ldots, n$，有
\begin{enumerate}
\item $c_i^2, c_i^3 \in A[c_1, \ldots, c_{i - 1}]$，或者
\item 存在素数 $p$，使得
$pc_i, c_i^p \in A[c_1, \ldots, c_{i - 1}]$。
\end{enumerate}
因此，若将 $B' \subset B$ 定义为满足下述条件的元素 $b \in B$ 的集合，
则性质 (2) 成立：包含于一个扩张
$
A[b_1, \ldots, b_n] \subset B
$
中，且满足 (*)：对 $i = 1, \ldots, n$，有
\begin{enumerate}
\item $b_i^2, b_i^3 \in A[b_1, \ldots, b_{i - 1}]$，或者
\item 存在素数 $p$，使得
$pb_i, b_i^p \in A[b_1, \ldots, b_{i - 1}]$。
\end{enumerate}
只需检查两件事：(a) $B'$ 是 $A$-子代数；(b)
$\Spec(B') \to \Spec(A)$ 是泛同胚。
对 (a)，若 $n \geq 0$ 且 $b_1, \ldots, b_n \in B$ 满足 (*)，
并且 $m \geq 0$ 且 $b'_1, \ldots, b'_m \in B$ 满足 (*)，
则整数 $n + m$ 以及
$b_1, \ldots, b_n, b'_1, \ldots, b'_m \in B$ 也满足 (*)，所以 (a) 成立。
最后，(b) 由命题 \ref{morphisms-proposition-universal-homeomorphism}
以及 $B'$ 的构造得到。
\end{proof}

\begin{lemma}
\label{morphisms-lemma-relative-weak-normalization-localization}
设环映射 $A \to B$ 诱导谱的支配态射
$\Spec(B) \to \Spec(A)$。引理
\ref{morphisms-lemma-relative-weak-normalization-algebra} 中 $A$-子代数
$B' \subset B$ 的构造与局部化相容（解释见证明）。
\end{lemma}

\begin{proof}
令 $S \subset A$ 为乘法子集。于是 $S^{-1}A \to S^{-1}B$ 是一个
环映射，并且同样诱导支配态射
$\Spec(S^{-1}B) \to \Spec(S^{-1}A)$
（见引理 \ref{morphisms-lemma-base-change-dominant} 和
\ref{morphisms-lemma-generalizations-lift-flat}）。因此，引理
\ref{morphisms-lemma-relative-weak-normalization-algebra} 产生一个
$S^{-1}A$-子代数 $(S^{-1}B)' \subset S^{-1}B$。
所述结论是指 $S^{-1}B' = (S^{-1}B)'$，并且它们是
$S^{-1}A$-子代数，且属于 $S^{-1}B$。

\medskip\noindent
为说明这一点，我们使用引理
\ref{morphisms-lemma-relative-weak-normalization-algebra} 的证明中对 $B'$
和 $(S^{-1}B)'$ 的构造。第一步中可见，$B'$ 是
$A_{red}$-子代数 $B'_{red} \subset B_{red}$ 的逆像，
该子代数由环映射 $A_{red} \to B_{red}$ 构造；$(S^{-1}B)'$ 同理。
注意 $S^{-1}B_{red} = (S^{-1}B)_{red}$，问题归约到下一段所讨论的情形。

\medskip\noindent
若 $A$ 与 $B$ 约化，则我们把 $B'$ 构造成下列子代数的并：
$A[b_1, \ldots, b_n]$，其中对 $i = 1, \ldots, n$ 有
\begin{enumerate}
\item $b_i^2, b_i^3 \in A[b_1, \ldots, b_{i - 1}]$，或者
\item 存在素数 $p$，使得
$pb_i, b_i^p \in A[b_1, \ldots, b_{i - 1}]$。
\end{enumerate}
对 $(S^{-1}B)' \subset S^{-1}B$ 也同理。因此显然，
$B' \to B \to S^{-1}B$ 的像包含于 $(S^{-1}B)'$ 中。
要证明相应的映射 $S^{-1}B' \to (S^{-1}B)'$ 是满射，
可使用引理 \ref{morphisms-lemma-special-elements-and-localization} 逐次清除分母；
细节略去。
\end{proof}

\begin{lemma}
\label{morphisms-lemma-relative-seminormalization-algebra}
设环映射 $A \to B$ 诱导谱的支配态射
$\Spec(B) \to \Spec(A)$。存在一个 $A$-子代数
$B' \subset B$，满足
\begin{enumerate}
\item $\Spec(B') \to \Spec(A)$ 是诱导剩余域同构的泛同胚；
\item 给定分解 $A \to C \to B$，其中 $\Spec(C) \to \Spec(A)$
是诱导剩余域同构的泛同胚，则 $C \to B$ 的像包含于 $B'$ 中。
\end{enumerate}
\end{lemma}

\begin{proof}
该证明与引理 \ref{morphisms-lemma-relative-weak-normalization-algebra} 的证明完全相同，
只是使用命题 \ref{morphisms-proposition-universal-homeomorphism-equal-residue-fields}，
而不是命题 \ref{morphisms-proposition-universal-homeomorphism}。
\end{proof}

\begin{lemma}
\label{morphisms-lemma-relative-seminormalization-localization}
设环映射 $A \to B$ 诱导谱的支配态射
$\Spec(B) \to \Spec(A)$。引理
\ref{morphisms-lemma-relative-seminormalization-algebra} 中 $A$-子代数
$B' \subset B$ 的构造与局部化相容（解释见证明）。
\end{lemma}

\begin{proof}
证明与引理 \ref{morphisms-lemma-relative-weak-normalization-localization} 的证明相同。
\end{proof}

\begin{lemma}
\label{morphisms-lemma-relative-seminormalization}
设 $f : Y \to X$ 为概形的拟紧、拟分离且支配的态射。
\begin{enumerate}
\item 由分解 $Y \to X' \to X$ 构成的范畴，其中 $X' \to X$ 为泛同胚，具有始对象
$Y \to X^{Y/wn} \to X$。
\item 由分解 $Y \to X' \to X$ 构成的范畴，其中 $X' \to X$ 为诱导剩余域同构的泛同胚，具有始对象
$Y \to X^{Y/sn} \to X$。
\end{enumerate}
此外，分解 $Y \to X^{Y/wn} \to X$ 和
$Y \to X^{Y/sn} \to X$ 的构造与到 $X$ 的开子概形的基变换相容。
\end{lemma}

\begin{proof}
我们证明 (1)，略去 (2) 的证明；最后的断言也将由分解的构造推出。
以下使用引理 \ref{morphisms-lemma-universal-homeomorphism}，不再另行说明。
首先，令 $Y \to X^{Y/n} \to X$ 为 $X$ 在 $Y$ 中的正规化，见定义
\ref{morphisms-definition-normalization-X-in-Y}。对于 (1) 中的分解
$Y \to X' \to X$，由引理 \ref{morphisms-lemma-characterize-normalization}
可得到一个与给定态射相容的唯一态射 $X^{Y/n} \to X'$。
因此，只需证明将 $f$ 替换为 $X^{Y/n} \to X$ 后的引理。
这就归约到下一段研究的情形。

\medskip\noindent
假设 $f$ 是整的（证明的其余部分在 $f$ 为仿射时也更一般地成立）。
令 $U = \Spec(A)$ 为 $X$ 的仿射开集，并令
$V = f^{-1}(U) = \Spec(B)$ 为其在 $Y$ 中的逆像。
则 $A \to B$ 是一个在谱上诱导支配态射的环映射。
由引理 \ref{morphisms-lemma-relative-weak-normalization-algebra}，得到
$A$-子代数 $B' \subset B$，使得置
$U^{V/wn} = \Spec(B')$ 后，分解 $V \to U^{V/wn} \to U$
是由 $V \to U' \to U$（其中 $U' \to U$ 为泛同胚）构成的范畴中的始对象。

\medskip\noindent
若 $U_1 \subset U_2 \subset X$ 是仿射开集，令
$V_i = f^{-1}(U_i)$，由此得到一个典范态射
$$
\rho_{U_1}^{U_2} : U_1^{V_1/wn} \to U_1 \times_{U_2} U_2^{V_2/wn}
$$
它在 $U_1$ 上，并由 $U_1^{V_1/wn}$ 的泛性质给出。
这些态射满足自然的函子性；其表述和证明留给读者。
此外，态射 $\rho_{U_1}^{U_2}$ 是同构；当
$U_1 \subset U_2$ 是标准开集时，这由引理
\ref{morphisms-lemma-relative-weak-normalization-localization} 得出，
一般情形则可由这些映射的函子性以及《概形》引理
\ref{schemes-lemma-standard-open-two-affines} 归约到标准开集情形（细节略去）。
因此，由相对粘合（《构造》引理
\ref{constructions-lemma-relative-glueing}）得到一个态射
$X^{Y/wn} \to X$，它限制为 $U^{V/wn} \to U$，并且在 $U$ 上与
$\rho_{U_1}^{U_2}$ 相容。
当然，态射 $V \to U^{V/wn}$ 粘合为态射
$Y \to X^{Y/wn}$（见《构造》注
\ref{constructions-remark-relative-glueing-functorial}），
从而得到分解 $Y \to X^{Y/wn} \to X$，其中第二个态射是泛同胚。

\medskip\noindent
最后，令 $Y \to X' \to X$ 为 (1) 中的分解。
如上，考虑 $V \to U^{V/wn} \to U \subset X$，
我们得到分解 $V \to U \times_X X' \to U$，其中第二个箭头是泛同胚，
并得到一个唯一态射
$g_U : U^{V/wn} \to U \times_X X'$，它在 $U$ 上由
$U^{V/wn}$ 的泛性质给出。
这些 $g_U$ 与态射 $\rho_{U_1}^{U_2}$ 相容；细节略去。
因此存在唯一态射 $g : X^{Y/wn} \to X'$，它在 $X$ 上与
$g_U$ 在 $U$ 上一致，见《构造》注
\ref{constructions-remark-relative-glueing-functorial}。
这证明 $Y \to X^{Y/wn} \to X$ 是我们范畴中的始对象，证明完成。
\end{proof}

\begin{definition}
\label{morphisms-definition-relative-seminormalization}
设 $f : Y \to X$ 为概形的拟紧、拟分离且支配的态射。
\begin{enumerate}
\item 引理 \ref{morphisms-lemma-relative-seminormalization} 的 (2) 中构造的分解
$Y \to X^{Y/sn} \to X$ 称为 $X$ 在 $Y$ 中的\textit{半正规化}。
\item 引理 \ref{morphisms-lemma-relative-seminormalization} 的 (1) 中构造的分解
$Y \to X^{Y/wn} \to X$ 称为 $X$ 在 $Y$ 中的\textit{弱正规化}。
\end{enumerate}
\end{definition}
\noindent
以下给出对局部具有有限多个不可约分支的概形的半正规化的一种重新解释。

\begin{lemma}
\label{morphisms-lemma-seminormal-and-normalization}
设 $X$ 为每个拟紧开集都有有限多个不可约分支的概形。令
$\nu : X^\nu \to X$ 为 $X$ 的正规化。则 $X$ 在 $X^\nu$ 中的半正规化
就是 $X$ 的半正规化。公式为：
$X^{sn} = X^{X^\nu/sn}$。
\end{lemma}

\begin{proof}
令 $f : Y \to X$ 如（\ref{morphisms-equation-generic-points}）中所述，
使得 $X^\nu$ 是 $X$ 在 $Y$ 中的正规化。
半正规化 $X^{sn} \to X$ 是概形 $X$ 的由诱导剩余域同构的泛同胚
$X' \to X$ 构成的范畴中的始对象。由于 $Y$ 是
$X$ 的不可约分支的泛点处的剩余域的谱的不交并，
可见对于该范畴中的任意 $X' \to X$，我们都得到一个典范提升
$f' : Y \to X'$，它是 $f$ 的提升。于是由引理
\ref{morphisms-lemma-characterize-normalization}，得到一个典范态射 $X^\nu \to X'$。
从而再得到一个典范态射 $X^{X^\nu/sn} \to X'$，
它由 $X^{X^\nu/sn}$ 的泛性质给出。
这样，我们看到 $X^{X^\nu/sn}$ 满足与 $X^{sn}$ 相同的泛性质，结论成立。
\end{proof}

\noindent
引理 \ref{morphisms-lemma-seminormal-and-normalization} 导出以下定义。
由于我们只在 $X$ 局部具有有限多个不可约分支时构造正规化，
对于弱正规化也限制在这种情形。

\begin{definition}
\label{morphisms-definition-weak-normalization}
设 $X$ 为每个拟紧开集都有有限多个不可约分支的概形。我们将 $X$ 的
\textit{弱正规化}定义为其弱正规化
$$
X^\nu \longrightarrow X^{wn} \longrightarrow X
$$
即 $X$ 在正规化 $X^\nu$（即 $X$ 的正规化）中的弱正规化（定义 \ref{morphisms-definition-normalization}）。公式为：
$X^{wn} = X^{X^\nu/wn}$。
\end{definition}

\noindent
结合引理 \ref{morphisms-lemma-seminormal-and-normalization}，可见对于一个
局部具有有限多个不可约分支的概形 $X$，存在典范态射
$$
X^\nu \to X^{wn} \to X^{sn} \to X
$$
有了这个定义，我们就可以说明概形何时是弱正规的
（前提是它局部具有有限多个不可约分支）。

\begin{definition}
\label{morphisms-definition-weakly-normal}
设 $X$ 为每个拟紧开集都有有限多个不可约分支的概形。
称 $X$ 为\textit{弱正规}，如果弱正规化 $X^{wn} \to X$ 是同构（定义
\ref{morphisms-definition-weak-normalization}）。
\end{definition}

\noindent
由定义立即得到，对于每个拟紧开集都有有限多个不可约分支的概形 $X$，有
$$
X\text{ normal} \Rightarrow
X\text{ weakly normal} \Rightarrow
X\text{ seminormal}
$$
下面在仿射情形中说明弱正规的含义。

\begin{lemma}
\label{morphisms-lemma-affine-weakly-normal}
设 $X = \Spec(A)$ 为具有有限多个不可约分支的仿射概形。
则 $X$ 弱正规，当且仅当
\begin{enumerate}
\item $A$ 是半正规环（定义 \ref{morphisms-definition-seminormal-ring}）；
\item 对于素数 $p$ 以及满足下列条件的 $z, w \in A$：
(a) $z$ 是非零因子；(b) $w^p$ 可被 $z^p$ 整除；(c) $pw$ 可被 $z$ 整除，
都有 $w$ 可被 $z$ 整除。
\end{enumerate}
\end{lemma}

\begin{proof}
假设 $X$ 弱正规。由于弱正规概形是半正规概形，
由弱正规概形的定义可知 (1) 成立。
特别地，$A$ 是约化的。令 $p, z, w$ 如 (2) 中所述。
取 $x, y \in A$，使得 $z^p x = w^p$ 且 $zy = pw$。
于是 $p^p x = y^p$。环映射
$A \to C = A[t]/(t^p - x, pt - y)$ 在谱上诱导泛同胚。
$X^\nu$ 是 $X$ 的整闭包 $A'$ 的谱，其中它取自 $A$ 在 $A$ 的总分式环中的整闭包，见引理
\ref{morphisms-lemma-description-normalization}。
注意 $a = w/z \in A'$，因为 $a^p = x$。
因此有一个 $A$-代数同态 $A \to C \to A'$，将 $t$ 映为 $a$。
此时，弱正规的定义性质
$X = X^{wn} = X^{X^\nu/wn}$ 告诉我们，$C \to A'$ 的像包含于 $A$ 中。
于是得到所需的 $a \in A$。

\medskip\noindent
反之，假设 (1) 和 (2) 成立。令 $A'$ 如上一段中所述。
我们需要证明 $X^{X^\nu/wn} = X$。根据引理
\ref{morphisms-lemma-relative-weak-normalization-algebra} 的证明中的构造，
$X^{X^\nu/wn}$ 是 $A'$ 的子环的谱，该子环是下列子环
$A[a_1, \ldots, a_n] \subset A'$ 的并，其中对 $i = 1, \ldots, n$ 有
\begin{enumerate}
\item[(a)] $a_i^2, a_i^3 \in A[a_1, \ldots, a_{i - 1}]$，或者
\item[(b)] 存在素数 $p$，使得
$pa_i, a_i^p \in A[a_1, \ldots, a_{i - 1}]$。
\end{enumerate}
于是可以使用 (1) 和 (2) 归纳地看到
$a_1, \ldots, a_n \in A$；细节略去。
因此 $X = X^{X^\nu/wn}$，从而 $X$ 弱正规。
\end{proof}

\noindent
下面是一个必需的引理。

\begin{lemma}
\label{morphisms-lemma-locally-weakly-normal}
设 $X$ 为每个拟紧开集都有有限多个不可约分支的概形。以下条件等价：
\begin{enumerate}
\item 概形 $X$ 弱正规。
\item 对每个仿射开集 $U \subset X$，环 $\mathcal{O}_X(U)$
满足引理 \ref{morphisms-lemma-affine-weakly-normal} 的条件 (1) 和 (2)。
\item 存在仿射开覆盖 $X = \bigcup U_i$，使得每个环 $\mathcal{O}_X(U_i)$
满足引理 \ref{morphisms-lemma-affine-weakly-normal} 的条件 (1) 和 (2)。
\item 存在开覆盖 $X = \bigcup X_j$，使得每个开子概形 $X_j$ 都弱正规。
\end{enumerate}
此外，若 $X$ 弱正规，则每个开子概形也弱正规。
\end{lemma}

\begin{proof}
$X$ 弱正规的条件是态射
$X^{wn} = X^{X^\nu/wn} \to X$ 为同构。
由于 $X^\nu \to X$ 的构造与到开子概形的基变换相容，
并且 $X^{X^\nu/wn}$ 的构造与到 $X$ 的开子概形的基变换相容
（引理 \ref{morphisms-lemma-relative-seminormalization}），引理显然成立。
\end{proof}

\section{Zariski 主定理（代数版本）}
\label{morphisms-section-Zariski}

\noindent
这是可以使用纯代数方法证明的版本。
在我们能够证明更强的版本（对于非仿射态射）之前，
还需要发展更多工具。见《概形的上同调》一节
\ref{coherent-section-applications-formal-functions}
以及《更多关于态射》一节
\ref{more-morphisms-section-application-etale-neighbourhoods}。

\begin{theorem}[Zariski 主定理的代数版本]
\label{morphisms-theorem-main-theorem}
设 $f : Y \to X$ 为概形的仿射态射。
假设 $f$ 是有限型的。
令 $X'$ 为 $X$ 在 $Y$ 中的正规化。图示为：
$$
\xymatrix{
Y \ar[rd]_f \ar[rr]_{f'} & & X' \ar[ld]^\nu \\
& X &
}
$$
则存在开子概形 $U' \subset X'$，使得
\begin{enumerate}
\item $(f')^{-1}(U') \to U'$ 是同构，并且
\item $(f')^{-1}(U') \subset Y$ 正是 $f$ 为拟有限的点集。
\end{enumerate}
\end{theorem}

\begin{proof}
可以立即归约到 $X$、从而 $Y$ 都是仿射的情形。
设 $X = \Spec(R)$ 且 $Y = \Spec(A)$。
于是 $X' = \Spec(A')$，其中 $A'$ 是 $R$ 在 $A$ 中的整闭包，见定义
\ref{morphisms-definition-integral-closure} 和 \ref{morphisms-definition-normalization-X-in-Y}。
由《代数》定理 \ref{algebra-theorem-main-theorem}，
对于每个拟有限点 $y \in Y$（即 $f$ 在该点拟有限），存在开集
$U'_y \subset X'$，使得 $(f')^{-1}(U'_y) \to U'_y$ 是同构。
令 $U' = \bigcup U'_y$，其中 $y \in Y$ 遍历 $f$ 为拟有限的所有点。
还需证明 $f$ 在 $(f')^{-1}(U')$ 的所有点处都是拟有限的。
若 $y \in (f')^{-1}(U')$ 的像为 $x \in X$，则可见
$Y_x \to X'_x$ 在 $y$ 的某个邻域中是同构。因此 $Y_x$ 中没有点专化到 $y$，
因为 $f'(y)$ 在 $X'_x$ 中具有同样性质，见引理
\ref{morphisms-lemma-integral-fibres}。由引理
\ref{morphisms-lemma-quasi-finite-at-point-characterize} 的 (3)，
这意味着 $f$ 在 $y$ 处拟有限。
\end{proof}

\noindent
我们可以使用 Zariski 主定理的代数版本证明，态射拟有限的点集是开集。

\begin{lemma}
\label{morphisms-lemma-quasi-finite-points-open}
\begin{slogan}
态射的局部拟有限轨迹是开集
\end{slogan}
设 $f : X \to S$ 为概形的态射。
$X$ 中 $f$ 拟有限的点集是一个开集 $U \subset X$。
诱导的态射 $U \to S$ 局部拟有限。
\end{lemma}

\begin{proof}
假设 $f$ 在 $x$ 处拟有限。
令 $x \in U = \Spec(A) \subset X$、$V = \Spec(R) \subset S$
为定义 \ref{morphisms-definition-quasi-finite} 中的仿射开集。
由上面的定理 \ref{morphisms-theorem-main-theorem} 或《代数》引理
\ref{algebra-lemma-quasi-finite-open}，素理想 $\mathfrak q$ 中满足
$R \to A$ 拟有限者的集合是 $\Spec(A)$ 中的开集。
由于这些素理想都对应于 $X$ 中 $f$ 拟有限的点，得到第一个断言。
第二个断言显然。
\end{proof}

\noindent
稍后我们将把以下引理推广到一般的拟有限分离态射，见《更多关于态射》引理
\ref{more-morphisms-lemma-quasi-finite-separated-pass-through-finite}。

\begin{lemma}
\label{morphisms-lemma-quasi-finite-affine}
设 $f : Y \to X$ 为概形的态射。
假设
\begin{enumerate}
\item $X$ 与 $Y$ 都是仿射的，并且
\item $f$ 是拟有限的。
\end{enumerate}
则存在图
$$
\xymatrix{
Y \ar[rd]_f \ar[rr]_j & & Z \ar[ld]^\pi \\
& X &
}
$$
其中 $Z$ 是仿射的，$\pi$ 有限，且 $j$ 是开浸入。
\end{lemma}

\begin{proof}
这就是《代数》引理
\ref{algebra-lemma-quasi-finite-open-integral-closure} 在概形语言中的改写。
\end{proof}

\begin{lemma}
\label{morphisms-lemma-image-nowhere-dense-quasi-finite}
设 $f : Y \to X$ 为概形的拟有限态射。
设 $T \subset Y$ 为 $Y$ 的闭无处稠密子集。
则 $f(T) \subset X$ 是 $X$ 的无处稠密子集。
\end{lemma}

\begin{proof}
与引理 \ref{morphisms-lemma-image-nowhere-dense-finite} 的证明一样，
这立即归约到基概形 $X$ 为仿射的情形。
在这种情形下，$Y = \bigcup_{i = 1, \ldots, n} Y_i$ 是仿射开集的有限并
（因为 $f$ 拟紧）。由于每个 $T \cap Y_i$ 都无处稠密，
而且无处稠密集的有限并仍无处稠密（见《拓扑》引理
\ref{topology-lemma-nowhere-dense}），只需证明像
$f(T \cap Y_i)$ 在 $X$ 中无处稠密。
这又把我们归约到 $X$ 与 $Y$ 都为仿射的情形。
此时应用上面的引理 \ref{morphisms-lemma-quasi-finite-affine}，得到图
$$
\xymatrix{
Y \ar[rd]_f \ar[rr]_j & & Z \ar[ld]^\pi \\
& X &
}
$$
其中 $Z$ 仿射，$\pi$ 有限，且 $j$ 是开浸入。
置 $\overline{T} = \overline{j(T)} \subset Z$。
由《拓扑》引理 \ref{topology-lemma-image-nowhere-dense-open}，
$\overline{T}$ 在 $Z$ 中无处稠密。
由于 $f(T) \subset \pi(\overline{T})$，引理由有限情形中的相应结果得出，
见引理 \ref{morphisms-lemma-image-nowhere-dense-finite}。
\end{proof}

\section{泛一致有界的纤维}
\label{morphisms-section-universally-bounded}

\noindent
设 $X$ 为域 $k$ 上的概形。若 $X$ 在 $k$ 上有限，
则 $X = \Spec(A)$，其中 $A$ 是有限 $k$-代数。
另一种说法是 $X$ 在 $\Spec(k)$ 上有限局部自由，见定义
\ref{morphisms-definition-finite-locally-free}。因此
$X \to \Spec(k)$ 有一个\textit{次数}，它是整数 $d \geq 0$，
即 $d = \dim_k(A)$。我们有时称它为（有限）概形 $X$ 在 $k$ 上的\textit{次数}。

\begin{definition}
\label{morphisms-definition-universally-bounded}
设 $f : X \to Y$ 为概形的态射。
\begin{enumerate}
\item 称整数 $n$\textit{界住了 $f$ 的纤维的次数}，如果对于所有
$y \in Y$，纤维 $X_y$ 是 $\kappa(y)$ 上的有限概形，
且它在 $\kappa(y)$ 上的次数为 $\leq n$。
\item 称\textit{$f$ 的纤维泛一致有界}\footnote{这很可能不是标准术语。}，
如果存在一个整数 $n$ 界住 $f$ 的纤维的次数。
\end{enumerate}
\end{definition}

\noindent
特别地，注意每个纤维中的点数也被 $n$ 界住。
（反之不成立，即使所有纤维都是有限约化概形。）

\begin{lemma}
\label{morphisms-lemma-characterize-universally-bounded}
设 $f : X \to Y$ 为概形的态射，并设 $n \geq 0$。以下条件等价：
\begin{enumerate}
\item 整数 $n$ 界住了 $f$ 的纤维的次数；
\item 对每个态射 $\Spec(k) \to Y$，其中 $k$ 是域，
纤维积 $X_k = \Spec(k) \times_Y X$ 在 $k$ 上有限，且次数为 $\leq n$。
\end{enumerate}
在这种情况下，$f$ 的纤维泛一致有界，并且概形 $X_k$ 至多有 $n$ 个点。
更精确地，若 $X_k = \{x_1, \ldots, x_t\}$，则有
$$
n \geq \sum\nolimits_{i = 1, \ldots, t} [\kappa(x_i) : k]
$$
\end{lemma}

\begin{proof}
(2) $\Rightarrow$ (1) 显然。反向蕴含成立，因为若
$\Spec(k) \to Y$ 的像为 $y$，则
$X_k = \Spec(k) \times_{\Spec(\kappa(y))} X_y$。
按照定义，$f$ 的纤维泛一致有界是指存在某个 $n$。
最后，设 $X_k = \Spec(A)$。则 $\dim_k A = n$。
于是 $A$ 是 Artin 环，所有素理想都是极大理想 $\mathfrak m_i$，
而且 $A$ 是它在这些极大理想处的局部化之积。
见《代数》引理 \ref{algebra-lemma-finite-dimensional-algebra}
和 \ref{algebra-lemma-artinian-finite-length}。
于是 $\mathfrak m_i$ 对应于 $x_i$，有
$A_{\mathfrak m_i} = \mathcal{O}_{X_k, x_i}$，
从而存在满射
$A \to  \bigoplus \kappa(\mathfrak m_i) = \bigoplus \kappa(x_i)$。
线性代数于是给出引理断言中的不等式。
\end{proof}

\begin{lemma}
\label{morphisms-lemma-finite-locally-free-universally-bounded}
若 $f$ 是次数为 $d$ 的有限局部自由态射，
则 $d$ 界住 $f$ 的纤维的次数。
\end{lemma}

\begin{proof}
这是因为 $f$ 的任意基变换仍为次数 $d$ 的有限局部自由态射
（引理 \ref{morphisms-lemma-base-change-finite-locally-free}），
因而 $f$ 的所有纤维的次数都是 $d$。
\end{proof}

\begin{lemma}
\label{morphisms-lemma-composition-universally-bounded}
纤维泛一致有界的态射的复合，其纤维仍泛一致有界。
更精确地，假设 $n$ 界住 $f : X \to Y$ 的纤维的次数，
且 $m$ 界住 $g : Y \to Z$ 的纤维的次数。
则 $nm$ 界住 $g \circ f : X \to Z$ 的纤维的次数。
\end{lemma}

\begin{proof}
设 $f : X \to Y$ 和 $g : Y \to Z$ 的纤维泛一致有界。
假设 $\deg(X_y/\kappa(y)) \leq n$ 对所有 $y \in Y$ 成立，
并且 $\deg(Y_z/\kappa(z)) \leq m$ 对所有 $z \in Z$ 成立。
令 $z \in Z$ 为一点。由假设，概形 $Y_z$ 在
$\Spec(\kappa(z))$ 上有限。特别地，$Y_z$ 的底拓扑空间是有限离散集。
态射 $f_z : X_z \to Y_z$ 的纤维是 $f$ 在 $Y$ 中相应点处的纤维；
由上面的论证，它们是有限离散集。
因此 $X_z$ 的底拓扑空间也是有限离散集。
于是 $X_z$ 是仿射概形（这是一个很好的练习；例如也可将《性质》引理
\ref{properties-lemma-maximal-points-affine} 应用于 $X_z$ 的所有点而得出）。
写 $X_z = \Spec(A)$、$Y_z = \Spec(B)$，并置 $k = \kappa(z)$。
于是有 $k \to B \to A$，并且我们知道：(a) $\dim_k(B) \leq m$；
(b) 对每个极大理想 $\mathfrak m \subset B$，有
$\dim_{\kappa(\mathfrak m)}(A/\mathfrak mA) \leq n$。
我们断言这意味着 $\dim_k(A) \leq nm$。
注意 $B$ 是它在各个极大理想 $B_{\mathfrak m}$ 处的局部化之积；
例如，这是因为 $Y_z$ 是 $1$-点概形的不交并，或者由《代数》引理
\ref{algebra-lemma-finite-dimensional-algebra} 和
\ref{algebra-lemma-artinian-finite-length} 得出。
因此
$\dim_k(B) = \sum_{\mathfrak m} \dim_k(B_{\mathfrak m})$，并且
$\dim_k(A) = \sum_{\mathfrak m} \dim_k(A_{\mathfrak m})$，
其中两式中的 $\mathfrak m$ 都遍历 $B$（而不是 $A$）的极大理想。
由上述结论和 Nakayama 引理（《代数》引理
\ref{algebra-lemma-NAK}），可见每个 $A_{\mathfrak m}$ 都是
$B_{\mathfrak m}^{\oplus n}$ 的商模（作为 $B_{\mathfrak m}$-模）。因此
$\dim_k(A_{\mathfrak m}) \leq n \dim_k(B_{\mathfrak m})$。
综合起来可得
$$
\dim_k(A) = \sum\nolimits_{\mathfrak m} \dim_ta	(A_{\mathfrak m})
\leq \sum\nolimits_{\mathfrak m} n \dim_k(B_{\mathfrak m})
= n \dim_k(B) \leq nm
$$
正是所需结论。
\end{proof}

\begin{lemma}
\label{morphisms-lemma-base-change-universally-bounded}
纤维泛一致有界的态射的基变换，其纤维仍泛一致有界。
更精确地，若 $n$ 界住 $f : X \to Y$ 的纤维的次数，且
$Y' \to Y$ 是任意态射，则基变换
$f' : Y' \times_Y X \to Y'$ 的纤维的次数也被 $n$ 界住。
\end{lemma}

\begin{proof}
这由引理 \ref{morphisms-lemma-characterize-universally-bounded} 的结论显然成立。
\end{proof}

\begin{lemma}
\label{morphisms-lemma-descent-universally-bounded}
设 $f : X \to Y$ 为概形的态射。设 $Y' \to Y$ 为概形的态射，并令
$f' : X' = X_{Y'} \to Y'$ 为 $f$ 的基变换。
若 $Y' \to Y$ 满射且 $f'$ 的纤维泛一致有界，
则 $f$ 的纤维也泛一致有界。更精确地，若 $n$ 界住 $f'$ 的纤维的次数，
则 $n$ 也界住 $f$ 的纤维的次数。
\end{lemma}

\begin{proof}
令 $n \geq 0$ 为界住 $f'$ 的纤维次数的整数。
我们断言 $n$ 对 $f$ 也适用。事实上，若 $y \in Y$ 为一点，
则取一点 $y' \in Y'$ 位于 $y$ 之上，并注意
$$
X'_{y'} = \Spec(\kappa(y')) \times_{\Spec(\kappa(y))} X_y.
$$
由于假设 $X'_{y'}$ 是次数为 $\leq n$ 的有限概形（在 $\kappa(y')$ 上），
所以 $X_y$ 也是次数为 $\leq n$ 的有限概形（在 $\kappa(y)$ 上）。
（略去一些细节。）
\end{proof}

\begin{lemma}
\label{morphisms-lemma-immersion-universally-bounded}
浸入的纤维泛一致有界。
\end{lemma}

\begin{proof}
定义中可取整数 $n = 1$。
\end{proof}

\begin{lemma}
\label{morphisms-lemma-etale-universally-bounded}
设 $f : X \to Y$ 为概形的 'etale 态射。设 $n \geq 0$。
以下条件等价：
\begin{enumerate}
\item 整数 $n$ 界住纤维的次数；
\item 对每个域 $k$ 和态射 $\Spec(k) \to Y$，基变换
$X_k = \Spec(k) \times_Y X$ 至多有 $n$ 个点；
\item 对每个 $y \in Y$ 和每个可分代数闭包
$\kappa(y) \subset \kappa(y)^{sep}$，概形
$X_{\kappa(y)^{sep}}$ 至多有 $n$ 个点。
\end{enumerate}
\end{lemma}

\begin{proof}
这由引理 \ref{morphisms-lemma-characterize-universally-bounded}
以及下述事实得出：纤维 $X_y$ 是 $\kappa(y)$ 的有限可分域扩张的谱的不交并，
见引理 \ref{morphisms-lemma-etale-over-field}。
\end{proof}

\noindent
纤维泛一致有界是绝对概念，而不是相对概念。
这就是下一个引理的条件要求 $X$ 拟紧，而不是要求 $f$ 拟紧的原因。

\begin{lemma}
\label{morphisms-lemma-locally-quasi-finite-qc-source-universally-bounded}
设 $f : X \to Y$ 为概形的态射。假设
\begin{enumerate}
\item $f$ 局部拟有限，并且
\item $X$ 拟紧。
\end{enumerate}
则 $f$ 的纤维泛一致有界。
\end{lemma}

\begin{proof}
由于 $X$ 拟紧，存在有限仿射开覆盖
$X = \bigcup_{i = 1, \ldots, n} U_i$ 以及仿射开集
$V_i \subset Y$，$i = 1, \ldots, n$，使得 $f(U_i) \subset V_i$。
由“局部拟有限”的局部性质（见引理
\ref{morphisms-lemma-quasi-finite-at-point-characterize} 的 (4)），可见态射
$f|_{U_i} : U_i \to V_i$ 是仿射概形之间的局部拟有限态射，
因而是拟有限的，见引理 \ref{morphisms-lemma-quasi-finite-locally-quasi-compact}。
对于 $y \in Y$，显然
$X_y = \bigcup_{y \in V_i} (U_i)_y$ 是开覆盖。
因此只需对仿射概形之间的拟有限态射证明该引理
（也就是说，若 $n_i$ 对态射 $f|_{U_i} : U_i \to V_i$ 适用，
则 $\sum n_i$ 对 $f$ 适用）。

\medskip\noindent
假设 $f : X \to Y$ 是仿射概形之间的拟有限态射。
由引理 \ref{morphisms-lemma-quasi-finite-affine}，可以找到图
$$
\xymatrix{
X \ar[rd]_f \ar[rr]_j & & Z \ar[ld]^\pi \\
& Y &
}
$$
其中 $Z$ 仿射，$\pi$ 有限，且 $j$ 是开浸入。
由于 $j$ 的纤维泛一致有界
（引理 \ref{morphisms-lemma-immersion-universally-bounded}），
问题归约为证明 $\pi$ 的纤维泛一致有界
（引理 \ref{morphisms-lemma-composition-universally-bounded}）。

\medskip\noindent
这又把我们归约到形如 $\Spec(B) \to \Spec(A)$ 的态射，
其中 $A \to B$ 有限。设 $B$ 由 $x_1, \ldots, x_n$
作为 $A$-模生成，于是
$$
A^{\oplus n} \longrightarrow B, \quad
(a_1, \ldots, a_n) \longmapsto \sum a_i x_i
$$
是满的 $A$-模映射。对任意素理想 $\mathfrak p \subset A$，
它诱导满的 $\kappa(\mathfrak p)$-向量空间映射
$$
\kappa(\mathfrak p)^{\oplus n} \longrightarrow B \otimes_A \kappa(\mathfrak p)
$$
换言之，定义纤维泛一致有界的态射时，可以取整数 $n$。
\end{proof}

\begin{lemma}
\label{morphisms-lemma-universally-bounded-permanence}
考虑概形态射的交换图
$$
\xymatrix{
X \ar[rd]_g \ar[rr]_f & & Y \ar[ld]^h \\
& Z &
}
$$
若 $g$ 的纤维泛一致有界，并且 $f$ 满且平坦，
则 $h$ 的纤维也泛一致有界。更精确地，若 $n$
界住 $g$ 的纤维的次数，则 $n$ 也界住 $h$ 的纤维的次数。
\end{lemma}

\begin{proof}
假设 $g$ 的纤维泛一致有界，并且 $f$ 满且平坦。
设 $g$ 的纤维的次数被 $n \in \mathbf{N}$ 界住。
我们断言 $n$ 对 $h$ 也适用。
令 $z \in Z$。考虑概形态射 $X_z \to Y_z$。
它平坦且满射。由假设，$X_z$ 是 $\kappa(z)$ 上的有限概形，
特别地，它是 Artin 环的谱（由《代数》引理
\ref{algebra-lemma-finite-dimensional-algebra}）。
由引理 \ref{morphisms-lemma-Artinian-affine}，态射 $X_z \to Y_z$ 是仿射的，
特别是拟紧的。由引理 \ref{morphisms-lemma-fpqc-quotient-topology}，
可知 $Y_z$ 也是有限离散的，因为 $X_z$ 具有这一性质。
因此 $Y_z$ 是仿射概形（这是一个很好的练习；例如也可将《性质》引理
\ref{properties-lemma-maximal-points-affine} 应用于 $Y_z$ 的所有点而得出）。
写 $Y_z = \Spec(B)$ 且 $X_z = \Spec(A)$。
则 $A$ 在 $B$ 上忠实平坦，因此 $B \subset A$。
所以 $\dim_k(B) \leq \dim_k(A) \leq n$，正是所需结论。
\end{proof}

\section{杂项}
\label{morphisms-section-misc}

\noindent
这里收录不适合放在其他地方的结果。

\begin{lemma}
\label{morphisms-lemma-factor-through-point}
设 $f : Y \to X$ 为概形的态射，并设 $x \in X$ 为一点。
假设 $Y$ 约化，并且在集合论意义下 $f(Y)$ 包含于 $\{x\}$ 中。
则 $f$ 经过典范态射 $x = \Spec(\kappa(x)) \to X$ 分解。
\end{lemma}

\begin{proof}
略。提示：在仿射局部工作，可将问题归约到一个交换代数引理。
给定环映射 $A \to B$，其中 $B$ 约化，并且存在唯一素理想
$\mathfrak p \subset A$ 位于 $\Spec(B) \to \Spec(A)$ 的像中，
则 $A \to B$ 经过 $\kappa(\mathfrak p)$ 分解。
这是一个很好的练习。
\end{proof}

\begin{lemma}
\label{morphisms-lemma-factor-through-set-of-points}
设 $f : Y \to X$ 为概形的态射，并设 $E \subset X$。
假设 $X$ 局部诺特，$E$ 的元素之间不存在非平凡专化，
$Y$ 约化，并且 $f(Y) \subset E$。
则 $f$ 经过 $\coprod_{x \in E} x \to X$ 分解。
\end{lemma}

\begin{proof}
当 $E$ 为单点集时，这由引理 \ref{morphisms-lemma-factor-through-point} 得出。
若 $E$ 有限，则由我们关于专化的假设，$E$（赋予从 $X$ 诱导的拓扑）
是有限离散空间。因此该情形归约到单点集情形。
一般地，可归约到 $X$ 与 $Y$ 都为仿射概形的情形。
设 $f : Y \to X$ 对应于环映射 $\varphi : A \to B$。
用 $A' \subset B$ 表示 $\varphi$ 的像。令
$E' \subset \Spec(A') \subset \Spec(A)$ 为 $A'$ 的极小素理想集。
由《代数》引理 \ref{algebra-lemma-injective-minimal-primes-in-image}，
集合 $E'$ 包含于
$\Spec(B) \to \Spec(A') \subset \Spec(A)$ 的像中。
因此 $E' \subset E$。由于 $A'$ 是诺特的，由《代数》引理
\ref{algebra-lemma-Noetherian-irreducible-components}，$E'$ 有限。
由于 $\Spec(B) \to \Spec(A)$ 的像中的任何其他点
都是 $E'$ 的某个元素的专化，并且属于 $E$，
可知该像包含于 $E'$ 中（由我们关于 $E$ 中各点之间专化的假设）。
于是归约到上面已经处理的 $E$ 有限的情形。
\end{proof}
% Generated mechanically from the frozen coherent.tex target.
% Do not edit this generated file; edit the bound locale source instead.

% OK, start here.
%

\setcounter{chapter}{29}
\chapter{概形的上同调}



\phantomsection
\label{coherent-section-phantom}


\section{引言}
\label{coherent-section-introduction}

\noindent
本章首先证明关于准凝聚层上同调的一系列结果。
一个基础参考文献是 \cite{EGA}。
在此基础上，我们将进一步讨论 Noether 情形下凝聚层的上同调。
见 \cite{FAC}。







\section{准凝聚层的 {\v C}ech 上同调}
\label{coherent-section-cech-quasi-coherent}

\noindent
设 $X$ 为概形。
设 $U \subset X$ 为仿射开集。
回忆，$U$ 的一个{\it 标准开覆盖}是形如
$\mathcal{U} : U = \bigcup_{i = 1}^n D(f_i)$ 的覆盖，其中
$f_1, \ldots, f_n \in \Gamma(U, \mathcal{O}_X)$ 生成单位理想；见
《概形》定义 \ref{schemes-definition-standard-covering}。

\begin{lemma}
\label{coherent-lemma-cech-cohomology-quasi-coherent-trivial}
设 $X$ 为概形。
设 $\mathcal{F}$ 为准凝聚 $\mathcal{O}_X$-模。
设 $\mathcal{U} : U = \bigcup_{i = 1}^n D(f_i)$ 为 $X$ 的一个仿射开集的
标准开覆盖。
则 $\check{H}^p(\mathcal{U}, \mathcal{F}) = 0$ 对所有 $p > 0$ 均成立。
\end{lemma}


\begin{proof}
写成 $U = \Spec(A)$，其中 $A$ 为某个环。
换言之，$f_1, \ldots, f_n$ 是 $A$ 中生成 $A$ 的单位理想的元素。
写成 $\mathcal{F}|_U = \widetilde{M}$，其中这个 $A$-模记作 $M$。
显然，{\v C}ech 复形
$\check{\mathcal{C}}^\bullet(\mathcal{U}, \mathcal{F})$
可与下列复形等同：
$$
\prod\nolimits_{i_0} M_{f_{i_0}} \to
\prod\nolimits_{i_0i_1} M_{f_{i_0}f_{i_1}} \to
\prod\nolimits_{i_0i_1i_2} M_{f_{i_0}f_{i_1}f_{i_2}} \to
\ldots
$$
我们需要证明扩充复形
\begin{equation}
\label{coherent-equation-extended}
0 \to
M \to
\prod\nolimits_{i_0} M_{f_{i_0}} \to
\prod\nolimits_{i_0i_1} M_{f_{i_0}f_{i_1}} \to
\prod\nolimits_{i_0i_1i_2} M_{f_{i_0}f_{i_1}f_{i_2}} \to
\ldots
\end{equation}

正合；我们已在《代数》引理 \ref{algebra-lemma-cover-module} 中研究过
它的截断。
只需证明复形 (\ref{coherent-equation-extended}) 在任意素理想 $\mathfrak p$
处局部化后正合；见《代数》引理
\ref{algebra-lemma-characterize-zero-local}。
事实上，我们将证明该扩充复形在 $\mathfrak p$ 处局部化后零伦。

\medskip\noindent
存在指标 $i$，使得 $f_i \not \in \mathfrak p$。
选定这样的指标，并记为 $i_{\text{fix}}$。注意
$M_{f_{i_{\text{fix}}}, \mathfrak p} = M_{\mathfrak p}$。类似地，
对乘积 $f_{i_0} \ldots f_{i_p}$ 再在 $\mathfrak p$ 处局部化时，
其中任一因子 $f_{i_j}$，只要 $i_j = i_{\text{fix}}$，均可略去。
定义同伦
$$
h :
\prod\nolimits_{i_0 \ldots i_{p + 1}}
M_{f_{i_0} \ldots f_{i_{p + 1}}, \mathfrak p}
\longrightarrow
\prod\nolimits_{i_0 \ldots i_p}
M_{f_{i_0} \ldots f_{i_p}, \mathfrak p}
$$
如下：
$$
h(s)_{i_0 \ldots i_p} = s_{i_{\text{fix}} i_0 \ldots i_p}
$$
（这与《上同调》引理
\ref{cohomology-lemma-homology-complex} 证明中的同伦“对偶”。）
换言之，$h : \prod_{i_0} M_{f_{i_0}, \mathfrak p} \to M_\mathfrak p$
是到因子
$M_{f_{i_{\text{fix}}}, \mathfrak p} = M_{\mathfrak p}$ 的投影；一般地，
映射 $h$ 是到诸因子
$M_{f_{i_{\text{fix}}} f_{i_1} \ldots f_{i_{p + 1}}, \mathfrak p}
= M_{f_{i_1} \ldots f_{i_{p + 1}}, \mathfrak p}$ 的投影。计算得
\begin{align*}
(dh + hd)(s)_{i_0 \ldots i_p}
& =
\sum\nolimits_{j = 0}^p
(-1)^j
h(s)_{i_0 \ldots \hat i_j \ldots i_p}
+
d(s)_{i_{\text{fix}} i_0 \ldots i_p}\\
& =
\sum\nolimits_{j = 0}^p
(-1)^j
s_{i_{\text{fix}} i_0 \ldots \hat i_j \ldots i_p}
+
s_{i_0 \ldots i_p}
+
\sum\nolimits_{j = 0}^p
(-1)^{j + 1}
s_{i_{\text{fix}} i_0 \ldots \hat i_j \ldots i_p} \\
& =
s_{i_0 \ldots i_p}
\end{align*}
这就按要求证明了恒等映射零伦。
\end{proof}


\noindent
下一个引理特别表明：对任意仿射概形 $X$ 以及任意准凝聚层
$\mathcal{F}$（定义在 $X$ 上），均有
$$
H^p(X, \mathcal{F}) = 0
$$
其中 $p > 0$。

\begin{lemma}
\label{coherent-lemma-quasi-coherent-affine-cohomology-zero}
\begin{slogan}
Serre 消失：仿射概形上准凝聚模的高阶上同调消失。
\end{slogan}
设 $X$ 为概形。
设 $\mathcal{F}$ 为准凝聚 $\mathcal{O}_X$-模。
对任意仿射开集 $U \subset X$，均有
$H^p(U, \mathcal{F}) = 0$，其中 $p > 0$。
\end{lemma}

\begin{proof}
我们应用《上同调》引理
\ref{cohomology-lemma-cech-vanish-basis}。
取拓扑基 $\mathcal{B}$；这里考虑 $X$ 的拓扑，并使用 $X$ 的全部仿射开集。
取开覆盖集 $\text{Cov}$ 为 $X$ 的仿射开集的全部标准开覆盖。
下面验证《上同调》引理
\ref{cohomology-lemma-cech-vanish-basis} 的条件 (1)、(2) 和 (3)。
注意，仿射概形中两个标准开集的交仍为标准开集，故性质 (1) 成立。
这些覆盖构成 $\mathcal{B}$ 中任一元素的开覆盖的共尾系；见
《概形》引理 \ref{schemes-lemma-standard-open}。
故 (2) 成立。
最后，条件 (3) 来自引理
\ref{coherent-lemma-cech-cohomology-quasi-coherent-trivial}。
\end{proof}

\noindent
下面给出仿射概形上准凝聚层上同调消失的相对版本。

\begin{lemma}
\label{coherent-lemma-relative-affine-vanishing}
设 $f : X \to S$ 为概形态射。
设 $\mathcal{F}$ 为准凝聚 $\mathcal{O}_X$-模。
若 $f$ 仿射，则 $R^if_*\mathcal{F} = 0$ 对所有 $i > 0$ 均成立。
\end{lemma}

\begin{proof}
按照《上同调》引理
\ref{cohomology-lemma-describe-higher-direct-images}，层
$R^if_*\mathcal{F}$ 是预层
$V \mapsto H^i(f^{-1}(V), \mathcal{F}|_{f^{-1}(V)})$ 的层化。
由假设，只要 $V$ 仿射，$f^{-1}(V)$ 就仿射；见《态射》定义
\ref{morphisms-definition-affine}。
由引理 \ref{coherent-lemma-quasi-coherent-affine-cohomology-zero} 可知，
$H^i(f^{-1}(V), \mathcal{F}|_{f^{-1}(V)}) = 0$
只要 $V$ 仿射就成立。
由于 $S$ 有一个由仿射开集构成的基，结论成立。
\end{proof}

\begin{lemma}
\label{coherent-lemma-relative-affine-cohomology}
设 $f : X \to S$ 为概形的仿射态射。
设 $\mathcal{F}$ 为准凝聚 $\mathcal{O}_X$-模。
则 $H^i(X, \mathcal{F}) = H^i(S, f_*\mathcal{F})$
对所有 $i \geq 0$ 均成立。
\end{lemma}

\begin{proof}
这来自引理 \ref{coherent-lemma-relative-affine-vanishing} 以及 Leray 谱序列。
见《上同调》引理 \ref{cohomology-lemma-apply-Leray}。
\end{proof}

\noindent
下面两个引理说明何时可用 {\v C}ech 上同调计算准凝聚模的上同调。

\begin{lemma}
\label{coherent-lemma-affine-diagonal}
设 $X$ 为概形。下列条件等价：
\begin{enumerate}
\item $X$ 的对角态射 $\Delta : X \to X \times X$ 仿射；
\item 对任意仿射开集 $U, V \subset X$，交 $U \cap V$ 仿射；以及
\item 存在开覆盖 $\mathcal{U} : X = \bigcup_{i \in I} U_i$，使得
$U_{i_0 \ldots i_p}$ 对所有 $p \ge 0$ 以及所有
$i_0, \ldots, i_p \in I$ 都是仿射开集。
\end{enumerate}
特别地，若 $X$ 分离，则这些条件成立。
\end{lemma}

\begin{proof}
假设 $X$ 的对角态射仿射。设 $U, V \subset X$ 为仿射开集。
则 $U \cap V = \Delta^{-1}(U \times V)$ 仿射，故 (2) 成立。
(2) 蕴含 (3) 是立即的。反之，若 $X$ 有如 (3) 所述的覆盖，则
$X \times X = \bigcup U_i \times U_{i'}$ 是仿射开覆盖，并且
$\Delta^{-1}(U_i \times U_{i'}) = U_i \cap U_{i'}$
仿射。于是由《态射》引理
\ref{morphisms-lemma-characterize-affine}，$\Delta$ 是仿射态射。
最后一个断言来自《概形》引理
\ref{schemes-lemma-characterize-separated}。
\end{proof}

\begin{lemma}
\label{coherent-lemma-cech-cohomology-quasi-coherent}
设 $X$ 为概形。
设 $\mathcal{U} : X = \bigcup_{i \in I} U_i$ 为开覆盖，并且
$U_{i_0 \ldots i_p}$ 对所有 $p \ge 0$ 以及所有
$i_0, \ldots, i_p \in I$ 都是仿射开集。
在这种情况下，对任意准凝聚层 $\mathcal{F}$，均有
$$
\check{H}^p(\mathcal{U}, \mathcal{F}) = H^p(X, \mathcal{F})
$$
作为 $\Gamma(X, \mathcal{O}_X)$-模的等式，并对所有 $p$ 成立。
\end{lemma}

\begin{proof}
由引理 \ref{coherent-lemma-quasi-coherent-affine-cohomology-zero}，这是
《上同调》引理
\ref{cohomology-lemma-cech-spectral-sequence-application} 的特殊情形。
\end{proof}







\section{上同调的消失}
\label{coherent-section-vanishing}

\noindent
我们已经看到，在仿射概形上，任意准凝聚层的高阶上同调群都消失
（引理 \ref{coherent-lemma-quasi-coherent-affine-cohomology-zero}）。
事实证明，这一性质也刻画仿射概形。我们给出两个版本。

\begin{lemma}
\label{coherent-lemma-quasi-compact-h1-zero-covering}
\begin{reference}
\cite{Serre-criterion}, \cite[II, Theorem 5.2.1 (d') and IV (1.7.17)]{EGA}
\end{reference}
\begin{slogan}
Serre 仿射性判别准则。
\end{slogan}
设 $X$ 为概形。假设：
\begin{enumerate}
\item $X$ 拟紧；
\item 对每个准凝聚理想层
$\mathcal{I} \subset \mathcal{O}_X$，均有 $H^1(X, \mathcal{I}) = 0$。
\end{enumerate}
则 $X$ 仿射。
\end{lemma}

\begin{proof}
设 $x \in X$ 为闭点。设 $U \subset X$ 为 $x$ 的一个仿射开邻域。
写成 $U = \Spec(A)$，并令 $\mathfrak m \subset A$ 为对应于 $x$ 的极大理想。
置 $Z = X \setminus U$ 以及 $Z' = Z \cup \{x\}$。
由《概形》引理 \ref{schemes-lemma-reduced-closed-subscheme}，存在准凝聚理想层
$\mathcal{I}$ 和 $\mathcal{I}'$，它们分别截出约化闭子概形
$Z$ 和 $Z'$。
考虑短正合列
$$
0 \to \mathcal{I}' \to \mathcal{I} \to \mathcal{I}/\mathcal{I}' \to 0.
$$
由于 $x$ 是 $X$ 的闭点且 $x \not \in Z$，
$\mathcal{I}/\mathcal{I}'$ 支持在 $x$ 上。事实上，
$\mathcal{I}/\mathcal{I'}$ 到 $U$ 的限制对应于 $A$-模
$A/\mathfrak m$。因此
$\Gamma(X, \mathcal{I}/\mathcal{I'})
= A/\mathfrak m$。
由于依假设 $H^1(X, \mathcal{I}') = 0$，存在整体截面
$f \in \Gamma(X, \mathcal{I})$，它映到元素 $1 \in A/\mathfrak m$；
这里把它视为 $\mathcal{I}/\mathcal{I'}$ 的截面。
显然 $x \in X_f \subset U$。这意味着
$X_f = D(f_A)$，其中 $f_A$ 是 $f$ 在
$A = \Gamma(U, \mathcal{O}_X)$ 中的像。特别地，$X_f$ 仿射。

\medskip\noindent
考虑并集 $W = \bigcup X_f$，其中遍历所有
$f \in \Gamma(X, \mathcal{O}_X)$ 且 $X_f$ 仿射。显然 $W$ 是 $X$ 的开集。
由上述论证，$X$ 的每个闭点都包含在 $W$ 中。
闭子集 $X \setminus W$（它是 $X$ 中的闭子集）也拟紧
（见《拓扑学》引理 \ref{topology-lemma-closed-in-quasi-compact}）。
因此，如果它非空，就有闭点（见《拓扑学》引理
\ref{topology-lemma-quasi-compact-closed-point}）。
这将与所有闭点都在 $W$ 中矛盾。故 $X = W$。

\medskip\noindent
选择有限多个 $f_1, \ldots, f_n \in \Gamma(X, \mathcal{O}_X)$，
使得 $X = X_{f_1} \cup \ldots \cup X_{f_n}$，并且每个
$X_{f_i}$ 都仿射。由上述论证，这样的选择是可能的。
由《性质》引理 \ref{properties-lemma-characterize-affine}，
为完成证明，只需证明 $f_1, \ldots, f_n$ 在
$\Gamma(X, \mathcal{O}_X)$ 中生成单位理想。考虑短正合列
$$
\xymatrix{
0 \ar[r] &
\mathcal{F} \ar[r] &
\mathcal{O}_X^{\oplus n} \ar[rr]^{f_1, \ldots, f_n} & &
\mathcal{O}_X \ar[r] &
0
}
$$
由 $f_1, \ldots, f_n$ 定义的箭头满射，因为诸开集
$X_{f_i}$ 覆盖 $X$。令 $\mathcal{F}$ 为该满射的核。
注意，$\mathcal{F}$ 有滤过
$$
0 = \mathcal{F}_0 \subset \mathcal{F}_1 \subset
\ldots \subset \mathcal{F}_n = \mathcal{F}
$$
使得每个子商 $\mathcal{F}_i/\mathcal{F}_{i - 1}$
都同构于一个准凝聚理想层。
具体而言，可取 $\mathcal{F}_i$ 为 $\mathcal{F}$ 与前 $i$ 个直和因子的交；
这些因子属于 $\mathcal{O}_X^{\oplus n}$。
引理的假设蕴含
$H^1(X, \mathcal{F}_i/\mathcal{F}_{i - 1}) = 0$
对所有 $i$ 都成立。于是 $H^1(X, \mathcal{F}_2) = 0$，
因为它夹在 $H^1(X, \mathcal{F}_1)$ 与
$H^1(X, \mathcal{F}_2/\mathcal{F}_1)$ 之间。
如此继续，可得 $H^1(X, \mathcal{F}) = 0$。
因此，映射
$$
\xymatrix{
\bigoplus\nolimits_{i = 1, \ldots, n} \Gamma(X, \mathcal{O}_X)
\ar[rr]^{f_1, \ldots, f_n} & &
\Gamma(X, \mathcal{O}_X)
}
$$
如所需为满射。
\end{proof}

\noindent
注意，若 $X$ 是 Noether 概形，则每个准凝聚理想层自动是凝聚理想层，
也是有限型准凝聚理想层。因此在这种情况下，前一引理与下一引理均适用。

\begin{lemma}
\label{coherent-lemma-quasi-separated-h1-zero-covering}
\begin{reference}
\cite{Serre-criterion}, \cite[II, Theorem 5.2.1]{EGA}
\end{reference}
\begin{slogan}
Serre 仿射性判别准则。
\end{slogan}
设 $X$ 为概形。假设：
\begin{enumerate}
\item $X$ 拟紧；
\item $X$ 拟分离；并且
\item $H^1(X, \mathcal{I}) = 0$ 对每个有限型准凝聚理想层
$\mathcal{I}$ 均成立。
\end{enumerate}
则 $X$ 仿射。
\end{lemma}

\begin{proof}
由《性质》引理
\ref{properties-lemma-quasi-coherent-colimit-finite-type}，
每个准凝聚理想层都是有限型准凝聚理想层的有向余极限。
由《上同调》引理
\ref{cohomology-lemma-quasi-separated-cohomology-colimit}，
在 $X$ 上取上同调与有向余极限可交换。
因此，$H^1(X, \mathcal{I}) = 0$ 对 $X$ 上的每个准凝聚理想层均成立。
换言之，引理
\ref{coherent-lemma-quasi-compact-h1-zero-covering} 适用。
\end{proof}

\noindent
可以利用上述论证给出一个判定可逆层何时丰沛的充分条件。
不过我们提醒读者，这个条件并非必要。

\begin{lemma}
\label{coherent-lemma-quasi-compact-h1-zero-invertible}
设 $X$ 为概形。设 $\mathcal{L}$ 为可逆 $\mathcal{O}_X$-模。
假设：
\begin{enumerate}
\item $X$ 拟紧；
\item 对每个准凝聚理想层
$\mathcal{I} \subset \mathcal{O}_X$，存在 $n \geq 1$，使得
$H^1(X, \mathcal{I} \otimes_{\mathcal{O}_X} \mathcal{L}^{\otimes n}) = 0$。
\end{enumerate}
则 $\mathcal{L}$ 丰沛。
\end{lemma}

\begin{proof}
证明与引理 \ref{coherent-lemma-quasi-compact-h1-zero-covering} 完全相同。
设 $x \in X$ 为闭点。设 $U \subset X$ 为 $x$ 的仿射开邻域，并且
$\mathcal{L}|_U \cong \mathcal{O}_U$。
写成 $U = \Spec(A)$，并令 $\mathfrak m \subset A$
为对应于 $x$ 的极大理想。
置 $Z = X \setminus U$ 以及 $Z' = Z \cup \{x\}$。
由《概形》引理 \ref{schemes-lemma-reduced-closed-subscheme}，存在准凝聚理想层
$\mathcal{I}$ 和 $\mathcal{I}'$，它们分别截出约化闭子概形
$Z$ 和 $Z'$。考虑短正合列
$$
0 \to \mathcal{I}' \to \mathcal{I} \to \mathcal{I}/\mathcal{I}' \to 0.
$$
对每个 $n \geq 1$，得到短正合列
$$
0 \to \mathcal{I}' \otimes_{\mathcal{O}_X} \mathcal{L}^{\otimes n}
\to \mathcal{I} \otimes_{\mathcal{O}_X} \mathcal{L}^{\otimes n} \to
\mathcal{I}/\mathcal{I}' \otimes_{\mathcal{O}_X} \mathcal{L}^{\otimes n} \to 0.
$$
由假设，可选择 $n$，使得
$H^1(X, \mathcal{I}' \otimes_{\mathcal{O}_X} \mathcal{L}^{\otimes n}) = 0$。
由于 $x$ 是 $X$ 的闭点且 $x \not \in Z$，
$\mathcal{I}/\mathcal{I}'$ 支持在 $x$ 上。事实上，
$\mathcal{I}/\mathcal{I'}$ 到 $U$ 的限制对应于 $A$-模
$A/\mathfrak m$。由于 $\mathcal{L}$ 在 $U$ 上平凡，
$\mathcal{I}/\mathcal{I}' \otimes_{\mathcal{O}_X} \mathcal{L}^{\otimes n}$
到 $U$ 的限制也对应于 $A$-模 $A/\mathfrak m$。
因此
$\Gamma(X, \mathcal{I}/\mathcal{I'} \otimes_{\mathcal{O}_X}
\mathcal{L}^{\otimes n}) = A/\mathfrak m$。
根据所选的 $n$，存在整体截面
$s \in \Gamma(X, \mathcal{I} \otimes_{\mathcal{O}_X} \mathcal{L}^{\otimes n})$，
它映到元素 $1 \in A/\mathfrak m$。显然
$x \in X_s \subset U$，因为 $s$ 在 $Z$ 的点处消失。
这意味着 $X_s = D(f)$，其中 $f \in A$ 是 $s$ 在
$A \cong \Gamma(U, \mathcal{L}^{\otimes n})$ 中的像。
特别地，$X_s$ 仿射。

\medskip\noindent
考虑并集 $W = \bigcup X_s$，其中遍历所有
$s \in \Gamma(X, \mathcal{L}^{\otimes n})$，取 $n \geq 1$，并要求
$X_s$ 仿射。显然 $W$ 是 $X$ 的开集。
由上述论证，$X$ 的每个闭点都包含在 $W$ 中。
闭子集 $X \setminus W$（它是 $X$ 中的闭子集）也拟紧
（见《拓扑学》引理 \ref{topology-lemma-closed-in-quasi-compact}）。
因此，如果它非空，就有闭点（见《拓扑学》引理
\ref{topology-lemma-quasi-compact-closed-point}）。
这将与所有闭点都在 $W$ 中矛盾。故 $X = W$。
由《性质》定义 \ref{properties-definition-ample}，
这意味着 $\mathcal{L}$ 丰沛。
\end{proof}

\noindent
引理 \ref{coherent-lemma-quasi-compact-h1-zero-invertible} 有一个使用有限型理想层的
变体；若以后需要，我们将在此陈述并证明。

\begin{lemma}
\label{coherent-lemma-criterion-affine-morphism}
设 $f : X \to Y$ 为拟紧态射，并且 $X$ 和 $Y$ 均拟分离。
若 $R^1f_*\mathcal{I} = 0$ 对每个准凝聚理想层
$\mathcal{I}$（定义在 $X$ 上）均成立，则 $f$ 仿射。
\end{lemma}

\begin{proof}
设 $V \subset Y$ 为仿射开子概形。我们需要证明
$U = f^{-1}(V)$ 仿射。由《概形》引理
\ref{schemes-lemma-quasi-compact-permanence}，包含态射 $V \to Y$ 拟紧。
因此基变换 $U \to X$ 拟紧；见《概形》引理
\ref{schemes-lemma-quasi-compact-preserved-base-change}。
于是，任意准凝聚理想层 $\mathcal{I}$（定义在 $U$ 上）
都可延拓为 $X$ 上的准凝聚理想层；见《性质》引理
\ref{properties-lemma-extend-trivial}。
由于 $R^1f_*$ 的构造在 $Y$ 上是局部的
（《上同调》节 \ref{cohomology-section-locality}），
由引理中的假设可得 $R^1(U \to V)_*\mathcal{I} = 0$。
因此由 Leray 谱序列
（《上同调》引理 \ref{cohomology-lemma-Leray}），可得
$H^1(U, \mathcal{I}) = H^1(V, (U \to V)_*\mathcal{I})$。
由《概形》引理
\ref{schemes-lemma-push-forward-quasi-coherent}，
$(U \to V)_*\mathcal{I}$ 准凝聚，故由引理
\ref{coherent-lemma-quasi-coherent-affine-cohomology-zero}，
$H^1(V, (U \to V)_*\mathcal{I}) = 0$。
于是由引理 \ref{coherent-lemma-quasi-compact-h1-zero-covering}，
$U$ 仿射。
\end{proof}















\section{高阶直像的准凝聚性}
\label{coherent-section-quasi-coherence}

\noindent
我们已经看到，仿射概形上一个准凝聚模的高阶上同调群为零。
对拟分离拟紧概形 $X$，这蕴含准凝聚层的上同调群在某一次数以上消失。
然而，$X$ 未必具有有限上同调维数，因为上同调维数的定义要求
{\it 所有} $\mathcal{O}_X$-模的上同调都消失。

\begin{lemma}[归纳原理]
\label{coherent-lemma-induction-principle}
\begin{reference}
\cite[Proposition 3.3.1]{BvdB}
\end{reference}
设 $X$ 为拟紧拟分离概形。设 $P$ 是 $X$ 的拟紧开集所具有的一个性质。
假设：
\begin{enumerate}
\item $P$ 对 $X$ 的每个仿射开集成立；
\item 若 $U$ 是拟紧开集，$V$ 是仿射开集，且 $P$ 对
$U$、$V$ 和 $U \cap V$ 成立，则 $P$ 对 $U \cup V$ 成立。
\end{enumerate}
则 $P$ 对 $X$ 的每个拟紧开集成立，特别地对 $X$ 成立。
\end{lemma}

\begin{proof}
首先用归纳法证明 $P$ 对{\it 分离的}拟紧开集 $W \subset X$ 成立。
事实上，可将这样的开集写成
$W = U_1 \cup \ldots \cup U_n$，并对 $n$ 归纳；
在性质 (2) 中取 $U = U_1 \cup \ldots \cup U_{n - 1}$ 和 $V = U_n$。
这是允许的，因为
$U \cap V = (U_1 \cap U_n) \cup \ldots \cup (U_{n - 1} \cap U_n)$
也是 $n - 1$ 个仿射开子概形之并；这里把《概形》引理
\ref{schemes-lemma-characterize-separated} 应用于仿射开集 $U_i$ 和 $U_n$，
它们均为 $W$ 的仿射开集。
由此，对任意拟紧开集 $W \subset X$，可对覆盖 $W$ 所需的仿射开集数目
归纳，使用与前面相同的办法，并利用如下事实：拟紧开集
$U_i \cap U_n$ 作为仿射概形 $U_n$ 的开子概形是分离的。
\end{proof}

\begin{lemma}
\label{coherent-lemma-vanishing-nr-affines}
\begin{slogan}
对具有仿射对角态射的概形，准凝聚模的上同调在次数大于某个覆盖所需的
仿射开集数目时消失。
\end{slogan}
设 $X$ 为具有仿射对角态射的拟紧概形（例如，若 $X$ 分离）。
令 $t = t(X)$ 为覆盖 $X$ 所需的仿射开集的最小数目。
则均有 $H^n(X, \mathcal{F}) = 0$，其中 $n \geq t$ 任意，
$\mathcal{F}$ 为任意准凝聚层。
\end{lemma}

\begin{proof}
第一种证明。
对 $t$ 归纳。
若 $t = 1$，结论来自引理
\ref{coherent-lemma-quasi-coherent-affine-cohomology-zero}。
若 $t > 1$，写成 $X = U \cup V$，其中 $V$ 为仿射开集，而
$U = U_1 \cup \ldots \cup U_{t - 1}$ 是 $t - 1$ 个仿射开集之并。
注意，此时
$U \cap V =  (U_1 \cap V) \cup \ldots (U_{t - 1} \cap V)$
也是 $t - 1$ 个仿射开子概形之并。
事实上，因为对角态射仿射，两个仿射开集的交仿射；见引理
\ref{coherent-lemma-affine-diagonal}。
应用 Mayer--Vietoris 长正合列
$$
0 \to
H^0(X, \mathcal{F}) \to
H^0(U, \mathcal{F}) \oplus H^0(V, \mathcal{F}) \to
H^0(U \cap V, \mathcal{F}) \to
H^1(X, \mathcal{F}) \to \ldots
$$
见《上同调》引理 \ref{cohomology-lemma-mayer-vietoris}。
由归纳可知，群 $H^i(U, \mathcal{F})$、
$H^i(V, \mathcal{F})$ 和 $H^i(U \cap V, \mathcal{F})$
在 $i \geq t - 1$ 时为零。因此立即得到
$H^i(X, \mathcal{F})$ 在 $i \geq t$ 时为零。

\medskip\noindent
第二种证明。
设 $\mathcal{U} : X = \bigcup_{i = 1}^t U_i$ 为有限仿射开覆盖。
由于 $X$ 具有仿射对角态射，多重交
$U_{i_0 \ldots i_p}$ 都仿射；见引理 \ref{coherent-lemma-affine-diagonal}。
由引理 \ref{coherent-lemma-cech-cohomology-quasi-coherent}，{\v C}ech 上同调群
$\check{H}^p(\mathcal{U}, \mathcal{F})$ 与上同调群相同。
由《上同调》引理 \ref{cohomology-lemma-alternating-usual}，
可用交错 {\v C}ech 复形
$\check{\mathcal{C}}_{alt}^\bullet(\mathcal{U}, \mathcal{F})$
计算 {\v C}ech 上同调群。由于覆盖由 $t$ 个元素组成，立即得到
$\check{\mathcal{C}}_{alt}^p(\mathcal{U}, \mathcal{F}) = 0$
对所有 $p \geq t$ 成立。因此结论成立。
\end{proof}

\begin{lemma}
\label{coherent-lemma-affine-diagonal-universal-delta-functor}
设 $X$ 为具有仿射对角态射的拟紧概形
（例如，若 $X$ 分离）。则：
\begin{enumerate}
\item 给定准凝聚 $\mathcal{O}_X$-模 $\mathcal{F}$，存在嵌入
$\mathcal{F} \to \mathcal{F}'$，其中后者是准凝聚
$\mathcal{O}_X$-模，
使得 $H^p(X, \mathcal{F}') = 0$ 对所有 $p \geq 1$ 成立；并且
\item $\{H^n(X, -)\}_{n \geq 0}$ 是泛 $\delta$-函子，
其源为 $\QCoh(\mathcal{O}_X)$，靶为 $\textit{Ab}$。
\end{enumerate}
\end{lemma}

\begin{proof}
设 $X = \bigcup U_i$ 为有限仿射开覆盖。
置 $U = \coprod U_i$，并以 $j : U \to X$ 表示诱导给定开浸入
$U_i \to X$ 的态射。
由于 $U$ 是仿射概形，且 $X$ 有仿射对角态射，态射 $j$ 仿射；见《态射》引理
\ref{morphisms-lemma-affine-permanence}。
对每个 $\mathcal{O}_X$-模 $\mathcal{F}$，存在典范映射
$\mathcal{F} \to j_*j^*\mathcal{F}$。
在茎上检验可知该映射单射：若 $x \in U_i$，则有分解
$$
\mathcal{F}_x \to (j_*j^*\mathcal{F})_x
\to (j^*\mathcal{F})_{x'} = \mathcal{F}_x
$$
其中 $x' \in U$ 是把点 $x$ 看成 $U_i \subset U$ 的点。
现在，若 $\mathcal{F}$ 准凝聚，则 $j^*\mathcal{F}$
在仿射概形 $U$ 上准凝聚，故由引理
\ref{coherent-lemma-quasi-coherent-affine-cohomology-zero}，其高阶上同调消失。
又有 $H^p(X, j_*j^*\mathcal{F}) = 0$ 对 $p > 0$ 成立；
这是由引理 \ref{coherent-lemma-relative-affine-cohomology} 以及 $j$ 仿射得到的。
这证明 (1)。
最后，映射
$H^p(X, \mathcal{F}) \to H^p(X, j_*j^*\mathcal{F})$
为零，而 (2) 来自《同调代数》引理
\ref{homology-lemma-efface-implies-universal}。
\end{proof}

\begin{lemma}
\label{coherent-lemma-vanishing-nr-affines-quasi-separated}
设 $X$ 为拟紧拟分离概形。
设 $X = U_1 \cup \ldots \cup U_n$ 为开覆盖，其中每个 $U_i$
都拟紧且分离（例如仿射）。置
$$
d = \max\nolimits_{I \subset \{1, \ldots, n\}}
\left(|I| + t(\bigcap\nolimits_{i \in I} U_i) - 1\right)
$$
其中 $t(U)$ 是覆盖概形 $U$ 所需的仿射开集的最小数目。
则均有 $H^p(X, \mathcal{F}) = 0$，其中 $p \geq d$ 任意，
$\mathcal{F}$ 为任意准凝聚层。
\end{lemma}

\begin{proof}
注意，由于 $X$ 拟分离且 $U_i$ 拟紧，数
$t(\bigcap_{i \in I} U_i)$ 有限。对 $n$ 归纳证明。
若 $n = 1$，结论来自引理 \ref{coherent-lemma-vanishing-nr-affines}。
若 $n > 1$，写成 $X = U \cup V$，其中
$U = U_1 \cup \ldots \cup U_{n - 1}$ 且 $V = U_n$。
应用 Mayer--Vietoris 长正合列
$$
0 \to
H^0(X, \mathcal{F}) \to
H^0(U, \mathcal{F}) \oplus H^0(V, \mathcal{F}) \to
H^0(U \cap V, \mathcal{F}) \to
H^1(X, \mathcal{F}) \to \ldots
$$
见《上同调》引理 \ref{cohomology-lemma-mayer-vietoris}。
为对 $q \geq d$ 完成证明，我们将证明
$H^q(V, \mathcal{F})$、$H^q(U, \mathcal{F})$ 和
$H^{q - 1}(U \cap V, \mathcal{F})$ 消失。
由 $n = 1$ 的情形，$H^q(V, \mathcal{F}) = 0$ 对
$q \geq t(V) = t(U_n)$ 成立。由于 $t(V) \leq d$，这就证明了所需断言。
由归纳假设，$H^q(U, \mathcal{F}) = 0$ 在下列范围成立：
$$
q \geq \max\nolimits_{I \subset \{1, \ldots, n - 1\}}
\left(|I| + t(\bigcap\nolimits_{i \in I} U_i) - 1\right)
$$
由于右端整数小于或等于 $d$，这就证明了所需断言。
最后，可对开集
$U \cap V = (U_1 \cap U_n) \cup \ldots \cup (U_{n - 1} \cap U_n)$
使用归纳假设，得到 $H^q(U \cap V, \mathcal{F}) = 0$ 在下列范围成立：
$$
q \geq \max\nolimits_{I \subset \{1, \ldots, n - 1\}}
\left(|I| + t(U_n \cap \bigcap\nolimits_{i \in I} U_i) - 1\right)
$$
由于右端整数严格小于 $d$，引理得证。
\end{proof}

\begin{lemma}
\label{coherent-lemma-quasi-coherence-higher-direct-images}
设 $f : X \to S$ 为概形态射。
假设 $f$ 拟分离且拟紧。
\begin{enumerate}
\item 对任意准凝聚 $\mathcal{O}_X$-模 $\mathcal{F}$，
高阶直像 $R^pf_*\mathcal{F}$ 在 $S$ 上准凝聚。
\item 若 $S$ 拟紧，则存在整数 $n = n(X, S, f)$，使得
$R^pf_*\mathcal{F} = 0$；这里 $p \geq n$ 任意，$\mathcal{F}$
为定义在 $X$ 上的任意准凝聚层。
\item 事实上，若 $S$ 拟紧，则可找到 $n = n(X, S, f)$，使得对每个
概形态射 $S' \to S$，都有 $R^p(f')_*\mathcal{F}' = 0$
对 $p \geq n$ 以及任意准凝聚层 $\mathcal{F}'$ 成立，后者定义在 $X'$ 上。
这里 $f' : X' = S' \times_S X \to S'$ 是 $f$ 的基变换。
\end{enumerate}
\end{lemma}

\begin{proof}
先证明 (1)。注意，在引理的假设下，层
$R^0f_*\mathcal{F} = f_*\mathcal{F}$ 由《概形》引理
\ref{schemes-lemma-push-forward-quasi-coherent} 是准凝聚的。
利用《上同调》引理
\ref{cohomology-lemma-localize-higher-direct-images}，
可知形成高阶直像与限制到开子概形可交换。
由于准凝聚性在 $S$ 上是局部性质，问题归约为下一段所讨论的情形。

\medskip\noindent
在 $S$ 仿射的情形下证明 (1)。我们使用归纳原理。
由于 $f$ 拟紧且拟分离，$X$ 拟紧且拟分离。
对拟紧开集 $U \subset X$，令 $a = f|_U$，并令 $P(U)$ 表示如下性质：
$R^pa_*\mathcal{F}$ 在 $S$ 上准凝聚；这对所有准凝聚模
$\mathcal{F}$（定义在 $U$ 上）以及所有 $p \geq 0$ 成立。
由于 $P(X)$ 就是 (1)，只需证明引理
\ref{coherent-lemma-induction-principle} 的条件 (1) 和 (2) 成立。
若 $U$ 仿射，则 $P(U)$ 成立，因为
$R^pa_*\mathcal{F} = 0$ 对 $p \geq 1$ 成立
（由引理 \ref{coherent-lemma-relative-affine-vanishing} 以及《态射》引理
\ref{morphisms-lemma-morphism-affines-affine}），并且我们在第一段已经观察到
$p = 0$ 时结论成立。
接下来，设 $U \subset X$ 为拟紧开集，$V \subset X$ 为仿射开集，并假设
$P(U)$、$P(V)$ 和 $P(U \cap V)$ 成立。
令 $a = f|_U$、$b = f|_V$、$c = f|_{U \cap V}$ 以及
$g = f|_{U \cup V}$。则对任意准凝聚
$\mathcal{O}_{U \cup V}$-模 $\mathcal{F}$，有相对
Mayer--Vietoris 序列
$$
0 \to
g_*\mathcal{F} \to
a_*(\mathcal{F}|_U) \oplus b_*(\mathcal{F}|_V) \to
c_*(\mathcal{F}|_{U \cap V}) \to
R^1g_*\mathcal{F} \to \ldots
$$
见《上同调》引理 \ref{cohomology-lemma-relative-mayer-vietoris}。
由 $P(U)$、$P(V)$ 和 $P(U \cap V)$ 可知，
$R^pa_*(\mathcal{F}|_U)$、$R^pb_*(\mathcal{F}|_V)$ 以及
$R^pc_*(\mathcal{F}|_{U \cap V})$ 都准凝聚。
使用《概形》节 \ref{schemes-section-quasi-coherent} 中关于准凝聚层的结果，
可知每个层 $R^pg_*\mathcal{F}$ 都准凝聚，因为它位于一个短正合列中间：
其左侧是准凝聚层之间一个映射的余核，右侧是准凝聚层之间一个映射的核。
故 $P(U \cup V)$ 成立，(1) 得证。

\medskip\noindent
下面证明 (3)，从而更证明 (2)。选取有限仿射开覆盖
$S = \bigcup_{j = 1, \ldots m} S_j$。对每个 $j$，选取有限仿射开覆盖
$f^{-1}(S_j) = \bigcup_{i = 1, \ldots t_j} U_{ji} $。
令
$$
d_j = \max\nolimits_{I \subset \{1, \ldots, t_j\}}
\left(|I| + t(\bigcap\nolimits_{i \in I} U_{ji})\right)
$$
为引理 \ref{coherent-lemma-vanishing-nr-affines-quasi-separated} 中所得到的整数。
我们断言 $n(X, S, f) = \max d_j$ 可用。

\medskip\noindent
事实上，设 $S' \to S$ 为概形态射，并设 $\mathcal{F}'$ 为
$X' = S' \times_S X$ 上的准凝聚层。
我们要证明 $R^pf'_*\mathcal{F}' = 0$ 对
$p \geq n(X, S, f)$ 成立。
由于该问题在 $S'$ 上是局部的，可假设 $S'$ 仿射，并映入某个
$S_j$，其中 $j$ 为相应指标。于是 $X' = S' \times_{S_j} f^{-1}(S_j)$
由仿射开集 $S' \times_{S_j} U_{ji}$（$i = 1, \ldots t_j$）覆盖，且诸交
$$
\bigcap\nolimits_{i \in I} S' \times_{S_j} U_{ji} =
S' \times_{S_j} \bigcap\nolimits_{i \in I} U_{ji}
$$
可由与基变换前相同数目的仿射开集覆盖。
应用引理 \ref{coherent-lemma-vanishing-nr-affines-quasi-separated}，得到
$H^p(X', \mathcal{F}') = 0$。
由证明的第一部分，已经知道每个 $R^qf'_*\mathcal{F}'$ 都准凝聚，
因而在仿射概形 $S'$ 上具有消失的高阶上同调群。
于是由《上同调》引理 \ref{cohomology-lemma-apply-Leray} 可得
$H^0(S', R^pf'_*\mathcal{F}') = H^p(X', \mathcal{F}') = 0$。
由于 $R^pf'_*\mathcal{F}'$ 准凝聚，结论是
$R^pf'_*\mathcal{F}' = 0$。
\end{proof}

\begin{lemma}
\label{coherent-lemma-quasi-coherence-higher-direct-images-application}
设 $f : X \to S$ 为概形态射。
假设 $f$ 拟分离且拟紧。
假设 $S$ 仿射。
对任意准凝聚 $\mathcal{O}_X$-模 $\mathcal{F}$，均有
$$
H^q(X, \mathcal{F}) = H^0(S, R^qf_*\mathcal{F})
$$
其中 $q \in \mathbf{Z}$。
\end{lemma}

\begin{proof}
考虑 Leray 谱序列 $E_2^{p, q} = H^p(S, R^qf_*\mathcal{F})$，
它收敛到 $H^{p + q}(X, \mathcal{F})$；见《上同调》引理
\ref{cohomology-lemma-Leray}。
由引理 \ref{coherent-lemma-quasi-coherence-higher-direct-images}，
层 $R^qf_*\mathcal{F}$ 准凝聚。
由引理 \ref{coherent-lemma-quasi-coherent-affine-cohomology-zero}，
有 $E_2^{p, q} = 0$，只要 $p > 0$。
因此谱序列在 $E_2$ 处退化，结论成立。
一般原理另见《上同调》引理
\ref{cohomology-lemma-apply-Leray} (2)。
\end{proof}


\section{上同调与基变换（一）}
\label{coherent-section-cohomology-and-base-change}

\noindent
设 $f : X \to S$ 为概形态射。
设 $\mathcal{F}$ 为 $X$ 上的准凝聚层。
再设 $g : S' \to S$ 为任意概形态射。记
$X' = X_{S'} = S' \times_S X$ 为 $X$ 的基变换，并记
$f' : X' \to S'$ 为 $f$ 的基变换。
还以 $g' : X' \to X$ 表示投影，并置
$\mathcal{F}' = (g')^*\mathcal{F}$。
下图表示这一情形：
\begin{equation}
\label{coherent-equation-base-change-diagram}
\vcenter{
\xymatrix{
\mathcal{F}' = (g')^*\mathcal{F} &
X' \ar[r]_{g'} \ar[d]_{f'} &
X \ar[d]^f &
\mathcal{F} \\
Rf'_*\mathcal{F}' &
S' \ar[r]^g &
S &
Rf_*\mathcal{F}
}
}
\end{equation}
下面是我们所考虑的基变换性质最简单的情形。

\begin{lemma}
\label{coherent-lemma-affine-base-change}
设 $f : X \to S$ 为概形态射。
设 $\mathcal{F}$ 为准凝聚 $\mathcal{O}_X$-模。
假设 $f$ 仿射。
此时 $f_*\mathcal{F} \cong Rf_*\mathcal{F}$ 是准凝聚层，并且对每个基变换图
(\ref{coherent-equation-base-change-diagram}) 都有
$$
g^*f_*\mathcal{F} = f'_*(g')^*\mathcal{F}.
$$
\end{lemma}

\begin{proof}
高阶直像的消失即引理 \ref{coherent-lemma-relative-affine-vanishing}。
该断言在 $S$ 和 $S'$ 上都是局部的。因此可假设
$X = \Spec(A)$、$S = \Spec(R)$、
$S' = \Spec(R')$ 且 $\mathcal{F} = \widetilde{M}$，
其中后者由某个 $A$-模 $M$ 给出。
我们用《概形》引理 \ref{schemes-lemma-widetilde-pullback}
描述 $\mathcal{F}$ 的拉回与直像。
具体地，$X' = \Spec(R' \otimes_R A)$，而
$\mathcal{F}'$ 是由 $(R' \otimes_R A) \otimes_A M$
所对应的准凝聚层。因此引理归结为 $R'$-模的等式
$$
(R' \otimes_R A) \otimes_A M = R' \otimes_R M
$$
。
\end{proof}

\noindent
在许多情形下，只需知道下面这个上同调与基变换的特殊情形。
它立即推出自节 \ref{coherent-section-cohomology-and-base-change-derived}
中的更强结果；但由于它如此重要，值得给出独立证明。

\begin{lemma}[平坦基变换]
\label{coherent-lemma-flat-base-change-cohomology}
考虑概形的笛卡尔图
$$
\xymatrix{
X' \ar[d]_{f'} \ar[r]_{g'} & X \ar[d]^f \\
S' \ar[r]^g & S
}
$$
设 $\mathcal{F}$ 为准凝聚 $\mathcal{O}_X$-模，其拉回为
$\mathcal{F}' = (g')^*\mathcal{F}$。
假设 $g$ 平坦，且 $f$ 拟紧拟分离。
对任意 $i \geq 0$，有：
\begin{enumerate}
\item 《上同调》引理 \ref{cohomology-lemma-base-change-map-flat-case}
中的基变换映射是同构
$$
g^*R^if_*\mathcal{F} \longrightarrow R^if'_*\mathcal{F}',
$$
\item 若 $S = \Spec(A)$ 且 $S' = \Spec(B)$，则
$H^i(X, \mathcal{F}) \otimes_A B = H^i(X', \mathcal{F}')$。
\end{enumerate}
\end{lemma}

\begin{proof}
在 (1) 中可以使用《上同调》引理
\ref{cohomology-lemma-base-change-map-flat-case}，因为由《态射》引理
\ref{morphisms-lemma-base-change-flat}，$g'$ 平坦。
说明这一点后，(1) 推出自 (2)。事实上，(1) 在 $S'$ 上是局部的，
故可假设 $S$ 和 $S'$ 都仿射。换言之，如 (2) 中那样，有
$S = \Spec(A)$ 和 $S' = \Spec(B)$。
于是，由于 $R^if_*\mathcal{F}$ 准凝聚
（引理 \ref{coherent-lemma-quasi-coherence-higher-direct-images}），
它是一个准凝聚 $\mathcal{O}_S$-模，该模与 $A$-模
$H^0(S, R^if_*\mathcal{F}) = H^i(X, \mathcal{F})$ 相对应
（等式来自引理
\ref{coherent-lemma-quasi-coherence-higher-direct-images-application}）。
类似地，$R^if'_*\mathcal{F}'$ 是一个准凝聚
$\mathcal{O}_{S'}$-模，该模与 $B$-模
$H^i(X', \mathcal{F}')$ 相对应。由于沿 $g$ 拉回在模上对应于
$- \otimes_A B$
（《概形》引理 \ref{schemes-lemma-widetilde-pullback}），
只需证明 (2)。

\medskip\noindent
设 $A \to B$ 为平坦环同态。
设 $X$ 为 $A$ 上的拟紧拟分离概形。
设 $\mathcal{F}$ 为准凝聚 $\mathcal{O}_X$-模。
置 $X_B = X \times_{\Spec(A)} \Spec(B)$，并以
$\mathcal{F}_B$ 表示 $\mathcal{F}$ 的拉回。
我们要证明映射
$$
H^i(X, \mathcal{F}) \otimes_A B \longrightarrow H^i(X_B, \mathcal{F}_B)
$$
（由引理陈述中的参考结果给出）是同构。

\medskip\noindent
若 $X$ 分离，选取仿射开覆盖
$\mathcal{U} : X = U_1 \cup \ldots \cup U_t$，并回忆
$$
\check{H}^p(\mathcal{U}, \mathcal{F}) = H^p(X, \mathcal{F}),
$$
见引理 \ref{coherent-lemma-cech-cohomology-quasi-coherent}。
若 $\mathcal{U}_B : X_B = (U_1)_B \cup \ldots \cup (U_t)_B$
是由基变换所得的覆盖，则每个 $(U_i)_B$ 仍仿射，且 $X_B$ 分离。
因此得到
$$
\check{H}^p(\mathcal{U}_B, \mathcal{F}_B) = H^p(X_B, \mathcal{F}_B).
$$
两个 {\v C}ech 复形之间有关系
$$
\check{\mathcal{C}}^\bullet(\mathcal{U}_B, \mathcal{F}_B) =
\check{\mathcal{C}}^\bullet(\mathcal{U}, \mathcal{F}) \otimes_A B
$$
这是引理 \ref{coherent-lemma-affine-base-change} 的推论。
由于 $A \to B$ 平坦，取上同调后同一结论仍成立。

\medskip\noindent
若 $X$ 拟分离，选取仿射开覆盖
$\mathcal{U} : X = U_1 \cup \ldots \cup U_t$。
我们将使用《上同调》引理
\ref{cohomology-lemma-cech-spectral-sequence} 的
{\v C}ech 到上同调谱序列。
希望避开这一谱序列的读者，可如引理
\ref{coherent-lemma-quasi-coherence-higher-direct-images} 的证明中那样，
使用 Mayer--Vietoris 并对 $t$ 归纳。
该谱序列的 $E_2$ 页为
$E_2^{p, q} = \check{H}^p(\mathcal{U}, \underline{H}^q(\mathcal{F}))$，
并收敛到 $H^{p + q}(X, \mathcal{F})$。
类似地，还有一个 $E_2$ 页为
$E_2^{p, q} = \check{H}^p(\mathcal{U}_B, \underline{H}^q(\mathcal{F}_B))$
的谱序列，它收敛到 $H^{p + q}(X_B, \mathcal{F}_B)$。
由于诸交 $U_{i_0 \ldots i_p}$ 拟紧且分离，
证明第二段的结果给出
$\check{H}^p(\mathcal{U}_B, \underline{H}^q(\mathcal{F}_B)) =
\check{H}^p(\mathcal{U}, \underline{H}^q(\mathcal{F})) \otimes_A B$。
利用 $A \to B$ 平坦，得到
$H^i(X, \mathcal{F}) \otimes_A B \to H^i(X_B, \mathcal{F}_B)$
对所有 $i$ 都是同构，证毕。
\end{proof}

\begin{lemma}[有限局部自由基变换]
\label{coherent-lemma-finite-locally-free-base-change-cohomology}
考虑概形的笛卡尔图
$$
\xymatrix{
Y \ar[d]_{g} \ar[r]_h & X \ar[d]^f \\
\Spec(B) \ar[r] & \Spec(A)
}
$$
设 $\mathcal{F}$ 为准凝聚 $\mathcal{O}_X$-模，其拉回为
$\mathcal{G} = h^*\mathcal{F}$。
若 $B$ 是有限局部自由 $A$-模，则
$H^i(X, \mathcal{F}) \otimes_A B = H^i(Y, \mathcal{G})$。
\end{lemma}

\noindent
{\bf 警告}：除非理解本引理与引理
\ref{coherent-lemma-flat-base-change-cohomology} 的差别，否则不要使用本引理。

\begin{proof}
若 $X$ 分离，选取仿射开覆盖
$\mathcal{U} : X = \bigcup_{i \in I} U_i$，并回忆
$$
\check{H}^p(\mathcal{U}, \mathcal{F}) = H^p(X, \mathcal{F}),
$$
见引理 \ref{coherent-lemma-cech-cohomology-quasi-coherent}。
令 $\mathcal{V} : Y = \bigcup_{i \in I} g^{-1}(U_i)$
为 $Y$ 的相应仿射开覆盖。
开集 $V_i = g^{-1}(U_i) = U_i \times_{\Spec(A)} \Spec(B)$
均仿射，且 $Y$ 分离。因此得到
$$
\check{H}^p(\mathcal{V}, \mathcal{G}) = H^p(Y, \mathcal{G}).
$$
我们断言 {\v C}ech 复形的映射
$$
\check{\mathcal{C}}^\bullet(\mathcal{U}, \mathcal{F}) \otimes_A B
\longrightarrow
\check{\mathcal{C}}^\bullet(\mathcal{V}, \mathcal{G})
$$
是同构。事实上，由于 $B$ 作为 $A$-模有限表示，
和 $B$ 张量在 $A$ 上与乘积可交换；见《代数》命题
\ref{algebra-proposition-fp-tensor}。
因此只需证明映射
$\Gamma(U_{i_0 \ldots i_p}, \mathcal{F}) \otimes_A B \to
\Gamma(V_{i_0 \ldots i_p}, \mathcal{G})$
是同构，而这来自引理 \ref{coherent-lemma-affine-base-change}。
由于 $A \to B$ 平坦，取上同调后同一结论仍成立。

\medskip\noindent
在一般情形下，论证完全相同：使用仿射开覆盖
$\mathcal{U} : X = \bigcup_{i \in I} U_i$ 及其相应覆盖
$\mathcal{V} : Y = \bigcup_{i \in I} V_i$，
其中 $V_i = g^{-1}(U_i)$ 如上。
我们使用《上同调》引理
\ref{cohomology-lemma-cech-spectral-sequence} 的
{\v C}ech 到上同调谱序列。
该谱序列的 $E_2$ 页为
$E_2^{p, q} = \check{H}^p(\mathcal{U}, \underline{H}^q(\mathcal{F}))$，
并收敛到 $H^{p + q}(X, \mathcal{F})$。
类似地，还有一个 $E_2$ 页为
$E_2^{p, q} = \check{H}^p(\mathcal{V}, \underline{H}^q(\mathcal{G}))$
的谱序列，它收敛到 $H^{p + q}(Y, \mathcal{G})$。
由于诸交 $U_{i_0 \ldots i_p}$ 分离，上一段的结果给出同构
$\Gamma(U_{i_0 \ldots i_p}, \underline{H}^q(\mathcal{F})) \otimes_A B
\to \Gamma(V_{i_0 \ldots i_p}, \underline{H}^q(\mathcal{G}))$。
利用 $- \otimes_A B$ 与乘积可交换且正合，得到
$\check{H}^p(\mathcal{U}, \underline{H}^q(\mathcal{F})) \otimes_A B
\to \check{H}^p(\mathcal{V}, \underline{H}^q(\mathcal{G}))$
是同构。再利用 $A \to B$ 平坦，得到
$H^i(X, \mathcal{F}) \otimes_A B \to H^i(Y, \mathcal{G})$
对所有 $i$ 都是同构，证毕。
\end{proof}


\section{余极限与高阶直像}
\label{coherent-section-colimits}

\noindent
这一性质的一般性结果见《上同调》节
\ref{cohomology-section-limits}、《层》引理
\ref{sheaves-lemma-directed-colimits-sections} 以及《模》引理
\ref{modules-lemma-finite-presentation-quasi-compact-colimit}。

\begin{lemma}
\label{coherent-lemma-colimit-cohomology}
设 $f : X \to S$ 为概形之间的拟紧拟分离态射。
设 $\mathcal{F} = \colim \mathcal{F}_i$ 为 $X$ 上阿贝尔层的滤过余极限。
则对任意 $p \geq 0$，有
$$
R^pf_*\mathcal{F} = \colim R^pf_*\mathcal{F}_i.
$$
\end{lemma}

\begin{proof}
应用《上同调》引理
\ref{cohomology-lemma-higher-direct-image-colimit}。
由于仿射开集构成 $S$ 的拓扑基，只需证明：对仿射开集
$U \subset S$，有
$H^p(f^{-1}U, \mathcal{F}) = \colim H^p(f^{-1}U, \mathcal{F}_i)$。
因为 $f^{-1}U$ 拟紧拟分离，结论来自《上同调》引理
\ref{cohomology-lemma-quasi-separated-cohomology-colimit}。
（这是因为 $f^{-1}U$ 中仿射开集所成的基满足该引理的假设。）
\end{proof}


\section{上同调与基变换（二）}
\label{coherent-section-cohomology-and-base-change-derived}

\noindent
设 $f : X \to S$ 为概形态射，并设 $\mathcal{F}$
为准凝聚 $\mathcal{O}_X$-模。若 $f$ 拟紧拟分离，我们希望将
$Rf_*\mathcal{F}$ 表示为 $S$ 上的准凝聚层复形。这由如下事实推出：
层 $R^if_*\mathcal{F}$ 在 $S$ 拟紧且具有仿射对角态射时准凝聚；
再利用 $D_\QCoh(S)$ 与 $D(\QCoh(\mathcal{O}_S))$ 的等价，
见《概形的导出范畴》命题
\ref{perfect-proposition-quasi-compact-affine-diagonal}。

\medskip\noindent
本节采用另一种方法，构造一个具有良好基变换性质的显式复形。
当 $f$ 与 $S$ 分离，或者更一般地具有仿射对角态射时，
这一构造尤其简单。由于这是迄今最常用的情形，我们将其单独处理。

\begin{lemma}
\label{coherent-lemma-separated-case-relative-cech}
设 $f : X \to S$ 为概形态射。
设 $\mathcal{F}$ 为准凝聚 $\mathcal{O}_X$-模。
假设 $X$ 拟紧，且 $X$ 与 $S$ 都具有仿射对角态射
（例如，当 $X$ 与 $S$ 都分离时）。
在此情形下，可按如下方式计算 $Rf_*\mathcal{F}$：
\begin{enumerate}
\item 选取有限仿射开覆盖
$\mathcal{U} : X = \bigcup_{i = 1, \ldots, n} U_i$。
\item 对 $i_0, \ldots, i_p \in \{1, \ldots, n\}$，记
$f_{i_0 \ldots i_p} : U_{i_0 \ldots i_p} \to S$ 为 $f$
在交集 $U_{i_0 \ldots i_p} = U_{i_0} \cap \ldots \cap U_{i_p}$
上的限制。
\item 令 $\mathcal{F}_{i_0 \ldots i_p}$ 等于 $\mathcal{F}$
在 $U_{i_0 \ldots i_p}$ 上的限制。
\item 置
$$
\check{\mathcal{C}}^p(\mathcal{U}, f, \mathcal{F}) =
\bigoplus\nolimits_{i_0 \ldots i_p}
f_{i_0 \ldots i_p *} \mathcal{F}_{i_0 \ldots i_p}
$$
并按《上同调》公式 (\ref{cohomology-equation-d-cech}) 定义微分
$d : \check{\mathcal{C}}^p(\mathcal{U}, f, \mathcal{F})
\to \check{\mathcal{C}}^{p + 1}(\mathcal{U}, f, \mathcal{F})$。
\end{enumerate}
则复形 $\check{\mathcal{C}}^\bullet(\mathcal{U}, f, \mathcal{F})$
是 $S$ 上的准凝聚层复形，并带有同构
$$
\check{\mathcal{C}}^\bullet(\mathcal{U}, f, \mathcal{F})
\longrightarrow
Rf_*\mathcal{F}
$$
它位于 $D^{+}(S)$ 中。该同构关于准凝聚层 $\mathcal{F}$ 具有函子性。
\end{lemma}

\begin{proof}
考虑解消
$\mathcal{F} \to {\mathfrak C}^\bullet(\mathcal{U}, \mathcal{F})$，
即《上同调》引理 \ref{cohomology-lemma-covering-resolution} 中的解消。
我们有复形等式
$\check{\mathcal{C}}^\bullet(\mathcal{U}, f, \mathcal{F}) =
f_*{\mathfrak C}^\bullet(\mathcal{U}, \mathcal{F})$；
二者都是准凝聚 $\mathcal{O}_S$-模的复形。
态射 $j_{i_0 \ldots i_p} : U_{i_0 \ldots i_p} \to X$
与态射 $f_{i_0 \ldots i_p} : U_{i_0 \ldots i_p} \to S$
由《态射》引理 \ref{morphisms-lemma-affine-permanence}
及引理 \ref{coherent-lemma-affine-diagonal} 都是仿射的。
故 $R^qj_{i_0 \ldots i_p *}\mathcal{F}_{i_0 \ldots i_p}$
以及 $R^qf_{i_0 \ldots i_p *}\mathcal{F}_{i_0 \ldots i_p}$
在 $q > 0$ 时均为零（引理 \ref{coherent-lemma-relative-affine-vanishing}）。
利用 $f \circ j_{i_0 \ldots i_p} = f_{i_0 \ldots i_p}$ 及
《上同调》引理 \ref{cohomology-lemma-relative-Leray} 中的谱序列，
可知
$R^qf_*(j_{i_0 \ldots i_p *}\mathcal{F}_{i_0 \ldots i_p}) = 0$
对 $q > 0$ 成立。
由于复形 ${\mathfrak C}^\bullet(\mathcal{U}, \mathcal{F})$ 的各项
是层 $j_{i_0 \ldots i_p *}\mathcal{F}_{i_0 \ldots i_p}$ 的有限直和，
由 Leray 无环性引理
（《导出范畴》引理 \ref{derived-lemma-leray-acyclicity}）得到
$$
Rf_* \mathcal{F} = f_*{\mathfrak C}^\bullet(\mathcal{U}, \mathcal{F}) =
\check{\mathcal{C}}^\bullet(\mathcal{U}, f, \mathcal{F})
$$
如所欲。
\end{proof}

\noindent
下面考察进行基变换时会发生什么。

\begin{lemma}
\label{coherent-lemma-base-change-complex}
采用图 (\ref{coherent-equation-base-change-diagram}) 中的记号。
假设 $f : X \to S$ 与 $\mathcal{F}$ 满足引理
\ref{coherent-lemma-separated-case-relative-cech} 的假设。选取有限仿射开覆盖
$\mathcal{U} : X = \bigcup U_i$ 作为 $X$ 的覆盖。
存在典范同构
$$
g^*\check{\mathcal{C}}^\bullet(\mathcal{U}, f, \mathcal{F})
\longrightarrow
Rf'_*\mathcal{F}'
$$
它位于 $D^{+}(S')$ 中。此外，若 $S' \to S$ 仿射，则实际上
$$
g^*\check{\mathcal{C}}^\bullet(\mathcal{U}, f, \mathcal{F})
=
\check{\mathcal{C}}^\bullet(\mathcal{U}', f', \mathcal{F}')
$$
其中 $\mathcal{U}' : X' = \bigcup U_i'$，而
$U_i' = (g')^{-1}(U_i) = U_{i, S'}$ 也仿射。
\end{lemma}

\begin{proof}
可定义 $U_i' = (g')^{-1}(U_i) = U_{i, S'}$，
而不论 $S'$ 在 $S$ 上是否仿射。
令 $\mathcal{U}' : X' = \bigcup U_i'$ 为由此得到的 $X'$ 的覆盖。
在此情形下，我们断言
$$
g^*\check{\mathcal{C}}^\bullet(\mathcal{U}, f, \mathcal{F})
=
\check{\mathcal{C}}^\bullet(\mathcal{U}', f', \mathcal{F}')
$$
其中 $\check{\mathcal{C}}^\bullet(\mathcal{U}', f', \mathcal{F}')$
的定义与引理 \ref{coherent-lemma-separated-case-relative-cech} 中完全相同。
这可通过在 $S'$ 上局部工作，由仿射态射的情形
（引理 \ref{coherent-lemma-affine-base-change}）立即得到。
此外，完全如引理 \ref{coherent-lemma-separated-case-relative-cech} 的证明，
可见存在同构
$$
\check{\mathcal{C}}^\bullet(\mathcal{U}', f', \mathcal{F}')
\longrightarrow
Rf'_*\mathcal{F}'
$$
它位于 $D^{+}(S')$ 中，因为态射 $U_i' \to X'$ 与 $U_i' \to S'$
仍然仿射（它们是仿射态射的基变换）。细节从略。
\end{proof}

\noindent
上述引理说明，复形
$$
\mathcal{K}^\bullet = \check{\mathcal{C}}^\bullet(\mathcal{U}, f, \mathcal{F})
$$
是 $S$ 上准凝聚层的下有界复形，它\emph{泛地}计算
$f : X \to S$ 的高阶直像。
这是该特定复形所具有的性质；一般而言，将
$\check{\mathcal{C}}^\bullet(\mathcal{U}, f, \mathcal{F})$
换成拟同构复形并不能保持这一性质！换言之，
这不是一个在导出范畴中有意义的陈述。
原因在于，拉回 $g^*\mathcal{K}^\bullet$ 一般\emph{不}等于
$Lg^*\mathcal{K}^\bullet$，后者是 $\mathcal{K}^\bullet$ 的导出拉回！

\medskip\noindent
下面给出可以证明该陈述的一个更一般情形。
我们指出，在大多数应用中，$S$ 分离这一条件无伤大雅，
因为这里通常是用它证明全导出像的某个局部性质。
证明要复杂得多，并使用超覆盖；这也是有时如何使用超覆盖的一个好例子。

\begin{lemma}
\label{coherent-lemma-hypercoverings}
设 $f : X \to S$ 为概形态射。
设 $\mathcal{F}$ 为 $X$ 上的准凝聚层。
假设 $f$ 拟紧拟分离，且 $S$ 拟紧并分离。
则存在下有界复形 $\mathcal{K}^\bullet$，
它由准凝聚 $\mathcal{O}_S$-模组成，并具有如下性质：对每个态射
$g : S' \to S$，复形 $g^*\mathcal{K}^\bullet$ 是
$Rf'_*\mathcal{F}'$ 的一个代表元，其中记号如图
(\ref{coherent-equation-base-change-diagram})。
\end{lemma}

\begin{proof}
（若 $f$ 也分离，请见引理 \ref{coherent-lemma-base-change-complex}。）
这些假设特别推出 $X$ 作为概形拟紧拟分离。
令 $\mathcal{B}$ 为 $X$ 的仿射开集所成的集合。由《超覆盖》引理
\ref{hypercovering-lemma-quasi-separated-quasi-compact-hypercovering}，
可找到超覆盖 $K = (I, \{U_i\})$，使每个 $I_n$ 有限，
且每个 $U_i$ 都是 $X$ 的仿射开集。由《超覆盖》引理
\ref{hypercovering-lemma-cech-spectral-sequence}，
存在一个 $E_2$ 页为
$$
E_2^{p, q} = \check{H}^p(K, \underline{H}^q(\mathcal{F}))
$$
且收敛到 $H^{p + q}(X, \mathcal{F})$ 的谱序列。注意，
$\check{H}^p(K, \underline{H}^q(\mathcal{F}))$ 是下列复形的第 $p$ 个上同调群：
$$
\prod\nolimits_{i \in I_0} H^q(U_i, \mathcal{F})
\to
\prod\nolimits_{i \in I_1} H^q(U_i, \mathcal{F})
\to
\prod\nolimits_{i \in I_2} H^q(U_i, \mathcal{F})
\to \ldots
$$
由于每个 $U_i$ 仿射，可见除非 $q = 0$，否则它为零；
而在后一情形得到
$$
\prod\nolimits_{i \in I_0} \mathcal{F}(U_i)
\to
\prod\nolimits_{i \in I_1} \mathcal{F}(U_i)
\to
\prod\nolimits_{i \in I_2} \mathcal{F}(U_i)
\to \ldots
$$
因此可知 $R\Gamma(X, \mathcal{F})$ 由这个复形计算。

\medskip\noindent
对任意 $n$ 与 $i \in I_n$，记 $f_i : U_i \to S$ 为 $f$
在 $U_i$ 上的限制。由于 $S$ 分离且 $U_i$ 仿射，该态射仿射。
考虑 $S$ 上的准凝聚层复形
$$
\mathcal{K}^\bullet = (
\prod\nolimits_{i \in I_0} f_{i, *}\mathcal{F}|_{U_i}
\to
\prod\nolimits_{i \in I_1} f_{i, *}\mathcal{F}|_{U_i}
\to
\prod\nolimits_{i \in I_2} f_{i, *}\mathcal{F}|_{U_i}
\to \ldots )
$$
与《超覆盖》引理 \ref{hypercovering-lemma-cech-spectral-sequence}
中一样，我们得到映射 $\mathcal{K}^\bullet \to Rf_*\mathcal{F}$，
它位于 $D(\mathcal{O}_S)$ 中；其构造是选取 $\mathcal{F}$ 的单射解消
（细节从略）。
考虑任意仿射概形 $V$ 及态射 $g : V \to S$。
则基变换 $X_V$ 具有由基变换得到的超覆盖
$K_V = (I, \{U_{i, V}\})$。此外，
$g^*f_{i, *}\mathcal{F} = f_{i, V, *}(g')^*\mathcal{F}|_{U_{i, V}}$。
因此上述论证证明，$\Gamma(V, g^*\mathcal{K}^\bullet)$
计算 $R\Gamma(X_V, (g')^*\mathcal{F})$。
要证明复形相等，只需在 $S'$ 上 Zariski 局部证明，故引理得证。
\end{proof}

\noindent
下面的引理是平坦基变换的一种变式。

\begin{lemma}
\label{coherent-lemma-base-change-tensor-with-flat}
考虑概形的笛卡尔图
$$
\xymatrix{
X' \ar[d]_{f'} \ar[r]_{g'} & X \ar[d]^f \\
S' \ar[r]^g & S
}
$$
设 $\mathcal{F}$ 为准凝聚 $\mathcal{O}_X$-模。
设 $\mathcal{G}$ 为准凝聚 $\mathcal{O}_{S'}$-模，
并且在 $S$ 上平坦。假设 $f$ 拟紧拟分离。
对任意 $i \geq 0$，有等同
$$
\mathcal{G} \otimes_{\mathcal{O}_{S'}} g^*R^if_*\mathcal{F} =
R^if'_*\left((f')^*\mathcal{G}
\otimes_{\mathcal{O}_{X'}} (g')^*\mathcal{F}\right)
$$
\end{lemma}

\begin{proof}
先构造从左到右的映射。首先，有基变换映射
$Lg^*Rf_*\mathcal{F} \to Rf'_*L(g')^*\mathcal{F}$。另有伴随映射
$\mathcal{G} \to Rf'_*L(f')^*\mathcal{G}$。利用相对杯积，得到
\begin{align*}
\mathcal{G}
\otimes_{\mathcal{O}_{S'}}^\mathbf{L}
Lg^*Rf_*\mathcal{F}
& \to
Rf'_*L(f')^*\mathcal{G}
\otimes_{\mathcal{O}_{S'}}^\mathbf{L}
Rf'_*L(g')^*\mathcal{F} \\
& \to
Rf'_*\left(L(f')^*\mathcal{G}
\otimes_{\mathcal{O}_{X'}}^\mathbf{L}
L(g')^*\mathcal{F}\right) \\
& \to
Rf'_*\left((f')^*\mathcal{G}
\otimes_{\mathcal{O}_{X'}} (g')^*\mathcal{F}\right)
\end{align*}
其中中间的箭头使用了相对杯积，见《上同调》备注
\ref{cohomology-remark-cup-product}。
该复合的源为
$$
\mathcal{G} \otimes_{\mathcal{O}_{S'}}^\mathbf{L} Lg^*Rf_*\mathcal{F} =
\mathcal{G} \otimes_{g^{-1}\mathcal{O}_S}^\mathbf{L} g^{-1}Rf_*\mathcal{F}
$$
这是《上同调》引理 \ref{cohomology-lemma-variant-derived-pullback}。
由于 $\mathcal{G}$ 作为 $g^{-1}\mathcal{O}_S$-模层平坦，
且 $g^{-1}$ 是正合函子，这是一个复形，其第 $i$ 个上同调层为
$\mathcal{G} \otimes_{g^{-1}\mathcal{O}_S} g^{-1}R^if_*\mathcal{F} =
\mathcal{G} \otimes_{\mathcal{O}_{S'}} g^*R^if_*\mathcal{F}$。
这样便得到引理等式中从左到右的全局映射。
为证明此映射是同构，可在 $S'$ 上局部工作。
因此可以并且确实假设 $S$ 与 $S'$ 都是仿射概形。

\medskip\noindent
现在证明 $S$ 与 $S'$ 仿射的情形。设 $S = \Spec(A)$ 且 $S' = \Spec(B)$，
并设 $\mathcal{G}$ 对应于 $B$-模 $N$，后者按假设在 $A$ 上平坦。
由于 $S$ 仿射，$X$ 拟紧拟分离。
我们用超覆盖论证完成证明；若 $X$ 分离或具有仿射对角态射，
则可使用一个 {\v C}ech 覆盖。
令 $\mathcal{B}$ 为 $X$ 的仿射开集所成的集合。由《超覆盖》引理
\ref{hypercovering-lemma-quasi-separated-quasi-compact-hypercovering}，
可找到超覆盖 $K = (I, \{U_i\})$，它是 $X$ 的超覆盖，且每个
$I_n$ 有限、每个 $U_i$ 都是 $X$ 的仿射开集。由《超覆盖》引理
\ref{hypercovering-lemma-cech-spectral-sequence}，
存在一个 $E_2$ 页为
$$
E_2^{p, q} = \check{H}^p(K, \underline{H}^q(\mathcal{F}))
$$
且收敛到 $H^{p + q}(X, \mathcal{F})$ 的谱序列。由于每个 $U_i$ 仿射，
且 $\mathcal{F}$ 准凝聚，$\underline{H}^q(\mathcal{F})$
在 $U_i$ 上的值对 $q > 0$ 为零。因此谱序列退化，
从而上同调模 $H^q(X, \mathcal{F})$ 由下列复形计算：
$$
\prod\nolimits_{i \in I_0} \mathcal{F}(U_i)
\to
\prod\nolimits_{i \in I_1} \mathcal{F}(U_i)
\to
\prod\nolimits_{i \in I_2} \mathcal{F}(U_i)
\to \ldots
$$
接着注意，我们的超覆盖在 $S'$ 上的基变换是
$X' = S' \times_S X$ 的一个超覆盖。
概形 $S' \times_S U_i$ 也都仿射，并且有
$$
\left((f')^*\mathcal{G} \otimes_{\mathcal{O}_{S'}} (g')^*\mathcal{F}\right)
(S' \times_S U_i) =
N \otimes_A \mathcal{F}(U_i)
$$
由此可知，上同调模
$H^q(X', (f')^*\mathcal{G} \otimes_{\mathcal{O}_{S'}} (g')^*\mathcal{F})$
由下列复形计算：
$$
N \otimes_A \left(
\prod\nolimits_{i \in I_0} \mathcal{F}(U_i)
\to
\prod\nolimits_{i \in I_1} \mathcal{F}(U_i)
\to
\prod\nolimits_{i \in I_2} \mathcal{F}(U_i)
\to \ldots
\right)
$$
由于 $N$ 在 $A$ 上平坦，可知
$$
H^q(X', (f')^*\mathcal{G} \otimes_{\mathcal{O}_{S'}} (g')^*\mathcal{F})
=
N \otimes_A H^q(X, \mathcal{F})
$$
这正是待证陈述的代数表述，故证明完成。
\end{proof}







\section{射影空间的上同调}
\label{coherent-section-cohomology-projective-space}

\noindent
本节计算 $\mathbf{P}^n_S$ 的结构层各扭转的上同调，其中基概形为 $S$。
回忆，$\mathbf{P}^n_S$ 定义为纤维积
$
\mathbf{P}^n_S = S \times_{\Spec(\mathbf{Z})} \mathbf{P}^n_{\mathbf{Z}}
$
见《构造》定义 \ref{constructions-definition-projective-space}。
《构造》引理 \ref{constructions-lemma-projective-space-bundle} 已证明它等于
$$
\mathbf{P}^n_S = \underline{\text{Proj}}_S(\mathcal{O}_S[T_0, \ldots, T_n])
$$
特别地，射影空间是射影丛的一个特殊情形。
若 $S = \Spec(R)$ 仿射，则有
$$
\mathbf{P}^n_S = \mathbf{P}^n_R = \text{Proj}(R[T_0, \ldots, T_n]).
$$
所有这些等同彼此相容，也与扭转结构层
$\mathcal{O}_{\mathbf{P}^n_S}(d)$ 的构造相容。

\medskip\noindent
在陈述结果前，需要引入一些记号。
设 $R$ 为环。
回忆，$R[T_0, \ldots, T_n]$ 是分次 $R$-代数，
其中每个 $T_i$ 都是次数为 $1$ 的齐次元。
以 $(R[T_0, \ldots, T_n])_d$ 表示次数为 $d$ 的分量。
它是有限自由 $R$-模，秩为 $\binom{n + d}{d}$，
当 $d \geq 0$；否则为零。它的一组基由单项式
$T_0^{e_0} \ldots T_n^{e_n}$ 组成，其中 $\sum e_i = d$。
还要使用如下记号：
$R[\frac{1}{T_0}, \ldots, \frac{1}{T_n}]$ 表示 $\mathbf{Z}$-分次环，
其中 $\frac{1}{T_i}$ 的次数为 $-1$。特别地，下述陈述中出现的
$\mathbf{Z}$-分次 $R[\frac{1}{T_0}, \ldots, \frac{1}{T_n}]$-模
$$
\frac{1}{T_0 \ldots T_n} R[\frac{1}{T_0}, \ldots, \frac{1}{T_n}]
$$
在次数 $\geq -n$ 时为零，以生成元 $\frac{1}{T_0 \ldots T_n}$
在次数 $-n - 1$ 时自由生成，并且是秩为
$(-1)^n\binom{n + d}{d}$ 的自由模，当 $d \leq -n - 1$。

\begin{lemma}
\label{coherent-lemma-cohomology-projective-space-over-ring}
\begin{reference}
\cite[III Proposition 2.1.12]{EGA}
\end{reference}
设 $R$ 为环。
设 $n \geq 0$ 为整数。
则
$$
H^q(\mathbf{P}^n, \mathcal{O}_{\mathbf{P}^n_R}(d)) =
\left\{
\begin{matrix}
(R[T_0, \ldots, T_n])_d & \text{若} & q = 0,\ d \geq 0 \\
\left(\frac{1}{T_0 \ldots T_n} R[\frac{1}{T_0}, \ldots, \frac{1}{T_n}]\right)_d
& \text{若} & q = n,\ d < 0 \\
0 & \text{若} & q, n, d\text{ not as above}
\end{matrix}
\right.
$$
这是 $R$-模的等式。
\end{lemma}

\begin{proof}
使用标准仿射开覆盖
$$
\mathcal{U} : \mathbf{P}^n_R = \bigcup\nolimits_{i = 0}^n D_{+}(T_i)
$$
并利用 {\v C}ech 复形计算上同调。
这是允许的，因为由引理 \ref{coherent-lemma-cech-cohomology-quasi-coherent}，
有限多个仿射开集 $D_{+}(T_i)$ 的任意交仍是标准仿射开集
（见《构造》节 \ref{constructions-section-proj}）。
事实上，根据《上同调》引理 \ref{cohomology-lemma-ordered-alternating}
及 \ref{cohomology-lemma-alternating-usual}，
可使用交错或有序 {\v C}ech 复形。

\medskip\noindent
在 $\{0, \ldots, n\}$ 上采用通常次序。
于是所考察复形的各项为
$$
\check{\mathcal{C}}_{ord}^p(\mathcal{U}, \mathcal{O}_{\mathbf{P}_R}(d))
=
\bigoplus\nolimits_{i_0 < \ldots < i_p}
(R[T_0, \ldots, T_n, \frac{1}{T_{i_0} \ldots T_{i_p}}])_d
$$
此外，映射由通常公式给出：
$$
d(s)_{i_0 \ldots i_{p + 1}} =
\sum\nolimits_{j = 0}^{p + 1} (-1)^j s_{i_0 \ldots \hat i_j \ldots i_{p + 1}}
$$
见《上同调》节 \ref{cohomology-section-alternating-cech}。
注意，此复形的每一项都有自然的 $\mathbf{Z}^{n + 1}$-分次。
具体而言，规定单项式 $T_0^{e_0} \ldots T_n^{e_n}$ 为权
$(e_0, \ldots, e_n) \in \mathbf{Z}^{n + 1}$ 的齐次元。
显然，上述微分保持该分次。用公式表示即
$$
\check{\mathcal{C}}_{ord}^\bullet(\mathcal{U}, \mathcal{O}_{\mathbf{P}_R}(d))
=
\bigoplus\nolimits_{\vec{e} \in \mathbf{Z}^{n + 1}}
\check{\mathcal{C}}^\bullet(\vec{e})
$$
其中右侧并非所有直和项都会出现（见下文）。
因此，为计算该复形的上同调模，只需计算各分次部分的上同调，
最后取直和。

\medskip\noindent
固定 $\vec{e} = (e_0, \ldots, e_n) \in \mathbf{Z}^{n + 1}$。
要使这个权出现在上述复形中，需要假设
$e_0 + \ldots + e_n = d$（否则它当然会出现在结构层的另一扭转中）。
在此假设下，置
$$
NEG(\vec{e}) = \{i \in \{0, \ldots, n\} \mid e_i < 0\}.
$$
采用此记号，上述 {\v C}ech 复形的权 $\vec{e}$ 直和项
$\check{\mathcal{C}}^\bullet(\vec{e})$ 的各项为
$$
\check{\mathcal{C}}^p(\vec{e})
=
\bigoplus\nolimits_{i_0 < \ldots < i_p,
\ NEG(\vec{e}) \subset \{i_0, \ldots, i_p\}}
R \cdot T_0^{e_0} \ldots T_n^{e_n}
$$
换言之，对满足 $i_0 < \ldots < i_p$ 但
$NEG(\vec{e})$ 不包含于 $\{i_0 \ldots i_p\}$ 的指标，
相应项为零。复形 $\check{\mathcal{C}}^\bullet(\vec{e})$ 的微分
仍由上述完全相同的公式给出。

\medskip\noindent
假设 $NEG(\vec{e}) = \{0, \ldots, n\}$，即所有指数
$e_i$ 都为负。
此时复形 $\check{\mathcal{C}}^\bullet(\vec{e})$ 只有一项，即
$\check{\mathcal{C}}^n(\vec{e}) =
R \cdot \frac{1}{T_0^{-e_0} \ldots T_n^{-e_n}}$。因此在此情形下
$$
H^q(\check{\mathcal{C}}^\bullet(\vec{e})) =
\left\{
\begin{matrix}
R \cdot \frac{1}{T_0^{-e_0} \ldots T_n^{-e_n}} & \text{若} & q = n \\
0 & \text{若} & \text{否则}
\end{matrix}
\right.
$$
所有这些项的直和显然给出
$$
\left(\frac{1}{T_0 \ldots T_n} R[\frac{1}{T_0}, \ldots, \frac{1}{T_n}]\right)_d
$$
在次数 $n$ 的值，如引理所述。此外，这些项对其他次数的上同调
没有贡献（也与引理的陈述一致）。

\medskip\noindent
假设 $NEG(\vec{e}) = \emptyset$。此时复形
$\check{\mathcal{C}}^\bullet(\vec{e})$ 有一个直和项 $R$，
对应于所有 $i_0 < \ldots < i_p$。
将复形 $\check{\mathcal{C}}^\bullet(\vec{e})$ 与另一个复形比较。
具体地，考虑仿射开覆盖
$$
\mathcal{V} : \Spec(R) = \bigcup\nolimits_{i \in \{0, \ldots, n\}} V_i
$$
其中 $V_i = \Spec(R)$ 对所有 $i$ 都成立。考虑交错 {\v C}ech 复形
$$
\check{\mathcal{C}}_{ord}^\bullet(\mathcal{V}, \mathcal{O}_{\Spec(R)})
$$
由与上面相同的论证，它计算 $\Spec(R)$ 上结构层的上同调。
因此可见\par\noindent
$H^p(
\check{\mathcal{C}}_{ord}^\bullet(\mathcal{V}, \mathcal{O}_{\Spec(R)})
) = R$\hfill\break
当 $p = 0$，而为 $0$ 当 $p > 0$。\par
这些事实见引理 \ref{coherent-lemma-cech-cohomology-quasi-coherent-trivial} 及其证明。
还要注意，
$\check{\mathcal{C}}_{ord}^\bullet(\mathcal{V}, \mathcal{O}_{\Spec(R)})$
也有一个直和项 $R$ 对应于每个 $i_0 < \ldots < i_p$，
并且其微分与 $\check{\mathcal{C}}^\bullet(\vec{e})$ 完全相同。
换言之，这些复形彼此同构，因而有相同的上同调。由此得到
$$
H^q(\check{\mathcal{C}}^\bullet(\vec{e})) =
\left\{
\begin{matrix}
R \cdot T_0^{e_0} \ldots T_n^{e_n} & \text{若} & q = 0 \\
0 & \text{若} & \text{否则}
\end{matrix}
\right.
$$
这里假设 $NEG(\vec{e}) = \emptyset$。
所有这些项的直和显然给出
$$
(R[T_0, \ldots, T_n])_d
$$
在次数 $0$ 的值，如引理所述。此外，这些项对其他次数的上同调
没有贡献（也与引理的陈述一致）。

\medskip\noindent
为完成引理的证明，还须说明：当
$\check{\mathcal{C}}^\bullet(\vec{e})$ 满足
$NEG(\vec{e})$ 既非空也不等于 $\{0, \ldots, n\}$ 时，它是无环的。
选取指标 $i_{\text{fix}} \not \in NEG(\vec{e})$（这样的指标存在）。
考虑映射
$$
h :
\check{\mathcal{C}}^{p + 1}(\vec{e})
\to
\check{\mathcal{C}}^p(\vec{e})
$$
其规则为：对 $i_0 < \ldots < i_p$，有
$$
h(s)_{i_0 \ldots i_p} =
\left\{
\begin{matrix}
0 & \text{若} & p \not \in \{0, \ldots, n - 1\} \\
0 & \text{若} & i_{\text{fix}} \in \{i_0, \ldots, i_p\} \\
s_{i_{\text{fix}} i_0 \ldots i_p} & \text{若} & i_{\text{fix}} < i_0 \\
(-1)^a s_{i_0 \ldots i_{a - 1} i_{\text{fix}} i_a \ldots i_p} &
\text{若} & i_{a - 1} < i_{\text{fix}} < i_a \\
(-1)^{p + 1} s_{i_0 \ldots i_p i_{\text{fix}}} &
\text{若} & i_p < i_{\text{fix}}
\end{matrix}
\right.
$$
请与引理 \ref{coherent-lemma-cech-cohomology-quasi-coherent-trivial} 的证明比较。
该定义有意义，因为
$$
NEG(\vec{e}) \subset \{i_0, \ldots, i_p\}
\Leftrightarrow
NEG(\vec{e}) \subset \{i_{\text{fix}}, i_0, \ldots, i_p\}
$$
与该引理证明中完全相同的组合计算\footnote{
例如，假设 $i_0 < \ldots < i_p$ 满足
$i_{\text{fix}} \not \in \{i_0, \ldots, i_p\}$，并且有
$i_{a - 1} < i_{\text{fix}} < i_a$，其中 $1 \leq a \leq p$。则
\begin{align*}
&
(dh + hd)(s)_{i_0 \ldots i_p} \\
& =
\sum\nolimits_{j = 0}^p
(-1)^j
h(s)_{i_0 \ldots \hat i_j \ldots i_p}
+
(-1)^a d(s)_{i_0 \ldots i_{a - 1} i_{\text{fix}} i_a \ldots i_p}\\
& =
\sum\nolimits_{j = 0}^{a - 1}
(-1)^{j + a - 1}
s_{i_0 \ldots \hat i_j \ldots i_{a - 1} i_{\text{fix}} i_a \ldots i_p} +
\sum\nolimits_{j = a}^p
(-1)^{j + a}
s_{i_0 \ldots i_{a - 1} i_{\text{fix}} i_a \ldots \hat i_j \ldots i_p}
+ \\
&
\sum\nolimits_{j = 0}^{a - 1}
(-1)^{a + j}
s_{i_0 \ldots \hat i_j \ldots i_{a - 1} i_{\text{fix}} i_a \ldots i_p} +
(-1)^{2a} s_{i_0 \ldots i_p} +
\sum\nolimits_{j = a}^p
(-1)^{a + j + 1}
s_{i_0 \ldots i_{a - 1} i_{\text{fix}} i_a \ldots \hat i_j \ldots i_p}
\\
& =
s_{i_0 \ldots i_p}
\end{align*}
如所欲。其他情形类似。}
与引理 \ref{coherent-lemma-cech-cohomology-quasi-coherent-trivial} 的证明一样，这说明
$$
(hd + dh)(s)_{i_0 \ldots i_p}
=
s_{i_0 \ldots i_p}
$$
因此复形 $\check{\mathcal{C}}^\bullet(\vec{e})$ 的恒等映射
零伦，从而该复形无环。
\end{proof}

\noindent
在下面的引理中，将使用自由 $R$-模之间的配对
$$
R[T_0, \ldots, T_n]
\times
\frac{1}{T_0 \ldots T_n} R[\frac{1}{T_0}, \ldots, \frac{1}{T_n}]
\longrightarrow
R
$$
它由如下规则定义：
$$
(f, g)
\longmapsto
\text{coefficient of }
\frac{1}{T_0 \ldots T_n}
\text{ 于 }fg.
$$
换言之，基元素 $T_0^{e_0} \ldots T_n^{e_n}$ 与基元素
$T_0^{d_0} \ldots T_n^{d_n}$ 配对得到 $1$，当且仅当
$e_i + d_i = -1$ 对所有 $i$ 成立；在其余所有情形下配对为零。
利用这个配对，得到等同
$$
\left(\frac{1}{T_0 \ldots T_n} R[\frac{1}{T_0}, \ldots, \frac{1}{T_n}]\right)_d
=
\Hom_R((R[T_0, \ldots, T_n])_{-n - 1 - d}, R)
$$
因此可将引理 \ref{coherent-lemma-cohomology-projective-space-over-ring} 的结果
重新表述为
\begin{equation}
\label{coherent-equation-identify}
H^q(\mathbf{P}^n, \mathcal{O}_{\mathbf{P}^n_R}(d)) =
\left\{
\begin{matrix}
(R[T_0, \ldots, T_n])_d & \text{若} & q = 0 \\
0 & \text{若} & q \not = 0, n \\
\Hom_R((R[T_0, \ldots, T_n])_{-n - 1 - d}, R)
& \text{若} & q = n
\end{matrix}
\right.
\end{equation}

\begin{lemma}
\label{coherent-lemma-identify-functorially}
公式 (\ref{coherent-equation-identify}) 中的等同与环同态 $R \to R'$ 的基变换相容。
此外，对任意齐次元 $f \in R[T_0, \ldots, T_n]$，若其次数为 $m$，
则乘以 $f$ 的映射
$$
\mathcal{O}_{\mathbf{P}^n_R}(d)
\longrightarrow
\mathcal{O}_{\mathbf{P}^n_R}(d + m)
$$
在公式 (\ref{coherent-equation-identify}) 的等同下，于上同调群上诱导的映射
为乘以 $f$（在 $H^0$ 上），以及乘以 $f$ 的对偶映射
$$
(R[T_0, \ldots, T_n])_{-n - 1 - (d + m)}
\longrightarrow
(R[T_0, \ldots, T_n])_{-n - 1 - d}
$$
（在 $H^n$ 上）。
\end{lemma}

\begin{proof}
假设 $R \to R'$ 为环同态。
令 $\mathcal{U}$ 为 $\mathbf{P}^n_R$ 的标准仿射开覆盖，
令 $\mathcal{U}'$ 为 $\mathbf{P}^n_{R'}$ 的标准仿射开覆盖。
注意，$\mathcal{U}'$ 是覆盖 $\mathcal{U}$ 沿典范态射
$\mathbf{P}^n_{R'} \to \mathbf{P}^n_R$ 的拉回。因此有 {\v C}ech 复形的映射
$$
\gamma :
\check{\mathcal{C}}_{ord}^\bullet(\mathcal{U},
\mathcal{O}_{\mathbf{P}_R}(d))
\longrightarrow
\check{\mathcal{C}}_{ord}^\bullet(\mathcal{U}',
\mathcal{O}_{\mathbf{P}_{R'}}(d))
$$
由《上同调》引理 \ref{cohomology-lemma-functoriality-cech}，
它与上同调上的映射相容。
从引理 \ref{coherent-lemma-cohomology-projective-space-over-ring} 证明中的计算可立即看出，
这个 {\v C}ech 复形映射与所论上同调群的等同相容。
（也就是说，$R$ 上 {\v C}ech 复形的基元素恰好映到
$R'$ 上 {\v C}ech 复形的相应基元素。）这证明了引理的第一项陈述。

\medskip\noindent
现在固定环 $R$，并考虑两个齐次多项式
$f, g \in R[T_0, \ldots, T_n]$，二者次数同为 $m$。
由于上同调是加性函子，乘以 $f + g$ 所诱导的映射显然等于
乘以 $f$ 所诱导的映射与乘以 $g$ 所诱导的映射之和。
此外，由于上同调是函子，乘以乘积 $fg$ 也有类似结论，
其中 $f, g$ 都齐次（但次数不必相同）。
因此，为验证引理的第二项陈述，只需证明 $f = x \in R$ 或 $f = T_i$ 的情形。
在乘以元素 $x \in R$ 的情形，结论成立，因为所涉及的每个上同调群或复形
都具有 $R$-模或 $R$-模复形的结构。
最后，考虑乘以 $T_i$，将其视为 $\mathcal{O}_{\mathbf{P}^n_R}$-线性映射
$$
\mathcal{O}_{\mathbf{P}^n_R}(d)
\longrightarrow
\mathcal{O}_{\mathbf{P}^n_R}(d + 1)
$$
关于 $H^0$ 的陈述显然。关于 $H^n$ 的陈述，
将乘以 $T_i$ 视为 {\v C}ech 复形上的映射
$$
\check{\mathcal{C}}_{ord}^\bullet(\mathcal{U},
\mathcal{O}_{\mathbf{P}_R}(d))
\longrightarrow
\check{\mathcal{C}}_{ord}^\bullet(\mathcal{U},
\mathcal{O}_{\mathbf{P}_R}(d + 1))
$$
将使用引理 \ref{coherent-lemma-cohomology-projective-space-over-ring} 的证明中引入的记号。
依据分解
$$
\check{\mathcal{C}}_{ord}^\bullet(\mathcal{U}, \mathcal{O}_{\mathbf{P}_R}(d))
=
\bigoplus\nolimits_{\vec{e} \in \mathbf{Z}^{n + 1}, \ \sum e_i = d}
\check{\mathcal{C}}^\bullet(\vec{e})
$$
以及
$$
\check{\mathcal{C}}_{ord}^\bullet(\mathcal{U},
\mathcal{O}_{\mathbf{P}_R}(d + 1))
=
\bigoplus\nolimits_{\vec{e} \in \mathbf{Z}^{n + 1}, \ \sum e_i = d + 1}
\check{\mathcal{C}}^\bullet(\vec{e})
$$
考察乘以 $T_i$ 的作用。
显然，它将子复形 $\check{\mathcal{C}}^\bullet(\vec{e})$ 映到子复形
$\check{\mathcal{C}}^\bullet(\vec{e} + \vec{b}_i)$，其中
$\vec{b}_i = (0, \ldots, 0, 1, 0, \ldots, 0))$ 是第 $i$ 个基向量。
换言之，它将 $H^n$ 中对应于 $\vec{e}$、满足 $e_i < 0$ 且
$\sum e_i = d$ 的直和项映到 $H^n$ 中对应于
$\vec{e} + \vec{b}_i$ 的直和项（若 $e_i + b_i \geq 0$，后者为零）。
容易看出，这恰好对应于乘以 $T_i$ 的对偶映射的作用，即映射
$$
(R[T_0, \ldots, T_n])_{-n - 1 - (d + 1)}
\longrightarrow
(R[T_0, \ldots, T_n])_{-n - 1 - d}
$$
引理得证。
\end{proof}

\noindent
在陈述相对版本前，需要引入一些记号。
回忆，$\mathcal{O}_S[T_0, \ldots, T_n]$ 是分次
$\mathcal{O}_S$-模，其中每个 $T_i$ 都是次数为 $1$ 的齐次元。
以 $(\mathcal{O}_S[T_0, \ldots, T_n])_d$ 表示次数为 $d$ 的分量。
它是秩为 $\binom{n + d}{d}$ 的有限局部自由层，定义在 $S$ 上。

\begin{lemma}
\label{coherent-lemma-cohomology-projective-space-over-base}
设 $S$ 为概形。
设 $n \geq 0$ 为整数。
考虑结构态射
$$
f : \mathbf{P}^n_S \longrightarrow S.
$$
则
$$
R^qf_*(\mathcal{O}_{\mathbf{P}^n_S}(d)) =
\left\{
\begin{matrix}
(\mathcal{O}_S[T_0, \ldots, T_n])_d & \text{若} & q = 0 \\
0 & \text{若} & q \not = 0, n \\
\SheafHom_{\mathcal{O}_S}(
(\mathcal{O}_S[T_0, \ldots, T_n])_{- n - 1 - d}, \mathcal{O}_S)
& \text{若} & q = n
\end{matrix}
\right.
$$
\end{lemma}

\begin{proof}
从略。提示：这由公式 (\ref{coherent-equation-identify}) 中的等同
与仿射基变换相容而得，见引理 \ref{coherent-lemma-identify-functorially}。
\end{proof}

\noindent
下面陈述与有限局部自由层相伴的射影丛版本。
设 $S$ 为概形。设 $\mathcal{E}$ 为有限局部自由 $\mathcal{O}_S$-模，
其秩恒为 $n + 1$；见《模》节 \ref{modules-section-locally-free}。
此时把 $\text{Sym}(\mathcal{E})$ 看作分次 $\mathcal{O}_S$-模，
其中 $\mathcal{E}$ 是次数为 $1$ 的分次部分。
而 $\text{Sym}^d(\mathcal{E})$ 是次数为 $d$ 的分量。
它是秩为 $\binom{n + d}{d}$ 的有限局部自由层，定义在 $S$ 上。
回忆，我们的规范化约定为
$$
\pi :
\mathbf{P}(\mathcal{E})
=
\underline{\text{Proj}}_S(\text{Sym}(\mathcal{E}))
\longrightarrow
S
$$
并且有自然映射
$\text{Sym}^d(\mathcal{E}) \to \pi_*\mathcal{O}_{\mathbf{P}(\mathcal{E})}(d)$。

\begin{lemma}
\label{coherent-lemma-cohomology-projective-bundle}
设 $S$ 为概形。设 $n \geq 1$。
设 $\mathcal{E}$ 为有限局部自由 $\mathcal{O}_S$-模，
其秩恒为 $n + 1$。
考虑结构态射
$$
\pi : \mathbf{P}(\mathcal{E}) \longrightarrow S.
$$
则
$$
R^q\pi_*(\mathcal{O}_{\mathbf{P}(\mathcal{E})}(d)) =
\left\{
\begin{matrix}
\text{Sym}^d(\mathcal{E}) & \text{若} & q = 0 \\
0 & \text{若} & q \not = 0, n \\
\SheafHom_{\mathcal{O}_S}(
\text{Sym}^{- n - 1 - d}(\mathcal{E})
\otimes_{\mathcal{O}_S}
\wedge^{n + 1}\mathcal{E},
\mathcal{O}_S)
& \text{若} & q = n
\end{matrix}
\right.
$$
这些等同与基变换以及局部自由层之间的同构相容。
\end{lemma}

\begin{proof}
考虑典范映射
$$
\pi^*\mathcal{E} \longrightarrow \mathcal{O}_{\mathbf{P}(\mathcal{E})}(1)
$$
向下扭转 $1$，得到
$$
\pi^*(\mathcal{E})(-1) \longrightarrow \mathcal{O}_{\mathbf{P}(\mathcal{E})}
$$
这是从秩为 $n + 1$ 的局部自由层到结构层的满射。
因此相应的 Koszul 复形正合（《更多代数》引理
\ref{more-algebra-lemma-homotopy-koszul-abstract}）。
换言之，存在正合复形
$$
0 \to
\pi^*(\wedge^{n + 1}\mathcal{E})(-n - 1) \to
\ldots \to
\pi^*(\wedge^i\mathcal{E})(-i) \to
\ldots \to
\pi^*\mathcal{E}(-1) \to
\mathcal{O}_{\mathbf{P}(\mathcal{E})} \to 0
$$
将项 $\pi^*(\wedge^i\mathcal{E})(-i)$ 视为位于次数 $-i$。
使用《导出范畴》引理 \ref{derived-lemma-two-ss-complex-functor}
中的第一谱序列，计算这个无环复形的高阶直像。
具体而言，存在各项为
$$
E_1^{p, q} = R^q\pi_*\left(\pi^*(\wedge^{-p}\mathcal{E})(p)\right)
$$
且收敛到零的谱序列！由投影公式
（《上同调》引理 \ref{cohomology-lemma-projection-formula}），有
$$
E_1^{p, q} = \wedge^{-p} \mathcal{E} \otimes_{\mathcal{O}_S}
R^q\pi_*\left(\mathcal{O}_{\mathbf{P}(\mathcal{E})}(p)\right).
$$
注意，在 $S$ 上局部，层 $\mathcal{E}$ 平凡，
即同构于 $\mathcal{O}_S^{\oplus n + 1}$；故在 $S$ 上局部，
态射 $\mathbf{P}(\mathcal{E}) \to S$ 可与
$\mathbf{P}^n_S \to S$ 等同。因此在 $S$ 上局部，
可使用引理 \ref{coherent-lemma-cohomology-projective-space-over-ring}、
\ref{coherent-lemma-identify-functorially} 或
\ref{coherent-lemma-cohomology-projective-space-over-base} 的结果。
由此，$E_1^{p, q} = 0$，除非 $(p, q)$ 为 $(0, 0)$
或 $(-n - 1, n)$。非零项为
\begin{align*}
E_1^{0, 0} & = \pi_*\mathcal{O}_{\mathbf{P}(\mathcal{E})} = \mathcal{O}_S \\
E_1^{-n - 1, n} & =
R^n\pi_*\left(\pi^*(\wedge^{n + 1}\mathcal{E})(-n - 1)\right) =
\wedge^{n + 1}\mathcal{E} \otimes_{\mathcal{O}_S}
R^n\pi_*\left(\mathcal{O}_{\mathbf{P}(\mathcal{E})}(-n - 1)\right)
\end{align*}
因此谱序列中只能有一个非零微分，即映射
$d_{n + 1}^{-n - 1, n} : E_{n + 1}^{-n - 1, n} \to E_{n + 1}^{0, 0}$；
它必为同构（因为谱序列收敛到 $0$ 层）。
于是 $E_1^{p, q} = E_{n + 1}^{p, q}$，并得到典范同构
$$
\wedge^{n + 1}\mathcal{E} \otimes_{\mathcal{O}_S}
R^n\pi_*\left(\mathcal{O}_{\mathbf{P}(\mathcal{E})}(-n - 1)\right) =
R^n\pi_*\left(\pi^*(\wedge^{n + 1}\mathcal{E})(-n - 1)\right)
\xrightarrow{d_{n + 1}^{-n - 1, n}}
\mathcal{O}_S
$$
由于 $\wedge^{n + 1}\mathcal{E}$ 是可逆层，这推出
$R^n\pi_*\mathcal{O}_{\mathbf{P}(\mathcal{E})}(-n - 1)$ 也可逆，
并且典范同构于 $\wedge^{n + 1}\mathcal{E}$ 的逆。
换言之，已经证明引理中 $d = - n - 1$ 的情形。

\medskip\noindent
在 $S$ 上局部工作，由引理
\ref{coherent-lemma-cohomology-projective-space-over-ring}、
\ref{coherent-lemma-identify-functorially} 或
\ref{coherent-lemma-cohomology-projective-space-over-base} 的上同调计算，
立即得到引理中的消失陈述。此外，关于 $R^0\pi_*$ 的结果也很清楚，
因为存在典范映射
$\text{Sym}^d(\mathcal{E}) \to \pi_* \mathcal{O}_{\mathbf{P}(\mathcal{E})}(d)$
对所有 $d$ 都成立。
还须证明 $R^n\pi_*\mathcal{O}_{\mathbf{P}(\mathcal{E})}(d)$ 的描述
对 $d < -n - 1$ 正确。为此考虑映射
$$
\pi^*(\text{Sym}^{-d - n - 1}(\mathcal{E}))
\otimes_{\mathcal{O}_{\mathbf{P}(\mathcal{E})}}
\mathcal{O}_{\mathbf{P}(\mathcal{E})}(d)
\longrightarrow
\mathcal{O}_{\mathbf{P}(\mathcal{E})}(-n - 1)
$$
应用 $R^n\pi_*$ 及投影公式（见上文），得到映射
$$
\text{Sym}^{-d - n - 1}(\mathcal{E})
\otimes_{\mathcal{O}_S}
R^n\pi_*(\mathcal{O}_{\mathbf{P}(\mathcal{E})}(d))
\longrightarrow
R^n\pi_*\mathcal{O}_{\mathbf{P}(\mathcal{E})}(-n - 1) =
(\wedge^{n + 1}\mathcal{E})^{\otimes -1}
$$
（最后一个等式已在上面证明）。
再次利用引理
\ref{coherent-lemma-cohomology-projective-space-over-ring}、
\ref{coherent-lemma-identify-functorially} 或
\ref{coherent-lemma-cohomology-projective-space-over-base} 的局部计算，
可知此映射在
$R^n\pi_*(\mathcal{O}_{\mathbf{P}(\mathcal{E})}(d))$ 与
$\text{Sym}^{-d - n - 1}(\mathcal{E}) \otimes \wedge^{n + 1}(\mathcal{E})$
之间诱导完美配对，如所欲。
\end{proof}

















\section{局部 Noether 概形上的凝聚层}
\label{coherent-section-coherent-sheaves}

\noindent
我们已在《模》节 \ref{modules-section-coherent} 中定义了
任意环层空间上凝聚模的概念。
尽管可以考虑非 Noether 概形上的凝聚层，
在考察凝聚层时，我们总是假设基概形局部 Noether。
下面给出局部 Noether 概形上凝聚层的一个刻画。

\begin{lemma}
\label{coherent-lemma-coherent-Noetherian}
设 $X$ 为局部 Noether 概形。
设 $\mathcal{F}$ 为 $\mathcal{O}_X$-模。
下列条件等价：
\begin{enumerate}
\item $\mathcal{F}$ 凝聚；
\item $\mathcal{F}$ 是拟凝聚有限型 $\mathcal{O}_X$-模；
\item $\mathcal{F}$ 是有限表示 $\mathcal{O}_X$-模；
\item 对任意仿射开集 $\Spec(A) = U \subset X$，都有
$\mathcal{F}|_U = \widetilde M$，其中 $M$ 为有限 $A$-模；
\item 存在仿射开覆盖 $X = \bigcup U_i$，
$U_i = \Spec(A_i)$，使得每个
$\mathcal{F}|_{U_i} = \widetilde M_i$，其中 $M_i$ 为有限 $A_i$-模。
\end{enumerate}
特别地，$\mathcal{O}_X$ 凝聚；任意可逆 $\mathcal{O}_X$-模凝聚；
更一般地，任意有限局部自由 $\mathcal{O}_X$-模凝聚。
\end{lemma}

\begin{proof}
蕴含 (1) $\Rightarrow$ (2) 与 (1) $\Rightarrow$ (3) 一般成立，见
《模》引理 \ref{modules-lemma-coherent-finite-presentation}。
若 $\mathcal{F}$ 有限表示，则 $\mathcal{F}$ 拟凝聚；见
《模》引理 \ref{modules-lemma-finite-presentation-quasi-coherent}。
故 (3) $\Rightarrow$ (2) 也成立。

\medskip\noindent
假设 $\mathcal{F}$ 是拟凝聚有限型 $\mathcal{O}_X$-模。
由《性质》引理 \ref{properties-lemma-finite-type-module}，
在任意仿射开集 $\Spec(A) = U \subset X$ 上都有
$\mathcal{F}|_U = \widetilde M$，其中 $M$ 为有限 $A$-模。
由于 $A$ Noether，$M$ 有有限解消
$$
A^{\oplus m} \to A^{\oplus n} \to M \to 0.
$$
因此由《性质》引理 \ref{properties-lemma-finite-presentation-module}，
$\mathcal{F}$ 有限表示。换言之，(2) $\Rightarrow$ (3)。

\medskip\noindent
由《模》引理 \ref{modules-lemma-coherent-structure-sheaf}，
为证明 (3) 推出 (1)，只需证明 $\mathcal{O}_X$ 凝聚。
因此须证明：给定任意开集 $U \subset X$ 与任意有限组截面
$f_i \in \mathcal{O}_X(U)$，$i = 1, \ldots, n$，映射
$\bigoplus_{i = 1, \ldots, n} \mathcal{O}_U \to \mathcal{O}_U$
的核为有限型。有限型是局部性质，故只需在任意 $x \in U$ 的邻域中检验。
于是可假设 $U = \Spec(A)$ 仿射。此时
$f_1, \ldots, f_n \in A$ 是 $A$ 的元素。由于 $A$ Noether，见
《性质》引理 \ref{properties-lemma-locally-Noetherian}，其核 $K$，即映射
$\bigoplus_{i = 1, \ldots, n} A \to A$ 的核，是有限 $A$-模；
例如见《代数》引理
\ref{algebra-lemma-Noetherian-basic}。
由于函子\ $\widetilde{ }$\ 正合，见《概形》引理
\ref{schemes-lemma-spec-sheaves}，得到正合序列
$$
\widetilde K \to
\bigoplus\nolimits_{i = 1, \ldots, n} \mathcal{O}_U \to
\mathcal{O}_U
$$
再由《性质》引理 \ref{properties-lemma-finite-type-module}，
可见 $\widetilde K$ 为有限型。因此 (1)、(2)、(3) 全都等价。

\medskip\noindent
由《性质》引理 \ref{properties-lemma-finite-type-module}，
(2) 推出 (4)；(4) 推出 (5) 是平凡的。
《概形》节 \ref{schemes-section-quasi-coherent} 中的讨论表明，
(5) 推出 $\mathcal{F}$ 拟凝聚；显然，(5) 还推出
$\mathcal{F}$ 为有限型。因此 (5) 推出 (2)，证毕。
\end{proof}

\begin{lemma}
\label{coherent-lemma-coherent-abelian-Noetherian}
设 $X$ 为局部 Noether 概形。
凝聚 $\mathcal{O}_X$-模所成的范畴是阿贝尔范畴。
更确切地，凝聚 $\mathcal{O}_X$-模之间映射的核与余核均凝聚。
凝聚层的任意扩张均凝聚。
\end{lemma}

\begin{proof}
这是《模》引理 \ref{modules-lemma-coherent-abelian}
在这一特殊情形下的重述。
\end{proof}

\noindent
下列引理对于凝聚 $\mathcal{O}_X$-模的范畴，
在一般环层空间 $X$ 上并不总是成立。

\begin{lemma}
\label{coherent-lemma-coherent-Noetherian-quasi-coherent-sub-quotient}
设 $X$ 为局部 Noether 概形。
设 $\mathcal{F}$ 为凝聚 $\mathcal{O}_X$-模。
$\mathcal{F}$ 的任意拟凝聚子模均凝聚。
$\mathcal{F}$ 的任意拟凝聚商模均凝聚。
\end{lemma}

\begin{proof}
可假设 $X$ 仿射，设 $X = \Spec(A)$。
《性质》引理 \ref{properties-lemma-locally-Noetherian} 推出 $A$ Noether。
引理 \ref{coherent-lemma-coherent-Noetherian} 将问题化为代数。
该引理的代数对应是：有限 $A$-模的商或子模仍为有限 $A$-模；
例如见《代数》引理 \ref{algebra-lemma-Noetherian-basic}。
\end{proof}

\begin{lemma}
\label{coherent-lemma-tensor-hom-coherent}
设 $X$ 为局部 Noether 概形。
设 $\mathcal{F}$、$\mathcal{G}$ 为凝聚 $\mathcal{O}_X$-模。
则 $\mathcal{O}_X$-模 $\mathcal{F} \otimes_{\mathcal{O}_X} \mathcal{G}$
与 $\SheafHom_{\mathcal{O}_X}(\mathcal{F}, \mathcal{G})$ 都凝聚。
\end{lemma}

\begin{proof}
《模》引理 \ref{modules-lemma-internal-hom-locally-kernel-direct-sum} 已证明
$\SheafHom_{\mathcal{O}_X}(\mathcal{F}, \mathcal{G})$ 凝聚。
张量积的结论是《模》引理 \ref{modules-lemma-tensor-product-permanence}。
\end{proof}

\begin{lemma}
\label{coherent-lemma-local-isomorphism}
设 $X$ 为局部 Noether 概形。
设 $\mathcal{F}$、$\mathcal{G}$ 为凝聚 $\mathcal{O}_X$-模。
设 $\varphi : \mathcal{G} \to \mathcal{F}$ 为 $\mathcal{O}_X$-模同态。
设 $x \in X$。
\begin{enumerate}
\item 若 $\mathcal{F}_x = 0$，则存在开邻域 $U \subset X$，
它是 $x$ 的邻域，且 $\mathcal{F}|_U = 0$。
\item 若 $\varphi_x : \mathcal{G}_x \to \mathcal{F}_x$ 单射，
则存在开邻域 $U \subset X$，它是 $x$ 的邻域，且 $\varphi|_U$ 单射。
\item 若 $\varphi_x : \mathcal{G}_x \to \mathcal{F}_x$ 满射，
则存在开邻域 $U \subset X$，它是 $x$ 的邻域，且 $\varphi|_U$ 满射。
\item 若 $\varphi_x : \mathcal{G}_x \to \mathcal{F}_x$ 双射，
则存在开邻域 $U \subset X$，它是 $x$ 的邻域，且 $\varphi|_U$ 为同构。
\end{enumerate}
\end{lemma}

\begin{proof}
见《模》引理
\ref{modules-lemma-finite-type-surjective-on-stalk}、
\ref{modules-lemma-finite-type-stalk-zero} 及
\ref{modules-lemma-finite-type-to-coherent-injective-on-stalk}。
\end{proof}

\begin{lemma}
\label{coherent-lemma-map-stalks-local-map}
设 $X$ 为局部 Noether 概形。
设 $\mathcal{F}$、$\mathcal{G}$ 为凝聚 $\mathcal{O}_X$-模。
设 $x \in X$。
假设 $\psi : \mathcal{G}_x \to \mathcal{F}_x$ 是
$\mathcal{O}_{X, x}$-模映射。
则存在开邻域 $U \subset X$（它是 $x$ 的邻域）及映射
$\varphi : \mathcal{G}|_U \to \mathcal{F}|_U$，使 $\varphi_x = \psi$。
\end{lemma}

\begin{proof}
根据引理 \ref{coherent-lemma-coherent-Noetherian}，
这是《模》引理 \ref{modules-lemma-stalk-internal-hom} 的重述。
\end{proof}

\begin{lemma}
\label{coherent-lemma-coherent-support-closed}
设 $X$ 为局部 Noether 概形。设 $\mathcal{F}$ 为凝聚
$\mathcal{O}_X$-模。则 $\text{Supp}(\mathcal{F})$ 闭，且
$\mathcal{F}$ 来自 $\mathcal{F}$ 的概形论支集上的凝聚层；见
《态射》定义 \ref{morphisms-definition-scheme-theoretic-support}。
\end{lemma}

\begin{proof}
设 $i : Z \to X$ 为 $\mathcal{F}$ 的概形论支集，
并设 $\mathcal{G}$ 为 $Z$ 上的有限型拟凝聚层，使
$i_*\mathcal{G} \cong \mathcal{F}$。
由于 $Z = \text{Supp}(\mathcal{F})$，支集闭。
由《态射》引理 \ref{morphisms-lemma-immersion-locally-finite-type}
及 \ref{morphisms-lemma-finite-type-noetherian}，概形 $Z$ 局部 Noether。
最后，由引理 \ref{coherent-lemma-coherent-Noetherian}，
$\mathcal{G}$ 是凝聚 $\mathcal{O}_Z$-模。
\end{proof}

\begin{lemma}
\label{coherent-lemma-i-star-equivalence}
设 $i : Z \to X$ 为局部 Noether 概形的闭浸入。
设 $\mathcal{I} \subset \mathcal{O}_X$ 为切出 $Z$ 的拟凝聚理想层。
函子 $i_*$ 在凝聚 $\mathcal{O}_X$-模（由 $\mathcal{I}$ 零化）所成的范畴与
凝聚 $\mathcal{O}_Z$-模范畴之间诱导等价。
\end{lemma}

\begin{proof}
由《态射》引理 \ref{morphisms-lemma-i-star-equivalence}，该函子全忠实。
设 $\mathcal{F}$ 为凝聚 $\mathcal{O}_X$-模，并由 $\mathcal{I}$ 零化。
由《态射》引理 \ref{morphisms-lemma-i-star-equivalence}，
可写成 $\mathcal{F} = i_*\mathcal{G}$，其中 $\mathcal{G}$
为 $Z$ 上的某个拟凝聚层。由《模》引理
\ref{modules-lemma-i-star-reflects-finite-type}，可见 $\mathcal{G}$ 有限型。
因此由引理 \ref{coherent-lemma-coherent-Noetherian}，$\mathcal{G}$ 凝聚。
故该函子也本质满，如所欲。
\end{proof}

\begin{lemma}
\label{coherent-lemma-finite-pushforward-coherent}
设 $f : X \to Y$ 为概形态射。
设 $\mathcal{F}$ 为拟凝聚 $\mathcal{O}_X$-模。
假设 $f$ 有限且 $Y$ 局部 Noether。
则 $R^pf_*\mathcal{F} = 0$ 对 $p > 0$ 成立，并且 $f_*\mathcal{F}$ 凝聚，
如果 $\mathcal{F}$ 凝聚。
\end{lemma}

\begin{proof}
高阶直像由引理 \ref{coherent-lemma-relative-affine-vanishing} 消失，
且有限态射按定义仿射。
注意，这些假设还推出 $X$ 局部 Noether
（见《态射》引理 \ref{morphisms-lemma-finite-type-noetherian}），
故陈述有意义。
设 $\Spec(A) = V \subset Y$ 为仿射开子集。
由《态射》定义 \ref{morphisms-definition-integral}，
有 $f^{-1}(V) = \Spec(B)$，其中 $A \to B$ 有限。
引理 \ref{coherent-lemma-coherent-Noetherian} 将引理的陈述化为如下代数事实：
若 $M$ 是有限 $B$-模，则把 $M$ 看作 $A$-模时也有限；
见《代数》引理 \ref{algebra-lemma-finite-module-over-finite-extension}。
\end{proof}

\noindent
在该引理的情形下，高阶直像也凝聚，因为它们消失。
我们将证明，对局部 Noether 概形之间的任意固有态射，
这一结论总是成立（命题 \ref{coherent-proposition-proper-pushforward-coherent}）。

\begin{lemma}
\label{coherent-lemma-coherent-support-dimension-0}
设 $X$ 为局部 Noether 概形。设 $\mathcal{F}$ 为凝聚层，且
$\dim(\text{Supp}(\mathcal{F})) \leq 0$。
则 $\mathcal{F}$ 由全局截面生成，并且
$H^i(X, \mathcal{F}) = 0$ 对 $i > 0$ 成立。
\end{lemma}

\begin{proof}
由引理 \ref{coherent-lemma-coherent-support-closed}，可写
$\mathcal{F} = i_*\mathcal{G}$，其中 $i : Z \to X$ 是
$\mathcal{F}$ 的概形论支集的包含，而 $\mathcal{G}$ 是凝聚
$\mathcal{O}_Z$-模。由于 $Z$ 的维数为 $0$，
可知 $Z$ 是仿射概形的不交并（《性质》引理
\ref{properties-lemma-locally-Noetherian-dimension-0}）。
因此 $\mathcal{G}$ 全局生成，且 $\mathcal{G}$ 的高阶上同调群为零
（引理 \ref{coherent-lemma-quasi-coherent-affine-cohomology-zero}）。
故 $\mathcal{F} = i_*\mathcal{G}$ 全局生成。
由于 $\mathcal{F}$ 与 $\mathcal{G}$ 的上同调相同
（闭浸入仿射，故可应用引理 \ref{coherent-lemma-relative-affine-cohomology}），
可知 $\mathcal{F}$ 的高阶上同调群为零。
\end{proof}

\begin{lemma}
\label{coherent-lemma-pushforward-coherent-on-open}
设 $X$ 为概形。设 $j : U \to X$ 为开集的包含。
设 $T \subset X$ 为包含于 $U$ 的闭子集。
若 $\mathcal{F}$ 是凝聚 $\mathcal{O}_U$-模，且
$\text{Supp}(\mathcal{F}) \subset T$，则
$j_*\mathcal{F}$ 是凝聚 $\mathcal{O}_X$-模。
\end{lemma}

\begin{proof}
考虑开覆盖 $X = U \cup (X \setminus T)$。
则 $j_*\mathcal{F}|_U = \mathcal{F}$ 凝聚，
且 $j_*\mathcal{F}|_{X \setminus T} = 0$ 也凝聚。
故 $j_*\mathcal{F}$ 凝聚。
\end{proof}

\section{Noether 概形上的凝聚层}
\label{coherent-section-coherent-quasi-compact}

\noindent
本节列举 Noether 概形上凝聚层的一些性质。

\begin{lemma}
\label{coherent-lemma-acc-coherent}
设 $X$ 为 Noether 概形。
设 $\mathcal{F}$ 为凝聚 $\mathcal{O}_X$-模。
$\mathcal{F}$ 的准凝聚子模满足升链条件。换言之，给定准凝聚子模的任意序列
$$
\mathcal{F}_1 \subset \mathcal{F}_2 \subset \ldots \subset \mathcal{F}
$$
则有 $\mathcal{F}_n = \mathcal{F}_{n + 1} = \ldots $
对某个 $n \geq 0$ 成立。
\end{lemma}

\begin{proof}
选取一个有限仿射开覆盖。
在覆盖的每个成员上，由《代数》引理
\ref{algebra-lemma-Noetherian-basic} 可知该链稳定。
故引理成立。
\end{proof}

\begin{lemma}
\label{coherent-lemma-power-ideal-kills-sheaf}
设 $X$ 为 Noether 概形。
设 $\mathcal{F}$ 为 $X$ 上的凝聚层。
设 $\mathcal{I} \subset \mathcal{O}_X$ 为对应于闭子概形
$Z \subset X$ 的准凝聚理想层。
则存在某个 $n \geq 0$ 使得 $\mathcal{I}^n\mathcal{F} = 0$，
当且仅当 $\text{Supp}(\mathcal{F}) \subset Z$（在集合论意义下）。
\end{lemma}

\begin{proof}
由于 $X$ 可由有限多个 Noether 环的谱覆盖，结论立即由
《代数》引理 \ref{algebra-lemma-Noetherian-power-ideal-kills-module}
推出。
\end{proof}

\begin{lemma}[Artin--Rees]
\label{coherent-lemma-Artin-Rees}
设 $X$ 为 Noether 概形。
设 $\mathcal{F}$ 为 $X$ 上的凝聚层。
设 $\mathcal{G} \subset \mathcal{F}$ 为准凝聚子层。
设 $\mathcal{I} \subset \mathcal{O}_X$ 为准凝聚理想层。
则存在 $c \geq 0$，使得对所有 $n \geq c$ 都有
$$
\mathcal{I}^{n - c}(\mathcal{I}^c\mathcal{F} \cap \mathcal{G})
=
\mathcal{I}^n\mathcal{F} \cap \mathcal{G}
$$
\end{lemma}

\begin{proof}
由于 $X$ 可由有限多个 Noether 环的谱覆盖，结论立即由
《代数》引理 \ref{algebra-lemma-Artin-Rees} 推出。
\end{proof}

\begin{lemma}
\label{coherent-lemma-directed-colimit-coherent}
设 $X$ 为 Noether 概形。每个准凝聚 $\mathcal{O}_X$-模
都是其凝聚子模的滤过余极限。
\end{lemma}

\begin{proof}
本结论只是《性质》引理
\ref{properties-lemma-quasi-coherent-colimit-finite-type}
的另一种表述；这里用到由
引理 \ref{coherent-lemma-coherent-Noetherian} 可知有限型准凝聚
$\mathcal{O}_X$-模是凝聚的这一事实。
\end{proof}

\begin{lemma}
\label{coherent-lemma-homs-over-open}
设 $X$ 为 Noether 概形。
设 $\mathcal{F}$ 为准凝聚 $\mathcal{O}_X$-模。
设 $\mathcal{G}$ 为凝聚 $\mathcal{O}_X$-模。
设 $\mathcal{I} \subset \mathcal{O}_X$ 为准凝聚理想层。
以 $Z \subset X$ 表示相应的闭子概形，并令
$U = X \setminus Z$。存在典范同构
$$
\colim_n \Hom_{\mathcal{O}_X}(\mathcal{I}^n\mathcal{G}, \mathcal{F})
\longrightarrow
\Hom_{\mathcal{O}_U}(\mathcal{G}|_U, \mathcal{F}|_U).
$$
特别地，存在同构
$$
\colim_n \Hom_{\mathcal{O}_X}(
\mathcal{I}^n, \mathcal{F})
\longrightarrow
\Gamma(U, \mathcal{F}).
$$
\end{lemma}

\begin{proof}
先证明第二个映射是同构。由《性质》引理
\ref{properties-lemma-sections-over-quasi-compact-open}，
它是单射。由于 $\mathcal{F}$ 是其凝聚子模的并，参见
《性质》引理 \ref{properties-lemma-quasi-coherent-colimit-finite-type}
（以及引理 \ref{coherent-lemma-coherent-Noetherian}），
为证明满射性，可以并且确实假设 $\mathcal{F}$ 凝聚。
以 $\mathcal{F}_n$ 表示 $\mathcal{F}$ 中由
$\mathcal{I}^n$ 零化的截面组成的准凝聚子层，参见
《性质》引理 \ref{properties-lemma-sections-over-quasi-compact-open}。
由于 $\mathcal{F}_1 \subset \mathcal{F}_2 \subset \ldots$，
可知 $\mathcal{F}_n = \mathcal{F}_{n + 1} = \ldots $
对某个 $n \geq 0$ 成立，这由引理
\ref{coherent-lemma-acc-coherent} 推出。
令 $\mathcal{H} = \mathcal{F}_n$，其中 $n$ 为上述整数。由 Artin--Rees
（引理 \ref{coherent-lemma-Artin-Rees}），存在 $c \geq 0$ 使得
$\mathcal{I}^m\mathcal{F} \cap \mathcal{H}
\subset \mathcal{I}^{m - c}\mathcal{H}$。
取 $m = n + c$，得到
$\mathcal{I}^m\mathcal{F} \cap \mathcal{H} \subset \mathcal{I}^n\mathcal{H}
= 0$。因此，若令 $\mathcal{F}' = \mathcal{I}^m\mathcal{F}$，
则有 $\mathcal{F}' \cap \mathcal{F}_n = 0$ 且
$\mathcal{F}'|_U = \mathcal{F}|_U$。特别注意到，子层
$(\mathcal{F}')_N$，即由 $\mathcal{I}^N$ 零化的截面所成子层，
对所有 $N \geq 0$ 均为零。因此由《性质》引理
\ref{properties-lemma-sections-over-quasi-compact-open}
可知，下列交换图中的上方水平箭头是双射：
$$
\xymatrix{
\colim_n \Hom_{\mathcal{O}_X}(
\mathcal{I}^n, \mathcal{F}')
\ar[r] \ar[d] &
\Gamma(U, \mathcal{F}') \ar[d] \\
\colim_n \Hom_{\mathcal{O}_X}(
\mathcal{I}^n, \mathcal{F})
\ar[r] &
\Gamma(U, \mathcal{F})
}
$$
右侧竖直箭头也是双射，故下方水平箭头如所要求是满射。

\medskip\noindent
接下来证明引理中的第一个箭头是双射。
由引理 \ref{coherent-lemma-coherent-Noetherian}，层 $\mathcal{G}$
是有限表示的，因而层
$\mathcal{H} = \SheafHom_{\mathcal{O}_X}(\mathcal{G}, \mathcal{F})$
是准凝聚的，参见《概形》第
\ref{schemes-section-quasi-coherent} 节。按定义有
$$
\mathcal{H}(U)
=
\Hom_{\mathcal{O}_U}(\mathcal{G}|_U, \mathcal{F}|_U)
$$
在引理第一个箭头的右端选取 $\psi$，即
$\psi \in \mathcal{H}(U)$。刚刚证明的结果适用于
$\mathcal{H}$，故存在 $n \geq 0$ 以及
$\varphi : \mathcal{I}^n \to \mathcal{H}$，它恢复
$\psi$（限制到 $U$ 后）。由《模》引理
\ref{modules-lemma-internal-hom}，$\varphi$ 对应于映射
$$
\varphi :
\mathcal{I}^n \otimes_{\mathcal{O}_X} \mathcal{G}
\longrightarrow
\mathcal{F}.
$$
这几乎就是所需结论，唯一差别是箭头的源是
$\mathcal{I}^n$ 与 $\mathcal{G}$ 的张量积，而非它们的积。
我们将说明，以增大 $n$ 为代价，这一区别无关紧要。
考虑短正合列
$$
0 \to \mathcal{K} \to
\mathcal{I}^n \otimes_{\mathcal{O}_X} \mathcal{G} \to
\mathcal{I}^n\mathcal{G} \to 0
$$
其中 $\mathcal{K}$ 定义为核。注意
$\mathcal{I}^n\mathcal{K} = 0$（证明略）。
再次应用 Artin--Rees 可知，对某个充分大的 $m$ 有
$$
\mathcal{K}
\cap
\mathcal{I}^m(\mathcal{I}^n \otimes_{\mathcal{O}_X} \mathcal{G})
=
0
$$
换言之，映射
$$
\mathcal{I}^m(\mathcal{I}^n \otimes_{\mathcal{O}_X} \mathcal{G})
\longrightarrow
\mathcal{I}^{n + m}\mathcal{G}
$$
是同构。设 $\varphi'$ 为 $\varphi$ 在该子模上的限制，
将它视为映射
$\mathcal{I}^{m + n}\mathcal{G} \to \mathcal{F}$。
于是 $\varphi'$ 给出引理第一个箭头左端的一个元素，
该元素经此箭头映为 $\psi$。换言之，我们已经证明该箭头满射。
其单射性的证明从略。
\end{proof}

\begin{lemma}
\label{coherent-lemma-extend-coherent}
设 $X$ 为局部 Noether 概形。设 $\mathcal{F}$、$\mathcal{G}$
为凝聚 $\mathcal{O}_X$-模。设 $U \subset X$ 为开集，并设
$\varphi : \mathcal{F}|_U \to \mathcal{G}|_U$ 为
$\mathcal{O}_U$-模映射。则存在凝聚子模
$\mathcal{F}' \subset \mathcal{F}$，它与
$\mathcal{F}$ 在 $U$ 上相同，且 $\varphi$ 可延拓为
$\varphi' : \mathcal{F}' \to \mathcal{G}$。
\end{lemma}

\begin{proof}
设 $\mathcal{I} \subset \mathcal{O}_X$ 为定义
$X \setminus U$ 上既约诱导概形结构的凝聚理想层。
若 $X$ 是 Noether 的，则引理 \ref{coherent-lemma-homs-over-open}
表明可取 $\mathcal{F}' = \mathcal{I}^n\mathcal{F}$，
其中 $n$ 为某个整数。
一般情形将借助 Zorn 引理由此推出。

\medskip\noindent
考虑所有三元组 $(U', \mathcal{F}', \varphi')$ 所成的集合，其中
$U \subset U' \subset X$ 为开集，
$\mathcal{F}' \subset \mathcal{F}|_{U'}$
是一个与 $\mathcal{F}$ 在 $U$ 上相同的凝聚子层，而
$\varphi' : \mathcal{F}' \to \mathcal{G}|_{U'}$
限制为 $\varphi$（在 $U$ 上）。我们约定
$(U'', \mathcal{F}'', \varphi'') \geq (U', \mathcal{F}', \varphi')$
当且仅当 $U'' \supset U'$、
$\mathcal{F}''|_{U'} = \mathcal{F}'$ 且
$\varphi''|_{U'} = \varphi'$。
显然，若有三元组的一个全序族
$(U_i, \mathcal{F}_i, \varphi_i)$，则可将
$\mathcal{F}_i$ 粘合成子层 $\mathcal{F}'$，它是
$\mathcal{F}$ 在 $U' = \bigcup U_i$ 上的子层，
并将 $\varphi$ 延拓为映射
$\varphi' : \mathcal{F}' \to \mathcal{G}|_{U'}$。
因此三元组的每个全序子集都有上界。
最后，设 $(U', \mathcal{F}', \varphi')$ 为任意三元组，
但 $U' \not = X$。则可选取仿射开集
$W \subset X$，它不包含于 $U'$。由第一段的结果，可以把子层
$\mathcal{F}'|_{W \cap U'}$ 及限制
$\varphi'|_{W \cap U'}$ 延拓为某个子层
$\mathcal{F}'' \subset \mathcal{F}|_W$ 及映射
$\varphi'' : \mathcal{F}'' \to \mathcal{G}|_W$。
当然，$(\mathcal{F}', \varphi')$ 与
$(\mathcal{F}'', \varphi'')$ 在 $W \cap U'$ 上相合，
恰好意味着可以将其延拓成三元组
$(U' \cup W, \mathcal{F}''', \varphi''')$。
因此任意极大三元组 $(U', \mathcal{F}', \varphi')$
（其存在性由 Zorn 引理保证）必有 $U' = X$，证明完毕。
\end{proof}

\section{深度}
\label{coherent-section-depth}

\noindent
本节简要讨论局部 Noether 概形上凝聚模的深度及性质
$(S_k)$。注意，我们已经在《性质》第
\ref{properties-section-Rk} 节讨论过局部 Noether 概形的这一概念。

\begin{definition}
\label{coherent-definition-depth}
设 $X$ 为局部 Noether 概形。
设 $\mathcal{F}$ 为凝聚 $\mathcal{O}_X$-模。
设 $k \geq 0$ 为整数。
\begin{enumerate}
\item 称 $\mathcal{F}$ 在一点处有{\it 深度 $k$}，若该点
$x$ 属于 $X$ 且
$\text{depth}_{\mathcal{O}_{X, x}}(\mathcal{F}_x) = k$。
\item 称 $X$ 在一点处有{\it 深度 $k$}，若该点
$x$ 属于 $X$ 且
$\text{depth}(\mathcal{O}_{X, x}) = k$。
\item 若 $\mathcal{F}$ 满足性质 {\it $(S_k)$}，是指
$$
\text{depth}_{\mathcal{O}_{X, x}}(\mathcal{F}_x)
\geq \min(k, \dim(\text{Supp}(\mathcal{F}_x)))
$$
对所有 $x \in X$ 成立。
\item 称 $X$ 满足性质 {\it $(S_k)$}，若 $\mathcal{O}_X$
满足性质 $(S_k)$。
\end{enumerate}
\end{definition}

\noindent
任意凝聚层都满足条件 $(S_0)$。
条件 $(S_1)$ 等价于没有嵌入伴随点，参见《除子》引理
\ref{divisors-lemma-S1-no-embedded}。

\begin{lemma}
\label{coherent-lemma-hom-into-depth}
设 $X$ 为局部 Noether 概形。设 $\mathcal{F}$、$\mathcal{G}$
为凝聚 $\mathcal{O}_X$-模，且 $x \in X$。
\begin{enumerate}
\item 若 $\mathcal{G}_x$ 的深度 $\geq 1$，则
$\SheafHom_{\mathcal{O}_X}(\mathcal{F}, \mathcal{G})_x$
的深度 $\geq 1$。
\item 若 $\mathcal{G}_x$ 的深度 $\geq 2$，则
$\Hom_{\mathcal{O}_X}(\mathcal{F}, \mathcal{G})_x$ 的深度 $\geq 2$。
\end{enumerate}
\end{lemma}

\begin{proof}
由引理 \ref{coherent-lemma-tensor-hom-coherent}，
$\SheafHom_{\mathcal{O}_X}(\mathcal{F}, \mathcal{G})$
是凝聚 $\mathcal{O}_X$-模。
凝聚模是有限表示的（引理 \ref{coherent-lemma-coherent-Noetherian}），
故取层茎与取 $\SheafHom$ 及 $\Hom$ 可交换，参见《模》引理
\ref{modules-lemma-stalk-internal-hom}。
因而可归结为局部环上的有限模情形，即《代数进阶》引理
\ref{more-algebra-lemma-hom-into-depth}。
\end{proof}

\begin{lemma}
\label{coherent-lemma-hom-into-S2}
设 $X$ 为局部 Noether 概形。设 $\mathcal{F}$、$\mathcal{G}$
为凝聚 $\mathcal{O}_X$-模。
\begin{enumerate}
\item 若 $\mathcal{G}$ 满足性质 $(S_1)$，则
$\SheafHom_{\mathcal{O}_X}(\mathcal{F}, \mathcal{G})$ 满足性质 $(S_1)$。
\item 若 $\mathcal{G}$ 满足性质 $(S_2)$，则
$\SheafHom_{\mathcal{O}_X}(\mathcal{F}, \mathcal{G})$ 满足性质 $(S_2)$。
\end{enumerate}
\end{lemma}

\begin{proof}
立即由引理 \ref{coherent-lemma-hom-into-depth} 和定义推出。
\end{proof}

\noindent
我们已在《性质》引理
\ref{properties-lemma-scheme-CM-iff-all-Sk} 中看到，
局部 Noether 概形是 Cohen--Macaulay 的，当且仅当
$(S_k)$ 对所有 $k$ 均成立。
因此引入下述定义是合理的；它等价于所有层茎均为
Cohen--Macaulay 模这一条件。

\begin{definition}
\label{coherent-definition-Cohen-Macaulay}
设 $X$ 为局部 Noether 概形。
设 $\mathcal{F}$ 为凝聚 $\mathcal{O}_X$-模。
称 $\mathcal{F}$ 是 {\it Cohen--Macaulay 的}，当且仅当
$(S_k)$ 对所有 $k \geq 0$ 均成立。
\end{definition}

\begin{lemma}
\label{coherent-lemma-Cohen-Macaulay-over-regular}
设 $X$ 为正则概形。设 $\mathcal{F}$ 为凝聚
$\mathcal{O}_X$-模。下列条件等价：
\begin{enumerate}
\item $\mathcal{F}$ 是 Cohen--Macaulay 的，且
$\text{Supp}(\mathcal{F}) = X$；
\item $\mathcal{F}$ 是秩 $> 0$ 的有限局部自由模。
\end{enumerate}
\end{lemma}

\begin{proof}
设 $x \in X$。若 (2) 成立，则 $\mathcal{F}_x$ 是自由
$\mathcal{O}_{X, x}$-模，且其秩 $> 0$。因此
$\text{depth}(\mathcal{F}_x) = \dim(\mathcal{O}_{X, x})$，
因为正则局部环是 Cohen--Macaulay 的（《代数》引理
\ref{algebra-lemma-regular-ring-CM}）。
反之，若 (1) 成立，则 $\mathcal{F}_x$ 是
$\mathcal{O}_{X, x}$ 上的极大 Cohen--Macaulay 模
（《代数》定义 \ref{algebra-definition-maximal-CM}）。
故由《代数》引理 \ref{algebra-lemma-regular-mcm-free}，
$\mathcal{F}_x$ 是自由的。
\end{proof}

\section{凝聚层的逐层剥离}
\label{coherent-section-devissage}

\noindent
设 $X$ 为 Noether 概形。考虑一个整的闭子概形
$i : Z \to X$。考察形如 $i_*\mathcal{G}$ 的凝聚层往往很方便，
其中 $\mathcal{G}$ 是 $Z$ 上的凝聚层。我们尤其关注
$\mathcal{G}$ 为无挠秩 $1$ 层的情形。例如，$\mathcal{G}$
可以是 $Z$ 上的非零理想层；更特殊地，也可取
$\mathcal{G} = \mathcal{O}_Z$。

\medskip\noindent
本节始终会用到：凝聚层与有限型准凝聚层是同一回事，
凝聚层的准凝聚子商仍然凝聚；参见第
\ref{coherent-section-coherent-sheaves} 节。
凝聚层的支集是闭的；参见《模》引理
\ref{modules-lemma-support-finite-type-closed}。

\begin{lemma}
\label{coherent-lemma-prepare-filter-support}
设 $X$ 为 Noether 概形。
设 $\mathcal{F}$ 为 $X$ 上的凝聚层。
假设 $\text{Supp}(\mathcal{F}) = Z \cup Z'$，其中
$Z$、$Z'$ 闭。则存在凝聚层的短正合列
$$
0 \to \mathcal{G}' \to \mathcal{F} \to \mathcal{G} \to 0
$$
使得 $\text{Supp}(\mathcal{G}') \subset Z'$ 且
$\text{Supp}(\mathcal{G}) \subset Z$。
\end{lemma}

\begin{proof}
设 $\mathcal{I} \subset \mathcal{O}_X$ 为定义
$Z$ 上既约诱导闭子概形结构的理想层，参见《概形》引理
\ref{schemes-lemma-reduced-closed-subscheme}。
考虑子层
$\mathcal{G}'_n = \mathcal{I}^n\mathcal{F}$ 以及商层
$\mathcal{G}_n = \mathcal{F}/\mathcal{I}^n\mathcal{F}$。
对每个 $n$，都有短正合列
$$
0 \to \mathcal{G}'_n \to \mathcal{F} \to \mathcal{G}_n \to 0
$$
对每个点 $x$（其属于 $Z' \setminus Z$），有
$\mathcal{I}_x = \mathcal{O}_{X, x}$，从而
$\mathcal{G}_{n, x} = 0$。故
$\text{Supp}(\mathcal{G}_n) \subset Z$。注意
$X \setminus Z'$ 是 Noether 概形。因此由引理
\ref{coherent-lemma-power-ideal-kills-sheaf}，存在 $n$ 使得
$\mathcal{G}'_n|_{X \setminus Z'} =
\mathcal{I}^n\mathcal{F}|_{X \setminus Z'} = 0$。
对这样的 $n$，有
$\text{Supp}(\mathcal{G}'_n) \subset Z'$。
因此取 $\mathcal{G}' = \mathcal{G}'_n$ 且
$\mathcal{G} = \mathcal{G}_n$ 即可。
\end{proof}

\begin{lemma}
\label{coherent-lemma-prepare-filter-irreducible}
设 $X$ 为 Noether 概形。
设 $i : Z \to X$ 为整的闭子概形。
设 $\xi \in Z$ 为其泛点。
设 $\mathcal{F}$ 为 $X$ 上的凝聚层。
假设 $\mathcal{F}_\xi$ 被 $\mathfrak m_\xi$ 零化。
则存在整数 $r \geq 0$、凝聚理想层
$\mathcal{I} \subset \mathcal{O}_Z$ 以及凝聚层的单射
$$
i_*\left(\mathcal{I}^{\oplus r}\right) \to \mathcal{F}
$$
它在 $\xi$ 的一个邻域上是同构。
\end{lemma}

\begin{proof}
设 $\mathcal{J} \subset \mathcal{O}_X$ 为 $Z$ 的理想层。
设 $\mathcal{F}' \subset \mathcal{F}$ 为
$\mathcal{F}$ 中被 $\mathcal{J}$ 零化的局部截面所成的子层。
由《性质》引理
\ref{properties-lemma-sections-annihilated-by-ideal}，它是准凝聚层。
此外，$\mathcal{F}'_\xi = \mathcal{F}_\xi$，因为
$\mathcal{J}_\xi = \mathfrak m_\xi$，并可应用《性质》引理
\ref{properties-lemma-sections-annihilated-by-ideal} 的第 (3) 部分。
由引理 \ref{coherent-lemma-local-isomorphism} 可知，
$\mathcal{F}' \to \mathcal{F}$ 在 $\xi$ 的一个邻域上诱导同构。
因此可将 $\mathcal{F}$ 替换为 $\mathcal{F}'$，并假设
$\mathcal{F}$ 被 $\mathcal{J}$ 零化。

\medskip\noindent
假设 $\mathcal{J}\mathcal{F} = 0$。由引理
\ref{coherent-lemma-i-star-equivalence}，可写成
$\mathcal{F} = i_*\mathcal{G}$，其中 $\mathcal{G}$ 是
$Z$ 上的某个凝聚层。
假设可以找到态射
$\mathcal{I}^{\oplus r} \to \mathcal{G}$，它在泛点
$\xi$（即 $Z$ 的泛点）的一个邻域上是同构。于是应用 $i_*$
（它左正合）便得到引理结论。因此已经归结为
$X = Z$ 的情形。

\medskip\noindent
假设 $Z = X$ 是泛点为 $\xi$ 的整 Noether 概形。
注意，此时 $\mathcal{O}_{X, \xi} = \kappa(\xi)$ 是 $X$ 的函数域。
由于 $\mathcal{F}_\xi$ 是有限 $\mathcal{O}_\xi$-模，
可知 $r = \dim_{\kappa(\xi)} \mathcal{F}_\xi$ 有限。
因此层 $\mathcal{O}_X^{\oplus r}$ 与 $\mathcal{F}$
在 $\xi$ 处的层茎同构。
由引理 \ref{coherent-lemma-map-stalks-local-map}，存在非空开集
$U \subset X$ 及态射
$\psi : \mathcal{O}_X^{\oplus r}|_U \to \mathcal{F}|_U$，
它在 $\xi$ 处是同构，进而由引理
\ref{coherent-lemma-local-isomorphism} 可知，它在 $\xi$ 的一个邻域上是同构。
由《概形》引理
\ref{schemes-lemma-reduced-closed-subscheme}，存在准凝聚理想层
$\mathcal{I} \subset \mathcal{O}_X$，其相伴闭子概形
$Y \subset X$ 是 $U$ 的补集。
由引理 \ref{coherent-lemma-homs-over-open}，存在 $n \geq 0$ 及态射
$\mathcal{I}^n(\mathcal{O}_X^{\oplus r}) \to \mathcal{F}$，
它恢复 $\psi$（在 $U$ 上）。由于
$\mathcal{I}^n(\mathcal{O}_X^{\oplus r}) = (\mathcal{I}^n)^{\oplus r}$，
便得到引理所述的映射。该映射是单射，因为 $X$ 是整的，
而且它在 $X$ 的泛点处是单射（简单证明略）。
\end{proof}

\begin{lemma}
\label{coherent-lemma-coherent-filter}
设 $X$ 为 Noether 概形。
设 $\mathcal{F}$ 为 $X$ 上的凝聚层。
存在由凝聚子层组成的滤过
$$
0 = \mathcal{F}_0 \subset \mathcal{F}_1 \subset
\ldots \subset \mathcal{F}_m = \mathcal{F}
$$
使得对每个 $j = 1, \ldots, m$，都存在整的闭子概形
$Z_j \subset X$ 及非零凝聚理想层
$\mathcal{I}_j \subset \mathcal{O}_{Z_j}$，满足
$$
\mathcal{F}_j/\mathcal{F}_{j - 1}
\cong (Z_j \to X)_* \mathcal{I}_j
$$
\end{lemma}

\begin{proof}
考虑集合
$$
\mathcal{T} =
\left\{
\begin{matrix}
Z \subset X
\text{ closed such that there exists a coherent sheaf }
\mathcal{F} \\
\text{ 其中 }
\text{Supp}(\mathcal{F}) = Z
\text{ for which the lemma is wrong}
\end{matrix}
\right\}
$$
我们要证明 $\mathcal{T}$ 为空。若非如此，由于 $X$ 是 Noether 的，
可选取极小元 $Z \in \mathcal{T}$。这意味着存在
凝聚层 $\mathcal{F}$，它在 $X$ 上且支集为 $Z$，
而引理对它不成立。
显然 $Z \not = \emptyset$，因为支集为空的层只有零层，
而引理对零层成立（取 $m = 0$）。

\medskip\noindent
若 $Z$ 不可约，则可写成 $Z = Z_1 \cup Z_2$，
其中 $Z_1, Z_2$ 闭且严格小于 $Z$。
于是应用引理 \ref{coherent-lemma-prepare-filter-support}，得到凝聚层的短正合列
$$
0 \to
\mathcal{G}_1 \to
\mathcal{F} \to
\mathcal{G}_2 \to 0
$$
满足 $\text{Supp}(\mathcal{G}_i) \subset Z_i$。
由 $Z$ 的极小性，每个 $\mathcal{G}_i$
都有引理陈述中的滤过。考察 $\mathcal{F}$ 上的诱导滤过便得到矛盾。
因此 $Z$ 必不可约。

\medskip\noindent
假设 $Z$ 不可约。设 $\mathcal{J}$ 为定义 $Z$
的既约诱导闭子概形结构的理想层，参见《概形》引理
\ref{schemes-lemma-reduced-closed-subscheme}。
由引理 \ref{coherent-lemma-power-ideal-kills-sheaf} 可知，存在
$n \geq 0$ 使得 $\mathcal{J}^n\mathcal{F} = 0$。
因而得到滤过
$$
0 = \mathcal{J}^n\mathcal{F} \subset \mathcal{J}^{n - 1}\mathcal{F}
\subset \ldots \subset \mathcal{J}\mathcal{F} \subset \mathcal{F}
$$
其每个相继子商都被 $\mathcal{J}$ 零化。
因此，若这些子商各自都有引理陈述中的滤过，
则 $\mathcal{F}$ 也有。换言之，可以假设
$\mathcal{J}$ 确实零化 $\mathcal{F}$。

\medskip\noindent
在 $Z$ 可约且 $\mathcal{J}\mathcal{F} = 0$ 的情形，
可应用引理 \ref{coherent-lemma-prepare-filter-irreducible}。
这给出短正合列
$$
0 \to
i_*(\mathcal{I}^{\oplus r}) \to
\mathcal{F} \to
\mathcal{Q} \to 0
$$
其中 $\mathcal{Q}$ 定义为该商。
由刚引用的引理，$\mathcal{Q}$ 在 $\xi$ 的一个邻域上为零，
故 $\mathcal{Q}$ 的支集严格小于 $Z$。于是
$\mathcal{Q}$ 由 $Z$ 的极小性具有所需类型的滤过。
但这样 $\mathcal{F}$ 显然也有，得到最终矛盾。
\end{proof}

\begin{lemma}
\label{coherent-lemma-property-initial}
设 $X$ 为 Noether 概形。
设 $\mathcal{P}$ 为 $X$ 上凝聚层的一种性质。假设：
\begin{enumerate}
\item 对凝聚层的任意短正合列
$$
0 \to \mathcal{F}_1 \to \mathcal{F} \to \mathcal{F}_2 \to 0
$$
若 $\mathcal{F}_i$（$i = 1, 2$）具有性质
$\mathcal{P}$，则 $\mathcal{F}$ 也具有该性质。
\item 对每个整的闭子概形 $Z \subset X$ 以及每个准凝聚理想层
$\mathcal{I} \subset \mathcal{O}_Z$，性质 $\mathcal{P}$
对 $i_*\mathcal{I}$ 成立。
\end{enumerate}
则性质 $\mathcal{P}$ 对 $X$ 上的每个凝聚层均成立。
\end{lemma}

\begin{proof}
先注意，若 $\mathcal{F}$ 是凝聚层，并有凝聚子层的滤过
$$
0 = \mathcal{F}_0 \subset \mathcal{F}_1 \subset
\ldots \subset \mathcal{F}_m = \mathcal{F}
$$
使每个 $\mathcal{F}_i/\mathcal{F}_{i - 1}$
都具有性质 $\mathcal{P}$，则 $\mathcal{F}$ 也具有该性质。
这由 $\mathcal{P}$ 的性质 (1) 推出。
另一方面，由引理 \ref{coherent-lemma-coherent-filter}，
可将任意 $\mathcal{F}$ 滤过，使相继子商具有 (2) 中的形式。
故引理成立。
\end{proof}

\begin{lemma}
\label{coherent-lemma-property-irreducible}
设 $X$ 为 Noether 概形。设 $Z_0 \subset X$ 为泛点
$\xi$ 的不可约闭子集。设 $\mathcal{P}$ 为
$X$ 上支集包含于 $Z_0$ 的凝聚层的一种性质，并且：
\begin{enumerate}
\item 对凝聚层的任意短正合列，若三项中有两项具有性质
$\mathcal{P}$，则第三项也具有该性质。
\item 对每个整的闭子概形 $Z \subset Z_0 \subset X$，
$Z \not = Z_0$，以及每个准凝聚理想层
$\mathcal{I} \subset \mathcal{O}_Z$，性质
$\mathcal{P}$ 对 $(Z \to X)_*\mathcal{I}$ 成立。
\item 存在某个凝聚层 $\mathcal{G}$，它在 $X$ 上且满足：
\begin{enumerate}
\item $\text{Supp}(\mathcal{G}) = Z_0$；
\item $\mathcal{G}_\xi$ 被 $\mathfrak m_\xi$ 零化；
\item $\dim_{\kappa(\xi)} \mathcal{G}_\xi = 1$；并且
\item 性质 $\mathcal{P}$ 对 $\mathcal{G}$ 成立。
\end{enumerate}
\end{enumerate}
则性质 $\mathcal{P}$ 对每个凝聚层 $\mathcal{F}$ 均成立，
只要该层在 $X$ 上且其支集包含于 $Z_0$。
\end{lemma}

\begin{proof}
先注意，若 $\mathcal{F}$ 是支集包含于 $Z_0$ 的凝聚层，
并有凝聚子层的滤过
$$
0 = \mathcal{F}_0 \subset \mathcal{F}_1 \subset
\ldots \subset \mathcal{F}_m = \mathcal{F}
$$
使每个 $\mathcal{F}_i/\mathcal{F}_{i - 1}$
都具有性质 $\mathcal{P}$，则 $\mathcal{F}$ 也具有该性质。
或者，若 $\mathcal{F}$ 具有性质 $\mathcal{P}$，且除一个之外的所有
$\mathcal{F}_i/\mathcal{F}_{i - 1}$ 都具有性质
$\mathcal{P}$，则最后一个也具有该性质。这由假设 (1) 推出。

\medskip\noindent
作为第一个应用，可知支集严格包含于 $Z_0$ 的任意凝聚层
都具有性质 $\mathcal{P}$。事实上，这样的层有一个滤过
（参见引理 \ref{coherent-lemma-coherent-filter}），其子商依 (2)
均具有性质 $\mathcal{P}$。

\medskip\noindent
设 $\mathcal{G}$ 如 (3) 所述。由引理
\ref{coherent-lemma-prepare-filter-irreducible}，存在理想层
$\mathcal{I}$（定义在 $Z_0$ 上）、整数 $r \geq 1$ 以及短正合列
$$
0 \to
\left((Z_0 \to X)_*\mathcal{I}\right)^{\oplus r} \to
\mathcal{G} \to
\mathcal{Q} \to 0
$$
其中 $\mathcal{Q}$ 的支集严格包含于 $Z_0$。
由 (3)(c) 可知 $r = 1$。由于 $\mathcal{Q}$
也具有性质 $\mathcal{P}$，可知
$(Z_0 \to X)_*\mathcal{I}$ 具有性质 $\mathcal{P}$。

\medskip\noindent
接着，设 $\mathcal{I}' \not = 0$ 是 $Z_0$
上的另一个准凝聚理想层。考虑交
$\mathcal{I}'' = \mathcal{I}' \cap \mathcal{I}$，得到两个短正合列
$$
0 \to
(Z_0 \to X)_*\mathcal{I}'' \to
(Z_0 \to X)_*\mathcal{I} \to
\mathcal{Q} \to 0
$$
和
$$
0 \to
(Z_0 \to X)_*\mathcal{I}'' \to
(Z_0 \to X)_*\mathcal{I}' \to
\mathcal{Q}' \to 0.
$$
注意，凝聚层 $\mathcal{Q}$ 与
$\mathcal{Q}'$ 的支集都严格包含于 $Z_0$。
因此 $\mathcal{Q}$ 与 $\mathcal{Q}'$ 均具有性质
$\mathcal{P}$（见上）。于是利用 (1) 可知
$(Z_0 \to X)_*\mathcal{I}''$ 与
$(Z_0 \to X)_*\mathcal{I}'$ 也都具有 $\mathcal{P}$。

\medskip\noindent
证明的最后一步是注意到，任意凝聚层 $\mathcal{F}$
只要在 $X$ 上且支集包含于 $Z_0$，便都有一个滤过（再次参见引理
\ref{coherent-lemma-coherent-filter}），依刚才所述，其子商均具有性质
$\mathcal{P}$。
\end{proof}

\begin{lemma}
\label{coherent-lemma-property}
设 $X$ 为 Noether 概形。
设 $\mathcal{P}$ 为 $X$ 上凝聚层的一种性质，并且：
\begin{enumerate}
\item 对凝聚层的任意短正合列，若三项中有两项具有性质
$\mathcal{P}$，则第三项也具有该性质。
\item 对每个整的闭子概形 $Z \subset X$（其泛点为
$\xi$），存在某个凝聚层 $\mathcal{G}$，满足：
\begin{enumerate}
\item $\text{Supp}(\mathcal{G}) = Z$；
\item $\mathcal{G}_\xi$ 被 $\mathfrak m_\xi$ 零化；
\item $\dim_{\kappa(\xi)} \mathcal{G}_\xi = 1$；并且
\item 性质 $\mathcal{P}$ 对 $\mathcal{G}$ 成立。
\end{enumerate}
\end{enumerate}
则性质 $\mathcal{P}$ 对 $X$ 上的每个凝聚层均成立。
\end{lemma}

\begin{proof}
依引理 \ref{coherent-lemma-property-initial}，只需证明：
对所有整的闭子概形 $Z \subset X$ 及所有准凝聚理想层
$\mathcal{I} \subset \mathcal{O}_Z$，性质 $\mathcal{P}$
对 $(Z \to X)_*\mathcal{I}$ 成立。若非如此，则由于
$X$ 是 Noether 的，存在极小的整闭子概形
$Z_0 \subset X$，使得性质 $\mathcal{P}$ 对
$(Z_0 \to X)_*\mathcal{I}_0$ 不成立，其中
$\mathcal{I}_0 \subset \mathcal{O}_{Z_0}$ 是某个准凝聚理想层；
但性质 $\mathcal{P}$ 对 $(Z \to X)_*\mathcal{I}$
成立，只要 $Z \subset Z_0$ 是整的闭子概形、
$Z \not = Z_0$ 且
$\mathcal{I} \subset \mathcal{O}_Z$ 是准凝聚理想层。
由第 (2) 部分，存在 $\mathcal{G}$（对应于 $Z_0$）；
依引理 \ref{coherent-lemma-property-irreducible}，这种情况不可能发生。
\end{proof}

\begin{lemma}
\label{coherent-lemma-property-irreducible-higher-rank-cohomological}
设 $X$ 为 Noether 概形。设 $Z_0 \subset X$ 为泛点
$\xi$ 的不可约闭子集。设 $\mathcal{P}$ 为 $X$
上凝聚层的一种性质，并且：
\begin{enumerate}
\item 对凝聚层的任意短正合列
$$
0 \to \mathcal{F}_1 \to \mathcal{F} \to \mathcal{F}_2 \to 0
$$
若 $\mathcal{F}_i$（$i = 1, 2$）具有性质
$\mathcal{P}$，则 $\mathcal{F}$ 也具有该性质。
\item 若性质 $\mathcal{P}$ 对 $\mathcal{F}^{\oplus r}$
成立，其中 $r \geq 1$，则它也对 $\mathcal{F}$ 成立。
\item 对每个整的闭子概形 $Z \subset Z_0 \subset X$、
$Z \not = Z_0$ 及每个准凝聚理想层
$\mathcal{I} \subset \mathcal{O}_Z$，性质
$\mathcal{P}$ 对 $(Z \to X)_*\mathcal{I}$ 成立。
\item 存在某个凝聚层 $\mathcal{G}$，满足：
\begin{enumerate}
\item $\text{Supp}(\mathcal{G}) = Z_0$；
\item $\mathcal{G}_\xi$ 被 $\mathfrak m_\xi$ 零化；并且
\item 对每个准凝聚理想层
$\mathcal{J} \subset \mathcal{O}_X$，若
$\mathcal{J}_\xi = \mathcal{O}_{X, \xi}$，则存在准凝聚子层
$\mathcal{G}' \subset \mathcal{J}\mathcal{G}$，使得
$\mathcal{G}'_\xi = \mathcal{G}_\xi$，且性质
$\mathcal{P}$ 对 $\mathcal{G}'$ 成立。
\end{enumerate}
\end{enumerate}
则性质 $\mathcal{P}$ 对每个凝聚层 $\mathcal{F}$ 均成立，
只要该层在 $X$ 上且其支集包含于 $Z_0$。
\end{lemma}

\begin{proof}
注意，若 $\mathcal{F}$ 是凝聚层，并有凝聚子层的滤过
$$
0 = \mathcal{F}_0 \subset \mathcal{F}_1 \subset
\ldots \subset \mathcal{F}_m = \mathcal{F}
$$
使每个 $\mathcal{F}_i/\mathcal{F}_{i - 1}$
都具有性质 $\mathcal{P}$，则 $\mathcal{F}$ 也具有该性质。
这由假设 (1) 推出。

\medskip\noindent
作为第一个应用，可知支集严格包含于 $Z_0$ 的任意凝聚层
都具有性质 $\mathcal{P}$。事实上，这样的层有一个滤过
（参见引理 \ref{coherent-lemma-coherent-filter}），其子商依 (3)
均具有性质 $\mathcal{P}$。

\medskip\noindent
以 $i : Z_0 \to X$ 表示该闭浸入。
考虑 (4) 中的凝聚层 $\mathcal{G}$。
由引理 \ref{coherent-lemma-prepare-filter-irreducible}，
存在理想层 $\mathcal{I}$（定义在 $Z_0$ 上）以及短正合列
$$
0 \to
i_*\mathcal{I}^{\oplus r} \to
\mathcal{G} \to
\mathcal{Q} \to 0
$$
其中 $\mathcal{Q}$ 的支集严格包含于 $Z_0$。
特别地，$r > 0$ 且 $\mathcal{I}$ 非零，因为
$\mathcal{G}$ 的支集等于 $Z_0$。
任取非零准凝聚理想层 $\mathcal{I}' \subset \mathcal{I}$；
它定义在 $Z_0$ 上且包含于 $\mathcal{I}$。于是也得到短正合列
$$
0 \to
i_*(\mathcal{I}')^{\oplus r} \to
\mathcal{G} \to
\mathcal{Q}' \to 0
$$
其中 $\mathcal{Q}'$ 的支集真包含于 $Z_0$。
设 $\mathcal{J} \subset \mathcal{O}_X$ 为定义
$\mathcal{Q}'$ 的支集的准凝聚理想层（例如，对应于
$\mathcal{Q}'$ 的支集上的既约诱导闭子概形结构的理想）。
则 $\mathcal{J}_\xi = \mathcal{O}_{X, \xi}$。由引理
\ref{coherent-lemma-power-ideal-kills-sheaf} 可知，
$\mathcal{J}^n\mathcal{Q}' = 0$ 对某个 $n$ 成立。
因此 $\mathcal{J}^n\mathcal{G} \subset i_*(\mathcal{I}')^{\oplus r}$。
由引理的假设 (4)(c)，可知存在准凝聚子层
$\mathcal{G}' \subset \mathcal{J}^n\mathcal{G}$，满足
$\mathcal{G}'_\xi = \mathcal{G}_\xi$，且性质
$\mathcal{P}$ 成立。
从而得到短正合列
$$
0 \to \mathcal{G}' \to
i_*(\mathcal{I}')^{\oplus r} \to
\mathcal{Q}'' \to 0
$$
其中 $\mathcal{Q}''$ 的支集真包含于 $Z_0$。
因此由开头的说明和引理的性质 (1)，可知
$i_*(\mathcal{I}')^{\oplus r}$ 满足 $\mathcal{P}$。
再由 (2) 可知 $i_*\mathcal{I}'$ 满足 $\mathcal{P}$。
最后，对任意准凝聚理想层
$\mathcal{I}'' \subset \mathcal{O}_{Z_0}$，令
$\mathcal{I}' = \mathcal{I}'' \cap \mathcal{I}$，便得到短正合列
$$
0 \to
i_*(\mathcal{I}') \to
i_*(\mathcal{I}'') \to
\mathcal{Q}''' \to 0
$$
其中 $\mathcal{Q}'''$ 的支集真包含于 $Z_0$。
故性质 $\mathcal{P}$ 对 $i_*\mathcal{I}''$ 也成立。

\medskip\noindent
证明的最后一步是注意到，任意凝聚层 $\mathcal{F}$
只要在 $X$ 上且支集包含于 $Z_0$，便都有一个滤过（再次参见引理
\ref{coherent-lemma-coherent-filter}），依刚才所述，其子商均具有性质
$\mathcal{P}$。
\end{proof}

\begin{lemma}
\label{coherent-lemma-property-higher-rank-cohomological}
设 $X$ 为 Noether 概形。
设 $\mathcal{P}$ 为 $X$ 上凝聚层的一种性质，并且：
\begin{enumerate}
\item 对凝聚层的任意短正合列
$$
0 \to \mathcal{F}_1 \to \mathcal{F} \to \mathcal{F}_2 \to 0
$$
若 $\mathcal{F}_i$（$i = 1, 2$）具有性质
$\mathcal{P}$，则 $\mathcal{F}$ 也具有该性质。
\item 若性质 $\mathcal{P}$ 对 $\mathcal{F}^{\oplus r}$
成立，其中 $r \geq 1$，则它也对 $\mathcal{F}$ 成立。
\item 对每个整的闭子概形 $Z \subset X$（其泛点为
$\xi$），存在某个凝聚层 $\mathcal{G}$，满足：
\begin{enumerate}
\item $\text{Supp}(\mathcal{G}) = Z$；
\item $\mathcal{G}_\xi$ 被 $\mathfrak m_\xi$ 零化；并且
\item 对每个准凝聚理想层
$\mathcal{J} \subset \mathcal{O}_X$，若
$\mathcal{J}_\xi = \mathcal{O}_{X, \xi}$，则存在准凝聚子层
$\mathcal{G}' \subset \mathcal{J}\mathcal{G}$，使得
$\mathcal{G}'_\xi = \mathcal{G}_\xi$，且性质
$\mathcal{P}$ 对 $\mathcal{G}'$ 成立。
\end{enumerate}
\end{enumerate}
则性质 $\mathcal{P}$ 对 $X$ 上的每个凝聚层均成立。
\end{lemma}

\begin{proof}
这由引理
\ref{coherent-lemma-property-irreducible-higher-rank-cohomological}
推出，其方式与引理 \ref{coherent-lemma-property} 由引理
\ref{coherent-lemma-property-irreducible} 推出的方式完全相同。
\end{proof}

\section{有限态射与仿射性}
\label{coherent-section-finite-affine}

\noindent
本节利用前几节的结果证明：Noether 仿射概形在有限态射下的像
是仿射的。稍后将看到，该结果在更一般的情形下成立
（参见《极限》引理 \ref{limits-lemma-affine} 及命题
\ref{limits-proposition-affine}）。

\begin{lemma}
\label{coherent-lemma-finite-morphism-Noetherian}
设 $f : Y \to X$ 为概形态射。
假设 $f$ 有限且满射，并且 $X$ 局部 Noether。
设 $Z \subset X$ 为泛点 $\xi$ 的整闭子概形。
则存在凝聚层 $\mathcal{F}$，它在 $Y$ 上，并使得
$f_*\mathcal{F}$ 的支集等于 $Z$，且
$(f_*\mathcal{F})_\xi$ 被 $\mathfrak m_\xi$ 零化。
\end{lemma}

\begin{proof}
注意，由《态射》引理
\ref{morphisms-lemma-finite-type-noetherian}，
$Y$ 局部 Noether。
由于 $f$ 满射，纤维 $Y_\xi$ 非空。
选取 $\xi' \in Y$，使其映到 $\xi$。令
$Z' = \overline{\{\xi'\}}$。
可将 $Z' \subset Y$ 视为既约闭子概形，参见《概形》引理
\ref{schemes-lemma-reduced-closed-subscheme}。
因此层 $\mathcal{F} = (Z' \to Y)_*\mathcal{O}_{Z'}$
是 $Y$ 上的凝聚层（参见引理
\ref{coherent-lemma-finite-pushforward-coherent}）。
考察交换图
$$
\xymatrix{
Z' \ar[r]_{i'} \ar[d]_{f'} &
Y \ar[d]^f \\
Z \ar[r]^i &
X
}
$$
可见 $f_*\mathcal{F} = i_*f'_*\mathcal{O}_{Z'}$。
因此 $f_*\mathcal{F}$ 在 $\xi$ 处的层茎就是
$f'_*\mathcal{O}_{Z'}$ 在 $\xi$ 处的层茎。注意，由于
$Z'$ 是泛点为 $\xi'$ 的整概形，$\xi'$ 是 $Z'$
中位于 $\xi$ 上方的唯一一点，参见《代数》引理
\ref{algebra-lemma-finite-is-integral} 与
\ref{algebra-lemma-integral-no-inclusion}。
故 $f'_*\mathcal{O}_{Z'}$ 在 $\xi$ 处的层茎等于
$\mathcal{O}_{Z', \xi'} = \kappa(\xi')$。特别地，
$f_*\mathcal{F}$ 在 $\xi$ 处的层茎非零。
结合 $f_*\mathcal{F}$ 具有
$i_*f'_*(\text{something})$ 的形式这一事实，即得引理。
\end{proof}

\begin{lemma}
\label{coherent-lemma-affine-morphism-projection-ideal}
设 $f : Y \to X$ 为概形态射。
设 $\mathcal{F}$ 为 $Y$ 上的准凝聚层。
设 $\mathcal{I}$ 为 $X$ 上的准凝聚理想层。
若态射 $f$ 仿射，则
$\mathcal{I}f_*\mathcal{F} = f_*(f^{-1}\mathcal{I}\mathcal{F})$。
\end{lemma}

\begin{proof}
该记号的含义如下。由于 $f^{-1}$ 是正合函子，
$f^{-1}\mathcal{I}$ 是 $f^{-1}\mathcal{O}_X$ 的理想层。
通过映射
$f^\sharp : f^{-1}\mathcal{O}_X \to \mathcal{O}_Y$，
它作用在 $\mathcal{F}$ 上。于是
$f^{-1}\mathcal{I}\mathcal{F}$ 是由形如 $as$
的局部截面的和所生成的子层，其中 $a$ 是
$f^{-1}\mathcal{I}$ 的局部截面，而 $s$ 是
$\mathcal{F}$ 的局部截面。它是 $\mathcal{O}_Y$-模
$\mathcal{F}$ 的准凝聚子模，因为它也是自然映射
$f^*\mathcal{I} \otimes_{\mathcal{O}_Y} \mathcal{F} \to \mathcal{F}$
的像。

\medskip\noindent
说明这些之后，证明是直接的。问题是局部的，故可假设
$X$ 仿射。由于 $f$ 仿射，可知 $Y$ 也仿射。
因此可写成
$Y = \Spec(B)$、$X = \Spec(A)$、
$\mathcal{F} = \widetilde{M}$ 以及
$\mathcal{I} = \widetilde{I}$。在此情形，引理断言归结为
$$
I(M_A) = ((IB)M)_A
$$
其中 $M_A$ 表示相应的 $A$-模；它由 $B$-模 $M$ 得到。
\end{proof}

\begin{lemma}
\label{coherent-lemma-image-affine-finite-morphism-affine-Noetherian}
设 $f : Y \to X$ 为概形态射。假设：
\begin{enumerate}
\item $f$ 有限；
\item $f$ 满射；
\item $Y$ 仿射；并且
\item $X$ Noether。
\end{enumerate}
则 $X$ 仿射。
\end{lemma}

\begin{proof}
我们将证明，在引理的假设下，对任意凝聚
$\mathcal{O}_X$-模 $\mathcal{F}$ 都有
$H^1(X, \mathcal{F}) = 0$。特别地，这蕴含对
$H^1(X, \mathcal{I}) = 0$ 对
$\mathcal{O}_X$ 的每个准凝聚理想层成立。于是由引理
\ref{coherent-lemma-quasi-compact-h1-zero-covering} 或引理
\ref{coherent-lemma-quasi-separated-h1-zero-covering} 可知 $X$ 仿射。

\medskip\noindent
设 $\mathcal{P}$ 为凝聚层 $\mathcal{F}$（在 $X$ 上）
的如下性质：
$$
\mathcal{P}(\mathcal{F}) \Leftrightarrow H^1(X, \mathcal{F}) = 0.
$$
我们将应用引理
\ref{coherent-lemma-property-higher-rank-cohomological}。
因此必须对 $\mathcal{P}$ 验证该引理的 (1)、(2)、(3)。
性质 (1) 由层的短正合列所伴随的上同调长正合列推出。
性质 (2) 由 $H^1(X, -)$ 是加性函子推出。
为验证 (3)，设 $Z \subset X$ 为泛点 $\xi$ 的整闭子概形。
设 $\mathcal{F}$ 为 $Y$ 上的凝聚层，使得
$f_*\mathcal{F}$ 的支集等于 $Z$，且
$(f_*\mathcal{F})_\xi$ 被 $\mathfrak m_\xi$ 零化；
参见引理 \ref{coherent-lemma-finite-morphism-Noetherian}。
我们断言取 $\mathcal{G} = f_*\mathcal{F}$ 即可。
只需验证引理
\ref{coherent-lemma-property-higher-rank-cohomological} 的 (3)(c)。
因此假设 $\mathcal{J} \subset \mathcal{O}_X$
是满足
$\mathcal{J}_\xi = \mathcal{O}_{X, \xi}$ 的准凝聚理想层。
有限态射是仿射的，故由引理
\ref{coherent-lemma-affine-morphism-projection-ideal} 可知
$\mathcal{J}\mathcal{G} = f_*(f^{-1}\mathcal{J}\mathcal{F})$。
此外，如引理
\ref{coherent-lemma-affine-morphism-projection-ideal} 的证明中所指出，
层 $f^{-1}\mathcal{J}\mathcal{F}$ 是准凝聚
$\mathcal{O}_Y$-模。
由于 $Y$ 仿射，有
$H^1(Y, f^{-1}\mathcal{J}\mathcal{F}) = 0$，
参见引理 \ref{coherent-lemma-quasi-coherent-affine-cohomology-zero}。
由于 $f$ 有限，因而仿射，所以
$$
H^1(X, \mathcal{J}\mathcal{G}) =
H^1(X, f_*(f^{-1}\mathcal{J}\mathcal{F})) =
H^1(Y, f^{-1}\mathcal{J}\mathcal{F}) = 0
$$
这由引理 \ref{coherent-lemma-relative-affine-cohomology} 得到。
因此准凝聚子层
$\mathcal{G}' = \mathcal{J}\mathcal{G}$ 满足
$\mathcal{P}$。这如所要求验证了引理
\ref{coherent-lemma-property-higher-rank-cohomological} 的性质 (3)(c)。
\end{proof}

\section{Proj 上的凝聚层，I}
\label{coherent-section-coherent-proj}

\noindent
本节讨论 $\text{Proj}(A)$ 上的凝聚层，其中 $A$
是由 $A_1$ 在 $A_0$ 上生成的 Noether 分次环。
下一节将讨论 $A$ 不由次数 $1$ 的元素生成时的情形。
首先，我们为 Noether 环上的射影空间表述一个汇总性结果。

\begin{lemma}
\label{coherent-lemma-coherent-projective}
设 $R$ 为 Noether 环。
设 $n \geq 0$ 为整数。
对每个凝聚层 $\mathcal{F}$（在 $\mathbf{P}^n_R$ 上），
有下列结论：
\begin{enumerate}
\item 存在 $r \geq 0$、
$d_1, \ldots, d_r \in \mathbf{Z}$ 以及满射
$$
\bigoplus\nolimits_{j = 1, \ldots, r}
\mathcal{O}_{\mathbf{P}^n_R}(d_j)
\longrightarrow
\mathcal{F}.
$$
\item 有 $H^i(\mathbf{P}^n_R, \mathcal{F}) = 0$，
除非 $0 \leq i \leq n$。
\item 对任意 $i$，上同调群
$H^i(\mathbf{P}^n_R, \mathcal{F})$ 是有限 $R$-模。
\item 若 $i > 0$，则
$H^i(\mathbf{P}^n_R, \mathcal{F}(d)) = 0$
对所有充分大的 $d$ 成立。
\item 对任意 $k \in \mathbf{Z}$，分次
$R[T_0, \ldots, T_n]$-模
$$
\bigoplus\nolimits_{d \geq k} H^0(\mathbf{P}^n_R, \mathcal{F}(d))
$$
是有限 $R[T_0, \ldots, T_n]$-模。
\end{enumerate}
\end{lemma}

\begin{proof}
我们将用到 $\mathcal{O}_{\mathbf{P}^n_R}(1)$ 是概形
$\mathbf{P}^n_R$ 上的丰沛可逆层。由于
$\mathbf{P}^n_R$ 被标准仿射开集 $D_{+}(T_i)$ 覆盖，
这直接由定义推出。因此由《性质》命题
\ref{properties-proposition-characterize-ample}，
每个有限型准凝聚 $\mathcal{O}_{\mathbf{P}^n_R}$-模
都是 $\mathcal{O}_{\mathbf{P}^n_R}(1)$ 的张量幂的
某个有限直和的商。另一方面，由引理
\ref{coherent-lemma-coherent-Noetherian}，在 $R$ 上的射影空间中，
凝聚层与有限型准凝聚层是同一回事。故得到 (1)。

\medskip\noindent
射影 $n$-空间 $\mathbf{P}^n_R$ 由 $n + 1$ 个仿射开集覆盖，
即标准开集 $D_{+}(T_i)$（$i = 0, \ldots, n$），参见《构造》引理
\ref{constructions-lemma-standard-covering-projective-space}。
因此，对任意准凝聚层 $\mathcal{F}$（在
$\mathbf{P}^n_R$ 上），有
$H^i(\mathbf{P}^n_R, \mathcal{F}) = 0$（当
$i \geq n + 1$ 时），参见引理
\ref{coherent-lemma-vanishing-nr-affines}。故 (2) 成立。

\medskip\noindent
我们同时为 $\mathbf{P}^n_R$ 上的所有凝聚层证明
(3) 与 (4)，并对 $i$ 作降归纳。由 (2)，结论对
$i \geq n + 1$ 显然成立。假设已知 $i + 1$ 的结论，
现在证明 $i$ 的结论。（若 $i = 0$，则 (4) 为空断言。）
设 $\mathcal{F}$ 为 $\mathbf{P}^n_R$ 上的凝聚层。
选取 (1) 中的满射，并以 $\mathcal{G}$ 表示其核，从而有短正合列
$$
0 \to \mathcal{G} \to
\bigoplus\nolimits_{j = 1, \ldots, r}
\mathcal{O}_{\mathbf{P}^n_R}(d_j)
\to
\mathcal{F} \to 0
$$
由引理 \ref{coherent-lemma-coherent-abelian-Noetherian} 可知
$\mathcal{G}$ 凝聚。上同调长正合列给出正合列
$$
H^i(\mathbf{P}^n_R, \bigoplus\nolimits_{j = 1, \ldots, r}
\mathcal{O}_{\mathbf{P}^n_R}(d_j))
\to
H^i(\mathbf{P}^n_R, \mathcal{F})
\to
H^{i + 1}(\mathbf{P}^n_R, \mathcal{G}).
$$
由归纳假设，右侧 $R$-模有限；由引理
\ref{coherent-lemma-cohomology-projective-space-over-ring}，
左侧 $R$-模有限。由于 $R$ Noether，立即可知
$H^i(\mathbf{P}^n_R, \mathcal{F})$ 是有限 $R$-模。
这证明了断言 (3) 的归纳步骤。
由于 $\mathcal{O}_{\mathbf{P}^n_R}(d)$ 可逆，
可知在 $\mathbf{P}^n_R$ 上作扭转是正合函子
（因为它由与可逆层张量得到，参见《构造》定义
\ref{constructions-definition-twist}）。
这意味着对所有 $d \in \mathbf{Z}$，序列
$$
0 \to \mathcal{G}(d) \to
\bigoplus\nolimits_{j = 1, \ldots, r}
\mathcal{O}_{\mathbf{P}^n_R}(d_j + d)
\to
\mathcal{F}(d) \to 0
$$
短正合。所得上同调序列为
$$
H^i(\mathbf{P}^n_R, \bigoplus\nolimits_{j = 1, \ldots, r}
\mathcal{O}_{\mathbf{P}^n_R}(d_j + d))
\to
H^i(\mathbf{P}^n_R, \mathcal{F}(d))
\to
H^{i + 1}(\mathbf{P}^n_R, \mathcal{G}(d)).
$$
由归纳假设，$d \gg 0$ 时右侧的模为零；由引理
\ref{coherent-lemma-cohomology-projective-space-over-ring} 中的计算，
一旦 $d \geq -\min\{d_j\}$ 且 $i \geq 1$，左侧的模即为零。
故得到断言 (4) 的归纳步骤。
这完成 (3) 与 (4) 的证明。

\medskip\noindent
为证明 (5)，注意对所有充分大的 $d$，映射
$$
H^0(\mathbf{P}^n_R, \bigoplus\nolimits_{j = 1, \ldots, r}
\mathcal{O}_{\mathbf{P}^n_R}(d_j + d))
\to
H^0(\mathbf{P}^n_R, \mathcal{F}(d))
$$
是满射，这是由刚刚证明的
$H^1(\mathbf{P}^n_R, \mathcal{G}(d))$ 的消失得到的。
换言之，当 $k$ 充分大时，模
$$
M_k
=
\bigoplus\nolimits_{d \geq k} H^0(\mathbf{P}^n_R, \mathcal{F}(d))
$$
是相应模
$$
N_k
=
\bigoplus\nolimits_{d \geq k} H^0(\mathbf{P}^n_R,
\bigoplus\nolimits_{j = 1, \ldots, r}
\mathcal{O}_{\mathbf{P}^n_R}(d_j + d)
)
$$
的商。当 $k$ 充分小时（例如有
$k < -d_j$ 对所有 $j$ 成立），由第
\ref{coherent-section-cohomology-projective-space} 节中的计算，有
$$
N_k = \bigoplus\nolimits_{j = 1, \ldots, r}
R[T_0, \ldots, T_n](d_j)
$$
特别地，它有限生成。
现在设 $k \in \mathbf{Z}$ 任意。
选取 $k_{-} \ll k \ll k_{+}$。
考虑交换图
$$
\xymatrix{
N_{k_{-}} & N_{k_{+}} \ar[d] \ar[l] \\
M_k & M_{k_{+}} \ar[l]
}
$$
其中竖直箭头是上述满射，而水平箭头是显然的包含映射。
由以上说明，$N_{k_{-}}$ 是有限生成
$R[T_0, \ldots, T_n]$-模。因此 $N_{k_{+}}$
也是有限生成 $R[T_0, \ldots, T_n]$-模，
因为它是有限生成模的子模，而环
$R[T_0, \ldots, T_n]$ 是 Noether 的。
由于竖直箭头满射，$M_{k_{+}}$ 是有限生成
$R[T_0, \ldots, T_n]$-模。商
$M_k/M_{k_{+}}$ 作为 $R$-模是有限的，因为它是有限个有限
$R$-模 $H^0(\mathbf{P}^n_R, \mathcal{F}(d))$
的直和，其中 $k \leq d < k_{+}$。
注意，此处使用了 (3) 中 $i = 0$ 的情形。因此
$M_k/M_{k_{+}}$ 更是有限
$R[T_0, \ldots, T_n]$-模。
换言之，我们把 $M_k$ 夹在两个有限
$R[T_0, \ldots, T_n]$-模之间，结论得证。
\end{proof}

\begin{lemma}
\label{coherent-lemma-coherent-on-proj}
设 $A$ 为分次环，使得 $A_0$ Noether，且 $A$
由 $A_1$ 的有限多个元素在 $A_0$ 上生成。
令 $X = \text{Proj}(A)$。则 $X$ 是 Noether 概形。
设 $\mathcal{F}$ 为凝聚 $\mathcal{O}_X$-模。
\begin{enumerate}
\item 存在 $r \geq 0$、
$d_1, \ldots, d_r \in \mathbf{Z}$ 以及满射
$$
\bigoplus\nolimits_{j = 1, \ldots, r} \mathcal{O}_X(d_j)
\longrightarrow \mathcal{F}.
$$
\item 对任意 $p$，上同调群 $H^p(X, \mathcal{F})$
是有限 $A_0$-模。
\item 若 $p > 0$，则
$H^p(X, \mathcal{F}(d)) = 0$ 对所有充分大的 $d$ 成立。
\item 对任意 $k \in \mathbf{Z}$，分次 $A$-模
$$
\bigoplus\nolimits_{d \geq k} H^0(X, \mathcal{F}(d))
$$
是有限 $A$-模。
\end{enumerate}
\end{lemma}

\begin{proof}
由假设，存在分次 $A_0$-代数的满射
$$
A_0[T_0, \ldots, T_n] \longrightarrow A
$$
其中 $\deg(T_j) = 1$ 对 $j = 0, \ldots, n$ 成立。
由《构造》引理
\ref{constructions-lemma-surjective-graded-rings-generated-degree-1-map-proj}，
这定义一个闭浸入 $i : X \to \mathbf{P}^n_{A_0}$，满足
$i^*\mathcal{O}_{\mathbf{P}^n_{A_0}}(1) = \mathcal{O}_X(1)$。
特别地，$X$ 作为 Noether 概形
$\mathbf{P}^n_{A_0}$ 的闭子概形，是 Noether 的。
我们断言，引理关于 $\mathcal{F}$ 的各结论可由引理
\ref{coherent-lemma-coherent-projective} 关于凝聚层
$i_*\mathcal{F}$（引理 \ref{coherent-lemma-i-star-equivalence}）在
$\mathbf{P}^n_{A_0}$ 上的相应结论推出。例如，由该引理存在满射
$$
\bigoplus\nolimits_{j = 1, \ldots, r} \mathcal{O}_{\mathbf{P}^n_{A_0}}(d_j)
\longrightarrow i_*\mathcal{F}.
$$
由伴随，这对应于映射
$\bigoplus_{j = 1, \ldots, r} \mathcal{O}_X(d_j) \longrightarrow \mathcal{F}$，
它也满射。关于上同调的断言由等式
$H^p(X, \mathcal{F}(d)) = H^p(\mathbf{P}^n_{A_0}, i_*\mathcal{F}(d))$
及引理 \ref{coherent-lemma-relative-affine-cohomology} 推出。
\end{proof}

\begin{lemma}
\label{coherent-lemma-recover-tail-graded-module}
设 $A$ 为分次环，使得 $A_0$ Noether，且 $A$
由 $A_1$ 的有限多个元素在 $A_0$ 上生成。设 $M$
为有限分次 $A$-模。令 $X = \text{Proj}(A)$，并令
$\widetilde{M}$ 为准凝聚 $\mathcal{O}_X$-模；它在
$X$ 上且与 $M$ 相伴。
《构造》引理 \ref{constructions-lemma-apply-modules} 中的映射
$$
M_n \longrightarrow \Gamma(X, \widetilde{M}(n))
$$
对所有充分大的 $n$ 都是同构。
\end{lemma}

\begin{proof}
由于 $M$ 是有限 $A$-模，可知 $\widetilde{M}$
是有限型 $\mathcal{O}_X$-模，即凝聚
$\mathcal{O}_X$-模。
令 $N = \bigoplus_{n \in \mathbf{Z}} \Gamma(X, \widetilde{M}(n))$。
必须证明映射 $M \to N$（分次 $A$-模的映射）
在所有充分大的次数上是同构。
由《性质》引理 \ref{properties-lemma-proj-quasi-coherent}，
存在典范同构 $\widetilde{N} \to \widetilde{M}$，使得诱导映射
$N_n \to N_n = \Gamma(X, \widetilde{M}(n))$
为恒等映射。因此有映射
$\widetilde{M} \to \widetilde{N} \to \widetilde{M}$，
使得对所有 $n$，交换图
$$
\xymatrix{
M_n \ar[d] \ar[r] & N_n \ar[d] \ar@{=}[rd] \\
\Gamma(X, \widetilde{M}(n)) \ar[r] &
\Gamma(X, \widetilde{N}(n)) \ar[r]^{\cong} &
\Gamma(X, \widetilde{M}(n))
}
$$
交换。这意味着复合
$$
M_n \to \Gamma(X, \widetilde{M}(n))
\to \Gamma(X, \widetilde{N}(n)) \to \Gamma(X, \widetilde{M}(n))
$$
等于典范映射
$M_n \to \Gamma(X, \widetilde{M}(n))$。
显然，这蕴含复合
$\widetilde{M} \to \widetilde{N} \to \widetilde{M}$
为恒等映射。因此 $\widetilde{M} \to \widetilde{N}$
是同构。令 $K = \Ker(M \to N)$ 且
$Q = \Coker(M \to N)$。
回忆函子 $M \mapsto \widetilde{M}$ 正合，参见《构造》引理
\ref{constructions-lemma-proj-sheaves}。
因此 $\widetilde{K} = 0$ 且 $\widetilde{Q} = 0$。
回忆 $A$ 是 Noether 环，$M$ 是有限生成 $A$-模，
而 $N$ 是分次 $A$-模，并且由引理
\ref{coherent-lemma-coherent-on-proj} 的最后一部分，
$N' = \bigoplus_{n \geq 0} N_n$ 有限生成。
因此 $K' = \bigoplus_{n \geq 0} K_n$ 及
$Q' = \bigoplus_{n \geq 0} Q_n$ 是有限 $A$-模。
注意 $\widetilde{Q} = \widetilde{Q'}$ 且
$\widetilde{K} = \widetilde{K'}$。
故为完成证明，只需说明：若有限 $A$-模 $K$ 满足
$\widetilde{K} = 0$，则它只有有限多个非零齐次分量
$K_d$，其中 $d \geq 0$。
为此，设 $x_1, \ldots, x_r \in K$ 为齐次生成元，
其次数分别为 $d_1, \ldots, d_r$。
设 $f_1, \ldots, f_n \in A_1$ 为生成
$A$ 的元素（在 $A_0$ 上）。对每个 $i$ 与 $j$，存在
$n_{ij} \geq 0$，使得
$f_i^{n_{ij}} x_j = 0$ 在 $K_{d_j + n_{ij}}$ 中成立：
否则 $x_i/f_i^{d_i} \in K_{(f_i)}$ 将不为零，即
$\widetilde{K}$ 不为零。
于是 $K_d$ 在 $d > \max_j(d_j + \sum_i n_{ij})$ 时
为零，因为 $K_d$ 的每个元素都是若干项之和，
每一项都是 $f_i$ 中的一个单项式乘以某个 $x_j$，
且总次数为 $d$。
\end{proof}

\noindent
设 $A$ 为分次环，使得 $A_0$ Noether，且 $A$
由 $A_1$ 的有限多个元素在 $A_0$ 上生成。回忆
$A_+ = \bigoplus_{n > 0} A_n$ 是无关理想。
设 $M$ 为分次 $A$-模。回忆，称 $M$ 是
$A_+$-幂挠模，是指对所有 $x \in M$ 都存在
$n \geq 1$ 使得 $(A_+)^n x = 0$；参见《代数进阶》定义
\ref{more-algebra-definition-f-power-torsion}。
若 $M$ 有限生成，则这等价于
$M_n = 0$ 对 $n \gg 0$ 成立。有时把 $A_+$-幂挠模简称为挠模。
有时称分次 $A$-模 $M$ 无挠，是指只要
$x \in M$ 满足 $(A_+)^n x = 0$ 且 $n > 0$，
便有 $x = 0$。
以 $\text{Mod}_A$ 表示分次 $A$-模范畴，
以 $\text{Mod}^{fg}_A$ 表示其中有限生成对象组成的满子范畴，
并以 $\text{Mod}^{fg}_{A, torsion}$ 表示由模
$M$ 组成的满子范畴；这些模满足
$M_n = 0$（当 $n \gg 0$）。

\begin{proposition}
\label{coherent-proposition-coherent-modules-on-proj}
设 $A$ 为分次环，使得 $A_0$ Noether，且 $A$
由 $A_1$ 的有限多个元素在 $A_0$ 上生成。
令 $X = \text{Proj}(A)$。函子 $M \mapsto \widetilde M$
诱导等价
$$
\text{Mod}^{fg}_A/\text{Mod}^{fg}_{A, torsion}
\longrightarrow
\textit{Coh}(\mathcal{O}_X)
$$
其拟逆由
$\mathcal{F} \longmapsto \bigoplus_{n \geq 0} \Gamma(X, \mathcal{F}(n))$
给出。
\end{proposition}

\begin{proof}
子范畴 $\text{Mod}^{fg}_{A, torsion}$ 是
$\text{Mod}^{fg}_A$ 的 Serre 子范畴，参见《同调》定义
\ref{homology-definition-serre-subcategory}。
这由上述对象描述显然可见，也可由《代数进阶》引理
\ref{more-algebra-lemma-I-power-torsion} 推出。
因此箭头左侧的商范畴已在《同调》引理
\ref{homology-lemma-serre-subcategory-is-kernel} 中定义。
为定义命题中的函子，只需说明函子
$M \mapsto \widetilde M$ 把挠模送到 $0$。
这是显然的，因为对任意齐次 $f \in A_+$，
模 $M_f$ 为零，因而值 $M_{(f)}$（$\widetilde M$
在 $D_+(f)$ 上的值）也为零。

\medskip\noindent
由引理 \ref{coherent-lemma-coherent-on-proj}，所提出的拟逆有意义。
具体而言，该引理表明
$\mathcal{F} \longmapsto \bigoplus_{n \geq 0} \Gamma(X, \mathcal{F}(n))$
是函子
$\textit{Coh}(\mathcal{O}_X) \to \text{Mod}^{fg}_A$；
可将它与商函子
$\text{Mod}^{fg}_A \to \text{Mod}^{fg}_A/\text{Mod}^{fg}_{A, torsion}$
复合。

\medskip\noindent
由引理 \ref{coherent-lemma-recover-tail-graded-module}，
从左到右再到左的复合同构于恒等函子。

\medskip\noindent
最后，设 $\mathcal{F}$ 为凝聚 $\mathcal{O}_X$-模。
令 $M = \bigoplus_{n \in \mathbf{Z}} \Gamma(X, \mathcal{F}(n))$，
将其视为分次 $A$-模，从而我们的函子把 $\mathcal{F}$ 送到
$M_{\geq 0} = \bigoplus_{n \geq 0} M_n$。
由《性质》引理 \ref{properties-lemma-proj-quasi-coherent}，
典范映射 $\widetilde M \to \mathcal{F}$ 是同构。
由于包含映射
$M_{\geq 0} \to M$ 定义同构
$\widetilde{M_{\geq 0}} \to \widetilde M$，
可知从右到左再到右的复合也同构于恒等函子。
\end{proof}





\section{Proj 上的凝聚层，II}
\label{coherent-section-coherent-proj-general}

\noindent
在本节中，我们讨论 $\text{Proj}(A)$ 上的凝聚层，
其中 $A$ 是 Noether 分次环。多数结果将借助一个技巧，
由前一节关于 $A$ 由次数 $1$ 元生成时的相应结果推出。

\begin{lemma}
\label{coherent-lemma-coherent-on-proj-general}
设 $A$ 为 Noether 分次环。令 $X = \text{Proj}(A)$。则 $X$
是 Noether 概形。设 $\mathcal{F}$ 为凝聚 $\mathcal{O}_X$-模。
\begin{enumerate}
\item 存在 $r \geq 0$、
$d_1, \ldots, d_r \in \mathbf{Z}$ 以及满射
$$
\bigoplus\nolimits_{j = 1, \ldots, r} \mathcal{O}_X(d_j)
\longrightarrow \mathcal{F}.
$$
\item 对任意 $p$，上同调群 $H^p(X, \mathcal{F})$ 是有限
$A_0$-模。
\item 若 $p > 0$，则 $H^p(X, \mathcal{F}(d)) = 0$
对所有充分大的 $d$ 成立。
\item 对任意 $k \in \mathbf{Z}$，分次 $A$-模
$$
\bigoplus\nolimits_{d \geq k} H^0(X, \mathcal{F}(d))
$$
是有限 $A$-模。
\end{enumerate}
\end{lemma}

\begin{proof}
我们把该断言归约到引理
\ref{coherent-lemma-coherent-on-proj} 来证明。
由《代数》引理 \ref{algebra-lemma-graded-Noetherian} 与
\ref{algebra-lemma-S-plus-generated}，环 $A_0$ 是 Noether 的，
并且 $A$ 在 $A_0$ 上由有限多个正次数齐次元
$f_1, \ldots, f_r$ 生成。
令 $d$ 为充分可除的整数。按《代数》第
\ref{algebra-section-graded} 节的记号，令 $A' = A^{(d)}$。
于是 $A'$ 在 $A'_0 = A_0$ 上由次数 $1$ 元生成；参见
《代数》引理 \ref{algebra-lemma-uple-generated-degree-1}。
因此引理 \ref{coherent-lemma-coherent-on-proj} 适用于
$X' = \text{Proj}(A')$。

\medskip\noindent
由《构造》引理 \ref{constructions-lemma-d-uple}，存在概形同构
$i : X \to X'$ 以及同构
$\mathcal{O}_X(nd) \to i^*\mathcal{O}_{X'}(n)$；
它们与映射 $A' \to A$ 以及映射
$A_n \to H^0(X, \mathcal{O}_X(n))$ 和 $A'_n \to H^0(X', \mathcal{O}_{X'}(n))$
相容。因此，引理 \ref{coherent-lemma-coherent-on-proj} 蕴含 $X$
为 Noether 概形，并且 (1) 与 (2) 成立。为证明 (3) 与 (4)，
可以利用如下事实：对任意固定的 $k$、$p$ 与 $q$，都有
$$
\bigoplus\nolimits_{dn + q \geq k} H^p(X, \mathcal{F}(dn + q)) =
\bigoplus\nolimits_{dn + q \geq k} H^p(X', (i_*\mathcal{F}(q))(n)
$$
这是由上述相容性得到的。若 $p > 0$，则由引理
\ref{coherent-lemma-coherent-on-proj}，对充分大的 $k$
（它可依赖于 $q$），右端为零。因为整数模 $d$
只有有限多个同余类，可知 (3) 对 $\mathcal{F}$ 在 $X$
上成立。若 $p = 0$，
则由引理 \ref{coherent-lemma-coherent-on-proj}，右端是有限
$A'$-模。再次利用同余类的有限性，可知
$\bigoplus_{n \geq k} H^0(X, \mathcal{F}(n))$ 也是有限
$A'$-模。由于该 $A'$-模结构来自 $A$-模结构
（由上述相容性），可知它也是有限 $A$-模。
\end{proof}

\begin{lemma}
\label{coherent-lemma-recover-tail-graded-module-general}
设 $A$ 为 Noether 分次环，并令 $d$ 为 $A$ 在 $A_0$
上生成元的最小公倍数。设 $M$ 为有限分次 $A$-模。
令 $X = \text{Proj}(A)$，并令 $\widetilde{M}$ 为准凝聚
$\mathcal{O}_X$-模；它在 $X$ 上且与 $M$ 相伴。
设 $k \in \mathbf{Z}$。
\begin{enumerate}
\item $N' = \bigoplus_{n \geq k} H^0(X, \widetilde{M(n)})$
是有限 $A$-模；
\item $N = \bigoplus_{n \geq k} H^0(X, \widetilde{M}(n))$
是有限 $A$-模；
\item 存在典范映射 $N \to N'$；
\item 若 $k$ 充分小，则存在典范映射 $M \to N'$；
\item 映射 $M_n \to N'_n$ 对 $n \gg 0$ 是同构；
\item $N_n \to N'_n$ 对 $d | n$ 是同构。
\end{enumerate}
\end{lemma}

\begin{proof}
(3) 中的映射 $N \to N'$ 来自《构造》等式
(\ref{constructions-equation-multiply-more-generally})，
取整体截面即可得到。

\medskip\noindent
由《构造》等式
(\ref{constructions-equation-global-sections-more-generally})，
存在分次 $A$-模的映射
$M \to \bigoplus_{n \in \mathbf{Z}} H^0(X, \widetilde{M(n)})$。
若 $M$ 的生成元所在次数均为 $\geq k$，则像包含在子模
$N' \subset \bigoplus_{n \in \mathbf{Z}} H^0(X, \widetilde{M(n)})$
中，从而得到 (4) 中的映射。

\medskip\noindent
由《代数》引理 \ref{algebra-lemma-graded-Noetherian} 与
\ref{algebra-lemma-S-plus-generated}，环 $A_0$ 是 Noether 的，
并且 $A$ 在 $A_0$ 上由有限多个正次数齐次元
$f_1, \ldots, f_r$ 生成。
令 $d = \text{lcm}(\deg(f_i))$。例如，由《构造》引理
\ref{constructions-lemma-where-invertible} 可知 (6) 成立。

\medskip\noindent
因为 $M$ 是有限 $A$-模，可知 $\widetilde{M}$ 是有限型
$\mathcal{O}_X$-模，即凝聚 $\mathcal{O}_X$-模。
因此 (2) 由引理 \ref{coherent-lemma-coherent-on-proj-general} 推出。

\medskip\noindent
我们用一个技巧由 (2) 推出 (1)。对
$q \in \{0, \ldots, d - 1\}$，记
$$
{}^qN = \bigoplus\nolimits_{n + q \geq k} H^0(X, \widetilde{M(q)}(n))
$$
由 (2)，这些都是有限 $A$-模。Noether 环 $A$
在 $A^{(d)} = \bigoplus_{n \geq 0} A_{dn}$ 上有限，
因为它由 $f_i$ 在 $A^{(d)}$ 上生成，且 $f_i^d \in A^{(d)}$。
因此 ${}^qN$ 是有限 $A^{(d)}$-模。
此外，$A^{(d)}$ 是 Noether 的（这由《代数》引理
\ref{algebra-lemma-dehomogenize-finite-type} 推出）。
于是 $A^{(d)}$-子模
${}^qN^{(d)} = \bigoplus_{n \in \mathbf{Z}} {}^qN_{dn}$
是有限 $A^{(d)}$-模。利用同构
$\widetilde{M(dn + q)} = \widetilde{M(q)}(dn)$，可写成
$$
N' =
\bigoplus\nolimits_{q \in \{0, \ldots, d - 1\}}
\bigoplus\nolimits_{dn + q \geq k}
H^0(X, \widetilde{M(q)}(dn)) =
\bigoplus\nolimits_{q \in \{0, \ldots, d - 1\}} {}^qN^{(d)}
$$
因此 $N'$ 在 $A^{(d)}$ 上有限，因而当然也在 $A$ 上有限。
故 (1) 成立。

\medskip\noindent
设 $K$ 为有限 $A$-模，且满足 $\widetilde{K} = 0$。
我们断言，$K_n = 0$ 对 $d|n$ 且 $n \gg 0$ 成立。
如上论证可知，$K^{(d)}$ 是有限 $A^{(d)}$-模。设
$x_1, \ldots, x_m \in K$ 是 $K^{(d)}$ 在 $A^{(d)}$
上的齐次生成元，其所在次数分别为
$d_1, \ldots, d_m$，并且 $d | d_j$。
对每个 $i$ 与 $j$，存在 $n_{ij} \geq 0$，使得
$f_i^{n_{ij}} x_j = 0$ 在 $K_{d_j + n_{ij}}$ 中成立；否则
$x_j/f_i^{d_i/\deg(f_i)} \in K_{(f_i)}$ 将不为零，即
$\widetilde{K}$ 不为零。这里使用了
$\deg(f_i) | d | d_j$ 对所有 $i, j$ 成立。
由此可知，$K_n$ 为零，其中 $n$ 满足 $d | n$ 且
$n > \max_j (d_j + \sum_i n_{ij} \deg(f_i))$，
因为 $K_n$ 的每个元素都是若干项之和，
其中每一项都是 $f_i$ 中的一个单项式乘以某个 $x_j$，
且总次数为 $n$。

\medskip\noindent
为完成证明，必须说明 $M \to N'$ 在所有充分大的次数上
是同构。映射 $N \to N'$ 诱导同构
$\widetilde{N} \to \widetilde{N'}$，因为在仿射开集
$D_+(f_i) = D_+(f_i^d)$ 上，相应的模由性质 (6) 给出同构
$N_{(f_i)} \cong N_{(f_i^d)} \cong N'_{(f_i^d)} \cong N'_{(f_i)}$。
由《性质》引理 \ref{properties-lemma-proj-quasi-coherent}，
存在典范同构 $\widetilde{N} \to \widetilde{M}$。
复合 $\widetilde{N} \to \widetilde{M} \to \widetilde{N'}$
就是上述同构（证明从略；提示：在标准仿射开集上检查）。
因此映射 $M \to N'$ 诱导同构
$\widetilde{M} \to \widetilde{N'}$。
令 $K = \Ker(M \to N')$ 且 $Q = \Coker(M \to N')$。
回忆函子 $M \mapsto \widetilde{M}$ 正合；参见《构造》引理
\ref{constructions-lemma-proj-sheaves}。
因此 $\widetilde{K} = 0$ 且 $\widetilde{Q} = 0$。
由上一段的结果可知，$K_n = 0$ 且 $Q_n = 0$
对 $d | n$ 且 $n \gg 0$ 成立。此时终于可见使用 $N'$ 而不是
$N$ 的优势：函子 $M \leadsto N'$ 与扭转相容
（由构造立即可见）。因此，把 $M$ 替换为 $M(q)$
并重复整个论证，可知 $K_n = 0$ 且 $Q_n = 0$
对 $n \equiv q \bmod d$ 且 $n \gg 0$ 成立。
因为模 $n$ 只有有限多个同余类，证明完成。
\end{proof}

\noindent
设 $A$ 为 Noether 分次环。回忆
$A_+ = \bigoplus_{n > 0} A_n$ 是无关理想。
由《代数》引理 \ref{algebra-lemma-graded-Noetherian} 与
\ref{algebra-lemma-S-plus-generated}，环 $A_0$ 是 Noether 的，
并且 $A$ 在 $A_0$ 上由有限多个正次数齐次元
$f_1, \ldots, f_r$ 生成。
令 $d = \text{lcm}(\deg(f_i))$。设 $M$ 为分次 $A$-模。
在这种情形下，若齐次元 $x \in M$ 满足
$$
(A_+ x)_{nd} = 0\text{ 对所有 }n \gg 0
$$
则称它是{\it 无关的}\footnote{这是非标准记号。}。
若 $x \in M$ 齐次且无关，而 $f \in A$ 齐次，
则 $fx$ 也无关。因此，无关元素的集合生成一个分次子模
$M_{irrelevant} \subset M$。若 $M$ 的每个齐次元都无关，
则称 $M$ 是{\it 无关的}，即 $M_{irrelevant} = M$。
若 $M$ 有限生成，则这等价于
$M_{nd} = 0$ 对 $n \gg 0$ 成立。
以 $\text{Mod}_A$ 表示分次 $A$-模范畴，
以 $\text{Mod}^{fg}_A$ 表示有限生成对象组成的满子范畴，
并以 $\text{Mod}^{fg}_{A, irrelevant}$ 表示无关模组成的满子范畴。

\begin{proposition}
\label{coherent-proposition-coherent-modules-on-proj-general}
设 $A$ 为 Noether 分次环。令 $X = \text{Proj}(A)$。
函子 $M \mapsto \widetilde M$ 诱导等价
$$
\text{Mod}^{fg}_A/\text{Mod}^{fg}_{A, irrelevant}
\longrightarrow
\textit{Coh}(\mathcal{O}_X)
$$
其拟逆由
$\mathcal{F} \longmapsto \bigoplus_{n \geq 0} \Gamma(X, \mathcal{F}(n))$
给出。
\end{proposition}

\begin{proof}
我们敦促读者先阅读 $A$ 由次数 $1$ 元生成的情形；
参见命题 \ref{coherent-proposition-coherent-modules-on-proj}。
设 $f_1, \ldots, f_r \in A$ 为 $A$ 的正次数齐次元，
它们在 $A_0$ 上生成该环。令 $d$ 为这些次数 $d_i$
的最小公倍数；这些是诸 $f_i$ 的次数。设 $M$ 为有限 $A$-模。
我们来说明，$\widetilde{M}$ 为零，当且仅当 $M$
是无关分次 $A$-模（按命题陈述上文的定义）。
为此，设 $x \in M$ 为齐次元。
选取充分小的 $k \in \mathbf{Z}$，并令
$N \to N'$ 与 $M \to N'$ 如引理
\ref{coherent-lemma-recover-tail-graded-module-general} 中所述。
还可以选取充分大的 $l$，使得
$M_n \to N_n$ 对 $n \geq l$ 为同构。
若 $\widetilde{M}$ 为零，则 $N = 0$。
因此，对任意齐次元 $f \in A_+$，若它满足
$\deg(f) + \deg(x) = nd$ 且 $nd > l$，则 $fx$ 为零，因为
$N_{nd} \to N'_{nd}$ 与 $M_{nd} \to N'_{nd}$ 都是同构。
故 $x$ 无关。
反之，设 $M$ 无关。则 $M_{nd}$ 对 $n \gg 0$ 为零
（参见命题前的讨论）。这显然蕴含
$M_{(f_i)} = M_{(f_i^{d/\deg(f_i)})} = 0$，
从而由构造可知 $\widetilde{M} = 0$。

\medskip\noindent
由此可知，子范畴 $\text{Mod}^{fg}_{A, irrelevant}$ 是
$\text{Mod}^{fg}_A$ 的 Serre 子范畴，因为它是正合函子
$M \mapsto \widetilde M$ 的核；参见《同调》引理
\ref{homology-lemma-kernel-exact-functor} 与《构造》引理
\ref{constructions-lemma-proj-sheaves}。
因此，箭头左侧的商范畴已在《同调》引理
\ref{homology-lemma-serre-subcategory-is-kernel} 中定义。
为定义命题中的函子，只需说明函子
$M \mapsto \widetilde M$ 把无关模送到 $0$；
这已在上文证明。

\medskip\noindent
由引理 \ref{coherent-lemma-coherent-on-proj-general}，所提出的拟逆有意义。
具体而言，该引理表明
$\mathcal{F} \longmapsto \bigoplus_{n \geq 0} \Gamma(X, \mathcal{F}(n))$
是函子 $\textit{Coh}(\mathcal{O}_X) \to \text{Mod}^{fg}_A$；
可将它与商函子
$\text{Mod}^{fg}_A \to \text{Mod}^{fg}_A/\text{Mod}^{fg}_{A, irrelevant}$
复合。

\medskip\noindent
由引理 \ref{coherent-lemma-recover-tail-graded-module-general}，
从左到右再到左的复合同构于恒等函子。
具体而言，设 $M$ 为有限分次 $A$-模，并取充分小的
$k \in \mathbf{Z}$，再令 $N \to N'$ 与 $M \to N'$
如引理 \ref{coherent-lemma-recover-tail-graded-module-general} 中所述。
于是 $M \to N'$ 的核与余核只在有限多个次数上非零，
因而都是无关的。此外，映射 $N \to N'$ 的核与余核
在所有充分大的、可被 $d$ 整除的次数上为零，
故它们也都是无关模。因此，$M \to N'$ 与 $N \to N'$
在商范畴中均为同构，正如所需。

\medskip\noindent
最后，设 $\mathcal{F}$ 为凝聚 $\mathcal{O}_X$-模。
令 $M = \bigoplus_{n \in \mathbf{Z}} \Gamma(X, \mathcal{F}(n))$，
将它视为分次 $A$-模，从而我们的函子把 $\mathcal{F}$ 送到
$M_{\geq 0} = \bigoplus_{n \geq 0} M_n$。
由《性质》引理 \ref{properties-lemma-proj-quasi-coherent}，
典范映射 $\widetilde M \to \mathcal{F}$ 是同构。
因为包含映射 $M_{\geq 0} \to M$ 定义同构
$\widetilde{M_{\geq 0}} \to \widetilde M$，可知
从右到左再到右的复合也同构于恒等函子。
\end{proof}










\section{沿射影态射的高阶直像}
\label{coherent-section-projective-pushforward}

\noindent
我们先陈述并证明基底仿射时的一个结果，
随后由此推出关于射影态射的一些结果。

\begin{lemma}
\label{coherent-lemma-coherent-proper-ample}
设 $R$ 为 Noether 环。设 $X \to \Spec(R)$ 为固有态射。
设 $\mathcal{L}$ 为 $X$ 上的丰沛可逆层。设 $\mathcal{F}$
为凝聚 $\mathcal{O}_X$-模。
\begin{enumerate}
\item 分次环
$A = \bigoplus_{d \geq 0} H^0(X, \mathcal{L}^{\otimes d})$
是有限生成 $R$-代数。
\item 存在 $r \geq 0$、
$d_1, \ldots, d_r \in \mathbf{Z}$ 以及满射
$$
\bigoplus\nolimits_{j = 1, \ldots, r} \mathcal{L}^{\otimes d_j}
\longrightarrow \mathcal{F}.
$$
\item 对任意 $p$，上同调群 $H^p(X, \mathcal{F})$ 是有限
$R$-模。
\item 若 $p > 0$，则
$H^p(X, \mathcal{F} \otimes_{\mathcal{O}_X} \mathcal{L}^{\otimes d}) = 0$
对所有充分大的 $d$ 成立。
\item 对任意 $k \in \mathbf{Z}$，分次 $A$-模
$$
\bigoplus\nolimits_{d \geq k}
H^0(X, \mathcal{F} \otimes_{\mathcal{O}_X} \mathcal{L}^{\otimes d})
$$
是有限 $A$-模。
\end{enumerate}
\end{lemma}

\begin{proof}
由《态射》引理
\ref{morphisms-lemma-finite-type-over-affine-ample-very-ample}，
存在 $d > 0$ 及浸入 $i : X \to \mathbf{P}^n_R$，使得
$\mathcal{L}^{\otimes d} \cong i^*\mathcal{O}_{\mathbf{P}^n_R}(1)$。
由于 $X$ 在 $R$ 上固有，态射 $i$ 是闭浸入
（《态射》引理 \ref{morphisms-lemma-image-proper-scheme-closed}）。
因此有
$H^i(X, \mathcal{G}) = H^i(\mathbf{P}^n_R, i_*\mathcal{G})$，
其中 $\mathcal{G}$ 是 $X$ 上任意准凝聚层
（由引理 \ref{coherent-lemma-relative-affine-cohomology} 以及闭浸入为仿射态射这一事实；
参见《态射》引理 \ref{morphisms-lemma-closed-immersion-affine}）。
此外，若 $\mathcal{G}$ 凝聚，则 $i_*\mathcal{G}$ 也凝聚
（引理 \ref{coherent-lemma-i-star-equivalence}）。
以下将不再特别提及这些事实。

\medskip\noindent
(1) 的证明。令 $S = R[T_0, \ldots, T_n]$，从而
$\mathbf{P}^n_R = \text{Proj}(S)$。
注意，$A$ 是 $S$-代数（但环映射 $S \to A$
不是分次环同态，因为 $S_n$ 映入 $A_{dn}$）。
由投影公式
（《上同调》引理 \ref{cohomology-lemma-projection-formula}），有
$$
i_*(\mathcal{L}^{\otimes nd + q}) =
i_*(\mathcal{L}^{\otimes q})
\otimes_{\mathcal{O}_{\mathbf{P}^n_R}}
\mathcal{O}_{\mathbf{P}^n_R}(n)
$$
对所有 $n \in \mathbf{Z}$ 成立。由引理
\ref{coherent-lemma-coherent-projective}，可知
$\bigoplus_{n \geq 0} A_{nd + q}$ 是有限分次 $S$-模。
由于
$A = \bigoplus_{q \in \{0, \ldots, d - 1} \bigoplus_{n \geq 0} A_{nd + q}$，
可知 $A$ 作为 $S$-代数是有限的，因而 (1) 成立。

\medskip\noindent
(2) 的证明。这由《性质》命题
\ref{properties-proposition-characterize-ample} 推出。

\medskip\noindent
(3) 的证明。应用引理 \ref{coherent-lemma-coherent-projective}，并使用
$H^p(X, \mathcal{F}) = H^p(\mathbf{P}^n_R, i_*\mathcal{F})$。

\medskip\noindent
(4) 的证明。固定 $p > 0$。由投影公式，有
$$
i_*(\mathcal{F} \otimes_{\mathcal{O}_X} \mathcal{L}^{\otimes nd + q}) =
i_*(\mathcal{F} \otimes_{\mathcal{O}_X} \mathcal{L}^{\otimes q})
\otimes_{\mathcal{O}_{\mathbf{P}^n_R}}
\mathcal{O}_{\mathbf{P}^n_R}(n)
$$
对所有 $n \in \mathbf{Z}$ 成立。由引理
\ref{coherent-lemma-coherent-projective} 可知
$H^p(X, \mathcal{F} \otimes \mathcal{L}^{nd + q}) = 0$
对 $n \gg 0$ 成立。因为整数模 $d$ 只有有限多个同余类，
这证明了 (4)。

\medskip\noindent
(5) 的证明。固定整数 $k$。令
{\scriptsize $M = \bigoplus_{n \geq k} H^0(X, \mathcal{F} \otimes \mathcal{L}^{\otimes n})$}。
如上论证可知，$\bigoplus_{nd + q \geq k} M_{nd + q}$
是有限分次 $S$-模。由于
$M = \bigoplus_{q \in \{0, \ldots, d - 1\}}
\bigoplus_{nd + q \geq k} M_{nd + q}$，可知 $M$
是有限 $S$-模。由于 $S$-模结构经由环映射 $S \to A$
分解，可知 $M$ 是有限 $A$-模。
\end{proof}

\begin{lemma}
\label{coherent-lemma-kill-by-twisting}
设 $f : X \to S$ 为概形态射。
设 $\mathcal{F}$ 为准凝聚 $\mathcal{O}_X$-模。
设 $\mathcal{L}$ 为 $X$ 上的可逆层。假设
\begin{enumerate}
\item $S$ 是 Noether 的；
\item $f$ 是固有的；
\item $\mathcal{F}$ 是凝聚的；并且
\item $\mathcal{L}$ 在 $X/S$ 上相对丰沛。
\end{enumerate}
则存在 $n_0$，使得对所有 $n \geq n_0$，都有
$$
R^pf_*\left(\mathcal{F} \otimes_{\mathcal{O}_X} \mathcal{L}^{\otimes n}\right)
=
0
$$
对所有 $p > 0$ 成立。
\end{lemma}

\begin{proof}
选取有限仿射开覆盖 $S = \bigcup V_j$，并令
$X_j = f^{-1}(V_j)$。
显然，若能对有限多个系统
$(X_j \to V_j, \mathcal{L}|_{X_j}, \mathcal{F}|_{V_j})$
中的每一个解决该问题，结果即告成立。因此可设 $S$ 仿射。
在这种情形下，
$R^pf_*(\mathcal{F} \otimes \mathcal{L}^{\otimes n})$
为零等价于
$H^p(X, \mathcal{F} \otimes \mathcal{L}^{\otimes n})$ 为零；参见引理
\ref{coherent-lemma-quasi-coherence-higher-direct-images-application}。
因此，所需消失由引理 \ref{coherent-lemma-coherent-proper-ample} 推出
（该引理适用，因为由《态射》引理
\ref{morphisms-lemma-finite-type-over-affine-ample-very-ample}，
$\mathcal{L}$ 在 $X$ 上丰沛）。
\end{proof}

\begin{lemma}
\label{coherent-lemma-locally-projective-pushforward}
设 $S$ 为局部 Noether 概形。
设 $f : X \to S$ 为局部射影态射。
设 $\mathcal{F}$ 为凝聚 $\mathcal{O}_X$-模。
则 $R^if_*\mathcal{F}$ 是凝聚 $\mathcal{O}_S$-模，
这对所有 $i \geq 0$ 成立。
\end{lemma}

\begin{proof}
首先注意，局部射影态射是固有的
（《态射》引理 \ref{morphisms-lemma-locally-projective-proper}），
因而是有限型的。特别地，$X$ 是局部 Noether 的
（《态射》引理 \ref{morphisms-lemma-finite-type-noetherian}），
故该陈述有意义。此外，由引理
\ref{coherent-lemma-quasi-coherence-higher-direct-images}，
诸层 $R^pf_*\mathcal{F}$ 都准凝聚。

\medskip\noindent
说明这些之后，该陈述在 $S$ 上是局部的（例如由《上同调》引理
\ref{cohomology-lemma-localize-higher-direct-images}）。
因此可设 $S = \Spec(R)$ 是 Noether 环的谱，且
$X$ 是 $\mathbf{P}^n_R$ 的闭子概形，其中 $n$ 取某个值；参见《态射》引理
\ref{morphisms-lemma-characterize-locally-projective}。
在这种情形下，诸层 $R^pf_*\mathcal{F}$ 是与 $R$-模
$H^p(X, \mathcal{F})$ 相伴的准凝聚层；参见引理
\ref{coherent-lemma-quasi-coherence-higher-direct-images-application}。
因此，只需说明 $R$-模 $H^p(X, \mathcal{F})$
是有限 $R$-模（引理 \ref{coherent-lemma-coherent-Noetherian}）。
这由引理 \ref{coherent-lemma-coherent-proper-ample} 推出
（因为 $\mathcal{O}_{\mathbf{P}^n_R}(1)$ 限制到 $X$
后在 $X$ 上丰沛）。
\end{proof}








\section{丰沛可逆层与上同调}
\label{coherent-section-ample-cohomology}

\noindent
下面给出仿射基底上的固有概形之丰沛性判据，
它以扭转后的上同调消失来表述。

\begin{lemma}
\label{coherent-lemma-vanshing-gives-ample}
\begin{reference}
\cite[III Proposition 2.6.1]{EGA}
\end{reference}
设 $R$ 为 Noether 环。设 $f : X \to \Spec(R)$ 为固有态射。
设 $\mathcal{L}$ 为可逆 $\mathcal{O}_X$-模。
下列条件等价：
\begin{enumerate}
\item $\mathcal{L}$ 在 $X$ 上丰沛（这与许多其他条件等价；参见
《性质》命题 \ref{properties-proposition-characterize-ample} 与
《态射》引理
\ref{morphisms-lemma-finite-type-over-affine-ample-very-ample}）；
\item 对每个凝聚 $\mathcal{O}_X$-模 $\mathcal{F}$，存在
$n_0 \geq 0$，使得
$H^p(X, \mathcal{F} \otimes \mathcal{L}^{\otimes n}) = 0$
对所有 $n \geq n_0$ 及 $p > 0$ 成立；并且
\item 对每个准凝聚理想层
$\mathcal{I} \subset \mathcal{O}_X$，存在 $n \geq 1$，
使得 $H^1(X, \mathcal{I} \otimes \mathcal{L}^{\otimes n}) = 0$。
\end{enumerate}
\end{lemma}

\begin{proof}
蕴含关系 (1) $\Rightarrow$ (2) 由引理
\ref{coherent-lemma-coherent-proper-ample} 推出。
蕴含关系 (2) $\Rightarrow$ (3) 是平凡的。
蕴含关系 (3) $\Rightarrow$ (1) 是引理
\ref{coherent-lemma-quasi-compact-h1-zero-invertible}。
\end{proof}

\begin{lemma}
\label{coherent-lemma-surjective-finite-morphism-ample}
设 $R$ 为 Noether 环。设 $f : Y \to X$ 为在 $R$
上固有的概形之间的态射。设 $\mathcal{L}$ 为可逆
$\mathcal{O}_X$-模。假设 $f$ 有限且满射。
则 $\mathcal{L}$ 丰沛，当且仅当 $f^*\mathcal{L}$ 丰沛。
\end{lemma}

\begin{proof}
丰沛可逆层经拟仿射态射拉回后仍丰沛；参见《态射》引理
\ref{morphisms-lemma-pullback-ample-tensor-relatively-ample}。
这证明了其中一个蕴含关系，因为有限态射按定义是仿射的。

\medskip\noindent
假设 $f^*\mathcal{L}$ 丰沛。令 $P$ 为凝聚
$\mathcal{O}_X$-模 $\mathcal{F}$ 的如下性质：存在 $n_0$，
使得 $H^p(X, \mathcal{F} \otimes \mathcal{L}^{\otimes n}) = 0$
对所有 $n \geq n_0$ 及 $p > 0$ 成立。我们将证明 $P$
对任意凝聚 $\mathcal{O}_X$-模 $\mathcal{F}$ 成立；由引理
\ref{coherent-lemma-vanshing-gives-ample}，这蕴含 $\mathcal{L}$ 丰沛。
我们要应用引理 \ref{coherent-lemma-property-higher-rank-cohomological}。
因此必须对 $P$ 验证该引理的 (1)、(2) 与 (3)。
性质 (1) 由与层的短正合列相伴的上同调长正合列，以及
与可逆层张量是正合函子这一事实推出。性质 (2) 由
$H^p(X, -)$ 是加性函子推出。为证明 (3)，设 $Z \subset X$
为泛点是 $\xi$ 的整闭子概形。
设 $\mathcal{F}$ 为 $Y$ 上的凝聚层，使得
$f_*\mathcal{F}$ 的支撑等于 $Z$，且
$(f_*\mathcal{F})_\xi$ 被 $\mathfrak m_\xi$ 零化；
参见引理 \ref{coherent-lemma-finite-morphism-Noetherian}。我们断言，
取 $\mathcal{G} = f_*\mathcal{F}$ 即可。只需验证引理
\ref{coherent-lemma-property-higher-rank-cohomological} 的 (3)(c)。
故假设 $\mathcal{J} \subset \mathcal{O}_X$ 为准凝聚理想层，
并满足 $\mathcal{J}_\xi = \mathcal{O}_{X, \xi}$。
有限态射是仿射的，故由引理
\ref{coherent-lemma-affine-morphism-projection-ideal} 可知
$\mathcal{J}\mathcal{G} = f_*(f^{-1}\mathcal{J}\mathcal{F})$。
此外，正如引理 \ref{coherent-lemma-affine-morphism-projection-ideal}
的证明所指出，层 $f^{-1}\mathcal{J}\mathcal{F}$ 是凝聚
$\mathcal{O}_Y$-模。由于 $\mathcal{L}$ 丰沛，由引理
\ref{coherent-lemma-vanshing-gives-ample} 可知存在 $n_0$，使得
$$
H^p(Y, f^{-1}\mathcal{J}\mathcal{F}
\otimes_{\mathcal{O}_Y} f^*\mathcal{L}^{\otimes n}) = 0,
$$
对 $n \geq n_0$ 及 $p > 0$ 成立。因为 $f$ 有限，因而仿射，故
\begin{align*}
H^p(X, \mathcal{J}\mathcal{G} \otimes_{\mathcal{O}_X}
\mathcal{L}^{\otimes n})
& =
H^p(X, f_*(f^{-1}\mathcal{J}\mathcal{F}) \otimes_{\mathcal{O}_X}
\mathcal{L}^{\otimes n}) \\
& =
H^p(X, f_*(f^{-1}\mathcal{J}\mathcal{F} \otimes_{\mathcal{O}_Y}
f^*\mathcal{L}^{\otimes n})) \\
& =
H^p(Y, f^{-1}\mathcal{J}\mathcal{F} \otimes_{\mathcal{O}_Y}
f^*\mathcal{L}^{\otimes n}) = 0
\end{align*}
这里使用了投影公式
（《上同调》引理 \ref{cohomology-lemma-projection-formula}）及引理
\ref{coherent-lemma-relative-affine-cohomology}。
因此，准凝聚子层 $\mathcal{G}' = \mathcal{J}\mathcal{G}$
满足 $P$。这正如所需地验证了引理
\ref{coherent-lemma-property-higher-rank-cohomological} 的性质 (3)(c)。
\end{proof}

\noindent
上同调具有函子性。特别地，给定环化空间 $X$、可逆
$\mathcal{O}_X$-模 $\mathcal{L}$ 以及截面
$s \in \Gamma(X, \mathcal{L})$，可得到映射
$$
H^p(X, \mathcal{F})
\longrightarrow
H^p(X, \mathcal{F} \otimes_{\mathcal{O}_X} \mathcal{L}), \quad
\xi \longmapsto s\xi
$$
它由映射
$\mathcal{F} \to \mathcal{F} \otimes_{\mathcal{O}_X} \mathcal{L}$
诱导；该映射是乘以 $s$。令
$\Gamma_*(X, \mathcal{L}) =
\bigoplus_{n \geq 0} \Gamma(X, \mathcal{L}^{\otimes n})$
为分次环；参见《模》定义 \ref{modules-definition-gamma-star}。
给定 $\mathcal{O}_X$-模层 $\mathcal{F}$ 及整数
$p \geq 0$，令
$$
H^p_*(X, \mathcal{L}, \mathcal{F}) =
\bigoplus\nolimits_{n \in \mathbf{Z}} H^p(X,
\mathcal{F} \otimes_{\mathcal{O}_X} \mathcal{L}^{\otimes n})
$$
由上述乘法，它是分次 $\Gamma_*(X, \mathcal{L})$-模。
警告：记号 $H^p_*(X, \mathcal{L}, \mathcal{F})$ 并非标准记号。

\begin{lemma}
\label{coherent-lemma-invert-s-cohomology}
设 $X$ 为概形。设 $\mathcal{L}$ 为 $X$ 上的可逆层。
设 $s \in \Gamma(X, \mathcal{L})$。
设 $\mathcal{F}$ 为准凝聚 $\mathcal{O}_X$-模。
若 $X$ 拟紧且拟分离，则典范映射
$$
H^p_*(X, \mathcal{L}, \mathcal{F})_{(s)}
\longrightarrow
H^p(X_s, \mathcal{F})
$$
把 $\xi/s^n$ 映到 $s^{-n}\xi$，并且是同构。
\end{lemma}

\begin{proof}
注意，当 $p = 0$ 时，这就是《性质》引理
\ref{properties-lemma-invert-s-sections}。
我们用归纳原理（引理 \ref{coherent-lemma-induction-principle}）证明该陈述；
对拟紧开集 $U \subset X$，令 $P(U)$ 为如下性质：
对所有 $p \geq 0$，映射
$$
H^p_*(U, \mathcal{L}, \mathcal{F})_{(s)}
\longrightarrow
H^p(U_s, \mathcal{F})
$$
是同构。

\medskip\noindent
若 $U$ 仿射，则由引理
\ref{coherent-lemma-quasi-coherent-affine-cohomology-zero} 与《性质》引理
\ref{properties-lemma-affine-cap-s-open}，
上述箭头两端对 $p > 0$ 均为零，故陈述成立。
若 $P$ 对 $U$、$V$ 及 $U \cap V$ 成立，
则可使用 Mayer--Vietoris 序列
（《上同调》引理 \ref{cohomology-lemma-mayer-vietoris}），
得到长正合列之间的映射
\begingroup\small
$$
\xymatrix{
H^{p - 1}_*(U \cap V, \mathcal{L}, \mathcal{F})_{(s)} \ar[r] \ar[d] &
H^p_*(U \cup V, \mathcal{L}, \mathcal{F})_{(s)} \ar[r] \ar[d] &
H^p_*(U, \mathcal{L}, \mathcal{F})_{(s)}
\oplus
H^p_*(V, \mathcal{L}, \mathcal{F})_{(s)} \ar[d] \\
H^{p - 1}(U_s \cap V_s, \mathcal{F}) \ar[r]&
H^p(U_s \cup V_s, \mathcal{F}) \ar[r] &
H^p(U_s, \mathcal{F})
\oplus
H^p(V_s, \mathcal{F})
}
$$
\endgroup
（这里只展示一个片段）。注意，$U_s \cap V_s = (U \cap V)_s$，
并且 $U_s \cup V_s = (U \cup V)_s$。因此左右两个竖直映射
都是同构（图中未画出的右侧再一个映射和左侧再一个映射也如此）。
由五引理（《同调》引理 \ref{homology-lemma-five-lemma}），
可知 $P(U \cup V)$ 成立。这就完成了证明。
\end{proof}

\begin{lemma}
\label{coherent-lemma-section-affine-open-kills-classes}
设 $X$ 为概形。
设 $\mathcal{L}$ 为可逆 $\mathcal{O}_X$-模。
设 $s \in \Gamma(X, \mathcal{L})$ 为截面。假设
\begin{enumerate}
\item $X$ 拟紧且拟分离；并且
\item $X_s$ 仿射。
\end{enumerate}
则对每个准凝聚 $\mathcal{O}_X$-模 $\mathcal{F}$、
每个 $p > 0$ 以及所有 $\xi \in H^p(X, \mathcal{F})$，
都存在 $n \geq 0$，使得 $s^n\xi = 0$ 在
$H^p(X, \mathcal{F} \otimes_{\mathcal{O}_X} \mathcal{L}^{\otimes n})$
中成立。
\end{lemma}

\begin{proof}
回忆，$H^p(X_s, \mathcal{G})$
对每个准凝聚模 $\mathcal{G}$ 都为零；这是引理
\ref{coherent-lemma-quasi-coherent-affine-cohomology-zero}。
因此本引理由引理 \ref{coherent-lemma-invert-s-cohomology} 推出。
\end{proof}

\noindent
下述引理的更一般版本见《极限》引理
\ref{limits-lemma-ample-on-reduction}。

\begin{lemma}
\label{coherent-lemma-ample-on-reduction}
设 $i : Z \to X$ 为 Noether 概形的闭浸入，
并诱导底拓扑空间的同胚。
设 $\mathcal{L}$ 为 $X$ 上的可逆层。
则 $i^*\mathcal{L}$ 在 $Z$ 上丰沛，当且仅当
$\mathcal{L}$ 在 $X$ 上丰沛。
\end{lemma}

\begin{proof}
若 $\mathcal{L}$ 丰沛，则例如由《态射》引理
\ref{morphisms-lemma-pullback-ample-tensor-relatively-ample}，
$i^*\mathcal{L}$ 丰沛。假设 $i^*\mathcal{L}$ 丰沛。
必须证明 $\mathcal{L}$ 在 $X$ 上丰沛。
设 $\mathcal{I} \subset \mathcal{O}_X$ 为定义闭子概形
$Z$ 的凝聚理想层。由于在集合意义下 $i(Z) = X$，
由引理 \ref{coherent-lemma-power-ideal-kills-sheaf} 可知，
$\mathcal{I}^n = 0$ 对某个 $n$ 成立。
考虑由闭子概形组成的序列
$$
X = Z_n \supset Z_{n - 1} \supset Z_{n - 2} \supset \ldots \supset Z_1 = Z
$$
它们分别由
$0 = \mathcal{I}^n \subset \mathcal{I}^{n - 1} \subset \ldots \subset
\mathcal{I}$ 定义。于是每个闭浸入 $Z_{i - 1} \to Z_{i}$
都由平方为零的凝聚理想层定义。这样便归约到
$\mathcal{I}^2 = 0$ 的情形。

\medskip\noindent
考虑准凝聚 $\mathcal{O}_X$-模的短正合列
$$
0 \to \mathcal{I} \to \mathcal{O}_X \to i_*\mathcal{O}_Z \to 0
$$
与 $\mathcal{L}^{\otimes n}$ 张量后，得到短正合列
\begin{equation}
\label{coherent-equation-ses}
0 \to \mathcal{I} \otimes_{\mathcal{O}_X} \mathcal{L}^{\otimes n}
\to \mathcal{L}^{\otimes n} \to i_*i^*\mathcal{L}^{\otimes n} \to 0
\end{equation}
因为 $\mathcal{I}^2 = 0$，可使用《态射》引理
\ref{morphisms-lemma-i-star-equivalence}，把 $\mathcal{I}$
看作准凝聚 $\mathcal{O}_Z$-模；于是
$\mathcal{I} \otimes_{\mathcal{O}_X} \mathcal{L}^{\otimes n} =
\mathcal{I} \otimes_{\mathcal{O}_Z} i^*\mathcal{L}^{\otimes n}$，
这里作了显然的记号滥用。此外，该层在 $Z$ 上的上同调
典范地等同于它在 $X$ 上的上同调（因为 $i$ 是同胚）。

\medskip\noindent
设 $x \in X$ 为一点，并以 $z \in Z$ 表示相应的点。
因为 $i^*\mathcal{L}$ 丰沛，存在 $n$ 及截面
$s \in \Gamma(Z, i^*\mathcal{L}^{\otimes n})$，使得
$z \in Z_s$ 且 $Z_s$ 仿射。把 $s$ 提升为
$\mathcal{L}^{\otimes n}$ 在 $X$ 上的截面之障碍，是来自层的短正合列
(\ref{coherent-equation-ses}) 的连接像
$$
\xi = \partial s \in
H^1(X, \mathcal{I} \otimes_{\mathcal{O}_X} \mathcal{L}^{\otimes n}) =
H^1(Z, \mathcal{I} \otimes_{\mathcal{O}_Z} i^*\mathcal{L}^{\otimes n})
$$
若把 $s$ 替换为 $s^{e + 1}$，则 $\xi$ 被替换为
$\partial(s^{e + 1}) = (e + 1) s^e \xi$；其所在群为
$H^1(Z, \mathcal{I} \otimes_{\mathcal{O}_Z} i^*\mathcal{L}^{\otimes (e + 1)n})$，
因为
$$
0 \to
\bigoplus\nolimits_{m \geq 0}
\mathcal{I} \otimes_{\mathcal{O}_X} \mathcal{L}^{\otimes m} \to
\bigoplus\nolimits_{m \geq 0}
\mathcal{L}^{\otimes m} \to
\bigoplus\nolimits_{m \geq 0}
i_*i^*\mathcal{L}^{\otimes m} \to 0
$$
的连接映射由《上同调》引理
\ref{cohomology-lemma-boundary-derivation-over-cup-product} 是导子。
由引理 \ref{coherent-lemma-section-affine-open-kills-classes}，
可知 $s^e \xi$ 对充分大的 $e$ 为零。
因此，把 $s$ 替换为某个幂后，可以假设 $s$ 是截面
$s' \in \Gamma(X, \mathcal{L}^{\otimes n})$ 的像。
于是 $X_{s'}$ 是开子概形，而 $Z_s \to X_{s'}$
是 Noether 概形之间的满射闭浸入，且 $Z_s$ 仿射。
因此，由引理 \ref{coherent-lemma-image-affine-finite-morphism-affine-Noetherian}，
$X_s$ 仿射，于是可知 $\mathcal{L}$ 丰沛。
\end{proof}

\noindent
下述引理的更一般版本见《极限》引理
\ref{limits-lemma-thickening-quasi-affine}。

\begin{lemma}
\label{coherent-lemma-thickening-quasi-affine}
设 $i : Z \to X$ 为 Noether 概形的闭浸入，
并诱导底拓扑空间的同胚。
则 $X$ 拟仿射，当且仅当 $Z$ 拟仿射。
\end{lemma}

\begin{proof}
回忆，概形拟仿射，当且仅当其结构层丰沛；参见《性质》引理
\ref{properties-lemma-quasi-affine-O-ample}。
因此，若 $Z$ 拟仿射，则 $\mathcal{O}_Z$ 丰沛；
由引理 \ref{coherent-lemma-ample-on-reduction}，$\mathcal{O}_X$ 丰沛；
从而 $X$ 拟仿射。反向阅读刚才给出的论证，
便得到逆命题的一个证明；它也可以用初等方法看出。
\end{proof}

\begin{lemma}
\label{coherent-lemma-affine-in-presence-ample}
设 $X$ 为概形。设 $\mathcal{L}$ 为丰沛可逆
$\mathcal{O}_X$-模。设 $n_0$ 为整数。
若 $H^p(X, \mathcal{L}^{\otimes -n}) = 0$
对 $n \geq n_0$ 及 $p > 0$ 成立，则 $X$ 仿射。
\end{lemma}

\begin{proof}
我们断言，$H^p(X, \mathcal{F}) = 0$ 对每个准凝聚
$\mathcal{O}_X$-模及 $p > 0$ 成立。由《性质》定义
\ref{properties-definition-ample}，$X$ 拟紧，
故该断言由引理 \ref{coherent-lemma-quasi-compact-h1-zero-covering}
完成证明。由《性质》引理
\ref{properties-lemma-ample-separated}，概形 $X$ 分离。
设 $X$ 被 $e + 1$ 个仿射开集覆盖。于是
$H^p(X, \mathcal{F}) = 0$ 对 $p > e$ 成立；参见引理
\ref{coherent-lemma-vanishing-nr-affines}。因此可对 $p$ 作降归纳
来证明该断言。把 $\mathcal{F}$ 写成有限型准凝聚模的滤过余极限
（《性质》引理
\ref{properties-lemma-quasi-coherent-colimit-finite-type}），
并使用《上同调》引理
\ref{cohomology-lemma-quasi-separated-cohomology-colimit}，
可设 $\mathcal{F}$ 为有限型。于是可以选取 $n > n_0$，
使得 $\mathcal{F} \otimes_{\mathcal{O}_X} \mathcal{L}^{\otimes n}$
由整体截面生成；参见《性质》命题
\ref{properties-proposition-characterize-ample}。
这意味着存在短正合列
$$
0 \to \mathcal{F}' \to
\bigoplus\nolimits_{i \in I} \mathcal{L}^{\otimes -n}
\to \mathcal{F} \to 0
$$
其中 $I$ 是某个集合（事实上可选取有限的 $I$）。
由归纳假设，$H^{p + 1}(X, \mathcal{F}') = 0$；
而由假设（结合已经使用过的上同调与余极限的交换性），有
$H^p(X, \bigoplus_{i \in I} \mathcal{L}^{\otimes -n}) = 0$。
由上同调长正合列，可知
$H^p(X, \mathcal{F}) = 0$，正如所需。
\end{proof}

\begin{lemma}
\label{coherent-lemma-affine-if-quasi-affine}
设 $X$ 为拟仿射概形。
若 $H^p(X, \mathcal{O}_X) = 0$ 对 $p > 0$ 成立，
则 $X$ 仿射。
\end{lemma}

\begin{proof}
由《性质》引理 \ref{properties-lemma-quasi-affine-O-ample}，
$\mathcal{O}_X$ 丰沛，故结论由引理
\ref{coherent-lemma-affine-in-presence-ample} 推出。
\end{proof}







\section{Chow 引理}
\label{coherent-section-chows-lemma}

\noindent
本节证明 Noether 情形的 Chow 引理
（引理 \ref{coherent-lemma-chow-Noetherian}）。
在《极限》第 \ref{limits-section-chows-lemma} 节中，
我们证明非 Noether 情形的一些变体。

\begin{lemma}
\label{coherent-lemma-chow-Noetherian}
\begin{reference}
\cite[II Theorem 5.6.1(a)]{EGA}
\end{reference}
设 $S$ 为 Noether 概形。
设 $f : X \to S$ 为有限型分离态射。
则存在 $n \geq 0$ 及交换图
$$
\xymatrix{
X \ar[rd] & X' \ar[d] \ar[l]^\pi \ar[r] & \mathbf{P}^n_S \ar[dl] \\
& S &
}
$$
其中 $X' \to \mathbf{P}^n_S$ 是浸入，且
$\pi : X' \to X$ 固有并满射。此外，还可以作如此安排：
存在稠密开子概形 $U \subset X$，使得
$\pi^{-1}(U) \to U$ 为同构。
\end{lemma}

\begin{proof}
在余下证明中遇到的所有概形，都将在 Noether 概形 $S$
上有限型，因而是 Noether 的
（参见《态射》引理 \ref{morphisms-lemma-finite-type-noetherian}）。
它们之间的所有态射都自动拟紧、局部有限型且拟分离；参见
《态射》引理 \ref{morphisms-lemma-permanence-finite-type} 及《性质》引理
\ref{properties-lemma-locally-Noetherian-quasi-separated} 与
\ref{properties-lemma-morphism-Noetherian-schemes-quasi-compact}。

\medskip\noindent
概形 $X$ 只有有限多个不可约分支
（《性质》引理 \ref{properties-lemma-Noetherian-irreducible-components}）。
设 $X = X_1 \cup \ldots \cup X_r$ 是 $X$
分解为不可约分支的表达式。
令 $\eta_i \in X_i$ 为泛点。对每个点 $x \in X$，
存在仿射开集 $U_x \subset X$，它包含 $x$ 以及所有泛点
$\eta_i$；参见《性质》引理
\ref{properties-lemma-point-and-maximal-points-affine}。
因为 $X$ 拟紧，可以找到有限仿射开覆盖
$X = U_1 \cup \ldots \cup U_m$，使得每个 $U_i$
都包含 $\eta_1, \ldots, \eta_r$。特别地，可知开集
$U = U_1 \cap \ldots \cap U_m \subset X$ 稠密。
以下只会用到这一点，以及诸 $U_i$ 是覆盖 $X$ 的仿射开集这一事实。

\medskip\noindent
令 $X^* \subset X$ 为 $U \to X$ 的概形论闭包；参见《态射》定义
\ref{morphisms-definition-scheme-theoretic-image}。
令 $U_i^* = X^* \cap U_i$。注意，$U_i^*$ 是 $U_i$
的闭子概形，因而 $U_i^*$ 仿射。由于 $U$ 在 $X$ 中稠密，
态射 $X^* \to X$ 是满射闭浸入。它在 $U$ 上是同构。
因此可以把 $X$ 替换为 $X^*$，把 $U_i$ 替换为 $U_i^*$，
并假设 $U$ 在 $X$ 中概形论稠密；参见《态射》定义
\ref{morphisms-definition-scheme-theoretically-dense}。

\medskip\noindent
由《态射》引理 \ref{morphisms-lemma-quasi-projective-finite-type-over-S}，
可以找到浸入 $j_i : U_i \to \mathbf{P}_S^{n_i}$，
这对每个 $i$ 成立。
由《态射》引理 \ref{morphisms-lemma-quasi-compact-immersion}，
可以找到闭子概形 $Z_i \subset \mathbf{P}_S^{n_i}$，
使得 $j_i : U_i \to Z_i$ 是概形论稠密开浸入。
注意，$Z_i \to S$ 固有；参见《态射》引理
\ref{morphisms-lemma-locally-projective-proper}。考虑态射
$$
j = (j_1|_U, \ldots, j_m|_U) : U \longrightarrow
\mathbf{P}_S^{n_1} \times_S \ldots \times_S \mathbf{P}_S^{n_m}.
$$
由上述引理，可以找到闭子概形 $Z$，它位于
$\mathbf{P}_S^{n_1} \times_S \ldots \times_S \mathbf{P}_S^{n_m}$
中，并且 $j : U \to Z$ 是开浸入，
并且 $U$ 在 $Z$ 中概形论稠密。态射 $Z \to S$ 固有。
考虑第 $i$ 个投影
$$
\text{pr}_i|_Z : Z \longrightarrow \mathbf{P}^{n_i}_S.
$$
该态射经由 $Z_i$ 分解（参见《态射》引理
\ref{morphisms-lemma-factor-factor}）。以 $p_i : Z \to Z_i$
表示诱导态射。例如由《态射》引理
\ref{morphisms-lemma-image-proper-scheme-closed}，这是固有态射。
此时已有 $U \subset U_i \subset Z_i$ 是概形论稠密开浸入。
此外，可把 $Z$ 看作“对角”态射
$U \to Z_1 \times_S \ldots \times_S Z_m$ 的概形论像。

\medskip\noindent
令 $V_i = p_i^{-1}(U_i)$。注意，
$p_i|_{V_i} : V_i \to U_i$ 固有。令
$X' = V_1 \cup \ldots \cup V_m$。由构造，$X'$ 可浸入概形
$\mathbf{P}^{n_1}_S \times_S \ldots \times_S \mathbf{P}^{n_m}_S$。
因此，由 Segre 嵌入（参见《构造》引理
\ref{constructions-lemma-segre-embedding}），可知 $X'$
可浸入 $S$ 上的某个射影空间。

\medskip\noindent
我们断言，诸态射 $p_i|_{V_i}: V_i \to U_i$
可粘合为态射 $X' \to X$。事实上，显然 $p_i|_U$
是从 $U$ 到 $U$ 的恒等映射。由构造，$U \subset X'$
概形论稠密，因而它在开子概形 $V_i \cap V_j$ 中也概形论稠密。
于是
$p_i|_{V_i \cap V_j} = p_j|_{V_i \cap V_j}$，
这里把两端视为到分离 $S$-概形 $X$ 的态射；参见《态射》引理
\ref{morphisms-lemma-equality-of-morphisms}。
把所得态射记为 $\pi : X' \to X$。

\medskip\noindent
我们断言 $\pi^{-1}(U_i) = V_i$。
由于 $\pi|_{V_i} = p_i|_{V_i}$，可知
$V_i \subset \pi^{-1}(U_i)$。考虑交换图
$$
\xymatrix{
V_i \ar[r] \ar[rd]_{p_i|_{V_i}} & \pi^{-1}(U_i) \ar[d] \\
& U_i
}
$$
因为 $V_i \to U_i$ 固有，由《态射》引理
\ref{morphisms-lemma-image-proper-scheme-closed}，
可知水平箭头的像闭。由于
$V_i \subset \pi^{-1}(U_i)$ 概形论稠密（因为它包含 $U$），
可知 $V_i = \pi^{-1}(U_i)$，正如所断言。

\medskip\noindent
这表明 $\pi^{-1}(U_i) \to U_i$ 可与固有态射
$p_i|_{V_i} : V_i \to U_i$ 识别。因此可知，$X$
有有限仿射覆盖 $X = \bigcup U_i$，使得 $\pi$
在该覆盖每个成员上的限制都固有。
于是由《态射》引理 \ref{morphisms-lemma-proper-local-on-the-base}，
可知 $\pi$ 固有。

\medskip\noindent
最后必须说明 $\pi^{-1}(U) = U$。为此，使用交换图
$$
\xymatrix{
U \ar[r] \ar[rd] & \pi^{-1}(U) \ar[d] \\
& U
}
$$
并使用 $\text{id}_U : U \to U$ 固有且 $U$
在 $\pi^{-1}(U)$ 中概形论稠密，便可如上作同样论证。
\end{proof}

\begin{remark}
\label{coherent-remark-chow-Noetherian}
在 Chow 引理 \ref{coherent-lemma-chow-Noetherian} 的情形中：
\begin{enumerate}
\item 态射 $\pi$ 实际上是 H-射影的（因而射影；参见《态射》引理
\ref{morphisms-lemma-H-projective}），因为态射
$X' \to \mathbf{P}^n_S \times_S X = \mathbf{P}^n_X$
是闭浸入（使用 $\pi$ 固有这一事实；参见《态射》引理
\ref{morphisms-lemma-image-proper-scheme-closed}）。
\item 可以假设 $\pi^{-1}(U)$ 在 $X'$ 中概形论稠密。
事实上，只需把 $X'$ 替换为 $\pi^{-1}(U)$ 的概形论闭包。
在这种情形下，可把 $U$ 看作 $X'$ 的概形论稠密开子概形。
参见《态射》第 \ref{morphisms-section-scheme-theoretic-image} 节。
\item 若 $X$ 既约，则可选取既约的 $X'$。这由 (2) 显然可见。
\end{enumerate}
\end{remark}






\section{凝聚层的高阶直像}
\label{coherent-section-proper-pushforward}

\noindent
本节证明一个基本事实：凝聚层沿固有态射的高阶直像仍凝聚。

\begin{proposition}
\label{coherent-proposition-proper-pushforward-coherent}
\begin{reference}
\cite[III Theorem 3.2.1]{EGA}
\end{reference}
设 $S$ 为局部 Noether 概形。
设 $f : X \to S$ 为固有态射。
设 $\mathcal{F}$ 为凝聚 $\mathcal{O}_X$-模。
则 $R^if_*\mathcal{F}$ 是凝聚 $\mathcal{O}_S$-模，
这对所有 $i \geq 0$ 成立。
\end{proposition}

\begin{proof}
由于该问题在 $S$ 上是局部的，可以假设 $S$ 是 Noether 概形。
由于固有态射是有限型的，可知在这种情形下 $X$ 也是 Noether 概形。
考虑性质 $\mathcal{P}$；它定义在 $X$ 上的凝聚层上，规则如下：
$$
\mathcal{P}(\mathcal{F}) \Leftrightarrow
R^pf_*\mathcal{F}\text{ is coherent for all }p \geq 0
$$
我们将用引理 \ref{coherent-lemma-property} 的结论证明
$\mathcal{P}$ 对 $X$ 上每个凝聚层都成立。

\medskip\noindent
设
$$
0 \to \mathcal{F}_1 \to \mathcal{F}_2 \to \mathcal{F}_3 \to 0
$$
为 $X$ 上凝聚层的短正合列。考虑高阶直像的长正合列
$$
R^{p - 1}f_*\mathcal{F}_3 \to
R^pf_*\mathcal{F}_1 \to
R^pf_*\mathcal{F}_2 \to
R^pf_*\mathcal{F}_3 \to
R^{p + 1}f_*\mathcal{F}_1
$$
显然，若诸层 $\mathcal{F}_i$ 中三项取二具有性质
$\mathcal{P}$，则第三个层的高阶直像在这个正合复形中夹在
两个凝聚层之间。因此，由引理
\ref{coherent-lemma-coherent-abelian-Noetherian} 与
\ref{coherent-lemma-coherent-Noetherian-quasi-coherent-sub-quotient}，
这些高阶直像也凝聚。故性质 $\mathcal{P}$ 对第三个层也成立。

\medskip\noindent
设 $Z \subset X$ 为整闭子概形。必须找到凝聚层
$\mathcal{F}$（在 $X$ 上），其支撑包含在 $Z$ 中；其在泛点
$\xi$（即 $Z$ 的泛点）处的茎是 $1$ 维向量空间，
标量域为 $\kappa(\xi)$，
并且 $\mathcal{P}$ 对 $\mathcal{F}$ 成立。
以 $g = f|_Z : Z \to S$ 表示 $f$ 的限制。
假设能找到凝聚层 $\mathcal{G}$（在 $Z$ 上），使得
(a) $\mathcal{G}_\xi$ 是 $1$ 维向量空间，其标量域为 $\kappa(\xi)$，
(b) $R^pg_*\mathcal{G} = 0$ 对 $p > 0$ 成立，并且
(c) $g_*\mathcal{G}$ 凝聚。于是可考虑
$\mathcal{F} = (Z \to X)_*\mathcal{G}$。因为 $Z \to X$ 是闭浸入，
可知 $(Z \to X)_*\mathcal{G}$ 在 $X$ 上凝聚，并且
$R^p(Z \to X)_*\mathcal{G} = 0$ 对 $p > 0$ 成立
（引理 \ref{coherent-lemma-finite-pushforward-coherent}）。
因此，由相对 Leray 谱序列
（《上同调》引理 \ref{cohomology-lemma-relative-Leray}），
将有 $R^pf_*\mathcal{F} = R^pg_*\mathcal{G} = 0$ 对 $p > 0$ 成立，
且 $f_*\mathcal{F} = g_*\mathcal{G}$ 凝聚。
最后，$\mathcal{F}_\xi = ((Z \to X)_*\mathcal{G})_\xi = \mathcal{G}_\xi$，
这验证了关于 $\xi$ 处之茎的条件。
因此，一切归结为寻找具有 (a)、(b)、(c) 的凝聚层
$\mathcal{G}$（在 $Z$ 上）。

\medskip\noindent
可以把 Chow 引理 \ref{coherent-lemma-chow-Noetherian}
应用于态射 $Z \to S$。由此得到 Chow 引理陈述中的交换图
$$
\xymatrix{
Z \ar[rd]_g & Z' \ar[d]^-{g'} \ar[l]^\pi \ar[r]_i & \mathbf{P}^m_S \ar[dl] \\
& S &
}
$$
此外，令 $U \subset Z$ 为稠密开子概形，使得
$\pi^{-1}(U) \to U$ 为同构。由评注
\ref{coherent-remark-chow-Noetherian} 中的讨论，可知
$i' = (i, \pi) : Z' \to \mathbf{P}^m_Z$ 是闭浸入。因此
$$
\mathcal{L} = i^*\mathcal{O}_{\mathbf{P}^m_S}(1) \cong
(i')^*\mathcal{O}_{\mathbf{P}^m_Z}(1)
$$
既 $g'$-相对丰沛，也 $\pi$-相对丰沛（例如由《态射》引理
\ref{morphisms-lemma-characterize-ample-on-finite-type}）。
因此，由引理 \ref{coherent-lemma-kill-by-twisting}，存在 $n \geq 0$，
使得 $R^p\pi_*\mathcal{L}^{\otimes n} = 0$ 对所有 $p > 0$ 成立，
且 $R^p(g')_*\mathcal{L}^{\otimes n} = 0$ 对所有 $p > 0$ 成立。
令 $\mathcal{G} = \pi_*\mathcal{L}^{\otimes n}$。
性质 (a) 成立，因为 $\pi_*\mathcal{L}^{\otimes n}|_U$
是可逆层（因为 $\pi^{-1}(U) \to U$ 是同构）。
性质 (b) 与 (c) 成立，因为由相对 Leray 谱序列
（《上同调》引理 \ref{cohomology-lemma-relative-Leray}），有
$$
E_2^{p, q} = R^pg_* R^q\pi_*\mathcal{L}^{\otimes n}
\Rightarrow
R^{p + q}(g')_*\mathcal{L}^{\otimes n}
$$
而由 $n$ 的选取，$E_2^{p, q}$ 中仅 $q = 0$ 的项可能非零，
且 $R^{p + q}(g')_*\mathcal{L}^{\otimes n}$ 中仅
$p = q = 0$ 的项可能非零。这蕴含
$R^pg_*\mathcal{G} = 0$ 对 $p > 0$ 成立，并且
$g_*\mathcal{G} = (g')_*\mathcal{L}^{\otimes n}$。
最后，应用前面的引理 \ref{coherent-lemma-locally-projective-pushforward}，
可知 $g_*\mathcal{G} = (g')_*\mathcal{L}^{\otimes n}$
如所需地凝聚。
\end{proof}

\begin{lemma}
\label{coherent-lemma-proper-over-affine-cohomology-finite}
设 $S = \Spec(A)$，其中 $A$ 为 Noether 环。
设 $f : X \to S$ 为固有态射。
设 $\mathcal{F}$ 为凝聚 $\mathcal{O}_X$-模。
则 $H^i(X, \mathcal{F})$ 是有限 $A$-模，
这对所有 $i \geq 0$ 成立。
\end{lemma}

\begin{proof}
这正是命题 \ref{coherent-proposition-proper-pushforward-coherent} 的仿射情形。
具体而言，由引理 \ref{coherent-lemma-quasi-coherence-higher-direct-images} 与
\ref{coherent-lemma-quasi-coherence-higher-direct-images-application}，已知
$R^if_*\mathcal{F}$ 是与 $A$-模 $H^i(X, \mathcal{F})$
相伴的准凝聚层；而由引理 \ref{coherent-lemma-coherent-Noetherian}，
该层凝聚，当且仅当 $H^i(X, \mathcal{F})$
是有限型 $A$-模。
\end{proof}

\begin{lemma}
\label{coherent-lemma-graded-finiteness}
设 $A$ 为 Noether 环。
设 $B$ 为有限生成分次 $A$-代数。
设 $f : X \to \Spec(A)$ 为固有态射。
令 $\mathcal{B} = f^*\widetilde B$。
设 $\mathcal{F}$ 为有限型准凝聚分次 $\mathcal{B}$-模。
\begin{enumerate}
\item 对每个 $p \geq 0$，分次 $B$-模 $H^p(X, \mathcal{F})$
是有限 $B$-模。
\item 若 $\mathcal{L}$ 为丰沛可逆 $\mathcal{O}_X$-模，
则存在整数 $d_0$，使得
$H^p(X, \mathcal{F} \otimes \mathcal{L}^{\otimes d}) = 0$
对所有 $p > 0$ 及 $d \geq d_0$ 成立。
\end{enumerate}
\end{lemma}

\begin{proof}
为证明该引理，考虑纤维积图
$$
\xymatrix{
X' = \Spec(B) \times_{\Spec(A)} X
\ar[r]_-\pi \ar[d]_{f'} &
X \ar[d]^f \\
\Spec(B) \ar[r] &
\Spec(A)
}
$$
注意，$f'$ 是固有态射；参见《态射》引理
\ref{morphisms-lemma-base-change-proper}。
此外，$B$ 是有限生成 $A$-代数，因而是 Noether 的
（《代数》引理 \ref{algebra-lemma-Noetherian-permanence}）。
这蕴含 $X'$ 是 Noether 概形
（《态射》引理 \ref{morphisms-lemma-finite-type-noetherian}）。
由《构造》引理 \ref{constructions-lemma-spec-properties}，
注意 $X'$ 是准凝聚 $\mathcal{O}_X$-代数 $\mathcal{B}$ 的相对谱。
因为 $\mathcal{F}$ 是准凝聚 $\mathcal{B}$-模，
可知存在唯一的准凝聚 $\mathcal{O}_{X'}$-模 $\mathcal{F}'$，
使得 $\pi_*\mathcal{F}' = \mathcal{F}$；参见《态射》引理
\ref{morphisms-lemma-affine-equivalence-modules}
由于 $\mathcal{F}$ 作为 $\mathcal{B}$-模是有限型的，
可知 $\mathcal{F}'$ 是有限型 $\mathcal{O}_{X'}$-模
（细节从略）。换言之，$\mathcal{F}'$ 是凝聚
$\mathcal{O}_{X'}$-模（引理 \ref{coherent-lemma-coherent-Noetherian}）。
由于态射 $\pi : X' \to X$ 仿射，由引理
\ref{coherent-lemma-relative-affine-cohomology} 有
$$
H^p(X, \mathcal{F}) = H^p(X', \mathcal{F}')
$$
因此 (1) 由引理 \ref{coherent-lemma-proper-over-affine-cohomology-finite} 推出。
给定 (2) 中的 $\mathcal{L}$，令
$\mathcal{L}' = \pi^*\mathcal{L}$。注意，由《态射》引理
\ref{morphisms-lemma-pullback-ample-tensor-relatively-ample}，
$\mathcal{L}'$ 在 $X'$ 上丰沛。由投影公式
（《上同调》引理 \ref{cohomology-lemma-projection-formula}），有
$\pi_*(\mathcal{F}' \otimes \mathcal{L}') = \mathcal{F} \otimes \mathcal{L}$。
因此，按上述相同论证，(2) 由引理
\ref{coherent-lemma-kill-by-twisting} 推出。
\end{proof}

\section{形式函数定理}
\label{coherent-section-theorem-formal-functions}

\noindent
本节研究形如 $\{I^n\mathcal{F}\}_{n \geq 0}$ 或
$\{\mathcal{F}/I^n\mathcal{F}\}_{n \geq 0}$ 的层序列在
$n$ 变化时其上同调的行为。

\medskip\noindent
这里及下文采用如下记号。给定概形态射 $f : X \to Y$、
准凝聚层 $\mathcal{F}$（在 $X$ 上），以及准凝聚理想层
$\mathcal{I} \subset \mathcal{O}_Y$，以
$\mathcal{I}^n\mathcal{F}$ 表示由 $f^{-1}(\mathcal{I}^n)$
的局部截面与 $\mathcal{F}$ 的局部截面之乘积生成的准凝聚子层。用公式表示为
$$
\mathcal{I}^n\mathcal{F}
=
\Im\left(
f^*(\mathcal{I}^n) \otimes_{\mathcal{O}_X} \mathcal{F}
\longrightarrow
\mathcal{F}
\right).
$$
注意有自然映射
$$
f^{-1}(\mathcal{I}^n) \otimes_{f^{-1}\mathcal{O}_Y} \mathcal{I}^m\mathcal{F}
\longrightarrow
f^*(\mathcal{I}^n) \otimes_{\mathcal{O}_X} \mathcal{I}^m\mathcal{F}
\longrightarrow
\mathcal{I}^{n + m}\mathcal{F}
$$
因此，$\mathcal{I}^n$ 的一个截面借助高阶直像的函子性给出映射
{\small $R^pf_*(\mathcal{I}^m\mathcal{F}) \to
R^pf_*(\mathcal{I}^{n + m}\mathcal{F})$}。
先局部化再层化，可见存在 $\mathcal{O}_Y$-模映射
$$
\mathcal{I}^n \otimes_{\mathcal{O}_Y} R^pf_*(\mathcal{I}^m\mathcal{F})
\longrightarrow
R^pf_*(\mathcal{I}^{n + m}\mathcal{F}).
$$
换言之，$\bigoplus_{n \geq 0} R^pf_*(\mathcal{I}^n\mathcal{F})$
是分次 $\bigoplus_{n \geq 0} \mathcal{I}^n$-模。

\medskip\noindent
若 $Y = \Spec(A)$ 且 $\mathcal{I} = \widetilde{I}$，则将
$\mathcal{I}^n\mathcal{F}$ 简记为 $I^n\mathcal{F}$。上面引入的映射赋予
$M = \bigoplus H^p(X, I^n\mathcal{F})$ 分次
$S = \bigoplus I^n$-模的结构。若 $f$ 固有，$A$ 是 Noether 环且
$\mathcal{F}$ 凝聚，则该模事实上是有限型的。

\begin{lemma}
\label{coherent-lemma-cohomology-powers-ideal-times-F}
\begin{reference}
\cite[III Cor 3.3.2]{EGA}
\end{reference}
设 $A$ 为 Noether 环，$I \subset A$ 为理想，并令
$B = \bigoplus_{n \geq 0} I^n$。设
$f : X \to \Spec(A)$ 为固有态射，并设凝聚层
$\mathcal{F}$（在 $X$ 上）。
则对每个 $p \geq 0$，分次 $B$-模
$\bigoplus_{n \geq 0} H^p(X, I^n\mathcal{F})$ 是有限 $B$-模。
\end{lemma}

\begin{proof}
令 $\mathcal{B} = f^*\widetilde{B}$。由于 $\mathcal{B}$ 满射到
$\bigoplus I^n\mathcal{O}_X$，显然
$\bigoplus I^n\mathcal{F}$ 是有限型分次 $\mathcal{B}$-模。
故结论由引理 \ref{coherent-lemma-graded-finiteness} 的 (1) 推出。
\end{proof}

\begin{lemma}
\label{coherent-lemma-cohomology-powers-ideal-times-sheaf}
给定概形态射 $f : X \to Y$、准凝聚层
$\mathcal{F}$（在 $X$ 上），以及准凝聚理想层
$\mathcal{I} \subset \mathcal{O}_Y$。假设 $Y$ 局部 Noether，
$f$ 固有且 $\mathcal{F}$ 凝聚。则
$$
\mathcal{M} =
\bigoplus\nolimits_{n \geq 0} R^pf_*(\mathcal{I}^n\mathcal{F})
$$
是准凝聚有限型分次
$\mathcal{A} = \bigoplus_{n \geq 0} \mathcal{I}^n$-模。
\end{lemma}

\begin{proof}
该陈述在 $Y$ 上是局部的，故可归约到 $Y$ 仿射的情形。
在仿射情形下，结论由引理
\ref{coherent-lemma-cohomology-powers-ideal-times-F} 推出。细节从略。
\end{proof}

\begin{lemma}
\label{coherent-lemma-cohomology-powers-ideal-application}
设 $A$ 为 Noether 环，$I \subset A$ 为理想，
$f : X \to \Spec(A)$ 为固有态射，并设
凝聚层 $\mathcal{F}$（在 $X$ 上）。则对每个 $p \geq 0$，
存在整数 $c \geq 0$，使得：
\begin{enumerate}
\item 乘法映射
$I^{n - c} \otimes H^p(X, I^c\mathcal{F}) \to H^p(X, I^n\mathcal{F})$
对所有 $n \geq c$ 都是满射；
\item 映射
$H^p(X, I^{n + m}\mathcal{F}) \to H^p(X, I^n\mathcal{F})$
之像包含于子模 $I^{m - e} H^p(X, I^n\mathcal{F})$，其中
$e = \max(0, c - n)$；这里 $n + m \geq c$ 且 $n, m \geq 0$；
\item 对 $n \geq c$，有
$$
\Ker(H^p(X, I^n\mathcal{F}) \to H^p(X, \mathcal{F})) =
\Ker(H^p(X, I^n\mathcal{F}) \to H^p(X, I^{n - c}\mathcal{F}))
$$
\item 存在映射
$I^nH^p(X, \mathcal{F}) \to H^p(X, I^{n - c}\mathcal{F})$（对 $n \geq c$），
使得当 $n \geq 2c$ 时，复合映射
$$
H^p(X, I^n\mathcal{F}) \to 
I^{n - c}H^p(X, \mathcal{F}) \to
H^p(X, I^{n - 2c}\mathcal{F})
$$
与
$$
I^nH^p(X, \mathcal{F}) \to
H^p(X, I^{n - c}\mathcal{F}) \to
I^{n - 2c}H^p(X, \mathcal{F})
$$
均为典范映射；并且
\item 逆向系统 $(H^p(X, I^n\mathcal{F}))$ 与
$(I^nH^p(X, \mathcal{F}))$ 是 pro-同构的。
\end{enumerate}
\end{lemma}

\begin{proof}
记 $M_n = H^p(X, I^n\mathcal{F})$（对 $n \geq 1$），并令
$M_0 = H^p(X, \mathcal{F})$，从而有映射
$\ldots \to M_3 \to M_2 \to M_1 \to M_0$。令
$B = \bigoplus_{n \geq 0} I^n$，则
$M = \bigoplus_{n \geq 0} M_n$ 是有限分次 $B$-模；参见引理
\ref{coherent-lemma-cohomology-powers-ideal-times-F}。
注意，乘积
$B_n \otimes M_m \to M_{m + n}$，$a \otimes m \mapsto a \cdot m$
与该逆向系统中的映射相容；更确切地说，对
$n, m' \geq 0$ 且 $m \geq m'$，交换图
$$
\xymatrix{
B_n \otimes_A M_m \ar[r] \ar[d] & M_{n + m} \ar[d] \\
B_n \otimes_A M_{m'} \ar[r] & M_{n + m'}
}
$$
交换。

\medskip\noindent
证明 (1)。选取 $d_1, \ldots, d_t \geq 0$ 及 $x_i \in M_{d_i}$，
使得 $M$ 由 $x_1, \ldots, x_t$ 在 $B$ 上生成。
对任意 $c \geq \max\{d_i\}$，可得
$B_{n - c} \cdot M_c = M_n$ 对 $n \geq c$ 成立，因而 (1) 成立。

\medskip\noindent
证明 (2)。令 $c$ 如 (1) 的证明中所取，并设 $n + m \geq c$。
有 $M_{n + m} = B_{n + m - c} \cdot M_c$。
若 $c > n$，则利用 $M_c \to M_n$ 以及上面指出的乘积与转移映射的相容性，
可知 $M_{n + m} \to M_n$ 之像包含于 $I^{n + m - c}M_n$。
若 $c \leq n$，则写成
$M_{n + m} = B_m \cdot B_{n - c} \cdot M_c =
B_m \cdot M_n$，从而看出其像包含于 $I^m M_n$。这证明了 (2)。

\medskip\noindent
令 $K_n \subset M_n$ 为映射 $M_n \to M_0$ 的核。
上面指出的乘积与转移映射的相容性表明
$K = \bigoplus K_n \subset M$ 是分次 $B$-子模。
由于 $B$ 是 Noether 的，且 $M$ 是有限生成分次 $B$-模，
$K$ 是有限生成分次 $B$-模。选取
$d'_1, \ldots, d'_{t'} \geq 0$ 及 $y_i \in K_{d'_i}$，
使得 $K$ 由 $y_1, \ldots, y_{t'}$ 在 $B$ 上生成。
令 $c = \max(d'_i, d'_j)$。由于
$y_i \in \Ker(M_{d'_i} \to M_0)$，可知
$B_n \cdot y_i \subset \Ker(M_{n + d'_i} \to M_n)$。
由此可见，$K_n = \Ker(M_n \to M_{n - c})$
对 $n \geq c$ 成立。这证明了 (3)。

\medskip\noindent
考虑下列实线部分交换的图
$$
\xymatrix{
I^n \otimes_A M_0 \ar[r] \ar[d] &
I^nM_0 \ar[r] \ar@{..>}[d] &
M_0 \ar[d] \\
M_n \ar[r] &
M_{n - c} \ar[r] &
M_0
}
$$
由上可知，满射 $I^n \otimes_A M_0 \to I^nM_0$ 的核映入
$K_n$，故得到虚线箭头，并且复合映射
$I^nM_0 \to M_{n - c} \to M_0$ 是典范映射。
于是，复合映射
$I^nM_0 \to M_{n - c} \to I^{n - 2c}M_0$
对 $n \geq 2c$ 显然是典范映射。
考虑复合映射 $M_n \to I^{n - c}M_0 \to M_{n - 2c}$。
第一映射把形如 $a \cdot m$ 的元素（其中
$a \in I^{n - c}$ 且 $m \in M_c$）送到 $a m'$，这里
$m'$ 是 $m$ 在 $M_0$ 中的像。第二映射再把它送到
$a \cdot m'$（位于 $M_{n - 2c}$ 中），故 (4) 成立。

\medskip\noindent
(5) 是 (4) 与 pro-对象态射之定义的直接推论。
\end{proof}

\noindent
在引理 \ref{coherent-lemma-cohomology-powers-ideal-times-F} 与
\ref{coherent-lemma-cohomology-powers-ideal-application} 的情形下，考虑逆向系统
$$
\mathcal{F}/I\mathcal{F} \leftarrow
\mathcal{F}/I^2\mathcal{F} \leftarrow
\mathcal{F}/I^3\mathcal{F} \leftarrow \ldots
$$
我们希望知道上同调群会发生什么。下面是第一个结论。

\begin{lemma}
\label{coherent-lemma-ML-cohomology-powers-ideal}
设 $A$ 为 Noether 环，$I \subset A$ 为理想，
$f : X \to \Spec(A)$ 为固有态射，并设凝聚层
$\mathcal{F}$（在 $X$ 上）。固定 $p \geq 0$。
存在 $c \geq 0$，使得：
\begin{enumerate}
\item 对所有 $n \geq c$，有
$$
\Ker(H^p(X, \mathcal{F}) \to H^p(X, \mathcal{F}/I^n\mathcal{F})) \subset
I^{n - c}H^p(X, \mathcal{F}).
$$
\item 逆向系统
$$
\left(H^p(X, \mathcal{F}/I^n\mathcal{F})\right)_{n \in \mathbf{N}}
$$
满足 Mittag--Leffler 条件（参见《同调》定义
\ref{homology-definition-Mittag-Leffler}）；并且
\item 有
$$
\Im(H^p(X, \mathcal{F}/I^k\mathcal{F})
\to H^p(X, \mathcal{F}/I^n\mathcal{F}))
=
\Im(H^p(X, \mathcal{F})
\to H^p(X, \mathcal{F}/I^n\mathcal{F}))
$$
这对所有 $k \geq n + c$ 成立。
\end{enumerate}
\end{lemma}

\begin{proof}
令 $c = \max\{c_p, c_{p + 1}\}$，其中 $c_p, c_{p + 1}$ 是引理
\ref{coherent-lemma-cohomology-powers-ideal-application} 分别对
$H^p$ 与 $H^{p + 1}$ 给出的整数。

\medskip\noindent
先证明 (1)。考虑短正合列
$$
0 \to I^n\mathcal{F} \to \mathcal{F} \to \mathcal{F}/I^n\mathcal{F} \to 0
$$
由上同调长正合列可知
$$
\Ker(
H^p(X, \mathcal{F}) \to H^p(X, \mathcal{F}/I^n\mathcal{F})
)
=
\Im(
H^p(X, I^n\mathcal{F}) \to H^p(X, \mathcal{F})
)
$$
因此，由引理 \ref{coherent-lemma-cohomology-powers-ideal-application} 的 (2)，
它包含于 $I^{n - c}H^p(X, \mathcal{F})$，这对 $n \geq c$ 成立。

\medskip\noindent
注意，(3) 按 Mittag--Leffler 系统的定义蕴含 (2)。

\medskip\noindent
下面证明 (3)。固定 $n$。考虑交换图
$$
\xymatrix{
0 \ar[r] &
I^n\mathcal{F} \ar[r] &
\mathcal{F} \ar[r] &
\mathcal{F}/I^n\mathcal{F} \ar[r] &
0 \\
0 \ar[r] &
I^{n + m}\mathcal{F} \ar[r] \ar[u] &
\mathcal{F} \ar[r] \ar[u] &
\mathcal{F}/I^{n + m}\mathcal{F} \ar[r] \ar[u] &
0
}
$$
它诱导出如下交换图
$$
\xymatrix{
H^p(X, \mathcal{F}) \ar[r] &
H^p(X, \mathcal{F}/I^n\mathcal{F}) \ar[r]_\delta &
H^{p + 1}(X, I^n\mathcal{F}) \ar[r] &
H^{p + 1}(X, \mathcal{F}) \\
H^p(X, \mathcal{F}) \ar[r] \ar[u]^1 &
H^p(X, \mathcal{F}/I^{n + m}\mathcal{F}) \ar[r] \ar[u]^\gamma &
H^{p + 1}(X, I^{n + m}\mathcal{F}) \ar[u]^\alpha \ar[r]^-\beta &
H^{p + 1}(X, \mathcal{F}) \ar[u]_1
}
$$
其两行正合。由引理
\ref{coherent-lemma-cohomology-powers-ideal-application} 的 (4)，$\beta$ 的核等于
$\alpha$ 的核，这对 $m \geq c$ 成立。
追图可知 $\gamma$ 之像包含于 $\delta$ 之核，故 (3) 成立
（令 $k = n + m$ 即得）。
\end{proof}

\begin{theorem}[形式函数定理]
\label{coherent-theorem-formal-functions}
设 $A$ 为 Noether 环，$I \subset A$ 为理想，
$f : X \to \Spec(A)$ 为固有态射，并设凝聚层
$\mathcal{F}$（在 $X$ 上）。固定 $p \geq 0$。
映射系
$$
H^p(X, \mathcal{F})/I^nH^p(X, \mathcal{F})
\longrightarrow
H^p(X, \mathcal{F}/I^n\mathcal{F})
$$
定义极限之间的同构
$$
H^p(X, \mathcal{F})^\wedge
\longrightarrow
\lim_n H^p(X, \mathcal{F}/I^n\mathcal{F})
$$
其中左端是 $A$-模 $H^p(X, \mathcal{F})$ 关于理想 $I$ 的完备化；
参见《代数》第 \ref{algebra-section-completion} 节。
而且，对极限拓扑而言，这事实上是同胚。
\end{theorem}

\begin{proof}
这可由引理 \ref{coherent-lemma-ML-cohomology-powers-ideal} 推出如下。
令 $M = H^p(X, \mathcal{F})$、
$M_n = H^p(X, \mathcal{F}/I^n\mathcal{F})$，并记
$N_n = \Im(M \to M_n)$。由引理
\ref{coherent-lemma-ML-cohomology-powers-ideal} 的 (2) 与 (3)，可知
$(M_n)$ 是 Mittag--Leffler 系统，并且 $N_n \subset M_n$
等于 $M_k$ 的像，这对所有 $k \gg n$ 成立。因此，作为具有极限拓扑的拓扑模，
$\lim M_n = \lim N_n$。另一方面，诸 $N_n$ 构成模 $M$
之商模的逆向系统，故 $\lim N_n$ 是 $M$ 关于诸核
$K_n = \Ker(M \to N_n)$ 所给拓扑的完备化。
由引理 \ref{coherent-lemma-ML-cohomology-powers-ideal} 的 (1)，有
$K_n \subset I^{n - c}M$；又由于 $N_n \subset M_n$ 被 $I^n$ 零化，
故 $I^n M \subset K_n$。因此，以子模 $K_n$ 作为 $0$ 的开邻域基本系
所定义的拓扑与 $I$-进拓扑相同，并得到诱导映射
$M^\wedge = \lim M/I^nM \to \lim N_n = \lim M_n$
是拓扑模的同构\footnote{
为确保无误，$M^\wedge$ 上的极限拓扑与其 $I$-进拓扑相同；
这由《代数》引理 \ref{algebra-lemma-hathat-finitely-generated} 推出。
另见《代数进阶》第 \ref{more-algebra-section-topological-ring} 节。}。
\end{proof}

\begin{lemma}
\label{coherent-lemma-spell-out-theorem-formal-functions}
设 $A$ 为环，$I \subset A$ 为理想。假设 $A$ 是 Noether 的且关于
$I$ 完备。设 $f : X \to \Spec(A)$ 为固有态射，
并设凝聚层 $\mathcal{F}$（在 $X$ 上）。则
$$
H^p(X, \mathcal{F}) = \lim_n H^p(X, \mathcal{F}/I^n\mathcal{F})
$$
对所有 $p \geq 0$ 成立。
\end{lemma}

\begin{proof}
这是形式函数定理（定理 \ref{coherent-theorem-formal-functions}）在完备 Noether
基环情形下的重新表述。事实上，在此情形下，$A$-模
$H^p(X, \mathcal{F})$ 是有限的
（引理 \ref{coherent-lemma-proper-over-affine-cohomology-finite}），因而是
$I$-进完备的（《代数》引理 \ref{algebra-lemma-completion-tensor}）；
由此可见左端不必再取完备化。
\end{proof}

\begin{lemma}
\label{coherent-lemma-formal-functions-stalk}
给定概形态射 $f : X \to Y$ 以及准凝聚层
$\mathcal{F}$（在 $X$ 上）。假设
\begin{enumerate}
\item $Y$ 局部 Noether；
\item $f$ 固有；并且
\item $\mathcal{F}$ 凝聚。
\end{enumerate}
设 $y \in Y$ 为一点。考虑纤维 $X_1 = X_y$ 的无穷小邻域
$$
\xymatrix{
X_n =
\Spec(\mathcal{O}_{Y, y}/\mathfrak m_y^n) \times_Y X
\ar[r]_-{i_n} \ar[d]_{f_n} &
X \ar[d]^f \\
\Spec(\mathcal{O}_{Y, y}/\mathfrak m_y^n) \ar[r]^-{c_n} & Y
}
$$
并令 $\mathcal{F}_n = i_n^*\mathcal{F}$。则有
$$
\left(R^pf_*\mathcal{F}\right)_y^\wedge
\cong
\lim_n H^p(X_n, \mathcal{F}_n)
$$
这是 $\mathcal{O}_{Y, y}^\wedge$-模的同构。
\end{lemma}

\begin{proof}
这只是形式函数定理，即定理 \ref{coherent-theorem-formal-functions}，
一个特殊情形的重新表述。下面详述之。注意
$\mathcal{O}_{Y, y}$ 是 Noether 局部环。考虑典范态射
$c : \Spec(\mathcal{O}_{Y, y}) \to Y$；参见《概形》公式
(\ref{schemes-equation-canonical-morphism})。该态射是平坦的，
因为它给出局部环的等同。暂记
$f' : X' \to \Spec(\mathcal{O}_{Y, y})$ 为 $f$ 到这个局部环的基变换。
由引理 \ref{coherent-lemma-flat-base-change-cohomology} 可知
$c^*R^pf_*\mathcal{F} = R^pf'_*\mathcal{F}'$。
此外，纤维 $X_y$ 与 $X'_y$ 的无穷小邻域彼此等同
（验证从略；提示：态射 $c_n$ 经由 $c$ 分解）。

\medskip\noindent
因此可以假设 $Y = \Spec(A)$ 是 Noether 局部环 $A$ 的谱，
其极大理想为 $\mathfrak m$，并且 $y \in Y$ 对应于闭点
（即对应于 $\mathfrak m$）。特别地，随即有
$$
\left(R^pf_*\mathcal{F}\right)_y =
\Gamma(Y, R^pf_*\mathcal{F}) =
H^p(X, \mathcal{F}).
$$

\medskip\noindent
在这种情形下，诸态射 $c_n$ 也全都是闭浸入，
故其基变换 $i_n$ 也都是闭浸入。注意
$i_{n, *}\mathcal{F}_n = i_{n, *}i_n^*\mathcal{F}
= \mathcal{F}/\mathfrak m^n\mathcal{F}$。由 $i_n$ 的 Leray 谱序列及引理
\ref{coherent-lemma-finite-pushforward-coherent}，可知
$$
H^p(X_n, \mathcal{F}_n) =
H^p(X, i_{n, *}\mathcal{F}_n) =
H^p(X, \mathcal{F}/\mathfrak m^n\mathcal{F})
$$
因此确可应用形式函数定理计算该引理陈述中的极限，结论得证。
\end{proof}

\noindent
下面这个引理稍后会推广到维数 $ > 0$ 的纤维，即再下一个引理。

\begin{lemma}
\label{coherent-lemma-higher-direct-images-zero-finite-fibre}
设 $f : X \to Y$ 为概形态射，并设 $y \in Y$。假设
\begin{enumerate}
\item $Y$ 局部 Noether；
\item $f$ 固有；并且
\item $f^{-1}(\{y\})$ 有限。
\end{enumerate}
则对任意凝聚层 $\mathcal{F}$（在 $X$ 上），都有
$(R^pf_*\mathcal{F})_y = 0$ 对所有 $p > 0$ 成立。
\end{lemma}

\begin{proof}
纤维 $X_y$ 是有限的，并且由《态射》引理
\ref{morphisms-lemma-finite-fibre}，它是有限离散空间。
而且，每个无穷小邻域 $X_n$ 的底拓扑空间都与之相同。
因此，由《概形》引理
\ref{schemes-lemma-scheme-finite-discrete-affine}，每个概形 $X_n$ 都仿射。
所以 $H^p(X_n, \mathcal{F}_n) = 0$ 对所有 $p > 0$ 成立。
于是由引理 \ref{coherent-lemma-formal-functions-stalk} 可知
$(R^pf_*\mathcal{F})_y^\wedge = 0$。注意，由命题
\ref{coherent-proposition-proper-pushforward-coherent}，
$R^pf_*\mathcal{F}$ 是凝聚的，因而 $R^pf_*\mathcal{F}_y$ 是有限
$\mathcal{O}_{Y, y}$-模。由 Nakayama 引理
（《代数》引理 \ref{algebra-lemma-NAK}），若局部环上有限模的完备化为零，
则该模为零。因此 $(R^pf_*\mathcal{F})_y = 0$。
\end{proof}


\begin{lemma}
\label{coherent-lemma-higher-direct-images-zero-above-dimension-fibre}
设 $f : X \to Y$ 为概形态射，并设 $y \in Y$。假设
\begin{enumerate}
\item $Y$ 局部 Noether；
\item $f$ 固有；并且
\item $\dim(X_y) = d$。
\end{enumerate}
则对任意凝聚层 $\mathcal{F}$（在 $X$ 上），都有
$(R^pf_*\mathcal{F})_y = 0$ 对所有 $p > d$ 成立。
\end{lemma}

\begin{proof}
纤维 $X_y$ 在 $\Spec(\kappa(y))$ 上是有限型的。
故由《态射》引理 \ref{morphisms-lemma-finite-type-noetherian}，
$X_y$ 是 Noether 概形。因此，由《性质》引理
\ref{properties-lemma-Noetherian-topology}，$X_y$ 的底拓扑空间是 Noether 的。
而且，每个无穷小邻域 $X_n$ 的底拓扑空间与 $X_y$ 的相同。
故由《上同调》命题 \ref{cohomology-proposition-vanishing-Noetherian}，
$H^p(X_n, \mathcal{F}_n) = 0$ 对所有 $p > d$ 成立。
于是可知 $(R^pf_*\mathcal{F})_y^\wedge = 0$；
由引理 \ref{coherent-lemma-formal-functions-stalk}，这对 $p > d$ 成立。注意，由命题
\ref{coherent-proposition-proper-pushforward-coherent}，
$R^pf_*\mathcal{F}$ 是凝聚的，因而 $R^pf_*\mathcal{F}_y$ 是有限
$\mathcal{O}_{Y, y}$-模。由 Nakayama 引理
（《代数》引理 \ref{algebra-lemma-NAK}），若局部环上有限模的完备化为零，
则该模为零。因此 $(R^pf_*\mathcal{F})_y = 0$。
\end{proof}

\section{形式函数定理的应用}
\label{coherent-section-applications-formal-functions}


\noindent
我们会按需在这里增补更多内容。目前需要 Noether 情形下有限态射的如下刻画。

\begin{lemma}
\label{coherent-lemma-characterize-finite}
（更一般的版本见《态射进阶》引理
\ref{more-morphisms-lemma-characterize-finite}。）
设 $f : X \to S$ 为概形态射，并假设 $S$ 局部 Noether。
则下列条件等价：
\begin{enumerate}
\item $f$ 有限；
\item $f$ 固有且纤维有限。
\end{enumerate}
\end{lemma}

\begin{proof}
由《态射》引理 \ref{morphisms-lemma-finite-proper}，有限态射固有。
由《态射》引理 \ref{morphisms-lemma-finite-quasi-finite}，有限态射拟有限。
拟有限态射有有限纤维；参见《态射》引理
\ref{morphisms-lemma-quasi-finite}。故有限态射固有且有有限纤维。

\medskip\noindent
假设 $f$ 固有且有有限纤维。要证明 $f$ 有限。
事实上，只须证明 $f$ 仿射。确实，若 $f$ 仿射，则由《态射》引理
\ref{morphisms-lemma-integral-universally-closed} 可知 $f$ 整；
继而由《态射》引理 \ref{morphisms-lemma-finite-integral} 可知 $f$ 有限。

\medskip\noindent
为证明 $f$ 仿射，可以假设 $S$ 仿射；目标是证明 $X$ 也仿射。
由于 $f$ 固有，可知 $X$ 分离且拟紧。因此可以用引理
\ref{coherent-lemma-quasi-separated-h1-zero-covering} 的判据证明 $X$ 仿射。
为此，设 $\mathcal{I} \subset \mathcal{O}_X$ 为有限型理想层。
特别地，$\mathcal{I}$ 是 $X$ 上的凝聚层。由引理
\ref{coherent-lemma-higher-direct-images-zero-finite-fibre}，可得
$R^1f_*\mathcal{I}_s = 0$ 对所有 $s \in S$ 成立。
换言之，$R^1f_*\mathcal{I} = 0$。于是由 $f$ 的 Leray 谱序列可见
$H^1(X , \mathcal{I}) = H^1(S, f_*\mathcal{I})$。
由于 $S$ 仿射，且 $f_*\mathcal{I}$ 准凝聚
（《概形》引理 \ref{schemes-lemma-push-forward-quasi-coherent}），
由引理 \ref{coherent-lemma-quasi-coherent-affine-cohomology-zero} 可得所需的
$H^1(S, f_*\mathcal{I}) = 0$。故如所需地，
$H^1(X, \mathcal{I}) = 0$。
\end{proof}

\noindent
作为推论，有如下有用结论。

\begin{lemma}
\label{coherent-lemma-proper-finite-fibre-finite-in-neighbourhood}
\begin{slogan}
固有态射在有限纤维的一个邻域上有限。
\end{slogan}
（更一般的版本见《态射进阶》引理
\ref{more-morphisms-lemma-proper-finite-fibre-finite-in-neighbourhood}。）
设 $f : X \to S$ 为概形态射，并设 $s \in S$。假设
\begin{enumerate}
\item $S$ 局部 Noether；
\item $f$ 固有；并且
\item $f^{-1}(\{s\})$ 是有限集。
\end{enumerate}
则存在开邻域 $V \subset S$（为 $s$ 的邻域），使得
$f|_{f^{-1}(V)} : f^{-1}(V) \to V$ 有限。
\end{lemma}

\begin{proof}
由《态射》引理 \ref{morphisms-lemma-finite-fibre}，态射 $f$
在 $f^{-1}(\{s\})$ 的所有点处拟有限。由《态射》引理
\ref{morphisms-lemma-quasi-finite-points-open}，$f$ 拟有限的点构成开集
$U \subset X$。令 $Z = X \setminus U$。则 $s \not \in f(Z)$。
由于 $f$ 固有，集合 $f(Z) \subset S$ 闭。选取任意开邻域
$V \subset S$（为 $s$ 的邻域），使得 $Z \cap V = \emptyset$。于是
$f^{-1}(V) \to V$ 局部拟有限且固有，故其拟有限
（《态射》引理 \ref{morphisms-lemma-quasi-finite-locally-quasi-compact}），
从而有有限纤维（《态射》引理 \ref{morphisms-lemma-quasi-finite}），
继而由引理 \ref{coherent-lemma-characterize-finite} 可知它有限。
\end{proof}

\begin{lemma}
\label{coherent-lemma-ample-on-fibre}
设 $f : X \to Y$ 为概形的固有态射，其中 $Y$ Noether。
设 $\mathcal{L}$ 为可逆 $\mathcal{O}_X$-模，
$\mathcal{F}$ 为凝聚 $\mathcal{O}_X$-模。
设 $y \in Y$ 为一点，使得 $\mathcal{L}_y$ 在 $X_y$ 上丰沛。
则存在 $d_0$，使得对所有 $d \geq d_0$，都有
$$
R^pf_*(\mathcal{F} \otimes_{\mathcal{O}_X} \mathcal{L}^{\otimes d})_y = 0
\text{ 对 }p > 0
$$
且映射
$$
f_*(\mathcal{F} \otimes_{\mathcal{O}_X} \mathcal{L}^{\otimes d})_y
\longrightarrow
H^0(X_y, \mathcal{F}_y \otimes_{\mathcal{O}_{X_y}} \mathcal{L}_y^{\otimes d})
$$
为满射。
\end{lemma}

\begin{proof}
注意，$\mathcal{O}_{Y, y}$ 是 Noether 局部环。考虑典范态射
$c : \Spec(\mathcal{O}_{Y, y}) \to Y$；参见《概形》公式
(\ref{schemes-equation-canonical-morphism})。该态射是平坦的，
因为它给出局部环的等同。暂记
$f' : X' \to \Spec(\mathcal{O}_{Y, y})$ 为 $f$ 到这个局部环的基变换。
由引理 \ref{coherent-lemma-flat-base-change-cohomology} 可知
$c^*R^pf_*\mathcal{F} = R^pf'_*\mathcal{F}'$。
此外，纤维 $X_y$ 与 $X'_y$ 彼此等同。
因此可以假设 $Y = \Spec(A)$ 是 Noether 局部环
$(A, \mathfrak m, \kappa)$ 的谱，且 $y \in Y$ 对应于
$\mathfrak m$。在这种情形下，
$R^pf_*(\mathcal{F} \otimes_{\mathcal{O}_X} \mathcal{L}^{\otimes d})_y =
H^p(X, \mathcal{F} \otimes_{\mathcal{O}_X} \mathcal{L}^{\otimes d})$
对所有 $p \geq 0$ 成立。记投影为
$f_y : X_y \to \Spec(\kappa)$。

\medskip\noindent
令 $B = \text{Gr}_\mathfrak m(A) =
\bigoplus_{n \geq 0} \mathfrak m^n/\mathfrak m^{n + 1}$。
考虑层 $\mathcal{B} = f_y^*\widetilde{B}$；它是准凝聚分次
$\mathcal{O}_{X_y}$-代数层。采用第
\ref{coherent-section-theorem-formal-functions} 节中的记号，并以
$I$ 由 $\mathfrak m$ 代替。由于 $X_y$ 是 $X$ 中由
$\mathfrak m\mathcal{O}_X$ 截出的闭子概形，可以把
$\mathfrak m^n\mathcal{F}/\mathfrak m^{n + 1}\mathcal{F}$
视为凝聚 $\mathcal{O}_{X_y}$-模；参见引理
\ref{coherent-lemma-i-star-equivalence}。于是
$\bigoplus_{n \geq 0} \mathfrak m^n\mathcal{F}/\mathfrak m^{n + 1}\mathcal{F}$
是有限型准凝聚分次 $\mathcal{B}$-模；这是因为它在
$\mathcal{B}$ 上由零次部分生成，并且其零次部分是
$\mathcal{F}_y = \mathcal{F}/\mathfrak m \mathcal{F}$，
这是凝聚 $\mathcal{O}_{X_y}$-模。故由引理
\ref{coherent-lemma-graded-finiteness} 的 (2)，可知
$$
H^p(X_y, \mathfrak m^n \mathcal{F}/ \mathfrak m^{n + 1}\mathcal{F}
\otimes_{\mathcal{O}_{X_y}} \mathcal{L}_y^{\otimes d}) = 0
$$
对所有 $p > 0$、$d \geq d_0$ 及 $n \geq 0$ 成立。
由引理 \ref{coherent-lemma-relative-affine-cohomology}，这等价于说
{\small
$
H^p(X, \mathfrak m^n \mathcal{F}/ \mathfrak m^{n + 1}\mathcal{F}
\otimes_{\mathcal{O}_X} \mathcal{L}^{\otimes d}) = 0
$}
对所有 $p > 0$、$d \geq d_0$ 及 $n \geq 0$ 成立。

\medskip\noindent
考虑凝聚 $\mathcal{O}_X$-模的短正合列
$$
0 \to \mathfrak m^n\mathcal{F}/\mathfrak m^{n + 1} \mathcal{F}
\to \mathcal{F}/\mathfrak m^{n + 1} \mathcal{F}
\to \mathcal{F}/\mathfrak m^n \mathcal{F} \to 0
$$
与 $\mathcal{L}^{\otimes d}$ 张量是正合函子，因而得到短正合列
$$
0 \to
\mathfrak m^n\mathcal{F}/\mathfrak m^{n + 1} \mathcal{F}
\otimes_{\mathcal{O}_X} \mathcal{L}^{\otimes d}
\to \mathcal{F}/\mathfrak m^{n + 1} \mathcal{F}
\otimes_{\mathcal{O}_X} \mathcal{L}^{\otimes d}
\to \mathcal{F}/\mathfrak m^n \mathcal{F}
\otimes_{\mathcal{O}_X} \mathcal{L}^{\otimes d} \to 0
$$
利用上同调长正合列和上面的消失，通过归纳可得：
\begin{enumerate}
\item $H^p(X, \mathcal{F}/\mathfrak m^n \mathcal{F}
\otimes_{\mathcal{O}_X} \mathcal{L}^{\otimes d}) = 0$
对所有 $p > 0$、$d \geq d_0$ 及 $n \geq 0$ 成立；并且
\item $H^0(X, \mathcal{F}/\mathfrak m^n \mathcal{F}
\otimes_{\mathcal{O}_X} \mathcal{L}^{\otimes d}) \to
H^0(X_y, \mathcal{F}_y \otimes_{\mathcal{O}_{X_y}} \mathcal{L}_y^{\otimes d})$
对所有 $d \geq d_0$ 及 $n \geq 1$ 都是满射。
\end{enumerate}
由形式函数定理（定理 \ref{coherent-theorem-formal-functions}），可知
$\mathfrak m$-进完备化（对
$H^p(X, \mathcal{F} \otimes_{\mathcal{O}_X} \mathcal{L}^{\otimes d})$）
在所有 $d \geq d_0$ 与 $p > 0$ 时都为零。
由于 $H^p(X, \mathcal{F} \otimes_{\mathcal{O}_X} \mathcal{L}^{\otimes d})$
是有限 $A$-模（引理 \ref{coherent-lemma-proper-over-affine-cohomology-finite}），
由 Nakayama 引理（《代数》引理 \ref{algebra-lemma-NAK}）可知
$H^p(X, \mathcal{F} \otimes_{\mathcal{O}_X} \mathcal{L}^{\otimes d})$
对所有 $d \geq d_0$ 与 $p > 0$ 都为零。
对 $p = 0$，由引理 \ref{coherent-lemma-ML-cohomology-powers-ideal} 的 (3) 可得
$H^0(X, \mathcal{F} \otimes_{\mathcal{O}_X} \mathcal{L}^{\otimes d}) \to
H^0(X_y, \mathcal{F}_y \otimes_{\mathcal{O}_{X_y}} \mathcal{L}_y^{\otimes d})$
为满射，这就给出了引理的最后一项陈述。
\end{proof}

\begin{lemma}
\label{coherent-lemma-ample-in-neighbourhood}
（更一般的版本见《态射进阶》引理
\ref{more-morphisms-lemma-ample-in-neighbourhood}。）
设 $f : X \to Y$ 为概形的固有态射，其中 $Y$ Noether。
设 $\mathcal{L}$ 为可逆 $\mathcal{O}_X$-模。
设 $y \in Y$ 为一点，使得 $\mathcal{L}_y$ 在 $X_y$ 上丰沛。
则存在开邻域 $V \subset Y$（为 $y$ 的邻域），使得
$\mathcal{L}|_{f^{-1}(V)}$ 在 $f^{-1}(V)/V$ 上丰沛。
\end{lemma}

\begin{proof}
选取 $d_0$，如引理 \ref{coherent-lemma-ample-on-fibre} 在
$\mathcal{F} = \mathcal{O}_X$ 时所给。选取 $d \geq d_0$，
使得可以找到 $r \geq 0$ 及截面
$s_{y, 0}, \ldots, s_{y, r} \in H^0(X_y, \mathcal{L}_y^{\otimes d})$，
它们定义闭浸入
$$
\varphi_y =
\varphi_{\mathcal{L}_y^{\otimes d}, (s_{y, 0}, \ldots, s_{y, r})} :
X_y \to \mathbf{P}^r_{\kappa(y)}.
$$
这可由《态射》引理
\ref{morphisms-lemma-finite-type-over-affine-ample-very-ample} 做到；
还要用《态射》引理
\ref{morphisms-lemma-image-proper-scheme-closed} 看出 $\varphi_y$
是闭浸入，并用《构造》第
\ref{constructions-section-projective-space} 节关于以可逆层与截面描述
到射影空间之态射的结论。由 $d_0$ 的选取，以 $Y$ 的一个开邻域
（为 $y$ 的邻域）替换后，可以选取
$s_0, \ldots, s_r \in H^0(X, \mathcal{L}^{\otimes d})$，
它们分别映到 $s_{y, 0}, \ldots, s_{y, r}$。
令 $X_{s_i} \subset X$ 为 $s_i$ 生成
$\mathcal{L}^{\otimes d}$ 的开子集。由于诸
$s_{y, i}$ 生成 $\mathcal{L}_y^{\otimes d}$，可见
$X_y \subset U = \bigcup X_{s_i}$。由于 $X \to Y$ 是闭映射，
可知存在开邻域 $y \in V \subset Y$，使得
$f^{-1}(V) \subset U$。将 $Y$ 替换为 $V$ 后，可以假设诸
$s_i$ 生成 $\mathcal{L}^{\otimes d}$。于是得到态射
$$
\varphi = \varphi_{\mathcal{L}^{\otimes d}, (s_0, \ldots, s_r)} :
X \longrightarrow \mathbf{P}^r_Y
$$
它满足 $\mathcal{L}^{\otimes d} \cong \varphi^*\mathcal{O}_{\mathbf{P}^r_Y}(1)$，
且其到 $y$ 的基变换给出 $\varphi_y$。

\medskip\noindent
我们以一个小技巧完成证明；“正确的”证明是直接说明，$\varphi$
在基变换到 $y$ 的某个开邻域后是闭浸入。具体而言，由引理
\ref{coherent-lemma-proper-finite-fibre-finite-in-neighbourhood}，可知 $\varphi$
在纤维 $\mathbf{P}^r_{\kappa(y)}$ 的一个开邻域上有限；这里该纤维是
$\mathbf{P}^r_Y \to Y$ 位于 $y$ 上方的纤维。
利用 $\mathbf{P}^r_Y \to Y$ 是闭映射，缩小 $Y$ 后，
可以假设 $\varphi$ 有限。于是
$\mathcal{L}^{\otimes d} \cong \varphi^*\mathcal{O}_{\mathbf{P}^r_Y}(1)$
由非常一般的《态射》引理
\ref{morphisms-lemma-pullback-ample-tensor-relatively-ample} 可知是丰沛的。
\end{proof}

\section{上同调与基变换（三）}
\label{coherent-section-cohomology-and-base-change-perfect}

\noindent
本节证明一个十分一般的现象之最简单情形；该现象将在《概形的导出范畴》第
\ref{perfect-section-cohomology-and-base-change-perfect} 节讨论。
关于把下列引理翻译成代数语言的表述，见评注
\ref{coherent-remark-explain-perfect-direct-image}。

\begin{lemma}
\label{coherent-lemma-perfect-direct-image}
设 $A$ 为 Noether 环，并令 $S = \Spec(A)$。设
$f : X \to S$ 为概形的固有态射。设 $\mathcal{F}$ 为凝聚
$\mathcal{O}_X$-模，并且在 $S$ 上平坦。则
\begin{enumerate}
\item $R\Gamma(X, \mathcal{F})$ 是 $D(A)$ 的完美对象；并且
\item 对任意环同态 $A \to A'$，基变换映射
$$
R\Gamma(X, \mathcal{F}) \otimes_A^{\mathbf{L}} A'
\longrightarrow
R\Gamma(X_{A'}, \mathcal{F}_{A'})
$$
是同构。
\end{enumerate}
\end{lemma}

\begin{proof}
选取有限仿射开覆盖 $X = \bigcup_{i = 1, \ldots, n} U_i$。
由引理 \ref{coherent-lemma-separated-case-relative-cech} 与
\ref{coherent-lemma-base-change-complex}，{\v C}ech 复形
$K^\bullet = {\check C}^\bullet(\mathcal{U}, \mathcal{F})$ 满足
$$
K^\bullet \otimes_A A' = R\Gamma(X_{A'}, \mathcal{F}_{A'})
$$
这对所有环同态 $A \to A'$ 成立。令
$K_{alt}^\bullet = {\check C}_{alt}^\bullet(\mathcal{U}, \mathcal{F})$
为交错 {\v C}ech 复形。由《上同调》引理
\ref{cohomology-lemma-alternating-usual}，存在同伦等价
$K_{alt}^\bullet \to K^\bullet$（在 $A$-模范畴中）。特别地，同样有
$$
K_{alt}^\bullet \otimes_A A' = R\Gamma(X_{A'}, \mathcal{F}_{A'})
$$
由于 $\mathcal{F}$ 在 $A$ 上平坦，可见每个 $K_{alt}^n$
在 $A$ 上平坦（参见《态射》引理
\ref{morphisms-lemma-flat-module-characterize}）。而且，
$K_{alt}^\bullet$ 有上界（这正是改用交错 {\v C}ech 复形的原因），故
$K_{alt}^\bullet \otimes_A A' = K_{alt}^\bullet \otimes_A^{\mathbf{L}} A'$；
这是导出张量积定义的结果（参见《代数进阶》第
\ref{more-algebra-section-derived-tensor-product} 节）。由引理
\ref{coherent-lemma-proper-over-affine-cohomology-finite}，上同调群
$H^i(K_{alt}^\bullet)$ 是有限 $A$-模。由于
$K_{alt}^\bullet$ 有界，可得 $K_{alt}^\bullet$ 是伪凝聚的；
参见《代数进阶》引理
\ref{more-algebra-lemma-Noetherian-pseudo-coherent}。
给定任意 $A$-模 $M$，令 $A' = A \oplus M$，其中 $M$ 是平方零理想，
即 $(a, m) \cdot (a', m') = (aa', am' + a'm)$。
由上可见，$K_{alt}^\bullet \otimes_A^\mathbf{L} A'$
仅在次数 $0, \ldots, n$ 有上同调。因此，
$K_{alt}^\bullet \otimes_A^\mathbf{L} M$ 仅在次数
$0, \ldots, n$ 有上同调。故 $K_{alt}^\bullet$ 有有限 Tor 维数；
参见《代数进阶》定义 \ref{more-algebra-definition-tor-amplitude}。
最后应用《代数进阶》引理 \ref{more-algebra-lemma-perfect} 即得结论。
\end{proof}

\begin{remark}
\label{coherent-remark-explain-perfect-direct-image}
引理 \ref{coherent-lemma-perfect-direct-image} 的一个推论是：存在由有限射影
$A$-模组成的有限复形 $M^\bullet$，使得
$$
H^i(X_{A'}, \mathcal{F}_{A'}) = H^i(M^\bullet \otimes_A A')
$$
关于 $A'$ 函子性地成立。$\mathcal{F}$ 在 $A$ 上平坦这一条件是本质的；
参见 \cite{Hartshorne}。
\end{remark}

\section{凝聚形式模}
\label{coherent-section-coherent-formal}

\noindent
由于尚无形式概形理论可供使用，我们发展一些语言来代替
“Noether 进形式概形上的凝聚模”这一概念。

\medskip\noindent
设 $X$ 为 Noether 概形，并设 $\mathcal{I} \subset \mathcal{O}_X$
为准凝聚理想层。我们考虑逆向系统 $(\mathcal{F}_n)$，
其各项是凝聚 $\mathcal{O}_X$-模，并满足：
\begin{enumerate}
\item $\mathcal{F}_n$ 被 $\mathcal{I}^n$ 零化；并且
\item 转移映射诱导同构
$\mathcal{F}_{n + 1}/\mathcal{I}^n\mathcal{F}_{n + 1} \to \mathcal{F}_n$。
\end{enumerate}
这类逆向系统的态射按通常方式定义。以
$\textit{Coh}(X, \mathcal{I})$ 表示这些逆向系统所成的范畴。
下面将证明一系列关于该范畴中对象的引理。事实上，后续大多数引理
都是下述仿射情形范畴描述的直接推论。

\begin{lemma}
\label{coherent-lemma-inverse-systems-affine}
若 $X = \Spec(A)$ 是 Noether 环的谱，且 $\mathcal{I}$ 是与理想
$I \subset A$ 相伴的准凝聚理想层，则
$\textit{Coh}(X, \mathcal{I})$ 等价于有限 $A^\wedge$-模范畴，
其中 $A^\wedge$ 是 $A$ 关于 $I$ 的完备化。
\end{lemma}

\begin{proof}
以 $\text{Mod}^{fg}_{A, I}$ 表示逆向系统 $(M_n)$ 所成的范畴，
其中各项是有限 $A$-模；这些系统满足：(1) $M_n$ 被 $I^n$ 零化；
(2) $M_{n + 1}/I^nM_{n + 1} = M_n$。由 $X$ 上凝聚层与有限
$A$-模之间的对应（引理 \ref{coherent-lemma-coherent-Noetherian}），只须证明
$\text{Mod}^{fg}_{A, I}$ 等价于有限 $A^\wedge$-模范畴。
为此，只须证明：给定对象 $(M_n)$（属于
$\text{Mod}^{fg}_{A, I}$），模
$$
M = \lim M_n
$$
是有限 $A^\wedge$-模，且 $M/I^nM = M_n$。
由于转移映射都是满射，可见 $M \to M_1$ 是满射。
选取 $x_1, \ldots, x_t \in M$，使其像生成 $M_1$。
这诱导系统间的映射 $(A/I^n)^{\oplus t} \to M_n$。
由 Nakayama 引理（《代数》引理 \ref{algebra-lemma-NAK}），
这些映射均为满射。令 $K_n \subset (A/I^n)^{\oplus t}$ 为其核。
性质 (2) 蕴含 $K_{n + 1} \to K_n$ 是满射；特别地，系统
$(K_n)$ 满足 Mittag--Leffler 条件。由《同调》引理
\ref{homology-lemma-Mittag-Leffler}，得到正合列
$0 \to K \to (A^\wedge)^{\oplus t} \to M \to 0$，
其中 $K = \lim K_n$。故 $M$ 是有限 $A^\wedge$-模。
由于 $K \to K_n$ 为满射，得到
$$
M/I^nM = \Coker(K \to (A/I^n)^{\oplus t}) =
(A/I^n)^{\oplus t}/K_n = M_n
$$
如所需。
\end{proof}

\begin{lemma}
\label{coherent-lemma-inverse-systems-abelian}
设 $X$ 为 Noether 概形，并设 $\mathcal{I} \subset \mathcal{O}_X$
为准凝聚理想层。
\begin{enumerate}
\item 范畴 $\textit{Coh}(X, \mathcal{I})$ 是阿贝尔范畴。
\item 对开集 $U \subset X$，限制函子
$\textit{Coh}(X, \mathcal{I}) \to \textit{Coh}(U, \mathcal{I}|_U)$
是正合的。
\item $\textit{Coh}(X, \mathcal{I})$ 中的正合性可通过限制到
$X$ 的一个开覆盖之各成员上来检验。
\end{enumerate}
\end{lemma}

\begin{proof}
设 $\alpha =(\alpha_n) : (\mathcal{F}_n) \to (\mathcal{G}_n)$ 为
$\textit{Coh}(X, \mathcal{I})$ 的态射。$\alpha$ 的余核是逆向系统
$(\Coker(\alpha_n))$（细节从略）。为描述其核，对 $l \geq m$ 令
$$
\mathcal{K}'_{l, m} = \Im(\Ker(\alpha_l) \to \mathcal{F}_m)
$$
我们断言：
\begin{enumerate}
\item[(a)] 逆向系统 $(\mathcal{K}'_{l, m})_{l \geq m}$ 最终常值，
设其最终值为 $\mathcal{K}'_m$；
\item[(b)] 系统
$(\mathcal{K}'_m/\mathcal{I}^n\mathcal{K}'_m)_{m \geq n}$ 最终常值，
设其最终值为 $\mathcal{K}_n$；
\item[(c)] 系统 $(\mathcal{K}_n)$ 构成
$\textit{Coh}(X, \mathcal{I})$ 的一个对象；并且
\item[(d)] 该对象是 $\alpha$ 的核。
\end{enumerate}
为证明 (a)、(b)、(c)，可以在仿射局部工作；设
$X = \Spec(A)$，且 $\mathcal{I}$ 对应于理想 $I \subset A$。
由引理 \ref{coherent-lemma-inverse-systems-affine}，$\alpha$ 对应于映射
$f : M \to N$（在有限 $A^\wedge$-模之间）。记
$K = \Ker(f)$。注意 $A^\wedge$ 是 Noether 环
（《代数》引理 \ref{algebra-lemma-completion-Noetherian-Noetherian}）。
选取整数 $c \geq 0$，使得
$K \cap I^n M \subset I^{n - c}K$ 对 $n \geq c$ 成立
（《代数》引理 \ref{algebra-lemma-Artin-Rees}），并且使它对映射
$f$ 与理想 $I^\wedge = IA^\wedge$ 满足《代数》引理
\ref{algebra-lemma-map-AR}。于是，$\mathcal{K}'_{l, m}$ 对应于 $A$-模
$$
K'_{l, m} = \frac{a^{-1}(I^lN) + I^mM}{I^mM} =
\frac{K + I^{l - c}f^{-1}(I^cN) + I^mM}{I^mM} =
\frac{K + I^mM}{I^mM}
$$
其中最后一个等号在 $l \geq m + c$ 时成立。因此
$\mathcal{K}'_m$ 对应于 $A$-模 $K/K \cap I^mM$，且
$\mathcal{K}'_m/\mathcal{I}^n\mathcal{K}'_m$ 对应于
$$
\frac{K}{K \cap I^mM + I^nK} = \frac{K}{I^nK}
$$
当 $m \geq n + c$ 时，这由上面对 $c$ 的选取成立。
故 $\mathcal{K}_n$ 对应于 $K/I^nK$。

\medskip\noindent
下面证明 (d)。由上面的仿射描述，复合映射
$(\mathcal{K}_n) \to (\mathcal{F}_n) \to (\mathcal{G}_n)$
显然为零。设 $\beta : (\mathcal{H}_n) \to (\mathcal{F}_n)$
为满足 $\alpha \circ \beta = 0$ 的态射。则
$\mathcal{H}_l \to \mathcal{F}_l$ 的像包含于 $\Ker(\alpha_l)$。
由于 $\mathcal{H}_m = \mathcal{H}_l/\mathcal{I}^m\mathcal{H}_l$
对 $l \geq m$ 成立，
得到一族映射 $\mathcal{H}_m \to \mathcal{K}'_{l, m}$，
从而得到映射 $\mathcal{H}_m \to \mathcal{K}_m'$。
又由于
$\mathcal{H}_n = \mathcal{H}_m/\mathcal{I}^n\mathcal{H}_m$，
得到一族映射
$\mathcal{H}_n \to \mathcal{K}'_m/\mathcal{I}^n\mathcal{K}'_m$，
因而如所需地得到映射 $\mathcal{H}_n \to \mathcal{K}_n$。

\medskip\noindent
为完成 (1) 的证明，还须证明 $\Coim = \Im$
（在 $\textit{Coh}(X, \mathcal{I})$ 中）。上面已经看出，
在仿射情形下，取核与余核和把
$\textit{Coh}(X, \mathcal{I})$ 描述为模范畴相交换。
由于模范畴中 $\Im = \Coim$，这给出
$\Coim = \Im$（在 $\textit{Coh}(X, \mathcal{I})$ 中）。
引理的 (2)、(3) 由我们对核与余核的构造立即推出。
\end{proof}

\begin{lemma}
\label{coherent-lemma-inverse-systems-surjective}
设 $X$ 为 Noether 概形，并设 $\mathcal{I} \subset \mathcal{O}_X$
为准凝聚理想层。映射
$(\mathcal{F}_n) \to (\mathcal{G}_n)$ 在
$\textit{Coh}(X, \mathcal{I})$ 中为满射，当且仅当
$\mathcal{F}_1 \to \mathcal{G}_1$ 为满射。
\end{lemma}

\begin{proof}
从略。提示：在仿射开集上考察，使用引理
\ref{coherent-lemma-inverse-systems-affine} 及《代数》引理
\ref{algebra-lemma-NAK}。
\end{proof}

\noindent
设 $X$ 为 Noether 概形，并设 $\mathcal{I} \subset \mathcal{O}_X$
为准凝聚理想层。有函子
\begin{equation}
\label{coherent-equation-completion-functor}
\textit{Coh}(\mathcal{O}_X) \longrightarrow \textit{Coh}(X, \mathcal{I}), \quad
\mathcal{F} \longmapsto \mathcal{F}^\wedge
\end{equation}
它把凝聚 $\mathcal{O}_X$-模 $\mathcal{F}$ 送到对象
$\mathcal{F}^\wedge = (\mathcal{F}/\mathcal{I}^n\mathcal{F})$；
该对象属于 $\textit{Coh}(X, \mathcal{I})$。

\begin{lemma}
\label{coherent-lemma-exact}
函子 (\ref{coherent-equation-completion-functor}) 是正合的。
\end{lemma}

\begin{proof}
只须在 $X$ 上局部检验。因此可以假设 $X$ 仿射，
即处于引理 \ref{coherent-lemma-inverse-systems-affine} 的情形。
该函子就是函子
$\text{Mod}^{fg}_A \to \text{Mod}^{fg}_{A^\wedge}$，
它把有限 $A$-模 $M$ 送到完备化 $M^\wedge$。
故结论由《代数》引理 \ref{algebra-lemma-completion-flat} 推出。
\end{proof}

\begin{lemma}
\label{coherent-lemma-completion-internal-hom}
设 $X$ 为 Noether 概形，并设 $\mathcal{I} \subset \mathcal{O}_X$
为准凝聚理想层。设 $\mathcal{F}$、$\mathcal{G}$ 为凝聚
$\mathcal{O}_X$-模。令
$\mathcal{H} = \SheafHom_{\mathcal{O}_X}(\mathcal{G}, \mathcal{F})$。
则
$$
\lim H^0(X, \mathcal{H}/\mathcal{I}^n\mathcal{H}) =
\Mor_{\textit{Coh}(X, \mathcal{I})}
(\mathcal{G}^\wedge, \mathcal{F}^\wedge).
$$
\end{lemma}

\begin{proof}
为证明该结论，可以在 $X$ 上仿射局部地工作。
故可假设 $X = \Spec(A)$，并且 $\mathcal{F}$、$\mathcal{G}$
由有限 $A$-模 $M$、$N$ 给出。于是 $\mathcal{H}$ 对应于有限
$A$-模 $H = \Hom_A(M, N)$。经由引理
\ref{coherent-lemma-inverse-systems-affine} 的等价，引理的陈述化为
$$
H^\wedge = \Hom_{A^\wedge}(M^\wedge, N^\wedge)
$$
由《代数》引理 \ref{algebra-lemma-completion-flat}（使用三次），有
$$
H^\wedge = \Hom_A(M, N) \otimes_A A^\wedge =
\Hom_{A^\wedge}(M \otimes_A A^\wedge, N \otimes_A A^\wedge) =
\Hom_{A^\wedge}(M^\wedge, N^\wedge)
$$
其中第二个等号使用了 $A^\wedge$ 在 $A$ 上平坦这一事实
（参见《代数进阶》引理
\ref{more-algebra-lemma-pseudo-coherence-and-base-change-ext}）。
引理得证。
\end{proof}

\noindent
设 $X$ 为 Noether 概形。设 $\mathcal{I} \subset \mathcal{O}_X$
为准凝聚理想层。称对象 $(\mathcal{F}_n)$（属于
$\textit{Coh}(X, \mathcal{I})$）是 {\it $\mathcal{I}$-幂挠的}，或称它
{\it 被 $\mathcal{I}$ 的某个幂零化}，若存在 $c \geq 1$，使得
$\mathcal{F}_n = \mathcal{F}_c$ 对所有 $n \geq c$ 成立。
若是如此，就称 $(\mathcal{F}_n)$ {\it 被 $\mathcal{I}^c$ 零化}。
若 $X = \Spec(A)$ 仿射，则经由引理
\ref{coherent-lemma-inverse-systems-affine} 的等价，这些对象恰对应于有限
$A$-模；这些模被 $I$ 的某个幂或被 $I^c$ 零化。

\begin{lemma}
\label{coherent-lemma-existence-easy}
设 $X$ 为 Noether 概形，并设 $\mathcal{I} \subset \mathcal{O}_X$
为准凝聚理想层。设 $\mathcal{G}$ 为凝聚 $\mathcal{O}_X$-模。
设 $(\mathcal{F}_n)$ 为 $\textit{Coh}(X, \mathcal{I})$ 的对象。
\begin{enumerate}
\item 若 $\alpha : (\mathcal{F}_n) \to \mathcal{G}^\wedge$ 是一个映射，
其核与余核均被 $\mathcal{I}$ 的某个幂零化，则存在三元组
$(\mathcal{F}, a, \beta)$，且它在唯一同构意义下唯一；其中
\begin{enumerate}
\item $\mathcal{F}$ 是凝聚 $\mathcal{O}_X$-模；
\item $a : \mathcal{F} \to \mathcal{G}$ 是 $\mathcal{O}_X$-模映射，
其核与余核均被 $\mathcal{I}$ 的某个幂零化；
\item $\beta : (\mathcal{F}_n) \to \mathcal{F}^\wedge$ 是同构；并且
\item $\alpha = a^\wedge \circ \beta$。
\end{enumerate}
\item 若 $\alpha : \mathcal{G}^\wedge \to (\mathcal{F}_n)$ 是一个映射，
其核与余核均被 $\mathcal{I}$ 的某个幂零化，则存在三元组
$(\mathcal{F}, a, \beta)$，且它在唯一同构意义下唯一；其中
\begin{enumerate}
\item $\mathcal{F}$ 是凝聚 $\mathcal{O}_X$-模；
\item $a : \mathcal{G} \to \mathcal{F}$ 是 $\mathcal{O}_X$-模映射，
其核与余核均被 $\mathcal{I}$ 的某个幂零化；
\item $\beta : \mathcal{F}^\wedge \to (\mathcal{F}_n)$ 是同构；并且
\item $\alpha = \beta \circ a^\wedge$。
\end{enumerate}
\end{enumerate}
\end{lemma}

\begin{proof}
证明 (1)。唯一性表明，只须构造 $(\mathcal{F}, a, \beta)$，
并可在 $X$ 上 Zariski 局部地完成构造。因此可以假设 $X = \Spec(A)$，
且 $\mathcal{I}$ 对应于理想 $I \subset A$。在此情形下可应用引理
\ref{coherent-lemma-inverse-systems-affine}。设 $M'$ 为有限 $A^\wedge$-模，
它对应于 $(\mathcal{F}_n)$。设 $N$ 为有限 $A$-模，
它对应于 $\mathcal{G}$。于是 $\alpha$ 对应于映射
$$
\varphi : M' \longrightarrow N^\wedge
$$
其核与余核被某个 $I^t$ 零化，其中 $t$ 为某个整数。回忆
$N^\wedge = N \otimes_A A^\wedge$
（《代数》引理 \ref{algebra-lemma-completion-tensor}）。
由《代数进阶》引理
\ref{more-algebra-lemma-application-formal-glueing}，存在 $A$-模映射
$\psi : M \to N$，其核与余核是 $I$-幂挠的，并且存在同构
$M \otimes_A A^\wedge = M'$，它与 $\varphi$ 相容。
由于 $N$ 与 $M'$ 都是有限模，可知 $M$ 是有限 $A$-模；
参见《代数进阶》评注
\ref{more-algebra-remark-formal-glueing-algebras}。故
$M \otimes_A A^\wedge = M^\wedge$。我们略去验证由此得到的三元组
$(M, N \to M, M^\wedge \to M')$ 在唯一同构意义下唯一。

\medskip\noindent
(2) 的证明完全相同，从略。
\end{proof}

\begin{lemma}
\label{coherent-lemma-torsion-hom-ext}
设 $X$ 为 Noether 概形，并设 $\mathcal{I} \subset \mathcal{O}_X$
为准凝聚理想层。$\textit{Coh}(X, \mathcal{I})$ 中任何被
$\mathcal{I}$ 的某个幂零化的对象，都在
(\ref{coherent-equation-completion-functor}) 的本质像中。
此外，若 $\mathcal{F}$、$\mathcal{G}$ 属于
$\textit{Coh}(\mathcal{O}_X)$，且 $\mathcal{F}$ 或 $\mathcal{G}$
之一被 $\mathcal{I}$ 的某个幂零化，则映射
$$
\xymatrix{
\Hom_X(\mathcal{F}, \mathcal{G}) \ar[d] &
\Ext_X(\mathcal{F}, \mathcal{G}) \ar[d] \\
\Hom_{\textit{Coh}(X, \mathcal{I})}(\mathcal{F}^\wedge, \mathcal{G}^\wedge) &
\Ext_{\textit{Coh}(X, \mathcal{I})}(\mathcal{F}^\wedge, \mathcal{G}^\wedge)
}
$$
均为同构。
\end{lemma}

\begin{proof}
设 $(\mathcal{F}_n)$ 是 $\textit{Coh}(X, \mathcal{I})$ 的对象，
且被 $\mathcal{I}^c$ 零化，其中 $c \geq 1$。于是
$\mathcal{F}_n \to \mathcal{F}_c$ 对 $n \geq c$ 是同构。
故若令 $\mathcal{F} = \mathcal{F}_c$，则可见
$\mathcal{F}^\wedge \cong (\mathcal{F}_n)$。这证明了第一项断言。

\medskip\noindent
设 $\mathcal{F}$、$\mathcal{G}$ 为 $\textit{Coh}(\mathcal{O}_X)$
的对象，使得 $\mathcal{F}$ 或 $\mathcal{G}$ 之一被
$\mathcal{I}^c$ 零化，其中 $c \geq 1$。则
$\mathcal{H} = \SheafHom_{\mathcal{O}_X}(\mathcal{G}, \mathcal{F})$
是凝聚 $\mathcal{O}_X$-模，并被 $\mathcal{I}^c$ 零化。因此
$$
\Hom_X(\mathcal{G}, \mathcal{F}) =
H^0(X, \mathcal{H}) =
\lim H^0(X, \mathcal{H}/\mathcal{I}^n\mathcal{H}) =
\Mor_{\textit{Coh}(X, \mathcal{I})}
(\mathcal{G}^\wedge, \mathcal{F}^\wedge).
$$
参见引理 \ref{coherent-lemma-completion-internal-hom}。
这证明了关于同态的陈述。

\medskip\noindent
记号 $\Ext$ 指《同调》第 \ref{homology-section-extensions} 节定义的扩张。
$\Ext$ 上映射的单射性立即由 $\Hom$ 上映射的双射性推出。
为证明满射性，假设 $\mathcal{F}$ 被 $I$ 的某个幂零化。
引理 \ref{coherent-lemma-existence-easy} 的 (1) 表明：给定
$\textit{Coh}(U, I\mathcal{O}_U)$ 中的扩张
$$
0 \to \mathcal{G}^\wedge \to (\mathcal{E}_n) \to \mathcal{F}^\wedge \to 0
$$
态射 $\mathcal{G}^\wedge \to (\mathcal{E}_n)$ 同构于
$\mathcal{G} \to \mathcal{E}^\wedge$，其中
$\mathcal{G} \to \mathcal{E}$ 是
$\textit{Coh}(\mathcal{O}_U)$ 中的某个态射。另一种情形同理。
\end{proof}

\begin{lemma}
\label{coherent-lemma-finite-over-rees-algebra}
设 $X$ 为 Noether 概形，并设 $\mathcal{I} \subset \mathcal{O}_X$
为准凝聚理想层。若 $(\mathcal{F}_n)$ 是
$\textit{Coh}(X, \mathcal{I})$ 的对象，则
$\bigoplus \Ker(\mathcal{F}_{n + 1} \to \mathcal{F}_n)$ 是有限型分次
准凝聚 $\bigoplus \mathcal{I}^n/\mathcal{I}^{n + 1}$-模。
\end{lemma}

\begin{proof}
该问题在 $X$ 上是局部的，故可假设 $X$ 仿射，
即处于引理 \ref{coherent-lemma-inverse-systems-affine} 的情形。
在这种情形下，若 $(\mathcal{F}_n)$ 对应于有限 $A^\wedge$-模
$M$，则 $\bigoplus \Ker(\mathcal{F}_{n + 1} \to \mathcal{F}_n)$
对应于 $\bigoplus I^nM/I^{n + 1}M$；后者显然是
$\bigoplus I^n/I^{n + 1}$ 上的有限模。
\end{proof}

\begin{lemma}
\label{coherent-lemma-inverse-systems-pullback}
设 $f : X \to Y$ 为 Noether 概形之间的态射。
设 $\mathcal{J} \subset \mathcal{O}_Y$ 为准凝聚理想层，并令
$\mathcal{I} = f^{-1}\mathcal{J} \mathcal{O}_X$。则有右正合函子
$$
f^* : \textit{Coh}(Y, \mathcal{J}) \longrightarrow \textit{Coh}(X, \mathcal{I})
$$
它把 $(\mathcal{G}_n)$ 送到 $(f^*\mathcal{G}_n)$。若 $f$ 平坦，
则 $f^*$ 是正合函子。
\end{lemma}

\begin{proof}
由于 $f^* : \textit{Coh}(\mathcal{O}_Y) \to \textit{Coh}(\mathcal{O}_X)$
右正合，有
$$
f^*\mathcal{G}_n =
f^*(\mathcal{G}_{n + 1}/\mathcal{I}^n\mathcal{G}_{n + 1}) =
f^*\mathcal{G}_{n + 1}/f^{-1}\mathcal{I}^nf^*\mathcal{G}_{n + 1} =
f^*\mathcal{G}_{n + 1}/\mathcal{J}^nf^*\mathcal{G}_{n + 1}
$$
故一个系统的拉回仍是系统。引理
\ref{coherent-lemma-inverse-systems-abelian} 的证明中关于余核的构造表明，
$f^* : \textit{Coh}(Y, \mathcal{J}) \to \textit{Coh}(X, \mathcal{I})$
总是右正合的。若 $f$ 平坦，则
$f^* : \textit{Coh}(\mathcal{O}_Y) \to \textit{Coh}(\mathcal{O}_X)$
是正合函子。由引理 \ref{coherent-lemma-inverse-systems-abelian}
的证明中对核的构造可知，在这种情形下，
$f^* : \textit{Coh}(Y, \mathcal{J}) \to \textit{Coh}(X, \mathcal{I})$
也把核变成核。
\end{proof}

\begin{lemma}
\label{coherent-lemma-inverse-systems-pullback-equivalence}
设 $f : X' \to X$ 为 Noether 概形之间的态射。设 $Z \subset X$
为闭子概形，并以 $Z' = f^{-1}Z$ 表示其概形论逆像。
设 $\mathcal{I} \subset \mathcal{O}_X$、
$\mathcal{I}' \subset \mathcal{O}_{X'}$ 为对应的准凝聚理想层。
若 $f$ 平坦且诱导态射 $Z' \to Z$ 是同构，则拉回函子
$f^* : \textit{Coh}(X, \mathcal{I}) \to \textit{Coh}(X', \mathcal{I}')$
（引理 \ref{coherent-lemma-inverse-systems-pullback}）是等价。
\end{lemma}

\begin{proof}
若 $X$ 与 $X'$ 仿射，则结论立即由《代数进阶》引理
\ref{more-algebra-lemma-neighbourhood-equivalence} 推出。
为在一般情形下证明它，令 $Z_n \subset X$、$Z'_n \subset X'$
为第 $n$ 个无穷小邻域，分别对应于 $Z$、$Z'$。
诱导态射 $Z_n \to Z'_n$ 在底拓扑空间上是同胚。
另一方面，若 $z' \in Z'$ 映到 $z \in Z$，则环同态
$\mathcal{O}_{X, z} \to \mathcal{O}_{X', z'}$ 平坦，且诱导同构
$\mathcal{O}_{X, z}/\mathcal{I}_z \to \mathcal{O}_{X', z'}/\mathcal{I}'_{z'}$。
因此，它诱导同构
$\mathcal{O}_{X, z}/\mathcal{I}_z^n \to
\mathcal{O}_{X', z'}/(\mathcal{I}'_{z'})^n$，
这对所有 $n \geq 1$ 成立；例如见《代数进阶》引理
\ref{more-algebra-lemma-neighbourhood-isomorphism}。
故 $Z'_n \to Z_n$ 是概形的同构。
因此，$f^*$ 在凝聚 $\mathcal{O}_X$-模范畴（其中对象被
$\mathcal{I}^n$ 零化）与凝聚 $\mathcal{O}_{X'}$-模范畴（其中对象被
$(\mathcal{I}')^n$ 零化）之间诱导等价；参见引理
\ref{coherent-lemma-i-star-equivalence}。这显然蕴含本引理。
\end{proof}

\begin{lemma}
\label{coherent-lemma-inverse-systems-ideals-equivalence}
设 $X$ 为 Noether 概形。设
$\mathcal{I}, \mathcal{J} \subset \mathcal{O}_X$ 为准凝聚理想层。
若 $V(\mathcal{I}) = V(\mathcal{J})$ 是 $X$ 的同一个闭子集，
则 $\textit{Coh}(X, \mathcal{I})$ 与
$\textit{Coh}(X, \mathcal{J})$ 等价。
\end{lemma}

\begin{proof}
首先假设 $X = \Spec(A)$ 仿射。设 $I, J \subset A$ 为与
$\mathcal{I}, \mathcal{J}$ 对应的理想。则
$V(I) = V(J)$ 蕴含 $I^c \subset J$ 且 $J^d \subset I$，
其中 $c, d \geq 1$ 为某些整数；这是 Zariski 拓扑的初等性质
（参见《代数》第 \ref{algebra-section-spectrum-ring} 节及引理
\ref{algebra-lemma-Noetherian-power}）。故 $I$-进与 $J$-进完备化
（对 $A$）相同；参见《代数》引理
\ref{algebra-lemma-change-ideal-completion}。因此，在此情形下，
所求等价由引理 \ref{coherent-lemma-inverse-systems-affine} 推出。

\medskip\noindent
一般而言，利用上面的论述及 $X$ 拟紧这一事实，选取
$c, d \geq 1$，使得 $\mathcal{I}^c \subset \mathcal{J}$ 且
$\mathcal{J}^d \subset \mathcal{I}$。给定对象 $(\mathcal{F}_n)$
（属于 $\textit{Coh}(X, \mathcal{I})$），
断言逆向系统
$$
(\mathcal{F}_{cn}/\mathcal{J}^n\mathcal{F}_{cn})
$$
属于 $\textit{Coh}(X, \mathcal{J})$。这可在一个仿射覆盖的各成员上检验；
细节从略。以同样方式，可以构造 $\textit{Coh}(X, \mathcal{I})$
的对象，其起点是 $\textit{Coh}(X, \mathcal{J})$ 的一个对象。
我们略去验证这些构造定义互为拟逆的函子。
\end{proof}

\section{Grothendieck 存在定理（一）}
\label{coherent-section-existence}

\noindent
本节讨论射影情形下的 Grothendieck 存在定理。我们将使用第
\ref{coherent-section-coherent-formal} 节发展的凝聚形式模概念。
建议熟悉形式概形的读者阅读 \cite{EGA} 中该定理的陈述与证明。

\begin{lemma}
\label{coherent-lemma-fully-faithful}
设 $A$ 为关于理想 $I$ 完备的 Noether 环。
设 $f : X \to \Spec(A)$ 为固有态射，并令
$\mathcal{I} = I\mathcal{O}_X$。则函子
(\ref{coherent-equation-completion-functor}) 是全忠实的。
\end{lemma}

\begin{proof}
设 $\mathcal{F}$、$\mathcal{G}$ 为凝聚 $\mathcal{O}_X$-模。
则 $\mathcal{H} = \SheafHom_{\mathcal{O}_X}(\mathcal{G}, \mathcal{F})$
是凝聚 $\mathcal{O}_X$-模；参见《模》引理
\ref{modules-lemma-internal-hom-locally-kernel-direct-sum}。
由引理 \ref{coherent-lemma-completion-internal-hom}，映射
$$
\lim_n H^0(X, \mathcal{H}/\mathcal{I}^n\mathcal{H})
\to
\Mor_{\textit{Coh}(X, \mathcal{I})}
(\mathcal{G}^\wedge, \mathcal{F}^\wedge)
$$
是双射。因此，函子 (\ref{coherent-equation-completion-functor}) 的全忠实性
由应用于凝聚层 $\mathcal{H}$ 的形式函数定理
（引理 \ref{coherent-lemma-spell-out-theorem-formal-functions}）推出。
\end{proof}

\begin{lemma}
\label{coherent-lemma-vanishing-projective}
设 $A$ 为 Noether 环，$I \subset A$ 为理想。
设 $f : X \to \Spec(A)$ 为固有态射，并设 $\mathcal{L}$
为 $f$-丰沛可逆层。令 $\mathcal{I} = I\mathcal{O}_X$。
设 $(\mathcal{F}_n)$ 为 $\textit{Coh}(X, \mathcal{I})$ 的对象。
则存在整数 $d_0$，使得
$$
H^1(X, \Ker(\mathcal{F}_{n + 1} \to \mathcal{F}_n)
\otimes \mathcal{L}^{\otimes d} )
= 0
$$
对所有 $n \geq 0$ 及所有 $d \geq d_0$ 成立。
\end{lemma}

\begin{proof}
令 $B = \bigoplus I^n/I^{n + 1}$，并令
$\mathcal{B} = \bigoplus \mathcal{I}^n/\mathcal{I}^{n + 1} = f^*\widetilde{B}$。
由引理 \ref{coherent-lemma-finite-over-rees-algebra}，分次准凝聚
$\mathcal{B}$-模
$\mathcal{G} = \bigoplus \Ker(\mathcal{F}_{n + 1} \to \mathcal{F}_n)$
是有限型的。故本引理由引理 \ref{coherent-lemma-graded-finiteness} 的 (2) 推出。
\end{proof}

\begin{lemma}
\label{coherent-lemma-existence-projective}
设 $A$ 为关于理想 $I$ 完备的 Noether 环。
设 $f : X \to \Spec(A)$ 为射影态射，并令
$\mathcal{I} = I\mathcal{O}_X$。则函子
(\ref{coherent-equation-completion-functor}) 是等价。
\end{lemma}

\begin{proof}
(\ref{coherent-equation-completion-functor}) 是全忠实的；这已在引理
\ref{coherent-lemma-fully-faithful} 中看到。因此只须证明该函子
本质满射。

\medskip\noindent
先证明每个对象 $(\mathcal{F}_n)$（属于
$\textit{Coh}(X, \mathcal{I})$）都是函子
(\ref{coherent-equation-completion-functor}) 之像中某个对象的商。
设 $\mathcal{L}$ 为 $f$-丰沛可逆层（在 $X$ 上）。
按引理 \ref{coherent-lemma-vanishing-projective} 选取
$d_0$。选取 $d \geq d_0$，使得
$\mathcal{F}_1 \otimes \mathcal{L}^{\otimes d}$
由某些整体截面 $s_{1, 1}, \ldots, s_{t, 1}$ 生成。
由于系统的转移映射
$$
H^0(X, \mathcal{F}_{n + 1} \otimes \mathcal{L}^{\otimes d})
\longrightarrow
H^0(X, \mathcal{F}_n \otimes \mathcal{L}^{\otimes d})
$$
因 $H^1$ 的消失而成为满射，可以把
$s_{1, 1}, \ldots, s_{t, 1}$ 提升为相容整体截面系
$s_{1, n}, \ldots, s_{t, n}$；这些是
$\mathcal{F}_n \otimes \mathcal{L}^{\otimes d}$ 的截面。
它们确定相容映射系
$$
(s_{1, n}, \ldots, s_{t, n}) :
(\mathcal{L}^{\otimes -d})^{\oplus t} \longrightarrow \mathcal{F}_n
$$
由引理 \ref{coherent-lemma-inverse-systems-surjective}，得到所需的满射
$$
\left((\mathcal{L}^{\otimes -d})^{\oplus t}\right)^\wedge
\longrightarrow
(\mathcal{F}_n)
$$

\medskip\noindent
上一段的结论以及 $\textit{Coh}(X, \mathcal{I})$ 是阿贝尔范畴这一事实
（引理 \ref{coherent-lemma-inverse-systems-abelian}）蕴含：
$\textit{Coh}(X, \mathcal{I})$ 的每个对象，都是来自
$\textit{Coh}(\mathcal{O}_X)$ 的两个对象之间某个映射的余核。
由于 (\ref{coherent-equation-completion-functor}) 全忠实且正合
（引理 \ref{coherent-lemma-fully-faithful} 与 \ref{coherent-lemma-exact}），结论成立。
\end{proof}

\section{Grothendieck 存在定理（二）}
\label{coherent-section-existence-proper}

\noindent
本节讨论固有情形下的 Grothendieck 存在定理。在给出其陈述与证明前，
还需要进一步发展第 \ref{coherent-section-coherent-formal} 节引入的凝聚形式模范畴
$\textit{Coh}(X, \mathcal{I})$ 的一些理论。

\begin{remark}
\label{coherent-remark-inverse-systems-kernel-cokernel-annihilated-by}
设 $X$ 为 Noether 概形，并设
$\mathcal{I}, \mathcal{K} \subset \mathcal{O}_X$ 为准凝聚理想层。
设 $\alpha : (\mathcal{F}_n) \to (\mathcal{G}_n)$ 为
$\textit{Coh}(X, \mathcal{I})$ 的态射。
给定仿射开集 $\Spec(A) = U \subset X$，其中
$\mathcal{I}|_U, \mathcal{K}|_U$ 对应于理想 $I, K \subset A$，
以 $\alpha_U : M \to N$ 表示有限 $A^\wedge$-模之间的映射，
它经由引理 \ref{coherent-lemma-inverse-systems-affine} 对应于 $\alpha|_U$。
我们断言下列条件等价：
\begin{enumerate}
\item 存在整数 $t \geq 1$，使得
$\Ker(\alpha_n)$ 与 $\Coker(\alpha_n)$ 被 $\mathcal{K}^t$ 零化，
这对所有 $n \geq 1$ 成立；
\item 对任意上述仿射开集 $\Spec(A) = U \subset X$，模
$\Ker(\alpha_U)$ 与 $\Coker(\alpha_U)$ 被某个 $K^t$ 零化，
其中 $t \geq 1$ 为某个整数；并且
\item 存在有限仿射开覆盖 $X = \bigcup U_i$，使得 (2) 的结论
对 $\alpha_{U_i}$ 成立。
\end{enumerate}
若这些等价条件成立，就称 $\alpha$ 为
{\it 核与余核被 $\mathcal{K}$ 的某个幂零化的映射}。
为看出其等价性，使用如下交换代数事实：假设给定 $A$-模的正合列
$$
0 \to T \to M \to N \to Q \to 0
$$
其中 $T$ 与 $Q$ 被 $K^t$ 零化，这里 $K \subset A$ 为某个理想。
则对每个 $f, g \in K^t$，存在典范映射
$"fg": N \to M$，使得 $M \to N \to M$ 等于乘以 $fg$。
事实上，对 $y \in N$，可以选取 $x \in M$，其像为
$fy$（在 $N$ 中），再令 $"fg"(y) = gx$。由此显然可见，
$\Ker(M/JM \to N/JN)$ 与 $\Coker(M/JM \to N/JN)$
均被 $K^{2t}$ 零化，这对任意理想 $J \subset A$ 成立。

\medskip\noindent
把该交换代数事实应用于 $\alpha_{U_i}$ 与 $J = I^n$，
可见 (3) 蕴含 (1)。反之，假设 (1) 成立且
$M \to N$ 等于 $\alpha_U$。则存在 $t \geq 1$，使得
$\Ker(M/I^nM \to N/I^nN)$ 与 $\Coker(M/I^nM \to N/I^nN)$
被 $K^t$ 零化，这对所有 $n$ 成立。得到映射
$"fg" : N/I^nN \to M/I^nM$，它们在极限中诱导映射
$N \to M$；这里 $N$ 与 $M$ 都是 $I$-进完备的。因为与
$N \to M \to N$ 的复合是乘以 $fg$，可知 $fg$ 零化 $T$ 与 $Q$。
换言之，$T$ 与 $Q$ 如所需地被 $K^{2t}$ 零化。
\end{remark}

\begin{lemma}
\label{coherent-lemma-existence-tricky}
设 $X$ 为 Noether 概形。设
$\mathcal{I}, \mathcal{K} \subset \mathcal{O}_X$ 为准凝聚理想层。
设 $X_e \subset X$ 为由 $\mathcal{K}^e$ 截出的闭子概形。
令 $\mathcal{I}_e = \mathcal{I}\mathcal{O}_{X_e}$。
设 $(\mathcal{F}_n)$ 为 $\textit{Coh}(X, \mathcal{I})$ 的对象。
假设：
\begin{enumerate}
\item 函子
$\textit{Coh}(\mathcal{O}_{X_e}) \to \textit{Coh}(X_e, \mathcal{I}_e)$
对所有 $e \geq 1$ 都是等价；并且
\item 存在凝聚层 $\mathcal{H}$（在 $X$ 上）及映射
$\alpha : (\mathcal{F}_n) \to \mathcal{H}^\wedge$，其核与余核
被 $\mathcal{K}$ 的某个幂零化。
\end{enumerate}
则 $(\mathcal{F}_n)$ 在 (\ref{coherent-equation-completion-functor}) 的本质像中。
\end{lemma}

\begin{proof}
在本证明中，将不再特别说明地使用如下事实：对闭浸入
$i : Z \to X$，函子 $i_*$ 在 $Z$ 上凝聚模范畴与 $X$ 上被
$Z$ 的理想层零化的凝聚模范畴之间给出等价；参见引理
\ref{coherent-lemma-i-star-equivalence}。特别地，可以把
$\textit{Coh}(\mathcal{O}_{X_e})$ 等同于凝聚 $\mathcal{O}_X$-模范畴
（其中对象被 $\mathcal{K}^e$ 零化），并把
$\textit{Coh}(X_e, \mathcal{I}_e)$
等同于 $\textit{Coh}(X, \mathcal{I})$ 中由被 $\mathcal{K}^e$
零化的对象组成的满子范畴。而且，(1) 告诉我们，这两个范畴在完备化函子
(\ref{coherent-equation-completion-functor}) 下等价。

\medskip\noindent
应用该等价，得到凝聚 $\mathcal{O}_X$-模 $\mathcal{G}_e$，
它被 $\mathcal{K}^e$ 零化，并对应于系统
$(\mathcal{F}_n/\mathcal{K}^e\mathcal{F}_n)$；该系统属于
$\textit{Coh}(X, \mathcal{I})$。映射
$\mathcal{F}_n/\mathcal{K}^{e + 1}\mathcal{F}_n \to
\mathcal{F}_n/\mathcal{K}^e\mathcal{F}_n$ 对应于典范映射
$\mathcal{G}_{e + 1} \to \mathcal{G}_e$，它们诱导同构
$\mathcal{G}_{e + 1}/\mathcal{K}^e\mathcal{G}_{e + 1} \to \mathcal{G}_e$。
故 $(\mathcal{G}_e)$ 是 $\textit{Coh}(X, \mathcal{K})$ 的对象。
映射 $\alpha$ 诱导映射系
$$
\mathcal{F}_n/\mathcal{K}^e\mathcal{F}_n
\longrightarrow
\mathcal{H}/(\mathcal{I}^n + \mathcal{K}^e)\mathcal{H}
$$
从而得到映射
$\mathcal{G}_e \to \mathcal{H}/\mathcal{K}^e\mathcal{H}$
（再次使用范畴等价）。由假设 (2)，存在整数 $t \geq 1$，
使得 $\mathcal{K}^t$ 零化所有映射
$\mathcal{F}_n \to \mathcal{H}/\mathcal{I}^n\mathcal{H}$ 的核与余核。
于是，$\mathcal{K}^{2t}$ 零化映射
$\mathcal{F}_n/\mathcal{K}^e\mathcal{F}_n \to
\mathcal{H}/(\mathcal{I}^n + \mathcal{K}^e)\mathcal{H}$ 的核与余核；
参见评注 \ref{coherent-remark-inverse-systems-kernel-cokernel-annihilated-by}。
继而可知 $\mathcal{K}^{4t}$ 零化下列映射的核与余核：
$$
\mathcal{G}_e
\longrightarrow
\mathcal{H}/\mathcal{K}^e\mathcal{H},
$$
参见评注 \ref{coherent-remark-inverse-systems-kernel-cokernel-annihilated-by}。
应用引理 \ref{coherent-lemma-existence-easy}，得到凝聚 $\mathcal{O}_X$-模
$\mathcal{F}$、映射 $a : \mathcal{F} \to \mathcal{H}$，以及同构
$\beta : (\mathcal{G}_e) \to (\mathcal{F}/\mathcal{K}^e\mathcal{F})$
（在 $\textit{Coh}(X, \mathcal{K})$ 中）。
反向追溯，对给定的 $n$，三元组
$(\mathcal{F}/\mathcal{I}^n\mathcal{F}, a \bmod \mathcal{I}^n, \beta
\bmod \mathcal{I}^n)$ 是该引理中对应于态射
$\alpha_n \bmod \mathcal{K}^e :
(\mathcal{F}_n/\mathcal{K}^e\mathcal{F}_n) \to
(\mathcal{H}/(\mathcal{I}^n + \mathcal{K}^e)\mathcal{H})$
（作为 $\textit{Coh}(X, \mathcal{K})$ 的态射）的三元组。
因此，引理 \ref{coherent-lemma-existence-easy} 的唯一性给出典范同构
$\mathcal{F}/\mathcal{I}^n\mathcal{F} \to \mathcal{F}_n$，
并与涉及的所有态射相容。这就完成了引理的证明。
\end{proof}

\begin{lemma}
\label{coherent-lemma-inverse-systems-push-pull}
设 $Y$ 为 Noether 概形。设
$\mathcal{J}, \mathcal{K} \subset \mathcal{O}_Y$ 为准凝聚理想层。
设 $f : X \to Y$ 为固有态射，并且它在
$V = Y \setminus V(\mathcal{K})$ 上为同构。
令 $\mathcal{I} = f^{-1}\mathcal{J} \mathcal{O}_X$。
设 $(\mathcal{G}_n)$ 为 $\textit{Coh}(Y, \mathcal{J})$ 的对象，
$\mathcal{F}$ 为凝聚 $\mathcal{O}_X$-模，并设
$\beta : (f^*\mathcal{G}_n)  \to \mathcal{F}^\wedge$ 为
$\textit{Coh}(X, \mathcal{I})$ 中的同构。则存在映射
$$
\alpha :
(\mathcal{G}_n)
\longrightarrow
(f_*\mathcal{F})^\wedge
$$
（在 $\textit{Coh}(Y, \mathcal{J})$ 中），其核与余核被
$\mathcal{K}$ 的某个幂零化。
\end{lemma}

\begin{proof}
由于 $f$ 是固有态射，可知 $f_*\mathcal{F}$ 是凝聚
$\mathcal{O}_Y$-模（命题
\ref{coherent-proposition-proper-pushforward-coherent}）。故引理的陈述有意义。
考虑复合映射
$$
\gamma_n : \mathcal{G}_n \to
f_*f^*\mathcal{G}_n \to
f_*(\mathcal{F}/\mathcal{I}^n\mathcal{F}).
$$
这里第一个映射是伴随映射，第二个是 $f_*\beta_n$。
我们断言：存在唯一的引理中所述 $\alpha$，使得复合映射
$$
\mathcal{G}_n \xrightarrow{\alpha_n}
f_*\mathcal{F}/\mathcal{J}^nf_*\mathcal{F} \to
f_*(\mathcal{F}/\mathcal{I}^n\mathcal{F})
$$
都等于 $\gamma_n$，这对所有 $n$ 成立。由于唯一性，可以假设
$Y = \Spec(B)$ 仿射。令 $J \subset B$ 为与理想层
$\mathcal{J}$ 对应的理想。令
$$
M_n = H^0(X, \mathcal{F}/\mathcal{I}^n\mathcal{F})
\quad\text{且}\quad
M = H^0(X, \mathcal{F})
$$
由引理 \ref{coherent-lemma-ML-cohomology-powers-ideal} 与定理
\ref{coherent-theorem-formal-functions}，诸模 $M_n$ 的逆极限等于完备化
$M^\wedge = \lim M/J^nM$。令
$N_n = H^0(Y, \mathcal{G}_n)$ 且 $N = \lim N_n$。
经由引理 \ref{coherent-lemma-inverse-systems-affine} 的范畴等价，有限
$B^\wedge$-模 $N$ 与 $M^\wedge$ 分别对应于
$(\mathcal{G}_n)$ 与 $f_*\mathcal{F}^\wedge$。
由此，$\alpha$ 必须是 $\textit{Coh}(Y, \mathcal{J})$ 的态射，
它对应于同态
$$
\lim \gamma_n : N = \lim_n N_n \longrightarrow \lim M_n = M^\wedge
$$
（在有限 $B^\wedge$-模之间）。

\medskip\noindent
还须证明 $\alpha$ 的核与余核被 $\mathcal{K}$ 的某个幂零化。
令 $Y' = \Spec(B^\wedge)$，并令 $X' = Y' \times_Y X$。
令 $\mathcal{K}'$、$\mathcal{J}'$、$\mathcal{G}'_n$ 以及
$\mathcal{I}'$、$\mathcal{F}'$ 分别为 $\mathcal{K}$、$\mathcal{J}$、
$\mathcal{G}_n$ 以及 $\mathcal{I}$、$\mathcal{F}$ 到
$Y'$ 与 $X'$ 的拉回。投影态射 $f' : X' \to Y'$ 是 $f$
沿 $Y' \to Y$ 的基变换。注意 $Y' \to Y$ 是平坦概形态射，
因为由《代数》引理 \ref{algebra-lemma-completion-flat}，
$B \to B^\wedge$ 平坦。因此，由引理
\ref{coherent-lemma-flat-base-change-cohomology}，$f'_*\mathcal{F}'$ 以及
$f'_*(f')^*\mathcal{G}_n'$ 分别是 $f_*\mathcal{F}$ 以及
$f_*f^*\mathcal{G}_n$ 到 $Y'$ 的拉回。
构造的唯一性表明，$\alpha$ 到 $Y'$ 的拉回就是对应映射
$\alpha'$，它为 $Y'$ 上情形构造。而且，为检验 $\alpha$ 的核与余核被
$\mathcal{K}^t$ 零化，只须检验 $\alpha'$ 的核与余核被
$(\mathcal{K}')^t$ 零化。确实，为此须检验映射
$\alpha_n$ 与 $\alpha'_n$ 的核与余核（见评注
\ref{coherent-remark-inverse-systems-kernel-cokernel-annihilated-by}），
而环同态 $B \to B^\wedge$ 在被 $J^n$ 与被 $(J')^n$ 零化的模
之间诱导范畴等价；参见《代数进阶》引理
\ref{more-algebra-lemma-neighbourhood-equivalence}。
故可以假设 $B$ 关于 $J$ 完备。

\medskip\noindent
假设 $Y = \Spec(B)$ 仿射，$\mathcal{J}$ 对应于理想
$J \subset B$，且 $B$ 关于 $J$ 完备。
在这种情形下，$(\mathcal{G}_n)$ 位于函子
$\textit{Coh}(\mathcal{O}_Y) \to \textit{Coh}(Y, \mathcal{J})$
的本质像中。设 $\mathcal{G}$ 是凝聚 $\mathcal{O}_Y$-模，满足
$(\mathcal{G}_n) = \mathcal{G}^\wedge$。注意
$f^*(\mathcal{G}^\wedge) = (f^*\mathcal{G})^\wedge$。
因此，引理 \ref{coherent-lemma-fully-faithful} 告诉我们，$\beta$ 来自同构
$b : f^*\mathcal{G} \to \mathcal{F}$，且 $\alpha$ 是把完备化函子
应用于下列映射所得的态射：
$$
\mathcal{G} \to f_*f^*\mathcal{G} \cong f_*\mathcal{F}
$$
故要验证伴随映射 $c : \mathcal{G} \to f_*f^*\mathcal{G}$ 的核与余核
被 $\mathcal{K}$ 的某个幂零化。然而，由于限制
$f|_{f^{-1}(V)} : f^{-1}(V) \to V$ 是同构，可知 $c|_V$ 是同构。
因此，凝聚层 $\Ker(c)$ 与 $\Coker(c)$ 支撑在
$V(\mathcal{K})$ 上，故被 $\mathcal{K}$ 的某个幂零化
（引理 \ref{coherent-lemma-power-ideal-kills-sheaf}），如所需。
\end{proof}

\noindent
下述命题是实践中最常用的 Grothendieck 存在定理形式。

\begin{proposition}
\label{coherent-proposition-existence-proper}
设 $A$ 为关于理想 $I$ 完备的 Noether 环。
设 $f : X \to \Spec(A)$ 为固有概形态射。
令 $\mathcal{I} = I\mathcal{O}_X$。则函子
(\ref{coherent-equation-completion-functor}) 是等价。
\end{proposition}

\begin{proof}
已经看到，函子 (\ref{coherent-equation-completion-functor}) 是全忠实的；
见引理 \ref{coherent-lemma-fully-faithful}。因此只须证明该函子本质满射。

\medskip\noindent
考虑集合 $\Xi$，其元素是准凝聚理想层
$\mathcal{K} \subset \mathcal{O}_X$，且满足：每个对象
$(\mathcal{F}_n)$（被 $\mathcal{K}$ 零化）均位于本质像中。
要证明 $(0)$ 属于 $\Xi$。
若不然，由于 $X$ 是 Noether 概形，存在极大准凝聚理想层
$\mathcal{K}$，它不属于 $\Xi$；参见引理 \ref{coherent-lemma-acc-coherent}。
将 $X$ 替换为 $X$ 中与 $\mathcal{K}$ 对应的闭子概形后，
可以假设
每个非零 $\mathcal{K}$ 均属于 $\Xi$。（这里使用了被
$\mathcal{K}$ 零化的凝聚模与 $\mathcal{K}$ 所对应闭子概形上的
凝聚模之间的对应；参见引理 \ref{coherent-lemma-i-star-equivalence}。）
设 $(\mathcal{F}_n)$ 为 $\textit{Coh}(X, \mathcal{I})$ 的对象。
将证明该对象位于函子 (\ref{coherent-equation-completion-functor}) 的本质像中，
从而完成命题的证明。

\medskip\noindent
应用 Chow 引理（引理 \ref{coherent-lemma-chow-Noetherian}），得到固有满射
$f : X' \to X$，它在稠密开集 $U \subset X$ 上为同构，
且 $X'$ 在 $A$ 上射影。设 $\mathcal{K}$ 为截出约化补集
$X \setminus U$ 的准凝聚理想层。由射影情形的 Grothendieck 存在定理
（引理 \ref{coherent-lemma-existence-projective}），存在凝聚模
$\mathcal{F}'$（在 $X'$ 上），使得
$(\mathcal{F}')^\wedge \cong (f^*\mathcal{F}_n)$。
由命题 \ref{coherent-proposition-proper-pushforward-coherent}，
$\mathcal{O}_X$-模 $\mathcal{H} = f_*\mathcal{F}'$ 凝聚；
由引理 \ref{coherent-lemma-inverse-systems-push-pull}，存在态射
$(\mathcal{F}_n) \to \mathcal{H}^\wedge$
（在 $\textit{Coh}(X, \mathcal{I})$ 中），其核与余核被
$\mathcal{K}$ 的某个幂零化。诸幂 $\mathcal{K}^e$ 都属于
$\Xi$，因而 (\ref{coherent-equation-completion-functor}) 对闭子概形
$X_e = V(\mathcal{K}^e)$ 是等价。最后应用引理
\ref{coherent-lemma-existence-tricky} 即得结论。
\end{proof}

\section{在基底上为固有}
\label{coherent-section-proper-over-base}

\noindent
本节很短，只是指出在基底上固有的闭子集，以及支撑在基底上固有的
有限型准凝聚模所具有的一些有用性质。

\begin{lemma}
\label{coherent-lemma-closed-proper-over-base}
设 $f : X \to S$ 为局部有限型概形态射，并设 $Z \subset X$ 为闭子集。
下列条件等价：
\begin{enumerate}
\item 态射 $Z \to S$ 在赋予 $Z$ 诱导的约化闭子概形结构时固有
（《概形》定义 \ref{schemes-definition-reduced-induced-scheme}）；
\item 对 $Z$ 的某个闭子概形结构，态射 $Z \to S$ 固有；
\item 对 $Z$ 的任意闭子概形结构，态射 $Z \to S$ 固有。
\end{enumerate}
\end{lemma}

\begin{proof}
(3) $\Rightarrow$ (1) 与 (1) $\Rightarrow$ (2) 都是立即的，
故只须证明 (2) 蕴含 (3)。我们建议读者自行找出这个事实的证明。
设 $Z'$ 与 $Z''$ 是 $Z$ 上的闭子概形结构，且 $Z' \to S$ 固有。
须证明 $Z'' \to S$ 固有。令 $Z''' = Z' \cup Z''$ 为概形论并；
参见《态射》定义
\ref{morphisms-definition-scheme-theoretic-intersection-union}。
则 $Z'''$ 是 $Z$ 上的另一个闭子概形结构。例如，这可由
《态射》引理 \ref{morphisms-lemma-scheme-theoretic-union}
对概形论并的描述推出。由于 $Z'' \to Z'''$ 是闭浸入，
只须证明 $Z''' \to S$ 固有（见《态射》引理
\ref{morphisms-lemma-closed-immersion-proper} 与
\ref{morphisms-lemma-composition-proper}）。
态射 $Z' \to Z'''$ 是双射闭浸入，特别地，它满且万有闭。
于是，$Z' \to S$ 分离蕴含 $Z''' \to S$ 分离；见《态射》引理
\ref{morphisms-lemma-image-universally-closed-separated}。
此外，$Z''' \to S$ 局部有限型，因为 $X \to S$ 局部有限型
（《态射》引理 \ref{morphisms-lemma-immersion-locally-finite-type}
与 \ref{morphisms-lemma-composition-finite-type}）。
由于 $Z' \to S$ 拟紧，且 $Z' \to Z'''$ 是同胚，可知
$Z''' \to S$ 拟紧。最后，由于 $Z' \to S$ 万有闭，
《态射》引理 \ref{morphisms-lemma-image-proper-is-proper}
表明 $Z''' \to S$ 也万有闭。这就完成了证明。
\end{proof}

\begin{definition}
\label{coherent-definition-proper-over-base}
设 $f : X \to S$ 为局部有限型概形态射，并设 $Z \subset X$ 为闭子集。
若引理 \ref{coherent-lemma-closed-proper-over-base} 的等价条件成立，
则称{\it $Z$ 在 $S$ 上固有}。
\end{definition}

\noindent
如果态射 $f : X \to S$ 不是局部有限型的，则上述定义所用的引理不成立。
因此，当 $f$ 不是局部有限型时，我们建议读者不要使用这一术语。

\begin{lemma}
\label{coherent-lemma-closed-closed-proper-over-base}
设 $f : X \to S$ 为局部有限型概形态射，并设
$Y \subset Z \subset X$ 为闭子集。若 $Z$ 在 $S$ 上固有，
则 $Y$ 也具有同一性质。
\end{lemma}

\begin{proof}
略。
\end{proof}

\begin{lemma}
\label{coherent-lemma-base-change-closed-proper-over-base}
考虑概形的笛卡尔图
$$
\xymatrix{
X' \ar[d]_{f'} \ar[r]_{g'} & X \ar[d]^f \\
S' \ar[r]^g & S
}
$$
其中 $f$ 局部有限型。若 $Z$ 是 $X$ 的闭子集并在 $S$ 上固有，
则 $(g')^{-1}(Z)$ 是 $X'$ 的闭子集并在 $S'$ 上固有。
\end{lemma}

\begin{proof}
由《态射》引理 \ref{morphisms-lemma-base-change-finite-type}，
$f'$ 局部有限型，故陈述有意义。赋予 $Z$ 诱导的约化闭子概形结构。
以 $Z' = (g')^{-1}(Z)$ 表示概形论逆像
（《概形》定义 \ref{schemes-definition-inverse-image-closed-subscheme}）。
则
$Z' = X' \times_X Z = (S' \times_S X) \times_X Z = S' \times_S Z$
在 $S'$ 上固有，因为它是 $Z$（在 $S$ 上）的基变换
（《态射》引理 \ref{morphisms-lemma-base-change-proper}）。
\end{proof}

\begin{lemma}
\label{coherent-lemma-functoriality-closed-proper-over-base}
设 $S$ 为概形。设 $f : X \to Y$ 为两个在 $S$ 上局部有限型的
概形之间的态射。
\begin{enumerate}
\item 若 $Y$ 在 $S$ 上分离，且 $Z \subset X$ 是在 $S$ 上固有的
闭子集，则 $f(Z)$ 是 $Y$ 的闭子集并在 $S$ 上固有。
\item 若 $f$ 万有闭，且 $Z \subset X$ 是在 $S$ 上固有的闭子集，
则 $f(Z)$ 是 $Y$ 的闭子集并在 $S$ 上固有。
\item 若 $f$ 固有，且 $Z \subset Y$ 是在 $S$ 上固有的闭子集，
则 $f^{-1}(Z)$ 是 $X$ 的闭子集并在 $S$ 上固有。
\end{enumerate}
\end{lemma}

\begin{proof}
证明 (1)。假设 $Y$ 在 $S$ 上分离，且 $Z \subset X$
是一个在 $S$ 上固有的闭子集。赋予 $Z$ 诱导的约化闭子概形结构，
并把《态射》引理
\ref{morphisms-lemma-scheme-theoretic-image-is-proper}
应用于态射 $Z \to Y$（在 $S$ 上），即得结论。

\medskip\noindent
证明 (2)。假设 $f$ 万有闭，且 $Z \subset X$ 是在 $S$ 上固有的
闭子集。赋予 $Z$ 与 $Z' = f(Z)$ 各自诱导的约化闭子概形结构，
从而得到诱导态射 $Z \to Z'$。
以 $Z'' = f^{-1}(Z')$ 表示概形论逆像
（《概形》定义 \ref{schemes-definition-inverse-image-closed-subscheme}）。
则 $Z'' \to Z'$ 作为 $f$ 的基变换是万有闭的
（《态射》引理 \ref{morphisms-lemma-base-change-proper}）。
于是，$Z \to Z'$ 作为闭浸入 $Z \to Z''$ 与 $Z'' \to Z'$ 的复合
也是万有闭的（《态射》引理
\ref{morphisms-lemma-closed-immersion-proper} 与
\ref{morphisms-lemma-composition-proper}）。
由《态射》引理 \ref{morphisms-lemma-image-universally-closed-separated}，
$Z' \to S$ 分离。由于 $Z \to S$ 拟紧且 $Z \to Z'$ 满，
可知 $Z' \to S$ 拟紧。由于 $Z' \to S$ 是 $Z' \to Y$ 与
$Y \to S$ 的复合，可知 $Z' \to S$ 局部有限型
（《态射》引理 \ref{morphisms-lemma-immersion-locally-finite-type}
与 \ref{morphisms-lemma-composition-finite-type}）。
最后，由于 $Z \to S$ 万有闭，《态射》引理
\ref{morphisms-lemma-image-proper-is-proper}
表明 $Z' \to S$ 也万有闭。这就完成了证明。

\medskip\noindent
证明 (3)。假设 $f$ 固有，且 $Z \subset Y$ 是在 $S$ 上固有的闭子集。
赋予 $Z$ 诱导的约化闭子概形结构。
以 $Z' = f^{-1}(Z)$ 表示概形论逆像
（《概形》定义 \ref{schemes-definition-inverse-image-closed-subscheme}）。
则 $Z' \to Z$ 作为 $f$ 的基变换是固有的
（《态射》引理 \ref{morphisms-lemma-base-change-proper}）。
因此，$Z' \to S$ 作为 $Z' \to Z$ 与 $Z \to S$ 的复合是固有的
（《态射》引理 \ref{morphisms-lemma-composition-proper}）。
这就完成了证明。
\end{proof}

\begin{lemma}
\label{coherent-lemma-union-closed-proper-over-base}
设 $f : X \to S$ 为局部有限型概形态射。设
$Z_i \subset X$（$i = 1, \ldots, n$）为闭子集。
若 $Z_i$（$i = 1, \ldots, n$）都在 $S$ 上固有，则
$Z_1 \cup \ldots \cup Z_n$ 也具有同一性质。
\end{lemma}

\begin{proof}
赋予 $Z_i$ 各自诱导的约化闭子概形结构。态射
$$
Z_1 \amalg \ldots \amalg Z_n \longrightarrow X
$$
由《态射》引理 \ref{morphisms-lemma-closed-immersion-finite}
与 \ref{morphisms-lemma-finite-union-finite} 是有限的。
由于有限态射万有闭（《态射》引理
\ref{morphisms-lemma-finite-proper}），并且
$Z_1 \amalg \ldots \amalg Z_n$ 在 $S$ 上固有，
由引理 \ref{coherent-lemma-functoriality-closed-proper-over-base} 的 (2)
可知其像 $Z_1 \cup \ldots \cup Z_n$ 在 $S$ 上固有。
\end{proof}

\noindent
设 $f : X \to S$ 为局部有限型概形态射，并设 $\mathcal{F}$ 为有限型
准凝聚 $\mathcal{O}_X$-模。则支撑
$\text{Supp}(\mathcal{F})$（即 $\mathcal{F}$ 的支撑）是 $X$ 的闭子集；
见《态射》引理
\ref{morphisms-lemma-support-finite-type}。因此，说
“$\mathcal{F}$ 的支撑在 $S$ 上固有”是有意义的。

\begin{lemma}
\label{coherent-lemma-module-support-proper-over-base}
设 $f : X \to S$ 为局部有限型概形态射，并设 $\mathcal{F}$ 为有限型
准凝聚 $\mathcal{O}_X$-模。下列条件等价：
\begin{enumerate}
\item $\mathcal{F}$ 的支撑在 $S$ 上固有；
\item $\mathcal{F}$ 的概形论支撑
（《态射》定义 \ref{morphisms-definition-scheme-theoretic-support}）
在 $S$ 上固有；
\item 存在闭子概形 $Z \subset X$ 及有限型准凝聚
$\mathcal{O}_Z$-模 $\mathcal{G}$，使得 (a) $Z \to S$ 固有，
且 (b) $(Z \to X)_*\mathcal{G} = \mathcal{F}$。
\end{enumerate}
\end{lemma}

\begin{proof}
$\text{Supp}(\mathcal{F})$（即 $\mathcal{F}$ 的支撑）是 $X$ 的闭子集；
见《态射》引理 \ref{morphisms-lemma-support-finite-type}。
故可应用定义 \ref{coherent-definition-proper-over-base}。
由于 $\mathcal{F}$ 的概形论支撑是一个闭子概形，其底层闭子集为
$\text{Supp}(\mathcal{F})$，定义 \ref{coherent-definition-proper-over-base}
表明 (1) 与 (2) 等价。(2) 蕴含 (3) 是清楚的。
反之，若 (3) 成立，则
$\text{Supp}(\mathcal{F}) \subset Z$
（这是 $X$ 的闭子集之间的包含），因而例如由引理
\ref{coherent-lemma-closed-closed-proper-over-base}，
$\text{Supp}(\mathcal{F})$ 在 $S$ 上固有。
\end{proof}

\begin{lemma}
\label{coherent-lemma-base-change-module-support-proper-over-base}
考虑概形的笛卡尔图
$$
\xymatrix{
X' \ar[d]_{f'} \ar[r]_{g'} & X \ar[d]^f \\
S' \ar[r]^g & S
}
$$
其中 $f$ 局部有限型。设 $\mathcal{F}$ 为有限型准凝聚
$\mathcal{O}_X$-模。若 $\mathcal{F}$ 的支撑在 $S$ 上固有，
则 $(g')^*\mathcal{F}$ 的支撑在 $S'$ 上固有。
\end{lemma}

\begin{proof}
陈述有意义，因为由《模》引理
\ref{modules-lemma-pullback-finite-type}，
$(g')*\mathcal{F}$ 是有限型的。由《态射》引理
\ref{morphisms-lemma-support-finite-type}，
有 $\text{Supp}((g')^*\mathcal{F}) = (g')^{-1}(\text{Supp}(\mathcal{F}))$。
故本引理由引理 \ref{coherent-lemma-base-change-closed-proper-over-base} 推出。
\end{proof}

\begin{lemma}
\label{coherent-lemma-cat-module-support-proper-over-base}
设 $f : X \to S$ 为局部有限型概形态射。设 $\mathcal{F}$、$\mathcal{G}$
为有限型准凝聚 $\mathcal{O}_X$-模。
\begin{enumerate}
\item 若 $\mathcal{F}$、$\mathcal{G}$ 的支撑都在 $S$ 上固有，
则 $\mathcal{F} \oplus \mathcal{G}$、$\mathcal{G}$ 被
$\mathcal{F}$ 的任意扩张、$\Im(u)$ 与 $\Coker(u)$
（这里所取的是任意 $\mathcal{O}_X$-模映射
$u : \mathcal{F} \to \mathcal{G}$），以及 $\mathcal{F}$ 或 $\mathcal{G}$
的任意准凝聚商也具有同一性质。
\item 若 $S$ 局部 Noether，则凝聚 $\mathcal{O}_X$-模中支撑在
$S$ 上固有者所成范畴，是凝聚 $\mathcal{O}_X$-模所成
阿贝尔范畴的 Serre 子范畴（《同调》定义
\ref{homology-definition-serre-subcategory}）。
\end{enumerate}
\end{lemma}

\begin{proof}
证明 (1)。设 $Z$、$Z'$ 分别为 $\mathcal{F}$ 与 $\mathcal{G}$ 的支撑。
则 (1) 中提到的所有层，其支撑都包含于 $Z \cup Z'$。
只要检验这些层为有限型且准凝聚，引理
\ref{coherent-lemma-closed-closed-proper-over-base} 与
\ref{coherent-lemma-union-closed-proper-over-base}
便表明断言成立。关于准凝聚性，参见《概形》第
\ref{schemes-section-quasi-coherent} 节。关于“有限型”，
建议读者参阅《模》第 \ref{modules-section-finite-type} 节。

\medskip\noindent
证明 (2)。证明与 (1) 的证明相同。注意这些断言有意义：
由《态射》引理 \ref{morphisms-lemma-finite-type-noetherian}，
$X$ 局部 Noether；还要使用第 \ref{coherent-section-coherent-sheaves} 节
对凝聚模范畴的描述。
\end{proof}

\begin{lemma}
\label{coherent-lemma-support-proper-over-base-pushforward}
设 $S$ 为局部 Noether 概形。设 $f : X \to S$ 为局部有限型概形态射。
设 $\mathcal{F}$ 为凝聚 $\mathcal{O}_X$-模，其支撑在 $S$ 上固有。
则 $R^pf_*\mathcal{F}$ 是凝聚 $\mathcal{O}_S$-模，
这对所有 $p \geq 0$ 成立。
\end{lemma}

\begin{proof}
由引理 \ref{coherent-lemma-module-support-proper-over-base}，存在闭浸入
$i : Z \to X$ 及有限型准凝聚 $\mathcal{O}_Z$-模 $\mathcal{G}$，
使得 (a) $g = f \circ i : Z \to S$ 固有，且 (b)
$i_*\mathcal{G} = \mathcal{F}$。由命题
\ref{coherent-proposition-proper-pushforward-coherent}，
$R^pg_*\mathcal{G}$ 在 $S$ 上凝聚。另一方面，
$R^qi_*\mathcal{G} = 0$ 对 $q > 0$ 成立（引理
\ref{coherent-lemma-finite-pushforward-coherent}）。
由《上同调》引理 \ref{cohomology-lemma-relative-Leray}，
$R^pf_*\mathcal{F} = R^pg_*\mathcal{G}$，故结论成立。
\end{proof}

\begin{lemma}
\label{coherent-lemma-systems-with-proper-support}
设 $S$ 为 Noether 概形，$f : X \to S$ 为有限型态射，并设
$\mathcal{I} \subset \mathcal{O}_X$ 为准凝聚理想层。下列范畴是
$\textit{Coh}(X, \mathcal{I})$ 的 Serre 子范畴：
\begin{enumerate}
\item $\textit{Coh}(X, \mathcal{I})$ 的满子范畴，其中的对象
$(\mathcal{F}_n)$ 满足 $\mathcal{F}_1$ 的支撑在 $S$ 上固有；
\item $\textit{Coh}(X, \mathcal{I})$ 的满子范畴，其中的对象
$(\mathcal{F}_n)$ 满足：存在闭子概形 $Z \subset X$，
它在 $S$ 上固有，并且
$\mathcal{I}_Z \mathcal{F}_n = 0$ 对所有 $n \geq 1$ 成立。
\end{enumerate}
\end{lemma}

\begin{proof}
将使用《同调》引理 \ref{homology-lemma-characterize-serre-subcategory}
的判据。此外，还将使用如下事实：若
$0 \to (\mathcal{G}_n) \to (\mathcal{F}_n) \to (\mathcal{H}_n) \to 0$
是 $\textit{Coh}(X, \mathcal{I})$ 的短正合列，则
(a) $\mathcal{G}_n \to \mathcal{F}_n \to \mathcal{H}_n \to 0$
对所有 $n \geq 1$ 正合；并且 (b) $\mathcal{G}_n$ 是
$\Ker(\mathcal{F}_m \to \mathcal{H}_m)$ 的商，其中某个 $m \geq n$。
见引理 \ref{coherent-lemma-inverse-systems-abelian} 的证明。

\medskip\noindent
证明 (1)。设 $(\mathcal{F}_n)$ 为
$\textit{Coh}(X, \mathcal{I})$ 的对象。则
$\text{Supp}(\mathcal{F}_n) = \text{Supp}(\mathcal{F}_1)$
对所有 $n \geq 1$ 成立。
因此，由上述 (a) 与 (b)，对
任意短正合列
$0 \to (\mathcal{G}_n) \to (\mathcal{F}_n) \to (\mathcal{H}_n) \to 0$
（在 $\textit{Coh}(X, \mathcal{I})$ 中），
都有
$\text{Supp}(\mathcal{G}_1) \cup \text{Supp}(\mathcal{H}_1) =
\text{Supp}(\mathcal{F}_1)$。
这证明 (1) 所定义的范畴是
$\textit{Coh}(X, \mathcal{I})$ 的 Serre 子范畴。

\medskip\noindent
证明 (2)。这里作同样的论证。设
$0 \to (\mathcal{G}_n) \to (\mathcal{F}_n) \to (\mathcal{H}_n) \to 0$
为 $\textit{Coh}(X, \mathcal{I})$ 的短正合列。
若 $Z \subset X$ 为闭子概形，且 $\mathcal{I}_Z$
零化 $\mathcal{F}_n$（对所有 $n$），则由上述 (a) 与 (b)，
$\mathcal{I}_Z$ 零化 $\mathcal{G}_n$ 与
$\mathcal{H}_n$（对所有 $n$）。因此，若 $Z \to S$ 固有，便可知 (2)
所定义的范畴在 $\textit{Coh}(X, \mathcal{I})$ 中对取子对象和
商对象封闭。最后，设 $Z \subset X$ 与 $Y \subset X$ 是在 $S$
上固有的闭子概形，并且
$\mathcal{I}_Z \mathcal{G}_n = 0$ 且
$\mathcal{I}_Y \mathcal{H}_n = 0$，这对所有 $n \geq 1$ 成立。
由上述 (a) 可知
$\mathcal{I}_{Z \cup Y} = \mathcal{I}_Z \cdot \mathcal{I}_Y$
零化 $\mathcal{F}_n$（对所有 $n$）。由引理
\ref{coherent-lemma-union-closed-proper-over-base}
（以及定义 \ref{coherent-definition-proper-over-base}，它表明在该并集上
可以任选概形结构），$Z \cup Y \to S$ 固有，证明完成。
\end{proof}

\section{Grothendieck 存在定理（三）}
\label{coherent-section-existence-proper-support}

\noindent
为陈述 Grothendieck 存在定理的一般形式，先引入一些记号。
设 $A$ 为关于理想 $I$ 完备的 Noether 环。设
$f : X \to \Spec(A)$ 为分离有限型概形态射，并令
$\mathcal{I} = I\mathcal{O}_X$。在此情形下，以
$$
\textit{Coh}_{\text{support proper over } A}(\mathcal{O}_X)
$$
表示 $\textit{Coh}(\mathcal{O}_X)$ 的满子范畴，其中的对象是凝聚
$\mathcal{O}_X$-模，且其支撑在 $\Spec(A)$ 上固有。由引理
\ref{coherent-lemma-cat-module-support-proper-over-base}，这是
$\textit{Coh}(\mathcal{O}_X)$ 的 Serre 子范畴。类似地，以
$$
\textit{Coh}_{\text{support proper over } A}(X, \mathcal{I})
$$
表示 $\textit{Coh}(X, \mathcal{I})$ 的满子范畴，其中的对象
$(\mathcal{F}_n)$ 满足：$\mathcal{F}_1$ 的支撑在
$\Spec(A)$ 上固有。由引理
\ref{coherent-lemma-systems-with-proper-support} 的 (1)，这是
$\textit{Coh}(X, \mathcal{I})$ 的 Serre 子范畴。
由于商模的支撑包含在原模的支撑中，
(\ref{coherent-equation-completion-functor}) 诱导函子
\begin{equation}
\label{coherent-equation-completion-functor-proper-over-A}
\textit{Coh}_{\text{support proper over }A}(\mathcal{O}_X)
\longrightarrow
\textit{Coh}_{\text{support proper over }A}(X, \mathcal{I})
\end{equation}
现在可以陈述本节的主要定理。

\begin{theorem}[Grothendieck 存在定理]
\label{coherent-theorem-grothendieck-existence}
\begin{reference}
\cite[III Theorem 5.1.5]{EGA}
\end{reference}
设 $A$ 为关于理想 $I$ 完备的 Noether 环。设 $X$ 为 $A$ 上分离的
有限型概形。则函子
(\ref{coherent-equation-completion-functor-proper-over-A})
$$
\textit{Coh}_{\text{support proper over }A}(\mathcal{O}_X)
\longrightarrow
\textit{Coh}_{\text{support proper over }A}(X, \mathcal{I})
$$
是等价。
\end{theorem}

\begin{proof}
将不再特别说明地使用引理 \ref{coherent-lemma-i-star-equivalence}
给出的范畴等价。在本证明中，对闭子概形
$Z \subset X$（它在 $A$ 上固有），若 $X$ 上的凝聚模被 $Z$ 的理想层零化，
或者等价地，它是 $Z$ 上凝聚模的直像，就称该模“支撑在 $Z$ 上”。
由命题 \ref{coherent-proposition-existence-proper}，已知凝聚模与支撑在
$Z$ 上的凝聚模系统之间的函子具有所述结论。因此，只须证明
$\textit{Coh}_{\text{support proper over }A}(\mathcal{O}_X)$ 的每个对象
以及 $\textit{Coh}_{\text{support proper over }A}(X, \mathcal{I})$
的每个对象，都支撑在某个闭子概形
$Z \subset X$ 上，该子概形在 $A$ 上固有。对
$\textit{Coh}_{\text{support proper over }A}(\mathcal{O}_X)$ 的对象，
这由定义成立。下面将依照命题 \ref{coherent-proposition-existence-proper}
的证明方法，对
$\textit{Coh}_{\text{support proper over }A}(X, \mathcal{I})$
的对象证明该断言。建议读者先阅读该命题的证明。

\medskip\noindent
考虑集合 $\Xi$，其元素为准凝聚理想层
$\mathcal{K} \subset \mathcal{O}_X$，并具有如下性质：对每个对象
$(\mathcal{F}_n)$（属于
$\textit{Coh}_{\text{support proper over }A}(X, \mathcal{I})$）
若它被 $\mathcal{K}$ 零化，
上述断言均成立。要证明 $(0)$ 属于 $\Xi$。
若不然，由于 $X$ 是 Noether 概形，存在极大准凝聚理想层
$\mathcal{K}$，它不属于 $\Xi$；见引理 \ref{coherent-lemma-acc-coherent}。
将 $X$ 替换为 $X$ 中与 $\mathcal{K}$ 对应的闭子概形后，
可以假设每个非零 $\mathcal{K}$ 均属于 $\Xi$。
设 $(\mathcal{F}_n)$ 为
$\textit{Coh}_{\text{support proper over }A}(X, \mathcal{I})$ 的对象。
将证明该对象支撑在某个闭子概形
$Z \subset X$ 上，该子概形在 $A$ 上固有，从而完成定理的证明。

\medskip\noindent
应用 Chow 引理（引理 \ref{coherent-lemma-chow-Noetherian}），得到固有满射
$f : Y \to X$，它在稠密开集 $U \subset X$ 上为同构，且 $Y$
在 $A$ 上 H-拟射影。选取开浸入 $j : Y \to Y'$，其中 $Y'$
在 $A$ 上射影；见《态射》引理
\ref{morphisms-lemma-H-quasi-projective-open-H-projective}。
注意
$$
\text{Supp}(f^*\mathcal{F}_n) = f^{-1}\text{Supp}(\mathcal{F}_n) =
f^{-1}\text{Supp}(\mathcal{F}_1)
$$
第一个等式来自《态射》引理
\ref{morphisms-lemma-support-finite-type}。由假设以及引理
\ref{coherent-lemma-functoriality-closed-proper-over-base} 的 (3)，
$f^{-1}\text{Supp}(\mathcal{F}_1)$ 在 $A$ 上固有。
因此，由引理 \ref{coherent-lemma-functoriality-closed-proper-over-base} 的 (1)，
$f^{-1}\text{Supp}(\mathcal{F}_1)$ 在 $j$ 下的像是 $Y'$ 中的闭集。
故由引理 \ref{coherent-lemma-pushforward-coherent-on-open}，
$\mathcal{F}'_n = j_*f^*\mathcal{F}_n$ 在 $Y'$ 上凝聚。
于是 $(\mathcal{F}_n')$ 是
$\textit{Coh}(Y', I\mathcal{O}_{Y'})$ 的对象。
由射影情形的 Grothendieck 存在定理
（引理 \ref{coherent-lemma-existence-projective}），存在凝聚
$\mathcal{O}_{Y'}$-模 $\mathcal{F}'$ 及同构
$(\mathcal{F}')^\wedge \cong (\mathcal{F}'_n)$
（在 $\textit{Coh}(Y', I\mathcal{O}_{Y'})$ 中）。
由于 $\mathcal{F}'/I\mathcal{F}' = \mathcal{F}'_1$，可知
$$
\text{Supp}(\mathcal{F}') \cap V(I\mathcal{O}_{Y'}) =
\text{Supp}(\mathcal{F}'_1) = j(f^{-1}\text{Supp}(\mathcal{F}_1))
$$
结构态射 $p' : Y' \to \Spec(A)$ 固有，故
$p'(\text{Supp}(\mathcal{F}') \setminus j(Y))$
是 $\Spec(A)$ 中的闭集。$\Spec(A)$ 的每个非空闭子集都包含
$V(I)$ 的一点，因为由《代数》引理
\ref{algebra-lemma-radical-completion}，$I$ 包含于 $A$ 的 Jacobson 根。
上述等式表明
$\text{Supp}(\mathcal{F}') \cap (p')^{-1}V(I) \subset j(Y)$，
从而 $\text{Supp}(\mathcal{F}') \subset j(Y)$。
因此，$\mathcal{F}'|_Y = j^*\mathcal{F}'$
支撑在某个闭子概形 $Z'$ 上，该子概形位于 $Y$ 中且在 $A$ 上固有，并且
$(\mathcal{F}'|_Y)^\wedge = (f^*\mathcal{F}_n)$。

\medskip\noindent
设 $\mathcal{K}$ 为截出约化补集 $X \setminus U$ 的准凝聚理想层。
由命题 \ref{coherent-proposition-proper-pushforward-coherent}，
$\mathcal{O}_X$-模 $\mathcal{H} = f_*(\mathcal{F}'|_Y)$ 凝聚；
由引理 \ref{coherent-lemma-inverse-systems-push-pull}，存在态射
$\alpha : (\mathcal{F}_n) \to \mathcal{H}^\wedge$
（在 $\textit{Coh}(X, \mathcal{I})$ 中），其核与余核被幂
$\mathcal{K}^t$ 零化，其中该幂来自 $\mathcal{K}$。得到正合列
$$
0 \to \Ker(\alpha) \to (\mathcal{F}_n) \to
\mathcal{H}^\wedge \to \Coker(\alpha) \to 0
$$
（在 $\textit{Coh}(X, \mathcal{I})$ 中）。
若 $Z_0 \subset X$ 是 $\mathcal{H}$ 的概形论支撑，则在集合意义下，
$Z_0 \subset f(Z')$ 是清楚的。于是，由引理
\ref{coherent-lemma-closed-closed-proper-over-base} 与引理
\ref{coherent-lemma-functoriality-closed-proper-over-base} 的 (2)，
$Z_0$ 在 $A$ 上固有。故 $\mathcal{H}^\wedge$
属于引理 \ref{coherent-lemma-systems-with-proper-support} 的 (2)
所定义的子范畴，从而尤其属于
$\textit{Coh}_{\text{support proper over }A}(X, \mathcal{I})$。
由引理 \ref{coherent-lemma-systems-with-proper-support} 的 (1)，
可知 $\Ker(\alpha)$ 与 $\Coker(\alpha)$ 属于
$\textit{Coh}_{\text{support proper over }A}(X, \mathcal{I})$。
由归纳假设，更确切地说，因为 $\mathcal{K}^t$ 属于 $\Xi$，
$\Ker(\alpha)$ 与 $\Coker(\alpha)$ 属于引理
\ref{coherent-lemma-systems-with-proper-support} 的 (2) 所定义的子范畴。
由于按该引理这是 Serre 子范畴，$(\mathcal{F}_n)$ 也属于其中，
这正是所要证明的。
\end{proof}

\begin{remark}[展开 Grothendieck 存在定理]
\label{coherent-remark-reformulate-existence-theorem}
设 $A$ 为关于理想 $I$ 完备的 Noether 环。记
$S = \Spec(A)$ 与 $S_n = \Spec(A/I^n)$。
设 $X \to S$ 为分离有限型态射。对 $n \geq 1$，令
$X_n = X \times_S S_n$。图示如下：
$$
\xymatrix{
X_1 \ar[r]_{i_1} \ar[d] & X_2 \ar[r]_{i_2} \ar[d] & X_3 \ar[r] \ar[d] &
\ldots & X \ar[d] \\
S_1 \ar[r] & S_2 \ar[r] & S_3 \ar[r] & \ldots & S
}
$$
在此情形下，考虑系统 $(\mathcal{F}_n, \varphi_n)$，其中
\begin{enumerate}
\item $\mathcal{F}_n$ 是凝聚 $\mathcal{O}_{X_n}$-模；
\item $\varphi_n : i_n^*\mathcal{F}_{n + 1} \to \mathcal{F}_n$
是同构；并且
\item $\text{Supp}(\mathcal{F}_1)$ 在 $S_1$ 上固有。
\end{enumerate}
定理 \ref{coherent-theorem-grothendieck-existence} 断言，完备化函子
$$
\begin{matrix}
\text{coherent }\mathcal{O}_X\text{-模 }\mathcal{F} \\
\text{with support proper over }A
\end{matrix}
\quad
\longrightarrow
\quad
\begin{matrix}
\text{systems }(\mathcal{F}_n) \\
\text{如上}
\end{matrix}
$$
是范畴等价。在 $X$ 在 $A$ 上固有这一特殊情形下，可以省略支撑条件。
\end{remark}

\section{Grothendieck 代数化定理}
\label{coherent-section-algebraization}

\noindent
第一个结论是用闭子概形和有限态射来表述 Grothendieck 存在定理。

\begin{lemma}
\label{coherent-lemma-algebraize-formal-closed-subscheme}
设 $A$ 为关于理想 $I$ 完备的 Noether 环。记
$S = \Spec(A)$ 与 $S_n = \Spec(A/I^n)$。
设 $X \to S$ 为分离有限型态射。对 $n \geq 1$，令
$X_n = X \times_S S_n$。假设给定交换图
$$
\xymatrix{
Z_1 \ar[r] \ar[d] & Z_2 \ar[r] \ar[d] & Z_3 \ar[r] \ar[d] & \ldots \\
X_1 \ar[r]^{i_1} & X_2 \ar[r]^{i_2} & X_3 \ar[r] & \ldots
}
$$
（诸方块均为笛卡尔方块），并假设：
\begin{enumerate}
\item $Z_1 \to X_1$ 是闭浸入；并且
\item $Z_1 \to S_1$ 固有。
\end{enumerate}
则存在概形的闭浸入 $Z \to X$，使得
$Z_n = Z \times_S S_n$。而且，$Z$ 在 $S$ 上固有。
\end{lemma}

\begin{proof}
以 $j_n : Z_n \to X_n$ 表示竖直态射。由于陈述中的方块是笛卡尔的，
$j_n$ 到 $X_1$ 的基变换就是 $j_1$。因此，《态射》引理
\ref{morphisms-lemma-check-closed-infinitesimally} 表明 $j_n$ 是闭浸入。
令 $\mathcal{F}_n = j_{n, *}\mathcal{O}_{Z_n}$，从而
$j_n^\sharp$ 是满射 $\mathcal{O}_{X_n} \to \mathcal{F}_n$。
再次利用方块的笛卡尔性可知，$\mathcal{F}_{n + 1}$ 到 $X_n$
的拉回就是 $\mathcal{F}_n$。故按评注
\ref{coherent-remark-reformulate-existence-theorem} 中的重新表述，
Grothendieck 存在定理告诉我们，存在映射
$\mathcal{O}_X \to \mathcal{F}$（在凝聚 $\mathcal{O}_X$-模之间），
其到 $X_n$ 的限制恢复
$\mathcal{O}_{X_n} \to \mathcal{F}_n$。而且，$\mathcal{F}$
的支撑在 $S$ 上固有。由于完备化函子正合（引理 \ref{coherent-lemma-exact}），
余核 $\mathcal{Q}$（对应于 $\mathcal{O}_X \to \mathcal{F}$）
具有零完备化。
由于 $\mathcal{F}$ 的支撑在 $S$ 上固有，$\mathcal{Q}$ 也如此，
这蕴含 $\mathcal{Q} = 0$；例如，这是因为
(\ref{coherent-equation-completion-functor-proper-over-A}) 按
Grothendieck 存在定理是等价。因此，有
$\mathcal{F} = \mathcal{O}_X/\mathcal{J}$，其中 $\mathcal{J}$
是某个准凝聚理想层。
令 $Z = V(\mathcal{J})$ 即完成证明。
\end{proof}

\noindent
在下述引理中，实际上只须假设 $Y_1 \to X_1$ 有限，因为这将蕴含
$Y_n \to X_n$ 也有限（见《态射进阶》引理
\ref{more-morphisms-lemma-thicken-property-morphisms-cartesian}）。

\begin{lemma}
\label{coherent-lemma-algebraize-formal-scheme-finite-over-proper}
设 $A$ 为关于理想 $I$ 完备的 Noether 环。记
$S = \Spec(A)$ 与 $S_n = \Spec(A/I^n)$。
设 $X \to S$ 为分离有限型态射。对 $n \geq 1$，令
$X_n = X \times_S S_n$。假设给定交换图
$$
\xymatrix{
Y_1 \ar[r] \ar[d] & Y_2 \ar[r] \ar[d] & Y_3 \ar[r] \ar[d] & \ldots \\
X_1 \ar[r]^{i_1} & X_2 \ar[r]^{i_2} & X_3 \ar[r] & \ldots
}
$$
（诸方块均为笛卡尔方块），并假设：
\begin{enumerate}
\item $Y_n \to X_n$ 是有限态射；并且
\item $Y_1 \to S_1$ 固有。
\end{enumerate}
则存在概形的有限态射 $Y \to X$，使得
$Y_n = Y \times_S S_n$。而且，$Y$ 在 $S$ 上固有。
\end{lemma}

\begin{proof}
以 $f_n : Y_n \to X_n$ 表示竖直态射。令
$\mathcal{F}_n = f_{n, *}\mathcal{O}_{Y_n}$。这是凝聚
$\mathcal{O}_{X_n}$-模，因为 $f_n$ 有限（引理
\ref{coherent-lemma-finite-pushforward-coherent}）。
利用方块的笛卡尔性可知，$\mathcal{F}_{n + 1}$ 到 $X_n$
的拉回就是 $\mathcal{F}_n$。故按评注
\ref{coherent-remark-reformulate-existence-theorem} 中的重新表述，
Grothendieck 存在定理告诉我们，存在凝聚 $\mathcal{O}_X$-模
$\mathcal{F}$，其到 $X_n$ 的限制恢复 $\mathcal{F}_n$。
而且，$\mathcal{F}$ 的支撑在 $S$ 上固有。
由于完备化函子全忠实（定理
\ref{coherent-theorem-grothendieck-existence}），诸乘法映射
$\mathcal{F}_n \otimes_{\mathcal{O}_{X_n}} \mathcal{F}_n \to
\mathcal{F}_n$ 相容地给出 $\mathcal{F}$ 上的代数结构。
令 $Y = \underline{\Spec}_X(\mathcal{F})$ 即完成证明。
\end{proof}

\begin{lemma}
\label{coherent-lemma-algebraize-morphism}
设 $A$ 为关于理想 $I$ 完备的 Noether 环。记
$S = \Spec(A)$ 与 $S_n = \Spec(A/I^n)$。设 $X$、$Y$ 为 $S$
上的概形。对 $n \geq 1$，令 $X_n = X \times_S S_n$ 及
$Y_n = Y \times_S S_n$。假设给定相容的交换图系统
$$
\xymatrix{
& & X_{n + 1} \ar[rd] \ar[rr]_{g_{n + 1}} & & Y_{n + 1} \ar[ld] \\
X_n \ar[rru] \ar[rd] \ar[rr]_{g_n} & & Y_n \ar[rru] \ar[ld] & S_{n + 1} \\
& S_n \ar[rru]
}
$$
并假设：
\begin{enumerate}
\item $X \to S$ 固有；并且
\item $Y \to S$ 分离且有限型。
\end{enumerate}
则存在唯一的概形态射 $g : X \to Y$（在 $S$ 上），使得
$g_n$ 是 $g$ 到 $S_n$ 的基变换。
\end{lemma}

\begin{proof}
态射 $(1, g_n) : X_n \to X_n \times_S Y_n$ 是闭浸入，因为
$Y_n \to S_n$ 分离（《概形》引理
\ref{schemes-lemma-section-immersion}）。因此，由引理
\ref{coherent-lemma-algebraize-formal-closed-subscheme}，存在闭子概形
$Z \subset X \times_S Y$，它在 $S$ 上固有，且到 $S_n$ 的基变换恢复
$X_n \subset X_n \times_S Y_n$。第一投影 $p : Z \to X$ 是固有态射
（因为 $Z$ 在 $S$ 上固有；见《态射》引理
\ref{morphisms-lemma-image-proper-scheme-closed}），其到 $S_n$ 的基变换
是同构，这对所有 $n$ 成立。特别地，由引理
\ref{coherent-lemma-proper-finite-fibre-finite-in-neighbourhood}，
$p : Z \to X$ 在 $X_0$ 的某个开邻域上有限。由于 $X$ 在 $S$ 上固有，
该开邻域就是整个 $X$，从而 $p : Z \to X$ 有限。
应用命题 \ref{coherent-proposition-existence-proper} 的等价可知
$p_*\mathcal{O}_Z = \mathcal{O}_X$，因为模 $I^n$ 后这一点对所有
$n$ 成立。因此 $p$ 是同构，并得到态射
$g$，它是复合 $X \cong Z \to Y$。唯一性的证明从略。
\end{proof}

\noindent
为了证明“抽象”的代数化定理，必须假设有一个丰沛可逆层；
没有这一假设，结论不成立。

\begin{theorem}[Grothendieck 代数化定理]
\label{coherent-theorem-algebraization}
设 $A$ 为关于理想 $I$ 完备的 Noether 环。令
$S = \Spec(A)$ 且 $S_n = \Spec(A/I^n)$。考虑交换图
$$
\xymatrix{
X_1 \ar[r]_{i_1} \ar[d] & X_2 \ar[r]_{i_2} \ar[d] & X_3 \ar[r] \ar[d] &
\ldots \\
S_1 \ar[r] & S_2 \ar[r] & S_3 \ar[r] & \ldots
}
$$
（诸方块均为笛卡尔方块）。假设给定 $(\mathcal{L}_n, \varphi_n)$，
其中每个 $\mathcal{L}_n$ 都是 $X_n$ 上的可逆层，并且
$\varphi_n : i_n^*\mathcal{L}_{n + 1} \to \mathcal{L}_n$
是同构。若
\begin{enumerate}
\item $X_1 \to S_1$ 固有；并且
\item $\mathcal{L}_1$ 在 $X_1$ 上丰沛，
\end{enumerate}
则存在概形的固有态射 $X \to S$、丰沛可逆
$\mathcal{O}_X$-模 $\mathcal{L}$，以及同构
$X_n \cong X \times_S S_n$ 与
$\mathcal{L}_n \cong \mathcal{L}|_{X_n}$，它们与态射
$i_n$ 和 $\varphi_n$ 相容。
\end{theorem}

\begin{proof}
由于图中的方块是笛卡尔的，且态射 $S_n \to S_{n + 1}$ 是闭浸入，
可知态射 $i_n$ 也是闭浸入。特别地，可以把
$X_m$ 看作 $X_n$ 的闭子概形（当 $m < n$ 时）。事实上，$X_m$ 是由准凝聚理想层
$I^m\mathcal{O}_{X_n}$ 截出的闭子概形。而且，概形
$X_1, X_2, X_3, \ldots$ 的底层拓扑空间都被等同，故可以（且确实）
把层 $\mathcal{O}_{X_n}$ 看作位于同一个底层拓扑空间上；
凝聚 $\mathcal{O}_{X_n}$-模亦然。令
$$
\mathcal{F}_n =
\Ker(\mathcal{O}_{X_{n + 1}} \to \mathcal{O}_{X_n})
$$
从而得到短正合列
$$
0 \to \mathcal{F}_n \to \mathcal{O}_{X_{n + 1}} \to \mathcal{O}_{X_n} \to 0
$$
由上可知 $\mathcal{F}_n = I^n\mathcal{O}_{X_{n + 1}}$。
于是 $\mathcal{F}_n$ 是 $X_{n + 1}$ 上被 $I$ 零化的凝聚层，
故可以（且确实）把它看作凝聚 $\mathcal{O}_{X_1}$-模。
注意，对 $m > n$，层
$$
I^n\mathcal{O}_{X_m}/I^{n + 1}\mathcal{O}_{X_m}
$$
同构地映到 $\mathcal{F}_n$，所用映射为
$\mathcal{O}_{X_m} \to \mathcal{O}_{X_{n + 1}}$。因此，给定
$n_1, n_2 \geq 0$，
可以选取 $m > n_1 + n_2$ 并考虑乘法映射
$$
I^{n_1}\mathcal{O}_{X_m} \times I^{n_2}\mathcal{O}_{X_m}
\longrightarrow
I^{n_1 + n_2}\mathcal{O}_{X_m} \to \mathcal{F}_{n_1 + n_2}
$$
它诱导 $\mathcal{O}_{X_1}$-双线性映射
$$
\mathcal{F}_{n_1} \times \mathcal{F}_{n_2} \longrightarrow
\mathcal{F}_{n_1 + n_2}
$$
继而定义分次 $\mathcal{O}_{X_1}$-代数结构，其底层分次模为
$\mathcal{F} = \bigoplus_{n \geq 0} \mathcal{F}_n$。

\medskip\noindent
令 $B = \bigoplus I^n/I^{n + 1}$；这是有限生成分次
$A/I$-代数。令 $\mathcal{B} = (X_1 \to S_1)^*\widetilde{B}$。
上述讨论给出分次 $\mathcal{O}_{X_1}$-代数的典范满射
$$
\mathcal{B} \longrightarrow \mathcal{F}
$$
特别地，$\mathcal{F}$ 是有限型准凝聚分次 $\mathcal{B}$-模。
由引理 \ref{coherent-lemma-graded-finiteness}，可找到整数 $d_0$，使得
$H^1(X_1, \mathcal{F} \otimes \mathcal{L}^{\otimes d}) = 0$
对所有 $d \geq d_0$ 成立。选取 $d \geq d_0$，使得存在截面
$s_{0, 1}, \ldots, s_{N, 1} \in \Gamma(X_1, \mathcal{L}_1^{\otimes d})$，
它们诱导浸入
$$
\psi_1 : X_1 \to \mathbf{P}^N_{S_1}
$$
（在 $S_1$ 上）；见《态射》引理
\ref{morphisms-lemma-finite-type-over-affine-ample-very-ample}。
由于 $X_1$ 在 $S_1$ 上固有，可知 $\psi_1$ 是闭浸入；见
《态射》引理 \ref{morphisms-lemma-image-proper-scheme-closed}
及《概形》引理 \ref{schemes-lemma-immersion-when-closed}。
将把 $\psi_1$“提升”为 $X_n$ 到 $\mathbf{P}^N$ 的相容闭浸入系统。

\medskip\noindent
把证明第一段中的短正合列与
$\mathcal{L}_{n + 1}^{\otimes d}$ 张量，得到短正合列
$$
0 \to \mathcal{F}_n \otimes \mathcal{L}_{n + 1}^{\otimes d} \to
\mathcal{L}_{n + 1}^{\otimes d} \to \mathcal{L}_{n + 1}^{\otimes d} \to 0
$$
使用同构 $\varphi_n$，得到同构
$\mathcal{L}_{n + 1} \otimes \mathcal{O}_{X_l} = \mathcal{L}_l$
（对 $l \leq n$）。因此上述序列变为
$$
0 \to \mathcal{F}_n \otimes \mathcal{L}_1^{\otimes d} \to
\mathcal{L}_{n + 1}^{\otimes d} \to \mathcal{L}_n^{\otimes d} \to 0
$$
$H^1(X, \mathcal{F}_n \otimes \mathcal{L}_1^{\otimes d})$ 的消失
使我们能够归纳地把
$s_{0, 1}, \ldots, s_{N, 1} \in \Gamma(X_1, \mathcal{L}_1^{\otimes d})$
提升为截面
$s_{0, n}, \ldots, s_{N, n} \in \Gamma(X_n, \mathcal{L}_n^{\otimes d})$。
从而得到交换图
$$
\xymatrix{
X_1 \ar[r]_{i_1} \ar[d]_{\psi_1} &
X_2 \ar[r]_{i_2} \ar[d]_{\psi_2} &
X_3 \ar[r] \ar[d]_{\psi_3} &
\ldots \\
\mathbf{P}^N_{S_1} \ar[r] &
\mathbf{P}^N_{S_2} \ar[r] &
\mathbf{P}^N_{S_3} \ar[r] & \ldots
}
$$
其中，按《构造》第 \ref{constructions-section-projective-space} 节的记号，
$\psi_n = \varphi_{(\mathcal{L}_n, (s_{0, n}, \ldots, s_{N, n}))}$。
由于定理陈述中的方块是笛卡尔的，上图中的方块也是笛卡尔的。
应用引理 \ref{coherent-lemma-algebraize-formal-closed-subscheme} 即得结论。
\end{proof}
% Generated mechanically from the frozen divisors.tex target.
% Do not edit this generated file; edit the bound locale source instead.

% OK, start here.
%
\raggedbottom

\setcounter{chapter}{30}
\chapter{除子}



\phantomsection
\label{divisors-section-phantom}


\section{引言}
\label{divisors-section-introduction}

\noindent
本章研究定义除子等问题所涉及的一些最基本问题。
基础参考文献是 \cite{EGA}。



\section{伴随点}
\label{divisors-section-associated}

\noindent
设 $R$ 为环，$M$ 为 $R$-模。
回忆，若素理想 $\mathfrak p \subset R$ 具有如下性质，
则称它是 $M$ 的{\it 伴随素理想}：$M$ 中存在某个元素，
其零化子为 $\mathfrak p$。
参见《代数》定义 \ref{algebra-definition-associated}。
下面给出概形上拟凝聚层的伴随点定义，
该定义取自 \cite[IV Definition 3.1.1]{EGA}。

\begin{definition}
\label{divisors-definition-associated}
设 $X$ 为概形。
设 $\mathcal{F}$ 为 $X$ 上的拟凝聚层。
\begin{enumerate}
\item 称 $x \in X$ 是 $\mathcal{F}$ 的{\it 伴随点}，
如果极大理想 $\mathfrak m_x$ 是 $\mathcal{O}_{X, x}$-模
$\mathcal{F}_x$ 的伴随素理想。
\item 以 $\text{Ass}(\mathcal{F})$ 或 $\text{Ass}_X(\mathcal{F})$
表示 $\mathcal{F}$ 的伴随点集。
\item 所谓 $X$ 的{\it 伴随点}，是指 $\mathcal{O}_X$ 的伴随点。
\end{enumerate}
\end{definition}

\noindent
当 $X$ 局部 Noether 且 $\mathcal{F}$ 为有限型时，
这些定义最为有用。
例如，$X$ 的某个不可约分支的泛点可能不是 $X$ 的伴随点，见
例 \ref{divisors-example-no-associated-prime}。
在非 Noether 情形下，使用弱伴随点有时更为方便，见
第 \ref{divisors-section-weakly-associated} 节。
下面把概形论概念与仿射开集上的代数概念联系起来；
注意，这一对应只有对局部 Noether 概形才完全成立。

\begin{lemma}
\label{divisors-lemma-associated-affine-open}
设 $X$ 为概形，$\mathcal{F}$ 为 $X$ 上的拟凝聚层。
设 $\Spec(A) = U \subset X$ 为仿射开集，并置
$M = \Gamma(U, \mathcal{F})$。
设 $x \in U$，并令 $\mathfrak p \subset A$ 为与之对应的素理想。
\begin{enumerate}
\item 若 $\mathfrak p$ 是 $M$ 的伴随素理想，则 $x$ 是
$\mathcal{F}$ 的伴随点。
\item 若 $\mathfrak p$ 有限生成，则逆命题也成立。
\end{enumerate}
特别地，若 $X$ 局部 Noether，则对上述所有配对
$(\mathfrak p, x)$ 都有等价关系
$$
\mathfrak p \in \text{Ass}(M) \Leftrightarrow x \in \text{Ass}(\mathcal{F})
$$
成立。
\end{lemma}

\begin{proof}
这由《代数》引理
\ref{algebra-lemma-associated-primes-localize} 推出。
也可如下直接论证。
假设 $\mathfrak p$ 是 $M$ 的伴随素理想。
则存在 $m \in M$，其零化子为 $\mathfrak p$。
由于局部化是正合的，可见
$\mathfrak pA_{\mathfrak p}$ 是
$m/1 \in M_{\mathfrak p}$ 的零化子。又因 $M_{\mathfrak p} = \mathcal{F}_x$
（《概形》引理 \ref{schemes-lemma-spec-sheaves}），
可知 $x$ 是 $\mathcal{F}$ 的伴随点。

\medskip\noindent
反之，假设 $x$ 是 $\mathcal{F}$ 的伴随点，且
$\mathfrak p$ 有限生成。
由 $x$ 是 $\mathcal{F}$ 的伴随点，
存在 $m' \in M_{\mathfrak p}$，其零化子为
$\mathfrak pA_{\mathfrak p}$。将其写成
$m' = m/f$，其中 $f \in A$ 且 $f \not \in \mathfrak p$。
设 $I$ 为 $m$ 的零化子。它是 $A$ 的理想，并满足
$IA_{\mathfrak p} = \mathfrak pA_{\mathfrak p}$。因此
$I \subset \mathfrak p$，且 $(\mathfrak p/I)_{\mathfrak p} = 0$。
由于 $\mathfrak p$ 有限生成，存在
$g \in A$ 且 $g \not \in \mathfrak p$，使得
$g(\mathfrak p/I) = 0$。因此 $gm$ 的零化子为
$\mathfrak p$，结论得证。

\medskip\noindent
若 $X$ 局部 Noether，则 $A$ 为 Noether 环
（《性质》引理 \ref{properties-lemma-locally-Noetherian}），
因而 $\mathfrak p$ 总是有限生成的。
\end{proof}

\begin{lemma}
\label{divisors-lemma-ass-support}
设 $X$ 为概形。
设 $\mathcal{F}$ 为拟凝聚 $\mathcal{O}_X$-模。
则 $\text{Ass}(\mathcal{F}) \subset \text{Supp}(\mathcal{F})$。
\end{lemma}

\begin{proof}
这由定义立即推出。
\end{proof}

\begin{lemma}
\label{divisors-lemma-ses-ass}
设 $X$ 为概形，并设
$0 \to \mathcal{F}_1 \to \mathcal{F}_2 \to \mathcal{F}_3 \to 0$
为 $X$ 上拟凝聚层的短正合列。则
$\text{Ass}(\mathcal{F}_2) \subset
\text{Ass}(\mathcal{F}_1) \cup \text{Ass}(\mathcal{F}_3)$，
并且
$\text{Ass}(\mathcal{F}_1) \subset \text{Ass}(\mathcal{F}_2)$。
\end{lemma}

\begin{proof}
对每个点 $x \in X$，茎的序列
$0 \to \mathcal{F}_{1, x} \to \mathcal{F}_{2, x} \to \mathcal{F}_{3, x} \to 0$
是 $\mathcal{O}_{X, x}$-模的短正合列。
因此本引理由《代数》引理 \ref{algebra-lemma-ass} 推出。
\end{proof}

\begin{lemma}
\label{divisors-lemma-finite-ass}
设 $X$ 为局部 Noether 概形。
设 $\mathcal{F}$ 为凝聚 $\mathcal{O}_X$-模。
则 $\text{Ass}(\mathcal{F}) \cap U$ 对每个拟紧开集
$U \subset X$ 均为有限集。
\end{lemma}

\begin{proof}
这是因为 Noether 环上的有限模只有有限多个伴随素理想；见
《代数》引理 \ref{algebra-lemma-finite-ass}。
为把概形陈述转化为代数陈述，使用如下事实：
$U$ 是有限多个仿射开集的并；每个这样的开集都是某个 Noether 环的谱
（《性质》引理 \ref{properties-lemma-locally-Noetherian}）；
$\mathcal{F}$ 对应于该环上的有限模
（《概形上同调》引理 \ref{coherent-lemma-coherent-Noetherian}）；
最后应用引理 \ref{divisors-lemma-associated-affine-open}。
\end{proof}

\begin{lemma}
\label{divisors-lemma-ass-zero}
设 $X$ 为局部 Noether 概形，$\mathcal{F}$ 为拟凝聚
$\mathcal{O}_X$-模。则
$$
\mathcal{F} = 0 \Leftrightarrow \text{Ass}(\mathcal{F}) = \emptyset.
$$
\end{lemma}

\begin{proof}
若 $\mathcal{F} = 0$，则由定义
$\text{Ass}(\mathcal{F}) = \emptyset$。
反之，若 $\text{Ass}(\mathcal{F}) = \emptyset$，则由
《代数》引理 \ref{algebra-lemma-ass-zero} 可知 $\mathcal{F} = 0$。
为把概形陈述转化为代数陈述，只需限制到任意仿射开集并应用
引理 \ref{divisors-lemma-associated-affine-open}。
\end{proof}

\begin{example}
\label{divisors-example-no-associated-prime}
设 $k$ 为域。环 $R = k[x_1, x_2, x_3, \ldots]/(x_i^2)$
是局部环，其极大理想 $\mathfrak m$ 局部幂零。
$R$ 中不存在零化子为 $\mathfrak m$ 的元素。
因此 $\text{Ass}(R) = \emptyset$，而 $X = \Spec(R)$
给出了一个没有伴随点的概形例子。
\end{example}

\begin{lemma}
\label{divisors-lemma-restriction-injective-open-contains-ass}
设 $X$ 为局部 Noether 概形，$\mathcal{F}$ 为拟凝聚
$\mathcal{O}_X$-模。若 $U \subset X$ 为开集且
$\text{Ass}(\mathcal{F}) \subset U$，则
$\Gamma(X, \mathcal{F}) \to \Gamma(U, \mathcal{F})$ 为单射。
\end{lemma}

\begin{proof}
设 $s \in \Gamma(X, \mathcal{F})$ 为限制到 $U$ 上等于零的截面。
令 $\mathcal{F}' \subset \mathcal{F}$ 为映射
$\mathcal{O}_X \to \mathcal{F}$ 的像，其中该映射由 $s$ 定义。于是
$\text{Supp}(\mathcal{F}') \cap U = \emptyset$。另一方面，由
引理 \ref{divisors-lemma-ses-ass} 有
$\text{Ass}(\mathcal{F}') \subset \text{Ass}(\mathcal{F})$。
又有
$\text{Ass}(\mathcal{F}') \subset \text{Supp}(\mathcal{F}')$
（引理 \ref{divisors-lemma-ass-support}），故
$\text{Ass}(\mathcal{F}') = \emptyset$。
最后由引理 \ref{divisors-lemma-ass-zero} 得 $\mathcal{F}' = 0$。
\end{proof}

\begin{lemma}
\label{divisors-lemma-minimal-support-in-ass}
设 $X$ 为局部 Noether 概形。
设 $\mathcal{F}$ 为拟凝聚 $\mathcal{O}_X$-模。
设 $x \in \text{Supp}(\mathcal{F})$ 为 $\mathcal{F}$ 的支撑中的点，
且它不是 $\text{Supp}(\mathcal{F})$ 中任何其他点的特殊化。
则 $x \in \text{Ass}(\mathcal{F})$。
特别地，$X$ 的每个不可约分支的泛点都是 $X$ 的伴随点。
\end{lemma}

\begin{proof}
由于 $x \in \text{Supp}(\mathcal{F})$，模 $\mathcal{F}_x$
非零。因此
$\text{Ass}(\mathcal{F}_x) \subset \Spec(\mathcal{O}_{X, x})$
由《代数》引理 \ref{algebra-lemma-ass-zero} 非空。
另一方面，由假设
$\text{Supp}(\mathcal{F}_x) = \{\mathfrak m_x\}$。
又因
$\text{Ass}(\mathcal{F}_x) \subset \text{Supp}(\mathcal{F}_x)$
（《代数》引理 \ref{algebra-lemma-ass-support}），
可知 $\mathfrak m_x$ 是 $\mathcal{F}_x$ 的伴随素理想，
结论得证。
\end{proof}

\noindent
下一个引理对应于《代数进阶》引理
\ref{more-algebra-lemma-check-injective-on-ass}。

\begin{lemma}
\label{divisors-lemma-check-injective-on-ass}
设 $X$ 为局部 Noether 概形，并设
$\varphi : \mathcal{F} \to \mathcal{G}$ 为拟凝聚
$\mathcal{O}_X$-模之间的映射。
假设对每个 $x \in X$，以下至少一种情形成立：
\begin{enumerate}
\item $\mathcal{F}_x \to \mathcal{G}_x$ 为单射；或
\item $x \not \in \text{Ass}(\mathcal{F})$。
\end{enumerate}
则 $\varphi$ 为单射。
\end{lemma}

\begin{proof}
这些假设蕴含 $\text{Ass}(\Ker(\varphi)) = \emptyset$，
故由引理 \ref{divisors-lemma-ass-zero} 有 $\Ker(\varphi) = 0$。
\end{proof}

\begin{lemma}
\label{divisors-lemma-check-isomorphism-via-depth-and-ass}
设 $X$ 为局部 Noether 概形，并设
$\varphi : \mathcal{F} \to \mathcal{G}$ 为拟凝聚
$\mathcal{O}_X$-模之间的映射。假设 $\mathcal{F}$ 凝聚，
并且对每个 $x \in X$，以下一种情形成立：
\begin{enumerate}
\item $\mathcal{F}_x \to \mathcal{G}_x$ 为同构；或
\item $\text{depth}(\mathcal{F}_x) \geq 2$ 且
$x \not \in \text{Ass}(\mathcal{G})$。
\end{enumerate}
则 $\varphi$ 为同构。
\end{lemma}

\begin{proof}
这只是把《代数进阶》引理
\ref{more-algebra-lemma-check-isomorphism-via-depth-and-ass}
改写成概形语言。
\end{proof}

\section{态射与伴随点}
\label{divisors-section-morphisms-associated}

\noindent
设 $f : X \to S$ 为概形的态射。
设 $\mathcal{F}$ 为 $\mathcal{O}_X$-模层。
若 $s \in S$ 是一点，则常用 $\mathcal{F}_s$ 表示下述
$\mathcal{O}_{X_s}$-模：它是把 $\mathcal{F}$ 沿态射
$i_s : X_s \to X$ 拉回所得的模。
这里 $X_s$ 是 $f$ 在 $s$ 上的概形论纤维。用公式表示即
$$
\mathcal{F}_s = i_s^*\mathcal{F}
$$
当然，这一记号会与 $\mathcal{F}$ 在点 $x \in X$ 处的茎所用的
既有记号冲突，这里指 $f = \text{id}_X$ 的情形。不过，这一记号通常很方便，
如下述引理的表述所示。

\begin{lemma}
\label{divisors-lemma-bourbaki}
设 $f : X \to S$ 为概形的态射。
设 $\mathcal{F}$ 为 $X$ 上在 $S$ 上平坦的拟凝聚层。
设 $\mathcal{G}$ 为 $S$ 上的拟凝聚层。则有
$$
\text{Ass}_X(\mathcal{F} \otimes_{\mathcal{O}_X} f^*\mathcal{G})
\supset
\bigcup\nolimits_{s \in \text{Ass}_S(\mathcal{G})}
\text{Ass}_{X_s}(\mathcal{F}_s)
$$
而当 $S$ 局部 Noether 时等号成立（记号 $\mathcal{F}_s$ 见上文）。
\end{lemma}

\begin{proof}
设 $x \in X$，并令 $s = f(x) \in S$。
置 $B = \mathcal{O}_{X, x}$、$A = \mathcal{O}_{S, s}$、
$N = \mathcal{F}_x$ 和 $M = \mathcal{G}_s$。
注意，$\mathcal{F} \otimes_{\mathcal{O}_X} f^*\mathcal{G}$ 在
$x$ 处的茎等于 $B$-模 $M \otimes_A N$。因此，
$x \in \text{Ass}_X(\mathcal{F} \otimes_{\mathcal{O}_X} f^*\mathcal{G})$
当且仅当 $\mathfrak m_B$ 属于 $\text{Ass}_B(M \otimes_A N)$。
类似地，$s \in \text{Ass}_S(\mathcal{G})$ 且
$x \in \text{Ass}_{X_s}(\mathcal{F}_s)$，当且仅当
$\mathfrak m_A \in \text{Ass}_A(M)$ 且
$\mathfrak m_B/\mathfrak m_A B \in
\text{Ass}_{B \otimes \kappa(\mathfrak m_A)}(N \otimes \kappa(\mathfrak m_A))$。
于是本引理由《代数》引理 \ref{algebra-lemma-bourbaki-fibres} 得出。
\end{proof}





\section{嵌入点}
\label{divisors-section-embedded}

\noindent
设 $R$ 为环，$M$ 为 $R$-模。
回忆，若素理想 $\mathfrak p \subset R$ 是 $M$ 的伴随素理想，
但在 $M$ 的伴随素理想中不是极小的，则称它为 $M$ 的
{\it 嵌入伴随素理想}。参见《代数》定义
\ref{algebra-definition-embedded-primes}。
下面给出概形上拟凝聚层的嵌入伴随点定义，
该定义取自 \cite[IV Definition 3.1.1]{EGA}。

\begin{definition}
\label{divisors-definition-embedded}
设 $X$ 为概形。
设 $\mathcal{F}$ 为 $X$ 上的拟凝聚层。
\begin{enumerate}
\item 称 $\mathcal{F}$ 的某个伴随点为{\it 嵌入伴随点}，
如果它在 $\mathcal{F}$ 的伴随点中不是极大的；换言之，
它是 $\mathcal{F}$ 的另一个伴随点的特化。
\item 设点 $x$ 所在的概形为 $X$。若 $x$ 是
$\mathcal{O}_X$ 的嵌入伴随点，则称其为{\it 嵌入点}。
\item $X$ 的一个{\it 嵌入分支}是不可约闭子集
$Z = \overline{\{x\}}$，其中 $x$ 是 $X$ 的嵌入点。
\end{enumerate}
\end{definition}

\noindent
在 Noether 情形下，当 $\mathcal{F}$ 凝聚时，有下述结论。

\begin{lemma}
\label{divisors-lemma-embedded}
设 $X$ 为局部 Noether 概形。
设 $\mathcal{F}$ 为凝聚 $\mathcal{O}_X$-模。则
\begin{enumerate}
\item $\text{Supp}(\mathcal{F})$ 的不可约分支的泛点是
$\mathcal{F}$ 的伴随点；
\item $\mathcal{F}$ 的伴随点是嵌入伴随点，当且仅当它不是
$\text{Supp}(\mathcal{F})$ 的某个不可约分支的泛点。
\end{enumerate}
特别地，$X$ 的嵌入点就是 $X$ 的伴随点中那些不是 $X$ 的
不可约分支泛点的点。
\end{lemma}

\begin{proof}
回忆，在此情形下 $Z = \text{Supp}(\mathcal{F})$ 是闭集，见
《态射》引理 \ref{morphisms-lemma-support-finite-type}；并且 $Z$
的不可约分支的泛点是 $\mathcal{F}$ 的伴随点，见引理
\ref{divisors-lemma-minimal-support-in-ass}。最后，由引理
\ref{divisors-lemma-ass-support} 有 $\text{Ass}(\mathcal{F}) \subset Z$。
结合这些结论以及 $Z$ 是 sober 拓扑空间这一事实（因而 $Z$ 的
每个点都是 $Z$ 的某个泛点的特化），便得到 (1) 和 (2)。
\end{proof}

\begin{lemma}
\label{divisors-lemma-S1-no-embedded}
设 $X$ 为局部 Noether 概形。
设 $\mathcal{F}$ 为 $X$ 上的凝聚层。则以下条件等价：
\begin{enumerate}
\item $\mathcal{F}$ 没有嵌入伴随点；
\item $\mathcal{F}$ 具有性质 $(S_1)$。
\end{enumerate}
\end{lemma}

\begin{proof}
这是《代数》引理 \ref{algebra-lemma-criterion-no-embedded-primes}
与上面的引理 \ref{divisors-lemma-associated-affine-open} 的结合。
\end{proof}

\begin{lemma}
\label{divisors-lemma-noetherian-dim-1-CM-no-embedded-points}
设 $X$ 为维数 $\leq 1$ 的局部 Noether 概形。以下条件等价：
\begin{enumerate}
\item $X$ 是 Cohen-Macaulay 概形；
\item $X$ 没有嵌入点。
\end{enumerate}
\end{lemma}

\begin{proof}
由引理 \ref{divisors-lemma-S1-no-embedded} 及各定义立即得到。
\end{proof}

\begin{lemma}
\label{divisors-lemma-scheme-theoretically-dense-contain-embedded-points}
设 $X$ 为局部 Noether 概形，并设 $U \subset X$ 为开子概形。
以下条件等价：
\begin{enumerate}
\item $U$ 在 $X$ 中概形论稠密
（《态射》定义 \ref{morphisms-definition-scheme-theoretically-dense}）；
\item $U$ 在 $X$ 中稠密，且 $U$ 包含 $X$ 的所有嵌入点；
\item $U$ 包含 $X$ 的所有伴随点。
\end{enumerate}
\end{lemma}

\begin{proof}
问题在 $X$ 上是局部的，故可假设 $X = \Spec(A)$，其中 $A$
为 Noether 环。此时 $U$ 拟紧（《性质》引理
\ref{properties-lemma-immersion-into-noetherian}），因而
$U = D(f_1) \cup \ldots \cup D(f_n)$
（《代数》引理 \ref{algebra-lemma-qc-open}）。在此情形下，
$U$ 在 $X$ 中概形论稠密，当且仅当
$A \to A_{f_1} \times \ldots \times A_{f_n}$ 是单射，见《态射》例
\ref{morphisms-example-scheme-theoretic-closure}。
条件 (3) 翻译为代数语言就是说：对每个伴随素理想
$\mathfrak p$，若它是 $A$ 的伴随素理想，则都存在某个 $i$
使得 $f_i \not \in \mathfrak p$。

\medskip\noindent
假设 (1)，即 $A \to A_{f_1} \times \ldots \times A_{f_n}$ 是单射。
设 $a \in A$ 的零化子是素理想 $\mathfrak p$。由于 $a$ 在某个
$A_{f_i}$ 中的像非零，便知对某个 $i$ 有
$f_i \not \in \mathfrak p$。故 (3) 成立。

\medskip\noindent
注意，$U$ 在 $\Spec(A)$ 中稠密，当且仅当它包含所有极小素理想。
由于 $A$ 的伴随素理想之集等于极小素理想之集与嵌入伴随素理想
之集的并（见《代数》命题
\ref{algebra-proposition-minimal-primes-associated-primes}），
可知 (2) 与 (3) 等价。

\medskip\noindent
假设 (3)。这表示每个伴随素理想 $\mathfrak p$，如果它是 $A$
的伴随素理想，便对应于某个 $A_{f_i}$ 的一个素理想，其中 $i$
可适当选取。因此
$A \to A_{f_1} \times \ldots \times A_{f_n}$ 是单射，因为由
《代数》引理 \ref{algebra-lemma-zero-at-ass-zero}，映射
$A \to \prod_{\mathfrak p \in \text{Ass}(A)} A_\mathfrak p$ 是单射。
由此可见 (3) 蕴含 (1)。
\end{proof}

\begin{lemma}
\label{divisors-lemma-remove-embedded-points}
设 $X$ 为局部 Noether 概形。
设 $\mathcal{F}$ 为 $X$ 上的凝聚层。满足下述条件的凝聚子层之集
$$
\{
\mathcal{K} \subset \mathcal{F}
\mid
\text{Supp}(\mathcal{K})\text{ is nowhere dense in }\text{Supp}(\mathcal{F})
\}
$$
有一个极大元 $\mathcal{K}$。令
$\mathcal{F}' = \mathcal{F}/\mathcal{K}$，则有
\begin{enumerate}
\item $\text{Supp}(\mathcal{F}') = \text{Supp}(\mathcal{F})$；
\item $\mathcal{F}'$ 没有嵌入伴随点；
\item 存在稠密开集 $U \subset X$，使得
$U \cap \text{Supp}(\mathcal{F})$ 在 $\text{Supp}(\mathcal{F})$
中稠密，且 $\mathcal{F}'|_U \cong \mathcal{F}|_U$。
\end{enumerate}
\end{lemma}

\begin{proof}
这由《代数》引理 \ref{algebra-lemma-remove-embedded-primes} 和
\ref{algebra-lemma-remove-embedded-primes-localize} 得出。
注意，$U$ 可取为 $\mathcal{F}$ 的嵌入伴随点之集的闭包的补集。
\end{proof}

\begin{lemma}
\label{divisors-lemma-no-embedded-points-endos}
设 $X$ 为局部 Noether 概形。
设 $\mathcal{F}$ 为没有嵌入伴随点的凝聚 $\mathcal{O}_X$-模。置
$$
\mathcal{I}
=
\Ker(\mathcal{O}_X
\longrightarrow
\SheafHom_{\mathcal{O}_X}(\mathcal{F}, \mathcal{F})).
$$
这是一个凝聚理想层，并定义了没有嵌入点的闭子概形
$Z \subset X$。此外，存在凝聚层 $\mathcal{G}$，它定义在 $Z$ 上，并且
(a) $\mathcal{F} = (Z \to X)_*\mathcal{G}$，
(b) $\mathcal{G}$ 没有嵌入伴随点，且
(c) $\text{Supp}(\mathcal{G}) = Z$（作为集合）。
\end{lemma}

\begin{proof}
其中一些断言已在《概形上同调》引理
\ref{coherent-lemma-coherent-support-closed} 的证明中见过，
其余断言由《代数》引理
\ref{algebra-lemma-no-embedded-primes-endos} 得出。
\end{proof}



\section{弱伴随点}
\label{divisors-section-weakly-associated}

\noindent
设 $R$ 为环，$M$ 为 $R$-模。
回忆，称素理想 $\mathfrak p \subset R$ 与 $M$ {\it 弱伴随}，
如果存在一个元素 $m$，它属于 $M$，并且 $\mathfrak p$
在包含 $m$ 的零化子的素理想中极小。参见《代数》定义
\ref{algebra-definition-weakly-associated}。
若 $R$ 是极大理想为 $\mathfrak m$ 的局部环，则
$\mathfrak m$ 与 $M$ 弱伴随，当且仅当存在元素 $m \in M$，
其零化子的根理想为 $\mathfrak m$，见《代数》引理
\ref{algebra-lemma-weakly-ass-local}。

\begin{definition}
\label{divisors-definition-weakly-associated}
设 $X$ 为概形。
设 $\mathcal{F}$ 为 $X$ 上的拟凝聚层。
\begin{enumerate}
\item 称 $x \in X$ 与 $\mathcal{F}$ {\it 弱伴随}，如果极大理想
$\mathfrak m_x$ 与 $\mathcal{O}_{X, x}$-模 $\mathcal{F}_x$ 弱伴随。
\item 用 $\text{WeakAss}(\mathcal{F})$ 表示 $\mathcal{F}$ 的
弱伴随点之集。
\item $X$ 的{\it 弱伴随点}是 $\mathcal{O}_X$ 的弱伴随点。
\end{enumerate}
\end{definition}

\noindent
在这种情形下，在任意仿射开集上，这一概念恰好对应于上面定义的
弱伴随素理想。精确表述如下。

\begin{lemma}
\label{divisors-lemma-weakly-associated-affine-open}
设 $X$ 为概形，$\mathcal{F}$ 为 $X$ 上的拟凝聚层。
设 $\Spec(A) = U \subset X$ 为仿射开集，并置
$M = \Gamma(U, \mathcal{F})$。
设 $x \in U$，并设 $\mathfrak p \subset A$ 为对应的素理想。
以下条件等价：
\begin{enumerate}
\item $\mathfrak p$ 与 $M$ 弱伴随；
\item $x$ 与 $\mathcal{F}$ 弱伴随。
\end{enumerate}
\end{lemma}

\begin{proof}
这由《代数》引理 \ref{algebra-lemma-weakly-ass-local} 得出。
\end{proof}

\begin{lemma}
\label{divisors-lemma-weakly-ass-support}
设 $X$ 为概形。
设 $\mathcal{F}$ 为拟凝聚 $\mathcal{O}_X$-模。则
$$
\text{Ass}(\mathcal{F}) \subset \text{WeakAss}(\mathcal{F}) \subset
\text{Supp}(\mathcal{F}).
$$
\end{lemma}

\begin{proof}
这由定义立即得到。
\end{proof}

\begin{lemma}
\label{divisors-lemma-ses-weakly-ass}
设 $X$ 为概形。设
$0 \to \mathcal{F}_1 \to \mathcal{F}_2 \to \mathcal{F}_3 \to 0$
为 $X$ 上拟凝聚层的短正合列。则
$\text{WeakAss}(\mathcal{F}_2) \subset
\text{WeakAss}(\mathcal{F}_1) \cup \text{WeakAss}(\mathcal{F}_3)$，且
$\text{WeakAss}(\mathcal{F}_1) \subset \text{WeakAss}(\mathcal{F}_2)$。
\end{lemma}

\begin{proof}
对每个点 $x \in X$，茎的序列
$0 \to \mathcal{F}_{1, x} \to \mathcal{F}_{2, x} \to \mathcal{F}_{3, x} \to 0$
是 $\mathcal{O}_{X, x}$-模的短正合列。因此，本引理由《代数》引理
\ref{algebra-lemma-weakly-ass} 得出。
\end{proof}

\begin{lemma}
\label{divisors-lemma-weakly-ass-zero}
设 $X$ 为概形。
设 $\mathcal{F}$ 为拟凝聚 $\mathcal{O}_X$-模。则
$$
\mathcal{F} = (0) \Leftrightarrow \text{WeakAss}(\mathcal{F}) = \emptyset
$$
\end{lemma}

\begin{proof}
由引理 \ref{divisors-lemma-weakly-associated-affine-open} 和《代数》引理
\ref{algebra-lemma-weakly-ass-zero} 得出。
\end{proof}

\begin{lemma}
\label{divisors-lemma-restriction-injective-open-contains-weakly-ass}
设 $X$ 为概形，$\mathcal{F}$ 为拟凝聚 $\mathcal{O}_X$-模。
若 $U \subset X$ 是开集，且
$\text{WeakAss}(\mathcal{F}) \subset U$，则映射
$\Gamma(X, \mathcal{F}) \to \Gamma(U, \mathcal{F})$ 是单射。
\end{lemma}

\begin{proof}
设 $s \in \Gamma(X, \mathcal{F})$ 为一个限制到 $U$ 上为零的截面。
令 $\mathcal{F}' \subset \mathcal{F}$ 为映射
$\mathcal{O}_X \to \mathcal{F}$ 的像，该映射由 $s$ 定义。于是
$\text{Supp}(\mathcal{F}') \cap U = \emptyset$。另一方面，由引理
\ref{divisors-lemma-ses-weakly-ass}，有
$\text{WeakAss}(\mathcal{F}') \subset \text{WeakAss}(\mathcal{F})$。
又有 $\text{WeakAss}(\mathcal{F}') \subset \text{Supp}(\mathcal{F}')$
（引理 \ref{divisors-lemma-weakly-ass-support}），故
$\text{WeakAss}(\mathcal{F}') = \emptyset$。
再由引理 \ref{divisors-lemma-weakly-ass-zero} 得 $\mathcal{F}' = 0$。
\end{proof}

\begin{lemma}
\label{divisors-lemma-minimal-support-in-weakly-ass}
设 $X$ 为概形。
设 $\mathcal{F}$ 为拟凝聚 $\mathcal{O}_X$-模。
设 $x \in \text{Supp}(\mathcal{F})$ 是 $\mathcal{F}$ 的支撑中的一点，
并且它不是 $\text{Supp}(\mathcal{F})$ 中另一个点的特化。则
$x \in \text{WeakAss}(\mathcal{F})$。
特别地，$X$ 的任一不可约分支的泛点都与 $\mathcal{O}_X$ 弱伴随。
\end{lemma}

\begin{proof}
由于 $x \in \text{Supp}(\mathcal{F})$，模 $\mathcal{F}_x$ 非零。
因此，由《代数》引理 \ref{algebra-lemma-weakly-ass-zero}，集合
$\text{WeakAss}(\mathcal{F}_x) \subset \Spec(\mathcal{O}_{X, x})$
非空。另一方面，根据假设，
$\text{Supp}(\mathcal{F}_x) = \{\mathfrak m_x\}$。
又因为
$\text{WeakAss}(\mathcal{F}_x) \subset \text{Supp}(\mathcal{F}_x)$
（《代数》引理 \ref{algebra-lemma-weakly-ass-support}），可见
$\mathfrak m_x$ 与 $\mathcal{F}_x$ 弱伴随，结论成立。
\end{proof}

\begin{lemma}
\label{divisors-lemma-ass-weakly-ass}
设 $X$ 为概形。
设 $\mathcal{F}$ 为拟凝聚 $\mathcal{O}_X$-模。
若 $\mathfrak m_x$ 是 $\mathcal{O}_{X, x}$ 的有限生成理想，则
$$
x \in \text{Ass}(\mathcal{F}) \Leftrightarrow
x \in \text{WeakAss}(\mathcal{F}).
$$
特别地，若 $X$ 局部 Noether，则
$\text{Ass}(\mathcal{F}) = \text{WeakAss}(\mathcal{F})$。
\end{lemma}

\begin{proof}
见《代数》引理 \ref{algebra-lemma-ass-weakly-ass}。
\end{proof}

\begin{lemma}
\label{divisors-lemma-weakass-pushforward}
设 $f : X \to S$ 为拟紧拟分离的概形态射。
设 $\mathcal{F}$ 为拟凝聚 $\mathcal{O}_X$-模。
设 $s \in S$ 是一个不属于 $f$ 的像的点。则
$s$ 不与 $f_*\mathcal{F}$ 弱伴随。
\end{lemma}

\begin{proof}
考虑基变换 $f' : X' \to \Spec(\mathcal{O}_{S, s})$；它由 $f$
沿态射 $g : \Spec(\mathcal{O}_{S, s}) \to S$ 得到。记另一个投影为
$g' : X' \to X$。于是
$$
(f_*\mathcal{F})_s = (g^*f_*\mathcal{F})_s = (f'_*(g')^*\mathcal{F})_s
$$
第一个等式成立，是因为 $g$ 在 $s$ 处的局部环上诱导同构；
第二个等式由平坦基变换得到（《概形上同调》引理
\ref{coherent-lemma-flat-base-change-cohomology}）。当然，
$s \in \Spec(\mathcal{O}_{S, s})$ 不属于 $f'$ 的像。
因此可假设 $S$ 是局部环 $(A, \mathfrak m)$ 的谱，并且 $s$
对应于 $\mathfrak m$。由《概形》引理
\ref{schemes-lemma-push-forward-quasi-coherent}，层
$f_*\mathcal{F}$ 拟凝聚，设它对应于 $A$-模 $M$。
由于 $s$ 不属于 $f$ 的像，可见
$X = \bigcup_{a \in \mathfrak m} f^{-1}D(a)$ 是一个开覆盖。
因为 $X$ 拟紧，可找到 $a_1, \ldots, a_n \in \mathfrak m$，使得
$X = f^{-1}D(a_1) \cup \ldots \cup f^{-1}D(a_n)$。由此
$$
M \to M_{a_1} \oplus \ldots \oplus M_{a_r}
$$
是单射。因此，对任意非零元素 $m$，若把它视作茎 $M_\mathfrak p$ 的元素，
存在一个 $i$，使得 $a_i^n m$ 对所有 $n \geq 0$ 都非零。
故 $\mathfrak m$ 不与 $M$ 弱伴随。
\end{proof}

\begin{lemma}
\label{divisors-lemma-check-injective-on-weakass}
设 $X$ 为概形。设 $\varphi : \mathcal{F} \to \mathcal{G}$ 为拟凝聚
$\mathcal{O}_X$-模之间的映射。假设对每个 $x \in X$，
以下至少一种情形成立：
\begin{enumerate}
\item $\mathcal{F}_x \to \mathcal{G}_x$ 是单射；
\item $x \not \in \text{WeakAss}(\mathcal{F})$。
\end{enumerate}
则 $\varphi$ 是单射。
\end{lemma}

\begin{proof}
这些假设蕴含 $\text{WeakAss}(\Ker(\varphi)) = \emptyset$，
故由引理 \ref{divisors-lemma-weakly-ass-zero} 有 $\Ker(\varphi) = 0$。
\end{proof}

\begin{lemma}
\label{divisors-lemma-depth-2-hartog}
设 $X$ 为局部 Noether 概形，$\mathcal{F}$ 为凝聚
$\mathcal{O}_X$-模。设 $j : U \to X$ 为开子概形，并且对
$x \in X \setminus U$ 有 $\text{depth}(\mathcal{F}_x) \geq 2$。则
$$
\mathcal{F} \longrightarrow j_*(\mathcal{F}|_U)
$$
是同构，从而
$\Gamma(X, \mathcal{F}) \to \Gamma(U, \mathcal{F})$ 也是同构。
\end{lemma}

\begin{proof}
我们断言，引理 \ref{divisors-lemma-check-isomorphism-via-depth-and-ass}
可用于本引理中展示的映射。
设 $x \in X$。若 $x \in U$，则该映射在茎上是同构，因为
$j_*(\mathcal{F}|_U)|_U = \mathcal{F}|_U$。
若 $x \in X \setminus U$，则
$x \not \in \text{Ass}(j_*(\mathcal{F}|_U))$
（引理 \ref{divisors-lemma-weakass-pushforward} 和
\ref{divisors-lemma-weakly-ass-support}）。再结合假设
$\text{depth}(\mathcal{F}_x) \geq 2$，证明即告完成。
\end{proof}

\begin{lemma}
\label{divisors-lemma-weakass-reduced}
设 $X$ 为约化概形。则 $X$ 的弱伴随点恰好是 $X$ 的
不可约分支的泛点。
\end{lemma}

\begin{proof}
由《代数》引理 \ref{algebra-lemma-reduced-weakly-ass-minimal} 得出。
\end{proof}



\section{态射与弱伴随点}
\label{divisors-section-morphisms-weakly-associated}

\begin{lemma}
\label{divisors-lemma-weakly-ass-reverse-functorial}
设 $f : X \to S$ 为概形的仿射态射。
设 $\mathcal{F}$ 为拟凝聚 $\mathcal{O}_X$-模。则
$$
\text{WeakAss}_S(f_*\mathcal{F}) \subset f(\text{WeakAss}_X(\mathcal{F}))
$$
\end{lemma}

\begin{proof}
可以假设 $X$ 和 $S$ 仿射，故 $X \to S$ 来自环映射
$A \to B$。于是 $\mathcal{F} = \widetilde M$，这里所用的是某个
$B$-模 $M$。由引理 \ref{divisors-lemma-weakly-associated-affine-open}，
$\mathcal{F}$ 的弱伴随点恰好对应于 $M$ 的弱伴随素理想。
类似地，$f_*\mathcal{F}$ 的弱伴随点恰好对应于把 $M$ 看作
$A$-模时的弱伴随素理想。因此，本引理由《代数》引理
\ref{algebra-lemma-weakly-ass-reverse-functorial} 得出。
\end{proof}

\begin{lemma}
\label{divisors-lemma-ass-functorial-equal}
设 $f : X \to S$ 为概形的仿射态射。
设 $\mathcal{F}$ 为拟凝聚 $\mathcal{O}_X$-模。
若 $X$ 局部 Noether，则
$$
f(\text{Ass}_X(\mathcal{F})) =
\text{Ass}_S(f_*\mathcal{F}) =
\text{WeakAss}_S(f_*\mathcal{F}) =
f(\text{WeakAss}_X(\mathcal{F}))
$$
\end{lemma}

\begin{proof}
可以假设 $X$ 和 $S$ 仿射，故 $X \to S$ 来自环映射
$A \to B$。由于 $X$ 局部 Noether，环 $B$ 是 Noether 环，见
《性质》引理 \ref{properties-lemma-locally-Noetherian}。
写 $\mathcal{F} = \widetilde M$，这里所用的是某个 $B$-模 $M$。
由引理 \ref{divisors-lemma-associated-affine-open}，$\mathcal{F}$ 的伴随点
恰好对应于 $M$ 的伴随素理想；而把 $M$ 看作 $A$-模时的任一
伴随素理想，都是 $f_*\mathcal{F}$ 的一个伴随点。因此，由《代数》引理
\ref{algebra-lemma-ass-functorial-Noetherian} 得到包含关系
$$
f(\text{Ass}_X(\mathcal{F})) \subset \text{Ass}_S(f_*\mathcal{F})
$$
由引理 \ref{divisors-lemma-weakly-ass-support}，有包含关系
$$
\text{Ass}_S(f_*\mathcal{F}) \subset \text{WeakAss}_S(f_*\mathcal{F})
$$
由引理 \ref{divisors-lemma-weakly-ass-reverse-functorial}，有包含关系
$$
\text{WeakAss}_S(f_*\mathcal{F}) \subset f(\text{WeakAss}_X(\mathcal{F}))
$$
由引理 \ref{divisors-lemma-ass-weakly-ass}，最外侧的两个集合相等，
故处处均为等号。
\end{proof}

\begin{lemma}
\label{divisors-lemma-weakly-associated-finite}
设 $f : X \to S$ 为概形的有限态射。
设 $\mathcal{F}$ 为拟凝聚 $\mathcal{O}_X$-模。则
$\text{WeakAss}(f_*\mathcal{F}) = f(\text{WeakAss}(\mathcal{F}))$。
\end{lemma}

\begin{proof}
可以假设 $X$ 和 $S$ 仿射，故 $X \to S$ 来自有限环映射
$A \to B$。写 $\mathcal{F} = \widetilde M$，这里所用的是某个
$B$-模 $M$。由引理 \ref{divisors-lemma-weakly-associated-affine-open}，
$\mathcal{F}$ 的弱伴随点恰好对应于 $M$ 的弱伴随素理想。
类似地，$f_*\mathcal{F}$ 的弱伴随点恰好对应于把 $M$ 看作
$A$-模时的弱伴随素理想。因此，本引理由《代数》引理
\ref{algebra-lemma-weakly-ass-finite-ring-map} 得出。
\end{proof}

\begin{lemma}
\label{divisors-lemma-weakly-ass-pullback}
设 $f : X \to S$ 为概形的态射，$\mathcal{G}$ 为拟凝聚
$\mathcal{O}_S$-模。设 $x \in X$ 且 $s = f(x)$。
若 $f$ 在 $x$ 处平坦，点 $x$ 是纤维 $X_s$ 的泛点，并且
$s \in \text{WeakAss}_S(\mathcal{G})$，则
$x \in \text{WeakAss}(f^*\mathcal{G})$。
\end{lemma}

\begin{proof}
置 $A = \mathcal{O}_{S, s}$、$B = \mathcal{O}_{X, x}$ 和
$M = \mathcal{G}_s$。设 $m \in M$ 为一个元素，其零化子
$I = \{a \in A \mid am = 0\}$ 的根理想为 $\mathfrak m_A$。
$m \otimes 1$ 的零化子为 $I B$，因为 $A \to B$ 忠实平坦。
因此，只须证明 $\sqrt{I B} = \mathfrak m_B$。这由以下两点得出：
$B/\mathfrak m_AB$ 的极大理想局部幂零（见《代数》引理
\ref{algebra-lemma-minimal-prime-reduced-ring}），以及假设
$\sqrt{I} = \mathfrak m_A$。略去一些细节。
\end{proof}

\begin{lemma}
\label{divisors-lemma-weakly-ass-change-fields}
设 $K/k$ 为域扩张，$X$ 为 $k$ 上的概形。
设 $\mathcal{F}$ 为拟凝聚 $\mathcal{O}_X$-模。
设 $y \in X_K$ 的像为 $x \in X$。若 $y$ 是拉回
$\mathcal{F}_K$ 的弱伴随点，则 $x$ 是 $\mathcal{F}$ 的弱伴随点。
\end{lemma}

\begin{proof}
这只是把《代数》引理 \ref{algebra-lemma-weakly-ass-change-fields}
改写成概形语言。
\end{proof}

\noindent
下面是一个简单引理，它表明直像常常具有至少为 2 的深度。

\begin{lemma}
\label{divisors-lemma-depth-pushforward}
设 $f : X \to S$ 为拟紧拟分离的概形态射。
设 $\mathcal{F}$ 为拟凝聚 $\mathcal{O}_X$-模。
设 $s \in S$。
\begin{enumerate}
\item 若 $s \not \in f(X)$，则 $s$ 不与 $f_*\mathcal{F}$ 弱伴随。
\item 若 $s \not \in f(X)$ 且 $\mathcal{O}_{S, s}$ 是 Noether 环，
则 $s$ 不是 $f_*\mathcal{F}$ 的伴随点。
\item 若 $s \not \in f(X)$，$(f_*\mathcal{F})_s$ 是有限
$\mathcal{O}_{S, s}$-模，并且 $\mathcal{O}_{S, s}$ 是 Noether 环，
则 $\text{depth}((f_*\mathcal{F})_s) \geq 2$。
\item 若 $\mathcal{F}$ 在 $S$ 上平坦，并且 $a \in \mathfrak m_s$
是非零因子，则 $a$ 在 $(f_*\mathcal{F})_s$ 上是非零因子。
\item 若 $\mathcal{F}$ 在 $S$ 上平坦，并且
$a, b \in \mathfrak m_s$ 是正则序列，则 $a$ 在
$(f_*\mathcal{F})_s$ 上是非零因子，且 $b$ 在
$(f_*\mathcal{F})_s/a(f_*\mathcal{F})_s$ 上是非零因子。
\item 若 $\mathcal{F}$ 在 $S$ 上平坦，并且 $(f_*\mathcal{F})_s$
是有限 $\mathcal{O}_{S, s}$-模，则
$\text{depth}((f_*\mathcal{F})_s) \geq
\min(2, \text{depth}(\mathcal{O}_{S, s}))$。
\end{enumerate}
\end{lemma}

\begin{proof}
(1) 是引理 \ref{divisors-lemma-weakass-pushforward}。
(2) 由 (1) 和引理 \ref{divisors-lemma-ass-weakly-ass} 得出。

{\small
\medskip\noindent
证明 (3)。要证明深度 $\geq 2$，只须证明
$\Hom_{\mathcal{O}_{S, s}}(\kappa(s), (f_*\mathcal{F})_s) = 0$ 及
$\Ext^1_{\mathcal{O}_{S, s}}(\kappa(s), (f_*\mathcal{F})_s) = 0$，
见《代数》引理 \ref{algebra-lemma-depth-ext}。\par}
利用正合列
$0 \to \mathfrak m_s \to \mathcal{O}_{S, s} \to \kappa(s) \to 0$，
只须证明映射
$$
\Hom_{\mathcal{O}_{S, s}}(\mathcal{O}_{S, s}, (f_*\mathcal{F})_s)
\to
\Hom_{\mathcal{O}_{S, s}}(\mathfrak m_s, (f_*\mathcal{F})_s)
$$
是同构。由平坦基变换（《概形上同调》引理
\ref{coherent-lemma-flat-base-change-cohomology}），可以把 $S$ 替换为
$\Spec(\mathcal{O}_{S, s})$，并把 $X$ 替换为
$\Spec(\mathcal{O}_{S, s}) \times_S X$。
用 $\mathfrak m \subset \mathcal{O}_S$ 表示 $s$ 的理想层。于是
$$
\Hom_{\mathcal{O}_{S, s}}(\mathfrak m_s, (f_*\mathcal{F})_s) =
\Hom_{\mathcal{O}_S}(\mathfrak m, f_*\mathcal{F}) =
\Hom_{\mathcal{O}_X}(f^*\mathfrak m, \mathcal{F})
$$
第一个等式成立，是因为 $S$ 是以 $s$ 为闭点的局部概形；
第二个等式由拟凝聚模上 $f^*, f_*$ 的伴随性得到。然而，由于
$s \not \in f(X)$，可见 $f^*\mathfrak m = \mathcal{O}_X$。
沿上述论证逆推，便得到所需等式。

\medskip\noindent
为证明 (4)、(5) 和 (6)，利用平坦基变换（《概形上同调》引理
\ref{coherent-lemma-flat-base-change-cohomology}），可约化到 $S$ 是
$\mathcal{O}_{S, s}$ 的谱的情形。此时非零因子
$a \in \mathcal{O}_{S, s}$ 确定短正合列
$$
0 \to \mathcal{O}_S \xrightarrow{a} \mathcal{O}_S \to
\mathcal{O}_S/a \mathcal{O}_S \to 0
$$
由于 $\mathcal{F}$ 在 $S$ 上平坦，得到正合列
$$
0 \to \mathcal{F} \xrightarrow{a} \mathcal{F} \to
\mathcal{F}/a\mathcal{F} \to 0
$$
取直像，得到正合列
$$
0 \to f_*\mathcal{F} \xrightarrow{a} f_*\mathcal{F} \to
f_*(\mathcal{F}/a\mathcal{F})
$$
这证明了 (4)，并且说明
$f_*\mathcal{F}/ af_*\mathcal{F} \subset f_*(\mathcal{F}/a\mathcal{F})$。
若 $b$ 在 $\mathcal{O}_{S, s}/a\mathcal{O}_{S, s}$ 上是非零因子，
则完全相同的论证说明
$b : \mathcal{F}/a\mathcal{F} \to \mathcal{F}/a\mathcal{F}$ 是单射。
取直像可知
$$
b : f_*(\mathcal{F}/a\mathcal{F}) \to f_*(\mathcal{F}/a\mathcal{F})
$$
是单射，因而
$b : f_*\mathcal{F}/ af_*\mathcal{F} \to f_*\mathcal{F}/ af_*\mathcal{F}$
也为单射。这证明了 (5)。(6) 由 (4)、(5) 及各定义得出。
\end{proof}


\section{相对伴随点集}
\label{divisors-section-relative-assassin}

\noindent
设 $A \to B$ 为环映射，$N$ 为 $B$-模。回忆，若素理想
$\mathfrak q \subset B$ 满足如下条件，则称它属于 $N$ 在
$B/A$ 上的相对伴随素理想集：$\mathfrak q$ 是
$N \otimes_A \kappa(\mathfrak p)$ 的伴随素理想。这里
$\mathfrak p = A \cap \mathfrak q$。见《代数》定义
\ref{algebra-definition-relative-assassin}。
下面定义概形态射上拟凝聚层的相对伴随点集。

\begin{definition}
\label{divisors-definition-relative-assassin}
设 $f : X \to S$ 为概形的态射。
设 $\mathcal{F}$ 为拟凝聚 $\mathcal{O}_X$-模。
$\mathcal{F}$ 在 $X$ 中、$S$ 上的{\it 相对伴随点集}是集合
$$
\text{Ass}_{X/S}(\mathcal{F}) =
\bigcup\nolimits_{s \in S} \text{Ass}_{X_s}(\mathcal{F}_s)
$$
其中 $\mathcal{F}_s = (X_s \to X)^*\mathcal{F}$ 是 $\mathcal{F}$
在 $f$ 于 $s$ 处的纤维上的限制。
\end{definition}

\noindent
这里仍有一个注意事项：最好只在 $f$ 的纤维局部 Noether 且
$\mathcal{F}$ 为有限型时使用这一概念。一般情形下，或许应当使用
相对弱伴随点集（定义见下一节）。下面把概形论概念与仿射开集上的
代数概念联系起来；注意，这一对应只有在概形态射的纤维局部
Noether 时才完全成立。

\begin{lemma}
\label{divisors-lemma-relative-assassin-affine-open}
设 $f : X \to S$ 为概形的态射。
设 $\mathcal{F}$ 为 $X$ 上的拟凝聚层。
设 $U \subset X$ 和 $V \subset S$ 为仿射开集，且
$f(U) \subset V$。写 $U = \Spec(A)$、$V = \Spec(R)$，并置
$M = \Gamma(U, \mathcal{F})$。
设 $x \in U$，并设 $\mathfrak p \subset A$ 为对应的素理想。则
$$
\mathfrak p \in \text{Ass}_{A/R}(M) \Rightarrow
x \in \text{Ass}_{X/S}(\mathcal{F})
$$
若所有纤维 $X_s$（即 $f$ 的纤维）都局部 Noether，则等价式
$\mathfrak p \in \text{Ass}_{A/R}(M) \Leftrightarrow
x \in \text{Ass}_{X/S}(\mathcal{F})$
对上述每一对 $(\mathfrak p, x)$ 都成立。
\end{lemma}

\begin{proof}
集合 $\text{Ass}_{A/R}(M)$ 定义于《代数》定义
\ref{algebra-definition-relative-assassin}。选取一对
$(\mathfrak p, x)$，并令 $s = f(x)$。
设 $\mathfrak r \subset R$ 为位于 $\mathfrak p$ 之下的素理想，
即对应于 $s$ 的素理想。设
$\mathfrak p' \subset A \otimes_R \kappa(\mathfrak r)$ 为逆像等于
$\mathfrak p$ 的素理想；换言之，它对应于把 $x$ 看作其纤维
$X_s$ 的点时的那个点。于是，
$\mathfrak p \in \text{Ass}_{A/R}(M)$ 当且仅当 $\mathfrak p'$ 是
$M \otimes_R \kappa(\mathfrak r)$ 的伴随素理想，见《代数》引理
\ref{algebra-lemma-compare-relative-assassins}。
注意，环 $A \otimes_R \kappa(\mathfrak r)$ 对应于 $U_s$，而模
$M \otimes_R \kappa(\mathfrak r)$ 对应于拟凝聚层
$\mathcal{F}_s|_{U_s}$。因此，由引理
\ref{divisors-lemma-associated-affine-open}，$x$ 是 $\mathcal{F}_s$ 的伴随点。
若 $\mathfrak p'$ 有限生成，则逆命题也成立；这便说明了最后一句。
\end{proof}

\begin{lemma}
\label{divisors-lemma-base-change-relative-assassin}
设 $f : X \to S$ 为概形的态射。
设 $\mathcal{F}$ 为拟凝聚 $\mathcal{O}_X$-模。
设 $g : S' \to S$ 为概形的态射。考虑基变换图
$$
\xymatrix{
X' \ar[d] \ar[r]_{g'} & X \ar[d] \\
S' \ar[r]^g & S
}
$$
并置 $\mathcal{F}' = (g')^*\mathcal{F}$。设点 $x' \in X'$ 的像
分别为 $x \in X$、$s' \in S'$ 和 $s \in S$。假设 $f$ 局部有限型。
则 $x' \in \text{Ass}_{X'/S'}(\mathcal{F}')$ 当且仅当
$x \in \text{Ass}_{X/S}(\mathcal{F})$，并且 $x'$ 对应于
$\Spec(\kappa(s') \otimes_{\kappa(s)} \kappa(x))$ 的一个不可约分支的泛点。
\end{lemma}

\begin{proof}
考虑纤维态射 $X'_{s'} \to X_s$。由于
$X_{s'} = X_s \times_{\Spec(\kappa(s))} \Spec(\kappa(s'))$，
这是平坦态射。此外，$\mathcal{F}'_{s'}$ 是 $\mathcal{F}_s$ 沿此态射的拉回。
由于 $X_s$ 在 Noether 概形 $\Spec(\kappa(s))$ 上局部有限型，
故 $X_s$ 局部 Noether，见《态射》引理
\ref{morphisms-lemma-finite-type-noetherian}。因此可应用引理
\ref{divisors-lemma-bourbaki}，得到
$$
\text{Ass}_{X'_{s'}}(\mathcal{F}'_{s'}) =
\bigcup\nolimits_{x \in \text{Ass}(\mathcal{F}_s)} \text{Ass}((X'_{s'})_x).
$$
于是，为证明本引理，只须说明纤维 $(X'_{s'})_x$ 的伴随点就是其泛点；
这里该纤维是态射 $X'_{s'} \to X_s$ 在 $x$ 上的纤维。注意，作为概形有
$(X'_{s'})_x = \Spec(\kappa(s') \otimes_{\kappa(s)} \kappa(x))$。
由《代数》引理 \ref{algebra-lemma-tensor-fields-CM}，环
$\kappa(s') \otimes_{\kappa(s)} \kappa(x)$ 是 Noether
Cohen-Macaulay 环。因此，它的伴随素理想就是其极小素理想；见
《代数》命题 \ref{algebra-proposition-minimal-primes-associated-primes}
（极小素理想是伴随素理想）和《代数》引理
\ref{algebra-lemma-criterion-no-embedded-primes}（没有嵌入素理想）。
\end{proof}

\begin{remark}
\label{divisors-remark-base-change-relative-assassin}
采用引理 \ref{divisors-lemma-base-change-relative-assassin} 的记号和假设，
总有
$(g')^{-1}(\text{Ass}_{X/S}(\mathcal{F})) \supset
\text{Ass}_{X'/S'}(\mathcal{F}')$。
若态射 $S' \to S$ 局部拟有限，则事实上有
$$
(g')^{-1}(\text{Ass}_{X/S}(\mathcal{F}))
=
\text{Ass}_{X'/S'}(\mathcal{F}')
$$
因为此时域扩张 $\kappa(s')/\kappa(s)$ 总是有限的。事实上，
对任意满足下述条件的态射 $g : S' \to S$，这一结论仍成立：
所有域扩张 $\kappa(s')/\kappa(s)$ 都是代数扩张。这是因为此时
$\kappa(s') \otimes_{\kappa(s)} \kappa(x)$ 的所有素理想都是
极大（且极小）素理想，见《代数》引理
\ref{algebra-lemma-integral-over-field}。
\end{remark}




\section{相对弱伴随点集}
\label{divisors-section-relative-weak-assassin}

\begin{definition}
\label{divisors-definition-relative-weak-assassin}
设 $f : X \to S$ 为概形的态射。
设 $\mathcal{F}$ 为拟凝聚 $\mathcal{O}_X$-模。
$\mathcal{F}$ 在 $X$ 中、$S$ 上的{\it 相对弱伴随点集}是集合
$$
\text{WeakAss}_{X/S}(\mathcal{F}) =
\bigcup\nolimits_{s \in S} \text{WeakAss}(\mathcal{F}_s)
$$
其中 $\mathcal{F}_s = (X_s \to X)^*\mathcal{F}$ 是 $\mathcal{F}$
在 $f$ 于 $s$ 处的纤维上的限制。
\end{definition}

\begin{lemma}
\label{divisors-lemma-relative-weak-assassin-assassin-finite-type}
设 $f : X \to S$ 为局部有限型的概形态射。
设 $\mathcal{F}$ 为拟凝聚 $\mathcal{O}_X$-模。则
$\text{WeakAss}_{X/S}(\mathcal{F}) = \text{Ass}_{X/S}(\mathcal{F})$。
\end{lemma}

\begin{proof}
这是因为 $f$ 的纤维都是局部 Noether 概形，并且在局部 Noether
概形上，伴随点与弱伴随点一致，见引理
\ref{divisors-lemma-ass-weakly-ass}。
\end{proof}

\begin{lemma}
\label{divisors-lemma-relative-weak-assassin-finite}
设 $f : X \to S$ 为概形的态射。
设 $i : Z \to X$ 为有限态射。
设 $\mathcal{F}$ 为拟凝聚 $\mathcal{O}_Z$-模。则
$\text{WeakAss}_{X/S}(i_*\mathcal{F}) =
i(\text{WeakAss}_{Z/S}(\mathcal{F}))$。
\end{lemma}

\begin{proof}
设 $i_s : Z_s \to X_s$ 为纤维间的诱导态射。于是
$(i_*\mathcal{F})_s = i_{s, *}(\mathcal{F}_s)$；这是《概形上同调》引理
\ref{coherent-lemma-affine-base-change} 与 $i$ 是仿射态射这一事实的结果。
因此可应用引理
\ref{divisors-lemma-weakly-associated-finite} 得出结论。
\end{proof}





\section{Fitting 理想}
\label{divisors-section-fitting-ideals}

\noindent
本节继续《代数进阶》第
\ref{more-algebra-section-fitting-ideals} 节中的讨论。
设 $S$ 为概形，$\mathcal{F}$ 为有限型拟凝聚 $\mathcal{O}_S$-模。
在这种情形下，可以构造 Fitting 理想
$$
0 = \text{Fit}_{-1}(\mathcal{F}) \subset \text{Fit}_0(\mathcal{F}) \subset
\text{Fit}_1(\mathcal{F}) \subset \ldots \subset \mathcal{O}_S
$$
它们是由下述性质刻画的拟凝聚理想序列：对每个仿射开集
$U = \Spec(A)$，若它是 $S$ 的开集且 $\mathcal{F}|_U$ 对应于 $A$-模 $M$，则
$\text{Fit}_i(\mathcal{F})|_U$ 对应于理想
$\text{Fit}_i(M) \subset A$。这一定义良好，并给出拟凝聚理想层，
因为若 $f \in A$，则第 $i$ 个 Fitting 理想（即 $M_f$ 在 $A_f$
上的相应理想）等于 $\text{Fit}_i(M) A_f$，见《代数进阶》引理
\ref{more-algebra-lemma-fitting-ideal-basics}。

\medskip\noindent
也可以用 $\mathcal{F}$ 的局部展示来构造 Fitting 理想。具体地，
若 $U \subset X$ 为开集，并且
$$
\bigoplus\nolimits_{i \in I} \mathcal{O}_U \to
\mathcal{O}_U^{\oplus n} \to \mathcal{F}|_U \to 0
$$
是 $\mathcal{F}$ 在 $U$ 上的一个展示，则
$\text{Fit}_r(\mathcal{F})|_U$ 由定义该展示中第一个箭头的矩阵的
所有 $(n - r) \times (n - r)$ 子式生成。这与上述构造相容，
因为环上模的 Fitting 理想本来就是这样定义的。略去一些细节。

\begin{lemma}
\label{divisors-lemma-base-change-fitting-ideal}
设 $f : T \to S$ 为概形的态射。
设 $\mathcal{F}$ 为有限型拟凝聚 $\mathcal{O}_S$-模。则
$f^{-1}\text{Fit}_i(\mathcal{F}) \cdot \mathcal{O}_T =
\text{Fit}_i(f^*\mathcal{F})$。
\end{lemma}

\begin{proof}
由《代数进阶》引理 \ref{more-algebra-lemma-fitting-ideal-basics}
的 (3) 立即得到。
\end{proof}

\begin{lemma}
\label{divisors-lemma-fitting-ideal-of-finitely-presented}
设 $S$ 为概形。
设 $\mathcal{F}$ 为有限展示 $\mathcal{O}_S$-模。则
$\text{Fit}_r(\mathcal{F})$ 是有限型拟凝聚理想。
\end{lemma}

\begin{proof}
由《代数进阶》引理 \ref{more-algebra-lemma-fitting-ideal-basics}
的 (4) 立即得到。
\end{proof}

\begin{lemma}
\label{divisors-lemma-on-subscheme-cut-out-by-Fit-0}
设 $S$ 为概形。
设 $\mathcal{F}$ 为有限型拟凝聚 $\mathcal{O}_S$-模。
设 $Z_0 \subset S$ 为由 $\text{Fit}_0(\mathcal{F})$ 截出的闭子概形。
设 $Z \subset S$ 为 $\mathcal{F}$ 的概形论支撑。则
\begin{enumerate}
\item 作为闭子概形，有 $Z \subset Z_0 \subset S$；
\item 作为闭子集，有 $Z = Z_0 = \text{Supp}(\mathcal{F})$；
\item 存在有限型拟凝聚 $\mathcal{O}_{Z_0}$-模 $\mathcal{G}_0$，使得
$$
(Z_0 \to X)_*\mathcal{G}_0 = \mathcal{F}.
$$
\end{enumerate}
\end{lemma}

\begin{proof}
回忆，$Z$ 局部由 $\mathcal{F}$ 的零化子截出，见《态射》定义
\ref{morphisms-definition-scheme-theoretic-support}（该定义使用《态射》引理
\ref{morphisms-lemma-scheme-theoretic-support} 来定义 $Z$）。因此，
由《代数进阶》引理 \ref{more-algebra-lemma-fitting-ideal-basics} 的 (6)，
可见在概形论意义下 $Z \subset Z_0$。另一方面，由《态射》引理
\ref{morphisms-lemma-scheme-theoretic-support}，集合论上有
$Z = \text{Supp}(\mathcal{F})$；又由《代数进阶》引理
\ref{more-algebra-lemma-fitting-ideal-basics} 的 (7)，集合论上有
$Z_0 = Z$。最后，为得到 (3) 中的 $\mathcal{G}_0$，可以使用
《态射》引理 \ref{morphisms-lemma-scheme-theoretic-support} 给出的
$\mathcal{G}$（它定义在 $Z$ 上），并置
$\mathcal{G}_0 = (Z \to Z_0)_*\mathcal{G}$；也可以使用《态射》引理
\ref{morphisms-lemma-i-star-equivalence}，以及由《代数进阶》引理
\ref{more-algebra-lemma-fitting-ideal-basics} 的 (6) 得到的事实：
$\text{Fit}_0(\mathcal{F})$ 零化 $\mathcal{F}$。
\end{proof}

\begin{lemma}
\label{divisors-lemma-fitting-ideal-generate-locally}
设 $S$ 为概形，$\mathcal{F}$ 为有限型拟凝聚 $\mathcal{O}_S$-模，
并设 $s \in S$。则 $\mathcal{F}$ 能由 $r$ 个元素生成，
并且可在 $s$ 的一个邻域中如此生成，当且仅当
$\text{Fit}_r(\mathcal{F})_s = \mathcal{O}_{S, s}$。
\end{lemma}

\begin{proof}
由《代数进阶》引理
\ref{more-algebra-lemma-fitting-ideal-generate-locally} 立即得到。
\end{proof}

\begin{lemma}
\label{divisors-lemma-fitting-ideal-finite-locally-free}
设 $S$ 为概形，$\mathcal{F}$ 为有限型拟凝聚 $\mathcal{O}_S$-模，
并设 $r \geq 0$。以下条件等价：
\begin{enumerate}
\item $\mathcal{F}$ 是秩为 $r$ 的有限局部自由模；
\item $\text{Fit}_{r - 1}(\mathcal{F}) = 0$ 且
$\text{Fit}_r(\mathcal{F}) = \mathcal{O}_S$；
\item $\text{Fit}_k(\mathcal{F}) = 0$ 对 $k < r$ 成立，而
$\text{Fit}_k(\mathcal{F}) = \mathcal{O}_S$ 对 $k \geq r$ 成立。
\end{enumerate}
\end{lemma}

\begin{proof}
由《代数进阶》引理
\ref{more-algebra-lemma-fitting-ideal-finite-locally-free} 立即得到。
\end{proof}

\begin{lemma}
\label{divisors-lemma-locally-free-rank-r-pullback}
设 $S$ 为概形，$\mathcal{F}$ 为有限型拟凝聚 $\mathcal{O}_S$-模。
由 $\mathcal{F}$ 的 Fitting 理想定义的闭子概形
$$
S = Z_{-1} \supset Z_0 \supset Z_1 \supset Z_2 \ldots
$$
具有下述性质：
\begin{enumerate}
\item 交集 $\bigcap Z_r$ 为空。
\item 由规则
$$
T \longmapsto
\left\{
\begin{matrix}
\{*\} & \text{若 }\mathcal{F}_T\text{ is locally generated by }
\leq r\text{ sections} \\
\emptyset & \text{otherwise}
\end{matrix}
\right.
$$
定义的函子 $(\Sch/S)^{opp} \to \textit{Sets}$ 由开子概形
$S \setminus Z_r$ 表示。
\item 由规则
$$
T \longmapsto
\left\{
\begin{matrix}
\{*\} & \text{若 }\mathcal{F}_T\text{ locally free rank }r\\
\emptyset & \text{otherwise}
\end{matrix}
\right.
$$
定义的函子 $F_r : (\Sch/S)^{opp} \to \textit{Sets}$ 由局部闭子概形
$Z_{r - 1} \setminus Z_r$ 表示；它是 $S$ 的局部闭子概形。
\end{enumerate}
若 $\mathcal{F}$ 有限展示，则态射 $Z_r \to S$、
$S \setminus Z_r \to S$ 和 $Z_{r - 1} \setminus Z_r \to S$
都是有限展示的。
\end{lemma}

\begin{proof}
(1) 成立，是因为在每个仿射开集 $U$ 上，都存在整数 $n$，使得
$\text{Fit}_n(\mathcal{F})|_U = \mathcal{O}_U$。事实上，可取 $n$
为 $\mathcal{F}$ 在 $U$ 上的生成元个数；见《代数进阶》第
\ref{more-algebra-section-fitting-ideals} 节。

\medskip\noindent
对任意态射 $g : T \to S$，由引理 \ref{divisors-lemma-base-change-fitting-ideal}
和 \ref{divisors-lemma-fitting-ideal-generate-locally} 可见，$\mathcal{F}_T$
局部由 $\leq r$ 个截面生成，当且仅当
$\text{Fit}_r(\mathcal{F}) \cdot \mathcal{O}_T = \mathcal{O}_T$。
这证明了 (2)。

\medskip\noindent
对任意态射 $g : T \to S$，由引理 \ref{divisors-lemma-base-change-fitting-ideal}
和 \ref{divisors-lemma-fitting-ideal-finite-locally-free} 可见，$\mathcal{F}_T$
是秩为 $r$ 的自由模，当且仅当
$\text{Fit}_r(\mathcal{F}) \cdot \mathcal{O}_T = \mathcal{O}_T$ 且
$\text{Fit}_{r - 1}(\mathcal{F}) \cdot \mathcal{O}_T = 0$。
这证明了 (3)。

\medskip\noindent
假设 $\mathcal{F}$ 有限展示。则每个态射 $Z_r \to S$ 都是有限展示的，
因为 $\text{Fit}_r(\mathcal{F})$ 有限型（引理
\ref{divisors-lemma-fitting-ideal-of-finitely-presented} 和《态射》引理
\ref{morphisms-lemma-closed-immersion-finite-presentation}）。这说明
$Z_{r - 1} \setminus Z_r$ 是 $Z_r$ 中的逆紧开集（《性质》引理
\ref{properties-lemma-quasi-coherent-finite-type-ideals}），因而态射
$Z_{r - 1} \setminus Z_r \to Z_r$ 也是有限展示的。
\end{proof}

\noindent
尽管有引理 \ref{divisors-lemma-locally-free-rank-r-pullback}，若 $\mathcal{F}$
只是有限型拟凝聚模，下述引理并不成立。确实，此时分层仍然存在，
但一般而言它并不表示函子 $F_{flat}$。

\begin{lemma}
\label{divisors-lemma-finite-presentation-module}
设 $S$ 为概形，$\mathcal{F}$ 为有限展示 $\mathcal{O}_S$-模。
设 $S = Z_{-1} \supset Z_0 \supset Z_1 \supset \ldots$ 如引理
\ref{divisors-lemma-locally-free-rank-r-pullback} 所述。置
$S_r = Z_{r - 1} \setminus Z_r$。则
$S' = \coprod_{r \geq 0} S_r$ 表示函子
$$
F_{flat} : \Sch/S \longrightarrow \textit{Sets},\quad\quad
T \longmapsto
\left\{
\begin{matrix}
\{*\} & \text{若 }\mathcal{F}_T\text{ flat over }T\\
\emptyset & \text{otherwise}
\end{matrix}
\right.
$$
此外，$\mathcal{F}|_{S_r}$ 是秩为 $r$ 的局部自由模，且态射
$S_r \to S$ 和 $S' \to S$ 都是有限展示的。
\end{lemma}

\begin{proof}
设 $g : T \to S$ 是概形态射，并且拉回
$\mathcal{F}_T = g^*\mathcal{F}$ 平坦。则 $\mathcal{F}_T$ 是有限展示
平坦 $\mathcal{O}_T$-模。因此，由《性质》引理
\ref{properties-lemma-finite-locally-free}，$\mathcal{F}_T$ 有限局部自由。
于是 $T = \coprod_{r \geq 0} T_r$，其中 $\mathcal{F}_T|_{T_r}$
局部自由且秩为 $r$。这说明在 $\Sch/S$ 上的 Zariski 层范畴中
$$
F_{flat} = \coprod\nolimits_{r \geq 0} F_r
$$
其中 $F_r$ 如引理 \ref{divisors-lemma-locally-free-rank-r-pullback} 所述。
因此，$F_{flat}$ 由
$\coprod_{r \geq 0} (Z_{r - 1} \setminus Z_r)$ 表示，其中 $Z_r$
如引理 \ref{divisors-lemma-locally-free-rank-r-pullback} 所述。
其他断言也由该引理得出。
\end{proof}

\begin{example}
\label{divisors-example-not-fp-MB}
设 $R = \prod_{n \in \mathbf{N}} \mathbf{F}_2$。设 $I \subset R$
为由这样的元素 $a = (a_n)_{n \in \mathbf{N}}$ 构成的理想：
它们只有有限多个分量非零。设 $\mathfrak m$ 为包含 $I$ 的极大理想。
则 $M = R/\mathfrak m$ 是有限平坦 $R$-模，因为 $R$ 绝对平坦
（《代数进阶》引理
\ref{more-algebra-lemma-product-fields-absolutely-flat}）。
置 $S = \Spec(R)$ 和 $\mathcal{F} = \widetilde{M}$。
引理 \ref{divisors-lemma-locally-free-rank-r-pullback} 中的闭子概形为
$S = Z_{-1}$、$Z_0 = \Spec(R/\mathfrak m)$，以及
$Z_i = \emptyset$（当 $i > 0$ 时）。但是，$\text{id} : S \to S$ 不通过
$(S \setminus Z_0) \amalg Z_0$ 分解，因为 $\mathfrak m$ 是 $S$
的非孤立点。因此，引理 \ref{divisors-lemma-finite-presentation-module}
对有限型模不成立。
\end{example}







\section{态射的奇异点集}
\label{divisors-section-singular-locus-morphism}

\noindent
设 $f : X \to S$ 为概形的有限型态射。令点集 $U$ 由 $f$ 光滑的点组成；
它是 $X$ 的开集（由《态射》定义 \ref{morphisms-definition-smooth}）。
在许多情形下，拥有一个典范闭子概形
$\text{Sing}(f) \subset X$ 很有用：它的补集为 $U$，并且它的构造
与任意基变换交换。

\medskip\noindent
若 $f$ 有限展示，一种选择是考虑由仿射局部意义下“严格标准”的
函数截出的闭子概形 $Z$；“严格标准”取《环映射的光滑化》定义
\ref{smoothing-definition-strictly-standard} 的意义。
由《环映射的光滑化》引理
\ref{smoothing-lemma-strictly-standard-base-change}，若
$f' : X' \to S'$ 是 $f$ 沿态射 $S' \to S$ 的基变换，则
$Z' \subset S' \times_S Z$，其中 $Z'$ 是 $X'$ 中由仿射局部严格标准
函数截出的闭子概形。然而，等号是否成立并不清楚。
严格标准元这一概念在 Popescu 定理一章中很有用。据我们所知，
由这些元素定义的闭子概形在文献中并未使用\footnote{若 $f$ 是
局部完全交态射（《态射进阶》定义
\ref{more-morphisms-definition-lci}），则由局部严格标准元截出的闭子概形
正是应当考察的对象。}。

\medskip\noindent
若 $f$ 平坦且有限展示，并且 $f$ 的所有纤维都等维且维数为 $d$，
则可用第 $d$ 个 Fitting 理想（即 $\Omega_{X/S}$ 的相应理想）得到一个
良好的闭子概形。对任意有限型态射，$X$ 的各闭子概形若由
$\Omega_{X/S}$ 的 Fitting 理想定义，则给出 $X$ 的一个分层；
该分层按照 $\Omega_{X/S}$ 的秩划分，其构造与基变换交换。
这一分层可能很有帮助，并与纤维的
嵌入维数有关，见《簇》第 \ref{varieties-section-embedding-dimension} 节。

\begin{lemma}
\label{divisors-lemma-base-change-and-fitting-ideal-omega}
设 $f : X \to S$ 为局部有限型的概形态射。
设 $X = Z_{-1} \supset Z_0 \supset Z_1 \supset \ldots$ 为由
$\Omega_{X/S}$ 的 Fitting 理想定义的闭子概形。则 $Z_i$ 的构造
与任意基变换交换。
\end{lemma}

\begin{proof}
注意，$\Omega_{X/S}$ 是有限型拟凝聚 $\mathcal{O}_X$-模
（《态射》引理 \ref{morphisms-lemma-finite-type-differentials}），
故其 Fitting 理想有定义。若 $f' : X' \to S'$ 是 $f$ 沿
$g : S' \to S$ 的基变换，则
$\Omega_{X'/S'} = (g')^*\Omega_{X/S}$，其中 $g' : X' \to X$ 为投影
（《态射》引理 \ref{morphisms-lemma-base-change-differentials}）。因此
$(g')^{-1}\text{Fit}_i(\Omega_{X/S}) \cdot \mathcal{O}_{X'} =
\text{Fit}_i(\Omega_{X'/S'})$。这表示在概形论意义下
$$
Z'_i = (g')^{-1}(Z_i) = Z_i \times_X X'
$$
而这正是本引理断言的含义。
\end{proof}

\noindent
第 $0$ 个 Fitting 理想（关于 $\Omega$）截出态射的“分歧点集”。

\begin{lemma}
\label{divisors-lemma-zero-fitting-ideal-omega-unramified}
设 $f : X \to S$ 为局部有限型的概形态射。
闭子概形 $Z \subset X$ 由第 $0$ 个 Fitting 理想（关于
$\Omega_{X/S}$）截出；它恰好是 $f$ 不为非分歧的点之集。
\end{lemma}

\begin{proof}
由引理 \ref{divisors-lemma-on-subscheme-cut-out-by-Fit-0}，$Z$ 的补集恰好是
$\Omega_{X/S}$ 为零的点集。由《态射》引理
\ref{morphisms-lemma-unramified-omega-zero}，这恰好是 $f$ 非分歧的点集。
\end{proof}

\begin{lemma}
\label{divisors-lemma-d-fitting-ideal-omega-smooth}
设 $f : X \to S$ 为概形的态射，并设 $d \geq 0$ 为整数。假设
\begin{enumerate}
\item $f$ 平坦；
\item $f$ 局部有限展示；
\item $f$ 的每个非空纤维都等维且维数为 $d$。
\end{enumerate}
设 $Z \subset X$ 为第 $d$ 个 Fitting 理想（关于 $\Omega_{X/S}$）
截出的闭子概形。则 $Z$ 恰好是 $f$ 不光滑的点之集。
\end{lemma}

\begin{proof}
由引理 \ref{divisors-lemma-locally-free-rank-r-pullback}，$Z$ 的补集恰好是
$\Omega_{X/S}$ 可由至多 $d$ 个元素生成的点集。因此，本引理由
《态射》引理 \ref{morphisms-lemma-smooth-at-point} 得出。
\end{proof}







\section{无挠模}
\label{divisors-section-torsion-free}

\noindent
本节是《代数进阶》第
\ref{more-algebra-section-torsion-flat} 节关于拟凝聚模的对应版本。

\begin{lemma}
\label{divisors-lemma-torsion-sections}
设 $X$ 为泛点是 $\eta$ 的整概形，$\mathcal{F}$ 为拟凝聚
$\mathcal{O}_X$-模。设 $U \subset X$ 为非空开集，且
$s \in \mathcal{F}(U)$。则下列条件等价：
\begin{enumerate}
\item 对某个 $x \in U$，$s$ 在 $\mathcal{F}_x$ 中的像是挠元；
\item 对所有 $x \in U$，$s$ 在 $\mathcal{F}_x$ 中的像都是挠元；
\item $s$ 在 $\mathcal{F}_\eta$ 中的像为零；
\item $s$ 在 $j_*\mathcal{F}_\eta$ 中的像为零，其中
$j : \eta \to X$ 是包含态射。
\end{enumerate}
\end{lemma}

\begin{proof}
略。
\end{proof}

\begin{definition}
\label{divisors-definition-torsion}
设 $X$ 为整概形，$\mathcal{F}$ 为拟凝聚 $\mathcal{O}_X$-模。
\begin{enumerate}
\item 若 $\mathcal{F}$ 的一个局部截面满足引理
\ref{divisors-lemma-torsion-sections} 中的等价条件，则称它是{\it 挠截面}。
\item 称 $\mathcal{F}$ 是{\it 无挠的}，如果 $\mathcal{F}$ 的每个
挠截面都是 $0$。
\end{enumerate}
\end{definition}

\noindent
下面给出按例必须有的引理，将此定义与通常的代数定义作比较。

\begin{lemma}
\label{divisors-lemma-check-torsion-on-affines}
设 $X$ 为整概形，$\mathcal{F}$ 为拟凝聚 $\mathcal{O}_X$-模。
则下列条件等价：
\begin{enumerate}
\item $\mathcal{F}$ 无挠；
\item 对每个仿射开集 $U \subset X$，$\mathcal{F}(U)$ 都是无挠
$\mathcal{O}(U)$-模。
\end{enumerate}
\end{lemma}

\begin{proof}
略。
\end{proof}

\begin{lemma}
\label{divisors-lemma-torsion}
设 $X$ 为整概形，$\mathcal{F}$ 为拟凝聚 $\mathcal{O}_X$-模。
$\mathcal{F}$ 的挠截面构成一个拟凝聚 $\mathcal{O}_X$-子模
$\mathcal{F}_{tors} \subset \mathcal{F}$。
商模 $\mathcal{F}/\mathcal{F}_{tors}$ 无挠。
\end{lemma}

\begin{proof}
略。代数版本见《代数进阶》引理 \ref{more-algebra-lemma-torsion}。
\end{proof}

\begin{lemma}
\label{divisors-lemma-flat-torsion-free}
设 $X$ 为整概形。任何平坦的拟凝聚 $\mathcal{O}_X$-模都是无挠的。
\end{lemma}

\begin{proof}
略。参见《代数进阶》引理
\ref{more-algebra-lemma-flat-torsion-free}。
\end{proof}

\begin{lemma}
\label{divisors-lemma-flat-pullback-torsion}
设 $f : X \to Y$ 为整概形之间的平坦态射，$\mathcal{G}$ 为无挠的
拟凝聚 $\mathcal{O}_Y$-模。则 $f^*\mathcal{G}$ 是无挠的
$\mathcal{O}_X$-模。
\end{lemma}

\begin{proof}
略。代数版本见《代数进阶》引理
\ref{more-algebra-lemma-flat-pullback-torsion}。
\end{proof}

\begin{lemma}
\label{divisors-lemma-flat-over-integral-integral-fibre}
设 $f : X \to Y$ 为概形的平坦态射。若 $Y$ 是整概形，且 $f$ 的泛纤维
是整概形，则 $X$ 是整概形。
\end{lemma}

\begin{proof}
其代数版本如下：设 $A$ 为整环，分式域为 $K$；设 $B$ 为平坦
$A$-代数，且 $B \otimes_A K$ 是整环。则 $B$ 是整环。这是因为由
《代数进阶》引理 \ref{more-algebra-lemma-flat-torsion-free}，
$B$ 无挠，因而 $B \subset B \otimes_A K$。
\end{proof}

\begin{lemma}
\label{divisors-lemma-check-torsion}
设 $X$ 为整概形，$\mathcal{F}$ 为拟凝聚 $\mathcal{O}_X$-模。
则 $\mathcal{F}$ 无挠，当且仅当 $\mathcal{F}_x$ 是无挠的
$\mathcal{O}_{X, x}$-模；这一条件须对所有 $x \in X$ 成立。
\end{lemma}

\begin{proof}
略。参见《代数进阶》引理
\ref{more-algebra-lemma-check-torsion}。
\end{proof}

\begin{lemma}
\label{divisors-lemma-extension-torsion-free}
设 $X$ 为整概形，并设
$0 \to \mathcal{F} \to \mathcal{F}' \to \mathcal{F}'' \to 0$
为拟凝聚 $\mathcal{O}_X$-模的短正合列。若 $\mathcal{F}$ 和
$\mathcal{F}''$ 无挠，则 $\mathcal{F}'$ 无挠。
\end{lemma}

\begin{proof}
略。代数版本见《代数进阶》引理
\ref{more-algebra-lemma-extension-torsion-free}。
\end{proof}

\begin{lemma}
\label{divisors-lemma-torsion-free-finite-noetherian-domain}
设 $X$ 为泛点是 $\eta$ 的局部诺特整概形，$\mathcal{F}$ 为非零凝聚
$\mathcal{O}_X$-模。则下列条件等价：
\begin{enumerate}
\item $\mathcal{F}$ 无挠；
\item $\eta$ 是 $\mathcal{F}$ 的唯一伴随素理想；
\item $\eta$ 属于 $\mathcal{F}$ 的支撑，且 $\mathcal{F}$ 具有性质
$(S_1)$；
\item $\eta$ 属于 $\mathcal{F}$ 的支撑，且 $\mathcal{F}$
没有嵌入伴随素理想。
\end{enumerate}
\end{lemma}

\begin{proof}
这是把《代数进阶》引理
\ref{more-algebra-lemma-torsion-free-finite-noetherian-domain}
翻译成概形语言的结果。略去该翻译。
\end{proof}

\begin{lemma}
\label{divisors-lemma-torsion-free-over-regular-dim-1}
设 $X$ 为维数 $\leq 1$ 的整正则概形，$\mathcal{F}$ 为凝聚
$\mathcal{O}_X$-模。则下列条件等价：
\begin{enumerate}
\item $\mathcal{F}$ 无挠；
\item $\mathcal{F}$ 有限局部自由。
\end{enumerate}
\end{lemma}

\begin{proof}
有限局部自由模显然无挠。反过来，我们将说明：若 $\mathcal{F}$ 无挠，
则 $\mathcal{F}_x$ 是自由 $\mathcal{O}_{X, x}$-模；这一结论对所有
$x \in X$ 成立。由《代数》引理
\ref{algebra-lemma-finite-projective} 以及 $\mathcal{F}$ 凝聚这一事实，
这就足够了。若 $\dim(\mathcal{O}_{X, x}) = 0$，则
$\mathcal{O}_{X, x}$ 是域，结论显然。若
$\dim(\mathcal{O}_{X, x}) = 1$，则 $\mathcal{O}_{X, x}$ 是离散赋值环
（《代数》引理 \ref{algebra-lemma-characterize-dvr}），并且
$\mathcal{F}_x$ 无挠。因此，由《代数进阶》引理
\ref{more-algebra-lemma-dedekind-torsion-free-flat}，
$\mathcal{F}_x$ 是自由的。
\end{proof}

\begin{lemma}
\label{divisors-lemma-hom-into-torsion-free}
设 $X$ 为整概形，$\mathcal{F}$、$\mathcal{G}$ 为拟凝聚
$\mathcal{O}_X$-模。若 $\mathcal{G}$ 无挠，且 $\mathcal{F}$ 有限展示，
则 $\SheafHom_{\mathcal{O}_X}(\mathcal{F}, \mathcal{G})$ 无挠。
\end{lemma}

\begin{proof}
该断言有意义，因为
$\SheafHom_{\mathcal{O}_X}(\mathcal{F}, \mathcal{G})$
由《概形》第 \ref{schemes-section-quasi-coherent} 节可知是拟凝聚的。
断言的证明见《代数进阶》引理
\ref{more-algebra-lemma-hom-into-torsion-free}。略去一些细节。
\end{proof}

\begin{lemma}
\label{divisors-lemma-isom-depth-2-torsion-free}
设 $X$ 为整的局部诺特概形，并设
$\varphi : \mathcal{F} \to \mathcal{G}$ 为拟凝聚
$\mathcal{O}_X$-模的映射。假设 $\mathcal{F}$ 凝聚、$\mathcal{G}$ 无挠，
且对每个 $x \in X$，下列情形之一成立：
\begin{enumerate}
\item $\mathcal{F}_x \to \mathcal{G}_x$ 是同构；或
\item $\text{depth}(\mathcal{F}_x) \geq 2$。
\end{enumerate}
则 $\varphi$ 是同构。
\end{lemma}

\begin{proof}
这是把《代数进阶》引理
\ref{more-algebra-lemma-isom-depth-2-torsion-free}
翻译成概形语言的结果。
\end{proof}










\section{模的秩}
\label{divisors-section-rank}

\noindent
本节是《代数进阶》第 \ref{more-algebra-section-rank} 节关于拟凝聚模的
对应版本。回忆，若 $X$ 是泛点为 $\eta$ 的整概形，则局部环
$\mathcal{O}_{X, \eta}$ 等于剩余域 $\kappa(\eta)$，见《态射》引理
\ref{morphisms-lemma-integral-scheme-rational-functions}。

\begin{definition}
\label{divisors-definition-rank}
设 $X$ 为泛点是 $\eta$ 的整概形，$\mathcal{F}$ 为拟凝聚
$\mathcal{O}_X$-模。$\mathcal{F}$ 的{\it 秩}定义为
$\dim_{\kappa(\eta)} \mathcal{F}_\eta$。
\end{definition}

\noindent
当两个定义都适用时，这与局部自由模之秩的定义并不冲突。若概形 $X$
不是不可约的，或者泛点处的局部环 $\mathcal{O}_{X, \eta}$ 不是域，则给
拟凝聚模定义秩看来不是好主意\footnote{若 $X$ 不可约，且
$\mathcal{F}_\eta$ 在 $\mathcal{O}_{X, \eta}$ 上自由，则采用该自由模的
秩似乎很自然。须注意，有些文献用 $\mathcal{F}_\eta$ 的长度定义秩。}

\begin{lemma}
\label{divisors-lemma-check-rank-on-affines}
设 $X$ 为整概形，$\mathcal{F}$ 为拟凝聚 $\mathcal{O}_X$-模，并设
$r$ 为一个基数。则下列条件等价：
\begin{enumerate}
\item $\mathcal{F}$ 的秩为 $r$；
\item 对所有非空仿射开集 $U \subset X$，$\mathcal{F}(U)$ 都是
$\mathcal{O}_X(U)$-模，其秩为 $r$；
\item 对某个非空仿射开集 $U \subset X$，$\mathcal{F}(U)$ 是
$\mathcal{O}_X(U)$-模，其秩为 $r$。
\end{enumerate}
\end{lemma}

\begin{proof}
略。
\end{proof}

\begin{lemma}
\label{divisors-lemma-dominant-pullback-rank}
设 $f : X \to Y$ 为整概形之间的支配态射，$\mathcal{G}$ 为拟凝聚
$\mathcal{O}_Y$-模。则 $\mathcal{G}$ 在 $Y$ 上的秩等于
$f^*\mathcal{G}$ 在 $X$ 上的秩。
\end{lemma}

\begin{proof}
使用引理 \ref{divisors-lemma-check-rank-on-affines}、《态射》引理
\ref{morphisms-lemma-dominant-between-integral} 以及《代数进阶》引理
\ref{more-algebra-lemma-dominant-pullback-rank}。
\end{proof}

\begin{lemma}
\label{divisors-lemma-rank-additive}
设 $X$ 为整概形，并设
$0 \to \mathcal{F} \to \mathcal{F}' \to \mathcal{F}'' \to 0$
为拟凝聚 $\mathcal{O}_X$-模的短正合列。则 $\mathcal{F}'$ 的秩等于
$\mathcal{F}$ 与 $\mathcal{F}''$ 的秩之和。
\end{lemma}

\begin{proof}
略。参见《代数进阶》引理 \ref{more-algebra-lemma-rank-additive}。
\end{proof}

\begin{lemma}
\label{divisors-lemma-rank-tensor}
设 $X$ 为整概形，$\mathcal{F}$ 和 $\mathcal{G}$ 为拟凝聚
$\mathcal{O}_X$-模。
\begin{enumerate}
\item $\mathcal{F} \otimes_{\mathcal{O}_X} \mathcal{G}$ 的秩等于
$\mathcal{F}$ 与 $\mathcal{G}$ 的秩之积。
\item 若 $\mathcal{F}$ 有限展示，则
$\SheafHom_{\mathcal{O}_X}(\mathcal{F}, \mathcal{G})$ 的秩等于
$\mathcal{F}$ 与 $\mathcal{G}$ 的秩之积。
\end{enumerate}
\end{lemma}

\begin{proof}
略。注意，(2) 是有意义的，因为
$\SheafHom_{\mathcal{O}_X}(\mathcal{F}, \mathcal{G})$
由《概形》第 \ref{schemes-section-quasi-coherent} 节的 (10)
可知是拟凝聚的。本引理的代数版本是《代数进阶》引理
\ref{more-algebra-lemma-rank-tensor}。
\end{proof}

\begin{lemma}
\label{divisors-lemma-finite-module-rank}
设 $X$ 为整概形，$\mathcal{F}$ 为有限型拟凝聚 $\mathcal{O}_X$-模。
则
\begin{enumerate}
\item $\mathcal{F}$ 具有有限秩 $r \geq 0$；
\item 存在非空开集 $U \subset X$，使得 $\mathcal{F}|_U$ 是秩为
$r$ 的自由模。
\end{enumerate}
\end{lemma}

\begin{proof}
代数版本见《代数进阶》引理
\ref{more-algebra-lemma-finite-module-rank}。
\end{proof}










\section{自反模}
\label{divisors-section-reflexive}

\noindent
本节是《代数进阶》第 \ref{more-algebra-section-reflexive} 节关于局部诺特
概形上凝聚模的对应版本。之所以考察凝聚模，是因为
$\SheafHom_{\mathcal{O}_X}(\mathcal{F}, \mathcal{G})$ 对任意一对凝聚
$\mathcal{O}_X$-模 $\mathcal{F}, \mathcal{G}$ 都是凝聚的，
见《模》引理
\ref{modules-lemma-internal-hom-locally-kernel-direct-sum}。

\begin{definition}
\label{divisors-definition-reflexive}
设 $X$ 为整的局部诺特概形，$\mathcal{F}$ 为凝聚
$\mathcal{O}_X$-模。$\mathcal{F}$ 的{\it 自反包}是
$\mathcal{O}_X$-模
$$
\mathcal{F}^{**} = \SheafHom_{\mathcal{O}_X}(
\SheafHom_{\mathcal{O}_X}(\mathcal{F}, \mathcal{O}_X), \mathcal{O}_X)
$$
称 $\mathcal{F}$ 是{\it 自反的}，如果自然映射
$j : \mathcal{F} \longrightarrow \mathcal{F}^{**}$ 是同构。
\end{definition}

\noindent
由引理 \ref{divisors-lemma-dual-reflexive}，自反包是自反
$\mathcal{O}_X$-模。在更一般的情形中也可以用同一定义来定义自反模，
但这看来并不十分有用。下面给出按例必须有的引理，将此定义与通常的
代数定义作比较。

\begin{lemma}
\label{divisors-lemma-check-reflexive-on-affines}
设 $X$ 为整的局部诺特概形，$\mathcal{F}$ 为凝聚
$\mathcal{O}_X$-模。则下列条件等价：
\begin{enumerate}
\item $\mathcal{F}$ 自反；
\item 对仿射开集 $U \subset X$，$\mathcal{F}(U)$ 是自反
$\mathcal{O}(U)$-模。
\end{enumerate}
\end{lemma}

\begin{proof}
略。
\end{proof}

\begin{remark}
\label{divisors-remark-different-reflexive}
若 $X$ 是域上的有限型概形，则有时会使用另一种自反模的概念
（例如见 \cite[第 5 页底部及定义 1.1.9]{HL}）。这种概念使用
$R\SheafHom$ 取值于对偶化复形 $\omega_X^\bullet$，而不是取值于
$\mathcal{O}_X$ 的版本；它也许应当另取名称，因为当 $X$ 不是
Gorenstein 概形时，两者可能不同。例如，若
$X = \Spec(k[t^3, t^4, t^5])$，则计算表明，对偶化层 $\omega_X$
按我们的定义并不自反，但按另一种定义是自反的，因为
$\omega_X \to \SheafHom(\SheafHom(\omega_X, \omega_X), \omega_X)$
是同构。
\end{remark}

\begin{lemma}
\label{divisors-lemma-reflexive-torsion-free}
设 $X$ 为整的局部诺特概形，$\mathcal{F}$ 为凝聚
$\mathcal{O}_X$-模。
\begin{enumerate}
\item 若 $\mathcal{F}$ 自反，则 $\mathcal{F}$ 无挠。
\item 映射 $j : \mathcal{F} \longrightarrow \mathcal{F}^{**}$ 是单射，
当且仅当 $\mathcal{F}$ 无挠。
\end{enumerate}
\end{lemma}

\begin{proof}
略。参见《代数进阶》引理
\ref{more-algebra-lemma-reflexive-torsion-free}。
\end{proof}

\begin{lemma}
\label{divisors-lemma-check-reflexive}
设 $X$ 为整的局部诺特概形，$\mathcal{F}$ 为凝聚
$\mathcal{O}_X$-模。则下列条件等价：
\begin{enumerate}
\item $\mathcal{F}$ 自反；
\item $\mathcal{F}_x$ 是自反 $\mathcal{O}_{X, x}$-模，这对所有
$x \in X$ 都成立；
\item $\mathcal{F}_x$ 是自反 $\mathcal{O}_{X, x}$-模，这对所有闭点
$x \in X$ 都成立。
\end{enumerate}
\end{lemma}

\begin{proof}
由《模》引理 \ref{modules-lemma-stalk-internal-hom}，(1) 与 (2)
等价。由于 $X$ 的每个点都特化到某个闭点（《性质》引理
\ref{properties-lemma-locally-Noetherian-closed-point}），
(2) 与 (3) 等价。
\end{proof}

\begin{lemma}
\label{divisors-lemma-flat-pullback-reflexive}
设 $f : X \to Y$ 为整的局部诺特概形之间的平坦态射，$\mathcal{G}$
为凝聚自反 $\mathcal{O}_Y$-模。则 $f^*\mathcal{G}$ 是凝聚自反
$\mathcal{O}_X$-模。
\end{lemma}

\begin{proof}
略。代数版本见《代数进阶》引理
\ref{more-algebra-lemma-flat-pullback-reflexive}。
\end{proof}

\begin{lemma}
\label{divisors-lemma-sequence-reflexive}
设 $X$ 为整的局部诺特概形，并设
$0 \to \mathcal{F} \to \mathcal{F}' \to \mathcal{F}''$
为凝聚 $\mathcal{O}_X$-模的正合列。若 $\mathcal{F}'$ 自反且
$\mathcal{F}''$ 无挠，则 $\mathcal{F}$ 自反。
\end{lemma}

\begin{proof}
略。参见《代数进阶》引理
\ref{more-algebra-lemma-sequence-reflexive}。
\end{proof}

\begin{lemma}
\label{divisors-lemma-dual-reflexive}
设 $X$ 为整的局部诺特概形，$\mathcal{F}$、$\mathcal{G}$ 为凝聚
$\mathcal{O}_X$-模。若 $\mathcal{G}$ 自反，则
$\SheafHom_{\mathcal{O}_X}(\mathcal{F}, \mathcal{G})$ 自反。
\end{lemma}

\begin{proof}
该断言有意义，因为
$\SheafHom_{\mathcal{O}_X}(\mathcal{F}, \mathcal{G})$
由《概形的上同调》引理
\ref{coherent-lemma-tensor-hom-coherent} 可知是凝聚的。断言的证明见
《代数进阶》引理 \ref{more-algebra-lemma-dual-reflexive}。
略去一些细节。
\end{proof}

\begin{remark}
\label{divisors-remark-tensor}
设 $X$ 为整的局部诺特概形。由引理 \ref{divisors-lemma-dual-reflexive}，
我们知道自反包 $\mathcal{F}^{**}$ 对凝聚 $\mathcal{O}_X$-模而言
是凝聚自反的。考虑范畴 $\mathcal{C}$，其对象为凝聚自反
$\mathcal{O}_X$-模。取自反包给出包含函子
$\mathcal{C} \to \textit{Coh}(\mathcal{O}_X)$ 的左伴随。
注意，$\mathcal{C}$ 是具有核和余核的加性范畴。事实上，给定
$\varphi : \mathcal{F} \to \mathcal{G}$，并设它是 $\mathcal{C}$ 中的态射，
通常的核 $\Ker(\varphi)$ 是自反的（引理
\ref{divisors-lemma-sequence-reflexive}），而通常余核的自反包
$\Coker(\varphi)^{**}$ 是 $\mathcal{C}$ 中的余核。此外，
$\mathcal{C}$ 继承了张量积
$$
\mathcal{F} \otimes_\mathcal{C} \mathcal{G} =
(\mathcal{F} \otimes_{\mathcal{O}_X} \mathcal{G})^{**}
$$
它是结合且对称的。还存在内 Hom：对任意三个对象
$\mathcal{F}, \mathcal{G}, \mathcal{H}$（它们属于 $\mathcal{C}$），有恒等式
$$
\Hom_\mathcal{C}(\mathcal{F} \otimes_\mathcal{C} \mathcal{G}, \mathcal{H}) =
\Hom_\mathcal{C}(\mathcal{F},
\SheafHom_{\mathcal{O}_X}(\mathcal{G}, \mathcal{H}))
$$
见《模》引理 \ref{modules-lemma-internal-hom}。在 $\mathcal{C}$ 中，
每个对象 $\mathcal{F}$ 都有一个{\it 对偶对象}
$\SheafHom_{\mathcal{O}_X}(\mathcal{F}, \mathcal{O}_X)$。
若不对 $X$ 加进一步条件，则可能有
$$
\SheafHom_{\mathcal{O}_X}(\mathcal{F}, \mathcal{G}) \not \cong
\SheafHom_{\mathcal{O}_X}(\mathcal{F}, \mathcal{O}_X)
\otimes_\mathcal{C} \mathcal{G}
\quad\text{且}\quad
\mathcal{F} \otimes_\mathcal{C}
\SheafHom_{\mathcal{O}_X}(\mathcal{F}, \mathcal{O}_X)
\not \cong \mathcal{O}_X
$$
其中 $\mathcal{F}, \mathcal{G}$ 的秩为 $1$，且它们是 $\mathcal{C}$ 中的对象。
为构造一个例子，令 $X = \Spec(R)$，其中 $R$ 如《代数进阶》例
\ref{more-algebra-example-ring-not-S2} 所述，并令 $\mathcal{F}, \mathcal{G}$
为与 $M$ 对应的模。略去计算。
\end{remark}

\begin{lemma}
\label{divisors-lemma-reflexive-depth-2}
设 $X$ 为整的局部诺特概形，$\mathcal{F}$ 为凝聚
$\mathcal{O}_X$-模。则下列条件等价：
\begin{enumerate}
\item $\mathcal{F}$ 自反；
\item 对每个 $x \in X$，下列情形之一成立：
\begin{enumerate}
\item $\mathcal{F}_x$ 是自反 $\mathcal{O}_{X, x}$-模；或
\item $\text{depth}(\mathcal{F}_x) \geq 2$。
\end{enumerate}
\end{enumerate}
\end{lemma}

\begin{proof}
略。参见《代数进阶》引理
\ref{more-algebra-lemma-reflexive-depth-2}。
\end{proof}

\begin{lemma}
\label{divisors-lemma-reflexive-S2}
设 $X$ 为整的局部诺特概形，$\mathcal{F}$ 为凝聚自反
$\mathcal{O}_X$-模，并设 $x \in X$。
\begin{enumerate}
\item 若 $\text{depth}(\mathcal{O}_{X, x}) \geq 2$，则
$\text{depth}(\mathcal{F}_x) \geq 2$。
\item 若 $X$ 具有 $(S_2)$ 性质，则 $\mathcal{F}$ 具有 $(S_2)$ 性质。
\end{enumerate}
\end{lemma}

\begin{proof}
略。参见《代数进阶》引理
\ref{more-algebra-lemma-reflexive-S2}。
\end{proof}

\begin{lemma}
\label{divisors-lemma-reflexive-S2-extend}
设 $X$ 为整的局部诺特概形，$j : U \to X$ 为补集是 $Z$ 的开子概形。
假设 $\mathcal{O}_{X, z}$ 的深度 $\geq 2$，且这一条件对所有
$z \in Z$ 成立。则 $j^*$ 与 $j_*$ 给出凝聚自反 $\mathcal{O}_X$-模范畴
同凝聚自反 $\mathcal{O}_U$-模范畴之间的等价。
\end{lemma}

\begin{proof}
设 $\mathcal{F}$ 为凝聚自反 $\mathcal{O}_X$-模。对 $z \in Z$，
由引理 \ref{divisors-lemma-reflexive-S2}，茎 $\mathcal{F}_z$ 的深度
$\geq 2$。因此，由引理 \ref{divisors-lemma-depth-2-hartog}，
$\mathcal{F} \to j_*j^*\mathcal{F}$ 是同构。反之，设 $\mathcal{G}$
为凝聚自反 $\mathcal{O}_U$-模。只须证明 $j_*\mathcal{G}$ 是凝聚自反
$\mathcal{O}_X$-模。为此，可以假设 $X$ 仿射。由《性质》引理
\ref{properties-lemma-lift-finite-presentation}，存在凝聚
$\mathcal{O}_X$-模 $\mathcal{F}$，使得
$\mathcal{G} = j^*\mathcal{F}$。取 $\mathcal{F}$ 的自反包后，
可以假设 $\mathcal{F}$ 自反
（见上文讨论，特别是引理 \ref{divisors-lemma-dual-reflexive}）。
由上述论证，$j_*\mathcal{G} = j_*j^*\mathcal{F} = \mathcal{F}$，
这正是所需结论。
\end{proof}

\noindent
若概形正规，则自反等价于无挠并满足 $(S_2)$。

\begin{lemma}
\label{divisors-lemma-reflexive-over-normal}
设 $X$ 为整的局部诺特正规概形，$\mathcal{F}$ 为凝聚
$\mathcal{O}_X$-模。则下列条件等价：
\begin{enumerate}
\item $\mathcal{F}$ 自反；
\item $\mathcal{F}$ 无挠并具有性质 $(S_2)$；
\item 存在开子概形 $j : U \to X$，使得
\begin{enumerate}
\item $X \setminus U$ 的每个不可约分支的余维都 $\geq 2$
（余维在 $X$ 中计算）；
\item $j^*\mathcal{F}$ 有限局部自由；
\item $\mathcal{F} = j_*j^*\mathcal{F}$。
\end{enumerate}
\end{enumerate}
\end{lemma}

\begin{proof}
使用引理 \ref{divisors-lemma-check-reflexive-on-affines}，(1) 与 (2) 的等价性
由《代数进阶》引理
\ref{more-algebra-lemma-reflexive-over-normal} 得出。设
$U \subset X$ 如 (3) 所述。由《性质》引理
\ref{properties-lemma-criterion-normal}，有
$\text{depth}(\mathcal{O}_{X, x}) \geq 2$，其中 $x \not \in U$。
由于有限局部自由模自反，
引理 \ref{divisors-lemma-reflexive-S2-extend} 表明 (3) 蕴含 (1)。

\medskip\noindent
假设 (1)。设 $U \subset X$ 为使
$j^*\mathcal{F} = \mathcal{F}|_U$ 有限局部自由的极大开子概形。
于是 (3)(b) 成立。设 $x \in X$。若 $\mathcal{F}_x$ 是自由
$\mathcal{O}_{X, x}$-模，则由《模》引理
\ref{modules-lemma-finite-presentation-stalk-free} 有 $x \in U$。
若 $\dim(\mathcal{O}_{X, x}) \leq 1$，则
$\mathcal{O}_{X, x}$ 是域或离散赋值环（《性质》引理
\ref{properties-lemma-criterion-normal}），从而 $\mathcal{F}_x$ 是自由的
（《代数进阶》引理
\ref{more-algebra-lemma-dedekind-torsion-free-flat}）。因此
$x \not \in U \Rightarrow \dim(\mathcal{O}_{X, x}) \geq 2$。
此时《性质》引理
\ref{properties-lemma-codimension-local-ring} 给出 (3)(a)。
再次使用前述《性质》引理
\ref{properties-lemma-criterion-normal}，还可知
$\text{depth}(\mathcal{O}_{X, x}) \geq 2$，其中 $x \not \in U$；
于是 (3)(c) 由引理
\ref{divisors-lemma-reflexive-S2-extend} 得出。
\end{proof}

\begin{lemma}
\label{divisors-lemma-describe-reflexive-hull}
设 $X$ 为泛点是 $\eta$ 的整局部诺特正规概形，$\mathcal{F}$、
$\mathcal{G}$ 为凝聚 $\mathcal{O}_X$-模。设
$T : \mathcal{G}_\eta \to \mathcal{F}_\eta$ 为线性映射。则 $T$
可延拓为映射 $\mathcal{G} \to \mathcal{F}^{**}$，且该映射为
$\mathcal{O}_X$-模映射，当且仅当
\begin{itemize}
\item[(*)] 对每个 $x \in X$，若 $\dim(\mathcal{O}_{X, x}) = 1$，则有
$$
T\left(\Im(\mathcal{G}_x \to \mathcal{G}_\eta)\right) \subset
\Im(\mathcal{F}_x \to \mathcal{F}_\eta).
$$
\end{itemize}
\end{lemma}

\begin{proof}
由于 $\mathcal{F}^{**}$ 无挠，且
$\mathcal{F}_\eta = \mathcal{F}^{**}_\eta$，所以延拓若存在就唯一。
因此，只须在 $X$ 的一个开覆盖的各成员上证明本引理；也就是说，可以假设
$X$ 仿射。在这种情形下，我们问的是如下代数问题：设 $R$ 为分式域是
$K$ 的诺特正规整环，$M$、$N$ 为有限 $R$-模，并设
$T : M \otimes_R K \to N \otimes_R K$ 为 $K$-线性映射。$T$
何时能延拓为映射 $N \to M^{**}$？由《代数进阶》引理
\ref{more-algebra-lemma-describe-reflexive-hull}，这当且仅当
$N_\mathfrak p$ 映入 $(M/M_{tors})_\mathfrak p$，而这对每个高为
$1$ 的素理想 $\mathfrak p$（属于 $R$）都成立。这恰好是本引理的
条件 $(*)$。
\end{proof}

\begin{lemma}
\label{divisors-lemma-reflexive-over-regular-dim-2}
设 $X$ 为维数 $\leq 2$ 的正则概形，$\mathcal{F}$ 为凝聚
$\mathcal{O}_X$-模。则下列条件等价：
\begin{enumerate}
\item $\mathcal{F}$ 自反；
\item $\mathcal{F}$ 有限局部自由。
\end{enumerate}
\end{lemma}

\begin{proof}
有限局部自由模显然自反。反过来，我们将说明：若 $\mathcal{F}$ 自反，
则 $\mathcal{F}_x$ 是自由 $\mathcal{O}_{X, x}$-模，这对所有
$x \in X$ 都成立。由《代数》引理
\ref{algebra-lemma-finite-projective} 以及 $\mathcal{F}$ 凝聚这一事实，
这就足够了。若 $\dim(\mathcal{O}_{X, x}) = 0$，则
$\mathcal{O}_{X, x}$ 是域，结论显然。若
$\dim(\mathcal{O}_{X, x}) = 1$，则 $\mathcal{O}_{X, x}$ 是离散赋值环
（《代数》引理 \ref{algebra-lemma-characterize-dvr}），并且
$\mathcal{F}_x$ 无挠。因此，由《代数进阶》引理
\ref{more-algebra-lemma-dedekind-torsion-free-flat}，
$\mathcal{F}_x$ 是自由的。若 $\dim(\mathcal{O}_{X, x}) = 2$，
则 $\mathcal{O}_{X, x}$ 是维数为 $2$ 的正则局部环。由《代数进阶》引理
\ref{more-algebra-lemma-reflexive-over-normal}，
$\mathcal{F}_x$ 的深度 $\geq 2$。因此，由《代数》引理
\ref{algebra-lemma-regular-mcm-free}，$\mathcal{F}_x$ 是自由的。
\end{proof}






\section{有效 Cartier 除子}
\label{divisors-section-effective-Cartier-divisors}

\noindent
在定义其他类型的除子之前，我们先定义有效 Cartier 除子。

\begin{definition}
\label{divisors-definition-effective-Cartier-divisor}
设 $S$ 为概形。
\begin{enumerate}
\item $S$ 的一个{\it 局部主闭子概形}，是其理想层可局部由单个元素
生成的闭子概形。
\item $S$ 上的一个{\it 有效 Cartier 除子}，是闭子概形
$D \subset S$，其理想层 $\mathcal{I}_D \subset \mathcal{O}_S$
是可逆 $\mathcal{O}_S$-模。
\end{enumerate}
\end{definition}

\noindent
因此，有效 Cartier 除子是局部主闭子概形，但反之并非总是成立。
在最强的可能意义下，有效 Cartier 除子是纯余维 $1$ 的闭子概形。
具体而言，它们局部由单个非零因子截出。特别地，它们无处稠密。

\begin{lemma}
\label{divisors-lemma-characterize-effective-Cartier-divisor}
设 $S$ 为概形，$D \subset S$ 为闭子概形。则下列条件等价：
\begin{enumerate}
\item 子概形 $D$ 是 $S$ 上的有效 Cartier 除子。
\item 对每个 $x \in D$，都存在仿射开邻域
$\Spec(A) = U \subset S$，它是 $x$ 的邻域，并且
$U \cap D = \Spec(A/(f))$，其中 $f \in A$ 是非零因子。
\end{enumerate}
\end{lemma}

\begin{proof}
假设 (1)。对每个 $x \in D$，都存在仿射开邻域
$\Spec(A) = U \subset S$，它是 $x$ 的邻域，并且
$\mathcal{I}_D|_U \cong \mathcal{O}_U$。换言之，存在截面
$f \in \Gamma(U, \mathcal{I}_D)$，它自由生成限制
$\mathcal{I}_D|_U$。因此 $f \in A$，且乘法映射
$f : A \to A$ 是单射。又因为 $\mathcal{I}_D$ 拟凝聚，所以
$D \cap U = \Spec(A/(f))$。

\medskip\noindent
假设 (2)。设 $x \in D$。根据假设，存在仿射开邻域
$\Spec(A) = U \subset S$，它是 $x$ 的邻域，并且
$U \cap D = \Spec(A/(f))$，其中 $f \in A$ 是非零因子。
于是 $\mathcal{I}_D|_U \cong \mathcal{O}_U$，因为它等于
$\widetilde{(f)} \cong \widetilde{A} \cong \mathcal{O}_U$。
当然，$\mathcal{I}_D$ 在开子概形 $S \setminus D$ 上的限制同构于
$\mathcal{O}_{S \setminus D}$。因此 $\mathcal{I}_D$ 是可逆
$\mathcal{O}_S$-模。
\end{proof}

\begin{lemma}
\label{divisors-lemma-complement-locally-principal-closed-subscheme}
设 $S$ 为概形，$Z \subset S$ 为局部主闭子概形，并设
$U = S \setminus Z$。则 $U \to S$ 是仿射态射。
\end{lemma}

\begin{proof}
该问题在 $S$ 上是局部的，见《态射》引理
\ref{morphisms-lemma-characterize-affine}。因此可以假设
$S = \Spec(A)$ 且 $Z = V(f)$，其中 $f \in A$。此时
$U = D(f) = \Spec(A_f)$ 仿射，故 $U \to S$ 仿射。
\end{proof}

\begin{lemma}
\label{divisors-lemma-complement-effective-Cartier-divisor}
设 $S$ 为概形，$D \subset S$ 为有效 Cartier 除子，并设
$U = S \setminus D$。则 $U \to S$ 是仿射态射，且 $U$ 在 $S$ 中
概形论稠密。
\end{lemma}

\begin{proof}
仿射性由引理
\ref{divisors-lemma-complement-locally-principal-closed-subscheme} 给出。
稠密性问题在 $S$ 上是局部的，见《态射》引理
\ref{morphisms-lemma-characterize-scheme-theoretically-dense}。
因此可以假设 $S = \Spec(A)$，而 $D$ 对应于非零因子 $f \in A$，
见引理 \ref{divisors-lemma-characterize-effective-Cartier-divisor}。于是
$A \subset A_f$，这说明 $U \subset S$ 概形论稠密，见《态射》例
\ref{morphisms-example-scheme-theoretic-closure}。
\end{proof}

\begin{lemma}
\label{divisors-lemma-effective-Cartier-makes-dimension-drop}
设 $S$ 为概形，$D \subset S$ 为有效 Cartier 除子，并设
$s \in D$。若 $\dim_s(S) < \infty$，则
$\dim_s(D) < \dim_s(S)$。
\end{lemma}

\begin{proof}
假设 $\dim_s(S) < \infty$。设 $U = \Spec(A) \subset S$ 为 $s$
的仿射开邻域，使得 $\dim(U) = \dim_s(S)$，并使得
$D = V(f)$，其中 $f \in A$ 是非零因子（见引理
\ref{divisors-lemma-characterize-effective-Cartier-divisor}）。回忆，
$\dim(U)$ 是环 $A$ 的 Krull 维数，而 $\dim(U \cap D)$ 是环
$A/(f)$ 的 Krull 维数。于是 $f$ 不属于 $A$ 的任何极小素理想。
因此，把 $A/(f)$ 中任一极大素理想链看作 $A$ 中的素理想链后，
都可以添加一个极小素理想来延长它。
\end{proof}

\begin{definition}
\label{divisors-definition-sum-effective-Cartier-divisors}
设 $S$ 为概形。给定有效 Cartier 除子 $D_1$、$D_2$，它们定义在 $S$ 上，
令 $D = D_1 + D_2$ 为 $S$ 的闭子概形，它对应于拟凝聚理想层
$\mathcal{I}_{D_1}\mathcal{I}_{D_2} \subset \mathcal{O}_S$。
称它为{\it 有效 Cartier 除子 $D_1$ 与 $D_2$ 的和}。
\end{definition}

\noindent
显然，可以定义和 $\sum n_iD_i$，只要给定有限多个有效 Cartier
除子 $D_i$（定义在 $X$ 上）和非负整数 $n_i$。

\begin{lemma}
\label{divisors-lemma-sum-effective-Cartier-divisors}
两个有效 Cartier 除子的和仍是有效 Cartier 除子。
\end{lemma}

\begin{proof}
略。局部地，$f_1, f_2 \in A$ 都是非零因子，因而
$f_1f_2 \in A$ 也是非零因子。
\end{proof}

\begin{lemma}
\label{divisors-lemma-difference-effective-Cartier-divisors}
设 $X$ 为概形，$D, D'$ 为 $X$ 上两个有效 Cartier 除子。
若 $D \subset D'$（作为 $X$ 的闭子概形），则存在有效 Cartier
除子 $D''$，使得 $D' = D + D''$。
\end{lemma}

\begin{proof}
略。
\end{proof}

\begin{lemma}
\label{divisors-lemma-sum-closed-subschemes-effective-Cartier}
设 $X$ 为概形，$Z, Y$ 为 $X$ 的两个闭子概形，其理想层分别为
$\mathcal{I}$ 和 $\mathcal{J}$。若 $\mathcal{I}\mathcal{J}$
定义有效 Cartier 除子 $D \subset X$，则 $Z$ 和 $Y$ 都是
有效 Cartier 除子，且 $D = Z + Y$。
\end{lemma}

\begin{proof}
应用引理 \ref{divisors-lemma-characterize-effective-Cartier-divisor}，得到如下
代数情形：$A$ 是环，$I, J \subset A$ 是理想，$f \in A$ 是非零因子，
且 $IJ = (f)$。因此，结论由《代数》引理
\ref{algebra-lemma-product-ideals-principal} 得出。
\end{proof}

\begin{lemma}
\label{divisors-lemma-sum-effective-Cartier-divisors-union}
设 $X$ 为概形，$D, D' \subset X$ 为有效 Cartier 除子，并假设
概形论交 $D \cap D'$ 是 $D'$ 上的有效 Cartier 除子。则
$D + D'$ 是 $D$ 与 $D'$ 的概形论并。
\end{lemma}

\begin{proof}
概形论交与并的定义见《态射》定义
\ref{morphisms-definition-scheme-theoretic-intersection-union}。
为证明本引理，局部工作（使用引理
\ref{divisors-lemma-characterize-effective-Cartier-divisor}）后得到如下代数问题：
给定环 $A$ 与非零因子 $f_1, f_2 \in A$，且 $f_1$ 在
$A/f_2A$ 中的像是非零因子，证明
$f_1A \cap f_2A = f_1f_2A$。略去这个直接的论证。
\end{proof}

\noindent
回忆，我们已在《概形》定义
\ref{schemes-definition-inverse-image-closed-subscheme} 中定义了闭子概形
在任意概形态射下的逆像。

\begin{lemma}
\label{divisors-lemma-pullback-locally-principal}
设 $f : S' \to S$ 为概形态射，$Z \subset S$ 为局部主闭子概形。
则逆像 $f^{-1}(Z)$ 是 $S'$ 的局部主闭子概形。
\end{lemma}

\begin{proof}
略。
\end{proof}

\begin{definition}
\label{divisors-definition-pullback-effective-Cartier-divisor}
设 $f : S' \to S$ 为概形态射，$D \subset S$ 为有效 Cartier 除子。
称{\it $D$ 沿 $f$ 的拉回有定义}，如果闭子概形
$f^{-1}(D) \subset S'$ 是有效 Cartier 除子。此时将它记作 $f^*D$ 或
$f^{-1}(D)$，并称之为{\it 有效 Cartier 除子的拉回}。
\end{definition}

\noindent
$f^{-1}(D)$ 是有效 Cartier 除子这一条件在实践中经常满足。
下面给出一个示例性引理。

\begin{lemma}
\label{divisors-lemma-pullback-effective-Cartier-defined}
设 $f : X \to Y$ 为概形态射，$D \subset Y$ 为有效 Cartier 除子。
在下列每种情形中，$D$ 沿 $f$ 的拉回都有定义：
\begin{enumerate}
\item $f(x) \not \in D$ 对每个弱伴随点 $x$（属于 $X$）都成立；
\item $X$、$Y$ 均为整概形，且 $f$ 支配；
\item $X$ 约化，并且 $f(\xi) \not \in D$ 对任一泛点 $\xi$
（属于 $X$ 的任一不可约分支）都成立；
\item $X$ 局部诺特，并且 $f(x) \not \in D$ 对任一伴随点 $x$
（属于 $X$）都成立；
\item $X$ 局部诺特且没有嵌入点，并且 $f(\xi) \not \in D$ 对任一
泛点 $\xi$（属于 $X$ 的某个不可约分支）都成立；
\item $f$ 平坦；
\item 可按需在此添加更多情形。
\end{enumerate}
\end{lemma}

\begin{proof}
该问题在 $X$ 上是局部的，因此可约化到如下情形：
$X = \Spec(A)$，$Y = \Spec(R)$，$f$ 由
$\varphi : R \to A$ 给出，并且 $D = \Spec(R/(t))$，其中
$t \in R$ 是非零因子。每种情形中的目标都是证明
$\varphi(t) \in A$ 是非零因子。

\medskip\noindent
情形 (1) 由《代数》引理
\ref{algebra-lemma-weakly-ass-zero-divisors} 得出。由引理
\ref{divisors-lemma-ass-weakly-ass}，情形 (4) 是 (1) 的特例。情形 (5)
由 (4) 和定义得出。由引理 \ref{divisors-lemma-weakass-reduced}，情形 (3)
是 (1) 的特例。情形 (2) 是 (3) 的特例。若 $R \to A$ 平坦，则
$t : R \to R$ 是单射这一事实表明 $t : A \to A$ 是单射。
这证明了 (6)。
\end{proof}

\begin{lemma}
\label{divisors-lemma-pullback-effective-Cartier-divisors-additive}
设 $f : S' \to S$ 为概形态射，$D_1$、$D_2$ 为 $S$ 上的
有效 Cartier 除子。若 $D_1$ 与 $D_2$ 的拉回都有定义，则
$D = D_1 + D_2$ 的拉回有定义，且
$f^*D = f^*D_1 + f^*D_2$。
\end{lemma}

\begin{proof}
略。
\end{proof}





\section{有效 Cartier 除子与可逆层}
\label{divisors-section-effective-Cartier-invertible}

\noindent
由于有效 Cartier 除子具有可逆理想层（定义
\ref{divisors-definition-effective-Cartier-divisor}），下述定义是有意义的。

\begin{definition}
\label{divisors-definition-invertible-sheaf-effective-Cartier-divisor}
设 $S$ 为概形，$D \subset S$ 为有效 Cartier 除子，其理想层为
$\mathcal{I}_D$。
\begin{enumerate}
\item {\it 可逆层 $\mathcal{O}_S(D)$（与 $D$ 伴随）}定义为
$$
\mathcal{O}_S(D) =
\SheafHom_{\mathcal{O}_S}(\mathcal{I}_D, \mathcal{O}_S) =
\mathcal{I}_D^{\otimes -1}.
$$
\item {\it 典范截面}通常记作 $1$ 或 $1_D$；它是
$\mathcal{O}_S(D)$ 的整体截面，对应于包含映射
$\mathcal{I}_D \to \mathcal{O}_S$。
\item 记
$\mathcal{O}_S(-D) = \mathcal{O}_S(D)^{\otimes -1} = \mathcal{I}_D$。
\item 给定第二个有效 Cartier 除子 $D' \subset S$，定义
$\mathcal{O}_S(D - D') =
\mathcal{O}_S(D) \otimes_{\mathcal{O}_S} \mathcal{O}_S(-D')$。
\end{enumerate}
\end{definition}

\noindent
作几点说明。下文将看到，赋值 $D \mapsto \mathcal{O}_S(D)$ 把有效
Cartier 除子的加法（定义
\ref{divisors-definition-sum-effective-Cartier-divisors}）变为 $S$ 的 Picard 群
中的加法（引理
\ref{divisors-lemma-invertible-sheaf-sum-effective-Cartier-divisors}）。
然而，上述定义中的表达式 $D - D'$ 没有任何几何含义。更确切地说，
可以把 $S$ 上的有效 Cartier 除子之集看作交换幺半群
$\text{EffCart}(S)$，其零元为空的有效 Cartier 除子。于是赋值
$(D, D') \mapsto \mathcal{O}_S(D - D')$ 定义群同态
$$
\text{EffCart}(S)^{gp} \longrightarrow \Pic(S)
$$
其中左端是 $\text{EffCart}(S)$ 的群完备化。换言之，当写
$\mathcal{O}_S(D - D')$ 时，可以把 $D - D'$ 看作
$\text{EffCart}(S)^{gp}$ 的元素。

\begin{lemma}
\label{divisors-lemma-conormal-effective-Cartier-divisor}
设 $S$ 为概形，$D \subset S$ 为有效 Cartier 除子。则余法层为
$\mathcal{C}_{D/S} = \mathcal{I}_D|_D =
\mathcal{O}_S(-D)|_D$，而法层为
$\mathcal{N}_{D/S} = \mathcal{O}_S(D)|_D$。
\end{lemma}

\begin{proof}
这由《态射》引理 \ref{morphisms-lemma-affine-conormal} 得出。
\end{proof}

\begin{lemma}
\label{divisors-lemma-ses-add-divisor}
设 $X$ 为概形，$D, D' \subset X$ 为有效 Cartier 除子，且
$D \subset D'$；再设 $C \subset X$ 为满足 $D' = D + C$ 的有效
Cartier 除子（其存在性由引理
\ref{divisors-lemma-difference-effective-Cartier-divisors} 保证）。则有短正合列
$$
0 \to \mathcal{O}_X(-D)|_C \to \mathcal{O}_{D'} \to \mathcal{O}_D \to 0
$$
它由 $\mathcal{O}_X$-模组成。
\end{lemma}

\begin{proof}
在本引理的陈述和证明中，我们使用《态射》引理
\ref{morphisms-lemma-i-star-equivalence} 的等价，把 $X$ 的闭子概形
上的拟凝聚模视为 $X$ 上的拟凝聚模。设 $\mathcal{I}$ 为 $D$ 在
$D'$ 中的理想层，并设 $\mathcal{I}_D = \mathcal{O}_X(-D)$ 为
$D$ 在 $X$ 中的理想层。于是有典范满射
$\mathcal{I}_D \to \mathcal{I}$。因此，只须证明该映射的核等于
$\mathcal{I}_C\mathcal{I}_D$（采用显然的记号）。这可局部检验，
故可假设 $X = \Spec(A)$、$D = \Spec(A/fA)$ 且
$C = \Spec(A/gA)$，其中 $f, g \in A$ 为非零因子
（引理 \ref{divisors-lemma-characterize-effective-Cartier-divisor}）。
于是 $\mathcal{I}_D$ 对应于 $fA$，非零因子 $fg \in A$ 生成
$D'$ 的理想，而 $\mathcal{I}$ 对应于 $I = fA/fgA$。结论随即清楚。
\end{proof}

\begin{lemma}
\label{divisors-lemma-invertible-sheaf-sum-effective-Cartier-divisors}
设 $S$ 为概形，$D_1$、$D_2$ 为 $S$ 上的有效 Cartier 除子，并设
$D = D_1 + D_2$。则存在唯一同构
$$
\mathcal{O}_S(D_1) \otimes_{\mathcal{O}_S} \mathcal{O}_S(D_2)
\longrightarrow
\mathcal{O}_S(D)
$$
它把 $1_{D_1} \otimes 1_{D_2}$ 映到 $1_D$。
\end{lemma}

\begin{proof}
略。
\end{proof}

\begin{lemma}
\label{divisors-lemma-pullback-effective-Cartier-divisors}
设 $f : S' \to S$ 为概形态射，$D$ 为 $S$ 上的有效 Cartier 除子。
若 $D$ 的拉回有定义，则
$f^*\mathcal{O}_S(D) = \mathcal{O}_{S'}(f^*D)$，并且典范截面
$1_D$ 拉回为典范截面 $1_{f^*D}$。
\end{lemma}

\begin{proof}
略。
\end{proof}

\begin{definition}
\label{divisors-definition-regular-section}
设 $(X, \mathcal{O}_X)$ 为局部赋环空间，$\mathcal{L}$ 为 $X$ 上的
可逆层。若整体截面 $s \in \Gamma(X, \mathcal{L})$ 所给出的映射
$\mathcal{O}_X \to \mathcal{L}$（$f \mapsto fs$）是单射，则称它为
{\it 正则截面}。
\end{definition}

\begin{lemma}
\label{divisors-lemma-regular-section-structure-sheaf}
设 $X$ 为局部赋环空间，$f \in \Gamma(X, \mathcal{O}_X)$。
则下列条件等价：
\begin{enumerate}
\item $f$ 是正则截面；
\item 对任意 $x \in X$，像 $f \in \mathcal{O}_{X, x}$ 是非零因子。
\end{enumerate}
若 $X$ 是概形，则它们还等价于
\begin{enumerate}
\item[(3)] 对任意仿射开集 $\Spec(A) = U \subset X$，像 $f \in A$
是非零因子；
\item[(4)] 存在仿射开覆盖 $X = \bigcup \Spec(A_i)$，使得 $f$ 在
$A_i$ 中的像对所有 $i$ 都是非零因子。
\end{enumerate}
\end{lemma}

\begin{proof}
略。
\end{proof}

\noindent
注意，整体截面 $s$ 属于某个可逆 $\mathcal{O}_X$-模
$\mathcal{L}$；它可以视为 $\mathcal{O}_X$-模映射
$s : \mathcal{O}_X \to \mathcal{L}$。因此，其对偶是映射
$s : \mathcal{L}^{\otimes -1} \to \mathcal{O}_X$。
（对偶可逆层的定义见《模》定义
\ref{modules-definition-powers}。）

\begin{definition}
\label{divisors-definition-zero-scheme-s}
设 $X$ 为概形，$\mathcal{L}$ 为可逆层，并设
$s \in \Gamma(X, \mathcal{L})$ 为整体截面。$s$ 的{\it 零点概形}
是闭子概形 $Z(s) \subset X$，它由拟凝聚理想层
$\mathcal{I} \subset \mathcal{O}_X$ 定义；该理想层是映射
$s : \mathcal{L}^{\otimes -1} \to \mathcal{O}_X$ 的像。
\end{definition}

\begin{lemma}
\label{divisors-lemma-zero-scheme}
设 $X$ 为概形，$\mathcal{L}$ 为可逆层，并设
$s \in \Gamma(X, \mathcal{L})$。
\begin{enumerate}
\item 考虑闭浸入 $i : Z \to X$，要求
$i^*s \in \Gamma(Z, i^*\mathcal{L})$ 为零，并按包含关系排序。
零点概形 $Z(s)$ 是该有序集的极大元。
\item 对任意概形态射 $f : Y \to X$，有
$f^*s = 0$（在 $\Gamma(Y, f^*\mathcal{L})$ 中），当且仅当 $f$
经由 $Z(s)$ 分解。
\item 零点概形 $Z(s)$ 是局部主闭子概形。
\item 零点概形 $Z(s)$ 是有效 Cartier 除子，当且仅当 $s$ 是
$\mathcal{L}$ 的正则截面。
\end{enumerate}
\end{lemma}

\begin{proof}
略。
\end{proof}

\begin{lemma}
\label{divisors-lemma-characterize-OD}
\begin{slogan}
概形上的有效 Cartier 除子等同于带有指定正则整体截面的可逆层。
\end{slogan}
设 $X$ 为概形。
\begin{enumerate}
\item 若 $D \subset X$ 是有效 Cartier 除子，则典范截面 $1_D$
（属于 $\mathcal{O}_X(D)$）是正则的。
\item 反之，若 $s$ 是可逆层 $\mathcal{L}$ 的正则截面，则存在唯一
有效 Cartier 除子 $D = Z(s) \subset X$ 和唯一同构
$\mathcal{O}_X(D) \to \mathcal{L}$，它把 $1_D$ 映到 $s$。
\end{enumerate}
构造 $D \mapsto (\mathcal{O}_X(D), 1_D)$ 和
$(\mathcal{L}, s) \mapsto Z(s)$ 给出互逆的映射
$$
\left\{
\begin{matrix}
\text{effective Cartier divisors on }X
\end{matrix}
\right\}
\leftrightarrow
\left\{
\begin{matrix}
\text{isomorphism classes of pairs }(\mathcal{L}, s)\\
\text{consisting of an invertible }
\mathcal{O}_X\text{-模}\\
\mathcal{L}\text{ and a regular global section }s
\end{matrix}
\right\}
$$
\end{lemma}

\begin{proof}
略。
\end{proof}

\begin{remark}
\label{divisors-remark-ses-regular-section}
设 $X$ 为概形，$\mathcal{L}$ 为可逆 $\mathcal{O}_X$-模，$s$ 为
$\mathcal{L}$ 的正则截面。则零点概形 $D = Z(s)$ 是 $X$ 上的
有效 Cartier 除子，并有短正合列
$$
0 \to \mathcal{O}_X \to \mathcal{L} \to i_*(\mathcal{L}|_D) \to 0
\quad\text{且}\quad
0 \to \mathcal{L}^{\otimes -1} \to \mathcal{O}_X \to i_*\mathcal{O}_D \to 0.
$$
给定有效 Cartier 除子 $D \subset X$，使用引理
\ref{divisors-lemma-characterize-OD} 和
\ref{divisors-lemma-conormal-effective-Cartier-divisor}，得到
$$
0 \to \mathcal{O}_X \to \mathcal{O}_X(D) \to i_*(\mathcal{N}_{D/X}) \to 0
\quad\text{且}\quad
0 \to \mathcal{O}_X(-D) \to \mathcal{O}_X \to i_*(\mathcal{O}_D) \to 0
$$
\end{remark}






\section{诺特概形上的有效 Cartier 除子}
\label{divisors-section-Noetherian-effective-Cartier}

\noindent
在局部诺特的情形中，关于有效 Cartier 除子和正则截面的许多讨论
都有所简化。

\begin{lemma}
\label{divisors-lemma-regular-section-associated-points}
设 $X$ 为局部诺特概形，$\mathcal{L}$ 为可逆
$\mathcal{O}_X$-模，并设 $s \in \Gamma(X, \mathcal{L})$。则 $s$
是正则截面，当且仅当 $s$ 在 $X$ 的伴随点处不消失。
\end{lemma}

\begin{proof}
略。提示：约化到仿射情形，并令 $\mathcal{L}$ 平凡，然后使用引理
\ref{divisors-lemma-regular-section-structure-sheaf} 和《代数》引理
\ref{algebra-lemma-ass-zero-divisors}。
\end{proof}

\begin{lemma}
\label{divisors-lemma-effective-Cartier-in-points}
设 $X$ 为局部诺特概形，$D \subset X$ 为闭子概形，对应于拟凝聚理想层
$\mathcal{I} \subset \mathcal{O}_X$。
\begin{enumerate}
\item 若对每个 $x \in D$，理想
$\mathcal{I}_x \subset \mathcal{O}_{X, x}$ 都可以由一个元素生成，
则 $D$ 是局部主的。
\item 若对每个 $x \in D$，理想
$\mathcal{I}_x \subset \mathcal{O}_{X, x}$ 都可以由一个非零因子生成，
则 $D$ 是有效 Cartier 除子。
\end{enumerate}
\end{lemma}

\begin{proof}
设 $\Spec(A)$ 为点 $x \in D$ 的仿射邻域。设
$\mathfrak p \subset A$ 为对应于 $x$ 的素理想。设 $I \subset A$
为定义 $D$ 在 $\Spec(A)$ 上之迹的理想。因为 $A$ 是诺特的
（由于 $X$ 局部诺特），理想 $I$ 可由有限多个元素生成，记为
$I = (f_1, \ldots, f_r)$。在 (1) 的假设下，有
$I_\mathfrak p = (f)$，其中 $f \in A_\mathfrak p$。于是有
$f_i = g_i f$，其中 $g_i \in A_\mathfrak p$。写
$g_i = a_i/h_i$ 且 $f = f'/h$，其中
$a_i, h_i, f', h \in A$，并且 $h_i, h \not \in \mathfrak p$。
于是 $I_{h_1 \ldots h_r h} \subset A_{h_1 \ldots h_r h}$ 是主理想，
因为它由 $f'$ 生成。这证明了 (1)。对于 (2)，可以假设
$I = (f)$。假设意味着 $f$ 在 $A_\mathfrak p$ 中的像是非零因子。
由《代数》引理
\ref{algebra-lemma-regular-sequence-in-neighbourhood}，$f$ 在 $x$
的某个邻域上是非零因子。这证明了 (2)。
\end{proof}

\begin{lemma}
\label{divisors-lemma-effective-Cartier-codimension-1}
设 $X$ 为局部诺特概形。
\begin{enumerate}
\item 设 $D \subset X$ 为局部主闭子概形，$\xi \in D$ 为 $D$
某个不可约分支的泛点。则 $\dim(\mathcal{O}_{X, \xi}) \leq 1$。
\item 设 $D \subset X$ 为有效 Cartier 除子，$\xi \in D$ 为 $D$
某个不可约分支的泛点。则 $\dim(\mathcal{O}_{X, \xi}) = 1$。
\end{enumerate}
\end{lemma}

\begin{proof}
证明 (1)。根据假设，可以设 $X = \Spec(A)$ 且
$D = \Spec(A/(f))$，其中 $A$ 为诺特环，$f \in A$。设 $\xi$
对应于素理想 $\mathfrak p \subset A$。$\xi$ 是 $D$ 某个不可约分支
的泛点这一假设，意味着 $\mathfrak p$ 是包含 $(f)$ 的极小素理想。
于是由《代数》引理 \ref{algebra-lemma-minimal-over-1}，
$\dim(A_\mathfrak p) \leq 1$。

\medskip\noindent
证明 (2)。由 (1) 可知 $\dim(\mathcal{O}_{X, \xi}) \leq 1$。
另一方面，由引理
\ref{divisors-lemma-characterize-effective-Cartier-divisor}，局部方程 $f$
在 $A_\mathfrak p$ 中是非零因子，这说明维数至少为 $1$
（因为由初等的《代数》引理
\ref{algebra-lemma-Zariski-topology}，$A_\mathfrak p$ 中必有一个
不包含 $f$ 的素理想）。
\end{proof}

\begin{lemma}
\label{divisors-lemma-integral-effective-Cartier-divisor-dvr}
设 $X$ 为诺特概形，$D \subset X$ 为整闭子概形，同时也是有效
Cartier 除子。则 $X$ 在 $D$ 的泛点处的局部环是离散赋值环。
\end{lemma}

\begin{proof}
由引理 \ref{divisors-lemma-characterize-effective-Cartier-divisor}，可以假设
$X = \Spec(A)$ 且 $D = \Spec(A/(f))$，其中 $A$ 为诺特环，
$f \in A$ 为非零因子。$D$ 为整概形这一假设意味着 $(f)$ 是素理想。
因此，$X$ 在泛点处的局部环为 $A_{(f)}$；这是诺特局部环，其极大理想
由一个非零因子生成。于是由《代数》引理
\ref{algebra-lemma-characterize-dvr}，它是离散赋值环。
\end{proof}

\begin{lemma}
\label{divisors-lemma-effective-Cartier-divisor-Sk}
设 $X$ 为局部诺特概形，$D \subset X$ 为有效 Cartier 除子。
若 $X$ 具有 $(S_k)$ 性质，则 $D$ 具有 $(S_{k - 1})$ 性质。
\end{lemma}

\begin{proof}
设 $x \in D$。则
$\mathcal{O}_{D, x} = \mathcal{O}_{X, x}/(f)$，其中
$f \in \mathcal{O}_{X, x}$ 是非零因子。根据假设，有
$\text{depth}(\mathcal{O}_{X, x}) \geq \min(\dim(\mathcal{O}_{X, x}), k)$。
由《代数》引理 \ref{algebra-lemma-depth-drops-by-one}，有
$\text{depth}(\mathcal{O}_{D, x}) = \text{depth}(\mathcal{O}_{X, x}) - 1$；
由《代数》引理 \ref{algebra-lemma-one-equation}，有
$\dim(\mathcal{O}_{D, x}) = \dim(\mathcal{O}_{X, x}) - 1$。
因此
$\text{depth}(\mathcal{O}_{D, x}) \geq \min(\dim(\mathcal{O}_{D, x}), k - 1)$，
这正是所需结论。
\end{proof}

\begin{lemma}
\label{divisors-lemma-normal-effective-Cartier-divisor-S1}
设 $X$ 为局部诺特正规概形，$D \subset X$ 为有效 Cartier 除子。
则 $D$ 具有 $(S_1)$ 性质。
\end{lemma}

\begin{proof}
由《性质》引理 \ref{properties-lemma-criterion-normal}，可知 $X$
具有 $(S_2)$ 性质。因此，由引理
\ref{divisors-lemma-effective-Cartier-divisor-Sk} 得出结论。
\end{proof}

\begin{lemma}
\label{divisors-lemma-weil-divisor-is-cartier-UFD}
设 $X$ 为局部诺特概形，$D \subset X$ 为整闭子概形。假设
\begin{enumerate}
\item $D$ 的余维为 $1$（余维在 $X$ 中计算）；
\item $\mathcal{O}_{X, x}$ 是 UFD，且这对所有 $x \in D$ 都成立。
\end{enumerate}
则 $D$ 是有效 Cartier 除子。
\end{lemma}

\begin{proof}
设 $x \in D$，并令 $A = \mathcal{O}_{X, x}$。设
$\mathfrak p \subset A$ 对应于 $D$ 的泛点。根据 (1)，
$A_\mathfrak p$ 的维数为 $1$。因此 $\mathfrak p$ 是高为 $1$
的素理想。由于 $A$ 是 UFD，这意味着
$\mathfrak p = (f)$，其中 $f \in A$。当然，$f$ 是非零因子，
于是结论由引理
\ref{divisors-lemma-effective-Cartier-in-points} 得出。
\end{proof}

\begin{lemma}
\label{divisors-lemma-codim-1-part}
设 $X$ 为诺特概形，$Z \subset X$ 为闭子概形。假设存在
\begin{enumerate}
\item 一族整有效 Cartier 除子 $D_i \subset X$，$i \in I$；
\item 闭子集 $Z' \subset X$，其所有不可约分支的余维都
$\geq 2$（余维在 $X$ 中计算；《拓扑》定义
\ref{topology-definition-codimension}）。
\end{enumerate}
并且集合论上有 $Z \subset Z' \cup \bigcup_{i \in I} D_i$。
则存在整数 $a_i \geq 0$，其中 $a_i = 0$ 对几乎所有
$i \in I$ 都成立，使得有效 Cartier 除子
$$
D = \sum a_i D_i
$$
包含在 $Z$ 中，并且包含态射 $D \to Z$ 在余维 $2$ 之外
是同构（余维在 $X$ 中计算；其含义是：存在开集 $U \subset Z$，使
$D \cap U \to Z \cap U$ 为同构，并且 $Z \setminus U$ 的每个
不可约分支的余维都 $\geq 2$，余维在 $X$ 中计算）。当 $Z$ 在 $X$
中无处稠密时，$D_i$、$i \in I$ 和 $Z'$ 的存在性得到保证，只要
$\mathcal{O}_{X, x}$ 是 UFD 且这对所有 $x \in Z$ 都成立，或者
$X$ 正则。
\end{lemma}

\begin{proof}
设 $\xi_i \in D_i$ 为泛点，并令
$\mathcal{O}_i = \mathcal{O}_{X, \xi_i}$ 为局部环；由引理
\ref{divisors-lemma-integral-effective-Cartier-divisor-dvr}，它是离散赋值环。
设 $a_i \geq 0$ 为
$\mathcal{I}_{Z, \xi_i} \subset \mathcal{O}_i$ 中元素的最小赋值。
当然，$a_i > 0$ 仅可能在 $D_i$ 是 $Z$ 的不可约分支时成立，因此
$a_i > 0$ 只对有限多个 $i \in I$ 成立。我们断言有效 Cartier
除子 $D = \sum a_i D_i$ 符合要求。

\medskip\noindent
具体地，设 $x \in X$，并令 $A = \mathcal{O}_{X, x}$。设
$D_1, \ldots, D_n$ 为两两不同的除子 $D_i$，满足
$x \in D_i$ 且 $a_i > 0$。对 $1 \leq i \leq n$，令
$f_i \in A$ 为 $D_i$ 的局部方程。则 $f_i$ 是 $A$ 的素元，且
$\mathcal{O}_i = A_{(f_i)}$。设
$I = \mathcal{I}_{Z, x} \subset A$ 为 $Z$ 的理想层之茎。
根据 $a_i$ 的选择，有
$I A_{(f_i)} = f_i^{a_i}A_{(f_i)}$。我们断言
$I \subset (\prod_{i = 1, \ldots, n} f_i^{a_i})$。

\medskip\noindent
证明该断言。局部化映射
$\varphi : A/(f_i) \to A_{(f_i)}/f_iA_{(f_i)}$ 是单射，因为素理想
$(f_i)$ 是极大理想 $f_iA_{(f_i)}$ 的逆像。对 $n$ 归纳可知
$\varphi_n : A/(f_i^n)\to A_{(f_i)}/f_i^nA_{(f_i)}$ 也为单射。
由于 $\varphi_{a_i}(I) = 0$，有 $I \subset (f_i^{a_i})$。
因此，对任意 $x \in I$，可写成 $x = f_1^{a_1}x_1$，其中
$x_1 \in A$。由于 $D_1, \ldots, D_n$ 两两不同，$f_i$ 在
$A_{(f_j)}$ 中是单位，只要 $i \not = j$。比较 $x$ 与 $x_1$ 在
$A_{(f_i)}$ 中的像，对 $n \geq i > 1$ 仍有
$x_1 \in (f_i^{a_i})$。重复上述过程，归纳地写出
$x_i = f_{i + 1}^{a_{i + 1}}x_{i + 1}$，这对任意
$n > i \geq 1$ 成立。总之，有
$x \in (\prod_{i = 1, \ldots n} f_i^{a_i})$，这对任意
$x \in I$ 成立，证毕。

\medskip\noindent
该断言表明 $\mathcal{I}_Z \subset \mathcal{I}_D$，即
$D \subset Z$。此外，还可看出 $D$ 与 $Z$ 在 $\xi_i$ 处相同，
这证明 $D \to Z$ 在余维 $2$ 之外是同构（余维在 $X$ 中计算）。

\medskip\noindent
最后几个断言可如下证明。正则局部环是 UFD（《代数进阶》引理
\ref{more-algebra-lemma-regular-local-UFD}），故只须处理 UFD 情形。
在该情形中，令 $D_i$ 为 $Z$ 的余维为 $1$ 的不可约分支
（余维在 $X$ 中计算）。由引理
\ref{divisors-lemma-weil-divisor-is-cartier-UFD}，每个 $D_i$ 都是有效
Cartier 除子。
\end{proof}

\begin{lemma}
\label{divisors-lemma-codimension-1-is-effective-Cartier}
设 $Z \subset X$ 为诺特概形的闭子概形。假设
\begin{enumerate}
\item $Z$ 没有嵌入点；
\item $Z$ 的每个不可约分支的余维均为 $1$（余维在 $X$ 中计算）；
\item 每个局部环 $\mathcal{O}_{X, x}$（$x \in Z$）都是
UFD，或者 $X$ 正则。
\end{enumerate}
则 $Z$ 是有效 Cartier 除子。
\end{lemma}

\begin{proof}
设 $D = \sum a_i D_i$ 如引理 \ref{divisors-lemma-codim-1-part} 中所述，
其中 $D_i \subset Z$ 是 $Z$ 的不可约分支。
若 $D \to Z$ 不是同构，则 $\mathcal{O}_Z \to \mathcal{O}_D$
有一个非零核，其支集位于余维 $\geq 2$ 处。这将意味着
$Z$ 存在嵌入点，与假设 (1) 矛盾。因此如所需有 $D \cong Z$。
\end{proof}

\begin{lemma}
\label{divisors-lemma-UFD-one-equation-CM}
设 $R$ 为诺特 UFD。设 $I \subset R$ 为一个理想，
使得 $R/I$ 没有嵌入素理想，且包含 $I$
的每个极小素理想的高均为 $1$。
则 $I = (f)$，其中 $f \in R$。
\end{lemma}

\begin{proof}
由引理 \ref{divisors-lemma-codimension-1-is-effective-Cartier}，
理想层 $\tilde I$ 在 $\Spec(R)$ 上可逆。
由《代数进阶》引理 \ref{more-algebra-lemma-UFD-Pic-trivial}，
它由一个元生成。
\end{proof}

\begin{lemma}
\label{divisors-lemma-effective-Cartier-divisor-is-a-sum}
设 $X$ 为诺特概形，$D \subset X$ 为有效 Cartier 除子。
假设存在整有效 Cartier 除子 $D_i \subset X$，使得集合论上有
$D \subset \bigcup D_i$。则 $D = \sum a_i D_i$，其中 $a_i \geq 0$。
$D_i$ 的存在性得到保证，只要 $\mathcal{O}_{X, x}$ 是 UFD，
且这对所有 $x \in D$ 都成立，或者 $X$ 正则。
\end{lemma}

\begin{proof}
如引理 \ref{divisors-lemma-codim-1-part} 中选取 $a_i$，并令
$D' = \sum a_i D_i$。则 $D' \to D$ 是有效 Cartier 除子的包含，
且在余维 $2$ 之外是同构（余维在 $X$ 上计算）。取 $x \in X$。
令 $A = \mathcal{O}_{X, x}$，并设 $f, f' \in A$ 分别生成
$D, D'$ 之理想的非零因子（在 $A$ 中）。则 $f = gf'$，
其中 $g \in A$。此外，对每个素理想 $\mathfrak p$，若其高
$\leq 1$（在 $A$ 中计算），则 $g$ 在 $A_\mathfrak p$ 中的像是单位。这意味着
$g$ 是单位，因为包含 $(g)$ 的极小素理想的高为 $1$
（《代数》引理 \ref{algebra-lemma-minimal-over-1}）。
\end{proof}

\begin{lemma}
\label{divisors-lemma-quasi-projective-Noetherian-pic-effective-Cartier}
\begin{slogan}
在射影概形上，每个线丛都有正则亚纯截面。
\end{slogan}
设 $X$ 为具有丰富可逆层的诺特概形。
则每个可逆 $\mathcal{O}_X$-模都同构于
$$
\mathcal{O}_X(D - D') =
\mathcal{O}_X(D) \otimes_{\mathcal{O}_X} \mathcal{O}_X(D')^{\otimes -1}
$$
其中 $D, D'$ 是 $X$ 中的某些有效 Cartier 除子。此外，给定
有限子集 $E \subset X$，可以选取 $D, D'$ 使得
$E \cap D = \emptyset$ 且 $E \cap D' = \emptyset$。若
$X$ 拟仿射，则可选取 $D' = \emptyset$。
\end{lemma}

\begin{proof}
设 $x_1, \ldots, x_n$ 为 $X$ 的伴随点
（引理 \ref{divisors-lemma-finite-ass}）。

\medskip\noindent
若 $X$ 拟仿射，而 $\mathcal{N}$ 是任意可逆
$\mathcal{O}_X$-模，则可以选取一个截面 $t$ 属于
$\mathcal{N}$，
使其在 $E \cup \{x_1, \ldots, x_n\}$ 的任一点处都不消失，见《性质》引理
\ref{properties-lemma-quasi-affine-invertible-nonvanishing-section}。
由引理 \ref{divisors-lemma-regular-section-associated-points}，$t$ 是
$\mathcal{N}$ 的正则截面。因此 $\mathcal{N} \cong \mathcal{O}_X(D)$，其中
$D = Z(t)$ 是对应于 $t$ 的有效 Cartier 除子，见引理
\ref{divisors-lemma-characterize-OD}。由构造可知 $E \cap D = \emptyset$，
因而该情形已证。

\medskip\noindent
回到一般情形，设 $\mathcal{L}$ 为 $X$ 上的丰富可逆层。
存在 $m > 0$ 以及截面 $s \in \Gamma(X, \mathcal{L}^{\otimes m})$，
使得 $X_s$ 为仿射概形，并且
$E \cup \{x_1, \ldots, x_n\} \subset X_s$
（《性质》引理 \ref{properties-lemma-ample-finite-set-in-principal-affine}）。

\medskip\noindent
设 $\mathcal{N}$ 为任意可逆 $\mathcal{O}_X$-模。
由拟仿射情形，可以找到截面 $t \in \mathcal{N}(X_s)$，
它在 $E \cup \{x_1, \ldots, x_n\}$ 的任一点处均不消失。
由《性质》引理 \ref{properties-lemma-invert-s-sections}，
对某个 $e \geq 0$，截面 $s^e|_{X_s} t$ 可以延拓为
一个整体截面 $\tau$，它是 $\mathcal{L}^{\otimes me} \otimes \mathcal{N}$ 的截面。
因此 $\mathcal{L}^{\otimes me} \otimes \mathcal{N}$ 与
$\mathcal{L}^{\otimes me}$ 都是可逆层，并且各自有一个整体截面，
该截面在 $E \cup \{x_1, \ldots, x_n\}$ 的任一点处均不消失。
由引理 \ref{divisors-lemma-regular-section-associated-points}，这些都是正则截面。
故有 $\mathcal{L}^{\otimes me} \otimes \mathcal{N} \cong \mathcal{O}_X(D)$
以及 $\mathcal{L}^{\otimes me} \cong \mathcal{O}_X(D')$，其中 $D$ 与 $D'$
是某些有效 Cartier 除子，见引理
\ref{divisors-lemma-characterize-OD}。由构造可知 $E \cap D = \emptyset$ 且
$E \cap D' = \emptyset$，证毕。
\end{proof}

\begin{lemma}
\label{divisors-lemma-wedge-product-ses}
设 $X$ 为整正则概形，其维数为 $2$。
设 $i : D \to X$ 为一个有效 Cartier 除子的浸入。
设 $\mathcal{F} \to \mathcal{F}' \to i_*\mathcal{G} \to 0$
为凝聚 $\mathcal{O}_X$-模的正合列。假设
\begin{enumerate}
\item $\mathcal{F}, \mathcal{F}'$ 是秩为 $r$ 的局部自由模，这在
$X$ 的一个非空开集上成立；
\item $D$ 是整概形；
\item $\mathcal{G}$ 是有限局部自由 $\mathcal{O}_D$-模，秩为 $s$。
\end{enumerate}
则 $\mathcal{L} = (\wedge^r\mathcal{F})^{**}$ 与
$\mathcal{L}' = (\wedge^r \mathcal{F}')^{**}$ 是可逆
$\mathcal{O}_X$-模，且 $\mathcal{L}' \cong \mathcal{L}(k D)$，其中
$k \in \{0, \ldots, \min(s, r)\}$。
\end{lemma}

\begin{proof}
第一个断言由引理 \ref{divisors-lemma-reflexive-over-regular-dim-2} 得出，
因为假设 (1) 意味着 $\mathcal{L}$ 与 $\mathcal{L}'$ 的秩都为 $1$。
取 $\wedge^r$ 和取双对偶都是函子，因此得到一个典范映射
$\sigma : \mathcal{L} \to \mathcal{L}'$。它在 (1) 中的非空开集上是同构，
因而非零。为了完成证明，只须证明：将 $\sigma$ 视为
$\mathcal{L}' \otimes \mathcal{L}^{\otimes -1}$ 的整体截面时，
它在 $X$ 的任意余维点处都不消失，除了
$D$ 的泛点；而在该点处，其消失阶至多为 $\min(s, r)$。

\medskip\noindent
翻译成代数语言，我们得到如下问题：
设 $(A, \mathfrak m, \kappa)$ 为离散赋值环，分式域为 $K$。
设 $M \to M' \to N \to 0$ 为有限 $A$-模的正合列，满足
$\dim_K(M \otimes K) = \dim_K(M' \otimes K) = r$，并且
$N \cong \kappa^{\oplus s}$。证明诱导映射
$L = \wedge^r(M)^{**} \to L' = \wedge^r(M')^{**}$ 的消失阶
至多为 $\min(s, r)$。我们将使用 $A$ 上模的结构定理，参见
《代数进阶》引理
\ref{more-algebra-lemma-generalized-valuation-ring-modules} 或
\ref{more-algebra-lemma-modules-PID}。
有限 $A$-模对其一个挠子模取商，不会改变其双对偶。
因此，可以将 $M$ 替换为 $M/M_{tors}$，并将 $M'$
替换为 $M'/\Im(M_{tors} \to M')$，从而假设 $M$ 无挠。
此时 $M \to M'$ 是单射，且 $M'_{tors} \to N$ 是单射。
因此，可以将 $M'$ 替换为 $M'/M'_{tors}$，并将 $N$
替换为 $N/M'_{tors}$。于是问题归约到 $M$ 和 $M'$ 都是秩为 $r$
的自由模且 $N \cong \kappa^{\oplus s}$ 的情形。此时 $\sigma$
是 $M \to M'$ 的行列式，并且其消失阶为 $s$；例如这可由
《代数》引理 \ref{algebra-lemma-order-vanishing-determinant} 得出。
\end{proof}

\section{仿射开集的补集}
\label{divisors-section-complement-affine-open}

\noindent
本节讨论如下结果：簇中一个仿射开集的补集具有纯余维 $1$。

\begin{lemma}
\label{divisors-lemma-affine-punctured-spec}
设 $(A, \mathfrak m)$ 为诺特局部环。
记 $U = \Spec(A) \setminus \{\mathfrak m\}$；这是 $A$ 的穿孔谱，并且它是仿射的，
当且仅当 $\dim(A) \leq 1$。
\end{lemma}

\begin{proof}
若 $\dim(A) = 0$，则 $U$ 为空，因而仿射（它等于 $0$ 环的谱）。
若 $\dim(A) = 1$，则可以选择元素 $f \in \mathfrak m$，使其不包含于
$A$ 的有限多个极小素理想中的任何一个
（《代数》引理 \ref{algebra-lemma-Noetherian-irreducible-components} 与
\ref{algebra-lemma-silly}）。于是 $U = \Spec(A_f)$ 是仿射的。

\medskip\noindent
逆命题更有意思。我们将给出一个稍微非标准的证明，并在下面的评注中讨论
标准论证。假设 $U = \Spec(B)$ 是仿射的。由于取约化不影响仿射性与维数，
可以用 $A$ 对其幂零元理想所取的商替换原环，从而假设 $A$ 约化。
令 $Q = B/A$，将其视为 $A$-模。$Q$ 的支集是
$\{\mathfrak m\}$，因为有 $A_\mathfrak p = B_\mathfrak p$，这对每个非极大素理想
$\mathfrak p$（属于 $A$）都成立。
可以假设 $\dim(A) \geq 1$；于是如上可选取 $f \in \mathfrak m$，
使其不包含于 $A$ 的任何极小理想。由于 $A$ 约化，这意味着 $f$ 在 $A$
中是非零因子。同一论证表明 $f : B \to B$ 是单射，因为 $\Spec(B)$
是 $\Spec(A)$ 的开集。注意，由《代数》引理
\ref{algebra-lemma-one-equation}，有
$\dim(A/fA) = \dim(A) - 1$。
对乘以 $f$ 的映射应用蛇形引理，其中该映射作用于短正合列
$0 \to A \to B \to Q \to 0$，
得到
$$
0 \to Q[f] \to A/fA \to B/fB \to Q/fQ \to 0
$$
其中 $Q[f] = \Ker(f : Q \to Q)$。
这说明 $Q[f]$ 是有限 $A$-模。由于 $Q[f]$ 的支集是
$\{\mathfrak m\}$，可知
$l = \text{length}_A(Q[f]) < \infty$
（《代数》引理 \ref{algebra-lemma-support-point}）。
令 $l_n = \text{length}_A(Q[f^n])$。正合列
$$
0 \to Q[f^n] \to Q[f^{n + 1}] \xrightarrow{f^n} Q[f]
$$
归纳地表明 $l_n < \infty$，且 $l_n \leq l_{n + 1}$。
考察正合列
$$
0 \to Q[f] \to Q[f^{n + 1}] \xrightarrow{f} Q[f^n] \to Q/fQ
$$
可知 $Q[f^n]$ 在 $Q/fQ$ 中之像的长度为
$l_n - l_{n + 1} + l \leq l$。由于 $Q = \bigcup Q[f^n]$，
可知 $Q/fQ$ 的长度至多为 $l$，即有界。
因此 $Q/fQ$ 是有限 $A$-模。故 $A/fA \to B/fB$ 是有限环同态，
特别地，它在谱上诱导闭映射
（《代数》引理 \ref{algebra-lemma-integral-going-up} 与
\ref{algebra-lemma-going-up-closed}）。
另一方面，$\Spec(B/fB)$ 是 $\Spec(A/fA)$ 的穿孔谱。
除非 $\Spec(B/fB) = \emptyset$，否则就产生矛盾；而该等式意味着
$\dim(A/fA) = 0$，正是所需结论。
\end{proof}

\begin{remark}
\label{divisors-remark-affine-punctured-spectrum-standard-proof}
设 $(A, \mathfrak m)$ 为维数 $\geq 2$ 的诺特局部正规整环，并设
$U$ 为 $A$ 的穿孔谱，则 $\Gamma(U, \mathcal{O}_U) = A$。
Hartogs 定理的这一代数版本来自如下事实：
$A = \bigcap_{\text{高度}(\mathfrak p) = 1} A_\mathfrak p$，
我们已在《代数》引理
\ref{algebra-lemma-normal-domain-intersection-localizations-height-1}
中见过它。因此在该情形中，$U$ 不可能仿射（否则会迫使 $\mathfrak m$
成为 $U$ 的一个点）。这常被用作引理
\ref{divisors-lemma-affine-punctured-spec} 之证明的出发点。
为把一般诺特局部环的情形约化到该情形，首先作完备化
（以得到 Nagata 局部环），然后对某个适当的极小素理想将 $A$
替换为 $A/\mathfrak q$，再作正规化。这些步骤都不改变维数，并由此得到矛盾。
也可以跳过完备化步骤，但此时正规化一般不是诺特整环。不过，它仍然是
同维数的 Krull 整环（这利用 Krull--Akizuki 定理证明），因而可以应用
同一论证。
\end{remark}

\begin{remark}
\label{divisors-remark-affine-puctured-spectrum-general}
尚不清楚如何刻画穿孔谱为仿射的非诺特局部环 $(A, \mathfrak m)$。
这样的环有一个有限生成理想 $I$，满足
$\mathfrak m = \sqrt{I}$。当然，若可以取 $I$ 由 $1$ 个元素生成，
则 $A$ 的穿孔谱是仿射的；这给出许多非诺特例子。
反过来，由引理 \ref{divisors-lemma-affine-punctured-spec} 之证明中的论证可知，
这样的环不可能有非零因子 $f \in \mathfrak m$ 满足
$H^0_I(A/fA) = 0$（所以 $A$ 不可能有长度为 $2$ 的正则列）。
此外，对任何环 $A'$，若它是局部环同态 $A \to A'$ 的目标环，只要
$\mathfrak m_{A'} = \sqrt{\mathfrak mA'}$，同一结论也成立。
\end{remark}

\begin{lemma}
\label{divisors-lemma-complement-affine-open-immersion}
\begin{reference}
\cite[EGA IV, Corollaire 21.12.7]{EGA4}
\end{reference}
设 $X$ 为局部诺特概形。设 $U \subset X$ 为开子概形，并且包含态射
$U \to X$ 是仿射的。对每个泛点 $\xi$，若它来自
$X \setminus U$ 的某个不可约分支，则局部环 $\mathcal{O}_{X, \xi}$ 的维数
$\leq 1$。若 $U$ 稠密，或 $\xi$ 位于 $U$ 的闭包中，则
$\dim(\mathcal{O}_{X, \xi}) = 1$。
\end{lemma}

\begin{proof}
由于 $\xi$ 是 $X \setminus U$ 的泛点，可知
$$
U_\xi = U \times_X \Spec(\mathcal{O}_{X, \xi}) \subset
\Spec(\mathcal{O}_{X, \xi})
$$
是 $\mathcal{O}_{X, \xi}$ 的穿孔谱（提示：使用《概形》引理
\ref{schemes-lemma-specialize-points}）。
因为 $U \to X$ 是仿射的，故
$U_\xi \to \Spec(\mathcal{O}_{X, \xi})$ 是仿射的
（《态射》引理 \ref{morphisms-lemma-base-change-affine}），
从而 $U_\xi$ 是仿射的。因此由引理
\ref{divisors-lemma-affine-punctured-spec}，有
$\dim(\mathcal{O}_{X, \xi}) \leq 1$。
若 $\xi \in \overline{U}$，则存在特化 $\eta \to \xi$，其中
$\eta \in U$（只需取泛点 $\eta$，它属于 $\overline{U}$ 的一个包含
$\xi$ 的不可约分支；由于 $\overline{U}$ 局部诺特，因而局部只有有限多个
不可约分支，可知 $\eta \in U$）。于是
$\eta \in \Spec(\mathcal{O}_{X, \xi})$，故维数不可能为 $0$。
\end{proof}

\begin{lemma}
\label{divisors-lemma-complement-affine-open}
设 $X$ 为分离的局部诺特概形。设 $U \subset X$ 为仿射开集。
对每个泛点 $\xi$，若它来自 $X \setminus U$ 的某个不可约分支，则局部环
$\mathcal{O}_{X, \xi}$ 的维数均 $\leq 1$。若 $U$ 稠密，或
$\xi$ 位于 $U$ 的闭包中，则
$\dim(\mathcal{O}_{X, \xi}) = 1$。
\end{lemma}

\begin{proof}
这由引理 \ref{divisors-lemma-complement-affine-open-immersion} 得出，因为
由《态射》引理 \ref{morphisms-lemma-affine-permanence}，
态射 $U \to X$ 是仿射的。
\end{proof}

\noindent
下面的引理有时可以用来构造有效 Cartier 除子。

\begin{lemma}
\label{divisors-lemma-complement-open-affine-effective-cartier-divisor}
设 $X$ 为诺特分离概形。设 $U \subset X$ 为稠密仿射开集。
若 $\mathcal{O}_{X, x}$ 是 UFD，且这对所有 $x \in X \setminus U$ 都成立，
则存在有效 Cartier 除子 $D \subset X$，使得 $U = X \setminus D$。
\end{lemma}

\begin{proof}
由于 $X$ 诺特，补集 $X \setminus U$ 只有有限多个不可约分支
$D_1, \ldots, D_r$
（将《性质》引理 \ref{properties-lemma-Noetherian-irreducible-components}
应用于 $X \setminus U$ 上的既约诱导子概形结构）。
由引理 \ref{divisors-lemma-complement-affine-open}
（以及《性质》引理 \ref{properties-lemma-codimension-local-ring}），
每个 $D_i \subset X$ 的余维均为 $1$。
于是由引理 \ref{divisors-lemma-weil-divisor-is-cartier-UFD}，
$D_i$ 是有效 Cartier 除子。因此可以取
$D = D_1 + \ldots + D_r$。
\end{proof}

\begin{lemma}
\label{divisors-lemma-complement-open-affine-effective-cartier-divisor-bis}
设 $X$ 为具有仿射对角的诺特概形。设 $U \subset X$ 为稠密仿射开集。
若 $\mathcal{O}_{X, x}$ 是 UFD，且这对所有 $x \in X \setminus U$ 都成立，
则存在有效 Cartier 除子 $D \subset X$，使得 $U = X \setminus D$。
\end{lemma}

\begin{proof}
由于 $X$ 诺特，补集 $X \setminus U$ 只有有限多个不可约分支
$D_1, \ldots, D_r$
（将《性质》引理 \ref{properties-lemma-Noetherian-irreducible-components}
应用于 $X \setminus U$ 上的既约诱导子概形结构）。
将 $D_i$ 视为 $X$ 的既约闭子概形。
设 $X = \bigcup_{j \in J} X_j$ 为 $X$ 的仿射开覆盖。对
所有 $j$（其属于 $J$），令 $U_j = U \cap X_j$。由于 $X$ 具有仿射对角，
概形
$$
U_j = X \times_{(X \times X)} (U \times X_j)
$$
是仿射的。因此，由于 $X_j$ 分离，根据引理
\ref{divisors-lemma-complement-open-affine-effective-cartier-divisor}
及其证明，对所有 $j \in J$ 与 $1 \leq i \leq r$，交
$D_i \cap X_j$ 或为空，或为 $X_j$ 中的有效 Cartier 除子。
因此 $D_i \subset X$ 是有效 Cartier 除子（因为这是局部性质）。
故可以取 $D = D_1 + \ldots + D_r$。
\end{proof}

\begin{lemma}
\label{divisors-lemma-regular-ample-family}
设 $X$ 为具有仿射对角的拟紧正则概形。则 $X$ 具有一个可逆模的丰富族
（《态射》定义
\ref{morphisms-definition-family-ample-invertible-modules}）。
\end{lemma}

\begin{proof}
注意 $X$ 是整概形的有限不交并
（《性质》引理
\ref{properties-lemma-regular-normal} 与
\ref{properties-lemma-normal-Noetherian}）。
因此可以进一步假设 $X$ 是整的、诺特的、正则的，并且具有仿射对角。
设 $x \in X$。选择仿射开邻域 $U \subset X$，它包含 $x$。
由于 $X$ 是整概形，$U$ 在 $X$ 中稠密。
由《代数进阶》引理 \ref{more-algebra-lemma-regular-local-UFD}，
$X$ 的局部环都是 UFD。因此由引理
\ref{divisors-lemma-complement-open-affine-effective-cartier-divisor-bis}，
可以找到有效 Cartier 除子 $D \subset X$，其补集为 $U$。
于是典范截面 $s = 1_D$ 属于 $\mathcal{L} = \mathcal{O}_X(D)$，
见定义 \ref{divisors-definition-invertible-sheaf-effective-Cartier-divisor}，
恰在 $D$ 上消失，故 $U = X_s$。
因此《态射》定义
\ref{morphisms-definition-family-ample-invertible-modules}
中的两个条件都成立，证毕。
\end{proof}

\section{范数}
\label{divisors-section-norms}

\noindent
设 $\pi : X \to Y$ 为概形的有限态射，并设 $d \geq 1$ 为整数。
如果存在一个乘法映射，则称存在{\it 次数为 $d$ 的 $\pi$ 的范数}
\footnote{这是非标准记号。}
$$
\text{Norm}_\pi : \pi_*\mathcal{O}_X \to \mathcal{O}_Y
$$
它是层的映射，并且满足
\begin{enumerate}
\item 复合
$\mathcal{O}_Y \xrightarrow{\pi^\sharp} \pi_*\mathcal{O}_X
\xrightarrow{\text{Norm}_\pi} \mathcal{O}_Y$ 等于 $g \mapsto g^d$；
\item 对开集 $V \subset Y$，若 $f \in \mathcal{O}_X(\pi^{-1}V)$
在 $x \in \pi^{-1}(V)$ 处为零，则 $\text{Norm}_\pi(f)$
在 $\pi(x)$ 处为零。
\end{enumerate}
注意条件 (1) 迫使 $\pi$ 满射。由于 $\text{Norm}_\pi$ 是乘法的，
它把单位映到单位；因此，给定 $y \in Y$，若 $f$ 是 $X$ 上的正则函数，
它在每个 $x \in X$ 处有定义但不消失，并且 $\pi(x) = y$，则
$\text{Norm}_\pi(f)$ 有定义且在 $y$ 处不消失。该结论不需要条件 (2)；
事实上，本节中的构造只需要条件 (1)，只有某些消失性质
（特别是在引理 \ref{divisors-lemma-norm-ample} 的证明中使用）才需要性质 (2)。

\begin{lemma}
\label{divisors-lemma-finite-trivialize-invertible-upstairs}
设 $\pi : X \to Y$ 为概形的有限态射。设 $\mathcal{L}$ 为可逆
$\mathcal{O}_X$-模。设 $y \in Y$。则存在开邻域
$V \subset Y$，它包含 $y$，并使得 $\mathcal{L}|_{\pi^{-1}(V)}$ 平凡。
\end{lemma}

\begin{proof}
显然可以假设 $Y$ 仿射，因而 $X$ 也仿射。由于 $\pi$ 有限，
纤维 $\pi^{-1}(\{y\})$（位于 $y$ 上）是有限的。
因为 $X$ 仿射，可以选取 $s \in \Gamma(X, \mathcal{L})$，使其在
$\pi^{-1}(\{y\})$ 的任何点处都不消失。这由《性质》引理
\ref{properties-lemma-quasi-affine-invertible-nonvanishing-section}
得出，不过我们也给出直接论证。具体地，可以选取闭点的有限集
$E \subset X$，使每个 $x \in \pi^{-1}(\{y\})$ 都特化到
$E$ 的某个点。对 $x \in E$，用 $i_x : x \to X$ 表示闭浸入。
于是
$\mathcal{L} \to \bigoplus_{x \in E} i_{x, *}i_x^*\mathcal{L}$
是拟凝聚 $\mathcal{O}_X$-模的满射，故映射
$$
\Gamma(X, \mathcal{L}) \to
\bigoplus\nolimits_{x \in E} \mathcal{L}_x/\mathfrak m_x\mathcal{L}_x
$$
是满射（因为在拟凝聚 $\mathcal{O}_X$-模范畴上取整体截面是正合函子，
见《概形》引理 \ref{schemes-lemma-equivalence-quasi-coherent}）。
因此，可以找到 $s \in \Gamma(X, \mathcal{L})$，使其在特化到
$E$ 中某点的任何点处都不消失。于是 $X_s \subset X$ 是
$\pi^{-1}(\{y\})$ 的开邻域。由于 $\pi$ 有限，因而是闭态射，
可知存在开邻域 $V \subset Y$，它包含 $y$，且其逆像包含于 $X_s$，正是所需。
\end{proof}

\begin{lemma}
\label{divisors-lemma-norm-invertible}
设 $\pi : X \to Y$ 为概形的有限态射。若存在次数为 $d$ 的 $\pi$ 的范数，
则存在阿贝尔群同态
$$
\text{Norm}_\pi : \Pic(X) \to \Pic(Y)
$$
使得 $\text{Norm}_\pi(\pi^*\mathcal{N}) \cong \mathcal{N}^{\otimes d}$
对所有可逆 $\mathcal{O}_Y$-模 $\mathcal{N}$ 都成立。
\end{lemma}

\begin{proof}
我们将使用可逆 $\mathcal{O}_X$-模的同构类与
$H^1(X, \mathcal{O}_X^*)$ 的元素之间的对应，而不再另作说明；
该对应见《上同调》引理 \ref{cohomology-lemma-h1-invertible}。
现说明如何对可逆 $\mathcal{O}_X$-模 $\mathcal{L}$ 取范数。
由引理 \ref{divisors-lemma-finite-trivialize-invertible-upstairs}，存在开覆盖
$Y = \bigcup V_j$，使得 $\mathcal{L}|_{\pi^{-1}V_j}$ 平凡。
选择一个生成截面 $s_j \in \mathcal{L}(\pi^{-1}V_j)$，对每个 $j$ 都如此。
在交集 $\pi^{-1}V_j \cap \pi^{-1}V_{j'}$ 上，可以写成
$$
s_j = u_{jj'} s_{j'}
$$
其中存在唯一的
$u_{jj'} \in \mathcal{O}^*_X(\pi^{-1}V_j \cap \pi^{-1}V_{j'})$。
于是可以考察元素
$$
v_{jj'} = \text{Norm}_\pi(u_{jj'}) \in \mathcal{O}_Y^*(V_j \cap V_{j'})
$$
这些元素满足余循环条件（因为 $u_{jj'}$ 满足该条件，且
$\text{Norm}_\pi$ 是乘法的），因此定义了一个可逆
$\mathcal{O}_Y$-模。我们略去如下验证：该构造是良定义的，在 Picard 群上
可加，并且满足
$\text{Norm}_\pi(\pi^*\mathcal{N}) \cong \mathcal{N}^{\otimes d}$，
这对所有可逆 $\mathcal{O}_Y$-模 $\mathcal{N}$ 都成立。
\end{proof}

\begin{lemma}
\label{divisors-lemma-norm-map-invertible}
设 $\pi : X \to Y$ 为概形的有限态射。假设存在次数为 $d$ 的 $\pi$ 的范数。
对任意 $\mathcal{O}_X$-线性映射
$\varphi : \mathcal{L} \to \mathcal{L}'$，若它连接两个可逆
$\mathcal{O}_X$-模，则存在 $\mathcal{O}_Y$-线性映射
$$
\text{Norm}_\pi(\varphi) :
\text{Norm}_\pi(\mathcal{L})
\longrightarrow
\text{Norm}_\pi(\mathcal{L}')
$$
其中 $\text{Norm}_\pi(\mathcal{L})$、$\text{Norm}_\pi(\mathcal{L}')$
如引理 \ref{divisors-lemma-norm-invertible} 所述。此外，对 $y \in Y$，下列条件等价：
\begin{enumerate}
\item $\varphi$ 在某个点 $x \in X$ 处为零，且 $\pi(x) = y$；
\item $\text{Norm}_\pi(\varphi)$ 在 $y$ 处为零。
\end{enumerate}
\end{lemma}

\begin{proof}
选择开覆盖 $Y = \bigcup V_j$，使 $\mathcal{L}$ 和 $\mathcal{L}'$
在开集 $\pi^{-1}V_j$ 上平凡。由引理
\ref{divisors-lemma-finite-trivialize-invertible-upstairs}，这可以做到。
选择生成截面 $s_j$ 与 $s'_j$；它们分别是 $\mathcal{L}$ 与
$\mathcal{L}'$ 在开集 $\pi^{-1}V_j$ 上的生成截面。于是
$\varphi(s_j) = f_js'_j$，其中
$f_j \in \mathcal{O}_X(\pi^{-1}V_j)$。
把 $\text{Norm}_\pi(\varphi)$ 定义为乘以
$\text{Norm}_\pi(f_j)$；此定义在 $V_j$ 上给出。对引理
\ref{divisors-lemma-norm-invertible} 之证明中
用来构造 $\text{Norm}_\pi(\mathcal{L})$、$\text{Norm}_\pi(\mathcal{L}')$
的余循环作简单计算，即可表明这定义了引理所述映射。
最后的断言来自次数为 $d$ 的范数映射之定义中的条件 (2)。
略去一些细节。
\end{proof}

\begin{lemma}
\label{divisors-lemma-norm-ample}
设 $\pi : X \to Y$ 为概形的有限态射。假设 $X$ 具有丰富可逆层，
并且存在次数为 $d$ 的 $\pi$ 的范数。则 $Y$ 具有丰富可逆层。
\end{lemma}

\begin{proof}
设 $\mathcal{L}$ 为假设所给的 $X$ 上的丰富可逆层。我们将证明
$\mathcal{N} = \text{Norm}_\pi(\mathcal{L})$ 在 $Y$ 上丰富。

\medskip\noindent
由于 $X$ 拟紧（《性质》定义
\ref{properties-definition-ample}），而 $X \to Y$ 满射
（因为 $\text{Norm}_\pi$ 存在），可知 $Y$ 拟紧。
取一点 $y \in Y$。为完成证明，我们将证明：存在截面 $t$，它属于
$\mathcal{N}$ 的某个正张量幂，在 $y$ 处不消失，并且 $Y_t$ 仿射。
为此，选择仿射开邻域 $V \subset Y$，它包含 $y$。
选择 $n \gg 0$ 以及截面
$s \in \Gamma(X, \mathcal{L}^{\otimes n})$，使得
$$
\pi^{-1}(\{y\}) \subset X_s \subset \pi^{-1}V
$$
这可由《性质》引理
\ref{properties-lemma-ample-finite-set-in-principal-affine} 做到。
于是 $t = \text{Norm}_\pi(s)$ 是 $\mathcal{N}^{\otimes n}$ 的截面，
它在 $x$ 处不消失，并且 $Y_t \subset V$，见引理
\ref{divisors-lemma-norm-map-invertible}。再由《性质》引理
\ref{properties-lemma-affine-cap-s-open}，$Y_t$ 是仿射的。
\end{proof}

\begin{lemma}
\label{divisors-lemma-norm-quasi-affine}
设 $\pi : X \to Y$ 为概形的有限态射。假设 $X$ 拟仿射，并且存在
次数为 $d$ 的 $\pi$ 的范数。则 $Y$ 拟仿射。
\end{lemma}

\begin{proof}
由《性质》引理 \ref{properties-lemma-quasi-affine-O-ample}，
$\mathcal{O}_X$ 是 $X$ 上的丰富可逆层。
引理 \ref{divisors-lemma-norm-ample} 的证明表明
$\text{Norm}_\pi(\mathcal{O}_X) = \mathcal{O}_Y$
是丰富可逆 $\mathcal{O}_Y$-模。因此《性质》引理
\ref{properties-lemma-quasi-affine-O-ample} 表明 $Y$ 拟仿射。
\end{proof}

\begin{lemma}
\label{divisors-lemma-finite-locally-free-has-norm}
设 $\pi : X \to Y$ 为次数 $d \geq 1$ 的有限局部自由态射。
则存在次数为 $d$ 的典范范数，其构造与任意基变换交换。
\end{lemma}

\begin{proof}
设 $V \subset Y$ 为仿射开集，使得 $(\pi_*\mathcal{O}_X)|_V$
是秩为 $d$ 的有限自由模。选择一组基，得到同构
$$
\mathcal{O}_V^{\oplus d} \cong (\pi_*\mathcal{O}_X)|_V
$$
对每个 $f \in \pi_*\mathcal{O}_X(V) = \mathcal{O}_X(\pi^{-1}(V))$，
乘以 $f$ 定义了所示自由向量丛的 $\mathcal{O}_V$-线性自同态 $m_f$。
因此得到一个 $d \times d$ 矩阵
$M_f \in \text{Mat}(d \times d, \mathcal{O}_Y(V))$，并可定义
$$
\text{Norm}_\pi(f) = \det(M_f)
$$
由于矩阵的行列式与所选基无关，可知该定义良好；这也意味着该构造可以
粘合成所需的整体映射。构造显然与基变换相容。

\medskip\noindent
性质 (1) 来自如下事实：一个 $d \times d$ 对角矩阵的对角元为
$g, g, \ldots, g$ 时，其行列式为 $g^d$。
为证明性质 (2)，可以作基变换，并假设 $Y$ 是域 $k$ 的谱。
于是 $X = \Spec(A)$，其中 $A$ 为 $k$-代数，并且
$\dim_k(A) = d$。若存在 $x \in X$ 使 $f \in A$ 在 $x$ 处消失，
则存在到某个域的映射 $A \to \kappa$，使 $f$ 在 $\kappa$ 中的像为零。
于是 $f : A \to A$ 不可能满射，因而
$\det(f : A \to A) = 0$，正是所需。
\end{proof}

\begin{lemma}
\label{divisors-lemma-norm-in-normal-case}
设 $\pi : X \to Y$ 为有限满态射，其中 $X$ 与 $Y$ 都是整概形，
且 $Y$ 正规。则存在次数为 $[R(X) : R(Y)]$ 的 $\pi$ 的范数。
\end{lemma}

\begin{proof}
设 $\Spec(B) \subset Y$ 为仿射开子集，并设
$\Spec(A) \subset X$ 为其逆像。则 $A$ 与 $B$ 都是整环。
设 $K$ 为 $A$ 的分式域，$L$ 为 $B$ 的分式域。图示如下：
$$
\xymatrix{
L \ar[r] & K \\
B \ar[u] \ar[r] & A \ar[u]
}
$$
由于 $K/L$ 是有限扩张，存在范数映射
$\text{Norm}_{K/L} : K^* \to L^*$，其次数为 $d = [K : L]$；如引理
\ref{divisors-lemma-finite-locally-free-has-norm} 的证明所述，它把
$f \in K$ 映到 $\det_L(f : K \to K)$。
注意 $f : K \to K$ 的特征多项式是 $f$ 在 $L$ 上的极小多项式的一个幂；
特别地，$\text{Norm}_{K/L}(f)$ 是 $f$ 在 $L$ 上之极小多项式的常数项的一个幂。
因此，由《代数》引理
\ref{algebra-lemma-minimal-polynomial-normal-domain}，
$\text{Norm}_{K/L}$ 把 $A$ 映入 $B$。
这在仿射开集的截面上确定了一个相容映射系，从而确定整体范数映射
$\text{Norm}_\pi$，其次数为 $d$。

\medskip\noindent
性质 (1) 由构造立即得出。
为证明性质 (2)，设 $f \in A$ 包含于素理想 $\mathfrak p \subset A$。
设 $f^m + b_1 f^{m - 1} + \ldots + b_m$ 为 $f$ 在 $L$ 上的极小多项式。
由《代数》引理 \ref{algebra-lemma-minimal-polynomial-normal-domain}，
有 $b_i \in B$。因此 $b_0 \in B \cap \mathfrak p$。
由于 $\text{Norm}_{K/L}(f) = b_0^{d/m}$（见上文），可知范数在
$\mathfrak p$ 的像点处消失。
\end{proof}

\begin{lemma}
\label{divisors-lemma-Frobenius-gives-norm-for-reduction}
设 $X$ 为诺特概形。设 $p$ 为满足 $p\mathcal{O}_X = 0$ 的素数。
则对某个 $e > 0$，存在次数为 $p^e$ 的范数，它对应态射 $X_{red} \to X$，
其中 $X_{red}$ 是 $X$ 的约化。
\end{lemma}

\begin{proof}
设 $A$ 为满足 $pA = 0$ 的诺特环。设 $I \subset A$ 为幂零元的理想。
则有 $I^n = 0$，其中 $n$ 为某个整数（《代数》引理
\ref{algebra-lemma-Noetherian-power}）。
选择 $e$ 使 $p^e \geq n$。于是
$$
A/I \longrightarrow A,\quad
f \bmod I \longmapsto f^{p^e}
$$
是良定义的。这给出次数为 $p^e$ 的范数，它对应
$\Spec(A/I) \to \Spec(A)$。
现在，若 $X$ 由某些仿射概形 $\Spec(A_i)$ 粘合而成，则对某个
$e \gg 0$，这些映射粘合成 $X_{red} \to X$ 的范数映射。略去细节。
\end{proof}

\begin{proposition}
\label{divisors-proposition-push-down-ample}
设 $\pi : X \to Y$ 为概形的有限满态射。假设 $X$ 具有丰富可逆
$\mathcal{O}_X$-模。若
\begin{enumerate}
\item $\pi$ 有限局部自由；或
\item $Y$ 为整正规概形；或
\item $Y$ 诺特，$p\mathcal{O}_Y = 0$，并且 $X = Y_{red}$，
\end{enumerate}
则 $Y$ 具有丰富可逆 $\mathcal{O}_Y$-模。
\end{proposition}

\begin{proof}
情形 (1) 由引理 \ref{divisors-lemma-finite-locally-free-has-norm} 与
\ref{divisors-lemma-norm-ample} 结合得出。
情形 (3) 由引理 \ref{divisors-lemma-Frobenius-gives-norm-for-reduction} 与
\ref{divisors-lemma-norm-ample} 结合得出。
在情形 (2) 中，首先替换 $X$；所取替代者是 $X$ 的一个支配 $Y$
的不可约分支（把它视为 $X$ 的既约闭子概形）。
然后可以应用引理 \ref{divisors-lemma-norm-in-normal-case}。
\end{proof}

\begin{lemma}
\label{divisors-lemma-push-down-quasi-affine}
设 $\pi : X \to Y$ 为概形的有限满态射。假设 $X$ 拟仿射。若下列任一条件成立：
\begin{enumerate}
\item $\pi$ 有限局部自由；或
\item $Y$ 为整正规概形，
\end{enumerate}
则 $Y$ 拟仿射。
\end{lemma}

\begin{proof}
情形 (1) 由引理
\ref{divisors-lemma-finite-locally-free-has-norm} 与 \ref{divisors-lemma-norm-quasi-affine}
结合得出。
在情形 (2) 中，首先替换 $X$；所取替代者是 $X$ 的一个支配 $Y$
的不可约分支（把它视为 $X$ 的既约闭子概形）。
然后可以应用引理 \ref{divisors-lemma-norm-in-normal-case}。
\end{proof}

\section{相对有效 Cartier 除子}
\label{divisors-section-effective-Cartier-morphisms}

\noindent
下面的引理表明，在基上平坦的有效 Cartier 除子确实是基上的一个
“有效 Cartier 除子族”。例如，它在任意纤维上的限制都是有效
Cartier 除子。

\begin{lemma}
\label{divisors-lemma-relative-Cartier}
设 $f : X \to S$ 为概形的态射。设 $D \subset X$ 为闭子概形。
假设
\begin{enumerate}
\item $D$ 是有效 Cartier 除子；并且
\item $D \to S$ 是平坦态射。
\end{enumerate}
则对概形的每个态射 $g : S' \to S$，拉回
$(g')^{-1}D$ 都是 $X' = S' \times_S X$ 上的有效 Cartier 除子，
其中 $g' : X' \to X$ 是投影。
\end{lemma}

\begin{proof}
使用引理 \ref{divisors-lemma-characterize-effective-Cartier-divisor}，
可把它翻译成如下代数陈述。设 $A \to B$ 为环同态，且 $h \in B$。
假设 $h$ 是非零因子，并且 $B/hB$ 在 $A$ 上平坦。则
$$
0 \to B \xrightarrow{h} B \to B/hB \to 0
$$
是 $A$-模的短正合列，其中 $B/hB$ 在 $A$ 上平坦。
由《代数》引理 \ref{algebra-lemma-flat-tor-zero}，该列与任意模在
$A$ 上张量后仍然正合，特别地，也可以与任意 $A$-代数 $A'$ 张量。
\end{proof}

\noindent
该引理是下述定义的动机。

\begin{definition}
\label{divisors-definition-relative-effective-Cartier-divisor}
设 $f : X \to S$ 为概形的态射。$X/S$ 上的一个{\it 相对有效 Cartier 除子}
是有效 Cartier 除子 $D \subset X$，使得 $D \to S$ 是概形的平坦态射。
\end{definition}

\noindent
提醒读者，这可能是非标准记号。特别地，在
\cite[IV, Section 21.15]{EGA} 中，只有当 $X \to S$ 平坦且局部有限表示时，
才讨论相对除子的概念。我们的定义稍微更一般。不过，事实证明，在许多情形
（但并非总是如此）下，若 $x \in D$，则 $X \to S$ 在 $x$ 处平坦。

\begin{lemma}
\label{divisors-lemma-sum-relative-effective-Cartier-divisor}
设 $f : X \to S$ 为概形的态射。若 $D_1, D_2 \subset X$ 是
$X/S$ 上的相对有效 Cartier 除子，则 $D_1 + D_2$ 也是
（定义 \ref{divisors-definition-sum-effective-Cartier-divisors}）。
\end{lemma}

\begin{proof}
这翻译成如下代数事实：
设 $A \to B$ 为环同态，且 $h_1, h_2 \in B$。
假设 $h_i$ 都是非零因子，并且 $B/h_iB$ 在 $A$ 上平坦。
则 $h_1h_2$ 是非零因子，并且 $B/h_1h_2B$ 在 $A$ 上平坦。
理由是存在短正合列
$$
0 \to B/h_1B \to B/h_1h_2B \to B/h_2B \to 0
$$
其中第一个箭头由乘以 $h_2$ 给出。由于两端的两项在 $A$ 上平坦，
中间项也平坦，见《代数》引理 \ref{algebra-lemma-flat-ses}。
\end{proof}

\begin{lemma}
\label{divisors-lemma-difference-relative-effective-Cartier-divisor}
设 $f : X \to S$ 为概形的态射。若 $D_1, D_2 \subset X$ 是
$X/S$ 上的相对有效 Cartier 除子，且作为闭子概形有 $D_1 \subset D_2$，
则有效 Cartier 除子 $D$，它满足 $D_2 = D_1 + D$
（引理 \ref{divisors-lemma-difference-effective-Cartier-divisors}）
是 $X/S$ 上的相对有效 Cartier 除子。
\end{lemma}

\begin{proof}
这翻译成如下代数事实：
设 $A \to B$ 为环同态，且 $h_1, h_2 \in B$。
假设 $h_i$ 都是非零因子，$B/h_iB$ 在 $A$ 上平坦，并且
$(h_2) \subset (h_1)$。于是可以写成 $h_2 = h h_1$，
其中 $h \in B$ 是非零因子。得到短正合列
$$
0 \to B/hB \to B/h_2B \to B/h_1B \to 0
$$
其中第一个箭头由乘以 $h_1$ 给出。由于右边两项在 $A$ 上平坦，
中间项也平坦，见《代数》引理 \ref{algebra-lemma-flat-ses}。
\end{proof}

\begin{lemma}
\label{divisors-lemma-flat-at-x}
设 $f : X \to S$ 为概形的态射。设 $D \subset X$ 为
$X/S$ 上的相对有效 Cartier 除子。若 $x \in D$ 且
$\mathcal{O}_{X, x}$ 诺特，则 $f$ 在 $x$ 处平坦。
\end{lemma}

\begin{proof}
令 $A = \mathcal{O}_{S, f(x)}$ 且 $B = \mathcal{O}_{X, x}$。
设 $h \in B$ 为生成 $D$ 之理想的元素。则 $h$ 是 $B$ 中的非零因子，
并且 $B/hB$ 是平坦局部 $A$-代数。设 $I \subset A$ 为有限生成理想。
考察交换图
$$
\xymatrix{
0 \ar[r] &
B \ar[r]_h &
B \ar[r] &
B/hB \ar[r] & 0 \\
0 \ar[r] &
B \otimes_A I \ar[r]^h \ar[u] &
B \otimes_A I \ar[r] \ar[u] &
B/hB \otimes_A I \ar[r] \ar[u] & 0
}
$$
由于 $B/hB$ 在 $A$ 上平坦，下面一列是短正合列，见《代数》引理
\ref{algebra-lemma-flat-tor-zero}。
由于 $B/hB$ 在 $A$ 上平坦，右侧竖直箭头是单射，见《代数》引理
\ref{algebra-lemma-flat}。
因此，乘以 $h$ 在中间竖直箭头的核 $K$ 上是满射。
由 Nakayama 引理，见《代数》引理 \ref{algebra-lemma-NAK}，
可知 $K= 0$。因此 $B$ 在 $A$ 上平坦，见《代数》引理
\ref{algebra-lemma-flat}。
\end{proof}

\noindent
下面的引理依赖平坦点集为开的代数版本。其概形论版本见
《态射进阶》第 \ref{more-morphisms-section-open-flat} 节。

\begin{lemma}
\label{divisors-lemma-flat-relative-Cartier-divisor}
设 $f : X \to S$ 为概形的态射。设 $D \subset X$ 为相对有效
Cartier 除子。若 $f$ 局部有限表示，则存在开子概形 $U \subset X$，
使得 $D \subset U$ 且 $f|_U : U \to S$ 平坦。
\end{lemma}

\begin{proof}
取 $x \in D$。只须找到开邻域 $U \subset X$，它包含 $x$，并使得
$f|_U$ 平坦。因此，该引理约化到如下情形：
$X = \Spec(B)$ 与 $S = \Spec(A)$ 都仿射，并且 $D$ 由非零因子
$h \in B$ 给出。根据假设，$B$ 是有限表示 $A$-代数，而
$B/hB$ 是平坦 $A$-代数。我们将使用绝对诺特逼近。

\medskip\noindent
写成 $B = A[x_1, \ldots, x_n]/(g_1, \ldots, g_m)$。
假设 $h$ 是 $h' \in A[x_1, \ldots, x_n]$ 的像。
选择一个在 $\mathbf{Z}$ 上有限型的子代数 $A_0 \subset A$，使多项式
$h', g_1, \ldots, g_m$ 的所有系数都属于 $A_0$。于是可以令
$B_0 = A_0[x_1, \ldots, x_n]/(g_1, \ldots, g_m)$，并令 $h_0$
为 $h'$ 在 $B_0$ 中的像。此时
$B = B_0 \otimes_{A_0} A$，且
$B/hB = B_0/h_0B_0 \otimes_{A_0} A$。由《代数》引理
\ref{algebra-lemma-flat-finite-presentation-limit-flat}，扩大 $A_0$ 后，
可以假设 $B_0/h_0B_0$ 在 $A_0$ 上平坦。
令 $K_0 = \Ker(h_0 : B_0 \to B_0)$。
由于 $B_0$ 在 $\mathbf{Z}$ 上有限型，可知 $K_0$ 是有限生成理想。
设 $A_1 \subset A$ 为有限型 $\mathbf{Z}$-子代数，它包含 $A_0$，
并用 $B_1$、$h_1$、$K_1$ 表示在 $A_1$ 上的相应对象。
由《代数进阶》引理
\ref{more-algebra-lemma-base-change-H1-regular}，映射
$K_0 \otimes_{A_0} A_1 \to K_1$ 是满射。另一方面，根据假设，
$h : B \to B$ 的核为零。因此，$K_0$ 的每个元素在 $K_1$ 中的像
都为零，这对充分大的子环 $A_1 \subset A$ 成立。由于 $K_0$ 有限生成，
可知 $K_1 = 0$，只须适当选择 $A_1$。

\medskip\noindent
令 $f_1 : X_1 \to S_1$ 等于 $\Spec$，它取自环同态 $A_1 \to B_1$。
令 $D_1 = \Spec(B_1/h_1B_1)$。由于
$B = B_1 \otimes_{A_1} A$，即 $X = X_1 \times_{S_1} S$，
现在只须对 $X_1 \to S_1$ 与相对有效 Cartier 除子 $D_1$
证明该引理，见《态射》引理
\ref{morphisms-lemma-base-change-module-flat}。
因此，问题已约化到 $A$ 为诺特环的情形。在该情形中，由引理
\ref{divisors-lemma-flat-at-x}，环同态 $A \to B$ 在每个素理想
$\mathfrak q$ 处平坦，这些素理想属于 $V(h)$。结合此情形下平坦点集为开的事实，见《代数》定理
\ref{algebra-theorem-openness-flatness}，即得结论。
\end{proof}

\noindent
还有下面的引理（其思想显然归于 Michael Artin，见 \cite{Nobile}），
它完全不需要任何有限性假设。

\begin{lemma}
\label{divisors-lemma-michael-artin}
设 $f : X \to S$ 为概形的态射。设 $D \subset X$ 为
$X/S$ 上的相对有效 Cartier 除子。若 $f$ 在 $X \setminus D$
的所有点处平坦，则 $f$ 平坦。
\end{lemma}

\begin{proof}
这翻译成如下代数事实：
设 $A \to B$ 为环同态，且 $h \in B$。
假设 $h$ 是非零因子，$B/hB$ 在 $A$ 上平坦，并且局部化
$B_h$ 在 $A$ 上平坦。则 $B$ 在 $A$ 上平坦。
理由是存在短正合列
$$
0 \to B \to B_h \to \colim_n (1/h^n)B/B \to 0
$$
并且第二项与第三项在 $A$ 上平坦，这意味着 $B$ 在 $A$ 上平坦
（见《代数》引理 \ref{algebra-lemma-flat-ses}）。
注意，平坦模的滤过余极限是平坦的
（见《代数》引理 \ref{algebra-lemma-colimit-flat}），
而对 $n$ 归纳可知，每个 $(1/h^n)B/B \cong B/h^nB$ 都在 $A$ 上平坦，
因为它嵌入短正合列
$$
0 \to B/h^{n - 1}B \xrightarrow{h} B/h^nB \to B/hB \to 0
$$
略去一些细节。
\end{proof}

\begin{example}
\label{divisors-example-relative-cartier-ambient-space-not-flat}
下面给出一个相对有效 Cartier 除子 $D$ 的例子，其中环境概形在 $D$
的一个邻域上不平坦。具体地，令 $A = k[t]$，并令
$$
B = k[t, x, y, x^{-1}y, x^{-2}y, \ldots]/(ty, tx^{-1}y, tx^{-2}y, \ldots)
$$
则 $B$ 在 $A$ 上不平坦，但 $B/xB \cong A$ 在 $A$ 上平坦。
此外，$x$ 是非零因子，因而在 $\Spec(B)$ 中定义了一个位于
$\Spec(A)$ 上的相对有效 Cartier 除子。
\end{example}

\noindent
若环境概形在基上平坦且局部有限表示，则可以用纤维来刻画相对有效
Cartier 除子。关于该引理稍有不同的表述，另见《态射进阶》引理
\ref{more-morphisms-lemma-slice-given-element}。

\begin{lemma}
\label{divisors-lemma-fibre-Cartier}
设 $\varphi : X \to S$ 为局部有限表示的平坦态射。设
$Z \subset X$ 为闭子概形。设 $x \in Z$，其像为 $s \in S$。
\begin{enumerate}
\item 若 $Z_s \subset X_s$ 在 $x$ 的一个邻域中是 Cartier 除子，
则存在开集 $U \subset X$ 以及相对有效 Cartier 除子 $D \subset U$，
使得 $Z \cap U \subset D$ 且 $Z_s \cap U = D_s$。
\item 若 $Z_s \subset X_s$ 在 $x$ 的一个邻域中是 Cartier 除子，
态射 $Z \to X$ 有限表示，并且 $Z \to S$ 在 $x$ 处平坦，
则可以选择 $U$ 与 $D$ 使得 $Z \cap U = D$。
\item 若 $Z_s \subset X_s$ 在 $x$ 的一个邻域中是 Cartier 除子，
并且 $Z$ 是 $X$ 的局部主闭子概形，这在 $x$ 的一个邻域中成立，
则可以选择 $U$ 与 $D$ 使得 $Z \cap U = D$。
\end{enumerate}
特别地，若 $Z \to S$ 局部有限表示且平坦，并且所有纤维
$Z_s \subset X_s$ 都是有效 Cartier 除子，则 $Z$ 是相对有效
Cartier 除子。类似地，若 $Z$ 是 $X$ 的局部主闭子概形，且所有纤维
$Z_s \subset X_s$ 都是有效 Cartier 除子，则 $Z$ 是相对有效
Cartier 除子。
\end{lemma}

\begin{proof}
选择仿射开邻域 $\Spec(A)$，它包含 $s$；以及仿射开邻域
$\Spec(B)$，它包含 $x$，并使得
$\varphi(\Spec(B)) \subset \Spec(A)$。
设 $\mathfrak p \subset A$ 为对应于 $s$ 的素理想。
设 $\mathfrak q \subset B$ 为对应于 $x$ 的素理想。
设 $I \subset B$ 为对应于 $Z$ 的理想。
由引理的初始假设可知 $A \to B$ 平坦且有限表示。
(1) 中的假设意味着，缩小 $\Spec(B)$ 后，可以假设
$I(B \otimes_A \kappa(\mathfrak p))$ 由一个元素生成，并且该元素在
$B \otimes_A \kappa(\mathfrak p)$ 中是非零因子。
设 $f \in I$ 映到这个生成元。我们断言，在可逆某个
$g \in B$（其中 $g \not \in \mathfrak q$）后，闭子概形
$D = V(f) \subset \Spec(B_g)$ 是相对有效 Cartier 除子。

\medskip\noindent
由《代数》引理
\ref{algebra-lemma-flat-finite-presentation-limit-flat}，
可以找到诺特环之间的平坦有限型环同态 $A_0 \to B_0$、一个元素
$f_0 \in B_0$、一个环同态 $A_0 \to A$ 以及同构
$A \otimes_{A_0} B_0 \cong B$。若
$\mathfrak p_0 = A_0 \cap \mathfrak p$，则可知
$$
B \otimes_A \kappa(\mathfrak p) =
\left(B_0 \otimes_{A_0} \kappa(\mathfrak p_0)\right)
\otimes_{\kappa(\mathfrak p_0))} \kappa(\mathfrak p)
$$
因此，$f_0$ 在 $B_0 \otimes_{A_0} \kappa(\mathfrak p_0)$ 中是非零因子。
由《代数》引理 \ref{algebra-lemma-grothendieck}，可知 $f_0$ 在
$(B_0)_{\mathfrak q_0}$ 中是非零因子，其中
$\mathfrak q_0 = B_0 \cap \mathfrak q$，并且
$(B_0/f_0B_0)_{\mathfrak q_0}$ 在 $A_0$ 上平坦。
因此，由《代数》引理
\ref{algebra-lemma-regular-sequence-in-neighbourhood}
与《代数》定理 \ref{algebra-theorem-openness-flatness}，
存在 $g_0 \in B_0$，其中 $g_0 \not \in \mathfrak q_0$，使得
$f_0$ 在 $(B_0)_{g_0}$ 中是非零因子，并且
$(B_0/f_0B_0)_{g_0}$ 在 $A_0$ 上平坦。因此可知
$D_0 = V(f_0) \subset \Spec((B_0)_{g_0})$ 是相对有效 Cartier 除子。
由于该性质在基变换下保持，见引理
\ref{divisors-lemma-relative-Cartier}，可得上面所述的断言，其中 $g$ 等于
$g_0$ 在 $B$ 中的像。

\medskip\noindent
至此已证明 (1)。为证明 (2)，考察闭浸入 $Z \to D$。
满环同态
$u : \mathcal{O}_{D, x} \to \mathcal{O}_{Z, x}$
是本质有限表示的平坦局部 $\mathcal{O}_{S, s}$-代数之间的映射，
并且模掉 $\mathfrak m_s$ 后成为同构。因此它是同构，见《代数》引理
\ref{algebra-lemma-mod-injective-general}。
由此可知 $Z \to D$ 在 $x$ 的一个邻域中是同构，见《代数》引理
\ref{algebra-lemma-local-isomorphism}。
为证明 (3)，必要时缩小 $U$ 后，可以假设 $D$ 的理想由单个非零因子
$f$ 生成，而 $Z$ 的理想由元素 $g$ 生成。于是 $f = gh$。
但是 $g|_{U_s}$ 与 $f|_{U_s}$ 在 $x$ 的一个邻域中截出同一个
有效 Cartier 除子。因此 $h|_{X_s}$ 是 $\mathcal{O}_{X_s, x}$ 中的单位，
从而 $h$ 是 $\mathcal{O}_{X, x}$ 中的单位，进而 $h$ 在 $x$
的一个开邻域中是单位。也就是说，有 $Z \cap U = D$，这是在缩小 $U$ 后成立。

\medskip\noindent
引理最后的断言立即由 (2)、(3) 以及如下事实得出：
$Z \to S$ 局部有限表示，当且仅当 $Z \to X$ 有限表示，见《态射》引理
\ref{morphisms-lemma-composition-finite-presentation} 与
\ref{morphisms-lemma-finite-presentation-permanence}。
\end{proof}

\section{浸入的法锥}
\label{divisors-section-normal-cone}

\noindent
设 $i : Z \to X$ 为闭浸入。设
$\mathcal{I} \subset \mathcal{O}_X$ 为相应的拟凝聚理想层。
考察分次 $\mathcal{O}_X$-代数的拟凝聚层
$\bigoplus_{n \geq 0} \mathcal{I}^n/\mathcal{I}^{n + 1}$。
由于每个层 $\mathcal{I}^n/\mathcal{I}^{n + 1}$ 都被
$\mathcal{I}$ 零化，由《态射》引理
\ref{morphisms-lemma-i-star-equivalence}，这个分次代数对应于一个
分次 $\mathcal{O}_Z$-代数的拟凝聚层。该拟凝聚分次
$\mathcal{O}_Z$-代数称为{\it $Z$ 在 $X$ 中的余法代数}，并且按照
《态射》第 \ref{morphisms-section-closed-immersions-quasi-coherent} 节
所述的记号滥用，常简记为
$\bigoplus_{n \geq 0} \mathcal{I}^n/\mathcal{I}^{n + 1}$。

\medskip\noindent
设 $f : Z \to X$ 为浸入。把 $f$ 的余法代数定义为闭浸入
$i : Z \to X \setminus \partial Z$ 的余法代数，其中
$\partial Z = \overline{Z} \setminus Z$。它常记为
$\bigoplus_{n \geq 0} \mathcal{I}^n/\mathcal{I}^{n + 1}$，
其中 $\mathcal{I}$ 是闭浸入
$i : Z \to X \setminus \partial Z$ 的理想层。

\begin{definition}
\label{divisors-definition-conormal-sheaf}
设 $f : Z \to X$ 为浸入。所谓{\it 余法代数
$\mathcal{C}_{Z/X, *}$，即 $Z$ 在 $X$ 中的余法代数}，或所谓
{\it $f$ 的余法代数}，是上面所述的分次 $\mathcal{O}_Z$-代数拟凝聚层
$\bigoplus_{n \geq 0} \mathcal{I}^n/\mathcal{I}^{n + 1}$。
\end{definition}

\noindent
因此 $\mathcal{C}_{Z/X, 1} = \mathcal{C}_{Z/X}$ 是该浸入的余法层。
此外，$\mathcal{C}_{Z/X, 0} = \mathcal{O}_Z$，而
$\mathcal{C}_{Z/X, n}$ 是由下述性质刻画的拟凝聚 $\mathcal{O}_Z$-模：
\begin{equation}
\label{divisors-equation-conormal-in-degree-n}
i_*\mathcal{C}_{Z/X, n} = \mathcal{I}^n/\mathcal{I}^{n + 1}
\end{equation}
其中 $i : Z \to X \setminus \partial Z$，且 $\mathcal{I}$ 如上是
$i$ 的理想层。最后，注意存在拟凝聚分次 $\mathcal{O}_Z$-代数的典范满射
\begin{equation}
\label{divisors-equation-conormal-algebra-quotient}
\text{Sym}^*(\mathcal{C}_{Z/X}) \longrightarrow \mathcal{C}_{Z/X, *}
\end{equation}
它在次数 $0$ 与 $1$ 上是同构。

\begin{lemma}
\label{divisors-lemma-affine-conormal-sheaf}
设 $i : Z \to X$ 为浸入。$i$ 的余法代数具有下列性质：
\begin{enumerate}
\item 设 $U \subset X$ 为任意开集，使得 $i(Z)$ 是 $U$ 的闭子集。
设 $\mathcal{I} \subset \mathcal{O}_U$ 为对应于闭子概形
$i(Z) \subset U$ 的理想层。则
$$
\mathcal{C}_{Z/X, *} =
i^*\left(\bigoplus\nolimits_{n \geq 0} \mathcal{I}^n\right) =
i^{-1}\left(
\bigoplus\nolimits_{n \geq 0} \mathcal{I}^n/\mathcal{I}^{n + 1}
\right)
$$
\item 对任意仿射开集 $\Spec(R) = U \subset X$，若
$Z \cap U = \Spec(R/I)$，则存在典范同构
$\Gamma(Z \cap U, \mathcal{C}_{Z/X, *}) = \bigoplus_{n \geq 0} I^n/I^{n + 1}$。
\end{enumerate}
\end{lemma}

\begin{proof}
大多由定义显然可得。注意，给定环 $R$ 以及一个理想 $I$（属于 $R$），有
$I^n/I^{n + 1} = I^n \otimes_R R/I$。略去细节。
\end{proof}

\begin{lemma}
\label{divisors-lemma-conormal-algebra-functorial}
设
$$
\xymatrix{
Z \ar[r]_i \ar[d]_f & X \ar[d]^g \\
Z' \ar[r]^{i'} & X'
}
$$
为概形范畴中的交换图。假设 $i$、$i'$ 都是浸入。存在分次
$\mathcal{O}_Z$-代数的典范映射
$$
f^*\mathcal{C}_{Z'/X', *}
\longrightarrow
\mathcal{C}_{Z/X, *}
$$
它由下列性质刻画：对每一对仿射开集
$(\Spec(R) = U \subset X, \Spec(R') = U' \subset X')$，若
$g(U) \subset U'$，且
$Z \cap U = \Spec(R/I)$、$Z' \cap U' = \Spec(R'/I')$，
则诱导映射
$$
\Gamma(Z' \cap U', \mathcal{C}_{Z'/X', *}) =
\bigoplus\nolimits (I')^n/(I')^{n + 1}
\longrightarrow
\bigoplus\nolimits_{n \geq 0} I^n/I^{n + 1} =
\Gamma(Z \cap U, \mathcal{C}_{Z/X, *})
$$
是由环同态 $f^\sharp : R' \to R$ 诱导的映射，该同态满足
$f^\sharp(I') \subset I$。
\end{lemma}

\begin{proof}
令 $\partial Z' = \overline{Z'} \setminus Z'$ 且
$\partial Z = \overline{Z} \setminus Z$。它们分别是 $X'$ 与 $X$
的闭子集。替换 $X'$，使用 $X' \setminus \partial Z'$；并替换 $X$，
使用 $X \setminus \big(g^{-1}(\partial Z') \cup \partial Z\big)$，
可知可以假设 $i$ 与 $i'$ 都是闭浸入。

\medskip\noindent
{\small
$g \circ i$ 经过 $i'$ 分解这一事实意味着，$g^*\mathcal{I}'$ 映入
$\mathcal{I}$，这是在典范映射 $g^*\mathcal{I}' \to \mathcal{O}_X$ 下成立，见《概形》引理
\ref{schemes-lemma-characterize-closed-subspace} 与
\ref{schemes-lemma-restrict-map-to-closed}。
因此得到拟凝聚层的诱导映射
$g^*((\mathcal{I}')^n/(\mathcal{I}')^{n + 1}) \to
\mathcal{I}^n/\mathcal{I}^{n + 1}$。
再由 $i$ 拉回，得到
$i^*g^*((\mathcal{I}')^n/(\mathcal{I}')^{n + 1}) \to
i^*(\mathcal{I}^n/\mathcal{I}^{n + 1})$。
注意
$i^*(\mathcal{I}^n/\mathcal{I}^{n + 1}) = \mathcal{C}_{Z/X, n}$。
另一方面，
$i^*g^*((\mathcal{I}')^n/(\mathcal{I}')^{n + 1}) =
f^*(i')^*((\mathcal{I}')^n/(\mathcal{I}')^{n + 1}) =
f^*\mathcal{C}_{Z'/X', n}$。
这给出所需映射。\par}

\medskip\noindent
检验该映射局部地由给定映射
$(I')^n/(I')^{n + 1} \to I^n/I^{n + 1}$ 描述，只须展开定义，
故略去。另一个观察是：给定任意 $x \in i(Z)$，确实存在仿射开邻域
$U$、$U'$，满足 $f(U) \subset U'$，并使 $Z \cap U$ 以及
$U' \cap Z'$ 都是闭的，且 $x \in U$。略去证明。因此，引理中的要求
确实刻画了该映射（也本可以用来定义它）。
\end{proof}

\begin{lemma}
\label{divisors-lemma-conormal-algebra-functorial-flat}
设
$$
\xymatrix{
Z \ar[r]_i \ar[d]_f & X \ar[d]^g \\
Z' \ar[r]^{i'} & X'
}
$$
{\small 为概形范畴中的纤维积图，其中 $i$、$i'$ 都是浸入。则引理
\ref{divisors-lemma-conormal-algebra-functorial} 的典范映射
$f^*\mathcal{C}_{Z'/X', *} \to \mathcal{C}_{Z/X, *}$ 是满射。
若 $g$ 平坦，则它是同构。\par}
\end{lemma}

\begin{proof}
设 $R' \to R$ 为环同态，且 $I' \subset R'$ 为理想。
令 $I = I'R$。则
$(I')^n/(I')^{n + 1} \otimes_{R'} R \to I^n/I^{n + 1}$
是满射。若 $R' \to R$ 平坦，则
$I^n = (I')^n \otimes_{R'} R$，从而该映射是同构。
\end{proof}

\begin{definition}
\label{divisors-definition-normal-cone}
设 $i : Z \to X$ 为概形的浸入。所谓{\it 法锥 $C_ZX$，即 $Z$ 在 $X$
中的法锥}，是
$$
C_ZX = \underline{\Spec}_Z(\mathcal{C}_{Z/X, *})
$$
见《构造》定义 \ref{constructions-definition-cone} 与
\ref{constructions-definition-abstract-cone}。$Z$ 在 $X$ 中的{\it 法丛}
是向量丛
$$
N_ZX = \underline{\Spec}_Z(\text{Sym}(\mathcal{C}_{Z/X}))
$$
见《构造》定义 \ref{constructions-definition-vector-bundle} 与
\ref{constructions-definition-abstract-vector-bundle}。
\end{definition}

\noindent
因此 $C_ZX \to Z$ 是 $Z$ 上的锥，$N_ZX \to Z$ 是 $Z$ 上的向量丛
（回忆：在我们的术语中，这并不意味着余法层是有限局部自由层）。
此外，分次代数的典范满射
(\ref{divisors-equation-conormal-algebra-quotient}) 定义了 $Z$ 上锥的典范闭浸入
\begin{equation}
\label{divisors-equation-normal-cone-in-normal-bundle}
C_ZX \longrightarrow N_ZX
\end{equation}

\section{正则理想层}
\label{divisors-section-regular-ideal-sheaves}

\noindent
本节把有效 Cartier 除子的概念推广到更高余维。回忆：元素列
$f_1, \ldots, f_r$（属于环 $R$）称为{\it 正则列}，如果对每个
$i = 1, \ldots, r$，元素 $f_i$ 在
$R/(f_1, \ldots, f_{i - 1})$ 上都是非零因子，并且
$R/(f_1, \ldots, f_r) \not = 0$，见《代数》定义
\ref{algebra-definition-regular-sequence}。
还可以施加三个密切相关的较弱条件。第一个条件是假设
$f_1, \ldots, f_r$ 是{\it Koszul-正则列}，即
$H_i(K_\bullet(f_1, \ldots, f_r)) = 0$ 对所有 $i > 0$ 成立，见《代数进阶》定义
\ref{more-algebra-definition-koszul-regular-sequence}。
称该列为{\it $H_1$-正则列}，如果
$H_1(K_\bullet(f_1, \ldots, f_r)) = 0$。另一个可以施加的条件是：令
$J = (f_1, \ldots, f_r)$，映射
$$
R/J[T_1, \ldots, T_r]
\longrightarrow
\bigoplus\nolimits_{n \geq 0}
J^n/J^{n + 1}
$$
把 $T_i$ 映到 $f_i \bmod J^2$，并且它是同构。在这种情形下，称
$f_1, \ldots, f_r$ 为{\it 拟正则列}，见《代数》定义
\ref{algebra-definition-quasi-regular-sequence}。
给定 $R$-模 $M$，还有 $M$-正则列与 $M$-拟正则列的概念。

\medskip\noindent
可以如下把这些概念推广到带环空间。设 $X$ 为带环空间，并设
$f_1, \ldots, f_r \in \Gamma(X, \mathcal{O}_X)$。
称 $f_1, \ldots, f_r$ 为{\it 正则列}，如果对每个
$i = 1, \ldots, r$，映射
\begin{equation}
\label{divisors-equation-map-regular}
f_i :
\mathcal{O}_X/(f_1, \ldots, f_{i - 1})
\longrightarrow
\mathcal{O}_X/(f_1, \ldots, f_{i - 1})
\end{equation}
都是层的单射。称 $f_1, \ldots, f_r$ 为{\it Koszul-正则列}，如果
Koszul 复形
\begin{equation}
\label{divisors-equation-koszul}
K_\bullet(\mathcal{O}_X, f_\bullet),
\end{equation}
见《模》定义 \ref{modules-definition-koszul-complex}，
在次数 $> 0$ 上无环。称 $f_1, \ldots, f_r$ 为
{\it $H_1$-正则列}，如果 Koszul 复形
$K_\bullet(\mathcal{O}_X, f_\bullet)$ 在次数 $1$ 上正合。
最后，称 $f_1, \ldots, f_r$ 为{\it 拟正则列}，如果映射
\begin{equation}
\label{divisors-equation-map-quasi-regular}
\mathcal{O}_X/\mathcal{J}[T_1, \ldots, T_r]
\longrightarrow
\bigoplus\nolimits_{d \geq 0} \mathcal{J}^d/\mathcal{J}^{d + 1}
\end{equation}
是层的同构，其中 $\mathcal{J} \subset \mathcal{O}_X$ 是由
$f_1, \ldots, f_r$ 生成的理想层。（还有 $\mathcal{F}$-正则列与
$\mathcal{F}$-拟正则列的概念，其中给定的 $\mathcal{O}_X$-模为
$\mathcal{F}$；若以后需要，我们将在那里引入。）

\begin{lemma}
\label{divisors-lemma-types-regular-sequences-implications}
设 $X$ 为带环空间。设
$f_1, \ldots, f_r \in \Gamma(X, \mathcal{O}_X)$。
有下列蕴含：
$f_1, \ldots, f_r$ 是正则列 $\Rightarrow$
$f_1, \ldots, f_r$ 是 Koszul-正则列 $\Rightarrow$
$f_1, \ldots, f_r$ 是 $H_1$-正则列 $\Rightarrow$
$f_1, \ldots, f_r$ 是拟正则列。
\end{lemma}

\begin{proof}
由于正合性可以在茎上检验，列 $f_1, \ldots, f_r$ 是正则列，
当且仅当映射
$$
f_i :
\mathcal{O}_{X, x}/(f_1, \ldots, f_{i - 1})
\longrightarrow
\mathcal{O}_{X, x}/(f_1, \ldots, f_{i - 1})
$$
对所有 $x \in X$ 都是单射。其他类型的正则性也都可以逐茎检验
（略去细节）。因此，这些蕴含由《代数进阶》引理
\ref{more-algebra-lemma-regular-koszul-regular}、
\ref{more-algebra-lemma-koszul-regular-H1-regular} 与
\ref{more-algebra-lemma-H1-regular-quasi-regular} 得出。
\end{proof}

\begin{definition}
\label{divisors-definition-regular-ideal-sheaf}
\begin{reference}
Koszul-正则理想层的概念由
\cite[Expose VII, Definition 1.4]{SGA6} 引入，其中称为正则理想层。
\end{reference}
设 $X$ 为带环空间。设 $\mathcal{J} \subset \mathcal{O}_X$ 为理想层。
\begin{enumerate}
\item 称 $\mathcal{J}$ {\it 正则}，如果对每个
$x \in \text{Supp}(\mathcal{O}_X/\mathcal{J})$，都存在开邻域
$x \in U \subset X$ 以及正则列
$f_1, \ldots, f_r \in \mathcal{O}_X(U)$，使得 $\mathcal{J}|_U$
由 $f_1, \ldots, f_r$ 生成。
\item 称 $\mathcal{J}$ {\it Koszul-正则}，如果对每个
$x \in \text{Supp}(\mathcal{O}_X/\mathcal{J})$，都存在开邻域
$x \in U \subset X$ 以及 Koszul-正则列
$f_1, \ldots, f_r \in \mathcal{O}_X(U)$，使得 $\mathcal{J}|_U$
由 $f_1, \ldots, f_r$ 生成。
\item 称 $\mathcal{J}$ {\it $H_1$-正则}，如果对每个
$x \in \text{Supp}(\mathcal{O}_X/\mathcal{J})$，都存在开邻域
$x \in U \subset X$ 以及 $H_1$-正则列
$f_1, \ldots, f_r \in \mathcal{O}_X(U)$，使得 $\mathcal{J}|_U$
由 $f_1, \ldots, f_r$ 生成。
\item 称 $\mathcal{J}$ {\it 拟正则}，如果对每个
$x \in \text{Supp}(\mathcal{O}_X/\mathcal{J})$，都存在开邻域
$x \in U \subset X$ 以及拟正则列
$f_1, \ldots, f_r \in \mathcal{O}_X(U)$，使得 $\mathcal{J}|_U$
由 $f_1, \ldots, f_r$ 生成。
\end{enumerate}
\end{definition}

\noindent
该概念的许多性质立即由环中正则列与拟正则列的相应概念得出。

\begin{lemma}
\label{divisors-lemma-regular-quasi-regular-scheme}
设 $X$ 为带环空间。设 $\mathcal{J}$ 为理想层。则有下列蕴含：
$\mathcal{J}$ 正则 $\Rightarrow$
$\mathcal{J}$ Koszul-正则 $\Rightarrow$
$\mathcal{J}$ 是 $H_1$-正则 $\Rightarrow$
$\mathcal{J}$ 拟正则。
\end{lemma}

\begin{proof}
该引理立即约化为引理
\ref{divisors-lemma-types-regular-sequences-implications}。
\end{proof}

\begin{lemma}
\label{divisors-lemma-quasi-regular-ideal}
设 $X$ 为局部带环空间。设 $\mathcal{J} \subset \mathcal{O}_X$ 为理想层。
则 $\mathcal{J}$ 拟正则，当且仅当下列条件成立：
\begin{enumerate}
\item $\mathcal{J}$ 是有限型 $\mathcal{O}_X$-模；
\item $\mathcal{J}/\mathcal{J}^2$ 是有限局部自由
$\mathcal{O}_X/\mathcal{J}$-模；并且
\item 典范映射
$$
\text{Sym}^n_{\mathcal{O}_X/\mathcal{J}}(\mathcal{J}/\mathcal{J}^2)
\longrightarrow
\mathcal{J}^n/\mathcal{J}^{n + 1}
$$
对所有 $n \geq 0$ 都是同构。
\end{enumerate}
\end{lemma}

\begin{proof}
显然，若 $U \subset X$ 是开集，并且 $\mathcal{J}|_U$ 由拟正则列
$f_1, \ldots, f_r \in \mathcal{O}_X(U)$ 生成，则 $\mathcal{J}|_U$
有限型，$\mathcal{J}|_U/\mathcal{J}^2|_U$ 自由且具有基
$f_1, \ldots, f_r$，而 (3) 中的映射是同构，因为它们是
(\ref{divisors-equation-map-quasi-regular}) 中次数 $n$ 部分的无坐标表述。
因此显然，拟正则蕴含条件 (1)、(2)、(3)。

\medskip\noindent
反过来，假设 (1)、(2)、(3) 成立。取一点
$x \in \text{Supp}(\mathcal{O}_X/\mathcal{J})$。则存在邻域
$U \subset X$，它包含 $x$，使得 $\mathcal{J}|_U/\mathcal{J}^2|_U$
是秩为 $r$、底环为 $\mathcal{O}_U/\mathcal{J}|_U$ 的自由模。
必要时缩小 $U$ 后，可以假设存在
$f_1, \ldots, f_r \in \mathcal{J}(U)$，它们映为
$\mathcal{J}|_U/\mathcal{J}^2|_U$ 作为
$\mathcal{O}_U/\mathcal{J}|_U$-模的一组基。
特别地，可知 $f_1, \ldots, f_r$ 在
$\mathcal{J}_x/\mathcal{J}^2_x$ 中的像生成该模。因此，由 Nakayama 引理
（《代数》引理 \ref{algebra-lemma-NAK}），可知
$f_1, \ldots, f_r$ 生成茎 $\mathcal{J}_x$。
于是，由于 $\mathcal{J}$ 有限型，根据《模》引理
\ref{modules-lemma-finite-type-surjective-on-stalk}，缩小 $U$ 后，
可以假设 $f_1, \ldots, f_r$ 生成 $\mathcal{J}$。
最后，由 (3) 与同构
$\mathcal{J}|_U/\mathcal{J}^2|_U = \bigoplus \mathcal{O}_U/\mathcal{J}|_U f_i$，显然
$f_1, \ldots, f_r \in \mathcal{O}_X(U)$ 是拟正则列。
\end{proof}

\begin{lemma}
\label{divisors-lemma-generate-regular-ideal}
设 $(X, \mathcal{O}_X)$ 为局部带环空间。设
$\mathcal{J} \subset \mathcal{O}_X$ 为理想层。设 $x \in X$，并设
$f_1, \ldots, f_r \in \mathcal{J}_x$ 的像给出
$\kappa(x)$-向量空间 $\mathcal{J}_x/\mathfrak m_x\mathcal{J}_x$ 的一组基。
\begin{enumerate}
\item 若 $\mathcal{J}$ 拟正则，则存在一个开邻域，使得
$f_1, \ldots, f_r \in \mathcal{O}_X(U)$ 构成生成
$\mathcal{J}|_U$ 的拟正则列。
\item 若 $\mathcal{J}$ 是 $H_1$-正则，则存在一个开邻域，使得
$f_1, \ldots, f_r \in \mathcal{O}_X(U)$ 构成 $H_1$-正则列并生成
$\mathcal{J}|_U$。
\item 若 $\mathcal{J}$ Koszul-正则，则存在一个开邻域，使得
$f_1, \ldots, f_r \in \mathcal{O}_X(U)$ 构成生成
$\mathcal{J}|_U$ 的 Koszul-正则列。
\end{enumerate}
\end{lemma}

\begin{proof}
首先假设 $\mathcal{J}$ 拟正则。可以选择开邻域
$U \subset X$（它包含 $x$）以及拟正则列
$g_1, \ldots, g_s \in \mathcal{O}_X(U)$，后者生成 $\mathcal{J}|_U$。
注意，这意味着 $\mathcal{J}/\mathcal{J}^2$ 是秩为 $s$、底环为
$\mathcal{O}_U/\mathcal{J}|_U$ 的自由模
（见引理 \ref{divisors-lemma-quasi-regular-ideal} 及其证明），因而 $r = s$。
可以缩小 $U$，并假设
$f_1, \ldots, f_r \in \mathcal{J}(U)$。于是可以写成
$$
f_i = \sum a_{ij} g_j
$$
其中 $a_{ij} \in \mathcal{O}_X(U)$。根据假设，矩阵
$A = (a_{ij})$ 映为 $\kappa(x)$ 上的可逆矩阵。
因此，再次缩小 $U$ 后，可以假设 $(a_{ij})$ 可逆。于是可知
$f_1, \ldots, f_r$ 给出 $(\mathcal{J}/\mathcal{J}^2)|_U$ 的一组基，
这证明 $f_1, \ldots, f_r$ 是 $U$ 上的拟正则列。

\medskip\noindent
注意，为证明 (2)、(3)，由于其假设比 (1) 的假设更强，可以直接假设
$f_1, \ldots, f_r \in \mathcal{J}(U)$ 且
$f_i = \sum a_{ij}g_j$，其中 $(a_{ij})$ 如上可逆，而现在
$g_1, \ldots, g_r$ 是 $H_1$-正则列或 Koszul-正则列。
由于 $f_1, \ldots, f_r$ 上的 Koszul 复形同构于
$g_1, \ldots, g_r$ 上的 Koszul 复形，此同构由矩阵 $(a_{ij})$ 给出
（见《代数进阶》引理 \ref{more-algebra-lemma-change-basis}），
可知 $f_1, \ldots, f_r$ 按所需是 $H_1$-正则或 Koszul-正则的。
\end{proof}

\begin{lemma}
\label{divisors-lemma-regular-ideal-sheaf-quasi-coherent}
概形上的任何正则、Koszul-正则、$H_1$-正则或拟正则理想层，
都是有限型拟凝聚理想层。
\end{lemma}

\begin{proof}
这是因为这样的理想层局部地由有限多个截面生成。概形上任何局部地由
截面生成的理想层都是拟凝聚的，见《概形》引理
\ref{schemes-lemma-closed-subspace-scheme}。
\end{proof}

\begin{lemma}
\label{divisors-lemma-regular-ideal-sheaf-scheme}
设 $X$ 为概形。设 $\mathcal{J}$ 为理想层。则 $\mathcal{J}$ 正则
（相应地，Koszul-正则、$H_1$-正则、拟正则），当且仅当对每个
$x \in \text{Supp}(\mathcal{O}_X/\mathcal{J})$，都存在仿射开邻域
$x \in U \subset X$，$U = \Spec(A)$，使得
$\mathcal{J}|_U = \widetilde{I}$，并且 $I$ 由正则列
（相应地，Koszul-正则列、$H_1$-正则列、拟正则列）
$f_1, \ldots, f_r \in A$ 生成。
\end{lemma}

\begin{proof}
根据假设，可以找到开邻域 $U$（它包含 $x$），使得 $\mathcal{J}$ 在其上由正则列
（相应地，Koszul-正则列、$H_1$-正则列、拟正则列）
$f_1, \ldots, f_r \in \mathcal{O}_X(U)$ 生成。缩小 $U$ 后，
可以假设 $U$ 仿射，记为 $U = \Spec(A)$。
由于由引理 \ref{divisors-lemma-regular-ideal-sheaf-quasi-coherent}，
$\mathcal{J}$ 拟凝聚，可知有 $\mathcal{J}|_U = \widetilde{I}$，其中
$I \subset A$ 为某个理想。现在可以使用如下事实：
$$
\widetilde{\ } : \text{Mod}_A \longrightarrow \QCoh(\mathcal{O}_U)
$$
是保持正合性的范畴等价。例如，函数 $f_i$ 生成 $\mathcal{J}$
这一事实，意味着把 $f_i$ 视为 $A$ 的元素时，它们生成 $I$。
映射 (\ref{divisors-equation-map-regular}) 为单射
（相应地，(\ref{divisors-equation-koszul}) 正合，(\ref{divisors-equation-koszul})
在次数 $1$ 上正合，(\ref{divisors-equation-map-quasi-regular}) 为同构）
这一事实，蕴含映射
$A/(f_1, \ldots, f_{i - 1}) \to A/(f_1, \ldots, f_{i - 1})$
（相应地，复形 $K_\bullet(A, f_1, \ldots, f_r)$，映射
$A/I[T_1, \ldots, T_r] \to \bigoplus I^n/I^{n + 1}$）
具有相应性质。因此，$f_1, \ldots, f_r \in A$ 是正则列
（相应地，Koszul-正则列、$H_1$-正则列、拟正则列），所在的环为 $A$。
\end{proof}

\begin{lemma}
\label{divisors-lemma-Noetherian-scheme-regular-ideal}
设 $X$ 为局部诺特概形。设 $\mathcal{J} \subset \mathcal{O}_X$
为拟凝聚理想层。设 $x$ 为 $\mathcal{O}_X/\mathcal{J}$ 之支集中的点。
下列条件等价：
\begin{enumerate}
\item $\mathcal{J}_x$ 在 $\mathcal{O}_{X, x}$ 中由正则列生成；
\item $\mathcal{J}_x$ 在 $\mathcal{O}_{X, x}$ 中由 Koszul-正则列生成；
\item $\mathcal{J}_x$ 由 $H_1$-正则列生成，该列位于
$\mathcal{O}_{X, x}$ 中；
\item $\mathcal{J}_x$ 在 $\mathcal{O}_{X, x}$ 中由拟正则列生成；
\item 存在仿射邻域 $U = \Spec(A)$，它包含 $x$，使得
$\mathcal{J}|_U = \widetilde{I}$，并且 $I$ 在 $A$ 中由正则列生成；
\item 存在仿射邻域 $U = \Spec(A)$，它包含 $x$，使得
$\mathcal{J}|_U = \widetilde{I}$，并且 $I$ 在 $A$ 中由 Koszul-正则列生成；
\item 存在仿射邻域 $U = \Spec(A)$，它包含 $x$，使得
$\mathcal{J}|_U = \widetilde{I}$，并且 $I$ 由 $H_1$-正则列生成，
该列位于 $A$ 中；
\item 存在仿射邻域 $U = \Spec(A)$，它包含 $x$，使得
$\mathcal{J}|_U = \widetilde{I}$，并且 $I$ 在 $A$ 中由拟正则列生成；
\item 存在邻域 $U$（它包含 $x$），使得 $\mathcal{J}|_U$ 正则；
\item 存在邻域 $U$（它包含 $x$），使得 $\mathcal{J}|_U$ Koszul-正则；
\item 存在邻域 $U$（它包含 $x$），使得 $\mathcal{J}|_U$ 是 $H_1$-正则；
\item 存在邻域 $U$（它包含 $x$），使得 $\mathcal{J}|_U$ 拟正则。
\end{enumerate}
特别地，在局部诺特概形上，正则、Koszul-正则、$H_1$-正则与拟正则
理想层的概念全都一致。
\end{lemma}

\begin{proof}
由引理 \ref{divisors-lemma-regular-ideal-sheaf-scheme} 可知
(5) $\Leftrightarrow$ (9)、(6) $\Leftrightarrow$ (10)、
(7) $\Leftrightarrow$ (11)，以及 (8) $\Leftrightarrow$ (12)。
显然，(5) $\Rightarrow$ (1)、(6) $\Rightarrow$ (2)、
(7) $\Rightarrow$ (3)，以及 (8) $\Rightarrow$ (4)。
由《代数》引理 \ref{algebra-lemma-regular-sequence-in-neighbourhood}，
有 (1) $\Rightarrow$ (5)。
由引理 \ref{divisors-lemma-regular-quasi-regular-scheme}，有
(9) $\Rightarrow$ (10) $\Rightarrow$ (11) $\Rightarrow$ (12)。
最后，由《代数》引理 \ref{algebra-lemma-quasi-regular-regular}，
有 (4) $\Rightarrow$ (1)。现在全部 12 个陈述都等价。
\end{proof}

\section{正则浸入}
\label{divisors-section-regular-immersions}

\noindent
设 $i : Z \to X$ 为概形的浸入。按定义，这意味着存在开子概形
$U \subset X$，使得 $Z$ 可与 $U$ 的一个闭子概形等同。令
$\mathcal{I} \subset \mathcal{O}_U$ 为相应的拟凝聚理想层。假设
$U' \subset X$ 是另一个这样的开子概形，并记
$\mathcal{I}' \subset \mathcal{O}_{U'}$ 为相应的拟凝聚理想层。则
$\mathcal{I}|_{U \cap U'} = \mathcal{I}'|_{U \cap U'}$。
此外，$\mathcal{O}_U/\mathcal{I}$ 的支集是 $Z$；它包含于
$U \cap U'$，并且也是 $\mathcal{O}_{U'}/\mathcal{I}'$ 的支集。
因此，由定义 \ref{divisors-definition-regular-ideal-sheaf} 可知，
$\mathcal{I}$ 是正则理想，当且仅当 $\mathcal{I}'$ 是正则理想。
对 Koszul-正则、$H_1$-正则或拟正则也同样成立。

\begin{definition}
\label{divisors-definition-regular-immersion}
\begin{reference}
Koszul-正则浸入的概念由
\cite[Expose VII, Definition 1.4]{SGA6} 引入，其中称为正则浸入。
\end{reference}
设 $i : Z \to X$ 为概形的浸入。选择开子概形
$U \subset X$，使得 $i$ 把 $Z$ 与 $U$ 的一个闭子概形等同，并记
$\mathcal{I} \subset \mathcal{O}_U$ 为相应的拟凝聚理想层。
\begin{enumerate}
\item 称 $i$ 为{\it 正则浸入}，如果 $\mathcal{I}$ 正则。
\item 称 $i$ 为{\it Koszul-正则浸入}，如果 $\mathcal{I}$ Koszul-正则。
\item 称 $i$ 为{\it $H_1$-正则浸入}，如果 $\mathcal{I}$ 是 $H_1$-正则。
\item 称 $i$ 为{\it 拟正则浸入}，如果 $\mathcal{I}$ 拟正则。
\end{enumerate}
\end{definition}

\noindent
上面的讨论表明，这一定义不依赖于 $U$ 的选择。这些条件按强弱递减的
顺序列出，见引理 \ref{divisors-lemma-regular-quasi-regular-immersion}。
Koszul-正则闭浸入经平滑局部化后是正则浸入，见引理
\ref{divisors-lemma-koszul-regular-smooth-locally-regular}。
在局部 Noether 情形，四种概念全都一致，见引理
\ref{divisors-lemma-Noetherian-scheme-regular-ideal}。

\begin{lemma}
\label{divisors-lemma-regular-quasi-regular-immersion}
设 $i : Z \to X$ 为概形的浸入。则有下列蕴含：
$i$ 正则 $\Rightarrow$
$i$ Koszul-正则 $\Rightarrow$
$i$ 是 $H_1$-正则 $\Rightarrow$
$i$ 拟正则。
\end{lemma}

\begin{proof}
该引理立即约化为引理
\ref{divisors-lemma-regular-quasi-regular-scheme}。
\end{proof}

\begin{lemma}
\label{divisors-lemma-regular-immersion-noetherian}
设 $i : Z \to X$ 为概形的浸入。假设 $X$ 局部 Noether。则
$i$ 正则 $\Leftrightarrow$
$i$ Koszul-正则 $\Leftrightarrow$
$i$ 是 $H_1$-正则 $\Leftrightarrow$
$i$ 拟正则。
\end{lemma}

\begin{proof}
这立即由引理
\ref{divisors-lemma-regular-quasi-regular-immersion}
与引理
\ref{divisors-lemma-Noetherian-scheme-regular-ideal}
得出。
\end{proof}

\begin{lemma}
\label{divisors-lemma-flat-base-change-regular-immersion}
设 $i : Z \to X$ 为正则浸入（相应地，Koszul-正则、
$H_1$-正则、拟正则浸入）。设 $X' \to X$ 为平坦态射。则基变换
$i' : Z \times_X X' \to X'$ 是正则浸入（相应地，Koszul-正则、
$H_1$-正则、拟正则浸入）。
\end{lemma}

\begin{proof}
借助引理
\ref{divisors-lemma-regular-ideal-sheaf-scheme}，
这转化为《代数》引理 \ref{algebra-lemma-flat-increases-depth} 与
\ref{algebra-lemma-flat-base-change-quasi-regular}
以及《代数进阶》引理
\ref{more-algebra-lemma-koszul-regular-flat-base-change}
中的代数陈述。
\end{proof}

\begin{lemma}
\label{divisors-lemma-quasi-regular-immersion}
设 $i : Z \to X$ 为概形的浸入。则 $i$ 是拟正则浸入，当且仅当
下列条件成立：
\begin{enumerate}
\item $i$ 局部有限表示；
\item 余法层 $\mathcal{C}_{Z/X}$ 有限局部自由；并且
\item 映射 (\ref{divisors-equation-conormal-algebra-quotient}) 是同构。
\end{enumerate}
\end{lemma}

\begin{proof}
开浸入局部有限表示。因此，可以把 $X$ 替换为开子概形
$U \subset X$，使得 $i$ 把 $Z$ 与 $U$ 的一个闭子概形等同；换言之，可以假设
$i$ 是闭浸入。令 $\mathcal{I} \subset \mathcal{O}_X$ 为相应的
拟凝聚理想层。回忆《态射》引理
\ref{morphisms-lemma-closed-immersion-finite-presentation}：
$\mathcal{I}$ 有限型，当且仅当 $i$ 局部有限表示。因此，所述等价性
由引理 \ref{divisors-lemma-quasi-regular-ideal} 和展开定义得出。
\end{proof}

\begin{lemma}
\label{divisors-lemma-transitivity-conormal-quasi-regular}
设 $Z \to Y \to X$ 为概形的浸入。假设 $Z \to Y$ 是
$H_1$-正则的。则《态射》引理
\ref{morphisms-lemma-transitivity-conormal} 的典范序列
$$
0 \to i^*\mathcal{C}_{Y/X} \to
\mathcal{C}_{Z/X} \to
\mathcal{C}_{Z/Y} \to 0
$$
正合且局部分裂。
\end{lemma}

\begin{proof}
由于 $\mathcal{C}_{Z/Y}$ 有限局部自由（见引理
\ref{divisors-lemma-quasi-regular-immersion}
与引理
\ref{divisors-lemma-regular-quasi-regular-scheme}），
只需证明该序列正合。由《态射》引理
\ref{morphisms-lemma-transitivity-conormal} 中已经证明的结论，
只需说明第一个映射是单射。仿射局部地工作，可把问题约化如下：
假设有环 $A$ 及理想 $I \subset J \subset A$。假设
$J/I \subset A/I$ 由一个 $H_1$-正则列生成。这是否蕴含
$I/I^2 \otimes_A A/J \to J/J^2$ 为单射？
注意 $I/I^2 \otimes_A A/J = I/IJ$。因此，需要证明
$I \cap J^2 = IJ$。这正是《代数进阶》引理
\ref{more-algebra-lemma-conormal-sequence-H1-regular} 的结论。
\end{proof}

\noindent
拟正则浸入的复合未必拟正则，见《代数》评注
\ref{algebra-remark-join-quasi-regular-sequences}。
其他类型的正则浸入在复合下保持。

\begin{lemma}
\label{divisors-lemma-composition-regular-immersion}
设 $i : Z \to Y$ 与 $j : Y \to X$ 为概形的浸入。
\begin{enumerate}
\item 若 $i$ 与 $j$ 都是正则浸入，则 $j \circ i$ 也是。
\item 若 $i$ 与 $j$ 都是 Koszul-正则浸入，则 $j \circ i$ 也是。
\item 若 $i$ 与 $j$ 都是 $H_1$-正则浸入，则 $j \circ i$ 也是。
\item 若 $i$ 是 $H_1$-正则浸入而 $j$ 是拟正则浸入，则
$j \circ i$ 是拟正则浸入。
\end{enumerate}
\end{lemma}

\begin{proof}
(1) 的代数版本是《代数》引理
\ref{algebra-lemma-join-regular-sequences}。
(2) 的代数版本是《代数进阶》引理
\ref{more-algebra-lemma-join-koszul-regular-sequences}。
(3) 的代数版本是《代数进阶》引理
\ref{more-algebra-lemma-join-H1-regular-sequences}。
(4) 的代数版本是《代数进阶》引理
\ref{more-algebra-lemma-join-quasi-regular-H1-regular}。
\end{proof}

\begin{lemma}
\label{divisors-lemma-permanence-regular-immersion}
设 $i : Z \to Y$ 与 $j : Y \to X$ 为概形的浸入。假设
$j$ 局部有限表示，并且《态射》引理
\ref{morphisms-lemma-transitivity-conormal} 的序列
$$
0 \to i^*\mathcal{C}_{Y/X} \to
\mathcal{C}_{Z/X} \to
\mathcal{C}_{Z/Y} \to 0
$$
正合且局部分裂。
\begin{enumerate}
\item 若 $j \circ i$ 是拟正则浸入，则 $i$ 也是。
\item 若 $j \circ i$ 是 $H_1$-正则浸入，则 $i$ 也是。
\item 若 $j$ 与 $j \circ i$ 都是 Koszul-正则浸入，则 $i$ 也是。
\end{enumerate}
\end{lemma}

\begin{proof}
缩小 $Y$ 与 $X$ 后，可以假设 $i$ 与 $j$ 都是闭浸入。记
$\mathcal{I} \subset \mathcal{O}_X$ 为 $Y$ 的理想层，并记
$\mathcal{J} \subset \mathcal{O}_X$ 为 $Z$ 的理想层。
余法序列是 $0 \to \mathcal{I}/\mathcal{I}\mathcal{J}
\to \mathcal{J}/\mathcal{J}^2 \to
\mathcal{J}/(\mathcal{I} + \mathcal{J}^2) \to 0$。
设 $z \in Z$，并置 $y = i(z)$、$x = j(y) = j(i(z))$。
选择 $f_1, \ldots, f_n \in \mathcal{I}_x$，其像构成
$\mathcal{I}_x/\mathfrak m_z\mathcal{I}_x$ 的一组基。把它扩充为
$f_1, \ldots, f_n, g_1, \ldots, g_m \in \mathcal{J}_x$，使其像构成
$\mathcal{J}_x/\mathfrak m_z\mathcal{J}_x$ 的一组基。这是可能的，
因为已经假设余法层的序列在 $z$ 的某个邻域中分裂，从而
$\mathcal{I}_x/\mathfrak m_x\mathcal{I}_x \to
\mathcal{J}_x/\mathfrak m_x\mathcal{J}_x$ 为单射。

\medskip\noindent
{\small (1) 的证明。由引理
\ref{divisors-lemma-generate-regular-ideal}，可以找到仿射开邻域
$U$（它包含 $x$），使得 $f_1, \ldots, f_n, g_1, \ldots, g_m$ 构成生成
$\mathcal{J}$ 的拟正则列。由于 $i$ 局部有限表示，还可以假设
$\mathcal{I}$ 由 $f_1, \ldots, f_n$ 生成（使用《态射》引理
\ref{morphisms-lemma-closed-immersion-finite-presentation}、
Nakayama 引理以及《模》引理
\ref{modules-lemma-finite-type-surjective-on-stalk}）。
因此，由《代数》引理
\ref{algebra-lemma-truncate-quasi-regular} 可知，
$g_1, \ldots, g_m$ 在 $Y \cap U$ 上诱导一个切出 $Z$ 的拟正则列。\par}

\medskip\noindent
(2) 的证明。与 (1) 的证明完全相同，只需改用《代数进阶》引理
\ref{more-algebra-lemma-truncate-H1-regular}。

\medskip\noindent
(3) 的证明。将引理
\ref{divisors-lemma-generate-regular-ideal}
应用两次，可找到仿射开邻域 $U$（它包含 $x$），使得
$f_1, \ldots, f_n$ 构成生成 $\mathcal{I}$ 的 Koszul-正则列，并且
$f_1, \ldots, f_n, g_1, \ldots, g_m$ 构成生成 $\mathcal{J}$ 的
Koszul-正则列。因此，由《代数进阶》引理
\ref{more-algebra-lemma-truncate-koszul-regular} 可知，
$g_1, \ldots, g_m$ 在 $Y \cap U$ 上诱导一个切出 $Z$ 的
Koszul-正则列。
\end{proof}

\begin{lemma}
\label{divisors-lemma-extra-permanence-regular-immersion-noetherian}
设 $i : Z \to Y$ 与 $j : Y \to X$ 为概形的浸入。取 $z \in Z$，
并记 $y \in Y$、$x \in X$ 为相应的点。假设 $X$ 局部 Noether。
下列条件等价：
\begin{enumerate}
\item $i$ 在 $z$ 的一个邻域中是正则浸入，并且 $j$ 在
$y$ 的一个邻域中是正则浸入；
\item $i$ 与 $j \circ i$ 在 $z$ 的一个邻域中都是正则浸入；
\item $j \circ i$ 在 $z$ 的一个邻域中是正则浸入，并且余法序列
$$
0 \to i^*\mathcal{C}_{Y/X} \to
\mathcal{C}_{Z/X} \to
\mathcal{C}_{Z/Y} \to 0
$$
在 $z$ 的一个邻域中分裂正合。
\end{enumerate}
\end{lemma}

\begin{proof}
由于 $X$（因而 $Y$）局部 Noether，四种正则浸入的概念全都一致；
此外，可以在局部环层面检验一个态射是否为正则浸入，见引理
\ref{divisors-lemma-Noetherian-scheme-regular-ideal}。
蕴含 (1) $\Rightarrow$ (2) 是引理
\ref{divisors-lemma-composition-regular-immersion}。
蕴含 (2) $\Rightarrow$ (3) 是引理
\ref{divisors-lemma-transitivity-conormal-quasi-regular}。
因此，只需证明 (3) 蕴含 (1)。

\medskip\noindent
假设 (3)。置 $A = \mathcal{O}_{X, x}$。记 $I \subset A$ 为满射
$\mathcal{O}_{X, x} \to \mathcal{O}_{Y, y}$ 的核，并记
$J \subset A$ 为满射
$\mathcal{O}_{X, x} \to \mathcal{O}_{Z, z}$ 的核。
注意，生成 $J$ 的任意极小元素列（位于 $A$ 中）都是拟正则列，因而是正则列，
见引理 \ref{divisors-lemma-generate-regular-ideal}。由假设，余法序列
$$
0 \to I/IJ \to J/J^2 \to J/(I + J^2) \to 0
$$
作为 $A/J$-模的序列分裂正合。因此，可以选择极小生成系
$f_1, \ldots, f_n, g_1, \ldots, g_m$，它生成 $J$，并使得
$f_1, \ldots, f_n \in I$ 是 $I$ 的一个极小生成系。如上所述，
$f_1, \ldots, f_n, g_1, \ldots, g_m$ 是 $A$ 中的正则列。
由正则列的定义直接可知，$f_1, \ldots, f_n$ 是 $A$ 中的正则列，
而 $\overline{g}_1, \ldots, \overline{g}_m$ 是 $A/I$ 中的正则列。
因此，$j$ 在 $y$ 处是正则浸入，而 $i$ 在 $z$ 处是正则浸入。
\end{proof}

\begin{remark}
\label{divisors-remark-not-always-extra-permanence}
在引理
\ref{divisors-lemma-extra-permanence-regular-immersion-noetherian}
的情形中，(1)、(2)、(3) 与下述条件{\bf 不}等价：
“$j \circ i$ 与 $j$ 分别在 $z$ 与 $y$ 处为正则浸入”。
一个例子是 $X = \mathbf{A}^1_k = \Spec(k[x])$、
$Y = \Spec(k[x]/(x^2))$ 与 $Z = \Spec(k[x]/(x))$。
\end{remark}

\begin{lemma}
\label{divisors-lemma-koszul-regular-smooth-locally-regular}
设 $i : Z \to X$ 为 Koszul 正则闭浸入。则存在满的平滑态射
$X' \to X$，使得基变换
$i' : Z \times_X X' \to X'$（它来自 $i$）是正则浸入。
\end{lemma}

\begin{proof}
可以假设 $X$ 仿射，并且 $Z$ 的理想由 Koszul-正则列生成：为此，
用某个适当仿射开覆盖的成员替换 $X$（仿射开集取引理
\ref{divisors-lemma-regular-ideal-sheaf-scheme} 中的那些）。仿射情形是
《代数进阶》引理
\ref{more-algebra-lemma-Koszul-regular-flat-locally-regular}。
\end{proof}

\begin{lemma}
\label{divisors-lemma-immersion-regular-regular-immersion}
设 $i : Z \to X$ 为浸入。若 $Z$ 与 $X$ 都是正则概形，则 $i$ 是
正则浸入。
\end{lemma}

\begin{proof}
设 $z \in Z$。由引理 \ref{divisors-lemma-Noetherian-scheme-regular-ideal}，
只需证明映射
$\mathcal{O}_{X, z} \to \mathcal{O}_{Z, z}$
的核由正则列生成。这由《代数》引理
\ref{algebra-lemma-regular-quotient-regular} 与
\ref{algebra-lemma-regular-ring-CM}
得出。
\end{proof}





\section{相对正则浸入}
\label{divisors-section-relative-regular-immersion}

\noindent
本节考察正则浸入的基变换性质。下述引理对正则浸入或 Koszul 浸入
并不成立，见《例》引理
\ref{examples-lemma-base-change-regular-sequence}。

\begin{lemma}
\label{divisors-lemma-relative-regular-immersion}
设 $f : X \to S$ 为概形的态射。设 $i : Z \subset X$ 为浸入。假设
\begin{enumerate}
\item $i$ 是 $H_1$-正则浸入（相应地，拟正则浸入）；并且
\item $Z \to S$ 是平坦态射。
\end{enumerate}
则对概形的每个态射 $g : S' \to S$，基变换
$Z' = S' \times_S Z \to X' = S' \times_S X$
是 $H_1$-正则浸入（相应地，拟正则浸入）。
\end{lemma}

\begin{proof}
展开定义并使用引理
\ref{divisors-lemma-regular-ideal-sheaf-scheme}，
可把它转化为《代数进阶》引理
\ref{more-algebra-lemma-relative-regular-immersion-algebra}。
\end{proof}

\noindent
这个引理引出了下面的定义。

\begin{definition}
\label{divisors-definition-relative-H1-regular-immersion}
设 $f : X \to S$ 为概形的态射。设 $i : Z \to X$ 为浸入。
\begin{enumerate}
\item 称 $i$ 为{\it 相对拟正则浸入}，如果 $Z \to S$ 平坦且
$i$ 是拟正则浸入。
\item 称 $i$ 为{\it 相对 $H_1$-正则浸入}，如果 $Z \to S$ 平坦且
$i$ 是 $H_1$-正则浸入。
\end{enumerate}
\end{definition}

\noindent
提醒读者，这可能不是标准记号。引理
\ref{divisors-lemma-relative-regular-immersion}
保证相对拟正则浸入（相应地，$H_1$-正则浸入）在任意基变换下保持。
相对 $H_1$-正则浸入是相对拟正则浸入，见引理
\ref{divisors-lemma-regular-quasi-regular-immersion}。请参见引理
\ref{divisors-lemma-flat-relative-H1-regular}
（或引理
\ref{divisors-lemma-relative-regular-immersion-flat-in-neighbourhood}）：
它表明，若 $Z \to X$ 是相对 $H_1$-正则浸入（或拟正则浸入），
且环境概形在 $S$ 上（平坦并）局部有限表示，则 $Z \to X$
实际上是正则浸入，并且在任意基变换后仍然如此。

\begin{lemma}
\label{divisors-lemma-quasi-regular-immersion-flat-at-x}
设 $f : X \to S$ 为概形的态射。设 $Z \to X$ 为相对拟正则浸入。
若 $x \in Z$ 且 $\mathcal{O}_{X, x}$ 是 Noether 环，则 $f$ 在 $x$ 处平坦。
\end{lemma}

\begin{proof}
设 $f_1, \ldots, f_r \in \mathcal{O}_{X, x}$ 为切出
$Z$ 在 $x$ 处之理想的拟正则列。由《代数》引理
\ref{algebra-lemma-quasi-regular-regular} 可知
$f_1, \ldots, f_r$ 是正则列。因此，$f_r$ 在
$\mathcal{O}_{X, x}/(f_1, \ldots, f_{r - 1})$ 上是非零因子，
并且其商是平坦 $\mathcal{O}_{S, f(x)}$-模。由引理
\ref{divisors-lemma-flat-at-x} 可知
$\mathcal{O}_{X, x}/(f_1, \ldots, f_{r - 1})$
是平坦 $\mathcal{O}_{S, f(x)}$-模。继续归纳，可得
$\mathcal{O}_{X, x}$ 是平坦 $\mathcal{O}_{S, s}$-模。
\end{proof}

\begin{lemma}
\label{divisors-lemma-relative-regular-immersion-flat-in-neighbourhood}
设 $X \to S$ 为概形的态射。设 $Z \to X$ 为浸入。假设
\begin{enumerate}
\item $X \to S$ 平坦且局部有限表示；
\item $Z \to X$ 是相对拟正则浸入。
\end{enumerate}
则 $Z \to X$ 是正则浸入，并且在任意基变换后仍然如此。
\end{lemma}

\begin{proof}
取 $x \in Z$，其像为 $s \in S$。为证明此结论，只需找到 $x$ 的
一个包含于 $U$ 的仿射邻域，使结论在该仿射开集上成立。因此，可以
假设 $X$ 仿射，并且存在拟正则列
$f_1, \ldots, f_r \in \Gamma(X, \mathcal{O}_X)$，使得
$Z = V(f_1, \ldots, f_r)$。由《代数进阶》引理
\ref{more-algebra-lemma-relative-regular-immersion-algebra}，
列 $f_1|_{X_s}, \ldots, f_r|_{X_s}$ 是
$\Gamma(X_s, \mathcal{O}_{X_s})$ 中的拟正则列。由于 $X_s$
Noether，这蕴含在必要时稍微缩小 $X$ 后，
$f_1|_{X_s}, \ldots, f_r|_{X_s}$ 是正则列；见《代数》引理
\ref{algebra-lemma-quasi-regular-regular} 与
\ref{algebra-lemma-regular-sequence-in-neighbourhood}。
由引理
\ref{divisors-lemma-fibre-Cartier} 可知，
$Z_1 = V(f_1) \subset X$ 是相对有效 Cartier 除子；这里必要时再次
稍微缩小 $X$。再次应用同一引理，
这次应用于 $Z_2 = V(f_1, f_2) \subset Z_1$，可知
$Z_2 \subset Z_1$ 是相对有效 Cartier 除子。依此类推，直到到达
$Z = Z_n = V(f_1, \ldots, f_n)$。由于相对有效 Cartier 除子的性质
在任意基变换下保持，见引理
\ref{divisors-lemma-relative-Cartier}，引理的最后一个陈述也随之成立。
\end{proof}

\begin{remark}
\label{divisors-remark-relative-regular-immersion-elements}
相对拟正则浸入的余维若为常数，则在基变换后不变。事实上，若有环映射
$A \to B$ 以及拟正则列 $f_1, \ldots, f_r \in B$，使得
$B/(f_1, \ldots, f_r)$ 在 $A$ 上平坦，则对任意环映射
$A \to A'$，由《代数进阶》引理
\ref{more-algebra-lemma-relative-regular-immersion-algebra}，
$f_1 \otimes 1, \ldots, f_r \otimes 1$ 在
$B' = B \otimes_A A'$ 中是拟正则列
（该引理也用于上面引理
\ref{divisors-lemma-relative-regular-immersion} 的证明）。
现在，引理
\ref{divisors-lemma-relative-regular-immersion-flat-in-neighbourhood}
的证明表明，若 $A \to B$ 平坦且局部有限表示，则对每个素理想
$\mathfrak q' \subset B'$，列
$f_1 \otimes 1, \ldots, f_r \otimes 1$ 甚至是局部环
$B'_{\mathfrak q'}$ 中的正则列。
\end{remark}

\begin{lemma}
\label{divisors-lemma-flat-relative-H1-regular}
设 $X \to S$ 为概形的态射。设 $Z \to X$ 为相对 $H_1$-正则浸入。
假设 $X \to S$ 局部有限表示。则
\begin{enumerate}
\item 存在开子概形 $U \subset X$，使得 $Z \subset U$ 且
$U \to S$ 平坦；并且
\item $Z \to X$ 是正则浸入，并且在任意基变换后仍然如此。
\end{enumerate}
\end{lemma}

\begin{proof}
取 $x \in Z$。为证明 (1)，只需找到开邻域
$U \subset X$，它包含 $x$，使得 $U \to S$ 平坦。因此，引理约化到如下情形：
$X = \Spec(B)$ 与 $S = \Spec(A)$ 仿射，并且 $Z$ 由
$H_1$-正则列 $f_1, \ldots, f_r \in B$ 给出。由假设，
$B$ 是有限表示 $A$-代数，并且
$B/(f_1, \ldots, f_r)B$ 是平坦 $A$-代数。下面使用绝对 Noether 逼近。

\medskip\noindent
写成 $B = A[x_1, \ldots, x_n]/(g_1, \ldots, g_m)$。假设
$f_i$ 是 $f_i' \in A[x_1, \ldots, x_n]$ 的像。选择有限型
$\mathbf{Z}$-子代数 $A_0 \subset A$，使得多项式
$f_1', \ldots, f_r', g_1, \ldots, g_m$ 的所有系数都属于 $A_0$。
置 $B_0 = A_0[x_1, \ldots, x_n]/(g_1, \ldots, g_m)$，并记
$f_{i, 0}$ 为 $f_i'$ 在 $B_0$ 中的像。则
$B = B_0 \otimes_{A_0} A$，并且
$$
B/(f_1, \ldots, f_r) =
B_0/(f_{0, 1}, \ldots, f_{0, r}) \otimes_{A_0} A.
$$
由《代数》引理
\ref{algebra-lemma-flat-finite-presentation-limit-flat}，扩大 $A_0$ 后，
可以假设 $B_0/(f_{0, 1}, \ldots, f_{0, r})$ 在 $A_0$ 上平坦。
此时 Koszul 同调群
$H_1(K_\bullet(B_0, f_{0, 1}, \ldots, f_{0, r}))$ 未必为零。
另一方面，由于 $B_0$ Noether，它是有限生成 $B_0$-模。令
$\xi_1, \ldots, \xi_n \in H_1(K_\bullet(B_0, f_{0, 1}, \ldots, f_{0, r}))$
为生成元。设 $A_0 \subset A_1 \subset A$ 为一个更大的有限型
$\mathbf{Z}$-子代数，它包含在 $A$ 中。记 $f_{1, i}$ 为 $f_{0, i}$ 在
$B_1 = B_0 \otimes_{A_0} A_1$ 中的像。由《代数进阶》引理
\ref{more-algebra-lemma-base-change-H1-regular}，映射
$$
H_1(K_\bullet(B_0, f_{0, 1}, \ldots, f_{0, r})) \otimes_{A_0} A_1
\longrightarrow
H_1(K_\bullet(B_1, f_{1, 1}, \ldots, f_{1, r}))
$$
是满射。此外，显然，对上述所有 $A_1$ 取法，复形
$K_\bullet(B_1, f_{1, 1}, \ldots, f_{1, r})$ 的余极限是复形
$K_\bullet(B, f_1, \ldots, f_r)$，后者在次数 $1$ 上无环。因此，由
《代数》引理 \ref{algebra-lemma-directed-colimit-exact}，
$$
\colim_{A_0 \subset A_1 \subset A}
H_1(K_\bullet(B_1, f_{1, 1}, \ldots, f_{1, r}))
= 0
$$
于是可选择 $A_1$，使得 $\xi_1, \ldots, \xi_n$ 全部映到
$H_1(K_\bullet(B_1, f_{1, 1}, \ldots, f_{1, r}))$ 中的零。
换言之，Koszul 同调群
$H_1(K_\bullet(B_1, f_{1, 1}, \ldots, f_{1, r}))$
为零。

\medskip\noindent
考察仿射概形的态射 $X_1 \to S_1$，它等于 $\Spec$，取自环映射
$A_1 \to B_1$，并令
$Z_1 = \Spec(B_1/(f_{1, 1}, \ldots, f_{1, r}))$。
由于 $B = B_1 \otimes_{A_1} A$，即 $X = X_1 \times_{S_1} S$，
并且类似地 $Z = Z_1 \times_S S_1$，现在只需对
$X_1 \to S_1$ 与相对 $H_1$-正则浸入 $Z_1 \to X_1$ 证明 (1)，
见《态射》引理
\ref{morphisms-lemma-base-change-module-flat}。
于是已约化到 $X \to S$ 为 Noether 概形之间的有限型态射的情形。
在此情形，由引理
\ref{divisors-lemma-quasi-regular-immersion-flat-at-x}，$X \to S$ 在
$Z$ 的每个点处平坦。结合此时平坦点集为开集这一事实，见《代数》定理
\ref{algebra-theorem-openness-flatness}，可知 (1) 成立。然后 (2) 由应用引理
\ref{divisors-lemma-relative-regular-immersion-flat-in-neighbourhood}
得出。
\end{proof}

\noindent
如果环境概形在基上平坦且局部有限表示，则可以用纤维刻画相对拟正则浸入。

\begin{lemma}
\label{divisors-lemma-fibre-quasi-regular}
设 $\varphi : X \to S$ 为局部有限表示的平坦态射。设
$T \subset X$ 为闭子概形。设 $x \in T$，其像为 $s \in S$。
\begin{enumerate}
\item 若 $T_s \subset X_s$ 在 $x$ 的一个邻域中是拟正则浸入，
则存在开集 $U \subset X$ 以及相对拟正则浸入 $Z \subset U$，
使得 $Z_s = T_s \cap U_s$ 且 $T \cap U \subset Z$。
\item 若 $T_s \subset X_s$ 在 $x$ 的一个邻域中是拟正则浸入，
态射 $T \to X$ 有限表示，并且 $T \to S$ 在 $x$ 处平坦，则可以选择
(1) 中的 $U$ 与 $Z$，使得 $T \cap U = Z$。
\item 若 $T_s \subset X_s$ 在 $x$ 的一个邻域中是拟正则浸入，
并且 $T$ 由 $c$ 个方程在 $x$ 的一个邻域中切出，其中
$c = \dim_x(X_s) - \dim_x(T_s)$，则可以选择 (1) 中的 $U$ 与 $Z$，
使得 $T \cap U = Z$。
\end{enumerate}
每种情形中，由引理
\ref{divisors-lemma-relative-regular-immersion-flat-in-neighbourhood}，
$Z \to U$ 都是正则浸入。特别地，若 $T \to S$ 局部有限表示且平坦，
并且所有纤维 $T_s \subset X_s$ 都是拟正则浸入，则 $T \to X$
是相对拟正则浸入。
\end{lemma}

\begin{proof}
选择仿射开邻域 $\Spec(A)$（它包含 $s$）与仿射开邻域
$\Spec(B)$（它包含 $x$），使得
$\varphi(\Spec(B)) \subset \Spec(A)$。令
$\mathfrak p \subset A$ 为对应于 $s$ 的素理想，令
$\mathfrak q \subset B$ 为对应于 $x$ 的素理想，并令
$I \subset B$ 为对应于 $T$ 的理想。由引理的初始假设可知
$A \to B$ 平坦且有限表示。(1) 中的假设意味着，在缩小
$\Spec(B)$ 后，可以假设 $I(B \otimes_A \kappa(\mathfrak p))$
由一个元素拟正则列生成。必要时把 $B$ 在某个
$g \in B$（且 $g \not \in \mathfrak q$）处局部化后，可以假设存在
$f_1, \ldots, f_r \in I$，它们在
$B \otimes_A \kappa(\mathfrak p)$ 中的像构成一个拟正则列，并生成
$I(B \otimes_A \kappa(\mathfrak p))$。由《代数》引理
\ref{algebra-lemma-quasi-regular-regular} 与
\ref{algebra-lemma-regular-sequence-in-neighbourhood}，再次局部化后，
可以假设 $f_1, \ldots, f_r \in I$ 在
$B \otimes_A \kappa(\mathfrak p)$ 中构成正则列。由引理
\ref{divisors-lemma-fibre-Cartier}，
$Z_1 = V(f_1) \subset \Spec(B)$ 是相对有效 Cartier 除子；这里必要时
再次局部化 $B$。
再次应用同一引理，这次应用于
$Z_2 = V(f_1, f_2) \subset Z_1$，可知 $Z_2 \subset Z_1$
是相对有效 Cartier 除子。依此类推，直到到达
$Z = Z_n = V(f_1, \ldots, f_n)$。于是
$Z \to \Spec(B)$ 是正则浸入，并且 $Z$ 在 $S$ 上平坦；特别地，
$Z \to \Spec(B)$ 在 $\Spec(A)$ 上是相对拟正则浸入。
这证明了 (1)。

\medskip\noindent
为证明 (2)，考察闭浸入 $Z \to D$。满环映射
$u : \mathcal{O}_{D, x} \to \mathcal{O}_{Z, x}$
是本质有限表示的平坦局部 $\mathcal{O}_{S, s}$-代数之间的映射，
并且模掉 $\mathfrak m_s$ 后成为同构。因此，由《代数》引理
\ref{algebra-lemma-mod-injective-general}，它是同构。
由此，$Z \to D$ 在 $x$ 的一个邻域中是同构，见《代数》引理
\ref{algebra-lemma-local-isomorphism}。

\medskip\noindent
为证明 (3)，必要时缩小 $U$ 后，可以假设 $Z$ 的理想由正则列
$f_1, \ldots, f_r$ 生成（见上面对 $Z$ 的构造），并且 $T$ 的理想由
$g_1, \ldots, g_c$ 生成。断言 $c = r$。事实上，
\begin{align*}
\dim_x(X_s) & = \dim(\mathcal{O}_{X_s, x}) +
\text{trdeg}_{\kappa(s)}(\kappa(x)), \\
\dim_x(T_s) & = \dim(\mathcal{O}_{T_s, x}) +
\text{trdeg}_{\kappa(s)}(\kappa(x)), \\
\dim(\mathcal{O}_{X_s, x}) & = \dim(\mathcal{O}_{T_s, x}) + r
\end{align*}
前两个等式来自《代数》引理
\ref{algebra-lemma-dimension-at-a-point-finite-type-field}，第二个等式
来自将《代数》引理
\ref{algebra-lemma-one-equation}
应用 $r$ 次。由于 $T \subset Z$，可知
$f_i = \sum b_{ij} g_j$。但是，$Z$ 与 $T$ 的理想切出
$X_s$ 的同一个拟正则闭子概形，而这一点在 $x$ 的一个邻域中成立。因此，矩阵
$(b_{ij}) \bmod \mathfrak m_x$ 可逆（略去一些细节）。于是
$(b_{ij})$ 在 $x$ 的一个开邻域中可逆。换言之，
$T \cap U = Z$，这是在缩小 $U$ 后成立。

\medskip\noindent
引理的最后几个陈述立即由 (2) 得出，并结合如下事实：
$Z \to S$ 局部有限表示，当且仅当 $Z \to X$ 有限表示；见《态射》引理
\ref{morphisms-lemma-composition-finite-presentation} 与
\ref{morphisms-lemma-finite-presentation-permanence}。
\end{proof}

\noindent
下述引理加强了《态射》引理
\ref{morphisms-lemma-section-smooth-morphism}。

\begin{lemma}
\label{divisors-lemma-section-smooth-regular-immersion}
设 $f : X \to S$ 为概形的平滑态射。设
$\sigma : S \to X$ 为 $f$ 的截面。则 $\sigma$ 是正则浸入。
\end{lemma}

\begin{proof}
由《概形》引理
\ref{schemes-lemma-semi-diagonal}，态射 $\sigma$ 是浸入。
用 $X$ 的一个开子概形替换它，该开子概形是 $\sigma(S)$ 的邻域；
于是可以假设 $\sigma$ 是闭浸入。
令 $T = \sigma(S)$ 为 $X$ 的相应闭子概形。由于 $T \to S$ 是同构，
它平坦且有限表示。平滑态射也平坦且局部有限表示，见《态射》引理
\ref{morphisms-lemma-smooth-flat} 与
\ref{morphisms-lemma-smooth-locally-finite-presentation}。
因此，根据引理
\ref{divisors-lemma-fibre-quasi-regular}，只需证明
$T_s \subset X_s$ 是拟正则闭子概形。这立即由《态射》引理
\ref{morphisms-lemma-section-smooth-morphism} 得出，但也可以如下直接说明。
设 $k$ 为域，并设 $A$ 为平滑 $k$-代数。设
$\mathfrak m \subset A$ 为剩余域等于 $k$ 的极大理想。则
$\mathfrak m$ 由拟正则列生成；这里必要时先将 $A$ 替换为某个
$A_g$，其中 $g \in A$、$g \not \in \mathfrak m$。在《代数》引理
\ref{algebra-lemma-characterize-smooth-over-field} 中已证明
$A_{\mathfrak m}$ 是正则局部环，故
$\mathfrak mA_{\mathfrak m}$ 由正则列生成。这的确蕴含
$\mathfrak m$ 由正则列生成（在必要时将 $A$ 替换为某个
$A_g$，其中 $g \in A$、$g \not \in \mathfrak m$），见《代数》引理
\ref{algebra-lemma-regular-sequence-in-neighbourhood}。
\end{proof}

\noindent
下述引理有某种逆，见引理
\ref{divisors-lemma-push-regular-immersion-thru-smooth}。

\begin{lemma}
\label{divisors-lemma-lift-regular-immersion-to-smooth}
设
$$
\xymatrix{
Y \ar[rd]_j \ar[rr]_i & & X \ar[ld] \\
& S
}
$$
为概形态射的交换图。假设 $X \to S$ 平滑，并且 $i$、$j$ 是浸入。
若 $j$ 是正则浸入（相应地，Koszul-正则、$H_1$-正则、拟正则浸入），
则 $i$ 也是。
\end{lemma}

\begin{proof}
可以把 $i$ 写成复合
$$
Y \to Y \times_S X \to X
$$
由引理
\ref{divisors-lemma-section-smooth-regular-immersion}，第一个箭头是正则浸入。
第二个箭头是 $Y \to S$ 的平坦基变换，因而是正则浸入（相应地，
Koszul-正则、$H_1$-正则、拟正则浸入），见引理
\ref{divisors-lemma-flat-base-change-regular-immersion}。
应用引理
\ref{divisors-lemma-composition-regular-immersion}
即得结论。
\end{proof}

\begin{lemma}
\label{divisors-lemma-immersion-lci-into-smooth-regular-immersion}
设
$$
\xymatrix{
Y \ar[rd] \ar[rr]_i & & X \ar[ld] \\
& S
}
$$
为概形态射的交换图。假设 $Y \to S$ 是 syntomic 态射，
$X \to S$ 平滑，并且 $i$ 是浸入。则 $i$ 是正则浸入。
\end{lemma}

\begin{proof}
用 $X$ 的一个开子概形替换它，该开子概形是 $i(Y)$ 的邻域；
于是可以假设 $i$ 是闭浸入。
令 $T = i(Y)$ 为 $X$ 的相应闭子概形。由于 $T \cong Y$，态射
$T \to S$ 平坦且有限表示（《态射》引理
\ref{morphisms-lemma-syntomic-locally-finite-presentation} 与
\ref{morphisms-lemma-syntomic-flat}）。
平滑态射也平坦且局部有限表示（《态射》引理
\ref{morphisms-lemma-smooth-flat} 与
\ref{morphisms-lemma-smooth-locally-finite-presentation}）。
因此，根据引理
\ref{divisors-lemma-fibre-quasi-regular}，只需证明
$T_s \subset X_s$ 是拟正则闭子概形。由于 $X_s$ 在域上局部有限型，
它是 Noether 概形（《态射》引理
\ref{morphisms-lemma-finite-type-noetherian}）。因此，可以逐点检验
$T_s \subset X_s$ 是拟正则浸入，见引理
\ref{divisors-lemma-Noetherian-scheme-regular-ideal}。取 $t \in T_s$。
由《态射》引理
\ref{morphisms-lemma-local-complete-intersection}，局部环
$\mathcal{O}_{T_s, t}$ 是 $\kappa(s)$ 上的局部完全交。局部环
$\mathcal{O}_{X_s, t}$ 正则，见《代数》引理
\ref{algebra-lemma-characterize-smooth-over-field}。由《代数》引理
\ref{algebra-lemma-lci-local} 可知，满射
$\mathcal{O}_{X_s, t} \to \mathcal{O}_{T_s, t}$
的核由正则列生成，这正是所需结论。
\end{proof}

\begin{lemma}
\label{divisors-lemma-immersion-smooth-into-smooth-regular-immersion}
设
$$
\xymatrix{
Y \ar[rd] \ar[rr]_i & & X \ar[ld] \\
& S
}
$$
为概形态射的交换图。假设 $Y \to S$ 平滑、$X \to S$ 平滑，并且
$i$ 是浸入。则 $i$ 是正则浸入。
\end{lemma}

\begin{proof}
这是引理
\ref{divisors-lemma-immersion-lci-into-smooth-regular-immersion}
的特例，因为平滑态射是 syntomic 态射，见《态射》引理
\ref{morphisms-lemma-smooth-syntomic}。
\end{proof}

\begin{lemma}
\label{divisors-lemma-push-regular-immersion-thru-smooth}
设
$$
\xymatrix{
Y \ar[rd]_j \ar[rr]_i & & X \ar[ld] \\
& S
}
$$
为概形态射的交换图。假设 $X \to S$ 平滑，并且 $i$ 与 $j$ 是浸入。
若 $i$ 是 Koszul-正则浸入（相应地，$H_1$-正则、拟正则浸入），
则 $j$ 也是。
\end{lemma}

\begin{proof}
以下不再另行说明地使用引理
\ref{divisors-lemma-regular-quasi-regular-immersion}。
任取 $y \in Y$。置 $x = i(y)$，并置 $s = j(y)$。
只需在用 $X$ 与 $S$ 的开邻域 $U$ 与 $V$（分别包含 $x$ 与 $s$）
替换它们，并用 $Y$ 中 $y$ 的一个开邻域（包含于
$i^{-1}(U) \cap j^{-1}(V)$）替换它后证明结论。

\medskip\noindent
先对 $X = \mathbf{A}^n_S$ 证明结论。用 $S$ 的仿射开子集 $V$
替换它，并用 $Y$ 的开子集 $j^{-1}(V)$ 替换它后，可以假设
$j$ 是闭浸入且 $S$ 仿射。
写成 $S = \Spec(A)$。于是 $j : Y \to S$ 把 $Y$ 同构到闭子概形
$\Spec(A/I)$，其中 $I \subset A$ 为某个理想。映射
$i : Y = \Spec(A/I) \to
\mathbf{A}^n_S = \Spec(A[x_1, \ldots, x_n])$
对应于 $A$-代数同态
$i^\sharp : A[x_1, \ldots, x_n] \to A/I$。
选择 $a_i \in A$，使其像为 $i^\sharp(x_i)$，此像位于 $A/I$ 中。
注意，闭浸入 $i$ 的理想是
$$
J = (x_1 - a_1, \ldots, x_n - a_n) + IA[x_1, \ldots, x_n].
$$
置 $K = (x_1 - a_1, \ldots, x_n - a_n)$。断言序列
$$
0 \to K/KJ \to J/J^2 \to J/(K + J^2) \to 0
$$
分裂正合。为说明这一点，注意 $K/K^2$ 以 $x_i - a_i$ 为基，
并且在环 $A[x_1, \ldots, x_n]/K \cong A$ 上自由。
因此，$K/KJ$ 在环 $A[x_1, \ldots, x_n]/J \cong A/I$ 上自由，
基仍相同。另一方面，取微分给出映射
$$
\text{d}_{A[x_1, \ldots, x_n]/A} :
J/J^2
\longrightarrow
\Omega_{A[x_1, \ldots, x_n]/A} \otimes_{A[x_1, \ldots, x_n]}
A[x_1, \ldots, x_n]/J
$$
它把生成元 $x_i - a_i$ 映到右边自由模的基元素 $\text{d}x_i$。
断言随之成立。此外，注意
$x_1 - a_1, \ldots, x_n - a_n$ 是
$A[x_1, \ldots, x_n]$ 中的正则列，其商环为
$A[x_1, \ldots, x_n]/(x_1 - a_1, \ldots, x_n - a_n) \cong A$。
于是，闭浸入 $i$ 有分解
$$
Y \to V(x_1 - a_1, \ldots, x_n - a_n) \to \mathbf{A}^n_S
$$
其中复合是 Koszul-正则浸入（相应地，$H_1$-正则、拟正则浸入），
第二个箭头是正则浸入，并且相应的余法序列分裂。现在由引理
\ref{divisors-lemma-permanence-regular-immersion}
得出结论。

\medskip\noindent
接着，在 $i$ 是 $H_1$-正则的情形证明结论。也就是说，按证明第一段
那样缩小后，可以假设 $Y$、$X$ 与 $S$ 都仿射。此时，可以选择闭浸入
$h : X \to \mathbf{A}^n_S$，它在 $S$ 上，其中 $n$ 为某个整数。由引理
\ref{divisors-lemma-immersion-smooth-into-smooth-regular-immersion}，
$h$ 是正则浸入。因此，$h \circ i$ 是 $H_1$-正则浸入，见引理
\ref{divisors-lemma-composition-regular-immersion}（注意，这一步在
“拟正则情形”中不成立）。于是约化到上面已经证明的
$X = \mathbf{A}^n_S$ 且 $S$ 仿射的情形。

\medskip\noindent
最后，假设 $i$ 拟正则。按证明第一段那样缩小后，可以使用《态射》引理
\ref{morphisms-lemma-smooth-etale-over-affine-space}，把 $f$ 分解为
$X \to \mathbf{A}^n_S \to S$，其中第一个态射
$X \to \mathbf{A}^n_S$ 是 \'etale 的。这把问题约化到两个情形：
(a) $X = \mathbf{A}^n_S$，以及 (b) $f$ 是 \'etale 的。
情形 (a) 已在证明第二段处理。情形 (b) 由下一段处理。

\medskip\noindent
假设 $f$ 是 \'etale 的。缩小后，可以假设 $X$、$Y$ 与 $S$ 仿射，
$i$ 与 $j$ 是闭浸入（略去一个小细节）。记
$S = \Spec(A)$、$X = \Spec(B)$ 且
$Y = \Spec(B/J) = \Spec(A/I)$。进一步缩小后，可以假设 $J$
由拟正则列生成。环映射 $A \to B$ 是 \'etale 的，因而形式 \'etale
（《代数》引理 \ref{algebra-lemma-formally-etale-etale}）。
因此，由《代数》引理
\ref{algebra-lemma-formally-etale-lift-infinitesimal}，
$\bigoplus I^n/I^{n + 1} \cong \bigoplus J^n/J^{n + 1}$。
由于 $J$ 由拟正则列生成，$I$ 也如此。这就完成了证明。
\end{proof}









\section{亚纯函数与亚纯截面}
\label{divisors-section-meromorphic-functions}

\noindent
本节只包含一般定义和一些初等结果。若要了解可能遇到的陷阱，见
\cite{misconceptions}\footnote{危险，威尔·罗宾逊！}。

\medskip\noindent
设 $(X, \mathcal{O}_X)$ 为局部带环空间。对任意开集
$U \subset X$，已经定义了集合
$\mathcal{S}(U) \subset \mathcal{O}_X(U)$，即 $\mathcal{O}_X$ 在 $U$ 上的
正则截面集合，见定义
\ref{divisors-definition-regular-section}。正则截面限制到更小开集后仍然正则。
因此，$\mathcal{S} : U \mapsto \mathcal{S}(U)$ 是
$\mathcal{O}_X$ 的子层（集合层）。有时记
$\mathcal{S} = \mathcal{S}_X$，以表明其对 $X$ 的依赖。此外，
$\mathcal{S}(U)$ 是环 $\mathcal{O}_X(U)$ 的乘法子集；这对每个 $U$
都成立。因此，可以考察环的预层
$$
U \longmapsto \mathcal{S}(U)^{-1} \mathcal{O}_X(U),
$$
见《模》引理
\ref{modules-lemma-simple-invert}。

\begin{definition}
\label{divisors-definition-sheaf-meromorphic-functions}
设 $(X, \mathcal{O}_X)$ 为局部带环空间。$X$ 上的{\it 亚纯函数层}
是与上面所示预层相伴的层{\it $\mathcal{K}_X$}。$X$ 上的
{\it 亚纯函数}是 $\mathcal{K}_X$ 的整体截面。
\end{definition}

\noindent
由于每个 $\mathcal{S}(U)$ 的每个元素都是
$\mathcal{O}_X(U)$ 上的非零因子，可知环层的自然映射
$\mathcal{O}_X \to \mathcal{K}_X$ 是单射。

\begin{example}
\label{divisors-example-no-change}
设
$A = \mathbf{C}[x, \{y_\alpha\}_{\alpha \in \mathbf{C}}]/
((x - \alpha)y_\alpha, y_\alpha y_\beta)$。$A$ 的每个元素都可唯一写成
$f(x) + \sum \lambda_\alpha y_\alpha$，其中
$f(x) \in \mathbf{C}[x]$ 且 $\lambda_\alpha \in \mathbf{C}$。
令 $X = \Spec(A)$。在此情形，$\mathcal{O}_X = \mathcal{K}_X$，
因为在任意仿射开集 $D(f)$ 上，环 $A_f$ 中的每个非零因子都是单位
（略去证明）。
\end{example}

\noindent
设 $(X, \mathcal{O}_X)$ 为局部带环空间。设 $\mathcal{F}$ 为
$\mathcal{O}_X$-模层。考察预层
$U \mapsto \mathcal{S}(U)^{-1}\mathcal{F}(U)$。它的层化是层
$\mathcal{F} \otimes_{\mathcal{O}_X} \mathcal{K}_X$，见《模》引理
\ref{modules-lemma-simple-invert-module}。

\begin{definition}
\label{divisors-definition-meromorphic-section}
设 $X$ 为局部带环空间。设 $\mathcal{F}$ 为 $\mathcal{O}_X$-模层。
\begin{enumerate}
\item 记 $\mathcal{K}_X(\mathcal{F})$ 为如下
$\mathcal{K}_X$-模层：它是预层
$U \mapsto \mathcal{S}(U)^{-1}\mathcal{F}(U)$ 的层化。等价地，
$\mathcal{K}_X(\mathcal{F}) =
\mathcal{F} \otimes_{\mathcal{O}_X} \mathcal{K}_X$（见上文）。
\item $\mathcal{F}$ 的{\it 亚纯截面}是
$\mathcal{K}_X(\mathcal{F})$ 的整体截面。
\end{enumerate}
\end{definition}

\noindent
特别地，对任意点 $x \in X$，有
$$
\mathcal{K}_X(\mathcal{F})_x
=
\mathcal{F}_x \otimes_{\mathcal{O}_{X, x}} \mathcal{K}_{X, x}
=
\mathcal{S}_x^{-1}\mathcal{F}_x
$$
不过必须谨慎，因为 $\mathcal{S}_x$ 未必等于局部环
$\mathcal{O}_{X, x}$ 中的非零因子集合。具体地，总有到全商环的单射
$$
\mathcal{K}_{X, x} \longrightarrow Q(\mathcal{O}_{X, x})
$$
它为满射，当且仅当 $\mathcal{S}_x$ 是
$\mathcal{O}_{X, x}$ 中的非零因子集合。亚纯截面层一般不是拟凝聚模，
但它们与拟凝聚模有一些共同性质。

\begin{definition}
\label{divisors-definition-pullback-meromorphic-sections}
设 $f : (X, \mathcal{O}_X) \to (Y, \mathcal{O}_Y)$ 为局部带环空间的态射。
称{\it 亚纯函数对 $f$ 的拉回有定义}，如果对每一对开集
$U \subset X$、$V \subset Y$，只要
$f(U) \subset V$，并对任意截面 $s \in \Gamma(V, \mathcal{S}_Y)$，
拉回 $f^\sharp(s) \in \Gamma(U, \mathcal{O}_X)$ 都属于
$\Gamma(U, \mathcal{S}_X)$。
\end{definition}

\noindent
在这种情形，存在诱导映射
$f^\sharp : f^{-1}\mathcal{K}_Y \to \mathcal{K}_X$；换言之，得到
带环空间态射的交换图
$$
\xymatrix{
(X, \mathcal{K}_X) \ar[r] \ar[d]^f &
(X, \mathcal{O}_X) \ar[d]^f \\
(Y, \mathcal{K}_Y) \ar[r] &
(Y, \mathcal{O}_Y)
}
$$
有时记 $f^*(s) = f^\sharp(s)$，其中
$s \in \Gamma(Y, \mathcal{K}_Y)$。

\begin{lemma}
\label{divisors-lemma-pullback-meromorphic-sections-defined}
设 $f : X \to Y$ 为概形的态射。在下列每种情形中，亚纯函数拉回有定义：
\begin{enumerate}
\item $X$ 的每个弱伴随点都映到 $Y$ 的某个不可约分支的泛点；
\item $X$、$Y$ 整，且 $f$ 支配；
\item $X$ 整，且 $X$ 的泛点映到 $Y$ 的某个不可约分支的泛点；
\item $X$ 约化，且 $X$ 的每个不可约分支的每个泛点都映到
$Y$ 的某个不可约分支的泛点；
\item $X$ 局部 Noether，且 $X$ 的任意伴随点都映到
$Y$ 的某个不可约分支的泛点；
\item $X$ 局部 Noether、没有嵌入点，并且 $X$ 的每个不可约分支的
泛点都映到 $Y$ 的某个不可约分支的泛点；以及
\item $f$ 平坦。
\end{enumerate}
\end{lemma}

\begin{proof}
问题在 $X$ 与 $Y$ 上都是局部的。因此，约化到
$X = \Spec(A)$、$Y = \Spec(R)$，并且 $f$ 由环映射
$\varphi : R \to A$ 给出的情形。由引理
\ref{divisors-lemma-regular-section-structure-sheaf}
中对结构层正则截面的刻画，只需证明 $R \to A$ 把非零因子映到
非零因子。设 $t \in R$ 为非零因子。

\medskip\noindent
若 $R \to A$ 平坦，则 $t : R \to R$ 为单射这一事实表明
$t : A \to A$ 为单射。这证明了 (7)。

\medskip\noindent
在其他情形中，注意 $t$ 不属于 $R$ 的任何极小素理想
（因为环的极小素理想中的每个元素都是零因子）。
因此，在情形 (1)，$\varphi(t)$ 不属于 $A$ 的任何弱伴随素理想。
此情形遂由《代数》引理
\ref{algebra-lemma-weakly-ass-zero-divisors}
得出。由引理 \ref{divisors-lemma-ass-weakly-ass}，情形 (5) 是 (1) 的特例。
情形 (6) 由 (5) 和定义得出。由引理
\ref{divisors-lemma-weakass-reduced}，情形 (4) 是 (1) 的特例。
情形 (2)、(3) 是 (4) 的特例。
\end{proof}


\begin{lemma}
\label{divisors-lemma-meromorphic-weakass-finite}
设 $X$ 为满足下列条件的概形：
\begin{enumerate}
\item[(a)] $X$ 的每个弱伴随点都是 $X$ 的某个不可约分支的泛点；并且
\item[(b)] 每个拟紧开集只有有限多个不可约分支。
\end{enumerate}
令 $X^0$ 为 $X$ 的不可约分支的泛点集合。则有
$$
\mathcal{K}_X =
\bigoplus\nolimits_{\eta \in X^0} j_{\eta, *}\mathcal{O}_{X, \eta} =
\prod\nolimits_{\eta \in X^0} j_{\eta, *}\mathcal{O}_{X, \eta}
$$
其中 $j_\eta : \Spec(\mathcal{O}_{X, \eta}) \to X$ 是《概形》第
\ref{schemes-section-points} 节的典范映射。此外：
\begin{enumerate}
\item $\mathcal{K}_X$ 是拟凝聚 $\mathcal{O}_X$-代数层；
\item 对每个拟凝聚 $\mathcal{O}_X$-模 $\mathcal{F}$，$\mathcal{F}$ 的亚纯截面层
$$
\mathcal{K}_X(\mathcal{F}) =
\bigoplus\nolimits_{\eta \in X^0} j_{\eta, *}\mathcal{F}_\eta =
\prod\nolimits_{\eta \in X^0} j_{\eta, *}\mathcal{F}_\eta
$$
是拟凝聚的；
\item $\mathcal{S}_x \subset \mathcal{O}_{X, x}$ 是非零因子集合，
对任意 $x \in X$；
\item $\mathcal{K}_{X, x}$ 是 $\mathcal{O}_{X, x}$ 的全商环，
对任意 $x \in X$；
\item $\mathcal{K}_X(U)$ 等于 $\mathcal{O}_X(U)$ 的全商环，
对任意仿射开集 $U \subset X$；
\item $X$ 的有理函数环（《态射》定义
\ref{morphisms-definition-rational-function}）就是 $X$ 的亚纯函数环；
写成公式即 $R(X) = \Gamma(X, \mathcal{K}_X)$。
\end{enumerate}
\end{lemma}

\begin{proof}
注意，局部有限的模层直和等于其直积；例如，这可以在茎上检验。又因
$\mathcal{K}_X(\mathcal{F}) =
\mathcal{F} \otimes_{\mathcal{O}_X} \mathcal{K}_X$，
可知 (2) 由其他陈述推出。还要注意，当 $X^0$ 有限时，(6) 由引理的
初始陈述与《态射》引理
\ref{morphisms-lemma-integral-scheme-rational-functions}
得出；一般情形的 (6) 再由此粘合而得（略去论证）。

\medskip\noindent
令
$j : Y = \coprod\nolimits_{\eta \in X^0} \Spec(\mathcal{O}_{X, \eta}) \to X$
为态射 $j_\eta$ 的余积。需要证明
$\mathcal{K}_X = j_*\mathcal{O}_Y$。首先注意，
$\mathcal{K}_Y = \mathcal{O}_Y$，因为 $Y$ 是 $0$ 维局部环之谱的不交并：
在零维局部环中，每个非零因子都是单位。其次注意，由引理
\ref{divisors-lemma-pullback-meromorphic-sections-defined}，$j$ 的亚纯函数拉回
有定义。这给出映射
$$
\mathcal{K}_X \longrightarrow j_*\mathcal{O}_Y.
$$
设 $\Spec(A) = U \subset X$ 为仿射开集。则 $A$ 是只有有限多个
极小素理想 $\mathfrak q_1, \ldots, \mathfrak q_t$ 的环，并且
$A$ 的每个弱伴随素理想都是某个 $\mathfrak q_i$。由《代数》引理
\ref{algebra-lemma-total-ring-fractions-no-embedded-points}
与 \ref{algebra-lemma-weakly-ass-zero-divisors}，得到
$Q(A) = \prod A_{\mathfrak q_i}$。换言之，预层的取值
$U \mapsto \mathcal{S}(U)^{-1}\mathcal{O}_X(U)$ 就与
$j_*\mathcal{O}_Y(U)$ 在仿射开集 $U$ 上一致。因此，上示映射是同构，
从而证明引理陈述中
第一个上示等式。

\medskip\noindent
最后，证明 (1)、(3)、(4)、(5)。在证明
$\mathcal{K}_X = j_*\mathcal{O}_Y$ 的过程中，已经看到了 (5)。
态射 $j$ 因我们关于 $X$ 的不可约分支集合局部有限这一假设而拟紧。
因此，$j$ 拟紧且拟分离（因为 $Y$ 分离）。由《概形》引理
\ref{schemes-lemma-push-forward-quasi-coherent}，
$j_*\mathcal{O}_Y$ 拟凝聚。这证明了 (1)。设 $x \in X$。
可以选择仿射开邻域 $U = \Spec(A)$ 作为 $x$ 的邻域，使其所有
不可约分支都经过 $x$。于是 $A \subset A_\mathfrak p$，因为 $A$ 的每个弱伴随素理想
都包含于 $\mathfrak p$，故由《代数》引理
\ref{algebra-lemma-weakly-ass-zero-divisors}，
$A \setminus \mathfrak p$ 中的元素都是非零因子。
容易推出，$A_\mathfrak p$ 的任意非零因子都是 $x$ 的某个
（可能更小的）仿射开邻域上的非零因子的像。这证明了 (3)。
(4) 由 (3) 通过计算茎得出。
\end{proof}

\begin{definition}
\label{divisors-definition-regular-meromorphic-section}
设 $X$ 为局部带环空间。设 $\mathcal{L}$ 为可逆
$\mathcal{O}_X$-模。设 $s$ 是 $\mathcal{L}$ 的一个亚纯截面。若其所诱导的映射
$\mathcal{K}_X \to \mathcal{K}_X(\mathcal{L})$ 为单射，则称其
{\it 正则}。换言之，$s$ 是可逆
$\mathcal{K}_X$-模 $\mathcal{K}_X(\mathcal{L})$ 的正则截面，见定义
\ref{divisors-definition-regular-section}。
\end{definition}

\noindent
下面明确说明何时可以拉回（正则）亚纯截面。

\begin{lemma}
\label{divisors-lemma-meromorphic-sections-pullback}
设 $f : X \to Y$ 为局部带环空间的态射。假设 $f$ 的亚纯函数拉回有定义
（见定义
\ref{divisors-definition-pullback-meromorphic-sections}）。
\begin{enumerate}
\item 设 $\mathcal{F}$ 为 $\mathcal{O}_Y$-模层。存在典范拉回映射
$f^* : \Gamma(Y, \mathcal{K}_Y(\mathcal{F})) \to
\Gamma(X, \mathcal{K}_X(f^*\mathcal{F}))$，用于 $\mathcal{F}$ 的亚纯截面。
\item 设 $\mathcal{L}$ 为可逆 $\mathcal{O}_Y$-模。
设 $s$ 是 $\mathcal{L}$ 的正则亚纯截面。则 $f^*s$ 是
$f^*\mathcal{L}$ 的正则亚纯截面。
\end{enumerate}
\end{lemma}

\begin{proof}
略。
\end{proof}

\begin{lemma}
\label{divisors-lemma-regular-meromorphic-ideal-denominators}
设 $X$ 为概形。设 $\mathcal{L}$ 为可逆 $\mathcal{O}_X$-模。
设 $s$ 为 $\mathcal{L}$ 的正则亚纯截面。记
$\mathcal{I} \subset \mathcal{O}_X$ 为由下述规则定义的理想层：
$$
\mathcal{I}(V)
=
\{f \in \mathcal{O}_X(V) \mid fs \in \mathcal{L}(V)\}.
$$
由于 $\mathcal{L}(V) \subset \mathcal{K}_X(\mathcal{L})(V)$，
该公式有意义。则 $\mathcal{I}$ 是拟凝聚理想层，并且有单射
$$
1 : \mathcal{I} \longrightarrow \mathcal{O}_X,
\quad
s : \mathcal{I} \longrightarrow \mathcal{L}
$$
二者的余核都支撑在 $X$ 的闭无处稠密子集上。
\end{lemma}

\begin{proof}
问题在 $X$ 上是局部的。因此，可以假设 $X = \Spec(A)$ 且
$\mathcal{L} = \mathcal{O}_X$。进一步缩小后，可以假设
$s = a/b$，其中 $a, b \in A$ 在 $A$ 中{\it 都}是非零因子。
置 $I = \{x \in A \mid x(a/b) \in A\}$。

\medskip\noindent
为证明 $\mathcal{I}$ 拟凝聚，需要证明
$I_f = \{x \in A_f \mid x(a/b) \in A_f\}$，对每个 $f \in A$。若
$c/f^n \in A_f$ 且 $(c/f^n)(a/b) \in A_f$，则
$f^mc(a/b) \in A$ 对某个 $m$ 成立，故 $c/f^n \in I_f$。反过来，容易看出
$I_f$ 包含于 $\{x \in A_f \mid x(a/b) \in A_f\}$。
这证明了拟凝聚性。

\medskip\noindent
下面证明最后一个陈述。显然 $(b) \subset I$。因此
$V(I) \subset V(b)$ 是无处稠密子集，因为 $b$ 是非零因子。
所以 $1$ 的余核支撑在某个闭无处稠密集合中。对 $s$ 的余核，
同样的论证成立，因为 $s(b) = (a) \subset sI \subset A$。
\end{proof}

\begin{definition}
\label{divisors-definition-regular-meromorphic-ideal-denominators}
设 $X$ 为概形。设 $\mathcal{L}$ 为可逆 $\mathcal{O}_X$-模。
设 $s$ 为 $\mathcal{L}$ 的正则亚纯截面。引理
\ref{divisors-lemma-regular-meromorphic-ideal-denominators}
中构造的理想层 $\mathcal{I}$ 称为 $s$ 的{\it 分母理想层}。
\end{definition}




\section{亚纯函数与亚纯截面：Noether 情形}
\label{divisors-section-meromorphic-noetherian}

\noindent
对于局部 Noether 概形，可以证明亚纯函数层的一些结果。然而，
\cite{misconceptions} 中有一个例子表明，$\mathcal{K}_X$ 对 Noether 概形
$X$ 也未必拟凝聚。

\begin{lemma}
\label{divisors-lemma-meromorphic-section-restricts-to-zero}
设 $X$ 为拟紧概形。设 $h \in \Gamma(X, \mathcal{O}_X)$ 且
$f \in \Gamma(X, \mathcal{K}_X)$，并假设 $f$ 在 $X_h$ 上的限制为零。
则 $h^n f = 0$ 对某个 $n \gg 0$ 成立。
\end{lemma}

\begin{proof}
可找到 $X$ 的一个由仿射开集 $U$ 构成的覆盖，使得
$f|_U = s^{-1}a$，其中 $a \in \mathcal{O}_X(U)$ 且
$s \in \mathcal{S}(U)$。由于 $X$ 拟紧，可以取这种形式的有限仿射开覆盖。
因此，只需在 $X = \Spec(A)$ 且 $f = s^{-1}a$ 时证明引理。
注意，$s \in A$ 是非零因子，故只需证明 $f = a$ 的情形。条件
$f|_{X_h} = 0$ 蕴含 $a$ 在
$A_h = \mathcal{O}_X(X_h)$ 中的像为零，因为
$\mathcal{O}_X \subset \mathcal{K}_X$。因此，$h^na = 0$ 对某个
$n > 0$ 成立，这正是所需结论。
\end{proof}

\begin{lemma}
\label{divisors-lemma-locally-Noetherian-K}
设 $X$ 为局部 Noether 概形。
\begin{enumerate}
\item 对任意 $x \in X$，$\mathcal{S}_x \subset \mathcal{O}_{X, x}$
是非零因子集合，因而 $\mathcal{K}_{X, x}$ 是
$\mathcal{O}_{X, x}$ 的全商环。
\item 对任意仿射开集 $U \subset X$，环
$\mathcal{K}_X(U)$ 等于 $\mathcal{O}_X(U)$ 的全商环。
\end{enumerate}
\end{lemma}

\begin{proof}
为证明本引理，可以假设 $X$ 是某个 Noether 环 $A$ 的谱。设
$x \in X$ 对应于 $\mathfrak p \subset A$。

\medskip\noindent
证明 (1)。显然，$\mathcal{S}_x$ 包含于
$\mathcal{O}_{X, x} = A_\mathfrak p$ 的非零因子集合。反过来，设
$f, g \in A$、$g \not \in \mathfrak p$，并假设
$f/g$ 是 $A_{\mathfrak p}$ 中的非零因子。令
$I = \{a \in A \mid af = 0\}$。由局部化的正合性可知
$I_{\mathfrak p} = 0$。由于 $A$ Noether，$I$ 有限生成，故
$g'I = 0$，其中存在元素 $g' \in A$，且 $g' \not \in \mathfrak p$。
因此，$f$ 在 $A_{g'}$ 中是非零因子，也就是在
$\mathfrak p$ 的某个 Zariski 开邻域中是非零因子。于是
$f/g$ 是 $\mathcal{S}_x$ 的元素。

\medskip\noindent
证明 (2)。设 $f \in \Gamma(X, \mathcal{K}_X)$ 为亚纯函数。令
$I = \{a \in A \mid af \in A\}$。固定素理想
$\mathfrak p \subset A$，它对应于点 $x \in X$。由 (1)，可以把 $f$
在 $\mathfrak p$ 处茎中的像写为 $a/b$，其中
$a, b \in A_{\mathfrak p}$，且 $b \in A_{\mathfrak p}$ 不是零因子。
写 $b = c/d$，其中 $c, d \in A$、$d \not \in \mathfrak p$。
则 $ad - cf$ 是 $\mathcal{K}_X$ 的一个截面，并在 $x$ 的某个
开邻域上消失。设它在 $D(e)$ 上消失，其中
$e \in A$、$e \not \in \mathfrak p$。由引理
\ref{divisors-lemma-meromorphic-section-restricts-to-zero}，有
$e^n(ad - cf) = 0$，其中 $n \gg 0$ 是某个整数。因此
$e^nc \in I$，且 $e^nc$ 在 $A_{\mathfrak p}$ 中的像是非零因子。令
$\text{Ass}(A) = \{\mathfrak q_1, \ldots, \mathfrak q_t\}$ 为
$A$ 的伴随素理想。考察 $IA_{\mathfrak q_i}$ 并使用《代数》引理
\ref{algebra-lemma-associated-primes-localize}，上述论证表明，
$I \not \subset \mathfrak q_i$ 对每个 $i$ 成立。由《代数》引理
\ref{algebra-lemma-silly}，存在元素 $z \in I$，使得
$z \not \in \bigcup \mathfrak q_i$。由《代数》引理
\ref{algebra-lemma-ass-zero-divisors}，可知 $z$ 不是 $A$ 上的零因子。
因此，$f = (zf)/z$ 是 $A$ 的全商环的元素。这证明了 (2)。
\end{proof}

\begin{lemma}
\label{divisors-lemma-quasi-coherent-K}
设 $X$ 为没有嵌入点的局部 Noether 概形。令 $X^0$ 为 $X$ 的
不可约分支的泛点集合。则有
$$
\mathcal{K}_X =
\bigoplus\nolimits_{\eta \in X^0} j_{\eta, *}\mathcal{O}_{X, \eta} =
\prod\nolimits_{\eta \in X^0} j_{\eta, *}\mathcal{O}_{X, \eta}
$$
其中 $j_\eta : \Spec(\mathcal{O}_{X, \eta}) \to X$ 是《概形》第
\ref{schemes-section-points} 节的典范映射。此外：
\begin{enumerate}
\item $\mathcal{K}_X$ 是拟凝聚 $\mathcal{O}_X$-代数层；
\item 对每个拟凝聚 $\mathcal{O}_X$-模 $\mathcal{F}$，层
$$
\mathcal{K}_X(\mathcal{F}) =
\bigoplus\nolimits_{\eta \in X^0} j_{\eta, *}\mathcal{F}_\eta =
\prod\nolimits_{\eta \in X^0} j_{\eta, *}\mathcal{F}_\eta
$$
即 $\mathcal{F}$ 的亚纯截面层，是拟凝聚的；并且
\item $X$ 的有理函数环就是 $X$ 的亚纯函数环；写成公式即
$R(X) = \Gamma(X, \mathcal{K}_X)$。
\end{enumerate}
\end{lemma}

\begin{proof}
本引理是引理 \ref{divisors-lemma-meromorphic-weakass-finite} 的特例，因为在
局部 Noether 情形，由引理 \ref{divisors-lemma-ass-weakly-ass} 可知，
弱伴随点与伴随点相同。
\end{proof}

\begin{lemma}
\label{divisors-lemma-regular-meromorphic-section-exists-noetherian}
设 $X$ 为没有嵌入点的局部 Noether 概形。设 $\mathcal{L}$ 为可逆
$\mathcal{O}_X$-模。则 $\mathcal{L}$ 有一个正则亚纯截面。
\end{lemma}

\begin{proof}
对每个泛点 $\eta$（作为 $X$ 的泛点），选取生成元 $s_\eta$；它生成秩为
$1$ 的自由模 $\mathcal{L}_\eta$，后者是 Artin 局部环
$\mathcal{O}_{X, \eta}$ 上的模。由引理 \ref{divisors-lemma-quasi-coherent-K} 中对
$\mathcal{K}_X$ 与 $\mathcal{K}_X(\mathcal{L})$ 的描述，立即可知
$s = \prod s_\eta$ 是 $\mathcal{L}$ 的一个正则亚纯截面。
\end{proof}

\begin{lemma}
\label{divisors-lemma-make-maps-regular-section}
假设给定：
\begin{enumerate}
\item 局部 Noether 概形 $X$；
\item 可逆模 $\mathcal{L}$（作为 $\mathcal{O}_X$-模）；
\item 正则亚纯截面 $s$（取值于 $\mathcal{L}$）；以及
\item 凝聚层 $\mathcal{F}$（定义在 $X$ 上），它没有嵌入伴随点且
$\text{Supp}(\mathcal{F}) = X$。
\end{enumerate}
设 $\mathcal{I} \subset \mathcal{O}_X$ 为 $s$ 的分母理想层。
令 $T \subset X$ 为 $\mathcal{O}_X/\mathcal{I}$ 与
$\mathcal{L}/s(\mathcal{I})$ 的支集之并；根据引理
\ref{divisors-lemma-regular-meromorphic-ideal-denominators}，它是无处稠密闭子集
$T \subset X$。则存在典范单射
$$
1 : \mathcal{I}\mathcal{F} \to \mathcal{F}, \quad
s : \mathcal{I}\mathcal{F} \to \mathcal{F} \otimes_{\mathcal{O}_X}\mathcal{L}
$$
二者的余核都支撑在 $T$ 上。
\end{lemma}

\begin{proof}
约化到仿射情形，其中 $\mathcal{L} \cong \mathcal{O}_X$，且
$s = a/b$，这里 $a, b \in A$ 都是非零因子。略去约化步骤的证明。
写 $\mathcal{F} = \widetilde{M}$。令
$I = \{x \in A \mid x(a/b) \in A\}$，于是
$\mathcal{I} = \widetilde{I}$（见引理
\ref{divisors-lemma-regular-meromorphic-ideal-denominators} 的证明）。
注意，$T = V(I) \cup V((a/b)I)$。对任意 $A$-模 $M$，考察映射
$1 : IM \to M$；它给出引理中的映射 $1$。另一方面，考察映射
$\sigma : IM \to M_b, x \mapsto ax/b$。由于 $b$ 不是 $A$ 中的零因子，
而且 $M$ 的支集是 $\Spec(A)$ 且没有嵌入素理想，可知
$b$ 也是 $M$ 上的非零因子。因此，$M \subset M_b$。
由 $I$ 的定义，有 $\sigma(IM) \subset M$；这里把两者看作 $M_b$ 的子模。
因此，得到 $A$-模映射 $s : IM \to M$（即满足
$s(z)/1 = \sigma(z)$ 在 $M_b$ 中对所有 $z \in IM$ 成立的唯一映射）。
它是单射，因为 $a$ 也是非零因子（在 $A$ 与 $M$ 上都是）。
显然，$M/IM$ 被 $I$ 零化，而 $M/s(IM)$ 被 $(a/b)I$ 零化。
引理得证。
\end{proof}




\section{亚纯函数与亚纯截面：约化情形}
\label{divisors-section-meromorphic-reduced}

\noindent
对于约化且局部只有有限多个不可约分支的概形，亚纯函数层是拟凝聚的。

\begin{lemma}
\label{divisors-lemma-reduced-finite-irreducible}
设 $X$ 为约化概形，并且每个拟紧开集都只有有限多个不可约分支。
令 $X^0$ 为 $X$ 的不可约分支的泛点集合。则有
$$
\mathcal{K}_X =
\bigoplus\nolimits_{\eta \in X^0} j_{\eta, *}\kappa(\eta) =
\prod\nolimits_{\eta \in X^0} j_{\eta, *}\kappa(\eta)
$$
其中 $j_\eta : \Spec(\kappa(\eta)) \to X$ 是《概形》第
\ref{schemes-section-points} 节的典范映射。此外：
\begin{enumerate}
\item $\mathcal{K}_X$ 是拟凝聚 $\mathcal{O}_X$-代数层；
\item 对每个拟凝聚 $\mathcal{O}_X$-模 $\mathcal{F}$，层
$$
\mathcal{K}_X(\mathcal{F}) =
\bigoplus\nolimits_{\eta \in X^0} j_{\eta, *}\mathcal{F}_\eta =
\prod\nolimits_{\eta \in X^0} j_{\eta, *}\mathcal{F}_\eta
$$
即 $\mathcal{F}$ 的亚纯截面层，是拟凝聚的；
\item $\mathcal{S}_x \subset \mathcal{O}_{X, x}$ 是非零因子集合，
对任意 $x \in X$；
\item $\mathcal{K}_{X, x}$ 是 $\mathcal{O}_{X, x}$ 的全商环，
对任意 $x \in X$；
\item $\mathcal{K}_X(U)$ 等于 $\mathcal{O}_X(U)$ 的全商环，
对任意仿射开集 $U \subset X$；
\item $X$ 的有理函数环就是 $X$ 的亚纯函数环；写成公式即
$R(X) = \Gamma(X, \mathcal{K}_X)$。
\end{enumerate}
\end{lemma}

\begin{proof}
本引理是引理 \ref{divisors-lemma-meromorphic-weakass-finite} 的特例，因为由引理
\ref{divisors-lemma-weakass-reduced}，约化概形的弱伴随点就是各不可约分支的泛点。
\end{proof}

\begin{lemma}
\label{divisors-lemma-reduced-normalization}
设 $X$ 为概形。假设 $X$ 约化，并且每个拟紧开集
$U \subset X$ 都只有有限多个不可约分支。则正规化态射
$\nu : X^\nu \to X$ 是态射
$$
\underline{\Spec}_X(\mathcal{O}') \longrightarrow X
$$
其中 $\mathcal{O}' \subset \mathcal{K}_X$ 是 $\mathcal{O}_X$
在亚纯函数层中的整闭包。
\end{lemma}

\begin{proof}
将正规化态射 $\nu : X^\nu \to X$ 的定义（见《态射》定义
\ref{morphisms-definition-normalization}）与上面引理
\ref{divisors-lemma-reduced-finite-irreducible} 中对 $\mathcal{K}_X$ 的描述比较即可。
\end{proof}

\begin{lemma}
\label{divisors-lemma-meromorphic-functions-integral-scheme}
设 $X$ 为整概形，泛点为 $\eta$。则：
\begin{enumerate}
\item 亚纯函数层同构于取值为 $X$ 的函数域（见《态射》定义
\ref{morphisms-definition-function-field}）的常值层。
\item 对任意拟凝聚层 $\mathcal{F}$（定义在 $X$ 上），层
$\mathcal{K}_X(\mathcal{F})$ 同构于取值为 $\mathcal{F}_\eta$ 的常值层。
\end{enumerate}
\end{lemma}

\begin{proof}
略。
\end{proof}

\noindent
在某些情形，可以证明正则亚纯截面存在。

\begin{lemma}
\label{divisors-lemma-regular-meromorphic-section-exists}
设 $X$ 为概形。设 $\mathcal{L}$ 为可逆 $\mathcal{O}_X$-模。
在下列每种情形，$\mathcal{L}$ 都有一个正则亚纯截面：
\begin{enumerate}
\item $X$ 整；
\item $X$ 约化，且每个拟紧开集都只有有限多个不可约分支；
\item $X$ 局部 Noether 且没有嵌入点。
\end{enumerate}
\end{lemma}

\begin{proof}
在情形 (1)，令 $\eta \in X$ 为泛点。由引理
\ref{divisors-lemma-meromorphic-functions-integral-scheme} 已知，
$\mathcal{K}_X$ 与 $\mathcal{K}_X(\mathcal{L})$ 分别是取值为
$\kappa(\eta)$ 与 $\mathcal{L}_\eta$ 的常值层。由于
$\dim_{\kappa(\eta)} \mathcal{L}_\eta = 1$，可以选取非零元
$s \in \mathcal{L}_\eta$。显然，$s$ 是 $\mathcal{L}$ 的正则亚纯截面。
在情形 (2)，选取非零元 $s_\eta \in \mathcal{L}_\eta$；这里
$\eta$ 遍历 $X$ 的所有泛点。这是可能的，因为 $\mathcal{L}_\eta$ 是
$1$ 维 $\kappa(\eta)$-向量空间。由引理
\ref{divisors-lemma-reduced-finite-irreducible} 中对
$\mathcal{K}_X$ 与 $\mathcal{K}_X(\mathcal{L})$ 的描述，立即可知
$s = \prod s_\eta$ 是 $\mathcal{L}$ 的正则亚纯截面。情形 (3) 是引理
\ref{divisors-lemma-regular-meromorphic-section-exists-noetherian}。
\end{proof}




\section{Weil 除子}
\label{divisors-section-Weil-divisors}

\noindent
下面对局部 Noether 整概形引入 Weil 除子以及 Weil 除子的有理等价。
由于不假设概形拟紧，在定义 Weil 除子时必须稍加谨慎。必须允许
素除子的无穷和，例如，有理函数可能有无穷多个极点。对于拟紧概形，
这里的 Weil 除子仍像通常一样是有限和。下面这个基本引理常用于证明
闭子概形族局部有限。

\begin{lemma}
\label{divisors-lemma-components-locally-finite}
设 $X$ 为局部 Noether 概形。设 $Z \subset X$ 为闭子概形。
$Z$ 的不可约分支族在 $X$ 中局部有限。
\end{lemma}

\begin{proof}
设 $U \subset X$ 为拟紧开子概形。则 $U$ 是 Noether 概形，因而其
底拓扑空间 Noether（《性质》引理
\ref{properties-lemma-Noetherian-topology}）。因此，每个子空间都
Noether，并且只有有限多个不可约分支（见《拓扑》引理
\ref{topology-lemma-Noetherian}）。
\end{proof}

\noindent
回忆，若 $Z$ 是概形 $X$ 的不可约闭子集，则 $Z$ 在 $X$ 中的余维
等于局部环 $\mathcal{O}_{X, \xi}$ 的维数，其中 $\xi \in Z$ 是泛点。
见《性质》引理 \ref{properties-lemma-codimension-local-ring}。

\begin{definition}
\label{divisors-definition-Weil-divisor}
设 $X$ 为局部 Noether 整概形。
\begin{enumerate}
\item {\it 素除子}是整闭子概形 $Z \subset X$，其余维为 $1$。
\item {\it Weil 除子}是形式和 $D = \sum n_Z Z$，其中求和遍历
$X$ 的素除子，并且族 $\{Z \mid n_Z \not = 0\}$ 局部有限
（《拓扑》定义 \ref{topology-definition-locally-finite}）。
\end{enumerate}
$X$ 上所有 Weil 除子组成的群记作 $\text{Div}(X)$。
\end{definition}

\noindent
下一步是定义与有理函数相伴的 Weil 除子。为此，使用有理函数沿
素除子的消失阶，其定义如下。

\begin{definition}
\label{divisors-definition-order-vanishing}
设 $X$ 为局部 Noether 整概形。设 $f \in R(X)^*$。
对每个素除子 $Z \subset X$，把{\it $f$ 沿 $Z$ 的消失阶}定义为整数
$$
\text{ord}_Z(f) = \text{ord}_{\mathcal{O}_{X, \xi}}(f)
$$
其中右端是《代数》定义 \ref{algebra-definition-ord} 中的概念，
而 $\xi$ 是 $Z$ 的泛点。
\end{definition}

\noindent
注意，对 $f, g \in R(X)^*$，有
$$
\text{ord}_Z(fg) = \text{ord}_Z(f) + \text{ord}_Z(g).
$$
当然可能出现 $\text{ord}_Z(f) < 0$。在这种情形，称 $f$ 沿 $Z$
有一个{\it 极点}，并称 $-\text{ord}_Z(f) > 0$ 为
{\it $f$ 沿 $Z$ 的极点阶}。务必注意，条件
$\text{ord}_Z(f) \geq 0$ 与条件 $f \in \mathcal{O}_{X, \xi}$
并不等价，除非局部环 $\mathcal{O}_{X, \xi}$ 是离散赋值环。

\begin{lemma}
\label{divisors-lemma-divisor-locally-finite}
设 $X$ 为局部 Noether 整概形。设 $f \in R(X)^*$。则族
$$
\{Z \subset X \mid Z\text{ a prime divisor with generic point }\xi
\text{ 且 }f\text{ not in }\mathcal{O}_{X, \xi}\}
$$
与族
$$
\{Z \subset X \mid Z \text{ a prime divisor and }\text{ord}_Z(f) \not = 0\}
$$
在 $X$ 中局部有限。
\end{lemma}

\begin{proof}
存在非空开子概形 $U \subset X$，使得 $f$ 对应于
$\Gamma(U, \mathcal{O}_X^*)$ 的一个截面。因此，引理中两个集合里
可能出现的素除子全都是 $X \setminus U$ 的不可约分支。
于是引理 \ref{divisors-lemma-components-locally-finite} 给出所需结论。
\end{proof}

\noindent
该引理使下述定义有意义。

\begin{definition}
\label{divisors-definition-principal-divisor}
设 $X$ 为局部 Noether 整概形。设 $f \in R(X)^*$。
与 $f$ 相伴的{\it 主 Weil 除子}是 Weil 除子
$$
\text{div}(f) = \text{div}_X(f) = \sum \text{ord}_Z(f) [Z]
$$
其中求和遍历素除子，$\text{ord}_Z(f)$ 如定义
\ref{divisors-definition-order-vanishing} 所述。由引理
\ref{divisors-lemma-divisor-locally-finite}，这是有意义的。
\end{definition}

\begin{lemma}
\label{divisors-lemma-div-additive}
设 $X$ 为局部 Noether 整概形。设 $f, g \in R(X)^*$。则
$$
\text{div}_X(fg) = \text{div}_X(f) + \text{div}_X(g)
$$
作为 $X$ 上的 Weil 除子成立。
\end{lemma}

\begin{proof}
这由 $\text{ord}$ 函数的可加性立即得出。
\end{proof}

\noindent
由上引理可知，主 Weil 除子族构成全体 Weil 除子群的一个子群。
这引出下述定义。

\begin{definition}
\label{divisors-definition-class-group}
设 $X$ 为局部 Noether 整概形。$X$ 的{\it Weil 除子类群}是
Weil 除子群对主 Weil 除子子群的商。记作 $\text{Cl}(X)$。
\end{definition}

\noindent
由构造得到正合复形
\begin{equation}
\label{divisors-equation-Weil-divisor-class}
R(X)^* \xrightarrow{\text{div}} \text{Div}(X) \to \text{Cl}(X) \to 0
\end{equation}
可以把它看作 $\text{Cl}(X)$ 的一个呈示。下一步将把 Weil 除子类群
与 Picard 群联系起来。




\section{与可逆模相伴的 Weil 除子类}
\label{divisors-section-c1}

\noindent
本节将完全沿用第 \ref{divisors-section-Weil-divisors} 节的步骤，为局部 Noether
整概形定义典范映射
$\Pic(X) \to \text{Cl}(X)$。

\medskip\noindent
设 $X$ 为概形。设 $\mathcal{L}$ 为可逆 $\mathcal{O}_X$-模。
设 $\xi \in X$ 为一点。若
$s_\xi, s'_\xi \in \mathcal{L}_\xi$ 生成 $\mathcal{L}_\xi$（这里把它视为
$\mathcal{O}_{X, \xi}$-模），则存在单位 $u \in \mathcal{O}_{X, \xi}^*$，使得
$s_\xi = u s'_\xi$。亚纯截面层 $\mathcal{K}_X(\mathcal{L})$
（即 $\mathcal{L}$ 的亚纯截面层）在 $x$ 处的茎等于
$\mathcal{K}_{X, x} \otimes_{\mathcal{O}_{X, x}} \mathcal{L}_x$。
因此，任意亚纯截面 $s$（取值于 $\mathcal{L}$）在 $x$ 处茎中的像都可写成
$s = fs_\xi$，其中 $f \in \mathcal{K}_{X, x}$。下文将把这简写为
$f = s/s_\xi$。此外，若 $X$ 整，则
$\mathcal{K}_{X, x} = R(X)$ 等于 $X$ 的函数域，所以
$s/s_\xi \in R(X)$。若 $s$ 是正则亚纯截面，则事实上
$s/s_\xi \in R(X)^*$。在整概形上，正则亚纯截面恰好就是非零亚纯截面。
最后可知，$s/s_\xi$ 不依赖于 $s_\xi$ 的选择，至多相差局部环
$\mathcal{O}_{X, x}$ 中一个单位的倍数。综合以上内容，下述定义有意义。

\begin{definition}
\label{divisors-definition-order-vanishing-meromorphic}
设 $X$ 为局部 Noether 整概形。设 $\mathcal{L}$ 为可逆
$\mathcal{O}_X$-模。设
$s \in \Gamma(X, \mathcal{K}_X(\mathcal{L}))$ 为 $\mathcal{L}$ 的
正则亚纯截面。对每个素除子 $Z \subset X$，把
{\it $s$ 沿 $Z$ 的消失阶}定义为整数
$$
\text{ord}_{Z, \mathcal{L}}(s)
= \text{ord}_{\mathcal{O}_{X, \xi}}(s/s_\xi)
$$
其中右端是《代数》定义 \ref{algebra-definition-ord} 中的概念，
$\xi \in Z$ 是泛点，而
$s_\xi \in \mathcal{L}_\xi$ 是一个生成元。
\end{definition}

\noindent
与主除子的情形一样，有下述引理。

\begin{lemma}
\label{divisors-lemma-divisor-meromorphic-locally-finite}
设 $X$ 为局部 Noether 整概形。设 $\mathcal{L}$ 为可逆
$\mathcal{O}_X$-模。设 $s \in \mathcal{K}_X(\mathcal{L})$ 为
$\mathcal{L}$ 的正则（即非零）亚纯截面。则集合
$$
\{Z \subset X \mid Z \text{ a prime divisor with generic point }\xi
\text{ 且 }s\text{ not in }\mathcal{L}_\xi\}
$$
与集合
$$
\{Z \subset X \mid Z \text{ is a prime divisor and }
\text{ord}_{Z, \mathcal{L}}(s) \not = 0\}
$$
在 $X$ 中局部有限。
\end{lemma}

\begin{proof}
存在非空开子概形 $U \subset X$，使得 $s$ 对应于
$\Gamma(U, \mathcal{L})$ 的一个截面，并且该截面生成
$\mathcal{L}$ 在 $U$ 上的限制。因此，引理中的集合里可能出现的素除子全都是
$X \setminus U$ 的不可约分支。于是引理
\ref{divisors-lemma-components-locally-finite} 给出所需结论。
\end{proof}

\begin{lemma}
\label{divisors-lemma-divisor-meromorphic-well-defined}
设 $X$ 为局部 Noether 整概形。设 $\mathcal{L}$ 为可逆
$\mathcal{O}_X$-模。设 $s, s' \in \mathcal{K}_X(\mathcal{L})$ 为
$\mathcal{L}$ 的非零亚纯截面。则 $f = s/s'$ 是 $R(X)^*$ 的元素，并且
$$
\sum \text{ord}_{Z, \mathcal{L}}(s)[Z]
=
\sum \text{ord}_{Z, \mathcal{L}}(s')[Z]
+
\text{div}(f)
$$
作为 Weil 除子成立。
\end{lemma}

\begin{proof}
这由定义立即得出。注意，引理
\ref{divisors-lemma-divisor-meromorphic-locally-finite}
保证这些和确实是 Weil 除子。
\end{proof}

\begin{definition}
\label{divisors-definition-divisor-invertible-sheaf}
设 $X$ 为局部 Noether 整概形。设 $\mathcal{L}$ 为可逆
$\mathcal{O}_X$-模。
\begin{enumerate}
\item 对任意非零亚纯截面 $s$（取值于 $\mathcal{L}$），把
{\it 与 $s$ 相伴的 Weil 除子}定义为
$$
\text{div}_\mathcal{L}(s) =
\sum \text{ord}_{Z, \mathcal{L}}(s) [Z] \in \text{Div}(X)
$$
其中求和遍历素除子。
\item 把{\it 与 $\mathcal{L}$ 相伴的 Weil 除子类}定义为
$\text{div}_\mathcal{L}(s)$ 在 $\text{Cl}(X)$ 中的像，其中
$s$ 是 $\mathcal{L}$ 在 $X$ 上的任意非零亚纯截面。由引理
\ref{divisors-lemma-divisor-meromorphic-well-defined}，这是良定义的。
\end{enumerate}
\end{definition}

\noindent
正如所料，该构造对可逆模可加。

\begin{lemma}
\label{divisors-lemma-c1-additive}
设 $X$ 为局部 Noether 整概形。设 $\mathcal{L}$、$\mathcal{N}$ 为可逆
$\mathcal{O}_X$-模。设 $s$ 与 $t$ 分别为
$\mathcal{L}$ 与 $\mathcal{N}$ 的非零亚纯截面。则 $st$ 是
$\mathcal{L} \otimes \mathcal{N}$ 的非零亚纯截面，并且
$$
\text{div}_{\mathcal{L} \otimes \mathcal{N}}(st)
=
\text{div}_\mathcal{L}(s) + \text{div}_\mathcal{N}(t)
$$
在 $\text{Div}(X)$ 中成立。特别地，
$\mathcal{L} \otimes_{\mathcal{O}_X} \mathcal{N}$ 的 Weil 除子类等于
$\mathcal{L}$ 与 $\mathcal{N}$ 的 Weil 除子类之和。
\end{lemma}

\begin{proof}
设 $s$ 与 $t$ 分别为 $\mathcal{L}$ 与 $\mathcal{N}$ 的非零亚纯截面。
则 $st$ 是 $\mathcal{L} \otimes \mathcal{N}$ 的非零亚纯截面。
设 $Z \subset X$ 为素除子。设 $\xi \in Z$ 为其泛点。选取生成元
$s_\xi \in \mathcal{L}_\xi$ 与 $t_\xi \in \mathcal{N}_\xi$。则
$s_\xi t_\xi$ 是 $(\mathcal{L} \otimes \mathcal{N})_\xi$ 的一个生成元。
所以 $st/(s_\xi t_\xi) = (s/s_\xi)(t/t_\xi)$。因而
$$
\text{ord}_{Z, \mathcal{L} \otimes \mathcal{N}}(st)
=
\text{ord}_{Z, \mathcal{L}}(s) + \text{ord}_{Z, \mathcal{N}}(t)
$$
这由 $\text{ord}_Z$ 函数的可加性得出。
\end{proof}

\noindent
这样得到交换群的同态
\begin{equation}
\label{divisors-equation-c1}
\Pic(X) \longrightarrow \text{Cl}(X)
\end{equation}
它把可逆模映到其 Weil 除子类。

\begin{lemma}
\label{divisors-lemma-normal-c1-injective}
设 $X$ 为局部 Noether 整概形。若 $X$ 正规，则映射
(\ref{divisors-equation-c1}) $\Pic(X) \to \text{Cl}(X)$ 是单射。
\end{lemma}

\begin{proof}
设 $\mathcal{L}$ 为可逆 $\mathcal{O}_X$-模，并且其相伴的 Weil 除子类
平凡。设 $s$ 为 $\mathcal{L}$ 的正则亚纯截面。该假设意味着
$\text{div}_\mathcal{L}(s) = \text{div}(f)$，其中
$f \in R(X)^*$。于是
$t = f^{-1}s$ 是 $\mathcal{L}$ 的正则亚纯截面，且
$\text{div}_\mathcal{L}(t) = 0$，见引理
\ref{divisors-lemma-divisor-meromorphic-well-defined}。下面证明 $t$ 给出
$\mathcal{L}$ 的平凡化，从而完成引理的证明。
为此，可以在 $X$ 上局部工作。因此，可以假设
$X = \Spec(A)$ 仿射且 $\mathcal{L}$ 平凡。于是 $A$ 是 Noether
正规整环，而 $t$ 是其分式域中的元素，并且满足
$\text{ord}_{A_\mathfrak p}(t) = 0$；所取素理想的高度为 $1$，记作
$\mathfrak p$，且它是 $A$ 的素理想。目标是证明 $t$ 是 $A$ 的单位。
由于 $A_\mathfrak p$ 对 $A$ 的高度一素理想而言是离散赋值环
（《代数》引理 \ref{algebra-lemma-criterion-normal}），所以该条件表示
$t \in A_\mathfrak p^*$ 对所有素理想 $\mathfrak p$（高度为 $1$）均成立。
由《代数》引理
\ref{algebra-lemma-normal-domain-intersection-localizations-height-1}，
这蕴含 $t \in A$ 且 $t^{-1} \in A$。证明完成。
\end{proof}

\begin{lemma}
\label{divisors-lemma-local-rings-UFD-c1-bijective}
设 $X$ 为局部 Noether 整概形。考察映射
(\ref{divisors-equation-c1}) $\Pic(X) \to \text{Cl}(X)$。下列条件等价：
\begin{enumerate}
\item $X$ 的局部环都是 UFD；并且
\item $X$ 正规，且 $\Pic(X) \to \text{Cl}(X)$ 为满射。
\end{enumerate}
在这种情形，$\Pic(X) \to \text{Cl}(X)$ 是同构。
\end{lemma}

\begin{proof}
若 (1) 成立，则由《代数》引理 \ref{algebra-lemma-UFD-normal}，
$X$ 正规。因此，映射 (\ref{divisors-equation-c1}) 由引理
\ref{divisors-lemma-normal-c1-injective} 可知为单射。此外，每个素除子 $D \subset X$ 都由引理
\ref{divisors-lemma-weil-divisor-is-cartier-UFD} 成为有效 Cartier 除子。在此情形，
典范截面 $1_D$（它是 $\mathcal{O}_X(D)$ 的截面）
（定义 \ref{divisors-definition-invertible-sheaf-effective-Cartier-divisor}）
恰好沿 $D$ 消失，可知 $D$ 的类是 $\mathcal{O}_X(D)$ 在映射
(\ref{divisors-equation-c1}) 下的像。因此，该映射也为满射。

\medskip\noindent
假设 (2) 成立。选取素除子 $D \subset X$。由于
(\ref{divisors-equation-c1}) 为满射，存在可逆层 $\mathcal{L}$、正则亚纯截面
$s$ 以及 $f \in R(X)^*$，使得
$\text{div}_\mathcal{L}(s) + \text{div}(f) = [D]$。换言之，
$\text{div}_\mathcal{L}(fs) = [D]$。设 $x \in X$，并令
$A = \mathcal{O}_{X, x}$。则 $A$ 是 Noether 局部正规整环，分式域为
$K = R(X)$。每个高度为 $1$ 的素理想（在 $A$ 中）都对应于 $X$ 上的素除子，
而每个可逆 $\mathcal{O}_X$-模限制到 $\Spec(A)$ 后都是平凡可逆模。
由此，对每个高度为 $1$ 的素理想 $\mathfrak p \subset A$，存在元素
$f \in K$，使得 $\text{ord}_{A_\mathfrak p}(f) = 1$，且
$\text{ord}_{A_{\mathfrak p'}}(f) = 0$ 对每个其他高度一素理想
$\mathfrak p'$ 成立。于是由《代数》引理
\ref{algebra-lemma-normal-domain-intersection-localizations-height-1}，
$f \in A$。同理可知，每个元素 $g \in \mathfrak p$ 都可写成
$g = af$，其中 $a \in A$。因此，$A$ 的每个高度一素理想都是主理想，
由《代数》引理 \ref{algebra-lemma-characterize-UFD-height-1}，$A$ 是 UFD。
\end{proof}




\section{关于可逆模的进一步结果}
\label{divisors-section-invertible-modules}

\noindent
本节讨论可逆模的一些性质。

\begin{lemma}
\label{divisors-lemma-in-image-pullback}
设 $\varphi : X \to Y$ 为概形态射。设 $\mathcal{L}$ 为可逆
$\mathcal{O}_X$-模。假设：
\begin{enumerate}
\item $X$ 局部 Noether；
\item $Y$ 局部 Noether、整且正规；
\item $\varphi$ 平坦，且其纤维整（因而非空）；
\item $\varphi$ 或拟紧，或局部有限型；
\item $\mathcal{L}$ 限制到 $\varphi$ 的泛纤维后平凡。
\end{enumerate}
则 $\mathcal{L} \cong \varphi^*\mathcal{N}$，其中右端由某个可逆
$\mathcal{O}_Y$-模 $\mathcal{N}$ 给出。
\end{lemma}

\begin{proof}
设 $\xi \in Y$ 为泛点。设 $X_\xi$ 为 $\varphi$ 在 $\xi$ 上的概形论纤维。
用 $\mathcal{L}_\xi$ 表示 $\mathcal{L}$ 到 $X_\xi$ 的拉回。假设 (5) 意味着
$\mathcal{L}_\xi$ 平凡。选取平凡化截面
$s \in \Gamma(X_\xi, \mathcal{L}_\xi)$。由引理
\ref{divisors-lemma-flat-over-integral-integral-fibre}，$X$ 是整概形。因此可以把
$s$ 看作 $\mathcal{L}$ 的正则亚纯截面。由引理
\ref{divisors-lemma-pullback-meromorphic-sections-defined}，亚纯函数沿
$\varphi$ 的拉回有定义。令 $\mathcal{N} \subset \mathcal{K}_Y$ 为如下
$\mathcal{O}_Y$-模：其在开集 $V \subset Y$ 上的截面，是那些亚纯函数
$g \in \mathcal{K}_Y(V)$，它们满足
$\varphi^*(g)s \in \mathcal{L}(\varphi^{-1}V)$。先验地，$\varphi^*(g)s$ 是
$\mathcal{K}_X(\mathcal{L})$ 在 $\varphi^{-1}V$ 上的截面。我们断言
$\mathcal{N}$ 是可逆 $\mathcal{O}_Y$-模，并且映射
$$
\varphi^*\mathcal{N} \longrightarrow \mathcal{L},\quad
g \longmapsto gs
$$
是同构。

\medskip\noindent
先在下述情形证明该断言：$X$ 与 $Y$ 均仿射，且 $\mathcal{L}$ 平凡。
记 $Y = \Spec(R)$、$X = \Spec(A)$，而 $s$ 由元素
$s \in A \otimes_R K$ 给出，其中 $K$ 是 $R$ 的分式域。可写
$s = a/r$，其中 $r \in R$ 非零且 $a \in A$。由于 $s$ 在泛纤维上生成
$\mathcal{L}$，存在 $s' \in A \otimes_R K$，使得 $ss' = 1$。所以
$s = r'/a'$，其中 $r' \in R$ 非零且 $a' \in A$。设
$\mathfrak p_1, \ldots, \mathfrak p_n \subset R$ 为包含 $rr'$ 的极小素理想。
每个 $R_{\mathfrak p_i}$ 都是离散赋值环（《代数》引理
\ref{algebra-lemma-minimal-over-1} 与
\ref{algebra-lemma-criterion-normal}）。由假设，
$\mathfrak q_i = \mathfrak p_i A$ 是素理想。因此
$\mathfrak q_i A_{\mathfrak q_i}$ 由单个元素生成，进而
$A_{\mathfrak q_i}$ 也是离散赋值环（《代数》引理
\ref{algebra-lemma-characterize-dvr}）。当然，
$R_{\mathfrak p_i} \to A_{\mathfrak q_i}$ 的分歧指数为 $1$。设
$e_i, e'_i \geq 0$ 分别为 $a, a'$ 在 $A_{\mathfrak q_i}$ 中的赋值。
则 $e_i + e'_i$ 是 $rr'$ 在 $R_{\mathfrak p_i}$ 中的赋值。注意，由
《代数》引理
\ref{algebra-lemma-normal-domain-intersection-localizations-height-1}，在
$R$ 中有
$$
\mathfrak p_1^{(e_1 + e'_1)} \cap \ldots \cap \mathfrak p_i^{(e_n + e'_n)} =
(rr')
$$
令
$$
I = \mathfrak p_1^{(e_1)} \cap \ldots \cap \mathfrak p_i^{(e_n)}
\quad\text{且}\quad
I' = \mathfrak p_1^{(e'_1)} \cap \ldots \cap \mathfrak p_i^{(e'_n)}
$$
于是 $II' \subset (rr')$。注意，由《代数》引理
\ref{algebra-lemma-symbolic-power-flat-extension} 与
\ref{algebra-lemma-flat-intersect-ideals}，有
$$
IA =
(\mathfrak p_1^{(e_1)} \cap \ldots \cap \mathfrak p_i^{(e_n)})A =
(\mathfrak p_1A)^{(e_1)} \cap \ldots \cap (\mathfrak p_i A)^{(e_n)}
$$
$I'A$ 也同样如此。因此 $a \in IA$ 且 $a' \in I'A$。由此，
$IA \otimes_A I'A \to rr'A$ 为满射。利用 $R \to A$ 的忠实平坦性，
$I \otimes_R I' \to (rr')$ 也为满射。于是 $II' = (rr')$，并且
$I$ 与 $I'$ 都是秩为 $1$ 的有限局部自由模，见《代数》引理
\ref{algebra-lemma-product-ideals-principal}。因此在 $R$ 上 Zariski 局部地，
可写成 $I = (g)$ 与 $I' = (g')$，其中 $gg' = rr'$。于是
$a = ug$ 且 $a' = u'g'$，其中 $u, u' \in A$。由此可知 $u, u'$ 都是单位。
所以在 $R$ 上 Zariski 局部地，有 $s = ug/r$，该情形的断言得证。

\medskip\noindent
设 $y \in Y$ 为一点。选取 $x \in X$，使它映到 $y$。可以把上一段的结果
应用于 $\Spec(\mathcal{O}_{X, x}) \to \Spec(\mathcal{O}_{Y, y})$。
由此存在元素 $g \in R(Y)^*$；它在相差
$\mathcal{O}_{Y, y}^*$ 中元素之乘法的意义下是良定义的，并且
$\varphi^*(g)s$ 生成 $\mathcal{L}_x$。因而 $\varphi^*(g)s$ 生成
$\mathcal{L}$；这在某个邻域 $U$ 上成立，且 $x$ 属于该邻域。设 $x'$ 是位于
$y$ 上方的另一个点，并设 $g' \in R(Y)^*$ 使得 $\varphi^*(g')s$ 生成
$\mathcal{L}$；这在某个开邻域 $U'$ 上成立，且 $x'$ 属于该邻域。由于纤维不可约，可以选取一点
$x''$，它位于 $U \cap U' \cap \varphi^{-1}(\{y\})$ 中。由环映射
$\mathcal{O}_{Y, y} \to \mathcal{O}_{X, x''}$ 所对应的唯一性，可知
$g$ 与 $g'$ 相差 $\mathcal{O}_{Y, y}^*$ 中元素的乘法。因此可知
$\varphi^*(g)s$ 生成 $\mathcal{L}$；这在 $\varphi^{-1}(y)$ 的某个开邻域上成立。
令 $Z \subset X$ 为如下点的集合：这些点
$z \in X$ 满足 $\varphi^*(g)s$ 不生成 $\mathcal{L}_z$。以上论证表明
$Z$ 是闭集，且有 $Z = \varphi^{-1}(T)$，其中 $T \subset Y$ 是某个子集，并且
$y \not \in T$。若能证明 $T$ 为闭集，则 $g$ 将生成 $\mathcal{N}$（作为
$\mathcal{O}_Y$-模）；这在开邻域 $Y \setminus T$ 上成立，而该邻域包含 $y$，从而完成证明
（略去一些细节）。

\medskip\noindent
若 $\varphi$ 拟紧，则由《态射》引理
\ref{morphisms-lemma-fpqc-quotient-topology}，$T$ 为闭集。若
$\varphi$ 局部有限型，则由《态射》引理
\ref{morphisms-lemma-fppf-open}，$\varphi$ 是开映射。于是
$Y \setminus T$ 作为开集 $X \setminus Z$ 的像是开集。
\end{proof}

\begin{lemma}
\label{divisors-lemma-closure-effective-cartier-divisor}
设 $X$ 为局部 Noether 概形。设 $U \subset X$ 为开集，并设
$D \subset U$ 为有效 Cartier 除子。若 $\mathcal{O}_{X, x}$ 对所有
$x \in X \setminus U$ 都是 UFD，则存在有效 Cartier 除子
$D' \subset X$，使得 $D = U \cap D'$。
\end{lemma}

\begin{proof}
设 $D' \subset X$ 为态射 $D \to X$ 的概形论像。由于 $X$ 局部 Noether，
态射 $D \to X$ 拟紧，见《性质》引理
\ref{properties-lemma-immersion-into-noetherian}。因此由《态射》引理
\ref{morphisms-lemma-quasi-compact-scheme-theoretic-image}，$D'$ 的形成与
取 $X$ 中的开集可交换。故可假设 $X = \Spec(A)$ 仿射。设
$I \subset A$ 为与 $D'$ 对应的理想。设素理想 $\mathfrak p \subset A$
对应于 $X \setminus U$ 中的一点。为完成证明，只需证明
$I_\mathfrak p$ 由单个元素生成，见引理
\ref{divisors-lemma-effective-Cartier-in-points}。因此由《态射》引理
\ref{morphisms-lemma-flat-base-change-scheme-theoretic-image}，可将 $X$ 替换为
$\Spec(A_\mathfrak p)$。换言之，可以假设 $X$ 是局部 UFD $A$ 的谱。
于是 $A$ 的所有局部环都是 UFD。由引理
\ref{divisors-lemma-effective-Cartier-divisor-is-a-sum}，有
$D = \sum a_i D_i$，其中 $D_i \subset U$ 是整的有效 Cartier 除子。
泛点 $\xi_i$（属于 $D_i$）对应于素理想
$\mathfrak p_i \subset A$，其高度为 $1$，见引理
\ref{divisors-lemma-effective-Cartier-codimension-1}。于是有
$\mathfrak p_i = (f_i)$，其中 $f_i \in A$ 是某个素元；从而 $D'$ 由
$\prod f_i^{a_i}$ 截出，正如所需。
\end{proof}

\begin{lemma}
\label{divisors-lemma-extend-invertible-module}
设 $X$ 为局部 Noether 概形。设 $U \subset X$ 为开集，并设
$\mathcal{L}$ 为可逆 $\mathcal{O}_U$-模。若 $\mathcal{O}_{X, x}$ 对所有
$x \in X \setminus U$ 都是 UFD，则存在可逆
$\mathcal{O}_X$-模 $\mathcal{L}'$，使得
$\mathcal{L} \cong \mathcal{L}'|_U$。
\end{lemma}

\begin{proof}
选取 $x \in X$，且 $x \not \in U$。我们将证明存在仿射开邻域
$W \subset X$，使得 $\mathcal{L}|_{W \cap U}$ 可延拓为 $W$ 上的可逆层。
由层的黏合（《层》节 \ref{sheaves-section-glueing-sheaves}），这蕴含可把
$\mathcal{L}$ 延拓到严格更大的开集 $U \cup W$。令
$W = \Spec(A)$ 为一个仿射开邻域。由于 $U \cap W$ 拟仿射，由引理
\ref{divisors-lemma-quasi-projective-Noetherian-pic-effective-Cartier}，可将
$\mathcal{L}|_{W \cap U}$ 写成
$\mathcal{O}(D_1) \otimes \mathcal{O}(D_2)^{\otimes -1}$，其中
$D_1, D_2 \subset W \cap U$ 是某些有效 Cartier 除子。再由引理
\ref{divisors-lemma-closure-effective-cartier-divisor}，$D_1$ 与 $D_2$ 可延拓为
$W$ 上的有效 Cartier 除子，这就给出可逆层的延拓。

\medskip\noindent
若 $X$ Noether（这是实践中最常用的情形），上述结论结合 Noether 归纳法
即完成证明。一般情形如下处理。首先，由于 $U$ 外每一点的局部环都是整环，
且 $X$ 局部 Noether，可知 $U$ 在 $X$ 中的闭包是开集。因此可以假设
$U \subset X$ 稠密且概形论稠密。现在考虑集合 $T$，其元素为三元组
$(U', \mathcal{L}', \alpha)$，其中
$U \subset U' \subset X$ 是开子概形，$\mathcal{L}'$ 是可逆
$\mathcal{O}_{U'}$-模，而
$\alpha : \mathcal{L}'|_U \to \mathcal{L}$ 是同构。在 $T$ 上赋予偏序
$\leq$，定义规则为：
$(U', \mathcal{L}', \alpha) \leq (U'', \mathcal{L}'', \alpha')$
当且仅当 $U' \subset U''$，且存在同构
$\beta : \mathcal{L}''|_{U'} \to \mathcal{L}'$，它与 $\alpha$ 和 $\alpha'$ 相容。
注意，$\beta$（若存在）是唯一的，因为 $U \subset X$ 稠密。证明的第一部分
表明，对任意元素 $t = (U', \mathcal{L}', \alpha)$（属于 $T$），若
$U' \not = X$，则存在 $t' \in T$ 使得 $t' > t$。因此为完成证明，只需说明
Zorn 引理适用。为此，考虑全序子集 $I \subset T$。若 $i \in I$ 对应于
三元组 $(U_i, \mathcal{L}_i, \alpha_i)$，则可如下构造可逆模
$\mathcal{L}'$，它定义在 $U' = \bigcup U_i$ 上。对开且拟紧的
$W \subset U'$，可知 $W \subset U_i$ 对某个 $i$ 成立，并置
$$
\mathcal{L}'(W) = \mathcal{L}_i(W)
$$
转移映射使用各个 $\beta$（它们唯一，因而复合相容）。这就定义了可逆
$\mathcal{O}$-模 $\mathcal{L}'$，定义所在的基是 $U'$ 的拟紧开集；这足以定义一个可逆模
（《层》节 \ref{sheaves-section-bases}）。细节从略。
\end{proof}

\begin{lemma}
\label{divisors-lemma-open-subscheme-UFD}
设 $R$ 为 UFD。下列概形的 Picard 群平凡：
\begin{enumerate}
\item $\Spec(R)$ 及其任意开子概形；
\item $\mathbf{A}^n_R = \Spec(R[x_1, \ldots, x_n])$ 及其任意开子概形。
\end{enumerate}
特别地，仿射 $n$-空间 $\mathbf{A}^n_k$ 在域 $k$ 上的任意开子概形之
Picard 群平凡。
\end{lemma}

\begin{proof}
由《代数》引理 \ref{algebra-lemma-polynomial-ring-UFD}，$R$ 上的任意多项式环
都是 UFD。因此，问题归结为证明 $\Pic(U)$ 平凡，其中
$U \subset \Spec(R)$ 是（非空）开集。

\medskip\noindent
在 Noether 情形，可用引理 \ref{divisors-lemma-extend-invertible-module} 将问题约化为证明
$\Pic(\Spec(R))$ 平凡；而后者已在《代数续篇》引理
\ref{more-algebra-lemma-UFD-Pic-trivial} 中证明。我们建议读者跳过证明的其余部分。

\medskip\noindent
设 $U \subset \Spec(R)$ 为非空开集。选取非零 $f \in R$，使得
$D(f) \subset U$。由于 $f$ 是 $R$ 中素元的乘积，可知只有有限多个素元
$a_1, \ldots, a_n \in R$ 满足 $(a_i) \not \in U$（或等价地，
$U \cap V(a_i) = \emptyset$）。将 $R$ 替换为
$R_{a_1 \ldots a_n}$ 后，可假设 $U$ 包含所有主素理想。（注意，UFD 的局部化
仍是 UFD，例如由《代数》引理
\ref{algebra-lemma-invert-prime-elements} 可知。）

\medskip\noindent
假设 $U$ 包含所有主素理想。写
$U = \Spec(R) \setminus V(I)$，其中 $I \subset R$ 是某个理想。选取非零 $f \in I$，并把
$f = a_1 \ldots a_n$ 写成素元的乘积。由于
$(a_i) \not \in V(I)$，由素理想回避，可选取 $g \in I$，使得
$g \not \in (a_i)$ 对所有 $i$ 均成立。图示为
$$
U' = D(f) \cup D(g) \subset U \subset \Spec(R)
$$
由《代数续篇》引理 \ref{more-algebra-lemma-UFD-Pic-trivial}，$D(f)$ 与
$D(g)$ 的 Picard 群都平凡；而由关于 UFD 的初等论证，映射
$R_f^* \times R_g^* \to R_{fg}^*$，$(u, v) \mapsto uv^{-1}$ 是满射。
由此可知 $\Pic(U')$ 平凡。

\medskip\noindent
因此只需证明：可逆模 $\mathcal{L}$（定义在 $U$ 上）若沿
$j : U' \to U'$ 拉回后平凡，则它平凡。为此，只需证明映射
$\mathcal{L} \to j_*j^*\mathcal{L}$ 是同构。进一步，设
$u \in U \setminus U'$。则
$f, g \in \mathfrak m_u \subset \mathcal{O}_{U, u} = A$ 是局部 UFD 中的元素，
且没有共同素因子。读者容易由考察素因子证明，这意味着
$$
0 \to A \to A_f \times A_g \to A_{fg}
$$
正合。由于 $A = \mathcal{L}_u$（这里使用 $\mathcal{L}$ 可逆），又因为
$(j_*j^*\mathcal{L})_u$ 是 $A_f \times A_g \to A_{fg}$ 的核，结论随即得出。
\end{proof}

\begin{lemma}
\label{divisors-lemma-Pic-projective-space-UFD}
设 $R$ 为 UFD。$\mathbf{P}^n_R$ 的 Picard 群是 $\mathbf{Z}$。更精确地，
存在同构
$$
\mathbf{Z} \longrightarrow \Pic(\mathbf{P}^n_R),\quad
m \longmapsto \mathcal{O}_{\mathbf{P}^n_R}(m)
$$
特别地，射影空间 $\mathbf{P}^n_k$ 在域 $k$ 上的 Picard 群是
$\mathbf{Z}$。
\end{lemma}

\begin{proof}
由引理 \ref{divisors-lemma-open-subscheme-UFD}，开集
$D_+(T_i) \cong \mathbf{A}^n_R$ 的 Picard 群平凡。因此若
$\mathcal{L}$ 是 $\mathbf{P}^n_R$ 上的可逆模，则对于开覆盖
$$
\mathbf{P}^n_R = \bigcup\nolimits_{i = 0, \ldots, n} D_+(T_i)
$$
它由一个 $1$-上闭链给出，该上闭链取值于可逆函子层 $\mathcal{O}^*$。注意，对
$i \not = j$，有
$$
\mathcal{O}^*(D_+(T_i) \cap D_+(T_j)) =
\mathcal{O}^*(D_+(T_iT_j)) = \left(R[T_0, \ldots, T_n]_{(T_iT_j)}\right)^* =
R^* \times (T_i/T_j)^\mathbf{Z}
$$
因此，这样的上闭链 $(g_{ij})$ 由单位
$$
g_{ij} = u_{ij} (T_i/T_j)^{e_{ij}}
$$
给出，其中 $u_{ij} \in R^*$ 且 $e_{ij} \in \mathbf{Z}$，并满足上闭链条件。
$D_+(T_iT_jT_k)$ 上的上闭链条件在 $\#\{i, j, k\} = 3$ 时告诉我们
$$
u_{ik} = u_{ij} u_{jk} \quad\text{且}\quad
e_{ik} = e_{ij} = e_{jk}
$$
于是 $u_{ij} = u_{i1} u_{j1}^{-1}$ 是一个上边界。因此，可逆模的所有同构类
都由取如下上闭链得到：
$$
g_{ij} = (T_i/T_j)^e
$$
其中 $e \in \mathbf{Z}$。由于 $\mathcal{O}(n)$ 有平凡化截面
$T_i^n$（定义在 $D_+(T_i)$ 上），可知 $\mathcal{O}(n)$ 所对应的上闭链是
$(T_i/T_j)^n$。证明完成。
\end{proof}


\section{正规概形上的 Weil 除子}
\label{divisors-section-weil-divisors-normal}

\noindent
首先讨论自反模的性质。

\begin{lemma}
\label{divisors-lemma-reflexive-normal}
设 $X$ 为整的局部 Noether 正规概形。对 $\mathcal{F}$ 与 $\mathcal{G}$（它们是
相干自反 $\mathcal{O}_X$-模），映射
$$
(\SheafHom_{\mathcal{O}_X}(\mathcal{F}, \mathcal{O}_X)
\otimes_{\mathcal{O}_X} \mathcal{G})^{**} \to
\SheafHom_{\mathcal{O}_X}(\mathcal{F}, \mathcal{G})
$$
是同构。规则 $\mathcal{F}, \mathcal{G} \mapsto
(\mathcal{F} \otimes_{\mathcal{O}_X} \mathcal{G})^{**}$
在秩为 $1$ 的相干自反 $\mathcal{O}_X$-模之同构类集合上定义一个交换群律。
\end{lemma}

\begin{proof}
虽非严格必要，但建议先阅读注
\ref{divisors-remark-tensor}，再继续此证明。选取开子概形 $j : U \to X$，使得
$X \setminus U$ 的每个不可约分支之余维均为 $\geq 2$（余维在 $X$ 中取），并且
$j^*\mathcal{F}$ 与 $j^*\mathcal{G}$ 都有限局部自由，见引理
\ref{divisors-lemma-reflexive-over-normal}。映射
$$
\SheafHom_{\mathcal{O}_U}(j^*\mathcal{F}, \mathcal{O}_U)
\otimes_{\mathcal{O}_U} j^*\mathcal{G} \to
\SheafHom_{\mathcal{O}_U}(j^*\mathcal{F}, j^*\mathcal{G})
$$
是同构，因为可在局部检验，而且当这些模有限自由时结论显然。注意，对引理中
显示的箭头应用 $j^*$，所得正是刚刚证明为同构的箭头（略去一个小细节）。由于
$j^*$ 在 $U$ 上的相干自反模与 $X$ 上的相干自反模之间定义范畴等价
（由引理 \ref{divisors-lemma-reflexive-S2-extend} 及 Serre 判别法《性质》引理
\ref{properties-lemma-criterion-normal}），可知引理中的箭头也是同构。
若 $\mathcal{F}$ 的秩为 $1$，则 $j^*\mathcal{F}$ 是可逆
$\mathcal{O}_U$-模，而自反模
$\mathcal{F}^\vee = \SheafHom(\mathcal{F}, \mathcal{O}_X)$
限制后是它的逆。与前面同理可得
$(\mathcal{F} \otimes_{\mathcal{O}_X} \mathcal{F}^\vee)^{**} = \mathcal{O}_X$。
这样便得到引理陈述中群律的逆元。
\end{proof}

\begin{lemma}
\label{divisors-lemma-normal-class-group}
设 $X$ 为整的局部 Noether 正规概形。秩为 $1$ 的相干自反
$\mathcal{O}_X$-模所成的群同构于 Weil 除子类群 $\text{Cl}(X)$（属于 $X$）。
\end{lemma}

\begin{proof}
设 $\mathcal{F}$ 为秩为 $1$ 的相干自反 $\mathcal{O}_X$-模。选取开集
$U \subset X$，使得 $X \setminus U$ 的每个不可约分支之余维均为
$\geq 2$（余维在 $X$ 中取），并且 $\mathcal{F}|_U$ 可逆，见引理
\ref{divisors-lemma-reflexive-over-normal}。注意，$\text{Cl}(U) = \text{Cl}(X)$，因为
$X$ 的 Weil 除子类群只依赖于其有理函数域、余维为 $1$ 的点及其局部环。
因此可以把 $\mathcal{F}$ 的 Weil 除子类定义为 $\mathcal{F}|_U$ 在
$\text{Cl}(U)$ 中的 Weil 除子类。我们略去它不依赖于 $U$ 之选择的验证。

\medskip\noindent
用 $\text{Cl}'(X)$ 表示秩为 $1$ 的相干自反 $\mathcal{O}_X$-模之同构类集合。
上述构造给出群同态
$$
\text{Cl}'(X) \longrightarrow \text{Cl}(X)
$$
这是因为，对任意一对 $\mathcal{F}, \mathcal{G}$（它们属于 $\text{Cl}'(X)$），
可以选取一个对二者都适用的 $U$，而将可逆模映到其 Weil 除子类的赋值
(\ref{divisors-equation-c1}) 是同态。若 $\mathcal{F}$ 属于该映射的核，则由引理
\ref{divisors-lemma-normal-c1-injective}，$\mathcal{F}|_U$ 平凡，继而由引理
\ref{divisors-lemma-reflexive-S2-extend} 及 Serre 判别法《性质》引理
\ref{properties-lemma-criterion-normal}，$\mathcal{F}$ 也平凡。为完成证明，
只需检验该映射为满射。

\medskip\noindent
设 $D = \sum n_Z Z$ 为 $X$ 上的 Weil 除子。我们断言存在开集
$U \subset X$，使得 $X \setminus U$ 的每个不可约分支之余维均为
$\geq 2$（余维在 $X$ 中取），并且 $Z|_U$ 在 $n_Z \not = 0$ 时是有效 Cartier 除子。
为证明该断言，可以假设 $X$ 仿射。于是可以假设
$D = n_1 Z_1 + \ldots + n_r Z_r$ 是有限和，其中
$Z_1, \ldots, Z_r$ 两两不同。删去 $Z_i \cap Z_j$（$i \not = j$）后，
可以假设 $Z_1, \ldots, Z_r$ 两两不交。这把问题约化到单个素除子
$Z$（在 $X$ 上）的情形。由《性质》引理 \ref{properties-lemma-criterion-normal}，
$X$ 满足 $(R_1)$，故局部环
$\mathcal{O}_{X, \xi}$ 是离散赋值环；这里 $\xi$ 是 $Z$ 的泛点。设
$f \in \mathcal{O}_{X, \xi}$ 为一个素元。设 $V \subset X$ 为 $\xi$ 的开邻域，
使得 $f$ 是某个元素 $f \in \mathcal{O}_X(V)$ 的像。缩小 $V$ 后，可假设
$Z \cap V = V(f)$ 概形论地成立，因为这在 $\xi$ 处的局部环中成立。此时取
$$
U = X \setminus (Z \setminus V) = (X \setminus Z) \cup V
$$
便得到所需开集，从而证明该断言。

\medskip\noindent
为证明 $D$ 的除子类在像中，可以写成
$D = \sum_{n_Z < 0} n_Z Z - \sum_{n_Z > 0} (-n_Z) Z$。由上述映射的可加性，
可以而且将假设 $n_Z \leq 0$ 对所有素除子 $Z$ 均成立
（若读者愿意，可以避开这一步）。设 $U \subset X$ 如上述断言所述。
若 $U$ 拟紧，则写
$D|_U = -n_1 Z_1 - \ldots - n_r Z_r$，其中 $Z_i$ 是两两不同的素除子且
$n_i > 0$；并考虑可逆 $\mathcal{O}_U$-模
$$
\mathcal{L} =
\mathcal{I}_1^{n_1} \ldots \mathcal{I}_r^{n_r} \subset \mathcal{O}_U
$$
其中 $\mathcal{I}_i$ 是 $Z_i$ 的理想层。由 $U$ 的选取及引理
\ref{divisors-lemma-sum-effective-Cartier-divisors}，这是可逆的。此外，
$\text{div}_\mathcal{L}(1) = D|_U$。由引理
\ref{divisors-lemma-reflexive-S2-extend}，有 $\mathcal{L} = \mathcal{F}|_U$，其中存在某个
秩为 $1$ 的相干自反 $\mathcal{O}_X$-模 $\mathcal{F}$；因而
$D$ 在所述映射的像中。

\medskip\noindent
若 $U$ 不拟紧，则按上面的显示公式局部定义
$\mathcal{L} \subset \mathcal{O}_U$。读者可验证该构造能够黏合，随后与前面
完全相同地完成证明。细节从略。
\end{proof}

\begin{lemma}
\label{divisors-lemma-structure-sheaf-Xs}
设 $X$ 为整的局部 Noether 正规概形。设 $\mathcal{F}$ 为秩为 1 的相干自反
$\mathcal{O}_X$-模。设 $s \in \Gamma(X, \mathcal{F})$。令
$$
U = \{x \in X \mid s : \mathcal{O}_{X, x} \to \mathcal{F}_x
\text{ 为同构}\}
$$
则 $j : U \to X$ 是 $X$ 的开子概形，并且
$$
j_*\mathcal{O}_U =
\colim (\mathcal{O}_X \xrightarrow{s} \mathcal{F}
\xrightarrow{s} \mathcal{F}^{[2]}
\xrightarrow{s} \mathcal{F}^{[3]}
\xrightarrow{s} \ldots)
$$
其中 $\mathcal{F}^{[1]} = \mathcal{F}$，并且归纳地定义
$\mathcal{F}^{[n + 1]} =
(\mathcal{F} \otimes_{\mathcal{O}_X} \mathcal{F}^{[n]})^{**}$。
\end{lemma}

\begin{proof}
由《模》引理 \ref{modules-lemma-finite-type-surjective-on-stalk} 与
\ref{modules-lemma-finite-type-to-coherent-injective-on-stalk}，集合 $U$ 是开集。
注意，由《性质》引理 \ref{properties-lemma-immersion-into-noetherian}，$j$ 拟紧。
为证明最后的陈述，只需证明：对每个拟紧开集 $W \subset X$，存在同构
$$
\colim \Gamma(W, \mathcal{F}^{[n]})
\longrightarrow
\Gamma(U \cap W, \mathcal{O}_U)
$$
它是 $\mathcal{O}_X(W)$-模同构，且与限制映射相容。我们略去相容性的验证。
将 $X$ 替换为 $W$，并用 Hom 改写上式后，可知只需构造同构
$$
\colim \Hom_{\mathcal{O}_X}(\mathcal{O}_X, \mathcal{F}^{[n]})
\longrightarrow
\Hom_{\mathcal{O}_U}(\mathcal{O}_U, \mathcal{O}_U)
$$
选取开集 $V \subset X$，使得 $X \setminus V$ 的每个不可约分支之余维均为
$\geq 2$（余维在 $X$ 中取），并且 $\mathcal{F}|_V$ 可逆，见引理
\ref{divisors-lemma-reflexive-over-normal}。于是，限制定义了秩为 $1$ 的相干自反模范畴
之间的等价（这两个范畴分别在 $X$ 与 $V$ 上），也定义了秩为 $1$ 的相干
自反模范畴之间的等价（这两个范畴分别在 $U$ 与 $V \cap U$ 上）。见引理
\ref{divisors-lemma-reflexive-S2-extend} 及 Serre 判别法
《性质》引理 \ref{properties-lemma-criterion-normal}。因此只需构造同构
$$
\colim \Gamma(V, (\mathcal{F}|_V)^{\otimes n}) \longrightarrow
\Gamma(V \cap U, \mathcal{O}_U)
$$
由于 $\mathcal{F}|_V$ 可逆，而且 $U \cap V$ 正是 $s|_V$ 生成这个可逆模的
点所成的集合，这是《性质》引理
\ref{properties-lemma-invert-s-sections} 的特殊情形（该映射也有显式公式）。
\end{proof}

\begin{lemma}
\label{divisors-lemma-Xs-codim-complement}
假设与记号同引理 \ref{divisors-lemma-structure-sheaf-Xs}。若 $s$ 非零，则
$X \setminus U$ 的每个不可约分支之余维均为 $1$（余维在 $X$ 中取）。
\end{lemma}

\begin{proof}
设 $\xi \in X$ 是一个不可约分支 $Z$ 的泛点，而该分支属于 $X \setminus U$。将 $X$
替换为 $\xi$ 的一个开邻域后，可以假设 $Z = X \setminus U$ 不可约。由于
$s : \mathcal{O}_U \to \mathcal{F}|_U$ 是同构，若 $Z$ 在 $X$ 中的余维为
$\geq 2$，则由引理 \ref{divisors-lemma-reflexive-S2-extend} 及 Serre 判别法《性质》引理
\ref{properties-lemma-criterion-normal}，
$s : \mathcal{O}_X \to \mathcal{F}$ 是同构。这将意味着
$Z = \emptyset$，矛盾。
\end{proof}

\begin{remark}
\label{divisors-remark-structure-sheaf-Xs}
设 $A$ 为 Noether 正规整环。设 $M$ 为秩为 $1$ 的有限自反 $A$-模。设
$s \in M$ 非零。设 $\mathfrak p_1, \ldots, \mathfrak p_r$ 为高度为
$1$ 的素理想（属于 $A$），它们位于 $M/As$ 的支集中。由引理
\ref{divisors-lemma-structure-sheaf-Xs}，开集 $U$ 为
$$
U = \Spec(A) \setminus
\left(V(\mathfrak p_1) \cup \ldots \cup V(\mathfrak p_r)\right)
$$
这是引理 \ref{divisors-lemma-Xs-codim-complement} 的结论。此外，若 $M^{[n]}$ 表示
$M \otimes_A \ldots \otimes_A M$（$n$ 个因子）的自反包，则由引理
\ref{divisors-lemma-structure-sheaf-Xs}，
$$
\Gamma(U, \mathcal{O}_U) = \colim M^{[n]}
$$
\end{remark}

\begin{lemma}
\label{divisors-lemma-affine-Xs}
假设与记号同引理 \ref{divisors-lemma-structure-sheaf-Xs}。下列条件等价：
\begin{enumerate}
\item 包含态射 $j : U \to X$ 是仿射的；并且
\item 对每个 $x \in X \setminus U$，存在 $n > 0$，使得
$s^n \in \mathfrak m_x \mathcal{F}^{[n]}_x$。
\end{enumerate}
\end{lemma}

\begin{proof}
假设 (1)。对 $x \in X \setminus U$，逆像 $U_x$（它是 $U$ 在典范态射
$f_x : \Spec(\mathcal{O}_{X, x}) \to X$ 下的逆像）是仿射的，且不包含
$x$。因此 $\mathfrak m_x \Gamma(U_x, \mathcal{O}_{U_x})$ 是单位理想。特别地，
可写成
$$
1 = \sum f_i g_i
$$
其中 $f_i \in \mathfrak m_x$ 且 $g_i \in \Gamma(U_x, \mathcal{O}_{U_x})$。
由引理 \ref{divisors-lemma-structure-sheaf-Xs}，有
$\Gamma(U_x, \mathcal{O}_{U_x}) = \colim \mathcal{F}^{[n]}_x$，其转移映射由乘以
$s$ 给出。因此对某个 $n > 0$，有
$$
s^n = \sum f_i t_i
$$
其中某些 $t_i = s^ng_i \in \mathcal{F}^{[n]}_x$。故 (2) 成立。

\medskip\noindent
反之，假设 (2) 成立。$j$ 是否仿射是 $X$ 上的局部问题，见《态射》引理
\ref{morphisms-lemma-characterize-affine}。因此可以而且将假设 $X$ 仿射。
目标是证明 $U$ 仿射。由《概形上同调》引理
\ref{coherent-lemma-affine-if-quasi-affine}，只需证明
$H^p(U, \mathcal{O}_U) = 0$ 对 $p > 0$ 成立。由于
$H^p(U, \mathcal{O}_U) = H^0(X, R^pj_*\mathcal{O}_U)$
（《概形上同调》引理
\ref{coherent-lemma-quasi-coherence-higher-direct-images-application}），并且
$R^pj_*\mathcal{O}_U$ 拟相干（《概形上同调》引理
\ref{coherent-lemma-quasi-coherence-higher-direct-images}），只需证明茎
$(R^pj_*\mathcal{O}_U)_x$ 在点 $x \in X$ 处为零。考虑基变换图
$$
\xymatrix{
U_x \ar[d]_{j_x} \ar[r] & U \ar[d]^j \\
\Spec(\mathcal{O}_{X, x}) \ar[r] & X
}
$$
由《概形上同调》引理
\ref{coherent-lemma-flat-base-change-cohomology}，有
$(R^pj_*\mathcal{O}_U)_x = R^pj_{x, *}\mathcal{O}_{U_x}$。因此可假设 $X$ 是以
$x$ 为闭点的局部概形，并且必须证明 $U$ 仿射（因为由上面给出的引用，这与
所需消失等价）。特别地，$d = \dim(X)$ 有限（《代数》命题
\ref{algebra-proposition-dimension}）。若 $x \in U$，则 $U = X$，结论显然。
若 $d = 0$ 且 $x \not \in U$，则 $U = \emptyset$，结论也显然。现在假设
$d > 0$ 且 $x \not \in U$。由于
$j_*\mathcal{O}_U = \colim \mathcal{F}^{[n]}$，假设意味着可写成
$$
1 = \sum f_i g_i
$$
其中对某个 $n > 0$，有 $f_i \in \mathfrak m_x$ 且
$g_i \in \mathcal{O}(U)$。由对 $d$ 的归纳可知，
$D(f_i) \cap U$ 对所有 $i$ 都仿射：把刚才的整个论证中的 $X$ 替换为 $D(f_i)$，最终会
得到维数严格小于 $d$ 的 Noether 局部环。因此由《性质》引理
\ref{properties-lemma-characterize-affine}，$U$ 仿射，正如所需。
\end{proof}


\section{相对 Proj}
\label{divisors-section-relative-proj}

\noindent
这里给出相对 Proj 的一些结果。先从若干非常基本的结果开始。回忆，相对 Proj
总是在基底上分离，见《构造》引理
\ref{constructions-lemma-relative-proj-separated}。

\begin{lemma}
\label{divisors-lemma-relative-proj-quasi-compact}
设 $S$ 为概形。设 $\mathcal{A}$ 为拟相干分次 $\mathcal{O}_S$-代数。设
$p : X = \underline{\text{Proj}}_S(\mathcal{A}) \to S$ 为
$\mathcal{A}$ 的相对 Proj。若下列条件之一成立：
\begin{enumerate}
\item $\mathcal{A}$ 作为 $\mathcal{A}_0$-代数层是有限型的；
\item $\mathcal{A}$ 由 $\mathcal{A}_1$ 生成（作为 $\mathcal{A}_0$-代数），且
$\mathcal{A}_1$ 是有限型 $\mathcal{A}_0$-模；
\item 存在有限型拟相干 $\mathcal{A}_0$-子模
$\mathcal{F} \subset \mathcal{A}_{+}$，使得
$\mathcal{A}_{+}/\mathcal{F}\mathcal{A}$ 是
$\mathcal{A}/\mathcal{F}\mathcal{A}$ 的局部幂零理想层；
\end{enumerate}
则 $p$ 拟紧。
\end{lemma}

\begin{proof}
该问题在基底上是局部的，见《概形》引理
\ref{schemes-lemma-quasi-compact-affine}。因此可假设 $S$ 仿射。记
$S = \Spec(R)$，且 $\mathcal{A}$ 对应于分次 $R$-代数 $A$。于是
$X = \text{Proj}(A)$，见《构造》节
\ref{constructions-section-relative-proj-via-glueing}。在情形 (1)，进一步局部化后
可假设 $A$ 由齐次元素 $f_1, \ldots, f_n \in A_{+}$ 生成（在 $A_0$ 上）。于是由
《代数》引理 \ref{algebra-lemma-S-plus-generated}，有
$A_{+} = (f_1, \ldots, f_n)$。在情形 (3)，可知
$\mathcal{F} = \widetilde{M}$，其中 $A_0$-模
$M \subset A_{+}$ 是有限型的。记
$M = \sum A_0f_i$。把 $f_i = \sum f_{i, j}$ 写成齐次分量之分解。条件 (3)
表示 $A_{+} \subset \sqrt{(f_{i, j})}$。因此在这两种情形，由《构造》引理
\ref{constructions-lemma-proj-quasi-compact}，可知 $\text{Proj}(A)$ 拟紧。
最后，(2) 由 (1) 得出。
\end{proof}

\begin{lemma}
\label{divisors-lemma-relative-proj-finite-type}
设 $S$ 为概形。设 $\mathcal{A}$ 为拟相干分次 $\mathcal{O}_S$-代数。设
$p : X = \underline{\text{Proj}}_S(\mathcal{A}) \to S$ 为
$\mathcal{A}$ 的相对 Proj。若 $\mathcal{A}$ 作为 $\mathcal{O}_S$-代数层是
有限型的，则 $p$ 是有限型态射，且 $\mathcal{O}_X(d)$ 是有限型
$\mathcal{O}_X$-模。
\end{lemma}

\begin{proof}
该假设蕴含 $p$ 拟紧，见引理 \ref{divisors-lemma-relative-proj-quasi-compact}。因此只需
证明 $p$ 局部有限型。故该问题在基底和目标上是局部的，见《态射》引理
\ref{morphisms-lemma-locally-finite-type-characterize}。记 $S = \Spec(R)$，且
$\mathcal{A}$ 对应于分次 $R$-代数 $A$。在 $S$ 上进一步局部化后，可假设
$A$ 是有限型 $R$-代数。概形 $X$ 由环 $A_{(f)}$ 的谱黏合而成，其中
$f \in A_{+}$ 齐次。由《代数》引理
\ref{algebra-lemma-dehomogenize-finite-type} 的 (1)，每个这样的环在 $R$ 上都是
有限型的。因此 $\text{Proj}(A)$ 在 $R$ 上有限型。关于
$\mathcal{O}_X(d)$ 的陈述使用《代数》引理
\ref{algebra-lemma-dehomogenize-finite-type} 的 (2)。
\end{proof}

\begin{lemma}
\label{divisors-lemma-relative-proj-universally-closed}
设 $S$ 为概形。设 $\mathcal{A}$ 为拟相干分次 $\mathcal{O}_S$-代数。设
$p : X = \underline{\text{Proj}}_S(\mathcal{A}) \to S$ 为
$\mathcal{A}$ 的相对 Proj。若 $\mathcal{O}_S \to \mathcal{A}_0$ 是整代数映射
\footnote{换言之，$\mathcal{O}_S$ 在 $\mathcal{A}_0$ 中的整闭包（见《态射》定义
\ref{morphisms-definition-integral-closure}）等于 $\mathcal{A}_0$。}，并且
$\mathcal{A}$ 作为 $\mathcal{A}_0$-代数有限型，则 $p$ 万有闭。
\end{lemma}

\begin{proof}
该问题在基底上是局部的。因此可假设 $X = \Spec(R)$ 仿射。设
$\mathcal{A}$ 为拟相干 $\mathcal{O}_X$-代数，它与分次 $R$-代数 $A$ 相伴。
假设是 $R \to A_0$ 为整映射，且 $A$ 在 $A_0$ 上有限型。把
$X \to \Spec(R)$ 写成复合 $X \to \Spec(A_0) \to \Spec(R)$。由于
$R \to A_0$ 是整环映射，由《态射》引理
\ref{morphisms-lemma-integral-universally-closed}，可知
$\Spec(A_0) \to \Spec(R)$ 万有闭。拟紧态射（见《构造》引理
\ref{constructions-lemma-proj-quasi-compact}）
$$
X = \text{Proj}(A) \to \Spec(A_0)
$$
由《构造》引理 \ref{constructions-lemma-proj-valuative-criterion} 满足赋值判别法的
存在性部分，因而由《概形》命题
\ref{schemes-proposition-characterize-universally-closed} 万有闭。所以
$X \to \Spec(R)$ 作为万有闭态射的复合是万有闭的。
\end{proof}

\begin{lemma}
\label{divisors-lemma-relative-proj-proper}
设 $S$ 为概形。设 $\mathcal{A}$ 为拟相干分次 $\mathcal{O}_S$-代数。设
$p : X = \underline{\text{Proj}}_S(\mathcal{A}) \to S$ 为
$\mathcal{A}$ 的相对 Proj。下列条件等价：
\begin{enumerate}
\item $\mathcal{A}_0$ 是有限型 $\mathcal{O}_S$-模，且 $\mathcal{A}$ 作为
$\mathcal{A}_0$-代数有限型；
\item $\mathcal{A}_0$ 是有限型 $\mathcal{O}_S$-模，且 $\mathcal{A}$ 作为
$\mathcal{O}_S$-代数有限型。
\end{enumerate}
若这些条件成立，则 $p$ 局部射影，特别地它是固有态射。
\end{lemma}

\begin{proof}
假设 $\mathcal{A}_0$ 是有限型 $\mathcal{O}_S$-模。选取仿射开集
$U = \Spec(R) \subset X$，使得 $\mathcal{A}$ 对应于分次 $R$-代数 $A$，且
$A_0$ 是有限 $R$-模。条件 (1) 意味着（在 $S$ 上进一步局部化后）$A$ 是
有限型 $A_0$-代数，而条件 (2) 意味着（在 $S$ 上进一步局部化后）$A$ 是
有限型 $R$-代数。因此由《代数》引理
\ref{algebra-lemma-compose-finite-type}，这两个条件相互蕴含。

\medskip\noindent
局部射影态射是固有的，见《态射》引理
\ref{morphisms-lemma-locally-projective-proper}。因此现在可假设
$S = \Spec(R)$ 且 $X = \text{Proj}(A)$，并且 $A_0$ 在 $R$ 上有限而 $A$ 在
$R$ 上有限型。我们将证明 $X = \text{Proj}(A) \to \Spec(R)$ 是射影态射。
我们建议读者自己证明这一点：直接构造 $X$ 到 $R$ 上射影空间中的闭浸入，
而不要阅读下面给出的论证。

\medskip\noindent
由引理 \ref{divisors-lemma-relative-proj-finite-type}，可知 $X$ 在 $\Spec(R)$ 上有限型。
《构造》引理 \ref{constructions-lemma-ample-on-proj} 告诉我们，
$\mathcal{O}_X(d)$ 在 $X$ 上对某个 $d \geq 1$ 丰沛（见《性质》节
\ref{properties-section-ample}）。因此 $X \to \Spec(R)$ 是拟射影态射（由
《态射》定义 \ref{morphisms-definition-quasi-projective}）。由《态射》引理
\ref{morphisms-lemma-quasi-projective-open-projective}，可知 $X$ 同构于某个在
$\Spec(R)$ 上射影的概形之开子概形。因此，为完成证明，只需证明
$X \to \Spec(R)$ 万有闭（使用《态射》引理
\ref{morphisms-lemma-image-proper-scheme-closed}）。这由引理
\ref{divisors-lemma-relative-proj-universally-closed} 得出。
\end{proof}

\begin{lemma}
\label{divisors-lemma-relative-proj-projective}
设 $S$ 为概形。设 $\mathcal{A}$ 为拟相干分次 $\mathcal{O}_S$-代数。设
$p : X = \underline{\text{Proj}}_S(\mathcal{A}) \to S$ 为
$\mathcal{A}$ 的相对 Proj。若 $\mathcal{A}$ 由 $\mathcal{A}_1$ 生成（在
$\mathcal{A}_0$ 上），且 $\mathcal{A}_1$ 是有限型 $\mathcal{O}_S$-模，则
$p$ 是射影态射。
\end{lemma}

\begin{proof}
具体而言，与分次 $\mathcal{O}_S$-代数映射
$$
\text{Sym}_{\mathcal{O}_X}^*(\mathcal{A}_1)
\longrightarrow
\mathcal{A}
$$
相伴的态射是闭浸入 $X \to \mathbf{P}(\mathcal{A}_1)$，见《构造》引理
\ref{constructions-lemma-surjective-generated-degree-1-map-relative-proj}。
\end{proof}

\begin{lemma}
\label{divisors-lemma-relative-proj-flat}
设 $S$ 为概形。设 $\mathcal{A}$ 为拟相干分次 $\mathcal{O}_S$-代数。设
$p : X = \underline{\text{Proj}}_S(\mathcal{A}) \to S$ 为
$\mathcal{A}$ 的相对 Proj。若 $\mathcal{A}_d$ 是平坦 $\mathcal{O}_S$-模
（当 $d \gg 0$ 时），则 $p$ 平坦，且 $\mathcal{O}_X(d)$ 在 $S$ 上平坦。
\end{lemma}

\begin{proof}
在仿射局部，$X$ 在 $S$ 上的平坦性约化为下述陈述：设 $R$ 为环，设 $A$ 为
分次 $R$-代数，且 $A_d$ 在 $R$ 上平坦（当 $d \gg 0$ 时）；设
$f \in A_d$，其中 $d > 0$。则 $A_{(f)}$ 在 $R$ 上平坦。由于
$A_{(f)} = \colim A_{nd}$，且转移映射由乘以 $f$ 给出，该结论由《代数》引理
\ref{algebra-lemma-colimit-flat} 得出。用同样的论证可得
$\mathcal{O}_X(d)$ 在 $S$ 上的平坦性。
\end{proof}

\begin{lemma}
\label{divisors-lemma-relative-proj-finite-presentation}
设 $S$ 为概形。设 $\mathcal{A}$ 为拟相干分次 $\mathcal{O}_S$-代数。设
$p : X = \underline{\text{Proj}}_S(\mathcal{A}) \to S$ 为
$\mathcal{A}$ 的相对 Proj。若 $\mathcal{A}$ 是有限呈示
$\mathcal{O}_S$-代数，则 $p$ 是有限呈示态射，且 $\mathcal{O}_X(d)$ 是有限呈示
$\mathcal{O}_X$-模。
\end{lemma}

\begin{proof}
在仿射局部，这约化为下述陈述：设 $R$ 为环，且设 $A$ 为有限呈示分次
$R$-代数。则 $\text{Proj}(A) \to \Spec(R)$ 是有限呈示态射，而且
$\mathcal{O}_{\text{Proj}(A)}(d)$ 是有限呈示
$\mathcal{O}_{\text{Proj}(A)}$-模。有限呈示条件蕴含可以选取呈示
$$
A = R[X_1, \ldots, X_n]/(F_1, \ldots, F_m)
$$
其中 $R[X_1, \ldots, X_n]$ 是如下分次多项式环：赋予权 $d_i$ 给 $X_i$；而
$F_1, \ldots, F_m$ 是次数为 $e_j$ 的齐次多项式。设 $R_0 \subset R$ 为由
多项式 $F_1, \ldots, F_m$ 的系数生成的子环。然后置
$A_0 = R_0[X_1, \ldots, X_n]/(F_1, \ldots, F_m)$。按构造，
$A = A_0 \otimes_{R_0} R$。因此由《构造》引理
\ref{constructions-lemma-base-change-map-proj}，只需对
$X_0 = \text{Proj}(A_0)$（在 $R_0$ 上）证明该结果。由引理
\ref{divisors-lemma-relative-proj-finite-type}，已知 $X_0$ 在 $R_0$ 上有限型，且
$\mathcal{O}_{X_0}(d)$ 是有限型拟相干 $\mathcal{O}_{X_0}$-模。由于 $R_0$
是 Noether 环（它是有限生成 $\mathbf{Z}$-代数），可知 $X_0$ 在 $R_0$ 上有限呈示
（《态射》引理 \ref{morphisms-lemma-noetherian-finite-type-finite-presentation}），
而由《概形上同调》引理 \ref{coherent-lemma-coherent-Noetherian}，
$\mathcal{O}_{X_0}(d)$ 是有限呈示模。证明完成。
\end{proof}


\section{相对 Proj 的闭子概形}
\label{divisors-section-closed-in-relative-proj}

\noindent
这里给出关于相对 Proj 的闭子概形的若干辅助引理。

\begin{lemma}
\label{divisors-lemma-closed-subscheme-proj}
设 $S$ 为概形。设 $\mathcal{A}$ 为拟相干分次
$\mathcal{O}_S$-代数。设
$p : X = \underline{\text{Proj}}_S(\mathcal{A}) \to S$ 为
$\mathcal{A}$ 的相对 Proj。设 $i : Z \to X$ 为闭子概形。记
$\mathcal{I} \subset \mathcal{A}$；它是下述典范映射的核：
$$
\mathcal{A}
\longrightarrow
\bigoplus\nolimits_{d \geq 0} p_*\left((i_*\mathcal{O}_Z)(d)\right).
$$
若 $p$ 拟紧，则有同构
$Z = \underline{\text{Proj}}_S(\mathcal{A}/\mathcal{I})$。
\end{lemma}

\begin{proof}
由《构造》引理 \ref{constructions-lemma-relative-proj-separated}，态射 $p$
是分离的。由于 $p$ 拟紧，$p_*$ 把拟相干模变为拟相干模，见《概形》引理
\ref{schemes-lemma-push-forward-quasi-coherent}。因此 $\mathcal{I}$ 是拟相干
$\mathcal{O}_S$-模。特别地，$\mathcal{B} = \mathcal{A}/\mathcal{I}$ 是拟相干
分次 $\mathcal{O}_S$-代数。由函子性得到的态射
$Z' = \underline{\text{Proj}}_S(\mathcal{B}) \to
\underline{\text{Proj}}_S(\mathcal{A})$ 处处有定义且为闭浸入，见《构造》引理
\ref{constructions-lemma-surjective-graded-rings-map-relative-proj}。
因此只需证明 $Z = Z'$ 作为 $X$ 的闭子概形相等。

\medskip\noindent
由此，该问题在基底上是局部的，故可假设
$S = \Spec(R)$ 且 $X = \text{Proj}(A)$，其中存在一个在 $R$ 上的分次代数
$A$。假设 $\mathcal{I} = \widetilde{I}$，其中
$I \subset A$ 是分次理想。由《构造》引理
\ref{constructions-lemma-proj-quasi-compact}，存在
$f_0, \ldots, f_n \in A_{+}$，使得
$A_{+} \subset \sqrt{(f_0, \ldots, f_n)}$；换言之，
$X = \bigcup D_{+}(f_i)$。因此只需验证
$Z \cap D_{+}(f_i) = Z' \cap D_{+}(f_i)$ 对每个 $i$ 都成立。重新编号后可假设
$i = 0$。设 $Z \cap D_{+}(f_0)$ 与 $Z' \cap D_{+}(f_0)$
分别由理想 $J$ 与 $J'$ 截出，这两个理想属于 $A_{(f_0)}$。

\medskip\noindent
先证包含关系 $J' \subset J$。设 $d$ 是
$\deg(f_0), \ldots, \deg(f_n)$ 的最小公倍数。注意，每个扭曲
$\mathcal{O}_X(nd)$ 都可逆，并且由 $f_i^{nd/\deg(f_i)}$ 在
$D_{+}(f_i)$ 上将其平凡化；对于任意拟相干模 $\mathcal{F}$（在 $X$ 上），乘法映射
$\mathcal{O}_X(nd) \otimes_{\mathcal{O}_X} \mathcal{F}(m)
\to \mathcal{F}(nd + m)$
都是同构，见《构造》引理
\ref{constructions-lemma-when-invertible}。注意，$J'$ 是由元素
$g/f_0^e$ 生成的理想，其中 $g \in I$ 是次数为
$e\deg(f_0)$ 的齐次元素（见《构造》引理
\ref{constructions-lemma-surjective-graded-rings-map-proj} 的证明）。
当然，把 $g$ 替换为 $f_0^lg$（这里 $l$ 取适当值），则总可假设
$d | e$。由于 $g$ 作为 $\mathcal{O}_X(e\deg(f_0))$ 的截面限制到
$Z$ 后消失，可知 $g/f_0^d$ 是 $J$ 的元素。因此 $J' \subset J$。

\medskip\noindent
反之，假设 $g/f_0^e \in J$。仍可假设 $d | e$。选取
$i \in \{1, \ldots, n\}$。此时 $Z \cap D_{+}(f_i)$ 由某个理想
$J_i \subset A_{(f_i)}$ 截出，并且
$$
J \cdot A_{(f_0f_i)} = J_i \cdot A_{(f_0f_i)}.
$$
右端是 $J_i$ 关于
$f_0^{\deg(f_i)}/f_i^{\deg(f_0)}$ 的局部化。因此，对某个充分可除的
$e_i \gg 0$，有
$$
f_0^{e_i}g/f_i^{(e_i + e)\deg(f_0)/\deg(f_i)} \in J_i
$$
这证明 $f_0^{\max(e_i)}g$ 是 $I$ 的元素，因为它在每个仿射开集
$D_{+}(f_i)$ 上的限制都在闭子概形
$Z \cap D_{+}(f_i)$ 上消失。所以 $g/f_0^e \in J'$，从而如所需有
$J \subset J'$。
\end{proof}

\begin{example}
\label{divisors-example-closed-subscheme-of-proj}
设 $A$ 为分次环。设 $X = \text{Proj}(A)$ 且 $S = \Spec(A_0)$。
给定分次理想 $I \subset A$，由《构造》引理
\ref{constructions-lemma-surjective-graded-rings-map-proj} 得到闭子概形
$V_+(I) = \text{Proj}(A/I) \to X$。把引理
\ref{divisors-lemma-closed-subscheme-proj} 的结果翻译到这一情形，可知若 $X$ 拟紧，
则任意闭子概形 $Z$ 都具有形式 $V_+(I(Z))$，其中分次理想
$I(Z) \subset A$ 由下式给出：
$$
I(Z) = \Ker(A \longrightarrow
\bigoplus\nolimits_{n \geq 0} \Gamma(Z, \mathcal{O}_Z(n)))
$$
于是可以提出以下两个自然问题：
\begin{enumerate}
\item 哪些理想 $I$ 具有形式 $I(Z)$？
\item 如何描述运算 $I \mapsto I(V_+(I))$？
\end{enumerate}
我们将在 $A$ 为 Noether 环时回答这些问题。

\medskip\noindent
首先假设 $A$ 由 $A_1$ 生成（在 $A_0$ 上）。在此情形，对任意理想
$I \subset A$，映射
$A/I \to \bigoplus \Gamma(\text{Proj}(A/I), \mathcal{O})$
的核是 $A/I$ 的挠元素集合，见《概形上同调》命题
\ref{coherent-proposition-coherent-modules-on-proj}。因此
$$
I(V_+(I)) = \{x \in A \mid A_n x \subset I\text{ 对某个 }n \geq 0\}
$$
右端的理想有时称为 $I$ 的饱和化。这回答了 (2)；对 (1) 的回答是：
一个理想具有形式 $I(Z)$，当且仅当它是饱和的，即等于自身的饱和化。

\medskip\noindent
若 $A$ 是一般的 Noether 分次环，则使用《概形上同调》命题
\ref{coherent-proposition-coherent-modules-on-proj-general}。因此，若 $d$
等于 $A$ 在 $A_0$ 上一组生成元的次数之最小公倍数，则
$$
I(V_+(I)) = \{x \in A \mid (Ax)_{nd} \subset I\text{ 对所有 }n \gg 0\}
$$
这可能不同于 $I$ 的饱和化；当 $d \not = 1$ 时会出现这种情形。例如，假设
$A = \mathbf{Q}[x, y]$，其中 $\deg(x) = 2$ 且 $\deg(y) = 3$。
于是 $d = 6$。令 $I = (y^2)$。可知 $y \in I(V_+(I))$，因为对于任意
齐次 $f \in A$，若 $6 | \deg(fy)$，则 $y | f$，故 $fy \in I$。
由此 $I(V_+(I)) = (y)$，但 $x^n y \not \in I$ 对所有 $n$
都成立；因此 $I(V_+(I))$ 不等于该饱和化。
\end{example}

\begin{lemma}
\label{divisors-lemma-equation-codim-1-in-projective-space}
设 $R$ 为 Noether UFD。设 $Z \subset \mathbf{P}^n_R$ 为闭子概形，
它不含嵌入点，并且 $Z$ 的每个不可约分支的余维均为 $1$（在
$\mathbf{P}^n_R$ 中）。则理想
$I(Z) \subset R[T_0, \ldots, T_n]$（它对应于 $Z$）为主理想。
\end{lemma}

\begin{proof}
注意，$X = \mathbf{P}^n_R$ 的局部环都是 UFD，因为 $X$ 被同构于
$\mathbf{A}^n_R$ 的仿射片覆盖，而且 $R[x_1, \ldots, x_n]$ 是 UFD
（《代数》引理 \ref{algebra-lemma-polynomial-ring-UFD}）。因此由引理
\ref{divisors-lemma-codimension-1-is-effective-Cartier}，$Z$ 是有效 Cartier 除子。
设 $\mathcal{I} \subset \mathcal{O}_X$ 为与 $Z$ 对应的拟相干理想层。
选取同构 $\mathcal{O}(m) \to \mathcal{I}$，其中
$m \in \mathbf{Z}$，见引理
\ref{divisors-lemma-Pic-projective-space-UFD}。于是复合
$$
\mathcal{O}_X(m) \to \mathcal{I} \to \mathcal{O}_X
$$
非零。可知 $m \leq 0$，且 $\mathcal{O}_X(m)^{\otimes -1} = \mathcal{O}_X(-m)$
的相应截面由某个
$F \in R[T_0, \ldots, T_n]$ 给出，其次数为 $-m$，见《概形上同调》引理
\ref{coherent-lemma-cohomology-projective-space-over-ring}。因此，在第
$i$ 个标准开集 $U_i = D_+(T_i)$ 上，闭子概形
$Z \cap U_i$ 由理想
$$
(F(T_0/T_i, \ldots, T_n/T_i)) \subset R[T_0/T_i, \ldots, T_n/T_i]
$$
截出。因此分次理想
$I(Z) = \Ker(R[T_0, \ldots, T_n] \to \bigoplus \Gamma(\mathcal{O}_Z(m)))$
的齐次元素是满足下式的齐次多项式 $G$：
$$
G(T_0/T_i, \ldots, T_n/T_i) \in (F(T_0/T_i, \ldots, T_n/T_i))
$$
其中 $i = 0, \ldots, n$。清去分母可知，存在 $e_i \geq 0$，使得
$$
T_i^{e_i}G \in (F)
$$
对 $i = 0, \ldots, n$ 成立。由于 $R$ 是 UFD，
$R[T_0, \ldots, T_n]$ 也是 UFD。因此
$F | T_0^{e_0}G$ 与 $F | T_1^{e_1}G$ 蕴含 $F | G$，因为
$T_0^{e_0}$ 与 $T_1^{e_1}$ 没有公因子。故 $I(Z) = (F)$。
\end{proof}

\noindent
当闭子概形在局部由有限多个方程截出时，可以用 $\mathcal{A}$ 的有限型理想层
定义它。

\begin{lemma}
\label{divisors-lemma-closed-subscheme-proj-finite}
设 $S$ 为拟紧拟分离概形。设 $\mathcal{A}$ 为拟相干分次
$\mathcal{O}_S$-代数。设
$p : X = \underline{\text{Proj}}_S(\mathcal{A}) \to S$ 为
$\mathcal{A}$ 的相对 Proj。设 $i : Z \to X$ 为闭子概形。
若 $p$ 拟紧且 $i$ 是有限呈示态射，则存在 $d > 0$ 和有限型拟相干
$\mathcal{O}_S$-子模 $\mathcal{F} \subset \mathcal{A}_d$，使得
$Z = \underline{\text{Proj}}_S(\mathcal{A}/\mathcal{F}\mathcal{A})$。
\end{lemma}

\begin{proof}
由引理 \ref{divisors-lemma-closed-subscheme-proj}，已知存在拟相干分次理想层
$\mathcal{I} \subset \mathcal{A}$，使得
$Z = \underline{\text{Proj}}(\mathcal{A}/\mathcal{I})$。由于 $S$ 拟紧，
可以选取有限仿射开覆盖 $S = U_1 \cup \ldots \cup U_n$。记
$U_i = \Spec(R_i)$。设 $\mathcal{A}|_{U_i}$ 对应于分次
$R_i$-代数 $A_i$，而 $\mathcal{I}|_{U_i}$ 对应于分次理想
$I_i \subset A_i$。注意，$p^{-1}(U_i) = \text{Proj}(A_i)$ 是
$R_i$ 上的概形。由于 $p$ 拟紧，可以选取有限多个齐次元素
$f_{i, j} \in A_{i, +}$，使得
$p^{-1}(U_i) = \bigcup_j D_{+}(f_{i, j})$。关于
$Z \to X$ 的条件意味着 $Z$ 的理想层（作为 $\mathcal{O}_X$ 的理想）有限型，见
《态射》引理 \ref{morphisms-lemma-closed-immersion-finite-presentation}。
因此可以找到有限多个齐次元素
$h_{i, j, k} \in I_i \cap A_{i, +}$，使得
$Z \cap D_{+}(f_{i, j})$ 的理想由元素
$h_{i, j, k}/f_{i, j}^{e_{i, j, k}}$ 生成。选取 $d > 0$，使其为所有整数
$\deg(f_{i, j})$ 与 $\deg(h_{i, j, k})$ 的公倍数。由《性质》引理
\ref{properties-lemma-quasi-coherent-colimit-finite-type}，存在有限型拟相干子模
$\mathcal{F} \subset \mathcal{I}_d$，使得所有局部截面
$$
h_{i, j, k}f_{i, j}^{(d - \deg(h_{i, j, k}))/\deg(f_{i, j})}
$$
都是 $\mathcal{F}$ 的截面。按构造，$\mathcal{F}$ 即为所求。
\end{proof}

\noindent
引理 \ref{divisors-lemma-closed-subscheme-proj-finite} 的下列版本将在引理
\ref{divisors-lemma-composition-admissible-blowups} 的证明中使用。

\begin{lemma}
\label{divisors-lemma-closed-subscheme-proj-finite-type}
设 $S$ 为拟紧拟分离概形。设 $\mathcal{A}$ 为拟相干分次
$\mathcal{O}_S$-代数。设
$p : X = \underline{\text{Proj}}_S(\mathcal{A}) \to S$ 为
$\mathcal{A}$ 的相对 Proj。设 $i : Z \to X$ 为闭子概形。
设 $U \subset S$ 为开集。假设：
\begin{enumerate}
\item $p$ 拟紧；
\item $i$ 是有限呈示态射；
\item $U \cap p(i(Z)) = \emptyset$；
\item $U$ 拟紧；
\item $\mathcal{A}_n$ 都是有限型 $\mathcal{O}_S$-模（对所有 $n$）。
\end{enumerate}
则存在 $d > 0$ 和有限型拟相干 $\mathcal{O}_S$-子模
$\mathcal{F} \subset \mathcal{A}_d$，使得 (a)
$Z = \underline{\text{Proj}}_S(\mathcal{A}/\mathcal{F}\mathcal{A})$，且 (b)
$\mathcal{A}_d/\mathcal{F}$ 的支集与 $U$ 不交。
\end{lemma}

\begin{proof}
设 $\mathcal{I} \subset \mathcal{A}$ 为引理
\ref{divisors-lemma-closed-subscheme-proj} 中构造的拟相干分次理想层。令
$U_i$、$R_i$、$A_i$、$I_i$、$f_{i, j}$、$h_{i, j, k}$ 与 $d$
如引理 \ref{divisors-lemma-closed-subscheme-proj-finite} 的证明中所构造。由于
$U \cap p(i(Z)) = \emptyset$，可知
$\mathcal{I}_d|_U = \mathcal{A}_d|_U$（由我们把 $\mathcal{I}$
构造为核的方式）。由于 $U$ 拟紧，可选取有限仿射开覆盖
$U = W_1 \cup \ldots \cup W_m$。由于 $\mathcal{A}_d$ 有限型，可以找到
有限多个截面 $g_{t, s} \in \mathcal{A}_d(W_t)$，它们生成
$\mathcal{A}_d|_{W_t} = \mathcal{I}_d|_{W_t}$（作为
$\mathcal{O}_{W_t}$-模）。为完成证明，注意由《性质》引理
\ref{properties-lemma-quasi-coherent-colimit-finite-type}，存在有限型子模
$\mathcal{F} \subset \mathcal{I}_d$，使得所有局部截面
$$
h_{i, j, k}f_{i, j}^{(d - \deg(h_{i, j, k}))/\deg(f_{i, j})}
\quad\text{且}\quad
g_{t, s}
$$
都是 $\mathcal{F}$ 的截面。按构造，$\mathcal{F}$ 即为所求。
\end{proof}

\begin{lemma}
\label{divisors-lemma-conormal-sheaf-section-projective-bundle}
设 $X$ 为概形。设 $\mathcal{E}$ 为拟相干
$\mathcal{O}_X$-模。存在双射
$$
\left\{
\begin{matrix}
\text{sections }\sigma\text{ of the } \\
\text{morphism } \mathbf{P}(\mathcal{E}) \to X
\end{matrix}
\right\}
\leftrightarrow
\left\{
\begin{matrix}
\text{surjections }\mathcal{E} \to \mathcal{L}\text{ 其中} \\
\mathcal{L}\text{ is an invertible }\mathcal{O}_X\text{-模}
\end{matrix}
\right\}
$$
在此情形，$\sigma$ 是闭浸入，并且存在典范同构
$$
\Ker(\mathcal{E} \to \mathcal{L})
\otimes_{\mathcal{O}_X} \mathcal{L}^{\otimes -1}
\longrightarrow
\mathcal{C}_{\sigma(X)/\mathbf{P}(\mathcal{E})}
$$
该双射与同构都和基变换相容。
\end{lemma}

\begin{proof}
回忆，$\pi : \mathbf{P}(\mathcal{E}) \to X$ 是 $\mathcal{E}$
上的对称代数之相对 Proj，见《构造》定义
\ref{constructions-definition-projective-bundle}。因此立即得到截面 $\sigma$
的描述；这来自《构造》引理 \ref{constructions-lemma-apply-relative} 中对
$\mathbf{P}(\mathcal{E})$ 的点函子的描述。
由于 $\pi$ 分离，任意截面都是闭浸入（《构造》引理
\ref{constructions-lemma-relative-proj-separated} 和《概形》引理
\ref{schemes-lemma-section-immersion}）。
设 $U \subset X$ 为仿射开集，并设
$k \in \mathcal{E}(U)$ 与 $s \in \mathcal{E}(U)$ 为局部截面，使得 $k$
在 $\mathcal{L}$ 中映为零，而 $s$ 映为一个生成元 $\overline{s}$
（属于 $\mathcal{L}$）。于是 $f = k/s$ 是
$\mathcal{O}_{\mathbf{P}(\mathcal{E})}$ 的截面，它定义在开邻域
$D_+(s)$ 上；这是 $s(U)$ 在 $\pi^{-1}(U)$ 中的一个开邻域。此外，由于 $k$
在 $\mathcal{L}$ 中映为零，可知 $f$ 是 $s(U)$ 在 $\pi^{-1}(U)$ 中的
理想层的截面。因此可以取像 $\overline{f}$，即 $f$ 在
$\mathcal{C}_{\sigma(X)/\mathbf{P}(\mathcal{E})}(U)$ 中的像。
我们断言：(1) 像 $\overline{f}$ 只依赖于截面 $k$
与 $\overline{s}$，而不依赖于 $s$ 的选取；(2) 以这种方式在 $U$
上得到一个同构（见下文）。一旦 (1) 与 (2) 成立，便可看出该构造与
$U' \to U$ 给出的基变换相容，其中 $U'$ 仿射；这就证明这些局部映射可以
黏合，并且与任意基变换相容。

\medskip\noindent
为证明 (1) 与 (2)，我们把具体情形写出来。即设
$U = \Spec(A)$，并设 $\mathcal{E} \to \mathcal{L}$ 对应于
$A$-模映射 $M \to N$。于是
$k \in K = \Ker(M \to N)$，并且 $s \in M$ 映为生成元
$\overline{s}$（属于 $N$）。因此 $M = K \oplus A s$。于是
$$
\text{Sym}(M) = \text{Sym}(K)[s]
$$
考虑等同 $\text{Sym}(K) \to \text{Sym}(M)_{(s)}$，它由规则
$g \mapsto g/s^n$ 给出，其中 $g \in \text{Sym}^n(K)$。这给出同构
$D_+(s) = \Spec(\text{Sym}(K))$，使得 $\sigma$ 对应于环映射
$\text{Sym}(K) \to A$，它把 $K$ 映为零。通过该同构可见，与 $K$
对应的拟相干模被等同于
$\mathcal{C}_{\sigma(U)/D_+(s)}$，这就证明了 (2)。
最后，假设 $s' = k' + s$，其中 $k' \in K$。则
$$
k/s' = (k/s) (s/s') = (k/s) (s'/s)^{-1} = (k/s) (1 + k'/s)^{-1}
$$
在 $\sigma(U)$ 的一个开邻域内成立，该邻域位于 $D_+(s)$ 中。因此可见
$s'/s$ 的限制等于 $1$；这里限制在 $\sigma(U)$ 上。并且 $k/s'$ 与 $k/s$ 在余法层中
映为同一元素，从而证明了 (1)。
\end{proof}

\section{爆破}
\label{divisors-section-blowing-up}

\noindent
爆破是代数几何中的重要工具。

\begin{definition}
\label{divisors-definition-blow-up}
设 $X$ 为概形。设 $\mathcal{I} \subset \mathcal{O}_X$ 为拟相干理想层，
并设 $Z \subset X$ 为与 $\mathcal{I}$ 对应的闭子概形，见《概形》定义
\ref{schemes-definition-immersion}。所谓{\it $X$ 沿 $Z$ 的爆破}，或
所谓{\it $X$ 在理想层 $\mathcal{I}$ 处的爆破}，是态射
$$
b :
\underline{\text{Proj}}_X
\left(\bigoplus\nolimits_{n \geq 0} \mathcal{I}^n\right)
\longrightarrow
X
$$
爆破的{\it 例外除子}是逆像 $b^{-1}(Z)$。有时把 $Z$ 称为爆破的
{\it 中心}。
\end{definition}

\noindent
我们稍后将看到，例外除子是有效 Cartier 除子。此外，爆破可刻画为
$X$ 上使 $Z$ 的逆像成为有效 Cartier 除子的“最小”概形。

\medskip\noindent
若 $b : X' \to X$ 是 $X$ 在 $Z$ 处的爆破，则常以
$\mathcal{O}_{X'}(n)$ 表示结构层的扭曲。注意，这些都是可逆
$\mathcal{O}_{X'}$-模，并且
$\mathcal{O}_{X'}(n) = \mathcal{O}_{X'}(1)^{\otimes n}$，因为 $X'$
是一个拟相干分次 $\mathcal{O}_X$-代数的相对 Proj，而该代数由次数
$1$ 的部分生成，见《构造》引理
\ref{constructions-lemma-apply-relative}。注意，$\mathcal{O}_{X'}(1)$
相对于 $b$ 甚丰沛；尽管 $b$ 未必有限型，甚至未必拟紧，这是因为 $X'$
自带到 $\mathbf{P}(\mathcal{I})$ 中的闭浸入，见《态射》例
\ref{morphisms-example-very-ample}。

\begin{lemma}
\label{divisors-lemma-blowing-up-affine}
设 $X$ 为概形。设 $\mathcal{I} \subset \mathcal{O}_X$ 为拟相干理想层。
设 $U = \Spec(A)$ 为 $X$ 的仿射开子概形，并设 $I \subset A$ 为与
$\mathcal{I}|_U$ 对应的理想。若 $b : X' \to X$ 是 $X$ 在
$\mathcal{I}$ 处的爆破，则存在典范同构
$$
b^{-1}(U) = \text{Proj}(\bigoplus\nolimits_{d \geq 0} I^d)
$$
把 $b^{-1}(U)$ 等同于 $I$ 在 $A$ 中的 Rees 代数之齐次谱。此外，
$b^{-1}(U)$ 具有由仿射爆破代数 $A[\frac{I}{a}]$ 的谱组成的仿射开覆盖。
\end{lemma}

\begin{proof}
第一条陈述由相对 Proj 的黏合构造直接得出，见《构造》节
\ref{constructions-section-relative-proj-via-glueing}。对 $a \in I$，记
$a^{(1)}$ 为 $a$ 视作次数 $1$ 的元素（在 Rees 代数
$\bigoplus_{n \geq 0} I^n$ 中）。由于这些元素在 $A$ 上生成
Rees 代数，可知 $\text{Proj}(\bigoplus_{d \geq 0} I^d)$ 被仿射开集
$D_{+}(a^{(1)})$ 覆盖。仿射概形 $D_{+}(a^{(1)})$ 是仿射爆破代数
$A' = A[\frac{I}{a}]$ 的谱，见《代数》定义
\ref{algebra-definition-blow-up}。证明完成。
\end{proof}

\begin{lemma}
\label{divisors-lemma-flat-base-change-blowing-up}
\begin{slogan}
爆破与平坦基变换可交换。
\end{slogan}
设 $X_1 \to X_2$ 为概形的平坦态射。设 $Z_2 \subset X_2$ 为闭子概形。
设 $Z_1$ 为 $Z_2$ 在 $X_1$ 中的逆像。设 $X'_i$ 为以 $Z_i$ 为中心对
$X_i$ 的爆破。则存在概形的笛卡尔图
$$
\xymatrix{
X_1' \ar[r] \ar[d] & X_2' \ar[d] \\
X_1 \ar[r] & X_2
}
$$
\end{lemma}

\begin{proof}
设 $\mathcal{I}_2$ 为 $Z_2$ 在 $X_2$ 中的理想层。记给定态射为
$g : X_1 \to X_2$。则理想层 $\mathcal{I}_1$（对应于 $Z_1$）是映射
$g^*\mathcal{I}_2 \to \mathcal{O}_{X_1}$ 的像（由逆像的定义，见
《概形》定义 \ref{schemes-definition-inverse-image-closed-subscheme}）。
由《构造》引理 \ref{constructions-lemma-relative-proj-base-change}，
$X_1 \times_{X_2} X_2'$ 是
$\bigoplus_{n \geq 0} g^*\mathcal{I}_2^n$ 的相对 Proj。由于 $g$ 平坦，
映射 $g^*\mathcal{I}_2^n \to \mathcal{O}_{X_1}$ 为单射，其像为
$\mathcal{I}_1^n$。因此 $X_1 \times_{X_2} X_2' = X_1'$。
\end{proof}

\begin{lemma}
\label{divisors-lemma-blowing-up-gives-effective-Cartier-divisor}
设 $X$ 为概形。设 $Z \subset X$ 为闭子概形。爆破
$b : X' \to X$（即以 $Z$ 为中心对 $X$ 的爆破）具有下列性质：
\begin{enumerate}
\item $b|_{b^{-1}(X \setminus Z)} : b^{-1}(X \setminus Z) \to X \setminus Z$
是同构；
\item 例外除子 $E = b^{-1}(Z)$ 是 $X'$ 上的有效 Cartier 除子；
\item 存在典范同构
$\mathcal{O}_{X'}(-1) = \mathcal{O}_{X'}(E)$。
\end{enumerate}
\end{lemma}

\begin{proof}
由于爆破与到开子概形的限制可交换（引理
\ref{divisors-lemma-flat-base-change-blowing-up}），第一条陈述只不过是说：
有 $X' = X$；这里 $Z = \emptyset$。在此情形，我们在理想层
$\mathcal{I} = \mathcal{O}_X$ 处爆破，结论由《构造》例
\ref{constructions-example-trivial-proj} 得出。

\medskip\noindent
第二条陈述在 $X$ 上是局部的，故可假设 $X$ 仿射。记
$X = \Spec(A)$ 且 $Z = \Spec(A/I)$。由引理
\ref{divisors-lemma-blowing-up-affine}，$X'$ 被仿射爆破代数
$A' = A[\frac{I}{a}]$ 的谱覆盖。此时 $IA' = aA'$，并且按《代数》引理
\ref{algebra-lemma-affine-blowup}，$a$ 在 $A'$ 中的像是非零因子。这证明了
引理，因为 $Z$ 在 $\Spec(A')$ 中的逆像对应于
$\Spec(A'/IA') \subset \Spec(A')$。

\medskip\noindent
考虑典范映射
$\psi_{univ, 1} : b^*\mathcal{I} \to \mathcal{O}_{X'}(1)$，见《构造》定义
\ref{constructions-definition-relative-proj} 后面的讨论。我们断言，它分解为同构
$\mathcal{I}_E \to \mathcal{O}_{X'}(1)$（这证明最后一条断言）。
具体而言，在上面提到的、与爆破代数
$A' = A[\frac{I}{a}]$ 对应的仿射开集上，$\psi_{univ, 1}$ 对应于
$A'$-模映射
$$
I \otimes_A A'
\longrightarrow
\left(\Big(\bigoplus\nolimits_{d \geq 0} I^d\Big)_{a^{(1)}}\right)_1
$$
其中 $a^{(1)}$ 如《代数》定义 \ref{algebra-definition-blow-up} 所述。
我们略去验证它正是映射 $I \otimes_A A' \to IA' = aA'$。
\end{proof}

\begin{lemma}[爆破的泛性质]
\label{divisors-lemma-universal-property-blowing-up}
设 $X$ 为概形。设 $Z \subset X$ 为闭子概形。设 $\mathcal{C}$ 为
$(\Sch/X)$ 的如下全子范畴：其对象是满足下列条件的 $Y \to X$，即
$Z$ 在 $Y$ 中的逆像是有效 Cartier 除子。则爆破
$b : X' \to X$（即以 $Z$ 为中心对 $X$ 作爆破）是 $\mathcal{C}$ 的终对象。
\end{lemma}

\begin{proof}
由引理 \ref{divisors-lemma-blowing-up-gives-effective-Cartier-divisor}，
$b : X' \to X$ 是 $\mathcal{C}$ 的对象。设 $f : Y \to X$ 是
$\mathcal{C}$ 的对象。必须证明存在唯一态射 $Y \to X'$，且该态射在 $X$ 上。令
$D = f^{-1}(Z)$。设 $\mathcal{I} \subset \mathcal{O}_X$ 为 $Z$ 的理想层，
并设 $\mathcal{I}_D$ 为 $D$ 的理想层。于是
$f^*\mathcal{I} \to \mathcal{I}_D$ 是到一个可逆
$\mathcal{O}_Y$-模的满射。它扩张为映射
$\psi : \bigoplus f^*\mathcal{I}^d \to \bigoplus \mathcal{I}_D^d$，
这是分次 $\mathcal{O}_Y$-代数映射。（注意，有
$\mathcal{I}_D^d = \mathcal{I}_D^{\otimes d}$，因为 $D$ 是有效 Cartier
除子。）由《构造》节
\ref{constructions-section-relative-proj} 的内容，三元组
$(1, f : Y \to X, \psi)$ 定义了态射 $Y \to X'$，它在 $X$ 上。其限制
$$
Y \setminus D \longrightarrow X' \setminus b^{-1}(Z) = X \setminus Z
$$
是唯一的。由引理 \ref{divisors-lemma-complement-effective-Cartier-divisor}，开集
$Y \setminus D$ 在 $Y$ 中概形论稠密。因此由《态射》引理
\ref{morphisms-lemma-equality-of-morphisms}，态射 $Y \to X'$ 唯一
（并且由《构造》引理
\ref{constructions-lemma-relative-proj-separated}，$b$ 也分离）。
\end{proof}

\begin{lemma}
\label{divisors-lemma-characterize-affine-blowup}
设 $b : X' \to X$ 为概形 $X$ 沿闭子概形 $Z$ 的爆破。设
$U = \Spec(A)$ 为 $X$ 的仿射开集，并设 $I \subset A$ 为与
$Z \cap U$ 对应的理想。设 $a \in I$，并设 $x' \in X'$ 为映到
$U$ 中某点的点。则 $x'$ 属于仿射开集
$U' = \Spec(A[\frac{I}{a}])$，当且仅当 $a$ 在
$\mathcal{O}_{X', x'}$ 中的像截出例外除子。
\end{lemma}

\begin{proof}
由于 $U'$ 上的例外除子由 $a$ 在
$A' = A[\frac{I}{a}]$ 中的像截出，一个方向是清楚的。反之，假设
$a$ 在 $\mathcal{O}_{X', x'}$ 中的像截出 $E$。由于 $I$ 的每个元素在
定义 $E$ 的理想中，这是在 $b^{-1}(U)$ 上而言；可知 $I$ 的元素都可被
$a$ 整除（在 $\mathcal{O}_{X', x'}$ 中）。因此对 $f \in I^n$，可写成
$f = \psi(f) a^n$，其中 $\psi(f) \in \mathcal{O}_{X', x'}$ 为某个元素。
注意，由于 $a$ 映到 $\mathcal{O}_{X', x'}$ 的非零因子，元素
$\psi(f)$ 由此唯一确定。于是定义
$$
A' \longrightarrow \mathcal{O}_{X', x'},\quad
f/a^n \longmapsto \psi(f)
$$
这里使用《代数》定义 \ref{divisors-definition-blow-up} 后给出的爆破代数描述。
上面提到的唯一性表明，这是 $A$-代数同态。它给出态射
$\Spec(\mathcal{O}_{X', x"}) \to \Spec(A') = U'$。由爆破的泛性质
（引理 \ref{divisors-lemma-universal-property-blowing-up}），这是 $X'$ 上的态射，
当然这蕴含 $x' \in U'$。
\end{proof}

\begin{lemma}
\label{divisors-lemma-blow-up-effective-Cartier-divisor}
设 $X$ 为概形。设 $Z \subset X$ 为有效 Cartier 除子。$X$ 在 $Z$
处的爆破是 $X$ 的恒等态射。
\end{lemma}

\begin{proof}
由爆破的泛性质（引理 \ref{divisors-lemma-universal-property-blowing-up}）立即得出。
\end{proof}

\begin{lemma}
\label{divisors-lemma-blow-up-reduced-scheme}
设 $X$ 为概形。设 $\mathcal{I} \subset \mathcal{O}_X$ 为拟相干理想层。
若 $X$ 约化，则爆破 $X'$（即 $X$ 在 $\mathcal{I}$ 处的爆破）约化。
\end{lemma}

\begin{proof}
结合引理 \ref{divisors-lemma-blowing-up-affine} 与《代数》引理
\ref{algebra-lemma-blowup-reduced}。
\end{proof}

\begin{lemma}
\label{divisors-lemma-blow-up-integral-scheme}
设 $X$ 为概形。设 $\mathcal{I} \subset \mathcal{O}_X$ 为非零拟相干理想层。
若 $X$ 整，则爆破 $X'$（即 $X$ 在 $\mathcal{I}$ 处的爆破）整。
\end{lemma}

\begin{proof}
结合引理 \ref{divisors-lemma-blowing-up-affine} 与《代数》引理
\ref{algebra-lemma-blowup-domain}。
\end{proof}

\begin{lemma}
\label{divisors-lemma-blow-up-and-irreducible-components}
设 $X$ 为概形。设 $Z \subset X$ 为闭子概形。设 $b : X' \to X$ 为
$X$ 沿 $Z$ 的爆破。则 $b$ 诱导双射：从 $X'$ 的不可约分支的泛点集合，
到 $X$ 的不可约分支中那些不属于 $Z$ 的泛点集合。
\end{lemma}

\begin{proof}
例外除子 $E \subset X'$ 是有效 Cartier 除子，并且
$X' \setminus E \to X \setminus Z$ 是同构，见引理
\ref{divisors-lemma-blowing-up-gives-effective-Cartier-divisor}。因此只需证明下述事实：
给定有效 Cartier 除子 $D \subset S$，其中 $S$ 是概形；$S$ 的不可约分支的泛点
没有一个属于 $D$。为此，可以用仿射开覆盖的成员替换 $S$。所以由引理
\ref{divisors-lemma-characterize-effective-Cartier-divisor}，可假设
$S = \Spec(A)$ 且 $D = V(f)$，其中 $f \in A$ 为非零因子。于是必须证明
$f$ 不属于任何极小素理想 $\mathfrak p \subset A$。否则，$f$ 将映为
$R_\mathfrak p$ 的极大理想中所含的非零因子，这与《代数》引理
\ref{algebra-lemma-minimal-prime-reduced-ring} 矛盾。
\end{proof}

\begin{lemma}
\label{divisors-lemma-blow-up-pullback-effective-Cartier}
设 $X$ 为概形。设 $b : X' \to X$ 为 $X$ 沿某闭子概形的爆破。拉回
$b^{-1}D$ 对所有有效 Cartier 除子 $D \subset X$ 都有定义，并且亚纯函数沿
$b$ 的拉回也有定义（定义
\ref{divisors-definition-pullback-effective-Cartier-divisor} 与
\ref{divisors-definition-pullback-meromorphic-sections}）。
\end{lemma}

\begin{proof}
由引理 \ref{divisors-lemma-blowing-up-affine} 与
\ref{divisors-lemma-characterize-effective-Cartier-divisor}，这约化为下述代数事实：
设 $A$ 为环，$I \subset A$ 为理想，$a \in I$，且 $x \in A$ 为非零因子。
则 $x$ 在 $A[\frac{I}{a}]$ 中的像是非零因子。事实上，假设
$x (y/a^n) = 0$ 在 $A[\frac{I}{a}]$ 中成立。则
$a^mxy = 0$ 在 $A$ 中成立，其中 $m$ 为某个整数。因此 $a^my = 0$，
因为 $x$ 是非零因子。
所以如所需，$y/a^n$ 在 $A[\frac{I}{a}]$ 中为零。
\end{proof}

\begin{lemma}
\label{divisors-lemma-blowing-up-two-ideals}
设 $X$ 为概形。设 $\mathcal{I}, \mathcal{J} \subset \mathcal{O}_X$
为拟相干理想层。设 $b : X' \to X$ 为 $X$ 在 $\mathcal{I}$ 处的爆破。
设 $b' : X'' \to X'$ 为 $X'$ 在
$b^{-1}\mathcal{J} \mathcal{O}_{X'}$ 处的爆破。则 $X'' \to X$
典范同构于 $X$ 在 $\mathcal{I}\mathcal{J}$ 处的爆破。
\end{lemma}

\begin{proof}
设 $E \subset X'$ 为 $b$ 的例外除子；由引理
\ref{divisors-lemma-blowing-up-gives-effective-Cartier-divisor}，它是有效 Cartier 除子。
由引理 \ref{divisors-lemma-blow-up-pullback-effective-Cartier}，
$(b')^{-1}E$ 是 $X''$ 上的有效 Cartier 除子。设 $E' \subset X''$ 为
$b'$ 的例外除子（它也是有效 Cartier 除子）。考虑有效 Cartier 除子
$E'' = E' + (b')^{-1}E$。按构造，$E''$ 的理想为
$(b \circ b')^{-1}\mathcal{I} (b \circ b')^{-1}\mathcal{J} \mathcal{O}_{X''}$。
因此由引理 \ref{divisors-lemma-universal-property-blowing-up}，存在从 $X''$ 到爆破
$c : Y \to X$ 的典范态射，其中后者是 $X$ 在
$\mathcal{I}\mathcal{J}$ 处的爆破。反之，由于
$\mathcal{I}\mathcal{J}$ 拉回为可逆理想，可知
$c^{-1}\mathcal{I}\mathcal{O}_Y$ 定义了有效 Cartier 除子，见引理
\ref{divisors-lemma-sum-closed-subschemes-effective-Cartier}。所以由引理
\ref{divisors-lemma-universal-property-blowing-up}，存在态射
$c' : Y \to X'$，它在 $X$ 上。于是
$(c')^{-1}b^{-1}\mathcal{J}\mathcal{O}_Y = c^{-1}\mathcal{J}\mathcal{O}_Y$，
它也定义了有效 Cartier 除子。因此存在态射
$c'' : Y \to X''$，它在 $X'$ 上。我们略去验证该态射与先前构造的态射
$X'' \to Y$ 互为逆态射。
\end{proof}

\begin{lemma}
\label{divisors-lemma-blowing-up-projective}
设 $X$ 为概形。设 $\mathcal{I} \subset \mathcal{O}_X$ 为拟相干理想层。
设 $b : X' \to X$ 为 $X$ 在理想层 $\mathcal{I}$ 处的爆破。若
$\mathcal{I}$ 有限型，则
\begin{enumerate}
\item $b : X' \to X$ 是射影态射；
\item $\mathcal{O}_{X'}(1)$ 是相对于 $b$ 丰沛的可逆层。
\end{enumerate}
\end{lemma}

\begin{proof}
分次 $\mathcal{O}_X$-代数的满射
$$
\text{Sym}_{\mathcal{O}_X}^*(\mathcal{I})
\longrightarrow
\bigoplus\nolimits_{d \geq 0} \mathcal{I}^d
$$
由《构造》引理
\ref{constructions-lemma-surjective-generated-degree-1-map-relative-proj}
定义闭浸入
$$
X' = \underline{\text{Proj}}_X (\bigoplus\nolimits_{d \geq 0} \mathcal{I}^d)
\longrightarrow
\mathbf{P}(\mathcal{I}).
$$
因此由《态射》定义 \ref{morphisms-definition-projective}，$b$ 是射影态射。
第二条陈述例如由《态射》引理
\ref{morphisms-lemma-characterize-relatively-ample} 中对相对丰沛可逆层的刻画
得出。略去若干细节。
\end{proof}

\begin{lemma}
\label{divisors-lemma-composition-finite-type-blowups}
\begin{slogan}
爆破的复合仍为爆破。
\end{slogan}
设 $X$ 为拟紧拟分离概形。设 $Z \subset X$ 为有限呈示闭子概形。
设 $b : X' \to X$ 为以 $Z$ 为中心的爆破。设 $Z' \subset X'$ 为有限呈示
闭子概形。设 $X'' \to X'$ 为以 $Z'$ 为中心的爆破。则存在有限呈示闭子概形
$Y \subset X$，使得
\begin{enumerate}
\item 在集合论意义下，$Y = Z \cup b(Z')$；
\item 复合 $X'' \to X$ 同构于 $X$ 在 $Y$ 处的爆破。
\end{enumerate}
\end{lemma}

\begin{proof}
条件 $Z \to X$ 有限呈示意味着 $Z$ 由有限型拟相干理想层
$\mathcal{I} \subset \mathcal{O}_X$ 截出，见《态射》引理
\ref{morphisms-lemma-closed-immersion-finite-presentation}。写成
$\mathcal{A} = \bigoplus_{n \geq 0} \mathcal{I}^n$，于是
$X' = \underline{\text{Proj}}(\mathcal{A})$。注意，由《性质》引理
\ref{properties-lemma-quasi-coherent-finite-type-ideals}，
$X \setminus Z$ 是 $X$ 的拟紧开集。由于
$b^{-1}(X \setminus Z) \to X \setminus Z$ 是同构（引理
\ref{divisors-lemma-blowing-up-gives-effective-Cartier-divisor}），同一结果表明
$b^{-1}(X \setminus Z) \setminus Z'$ 是 $X'$ 的拟紧开集。因此
$U = X \setminus (Z \cup b(Z'))$ 是 $X$ 的拟紧开集。由引理
\ref{divisors-lemma-closed-subscheme-proj-finite-type}，存在 $d > 0$ 和有限型
$\mathcal{O}_X$-子模 $\mathcal{F} \subset \mathcal{I}^d$，使得
$Z' = \underline{\text{Proj}}(\mathcal{A}/\mathcal{F}\mathcal{A})$，
并且 $\mathcal{I}^d/\mathcal{F}$ 的支集包含在 $X \setminus U$ 中。

\medskip\noindent
由于 $\mathcal{F} \subset \mathcal{I}^d$ 是 $\mathcal{O}_X$-子模，
可以把 $\mathcal{F} \subset \mathcal{I}^d \subset \mathcal{O}_X$
视为 $X$ 上的有限型拟相干理想层。为避免混淆，将其记作
$\mathcal{J} \subset \mathcal{O}_X$。由于
$\mathcal{I}^d / \mathcal{J}$ 与 $\mathcal{O}/\mathcal{I}^d$
的支集都在 $X \setminus U$ 上，可知 $V(\mathcal{J})$ 包含在
$X \setminus U$ 中。反之，由于 $\mathcal{J} \subset \mathcal{I}^d$，
可知 $Z \subset V(\mathcal{J})$。在
$X \setminus Z \cong X' \setminus b^{-1}(Z)$ 上，理想层
$\mathcal{J}$ 截出 $Z'$（见下面的显示公式）。所以
$V(\mathcal{J})$ 等于 $Z \cup b(Z')$。从而在集合论意义下，也有
$V(\mathcal{I}\mathcal{J}) = Z \cup b(Z')$。此外，
$\mathcal{I}\mathcal{J}$ 作为两个有限型理想之积仍有限型。我们断言
$X'' \to X$ 同构于 $X$ 在 $\mathcal{I}\mathcal{J}$ 处的爆破；令
$Y = V(\mathcal{I}\mathcal{J})$ 即完成引理的证明。

\medskip\noindent
首先回忆，$X$ 在 $\mathcal{I}\mathcal{J}$ 处的爆破与 $X'$ 在
$b^{-1}\mathcal{J} \mathcal{O}_{X'}$ 处的爆破相同，见引理
\ref{divisors-lemma-blowing-up-two-ideals}。因此只需证明 $X'$ 在
$b^{-1}\mathcal{J} \mathcal{O}_{X'}$ 处的爆破与 $X'$ 在 $Z'$ 处的爆破
一致。我们将证明
$$
b^{-1}\mathcal{J} \mathcal{O}_{X'} = \mathcal{I}_E^d \mathcal{I}_{Z'}
$$
作为 $X''$ 上的理想层成立。这将证明所需结论，因为
$\mathcal{I}_E^d$ 截出有效 Cartier 除子 $dE$，并且可以使用引理
\ref{divisors-lemma-blow-up-effective-Cartier-divisor} 与
\ref{divisors-lemma-blowing-up-two-ideals}。

\medskip\noindent
为验证所显示的理想等式，可以作局部计算。采用引理
\ref{divisors-lemma-blowing-up-affine} 中的记号 $A$、$I$、$a \in I$，可知
$\mathcal{F}$ 对应于 $R$-子模 $M \subset I^d$，后者同构地映到理想
$J \subset R$。条件
$Z' = \underline{\text{Proj}}(\mathcal{A}/\mathcal{F}\mathcal{A})$
意味着 $Z' \cap \Spec(A[\frac{I}{a}])$ 由元素
$m/a^d$（其中 $m \in M$）生成的理想截出。设元素 $m \in M$ 对应于函数
$f \in J$。则在仿射爆破代数 $A' = A[\frac{I}{a}]$ 中，有
$f = (a^dm)/a^d = a^d (m/a^d)$。因此等式成立。
\end{proof}

\section{严格变换}
\label{divisors-section-strict-transform}

\noindent
本节简要讨论爆破下的严格变换。设 $S$ 为概形，设 $Z \subset S$ 为闭子概形。
设 $b : S' \to S$ 为 $S$ 在 $Z$ 处的爆破，并记 $E \subset S'$ 为例外除子
$E = b^{-1}Z$。下面常常考虑概形 $X$（它在 $S$ 上），并作笛卡尔图
$$
\xymatrix{
\text{pr}_{S'}^{-1}E \ar[r] \ar[d] &
X \times_S S' \ar[r]_-{\text{pr}_X} \ar[d]_{\text{pr}_{S'}} &
X \ar[d]^f \\
E \ar[r] & S' \ar[r] & S
}
$$
由于 $E$ 是有效 Cartier 除子（引理
\ref{divisors-lemma-blowing-up-gives-effective-Cartier-divisor}），可知
$\text{pr}_{S'}^{-1}E \subset X \times_S S'$ 局部为主理想所定义
（引理 \ref{divisors-lemma-pullback-locally-principal}）。因此
$\text{pr}_{S'}^{-1}E$ 在 $X \times_S S'$ 中的补集是逆紧的
（引理 \ref{divisors-lemma-complement-locally-principal-closed-subscheme}）。
所以，对拟相干 $\mathcal{O}_{X \times_S S'}$-模 $\mathcal{G}$，
由支集在 $\text{pr}_{S'}^{-1}E$ 上的截面组成的子层是拟相干子模，见《性质》引理
\ref{properties-lemma-sections-supported-on-closed-subset}。若 $\mathcal{G}$
是拟相干代数层，例如 $\mathcal{G} = \mathcal{O}_{X \times_S S'}$，
则该子层是 $\mathcal{G}$ 的理想。

\begin{definition}
\label{divisors-definition-strict-transform}
沿用上面的 $Z \subset S$ 与 $f : X \to S$。
\begin{enumerate}
\item 给定拟相干 $\mathcal{O}_X$-模 $\mathcal{F}$，$\mathcal{F}$
关于 $S$ 在 $Z$ 处之爆破的{\it 严格变换} $\mathcal{F}'$，是
$\text{pr}_X^*\mathcal{F}$ 对支集在 $\text{pr}_{S'}^{-1}E$ 上的截面子模
取商。
\item $X$ 的{\it 严格变换}是闭子概形
$X' \subset X \times_S S'$，它由 $\mathcal{O}_{X \times_S S'}$
中支集在 $\text{pr}_{S'}^{-1}E$ 上的截面组成的拟相干理想截出。
\end{enumerate}
\end{definition}

\noindent
注意，沿一次爆破取严格变换依赖于用来作爆破的闭子概形
（而不只依赖态射 $S' \to S$）。这个概念常用于 $S$ 的闭子概形。
事实表明，$X$ 的严格变换是 $X$ 的一次爆破。

\begin{lemma}
\label{divisors-lemma-strict-transform}
在定义 \ref{divisors-definition-strict-transform} 的情形中：
\begin{enumerate}
\item 严格变换 $X'$（属于 $X$）是 $X$ 的爆破，其中心是闭子概形
$f^{-1}Z$（位于 $X$ 中）；
\item 对拟相干 $\mathcal{O}_X$-模 $\mathcal{F}$，严格变换
$\mathcal{F}'$ 典范同构于沿 $X' \to X \times_S S'$ 的下述正像：
$\mathcal{F}$ 关于爆破 $X' \to X$ 的严格变换。
\end{enumerate}
\end{lemma}

\begin{proof}
设 $X'' \to X$ 为 $X$ 在 $f^{-1}Z$ 处的爆破。由爆破的泛性质
（引理 \ref{divisors-lemma-universal-property-blowing-up}），存在交换图
$$
\xymatrix{
X'' \ar[r] \ar[d] & X \ar[d] \\
S' \ar[r] & S
}
$$
从而有态射 $X'' \to X \times_S S'$。所以第一条断言是说，该态射是以
$X'$ 为像的闭浸入。问题在 $X$ 上是局部的，故可假设 $X$ 与 $S$ 都仿射。
记 $S = \Spec(A)$、$X = \Spec(B)$，并设 $Z$ 由理想
$I \subset A$ 截出。置 $J = IB$。映射
$B \otimes_A \bigoplus_{n \geq 0} I^n \to \bigoplus_{n \geq 0} J^n$
定义了闭浸入 $X'' \to X \times_S S'$，见《构造》引理
\ref{constructions-lemma-base-change-map-proj} 与
\ref{constructions-lemma-surjective-graded-rings-generated-degree-1-map-proj}。
我们略去验证该态射与上面由泛性质构造的态射相同。选取 $a \in I$，它对应于
仿射开集 $\Spec(A[\frac{I}{a}]) \subset S'$，见引理
\ref{divisors-lemma-blowing-up-affine}。$\Spec(A[\frac{I}{a}])$ 在严格变换
$X'$ 中的逆像（这里严格变换属于 $X$）是下述环的谱：
$$
B' = (B \otimes_A A[\textstyle{\frac{I}{a}}])/a\text{-power-torsion}
$$
见《性质》引理
\ref{properties-lemma-sections-supported-on-closed-subset}。另一方面，令
$b \in J$ 为 $a$ 的像，则 $\Spec(B[\frac{J}{b}])$ 是
$\Spec(A[\frac{I}{a}])$ 在 $X''$ 中的逆像。由《代数》引理
\ref{algebra-lemma-blowup-base-change}，开集
$\Spec(B[\frac{J}{b}])$ 同构地映到开子概形
$\text{pr}_{S'}^{-1}(\Spec(A[\frac{I}{a}]))$；后者位于 $X'$ 中。因此
$X'' \to X'$ 是同构。

\medskip\noindent
在上述记号下，设 $\mathcal{F}$ 对应于 $B$-模 $N$。$\mathcal{F}$
的严格变换对应于 $B \otimes_A A[\frac{I}{a}]$-模
$$
N' = (N \otimes_A A[\textstyle{\frac{I}{a}}])/a\text{-power-torsion}
$$
见《性质》引理
\ref{properties-lemma-sections-supported-on-closed-subset}。
$\mathcal{F}$ 关于 $X$ 在 $f^{-1}Z$ 处之爆破的严格变换，对应于
$B[\frac{J}{b}]$-模
$N \otimes_B B[\frac{J}{b}]/b\text{-power-torsion}$。完全仿照上面即可证明
这两个模同构。略去细节。
\end{proof}

\begin{lemma}
\label{divisors-lemma-strict-transform-flat}
在定义 \ref{divisors-definition-strict-transform} 的情形中：
\begin{enumerate}
\item 如果 $X$ 平坦于 $S$，至少在映到 $Z$ 的每一点处如此，则
$X$ 的严格变换等于基变换 $X \times_S S'$。
\item 设 $\mathcal{F}$ 为拟相干 $\mathcal{O}_X$-模。
如果 $\mathcal{F}$ 平坦于 $S$，至少在映到 $Z$ 的每一点处如此，则
严格变换 $\mathcal{F}'$（属于 $\mathcal{F}$）等于拉回
$\text{pr}_X^*\mathcal{F}$。
\end{enumerate}
\end{lemma}

\begin{proof}
我们证明第 (2) 条；由 $S$ 上概形的严格变换之定义，它蕴含第 (1) 条。
问题在 $X$ 上是局部的。因此可假设 $S = \Spec(A)$、$X = \Spec(B)$，并且
$\mathcal{F}$ 对应于 $B$-模 $N$。于是，$\mathcal{F}'$ 在开集
$\Spec(B \otimes_A A[\frac{I}{a}])$ 上（该开集位于 $X \times_S S'$ 中）对应于模
$$
N' = (N \otimes_A A[\textstyle{\frac{I}{a}}])/a\text{-power-torsion}
$$
见《性质》引理
\ref{properties-lemma-sections-supported-on-closed-subset}。
所以只须证明 $a$-幂挠在 $N \otimes_A A[\frac{I}{a}]$ 中为零。
设 $y \in N \otimes_A A[\frac{I}{a}]$ 且 $a^n y = 0$。如果
$\mathfrak q \subset B$ 是素理想且 $a \not \in \mathfrak q$，那么 $y$ 在
$(N \otimes_A A[\frac{I}{a}])_\mathfrak q$ 中的像为零。另一方面，如果
$a \in \mathfrak q$，则 $N_\mathfrak q$ 是平坦 $A$-模，因而
$N_\mathfrak q \otimes_A A[\frac{I}{a}]
=(N \otimes_A A[\frac{I}{a}])_\mathfrak q$
没有 $a$-幂挠（因为 $A[\frac{I}{a}]$ 没有）。故 $y$ 在这个局部化中的像
也为零。由《代数》引理 \ref{algebra-lemma-characterize-zero-local}，可知
$y$ 为零。
\end{proof}

\begin{lemma}
\label{divisors-lemma-strict-transform-affine}
设 $S$ 为概形。设 $Z \subset S$ 为闭子概形。设 $b : S' \to S$ 为
沿 $Z$ 对 $S$ 作的爆破。设 $g : X \to Y$ 为 $S$ 上概形之间的仿射态射。
设 $\mathcal{F}$ 为 $X$ 上的拟相干层。设
$g' : X \times_S S' \to Y \times_S S'$ 为 $g$ 的基变换。
设 $\mathcal{F}'$ 为 $\mathcal{F}$ 关于 $b$ 的严格变换。那么
$g'_*\mathcal{F}'$ 是 $g_*\mathcal{F}$ 的严格变换。
\end{lemma}

\begin{proof}
由《概形上同调》引理 \ref{coherent-lemma-affine-base-change}，有
$g'_*\text{pr}_X^*\mathcal{F} = \text{pr}_Y^*g_*\mathcal{F}$。
设 $\mathcal{K} \subset \text{pr}_X^*\mathcal{F}$ 为由支集在
$Z$ 的逆像上的截面组成的子层；该逆像位于 $X \times_S S'$ 中。由《性质》引理
\ref{properties-lemma-push-sections-supported-on-closed-subset}，正像
$g'_*\mathcal{K}$ 是 $\text{pr}_Y^*g_*\mathcal{F}$ 的截面子层，
其支集在 $Z$ 的逆像上；该逆像位于 $Y \times_S S'$ 中。由于 $g'$ 是仿射的
（《态射》引理 \ref{morphisms-lemma-base-change-affine}），函子 $g'_*$ 正合，
结论随即成立。
\end{proof}

\begin{lemma}
\label{divisors-lemma-strict-transform-different-centers}
设 $S$ 为概形。设 $Z \subset S$ 为闭子概形。设 $D \subset S$ 为
有效 Cartier 除子。设 $Z' \subset S$ 为由 $Z$ 与 $D$ 的理想层之积
截出的闭子概形。设 $S' \to S$ 为 $S$ 沿 $Z$ 的爆破。
\begin{enumerate}
\item $S$ 沿 $Z'$ 的爆破同构于 $S' \to S$。
\item 设 $f : X \to S$ 为概形态射，设 $\mathcal{F}$ 为拟相干
$\mathcal{O}_X$-模。如果 $\mathcal{F}$ 没有支集在 $f^{-1}D$ 上的非零
局部截面，那么 $\mathcal{F}$ 关于沿 $Z$ 爆破的严格变换，与
$\mathcal{F}$ 关于 $S$ 沿 $Z'$ 爆破的严格变换相同。
\end{enumerate}
\end{lemma}

\begin{proof}
第一条由引理 \ref{divisors-lemma-blowing-up-two-ideals} 和
\ref{divisors-lemma-blow-up-effective-Cartier-divisor} 合并即得。
利用引理 \ref{divisors-lemma-blowing-up-affine}，第二条转化为下面的代数问题。
设 $A$ 为环，$I \subset A$ 为理想，$x \in A$ 为非零因子，且
$a \in I$。设 $M$ 为 $A$-模，其 $x$-挠为零。要证：
$a$-幂挠在 $M \otimes_A A[\frac{I}{a}]$ 中等于 $xa$-幂挠。
理由如下：映射 $A \to A[\frac{I}{a}]$ 的核与余核都是 $a$-幂挠，
所以将 $a$ 可逆化之后，该映射成为同构。因此
$M \to M \otimes_A A[\frac{I}{a}]$ 的核与余核也都是 $a$-幂挠。
这就蕴含所求结论。
\end{proof}

\begin{lemma}
\label{divisors-lemma-strict-transform-composition-blowups}
设 $S$ 为概形。设 $Z \subset S$ 为闭子概形。设 $b : S' \to S$ 为
中心是 $Z$ 的爆破。设 $Z' \subset S'$ 为闭子概形。设 $S'' \to S'$ 为
中心是 $Z'$ 的爆破。设 $Y \subset S$ 为闭子概形，使得集合论上
$Y = Z \cup b(Z')$，并且复合 $S'' \to S$ 同构于 $S$ 沿 $Y$ 的爆破。
在这种情形中，给定任意概形 $X$（定义在 $S$ 上）以及
$\mathcal{F} \in \QCoh(\mathcal{O}_X)$，有：
\begin{enumerate}
\item $\mathcal{F}$ 关于 $S$ 沿 $Y$ 爆破的严格变换，等于下述严格变换：
对于 $S'' \to S'$ 沿 $Z'$ 的爆破，取 $\mathcal{F}$ 关于
$S' \to S$（即 $S$ 沿 $Z$ 的爆破）的严格变换之严格变换；并且
\item $X$ 关于 $S$ 沿 $Y$ 爆破的严格变换，等于下述严格变换：
对于 $S'' \to S'$ 沿 $Z'$ 的爆破，取 $X$ 关于
$S' \to S$（即 $S$ 沿 $Z$ 的爆破）的严格变换之严格变换。
\end{enumerate}
\end{lemma}

\begin{proof}
设 $\mathcal{F}'$ 为 $\mathcal{F}$ 关于 $S' \to S$（即 $S$ 沿 $Z$ 的爆破）
的严格变换。设 $\mathcal{F}''$ 为 $\mathcal{F}'$ 关于
$S'' \to S'$（即 $S'$ 沿 $Z'$ 的爆破）的严格变换。
设 $\mathcal{G}$ 为 $\mathcal{F}$ 关于 $S'' \to S$（即 $S$ 沿 $Y$ 的爆破）
的严格变换。我们还如下标记各态射：
$$
\xymatrix{
X \times_S S'' \ar[r]_q \ar[d]^{f''} &
X \times_S S' \ar[r]_p \ar[d]^{f'} &
X \ar[d]^f \\
S'' \ar[r] & S' \ar[r] & S
}
$$
由定义有满射 $p^*\mathcal{F} \to \mathcal{F}'$ 和满射
$q^*\mathcal{F}' \to \mathcal{F}''$；利用 $q^*$ 的右正合性，可将二者合成满射
$(p \circ q)^*\mathcal{F} \to \mathcal{F}''$。另有满射
$(p \circ q)^*\mathcal{F} \to \mathcal{G}$。所以只须证明这两个满射的核相同。

\medskip\noindent
满射 $p^*\mathcal{F} \to \mathcal{F}'$ 的核支集在
$(f \circ p)^{-1}Z$ 上，故这个映射在其补集的各点处是同构。因此
$q^*p^*\mathcal{F} \to q^*\mathcal{F}'$ 的核支集在
$(f \circ p \circ q)^{-1}Z$ 上。映射
$q^*\mathcal{F}' \to \mathcal{F}''$ 的核支集在 $(f' \circ q)^{-1}Z'$ 上。
合在一起可知，$(p \circ q)^*\mathcal{F} \to \mathcal{F}''$ 的核支集在
$(f \circ p \circ q)^{-1}Z \cup (f' \circ q)^{-1}Z' =
(f \circ p \circ q)^{-1}Y$ 上。由 $\mathcal{G}$ 的构造可得分解
$(p \circ q)^*\mathcal{F} \to \mathcal{F}'' \to \mathcal{G}$。
为完成证明，只须说明 $\mathcal{F}''$ 没有支集在
$(f \circ p \circ q)^{-1}(Y) =
(f \circ p \circ q)^{-1}Z \cup (f' \circ q)^{-1}Z'$ 上的非零（局部）截面。
把引理 \ref{divisors-lemma-strict-transform-different-centers} 应用于
$\mathcal{F}'$（它在 $X \times_S S'$ 上，而后者定义在 $S'$ 上）、
闭子概形 $Z'$ 以及有效 Cartier 除子 $b^{-1}Z$，即得此结论。
\end{proof}

\begin{lemma}
\label{divisors-lemma-strict-transform-universally-injective}
在定义 \ref{divisors-definition-strict-transform} 的情形中。假设
$$
0 \to \mathcal{F}_1 \to \mathcal{F}_2 \to \mathcal{F}_3 \to 0
$$
是 $X$ 上拟相干层的正合列，并且在任意基变换 $T \to S$ 后仍正合。
那么，严格变换 $\mathcal{F}_i'$（关于任意爆破 $S' \to S$）也组成短正合列
$0 \to \mathcal{F}'_1 \to \mathcal{F}'_2 \to \mathcal{F}'_3 \to 0$。
\end{lemma}

\begin{proof}
可在 $S$ 与 $X$ 上局部化，并假设二者都仿射。然后可把
$\mathcal{F}_i$ 正像到 $S$，见引理 \ref{divisors-lemma-strict-transform-affine}。
可以假设我们的爆破是态射 $1 : S \to S$，它与有效 Cartier 除子
$D \subset S$ 相伴。此时，代数表述如下：设 $A$ 为环，并设
$0 \to M_1 \to M_2 \to M_3 \to 0$ 是 $A$-模的泛正合列。
令 $a\in A$。则序列
$$
0 \to
M_1/a\text{-幂挠子} \to
M_2/a\text{-幂挠子} \to
M_3/a\text{-幂挠子} \to 0
$$
也正合。事实上，最后一个映射的满射性与第一个映射的单射性是立即的。
问题在于中间位置的正合性。设 $x \in M_2$ 在
$M_3/a\text{-幂挠子}$ 中的像为零。则
$y = a^n x \in M_1$，且这对某个 $n$ 成立。于是 $y$ 在 $M_2/a^nM_2$ 中的像为零。
由于 $M_1 \to M_2$ 是泛单射，可知 $y$ 在 $M_1/a^nM_1$ 中的像为零。
因而有 $y = a^n z$，其中 $z \in M_1$ 是某个元素。所以
$a^n(x - y) = 0$。故 $y$ 映到 $x$ 的类；该类位于
$M_2/a\text{-幂挠子}$ 中，
这正是所需结论。
\end{proof}

\section{容许爆破}
\label{divisors-section-admissible-blowups}

\noindent
为了对所用的爆破有更多控制，我们引入下面的标准术语。

\begin{definition}
\label{divisors-definition-admissible-blowup}
设 $X$ 为概形。设 $U \subset X$ 为开子概形。称态射
$X' \to X$ 为一个{\it $U$-容许爆破}，如果存在有限呈示的闭浸入
$Z \to X$，其中 $Z$ 与 $U$ 不交，并且 $X'$ 同构于 $X$ 沿 $Z$ 的爆破。
\end{definition}

\noindent
回忆，$Z \to X$ 有限呈示，当且仅当理想层
$\mathcal{I}_Z \subset \mathcal{O}_X$ 为有限型，见《态射》引理
\ref{morphisms-lemma-closed-immersion-finite-presentation}。
特别地，$U$-容许爆破是射影态射，见引理
\ref{divisors-lemma-blowing-up-projective}。
注意，产生同一态射的中心可以不止一个。因此，要求仅仅是存在某个与
$U$ 不交的中心，它产生 $X'$。最后，由于态射 $b : X' \to X$ 在 $U$ 上
是同构（见引理 \ref{divisors-lemma-blowing-up-gives-effective-Cartier-divisor}），
我们常常滥用记号，也把 $U$ 看作 $X'$ 的开子概形。

\begin{lemma}
\label{divisors-lemma-composition-admissible-blowups}
\begin{slogan}
容许爆破对复合稳定。
\end{slogan}
设 $X$ 为拟紧且拟分离概形。设 $U \subset X$ 为拟紧开子概形。
设 $b : X' \to X$ 为 $U$-容许爆破。设 $X'' \to X'$ 为
$U$-容许爆破。那么复合 $X'' \to X$ 是 $U$-容许爆破。
\end{lemma}

\begin{proof}
这是更精确的引理 \ref{divisors-lemma-composition-finite-type-blowups} 的直接推论。
\end{proof}

\begin{lemma}
\label{divisors-lemma-extend-admissible-blowups}
设 $X$ 为拟紧且拟分离概形。设 $U, V \subset X$ 为拟紧开子概形。
设 $b : V' \to V$ 为 $U \cap V$-容许爆破。则存在 $U$-容许爆破
$X' \to X$，它在 $V$ 上的限制为 $V'$。
\end{lemma}

\begin{proof}
设 $\mathcal{I} \subset \mathcal{O}_V$ 为有限型拟相干理想层，满足
$V(\mathcal{I})$ 与 $U \cap V$ 不交，并使 $V'$ 同构于 $V$ 沿
$\mathcal{I}$ 的爆破。设
$\mathcal{I}' \subset \mathcal{O}_{U \cup V}$ 为拟相干理想层，它在
$U$ 上的限制是 $\mathcal{O}_U$，在 $V$ 上的限制是 $\mathcal{I}$
（见《层》\ref{sheaves-section-glueing-sheaves} 节）。由《性质》引理
\ref{properties-lemma-extend}，存在有限型拟相干理想层
$\mathcal{J} \subset \mathcal{O}_X$，它在 $U \cup V$ 上的限制是
$\mathcal{I}'$。引理得证。
\end{proof}

\begin{lemma}
\label{divisors-lemma-dominate-admissible-blowups}
设 $X$ 为拟紧且拟分离概形。设 $U \subset X$ 为拟紧开子概形。
设 $b_i : X_i \to X$（$i = 1, \ldots, n$）为 $U$-容许爆破。
存在 $U$-容许爆破 $b : X' \to X$，使得
(a) $b$ 分解为 $X' \to X_i \to X$，其中 $i = 1, \ldots, n$；且
(b) 每个态射 $X' \to X_i$ 都是 $U$-容许爆破。
\end{lemma}

\begin{proof}
设 $\mathcal{I}_i \subset \mathcal{O}_X$ 为有限型拟相干理想层，满足
$V(\mathcal{I}_i)$ 与 $U$ 不交，并使 $X_i$ 同构于 $X$ 沿
$\mathcal{I}_i$ 的爆破。置
$\mathcal{I} = \mathcal{I}_1 \cdot \ldots \cdot \mathcal{I}_n$，并令
$X'$ 为 $X$ 沿 $\mathcal{I}$ 的爆破。由引理
\ref{divisors-lemma-blowing-up-two-ideals}，$X' \to X$ 通过 $b_i$ 分解。
\end{proof}

\begin{lemma}
\label{divisors-lemma-separate-disjoint-opens-by-blowing-up}
\begin{slogan}
通过爆破分离不可约分支。
\end{slogan}
设 $X$ 为拟紧且拟分离概形。设 $U, V$ 为 $X$ 的两个拟紧且不交的
开子概形。则存在一个 $U \cup V$-容许爆破 $b : X' \to X$，使得
$X'$ 是开子概形的不交并 $X' = X'_1 \amalg X'_2$，并且
$b^{-1}(U) \subset X'_1$ 及 $b^{-1}(V) \subset X'_2$。
\end{lemma}

\begin{proof}
选取有限型拟相干理想层 $\mathcal{I}$ 以及相应的 $\mathcal{J}$，使得
$X \setminus U = V(\mathcal{I})$ 以及相应的
$X \setminus V = V(\mathcal{J})$，见《性质》引理
\ref{properties-lemma-quasi-coherent-finite-type-ideals}。
于是集合论上 $V(\mathcal{I}\mathcal{J}) = X$，因而
$\mathcal{I}\mathcal{J}$ 是局部幂零理想层。由于 $\mathcal{I}$ 与
$\mathcal{J}$ 都是有限型的，并且 $X$ 拟紧，存在 $n > 0$ 使得
$\mathcal{I}^n \mathcal{J}^n = 0$。我们可以而且确实把 $\mathcal{I}$
换成 $\mathcal{I}^n$，把 $\mathcal{J}$ 换成 $\mathcal{J}^n$。于是
$\mathcal{I} \mathcal{J} = 0$。设 $b : X' \to X$ 为沿
$\mathcal{I} + \mathcal{J}$ 的爆破。它是 $U \cup V$-容许的，因为
$V(\mathcal{I} + \mathcal{J}) = X \setminus U \cup V$。我们将证明
$X'$ 是开子概形的不交并 $X' = X'_1 \amalg X'_2$，并且
$b^{-1}\mathcal{I}|_{X'_2} = 0$ 与 $b^{-1}\mathcal{J}|_{X'_1} = 0$；
这就证明了引理。

\medskip\noindent
我们将使用引理 \ref{divisors-lemma-blowing-up-affine} 中对爆破的描述。
假设 $U = \Spec(A) \subset X$ 是仿射开集，使得
$\mathcal{I}|_U$ 以及相应的 $\mathcal{J}|_U$，分别对应于有限生成理想
$I \subset A$ 以及相应的 $J \subset A$。于是
$$
b^{-1}(U) = \text{Proj}(A \oplus (I + J) \oplus (I + J)^2 \oplus \ldots)
$$
它由仿射开子集 $A[\frac{I + J}{x}]$ 与
$A[\frac{I + J}{y}]$ 覆盖，其中 $x \in I$ 且 $y \in J$。由于
$x \in I$ 在 $A[\frac{I + J}{x}]$ 中是非零因子，并且 $IJ = 0$，可知
$J A[\frac{I + J}{x}] = 0$。由于 $y \in J$ 在
$A[\frac{I + J}{y}]$ 中是非零因子，并且 $IJ = 0$，可知
$I A[\frac{I + J}{y}] = 0$。此外，
$$
\Spec(A[\textstyle{\frac{I + J}{x}}]) \cap
\Spec(A[\textstyle{\frac{I + J}{y}}]) =
\Spec(A[\textstyle{\frac{I + J}{xy}}]) = \emptyset
$$
因为 $xy$ 既是非零因子又为零。所以 $b^{-1}(U)$ 是下列两个开子概形的
不交并：$U_1$ 定义为标准开集
$\Spec(A[\frac{I + J}{x}])$（$x \in I$）之并；而 $U_2$ 是仿射开集
$\Spec(A[\frac{I + J}{y}])$（$y \in J$）之并。我们已经看到，
$b^{-1}\mathcal{I}\mathcal{O}_{X'}$ 在 $U_2$ 上的限制为零，并且
$b^{-1}\mathcal{I}\mathcal{O}_{X'}$ 在 $U_1$ 上的限制为零。
我们略去验证这些开子概形可以黏合成全局开子概形 $X'_1$ 与 $X'_2$。
\end{proof}

\begin{lemma}
\label{divisors-lemma-blowing-up-denominators}
设 $X$ 为局部 Noether 概形。设 $\mathcal{L}$ 为可逆
$\mathcal{O}_X$-模。设 $s$ 为 $\mathcal{L}$ 的正则亚纯截面。设
$U \subset X$ 为最大的开子概形，使得 $s$ 对应于 $\mathcal{L}$ 在 $U$ 上的
一个截面。爆破 $b : X' \to X$ 是沿 $s$ 的分母理想作的；它是
$U$-容许的。存在有效 Cartier 除子 $D \subset X'$ 以及同构
$$
b^*\mathcal{L} = \mathcal{O}_{X'}(D - E),
$$
其中 $E \subset X'$ 是例外除子，并且经由这个同构，亚纯截面
$b^*s$ 对应于亚纯截面 $1_D \otimes (1_E)^{-1}$。
\end{lemma}

\begin{proof}
由定义 \ref{divisors-definition-regular-meromorphic-ideal-denominators} 中分母理想的定义，
立即可见 $b$ 是 $U$-容许爆破。关于记号 $1_D$、$1_E$ 和
$\mathcal{O}_{X'}(D - E)$，请见定义
\ref{divisors-definition-invertible-sheaf-effective-Cartier-divisor}。
拉回 $b^*s$ 由引理 \ref{divisors-lemma-blow-up-pullback-effective-Cartier} 和
\ref{divisors-lemma-meromorphic-sections-pullback} 定义。因此引理的陈述是有意义的。
可把最后一个断言重新解释为：$b^*s$ 是
$b^*\mathcal{L}(E)$ 的整体正则截面，其零点概形为 $D$。
这唯一确定了 $D$，所以为了证明引理，可以在 $X$ 与 $X'$ 上仿射局部地工作。
假设 $X = \Spec(A)$ 为仿射概形，并且
$\mathcal{L} = \mathcal{O}_X$。于是 $s$ 是正则亚纯函数；再作缩小后，
可假设 $s = a'/a$，其中 $a', a \in A$ 都是非零因子。
于是 $s$ 的分母理想对应于理想
$I = \{x \in A \mid xa' \in aA\}$。回忆，$X'$ 由仿射爆破代数的谱覆盖，
这些代数为 $A' = A[\frac{I}{x}]$，其中 $x \in I$
（引理 \ref{divisors-lemma-blowing-up-affine}）。固定 $x \in I$，并写成
$xa' = a a''$，其中某个 $a'' \in A$。除子 $E \subset X'$ 由
$x \in A'$ 截出（这是在 $A'$ 的谱上），因而 $1/x$ 是
$\mathcal{O}_{X'}(E)$ 在 $\Spec(A')$ 上的生成元。最后，在 $A'$ 的全商环中，
有 $a'/a = a''/x$。所以 $b^*s = a'/a$ 限制为
$\mathcal{O}_{X'}(E)$ 的正则截面；在 $\Spec(A')$ 上，这个截面由
$a''/x$ 给出。这就完成了证明。
（除子 $D \cap \Spec(A')$ 由 $a''$ 在 $A'$ 中的像截出。）
\end{proof}

\section{爆破与平坦性}
\label{divisors-section-blowup-flat}

\noindent
我们继续《代数进阶》\ref{more-algebra-section-blowup-flat} 节中开始的讨论。
我们还将在《平坦性进阶》\ref{flat-section-blowup-flat} 节中证明进一步的结果。


\begin{lemma}
\label{divisors-lemma-strict-transform-blowup-fitting-ideal}
设 $S$ 为概形。设 $\mathcal{F}$ 为有限型拟相干
$\mathcal{O}_S$-模。设 $Z_k \subset S$ 为由
$\text{Fit}_k(\mathcal{F})$ 截出的闭子概形，见
\ref{divisors-section-fitting-ideals} 节。设 $S' \to S$ 为 $S$ 沿 $Z_k$ 的爆破，
并设 $\mathcal{F}'$ 为 $\mathcal{F}$ 的严格变换。那么
$\mathcal{F}'$ 局部可由 $\leq k$ 个截面生成。
\end{lemma}

\begin{proof}
回忆，$\mathcal{F}'$ 局部可由 $\leq k$ 个截面生成，当且仅当
$\text{Fit}_k(\mathcal{F}') = \mathcal{O}_{S'}$，见引理
\ref{divisors-lemma-fitting-ideal-generate-locally}。所以本引理是《代数进阶》引理
\ref{more-algebra-lemma-blowup-fitting-ideal} 的翻译。
\end{proof}

\begin{lemma}
\label{divisors-lemma-strict-transform-blowup-fitting-ideal-locally-free}
设 $S$ 为概形。设 $\mathcal{F}$ 为有限型拟相干
$\mathcal{O}_S$-模。设 $Z_k \subset S$ 为由
$\text{Fit}_k(\mathcal{F})$ 截出的闭子概形，见
\ref{divisors-section-fitting-ideals} 节。假设 $\mathcal{F}$ 局部自由且秩为 $k$，
这是在 $S \setminus Z_k$ 上。设 $S' \to S$ 为 $S$ 沿 $Z_k$ 的爆破，并设
$\mathcal{F}'$ 为 $\mathcal{F}$ 的严格变换。那么 $\mathcal{F}'$ 局部自由且
秩为 $k$。
\end{lemma}

\begin{proof}
这是《代数进阶》引理
\ref{more-algebra-lemma-blowup-fitting-ideal-locally-free} 的翻译。
\end{proof}

\begin{lemma}
\label{divisors-lemma-blowup-fitting-ideal}
设 $X$ 为概形。设 $\mathcal{F}$ 为有限呈示的
$\mathcal{O}_X$-模。设 $U \subset X$ 为概形论稠密开集，使得
$\mathcal{F}|_U$ 为常秩 $r$ 的有限局部自由模。则：
\begin{enumerate}
\item 爆破 $b : X' \to X$ 是对 $X$ 沿第 $r$ 个 Fitting 理想
（属于 $\mathcal{F}$）作的，并且是 $U$-容许的；
\item 严格变换 $\mathcal{F}'$（取自 $\mathcal{F}$，关于 $b$）局部自由且秩为 $r$；
\item 核 $\mathcal{K}$（对应于满射 $b^*\mathcal{F} \to \mathcal{F}'$）
有限呈示，并且 $\mathcal{K}|_U = 0$；
\item $b^*\mathcal{F}$ 与 $\mathcal{K}$ 都是完美
$\mathcal{O}_{X'}$-模，其 Tor 维数 $\leq 1$。
\end{enumerate}
\end{lemma}

\begin{proof}
由引理 \ref{divisors-lemma-fitting-ideal-of-finitely-presented}，理想
$\text{Fit}_r(\mathcal{F})$ 是有限型的；由引理
\ref{divisors-lemma-fitting-ideal-finite-locally-free}，它在 $U$ 上的限制等于
$\mathcal{O}_U$。所以 $b : X' \to X$ 是 $U$-容许的，见定义
\ref{divisors-definition-admissible-blowup}。

\medskip\noindent
由引理 \ref{divisors-lemma-fitting-ideal-finite-locally-free}，
$\text{Fit}_{r - 1}(\mathcal{F})$ 在 $U$ 上的限制为零；由于 $U$ 概形论稠密，
可知 $\text{Fit}_{r - 1}(\mathcal{F}) = 0$ 在整个 $X$ 上成立。因此，再由引理
\ref{divisors-lemma-fitting-ideal-finite-locally-free}，$\mathcal{F}$ 局部自由且秩为 $r$，
这是在下述子概形的补集上：该子概形由第 $r$ 个 Fitting 理想
（属于 $\mathcal{F}$）截出（这个补集可能大于 $U$，这正是论证中必须完成这一步的原因）。
所以，由引理 \ref{divisors-lemma-strict-transform-blowup-fitting-ideal-locally-free}，
严格变换
$$
b^*\mathcal{F} \longrightarrow \mathcal{F}'
$$
局部自由且秩为 $r$。这个映射的核 $\mathcal{K}$ 支集在爆破 $b$ 的
例外除子上，因而 $\mathcal{K}|_U = 0$。最后，由于
$\mathcal{F}'$ 有限局部自由，并且所示箭头是满射，可以在 $X'$ 上局部地把
$b^*\mathcal{F}$ 写成 $\mathcal{K}$ 与 $\mathcal{F}'$ 的直和。
由于 $b^*\mathcal{F}'$ 是有限呈示的
（《模》引理 \ref{modules-lemma-pullback-finite-presentation}），
$\mathcal{K}$ 也如此。

\medskip\noindent
关于 Tor 维数的陈述由《代数进阶》引理
\ref{more-algebra-lemma-fitting-ideals-and-pd1} 得出。
\end{proof}

\section{修改}
\label{divisors-section-modifications}

\noindent
本节将汇集如下类型的结果：作一次修改之后，某某性质成立。以后将看到，
修改可由一次爆破支配（《平坦性进阶》引理
\ref{flat-lemma-dominate-modification-by-blowup}）。

\begin{lemma}
\label{divisors-lemma-filter-after-modification}
设 $X$ 为整概形。设 $\mathcal{E}$ 为有限局部自由
$\mathcal{O}_X$-模。存在修改 $f : X' \to X$，使得
$f^*\mathcal{E}$ 有一个滤过，其逐次商都是可逆 $\mathcal{O}_{X'}$-模。
\end{lemma}

\begin{proof}
关于秩 $r$ 作归纳证明，其中这是 $\mathcal{E}$ 的秩。如果 $r = 1$ 或 $r = 0$，
引理显然成立。假设 $r > 1$。令
$P = \mathbf{P}(\mathcal{E})$，其结构态射为 $\pi : P \to X$，见《构造》
\ref{constructions-section-projective-bundle} 节。于是 $\pi$ 是固有的
（引理 \ref{divisors-lemma-relative-proj-proper}）。有典范满射
$$
\pi^*\mathcal{E} \to \mathcal{O}_P(1)
$$
其核有限局部自由且秩为 $r - 1$。选取非空开子概形
$U \subset X$，使得 $\mathcal{E}|_U \cong \mathcal{O}_U^{\oplus r}$。
于是 $P_U = \pi^{-1}(U)$ 同构于 $\mathbf{P}^{r - 1}_U$。
特别地，存在截面 $s : U \to P_U$（属于 $\pi$）。设
$X' \subset P$ 为态射 $U \to P_U \to P$ 的概形论像。那么
$X'$ 是整概形（《态射》引理
\ref{morphisms-lemma-scheme-theoretic-image-reduced}），态射
$f = \pi|_{X'} : X' \to X$ 是固有的（《态射》引理
\ref{morphisms-lemma-closed-immersion-proper} 和
\ref{morphisms-lemma-composition-proper}），并且
$f^{-1}(U) \to U$ 是同构。所以 $f$ 是修改
（《态射》定义 \ref{morphisms-definition-modification}）。
由构造，拉回 $f^*\mathcal{E}$ 有一个二步滤过；它的商可逆，因为该商等于
$\mathcal{O}_P(1)|_{X'}$；它的子模 $\mathcal{E}'$ 局部自由且秩为
$r - 1$。由归纳假设，可以找到修改 $g : X'' \to X'$，使得
$g^*\mathcal{E}'$ 有如引理陈述中的滤过。因此
$f \circ g : X'' \to X$ 就是所求的修改。
\end{proof}

\begin{lemma}
\label{divisors-lemma-extend-rational-map-after-modification}
设 $S$ 为概形。设 $X$、$Y$ 为 $S$ 上概形。假设 $X$ 是 Noether 的，
并且 $Y$ 在 $S$ 上固有。给定 $S$-有理映射
$f : U \to Y$，其源概形为 $X$、目标为 $Y$，则存在态射 $p : X' \to X$ 与
$S$-态射 $f' : X' \to Y$，使得：
\begin{enumerate}
\item $p$ 是固有的，并且 $p^{-1}(U) \to U$ 是同构；
\item $f'|_{p^{-1}(U)}$ 等于 $f \circ p|_{p^{-1}(U)}$。
\end{enumerate}
\end{lemma}

\begin{proof}
记包含态射 $j : U \to X$。设 $X' \subset Y \times_S X$ 为
$(f, j) : U \to Y \times_S X$ 的概形论像（《态射》定义
\ref{morphisms-definition-scheme-theoretic-image}）。投影
$g : Y \times_S X \to X$ 是固有的（《态射》引理
\ref{morphisms-lemma-base-change-proper}）。态射 $p : X' \to X$ 是
$X' \to Y \times_S X$ 与 $g$ 的复合；它是固有的（《态射》引理
\ref{morphisms-lemma-closed-immersion-proper} 和
\ref{morphisms-lemma-composition-proper}）。由于 $g$ 是分离的，并且
$U \subset X$ 是逆紧的（因为 $X$ 是 Noether 的），由《态射》引理
\ref{morphisms-lemma-scheme-theoretic-image-of-partial-section} 可知
$p^{-1}(U) \to U$ 是同构。另一方面，$f' : X' \to Y$ 是
$X' \to Y \times_S X$ 与投影 $Y \times_S X \to Y$ 的复合；它与
$f$ 一致，这是在 $p^{-1}(U)$ 上。
\end{proof}
% Generated mechanically from the frozen limits.tex target.
% Do not edit this generated file; edit the bound locale source instead.

% OK, start here.
%

\raggedbottom

\setcounter{chapter}{31}
\chapter{概形的极限}



\phantomsection
\label{limits-section-phantom}


\section{引言}
\label{limits-section-introduction}

\noindent
本章汇集与概形极限有关的材料。我们主要研究以有向集
（《范畴》定义 \ref{categories-definition-directed-set}）为指标、
过渡态射为仿射态射的逆系之极限。我们讨论绝对 Noether 逼近。
我们把基上局部有限呈示概形刻画为如下概形：其相伴点函子保持极限。
作为绝对 Noether 逼近的一个应用，我们证明整态射下仿射概形的像是仿射的。
此外，我们证明 Chow 引理的一些非常一般的变体。
基础参考文献是 \cite{EGA}。




\section{带仿射过渡态射的概形有向极限}
\label{limits-section-limits}

\noindent
本节构造这种极限。

\begin{lemma}
\label{limits-lemma-directed-inverse-system-affine-schemes-has-limit}
设 $I$ 为有向集。设 $(S_i, f_{ii'})$ 为以 $I$ 为指标的概形逆系。
如果所有概形 $S_i$ 都是仿射的，那么极限 $S = \lim_i S_i$ 在概形范畴中存在。
事实上，$S$ 是仿射的，并且 $S = \Spec(\colim_i R_i)$，其中
$R_i = \Gamma(S_i, \mathcal{O})$。
\end{lemma}

\begin{proof}
直接定义 $S = \Spec(\colim_i R_i)$。由《概形》引理
\ref{schemes-lemma-morphism-into-affine}，可知 $S$ 即使在局部环空间范畴中
也是该极限。
\end{proof}

\begin{lemma}
\label{limits-lemma-directed-inverse-system-has-limit}
设 $I$ 为有向集。设 $(S_i, f_{ii'})$ 为以 $I$ 为指标的概形逆系。
如果所有态射 $f_{ii'} : S_i \to S_{i'}$ 都是仿射的，那么极限
$S = \lim_i S_i$ 在概形范畴中存在。此外：
\begin{enumerate}
\item 每个态射 $f_i : S \to S_i$ 都是仿射的；
\item 对元素 $0 \in I$ 以及任意开子概形 $U_0 \subset S_0$，有
$$
f_0^{-1}(U_0) = \lim_{i \geq 0} f_{i0}^{-1}(U_0)
$$
这是概形范畴中的等式。
\end{enumerate}
\end{lemma}

\begin{proof}
选取元素 $0 \in I$。注意，$I$ 非空，因为所述极限是有向的。
对每个 $i \geq 0$，考虑拟相干 $\mathcal{O}_{S_0}$-代数层
$\mathcal{A}_i = f_{i0, *}\mathcal{O}_{S_i}$。回忆，
$S_i = \underline{\Spec}_{S_0}(\mathcal{A}_i)$，见《态射》引理
\ref{morphisms-lemma-characterize-affine}。置
$\mathcal{A} = \colim_{i \geq 0} \mathcal{A}_i$。这是拟相干
$\mathcal{O}_{S_0}$-代数层，见《概形》\ref{schemes-section-quasi-coherent} 节。
置 $S = \underline{\Spec}_{S_0}(\mathcal{A})$。由《态射》引理
\ref{morphisms-lemma-affine-equivalence-algebras}，对于 $i \geq 0$，
得到与过渡态射相容的态射 $f_i : S \to S_i$。例如，由《态射》引理
\ref{morphisms-lemma-affine-permanence}，这些态射 $f_i$ 是仿射的。
由上面的引理 \ref{limits-lemma-directed-inverse-system-affine-schemes-has-limit}，
可知对于任意仿射开集 $U_0 \subset S_0$，逆像
$U = f_0^{-1}(U_0) \subset S$ 是开集逆系
$U_i = f_{i0}^{-1}(U_0)$（$i \geq 0$）在概形范畴中的极限。

\medskip\noindent
设 $T$ 为概形。设 $g_i : T \to S_i$ 为相容态射系。为了证明
$S = \lim_i S_i$，须证存在唯一态射 $g : T \to S$，使得都有
$g_i = f_i \circ g$，这是对所有 $i \in I$。对每个 $t \in T$，存在
仿射开集 $U_0 \subset S_0$，它包含 $g_0(t)$。设
$V \subset g_0^{-1}(U_0)$ 为包含 $t$ 的仿射开邻域。由上述观察，得到唯一态射
$g_V : V \to U = f_0^{-1}(U_0)$，使得都有
$f_i \circ g_V = g_i|_{U_i}$，这是对所有 $i$。如此构造的开集 $V \subset T$ 构成
$T$ 的拓扑之一个基。由于唯一性，态射 $g_V$ 黏合成态射
$g : T \to S$。这就给出了所需态射 $g : T \to S$。

\medskip\noindent
最后一个陈述由上面的极限构造显然成立。
\end{proof}

\begin{lemma}
\label{limits-lemma-scheme-over-limit}
设 $I$ 为有向集。设 $(S_i, f_{ii'})$ 为以 $I$ 为指标的概形逆系。
假设所有态射 $f_{ii'} : S_i \to S_{i'}$ 都是仿射的，令
$S = \lim_i S_i$。令 $0 \in I$。设 $T$ 为 $S_0$ 上概形。则
$$
T \times_{S_0} S = \lim_{i \geq 0} T \times_{S_0} S_i
$$
\end{lemma}

\begin{proof}
由引理 \ref{limits-lemma-directed-inverse-system-has-limit}，右端是概形。
该等式是形式的，见《范畴》引理
\ref{categories-lemma-colimits-commute}。
\end{proof}

\section{无限乘积}
\label{limits-section-inifinite-products}

\noindent
概形的无限乘积通常不存在。例如，《例》\ref{examples-section-not-algebraic} 节
证明了 $\mathbf{P}^1$ 的无限多个副本之乘积甚至不是代数空间。

\medskip\noindent
另一方面，仿射概形的无限乘积确实存在，并且仍为仿射概形。利用《概形》引理
\ref{schemes-lemma-morphism-into-affine}，这对应于环范畴中存在无限余积这一事实：
若 $I$ 是集合，且 $R_i$ 对每个 $i$ 都是环，那么可以考虑环
$$
R = \otimes R_i =
\colim_{\{i_1, \ldots, i_n\} \subset I}
R_{i_1} \otimes_\mathbf{Z} \ldots \otimes_\mathbf{Z} R_{i_n}
$$
给定另一个环 $A$，映射 $R \to A$ 等同于给出一族环映射
$R_i \to A$（$i \in I$）；这由有限张量积的相应性质推出。

\begin{lemma}
\label{limits-lemma-infinite-product}
\begin{slogan}
仿射概形的无限乘积存在且为仿射概形。
\end{slogan}
设 $S$ 为概形。设 $I$ 为集合，并且对每个 $i \in I$，设
$f_i : T_i \to S$ 为仿射态射。则乘积 $T = \prod T_i$ 在 $S$ 上概形范畴中存在。
事实上，有
$$
T = \lim_{\{i_1, \ldots, i_n\} \subset I}
T_{i_1} \times_S \ldots \times_S T_{i_n}
$$
并且投影态射 $T \to T_{i_1} \times_S \ldots \times_S T_{i_n}$ 是仿射的。
\end{lemma}

\begin{proof}
略。提示：仿照引理前的讨论论证，并用引理
\ref{limits-lemma-directed-inverse-system-has-limit} 证明极限存在。
\end{proof}

\begin{lemma}
\label{limits-lemma-infinite-product-surjective}
设 $S$ 为概形。设 $I$ 为集合，并且对每个 $i \in I$，设
$f_i : T_i \to S$ 为满射仿射态射。则乘积 $T = \prod T_i$
在 $S$ 上概形范畴中（引理 \ref{limits-lemma-infinite-product}）满射到 $S$。
\end{lemma}

\begin{proof}
设 $s \in S$。选取 $t_i \in T_i$ 使其映到 $s$。选取足够大的域扩张
$K/\kappa(s)$，使得 $\kappa(s_i)$ 嵌入 $K$，这是对每个 $i$。于是得到态射
$\Spec(K) \to T_i$，其像为 $s_i$，并且它们作为到 $S$ 的态射彼此相同。
因而得到态射 $\Spec(K) \to T$，这证明存在 $T$ 中映到 $s$ 的点。
\end{proof}

\begin{lemma}
\label{limits-lemma-infinite-product-integral}
设 $S$ 为概形。设 $I$ 为集合，并且对每个 $i \in I$，设
$f_i : T_i \to S$ 为整态射。则乘积 $T = \prod T_i$
在 $S$ 上概形范畴中（引理 \ref{limits-lemma-infinite-product}）在 $S$ 上是整的。
\end{lemma}

\begin{proof}
略。提示：在仿射片上，这归结为下述代数事实：若 $A \to B_i$
对所有 $i$ 都是整的，则 $A \to \otimes_A B_i$ 是整的。
\end{proof}


\section{性质的下降}
\label{limits-section-descent}

\noindent
先给出若干描述极限拓扑的基本引理。

\begin{lemma}
\label{limits-lemma-inverse-limit-sets}
设 $S = \lim S_i$ 是一个以仿射态射为过渡态射的概形有向逆系之极限
（引理 \ref{limits-lemma-directed-inverse-system-has-limit}）。则
$S_{set} = \lim_i S_{i, set}$，其中 $S_{set}$ 表示概形 $S$ 的底集。
\end{lemma}

\begin{proof}
选取 $i \in I$。取仿射开集 $U_i \subset S_i$。记
$U_{i'} = f_{i'i}^{-1}(U_i)$ 以及 $U = f_i^{-1}(U_i)$。
这里 $f_{i'i} : S_{i'} \to S_i$ 是过渡态射，而
$f_i : S \to S_i$ 是投影。由引理
\ref{limits-lemma-directed-inverse-system-has-limit}，有
$U = \lim_{i' \geq i} U_i$。假设可以证明
$U_{set} = \lim_{i' \geq i} U_{i', set}$。那么用 $S_i$ 的一个仿射覆盖作简单论证，
即可推出该引理。因此可以假设所有 $S_i$ 以及 $S$ 都是仿射的。
这把问题化为下一段所考察的代数问题。

\medskip\noindent
设给定环系 $(A_i, \varphi_{ii'})$，其指标集为 $I$。置
$A = \colim_i A_i$，其典范映射为 $\varphi_i : A_i \to A$。则
$$
\Spec(A) = \lim_i \Spec(A_i)
$$
事实上，设给定素理想 $\mathfrak p_i \subset A_i$，且等式
$\mathfrak p_i = \varphi_{ii'}^{-1}(\mathfrak p_{i'})$ 对所有 $i' \geq i$ 都成立。
只需置
$$
\mathfrak p =
\{x \in A
\mid
\exists i, x_i \in \mathfrak p_i \text{ 其中 }\varphi_i(x_i) = x\}
$$
显然这是一个理想，并且满足
$\varphi_i^{-1}(\mathfrak p) = \mathfrak p_i$。由此也很容易看出它是素理想。
\end{proof}

\begin{lemma}
\label{limits-lemma-inverse-limit-top}
\begin{reference}
\cite[IV, Proposition 8.2.9]{EGA}
\end{reference}
设 $S = \lim S_i$ 是一个以仿射态射为过渡态射的概形有向逆系之极限
（引理 \ref{limits-lemma-directed-inverse-system-has-limit}）。则
$S_{top} = \lim_i S_{i, top}$，其中 $S_{top}$ 表示概形 $S$ 的底拓扑空间。
\end{lemma}

\begin{proof}
我们使用《拓扑》引理 \ref{topology-lemma-characterize-limit} 的判别准则。
在引理 \ref{limits-lemma-inverse-limit-sets} 中已经看到
$S_{set} = \lim_i S_{i, set}$。映射 $f_i : S \to S_i$ 是概形态射，
因而连续。因此，$f_i^{-1}(U_i)$ 对每个开集 $U_i \subset S_i$ 都是开集。
最后，设 $s \in S$，并设 $s \in V \subset S$ 是一个开邻域。
选取 $0 \in I$，并选取仿射开邻域 $U_0 \subset S_0$，它包含 $s$ 之像。于是
$f_0^{-1}(U_0) = \lim_{i \geq 0} f_{i0}^{-1}(U_0)$，见引理
\ref{limits-lemma-directed-inverse-system-has-limit}。此时
$f_0^{-1}(U_0)$ 以及 $f_{i0}^{-1}(U_0)$ 都是仿射的，并且
$$
\mathcal{O}_S(f_0^{-1}(U_0)) =
\colim_{i \geq 0} \mathcal{O}_{S_i}(f_{i0}^{-1}(U_0))
$$
这可由引理 \ref{limits-lemma-directed-inverse-system-has-limit} 的证明得到，
也可由引理 \ref{limits-lemma-directed-inverse-system-affine-schemes-has-limit} 得到。
选取 $a \in \mathcal{O}_S(f_0^{-1}(U_0))$，使得
$s \in D(a) \subset V$。这是可能的，因为主开集构成仿射概形
$f_0^{-1}(U_0)$ 的拓扑之一个基。于是可以选取 $i \geq 0$ 以及
$a_i \in \mathcal{O}_{S_i}(f_{i0}^{-1}(U_0))$，使后者映到 $a$。
由此，$D(a_i) \subset f_{i0}^{-1}(U_0) \subset S_i$ 是一个开子集，
其在 $S$ 中的逆像为 $D(a)$。证明完毕。
\end{proof}

\begin{lemma}
\label{limits-lemma-limit-nonempty}
设 $S = \lim S_i$ 是一个以仿射态射为过渡态射的概形有向逆系之极限
（引理 \ref{limits-lemma-directed-inverse-system-has-limit}）。如果所有概形
$S_i$ 都非空且拟紧，那么极限 $S = \lim_i S_i$ 非空。
\end{lemma}

\begin{proof}
选取 $0 \in I$。注意，$I$ 非空，因为所述极限是有向的。
选取仿射开覆盖 $S_0 = \bigcup_{j = 1, \ldots, m} U_j$。
由于 $I$ 是有向的，存在 $j \in \{1, \ldots, m\}$，使得
$f_{i0}^{-1}(U_j) \not = \emptyset$ 对所有 $i \geq 0$ 都成立。因此
$\lim_{i \geq 0} f_{i0}^{-1}(U_j)$ 非空，因为非零环的有向余极限非零
（这是由于 $1 \not = 0$）。由于
$\lim_{i \geq 0} f_{i0}^{-1}(U_j)$ 是该极限的开子概形，故得证。
\end{proof}

\begin{lemma}
\label{limits-lemma-inverse-limit-irreducibles}
设 $S = \lim S_i$ 是一个以仿射态射为过渡态射的概形有向逆系之极限
（引理 \ref{limits-lemma-directed-inverse-system-has-limit}）。设
$s \in S$，其像为 $s_i \in S_i$。则：
\begin{enumerate}
\item 作为概形，$s = \lim s_i$，即 $\kappa(s) = \colim \kappa(s_i)$；
\item 作为集合，$\overline{\{s\}} = \lim \overline{\{s_i\}}$；
\item 作为概形，$\overline{\{s\}} = \lim \overline{\{s_i\}}$，
其中 $\overline{\{s\}}$ 以及 $\overline{\{s_i\}}$ 都赋予约化诱导概形结构。
\end{enumerate}
\end{lemma}

\begin{proof}
选取 $0 \in I$ 以及一个仿射开覆盖
$S_0 = \bigcup_{j \in J} U_{0, j}$。对 $i \geq 0$，令
$U_{i, j} = f_{i, 0}^{-1}(U_{0, j})$，并置
$U_j = f_0^{-1}(U_{0, j})$。这里 $f_{i'i} : S_{i'} \to S_i$ 是过渡态射，
而 $f_i : S \to S_i$ 是投影。对 $j \in J$，下列条件等价：
(a) $s \in U_j$；(b) $s_0 \in U_{0, j}$；
(c) $s_i \in U_{i, j}$ 对所有 $i \geq 0$ 都成立。
设 $J' \subset J$ 是使 (a)、(b)、(c) 成立的指标集合。则
$\overline{\{s\}} = \bigcup_{j \in J'} (\overline{\{s\}} \cap U_j)$，
并且 $\overline{\{s_i\}}$ 对 $i \geq 0$ 也有类似等式。
注意，$\overline{\{s\}} \cap U_j$ 是集合 $\{s\}$ 在拓扑空间
$U_j$ 中的闭包。类似地，$\overline{\{s_i\}} \cap U_{i, j}$ 对
$i \geq 0$ 也是如此。因此，只需证明 $S$ 和 $S_i$ 对所有 $i$
都是仿射时的情形。这把问题化为下一段所考察的代数问题。

\medskip\noindent
设给定环系 $(A_i, \varphi_{ii'})$，其指标集为 $I$。置
$A = \colim_i A_i$，其典范映射为 $\varphi_i : A_i \to A$。
设 $\mathfrak p \subset A$ 为素理想，并置
$\mathfrak p_i = \varphi_i^{-1}(\mathfrak p)$。则
$$
V(\mathfrak p) = \lim_i V(\mathfrak p_i)
$$
这由引理 \ref{limits-lemma-inverse-limit-sets} 推出，因为
$A/\mathfrak p = \colim A_i/\mathfrak p_i$。该环等式还表明，
关于约化诱导概形结构的最后一个陈述成立。等式
$\kappa(\mathfrak p) = \colim \kappa(\mathfrak p_i)$ 也由此推出。
\end{proof}

\noindent
本节其余部分都在下述情形中工作。

\begin{situation}
\label{limits-situation-descent}
设 $S = \lim_{i \in I} S_i$ 是一个以仿射态射
$f_{i'i} : S_{i'} \to S_i$ 为过渡态射的概形有向系之极限
（引理 \ref{limits-lemma-directed-inverse-system-has-limit}）。假设 $S_i$
对所有 $i \in I$ 都是拟紧且拟分离的。我们用
$f_i : S \to S_i$ 表示投影。还选定一个元素 $0 \in I$。
\end{situation}

\noindent
在此情形下，态射 $S \to S_0$ 是仿射的。由此可知，$S$ 拟紧且拟分离
\footnote{这由《态射》引理 \ref{morphisms-lemma-affine-separated}、
《拓扑》定义 \ref{topology-definition-quasi-compact} 以及
《概形》引理 \ref{schemes-lemma-separated-permanence} 推出。}。
我们所寻求的结果具有如下形式：若在 $S$ 上有一个对象，那么对某个 $i$，
在 $S_i$ 上存在一个类似对象。

\begin{lemma}
\label{limits-lemma-topology-limit}
在情形 \ref{limits-situation-descent} 中：
\begin{enumerate}
\item 有 $S_{set} = \lim_i S_{i, set}$，其中 $S_{set}$ 表示概形 $S$ 的底集。
\item 有 $S_{top} = \lim_i S_{i, top}$，其中 $S_{top}$ 表示概形 $S$ 的底拓扑空间。
\item 若 $s, s' \in S$，且 $s'$ 不是 $s$ 的特化，则对某个 $i \in I$，
像 $s'_i \in S_i$（$s'$ 的像）不是像 $s_i \in S_i$（$s$ 的像）的特化。
\item 在此加入更多关于 $S$ 拓扑的简单事实。
（要求：所加入的任何内容在仿射情形下都应易于证明。）
\end{enumerate}
\end{lemma}

\begin{proof}
第 (1) 部分是引理 \ref{limits-lemma-inverse-limit-sets} 的特殊情形。

\medskip\noindent
第 (2) 部分是引理 \ref{limits-lemma-inverse-limit-top} 的特殊情形。

\medskip\noindent
第 (3) 部分是引理 \ref{limits-lemma-inverse-limit-irreducibles} 的特殊情形。
\end{proof}

\begin{lemma}
\label{limits-lemma-descend-section}
在情形 \ref{limits-situation-descent} 中，设 $\mathcal{F}_0$ 是 $S_0$ 上的拟相干层。
置 $\mathcal{F}_i = f_{i0}^*\mathcal{F}_0$（$i \geq 0$），并置
$\mathcal{F} = f_0^*\mathcal{F}_0$。则
$$
\Gamma(S, \mathcal{F}) = \colim_{i \geq 0} \Gamma(S_i, \mathcal{F}_i)
$$
\end{lemma}

\begin{proof}
记 $\mathcal{A}_j = f_{i0, *} \mathcal{O}_{S_i}$。这是一个拟相干
$\mathcal{O}_{S_0}$-代数层（见《态射》引理
\ref{morphisms-lemma-affine-equivalence-algebras}），并且 $S_i$ 是
$\mathcal{A}_i$ 在 $S_0$ 上的相对谱。在引理
\ref{limits-lemma-directed-inverse-system-has-limit} 的证明中，我们把 $S$ 构造为
$\mathcal{A} = \colim_{i \geq 0} \mathcal{A}_i$ 在 $S_0$ 上的相对谱。置
$$
\mathcal{M}_i = \mathcal{F}_0 \otimes_{\mathcal{O}_{S_0}} \mathcal{A}_i
$$
以及
$$
\mathcal{M} = \mathcal{F}_0 \otimes_{\mathcal{O}_{S_0}} \mathcal{A}.
$$
则有 $f_{i0, *} \mathcal{F}_i = \mathcal{M}_i$ 以及
$f_{0, *}\mathcal{F} = \mathcal{M}$。由于 $\mathcal{A}$ 是层
$\mathcal{A}_i$ 的余极限，并且张量积与有向余极限交换，故有
$\mathcal{M} = \colim_{i \geq 0} \mathcal{M}_i$。由于 $S_0$ 拟紧且拟分离，
可见
\begin{eqnarray*}
\Gamma(S, \mathcal{F})
& = &
\Gamma(S_0, \mathcal{M}) \\
& = &
\Gamma(S_0, \colim_{i \geq 0} \mathcal{M}_i) \\
& = &
\colim_{i \geq 0} \Gamma(S_0, \mathcal{M}_i) \\
& = &
\colim_{i \geq 0} \Gamma(S_i, \mathcal{F}_i)
\end{eqnarray*}
中间的等号见《层》引理 \ref{sheaves-lemma-directed-colimits-sections} 以及
《拓扑》引理 \ref{topology-lemma-topology-quasi-separated-scheme}。
\end{proof}

\begin{lemma}
\label{limits-lemma-limit-closed-nonempty}
在情形 \ref{limits-situation-descent} 中，设对每个 $i$ 都给定一个非空闭子集
$Z_i \subset S_i$，并且有 $f_{i'i}(Z_{i'}) \subset Z_i$ 对所有
$i' \geq i$ 成立。则存在一点 $s \in S$，使得 $f_i(s) \in Z_i$
对所有 $i$ 成立。
\end{lemma}

\begin{proof}
仍用 $Z_i \subset S_i$ 表示与 $Z_i$ 相关联的既约闭子概形，见《概形》
定义 \ref{schemes-definition-reduced-induced-scheme}。
闭浸入是仿射的，而仿射态射的复合仍是仿射的（见《态射》引理
\ref{morphisms-lemma-closed-immersion-affine} 和
\ref{morphisms-lemma-composition-affine}），所以 $Z_{i'} \to S_i$ 在
$i' \geq i$ 时是仿射的。由《态射》引理
\ref{morphisms-lemma-affine-permanence} 可知态射
$f_{i'i} : Z_{i'} \to Z_i$ 是仿射的。每个概形 $Z_i$ 都是拟紧概形的闭子概形，
因而是拟紧的。因此可以应用引理 \ref{limits-lemma-limit-nonempty}，得到
$Z = \lim_i Z_i$ 非空。由于有典范态射 $Z \to S$，结论成立。
\end{proof}

\begin{lemma}
\label{limits-lemma-limit-fibre-product-empty}
在情形 \ref{limits-situation-descent} 中，设给定某个 $i$ 以及态射 $T \to S_i$，满足
\begin{enumerate}
\item $T \times_{S_i} S = \emptyset$；
\item $T$ 拟紧。
\end{enumerate}
则 $T \times_{S_i} S_{i'} = \emptyset$ 对所有充分大的 $i'$ 成立。
\end{lemma}

\begin{proof}
由引理 \ref{limits-lemma-scheme-over-limit} 可知
$T \times_{S_i} S = \lim_{i' \geq i} T \times_{S_i} S_{i'}$。
于是结论由引理 \ref{limits-lemma-limit-nonempty} 得出。
\end{proof}

\begin{lemma}
\label{limits-lemma-limit-contained-in-constructible}
在情形 \ref{limits-situation-descent} 中，设给定某个 $i$ 以及一个局部可构造子集
$E \subset S_i$，并且 $f_i(S) \subset E$。则
$f_{i'i}(S_{i'}) \subset E$ 对所有充分大的 $i'$ 成立。
\end{lemma}

\begin{proof}
把 $S_i$ 写成有限个仿射开子概形之并，就把问题约化到 $S_i$ 仿射且 $E$
可构造的情形，见引理 \ref{limits-lemma-directed-inverse-system-has-limit} 以及
《性质》引理 \ref{properties-lemma-locally-constructible}。
在此情形下，补集 $S_i \setminus E$ 也是可构造的。因此存在一个仿射概形 $T$
以及一个态射 $T \to S_i$，其像为 $S_i \setminus E$，见《代数》引理
\ref{algebra-lemma-constructible-is-image}。由引理
\ref{limits-lemma-limit-fibre-product-empty} 可知，$T \times_{S_i} S_{i'}$
对所有充分大的 $i'$ 都为空，因而 $f_{i'i}(S_{i'}) \subset E$
对所有充分大的 $i'$ 成立。
\end{proof}

\begin{lemma}
\label{limits-lemma-descend-opens}
在情形 \ref{limits-situation-descent} 中有以下结论：
\begin{enumerate}
\item 给定任意拟紧开集 $V \subset S = \lim_i S_i$，存在 $i \in I$ 以及
拟紧开集 $V_i \subset S_i$，使得 $f_i^{-1}(V_i) = V$。
\item 设 $V_i \subset S_i$ 和 $V_{i'} \subset S_{i'}$ 是拟紧开集，且
$f_i^{-1}(V_i) = f_{i'}^{-1}(V_{i'})$。则存在指标
$i'' \geq i, i'$，使得
$f_{i''i}^{-1}(V_i) = f_{i''i'}^{-1}(V_{i'})$。
\item 若 $V_{1, i}, \ldots, V_{n, i} \subset S_i$ 是拟紧开集，且
$S = f_i^{-1}(V_{1, i}) \cup \ldots \cup f_i^{-1}(V_{n, i})$，则有
$S_{i'} = f_{i'i}^{-1}(V_{1, i}) \cup \ldots \cup f_{i'i}^{-1}(V_{n, i})$
对某个 $i' \geq i$ 成立。
\end{enumerate}
\end{lemma}

\begin{proof}
选取 $i_0 \in I$。注意，因为极限是有向的，所以 $I$ 非空。为方便起见，记
$S_0 = S_{i_0}$ 且记 $i_0 = 0$。选取仿射开覆盖
$S_0 = U_{1, 0} \cup \ldots \cup U_{m, 0}$。记
$U_{j, i} \subset S_i$ 为 $U_{j, 0}$ 在过渡态射下的逆像（$i \geq 0$），
并记 $U_j$ 为
$U_{j, 0}$ 在 $S$ 中的逆像。注意，$U_j = \lim_i U_{j, i}$ 是仿射概形的极限。

\medskip\noindent
先证明唯一性陈述。设 $V_i \subset S_i$ 和 $V_{i'} \subset S_{i'}$ 是拟紧开集，
且 $f_i^{-1}(V_i) = f_{i'}^{-1}(V_{i'})$。只需证明
$f_{i''i}^{-1}(V_i \cap U_{j, i''})$ 和
$f_{i''i'}^{-1}(V_{i'} \cap U_{j, i''})$ 当 $i''$ 充分大时相等。
于是问题约化为仿射概形极限的情形。在此情形下，写
$S = \Spec(R)$ 和 $S_i = \Spec(R_i)$（对所有 $i \in I$）。可写成
$V_i = S_i \setminus V(h_1, \ldots, h_m)$ 以及
$V_{i'} = S_{i'} \setminus V(g_1, \ldots, g_n)$。
假设意味着理想 $\sum g_jR$ 和 $\sum h_jR$ 在 $R$ 中具有相同的根基。
这意味着有 $g_j^N = \sum a_{jj'}h_{j'}$ 及
$h_j^N = \sum b_{jj'} g_{j'}$，其中某个 $N \gg 0$，且
$a_{jj'}$ 和 $b_{jj'}$ 属于 $R$。
由于 $R = \colim_i R_i$，可以选取指标 $i'' \geq i$，使得等式
$g_j^N = \sum a_{jj'}h_{j'}$ 和
$h_j^N = \sum b_{jj'} g_{j'}$ 在 $R_{i''}$ 中成立，其中某些
$a_{jj'}$ 和 $b_{jj'}$ 属于 $R_{i''}$。这说明理想
$\sum g_jR_{i''}$ 和 $\sum h_jR_{i''}$ 在 $R_{i''}$ 中具有相同的根基，
正如所需。

\medskip\noindent
下面证明存在性。若 $S_0$ 仿射，则有 $S_i = \Spec(R_i)$ 对所有
$i \geq 0$ 成立，并且 $S = \Spec(R)$，其中 $R = \colim R_i$。
于是有 $V = S \setminus V(g_1, \ldots, g_n)$，其中某些
$g_1, \ldots, g_n \in R$。选取充分大的 $i$，使每个 $g_j$
都来自某个元素 $g_{j, i} \in R_i$，并取
$V_i = S_i \setminus V(g_{1, i}, \ldots, g_{n, i})$。
若 $S_0$ 一般，则开集 $V \cap U_j$ 是拟紧的，因为 $S$ 拟分离。
因此由仿射情形可知，对每个 $j = 1, \ldots, m$，存在 $i_j \in I$
以及拟紧开集 $V_{i_j} \subset U_{j, i_j}$，其在 $U_j$ 中的逆像为
$V \cap U_j$。置 $i = \max(i_1, \ldots, i_m)$，并令
$V_i = \bigcup f_{ii_j}^{-1}(V_{i_j})$。

\medskip\noindent
关于覆盖的陈述由唯一性陈述应用于两个开集
$V_{1, i} \cup \ldots \cup V_{n, i}$ 和 $S_i$（它们都是 $S_i$ 的开集）得出。
\end{proof}

\begin{lemma}
\label{limits-lemma-limit-quasi-affine}
在情形 \ref{limits-situation-descent} 中，若 $S$ 拟仿射，则对某个 $i_0 \in I$，
概形 $S_i$ 在 $i \geq i_0$ 时都拟仿射。
\end{lemma}

\begin{proof}
选取 $i_0 \in I$。注意，因为极限是有向的，所以 $I$ 非空。为方便起见，记
$S_0 = S_{i_0}$ 且记 $i_0 = 0$。取 $s \in S$。可以选取一个仿射开集
$U_0 \subset S_0$，它包含 $f_0(s)$。由于 $S$ 拟仿射，可以选取元素
$a \in \Gamma(S, \mathcal{O}_S)$，使得
$s \in D(a) \subset f_0^{-1}(U_0)$，并且 $D(a)$ 仿射。
由引理 \ref{limits-lemma-descend-section}，存在 $i \geq 0$，使得 $a$ 来自某个元素
$a_i \in \Gamma(S_i, \mathcal{O}_{S_i})$。对任意指标 $j \geq i$，
记 $a_j$ 为 $a_i$ 在 $S_j$ 的结构层整体截面中的像。考虑开集
$D(a_j) \subset S_j$ 和 $U_j = f_{j0}^{-1}(U_0)$。注意，$U_j$ 仿射，
而 $D(a_j)$ 是 $S_j$ 的拟紧开集；例如见《性质》引理
\ref{properties-lemma-affine-cap-s-open}。因此可将引理
\ref{limits-lemma-descend-opens} 应用于开集 $U_j$ 和 $U_j \cup D(a_j)$，
从而得到 $D(a_j) \subset U_j$ 对某个 $j \geq i$ 成立。

对这样的指标 $j$，$D(a_j) \subset S_j$ 是仿射开集
（因为 $D(a_j)$ 是仿射开集 $U_j$ 的标准仿射开集），并且它包含像 $f_j(s)$。

\medskip\noindent
综上，对每个 $s \in S$，都存在指标 $i \in I$ 以及整体截面
$a \in \Gamma(S_i, \mathcal{O}_{S_i})$，使得 $D(a) \subset S_i$
是一个包含 $f_i(s)$ 的仿射开集。由于 $S$ 拟紧，可以选取单个指标
$i \in I$ 以及整体截面
$a_1, \ldots, a_m \in \Gamma(S_i, \mathcal{O}_{S_i})$，使每个
$D(a_j) \subset S_i$ 都是仿射开集，并且 $f_i : S \to S_i$ 的像包含在并集
$W_i = \bigcup_{j = 1, \ldots, m} D(a_j)$ 中。对 $i' \geq i$，置
$W_{i'} = f_{i'i}^{-1}(W_i)$。因为 $f_i^{-1}(W_i)$ 是整个 $S$，再次应用引理
\ref{limits-lemma-descend-opens} 可知，对适当的 $i' \geq i$ 有
$S_{i'} = W_{i'}$。因此可将 $i$ 替换为 $i'$，并假设
$S_i = \bigcup_{j = 1, \ldots, m} D(a_j)$。这说明
$\mathcal{O}_{S_i}$ 是 $S_i$ 上的丰沛可逆层（见《性质》定义
\ref{properties-definition-ample}），因而 $S_i$ 拟仿射，见《性质》引理
\ref{properties-lemma-quasi-affine-O-ample}。结论成立。
\end{proof}

\begin{lemma}
\label{limits-lemma-limit-affine}
在情形 \ref{limits-situation-descent} 中，若 $S$ 仿射，则对某个 $i_0 \in I$，
概形 $S_i$ 在 $i \geq i_0$ 时都仿射。
\end{lemma}

\begin{proof}
由引理 \ref{limits-lemma-limit-quasi-affine}，可假设 $S_0$ 对某个 $0 \in I$ 拟仿射。
置 $R_0 = \Gamma(S_0, \mathcal{O}_{S_0})$。则 $S_0$ 是
$T_0 = \Spec(R_0)$ 的一个拟紧开集。记
$j_0 : S_0 \to T_0$ 为相应的拟紧开浸入。对 $i \geq 0$，置
$\mathcal{A}_i = f_{i0, *}\mathcal{O}_{S_i}$。因为 $f_{i0}$ 仿射，所以
$S_i = \underline{\Spec}_{S_0}(\mathcal{A}_i)$。置
$T_i = \underline{\Spec}_{T_0}(j_{0, *}\mathcal{A}_i)$。
则 $T_i \to T_0$ 仿射，因而 $T_i$ 仿射。因此 $T_i$ 是以下环的谱：
$$
R_i = \Gamma(T_0, j_{0, *}\mathcal{A}_i) = \Gamma(S_0, \mathcal{A}_i) =
\Gamma(S_i, \mathcal{O}_{S_i}).
$$
写 $S = \Spec(R)$。由引理 \ref{limits-lemma-descend-section} 有
$R = \colim_i R_i$。因此也有 $S = \lim_i T_i$。由于相对谱的形成与基变换交换，
开集 $S_0 \subset T_0$ 在 $T_i$ 中的逆像是 $S_i$。令
$Z_0 = T_0 \setminus S_0$，并令 $Z_i \subset T_i$ 为 $Z_0$ 的逆像。
由于 $S_i = T_i \setminus Z_i$，只需证明 $Z_i$ 对某个 $i$ 为空。
反设 $Z_i$ 对所有 $i$ 都非空。由引理 \ref{limits-lemma-limit-closed-nonempty}，存在一点
$s$ 属于 $S = \lim T_i$，它映到 $Z_i$ 中的一点（对每个 $i$）。
但 $S = \lim_i S_i$，故由引理 \ref{limits-lemma-topology-limit} 得到矛盾。
\end{proof}

\begin{lemma}
\label{limits-lemma-limit-separated}
在情形 \ref{limits-situation-descent} 中，若 $S$ 分离，则对某个 $i_0 \in I$，
概形 $S_i$ 在 $i \geq i_0$ 时都分离。
\end{lemma}

\begin{proof}
选取有限仿射开覆盖
$S_0 = U_{0, 1} \cup \ldots \cup U_{0, m}$。令
$U_{i, j} \subset S_i$ 和 $U_j \subset S$ 分别为 $U_{0, j}$ 的逆像。
注意，$U_{i, j}$ 和 $U_j$ 都仿射。由于 $S$ 分离，交集
$U_{j_1} \cap U_{j_2}$ 仿射。因为
$U_{j_1} \cap U_{j_2} = \lim_{i \geq 0} U_{i, j_1} \cap U_{i, j_2}$，
由引理 \ref{limits-lemma-limit-affine} 可知，$U_{i, j_1} \cap U_{i, j_2}$
在 $i$ 充分大时仿射。为了证明 $S_i$ 在 $i$ 充分大时分离，
现在只需证明
$$
\mathcal{O}_{S_i}(U_{i, j_1})
\otimes_{\mathcal{O}_S(S)}
\mathcal{O}_{S_i}(U_{i, j_2})
\longrightarrow
\mathcal{O}_{S_i}(U_{i, j_1} \cap U_{i, j_2})
$$
当 $i$ 充分大时满射（《概形》引理
\ref{schemes-lemma-characterize-separated}）。

\medskip\noindent
为了去掉烦琐的指标，设有仿射开集 $U, V \subset S_0$，且 $U \cap V$ 也仿射。
令 $U_i, V_i \subset S_i$ 以及 $U, V \subset S$ 分别为它们的逆像。
需要证明 $\mathcal{O}(U_i) \otimes \mathcal{O}(V_i) \to
\mathcal{O}(U_i \cap V_i)$ 在 $i$ 充分大时满射；已知
$\mathcal{O}(U) \otimes \mathcal{O}(V) \to \mathcal{O}(U \cap V)$ 满射。
注意，
$\mathcal{O}(U_0) \otimes \mathcal{O}(V_0) \to
\mathcal{O}(U_0 \cap V_0)$ 是有限型的，因为对角态射
$S_i \to S_i \times S_i$ 是浸入（《概形》引理
\ref{schemes-lemma-diagonal-immersion}），因而是局部有限型的
（《态射》引理 \ref{morphisms-lemma-locally-finite-type-characterize} 和
\ref{morphisms-lemma-immersion-locally-finite-type}）。
因此可以选取元素
$f_{0, 1}, \ldots, f_{0, n} \in \mathcal{O}(U_0 \cap V_0)$，
它们生成 $\mathcal{O}(U_0 \cap V_0)$ 作为
$\mathcal{O}(U_0) \otimes \mathcal{O}(V_0)$-代数。注意，对 $i \geq 0$，概形图表
$$
\xymatrix{
U_i \cap V_i \ar[r] \ar[d] & U_i \ar[d] \\
U_0 \cap V_0 \ar[r] & U_0
}
$$
是笛卡尔的。因此，像
$f_{i, 1}, \ldots, f_{i, n} \in \mathcal{O}(U_i \cap V_i)$
生成 $\mathcal{O}(U_i \cap V_i)$ 作为
$\mathcal{O}(U_i) \otimes \mathcal{O}(V_0)$-代数，从而更是作为
$\mathcal{O}(U_i) \otimes \mathcal{O}(V_i)$-代数生成它。
由假设，像 $f_1, \ldots, f_n \in \mathcal{O}(U \otimes V)$ 属于映射
$\mathcal{O}(U) \otimes \mathcal{O}(V) \to \mathcal{O}(U \cap V)$ 的像。
由于
$\mathcal{O}(U) \otimes \mathcal{O}(V) =
\colim \mathcal{O}(U_i) \otimes \mathcal{O}(V_i)$，
它们在某个有限层已经属于相应映射的像，引理得证。
\end{proof}

\begin{lemma}
\label{limits-lemma-limit-ample}
在情形 \ref{limits-situation-descent} 中，设 $\mathcal{L}_0$ 是 $S_0$ 上的可逆模层。
若拉回 $\mathcal{L}$ 到 $S$ 是丰沛的，则对某个 $i \in I$，
拉回 $\mathcal{L}_i$ 到 $S_i$ 是丰沛的。
\end{lemma}

\begin{proof}
该假设意味着存在有限多个截面
$s_1, \ldots, s_m \in \Gamma(S, \mathcal{L})$，使得每个 $S_{s_j}$ 仿射，
并且 $S = \bigcup S_{s_j}$，见《性质》定义
\ref{properties-definition-ample}。由引理 \ref{limits-lemma-descend-section}，
可以找到 $i \in I$ 以及截面
$s_{i, j} \in \Gamma(S_i, \mathcal{L}_i)$，它们映到 $s_j$。
由引理 \ref{limits-lemma-limit-affine}，增大 $i$ 后可假设
$(S_i)_{s_{i, j}}$ 对 $j = 1, \ldots, m$ 仿射。最后再次增大 $i$，
由引理 \ref{limits-lemma-descend-opens} 可假设
$S_i = \bigcup (S_i)_{s_{i, j}}$。于是按定义，$\mathcal{L}_i$ 是丰沛的。
\end{proof}

\begin{lemma}
\label{limits-lemma-finite-type-eventually-closed}
设 $S$ 是概形。设 $X = \lim X_i$ 是 $S$ 上概形的有向极限，且过渡态射仿射。
设 $Y \to X$ 是 $S$ 上概形的态射。
\begin{enumerate}
\item 若 $Y \to X$ 是闭浸入，$X_i$ 拟紧，并且 $Y$ 在 $S$ 上局部有限型，
则 $Y \to X_i$ 当 $i$ 充分大时是闭浸入。
\item 若 $Y \to X$ 是浸入，$X_i$ 拟分离，$Y \to S$ 局部有限型，
并且 $Y$ 拟紧，则 $Y \to X_i$ 当 $i$ 充分大时是浸入。
\item 若 $Y \to X$ 是同构，$X_i$ 拟紧，$X_i \to S$ 局部有限型，
过渡态射 $X_{i'} \to X_i$ 是闭浸入，并且 $Y \to S$ 局部有限表示，
则 $Y \to X_i$ 当 $i$ 充分大时是同构。
\end{enumerate}
\end{lemma}

\begin{proof}
证明 (1)。选取 $0 \in I$，并选取有限仿射开覆盖
$X_0 = U_{0, 1} \cup \ldots \cup U_{0, m}$，使每个 $U_{0, j}$ 都映到
某个仿射开集 $W_j \subset S$ 中。令
$V_j \subset Y$、$U_{i, j} \subset X_i$（$i \geq 0$）以及
$U_j \subset X$ 分别为 $U_{0, j}$ 的逆像。只需证明
$V_j \to U_{i, j}$ 当 $i$ 充分大时是闭浸入，而我们已知 $V_j \to U_j$ 是闭浸入。
因此问题约化为以下代数事实：若 $A = \colim A_i$ 是 $R$-代数的有向余极限，
$A \to B$ 是 $R$-代数的满射，并且 $B$ 是有限生成 $R$-代数，
则 $A_i \to B$ 当 $i$ 充分大时满射。

\medskip\noindent
证明 (2)。选取 $0 \in I$。选取一个拟紧开集 $X'_0 \subset X_0$，使得
$Y \to X_0$ 分解通过 $X'_0$。把 $X_i$ 替换为 $X'_0$ 的逆像
（$i \geq 0$）后，可以假设所有 $X_i'$ 都拟紧且拟分离。
取拟紧开集 $U \subset X$，使得 $Y \to X$ 分解通过闭浸入 $Y \to U$
（这样的 $U$ 存在，因为 $Y$ 拟紧）。
由引理 \ref{limits-lemma-descend-opens}，可假设
$U = \lim U_i$，其中 $U_i \subset X_i$ 是拟紧开集。由第 (1) 部分，
$Y \to U_i$ 对某个 $i$ 是闭浸入。因此 (2) 成立。

\medskip\noindent
证明 (3)。像 (1) 的证明那样，在 $X_0$ 上对某个 $0 \in I$ 仿射局部地工作，
便将问题约化为以下代数事实：若 $A = \lim A_i$ 是 $R$-代数的有向余极限，
过渡映射均满射，并且 $A$ 在 $A_0$ 上有限表示，则有 $A = A_i$ 对某个 $i$ 成立。
事实上，写 $A = A_0/(f_1, \ldots, f_n)$。选取 $i$，使得
$f_1, \ldots, f_n$ 在满射 $A_0 \to A_i$ 下都映到零。
\end{proof}

\begin{lemma}
\label{limits-lemma-eventually-separated}
设 $S$ 是概形。设 $X = \lim X_i$ 是 $S$ 上概形的有向极限，且过渡态射仿射。
假设
\begin{enumerate}
\item $S$ 拟分离；
\item $X_i$ 拟紧且拟分离；
\item $X \to S$ 分离。
\end{enumerate}
则 $X_i \to S$ 对所有充分大的 $i$ 都分离。
\end{lemma}

\begin{proof}
取 $0 \in I$。注意，因为极限是有向的，所以 $I$ 非空。由于 $X_0$ 拟紧，
可以找到有限多个仿射开集 $U_1, \ldots, U_n \subset S$，使得
$X_0 \to S$ 的像包含在 $U_1 \cup \ldots \cup U_n$ 中。记
$h_i : X_i \to S$ 为结构态射。只需检验对某个 $i \geq 0$，态射
$h_i^{-1}(U_j) \to U_j$ 对 $j = 1, \ldots,  n$ 都分离。
由于 $S$ 拟分离，态射 $U_j \to S$ 拟紧。因此 $h_i^{-1}(U_j)$ 拟紧且拟分离。
这样便约化到 $S$ 仿射的情形。在此情形下，需要证明 $X_i$ 分离，而已知 $X$ 分离。
因此引理由引理 \ref{limits-lemma-limit-separated} 得出。
\end{proof}

\begin{lemma}
\label{limits-lemma-eventually-affine}
设 $S$ 是概形。设 $X = \lim X_i$ 是 $S$ 上概形的有向极限，且过渡态射仿射。
假设
\begin{enumerate}
\item $S$ 拟紧且拟分离；
\item $X_i$ 拟紧且拟分离；
\item $X \to S$ 仿射。
\end{enumerate}
则 $X_i \to S$ 当 $i$ 充分大时仿射。
\end{lemma}

\begin{proof}
选取有限仿射开覆盖 $S = \bigcup_{j = 1, \ldots, n} V_j$。记
$f : X \to S$ 和 $f_i : X_i \to S$ 为结构态射。对每个 $j$，概形
$f^{-1}(V_j) = \lim_i f_i^{-1}(V_j)$ 是仿射的（因为有限态射按定义是仿射的）。
因此由引理 \ref{limits-lemma-limit-affine}，存在 $i \in I$，使每个
$f_i^{-1}(V_j)$ 都仿射。换言之，$f_i : X_i \to S$ 当 $i$ 充分大时仿射，
见《态射》引理 \ref{morphisms-lemma-characterize-affine}。
\end{proof}

\begin{lemma}
\label{limits-lemma-eventually-finite}
设 $S$ 是概形。设 $X = \lim X_i$ 是 $S$ 上概形的有向极限，且过渡态射仿射。
假设
\begin{enumerate}
\item $S$ 拟紧且拟分离；
\item $X_i$ 拟紧且拟分离；
\item 过渡态射 $X_{i'} \to X_i$ 有限；
\item $X_i \to S$ 局部有限型；
\item $X \to S$ 整。
\end{enumerate}
则 $X_i \to S$ 当 $i$ 充分大时有限。
\end{lemma}

\begin{proof}
由引理 \ref{limits-lemma-eventually-affine}，可假设 $X_i \to S$ 对所有 $i$ 都仿射。
选取有限仿射开覆盖 $S = \bigcup_{j = 1, \ldots, n} V_j$。记
$f : X \to S$ 和 $f_i : X_i \to S$ 为结构态射。只需证明存在 $i$，使得
$f_i^{-1}(V_j)$ 在 $V_j$ 上有限（$j = 1, \ldots, m$）
（《态射》引理 \ref{morphisms-lemma-finite-local}）。
事实上，对 $i' \geq i$，复合 $X_{i'} \to X_i \to S$ 是有限态射的复合，
因而有限（《态射》引理 \ref{morphisms-lemma-composition-finite}）。
这把问题约化到仿射情形：设 $R$ 是环，$A = \colim A_i$，其中
$R \to A$ 整，并且 $A_i \to A_{i'}$ 对所有 $i \leq i'$ 都有限。
此外，$R \to A_i$ 对所有 $i$ 都有限型。目标是证明 $A_i$ 在 $R$ 上
对某个 $i$ 有限。为此，选取 $i \in I$，并选取元素
$x_1, \ldots, x_m \in A_i$，它们生成 $A_i$ 作为 $R$-代数。
因为 $A$ 在 $R$ 上整，可以找到首一多项式 $P_j \in R[T]$，满足
$P_j(x_j) = 0$ 于 $A$ 中。因此存在 $i' \geq i$，使得
$P_j(x_j) = 0$ 在 $A_{i'}$ 中对 $j = 1, \ldots, m$ 都成立。
于是由《代数》引理
\ref{algebra-lemma-characterize-finite-in-terms-of-integral}，像 $A'_i$
（即 $A_i$ 在 $A_{i'}$ 中的像）在 $R$ 上有限。又因为
$A'_i \subset A_{i'}$ 也有限，
由《代数》引理 \ref{algebra-lemma-finite-transitive} 可知 $A_{i'}$ 在 $R$ 上有限。
\end{proof}

\begin{lemma}
\label{limits-lemma-eventually-closed-immersion}
设 $S$ 是概形。设 $X = \lim X_i$ 是 $S$ 上概形的有向极限，且过渡态射仿射。
假设
\begin{enumerate}
\item $S$ 拟紧且拟分离；
\item $X_i$ 拟紧且拟分离；
\item 过渡态射 $X_{i'} \to X_i$ 是闭浸入；
\item $X_i \to S$ 局部有限型；
\item $X \to S$ 是闭浸入。
\end{enumerate}
则 $X_i \to S$ 当 $i$ 充分大时是闭浸入。
\end{lemma}

\begin{proof}
由引理 \ref{limits-lemma-eventually-affine}，可假设 $X_i \to S$ 对所有 $i$ 都仿射。
选取有限仿射开覆盖 $S = \bigcup_{j = 1, \ldots, n} V_j$。记
$f : X \to S$ 和 $f_i : X_i \to S$ 为结构态射。只需证明存在 $i$，使得
$f_i^{-1}(V_j)$ 是 $V_j$ 的闭子概形（$j = 1, \ldots, m$）
（《态射》引理 \ref{morphisms-lemma-closed-immersion}）。
这把问题约化到仿射情形：设 $R$ 是环，$A = \colim A_i$，其中
$R \to A$ 满射，并且 $A_i \to A_{i'}$ 对所有 $i \leq i'$ 都满射。
此外，$R \to A_i$ 对所有 $i$ 都有限型。目标是证明 $R \to A_i$
对某个 $i$ 满射。为此，选取 $i \in I$，并选取元素
$x_1, \ldots, x_m \in A_i$，它们生成 $A_i$ 作为 $R$-代数。
由于 $R \to A$ 满射，可以找到 $r_j \in R$，使得 $r_j$ 映到 $x_j$
于 $A$ 中。因此存在 $i' \geq i$，使得 $r_j$ 映到 $x_j$ 的像于
$A_{i'}$ 中（$j = 1, \ldots, m$）。
由于 $A_i \to A_{i'}$ 满射，这说明 $R \to A_{i'}$ 满射。
\end{proof}

\begin{lemma}
\label{limits-lemma-eventually-immersion}
设 $S$ 是概形。设 $X = \lim X_i$ 是 $S$ 上概形的有向极限，且过渡态射仿射。
假设
\begin{enumerate}
\item $S$ 拟分离；
\item $X_i$ 拟紧且拟分离；
\item 过渡态射 $X_{i'} \to X_i$ 是闭浸入；
\item $X_i \to S$ 局部有限型；
\item $X \to S$ 是浸入。
\end{enumerate}
则 $X_i \to S$ 当 $i$ 充分大时是浸入。
\end{lemma}

\begin{proof}
选取开子概形 $U \subset S$，使得 $X \to S$ 分解为闭浸入
$X \to U$ 与包含态射 $U \to S$ 的复合。由于 $X$ 拟紧，可以缩小 $U$
并假设 $U$ 拟紧。记 $V_i \subset X_i$ 为 $U$ 的逆像。由于 $V_i$ 拉回为
$X$，由引理 \ref{limits-lemma-descend-opens} 可知，$V_i = X_i$ 对所有充分大的 $i$ 成立。
因此可以假设 $X = \lim X_i$ 是 $U$ 上概形范畴中的极限。
于是由引理 \ref{limits-lemma-eventually-closed-immersion} 可知，$X_i \to U$
当 $i$ 充分大时是闭浸入。这就证明了引理。
\end{proof}


\section{绝对 Noether 逼近}
\label{limits-section-approximation}

\noindent
本节的一个很好的参考文献是 Thomason 与 Trobaugh 的论文
\cite{TT} 的附录 C。
关于有向系的约定，见《范畴》第 \ref{categories-section-posets-limits} 节。
下文将直接使用第 \ref{limits-section-limits} 节给出的极限存在性结果及其性质，
不再另行说明。

\begin{lemma}
\label{limits-lemma-quasi-affine-finite-type-over-Z}
设 $W$ 是 $\mathbf{Z}$ 上有限型的拟仿射概形。假设
$W \to \Spec(R)$ 是到一个仿射概形的开浸入。则存在有限型
$\mathbf{Z}$-子代数 $A \subset R$，使其诱导开浸入
$W \to \Spec(A)$。此外，$R$ 是所有这类子代数的有向余极限。
\end{lemma}

\begin{proof}
选取仿射开覆盖 $W = \bigcup_{i = 1, \ldots, n} W_i$，使每个 $W_i$
都是 $\Spec(R)$ 中的标准仿射开集。换言之，若写
$W_i = \Spec(R_i)$，则 $R_i = R_{f_i}$，其中 $f_i \in R$。
选取有限多个 $x_{ij} \in R_i$，使它们生成 $R_i$ 作为 $\mathbf{Z}$-代数。
取 $N \gg 0$，使每个 $f_i^Nx_{ij}$ 都来自 $R$ 的某个元素，
记为 $y_{ij} \in R$。令 $A$ 等于如下 $\mathbf{Z}$-代数：
它由各 $f_i$、各 $y_{ij}$ 以及（可选的）$R$ 中另外有限多个元素生成。
则 $A$ 满足要求。细节从略。
\end{proof}

\begin{lemma}
\label{limits-lemma-diagram}
设给定如下环的笛卡尔图
$$
\xymatrix{
B \ar[r]_s & R \\
B'\ar[u] \ar[r] & R' \ar[u]_t
}
$$
设 $W' \subset \Spec(R')$ 为形如
$W' = D(f_1) \cup \ldots \cup D(f_n)$ 的开集，并且
有 $t(f_i) = s(g_i)$（其中 $g_i \in B$）及
$B_{g_i} \cong R_{s(g_i)}$。则 $B' \to R'$
诱导从 $W'$ 到 $\Spec(B')$ 的开浸入。
\end{lemma}

\begin{proof}
令 $h_i = (g_i, f_i) \in B'$。《高等代数》引理
\ref{more-algebra-lemma-diagram-localize} 表明
$(B')_{h_i} \cong (R')_{f_i}$，正合所需。
\end{proof}

\noindent
下面的引理给出 Noether 逼近的一个精确表述。

\begin{lemma}
\label{limits-lemma-approximate}
设 $S$ 是拟紧且拟分离的概形，$V \subset S$ 是拟紧开集。设 $I$ 是有向集，
并设 $(V_i, f_{ii'})$ 是以 $I$ 为指标、过渡映射仿射的概形逆系，
每个 $V_i$ 都是 $\mathbf{Z}$ 上有限型的，并且 $V = \lim V_i$。
则存在
\begin{enumerate}
\item 一个有向集 $J$；
\item 一个概形逆系 $(S_j, g_{jj'})$，以 $J$ 为指标；
\item 一个保序映射 $\alpha : J \to I$；
\item 开子概形 $V'_j \subset S_j$；以及
\item 同构 $V'_j \to V_{\alpha(j)}$，
\end{enumerate}
使得
\begin{enumerate}
\item 过渡态射 $g_{jj'} : S_j \to S_{j'}$ 仿射；
\item 每个 $S_j$ 都是 $\mathbf{Z}$ 上有限型的；
\item $g_{jj'}^{-1}(V'_{j'}) = V'_j$；
\item $S = \lim S_j$ 且 $V = \lim V'_j$；并且
\item 图
$$
\vcenter{
\xymatrix{
V \ar[d] \ar[rd] \\
V'_j \ar[r] & V_{\alpha(j)}
}
}
\quad\text{且}\quad
\vcenter{
\xymatrix{
V'_j \ar[r] \ar[d] & V_{\alpha(j)} \ar[d] \\
V'_{j'} \ar[r] & V_{\alpha(j')}
}
}
$$
交换。
\end{enumerate}
\end{lemma}

\begin{proof}
令 $Z = S \setminus V$。选取仿射开集 $U_1, \ldots, U_m \subset S$，
使得 $Z \subset \bigcup_{l = 1, \ldots, m} U_l$。考察开集
$$
V \subset V \cup U_1 \subset V \cup U_1 \cup U_2 \subset
\ldots \subset V \cup \bigcup\nolimits_{l = 1, \ldots, m} U_l = S
$$
如果能对每一种情形
$$
V \cup U_1 \cup \ldots \cup U_l
\subset
V \cup U_1 \cup \ldots \cup U_{l + 1}
$$
依次证明此引理，那么 $V \subset S$ 的结论随即成立。每一步只添加一个仿射开集。
因此可以假设
\begin{enumerate}
\item $S = U \cup V$；
\item $U$ 是 $S$ 的仿射开集；
\item $V$ 是 $S$ 的拟紧开集；并且
\item $V = \lim_i V_i$，其中 $(V_i, f_{ii'})$ 是以有向集 $I$ 为指标的逆系，
每个 $f_{ii'}$ 仿射，且每个 $V_i$ 都是 $\mathbf{Z}$ 上有限型的。
\end{enumerate}
记 $f_i : V \to V_i$ 为各投影。令 $W = U \cap V$。由于 $S$ 拟分离，
这是 $V$ 的拟紧开集。由引理 \ref{limits-lemma-descend-opens}
（并在缩小 $I$ 后），可以假设存在开集 $W_i \subset V_i$，使得
$f_{ii'}^{-1}(W_{i'}) = W_i$ 且 $f_i^{-1}(W_i) = W$。由于 $W$
是 $U$ 的拟紧开集，它是拟仿射的。因此可以再次缩小 $I$，并假设
$W_i$ 对所有 $i$ 都拟仿射；见引理 \ref{limits-lemma-limit-quasi-affine}。

\medskip\noindent
写成 $U = \Spec(B)$。令 $R = \Gamma(W, \mathcal{O}_W)$，
并令 $R_i = \Gamma(W_i, \mathcal{O}_{W_i})$。由引理
\ref{limits-lemma-descend-section}，有 $R = \colim_i R_i$。
现在得到环映射
$$
\xymatrix{
B \ar[r]_s & R \\
& R_i \ar[u]_{t_i}
}
$$
令 $B_i = \{(b, r) \in B \times R_i \mid s(b) = t_i(r)\}$，于是对每个 $i$
都有笛卡尔图
$$
\xymatrix{
B \ar[r]_s & R \\
B_i \ar[u] \ar[r] & R_i \ar[u]_{t_i}
}
$$
过渡映射 $R_i \to R_{i'}$ 诱导映射 $B_i \to B_{i'}$。显然
$B = \colim_i B_i$。下一段将证明，对所有充分大的 $i$，复合
$W_i \to \Spec(R_i) \to \Spec(B_i)$ 是开浸入。

\medskip\noindent
由于 $W$ 是 $U = \Spec(B)$ 的拟紧开集，可以找到有限多个元素
$g_l \in B$（$l = 1, \ldots, m$），使得 $D(g_l) \subset W$ 且
$W = \bigcup_{l = 1, \ldots, m} D(g_l)$。注意，这蕴含
$D(g_l) = W_{s(g_l)}$（作为 $U$ 的开子集），其中 $W_{s(g_l)}$
表示 $W$ 中使 $s(g_l)$ 可逆的最大开子集。因此
$$
B_{g_l} =
\Gamma(D(g_l), \mathcal{O}_U) =
\Gamma(W_{s(g_l)}, \mathcal{O}_W) = R_{s(g_l)},
$$
这里最后一个等式来自《性质》引理
\ref{properties-lemma-invert-f-sections}。
由于 $W_{s(g_l)}$ 仿射，这还蕴含
$D(s(g_l)) = W_{s(g_l)}$（作为 $\Spec(R)$ 的开子集）。
由于 $R = \colim_i R_i$，可以（缩小 $I$ 后）假设存在
$g_{l, i} \in R_i$（对所有 $i \in I$），使得
$s(g_l) = t_i(g_{l, i})$。当然，我们选取 $g_{l, i}$，使得
$g_{l, i}$ 映到 $g_{l, i'}$，所经的过渡映射为
$R_i \to R_{i'}$。于是由引理
\ref{limits-lemma-descend-opens}，可以（再次缩小 $I$ 后）假设相应开集
$D(g_{l, i}) \subset \Spec(R_i)$ 包含在 $W_i$ 中（$l = 1, \ldots, m$），
并且它们覆盖 $W_i$。由引理 \ref{limits-lemma-diagram} 可知，态射
$W_i \to \Spec(R_i) \to \Spec(B_i)$ 是开浸入。

\medskip\noindent
由引理 \ref{limits-lemma-quasi-affine-finite-type-over-Z}，可将 $B_i$
写成子代数 $A_{i, p} \subset B_i$（$p \in P_i$）的有向余极限，
其中每个子代数都是 $\mathbf{Z}$ 上有限型的，并且 $W_i$
被等同于 $\Spec(A_{i, p})$ 的一个开子概形。令 $S_{i, p}$
为将 $V_i$ 与 $\Spec(A_{i, p})$ 沿开集 $W_i$ 黏合所得的概形；
见《概形》第 \ref{schemes-section-glueing-schemes} 节。
由此得到如下概形交换图：
$$
\xymatrix{
& & V \ar[lld] \ar[d] & W \ar[l] \ar[lld] \ar[d] \\
V_i \ar[d] & W_i \ar[l] \ar[d] & S \ar[lld] & U \ar[lld] \ar[l] \\
S_{i, p} & \Spec(A_{i, p}) \ar[l]
}
$$
态射 $S \to S_{i, p}$ 存在，是因为右上方的方块在概形范畴中是推出。
注意，$S_{i, p}$ 是 $\mathbf{Z}$ 上有限型的，因为它有一个有限仿射开覆盖，
其各成员都是有限型 $\mathbf{Z}$-代数的谱。
在 $J = \coprod_{i \in I} P_i$ 上定义预序如下：规定
$(i', p') \geq (i, p)$ 当且仅当 $i' \geq i$ 且映射
$B_i \to B_{i'}$ 把 $A_{i, p}$ 映入 $A_{i', p'}$。
这恰好是定义态射
$S_{i', p'} \to S_{i, p}$ 所需的条件：即利用过渡态射 $V_{i'} \to V_i$、
$W_{i'} \to W_i$ 以及态射
$\Spec(A_{i', p'}) \to \Spec(A_{i, p})$（它由环映射
$A_{i, p} \to A_{i', p'}$ 诱导），构造如上的交换图。
相关的交换性已经包含在这些构造中。
我们断言 $S$ 是概形 $S_{i, p}$ 的有向极限。由于概形 $V_i$
按构造以 $V$ 为极限，这归结为 $B$ 是环 $A_{i, p}$ 的极限；
此事由构造即知。映射 $\alpha : J \to I$ 由规则
$j = (i, p) \mapsto i$ 给出。开子概形 $V'_j$ 就是上面
$V_i \to S_{i, p}$ 的像。第 (5) 项各图的交换性由构造显然可知。
引理得证。
\end{proof}

\begin{proposition}
\label{limits-proposition-approximate}
设 $S$ 是拟紧且拟分离的概形。则存在有向集 $I$ 以及概形逆系
$(S_i, f_{ii'})$，后者以 $I$ 为指标，并且满足
\begin{enumerate}
\item 过渡态射 $f_{ii'}$ 仿射；
\item 每个 $S_i$ 都是 $\mathbf{Z}$ 上有限型的；并且
\item $S = \lim_i S_i$。
\end{enumerate}
\end{proposition}

\begin{proof}
这是引理 \ref{limits-lemma-approximate} 在 $V = \emptyset$ 时的特例。
\end{proof}


\section{极限与有限呈示态射}
\label{limits-section-finite-presentation}

\noindent
下面是《代数》引理
\ref{algebra-lemma-characterize-finite-presentation} 的推广。

\begin{proposition}
\label{limits-proposition-characterize-locally-finite-presentation}
\begin{reference}
\cite[IV, Proposition 8.14.2]{EGA}
\end{reference}
设 $f : X \to S$ 是概形态射。下列条件等价：
\begin{enumerate}
\item 态射 $f$ 局部有限呈示。
\item 对任意有向集 $I$，以及任意逆系 $(T_i, f_{ii'})$，
若它由 $S$-概形组成、以 $I$ 为指标，并且每个 $T_i$ 都仿射，则有
$$
\Mor_S(\lim_i T_i, X) =
\colim_i \Mor_S(T_i, X)
$$
\item 对任意有向集 $I$，以及任意逆系 $(T_i, f_{ii'})$，
若它由 $S$-概形组成、以 $I$ 为指标，每个 $f_{ii'}$ 都仿射，
并且每个 $T_i$ 作为概形均拟紧且拟分离，则有
$$
\Mor_S(\lim_i T_i, X) =
\colim_i \Mor_S(T_i, X)
$$
\end{enumerate}
\end{proposition}

\begin{proof}
(3) 蕴含 (2) 是显然的。

\medskip\noindent
下面证明 (2) 蕴含 (1)。假设 (2) 成立。任取仿射开集
$U \subset X$ 与 $V \subset S$，使得 $f(U) \subset V$。
必须证明 $\mathcal{O}_S(V) \to \mathcal{O}_X(U)$ 有限呈示。
设 $(A_i, \varphi_{ii'})$ 是 $\mathcal{O}_S(V)$-代数的有向系。
令 $A = \colim_i A_i$。根据《代数》引理
\ref{algebra-lemma-characterize-finite-presentation}，必须证明
$$
\Hom_{\mathcal{O}_S(V)}(\mathcal{O}_X(U), A) =
\colim_i \Hom_{\mathcal{O}_S(V)}(\mathcal{O}_X(U), A_i)
$$
考察概形 $T_i = \Spec(A_i)$。它们组成 $V$-概形逆系，以 $I$ 为指标，
其过渡态射 $f_{ii'} : T_i \to T_{i'}$ 由
$\mathcal{O}_S(V)$-代数映射 $\varphi_{i'i}$ 诱导。令
$T := \Spec(A) = \lim_i T_i$。用概形态射集表示，上式化为
$$
\Mor_V(\lim_i T_i, U) =
\colim_i \Mor_V(T_i, U).
$$
首先注意到
$\Mor_V(T_i, U) = \Mor_S(T_i, U)$
以及
$\Mor_V(T, U) = \Mor_S(T, U)$。
因此必须证明
$$
\Mor_S(\lim_i T_i, U) =
\colim_i \Mor_S(T_i, U)
$$
而已知
$$
\Mor_S(\lim_i T_i, X) =
\colim_i \Mor_S(T_i, X).
$$
所以只需证明：给定态射 $g_i : T_i \to X$（在 $S$ 上），若复合
$T \to T_i \to X$ 的像落在 $U$ 中，则存在某个 $i' \geq i$，
使复合 $g_{i'} : T_{i'} \to T_i \to X$ 的像落在 $U$ 中。
记 $Z_{i'} = g_{i'}^{-1}(X \setminus U)$。假设每个 $Z_{i'}$ 都非空，
以导出矛盾。由引理 \ref{limits-lemma-limit-closed-nonempty}，存在一点 $t$，
它属于 $T$，并且映入 $Z_{i'}$（对所有 $i' \geq i$）。
这样的点不映入 $U$，矛盾。

\medskip\noindent
最后证明 (1) 蕴含 (3)。假设 (1) 成立。给定有向逆系
$(T_i, f_{ii'})$，其对象是 $S$-概形。假设态射 $f_{ii'}$ 仿射，并且每个 $T_i$
作为概形均拟紧且拟分离。令 $T = \lim_i T_i$。记
$f_i : T \to T_i$ 为投影态射。需要证明：
\begin{enumerate}
\item[(a)] 给定态射 $g_i, g'_i : T_i \to X$（在 $S$ 上），如果
$g_i \circ f_i = g'_i \circ f_i$，则存在 $i' \geq i$，使得
$g_i \circ f_{i'i} = g'_i \circ f_{i'i}$。
\item[(b)] 给定任意态射 $g : T \to X$（在 $S$ 上），存在 $i \in I$
及态射 $g_i : T_i \to X$，使得 $g = f_i \circ g_i$。
\end{enumerate}

\noindent
先证明唯一性部分 (a)。设 $g_i, g'_i : T_i \to X$ 是满足
$g_i \circ f_i = g'_i \circ f_i$ 的态射。对任意 $i' \geq i$，
令 $g_{i'} = g_i \circ f_{i'i}$ 且
$g'_{i'} = g'_i \circ f_{i'i}$。还令
$g = g_i \circ f_i = g'_i \circ f_i$。
考察态射
$(g_i, g'_i) : T_i \to X \times_S X$。令
$$
W =
\bigcup\nolimits_{U \subset X\text{ affine open},
V \subset S\text{ affine open}, f(U) \subset V}
U \times_V U.
$$
这是 $X \times_S X$ 的一个开集，且态射 $\Delta_{X/S}$ 可通过到 $W$
的闭浸入分解；见《概形》引理
\ref{schemes-lemma-diagonal-immersion} 的证明。
注意，复合 $(g_i, g'_i) \circ f_i : T \to X \times_S X$ 是到 $W$
的态射，因为按假设它通过对角态射分解。
令 $Z_{i'} = (g_{i'}, g'_{i'})^{-1}(X \times_S X \setminus W)$。
如果每个 $Z_{i'}$ 都非空，则由引理
\ref{limits-lemma-limit-closed-nonempty}，存在点 $t \in T$，使其映入
$Z_{i'}$（对所有 $i' \geq i$）。这与 $T$ 映入 $W$ 的事实矛盾。
因此可以增大 $i$，并假设 $(g_i, g'_i) : T_i \to X \times_S X$
是到 $W$ 的态射。由 $W$ 的构造以及 $T_i$ 的拟紧性，可以找到有限仿射开覆盖
$T_i = T_{1, i} \cup \ldots \cup T_{n, i}$，使得
$(g_i, g'_i)|_{T_{j, i}}$ 是到 $U \times_V U$ 的态射，其中
$(U, V)$ 是上述 $W$ 的定义中出现的一对开集。只需证明 $g_{i'}$
与 $g'_{i'}$ 在每个 $f_{i'i}^{-1}(T_{j, i})$ 上相等，
因而问题约化到仿射情形。仿射情形由《代数》引理
\ref{algebra-lemma-characterize-finite-presentation} 以及环映射
$\mathcal{O}_S(V) \to \mathcal{O}_X(U)$ 有限呈示这一事实得出
（见《态射》引理
\ref{morphisms-lemma-locally-finite-presentation-characterize}）。

\medskip\noindent
最后证明存在性部分 (b)。设 $g : T \to X$ 是 $S$ 上的概形态射。
可以找到有限仿射开覆盖
$T = W_1 \cup \ldots \cup W_n$，使得对每个
$j \in \{1, \ldots, n\}$，存在仿射开集 $U_j \subset X$ 与
$V_j \subset S$，满足 $f(U_j) \subset V_j$ 和 $g(W_j) \subset U_j$。
由引理 \ref{limits-lemma-descend-opens} 及 \ref{limits-lemma-limit-affine}，
可以（在必要时缩小 $I$ 后）假设存在与过渡映射相容的仿射开覆盖
$T_i = W_{1, i} \cup \ldots \cup W_{n, i}$，并且
$W_j = \lim_i W_{j, i}$。把《代数》引理
\ref{algebra-lemma-characterize-finite-presentation} 应用于仿射概形
$U_j$、$V_j$、$W_{j, i}$ 与 $W_j$ 对应的环；这里使用
$\mathcal{O}_S(V_j) \to \mathcal{O}_X(U_j)$ 有限呈示
（见《态射》引理
\ref{morphisms-lemma-locally-finite-presentation-characterize}）。
于是对每个 $j$，可以找到指标 $i_j \in I$ 及态射
$g_{j, i_j} : W_{j, i_j} \to X$，使得
$g_{j, i_j} \circ f_i|_{W_j} : W_j \to W_{j, i} \to X$
等于 $g|_{W_j}$。利用已证的 (a) 以及
$W_{j_1, i} \cap W_{j_2, i}$ 的拟紧性（它由 $T_i$ 拟分离得出），
可以找到指标 $i' \in I$，它大于所有 $i_j$，并使得
$$
g_{j_1, i_{j_1}} \circ f_{i'i_{j_1}}|_{W_{j_1, i'} \cap W_{j_2, i'}} =
g_{j_2, i_{j_2}} \circ f_{i'i_{j_2}}|_{W_{j_1, i'} \cap W_{j_2, i'}}
$$
对所有 $j_1, j_2 \in \{1, \ldots, n\}$ 成立。因此各态射
$g_{j, i_j} \circ f_{i'i_j}|_{W_{j, i'}}$ 可黏合成所需的态射
$T_{i'} \to X$。
\end{proof}

\begin{remark}
\label{limits-remark-limit-preserving}
设 $S$ 是概形。称函子
$F : (\Sch/S)^{opp} \to \textit{Sets}$ {\it 保持极限}，如果对每个
仿射概形有向逆系 $\{T_i\}_{i \in I}$，其极限为 $T$ 时，都有
$F(T) = \colim_i F(T_i)$。设 $X$ 是 $S$ 上的概形，并设
$h_X : (\Sch/S)^{opp} \to \textit{Sets}$ 是其点函子；见《概形》第
\ref{schemes-section-representable} 节。用这种术语，命题
\ref{limits-proposition-characterize-locally-finite-presentation} 断言：
概形 $X$ 在 $S$ 上局部有限呈示，当且仅当 $h_X$ 保持极限。
\end{remark}

\begin{lemma}
\label{limits-lemma-surjection-is-enough}
设 $f : X \to S$ 是概形态射。如果对每个有向极限
$T = \lim_{i \in I} T_i$（由 $S$ 上的仿射概形组成），映射
$$
\colim \Mor_S(T_i, X) \longrightarrow \Mor_S(T, X)
$$
都满射，则 $f$ 局部有限呈示。换言之，在命题
\ref{limits-proposition-characterize-locally-finite-presentation} 的
(2) 和 (3) 中，只需检验该映射的满射性。
\end{lemma}

\begin{proof}
证明与命题 \ref{limits-proposition-characterize-locally-finite-presentation}
中“(2) 蕴含 (1)”的证明完全相同。任取仿射开集 $U \subset X$
与 $V \subset S$，使得 $f(U) \subset V$。必须证明
$\mathcal{O}_S(V) \to \mathcal{O}_X(U)$ 有限呈示。
设 $(A_i, \varphi_{ii'})$ 是 $\mathcal{O}_S(V)$-代数的有向系。
令 $A = \colim_i A_i$。根据《代数》引理
\ref{algebra-lemma-characterize-finite-presentation}，只需证明
$$
\colim_i \Hom_{\mathcal{O}_S(V)}(\mathcal{O}_X(U), A_i) \to
\Hom_{\mathcal{O}_S(V)}(\mathcal{O}_X(U), A)
$$
满射。考察概形 $T_i = \Spec(A_i)$。它们组成
$V$-概形逆系，以 $I$ 为指标，其过渡态射 $f_{ii'} : T_i \to T_{i'}$
由 $\mathcal{O}_S(V)$-代数映射 $\varphi_{i'i}$ 诱导。令
$T := \Spec(A) = \lim_i T_i$。用概形态射集表示，上式化为
$$
\colim_i \Mor_V(T_i, U) \to \Mor_V(\lim_i T_i, U)
$$
首先注意到
$\Mor_V(T_i, U) = \Mor_S(T_i, U)$
以及
$\Mor_V(T, U) = \Mor_S(T, U)$。
因此必须证明
$$
\colim_i \Mor_S(T_i, U) \to
\Mor_S(\lim_i T_i, U)
$$
满射，而已知
$$
\colim_i \Mor_S(T_i, X) \to
\Mor_S(\lim_i T_i, X)
$$
满射。所以只需证明：给定态射 $g_i : T_i \to X$（在 $S$ 上），若复合
$T \to T_i \to X$ 的像落在 $U$ 中，则存在某个 $i' \geq i$，
使复合 $g_{i'} : T_{i'} \to T_i \to X$ 的像落在 $U$ 中。
记 $Z_{i'} = g_{i'}^{-1}(X \setminus U)$。假设每个 $Z_{i'}$ 都非空，
以导出矛盾。由引理 \ref{limits-lemma-limit-closed-nonempty}，存在一点 $t$，
它属于 $T$，并且映入 $Z_{i'}$（对所有 $i' \geq i$）。
这样的点不映入 $U$，矛盾。
\end{proof}

\noindent
下面是命题 \ref{limits-proposition-characterize-locally-finite-presentation}
的一个应用示例。

\begin{lemma}
\label{limits-lemma-morphism-glueing-near-closed-point}
设 $S$ 是概形，$X$ 和 $Y$ 是 $S$ 上的概形。假设 $Y$
在 $S$ 上局部有限呈示。设 $x \in X$ 是闭点，并且
$U = X \setminus \{x\} \to X$ 拟紧。令
$V = \Spec(\mathcal{O}_{X, x}) \setminus \{x\}$，则有双射
{\small
$$
\left\{
\begin{matrix}
\text{morphisms }X \to Y\text{ 相对于 }S
\end{matrix}
\right\}
\longrightarrow
\left\{
\begin{matrix}
(a, b)\text{ 其中 }
a : U \to Y\text{ 且 }b : \Spec(\mathcal{O}_{X, x}) \to Y\\
\text{ are morphisms over }S
\text{ which agree over }V
\end{matrix}
\right\}
$$
}
\end{lemma}

\begin{proof}
设 $W \subset X$ 是 $x$ 的开邻域。由概形的黏合（见《概形》第
\ref{schemes-section-glueing-schemes} 节），如果考察态射对
$a : U \to Y$ 与 $c : W \to Y$，并要求它们在 $U \cap W$ 上相同，
则结论成立。有
$\mathcal{O}_{X, x} = \colim \mathcal{O}_W(W)$，其中 $W$
遍历 $x$ 在 $X$ 中的仿射开邻域。因此
$\Spec(\mathcal{O}_{X, x}) = \lim W$，其中 $W$ 遍历 $s$ 的仿射开邻域。
故由命题 \ref{limits-proposition-characterize-locally-finite-presentation}，
任意态射 $b : \Spec(\mathcal{O}_{X, x}) \to Y$（在 $S$ 上）
都来自态射 $c : W \to Y$，其中 $W$ 如上
（并且 $c$ 在进一步缩小 $W$ 后唯一）。
对每个满足 $x \in W$ 的仿射开集，可知 $U \cap W$ 拟紧，
因为 $U \to X$ 拟紧。因此
$V = \lim W \cap U = \lim W \setminus \{x\}$ 是拟紧且拟分离概形的极限
（见引理 \ref{limits-lemma-directed-inverse-system-has-limit}）。
所以，如果 $a$ 与 $b$ 在 $V$ 上相同，那么在缩小 $W$ 后，
$a$ 与 $c$ 在 $U \cap W$ 上相同（由同一命题）。引理得证。
\end{proof}


\section{相对逼近}
\label{limits-section-relative-approximation}

\noindent
我们讨论命题 \ref{limits-proposition-approximate} 在给定基底上的各种变体。

\begin{lemma}
\label{limits-lemma-approximate-morphism}
设 $f : X \to S$ 是拟紧且拟分离概形之间的态射。则存在有向集 $I$
以及态射逆系 $(f_i : X_i \to S_i)$，后者以 $I$ 为指标，使得过渡态射
$X_i \to X_{i'}$ 和 $S_i \to S_{i'}$ 仿射，$X_i$ 与 $S_i$
均为 $\mathbf{Z}$ 上有限型的，并且
$(X \to S) = \lim (X_i \to S_i)$。
\end{lemma}

\begin{proof}
按命题 \ref{limits-proposition-approximate}，写成
$X = \lim_{a \in A} X_a$ 以及 $S = \lim_{b \in B} S_b$；
换言之，$X_a$ 与 $S_b$ 均为 $\mathbf{Z}$ 上有限型的，并且过渡态射仿射。

\medskip\noindent
固定 $b \in B$。把命题
\ref{limits-proposition-characterize-locally-finite-presentation}
应用于 $S_b$ 和 $X = \lim X_a$（均在 $\mathbf{Z}$ 上），可找到
$a \in A$ 及态射 $f_{a, b} : X_a \to S_b$，使图
$$
\xymatrix{
X \ar[d] \ar[r] & S \ar[d] \\
X_a \ar[r] & S_b
}
$$
交换。令 $I$ 为以这种方式得到的三元组 $(a, b, f_{a, b})$ 之集合。

\medskip\noindent
设 $(a, b, f_{a, b})$ 与 $(a', b', f_{a', b'})$ 属于 $I$。
取 $b'' \leq \min(b, b')$。再次利用命题
\ref{limits-proposition-characterize-locally-finite-presentation}，存在
$a'' \geq \max(a, a')$，使复合
$X_{a''} \to X_a \to S_b \to S_{b''}$ 与
$X_{a''} \to X_{a'} \to S_{b'} \to S_{b''}$ 相等。
在 $I$ 上赋予如下预序：
$$
(a, b, f_{a, b}) \geq (a', b', f_{a', b'})
\Leftrightarrow
a \geq a',\ b \geq b',\text{ 且 }
g_{b, b'} \circ f_{a, b} = f_{a', b'} \circ h_{a, a'}
$$
其中 $h_{a, a'} : X_a \to X_{a'}$ 与
$g_{b, b'} : S_b \to S_{b'}$ 是过渡态射。上面的论述表明 $I$ 有向，
且映射 $I \to A$、$(a, b, f_{a, b}) \mapsto a$ 以及
$I \to B$、$(a, b, f_{a, b})$ 都是共尾的。如果对
$i = (a, b, f_{a, b})$ 令 $X_i = X_a$、$S_i = S_b$ 且
$f_i = f_{a, b}$，便得到以 $I$ 为指标的态射逆系，并且
$$
\lim_{i \in I} X_i = \lim_{a \in A} X_a = X
\quad\text{且}\quad
\lim_{i \in I} S_i = \lim_{b \in B} S_b = S
$$
这由《范畴》引理 \ref{categories-lemma-initial} 得出
（回忆：$I$ 上的极限实际上是与 $I$ 相伴的反范畴上的极限，
所以共尾变为始）。证明完毕。
\end{proof}

\begin{lemma}
\label{limits-lemma-relative-approximation}
设 $f : X \to S$ 是概形态射。假设
\begin{enumerate}
\item $X$ 拟紧且拟分离；并且
\item $S$ 拟分离。
\end{enumerate}
则 $X = \lim X_i$ 是概形有向系的极限，其中 $X_i$ 在 $S$ 上有限呈示，
且过渡态射在 $S$ 上仿射。
\end{lemma}

\begin{proof}
由于 $f(X)$ 拟紧，可以把 $S$ 替换为一个包含 $f(X)$ 的拟紧开集。
因此可以假设 $S$ 拟紧。由引理 \ref{limits-lemma-approximate-morphism}，可将
$(X \to S) = \lim (X_i \to S_i)$ 写成某个有向逆系，其中各态射涉及的概形
均为 $\mathbf{Z}$ 上有限型，且过渡态射仿射。由于极限与极限交换
（《范畴》引理 \ref{categories-lemma-colimits-commute}），有
$X = \lim X_i \times_{S_i} S$。设 $i \geq i'$，其中二者属于 $I$。态射
$X_i \times_{S_i} S \to X_{i'} \times_{S_{i'}} S$ 是仿射的，因为它是复合
$$
X_i \times_{S_i} S \to X_i \times_{S_{i'}} S \to X_{i'} \times_{S_{i'}} S
$$
其中第一个态射是闭浸入（由《概形》引理
\ref{schemes-lemma-fibre-product-after-map}），第二个态射是仿射态射的基变换
（《态射》引理 \ref{morphisms-lemma-base-change-affine}），
而仿射态射的复合仿射（《态射》引理
\ref{morphisms-lemma-composition-affine}）。
态射 $f_i$ 有限呈示（《态射》引理
\ref{morphisms-lemma-noetherian-finite-type-finite-presentation} 及
\ref{morphisms-lemma-finite-presentation-permanence}），
所以其基变换 $X_i \times_{f_i, S_i} S \to S$ 也有限呈示
（《态射》引理
\ref{morphisms-lemma-base-change-finite-presentation}）。
\end{proof}

\begin{lemma}
\label{limits-lemma-integral-limit-finite-and-finite-presentation}
设 $X \to S$ 是整态射，且 $S$ 拟紧并拟分离。则
$X = \lim X_i$，其中 $X_i \to S$ 有限且有限呈示。
\end{lemma}

\begin{proof}
考察层 $\mathcal{A} = f_*\mathcal{O}_X$。这是拟相干
$\mathcal{O}_S$-代数层；见《概形》引理
\ref{schemes-lemma-push-forward-quasi-coherent}。
结合《性质》引理
\ref{properties-lemma-integral-algebra-directed-colimit-finite}，
可将 $\mathcal{A} = \colim_i \mathcal{A}_i$ 写成有限且有限呈示
$\mathcal{O}_S$-代数的滤过余极限。于是
$$
X_i = \underline{\Spec}_S(\mathcal{A}_i)
\longrightarrow
S
$$
是有限且有限呈示的概形态射。按构造，$X = \lim_i X_i$，引理得证。
\end{proof}


\section{态射性质的下降}
\label{limits-section-descent-of-properties}

\noindent
本节是第 \ref{limits-section-descent} 节对 $S$ 上态射性质的对应版本。
我们将在下列情形中工作。

\begin{situation}
\label{limits-situation-descent-property}
设 $S = \lim S_i$ 是带仿射过渡态射的概形有向系之极限
（引理 \ref{limits-lemma-directed-inverse-system-has-limit}）。
设 $0 \in I$，并设 $f_0 : X_0 \to Y_0$ 是 $S_0$ 上的概形态射。
假设 $S_0$、$X_0$、$Y_0$ 拟紧且拟分离。设
$f_i : X_i \to Y_i$ 是 $f_0$ 到 $S_i$ 的基变换，并设
$f : X \to Y$ 是 $f_0$ 到 $S$ 的基变换。
\end{situation}

\begin{lemma}
\label{limits-lemma-descend-affine-finite-presentation}
记号与假设同情形 \ref{limits-situation-descent-property}。如果 $f$ 仿射，
则存在指标 $i \geq 0$，使得 $f_i$ 仿射。
\end{lemma}

\begin{proof}
设 $Y_0 = \bigcup_{j = 1, \ldots, m} V_{j, 0}$ 是有限仿射开覆盖。
令 $U_{j, 0} = f_0^{-1}(V_{j, 0})$。对 $i \geq 0$，记
$V_{j, i}$ 为 $V_{j, 0}$ 在 $Y_i$ 中的逆像，并令
$U_{j, i} = f_i^{-1}(V_{j, i})$。类似地，有
$U_j = f^{-1}(V_j)$。于是 $U_j = \lim_{i \geq 0} U_{j, i}$
（见引理 \ref{limits-lemma-directed-inverse-system-has-limit}）。
由于按假设 $U_j$ 仿射，由引理 \ref{limits-lemma-limit-affine} 可知，
$U_{j, i}$ 对充分大的 $i$ 都仿射。因为 $j$ 只有有限多个，
可以选取一个 $i$，它对所有 $j$ 都适用。因此 $f_i$
对充分大的 $i$ 仿射；见《态射》引理
\ref{morphisms-lemma-characterize-affine}。
\end{proof}

\begin{lemma}
\label{limits-lemma-descend-finite-finite-presentation}
记号与假设同情形 \ref{limits-situation-descent-property}。如果
\begin{enumerate}
\item $f$ 是有限态射；并且
\item $f_0$ 局部有限型，
\end{enumerate}
则存在 $i \geq 0$，使得 $f_i$ 有限。
\end{lemma}

\begin{proof}
有限态射仿射；见《态射》定义
\ref{morphisms-definition-integral}。因此由上面的引理
\ref{limits-lemma-descend-affine-finite-presentation}，增大 $0$ 后可假设
$f_0$ 仿射。把 $Y_0$ 写成有限多个仿射开集之并，便约化到
$X_0$ 与 $Y_0$ 仿射且映入同一个仿射开集 $W \subset S_0$ 的情形。
相应的代数陈述来自《代数》引理
\ref{algebra-lemma-colimit-finite}。
\end{proof}

\begin{lemma}
\label{limits-lemma-descend-unramified}
记号与假设同情形 \ref{limits-situation-descent-property}。如果
\begin{enumerate}
\item $f$ 非分歧；并且
\item $f_0$ 局部有限型，
\end{enumerate}
则存在 $i \geq 0$，使得 $f_i$ 非分歧。
\end{lemma}

\begin{proof}
选取有限仿射开覆盖
$Y_0 = \bigcup_{j = 1, \ldots, m} Y_{j, 0}$，
使每个 $Y_{j, 0}$ 都映入仿射开集 $S_{j, 0} \subset S_0$。
对每个 $j$，设
$f_0^{-1}Y_{j, 0} = \bigcup_{k = 1, \ldots, n_j} X_{k, 0}$
为有限仿射开覆盖。因为非分歧性是局部的，只需对仿射概形的态射
$X_{k, i} \to Y_{j, i} \to S_{j, i}$ 证明此引理；这些态射是
$X_{k, 0} \to Y_{j, 0} \to S_{j, 0}$ 到 $S_i$ 的基变换。
因此约化到 $X_0, Y_0, S_0$ 均仿射的情形。

\medskip\noindent
在仿射情形，问题约化为以下代数结果。假设
$R = \colim_{i \in I} R_i$。对某个 $0 \in I$，设给定有限型
$R_0$-代数映射 $A_i \to B_i$。如果
$R \otimes_{R_0} A_0 \to R \otimes_{R_0} B_0$ 非分歧，
则对某个 $i \geq 0$，映射
$R_i \otimes_{R_0} A_0 \to R_i \otimes_{R_0} B_0$ 非分歧。
这来自《代数》引理
\ref{algebra-lemma-colimit-unramified}。
\end{proof}

\begin{lemma}
\label{limits-lemma-descend-closed-immersion-finite-presentation}
记号与假设同情形 \ref{limits-situation-descent-property}。如果
\begin{enumerate}
\item $f$ 是闭浸入；并且
\item $f_0$ 局部有限型，
\end{enumerate}
则存在 $i \geq 0$，使得 $f_i$ 是闭浸入。
\end{lemma}

\begin{proof}
闭浸入仿射；见《态射》引理
\ref{morphisms-lemma-closed-immersion-affine}。因此由上面的引理
\ref{limits-lemma-descend-affine-finite-presentation}，增大 $0$ 后可假设
$f_0$ 仿射。把 $Y_0$ 写成有限多个仿射开集之并，便约化到
$X_0$ 与 $Y_0$ 仿射且映入同一个仿射开集 $W \subset S_0$ 的情形。
相应的代数陈述是《代数》引理
\ref{algebra-lemma-colimit-surjective} 的推论。
\end{proof}

\begin{lemma}
\label{limits-lemma-descend-separated-finite-presentation}
记号与假设同情形 \ref{limits-situation-descent-property}。如果 $f$ 分离，
则 $f_i$ 对某个 $i \geq 0$ 分离。
\end{lemma}

\begin{proof}
把引理 \ref{limits-lemma-descend-closed-immersion-finite-presentation}
应用于对角态射
$\Delta_{X_0/S_0} : X_0 \to X_0 \times_{S_0} X_0$。
这是允许的，因为对角态射局部有限型，并且纤维积
$X_0 \times_{S_0} X_0$ 拟紧且拟分离；见《概形》引理
\ref{schemes-lemma-diagonal-immersion}、《态射》引理
\ref{morphisms-lemma-immersion-locally-finite-type} 以及《概形》注记
\ref{schemes-remark-quasi-compact-and-quasi-separated}。
\end{proof}

\begin{lemma}
\label{limits-lemma-descend-flat-finite-presentation}
记号与假设同情形 \ref{limits-situation-descent-property}。如果
\begin{enumerate}
\item $f$ 平坦；
\item $f_0$ 局部有限呈示，
\end{enumerate}
则 $f_i$ 对某个 $i \geq 0$ 平坦。
\end{lemma}

\begin{proof}
选取有限仿射开覆盖
$Y_0 = \bigcup_{j = 1, \ldots, m} Y_{j, 0}$，
使每个 $Y_{j, 0}$ 都映入仿射开集 $S_{j, 0} \subset S_0$。
对每个 $j$，设
$f_0^{-1}Y_{j, 0} = \bigcup_{k = 1, \ldots, n_j} X_{k, 0}$
为有限仿射开覆盖。因为平坦性是局部的，只需对仿射概形的态射
$X_{k, i} \to Y_{j, i} \to S_{j, i}$ 证明此引理；这些态射是
$X_{k, 0} \to Y_{j, 0} \to S_{j, 0}$ 到 $S_i$ 的基变换。
因此约化到 $X_0, Y_0, S_0$ 均仿射的情形。

\medskip\noindent
在仿射情形，问题约化为以下代数结果。假设
$R = \colim_{i \in I} R_i$。对某个 $0 \in I$，设给定有限呈示
$R_0$-代数映射 $A_i \to B_i$。如果
$R \otimes_{R_0} A_0 \to R \otimes_{R_0} B_0$ 平坦，
则对某个 $i \geq 0$，映射
$R_i \otimes_{R_0} A_0 \to R_i \otimes_{R_0} B_0$ 平坦。
这来自《代数》引理
\ref{algebra-lemma-flat-finite-presentation-limit-flat} 的第 (3) 部分。
\end{proof}


\begin{lemma}
\label{limits-lemma-descend-finite-locally-free}
记号与假设同情形 \ref{limits-situation-descent-property}。如果
\begin{enumerate}
\item $f$ 有限局部自由（次数为 $d$）；
\item $f_0$ 局部有限呈示，
\end{enumerate}
则 $f_i$ 有限局部自由（次数为 $d$），其中 $i \geq 0$ 为某个指标。
\end{lemma}

\begin{proof}
由引理 \ref{limits-lemma-descend-flat-finite-presentation} 及
\ref{limits-lemma-descend-finite-finite-presentation}，可找到 $i$，使 $f_i$
平坦且有限。另一方面，$f_i$ 局部有限呈示。因此由《态射》引理
\ref{morphisms-lemma-finite-flat}，$f_i$ 有限局部自由。
如果进一步假设 $f$ 有限局部自由且次数为 $d$，则 $Y \to Y_i$ 的像
包含在开闭轨迹 $W_d \subset Y_i$ 中；在该轨迹上，$f_i$ 的次数为 $d$。
由引理 \ref{limits-lemma-limit-contained-in-constructible}，对某个
$i' \geq i$，$Y_{i'} \to Y_i$ 的像包含在 $W_d$ 中。
于是 $f_{i'}$ 有限局部自由且次数为 $d$。
\end{proof}

\begin{lemma}
\label{limits-lemma-descend-smooth}
记号与假设同情形 \ref{limits-situation-descent-property}。如果
\begin{enumerate}
\item $f$ 光滑；
\item $f_0$ 局部有限呈示，
\end{enumerate}
则 $f_i$ 对某个 $i \geq 0$ 光滑。
\end{lemma}

\begin{proof}
光滑性在源和靶上都是局部的（《态射》引理
\ref{morphisms-lemma-smooth-characterize}），因此可以假设
$S_0, X_0, Y_0$ 仿射（细节从略）。相应的代数事实是《代数》引理
\ref{algebra-lemma-colimit-smooth}。
\end{proof}

\begin{lemma}
\label{limits-lemma-descend-etale}
记号与假设同情形 \ref{limits-situation-descent-property}。如果
\begin{enumerate}
\item $f$ 平展；
\item $f_0$ 局部有限呈示，
\end{enumerate}
则 $f_i$ 对某个 $i \geq 0$ 平展。
\end{lemma}

\begin{proof}
平展性在源和靶上都是局部的（《态射》引理
\ref{morphisms-lemma-etale-characterize}），因此可以假设
$S_0, X_0, Y_0$ 仿射（细节从略）。相应的代数事实是《代数》引理
\ref{algebra-lemma-colimit-etale}。
\end{proof}

\begin{lemma}
\label{limits-lemma-descend-isomorphism}
记号与假设同情形 \ref{limits-situation-descent-property}。如果
\begin{enumerate}
\item $f$ 是同构；并且
\item $f_0$ 局部有限呈示，
\end{enumerate}
则 $f_i$ 对某个 $i \geq 0$ 是同构。
\end{lemma}

\begin{proof}
由引理 \ref{limits-lemma-descend-etale} 及
\ref{limits-lemma-descend-closed-immersion-finite-presentation}，可找到 $i$，
使 $f_i$ 平坦且为闭浸入。于是 $f_i$ 把 $X_i$ 等同于 $Y_i$
的一个开闭子概形；见《态射》引理
\ref{morphisms-lemma-flat-closed-immersions-finite-presentation}。
按假设，$Y \to Y_i$ 的像落在 $f_i(X_i)$ 中。因此由引理
\ref{limits-lemma-limit-contained-in-constructible}，可知 $Y_{i'}$
映入 $f_i(X_i)$，其中 $i' \geq i$ 为某个指标。这蕴含
$X_{i'} \to Y_{i'}$ 满射，
结论得证。
\end{proof}

\begin{lemma}
\label{limits-lemma-descend-open-immersion}
记号与假设同情形 \ref{limits-situation-descent-property}。如果
\begin{enumerate}
\item $f$ 是开浸入；并且
\item $f_0$ 局部有限呈示，
\end{enumerate}
则 $f_i$ 对某个 $i \geq 0$ 是开浸入。
\end{lemma}

\begin{proof}
由引理 \ref{limits-lemma-descend-etale}，可找到 $i$，使 $f_i$ 平展。
于是 $V_i = f_i(X_i)$ 是 $Y_i$ 的拟紧开子概形
（《态射》引理 \ref{morphisms-lemma-etale-open}）。
令 $V$ 以及 $V_{i'}$（$i' \geq i$）分别为 $V_i$ 在
$Y$ 及 $Y_{i'}$ 中的逆像。于是 $f : X \to V$ 是同构
（即它是满的开浸入）。所以由引理
\ref{limits-lemma-descend-isomorphism}，$X_{i'} \to V_{i'}$
对某个 $i' \geq i$ 是同构，正合所需。
\end{proof}

\begin{lemma}
\label{limits-lemma-descend-immersion}
记号与假设同情形 \ref{limits-situation-descent-property}。如果
\begin{enumerate}
\item $f$ 是浸入；并且
\item $f_0$ 局部有限型，
\end{enumerate}
则 $f_i$ 对某个 $i \geq 0$ 是浸入。
\end{lemma}

\begin{proof}
存在开集 $V \subset Y$，使态射 $f$ 分解为 $X \to V \to Y$，
并且 $X \to V$ 是闭浸入；见《概形》第
\ref{schemes-section-immersions} 节中的讨论。
由于 $X$ 拟紧，可以且确实假设 $V$ 是 $Y$ 的拟紧开集。
由引理 \ref{limits-lemma-descend-opens}，增大 $0$ 后，可找到拟紧开集
$V_0 \subset Y_0$，使 $V$ 是 $V_0$ 的逆像。于是 $V_0$ 在 $X_0$
中的逆像是拟紧开集，其在 $X$ 中的逆像为 $X$。因此，把同一引理应用于
$X = \lim X_i$，增大 $0$ 后可假设有分解
$X_0 \to V_0 \to Y_0$。于是对充分大的 $i \geq 0$，态射
$X_i \to V_i$（其中 $V_i = Y_i \times_{Y_0} V_0$）由引理
\ref{limits-lemma-descend-closed-immersion-finite-presentation} 可知是闭浸入，
证明完成。
\end{proof}

\begin{lemma}
\label{limits-lemma-descend-monomorphism}
记号与假设同情形 \ref{limits-situation-descent-property}。如果
\begin{enumerate}
\item $f$ 是单态射；并且
\item $f_0$ 局部有限型，
\end{enumerate}
则 $f_i$ 对某个 $i \geq 0$ 是单态射。
\end{lemma}

\begin{proof}
回忆，概形态射 $V \to W$ 是单态射，当且仅当对角态射
$V \to V \times_W V$ 是同构（《概形》引理
\ref{schemes-lemma-monomorphism}）。
态射 $X_0 \to X_0 \times_{Y_0} X_0$ 由《态射》引理
\ref{morphisms-lemma-diagonal-morphism-finite-type} 可知局部有限呈示。
由于 $X_0 \times_{Y_0} X_0$ 拟紧且拟分离
（《概形》注记
\ref{schemes-remark-quasi-compact-and-quasi-separated}），
由引理 \ref{limits-lemma-descend-isomorphism} 可知，
$\Delta_i : X_i \to X_i \times_{Y_i} X_i$ 对某个 $i \geq 0$ 是同构。
对这个 $i$，
态射 $f_i$ 是单态射。
\end{proof}

\begin{lemma}
\label{limits-lemma-descend-surjective}
记号与假设同情形 \ref{limits-situation-descent-property}。如果
\begin{enumerate}
\item $f$ 满射；并且
\item $f_0$ 局部有限呈示，
\end{enumerate}
则存在 $i \geq 0$，使得 $f_i$ 满射。
\end{lemma}

\begin{proof}
态射 $f_0$ 有限呈示。因此 $E = f_0(X_0)$ 是 $Y_0$ 的可构造子集；
见《态射》引理 \ref{morphisms-lemma-chevalley}。
由于 $f_i$ 是 $f_0$ 沿 $Y_i \to Y_0$ 的基变换，$f_i$ 的像就是
$E$ 在 $Y_i$ 中的逆像。此外，已知 $Y \to Y_0$ 映入 $E$。
因此由引理 \ref{limits-lemma-limit-contained-in-constructible} 即得结论。
\end{proof}

\begin{lemma}
\label{limits-lemma-descend-syntomic}
记号与假设同情形 \ref{limits-situation-descent-property}。如果
\begin{enumerate}
\item $f$ 是 syntomic 态射；并且
\item $f_0$ 局部有限呈示，
\end{enumerate}
则存在 $i \geq 0$，使得 $f_i$ 是 syntomic 态射。
\end{lemma}

\begin{proof}
选取有限仿射开覆盖
$Y_0 = \bigcup_{j = 1, \ldots, m} Y_{j, 0}$，
使每个 $Y_{j, 0}$ 都映入仿射开集 $S_{j, 0} \subset S_0$。
对每个 $j$，设
$f_0^{-1}Y_{j, 0} = \bigcup_{k = 1, \ldots, n_j} X_{k, 0}$
为有限仿射开覆盖。因为 syntomic 性是局部的，只需对仿射概形的态射
$X_{k, i} \to Y_{j, i} \to S_{j, i}$ 证明此引理；这些态射是
$X_{k, 0} \to Y_{j, 0} \to S_{j, 0}$ 到 $S_i$ 的基变换。
因此约化到 $X_0, Y_0, S_0$ 均仿射的情形。

\medskip\noindent
在仿射情形，问题约化为以下代数结果。假设
$R = \colim_{i \in I} R_i$。对某个 $0 \in I$，设给定有限呈示
$R_0$-代数映射 $A_i \to B_i$。如果
$R \otimes_{R_0} A_0 \to R \otimes_{R_0} B_0$ 是 syntomic 态射，
则对某个 $i \geq 0$，映射
$R_i \otimes_{R_0} A_0 \to R_i \otimes_{R_0} B_0$ 是 syntomic 态射。
这来自《代数》引理
\ref{algebra-lemma-colimit-lci}。
\end{proof}


\section{有限型概形闭嵌入有限呈示概形}
\label{limits-section-finite-type-closed-in-finite-presentation}

\noindent
这类结果可见 \cite[Satz 2.10]{Kiehl}。另一个参考文献是
\cite{Conrad-Nagata}。

\begin{lemma}
\label{limits-lemma-locally-finite-type-in-finite-presentation}
设 $f : X \to S$ 是概形态射。假设：
\begin{enumerate}
\item 态射 $f$ 局部有限型。
\item 概形 $X$ 拟紧且拟分离。
\end{enumerate}
则存在有限呈示态射 $f' : X' \to S$，以及浸入
$X \to X'$，后者是 $S$ 上的概形态射。
\end{lemma}

\begin{proof}
由命题 \ref{limits-proposition-approximate}，可写成
$X = \lim_i X_i$，其中每个 $X_i$ 都是 $\mathbf{Z}$ 上有限型的，
过渡态射 $f_{ii'} : X_i \to X_{i'}$ 仿射。考察交换图
$$
\xymatrix{
X \ar[r] \ar[rd] & X_{i, S} \ar[r] \ar[d] & X_i \ar[d] \\
& S \ar[r] & \Spec(\mathbf{Z})
}
$$
注意，$X_i$ 在 $\Spec(\mathbf{Z})$ 上有限呈示；见《态射》引理
\ref{morphisms-lemma-noetherian-finite-type-finite-presentation}。
所以由《态射》引理
\ref{morphisms-lemma-base-change-finite-presentation}，
基变换 $X_{i, S} \to S$ 有限呈示。因此只需证明，箭头
$X \to X_{i, S}$ 对充分大的 $i$ 是浸入。

\medskip\noindent
为此，选取有限仿射开覆盖 $X = V_1 \cup \ldots \cup V_n$，
使 $f$ 把每个 $V_j$ 映入某个仿射开集 $U_j \subset S$。
设 $h_{j, a} \in \mathcal{O}_X(V_j)$ 是有限组元素，它们生成
$\mathcal{O}_X(V_j)$ 作为 $\mathcal{O}_S(U_j)$-代数；见《态射》引理
\ref{morphisms-lemma-locally-finite-type-characterize}。
由引理 \ref{limits-lemma-descend-opens} 及 \ref{limits-lemma-limit-affine}，
可以（在必要时缩小 $I$ 后）假设存在与过渡映射相容的仿射开覆盖
$X_i = V_{1, i} \cup \ldots \cup V_{n, i}$，使
$V_j = \lim_i V_{j, i}$。由引理 \ref{limits-lemma-descend-section}，
可以把 $i$ 取得充分大，使每个 $h_{j, a}$ 都来自某个元素
$h_{j, a, i} \in \mathcal{O}_{X_i}(V_{j, i})$。
于是下图中的箭头
$$
V_j \longrightarrow U_j \times_{\Spec(\mathbf{Z})} V_{j, i} =
(V_{j, i})_{U_j} \subset (V_{j, i})_S \subset X_{i, S}
$$
是闭浸入。由于 $\bigcup (V_{j, i})_{U_j}$ 构成 $X_{i, S}$
的一个开集，并且 $(V_{j, i})_{U_j}$ 在 $X$ 中的逆像是 $V_j$，
可知 $X \to X_{i, S}$ 是浸入。
\end{proof}

\begin{remark}
\label{limits-remark-cannot-do-better}
如果不对 $S$ 和态射 $f : X \to S$ 作更多假设，就不能把结论加强。
例如，一般不能找到如引理所述的{\it 闭}浸入 $X \to X'$。
原因是这将蕴含 $f$ 拟紧，而事实未必如此。一个例子是：令 $S$ 为把无限维
仿射空间的 $0$ 加倍所得的概形，并取 $X$ 为这两个无限维仿射空间之一。
\end{remark}

\begin{lemma}
\label{limits-lemma-finite-type-closed-in-finite-presentation}
设 $f : X \to S$ 是概形态射。假设：
\begin{enumerate}
\item 态射 $f$ 属于局部有限型。
\item 概形 $X$ 拟紧且拟分离；并且
\item 概形 $S$ 拟分离。
\end{enumerate}
则存在有限呈示态射 $f' : X' \to S$，以及闭浸入
$X \to X'$，后者是 $S$ 上的概形态射。
\end{lemma}

\begin{proof}
由上面的引理 \ref{limits-lemma-locally-finite-type-in-finite-presentation}，
存在有限呈示态射 $Y \to S$ 及浸入
$i : X \to Y$，后者是 $S$ 上的概形态射。对每个点 $x \in X$，存在仿射开集
$V_x \subset Y$，使 $i^{-1}(V_x) \to V_x$ 是闭浸入。
由于 $X$ 拟紧，可以找到有限多个仿射开集
$V_1, \ldots, V_n \subset Y$，使得
$i(X) \subset V_1 \cup \ldots \cup V_n$，且
$i^{-1}(V_j) \to V_j$ 是闭浸入。换言之，
$i : X \to X' = V_1 \cup \ldots \cup V_n$ 是 $S$ 上概形的闭浸入。
由于 $S$ 拟分离，且 $Y$ 在 $S$ 上拟分离，由《概形》引理
\ref{schemes-lemma-separated-permanence} 可知 $Y$ 拟分离。
所以开浸入 $X' = V_1 \cup \ldots \cup V_n \to Y$ 拟紧。
这蕴含 $X' \to Y$ 有限呈示；见《态射》引理
\ref{morphisms-lemma-quasi-compact-open-immersion-finite-presentation}。
由于 $X' \to Y \to S$ 是有限呈示态射的复合，故它有限呈示
（见《态射》引理
\ref{morphisms-lemma-composition-finite-presentation}），从而得出结论。
\end{proof}

\begin{lemma}
\label{limits-lemma-closed-is-limit-closed-and-finite-presentation}
设 $X \to Y$ 是概形的闭浸入。假设 $Y$ 拟紧且拟分离。
则 $X$ 可写成有向极限 $X = \lim X_i$，其中各项均为 $Y$ 上的概形，
且 $X_i \to Y$ 是有限呈示的闭浸入。
\end{lemma}

\begin{proof}
设 $\mathcal{I} \subset \mathcal{O}_Y$ 是把 $X$ 定义为 $Y$
的闭子概形的拟相干理想层。由《性质》引理
\ref{properties-lemma-quasi-coherent-colimit-finite-type}，
可将 $\mathcal{I}$ 写成它的有限型拟相干理想层之有向余极限
$\mathcal{I} = \colim_{i \in I} \mathcal{I}_i$。
设 $X_i \subset Y$ 是由 $\mathcal{I}_i$ 定义的闭子概形。
它们组成以 $I$ 为指标的概形逆系。过渡态射 $X_i \to X_{i'}$
仿射，因为它们是闭浸入。每个 $X_i$ 拟紧且拟分离，因为它是 $Y$
的闭子概形，而按假设 $Y$ 拟紧且拟分离。
有 $X = \lim_i X_i$；这直接来自
$\mathcal{I} = \colim_{i \in I} \mathcal{I}_a$ 这一事实。
每个态射 $X_i \to Y$ 都有限呈示；见《态射》引理
\ref{morphisms-lemma-closed-immersion-finite-presentation}。
\end{proof}


\begin{lemma}
\label{limits-lemma-finite-type-is-limit-finite-presentation}
设 $f : X \to S$ 是概形态射。假设
\begin{enumerate}
\item 态射 $f$ 属于局部有限型。
\item 概形 $X$ 拟紧且拟分离；并且
\item 概形 $S$ 拟分离。
\end{enumerate}
则 $X = \lim X_i$，其中 $X_i \to S$ 有限呈示，$X_i$
拟紧且拟分离，而过渡态射 $X_{i'} \to X_i$ 是闭浸入
（这蕴含 $X \to X_i$ 对所有 $i$ 都是闭浸入）。
\end{lemma}

\begin{proof}
由引理 \ref{limits-lemma-finite-type-closed-in-finite-presentation}，
存在闭浸入 $X \to Y$，其中 $Y \to S$ 有限呈示。
由《概形》引理 \ref{schemes-lemma-separated-permanence}，$Y$ 拟分离。
由于 $X$ 拟紧，可以假设 $Y$ 拟紧：只需把 $Y$ 替换为一个包含 $X$
的拟紧开集。由引理
\ref{limits-lemma-closed-is-limit-closed-and-finite-presentation}，可知
$X = \lim X_i$，其中 $X_i \to Y$ 是有限呈示的闭浸入。
由《态射》引理 \ref{morphisms-lemma-composition-finite-presentation}，
态射 $X_i \to S$ 有限呈示。
\end{proof}

\begin{proposition}
\label{limits-proposition-separated-closed-in-finite-presentation}
设 $f : X \to S$ 是概形态射。假设
\begin{enumerate}
\item $f$ 有限型且分离；并且
\item $S$ 拟紧且拟分离。
\end{enumerate}
则存在分离的有限呈示态射 $f' : X' \to S$，以及闭浸入
$X \to X'$，后者是 $S$ 上的概形态射。
\end{proposition}

\begin{proof}
应用引理 \ref{limits-lemma-finite-type-is-limit-finite-presentation}。
由引理 \ref{limits-lemma-eventually-separated}，$X_i \to S$
对充分大的 $i$ 分离，因为已经假设 $X \to S$ 分离。
\end{proof}

\begin{lemma}
\label{limits-lemma-finite-closed-in-finite-finite-presentation}
设 $f : X \to S$ 是概形态射。假设
\begin{enumerate}
\item $f$ 有限；并且
\item $S$ 拟紧且拟分离。
\end{enumerate}
则存在有限且有限呈示的态射 $f' : X' \to S$，以及闭浸入
$X \to X'$，后者是 $S$ 上的概形态射。
\end{lemma}

\begin{proof}
可以按引理 \ref{limits-lemma-finite-type-is-limit-finite-presentation}
写成 $X = \lim X_i$。应用引理 \ref{limits-lemma-eventually-finite} 可知，
$X_i \to S$ 对充分大的 $i$ 有限。
\end{proof}

\begin{lemma}
\label{limits-lemma-finite-in-finite-and-finite-presentation}
设 $f : X \to S$ 是概形态射。假设
\begin{enumerate}
\item $f$ 有限；并且
\item $S$ 拟紧且拟分离。
\end{enumerate}
则 $X$ 是有向极限 $X = \lim X_i$，其中过渡映射是闭浸入，
各对象 $X_i$ 在 $S$ 上有限且有限呈示。
\end{lemma}

\begin{proof}
可以按引理 \ref{limits-lemma-finite-type-is-limit-finite-presentation}
写成 $X = \lim X_i$。应用引理 \ref{limits-lemma-eventually-finite} 可知，
$X_i \to S$ 对充分大的 $i$ 有限。
\end{proof}










\section{相对对象的下降}
\label{limits-section-descending-relative}

\noindent
下面的引理是本节这类结果的典型例子。我们完整写出这个``标准''证明。
与其阅读此证明，读者或许能更快地自行确信该结果成立。

\begin{lemma}
\label{limits-lemma-descend-finite-presentation}
设 $I$ 是有向集。设 $(S_i, f_{ii'})$ 是以 $I$ 为指标的概形逆系。假设
\begin{enumerate}
\item 态射 $f_{ii'} : S_i \to S_{i'}$ 仿射；
\item 概形 $S_i$ 拟紧且拟分离。
\end{enumerate}
令 $S = \lim_i S_i$。则有：
\begin{enumerate}
\item 对任意有限呈示态射 $X \to S$，存在指标 $i \in I$ 和有限呈示态射
$X_i \to S_i$，使 $X \cong X_{i, S}$ 是 $S$ 上概形的同构。
\item 给定指标 $i \in I$、概形 $X_i$、$Y_i$（它们在 $S_i$ 上有限呈示），
以及态射 $\varphi : X_{i, S} \to Y_{i, S}$（它在 $S$ 上），存在指标
$i' \geq i$ 和态射
$\varphi_{i'} : X_{i, S_{i'}} \to Y_{i, S_{i'}}$，其到 $S$ 的基变换为
$\varphi$。
\item 给定指标 $i \in I$、概形 $X_i$、$Y_i$（它们在 $S_i$ 上有限呈示），
以及一对态射 $\varphi_i, \psi_i : X_i \to Y_i$；若其基变换相等，
即 $\varphi_{i, S} = \psi_{i, S}$，则存在指标 $i' \geq i$，使得
$\varphi_{i, S_{i'}} = \psi_{i, S_{i'}}$。
\end{enumerate}
换言之，$S$ 上有限呈示概形的范畴，是在 $I$ 上取余极限所得；
该图中的各项为 $S_i$ 上有限呈示概形的范畴。
\end{lemma}

\begin{proof}
如果每个概形 $S_i$ 都仿射，并且只考虑 $S_i$ 上（相应地，$S$ 上）
有限呈示的仿射概形，则此引理等价于《代数》引理
\ref{algebra-lemma-colimit-category-fp-algebras}。我们断言，仿射情形蕴含
一般情形的引理。

\medskip\noindent
先证明 (3)。设给定指标 $i \in I$、概形 $X_i$、$Y_i$（它们在
$S_i$ 上有限呈示）及一对态射 $\varphi_i, \psi_i : X_i \to Y_i$。
假设基变换相等：$\varphi_{i, S} = \psi_{i, S}$。采用记号
$X_{i'} = X_{i, S_{i'}}$ 和 $Y_{i'} = Y_{i, S_{i'}}$，其中 $i' \geq i$。
还令 $X = X_{i, S}$ 和 $Y = Y_{i, S}$。由引理
\ref{limits-lemma-scheme-over-limit}，有 $X = \lim_{i' \geq i} X_{i'}$，
对 $Y$ 也同样成立。此外，用 $\varphi_{i'}$ 和 $\psi_{i'}$
（相应地，$\varphi$ 和 $\psi$）表示 $\varphi_i$ 和 $\psi_i$ 到
$S_{i'}$（相应地，到 $S$）的基变换。因此假设意味着
$\varphi = \psi$。由于 $Y_i$ 和 $X_i$ 在 $S_i$ 上有限呈示，而
$S_i$ 拟紧且拟分离，所以 $X_i$ 和 $Y_i$ 也拟紧且拟分离
（见《态射》引理
\ref{morphisms-lemma-finite-presentation-quasi-compact-quasi-separated}）。
因此可选取有限仿射开覆盖 $Y_i = \bigcup V_{j, i}$，使每个
$V_{j, i}$ 都映入 $S_i$ 的一个仿射开集。如上，用 $V_{j, i'}$
表示 $V_{j, i}$ 在 $Y_{i'}$ 中的逆像，用 $V_j$ 表示它在 $Y$ 中的逆像。
浸入 $V_{j, i'} \to Y_{i'}$ 拟紧，而逆像
$U_{j, i'} = \varphi_i^{-1}(V_{j, i'})$ 和
$U_{j, i'}' = \psi_i^{-1}(V_{j, i'})$ 是 $X_{i'}$ 的拟紧开集。
按假设，$V_j$ 在 $\varphi$ 和 $\psi$ 下于 $X$ 中的逆像相等。
因此由引理 \ref{limits-lemma-descend-opens}，存在指标 $i' \geq i$，使得
$U_{j, i'} = U_{j, i'}'$ 在 $X_{i'}$ 中成立。选取有限仿射开覆盖
$U_{j, i'} = U_{j, i'}' = \bigcup W_{j, k, i'}$；它诱导覆盖
$U_{j, i''} = U_{j, i''}' = \bigcup W_{j, k, i''}$，对所有
$i'' \geq i'$ 均如此。
由仿射情形，存在指标 $i''$，使
$\varphi_{i''}|_{W_{j, k, i''}} = \psi_{i''}|_{W_{j, k, i''}}$
对所有 $j, k$ 都成立。
于是 $i''$ 是使 $\varphi_{i''} = \psi_{i''}$ 成立的指标，(3) 得证。

\medskip\noindent
再证明 (2)。设给定指标 $i \in I$、概形 $X_i$、$Y_i$（它们在
$S_i$ 上有限呈示）及态射 $\varphi : X_{i, S} \to Y_{i, S}$。
采用记号 $X_{i'} = X_{i, S_{i'}}$ 和 $Y_{i'} = Y_{i, S_{i'}}$，
其中 $i' \geq i$。还令 $X = X_{i, S}$ 和 $Y = Y_{i, S}$。
由引理 \ref{limits-lemma-scheme-over-limit}，有
$X = \lim_{i' \geq i} X_{i'}$，对 $Y$ 也同样成立。
由于 $Y_i$ 和 $X_i$ 在 $S_i$ 上有限呈示，而 $S_i$ 拟紧且拟分离，
所以 $X_i$ 和 $Y_i$ 也拟紧且拟分离（见《态射》引理
\ref{morphisms-lemma-finite-presentation-quasi-compact-quasi-separated}）。
因此可选取有限仿射开覆盖 $Y_i = \bigcup V_{j, i}$，使每个
$V_{j, i}$ 都映入 $S_i$ 的一个仿射开集。如上，用 $V_{j, i'}$
表示 $V_{j, i}$ 在 $Y_{i'}$ 中的逆像，用 $V_j$ 表示它在 $Y$ 中的逆像。
浸入 $V_j \to Y$ 拟紧，而逆像 $U_j = \varphi^{-1}(V_j)$ 是 $X$
的拟紧开集。因此由引理 \ref{limits-lemma-descend-opens}，存在指标
$i' \geq i$ 以及拟紧开集 $U_{j, i'}$（它属于 $X_{i'}$），其在 $X$
中的逆像是 $U_j$。选取有限仿射开覆盖
$U_{j, i'} = \bigcup W_{j, k, i'}$；它诱导仿射开覆盖
$U_{j, i''} = \bigcup W_{j, k, i''}$（对所有 $i'' \geq i'$），并诱导仿射开覆盖
$U_j = \bigcup W_{j, k}$。由仿射情形，存在指标 $i''$ 和态射
$\varphi_{j, k, i''} : W_{j, k, i''} \to V_{j, i''}$，使
$\varphi|_{W_{j, k}} = \varphi_{j, k, i'', S}$ 对所有 $j, k$ 都成立。
由上面已证的 (3)，存在更大的指标 $i''' \geq i''$，使得
$$
\varphi_{j_1, k_1, i'', S_{i'''}}|_{W_{j_1, k_1, i'''} \cap W_{j_2, k_2, i'''}}
=
\varphi_{j_2, k_2, i'', S_{i'''}}|_{W_{j_1, k_1, i'''} \cap W_{j_2, k_2, i'''}}
$$
对所有 $j_1, j_2, k_1, k_2$ 成立。于是 $i'''$ 是这样的指标：存在态射
$\varphi_{i'''} : X_{i'''} \to Y_{i'''}$，其到 $S$ 的基变换给出
$\varphi$。故 (2) 成立。

\medskip\noindent
最后证明 (1)。设给定概形 $X$，它在 $S$ 上有限呈示。由于 $X$ 在
$S$ 上有限呈示，而 $S$ 拟紧且拟分离，所以 $X$ 也拟紧且拟分离
（见《态射》引理
\ref{morphisms-lemma-finite-presentation-quasi-compact-quasi-separated}）。
选取有限仿射开覆盖 $X = \bigcup U_j$，使每个 $U_j$ 都映入仿射开集
$V_j \subset S$。记 $U_{j_1j_2} = U_{j_1} \cap U_{j_2}$ 及
$U_{j_1j_2j_3} = U_{j_1} \cap U_{j_2} \cap U_{j_3}$。
由引理 \ref{limits-lemma-descend-opens} 及 \ref{limits-lemma-limit-affine}，
可以找到指标 $i_1$ 和仿射开集 $V_{j, i_1} \subset S_{i_1}$，
使每个 $V_j$ 都是它在 $S$ 中的逆像。令 $V_{j, i}$ 为
$V_{j, i_1}$ 在 $S_i$ 中的逆像，其中 $i \geq i_1$。由仿射情形，可以找到
指标 $i_2 \geq i_1$ 和仿射概形 $U_{j, i_2} \to V_{j, i_2}$，使
$U_j = S \times_{S_{i_2}} U_{j, i_2}$ 是其基变换。记
$U_{j, i} = S_i \times_{S_{i_2}} U_{j, i_2}$，其中 $i \geq i_2$。
由引理 \ref{limits-lemma-descend-opens}，存在指标 $i_3 \geq i_2$ 及开子概形
$W_{j_1, j_2, i_3} \subset U_{j_1, i_3}$，其到 $S$ 的基变换等于
$U_{j_1j_2}$。记
$W_{j_1, j_2, i} = S_i \times_{S_{i_3}} W_{j_1, j_2, i_3}$，其中
$i \geq i_3$。
由上面已证的 (2)，存在指标 $i_4 \geq i_3$ 和态射
$\varphi_{j_1, j_2, i_4} : W_{j_1, j_2, i_4} \to W_{j_2, j_1, i_4}$，
其到 $S$ 的基变换给出恒等态射
$U_{j_1j_2} = U_{j_2j_1}$，对所有 $j_1, j_2$ 都如此。对所有 $i \geq i_4$，记其基变换为
$\varphi_{j_1, j_2, i} = \text{id}_S \times \varphi_{j_1, j_2, i_4}$。
我们断言，对某个 $i_5 \geq i_4$，系统
$((U_{j, i_5})_j, (W_{j_1, j_2, i_5})_{j_1, j_2},
(\varphi_{j_1, j_2, i_5})_{j_1, j_2})$ 构成《概形》第
\ref{schemes-section-glueing-schemes} 节所述的黏合资料。
为此，需要验证：对充分大的 $i$，有
$$
\varphi_{j_1, j_2, i}^{-1}(W_{j_2, j_1, i} \cap W_{j_2, j_3, i})
=
W_{j_1, j_2, i} \cap W_{j_1, j_3, i}
$$
并且对充分大的 $i$，余圈条件成立。第一个条件来自引理
\ref{limits-lemma-descend-opens} 及事实 $U_{j_2j_1j_3} = U_{j_1j_2j_3}$。
第二个条件来自上面已证的引理 (3)，以及恒等映射
$\text{id} : U_{j_1j_2} \to U_{j_2j_1}$ 满足余圈条件这一事实。
现在可以应用《概形》引理 \ref{schemes-lemma-glue-schemes}，黏合系统
{\footnotesize\par\noindent\centering
$((U_{j, i_5})_j, (W_{j_1, j_2, i_5})_{j_1, j_2},
(\varphi_{j_1, j_2, i_5})_{j_1, j_2})$\par}
由此得到概形
$X_{i_5} \to S_{i_5}$。按构造，$X_{i_5}$ 到 $S$ 的基变换，是把仿射开集
$U_j$ 沿开集 $U_{j_1} \leftarrow U_{j_1j_2} \rightarrow U_{j_2}$
黏合而成的。因此如所需，有 $S \times_{S_{i_5}} X_{i_5} \cong X$。
\end{proof}

\begin{lemma}
\label{limits-lemma-descend-modules-finite-presentation}
设 $I$ 是有向集。设 $(S_i, f_{ii'})$ 是以 $I$ 为指标的概形逆系。假设
\begin{enumerate}
\item 所有态射 $f_{ii'} : S_i \to S_{i'}$ 都仿射；
\item 所有概形 $S_i$ 都拟紧且拟分离。
\end{enumerate}
令 $S = \lim_i S_i$。则有：
\begin{enumerate}
\item 对任意有限呈示的 $\mathcal{O}_S$-模层 $\mathcal{F}$，存在指标
$i \in I$ 和有限呈示的 $\mathcal{O}_{S_i}$-模层 $\mathcal{F}_i$，使得
$\mathcal{F} \cong f_i^*\mathcal{F}_i$。
\item 设给定指标 $i \in I$、有限呈示的 $\mathcal{O}_{S_i}$-模层
$\mathcal{F}_i$、$\mathcal{G}_i$，以及态射
$\varphi : f_i^*\mathcal{F}_i \to f_i^*\mathcal{G}_i$（它在 $S$ 上）。则存在指标
$i' \geq i$ 和态射
$\varphi_{i'} : f_{i'i}^*\mathcal{F}_i \to f_{i'i}^*\mathcal{G}_i$，
其到 $S$ 的基变换是 $\varphi$。
\item 设给定指标 $i \in I$、有限呈示的 $\mathcal{O}_{S_i}$-模层
$\mathcal{F}_i$、$\mathcal{G}_i$，以及一对态射
$\varphi_i, \psi_i : \mathcal{F}_i \to \mathcal{G}_i$。假设其基变换相等：
$f_i^*\varphi_i = f_i^*\psi_i$。则存在指标 $i' \geq i$，使得
$f_{i'i}^*\varphi_i = f_{i'i}^*\psi_i$。
\end{enumerate}
换言之，$S$ 上有限呈示模的范畴，是在 $I$ 上取余极限所得；
该图中的各项为 $S_i$ 上有限呈示模的范畴。
\end{lemma}

\begin{proof}
我们概述两个证明，但省略细节。

\medskip\noindent
第一种证明。若 $S$ 和 $S_i$ 都是仿射概形，则此引理等价于《代数》引理
\ref{algebra-lemma-colimit-category-fp-modules}。一般情形用 Zariski 黏合
从仿射情形推出。

\medskip\noindent
第二种证明。我们使用：
\begin{enumerate}
\item 拟相干 $\mathcal{O}_S$-模与 $S$ 上向量丛的范畴之间存在等价；见
《构造》第 \ref{constructions-section-vector-bundle} 节；并且
\item 向量丛 $\mathbf{V}(\mathcal{F}) \to S$ 在 $S$ 上有限呈示，当且仅当
$\mathcal{F}$ 是有限呈示的 $\mathcal{O}_S$-模。
\end{enumerate}
由此，可以使用引理 \ref{limits-lemma-descend-finite-presentation} 证明：
$S$ 上有限呈示向量丛的范畴，是在 $I$ 上取余极限所得；
该图中的各项为 $S_i$ 上的向量丛范畴。
\end{proof}

\begin{lemma}
\label{limits-lemma-descend-invertible-modules}
设 $S = \lim S_i$ 是拟紧且拟分离概形 $S_i$ 所成有向系的极限，
过渡态射仿射。则
\begin{enumerate}
\item 任意有限局部自由 $\mathcal{O}_S$-模都是某个有限局部自由
$\mathcal{O}_{S_i}$-模的拉回，其中指标为 $i$；
\item 任意可逆 $\mathcal{O}_S$-模都是某个可逆
$\mathcal{O}_{S_i}$-模的拉回，其中指标为 $i$；并且
\item 任意有限型拟相干理想 $\mathcal{I} \subset \mathcal{O}_S$ 都具有形式
$\mathcal{I}_i \cdot \mathcal{O}_S$，其中某个 $i$ 的
$\mathcal{I}_i \subset \mathcal{O}_{S_i}$ 是有限型拟相干理想。
\end{enumerate}
\end{lemma}

\begin{proof}
设 $\mathcal{E}$ 是有限局部自由 $\mathcal{O}_S$-模。由于有限局部自由模
有限呈示，可以找到 $i$ 和有限呈示的 $\mathcal{O}_{S_i}$-模
$\mathcal{E}_i$，使得 $f_i^*\mathcal{E}_i \cong \mathcal{E}$；见引理
\ref{limits-lemma-descend-modules-finite-presentation}。增大 $i$ 后，可以假设
$\mathcal{E}_i$ 是平坦 $\mathcal{O}_{S_i}$-模；见《代数》引理
\ref{algebra-lemma-flat-finite-presentation-limit-flat}。（使用此引理并非必要，
但较为方便。）于是由《代数》引理
\ref{algebra-lemma-finite-projective}，$\mathcal{E}_i$ 有限局部自由。

\medskip\noindent
如果 $\mathcal{L}$ 是可逆 $\mathcal{O}_S$-模，则由上可找到 $i$ 以及
有限局部自由 $\mathcal{O}_{S_i}$-模 $\mathcal{L}_i$ 和 $\mathcal{N}_i$，
它们分别拉回为 $\mathcal{L}$ 和 $\mathcal{L}^{\otimes -1}$。
必要时增大 $i$ 后，映射
$\mathcal{L} \otimes_{\mathcal{O}_X} \mathcal{L}^{\otimes -1}
\to \mathcal{O}_X$ 下降为映射
$\mathcal{L}_i \otimes_{\mathcal{O}_{S_i}} \mathcal{N}_i \to
\mathcal{O}_{S_i}$。进一步增大 $i$ 后，可以假设它是同构。
于是 $\mathcal{L}_i$ 是可逆模（《模》引理
\ref{modules-lemma-invertible}），(2) 得证。

\medskip\noindent
给定 (3) 中的 $\mathcal{I}$，可见
$\mathcal{O}_S \to \mathcal{O}_S/\mathcal{I}$ 是有限呈示
$\mathcal{O}_S$-模之间的映射。因此由引理
\ref{limits-lemma-descend-modules-finite-presentation}，它是某个映射
$\mathcal{O}_{S_i} \to \mathcal{F}_i$ 的拉回；该映射位于有限呈示
$\mathcal{O}_{S_i}$-模之间。增大 $i$ 后，可以假设
此映射满射（略去细节；提示：在仿射开覆盖上使用《代数》引理
\ref{algebra-lemma-module-map-property-in-colimit}）。于是
$\mathcal{O}_{S_i} \to \mathcal{F}_i$ 的核是 $\mathcal{O}_{S_i}$ 中的
有限型拟相干理想，其拉回给出 $\mathcal{I}$。
\end{proof}

\begin{lemma}
\label{limits-lemma-descend-module-flat-finite-presentation}
记号与假设同引理 \ref{limits-lemma-descend-finite-presentation}。令 $i \in I$。
设 $\varphi_i : X_i \to Y_i$ 是 $S_i$ 上有限呈示概形之间的态射，且
$\mathcal{F}_i$ 是有限呈示拟相干 $\mathcal{O}_{X_i}$-模。
如果 $\mathcal{F}_i$ 到 $X_i \times_{S_i} S$ 的拉回在
$Y_i \times_{S_i} S$ 上平坦，则存在指标 $i' \geq i$，使得
$\mathcal{F}_i$ 到 $X_i \times_{S_i} S_{i'}$ 的拉回在
$Y_i \times_{S_i} S_{i'}$ 上平坦。
\end{lemma}

\begin{proof}
（此引理是引理 \ref{limits-lemma-descend-flat-finite-presentation} 的模版本。）
对 $i' \geq i$，记 $X_{i'} = S_{i'} \times_{S_i} X_i$、
$\mathcal{F}_{i'} = (X_{i'} \to X_i)^*\mathcal{F}_i$；对 $Y_{i'}$
也采用类似记号。用 $\varphi_{i'}$ 表示 $\varphi_i$ 到 $S_{i'}$ 的基变换。
还令 $X = S \times_{S_i} X_i$、$Y =S \times_{S_i} X_i$、
$\mathcal{F} = (X \to X_i)^*\mathcal{F}_i$，并用 $\varphi$ 表示
$\varphi_i$ 到 $S$ 的基变换。取有限仿射开覆盖
$Y_i = \bigcup_{j = 1, \ldots, m} V_{j, i}$，使每个 $V_{j, i}$
都映入 $S_i$ 的某个仿射开集。对每个 $j = 1, \ldots m$，取有限仿射开覆盖
$\varphi_i^{-1}(V_{j, i}) = \bigcup_{k = 1, \ldots, m(j)} U_{k, j, i}$。
对 $i' \geq i$，用 $V_{j, i'}$ 表示 $V_{j, i}$ 在 $Y_{i'}$ 中的逆像，
用 $U_{k, j, i'}$ 表示 $U_{k, j, i}$ 在 $X_{i'}$ 中的逆像。
类似地，有 $U_{k, j} \subset X$ 和 $V_j \subset Y$。于是
$U_{k, j} = \lim_{i' \geq i} U_{k, j, i'}$，且
$V_j = \lim_{i' \geq i} V_j$（见引理
\ref{limits-lemma-directed-inverse-system-has-limit}）。由于
$X_{i'} = \bigcup_{k, j} U_{k, j, i'}$ 是有限开覆盖，只需对每个态射
$U_{k, j, i} \to V_{j, i}$ 及层 $\mathcal{F}_i|_{U_{k, j, i}}$
证明此引理。因此引理约化到 $X_i$ 和 $Y_i$ 均仿射且映入 $S_i$
的一个仿射开集的情形；也就是说，还可假设 $S$ 仿射。

\medskip\noindent
在仿射情形，问题约化为以下代数结果。假设
$R = \colim_{i \in I} R_i$。对某个 $i \in I$，设给定映射
$A_i \to B_i$，其两端是有限呈示 $R_i$-代数。设 $N_i$ 是有限呈示
$B_i$-模。那么，如果 $R \otimes_{R_i} N_i$ 在
$R \otimes_{R_i} A_i$ 上平坦，则对某个 $i' \geq i$，模
$R_{i'} \otimes_{R_i} N_i$ 在 $R_{i'} \otimes_{R_i} A$ 上平坦。
这正是《代数》引理
\ref{algebra-lemma-flat-finite-presentation-limit-flat} (3) 中证明的结果。
\end{proof}

\begin{lemma}
\label{limits-lemma-descend-finite-presentation-variant}
对概形 $T$，用 $\mathcal{C}_T$ 表示概形 $W$ 所成的全子范畴；这些概形
在 $T$ 上，$W$ 拟紧且拟分离，并且结构态射 $W \to T$ 局部有限呈示。
设 $S = \lim S_i$ 是具有仿射过渡态射的概形有向极限。则由基变换函子给出
范畴等价
$$
\colim \mathcal{C}_{S_i} \longrightarrow \mathcal{C}_S
$$
。
\end{lemma}

\noindent
警告：如果不了解此引理与引理 \ref{limits-lemma-descend-finite-presentation}
之间的差别，请勿使用此引理。

\begin{proof}
全忠实性。设有 $i \in I$ 及对象 $X_i$、$Y_i$，它们属于
$\mathcal{C}_{S_i}$。
记 $X = X_i \times_{S_i} S$ 和 $Y = Y_i \times_{S_i} S$。
设给定态射 $f : X \to Y$，它在 $S$ 上。可以选取有限仿射开覆盖
$Y_i = V_{i, 1} \cup \ldots \cup V_{i, m}$，使
$V_{i, j} \to Y_i \to S_i$ 映入某个仿射开集 $W_{i, j}$，后者属于 $S_i$。
用 $Y = V_1 \cup \ldots \cup V_m$ 表示 $Y$ 上诱导出的仿射开覆盖。
由于 $f : X \to Y$ 拟紧（《概形》引理
\ref{schemes-lemma-quasi-compact-permanence}），增大 $i$ 后，可以假设存在
由拟紧开集组成的有限开覆盖
$X_i = U_{i, 1} \cup \ldots \cup U_{i, m}$，使 $U_{i, j}$ 在 $Y$
中的逆像为 $f^{-1}(V_j)$；见引理 \ref{limits-lemma-descend-opens}。
对 $f|_{f^{-1}(V_j)}$ 应用引理 \ref{limits-lemma-descend-finite-presentation}；
这里的底为 $W_j$。增大 $i$ 后，可以假设存在态射
$f_{i, j} : V_{i, j} \to U_{i, j}$，它在 $S$ 上，且其到 $S$ 的基变换为
$f|_{f^{-1}(V_j)}$。进一步增大 $i$ 后，可以假设 $f_{i, j}$ 和
$f_{i, j'}$ 在拟紧开集 $U_{i, j} \cap U_{i, j'}$ 上一致。
于是可以黏合这些态射，得到所需态射 $f_i : X_i \to Y_i$。
此态射唯一（允许增大 $i$），因为态射 $f_{i, j}$ 具有该唯一性。

\medskip\noindent
为证明该函子本质满射，以完全相同的方式论证。具体地，设 $X$ 是
$\mathcal{C}_S$ 的对象。选取 $i \in I$。可以选取有限仿射开覆盖
$X = U_1 \cup \ldots \cup U_m$，使
$U_j \to X \to S \to S_i$ 通过仿射开集 $W_{i, j} \subset S_i$ 分解。
令 $W_j = W_{i, j} \times_{S_i} S$。这是 $S$ 的仿射开集。
由引理 \ref{limits-lemma-descend-finite-presentation}，增大 $i$ 后，可以假设存在
有限呈示态射 $U_{i, j} \to W_{i, j}$，其到 $W_j$ 的基变换为 $U_j$。
增大 $i$ 后，可以假设存在拟紧开集
$U_{i, j, j'} \subset U_{i, j}$，其到 $S$ 的基变换等于
$U_j \cap U_{j'}$。断言：增大 $i$ 后，可以假设态射
$U_{i, j, j'} \to U_{i, j} \to W_{i, j}$ 的像落在
$W_{i, j} \cap W_{i, j'}$ 中。事实上，
$W_{i, j} \cap W_{i, j'}$ 的补集在仿射概形 $W_{i, j}$ 中闭，因而仿射。
由于 $U_j \cap U_{j'} = \lim U_{i, j, j'}$ 的确映入
$W_{i, j} \cap W_{i, j'}$，可以应用引理
\ref{limits-lemma-limit-fibre-product-empty} 得到断言。因此可将二者
$$
U_{i, j, j'} \quad\text{且}\quad U_{i, j', j}
$$
视为 $W_{i, j'}$ 上的概形；它们到 $W_{j'}$ 的基变换都恢复
$U_j \cap U_{j'}$。所以增大 $i$ 后，利用引理
\ref{limits-lemma-descend-finite-presentation}，可以假设存在同构
$U_{i, j, j'} \to U_{i, j', j}$，它在 $W_{i, j'}$ 上，从而也在 $S_i$ 上。
进一步增大 $i$ 后（略去细节），可以假设这些同构满足《概形》第
\ref{schemes-section-glueing-schemes} 节所述的余圈条件。
应用《概形》引理 \ref{schemes-lemma-glue}，得到对象 $X_i$，它属于
$\mathcal{C}_{S_i}$；其到 $S$ 的基变换同构于 $X$，我们略去一些验证。
\end{proof}

\section{仿射概形的刻画}
\label{limits-section-affine}

\noindent
如果 $f : X \to S$ 是概形的满整态射且 $X$ 是仿射概形，则 $S$
也仿射。见 \cite[A.2]{Conrad-Nagata}。我们的证明依赖 Noether 情形；
该情形已在《概形上同调》引理
\ref{coherent-lemma-image-affine-finite-morphism-affine-Noetherian}
中陈述并证明。另见 \cite[II 6.7.1]{EGA}。

\begin{lemma}
\label{limits-lemma-affine}
\begin{slogan}
若一个概形容许从仿射概形出发的有限满射，则它仿射。
\end{slogan}
设 $f : X \to S$ 是概形态射。假设 $f$ 满且有限，并假设 $X$ 仿射。
则 $S$ 仿射。
\end{lemma}

\begin{proof}
由于 $f$ 满且 $X$ 拟紧，可知 $S$ 拟紧。由于 $X$ 分离，而 $f$
满且万有闭（《态射》引理
\ref{morphisms-lemma-integral-universally-closed}），可知 $S$ 分离
（《态射》引理
\ref{morphisms-lemma-image-universally-closed-separated}）。

\medskip\noindent
由引理 \ref{limits-lemma-finite-in-finite-and-finite-presentation}，可写成
$X = \lim_a X_a$，其中 $X_a \to S$ 有限且有限呈示。
由引理 \ref{limits-lemma-limit-affine}，$X_a$ 对某个 $a \in A$ 仿射。
把 $X$ 替换为 $X_a$ 后，可以假设 $X \to S$ 满、有限且有限呈示，
并且 $X$ 仿射。

\medskip\noindent
由命题 \ref{limits-proposition-approximate}，可以把
$S = \lim_{i \in I} S_i$ 写成 $\mathbf{Z}$ 上有限型概形的有向极限。
由引理 \ref{limits-lemma-descend-finite-presentation}，缩小 $I$ 后，可以假设存在
有限呈示概形 $X_i \to S_i$，使
$X_{i'} = X_i \times_S S_{i'}$ 对 $i' \geq i$ 成立，且 $X = \lim_i X_i$。
再由引理 \ref{limits-lemma-descend-finite-finite-presentation}，还可假设
$X_i \to S_i$ 对所有 $i \in I$ 都有限。再次应用引理
\ref{limits-lemma-limit-affine}，可以假设 $X_i$ 对所有 $i \in I$ 都仿射。
因此结果来自 Noether 情形；见《概形上同调》引理
\ref{coherent-lemma-image-affine-finite-morphism-affine-Noetherian}。
\end{proof}

\begin{proposition}
\label{limits-proposition-affine}
\begin{slogan}
若一个概形容许从仿射概形出发的满整态射，则它仿射。
\end{slogan}
设 $f : X \to S$ 是概形态射。假设 $X$ 仿射，且 $f$ 满并万有闭
\footnote{整态射万有闭；见《态射》引理
\ref{morphisms-lemma-integral-universally-closed}。}。则 $S$ 仿射。
\end{proposition}

\begin{proof}
由《态射》引理 \ref{morphisms-lemma-image-universally-closed-separated}，
概形 $S$ 分离。再由《态射》引理
\ref{morphisms-lemma-affine-permanence}，可知 $f$ 仿射。于是由《态射》引理
\ref{morphisms-lemma-integral-universally-closed}，可知 $f$ 整。

\medskip\noindent
由前一段，可以假设 $f : X \to S$ 满且整，$X$ 仿射，并且 $S$ 分离。
由于 $f$ 满且 $X$ 拟紧，还可推出 $S$ 拟紧。

\medskip\noindent
由引理 \ref{limits-lemma-integral-limit-finite-and-finite-presentation}，可写成
$X = \lim_i X_i$，其中 $X_i \to S$ 有限。由引理
\ref{limits-lemma-limit-affine}，对充分大的 $i$，概形 $X_i$ 仿射。
此外，由于 $X \to S$ 通过每个 $X_i$ 分解，可知 $X_i \to S$ 满。
故由引理 \ref{limits-lemma-affine} 可得 $S$ 仿射。
\end{proof}

\begin{lemma}
\label{limits-lemma-affines-glued-in-closed-affine}
设概形 $X$ 在集合论意义下是有限多个仿射闭子概形的并。则 $X$ 仿射。
\end{lemma}

\begin{proof}
设 $Z_i \subset X$（$i = 1, \ldots, n$）是仿射闭子概形，并且在集合论意义下
$X = \bigcup Z_i$。则 $\coprod Z_i \to X$ 满且整，其源仿射。
故由命题 \ref{limits-proposition-affine}，$X$ 仿射。
\end{proof}

\begin{lemma}
\label{limits-lemma-ample-on-reduction}
设 $i : Z \to X$ 是概形的闭浸入，并诱导底拓扑空间之间的同胚。
设 $\mathcal{L}$ 是 $X$ 上的可逆层。则 $i^*\mathcal{L}$ 在 $Z$ 上丰沛，
当且仅当 $\mathcal{L}$ 在 $X$ 上丰沛。
\end{lemma}

\begin{proof}
如果 $\mathcal{L}$ 丰沛，则 $i^*\mathcal{L}$ 丰沛；例如这来自《态射》引理
\ref{morphisms-lemma-pullback-ample-tensor-relatively-ample}。
假设 $i^*\mathcal{L}$ 丰沛。则 $Z$ 拟紧
（《性质》定义 \ref{properties-definition-ample}）且分离
（《性质》引理 \ref{properties-lemma-ample-separated}）。
由于 $i$ 满，可知 $X$ 拟紧。由于 $i$ 万有闭且满，可知 $X$ 分离
（《态射》引理 \ref{morphisms-lemma-image-universally-closed-separated}）。

\medskip\noindent
由命题 \ref{limits-proposition-approximate}，可将
$X = \lim X_i$ 写成 $\mathbf{Z}$ 上有限型概形的有向极限，过渡态射仿射。
可以找到 $i$ 和可逆层 $\mathcal{L}_i$（它在 $X_i$ 上），其到 $X$ 的拉回同构于
$\mathcal{L}$；见引理 \ref{limits-lemma-descend-modules-finite-presentation}。

\medskip\noindent
对每个 $i$，设 $Z_i \subset X_i$ 是态射 $Z \to X_i$ 的概形论像。
如果 $\Spec(A_i) \subset X_i$ 是仿射开子概形，其逆像为 $\Spec(A)$
（逆像取在 $X$ 中），并且 $Z \cap \Spec(A)$ 由理想 $I \subset A$ 定义，
则 $Z_i \cap \Spec(A_i)$ 由理想 $I_i \subset A_i$ 定义；后者是 $I$
在 $A_i$ 中关于环映射 $A_i \to A$ 的逆像。见《态射》例
\ref{morphisms-example-scheme-theoretic-image}。
由于 $\colim A_i/I_i = A/I$，可知 $\lim Z_i = Z$。
由引理 \ref{limits-lemma-limit-ample}，$\mathcal{L}_i|_{Z_i}$ 对某个 $i$ 丰沛。
由于 $Z$，从而 $X$，在集合论意义下映入 $Z_i$，由引理
\ref{limits-lemma-limit-contained-in-constructible}，$X_{i'} \to X_i$
在集合论意义下映入 $Z_i$，对某个 $i' \geq i$ 如此。
（注意，由于 $X_i$ Noether，$X_i$ 的每个闭子集都可构造。）
设 $T \subset X_{i'}$ 是 $Z_i$ 在 $X_{i'}$ 中的概形论逆像。
注意，$\mathcal{L}_{i'}|_T$ 是 $\mathcal{L}_i|_{Z_i}$ 的拉回，故由《态射》引理
\ref{morphisms-lemma-pullback-ample-tensor-relatively-ample} 及
$T \to Z_i$ 是仿射态射这一事实，它丰沛。
于是由《概形上同调》引理 \ref{coherent-lemma-ample-on-reduction}，
$\mathcal{L}_{i'}$ 在 $X_{i'}$ 上丰沛。将其拉回到 $X$
（使用与上面相同的引理），可知 $\mathcal{L}$ 丰沛。
\end{proof}

\begin{lemma}
\label{limits-lemma-thickening-quasi-affine}
设 $i : Z \to X$ 是概形的闭浸入，并诱导底拓扑空间之间的同胚。
则 $X$ 拟仿射，当且仅当 $Z$ 拟仿射。
\end{lemma}

\begin{proof}
回忆，概形拟仿射，当且仅当其结构层丰沛；见《性质》引理
\ref{properties-lemma-quasi-affine-O-ample}。因此，如果 $Z$ 拟仿射，
则 $\mathcal{O}_Z$ 丰沛；由引理 \ref{limits-lemma-ample-on-reduction}，
$\mathcal{O}_X$ 丰沛，因而 $X$ 拟仿射。反过来的证明也可用初等方式看出，
只需逆向阅读刚才的论证。
\end{proof}

\noindent
下面的引理其实不属于本节。

\begin{lemma}
\label{limits-lemma-ample-profinite-set-in-principal-affine}
设 $X$ 是概形。设 $\mathcal{L}$ 是 $X$ 上的丰沛可逆层。
假设有概形态射
$$
\Spec(k) \leftarrow \Spec(A) \to W \subset X
$$
，其中 $k$ 是域，$A$ 是整 $k$-代数，$W$ 在 $X$ 中开。
则存在 $n > 0$ 及截面 $s \in \Gamma(X, \mathcal{L}^{\otimes n})$，
使得 $X_s$ 仿射、$X_s \subset W$，并且 $\Spec(A) \to W$ 通过 $X_s$ 分解。
\end{lemma}

\begin{proof}
由于 $\Spec(A)$ 拟紧，可以把 $W$ 替换为一个仍包含
$\Spec(A) \to X$ 之像的拟紧开集。回忆，由于存在丰沛可逆层，$X$
拟分离且拟紧；见《性质》定义 \ref{properties-definition-ample} 及引理
\ref{properties-lemma-affine-s-opens-cover-quasi-separated}。
由命题 \ref{limits-proposition-approximate}，可将 $X = \lim X_i$ 写成
$\mathbf{Z}$ 上有限型概形的有向系之极限，过渡态射仿射。
对某个 $i$，丰沛可逆层 $\mathcal{L}$（它在 $X$ 上）下降为丰沛可逆层
$\mathcal{L}_i$（它在 $X_i$ 上），而开集 $W$ 是某个拟紧开集
$W_i \subset X_i$ 的逆像；见引理 \ref{limits-lemma-limit-ample}、
\ref{limits-lemma-descend-invertible-modules} 及 \ref{limits-lemma-descend-opens}。
可以把 $X, W, \mathcal{L}$ 替换为 $X_i, W_i, \mathcal{L}_i$，并假设
$X$ 在 $\mathbf{Z}$ 上有限呈示。把 $A = \colim A_j$ 写成其有限
$k$-子代数的余极限。于是对某个 $j$，态射 $\Spec(A) \to X$
通过态射 $\Spec(A_j) \to X$ 分解；见命题
\ref{limits-proposition-characterize-locally-finite-presentation}。
由于 $\Spec(A_j)$ 有限，问题约化为《性质》引理
\ref{properties-lemma-ample-finite-set-in-principal-affine}。
\end{proof}

\section{Chow 引理的变体}
\label{limits-section-chows-lemma}

\noindent
本节证明 Chow 引理的若干变体。最有意义的版本大概就是 Noether 情形；
我们已在《概形上同调》第 \ref{coherent-section-chows-lemma} 节中陈述并证明。

\begin{lemma}
\label{limits-lemma-chow-finite-type}
设 $S$ 是拟紧且拟分离概形。设 $f : X \to S$ 是有限型分离态射。
则存在 $n \geq 0$ 及图
$$
\xymatrix{
X \ar[rd] & X' \ar[d] \ar[l]^\pi \ar[r] & \mathbf{P}^n_S \ar[dl] \\
& S &
}
$$
，其中 $X' \to \mathbf{P}^n_S$ 是浸入，而
$\pi : X' \to X$ 固有且满。
\end{lemma}

\begin{proof}
由命题 \ref{limits-proposition-separated-closed-in-finite-presentation}，可以找到闭浸入
$X \to Y$，其中 $Y$ 在 $S$ 上分离且有限呈示。显然，如果对 $Y$
证明了断言，则 $X$ 的结果也随之成立。因此可以假设 $X$ 在 $S$ 上有限呈示。

\medskip\noindent
将 $S = \lim_i S_i$ 写成 Noether 概形的有向极限；见命题
\ref{limits-proposition-approximate}。由引理 \ref{limits-lemma-descend-finite-presentation}，
可以找到指标 $i \in I$ 和有限呈示概形 $X_i \to S_i$，使得
$X = S \times_{S_i} X_i$。由引理
\ref{limits-lemma-descend-separated-finite-presentation}，可以假设
$X_i \to S_i$ 分离。显然，如果对 $X_i$（它在 $S_i$ 上）证明了断言，
则 $X$ 的断言成立。态射 $X_i \to S_i$ 的情形由《概形上同调》引理
\ref{coherent-lemma-chow-Noetherian} 处理。
\end{proof}

\begin{remark}
\label{limits-remark-chow-finite-type}
在 Chow 引理 \ref{limits-lemma-chow-finite-type} 的情形下：
\begin{enumerate}
\item 态射 $\pi$ 事实上是 H-射影的（因而射影；见《态射》引理
\ref{morphisms-lemma-H-projective}），因为态射
$X' \to \mathbf{P}^n_S \times_S X = \mathbf{P}^n_X$ 是闭浸入
（使用 $\pi$ 固有这一事实；见《态射》引理
\ref{morphisms-lemma-image-proper-scheme-closed}）。
\item 可以假设 $X'$ 既约，因为可将 $X'$ 替换为它的既约化，
而不改变引理的其他断言。
\item 可以假设 $X' \to X$ 有限呈示，而不改变引理的其他断言。
这可从引理 \ref{limits-lemma-chow-finite-type} 的证明推出，也可直接证明如下。
由 (1)，有闭浸入 $X' \to \mathbf{P}^n_X$。由引理
\ref{limits-lemma-closed-is-limit-closed-and-finite-presentation}，可写成
$X' = \lim X'_i$，其中 $X'_i \to \mathbf{P}^n_X$ 是有限呈示闭浸入。
特别地，$X'_i \to X$ 有限呈示、固有且满。由引理
\ref{limits-lemma-finite-type-eventually-closed}，对充分大的 $i$，态射
$X'_i \to \mathbf{P}^n_S$ 是浸入。把 $X'$ 替换为 $X'_i$，即得所需结论。
\end{enumerate}
当然，一般不能同时实现 (2) 和 (3)。
\end{remark}

\noindent
下面是 Chow 引理的一个变体，其中假设上方的概形只有有限多个不可约分支。

\begin{lemma}
\label{limits-lemma-chow-EGA}
设 $S$ 是拟紧且拟分离概形。设 $f : X \to S$ 是有限型分离态射。
假设 $X$ 有有限多个不可约分支。则存在 $n \geq 0$ 及图
$$
\xymatrix{
X \ar[rd] & X' \ar[d] \ar[l]^\pi \ar[r] & \mathbf{P}^n_S \ar[dl] \\
& S &
}
$$
，其中 $X' \to \mathbf{P}^n_S$ 是浸入，而
$\pi : X' \to X$ 固有且满。此外，存在稠密开子概形 $U \subset X$，
使得 $\pi^{-1}(U) \to U$ 是概形同构。
\end{lemma}

\begin{proof}
设 $X = Z_1 \cup \ldots \cup Z_n$ 是 $X$ 的不可约分支分解。
设 $\eta_j \in Z_j$ 为一般点。

\medskip\noindent
证明至少有两种进行方式。第一种是重做《概形上同调》引理
\ref{coherent-lemma-chow-Noetherian} 的证明，使用一般的《性质》引理
\ref{properties-lemma-point-and-maximal-points-affine} 在 $X$ 中寻找合适的
仿射开集。（这是``标准''证明。）第二种是像上面引理
\ref{limits-lemma-chow-finite-type} 的证明那样，使用绝对 Noether 逼近。
这里采用第二种方式。

\medskip\noindent
由命题 \ref{limits-proposition-separated-closed-in-finite-presentation}，可以找到闭浸入
$X \to Y$，其中 $Y$ 在 $S$ 上分离且有限呈示。
将 $S = \lim_i S_i$ 写成 Noether 概形的有向极限；见命题
\ref{limits-proposition-approximate}。由引理 \ref{limits-lemma-descend-finite-presentation}，
可以找到指标 $i \in I$ 和有限呈示概形 $Y_i \to S_i$，使得
$Y = S \times_{S_i} Y_i$。由引理
\ref{limits-lemma-descend-separated-finite-presentation}，可以假设
$Y_i \to S_i$ 分离。有图
$$
\xymatrix{
\eta_j \in Z_j \ar[r] & X \ar[r] \ar[rd] & Y \ar[r] \ar[d] & Y_i \ar[d] \\
& & S \ar[r] & S_i
}
$$
。用 $h : X \to Y_i$ 表示该复合。

\medskip\noindent
对 $i' \geq i$，记 $Y_{i'} = S_{i'} \times_{S_i} Y_i$。
则 $Y = \lim_{i' \geq i} Y_{i'}$；见引理
\ref{limits-lemma-scheme-over-limit}。选取 $j, j' \in \{1, \ldots, n\}$，
$j \not = j'$。注意，$\eta_j$ 不是 $\eta_{j'}$ 的特殊化。
由引理 \ref{limits-lemma-topology-limit}，可将 $i$ 替换为更大的指标，并假设
$h(\eta_j)$ 不是 $h(\eta_{j'})$ 的特殊化，对所有上述指标对 $(j, j')$ 都如此。
对这样的指标，设 $Y' \subset Y_i$ 是 $h : X \to Y_i$ 的概形论像；
见《态射》定义 \ref{morphisms-definition-scheme-theoretic-image}。
态射 $h$ 拟紧，因为它是拟紧态射 $X \to Y$ 和 $Y \to Y_i$
（后者仿射）的复合。因此由《态射》引理
\ref{morphisms-lemma-quasi-compact-scheme-theoretic-image}，
态射 $X \to Y'$ 支配。故 $Y'$ 的一般点全都包含在集合
$\{h(\eta_1), \ldots, h(\eta_n)\}$ 中；见《态射》引理
\ref{morphisms-lemma-quasi-compact-dominant}。
由于没有一个 $h(\eta_j)$ 是另一个的特殊化，可知点
$h(\eta_1), \ldots, h(\eta_n)$ 两两不同，且每个都是 $Y'$ 的一般点。

\medskip\noindent
对态射 $Y' \to S_i$ 应用上面的《概形上同调》引理
\ref{coherent-lemma-chow-Noetherian}。得到图
$$
\xymatrix{
Y' \ar[rd] & Y^* \ar[d] \ar[l]^\pi \ar[r] & \mathbf{P}^n_{S_i} \ar[dl] \\
& S_i &
}
$$
，使得 $\pi$ 固有且满，并且在稠密开子概形 $V \subset Y'$ 上是同构。
由上面对 $i$ 的选取，知道 $h(\eta_1), \ldots, h(\eta_n) \in V$。
考察交换图
$$
\xymatrix{
X' \ar@{=}[r] &
X \times_{Y'} Y^* \ar[r] \ar[d] &
Y^* \ar[r] \ar[d] &
\mathbf{P}^n_{S_i} \ar[ddl] \\
& X \ar[r] \ar[d] & Y' \ar[d] & \\
& S \ar[r] & S_i &
}
$$
。注意，$X' \to X$ 在开子概形 $U = h^{-1}(V)$ 上是同构；后者包含每个
$\eta_j$，因而在 $X$ 中稠密。由此得出
$X \leftarrow X' \rightarrow \mathbf{P}^n_S$ 是引理中问题的一个解。
\end{proof}

\section{Chow 引理的应用}
\label{limits-section-apply-chow}

\noindent
下面是 Chow 引理的第一个应用。

\begin{lemma}
\label{limits-lemma-eventually-proper}
\begin{slogan}
如果一个概形到极限的基变换固有，则它在某个有限层级的基变换已经固有。
\end{slogan}
假设与记号同情形 \ref{limits-situation-descent-property}。如果
\begin{enumerate}
\item $f$ 固有；并且
\item $f_0$ 局部有限型，
\end{enumerate}
则存在 $i$，使得 $f_i$ 固有。
\end{lemma}

\begin{proof}
由引理 \ref{limits-lemma-descend-separated-finite-presentation}，$f_i$ 对某个
$i \geq 0$ 分离。把 $0$ 替换为 $i$ 后，可以假设 $f_0$ 分离。
注意，$f_0$ 拟紧；见《概形》引理
\ref{schemes-lemma-quasi-compact-permanence}。由引理
\ref{limits-lemma-chow-finite-type}，可选取图
$$
\xymatrix{
X_0 \ar[rd] & X_0' \ar[d] \ar[l]^\pi \ar[r] & \mathbf{P}^n_{Y_0} \ar[dl] \\
& Y_0 &
}
$$
，其中 $X_0' \to \mathbf{P}^n_{Y_0}$ 是浸入，而
$\pi : X_0' \to X_0$ 固有且满。引入
$X' = X_0' \times_{Y_0} Y$ 和 $X_i' = X_0' \times_{Y_0} Y_i$。
由《态射》引理 \ref{morphisms-lemma-composition-proper} 及
\ref{morphisms-lemma-base-change-proper}，可知 $X' \to Y$ 固有。
所以 $X' \to \mathbf{P}^n_Y$ 是闭浸入（《态射》引理
\ref{morphisms-lemma-image-proper-scheme-closed}）。由《态射》引理
\ref{morphisms-lemma-image-proper-is-proper}，只需证明 $X'_i \to Y_i$
对某个 $i$ 固有。由引理
\ref{limits-lemma-descend-closed-immersion-finite-presentation}，
$X'_i \to \mathbf{P}^n_{Y_i}$ 对充分大的 $i$ 是闭浸入。
于是 $X'_i \to Y_i$ 固有，证毕。
\end{proof}

\begin{lemma}
\label{limits-lemma-proper-limit-of-proper-finite-presentation}
设 $f : X \to S$ 是固有态射，且 $S$ 拟紧且拟分离。则
$X = \lim X_i$ 是概形的有向极限，其中 $X_i$ 在 $S$ 上固有且有限呈示，
所有过渡态射及态射 $X \to X_i$ 都是闭浸入。
\end{lemma}

\begin{proof}
由命题 \ref{limits-proposition-separated-closed-in-finite-presentation}，可以找到闭浸入
$X \to Y$，其中 $Y$ 在 $S$ 上分离且有限呈示。由引理
\ref{limits-lemma-chow-finite-type}，可以找到图
$$
\xymatrix{
Y \ar[rd] & Y' \ar[d] \ar[l]^\pi \ar[r] & \mathbf{P}^n_S \ar[dl] \\
& S &
}
$$
，其中 $Y' \to \mathbf{P}^n_S$ 是浸入，而
$\pi : Y' \to Y$ 固有且满。由引理
\ref{limits-lemma-closed-is-limit-closed-and-finite-presentation}，可写成
$X = \lim X_i$，其中 $X_i \to Y$ 是有限呈示闭浸入。
用 $X'_i \subset Y'$（相应地，$X' \subset Y'$）表示
$X_i \subset Y$（相应地，$X \subset Y$）的概形论逆像。
则 $\lim X'_i = X'$。由于 $X' \to S$ 固有
（《态射》引理 \ref{morphisms-lemma-composition-proper}），可知
$X' \to \mathbf{P}^n_S$ 是闭浸入（《态射》引理
\ref{morphisms-lemma-image-proper-scheme-closed}）。因此由引理
\ref{limits-lemma-eventually-closed-immersion}，对充分大的 $i$，
$X'_i \to \mathbf{P}^n_S$ 是闭浸入。于是 $X'_i$ 在 $S$ 上固有。
对这样的 $i$，由《态射》引理
\ref{morphisms-lemma-image-proper-is-proper}，态射 $X_i \to S$ 固有。
\end{proof}

\begin{lemma}
\label{limits-lemma-proper-limit-of-proper-finite-presentation-noetherian}
设 $f : X \to S$ 是固有态射，且 $S$ 拟紧且拟分离。则存在有向集 $I$
及概形态射逆系 $(f_i : X_i \to S_i)$，它以 $I$ 为指标，并使得过渡态射
$X_i \to X_{i'}$ 和 $S_i \to S_{i'}$ 仿射，$f_i$ 固有，$S_i$
在 $\mathbf{Z}$ 上有限型，并且
$(X \to S) = \lim (X_i \to S_i)$。
\end{lemma}

\begin{proof}
由引理 \ref{limits-lemma-proper-limit-of-proper-finite-presentation}，可写成
$X = \lim_{k \in K} X_k$，其中 $X_k \to S$ 固有且有限呈示。
接着，由绝对 Noether 逼近（命题 \ref{limits-proposition-approximate}），可写成
$S = \lim_{j \in J} S_j$，其中 $S_j$ 在 $\mathbf{Z}$ 上有限型。
对每个 $k$，存在 $j$ 和有限呈示态射 $X_{k, j} \to S_j$，使
$X_k \cong S \times_{S_j} X_{k, j}$ 是 $S$ 上概形的同构；见引理
\ref{limits-lemma-descend-finite-presentation}。增大 $j$ 后，可以假设
$X_{k, j} \to S_j$ 固有；见引理 \ref{limits-lemma-eventually-proper}。
集合 $I$ 将由这些指标对 $(k, j)$ 组成，而相应的态射为
$X_{k, j} \to S_j$。对每个 $k' \geq k$，可以找到 $j' \geq j$
以及态射 $X_{j', k'} \to X_{j, k}$，它位于 $S_{j'} \to S_j$ 上，
且其到 $S$ 的基变换给出态射
$X_{k'} \to X_k$（再次由引理 \ref{limits-lemma-descend-finite-presentation}）。
这些态射构成该系统的过渡态射。略去一些细节。
\end{proof}

\begin{lemma}
\label{limits-lemma-finite-type-eventually-proper}
设 $S$ 是概形。设 $X = \lim X_i$ 是具有仿射过渡态射的 $S$ 上概形有向极限。
设 $Y \to X$ 是 $S$ 上概形的态射。如果 $Y \to X$ 固有、$X_i$
拟紧且拟分离，并且 $Y$ 在 $S$ 上局部有限型，则 $Y \to X_i$
对充分大的 $i$ 固有。
\end{lemma}

\begin{proof}
选取闭浸入 $Y \to Y'$，其中 $Y'$ 在 $X$ 上固有且有限呈示；见引理
\ref{limits-lemma-proper-limit-of-proper-finite-presentation}。然后选取 $i$ 及固有态射
$Y'_i \to X_i$，使 $Y' = X \times_{X_i} Y'_i$。由引理
\ref{limits-lemma-descend-finite-presentation} 及 \ref{limits-lemma-eventually-proper}，
这是可能的。然后用更大的指标替换 $i$ 后，$Y \to Y'_i$ 是闭浸入；
见引理 \ref{limits-lemma-finite-type-eventually-closed}。
\end{proof}

\noindent
回忆有限型拟相干模的概形论支集；见《态射》定义
\ref{morphisms-definition-scheme-theoretic-support}。

\begin{lemma}
\label{limits-lemma-eventually-proper-support}
假设与记号同情形 \ref{limits-situation-descent-property}。设 $\mathcal{F}_0$
是拟相干 $\mathcal{O}_{X_0}$-模。用 $\mathcal{F}$ 和 $\mathcal{F}_i$
表示 $\mathcal{F}_0$ 到 $X$ 和 $X_i$ 的拉回。假设
\begin{enumerate}
\item $f_0$ 局部有限型；
\item $\mathcal{F}_0$ 有限型；
\item $\mathcal{F}$ 的概形论支集在 $Y$ 上固有。
\end{enumerate}
则 $\mathcal{F}_i$ 的概形论支集在 $Y_i$ 上固有，对某个 $i$ 如此。
\end{lemma}

\begin{proof}
可以把 $X_0$ 替换为 $\mathcal{F}_0$ 的概形论支集。
由《态射》引理 \ref{morphisms-lemma-support-finite-type}，这保证 $X_i$
是 $\mathcal{F}_i$ 的支集，而 $X$ 是 $\mathcal{F}$ 的支集。
于是，若 $Z \subset X$ 表示 $\mathcal{F}$ 的概形论支集，则可知
$Z \to X$ 是泛同胚。由《态射》引理
\ref{morphisms-lemma-image-proper-is-proper}，可得 $X \to Y$ 固有，
因为按假设 $Z \to Y$ 固有。
由引理 \ref{limits-lemma-eventually-proper}，$X_i \to Y$ 对某个 $i$ 固有。
于是由《态射》引理 \ref{morphisms-lemma-closed-immersion-proper} 及
\ref{morphisms-lemma-composition-proper}，概形论支集 $Z_i$（属于
$\mathcal{F}_i$）在 $Y$ 上固有。
\end{proof}

\section{万有闭态射}
\label{limits-section-universally-closed}

\noindent
本节讨论拟紧但未必分离的态射何时万有闭。先证明一个引理，借此可在
局部有限呈示的基变换之后检验万有闭性。

\begin{lemma}
\label{limits-lemma-separate}
设 $f : X \to S$ 是概形的拟紧态射。设 $g : T \to S$ 是概形态射。
设 $t \in T$ 为一点，而 $Z \subset X_T$ 是闭子概形，满足
$Z \cap X_t = \emptyset$。则存在开邻域 $V \subset T$（它包含点 $t$）、交换图
$$
\xymatrix{
V \ar[d] \ar[r]_a & T' \ar[d]^b \\
T \ar[r]^g & S,
}
$$
以及闭子概形 $Z' \subset X_{T'}$，使得
\begin{enumerate}
\item 态射 $b : T' \to S$ 局部有限呈示；
\item 令 $t' = a(t)$，则 $Z' \cap X_{t'} = \emptyset$；并且
\item $Z \cap X_V$ 映入 $Z'$，所经态射为 $X_V \to X_{T'}$。
\end{enumerate}
此外，可以假设 $V$ 和 $T'$ 仿射。
\end{lemma}

\begin{proof}
令 $s = g(t)$。在证明过程中，总可以把 $T$ 替换为 $t$ 的开邻域，
因而也可把 $S$ 替换为 $s$ 的开邻域。因此可以并且确实假设 $T$ 和
$S$ 仿射。记 $S = \Spec(A)$、$T = \Spec(B)$，$g$ 由环映射
$A \to B$ 给出，而 $t$ 对应于素理想 $\mathfrak q \subset B$。

\medskip\noindent
由于 $X \to S$ 拟紧且 $S$ 仿射，可以把
$X = \bigcup_{i = 1, \ldots, n} U_i$ 写成有限多个仿射开集的并。
写 $U_i = \Spec(C_i)$。特别地，
$X_T = \bigcup_{i = 1, \ldots, n} U_{i, T} =
\bigcup_{i = 1, \ldots n} \Spec(C_i \otimes_A B)$。
设 $I_i \subset C_i \otimes_A B$ 是对应于闭子概形
$Z \cap U_{i, T}$ 的理想。条件 $Z \cap X_t = \emptyset$ 表示 $I_i$
在环
$$
C_i \otimes_A \kappa(\mathfrak q) =
(B \setminus \mathfrak q)^{-1}\left(
C_i \otimes_A B/\mathfrak q C_i \otimes_A B \right)
$$
中生成单位理想。由于
$I_i (B \setminus \mathfrak q)^{-1}(C_i \otimes_A B) =
(B \setminus \mathfrak q)^{-1} I_i$，这意味着
$1 = x_i/g_i$，其中某个 $x_i \in I_i$ 和 $g_i \in B$ 满足
$g_i \not \in \mathfrak q$。所以清除分母后，可以找到形如
$$
x_i + \sum\nolimits_j f_{i, j}c_{i, j} = g_i
$$
的关系，其中 $x_i \in I_i$、$f_{i, j} \in \mathfrak q$、
$c_{i, j} \in C_i \otimes_A B$，并且 $g_i \in B$、
$g_i \not \in \mathfrak q$。把 $B$ 替换为 $B_{g_1 \ldots g_n}$，
即把 $T$ 替换为 $t$ 的更小仿射邻域后，可以假设方程为
$$
x_i + \sum\nolimits_j f_{i, j}c_{i, j} = 1
$$
，其中 $x_i \in I_i$、$f_{i, j} \in \mathfrak q$、
$c_{i, j} \in C_i \otimes_A B$。

\medskip\noindent
为完成论证，把 $B$ 写成有限呈示 $A$-代数 $B_\lambda$ 在有向集
$\Lambda$ 上的余极限。对每个 $\lambda$，令
$\mathfrak q_\lambda = (B_\lambda \to B)^{-1}(\mathfrak q)$。
对充分大的 $\lambda \in \Lambda$，可以找到
\begin{enumerate}
\item 元素 $x_{i, \lambda} \in C_i \otimes_A B_\lambda$，它映为 $x_i$；
\item 元素 $f_{i, j, \lambda} \in \mathfrak q_{i, \lambda}$，它们映为
$f_{i, j}$；并且
\item 元素 $c_{i, j, \lambda} \in C_i \otimes_A B_\lambda$，它们映为
$c_{i, j}$。
\end{enumerate}
再增大一点 $\lambda$ 后，方程
$$
x_{i, \lambda} + \sum\nolimits_j f_{i, j, \lambda}c_{i, j, \lambda} = 1
$$
成立。固定这样的 $\lambda$，并令 $T' = \Spec(B_\lambda)$。
则 $t' \in T'$ 是对应于素理想 $\mathfrak q_\lambda$ 的点。
最后，设 $Z' \subset X_{T'}$ 为 $Z \to X_T \to X_{T'}$ 的概形论像。
由于 $X_T \to X_{T'}$ 仿射，可以把 $Z'$ 在仿射开片 $U_{i, T'}$ 上
计算为对应于
$\Ker(C_i \otimes_A B_\lambda \to C_i \otimes_A B/I_i)$ 的闭子概形；
见《态射》例 \ref{morphisms-example-scheme-theoretic-image}。
因此 $x_{i, \lambda}$ 属于定义 $Z'$ 的理想。故最后一个显示方程表明
$Z' \cap X_{t'}$ 为空。
\end{proof}

\begin{lemma}
\label{limits-lemma-test-universally-closed}
设 $f : X \to S$ 是概形的拟紧态射。下列条件等价：
\begin{enumerate}
\item $f$ 万有闭；
\item 对每个局部有限呈示的态射 $S' \to S$，基变换
$X_{S'} \to S'$ 都是闭映射；并且
\item 对每个 $n$，态射
$\mathbf{A}^n \times X \to \mathbf{A}^n \times S$
都是闭映射。
\end{enumerate}
\end{lemma}

\begin{proof}
显然 (1) 蕴含 (2)。下面证明 (2) 蕴含 (1)。假设基变换 $X_T \to T$
对某个概形 $T$（它在 $S$ 上）不是闭映射。由《概形》引理
\ref{schemes-lemma-quasi-compact-closed}，这意味着存在某个特化
$t_1 \leadsto t$ 于 $T$ 中，以及点 $\xi \in X_T$，它映到 $t_1$，使得 $\xi$ 不能特化到
$t$ 上纤维中的点。令 $Z = \overline{\{\xi\}} \subset X_T$。则
$Z \cap X_t = \emptyset$。应用引理 \ref{limits-lemma-separate}，得到开邻域
$V \subset T$（它包含点 $t$）、交换图
$$
\xymatrix{
V \ar[d] \ar[r]_a & T' \ar[d]^b \\
T \ar[r]^g & S,
}
$$
以及闭子概形 $Z' \subset X_{T'}$，使得
\begin{enumerate}
\item 态射 $b : T' \to S$ 局部有限呈示；
\item 令 $t' = a(t)$，则 $Z' \cap X_{t'} = \emptyset$；并且
\item $Z \cap X_V$ 映入 $Z'$，所经态射为 $X_V \to X_{T'}$。
\end{enumerate}
这显然意味着 $X_{T'} \to T'$ 把闭子集 $Z'$ 映到 $T'$ 的一个子集，
后者包含 $a(t_1)$，却不包含 $t' = a(t)$。由于
$a(t_1) \leadsto a(t) = t'$，可知 $X_{T'} \to T'$ 不是闭映射。因此已经证明：
$X \to S$ 非万有闭蕴含 $X_{T'} \to T'$ 对某个局部有限呈示的
$T' \to S$ 不是闭映射。换言之，(2) 蕴含 (1)。

\medskip\noindent
假设 $\mathbf{A}^n \times X \to \mathbf{A}^n \times S$ 对每个整数 $n$
都是闭映射。要证明：$X_T \to T$ 对每个概形 $T$（它在 $S$ 上局部有限呈示）
都是闭映射。当然，可以假设 $T$ 仿射，并且映入一个仿射开集 $V$，后者是
$S$ 的开集（因为 $X_T \to T$ 为闭映射是 $T$ 上的局部性质）。此时存在
闭浸入 $T \to \mathbf{A}^n \times V$，因为
$\mathcal{O}_T(T)$ 是有限呈示 $\mathcal{O}_S(V)$-代数；见《态射》引理
\ref{morphisms-lemma-locally-finite-presentation-characterize}。
于是 $T \to \mathbf{A}^n \times S$ 是局部闭浸入。因此得到概形的笛卡尔图
$$
\xymatrix{
X_T \ar[d]_{f_T} \ar[r] & \mathbf{A}^n \times X \ar[d]^{f_n} \\
T \ar[r] & \mathbf{A}^n \times S
}
$$
，其中水平箭头都是局部闭浸入。因此任意闭子集 $Z \subset X_T$ 都可写成
$X_T \cap Z'$，其中 $Z' \subset \mathbf{A}^n \times X$ 为闭子集。
于是 $f_T(Z) = T \cap f_n(Z')$；可见若 $f_n$ 是闭映射，则 $f_T$ 也是闭映射。
\end{proof}

\begin{lemma}
\label{limits-lemma-limited-base-change}
设 $S$ 是概形。设 $f : X \to S$ 是分离有限型态射。下列条件等价：
\begin{enumerate}
\item 态射 $f$ 固有。
\item 对任意局部有限型的态射 $S' \to S$，基变换
$X_{S'} \to S'$ 是闭映射。
\item 对每个 $n \geq 0$，态射
$\mathbf{A}^n \times X \to \mathbf{A}^n \times S$ 是闭映射。
\end{enumerate}
\end{lemma}

\begin{proof}[第一种证明]
固有态射恰好就是分离、有限型且万有闭的态射；据此，本引理是引理
\ref{limits-lemma-test-universally-closed} 的一个特例。
\end{proof}

\begin{proof}[第二种证明]
显然 (1) 蕴含 (2)，而 (2) 蕴含 (3)，故只需证明 (3) 蕴含 (1)。
先归约到 $S$ 仿射的情形。假设当基概形仿射时 (3) 蕴含 (1)。现在设
$f: X \to S$ 是分离有限型态射。固有性在基概形上是局部的（见《态射》引理
\ref{morphisms-lemma-proper-local-on-the-base}）；因此，若
$S = \bigcup_\alpha S_\alpha$ 是仿射开覆盖，并记
$X_\alpha := f^{-1}(S_\alpha)$，则只需证明
$f|_{X_\alpha}: X_\alpha \to S_\alpha$ 对所有 $\alpha$ 都固有。由于 $S_\alpha$ 仿射，若映射
$f|_{X_\alpha}$ 满足 (3)，则依假设它满足 (1)，从而固有。为完成向
$S$ 仿射情形的归约，必须证明：若分离有限型态射 $f: X \to S$ 满足 (3)，
则 $f|_{X_\alpha} : X_\alpha \to S_\alpha$ 是满足 (3) 的分离有限型态射。
分离性和有限型性是显然的。为验证 (3)，注意到
$\mathbf{A}^n \times X_\alpha$ 是 $\mathbf{A}^n \times S_\alpha$ 在映射
$1 \times f$ 下的开逆像。固定闭集
$Z \subset \mathbf A^n \times X_\alpha$。令 $\bar Z$ 为 $Z$ 在
$\mathbf{A}^n \times X$ 中的闭包。则由拓扑理由，
$$
1 \times f(\bar Z) \cap \mathbf{A}^n \times S_\alpha  = 1 \times f(Z).
$$
所以 $1 \times f(Z)$ 闭，而 (3) $\Rightarrow$ (1) 的证明已经归约到仿射情形。

\medskip\noindent
假设 $S$ 仿射，而 $f : X \to S$ 分离且有限型。应用 Chow 引理
\ref{limits-lemma-chow-finite-type}，得到固有满射 $\pi : X' \to X$，以及浸入
$X' \to \mathbf{P}^n_S$。若 $X$ 在 $S$ 上固有，则 $X' \to S$ 固有
（《态射》引理 \ref{morphisms-lemma-composition-proper}）。由于
$\mathbf{P}^n_S \to S$ 分离，可知 $X' \to
\mathbf{P}^n_S$ 固有（《态射》引理
\ref{morphisms-lemma-image-proper-scheme-closed}），因而是闭浸入
（《概形》引理 \ref{schemes-lemma-immersion-when-closed}）。反之，假设
$X' \to \mathbf{P}^n_S$ 是闭浸入。考虑图：
\begin{equation}
\label{limits-equation-check-proper}
\xymatrix{
X' \ar[r] \ar@{->>}[d]_{\pi} &
\mathbf{P}^n_S \ar[d] \\
X \ar[r]^f & S
}
\end{equation}
除 $X \to S$ 外，所有映射先验地都是固有的。因而由《态射》引理
\ref{morphisms-lemma-image-proper-is-proper} 可知 $X \to S$ 固有。
所以已经证明：$X \to S$ 固有，当且仅当
$X' \to \mathbf{P}^n_S$ 是闭浸入。

\medskip\noindent
假设 $S$ 仿射且 (3) 成立，并取如上的 $n, X', \pi$。由于一个态射为闭映射
在基概形上是局部的，映射 $X \times \mathbf{P}^n \to S \times \mathbf{P}^n$
是闭映射：由 (3)，$X \times \mathbf{A}^n \to S \times \mathbf{A}^n$
是闭映射，而射影空间由仿射 $n$-空间的若干副本覆盖；见《构造》引理
\ref{constructions-lemma-standard-covering-projective-space}。由《态射》引理
\ref{morphisms-lemma-base-change-proper}，态射
$$
X' \times_S \mathbf{P}^n_S
\to
X \times_S \mathbf{P}^n_S =
X \times \mathbf{P}^n
$$
固有。由于 $\mathbf{P}^n$ 分离，投影
$$
X' \times_S \mathbf{P}^n_S = \mathbf{P}^n_{X'} \to X'
$$
是分离的，因为它只是一个分离态射的基变换。因此映射
$X' \to X' \times_S \mathbf{P}^n_S$ 固有，因为它是一个分离映射的截面
（见《概形》引理 \ref{schemes-lemma-section-immersion}）。复合这些态射
$$
X' \to X' \times_S \mathbf{P}^n_S \to X \times_S \mathbf{P}^n_S
= X \times \mathbf{P}^n \to S \times \mathbf{P}^n = \mathbf{P}^n_S
$$
，便发现浸入 $X' \to \mathbf{P}^n_S$ 是闭映射，因而是闭浸入。
\end{proof}

\section{Noether 型赋值判据}
\label{limits-section-Noetherian-valuative-criterion}

\noindent
若基概形为 Noether 概形，则只用离散赋值环便可证明赋值判据成立。

\medskip\noindent
本节的许多结果都可以（或许也应该）借助下述引理证明，尽管我们并非
总是这样做。

\begin{lemma}
\label{limits-lemma-reach-point-closure-Noetherian}
设 $f : X \to Y$ 是概形态射。假设 $f$ 有限型且 $Y$ 局部 Noether。
设 $y \in Y$ 是 $f$ 的像的闭包中的一点。则存在交换图
$$
\xymatrix{
\Spec(K) \ar[r] \ar[d] & X \ar[d]^f \\
\Spec(A) \ar[r] & Y
}
$$
，其中 $A$ 是离散赋值环，而 $K$ 是它的分式域；该图把 $\Spec(A)$ 的闭点
映到 $y$。此外，可以假设 $\Spec(K) \to X$ 的像点是某个一般点 $\eta$，
它属于 $X$ 的一个不可约分支，且 $K = \kappa(\eta)$。
\end{lemma}

\begin{proof}
由本引理的非 Noether 版本（《态射》引理
\ref{morphisms-lemma-reach-points-scheme-theoretic-image}），存在一点
$x \in X$，使 $f(x)$ 特化到 $y$。可以用任意特化到 $x$ 的点替换 $x$，
因而可假设 $x$ 是 $X$ 的某个不可约分支的一般点。这给出环映射
$\mathcal{O}_{Y, y} \to \kappa(x)$（见《概形》第
\ref{schemes-section-points} 节）。令 $R \subset \kappa(x)$ 为其像。
则 $R$ 是 Noether 局部环 $\mathcal{O}_{Y, y}$ 的商，故为 Noether 环。
另一方面，$\kappa(x)$ 是 $R$ 的分式域的有限生成扩张，因为 $f$ 有限型。
因此，由《代数》引理 \ref{algebra-lemma-exists-dvr}，存在离散赋值环
$A \subset \kappa(x)$，其分式域为 $\kappa(x)$，且支配 $R$。于是
$$
\xymatrix{
\Spec(\kappa(x)) \ar[d] \ar[rrr] & & & X \ar[d] \\
\Spec(A) \ar[r] & \Spec(R) \ar[r] & \Spec(\mathcal{O}_{Y, y}) \ar[r] & Y
}
$$
给出所需的图。
\end{proof}

\noindent
先陈述关于分离性的结果。我们将经常使用如下形状的、由概形态射构成的实线交换图
\begin{equation}
\label{limits-equation-valuative}
\vcenter{
\xymatrix{
\Spec(K) \ar[r] \ar[d] & X \ar[d] \\
\Spec(A) \ar[r] \ar@{-->}[ru] & S
}
}
\end{equation}
，其中 $A$ 是赋值环，而 $K$ 是它的分式域。

\begin{lemma}
\label{limits-lemma-Noetherian-dvr-valuative-separation}
设 $S$ 是局部 Noether 概形。设 $f : X \to S$ 是概形态射。假设 $f$
局部有限型。下列条件等价：
\begin{enumerate}
\item 态射 $f$ 分离。
\item 对任意图 (\ref{limits-equation-valuative})，虚线箭头至多有一条。
\item 对所有图 (\ref{limits-equation-valuative})，若 $A$ 是离散赋值环，则虚线箭头
至多有一条。
\item 对任意不可约分支 $X_0$（它属于 $X$）及其一般点 $\eta \in X_0$，
对任意离散赋值环 $A \subset K = \kappa(\eta)$（其分式域为 $K$），以及
任意图 (\ref{limits-equation-valuative})，若态射 $\Spec(K) \to X$ 是典范态射
（见《概形》第 \ref{schemes-section-points} 节），则虚线箭头至多有一条。
\end{enumerate}
\end{lemma}

\begin{proof}
显然 (1) 蕴含 (2)，(2) 蕴含 (3)，而 (3) 蕴含 (4)。只须证明 (4)
蕴含 (1)。假设 (4)。

先归约到 $S$ 仿射的情形。分离性在基概形上是局部的（见《概形》引理
\ref{schemes-lemma-characterize-separated}）。因此，只要能证明：每当
$X \to S$ 满足 (4) 时，限制 $X_\alpha \to S_\alpha$ 也满足 (4)，其中
$S_\alpha \subset S$ 是开子集且 $X_\alpha :=
f^{-1}(S_\alpha)$，便得证。$X_\alpha$ 的不可约分支的一般点将是 $X$ 的不可约分支的一般点，
因为 $X_\alpha$ 在 $X$ 中开。因此，图
\begin{equation}
\label{limits-equation-valuative-alpha}
\xymatrix{
\Spec(K) \ar[r] \ar[d] & X_\alpha \ar[d] \\
\Spec(A) \ar[r] \ar@{-->}[ru] & S_\alpha
}
\end{equation}
中的任意两条不同虚线箭头，经由映射 $X_\alpha \to X$ 和
$S_\alpha \to S$，便会在图 (\ref{limits-equation-valuative}) 中给出两条不同箭头，
这与假设矛盾。于是已经归约到 $S$ 仿射的情形。顺便指出，在这一归约中，
我们证明了：若 $X \to S$ 满足 (4)，则限制 $U
\to V$ 也满足 (4)，其中 $U \subset X$ 和 $V \subset S$ 是满足
$f(U) \subset V$ 的开集。

\medskip\noindent
其次要归约到 $X \to S$ 有限型的情形。假设已经知道当 $X$ 在该基上有限型时
(4) 蕴含 (1)。由于 $S$ 是 Noether 概形，而 $X$ 在 $S$ 上局部有限型，
可知 $X$ 也局部 Noether（见《态射》引理
\ref{morphisms-lemma-finite-type-noetherian}）。因此 $X \to S$ 拟分离
（见《性质》引理 \ref{properties-lemma-locally-Noetherian-quasi-separated}），
故可用赋值判据检验 $X$ 是否分离（见《概形》引理
\ref{schemes-lemma-valuative-criterion-separatedness}）。令
$X = \bigcup_\alpha X_\alpha$ 是 $X$ 的仿射开覆盖。在图
(\ref{limits-equation-valuative}) 中给定任意两条虚线箭头，$\Spec A$ 的闭点的像
将分别落入两个集合 $X_\alpha$ 和 $X_\beta$。由于
$X_\alpha \cup
X_\beta$ 是开集，出于拓扑理由，它必定包含两个映射下
$\Spec(A)$ 的像。因此两条虚线箭头都分解通过
$X_\alpha \cup X_\beta \to X$，而这个概形在 $S$ 上有限型。由于
$X_\alpha \cup X_\beta$ 是 $X$ 的开子集，依前面的说明，
$X_\alpha \cup X_\beta$ 满足 (4)，故依假设它分离。这说明给定的两条
虚线箭头相同。因此已经归约到 $X \to S$ 有限型的情形。

\medskip\noindent
假设 $X \to S$ 有限型并满足 (4)。由于 $X \to S$ 有限型，而 $S$ 是仿射
Noether 概形，$X$ 也为 Noether（见《态射》引理
\ref{morphisms-lemma-finite-type-noetherian}）。因此
$X \to X \times_S X$ 是 Noether 概形间的拟紧浸入。用反证法。假设
$X \to X \times_S X$ 不闭，则存在某个
$y \in X \times_S X$，它属于像的闭包但不属于像。由于 $X$ 是 Noether
概形，它只有有限多个不可约分支。因此 $y$ 属于某个不可约分支
$X_0 \subset X$ 的像的闭包。赋予 $X_0$ 既约诱导结构。复合态射
$X_0 \to X \to X \times_S X$ 分解通过闭子概形
$X_0 \times_S X_0 \subset X \times_S X$。将 $\Delta(X_0)$ 在
$X_0 \times_S X_0$ 中的闭包记为 $\bar X_0$（仍取既约闭子概形）。于是
$y \in \bar X_0$。由于 $X_0 \to X_0 \times_S X_0$ 是浸入，$X_0$ 的像
在 $\bar X_0$ 中开。因此 $X_0$ 与 $\bar X_0$ 双有理。由于
$\bar{X}_0$ 是 Noether 概形的闭子概形，它是 Noether 概形。所以局部环
$\mathcal O_{{\bar X_0, y}}$ 是局部 Noether 整环，其分式域 $K$ 等于
$X_0$ 的函数域。由 Krull--Akizuki 定理（见《代数》引理
\ref{algebra-lemma-exists-dvr}），存在离散赋值环 $A$，它支配
$\mathcal O_{{\bar X_0, y}}$，且分式域为 $K$。这使我们能构造图
\begin{equation}
\label{limits-equation-valuative-generic}
\xymatrix{
\Spec(K) \ar[r] \ar[d] & X_0 \ar[d]^{\Delta} \\
\Spec(A) \ar[r] \ar@{-->}[ur]& X_0 \times_S X_0 \\
}
\end{equation}
，它把 $\Spec K$ 映到 $\Delta(X_0)$ 的一般点，并把 $A$ 的闭点映到
$y \in X_0 \times_S X_0$（使用《概形》第
\ref{schemes-section-points} 节的材料构造这些箭头）。甚至不存在集合论意义的
虚线箭头，因为 $y$ 不在 $\Delta$ 的像中，这是由 $y$ 的选取所保证的。由范畴方法，
上图中存在虚线箭头，等价于下图中虚线箭头的唯一性：
\begin{equation}
\label{limits-equation-valuative-nonexistent}
\xymatrix{
\Spec(K) \ar[r] \ar[d] & X_0 \ar[d]\\
\Spec(A) \ar[r] \ar@{-->}[ur] & S \\
}
\end{equation}
所以，第一幅图中不存在虚线箭头，便使后一幅图不满足唯一性。因此 $X_0$
不满足离散赋值环的唯一性；又因为 $X_0$ 是 $X$ 的不可约分支，所以
$X \to S$ 不满足 (4)。由此证明了 (4) 蕴含 (1)。
\end{proof}

\begin{lemma}
\label{limits-lemma-Noetherian-dvr-valuative-proper}
设 $S$ 是局部 Noether 概形。设 $f : X \to S$ 是有限型态射。下列条件等价：
\begin{enumerate}
\item 态射 $f$ 固有。
\item 对任意图 (\ref{limits-equation-valuative})，正好存在一条虚线箭头。
\item 对所有图 (\ref{limits-equation-valuative})，若 $A$ 是离散赋值环，则正好存在
一条虚线箭头。
\item 对任意不可约分支 $X_0$（它属于 $X$）及其一般点 $\eta \in X_0$，
对任意离散赋值环 $A \subset K = \kappa(\eta)$（其分式域为 $K$），以及
任意图 (\ref{limits-equation-valuative})，若态射 $\Spec(K) \to X$ 是典范态射
（见《概形》第 \ref{schemes-section-points} 节），则正好存在一条虚线箭头。
\end{enumerate}
\end{lemma}

\begin{proof}
(1) 蕴含 (2)，(2) 蕴含 (3)，而 (3) 蕴含 (4)。现在证明 (4) 蕴含 (1)。
如引理 \ref{limits-lemma-Noetherian-dvr-valuative-separation} 的证明所示，可以归约到
$S$ 仿射的情形，因为固有性在基概形上是局部的；而若 $X
\to S$ 满足 (4)，则 $X_\alpha \to S_\alpha$ 也满足 (4)，其中
$S_\alpha \subset S$ 是开集且 $X_\alpha = f^{-1}(S_\alpha)$。

\medskip\noindent
现在 $S$ 是 Noether 概形，所以 $X$ 也是 Noether 概形，因为 $X \to
S$ 有限型。
现在可使用 Chow 引理（《概形上同调》引理
\ref{coherent-lemma-chow-Noetherian}），得到满射、固有、双有理态射
$X' \to X$ 以及浸入 $X' \to \mathbf{P}^n_S$。我们要证明
$X \to S$ 万有闭。如引理 \ref{limits-lemma-limited-base-change} 的证明所示，
只须检验 $X' \to \mathbf{P}^n_S$ 是闭浸入。
用反证法，假设 $X' \to
\mathbf{P}^n_S$ 不是闭浸入。则存在某个 $y
\in \mathbf{P}^n_S$，它在 $X'$ 的像的闭包中，却不在像中。因此 $y$ 在
$X_0'$ 的像的闭包中，其中该分支是 $X'$ 的某个不可约分支；但该点不在像中。令
$\bar X_0' \subset \mathbf{P}^n_S$ 是 $X_0'$ 的像的闭包。由于
$X' \to \mathbf{P}^n_S$ 是 Noether 概形间的浸入，态射
$X'_0 \to \bar X_0'$ 开且稠密。由《代数》引理
\ref{algebra-lemma-exists-dvr} 或《性质》引理
\ref{properties-lemma-locally-Noetherian-specialization-dvr}，可以找到离散赋值环
$A$，它支配 $\mathcal{O}_{\bar X_0', y}$ 且具有相同的分式域 $K$。
显然 $K$ 是 $X_0'$ 的一般点处的剩余域。由此得到实线交换图
\begin{equation}
\label{limits-equation-solid}
\xymatrix{
\Spec K \ar[r] \ar[d] & X' \ar [r] \ar[d] &
\mathbf{P}^n_S \ar[d] \\
\Spec A \ar@{-->}[r] \ar@{-->}[ru] \ar[urr] & X \ar[r] & S\\
}
\end{equation}
注意 $A$ 的闭点映到 $y \in \mathbf{P}^n_S$。依构造，不存在集合论意义的
到 $X'$ 的提升。由于 $X' \to X$ 双有理，$X'_0$ 在 $X$ 中的像是一个不可约
分支 $X_0$，它属于 $X$，而 $K$ 也等同于 $X_0$ 的函数域。于是，因为假设
$X \to S$ 满足 (4)，虚线箭头 $\Spec(A) \to X$ 存在。由于
$X' \to X$ 固有，这条虚线箭头提升为虚线箭头 $\Spec(A) \to X'$（使用
《概形》命题 \ref{schemes-proposition-characterize-universally-closed}）。
将它与浸入 $X' \to \mathbf{P}^n_S$ 复合，得到另一个未在图中绘出的态射
$\Spec(A) \to \mathbf{P}^n_S$。由于 $\mathbf{P}^n_S$ 在 $S$ 上固有，
它满足 (2)，故这两个态射相同。这与我们已经构造出原映射
$\Spec(A) \to \mathbf{P}^n_S$ 到 $X'$ 的、原本不可能存在的提升矛盾。
\end{proof}

\begin{lemma}
\label{limits-lemma-check-universally-closed-Noetherian}
设 $f : X \to S$ 是概形的有限型态射。假设 $S$ 局部 Noether。则下列条件
等价：
\begin{enumerate}
\item $f$ 万有闭；
\item 对每个 $n$，态射
$\mathbf{A}^n \times X \to \mathbf{A}^n \times S$ 是闭映射；
\item 对任意图 (\ref{limits-equation-valuative})，存在某条虚线箭头；
\item 对所有图 (\ref{limits-equation-valuative})，若 $A$ 是离散赋值环，则存在某条
虚线箭头。
\end{enumerate}
\end{lemma}

\begin{proof}
(1) 与 (2) 的等价性是引理 \ref{limits-lemma-test-universally-closed} 的特例。
(1) 与 (3) 的等价性是《概形》命题
\ref{schemes-proposition-characterize-universally-closed} 的特例。显然 (3)
蕴含 (4)。所以只须证明 (4) 蕴含 (2)。我们将用《概形》引理
\ref{schemes-lemma-quasi-compact-closed} 的判据证明
$\mathbf{A}^n \times X \to \mathbf{A}^n \times S$ 是闭映射。
选取 $n$，再选取一个非平凡特化 $z \leadsto z'$，其两点均位于
$\mathbf{A}^n \times S$ 中。再选取一点
$y \in \mathbf{A}^n \times X$，它位于 $z$ 上方。注意 $\kappa(y)$ 是 $\kappa(z)$ 的
有限生成域扩张，因为 $\mathbf{A}^n \times X \to \mathbf{A}^n \times S$
有限型。因此，由《性质》引理
\ref{properties-lemma-locally-Noetherian-specialization-dvr} 或《代数》引理
\ref{algebra-lemma-exists-dvr}，存在离散赋值环 $A \subset \kappa(y)$，
其分式域为 $\kappa(z)$，并支配 $\mathcal{O}_{\mathbf{A}^n \times S, z'}$
在 $\kappa(z)$ 中的像。这给出交换图
$$
\xymatrix{
\Spec(\kappa(y)) \ar[r] \ar[d] &
\mathbf{A}^n \times X \ar[d] \ar[r] & X \ar[d] \\
\Spec(A) \ar[r] & \mathbf{A}^n \times S \ar[r] & S
}
$$
现在性质 (4) 蕴含存在使此图交换的态射 $\Spec(A) \to X$。由于从左下水平箭头
已经得到态射 $\Spec(A) \to \mathbf{A}^n$，也就得到使左方正方形交换的态射
$\Spec(A) \to \mathbf{A}^n \times X$。因此闭点的像
$y' \in \mathbf{A}^n \times X$ 是 $y$ 的一个特化，并位于 $z'$ 上方。
这证明特化可沿 $\mathbf{A}^n \times X \to \mathbf{A}^n \times S$
提升，从而得证。
\end{proof}

\section{精细化的 Noether 型赋值判据}
\label{limits-section-refined-valuative-criteria}

\noindent
检验赋值判据时，通常不必考虑赋值环所给出的所有可能图。《态射》引理
\ref{morphisms-lemma-refined-valuative-criterion-universally-closed}
给出了一个例子。在 Noether 情形，我们也已在引理
\ref{limits-lemma-Noetherian-dvr-valuative-separation} 和
\ref{limits-lemma-Noetherian-dvr-valuative-proper} 中见过这一点。下面是另一种变体。

\begin{lemma}
\label{limits-lemma-refined-valuative-criterion-proper}
设 $f : X \to S$ 和 $h : U \to X$ 是概形态射。假设 $S$ 局部 Noether，
$f$ 和 $h$ 有限型，$f$ 分离，并且 $h(U)$ 在 $X$ 中稠密。若对任意交换实线图
$$
\xymatrix{
\Spec(K) \ar[r] \ar[d] & U \ar[r]^h & X \ar[d]^f \\
\Spec(A) \ar[rr] \ar@{-->}[rru] & & S
}
$$
，其中 $A$ 是以 $K$ 为分式域的离散赋值环，都存在一条使该图交换的虚线箭头，
则 $f$ 固有。
\end{lemma}

\begin{proof}
可以立即归约到 $S$ 仿射的情形。此时 $U$ 拟紧。令
$U = U_1 \cup \ldots \cup U_n$ 为仿射开覆盖。将 $U$ 替换为
$U_1 \amalg \ldots \amalg U_n$ 不会改变假设，故可假设 $U$ 仿射。
于是可以找到开浸入 $U \to Y$，它位于 $X$ 上，并使 $Y$ 在 $X$ 上固有。（先用《态射》
引理 \ref{morphisms-lemma-quasi-affine-finite-type-over-S} 将 $U$ 放入
$\mathbf{A}^n_X$，再在 $\mathbf{P}^n_X$ 中取闭包；或者直接使用《态射》引理
\ref{morphisms-lemma-quasi-projective-open-projective}。）可以假设 $U$ 在 $Y$
中稠密（必要时将 $Y$ 替换为 $U$ 的概形论闭包；见《态射》第
\ref{morphisms-section-scheme-theoretic-closure} 节）。注意
$g : Y \to X$ 满射，因为它的像闭且包含稠密子集 $h(U)$。
我们将证明 $Y \to S$ 固有。由《态射》引理
\ref{morphisms-lemma-image-proper-is-proper}，这将蕴含 $X \to S$ 固有，
从而完成证明。为证明 $Y \to S$ 固有，我们将使用引理
\ref{limits-lemma-Noetherian-dvr-valuative-proper} 的第 (4) 部分。为此考虑图
$$
\xymatrix{
\Spec(K) \ar[r]_y \ar[d] & Y \ar[d]^{f \circ g} \\
\Spec(A) \ar[r] \ar@{..>}[ru] & S
}
$$
，其中 $A$ 是以 $K$ 为分式域的离散赋值环，而
$y : \Spec(K) \to Y$ 是某个一般点的包含映射。必须证明正好存在一条虚线箭头。
唯一性来自分离性赋值判据的逆向结论（《概形》引理
\ref{schemes-lemma-separated-implies-valuative}），因为 $Y \to S$ 是分离态射
$Y \to X$ 与 $X \to S$ 的复合，故分离（《概形》引理
\ref{schemes-lemma-separated-permanence}）。存在性可如下看出。由于 $y$ 是
$Y$ 的一般点，它属于 $U$。由本引理的假设，存在态射
$a : \Spec(A) \to X$，使
$$
\xymatrix{
\Spec(K) \ar[r]_y \ar[d] & U \ar[r] & X \ar[d]^f \\
\Spec(A) \ar[rr] \ar[rru]^a & & S
}
$$
交换。由于 $Y \to X$ 固有，可以应用固有性的赋值判据（《态射》引理
\ref{morphisms-lemma-characterize-proper}），找到态射
$b : \Spec(A) \to Y$，使
$$
\xymatrix{
\Spec(K) \ar[r]_y \ar[d] & Y \ar[d]^g \\
\Spec(A) \ar[r]^a \ar[ru]^b & X
}
$$
交换。以 $b$ 充当上面的虚线箭头，证明即告完成。
\end{proof}

\begin{lemma}
\label{limits-lemma-refined-valuative-criterion-separated}
设 $f : X \to S$ 和 $h : U \to X$ 是概形态射。假设 $S$ 局部 Noether，
$f$ 局部有限型，$h$ 有限型，并且 $h(U)$ 在 $X$ 中稠密。若对任意交换实线图
$$
\xymatrix{
\Spec(K) \ar[r] \ar[d] & U \ar[r]^h & X \ar[d]^f \\
\Spec(A) \ar[rr] \ar@{-->}[rru] & & S
}
$$
，其中 $A$ 是以 $K$ 为分式域的离散赋值环，至多存在一条使该图交换的
虚线箭头，则 $f$ 分离。
\end{lemma}

\begin{proof}
将引理 \ref{limits-lemma-refined-valuative-criterion-proper} 应用于态射
$U \to X$ 和 $\Delta : X \to X \times_S X$。逐一检验其条件。由《性质》引理
\ref{properties-lemma-locally-Noetherian-quasi-separated}（以及《概形》引理
\ref{schemes-lemma-compose-after-separated}），$\Delta$ 拟紧。当然，
$\Delta$ 局部有限型且分离（任意对角态射都是如此）。最后，假设给定交换实线图
$$
\xymatrix{
\Spec(K) \ar[r] \ar[d] & U \ar[r]^h & X \ar[d]^\Delta \\
\Spec(A) \ar[rr]^{(a, b)} \ar@{-->}[rru] & & X \times_S X
}
$$
，其中 $A$ 是以 $K$ 为分式域的离散赋值环。则 $a$ 和 $b$ 在本引理的图中
给出两条虚线箭头，因而必须相等。因此可用 $a = b$ 作为虚线箭头，这就给出
存在性。证明完成。
\end{proof}

\begin{lemma}
\label{limits-lemma-refined-valuative-criterion-universally-closed}
设 $f : X \to S$ 和 $h : U \to X$ 是概形态射。假设 $S$ 局部 Noether，
$f$ 和 $h$ 有限型，并且 $h(U)$ 在 $X$ 中稠密。若对任意交换实线图
$$
\xymatrix{
\Spec(K) \ar[r] \ar[d] & U \ar[r]^h & X \ar[d]^f \\
\Spec(A) \ar[rr] \ar@{-->}[rru] & & S
}
$$
，其中 $A$ 是以 $K$ 为分式域的离散赋值环，正好存在一条使该图交换的
虚线箭头，则 $f$ 固有。
\end{lemma}

\begin{proof}
合并引理 \ref{limits-lemma-refined-valuative-criterion-separated} 与
\ref{limits-lemma-refined-valuative-criterion-proper}。
\end{proof}

\section{Nagata 基上的赋值判据}
\label{limits-section-nagata-valuative}

\noindent
处理在 Nagata 基上局部有限型的概形时，可以归约到在该基上本质有限型的
离散赋值环。下面只是由此可得的一些示例性结果。

\begin{lemma}
\label{limits-lemma-essentially-finite-type-criterion-universally-closed}
设 $S$ 是 Nagata 概形（因而特别地局部 Noether）。设
$f : X \to Y$ 是在 $S$ 上局部有限型的概形间的拟紧态射。下列条件等价：
\begin{enumerate}
\item $f$ 万有闭；
\item 对每个 $n$，态射
$\mathbf{A}^n \times X \to \mathbf{A}^n \times Y$ 是闭映射；
\item 对任意 $S$ 上的概形交换图
$$
\xymatrix{
U \ar[r] \ar[d] & X \ar[d]^f \\
C \ar[r] \ar@{..>}[ru] & Y
}
$$
，若
\begin{enumerate}
\item $C$ 是在 $S$ 上有限型的正规整概形；
\item 有 $U = C \setminus \{c\}$，其中 $c \in C$ 是某个闭点；
\item $A = \mathcal{O}_{C, c}$ 的维数为 $1$\footnote{由此可知
$A$ 是离散赋值环；见《代数》引理 \ref{algebra-lemma-characterize-dvr}。
此外，$c$ 映到一个有限型点 $s \in S$，并且 $A$ 在
$\mathcal{O}_{S, s}$ 上本质有限型。}
\end{enumerate}
，则在交换图
$$
\xymatrix{
\Spec(K) \ar[r] \ar[d] & X \ar[d]^f \\
\Spec(A) \ar[r] \ar@{-->}[ru] & Y
}
$$
中，其中 $K = \text{Frac}(A)$，存在某条使该图交换的虚线箭头\footnote{由引理
\ref{limits-lemma-morphism-glueing-near-closed-point}，这等价于要求存在一条
使第一个交换图交换的虚线箭头。}。
\end{enumerate}
\end{lemma}

\begin{proof}
引理 \ref{limits-lemma-check-universally-closed-Noetherian} 已经给出了 (1) 与 (2)
的等价性以及它们蕴含 (3) 这一事实。所以只须证明 (3) 蕴含 (2)。注意，若
$f : X \to Y$ 满足条件 (3)，则
$1 \times f : \mathbf{A}^n \times X \to \mathbf{A}^n \times Y$
也满足条件 (3)（见引理 \ref{limits-lemma-check-universally-closed-Noetherian}
的证明中的论证）。因此只须证明 (3) 蕴含 $f$ 是闭映射。

\medskip\noindent
先归约到 $Y$ 和 $S$ 都仿射的情形；建议略过本段。令
$S' \subset S$ 是仿射开集，并令 $Y' \subset Y$ 是一个映入 $S'$ 的仿射开集。
置 $X' = f^{-1}(Y')$。我们断言，限制 $f' : X' \to Y'$（即 $f$ 的限制）
在视为 $S'$ 上的概形态射后仍有性质 (3)。略去细节。现在，若能证明 $f'$
对 $S'$ 和 $Y'$ 的所有选取都是闭映射，则 $f$ 是闭映射。由此归约到下一段讨论的情形。

\medskip\noindent
假设 $S$ 和 $Y$ 仿射。令 $Z \subset X$ 是闭子集。我们可以且确实把 $Z$
视为 $X$ 的既约闭子概形。必须证明 $E = f(Z)$ 闭。取闭点
$y \in Y$，它属于 $f(Z)$ 的闭包。只须证明 $y \in E$。反设
$y \not \in E$，以求矛盾。$s \in S$ 是 $y$ 的像；它是 $S$ 的有限型点；
见《态射》引理 \ref{morphisms-lemma-finite-type-points-morphism}。回忆
$E$ 可构（《态射》引理 \ref{morphisms-lemma-chevalley}）。考虑交集
$\Spec(\mathcal{O}_{Y, y}) \cap E$。这是该谱的一个不含闭点的可构子集
（《态射》引理 \ref{morphisms-lemma-inverse-image-constructible}）。
由于穿孔谱 $\Spec(\mathcal{O}_{Y, y}) \setminus \{y\}$ 是 Jacobson 空间
（《态射》引理 \ref{morphisms-lemma-ubiquity-Jacobson-schemes}），可找到闭点
$t \in \Spec(\mathcal{O}_{Y, y}) \setminus \{y\}$，使 $t \in E$
（见《拓扑》引理 \ref{topology-lemma-jacobson-inherited}）。换言之，
$t \in E$ 是 $Y$ 中的一个点，且有直接特化 $t \leadsto y$。由于
$t \in E$，概形论纤维 $Z_t$ 非空。取闭点 $x \in Z_t$。特别地，由 Hilbert
零点定理可知 $[\kappa(x) : \kappa(t)] < \infty$（《态射》引理
\ref{morphisms-lemma-closed-point-fibre-locally-finite-type}）。

\medskip\noindent
记 $T = \overline{\{t\}} \subset Y$ 为一个整闭子概形，其底拓扑空间就是所示闭包
（《概形》定义 \ref{schemes-definition-reduced-induced-scheme}）。于是
$t \in T$ 是一般点。记 $C \to T$ 为 $T$ 在 $\kappa(x)$ 中的正规化；见《态射》第
\ref{morphisms-section-normalization-X-in-Y} 节（更准确地说，$C \to T$ 是
$T$ 在 $x$ 中的正规化，这里将 $x = \Spec(\kappa(x)) \to T$ 视为 $T$
上的概形）。因为 $S$ 是 Nagata 概形，$T$ 也是 Nagata 概形（《态射》引理
\ref{morphisms-lemma-finite-type-nagata}）。所以 $C \to T$ 有限
（《态射》引理 \ref{morphisms-lemma-nagata-normalization-finite-general}）。
由于 $t$ 在像中，$C \to T$ 满射（因为像闭，而 $T$ 是 $t$ 在 $Y$ 中的闭包）。
选取一点 $c \in C$，它映到 $y \in T$。由于 $y$ 是 $T$ 的闭点，$c$ 是 $C$
的闭点。由于 $\dim(\mathcal{O}_{T, y}) = 1$，可知
$\dim(\mathcal{O}_{C, c}) = 1$（维数至少为 $1$，因为 $c$ 不是 $C$ 的一般点；
至多为 $1$，因为 $C \to T$ 有限）。因为 $C$ 的函数域为 $\kappa(x)$，而 $x$
是 $X$ 的点，存在一个 $Y$-有理映射，其源为 $C$、靶为 $X$（例如见《态射》引理
\ref{morphisms-lemma-rational-map-finite-presentation}）。令
$C \supset U \to X$ 是一个代表元（特别地，$U$ 非空）。可以假设
$c \not \in U$（将 $U$ 替换为 $U \setminus \{c\}$）。由于 $c$ 是余维数为
$1$ 的闭点，并且位于整概形 $C$ 中，有
$C = U \amalg \{c\} \amalg \Sigma$，其中 $\Sigma \subset C$ 是某个真闭子集。
将 $C$ 替换为 $C \setminus \Sigma$ 后，
便构造出第 (3) 部分中的交换图。由引理陈述中的第二个脚注，虚线箭头的存在性
使有理映射延拓到整个 $C$；于是得到矛盾，因为 $c$ 的像将是 $Z$ 中映到 $y$
的一点。
\end{proof}

\begin{lemma}
\label{limits-lemma-essentially-finite-type-criterion-separation}
\enlargethispage{3pt}
设 $S$ 是 Nagata 概形（因而特别地局部 Noether）。设
$f : X \to Y$ 是在 $S$ 上局部有限型的概形间的态射。下列条件等价：
\begin{enumerate}
\item $f$ 分离；
\item 对任意 $S$ 上的概形交换图
$$
\xymatrix{
U \ar[r] \ar[d] & X \ar[d]^f \\
C \ar[r] \ar@{..>}[ru] & Y
}
$$
，若
\begin{enumerate}
\item $C$ 是在 $S$ 上有限型的正规整概形；
\item 有 $U = C \setminus \{c\}$，其中 $c \in C$ 是某个闭点；
\item $A = \mathcal{O}_{C, c}$ 的维数为 $1$\footnote{由此可知
$A$ 是离散赋值环；见《代数》引理 \ref{algebra-lemma-characterize-dvr}。
此外，$c$ 映到一个有限型点 $s \in S$，并且 $A$ 在
$\mathcal{O}_{S, s}$ 上本质有限型。}
\end{enumerate}
，则在交换图
$$
\xymatrix{
\Spec(K) \ar[r] \ar[d] & X \ar[d]^f \\
\Spec(A) \ar[r] \ar@{-->}[ru] & Y
}
$$
中，其中 $K = \text{Frac}(A)$，至多存在一条使该图交换的虚线箭头\footnote{由
引理 \ref{limits-lemma-morphism-glueing-near-closed-point}，这等价于要求至多
存在一条使第一个交换图交换的虚线箭头。}。
\end{enumerate}
\end{lemma}

\begin{proof}
由引理 \ref{limits-lemma-Noetherian-dvr-valuative-separation} 可知 (1) 蕴含 (2)。
假设 (2)。为证明 $f$ 分离，必须证明
$\Delta : X \to X \times_Y X$ 闭。由《态射》引理
\ref{morphisms-lemma-finite-type-Noetherian-quasi-separated}，态射 $\Delta$
拟紧。由引理 \ref{limits-lemma-essentially-finite-type-criterion-universally-closed}，
只须证明：对任意 $S$ 上的概形交换图
$$
\xymatrix{
U \ar[rr] \ar[d] & & X \ar[d]^\Delta \\
C \ar[rr]^{(a_1, a_2)} \ar@{..>}[rru] & & X \times_Y X
}
$$
，若
\begin{enumerate}
\item $C$ 是在 $S$ 上有限型的正规整概形；
\item 有 $U = C \setminus \{c\}$，其中 $c \in C$ 是某个闭点；
\item $A = \mathcal{O}_{C, c}$ 的维数为 $1$；
\end{enumerate}
则在交换图
$$
\xymatrix{
\Spec(K) \ar[r] \ar[d] & X \ar[d]^\Delta \\
\Spec(A) \ar[r] \ar@{-->}[ru] & X \times_Y X
}
$$
中，其中 $K = \text{Frac}(A)$，存在某条使该图交换的虚线箭头。由引理
\ref{limits-lemma-morphism-glueing-near-closed-point}，第二幅图中虚线箭头的存在性
等价于第一幅图中虚线箭头的存在性。而后一种存在性又等价于要求
$a_1 = a_2$。不过 $a_1|_U = a_2|_U$，所以由唯一性假设 (2)，该等式确实成立，
证明完成。
\end{proof}

\begin{lemma}
\label{limits-lemma-essentially-finite-type-criterion-proper}
设 $S$ 是 Nagata 概形（因而特别地局部 Noether）。设
$f : X \to Y$ 是在 $S$ 上局部有限型的概形间的拟紧态射。下列条件等价：
\begin{enumerate}
\item $f$ 固有；
\item 对任意 $S$ 上的概形交换图
$$
\xymatrix{
U \ar[r] \ar[d] & X \ar[d]^f \\
C \ar[r] \ar@{..>}[ru] & Y
}
$$
，若
\begin{enumerate}
\item $C$ 是在 $S$ 上有限型的正规整概形；
\item 有 $U = C \setminus \{c\}$，其中 $c \in C$ 是某个闭点；
\item $A = \mathcal{O}_{C, c}$ 的维数为 $1$\footnote{由此可知
$A$ 是离散赋值环；见《代数》引理 \ref{algebra-lemma-characterize-dvr}。
此外，$c$ 映到一个有限型点 $s \in S$，并且 $A$ 在
$\mathcal{O}_{S, s}$ 上本质有限型。}
\end{enumerate}
，则在交换图
$$
\xymatrix{
\Spec(K) \ar[r] \ar[d] & X \ar[d]^f \\
\Spec(A) \ar[r] \ar@{-->}[ru] & Y
}
$$
中，其中 $K = \text{Frac}(A)$，正好存在一条使该图交换的虚线箭头\footnote{由
引理 \ref{limits-lemma-morphism-glueing-near-closed-point}，这等价于要求正好
存在一条使第一个交换图交换的虚线箭头。}。
\end{enumerate}
\end{lemma}

\begin{proof}
这由引理 \ref{limits-lemma-essentially-finite-type-criterion-universally-closed}、
\ref{limits-lemma-essentially-finite-type-criterion-separation}，以及固有态射定义为
有限型、分离且万有闭的态射形式地得出。
\end{proof}

\section{极限与纤维维数}
\label{limits-section-limits-dimension}

\noindent
下述引理最常用于引理 \ref{limits-lemma-descend-finite-presentation} 的情形，以确保：
若极限的纤维维数 $\leq d$，则在某个有限阶段，纤维维数也 $\leq d$。

\begin{lemma}
\label{limits-lemma-limit-dimension}
设 $I$ 是有向集。设 $(f_i : X_i \to S_i)$ 是指标集为 $I$ 的概形态射逆系统。
假设
\begin{enumerate}
\item 所有态射 $S_{i'} \to S_i$ 都仿射；
\item 所有概形 $S_i$ 都拟紧且拟分离；
\item 态射 $f_i$ 都有限型；并且
\item 态射 $X_{i'} \to X_i \times_{S_i} S_{i'}$ 都是闭浸入。
\end{enumerate}
令 $f : X = \lim_i X_i \to S = \lim_i S_i$ 为极限。令 $d \geq 0$。
若 $f$ 的每个纤维维数都 $\leq d$，则对某个 $i$，$f_i$ 的每个纤维维数
都 $\leq d$。
\end{lemma}

\begin{proof}
对每个 $i$，令
$U_i = \{x \in X_i \mid \dim_x((X_i)_{f_i(x)}) \leq d\}$。
这是 $X_i$ 的开子集；见《态射》引理
\ref{morphisms-lemma-openness-bounded-dimension-fibres}。置
$Z_i = X_i \setminus U_i$（赋予既约诱导概形结构）。必须证明
$Z_i = \emptyset$ 对某个 $i$ 成立。否则，由引理 \ref{limits-lemma-limit-nonempty}，
$Z = \lim Z_i \not = \emptyset$。设 $z \in Z$ 是一点。注意
$Z \subset X$ 是闭子概形。置 $s = f(z)$。对每个 $i$，令 $s_i \in S_i$
为 $s$ 的像。指出：$Z_s$ 是概形 $(Z_i)_{s_i}$ 的极限，而 $Z_s$ 也是将概形
$(Z_i)_{s_i}$ 基变换到 $\kappa(s)$ 后所得诸概形的极限。此外，由假设 (4)，
所有态射
$$
Z_s
\longrightarrow
(Z_{i'})_{s_{i'}} \times_{\Spec(\kappa(s_{i'}))} \Spec(\kappa(s))
\longrightarrow
(Z_i)_{s_i} \times_{\Spec(\kappa(s_i))} \Spec(\kappa(s))
\longrightarrow
X_s
$$
都是闭浸入。因此 $Z_s$ 是诸闭子概形
$(Z_i)_{s_i} \times_{\Spec(\kappa(s_i))} \Spec(\kappa(s))$
在 $X_s$ 中的概形论交。由于概形
$(Z_i)_{s_i} \times_{\Spec(\kappa(s_i))} \Spec(\kappa(s))$
的所有不可约分支维数都 $> d$ 且包含 $z$，可知
$Z_s$ 含有一个维数 $> d$、通过 $z$ 的不可约分支。这与
$Z_s \subset X_s$ 及 $\dim(X_s) \leq d$ 矛盾。
\end{proof}

\begin{lemma}
\label{limits-lemma-descend-quasi-finite}
记号和假设同情形 \ref{limits-situation-descent-property}。若
\begin{enumerate}
\item $f$ 是拟有限态射；并且
\item $f_0$ 局部有限型，
\end{enumerate}
则存在 $i \geq 0$，使 $f_i$ 拟有限。
\end{lemma}

\begin{proof}
立即由引理 \ref{limits-lemma-limit-dimension} 得出。
\end{proof}

\begin{lemma}
\label{limits-lemma-eventually-relative-dimension}
假设和记号同情形 \ref{limits-situation-descent-property}。令 $d \geq 0$。若
\begin{enumerate}
\item $f$ 的相对维数 $\leq d$
（《态射》定义 \ref{morphisms-definition-relative-dimension-d}）；并且
\item $f_0$ 局部有限型，
\end{enumerate}
则存在某个 $i$，使 $f_i$ 的相对维数 $\leq d$。
\end{lemma}

\begin{proof}
立即由引理 \ref{limits-lemma-limit-dimension} 得出。
\end{proof}

\begin{lemma}
\label{limits-lemma-descend-dimension-d}
记号和假设同情形 \ref{limits-situation-descent-property}。若
\begin{enumerate}
\item $f$ 的相对维数为 $d$；并且
\item $f_0$ 局部有限呈示，
\end{enumerate}
则存在 $i \geq 0$，使 $f_i$ 的相对维数为 $d$。
\end{lemma}

\begin{proof}
由引理 \ref{limits-lemma-limit-dimension}，可假设 $f_0$ 的所有纤维维数都
$\leq d$。由《态射》引理
\ref{morphisms-lemma-openness-bounded-dimension-fibres-finite-presentation}，
集合 $U_0 \subset X_0$ 由满足下述条件的点 $x \in X_0$ 构成：
$X_0 \to Y_0$ 的纤维在 $x$ 处维数 $\leq d - 1$；该集合在 $X_0$ 中开且逆紧。因此其补集
$E = X_0 \setminus U_0$ 可构。此外，由《态射》引理
\ref{morphisms-lemma-dimension-fibre-after-base-change}，$X \to X_0$
的像包含在 $E$ 中。所以当 $i \gg 0$ 时，$X_i \to X_0$ 的像包含在 $E$ 中
（引理 \ref{limits-lemma-limit-contained-in-constructible}）。于是再由上述《态射》引理
\ref{morphisms-lemma-dimension-fibre-after-base-change}，$X_i \to Y_i$
的所有纤维维数均为 $d$。
\end{proof}

\begin{lemma}
\label{limits-lemma-approximate-given-relative-dimension}
设 $S$ 是拟紧拟分离概形。设 $f : X \to S$ 是有限呈示态射。设
$d \geq 0$ 是整数。若 $Z \subset X$ 是闭子概形，且
$\dim(Z_s) \leq d$ 对所有 $s \in S$ 成立，则存在闭子概形 $Z' \subset X$，满足
\begin{enumerate}
\item $Z \subset Z'$；
\item $Z' \to X$ 有限呈示；并且
\item $\dim(Z'_s) \leq d$ 对所有 $s \in S$ 成立。
\end{enumerate}
\end{lemma}

\begin{proof}
由命题 \ref{limits-proposition-approximate}，可将 $S = \lim S_i$ 写成带仿射过渡态射的
Noether 概形有向逆系统的极限。由引理 \ref{limits-lemma-descend-finite-presentation}，
可以假设存在有限呈示态射系统 $f_i : X_i \to S_i$，使得
$X_{i'} = X_i \times_{S_i} S_{i'}$ 对所有 $i' \geq i$ 成立，并且
$X = X_i \times_{S_i} S$。令 $Z_i \subset X_i$ 是
$Z \to X \to X_i$ 的概形论像。则当 $i' \geq i$ 时，态射
$X_{i'} \to X_i$ 将 $Z_{i'}$ 映入 $Z_i$，且诱导态射
$Z_{i'} \to Z_i \times_{S_i} S_{i'}$ 是闭浸入。由引理
\ref{limits-lemma-limit-dimension} 可知，$Z_i \to S_i$ 的所有纤维维数均
$\leq d$，这里 $i \in I$ 取为某个适当指标。固定这样的 $i$，并置
$Z' = Z_i \times_{S_i} S \subset X$。由于 $S_i$ 是 Noether 概形，
$X_i$ 也是 Noether 概形，因而态射 $Z_i \to X_i$ 有限呈示。所以其基变换
$Z' \to X$ 也有限呈示。此外，$Z' \to S$ 的纤维都是 $Z_i \to S_i$
相应纤维的基变换，因而维数 $\leq d$。
\end{proof}

\section{最高次的基变换}
\label{limits-section-top-degree}

\noindent
对于固有态射和有限型拟凝聚模，基变换映射在最高次是同构。

\begin{lemma}
\label{limits-lemma-top-cohomology-functor}
设 $f : X \to Y$ 是概形态射。令 $d \geq 0$。假设
\begin{enumerate}
\item $X$ 和 $Y$ 都拟紧且拟分离；并且
\item $R^if_*\mathcal{F} = 0$ 对 $i > d$ 以及每个拟凝聚
$\mathcal{O}_X$-模 $\mathcal{F}$ 成立。
\end{enumerate}
则有：
\begin{enumerate}
\item[(a)] 对任意基变换图
$$
\xymatrix{
X' \ar[d]_{f'} \ar[r]_{g'} & X \ar[d]^f \\
Y' \ar[r]^g & Y
}
$$
，都有 $R^if'_*\mathcal{F}' = 0$，这里 $i > d$，且任意拟凝聚
$\mathcal{O}_{X'}$-模 $\mathcal{F}'$ 均可取；
\item[(b)]
$R^df'_*(\mathcal{F}' \otimes_{\mathcal{O}_{X'}} (f')^*\mathcal{G}') =
R^df'_*\mathcal{F}' \otimes_{\mathcal{O}_{Y'}} \mathcal{G}'$
对任意拟凝聚 $\mathcal{O}_{Y'}$-模 $\mathcal{G}'$ 成立；
\item[(c)] $R^df'_*\mathcal{F}'$ 的构造与任意进一步的基变换可交换
（说明见证明）。
\end{enumerate}
\end{lemma}

\begin{proof}
在给出证明之前，先解释 (c) 的含义。假设还有一个接在给定图后的笛卡尔方块
$$
\xymatrix{
X'' \ar[d]_{f''} \ar[r]_{h'} &
X' \ar[d]_{f'} \ar[r]_{g'} &
X \ar[d]^f \\
Y'' \ar[r]^h &
Y' \ar[r]^g &
Y
}
$$
。若 (a) 成立，则有典范映射
$\gamma : h^*R^df'_*\mathcal{F}' \to R^df''_*(h')^*\mathcal{F}'$。
也就是说，$\gamma$ 是由复合
$$
Lh^*Rf'_*\mathcal{F}' \longrightarrow
Rf''_*L(h')^*\mathcal{F}' \longrightarrow Rf''_*(h')^*\mathcal{F}'
$$
在第 $d$ 次上同调层上诱导的映射。这里第一箭头是基变换映射
（《上同调》注记 \ref{cohomology-remark-base-change}），第二箭头来自典范映射
$L(g')^*\mathcal{F} \to (g')^*\mathcal{F}$。同样，由 (a) 可知
$Rf'_*\mathcal{F}$ 在次数 $> d$ 没有非零上同调层，故
$H^d(Lh^*Rf_*\mathcal{F}') = h^*R^df_*\mathcal{F}$。
(c) 的内容就是 $\gamma$ 为同构。

\medskip\noindent
作此说明后，可以在 $Y'$ 和 $Y''$ 上局部检验 (a)、(b) 和 (c)。
假设 $V \subset Y$ 是拟紧开子概形。则我们断言，(1) 和 (2) 对
$f|_{f^{-1}(V)} : f^{-1}(V) \to V$ 成立。事实上，(1) 是立即的；(2) 成立，
是因为 $f^{-1}(V)$ 上任意拟凝聚模都是 $X$ 上某个拟凝聚模的限制
（《性质》引理 \ref{properties-lemma-extend-trivial}），并且高阶直像的构造
与对开集的限制可交换。因此也可在 $Y$ 上局部工作。换言之，可以假设
$Y''$、$Y'$ 和 $Y$ 都是仿射概形。

\medskip\noindent
当 $Y'$ 和 $Y$ 仿射时证明 (a)。此时态射 $g$ 和 $g'$ 都仿射。因此
$g_* = Rg_*$ 且 $g'_* = Rg'_*$（《概形上同调》引理
\ref{coherent-lemma-relative-affine-vanishing}），并且 $g_*$ 与模的限制标量函子
等同（《概形》引理 \ref{schemes-lemma-widetilde-pullback}）。于是
$$
g_*(R^if'_*\mathcal{F}') = H^i(Rg_*Rf'_*\mathcal{F}') =
H^i(Rf_*Rg'_*\mathcal{F}') =
H^i(Rf_*g'_*\mathcal{F}') =
Rf^i_*g'_*\mathcal{F}'
$$
由假设 (2)，这等于零。结合上述对 $g_*$ 的刻画即得 (a)。

\medskip\noindent
当 $Y'$ 仿射时证明 (b)，设 $Y' = \Spec(R')$。由 (a)，有
$H^{d + 1}(X', \mathcal{F}') = 0$ 对任意拟凝聚
$\mathcal{O}_{X'}$-模 $\mathcal{F}'$ 成立；见《概形上同调》引理
\ref{coherent-lemma-quasi-coherence-higher-direct-images-application}。
考虑函子 $F$，其定义域为 $R'$-模：
$$
F(M) = H^d(X', \mathcal{F}' \otimes_{\mathcal{O}_{X'}} (f')^*\widetilde{M})
$$
由《上同调》引理 \ref{cohomology-lemma-quasi-separated-cohomology-colimit}，
该函子与直和可交换（这里用到 $X$，从而 $X'$ 拟紧且拟分离）。
另一方面，若 $M_1 \to M_2 \to M_3 \to 0$ 是正合列，则
$$
\mathcal{F}' \otimes_{\mathcal{O}_{X'}} (f')^*\widetilde{M}_1 \to
\mathcal{F}' \otimes_{\mathcal{O}_{X'}} (f')^*\widetilde{M}_2 \to
\mathcal{F}' \otimes_{\mathcal{O}_{X'}} (f')^*\widetilde{M}_3 \to 0
$$
是 $X'$ 上拟凝聚模的正合列；由上面的高阶上同调消失性，得到正合列
$$
F(M_1) \to F(M_2) \to F(M_3) \to 0
$$
。换言之，$F$ 右正合。任意与直和可交换的右正合 $R'$-线性函子
$F : \text{Mod}_{R'} \to \text{Mod}_{R'}$
都由与某个 $R'$-模张量给出（证明略去，留作读者练习）。
所以得到 $F(M) = H^d(X', \mathcal{F}') \otimes_{R'} M$。
由于 $R^d(f')_*\mathcal{F}'$ 和
$R^d(f')_*(\mathcal{F}' \otimes_{\mathcal{O}_{X'}} (f')^*\widetilde{M})$
都是拟凝聚的（《概形上同调》引理
\ref{coherent-lemma-quasi-coherence-higher-direct-images}），等式
$F(M) = H^d(X', \mathcal{F}') \otimes_{R'} M$ 正好转化为 (b) 中的陈述。

\medskip\noindent
当 $Y'' \to Y' \to Y$ 是仿射概形态射时证明 (c)。设
$Y'' = \Spec(R'')$ 且 $Y' = \Spec(R')$。于是
$R^df''_*(h')^*\mathcal{F}'$ 是 $Y'$ 上由 $R''$-模
$H^d(X'', (h')^*\mathcal{F}')$ 所伴随的拟凝聚模。现在
$h' : X'' \to X'$ 仿射，因而由已经使用过的《概形上同调》引理
\ref{coherent-lemma-relative-affine-cohomology}，
$H^d(X'', (h')^*\mathcal{F}') = H^d(X, h'_*(h')^*\mathcal{F}')$。
有
$$
h'_*(h')^*\mathcal{F}' =
\mathcal{F}' \otimes_{\mathcal{O}_{X'}} (f')^*\widetilde{R''}
$$
；读者可在某个仿射开覆盖上检验这一点。因此
$H^d(X'', (h')^*\mathcal{F}') = H^d(X', \mathcal{F}') \otimes_{R'} R''$；
这是由 (b) 应用于 $f'$ 得到的。证明完成。
\end{proof}

\begin{lemma}
\label{limits-lemma-higher-direct-images-zero-above-dimension-fibre}
设 $f : X \to Y$ 是概形态射。设 $y \in Y$。假设 $f$ 固有且
$\dim(X_y) = d$。则
\begin{enumerate}
\item 对 $\mathcal{F} \in \QCoh(\mathcal{O}_X)$，有
$(R^if_*\mathcal{F})_y = 0$ 对所有 $i > d$ 成立；
\item 存在仿射开邻域 $V \subset Y$，它包含 $y$，并使
$f^{-1}(V) \to V$ 和 $d$ 满足引理
\ref{limits-lemma-top-cohomology-functor} 的假设及结论。
\end{enumerate}
\end{lemma}

\begin{proof}
由《态射》引理 \ref{morphisms-lemma-openness-bounded-dimension-fibres}
以及 $f$ 是闭映射这一事实，可找到仿射开邻域 $V$，它包含 $y$，并使 $V$ 上各点的
纤维维数均 $\leq d$。因此可以假设 $X \to Y$ 是固有态射，其所有纤维维数
都 $\leq d$，且 $Y$ 仿射。我们将证明 (2)；这立即蕴含 (1) 对所有 $y \in Y$ 成立。

\medskip\noindent
由引理 \ref{limits-lemma-proper-limit-of-proper-finite-presentation}，可以把
$X = \lim X_i$ 写成余滤极限，其中 $X_i \to Y$ 固有且有限呈示，并且
$X \to X_i$ 及过渡态射都是闭浸入。由引理 \ref{limits-lemma-limit-dimension}，
对某个 $i$，$X_i \to Y$ 的纤维维数都 $\leq d$。对拟凝聚
$\mathcal{O}_X$-模 $\mathcal{F}$，由《概形上同调》引理
\ref{coherent-lemma-relative-affine-vanishing} 及 Leray
（《上同调》引理 \ref{cohomology-lemma-relative-Leray}），有
$R^pf_*\mathcal{F} = R^pf_{i, *}(X \to X_i)_*\mathcal{F}$。
所以可将 $X$ 替换为 $X_i$，归约到下一段讨论的情形。

\medskip\noindent
假设 $Y$ 仿射，$f : X \to Y$ 固有且有限呈示，并且所有纤维维数都
$\leq d$。只须证明 $H^p(X, \mathcal{F}) = 0$ 对 $p > d$ 成立。事实上，由
《概形上同调》引理
\ref{coherent-lemma-quasi-coherence-higher-direct-images-application}，
有 $H^p(X, \mathcal{F}) = H^0(Y, R^pf_*\mathcal{F})$。另一方面，由
《概形上同调》引理
\ref{coherent-lemma-quasi-coherence-higher-direct-images}，
$R^pf_*\mathcal{F}$ 在 $Y$ 上拟凝聚，所以整体截面消失蕴含该层消失。
将 $Y = \lim_{i \in I} Y_i$ 写成仿射概形的余滤极限，其中 $Y_i$
是某个 Noether 环（例如有限型 $\mathbf{Z}$-代数）的谱。由引理
\ref{limits-lemma-descend-finite-presentation}，可以选取元素 $0 \in I$ 和有限型态射
$X_0 \to Y_0$，使 $X \cong Y \times_{Y_0} X_0$。增大 $0$ 后，可以假设
$X_0 \to Y_0$ 固有（引理 \ref{limits-lemma-eventually-proper}），且
$X_0 \to Y_0$ 的纤维维数都 $\leq d$（引理 \ref{limits-lemma-limit-dimension}）。
由于 $X \to X_0$ 仿射，由《概形上同调》引理
\ref{coherent-lemma-relative-affine-cohomology} 可知
$H^p(X, \mathcal{F}) = H^p(X_0, (X \to X_0)_*\mathcal{F})$。
这就归约到下一段讨论的情形。

\medskip\noindent
假设 $Y$ 是仿射 Noether 概形，$f : X \to Y$ 固有，并且所有纤维维数都
$\leq d$。此时，由《性质》引理
\ref{properties-lemma-directed-colimit-finite-presentation}，可将
$\mathcal{F} = \colim \mathcal{F}_i$ 写成凝聚 $\mathcal{O}_X$-模的滤过余极限。
由《上同调》引理 \ref{cohomology-lemma-quasi-separated-cohomology-colimit}，
$H^p(X, \mathcal{F}) = \colim H^p(X, \mathcal{F}_i)$。所以可以假设
$\mathcal{F}$ 凝聚。在此情形下，由《概形上同调》引理
\ref{coherent-lemma-higher-direct-images-zero-above-dimension-fibre}，有
$(R^pf_*\mathcal{F})_y = 0$ 对所有 $y \in Y$ 成立。因此
$R^pf_*\mathcal{F} = 0$，从而 $H^p(X, \mathcal{F}) = 0$（见上），
结论成立。
\end{proof}

\begin{lemma}
\label{limits-lemma-proper-top-cohomology-finite-type}
设 $f : X \to Y$ 是概形态射。令 $d \geq 0$。设 $\mathcal{F}$ 是
$\mathcal{O}_X$-模。假设
\begin{enumerate}
\item $f$ 是固有态射，且其所有纤维维数都 $\leq d$；
\item $\mathcal{F}$ 是有限型拟凝聚 $\mathcal{O}_X$-模。
\end{enumerate}
则 $R^df_*\mathcal{F}$ 是有限型拟凝聚 $\mathcal{O}_X$-模。
\end{lemma}

\begin{proof}
由《概形上同调》引理
\ref{coherent-lemma-quasi-coherence-higher-direct-images}，
模 $R^df_*\mathcal{F}$ 拟凝聚。问题在 $Y$ 上是局部的，所以可以假设
$Y$ 仿射。设 $Y = \Spec(R)$。于是只须证明
$H^d(X, \mathcal{F})$ 是有限 $R$-模。

\medskip\noindent
由引理 \ref{limits-lemma-proper-limit-of-proper-finite-presentation}，可以把
$X = \lim X_i$ 写成余滤极限，其中 $X_i \to Y$ 固有且有限呈示，并且
$X \to X_i$ 及过渡态射都是闭浸入。由引理 \ref{limits-lemma-limit-dimension}，
对某个 $i$，$X_i \to Y$ 的纤维维数都 $\leq d$。由《概形上同调》引理
\ref{coherent-lemma-relative-affine-vanishing} 及 Leray
（《上同调》引理 \ref{cohomology-lemma-relative-Leray}），有
$R^pf_*\mathcal{F} = R^pf_{i, *}(X \to X_i)_*\mathcal{F}$。
所以可以将 $X$ 替换为 $X_i$，归约到下一段讨论的情形。

\medskip\noindent
假设 $Y$ 仿射，$f : X \to Y$ 固有且有限呈示，并且所有纤维维数都
$\leq d$。由《性质》引理
\ref{properties-lemma-finite-directed-colimit-surjective-maps}，
可将 $\mathcal{F}$ 写成某个有限呈示 $\mathcal{O}_X$-模
$\mathcal{F}'$ 的商。映射
$H^d(X, \mathcal{F}') \to H^d(X, \mathcal{F})$ 满射，因为由引理
\ref{limits-lemma-higher-direct-images-zero-above-dimension-fibre}
中的高阶上同调消失性（或其证明），有
$H^{d + 1}(X, \Ker(\mathcal{F}' \to \mathcal{F})) = 0$。
于是归约到下一段讨论的情形。

\medskip\noindent
假设 $Y = \Spec(R)$ 仿射，$f : X \to Y$ 固有且有限呈示，所有纤维维数
都 $\leq d$，并且 $\mathcal{F}$ 是有限呈示 $\mathcal{O}_X$-模。
将 $Y = \lim_{i \in I} Y_i$ 写成仿射概形的余滤极限，其中
$Y_i = \Spec(R_i)$ 是某个 Noether 环（例如有限型 $\mathbf{Z}$-代数）的谱。
由引理 \ref{limits-lemma-descend-finite-presentation}，可以选取元素 $0 \in I$
和有限型态射 $X_0 \to Y_0$，使 $X \cong Y \times_{Y_0} X_0$。
增大 $0$ 后，可以假设 $X_0 \to Y_0$ 固有
（引理 \ref{limits-lemma-eventually-proper}），且 $X_0 \to Y_0$ 的纤维维数
都 $\leq d$（引理 \ref{limits-lemma-limit-dimension}）。再次增大 $0$ 后，
可以假设存在凝聚 $\mathcal{O}_{X_0}$-模 $\mathcal{F}_0$，其拉回为
$\mathcal{F}$；见引理 \ref{limits-lemma-descend-modules-finite-presentation}。
由引理 \ref{limits-lemma-top-cohomology-functor}，有
$$
H^d(X, \mathcal{F}) = H^d(X_0, \mathcal{F}_0) \otimes_{R_0} R
$$
。由《概形上同调》引理
\ref{coherent-lemma-proper-over-affine-cohomology-finite}，
上同调模 $H^d(X_0, \mathcal{F}_0)$ 有限，故证明完成。
\end{proof}

\begin{lemma}
\label{limits-lemma-proper-top-cohomology-finite-presentation}
设 $f : X \to Y$ 是概形态射。令 $d \geq 0$。设 $\mathcal{F}$ 是
$\mathcal{O}_X$-模。假设
\begin{enumerate}
\item $f$ 是有限呈示的固有态射，且其所有纤维维数都 $\leq d$；
\item $\mathcal{F}$ 是有限呈示 $\mathcal{O}_X$-模。
\end{enumerate}
则 $R^df_*\mathcal{F}$ 是有限呈示 $\mathcal{O}_X$-模。
\end{lemma}

\begin{proof}
证明与引理 \ref{limits-lemma-proper-top-cohomology-finite-type} 的证明完全相同，
但可略去第三段。细节从略。
\end{proof}

\section{闭纤维上的黏合}
\label{limits-section-change-over-closed-points}

\noindent
将上述理论应用于局部环的谱，得到下述关于相对概形的优美黏合结果。

\begin{lemma}
\label{limits-lemma-glueing-near-closed-point}
设 $S$ 是概形。设 $s \in S$ 是闭点，并且
$U = S \setminus \{s\} \to S$ 拟紧。置
$V = \Spec(\mathcal{O}_{S, s}) \setminus \{s\}$，则有范畴等价
$$
\left\{
\begin{matrix}
X \to S\text{ of finite presentation}
\end{matrix}
\right\}
\longrightarrow
\left\{
\vcenter{
\xymatrix{
X' \ar[d] & Y' \ar[d] \ar[l] \ar[r] & Y \ar[d] \\
U & V \ar[l] \ar[r] & \Spec(\mathcal{O}_{S, s})
}
}
\right\}
$$
。右端考虑这样的交换图：其中各方块都是笛卡尔方块，各竖直箭头都是
有限呈示态射。
\end{lemma}

\begin{proof}
设 $W \subset S$ 是 $s$ 的开邻域。由相对概形的黏合（见《构造》第
\ref{constructions-section-relative-glueing} 节），函子
$$
\left\{
\begin{matrix}
X \to S\text{ of finite presentation}
\end{matrix}
\right\}
\longrightarrow
\left\{
\vcenter{
\xymatrix{
X' \ar[d] & Y' \ar[d] \ar[l] \ar[r] & Y \ar[d] \\
U & W \setminus \{s\} \ar[l] \ar[r] & W
}
}
\right\}
$$
是范畴等价。有
$\mathcal{O}_{S, s} = \colim \mathcal{O}_W(W)$，其中 $W$ 遍历
$s$ 的仿射开邻域。所以
$\Spec(\mathcal{O}_{S, s}) = \lim W$，其中 $W$ 遍历 $s$ 的仿射开邻域。
因此，$\Spec(\mathcal{O}_{S, s})$ 上有限呈示概形的范畴，是各
$W$ 上有限呈示概形范畴的极限，其中 $W$ 遍历 $s$ 的仿射开邻域；见引理
\ref{limits-lemma-descend-finite-presentation}。
对每个仿射开集 $s \in W$，可见 $U \cap W$ 拟紧，因为 $U \to S$ 拟紧。
所以 $V = \lim W \cap U = \lim W \setminus \{s\}$ 是拟紧拟分离概形的极限
（见引理 \ref{limits-lemma-directed-inverse-system-has-limit}）。因此，$V$
上有限呈示概形的范畴也是各 $W \cap U$ 上有限呈示概形范畴的极限，
其中 $W$ 遍历 $s$ 的仿射开邻域。合并这些结果即可形式地得出本引理。
\end{proof}

\begin{lemma}
\label{limits-lemma-glueing-near-closed-point-modules}
设 $S$ 是概形。设 $s \in S$ 是闭点，并且
$U = S \setminus \{s\} \to S$ 拟紧。置
$V = \Spec(\mathcal{O}_{S, s}) \setminus \{s\}$，则有范畴等价
$$
\left\{
\mathcal{O}_S\text{-模 }\mathcal{F}\text{ of finite presentation}
\right\}
\longrightarrow
\left\{
(\mathcal{G}, \mathcal{H}, \alpha)
\right\}
$$
。右端考虑三元组，其组成如下：有限呈示 $\mathcal{O}_U$-模
$\mathcal{G}$；有限呈示
$\mathcal{O}_{\Spec(\mathcal{O}_{S, s})}$-模 $\mathcal{H}$；以及同构
$\alpha : \mathcal{G}|_V \to \mathcal{H}|_V$，它是
$\mathcal{O}_V$-模同构。
\end{lemma}

\begin{proof}
可以重做引理 \ref{limits-lemma-glueing-near-closed-point} 的证明，其中使用引理
\ref{limits-lemma-descend-modules-finite-presentation}；也可以从引理
\ref{limits-lemma-glueing-near-closed-point} 推出本结论，所用的是拟凝聚模与
《构造》第 \ref{constructions-section-vector-bundle} 节所述“向量丛”之间的等价。
细节从略。
\end{proof}

\begin{lemma}
\label{limits-lemma-glueing-near-point}
设 $S$ 是概形。设 $U \subset S$ 是逆紧开集。设 $s \in S$ 是
$U$ 的补集中的一点。置 $V = \Spec(\mathcal{O}_{S, s}) \cap U$，则有范畴等价
$$
\colim_{s \in U' \supset U\text{ 开}}
\left\{
\vcenter{
\xymatrix{
X \ar[d] \\
U'
}
}
\right\}
\longrightarrow
\left\{
\vcenter{
\xymatrix{
X' \ar[d] & Y' \ar[d] \ar[l] \ar[r] & Y \ar[d] \\
U & V \ar[l] \ar[r] & \Spec(\mathcal{O}_{S, s})
}
}
\right\}
$$
。左端的竖直箭头是有限呈示态射；右端考虑这样的交换图：其中各方块
都是笛卡尔方块，各竖直箭头都是有限呈示态射。
\end{lemma}

\begin{proof}
设 $W \subset S$ 是 $s$ 的开邻域。由相对概形的黏合（见《构造》第
\ref{constructions-section-relative-glueing} 节），函子
$$
\left\{
\begin{matrix}
X \to U' = U \cup W \text{ of finite presentation}
\end{matrix}
\right\}
\longrightarrow
\left\{
\vcenter{
\xymatrix{
X' \ar[d] & Y' \ar[d] \ar[l] \ar[r] & Y \ar[d] \\
U & W \cap U \ar[l] \ar[r] & W
}
}
\right\}
$$
是范畴等价。有
$\mathcal{O}_{S, s} = \colim \mathcal{O}_W(W)$，其中 $W$ 遍历
$s$ 的仿射开邻域。所以
$\Spec(\mathcal{O}_{S, s}) = \lim W$，其中 $W$ 遍历 $s$ 的仿射开邻域。
因此，$\Spec(\mathcal{O}_{S, s})$ 上有限呈示概形的范畴，是各
$W$ 上有限呈示概形范畴的极限，其中 $W$ 遍历 $s$ 的仿射开邻域；见引理
\ref{limits-lemma-descend-finite-presentation}。
对每个仿射开集 $s \in W$，可见 $U \cap W$ 拟紧，因为 $U \to S$ 拟紧。
所以 $V = \lim W \cap U$ 是拟紧拟分离概形的极限
（见引理 \ref{limits-lemma-directed-inverse-system-has-limit}）。因此，$V$
上有限呈示概形的范畴也是各 $W \cap U$ 上有限呈示概形范畴的极限，
其中 $W$ 遍历 $s$ 的仿射开邻域。合并这些结果即可形式地得出本引理。
\end{proof}

\begin{lemma}
\label{limits-lemma-glueing-near-point-properties}
记号和假设同引理 \ref{limits-lemma-glueing-near-point}。设
$U \subset U' \subset X$ 是含 $s$ 的开集。
\begin{enumerate}
\item 设 $f' : X \to U'$ 通过该等价对应于 $f : X' \to U$ 和
$g : Y \to \Spec(\mathcal{O}_{S, s})$。若 $f$ 和 $g$ 具有下列任一性质：
分离、固有、有限或 \'etale，则在可能缩小 $U'$ 后，态射 $f'$ 具有同一性质。
\item 设 $a : X_1 \to X_2$ 是 $U'$ 上有限呈示概形的态射，其基变换
$a' : X'_1 \to X'_2$ 位于 $U$ 上，而
$b : Y_1 \to Y_2$ 位于 $\Spec(\mathcal{O}_{S, s})$ 上。若 $a'$ 和 $b$
具有下列任一性质：分离、固有、有限或 \'etale，则在可能缩小 $U'$ 后，态射 $a$
具有同一性质。
\end{enumerate}
\end{lemma}

\begin{proof}
证明 (1)。回忆 $\Spec(\mathcal{O}_{S, s})$ 是 $s$ 在 $S$ 中的各仿射开邻域
的极限。由于 $g$ 具有所论性质，$f'$ 对其中某个仿射开邻域的限制也具有该性质；
见引理 \ref{limits-lemma-descend-separated-finite-presentation}、
\ref{limits-lemma-eventually-proper}、
\ref{limits-lemma-descend-finite-finite-presentation} 和
\ref{limits-lemma-descend-etale}。由于 $f'$ 在 $U$ 上具有所给性质，正如 $f$
所具有的一样，又因为该性质可在基上局部检验，结论成立。

\medskip\noindent
证明 (2)。若写
$\Spec(\mathcal{O}_{S, s}) = \lim W$，其中 $W$ 遍历 $s$ 在 $S$ 中的仿射开邻域，
则有 $Y_i = \lim W \times_S X_i$。因此可以使用与 (1) 的证明完全相同的论证。
\end{proof}

\begin{lemma}
\label{limits-lemma-glueing-near-multiple-closed-points}
设 $S$ 是概形。设 $s_1, \ldots, s_n \in S$ 是两两不同的闭点，并且
$U = S \setminus \{s_1, \ldots, s_n\} \to S$ 拟紧。置
$S_i = \Spec(\mathcal{O}_{S, s_i})$ 及 $U_i = S_i \setminus \{s_i\}$，
则有范畴等价
$$
FP_S \longrightarrow
FP_U \times_{(FP_{U_1} \times \ldots \times FP_{U_n})}
(FP_{S_1} \times \ldots \times FP_{S_n})
$$
，其中 $FP_T$ 是概形 $T$ 上有限呈示概形的范畴。
\end{lemma}

\begin{proof}
当 $n = 1$ 时，这就是引理 \ref{limits-lemma-glueing-near-closed-point}。
当 $n > 1$ 时，可以用完全相同的方法证明本引理，也可以由该引理推出。
例如，假设 $f_i : X_i \to S_i$ 是 $FP_{S_i}$ 的对象，
$f : X \to U$ 是 $FP_U$ 的对象，并且给定同构
$X_i \times_{S_i} U_i = X \times_U U_i$。
由引理 \ref{limits-lemma-glueing-near-closed-point}，可找到有限呈示态射
$f' : X' \to U' = S \setminus \{s_1, \ldots, s_{n - 1}\}$；
它与 $X_i$ 在 $S_i$ 上同构，与 $X$ 在 $U$ 上同构，并且这些同构
与给定同构 $X_i \times_{S_n} U_n = X \times_U U_n$ 相容。
然后可对
$f_i : X_i \to S_i$，$i \leq n - 1$，
$f' : X' \to U'$，以及诱导同构
$X_i \times_{S_i} U_i = X' \times_{U'} U_i$，$i \leq n - 1$
应用归纳法。这证明了本质满射性。全忠实性的证明从略。
\end{proof}

\section{修改的应用}
\label{limits-section-modifications-at-a-point}

\noindent
利用第 \ref{limits-section-change-over-closed-points} 节的结果，
可以用局部环来描述概形在一个闭点上的修改范畴。

\begin{lemma}
\label{limits-lemma-modifications}
设 $S$ 是概形。设 $s \in S$ 是闭点，且
$U = S \setminus \{s\} \to S$ 拟紧。置
$V = \Spec(\mathcal{O}_{S, s}) \setminus \{s\}$，则基变换函子
{\small
$$
\left\{
\begin{matrix}
f : X \to S\text{ of finite presentation} \\
f^{-1}(U) \to U\text{ 为同构}
\end{matrix}
\right\}
\longrightarrow
\left\{
\begin{matrix}
g : Y \to \Spec(\mathcal{O}_{S, s})\text{ of finite presentation} \\
g^{-1}(V) \to V\text{ 为同构}
\end{matrix}
\right\}
$$
}
是范畴等价。
\end{lemma}

\begin{proof}
这是引理 \ref{limits-lemma-glueing-near-closed-point} 的特例。
\end{proof}

\begin{lemma}
\label{limits-lemma-modifications-properties}
记号和假设同引理 \ref{limits-lemma-modifications}。
设 $f : X \to S$ 通过该等价对应于 $g : Y \to \Spec(\mathcal{O}_{S, s})$。
则 $f$ 是分离的、固有的、有限的、\'etale 的，且还可在此加入更多性质，
当且仅当 $g$ 也如此。
\end{lemma}

\begin{proof}
分离、固有、整、有限等性质在基变换下稳定。见
《概形》引理 \ref{schemes-lemma-separated-permanence}
以及《态射》引理 \ref{morphisms-lemma-base-change-proper} 和
\ref{morphisms-lemma-base-change-finite}。
因此，若 $f$ 具有该性质，则 $g$ 也具有该性质。
反向结论由引理 \ref{limits-lemma-glueing-near-point-properties} 得出，
不过这里也给出一个直接证明。
具体而言，若 $g$ 具有该性质，则由引理
\ref{limits-lemma-descend-separated-finite-presentation}、
\ref{limits-lemma-eventually-proper}、
\ref{limits-lemma-descend-finite-finite-presentation} 和
\ref{limits-lemma-descend-etale}，$f$ 在 $s$ 的某个邻域上具有该性质。
由于 $f$ 在 $S \setminus \{s\}$ 上显然具有所给性质，
而该性质可在基上局部检验，故结论成立。
\end{proof}

\begin{remark}
\label{limits-remark-more-general-modification}
上述引理可推广如下。设 $S$ 是概形，$T \subset S$ 是闭子集。
假设存在开邻域的共尾系 $T \subset W_i$，使得
(1) $W_i \setminus T$ 拟紧，并且
(2) $W_i \subset W_j$ 是仿射态射。
则 $W = \lim W_i$ 是包含 $T$ 作为闭子概形的概形。置
$U = X \setminus T$ 和 $V = W \setminus T$。则基变换函子
$$
\left\{
\begin{matrix}
f : X \to S\text{ of finite presentation} \\
f^{-1}(U) \to U\text{ 为同构}
\end{matrix}
\right\}
\longrightarrow
\left\{
\begin{matrix}
g : Y \to W\text{ of finite presentation} \\
g^{-1}(V) \to V\text{ 为同构}
\end{matrix}
\right\}
$$
是范畴等价。若日后需要此结论，我们会把这条注记改为引理并给出详细证明。
\end{remark}

\section{有限型概形的下降}
\label{limits-section-finite-type-quasi-separated}

\noindent
本节继续第 \ref{limits-section-finite-type-closed-in-finite-presentation} 节的主题，
其精神与第 \ref{limits-section-descending-relative} 节所讨论的结果相同。

\begin{situation}
\label{limits-situation-limit-noetherian}
设 $S = \lim_{i \in I} S_i$ 是 Noether 概形的有向系统的极限，
其过渡态射 $S_{i'} \to S_i$ 在 $i' \geq i$ 时为仿射态射。
\end{situation}

\begin{lemma}
\label{limits-lemma-good-diagram}
在情形 \ref{limits-situation-limit-noetherian} 中。
设 $X \to S$ 拟分离且有限型。
则存在 $i \in I$ 及交换图
\begin{equation}
\label{limits-equation-good-diagram}
\vcenter{
\xymatrix{
X \ar[r] \ar[d] & W \ar[d] \\
S \ar[r] & S_i
}
}
\end{equation}
使得 $W \to S_i$ 有限型，并且诱导态射
$X \to S \times_{S_i} W$ 是闭浸入。
\end{lemma}

\begin{proof}
由引理 \ref{limits-lemma-finite-type-closed-in-finite-presentation}，
可以找到闭浸入 $X \to X'$；它是 $S$ 上的态射，其中 $X'$ 是 $S$ 上有限呈示概形。
由引理 \ref{limits-lemma-descend-finite-presentation}，可以找到指标 $i$ 和有限呈示态射
$X'_i \to S_i$，其拉回为 $X'$。置 $W = X'_i$。
\end{proof}

\begin{lemma}
\label{limits-lemma-limit-from-good-diagram}
在情形 \ref{limits-situation-limit-noetherian} 中。
设 $X \to S$ 拟分离且有限型。
给定 $i \in I$ 及如式 (\ref{limits-equation-good-diagram}) 的交换图
$$
\vcenter{
\xymatrix{
X \ar[r] \ar[d] & W \ar[d] \\
S \ar[r] & S_i
}
}
$$
，对 $i' \geq i$，设 $X_{i'}$ 为态射
$X \to S_{i'} \times_{S_i} W$ 的概形论像。则
$X = \lim_{i' \geq i} X_{i'}$。
\end{lemma}

\begin{proof}
由于 $X$ 拟紧且拟分离，取态射
$X \to S_{i'} \times_{S_i} W$ 的概形论像与限制到开子概形可交换
（《态射》引理 \ref{morphisms-lemma-quasi-compact-scheme-theoretic-image}）。
因此可以并且确实假设 $W$ 仿射，且映入仿射开集 $U_i$，后者位于
$S_i$ 中。设 $U \subset S$、$U_{i'} \subset S_{i'}$ 是 $U_i$ 的逆像。
则 $U$、$U_{i'}$、
$S_{i'} \times_{S_i} W = U_{i'} \times_{U_i} W$ 以及
$S \times_{S_i} W = U \times_{U_i} W$ 都是仿射的。
这蕴含 $X$ 仿射，因为 $X \to S \times_{S_i} W$ 是闭浸入。
这还说明环同态
$$
\mathcal{O}(U) \otimes_{\mathcal{O}(U_i)} \mathcal{O}(W) \to
\mathcal{O}(X)
$$
是满射。设 $I$ 为其核。于是 $X_{i'}$ 是下述环的谱：
$$
\mathcal{O}(X_{i'}) =
\mathcal{O}(U_{i'}) \otimes_{\mathcal{O}(U_i)} \mathcal{O}(W)/I_{i'}
$$
，其中 $I_{i'}$ 是理想 $I$ 的逆像（见《态射》例
\ref{morphisms-example-scheme-theoretic-image}）。
由于 $\mathcal{O}(U) = \colim \mathcal{O}(U_{i'})$，
可见 $I = \colim I_{i'}$，从而
$\colim \mathcal{O}(X_{i'}) = \mathcal{O}(X)$。
\end{proof}

\begin{lemma}
\label{limits-lemma-morphism-good-diagram}
在情形 \ref{limits-situation-limit-noetherian} 中。
设 $f : X \to Y$ 是 $S$ 上拟分离有限型概形之间的态射。设
$$
\vcenter{
\xymatrix{
X \ar[r] \ar[d] & W \ar[d] \\
S \ar[r] & S_{i_1}
}
}
\quad\text{且}\quad
\vcenter{
\xymatrix{
Y \ar[r] \ar[d] & V \ar[d] \\
S \ar[r] & S_{i_2}
}
}
$$
是如式 (\ref{limits-equation-good-diagram}) 的交换图。设
$X = \lim_{i \geq i_1} X_i$ 和
$Y = \lim_{i \geq i_2} Y_i$ 是引理 \ref{limits-lemma-limit-from-good-diagram}
给出的相应极限表示。则存在 $i_0 \geq \max(i_1, i_2)$ 以及
$(S_i)_{i \geq i_0}$ 上逆系统的态射
$$
(f_i)_{i \geq i_0} : (X_i)_{i \geq i_0} \to (Y_i)_{i \geq i_0}
$$
使得 $f = \lim_{i \geq i_0} f_i$。
若 $(g_i)_{i \geq i_0} : (X_i)_{i \geq i_0} \to (Y_i)_{i \geq i_0}$
是 $(S_i)_{i \geq i_0}$ 上逆系统的另一个态射，并且
$f = \lim_{i \geq i_0} g_i$，则 $f_i = g_i$ 对所有 $i \gg i_0$ 都成立。
\end{lemma}

\begin{proof}
由于 $V \to S_{i_2}$ 有限呈示且
$X = \lim_{i \geq i_1} X_i$，可以应用命题
\ref{limits-proposition-characterize-locally-finite-presentation}，找到
$i_0 \geq \max(i_1, i_2)$ 以及态射 $h : X_{i_0} \to V$，它位于 $S_{i_2}$ 上，
使得 $X \to X_{i_0} \to V$ 等于 $X \to Y \to V$。
对 $i \geq i_0$，得到实线部分交换的图
$$
\xymatrix{
X \ar[d] \ar[r] &
X_i \ar[r] \ar@{..>}[d] \ar@/_2pc/[dd] |!{[d];[ld]}\hole &
X_{i_0} \ar[d]^h \\
Y \ar[r] \ar[d] & Y_i \ar[r] \ar[d] & V \ar[d] \\
S \ar[r] & S_i \ar[r] & S_{i_0}
}
$$
由于 $X \to X_i$ 的像概形论稠密，而 $Y_i$ 是态射
$Y \to S_i \times_{S_{i_2}} V$ 的概形论像，由该图诱导的态射
$X_i \to S_i \times_{S_{i_2}} V$ 经 $Y_i$ 分解
（《态射》引理 \ref{morphisms-lemma-factor-factor}）。这证明了存在性。

\medskip\noindent
唯一性。设 $E_i \subset X_i$ 是 $f_i$ 与 $g_i$ 的等化子，其中 $i \geq i_0$。
由《概形》引理 \ref{schemes-lemma-where-are-they-equal}，
$E_i$ 是 $X_i$ 的局部闭子概形。由于 $X_i$ 是
$S_i \times_{S_{i_0}} X_{i_0}$ 的闭子概形，而 $Y_i$ 也类似，可见
$$
E_i = X_i \times_{(S_i \times_{S_{i_0}} X_{i_0})} (S_i \times_{S_{i_0}} E_{i_0})
$$
因此，要完成证明，只须说明态射 $X_i \to X_{i_0}$ 经 $E_{i_0}$ 分解，
对某个 $i \geq i_0$ 如此。为此，将使用如下事实：
$X \to X_{i_0}$ 经 $E_{i_0}$ 分解，因为 $f_{i_0}$ 和 $g_{i_0}$
都与 $f$ 相容。由于 $X_i$ 是 Noether 概形，可见底层拓扑空间
$|E_{i_0}|$ 是 $|X_{i_0}|$ 的可构造子集
（《拓扑》引理 \ref{topology-lemma-constructible-Noetherian-space}）。
因此，由引理 \ref{limits-lemma-limit-contained-in-constructible}，
$X_i \to X_{i_0}$ 在集合论意义下经 $E_{i_0}$ 分解，对充分大的 $i$ 如此。
对这样的 $i$，概形论逆像
$(X_i \to X_{i_0})^{-1}(E_{i_0})$ 是 $X_i$ 的闭子概形，并且 $X$
经它分解；该闭子概形因而等于 $X_i$，因为 $X \to X_i$ 的像按构造概形论稠密。
证明完毕。
\end{proof}

\begin{remark}
\label{limits-remark-finite-type-gives-well-defined-system}
在情形 \ref{limits-situation-limit-noetherian} 中，引理 \ref{limits-lemma-good-diagram}、
\ref{limits-lemma-limit-from-good-diagram} 和
\ref{limits-lemma-morphism-good-diagram} 表明，$S$ 上拟分离有限型概形的范畴
等价于 $(S_i)_{i \in I}$ 上某些类型的概形逆系统；具体而言，就是把引理
\ref{limits-lemma-limit-from-good-diagram} 应用于形如
(\ref{limits-equation-good-diagram}) 的图所产生的系统。
例如，给定有限型拟分离态射 $X \to S$，若选取如式
(\ref{limits-equation-good-diagram}) 的两个不同交换图 $X \to V_1 \to S_{i_1}$
和 $X \to V_2 \to S_{i_2}$，则在两个方向上把引理
\ref{limits-lemma-morphism-good-diagram} 应用于 $\text{id}_X$，可知 $X$
的相应极限表示典范同构（允许缩小有向集 $I$）。其余依此类推。
\end{remark}

\begin{lemma}
\label{limits-lemma-morphism-good-diagram-flat}
记号和假设同引理 \ref{limits-lemma-morphism-good-diagram}。
若 $f$ 平坦且有限呈示，则存在 $i_3 \geq i_0$，使得对 $i \geq i_3$，
$f_i$ 平坦，$X_i = Y_i \times_{Y_{i_3}} X_{i_3}$，并且
$X = Y \times_{Y_{i_3}} X_{i_3}$。
\end{lemma}

\begin{proof}
由引理 \ref{limits-lemma-descend-finite-presentation}，可以选取 $i \geq i_2$
以及有限呈示态射 $U \to Y_i$，使得 $X = Y \times_{Y_i} U$
（此处用到了 $f$ 有限呈示）。增大 $i$ 后，可以假设 $U \to Y_i$ 平坦，
见引理 \ref{limits-lemma-descend-flat-finite-presentation}。
如注记 \ref{limits-remark-finite-type-gives-well-defined-system} 所讨论的，
可以并且确实替换定义系统 $(X_i)_{i \geq i_1}$ 时使用的初始图，
替代者是对应于 $X \to U \to S_i$ 的系统。因此，$X_{i'}$ 在
$i' \geq i$ 时定义为态射 $X \to S_{i'} \times_{S_i} U$ 的概形论像。

\medskip\noindent
因为 $U \to Y_i$ 平坦（此处用到了 $f$ 平坦），并且
$X = Y \times_{Y_i} U$，还因为态射 $Y \to Y_i$ 的概形论像是 $Y_i$，
所以态射 $X \to U$ 的概形论像是 $U$
（《态射》引理
\ref{morphisms-lemma-flat-base-change-scheme-theoretic-image}）。
注意，态射 $Y_{i'} \to S_{i'} \times_{S_i} Y_i$ 在 $i' \geq i$ 时
按系统 $Y_j$ 的构造是闭浸入。
于是与上面相同的论证说明，态射 $X \to S_{i'} \times_{S_i} U$
的概形论像等于闭子概形 $Y_{i'} \times_{Y_i} U$。
因此，$X_{i'} = Y_{i'} \times_{Y_i} U$ 对所有 $i' \geq i$ 都成立，
故取 $i_3 = i$ 即得本引理。
\end{proof}

\begin{lemma}
\label{limits-lemma-morphism-good-diagram-smooth}
记号和假设同引理 \ref{limits-lemma-morphism-good-diagram}。
若 $f$ 光滑，则存在 $i_3 \geq i_0$，使得对 $i \geq i_3$，
$f_i$ 光滑。
\end{lemma}

\begin{proof}
合并引理 \ref{limits-lemma-morphism-good-diagram-flat} 和
\ref{limits-lemma-descend-smooth}。
\end{proof}

\begin{lemma}
\label{limits-lemma-morphism-good-diagram-proper}
记号和假设同引理 \ref{limits-lemma-morphism-good-diagram}。
若 $f$ 固有，则存在 $i_3 \geq i_0$，使得对 $i \geq i_3$，
$f_i$ 固有。
\end{lemma}

\begin{proof}
由注记 \ref{limits-remark-finite-type-gives-well-defined-system} 中的讨论，
对本引理是否成立而言，选择适配如式 (\ref{limits-equation-good-diagram}) 之图的
$i_1$ 与 $W$ 并不重要。因此按如下方式选择 $W$。
首先选取闭浸入 $X \to X'$，使得 $X' \to S$ 固有且有限呈示；见引理
\ref{limits-lemma-proper-limit-of-proper-finite-presentation}。
然后选取 $i_3 \geq i_2$ 以及固有态射 $W \to Y_{i_3}$，使得
$X' = Y \times_{Y_{i_3}} W$。这是可能的，因为
$Y = \lim_{i \geq i_2} Y_i$，并可应用引理
\ref{limits-lemma-descend-finite-presentation} 和
\ref{limits-lemma-eventually-proper}。
对这个 $W$，由构造立即可见：当 $i \geq i_3$ 时，概形 $X_i$ 是
$Y_i \times_{Y_{i_3}} W \subset S_i \times_{S_{i_3}} W$ 的闭子概形，
因而在 $Y_i$ 上固有。
\end{proof}

\begin{lemma}
\label{limits-lemma-good-diagram-fibre-product}
在情形 \ref{limits-situation-limit-noetherian} 中，假设有笛卡尔图
$$
\xymatrix{
X^1 \ar[r]_p \ar[d]_q & X^3 \ar[d]^a \\
X^2 \ar[r]^b & X^4
}
$$
，其中各概形在 $S$ 上拟分离且有限型。对每个 $j = 1, 2, 3, 4$，
选取 $i_j \in I$ 以及如式 (\ref{limits-equation-good-diagram}) 的图
$$
\xymatrix{
X^j \ar[r] \ar[d] & W^j \ar[d] \\
S \ar[r] & S_{i_j}
}
$$
。设 $X^j = \lim_{i \geq i_j} X^j_i$ 是引理
\ref{limits-lemma-morphism-good-diagram} 所述的相应极限表示。
设 $(a_i)_{i \geq i_5}$、$(b_i)_{i \geq i_6}$、$(p_i)_{i \geq i_7}$ 和
$(q_i)_{i \geq i_8}$ 是引理 \ref{limits-lemma-morphism-good-diagram}
所构造的相应系统态射。则存在
$i_9 \geq \max(i_5, i_6, i_7, i_8)$，使得对 $i \geq i_9$ 有
$a_i \circ p_i = b_i \circ q_i$，并且
$$
(q_i, p_i) : X^1_i \longrightarrow X^2_i \times_{b_i, X^4_i, a_i} X^3_i
$$
是闭浸入。
若 $a$ 与 $b$ 平坦且有限呈示，则存在
$i_{10} \geq \max(i_5, i_6, i_7, i_8, i_9)$，使得对 $i \geq i_{10}$，
最后显示的态射是同构。
\end{lemma}

\begin{proof}
由注记 \ref{limits-remark-finite-type-gives-well-defined-system} 中的讨论，
对本引理是否成立而言，选择适配如式 (\ref{limits-equation-good-diagram}) 之图的
$W^1$ 并不重要。因此可以选取 $W^1 = W^2 \times_{W^4} W^3$。
由 $X^1_i$ 的构造立即可见 $a_i \circ p_i = b_i \circ q_i$，并且
$$
(q_i, p_i) : X^1_i \longrightarrow X^2_i \times_{b_i, X^4_i, a_i} X^3_i
$$
是闭浸入。

\medskip\noindent
若 $a$ 与 $b$ 平坦且有限呈示，则 $p$ 与 $q$ 也平坦且有限呈示，
因为它们是 $a$ 与 $b$ 的基变换。因此可以把引理
\ref{limits-lemma-morphism-good-diagram-flat} 分别应用于 $a$、$b$、$p$、$q$
以及 $a \circ p = b \circ q$。由此存在 $i_9 \in I$，使得
$$
(q_i, p_i) : X^1_i \to X^2_i \times_{X^4_i} X^3_i
$$
是 $(q_{i_9}, p_{i_9})$ 沿态射 $X^4_i \to X^4_{i_9}$ 的基变换，
对所有 $i \geq i_9$ 如此。由引理 \ref{limits-lemma-descend-isomorphism}，
可知 $(q_i, p_i)$ 对所有充分大的 $i$ 都是同构。
\end{proof}
% Generated mechanically from the frozen varieties.tex target.
% Do not edit this generated file; edit the bound locale source instead.

% OK, start here.
%

\setcounter{chapter}{32}
\chapter{簇}



\phantomsection
\label{varieties-section-phantom}


\section{引言}
\label{varieties-section-introduction}

\noindent
本章开始研究簇，并更一般地研究域上的概形。一个基本参考文献是
\cite{EGA}。








\section{记号}
\label{varieties-section-notation}

\noindent
本章始终用字母 $k$ 表示基域。








\section{簇}
\label{varieties-section-varieties}

\noindent
在 Stacks 计划中，我们采用如下簇的定义。

\begin{definition}
\label{varieties-definition-variety}
设 $k$ 为域。所谓一个{\it 簇}，是指 $X$ 是域 $k$ 上的概形，且
$X$ 为整概形，且结构态射
$X \to \Spec(k)$ 分离并且是有限型的。
\end{definition}

\noindent
这个定义有如下缺点。设 $k'/k$ 是域扩张，并设 $X$ 是域 $k$ 上的簇。
则基变换
$X_{k'} = X \times_{\Spec(k)} \Spec(k')$
未必是域 $k'$ 上的簇。后续各节将以更一般的形式详细讨论这一现象。
两个簇的积未必是簇（本质上这正是同一现象）。下面给出一个例子。

\begin{example}
\label{varieties-example-product-not-a-variety}
令 $k = \mathbf{Q}$。令 $X = \Spec(\mathbf{Q}(i))$ 且
$Y = \Spec(\mathbf{Q}(i))$。则积
$X \times_{\Spec(k)} Y$，亦即簇 $X$ 与 $Y$ 的积，不是簇，因为它是可约的。（它同构于两个
$X$ 的不交并。）
\end{example}

\noindent
然而，如果基域是代数闭域，那么簇的积仍是簇。这可由《代数》章中的
结果推出，不过那里处理的是远为一般的情形。这里也有一个简单的直接
证明，我们将其写出。

\begin{lemma}
\label{varieties-lemma-product-varieties}
\begin{slogan}
代数闭域上簇的积仍是簇。
\end{slogan}
设 $k$ 为代数闭域。设 $X$, $Y$ 为域 $k$ 上的簇。则
$X \times_{\Spec(k)} Y$ 是域 $k$ 上的簇。
\end{lemma}

\begin{proof}
态射 $X \times_{\Spec(k)} Y \to \Spec(k)$ 是有限型且分离的，
因为它是态射
$X \times_{\Spec(k)} Y \to Y \to \Spec(k)$
的复合，而这两个态射都分离且为有限型；见《态射》引理
\ref{morphisms-lemma-base-change-finite-type} 和
\ref{morphisms-lemma-composition-finite-type}，以及《概形》引理
\ref{schemes-lemma-separated-permanence}。
为完成证明，只需证明 $X \times_{\Spec(k)} Y$ 是整概形。
取有限仿射开覆盖
$X = \bigcup_{i = 1, \ldots, n} U_i$、
$Y = \bigcup_{j = 1, \ldots, m} V_j$。如果能证明每个
$U_i \times_{\Spec(k)} V_j$ 都是整概形，则由《性质》引理
\ref{properties-lemma-characterize-reduced}、
\ref{properties-lemma-characterize-irreducible} 和
\ref{properties-lemma-characterize-integral} 即得结论。
所以问题约化到仿射情形。

\medskip\noindent
仿射情形可转化为如下代数命题：设 $A$, $B$ 是整环且为有限生成
$k$-代数，则 $A \otimes_k B$ 是整环。反设不然，并设
$$
(\sum\nolimits_{i = 1, \ldots, n} a_i \otimes b_i)
(\sum\nolimits_{j = 1, \ldots, m} c_j \otimes d_j) = 0
$$
在 $A \otimes_k B$ 中成立，且两个因子在 $A \otimes_k B$ 中均非零。
不妨设 $b_1, \ldots, b_n$ 关于 $k$ 在 $B$ 中线性无关，并且
$d_1, \ldots, d_m$ 关于 $k$ 在 $B$ 中线性无关。当然，也可设
$a_1$ 与 $c_1$ 在 $A$ 中非零。因此
$D(a_1c_1) \subset \Spec(A)$ 非空。由 Hilbert 零点定理
（《代数》定理 \ref{algebra-theorem-nullstellensatz}），可找到极大理想
$\mathfrak m \subset A$，它包含于 $D(a_1c_1)$，且
$A/\mathfrak m = k$；这是因为 $k$ 代数闭。用
$\overline{a}_i, \overline{c}_j$ 表示 $a_i, c_j$ 在
$A/\mathfrak m = k$ 中的剩余类。上述等式变为
$$
(\sum\nolimits_{i = 1, \ldots, n} \overline{a}_i b_i)
(\sum\nolimits_{j = 1, \ldots, m} \overline{c}_j d_j) = 0
$$
，这与 $\mathfrak m \in D(a_1c_1)$、
$b_1, \ldots, b_n$ 和 $d_1, \ldots, d_m$ 的线性无关性以及
$B$ 是整环这一事实矛盾。
\end{proof}






\section{簇与有理映射}
\label{varieties-section-varieties-rational-maps}

\noindent
设 $k$ 为域。设 $X$ 和 $Y$ 为域 $k$ 上的簇。我们所说的
{\it 从 $X$ 到 $Y$ 的簇有理映射}，是指按照《态射》定义
\ref{morphisms-definition-rational-map} 的 $\Spec(k)$-有理映射，
其源概形为 $X$，目标概形为 $Y$。依照惯例，“簇有理映射”这一说法并不指明这些簇
（共同）的基域；但对于一般概形，我们会区分有理映射和给定基上的
有理映射。

\medskip\noindent
本节标题所指的是下面这个基本定理。

\begin{theorem}
\label{varieties-theorem-varieties-rational-maps}
设 $k$ 为域。由簇及占优有理映射构成的范畴等价于有限生成域扩张
$K/k$ 构成的范畴。
\end{theorem}

\begin{proof}
设 $X$ 和 $Y$ 为簇，其一般点分别为 $x \in X$ 和 $y \in Y$。
回忆，从 $X$ 到 $Y$ 的占优有理映射，恰好是把 $x$ 映到 $y$ 的那些
有理映射（见《态射》定义 \ref{morphisms-definition-dominant-rational}
及其后的讨论）。因此，给定占优有理映射 $X \supset U \to Y$，便得到
函数域的映射
$$
k(Y) = \kappa(y) = \mathcal{O}_{Y, y}
\longrightarrow
\mathcal{O}_{X, x} = \kappa(x) = k(X)
$$
反之，这样的 $k$-代数映射（由于源和目标都是域，它自动是局部的）
由《态射》引理 \ref{morphisms-lemma-rational-map-finite-presentation}
唯一确定一个占优有理映射。由此得到一个全忠实函子。为完成证明，
只需证明每个有限生成域扩张 $K/k$ 都在其本质像中。由于 $K/k$ 有限生成，
存在有限型 $k$-代数 $A \subset K$，使 $K$ 是 $A$ 的分式域。于是
$X = \Spec(A)$ 是函数域为 $K$ 的簇。
\end{proof}

\noindent
设 $k$ 为域。设 $X$ 和 $Y$ 为域 $k$ 上的簇。我们所说的
{\it $X$ 与 $Y$ 是双有理簇}，是指按照《态射》定义
\ref{morphisms-definition-birational}，$X$ 与 $Y$ 是
$\Spec(k)$-双有理的。依照惯例，“双有理簇”这一说法并不指明这些簇
（共同）的基域；但对于一般不可约概形，我们会区分双有理和给定基上
双有理。

\begin{lemma}
\label{varieties-lemma-birational-varieties}
设 $X$ 与 $Y$ 是域 $k$ 上的簇。下列条件等价：
\begin{enumerate}
\item $X$ 与 $Y$ 是双有理簇；
\item 函数域 $k(X)$ 与 $k(Y)$ 同构；
\item $X$ 与 $Y$ 各有一个非空开集，二者作为簇同构；
\item 存在开集 $U \subset X$ 以及簇的双有理态射 $U \to Y$。
\end{enumerate}
\end{lemma}

\begin{proof}
这是《态射》引理
\ref{morphisms-lemma-criterion-birational-finite-presentation} 的特殊情形。
\end{proof}






\section{域变换与局部环}
\label{varieties-section-local-rings}

\noindent
先给出基域扩张时局部环所发生变化的一些预备结果。

\begin{lemma}
\label{varieties-lemma-change-fields-flat}
设 $K/k$ 为域扩张。设 $X$ 为域 $k$ 上的概形，并令 $Y = X_K$。
若 $y \in Y$ 的像为 $x \in X$，则：
\begin{enumerate}
\item $\mathcal{O}_{X, x} \to \mathcal{O}_{Y, y}$ 是忠实平坦的局部环同态；
\item 令 $\mathfrak p_0 = \Ker(\kappa(x) \otimes_k K \to \kappa(y))$，
则 $\kappa(y) = \kappa(\mathfrak p_0)$；
\item $\mathcal{O}_{Y, y} = (\mathcal{O}_{X, x} \otimes_k K)_\mathfrak p$，
其中 $\mathfrak p \subset \mathcal{O}_{X, x} \otimes_k K$ 是
$\mathfrak p_0$ 的逆像；
\item 有
$\mathcal{O}_{Y, y}/\mathfrak m_x\mathcal{O}_{Y, y} =
(\kappa(x) \otimes_k K)_{\mathfrak p_0}$。
\end{enumerate}
\end{lemma}

\begin{proof}
不妨设 $X = \Spec(A)$ 为仿射概形。于是 $Y = \Spec(A \otimes_k K)$。
由于 $K$ 在 $k$ 上平坦，$A \to A \otimes_k K$ 平坦。因此
$Y \to X$ 平坦；再用《代数》引理
\ref{algebra-lemma-local-flat-ff}，便得到第一项。第二项由《概形》引理
\ref{schemes-lemma-points-fibre-product} 得出。现在 $y$ 对应于素理想
$\mathfrak q \subset A \otimes_k K$，而 $x$ 对应于
$\mathfrak r = A \cap \mathfrak q$。于是 $\mathfrak p_0$ 是诱导映射
$\kappa(\mathfrak r) \otimes_k K \to \kappa(\mathfrak q)$ 的核。
局部环上的映射为
$$
A_\mathfrak r \longrightarrow (A \otimes_k K)_\mathfrak q
$$
我们可将此映射分解经过
$A_\mathfrak r \otimes_k K = (A \otimes_k K)_{\mathfrak r}$，得到
$$
A_\mathfrak r \longrightarrow A_\mathfrak r \otimes_k K
\longrightarrow (A \otimes_k K)_\mathfrak q
$$
而第二个箭头是关于某个素理想的局部化。这个素理想正是
$\mathfrak p_0$ 的逆像（略去细节），这就证明了 (3)。为证明 (4)，
使用 (3)，以及局部化和 $- \otimes_k K$ 都是正合函子这一事实。
\end{proof}

\begin{lemma}
\label{varieties-lemma-change-fields-algebraic-dim}
采用引理 \ref{varieties-lemma-change-fields-flat} 的记号。设 $X$ 在 $k$ 上局部有限型。
则
$$
\dim(\mathcal{O}_{Y, y}/\mathfrak m_x\mathcal{O}_{Y, y}) =
\text{trdeg}_k(\kappa(x)) - \text{trdeg}_K(\kappa(y)) =
\dim(\mathcal{O}_{Y, y}) - \dim(\mathcal{O}_{X, x})
$$
\end{lemma}

\begin{proof}
这只是《代数》引理
\ref{algebra-lemma-inequalities-under-field-extension} 的改述。
\end{proof}

\begin{lemma}
\label{varieties-lemma-change-fields-algebraic-unramified}
采用引理 \ref{varieties-lemma-change-fields-flat} 的记号。设 $X$ 在 $k$ 上局部有限型，
并且 $\dim(\mathcal{O}_{X, x}) = \dim(\mathcal{O}_{Y, y})$，且
$\kappa(x) \otimes_k K$ 约化（例如，当 $\kappa(x)/k$ 可分或
$K/k$ 可分时）。则 $\mathfrak m_x \mathcal{O}_{Y, y} = \mathfrak m_y$。
\end{lemma}

\begin{proof}
（括号内的断言由《代数》引理
\ref{algebra-lemma-separable-extension-preserves-reducedness} 得出。）
结合引理 \ref{varieties-lemma-change-fields-flat} 与
\ref{varieties-lemma-change-fields-algebraic-dim}，可知
$\mathcal{O}_{Y, y}/\mathfrak m_x \mathcal{O}_{Y, y}$ 的维数为 $0$，
并且它约化。因此它是域。
\end{proof}





\section{几何约化概形}
\label{varieties-section-geometrically-reduced}

\noindent
如果 $X$ 是域上的约化概形，那么扩张基域之后，$X$ 可能变成非约化的。
对于几何约化概形，这种情况不会发生。

\begin{definition}
\label{varieties-definition-geometrically-reduced}
设 $k$ 为域。设 $X$ 为 $k$ 上的概形。
\begin{enumerate}
\item 设 $x \in X$ 为一点。称 $X$ 在 $x$ 处{\it 几何约化}，如果对于
任意域扩张 $k'/k$ 以及任意点 $x' \in X_{k'}$，只要它位于 $x$ 上方，
局部环 $\mathcal{O}_{X_{k'}, x'}$ 就约化。
\item 称 $X$ 在 $k$ 上{\it 几何约化}，如果 $X$ 在 $X$ 的每一点处
都几何约化。
\end{enumerate}
\end{definition}

\noindent
初看起来，这个概念可能略显神秘，但它实际上正是《代数》章中讨论的
概念。下面给出一些说明其联系的基本结果。另见引理
\ref{varieties-lemma-geometrically-reduced-dense-smooth-open}；该引理说明，对域上的簇，
“几何约化”等价于“一般光滑”。

\begin{lemma}
\label{varieties-lemma-geometrically-reduced-at-point}
\begin{slogan}
几何约化性可在局部环上检验。
\end{slogan}
设 $k$ 为域。设 $X$ 为 $k$ 上的概形。设 $x \in X$。下列条件等价：
\begin{enumerate}
\item $X$ 在 $x$ 处几何约化；
\item 环 $\mathcal{O}_{X, x}$ 在 $k$ 上几何约化（见《代数》定义
\ref{algebra-definition-geometrically-reduced}）。
\end{enumerate}
\end{lemma}

\begin{proof}
设 (1) 成立。这尤其说明 $\mathcal{O}_{X, x}$ 约化。设 $k'/k$ 为有限
纯不可分域扩张。考虑环 $\mathcal{O}_{X, x} \otimes_k k'$。由《代数》引理
\ref{algebra-lemma-p-ring-map}，它的谱与 $\mathcal{O}_{X, x}$ 的谱相同，
故它也是局部环（《代数》引理
\ref{algebra-lemma-characterize-local-ring}）。因此存在唯一的点
$x' \in X_{k'}$ 位于 $x$ 上方，并且
$\mathcal{O}_{X_{k'}, x'} \cong \mathcal{O}_{X, x} \otimes_k k'$；见引理
\ref{varieties-lemma-change-fields-flat}。由假设，这是约化环。因此，由《代数》引理
\ref{algebra-lemma-geometrically-reduced-finite-purely-inseparable-extension}
推出 (2)。

\medskip\noindent
设 (2) 成立。设 $k'/k$ 为域扩张。由于
$\Spec(k') \to \Spec(k)$ 满射，故
$X_{k'} \to X$ 也满射（《态射》引理
\ref{morphisms-lemma-base-change-surjective}）。设 $x' \in X_{k'}$ 是位于
$x$ 上方的任意点。由引理 \ref{varieties-lemma-change-fields-flat}，局部环
$\mathcal{O}_{X_{k'}, x'}$ 是环
$\mathcal{O}_{X, x} \otimes_k k'$ 的局部化。因此由假设它约化，(1) 得证。
\end{proof}

\noindent
这个概念在特征零情形没有新的内容。

\begin{lemma}
\label{varieties-lemma-perfect-reduced}
设 $X$ 是完全域 $k$（例如，$k$ 的特征为零）上的概形。设 $x \in X$。
若 $\mathcal{O}_{X, x}$ 约化，则 $X$ 在 $x$ 处几何约化。
若 $X$ 约化，则 $X$ 在 $k$ 上几何约化。
\end{lemma}

\begin{proof}
第一条断言由引理 \ref{varieties-lemma-geometrically-reduced-at-point}、《代数》引理
\ref{algebra-lemma-separable-extension-preserves-reducedness} 以及完全域的定义
（《代数》定义 \ref{algebra-definition-perfect}）得出。第二条断言由第一条得出。
\end{proof}

\begin{lemma}
\label{varieties-lemma-geometrically-reduced}
设 $k$ 为特征 $p > 0$ 的域。设 $X$ 为 $k$ 上的概形。下列条件等价：
\begin{enumerate}
\item $X$ 几何约化；
\item $X_{k'}$ 对每个域扩张 $k'/k$ 都约化；
\item $X_{k'}$ 对每个有限纯不可分域扩张 $k'/k$ 都约化；
\item $X_{k^{1/p}}$ 约化；
\item $X_{k^{perf}}$ 约化；
\item $X_{\bar k}$ 约化；
\item 对每个仿射开集 $U \subset X$，环 $\mathcal{O}_X(U)$ 都几何约化
（见《代数》定义 \ref{algebra-definition-geometrically-reduced}）。
\end{enumerate}
\end{lemma}

\begin{proof}
设 (1) 成立。于是，对每个域扩张 $k'/k$ 和每个点
$x' \in X_{k'}$，$X_{k'}$ 在 $x'$ 处的局部环都约化。换言之，
$X_{k'}$ 约化。因此 (2) 成立。

\medskip\noindent
设 (2) 成立。令 $U \subset X$ 为仿射开集。于是，对每个域扩张
$k'/k$，概形 $X_{k'}$ 都约化，因而
$U_{k'} = \Spec(\mathcal{O}(U)\otimes_k k')$ 约化，继而
$\mathcal{O}(U)\otimes_k k'$ 约化（见《性质》第
\ref{properties-section-integral} 节）。换言之，$\mathcal{O}(U)$ 几何约化，
故 (7) 成立。

\medskip\noindent
设 (7) 成立。对任意域扩张 $k'/k$，基变换 $X_{k'}$ 可由环
$\mathcal{O}_X(U) \otimes_k k'$ 的谱粘合而得，其中 $U$ 遍历 $X$ 中的
仿射开集（见《概形》第 \ref{schemes-section-fibre-products} 节）。因此
$X_{k'}$ 约化，故 (1) 成立。

\medskip\noindent
这证明了 (1)、(2) 与 (7) 等价。它们又等价于 (3)、(4)、(5) 与 (6)，
因为可以对 $\mathcal{O}_X(U)$ 应用《代数》引理，其中 $U \subset X$
为仿射开集：
\ref{algebra-lemma-geometrically-reduced-finite-purely-inseparable-extension}。
\end{proof}

\begin{lemma}
\label{varieties-lemma-check-only-finite-inseparable-extensions}
设 $k$ 为特征 $p > 0$ 的域。设 $X$ 为 $k$ 上的概形。设 $x \in X$。
下列条件等价：
\begin{enumerate}
\item $X$ 在 $x$ 处几何约化；
\item 局部环 $\mathcal{O}_{X_{k'}, x'}$ 约化，其中 $k'$ 是 $k$ 的任一
有限纯不可分域扩张，且 $x' \in X_{k'}$ 是位于 $x$ 上方的唯一点；
\item 局部环 $\mathcal{O}_{X_{k^{1/p}}, x'}$ 约化，其中
$x' \in X_{k^{1/p}}$ 是位于 $x$ 上方的唯一点；
\item 局部环 $\mathcal{O}_{X_{k^{perf}}, x'}$ 约化，其中
$x' \in X_{k^{perf}}$ 是位于 $x$ 上方的唯一点。
\end{enumerate}
\end{lemma}

\begin{proof}
注意，若 $k'/k$ 纯不可分，则 $X_{k'} \to X$ 在底拓扑空间上诱导同胚；
见《代数》引理 \ref{algebra-lemma-p-ring-map}。这说明了命题中 $x'$ 位于
$x$ 上方时的唯一性。此外，在这种情况下，
$\mathcal{O}_{X_{k'}, x'} = \mathcal{O}_{X, x} \otimes_k k'$。
所以本引理由上面的引理 \ref{varieties-lemma-geometrically-reduced-at-point} 以及
《代数》引理
\ref{algebra-lemma-geometrically-reduced-finite-purely-inseparable-extension}
得出。
\end{proof}

\begin{lemma}
\label{varieties-lemma-geometrically-reduced-upstairs}
设 $k$ 为域。设 $X$ 为 $k$ 上的概形。设 $k'/k$ 为域扩张。
设 $x \in X$ 为一点，并设 $x' \in X_{k'}$ 为位于 $x$ 上方的一点。
下列条件等价：
\begin{enumerate}
\item $X$ 在 $x$ 处几何约化；
\item $X_{k'}$ 在 $x'$ 处几何约化。
\end{enumerate}
特别地，$X$ 在 $k$ 上几何约化，当且仅当 $X_{k'}$ 在 $k'$ 上几何约化。
\end{lemma}

\begin{proof}
显然 (1) 蕴含 (2)。设 (2) 成立。令 $k''/k$ 为有限纯不可分域扩张，
并令 $x'' \in X_{k''}$ 为位于 $x$ 上方的一点（事实上，它是唯一的）。
选取 $k'''/k$，使它是 $k''/k$ 与 $k'/k$ 的公共域扩张；见《域》引理
\ref{fields-lemma-common-extension-field}。选取一点 $x''' \in X_{k'''}$，
它同时位于 $x'$ 与 $x''$ 上方。考虑局部环映射
$$
\mathcal{O}_{X_{k''}, x''} \longrightarrow \mathcal{O}_{X_{k'''}, x'''}.
$$
这是平坦局部环同态，因而忠实平坦。由 (2)，右侧局部环约化。
于是由《代数》引理 \ref{algebra-lemma-descent-reduced} 可知
$\mathcal{O}_{X_{k''}, x''}$ 约化。因此，由引理
\ref{varieties-lemma-check-only-finite-inseparable-extensions} 可知 $X$ 在 $x$ 处
几何约化。
\end{proof}

\begin{lemma}
\label{varieties-lemma-geometrically-reduced-any-base-change}
设 $k$ 为域。设 $X$、$Y$ 为 $k$ 上的概形。
\begin{enumerate}
\item 若 $X$ 在 $x$ 处几何约化，且 $Y$ 约化，则 $X \times_k Y$ 在每个
位于 $x$ 上方的点处都约化。
\item 若 $X$ 在 $k$ 上几何约化，且 $Y$ 约化，则 $X \times_k Y$ 约化。
\item 若 $X$ 与 $Y$ 在 $k$ 上几何约化，则 $X \times_k Y$ 几何约化。
\item 若 $k$ 为完全域，且 $X$ 与 $Y$ 约化，则 $X \times_k Y$ 约化。
\item 此处补充更多内容。
\end{enumerate}
\end{lemma}

\begin{proof}
为证明 (1)，结合引理 \ref{varieties-lemma-geometrically-reduced-at-point} 与《代数》引理
\ref{algebra-lemma-geometrically-reduced-any-reduced-base-change}。
为证明 (2)，结合引理 \ref{varieties-lemma-geometrically-reduced} 与《代数》引理
\ref{algebra-lemma-geometrically-reduced-any-reduced-base-change}。
为证明 (3)，注意
$(X \times_k Y)_{\overline{k}} =
X_{\overline{k}} \times_{\overline{k}} Y_{\overline{k}}$，
并使用 (2) 及引理 \ref{varieties-lemma-geometrically-reduced}。
为证明 (4)，结合使用 (3) 与引理 \ref{varieties-lemma-perfect-reduced}。
\end{proof}

\begin{lemma}
\label{varieties-lemma-generic-points-geometrically-reduced}
设 $k$ 为域。设 $X$ 为 $k$ 上的概形。
\begin{enumerate}
\item 若 $x' \leadsto x$ 是一次特化，且 $X$ 在 $x$ 处几何约化，则 $X$
在 $x'$ 处几何约化。
\item 设 $x \in X$ 满足：(a) $\mathcal{O}_{X, x}$ 约化；并且 (b) 对每个
特化 $x' \leadsto x$，若 $x'$ 是 $X$ 的一个不可约分支的泛点，则 $X$
在 $x'$ 处几何约化。那么 $X$ 在 $x$ 处几何约化。
\item 若 $X$ 约化，且在 $X$ 的不可约分支的所有泛点处都几何约化，
则 $X$ 几何约化。
\item 若 $X$ 是 $k$ 上的簇，则 $X$ 几何约化，当且仅当其函数域是 $k$
的可分扩张。
\end{enumerate}
\end{lemma}

\begin{proof}
第 (1) 部分由引理 \ref{varieties-lemma-geometrically-reduced-at-point} 及以下事实得出：
若 $A$ 为几何约化 $k$-代数，则 $S^{-1}A$ 是几何约化 $k$-代数，
其中 $S$ 是 $A$ 的任意乘法子集；见《代数》引理
\ref{algebra-lemma-geometrically-reduced-permanence}。

\medskip\noindent
令 $A = \mathcal{O}_{X, x}$。第 (2) 部分的假设 (a) 与 (b) 蕴含 $A$
约化，并且 $A_{\mathfrak q}$ 在 $k$ 上几何约化，其中 $\mathfrak q$
遍历 $A$ 的极小素理想。因此 $A$ 在 $k$ 上几何约化；见《代数》引理
\ref{algebra-lemma-generic-points-geometrically-reduced}。由此，$X$ 在 $x$
处几何约化；见引理 \ref{varieties-lemma-geometrically-reduced-at-point}。

\medskip\noindent
第 (3) 部分显然由第 (2) 部分得出。最后，第 (4) 部分由第 (3) 部分及
《代数》引理 \ref{algebra-lemma-characterize-separable-field-extensions}
得出。
\end{proof}

\begin{lemma}
\label{varieties-lemma-Noetherian-geometrically-reduced-at-point}
设 $k$ 为域。设 $X$ 为 $k$ 上的概形。设 $x \in X$。假设 $X$ 局部
Noether，且在 $x$ 处几何约化。那么存在开邻域 $U \subset X$，它包含
$x$，并且在 $k$ 上几何约化。
\end{lemma}

\begin{proof}
假设 $X$ 局部 Noether，且在 $x$ 处几何约化。由《性质》引理
\ref{properties-lemma-ring-affine-open-injective-local-ring}，可以找到
仿射开邻域 $U \subset X$，它包含 $x$，并且使得
$R = \mathcal{O}_X(U) \to \mathcal{O}_{X, x}$ 为单射。由引理
\ref{varieties-lemma-geometrically-reduced-at-point}，假设意味着
$\mathcal{O}_{X, x}$ 在 $k$ 上几何约化。由《代数》引理
\ref{algebra-lemma-subalgebra-separable}，这蕴含 $R$ 在 $k$ 上几何约化，
进而蕴含 $U$ 几何约化。
\end{proof}

\begin{example}
\label{varieties-example-not-geometrically-reduced}
设 $k = \mathbf{F}_p(s, t)$，即素域的纯超越扩张。考虑簇
$X = \Spec(k[x, y]/(1 + sx^p + ty^p))$。设 $k'/k$ 为任一扩张，
使 $s$ 与 $t$ 都有 $p$ 次根，且这些根属于 $k'$。那么基变换 $X_{k'}$
非约化。
事实上，环 $k'[x, y]/(1 + s x^p + ty^p)$ 包含元素
$1 + s^{1/p}x + t^{1/p}y$；它的 $p$ 次幂为零，但它本身不为零
（因为理想 $(1 + sx^p + ty^p)$ 显然不含任何次数 $< p$ 的非零元素）。
\end{example}

\begin{lemma}
\label{varieties-lemma-finite-extension-geometrically-reduced}
设 $k$ 为域。设 $X \to \Spec(k)$ 局部有限型。假设 $X$ 只有有限多个
不可约分支。那么存在有限纯不可分扩张 $k'/k$，使得 $(X_{k'})_{red}$
在 $k'$ 上几何约化。
\end{lemma}

\begin{proof}
为证明本引理，可以将 $X$ 替换为其约化 $X_{red}$。因此可以假设 $X$
约化且在 $k$ 上局部有限型。令 $x_1, \ldots, x_n \in X$ 为 $X$ 的
各不可约分支的泛点。注意，对每个纯不可分代数扩张 $k'/k$，态射
$(X_{k'})_{red} \to X$ 都是同胚；见《代数》引理
\ref{algebra-lemma-p-ring-map}。因此，点 $x'_1, \ldots, x'_n$（分别位于
$x_1, \ldots, x_n$ 上方）是 $(X_{k'})_{red}$ 的各不可约分支的泛点。
由于 $X$ 约化，局部环 $K_i = \mathcal{O}_{X, x_i}$ 都是域；见《代数》引理
\ref{algebra-lemma-minimal-prime-reduced-ring}。由于 $X$ 在 $k$ 上局部有限型，
域扩张 $K_i/k$ 都有限生成。最后，局部环
$\mathcal{O}_{(X_{k'})_{red}, x'_i}$ 是域 $(K_i \otimes_k k')_{red}$。
由《代数》引理 \ref{algebra-lemma-make-separable}，可以找到有限纯不可分扩张
$k'/k$，使得 $(K_i \otimes_k k')_{red}$ 都是 $k'$ 的可分域扩张。
特别地，由《代数》引理
\ref{algebra-lemma-characterize-separable-field-extensions}，每个
$(K_i \otimes_k k')_{red}$ 都在 $k'$ 上几何约化。此时，引理
\ref{varieties-lemma-generic-points-geometrically-reduced} 的第 (3) 部分蕴含
$(X_{k'})_{red}$ 几何约化。
\end{proof}

\section{几何连通概形}
\label{varieties-section-geometrically-connected}

\noindent
如果 $X$ 是域上的连通概形，那么扩张基域之后，$X$ 可能变成不连通的。
对于几何连通概形，这种情况不会发生。

\begin{definition}
\label{varieties-definition-geometrically-connected}
设 $X$ 为域 $k$ 上的概形。称 $X$ 在 $k$ 上{\it 几何连通}，如果概形
$X_{k'}$ 对每个域 $k'$（作为 $k$ 的扩张）都连通。
\end{definition}

\noindent
按照约定，连通拓扑空间非空；因此更不必说，几何连通概形非空。
下面给出一个并非几何连通的簇的例子。

\begin{example}
\label{varieties-example-not-geometrically-irreducible}
设 $k = \mathbf{Q}$。概形
$X = \Spec(\mathbf{Q}(i))$ 是 $\Spec(\mathbf{Q})$ 上的簇。
但是基变换 $X_{\mathbf{C}}$ 是
$\mathbf{C} \otimes_{\mathbf{Q}} \mathbf{Q}(i) \cong
\mathbf{C} \times \mathbf{C}$ 的谱，也就是两个 $\Spec(\mathbf{C})$
的不交并。因此，这确实是一个非几何连通簇的例子。
\end{example}

\begin{lemma}
\label{varieties-lemma-geometrically-connected-check-after-extension}
设 $X$ 为域 $k$ 上的概形。设 $k'/k$ 为域扩张。那么 $X$ 在 $k$ 上
几何连通，当且仅当 $X_{k'}$ 在 $k'$ 上几何连通。
\end{lemma}

\begin{proof}
如果 $X$ 在 $k$ 上几何连通，那么显然 $X_{k'}$ 在 $k'$ 上几何连通。
反之，对任意域扩张 $k''/k$，存在公共域扩张 $k'''/k'$ 与 $k'''/k''$；
见《域》引理 \ref{fields-lemma-common-extension-field}。由于态射
$X_{k'''} \to X_{k''}$ 满射（它是域的谱之间满射的基变换），
$X_{k'''}$ 的连通性蕴含 $X_{k''}$ 的连通性。因此，如果 $X_{k'}$ 在
$k'$ 上几何连通，那么 $X$ 在 $k$ 上几何连通。
\end{proof}

\begin{lemma}
\label{varieties-lemma-bijection-connected-components}
设 $k$ 为域。设 $X$、$Y$ 为 $k$ 上的概形。假设 $X$ 在 $k$ 上
几何连通。那么投影态射
$$
p : X \times_k Y \longrightarrow Y
$$
诱导连通分支之间的双射。
\end{lemma}

\begin{proof}
$p$ 的概形论纤维连通，因为它们是几何连通概形 $X$ 沿域扩张的
基变换。此外，概形论纤维与集合论纤维同胚；见《概形》引理
\ref{schemes-lemma-fibre-topological}。由《态射》引理
\ref{morphisms-lemma-scheme-over-field-universally-open}，映射 $p$ 是开的。
因此可以应用《拓扑学》引理
\ref{topology-lemma-connected-fibres-connected-components} 得出结论。
\end{proof}

\begin{lemma}
\label{varieties-lemma-affine-geometrically-connected}
设 $k$ 为域。设 $A$ 为 $k$-代数。那么 $X = \Spec(A)$ 在 $k$ 上
几何连通，当且仅当 $A$ 在 $k$ 上几何连通（见《代数》定义
\ref{algebra-definition-geometrically-connected}）。
\end{lemma}

\begin{proof}
这直接来自定义。
\end{proof}

\begin{lemma}
\label{varieties-lemma-separably-closed-field-connected-components}
设 $k'/k$ 为域扩张。设 $X$ 为 $k$ 上的概形。假设 $k$ 可分代数闭。
那么态射 $X_{k'} \to X$ 诱导连通分支之间的双射。特别地，$X$ 在 $k$
上几何连通，当且仅当 $X$ 连通。
\end{lemma}

\begin{proof}
由于 $k$ 可分代数闭，可知 $k'$ 在 $k$ 上几何连通；见《代数》引理
\ref{algebra-lemma-separably-closed-connected-implies-geometric}。
因此，由上面的引理 \ref{varieties-lemma-affine-geometrically-connected}，
$Z = \Spec(k')$ 在 $k$ 上几何连通。由于
$X_{k'} = Z \times_k X$，结论是引理
\ref{varieties-lemma-bijection-connected-components} 的特殊情形。
\end{proof}

\begin{lemma}
\label{varieties-lemma-characterize-geometrically-connected}
设 $k$ 为域。设 $X$ 为 $k$ 上的概形。设 $\overline{k}$ 为 $k$ 的
可分代数闭包。那么 $X$ 几何连通，当且仅当基变换
$X_{\overline{k}}$ 连通。
\end{lemma}

\begin{proof}
假设 $X_{\overline{k}}$ 连通。设 $k'/k$ 为域扩张。存在域扩张
$\overline{k}'/\overline{k}$，使得 $k'$ 嵌入 $\overline{k}'$，并且该嵌入
是 $k$-扩张之间的嵌入。由引理
\ref{varieties-lemma-separably-closed-field-connected-components} 可知
$X_{\overline{k}'}$ 连通。由于 $X_{\overline{k}'} \to X_{k'}$ 满射，
可知 $X_{k'}$ 连通，如所需。
\end{proof}

\begin{lemma}
\label{varieties-lemma-descend-open}
设 $k$ 为域。设 $X$ 为 $k$ 上的概形。设 $A$ 为 $k$-代数。
设 $V \subset X_A$ 为拟紧开集。那么存在有限生成 $k$-子代数
$A' \subset A$ 以及拟紧开集 $V' \subset X_{A'}$，使得 $V = V'_A$。
\end{lemma}

\begin{proof}
顺带指出，若 $X$ 还拟分离，则该结论由《极限》引理
\ref{limits-lemma-descend-opens} 得出。令 $U_1, \ldots, U_n$ 为 $X$ 的
有限多个仿射开集，使得 $V \subset \bigcup U_{i, A}$。记
$U_i = \Spec(R_i)$。由于 $V$ 拟紧，可以找到有限多个
$f_{ij} \in R_i \otimes_k A$，$j = 1, \ldots, n_i$，使得
$V = \bigcup_i \bigcup_{j = 1, \ldots, n_i} D(f_{ij})$，其中
$D(f_{ij}) \subset U_{i, A}$ 是相应的标准开集。（我们并未断言
$V \cap U_{i, A}$ 是 $D(f_{ij})$（$j = 1, \ldots, n_i$）的并。）
显然可以找到有限生成 $k$-子代数 $A' \subset A$，使得 $f_{ij}$ 是某个
$f_{ij}' \in R_i \otimes_k A'$ 的像。置
$V' = \bigcup D(f_{ij}')$；这是 $X_{A'}$ 的拟紧开集。用
$\pi : X_A \to X_{A'}$ 表示典范态射。有 $\pi(V) \subset V'$，因为
$\pi(D(f_{ij})) \subset D(f_{ij}')$。若 $x \in X_A$ 且
$\pi(x) \in V'$，则 $\pi(x) \in D(f_{ij}')$ 对某个 $i, j$ 成立；于是
$x \in D(f_{ij})$，因为 $f_{ij}'$ 映到 $f_{ij}$。因此
$V = \pi^{-1}(V')$，如所需。
\end{proof}

\noindent
设 $k$ 为域。设 $\overline{k}/k$ 为（可能无限的）Galois 扩张。例如，
$\overline{k}$ 可以是 $k$ 的可分代数闭包。对任意
$\sigma \in \text{Gal}(\overline{k}/k)$，得到相应的自同构
$
\Spec(\sigma) :
\Spec(\overline{k})
\longrightarrow
\Spec(\overline{k})
$。
注意
$\Spec(\sigma) \circ \Spec(\tau) = \Spec(\tau \circ \sigma)$。
因此得到作用
\begingroup\small
$$
\text{Gal}(\overline{k}/k)^{opp} \times \Spec(\overline{k})
\longrightarrow
\Spec(\overline{k})
$$
\endgroup
即反群在概形 $\Spec(\overline{k})$ 上的作用。\par
设 $X$ 为 $k$ 上的概形。
按照定义，
$X_{\overline{k}} =
\Spec(\overline{k}) \times_{\Spec(k)} X$，
所以以上作用诱导典范作用
\begin{equation}
\label{varieties-equation-galois-action-base-change-kbar}
\text{Gal}(\overline{k}/k)^{opp} \times X_{\overline{k}}
\longrightarrow
X_{\overline{k}}.
\end{equation}

\begin{lemma}
\label{varieties-lemma-Galois-action-quasi-compact-open}
设 $k$ 为域。设 $X$ 为 $k$ 上的概形。设 $\overline{k}$ 为 $k$ 的
（可能无限的）Galois 扩张。设 $V \subset X_{\overline{k}}$ 为拟紧开集。
那么：
\begin{enumerate}
\item 存在有限子扩张 $\overline{k}/k'/k$ 以及拟紧开集
$V' \subset X_{k'}$，使得 $V = (V')_{\overline{k}}$；
\item 存在开子群 $H \subset \text{Gal}(\overline{k}/k)$，使得
$\sigma(V) = V$ 对所有 $\sigma \in H$ 都成立。
\end{enumerate}
\end{lemma}

\begin{proof}
由引理 \ref{varieties-lemma-descend-open}，存在有限子扩张
$k'/k \subset \overline{k}$ 以及开集 $V' \subset X_{k'}$，其拉回为 $V$。
这证明了 (1)。由于 $\text{Gal}(\overline{k}/k')$ 是
$\text{Gal}(\overline{k}/k)$ 的开子群，第 (2) 部分也很清楚。
\end{proof}

\begin{lemma}
\label{varieties-lemma-closed-fixed-by-Galois}
设 $k$ 为域。设 $\overline{k}/k$ 为（可能无限的）Galois 扩张。
设 $X$ 为 $k$ 上的概形。设 $\overline{T} \subset X_{\overline{k}}$ 满足：
\begin{enumerate}
\item $\overline{T}$ 是 $X_{\overline{k}}$ 的闭子集；
\item 对每个 $\sigma \in \text{Gal}(\overline{k}/k)$，都有
$\sigma(\overline{T}) = \overline{T}$。
\end{enumerate}
那么存在闭子集 $T \subset X$，它在 $X_{\overline{k}}$ 中的逆像是
$\overline{T}$。
\end{lemma}

\begin{proof}
本引理立即归约到 $X = \Spec(A)$ 为仿射的情形。在这种情况下，令
$\overline{I} \subset A \otimes_k \overline{k}$ 为对应于 $\overline{T}$ 的
根理想。假设 (2) 蕴含 $\sigma(\overline{I}) = \overline{I}$ 对所有
$\sigma \in \text{Gal}(\overline{k}/k)$ 都成立。取 $x \in \overline{I}$。
存在有限 Galois 扩张 $k'/k$，它包含于 $\overline{k}$ 中，并且使得
$x \in A \otimes_k k'$。置 $G = \text{Gal}(k'/k)$。置
$$
P(T) = \prod\nolimits_{\sigma \in G} (T - \sigma(x)) \in (A \otimes_k k')[T]
$$
显然，$P(T)$ 首一，并且它实际上是
$(A \otimes_k k')^G[T] = A[T]$ 的元素（由基本 Galois 理论）。此外，
若写成 $P(T) = T^d + a_1T^{d - 1} + \ldots + a_d$，则可见
$a_i \in I := A \cap \overline{I}$。结合 $P(x) = 0$ 与 $a_i \in I$，得到
$x^d = - a_1 x^{d - 1} - \ldots  - a_d \in I(A \otimes_k \overline{k})$。
因此 $x$ 属于 $I(A \otimes_k \overline{k})$ 的根。故
$\overline{I}$ 是 $I(A \otimes_k \overline{k})$ 的根，而置 $T = V(I)$
便得到一个解。
\end{proof}

\begin{lemma}
\label{varieties-lemma-characterize-geometrically-disconnected}
设 $k$ 为域。设 $X$ 为 $k$ 上的概形。下列条件等价：
\begin{enumerate}
\item $X$ 几何连通；
\item 对每个有限可分域扩张 $k'/k$，概形 $X_{k'}$ 都连通。
\end{enumerate}
\end{lemma}

\begin{proof}
由定义立即可知 (1) 蕴含 (2)。假设 $X$ 不几何连通。设
$k \subset \overline{k}$ 为 $k$ 的可分代数闭包。由引理
\ref{varieties-lemma-characterize-geometrically-connected}，可知
$X_{\overline{k}}$ 不连通。设
$X_{\overline{k}} = \overline{U} \amalg \overline{V}$，其中
$\overline{U}$ 与 $\overline{V}$ 均开、闭且非空。

\medskip\noindent
设 $W \subset X$ 是任意拟紧开集。那么
$W_{\overline{k}} \cap \overline{U}$ 与
$W_{\overline{k}} \cap \overline{V}$ 在 $W_{\overline{k}}$ 中都既开又闭。
特别地，$W_{\overline{k}} \cap \overline{U}$ 与
$W_{\overline{k}} \cap \overline{V}$ 都拟紧，并且由引理
\ref{varieties-lemma-Galois-action-quasi-compact-open}，
$W_{\overline{k}} \cap \overline{U}$ 与
$W_{\overline{k}} \cap \overline{V}$ 都可在一个有限子扩张上定义，
并且在 $\text{Gal}(\overline{k}/k)$ 的一个开子群下不变。下文将不再
另行说明而直接使用这一点。

\medskip\noindent
选取拟紧开集 $W_0 \subset X$，使
$W_{0, \overline{k}} \cap \overline{U}$ 与
$W_{0, \overline{k}} \cap \overline{V}$ 都非空。选取有限子扩张
$\overline{k}/k'/k$ 以及分解 $W_{0, k'} = U_0' \amalg V_0'$，其两项
都是既开又闭的子集，并且使得
$W_{0, \overline{k}} \cap \overline{U} = (U'_0)_{\overline{k}}$ 且
$W_{0, \overline{k}} \cap \overline{V} = (V'_0)_{\overline{k}}$。
令 $H = \text{Gal}(\overline{k}/k') \subset \text{Gal}(\overline{k}/k)$。
特别地，
$\sigma(W_{0, \overline{k}} \cap \overline{U}) =
W_{0, \overline{k}} \cap \overline{U}$，对 $\overline{V}$ 也类似。

\medskip\noindent
如上选定 $W_0$、$k'$ 后，对每个拟紧开集 $W \subset X$，置
$$
U_W =
\bigcap\nolimits_{\sigma \in H} \sigma(W_{\overline{k}} \cap \overline{U}),
\quad
V_W =
\bigcup\nolimits_{\sigma \in H} \sigma(W_{\overline{k}} \cap \overline{V}).
$$
由于 $W_{\overline{k}} \cap \overline{U}$ 与
$W_{\overline{k}} \cap \overline{V}$ 在
$\text{Gal}(\overline{k}/k)$ 的一个开子群下固定，上面的并与交都是
有限的。因此 $U_W$ 与 $V_W$ 都既开又闭。此外，按构造有
$W_{\bar k} = U_W \amalg V_W$。

\medskip\noindent
我们断言，若 $W \subset W' \subset X$ 为拟紧开集，则
$W_{\overline{k}} \cap U_{W'} = U_W$ 且
$W_{\overline{k}} \cap V_{W'} = V_W$。验证从略。因此，在定义
$U = \bigcup_{W \subset X} U_W$ 与
$V = \bigcup_{W \subset X} V_W$ 后，得到
$X_{\overline{k}} = U \amalg V$，这是既开又闭子集的不交并。
显然 $V$ 非空，因为它是（局部地）取并构造出的。另一方面，$U$ 非空，
因为按构造它包含 $W_0 \cap \overline{U}$。最后，
$U, V \subset X_{\bar k}$ 按构造闭且 $H$-不变。因此，由引理
\ref{varieties-lemma-closed-fixed-by-Galois}，有 $U = (U')_{\bar k}$ 且
$V = (V')_{\bar k}$，其中 $U', V' \subset X_{k'}$ 为闭子集。显然
$X_{k'} = U' \amalg V'$，由此可见
$X_{k'}$ 不连通，如所需。
\end{proof}

\begin{lemma}
\label{varieties-lemma-tricky}
设 $k$ 为域。设 $\overline{k}/k$ 为（可能无限的）Galois 扩张。
设 $f : T \to X$ 为 $k$ 上概形的态射。假设 $T_{\overline{k}}$ 连通且
$X_{\overline{k}}$ 不连通。那么 $X$ 不连通。
\end{lemma}

\begin{proof}
写成 $X_{\overline{k}} = \overline{U} \amalg \overline{V}$，其中
$\overline{U}$ 与 $\overline{V}$ 既开又闭。记
$\overline{f} : T_{\overline{k}} \to X_{\overline{k}}$ 为 $f$ 的基变换。
由于 $T_{\overline{k}}$ 连通，可知 $T_{\overline{k}}$ 包含于
$\overline{f}^{-1}(\overline{U})$ 或 $\overline{f}^{-1}(\overline{V})$
二者之一。不妨设
$T_{\overline{k}} \subset \overline{f}^{-1}(\overline{U})$。

\medskip\noindent
固定拟紧开集 $W \subset X$。存在有限 Galois 子扩张
$\overline{k}/k'/k$，使得 $\overline{U} \cap W_{\overline{k}}$ 与
$\overline{V} \cap W_{\overline{k}}$ 来自拟紧开集
$U', V' \subset W_{k'}$。于是也有 $W_{k'} = U' \amalg V'$。考虑
$$
U'' = \bigcap\nolimits_{\sigma \in \text{Gal}(k'/k)} \sigma(U'),
\quad
V'' = \bigcup\nolimits_{\sigma \in \text{Gal}(k'/k)} \sigma(V').
$$
它们 Galois 不变、既开又闭，并且 $W_{k'} = U'' \amalg V''$。
由引理 \ref{varieties-lemma-closed-fixed-by-Galois}，得到既开又闭的子集
$U_W, V_W \subset W$，使得
$U'' = (U_W)_{k'}$、$V'' = (V_W)_{k'}$ 且
$W = U_W \amalg V_W$。

\medskip\noindent
我们断言，若 $W \subset W' \subset X$ 为拟紧开集，则
$W \cap U_{W'} = U_W$ 且 $W \cap V_{W'} = V_W$。验证从略。因此，在定义
$U = \bigcup_{W \subset X} U_W$ 与
$V = \bigcup_{W \subset X} V_W$ 后，得到 $X = U \amalg V$。
显然 $V$ 非空，因为它是（局部地）取并构造出的。另一方面，$U$ 非空，
因为按构造它包含 $f(T)$。
\end{proof}

\begin{lemma}
\label{varieties-lemma-geometrically-connected-criterion}
\begin{reference}
\cite[IV Corollary 4.5.13.1(i)]{EGA}
\end{reference}
设 $k$ 为域。设 $T \to X$ 为 $k$ 上概形的态射。假设 $T$ 几何连通且
$X$ 连通。那么 $X$ 几何连通。
\end{lemma}

\begin{proof}
这是引理 \ref{varieties-lemma-tricky} 的改述。
\end{proof}

\begin{lemma}
\label{varieties-lemma-geometrically-connected-if-connected-and-point}
设 $k$ 为域。设 $X$ 为 $k$ 上的概形。假设 $X$ 连通，并且有一点 $x$，
使得 $k$ 在 $\kappa(x)$ 中代数闭。那么 $X$ 几何连通。特别地，如果
$X$ 有 $k$-有理点且 $X$ 连通，那么 $X$ 几何连通。
\end{lemma}

\begin{proof}
置 $T = \Spec(\kappa(x))$。设 $\overline{k}$ 为 $k$ 的可分代数闭包。
关于 $\kappa(x)/k$ 的假设蕴含 $T_{\overline{k}}$ 不可约；见《代数》引理
\ref{algebra-lemma-field-extension-geometrically-irreducible}。因此，由引理
\ref{varieties-lemma-geometrically-connected-criterion} 可知 $X_{\overline{k}}$ 连通。
由引理 \ref{varieties-lemma-characterize-geometrically-connected}，可知 $X$ 几何连通。
\end{proof}

\begin{lemma}
\label{varieties-lemma-inverse-image-connected-component}
设 $K/k$ 为域扩张。设 $X$ 为 $k$ 上的概形。对每个连通分支 $T$
（属于 $X$），逆像 $T_K \subset X_K$ 是 $X_K$ 的若干连通分支之并。
\end{lemma}

\begin{proof}
这是纯拓扑的陈述。记 $p : X_K \to X$ 为投影态射。设
$T \subset X$ 为 $X$ 的一个连通分支。设
$t \in T_K = p^{-1}(T)$。设 $C \subset X_K$ 为包含 $t$ 的连通分支。
那么 $p(C)$ 是 $X$ 的一个与 $T$ 相交的连通子集，故
$p(C) \subset T$。因此 $C \subset T_K$。
\end{proof}

\noindent
下面的引理将被后面更强的引理 \ref{varieties-lemma-image-connected-component}
所取代。

\begin{lemma}
\label{varieties-lemma-image-connected-component-finite-extension}
设 $K/k$ 为有限域扩张，并设 $X$ 为 $k$ 上的概形。用
$p : X_K \to X$ 表示投影态射。对每个连通分支 $T$（属于 $X_K$），像
$p(T)$ 是 $X$ 的一个连通分支。
\end{lemma}

\begin{proof}
像 $p(T)$ 包含于某个连通分支 $X'$ 中，而该分支属于 $X$。以任意方式
把 $X'$ 视为 $X$ 的闭子概形。那么 $T$ 也是
$X'_K = p^{-1}(X')$ 的连通分支，因此可以假设 $X$ 连通。态射 $p$ 是
开映射（《态射》引理
\ref{morphisms-lemma-scheme-over-field-universally-open}）、闭映射
（《态射》引理 \ref{morphisms-lemma-integral-universally-closed}），并且
$p$ 的纤维是有限集（《态射》引理
\ref{morphisms-lemma-finite-quasi-finite}）。所以可以应用《拓扑学》引理
\ref{topology-lemma-finite-fibre-connected-components} 得出结论。
\end{proof}

\begin{lemma}[Gabber]
\label{varieties-lemma-image-connected-component}
\begin{reference}
Ofer Gabber 于 2016 年 6 月 4 日发送的电子邮件
\end{reference}
设 $K/k$ 为域扩张。设 $X$ 为 $k$ 上的概形。记
$p : X_K \to X$ 为投影态射。设 $\overline{T} \subset X_K$ 为连通分支。
那么 $p(\overline{T})$ 是 $X$ 的连通分支。
\end{lemma}

\begin{proof}
当 $K/k$ 有限时，这就是引理
\ref{varieties-lemma-image-connected-component-finite-extension}。一般情况下，
证明较为困难。

\medskip\noindent
令 $T \subset X$ 为 $X$ 的包含 $\overline{T}$ 之像的连通分支。可以将
$X$ 替换为 $T$（赋予诱导的约化子概形结构）。因此可以假设 $X$ 连通。
令 $A = H^0(X, \mathcal{O}_X)$。设 $L \subset A$ 为极大的弱
\'etale $k$-子代数；见《更多代数》引理
\ref{more-algebra-lemma-max-weakly-etale-subalgebra}。由于 $A$ 没有任何
非平凡幂等元，由《更多代数》引理
\ref{more-algebra-lemma-class-weakly-etale-over-field} 可知，$L$ 是域，
并且是 $k$ 的可分代数扩张。注意，$L$ 也是极大弱 \'etale
$L$-子代数；这里所指的母代数为 $A$。这是因为《更多代数》引理
\ref{more-algebra-lemma-composition-weakly-etale} 表明，任意弱 \'etale
$L$-代数在 $k$ 上都是弱 \'etale 的。由《概形》引理
\ref{schemes-lemma-morphism-into-affine}，得到结构态射的分解
$X \to \Spec(L) \to \Spec(k)$。

\medskip\noindent
设 $L'/L$ 为有限可分扩张。由《概形的上同调》引理
\ref{coherent-lemma-finite-locally-free-base-change-cohomology}，有
$$
A \otimes_L L' =
H^0(X \times_{\Spec(L)} \Spec(L'), \mathcal{O}_{X \times_{\Spec(L)} \Spec(L')})
$$
极大弱 \'etale $L'$-子代数（在 $A \otimes_L L'$ 中）是
$L \otimes_L L' = L'$；见《更多代数》引理
\ref{more-algebra-lemma-change-fields-max-weakly-etale-subalgebra}。特别地，
$A \otimes_L L'$ 没有非平凡幂等元（这样的幂等元会生成一个弱
\'etale 子代数），所以 $X \times_{\Spec(L)} \Spec(L')$ 连通。由引理
\ref{varieties-lemma-characterize-geometrically-disconnected}，可知 $X$ 在 $L$ 上
几何连通。

\medskip\noindent
赋予 $\overline{T}$ 诱导的约化概形结构，并考虑复合
$$
\overline{T} \xrightarrow{i} X_K = X \times_{\Spec(k)} \Spec(K)
\xrightarrow{\pi}
\Spec(L \otimes_k K)
$$
其像包含于 $\Spec(L \otimes_k K)$ 的一个连通分支中。由于
$K \to L \otimes_k K$ 整，可知 $\Spec(L \otimes_k K)$ 的各连通分支
都是点，并且所有点都闭；见《代数》引理
\ref{algebra-lemma-integral-over-field}。因此得到商域
$L \otimes_k K \to E$，使 $\overline{T}$ 映入
$\Spec(E) \subset \Spec(L \otimes_k K)$。故
$i(\overline{T}) \subset \pi^{-1}(\Spec(E))$。但是
$$
\pi^{-1}(\Spec(E)) =
(X \times_{\Spec(k)} \Spec(K)) \times_{\Spec(L \otimes_k K)} \Spec(E) =
X \times_{\Spec(L)} \Spec(E)
$$
它连通，因为 $X$ 在 $L$ 上几何连通。于是得到等式
$\overline{T} = X \times_{\Spec(L)} \Spec(E)$（在集合论意义下），
并由此得出 $\overline{T} \to X$ 满射，如所需。
\end{proof}

\noindent
设 $X$ 为概形。用 $\pi_0(X)$ 表示 $X$ 的连通分支集合。

\begin{lemma}
\label{varieties-lemma-galois-action-connected-components}
设 $k$ 为域，并设 $\overline{k}$ 为其可分代数闭包。设 $X$ 为 $k$ 上的
概形。存在作用
$$
\text{Gal}(\overline{k}/k)^{opp} \times \pi_0(X_{\overline{k}})
\longrightarrow
\pi_0(X_{\overline{k}})
$$
并具有下列性质：
\begin{enumerate}
\item 元素 $\overline{T} \in \pi_0(X_{\overline{k}})$ 被该作用固定，
当且仅当存在连通分支 $T \subset X$，它在 $k$ 上几何连通，并且满足
$T_{\overline{k}} = \overline{T}$。
\item 对任意域扩张 $k'/k$ 及其可分代数闭包 $\overline{k}'$，图表
$$
\xymatrix{
\text{Gal}(\overline{k}'/k') \times \pi_0(X_{\overline{k}'})
\ar[r] \ar[d] &
\pi_0(X_{\overline{k}'}) \ar[d] \\
\text{Gal}(\overline{k}/k) \times \pi_0(X_{\overline{k}})
\ar[r] &
\pi_0(X_{\overline{k}})
}
$$
交换（由引理 \ref{varieties-lemma-separably-closed-field-connected-components}，
右侧竖直箭头为双射）。
\end{enumerate}
\end{lemma}

\begin{proof}
作用（\ref{varieties-equation-galois-action-base-change-kbar}）即
$\text{Gal}(\overline{k}/k)$ 在 $X_{\overline{k}}$ 上的作用，它诱导
在连通分支上的作用。连通分支总是闭的（《拓扑学》引理
\ref{topology-lemma-connected-components}）。因此，若 $\overline{T}$ 如
(1) 所述，则由引理 \ref{varieties-lemma-closed-fixed-by-Galois}，存在闭子集
$T \subset X$，使得 $\overline{T} = T_{\overline{k}}$。注意，$T$ 在
$k$ 上几何连通；见引理 \ref{varieties-lemma-characterize-geometrically-connected}。
为说明 $T$ 是 $X$ 的一个连通分支，假设
$T \subset T'$、$T \not = T'$，其中 $T'$ 是 $X$ 的一个连通分支。
在这种情况下，$T'_{k'}$ 严格包含 $\overline{T}$，因而不连通。由引理
\ref{varieties-lemma-tricky}，这意味着 $T'$ 不连通！矛盾。

\medskip\noindent
我们略去 (2) 中函子性的证明。
\end{proof}

\begin{lemma}
\label{varieties-lemma-galois-action-connected-components-continuous}
设 $k$ 为域，并设 $\overline{k}$ 为其可分代数闭包。设 $X$ 为 $k$ 上的
概形。假设：
\begin{enumerate}
\item $X$ 拟紧；
\item $X_{\overline{k}}$ 的连通分支都是开集。
\end{enumerate}
那么：
\begin{enumerate}
\item[(a)] $\pi_0(X_{\overline{k}})$ 有限；
\item[(b)] $\text{Gal}(\overline{k}/k)$ 在
$\pi_0(X_{\overline{k}})$ 上的作用连续。
\end{enumerate}
此外，当 $X$ 在 $k$ 上有限型时，假设 (1) 与 (2) 均成立。
\end{lemma}

\begin{proof}
由于各连通分支都是开集、覆盖 $X_{\overline{k}}$（《拓扑学》引理
\ref{topology-lemma-connected-components}），而且 $X_{\overline{k}}$ 拟紧，
可知它们只有有限多个。因此 (a) 成立。由引理 \ref{varieties-lemma-descend-open}，
这些连通分支中的每一个都可以在 $\overline{k}/k$ 的某个有限子扩张上
定义，故得到 (b)。若 $X$ 在 $k$ 上有限型，则 $X_{\overline{k}}$ 在
$\overline{k}$ 上有限型（《态射》引理
\ref{morphisms-lemma-base-change-finite-type}）。因此
$X_{\overline{k}}$ 是 Noether 概形（《态射》引理
\ref{morphisms-lemma-finite-type-noetherian}）。所以
$X_{\overline{k}}$ 只有有限多个不可约分支（《性质》引理
\ref{properties-lemma-Noetherian-irreducible-components}），更不必说，它只有
有限多个连通分支（因而这些分支都是开的）。
\end{proof}

\section{几何不可约概形}
\label{varieties-section-geometrically-irreducible}

\noindent
若 $X$ 是域上的不可约概形，则在扩张基域之后，$X$ 可能变为可约概形。
几何不可约概形不会发生这种情况。

\begin{definition}
\label{varieties-definition-geometrically-irreducible}
设 $X$ 为域 $k$ 上的概形。称 $X$ 在 $k$ 上{\it 几何不可约}，如果概形
$X_{k'}$ 不可约\footnote{不可约空间非空。}，其中 $k'$ 是 $k$ 的任意域扩张。
\end{definition}

\begin{lemma}
\label{varieties-lemma-geometrically-irreducible-check-after-extension}
设 $X$ 为域 $k$ 上的概形，并设 $k'/k$ 为域扩张。那么，$X$ 在
$k$ 上几何不可约，当且仅当 $X_{k'}$ 在 $k'$ 上几何不可约。
\end{lemma}

\begin{proof}
若 $X$ 在 $k$ 上几何不可约，则 $X_{k'}$ 在 $k'$ 上几何不可约是显然的。
反过来，对任意域扩张 $k''/k$，存在公共域扩张 $k'''/k'$ 和
$k'''/k''$；见《域》引理 \ref{fields-lemma-common-extension-field}。
态射 $X_{k'''} \to X_{k''}$ 是满的（因为它是域谱之间的满态射的基变换），
故 $X_{k'''}$ 不可约蕴含 $X_{k''}$ 不可约。因此，若 $X_{k'}$ 在
$k'$ 上几何不可约，则 $X$ 在 $k$ 上几何不可约。
\end{proof}

\begin{lemma}
\label{varieties-lemma-separably-closed-irreducible}
设 $X$ 为一个概形，它定义在可分代数闭的域 $k$ 上。若 $X$ 不可约，则 $X_K$ 对任意
域扩张 $K/k$ 都不可约。换言之，$X$ 在 $k$ 上几何不可约。
\end{lemma}

\begin{proof}
使用《性质》引理 \ref{properties-lemma-characterize-irreducible}
和《代数》引理 \ref{algebra-lemma-separably-closed-irreducible}。
\end{proof}

\begin{lemma}
\label{varieties-lemma-bijection-irreducible-components}
设 $k$ 为域，设 $X$、$Y$ 为 $k$ 上的概形。假设 $X$ 在 $k$ 上几何不可约。
那么投影态射
$$
p : X \times_k Y \longrightarrow Y
$$
诱导不可约分支之间的双射。
\end{lemma}

\begin{proof}
首先，$p$ 的概形论纤维不可约，因为它们是几何不可约概形 $X$ 沿域扩张
的基变换。此外，概形论纤维与集合论纤维同胚；见《概形》引理
\ref{schemes-lemma-fibre-topological}。由《态射》引理
\ref{morphisms-lemma-scheme-over-field-universally-open}，映射 $p$ 是开的。
因此可以应用《拓扑学》引理
\ref{topology-lemma-irreducible-fibres-irreducible-components} 得出结论。
\end{proof}

\begin{lemma}
\label{varieties-lemma-geometrically-irreducible-local}
\begin{slogan}
除连通性外，几何不可约性是 Zariski 局部的。
\end{slogan}
设 $k$ 为域，设 $X$ 为 $k$ 上的概形。下列条件等价：
\begin{enumerate}
\item $X$ 在 $k$ 上几何不可约；
\item 对每个非空仿射开集 $U$，$k$-代数 $\mathcal{O}_X(U)$ 在
$k$ 上几何不可约（见《代数》定义
\ref{algebra-definition-geometrically-irreducible}）；
\item $X$ 不可约，并且存在仿射开覆盖 $X = \bigcup U_i$，使得每个
$k$-代数 $\mathcal{O}_X(U_i)$ 都几何不可约；
\item 存在开覆盖 $X = \bigcup_{i \in I} X_i$，其中
$I \not = \emptyset$，使得 $X_i$ 对每个 $i$ 都几何不可约，并且
$X_i \cap X_j \not = \emptyset$ 对所有 $i, j \in I$ 都成立。
\end{enumerate}
此外，若 $X$ 几何不可约，则 $X$ 的每个非空开子概形也几何不可约。
\end{lemma}

\begin{proof}
仿射概形 $\Spec(A)$ 在 $k$ 上几何不可约，当且仅当 $A$ 在 $k$ 上
几何不可约；这直接来自定义。回忆，若 $X$ 不可约，则它的每个非空开子概形
也不可约，而且任意两个非空开子集的交非空。又若每个仿射开集都不可约，
则该概形不可约；见《性质》引理
\ref{properties-lemma-characterize-irreducible}。因此，引理的最后一个断言，
以及蕴含 (1) $\Rightarrow$ (2)、(2) $\Rightarrow$ (3) 和
(3) $\Rightarrow$ (4) 都是显然的。若 (4) 成立，则对任意域扩张
$k'/k$，概形 $X_{k'}$ 有一个由两两相交的不可约开集组成的覆盖。因此
$X_{k'}$ 不可约，故 (4) 蕴含 (1)。
\end{proof}

\begin{lemma}
\label{varieties-lemma-geometrically-irreducible-function-field}
设 $X$ 为域 $k$ 上的不可约概形，设 $\xi \in X$ 为其泛点。下列条件等价：
\begin{enumerate}
\item $X$ 在 $k$ 上几何不可约；
\item $\kappa(\xi)$ 在 $k$ 上几何不可约；
\item $k$ 在 $\kappa(\xi)$ 中的可分代数闭包等于 $k$。
\end{enumerate}
\end{lemma}

\begin{proof}
由《代数》引理
\ref{algebra-lemma-geometrically-irreducible-separable-elements}，
(2) 与 (3) 等价。

\medskip\noindent
假设 (1)。回忆，$\mathcal{O}_{X, \xi}$ 是 $\mathcal{O}_X(U)$ 的滤过
余极限，其中 $U$ 遍历 $X$ 的非空仿射开子概形。结合引理
\ref{varieties-lemma-geometrically-irreducible-local} 与《代数》引理
\ref{algebra-lemma-subalgebra-geometrically-irreducible}，可知
$\mathcal{O}_{X, \xi}$ 在 $k$ 上几何不可约。由于
$\mathcal{O}_{X, \xi} \to \kappa(\xi)$ 是具有局部幂零核的满射
（见《代数》引理 \ref{algebra-lemma-minimal-prime-reduced-ring}），由《代数》
引理 \ref{algebra-lemma-p-ring-map} 可知 $\kappa(\xi)$ 几何不可约。

\medskip\noindent
假设 (2)。可以假设 $X$ 约化。设 $U \subset X$ 为非空仿射开集。则
$U = \Spec(A)$，其中 $A$ 是以 $\kappa(\xi)$ 为分式域的整环。因此，
$A$ 是一个几何不可约 $k$-代数的 $k$-子代数。于是由《代数》引理
\ref{algebra-lemma-subalgebra-geometrically-irreducible}，$A$ 在 $k$ 上
几何不可约。由引理 \ref{varieties-lemma-geometrically-irreducible-local} 得出
$X$ 在 $k$ 上几何不可约。
\end{proof}

\begin{lemma}
\label{varieties-lemma-geometrically-irreducible-self-product}
概形 $X$ 定义在域 $k$ 上时，它在 $k$ 上几何不可约，当且仅当
$X\times_k X$ 不可约。
\end{lemma}

\begin{proof}
若 $X$ 在 $k$ 上几何不可约，则由引理
\ref{varieties-lemma-bijection-irreducible-components}，$X\times_k X$ 不可约。

\medskip\noindent
反过来，假设 $X \times_k X$ 不可约。把 $X$ 替换为其底层约化概形，可以
在不改变问题的前提下假设 $X$ 约化。由于第一投影
$X \times_k X \to X$ 满射，$X$ 也不可约，因而是整概形。假设存在元素
$\alpha$，它属于函数域 $k(X)$，是 $k$ 上的可分代数元且不属于 $k$。
令 $E = k(\alpha)$。

\medskip\noindent
那么，存在非空仿射开子概形 $U$，它是 $X$ 的开子概形，并且使得
$k \subset E \subset O(U)$，因而
$E \otimes_k E \subset
O(U)\otimes_k O(U) = O(U\times_k U)$。
于是有支配态射 $U \times_k U \to \Spec(E\otimes_k E)$。由于
$X\times_k X$ 不可约，$U \times_k U$ 也不可约，因而
$\Spec(E \otimes_k E)$ 连通。若 $E$ 严格包含 $k$，则
$E \otimes_k E$ 是若干域的直积，其中一个因子为 $E$。它的维数为
$(\dim_k(E))^2>\dim_k(E)$；这里维数是作为 $k$-向量空间计算的。故 $E\otimes_k E$
是至少两个域的直积，这与 $\Spec(E \otimes_k E)$ 连通矛盾。因此事实上
$E$ 等于 $k$。由引理 \ref{varieties-lemma-geometrically-irreducible-function-field}，
$X$ 在 $k$ 上几何不可约。
\end{proof}

\begin{lemma}
\label{varieties-lemma-separably-closed-field-irreducible-components}
设 $k'/k$ 为域扩张，设 $X$ 为 $k$ 上的概形。置 $X' = X_{k'}$。
假设 $k$ 可分代数闭。那么态射 $X' \to X$ 诱导不可约分支之间的双射。
\end{lemma}

\begin{proof}
由于 $k$ 可分代数闭，可知 $k'$ 在 $k$ 上几何不可约；见《代数》引理
\ref{algebra-lemma-separably-closed-irreducible-implies-geometric}。
因此 $Z = \Spec(k')$ 在 $k$ 上几何不可约；见上面的引理
\ref{varieties-lemma-geometrically-irreducible-local}。由于
$X' = Z \times_k X$，结论是引理
\ref{varieties-lemma-bijection-irreducible-components} 的特例。
\end{proof}

\begin{lemma}
\label{varieties-lemma-characterize-geometrically-irreducible}
\begin{slogan}
几何不可约性可以在基域的一个可分代数闭包上检验。
\end{slogan}
设 $k$ 为域，设 $X$ 为 $k$ 上的概形。下列条件等价：
\begin{enumerate}
\item $X$ 在 $k$ 上几何不可约；
\item 对每个有限可分域扩张 $k'/k$，概形 $X_{k'}$ 都不可约；
\item $X_{\overline{k}}$ 不可约，其中 $k \subset \overline{k}$ 是
$k$ 的一个可分代数闭包。
\end{enumerate}
\end{lemma}

\begin{proof}
假设 $X_{\overline{k}}$ 不可约，即假设 (3)。设 $k'/k$ 为域扩张。
存在域扩张 $\overline{k}'/\overline{k}$，使得 $k'$ 嵌入
$\overline{k}'$，并且该嵌入是 $k$-扩张的嵌入。由引理
\ref{varieties-lemma-separably-closed-field-irreducible-components} 可知
$X_{\overline{k}'}$ 不可约。由于 $X_{\overline{k}'} \to X_{k'}$ 满射，
可知 $X_{k'}$ 不可约。因此 (1) 成立。

\medskip\noindent
设 $k \subset \overline{k}$ 为 $k$ 的一个可分代数闭包。假设 (3) 不成立，
即假设 $X_{\overline{k}}$ 可约。我们的目标是证明，$X_{k'}$ 对某个有限
子扩张 $\overline{k}/k'/k$ 也可约。设
$X = \bigcup_{i \in I} U_i$ 为仿射开覆盖，其中 $U_i$ 非空。若对某个
$i$，概形 $U_i$ 可约，或者对某对 $i \not = j$，交
$U_i \cap U_j$ 为空，则 $X$ 可约（《性质》引理
\ref{properties-lemma-characterize-irreducible}），于是结论成立。特别地，
可以假设 $U_{i, \overline{k}} \cap U_{j, \overline{k}}$ 对所有
$i, j \in I$ 都非空；由此可知，$U_{i, \overline{k}}$ 对某个 $i$
必须可约。根据《代数》引理
\ref{algebra-lemma-geometrically-irreducible}，这意味着 $U_{i, k'}$ 对某个
有限可分域扩张 $k'/k$ 可约。
因此 $X_{k'}$ 也可约。于是 (2) 蕴含 (3)。

\medskip\noindent
蕴含 (1) $\Rightarrow$ (2) 是直接的。这就证明了引理。
\end{proof}

\begin{lemma}
\label{varieties-lemma-inverse-image-irreducible}
设 $K/k$ 为域扩张，设 $X$ 为 $k$ 上的概形。对每个不可约分支 $T$
（它是 $X$ 的不可约分支），逆像 $T_K \subset X_K$ 是 $X_K$
的若干不可约分支之并。
\end{lemma}

\begin{proof}
设 $T \subset X$ 为 $X$ 的一个不可约分支。态射 $T_K \to T$ 平坦，
故泛化可沿 $T_K \to T$ 提升。因此，每个 $\xi \in T_K$，只要它是
$T_K$ 的一个不可约分支的泛点，就映到泛点 $\eta$；后者是 $T$ 的泛点。若
$\xi' \leadsto \xi$ 是 $X_K$ 中的一次特化，则 $\xi'$ 映到 $\eta$，
因为没有特化到 $\eta$ 的点存在于 $X$ 中。故 $\xi' \in T_K$，由此得出
$\xi = \xi'$。换言之，$\xi$ 是 $X_K$ 的一个不可约分支的泛点。
这意味着 $T_K$ 的所有不可约分支也都是 $X_K$ 的不可约分支。
\end{proof}

\noindent
对概形 $X$，用 $\text{IrredComp}(X)$ 表示 $X$ 的不可约分支集合。

\begin{lemma}
\label{varieties-lemma-image-irreducible}
设 $K/k$ 为域扩张，设 $X$ 为 $k$ 上的概形。对每个不可约分支
$\overline{T} \subset X_K$，$\overline{T}$ 在 $X$ 中的像是 $X$ 的一个
不可约分支。这定义了典范映射
$$
\text{IrredComp}(X_K)
\longrightarrow
\text{IrredComp}(X)
$$
并且该映射满射。
\end{lemma}

\begin{proof}
考虑图表
$$
\xymatrix{
X_K \ar[d] & X_{\overline{K}} \ar[d] \ar[l] \\
X & X_{\overline{k}} \ar[l]
}
$$
其中 $\overline{K}$ 是 $K$ 的可分代数闭包，而 $\overline{k}$ 是
$k$ 的可分代数闭包。由引理
\ref{varieties-lemma-separably-closed-field-irreducible-components}，态射
$X_{\overline{K}} \to X_{\overline{k}}$ 诱导不可约分支之间的双射。因此，
只需对态射 $X_{\overline{k}} \to X$ 和
$X_{\overline{K}} \to X_K$ 证明本引理。换言之，可以假设
$K = \overline{k}$。

\medskip\noindent
态射 $p : X_{\overline{k}} \to X$ 是整、平坦且满的。平坦性蕴含泛化可沿
$p$ 提升；见《态射》引理
\ref{morphisms-lemma-generalizations-lift-flat}。因此，
$X_{\overline{k}}$ 的不可约分支的泛点映到 $X$ 的不可约分支的泛点。
整性蕴含 $p$ 万有闭；见《态射》引理
\ref{morphisms-lemma-integral-universally-closed}。所以，不可约分支的像
$p(\overline{T})$ 是包含 $X$ 的某个不可约分支的泛点的闭不可约子集；
因而 $p(\overline{T})$ 是 $X$ 的一个不可约分支。这证明了第一个断言。
若 $T \subset X$ 是不可约分支，则 $p^{-1}(T) =T_K$ 是不可约分支的
非空并；见引理 \ref{varieties-lemma-inverse-image-irreducible}。根据第一部分，
这些分支中的每一个必然满射到 $T$。故该映射是满的。
\end{proof}

\begin{lemma}
\label{varieties-lemma-irreducible-dense-rational-points-geometrically-irreducible}
设 $k$ 为域，设 $X$ 为 $k$ 上的概形。若 $X$ 不可约且具有稠密的
$k$-有理点集，则 $X$ 几何不可约。
\end{lemma}

\begin{proof}
设 $k'/k$ 为有限域扩张，并设 $Z, Z' \subset X_{k'}$ 为不可约分支。
只需证明 $Z = Z'$；见引理
\ref{varieties-lemma-characterize-geometrically-irreducible}。由引理
\ref{varieties-lemma-image-irreducible}，有 $p(Z) = p(Z') = X$，其中
$p : X_{k'} \to X$ 是投影。若 $Z \not = Z'$，则
$Z \cap Z'$ 在 $X_{k'}$ 中无处稠密，因而由《态射》引理
\ref{morphisms-lemma-image-nowhere-dense-finite}，$p(Z \cap Z')$
不稠密；这里还使用了 $p$ 是有限态射，因为它是有限态射
$\Spec(k') \to \Spec(k)$ 的基变换；见《态射》引理
\ref{morphisms-lemma-base-change-finite}。因此，可以选取一个
$k$-有理点 $x \in X$，使得 $x \not \in p(Z \cap Z')$。由于 $x$ 的
剩余域是 $k$，可知 $p^{-1}(\{x\}) = \{x'\}$，其中
$x' \in X_{k'}$ 是剩余域为 $k'$ 的点。由于
$x \in p(Z) = p(Z')$，可知 $x' \in Z \cap Z'$；这就是所需的矛盾。
\end{proof}

\begin{lemma}
\label{varieties-lemma-galois-action-irreducible-components}
设 $k$ 为域，并设 $\overline{k}$ 为其可分代数闭包。设 $X$ 为 $k$ 上的
概形。存在作用
$$
\text{Gal}(\overline{k}/k)^{opp} \times \text{IrredComp}(X_{\overline{k}})
\longrightarrow
\text{IrredComp}(X_{\overline{k}})
$$
并具有下列性质：
\begin{enumerate}
\item 元素 $\overline{T} \in \text{IrredComp}(X_{\overline{k}})$ 被该作用
固定，当且仅当存在不可约分支 $T \subset X$，它在 $k$ 上几何不可约，
并且满足 $T_{\overline{k}} = \overline{T}$。
\item 对任意域扩张 $k'/k$ 及其可分代数闭包 $\overline{k}'$，图表
$$
\xymatrix{
\text{Gal}(\overline{k}'/k') \times \text{IrredComp}(X_{\overline{k}'})
\ar[r] \ar[d] &
\text{IrredComp}(X_{\overline{k}'}) \ar[d] \\
\text{Gal}(\overline{k}/k) \times \text{IrredComp}(X_{\overline{k}})
\ar[r] &
\text{IrredComp}(X_{\overline{k}})
}
$$
交换（由引理 \ref{varieties-lemma-separably-closed-field-irreducible-components}，
右侧竖直箭头为双射）。
\end{enumerate}
\end{lemma}

\begin{proof}
$\text{Gal}(\overline{k}/k)$ 在 $X_{\overline{k}}$ 上的作用
（\ref{varieties-equation-galois-action-base-change-kbar}）诱导其在不可约分支上的
作用。不可约分支总是闭的（《拓扑学》引理
\ref{topology-lemma-connected-components}）。因此，若 $\overline{T}$
如 (1) 所述，则由引理 \ref{varieties-lemma-closed-fixed-by-Galois}，存在闭子集
$T \subset X$，使得 $\overline{T} = T_{\overline{k}}$。注意，$T$ 在
$k$ 上几何不可约；见引理
\ref{varieties-lemma-characterize-geometrically-irreducible}。为说明 $T$ 是
$X$ 的一个不可约分支，假设 $T \subset T'$、$T \not = T'$，其中
$T'$ 是 $X$ 的一个不可约分支。设 $\overline{\eta}$ 为
$\overline{T}$ 的泛点。它映到泛点 $\eta$，后者是 $T$ 的泛点。于是泛点
$\xi \in T'$ 特化到 $\eta$。由于 $X_{\overline{k}} \to X$ 平坦，
存在点 $\overline{\xi} \in X_{\overline{k}}$，它映到 $\xi$ 并特化到
$\overline{\eta}$。于是，单点集 $\{\overline{\xi}\}$ 在
$X_{\overline{\xi}}$ 中的闭包是一个严格包含 $\overline{T}$ 的
不可约闭子集。这就是所需的矛盾。

\medskip\noindent
我们略去 (2) 中函子性的证明。
\end{proof}

\begin{lemma}
\label{varieties-lemma-orbit-irreducible-components}
设 $k$ 为域，并设 $\overline{k}$ 为其可分代数闭包。设 $X$ 为 $k$ 上的
概形。引理 \ref{varieties-lemma-image-irreducible} 的映射
$$
\text{IrredComp}(X_{\overline{k}})
\longrightarrow
\text{IrredComp}(X)
$$
的各纤维，恰好是 $\text{Gal}(\overline{k}/k)$ 在引理
\ref{varieties-lemma-galois-action-irreducible-components} 的作用下的轨道。
\end{lemma}

\begin{proof}
设 $T \subset X$ 为 $X$ 的一个不可约分支，设 $\eta \in T$ 为其泛点。
由引理 \ref{varieties-lemma-inverse-image-irreducible} 和
\ref{varieties-lemma-image-irreducible}，$\overline{T}$ 的各不可约分支中映入
$T$ 的那些分支的泛点都映到 $\eta$。由《代数》引理
\ref{algebra-lemma-Galois-orbit}，Galois 群在 $X_{\overline{k}}$ 中所有
映到 $\eta$ 的点上可迁地作用。因此引理成立。
\end{proof}

\begin{lemma}
\label{varieties-lemma-galois-action-irreducible-components-locally-finite-type}
设 $k$ 为域。假设 $X \to \Spec(k)$ 局部有限型。在这种情况下：
\begin{enumerate}
\item 若赋予 $\text{IrredComp}(X_{\overline{k}})$ 离散拓扑，则作用
$$
\text{Gal}(\overline{k}/k)^{opp} \times \text{IrredComp}(X_{\overline{k}})
\longrightarrow
\text{IrredComp}(X_{\overline{k}})
$$
是连续的；
\item $X_{\overline{k}}$ 的每个不可约分支都可以在 $k$ 的某个有限扩张上
定义；
\item 给定任意不可约分支 $T \subset X$，概形
$T_{\overline{k}}$ 是 $X_{\overline{k}}$ 的有限多个不可约分支之并，
而且这些分支都属于同一个 $\text{Gal}(\overline{k}/k)$-轨道。
\end{enumerate}
\end{lemma}

\begin{proof}
设 $\overline{T}$ 为 $X_{\overline{k}}$ 的一个不可约分支。可以选取仿射
开集 $U \subset X$，使得 $\overline{T} \cap U_{\overline{k}}$ 非空。
写成 $U = \Spec(A)$；于是 $A$ 是有限型 $k$-代数，见《态射》引理
\ref{morphisms-lemma-locally-finite-type-characterize}。因此
$A_{\overline{k}}$ 是有限型 $\overline{k}$-代数，特别地，它是 Noether
环。设 $\mathfrak p = (f_1, \ldots, f_n)$ 为对应于
$\overline{T} \cap U_{\overline{k}}$ 的素理想。由于
$A_{\overline{k}} = A \otimes_k \overline{k}$，可知存在有限子扩张
$\overline{k}/k'/k$，使得每个 $f_i \in A_{k'}$。显然，
$\text{Gal}(\overline{k}/k')$ 固定 $\overline{T}$，这证明了 (1)。

\medskip\noindent
把引理 \ref{varieties-lemma-galois-action-irreducible-components} (1) 应用于
$k'$ 上的情形，可知不可约分支 $\overline{T}$ 形如
$T'_{\overline{k}}$，其中 $T' \subset X_{k'}$ 不可约，故 (2) 成立。

\medskip\noindent
为证明 (3)，设 $T \subset X$ 为一个不可约分支。选取不可约分支
$\overline{T} \subset X_{\overline{k}}$，它映到 $T$；见引理
\ref{varieties-lemma-image-irreducible}。由上可知，$\overline{T}$ 的轨道有限，
记为 $\overline{T}_1, \ldots, \overline{T}_n$。那么
$\overline{T}_1 \cup \ldots \cup \overline{T}_n$ 是一个
$\text{Gal}(\overline{k}/k)$-不变闭子集，它位于 $X_{\overline{k}}$ 中，故由引理
\ref{varieties-lemma-closed-fixed-by-Galois}，它形如 $W_{\overline{k}}$，其中
$W \subset X$ 是闭子集。显然 $W = T$，故结论成立。
\end{proof}

\begin{lemma}
\label{varieties-lemma-finite-extension-geometrically-irreducible-components}
设 $k$ 为域，设 $X \to \Spec(k)$ 局部有限型。假设 $X$ 只有有限多个
不可约分支。那么存在有限可分扩张 $k'/k$，使得 $X_{k'}$ 的每个
不可约分支都在 $k'$ 上几何不可约。
\end{lemma}

\begin{proof}
设 $\overline{k}$ 为 $k$ 的可分代数闭包。$X$ 只有有限多个不可约分支
这一假设，与引理
\ref{varieties-lemma-galois-action-irreducible-components-locally-finite-type} (3)
合起来表明，$X_{\overline{k}}$ 只有有限多个不可约分支
$\overline{T}_1, \ldots, \overline{T}_n$。由引理
\ref{varieties-lemma-galois-action-irreducible-components-locally-finite-type} (2)，
存在有限扩张 $\overline{k}/k'/k$ 和不可约分支 $T_i \subset X_{k'}$，
使得 $\overline{T}_i = T_{i, \overline{k}}$，故结论成立。
\end{proof}

\begin{lemma}
\label{varieties-lemma-irreducible-components-geometrically-irreducible}
设 $X$ 为域 $k$ 上的概形。假设 $X$ 只有有限多个不可约分支，并且它们
全都几何不可约。那么 $X$ 只有有限多个连通分支，并且每个连通分支都
几何连通。
\end{lemma}

\begin{proof}
这是显然的，因为一个连通分支是若干不可约分支之并。细节略去。
\end{proof}


\section{几何整概形}
\label{varieties-section-geometrically-integral}

\noindent
若 $X$ 是域上的整概形，则在扩张基域后，$X$ 可能变得非约化或可约。
几何整概形不会发生这种情况。

\begin{definition}
\label{varieties-definition-geometrically-integral}
设 $X$ 为域 $k$ 上的概形。
\begin{enumerate}
\item 设 $x \in X$。称 $X$ 在 $x$ 处{\it 几何逐点整}，如果对每个域扩张
$k'/k$，以及每个 $x' \in X_{k'}$，只要它位于 $x$ 上方，局部环
$\mathcal{O}_{X_{k'}, x'}$ 都是整环。
\item 若 $X$ 在每一点处都几何逐点整，则称 $X$ {\it 几何逐点整}。
\item 称 $X$ 在 $k$ 上{\it 几何整}，如果概形 $X_{k'}$ 对每个域扩张
$k'$（它是 $k$ 的扩张）都整。
\end{enumerate}
\end{definition}

\noindent
概念 (2) 与 (3) 之间的区别是必要的。例如，若
$k = \mathbf{R}$ 且 $X = \Spec(\mathbf{C}[x])$，则 $X$ 在
$\mathbf{R}$ 上几何逐点整，但当然不是几何整的。

\begin{lemma}
\label{varieties-lemma-geometrically-integral}
设 $k$ 为域，设 $X$ 为 $k$ 上的概形。那么，$X$ 在 $k$ 上几何整，
当且仅当 $X$ 在 $k$ 上既几何约化又几何不可约。
\end{lemma}

\begin{proof}
见《性质》引理 \ref{properties-lemma-characterize-integral}。
\end{proof}

\begin{lemma}
\label{varieties-lemma-proper-geometrically-reduced-global-sections}
设 $k$ 为域，设 $X$ 为 $k$ 上的固有概形。
\begin{enumerate}
\item $A = H^0(X, \mathcal{O}_X)$ 是有限维 $k$-代数；
\item $A = \prod_{i = 1, \ldots, n} A_i$ 是 Artin 局部
$k$-代数的直积，$X$ 的每个连通分支对应一个因子；
\item 若 $X$ 约化，则 $A = \prod_{i = 1, \ldots, n} k_i$
是若干域的直积，其中每个域都是 $k$ 的有限扩张；
\item 若 $X$ 几何约化，则 $k_i$ 是 $k$ 的有限可分扩张；
\item 若 $X$ 几何连通，则 $A$ 在 $k$ 上几何不可约；
\item 若 $X$ 几何不可约，则 $A$ 在 $k$ 上几何不可约；
\item 若 $X$ 几何约化且几何连通，则 $A = k$；
\item 若 $X$ 几何整，则 $A = k$。
\end{enumerate}
\end{lemma}

\begin{proof}
由《概形的上同调》引理
\ref{coherent-lemma-proper-over-affine-cohomology-finite}，可知
$A = H^0(X, \mathcal{O}_X)$ 是有限维 $k$-代数。这证明了 (1)。

\medskip\noindent
于是，由《代数》引理 \ref{algebra-lemma-finite-dimensional-algebra}
和命题 \ref{algebra-proposition-dimension-zero-ring}，$A$ 是若干 Artin
局部 $k$-代数的直积。若 $X = Y \amalg Z$，其中 $Y$ 与 $Z$ 在 $X$
中都是开集，则可得到幂等元 $e \in A$：取 $\mathcal{O}_X$ 的一个截面，
使其取值为 $1$（在 $Y$ 上）、取值为 $0$（在 $Z$ 上）。反过来，若 $e \in A$
是幂等元，则得到 $X$ 的相应分解。最后，由于 $X$ 的底层拓扑空间是
Noether 空间，它的连通分支都是开的。因此，$X$ 的连通分支以
$1$-to-$1$ 的方式对应于 $A$ 的本原幂等元。这证明了 (2)。

\medskip\noindent
若 $X$ 约化，则 $A$ 约化。因此，局部环 $A_i = k_i$ 约化，从而是域
（例如可用《代数》引理
\ref{algebra-lemma-minimal-prime-reduced-ring}）。这证明了 (3)。

\medskip\noindent
若 $X$ 几何约化，则
$A \otimes_k \overline{k} =
H^0(X_{\overline{k}}, \mathcal{O}_{X_{\overline{k}}})$
（该等式来自《概形的上同调》引理
\ref{coherent-lemma-flat-base-change-cohomology}）是约化的。这蕴含
$k_i \otimes_k \overline{k}$ 是若干域的直积，因而 $k_i/k$ 可分；
例如可用《代数》引理
\ref{algebra-lemma-characterize-separable-field-extensions} 和
\ref{algebra-lemma-geometrically-reduced-finite-purely-inseparable-extension}。
这证明了 (4)。

\medskip\noindent
若 $X$ 几何连通，则
$A \otimes_k \overline{k} =
H^0(X_{\overline{k}}, \mathcal{O}_{X_{\overline{k}}})$
由 (2) 是零维局部环，故它的谱只有一个点，特别地，它不可约。因此
$A$ 几何不可约。这证明了 (5)。当然，(5) 蕴含 (6)。

\medskip\noindent
若 $X$ 几何约化且几何连通，则 $A = k_1$ 是域，而扩张 $k_1/k$
有限可分且几何不可约。然而，此时 $k_1 \otimes_k \overline{k}$
是 $[k_1 : k]$ 份 $\overline{k}$ 的直积，故可知 $k_1 = k$。
这证明了 (7)。当然，(7) 蕴含 (8)。
\end{proof}

\noindent
下面给出 Stein 分解的一个初步版本；真正的 Stein 分解将在《态射续篇》第
\ref{more-morphisms-section-stein-factorization} 节讨论。

\begin{lemma}
\label{varieties-lemma-baby-stein}
设 $X$ 为域 $k$ 上的固有概形。置
$A = H^0(X, \mathcal{O}_X)$。典范态射 $X \to \Spec(A)$ 的纤维几何连通。
\end{lemma}

\begin{proof}
置 $S = \Spec(A)$。典范态射 $X \to S$ 是通过《概形》引理
\ref{schemes-lemma-morphism-into-affine} 与等式
$\Gamma(S, \mathcal{O}_S) = A = \Gamma(X, \mathcal{O}_X)$
相对应的态射。$k$-代数 $A$ 是有限直积
$A = \prod A_i$，其因子是 $k$ 上有限的 Artin 局部 $k$-代数；见引理
\ref{varieties-lemma-proper-geometrically-reduced-global-sections}。用
$s_i \in S$ 表示对应于 $A_i$ 的极大理想的点。选取代数闭包
$\overline{k}$，它是 $k$ 的一个代数闭包，并置
$\overline{A} = A \otimes_k \overline{k}$。
选取嵌入 $\kappa(s_i) \to \overline{k}$，它在 $k$ 上；该嵌入确定
$\overline{k}$-代数映射
$$
\sigma_i : \overline{A} = A \otimes_k \overline{k} \to
\kappa(s_i) \otimes_k \overline{k} \to \overline{k}
$$
考虑基变换
$$
\xymatrix{
\overline{X} \ar[r] \ar[d] & X \ar[d] \\
\overline{S} \ar[r] & S
}
$$
它把 $X$ 基变换到 $\overline{S} = \Spec(\overline{A})$。由
《概形的上同调》引理
\ref{coherent-lemma-flat-base-change-cohomology}，有
$\Gamma(\overline{X}, \mathcal{O}_{\overline{X}}) = \overline{A}$。
若 $\overline{s}_i \in \Spec(\overline{A})$ 表示一个
$\overline{k}$-有理点，它对应于 $\sigma_i$，则 $\overline{s}_i$ 映到
$s_i \in S$，并且
$\overline{X}_{\overline{s}_i}$ 是 $X_{s_i}$ 沿 $\Spec(\sigma_i)$
的基变换。因此，只需在 $k$ 代数闭的情形下证明本引理。

\medskip\noindent
假设 $k$ 代数闭。在这种情况下，$\kappa(s_i)$ 代数闭，而我们需要证明
$X_{s_i}$ 连通。直积分解 $A = \prod A_i$ 对应于不交并分解
$\Spec(A) = \coprod \Spec(A_i)$；见《代数》引理
\ref{algebra-lemma-spec-product}。用 $X_i$ 表示 $\Spec(A_i)$ 的逆像。
由引理 \ref{varieties-lemma-proper-geometrically-reduced-global-sections} (2)，有
$A_i = \Gamma(X_i, \mathcal{O}_{X_i})$。注意，
$X_{s_i} \to X_i$ 是闭浸入，并在底层拓扑空间上诱导同构（因为
$\Spec(A_i)$ 是单点集）。因此，若 $X_{s_i}$ 不连通，则 $X_i$
也不连通。所以，要么 $X_i$ 为空且 $A_i = 0$，要么 $X_i$
可写成 $U \amalg V$，其中 $U$ 与 $V$ 是非空开集；这将蕴含 $A_i$
有非平凡幂等元。由于 $A_i$ 是局部环，这构成矛盾，证明完成。
\end{proof}

\begin{lemma}
\label{varieties-lemma-geometrically-reduced-stein}
设 $k$ 为域，设 $X$ 为 $k$ 上固有且几何约化的概形。下列条件等价：
\begin{enumerate}
\item $H^0(X, \mathcal{O}_X) = k$；
\item $X$ 几何连通。
\end{enumerate}
\end{lemma}

\begin{proof}
由引理 \ref{varieties-lemma-baby-stein}，有 (1) $\Rightarrow$ (2)。
由引理 \ref{varieties-lemma-proper-geometrically-reduced-global-sections}，
有 (2) $\Rightarrow$ (1)。
\end{proof}


\section{几何正规概形}
\label{varieties-section-geometrically-normal}

\noindent
我们已在《性质》定义 \ref{properties-definition-normal} 中定义了正规概形的
概念。这个概念甚至对非 Noether 概形也有定义。因此，与我们对
“几何正则”概形的讨论不同，这里要考虑基域的所有域扩张。

\begin{definition}
\label{varieties-definition-geometrically-normal}
设 $X$ 为域 $k$ 上的概形。
\begin{enumerate}
\item 设 $x \in X$。称 $X$ 在 $x$ 处{\it 几何正规}，如果对每个域扩张
$k'/k$，以及每个 $x' \in X_{k'}$，只要它位于 $x$ 上方，局部环
$\mathcal{O}_{X_{k'}, x'}$ 都正规。
\item 称 $X$ 在 $k$ 上{\it 几何正规}，如果 $X$ 在每个 $x \in X$
处都几何正规。
\end{enumerate}
\end{definition}

\begin{lemma}
\label{varieties-lemma-geometrically-normal-at-point}
设 $k$ 为域，设 $X$ 为 $k$ 上的概形，并设 $x \in X$。下列条件等价：
\begin{enumerate}
\item $X$ 在 $x$ 处几何正规；
\item 对每个有限纯不可分域扩张 $k'$（它是 $k$ 的扩张），以及每个
$x' \in X_{k'}$，只要它位于 $x$ 上方，局部环
$\mathcal{O}_{X_{k'}, x'}$ 就正规；
\item 环 $\mathcal{O}_{X, x}$ 在 $k$ 上几何正规（见《代数》定义
\ref{algebra-definition-geometrically-normal}）。
\end{enumerate}
\end{lemma}

\begin{proof}
(1) 蕴含 (2) 是显然的。假设 (2)。设 $k'/k$ 为有限纯不可分域扩张
（例如 $k = k'$）。考虑环 $\mathcal{O}_{X, x} \otimes_k k'$。
由《代数》引理 \ref{algebra-lemma-p-ring-map}，它的谱与
$\mathcal{O}_{X, x}$ 的谱相同。因此，它也是局部环（《代数》引理
\ref{algebra-lemma-characterize-local-ring}）。所以，存在唯一的点
$x' \in X_{k'}$ 位于 $x$ 上方，并且
$\mathcal{O}_{X_{k'}, x'} \cong \mathcal{O}_{X, x} \otimes_k k'$。
由假设，这是正规环。因此，由《代数》引理
\ref{algebra-lemma-geometrically-normal} 得到 (3)。

\medskip\noindent
假设 (3)。设 $k'/k$ 为域扩张。由于
$\Spec(k') \to \Spec(k)$ 满射，$X_{k'} \to X$ 也满射
（《态射》引理 \ref{morphisms-lemma-base-change-surjective}）。
设 $x' \in X_{k'}$ 为位于 $x$ 上方的任意一点。局部环
$\mathcal{O}_{X_{k'}, x'}$ 是环
$\mathcal{O}_{X, x} \otimes_k k'$ 的一个局部化。因此，它由假设是正规
的，从而证明了 (1)。
\end{proof}

\begin{lemma}
\label{varieties-lemma-geometrically-normal}
设 $k$ 为域，设 $X$ 为 $k$ 上的概形。下列条件等价：
\begin{enumerate}
\item $X$ 几何正规；
\item $X_{k'}$ 对每个域扩张 $k'/k$ 都是正规概形；
\item $X_{k'}$ 对每个有限生成域扩张 $k'/k$ 都是正规概形；
\item $X_{k'}$ 对每个有限纯不可分域扩张 $k'/k$ 都是正规概形；
\item 对每个仿射开集 $U \subset X$，环 $\mathcal{O}_X(U)$ 都几何正规
（见《代数》定义 \ref{algebra-definition-geometrically-normal}）；
\item $X_{k^{perf}}$ 是正规概形。
\end{enumerate}
\end{lemma}

\begin{proof}
假设 (1)。那么，对每个域扩张 $k'/k$ 和每个点
$x' \in X_{k'}$，$X_{k'}$ 在 $x'$ 处的局部环都正规。按定义，这意味着
$X_{k'}$ 正规。因此 (2) 成立。

\medskip\noindent
(2) 蕴含 (3)、(3) 蕴含 (4)，这都是显然的。

\medskip\noindent
假设 (4)，并设 $U \subset X$ 为仿射开子概形。那么，$U_{k'}$ 对任意
有限纯不可分扩张 $k'/k$ 都是正规概形（包括 $k = k'$）。这意味着，环
$k' \otimes_k \mathcal{O}(U)$ 对所有有限纯不可分扩张 $k'/k$ 都正规。因此，
$\mathcal{O}(U)$ 按定义是几何正规 $k$-代数。故 (4) 蕴含 (5)。

\medskip\noindent
假设 (5)。对任意域扩张 $k'/k$，基变换 $X_{k'}$ 由各个环
$\mathcal{O}_X(U) \otimes_k k'$ 的谱粘合而得，其中 $U$ 是 $X$ 中的
仿射开集（见《概形》第 \ref{schemes-section-fibre-products} 节）。
因此 $X_{k'}$ 正规，所以 (1) 成立。

\medskip\noindent
(5) 与 (6) 的等价性来自几何正规代数的定义，以及刚刚证明的 (3) 与
(4) 的等价性。
\end{proof}

\begin{lemma}
\label{varieties-lemma-geometrically-normal-upstairs}
设 $k$ 为域，设 $X$ 为 $k$ 上的概形，设 $k'/k$ 为域扩张。设
$x \in X$ 为一点，并设 $x' \in X_{k'}$ 为位于 $x$ 上方的一点。
下列条件等价：
\begin{enumerate}
\item $X$ 在 $x$ 处几何正规；
\item $X_{k'}$ 在 $x'$ 处几何正规。
\end{enumerate}
特别地，$X$ 在 $k$ 上几何正规，当且仅当 $X_{k'}$ 在 $k'$ 上几何正规。
\end{lemma}

\begin{proof}
(1) 蕴含 (2) 是显然的。假设 (2)。设 $k''/k$ 为有限纯不可分域扩张，
并设 $x'' \in X_{k''}$ 为位于 $x$ 上方的一点（实际上它是唯一的）。
选取公共域扩张 $k'''/k$，它同时扩张 $k'/k$ 与 $k''/k$；见《域》引理
\ref{fields-lemma-common-extension-field}。选取点
$x''' \in X_{k'''}$，它同时位于 $x'$ 与 $x''$ 上方。考虑局部环映射
$$
\mathcal{O}_{X_{k''}, x''} \longrightarrow \mathcal{O}_{X_{k'''}, x''''}.
$$
这是平坦局部环同态，因而忠实平坦。由 (2)，右侧的局部环正规。因此由
《代数》引理 \ref{algebra-lemma-descent-normal}，可知
$\mathcal{O}_{X_{k''}, x''}$ 正规。由引理
\ref{varieties-lemma-geometrically-normal-at-point}，可知 $X$ 在 $x$ 处几何正规。
\end{proof}

\begin{lemma}
\label{varieties-lemma-fibre-product-normal}
设 $k$ 为域，设 $X$ 为 $k$ 上几何正规的概形，并设 $Y$ 为 $k$ 上正规的
概形。那么 $X \times_k Y$ 是正规概形。
\end{lemma}

\begin{proof}
这归结为《代数》引理
\ref{algebra-lemma-geometrically-normal-tensor-normal}；所用的归约来自引理
\ref{varieties-lemma-geometrically-normal}。
\end{proof}

\begin{lemma}
\label{varieties-lemma-base-change-normal-by-separable}
设 $k$ 为域，设 $X$ 为 $k$ 上的正规概形。设 $K/k$ 为可分域扩张。
那么 $X_K$ 是正规概形。
\end{lemma}

\begin{proof}
这来自引理 \ref{varieties-lemma-fibre-product-normal} 和《代数》引理
\ref{algebra-lemma-separable-field-extension-geometrically-normal}。
\end{proof}

\begin{lemma}
\label{varieties-lemma-geometrically-normal-stein}
设 $k$ 为域，设 $X$ 为 $k$ 上固有且几何正规的概形。下列条件等价：
\begin{enumerate}
\item $H^0(X, \mathcal{O}_X) = k$；
\item $X$ 几何连通；
\item $X$ 几何不可约；
\item $X$ 几何整。
\end{enumerate}
\end{lemma}

\begin{proof}
由引理 \ref{varieties-lemma-geometrically-reduced-stein}，(1) 与 (2) 等价。局部
Noether 正规概形（例如 $X_{\overline{k}}$）是其不可约分支的不交并
（《性质》引理 \ref{properties-lemma-normal-Noetherian}）。因此 (2) 与
(3) 等价。由于假设 $X_{\overline{k}}$ 约化，(3) 与 (4) 也等价。
\end{proof}


\section{域变换与局部 Noether 概形}
\label{varieties-section-locally-Noetherian}

\noindent
设 $X$ 为域 $k$ 上的局部 Noether 概形。$X_{k'}$ 不一定仍是局部
Noether 概形。例如，若 $X = \Spec(\overline{\mathbf{Q}})$ 且
$k = \mathbf{Q}$，则 $X_{\overline{\mathbf{Q}}}$ 是
$\overline{\mathbf{Q}} \otimes_{\mathbf{Q}} \overline{\mathbf{Q}}$
的谱，而后者不是 Noether 环。（提示：它有太多幂等元。）但是，如果只沿
有限生成域扩张作基变换，则 Noether 性得以保持。（或者，若 $X$ 在
$k$ 上局部有限型，结论也成立，因为这个性质在基变换下保持。）

\begin{lemma}
\label{varieties-lemma-locally-Noetherian-base-change}
设 $k$ 为域，设 $X$ 为 $k$ 上的概形，并设 $k'/k$ 为有限生成域扩张。
那么，$X$ 局部 Noether，当且仅当 $X_{k'}$ 局部 Noether。
\end{lemma}

\begin{proof}
利用《性质》引理 \ref{properties-lemma-locally-Noetherian}，可归约到
$X$ 仿射的情形，比如说 $X = \Spec(A)$。在这种情况下，需要证明
$A$ 是 Noether 环，当且仅当 $A_{k'}$ 是 Noether 环。由于
$A \to A_{k'} = k' \otimes_k A$ 忠实平坦，由《代数》引理
\ref{algebra-lemma-descent-Noetherian}，若 $A_{k'}$ 是 Noether 环，
则 $A$ 也是 Noether 环。反过来，若 $A$ 是 Noether 环，则由《代数》引理
\ref{algebra-lemma-Noetherian-field-extension}，$A_{k'}$ 是 Noether 环。
\end{proof}


\section{几何正则概形}
\label{varieties-section-geometrically-regular}

\noindent
域 $k$ 上的几何正则概形，是域 $k$ 上的局部 Noether 概形，并且在适当的
基域变换后仍为正则概形。域 $k$ 上的有限型概形几何正则，当且仅当它在
$k$ 上光滑（见引理 \ref{varieties-lemma-geometrically-regular-smooth}）。几何正则性
这个概念最有意义的情形，是无法使用光滑性的情形，例如形式纤维（此处以后
插入引用）。

\medskip\noindent
在下面的定义中，我们限于局部 Noether 概形，因为正则局部环这个性质只对
Noether 局部环定义。由上面的引理
\ref{varieties-lemma-locally-Noetherian-base-change}，如果只考虑有限生成域扩张，
这个性质在基域变换下保持。本节将直接使用这一说明，不再另行引用。特别地，
下面的定义是有意义的。

\begin{definition}
\label{varieties-definition-geometrically-regular}
设 $k$ 为域，设 $X$ 为 $k$ 上的局部 Noether 概形。
\begin{enumerate}
\item 设 $x \in X$。称 $X$ 在 $x$ 处相对于 $k$ {it 几何正则}，如果对
每个有限生成域扩张 $k'/k$ 以及任意 $x' \in X_{k'}$，只要后者位于 $x$
上方，局部环 $\mathcal{O}_{X_{k'}, x'}$ 就正则。
\item 称 $X$ {it 在 $k$ 上几何正则}，如果 $X$ 在其所有点处都几何正则。
\end{enumerate}
\end{definition}

\noindent
类似的定义也可用于定义域上的几何 Cohen--Macaulay、$(R_k)$ 与 $(S_k)$
概形。需要时，我们会另设一节讨论这些性质。

\begin{lemma}
\label{varieties-lemma-geometrically-regular-at-point}
设 $k$ 为域，设 $X$ 为 $k$ 上的局部 Noether 概形，并设 $x \in X$。
下列条件等价：
\begin{enumerate}
\item $X$ 在 $x$ 处几何正则；
\item 对每个有限纯不可分扩域 $k'$（其底域为 $k$）以及
$x' \in X_{k'}$，只要后者位于 $x$ 上方，局部环
$\mathcal{O}_{X_{k'}, x'}$ 就正则；
\item 环 $\mathcal{O}_{X, x}$ 在 $k$ 上几何正则（见《代数》定义
\ref{algebra-definition-geometrically-regular}）。
\end{enumerate}
\end{lemma}

\begin{proof}
(1) 蕴含 (2) 是清楚的。假设 (2)。特别地，这蕴含 $\mathcal{O}_{X, x}$
是正则局部环。设 $k'/k$ 为有限纯不可分域扩张。考虑环
$\mathcal{O}_{X, x} \otimes_k k'$。由《代数》引理
\ref{algebra-lemma-p-ring-map}，它的谱与 $\mathcal{O}_{X, x}$ 的谱相同，
因而它也是局部环（《代数》引理
\ref{algebra-lemma-characterize-local-ring}）。所以存在唯一一点
$x' \in X_{k'}$ 位于 $x$ 上方，且
$\mathcal{O}_{X_{k'}, x'} \cong \mathcal{O}_{X, x} \otimes_k k'$。
根据假设，这是正则环。因此，由几何正则环的定义得到 (3)。

\medskip\noindent
假设 (3)。设 $k'/k$ 为域扩张。由于
$\Spec(k') \to \Spec(k)$ 满射，
$X_{k'} \to X$ 也满射（《态射》引理
\ref{morphisms-lemma-base-change-surjective}）。任取一点
$x' \in X_{k'}$，使它位于 $x$ 上方。局部环
$\mathcal{O}_{X_{k'}, x'}$ 是环
$\mathcal{O}_{X, x} \otimes_k k'$ 的局部化。因此，根据假设它正则，
从而 (1) 得证。
\end{proof}

\begin{lemma}
\label{varieties-lemma-geometrically-regular}
设 $k$ 为域，设 $X$ 为 $k$ 上的局部 Noether 概形。下列条件等价：
\begin{enumerate}
\item $X$ 几何正则；
\item $X_{k'}$ 对每个有限生成域扩张 $k'/k$ 都是正则概形；
\item $X_{k'}$ 对每个有限纯不可分域扩张 $k'/k$ 都是正则概形；
\item 对每个仿射开集 $U \subset X$，环 $\mathcal{O}_X(U)$ 几何正则
（见《代数》定义 \ref{algebra-definition-geometrically-regular}）；
\item 存在仿射开覆盖 $X = \bigcup U_i$，使每个 $\mathcal{O}_X(U_i)$
都在 $k$ 上几何正则。
\end{enumerate}
\end{lemma}

\begin{proof}
假设 (1)。那么，对每个有限生成域扩张 $k'/k$ 以及每个点
$x' \in X_{k'}$，$X_{k'}$ 在 $x'$ 处的局部环都正则。由《性质》引理
\ref{properties-lemma-characterize-regular}，这意味着 $X_{k'}$ 正则，
所以 (2) 成立。

\medskip\noindent
(2) 蕴含 (3) 是清楚的。

\medskip\noindent
假设 (3)，并设 $U \subset X$ 为仿射开子概形。那么 $U_{k'}$ 对任意有限
纯不可分扩张 $k'/k$（包括 $k = k'$）都是正则概形。这意味着
$k' \otimes_k \mathcal{O}(U)$ 对所有有限纯不可分扩张 $k'/k$ 都是正则环。因此，
$\mathcal{O}(U)$ 是几何正则 $k$-代数，从而 (4) 成立。

\medskip\noindent
(4) 蕴含 (5) 是清楚的。设 $X = \bigcup U_i$ 为 (5) 中的仿射开覆盖。
对任意域扩张 $k'/k$，基变换 $X_{k'}$ 由各环
$\mathcal{O}_X(U_i) \otimes_k k'$ 的谱粘合而成（见《概形》第
\ref{schemes-section-fibre-products} 节）。因此 $X_{k'}$ 正则，所以 (1)
成立。
\end{proof}

\begin{lemma}
\label{varieties-lemma-geometrically-regular-upstairs}
设 $k$ 为域，设 $X$ 为 $k$ 上的概形。设 $k'/k$ 为有限生成域扩张。
设 $x \in X$，并设 $x' \in X_{k'}$ 为位于 $x$ 上方的点。下列条件等价：
\begin{enumerate}
\item $X$ 在 $x$ 处几何正则；
\item $X_{k'}$ 在 $x'$ 处几何正则。
\end{enumerate}
特别地，$X$ 在 $k$ 上几何正则，当且仅当 $X_{k'}$ 在 $k'$ 上几何正则。
\end{lemma}

\begin{proof}
(1) 蕴含 (2) 是清楚的。假设 (2)。设 $k''/k$ 为有限纯不可分域扩张，
并设 $x'' \in X_{k''}$ 为位于 $x$ 上方的点（事实上它是唯一的）。我们可以
找到一个公共有限生成域扩张 $k'''/k$，它同时扩张 $k'/k$ 和 $k''/k$，见
《域》引理 \ref{fields-lemma-common-extension-field} 及其证明。并且可以找到
一点 $x''' \in X_{k'''}$，它同时位于 $x'$ 与 $x''$ 上方。考虑局部环映射
$$
\mathcal{O}_{X_{k''}, x''} \longrightarrow \mathcal{O}_{X_{k'''}, x''''}.
$$
这是 Noether 局部环之间的平坦局部同态，因而忠实平坦。由 (2)，右侧局部环
正则。于是由《代数》引理 \ref{algebra-lemma-flat-under-regular}，可得
$\mathcal{O}_{X_{k''}, x''}$ 正则。由引理
\ref{varieties-lemma-geometrically-regular-at-point}，$X$ 在 $x$ 处几何正则。
\end{proof}

\noindent
下一个引理是《代数》引理
\ref{algebra-lemma-geometrically-regular-descent} 的几何版本。

\begin{lemma}
\label{varieties-lemma-flat-under-geometrically-regular}
设 $k$ 为域，设 $f : X \to Y$ 为 $k$ 上局部 Noether 概形之间的态射。
设 $x \in X$，并置 $y = f(x)$。如果 $X$ 在 $x$ 处几何正则，且 $f$
在 $x$ 处平坦，那么 $Y$ 在 $y$ 处几何正则。特别地，如果 $X$ 在 $k$
上几何正则，且 $f$ 平坦满射，那么 $Y$ 在 $k$ 上几何正则。
\end{lemma}

\begin{proof}
设 $k'$ 为 $k$ 的有限纯不可分扩张。设 $f' : X_{k'} \to Y_{k'}$ 为 $f$
的基变换，并设 $x' \in X_{k'}$ 为位于 $x$ 上方的唯一一点。如果证明
$Y_{k'}$ 在 $y' = f'(x')$ 处正则，那么由引理
\ref{varieties-lemma-geometrically-regular}，$Y$ 在 $k$ 上于 $y'$ 处几何正则。
由《态射》引理 \ref{morphisms-lemma-base-change-module-flat}，态射
$X_{k'} \to Y_{k'}$ 在 $x'$ 处平坦。因此环映射
$$
\mathcal{O}_{Y_{k'}, y'}
\longrightarrow
\mathcal{O}_{X_{k'}, x'}
$$
是局部 Noether 环之间的平坦局部同态，而其右侧根据假设正则。因此由
《代数》引理 \ref{algebra-lemma-flat-under-regular}，其左侧是正则局部环。
\end{proof}

\begin{lemma}
\label{varieties-lemma-geometrically-regular-smooth}
设 $k$ 为域，设 $X$ 为 $k$ 上局部有限型的概形，并设 $x \in X$。
那么，$X$ 在 $x$ 处几何正则，当且仅当 $X \to \Spec(k)$ 在 $x$ 处光滑
（《态射》定义 \ref{morphisms-definition-smooth}）。
\end{lemma}

\begin{proof}
问题在 $x$ 附近是局部的，因而可以假设 $X = \Spec(A)$，其中所用代数是
某个有限型 $k$-代数。设 $x$ 对应于素理想 $\mathfrak p$。

\medskip\noindent
如果 $A$ 在 $k$ 上于 $\mathfrak p$ 处光滑，那么可以局部化 $A$，并假设
$A$ 在 $k$ 上光滑。在这种情况下，$k' \otimes_k A$ 对所有扩域都在
$k'$ 上光滑，其中 $k'/k$ 任意；由《代数》引理
\ref{algebra-lemma-characterize-smooth-over-field}，这些 Noether 环都正则。

\medskip\noindent
假设 $X$ 在 $x$ 处几何正则。考虑剩余域
$K := \kappa(x) = \kappa(\mathfrak p)$，它属于点 $x$。这是 $k$ 的有限生成扩张。由
《代数》引理 \ref{algebra-lemma-make-separable}，存在有限纯不可分扩张
$k'/k$，使合成域 $k'K$ 是 $k'$ 的可分域扩张。设
$\mathfrak p' \subset A' = k' \otimes_k A$ 为位于 $\mathfrak p$ 上方的
素理想。由《代数》引理 \ref{algebra-lemma-p-ring-map}，它是位于
$\mathfrak p$ 上方的唯一素理想。因此剩余域
$K' := \kappa(\mathfrak p')$ 就是合成域 $k'K$。根据假设，局部环
$(A')_{\mathfrak p'}$ 正则。于是由《代数》引理
\ref{algebra-lemma-separable-smooth}，$k' \to A'$ 在 $\mathfrak p'$ 处光滑。
进而由《代数》引理 \ref{algebra-lemma-smooth-field-change-local}，
$k \to A$ 在 $\mathfrak p$ 处光滑。引理得证。
\end{proof}


\begin{example}
\label{varieties-example-geometrically-reduced-not-normal}
设 $k =\mathbf{F}_p(t)$。很容易给出一个正则但非几何约化的簇 $V$，它在
$k$ 上定义。例如可取 $\Spec(k[x]/(x^p - t))$。事实上，还存在正则且几何约化、
但甚至不几何正规的簇 $V$。具体地，当 $p > 2$ 时取
$V = \Spec(k[x, y]/(y^2 - x^p + t))$。由于多项式
$y^2 - x^p + t \in k[x, y]$ 不可约，这是一个簇。态射
$V \to \Spec(k)$ 在所有点处都光滑，唯独可能在点 $v_0 \in V$ 处不光滑；
该点对应于极大理想 $(y, x^p - t)$（因为 $2y$ 可逆）。特别地，
我们看到 $V$ 在所有点处都（几何）正则，唯独 $v_0$ 处尚有可能例外。
局部环
$$
\mathcal{O}_{V, v_0} = \left(k[x, y]/(y^2 - x^p + t)\right)_{(y, x^p - t)}
$$
是维数为 $1$ 的整环。它的极大理想由 $1$ 个元素生成，即由 $y$ 生成。
所以它是离散赋值环，从而正则。设 $k' = k[t^{1/p}]$。记
$t' = t^{1/p} \in k'$、$V' = V_{k'}$，并记 $v'_0 \in V'$ 为位于 $v_0$
上方的唯一一点。在 $k'$ 上，可写 $x^p - t = (x - t')^p$，但多项式
$y^2 - (x - t')^p$ 仍不可约，且 $V'$ 仍是簇。然而元素
$$
\frac{y}{x - t'} \in (\text{fraction field of }\mathcal{O}_{V', v'_0})
$$
在 $\mathcal{O}_{V', v'_0}$ 上整（只须计算它的平方），却不属于该环；
所以 $V'$ 在 $v'_0$ 处不正规。这就完成了这个例子。
\end{example}


\section{域变换与 Cohen--Macaulay 性}
\label{varieties-section-CM}

\noindent
下一个引理说明，定义几何 Cohen--Macaulay 概形没有意义，因为这样的概形
恰好就是 Cohen--Macaulay 概形。

\begin{lemma}
\label{varieties-lemma-CM-base-change}
设 $X$ 为域 $k$ 上的局部 Noether 概形。设 $k'/k$ 为有限生成域扩张。
设 $x \in X$，并设 $x' \in X_{k'}$ 为位于 $x$ 上方的点。那么
$$
\mathcal{O}_{X, x}\text{ 为 Cohen--Macaulay}
\Leftrightarrow
\mathcal{O}_{X_{k'}, x'}\text{ 为 Cohen--Macaulay}
$$
如果 $X$ 在 $k$ 上局部有限型，那么对任意域扩张 $k'/k$，同一结论成立。
\end{lemma}

\begin{proof}
引理的第一种情形来自《代数》引理
\ref{algebra-lemma-CM-geometrically-CM}。引理的第二种情形等价于《代数》引理
\ref{algebra-lemma-extend-field-CM-locus}。
\end{proof}


\section{域变换与 Jacobson 性}
\label{varieties-section-overfield}

\noindent
域上局部有限型的概形有充足的闭点，也就是说，它是 Jacobson 概形。此外，
这些剩余域都是基域的有限扩张。

\begin{lemma}
\label{varieties-lemma-locally-finite-type-Jacobson}
设 $X$ 为在 $k$ 上局部有限型的概形。那么：
\begin{enumerate}
\item 对任意闭点 $x \in X$，扩张 $\kappa(x)/k$ 是代数扩张；
\item $X$ 是 Jacobson 概形（《性质》定义
\ref{properties-definition-jacobson}）。
\end{enumerate}
\end{lemma}

\begin{proof}
一个概形是 Jacobson 概形，当且仅当它有一个由 Jacobson 概形组成的仿射
开覆盖，见《性质》引理 \ref{properties-lemma-locally-jacobson}。闭点处剩余域
的性质对 $X$ 也是局部的。因此可以假设 $X$ 仿射。在这种情况下，结论是
Hilbert 零点定理的推论，见《代数》定理
\ref{algebra-theorem-nullstellensatz}。它也可由《态射》引理
\ref{morphisms-lemma-jacobson-finite-type-points}、
\ref{morphisms-lemma-Jacobson-universally-Jacobson} 和
\ref{morphisms-lemma-ubiquity-Jacobson-schemes} 合并推出。
\end{proof}

\noindent
事实证明，如果 $X$ 不局部有限型，那么在作足够大的基域扩张后，可以得到
同样的结论。

\begin{lemma}
\label{varieties-lemma-make-Jacobson}
设 $X$ 为域 $k$ 上的概形。对任意基数足够大的域扩张 $K/k$，有：
\begin{enumerate}
\item 对任意闭点 $x \in X_K$，扩张 $\kappa(x)/K$ 是代数扩张；
\item $X_K$ 是 Jacobson 概形（《性质》定义
\ref{properties-definition-jacobson}）。
\end{enumerate}
\end{lemma}

\begin{proof}
选取仿射开覆盖 $X = \bigcup U_i$。由《代数》引理
\ref{algebra-lemma-base-change-Jacobson} 和《性质》引理
\ref{properties-lemma-affine-jacobson}，存在基数 $\kappa_i$，使得只要
$U_{i, K}$ 是在 $K$ 上考虑的，且 $\#(K) \geq \kappa_i$，它就具有所需性质。
置 $\kappa = \max\{\kappa_i\}$。那么，如果 $K$ 的基数大于 $\kappa$，
每个 $U_{i, K}$ 都满足引理的结论。因此，由《性质》引理
\ref{properties-lemma-locally-jacobson}，$X_K$ 是 Jacobson 概形。关于
$X_K$ 的闭点处剩余域的断言，来自 $U_{i, K}$ 的闭点处剩余域的相应断言。
\end{proof}


\section{域变换与丰沛可逆层}
\label{varieties-section-change-fields-ample}

\noindent
下述结果是本节各结果的典型例子。

\begin{lemma}
\label{varieties-lemma-ample-after-field-extension}
设 $k$ 为域，设 $X$ 为 $k$ 上的概形。如果 $X_K$ 对某个域扩张 $K/k$
有丰沛可逆层，那么 $X$ 有丰沛可逆层。
\end{lemma}

\begin{proof}
设 $K/k$ 为域扩张，使得 $X_K$ 有丰沛可逆层 $\mathcal{L}$。态射
$X_K \to X$ 满射。因此，$X$ 作为拟紧概形的像而拟紧（《性质》定义
\ref{properties-definition-ample}）。由于 $X_K$ 拟分离（由《性质》引理
\ref{properties-lemma-affine-s-opens-cover-quasi-separated}），可知 $X$
拟分离：如果 $U, V \subset X$ 是仿射开集，那么
$(U \cap V)_K = U_K \cap V_K$ 拟紧，而
$(U \cap V)_K \to U \cap V$ 满射。因此《概形》引理
\ref{schemes-lemma-characterize-quasi-separated} 适用。

\medskip\noindent
将 $K = \colim A_i$ 写成 $K$ 的各有限型 $k$-子代数的余极限。记
$X_i = X \times_{\Spec(k)} \Spec(A_i)$。由于 $X_K = \lim X_i$，可找到
一个 $i$ 和一个可逆层 $\mathcal{L}_i$，后者定义在 $X_i$ 上，使它在 $X_K$ 上的拉回
是 $\mathcal{L}$（《极限》引理 \ref{limits-lemma-descend-invertible-modules}；
这里以及下文都使用刚刚证明的 $X$ 拟紧且拟分离）。由《极限》引理
\ref{limits-lemma-limit-ample}，可以假设 $\mathcal{L}_i$ 丰沛，只需在必要时
增大 $i$。固定这样的 $i$，并设
$\mathfrak m \subset A_i$ 为极大理想。由 Hilbert 零点定理（《代数》定理
\ref{algebra-theorem-nullstellensatz}），剩余域
$k' = A_i/\mathfrak m$ 是 $k$ 的有限扩张。因此
$X_{k'} \subset X_i$ 是闭子概形，从而有丰沛可逆层（《性质》引理
\ref{properties-lemma-ample-on-closed}）。由于 $X_{k'} \to X$ 有限局部自由，
由《除子》命题 \ref{divisors-proposition-push-down-ample}，可得 $X$ 有丰沛
可逆层。
\end{proof}

\begin{lemma}
\label{varieties-lemma-quasi-affine-after-field-extension}
设 $k$ 为域，设 $X$ 为 $k$ 上的概形。如果 $X_K$ 对某个域扩张 $K/k$
拟仿射，那么 $X$ 拟仿射。
\end{lemma}

\begin{proof}
设 $K/k$ 为域扩张，使得 $X_K$ 拟仿射。态射 $X_K \to X$ 满射。因此，
$X$ 作为拟紧概形的像而拟紧（《性质》定义
\ref{properties-definition-quasi-affine}）。由于 $X_K$ 拟分离（因为它是
仿射概形的开子概形），可知 $X$ 拟分离：如果 $U, V \subset X$ 是仿射
开集，那么 $(U \cap V)_K = U_K \cap V_K$ 拟紧，而
$(U \cap V)_K \to U \cap V$ 满射。因此《概形》引理
\ref{schemes-lemma-characterize-quasi-separated} 适用。

\medskip\noindent
将 $K = \colim A_i$ 写成 $K$ 的各有限型 $k$-子代数的余极限。记
$X_i = X \times_{\Spec(k)} \Spec(A_i)$。由于 $X_K = \lim X_i$，可找到
一个 $i$，使 $X_i$ 拟仿射（《极限》引理
\ref{limits-lemma-limit-quasi-affine}；这里使用刚刚证明的 $X$ 拟紧且拟分离）。
由 Hilbert 零点定理（《代数》定理
\ref{algebra-theorem-nullstellensatz}），剩余域 $k' = A_i/\mathfrak m$ 是
$k$ 的有限扩张。因此 $X_{k'} \subset X_i$ 是闭子概形，从而拟仿射
（《性质》引理 \ref{properties-lemma-quasi-affine-locally-closed}）。由于
$X_{k'} \to X$ 有限局部自由，结论来自《除子》引理
\ref{divisors-lemma-push-down-quasi-affine}。
\end{proof}

\begin{lemma}
\label{varieties-lemma-quasi-projective-after-field-extension}
设 $k$ 为域，设 $X$ 为 $k$ 上的概形。如果 $X_K$ 在 $K$ 上拟射影，其中
$K/k$ 是某个域扩张，那么 $X$ 在 $k$ 上拟射影。
\end{lemma}

\begin{proof}
根据定义，概形态射 $g : Y \to T$ 拟射影，是指它局部有限型、拟紧，并且
存在一个相对于 $g$ 丰沛的可逆层，定义在 $Y$ 上。设 $K/k$ 为域扩张，
使得 $X_K$ 在 $K$ 上拟射影。设 $\Spec(A) \subset X$ 为仿射开集。那么
$U_K$ 是 $X_K$ 的仿射开子概形，因而 $A_K$ 是有限型 $K$-代数。由《代数》
引理 \ref{algebra-lemma-finite-type-descends}，$A$ 是有限型 $k$-代数。因此
$X \to \Spec(k)$ 局部有限型。由于 $X_K \to \Spec(K)$ 拟紧，可知
$X_K$ 拟紧，进而 $X$ 拟紧，进而 $X \to \Spec(k)$ 有限型。由《态射》引理
\ref{morphisms-lemma-finite-type-over-affine-ample-very-ample}，可知 $X_K$
有丰沛可逆层。再由引理 \ref{varieties-lemma-ample-after-field-extension}，$X$ 有丰沛
可逆层。因此，由《态射》引理
\ref{morphisms-lemma-finite-type-over-affine-ample-very-ample}，
$X \to \Spec(k)$ 拟射影。
\end{proof}

\noindent
下一个引理是《下降》引理
\ref{descent-lemma-descending-property-proper} 的特殊情形。

\begin{lemma}
\label{varieties-lemma-proper-after-field-extension}
设 $k$ 为域，设 $X$ 为 $k$ 上的概形。如果 $X_K$ 在 $K$ 上固有，其中
$K/k$ 是某个域扩张，那么 $X$ 在 $k$ 上固有。
\end{lemma}

\begin{proof}
设 $K/k$ 为域扩张，使得 $X_K$ 在 $K$ 上固有。回忆，这蕴含 $X_K$
分离且拟紧（《态射》定义 \ref{morphisms-definition-proper}）。态射
$X_K \to X$ 满射。因此，$X$ 作为拟紧概形的像而拟紧（《性质》定义
\ref{properties-definition-ample}）。由于 $X_K$ 分离，可知 $X$ 拟分离：
如果 $U, V \subset X$ 是仿射开集，那么
$(U \cap V)_K = U_K \cap V_K$ 拟紧，而
$(U \cap V)_K \to U \cap V$ 满射。因此《概形》引理
\ref{schemes-lemma-characterize-quasi-separated} 适用。

\medskip\noindent
将 $K = \colim A_i$ 写成 $K$ 的各有限型 $k$-子代数的余极限。记
$X_i = X \times_{\Spec(k)} \Spec(A_i)$。由《极限》引理
\ref{limits-lemma-eventually-proper}，存在一个 $i$，使得
$X_i \to \Spec(A_i)$ 固有。这里使用刚刚证明的 $X$ 拟紧且拟分离。选取极大
理想 $\mathfrak m \subset A_i$。由 Hilbert 零点定理（《代数》定理
\ref{algebra-theorem-nullstellensatz}），剩余域 $k' = A_i/\mathfrak m$ 是
$k$ 的有限扩张。基变换 $X_{k'} \to \Spec(k')$ 固有（《态射》引理
\ref{morphisms-lemma-base-change-proper}）。由于 $k'/k$ 有限，
$X_{k'} \to X$ 以及复合态射 $X_{k'} \to \Spec(k)$ 也都固有（《态射》引理
\ref{morphisms-lemma-finite-proper}、
\ref{morphisms-lemma-base-change-proper} 和
\ref{morphisms-lemma-composition-proper}）。第一个断言蕴含 $X$ 在 $k$ 上分离，
因为 $X_{k'}$ 分离（《态射》引理
\ref{morphisms-lemma-image-universally-closed-separated}）。第二个断言由《态射》
引理 \ref{morphisms-lemma-image-proper-is-proper} 蕴含
$X \to \Spec(k)$ 固有。
\end{proof}

\begin{lemma}
\label{varieties-lemma-projective-after-field-extension}
设 $k$ 为域，设 $X$ 为 $k$ 上的概形。如果 $X_K$ 在 $K$ 上射影，其中
$K/k$ 是某个域扩张，那么 $X$ 在 $k$ 上射影。
\end{lemma}

\begin{proof}
域 $k$ 上的概形在 $k$ 上射影，当且仅当它在 $k$ 上拟射影且固有。见
《态射》引理 \ref{morphisms-lemma-projective-is-quasi-projective-proper}。
因此，引理由引理 \ref{varieties-lemma-quasi-projective-after-field-extension} 和
\ref{varieties-lemma-proper-after-field-extension} 推出。
\end{proof}


\section{切空间}
\label{varieties-section-tangent-spaces}

\noindent
本节利用取值于对偶数的点，定义概形态射在源概形一点处的切空间。

\begin{definition}
\label{varieties-definition-dual-numbers}
对任意环 $R$，$R$ 上的{it 对偶数}是一个 $R$-代数，记作
$R[\epsilon]$。作为 $R$-模，它是自由的，基为 $1$、$\epsilon$；其
$R$-代数结构来自规定 $\epsilon^2 = 0$。
\end{definition}

\noindent
设 $f : X \to S$ 为概形态射。设 $x \in X$ 为一点，它的像为
$s = f(x)$，这是 $S$ 中的一点。考虑实线组成的交换图
\begin{equation}
\label{varieties-equation-tangent-space}
\vcenter{
\xymatrix{
\Spec(\kappa(x)) \ar[r] \ar[dr] \ar@/^1pc/[rr] &
\Spec(\kappa(x)[\epsilon]) \ar@{.>}[r] \ar[d]&
X \ar[d] \\
&
\Spec(\kappa(s)) \ar[r] &
S
}
}
\end{equation}
其中弯曲箭头是从 $\Spec(\kappa(x))$ 到 $X$ 的典范态射。

\begin{lemma}
\label{varieties-lemma-tangent-space}
使图 (\ref{varieties-equation-tangent-space}) 交换的虚线箭头之集合，具有典范的
$\kappa(x)$-向量空间结构。
\end{lemma}

\begin{proof}
置 $\kappa = \kappa(x)$。注意，在概形范畴中有推出
$$
\Spec(\kappa[\epsilon]) \amalg_{\Spec(\kappa)} \Spec(\kappa[\epsilon])
= \Spec(\kappa[\epsilon_1, \epsilon_2])
$$
其中 $\kappa[\epsilon_1, \epsilon_2]$ 是一个 $\kappa$-代数，以
$1, \epsilon_1, \epsilon_2$ 为基，并且
$\epsilon_1^2 = \epsilon_1\epsilon_2 = \epsilon_2^2 = 0$。这立即来自环的相应
结论，以及《概形》引理 \ref{schemes-lemma-morphism-from-spec-local-ring}
中从局部环的谱到概形之态射的描述。给定两个箭头
$\theta_1, \theta_2 : \Spec(\kappa[\epsilon]) \to X$，可以考虑态射
$$
\theta_1 + \theta_2 :
\Spec(\kappa[\epsilon]) \to
\Spec(\kappa[\epsilon_1, \epsilon_2])
\xrightarrow{\theta_1, \theta_2} X
$$
其中第一个箭头由 $\epsilon_i \mapsto \epsilon$ 给出。另一方面，给定
$\lambda \in \kappa$，存在 $\Spec(\kappa[\epsilon])$ 的一个自映射，它对应于
$\kappa$-代数 $\kappa[\epsilon]$ 的自同态，而该自同态将 $\epsilon$ 映为
$\lambda \epsilon$。用这个自映射预合成
$\theta : \Spec(\kappa[\epsilon]) \to X$，便得到 $\lambda \theta$。读者可以
通过验证适当概形交换图的存在性来验证向量空间公理。细节从略。（另一种证明
是把一切都用局部环表述，再在环映射层面验证向量空间公理。）
\end{proof}

\begin{definition}
\label{varieties-definition-tangent-space}
设 $f : X \to S$ 为概形态射，并设 $x \in X$。使图
(\ref{varieties-equation-tangent-space}) 交换的虚线箭头之集合，连同其典范的
$\kappa(x)$-向量空间结构，称为{it $X$ 在 $S$ 上于 $x$ 处的切空间}，
记作 $T_{X/S, x}$。这个空间的元素称为 $X/S$ 在 $x$ 处的{it 切向量}。
\end{definition}

\noindent
由于 $x \in X$ 处的切向量存在于概形论纤维 $X_s$ 中，而这个纤维来自
$f : X \to S$ 且位于 $s = f(x)$ 上，我们得到典范等同
\begin{equation}
\label{varieties-equation-tangent-space-fibre}
T_{X/S, x} = T_{X_s/s, x}
\end{equation}
这个涉及点函子的优美定义有如下代数描述；它提示我们把{it $X$ 在 $S$ 上于
$x$ 处的余切空间}定义为 $\kappa(x)$-向量空间
$$
T^*_{X/S, x} = \Omega_{X/S, x} \otimes_{\mathcal{O}_{X, x}} \kappa(x)
$$
其理由就在于，它典范地是 $\kappa(x)$ 上的对偶，而所对偶的是 $X$ 在 $S$
上于 $x$ 处的切空间。

\begin{lemma}
\label{varieties-lemma-tangent-space-cotangent-space}
设 $f : X \to S$ 为概形态射，并设 $x \in X$。存在典范同构
$$
T_{X/S, x} = \Hom_{\mathcal{O}_{X, x}}(\Omega_{X/S, x}, \kappa(x))
$$
它是 $\kappa(x)$ 上的向量空间同构。
\end{lemma}

\begin{proof}
置 $\kappa = \kappa(x)$。给定 $\theta \in T_{X/S, x}$，得到映射
$$
\theta^*\Omega_{X/S} \to
\Omega_{\Spec(\kappa[\epsilon])/\Spec(\kappa(s))} \to
\Omega_{\Spec(\kappa[\epsilon])/\Spec(\kappa)}
$$
取截面得到一个 $\mathcal{O}_{X, x}$-线性映射
$\xi_\theta : \Omega_{X/S, x} \to \kappa \text{d}\epsilon$，也就是引理公式
右侧的一个元素。为了证明 $\theta \mapsto \xi_\theta$ 是同构，可以把 $S$
替换为 $s$，并把 $X$ 替换为概形论纤维 $X_s$。事实上，公式两侧都只依赖于
概形论纤维；对 $T_{X/S, x}$ 这是清楚的，右侧则见《态射》引理
\ref{morphisms-lemma-base-change-differentials}。还可以把 $X$ 替换为
$\mathcal{O}_{X, x}$ 的谱，因为这既不改变 $T_{X/S, x}$（《概形》
引理 \ref{schemes-lemma-morphism-from-spec-local-ring}），也不改变
$\Omega_{X/S, x}$（《模》引理
\ref{modules-lemma-stalk-module-differentials}）。

\medskip\noindent
设 $(A, \mathfrak m, \kappa)$ 为域 $k$ 上的局部环。为完成证明，需要说明
任意 $A$-线性映射 $\xi : \Omega_{A/k} \to \kappa$ 都来自唯一的 $k$-代数映射
$\varphi : A \to \kappa[\epsilon]$，且后者与典范映射
$c : A \to \kappa$ 在模 $\epsilon$ 后一致。写成
$\varphi(a) = c(a) + D(a) \epsilon$，读者可见
$a \mapsto D(a)$ 是 $k$-导子。利用 $\Omega_{A/k}$ 的泛性质，可知每个
$D$ 对应唯一的 $\xi$，反之亦然。证明完成。
\end{proof}

\begin{lemma}
\label{varieties-lemma-tangent-space-rational-point}
设 $f : X \to S$ 为概形态射。设 $x \in X$，并设 $s = f(x) \in S$。
假设 $\kappa(x) = \kappa(s)$。那么存在典范同构
$$
\mathfrak m_x/(\mathfrak m_x^2 + \mathfrak m_s\mathcal{O}_{X, x})
=
\Omega_{X/S, x} \otimes_{\mathcal{O}_{X, x}} \kappa(x)
$$
和
$$
T_{X/S, x} =
\Hom_{\kappa(x)}(
\mathfrak m_x/(\mathfrak m_x^2 + \mathfrak m_s\mathcal{O}_{X, x}),
\kappa(x))
$$
更一般地，当 $\kappa(x)/\kappa(s)$ 是可分代数扩张时，结论仍成立。
\end{lemma}

\begin{proof}
第二个同构由第一个同构和引理
\ref{varieties-lemma-tangent-space-cotangent-space} 推出。对于第一个同构，可以把 $S$
替换为 $s$，并把 $X$ 替换为 $X_s$，见《态射》引理
\ref{morphisms-lemma-base-change-differentials}。还可以把 $X$ 替换为
$\mathcal{O}_{X, x}$ 的谱，见《模》引理
\ref{modules-lemma-stalk-module-differentials}。因此，只须证明如下代数事实：设
$(A, \mathfrak m, \kappa)$ 为域 $k$ 上的局部环，并且 $\kappa/k$ 是可分
代数扩张。那么典范映射
$$
\mathfrak m/\mathfrak m^2
\longrightarrow
\Omega_{A/k} \otimes \kappa
$$
是同构。注意，$\mathfrak m/\mathfrak m^2 = H_1(\NL_{\kappa/A})$。由《代数》
引理 \ref{algebra-lemma-exact-sequence-NL}，只须证明
$\Omega_{\kappa/k} = 0$ 和 $H_1(\NL_{\kappa/k}) = 0$。由于 $\kappa$ 是
$k$ 中各有限可分扩张的并，只须证明 $\kappa$ 是 $k$ 的有限可分扩张时的情形
（《代数》引理 \ref{algebra-lemma-colimits-NL}）。在这种情况下，环映射
$k \to \kappa$ 是 \'etale 的，因而 $\NL_{\kappa/k} = 0$（这基本上就是定义，
见《代数》第 \ref{algebra-section-etale} 节）。
\end{proof}

\begin{lemma}
\label{varieties-lemma-map-tangent-spaces}
设 $f : X \to Y$ 为基概形 $S$ 上的概形态射。设 $x \in X$，并置
$y = f(x)$。如果 $\kappa(y) = \kappa(x)$，那么 $f$ 诱导自然线性映射
$$
\text{d}f : T_{X/S, x} \longrightarrow T_{Y/S, y}
$$
经由引理 \ref{varieties-lemma-tangent-space-cotangent-space} 的等同，它对偶于线性映射
$\Omega_{Y/S, y} \otimes \kappa(y) \to \Omega_{X/S, x} \otimes \kappa(x)$。
\end{lemma}

\begin{proof}
从略。
\end{proof}

\begin{lemma}
\label{varieties-lemma-tangent-space-product}
设 $X$、$Y$ 为基 $S$ 上的概形。设 $x \in X$ 与 $y \in Y$ 有同一
像点 $s \in S$，且 $\kappa(s) = \kappa(x)$ 和
$\kappa(s) = \kappa(y)$。存在典范同构
$$
T_{X \times_S Y/S, (x, y)} = T_{X/S, x} \oplus T_{Y/S, y}
$$
从左到右的映射由投影 $X \times_S Y \to X$ 与
$X \times_S Y \to Y$ 在切空间上诱导的映射给出。从右到左的映射由
$1 \times y : X_s \to X_s \times_s Y_s$ 与
$x \times 1 : Y_s \to X_s \times_s Y_s$ 给出，其中使用等同
(\ref{varieties-equation-tangent-space-fibre}) 将切空间等同于纤维的切空间。
\end{lemma}

\begin{proof}
直和分解来自《态射》引理 \ref{morphisms-lemma-differential-product}，其中
使用引理 \ref{varieties-lemma-tangent-space-rational-point}。它与这些映射的相容性来自引理
\ref{varieties-lemma-map-tangent-spaces}。
\end{proof}

\begin{lemma}
\label{varieties-lemma-injective-tangent-spaces-unramified}
设 $f : X \to Y$ 为基概形 $S$ 上局部有限型的概形态射。设 $x \in X$，
置 $y = f(x)$，并假设 $\kappa(y) = \kappa(x)$。那么下列条件等价：
\begin{enumerate}
\item $\text{d}f : T_{X/S, x} \longrightarrow T_{Y/S, y}$ 是单射；
\item $f$ 在 $x$ 处非分歧。
\end{enumerate}
\end{lemma}

\begin{proof}
由《态射》引理 \ref{morphisms-lemma-permanence-finite-type}，态射 $f$ 局部
有限型。映射 $\text{d}f$ 是单射，当且仅当
$\Omega_{Y/S, y} \otimes \kappa(y) \to \Omega_{X/S, x} \otimes \kappa(x)$
是满射（引理 \ref{varieties-lemma-map-tangent-spaces}）。正合列
$f^*\Omega_{Y/S} \to \Omega_{X/S} \to \Omega_{X/Y} \to 0$
（《态射》引理 \ref{morphisms-lemma-triangle-differentials}）于是说明，发生这种
情况当且仅当 $\Omega_{X/Y, x} \otimes \kappa(x) = 0$。因此结论来自《态射》
引理 \ref{morphisms-lemma-unramified-at-point}。
\end{proof}


\section{一般有限态射}
\label{varieties-section-generically-finite}

\noindent
本节重新讨论《态射》第 \ref{morphisms-section-generically-finite} 节研究过的
概形一般有限态射。

\begin{lemma}
\label{varieties-lemma-quasi-finite-in-codim-1}
设 $f : X \to Y$ 局部有限型。设 $y \in Y$ 为一点，使
$\mathcal{O}_{Y, y}$ 是维数 $\leq 1$ 的 Noether 环。再假设下列条件之一成立：
\begin{enumerate}
\item 对每个泛点 $\eta$，若它属于 $X$ 的某个不可约分支，则域扩张
$\kappa(\eta)/\kappa(f(\eta))$ 有限（或为代数扩张）；
\item 对每个泛点 $\eta$，若它属于 $X$ 的某个不可约分支，且
$f(\eta) \leadsto y$，域扩张 $\kappa(\eta)/\kappa(f(\eta))$ 就有限
（或为代数扩张）；
\item $f$ 在 $X$ 的不可约分支的每个泛点处拟有限；
\item $Y$ 局部 Noether，且 $f$ 在 $X$ 的一个稠密点集上拟有限；
\item 此处添加更多条件。
\end{enumerate}
那么，$f$ 在 $X$ 中位于 $y$ 上方的每一点处都拟有限。
\end{lemma}

\begin{proof}
条件 (4) 蕴含 $X$ 局部 Noether（《态射》引理
\ref{morphisms-lemma-finite-type-noetherian}）。态射拟有限的点集是开集
（《态射》引理 \ref{morphisms-lemma-quasi-finite-points-open}）。局部
Noether 概形的稠密开集包含不可约分支的所有泛点，因而 (4) 蕴含 (3)。
由《态射》引理 \ref{morphisms-lemma-residue-field-quasi-finite}，条件 (3)
蕴含条件 (1)。条件 (1) 蕴含条件 (2)。所以只须在 (2) 成立时证明引理。

\medskip\noindent
假设 (2) 成立。回忆，$\Spec(\mathcal{O}_{Y, y})$ 是 $Y$ 中特化到 $y$
的点集，见《概形》引理 \ref{schemes-lemma-specialize-points}。结合《态射》引理
\ref{morphisms-lemma-base-change-quasi-finite}，这说明可以把 $Y$ 替换为
$\Spec(\mathcal{O}_{Y, y})$。因此可以假设 $Y = \Spec(B)$，其中 $B$ 是
维数 $\leq 1$ 的 Noether 局部环，而 $y$ 是闭点。

\medskip\noindent
设 $X = \bigcup X_i$ 是 $X$ 的不可约分支，并把它们看作约化闭子概形。
如果能证明每个纤维 $X_{i, y}$ 都是离散空间，那么
$X_y = \bigcup X_{i, y}$ 也离散；由《态射》引理
\ref{morphisms-lemma-quasi-finite-at-point-characterize}，可得 $X \to Y$
在 $X_y$ 的所有点处拟有限。因此可以假设 $X$ 是整概形。

\medskip\noindent
如果 $X \to Y$ 把泛点 $\eta$（它是 $X$ 的泛点）映到 $y$，那么 $X$ 是
$\kappa(y)$ 的一个有限扩张的谱，结论成立。假设 $X$ 把 $\eta$ 映到
一个极小素理想 $\mathfrak q$（它属于 $B$）所对应的点，并且该点不同于
$\mathfrak m_B$。我们得到分解
$X \to \Spec(B/\mathfrak q) \to \Spec(B)$。设 $x \in X$ 为位于 $y$
上方的点。由维数公式（《态射》引理
\ref{morphisms-lemma-dimension-formula}），有
$$
\dim(\mathcal{O}_{X, x}) \leq \dim(B/\mathfrak q) +
\text{trdeg}_{\kappa(\mathfrak q)}(R(X)) - \text{trdeg}_{\kappa(y)} \kappa(x)
$$
我们知道 $\dim(B/\mathfrak q) = 1$；$X$ 的泛点不等于 $x$ 且特化到 $x$；
并且 $R(X)$ 是 $\kappa(\mathfrak q)$ 的代数扩张。因此得到
$$
1 \leq 1 - \text{trdeg}_{\kappa(y)} \kappa(x)
$$
于是，每一点 $x$（属于 $X_y$）都在 $X_y$ 中闭，这是由《态射》引理
\ref{morphisms-lemma-algebraic-residue-field-extension-closed-point-fibre} 得到的；
再由《态射》引理 \ref{morphisms-lemma-quasi-finite-at-point-characterize}，
$X \to Y$ 在每一点 $x$（属于 $X_y$）处拟有限（该引理也蕴含 $X_y$
是离散拓扑空间）。
\end{proof}

\begin{lemma}
\label{varieties-lemma-finite-in-codim-1}
设 $f : X \to Y$ 为固有态射。设 $y \in Y$ 为一点，使
$\mathcal{O}_{Y, y}$ 是维数 $\leq 1$ 的 Noether 环。再假设下列条件之一成立：
\begin{enumerate}
\item 对每个泛点 $\eta$，若它属于 $X$ 的某个不可约分支，则域扩张
$\kappa(\eta)/\kappa(f(\eta))$ 有限（或为代数扩张）；
\item 对每个泛点 $\eta$，若它属于 $X$ 的某个不可约分支，且
$f(\eta) \leadsto y$，域扩张 $\kappa(\eta)/\kappa(f(\eta))$ 就有限
（或为代数扩张）；
\item $f$ 在 $X$ 的每个泛点处拟有限；
\item $Y$ 局部 Noether，且 $f$ 在 $X$ 的一个稠密点集上拟有限；
\item 此处添加更多条件。
\end{enumerate}
那么，存在开邻域 $V \subset Y$，它包含 $y$，并使得
$f^{-1}(V) \to V$ 有限。
\end{lemma}

\begin{proof}
由引理 \ref{varieties-lemma-quasi-finite-in-codim-1}，态射 $f$ 在纤维 $X_y$ 的每一点
处拟有限。因此 $X_y$ 是离散拓扑空间（《态射》引理
\ref{morphisms-lemma-quasi-finite-at-point-characterize}）。由于 $f$ 固有，纤维
$X_y$ 拟紧，也就是说，它是有限集。因此可以应用《概形上同调》引理
\ref{coherent-lemma-proper-finite-fibre-finite-in-neighbourhood} 得出结论。
\end{proof}

\begin{lemma}
\label{varieties-lemma-modification-normal-iso-over-codimension-1}
设 $X$ 为 Noether 概形。设 $f : Y \to X$ 为双有理固有态射，且 $Y$ 约化。
设 $U \subset X$ 为 $f$ 在其上为同构的最大开集。那么 $U$ 包含：
\begin{enumerate}
\item 每个余维为 $0$ 的点，只要它属于 $X$；
\item 每个 $x \in X$，只要它余维为 $1$（这是在 $X$ 上计算的），且 $\mathcal{O}_{X, x}$
是离散赋值环；
\item 每个 $x \in X$，只要 $Y \to X$ 在 $x$ 上的纤维有限，且
$\mathcal{O}_{X, x}$ 正规；
\item 每个 $x \in X$，只要 $f$ 在某个 $y \in Y$ 处拟有限，该点
位于 $x$ 上方，并且 $\mathcal{O}_{X, x}$ 正规。
\end{enumerate}
\end{lemma}

\begin{proof}
(1) 来自《态射》引理
\ref{morphisms-lemma-birational-isomorphism-over-dense-open}。(2) 来自 (3)
以及引理 \ref{varieties-lemma-finite-in-codim-1}（并使用有限态射有有限纤维这一事实）。

\medskip\noindent
(3) 来自 (4) 和《态射》引理 \ref{morphisms-lemma-finite-fibre}，但我们也给出
一个直接证明。设 $x \in X$ 如 (3) 所述。由《概形上同调》引理
\ref{coherent-lemma-proper-finite-fibre-finite-in-neighbourhood}，可以假设 $f$
有限。可以假设 $X$ 仿射。这把问题归约为 Noether 仿射概形之间的有限双有理
态射 $Y \to X$ 的情形，其中 $x \in X$ 且 $\mathcal{O}_{X, x}$ 是正规整环。
由于 $\mathcal{O}_{X, x}$ 是整环且 $X$ Noether，可以把 $X$ 替换为 $x$
的某个仿射开集，使之整。于是，由于 $Y \to X$ 双有理且 $Y$ 约化，可知
$Y$ 整。写成 $X = \Spec(A)$ 和 $Y = \Spec(B)$，可知 $A \subset B$ 是
分式域相同的整环之间的有限包含。如果 $\mathfrak p \subset A$ 是对应于
$x$ 的素理想，那么 $A_\mathfrak p$ 正规蕴含
$A_\mathfrak p \subset B_\mathfrak p$ 为等式。由于 $B$ 是有限 $A$-模，
可知存在 $a \in A$、$a \not \in \mathfrak p$，使 $A_a \to B_a$ 是同构。

\medskip\noindent
设 $x \in X$ 与 $y \in Y$ 如 (4) 所述。用仿射开邻域替换 $X$ 后，
可以假设 $X = \Spec(A)$ 且 $A \subset \mathcal{O}_{X, x}$，见《性质》引理
\ref{properties-lemma-ring-affine-open-injective-local-ring}。于是 $A$ 是整环，
从而 $X$ 整。由于 $f$ 双有理且 $Y$ 约化，$Y$ 也整。考虑环映射
$\mathcal{O}_{X, x} \to \mathcal{O}_{Y, y}$。这是本质有限型的环映射，剩余
域扩张有限，并且
$\dim(\mathcal{O}_{Y, y}/\mathfrak m_x\mathcal{O}_{Y, y}) = 0$（要看出这一点，
可逐项追溯《态射》定义 \ref{morphisms-definition-quasi-finite} 与《代数》定义
\ref{algebra-definition-quasi-finite} 中拟有限映射的定义）。由《代数》引理
\ref{algebra-lemma-essentially-finite-type-fibre-dim-zero}，
$\mathcal{O}_{Y, y}$ 是某个有限 $\mathcal{O}_{X, x}$-代数 $B$ 的局部化。
当然，可以用 $B$ 的像替换 $B$，该像取在 $\mathcal{O}_{Y, y}$ 中；并假设 $B$
是与 $\mathcal{O}_{Y, y}$ 有相同分式域的整环。那么
$\mathcal{O}_{X, x} \subset B$ 与之有相同分式域，因为 $f$ 双有理。由于
$\mathcal{O}_{X, x}$ 正规，可得 $\mathcal{O}_{X, x} = B$（因为有限蕴含整），
特别地，$\mathcal{O}_{X, x} = \mathcal{O}_{Y, y}$。由《态射》引理
\ref{morphisms-lemma-morphism-defined-local-ring}，缩小 $X$ 后，可以假设存在
截面 $X \to Y$，它是 $f$ 的截面，把 $x$ 映到 $y$，并在局部环上诱导给定同构。
由于 $X \to Y$ 是闭的（由《概形》引理
\ref{schemes-lemma-section-immersion}），并且必然把 $X$ 的泛点映到 $Y$
的泛点，可知 $X \to Y$ 的像就是 $Y$。于是 $Y = X$，所需结论得证。
\end{proof}


\section{Noether 正规化的变式}
\label{varieties-section-noether-normalization}

\noindent
Noether 正规化定理断言：若 $k$ 是域，$A$ 是有限型 $k$-代数，且维数为
$d$，那么存在有限单射 $k$-代数同态
$k[x_1, \ldots, x_d] \to A$。见《代数》引理
\ref{algebra-lemma-Noether-normalization}。在几何上，这意味着存在有限满态射
$\Spec(A) \to \mathbf{A}^d_k$，它定义在 $\Spec(k)$ 上。

\begin{lemma}
\label{varieties-lemma-noether-normalization}
设 $f : X \to S$ 为概形态射。设 $x \in X$，其像为 $s \in S$。设
$V \subset S$ 是 $s$ 的仿射开邻域。若 $f$ 局部有限型且
$\dim_x(X_s) = d$，那么存在仿射开集 $U \subset X$，使得
$x \in U$ 且 $f(U) \subset V$，并且 $f|_U : U \to V$ 有分解
$$
U \xrightarrow{\pi} \mathbf{A}^d_V \to V
$$
其中 $\pi$ 拟有限。
\end{lemma}

\begin{proof}
这来自《代数》引理
\ref{algebra-lemma-quasi-finite-over-polynomial-algebra}。
\end{proof}

\begin{lemma}
\label{varieties-lemma-noether-normalization-affine}
设 $f : X \to S$ 为仿射概形之间的有限型态射。设 $s \in S$。若
$\dim(X_s) = d$，那么 $f$ 存在分解
$$
X \xrightarrow{\pi} \mathbf{A}^d_S \to S
$$
使得态射 $\pi_s : X_s \to \mathbf{A}^d_{\kappa(s)}$（即在 $s$ 上的纤维态射）有限。
\end{lemma}

\begin{proof}
写成 $S = \Spec(A)$ 与 $X = \Spec(B)$，并设 $A \to B$ 是对应于
$f$ 的环映射。设 $\mathfrak p \subset A$ 是对应于 $s$ 的素理想。可以选取满射
$A[x_1, \ldots, x_r] \to B$。由《代数》引理
\ref{algebra-lemma-Noether-normalization}，存在元
$y_1, \ldots, y_d \in A$，它们位于一个 $\mathbf{Z}$-子代数中；该子代数是在
$A$ 中由 $x_1, \ldots, x_r$ 生成的，并且使 $A$-代数同态
$A[t_1, \ldots, t_d] \to B$（把 $t_i$ 送到 $y_i$）诱导有限
$\kappa(\mathfrak p)$-代数同态
$\kappa(\mathfrak p)[t_1, \ldots, t_d] \to B \otimes_A \kappa(\mathfrak p)$。
这证明了引理。
\end{proof}

\begin{lemma}
\label{varieties-lemma-geometric-structure-unramified}
设 $f : X \to S$ 为概形态射。设 $x \in X$。设
$V = \Spec(A)$ 是 $f(x)$ 的仿射开邻域，并且位于 $S$ 中。若 $f$ 在 $x$
处非分歧，那么存在仿射开集 $U \subset X$，使 $x \in U$ 且
$f(U) \subset V$，并且存在交换图
$$
\xymatrix{
X \ar[d] & U \ar[l] \ar[rd] \ar[r]^-j &
\Spec(A[t]_{g'}/(g)) \ar[d] \ar[r] &
\Spec(A[t]) = \mathbf{A}^1_V \ar[ld] \\
Y & & V \ar[ll]
}
$$
其中 $j$ 是浸入，$g \in A[t]$ 是首一多项式，而 $g'$ 是 $g$ 对 $t$
的导数。若 $f$ 在 $x$ 处是 \'etale 的，那么可选取该图，使 $j$
是开浸入。
\end{lemma}

\begin{proof}
非分歧情形是《代数》命题
\ref{algebra-proposition-unramified-locally-standard} 的几何翻译。在 \'etale
情形，这是《代数》命题
\ref{algebra-proposition-etale-locally-standard} 的几何翻译；或者等价地，它来自
《态射》引理 \ref{morphisms-lemma-etale-locally-standard-etale}，尽管两处的表述略有不同。
\end{proof}

\begin{lemma}
\label{varieties-lemma-unramfied-over-affine}
设 $f : X \to S$ 为仿射概形之间的有限型态射。设 $x \in X$，其像为
$s \in S$。设
$$
r =
\dim_{\kappa(x)} \Omega_{X/S, x} \otimes_{\mathcal{O}_{X, x}} \kappa(x) =
\dim_{\kappa(x)} \Omega_{X_s/s, x} \otimes_{\mathcal{O}_{X_s, x}} \kappa(x) =
\dim_{\kappa(x)} T_{X/S, x}
$$
那么 $f$ 存在分解
$$
X \xrightarrow{\pi} \mathbf{A}^r_S \to S
$$
使得 $\pi$ 在 $x$ 处非分歧。
\end{lemma}

\begin{proof}
由《态射》引理 \ref{morphisms-lemma-finite-type-differentials}，第一个维数是有限的。
第一个等式来自《态射》引理
\ref{morphisms-lemma-base-change-differentials}：$\Omega_{X/S}$ 到纤维上的限制就是微分模。
最后一个等式来自引理 \ref{varieties-lemma-tangent-space-cotangent-space}。因此该表述有意义。

\medskip\noindent
为证明引理，写成 $S = \Spec(A)$ 与 $X = \Spec(B)$，并设 $A \to B$
是对应于 $f$ 的环映射。设 $\mathfrak q \subset B$ 是对应于 $x$ 的素理想。
选取 $A$-代数的满射 $A[x_1, \ldots, x_t] \to B$。由于
$\Omega_{B/A}$ 由 $\text{d}x_1, \ldots, \text{d}x_t$ 生成，它们在
$\Omega_{X/S, x} \otimes_{\mathcal{O}_{X, x}} \kappa(x)$ 中的像生成该
$\kappa(x)$-向量空间。重新编号后，可以假设
$\text{d}x_1, \ldots, \text{d}x_r$ 的像构成
$\Omega_{X/S, x} \otimes_{\mathcal{O}_{X, x}} \kappa(x)$ 的一组基。我们断言
$P = A[x_1, \ldots, x_r] \to B$ 在 $\mathfrak q$ 处非分歧。为此，只须证明
$\Omega_{B/P, \mathfrak q} = 0$（《代数》引理 \ref{algebra-lemma-unramified}）。
注意，$\Omega_{B/P}$ 是 $\Omega_{B/A}$ 对由
$\text{d}x_1, \ldots, \text{d}x_r$ 生成的子模之商。因此有
$\Omega_{B/P, \mathfrak q} \otimes_{B_\mathfrak q} \kappa(\mathfrak q) = 0$；这是由我们对
$x_1, \ldots, x_r$ 的选取得到的。
由 Nakayama 引理，更精确地说，由《代数》引理 \ref{algebra-lemma-NAK}
的 (2)（它可用，因为 $\Omega_{B/P}$ 是有限模，见上面的引用），可得
$\Omega_{B/P, \mathfrak q} = 0$。
\end{proof}

\begin{lemma}
\label{varieties-lemma-immersion-into-affine}
设 $f : X \to S$ 为概形态射。设 $x \in X$，其像为 $s \in S$。设
$V \subset S$ 是 $s$ 的仿射开邻域。若 $f$ 局部有限型且
$$
r =
\dim_{\kappa(x)} \Omega_{X/S, x} \otimes_{\mathcal{O}_{X, x}} \kappa(x) =
\dim_{\kappa(x)} \Omega_{X_s/s, x} \otimes_{\mathcal{O}_{X_s, x}} \kappa(x) =
\dim_{\kappa(x)} T_{X/S, x}
$$
那么存在下列两种情形：
\begin{enumerate}
\item 存在仿射开集 $U \subset X$，使 $x \in U$ 且 $f(U) \subset V$，并且
$f|_U$ 存在分解
$$
U \xrightarrow{j} \mathbf{A}^{r + 1}_V \to V
$$
其中 $j$ 是浸入；或者
\item 存在仿射开集 $U \subset X$，使 $x \in U$ 且 $f(U) \subset V$，并且
$f|_U$ 存在分解
$$
U \xrightarrow{j} D \to V
$$
其中 $j$ 是闭浸入，而 $D \to V$ 光滑且相对维数为 $r$。
\end{enumerate}
\end{lemma}

\begin{proof}
任取仿射开集 $U \subset X$，使 $x \in U$ 且 $f(U) \subset V$。对
$U \to V$ 应用引理 \ref{varieties-lemma-unramfied-over-affine}，得到该引理中的
$U \to \mathbf{A}^r_V \to V$。由引理 \ref{varieties-lemma-geometric-structure-unramified}，
得到分解
$$
U \xrightarrow{j} D \xrightarrow{j'} \mathbf{A}^{r + 1}_V
\xrightarrow{p} \mathbf{A}^r_V \to V
$$
其中 $j$ 与 $j'$ 是浸入，$p$ 是投影，而
$p \circ j'$ 是标准 \'etale 态射。特别地，这说明 (1) 与 (2) 成立。
\end{proof}


\section{纤维的维数}
\label{varieties-section-dimension-fibres}

\noindent
我们已经看到，有限型态射的纤维维数通常会跳升。本节讨论余维
$1$ 中不会发生这种现象，并更一般地讨论纤维维数能跳升多少。
下面列出一些相关结果：
\begin{enumerate}
\item 对有限型态射 $X \to S$，若点 $x \in X$ 满足
$\dim_x(X_{f(x)}) \leq d$，则所有这类点组成开集，见《代数》引理
\ref{algebra-lemma-dimension-fibres-bounded-open-upstairs} 和《态射》引理
\ref{morphisms-lemma-openness-bounded-dimension-fibres}。
\item 我们有维数公式，见《代数》引理
\ref{algebra-lemma-dimension-formula} 和《态射》引理
\ref{morphisms-lemma-dimension-formula}。
\item 对支配赋值环的整有限型概形，纤维维数为常数，见《代数》引理
\ref{algebra-lemma-finite-type-domain-over-valuation-ring-dim-fibres}。
\item 若 $X \to S$ 有限型，并且在 $X$ 的每个泛点处拟有限，那么
$X \to S$ 在余维 $1$ 中拟有限，见《代数》引理
\ref{algebra-lemma-finite-in-codim-1} 和引理
\ref{varieties-lemma-quasi-finite-in-codim-1}。
\end{enumerate}
上述最后一个结果可如下推广。

\begin{lemma}
\label{varieties-lemma-dimension-fibre-in-codim-1}
设 $f : X \to Y$ 局部有限型。设 $x \in X$ 是一点，其像为
$y \in Y$，并且 $\mathcal{O}_{Y, y}$ 是维数 $\leq 1$ 的 Noether 环。设
$d \geq 0$ 是整数，使得对每个泛点 $\eta$，若它属于 $X$ 的一个包含
$x$ 的不可约分支，则 $\dim_\eta(X_{f(\eta)}) = d$。那么
$\dim_x(X_y) = d$。
\end{lemma}

\begin{proof}
回忆，$\Spec(\mathcal{O}_{Y, y})$ 是 $Y$ 中特化到 $y$ 的点集，见《概形》引理
\ref{schemes-lemma-specialize-points}。因此可以替换 $Y$，把它取为
$\Spec(\mathcal{O}_{Y, y})$，并假设 $Y = \Spec(B)$，其中 $B$ 是维数 $\leq 1$ 的
Noether 局部环，而 $y$ 是闭点。还可以替换 $X$，把它取为 $x$ 的一个仿射邻域。

\medskip\noindent
设 $X = \bigcup X_i$ 是 $X$ 的不可约分支，并把它们视为约化闭子概形。
若能证明每个纤维 $X_{i, y}$ 的维数都为 $d$，那么
$X_y = \bigcup X_{i, y}$ 的维数也为 $d$。因此可以假设 $X$ 是整概形。

\medskip\noindent
若 $X \to Y$ 把泛点 $\eta$（它是 $X$ 的泛点）映到 $y$，那么 $X$ 是
$\kappa(y)$ 上的概形，结论由假设立即成立。假设 $X$ 把 $\eta$
映到点 $\xi \in Y$，它对应于极小素理想 $\mathfrak q$（它属于 $B$），
且它不同于 $\mathfrak m_B$。我们得到分解
$X \to \Spec(B/\mathfrak q) \to \Spec(B)$。由维数公式（《态射》引理
\ref{morphisms-lemma-dimension-formula}），有
$$
\dim(\mathcal{O}_{X, x}) + \text{trdeg}_{\kappa(y)} \kappa(x) \leq
\dim(B/\mathfrak q) + \text{trdeg}_{\kappa(\mathfrak q)}(R(X))
$$
我们有 $\dim(B/\mathfrak q) = 1$。又有
$\text{trdeg}_{\kappa(\mathfrak q)}(R(X)) = d$，这来自假设
$\dim_\eta(X_\xi) = d$，见《态射》引理
\ref{morphisms-lemma-dimension-fibre-at-a-point}。由于
$\mathcal{O}_{X, x} \to \mathcal{O}_{X_s, x}$ 有核（因为
$\eta \mapsto \xi \not = y$），并且 $\mathcal{O}_{X, x}$ 是 Noether 整环，可知
$\dim(\mathcal{O}_{X, x}) > \dim(\mathcal{O}_{X_y, x})$。从而
$$
\dim_x(X_s) =
\dim(\mathcal{O}_{X_s, x}) + \text{trdeg}_{\kappa(y)} \kappa(x) \leq d
$$
（《态射》引理 \ref{morphisms-lemma-dimension-fibre-at-a-point}）。另一方面，由《态射》引理
\ref{morphisms-lemma-openness-bounded-dimension-fibres}，有
$\dim_x(X_s) \geq \dim_\eta(X_{f(\eta)}) = d$。
\end{proof}

\begin{lemma}
\label{varieties-lemma-dominate-valuation-ring-dimension-fibres}
设 $f : X \to \Spec(R)$ 是从不可约概形到赋值环之谱的态射。若 $f$
局部有限型且满射，那么特殊纤维等维，且其维数等于一般纤维的维数。
\end{lemma}

\begin{proof}
可以用其约化替换 $X$，因为这不改变 $X$ 或特殊纤维的维数。
于是 $X$ 整，引理来自《代数》引理
\ref{algebra-lemma-finite-type-domain-over-valuation-ring-dim-fibres}。
\end{proof}

\noindent
下一个引理推广引理 \ref{varieties-lemma-dimension-fibre-in-codim-1}。

\begin{lemma}
\label{varieties-lemma-dimension-fibre-in-higher-codimension}
设 $f : X \to Y$ 局部有限型。设 $x \in X$ 是一点，其像为
$y \in Y$，并且 $\mathcal{O}_{Y, y}$ 是 Noether 环。设 $d \geq 0$
是整数，使得对每个泛点 $\eta$，若它属于 $X$ 的一个包含 $x$ 的不可约分支，则有
$f(\eta) \not = y$ 且 $\dim_\eta(X_{f(\eta)}) = d$。那么
$\dim_x(X_y) \leq d + \dim(\mathcal{O}_{Y, y}) - 1$。
\end{lemma}

\begin{proof}
与引理 \ref{varieties-lemma-dimension-fibre-in-codim-1} 的证明完全一样，我们把问题归约到
$X = \Spec(A)$ 的情形，其中 $A$ 是整环，而 $Y = \Spec(B)$，其中
$B$ 是局部 Noether 环，它的极大理想对应于 $y$。把 $B$
替换为 $B/\Ker(B \to A)$ 后，可以假设 $B$ 是整环且
$B \subset A$。然后使用维数公式（《态射》引理
\ref{morphisms-lemma-dimension-formula}）得到
$$
\dim(\mathcal{O}_{X, x}) + \text{trdeg}_{\kappa(y)} \kappa(x) \leq
\dim(B) + \text{trdeg}_B(A)
$$
我们有 $\text{trdeg}_B(A) = d$，这是因为假设 $\dim_\eta(X_\xi) = d$，见《态射》引理
\ref{morphisms-lemma-dimension-fibre-at-a-point}。由于
$\mathcal{O}_{X, x} \to \mathcal{O}_{X_y, x}$ 有核（因为 $f(\eta) \not = y$），
并且 $\mathcal{O}_{X, x}$ 是 Noether 整环，可知
$\dim(\mathcal{O}_{X, x}) > \dim(\mathcal{O}_{X_y, x})$。从而
$$
\dim_x(X_y) =
\dim(\mathcal{O}_{X_y, x}) + \text{trdeg}_{\kappa(y)} \kappa(x)
< \dim(B) + d
$$
（等式来自《态射》引理 \ref{morphisms-lemma-dimension-fibre-at-a-point}），
这就证明了所需结论。
\end{proof}


\section{代数概形}
\label{varieties-section-algebraic-schemes}

\noindent
下列定义取自 \cite[I Definition 6.4.1]{EGA}。

\begin{definition}
\label{varieties-definition-algebraic-scheme}
设 $k$ 为域。称 $k$ 上的概形 $X$ 为{it 代数 $k$-概形}，如果结构态射
$X \to \Spec(k)$ 有限型。称 $k$ 上的概形 $X$ 为{it 局部代数 $k$-概形}，
如果结构态射 $X \to \Spec(k)$ 局部有限型。
\end{definition}

\noindent
注意，每个（局部）代数 $k$-概形都是（局部）Noether 概形，见《态射》引理
\ref{morphisms-lemma-finite-type-noetherian}。代数 $k$-概形的范畴具有所有积和
纤维积（这不同于 $k$ 上簇的范畴）。局部代数 $k$-概形的范畴也一样。

\begin{lemma}
\label{varieties-lemma-algebraic-scheme-dim-0}
设 $k$ 为域。设 $X$ 是局部代数 $k$-概形，且维数为 $0$。那么 $X$ 是若干
局部 Artin $k$-代数 $A$ 的谱的不交并，并且 $\dim_k(A) < \infty$。若 $X$
是代数 $k$-概形，且维数为 $0$，那么此外 $X$ 是仿射的，并且态射
$X \to \Spec(k)$ 是有限态射。
\end{lemma}

\begin{proof}
设 $X$ 是局部代数 $k$-概形，且维数为 $0$。设
$U = \Spec(A) \subset X$ 是仿射开子概形。由于 $\dim(X) = 0$，可知
$\dim(A) = 0$。由 Noether 正规化（见《代数》引理
\ref{algebra-lemma-Noether-normalization}），存在有限单射 $k \to A$，即
$\dim_k(A) < \infty$。因此 $A$ 是 Artin 环，见《代数》引理
\ref{algebra-lemma-finite-dimensional-algebra}。这蕴涵
$A = A_1 \times \ldots \times A_r$ 是有限多个 Artin 局部环的积，见《代数》引理
\ref{algebra-lemma-artinian-finite-length}。当然，$\dim_k(A_i) < \infty$
对每个 $i$ 都成立，因为这些维数之和等于 $\dim_k(A)$。

\medskip\noindent
以上论证说明，$X$ 有一个开覆盖，其成员都是有限离散拓扑空间。因此 $X$
是离散拓扑空间。于是 $X$ 同构于其连通分支的不交并，而每个连通分支都是单点。
单点概形是仿射的，故由上一段的结果，每个这样的单点都是某个局部 Artin
$k$-代数 $A$ 的谱，并且 $\dim_k(A) < \infty$。

\medskip\noindent
最后，若 $X$ 是代数 $k$-概形且维数为 $0$，则 $X$ 拟紧，因而是有限不交并
$X = \Spec(A_1) \amalg \ldots \amalg \Spec(A_r)$，所以是仿射的（见《概形》引理
\ref{schemes-lemma-disjoint-union-affines}）；而证明的第一段已经说明
$X \to \Spec(k)$ 是有限态射。
\end{proof}

\noindent
下列引理汇集了局部代数概形维数论的一些陈述。

\begin{lemma}
\label{varieties-lemma-dimension-locally-algebraic}
设 $k$ 为域。设 $X$ 是局部代数 $k$-概形。
\begin{enumerate}
\item
\label{varieties-item-catenary}
$X$ 的拓扑空间是链状的（《拓扑》定义
\ref{topology-definition-catenary}）。
\item
\label{varieties-item-dimension-at-closed-point}
若 $x \in X$ 是闭点，则 $\dim_x(X) = \dim(\mathcal{O}_{X, x})$。
\item
\label{varieties-item-dimension-irreducible}
若 $X$ 不可约，则有 $\dim(X) = \dim(U)$，其中 $U \subset X$
为任意非空开集；并且有 $\dim(X) = \dim_x(X)$，其中 $x \in X$ 任意。
\item
\label{varieties-item-irreducible-maximal-chains}
若 $X$ 不可约，则不可约闭子集的任意链都可以延拓为极大链，并且不可约闭子集的
每条极大链的长度都等于 $\dim(X)$。
\item
\label{varieties-item-dimension-irreducibles-passing-through}
对 $x \in X$，有
$\dim_x(X) = \max \dim(Z) = \min \dim(\mathcal{O}_{X, x'})$，
其中极大值取遍不可约分支 $Z \subset X$，这些分支包含 $x$；极小值取遍满足
$x \leadsto x'$ 且 $x'$ 在 $X$ 中闭的特化。
\item
\label{varieties-item-dimension-irreducible-trdeg}
若 $X$ 不可约且泛点为 $x$，则
$\dim(X) = \text{trdeg}_k(\kappa(x))$。
\item
\label{varieties-item-immediate-specialization}
若 $x \leadsto x'$ 是 $X$ 中点的直接特化，则
$\text{trdeg}_k(\kappa(x)) = \text{trdeg}_k(\kappa(x')) + 1$。
\item
\label{varieties-item-dimension-sup-trdeg}
$X$ 的维数等于数 $\text{trdeg}_k(\kappa(x))$ 的上确界，其中 $x$
取遍 $X$ 的不可约分支的泛点。
\item
\label{varieties-item-specialization}
若 $x \leadsto x'$ 是 $X$ 中点的非平凡特化，则
\begin{enumerate}
\item $\dim_x(X) \leq \dim_{x'}(X)$，
\item $\dim(\mathcal{O}_{X, x}) < \dim(\mathcal{O}_{X, x'})$，
\item $\text{trdeg}_k(\kappa(x)) > \text{trdeg}_k(\kappa(x'))$，并且
\item 任意非平凡特化的极大链
$x = x_0 \leadsto x_1 \leadsto \ldots \leadsto x_n = x'$ 的长度为
$n = \text{trdeg}_k(\kappa(x)) - \text{trdeg}_k(\kappa(x'))$。
\end{enumerate}
\item
\label{varieties-item-dimension-formula}
对 $x \in X$，有
$\dim_x(X) = \text{trdeg}_k(\kappa(x)) + \dim(\mathcal{O}_{X, x})$。
\item
\label{varieties-item-immediate-specialization-local-ring}
若 $x \leadsto x'$ 是 $X$ 中点的直接特化，并且 $X$ 不可约或等维，则
$\dim(\mathcal{O}_{X, x'}) = \dim(\mathcal{O}_{X, x}) + 1$。
\end{enumerate}
\end{lemma}

\begin{proof}
我们不依赖前面证明的更一般结果，而把这些陈述归约为有限型 $k$-代数的相应陈述，
并引用交换代数章中的结果。

\medskip\noindent
(\ref{varieties-item-catenary}) 的证明。由《拓扑》引理
\ref{topology-lemma-catenary}，这是关于 $X$ 的局部问题。因此可以假设
$X = \Spec(A)$，其中 $A$ 是有限型 $k$-代数。我们必须证明 $A$ 是链环
（《代数》引理 \ref{algebra-lemma-catenary}）。利用《代数》引理
\ref{algebra-lemma-quotient-catenary}，可以归约到 $k[x_1, \ldots, x_n]$，
然后应用《代数》引理 \ref{algebra-lemma-dimension-height-polynomial-ring}。
另一种证明是：$k$ 显然是 Cohen--Macaulay 环，而 Cohen--Macaulay 环是普遍链环
（《代数》引理 \ref{algebra-lemma-CM-ring-catenary}）。

\medskip\noindent
(\ref{varieties-item-dimension-at-closed-point}) 的证明。选取仿射邻域
$U = \Spec(A)$，它包含 $x$。则 $\dim_x(X) = \dim_x(U)$。因此归约到仿射情形，这正是
《代数》引理 \ref{algebra-lemma-dimension-closed-point-finite-type-field}。

\medskip\noindent
(\ref{varieties-item-dimension-irreducible}) 的证明。只需证明任意两个非空仿射开集
$U, U' \subset X$ 具有相同维数（不可约子集的任意有限链都与某个仿射开集相交）。
选取闭点 $x$（它是 $X$ 的闭点），使得 $x \in U \cap U'$。这是可能的，因为 $X$ 不可约，
故 $U \cap U'$ 非空；又由引理 \ref{varieties-lemma-locally-finite-type-Jacobson}，
$X$ 是 Jacobson 概形，所以其中存在这样的闭点。由《代数》引理
\ref{algebra-lemma-dimension-spell-it-out}，有
$\dim(U) = \dim(\mathcal{O}_{X, x}) = \dim(U')$（严格地说，应用该引理前必须
用其约化替换 $X$）。

\medskip\noindent
(\ref{varieties-item-irreducible-maximal-chains}) 的证明。给定不可约闭子集的一条链，
可以找到与其中最小者相交的仿射开集 $U \subset X$。因此，结论来自《代数》引理
\ref{algebra-lemma-dimension-spell-it-out} 以及我们在
(\ref{varieties-item-dimension-irreducible}) 中看到的 $\dim(U) = \dim(X)$。

\medskip\noindent
(\ref{varieties-item-dimension-irreducibles-passing-through}) 的证明。选取仿射邻域
$U = \Spec(A)$，它包含 $x$。则 $\dim_x(X) = \dim_x(U)$。规则 $Z \mapsto Z \cap U$
给出 $X$ 的不可约分支（它们经过 $x$）与 $U$ 的不可约分支（它们经过 $x$）之间的双射。
并且由 (\ref{varieties-item-dimension-irreducible})，有
$\dim(Z \cap U) = \dim(Z)$，这对这样的 $Z$ 成立。因此结论来自《代数》引理
\ref{algebra-lemma-dimension-at-a-point-finite-type-over-field}。

\medskip\noindent
(\ref{varieties-item-dimension-irreducible-trdeg}) 的证明。由
(\ref{varieties-item-dimension-irreducible})，可以归约到 $X = \Spec(A)$ 为仿射的情形。
此时，把《代数》引理 \ref{algebra-lemma-dimension-prime-polynomial-ring}
应用于 $A_{red}$ 即得结论。

\medskip\noindent
(\ref{varieties-item-immediate-specialization}) 的证明。设
$Z = \overline{\{x\}} \supset Z' = \overline{\{x'\}}$。
由 (\ref{varieties-item-irreducible-maximal-chains})，$Z \supset Z'$ 是 $Z$
中不可约闭子概形的一条极大链的开端，因而 $\dim(Z) = \dim(Z') + 1$。
再由 (\ref{varieties-item-dimension-irreducible-trdeg}) 即得结论。

\medskip\noindent
(\ref{varieties-item-dimension-sup-trdeg}) 的证明。一个简单的拓扑论证说明
$\dim(X) = \sup \dim(Z)$，其中上确界取遍 $X$ 的不可约分支（提示：使用《拓扑》引理
\ref{topology-lemma-irreducible}）。因此结论来自
(\ref{varieties-item-dimension-irreducible-trdeg})。

\medskip\noindent
(\ref{varieties-item-specialization}) 的证明。(a) 来自如下事实：每个开集
$U \subset X$，若它包含 $x'$，则也包含 $x$。(b) 来自如下事实：$\mathcal{O}_{X, x}$ 是
$\mathcal{O}_{X, x'}$ 的一个局部化，因而 $\mathcal{O}_{X, x}$ 中的每条素理想链
都对应于 $\mathcal{O}_{X, x'}$ 中的一条素理想链，而且可在末端添加
$\mathfrak m_{x'}$ 来延长它。(c) 和 (d) 都形式地来自
(\ref{varieties-item-immediate-specialization})。

\medskip\noindent
(\ref{varieties-item-dimension-formula}) 的证明。选取仿射邻域
$U = \Spec(A)$，它包含 $x$。则 $\dim_x(X) = \dim_x(U)$。因此归约到仿射情形，这正是
《代数》引理 \ref{algebra-lemma-dimension-at-a-point-finite-type-field}。

\medskip\noindent
(\ref{varieties-item-immediate-specialization-local-ring}) 的证明。若 $X$ 等维
（《拓扑》定义 \ref{topology-definition-equidimensional}），则 $\dim(X)$
等于 $X$ 的每个不可约分支的维数，从而由
(\ref{varieties-item-dimension-irreducibles-passing-through})，有
$\dim_x(X) = \dim(X) = \dim_{x'}(X)$。因此结论来自
(\ref{varieties-item-immediate-specialization})。
\end{proof}

\begin{lemma}
\label{varieties-lemma-dimension-fibres-locally-algebraic}
设 $k$ 为域。设 $f : X \to Y$ 是局部代数 $k$-概形之间的态射。
\begin{enumerate}
\item 对 $y \in Y$，纤维 $X_y$ 是 $\kappa(y)$ 上的局部代数概形，因而引理
\ref{varieties-lemma-dimension-locally-algebraic} 的所有结果都适用。
\item 假设 $X$ 不可约。令 $Z = \overline{f(X)}$，并令
$d = \dim(X) - \dim(Z)$。则
\begin{enumerate}
\item 有 $\dim_x(X_{f(x)}) \geq d$，对所有 $x \in X$ 都成立，
\item 点 $x \in X$ 中满足 $\dim_x(X_{f(x)}) = d$ 的那些点组成稠密开集，
\item 若 $\dim(\mathcal{O}_{Z, f(x)}) \geq 1$，则
$\dim_x(X_{f(x)}) \leq d + \dim(\mathcal{O}_{Z, f(x)}) - 1$，
\item 若 $\dim(\mathcal{O}_{Z, f(x)}) = 1$，则
$\dim_x(X_{f(x)}) = d$。
\end{enumerate}
\item 对 $x \in X$，若 $y = f(x)$，则有
$\dim_x(X_y) \geq \dim_x(X) - \dim_y(Y)$。
\end{enumerate}
\end{lemma}

\begin{proof}
由《态射》引理 \ref{morphisms-lemma-permanence-finite-type}，态射 $f$
局部有限型。因此基变换 $X_y \to \Spec(\kappa(y))$ 局部有限型。这证明了 (1)。
在证明的其余部分，我们将自由地对 $X$、$Y$ 和 $f$ 的纤维使用引理
\ref{varieties-lemma-dimension-locally-algebraic} 的结果。

\medskip\noindent
(2) 的证明。设 $\eta \in X$ 为泛点，并令 $\xi = f(\eta)$。则
$Z = \overline{\{\xi\}}$。因此
$$
d = \dim(X) - \dim(Z) =
\text{trdeg}_k \kappa(\eta) - \text{trdeg}_k \kappa(\xi) =
\text{trdeg}_{\kappa(\xi)} \kappa(\eta) = \dim_\eta(X_\xi)
$$
于是 (2)(a) 和 (2)(b) 来自《态射》引理
\ref{morphisms-lemma-openness-bounded-dimension-fibres}。(2)(c) 和 (2)(d)
来自引理 \ref{varieties-lemma-dimension-fibre-in-higher-codimension} 和
\ref{varieties-lemma-dimension-fibre-in-codim-1}。

\medskip\noindent
(3) 的证明。设 $x \in X$。设 $X' \subset X$ 是 $X$ 的一个不可约分支，
它经过 $x$ 且维数为 $\dim_x(X)$。于是 (2) 蕴涵
$\dim_x(X_y) \geq \dim(X') - \dim(Z')$，其中 $Z' \subset Y$
是 $X'$ 的像的闭包。这证明了 (3)。
\end{proof}

\begin{lemma}
\label{varieties-lemma-dimension-product-locally-algebraic}
\begin{slogan}
积的维数等于各因子的维数之和。
\end{slogan}
设 $k$ 为域。设 $X$、$Y$ 是局部代数 $k$-概形。
\begin{enumerate}
\item 若 $z \in X \times Y$ 位于 $(x, y)$ 上方，则
$\dim_z(X \times Y) = \dim_x(X) + \dim_y(Y)$。
\item 有 $\dim(X \times Y) = \dim(X) + \dim(Y)$。
\end{enumerate}
\end{lemma}

\begin{proof}
(1) 的证明。考虑结构态射的分解
$$
X \times Y \longrightarrow Y \longrightarrow \Spec(k)
$$
第一个态射 $p : X \times Y \to Y$ 是平坦态射，因为它是平坦态射
$X \to \Spec(k)$ 的基变换，见《态射》引理
\ref{morphisms-lemma-base-change-flat}。此外，由《态射》引理
\ref{morphisms-lemma-dimension-fibre-after-base-change}，有
$\dim_z(p^{-1}(y)) = \dim_x(X)$。因此由《态射》引理
\ref{morphisms-lemma-dimension-fibre-at-a-point-additive}，有
$\dim_z(X \times Y) = \dim_x(X) + \dim_y(Y)$。(2) 是 (1) 的直接推论。
\end{proof}


\section{完备局部环}
\label{varieties-section-complete-local-rings}

\noindent
这里给出域上概形的完备局部环的一些结果。

\begin{lemma}
\label{varieties-lemma-complete-local-ring-structure-as-algebra}
设 $k$ 为域。设 $X$ 是 $k$ 上的局部 Noether 概形。设 $x \in X$
是一点，其剩余域为 $\kappa$。存在同构
\begin{equation}
\label{varieties-equation-complete-local-ring}
\kappa[[x_1, \ldots, x_n]]/I \longrightarrow \mathcal{O}_{X, x}^\wedge
\end{equation}
它在剩余域上诱导恒等映射。一般而言，不能选取
(\ref{varieties-equation-complete-local-ring}) 使其成为 $k$-代数同构。然而，若扩张
$\kappa/k$ 可分，则可以选取 (\ref{varieties-equation-complete-local-ring})
使其成为 $k$-代数同构。
\end{lemma}

\begin{proof}
该同构的存在性是 Cohen 结构定理的直接推论\footnote{注意，若 $\kappa$
的特征为 $p$，则该定理只说明存在满射
$\Lambda[[x_1, \ldots, x_n]] \to \mathcal{O}_{X, x}^\wedge$，其中
$\Lambda$ 是 $\kappa$ 的 Cohen 环。但在这种情形下，该映射当然经过
$\Lambda/p\Lambda[[x_1, \ldots, x_n]]$ 分解，并且
$\Lambda/p\Lambda = \kappa$。}（《代数》定理
\ref{algebra-theorem-cohen-structure-theorem}）。

\medskip\noindent
设 $p$ 为奇素数，设 $k = \mathbf{F}_p(t)$，并令
$A = k[x, y]/(y^2 + x^p - t)$。则完备化 $A^\wedge$（即 $A$
在极大理想 $\mathfrak m = (y)$ 处的完备化）作为环同构于
$k(t^{1/p})[[z]]$，但不作为 $k$-代数同构。原因是 $A^\wedge$
不含有 $p$ 次幂为 $t$ 的元素（读者可以通过模 $y^2$ 计算看出这一点）。
这也说明 (\ref{varieties-equation-complete-local-ring}) 的任何同构都不可能是
$k$-代数同构。

\medskip\noindent
若 $\kappa/k$ 可分，则由《代数进阶》引理
\ref{more-algebra-lemma-lift-residue-field}，存在 $k$-代数同态
$\kappa \to \mathcal{O}_{X, x}^\wedge$，它在剩余域上诱导恒等映射。
设 $f_1, \ldots, f_n \in \mathfrak m_x$ 为一组生成元。考虑映射
$$
\kappa[[x_1, \ldots, x_n]] \longrightarrow \mathcal{O}_{X, x}^\wedge,\quad
x_i \longmapsto f_i
$$
由于两边关于 $(x_1, \ldots, x_n)$ 都完备（右边的完备性见《代数》引理
\ref{algebra-lemma-hathat-finitely-generated}），而且按构造，该映射模
$(x_1, \ldots, x_n)$ 是满射，所以由《代数》引理
\ref{algebra-lemma-completion-generalities}，该映射是满射。
\end{proof}

\begin{lemma}
\label{varieties-lemma-base-change-complete-local-ring}
设 $K/k$ 是域扩张。设 $X$ 是局部代数 $k$-概形。令 $Y = X_K$。
设 $y \in Y$ 是一点，其像为 $x \in X$。假设
$\dim(\mathcal{O}_{X, x}) = \dim(\mathcal{O}_{Y, y})$，并且
$\kappa(x)/k$ 可分。选取如 (\ref{varieties-equation-complete-local-ring}) 中的
$k$-代数同构
$$
\kappa(x)[[x_1, \ldots, x_n]]/(g_1, \ldots, g_m) \longrightarrow
\mathcal{O}_{X, x}^\wedge
$$
则存在如 (\ref{varieties-equation-complete-local-ring}) 中的 $K$-代数同构
$$
\kappa(y)[[x_1, \ldots, x_n]]/(g_1, \ldots, g_m) \longrightarrow
\mathcal{O}_{Y, y}^\wedge
$$
这里使用 $\kappa(x) \to \kappa(y)$，把 $g_j$ 看作 $\kappa(y)$
上的幂级数。
\end{lemma}

\begin{proof}
局部环映射 $\mathcal{O}_{X, x} \to \mathcal{O}_{Y, y}$ 诱导局部环映射
$\mathcal{O}_{X, x}^\wedge \to \mathcal{O}_{Y, y}^\wedge$。诱导映射
$$
\kappa(x) \to \kappa(x)[[x_1, \ldots, x_n]]/(g_1, \ldots, g_m)
\to \mathcal{O}_{X, x}^\wedge \to \mathcal{O}_{Y, y}^\wedge
$$
与到 $\kappa(y)$ 的投影复合后，就是典范同态
$\kappa(x) \to \kappa(y)$。由引理 \ref{varieties-lemma-change-fields-flat}，
剩余域 $\kappa(y)$ 是 $\kappa(x) \otimes_k K$ 在核 $\mathfrak p_0$ 处的局部化；
这里所指的同态为 $\kappa(x) \otimes_k K \to \kappa(y)$。
另一方面，由引理 \ref{varieties-lemma-change-fields-algebraic-unramified}，局部环
$(\kappa(x) \otimes_k K)_{\mathfrak p_0}$ 等于 $\kappa(y)$。因此映射
$$
\kappa(x) \otimes_k K \to \mathcal{O}_{Y, y}^\wedge
$$
典范地经过 $\kappa(y)$ 分解。我们得到交换图
$$
\xymatrix{
\kappa(y) \ar[rr] & & \mathcal{O}_{Y, y}^\wedge \\
\kappa(x) \ar[r] \ar[u] &
\kappa(x)[[x_1, \ldots, x_n]]/(g_1, \ldots, g_m) \ar[r] &
\mathcal{O}_{X, x}^\wedge \ar[u]
}
$$
设 $f_i \in \mathfrak m_x^\wedge \subset \mathcal{O}_{X, x}^\wedge$
为 $x_i$ 的像。注意，由于该映射是满射，有
$\mathfrak  m_x^\wedge = (f_1, \ldots, f_n)$。考虑映射
$$
\kappa(y)[[x_1, \ldots, x_n]] \longrightarrow \mathcal{O}_{Y, y}^\wedge,\quad
x_i \longmapsto f_i
$$
这里的 $f_i$ 实际上是指 $f_i$ 在 $\mathfrak m_y^\wedge$ 中的像。
由引理 \ref{varieties-lemma-change-fields-algebraic-unramified}，有
$\mathfrak m_x \mathcal{O}_{Y, y} = \mathfrak m_y$，所以右边关于
$(x_1, \ldots, x_n)$ 完备（使用《代数》引理
\ref{algebra-lemma-hathat-finitely-generated} 可知它是完备局部环）。
由于两边关于 $(x_1, \ldots, x_n)$ 都完备，而该映射模
$(x_1, \ldots, x_n)$ 是满射，故由《代数》引理
\ref{algebra-lemma-completion-generalities}，我们的映射是满射。
当然，幂级数 $g_1, \ldots, g_m$ 在此映射下映为零，因为它们在
$\mathcal{O}_{X, x}^\wedge$ 中已经映为零。因此有交换图
$$
\xymatrix{
\kappa(y)[[x_1, \ldots, x_n]]/(g_1, \ldots, g_m) \ar[r] &
\mathcal{O}_{Y, y}^\wedge \\
\kappa(x)[[x_1, \ldots, x_n]]/(g_1, \ldots, g_m) \ar[r] \ar[u] &
\mathcal{O}_{X, x}^\wedge \ar[u]
}
$$
还需要证明上方的水平箭头是同构。我们已经知道它是满射。由引理
\ref{varieties-lemma-change-fields-flat}，$\mathcal{O}_{X, x} \to \mathcal{O}_{Y, y}$
是平坦的，这蕴涵 $\mathcal{O}_{X, x}^\wedge \to \mathcal{O}_{Y, y}^\wedge$
是平坦的（《代数进阶》引理
\ref{more-algebra-lemma-flat-completion}）。因此可以应用《代数》引理
\ref{algebra-lemma-mod-injective}，其中
$R = \kappa(x)[[x_1, \ldots, x_n]]/(g_1, \ldots, g_m)$，
$S = \kappa(y)[[x_1, \ldots, x_n]]/(g_1, \ldots, g_m)$，
$M = \mathcal{O}_{Y, y}^\wedge$，并且 $N = S$，从而得出该映射是单射。
\end{proof}


\section{全局生成}
\label{varieties-section-global-generation}

\noindent
这里给出与准凝聚模全局生成有关的一些引理。

\begin{lemma}
\label{varieties-lemma-globally-generated-base-change}
设 $X \to \Spec(A)$ 是概形的态射。设 $A \subset A'$ 是忠实平坦环映射。
设 $\mathcal{F}$ 是准凝聚 $\mathcal{O}_X$-模。那么 $\mathcal{F}$ 全局生成，
当且仅当基变换 $\mathcal{F}_{A'}$ 全局生成。
\end{lemma}

\begin{proof}
更准确地说，令 $X_{A'} = X \times_{\Spec(A)} \Spec(A')$。令
$\mathcal{F}_{A'} = p^*\mathcal{F}$，其中 $p : X_{A'} \to X$ 是投影。
由《概形上同调》引理 \ref{coherent-lemma-flat-base-change-cohomology}，有
$H^0(X_{k'}, \mathcal{F}_{A'}) = H^0(X, \mathcal{F}) \otimes_A A'$。
因此，若 $s_i$（$i \in I$）是 $H^0(X, \mathcal{F})$ 作为 $A$-模的一组生成元，
则它们在 $H^0(X_{A'}, \mathcal{F}_{A'})$ 中的像是
$H^0(X_{A'}, \mathcal{F}_{A'})$ 作为 $A'$-模的一组生成元。
所以只需证明映射
$\alpha : \bigoplus_{i \in I} \mathcal{O}_X \to \mathcal{F}$，
$(f_i) \mapsto \sum f_i s_i$ 是满射，当且仅当 $p^*\alpha$ 是满射。
这可以在仿射开集 $U = \Spec(B)$ 上检验；它是 $X$ 的开集。此时
$\mathcal{F}|_U$ 对应于 $B$-模 $M$，而 $s_i|_U$ 对应于元素
$x_i \in M$。因此只需证明 $\bigoplus_{i \in I} B \to M$ 是满射，
当且仅当基变换
$\bigoplus_{i \in I} B \otimes_A A' \to M \otimes_A A'$ 是满射。
这成立是因为 $A \to A'$ 忠实平坦。
\end{proof}

\begin{lemma}
\label{varieties-lemma-very-ample-vanish-at-point}
设 $k$ 是无限域。设 $X$ 是 $k$ 上的有限型概形。设 $\mathcal{L}$
是 $X$ 上的甚丰沛可逆层。设 $n \geq 0$，并设
$x, x_1, \ldots, x_n \in X$ 是若干点，其中 $x$ 是 $k$-有理点，即
$\kappa(x) = k$，并且有 $x \not = x_i$，其中 $i = 1, \ldots, n$。
那么存在 $s \in H^0(X, \mathcal{L})$，它在 $x$ 处为零，但在 $x_i$
处不为零。
\end{lemma}

\begin{proof}
若 $n = 0$，结论显然，故假设 $n > 0$。由甚丰沛可逆层的定义，
引理立即归约到 $X = \mathbf{P}^r_k$ 的情形，其中某个 $r > 0$，并且
$\mathcal{L} = \mathcal{O}_X(1)$。写成
$\mathbf{P}^r_k = \text{Proj}(k[T_0, \ldots, T_r])$。令
$V = H^0(X, \mathcal{L}) = kT_0 \oplus \ldots \oplus kT_r$。
由于 $x$ 是 $k$-有理点，元素 $s \in V$ 中在 $x$ 处为零的那些元素组成余维为 $1$
的子空间 $W \subset V$，并且 $W$ 生成与 $x$ 对应的齐次素理想。
由于 $x_i \not = x$，相应的齐次素理想
$\mathfrak p_i \subset k[T_0, \ldots, T_r]$ 不包含 $W$。由于 $k$ 无限，
可知 $W \not = \bigcup W \cap \mathfrak q_i$，证明完成。
\end{proof}

\begin{lemma}
\label{varieties-lemma-generated-by-dim-plus-1-sections}
设 $k$ 是无限域。设 $X$ 是代数 $k$-概形。设 $\mathcal{L}$ 是可逆
$\mathcal{O}_X$-模。设 $V \to \Gamma(X, \mathcal{L})$ 是 $k$-向量空间的
线性映射，其像生成 $\mathcal{L}$。那么存在子空间 $W \subset V$，满足
$\dim_k(W) \leq \dim(X) + 1$，并且它生成 $\mathcal{L}$。
\end{lemma}

\begin{proof}
在整个证明中，我们将使用如下事实：对每个 $x \in X$，线性映射
$$
\psi_x : V \to \Gamma(X, \mathcal{L}) \to \mathcal{L}_x \to
\mathcal{L}_x \otimes_{\mathcal{O}_{X, x}} \kappa(x)
$$
非零。证明对 $\dim(X)$ 使用归纳法。

\medskip\noindent
基始情形是 $\dim(X) = 0$。此时 $X$ 只有有限多个点
$X = \{x_1, \ldots, x_n\}$（例如见引理
\ref{varieties-lemma-algebraic-scheme-dim-0}）。由于 $k$ 无限，存在向量
$v \in V$，使得 $\psi_{x_i}(v) \not = 0$ 对所有 $i$ 都成立。
于是 $W = k\cdot v$ 满足要求。

\medskip\noindent
假设 $\dim(X) > 0$。设 $X_i \subset X$ 是维数等于 $\dim(X)$ 的不可约分支。
由于 $X$ 是 Noether 概形，这样的分支只有有限多个。对每个 $i$，选取一点
$x_i \in X_i$。如上选取 $v \in V$，使得 $\psi_{x_i}(v) \not = 0$
对所有 $i$ 都成立。设 $Z \subset X$ 是 $v$ 在
$\Gamma(X, \mathcal{L})$ 中之像的零概形，见《除子》定义
\ref{divisors-definition-zero-scheme-s}。按构造，$\dim(Z) < \dim(X)$。
由归纳假设，可以找到 $W \subset V$，满足 $\dim(W) \leq \dim(X)$，
使得 $W$ 生成 $\mathcal{L}|_Z$。于是 $W + k\cdot v$ 生成 $\mathcal{L}$。
\end{proof}


\section{分离点与切向量}
\label{varieties-section-separating-points-tangent-vectors}

\noindent
本节就是下列结果。

\begin{lemma}
\label{varieties-lemma-separate-points-tangent-vectors}
设 $k$ 是代数闭域。设 $X$ 是固有 $k$-概形。设 $\mathcal{L}$ 是可逆
$\mathcal{O}_X$-模。设 $V \subset H^0(X, \mathcal{L})$ 是 $k$-向量子空间。
如果
\begin{enumerate}
\item 对每一对互异闭点 $x, y \in X$，存在截面 $s \in V$，它在 $x$
处为零而在 $y$ 处不为零；并且
\item 对每个闭点 $x \in X$ 和非零切向量 $\theta \in T_{X/k, x}$，
存在截面 $s \in V$，它在 $x$ 处为零，但经 $\theta$ 拉回后非零，
\end{enumerate}
那么 $\mathcal{L}$ 甚丰沛，并且典范态射
$\varphi_{\mathcal{L}, V} : X \to \mathbf{P}(V)$ 是闭浸入。
\end{lemma}

\begin{proof}
条件 (1) 尤其蕴涵 $V$ 的元素生成 $\mathcal{L}$（作为 $X$ 上的层）。因此由《构造》例
\ref{constructions-example-projective-space}，得到典范态射
$$
\varphi = \varphi_{\mathcal{L}, V} :
X
\longrightarrow
\mathbf{P}(V)
$$
由《态射》引理 \ref{morphisms-lemma-image-proper-scheme-closed}，态射
$\varphi$ 是固有的。由 (1)，映射 $\varphi$ 在闭点上是单射（略去计算）。
特别地，$\mathbf{P}(V)$ 的任意闭点上方的纤维都是单点（略去一个小细节）。
因此 $\varphi$ 是有限态射；例如可使用《概形上同调》引理
\ref{coherent-lemma-proper-finite-fibre-finite-in-neighbourhood}。
为完成证明，只需说明映射
$$
\varphi^\sharp :
\mathcal{O}_{\mathbf{P}(V)}
\longrightarrow
\varphi_*\mathcal{O}_X
$$
是满射。这可以在闭点的茎上检验。设 $x \in X$ 是闭点，其像为闭点
$p = \varphi(x) \in \mathbf{P}(V)$。由 (1)，有
$\varphi^{-1}(\{p\}) = \{x\}$；又因为 $\varphi$ 固有（因而是闭映射），
$\varphi^{-1}(U)$ 取遍 $x$ 的开邻域的一个基本系，其中 $U$ 取遍
$p$ 的开邻域的一个基本系。因此在 $p$ 处取茎，得到映射
$$
\varphi^\sharp_x :
\mathcal{O}_{\mathbf{P}(V), p}
\longrightarrow
\mathcal{O}_{X, x}
$$
特别地，$\mathcal{O}_{X, x}$ 是有限 $\mathcal{O}_{\mathbf{P}(V), p}$-模。
此外，$x$ 和 $p$ 的剩余域都等于 $k$（因为 $k$ 代数闭；使用 Hilbert
零点定理）。最后，条件 (2) 蕴涵映射
$$
T_{X/k, x} \longrightarrow T_{\mathbf{P}(V)/k, p}
$$
是单射，因为该映射的核中的任何非零 $\theta$ 都不可能满足 (2) 的结论。
就上述局部环映射而言，这意味着
$$
\mathfrak m_p/\mathfrak m_p^2 \longrightarrow
\mathfrak m_x/\mathfrak m_x^2
$$
是满射，见引理 \ref{varieties-lemma-tangent-space-rational-point}。现在应用《代数》引理
\ref{algebra-lemma-when-surjective-local} 即完成证明。
\end{proof}

\begin{lemma}
\label{varieties-lemma-variant-separate-points-tangent-vectors}
设 $k$ 是代数闭域。设 $X$ 是固有 $k$-概形。设 $\mathcal{L}$ 是可逆
$\mathcal{O}_X$-模。假设对每个闭子概形 $Z \subset X$，只要它的维数为 $0$、
次数为 $2$（相对于 $k$），映射
$$
H^0(X, \mathcal{L}) \longrightarrow H^0(Z, \mathcal{L}|_Z)
$$
就是满射。那么 $\mathcal{L}$ 在 $X$ 上相对于 $k$ 甚丰沛。
\end{lemma}

\begin{proof}
这是引理 \ref{varieties-lemma-separate-points-tangent-vectors} 的改述。事实上，给定互异闭点
$x, y \in X$，取 $Z = x \cup y$（视为闭子概形），便得到该引理的条件 (1)。
给定非零切向量 $\theta \in T_{X/k, x}$，态射
$\theta : \Spec(k[\epsilon]) \to X$ 是闭浸入。令 $Z = \Im(\theta)$，
便得到该引理的条件 (2)。
\end{proof}


\section{积的闭包}
\label{varieties-section-closure-of-products}

\noindent
以下给出若干关于闭包与积之间关系的结果。

\begin{lemma}
\label{varieties-lemma-closure-of-product}
设 $k$ 是域。设 $X$、$Y$ 是 $k$-概形，并设
$A \subset X$、$B \subset Y$ 是子集。令
$$
AB =
\{z \in X \times_k Y \mid \text{pr}_X(z) \in A, \ \text{pr}_Y(z) \in B\}
\subset X \times_k Y
$$
则在集合论意义下有
$$
\overline{A} \times_k \overline{B} = \overline{AB}
$$
\end{lemma}

\begin{proof}
包含关系 $\overline{AB} \subset \overline{A} \times_k \overline{B}$ 是显然的。
我们可将 $X$ 和 $Y$ 分别替换为既约闭子概形 $\overline{A}$ 和
$\overline{B}$。设 $W \subset X \times_k Y$ 是非空开子集。由《态射》引理
\ref{morphisms-lemma-scheme-over-field-universally-open}，子集
$U = \text{pr}_X(W)$ 是 $X$ 中的非空开集。因此 $A \cap U$ 非空；取
$a \in A \cap U$。记
$Y_a = \{a\} \times_k Y \cong Y_{\kappa(a)}$ 为投影
$\text{pr}_X : X \times_k Y \to X$ 在 $a$ 上方的纤维。再次由《态射》引理
\ref{morphisms-lemma-scheme-over-field-universally-open}，态射 $Y_a \to Y$
是开映射，因为 $\Spec(\kappa(a)) \to \Spec(k)$ 是万有开态射。因此非空开子集
$W_a = W \times_{X \times_k Y} Y_a$ 映到 $Y$ 的非空开子集。于是该像中存在
$b \in B$。故 $AB \cap W \not = \emptyset$，正是所需结论。
\end{proof}

\begin{lemma}
\label{varieties-lemma-closure-image-product-map}
设 $k$ 是域。设 $f : A \to X$、$g : B \to Y$ 是 $k$-概形之间的态射。
则在集合论意义下有
$$
\overline{f(A)} \times_k \overline{g(B)} =
\overline{(f \times g)(A \times_k B)}
$$
\end{lemma}

\begin{proof}
这由引理 \ref{varieties-lemma-closure-of-product} 得出，因为在该引理的记号下，
$f \times g$ 的像为 $f(A)g(B)$。
\end{proof}

\begin{lemma}
\label{varieties-lemma-scheme-theoretic-image-product-map}
设 $k$ 是域。设 $f : A \to X$、$g : B \to Y$ 是 $k$-概形之间的拟紧态射。
设 $Z \subset X$ 是 $f$ 的概形论像，见《态射》定义
\ref{morphisms-definition-scheme-theoretic-image}。类似地，设 $Z' \subset Y$
是 $g$ 的概形论像。则 $Z \times_k Z'$ 是 $f \times g$ 的概形论像。
\end{lemma}

\begin{proof}
回忆 $Z$ 是 $X$ 的最小闭子概形，并且 $f$ 分解通过它；$Z'$ 也类似。
设 $W \subset X \times_k Y$ 是 $f \times g$ 的概形论像。由于 $f \times g$
分解通过 $Z \times_k Z'$，故 $W \subset Z \times_k Z'$。

\medskip\noindent
为证明反向包含，设 $U \subset X$ 和 $V \subset Y$ 是仿射开集。由《态射》引理
\ref{morphisms-lemma-quasi-compact-scheme-theoretic-image}，概形 $Z \cap U$
是 $f|_{f^{-1}(U)} : f^{-1}(U) \to U$ 的概形论像；$Z' \cap V$ 和
$W \cap U \times_k V$ 也类似。因此可假设 $X$ 和 $Y$ 是仿射的。由于 $f$
和 $g$ 拟紧，这蕴涵 $A = \bigcup U_i$ 是仿射开集的有限并，而
$B = \bigcup V_j$ 也是仿射开集的有限并。于是可将 $A$ 替换为
$\coprod U_i$，将 $B$ 替换为 $\coprod V_j$；也就是说，还可假设 $A$ 和 $B$
也是仿射的。在此情形，$Z$ 由
$\Ker(\Gamma(X, \mathcal{O}_X) \to \Gamma(A, \mathcal{O}_A))$ 截出，
$Z'$ 和 $W$ 也类似。因此结论由等式
$$
\Gamma(A \times_k B, \mathcal{O}_{A \times_k B})
=
\Gamma(A, \mathcal{O}_A) \otimes_k \Gamma(B, \mathcal{O}_B)
$$
得出；该等式成立是因为 $A$ 和 $B$ 是仿射的。略去细节。
\end{proof}


\section{域上的光滑概形}
\label{varieties-section-smooth}

\noindent
下面给出两个刻画域上光滑概形的引理。

\begin{lemma}
\label{varieties-lemma-char-zero-differentials-free-smooth}
设 $k$ 是域。设 $X$ 是 $k$-概形。假设
\begin{enumerate}
\item $X$ 在 $k$ 上局部有限型，
\item $\Omega_{X/k}$ 局部自由，并且
\item $k$ 的特征为零。
\end{enumerate}
那么结构态射 $X \to \Spec(k)$ 是光滑的。
\end{lemma}

\begin{proof}
这由《代数》引理 \ref{algebra-lemma-characteristic-zero-local-smooth} 得出。
\end{proof}

\noindent
在正特征中，存在微分层自由的有限型非既约概形，例如
$\Spec(\mathbf{F}_p[t]/(t^p))$ 在 $\Spec(\mathbf{F}_p)$ 上的情形。若基域 $k$
的特征为 $p$ 且非完全，则存在既约概形 $X/k$，其微分层 $\Omega_{X/k}$
自由却不光滑，例如 $\Spec(k[t]/(t^p-a)$，其中 $a \in k$ 不是 $p$ 次幂。

\begin{lemma}
\label{varieties-lemma-char-p-differentials-free-smooth}
设 $k$ 是域。设 $X$ 是 $k$-概形。假设
\begin{enumerate}
\item $X$ 在 $k$ 上局部有限型，
\item $\Omega_{X/k}$ 局部自由，
\item $X$ 既约，并且
\item $k$ 完全。
\end{enumerate}
那么结构态射 $X \to \Spec(k)$ 是光滑的。
\end{lemma}

\begin{proof}
设 $x \in X$ 是一点。因为 $X$ 局部 Noether（见《态射》引理
\ref{morphisms-lemma-finite-type-noetherian}），只有有限多个不可约分支
$X_1, \ldots, X_n$ 通过 $x$（见《性质》引理
\ref{properties-lemma-Noetherian-topology} 和《拓扑》引理
\ref{topology-lemma-Noetherian}）。设 $\eta_i \in X_i$ 是泛点。由于 $X$ 既约，
有 $\mathcal{O}_{X, \eta_i} = \kappa(\eta_i)$，见《代数》引理
\ref{algebra-lemma-minimal-prime-reduced-ring}。此外，$\kappa(\eta_i)$ 是完全域
$k$ 的有限生成域扩张，因而在 $k$ 上可分生成（见《代数》第
\ref{algebra-section-separability} 节）。所以
$\Omega_{X/k, \eta_i} = \Omega_{\kappa(\eta_i)/k}$ 是自由的，其秩等于
$\kappa(\eta_i)$ 在 $k$ 上的超越次数。由《态射》引理
\ref{morphisms-lemma-dimension-fibre-at-a-point}，得到
$\dim_{\eta_i}(X_i) = \text{rank}_{\eta_i}(\Omega_{X/k})$。由于
$x \in X_1 \cap \ldots \cap X_n$，可见
$$
\text{rank}_x(\Omega_{X/k}) = \text{rank}_{\eta_i}(\Omega_{X/k}) = \dim(X_i).
$$
因此 $\dim_x(X) = \text{rank}_x(\Omega_{X/k})$，见《代数》引理
\ref{algebra-lemma-dimension-at-a-point-finite-type-over-field}。于是例如由《代数》引理
\ref{algebra-lemma-characterize-smooth-over-field}，$X \to \Spec(k)$ 在 $x$
处光滑。
\end{proof}

\begin{lemma}
\label{varieties-lemma-smooth-regular}
\begin{slogan}
域上光滑蕴涵正则
\end{slogan}
设 $X \to \Spec(k)$ 是光滑态射，其中 $k$ 是域。则 $X$ 是正则概形。
\end{lemma}

\begin{proof}
（另见引理 \ref{varieties-lemma-geometrically-regular-smooth}。）由《代数》引理
\ref{algebra-lemma-characterize-smooth-over-field}，每个局部环
$\mathcal{O}_{X, x}$ 都正则。又因为 $X$ 在 $k$ 上局部有限型，所以它局部
Noether。因此由《性质》引理 \ref{properties-lemma-characterize-regular}，
$X$ 正则。
\end{proof}

\begin{lemma}
\label{varieties-lemma-smooth-geometrically-normal}
设 $X \to \Spec(k)$ 是光滑态射，其中 $k$ 是域。则 $X$ 在 $k$ 上几何正则、
几何正规且几何既约。
\end{lemma}

\begin{proof}
（另见引理 \ref{varieties-lemma-geometrically-regular-smooth}。）设 $k'$ 是 $k$ 的有限纯不可分
扩张。只需证明 $X_{k'}$ 正则、正规且既约，见引理
\ref{varieties-lemma-geometrically-regular}、\ref{varieties-lemma-geometrically-normal} 和
\ref{varieties-lemma-check-only-finite-inseparable-extensions}。由《态射》引理
\ref{morphisms-lemma-base-change-smooth}，态射 $X_{k'} \to \Spec(k')$ 也是光滑的。
因此只需说明域上的光滑概形 $X$ 正则、正规且既约。由引理
\ref{varieties-lemma-smooth-regular} 可知 $X$ 正则。于是《性质》引理
\ref{properties-lemma-regular-normal} 保证 $X$ 正规。
\end{proof}

\begin{lemma}
\label{varieties-lemma-affine-space-over-field}
设 $k$ 是域。设 $d \geq 0$。设 $W \subset \mathbf{A}^d_k$ 是非空开集。
则存在闭点 $w \in W$，使得 $k \subset \kappa(w)$ 是有限可分扩张。
\end{lemma}

\begin{proof}
必要时缩小 $W$ 后，可假设 $W = \mathbf{A}^d_k \setminus V(f)$，其中
$f \in k[x_1, \ldots, x_d]$。若引理不成立，则
$f(a_1, \ldots, a_d) = 0$ 对所有
$(a_1, \ldots, a_d) \in (k^{sep})^d$ 都成立。
这是荒谬的，因为 $k^{sep}$ 是无限域。
\end{proof}

\begin{lemma}
\label{varieties-lemma-smooth-separable-closed-points-dense}
设 $k$ 是域。若 $X$ 在 $\Spec(k)$ 上光滑，则集合
$$
\{x \in X\text{ closed such that }k \subset \kappa(x)
\text{ is finite separable}\}
$$
在 $X$ 中稠密。
\end{lemma}

\begin{proof}
只需说明：给定非空光滑概形 $X$（定义在 $k$ 上），至少存在一个剩余域在 $k$ 上有限可分
的闭点。为此，选取图表
$$
\xymatrix{
X & U \ar[l] \ar[r]^-\pi & \mathbf{A}^d_k
}
$$
其中 $\pi$ 为平展态射，见《态射》引理
\ref{morphisms-lemma-smooth-etale-over-affine-space}。态射
$\pi : U \to \mathbf{A}^d_k$ 是开映射，见《态射》引理
\ref{morphisms-lemma-etale-open}。由引理 \ref{varieties-lemma-affine-space-over-field}，
可选取闭点 $w \in \pi(U)$，其剩余域在 $k$ 上有限可分。任选
$x \in U$ 使得 $\pi(x) = w$。由《态射》引理
\ref{morphisms-lemma-etale-over-field}，域扩张 $\kappa(x)/\kappa(w)$ 有限可分。
故 $\kappa(x)/k$ 有限可分。由《态射》引理
\ref{morphisms-lemma-algebraic-residue-field-extension-closed-point-fibre}，
点 $x$ 是 $X$ 的闭点。
\end{proof}

\begin{lemma}
\label{varieties-lemma-geometrically-reduced-dense-smooth-open}
设 $X$ 是域 $k$ 上局部有限型的既约概形。则 $X$ 在 $k$ 上几何既约，当且仅当
$X$ 含有一个在 $k$ 上光滑的稠密开集。
\end{lemma}

\begin{proof}
问题关于 $X$ 是局部的，故可假设 $X$ 拟紧。设
$X = X_1 \cup \ldots \cup X_n$ 是 $X$ 的不可约分支。先假设 $X$ 含有一个
在 $k$ 上光滑的稠密开集。那么 $X$ 在每个 $X_i$ 的泛点处几何既约（例如由引理
\ref{varieties-lemma-smooth-geometrically-normal}）。由引理
\ref{varieties-lemma-generic-points-geometrically-reduced}，可知 $X$ 几何既约。

\medskip\noindent
反之，假设 $X$ 几何既约。由于
$Z = \bigcup_{i \not = j} X_i \cap X_j$ 在 $X$ 中无处稠密，可将
$X$ 替换为 $X \setminus Z$。因为 $X \setminus Z$ 是不可约概形的不交并，问题归结为
$X$ 不可约的情形。由于 $X$ 不可约且既约，它是整概形，见《性质》引理
\ref{properties-lemma-characterize-integral}。设 $\eta \in X$ 是其泛点。那么函数域
$K = k(X) = \kappa(\eta) = \mathcal{O}_{X, \eta}$ 在 $k$ 上几何既约，因而在
$k$ 上可分，见《代数》引理
\ref{algebra-lemma-characterize-separable-field-extensions}。任取非空仿射开集
$U = \Spec(A) \subset X$，使得 $K = A_{(0)}$ 是 $A$ 的分式域。应用《代数》引理
\ref{algebra-lemma-separable-smooth}，可知 $A$ 在 $(0)$ 处在 $k$ 上光滑。
按定义，这意味着 $A$ 的某个主局部化在 $k$ 上光滑，结论得证。
\end{proof}

\begin{lemma}
\label{varieties-lemma-dense-smooth-open-variety-over-perfect-field}
设 $k$ 是完全域。设 $X$ 是局部代数的既约 $k$-概形，例如 $k$ 上的簇。则
$$
\{x \in X \mid X \to \Spec(k)\text{ is smooth at }x\} =
\{x \in X \mid \mathcal{O}_{X, x}\text{ 正则}\}
$$
且这是 $X$ 的稠密开子概形。
\end{lemma}

\begin{proof}
两个集合相等立即由《代数》引理 \ref{algebra-lemma-separable-smooth} 和各定义得出
（完全域的定义见《代数》定义 \ref{algebra-definition-perfect}）。该集合是开的，因为
概形态射的光滑点集合是开的，见《态射》定义
\ref{morphisms-definition-smooth}。最后给出两种论证说明它稠密：
(1) $X$ 的泛点属于该集合，因为泛点处的局部环是域（《代数》引理
\ref{algebra-lemma-minimal-prime-reduced-ring}），因而正则。
(2) 由引理 \ref{varieties-lemma-perfect-reduced}，$X$ 几何既约，因此可应用引理
\ref{varieties-lemma-geometrically-reduced-dense-smooth-open}。
\end{proof}

\begin{lemma}
\label{varieties-lemma-flat-under-smooth}
设 $k$ 是域。设 $f : X \to Y$ 是 $k$ 上局部有限型概形之间的态射。设
$x \in X$ 是一点，并令 $y = f(x)$。若 $X \to \Spec(k)$ 在 $x$ 处光滑，
且 $f$ 在 $x$ 处平坦，则 $Y \to \Spec(k)$ 在 $y$ 处光滑。特别地，若 $X$
在 $k$ 上光滑，并且 $f$ 平坦且满射，则 $Y$ 在 $k$ 上光滑。
\end{lemma}

\begin{proof}
只需说明 $Y$ 在 $y$ 处几何正则，见引理
\ref{varieties-lemma-geometrically-regular-smooth}。这由引理
\ref{varieties-lemma-flat-under-geometrically-regular} 得出（并将引理
\ref{varieties-lemma-geometrically-regular-smooth} 应用于 $(X, x)$）。
\end{proof}

\begin{lemma}
\label{varieties-lemma-variety-with-smooth-rational-point}
设 $k$ 是域。设 $X$ 是 $k$ 上的簇，并有一个 $k$-有理点 $x$，使得 $X$
在 $x$ 处光滑。则 $X$ 在 $k$ 上几何整。
\end{lemma}

\begin{proof}
设 $U \subset X$ 是 $X$ 的光滑点集。由假设，$U$ 非空，因而稠密且概形论稠密。
于是 $U_{\overline{k}} \subset X_{\overline{k}}$ 也稠密且概形论稠密（略去若干细节）。
因此只需说明 $U$ 几何整。由于 $U$ 有一个 $k$-有理点，由引理
\ref{varieties-lemma-geometrically-connected-if-connected-and-point}，它几何连通。另一方面，
$U_{\overline{k}}$ 既约且正规（引理
\ref{varieties-lemma-smooth-geometrically-normal}）。由于连通正规 Noether 概形是整概形
（《性质》引理 \ref{properties-lemma-normal-Noetherian}），证明完成。
\end{proof}

\begin{lemma}
\label{varieties-lemma-smooth-stratification}
设 $X$ 是域 $k$ 上的有限型概形。存在有限纯不可分扩张 $k'/k$、整数
$t \geq 0$ 和闭子概形
$$
X_{k'} \supset Z_0 \supset Z_1 \supset \ldots \supset Z_t = \emptyset
$$
使得 $Z_0 = (X_{k'})_{red}$，并且 $Z_i \setminus Z_{i + 1}$ 在 $k'$ 上
对所有 $i$ 都光滑。
\end{lemma}

\begin{proof}
可对 $\dim(X)$ 使用归纳法。由引理
\ref{varieties-lemma-finite-extension-geometrically-reduced}，可找到有限纯不可分扩张
$k'/k$，使得 $(X_{k'})_{red}$ 在 $k'$ 上几何既约。由引理
\ref{varieties-lemma-geometrically-reduced-dense-smooth-open}，存在无处稠密闭子概形
$X' \subset (X_{k'})_{red}$，使得 $(X_{k'})_{red} \setminus X'$ 在 $k'$ 上光滑。
于是 $\dim(X') < \dim(X)$。由归纳假设，存在有限纯不可分扩张 $k''/k'$、整数
$t' \geq 0$ 和闭子概形
$$
X'_{k''} \supset Y_0 \supset Y_1 \supset \ldots \supset Y_{t'} = \emptyset
$$
使得 $Y_0 = (X'_{k''})_{red}$，并且 $Y_i \setminus Y_{i + 1}$ 在 $k''$ 上
对所有 $i$ 都光滑。令 $t = t' + 1$，并考虑
$$
X_{k''} \supset Z_0 \supset Z_1 \supset \ldots \supset Z_t = \emptyset
$$
其中 $Z_0 = (X_{k''})_{red}$，且 $Z_i = Y_{i - 1}$（当 $i > 0$）；这是有意义的，
因为 $X'_{k''}$ 是 $X_{k''}$ 的闭子概形。略去所有所述性质均成立的验证。
\end{proof}


\section{簇的类型}
\label{varieties-section-types}

\noindent
本短节讨论簇的一些初等整体性质。

\begin{definition}
\label{varieties-definition-variety-type}
设 $k$ 是域。设 $X$ 是 $k$ 上的簇。
\begin{enumerate}
\item 若 $X$ 是仿射概形，则称 $X$ 为{\it 仿射簇}。这等价于要求 $X$
同构于 $\mathbf{A}^n_k$ 的闭子概形，其中 $n$ 适当选取。
\item 称 $X$ 为{\it 射影簇}，如果结构态射 $X \to \Spec(k)$ 是射影态射。
由《态射》引理 \ref{morphisms-lemma-characterize-locally-projective}，这成立当且仅当
$X$ 同构于 $\mathbf{P}^n_k$ 的闭子概形，其中 $n$ 适当选取。
\item 称 $X$ 为{\it 拟射影簇}，如果结构态射 $X \to \Spec(k)$ 是拟射影态射。
由《态射》引理 \ref{morphisms-lemma-characterize-locally-quasi-projective}，这成立当且仅当
$X$ 同构于 $\mathbf{P}^n_k$ 的局部闭子概形，其中 $n$ 适当选取。
\item {\it 固有簇}是态射 $X \to \Spec(k)$ 固有的簇。
\item {\it 光滑簇}是态射 $X \to \Spec(k)$ 光滑的簇。
\end{enumerate}
\end{definition}

\noindent
注意，射影簇是固有簇，见《态射》引理
\ref{morphisms-lemma-locally-projective-proper}。此外，仿射簇是拟射影的，因为
$\mathbf{A}^n_k$ 同构于 $\mathbf{P}^n_k$ 的开子概形，见《构造》引理
\ref{constructions-lemma-standard-covering-projective-space}。

\begin{lemma}
\label{varieties-lemma-regular-functions-proper-variety}
设 $X$ 是 $k$ 上的固有簇。则
\begin{enumerate}
\item $K = H^0(X, \mathcal{O}_X)$ 是域，并且是域 $k$ 的有限扩张；
\item 若 $X$ 几何既约，则 $K/k$ 可分；
\item 若 $X$ 几何不可约，则 $K/k$ 纯不可分；
\item 若 $X$ 几何整，则 $K = k$。
\end{enumerate}
\end{lemma}

\begin{proof}
这是引理 \ref{varieties-lemma-proper-geometrically-reduced-global-sections} 的特殊情形。
\end{proof}


\section{正规化}
\label{varieties-section-normalization}

\noindent
下面讨论与正规化有关的一些问题。

\begin{lemma}
\label{varieties-lemma-normalization-locally-algebraic}
设 $k$ 是域。设 $X$ 是 $k$ 上的局部代数概形。设
$\nu : X^\nu \to X$ 是正规化态射，见《态射》定义
\ref{morphisms-definition-normalization}。则
\begin{enumerate}
\item $\nu$ 有限且支配，而 $X^\nu$ 是 $k$ 上正规不可约局部代数概形的不交并；
\item $\nu$ 分解为 $X^\nu \to X_{red} \to X$，且第一个态射是
$X_{red}$ 的正规化态射；
\item 若 $X$ 是既约代数概形，则 $\nu$ 是双有理的；
\item 若 $X$ 是簇，则 $X^\nu$ 是簇，且 $\nu$ 是簇的有限双有理态射。
\end{enumerate}
\end{lemma}

\begin{proof}
由于 $X$ 是域上的局部有限型概形，可知 $X$ 局部 Noether（《态射》引理
\ref{morphisms-lemma-finite-type-noetherian}），因此每个拟紧开集只有有限多个不可约分支
（《性质》引理 \ref{properties-lemma-Noetherian-irreducible-components}）。所以可应用
《态射》定义 \ref{morphisms-definition-normalization}。正规化 $X^\nu$ 总是正规整概形
的不交并，且正规化态射 $\nu$ 总是支配的，见《态射》引理
\ref{morphisms-lemma-normalization-normal}。由于 $X$ 万有 Nagata（《态射》引理
\ref{morphisms-lemma-ubiquity-nagata}），可知 $\nu$ 有限（《态射》引理
\ref{morphisms-lemma-nagata-normalization}）。故 $X^\nu$ 也是局部代数的。至此已证明 (1)。

\medskip\noindent
(2) 是《态射》引理 \ref{morphisms-lemma-normalization-reduced}。

\medskip\noindent
(3) 是《态射》引理 \ref{morphisms-lemma-normalization-birational}。

\medskip\noindent
(4) 由 (1)、(2)、(3) 以及如下事实得出：$X^\nu$ 作为分离概形上的有限概形是分离的。
\end{proof}

\begin{lemma}
\label{varieties-lemma-normalize-projective}
设 $k$ 是域。设 $X$ 是 $k$ 上的固有概形。设
$\nu : X^\nu \to X$ 是正规化态射，见《态射》定义
\ref{morphisms-definition-normalization}。则 $X^\nu$ 在 $k$ 上固有。若 $X$
在 $k$ 上射影，则 $X^\nu$ 在 $k$ 上射影。
\end{lemma}

\begin{proof}
由引理 \ref{varieties-lemma-normalization-locally-algebraic}，态射 $\nu$ 有限。故由《态射》引理
\ref{morphisms-lemma-finite-proper} 和
\ref{morphisms-lemma-composition-proper}，$X^\nu$ 在 $k$ 上固有。
由《态射》引理 \ref{morphisms-lemma-finite-projective}，态射 $\nu$ 是射影的；
所以若 $X$ 在 $k$ 上射影，则由《态射》引理
\ref{morphisms-lemma-composition-projective}，$X^\nu$ 在 $k$ 上射影。
\end{proof}

\begin{lemma}
\label{varieties-lemma-relative-normalization-finite}
设 $k$ 是域。设 $f : Y \to X$ 是 $k$ 上局部代数概形之间的拟紧态射。
设 $X'$ 是 $X$ 在 $Y$ 中的正规化。若 $Y$ 既约，则 $X' \to X$ 有限。
\end{lemma}

\begin{proof}
由于 $Y$ 拟分离（由《性质》引理
\ref{properties-lemma-locally-Noetherian-quasi-separated} 和《态射》引理
\ref{morphisms-lemma-finite-type-noetherian}），态射 $f$ 拟分离（《概形》引理
\ref{schemes-lemma-compose-after-separated}）。因此可应用《态射》定义
\ref{morphisms-definition-normalization-X-in-Y}。结论由《态射》引理
\ref{morphisms-lemma-nagata-normalization-finite-general} 得出。这里用到：局部代数概形
局部 Noether（因而局部只有有限多个不可约分支），并且局部代数概形是 Nagata 的
（《态射》引理 \ref{morphisms-lemma-ubiquity-nagata}）。略去一些小细节。
\end{proof}

\begin{lemma}
\label{varieties-lemma-finite-extension-geometrically-normal}
设 $k$ 是域。设 $X$ 是代数 $k$-概形。则存在有限纯不可分扩张 $k'/k$，
使得正规化 $Y$（取自 $X_{k'}$）在 $k'$ 上几何正规。
\end{lemma}

\begin{proof}
令 $K = k^{perf}$ 为完全闭包。设 $Y_K$ 是 $X_K$ 的正规化，见引理
\ref{varieties-lemma-normalization-locally-algebraic}。由《极限》引理
\ref{limits-lemma-descend-finite-presentation}，存在有限子扩张 $K/k'/k$ 和有限表示态射
$\nu : Y \to X_{k'}$，其到 $K$ 的基变换是正规化态射
$\nu_K : Y_K \to X_K$。注意 $Y$ 在 $k'$ 上几何正规（引理
\ref{varieties-lemma-geometrically-normal}）。增大 $k'$ 后，可假设 $Y \to X_{k'}$ 有限
（《极限》引理 \ref{limits-lemma-descend-finite-finite-presentation}）。由于
$\nu_K : Y_K \to X_K$ 是正规化态射，它诱导双有理态射
$Y_K \to (X_K)_{red}$。故存在稠密开集 $V_K \subset X_K$，使得
$\nu_K^{-1}(V_K) \to V_K$ 是闭浸入（并诱导 $\nu_K^{-1}(V_K)$ 与
$V_{K, red}$ 的同构；例如见《态射》引理
\ref{morphisms-lemma-birational-isomorphism-over-dense-open}）。增大 $k'$ 后，可找到
$V_K$ 是某个稠密开集 $V \subset Y$ 的基变换，并且态射
$\nu^{-1}(V) \to V$ 是闭浸入（《极限》引理
\ref{limits-lemma-descend-opens} 和
\ref{limits-lemma-descend-closed-immersion-finite-presentation}）。由此容易看出
$\nu$ 是正规化态射，证明完成。
\end{proof}

\begin{lemma}
\label{varieties-lemma-normalization-and-change-of-fields}
设 $k$ 是域。设 $X$ 是局部代数 $k$-概形。设 $K/k$ 是域扩张。设
$\nu : X^\nu \to X$ 是 $X$ 的正规化，并设 $Y^\nu \to X_K$ 是基变换的
正规化。那么典范态射
$$
Y^\nu \longrightarrow X^\nu \times_{\Spec(k)} \Spec(K)
$$
在 $K/k$ 可分时是同构，一般情形下是万有同胚。
\end{lemma}

\begin{proof}
令 $Y = X_K$。设 $X^{(0)}$ 和 $Y^{(0)}$ 分别是 $X$ 和 $Y$ 的不可约分支泛点集。
那么投影态射 $\pi : Y \to X$ 满足 $\pi(Y^{(0)}) = X^{(0)}$。这是因为
$\pi$ 满射、开且保泛化，见《态射》引理
\ref{morphisms-lemma-scheme-over-field-universally-open} 和
\ref{morphisms-lemma-open-generizing}。若将 $X^{(0)}$ 和 $Y^{(0)}$ 分别视为（既约）
概形，则 $X^\nu$ 和 $Y^\nu$ 分别是 $X$ 和 $Y$ 在 $X^{(0)}$ 和
$Y^{(0})$ 中的正规化。因此《态射》引理
\ref{morphisms-lemma-functoriality-normalization} 给出典范态射
$Y^\nu \to X^\nu$，它位于 $Y \to X$ 上方；再由纤维积的万有性质给出
本引理中的典范态射。

\medskip\noindent
为证明该态射具有引理所述性质，可假设 $X = \Spec(A)$ 是仿射的。设
$Q(A_{red})$ 为 $A_{red}$ 的总分式环。那么 $X^\nu$ 是整闭包 $A'$ 的谱，
该整闭包取自 $A$ 在 $Q(A_{red})$ 中的整闭包，见《态射》引理
\ref{morphisms-lemma-normalization-reduced} 和
\ref{morphisms-lemma-description-normalization}。类似地，$Y^\nu$ 是整闭包 $B'$ 的谱，
该整闭包取自 $A \otimes_k K$ 在 $Q((A \otimes_k K)_{red})$ 中的整闭包。
有典范映射 $Q(A_{red}) \to Q((A \otimes_k K)_{red})$、典范映射
$A' \to B'$，而引理中的态射对应于诱导映射
$$
A' \otimes_k K \longrightarrow B'
$$
这个 $K$-代数映射的核由幂零元组成，因为
$Q(A_{red}) \otimes_k K \to Q((A \otimes_k K)_{red})$ 的核就是幂零元集合。

\medskip\noindent
若 $K/k$ 可分，则由引理 \ref{varieties-lemma-base-change-normal-by-separable}，
$A' \otimes_k K$ 正规。特别地，它既约，因而
$Q((A \otimes_k K)_{red}) = Q(A' \otimes_k K)$，并且由《代数》引理
\ref{algebra-lemma-characterize-reduced-ring-normal}，
$B' = A' \otimes_k K$。

\medskip\noindent
假设 $K/k$ 不可分。那么 $k$ 的特征为 $p > 0$。我们将证明：对每个
$b \in B'$，都存在幂 $q$，它是 $p$ 的幂，使得 $b^q$ 属于
$A' \otimes_k K$ 的像。
由《代数》引理 \ref{algebra-lemma-p-ring-map}，这将证明上述映射是万有同胚。
给定 $b$，存在子域 $F \subset K$，其中 $F/k$ 有限生成，使得 $b$ 属于
$Q((A \otimes_k F)_{red})$，且整于 $A \otimes_k F$。选取首一多项式
$P = T^d + \alpha_1 T^{d - 1} + \ldots + \alpha _d$，使得
$P(b) = 0$ 且 $\alpha_i \in A \otimes_k F$。选取超越基
$t_1, \ldots, t_r$，它是 $F$ 在 $k$ 上的超越基。设
$F/F'/k(t_1, \ldots, t_r)$ 是最大可分子扩张
（《域》引理 \ref{fields-lemma-separable-first}）。由于 $F/F'$ 有限纯不可分，
存在 $q$，使得 $\lambda^q \in F'$ 对所有 $\lambda \in F$ 都成立。于是 $b^q$
属于 $Q((A \otimes_k F')_{red})$，并满足多项式
$T^d + \alpha_1^q T^{d - 1} + \ldots + \alpha _d^q$，其中
$\alpha_i^q \in A \otimes_k F'$。由可分情形可知
$b^q \in A' \otimes_k F'$，证明完成。
\end{proof}

\begin{lemma}
\label{varieties-lemma-geometrically-normal-in-codim-1}
设 $k$ 是域。设 $X$ 是局部代数 $k$-概形。设 $\nu : X^\nu \to X$ 是
$X$ 的正规化。设 $x \in X$ 是一点，满足：(a) $\mathcal{O}_{X, x}$ 既约；
(b) $\dim(\mathcal{O}_{X, x}) = 1$；并且 (c) 对每个 $x' \in X^\nu$，
若 $\nu(x') = x$，则扩张 $\kappa(x')/k$ 可分。那么 $X$ 在 $x$ 处几何既约，
并且 $X^\nu$ 在 $x'$ 处几何正则，只要 $\nu(x') = x$。
\end{lemma}

\begin{proof}
以下不再特别说明地使用引理 \ref{varieties-lemma-normalization-locally-algebraic} 的结果。
设 $x' \in X^\nu$ 是 $x$ 上方的一点。由维数理论（第
\ref{varieties-section-algebraic-schemes} 节），有
$\dim(\mathcal{O}_{X^\nu, x'}) = 1$。由于 $X^\nu$ 正规，可知
$\mathcal{O}_{X^\nu, x'}$ 是离散赋值环（《性质》引理
\ref{properties-lemma-criterion-normal}）。因此 $\mathcal{O}_{X^\nu, x'}$
是正则局部 $k$-代数，其剩余域在 $k$ 上可分。故
$k \to \mathcal{O}_{X^\nu, x'}$ 在 $\mathfrak m_{x'}$-进拓扑下形式光滑，见
《更多代数》引理 \ref{more-algebra-lemma-regular-implies-fs}。于是由《更多代数》定理
\ref{more-algebra-theorem-regular-fs}，$\mathcal{O}_{X^\nu, x'}$ 在 $k$ 上几何正则。
再由引理 \ref{varieties-lemma-geometrically-regular-at-point}，$X^\nu$ 在 $x'$ 处几何正则。

\medskip\noindent
由于 $\mathcal{O}_{X, x}$ 既约，映射族
$\mathcal{O}_{X, x} \to \mathcal{O}_{X^\nu, x'}$ 是单射。由于
$\mathcal{O}_{X^\nu, x'}$ 是几何既约 $k$-代数，立即可知
$\mathcal{O}_{X, x}$ 是几何既约 $k$-代数。故由引理
\ref{varieties-lemma-geometrically-reduced-at-point}，$X$ 在 $x$ 处几何既约。
\end{proof}


\section{可逆函数群}
\label{varieties-section-units}

\noindent
群 $\mathcal{O}^*(X)/k^*$ 往往（但不总是）是有限生成 Abel 群，只要 $X$ 是 $k$
上的簇。我们用一系列引理证明这一点。一切都建立在下述特殊情形上。

\begin{lemma}
\label{varieties-lemma-open-in-normal-proper}
设 $k$ 是代数闭域。设 $\overline{X}$ 是 $k$ 上的固有簇。设
$X \subset \overline{X}$ 是开子概形。假设 $X$ 正规。则
$\mathcal{O}^*(X)/k^*$ 是有限生成 Abel 群。
\end{lemma}

\begin{proof}
由于命题只涉及 $X$，可以将 $\overline{X}$ 替换为 $k$ 上另一个固有簇。设
$\nu : \overline{X}^\nu \to
\overline{X}$ 是正规化态射。由引理
\ref{varieties-lemma-normalization-locally-algebraic}，$\nu$ 有限，且
$\overline{X}^\nu$ 是簇。由于 $X$ 正规，可知
$\nu^{-1}(X) \to X$ 是同构（略去一个小细节）。最后，因为有限态射是固有的
（《态射》引理 \ref{morphisms-lemma-finite-proper}），且固有态射的复合是固有的
（《态射》引理 \ref{morphisms-lemma-composition-proper}），可知
$\overline{X}^\nu$ 在 $k$ 上固有。因此可以且将假设 $\overline{X}$ 正规。

\medskip\noindent
以下不再特别说明地使用：对任意仿射开集 $U$（它是 $\overline{X}$ 的开集），环
$\mathcal{O}(U)$ 是有限生成 $k$-代数，它是 Noether 环、整环且正规；见《代数》引理
\ref{algebra-lemma-Noetherian-permanence}、《性质》定义
\ref{properties-definition-integral}、《性质》引理
\ref{properties-lemma-locally-Noetherian} 和
\ref{properties-lemma-locally-normal}，以及《态射》引理
\ref{morphisms-lemma-locally-finite-type-characterize}。

\medskip\noindent
设 $\xi_1, \ldots, \xi_r$ 是 $X$ 在 $\overline{X}$ 中之补集的泛点。它们只有有限多个，
因为 $\overline{X}$ 的底拓扑空间是 Noether 的（见《态射》引理
\ref{morphisms-lemma-finite-type-noetherian}、《性质》引理
\ref{properties-lemma-Noetherian-topology} 和《拓扑》引理
\ref{topology-lemma-Noetherian}）。对每个 $i$，局部环
$\mathcal{O}_i = \mathcal{O}_{X, \xi_i}$ 是正规 Noether 局部整环（因为它是 Noether 正规整环的
局部化）。设 $J \subset \{1, \ldots, r\}$ 为所有指标 $i$ 所成的集合，这些指标满足
$\dim(\mathcal{O}_i) = 1$。对 $j \in J$，局部环
$\mathcal{O}_j$ 是离散赋值环，见《代数》引理
\ref{algebra-lemma-characterize-dvr}。因此得到赋值
$$
v_j : k(\overline{X})^* \longrightarrow \mathbf{Z}
$$
并且 $v_j(f) \geq 0 \Leftrightarrow f \in \mathcal{O}_j$。

\medskip\noindent
将 $\mathcal{O}(X)$ 视为一个 $k$-子代数，包含在 $k(X) = k(\overline{X})$ 中。我们断言映射
$$
\mathcal{O}(X)^* \longrightarrow
\prod\nolimits_{j \in J} \mathbf{Z},
\quad
f \longmapsto \prod v_j(f)
$$
的核是 $k^*$。显然，该断言证明引理。事实上，设
$f \in \mathcal{O}(X)$ 是核中的一个元素。设
$U = \Spec(B) \subset \overline{X}$ 是任意仿射开集。则 $B$ 是 Noether 正规整环。
对每个高为一的素理想 $\mathfrak q \subset B$ 及其对应点 $\xi \in X$，可知要么
$\xi = \xi_j$ 对某个 $j \in J$ 成立，要么 $\xi \in X$。理由是，由《性质》引理
\ref{properties-lemma-codimension-local-ring}，
$\text{codim}(\overline{\{\xi\}}, \overline{X}) = 1$；所以若
$\xi \in \overline{X} \setminus X$，则它必是 $\overline{X} \setminus X$ 的泛点，从而等于某个
$\xi_j$，$j \in J$。两种情形下都有
$f \in \mathcal{O}_{X, \xi} = B_{\mathfrak q}$，因为 $f$ 在该映射的核中。因此
$f \in \bigcap_{\text{ht}(\mathfrak q) = 1} B_{\mathfrak q} = B$，见《代数》引理
\ref{algebra-lemma-normal-domain-intersection-localizations-height-1}。换言之，可知
$f \in \Gamma(\overline{X}, \mathcal{O}_{\overline{X}})$。但由于 $k$ 代数闭，由引理
\ref{varieties-lemma-regular-functions-proper-variety} 可得 $f \in k$。
\end{proof}

\noindent
接下来，我们用一些初等论证推广上述情形，仍然保持域代数闭。

\begin{lemma}
\label{varieties-lemma-units-integral-finite-type-algebraically-closed}
设 $k$ 是代数闭域。设 $X$ 是 $k$ 上的局部有限型整概形。则
$\mathcal{O}^*(X)/k^*$ 是有限生成 Abel 群。
\end{lemma}

\begin{proof}
由于 $X$ 是整的，限制映射 $\mathcal{O}(X) \to \mathcal{O}(U)$ 对任意非空开子概形
$U \subset X$ 都是单射。因此可假设 $X$ 仿射。选取闭浸入
$X \to \mathbf{A}^n_k$，并用 $\overline{X}$ 表示 $X$ 在 $\mathbf{P}^n_k$ 中经由通常浸入
$\mathbf{A}^n_k \to \mathbf{P}^n_k$ 所得的闭包。从而可假设 $X$ 是某个射影簇
$\overline{X}$ 的仿射开集。

\medskip\noindent
设 $\nu : \overline{X}^\nu \to \overline{X}$ 是正规化态射，见《态射》定义
\ref{morphisms-definition-normalization}。我们知道 $\nu$ 有限且支配，并且
$\overline{X}^\nu$ 是正规不可约概形，见《态射》引理
\ref{morphisms-lemma-normalization-normal}、
\ref{morphisms-lemma-Japanese-normalization} 和
\ref{morphisms-lemma-ubiquity-nagata}。因此 $\overline{X}^\nu$ 是固有簇，因为
$\overline{X} \to \Spec(k)$ 作为有限态射与固有态射的复合是固有的（见《态射》第
\ref{morphisms-section-proper} 节和第
\ref{morphisms-section-integral} 节中的结果）。还可知 $\nu$ 是满射，因为
$\nu$ 的像是闭的，并包含 $\overline{X}$ 的泛点。因此令
$X^\nu = \nu^{-1}(X)$，可知只需对 $X^\nu$ 证明结论。换言之，可假设
$X$ 是正规固有簇 $\overline{X}$ 的非空开集。该情形由引理
\ref{varieties-lemma-open-in-normal-proper} 处理。
\end{proof}

\noindent
前一引理立即给出下述稍微的推广。

\begin{lemma}
\label{varieties-lemma-units-general-algebraically-closed}
设 $k$ 是代数闭域。设 $X$ 是 $k$ 上局部有限型的连通既约概形，并且只有
有限多个不可约分支。则 $\mathcal{O}^*(X)/k^*$ 是有限生成 Abel 群。
\end{lemma}

\begin{proof}
设 $X = \bigcup X_i$ 是不可约分支的并。由引理
\ref{varieties-lemma-units-integral-finite-type-algebraically-closed}，可知
$\mathcal{O}(X_i)^*/k^*$ 是有限生成 Abel 群。设
$f \in \mathcal{O}(X)^*$ 属于映射
$$
\mathcal{O}(X)^* \longrightarrow \prod \mathcal{O}(X_i)^*/k^*.
$$
的核。则对每个 $i$，存在元素 $\lambda_i \in k$，使得
$f|_{X_i} = \lambda_i$。限制到 $X_i \cap X_j$ 可得 $\lambda_i = \lambda_j$，只要
$X_i \cap X_j \not = \emptyset$。由于 $X$ 连通，可知所有
$\lambda_i$ 相等，因而 $f \in k^*$。这证明了
$$
\mathcal{O}(X)^*/k^* \subset \prod \mathcal{O}(X_i)^*/k^*
$$
而引理得证，因为右端是有限多个有限生成 Abel 群的积。
\end{proof}

\begin{lemma}
\label{varieties-lemma-integral-closure-ground-field}
设 $k$ 是域。设 $X$ 是 $k$ 上连通且既约的概形。则 $k$ 在
$\Gamma(X, \mathcal{O}_X)$ 中的整闭包是域。
\end{lemma}

\begin{proof}
设 $k' \subset \Gamma(X, \mathcal{O}_X)$ 是 $k$ 的整闭包。则
$X \to \Spec(k)$ 经由 $\Spec(k')$ 分解，见《概形》引理
\ref{schemes-lemma-morphism-into-affine}。由于 $X$ 既约，可知 $k'$ 没有非零幂零元。
由于 $k \to k'$ 整，可知 $k'$ 的每个素理想既是极大理想，又是极小素理想，
而 $\Spec(k')$ 全不连通；见《代数》引理
\ref{algebra-lemma-integral-no-inclusion} 和
\ref{algebra-lemma-ring-with-only-minimal-primes}。由于 $X$ 连通，态射
$X \to \Spec(k')$ 为常值态射，设其像为对应于
$\mathfrak p \subset k'$ 的点。那么任意 $f \in k'$ 若满足
$f \not \in \mathfrak p$，其像都是 $\mathcal{O}_X$ 的可逆元。由 $k'$ 的定义，这进而迫使
$f$ 是 $k'$ 的单位。所以 $k'$ 是以 $\mathfrak p$ 为极大理想的局部环，见《代数》引理
\ref{algebra-lemma-characterize-local-ring}。我们已经看到 $k'$ 既约，因此这意味着 $k'$ 是域，
见《代数》引理 \ref{algebra-lemma-minimal-prime-reduced-ring}。
\end{proof}

\begin{proposition}
\label{varieties-proposition-units-general}
设 $k$ 是域。设 $X$ 是 $k$ 上的概形。假设 $X$ 在 $k$ 上局部有限型、连通、既约，
并且只有有限多个不可约分支。若除上述条件外，下列条件之一成立，则
$\mathcal{O}(X)^*/k^*$ 是有限生成 Abel 群：
\begin{enumerate}
\item $k$ 在 $\Gamma(X, \mathcal{O}_X)$ 中的整闭包是 $k$；
\item $X$ 有一个 $k$-有理点；或
\item $X$ 几何整。
\end{enumerate}
\end{proposition}

\begin{proof}
设 $\overline{k}$ 是 $k$ 的代数闭包。设 $Y$ 是
$(X_{\overline{k}})_{red}$ 的一个连通分支。注意，典范态射
$p : Y \to X$ 是开的（由《态射》引理
\ref{morphisms-lemma-scheme-over-field-universally-open}），也是闭的（由《态射》引理
\ref{morphisms-lemma-integral-universally-closed}）。因此 $p(Y) = X$，因为假设 $X$ 连通。
特别地，由于 $X$ 既约，这意味着
$\mathcal{O}(X) \subset \mathcal{O}(Y)$。由引理
\ref{varieties-lemma-galois-action-irreducible-components-locally-finite-type}，可知 $Y$ 只有有限多个不可约分支
（该引理中的域 $\overline{k}$ 是可分代数闭包，但这无关紧要，因为由引理
\ref{varieties-lemma-separably-closed-field-irreducible-components}，两个基变换
$X_{\overline{k}}$ 有相同的不可约分支）。因此引理
\ref{varieties-lemma-units-general-algebraically-closed} 可应用于 $Y$。这意味着：若
$\mathcal{O}(X)^*/k^*$ 不是有限生成 Abel 群，则存在元素
$f \in \mathcal{O}(X)$、$f \not \in k$，其像属于 $\overline{k}$，这里使用映射
$\mathcal{O}(X) \to \mathcal{O}(Y)$。在这种情形下，
$f$ 在 $k$ 上代数，从而整于 $k$。因此若条件 (1) 成立，则这不可能发生。
为完成证明，我们证明条件 (2) 和 (3) 都推出 (1)。

\medskip\noindent
设 $k \subset k' \subset \Gamma(X, \mathcal{O}_X)$ 为 $k$ 在
$\Gamma(X, \mathcal{O}_X)$ 中的整闭包。由引理
\ref{varieties-lemma-integral-closure-ground-field}，可知 $k'$ 是域。若
$e : \Spec(k) \to X$ 是 $k$-有理点，则
$e^\sharp : \Gamma(X, \mathcal{O}_X) \to k$ 是包含映射
$k \to \Gamma(X, \mathcal{O}_X)$ 的一个截面。特别地，$e^\sharp$ 限制到 $k'$ 后是域映射
$k' \to k$，且位于 $k$ 上；这显然表明 (2) 推出 (1)。

\medskip\noindent
若整闭包 $k'$，即 $k$ 在 $\Gamma(X, \mathcal{O}_X)$ 中的整闭包，非平凡，则可知：
要么 $X$ 不几何连通（当 $k'/k$ 不是纯不可分的时候），要么 $X$ 不几何既约
（当 $k'/k$ 是非平凡纯不可分的时候）。略去细节。因此 (3) 推出 (1)。
\end{proof}

\begin{lemma}
\label{varieties-lemma-units-variety}
设 $k$ 是域。设 $X$ 是 $k$ 上的簇。若下列条件中至少一个成立，则群
$\mathcal{O}(X)^*/k^*$ 是有限生成 Abel 群：
\begin{enumerate}
\item $k$ 在 $\Gamma(X, \mathcal{O}_X)$ 中整闭；
\item $k$ 在 $k(X)$ 中代数闭；
\item $X$ 在 $k$ 上几何整；或
\item $k$ 是域扩张 $\kappa(x)/k$ 的“交”，其中 $x$ 遍历 $x$ 的闭点。
\end{enumerate}
\end{lemma}

\begin{proof}
由命题 \ref{varieties-proposition-units-general} 可知 (1) 足够。我们略去 (2)、(3)、(4) 各自推出 (1) 的验证。
\end{proof}


\section{Künneth 公式，I}
\label{varieties-section-kunneth}

\noindent
本节证明底是域且我们考虑拟凝聚模上同调时的 Künneth 公式。
更一般的版本见《概形的导出范畴》第
\ref{perfect-section-kunneth} 节。

\begin{lemma}
\label{varieties-lemma-kunneth}
设 $k$ 是域。设 $X$ 和 $Y$ 是 $k$ 上的概形，并设 $\mathcal{F}$ 和
$\mathcal{G}$ 分别是拟凝聚 $\mathcal{O}_X$-模和 $\mathcal{O}_Y$-模。则有典范同构
$$
H^n(X \times_{\Spec(k)} Y, \text{pr}_1^*\mathcal{F}
\otimes_{\mathcal{O}_{X \times_{\Spec(k)} Y}} \text{pr}_2^*\mathcal{G}) =
\bigoplus\nolimits_{p + q = n}
H^p(X, \mathcal{F}) \otimes_k H^q(Y, \mathcal{G})
$$
只要 $X$ 和 $Y$ 拟紧且具有仿射对角态射\footnote{
$X$ 和 $Y$ 拟分离的情形将在下文引理
\ref{varieties-lemma-kunneth-general} 中讨论。}（例如 $X$ 和 $Y$ 都分离）。
\end{lemma}

\begin{proof}
在本证明中，没有额外标注的积与张量积都取在 $k$ 上。作为映射
$$
H^p(X, \mathcal{F}) \otimes H^q(Y, \mathcal{G})
\longrightarrow
H^n(X \times Y, \text{pr}_1^*\mathcal{F}
\otimes_{\mathcal{O}_{X \times Y}} \text{pr}_2^*\mathcal{G})
$$
我们先利用上同调的函子性得到映射
$H^p(X, \mathcal{F}) \to H^p(X \times Y, \text{pr}_1^*\mathcal{F})$ 和
$H^p(Y, \mathcal{G}) \to H^p(X \times Y, \text{pr}_2^*\mathcal{G})$，
然后使用杯积
$$
\cup :
H^p(X \times Y, \text{pr}_1^*\mathcal{F})
\otimes H^q(X \times Y, \text{pr}_2^*\mathcal{G})
\longrightarrow
H^n(X \times Y, \text{pr}_1^*\mathcal{F}
\otimes_{\mathcal{O}_{X \times Y}} \text{pr}_2^*\mathcal{G})
$$
当 $X$ 和 $Y$ 仿射时，结论由仿射概形上高次上同调群的消失
（《概形的上同调》引理
\ref{coherent-lemma-quasi-coherent-affine-cohomology-zero}）以及定义
（拟凝聚模的拉回和拟凝聚模的张量积的定义）得出。

\medskip\noindent
选取有限仿射开覆盖
$\mathcal{U} : X = \bigcup_{i \in I} U_i$ 和
$\mathcal{V} : Y = \bigcup_{j \in J} V_j$。它们确定仿射开覆盖
$\mathcal{W} : X \times Y = \bigcup_{(i, j) \in I \times J} U_i \times V_j$。
注意，$\mathcal{W}$ 同时是 $\text{pr}_1^{-1}\mathcal{U}$ 和
$\text{pr}_2^{-1}\mathcal{V}$ 的加细。所以由《上同调》引理
\ref{cohomology-lemma-functoriality-cech}，得到映射
$$
\check{\mathcal{C}}^\bullet(\mathcal{U}, \mathcal{F}) \to
\check{\mathcal{C}}^\bullet(\mathcal{W}, \text{pr}_1^*\mathcal{F})
\quad\text{且}\quad
\check{\mathcal{C}}^\bullet(\mathcal{V}, \mathcal{G}) \to
\check{\mathcal{C}}^\bullet(\mathcal{W}, \text{pr}_2^*\mathcal{G})
$$
它们与上同调上的拉回映射相容。在《上同调》等式
(\ref{cohomology-equation-needs-signs}) 中，我们已构造复形映射
$$
\text{Tot}(
\check{\mathcal{C}}^\bullet(\mathcal{W}, \text{pr}_1^*\mathcal{F})
\otimes
\check{\mathcal{C}}^\bullet(\mathcal{W}, \text{pr}_2^*\mathcal{G}))
\longrightarrow
\check{\mathcal{C}}^\bullet(\mathcal{W},
\text{pr}_1^*\mathcal{F} \otimes_{\mathcal{O}_{X \times Y}}
\text{pr}_2^*\mathcal{G})
$$
它定义上同调上的杯积。结合上述映射，得到复形映射
\begin{equation}
\label{varieties-equation-kunneth-on-cech}
\text{Tot}(
\check{\mathcal{C}}^\bullet(\mathcal{U}, \mathcal{F})
\otimes
\check{\mathcal{C}}^\bullet(\mathcal{V}, \mathcal{G}))
\longrightarrow
\check{\mathcal{C}}^\bullet(\mathcal{W},
\text{pr}_1^*\mathcal{F} \otimes_{\mathcal{O}_{X \times Y}}
\text{pr}_2^*\mathcal{G})
\end{equation}
我们提醒读者，该映射不是复形的同构。回忆，可用这些覆盖计算我们的拟凝聚层的上同调
（《概形的上同调》引理
\ref{coherent-lemma-affine-diagonal} 和
\ref{coherent-lemma-cech-cohomology-quasi-coherent}）。因此，
(\ref{varieties-equation-kunneth-on-cech}) 在上同调上重现本引理的映射。

\medskip\noindent
考虑拟凝聚模的短正合列
$0 \to \mathcal{F} \to \mathcal{F}' \to \mathcal{F}'' \to 0$。由于
(\ref{varieties-equation-kunneth-on-cech}) 的构造对 $\mathcal{F}$ 具有函子性，并且相关
{\v C}ech 复形的构造对变量 $\mathcal{F}$ 正合（因为我们在仿射开集上取截面），
得到短正合复形列之间的映射
\begingroup\scriptsize
$$
\xymatrix{
\text{Tot}(
\check{\mathcal{C}}^\bullet(\mathcal{U}, \mathcal{F})
\otimes
\check{\mathcal{C}}^\bullet(\mathcal{V}, \mathcal{G})) \ar[r] \ar[d] &
\text{Tot}(
\check{\mathcal{C}}^\bullet(\mathcal{U}, \mathcal{F}')
\otimes
\check{\mathcal{C}}^\bullet(\mathcal{V}, \mathcal{G})) \ar[r] \ar[d] &
\text{Tot}(
\check{\mathcal{C}}^\bullet(\mathcal{U}, \mathcal{F}'')
\otimes
\check{\mathcal{C}}^\bullet(\mathcal{V}, \mathcal{G})) \ar[d] \\
\check{\mathcal{C}}^\bullet(\mathcal{W},
\text{pr}_1^*\mathcal{F} \otimes_{\mathcal{O}_{X \times Y}}
\text{pr}_2^*\mathcal{G}) \ar[r] &
\check{\mathcal{C}}^\bullet(\mathcal{W},
\text{pr}_1^*\mathcal{F}' \otimes_{\mathcal{O}_{X \times Y}}
\text{pr}_2^*\mathcal{G}) \ar[r] &
\check{\mathcal{C}}^\bullet(\mathcal{W},
\text{pr}_1^*\mathcal{F}'' \otimes_{\mathcal{O}_{X \times Y}}
\text{pr}_2^*\mathcal{G})
}
$$
\endgroup
（我们略去了外侧的零项）。考察上同调长正合列，可知：若引理的结论对 $2$-out-of-$3$ of
$\mathcal{F}, \mathcal{F}', \mathcal{F}''$ 成立，则它对第三个也成立。

\medskip\noindent
注意，$X$ 对拟凝聚模具有有限上同调维数，见《概形的上同调》引理
\ref{coherent-lemma-vanishing-nr-affines}。对
$d(\mathcal{F}) = \max \{d \mid H^d(X, \mathcal{F}) \not = 0\}$
作归纳，我们将化归到 $d(\mathcal{F}) = 0$ 的情形。假设
$d(\mathcal{F}) > 0$。由《概形的上同调》引理
\ref{coherent-lemma-affine-diagonal-universal-delta-functor}，我们已经看到，存在嵌入
$\mathcal{F} \to \mathcal{F}'$，使得 $H^p(X, \mathcal{F}') = 0$ 对所有 $p \geq 1$ 成立。令
$\mathcal{F}'' = \Coker(\mathcal{F} \to \mathcal{F}')$，则
$d(\mathcal{F}'') < d(\mathcal{F})$。因此可用前段的结果看出，只需对
$\mathcal{F}'$ 和 $\mathcal{F}''$ 证明引理，由此完成归纳步骤。

\medskip\noindent
对 $\mathcal{G}$ 作同样的论证，可假设 $\mathcal{F}$ 和 $\mathcal{G}$ 都只在次数
$0$ 具有非零上同调。设 $V \subset Y$ 是仿射开集。考虑仿射开覆盖
$\mathcal{U}_V : X \times V = \bigcup_{i \in I} U_i \times V$。立即有
$$
\check{\mathcal{C}}^\bullet(\mathcal{U}, \mathcal{F}) \otimes \mathcal{G}(V) =
\check{\mathcal{C}}^\bullet(\mathcal{U}_V,
\text{pr}_1^*\mathcal{F} \otimes_{\mathcal{O}_{X \times Y}}
\text{pr}_2^*\mathcal{G})
$$
（复形相等）。我们得出
$$
R\text{pr}_{2, *}(\text{pr}_1^*\mathcal{F} \otimes_{\mathcal{O}_{X \times Y}}
\text{pr}_2^*\mathcal{G})
\cong
\Gamma(X, \mathcal{F}) \otimes_k \mathcal{G} \cong
\bigoplus\nolimits_{\alpha \in A} \mathcal{G}
$$
在 $Y$ 上成立。这里 $A$ 是 $k$-向量空间 $\Gamma(X, \mathcal{F})$ 的一组基。
$Y$ 上的上同调与直和可交换（《上同调》引理
\ref{cohomology-lemma-quasi-separated-cohomology-colimit}）。对 $\text{pr}_2$ 使用 Leray 谱序列
（经由《上同调》引理 \ref{cohomology-lemma-apply-Leray}），我们得出，
$H^n(X \times Y, \text{pr}_1^*\mathcal{F} \otimes_{\mathcal{O}_{X \times Y}}
\text{pr}_2^*\mathcal{G})$
在 $n > 0$ 时为零，且同构于
$H^0(X, \mathcal{F}) \otimes H^0(Y, \mathcal{G})$，当 $n = 0$ 时。这完成了证明
（除了还应检查该同构确实由次数 $0$ 上的杯积给出；我们略去验证）。
\end{proof}

\begin{lemma}
\label{varieties-lemma-kunneth-general}
设 $k$ 是域。设 $X$ 和 $Y$ 是 $k$ 上的概形，并设 $\mathcal{F}$ 和
$\mathcal{G}$ 分别是拟凝聚 $\mathcal{O}_X$-模和 $\mathcal{O}_Y$-模。则有典范同构
$$
H^n(X \times_{\Spec(k)} Y, \text{pr}_1^*\mathcal{F}
\otimes_{\mathcal{O}_{X \times_{\Spec(k)} Y}} \text{pr}_2^*\mathcal{G}) =
\bigoplus\nolimits_{p + q = n}
H^p(X, \mathcal{F}) \otimes_k H^q(Y, \mathcal{G})
$$
只要 $X$ 和 $Y$ 拟紧且拟分离。
\end{lemma}

\begin{proof}
若 $X$ 和 $Y$ 分离，或更一般地具有仿射对角态射，请参见引理
\ref{varieties-lemma-kunneth} 中的“更好”证明（它胜过本证明之处在于：它把这些映射识别为先拉回再作杯积）。
设 $X'$ 和 $Y'$ 分别是 $X$ 和 $Y$ 的无穷小加厚，其结构层分别为
$\mathcal{O}_{X'} = \mathcal{O}_X \oplus \mathcal{F}$
和 $\mathcal{O}_{Y'} = \mathcal{O}_Y \oplus \mathcal{G}$，其中 $\mathcal{F}$ 和
$\mathcal{G}$ 分别是平方为零的理想。那么
$$
\mathcal{O}_{X' \times Y'} =
\mathcal{O}_{X \times Y}
\oplus \text{pr}_1^*\mathcal{F}
\oplus \text{pr}_2^*\mathcal{G}
\oplus \text{pr}_1^*\mathcal{F}
\otimes_{\mathcal{O}_{X \times Y}} \text{pr}_2^*\mathcal{G}
$$
作为 $X \times Y$ 上的层成立。由此可知，只需证明
$$
H^n(X \times Y, \mathcal{O}_{X \times Y}) =
\bigoplus\nolimits_{p + q = n}
H^p(X, \mathcal{O}_X) \otimes_k H^q(Y, \mathcal{O}_Y)
$$
对 $k$ 上任意一对拟紧拟分离概形成立。略去一些细节。

\medskip\noindent
为证明此断言，我们使用《概形的上同调》引理
\ref{coherent-lemma-hypercoverings} 所给出形式的上同调与基变换。该引理告诉我们，存在一个下有界的
$k$-向量空间复形，也就是复形 $\mathcal{K}^\bullet$，它由 $\Spec(k)$ 上的拟凝聚模组成，
它万有地计算 $Y$ 在 $\Spec(k)$ 上的上同调。特别地，可知
$$
R\text{pr}_{1, *}(\mathcal{O}_{X \times Y}) \cong
(X \to \Spec(k))^*\mathcal{K}^\bullet
$$
在 $D(\mathcal{O}_X)$ 中成立。在同伦意义下，复形 $\mathcal{K}^\bullet$ 同构于
$\bigoplus_{q \geq 0} H^q(Y, \mathcal{O}_Y)[-q]$，因为域上的每个向量空间复形都如此。
我们得出
$$
R\text{pr}_{1, *}(\mathcal{O}_{X \times Y}) \cong
\bigoplus\nolimits_{q \geq 0}
H^q(Y, \mathcal{O}_Y)[-q] \otimes_k \mathcal{O}_X
$$
在 $D(\mathcal{O}_X)$ 中成立。于是
\begin{align*}
R\Gamma(X \times Y, \mathcal{O}_{X \times Y})
& =
R\Gamma(X, R\text{pr}_{1, *}(\mathcal{O}_{X \times Y})) \\
& =
R\Gamma(X, \bigoplus\nolimits_{q \geq 0}
H^q(Y, \mathcal{O}_Y)[-q] \otimes_k \mathcal{O}_X) \\
& =
\bigoplus\nolimits_{q \geq 0}
R\Gamma(X,  H^q(Y, \mathcal{O}_Y) \otimes \mathcal{O}_X)[-q] \\
& =
\bigoplus\nolimits_{q \geq 0}
R\Gamma(X, \mathcal{O}_X) \otimes_k H^q(Y, \mathcal{O}_Y)[-q] \\
& =
\bigoplus\nolimits_{p, q \geq 0}
H^p(X, \mathcal{O}_X)[-p] \otimes_k H^q(Y, \mathcal{O}_Y)[-q]
\end{align*}
如所求。第一个等式由 $\text{pr}_1$ 的 Leray 谱序列得出
（《上同调》引理 \ref{cohomology-lemma-before-Leray}）。
第二个等式由上面给出的总直像分解得出。
第三个等式成立，因为上同调总与有限直和可交换（并且由《概形的上同调》引理
\ref{coherent-lemma-vanishing-nr-affines-quasi-separated}，$Y$ 的上同调在充分高次数上消失）。
第四个等式成立，因为由《上同调》引理
\ref{cohomology-lemma-quasi-separated-cohomology-colimit}，$X$ 上的上同调与无限直和可交换。
最后一个等式由上面对域的导出范畴的评论得出。
\end{proof}

\section{簇的 Picard 群}
\label{varieties-section-picard-groups}

\noindent
本节汇集代数簇的 Picard 群的一些初等结果。

\begin{lemma}
\label{varieties-lemma-change-rings-pic-pre}
设 $A \to B$ 是忠实平坦环同态。设 $X$ 是 $A$ 上的拟紧拟分离概形。
设 $\mathcal{L}$ 是可逆 $\mathcal{O}_X$-模，其到 $X_B$ 的拉回平凡。则
$H^0(X, \mathcal{L})$ 和 $H^0(X, \mathcal{L}^{\otimes -1})$ 是可逆
$H^0(X, \mathcal{O}_X)$-模，并且乘法映射诱导同构
$$
H^0(X, \mathcal{L}) \otimes_{H^0(X, \mathcal{O}_X)}
H^0(X, \mathcal{L}^{\otimes -1}) \longrightarrow
H^0(X, \mathcal{O}_X)
$$
\end{lemma}

\begin{proof}
以 $\mathcal{L}_B$ 表示 $\mathcal{L}$ 到 $X_B$ 的拉回。
选取同构 $\mathcal{L}_B \to \mathcal{O}_{X_B}$。
令 $R = H^0(X, \mathcal{O}_X)$、$M = H^0(X, \mathcal{L})$，并将
$M$ 视为 $R$-模。对每个拟凝聚 $\mathcal{O}_X$-模 $\mathcal{F}$，记
$\mathcal{F}_B$ 为它在 $X_B$ 上的拉回，则有典范同构
$H^0(X_B, \mathcal{F}_B) = H^0(X, \mathcal{F}) \otimes_A B$；见
《概形的上同调》引理 \ref{coherent-lemma-flat-base-change-cohomology}。
因此有
$$
M \otimes_R (R \otimes_A B) =
M \otimes_A B = H^0(X_B, \mathcal{L}_B) \cong
H^0(X_B, \mathcal{O}_{X_B}) = R \otimes_A B
$$
由于 $R \to R \otimes_A B$ 忠实平坦（它是忠实平坦同态
$A \to B$ 的基变换），由《代数》命题
\ref{algebra-proposition-ffdescent-finite-projectivity} 可知，$M$ 是可逆
$R$-模。类似地，$N = H^0(X, \mathcal{L}^{\otimes -1})$ 是可逆 $R$-模。
要看关于张量积的断言成立，只需利用它在拉回到 $X_B$ 后成立，以及
$R \to R \otimes_A B$ 的忠实平坦性。略去一些细节。
\end{proof}

\begin{lemma}
\label{varieties-lemma-change-rings-pic}
设 $A \to B$ 是忠实平坦环同态。设 $X$ 是 $A$ 上的概形，并且
\begin{enumerate}
\item $X$ 拟紧且拟分离；
\item $R = H^0(X, \mathcal{O}_X)$ 是半局部环。
\end{enumerate}
则拉回映射 $\Pic(X) \to \Pic(X_B)$ 是单射。
\end{lemma}

\begin{proof}
设 $\mathcal{L}$ 是可逆 $\mathcal{O}_X$-模，其拉回 $\mathcal{L}'$ 在
$X_B$ 上平凡。令 $M = H^0(X, \mathcal{L})$ 和
$N = H^0(X, \mathcal{L}^{\otimes - 1})$。由引理
\ref{varieties-lemma-change-rings-pic-pre}，$R$-模 $M$ 和 $N$ 可逆。由于 $R$ 半局部，
$M \cong R$ 且 $N \cong R$；见《代数》引理
\ref{algebra-lemma-locally-free-semi-local-free}。选取生成元 $s \in M$ 和
$t \in N$。由引理 \ref{varieties-lemma-change-rings-pic-pre} 的最后一部分，
$st \in R = H^0(X, \mathcal{O}_X)$ 是单位。我们得出，$s$ 和 $t$ 分别给出
$\mathcal{L}$ 和 $\mathcal{L}^{\otimes -1}$ 在 $X$ 上的平凡化。
\end{proof}

\begin{lemma}
\label{varieties-lemma-change-fields-pic}
设 $k'/k$ 是域扩张。设 $X$ 是 $k$ 上的概形，并且
\begin{enumerate}
\item $X$ 拟紧且拟分离；
\item $R = H^0(X, \mathcal{O}_X)$ 是半局部环，例如当 $\dim_k R < \infty$ 时如此。
\end{enumerate}
则拉回映射 $\Pic(X) \to \Pic(X_{k'})$ 是单射。
\end{lemma}

\begin{proof}
这是引理 \ref{varieties-lemma-change-rings-pic} 的特殊情形。若 $\dim_k R < \infty$，则
$R$ 是 Artin 环，因而是半局部环（《代数》引理
\ref{algebra-lemma-finite-dimensional-algebra} 和
\ref{algebra-lemma-artinian-finite-nr-max}）。
\end{proof}

\begin{example}
\label{varieties-example-need-condition}
若不对概形 $X$（视为域 $k$ 上的概形）施加某种条件，引理
\ref{varieties-lemma-change-fields-pic} 并不成立。下面给出一个例子。
设 $k$ 是域。设 $t \in \mathbf{P}^1_k$ 是闭点。令
$X = \mathbf{P}^1 \setminus \{t\}$。则有满射
$$
\mathbf{Z} = \Pic(\mathbf{P}^1_k) \longrightarrow \Pic(X)
$$
第一个等式由《除子》引理 \ref{divisors-lemma-Pic-projective-space-UFD} 得出，
满射性由《除子》引理 \ref{divisors-lemma-extend-invertible-module} 得出
（因为 $\mathbf{P}^1_k$ 的维数为 $1$，且它在 $k$ 上光滑，故其所有局部环都是离散赋值环）。
若 $\mathcal{L}$ 位于所示映射的核中，则有
$\mathcal{L} \cong \mathcal{O}_{\mathbf{P}^1_k}(nt)$，其中 $n \in \mathbf{Z}$。我们把下述验证留给读者：
$\mathcal{O}_{\mathbf{P}^1_k}(t) \cong \mathcal{O}_{\mathbf{P}^1_k}(d)$，
其中 $d = [\kappa(t) : k]$。因此
$$
\Pic(X) = \mathbf{Z}/d\mathbf{Z}
$$
所以，若 $t$ 不是 $k$-有理点，则 $d > 1$，这个 Picard 群非零。
另一方面，若把底域 $k$ 扩张为任意域扩张 $k'$，且存在 $k$-嵌入
$\kappa(t) \to k'$，则 $\mathbf{P}^1_{k'} \setminus X_{k'}$ 有一个
$k'$-有理点 $t'$。因此
$\mathcal{O}_{\mathbf{P}^1_{k'}}(1) = \mathcal{O}_{\mathbf{P}^1_{k'}}(t')$
属于映射 $\mathbf{Z} \to \Pic(X_{k'})$ 的核；按与上面相同的方法可得
$\Pic(X_{k'}) = 0$。
\end{example}

\noindent
下列引理告诉我们，在正规簇上，“有理等价的可逆模”彼此同构。

\begin{lemma}
\label{varieties-lemma-rational-equivalence-for-Pic}
设 $k$ 是域。设 $X$ 是 $k$ 上的正规簇。设
$U \subset \mathbf{A}^n_k$ 是开子概形，并具有 $k$-有理点
$p, q \in U(k)$。对每个可逆模 $\mathcal{L}$（定义在
$X \times_{\Spec(k)} U$ 上），
其限制 $\mathcal{L}|_{X \times p}$ 和 $\mathcal{L}|_{X \times q}$ 同构。
\end{lemma}

\begin{proof}
$X \times_{\Spec(k)} U \to X$ 的纤维是域上的仿射 $n$-空间的开子概形。
因此，由《除子》引理 \ref{divisors-lemma-open-subscheme-UFD}，这些纤维的
Picard 群平凡。应用《除子》引理 \ref{divisors-lemma-in-image-pullback}，
可见 $\mathcal{L}$ 是某个可逆模 $\mathcal{N}$ 的拉回，而该模定义在 $X$ 上。
\end{proof}

\section{底域的唯一性}
\label{varieties-section-base-field}

\noindent
“设 $X$ 是 $k$ 上的概形”这一说法意指 $X$ 是一个概形，并配备态射
$X \to \Spec(k)$。现在可以问，域 $k$ 是否由概形 $X$ 唯一确定。
当然并非如此。例如，我们通常把 $\mathbf{A}^1_{\mathbf{C}}$ 看作复数域
$\mathbf{C}$ 上的概形，但也可以把它看作 $\mathbf{Q}$ 上的概形。然而，若要求态射
$X \to \Spec(k)$ 不通过
$X \to \Spec(k') \to \Spec(k)$ 分解，其中 $k'/k$ 是任意非平凡域扩张，
情形又如何？换言之，我们要求 $k$ 在某种意义下是使 $X$ 定义在 $k$ 上的极大域。

\medskip\noindent
下列例子表明，这仍不能保证 $k$ 的唯一性。考虑概形
$$
X =
\Spec\left(
\mathbf{Q}(x)[y]\left[\frac{1}{P(y)}, P \in \mathbf{Q}[y], P \not = 0\right]
\right)
$$
乍看之下，它似乎是 $\mathbf{Q}(x)$ 上的概形；再看便会发现，它也是
$\mathbf{Q}(y)$ 上的概形。此外，域 $\mathbf{Q}(x)$ 和 $\mathbf{Q}(y)$ 都是
$R = \Gamma(X, \mathcal{O}_X)$ 的子域，并且在 $R$ 的所有子域中极大
（略去细节）。特别地，$\mathbf{Q}(x)$ 和 $\mathbf{Q}(y)$ 在上述意义下都是极大的。
注意，态射 $X \to \Spec(\mathbf{Q}(x))$ 和
$X \to \Spec(\mathbf{Q}(y))$ 都“本质有限型”
（即对应的环同态本质有限型）。因此 $X$ 是有限维 Noether 概形，也就是说，
它并非完全病态。

\medskip\noindent
另一个可能妨碍唯一性的问题是，概形 $X$ 可能非既约。在这种情形，从 $X$ 到
给定域的谱可以有许多不同的态射。作为一个显式例子，考虑对偶数
$D = \mathbf{C}[y]/(y^2) = \mathbf{C} \oplus \epsilon \mathbf{C}$。
给定任意导子 $\theta : \mathbf{C} \to \mathbf{C}$（定义在 $\mathbf{Q}$ 上），便得到环同态
$$
\mathbf{C} \longrightarrow D, \quad
c \longmapsto c + \epsilon \theta(c).
$$
$\mathbf{C}$ 中使所有这些映射都相同的子域是 $\mathbf{Q}$ 的代数闭包。
这意味着，若 $X$ 非既约，把 $X$ 能定义于其上的所有域取交，所得可能是很小的域。

\medskip\noindent
与此有关的一个观察如下：给定域 $k$ 以及两个子域 $k_1, k_2$（它们都是 $k$ 的子域），
若 $k$ 在 $k_1$ 和 $k_2$ 上都有限，一般而言 $k$ 在 $k_1 \cap k_2$ 上
{\it 并不}有限。例如，取域 $k = \mathbf{Q}(t)$ 及其子域
$k_1 = \mathbf{Q}(t^2)$ 和 $\mathbf{Q}((t + 1)^2)$。在此情形，
$k_1 \cap k_2 = \mathbf{Q}$。因此在下文取域的交时必须谨慎。

\medskip\noindent
说完这些，我们现在证明：若 $X$ 在某个域上局部有限型，则仍有某种唯一性。
精确结果如下。

\begin{proposition}
\label{varieties-proposition-unique-base-field}
设 $X$ 是概形。设 $a : X \to \Spec(k_1)$ 和
$b : X \to \Spec(k_2)$ 是从 $X$ 到域的谱的态射。
假设 $a, b$ 局部有限型，并且 $X$ 既约且连通。则
$k_1' = k_2'$，其中 $k_i' \subset \Gamma(X, \mathcal{O}_X)$ 是
$k_i$ 在 $\Gamma(X, \mathcal{O}_X)$ 中的整闭包。
\end{proposition}

\begin{proof}
首先假设该引理在 $X$ 拟紧的情形成立（我们将在下面处理拟紧情形）。
由于 $X$ 在域上局部有限型，它局部 Noether；见《态射》引理
\ref{morphisms-lemma-finite-type-noetherian}。特别地，这意味着它局部连通，
开子集的连通分支是开集，并且拟紧开集的交拟紧；见《性质》引理
\ref{properties-lemma-Noetherian-topology}、《拓扑》引理
\ref{topology-lemma-locally-connected}、《拓扑》第
\ref{topology-section-noetherian} 节以及《拓扑》引理
\ref{topology-lemma-constructible-Noetherian-space}。
选取开覆盖 $X = \bigcup_{i \in I} U_i$，使每个 $U_i$ 都拟紧且连通。
对每个 $i$，令 $K_i \subset \mathcal{O}_X(U_i)$ 为 $k_1$ 和 $k_2$ 的整闭包。
对每对 $i, j \in I$，把
$$
U_i \cap U_j = \coprod U_{i, j, l}
$$
分解为其有限多个连通分支。记
$K_{i, j, l} \subset \mathcal{O}(U_{i, j, l})$ 为 $k_1$ 和 $k_2$ 的整闭包。
由引理 \ref{varieties-lemma-integral-closure-ground-field}，环 $K_i$ 和
$K_{i, j, l}$ 都是域。现在我们断言，$k_1'$ 和 $k_2'$ 都等于下列映射的核：
$$
\prod K_i \longrightarrow \prod K_{i, j, l}, \quad
(x_i)_i \longmapsto x_i|_{U_{i, j, l}} - x_j|_{U_{i, j, l}}
$$
这便证明了所求。的确，显然 $k_1'$ 包含于这个核。反过来，设
$(x_i)_i$ 位于该核中。由层条件，$(x_i)_i$ 对应于
$f \in \mathcal{O}(X)$。选取某个 $i_0 \in I$，并令 $P(T) \in k_1[T]$
为满足 $P(x_{i_0}) = 0$ 的首一多项式。于是我们断言 $P(f) = 0$；这证明了
$f \in k_1$。为证明这一点，须证明 $P(x_i) = 0$ 对所有 $i$ 成立。
选取 $i \in I$。由于 $X$ 连通，存在序列
$i_0, i_1, \ldots, i_n = i \in I$，使得
$U_{i_t} \cap U_{i_{t + 1}} \not = \emptyset$。这意味着，对每个 $t$，存在
$l_t$，使 $x_{i_t}$ 和 $x_{i_{t + 1}}$ 映到域
$K_{i_t, i_{t + 1}, l_t}$ 的同一个元素。因此若 $P(x_{i_t}) = 0$，则
$P(x_{i_{t + 1}}) = 0$。从 $P(x_{i_0}) = 0$ 出发归纳，便得到
$P(x_i) = 0$，如所求。

\medskip\noindent
为完成该引理的证明，我们在附加假设 $X$ 拟紧的条件下证明该引理。
由引理 \ref{varieties-lemma-integral-closure-ground-field}，把 $k_i$ 替换为 $k_i'$ 后，
可以假设 $k_i$ 在 $\Gamma(X, \mathcal{O}_X)$ 中整闭。这蕴涵
$\mathcal{O}(X)^*/k_i^*$ 是有限生成 Abel 群；见命题
\ref{varieties-proposition-units-general}。令 $k_{12} = k_1 \cap k_2$，视其为
$\mathcal{O}(X)$ 的子环。注意，$k_{12}$ 是域。由于有映射
$$
k_1^*/k_{12}^* \longrightarrow \mathcal{O}(X)^*/k_2^*
$$
可见 $k_1^*/k_{12}^*$ 也是有限生成 Abel 群。因此存在
$\alpha_1, \ldots, \alpha_n \in k_1^*$，使每个元素 $\lambda \in k_1$ 都具有形式
$$
\lambda = c \alpha_1^{e_1} \ldots \alpha_n^{e_n}
$$
其中 $e_i \in \mathbf{Z}$ 且 $c \in k_{12}$。特别地，环同态
$$
k_{12}[x_1, \ldots, x_n, \frac{1}{x_1 \ldots x_n}] \longrightarrow k_1, \quad
x_i \longmapsto \alpha_i
$$
是满射。由 Hilbert 零点定理，即《代数》定理
\ref{algebra-theorem-nullstellensatz}，可知 $k_1$ 是 $k_{12}$ 的有限扩张。
同理可知 $k_2$ 是 $k_{12}$ 的有限扩张。特别地，$k_1$ 和 $k_2$ 都包含于
整闭包 $k_{12}'$；这是 $k_{12}$ 在 $\Gamma(X, \mathcal{O}_X)$ 中的整闭包。但由引理
\ref{varieties-lemma-integral-closure-ground-field}，$k_{12}'$ 是域；并且由于我们选取
$k_i$ 在 $\Gamma(X, \mathcal{O}_X)$ 中整闭，故得到
$k_1 = k_{12} = k_2$，如所求。
\end{proof}

\section{自同构}
\label{varieties-section-automorphisms}

\noindent
本节讨论域上概形的自同构。关于曲线的（无穷小）自同构的一些信息，见
《代数曲线》第 \ref{curves-section-vector-fields} 节和
《曲线的模》第 \ref{moduli-curves-section-finite-aut} 节。

\begin{lemma}
\label{varieties-lemma-automorphism}
设 $X$ 是域 $k$ 上有限型的既约概形。设 $f : X \to X$ 是 $k$ 上的自同构，
并且它在 $X$ 的底拓扑空间上诱导恒等映射。则
\begin{enumerate}
\item 有 $f^*\mathcal{F} \cong \mathcal{F}$，对任意凝聚
$\mathcal{O}_X$-模皆然；
\item 若每个不可约分支的维数都满足 $\dim(Z) > 0$（其中 $Z \subset X$），
则 $f$ 是恒等映射。
\end{enumerate}
\end{lemma}

\begin{proof}
(1) 由 (2) 以及 $X$ 的维数为 $0$ 的连通分支都是域的谱这一事实得出。

\medskip\noindent
设 $Z \subset X$ 是不可约分支，并把它看作整的闭子概形。显然
$f(Z) \subset Z$，且 $f|_Z : Z \to Z$ 是 $k$ 上的自同构，并在 $Z$ 的
底拓扑空间上诱导恒等映射。由于 $X$ 既约，只须证明箭头
$f|_Z : Z \to Z$ 是恒等映射。这把问题归约到下一段讨论的情形。

\medskip\noindent
假设 $X$ 不可约且维数 $> 0$。选取非空仿射开集 $U \subset X$。
由于 $f(U) \subset U$，并且 $U \subset X$ 概形论稠密，只须证明
$f|_U : U \to U$ 是恒等映射。

\medskip\noindent
假设 $X = \Spec(A)$ 仿射、不可约、维数 $> 0$，并且 $k$ 是无限域。
设 $g \in A$ 非常数。集合
$$
S = \bigcup\nolimits_{\lambda \in k} V(g - \lambda)
$$
在 $X$ 中稠密，因为它是稠密子集 $\mathbf{A}^1_k(k)$ 在非常数态射
$g : X \to \mathbf{A}^1_k$ 下的逆像。若 $x \in S$，则
$g(x)$（即 $g$ 在 $\kappa(x)$ 中的像）属于 $k \to \kappa(x)$ 的像。因此
$f^\sharp : \kappa(x) \to \kappa(x)$ 固定 $g(x)$。所以
$f^\sharp(g)$ 在 $\kappa(x)$ 中的像等于 $g(x)$。我们得出
$$
S \subset V(g - f^\sharp(g))
$$
而由于 $X$ 既约且 $S$ 稠密，便得到 $g=f^\sharp(g)$。这证明
$f^\sharp = \text{id}_A$，因为 $A$ 作为 $k$-代数由上述元素 $g$ 生成
（略去细节；提示：常数函数之集是有限维 $k$-向量子空间，位于 $A$ 中）。
我们得出 $f = \text{id}_X$。

\medskip\noindent
假设 $X = \Spec(A)$ 仿射、不可约、维数 $> 0$，并且 $k$ 是有限域。
若对每个 $1$ 维整的闭子概形 $C \subset X$，限制 $f|_C : C \to C$ 都是
恒等映射，则 $f$ 是恒等映射。这把问题归约到 $X$ 为曲线的情形。
有限域上的曲线具有有限自同构群（略去细节）。因此 $f$ 具有有限阶，记为
$n$。然后如上选取非常数态射 $g : X \to \mathbf{A}^1_k$，并考虑
$$
S = \{x \in X\text{ closed such that }[\kappa(g(x)) : k]
\text{ is prime to }n\}
$$
如前论证，可知 $S$ 在 $X$ 中稠密。对闭点 $x \in X$，映射
$f^\sharp : \kappa(x) \to \kappa(x)$ 是阶整除 $n$ 的自同构。因此对
$x \in S$，该自同构在由 $g$ 在 $\kappa(x)$ 中的像生成的子域上作用平凡。
于是 $S \subset V(g - f^\sharp(g))$，按前述论证即得结论。
\end{proof}


\section{欧拉示性数}
\label{varieties-section-euler}

\noindent
本节证明域上固有概形上的凝聚层的欧拉示性数的一些初等性质。

\begin{definition}
\label{varieties-definition-euler-characteristic}
设 $k$ 是域。设 $X$ 是 $k$ 上的固有概形。设 $\mathcal{F}$ 是凝聚
$\mathcal{O}_X$-模。在此情形下，{\it $\mathcal{F}$ 的欧拉示性数}是整数
$$
\chi(X, \mathcal{F}) = \sum\nolimits_i (-1)^i \dim_k H^i(X, \mathcal{F}).
$$
关于这个公式的合理性，见下文。
\end{definition}

\noindent
在该定义的情形下，向量空间 $H^i(X, \mathcal{F})$ 中只有有限多个非零
（《概形的上同调》引理
\ref{coherent-lemma-quasi-coherence-higher-direct-images}），
并且这些空间都是有限维的（《概形的上同调》引理
\ref{coherent-lemma-proper-over-affine-cohomology-finite}）。因此
$\chi(X, \mathcal{F}) \in \mathbf{Z}$ 有良好定义。注意，这个定义依赖于域
$k$，而不只依赖于二元组 $(X, \mathcal{F})$。

\begin{lemma}
\label{varieties-lemma-euler-characteristic-additive}
设 $k$ 是域。设 $X$ 是 $k$ 上的固有概形。设
$0 \to \mathcal{F}_1 \to \mathcal{F}_2 \to \mathcal{F}_3 \to 0$
是 $X$ 上凝聚模的短正合列。则
$$
\chi(X, \mathcal{F}_2) = \chi(X, \mathcal{F}_1) + \chi(X, \mathcal{F}_3)
$$
\end{lemma}

\begin{proof}
考虑与该引理中的短正合列相伴的上同调长正合列
$$
0 \to H^0(X, \mathcal{F}_1) \to H^0(X, \mathcal{F}_2) \to
H^0(X, \mathcal{F}_3) \to H^1(X, \mathcal{F}_1) \to \ldots
$$
线性代数中的秩-零化度定理表明
$$
0 = \dim H^0(X, \mathcal{F}_1) - \dim H^0(X, \mathcal{F}_2)
+ \dim H^0(X, \mathcal{F}_3) - \dim H^1(X, \mathcal{F}_1) + \ldots
$$
这立即蕴涵该引理。
\end{proof}

\begin{lemma}
\label{varieties-lemma-chi-tensor-finite}
设 $k$ 是域。设 $X$ 是 $k$ 上的固有概形。设 $\mathcal{F}$ 是满足
$\dim(\text{Supp}(\mathcal{F})) \leq 0$ 的凝聚层。则
\begin{enumerate}
\item $\mathcal{F}$ 由全局截面生成，
\item $H^0(X, \mathcal{F}) =
\bigoplus_{x \in \text{Supp}(\mathcal{F})} \mathcal{F}_x$，
\item $H^i(X, \mathcal{F}) = 0$ 对每个 $i > 0$ 成立，
\item $\chi(X, \mathcal{F}) = \dim_k H^0(X, \mathcal{F})$，并且
\item 等式
$\chi(X, \mathcal{F} \otimes \mathcal{E}) = n\chi(X, \mathcal{F})$
对每个局部自由模 $\mathcal{E}$（其秩为 $n$）成立。
\end{enumerate}
\end{lemma}

\begin{proof}
由《概形的上同调》引理 \ref{coherent-lemma-coherent-support-closed}，
可知 $\mathcal{F} = i_*\mathcal{G}$，其中 $i : Z \to X$ 是
$\mathcal{F}$ 的概形论支集的包含映射，而 $\mathcal{G}$ 是凝聚
$\mathcal{O}_Z$-模。由概形论支集的定义，$Z$ 的底拓扑空间是
$\text{Supp}(\mathcal{F})$。由于 $Z$ 的维数为 $0$，由《性质》引理
\ref{properties-lemma-locally-Noetherian-dimension-0} 可知 $Z$ 是仿射概形。
所以 $\mathcal{G}$ 是全局生成的，并且 $\mathcal{G}$ 的高阶上同调群为零
（《概形的上同调》引理
\ref{coherent-lemma-quasi-coherent-affine-cohomology-zero}）。
事实上，由引理 \ref{varieties-lemma-algebraic-scheme-dim-0}，概形 $Z$ 是局部
Artin 环的谱的有限不交并。因此相应地有
$H^0(Z, \mathcal{G}) = \bigoplus_{z \in Z} \mathcal{G}_z$。
由《概形的上同调》引理 \ref{coherent-lemma-relative-affine-cohomology}，
$\mathcal{F}$ 与 $\mathcal{G}$ 的上同调相同。因此
$H^i(X, \mathcal{F}) = 0$ 对每个 $i > 0$ 成立，并且
$H^0(X, \mathcal{F}) = H^0(Z, \mathcal{G})$。
特别地，(3) 成立。对 $z \in Z$，若它对应于
$x \in \text{Supp}(\mathcal{F})$，则有
$\mathcal{G}_z = (i_*\mathcal{G})_x = \mathcal{F}_x$。
我们得出 (2) 成立。当然，(2) 蕴涵 (1)。由欧拉示性数
$\chi(X, \mathcal{F})$ 的定义和 (3)，得到 (4)。由投影公式
（《上同调》引理 \ref{cohomology-lemma-projection-formula}），有
$$
i_*(\mathcal{G} \otimes i^*\mathcal{E}) = \mathcal{F} \otimes \mathcal{E}.
$$
由于 $Z$ 的维数为 $0$，局部自由层 $i^*\mathcal{E}$ 同构于
$\mathcal{O}_Z^{\oplus n}$；按上述方式论证即知 (5) 成立。
\end{proof}

\begin{lemma}
\label{varieties-lemma-euler-characteristic-extend-base-field}
设 $k'/k$ 是域扩张。设 $X$ 是 $k$ 上的固有概形。设 $\mathcal{F}$ 是
$X$ 上的凝聚层。设 $\mathcal{F}'$ 是 $\mathcal{F}$ 在 $X_{k'}$ 上的拉回。
则 $\chi(X, \mathcal{F}) = \chi(X', \mathcal{F}')$。
\end{lemma}

\begin{proof}
这是因为由平坦基变换有
$$
H^i(X_{k'}, \mathcal{F}') = H^i(X, \mathcal{F}) \otimes_k k'
$$
见《概形的上同调》引理
\ref{coherent-lemma-flat-base-change-cohomology}。
\end{proof}

\begin{lemma}
\label{varieties-lemma-euler-characteristic-morphism}
设 $k$ 是域。设 $f : Y \to X$ 是 $k$ 上的固有概形之间的态射。设
$\mathcal{G}$ 是凝聚 $\mathcal{O}_Y$-模。则
$$
\chi(Y, \mathcal{G}) = \sum (-1)^i \chi(X, R^if_*\mathcal{G})
$$
\end{lemma}

\begin{proof}
该公式是有意义的：层 $R^if_*\mathcal{G}$ 是凝聚的，并且其中只有有限多个非零；
见《概形的上同调》命题
\ref{coherent-proposition-proper-pushforward-coherent} 和引理
\ref{coherent-lemma-quasi-coherence-higher-direct-images}。
由《上同调》引理 \ref{cohomology-lemma-Leray}，存在以
$$
E_2^{p, q} = H^p(X, R^qf_*\mathcal{G})
$$
为第二页并收敛到 $H^{p + q}(Y, \mathcal{G})$ 的谱序列。由 $X$ 上上同调的
有限性可知，只有有限多个 $E_2^{p, q}$ 非零，并且每个 $E_2^{p, q}$ 都是
有限维向量空间。由此，$E_r^{p, q}$ 对每个 $r \geq 2$ 亦如此，并且
$$
\sum (-1)^{p + q} \dim_k E_r^{p, q}
$$
与 $r$ 无关。对充分大的 $r$，有 $E_r^{p, q} = E_\infty^{p, q}$；而收敛是指
$H^n(Y, \mathcal{G})$ 上存在一个滤过，其分次片为
$E_\infty^{p, q}$，其中 $p + q = n$（这就是谱序列收敛的含义），故得出结论。
另见更一般的《同调》引理
\ref{homology-lemma-biregular-ss-relation-in-K0}。
\end{proof}

\section{射影空间}
\label{varieties-section-projective-space}

\noindent
域上的射影空间的一些结果。

\begin{lemma}
\label{varieties-lemma-projective-space-smooth}
\begin{slogan}
射影空间是光滑的。
\end{slogan}
设 $k$ 是域且 $n \geq 0$。则 $\mathbf{P}^n_k$ 是维数为 $n$ 的
$k$ 上光滑射影簇。
\end{lemma}

\begin{proof}
从略。
\end{proof}

\begin{lemma}
\label{varieties-lemma-intersection-in-affine-space}
设 $k$ 是域且 $n \geq 0$。设 $X, Y \subset \mathbf{A}^n_k$
为闭子集。假设 $X$ 和 $Y$ 都是等维的，
$\dim(X) = r$ 且 $\dim(Y) = s$。
则 $X \cap Y$ 的每个不可约分支的维数都 $\geq r + s - n$。
\end{lemma}

\begin{proof}
考虑闭子概形 $X \times Y \subset \mathbf{A}^{2n}_k$，
这里使用坐标 $x_1, \ldots, x_n, y_1, \ldots, y_n$。于是
$X \cap Y = X \times Y \cap V(x_1 - y_1, \ldots, x_n - y_n)$。
设 $t \in X \cap Y \subset X \times Y$ 为闭点。
由引理 \ref{varieties-lemma-dimension-product-locally-algebraic}，有
$\dim_t(X \times Y) = \dim(X) + \dim(Y)$。
因而由引理 \ref{varieties-lemma-dimension-locally-algebraic}，有
$\dim(\mathcal{O}_{X \times Y, t}) = r + s$。
由《代数》引理 \ref{algebra-lemma-one-equation}，可得
$$
\dim(\mathcal{O}_{X \cap Y, t}) =
\dim(\mathcal{O}_{X \times Y, t}/(x_1 - y_1, \ldots, x_n - y_n)) \geq
r + s - n
$$
再由引理 \ref{varieties-lemma-dimension-locally-algebraic} 即得结论。
\end{proof}

\begin{lemma}
\label{varieties-lemma-intersection-in-projective-space}
设 $k$ 是域且 $n \geq 0$。设 $X, Y \subset \mathbf{P}^n_k$
为非空闭子集。若 $\dim(X) = r$ 且 $\dim(Y) = s$，并且
$r + s \geq n$，则 $X \cap Y$ 非空，且
$\dim(X \cap Y) \geq r + s - n$。
\end{lemma}

\begin{proof}
记 $\mathbf{A}^n = \Spec(k[x_0, \ldots, x_n])$ 且
$\mathbf{P}^n = \text{Proj}(k[T_0, \ldots, T_n])$。
考虑态射
$\pi : \mathbf{A}^{n + 1} \setminus \{0\} \to \mathbf{P}^n$，
它把 $(x_0, \ldots, x_n)$ 送到点 $[x_0 : \ldots : x_n]$。
更准确地说，它是由二元组
$(\mathcal{O}_{\mathbf{A}^{n + 1} \setminus \{0\}}, (x_0, \ldots, x_n))$
相伴的态射；见《构造》引理 \ref{constructions-lemma-projective-space}。
在标准仿射开集 $D_+(T_i)$ 上，得到由环同态
$$
k\left[\frac{T_0}{T_i}, \ldots, \frac{T_n}{T_i}\right]
\longrightarrow
k\left[T_0, \ldots, T_n, \frac{1}{T_i}\right] \cong
k\left[\frac{T_0}{T_i}, \ldots, \frac{T_n}{T_i}\right]
\left[T_i, \frac{1}{T_i}\right]
$$
相伴的态射；该环同态是满的，并且是相对维数为 $1$、纤维不可约的光滑同态
（略去细节）。
因此 $\pi^{-1}(X)$ 和 $\pi^{-1}(Y)$ 是维数分别为
$r + 1$ 和 $s + 1$ 的非空闭子集。选取不可约分支
$V \subset \pi^{-1}(X)$，其维数为 $r + 1$；再选取不可约分支
$W \subset \pi^{-1}(Y)$，其维数为 $s + 1$。注意，这蕴涵 $V$ 和 $W$ 包含它们所遇到的
$\pi$ 的每条纤维（因为 $\pi$ 的纤维是维数为 $1$ 的不可约集，
且引理 \ref{varieties-lemma-dimension-fibres-locally-algebraic} 表明
$V \to \pi(V)$ 和 $W \to \pi(W)$ 的纤维维数都 $\geq 1$）。
令 $\overline{V}$ 和 $\overline{W}$ 分别为
$V$ 和 $W$ 在 $\mathbf{A}^{n + 1}$ 中的闭包。因为
$0 \in \mathbf{A}^{n + 1}$ 属于 $\pi$ 的每条纤维的闭包，所以
$0 \in \overline{V} \cap \overline{W}$。由引理
\ref{varieties-lemma-intersection-in-affine-space}，有
$\dim(\overline{V} \cap \overline{W}) \geq r + s - n + 1$。
再次使用引理 \ref{varieties-lemma-dimension-fibres-locally-algebraic} 如上论证，
便得出 $\pi(V \cap W) \subset X \cap Y$ 的维数至少为
$r + s - n$，如所需。
\end{proof}

\begin{lemma}
\label{varieties-lemma-equation-codim-1-in-projective-space}
设 $k$ 是域。设 $Z \subset \mathbf{P}^n_k$ 是没有嵌入点的闭子概形，
且 $Z$ 的每个不可约分支的维数都是 $n - 1$。则理想
$I(Z) \subset k[T_0, \ldots, T_n]$（与 $Z$ 对应）是主理想。
\end{lemma}

\begin{proof}
这是《除子》引理
\ref{divisors-lemma-equation-codim-1-in-projective-space} 的特殊情形。
\end{proof}

\section{射影空间上的凝聚层}
\label{varieties-section-coherent-Pn}

\noindent
本节证明域上的 $\mathbf{P}^n$ 上凝聚层之上同调的一些结果；这些结果可见
\cite{Mum}。以后讨论 Quot 概形和 Hilbert 概形时会用到它们。

\subsection{预备知识}
\label{varieties-subsection-preliminaries}

\noindent
设 $k$ 是域，$n \geq 1$，$d \geq 1$，并设
$s \in \Gamma(\mathbf{P}_k^n, \mathcal{O}(d))$
为非零截面。本节用 $\mathcal{O}(d)$ 表示射影空间上结构层的第
$d$ 次扭转（《构造》定义 \ref{constructions-definition-twist} 和
\ref{constructions-definition-projective-space}）。
由于 $\mathbf{P}^n_k$ 是簇，该截面是正则的；因而
$s$ 是 $\mathcal{O}(d)$ 的正则截面，并定义有效 Cartier 除子
$H = Z(s) \subset \mathbf{P}^n_k$；见《除子》第
\ref{divisors-section-effective-Cartier-divisors} 节。
这样的除子 $H$ 称为{it 超曲面}；若 $d = 1$，则称为{it 超平面}。

\begin{lemma}
\label{varieties-lemma-hyperplane}
设 $k$ 是域。设 $n \geq 1$。
设 $i : H \to \mathbf{P}^n_k$ 是超平面。
则存在同构
$$
\varphi : \mathbf{P}^{n - 1}_k \longrightarrow H
$$
使得 $i^*\mathcal{O}(1)$ 拉回为 $\mathcal{O}(1)$。
\end{lemma}

\begin{proof}
有 $\mathbf{P}^n_k = \text{Proj}(k[T_0, \ldots, T_n])$。
截面 $s$ 对应于 $T_0, \ldots, T_n$ 中次数为 $1$ 的齐次形式；见
《概形的上同调》第
\ref{coherent-section-cohomology-projective-space} 节。
设 $s = \sum a_i T_i$。
由《构造》引理 \ref{constructions-lemma-closed-in-projective-space}，有
$H = \text{Proj}(k[T_0, \ldots, T_n]/I)$，其中分次理想 $I$ 的定义是：
$I_d$ 等于映射
$\Gamma(\mathbf{P}^n_k, \mathcal{O}(d)) \to \Gamma(H, i^*\mathcal{O}(d))$
的核。由《除子》定义 \ref{divisors-definition-zero-scheme-s} 中
$Z(s)$ 的构造可知，在 $D_{+}(T_j)$ 上，$H$ 的理想由
$\sum a_i T_i/T_j$ 生成；这是多项式环
$k[T_0/T_j, \ldots, T_n/T_j]$ 中的元素。因此很清楚，$I$ 是由
$\sum a_i T_i$ 生成的理想。注意
$$
k[T_0, \ldots, T_n]/I = k[T_0, \ldots, T_n]/(\sum a_i T_i) \cong
k[S_0, \ldots, S_{n - 1}]
$$
为分次环同构。例如，若 $a_n \not = 0$，则令
$S_i$ 映到 $T_i$ 的类即可。由 $\text{Proj}$ 的函子性得到所需同构。
结构层扭转的相等性例如由《构造》引理
\ref{constructions-lemma-surjective-graded-rings-generated-degree-1-map-proj}
推出。
\end{proof}

\begin{lemma}
\label{varieties-lemma-exact-sequence-induction}
设 $k$ 是无限域。设 $n \geq 1$。
设 $\mathcal{F}$ 是 $\mathbf{P}^n_k$ 上的凝聚模。
则存在非零截面
$s \in \Gamma(\mathbf{P}^n_k, \mathcal{O}(1))$
和短正合列
$$
0 \to \mathcal{F}(-1) \to \mathcal{F} \to i_*\mathcal{G} \to 0
$$
其中 $i : H \to \mathbf{P}^n_k$ 是超平面 $H$，它与 $s$ 相伴，且
$\mathcal{G} = i^*\mathcal{F}$。
\end{lemma}

\begin{proof}
映射 $\mathcal{F}(-1) \to \mathcal{F}$ 来自《构造》等式
(\ref{constructions-equation-multiply-on-sheaf})，其中 $n = 1$、
$m = -1$，并取截面 $s$（属于 $\mathcal{O}(1)$）。
具体考察该映射限制到 $D_{+}(T_0)$ 后的形式。记
$D_{+}(T_0) = \Spec(k[x_1, \ldots, x_n])$，其中
$x_i = T_i/T_0$。将 $\mathcal{O}(1)|_{D_{+}(T_0)}$ 与
$\mathcal{O}$ 等同，所用截面为 $T_0$。因此，若
$s = \sum a_iT_i$，则在上述等同下
$s|_{D_{+}(T_0)} = a_0 + \sum a_ix_i$。再假设
$\mathcal{F}|_{D_{+}(T_0)}$ 对应于有限
$k[x_1, \ldots, x_n]$-模 $M$。通过等同
$\mathcal{F}(-1) = \mathcal{F} \otimes \mathcal{O}(-1)$
及所选的 $\mathcal{O}(1)$ 平凡化，可知
$\mathcal{F}(-1)$ 也对应于 $M$。所以，把
$\mathcal{F}(-1) \to \mathcal{F}$ 限制到 $D_{+}(T_0)$，得到映射
$$
M \xrightarrow{a_0 + \sum a_ix_i} M
$$
为说明该箭头是单射，只须选取 $a_0 + \sum a_ix_i$，使它不属于
$M$ 的任何相伴素理想；见《代数》引理
\ref{algebra-lemma-ass-zero-divisors}。由《代数》引理
\ref{algebra-lemma-finite-ass}，集合 $\text{Ass}(M)$（即
$M$ 的相伴素理想集）是有限的。注意，对
$\mathfrak p \in \text{Ass}(M)$，交集
$\mathfrak p \cap \{a_0 + \sum a_i x_i\}$ 是真 $k$-向量子空间。
我们得出，存在真子向量空间的有限族
$V_1, \ldots, V_m \subset \Gamma(\mathbf{P}^n_k, \mathcal{O}(1))$，
使得若取 $s$ 不属于 $\bigcup V_i$，则乘以 $s$ 在
$D_{+}(T_0)$ 上是单射。对限制到 $D_{+}(T_j)$ 的情形也同样，
其中 $j = 1, \ldots, n$。由于 $k$ 是无限域，有限个真子向量空间的并
绝不等于整个空间，故可选取 $s$ 使该映射为单射。
$\mathcal{F}(-1) \to \mathcal{F}$ 的余核被
$\Im(s : \mathcal{O}(-1) \to \mathcal{O})$ 零化；由
《除子》定义 \ref{divisors-definition-zero-scheme-s}，后者是 $H$ 的理想层。
因此，由《概形的上同调》引理
\ref{coherent-lemma-i-star-equivalence}，得到 $\mathcal{G}$，它在 $H$ 上。
\end{proof}

\begin{remark}
\label{varieties-remark-exact-sequence-induction}
设 $k$ 是无限域。设 $n \geq 1$。给定有限多个凝聚模
$\mathcal{F}_i$，它们在 $\mathbf{P}^n_k$ 上，可选取同一个
$s \in \Gamma(\mathbf{P}^n_k, \mathcal{O}(1))$，
使引理 \ref{varieties-lemma-exact-sequence-induction} 的陈述对每个模都成立。
为证明这一点，只需把该引理应用于 $\bigoplus \mathcal{F}_i$。
\end{remark}

\begin{remark}
\label{varieties-remark-exact-sequence-induction-cohomology}
在引理 \ref{varieties-lemma-hyperplane} 和
\ref{varieties-lemma-exact-sequence-induction} 的情形下，有
$H \cong \mathbf{P}^{n - 1}_k$，且 Serre 扭转满足
$\mathcal{O}_H(d) = i^*\mathcal{O}_{\mathbf{P}^n_k}(d)$。
对每个 $d \in \mathbf{Z}$，有短正合列
$$
0 \to \mathcal{F}(d - 1) \to \mathcal{F}(d) \to i_*(\mathcal{G}(d)) \to 0
$$
事实上，与 $\mathcal{O}_{\mathbf{P}^n_k}(d)$ 作张量积是正合函子，
并且由投影公式（《上同调》引理
\ref{cohomology-lemma-projection-formula}），有
$i_*(\mathcal{G}(d)) = i_*\mathcal{G} \otimes \mathcal{O}_{\mathbf{P}^n_k}(d)$。
得到相应的长正合列
$$
H^i(\mathbf{P}^n_k, \mathcal{F}(d - 1)) \to
H^i(\mathbf{P}^n_k, \mathcal{F}(d)) \to
H^i(H, \mathcal{G}(d)) \to
H^{i + 1}(\mathbf{P}^n_k, \mathcal{F}(d - 1))
$$
这是由上述事实以及下列等式推出的：由《概形的上同调》引理
\ref{coherent-lemma-relative-affine-cohomology}（闭浸入是仿射的），有
$H^i(\mathbf{P}^n_k, i_*\mathcal{G}(d)) = H^i(H, \mathcal{G}(d))$。
\end{remark}





\subsection{正则性}
\label{varieties-subsection-regularity}

\noindent
定义如下。

\begin{definition}
\label{varieties-definition-regularity}
设 $k$ 是域。设 $n \geq 0$。设 $\mathcal{F}$ 是
$\mathbf{P}^n_k$ 上的凝聚层。称 $\mathcal{F}$ 是{it $m$-正则的}，若
$$
H^i(\mathbf{P}^n_k, \mathcal{F}(m - i)) = 0
$$
对 $i = 1, \ldots, n$ 成立。
\end{definition}

\noindent
注意，$\mathcal{F} = \mathcal{O}(d)$ 是 $m$-正则的，当且仅当
$d \geq -m$。这由《概形的上同调》等式
(\ref{coherent-equation-identify}) 中的上同调群计算推出。
也就是说，$H^n(\mathbf{P}^n_k, \mathcal{O}(d)) = 0$，当且仅当
$d \geq -n$。

\begin{lemma}
\label{varieties-lemma-m-regular-extend-base-field}
设 $k'/k$ 是域扩张。设 $n \geq 0$。
设 $\mathcal{F}$ 是 $\mathbf{P}^n_k$ 上的凝聚层。
设 $\mathcal{F}'$ 是 $\mathcal{F}$ 在 $\mathbf{P}^n_{k'}$ 上的拉回。
则 $\mathcal{F}$ 是 $m$-正则的，当且仅当 $\mathcal{F}'$ 是
$m$-正则的。
\end{lemma}

\begin{proof}
这是因为由平坦基变换有
$$
H^i(\mathbf{P}^n_{k'}, \mathcal{F}') =
H^i(\mathbf{P}^n_k, \mathcal{F}) \otimes_k k'
$$
见《概形的上同调》引理
\ref{coherent-lemma-flat-base-change-cohomology}。
\end{proof}

\begin{lemma}
\label{varieties-lemma-m-regular}
在引理 \ref{varieties-lemma-exact-sequence-induction} 的情形下，若
$\mathcal{F}$ 是 $m$-正则的，则 $\mathcal{G}$ 是 $m$-正则的；
这里它在 $H \cong \mathbf{P}^{n - 1}_k$ 上。
\end{lemma}

\begin{proof}
回忆，由《概形的上同调》引理
\ref{coherent-lemma-relative-affine-cohomology}，有
$H^i(\mathbf{P}^n_k, i_*\mathcal{G}) = H^i(H, \mathcal{G})$。
因此由评注 \ref{varieties-remark-exact-sequence-induction-cohomology} 可知，
对 $i \geq 1$ 得到
$$
H^i(\mathbf{P}^n_k, \mathcal{F}(m - i)) \to
H^i(H, \mathcal{G}(m - i)) \to
H^{i + 1}(\mathbf{P}^n_k, \mathcal{F}(m - 1 - i))
$$
该引理随即成立。
\end{proof}

\begin{lemma}
\label{varieties-lemma-m-regular-up}
设 $k$ 是域。设 $n \geq 0$。
设 $\mathcal{F}$ 是 $\mathbf{P}^n_k$ 上的凝聚层。
若 $\mathcal{F}$ 是 $m$-正则的，则 $\mathcal{F}$ 是
$(m + 1)$-正则的。
\end{lemma}

\begin{proof}
对 $n$ 作归纳证明。若 $n = 0$，则每个层对所有 $m$ 都是 $m$-正则的，
无需证明。由引理 \ref{varieties-lemma-m-regular-extend-base-field}，可把 $k$
替换为无限扩域，并假设 $k$ 是无限域。
于是可应用引理 \ref{varieties-lemma-exact-sequence-induction}。
由引理 \ref{varieties-lemma-m-regular}，知道 $\mathcal{G}$ 是
$m$-正则的。对 $n$ 使用归纳假设，便知 $\mathcal{G}$ 是
$(m + 1)$-正则的。考察与下列正合列相伴的上同调长正合列
$$
0 \to \mathcal{F}(m - i) \to \mathcal{F}(m + 1 - i)
\to i_*\mathcal{G}(m + 1 - i) \to 0
$$
如评注 \ref{varieties-remark-exact-sequence-induction-cohomology} 所述，读者容易推出：
对 $i \geq 1$，$H^i(\mathbf{P}^n_k, \mathcal{F}(m + 1 - i))$ 消失；
这是由 $H^i(\mathbf{P}^n_k, \mathcal{F}(m - i))$ 和
$H^i(\mathbf{P}^n_k, \mathcal{G}(m + 1 - i))$ 的（已知）消失得出的。
\end{proof}

\begin{lemma}
\label{varieties-lemma-m-regular-multiply}
设 $k$ 是域。设 $n \geq 0$。
设 $\mathcal{F}$ 是 $\mathbf{P}^n_k$ 上的凝聚层。
若 $\mathcal{F}$ 是 $m$-正则的，则乘法映射
$$
H^0(\mathbf{P}^n_k, \mathcal{F}(m)) \otimes_k
H^0(\mathbf{P}^n_k, \mathcal{O}(1))
\longrightarrow
H^0(\mathbf{P}^n_k, \mathcal{F}(m + 1))
$$
是满射。
\end{lemma}

\begin{proof}
设 $k'/k$ 是域扩张。设 $\mathcal{F}'$ 如引理
\ref{varieties-lemma-m-regular-extend-base-field} 中所定义。由《概形的上同调》引理
\ref{coherent-lemma-flat-base-change-cohomology}，该引理的线性映射基变换到
$k'$ 后，就是层 $\mathcal{F}'$ 的同一线性映射。由于
$k \to k'$ 忠实平坦，只须在 $k'$ 上证明该引理；也就是说，可以假设
$k$ 是无限域。

\medskip\noindent
假设 $k$ 是无限域。对 $n$ 作归纳证明。情形 $n = 0$ 是平凡的，因为
$\mathcal{O}(1) \cong \mathcal{O}$ 由 $T_0$ 生成。对 $n > 0$，应用
引理 \ref{varieties-lemma-exact-sequence-induction}，并将该正合列与
$\mathcal{O}(m + 1)$ 作张量积，得到
$$
0 \to \mathcal{F}(m) \xrightarrow{s} \mathcal{F}(m + 1) \to
i_*\mathcal{G}(m + 1) \to 0
$$
见评注 \ref{varieties-remark-exact-sequence-induction-cohomology}。
设 $t \in H^0(\mathbf{P}^n_k, \mathcal{F}(m + 1))$。
由归纳假设，像 $\overline{t} \in H^0(H, \mathcal{G}(m + 1))$
是 $\sum g_i \otimes \overline{s}_i$ 的像，其中
$\overline{s}_i \in \Gamma(H, \mathcal{O}(1))$ 且
$g_i \in H^0(H, \mathcal{G}(m))$。由于 $\mathcal{F}$ 是 $m$-正则的，
有 $H^1(\mathbf{P}^n_k, \mathcal{F}(m - 1)) = 0$；因此，与短正合列
$$
0 \to \mathcal{F}(m - 1) \xrightarrow{s} \mathcal{F}(m) \to
i_*\mathcal{G}(m) \to 0
$$
相伴的上同调长正合列表明，可把 $g_i$ 提升为
$f_i \in H^0(\mathbf{P}^n_k, \mathcal{F}(m))$。
也可以把 $\overline{s}_i$ 提升为
$s_i \in H^0(\mathbf{P}^n_k, \mathcal{O}(1))$（例如见引理
\ref{varieties-lemma-hyperplane} 的证明）。减去
$\sum f_i \otimes s_i$ 的像后，可把 $t$ 取为满足
$\overline{t} = 0$。但这恰好意味着 $t$ 是某个
$f \otimes s$ 的像，其中
$f \in H^0(\mathbf{P}^n_k, \mathcal{F}(m))$，如所需。
\end{proof}

\begin{lemma}
\label{varieties-lemma-m-regular-globally-generated}
设 $k$ 是域。设 $n \geq 0$。
设 $\mathcal{F}$ 是 $\mathbf{P}^n_k$ 上的凝聚层。
若 $\mathcal{F}$ 是 $m$-正则的，则 $\mathcal{F}(m)$ 是全局生成的。
\end{lemma}

\begin{proof}
对所有 $d \gg 0$，层 $\mathcal{F}(d)$ 都是全局生成的。
这例如由《概形的上同调》引理
\ref{coherent-lemma-coherent-projective} 的第一部分推出。
选取 $d \geq m$，使 $\mathcal{F}(d)$ 全局生成。
选取一组基 $f_1, \ldots, f_r \in H^0(\mathbf{P}^n_k, \mathcal{F})$。
由引理 \ref{varieties-lemma-m-regular-multiply}，每个元素
$f \in H^0(\mathbf{P}^n_k, \mathcal{F}(d))$ 都可写成
$f = \sum P_if_i$，其中某些
$P_i \in k[T_0, \ldots, T_n]$ 是次数为 $d - m$ 的齐次多项式。
由于截面 $f$ 生成 $\mathcal{F}(d)$，可知截面 $f_i$ 生成
$\mathcal{F}(m)$。
\end{proof}

\subsection{Hilbert 多项式}
\label{varieties-subsection-hilbert}

\noindent
下面的引理将被更一般的引理
\ref{varieties-lemma-numerical-polynomial-from-euler} 所取代。

\begin{lemma}
\label{varieties-lemma-hilbert-polynomial}
设 $k$ 是域。设 $n \geq 0$。设 $\mathcal{F}$ 是
$\mathbf{P}^n_k$ 上的凝聚层。函数
$$
d \longmapsto \chi(\mathbf{P}^n_k, \mathcal{F}(d))
$$
是多项式。
\end{lemma}

\begin{proof}
对 $n$ 作归纳证明。若 $n = 0$，则
$\mathbf{P}^n_k = \Spec(k)$ 且 $\mathcal{F}(d) = \mathcal{F}$。
因此在此情形下，该函数是常数，即次数为 $0$ 的多项式。
假设 $n > 0$。由引理
\ref{varieties-lemma-euler-characteristic-extend-base-field}，可假设 $k$ 是无限域。
应用引理 \ref{varieties-lemma-exact-sequence-induction}。
把引理 \ref{varieties-lemma-euler-characteristic-additive} 应用于扭转后的正合列
$0 \to \mathcal{F}(d - 1) \to \mathcal{F}(d) \to i_*\mathcal{G}(d) \to 0$，
得到
$$
\chi(\mathbf{P}^n_k, \mathcal{F}(d)) -
\chi(\mathbf{P}^n_k, \mathcal{F}(d - 1)) =
\chi(H, \mathcal{G}(d))
$$
见评注 \ref{varieties-remark-exact-sequence-induction-cohomology}。
由于 $H \cong \mathbf{P}^{n - 1}_k$，由归纳假设，右端是多项式。
最后注意：任何函数 $f : \mathbf{Z} \to \mathbf{Z}$，若映射
$d \mapsto f(d) - f(d - 1)$ 是多项式，则其自身也是多项式；这便完成证明。
略去该事实的证明（提示：与《代数》引理
\ref{algebra-lemma-numerical-polynomial} 比较）。
\end{proof}

\begin{definition}
\label{varieties-definition-hilbert-polynomial}
设 $k$ 是域。设 $n \geq 0$。设 $\mathcal{F}$ 是
$\mathbf{P}^n_k$ 上的凝聚层。函数
$d \mapsto \chi(\mathbf{P}^n_k, \mathcal{F}(d))$ 称为
$\mathcal{F}$ 的{it Hilbert 多项式}。
\end{definition}

\noindent
Hilbert 多项式的系数属于 $\mathbf{Q}$，一般并不属于 $\mathbf{Z}$。
例如，$\mathcal{O}_{\mathbf{P}^n_k}$ 的 Hilbert 多项式为
$$
d \longmapsto {d + n \choose n} = \frac{d^n}{n!} + \ldots
$$
这由下面的引理和下列事实推出：
$$
H^0(\mathbf{P}^n_k, \mathcal{O}_{\mathbf{P}^n_k}(d)) = k[T_0, \ldots, T_n]_d
$$
是次数为 $d$ 的部分，其在 $k$ 上的维数是 ${d + n \choose n}$。

\begin{lemma}
\label{varieties-lemma-hilbert-polynomial-H0}
设 $k$ 是域。设 $n \geq 0$。设 $\mathcal{F}$ 是
$\mathbf{P}^n_k$ 上的凝聚层，其 Hilbert 多项式为
$P \in \mathbf{Q}[t]$。则
$$
P(d) = \dim_k H^0(\mathbf{P}^n_k, \mathcal{F}(d))
$$
对所有 $d \gg 0$ 成立。
\end{lemma}

\begin{proof}
这由 $\mathcal{F}$ 的充分高次扭转之上同调消失推出。见
《概形的上同调》引理
\ref{coherent-lemma-coherent-projective}。
\end{proof}




\subsection{商的有界性}
\label{varieties-subsection-boundedness}

\noindent
本小节用 Hilbert 多项式来界定 $\mathbf{P}^n$ 上给定凝聚层之商的正则性。

\begin{lemma}
\label{varieties-lemma-bound-quotients-free}
设 $k$ 是域。设 $n \geq 0$。设 $r \geq 1$。设 $P \in \mathbf{Q}[t]$。
存在整数 $m$，它依赖于 $n$、$r$ 和 $P$，并具有如下性质：若
$$
0 \to \mathcal{K} \to \mathcal{O}^{\oplus r} \to \mathcal{F} \to 0
$$
是 $\mathbf{P}^n_k$ 上凝聚层的短正合列，并且 $\mathcal{F}$ 的
Hilbert 多项式为 $P$，则 $\mathcal{K}$ 是 $m$-正则的。
\end{lemma}

\begin{proof}
对 $n$ 作归纳证明。若 $n = 0$，则
$\mathbf{P}^n_k = \Spec(k)$，任何凝聚模都是 $0$-正则的，
而任何满射在全局截面上也是满射。
假设 $n > 0$。考虑引理中的正合列。
设 $P' \in \mathbf{Q}[t]$ 为多项式
$P'(t) = P(t) - P(t - 1)$。设 $m'$ 为对
$n - 1$、$r$ 和 $P'$ 有效的整数。
由引理 \ref{varieties-lemma-m-regular-extend-base-field} 和
\ref{varieties-lemma-euler-characteristic-extend-base-field}，可把 $k$ 替换为域扩张，
故可假设 $k$ 是无限域。把引理
\ref{varieties-lemma-exact-sequence-induction} 应用于凝聚层 $\mathcal{F}$。
$\mathcal{F}' = i^*\mathcal{F}$ 的 Hilbert 多项式为
$P'$（见引理 \ref{varieties-lemma-hilbert-polynomial} 的证明）。
由于 $i^*$ 右正合，可知 $\mathcal{F}'$ 是
$\mathcal{O}_H^{\oplus r} = i^*\mathcal{O}^{\oplus r}$ 的商。
因而归纳假设适用于 $\mathcal{F}'$；它在
$H \cong \mathbf{P}^{n - 1}_k$ 上（引理 \ref{varieties-lemma-hyperplane}）。
注意，映射 $\mathcal{K}(-1) \to \mathcal{K}$ 是单射，因为
$\mathcal{K} \subset \mathcal{O}^{\oplus r}$；其余核为
$i_*\mathcal{H}$，其中 $\mathcal{H} = i^*\mathcal{K}$。
由蛇形引理（《同调》引理 \ref{homology-lemma-snake}），得到行、列都正合的交换图
$$
\xymatrix{
& 0 \ar[d] & 0 \ar[d] & 0 \ar[d] \\
0 \ar[r] &
\mathcal{K}(-1) \ar[r] \ar[d] &
\mathcal{O}^{\oplus r}(-1) \ar[r] \ar[d] &
\mathcal{F}(-1) \ar[d] \ar[r] & 0 \\
0 \ar[r] &
\mathcal{K} \ar[r] \ar[d] &
\mathcal{O}^{\oplus r} \ar[r] \ar[d] &
\mathcal{F} \ar[d] \ar[r] & 0\\
0 \ar[r] &
i_*\mathcal{H} \ar[r] \ar[d] &
i_*\mathcal{O}_H^{\oplus r} \ar[r] \ar[d] &
i_*\mathcal{F}' \ar[r] \ar[d] & 0 \\
& 0 & 0 & 0
}
$$
因此，归纳假设适用于正合列
$0 \to \mathcal{H} \to \mathcal{O}_H^{\oplus r} \to \mathcal{F}' \to 0$；
该列在 $H \cong \mathbf{P}^{n - 1}_k$ 上（引理
\ref{varieties-lemma-hyperplane}），且 $\mathcal{H}$ 是 $m'$-正则的。
回忆，这蕴涵 $\mathcal{H}$ 是 $d$-正则的，对所有 $d \geq m'$ 都如此
（引理 \ref{varieties-lemma-m-regular-up}）。

\medskip\noindent
设 $i \geq 2$ 且 $d \geq m'$。由与上图左列相伴的上同调长正合列以及
$H^{i - 1}(H, \mathcal{H}(d))$ 的消失可知，映射
$$
H^i(\mathbf{P}^n_k, \mathcal{K}(d - 1))
\longrightarrow
H^i(\mathbf{P}^n_k, \mathcal{K}(d))
$$
是单射。由于这些群对 $d \gg 0$ 为零（《概形的上同调》引理
\ref{coherent-lemma-coherent-projective}），可知
$H^i(\mathbf{P}^n_k, \mathcal{K}(d))$ 对所有 $d \geq m'$ 和
$i \geq 2$ 都为零。

\medskip\noindent
还须控制 $H^1$。首先注意，所有映射
$$
H^1(\mathbf{P}^n_k, \mathcal{K}(m' - 1)) \to
H^1(\mathbf{P}^n_k, \mathcal{K}(m')) \to
H^1(\mathbf{P}^n_k, \mathcal{K}(m' + 1)) \to \ldots
$$
都是满射，因为 $H^1(H, \mathcal{H}(d))$ 对
$d \geq m'$ 消失。假设 $d > m'$，且
$$
H^1(\mathbf{P}^n_k, \mathcal{K}(d - 1))
\longrightarrow
H^1(\mathbf{P}^n_k, \mathcal{K}(d))
$$
是单射。则
$H^0(\mathbf{P}^n_k, \mathcal{K}(d)) \to H^0(H, \mathcal{H}(d))$
是满射。考虑交换图
$$
\xymatrix{
H^0(\mathbf{P}^n_k, \mathcal{K}(d)) \otimes_k
H^0(\mathbf{P}^n_k, \mathcal{O}(1))
\ar[r] \ar[d] &
H^0(\mathbf{P}^n_k, \mathcal{K}(d + 1)) \ar[d] \\
H^0(H, \mathcal{H}(d)) \otimes_k
H^0(H, \mathcal{O}_H(1))
\ar[r] &
H^0(H, \mathcal{H}(d + 1))
}
$$
由引理 \ref{varieties-lemma-m-regular-multiply}，下方水平箭头是满射。
因此右侧竖直箭头是满射。我们得出
$$
H^1(\mathbf{P}^n_k, \mathcal{K}(d))
\longrightarrow
H^1(\mathbf{P}^n_k, \mathcal{K}(d + 1))
$$
是单射。归纳可知
$$
H^1(\mathbf{P}^n_k, \mathcal{K}(d - 1)) \to
H^1(\mathbf{P}^n_k, \mathcal{K}(d)) \to
H^1(\mathbf{P}^n_k, \mathcal{K}(d + 1)) \to \ldots
$$
全都是单射；由这些群最终消失，得出
$H^1(\mathbf{P}^n_k, \mathcal{K}(d - 1)) = 0$。
所以，群 $H^1(\mathbf{P}^n_k, \mathcal{K}(d))$ 的维数在
$d \geq m'$ 时严格递减，直至变为零。
由此，$\mathcal{K}$ 的正则性被下式所界：
$m' + \dim_k H^1(\mathbf{P}^n_k, \mathcal{K}(m'))$。
另一方面，由高阶上同调群的消失，有
$$
\dim_k H^1(\mathbf{P}^n_k, \mathcal{K}(m')) = 
- \chi(\mathbf{P}^n_k, \mathcal{K}(m')) +
\dim_k H^0(\mathbf{P}^n_k, \mathcal{K}(m'))
$$
注意，$H^0$ 的维数被
$H^0(\mathbf{P}^n_k, \mathcal{O}^{\oplus r}(m'))$ 的维数所界；
后者至多为 $r{n + m' \choose n}$（若 $m' > 0$），否则为零。
最后，项 $\chi(\mathbf{P}^n_k, \mathcal{K}(m'))$ 等于
$r{n + m' \choose n} - P(m')$。这给出了所需类型的界，从而完成引理的证明。
\end{proof}







\section{Frobenius 态射}
\label{varieties-section-frobenius}

\noindent
设 $p$ 为素数。若 $X$ 是概形，则称
``$X$ 具有特征 $p$''，或称 ``$X$ 是特征 $p$ 的''，
或称 ``$X$ 处于特征 $p$''，若 $p$ 在 $\mathcal{O}_X$ 中为零。

\begin{definition}
\label{varieties-definition-absolute-frobenius}
设 $p$ 为素数。设 $X$ 是特征 $p$ 的概形。
$X$ 的\emph{绝对 Frobenius} 是态射 $F_X : X \to X$，
它在底层拓扑空间上为恒等映射，并且其
$F_X^\sharp : \mathcal{O}_X \to \mathcal{O}_X$ 由 $g \mapsto g^p$ 给出。
\end{definition}

\noindent
这是有意义的，因为对于任意环 $A$（其特征为 $p$），映射
$F_A : A \to A$，$a \mapsto a^p$ 是一个环自同态，并且在
$\Spec(A)$ 上诱导恒等映射。此外，若 $A$ 是局部环，则 $F_A$ 是局部同态。
因此，我们看到 $X$ 的绝对 Frobenius 是 $X$ 的一个自同态，
且属于概形范畴。事实证明，绝对 Frobenius 在特征 $p$
的概形范畴上定义了恒等函子的一个自映射。

\begin{lemma}
\label{varieties-lemma-frobenius-endomorphism-identity}
设 $p > 0$ 为素数。设 $f : X \to Y$ 是特征 $p$ 的概形之间的态射。
则下图交换：
$$
\xymatrix{
X \ar[d]_f \ar[r]_{F_X} & X \ar[d]^f \\
Y \ar[r]^{F_Y} & Y
}
$$
\end{lemma}

\begin{proof}
这直接由如下平凡的代数事实得到：若
$\varphi : A \to B$ 是特征 $p$ 的环同态，则
$\varphi(a^p) = \varphi(a)^p$。
\end{proof}

\begin{lemma}
\label{varieties-lemma-frobenius}
设 $p > 0$ 为素数。设 $X$ 是特征 $p$ 的概形。
则绝对 Frobenius $F_X : X \to X$ 是泛同胚，是整的，且
诱导纯不可分的剩余域扩张。
\end{lemma}

\begin{proof}
这由 Frobenius 自同态 $F_A : A \to A$ 的相应结果得到，其中 $A$
是特征 $p > 0$ 的环。例如参见《代数》一节
\ref{algebra-section-universal-homeomorphism} 的讨论，特别是引理
\ref{algebra-lemma-p-ring-map}。
\end{proof}

\noindent
若我们在固定的底概形上工作，则有 Frobenius 态射的相对版本。

\begin{definition}
\label{varieties-definition-relative-frobenius}
设 $p > 0$ 为素数。设 $S$ 是特征 $p$ 的概形。
设 $X$ 是 $S$ 上的概形。定义
$$
X^{(p)} = X^{(p/S)} = X \times_{S, F_S} S
$$
并将其视为 $S$ 上的概形。应用引理
\ref{varieties-lemma-frobenius-endomorphism-identity}，可知存在唯一的
态射 $F_{X/S} : X \longrightarrow X^{(p)}$，它在 $S$ 上且使下图交换：
$$
\xymatrix{
X \ar[rr]_{F_{X/S}} \ar[rrd] \ar@/^1em/[rrrr]^{F_X}
& & X^{(p)} \ar[rr] \ar[d] & & X \ar[d] \\
& & S \ar[rr]^{F_S} & & S
}
$$
其中右方的方形是笛卡尔的。态射 $F_{X/S}$ 称为
$X/S$ 的\emph{相对 Frobenius 态射}。
\end{definition}

\noindent
注意，$X \mapsto X^{(p)}$ 是一个函子；它是绝对 Frobenius 态射
$F_S : S \to S$ 的基变换函子。关于相对 Frobenius 态射，
我们有与此前相同的引理。

\begin{lemma}
\label{varieties-lemma-relative-frobenius-endomorphism-identity}
设 $p > 0$ 为素数。设 $S$ 是特征 $p$ 的概形。
设 $f : X \to Y$ 是 $S$ 上概形之间的态射。
则下图交换：
$$
\xymatrix{
X \ar[d]_f \ar[r]_{F_{X/S}} & X^{(p)} \ar[d]^{f^{(p)}} \\
Y \ar[r]^{F_{Y/S}} & Y^{(p)}
}
$$
\end{lemma}

\begin{proof}
这由引理 \ref{varieties-lemma-frobenius-endomorphism-identity} 及定义得到。
\end{proof}

\begin{lemma}
\label{varieties-lemma-relative-frobenius}
设 $p > 0$ 为素数。设 $S$ 是特征 $p$ 的概形。
设 $X$ 是 $S$ 上的概形。则相对 Frobenius
$F_{X/S} : X \to X^{(p)}$ 是泛同胚，是整的，且
诱导纯不可分的剩余域扩张。
\end{lemma}

\begin{proof}
由引理 \ref{varieties-lemma-frobenius}，态射 $F_X : X \to X$ 以及
基变换 $h : X^{(p)} \to X$（它是 $F_S$ 的基变换）都是泛同胚。
由于 $h \circ F_{X/S} = F_X$，我们得到 $F_{X/S}$ 是泛同胚
（《态射》引理 \ref{morphisms-lemma-permanence-universal-homeomorphism}）。
由《态射》引理
\ref{morphisms-lemma-universal-homeomorphism} 和
\ref{morphisms-lemma-universally-injective}，我们也得到
$F_{X/S}$ 具有其余性质。
\end{proof}

\begin{lemma}
\label{varieties-lemma-relative-frobenius-omega}
设 $p > 0$ 为素数。设 $S$ 是特征 $p$ 的概形。
设 $X$ 是 $S$ 上的概形。则 $\Omega_{X/S} = \Omega_{X/X^{(p)}}$。
\end{lemma}

\begin{proof}
这转化为如下代数事实。设 $A \to B$ 是特征 $p$ 的环同态。
置 $B' = B \otimes_{A, F_A} A$，并考虑环映射
$F_{B/A} : B' \to B$，$b \otimes a \mapsto b^pa$。
我们要断言的是 $\Omega_{B/A} = \Omega_{B/B'}$。
这是正确的，因为有 $\text{d}(b^pa) = 0$；其中
$\text{d} : B \to \Omega_{B/A}$ 是泛导子，因而 $\text{d}$ 是一个 $B'$-导子。
\end{proof}

\begin{lemma}
\label{varieties-lemma-relative-frobenius-finite}
设 $p > 0$ 为素数。设 $S$ 是特征 $p$ 的概形。
设 $X$ 是 $S$ 上的概形。若 $X \to S$ 局部有限型，则 $F_{X/S}$ 是有限的。
\end{lemma}

\begin{proof}
这转化为如下代数事实。设 $A \to B$ 是特征 $p$ 的有限型环同态。
置 $B' = B \otimes_{A, F_A} A$，并考虑环映射
$F_{B/A} : B' \to B$，$b \otimes a \mapsto b^pa$。
我们要断言的是 $F_{B/A}$ 是有限的。
事实上，若 $x_1, \ldots, x_n \in B$ 是在 $A$ 上的生成元，
则 $x_i$ 在 $B'$ 上整，因为 $x_i^p = F_{B/A}(x_i \otimes 1)$。
因此由《代数》引理
\ref{algebra-lemma-characterize-finite-in-terms-of-integral}，
$F_{B/A} : B' \to B$ 是有限的。
\end{proof}

\begin{lemma}
\label{varieties-lemma-geometrically-reduced-p}
设 $k$ 是特征 $p > 0$ 的域。设 $X$ 是 $k$ 上的概形。
则 $X$ 几何约化，当且仅当 $X^{(p)}$ 约化。
\end{lemma}

\begin{proof}
考虑绝对 Frobenius $F_k : k \to k$。则 $F_k(k) = k^p$；
换言之，$F_k : k \to k$ 同构于 $k$ 到 $k^{1/p}$ 的嵌入。
因此引理由引理 \ref{varieties-lemma-geometrically-reduced} 得到。
\end{proof}

\begin{lemma}
\label{varieties-lemma-inseparable-deg-p-smooth}
设 $k$ 是特征 $p > 0$ 的域。设 $X$ 是 $k$ 上的簇。
以下条件等价：
\begin{enumerate}
\item $X^{(p)}$ 约化；
\item $X$ 几何约化；
\item 存在非空开子集 $U \subset X$，其到 $k$ 光滑。
\end{enumerate}
在此情形，$X^{(p)}$ 是 $k$ 上的簇，且 $F_{X/k} : X \to X^{(p)}$
是次数为 $p^{\dim(X)}$ 的有限主导态射。
\end{lemma}

\begin{proof}
我们已经在引理 \ref{varieties-lemma-geometrically-reduced-p} 中证明了
(1) 与 (2) 等价。我们已经在引理
\ref{varieties-lemma-geometrically-reduced-dense-smooth-open} 中证明了
(2) 蕴含 (3)。若 (3) 成立，则 $U$ 几何约化
（例如见引理 \ref{varieties-lemma-geometrically-regular-smooth}），
从而由引理 \ref{varieties-lemma-generic-points-geometrically-reduced}，
$X$ 几何约化。由此可见 (1)、(2) 与 (3) 等价。

\medskip\noindent
假设 (1)、(2) 与 (3) 成立。由于 $F_{X/k}$ 是同胚
（引理 \ref{varieties-lemma-relative-frobenius}），我们看到 $X^{(p)}$ 是一个簇。
随后由引理 \ref{varieties-lemma-relative-frobenius-finite}，$F_{X/k}$ 是有限的。
它是主导的，因为它是满射。为计算次数
（《态射》定义 \ref{morphisms-definition-degree}），
只需计算 $F_{U/k} : U \to U^{(p)}$ 的次数
（因为由引理 \ref{varieties-lemma-relative-frobenius-endomorphism-identity}，
$F_{U/k} = F_{X/k}|_U$）。
缩小一点 $U$ 后，不妨设存在 'etale 态射
$h : U \to \mathbf{A}^n_k$，见《态射》引理
\ref{morphisms-lemma-smooth-etale-over-affine-space}。
当然 $n = \dim(U)$，因为 $\mathbf{A}^n_k \to \Spec(k)$
是相对维数为 $n$ 的光滑态射，'etale 态射 $h$ 是相对维数为 $0$ 的光滑态射，
而 $U \to \Spec(k)$ 是相对维数为 $\dim(U)$ 的光滑态射，且相对维数正确相加
（《态射》引理 \ref{morphisms-lemma-composition-relative-dimension-d}）。
注意，$h$ 是簇之间的一般有限主导态射，因而 $\deg(h)$ 有定义。
由引理 \ref{varieties-lemma-relative-frobenius-endomorphism-identity}，有交换图
$$
\xymatrix{
U \ar[rr]_{F_{X/k}} \ar[d]_h & & U^{(p)} \ar[d]^{h^{(p)}} \\
\mathbf{A}^n_k \ar[rr]^{F_{\mathbf{A}^n_k/k}} & &
(\mathbf{A}^n_k)^{(p)}
}
$$
由于 $h^{(p)}$ 是 $h$ 的基变换，它也是 'etale 的，
因此 $h^{(p)}$ 也是簇之间的一般有限主导态射。$h^{(p)}$ 的次数是扩张
$k(X^{(p)})/k((\mathbf{A}^n_k)^{(p)})$ 的次数；
该次数与扩张 $k(X)/k(\mathbf{A}^n_k)$ 的次数相同，因为 $h^{(p)}$ 是 $h$ 的基变换
（省略细节）。
由次数的乘法性（《态射》引理
\ref{morphisms-lemma-degree-composition}），只需证明
$F_{\mathbf{A}^n_k/k}$ 的次数为 $p^n$。为此注意
$(\mathbf{A}^n_k)^{(p)} = \mathbf{A}^n_k$，且
$F_{\mathbf{A}^n_k/k}$ 由将坐标映射为其 $p$ 次幂的映射给出。
\end{proof}

\begin{remark}
\label{varieties-remark-n-fold-relative-frobenius}
设 $p > 0$ 为素数。设 $S$ 是特征 $p$ 的概形。
设 $X$ 是 $S$ 上的概形。对 $n \geq 1$，
$$
X^{(p^n)} = X^{(p^n/S)} = X \times_{S, F_S^n} S
$$
并将其视为 $S$ 上的概形。注意，$X \mapsto X^{(p^n)}$
是一个函子。应用引理 \ref{varieties-lemma-frobenius-endomorphism-identity}，
可知 $F_{X/S, n} = (F_X^n, \text{id}_S) : X \longrightarrow X^{(p^n)}$
是一个在 $S$ 上的态射，并且适合下列交换图：
$$
\xymatrix{
X \ar[rr]_{F_{X/S, n}} \ar[rrd] \ar@/^1em/[rrrr]^{F_X^n}
& & X^{(p^n)} \ar[rr] \ar[d] & & X \ar[d] \\
& & S \ar[rr]^{F_S^n} & & S
}
$$
其中右方的方形是笛卡尔的。态射 $F_{X/S, n}$ 有时称为
\emph{$n$ 重相对 Frobenius 态射}（$X/S$）。
这是有意义的，因为我们有公式
$$
F_{X/S, n} =
F_{X^{(p^{n - 1})}/S} \circ \ldots \circ F_{X^{(p)}/S} \circ F_{X/S}
$$
这表明 $F_{X/S, n}$ 是 $n$ 个相对 Frobenius 的复合。由于
$$
F_{X^{(p^m)}/S} = F_{X^{(p^{m - 1})}/S}^{(p)} = \ldots = F_{X/S}^{(p^m)}
$$
（细节略去），我们还得到
$$
F_{X/S, n} =
F_{X/S}^{(p^{n - 1})} \circ \ldots \circ F_{X/S}^{(p)} \circ F_{X/S}
$$
\end{remark}




\section{粘合一维环}
\label{varieties-section-glueing-dim-1}

\noindent
本节给出一些代数预备，用于证明：分离概形的有限个余维 $1$ 点包含于某个仿射开集。

\begin{situation}
\label{varieties-situation-glue}
给定环的交换图
$$
\xymatrix{
A \ar[r] & K \\
R \ar[u] \ar[r] & B \ar[u]
}
$$
其中 $K$ 是域，$A$、$B$ 是 $K$ 的子环，且分式域都是 $K$。最后，
$R = A \times_K B = A \cap B$。
\end{situation}

\begin{lemma}
\label{varieties-lemma-glue-valuation-ring}
在情形 \ref{varieties-situation-glue} 中，假设 $B$ 是赋值环。
那么对每个单元 $u$（来自 $A$），或者 $u \in R$，或者 $u^{-1} \in R$。
\end{lemma}

\begin{proof}
事实上，若 $c$（$u$ 在 $K$ 中的像）属于 $B$，则 $u \in R$。
否则 $c^{-1} \in B$
（《代数》引理 \ref{algebra-lemma-valuation-ring-x-or-x-inverse}），
从而 $u^{-1} \in R$。
\end{proof}

\noindent
下面的引理解释稍后反复出现的条件
“$A \otimes B \to K$ 是满射”的含义。

\begin{lemma}
\label{varieties-lemma-glue-separated}
在情形 \ref{varieties-situation-glue} 中，假设 $A$ 是维数为 $1$ 的 Noether 环。
以下条件等价：
\begin{enumerate}
\item $A \otimes B \to K$ 不是满射；
\item 存在离散赋值环 $\mathcal{O} \subset K$，它包含 $A$ 与 $B$。
\end{enumerate}
\end{lemma}

\begin{proof}
(2) 蕴含 (1) 是显然的。另一方面，若 $A \otimes B \to K$
不是满射，则其像 $C \subset K$ 不是域，因而 $C$ 有非零极大理想
$\mathfrak m$。取 $\mathcal{O} \subset K$ 为支配 $C_\mathfrak m$ 的赋值环。
由《代数》引理
\ref{algebra-lemma-krull-akizuki} 应用于 $A \subset \mathcal{O}$，
环 $\mathcal{O}$ 是 Noether 环。因此，由《代数》引理
\ref{algebra-lemma-valuation-ring-Noetherian-discrete}，
$\mathcal{O}$ 是离散赋值环。
\end{proof}

\begin{lemma}
\label{varieties-lemma-semi-local}
在情形 \ref{varieties-situation-glue} 中，假设
\begin{enumerate}
\item $A$ 是维数为 $1$ 的 Noether 半局部整环；
\item $B$ 是离散赋值环；
\end{enumerate}
则有以下两种可能：
\begin{enumerate}
\item[(a)] 若 $A^*$ 不包含于 $R$，则
$\Spec(A) \to \Spec(R)$ 与 $\Spec(B) \to \Spec(R)$
都是覆盖 $\Spec(R)$ 的开浸入，并且 $K = A \otimes_R B$；
\item[(b)] 若 $A^*$ 包含于 $R$，则 $B$ 支配 $A$ 在某个极大理想处的
局部环，并且 $A \otimes B \to K$ 不是满射。
\end{enumerate}
\end{lemma}

\begin{proof}
假设 (a) 成立，则存在一个单元 $u$，它属于 $A$，其在 $K$ 中的像属于 $B$ 的极大理想。
于是 $u$ 是 $R$ 中的非零因子，并且对每个 $a \in A$，存在某个 $n$ 使
$u^n a \in R$。因此 $A = R_u$。

\medskip\noindent
令 $\mathfrak m_A$ 为 $A$ 的 Jacobson 根。取非零元
$x \in \mathfrak m_A$。由于 $\dim(A) = 1$，可知 $K = A_x$。
将 $x$ 替换为某个 $x^n u^m$（其中 $n \geq 1$ 且 $m \in \mathbf{Z}$）后，
可假设 $x$ 映到 $B$ 的单元。于是对每个 $b \in B$，都有 $x^nb$
属于 $R$ 的像，对某个 $n$ 成立。因此 $B = R_x$。

\medskip\noindent
设 $z \in R$。若 $z \not \in \mathfrak m_A$ 且 $z$ 不映到
$\mathfrak m_B$ 中的元素，则 $z$ 可逆。因此 $x + u$ 在 $R$ 中可逆。
于是 $\Spec(R) = D(x) \cup D(u)$。上面已经看到
$D(u) = \Spec(A)$ 且 $D(x) = \Spec(B)$。

\medskip\noindent
情形 (b)。若 $x \in \mathfrak m_A$，则 $1 + x$ 是单元，因而
$1 + x \in R$，即 $x \in R$。所以
$\mathfrak m_A \subset R \subset A$。事实上，在此情形 $A$ 在 $R$ 上整。
将
$A/\mathfrak m_A = \kappa_1 \times \ldots \times \kappa_n$
写成域的乘积。设
$x = (c_1, \ldots, c_r, 0, \ldots, 0)$
是一个其中 $c_i \not = 0$ 的元素。则
$$
x^2 - x(c_1, \ldots, c_r, 1, \ldots, 1)  = 0
$$
由于 $R$ 包含所有单元，可知 $A/\mathfrak m_A$ 在 $R$ 在其中的像上整，
从而 $A$ 在 $R$ 上整。由此得到 $R \subset A \subset B$，因为 $B$ 整闭。
此外，若非零元 $x \in \mathfrak m_A$，则
$K = A_x = \bigcup x^{-n}A = \bigcup x^{-n}R$。
因此 $x^{-1} \not \in B$，即 $x \in \mathfrak m_B$。
我们得到 $\mathfrak m_A \subset \mathfrak m_B$。
于是 $A \cap \mathfrak m_B$ 是 $A$ 的极大理想，证明完成。
\end{proof}

\begin{lemma}
\label{varieties-lemma-semi-local-dimension-one-conductor}
设 $B$ 是维数为 $1$ 的 Noether 半局部整环。
设 $B'$ 是 $B$ 在其分式域中的整闭包。
则 $B'$ 是半局部 Dedekind 整环。
设 $x$ 是 $B'$ 的 Jacobson 根中的非零元。
则对每个 $y \in B'$，存在某个 $n$ 使得 $x^n y \in B$。
\end{lemma}

\begin{proof}
令 $\mathfrak m_B$ 为 $B$ 的 Jacobson 根。
$B'$ 的结构由《代数》引理
\ref{algebra-lemma-integral-closure-Dedekind} 给出。
给定引理陈述中的 $x, y \in B'$，考虑子环
$B \subset A \subset B'$，它由 $x$ 与 $y$ 生成。则 $A$ 在 $B$ 上有限（《代数》引理
\ref{algebra-lemma-characterize-finite-in-terms-of-integral}）。
由于 $B$ 与 $A$ 的分式域相同，可知有限模 $A/B$ 支集于 $B$ 的闭点集。
因此有 $\mathfrak m_B^n A \subset B$，其中 $n$ 适当。
此外，$\Spec(B') \to \Spec(A)$ 是满射（《代数》引理
\ref{algebra-lemma-integral-overring-surjective}），所以 $A$ 也是半局部的。
同样可知 $x$ 属于 Jacobson 根 $\mathfrak m_A$，这是 $A$ 的 Jacobson 根。
注意 $\mathfrak m_A = \sqrt{\mathfrak m_B A}$。
因此有 $x^m y \in \mathfrak m_B A$，其中 $m$ 适当。
于是 $x^{nm} y \in B$。
\end{proof}

\begin{lemma}
\label{varieties-lemma-semi-local-both-side}
在情形 \ref{varieties-situation-glue} 中，假设
\begin{enumerate}
\item $A$ 是维数为 $1$ 的 Noether 半局部整环；
\item $B$ 是维数为 $1$ 的 Noether 半局部整环；
\item $A \otimes B \to K$ 是满射。
\end{enumerate}
则 $\Spec(A) \to \Spec(R)$ 与 $\Spec(B) \to \Spec(R)$
是覆盖 $\Spec(R)$ 的开浸入，并且 $K = A \otimes_R B$。
\end{lemma}

\begin{proof}
特殊情形：$B$ 在 $K$ 中整闭。这意味着 $B$ 是 Dedekind 整环
（《代数》引理 \ref{algebra-lemma-characterize-Dedekind}），
所以它在极大理想处的所有局部化都是离散赋值环。
令 $\mathfrak m_1, \ldots, \mathfrak m_r$ 为 $B$ 的极大理想。置
$$
R_1 = A \times_K B_{\mathfrak m_1}
$$
注意到 $A \otimes_{R_1} B_{\mathfrak m_1} \to K$ 是满射，
由引理 \ref{varieties-lemma-semi-local} 可知，$A$ 与
$B_{\mathfrak m_1}$ 定义覆盖 $\Spec(R_1)$ 的开子概形，并且
$K = A \otimes_{R_1} B_{\mathfrak m_1}$。特别地，$R_1$ 是维数为 $1$ 的
Noether 半局部环。归纳定义
$$
R_{i + 1} = R_i \times_K B_{\mathfrak m_{i + 1}}
$$
其中 $i = 1, \ldots, r - 1$。注意，$R = R_r$，因为
$B = B_{\mathfrak m_1} \cap \ldots \cap B_{\mathfrak m_r}$
（见《代数》引理
\ref{algebra-lemma-normal-domain-intersection-localizations-height-1}）。
由这一归纳过程，$R \to A$ 定义开浸入
$\Spec(A) \to \Spec(R)$。另一方面，极大理想
$\mathfrak n_i$（属于 $R$ 且不在这个开集内）对应于极大理想
$\mathfrak m_i$（属于 $B$）；事实上环同态
$R \to B$ 定义同构
$R_{\mathfrak n_i} \to B_{\mathfrak m_i}$（细节略去；提示：
在每一步中我们恰好向 $\Spec(R_i)$ 添加了一个极大理想）。
所以 $\Spec(B) \to \Spec(R)$ 是开浸入，如所需。

\medskip\noindent
一般情形。令 $B' \subset K$ 为 $B$ 的整闭包。
见引理 \ref{varieties-lemma-semi-local-dimension-one-conductor}。
于是特殊情形适用于 $R' = A \times_K B'$。
取 $x \in R'$，使它不属于 $A$ 的极大理想，而属于 $B'$ 的极大理想
（见《代数》引理 \ref{algebra-lemma-chinese-remainder}）。
由引理 \ref{varieties-lemma-semi-local-dimension-one-conductor}，
存在整数 $n$ 使得 $x^n \in R = A \times_K B$。
将 $x$ 替换为 $x^n$，于是 $x \in R$。对每个 $y \in R'$，
存在整数 $n$ 使得 $x^n y \in R$。另一方面，显然 $R'_x = A$。
因此 $R_x = A$。
交换 $A$ 与 $B$ 的角色，还可找到 $y \in R$ 使得 $B = R_y$。
注意，同时对 $x$ 与 $y$ 取逆后只剩下 $(0)$ 这一素理想。
因此 $K = R_{xy} = R_x \otimes_R R_y$。
证明完成。
\end{proof}

\begin{lemma}
\label{varieties-lemma-glue-a-bunch-of-local-rings}
设 $K$ 为域。设 $A_1, \ldots, A_r \subset K$ 是维数为 $1$ 的
Noether 半局部环，其分式域为 $K$。若
$A_i \otimes A_j \to K$ 对所有 $i \not = j$ 都是满射，
则存在 Noether 半局部整环 $A \subset K$，其维数为 $1$，并且它包含于
$A_1, \ldots, A_r$，并且满足：
\begin{enumerate}
\item $A \to A_i$ 诱导开浸入 $j_i : \Spec(A_i) \to \Spec(A)$；
\item $\Spec(A)$ 是各开集 $j_i(\Spec(A_i))$ 的并；
\item $\Spec(A)$ 的每个闭点恰好属于这些开集中的一个。
\end{enumerate}
\end{lemma}

\begin{proof}
事实上，可以取 $A = A_1 \cap \ldots \cap A_r$。首先注意，一旦证明了
(1) 和 (2)，则 (3) 由假设 $A_i \otimes A_j \to K$ 是满射推出，
因为若
$\mathfrak m \in j_i(\Spec(A_i)) \cap j_j(\Spec(A_j))$，则
$A_i \otimes A_j \to K$ 的像落在 $A_\mathfrak m$ 中。
为证明 (1) 和 (2)，对 $r$ 归纳。
若 $r > 1$，由归纳假设，(1) 和 (2) 对
$B = A_2 \cap \ldots \cap A_r$ 成立。再应用引理
\ref{varieties-lemma-semi-local-both-side}，可知它们对
$A = A_1 \cap B$ 也成立。
\end{proof}

\begin{lemma}
\label{varieties-lemma-create-globally-generated}
设 $A$ 为分式域是 $K$ 的整环。设 $B_1, \ldots, B_r \subset K$
为 Noether $1$-维半局部整环，其分式域均是 $K$。若
$A \otimes B_i \to K$ 对 $i = 1, \ldots, r$ 都是满射，
则存在 $x \in A$，使得 $x^{-1}$ 属于 $B_i$ 的 Jacobson 根，
对 $i = 1, \ldots, r$ 都成立。
\end{lemma}

\begin{proof}
设 $B_i'$ 为 $B_i$ 在 $K$ 中的整闭包。假设能找到非零元
$x \in A$，使得 $x^{-1}$ 属于 $B'_i$ 的 Jacobson 根，
对 $i = 1, \ldots, r$ 都成立。则由引理
\ref{varieties-lemma-semi-local-dimension-one-conductor}，
将 $x$ 替换为某个幂后，得到 $x^{-1} \in B_i$。
由于 $\Spec(B'_i) \to \Spec(B_i)$ 是满射，可知
$x^{-1}$ 也属于 $B_i$ 的 Jacobson 根。因此可以假设每个
$B_i$ 都是半局部 Dedekind 整环。

\medskip\noindent
若 $B_i$ 不是局部环，则从列表中删去 $B_i$，并加入有限个局部环
$(B_i)_\mathfrak m$。因此可以假设 $B_i$ 是离散赋值环，
对 $i = 1, \ldots, r$ 都成立。

\medskip\noindent
设 $v_i : K \to \mathbf{Z}$（$i = 1, \ldots, r$）为相应的离散赋值
（见《代数》引理 \ref{algebra-lemma-characterize-Dedekind}）。
我们要找非零元 $x \in A$，使得 $v_i(x) < 0$，
对 $i = 1, \ldots, r$ 都成立。对此对 $r$ 归纳。

\medskip\noindent
若 $r = 1$ 且结论不成立，则 $A \subset B$，并且映射
$A \otimes B \to K$ 不是满射，矛盾。

\medskip\noindent
若 $r > 1$，则由归纳假设，可以找到非零元 $x \in A$，
使得 $v_i(x) < 0$ 对 $i = 1, \ldots, r - 1$ 都成立。
若 $v_r(x) < 0$，则已经完成；所以可假设 $v_r(x) \geq 0$。
由基本情形，可找到非零元 $y \in A$，使得 $v_r(y) < 0$。
将 $x$ 替换为某个幂后，可假设 $v_i(x) < v_i(y)$，
对 $i = 1, \ldots, r - 1$ 都成立。于是 $x + y$ 就是所需元素。
\end{proof}

\begin{lemma}
\label{varieties-lemma-power-equal}
设 $A$ 是维数为 $1$ 的 Noether 局部环。
令 $L = \prod A_\mathfrak p$，其中乘积取遍 $A$ 的极小素理想。
设 $a_1, a_2 \in \mathfrak m_A$ 映到 $L$ 中的同一元素。
则有 $a_1^n = a_2^n$，对某个 $n > 0$ 成立。
\end{lemma}

\begin{proof}
写成 $a_1 = a_2 + x$。则 $x$ 在 $L$ 中的像为零。因此 $x$ 是
$A$ 的幂零元，因为 $\bigcap \mathfrak p$ 是 $(0)$ 的根基，而
零化理想 $I$（即 $x$ 的零化理想）包含极大理想的某个幂，这是由于对所有极小素理想都有
$\mathfrak p \not \in V(I)$。设 $x^k = 0$ 且
$\mathfrak m^n \subset I$。则
\begingroup
\small
$$
a_1^{k + n} = a_2^{k + n} + {n + k \choose 1} a_2^{n + k - 1} x +
{n + k \choose 2} a_2^{n + k - 2} x^2 + \ldots +
{n + k \choose k - 1} a_2^{n + 1} x^{k - 1} = a_2^{n + k}
$$
\endgroup
因为 $a_2 \in \mathfrak m_A$。
\end{proof}

\begin{lemma}
\label{varieties-lemma-power-works}
设 $A$ 是维数为 $1$ 的 Noether 局部环。
令 $L = \prod A_\mathfrak p$ 且 $I = \bigcap \mathfrak p$，
其中乘积与交均取遍 $A$ 的极小素理想。设 $f \in L$ 是形如
$f = i + a$ 的元素，其中 $a \in \mathfrak m_A$ 且
$i \in IL$。则 $f$ 的某个幂属于 $A \to L$ 的像。
\end{lemma}

\begin{proof}
由于 $A$ 是 Noether 环，有 $I^t = 0$，对某个 $t > 0$ 成立。
假设已知 $f = a + i$，其中 $i \in I^kL$。
则 $f^n = a^n + na^{n - 1}i \bmod I^{k + 1}L$。
所以只需证明：$na^{n - 1}i$ 属于 $I^k \to I^kL$ 的像，
对某个 $n \gg 0$ 成立。
为此，取 $g \in A$，使得 $\mathfrak m_A = \sqrt{(g)}$
（《代数》引理 \ref{algebra-lemma-height-1}）。
于是，例如由《代数》命题
\ref{algebra-proposition-dimension-zero-ring} 可知 $L = A_g$。
另一方面，存在某个 $n$ 使得 $a^n \in (g)$。
因此，可以清除 $L$ 中元素的分母，方法是乘以 $a$ 的高次幂。
\end{proof}

\begin{lemma}
\label{varieties-lemma-good-intersection}
设 $A$ 是维数为 $1$ 的 Noether 局部环。
令 $L = \prod A_\mathfrak p$，其中乘积取遍 $A$ 的极小素理想。
设 $K \to L$ 为整的环同态。
则存在 $a \in \mathfrak m_A$ 与 $x \in K$，它们映到
$L$ 中的同一元素，并且 $\mathfrak m_A = \sqrt{(a)}$。
\end{lemma}

\begin{proof}
由引理 \ref{varieties-lemma-power-works}，可以把 $A$ 替换为
$A/(\bigcap \mathfrak p)$，并假设 $A$ 是约化环
（略去一些细节）。
还可以把 $K$ 替换为 $K \to L$ 的像。
于是 $K$ 是约化环。映射 $\Spec(L) \to \Spec(K)$ 是满射且是闭映射
（细节略去）。因此 $\Spec(K)$ 是有限离散空间。
所以 $K$ 是有限多个域的积。

\medskip\noindent
令 $\mathfrak p_j$（$j = 1, \ldots, m$）为 $A$ 的极小素理想。
令 $L_j$ 为 $A_j$ 的分式域，从而
$L = \prod_{j = 1, \ldots, m} L_j$。
令 $A_j$ 为 $A/\mathfrak p_j$ 的正规化。则
$A_j$ 是至少有一个极大理想的半局部 Dedekind 整环，见《代数》引理
\ref{algebra-lemma-integral-closure-Dedekind}。
令 $n$ 为 $A_1, \ldots, A_m$ 的极大理想个数之和。
对这样的极大理想 $\mathfrak m \subset A_j$，考虑函数
$$
v_{\mathfrak m} : L \to L_j \to \mathbf{Z} \cup \{\infty\}
$$
其中第二个箭头是与离散赋值环 $(A_j)_{\mathfrak m}$ 相应的离散赋值，
并通过把 $0$ 映到 $\infty$ 来扩张。这样得到 $n$ 个函数
$v_1, \ldots, v_n : L \to \mathbf{Z} \cup \{\infty\}$。
我们将找出 $x \in K$，使得 $v_i(x) < 0$，
对所有 $i = 1, \ldots, n$ 成立。

\medskip\noindent
首先断言：对每个 $i$，存在 $x \in K$ 使得 $v_i(x) < 0$。
事实上，假设 $v_i$ 对应于 $\mathfrak m \subset A_j$。
若 $v_i(x) \geq 0$ 对所有 $x \in K$ 都成立，则 $K$ 映入
$(A_j)_{\mathfrak m}$，这是分式域 $L_j$（$A_j$ 的分式域）中的子环。
$K$ 在 $L_j$ 中的像是域；$L_j$ 在该域上是代数的，这是由《代数》引理
\ref{algebra-lemma-integral-under-field} 得到的。
结合起来，这与《代数》引理
\ref{algebra-lemma-valuation-ring-cap-field-finite} 矛盾。

\medskip\noindent
假设已经找到 $x \in K$，使得
$v_1(x) < 0, \ldots, v_r(x) < 0$，其中某个 $r < n$。
若 $v_{r + 1}(x) < 0$，则 $x$ 对 $r + 1$ 也适用。
否则，由上一段的结果，可以选取某个 $y \in K$，使
$v_{r + 1}(y) < 0$。将 $x$ 替换为某个 $x^n$（其中 $n > 0$）后，
可以假设 $v_i(x) < v_i(y)$ 对 $i = 1, \ldots, r$ 成立。
于是由赋值的性质，
$v_j(x + y) = v_j(x) < 0$ 对 $j = 1, \ldots, r$ 成立；类似地，
$v_{r + 1}(x + y) = v_{r + 1}(y) < 0$。
归纳可得 $x \in K$，使得 $v_i(x) < 0$
对 $i = 1, \ldots, n$ 成立。

\medskip\noindent
特别地，元素 $x \in K$ 在 $K$ 的每个因子中的投影均非零
（回忆 $K$ 是有限多个域的积；如果 $x$ 的某个分量为零，
则某个值 $v_i(x)$ 会是 $\infty$）。因此 $x$ 可逆，而
$x^{-1} \in K$ 满足
$\infty > v_i(x^{-1}) > 0$，对所有 $i$ 都成立。由引理
\ref{varieties-lemma-semi-local-dimension-one-conductor}，可知对某个
$e < 0$，元素 $x^e \in K$ 映到
$\mathfrak m_A/\mathfrak p_j \subset A/\mathfrak p_j$ 中，
对所有 $j = 1, \ldots, m$ 都成立。
注意，映射
$\mathfrak m_A \to \prod \mathfrak m_A/\mathfrak p_j$
的余核被 $\mathfrak m_A$ 的某个幂零化。因此，把 $e$ 替换为一个
更小的 $e$ 后，可找到 $a \in \mathfrak m_A$，使它在
$\mathfrak m_A/\mathfrak p_j$ 中的像等于 $x^e$ 的像。
于是二元组 $(a, x^e)$ 满足引理的结论。
\end{proof}

\begin{lemma}
\label{varieties-lemma-localization-semi-local}
设 $A$ 是环。设 $\mathfrak p_1, \ldots, \mathfrak p_r$
是 $A$ 的有限个素理想。令
$S = A \setminus \bigcup \mathfrak p_i$。则 $S$ 是乘法系，
而 $S^{-1}A$ 是半局部环，其极大理想对应于集合
$\{\mathfrak p_i\}$ 的极大元。
\end{lemma}

\begin{proof}
若 $a, b \in A$ 且 $a, b \in S$，则
$a, b \not \in \mathfrak p_i$，从而
$ab \not \in \mathfrak p_i$，所以 $ab \in S$。又有 $1 \in S$。
因此 $S$ 是 $A$ 的乘法子集。由《代数》引理
\ref{algebra-lemma-spec-localization} 中对 $\Spec(S^{-1}A)$ 的描述，
以及《代数》引理 \ref{algebra-lemma-silly}，可知
$S^{-1}A$ 的素理想对应于 $A$ 中包含于某个
$\mathfrak p_i$ 的素理想。
所以 $S^{-1}A$ 的极大理想一一对应于集合
$\{\mathfrak p_1, \ldots, \mathfrak p_r\}$ 中关于包含关系的极大元。
\end{proof}

\section{一维 Noether 概形}
\label{varieties-section-dimension-one}

\noindent
本节的主要结果是：维数为 $1$ 的 Noether 分离概形具有
丰沛可逆层。见命题
\ref{varieties-proposition-dim-1-noetherian-separated-has-ample}。

\begin{lemma}
\label{varieties-lemma-affine}
设 $X$ 是概形，其所有局部环都是维数
$\leq 1$ 的 Noether 环。设 $U \subset X$ 是逆紧开集，并以
$j : U \to X$ 表示包含态射。则
$R^pj_*\mathcal{F} = 0$，其中 $p > 0$；
这对每个拟凝聚 $\mathcal{O}_U$-模 $\mathcal{F}$ 都成立。
\end{lemma}

\begin{proof}
可以在茎上检验 $R^pj_*\mathcal{F}$ 的消失性。
$R^qj_*$ 的形成与平坦基变换可交换，见《概形上同调》引理
\ref{coherent-lemma-flat-base-change-cohomology}。
因此可假设 $X$ 是维数 $\leq 1$ 的 Noether 局部环的谱。
在此情形，$X$ 有一个闭点 $x$，以及有限个其他点
$x_1, \ldots, x_n$；它们特化到 $x$，但彼此之间不发生特化
（见《代数》引理
\ref{algebra-lemma-Noetherian-irreducible-components}）。
若 $x \in U$，则 $U = X$，结论显然。否则，在必要时重编号各点后，
有 $U = \{x_1, \ldots, x_r\}$，对某个 $r$ 成立。
于是 $U$ 是仿射的
（《概形》引理 \ref{schemes-lemma-scheme-finite-discrete-affine}）。
因此结论由《概形上同调》引理
\ref{coherent-lemma-relative-affine-vanishing} 得到。
\end{proof}

\begin{lemma}
\label{varieties-lemma-open-in-affine-curve-affine}
设 $X$ 是仿射概形，其所有局部环都是维数 $\leq 1$ 的 Noether 环。
则任意拟紧开集 $U \subset X$ 都是仿射的。
\end{lemma}

\begin{proof}
以 $j : U \to X$ 表示包含态射。设 $\mathcal{F}$
是拟凝聚 $\mathcal{O}_U$-模。
由引理 \ref{varieties-lemma-affine}，高阶正像
$R^pj_*\mathcal{F}$ 均为零。$\mathcal{O}_X$-模 $j_*\mathcal{F}$
是拟凝聚的
（《概形》引理 \ref{schemes-lemma-push-forward-quasi-coherent}）。
因此，由《概形上同调》引理
\ref{coherent-lemma-quasi-coherent-affine-cohomology-zero}，
它的高阶上同调群均为零。
由 Leray 谱序列（《上同调》引理
\ref{cohomology-lemma-apply-Leray}），有
$H^p(U, \mathcal{F}) = 0$，对所有 $p > 0$ 都成立。
所以 $U$ 是仿射的，例如可由《概形上同调》引理
\ref{coherent-lemma-quasi-compact-h1-zero-covering} 得到。
\end{proof}

\begin{lemma}
\label{varieties-lemma-complement-codim-1-closed-points}
设 $X$ 是概形，$U \subset X$ 是开集。假设
\begin{enumerate}
\item $U$ 是 $X$ 的逆紧开集；
\item $X \setminus U$ 是离散的；
\item 对 $x \in X \setminus U$，局部环
$\mathcal{O}_{X, x}$ 是维数 $\leq 1$ 的 Noether 环。
\end{enumerate}
则 (1) 存在可逆 $\mathcal{O}_X$-模 $\mathcal{L}$
以及一个截面 $s$，使得 $U = X_s$；并且 (2) 映射
$\Pic(X) \to \Pic(U)$ 是满射。
\end{lemma}

\begin{proof}
令 $X \setminus U = \{x_i; i \in I\}$。
选择仿射开集 $U_i \subset X$，使得 $x_i \in U_i$ 且
$x_j \not \in U_i$，其中 $j \not = i$。这由条件 (2) 可以做到。
写 $U_i = \Spec(A_i)$。令 $\mathfrak m_i \subset A_i$
为与 $x_i$ 对应的极大理想。由对局部环的假设，只有有限多个素理想
$\mathfrak q \subset \mathfrak m_i$ 满足
$\mathfrak q \not = \mathfrak m_i$（见《代数》引理
\ref{algebra-lemma-Noetherian-irreducible-components}）。
所以，由素理想回避（《代数》引理
\ref{algebra-lemma-silly}），可以找到 $f_i \in \mathfrak m_i$，
使它不属于上述任何素理想。于是
$V(f_i) = \{\mathfrak m_i\} \amalg Z_i$，其中某个闭子集
$Z_i \subset U_i$；这是因为 $Z_i$ 是
$V(f_i)$ 中逆紧且对特化封闭的开子集，见《代数》引理
\ref{algebra-lemma-constructible-stable-specialization-closed}。
缩小 $U_i$ 后，可假设 $V(f_i) = \{x_i\}$。于是
$$
\mathcal{U} : X = U \cup \bigcup U_i
$$
是 $X$ 的开覆盖。考虑 $2$-余循环，它取值于
$\mathcal{O}_X^*$：由 $f_i$ 在 $U \cap U_i$ 上给出，并由
$f_i/f_j$ 在 $U_i \cap U_j$ 上给出。它定义一个线丛
$\mathcal{L}$；而截面 $s$ 由 $1$ 在 $U$ 上定义、由 $f_i$
在 $U_i$ 上定义，它满足引理的陈述。

\medskip\noindent
设 $\mathcal{N}$ 是可逆 $\mathcal{O}_U$-模。
设 $N_i$ 为可逆 $(A_i)_{f_i}$-模，使得
$\mathcal{N}|_{U \cap U_i}$ 等于 $\tilde N_i$。
注意 $(A_{\mathfrak m_i})_{f_i}$ 是 Artin 环
（因为它是零维 Noether 环，见《代数》引理
\ref{algebra-lemma-Noetherian-dimension-0}）。
所以它是局部环的乘积
（《代数》引理 \ref{algebra-lemma-artinian-finite-length}），
因而其 Picard 群平凡。因此，缩小 $U_i$
（即把 $A_i$ 替换为某个 $(A_i)_g$，其中 $g \in A_i$ 且
$g \not \in \mathfrak m_i$）后，
可假设 $N_i = (A_i)_{f_i}$，即
$\mathcal{N}|_{U \cap U_i}$ 平凡。在这种情形，如何将
$\mathcal{N}$ 扩张为 $X$ 上的可逆层是显然的
（在每个 $U_i$ 上用平凡可逆模扩张）。
\end{proof}

\begin{lemma}
\label{varieties-lemma-find-globally-generated}
设 $X$ 是整分离概形。设 $U \subset X$ 是非空仿射开集，使得
$X \setminus U$ 是有限点集
$x_1, \ldots, x_r$，其中 $\mathcal{O}_{X, x_i}$
是维数为 $1$ 的 Noether 环。
则存在全局生成的可逆 $\mathcal{O}_X$-模
$\mathcal{L}$ 以及一个截面 $s$，使得 $U = X_s$。
\end{lemma}

\begin{proof}
写 $U = \Spec(A)$，并令 $K$ 为 $X$ 的函数域。
写 $B_i = \mathcal{O}_{X, x_i}$ 且
$\mathfrak m_i = \mathfrak m_{x_i}$。
由于 $x_i \not \in U$，可知开集
$U \times_X \Spec(B_i)$（属于 $\Spec(B_i)$）只有一个点，即
$U \times_X \Spec(B_i) = \Spec(K)$。
由于 $X$ 是分离的，可知 $\Spec(K)$ 是
$U \times \Spec(B_i)$ 的闭子概形，即映射
$A \otimes B_i \to K$ 是满射。
由引理 \ref{varieties-lemma-create-globally-generated}，可以找到非零元
$f \in A$，使得 $f^{-1} \in \mathfrak m_i$ 对
$i = 1, \ldots, r$ 都成立。
选择开集 $x_i \in U_i \subset X$，使得
$f^{-1} \in \mathcal{O}(U_i)$。于是
$$
\mathcal{U} : X = U \cup \bigcup U_i
$$
是 $X$ 的开覆盖。考虑 $2$-余循环，它取值于
$\mathcal{O}_X^*$：由 $f$ 在 $U \cap U_i$ 上给出，并由
$1$ 在 $U_i \cap U_j$ 上给出。它定义一个线丛
$\mathcal{L}$，并且该线丛有两个截面：
\begin{enumerate}
\item 截面 $s$ 由 $1$ 在 $U$ 上定义，由
$f^{-1}$ 在 $U_i$ 上定义；它满足引理的陈述；
\item 截面 $t$ 由 $f$ 在 $U$ 上定义，并由
$1$ 在 $U_i$ 上定义。
\end{enumerate}
注意 $X_t \supset U_1 \cup \ldots \cup U_r$。
因此 $s, t$ 生成 $\mathcal{L}$，引理得证。
\end{proof}

\begin{lemma}
\label{varieties-lemma-enough-globally-generated-ample}
设 $X$ 是拟紧概形。若对每个 $x \in X$ 都存在二元组
$(\mathcal{L}, s)$，其中 $\mathcal{L}$ 是全局生成的可逆层，
$s$ 是全局截面，且 $x \in X_s$、$X_s$ 是仿射的，
则 $X$ 具有丰沛可逆层。
\end{lemma}

\begin{proof}
由于 $X$ 拟紧，可以找到有限多个二元组
$(\mathcal{L}_i, s_i)$（$i = 1, \ldots, n$），使得
$\mathcal{L}_i$ 全局生成，$X_{s_i}$ 是仿射的，并且
$X = \bigcup X_{s_i}$。
再由 $X$ 拟紧，对每个 $i$ 可以找到有限多个截面
$t_{i, j}$，它们属于 $\mathcal{L}_i$（$j = 1, \ldots, m_i$），使得
$X = \bigcup X_{t_{i, j}}$。置 $t_{i, 0} = s_i$。
考虑可逆层
$$
\mathcal{L} = \mathcal{L}_1 
\otimes_{\mathcal{O}_X} \ldots
\otimes_{\mathcal{O}_X} \mathcal{L}_n
$$
及其全局截面
$$
\tau_J = t_{1, j_1} \otimes \ldots \otimes t_{n, j_n}
$$
由《性质》引理 \ref{properties-lemma-affine-cap-s-open}，
开集 $X_{\tau_J}$ 在只要有 $j_i = 0$（对某个 $i$）时就是仿射的。
不难看出这些开集覆盖 $X$。所以按定义，$\mathcal{L}$ 是丰沛的。
\end{proof}

\begin{lemma}
\label{varieties-lemma-dim-1-noetherian-integral-separated-has-ample}
设 $X$ 是维数为 $1$ 的 Noether 整分离概形。
则 $X$ 具有丰沛可逆层。
\end{lemma}

\begin{proof}
选择仿射开覆盖 $X = U_1 \cup \ldots \cup U_n$。
由于 $X$ 是 Noether 的，每个集合 $X \setminus U_i$ 都是有限的。
因此，由引理 \ref{varieties-lemma-find-globally-generated}，
可以找到二元组 $(\mathcal{L}_i, s_i)$，
其中 $\mathcal{L}_i$ 是全局生成的可逆层，
$s_i$ 是全局截面，且 $U_i = X_{s_i}$。
由引理 \ref{varieties-lemma-enough-globally-generated-ample}，
得到 $X$ 具有丰沛可逆层。
\end{proof}

\begin{lemma}
\label{varieties-lemma-surjective-pic-birational-finite}
设 $f : X \to Y$ 是概形的有限态射。假设存在开集
$V \subset Y$，使得 $f^{-1}(V) \to V$ 是同构，并且
$Y \setminus V$ 是离散空间。则每个可逆
$\mathcal{O}_X$-模都是某个可逆 $\mathcal{O}_Y$-模的拉回。
\end{lemma}

\begin{proof}
我们将使用 $\Pic(X) = H^1(X, \mathcal{O}_X^*)$，见《上同调》引理
\ref{cohomology-lemma-h1-invertible}。
考虑阿贝尔层 $\mathcal{O}_X^*$ 与 $f$ 的 Leray 谱序列，
见《上同调》引理 \ref{cohomology-lemma-Leray}。
考虑诱导映射
$$
H^1(X, \mathcal{O}_X^*) \longrightarrow H^0(Y, R^1f_*\mathcal{O}_X^*)
$$
《除子》引理
\ref{divisors-lemma-finite-trivialize-invertible-upstairs}
恰好说明此映射为零。因此 Leray 谱序列给出
$H^1(X, \mathcal{O}_X^*) = H^1(Y, f_*\mathcal{O}_X^*)$。
接着考虑映射
$$
f^\sharp : \mathcal{O}_Y^* \longrightarrow f_*\mathcal{O}_X^*
$$
由假设，此映射的核与余核都支集于闭子集
$T = Y \setminus V$（属于 $Y$）上。
因为按假设 $T$ 是离散拓扑空间，所以 $Y$ 上支集于
$T$ 的任意阿贝尔层的高阶上同调群都为零
（这由《上同调》引理
\ref{cohomology-lemma-cohomology-and-closed-immersions}、
《模》引理 \ref{modules-lemma-i-star-exact}，以及如下事实推出：
$H^i(T, \mathcal{F}) = 0$，对任意 $i > 0$ 与任意阿贝尔层
$\mathcal{F}$（在 $T$ 上）都成立）。
把上述映射分解为短正合列
$$
0 \to \Ker(f^\sharp) \to \mathcal{O}_Y^* \to \Im(f^\sharp) \to 0,\quad
0 \to \Im(f^\sharp) \to f_*\mathcal{O}_X^* \to \Coker(f^\sharp) \to 0
$$
先得到 $H^1(Y, \mathcal{O}_Y^*) \to H^1(Y, \Im(f^\sharp))$
是满射，再得到
$H^1(Y, \Im(f^\sharp)) \to H^1(Y, f_*\mathcal{O}_X^*)$ 是满射。
结合以上各点，得到
$H^1(Y, \mathcal{O}_Y^*) \to
H^1(X, \mathcal{O}_X^*)$ 是满射，如所需。
\end{proof}

\begin{lemma}
\label{varieties-lemma-glue-invertible-sheaves}
设 $X$ 是概形。设 $Z_1, \ldots, Z_n \subset X$ 为闭子概形。
设 $\mathcal{L}_i$ 为 $Z_i$ 上的可逆层。假设
\begin{enumerate}
\item $X$ 是约化的；
\item 在集合论意义下 $X = \bigcup Z_i$；
\item $Z_i \cap Z_j$ 是离散拓扑空间，其中 $i \not = j$。
\end{enumerate}
则存在可逆层 $\mathcal{L}$（在 $X$ 上），其在 $Z_i$ 上的限制是
$\mathcal{L}_i$。此外，若给定截面
$s_i \in \Gamma(Z_i, \mathcal{L}_i)$，它们在
$Z_i \cap Z_j$ 的各点处不消失，则可选择 $\mathcal{L}$，使得存在
$s \in \Gamma(X, \mathcal{L})$，并且有
$s|_{Z_i} = s_i$，对所有 $i$ 都成立。
\end{lemma}

\begin{proof}
$\mathcal{L}$ 的存在性可以由引理
\ref{varieties-lemma-surjective-pic-birational-finite} 推出；
但我们也将给出直接证明，并用该证明说明关于截面的陈述成立。
置 $T = \bigcup_{i \not = j} Z_i \cap Z_j$。由于 $X$ 约化，作为概形有
$$
X \setminus T = \bigcup (Z_i \setminus T)
$$
。假设 (3) 蕴含 $T$ 是 $X$ 的离散子集。
所以对每个 $t \in T$，可以找到开集 $U_t \subset X$，
使得 $t \in U_t$，但 $t' \not \in U_t$；这里
$t' \in T$ 且 $t' \not = t$。
必要时缩小 $U_t$，可假设存在同构
$\varphi_{t, i} : \mathcal{L}_i|_{U_t \cap Z_i} \to
\mathcal{O}_{U_t \cap Z_i}$。此外，对每个 $i$ 选择开覆盖
$$
Z_i \setminus T = \bigcup\nolimits_j U_{ij}
$$
使得存在同构
$\varphi_{i, j} : \mathcal{L}_i|_{U_{ij}} \cong \mathcal{O}_{U_{ij}}$。
注意
$$
\mathcal{U} : X = \bigcup U_t \cup \bigcup U_{ij}
$$
是 $X$ 的开覆盖。我们断言，可以用这些同构
$\varphi_{t, i}$ 和 $\varphi_{i, j}$ 为此覆盖定义一个
$2$-余循环，它取值于 $\mathcal{O}_X^*$，并定义引理陈述中的
$\mathcal{L}$。

\medskip\noindent
事实上，若 $i \not = i'$，则
$U_{i, j} \cap U_{i', j'} = \emptyset$，无事可做。
对 $U_{i, j} \cap U_{i, j'}$，由证明开头的说明，有
$\mathcal{O}_X(U_{i, j} \cap U_{i, j'}) =
\mathcal{O}_{Z_i}(U_{i, j} \cap U_{i, j'})$。
所以 $\mathcal{L}_i$ 的转移函数（更准确地说，
$\varphi_{i, j} \circ \varphi_{i, j'}^{-1}$）
定义了余循环在这一交集上的值。
对 $U_t \cap U_{i, j}$ 可以作同样处理。
最后，当 $t \not = t'$ 时，有
$$
U_t \cap U_{t'} = \coprod (U_t \cap U_{t'}) \cap Z_i
$$
，而且交集 $U_t \cap U_{t'} \cap Z_i$ 包含于
$Z_i \setminus T$。所以，由与前面相同的理由，有
$$
\mathcal{O}_X(U_t \cap U_{t'}) =
\prod \mathcal{O}_{Z_i}(U_t \cap U_{t'} \cap Z_i)
$$
，并且可以用 $\mathcal{L}_i$ 的转移函数（更准确地说，
$\varphi_{t, i} \circ \varphi_{t', i}^{-1}$）定义余循环在
$U_t \cap U_{t'}$ 上的值。这就完成了 $\mathcal{L}$ 的存在性证明。

\medskip\noindent
给定引理最后一个断言中的截面 $s_i$。在上述论证中，选择
$U_t$，使得 $s_i|_{U_t \cap Z_i}$ 不消失，并选择
$\varphi_{t, i}$，使得
$\varphi_{t, i}(s_i|_{U_t \cap Z_i}) = 1$。
于是，使用 $1$（在 $U_t$ 上）以及
$\varphi_{i, j}(s_i|_{U_{i, j}})$（在 $U_{i, j}$ 上），
将定义 $\mathcal{L}$ 的一个截面，
它限制为 $s_i$（在 $Z_i$ 上）。
\end{proof}

\begin{remark}
\label{varieties-remark-conductor}
设 $A$ 是约化环。设 $I, J$ 是 $A$ 的理想，并且
$V(I) \cup V(J) = \Spec(A)$。置 $B = A/J$。
则 $I \to IB$ 是 $A$-模的同构。事实上，有
$IB = I + J/J = I/(I \cap J)$；而 $I \cap J$ 为零，因为
$A$ 约化，并且
$\Spec(A) = V(I) \cup V(J) = V(I \cap J)$。
所以对任意投射 $A$-模 $P$，也有 $IP = I(P/JP)$。
\end{remark}

\begin{lemma}
\label{varieties-lemma-dim-1-noetherian-reduced-separated-has-ample}
设 $X$ 是维数为 $1$ 的 Noether 约化分离概形。
则 $X$ 具有丰沛可逆层。
\end{lemma}

\begin{proof}
令 $Z_i$（$i = 1, \ldots, n$）为 $X$ 的不可约分支。
把它们视为 $X$ 的约化闭子概形。
由引理 \ref{varieties-lemma-dim-1-noetherian-integral-separated-has-ample}，
存在丰沛可逆层 $\mathcal{L}_i$（在 $Z_i$ 上）。
置 $T = \bigcup_{i \not = j} Z_i \cap Z_j$。由于 $X$ 是维数
$1$ 的 Noether 概形，集合 $T$ 是有限的，并由 $X$ 的闭点组成。
对每个 $i$，在必要时把 $\mathcal{L}_i$ 替换为某个幂后，
可以选择 $s_i \in \Gamma(Z_i, \mathcal{L}_i)$，使得
$(Z_i)_{s_i}$ 是仿射的并包含 $T \cap Z_i$，见《性质》引理
\ref{properties-lemma-ample-finite-set-in-principal-affine}。

\medskip\noindent
由引理 \ref{varieties-lemma-glue-invertible-sheaves}，可以找到可逆层
$\mathcal{L}$（在 $X$ 上）以及 $s \in \Gamma(X, \mathcal{L})$，使得
$(\mathcal{L}, s)|_{Z_i} = (\mathcal{L}_i, s_i)$。
注意 $X_s$ 包含 $T$，而且在集合论意义下等于这些仿射闭子概形
$(Z_i)_{s_i}$ 的并。因此，由《极限》引理
\ref{limits-lemma-affines-glued-in-closed-affine}，它是仿射的。
为完成证明，只需对每个 $x \in X$（其中 $x \not \in T$）
找到整数 $m > 0$ 和截面
$t \in \Gamma(X, \mathcal{L}^{\otimes m})$，使得 $X_t$ 是仿射的且
$x \in X_t$。由于 $x \not \in T$，可知 $x \in Z_i$，其中
$i$ 唯一；设 $i = 1$。
令 $Z \subset X$ 为约化闭子概形，其底拓扑空间是
$Z_2 \cup \ldots \cup Z_n$。
令 $\mathcal{I} \subset \mathcal{O}_X$ 为 $Z$ 的理想层。
以 $\mathcal{I}_1 \subset \mathcal{O}_{Z_1}$ 表示该理想层在包含态射
$Z_1 \to X$ 下的逆像。注意
$$
\Gamma(X, \mathcal{I}\mathcal{L}^{\otimes m}) =
\Gamma(Z_1, \mathcal{I}_1 \mathcal{L}_1^{\otimes m})
$$
，见注 \ref{varieties-remark-conductor}。所以只需找到 $m > 0$ 以及
$t \in \Gamma(Z_1, \mathcal{I}_1 \mathcal{L}_1^{\otimes m})$，
使得 $x \in (Z_1)_t$ 且它是仿射的。
由于 $\mathcal{L}_1$ 丰沛，并且 $x$ 不在
$Z_1 \cap T = V(\mathcal{I}_1)$ 中，可以找到截面
$t_1 \in \Gamma(Z_1, \mathcal{I}_1 \mathcal{L}_1^{\otimes m_1})$，
使得 $x \in (Z_1)_{t_1}$，见《性质》命题
\ref{properties-proposition-characterize-ample}。
由于 $\mathcal{L}_1$ 丰沛，可以找到截面
$t_2 \in \Gamma(Z_1, \mathcal{L}_1^{\otimes m_2})$，
使得 $x \in (Z_1)_{t_2}$ 且 $(Z_1)_{t_2}$ 是仿射的，见《性质》定义
\ref{properties-definition-ample}。
置 $m = m_1 + m_2$ 且 $t = t_1 t_2$。则由构造，
$t \in \Gamma(Z_1, \mathcal{I}_1 \mathcal{L}_1^{\otimes m})$，
且 $x \in (Z_1)_t$；并且由《性质》引理
\ref{properties-lemma-affine-cap-s-open}，$(Z_1)_t$ 是仿射的。
\end{proof}

\begin{lemma}
\label{varieties-lemma-lift-line-bundle-from-reduction-dimension-1}
设 $i : Z \to X$ 是概形的闭浸入。
若 $X$ 的底拓扑空间是 Noether 的，并且
$\dim(X) \leq 1$，则 $\Pic(X) \to \Pic(Z)$ 是满射。
\end{lemma}

\begin{proof}
考虑 $X$ 上阿贝尔层的短正合列
$$
0 \to (1 + \mathcal{I}) \cap \mathcal{O}_X^* \to
\mathcal{O}^*_X \to i_*\mathcal{O}^*_Z \to 0
$$
，其中 $\mathcal{I}$ 是与 $Z$ 对应的拟凝聚理想层。
由于 $\dim(X) \leq 1$，有
$H^2(X, \mathcal{F}) = 0$；这对任意阿贝尔层 $\mathcal{F}$ 成立，见《上同调》命题
\ref{cohomology-proposition-vanishing-Noetherian}。
因此映射
$H^1(X, \mathcal{O}^*_X) \to H^1(X, i_*\mathcal{O}_Z^*)$
是满射。由《上同调》引理
\ref{cohomology-lemma-cohomology-and-closed-immersions}，有
$H^1(X, i_*\mathcal{O}_Z^*) = H^1(Z, \mathcal{O}_Z^*)$。
再由《上同调》引理
\ref{cohomology-lemma-h1-invertible}，这就证明了引理。
\end{proof}

\begin{proposition}
\label{varieties-proposition-dim-1-noetherian-separated-has-ample}
设 $X$ 是维数为 $1$ 的 Noether 分离概形。
则 $X$ 具有丰沛可逆层。
\end{proposition}

\begin{proof}
令 $Z \subset X$ 为 $X$ 的约化。由引理
\ref{varieties-lemma-dim-1-noetherian-reduced-separated-has-ample}，
概形 $Z$ 具有丰沛可逆层。
所以，由引理 \ref{varieties-lemma-lift-line-bundle-from-reduction-dimension-1}，
存在可逆 $\mathcal{O}_X$-模 $\mathcal{L}$（在 $X$ 上），
它在 $Z$ 上的限制是丰沛的。
于是，应用《概形上同调》引理
\ref{coherent-lemma-ample-on-reduction}，可知 $\mathcal{L}$ 是丰沛的。
\end{proof}

\begin{remark}
\label{varieties-remark-useless-generalization}
事实上，若 $X$ 是概形，其约化是维数为 $1$ 的 Noether 分离概形，
则 $X$ 具有丰沛可逆层。证明这一点的论证与命题
\ref{varieties-proposition-dim-1-noetherian-separated-has-ample} 的证明相同，
只是使用《极限》引理
\ref{limits-lemma-ample-on-reduction}
来代替《概形上同调》引理
\ref{coherent-lemma-ample-on-reduction}。
\end{remark}

\noindent
下面的引理实际上对拟有限分离态射也成立；读者可以使用 Zariski 主定理
（《关于态射的更多内容》引理
\ref{more-morphisms-lemma-quasi-finite-separated-pass-through-finite}）
以及引理 \ref{varieties-lemma-complement-codim-1-closed-points} 看出这一点。

\begin{lemma}
\label{varieties-lemma-surjection-on-pic-quasi-finite}
设 $f : X \to Y$ 是概形的态射。假设
$Y$ 是维数 $\leq 1$ 的 Noether 概形，$f$ 是有限的，并且
存在稠密开集 $V \subset Y$，使得
$f^{-1}(V) \to V$ 是闭浸入。则每个可逆
$\mathcal{O}_X$-模都是某个可逆 $\mathcal{O}_Y$-模的拉回。
\end{lemma}

\begin{proof}
把 $f$ 分解为 $X \to Z \to Y$，其中 $Z$ 是 $f$ 的概形论像。
则 $X \to Z$ 在 $V \cap Z$ 上是同构，并可应用引理
\ref{varieties-lemma-surjective-pic-birational-finite}。另一方面，
引理 \ref{varieties-lemma-lift-line-bundle-from-reduction-dimension-1}
适用于 $Z \to Y$。略去一些细节。
\end{proof}

\section{\texorpdfstring{\ensuremath{\delta}}{delta} 不变量}
\label{varieties-section-delta-invariant}

\noindent
本节定义奇异点的 $\delta$-不变量，其中该点位于约化
$1$ 维 Nagata 概形上。

\begin{lemma}
\label{varieties-lemma-pre-pre-delta-invariant}
设 $(A, \mathfrak m)$ 是 Noether $1$ 维局部环。
设 $f \in \mathfrak m$。下列条件等价：
\begin{enumerate}
\item $\mathfrak m = \sqrt{(f)}$；
\item $f$ 不包含在 $A$ 的任何极小素理想中；
\item 有 $A_f = \prod_{\mathfrak p\text{ minimal}} A_\mathfrak p$（作为 $A$-代数）。
\end{enumerate}
这样的 $f \in \mathfrak m$ 总是存在。若 $\text{depth}(A) = 1$（例如
$A$ 约化），则 (1)--(3) 还等价于：
\begin{enumerate}
\item[(4)] $f$ 是非零因子；
\item[(5)] $A_f$ 是 $A$ 的总分式环。
\end{enumerate}
若 $A$ 约化，则 (1)--(5) 还等价于：
\begin{enumerate}
\item[(6)] $A_f$ 是 $A$ 在各极小素理想处的剩余域之积。
\end{enumerate}
\end{lemma}

\begin{proof}
$A$ 的谱有有限多个素理想
$\mathfrak p_1, \ldots, \mathfrak p_n$ 位于 $\mathfrak m$ 之外，
并且它们全都极小；见
《代数》引理 \ref{algebra-lemma-Noetherian-irreducible-components}。
于是 (1) 与 (2) 的等价性由《代数》引理
\ref{algebra-lemma-Zariski-topology} 得出。
显然，(3) 蕴含 (2)。反之，若 (2) 成立，则
$A_f$ 的谱是子集 $\{\mathfrak p_1, \ldots, \mathfrak p_n\}$，
它位于 $\Spec(A)$ 中并赋予诱导拓扑；见
《代数》引理 \ref{algebra-lemma-spec-localization}。
这是有限离散拓扑空间。
故由《代数》命题 \ref{algebra-proposition-dimension-zero-ring}，
$A_f = \prod_{\mathfrak p\text{ minimal}} A_\mathfrak p$。
《代数》引理 \ref{algebra-lemma-height-1} 断言了这样的 $f$ 的存在性。

\medskip\noindent
假设 $A$ 的深度为 $1$。（由《代数》引理
\ref{algebra-lemma-bound-depth}，这是最大可能值；若 $A$ 约化，
则由《代数》引理 \ref{algebra-lemma-criterion-reduced} 可知它成立。）
于是 $\mathfrak m$ 不是 $A$ 的相伴素理想。
$A$ 的每个极小素理想都是相伴素理想
（《代数》命题
\ref{algebra-proposition-minimal-primes-associated-primes}）。
因此，由《代数》引理 \ref{algebra-lemma-ass-zero-divisors}，
$A$ 的非零因子集合恰为不包含在任何极小素理想中的元素集合。
所以 (4) 等价于 (2)。
由《代数》引理
\ref{algebra-lemma-total-ring-fractions-no-embedded-points}，
(5) 等价于 (3)。

\medskip\noindent
$A_\mathfrak p$ 是域，其中 $\mathfrak p \subset A$ 极小；见
《代数》引理 \ref{algebra-lemma-minimal-prime-reduced-ring}。
故 (3) 等价于 (6)。
\end{proof}

\begin{lemma}
\label{varieties-lemma-pre-delta-invariant}
设 $(A, \mathfrak m)$ 是约化 Nagata $1$ 维局部环。
设 $A'$ 是 $A$ 在 $A$ 的总分式环中的整闭包。则
$A'$ 是正规 Nagata 环，$A \to A'$ 有限，并且
$A'/A$ 作为 $A$-模具有有限长度。
\end{lemma}

\begin{proof}
总分式环在 $A$ 上本质有限型；因此 $A \to A'$ 有限，因为
$A$ 是 Nagata 环；见《代数》引理
\ref{algebra-lemma-nagata-in-reduced-finite-type-finite-integral-closure}。
例如，由《代数》引理
\ref{algebra-lemma-characterize-reduced-ring-normal} 和
\ref{algebra-lemma-Noetherian-irreducible-components} 可知环 $A'$ 正规。
例如，由《代数》引理 \ref{algebra-lemma-quasi-finite-over-nagata}
可知环 $A'$ 是 Nagata 环。
按引理 \ref{varieties-lemma-pre-pre-delta-invariant} 选取 $f \in \mathfrak m$。
由于 $A' \subset A_f$，显然 $A_f = A'_f$。因此有限
$A$-模 $A'/A$ 的支撑包含在 $\{\mathfrak m\}$ 中。
由《代数》引理 \ref{algebra-lemma-support-point}，它具有有限长度。
\end{proof}

\begin{definition}
\label{varieties-definition-delta-invariant-algebra}
设 $A$ 是维数为 $1$ 的约化 Nagata 局部环。
{\it $\delta$-不变量（关于 $A$）}定义为 $\text{length}_A(A'/A)$，
其中 $A'$ 如引理 \ref{varieties-lemma-pre-delta-invariant} 所述。
\end{definition}

\noindent
下面证明关于这个不变量之行为的一些引理。

\begin{lemma}
\label{varieties-lemma-delta-invariant-is-zero}
设 $A$ 是维数为 $1$ 的约化 Nagata 局部环。
$\delta$-不变量（关于 $A$）为 $0$，当且仅当
$A$ 是离散赋值环。
\end{lemma}

\begin{proof}
若 $A$ 是离散赋值环，则 $A$ 正规，并且环
$A'$ 等于 $A$。反之，若
$\delta$-不变量（关于 $A$）为 $0$，则 $A$ 在其总分式环中整闭，从而
$A$ 正规
（《代数》引理 \ref{algebra-lemma-characterize-reduced-ring-normal}）；
再由《代数》引理 \ref{algebra-lemma-characterize-dvr}，
这迫使 $A$ 成为离散赋值环。
\end{proof}

\begin{lemma}
\label{varieties-lemma-normalization-same-after-completion}
设 $A$ 是维数为 $1$ 的约化 Nagata 局部环。
令 $A \to A'$ 如引理 \ref{varieties-lemma-pre-delta-invariant} 所述。
令 $A^h$、$A^{sh}$ 以及相应的 $A^\wedge$
分别为 $A$ 的 Hensel 化、严格 Hensel 化以及完备化。
则 $A^h$、$A^{sh}$ 以及相应的 $A^\wedge$ 都是约化 Nagata
局部环，维数为 $1$，并且
$A' \otimes_A A^h$、$A' \otimes_A A^{sh}$ 以及相应的 $A' \otimes_A A^\wedge$
分别是 $A^h$、$A^{sh}$ 以及相应的 $A^\wedge$
在其总分式环中的整闭包。
\end{lemma}

\begin{proof}
注意 $A^\wedge$ 约化；见《代数续篇》引理
\ref{more-algebra-lemma-completion-reduced}。
由《代数续篇》引理 \ref{more-algebra-lemma-henselization-reduced}，
环 $A^h$ 与 $A^{sh}$ 约化。
由《代数续篇》引理
\ref{more-algebra-lemma-completion-dimension} 和
\ref{more-algebra-lemma-henselization-dimension}，
$A$、$A^h$、$A^{sh}$ 与 $A^\wedge$ 的维数相同。

\medskip\noindent
回顾：Noether 局部环是 Nagata 环，当且仅当
$A$ 的形式纤维几何约化；见《代数续篇》引理
\ref{more-algebra-lemma-Nagata-local-ring}。
这个性质由 $A^h$ 与 $A^{sh}$ 继承；见《代数续篇》第
\ref{more-algebra-section-properties-formal-fibres} 节中的材料，
尤其是引理 \ref{more-algebra-lemma-henselization-P-ring}。
由《代数》引理 \ref{algebra-lemma-Noetherian-complete-local-Nagata}，
完备化是 Nagata 环。

\medskip\noindent
现在证明关于整闭包的陈述。继续之前，按引理
\ref{varieties-lemma-pre-pre-delta-invariant} 选取 $f \in \mathfrak m$。
$f$ 在 $A^h$、$A^{sh}$ 与 $A^\wedge$ 中的像
显然都在相应环中满足性质 (1)--(6)。

\medskip\noindent
由于 $A \to A'$ 有限，可知
$A' \otimes_A A^h$ 与 $A' \otimes_A A^{sh}$
分别是若干 Hensel 局部环之积，而这些局部环分别在 $A^h$ 与
$A^{sh}$ 上有限；见《代数》引理
\ref{algebra-lemma-finite-over-henselian}。
其中每个局部环都是 $A'$ 在某个极大理想
$\mathfrak m' \subset A'$ 处的 Hensel 化，该理想位于 $\mathfrak m$
上方；见《代数》引理
\ref{algebra-lemma-quasi-finite-henselization} 或
\ref{algebra-lemma-quasi-finite-strict-henselization}。
故由《代数续篇》引理
\ref{more-algebra-lemma-henselization-normal}，这些局部环都是正规整环。
于是 $A' \otimes_A A^h$ 与 $A' \otimes_A A^{sh}$
是正规环。由于 $A^h \to A' \otimes_A A^h$ 与
$A^{sh} \to A' \otimes_A A^{sh}$ 有限（因而整），并且
$A' \otimes_A A^h \subset (A^h)_f = Q(A^h)$ 与
$A' \otimes_A A^{sh} \subset (A^{sh})_f = Q(A^{sh})$，
我们断定 $A' \otimes_A A^h$ 与 $A' \otimes_A A^{sh}$
就是所需的整闭包。

\medskip\noindent
对于完备化，论证完全相同。首先，由《代数》引理
\ref{algebra-lemma-completion-finite-extension} 有
$$
A' \otimes_A A^\wedge = (A')^\wedge =
\prod\nolimits (A'_{\mathfrak m'})^\wedge
$$
局部环 $A'_{\mathfrak m'}$ 正规且维数为 $1$
（例如由《代数》引理 \ref{algebra-lemma-finite-in-codim-1}，或由
《代数》第 \ref{algebra-section-homomorphism-dimension} 节中的讨论）。
故由《代数》引理 \ref{algebra-lemma-characterize-dvr}，
$A'_{\mathfrak m'}$ 是离散赋值环。
再由《代数续篇》引理 \ref{more-algebra-lemma-completion-dvr}，
$(A'_{\mathfrak m'})^\wedge$ 是离散赋值环。
于是 $A' \otimes_A A^\wedge$ 是正规环，并且
可以完全按前面的方式得出结论。
\end{proof}

\begin{lemma}
\label{varieties-lemma-delta-same-after-completion}
设 $A$ 是维数为 $1$ 的约化 Nagata 局部环。
$\delta$-不变量（关于 $A$）与 $\delta$-不变量（关于
$A$ 的 Hensel 化、严格 Hensel 化或完备化）相同。
\end{lemma}

\begin{proof}
我们证明完备化 $B = A^\wedge$ 的情形；其他情形的证明完全相同。
令 $A'$ 以及相应的 $B'$ 分别为 $A$ 以及相应的 $B$
在其总分式环中的整闭包。
由引理 \ref{varieties-lemma-normalization-same-after-completion}，
$B' = A' \otimes_A B$。
故 $B'/B = (A'/A) \otimes_A B$。
现在，等式由《代数》引理 \ref{algebra-lemma-pullback-module}
以及 $B \otimes_A \kappa_A = \kappa_B$ 这一事实得出。
\end{proof}

\begin{definition}
\label{varieties-definition-delta-invariant}
设 $k$ 是域。设 $X$ 是局部代数 $k$-概形。
设 $x \in X$ 是一点，使得 $\mathcal{O}_{X, x}$
约化且 $\dim(\mathcal{O}_{X, x}) = 1$。
{\it $\delta$-不变量（$X$ 在 $x$ 处）}定义为
$\delta$-不变量（关于 $\mathcal{O}_{X, x}$），如定义
\ref{varieties-definition-delta-invariant-algebra} 所述。
\end{definition}

\noindent
这个定义有意义，因为局部代数概形的局部环由
《代数》命题 \ref{algebra-proposition-ubiquity-nagata} 是 Nagata 环。
当然，更一般地，只要概形的一点 $x \in X$ 的局部环
$\mathcal{O}_{X, x}$ 是约化 Nagata 环且维数为 $1$，我们都可以作此定义。
由引理 \ref{varieties-lemma-delta-same-after-completion}，
$\delta$-不变量（$X$ 在 $x$ 处）为
$$
\delta\text{-invariant of }X\text{ 在 }x =
\delta\text{-invariant of }\mathcal{O}_{X, x}^h =
\delta\text{-invariant of }\mathcal{O}_{X, x}^\wedge
$$
由此可见，$\delta$-不变量是该点之完备局部环的不变量。

\begin{lemma}
\label{varieties-lemma-delta-invariant-and-change-of-fields}
设 $k$ 是域。设 $X$ 是局部代数 $k$-概形。
设 $K/k$ 是域扩张，并令 $Y = X_K$。
设 $y \in Y$ 的像为 $x \in X$。
假设 $X$ 在 $x$ 处几何约化，并且
$\dim(\mathcal{O}_{X, x}) = \dim(\mathcal{O}_{Y, y}) = 1$。
则
$$
\delta\text{-invariant of }X\text{ 在 }x \leq
\delta\text{-invariant of }Y\text{ 在 }y
$$
\end{lemma}

\begin{proof}
置 $A = \mathcal{O}_{X, x}$ 与 $B = \mathcal{O}_{Y, y}$。
由引理 \ref{varieties-lemma-geometrically-reduced-at-point} 可知
$A$ 几何约化。
因此，$B$ 是 $A \otimes_k K$ 的一个局部化。
令 $A \to A'$ 如引理 \ref{varieties-lemma-pre-delta-invariant} 所述。则
$$
B' = B \otimes_{(A \otimes_k K)} (A' \otimes_k K)
$$
在 $B$ 上有限，并且 $B \to B'$ 在总分式环上诱导同构。
事实上，选取满足引理 \ref{varieties-lemma-pre-pre-delta-invariant}
中 (1)--(6) 的 $f \in \mathfrak m_A$；由于
$\dim(B) = 1$，可知 $f \in \mathfrak m_B$
对 $B$ 起同样作用；又有 $B_f = B'_f$，因为 $A_f = A'_f$。
令 $B''$ 是 $B$ 在其总分式环中的整闭包，如引理
\ref{varieties-lemma-pre-delta-invariant} 所述。
则 $B' \subset B''$。因此，$\delta$-不变量（$Y$ 在 $y$ 处）为
$\text{length}_B(B''/B)$，并且
\begin{align*}
\text{length}_B(B''/B)
& \geq
\text{length}_B(B'/B) \\
& =
\text{length}_B((A'/A) \otimes_A B) \\
& =
\text{length}_B(B/\mathfrak m_A B) \text{length}_A(A'/A)
\end{align*}
这是因为 $A \to B$ 平坦（它是 $A \to A \otimes_k K$ 的一个局部化），
并可应用《代数》引理 \ref{algebra-lemma-pullback-module}。
由于 $\text{length}_A(A'/A)$ 是 $\delta$-不变量（$X$ 在 $x$ 处），
又由于 $\text{length}_B(B/\mathfrak m_A B) \geq 1$，引理得证。
\end{proof}

\begin{lemma}
\label{varieties-lemma-delta-invariant-and-change-of-fields-better}
设 $k$ 是域。设 $X$ 是局部代数 $k$-概形。
设 $K/k$ 是域扩张，并令 $Y = X_K$。
设 $y \in Y$ 的像为 $x \in X$。
假设引理 \ref{varieties-lemma-geometrically-normal-in-codim-1}
中的假设 (a)、(b)、(c) 对 $x \in X$ 成立，并且
$\dim(\mathcal{O}_{Y, y}) = 1$。
则 $\delta$-不变量（$X$ 在 $x$ 处）
等于 $\delta$-不变量（$Y$ 在 $y$ 处）。
\end{lemma}

\begin{proof}
置 $A = \mathcal{O}_{X, x}$ 与 $B = \mathcal{O}_{Y, y}$。
由引理 \ref{varieties-lemma-geometrically-normal-in-codim-1}
可知 $A$ 几何约化。
因此，$B$ 是 $A \otimes_k K$ 的一个局部化。
令 $A \to A'$ 如引理 \ref{varieties-lemma-pre-delta-invariant} 所述。
由引理 \ref{varieties-lemma-geometrically-normal-in-codim-1}
可知 $A' \otimes_k K$ 正规。
因此
$$
B' = B \otimes_{(A \otimes_k K)} (A' \otimes_k K)
$$
正规、在 $B$ 上有限，并且 $B \to B'$ 在总分式环上诱导同构。
事实上，选取满足引理 \ref{varieties-lemma-pre-pre-delta-invariant}
中 (1)--(6) 的 $f \in \mathfrak m_A$；由于
$\dim(B) = 1$，可知 $f \in \mathfrak m_B$
对 $B$ 起同样作用；又有 $B_f = B'_f$，因为 $A_f = A'_f$。
由此，$B \to B'$ 对 $B$ 而言如引理
\ref{varieties-lemma-pre-delta-invariant} 所述。因此只需证明
$\text{length}_A(A'/A) =
\text{length}_B(B'/B) = \text{length}_B((A'/A) \otimes_A B)$。
由于 $A \to B$ 平坦（它是 $A \to A \otimes_k K$ 的一个局部化），
并且 $\mathfrak m_B = \mathfrak m_A B$（因为
$B/\mathfrak m_A B$ 由上述说明是零维的，又是
$K \otimes_k \kappa(x)$ 的一个局部化；后者约化，因为
$\kappa(x)$ 在 $k$ 上可分），结论由《代数》引理
\ref{algebra-lemma-pullback-module} 得出。
\end{proof}





\section{分枝数}
\label{varieties-section-number-of-branches}

\noindent
我们已在《性质》第 \ref{properties-section-unibranch} 节中
定义了概形在一点处的分枝数。

\begin{lemma}
\label{varieties-lemma-number-of-branches}
设 $X$ 是概形。假设 $X$ 的每个拟紧开集都只有有限多个不可约分支。
设 $\nu : X^\nu \to X$ 是 $X$ 的正规化。设 $x \in X$。
\begin{enumerate}
\item $X$ 在 $x$ 处的分枝数等于 $x$ 在 $X^\nu$ 中的逆像数。
\item $X$ 在 $x$ 处的几何分枝数为
$\sum_{\nu(x^\nu) = x} [\kappa(x^\nu) : \kappa(x)]_s$。
\end{enumerate}
\end{lemma}

\begin{proof}
首先注意，对 $X$ 的假设恰好意味着正规化有定义；见《态射》定义
\ref{morphisms-definition-normalization}。
于是茎 $A' = (\nu_*\mathcal{O}_{X^\nu})_x$ 是
$A = \mathcal{O}_{X, x}$ 在 $A_{red}$ 的总分式环中的整闭包；见
《态射》引理 \ref{morphisms-lemma-stalk-normalization}。
由于 $\nu$ 是整态射，可知 $X^\nu$ 中位于 $x$ 上方的点
对应于 $A'$ 中位于极大理想 $\mathfrak m$ 上方的素理想，
而后者是 $A$ 的极大理想。
由于 $A \to A'$ 是整的，这也就是 $A'$ 的极大理想
（《代数》引理 \ref{algebra-lemma-integral-no-inclusion} 与
\ref{algebra-lemma-integral-going-up}）。
因此，本引理由其代数版本得出：
《代数续篇》引理 \ref{more-algebra-lemma-number-of-branches-1}。
\end{proof}

\begin{lemma}
\label{varieties-lemma-geometric-branches-and-change-of-fields}
设 $k$ 是域。设 $X$ 是局部代数 $k$-概形。
设 $K/k$ 是域扩张。设 $y \in X_K$ 是一点，其像 $x$ 位于
$X$ 中。则 $X$ 在 $x$ 处的几何分枝数等于
$X_K$ 在 $y$ 处的几何分枝数。
\end{lemma}

\begin{proof}
记 $Y = X_K$，并令 $X^\nu$ 以及相应的 $Y^\nu$ 为
$X$ 以及相应的 $Y$ 的正规化。考虑交换图
$$
\xymatrix{
Y^\nu \ar[r] \ar[d] & X^\nu_K \ar[r] \ar[d]_{\nu_K} & X^\nu \ar[d]_\nu \\
Y \ar@{=}[r] & Y \ar[r] & X
}
$$
由引理 \ref{varieties-lemma-normalization-and-change-of-fields}，左上方的水平箭头
是泛同胚。因此它诱导纯不可分剩余域扩张；见《态射》引理
\ref{morphisms-lemma-universal-homeomorphism} 与
\ref{morphisms-lemma-universally-injective}。
故由引理 \ref{varieties-lemma-number-of-branches}，$Y$ 在 $y$ 处的几何分枝数为
$\sum_{\nu_K(y') = y} [\kappa(y') : \kappa(y)]_s$。
类似地，$\sum_{\nu(x') = x} [\kappa(x') : \kappa(x)]_s$
是 $X$ 在 $x$ 处的几何分枝数。
利用《概形》引理 \ref{schemes-lemma-points-fibre-product}，
我们的陈述由下列代数事实得出：给定域扩张 $l/\kappa$
与代数域扩张 $m/\kappa$，则
$$
\sum\nolimits_{m \otimes_\kappa l \to m'} [m' : l']_s = [m : \kappa]_s
$$
其中求和遍历 $m \otimes_\kappa l$ 的商域。
可以用初等方法证明这一点，也可以把引理
\ref{varieties-lemma-separably-closed-field-connected-components} 应用于
\begingroup
\small
$$
\Spec(m \otimes_\kappa l) \times_{\Spec(l)} \Spec(\overline{l}) =
\Spec(m) \otimes_{\Spec(\kappa)} \Spec(\overline{l})
\longrightarrow
\Spec(m) \times_{\Spec(\kappa)} \Spec(\overline{\kappa})
$$
\endgroup
因为可以把 $[m : \kappa]_s$ 解释为右端的连通分支数，而把和
$\sum_{m \otimes_\kappa l \to m'} [m' : l']_s$
解释为左端的连通分支数。
\end{proof}

\begin{lemma}
\label{varieties-lemma-geometrically-unibranch-and-change-of-fields}
设 $k$ 是域。设 $X$ 是局部代数 $k$-概形。
设 $K/k$ 是域扩张。设 $y \in X_K$ 是一点，其像 $x$ 位于
$X$ 中。则 $X$ 在 $x$ 处几何单枝，当且仅当
$X_K$ 在 $y$ 处几何单枝。
\end{lemma}

\begin{proof}
这立即由引理 \ref{varieties-lemma-geometric-branches-and-change-of-fields}
和《代数续篇》引理
\ref{more-algebra-lemma-number-of-branches-1} 得出。
\end{proof}

\begin{definition}
\label{varieties-definition-wedge}
设 $A$ 以及 $A_i$（$1 \leq i \leq n$）都是局部环。称
{\it $A$ 是 $A_1, \ldots, A_n$ 的楔合}，若存在同构
$$
\kappa_{A_1} \to \kappa_{A_2} \to \ldots \to \kappa_{A_n}
$$
且 $A$ 同构于由那些 $n$-元组
$(a_1, \ldots, a_n) \in A_1 \times \ldots \times A_n$ 构成的环，
其中这些元组映到 $\kappa_{A_n}$ 的同一个元素。
\end{definition}

\noindent
若给定基环 $\Lambda$，并且 $A$ 与 $A_i$ 都是 $\Lambda$-代数，
则还要求 $\kappa_{A_i} \to \kappa_{A_{i + 1}}$
是 $\Lambda$-代数同构，并要求 $A$ 作为 $\Lambda$-代数同构于如下
$\Lambda$-代数：它由那些 $n$-元组
$(a_1, \ldots, a_n) \in A_1 \times \ldots \times A_n$ 构成，
且这些元组映到 $\kappa_{A_n}$ 的同一个元素。
特别地，若 $\Lambda = k$ 是域，且映射
$k \to \kappa_{A_i}$ 都是同构，则诸同构
$\kappa_{A_i} \to \kappa_{A_{i + 1}}$ 有唯一选择，
并且我们常说{\it $A_1, \ldots, A_n$ 的楔合}。

\begin{lemma}
\label{varieties-lemma-delta-number-branches-inequality-sh}
设 $(A, \mathfrak m)$ 是严格 Hensel 的
$1$ 维约化 Nagata 局部环。则
$$
\delta\text{-invariant of }A \geq \text{number of geometric branches of }A - 1
$$
若等号成立，则 $A$ 是 $n \geq 1$ 个严格 Hensel 离散赋值环的楔合。
\end{lemma}

\begin{proof}
几何分枝数等于 $A$ 的分枝数（这由《代数续篇》定义
\ref{more-algebra-definition-number-of-branches} 立即得出）。
令 $A \to A'$ 如引理 \ref{varieties-lemma-pre-delta-invariant} 所述。
注意，$A$ 的分枝数等于 $A'$ 的极大理想数；见
《代数续篇》引理 \ref{more-algebra-lemma-number-of-branches-1}。
有满射
$$
A'/A \longrightarrow
\left(\prod\nolimits_{\mathfrak m'} \kappa(\mathfrak m')\right)/
\kappa(\mathfrak m)
$$
由于 $\dim_{\kappa(\mathfrak m)} \prod \kappa(\mathfrak m')$
是 $\geq$ 分枝数，不等式显然成立。

\medskip\noindent
若等号成立，则 $\kappa(\mathfrak m') = \kappa(\mathfrak m)$
对所有 $\mathfrak m' \subset A'$ 都成立，并且上面的显示箭头是同构。
由于 $A$ 是 Hensel 环且 $A \to A'$ 有限，可知 $A'$ 是若干
Hensel 局部环之积；见《代数》引理
\ref{algebra-lemma-finite-over-henselian}。
这些因子就是局部环 $A'_{\mathfrak m'}$；由于
$A'$ 正规，这些因子是离散赋值环
（《代数》引理 \ref{algebra-lemma-characterize-dvr}）。
由于显示箭头是同构，可知 $A$ 确实是这些局部环的楔合。
\end{proof}

\begin{lemma}
\label{varieties-lemma-delta-number-branches-inequality}
设 $(A, \mathfrak m)$ 是 $1$ 维约化 Nagata 局部环。则
$$
\delta\text{-invariant of }A \geq \text{number of geometric branches of }A - 1
$$
\end{lemma}

\begin{proof}
可以把 $A$ 替换为 $A$ 的严格 Hensel 化，而不改变
$\delta$-不变量（引理 \ref{varieties-lemma-delta-same-after-completion}），
也不改变 $A$ 的几何分枝数（这由定义立即得出；见
《代数续篇》定义 \ref{more-algebra-definition-number-of-branches}）。
因此可以假设 $A$ 严格 Hensel，并应用引理
\ref{varieties-lemma-delta-number-branches-inequality-sh}。
\end{proof}






\section{一维概形的正规化}
\label{varieties-section-normalization-one-dimensional}

\noindent
由 Krull--Akizuki 结果，维数为 $1$ 的 Noether 概形的正规化态射
具有意想不到的良好性质。

\begin{lemma}
\label{varieties-lemma-normalize-noetherian-dim-1}
设 $X$ 是维数为 $1$ 的局部 Noether 概形。
设 $\nu : X^\nu \to X$ 是正规化。则
\begin{enumerate}
\item $\nu$ 是整的、满的，并在不可约分支上诱导双射；
\item 存在分解 $X^\nu \to X_{red} \to X$，且态射
$X^\nu \to X_{red}$ 是 $X_{red}$ 的正规化；
\item $X^\nu \to X_{red}$ 是双有理的；
\item 对每个闭点 $x \in X$，茎
$(\nu_*\mathcal{O}_{X^\nu})_x$ 是 $\mathcal{O}_{X, x}$
在 $(\mathcal{O}_{X, x})_{red} = \mathcal{O}_{X_{red}, x}$
的总分式环中的整闭包；
\item $\nu$ 的纤维有限，且剩余域扩张有限；
\item $X^\nu$ 是若干整正规局部 Noether 概形的不交并，
且每个仿射开集都是有限个 Dedekind 整环之积的谱。
\end{enumerate}
\end{lemma}

\begin{proof}
许多结论实际上都是正规化态射的一般性质；见《态射》引理
\ref{morphisms-lemma-normalization-reduced}、
\ref{morphisms-lemma-stalk-normalization}、
\ref{morphisms-lemma-normalization-normal} 和
\ref{morphisms-lemma-normalization-birational}。
并不明显的是：纤维有限，所诱导的剩余域扩张有限，以及
$X^\nu$ 局部看起来是 Dedekind 整环的谱（因而是 Noether 的）。
为证明这一点，可以假设 $X = \Spec(A)$ 仿射、Noether、维数为
$1$，并且 $A$ 约化。
然后可以使用《态射》引理
\ref{morphisms-lemma-description-normalization} 中的描述，
把问题约化到 $A$ 是维数为 $1$ 的 Noether 整环的情形。
此时所需性质由 Krull--Akizuki 定理得出，其形式见《代数》引理
\ref{algebra-lemma-integral-closure-Dedekind}。
\end{proof}

\noindent
当然，当 $X$ 不约化时，下列引理也有一个变体。

\begin{lemma}
\label{varieties-lemma-prepare-delta-invariant}
设 $X$ 是维数为 $1$ 的约化 Nagata 概形。设 $\nu : X^\nu \to X$
是正规化。设 $x \in X$ 表示一个闭点。则
\begin{enumerate}
\item $\nu : X^\nu \to X$ 有限、满且双有理；
\item $\mathcal{O}_X \subset \nu_*\mathcal{O}_{X^\nu}$，并且
$\nu_*\mathcal{O}_{X^\nu}/\mathcal{O}_X$ 是摩天层的直和，
其值为 $\mathcal{Q}_x$，位于 $x$ 处；后者非零当且仅当
$x$ 是 $X$ 的奇异点；
\item $A' = (\nu_*\mathcal{O}_{X^\nu})_x$ 是
$A = \mathcal{O}_{X, x}$ 在其总分式环中的整闭包；
\item $\mathcal{Q}_x = A'/A$ 具有有限长度，且该长度等于
$\delta$-不变量（$X$ 在 $x$ 处）；
\item $A'$ 是半局部环，它是有限个 Dedekind 整环之积；
\item $A^\wedge$ 是维数为 $1$ 的约化 Noether 完备局部环；
\item $(A')^\wedge$ 是 $A^\wedge$ 在其总分式环中的整闭包；
\item $(A')^\wedge$ 是有限个完备离散赋值环之积；
\item $A'/A \cong (A')^\wedge/A^\wedge$。
\end{enumerate}
\end{lemma}

\begin{proof}
我们可以并且将要使用引理 \ref{varieties-lemma-normalize-noetherian-dim-1}
的全部结论。
$\nu$ 的有限性由《态射》引理
\ref{morphisms-lemma-nagata-normalization} 得出。
由于 $X$ 约化、Nagata 且维数为 $1$，由《代数续篇》命题
\ref{more-algebra-proposition-ubiquity-J-2} 可知正则轨迹是
稠密开集 $U \subset X$。
由于正则概形正规，这说明 $\nu$ 在 $U$ 上是同构。
由于 $\dim(X) \leq 1$，这意味着 $\nu$ 在一离散闭点集
$x \in X$ 上不是同构。
特别地，有短正合列
$$
0 \to \mathcal{O}_X \to \nu_*\mathcal{O}_{X^\nu} \to
\bigoplus\nolimits_{x \in X \setminus U} \mathcal{Q}_x \to 0
$$
由于引理 \ref{varieties-lemma-normalize-noetherian-dim-1} 已描述
$\nu_*\mathcal{O}_{X^\nu}$ 的茎，我们断定 $Q_x = A'/A$
的长度确实等于 $\delta$-不变量（$X$ 在 $x$ 处）。
例如由引理 \ref{varieties-lemma-delta-invariant-is-zero} 可知，
$Q_x \not = 0$ 恰好当 $x$ 是奇异点。
$A'$ 是若干半局部 Dedekind 整环之积这一描述也由引理
\ref{varieties-lemma-normalize-noetherian-dim-1} 得出。
$A$、$A'$ 与 $(A')^\wedge$ 之间的关系，我们已在引理
\ref{varieties-lemma-normalization-same-after-completion}（及其证明）中见过。
\end{proof}




\section{寻找仿射开集}
\label{varieties-section-finding-affine-opens}

\noindent
我们继续《性质》第 \ref{properties-section-finding-affine-opens} 节中
开始的讨论。事实证明，在分离概形上，可以找到包含给定有限个余维
$1$ 的点的仿射开集。见命题
\ref{varieties-proposition-finite-set-of-points-of-codim-1-in-affine}。

\medskip\noindent
我们将在《下降》引理
\ref{descent-lemma-characterize-open-immersion} 中改进下列引理。

\begin{lemma}
\label{varieties-lemma-characterize-open-immersion}
设 $f : X \to Y$ 是概形态射。以 $X^0$ 表示
$X$ 的不可约分支的泛点集合。若
\begin{enumerate}
\item $f$ 分离；
\item 存在开覆盖 $X = \bigcup U_i$，使得
$f|_{U_i} : U_i \to Y$ 是开浸入；
\item 若 $\xi, \xi' \in X^0$ 且 $\xi \not = \xi'$，则
$f(\xi) \not = f(\xi')$，
\end{enumerate}
则 $f$ 是开浸入。
\end{lemma}

\begin{proof}
假设 $y = f(x) = f(x')$。选取特化 $y_0 \leadsto y$，
其中 $y_0$ 是 $Y$ 的某个不可约分支的泛点。
由于 $f$ 在源上局部为同构，可以选取特化
$x_0 \leadsto x$ 与 $x'_0 \leadsto x'$，它们映到
$y_0 \leadsto y$。
注意 $x_0, x'_0 \in X^0$。故由假设 (3)，$x_0 = x'_0$。
由于 $f$ 分离，我们断定 $x = x'$。因此 $f$ 是开浸入。
\end{proof}

\begin{lemma}
\label{varieties-lemma-local-isomorphism}
设 $X \to S$ 是概形态射。设 $x \in X$ 是一点，其像为
$s \in S$。若
\begin{enumerate}
\item $\mathcal{O}_{X, x} = \mathcal{O}_{S, s}$；
\item $S$ 约化；
\item $X \to S$ 是有限型的；
\item $S$ 有有限多个不可约分支，
\end{enumerate}
则存在开邻域 $U$（关于 $x$），使得 $f|_U$ 是开浸入。
\end{lemma}

\begin{proof}
可以移除 $S$ 中不包含 $s$ 的（有限多个）不可约分支。
可以用仿射开邻域替换 $S$，该邻域包含 $s$。
可以用仿射开邻域替换 $X$，该邻域包含 $x$。
记 $S = \Spec(A)$、$X = \Spec(B)$。令
$\mathfrak q \subset B$ 以及相应的 $\mathfrak p \subset A$
分别为对应于 $x$ 以及相应的 $s$ 的素理想。
由于 $A$ 约化，且 $A$ 的所有极小素理想都包含在
$\mathfrak p$ 中，可知 $A \subset A_\mathfrak p$。
由于 $X \to S$ 是有限型的，$B$ 在 $A$ 上有限型。
设 $b_1, \ldots, b_n \in B$ 是生成 $B$（在 $A$ 上）的元素。
由于 $A_\mathfrak p = B_\mathfrak q$，可以找到
$f \in A$、$f \not \in \mathfrak p$ 以及 $a_i \in A$，
使得 $b_i$ 与 $a_i/f$ 在 $B_\mathfrak q$ 中有相同的像。
因此可以找到 $g \in B$、$g \not \in \mathfrak q$，使得
$g(fb_i - a_i) = 0$ 于 $B$ 中。由此，
$A_f \to B_{fg}$ 的像包含 $b_1, \ldots, b_n$ 的像，
特别也包含 $g$ 的像。
选取 $n \geq 0$ 与 $f' \in A$，使得 $f'/f^n$
映到 $g$ 在 $B_{fg}$ 中的像。由于
$A_\mathfrak p = B_\mathfrak q$，可知 $f' \not \in \mathfrak p$。
我们断定 $A_{ff'} \to B_{fg}$ 满。
最后，由于 $A_{ff'} \subset A_\mathfrak p = B_\mathfrak q$
（见上文），映射 $A_{ff'} \to B_{fg}$ 单，因而是同构。
\end{proof}

\begin{lemma}
\label{varieties-lemma-points-in-affine}
设 $f : T \to X$ 是概形态射。以 $X^0$ 以及相应的 $T^0$
表示不可约分支的泛点集合。
设 $t_1, \ldots, t_m \in T$ 是有限点集，其像为
$x_j = f(t_j)$。若
\begin{enumerate}
\item $T$ 仿射；
\item $X$ 拟分离；
\item $X^0$ 有限；
\item $f(T^0) \subset X^0$，且 $f : T^0 \to X^0$ 单；
\item $\mathcal{O}_{X, x_j} = \mathcal{O}_{T, t_j}$，
\end{enumerate}
则存在 $X$ 的仿射开集，包含 $x_1, \ldots, x_r$。
\end{lemma}

\begin{proof}
使用《极限》命题 \ref{limits-proposition-affine}，
可立即约化到 $X$ 与 $T$ 都约化的情形。略去细节。

\medskip\noindent
假设 $X$ 与 $T$ 约化。可以把 $T = \lim_{i \in I} T_i$
写成 $X$ 上有限表示概形的有向极限，转移态射为仿射态射；见
《极限》引理 \ref{limits-lemma-relative-approximation}。
选取 $i \in I$ 使得 $T_i$ 仿射；见《极限》引理
\ref{limits-lemma-limit-affine}。
记 $T_i = \Spec(R_i)$、$T = \Spec(R)$。
令 $R' \subset R$ 是 $R_i \to R$ 的像。
则 $T' = \Spec(R')$ 仿射、约化、在 $X$ 上有限型，且
$T \to T'$ 支配。对 $j = 1, \ldots, r$，令
$t'_j \in T'$ 为 $t_j$ 的像。考虑局部环映射
$$
\mathcal{O}_{X, x_j} \to
\mathcal{O}_{T', t'_j} \to
\mathcal{O}_{T, t_j}
$$
以 $(T')^0$ 表示 $T'$ 的不可约分支的泛点集合。
设 $\xi \leadsto t'_j$ 是一个特化，其中 $\xi \in (T')^0$。
由于 $T \to T'$ 支配，可以选取 $\eta \in T^0$ 映到
$\xi$（警告：事先并不知道 $\eta$ 特化到 $t_j$）。
把假设 (3) 应用于 $\eta$，可知像 $\theta$（即
$\xi$ 在 $X$ 中的像）对应于 $\mathcal{O}_{X, x_j}$
的一个极小素理想。
通过 (5) 中的同构提升 $\xi$，得到特化
$\eta' \leadsto t_j$，其中 $\eta' \in T^0$ 映到
$\theta \leadsto x_j$。
(4) 的单射性表明 $\eta = \eta'$。因此
$\mathcal{O}_{T', t'_j}$ 的每个极小素理想都位于
$\mathcal{O}_{T, t_j}$ 的一个极小素理想之下。我们断定
$\mathcal{O}_{T', t'_j} \to \mathcal{O}_{T, t_j}$ 单，
因而上面两个映射都是同构。

\medskip\noindent
由引理 \ref{varieties-lemma-local-isomorphism}，存在开集
$U \subset T'$，它包含所有点 $t'_j$，使得
$U \to X$ 是引理 \ref{varieties-lemma-characterize-open-immersion}
意义下的局部同构。
由该引理可知 $U \to X$ 是开浸入。
最后，由《性质》引理
\ref{properties-lemma-ample-finite-set-in-affine}，
可以找到开集 $W \subset U \subset T'$，它包含所有 $t'_j$。
$W$ 在 $X$ 中的像就是所需的仿射开集。
\end{proof}

\begin{lemma}
\label{varieties-lemma-finite-set-codim-1-points-in-affine}
设 $X$ 是整分离概形。设 $x_1, \ldots, x_r \in X$
是有限点集，使得 $\mathcal{O}_{X, x_i}$ 是维数
$\leq 1$ 的 Noether 环。则存在 $X$ 的仿射开子概形，
它包含所有 $x_1, \ldots, x_r$。
\end{lemma}

\begin{proof}
设 $K$ 是 $X$ 的有理函数域。
置 $A_i = \mathcal{O}_{X, x_i}$。则 $A_i \subset K$，且 $K$
是 $A_i$ 的分式域。由于 $X$ 分离，并且
$x_i \not = x_j$，不可能存在赋值环 $\mathcal{O} \subset K$
同时支配 $A_i$ 与 $A_j$。事实上，考虑图
$$
\xymatrix{
\Spec(\mathcal{O}) \ar[r] \ar[d] & \Spec(A_1) \ar[d] \\
\Spec(A_2) \ar[r] & X
}
$$
并应用分离性的赋值判据
（《概形》引理 \ref{schemes-lemma-separated-implies-valuative}），
就会得到 $x_i = x_j$。因此，由引理
\ref{varieties-lemma-glue-separated} 可知，$A_i \otimes A_j \to K$
对所有 $i \not = j$ 都满。
由引理 \ref{varieties-lemma-glue-a-bunch-of-local-rings} 可知
$A = A_1 \cap \ldots \cap A_r$ 是 Noether 半局部环，
恰有 $r$ 个极大理想
$\mathfrak m_1, \ldots, \mathfrak m_r$，使得
$A_i = A_{\mathfrak m_i}$。此外，
$$
\Spec(A) = \Spec(A_1) \cup \ldots \cup \Spec(A_r)
$$
是开覆盖，并且该覆盖任意两块的交都是 $\Spec(K)$。
因此，给定态射 $\Spec(A_i) \to X$ 可粘合为概形态射
$$
\Spec(A) \longrightarrow X
$$
它把 $\mathfrak m_i$ 映到 $x_i$，并诱导局部环同构。
于是结论由引理 \ref{varieties-lemma-points-in-affine} 得出。
\end{proof}

\begin{lemma}
\label{varieties-lemma-extra-silly}
设 $A$ 是环，$I \subset A$ 是理想，
$\mathfrak p_1, \ldots, \mathfrak p_r$ 是 $A$ 的素理想，并且
$\overline{f} \in A/I$ 是元素。若
$I \not \subset \mathfrak p_i$ 对所有 $i$ 都成立，则存在
$f \in A$、$f \not \in \mathfrak p_i$，它映到
$\overline{f}$（在 $A/I$ 中）。
\end{lemma}

\begin{proof}
可以假设诸 $\mathfrak p_i$ 之间没有包含关系
（移除较小的素理想即可）。首先选取任意 $f \in A$ 提升
$\overline{f}$。令 $S$ 为那些 $s \in \{1, \ldots, r\}$ 的集合，
使得 $f \in \mathfrak p_s$。若 $S$ 为空，则已经完成。否则，
考虑理想 $J = I \prod_{i \not \in S} \mathfrak p_i$。
注意，$J$ 不包含在 $\mathfrak p_s$ 中，其中 $s \in S$，
因为诸 $\mathfrak p_i$ 之间没有包含关系，并且
$I$ 不包含在任何 $\mathfrak p_i$ 中。
故由《代数》引理 \ref{algebra-lemma-silly}，可以选取
$g \in J$、$g \not \in \mathfrak p_s$，其中 $s \in S$。
于是 $f + g$ 是引理所述问题的一个解。
\end{proof}

\begin{lemma}
\label{varieties-lemma-finite-set-codim-1-points-in-affine-per-component}
设 $X$ 是概形。设 $T \subset X$ 是有限点集。假设
\begin{enumerate}
\item $X$ 有有限多个不可约分支 $Z_1, \ldots, Z_t$；
\item $Z_i \cap T$ 包含在与 $Z_i$ 对应的赋予约化诱导结构的
子概形的一个仿射开集中。
\end{enumerate}
则存在 $X$ 的仿射开子概形，它包含 $T$。
\end{lemma}

\begin{proof}
使用《极限》命题 \ref{limits-proposition-affine}，
可立即约化到 $X$ 约化的情形。略去细节。在证明的其余部分，
对 $X$ 的每个闭子集都赋予诱导约化闭子概形结构。

\medskip\noindent
我们归纳证明，可以找到仿射开集
$U \subset Z_1 \cup \ldots \cup Z_r$，它包含
$T \cap (Z_1 \cup \ldots \cup Z_r)$。当 $r = 1$ 时，
这由假设成立。
设 $r > 1$，并令 $U \subset Z_1 \cup \ldots \cup Z_{r -  1}$
是包含 $T \cap (Z_1 \cup \ldots \cup Z_{r - 1})$ 的仿射开集。
设 $V \subset X_r$ 是包含 $T \cap Z_r$ 的仿射开集
（由假设存在）。则 $U \cap V$ 包含
$T \cap ( Z_1 \cup \ldots \cup Z_{r - 1} ) \cap Z_r$。
因此
$$
\Delta = (U \cap Z_r) \setminus (U \cap V)
$$
不包含 $T$ 的任何元素。注意 $\Delta$ 是 $U$ 的闭子集。
由素理想回避（《代数》引理 \ref{algebra-lemma-silly}），
可以找到标准开集 $U'$（属于 $U$），它包含 $T \cap U$ 并避开
$\Delta$，即 $U' \cap Z_r \subset U \cap V$。
把 $U$ 替换为 $U'$ 后，可以假设 $U \cap V$ 在 $U$ 中闭。

\medskip\noindent
与上面相同的论证还表明，集合
$\Delta' = (U \cap (Z_1 \cup \ldots \cup Z_{r - 1})) \setminus (U \cap V)$
不包含 $T$ 的任何元素。利用这一点，可以找到
$h \in \mathcal{O}(V)$，使得 $D(h) \subset V$ 包含
$T \cap V$，并且 $U \cap D(h) \subset U \cap V$。
利用 $U \cap V$ 在 $U$ 中闭这一事实，可以应用引理
\ref{varieties-lemma-extra-silly} 找到 $g \in \mathcal{O}(U)$，
它在 $U \cap V$ 上的限制等于 $h$ 在 $U \cap V$ 上的限制，
并且 $T \cap U \subset D(g)$。然后可以把 $U$ 替换为 $D(g)$，
把 $V$ 替换为 $D(h)$，从而达到 $U \cap V$ 在 $U$ 与 $V$
中都闭的情形。此时概形 $U \cup V$ 由《极限》引理
\ref{limits-lemma-affines-glued-in-closed-affine} 是仿射的。
这证明了归纳步骤，从而证明了引理。
\end{proof}

\noindent
下面是我们可以从上述材料中得出的结论。

\begin{proposition}
\label{varieties-proposition-finite-set-of-points-of-codim-1-in-affine}
设 $X$ 是分离概形，使得每个拟紧开集只有有限多个不可约分支。
设 $x_1, \ldots, x_r \in X$ 是点，使得
$\mathcal{O}_{X, x_i}$ 是维数 $\leq 1$ 的 Noether 环。
则存在 $X$ 的仿射开子概形，它包含所有 $x_1, \ldots, x_r$。
\end{proposition}

\begin{proof}
可以替换 $X$，所用拟紧开集包含 $x_1, \ldots, x_r$，
因而可以假设 $X$ 有有限多个不可约分支。
由引理 \ref{varieties-lemma-finite-set-codim-1-points-in-affine-per-component}，
可约化到 $X$ 为整概形的情形。这就是引理
\ref{varieties-lemma-finite-set-codim-1-points-in-affine}。
\end{proof}




\section{曲线}
\label{varieties-section-curves}

\noindent
在 Stacks 计划中，我们采用如下曲线定义。

\begin{definition}
\label{varieties-definition-curve}
设 $k$ 为域。所谓一条{\it 曲线}，是指维数为 $1$ 的、域 $k$ 上的簇。
\end{definition}

\noindent
域 $k$ 上曲线的两个标准例子是仿射直线 $\mathbf{A}^1_k$
和射影直线 $\mathbf{P}^1_k$。
\par\noindent
概形 $X = \Spec(k[x, y]/(f))$
是曲线，当且仅当 $f \in k[x, y]$ 不可约。

\medskip\noindent
曲线的定义具有与簇的定义相同的问题；见定义 \ref{varieties-definition-variety}
之后的讨论。此外，这意味着每条曲线都带有一个指定的定义域。
例如，$X = \Spec(\mathbf{C}[x])$ 是域 $\mathbf{C}$ 上的曲线，
但我们也可以把它看成域 $\mathbf{R}$ 上的曲线。概形
$\Spec(\mathbf{Z})$ 不是曲线，尽管概形 $\Spec(\mathbf{Z})$
与 $\mathbf{A}^1_{\mathbf{F}_p}$ 在许多方面表现相似。

\begin{lemma}
\label{varieties-lemma-proper-minus-point}
设 $X$ 是维数 $> 0$ 的、域 $k$ 上的分离不可约概形。
设 $x \in X$ 为闭点。开子概形 $X \setminus \{x\}$
在 $k$ 上不固有。
\end{lemma}

\begin{proof}
由于 $X$ 不可约，$U = X \setminus \{x\}$ 在 $X$ 中不是闭集。
特别地，浸入 $U \to X$ 不固有。
由《态射》引理 \ref{morphisms-lemma-image-proper-scheme-closed}
（这里使用了 $X$ 分离），$U \to \Spec(k)$ 也不固有。
\end{proof}

\begin{lemma}
\label{varieties-lemma-dim-1-quasi-projective}
设 $X$ 是域 $k$ 上分离的有限型概形。
若 $\dim(X) \leq 1$，则 $X$ 在 $k$ 上 H-拟射影。
\end{lemma}

\begin{proof}
由命题 \ref{varieties-proposition-dim-1-noetherian-separated-has-ample}，
概形 $X$ 具有丰沛可逆层 $\mathcal{L}$。
由《态射》引理 \ref{morphisms-lemma-quasi-projective-finite-type-over-S}，
可知 $X$ 同构于 $\mathbf{P}^n_k$ 在 $\Spec(k)$ 上的一个局部闭子概形。
这正是在 $k$ 上 H-拟射影的定义；见《态射》定义
\ref{morphisms-definition-quasi-projective}。
\end{proof}

\begin{lemma}
\label{varieties-lemma-dim-1-proper-projective}
设 $X$ 是域 $k$ 上的固有概形。
若 $\dim(X) \leq 1$，则 $X$ 在 $k$ 上 H-射影。
\end{lemma}

\begin{proof}
由引理 \ref{varieties-lemma-dim-1-quasi-projective}，可知
$X$ 是 $\mathbf{P}^n_k$ 的局部闭子概形，其中 $k$ 为某个域。
由于 $X$ 在 $k$ 上固有，可知 $X$ 是 $\mathbf{P}^n_k$
的闭子概形（《态射》引理
\ref{morphisms-lemma-image-proper-scheme-closed}）。
\end{proof}

\begin{lemma}
\label{varieties-lemma-dim-1-projective-completion}
设 $X$ 是 $k$ 上分离的有限型概形。
若 $\dim(X) \leq 1$，则存在开浸入
$j : X \to \overline{X}$，满足下列性质：
\begin{enumerate}
\item $\overline{X}$ 在 $k$ 上 H-射影，即
$\overline{X}$ 是 $\mathbf{P}^d_k$ 的闭子概形，对某个 $d$ 成立；
\item $j(X) \subset \overline{X}$ 稠密且概形论稠密；
\item $\overline{X} \setminus X = \{x_1, \ldots, x_n\}$，
其中 $x_i \in \overline{X}$ 是闭点。
\end{enumerate}
\end{lemma}

\begin{proof}
由引理 \ref{varieties-lemma-dim-1-quasi-projective}，可假设
$X$ 是 $\mathbf{P}^d_k$ 的局部闭子概形，其中 $d$ 为某个整数。令
$\overline{X} \subset \mathbf{P}^d_k$ 为 $X \to \mathbf{P}^d_k$
的概形论像；见《态射》定义
\ref{morphisms-definition-scheme-theoretic-image}。
《态射》引理 \ref{morphisms-lemma-quasi-compact-immersion}
中的描述给出性质 (1) 与 (2)。
例如考察泛点并应用引理 \ref{varieties-lemma-dimension-locally-algebraic}，可知
$\dim(X) = 1 \Rightarrow \dim(\overline{X}) = 1$。
由于 $\overline{X}$ 是 Noether 概形，因而
$\overline{X} \setminus X = \{x_1, \ldots, x_n\}$
是有限个闭点组成的集合。
\end{proof}

\begin{lemma}
\label{varieties-lemma-reduced-dim-1-projective-completion}
设 $X$ 是 $k$ 上分离的有限型概形。
若 $X$ 约化且 $\dim(X) \leq 1$，则存在开浸入
$j : X \to \overline{X}$，使得
\begin{enumerate}
\item $\overline{X}$ 在 $k$ 上 H-射影，即
$\overline{X}$ 是 $\mathbf{P}^d_k$ 的闭子概形，对某个 $d$ 成立；
\item $j(X) \subset \overline{X}$ 稠密且概形论稠密；
\item $\overline{X} \setminus X = \{x_1, \ldots, x_n\}$，
其中 $x_i \in \overline{X}$ 是闭点；
\item 局部环 $\mathcal{O}_{\overline{X}, x_i}$ 都是离散赋值环，
对 $i = 1, \ldots, n$ 成立。
\end{enumerate}
\end{lemma}

\begin{proof}
取引理 \ref{varieties-lemma-dim-1-projective-completion} 中的
$j : X \to \overline{X}$。考虑正规化 $X'$，即 $\overline{X}$ 在
$X$ 中的正规化。由引理
\ref{varieties-lemma-relative-normalization-finite}，态射
$X' \to \overline{X}$ 有限。
由《态射》引理 \ref{morphisms-lemma-finite-projective}，
$X' \to \overline{X}$ 射影。由《态射》引理
\ref{morphisms-lemma-projective-over-quasi-projective-is-H-projective}，
可知 $X' \to \overline{X}$ H-射影。
由《态射》引理 \ref{morphisms-lemma-H-projective-composition}，
可知 $X' \to \Spec(k)$ H-射影。
令 $\{x'_1, \ldots, x'_m\} \subset X'$ 为
$\{x_1, \ldots, x_n\} = \overline{X} \setminus X$ 的逆像。
则都有 $\dim(\mathcal{O}_{X', x'_i}) = 1$，
对所有 $1 \leq i \leq m$ 成立。
因此，由《态射》引理
\ref{morphisms-lemma-relative-normalization-normal-codim-1}，
局部环 $\mathcal{O}_{X', x'}$ 都是离散赋值环。
于是 $X \to X'$ 与 $\{x'_1, \ldots, x'_m\}$ 即为所求。
\end{proof}

\begin{lemma}
\label{varieties-lemma-dim-1-projective-completion-after-insep}
设 $X$ 是 $k$ 上分离的有限型概形，且 $\dim(X) \leq 1$。
则存在交换图
$$
\xymatrix{
\overline{Y}_1 \amalg \ldots \amalg \overline{Y}_n \ar[rd] &
Y_1 \amalg \ldots \amalg Y_n \ar[r]_-\nu \ar[d] \ar[l]^j &
X_{k'} \ar[r] \ar[d] &
X \ar[d]^f \\
& \Spec(k'_1) \amalg \ldots \amalg \Spec(k'_n) \ar[r] &
\Spec(k') \ar[r] &
\Spec(k)
}
$$
满足下列性质：
\begin{enumerate}
\item $k'/k$ 是域的有限纯不可分扩张；
\item $\nu$ 是 $X_{k'}$ 的正规化；
\item $j$ 是像稠密的开浸入；
\item $k'_i/k'$ 是有限可分扩张，对 $i = 1, \ldots, n$ 成立；
\item $\overline{Y}_i$ 光滑、射影、几何不可约，维数 $\leq 1$，
其基域为 $k'_i$。
\end{enumerate}
\end{lemma}

\begin{proof}
由于可以把 $X$ 替换为它的约化，我们可以并且确实假设 $X$ 约化。
选取引理 \ref{varieties-lemma-reduced-dim-1-projective-completion}
中的 $X \to \overline{X}$。
若能对 $\overline{X}$ 证明本引理，则对 $X$ 的结论也随之成立
（略去细节）。因此可以并且确实假设 $X$ 射影。

\medskip\noindent
选取有限纯不可分扩张 $k'/k$，使 $X_{k'}$ 的正规化
在 $k'$ 上几何正规；见引理
\ref{varieties-lemma-finite-extension-geometrically-normal}。
以 $Y = (X_{k'})^\nu$ 记该正规化；关于正规化的性质，
见第 \ref{varieties-section-normalization} 节。
$Y$ 几何正则，因为在维数 $\leq 1$ 时正规与正则等价；见《性质》引理
\ref{properties-lemma-normal-dimension-1-regular}。
因而由引理 \ref{varieties-lemma-geometrically-regular-smooth}，
$Y$ 在 $k'$ 上光滑。
设 $Y = Y_1 \amalg \ldots \amalg Y_n$ 为 $Y$ 的不可约分支分解。
置 $k'_i = \Gamma(Y_i, \mathcal{O}_{Y_i})$。
由引理 \ref{varieties-lemma-proper-geometrically-reduced-global-sections}，
这些都是 $k'$ 的有限可分扩张。
最后应用引理 \ref{varieties-lemma-baby-stein} 即完成证明。
\end{proof}

\begin{lemma}
\label{varieties-lemma-regular-point-on-curve}
设 $k$ 为域，$X$ 为域 $k$ 上的曲线，$x \in X$ 为闭点。
我们把 $x$ 看作 $X$ 的一个（约化）闭子概形，其理想层为
$\mathcal{I}$。下列条件等价：
\begin{enumerate}
\item $\mathcal{O}_{X, x}$ 正则；
\item $\mathcal{O}_{X, x}$ 正规；
\item $\mathcal{O}_{X, x}$ 是离散赋值环；
\item $\mathcal{I}$ 是可逆 $\mathcal{O}_X$-模；
\item $x$ 是 $X$ 上的有效 Cartier 除子。
\end{enumerate}
若 $k$ 完全，或者 $\kappa(x)$ 在 $k$ 上可分，
则这些条件还等价于
\begin{enumerate}
\item[(6)] $X \to \Spec(k)$ 在 $x$ 处光滑。
\end{enumerate}
\end{lemma}

\begin{proof}
由于 $X$ 是曲线，局部环 $\mathcal{O}_{X, x}$ 是维数为 $1$ 的
Noether 局部整环（引理 \ref{varieties-lemma-dimension-locally-algebraic}）。
条件 (4) 与 (5) 依定义等价，并且都等价于
$\mathcal{I}_x = \mathfrak m_x \subset \mathcal{O}_{X, x}$
可以由一个元素生成（《除子》引理
\ref{divisors-lemma-effective-Cartier-in-points}）。
因此，(1)、(2)、(3)、(4)、(5) 的等价性由《代数》引理
\ref{algebra-lemma-characterize-dvr} 得出。
当 $k$ 完全时，最后的陈述由引理
\ref{varieties-lemma-dense-smooth-open-variety-over-perfect-field} 得出。
若 $\kappa(x)/k$ 可分，则该等价性由《代数》引理
\ref{algebra-lemma-separable-smooth} 得出。
\end{proof}

\begin{remark}
\label{varieties-remark-smooth-projective-compactification}
设 $k$ 为域，$X$ 为域 $k$ 上的正则曲线。
由引理 \ref{varieties-lemma-regular-point-on-curve} 与
\ref{varieties-lemma-reduced-dim-1-projective-completion}，
存在非奇异射影曲线 $\overline{X}$，它是 $X$ 的紧化；换言之，存在开浸入
$j : X \to \overline{X}$，其补集由有限个闭点组成。
若 $k$ 完全，则 $X$ 与 $\overline{X}$ 都在 $k$ 上光滑，
而 $\overline{X}$ 是 $X$ 的光滑射影紧化。
\end{remark}

\noindent
注意，若仿射概形 $X$ 在 $k$ 上且在 $k$ 上固有，
则 $X$ 在 $k$ 上有限（《态射》引理
\ref{morphisms-lemma-finite-proper}），因而维数为 $0$
（《代数》引理 \ref{algebra-lemma-finite-dimensional-algebra} 与
命题 \ref{algebra-proposition-dimension-zero-ring}）。
因此，维数 $> 0$ 的、域 $k$ 上的概形不可能既仿射又在 $k$ 上固有。
所以，下面引理中的两种可能相互排斥。

\begin{lemma}
\label{varieties-lemma-curve-affine-projective}
设 $X$ 是域 $k$ 上的曲线。则 $X$ 或者是仿射概形，
或者 $X$ 在 $k$ 上 H-射影。
\end{lemma}

\begin{proof}
选取引理 \ref{varieties-lemma-reduced-dim-1-projective-completion} 中的
$X \to \overline{X}$，使得
$\overline{X} \setminus X = \{x_1, \ldots, x_r\}$。
于是 $\overline{X}$ 也是曲线。
若 $r = 0$，则 $X = \overline{X}$ 在 $k$ 上 H-射影。
因此可以假设 $r \geq 1$，我们的目标是证明 $X$ 仿射。
由引理 \ref{varieties-lemma-open-in-affine-curve-affine}，
只须证明 $\overline{X} \setminus \{x_1\}$ 仿射。
这把问题约化为下一段中的断言。

\medskip\noindent
设 $X$ 为域 $k$ 上的 H-射影曲线。设 $x \in X$ 为闭点，
使得 $\mathcal{O}_{X, x}$ 是离散赋值环。我们断言
$U = X \setminus \{x\}$ 仿射。
由引理 \ref{varieties-lemma-regular-point-on-curve}，
点 $x$ 定义了 $X$ 上的有效 Cartier 除子。
对 $n \geq 1$，以 $nx = x + \ldots + x$ 表示 $n$ 重和；见《除子》定义
\ref{divisors-definition-sum-effective-Cartier-divisors}。
以 $\mathcal{O}_{nx}$ 表示 $nx$ 的结构层，把它看作 $X$ 上的凝聚模。
由于局部概形 $nx$ 上的每个可逆模都是平凡的，
《除子》注 \ref{divisors-remark-ses-regular-section}
中的第一个短正合列在这里成为
$$
0 \to \mathcal{O}_X
\xrightarrow{1} \mathcal{O}_X(nx) \to \mathcal{O}_{nx} \to 0
$$
注意，$\dim_k H^0(X, \mathcal{O}_{nx}) \geq n$。
事实上，由引理 \ref{varieties-lemma-chi-tensor-finite}，有
$H^0(X, \mathcal{O}_{nx}) = \mathcal{O}_{X, x}/(\pi^n)$，其中
$\pi$ 是 $\mathcal{O}_{X, x}$ 中的素元，并且幂 $\pi^i$ 映到
关于 $k$ 线性无关的 $\mathcal{O}_{X, x}/(\pi^n)$ 中元素，
对 $i = 0, 1, \ldots, n - 1$ 成立。
由《概形上同调》引理
\ref{coherent-lemma-proper-over-affine-cohomology-finite}，
有 $\dim_k H^1(X, \mathcal{O}_X) < \infty$。
若 $n > \dim_k H^1(X, \mathcal{O}_X)$，
则由同调长正合列可知，存在
$s \in \Gamma(X, \mathcal{O}_X(nx))$，
它不是 $\mathcal{O}_X$ 的截面。
若取具有此性质的最小 $n$，则 $s$ 映到茎
$\left(\mathcal{O}_X(nx)\right)_x$ 的一个生成元；
否则它会定义 $\mathcal{O}_X((n - 1)x) \subset \mathcal{O}_X(nx)$
的一个截面。对这个 $n$，可知 $s_0 = 1$ 与 $s_1 = s$
生成可逆模 $\mathcal{L} = \mathcal{O}_X(nx)$。

\medskip\noindent
考虑《构造》第 \ref{constructions-section-projective-space} 节中的
相应态射
$f = \varphi_{\mathcal{L}, (s_0, s_1)} : X \to \mathbf{P}^1_k$。
注意，$D_{+}(T_0)$ 的逆像是 $U = X \setminus \{x\}$，
因为截面 $s_0$（作为 $\mathcal{L}$ 的截面）仅在 $x$ 处消失。
特别地，$f$ 不是常值态射，即 $\Im(f)$ 含有不止一个点。
因此，$f$ 必把泛点 $\eta$（$X$ 的泛点）映到
$\mathbf{P}^1_k$ 的泛点。所以，若
$y \in \mathbf{P}^1_k$ 是闭点，则 $f^{-1}(\{y\})$
是 $X$ 中不含 $\eta$ 的闭集，因而是有限集。
最后，$f$ 固有\footnote{事实上，由《态射》引理
\ref{morphisms-lemma-projective-is-quasi-projective-proper}，
H-射影簇是固有簇。由《态射》引理
\ref{morphisms-lemma-image-proper-scheme-closed}，
源为固有簇的簇之间的态射是固有态射。}。
由《概形上同调》引理
\ref{coherent-lemma-proper-finite-fibre-finite-in-neighbourhood}\footnote{可以避免使用这个依赖形式函数定理的引理。
事实上，$X$ 射影，因此只须证明：固有态射 $f : X \to Y$
若连接两个在 $k$ 上拟射影的概形且纤维有限，则是有限态射。为此，选取 $X$ 中包含
$f$ 在某点 $y$ 上之纤维的仿射开集；这里使用了拟射影概形的任意有限点集
在域 $k$ 上都包含于某个仿射开集这一事实。把 $Y$ 缩小到 $y$ 的一个小仿射
邻域后，问题约化为仿射概形之间的固有态射。由《态射》引理
\ref{morphisms-lemma-integral-universally-closed}，这种态射有限。}，
可知 $f$ 有限。因此 $U = f^{-1}(D_{+}(T_0))$ 仿射。
\end{proof}

\noindent
下面的引理与引理 \ref{varieties-lemma-proper-minus-point} 合用说明：
给定分离概形 $X$，其维数为 $1$ 且在 $k$ 上有限型，
那么 $X \setminus Z$ 是仿射概形，只要闭子集 $Z$
与 $X$ 的每个不可约分支相交。

\begin{lemma}
\label{varieties-lemma-dim-1-nonproper-affine}
设 $X$ 是 $k$ 上分离的有限型概形。
若 $\dim(X) \leq 1$，并且 $X$ 的不可约分支中没有维数为 $1$
的固有分支，则 $X$ 仿射。
\end{lemma}

\begin{proof}
令 $X = \bigcup X_i$ 为 $X$ 的不可约分支分解。
我们把 $X_i$ 看作整概形（赋予约化诱导概形结构；见《概形》定义
\ref{schemes-definition-reduced-induced-scheme}）。
特别地，$X_i$ 或者是单点集（因而仿射），或者是曲线，
从而由引理 \ref{varieties-lemma-curve-affine-projective} 仿射。
于是 $\coprod X_i \to X$ 是有限满射，而 $\coprod X_i$ 仿射。
因此，由《概形上同调》引理
\ref{coherent-lemma-image-affine-finite-morphism-affine-Noetherian}，
可知 $X$ 仿射。
\end{proof}




\section{曲线上的次数}
\label{varieties-section-divisors-curves}

\noindent
我们开始定义域上维数为 $1$ 的固有概形之可逆层，乃至更一般的
局部自由层的次数。在第 \ref{varieties-section-euler} 节中，
我们定义了凝聚层 $\mathcal{F}$（位于固有概形 $X$ 上，基域为
$k$）的欧拉示性数，公式为
$$
\chi(X, \mathcal{F}) = \sum (-1)^i \dim_k H^i(X, \mathcal{F}).
$$

\begin{definition}
\label{varieties-definition-degree-invertible-sheaf}
设 $k$ 为域，$X$ 为维数 $\leq 1$ 的固有概形，其基域为 $k$，
$\mathcal{L}$ 为可逆 $\mathcal{O}_X$-模。$\mathcal{L}$ 的{\it 次数}定义为
$$
\deg(\mathcal{L}) = \chi(X, \mathcal{L}) - \chi(X, \mathcal{O}_X)
$$
更一般地，若 $\mathcal{E}$ 是秩为 $n$ 的局部自由层，
则定义 $\mathcal{E}$ 的{\it 次数}为
$$
\deg(\mathcal{E}) = \chi(X, \mathcal{E}) - n\chi(X, \mathcal{O}_X)
$$
\end{definition}

\noindent
注意，这依赖于三元组 $\mathcal{E}/X/k$。
若 $X$ 不连通，而 $\mathcal{E}$ 有限局部自由（但秩不恒定），
则可以修改定义，对 $\mathcal{E}$ 在 $X$ 各连通分支上的限制之次数求和。
若 $\mathcal{E}$ 只是凝聚层，则有若干不同的方式扩展此定义\footnote{若
$X$ 是固有曲线，$\mathcal{F}$ 是 $X$ 上的凝聚层，
则次数常定义为 $\chi(X, \mathcal{F}) - r\chi(X, \mathcal{O}_X)$，其中
$r = \dim_{\kappa(\xi)} \mathcal{F}_\xi$ 是 $\mathcal{F}$
在泛点 $\xi$（$X$ 的泛点）处的秩。}。
下面用一系列引理证明，该定义具有次数应有的一切性质。

\begin{lemma}
\label{varieties-lemma-degree-base-change}
设 $k'/k$ 为域扩张。设 $X$ 是维数 $\leq 1$ 的固有概形，其基域为 $k$。
设 $\mathcal{E}$ 是局部自由 $\mathcal{O}_X$-模，其秩恒为 $n$。
则 $\mathcal{E}/X/k$ 的次数等于
$\mathcal{E}_{k'}/X_{k'}/k'$ 的次数。
\end{lemma}

\begin{proof}
更明确地，置 $X_{k'} = X \times_{\Spec(k)} \Spec(k')$。
令 $\mathcal{E}_{k'} = p^*\mathcal{E}$，其中
$p : X_{k'} \to X$ 是投影。由《概形上同调》引理
\ref{coherent-lemma-flat-base-change-cohomology}，有
$H^i(X_{k'}, \mathcal{E}_{k'}) = H^i(X, \mathcal{E}) \otimes_k k'$
以及
$H^i(X_{k'}, \mathcal{O}_{X_{k'}}) = H^i(X, \mathcal{O}_X) \otimes_k k'$。
因此欧拉示性数不变，次数也不变。
\end{proof}

\begin{lemma}
\label{varieties-lemma-degree-additive}
设 $k$ 为域。设 $X$ 是维数 $\leq 1$ 的固有概形，其基域为 $k$。
设 $0 \to \mathcal{E}_1 \to \mathcal{E}_2 \to \mathcal{E}_3 \to 0$
是有限恒秩局部自由 $\mathcal{O}_X$-模的短正合列。则
$$
\deg(\mathcal{E}_2) = \deg(\mathcal{E}_1) + \deg(\mathcal{E}_3)
$$
\end{lemma}

\begin{proof}
这立即由欧拉示性数的可加性（引理
\ref{varieties-lemma-euler-characteristic-additive}）与秩的可加性得出。
\end{proof}

\begin{lemma}
\label{varieties-lemma-degree-birational-pullback}
设 $k$ 为域。设 $f : X' \to X$ 是维数 $\leq 1$
的固有概形之间的双有理态射，这些概形定义在 $k$ 上。则
$$
\deg(f^*\mathcal{E}) = \deg(\mathcal{E})
$$
对每个恒秩有限局部自由层都成立。更一般地，只须 $f$
在维数为 $1$ 的不可约分支之间诱导双射，并在相应泛点处的局部环之间诱导同构。
\end{lemma}

\begin{proof}
态射 $f$ 固有（《态射》引理
\ref{morphisms-lemma-image-proper-scheme-closed}），
且其纤维维数 $\leq 0$。因此 $f$ 有限（《概形上同调》引理
\ref{coherent-lemma-proper-finite-fibre-finite-in-neighbourhood}）。
从而
$$
Rf_*f^*\mathcal{E} = f_*f^*\mathcal{E} =
\mathcal{E} \otimes_{\mathcal{O}_X} f_*\mathcal{O}_{X'}
$$
由于 $f$ 在所有维数为 $1$ 的不可约分支的泛点处的局部环上诱导同构，
可知核与余核
$$
0 \to \mathcal{K} \to \mathcal{O}_X \to f_*\mathcal{O}_{X'}
\to \mathcal{Q} \to 0
$$
的支撑维数 $\leq 0$。注意，由于 $\mathcal{E}$ 局部自由，
把该列与 $\mathcal{E}$ 张量后仍是正合列。于是得到
\begin{align*}
\chi(X, \mathcal{E}) - \chi(X', f^*\mathcal{E})
& =
\chi(X, \mathcal{E}) - \chi(X, f_*f^*\mathcal{E}) \\
& =
\chi(X, \mathcal{E}) - \chi(X, \mathcal{E} \otimes f_*\mathcal{O}_{X'}) \\
& =
\chi(X, \mathcal{K} \otimes \mathcal{E}) -
\chi(X, \mathcal{Q} \otimes \mathcal{E}) \\
& =
n\chi(X, \mathcal{K}) -
n\chi(X, \mathcal{Q}) \\
& =
n\chi(X, \mathcal{O}_X) - n\chi(X, f_*\mathcal{O}_{X'}) \\
& =
n\chi(X, \mathcal{O}_X) - n\chi(X', \mathcal{O}_{X'})
\end{align*}
这正是所需结论。第一个等号来自 $f$ 有限；见《概形上同调》引理
\ref{coherent-lemma-relative-affine-cohomology}。
第二个等号来自投影公式；见《同调》引理
\ref{cohomology-lemma-projection-formula}。
第三个等号来自欧拉示性数的可加性；见引理
\ref{varieties-lemma-euler-characteristic-additive}。
第四个等号来自引理 \ref{varieties-lemma-chi-tensor-finite}。
\end{proof}

\begin{lemma}
\label{varieties-lemma-degree-on-proper-curve}
设 $k$ 为域。设 $X$ 是 $k$ 上泛点为 $\xi$ 的固有曲线。
设 $\mathcal{E}$ 是局部自由 $\mathcal{O}_X$-模，其秩为 $n$，
$\mathcal{F}$ 是凝聚 $\mathcal{O}_X$-模。则
$$
\chi(X, \mathcal{E} \otimes \mathcal{F}) =
r \deg(\mathcal{E}) + n \chi(X, \mathcal{F})
$$
其中 $r = \dim_{\kappa(\xi)} \mathcal{F}_\xi$ 是 $\mathcal{F}$ 的秩。
\end{lemma}

\begin{proof}
令 $\mathcal{P}$ 为凝聚层 $\mathcal{F}$（位于 $X$ 上）的如下性质：
引理中的公式成立。我们断言，《概形上同调》引理
\ref{coherent-lemma-property} 的假设 (1) 与 (2) 对 $\mathcal{P}$ 成立。
事实上，(1) 成立，因为欧拉示性数与秩 $r$
在凝聚层的短正合列中都可加。(2) 也成立：若 $Z = X$，
则可取 $\mathcal{G} = \mathcal{O}_X$，而
$\mathcal{P}(\mathcal{O}_X)$ 依次数定义为真。
若 $i : Z \to X$ 是一个闭点的包含态射，则可取
$\mathcal{G} = i_*\mathcal{O}_Z$；而 $\mathcal{P}$ 成立，
这由引理 \ref{varieties-lemma-chi-tensor-finite} 以及此时 $r = 0$ 这一事实得出。
\end{proof}

\noindent
设 $k$ 为域。设 $X$ 是 $k$ 上维数 $\leq 1$ 的有限型概形。
设 $C_i \subset X$（$i = 1, \ldots, t$）为维数 $1$ 的不可约分支。
我们赋予 $C_i$ 约化诱导概形结构，从而把它看作概形。
令 $\xi_i \in C_i$ 为泛点。$C_i$ 在 $X$ 中的{\it 重数}定义为长度
$$
m_i = \text{length}_{\mathcal{O}_{X, \xi_i}} \mathcal{O}_{X, \xi_i}
$$
这是有意义的，因为 $\mathcal{O}_{X, \xi_i}$ 是零维 Noether 局部环，
因而作为自身上的模具有有限长度（《代数》命题
\ref{algebra-proposition-dimension-zero-ring}）。
更多信息见《Chow 同调》第
\ref{chow-section-cycle-of-closed-subscheme} 节。
事实上，局部自由层的次数只依赖于它在各不可约分支上的限制。

\begin{lemma}
\label{varieties-lemma-degree-in-terms-of-components}
设 $k$ 为域。设 $X$ 是维数 $\leq 1$ 的固有概形，其基域为 $k$。
设 $\mathcal{E}$ 是局部自由 $\mathcal{O}_X$-模，其秩为 $n$。则
$$
\deg(\mathcal{E}) = \sum m_i \deg(\mathcal{E}|_{C_i})
$$
其中 $C_i \subset X$（$i = 1, \ldots, t$）是赋予约化诱导概形结构的
维数 $1$ 的不可约分支，而 $m_i$ 是 $C_i$ 在 $X$ 中的重数。
\end{lemma}

\begin{proof}
先注意，该陈述有意义，因为 $C_i \to \Spec(k)$ 固有且维数为 $1$；
下面还将使用如下事实：由《态射》引理
\ref{morphisms-lemma-closed-immersion-proper} 与
\ref{morphisms-lemma-composition-proper}，$X$ 的任意闭子概形都在 $k$ 上固有。
考虑开子概形 $U_i = X \setminus (\bigcup_{j \not = i} C_j)$，
并令 $X_i \subset X$ 为 $U_i$ 的概形论闭包。
注意，$X_i \cap U_i = U_i$（概形论意义下），而
$X_i \cap U_j = \emptyset$（集合论意义下），对 $i \not = j$ 成立；这由《态射》引理
\ref{morphisms-lemma-quasi-compact-immersion} 中对概形论闭包的描述得出。
因此可把引理 \ref{varieties-lemma-degree-birational-pullback} 应用于态射
$X' = \coprod X_i \to X$。显然，$C_i \subset X_i$（概形论意义下），
且 $C_i$ 在 $X_i$ 中的重数等于 $C_i$ 在 $X$ 中的重数，
故问题约化为下一段所讨论的情形。

\medskip\noindent
假设 $X$ 不可约，泛点为 $\xi$。设
$C = X_{red}$ 的重数为 $m$。
我们必须证明 $\deg(\mathcal{E}) = m \deg(\mathcal{E}|_C)$。
令 $\mathcal{I} \subset \mathcal{O}_X$ 为定义闭子概形 $C$ 的理想。
取最小的 $e \geq 0$，使 $\mathcal{I}^{e + 1} = 0$
（《概形上同调》引理
\ref{coherent-lemma-power-ideal-kills-sheaf}）。
我们对 $e$ 作归纳。若 $e = 0$，则 $X = C$，结论显然。
否则，置 $\mathcal{F} = \mathcal{I}^e$，将其看作凝聚
$\mathcal{O}_C$-模（《概形上同调》引理
\ref{coherent-lemma-i-star-equivalence}）。
令 $X' \subset X$ 为凝聚理想 $\mathcal{I}^e$ 截出的闭子概形，
并令 $m'$ 为 $C$ 在 $X'$ 中的重数。取短正合列
$$
0 \to \mathcal{F} \to \mathcal{O}_X \to \mathcal{O}_{X'} \to 0
$$
在 $\xi$ 处的茎，可得（使用《代数》引理
\ref{algebra-lemma-length-additive}、
\ref{algebra-lemma-dimension-is-length} 与
\ref{algebra-lemma-length-independent}）
$$
m = \text{length}_{\mathcal{O}_{X, \xi}} \mathcal{O}_{X, \xi}
= \dim_{\kappa(\xi)} \mathcal{F}_\xi +
\text{length}_{\mathcal{O}_{X', \xi}} \mathcal{O}_{X', \xi}
= r + m'
$$
其中 $r$ 是 $\mathcal{F}$ 作为 $C$ 上凝聚层的秩。
与 $\mathcal{E}$ 张量后，得到短正合列
$$
0 \to \mathcal{E}|_C \otimes \mathcal{F} \to \mathcal{E} \to
\mathcal{E} \otimes \mathcal{O}_{X'} \to 0
$$
由归纳假设，有
$\deg(\mathcal{E}|_{X'}) = m' \deg(\mathcal{E}|_C)$。
由引理 \ref{varieties-lemma-degree-on-proper-curve}，有
$\chi(\mathcal{E}|_C \otimes \mathcal{F}) =
r \deg(\mathcal{E}|_C) + n \chi(\mathcal{F})$。
综合以上各式，即得结论。
\end{proof}




\begin{lemma}
\label{varieties-lemma-degree-tensor-product}
设 $k$ 为域，$X$ 为维数 $\leq 1$ 的固有概形，其基域为 $k$，
$\mathcal{E}$、$\mathcal{V}$ 为恒定有限秩的局部自由
$\mathcal{O}_X$-模。则
$$
\deg(\mathcal{E} \otimes \mathcal{V}) =
\text{rank}(\mathcal{E}) \deg(\mathcal{V}) +
\text{rank}(\mathcal{V}) \deg(\mathcal{E})
$$
\end{lemma}

\begin{proof}
由引理 \ref{varieties-lemma-degree-in-terms-of-components} 与初等算术，
问题约化为固有曲线的情形。此情形由引理
\ref{varieties-lemma-degree-on-proper-curve} 得出。
\end{proof}

\begin{lemma}
\label{varieties-lemma-degree-and-det}
设 $k$ 为域，$X$ 为维数 $\leq 1$ 的固有概形，其基域为 $k$，
$\mathcal{E}$ 为局部自由 $\mathcal{O}_X$-模，秩为 $n$。则
$$
\deg(\mathcal{E}) = \deg(\wedge^n(\mathcal{E})) = \deg(\det(\mathcal{E}))
$$
\end{lemma}

\begin{proof}
由引理 \ref{varieties-lemma-degree-in-terms-of-components} 与初等算术，
问题约化为固有曲线的情形。
此时存在修正 $f : X' \to X$，使得
$f^*\mathcal{E}$ 具有一个逐次商均为可逆模的滤过；见《除子》引理
\ref{divisors-lemma-filter-after-modification}。
由引理 \ref{varieties-lemma-degree-birational-pullback}，可在 $X'$ 上处理。
因此，可以假设存在由局部自由 $\mathcal{O}_X$-模组成的滤过
$$
0 = \mathcal{E}_0 \subset \mathcal{E}_1 \subset \mathcal{E}_2 \subset
\ldots \subset \mathcal{E}_n = \mathcal{E}
$$
使得 $\mathcal{L}_i = \mathcal{E}_i/\mathcal{E}_{i - 1}$ 可逆。
由《模》引理 \ref{modules-lemma-det-ses} 及归纳法，可得
$\det(\mathcal{E}) = \mathcal{L}_1 \otimes \ldots \otimes \mathcal{L}_n$。
于是等式由引理 \ref{varieties-lemma-degree-tensor-product}
与可加性（引理 \ref{varieties-lemma-degree-additive}）得出。
\end{proof}

\begin{lemma}
\label{varieties-lemma-degree-effective-Cartier-divisor}
设 $k$ 为域，$X$ 为维数 $\leq 1$ 的固有概形，其基域为 $k$。
设 $D$ 是 $X$ 上的有效 Cartier 除子。
则 $D$ 在 $\Spec(k)$ 上有限，次数为
$\deg(D) = \dim_k \Gamma(D, \mathcal{O}_D)$。对于局部自由层
$\mathcal{E}$，若其秩为 $n$，则有
$$
\deg(\mathcal{E}(D)) = n\deg(D) + \deg(\mathcal{E})
$$
其中 $\mathcal{E}(D) = \mathcal{E} \otimes_{\mathcal{O}_X} \mathcal{O}_X(D)$。
\end{lemma}

\begin{proof}
由于 $D$ 在 $X$ 中无处稠密（《除子》引理
\ref{divisors-lemma-complement-effective-Cartier-divisor}），
可知 $\dim(D) \leq 0$。因此由引理
\ref{varieties-lemma-algebraic-scheme-dim-0}，$D$ 在 $k$ 上有限。
由于 $k$ 为域，态射 $D \to \Spec(k)$ 有限局部自由，
因而具有次数（《态射》定义
\ref{morphisms-definition-finite-locally-free}）；该次数显然等于
$\dim_k \Gamma(D, \mathcal{O}_D)$，如引理所述。
由《除子》定义
\ref{divisors-definition-invertible-sheaf-effective-Cartier-divisor}，
存在短正合列
$$
0 \to \mathcal{O}_X \to \mathcal{O}_X(D) \to i_*i^*\mathcal{O}_X(D) \to 0
$$
其中 $i : D \to X$ 是闭浸入。与 $\mathcal{E}$ 张量，得到短正合列
$$
0 \to \mathcal{E} \to \mathcal{E}(D) \to i_*i^*\mathcal{E}(D) \to 0
$$
引理中的等式由欧拉示性数的可加性（引理
\ref{varieties-lemma-euler-characteristic-additive}）与引理
\ref{varieties-lemma-chi-tensor-finite} 得出。
\end{proof}

\begin{lemma}
\label{varieties-lemma-divisible}
设 $k$ 为域。设 $X$ 是 $k$ 上约化且连通的固有概形。
令 $\kappa = H^0(X, \mathcal{O}_X)$。则 $\kappa/k$
是有限域扩张，并且 $w = [\kappa : k]$ 整除下列各数：
\begin{enumerate}
\item $\deg(\mathcal{E})$，其中所取的是任意局部自由
$\mathcal{O}_X$-模 $\mathcal{E}$；
\item $[\kappa(x) : k]$，其中 $x \in X$ 为任意闭点；以及
\item $\deg(D)$，其中 $D \subset X$ 为任意零维闭子概形。
\end{enumerate}
\end{lemma}

\begin{proof}
关于 $\kappa$ 的断言见引理
\ref{varieties-lemma-proper-geometrically-reduced-global-sections}。
对每个拟凝聚 $\mathcal{O}_X$-模，$k$-向量空间
$H^i(X, \mathcal{F})$ 都是 $\kappa$-向量空间。
上述整除性很容易由这一陈述与各定义得出。
\end{proof}



\begin{lemma}
\label{varieties-lemma-degree-pullback-map-proper-curves}
设 $k$ 为域。设 $f : X \to Y$ 是 $k$ 上固有曲线之间的非常值态射。
设 $\mathcal{E}$ 是局部自由 $\mathcal{O}_Y$-模。则
$$
\deg(f^*\mathcal{E}) = \deg(X/Y) \deg(\mathcal{E})
$$
\end{lemma}

\begin{proof}
$X$ 在 $Y$ 上的次数定义见《态射》定义
\ref{morphisms-definition-degree}。
因此 $f_*\mathcal{O}_X$ 是凝聚 $\mathcal{O}_Y$-模，其秩为
$\deg(X/Y)$；换言之，
$\deg(X/Y) = \dim_{\kappa(\xi)} (f_*\mathcal{O}_X)_\xi$，其中
$\xi$ 是 $Y$ 的泛点。于是得到
\begin{align*}
\chi(X, f^*\mathcal{E})
& =
\chi(Y, f_*f^*\mathcal{E}) \\
& =
\chi(Y, \mathcal{E} \otimes f_*\mathcal{O}_X) \\
& =
\deg(X/Y) \deg(\mathcal{E}) + n \chi(Y, f_*\mathcal{O}_X) \\
& =
\deg(X/Y) \deg(\mathcal{E}) + n \chi(X, \mathcal{O}_X)
\end{align*}
这正是所需结论。第一个等号来自 $f$ 有限；见《概形上同调》引理
\ref{coherent-lemma-relative-affine-cohomology}。
第二个等号来自投影公式；见《同调》引理
\ref{cohomology-lemma-projection-formula}。
第三个等号来自引理 \ref{varieties-lemma-degree-on-proper-curve}。
\end{proof}

\noindent
下面是以上次数结果的一个平凡但重要的推论。

\begin{lemma}
\label{varieties-lemma-check-invertible-sheaf-trivial}
设 $k$ 为域。设 $X$ 是 $k$ 上的固有曲线。
设 $\mathcal{L}$ 是可逆 $\mathcal{O}_X$-模。
\begin{enumerate}
\item 若 $\mathcal{L}$ 有非零截面，则
$\deg(\mathcal{L}) \geq 0$。
\item 若 $\mathcal{L}$ 有一个在某点处消失的非零截面 $s$，则
$\deg(\mathcal{L}) > 0$。
\item 若 $\mathcal{L}$ 与 $\mathcal{L}^{-1}$ 都有非零截面，则
$\mathcal{L} \cong \mathcal{O}_X$。
\item 若 $\deg(\mathcal{L}) \leq 0$ 且 $\mathcal{L}$ 有非零截面，则
$\mathcal{L} \cong \mathcal{O}_X$。
\item 若 $\mathcal{N} \to \mathcal{L}$ 是可逆
$\mathcal{O}_X$-模之间的非零映射，则
$\deg(\mathcal{L}) \geq \deg(\mathcal{N})$；若等号成立，则该映射是同构。
\end{enumerate}
\end{lemma}

\begin{proof}
设 $s$ 是 $\mathcal{L}$ 的非零截面。由于 $X$ 是曲线，
可知 $s$ 是正则截面。因此存在有效 Cartier 除子 $D \subset X$
以及同构 $\mathcal{L} \to \mathcal{O}_X(D)$，它把 $s$
映到典范截面 $1$，即 $\mathcal{O}_X(D)$ 的典范截面；见《除子》引理
\ref{divisors-lemma-characterize-OD}。
于是由引理 \ref{varieties-lemma-degree-effective-Cartier-divisor}，
$\deg(\mathcal{L}) = \deg(D)$。
由于 $\deg(D) \geq 0$，且其 $= 0$ 当且仅当
$D = \emptyset$，这证明了 (1) 与 (2)。
在情形 (3) 中，可知 $\deg(\mathcal{L}) = 0$ 且
$D = \emptyset$。情形 (4) 同理。为证明 (5)，把 (1) 与 (4)
应用于可逆层
$$
\mathcal{L} \otimes_{\mathcal{O}_X} \mathcal{N}^{\otimes -1} =
\SheafHom_{\mathcal{O}_X}(\mathcal{N}, \mathcal{L})
$$
即可；由引理 \ref{varieties-lemma-degree-tensor-product}，该层的次数为
$\deg(\mathcal{L}) - \deg(\mathcal{N})$。
\end{proof}

\begin{lemma}
\label{varieties-lemma-no-sections-dual-nef}
设 $k$ 为域。设 $X$ 是 $k$ 上约化、连通且等维数的固有概形，
其维数为 $1$。设 $\mathcal{L}$ 是可逆 $\mathcal{O}_X$-模。
若 $\deg(\mathcal{L}|_C) \leq 0$ 对每个不可约分支 $C$
均成立（这些分支均取自 $X$），则或者
$H^0(X, \mathcal{L}) = 0$，或者
$\mathcal{L} \cong \mathcal{O}_X$。
\end{lemma}

\begin{proof}
设 $s \in H^0(X, \mathcal{L})$ 非零。由于 $X$ 约化，存在
不可约分支 $C$，它是 $X$ 的分支，并且 $s|_C \not = 0$。
然而，若 $s|_C$ 非零，则由引理
\ref{varieties-lemma-check-invertible-sheaf-trivial}，$s$
在 $C$ 上处处不消失。这又蕴含 $s$ 在 $X$ 的每个与 $C$
相交的不可约分支上都处处不消失。
由于 $X$ 连通，可知 $s$ 无处消失，引理得证。
\end{proof}

\begin{lemma}
\label{varieties-lemma-ample-curve}
设 $k$ 为域。设 $X$ 是 $k$ 上的固有曲线。
设 $\mathcal{L}$ 是可逆 $\mathcal{O}_X$-模。
则 $\mathcal{L}$ 丰沛，当且仅当 $\deg(\mathcal{L}) > 0$。
\end{lemma}

\begin{proof}
若 $\mathcal{L}$ 丰沛，则存在 $n > 0$ 以及截面
$s \in H^0(X, \mathcal{L}^{\otimes n})$，使 $X_s$ 仿射。
由于 $X$ 不仿射（否则由《态射》引理
\ref{morphisms-lemma-finite-proper}，$X$ 将是有限概形），
可知 $s$ 在某点处消失。因此由引理
\ref{varieties-lemma-check-invertible-sheaf-trivial}，有
$\deg(\mathcal{L}^{\otimes n}) > 0$。
由引理 \ref{varieties-lemma-degree-tensor-product}，可得
$\deg(\mathcal{L}) = 1/n\deg(\mathcal{L}^{\otimes n}) > 0$。

\medskip\noindent
假设 $\deg(\mathcal{L}) > 0$。则
$$
\dim_k H^0(X, \mathcal{L}^{\otimes n}) \geq \chi(X, \mathcal{L}^n)
= n\deg(\mathcal{L}) + \chi(X, \mathcal{O}_X)
$$
随 $n$ 线性增长。因此，对任意有限个闭点
$x_1, \ldots, x_t$（均在 $X$ 上），可以找到 $n$，使得
$\dim_k H^0(X, \mathcal{L}^{\otimes n}) > \sum \dim_k \kappa(x_i)$。
（回忆，由 Hilbert 零点定理，域扩张
$\kappa(x_i)/k$ 有限；例如见《态射》引理
\ref{morphisms-lemma-closed-point-fibre-locally-finite-type}。）
因此可以找到非零截面 $s \in H^0(X, \mathcal{L}^{\otimes n})$，
它在 $x_1, \ldots, x_t$ 处消失。特别地，若选取
$x_1, \ldots, x_t$，使 $X \setminus \{x_1, \ldots, x_t\}$ 仿射，
则 $X_s$ 也仿射（例如由《性质》引理
\ref{properties-lemma-affine-cap-s-open}；不过，若选取该有限集，使得
$\mathcal{L}|_{X \setminus \{x_1, \ldots, x_t\}}$ 平凡，
则结论是立即的）。所以可以找到 $n > 0$ 与非零截面
$s \in H^0(X, \mathcal{L}^{\otimes n})$，使 $X_s$ 仿射。

\medskip\noindent
我们将证明：对每个拟凝聚理想层 $\mathcal{I}$，
都存在 $m > 0$，使
$H^1(X, \mathcal{I} \otimes \mathcal{L}^{\otimes m})$ 为零。
由《概形上同调》引理
\ref{coherent-lemma-vanshing-gives-ample}，这将完成证明。
为此，考虑映射
$$
\mathcal{I} \xrightarrow{s}
\mathcal{I} \otimes \mathcal{L}^{\otimes n} \xrightarrow{s}
\mathcal{I} \otimes \mathcal{L}^{\otimes 2n} \xrightarrow{s} \ldots
$$
由于 $\mathcal{I}$ 无挠，这些映射是单射，并且在 $X_s$ 上是同构；
因而这些映射的余核的 $H^1$ 为零（例如由《概形上同调》引理
\ref{coherent-lemma-coherent-support-dimension-0}）。
可知向量空间的映射
$$
H^1(X, \mathcal{I}) \to
H^1(X, \mathcal{I} \otimes \mathcal{L}^{\otimes n}) \to
H^1(X, \mathcal{I} \otimes \mathcal{L}^{\otimes 2n}) \to \ldots
$$
均为满射。另一方面，$H^1(X, \mathcal{I})$ 的维数有限，
且由《概形上同调》引理
\ref{coherent-lemma-section-affine-open-kills-classes}，
每个元素最终都映到零。因此，对某个 $e > 0$，有
$H^1(X, \mathcal{I} \otimes \mathcal{L}^{\otimes en})$ 为零。
证明完成。
\end{proof}

\begin{lemma}
\label{varieties-lemma-ampleness-in-terms-of-degrees-components}
设 $k$ 为域。设 $X$ 是维数 $\leq 1$ 的固有概形，其基域为 $k$。
设 $\mathcal{L}$ 是可逆 $\mathcal{O}_X$-模。
设 $C_i \subset X$（$i = 1, \ldots, t$）为维数 $1$ 的不可约分支。
下列条件等价：
\begin{enumerate}
\item $\mathcal{L}$ 丰沛；以及
\item $\deg(\mathcal{L}|_{C_i}) > 0$ 对 $i = 1, \ldots, t$ 均成立。
\end{enumerate}
\end{lemma}

\begin{proof}
设 $x_1, \ldots, x_r \in X$ 为孤立闭点。
把 $x_i = \Spec(\kappa(x_i))$ 看作概形。考虑概形态射
$$
f : C_1 \amalg \ldots \amalg C_t \amalg x_1 \amalg \ldots \amalg x_r
\longrightarrow X
$$
这是 $k$ 上固有概形之间的有限满射（略去细节）。
因此，$\mathcal{L}$ 丰沛，当且仅当 $f^*\mathcal{L}$ 丰沛
（《概形上同调》引理
\ref{coherent-lemma-surjective-finite-morphism-ample}）。
于是结论由引理 \ref{varieties-lemma-ample-curve} 得出。
\end{proof}

\begin{lemma}
\label{varieties-lemma-general-degree-g-line-bundle}
设 $k$ 为代数闭域。设 $X$ 是 $k$ 上的固有曲线。则存在
\begin{enumerate}
\item 可逆 $\mathcal{O}_X$-模 $\mathcal{L}$，使
$\dim_k H^0(X, \mathcal{L}) = 1$ 且 $H^1(X, \mathcal{L}) = 0$；以及
\item 可逆 $\mathcal{O}_X$-模 $\mathcal{N}$，使
$\dim_k H^0(X, \mathcal{N}) = 0$ 且 $H^1(X, \mathcal{N}) = 0$。
\end{enumerate}
\end{lemma}

\begin{proof}
选取闭浸入 $i : X \to \mathbf{P}^n_k$
（引理 \ref{varieties-lemma-dim-1-proper-projective}）。
置 $\mathcal{L} = i^*\mathcal{O}_{\mathbf{P}^n}(d)$，其中
$d \gg 0$，
可知存在可逆层 $\mathcal{L}$，满足
$H^0(X, \mathcal{L}) \not = 0$ 且
$H^1(X, \mathcal{L}) = 0$（关于消失，见《概形上同调》引理
\ref{coherent-lemma-vanshing-gives-ample}；关于非消失，见该处所列参考）。
我们对 $t = \dim_k H^0(X, \mathcal{L})$ 作降归纳，以完成 (1) 的证明。
基础情形 $t = 1$ 显然。假设 $t > 1$。

\medskip\noindent
令 $U \subset X$ 为引理
\ref{varieties-lemma-dense-smooth-open-variety-over-perfect-field}
中所研究的非奇异点组成的非空开集。
设 $s \in H^0(X, \mathcal{L})$ 非零。存在闭点
$x \in U$，使得 $s$ 在 $x$ 处不消失。
设 $\mathcal{I}$ 为引理 \ref{varieties-lemma-regular-point-on-curve}
中 $i : x \to X$ 的理想层。考察短正合列
$$
0 \to \mathcal{I} \otimes_{\mathcal{O}_X} \mathcal{L} \to
\mathcal{L} \to i_*i^*\mathcal{L} \to 0
$$
注意，$H^0(X, i_*i^*\mathcal{L}) = H^0(x, i^*\mathcal{L})$
的维数为 $1$，因为 $x$ 是 $k$-有理点（$k$ 代数闭）。
由于 $s$ 在 $x$ 处不消失，可知
$$
H^0(X, \mathcal{L}) \longrightarrow H^0(X, i_*i^*\mathcal{L})
$$
满射。因此
$\dim_k H^0(X, \mathcal{I} \otimes_{\mathcal{O}_X} \mathcal{L}) = t - 1$。
最后，同调长正合列还表明
$H^1(X, \mathcal{I} \otimes_{\mathcal{O}_X} \mathcal{L}) = 0$，
从而完成归纳步骤的证明。

\medskip\noindent
为得到 (2) 中的可逆层，取 (1) 中的可逆层
$\mathcal{L}$，再作一次上一段的论证即可。
\end{proof}

\begin{lemma}
\label{varieties-lemma-vanishing-degree-2g-and-1-line-bundle}
设 $k$ 为代数闭域。设 $X$ 是 $k$ 上的固有曲线。
置 $g = \dim_k H^1(X, \mathcal{O}_X)$。
对每个可逆 $\mathcal{O}_X$-模 $\mathcal{L}$，若
$\deg(\mathcal{L}) \geq 2g - 1$，则
$H^1(X, \mathcal{L}) = 0$。
\end{lemma}

\begin{proof}
令 $\mathcal{N}$ 为引理
\ref{varieties-lemma-general-degree-g-line-bundle} 第 (2) 部分中找到的可逆模。
$\mathcal{N}$ 的次数为
$\chi(X, \mathcal{N}) - \chi(X, \mathcal{O}_X) = 0 - (1 - g) = g - 1$。
因此 $\mathcal{L} \otimes \mathcal{N}^{\otimes - 1}$ 的次数为
$\deg(\mathcal{L}) - (g - 1) \geq g$。从而
$\chi(X, \mathcal{L} \otimes \mathcal{N}^{\otimes -1}) \geq g + 1 - g = 1$。
所以存在非零整体截面 $s$，其零点概形是有效 Cartier 除子
$D$，且次数为 $\deg(\mathcal{L}) - (g - 1)$。
这给出短正合列
$$
0 \to \mathcal{N} \xrightarrow{s} \mathcal{L} \to i_*(\mathcal{L}|_D) \to 0
$$
其中 $i : D \to X$ 是包含态射。可知
$H^0(X, \mathcal{L})$ 同构地映到
$H^0(D, \mathcal{L}|_D)$，后者的维数为
$\deg(\mathcal{L}) - (g - 1)$。结论由次数定义得出。
\end{proof}



\section{数值相交}
\label{varieties-section-num}

\noindent
本节利用固有概形上凝聚层的欧拉示性数作一些运算，从而得到
可逆模的数值相交数。我们的主要工具是下述引理。

\begin{lemma}
\label{varieties-lemma-numerical-polynomial-from-euler}
设 $k$ 为域。设 $X$ 是 $k$ 上的固有概形。设 $\mathcal{F}$
是凝聚 $\mathcal{O}_X$-模。设
$\mathcal{L}_1, \ldots, \mathcal{L}_r$ 是可逆 $\mathcal{O}_X$-模。
映射
$$
(n_1, \ldots, n_r) \longmapsto
\chi(X, \mathcal{F} \otimes
\mathcal{L}_1^{\otimes n_1} \otimes \ldots \otimes
\mathcal{L}_r^{\otimes n_r})
$$
是关于 $n_1, \ldots, n_r$ 的数值多项式，其总次数不超过
$\mathcal{F}$ 的支撑的维数。
\end{lemma}

\begin{proof}
我们对 $\dim(\text{Supp}(\mathcal{F}))$ 作归纳。若该数为零，
则由引理 \ref{varieties-lemma-chi-tensor-finite}，此函数为常值函数，
其值为 $\dim_k \Gamma(X, \mathcal{F})$。假设
$\dim(\text{Supp}(\mathcal{F})) > 0$。

\medskip\noindent
若 $\mathcal{F}$ 有嵌入伴随点，则可考虑《除子》引理
\ref{divisors-lemma-remove-embedded-points} 中构造的短正合列
$0 \to \mathcal{K} \to \mathcal{F} \to \mathcal{F}' \to 0$。
由于 $\mathcal{K}$ 的支撑维数严格较小，由归纳假设，结论对
$\mathcal{K}$ 成立，且所得总次数严格较小。
由欧拉示性数的可加性（引理
\ref{varieties-lemma-euler-characteristic-additive}），只需对
$\mathcal{F}'$ 证明结论。因此可以假设 $\mathcal{F}$ 没有嵌入伴随点。

\medskip\noindent
若 $i : Z \to X$ 是闭浸入，且 $\mathcal{F} = i_*\mathcal{G}$，
则关于 $X$、$\mathcal{F}$、
$\mathcal{L}_1, \ldots, \mathcal{L}_r$ 的结论等价于关于
$Z$、$\mathcal{G}$、
$i^*\mathcal{L}_1, \ldots, i^*\mathcal{L}_r$ 的结论
（因为二者的上同调相同；见《概形上同调》引理
\ref{coherent-lemma-relative-affine-cohomology}）。
应用《除子》引理 \ref{divisors-lemma-no-embedded-points-endos}，
可以假设 $X$ 没有嵌入分支，并且
$X = \text{Supp}(\mathcal{F})$。

\medskip\noindent
选取正则亚纯截面 $s$，它是 $\mathcal{L}_1$ 的截面；见《除子》引理
\ref{divisors-lemma-regular-meromorphic-section-exists}。
令 $\mathcal{I} \subset \mathcal{O}_X$ 为 $s$ 的分母理想，
并考虑《除子》引理
\ref{divisors-lemma-make-maps-regular-section} 中的映射
$$
\mathcal{I}\mathcal{F} \to \mathcal{F},\quad
\mathcal{I}\mathcal{F} \to \mathcal{F} \otimes \mathcal{L}_1
$$
它们都是单射，余核分别为 $\mathcal{Q}$、$\mathcal{Q}'$，
这些余核支撑在
$X = \text{Supp}(\mathcal{F})$ 的无处稠密闭子概形上。
与可逆模
$\mathcal{L}_1^{\otimes n_1} \otimes \ldots \otimes \mathcal{L}_r^{\otimes n_r}$
张量是正合的；因此再次使用可加性可知
\begin{align*}
&\chi(X, \mathcal{F} \otimes \mathcal{L}_1^{\otimes n_1} \otimes \ldots \otimes
\mathcal{L}_r^{\otimes n_r}) -
\chi(X, \mathcal{F} \otimes \mathcal{L}_1^{\otimes n_1 + 1}
\otimes \ldots \otimes \mathcal{L}_r^{\otimes n_r}) \\
& =
\chi(\mathcal{Q} \otimes \mathcal{L}_1^{\otimes n_1} \otimes \ldots \otimes
\mathcal{L}_r^{\otimes n_r}) -
\chi(\mathcal{Q}' \otimes \mathcal{L}_1^{\otimes n_1} \otimes \ldots \otimes
\mathcal{L}_r^{\otimes n_r})
\end{align*}
因此，引理中的函数 $P(n_1, \ldots, n_r)$ 具有如下性质：
$$
P(n_1 + 1, n_2, \ldots, n_r) - P(n_1, \ldots, n_r)
$$
是总次数 $<$ $\mathcal{F}$ 支撑维数的数值多项式。
当然，由对称性，对任意 $i \in \{1, \ldots, r\}$，同一结论也适用于
$$
P(n_1, \ldots, n_{i - 1}, n_i + 1, n_{i + 1}, \ldots, n_r)
- P(n_1, \ldots, n_r)
$$
一个简单的算术论证表明，$P$ 是总次数不超过
$\dim(\text{Supp}(\mathcal{F}))$ 的数值多项式。
\end{proof}

\noindent
下述引理粗略地说明，首项系数只依赖于凝聚模在其支撑泛点处的长度。

\begin{lemma}
\label{varieties-lemma-numerical-polynomial-leading-term}
设 $k$ 为域。设 $X$ 是 $k$ 上的固有概形。设
$\mathcal{F}$ 是凝聚 $\mathcal{O}_X$-模。设
$\mathcal{L}_1, \ldots, \mathcal{L}_r$ 是可逆 $\mathcal{O}_X$-模。
令 $d = \dim(\text{Supp}(\mathcal{F}))$。
设 $Z_i \subset X$ 是 $\text{Supp}(\mathcal{F})$ 的维数为 $d$
的不可约分支。令 $\xi_i \in Z_i$ 为泛点，并置
$m_i = \text{length}_{\mathcal{O}_{X, \xi_i}}(\mathcal{F}_{\xi_i})$。
则
$$
\chi(X, \mathcal{F} \otimes \mathcal{L}_1^{\otimes n_1} \otimes \ldots \otimes
\mathcal{L}_r^{\otimes n_r}) -
\sum\nolimits_i
m_i\ \chi(Z_i, \mathcal{L}_1^{\otimes n_1} \otimes \ldots \otimes
\mathcal{L}_r^{\otimes n_r}|_{Z_i})
$$
是关于 $n_1, \ldots, n_r$ 的数值多项式，其总次数 $< d$。
\end{lemma}

\begin{proof}
考虑二元组 $(\xi , Z)$，其中 $Z \subset X$ 是维数为 $d$
的整闭子概形，而 $\xi$ 是其泛点。有限
$\mathcal{O}_{X, \xi}$-模 $\mathcal{F}_\xi$ 的支撑包含在
$\{\xi\}$ 中；因此长度
$m_Z = \text{length}_{\mathcal{O}_{X, \xi}}(\mathcal{F}_\xi)$
有限（《代数》引理 \ref{algebra-lemma-support-point}），
并且除非 $Z = Z_i$ 对某个 $i$ 成立，否则它为零。
所以引理中的表达式可写成
$$
E(\mathcal{F}) =
\chi(X, \mathcal{F} \otimes \mathcal{L}_1^{\otimes n_1} \otimes \ldots \otimes
\mathcal{L}_r^{\otimes n_r}) -
\sum\nolimits
m_Z\ \chi(Z, \mathcal{L}_1^{\otimes n_1} \otimes \ldots \otimes
\mathcal{L}_r^{\otimes n_r}|_Z)
$$
其中求和遍历所有整闭子概形 $Z \subset X$，其维数为 $d$。
赋值 $\mathcal{F} \mapsto E(\mathcal{F})$ 在短正合列
$0 \to \mathcal{F} \to \mathcal{F}' \to \mathcal{F}'' \to 0$
中可加；这些列由凝聚 $\mathcal{O}_X$-模组成，其支撑维数
$\leq d$。
这由欧拉示性数的可加性（引理
\ref{varieties-lemma-euler-characteristic-additive}）与长度的可加性
（《代数》引理 \ref{algebra-lemma-length-additive}）得出。
应用《概形上同调》引理
\ref{coherent-lemma-coherent-filter}，取滤过
$$
0 = \mathcal{F}_0 \subset \mathcal{F}_1 \subset
\ldots \subset \mathcal{F}_m = \mathcal{F}
$$
其中各项均为凝聚子层，并且对每个 $j = 1, \ldots, m$，
都存在整闭子概形 $V_j \subset X$ 与非零理想层
$\mathcal{I}_j \subset \mathcal{O}_{V_j}$，使得
$$
\mathcal{F}_j/\mathcal{F}_{j - 1}
\cong (V_j \to X)_* \mathcal{I}_j
$$
于是 $V_j \subset \text{Supp}(\mathcal{F})$，故
$\dim(V_j) \leq d$。
由上述可加性，只须对每个次商
$\mathcal{F}_j/\mathcal{F}_{j - 1}$ 证明结论。
所以只须考虑
$\mathcal{F} = (V \to X)_*\mathcal{I}$，其中
$V \subset X$ 是维数 $\leq d$ 的整闭子概形，而
$\mathcal{I} \subset \mathcal{O}_V$ 是非零凝聚理想层。
若 $\dim(V) < d$，更一般地，若 $\mathcal{F}$ 的支撑维数
$< d$，则由引理 \ref{varieties-lemma-numerical-polynomial-from-euler}，
$E(\mathcal{F})$ 中第一项的总次数 $< d$，第二项为零。
若 $\dim(V) = d$，则可使用短正合列
$$
0 \to (V \to X)_*\mathcal{I} \to (V \to X)_*\mathcal{O}_V
\to (V \to X)_*(\mathcal{O}_V/\mathcal{I}) \to 0
$$
结论对中间的层成立，因为求和中唯一出现的 $Z$ 是 $Z = V$，
且 $m_Z = 1$；此外，由投影公式
（《同调》第 \ref{cohomology-section-projection-formula} 节）
与《概形上同调》引理
\ref{coherent-lemma-relative-affine-cohomology}，有
$$
H^i(X, ((V \to X)_*\mathcal{O}_V) \otimes 
 \mathcal{L}_1^{\otimes n_1} \otimes \ldots \otimes
\mathcal{L}_r^{\otimes n_r}) =
H^i(V,  \mathcal{L}_1^{\otimes n_1} \otimes \ldots \otimes
\mathcal{L}_r^{\otimes n_r}|_V)
$$
因此在此情形中，实际上 $E(\mathcal{F}) = 0$。
结论对右端的层也成立，因为其支撑维数 $< d$。
所以结论对左端的层成立，引理得证。
\end{proof}

\begin{definition}
\label{varieties-definition-intersection-number}
设 $k$ 为域。设 $X$ 是 $k$ 上的固有概形。设
$i : Z \to X$ 是维数为 $d$ 的闭子概形。设
$\mathcal{L}_1, \ldots, \mathcal{L}_d$ 是可逆
$\mathcal{O}_X$-模。定义{\it 相交数}
$(\mathcal{L}_1 \cdots \mathcal{L}_d \cdot Z)$
为下述数值多项式中 $n_1 \ldots n_d$ 的系数：
$$
\chi(X, i_*\mathcal{O}_Z \otimes \mathcal{L}_1^{\otimes n_1} \otimes
\ldots \otimes \mathcal{L}_d^{\otimes n_d}) =
\chi(Z, \mathcal{L}_1^{\otimes n_1} \otimes
\ldots \otimes \mathcal{L}_d^{\otimes n_d}|_Z)
$$
在特殊情形
$\mathcal{L}_1 = \ldots = \mathcal{L}_d = \mathcal{L}$ 下，
记作 $(\mathcal{L}^d \cdot Z)$。
\end{definition}

\noindent
定义中的显示等式由投影公式
（《同调》第 \ref{cohomology-section-projection-formula} 节）
与《概形上同调》引理
\ref{coherent-lemma-relative-affine-cohomology} 得出。
下面证明关于这些相交数的几个引理。

\begin{lemma}
\label{varieties-lemma-intersection-number-integer}
在定义 \ref{varieties-definition-intersection-number} 的情形中，相交数
$(\mathcal{L}_1 \cdots \mathcal{L}_d \cdot Z)$ 是整数。
\end{lemma}

\begin{proof}
任意次数为 $e$、关于 $n_1, \ldots, n_d$ 的数值多项式，
都可以唯一地写成 $\mathbf{Z}$-线性组合，所用函数为
${n_1 \choose k_1}{n_2 \choose k_2} \ldots {n_d \choose k_d}$，其中
$k_1 + \ldots + k_d \leq e$。取 $e = d$ 即可。
细节留作练习。
\end{proof}

\begin{lemma}
\label{varieties-lemma-intersection-number-additive}
在定义 \ref{varieties-definition-intersection-number} 的情形中，相交数
$(\mathcal{L}_1 \cdots \mathcal{L}_d \cdot Z)$ 可加：
若 $\mathcal{L}_i = \mathcal{L}_i' \otimes \mathcal{L}_i''$，则
$$
(\mathcal{L}_1 \cdots \mathcal{L}_i \cdots \mathcal{L}_d \cdot Z) =
(\mathcal{L}_1 \cdots \mathcal{L}_i' \cdots \mathcal{L}_d \cdot Z) +
(\mathcal{L}_1 \cdots \mathcal{L}_i'' \cdots \mathcal{L}_d \cdot Z)
$$
\end{lemma}

\begin{proof}
这是因为，由引理 \ref{varieties-lemma-numerical-polynomial-from-euler}，函数
\begingroup\small
$$
(n_1, \ldots, n_{i - 1}, n_i', n_i'', n_{i + 1}, \ldots, n_d)
\mapsto
\chi(Z, \mathcal{L}_1^{\otimes n_1} \otimes
\ldots \otimes (\mathcal{L}_i')^{\otimes n_i'} \otimes
(\mathcal{L}_i'')^{\otimes n_i''} \otimes \ldots \otimes
\mathcal{L}_d^{\otimes n_d}|_Z)
$$
\endgroup
是总次数至多为 $d$ 的数值多项式，变量个数为 $d + 1$。
\end{proof}


\begin{lemma}
\label{varieties-lemma-intersection-number-in-terms-of-components}
在定义 \ref{varieties-definition-intersection-number} 的情形中，
设 $Z_i \subset Z$ 是维数为 $d$ 的不可约分支。令
$m_i = \text{length}_{\mathcal{O}_{X, \xi_i}}(\mathcal{O}_{Z, \xi_i})$，
其中 $\xi_i \in Z_i$ 为泛点。则
$$
(\mathcal{L}_1 \cdots \mathcal{L}_d \cdot Z) =
\sum m_i(\mathcal{L}_1 \cdots \mathcal{L}_d \cdot Z_i)
$$
\end{lemma}

\begin{proof}
这立即由引理 \ref{varieties-lemma-numerical-polynomial-leading-term} 与各定义得出。
\end{proof}

\begin{lemma}
\label{varieties-lemma-intersection-number-and-pullback}
设 $k$ 为域。设 $f : Y \to X$ 是 $k$ 上固有概形之间的态射。
设 $Z \subset Y$ 是维数为 $d$ 的整闭子概形，并设
$\mathcal{L}_1, \ldots, \mathcal{L}_d$ 是可逆 $\mathcal{O}_X$-模。
则
$$
(f^*\mathcal{L}_1 \cdots f^*\mathcal{L}_d \cdot Z) =
\deg(f|_Z : Z \to f(Z)) (\mathcal{L}_1 \cdots \mathcal{L}_d \cdot f(Z))
$$
其中 $\deg(Z \to f(Z))$ 的定义见《态射》定义
\ref{morphisms-definition-degree}；也可将其取为 $0$，条件是
$\dim(f(Z)) < d$。
\end{lemma}

\begin{proof}
左端由下述函数中 $n_1 \ldots n_d$ 的系数计算：
$$
\chi(Y, \mathcal{O}_Z \otimes f^*\mathcal{L}_1^{\otimes n_1} \otimes
\ldots \otimes f^*\mathcal{L}_d^{\otimes n_d}) =
\sum (-1)^i
\chi(X, R^if_*\mathcal{O}_Z \otimes
\mathcal{L}_1^{\otimes n_1} \otimes \ldots \otimes
\mathcal{L}_d^{\otimes n_d})
$$
该等式由引理 \ref{varieties-lemma-euler-characteristic-morphism} 与投影公式
（《同调》引理 \ref{cohomology-lemma-projection-formula}）得出。
若 $f(Z)$ 的维数 $< d$，则由引理
\ref{varieties-lemma-numerical-polynomial-from-euler}，右端是总次数
$< d$ 的多项式，故结论成立。假设 $\dim(f(Z)) = d$。
令 $\xi \in f(Z)$ 为泛点。由维数理论（见引理
\ref{varieties-lemma-dimension-locally-algebraic} 与
\ref{varieties-lemma-dimension-fibres-locally-algebraic}），
$Z$ 的泛点是 $Z$ 中唯一映到 $\xi$ 的点。
于是 $f : Z \to f(Z)$ 在 $f(Z)$ 的某个非空开集上有限；见《态射》引理
\ref{morphisms-lemma-generically-finite}。
因此 $\deg(f : Z \to f(Z))$ 有定义；事实上，它等于
$f_*\mathcal{O}_Z$ 在 $\xi$ 处的茎作为
$\mathcal{O}_{X, \xi}$-模的长度。此外，
$R^if_*\mathcal{O}_X$ 在 $\xi$ 处的茎对 $i > 0$ 为零，
因为刚刚已经看到
$f|_Z$ 在 $\xi$ 的一个邻域上有限
（因而《概形上同调》引理
\ref{coherent-lemma-finite-pushforward-coherent} 给出该消失）。
所以各项
$\chi(X, R^if_*\mathcal{O}_Z \otimes
\mathcal{L}_1^{\otimes n_1} \otimes \ldots \otimes
\mathcal{L}_d^{\otimes n_d})$（其中 $i > 0$）的总次数 $< d$，并且由引理
\ref{varieties-lemma-numerical-polynomial-leading-term}，模掉一个总次数
$< d$ 的多项式后，有
$$
\chi(X, f_*\mathcal{O}_Z \otimes
\mathcal{L}_1^{\otimes n_1} \otimes \ldots \otimes
\mathcal{L}_d^{\otimes n_d})
=
\deg(f : Z \to f(Z)) \chi(f(Z),
\mathcal{L}_1^{\otimes n_1} \otimes \ldots \otimes
\mathcal{L}_d^{\otimes n_d}|_{f(Z)})
$$
所需结论随即得出。
\end{proof}

\begin{lemma}
\label{varieties-lemma-numerical-intersection-effective-Cartier-divisor}
设 $k$ 为域。设 $X$ 在 $k$ 上固有。设 $Z \subset X$
是维数为 $d$ 的闭子概形。设
$\mathcal{L}_1, \ldots, \mathcal{L}_d$ 是可逆
$\mathcal{O}_X$-模。假设存在有效 Cartier 除子
$D \subset Z$，使得 $\mathcal{L}_1|_Z \cong \mathcal{O}_Z(D)$。
则
$$
(\mathcal{L}_1 \cdots \mathcal{L}_d \cdot Z) =
(\mathcal{L}_2 \cdots \mathcal{L}_d \cdot D)
$$
\end{lemma}

\begin{proof}
可以把 $X$ 替换为 $Z$，并把 $\mathcal{L}_i$ 替换为
$\mathcal{L}_i|_Z$。因此可以假设
$X = Z$ 且 $\mathcal{L}_1 = \mathcal{O}_X(D)$。
此时 $\mathcal{L}_1^{-1}$ 是 $D$ 的理想层，并可考虑短正合列
$$
0 \to \mathcal{L}_1^{\otimes -1} \to \mathcal{O}_X \to \mathcal{O}_D \to 0
$$
置
$P(n_1, \ldots, n_d) =
\chi(X, \mathcal{L}_1^{\otimes n_1} \otimes \ldots \otimes
\mathcal{L}_d^{\otimes n_d})$
以及
$Q(n_1, \ldots, n_d) =
\chi(D, \mathcal{L}_1^{\otimes n_1} \otimes \ldots \otimes
\mathcal{L}_d^{\otimes n_d}|_D)$。
由可加性可得
$$
P(n_1, \ldots, n_d) - P(n_1 - 1, n_2, \ldots, n_d) =
Q(n_1, \ldots, n_d)
$$
由于 $P$ 的总次数至多为 $d$，可知 $n_1 \ldots n_d$
在 $P$ 中的系数等于 $n_2 \ldots n_d$ 在 $Q$ 中的系数。
\end{proof}

\begin{lemma}
\label{varieties-lemma-ample-positive}
设 $k$ 为域。设 $X$ 在 $k$ 上固有。设 $Z \subset X$
是维数为 $d$ 的闭子概形。若
$\mathcal{L}_1, \ldots, \mathcal{L}_d$ 都丰沛，则
$(\mathcal{L}_1 \cdots \mathcal{L}_d \cdot Z)$ 为正。
\end{lemma}

\begin{proof}
我们对 $d$ 作归纳。$d = 0$ 的情形由引理
\ref{varieties-lemma-chi-tensor-finite} 得出。假设 $d > 0$。
由引理 \ref{varieties-lemma-intersection-number-in-terms-of-components}，
可以假设 $Z$ 是整闭子概形。事实上，可以把 $X$ 替换为 $Z$，
并把 $\mathcal{L}_i$ 替换为 $\mathcal{L}_i|_Z$，
从而约化到 $Z = X$ 是维数为 $d$ 的固有簇的情形。
由引理 \ref{varieties-lemma-intersection-number-additive}，可以把
$\mathcal{L}_1$ 替换为它的某个正张量幂。
所以可以假设存在非零截面
$s \in \Gamma(X, \mathcal{L}_1)$，使得
$X_s$ 仿射（这里使用丰沛可逆层的定义；见《性质》定义
\ref{properties-definition-ample}）。
注意，$X$ 不仿射，因为固有且仿射蕴含有限
（《态射》引理 \ref{morphisms-lemma-finite-proper}），
这与 $d > 0$ 矛盾。因此 $s$ 有非空消失概形
$Z(s) \subset X$。由于 $X$ 是簇，$s$ 是
$\mathcal{L}_1$ 的正则截面，所以 $Z(s)$ 是有效 Cartier 除子；
因而 $Z(s)$ 的余维为 $1$（这是在 $X$ 中而言），所以 $Z(s)$ 的维数为
$d - 1$（这里使用《除子》第
\ref{divisors-section-effective-Cartier-divisors}、
\ref{divisors-section-effective-Cartier-invertible}、
\ref{divisors-section-Noetherian-effective-Cartier} 节中的材料，
以及引理 \ref{varieties-lemma-dimension-locally-algebraic} 所用的维数理论）。
由引理 \ref{varieties-lemma-numerical-intersection-effective-Cartier-divisor}，有
$$
(\mathcal{L}_1 \cdots \mathcal{L}_d \cdot X) =
(\mathcal{L}_2 \cdots \mathcal{L}_d \cdot Z(s))
$$
由归纳假设，右端为正，证明完成。
\end{proof}

\begin{definition}
\label{varieties-definition-degree}
设 $k$ 为域。设 $X$ 是 $k$ 上的固有概形。设
$\mathcal{L}$ 是丰沛可逆 $\mathcal{O}_X$-模。
对任意闭子概形 $Z$，其关于 $\mathcal{L}$ 的{\it 次数}记作
$\deg_\mathcal{L}(Z)$，并定义为相交数
$(\mathcal{L}^d \cdot Z)$，其中 $d = \dim(Z)$。
\end{definition}

\noindent
由引理 \ref{varieties-lemma-ample-positive}，子概形的次数总是正整数。
注意，$\deg_\mathcal{L}(Z) = d$ 当且仅当
$$
\chi(Z, \mathcal{L}^{\otimes n}|_Z) = \frac{d}{\dim(Z)!} n^{\dim(Z)} + l.o.t
$$
这可由下式看出：
$$
(n_1 + \ldots + n_{\dim(Z)})^{\dim(Z)} =
\dim(Z)!\ n_1 \ldots n_{\dim(Z)} + \text{other terms}
$$

\begin{lemma}
\label{varieties-lemma-degree-finite-morphism-in-terms-degrees}
设 $k$ 为域。设 $f : Y \to X$ 是 $k$ 上固有簇之间的有限支配态射。
设 $\mathcal{L}$ 是丰沛可逆 $\mathcal{O}_X$-模。则
$$
\deg_{f^*\mathcal{L}}(Y) = \deg(f) \deg_\mathcal{L}(X)
$$
其中 $\deg(f)$ 的定义见《态射》定义
\ref{morphisms-definition-degree}。
\end{lemma}

\begin{proof}
该陈述有意义，因为由《态射》引理
\ref{morphisms-lemma-pullback-ample-tensor-relatively-ample}，
$f^*\mathcal{L}$ 丰沛。说明这一点后，结论就是引理
\ref{varieties-lemma-intersection-number-and-pullback} 的特殊情形。
\end{proof}

\noindent
最后，我们把与曲线的相交数同第
\ref{varieties-section-divisors-curves} 节中引入的曲线上可逆模次数联系起来。

\begin{lemma}
\label{varieties-lemma-intersection-numbers-and-degrees-on-curves}
设 $k$ 为域。设 $X$ 是 $k$ 上的固有概形。
设 $Z \subset X$ 是维数 $\leq 1$ 的闭子概形。
设 $\mathcal{L}$ 是可逆 $\mathcal{O}_X$-模。则
$$
(\mathcal{L} \cdot Z) = \deg(\mathcal{L}|_Z)
$$
其中 $\deg(\mathcal{L}|_Z)$ 的定义见定义
\ref{varieties-definition-degree-invertible-sheaf}。若 $\mathcal{L}$ 丰沛，则
$\deg_\mathcal{L}(Z) = \deg(\mathcal{L}|_Z)$。
\end{lemma}

\begin{proof}
这是因为函数
$n \mapsto \chi(Z, \mathcal{L}|_Z^{\otimes n})$ 的次数为 $1$，
故其首项系数等于相邻两值之差。
\end{proof}

\begin{proposition}[渐近 Riemann--Roch]
\label{varieties-proposition-asymptotic-riemann-roch}
设 $k$ 为域。设 $X$ 是 $k$ 上维数为 $d$ 的固有概形。
设 $\mathcal{L}$ 是丰沛可逆 $\mathcal{O}_X$-模。则
$$
\dim_k \Gamma(X, \mathcal{L}^{\otimes n}) \sim c n^d + l.o.t.
$$
其中 $c = \deg_\mathcal{L}(X)/d!$ 是正的常数。
\end{proposition}

\begin{proof}
这由各定义、引理 \ref{varieties-lemma-ample-positive}，以及《概形上同调》引理
\ref{coherent-lemma-vanshing-gives-ample} 中的高次上同调消失得出。
\end{proof}

\section{嵌入维数}
\label{varieties-section-embedding-dimension}

\noindent
嵌入维数有若干种定义方式，但对于代数闭域上的代数概形的闭点，
所有定义都与下述定义等价。

\begin{definition}
\label{varieties-definition-embed-dim}
设 $k$ 是代数闭域。设 $X$ 是局部代数 $k$-概形，并设 $x \in X$ 是闭点。
定义{\it $X$ 在 $x$ 处的嵌入维数}为
$\dim_k \mathfrak m_x/\mathfrak m_x^2$。
\end{definition}

\noindent
下面列出关于嵌入维数的一些事实。
设 $k, X, x$ 如定义 \ref{varieties-definition-embed-dim} 所述。
\begin{enumerate}
\item $X$ 在 $x$ 处的嵌入维数，等于切空间
$T_{X/k, x}$（定义 \ref{varieties-definition-tangent-space}）作为 $k$-向量空间的维数。
\item $X$ 在 $x$ 处的嵌入维数，是使得存在如下满射的最小整数
$d \geq 0$：
$$
k[[x_1, \ldots, x_d]] \longrightarrow \mathcal{O}_{X, x}^\wedge
$$
这里要求它是 $k$-代数满射。
\item $X$ 在 $x$ 处的嵌入维数，是最小整数 $d \geq 0$，使得存在开邻域
$U \subset X$（它是 $x$ 的邻域）以及闭浸入 $U \to Y$，其中 $Y$ 是维数为 $d$
的 $k$ 上光滑簇。
\item $X$ 在 $x$ 处的嵌入维数，是最小整数 $d \geq 0$，使得存在开邻域
$U \subset X$（它是 $x$ 的邻域）以及非分歧态射 $U \to \mathbf{A}^d_k$。
\item 若给定闭嵌入 $X \to Y$，且 $Y$ 在 $k$ 上光滑，则 $X$ 在 $x$ 处的嵌入维数
是满足下述条件的最小整数 $d \geq 0$：存在闭子概形 $Z \subset Y$，使得
$X \subset Z$，$Z \to \Spec(k)$ 在 $x$ 处光滑，并且 $\dim_x(Z) = d$。
\end{enumerate}
若以后需要这些事实，我们会陈述精确结果并给出证明。

\medskip\noindent
非代数闭的基域或非闭点。
设 $k$ 是域，且 $X$ 是局部代数 $k$-概形。
若 $x \in X$ 是一点，则定义 $X$ 在 $x$ 处的嵌入维数有若干种选择。具体说，可以采用
\begin{enumerate}
\item $\dim_{\kappa(x)}(\mathfrak m_x/\mathfrak m_x^2)$，
\item $\dim_{\kappa(x)}(T_{X/k, x}) =
\dim_{\kappa(x)}(\Omega_{X/k, x} \otimes_{\mathcal{O}_{X, x}} \kappa(x))$
（引理 \ref{varieties-lemma-tangent-space-cotangent-space}），
\item 最小整数 $d \geq 0$，使得存在开邻域 $U \subset X$（它是 $x$ 的邻域）
以及闭浸入 $U \to Y$，其中 $Y$ 是维数为 $d$ 的 $k$ 上光滑簇。
\end{enumerate}
在特征零的情形，(1) $=$ (2) 在 $x$ 是闭点时成立；
更一般地，若 $\kappa(x)$ 是 $k$ 上的可分代数扩张，这一等式也成立，
见引理 \ref{varieties-lemma-tangent-space-rational-point}。
几何定义 (3) 似乎最贴近“嵌入维数”这一措辞所唤起的几何直观。
由于可以证明 (3) 与 (2) 定义同一个数
（这由引理 \ref{varieties-lemma-immersion-into-affine} 得出），
我们将采用这个定义。在术语中，我们会明确说明所取的是相对于基域的嵌入维数。

\begin{definition}
\label{varieties-definition-embedding-dimension}
设 $k$ 是域。设 $X$ 是局部代数 $k$-概形。
设 $x \in X$ 是一点。定义{\it $X/k$ 在 $x$ 处的嵌入维数}为
$\dim_{\kappa(x)}(T_{X/k, x})$。
\end{definition}

\noindent
若 $(A, \mathfrak m, \kappa)$ 是 Noether 局部环，
有时将 $A$ 的{\it 嵌入维数}定义为 $\mathfrak m/\mathfrak m^2$ 在 $\kappa$ 上的维数。
上文已经看到，若 $A$ 具有域 $k$ 上的代数结构，采用
$\dim_\kappa(\Omega_{A/k} \otimes_A \kappa)$ 可能更合适。
我们将此量称为 $A/k$ 的{\it 嵌入维数}。采用这一术语，有
$$
\text{embed dim of }X/k\text{ 在 }x =
\text{embed dim of }\mathcal{O}_{X, x}/k =
\text{embed dim of }\mathcal{O}_{X, x}^\wedge/k
$$
其中 $k, X, x$ 如定义 \ref{varieties-definition-embedding-dimension} 所述。

\section{Bertini 定理}
\label{varieties-section-bertini}

\noindent
本节证明如下类型的结果：给定一个光滑射影簇 $X$，其基域为 $k$，
存在一个光滑的丰沛除子 $H \subset X$。

\begin{lemma}
\label{varieties-lemma-generate-over-complement}
设 $k$ 是域。设 $X$ 是 $k$ 上的本征概形。设 $\mathcal{L}$ 是丰沛可逆
$\mathcal{O}_X$-模。设 $Z \subset X$ 是闭子概形。则存在整数 $n_0$，使得对所有
$n \geq n_0$，令 $V_n$ 表示映射
$\Gamma(X, \mathcal{L}^{\otimes n}) \to \Gamma(Z, \mathcal{L}^{\otimes n}|_Z)$
的核；它生成 $\mathcal{L}^{\otimes n}|_{X \setminus Z}$，并且典范态射
$$
X \setminus Z \longrightarrow \mathbf{P}(V_n)
$$
是 $k$ 上概形的浸入。
\end{lemma}

\begin{proof}
设 $\mathcal{I} \subset \mathcal{O}_X$ 是 $Z$ 的拟凝聚理想层。注意，通过包含关系
$\mathcal{I} \otimes_{\mathcal{O}_X} \mathcal{L}^{\otimes n} \subset
\mathcal{L}^{\otimes n}$，有
$V_n = \Gamma(X, \mathcal{I} \otimes_{\mathcal{O}_X} \mathcal{L}^{\otimes n})$。
选取 $n_1$，使得对 $n \geq n_1$，层
$\mathcal{I} \otimes \mathcal{L}^{\otimes n}$ 是全局生成的，见
《性质》，命题 \ref{properties-proposition-characterize-ample}。
于是 $V_n$ 生成 $\mathcal{L}^{\otimes n}|_{X \setminus Z}$，只要 $n \geq n_1$。

\medskip\noindent
对 $n \geq n_1$，将限制映射记为
$\psi_n : V_n \to
\Gamma(X \setminus Z, \mathcal{L}^{\otimes n}|_{X \setminus Z})$。
由《构造》，例 \ref{constructions-example-projective-space}，得到典范态射
$$
\varphi = \varphi_{\mathcal{L}^{\otimes n}|_{X \setminus Z}, \psi_n} :
X \setminus Z
\longrightarrow
\mathbf{P}(V_n)
$$
。选取 $n_2$，使得对所有 $n \geq n_2$，可逆层
$\mathcal{L}^{\otimes n}$ 在 $X$ 上很丰沛。
我们断言 $n_0 = n_1 + n_2$ 可用。

\medskip\noindent
断言的证明。设 $n \geq n_0$，并写成 $n = n_1 + n'$。
对 $x \in X \setminus Z$，可选取 $s_1 \in V_1$，使它在 $x$ 处不消失。
置 $V' = \Gamma(X, \mathcal{L}^{\otimes n'})$。
由对 $n$ 与 $n'$ 的选取可知，相应态射
$\varphi' : X \to \mathbf{P}(V')$ 是闭浸入。因此，若选取
$s' \in \Gamma(X, \mathcal{L}^{\otimes n'})$ 在 $x$ 处不消失，则
$X_{s'} = (\varphi')^{-1}(D_+(s'))$（见《构造》，
引理 \ref{constructions-lemma-invertible-map-into-proj}）是仿射的，
且 $X_{s'} \to D_+(s')$ 是闭浸入。
于是 $s = s_1 \otimes s' \in V_n$ 在 $x$ 处不消失。
若 $D_+(s) \subset \mathbf{P}(V_n)$ 表示该射影空间中相应的开仿射空间，则
$\varphi^{-1}(D_+(s)) = X_s \subset X \setminus Z$（见上引文）。
开集 $X_s = X_{s'} \cap X_{s_1}$ 是仿射的，见
《性质》，引理 \ref{properties-lemma-affine-cap-s-open}。
考虑定义态射 $X_s \to D_+(s)$ 的环映射
$$
\text{Sym}(V)_{(s)} \longrightarrow \mathcal{O}_X(X_s)
$$
。因为 $X_{s'} \to D_+(s')$ 是闭浸入，以下元素的像
$$
\frac{s_1 \otimes t'}{s_1 \otimes s'}
$$
（其中 $t' \in V'$）生成映射
$\mathcal{O}_X(X_{s'}) \to \mathcal{O}_X(X_s)$ 的像。
由于 $X_s \to X_{s'}$ 是开浸入，这说明 $X_s \to D_+(s)$
是仿射概形的浸入（见下文）。从而由《态射》，
引理 \ref{morphisms-lemma-check-immersion}，$\varphi_n$ 是浸入。

\medskip\noindent
设 $a : A' \to A$ 与 $c : B \to A$ 是环映射，且
$\Spec(a)$ 是浸入并满足 $\Im(a) \subset \Im(c)$。
置 $B' = A' \times_A B$，其投影为 $b : B' \to B$ 与
$c' : B' \to A'$。由假设，$c'$ 是满射，因而
$\Spec(c')$ 是闭浸入。于是 $\Spec(c') \circ \Spec(a)$ 是浸入
（《概形》，引理 \ref{schemes-lemma-composition-immersion}）。
由于 $\Spec(c)$ 分解浸入
$\Spec(c') \circ \Spec(a) = \Spec(b) \circ \Spec(c)$，
故它必为浸入，见《态射》，引理 \ref{morphisms-lemma-immersion-permanence}。
\end{proof}

\begin{situation}
\label{varieties-situation-family-divisors}
设 $k$ 是域，
设 $X$ 是 $k$ 上的概形，
设 $\mathcal{L}$ 是可逆 $\mathcal{O}_X$-模，
设 $V$ 是有限维 $k$-向量空间，并且
设 $\psi : V \to \Gamma(X, \mathcal{L})$ 是 $k$-线性映射。
设 $\dim(V) = r$，且 $v_1, \ldots, v_r$ 是 $V$ 的一组基。
于是得到“泛除子”
$$
H_{univ} = Z(s_{univ}) \subset  \mathbf{A}^r \times_k X
$$
，它是截面
$$
s_{univ} = \sum\nolimits_{i = 1, \ldots, r} x_i \psi(v_i) \in
\Gamma(\mathbf{A}^r \times_k X, \text{pr}_2^*\mathcal{L})
$$
的零概形（《除子》，定义 \ref{divisors-definition-zero-scheme-s}）。
对域扩张 $k'/k$，$k'$-点 $v \in \mathbf{A}^r_k(k')$
对应于向量 $(a_1, \ldots, a_r)$，其分量均属于 $k'$。
因此，一方面可以把 $v$ 看作元素
$v = \sum_{i = 1, \ldots, r} a_i v_i \in V \otimes_k k'$，
另一方面可以向 $v$ 指派截面
$$
\psi(v) = \sum\nolimits_{i = 1, \ldots, r} a_i \psi(v_i) \in
\Gamma(X_{k'}, \mathcal{L}|_{X_{k'}})
$$
。按此记号，$H_{univ}$ 在 $v \in V \otimes k'$ 上的纤维显然是
$\psi(v)$ 的零概形。写成公式即
$$
H_v = H_{univ, v} = Z(\psi(v))
$$
。如式中所示，将这个公共值记作 $H_v$。最后，在此情形中，令 $P$ 是向量
$v \in V \otimes_k k'$ 的一个性质，其中 $k'/k$ 是任意域扩张\footnote{例如，
可以考虑 $H_v$ 在 $k'$ 上光滑，或在 $k'$ 上几何不可约这一条件。}。
称 $P$ 对{\it 一般的} $v \in V \otimes_k k'$ {\it 成立}，如果存在非空 Zariski 开集
$U \subset \mathbf{A}^r_k$，使得只要 $v$ 对应于一个 $k'$-点且该点属于 $U$，
则对任意 $k'/k$ 都有 $P(v)$ 成立。
\end{situation}

\begin{lemma}
\label{varieties-lemma-bertini}
在情形 \ref{varieties-situation-family-divisors} 中，假设
\begin{enumerate}
\item $X$ 在 $k$ 上光滑，
\item $\psi : V \to \Gamma(X, \mathcal{L})$ 的像生成 $\mathcal{L}$，
\item 相应态射
$\varphi_{\mathcal{L}, \psi} : X \to \mathbf{P}(V)$ 是浸入。
\end{enumerate}
则对一般的 $v \in V \otimes_k k'$，概形 $H_v$ 在 $k'$ 上光滑。
\end{lemma}

\begin{proof}
（注意，$X$ 作为射影空间的局部闭子概形，是分离且有限型的。）
采用引理陈述上方引入的记号。考虑投影
$$
\xymatrix{
\mathbf{A}^r_k \times_k X \ar[d] &
H_{univ} \ar[l] \ar[ld]^p \ar[r] \ar[rd]_q &
\mathbf{A}^r_k \times_k X \ar[d] \\
X & &
\mathbf{A}^r_k
}
$$
。设 $\Sigma \subset H_{univ}$ 是态射
$q : H_{univ} \to \mathbf{A}^r_k$ 的奇异轨迹，即 $q$ 不光滑的点集。
由于态射的光滑轨迹按定义是开的，$\Sigma$ 是闭的。
光滑态射的纤维是光滑的，故只需证明 $q(\Sigma)$ 包含在
$\mathbf{A}^r_k$ 的真闭子集中。
因为 $\Sigma$（赋予约化诱导概形结构）是 $k$ 上有限型概形，
只需证明 $\dim(\Sigma) < r$。这由引理
\ref{varieties-lemma-dimension-fibres-locally-algebraic} 得出。
以更大的域代替 $k$ 不改变维数
（《态射》，引理 \ref{morphisms-lemma-dimension-fibre-after-base-change}），
因此不妨假设 $k$ 是代数闭的。
由维数理论（引理 \ref{varieties-lemma-dimension-fibres-locally-algebraic}），
只需证明：对闭点 $x \in X \setminus Z$，有
$p^{-1}(\{x\}) \cap \Sigma$ 的维数 $< r - \dim(X')$，
其中 $X'$ 是 $X$ 中包含 $x$ 的唯一不可约分支。
由于 $X$ 在 $k$ 上光滑且 $x$ 是闭点，有
$\dim(X') = \dim \mathfrak m_x/\mathfrak m_x^2$
（《态射》，引理 \ref{morphisms-lemma-smooth-omega-finite-locally-free}
及《代数》，引理 \ref{algebra-lemma-rank-omega}）。
因此，只要对所有闭点 $x \in X$ 都有
$$
\dim p^{-1}(x) \cap \Sigma < r - \dim \mathfrak m_x/\mathfrak m_x^2
$$
，结论即成立。

\medskip\noindent
由于 $V$ 全局生成 $\mathcal{L}$，对每个不可约分支
$X'$（它属于 $X$），$\mathbf{A}^r$ 都有一个非空 Zariski 开集，使得 $q$ 在该开集上的纤维不包含
$X'$。（例如，若 $x' \in X'$ 是闭点，则可取对应于下述向量的开集：
$v \in V$，且 $\psi(v)$ 在 $x'$ 处不消失。这个开集是
$\mathbf{A}^r_k$ 中一个超平面的补集。）
令 $U \subset \mathbf{A}^r$ 为这些开集的（非空）交。
则 $q^{-1}(U) \to U$ 的纤维是 $U \times_k X \to U$ 的纤维上的有效 Cartier 除子
（因为整概形上可逆模的非零截面是正则截面）。因此由《除子》，
引理 \ref{divisors-lemma-fibre-Cartier}，态射 $q^{-1}(U) \to U$ 是平坦的。
从而，对闭点 $x \in X$ 与 $v \in V = \mathbf{A}^r_k(k)$，若
$(x, v) \in H_{univ}$，即 $x \in H_v$，则 $q$ 在 $(x, v)$ 处光滑当且仅当纤维
$H_v$ 在 $x$ 处光滑，见《态射》，引理 \ref{morphisms-lemma-smooth-at-point}。

\medskip\noindent
考虑像 $\psi(v)_x$；它位于茎 $\mathcal{L}_x$ 中，来自与 $v \in V$ 对应的截面。有
$$
x \in H_v \Leftrightarrow \psi(v)_x \in \mathfrak m_x\mathcal{L}_x
$$
。若此条件成立，则有
$$
H_v\text{ singular at }x \Leftrightarrow
\psi(v)_x \in \mathfrak m_x^2\mathcal{L}_x
$$
。事实上，$\psi(v)_x$ 不属于 $\mathfrak m_x^2\mathcal{L}_x$
$\Leftrightarrow$
$H_v \subset X$ 在 $x$ 处的局部方程不属于 $\mathfrak m_x^2$
$\Leftrightarrow$
$\mathcal{O}_{H_v, x}$ 是正则的
（《代数》，引理 \ref{algebra-lemma-regular-ring-CM}）
$\Leftrightarrow$
$H_v$ 在点 $x$ 处相对于 $k$ 光滑
（《代数》，引理 \ref{algebra-lemma-separable-smooth}）。
由此，$p^{-1}(x) \cap \Sigma$ 的闭点对应于这些 $v \in V$：它们满足
$\psi(v)_x \in \mathfrak m_x^2\mathcal{L}_x$。
然而，由于 $\varphi_{\mathcal{L}, \psi}$ 是浸入，映射
$$
V \longrightarrow \mathcal{L}_x/\mathfrak m_x^2\mathcal{L}_x
$$
是满射（略去一个小细节）。由上可知，将轨迹 $p^{-1}(x) \cap \Sigma$
视为 $V$ 的子空间时，其闭点构成此映射的核，因而其维数为
$r - \dim \mathfrak m_x/\mathfrak m_x^2 - 1$，如所欲证。
\end{proof}


\section{Enriques--Severi--Zariski 定理}
\label{varieties-section-vanishing-negative}

\noindent
本节证明若干如下形式的结果：以一个“非常负”的可逆模扭转会使低次上同调消失。
我们还推出，维数 $\geq 2$ 的正规本征概形的超曲面截面是连通的。

\begin{lemma}
\label{varieties-lemma-vanishin-h0-negative}
设 $k$ 是域。设 $X$ 是 $k$ 上的本征概形。设 $\mathcal{L}$ 是丰沛可逆
$\mathcal{O}_X$-模。设 $\mathcal{F}$ 是凝聚 $\mathcal{O}_X$-模。
若 $\text{Ass}(\mathcal{F})$ 不含任何闭点，则
$\Gamma(X, \mathcal{F} \otimes_{\mathcal{O}_X} \mathcal{L}^{\otimes n}) = 0$，
只要 $n \ll 0$。
\end{lemma}

\begin{proof}
对凝聚 $\mathcal{O}_X$-模 $\mathcal{F}$，令 $\mathcal{P}(\mathcal{F})$
表示如下性质：存在 $n_0 \in \mathbf{Z}$，使得对 $n \leq n_0$，
每个截面 $s$（它是
$\mathcal{F} \otimes_{\mathcal{O}_X} \mathcal{L}^{\otimes n}$ 的截面）
的支撑都只由闭点组成。
由于 $\text{Ass}(\mathcal{F}) =
\text{Ass}(\mathcal{F} \otimes_{\mathcal{O}_X} \mathcal{L}^{\otimes n})$，
只需证明 $\mathcal{P}$ 对 $X$ 上所有凝聚模都成立。
为此，证明《概形的上同调》，引理
\ref{coherent-lemma-property-higher-rank-cohomological}
中的条件 (1)、(2) 和 (3) 均成立。

\medskip\noindent
为验证条件 (1)，设
$$
0 \to \mathcal{F}_1 \to \mathcal{F} \to \mathcal{F}_2 \to 0
$$
是凝聚 $\mathcal{O}_X$-模的短正合列，并且已有 $\mathcal{P}$
对 $\mathcal{F}_i$ 成立，其中 $i = 1, 2$。
令 $n_1, n_2$ 为所得界限。
令 $\mathcal{F}'_2 \subset \mathcal{F}_2$ 为支撑是有限闭点集的最大凝聚子模。
令 $\mathcal{I} \subset \mathcal{O}_X$ 为 $\mathcal{F}'_2$ 的零化子。
由于 $\mathcal{L}$ 丰沛，可找到 $e > 0$，使
$\mathcal{I} \otimes_{\mathcal{O}_X} \mathcal{L}^{\otimes e}$
是全局生成的。置 $n_0 = \min(n_2, n_1 - e)$。
取 $n \leq n_0$，并令 $t$ 是
$\mathcal{F} \otimes \mathcal{L}^{\otimes n}$ 的一个全局截面。
截面 $t$ 在
$\mathcal{F}_2 \otimes \mathcal{L}^{\otimes n}$ 中的像落在
$\mathcal{F}'_2 \otimes \mathcal{L}^{\otimes n}$ 中，这是因为 $n \leq n_2$。
因此，对任意
$s \in \Gamma(X, \mathcal{I} \otimes_{\mathcal{O}_X} \mathcal{L}^{\otimes e})$，
乘积 $t \otimes s$ 属于
$\mathcal{F}_1 \otimes \mathcal{L}^{\otimes n + e}$。
$t \otimes s$ 的支撑包含在
$\text{Ass}(\mathcal{F}_1)$ 中的有限闭点集中，因为 $n + e \leq n_1$。
依照对 $e$ 的选取，可选取 $s$ 在任何不属于
$\mathcal{F}'_2$ 支撑的点处可逆；故 $t$ 的支撑包含在
$\text{Ass}(\mathcal{F}_1)$ 中的有限闭点集与
$\text{Ass}(\mathcal{F}_2)$ 中的有限闭点集之并中。
这就完成了条件 (1) 的证明。

\medskip\noindent
条件 (2) 是立即的。

\medskip\noindent
对条件 (3)，取 $\mathcal{G} = \mathcal{O}_Z$。
在这种情形下，若 $Z$ 是 $X$ 的闭点，则无须证明。
若 $\dim(Z) > 0$，则证明
$\Gamma(Z, \mathcal{L}^{\otimes n}|_Z) = 0$，只要 $n < 0$。
事实上，令 $s$ 是 $\mathcal{L}|_Z$ 的某个负次幂的非零截面。
选取非零截面 $t$，它属于 $\mathcal{L}|_Z$ 的某个正次幂
（这是可能的，因为 $\mathcal{L}$ 丰沛，见《性质》，命题
\ref{properties-proposition-characterize-ample}）。
则
$s^{\deg(t)} \otimes t^{\deg(s)}$
是 $\mathcal{O}_Z$ 的非零全局截面（因为 $Z$ 是整的），
从而是单位（引理 \ref{varieties-lemma-proper-geometrically-reduced-global-sections}）。
这说明 $t$ 是 $\mathcal{L}$ 的某个正次幂的平凡化截面。
于是函数 $n \mapsto \dim_k \Gamma(X, \mathcal{L}^{\otimes n})$
在一个无穷正整数集上有界；由于 $\dim(Z) > 0$，
这与渐近 Riemann--Roch（命题
\ref{varieties-proposition-asymptotic-riemann-roch}）矛盾。
\end{proof}

\begin{lemma}[Enriques--Severi--Zariski]
\label{varieties-lemma-vanishin-h1-negative}
设 $k$ 是域。设 $X$ 是 $k$ 上的本征概形。设 $\mathcal{L}$ 是丰沛可逆
$\mathcal{O}_X$-模。设 $\mathcal{F}$ 是凝聚 $\mathcal{O}_X$-模。
假设对闭点 $x \in X$ 有 $\text{depth}(\mathcal{F}_x) \geq 2$。
则
$H^1(X, \mathcal{F} \otimes_{\mathcal{O}_X} \mathcal{L}^{\otimes m}) = 0$，
只要 $m \ll 0$。
\end{lemma}

\begin{proof}
选取闭浸入 $i : X \to \mathbf{P}^n_k$，使得
$i^*\mathcal{O}(1) \cong \mathcal{L}^{\otimes e}$，其中 $e > 0$
（见《态射》，引理
\ref{morphisms-lemma-finite-type-over-affine-ample-very-ample}）。
只需在 $\mathbf{P}^n_k$ 上对
$$
\mathcal{G} =
i_*(\mathcal{F} \oplus \mathcal{F} \otimes \mathcal{L} \oplus \ldots
\oplus \mathcal{F} \otimes \mathcal{L}^{\otimes e - 1})
\quad\text{且}\quad \mathcal{O}(1)
$$
证明引理。事实上，由《概形的上同调》，引理
\ref{coherent-lemma-relative-affine-cohomology}，有
$$
H^1(\mathbf{P}^n_k, \mathcal{G}(m)) =
\bigoplus\nolimits_{j = 0, \ldots, e - 1}
H^1(X, \mathcal{F} \otimes \mathcal{L}^{\otimes j + me})
$$
。另外，若 $y \in \mathbf{P}^n_k$ 是闭点，则
$\text{depth}(\mathcal{G}_y) = \infty$，如果 $y \not \in i(X)$；
并且 $\text{depth}(\mathcal{G}_y) = \text{depth}(\mathcal{F}_x)$，
如果 $y = i(x)$。这是因为在后一种情形下，有
$\mathcal{G}_y \cong \mathcal{F}_x^{\oplus e}$，这是作为
$\mathcal{O}_{\mathbf{P}^n_k, x}$-模而言；并可用例如
《代数》，引理 \ref{algebra-lemma-depth-goes-down-finite}
得到该等式。

\medskip\noindent
假设 $X = \mathbf{P}^n_k$、$\mathcal{L} = \mathcal{O}(1)$，且
$k$ 是无限域。选取 $s \in H^0(\mathbf{P}^1_k, \mathcal{O}(1))$，
它确定正合列
$$
0 \to \mathcal{F}(-1) \xrightarrow{s} \mathcal{F} \to \mathcal{G} \to 0
$$
，如引理 \ref{varieties-lemma-exact-sequence-induction} 中所述。
由于映射 $\mathcal{F}(-1) \to \mathcal{F}$ 在仿射局部由
$\mathcal{F}$ 上的非零因子相乘给出，所以对闭点
$x \in \mathbf{P}^n_k$ 有 $\text{depth}(\mathcal{G}_x) \geq 1$，
见《代数》，引理 \ref{algebra-lemma-depth-drops-by-one}。
因此由引理 \ref{varieties-lemma-vanishin-h0-negative}，
$H^0(\mathcal{G}(m)) = 0$，只要 $m \ll 0$。
考察扭转后的上同调长正合列
（见注 \ref{varieties-remark-exact-sequence-induction-cohomology}），可知数列
$$
\dim H^1(\mathbf{P}^n_k, \mathcal{F}(m))
$$
在 $m \leq m_0$ 时稳定，其中 $m_0$ 是某个整数。
令 $N$ 为这些空间在 $m \leq m_0$ 时的公共维数。需要证明 $N = 0$。

\medskip\noindent
对 $d > 0$ 与 $m \leq m_0$，考虑双线性映射
$$
H^0(\mathbf{P}^n_k, \mathcal{O}(d)) \times
H^1(\mathbf{P}^n_k, \mathcal{F}(m - d))
\longrightarrow
H^1(\mathbf{P}^n_k, \mathcal{F}(m))
$$
。由线性代数，存在余维 $\leq N^2$ 的子空间
$V_m \subset H^0(\mathbf{P}^n_k, \mathcal{O}(d))$，使得由
$s' \in V_m$ 给出的乘法零化
$H^1(\mathbf{P}^n_k, \mathcal{F}(m - d))$。
注意，对 $m' < m \leq m_0$，图表
$$
\xymatrix{
H^0(\mathbf{P}^n_k, \mathcal{O}(d)) \times
H^1(\mathbf{P}^n_k, \mathcal{F}(m' - d)) \ar[r] \ar[d]^{1 \times s^{m' - m}} &
H^1(\mathbf{P}^n_k, \mathcal{F}(m')) \ar[d]^{s^{m' - m}}\\
H^0(\mathbf{P}^n_k, \mathcal{O}(d)) \times
H^1(\mathbf{P}^n_k, \mathcal{F}(m - d)) \ar[r] &
H^1(\mathbf{P}^n_k, \mathcal{F}(m))
}
$$
可交换，且竖直映射均为同构。因此 $V_m = V$ 与
$m \leq m_0$ 无关。对 $x \in \text{Ass}(\mathcal{F})$，
置 $Z = \overline{\{x\}}$。当 $d$ 足够大时，线性映射
$$
H^0(\mathbf{P}^n_k, \mathcal{O}(d)) \to H^0(Z, \mathcal{O}(d)|_Z)
$$
的秩 $> N^2$，因为 $\dim(Z) \geq 1$
（例如，这由渐近 Riemann--Roch 与 $\mathcal{O}(1)$ 的丰沛性得出；略去细节）。
因此可找到 $s' \in V$，使得 $s'$ 在
$\mathcal{F}$ 的任何伴随点都不消失（使用伴随点集有限这一事实）。
于是得到
$$
0 \to \mathcal{F}(-d) \xrightarrow{s'} \mathcal{F} \to \mathcal{G}' \to 0
$$
；与前面一样可知，乘以 $s'$ 在
$H^1(\mathbf{P}^n_k, \mathcal{F}(m - d))$ 上是单射，只要 $m \ll 0$。
这与对 $s'$ 的选取矛盾，除非 $N = 0$，故结论成立。

\medskip\noindent
还须处理 $k$ 是有限域的情形。
此时，令 $K/k$ 为任意无限代数域扩张。
记 $\mathcal{F}_K$ 与 $\mathcal{L}_K$ 为
$\mathcal{F}$ 与 $\mathcal{L}$ 到
$X_K = \Spec(K) \times_{\Spec(k)} X$ 上的拉回。
由《概形的上同调》，引理
\ref{coherent-lemma-flat-base-change-cohomology}，有
$$
H^1(X_K, \mathcal{F}_K \otimes \mathcal{L}_K^{\otimes m}) =
H^1(X, \mathcal{F} \otimes \mathcal{L}^{\otimes m}) \otimes_k K
$$
。另一方面，闭点 $x_K$（它属于 $X_K$）映到闭点
$x$（它属于 $X$），因为 $K/k$ 是代数扩张。
环映射 $\mathcal{O}_{X, x} \to \mathcal{O}_{X_K, x_K}$
是平坦的（引理 \ref{varieties-lemma-change-fields-flat}）。因此由《代数》，引理
\ref{algebra-lemma-apply-grothendieck-module}，有
$$
\text{depth}(\mathcal{F}_{x_K}) =
\text{depth}(\mathcal{F}_x \otimes_{\mathcal{O}_{X, x}}
\mathcal{O}_{X_K, x_K}) \geq
\text{depth}(\mathcal{F}_x)
$$
（事实上这里等号成立，但我们不需要它）。
所以 $k$ 上的结果由无限域 $K$ 上的结果得出，证明完成。
\end{proof}

\begin{lemma}
\label{varieties-lemma-connectedness-ample-divisor}
设 $k$ 是域。设 $X$ 是 $k$ 上的本征概形。设 $\mathcal{L}$ 是丰沛可逆
$\mathcal{O}_X$-模。设 $s \in \Gamma(X, \mathcal{L})$。假设
\begin{enumerate}
\item $s$ 是正则截面
（《除子》，定义 \ref{divisors-definition-regular-section}），
\item 对每个闭点 $x \in X$，有
$\text{depth}(\mathcal{O}_{X, x}) \geq 2$，并且
\item $X$ 是连通的。
\end{enumerate}
则 $Z(s)$（即 $s$ 的零概形）是连通的。
\end{lemma}

\begin{proof}
由于 $s$ 是正则截面，
$s^n \in \Gamma(X, \mathcal{L}^{\otimes n})$ 对所有 $n > 1$ 也是正则截面。
此外，包含态射 $Z(s) \to Z(s^n)$ 在底拓扑空间上是双射。
因此，若 $Z(s)$ 不连通，则 $Z(s^n)$ 也不连通。
现在考虑典范短正合列
$$
0 \to \mathcal{L}^{\otimes -n} \xrightarrow{s^n}
\mathcal{O}_X \to \mathcal{O}_{Z(s^n)} \to 0
$$
。考虑 $k$-代数 $R_n = \Gamma(X, \mathcal{O}_{Z(s^n)})$。
若 $Z(s)$ 不连通，即 $Z(s^n)$ 不连通，则或者
$R_n$ 为零（当 $Z(s^n) = \emptyset$ 时），或者
$R_n$ 含有非平凡幂等元（当
$Z(s^n) = U \amalg V$，其中 $U, V \subset Z(s^n)$ 开且非空时；
读者可查阅引理
\ref{varieties-lemma-proper-geometrically-reduced-global-sections}）。
所以映射 $\Gamma(X, \mathcal{O}_X) \to R_n$ 不可能是同构。
从而，要么 $H^0(X, \mathcal{L}^{\otimes -n})$ 非零，要么
$H^1(X, \mathcal{L}^{\otimes -n})$ 对无穷多个正整数 $n$ 非零。
这与引理 \ref{varieties-lemma-vanishin-h0-negative} 或
\ref{varieties-lemma-vanishin-h1-negative} 矛盾，证明完成。
\end{proof}
% Generated mechanically from the frozen topologies.tex target.
% Do not edit this generated file; edit the bound locale source instead.

% OK, start here.
%

\setcounter{chapter}{33}
\chapter{概形上的拓扑}


\phantomsection
\label{topologies-section-phantom}




\section{引言}
\label{topologies-section-introduction}

\noindent
本文说明概形范畴上的各种拓扑。可参阅 \cite{SGA1} 和 \cite{Ner}。
在此之前，我们想先指出：对于同一个底层范畴，可以选取许多不同的位点
（定义见 Sites, Definition \ref{sites-definition-site}），
而它们给出的层概念相同。因此，我们的选择可能与参考文献略有不同，
但最终得到的上同调群等对象是相同的。

\section{一般构造步骤}
\label{topologies-section-procedure}

\noindent
本节说明构造本文所用位点的一般步骤。设我们希望研究概形上相对于某个
拓扑 $\tau$ 的层。为了得到 Sites, Definition
\ref{sites-definition-site} 意义下配备该拓扑的概形位点，
还需要作一些处理。具体而言，我们不能简单地说
“考虑所有概形及其 Zariski 拓扑”，因为这样会得到一个“大”范畴。
本章各节将改按如下步骤进行：
\begin{enumerate}
\item 定义满足 Sites, Definition \ref{sites-definition-site}
中各条公理的概形覆盖类 $\text{Cov}_\tau$。概形的每个 Zariski 开覆盖
都将是一个 $\tau$-覆盖。
\item 在仿射概形范畴中指定“标准 $\tau$-覆盖”的概念。
\item 定义何谓“绝对”的大 $\tau$-位点 $\Sch_\tau$。
这类位点由适当选取的一组概形和一组覆盖构成。
\item 对 $\Sch_\tau$ 的任意对象 $S$，定义大 $\tau$-位点
$(\Sch/S)_\tau$；当 $\tau$ 合适时，还定义小\footnote{这里“大”和
“小”并非指相应范畴的大小。}$\tau$-位点 $S_\tau$。
\item 此外，利用仿射概形的标准 $\tau$-覆盖概念，还可得到位点
$(\textit{Aff}/S)_\tau$\footnote{对于 ph 拓扑，我们会略微偏离这一做法；
见 Definition \ref{topologies-definition-big-small-ph} 及其附近的讨论。}；
其层范畴与 $(\Sch/S)_\tau$ 上的层范畴等价。
\end{enumerate}
上述做法稍显繁琐，因为它并未给出概形的大 $\tau$-位点的典范选择，
甚至也未给出小 $\tau$-位点的典范选择。若愿意忽略集合论方面的困难，
则可以使用类来工作，从而得到典范的大位点和小位点……







\section{Zariski 拓扑}
\label{topologies-section-zariski}

\begin{definition}
\label{topologies-definition-zariski-covering}
设 $T$ 为概形。$T$ 的一个{\it Zariski 覆盖}是概形态射族
$\{f_i : T_i \to T\}_{i \in I}$，其中每个 $f_i$ 都是开浸入，
并且 $T = \bigcup f_i(T_i)$。
\end{definition}

\noindent
这定义了一个覆盖的（真）类。下面证明这一概念满足
Sites, Definition \ref{sites-definition-site} 的条件。

\begin{lemma}
\label{topologies-lemma-zariski}
设 $T$ 为概形。
\begin{enumerate}
\item 若 $T' \to T$ 是同构，则 $\{T' \to T\}$ 是 $T$ 的 Zariski 覆盖。
\item 若 $\{T_i \to T\}_{i\in I}$ 是 Zariski 覆盖，且对每个 $i$，
$\{T_{ij} \to T_i\}_{j\in J_i}$ 是 Zariski 覆盖，则
$\{T_{ij} \to T\}_{i \in I, j\in J_i}$ 是 Zariski 覆盖。
\item 若 $\{T_i \to T\}_{i\in I}$ 是 Zariski 覆盖，且
$T' \to T$ 是概形态射，则
$\{T' \times_T T_i \to T'\}_{i\in I}$ 是 Zariski 覆盖。
\end{enumerate}
\end{lemma}

\begin{proof}
略。
\end{proof}

\begin{lemma}
\label{topologies-lemma-zariski-affine}
设 $T$ 为仿射概形，且 $\{T_i \to T\}_{i \in I}$ 是 $T$ 的
Zariski 覆盖。则存在 Zariski 覆盖
$\{U_j \to T\}_{j = 1, \ldots, m}$，它是
$\{T_i \to T\}_{i \in I}$ 的一个加细，并且每个 $U_j$ 都是 $T$ 的
标准开集；见 Schemes, Definition
\ref{schemes-definition-standard-covering}。此外，可以选取每个 $U_j$，
使之成为某个 $T_i$ 的开子概形。
\end{lemma}

\begin{proof}
这是因为 $T$ 拟紧，而标准开集构成其拓扑的一组基。
Schemes, Lemma \ref{schemes-lemma-standard-open} 也证明了这一点。
\end{proof}

\noindent
于是，相应的仿射概形标准覆盖定义如下。

\begin{definition}
\label{topologies-definition-standard-Zariski}
比较 Schemes, Definition \ref{schemes-definition-standard-covering}。
设 $T$ 为仿射概形。$T$ 的一个{\it 标准 Zariski 覆盖}是 Zariski 覆盖
$\{U_j \to T\}_{j = 1, \ldots, m}$，其中每个 $U_j \to T$
都给出到 $T$ 的某个标准仿射开子概形的同构。
\end{definition}

\begin{definition}
\label{topologies-definition-big-zariski-site}
一个{\it 大 Zariski 位点}是 Sites, Definition
\ref{sites-definition-site} 意义下按如下方式构造的任一位点 $\Sch_{Zar}$：
\begin{enumerate}
\item 任取一组概形 $S_0$，以及这些概形之间的一组 Zariski 覆盖
$\text{Cov}_0$。
\item 取由集合 $S_0$ 出发、按 Sets, Lemma
\ref{sets-lemma-construct-category} 构造的任一范畴 $\Sch_\alpha$，
作为 $\Sch_{Zar}$ 的底层范畴。
\item 由范畴 $\Sch_\alpha$、Zariski 覆盖类以及上述集合
$\text{Cov}_0$ 出发，按 Sets, Lemma
\ref{sets-lemma-coverings-site} 选取任一组覆盖，作为 $\Sch_{Zar}$ 的覆盖。
\end{enumerate}
\end{definition}

\noindent
Sites, Lemma \ref{sites-lemma-choice-set-coverings-immaterial} 表明，
一旦选定范畴 $\Sch_\alpha$，$\Sch_\alpha$ 上的层范畴便不依赖于
上述第 (3) 项中所选的覆盖。换言之，拓扑斯 $\Sh(\Sch_{Zar})$
只依赖于范畴 $\Sch_\alpha$ 的选择。Sets, Lemma
\ref{sets-lemma-what-is-in-it} 表明，这些范畴对代数几何中的许多构造
封闭，例如纤维积以及取开、闭子概形。还可以证明，$\Sch_\alpha$
的具体选择并不十分重要；见 Section \ref{topologies-section-change-alpha}。

\medskip\noindent
另一种做法是假设存在强不可达基数，并将 $\Sch_{Zar}$ 定义为某个
给定宇宙中所含概形的范畴，而覆盖集则取同一宇宙中所含的 Zariski 覆盖。

\medskip\noindent
在继续引入概形 $S$ 的大 Zariski 位点之前，先指出：大 Zariski 位点
$\Sch_{Zar}$ 上的拓扑，在某种意义上由全体概形范畴上的 Zariski 拓扑诱导而来。

\begin{lemma}
\label{topologies-lemma-zariski-induced}
设 $\Sch_{Zar}$ 是 Definition \ref{topologies-definition-big-zariski-site}
所述的大 Zariski 位点，$T \in \Ob(\Sch_{Zar})$，并设
$\{T_i \to T\}_{i \in I}$ 是 $T$ 的任意 Zariski 覆盖。
则在位点 $\Sch_{Zar}$ 中存在 $T$ 的一个覆盖
$\{U_j \to T\}_{j \in J}$，它与 $\{T_i \to T\}_{i \in I}$
自明等价（见 Sites, Definition
\ref{sites-definition-combinatorial-tautological}）。
\end{lemma}

\begin{proof}
由于每个 $T_i \to T$ 都是开浸入，Sets, Lemma
\ref{sets-lemma-what-is-in-it} 表明，每个 $T_i$ 都同构于
$\Sch_{Zar}$ 的某个对象 $V_i$。覆盖 $\{V_i \to T\}_{i \in I}$
与 $\{T_i \to T\}_{i \in I}$ 自明等价（两个方向都使用 $I$ 上的恒等映射）。
此外，由 Sets, Lemma \ref{sets-lemma-coverings-site}，
$\{V_i \to T\}_{i \in I}$ 与位点 $\Sch_{Zar}$ 中 $T$ 的某个覆盖
$\{U_j \to T\}_{j \in J}$ 组合等价。
\end{proof}

\begin{definition}
\label{topologies-definition-big-small-Zariski}
设 $S$ 为概形，$\Sch_{Zar}$ 为包含 $S$ 的大 Zariski 位点。
\begin{enumerate}
\item $S$ 的{\it 大 Zariski 位点}记作 $(\Sch/S)_{Zar}$，
它是 Sites, Section \ref{sites-section-localize} 中引入的位点
$\Sch_{Zar}/S$。
\item $S$ 的{\it 小 Zariski 位点}记作 $S_{Zar}$，它是
$(\Sch/S)_{Zar}$ 的全子范畴，其对象为满足 $U \to S$ 是开浸入的
$U/S$。$S_{Zar}$ 的覆盖是 $(\Sch/S)_{Zar}$ 中任意满足
$U \in \Ob(S_{Zar})$ 的覆盖 $\{U_i \to U\}$。
\item $S$ 的{\it 大仿射 Zariski 位点}记作
$(\textit{Aff}/S)_{Zar}$，它是 $(\Sch/S)_{Zar}$ 的全子范畴，
由满足 $U$ 是仿射概形的对象 $U/S$ 构成。$(\textit{Aff}/S)_{Zar}$
的覆盖是 $(\Sch/S)_{Zar}$ 中任意满足
$U \in \Ob((\textit{Aff}/S)_{Zar})$ 且为标准 Zariski 覆盖的
$\{U_i \to U\}$。
\item $S$ 的{\it 小仿射 Zariski 位点}记作 $S_{affine, Zar}$，
它是 $S_{Zar}$ 的全子范畴，其对象为满足 $U$ 是仿射概形的 $U/S$。
$S_{affine, Zar}$ 的覆盖是 $S_{Zar}$ 中任意满足
$U \in \Ob(S_{affine, Zar})$ 且为标准 Zariski 覆盖的
$\{U_i \to U\}$。
\end{enumerate}
\end{definition}

\noindent
小 Zariski 位点、大仿射 Zariski 位点和小仿射 Zariski 位点确实是位点，
这一点并非完全显然。下面予以验证。

\begin{lemma}
\label{topologies-lemma-verify-site-Zariski}
设 $S$ 为概形，$\Sch_{Zar}$ 为包含 $S$ 的大 Zariski 位点。
上述定义的结构 $S_{Zar}$、$(\textit{Aff}/S)_{Zar}$ 和
$S_{affine, Zar}$ 都是位点。
\end{lemma}

\begin{proof}
先证明 $S_{Zar}$ 是位点。它是一个范畴，并配有一组具有固定目标的态射族。
因此必须验证 Sites, Definition \ref{sites-definition-site}
中的性质 (1)、(2)、(3)。由于 $(\Sch/S)_{Zar}$ 是位点，只需证明：
对 $(\Sch/S)_{Zar}$ 的任意覆盖 $\{U_i \to U\}$，若
$U \in \Ob(S_{Zar})$，则也有 $U_i \in \Ob(S_{Zar})$。
这由定义立即得到，因为开浸入的复合仍是开浸入。

\medskip\noindent
再证明 $(\textit{Aff}/S)_{Zar}$ 是位点。与上面相同，只需证明仿射概形的
标准 Zariski 覆盖族满足 Sites, Definition
\ref{sites-definition-site} 中的性质 (1)、(2)、(3)。设 $R$ 为环，
$f_1, \ldots, f_n \in R$ 生成单位理想。对每个
$i \in \{1, \ldots, n\}$，设 $g_{i1}, \ldots, g_{in_i} \in R_{f_i}$
生成 $R_{f_i}$ 的单位理想。可以写成
$g_{ij} = f_{ij}/f_i^{e_{ij}}$。必要时将 $f_{ij}$ 替换为
$f_i f_{ij}$，便有
$D(f_{ij}) \subset D(f_i) \cong \Spec(R_{f_i})$
等于 $D(g_{ij}) \subset \Spec(R_{f_i})$。因此，态射族
$\{D(g_{ij}) \to \Spec(R)\}$ 是标准 Zariski 覆盖。
由此可知，标准 Zariski 覆盖满足性质 (2)。性质 (1) 和 (3) 的验证从略。

\medskip\noindent
关于 $S_{affine, Zar}$ 是位点的证明从略。
\end{proof}

\begin{lemma}
\label{topologies-lemma-fibre-products-Zariski}
设 $S$ 为概形，$\Sch_{Zar}$ 为包含 $S$ 的大 Zariski 位点。
位点 $\Sch_{Zar}$、$(\Sch/S)_{Zar}$、$S_{Zar}$、
$(\textit{Aff}/S)_{Zar}$ 和 $S_{affine, Zar}$ 的底层范畴都具有纤维积。
在每种情形中，到全体概形范畴 $\Sch$ 的显然函子都与取纤维积相容。
范畴 $(\Sch/S)_{Zar}$ 和 $S_{Zar}$ 都有终对象，即 $S/S$。
\end{lemma}

\begin{proof}
对 $\Sch_{Zar}$ 而言，这是构造本身保证的；见 Sets, Lemma
\ref{sets-lemma-what-is-in-it}。设 $U \to S$、$V \to U$、$W \to U$
为概形态射，且 $U, V, W \in \Ob(\Sch_{Zar})$。
$\Sch_{Zar}$ 中的纤维积 $V \times_U W$ 也是 $\Sch$ 中的纤维积，
并且在 $S$ 上全体概形的范畴中，它是 $V/S$ 与 $W/S$ 在 $U/S$ 上的
纤维积，因而也是 $(\Sch/S)_{Zar}$ 中的纤维积。这证明了
$(\Sch/S)_{Zar}$ 的结论。若 $U \to S$、$V \to U$ 和 $W \to U$
都是开浸入，则 $V \times_U W \to S$ 也是开浸入，故得到
$S_{Zar}$ 的结论。若 $U, V, W$ 都是仿射概形，则 $V \times_U W$
也是仿射概形，从而得到 $(\textit{Aff}/S)_{Zar}$ 和
$S_{affine, Zar}$ 的结论。
\end{proof}

\noindent
下面验证大仿射位点（相应地，小仿射位点）与大位点（相应地，小位点）
定义同一个拓扑斯。

\begin{lemma}
\label{topologies-lemma-affine-big-site-Zariski}
设 $S$ 为概形，$\Sch_{Zar}$ 为包含 $S$ 的大 Zariski 位点。
函子 $(\textit{Aff}/S)_{Zar} \to (\Sch/S)_{Zar}$ 是特殊余连续函子，
因而诱导从 $\Sh((\textit{Aff}/S)_{Zar})$ 到
$\Sh((\Sch/S)_{Zar})$ 的拓扑斯等价。
\end{lemma}

\begin{proof}
特殊余连续函子的概念见 Sites, Definition
\ref{sites-definition-special-cocontinuous-functor}。因此需要验证
Sites, Lemma \ref{sites-lemma-equivalence} 的假设 (1)--(5)。
将包含函子记作
$u : (\textit{Aff}/S)_{Zar} \to (\Sch/S)_{Zar}$。
余连续性恰好意味着：当 $T$ 仿射时，$T/S$ 的每个 Zariski 覆盖
都可由 $T$ 的某个标准 Zariski 覆盖加细。这正是 Lemma
\ref{topologies-lemma-zariski-affine} 的内容，故 (1) 成立。由于标准 Zariski 覆盖
本身就是 Zariski 覆盖，$u$ 显然连续，故 (2) 成立。(3) 和 (4)
由 $u$ 全忠实立即得到。最后，每个概形都有仿射开覆盖，所以 (5) 成立。
\end{proof}

\begin{lemma}
\label{topologies-lemma-alternative-zariski}
设 $S$ 为概形，$\Sch_{Zar}$ 为包含 $S$ 的大 Zariski 位点。
函子 $S_{affine, Zar} \to S_{Zar}$ 是特殊余连续函子，因而诱导从
$\Sh(S_{affine, Zar})$ 到 $\Sh(S_{Zar})$ 的拓扑斯等价。
\end{lemma}

\begin{proof}
略。提示：与 Lemma \ref{topologies-lemma-affine-big-site-Zariski} 的证明比较。
\end{proof}

\noindent
下面验证小 Zariski 位点上的层概念与 $S$ 上的层概念相对应。

\begin{lemma}
\label{topologies-lemma-Zariski-usual}
$S_{Zar}$ 上的层范畴与 $S$ 的底层拓扑空间上的层范畴等价。
\end{lemma}

\begin{proof}
我们将反复使用如下事实：对 $S_{Zar}$ 的任意对象 $U/S$，态射
$U \to S$ 都给出到某个开子概形的同构。设 $\mathcal{F}$ 是
$S$ 上的层。按规则
$\mathcal{F}'(U/S) = \mathcal{F}(\Im(U \to S))$
定义 $S_{Zar}$ 上的层。反过来，对每个开子概形 $U \subset S$，
选取对象 $U'/S \in \Ob(S_{Zar})$，使得 $\Im(U' \to S) = U$
（这里要用 Sets, Lemma \ref{sets-lemma-what-is-in-it}）。
给定 $S_{Zar}$ 上的层 $\mathcal{G}$，令
$\mathcal{G}'(U) = \mathcal{G}(U'/S)$，从而定义 $S$ 上的层。
为了说明 $\mathcal{G}'$ 是层，对任意开覆盖
$U = \bigcup_{i \in I} U_i$，使用 Sets, Lemma
\ref{sets-lemma-coverings-site} 所给出的事实：覆盖
$\{U_i \to U\}_{i \in I}$ 与 $S_{Zar}$ 中的某个覆盖
$\{U_j' \to U'\}_{j \in J}$ 组合等价；再使用 Sites, Lemma
\ref{sites-lemma-tautological-same-sheaf}。细节从略。
\end{proof}

\noindent
从现在起，我们不再区分 $S_{Zar}$ 上的层与 $S$ 上的层；总是使用上述引理
证明中的步骤在两个概念之间转换。下面建立这些位点所对应拓扑斯之间的一些关系。

\begin{lemma}
\label{topologies-lemma-put-in-T}
设 $\Sch_{Zar}$ 为大 Zariski 位点，$f : T \to S$ 为
$\Sch_{Zar}$ 中的态射。函子 $T_{Zar} \to (\Sch/S)_{Zar}$
余连续，并诱导拓扑斯态射
$$
i_f :
\Sh(T_{Zar})
\longrightarrow
\Sh((\Sch/S)_{Zar})
$$
对 $(\Sch/S)_{Zar}$ 上的层 $\mathcal{G}$，有公式
$(i_f^{-1}\mathcal{G})(U/T) = \mathcal{G}(U/S)$。
函子 $i_f^{-1}$ 还有左伴随 $i_{f, !}$，后者与纤维积和等化子相容。
\end{lemma}

\begin{proof}
将该函子记作 $u : T_{Zar} \to (\Sch/S)_{Zar}$。换言之，给定与
$T_{Zar}$ 的某个对象对应的开浸入 $j : U \to T$，令
$u(U \to T) = (f \circ j : U \to S)$。由 Lemma
\ref{topologies-lemma-fibre-products-Zariski}，该函子与纤维积相容。
此外，$T_{Zar}$ 有等化子（因为具有相同源和目标的任意两个态射都相同），
且 $u$ 与等化子相容。它显然余连续。由于 $u$ 将覆盖变为覆盖并与纤维积相容，
它也连续。因此结论由 Sites, Lemmas \ref{sites-lemma-when-shriek}
和 \ref{sites-lemma-preserve-equalizers} 得到。
\end{proof}

\begin{lemma}
\label{topologies-lemma-at-the-bottom}
设 $S$ 为概形，$\Sch_{Zar}$ 为包含 $S$ 的大 Zariski 位点。
包含函子 $S_{Zar} \to (\Sch/S)_{Zar}$ 满足 Sites, Lemma
\ref{sites-lemma-bigger-site} 的假设，因而诱导位点态射
$$
\pi_S : (\Sch/S)_{Zar} \longrightarrow S_{Zar}
$$
以及拓扑斯态射
$$
i_S : \Sh(S_{Zar}) \longrightarrow \Sh((\Sch/S)_{Zar})
$$
并且 $\pi_S \circ i_S = \text{id}$。此外，$i_S = i_{\text{id}_S}$，
其中 $i_{\text{id}_S}$ 如 Lemma \ref{topologies-lemma-put-in-T} 所述。
特别地，函子 $i_S^{-1} = \pi_{S, *}$ 由规则
$i_S^{-1}(\mathcal{G})(U/S) = \mathcal{G}(U/S)$ 描述。
\end{lemma}

\begin{proof}
在此情形下，函子 $u : S_{Zar} \to (\Sch/S)_{Zar}$ 除了具有上述
Lemma \ref{topologies-lemma-put-in-T} 证明中所见的性质外，还是全忠实的，
并将终对象变为终对象。引理随即成立。
\end{proof}

\begin{definition}
\label{topologies-definition-restriction-small-zariski}
在 Lemma \ref{topologies-lemma-at-the-bottom} 的情形中，函子
$i_S^{-1} = \pi_{S, *}$ 常称为{\it 到小 Zariski 位点的限制}。
对大 Zariski 位点上的层 $\mathcal{F}$，将这一限制记作
$\mathcal{F}|_{S_{Zar}}$。
\end{definition}

\noindent
采用这一记号，对大位点上的层 $\mathcal{F}$ 和小位点上的层
$\mathcal{G}$，有
\begin{align*}
\Mor_{\Sh(S_{Zar})}(\mathcal{F}|_{S_{Zar}}, \mathcal{G})
& =
\Mor_{\Sh((\Sch/S)_{Zar})}(\mathcal{F},
i_{S, *}\mathcal{G}) \\
\Mor_{\Sh(S_{Zar})}(\mathcal{G}, \mathcal{F}|_{S_{Zar}})
& =
\Mor_{\Sh((\Sch/S)_{Zar})}(\pi_S^{-1}\mathcal{G},
\mathcal{F})
\end{align*}
此外，有 $(i_{S, *}\mathcal{G})|_{S_{Zar}} = \mathcal{G}$ 以及
$(\pi_S^{-1}\mathcal{G})|_{S_{Zar}} = \mathcal{G}$。

\begin{lemma}
\label{topologies-lemma-morphism-big}
设 $\Sch_{Zar}$ 为大 Zariski 位点，$f : T \to S$ 为
$\Sch_{Zar}$ 中的态射。函子
$$
u : (\Sch/T)_{Zar} \longrightarrow (\Sch/S)_{Zar},
\quad
V/T \longmapsto V/S
$$
余连续，并具有连续的右伴随
$$
v : (\Sch/S)_{Zar} \longrightarrow (\Sch/T)_{Zar},
\quad
(U \to S) \longmapsto (U \times_S T \to T).
$$
二者诱导同一个拓扑斯态射
$$
f_{big} :
\Sh((\Sch/T)_{Zar})
\longrightarrow
\Sh((\Sch/S)_{Zar})
$$
有 $f_{big}^{-1}(\mathcal{G})(U/T) = \mathcal{G}(U/S)$，以及
$f_{big, *}(\mathcal{F})(U/S) = \mathcal{F}(U \times_S T/T)$。
此外，$f_{big}^{-1}$ 有左伴随 $f_{big!}$，它与纤维积和等化子相容。
\end{lemma}

\begin{proof}
函子 $u$ 余连续、连续，并与纤维积和等化子相容（细节从略；比较
Lemma \ref{topologies-lemma-put-in-T} 的证明）。因此可以应用 Sites, Lemmas
\ref{sites-lemma-when-shriek} 和
\ref{sites-lemma-preserve-equalizers}，从而得到 $f_{big}^{-1}$
的公式以及 $f_{big!}$ 的存在性。此外，函子 $v$ 是右伴随，因为给定
$U/T$ 和 $V/S$，恰有
$\Mor_S(u(U), V) = \Mor_T(U, V \times_S T)$。于是可应用 Sites, Lemmas
\ref{sites-lemma-have-functor-other-way} 和
\ref{sites-lemma-have-functor-other-way-morphism}，得到 $f_{big, *}$ 的公式。
\end{proof}

\begin{lemma}
\label{topologies-lemma-morphism-big-small}
设 $\Sch_{Zar}$ 为大 Zariski 位点，$f : T \to S$ 为
$\Sch_{Zar}$ 中的态射。
\begin{enumerate}
\item 有 $i_f = f_{big} \circ i_T$，其中 $i_f$ 如 Lemma
\ref{topologies-lemma-put-in-T} 所述，$i_T$ 如 Lemma
\ref{topologies-lemma-at-the-bottom} 所述。
\item 函子 $S_{Zar} \to T_{Zar}$，
$(U \to S) \mapsto (U \times_S T \to T)$，是连续的，并诱导位点态射
$$
f_{small} :
T_{Zar}
\longrightarrow
S_{Zar}
$$
若通过 Lemma \ref{topologies-lemma-Zariski-usual} 将 $T_{Zar}$（相应地，
$S_{Zar}$）上的层与 $T$（相应地，$S$）上的层等同，则函子
$f_{small}^{-1}$ 和 $f_{small, *}$ 分别与通常的 $f^{-1}$ 和 $f_*$ 一致。
\item 有如下位点态射的交换图
$$
\xymatrix{
T_{Zar} \ar[d]_{f_{small}} &
(\Sch/T)_{Zar} \ar[d]^{f_{big}} \ar[l]^{\pi_T} \\
S_{Zar} &
(\Sch/S)_{Zar} \ar[l]_{\pi_S}
}
$$
因而作为拓扑斯态射，
$f_{small} \circ \pi_T = \pi_S \circ f_{big}$。
\item 有 $f_{small} = \pi_S \circ f_{big} \circ i_T = \pi_S \circ i_f$。
\end{enumerate}
\end{lemma}

\begin{proof}
等式 $i_f = f_{big} \circ i_T$ 由等式
$i_f^{-1} = i_T^{-1} \circ f_{big}^{-1}$ 得到；后者由上述函子描述显然成立。
这证明了 (1)。

\medskip\noindent
关于 (2)，见 Sites, Example \ref{sites-example-continuous-map}。

\medskip\noindent
(3) 成立，因为 $\pi_S$ 和 $\pi_T$ 由包含函子给出，而
$f_{small}$ 和 $f_{big}$ 由基变换函子 $U \mapsto U \times_S T$ 给出。

\medskip\noindent
(4) 由 (3) 两边预复合 $i_T$ 得到。
\end{proof}

\noindent
在该引理的情形中，采用 Definition
\ref{topologies-definition-restriction-small-zariski} 的术语，对 $T$ 的大
Zariski 位点上的层 $\mathcal{F}$，有
$$
(f_{big, *}\mathcal{F})|_{S_{Zar}} =
f_{small, *}(\mathcal{F}|_{T_{Zar}}),
$$
该等式由引理中的位点图交换立即得到，因为到 $T$（相应地，$S$）
的小 Zariski 位点的限制由 $\pi_{T, *}$（相应地，$\pi_{S, *}$）给出。
涉及拉回与限制的类似公式并不成立。

\begin{lemma}
\label{topologies-lemma-composition}
给定 $(\Sch/S)_{Zar}$ 中的概形 $X$、$Y$、$Z$ 以及态射
$f : X \to Y$、$g : Y \to Z$，有
$g_{big} \circ f_{big} = (g \circ f)_{big}$ and
$g_{small} \circ f_{small} = (g \circ f)_{small}$.
\end{lemma}

\begin{proof}
对大位点上的函子，这由 Lemma \ref{topologies-lemma-morphism-big} 中推前与拉回的
简单描述得到。对小位点上的函子，这是 Sheaves, Lemma
\ref{sheaves-lemma-pushforward-composition}，其中使用了 Lemma
\ref{topologies-lemma-Zariski-usual} 的等同。
\end{proof}

\begin{lemma}
\label{topologies-lemma-morphism-big-small-cartesian-diagram}
设 $\Sch_{Zar}$ 为大 Zariski 位点。考虑 $\Sch_{Zar}$ 中的笛卡儿图
$$
\xymatrix{
T' \ar[r]_{g'} \ar[d]_{f'} & T \ar[d]^f \\
S' \ar[r]^g & S
}
$$
则
$i_g^{-1} \circ f_{big, *} = f'_{small, *} \circ (i_{g'})^{-1}$
and $g_{big}^{-1} \circ f_{big, *} = f'_{big, *} \circ (g'_{big})^{-1}$.
\end{lemma}

\begin{proof}
由于该图是笛卡儿图，对 $U'/S'$ 有
$U' \times_{S'} T' = U' \times_S T$。因此，
$i_g^{-1} \circ f_{big, *}$ 和
$f'_{small, *} \circ (i_{g'})^{-1}$ 都将 $(\Sch/T)_{Zar}$ 上的层
$\mathcal{F}$ 送到 $S'_{Zar}$ 上的层
$U' \mapsto \mathcal{F}(U' \times_{S'} T')$
（使用 Lemmas \ref{topologies-lemma-put-in-T} 和
\ref{topologies-lemma-morphism-big-small}）。第二个等式可用同样方法证明，
也可由更一般的 Sites, Lemma \ref{sites-lemma-localize-morphism} 推出。
\end{proof}

\noindent
可以把 $S$ 的大 Zariski 位点上的一个层看成：在所有 $S$ 上概形上的
一族“通常”层。

\begin{lemma}
\label{topologies-lemma-characterize-sheaf-big}
设 $S$ 是大 Zariski 位点 $\Sch_{Zar}$ 中的概形。
大 Zariski 位点 $(\Sch/S)_{Zar}$ 上的层 $\mathcal{F}$
由如下数据给出：
\begin{enumerate}
\item 对每个 $T/S \in \Ob((\Sch/S)_{Zar})$，给定 $T$ 上的层
$\mathcal{F}_T$；
\item 对 $(\Sch/S)_{Zar}$ 中的每个 $f : T' \to T$，给定映射
$c_f : f^{-1}\mathcal{F}_T \to \mathcal{F}_{T'}$。
\end{enumerate}
这些数据满足如下条件：
\begin{enumerate}
\item[(a)] 对 $(\Sch/S)_{Zar}$ 中任意的 $f : T' \to T$ 和
$g : T'' \to T'$，复合 $c_g \circ g^{-1}c_f$ 等于 $c_{f \circ g}$；
\item[(b)] 若 $(\Sch/S)_{Zar}$ 中的 $f : T' \to T$ 是开浸入，
则 $c_f$ 是同构。
\end{enumerate}
\end{lemma}

\begin{proof}
该引理由 Sites, Remark
\ref{sites-remark-variant-glue-sheaves-absolute} 中讨论的一个纯层论命题得到。
这里也给出直接证明。

\medskip\noindent
给定 $\Sh((\Sch/S)_{Zar})$ 上的层 $\mathcal{F}$，令
$\mathcal{F}_T = i_p^{-1}\mathcal{F}$，其中 $p : T \to S$ 是结构态射。
注意，对任意开集 $U \subset T$ 以及 $(\Sch/T)_{Zar}$ 中像为 $U$
的开浸入 $U' \to T$，有
$\mathcal{F}_T(U) = \mathcal{F}(U'/S)$；见 Lemmas
\ref{topologies-lemma-Zariski-usual} 和 \ref{topologies-lemma-put-in-T}。因此，给定 $S$ 上的
$f : T' \to T$ 以及 $U, U' \to T$，得到典范映射
$\mathcal{F}_T(U) = \mathcal{F}(U'/S) \to \mathcal{F}(U'\times_T T'/S)
= \mathcal{F}_{T'}(f^{-1}(U))$。其中间的箭头是 $\mathcal{F}$ 关于
$S$ 上态射 $U' \times_T T' \to U'$ 的限制映射。这些映射与限制相容，
因而定义从 $\mathcal{F}_T$ 到 $\mathcal{F}_{T'}$ 的 $f$-映射 $c_f$；
见 Sheaves, Definition \ref{sheaves-definition-f-map} 及其附近的讨论。
显然，$c_{f \circ g}$ 是 $c_f$ 与 $c_g$ 的复合，因为
$\mathcal{F}$ 的限制映射的复合仍是限制映射。

\medskip\noindent
反过来，给定引理所述的系统 $(\mathcal{F}_T, c_f)$，只需令
$\mathcal{F}(T/S) = \mathcal{F}_T(T)$，便可在
$\Sh((\Sch/S)_{Zar})$ 上定义预层 $\mathcal{F}$。对于限制映射，
给定 $f : T' \to T$ 和 $s \in \mathcal{F}(T)$，令拉回
$f^*(s)$ 等于 $c_f(s)$（这里仍把 $c_f$ 看作一个 $f$-映射）。
关于 $c_f$ 的条件保证拉回满足所需的函子性。关于它是层的验证从略。
显然，这样定义的两个构造互为逆构造。
\end{proof}























\section{平展拓扑}
\label{topologies-section-etale}

\noindent
设 $S$ 为概形。我们希望在 $S$ 上概形的范畴上定义平展拓扑。
按照前述一般原则，先引入平展覆盖的概念。

\begin{definition}
\label{topologies-definition-etale-covering}
设 $T$ 为概形。$T$ 的一个{\it 平展覆盖}是概形态射族
$\{f_i : T_i \to T\}_{i \in I}$，其中每个 $f_i$ 都是平展的，
并且 $T = \bigcup f_i(T_i)$。
\end{definition}

\begin{lemma}
\label{topologies-lemma-zariski-etale}
任意 Zariski 覆盖都是平展覆盖。
\end{lemma}

\begin{proof}
这由定义以及开浸入是平展态射这一事实立即得到；见
Morphisms, Lemma \ref{morphisms-lemma-open-immersion-etale}。
\end{proof}

\noindent
下面证明这一概念满足 Sites, Definition \ref{sites-definition-site} 的条件。

\begin{lemma}
\label{topologies-lemma-etale}
设 $T$ 为概形。
\begin{enumerate}
\item 若 $T' \to T$ 是同构，则 $\{T' \to T\}$ 是 $T$ 的平展覆盖。
\item 若 $\{T_i \to T\}_{i\in I}$ 是平展覆盖，且对每个 $i$，
$\{T_{ij} \to T_i\}_{j\in J_i}$ 是平展覆盖，则
$\{T_{ij} \to T\}_{i \in I, j\in J_i}$ 是平展覆盖。
\item 若 $\{T_i \to T\}_{i\in I}$ 是平展覆盖，且 $T' \to T$
是概形态射，则 $\{T' \times_T T_i \to T'\}_{i\in I}$ 是平展覆盖。
\end{enumerate}
\end{lemma}

\begin{proof}
略。
\end{proof}

\begin{lemma}
\label{topologies-lemma-etale-affine}
设 $T$ 为仿射概形，$\{T_i \to T\}_{i \in I}$ 是 $T$ 的平展覆盖。
则存在平展覆盖 $\{U_j \to T\}_{j = 1, \ldots, m}$，它是
$\{T_i \to T\}_{i \in I}$ 的一个加细，并且每个 $U_j$ 都是仿射概形。
此外，可以选取每个 $U_j$，使之成为某个 $T_i$ 中的仿射开子概形。
\end{lemma}

\begin{proof}
略。
\end{proof}

\noindent
于是，相应的仿射概形标准覆盖定义如下。

\begin{definition}
\label{topologies-definition-standard-etale}
设 $T$ 为仿射概形。$T$ 的一个{\it 标准平展覆盖}是态射族
$\{f_j : U_j \to T\}_{j = 1, \ldots, m}$，其中每个 $U_j$ 都是
$T$ 上的平展仿射概形，并且 $T = \bigcup f_j(U_j)$。
\end{definition}

\noindent
在上述定义中，我们{\bf 不}假设态射 $f_j$ 是标准平展态射。
原因在于，若作此假设，标准平展覆盖便不能在 $\textit{Aff}/S$ 上定义位点；
例如可参见 Algebra, Lemma \ref{algebra-lemma-standard-etale} 的第 (4) 项。
另一方面，仿射概形之间的平展态射自动是标准光滑态射；见
Algebra, Lemma \ref{algebra-lemma-etale-standard-smooth}。
因此，标准平展覆盖既是标准光滑覆盖，也是标准 syntomic 覆盖。

\begin{definition}
\label{topologies-definition-big-etale-site}
一个{\it 大平展位点}是 Sites, Definition
\ref{sites-definition-site} 意义下按如下方式构造的任一位点 $\Sch_\etale$：
\begin{enumerate}
\item 任取一组概形 $S_0$，以及这些概形之间的一组平展覆盖
$\text{Cov}_0$。
\item 取由集合 $S_0$ 出发、按 Sets, Lemma
\ref{sets-lemma-construct-category} 构造的任一范畴 $\Sch_\alpha$
作为底层范畴。
\item 由范畴 $\Sch_\alpha$、平展覆盖类以及上述集合 $\text{Cov}_0$
出发，按 Sets, Lemma \ref{sets-lemma-coverings-site} 选取任一组覆盖。
\end{enumerate}
\end{definition}

\noindent
关于大位点定义的动机和解释，见 Definition
\ref{topologies-definition-big-zariski-site} 后面的评注。

\medskip\noindent
在继续引入概形 $S$ 的大平展位点之前，先指出：大平展位点
$\Sch_\etale$ 上的拓扑，在某种意义上由全体概形范畴上的平展拓扑诱导而来。

\begin{lemma}
\label{topologies-lemma-etale-induced}
设 $\Sch_\etale$ 是 Definition \ref{topologies-definition-big-etale-site}
所述的大平展位点，$T \in \Ob(\Sch_\etale)$，并设
$\{T_i \to T\}_{i \in I}$ 是 $T$ 的任意平展覆盖。
\begin{enumerate}
\item 位点 $\Sch_\etale$ 中存在 $T$ 的覆盖
$\{U_j \to T\}_{j \in J}$，它加细 $\{T_i \to T\}_{i \in I}$。
\item 若 $\{T_i \to T\}_{i \in I}$ 是标准平展覆盖，则它与
$\Sch_\etale$ 中的某个覆盖自明等价。
\item 若 $\{T_i \to T\}_{i \in I}$ 是 Zariski 覆盖，则它与
$\Sch_\etale$ 中的某个覆盖自明等价。
\end{enumerate}
\end{lemma}

\begin{proof}
对每个 $i$，选取仿射开覆盖
$T_i = \bigcup_{j \in J_i} T_{ij}$，使每个 $T_{ij}$ 都映入 $T$
的某个仿射开子概形。由 Lemma \ref{topologies-lemma-etale}，其加细
$\{T_{ij} \to T\}_{i \in I, j \in J_i}$ 仍是 $T$ 的平展覆盖。
因此可以假设每个 $T_i$ 都是仿射的，并映入 $T$ 的某个仿射开集 $W_i$。
应用 Sets, Lemma \ref{sets-lemma-what-is-in-it} 可知，$W_i$
同构于 $\Sch_\etale$ 的某个对象。于是再次应用 Sets, Lemma
\ref{sets-lemma-what-is-in-it}，作为 $W_i$
上有限型概形的 $T_i$ 同构于 $\Sch_\etale$ 的某个对象 $V_i$。
覆盖 $\{V_i \to T\}_{i \in I}$ 加细 $\{T_i \to T\}_{i \in I}$
（因为对应对象彼此同构）。此外，由 Sets, Lemma
\ref{sets-lemma-what-is-in-it}，$\{V_i \to T\}_{i \in I}$
与位点 $\Sch_\etale$ 中 $T$ 的某个覆盖
$\{U_j \to T\}_{j \in J}$ 组合等价。覆盖
$\{U_j \to T\}_{j \in J}$ 正是 (1) 所需的加细。
在 (2)、(3) 的情形中，由 Sets, Lemma
\ref{sets-lemma-what-is-in-it}，每个 $T_i$ 都同构于
$\Sch_\etale$ 的某个对象；再应用 Sets, Lemma
\ref{sets-lemma-coverings-site} 即得所需结论。
\end{proof}

\begin{definition}
\label{topologies-definition-big-small-etale}
设 $S$ 为概形，$\Sch_\etale$ 为包含 $S$ 的大平展位点。
\begin{enumerate}
\item $S$ 的{\it 大平展位点}记作 $(\Sch/S)_\etale$，它是
Sites, Section \ref{sites-section-localize} 中引入的位点 $\Sch_\etale/S$。
\item $S$ 的{\it 小平展位点}记作 $S_\etale$，它是
$(\Sch/S)_\etale$ 的全子范畴，其对象为满足 $U \to S$ 平展的 $U/S$。
$S_\etale$ 的覆盖是 $(\Sch/S)_\etale$ 中任意满足
$U \in \Ob(S_\etale)$ 的覆盖 $\{U_i \to U\}$。
\item $S$ 的{\it 大仿射平展位点}记作 $(\textit{Aff}/S)_\etale$，
它是 $(\Sch/S)_\etale$ 的全子范畴，其对象为满足 $U$ 是仿射概形的
$U/S$。$(\textit{Aff}/S)_\etale$ 的覆盖是 $(\Sch/S)_\etale$
中任意满足 $U \in \Ob((\textit{Aff}/S)_\etale)$ 且为标准平展覆盖的
$\{U_i \to U\}$。
\item $S$ 的{\it 小仿射平展位点}记作 $S_{affine, \etale}$，
它是 $S_\etale$ 的全子范畴，其对象为满足 $U$ 是仿射概形的 $U/S$。
$S_{affine, \etale}$ 的覆盖是 $S_\etale$ 中任意满足
$U \in \Ob(S_{affine, \etale})$ 且为标准平展覆盖的 $\{U_i \to U\}$。
\end{enumerate}
\end{definition}

\noindent
大仿射平展位点、小平展位点和小仿射平展位点确实是位点，
这一点并非完全显然。下面予以验证。

\begin{lemma}
\label{topologies-lemma-verify-site-etale}
设 $S$ 为概形，$\Sch_\etale$ 为包含 $S$ 的大平展位点。
结构 $S_\etale$、$(\textit{Aff}/S)_\etale$ 和
$S_{affine, \etale}$ 都是位点。
\end{lemma}

\begin{proof}
先证明 $S_\etale$ 是位点。它是一个范畴，并配有一组具有固定目标的态射族。
因此必须验证 Sites, Definition \ref{sites-definition-site}
中的性质 (1)、(2)、(3)。由于 $(\Sch/S)_\etale$ 是位点，只需证明：
对 $(\Sch/S)_\etale$ 的任意覆盖 $\{U_i \to U\}$，若
$U \in \Ob(S_\etale)$，则也有 $U_i \in \Ob(S_\etale)$。
这由定义立即得到，因为平展态射的复合仍是平展态射。

\medskip\noindent
再证明 $(\textit{Aff}/S)_\etale$ 是位点。与上面相同，只需证明
仿射概形的标准平展覆盖族满足 Sites, Definition
\ref{sites-definition-site} 中的性质 (1)、(2)、(3)。这是显然的。
例如，给定标准平展覆盖 $\{T_i \to T\}_{i\in I}$，且对每个 $i$，
$\{T_{ij} \to T_i\}_{j\in J_i}$ 是标准平展覆盖，则
$\{T_{ij} \to T\}_{i \in I, j\in J_i}$ 是标准平展覆盖，
因为 $\bigcup_{i\in I} J_i$ 有限且每个 $T_{ij}$ 都是仿射的。

\medskip\noindent
关于 $S_{affine, \etale}$ 是位点的证明从略。
\end{proof}

\begin{lemma}
\label{topologies-lemma-fibre-products-etale}
设 $S$ 为概形，$\Sch_\etale$ 为包含 $S$ 的大平展位点。
位点 $\Sch_\etale$、$(\Sch/S)_\etale$、$S_\etale$、
$(\textit{Aff}/S)_\etale$ 和 $S_{affine, \etale}$ 的底层范畴都具有纤维积。
在每种情形中，到全体概形范畴 $\Sch$ 的显然函子都与取纤维积相容。
范畴 $(\Sch/S)_\etale$ 和 $S_\etale$ 都有终对象，即 $S/S$。
\end{lemma}

\begin{proof}
对 $\Sch_\etale$ 而言，这是构造本身保证的；见 Sets, Lemma
\ref{sets-lemma-what-is-in-it}。设 $U \to S$、$V \to U$、$W \to U$
为概形态射，且 $U, V, W \in \Ob(\Sch_\etale)$。
$\Sch_\etale$ 中的纤维积 $V \times_U W$ 也是 $\Sch$ 中的纤维积，
并且在 $S$ 上全体概形的范畴中，它是 $V/S$ 与 $W/S$ 在 $U/S$ 上的
纤维积，因而也是 $(\Sch/S)_\etale$ 中的纤维积。这证明了
$(\Sch/S)_\etale$ 的结论。若 $U \to S$、$V \to U$ 和 $W \to U$
都平展，则 $V \times_U W \to S$ 也平展，故得到 $S_\etale$ 的结论。
若 $U, V, W$ 都仿射，则 $V \times_U W$ 也仿射，从而得到
$(\textit{Aff}/S)_\etale$ 和 $S_{affine, \etale}$ 的结论。
\end{proof}

\noindent
下面验证大仿射位点（相应地，小仿射位点）与大位点（相应地，小位点）
定义同一个拓扑斯。

\begin{lemma}
\label{topologies-lemma-affine-big-site-etale}
设 $S$ 为概形，$\Sch_\etale$ 为包含 $S$ 的大平展位点。
函子 $(\textit{Aff}/S)_\etale \to (\Sch/S)_\etale$ 特殊余连续，
并诱导从 $\Sh((\textit{Aff}/S)_\etale)$ 到
$\Sh((\Sch/S)_\etale)$ 的拓扑斯等价。
\end{lemma}

\begin{proof}
特殊余连续函子的概念见 Sites, Definition
\ref{sites-definition-special-cocontinuous-functor}。因此需要验证
Sites, Lemma \ref{sites-lemma-equivalence} 的假设 (1)--(5)。
将包含函子记作
$u : (\textit{Aff}/S)_\etale \to (\Sch/S)_\etale$。
余连续性恰好意味着：当 $T$ 仿射时，$T/S$ 的每个平展覆盖都可由
$T$ 的某个标准平展覆盖加细。这正是 Lemma
\ref{topologies-lemma-etale-affine} 的内容，故 (1) 成立。由于标准平展覆盖本身就是
平展覆盖，$u$ 显然连续，故 (2) 成立。(3) 和 (4) 由 $u$ 全忠实立即得到。
最后，每个概形都有仿射开覆盖，所以 (5) 成立。
\end{proof}

\begin{lemma}
\label{topologies-lemma-alternative}
设 $S$ 为概形，$\Sch_\etale$ 为包含 $S$ 的大平展位点。
函子 $S_{affine, \etale} \to S_\etale$ 特殊余连续，并诱导从
$\Sh(S_{affine, \etale})$ 到 $\Sh(S_\etale)$ 的拓扑斯等价。
\end{lemma}

\begin{proof}
略。提示：与 Lemma \ref{topologies-lemma-affine-big-site-etale} 的证明比较。
\end{proof}

\noindent
下面建立这些位点所对应拓扑斯之间的一些关系。

\begin{lemma}
\label{topologies-lemma-put-in-T-etale}
设 $\Sch_\etale$ 为大平展位点，$f : T \to S$ 为
$\Sch_\etale$ 中的态射。函子 $T_\etale \to (\Sch/S)_\etale$
余连续，并诱导拓扑斯态射
$$
i_f :
\Sh(T_\etale)
\longrightarrow
\Sh((\Sch/S)_\etale)
$$
对 $(\Sch/S)_\etale$ 上的层 $\mathcal{G}$，有公式
$(i_f^{-1}\mathcal{G})(U/T) = \mathcal{G}(U/S)$。
函子 $i_f^{-1}$ 还有左伴随 $i_{f, !}$，后者与纤维积和等化子相容。
\end{lemma}

\begin{proof}
将该函子记作 $u : T_\etale \to (\Sch/S)_\etale$。换言之，给定与
$T_\etale$ 的某个对象对应的平展态射 $j : U \to T$，令
$u(U \to T) = (f \circ j : U \to S)$。由 Lemma
\ref{topologies-lemma-fibre-products-etale}，该函子与纤维积相容。
设 $a, b : U \to V$ 是 $T_\etale$ 中的两个态射。此时，$a$ 与 $b$
在概形范畴中的等化子为
$$
V \times_{\Delta_{V/T}, V \times_T V, (a, b)} U
$$
它是 $T$ 上平展概形的纤维积，因而仍在 $T$ 上平展。因此
$T_\etale$ 有等化子，且 $u$ 与等化子相容。它显然余连续。
由于 $u$ 将覆盖变为覆盖并与纤维积相容，它也连续。因此引理由
Sites, Lemmas \ref{sites-lemma-when-shriek}
和 \ref{sites-lemma-preserve-equalizers} 得到。
\end{proof}

\begin{lemma}
\label{topologies-lemma-at-the-bottom-etale}
设 $S$ 为概形，$\Sch_\etale$ 为包含 $S$ 的大平展位点。
包含函子 $S_\etale \to (\Sch/S)_\etale$ 满足 Sites, Lemma
\ref{sites-lemma-bigger-site} 的假设，因而诱导位点态射
$$
\pi_S : (\Sch/S)_\etale \longrightarrow S_\etale
$$
以及拓扑斯态射
$$
i_S : \Sh(S_\etale) \longrightarrow \Sh((\Sch/S)_\etale)
$$
并且 $\pi_S \circ i_S = \text{id}$。此外，$i_S = i_{\text{id}_S}$，
其中 $i_{\text{id}_S}$ 如 Lemma \ref{topologies-lemma-put-in-T-etale} 所述。
特别地，函子 $i_S^{-1} = \pi_{S, *}$ 由规则
$i_S^{-1}(\mathcal{G})(U/S) = \mathcal{G}(U/S)$ 描述。
\end{lemma}

\begin{proof}
在此情形下，函子 $u : S_\etale \to (\Sch/S)_\etale$ 除了具有上述
Lemma \ref{topologies-lemma-put-in-T-etale} 证明中所见的性质外，还是全忠实的，
并将终对象变为终对象。结论由 Sites, Lemma
\ref{sites-lemma-bigger-site} 得到。
\end{proof}

\begin{definition}
\label{topologies-definition-restriction-small-etale}
在 Lemma \ref{topologies-lemma-at-the-bottom-etale} 的情形中，函子
$i_S^{-1} = \pi_{S, *}$ 常称为{\it 到小平展位点的限制}。
对大平展位点上的层 $\mathcal{F}$，将这一限制记作
$\mathcal{F}|_{S_\etale}$。
\end{definition}

\noindent
采用这一记号，对大位点上的层 $\mathcal{F}$ 和小位点上的层
$\mathcal{G}$，有
\begin{align*}
\Mor_{\Sh(S_\etale)}(
\mathcal{F}|_{S_\etale},
\mathcal{G})
& =
\Mor_{\Sh((\Sch/S)_\etale)}(
\mathcal{F},
i_{S, *}\mathcal{G}) \\
\Mor_{\Sh(S_\etale)}(
\mathcal{G},
\mathcal{F}|_{S_\etale})
& =
\Mor_{\Sh((\Sch/S)_\etale)}(
\pi_S^{-1}\mathcal{G},
\mathcal{F})
\end{align*}
此外，有 $(i_{S, *}\mathcal{G})|_{S_\etale} = \mathcal{G}$ 以及
$(\pi_S^{-1}\mathcal{G})|_{S_\etale} = \mathcal{G}$。

\begin{lemma}
\label{topologies-lemma-morphism-big-etale}
设 $\Sch_\etale$ 为大平展位点，$f : T \to S$ 为
$\Sch_\etale$ 中的态射。函子
$$
u :
(\Sch/T)_\etale
\longrightarrow
(\Sch/S)_\etale,
\quad
V/T \longmapsto V/S
$$
余连续，并具有连续的右伴随
$$
v :
(\Sch/S)_\etale
\longrightarrow
(\Sch/T)_\etale,
\quad
(U \to S) \longmapsto (U \times_S T \to T).
$$
二者诱导同一个拓扑斯态射
$$
f_{big} :
\Sh((\Sch/T)_\etale)
\longrightarrow
\Sh((\Sch/S)_\etale)
$$
有 $f_{big}^{-1}(\mathcal{G})(U/T) = \mathcal{G}(U/S)$，以及
$f_{big, *}(\mathcal{F})(U/S) = \mathcal{F}(U \times_S T/T)$。
此外，$f_{big}^{-1}$ 有左伴随 $f_{big!}$，它与纤维积和等化子相容。
\end{lemma}

\begin{proof}
函子 $u$ 余连续、连续，并与纤维积和等化子相容（细节从略；比较
Lemma \ref{topologies-lemma-put-in-T-etale} 的证明）。因此可以应用 Sites, Lemmas
\ref{sites-lemma-when-shriek} 和
\ref{sites-lemma-preserve-equalizers}，从而得到 $f_{big}^{-1}$
的公式以及 $f_{big!}$ 的存在性。此外，函子 $v$ 是右伴随，因为给定
$U/T$ 和 $V/S$，恰有
$\Mor_S(u(U), V) = \Mor_T(U, V \times_S T)$。于是可应用 Sites, Lemmas
\ref{sites-lemma-have-functor-other-way} 和
\ref{sites-lemma-have-functor-other-way-morphism}，得到 $f_{big, *}$ 的公式。
\end{proof}

\begin{lemma}
\label{topologies-lemma-morphism-big-small-etale}
设 $\Sch_\etale$ 为大平展位点，$f : T \to S$ 为
$\Sch_\etale$ 中的态射。
\begin{enumerate}
\item 有 $i_f = f_{big} \circ i_T$，其中 $i_f$ 如 Lemma
\ref{topologies-lemma-put-in-T-etale} 所述，$i_T$ 如 Lemma
\ref{topologies-lemma-at-the-bottom-etale} 所述。
\item 函子 $S_\etale \to T_\etale$，
$(U \to S) \mapsto (U \times_S T \to T)$，是连续的，并诱导位点态射
$$
f_{small} : T_\etale \longrightarrow S_\etale
$$
有 $f_{small, *}(\mathcal{F})(U/S) = \mathcal{F}(U \times_S T/T)$。
\item 有如下位点态射的交换图
$$
\xymatrix{
T_\etale \ar[d]_{f_{small}} &
(\Sch/T)_\etale \ar[d]^{f_{big}} \ar[l]^{\pi_T}\\
S_\etale &
(\Sch/S)_\etale \ar[l]_{\pi_S}
}
$$
因而作为拓扑斯态射，
$f_{small} \circ \pi_T = \pi_S \circ f_{big}$。
\item 有 $f_{small} = \pi_S \circ f_{big} \circ i_T = \pi_S \circ i_f$。
\end{enumerate}
\end{lemma}

\begin{proof}
等式 $i_f = f_{big} \circ i_T$ 由等式
$i_f^{-1} = i_T^{-1} \circ f_{big}^{-1}$ 得到；后者由上述函子描述显然成立。
这证明了 (1)。

\medskip\noindent
函子
$u :
S_\etale
\to
T_\etale$，$u(U \to S) = (U \times_S T \to T)$，将覆盖变为覆盖，
并与纤维积相容；见 Lemma \ref{topologies-lemma-etale} (3) 和
\ref{topologies-lemma-fibre-products-etale}。此外，$S_\etale$ 和 $T_\etale$
分别有终对象 $S/S$ 和 $T/T$，且 $u(S/S) = T/T$。因此由
Sites, Proposition \ref{sites-proposition-get-morphism}，函子 $u$
对应于位点态射 $T_\etale \to S_\etale$，后者又给出拓扑斯态射；
见 Sites, Lemma \ref{sites-lemma-morphism-sites-topoi}。
推前的描述由这些参考结果立即得到。

\medskip\noindent
(3) 成立，因为 $\pi_S$ 和 $\pi_T$ 由包含函子给出，而
$f_{small}$ 和 $f_{big}$ 由基变换函子 $U \mapsto U \times_S T$ 给出。

\medskip\noindent
(4) 由 (3) 两边预复合 $i_T$ 得到。
\end{proof}

\noindent
在该引理的情形中，采用 Definition
\ref{topologies-definition-restriction-small-etale} 的术语，对 $T$ 的大平展位点上的层
$\mathcal{F}$，有
$$
(f_{big, *}\mathcal{F})|_{S_\etale} =
f_{small, *}(\mathcal{F}|_{T_\etale}),
$$
该等式由引理中的位点图交换立即得到，因为到 $T$（相应地，$S$）
的小平展位点的限制由 $\pi_{T, *}$（相应地，$\pi_{S, *}$）给出。
涉及拉回与限制的类似公式并不成立。

\begin{lemma}
\label{topologies-lemma-composition-etale}
% P05-ERR-0001: frozen source repeats Y and omits Z; see part-local errata ledger.
给定 $\Sch_\etale$ 中的概形 $X$、$Y$、$Y$ 以及态射
$f : X \to Y$、$g : Y \to Z$，有
$g_{big} \circ f_{big} = (g \circ f)_{big}$ and
$g_{small} \circ f_{small} = (g \circ f)_{small}$.
\end{lemma}

\begin{proof}
对大位点上的函子，这由 Lemma \ref{topologies-lemma-morphism-big-etale}
中推前与拉回的简单描述得到。对小位点上的函子，这由 Lemma
\ref{topologies-lemma-morphism-big-small-etale} 中推前函子的描述得到。
\end{proof}

\begin{lemma}
\label{topologies-lemma-morphism-big-small-cartesian-diagram-etale}
设 $\Sch_\etale$ 为大平展位点。考虑 $\Sch_\etale$ 中的笛卡儿图
$$
\xymatrix{
T' \ar[r]_{g'} \ar[d]_{f'} & T \ar[d]^f \\
S' \ar[r]^g & S
}
$$
则
$i_g^{-1} \circ f_{big, *} = f'_{small, *} \circ (i_{g'})^{-1}$
and $g_{big}^{-1} \circ f_{big, *} = f'_{big, *} \circ (g'_{big})^{-1}$.
\end{lemma}

\begin{proof}
由于该图是笛卡儿图，对 $U'/S'$ 有
$U' \times_{S'} T' = U' \times_S T$。因此，
$i_g^{-1} \circ f_{big, *}$ 和
$f'_{small, *} \circ (i_{g'})^{-1}$ 都将 $(\Sch/T)_\etale$ 上的层
$\mathcal{F}$ 送到 $S'_\etale$ 上的层
$U' \mapsto \mathcal{F}(U' \times_{S'} T')$
（使用 Lemmas \ref{topologies-lemma-put-in-T-etale} 和
\ref{topologies-lemma-morphism-big-etale}）。第二个等式可用同样方法证明，
也可由更一般的 Sites, Lemma \ref{sites-lemma-localize-morphism} 推出。
\end{proof}

\noindent
可以把 $S$ 的大平展位点上的一个层看成：在所有 $S$ 上概形上的
一族“通常”层。

\begin{lemma}
\label{topologies-lemma-characterize-sheaf-big-etale}
设 $S$ 是大平展位点 $\Sch_\etale$ 中的概形。
大平展位点 $(\Sch/S)_\etale$ 上的层 $\mathcal{F}$ 由如下数据给出：
\begin{enumerate}
\item 对每个 $T/S \in \Ob((\Sch/S)_\etale)$，给定 $T_\etale$ 上的层
$\mathcal{F}_T$；
\item 对 $(\Sch/S)_\etale$ 中的每个 $f : T' \to T$，给定映射
$c_f : f_{small}^{-1}\mathcal{F}_T \to \mathcal{F}_{T'}$。
\end{enumerate}
这些数据满足如下条件：
\begin{enumerate}
\item[(a)] 对 $(\Sch/S)_\etale$ 中任意的 $f : T' \to T$ 和
$g : T'' \to T'$，复合 $c_g \circ g_{small}^{-1}c_f$ 等于
$c_{f \circ g}$；
\item[(b)] 若 $(\Sch/S)_\etale$ 中的 $f : T' \to T$ 平展，
则 $c_f$ 是同构。
\end{enumerate}
\end{lemma}

\begin{proof}
该引理由 Sites, Remark
\ref{sites-remark-variant-glue-sheaves-absolute} 中讨论的一个纯层论命题得到。
这里也给出直接证明。

\medskip\noindent
给定 $\Sh((\Sch/S)_\etale)$ 上的层 $\mathcal{F}$，令
$\mathcal{F}_T = i_p^{-1}\mathcal{F}$，其中 $p : T \to S$ 是结构态射。
注意，对 $T_\etale$ 中任意的 $U \to T$，有
$\mathcal{F}_T(U) = \mathcal{F}(U/S)$；见 Lemma
\ref{topologies-lemma-put-in-T-etale}。因此，给定 $S$ 上的 $f : T' \to T$
以及 $U \to T$，得到典范映射
$\mathcal{F}_T(U) = \mathcal{F}(U/S) \to \mathcal{F}(U \times_T T'/S)
= \mathcal{F}_{T'}(U \times_T T')$。其中间的箭头是 $\mathcal{F}$
关于 $S$ 上态射 $U \times_T T' \to U$ 的限制映射。这些映射与限制相容，
因而定义映射 $c'_f : \mathcal{F}_T \to f_{small, *}\mathcal{F}_{T'}$，
其中 $u : T_\etale \to T'_\etale$ 是与 $f$ 相关的基变换函子。
由 $f_{small, *}$ 与 $f_{small}^{-1}$ 的伴随关系（见 Sites, Section
\ref{sites-section-continuous-functors}），这等价于映射
$c_f : f_{small}^{-1}\mathcal{F}_T \to \mathcal{F}_{T'}$。
显然，$c'_{f \circ g}$ 是 $c'_f$ 与 $f_{small, *}c'_g$ 的复合，
因为 $\mathcal{F}$ 的限制映射的复合仍是限制映射；这给出了
$c_f$、$c_g$ 与 $c_{f \circ g}$ 之间所需的关系。

\medskip\noindent
反过来，给定引理所述的系统 $(\mathcal{F}_T, c_f)$，只需令
$\mathcal{F}(T/S) = \mathcal{F}_T(T)$，便可在
$\Sh((\Sch/S)_\etale)$ 上定义预层 $\mathcal{F}$。对于限制映射，
给定 $f : T' \to T$ 和 $s \in \mathcal{F}(T)$，令拉回 $f^*(s)$
等于 $c_f(s)$（这里仍把 $c_f$ 看作映射
$\mathcal{F}_T \to f_{small, *}\mathcal{F}_{T'}$）。关于 $c_f$ 的条件
保证拉回满足所需的函子性。关于它是层的验证从略。
显然，这样定义的两个构造互为逆构造。
\end{proof}























\section{光滑拓扑}
\label{topologies-section-smooth}

\noindent
本节定义光滑拓扑。这多少有些多余，因为以后会看到（见
More on Morphisms, Section \ref{more-morphisms-section-etale-over-smooth}），
这一拓扑与平展拓扑定义同一个拓扑斯。不过，这一定义仍然合理，
并且偶尔会用到。

\begin{definition}
\label{topologies-definition-smooth-covering}
设 $T$ 为概形。$T$ 的一个{\it 光滑覆盖}是概形态射族
$\{f_i : T_i \to T\}_{i \in I}$，其中每个 $f_i$ 都是光滑的，
并且 $T = \bigcup f_i(T_i)$。
\end{definition}

\begin{lemma}
\label{topologies-lemma-zariski-etale-smooth}
任意平展覆盖都是光滑覆盖；因此，任意 Zariski 覆盖更是光滑覆盖。
\end{lemma}

\begin{proof}
这由定义、平展态射是光滑态射这一事实（见 Morphisms, Definition
\ref{morphisms-definition-etale}）以及 Lemma \ref{topologies-lemma-zariski-etale}
立即得到。
\end{proof}

\noindent
下面证明这一概念满足 Sites, Definition \ref{sites-definition-site} 的条件。

\begin{lemma}
\label{topologies-lemma-smooth}
设 $T$ 为概形。
\begin{enumerate}
\item 若 $T' \to T$ 是同构，则 $\{T' \to T\}$ 是 $T$ 的光滑覆盖。
\item 若 $\{T_i \to T\}_{i\in I}$ 是光滑覆盖，且对每个 $i$，
$\{T_{ij} \to T_i\}_{j\in J_i}$ 是光滑覆盖，则
$\{T_{ij} \to T\}_{i \in I, j\in J_i}$ 是光滑覆盖。
\item 若 $\{T_i \to T\}_{i\in I}$ 是光滑覆盖，且 $T' \to T$
是概形态射，则 $\{T' \times_T T_i \to T'\}_{i\in I}$ 是光滑覆盖。
\end{enumerate}
\end{lemma}

\begin{proof}
略。
\end{proof}

\begin{lemma}
\label{topologies-lemma-smooth-affine}
设 $T$ 为仿射概形，$\{T_i \to T\}_{i \in I}$ 是 $T$ 的光滑覆盖。
则存在光滑覆盖 $\{U_j \to T\}_{j = 1, \ldots, m}$，它是
$\{T_i \to T\}_{i \in I}$ 的一个加细，每个 $U_j$ 都是仿射概形，
而且每个态射 $U_j \to T$ 都是标准光滑态射；见 Morphisms, Definition
\ref{morphisms-definition-smooth}。此外，可以选取每个 $U_j$，
使之成为某个 $T_i$ 中的仿射开子概形。
\end{lemma}

\begin{proof}
略；但可参见 Algebra, Lemma \ref{algebra-lemma-smooth-syntomic}。
\end{proof}

\noindent
于是，相应的仿射概形标准覆盖定义如下。

\begin{definition}
\label{topologies-definition-standard-smooth}
设 $T$ 为仿射概形。$T$ 的一个{\it 标准光滑覆盖}是态射族
$\{f_j : U_j \to T\}_{j = 1, \ldots, m}$，其中每个 $U_j$ 都是仿射的，
$U_j \to T$ 是标准光滑态射，并且 $T = \bigcup f_j(U_j)$。
\end{definition}

\begin{definition}
\label{topologies-definition-big-smooth-site}
一个{\it 大光滑位点}是 Sites, Definition
\ref{sites-definition-site} 意义下按如下方式构造的任一位点 $\Sch_{smooth}$：
\begin{enumerate}
\item 任取一组概形 $S_0$，以及这些概形之间的一组光滑覆盖
$\text{Cov}_0$。
\item 取由集合 $S_0$ 出发、按 Sets, Lemma
\ref{sets-lemma-construct-category} 构造的任一范畴 $\Sch_\alpha$
作为底层范畴。
\item 由范畴 $\Sch_\alpha$、光滑覆盖类以及上述集合 $\text{Cov}_0$
出发，按 Sets, Lemma \ref{sets-lemma-coverings-site} 选取任一组覆盖。
\end{enumerate}
\end{definition}

\noindent
关于大位点定义的动机和解释，见 Definition
\ref{topologies-definition-big-zariski-site} 后面的评注。

\medskip\noindent
在继续引入概形 $S$ 的大光滑位点之前，先指出：大光滑位点
$\Sch_{smooth}$ 上的拓扑，在某种意义上由全体概形范畴上的光滑拓扑诱导而来。

\begin{lemma}
\label{topologies-lemma-smooth-induced}
设 $\Sch_{smooth}$ 是 Definition \ref{topologies-definition-big-smooth-site}
所述的大光滑位点，\linebreak
$T \in \Ob(\Sch_{smooth})$，并设
$\{T_i \to T\}_{i \in I}$ 是 $T$ 的任意光滑覆盖。
\begin{enumerate}
\item 位点 $\Sch_{smooth}$ 中存在 $T$ 的覆盖
$\{U_j \to T\}_{j \in J}$，它加细 $\{T_i \to T\}_{i \in I}$。
\item 若 $\{T_i \to T\}_{i \in I}$ 是标准光滑覆盖，则它与
$\Sch_{smooth}$ 中的某个覆盖自明等价。
\item 若 $\{T_i \to T\}_{i \in I}$ 是 Zariski 覆盖，则它与
$\Sch_{smooth}$ 中的某个覆盖自明等价。
\end{enumerate}
\end{lemma}

\begin{proof}
对每个 $i$，选取仿射开覆盖
$T_i = \bigcup_{j \in J_i} T_{ij}$，使每个 $T_{ij}$ 都映入 $T$
的某个仿射开子概形。由 Lemma \ref{topologies-lemma-smooth}，其加细
$\{T_{ij} \to T\}_{i \in I, j \in J_i}$ 仍是 $T$ 的光滑覆盖。
因此可以假设每个 $T_i$ 都是仿射的，并映入 $T$ 的某个仿射开集 $W_i$。
应用 Sets, Lemma \ref{sets-lemma-what-is-in-it} 可知，$W_i$
同构于 $\Sch_{smooth}$ 的某个对象。于是再次应用 Sets, Lemma
\ref{sets-lemma-what-is-in-it}，作为 $W_i$
上有限型概形的 $T_i$ 同构于 $\Sch_{smooth}$ 的某个对象 $V_i$。
覆盖 $\{V_i \to T\}_{i \in I}$ 加细 $\{T_i \to T\}_{i \in I}$
（因为对应对象彼此同构）。此外，由 Sets, Lemma
\ref{sets-lemma-what-is-in-it}，$\{V_i \to T\}_{i \in I}$
与位点 $\Sch_{smooth}$ 中 $T$ 的某个覆盖
$\{U_j \to T\}_{j \in J}$ 组合等价。覆盖
$\{U_j \to T\}_{j \in J}$ 正是 (1) 所需的加细。
在 (2)、(3) 的情形中，由 Sets, Lemma
\ref{sets-lemma-what-is-in-it}，每个 $T_i$ 都同构于
$\Sch_{smooth}$ 的某个对象；再应用 Sets, Lemma
\ref{sets-lemma-coverings-site} 即得所需结论。
\end{proof}

\begin{definition}
\label{topologies-definition-big-small-smooth}
设 $S$ 为概形，$\Sch_{smooth}$ 为包含 $S$ 的大光滑位点。
\begin{enumerate}
\item $S$ 的{\it 大光滑位点}记作 $(\Sch/S)_{smooth}$，它是
Sites, Section \ref{sites-section-localize} 中引入的位点 $\Sch_{smooth}/S$。
\item $S$ 的{\it 大仿射光滑位点}记作 $(\textit{Aff}/S)_{smooth}$，
它是 $(\Sch/S)_{smooth}$ 的全子范畴，其对象是仿射的 $U/S$。
$(\textit{Aff}/S)_{smooth}$ 的覆盖是 $(\Sch/S)_{smooth}$ 中任意的
标准光滑覆盖 $\{U_i \to U\}$。
\end{enumerate}
\end{definition}

\noindent
下面验证大仿射位点与大位点定义同一个拓扑斯。

\begin{lemma}
\label{topologies-lemma-affine-big-site-smooth}
设 $S$ 为概形，$\Sch_{smooth}$ 为包含 $S$ 的大光滑位点。
函子\linebreak
$(\textit{Aff}/S)_{smooth} \to (\Sch/S)_{smooth}$ 特殊余连续，
并诱导从 $\Sh((\textit{Aff}/S)_{smooth})$ 到
$\Sh((\Sch/S)_{smooth})$ 的拓扑斯等价。
\end{lemma}

\begin{proof}
特殊余连续函子的概念见 Sites, Definition
\ref{sites-definition-special-cocontinuous-functor}。因此需要验证
Sites, Lemma \ref{sites-lemma-equivalence} 的假设 (1)--(5)。
将包含函子记作
$u : (\textit{Aff}/S)_{smooth} \to (\Sch/S)_{smooth}$。
余连续性恰好意味着：当 $T$ 仿射时，$T/S$ 的每个光滑覆盖都可由
$T$ 的某个标准光滑覆盖加细。这正是 Lemma
\ref{topologies-lemma-smooth-affine} 的内容，故 (1) 成立。由于标准光滑覆盖本身就是
光滑覆盖，$u$ 显然连续，故 (2) 成立。(3) 和 (4) 由 $u$ 全忠实立即得到。
最后，每个概形都有仿射开覆盖，所以 (5) 成立。
\end{proof}

\noindent
待续……

\begin{lemma}
\label{topologies-lemma-morphism-big-smooth}
设 $\Sch_{smooth}$ 为大光滑位点，$f : T \to S$ 为
$\Sch_{smooth}$ 中的态射。函子
$$
u : (\Sch/T)_{smooth} \longrightarrow (\Sch/S)_{smooth},
\quad
V/T \longmapsto V/S
$$
余连续，并具有连续的右伴随
$$
v : (\Sch/S)_{smooth} \longrightarrow (\Sch/T)_{smooth},
\quad
(U \to S) \longmapsto (U \times_S T \to T).
$$
二者诱导同一个拓扑斯态射
$$
f_{big} :
\Sh((\Sch/T)_{smooth})
\longrightarrow
\Sh((\Sch/S)_{smooth})
$$
有 $f_{big}^{-1}(\mathcal{G})(U/T) = \mathcal{G}(U/S)$，以及
$f_{big, *}(\mathcal{F})(U/S) = \mathcal{F}(U \times_S T/T)$。
此外，$f_{big}^{-1}$ 有左伴随 $f_{big!}$，它与纤维积和等化子相容。
\end{lemma}

\begin{proof}
函子 $u$ 余连续、连续，并与纤维积和等化子相容。因此可以应用
Sites, Lemmas \ref{sites-lemma-when-shriek} 和
\ref{sites-lemma-preserve-equalizers}，从而得到 $f_{big}^{-1}$
的公式以及 $f_{big!}$ 的存在性。此外，函子 $v$ 是右伴随，因为给定
$U/T$ 和 $V/S$，恰有
$\Mor_S(u(U), V) = \Mor_T(U, V \times_S T)$。于是可应用 Sites, Lemmas
\ref{sites-lemma-have-functor-other-way} 和
\ref{sites-lemma-have-functor-other-way-morphism}，得到 $f_{big, *}$ 的公式。
\end{proof}











\section{syntomic 拓扑}
\label{topologies-section-syntomic}

\noindent
本节定义 syntomic 拓扑。这一拓扑颇为有趣：它常常与 fppf 拓扑
具有相同的上同调群，但在技术上更容易处理。

\begin{definition}
\label{topologies-definition-syntomic-covering}
设 $T$ 为概形。$T$ 的一个{\it syntomic 覆盖}是概形态射族
$\{f_i : T_i \to T\}_{i \in I}$，其中每个 $f_i$ 都是 syntomic 态射，
并且 $T = \bigcup f_i(T_i)$。
\end{definition}

\begin{lemma}
\label{topologies-lemma-zariski-etale-smooth-syntomic}
任意光滑覆盖都是 syntomic 覆盖；因此，任意平展覆盖或 Zariski 覆盖
更是 syntomic 覆盖。
\end{lemma}

\begin{proof}
这由定义、光滑态射是 syntomic 态射这一事实（见 Morphisms, Lemma
\ref{morphisms-lemma-smooth-syntomic}）以及 Lemma
\ref{topologies-lemma-zariski-etale-smooth} 立即得到。
\end{proof}

\noindent
下面证明这一概念满足 Sites, Definition \ref{sites-definition-site} 的条件。

\begin{lemma}
\label{topologies-lemma-syntomic}
设 $T$ 为概形。
\begin{enumerate}
\item 若 $T' \to T$ 是同构，则 $\{T' \to T\}$ 是 $T$ 的 syntomic 覆盖。
\item 若 $\{T_i \to T\}_{i\in I}$ 是 syntomic 覆盖，且对每个 $i$，
$\{T_{ij} \to T_i\}_{j\in J_i}$ 是 syntomic 覆盖，则
$\{T_{ij} \to T\}_{i \in I, j\in J_i}$ 是 syntomic 覆盖。
\item 若 $\{T_i \to T\}_{i\in I}$ 是 syntomic 覆盖，且 $T' \to T$
是概形态射，则 $\{T' \times_T T_i \to T'\}_{i\in I}$ 是 syntomic 覆盖。
\end{enumerate}
\end{lemma}

\begin{proof}
略。
\end{proof}

\begin{lemma}
\label{topologies-lemma-syntomic-affine}
设 $T$ 为仿射概形，$\{T_i \to T\}_{i \in I}$ 是 $T$ 的 syntomic 覆盖。
则存在 syntomic 覆盖 $\{U_j \to T\}_{j = 1, \ldots, m}$，它是
$\{T_i \to T\}_{i \in I}$ 的一个加细，每个 $U_j$ 都是仿射概形，
而且每个态射 $U_j \to T$ 都是标准 syntomic 态射；见 Morphisms,
Definition \ref{morphisms-definition-syntomic}。此外，可以选取每个 $U_j$，
使之成为某个 $T_i$ 中的仿射开子概形。
\end{lemma}

\begin{proof}
略；但可参见 Algebra, Lemma \ref{algebra-lemma-syntomic}。
\end{proof}

\noindent
于是，相应的仿射概形标准覆盖定义如下。

\begin{definition}
\label{topologies-definition-standard-syntomic}
设 $T$ 为仿射概形。$T$ 的一个{\it 标准 syntomic 覆盖}是态射族
$\{f_j : U_j \to T\}_{j = 1, \ldots, m}$，其中每个 $U_j$ 都是仿射的，
$U_j \to T$ 是标准 syntomic 态射，并且 $T = \bigcup f_j(U_j)$。
\end{definition}

\begin{definition}
\label{topologies-definition-big-syntomic-site}
一个{\it 大 syntomic 位点}是 Sites, Definition
\ref{sites-definition-site} 意义下按如下方式构造的任一位点 $\Sch_{syntomic}$：
\begin{enumerate}
\item 任取一组概形 $S_0$，以及这些概形之间的一组 syntomic 覆盖
$\text{Cov}_0$。
\item 取由集合 $S_0$ 出发、按 Sets, Lemma
\ref{sets-lemma-construct-category} 构造的任一范畴 $\Sch_\alpha$
作为底层范畴。
\item 由范畴 $\Sch_\alpha$、syntomic 覆盖类以及上述集合 $\text{Cov}_0$
出发，按 Sets, Lemma \ref{sets-lemma-coverings-site} 选取任一组覆盖。
\end{enumerate}
\end{definition}

\noindent
关于大位点定义的动机和解释，见 Definition
\ref{topologies-definition-big-zariski-site} 后面的评注。

\medskip\noindent
在继续引入概形 $S$ 的大 syntomic 位点之前，先指出：大 syntomic 位点
$\Sch_{syntomic}$ 上的拓扑，在某种意义上由全体概形范畴上的
syntomic 拓扑诱导而来。

\begin{lemma}
\label{topologies-lemma-syntomic-induced}
设 $\Sch_{syntomic}$ 是 Definition \ref{topologies-definition-big-syntomic-site}
所述的大 syntomic 位点，$T \in \Ob(\Sch_{syntomic})$，并设
$\{T_i \to T\}_{i \in I}$ 是 $T$ 的任意 syntomic 覆盖。
\begin{enumerate}
\item 位点 $\Sch_{syntomic}$ 中存在 $T$ 的覆盖
$\{U_j \to T\}_{j \in J}$，它加细 $\{T_i \to T\}_{i \in I}$。
\item 若 $\{T_i \to T\}_{i \in I}$ 是标准 syntomic 覆盖，则它与
$\Sch_{syntomic}$ 中的某个覆盖自明等价。
\item 若 $\{T_i \to T\}_{i \in I}$ 是 Zariski 覆盖，则它与
$\Sch_{syntomic}$ 中的某个覆盖自明等价。
\end{enumerate}
\end{lemma}

\begin{proof}
对每个 $i$，选取仿射开覆盖
$T_i = \bigcup_{j \in J_i} T_{ij}$，使每个 $T_{ij}$ 都映入 $T$
的某个仿射开子概形。由 Lemma \ref{topologies-lemma-syntomic}，其加细
$\{T_{ij} \to T\}_{i \in I, j \in J_i}$ 仍是 $T$ 的 syntomic 覆盖。
因此可以假设每个 $T_i$ 都是仿射的，并映入 $T$ 的某个仿射开集 $W_i$。
应用 Sets, Lemma \ref{sets-lemma-what-is-in-it} 可知，$W_i$
同构于 $\Sch_{syntomic}$ 的某个对象。于是再次应用 Sets, Lemma
\ref{sets-lemma-what-is-in-it}，作为 $W_i$
上有限型概形的 $T_i$ 同构于 $\Sch_{syntomic}$ 的某个对象 $V_i$。
覆盖 $\{V_i \to T\}_{i \in I}$ 加细 $\{T_i \to T\}_{i \in I}$
（因为对应对象彼此同构）。此外，由 Sets, Lemma
\ref{sets-lemma-what-is-in-it}，$\{V_i \to T\}_{i \in I}$
与位点 $\Sch_{syntomic}$ 中 $T$ 的某个覆盖
$\{U_j \to T\}_{j \in J}$ 组合等价。覆盖
$\{U_j \to T\}_{j \in J}$ 正是 (1) 所需的覆盖。
在 (2)、(3) 的情形中，由 Sets, Lemma
\ref{sets-lemma-what-is-in-it}，每个 $T_i$ 都同构于
$\Sch_{syntomic}$ 的某个对象；再应用 Sets, Lemma
\ref{sets-lemma-coverings-site} 即得所需结论。
\end{proof}

\begin{definition}
\label{topologies-definition-big-small-syntomic}
设 $S$ 为概形，$\Sch_{syntomic}$ 为包含 $S$ 的大 syntomic 位点。
\begin{enumerate}
\item $S$ 的{\it 大 syntomic 位点}记作 $(\Sch/S)_{syntomic}$，
它是 Sites, Section \ref{sites-section-localize} 中引入的位点
$\Sch_{syntomic}/S$。
\item $S$ 的{\it 大仿射 syntomic 位点}记作
$(\textit{Aff}/S)_{syntomic}$，它是 $(\Sch/S)_{syntomic}$ 的全子范畴，
其对象是仿射的 $U/S$。$(\textit{Aff}/S)_{syntomic}$ 的覆盖是
$(\Sch/S)_{syntomic}$ 中任意的标准 syntomic 覆盖 $\{U_i \to U\}$。
\end{enumerate}
\end{definition}

\noindent
下面验证大仿射位点与大位点定义同一个拓扑斯。

\begin{lemma}
\label{topologies-lemma-affine-big-site-syntomic}
设 $S$ 为概形，$\Sch_{syntomic}$ 为包含 $S$ 的大 syntomic 位点。
函子 $(\textit{Aff}/S)_{syntomic} \to (\Sch/S)_{syntomic}$ 特殊余连续，
并诱导从 $\Sh((\textit{Aff}/S)_{syntomic})$ 到
$\Sh((\Sch/S)_{syntomic})$ 的拓扑斯等价。
\end{lemma}

\begin{proof}
特殊余连续函子的概念见 Sites, Definition
\ref{sites-definition-special-cocontinuous-functor}。因此需要验证
Sites, Lemma \ref{sites-lemma-equivalence} 的假设 (1)--(5)。
将包含函子记作
$u : (\textit{Aff}/S)_{syntomic} \to (\Sch/S)_{syntomic}$。
余连续性恰好意味着：当 $T$ 仿射时，$T/S$ 的每个 syntomic 覆盖
都可由 $T$ 的某个标准 syntomic 覆盖加细。这正是 Lemma
\ref{topologies-lemma-syntomic-affine} 的内容，故 (1) 成立。由于标准 syntomic 覆盖
本身就是 syntomic 覆盖，$u$ 显然连续，故 (2) 成立。(3) 和 (4)
由 $u$ 全忠实立即得到。最后，每个概形都有仿射开覆盖，所以 (5) 成立。
\end{proof}

\noindent
待续……

\begin{lemma}
\label{topologies-lemma-morphism-big-syntomic}
设 $\Sch_{syntomic}$ 为大 syntomic 位点，$f : T \to S$ 为
$\Sch_{syntomic}$ 中的态射。函子
$$
u : (\Sch/T)_{syntomic} \longrightarrow (\Sch/S)_{syntomic},
\quad
V/T \longmapsto V/S
$$
余连续，并具有连续的右伴随
$$
v : (\Sch/S)_{syntomic} \longrightarrow (\Sch/T)_{syntomic},
\quad
(U \to S) \longmapsto (U \times_S T \to T).
$$
二者诱导同一个拓扑斯态射
$$
f_{big} :
\Sh((\Sch/T)_{syntomic})
\longrightarrow
\Sh((\Sch/S)_{syntomic})
$$
有 $f_{big}^{-1}(\mathcal{G})(U/T) = \mathcal{G}(U/S)$，以及
$f_{big, *}(\mathcal{F})(U/S) = \mathcal{F}(U \times_S T/T)$。
此外，$f_{big}^{-1}$ 有左伴随 $f_{big!}$，它与纤维积和等化子相容。
\end{lemma}

\begin{proof}
函子 $u$ 余连续、连续，并与纤维积和等化子相容。因此可以应用
Sites, Lemmas \ref{sites-lemma-when-shriek} 和
\ref{sites-lemma-preserve-equalizers}，从而得到 $f_{big}^{-1}$
的公式以及 $f_{big!}$ 的存在性。此外，函子 $v$ 是右伴随，因为给定
$U/T$ 和 $V/S$，恰有
$\Mor_S(u(U), V) = \Mor_T(U, V \times_S T)$。于是可应用 Sites, Lemmas
\ref{sites-lemma-have-functor-other-way} 和
\ref{sites-lemma-have-functor-other-way-morphism}，得到 $f_{big, *}$ 的公式。
\end{proof}













\section{fppf 拓扑}
\label{topologies-section-fppf}

\noindent
设 $S$ 为概形。我们希望在 $S$ 上概形的范畴上定义 fppf 拓扑\footnote{
字母 fppf 代表法语 ``fid\`element plat de pr\'esentation finie''，
即“忠实平坦且有限表示”。}。按照前述一般原则，先引入 fppf 覆盖的概念。

\begin{definition}
\label{topologies-definition-fppf-covering}
设 $T$ 为概形。$T$ 的一个{\it fppf 覆盖}是概形态射族
$\{f_i : T_i \to T\}_{i \in I}$，其中每个 $f_i$ 都平坦且局部有限表示，
并且 $T = \bigcup f_i(T_i)$。
\end{definition}

\begin{lemma}
\label{topologies-lemma-zariski-etale-smooth-syntomic-fppf}
任意 syntomic 覆盖都是 fppf 覆盖；因此，任意光滑覆盖、平展覆盖或
Zariski 覆盖更是 fppf 覆盖。
\end{lemma}

\begin{proof}
这由定义、syntomic 态射平坦且局部有限表示这一事实（见 Morphisms,
Lemmas \ref{morphisms-lemma-syntomic-locally-finite-presentation} 和
\ref{morphisms-lemma-syntomic-flat}）以及 Lemma
\ref{topologies-lemma-zariski-etale-smooth-syntomic} 立即得到。
\end{proof}

\noindent
下面证明这一概念满足 Sites, Definition \ref{sites-definition-site} 的条件。

\begin{lemma}
\label{topologies-lemma-fppf}
设 $T$ 为概形。
\begin{enumerate}
\item 若 $T' \to T$ 是同构，则 $\{T' \to T\}$ 是 $T$ 的 fppf 覆盖。
\item 若 $\{T_i \to T\}_{i\in I}$ 是 fppf 覆盖，且对每个 $i$，
$\{T_{ij} \to T_i\}_{j\in J_i}$ 是 fppf 覆盖，则
$\{T_{ij} \to T\}_{i \in I, j\in J_i}$ 是 fppf 覆盖。
\item 若 $\{T_i \to T\}_{i\in I}$ 是 fppf 覆盖，且 $T' \to T$
是概形态射，则 $\{T' \times_T T_i \to T'\}_{i\in I}$ 是 fppf 覆盖。
\end{enumerate}
\end{lemma}

\begin{proof}
第一项显然。第二项成立，因为平坦态射的复合仍平坦（见 Morphisms,
Lemma \ref{morphisms-lemma-composition-flat}），有限表示态射的复合仍有限表示
（见 Morphisms, Lemma \ref{morphisms-lemma-composition-finite-presentation}）。
第三项成立，因为平坦态射的基变换仍平坦（见 Morphisms, Lemma
\ref{morphisms-lemma-base-change-flat}），有限表示态射的基变换仍有限表示
（见 Morphisms, Lemma \ref{morphisms-lemma-base-change-finite-presentation}）。
此外，满射态射族的基变换仍是满射的（证明从略）。
\end{proof}

\begin{lemma}
\label{topologies-lemma-fppf-affine}
设 $T$ 为仿射概形，$\{T_i \to T\}_{i \in I}$ 是 $T$ 的 fppf 覆盖。
则存在 fppf 覆盖 $\{U_j \to T\}_{j = 1, \ldots, m}$，它是
$\{T_i \to T\}_{i \in I}$ 的一个加细，并且每个 $U_j$ 都是仿射概形。
此外，可以选取每个 $U_j$，使之成为某个 $T_i$ 中的仿射开子概形。
\end{lemma}

\begin{proof}
这直接由定义以及“平坦且局部有限表示的态射是开映射”这一事实得到；
见 Morphisms, Lemma \ref{morphisms-lemma-fppf-open}。
\end{proof}

\noindent
于是，相应的仿射概形标准覆盖定义如下。

\begin{definition}
\label{topologies-definition-standard-fppf}
设 $T$ 为仿射概形。$T$ 的一个{\it 标准 fppf 覆盖}是态射族
$\{f_j : U_j \to T\}_{j = 1, \ldots, m}$，其中每个 $U_j$ 都是
$T$ 上平坦、有限表示的仿射概形，并且 $T = \bigcup f_j(U_j)$。
\end{definition}

\begin{definition}
\label{topologies-definition-big-fppf-site}
一个{\it 大 fppf 位点}是 Sites, Definition
\ref{sites-definition-site} 意义下按如下方式构造的任一位点 $\Sch_{fppf}$：
\begin{enumerate}
\item 任取一组概形 $S_0$，以及这些概形之间的一组 fppf 覆盖
$\text{Cov}_0$。
\item 取由集合 $S_0$ 出发、按 Sets, Lemma
\ref{sets-lemma-construct-category} 构造的任一范畴 $\Sch_\alpha$
作为底层范畴。
\item 由范畴 $\Sch_\alpha$、fppf 覆盖类以及上述集合 $\text{Cov}_0$
出发，按 Sets, Lemma \ref{sets-lemma-coverings-site} 选取任一组覆盖。
\end{enumerate}
\end{definition}

\noindent
关于大位点定义的动机和解释，见 Definition
\ref{topologies-definition-big-zariski-site} 后面的评注。

\medskip\noindent
在继续引入概形 $S$ 的大 fppf 位点之前，先指出：大 fppf 位点
$\Sch_{fppf}$ 上的拓扑，在某种意义上由全体概形范畴上的 fppf 拓扑诱导而来。

\begin{lemma}
\label{topologies-lemma-fppf-induced}
设 $\Sch_{fppf}$ 是 Definition \ref{topologies-definition-big-fppf-site}
所述的大 fppf 位点，$T \in \Ob(\Sch_{fppf})$，并设
$\{T_i \to T\}_{i \in I}$ 是 $T$ 的任意 fppf 覆盖。
\begin{enumerate}
\item 位点 $\Sch_{fppf}$ 中存在 $T$ 的覆盖
$\{U_j \to T\}_{j \in J}$，它加细 $\{T_i \to T\}_{i \in I}$。
\item 若 $\{T_i \to T\}_{i \in I}$ 是标准 fppf 覆盖，则它与
$\Sch_{fppf}$ 中的某个覆盖自明等价。
\item 若 $\{T_i \to T\}_{i \in I}$ 是 Zariski 覆盖，则它与
$\Sch_{fppf}$ 中的某个覆盖自明等价。
\end{enumerate}
\end{lemma}

\begin{proof}
对每个 $i$，选取仿射开覆盖
$T_i = \bigcup_{j \in J_i} T_{ij}$，使每个 $T_{ij}$ 都映入 $T$
的某个仿射开子概形。由 Lemma \ref{topologies-lemma-fppf}，其加细
$\{T_{ij} \to T\}_{i \in I, j \in J_i}$ 仍是 $T$ 的 fppf 覆盖。
因此可以假设每个 $T_i$ 都是仿射的，并映入 $T$ 的某个仿射开集 $W_i$。
应用 Sets, Lemma \ref{sets-lemma-what-is-in-it} 可知，$W_i$
同构于 $\Sch_{fppf}$ 的某个对象。于是再次应用 Sets, Lemma
\ref{sets-lemma-what-is-in-it}，作为 $W_i$
上有限型概形的 $T_i$ 同构于 $\Sch_{fppf}$ 的某个对象 $V_i$。
覆盖 $\{V_i \to T\}_{i \in I}$ 加细 $\{T_i \to T\}_{i \in I}$
（因为对应对象彼此同构）。此外，由 Sets, Lemma
\ref{sets-lemma-what-is-in-it}，$\{V_i \to T\}_{i \in I}$
与位点 $\Sch_{fppf}$ 中 $T$ 的某个覆盖
$\{U_j \to T\}_{j \in J}$ 组合等价。覆盖
$\{U_j \to T\}_{j \in J}$ 正是 (1) 所需的加细。
在 (2)、(3) 的情形中，由 Sets, Lemma
\ref{sets-lemma-what-is-in-it}，每个 $T_i$ 都同构于 $\Sch_{fppf}$
的某个对象；再应用 Sets, Lemma \ref{sets-lemma-coverings-site}
即得所需结论。
\end{proof}

\begin{definition}
\label{topologies-definition-big-small-fppf}
设 $S$ 为概形，$\Sch_{fppf}$ 为包含 $S$ 的大 fppf 位点。
\begin{enumerate}
\item $S$ 的{\it 大 fppf 位点}记作 $(\Sch/S)_{fppf}$，它是
Sites, Section \ref{sites-section-localize} 中引入的位点 $\Sch_{fppf}/S$。
\item $S$ 的{\it 大仿射 fppf 位点}记作 $(\textit{Aff}/S)_{fppf}$，
它是 $(\Sch/S)_{fppf}$ 的全子范畴，其对象是仿射的 $U/S$。
$(\textit{Aff}/S)_{fppf}$ 的覆盖是 $(\Sch/S)_{fppf}$ 中任意的
标准 fppf 覆盖 $\{U_i \to U\}$。
\end{enumerate}
\end{definition}

\noindent
大仿射 fppf 位点确实是位点，这一点并非完全显然。下面予以验证。

\begin{lemma}
\label{topologies-lemma-verify-site-fppf}
设 $S$ 为概形，$\Sch_{fppf}$ 为包含 $S$ 的大 fppf 位点。
则 $(\textit{Aff}/S)_{fppf}$ 是位点。
\end{lemma}

\begin{proof}
证明 $(\textit{Aff}/S)_{fppf}$ 是位点。与 Lemma
\ref{topologies-lemma-verify-site-etale} 的证明相同，只需证明仿射概形的标准 fppf
覆盖族满足 Sites, Definition \ref{sites-definition-site}
中的性质 (1)、(2)、(3)。这是显然的。例如，给定标准 fppf 覆盖
$\{T_i \to T\}_{i\in I}$，且对每个 $i$，
$\{T_{ij} \to T_i\}_{j\in J_i}$ 是标准 fppf 覆盖，则
$\{T_{ij} \to T\}_{i \in I, j\in J_i}$ 是标准 fppf 覆盖，
因为 $\bigcup_{i\in I} J_i$ 有限且每个 $T_{ij}$ 都是仿射的。
\end{proof}

\begin{lemma}
\label{topologies-lemma-fibre-products-fppf}
设 $S$ 为概形，$\Sch_{fppf}$ 为包含 $S$ 的大 fppf 位点。
位点 $\Sch_{fppf}$、$(\Sch/S)_{fppf}$ 和 $(\textit{Aff}/S)_{fppf}$
的底层范畴都具有纤维积。在每种情形中，到全体概形范畴 $\Sch$
的显然函子都与取纤维积相容。范畴 $(\Sch/S)_{fppf}$ 有终对象，即 $S/S$。
\end{lemma}

\begin{proof}
对 $\Sch_{fppf}$ 而言，这是构造本身保证的；见 Sets, Lemma
\ref{sets-lemma-what-is-in-it}。设 $U \to S$、$V \to U$、$W \to U$
为概形态射，且 $U, V, W \in \Ob(\Sch_{fppf})$。
$\Sch_{fppf}$ 中的纤维积 $V \times_U W$ 也是 $\Sch$ 中的纤维积，
并且在 $S$ 上全体概形的范畴中，它是 $V/S$ 与 $W/S$ 在 $U/S$ 上的
纤维积，因而也是 $(\Sch/S)_{fppf}$ 中的纤维积。这证明了
$(\Sch/S)_{fppf}$ 的结论。若 $U, V, W$ 都仿射，则 $V \times_U W$
也仿射，从而得到 $(\textit{Aff}/S)_{fppf}$ 的结论。
\end{proof}

\noindent
下面验证大仿射位点与大位点定义同一个拓扑斯。

\begin{lemma}
\label{topologies-lemma-affine-big-site-fppf}
设 $S$ 为概形，$\Sch_{fppf}$ 为包含 $S$ 的大 fppf 位点。
函子 $(\textit{Aff}/S)_{fppf} \to (\Sch/S)_{fppf}$ 余连续，
并诱导从 $\Sh((\textit{Aff}/S)_{fppf})$ 到
$\Sh((\Sch/S)_{fppf})$ 的拓扑斯等价。
\end{lemma}

\begin{proof}
特殊余连续函子的概念见 Sites, Definition
\ref{sites-definition-special-cocontinuous-functor}。因此需要验证
Sites, Lemma \ref{sites-lemma-equivalence} 的假设 (1)--(5)。
将包含函子记作
$u : (\textit{Aff}/S)_{fppf} \to (\Sch/S)_{fppf}$。
余连续性恰好意味着：当 $T$ 仿射时，$T/S$ 的每个 fppf 覆盖都可由
$T$ 的某个标准 fppf 覆盖加细。这正是 Lemma
\ref{topologies-lemma-fppf-affine} 的内容，故 (1) 成立。由于标准 fppf 覆盖本身就是
fppf 覆盖，$u$ 显然连续，故 (2) 成立。(3) 和 (4) 由 $u$ 全忠实立即得到。
最后，每个概形都有仿射开覆盖，所以 (5) 成立。
\end{proof}

\noindent
下面建立这些位点所对应拓扑斯之间的一些关系。

\begin{lemma}
\label{topologies-lemma-morphism-big-fppf}
设 $\Sch_{fppf}$ 为大 fppf 位点，$f : T \to S$ 为
$\Sch_{fppf}$ 中的态射。函子
$$
u : (\Sch/T)_{fppf} \longrightarrow (\Sch/S)_{fppf},
\quad
V/T \longmapsto V/S
$$
余连续，并具有连续的右伴随
$$
v : (\Sch/S)_{fppf} \longrightarrow (\Sch/T)_{fppf},
\quad
(U \to S) \longmapsto (U \times_S T \to T).
$$
二者诱导同一个拓扑斯态射
$$
f_{big} :
\Sh((\Sch/T)_{fppf})
\longrightarrow
\Sh((\Sch/S)_{fppf})
$$
有 $f_{big}^{-1}(\mathcal{G})(U/T) = \mathcal{G}(U/S)$，以及
$f_{big, *}(\mathcal{F})(U/S) = \mathcal{F}(U \times_S T/T)$。
此外，$f_{big}^{-1}$ 有左伴随 $f_{big!}$，它与纤维积和等化子相容。
\end{lemma}

\begin{proof}
函子 $u$ 余连续、连续，并与纤维积和等化子相容。因此可以应用
Sites, Lemmas \ref{sites-lemma-when-shriek} 和
\ref{sites-lemma-preserve-equalizers}，从而得到 $f_{big}^{-1}$
的公式以及 $f_{big!}$ 的存在性。此外，函子 $v$ 是右伴随，因为给定
$U/T$ 和 $V/S$，恰有
$\Mor_S(u(U), V) = \Mor_T(U, V \times_S T)$。于是可应用 Sites, Lemmas
\ref{sites-lemma-have-functor-other-way} 和
\ref{sites-lemma-have-functor-other-way-morphism}，得到 $f_{big, *}$ 的公式。
\end{proof}

\begin{lemma}
\label{topologies-lemma-composition-fppf}
给定 $(\Sch/S)_{fppf}$ 中的概形 $X$、$Y$、$Z$ 以及态射
$f : X \to Y$、$g : Y \to Z$，有
$g_{big} \circ f_{big} = (g \circ f)_{big}$.
\end{lemma}

\begin{proof}
这由 Lemma \ref{topologies-lemma-morphism-big-fppf} 中大位点上函子的推前与拉回的
简单描述得到。
\end{proof}



















































\section{ph 拓扑}
\label{topologies-section-ph}

\noindent
本节定义 ph 拓扑。它是由 Zariski 覆盖和固有满射态射生成的拓扑；
见 Lemma \ref{topologies-lemma-characterize-sheaf}。

\medskip\noindent
本文的记号和术语取自 Goodwillie 与 Lichtenbaum 的论文 \cite{ph}。
他们证明：若限制到 Noether 概形的子范畴，则 ph 拓扑与 Voevodsky
最初定义的“h 拓扑”相同；后者是由 Zariski 开覆盖和普遍商映射的有限型态射
生成的拓扑。他们还证明，这两个拓扑在非 Noether 概形上并不一致；
见 \cite[Example 4.5]{ph}。我们将在 More on Flatness, Section
\ref{flat-section-h} 中回到（本文版本的）h 拓扑。

\medskip\noindent
在定义这一拓扑的覆盖之前，还需要作一些准备。

\begin{definition}
\label{topologies-definition-standard-ph-covering}
设 $T$ 为仿射概形。一个{\it 标准 ph 覆盖}是如下构造的态射族
$\{f_j : U_j \to T\}_{j = 1, \ldots, m}$：取固有满射态射
$f : U \to T$ 和仿射开覆盖
$U = \bigcup_{j = 1, \ldots, m} U_j$，并令 $f_j = f|_{U_j}$。
\end{definition}

\noindent
由 Chow 引理立即可知，任意标准 ph 覆盖都可由对应于满射射影态射的
标准 ph 覆盖加细。

\begin{lemma}
\label{topologies-lemma-base-change-standard-ph}
设 $\{f_j : U_j \to T\}_{j = 1, \ldots, m}$ 是标准 ph 覆盖，
$T' \to T$ 是仿射概形态射。则
$\{U_j \times_T T' \to T'\}_{j = 1, \ldots, m}$ 是标准 ph 覆盖。
\end{lemma}

\begin{proof}
设固有满射态射 $f : U \to T$ 和仿射开覆盖
$U = \bigcup_{j = 1, \ldots, m} U_j$ 如 Definition
\ref{topologies-definition-standard-ph-covering} 所给。则
$U \times_T T' \to T'$ 是固有满射态射（Morphisms, Lemmas
\ref{morphisms-lemma-base-change-surjective} 和
\ref{morphisms-lemma-base-change-proper}）。此外，
$U \times_T T' = \bigcup_{j = 1, \ldots, m} U_j \times_T T'$
是仿射开覆盖。证明完毕。
\end{proof}

\begin{lemma}
\label{topologies-lemma-refine-by-standard-ph}
设 $T$ 为仿射概形。下列每类以 $T$ 为目标的映射族都可由标准 ph 覆盖加细：
\begin{enumerate}
\item $T$ 的任意 Zariski 开覆盖；
\item $\{W_{ji} \to T\}_{j = 1, \ldots, m, i = 1, \ldots n_j}$，
其中 $\{W_{ji} \to U_j\}_{i = 1, \ldots, n_j}$ 和
$\{U_j \to T\}_{j = 1, \ldots, m}$ 都是标准 ph 覆盖。
\end{enumerate}
\end{lemma}

\begin{proof}
(1) 成立，因为 $T$ 的任意 Zariski 开覆盖都可由有限仿射开覆盖加细。

\medskip\noindent
% P05-ERR-0002: frozen source says “Proof of (3)” in a two-item lemma.
第 (3) 项的证明。按 Definition \ref{topologies-definition-standard-ph-covering}
选取固有满射态射 $U \to T$ 以及
$U = \bigcup_{j = 1, \ldots, m} U_j$。再按 Definition
\ref{topologies-definition-standard-ph-covering} 选取固有满射态射
$W_j \to U_j$ 以及 $W_j = \bigcup W_{ji}$。由 Chow 引理
（Limits, Lemma \ref{limits-lemma-chow-finite-type}），可找到固有满射态射
$W'_j \to W_j$ 和闭浸入 $W'_j \to \mathbf{P}^{e_j}_{U_j}$。
因此，将 $W_j$ 替换为 $W'_j$，并将 $W_j = \bigcup W_{ji}$
替换为 $W'_j$ 的适当仿射开覆盖后，可以假设对所有
$j = 1, \ldots, m$ 都有闭浸入
$W_j \subset \mathbf{P}^{e_j}_{U_j}$。

\medskip\noindent
令 $\overline{W}_j \subset \mathbf{P}^{e_j}_U$ 为 $W_j$ 的概形论闭包。
则 $W_j \subset \overline{W}_j$ 是开子概形；事实上，在态射
$\overline{W}_j \to U$ 下，$W_j$ 是 $U_j \subset U$ 的逆像。
（为说明这一点，使用 $W_j \to \mathbf{P}^{e_j}_U$ 拟紧，因而概形论像的
形成与限制到开集相容；见 Morphisms, Section
\ref{morphisms-section-scheme-theoretic-image}。）令
$Z_j = U \setminus U_j$ 配备诱导的约化闭子概形结构。于是
$$
V_j = \overline{W}_j \amalg Z_j \to U
$$
是固有满射态射，而开子概形 $W_j \subset V_j$ 是 $U_j$ 的逆像。
因此，对 $v \in V_j$、$v \not \in W_j$，可以选取仿射开邻域
$v \in V_{j, v} \subset V_j$，使其映入某个满足
$1 \leq j' \leq m$ 的 $U_{j'}$。

\medskip\noindent
为完成证明，考虑固有满射态射
$$
V = V_1 \times_U V_2 \times_U \ldots \times_U V_m
\longrightarrow U \longrightarrow T
$$
以及由下列仿射开集组成的 $V$ 的覆盖
$$
V_{1, v_1} \times_U \ldots \times_U V_{j - 1, v_{j - 1}}
\times_U W_{j i} \times_U
V_{j + 1, v_{j + 1}} \times_U \ldots \times_U V_{m, v_m}
$$
这些开集确实构成覆盖，因为 $U$ 的每个点都属于某个 $U_j$，
而 $U_j$ 在 $V$ 中的逆像等于
$V_1 \times \ldots \times V_{j - 1} \times W_j \times V_{j + 1} \times
\ldots \times V_m$。注意，从上述仿射开集到 $T$ 的态射经过 $W_{ji}$
分解，因而得到一个加细。最后，只需要这些仿射开集中的有限多个，
因为 $V$ 拟紧（它在仿射概形 $T$ 上固有）。
\end{proof}

\begin{definition}
\label{topologies-definition-ph-covering}
设 $T$ 为概形。$T$ 的一个{\it ph 覆盖}是概形态射族
$\{f_i : T_i \to T\}_{i \in I}$，其中每个 $f_i$ 都局部有限型，
并且对每个仿射开集 $U \subset T$，都存在标准 ph 覆盖
$\{U_j \to U\}_{j = 1, \ldots, m}$，它加细态射族
$\{T_i \times_T U \to U\}_{i \in I}$。
\end{definition}

\noindent
由 Lemma \ref{topologies-lemma-base-change-standard-ph}，标准 ph 覆盖是 ph 覆盖。

\begin{lemma}
\label{topologies-lemma-zariski-ph}
Zariski 覆盖是 ph 覆盖\footnote{我们将在 More on Morphisms, Lemma
\ref{more-morphisms-lemma-fppf-ph} 中看到，fppf 覆盖（因而 syntomic、
光滑或平展覆盖）也都是 ph 覆盖。}。
\end{lemma}

\begin{proof}
这是因为，由 Lemma \ref{topologies-lemma-refine-by-standard-ph}，
仿射概形的 Zariski 覆盖可由标准 ph 覆盖加细。
\end{proof}

\begin{lemma}
\label{topologies-lemma-surjective-proper-ph}
设 $f : Y \to X$ 为概形之间的固有满射态射。则
$\{Y \to X\}$ 是 ph 覆盖。
\end{lemma}

\begin{proof}
略。
\end{proof}

\begin{lemma}
\label{topologies-lemma-refine-by-ph}
设 $T$ 为概形，$\{f_i : T_i \to T\}_{i \in I}$ 为态射族，
且对所有 $i$，$f_i$ 都局部有限型。下列条件等价：
\begin{enumerate}
\item $\{T_i \to T\}_{i \in I}$ 是 ph 覆盖；
\item 存在加细 $\{T_i \to T\}_{i \in I}$ 的 ph 覆盖；
\item $\{\coprod_{i \in I} T_i \to T\}$ 是 ph 覆盖。
\end{enumerate}
\end{lemma}

\begin{proof}
(1) 与 (2) 的等价性由 Definition \ref{topologies-definition-ph-covering}
以及“加细的加细仍是加细”立即得到。由于 (1) 与 (2) 等价，并且
$\{T_i \to T\}_{i \in I}$ 加细
$\{\coprod_{i \in I} T_i \to T\}$，可知 (1) 蕴含 (3)。
最后，设 (3) 成立。令 $U \subset T$ 为仿射开集，并令
$\{U_j \to U\}_{j = 1, \ldots, m}$ 为加细
$\{U \times_T \coprod_{i \in I} T_i \to U\}$ 的标准 ph 覆盖。
这意味着对每个 $j$，都有态射
$$
h_j :
U_j
\longrightarrow
U \times_T \coprod\nolimits_{i \in I} T_i =
\coprod\nolimits_{i \in I} U \times_T T_i
$$
位于 $U$ 上。由于 $U_j$ 拟紧，得到由开闭子概形组成的不交并分解
$U_j = \coprod_{i \in I} U_{j, i}$，其中几乎所有子概形都为空，
并且 $h_j|_{U_{j, i}}$ 将 $U_{j, i}$ 映入 $U \times_T T_i$。
于是
$$
\{U_{j, i} \to U\}_{j = 1, \ldots, m,\ i \in I,\ U_{j, i} \not = \emptyset}
$$
是加细 $\{U \times_T T_i \to U\}_{i \in I}$ 的标准 ph 覆盖
（一个小细节从略）。故 (1) 成立。
\end{proof}

\noindent
下面证明这一概念满足 Sites, Definition \ref{sites-definition-site} 的条件。

\begin{lemma}
\label{topologies-lemma-ph}
设 $T$ 为概形。
\begin{enumerate}
\item 若 $T' \to T$ 是同构，则 $\{T' \to T\}$ 是 $T$ 的 ph 覆盖。
\item 若 $\{T_i \to T\}_{i\in I}$ 是 ph 覆盖，且对每个 $i$，
$\{T_{ij} \to T_i\}_{j\in J_i}$ 是 ph 覆盖，则
$\{T_{ij} \to T\}_{i \in I, j\in J_i}$ 是 ph 覆盖。
\item 若 $\{T_i \to T\}_{i\in I}$ 是 ph 覆盖，且 $T' \to T$
是概形态射，则 $\{T' \times_T T_i \to T'\}_{i\in I}$ 是 ph 覆盖。
\end{enumerate}
\end{lemma}

\begin{proof}
(1) 显然。

\medskip\noindent
第 (3) 项的证明。由 Morphisms, Lemma
\ref{morphisms-lemma-base-change-finite-type}，基变换
$T_i \times_T T' \to T'$ 局部有限型。因此只需验证仿射开集上的条件。
设 $U' \subset T'$ 为仿射开子概形。由于 $U'$ 拟紧，可以找到有限仿射开覆盖
% P05-ERR-0003: frozen source ends this finite-cover equation with unindexed U'.
$U' = U'_1 \cup \ldots \cup U'$，使得 $U'_j \to T$ 映入某个仿射开集
$U_j \subset T$。选取加细 $\{T_i \times_T U_j \to U_j\}$ 的标准 ph 覆盖
$\{U_{jl} \to U_j\}_{l = 1, \ldots, n_j}$。由 Lemma
\ref{topologies-lemma-base-change-standard-ph}，基变换
$\{U_{jl} \times_{U_j} U'_j \to U'_j\}$ 是标准 ph 覆盖。
注意，$\{U'_j \to U'\}$ 也是标准 ph 覆盖。由 Lemma
\ref{topologies-lemma-refine-by-standard-ph}，态射族
$\{U_{jl} \times_{U_j} U'_j \to U'\}$ 可由标准 ph 覆盖加细。
由于 $\{U_{jl} \times_{U_j} U'_j \to U'\}$ 加细
$\{T_i \times_T U' \to U'\}$，结论成立。

\medskip\noindent
第 (2) 项的证明。复合保持局部有限型；见 Morphisms, Lemma
\ref{morphisms-lemma-composition-finite-type}。因此只需验证仿射开集上的条件。
设 $U \subset T$ 为仿射开集。先选取加细
$\{T_i \times_T U \to U\}$ 的标准 ph 覆盖
$\{U_k \to U\}_{k = 1, \ldots, m}$。设该加细由 $T$ 上的态射
$U_k \to T_{i_k}$ 给出。则
$$
\{T_{i_kj} \times_{T_{i_k}} U_k \to U_k\}_{j \in J_{i_k}}
$$
由第 (3) 项是 ph 覆盖。由于 $U_k$ 仿射，可找到加细该态射族的标准 ph 覆盖
$\{U_{ka} \to U_k\}_{a = 1, \ldots, b_k}$。再应用 Lemma
\ref{topologies-lemma-refine-by-standard-ph} 可知，$\{U_{ka} \to U\}$
可由标准 ph 覆盖加细。由于 $\{U_{ka} \to U\}$ 加细
$\{T_{ij} \times_T U \to U\}$，证明完毕。
\end{proof}

\begin{definition}
\label{topologies-definition-big-ph-site}
一个{\it 大 ph 位点}是 Sites, Definition
\ref{sites-definition-site} 意义下按如下方式构造的任一位点 $\Sch_{ph}$：
\begin{enumerate}
\item 任取一组概形 $S_0$，以及这些概形之间的一组 ph 覆盖 $\text{Cov}_0$。
\item 取由集合 $S_0$ 出发、按 Sets, Lemma
\ref{sets-lemma-construct-category} 构造的任一范畴 $\Sch_\alpha$
作为底层范畴。
\item 由范畴 $\Sch_\alpha$、ph 覆盖类以及上述集合 $\text{Cov}_0$
出发，按 Sets, Lemma \ref{sets-lemma-coverings-site} 选取任一组覆盖。
\end{enumerate}
\end{definition}

\noindent
关于大位点定义的动机和解释，见 Definition
\ref{topologies-definition-big-zariski-site} 后面的评注。

\medskip\noindent
在继续引入概形 $S$ 的大 ph 位点之前，先指出：大 ph 位点
$\Sch_{ph}$ 上的拓扑，在某种意义上由全体概形范畴上的 ph 拓扑诱导而来。

\begin{lemma}
\label{topologies-lemma-ph-induced}
设 $\Sch_{ph}$ 是 Definition \ref{topologies-definition-big-ph-site} 所述的大 ph 位点，
$T \in \Ob(\Sch_{ph})$，并设 $\{T_i \to T\}_{i \in I}$
是 $T$ 的任意 ph 覆盖。
\begin{enumerate}
\item 位点 $\Sch_{ph}$ 中存在 $T$ 的覆盖 $\{U_j \to T\}_{j \in J}$，
它加细 $\{T_i \to T\}_{i \in I}$。
\item 若 $\{T_i \to T\}_{i \in I}$ 是标准 ph 覆盖，则它与
$\Sch_{ph}$ 中的某个覆盖自明等价。
\item 若 $\{T_i \to T\}_{i \in I}$ 是 Zariski 覆盖，则它与
$\Sch_{ph}$ 中的某个覆盖自明等价。
\end{enumerate}
\end{lemma}

\begin{proof}
对每个 $i$，选取仿射开覆盖 $T_i = \bigcup_{j \in J_i} T_{ij}$，
使每个 $T_{ij}$ 都映入 $T$ 的某个仿射开子概形。由 Lemmas
\ref{topologies-lemma-zariski-ph} 和 \ref{topologies-lemma-ph}，其加细
$\{T_{ij} \to T\}_{i \in I, j \in J_i}$ 仍是 $T$ 的 ph 覆盖。
因此可以假设每个 $T_i$ 都是仿射的，并映入 $T$ 的某个仿射开集 $W_i$。
应用 Sets, Lemma \ref{sets-lemma-what-is-in-it} 可知，$W_i$
同构于 $\Sch_{ph}$ 的某个对象。于是再次应用 Sets, Lemma
\ref{sets-lemma-what-is-in-it}，作为 $W_i$ 上有限型概形的 $T_i$
同构于 $\Sch_{ph}$ 的某个对象 $V_i$。覆盖
$\{V_i \to T\}_{i \in I}$ 加细 $\{T_i \to T\}_{i \in I}$
（因为对应对象彼此同构）。此外，由 Sets, Lemma
\ref{sets-lemma-what-is-in-it}，$\{V_i \to T\}_{i \in I}$
与位点 $\Sch_{ph}$ 中 $T$ 的某个覆盖
$\{U_j \to T\}_{j \in J}$ 组合等价。覆盖
$\{U_j \to T\}_{j \in J}$ 正是 (1) 所需的加细。
在 (2)、(3) 的情形中，由 Sets, Lemma
\ref{sets-lemma-what-is-in-it}，每个 $T_i$ 都同构于 $\Sch_{ph}$
的某个对象；再应用 Sets, Lemma \ref{sets-lemma-coverings-site}
即得所需结论。
\end{proof}

\begin{definition}
\label{topologies-definition-big-small-ph}
设 $S$ 为概形，$\Sch_{ph}$ 为包含 $S$ 的大 ph 位点。
\begin{enumerate}
\item $S$ 的{\it 大 ph 位点}记作 $(\Sch/S)_{ph}$，它是 Sites,
Section \ref{sites-section-localize} 中引入的位点 $\Sch_{ph}/S$。
\item $S$ 的{\it 大仿射 ph 位点}记作 $(\textit{Aff}/S)_{ph}$，
它是 $(\Sch/S)_{ph}$ 的全子范畴，其对象是仿射的 $U/S$。
$(\textit{Aff}/S)_{ph}$ 的覆盖是 $(\Sch/S)_{ph}$ 中任意满足
$U_i$ 和 $U$ 都仿射的有限覆盖 $\{U_i \to U\}$。
\end{enumerate}
\end{definition}

\noindent
注意，$(\textit{Aff}/S)_{ph}$ 中的覆盖{\it 并非}由标准 ph 覆盖给出；
原因很简单：否则 Sites, Definition \ref{sites-definition-site}
的第二条公理将不成立。$(\textit{Aff}/S)_{ph}$ 中的覆盖其实是
$(\Sch/S)_{ph}$ 的仿射对象之间、可由标准 ph 覆盖加细的有限型态射的
有限族 $\{U_i \to U\}$。下面明确陈述并证明大仿射 ph 位点确实是位点。

\begin{lemma}
\label{topologies-lemma-verify-site-ph}
设 $S$ 为概形，$\Sch_{ph}$ 为包含 $S$ 的大 ph 位点。
则 $(\textit{Aff}/S)_{ph}$ 是位点。
\end{lemma}

\begin{proof}
与 Lemma \ref{topologies-lemma-verify-site-etale} 的证明相同，只需证明满足
$U$、$U_i$ 仿射的有限 ph 覆盖 $\{U_i \to U\}$ 所组成的集合满足
Sites, Definition \ref{sites-definition-site} 中的性质 (1)、(2)、(3)。
这是显然的。例如，给定 $T_i, T$ 仿射的有限 ph 覆盖
$\{T_i \to T\}_{i\in I}$，并且对每个 $i$，给定 $T_{ij}$ 仿射的有限
ph 覆盖 $\{T_{ij} \to T_i\}_{j\in J_i}$，则
$\{T_{ij} \to T\}_{i \in I, j\in J_i}$ 是 ph 覆盖
（Lemma \ref{topologies-lemma-ph}），$\bigcup_{i\in I} J_i$ 有限，且每个
$T_{ij}$ 都仿射。
\end{proof}

\begin{lemma}
\label{topologies-lemma-fibre-products-ph}
设 $S$ 为概形，$\Sch_{ph}$ 为包含 $S$ 的大 ph 位点。
位点 $\Sch_{ph}$、$(\Sch/S)_{ph}$ 和 $(\textit{Aff}/S)_{ph}$
的底层范畴都具有纤维积。在每种情形中，到全体概形范畴 $\Sch$
的显然函子都与取纤维积相容。范畴 $(\Sch/S)_{ph}$ 有终对象，即 $S/S$。
\end{lemma}

\begin{proof}
对 $\Sch_{ph}$ 而言，这是构造本身保证的；见 Sets, Lemma
\ref{sets-lemma-what-is-in-it}。设 $U \to S$、$V \to U$、$W \to U$
为概形态射，且 $U, V, W \in \Ob(\Sch_{ph})$。
$\Sch_{ph}$ 中的纤维积 $V \times_U W$ 也是 $\Sch$ 中的纤维积，
并且在 $S$ 上全体概形的范畴中，它是 $V/S$ 与 $W/S$ 在 $U/S$ 上的
纤维积，因而也是 $(\Sch/S)_{ph}$ 中的纤维积。这证明了
$(\Sch/S)_{ph}$ 的结论。若 $U, V, W$ 都仿射，则 $V \times_U W$
也仿射，从而得到 $(\textit{Aff}/S)_{ph}$ 的结论。
\end{proof}

\noindent
下面验证大仿射位点与大位点定义同一个拓扑斯。

\begin{lemma}
\label{topologies-lemma-affine-big-site-ph}
设 $S$ 为概形，$\Sch_{ph}$ 为包含 $S$ 的大 ph 位点。
函子 $(\textit{Aff}/S)_{ph} \to (\Sch/S)_{ph}$ 余连续，
并诱导从 $\Sh((\textit{Aff}/S)_{ph})$ 到
$\Sh((\Sch/S)_{ph})$ 的拓扑斯等价。
\end{lemma}

\begin{proof}
特殊余连续函子的概念见 Sites, Definition
\ref{sites-definition-special-cocontinuous-functor}。因此需要验证
Sites, Lemma \ref{sites-lemma-equivalence} 的假设 (1)--(5)。
将包含函子记作 $u : (\textit{Aff}/S)_{ph} \to (\Sch/S)_{ph}$。
余连续性由定义得到，因为当 $T$ 仿射时，$T/S$ 的任意 ph 覆盖都可由
$T$ 的某个标准 ph 覆盖加细，故 (1) 成立。由于仿射概形由仿射概形组成的
有限 ph 覆盖本身就是 ph 覆盖，$u$ 显然连续，故 (2) 成立。
(3) 和 (4) 由 $u$ 全忠实立即得到。最后，每个概形都有仿射开覆盖
（它是 ph 覆盖），所以 (5) 成立。
\end{proof}

\begin{lemma}
\label{topologies-lemma-characterize-sheaf}
设 $\mathcal{F}$ 是 $(\Sch/S)_{ph}$ 上的预层。则 $\mathcal{F}$ 是层，
当且仅当
\begin{enumerate}
\item $\mathcal{F}$ 对 Zariski 覆盖满足层条件；
\item 若 $f : V \to U$ 是固有满射态射，则 $\mathcal{F}(U)$
双射地映到两个映射 $\mathcal{F}(V) \to \mathcal{F}(V \times_U V)$
的等化子。
\end{enumerate}
此外，在 (1) 成立时，性质 (2) 等价于性质
\begin{enumerate}
\item[(2')] 当 $U$ 仿射时，对 (2) 所述的 $\{V \to U\}$ 满足层性质。
\end{enumerate}
\end{lemma}

\begin{proof}
证明若 (1) 和 (2) 成立，则 $\mathcal{F}$ 是层。设
$\{T_i \to T\}$ 为 ph 覆盖，即 $(\Sch/S)_{ph}$ 中的覆盖。
我们验证此覆盖的层条件。设 $s_i \in \mathcal{F}(T_i)$ 是一族截面，
它们限制到 $T_i \times_T T_{i'}$ 后相同。我们要证明存在唯一截面
$s \in \mathcal{F}(T)$，其在 $T_i$ 上限制为 $s_i$。
设 $T = \bigcup U_j$ 为仿射开覆盖。由性质 (1)，为构造 $s$，
只需构造截面 $s_j \in \mathcal{F}(U_j)$，使其在
$U_j \cap U_{j'}$ 上相同。考虑 ph 覆盖
$\{T_i \times_T U_j \to U_j\}$。则
$s_{ji} = s_i|_{T_i \times_T U_j}$ 是一族在
$(T_i \times_T U_j) \times_{U_j} (T_{i'} \times_T U_j)$ 上相同的截面。
选取固有满射态射 $V_j \to U_j$ 和有限仿射开覆盖
$V_j = \bigcup V_{jk}$，使标准 ph 覆盖 $\{V_{jk} \to U_j\}$ 加细
$\{T_i \times_T U_j \to U_j\}$。若
$s_{jk} \in \mathcal{F}(V_{jk})$ 表示由隐含态射将 $s_{ji}$ 拉回到
$V_{jk}$ 所得的截面，则 $s_{jk}$ 可黏合为截面
$s'_j \in \mathcal{F}(V_j)$。再次使用交叠处的一致性可知，$s'_j$
属于两个映射 $\mathcal{F}(V_j) \to \mathcal{F}(V_j \times_{U_j} V_j)$
的等化子。因此由 (2)，$s'_j$ 来自唯一截面
$s_j \in \mathcal{F}(U_j)$。关于这些截面 $s_j$ 具有全部所需性质的验证从略。

\medskip\noindent
在 (1) 成立时证明 (2) 与 (2') 等价。设 $V \to U$ 是
$(\Sch/S)_{ph}$ 中的固有满射态射。选取仿射开覆盖
$U = \bigcup U_i$，并令 $V_i = V \times_U U_i$。由 (2')，
$\mathcal{F}(U_i) \to \mathcal{F}(V_i)$ 是单射；由 (1)，
$\mathcal{F}(U) \to \prod \mathcal{F}(U_i)$ 是单射。因此
$\mathcal{F}(U) \to \mathcal{F}(V)$ 是单射。最后，设给定
$t \in \mathcal{F}(V)$，它属于两个映射
$\mathcal{F}(V) \to \mathcal{F}(V \times_U V)$ 的等化子。
则对所有 $i$，$t|_{V_i}$ 都属于两个映射
$\mathcal{F}(V_i) \to \mathcal{F}(V_i \times_{U_i} V_i)$ 的等化子。
因此由 (2')，对所有 $i$ 都得到唯一的截面
$s_i \in \mathcal{F}(U_i)$，它映到 $t|_{V_i}$。关于对所有 $i, j$
都有 $s_i|_{U_i \cap U_j} = s_j|_{U_i \cap U_j}$ 的验证从略；
这里使用刚刚证明的唯一性。由覆盖 $U = \bigcup U_i$ 的层性质，
得到截面 $s \in \mathcal{F}(U)$。关于 $s$ 映到
$\mathcal{F}(V)$ 中的 $t$ 的证明从略。
\end{proof}

\noindent
下面建立这些位点所对应拓扑斯之间的一些关系。

\begin{lemma}
\label{topologies-lemma-morphism-big-ph}
设 $\Sch_{ph}$ 为大 ph 位点，$f : T \to S$ 为 $\Sch_{ph}$ 中的态射。
函子
$$
u : (\Sch/T)_{ph} \longrightarrow (\Sch/S)_{ph},
\quad
V/T \longmapsto V/S
$$
余连续，并具有连续的右伴随
$$
v : (\Sch/S)_{ph} \longrightarrow (\Sch/T)_{ph},
\quad
(U \to S) \longmapsto (U \times_S T \to T).
$$
二者诱导同一个拓扑斯态射
$$
f_{big} :
\Sh((\Sch/T)_{ph})
\longrightarrow
\Sh((\Sch/S)_{ph})
$$
有 $f_{big}^{-1}(\mathcal{G})(U/T) = \mathcal{G}(U/S)$，以及
$f_{big, *}(\mathcal{F})(U/S) = \mathcal{F}(U \times_S T/T)$。
此外，$f_{big}^{-1}$ 有左伴随 $f_{big!}$，它与纤维积和等化子相容。
\end{lemma}

\begin{proof}
函子 $u$ 余连续、连续，并与纤维积和等化子相容。因此可以应用
Sites, Lemmas \ref{sites-lemma-when-shriek} 和
\ref{sites-lemma-preserve-equalizers}，从而得到 $f_{big}^{-1}$
的公式以及 $f_{big!}$ 的存在性。此外，函子 $v$ 是右伴随，因为给定
$U/T$ 和 $V/S$，恰有
$\Mor_S(u(U), V) = \Mor_T(U, V \times_S T)$。于是可应用 Sites, Lemmas
\ref{sites-lemma-have-functor-other-way} 和
\ref{sites-lemma-have-functor-other-way-morphism}，得到 $f_{big, *}$ 的公式。
\end{proof}

\begin{lemma}
\label{topologies-lemma-composition-ph}
% P05-ERR-0004: frozen source repeats Y and omits Z; see part-local errata ledger.
给定 $(\Sch/S)_{ph}$ 中的概形 $X$、$Y$、$Y$ 以及态射
$f : X \to Y$、$g : Y \to Z$，有
$g_{big} \circ f_{big} = (g \circ f)_{big}$.
\end{lemma}

\begin{proof}
这由 Lemma \ref{topologies-lemma-morphism-big-ph} 中大位点上函子的推前与拉回的
简单描述得到。
\end{proof}



































\section{fpqc 拓扑}
\label{topologies-section-fpqc}

\begin{definition}
\label{topologies-definition-fpqc-covering}
设 $T$ 为概形。$T$ 的一个{\it fpqc 覆盖}是概形态射族
$\{f_i : T_i \to T\}_{i \in I}$，其中每个 $f_i$ 都平坦，并且对每个
仿射开集 $U \subset T$，存在 $n \geq 0$、映射
$a : \{1, \ldots, n\} \to I$ 以及仿射开集
$V_j \subset T_{a(j)}$，$j = 1, \ldots, n$，使得
$\bigcup_{j = 1}^n f_{a(j)}(V_j) = U$。
\end{definition}

\noindent
这一条件当然蕴含 $T = \bigcup f_i(T_i)$。识别 fpqc 覆盖稍为困难，
所以下面给出若干判别引理。

\begin{lemma}
\label{topologies-lemma-recognize-fpqc-covering}
设 $T$ 为概形，$\{f_i : T_i \to T\}_{i \in I}$ 为以 $T$ 为目标的
概形态射族。下列条件等价：
\begin{enumerate}
\item $\{f_i : T_i \to T\}_{i \in I}$ 是 fpqc 覆盖；
\item 每个 $f_i$ 都平坦，并且对每个仿射开集 $U \subset T$，
存在几乎全为空的拟紧开集 $U_i \subset T_i$，使得
$U = \bigcup f_i(U_i)$；
\item 每个 $f_i$ 都平坦，并且存在仿射开覆盖
$T = \bigcup_{\alpha \in A} U_\alpha$，使得对每个 $\alpha \in A$，
都存在 $i_{\alpha, 1}, \ldots, i_{\alpha, n(\alpha)} \in I$
和拟紧开集 $U_{\alpha, j} \subset T_{i_{\alpha, j}}$，满足
$U_\alpha =
\bigcup_{j = 1, \ldots, n(\alpha)} f_{i_{\alpha, j}}(U_{\alpha, j})$。
\end{enumerate}
若 $T$ 拟分离，则这些条件还等价于
\begin{enumerate}
\item[(4)] 每个 $f_i$ 都平坦，并且对每个 $t \in T$，存在
$i_1, \ldots, i_n \in I$ 和拟紧开集 $U_j \subset T_{i_j}$，
使得 $\bigcup_{j = 1, \ldots, n} f_{i_j}(U_j)$ 是 $t$ 在 $T$
中的一个邻域（不必是开集）。
\end{enumerate}
\end{lemma}

\begin{proof}
关于 (1)、(2)、(3) 等价的证明从略。以下假设 $T$ 拟分离。
证明 (4) 蕴含 (2)。设 $U \subset T$ 为仿射开集。为证明 (2)，
只需说明：对每个 $t \in U$，存在有限多个拟紧开集
$U_j \subset T_{i_j}$，使得 $f_{i_j}(U_j) \subset U$，并且
$\bigcup f_{i_j}(U_j)$ 是 $t$ 在 $U$ 中的邻域。由假设，确实存在有限多个
拟紧开集 $U'_j \subset T_{i_j}$，使得
% P05-ERR-0005: frozen source duplicates “such that” at this locus.
$\bigcup f_{i_j}(U'_j)$ 是 $t$ 在 $T$ 中的邻域。由于 $T$ 拟分离，
$U_j = U'_j \cap f_{i_j}^{-1}(U) = (f_{i_j}|_{U'_j})^{-1}(U)$
由 Schemes, Lemma \ref{schemes-lemma-quasi-compact-permanence}，
它拟紧。故 (2) 成立。由于 (2) 显然蕴含 (4)，证明完毕。
\end{proof}

\begin{example}
\label{topologies-example-not-fpqc}
令 $X = X_1 \cup X_2$，其中
$X_1 = X_2 = \Spec(k[t_1,t_2,\dots])$，取 Schemes, Example
\ref{schemes-example-not-quasi-separated} 中的非拟分离概形，并令
$Y = X_1 \amalg \Spec(\mathcal{O}_{X_2, 0})$。显然的态射
$f : Y \to X$ 平坦且满射，并且（因为 $Y$ 拟紧）显然满足 Lemma
\ref{topologies-lemma-recognize-fpqc-covering} 的 (4)。但是
$\{f : Y \to X\}$ 不是 fpqc 覆盖。具体而言，我们断言：不存在 $Y$
的拟紧开集 $W$，使其在 $f$ 下的像为 $X_2$。事实上，
$f^{-1}(X_2) = X_1 \setminus \{0\} \amalg \Spec(\mathcal{O}_{X_2, 0})$。
因此，若 $W$ 存在，则对 $X_1 \setminus \{0\} = X_2 \setminus \{0\}$
的某个拟紧开集 $U$，必有
$W = U \amalg \Spec(\mathcal{O}_{X_2, 0})$。于是可写
$U = D(f_1) \cup \ldots \cup D(f_n)$，其中某些
$f_i \in k[t_1, t_2, \ldots]$ 在 $0$ 处为零。设
% P05-ERR-0006: frozen source switches from t-variables to x-variables.
$f_1, \ldots, f_n$ 只使用变量 $x_1, \ldots, x_m$。则
$p = (0, \ldots, 0, 1, 0, \ldots)$（其中 $1$ 位于第 $(m + 1)$ 个位置）
是 $X_2$ 的一个闭点，但不在 $W$ 的像中。
\end{example}

\begin{lemma}
\label{topologies-lemma-disjoint-union-is-fpqc-covering}
设 $T$ 为概形，$\{f_i : T_i \to T\}_{i \in I}$ 为以 $T$ 为目标的
概形态射族。下列条件等价：
\begin{enumerate}
\item $\{f_i : T_i \to T\}_{i \in I}$ 是 fpqc 覆盖；
\item 令 $T' = \coprod_{i \in I} T_i$ 以及 $f = \coprod_{i \in I} f_i$，
则态射族 $\{f : T' \to T\}$ 是 fpqc 覆盖。
\end{enumerate}
\end{lemma}

\begin{proof}
设 $U \subset T$ 为仿射开集。若 (1) 成立，则可找到
$i_1, \ldots, i_n \in I$ 和仿射开集 $U_j \subset T_{i_j}$，使得
$U = \bigcup_{j = 1, \ldots, n} f_{i_j}(U_j)$。于是
$U_1 \amalg \ldots \amalg U_n \subset T'$ 是满射到 $U$ 的拟紧开集。
因此由 Lemma \ref{topologies-lemma-recognize-fpqc-covering}，
$\{f : T' \to T\}$ 是 fpqc 覆盖。反过来，若 (2) 成立，则存在拟紧开集
$U' \subset T'$，使得 $U = f(U')$。此时 $U_j = U' \cap T_j$
是 $T_j$ 中的拟紧开集，并且对几乎所有 $j$ 都为空。由 Lemma
\ref{topologies-lemma-recognize-fpqc-covering} 可知 (1) 成立。
\end{proof}

\begin{lemma}
\label{topologies-lemma-family-flat-dominated-covering}
设 $T$ 为概形，$\{f_i : T_i \to T\}_{i \in I}$ 为一族以 $T$ 为目标的
概形态射。假设：
\begin{enumerate}
\item 每个 $f_i$ 都是平坦的；并且
\item 族 $\{f_i : T_i \to T\}_{i \in I}$ 可由 $T$ 的某个 fpqc 覆盖加细。
\end{enumerate}
则 $\{f_i : T_i \to T\}_{i \in I}$ 是 $T$ 的 fpqc 覆盖。
\end{lemma}

\begin{proof}
设 $\{g_j : X_j \to T\}_{j \in J}$ 是加细
$\{f_i : T_i \to T\}$ 的 fpqc 覆盖。设 $U \subset T$ 为仿射开集。
选取 $j_1, \ldots, j_m \in J$ 及仿射开集 $V_k \subset X_{j_k}$，使得
$U = \bigcup g_{j_k}(V_k)$。对每个 $j$，选取 $i_j \in I$ 以及态射
$h_j : X_j \to T_{i_j}$，使得 $g_j = f_{i_j} \circ h_j$。
由于 $h_{j_k}(V_k)$ 拟紧，可找到拟紧开集
$h_{j_k}(V_k) \subset U_k \subset f_{i_{j_k}}^{-1}(U)$。
于是 $U = \bigcup f_{i_{j_k}}(U_k)$。由 Lemma
\ref{topologies-lemma-recognize-fpqc-covering} 可知
$\{f_i : T_i \to T\}_{i \in I}$ 是 fpqc 覆盖。
\end{proof}

\begin{lemma}
\label{topologies-lemma-family-flat-fpqc-local-covering}
设 $T$ 为概形，$\{f_i : T_i \to T\}_{i \in I}$ 为一族以 $T$ 为目标的
概形态射。假设：
\begin{enumerate}
\item 每个 $f_i$ 都是平坦的；并且
\item 存在 fpqc 覆盖 $\{g_j : S_j \to T\}_{j \in J}$，使得每个
$\{S_j \times_T T_i \to S_j\}_{i \in I}$ 都是 fpqc 覆盖。
\end{enumerate}
则 $\{f_i : T_i \to T\}_{i \in I}$ 是 $T$ 的 fpqc 覆盖。
\end{lemma}

\begin{proof}
以下不再另行说明地使用 Lemma \ref{topologies-lemma-recognize-fpqc-covering}。
设 $U \subset T$ 为仿射开集。由 (2)，可对 $j \in J$ 找到拟紧开集
$V_j \subset S_j$，其中除有限多个外均为空，并且
$U = \bigcup g_j(V_j)$。于是对每个 $j$，可对 $i \in I$ 选取拟紧开集
$W_{ij} \subset S_j \times_T T_i$，其中除有限多个外均为空，并满足
$V_j = \bigcup_i \text{pr}_1(W_{ij})$。因此
$\{S_j \times_T T_i \to T\}$ 是 fpqc 覆盖。
此覆盖加细 $\{f_i : T_i \to T\}$，故结论由 Lemma
\ref{topologies-lemma-family-flat-dominated-covering} 得出。
\end{proof}

\begin{lemma}
\label{topologies-lemma-zariski-etale-smooth-syntomic-fppf-fpqc}
任意 fppf 覆盖都是 fpqc 覆盖；从而尤其，任意 syntomic、光滑、
\'etale 或 Zariski 覆盖都是 fpqc 覆盖。
\end{lemma}

\begin{proof}
我们证明 fppf 覆盖是 fpqc 覆盖；其余结论随即由 Lemma
\ref{topologies-lemma-zariski-etale-smooth-syntomic-fppf} 得出。
设 $\{f_i : U_i \to U\}_{i \in I}$ 为 fppf 覆盖。
按定义，这意味着各 $f_i$ 平坦，因而满足 Definition
\ref{topologies-definition-fpqc-covering} 的第一个条件。为验证第二个条件，
设 $V \subset U$ 为仿射开子集。对某些仿射开集 $V_{ij} \subset U_i$，写成
$f_i^{-1}(V) = \bigcup_{j \in J_i} V_{ij}$。由于每个 $f_i$ 都是开态射
（Morphisms, Lemma \ref{morphisms-lemma-fppf-open}），可见
$V = \bigcup_{i\in I} \bigcup_{j \in J_i} f_i(V_{ij})$
是 $V$ 的开覆盖。由于 $V$ 拟紧，此覆盖有有限加细。这就完成了证明。
\end{proof}

\noindent
fpqc\footnote{字母 fpqc 代表
``fid\`element plat quasi-compacte''.}
拓扑不能像 fppf 拓扑那样处理\footnote{更准确地说，Lemma
\ref{topologies-lemma-fppf-induced} 对 fpqc 拓扑的类似结论并不成立。}。
确切而言，设 $R$ 为非零环。我们将在 Lemma
\ref{topologies-lemma-no-set-of-fpqc-covers-is-initial} 中看到，不存在由
$\Spec(R)$ 的 fpqc 覆盖组成的集合 $A$，使得每个 fpqc 覆盖都可由
$A$ 的某个元素加细。若 $R = k$ 是域，则这种无界性的原因在于：
不存在 $k$ 的一个域扩张，其中包含 $k$ 的每个域扩张。

\medskip\noindent
若忽略集合论方面的困难，就会遇到无法层化的预层；见
\cite[Theorem 5.5]{Waterhouse-fpqc-sheafification}。
一个略有意思的选择，是只考虑那些忠实平坦环扩张 $R \to R'$，
其中 $R'$ 的基数受到适当限制。（若像 SGA4 那样考虑某个固定宇宙中的
所有概形，便是用一个强不可达基数来限制该基数。）然而，把这个基数
换成更大的基数后会发生什么，并不十分清楚。

\medskip\noindent
基于这些原因，我们不引入 fpqc 位点，也不考虑关于 fpqc 拓扑的上同调。

\medskip\noindent
另一方面，给定反变函子 $F : \Sch^{opp} \to \textit{Sets}$，
考察 $F$ 是否满足 fpqc 拓扑的层性质是有意义的，见下文。
此外，还可以研究 fpqc 拓扑中对象的下降等问题。简言之，对某些结果而言，
恰当的一般性正是使用 fpqc 覆盖。

\begin{lemma}
\label{topologies-lemma-fpqc}
设 $T$ 为概形。
\begin{enumerate}
\item 若 $T' \to T$ 是同构，则 $\{T' \to T\}$ 是 $T$ 的 fpqc 覆盖。
\item 若 $\{T_i \to T\}_{i\in I}$ 是 fpqc 覆盖，并且对每个 $i$ 都有
fpqc 覆盖 $\{T_{ij} \to T_i\}_{j\in J_i}$，则
$\{T_{ij} \to T\}_{i \in I, j\in J_i}$ 是 fpqc 覆盖。
\item 若 $\{T_i \to T\}_{i\in I}$ 是 fpqc 覆盖，且 $T' \to T$
是概形态射，则 $\{T' \times_T T_i \to T'\}_{i\in I}$ 是 fpqc 覆盖。
\end{enumerate}
\end{lemma}

\begin{proof}
第 (1) 项是立即的。回忆平坦态射的复合是平坦的，而平坦态射的基变换
也是平坦的（Morphisms, Lemmas \ref{morphisms-lemma-base-change-flat} and
\ref{morphisms-lemma-composition-flat}）。因此在每种情形下，均可应用
Lemma \ref{topologies-lemma-recognize-fpqc-covering} 验证所给态射族为 fpqc 覆盖。

\medskip\noindent
第 (2) 项的证明。设 $\{T_i \to T\}_{i\in I}$ 是 fpqc 覆盖，并且对每个
$i$ 都有 fpqc 覆盖 $\{f_{ij} : T_{ij} \to T_i\}_{j\in J_i}$。
设 $U \subset T$ 为仿射开集。可对 $i \in I$ 找到拟紧开集
$U_i \subset T_i$，其中除有限多个外均为空，并且
$U = \bigcup f_i(U_i)$。于是对每个 $i$，可对 $j \in J_i$ 选取拟紧开集
$U_{ij} \subset T_{ij}$，其中除有限多个外均为空，并满足
$U_i = \bigcup_j f_{ij}(U_{ij})$。因此
$\{T_{ij} \to T\}$ 是 fpqc 覆盖。

\medskip\noindent
第 (3) 项的证明。设 $\{T_i \to T\}_{i\in I}$ 是 fpqc 覆盖，且
$T' \to T$ 是概形态射。设 $U' \subset T'$ 为映入仿射开集
$U \subset T$ 的仿射开集。选取拟紧开集 $U_i \subset T_i$，其中除有限多个外
均为空，并使 $U = \bigcup f_i(U_i)$。则 $U' \times_U U_i$ 是
$T' \times_T T_i$ 的拟紧开集，并且
$U' = \bigcup \text{pr}_1(U' \times_U U_i)$。由于 $T'$ 可由这样的仿射开集
$U' \subset T'$ 覆盖，Lemma \ref{topologies-lemma-recognize-fpqc-covering} 表明
$\{T' \times_T T_i \to T'\}_{i\in I}$ 是 fpqc 覆盖。
\end{proof}

\begin{lemma}
\label{topologies-lemma-fpqc-affine}
设 $T$ 为仿射概形，$\{T_i \to T\}_{i \in I}$ 为 $T$ 的 fpqc 覆盖。
则存在加细 $\{T_i \to T\}_{i \in I}$ 的 fpqc 覆盖
$\{U_j \to T\}_{j = 1, \ldots, n}$，使每个 $U_j$ 都是仿射概形。
此外，还可选取每个 $U_j$，使其为某个 $T_i$ 中的仿射开集。
\end{lemma}

\begin{proof}
这直接由定义得出。
\end{proof}

\begin{definition}
\label{topologies-definition-standard-fpqc}
设 $T$ 为仿射概形。$T$ 的一个{\it 标准 fpqc 覆盖}，是一个态射族
$\{f_j : U_j \to T\}_{j = 1, \ldots, n}$，其中每个 $U_j$ 都是仿射的、
在 $T$ 上平坦，并且 $T = \bigcup f_j(U_j)$。
\end{definition}

\noindent
由于我们不引入仿射位点，必须直接证明所有标准 fpqc 覆盖的集合满足公理。

\begin{lemma}
\label{topologies-lemma-fpqc-affine-axioms}
设 $T$ 为仿射概形。
\begin{enumerate}
\item 若 $T' \to T$ 是同构，则 $\{T' \to T\}$ 是 $T$ 的标准 fpqc 覆盖。
\item 若 $\{T_i \to T\}_{i\in I}$ 是标准 fpqc 覆盖，并且对每个 $i$ 都有
标准 fpqc 覆盖 $\{T_{ij} \to T_i\}_{j\in J_i}$，则
$\{T_{ij} \to T\}_{i \in I, j\in J_i}$ 是标准 fpqc 覆盖。
\item 若 $\{T_i \to T\}_{i\in I}$ 是标准 fpqc 覆盖，且 $T' \to T$
是仿射概形态射，则 $\{T' \times_T T_i \to T'\}_{i\in I}$ 是标准
fpqc 覆盖。
\end{enumerate}
\end{lemma}

\begin{proof}
这形式地源于以下事实：平坦态射的复合和基变换都是平坦的
（Morphisms, Lemmas \ref{morphisms-lemma-base-change-flat} and
\ref{morphisms-lemma-composition-flat}），而仿射概形的纤维积仍为仿射概形
（Schemes, Lemma \ref{schemes-lemma-fibre-product-affines}）。
\end{proof}

\begin{lemma}
\label{topologies-lemma-fpqc-covering-affines-mapping-in}
设 $T$ 为概形，$\{f_i : T_i \to T\}_{i \in I}$ 为一族以 $T$ 为目标的
概形态射。假设：
\begin{enumerate}
\item 每个 $f_i$ 都是平坦的；并且
% P05-ERR-0007: frozen English omits the preposition before "every affine scheme".
\item 对每个仿射概形 $Z$ 及态射 $h : Z \to T$，都存在标准 fpqc 覆盖
$\{Z_j \to Z\}_{j = 1, \ldots, n}$，它加细态射族
$\{T_i \times_T Z \to Z\}_{i \in I}$。
\end{enumerate}
则 $\{f_i : T_i \to T\}_{i \in I}$ 是 $T$ 的 fpqc 覆盖。
\end{lemma}

\begin{proof}
设 $T = \bigcup U_\alpha$ 为仿射开覆盖。对每个 $\alpha$，拉回态射族
$\{T_i \times_T U_\alpha \to U_\alpha\}$ 可由标准 fpqc 覆盖加细，
故由 Lemma \ref{topologies-lemma-family-flat-dominated-covering}，它是 fpqc 覆盖。
又因 $\{U_\alpha \to T\}$ 是 fpqc 覆盖，Lemma
\ref{topologies-lemma-family-flat-fpqc-local-covering} 表明
$\{T_i \to T\}$ 是 fpqc 覆盖。
\end{proof}

\begin{definition}
\label{topologies-definition-sheaf-property-fpqc}
设 $F$ 是从概形范畴到集合范畴的反变函子。
\begin{enumerate}
\item 设 $\{U_i \to T\}_{i \in I}$ 为一族具有固定目标的概形态射。
若对任意一族元素 $\xi_i \in F(U_i)$，只要
$\xi_i|_{U_i \times_T U_j} = \xi_j|_{U_i \times_T U_j}$，就存在唯一元素
$\xi \in F(T)$，使得在 $F(U_i)$ 中有 $\xi_i = \xi|_{U_i}$，则称 $F$
{\it 对给定态射族满足层性质}。
\item 若该函子对任意 fpqc 覆盖都满足层性质，则称 $F$
{\it 满足 fpqc 拓扑的层性质}。
\end{enumerate}
\end{definition}

\noindent
在这种情形下，我们尽量避免使用“$F$ 是一个层”这一术语，因为如上所述，
我们并未定义 fpqc 层的范畴。

\begin{lemma}
\label{topologies-lemma-sheaf-property-fpqc}
设 $F$ 是从概形范畴到集合范畴的反变函子。则 $F$ 满足 fpqc 拓扑的层性质，
当且仅当它满足：
\begin{enumerate}
\item 对每个 Zariski 覆盖的层性质；以及
\item 对任意标准 fpqc 覆盖的层性质。
\end{enumerate}
此外，在 (1) 成立的前提下，性质 (2) 等价于下列性质：
\begin{enumerate}
\item[(2')] 对 $\{V \to U\}$ 的层性质，其中 $V$、$U$ 均为仿射概形，
且 $V \to U$ 忠实平坦。
\end{enumerate}
\end{lemma}

\begin{proof}
假设 (1) 和 (2) 成立。设 $\{f_i : T_i \to T\}_{i \in I}$ 为 fpqc 覆盖，
并设 $s_i \in F(T_i)$ 为一族元素，使得 $s_i$ 与 $s_j$ 映到
$F(T_i \times_T T_j)$ 中的同一元素。令 $W \subset T$ 为满足下列性质的
最大开子集：存在唯一的 $s \in F(W)$，使得对所有 $i$ 都有
$s|_{f_i^{-1}(W)} = s_i|_{f_i^{-1}(W)}$。这样的最大开集存在，因为 $F$
对 Zariski 覆盖满足层性质；事实上，$W$ 是所有具有此性质的开集之并。
设 $t \in T$。我们将证明 $t \in W$。为此，选取仿射开集
$t \in U \subset T$，并证明存在唯一的 $s \in F(U)$，使得对所有 $i$
都有 $s|_{f_i^{-1}(U)} = s_i|_{f_i^{-1}(U)}$。

\medskip\noindent
由 Lemma \ref{topologies-lemma-fpqc-affine}，可找到加细
$\{U \times_T T_i \to U\}$ 的标准 fpqc 覆盖
$\{U_j \to U\}_{j = 1, \ldots, n}$；设给出加细的态射为
$h_j : U_j \to T_{i_j}$。由 (2)，得到唯一元素 $s \in F(U)$，满足
$s|_{U_j} = F(h_j)(s_{i_j})$。注意，对于任意 $U$ 上的概形 $V \to U$，
都存在唯一截面 $s_V \in F(V)$，使其在 $V \times_U U_j$ 上的限制为
$F(h_j \circ \text{pr}_2)(s_{i_j})$，其中 $j = 1, \ldots, n$。
确实，若 $V$ 仿射，则因 $\{V \times_U U_j \to V\}$ 是标准 fpqc 覆盖，
这由 (2) 成立；一般情形则通过选取 $V$ 的仿射开覆盖，由 (1) 和仿射情形得出。
特别地，$s_V = s|_V$。现在取 $V = U \times_T T_i$，并利用
$s_{i_j}|_{T_{i_j} \times_T T_i} = s_i|_{T_{i_j} \times_T T_i}$，可得
$s|_{U \times_T T_i} = s_V = s_i|_{U \times_T T_i}$，这正是所需结论。

\medskip\noindent
在 (1) 成立的前提下证明 (2) 与 (2') 等价。设 $\{T_i \to T\}$ 是标准
fpqc 覆盖，则 $\coprod T_i \to T$ 是仿射概形之间的忠实平坦态射。
在 (1) 成立的前提下，有 $F(\coprod T_i) = \prod F(T_i)$，类似地还有
$F((\coprod T_i) \times_T (\coprod T_i)) = \prod F(T_i \times_T T_{i'})$。
因此，$\{T_i \to T\}$ 与 $\{\coprod T_i \to T\}$ 的层条件相同。
\end{proof}

\noindent
下面这个引理只是为了指出，集合论困难确实会出现；大多数读者可以忽略它。

\begin{lemma}
\label{topologies-lemma-no-set-of-fpqc-covers-is-initial}
设 $R$ 为非零环。不存在由 $\Spec(R)$ 的 fpqc 覆盖组成的集合 $A$，
使得每个 fpqc 覆盖都可由 $A$ 的某个元素加细。
\end{lemma}

\begin{proof}
先说明 $R = k$ 为域的情形。对任意集合 $I$，考虑纯超越域扩张
$k_I = k(\{t_i\}_{i \in I})/k$。由于 $k \to k_I$ 忠实平坦，
$\{\Spec(k_I) \to \Spec(k)\}$ 是 fpqc 覆盖。设 $A$ 为集合，并对每个
$\alpha \in A$，令
$\mathcal{U}_\alpha = \{S_{\alpha, j} \to \Spec(k)\}_{j \in J_\alpha}$
为 fpqc 覆盖。若 $\mathcal{U}_\alpha$ 加细
$\{\Spec(k_I) \to \Spec(k)\}$，则各态射 $S_{\alpha, j} \to \Spec(k)$
都经由 $\Spec(k_I)$ 分解。由于 $\mathcal{U}_\alpha$ 是覆盖，至少有某个
$S_{\alpha, j}$ 非空。选取一点 $s \in S_{\alpha, j}$。由分解
$S_{\alpha, j} \to \Spec(k_I) \to \Spec(k)$，得到域同态
$k_I \to \kappa(s)$。特别地，$\kappa(s)$ 的基数至少是 $I$ 的基数。
因此，若取集合 $I$ 的基数大于所有概形 $S_{\alpha, j}$ 的剩余域的基数，
则不存在这样的分解，引理对 $R = k$ 成立。

\medskip\noindent
一般情形。由于 $R$ 非零，它有一个极大素理想 $\mathfrak m$，其剩余域为
$\kappa$。设 $I$ 为集合，并考虑 $R_I = S_I^{-1} R[\{t_i\}_{i \in I}]$，
其中 $S_I \subset R[\{t_i\}_{i \in I}]$ 是由如下元素
$f \in R[\{t_i\}_{i \in I}]$ 组成的乘法子集：对 $R$ 的每个素理想
$\mathfrak p$，$f$ 在 $R/\mathfrak p[\{t_i\}_{i \in I}]$ 中的像均非零。
则 $R_I$ 是忠实平坦的 $R$-代数，而
$\{\Spec(R_I) \to \Spec(R)\}$ 是 fpqc 覆盖。留给读者作为练习证明，
沿用上述记号，有
$R_I \otimes_R \kappa \cong \kappa(\{t_i\}_{i \in I}) = \kappa_I$
（提示：利用 $R \to \kappa$ 是满射，以及任意
$f \in R[\{t_i\}_{i \in I}]$，只要其某个单项式的系数为 $1$，便属于
$S_I$）。设 $A$ 为集合，并对每个 $\alpha \in A$，令
$\mathcal{U}_\alpha = \{S_{\alpha, j} \to \Spec(R)\}_{j \in J_\alpha}$
为 fpqc 覆盖。若 $\mathcal{U}_\alpha$ 加细
$\{\Spec(R_I) \to \Spec(R)\}$，则基变换给出
$\{S_{\alpha, j} \times_{\Spec(R)} \Spec(\kappa) \to \Spec(\kappa)\}$
加细 $\{\Spec(\kappa_I) \to \Spec(\kappa)\}$。因此，由前一段的结果，
存在某个 $I$ 使其不成立，于是引理得证。
\end{proof}












\section{V 拓扑}
\label{topologies-section-V}

\noindent
V 拓扑强于本章中的所有其他拓扑。粗略地说，它由 Zariski 覆盖以及满足
特化提升性质的拟紧态射生成（Lemma \ref{topologies-lemma-refine-qcqs-V}）。
不过，我们用来定义 V 覆盖的步骤略有不同。我们将先定义仿射概形的标准
V 覆盖，再用它们定义一般的 V 覆盖。记号方面，在文献中有时用
“$v$-覆盖”代替“V 覆盖”。

\begin{definition}
\label{topologies-definition-standard-V-covering}
设 $T$ 为仿射概形。一个{\it 标准 V 覆盖}是有限族
$\{T_j \to T\}_{j = 1, \ldots, m}$，其中 $T_j$ 仿射，并且对每个态射
$g : \Spec(V) \to T$，只要 $V$ 是赋值环，就存在赋值环扩张
$V \subset W$（More on Algebra, Definition
\ref{more-algebra-definition-extension-valuation-rings}）、指标
$1 \leq j \leq m$ 以及交换图
$$
\xymatrix{
\Spec(W) \ar[r] \ar[d] & T_j \ar[d] \\
\Spec(V) \ar[r]^g & T
}
$$
\end{definition}

\noindent
我们先证明关于这一概念的几个基本引理。

\begin{lemma}
\label{topologies-lemma-standard-fpqc-standard-V}
标准 fpqc 覆盖是标准 V 覆盖。
\end{lemma}

\begin{proof}
设 $\{X_i \to X\}_{i = 1, \ldots, n}$ 为标准 fpqc 覆盖
（Definition \ref{topologies-definition-standard-fpqc}）。设
$g : \Spec(V) \to X$ 为态射，其中 $V$ 是赋值环。令 $x \in X$ 为
$\Spec(V)$ 的闭点之像。选取一个 $i$ 以及映到 $x$ 的点
$x_i \in X_i$。则 $\Spec(V) \times_X X_i$ 有一点 $x'_i$，它映到
$\Spec(V)$ 的闭点。由于 $\Spec(V) \times_X X_i \to \Spec(V)$ 平坦，
可找到 $\Spec(V) \times_X X_i$ 中点的一个特化
$x''_i \leadsto x'_i$，其中 $x''_i$ 映到 $\Spec(V)$ 的泛点；见
Morphisms, Lemma \ref{morphisms-lemma-generalizations-lift-flat}。
由 Schemes, Lemma \ref{schemes-lemma-points-specialize}，可选取赋值环
$W$ 以及态射 $h : \Spec(W) \to \Spec(V) \times_X X_i$，使得 $h$
将 $\Spec(W)$ 的泛点映到 $x''_i$，并将 $\Spec(W)$ 的闭点映到 $x'_i$。
由此得到交换图
$$
\xymatrix{
\Spec(W) \ar[r] \ar[d] & X_i \ar[d] \\
\Spec(V) \ar[r] & X
}
$$
其中 $V \to W$ 是赋值环扩张。这就证明了引理。
\end{proof}

\begin{lemma}
\label{topologies-lemma-standard-ph-standard-V}
标准 ph 覆盖是标准 V 覆盖。
\end{lemma}

\begin{proof}
设 $T$ 为仿射概形，$f : U \to T$ 为固有满射态射，并设
$U = \bigcup_{j = 1, \ldots, m} U_j$ 为有限仿射开覆盖。我们必须证明
$\{U_j \to T\}$ 是标准 V 覆盖；见
Definition \ref{topologies-definition-standard-ph-covering}。设
$g : \Spec(V) \to T$ 为态射，其中 $V$ 是以 $K$ 为分式域的赋值环。
由于 $U \to T$ 满射，可以选取域扩张 $L/K$ 以及交换图
$$
\xymatrix{
\Spec(L) \ar[rr] \ar[d] & & U \ar[d] \\
\Spec(K) \ar[r] & \Spec(V) \ar[r]^g & T
}
$$
由 Algebra, Lemma \ref{algebra-lemma-dominate}，可选取支配 $V$ 的赋值环
$W \subset L$。再由固有性的赋值判别准则（Morphisms, Lemma
\ref{morphisms-lemma-characterize-proper}），可在下列交换图中找到态射 $h$：
$$
\xymatrix{
\Spec(L) \ar[r] \ar[d] & \Spec(W) \ar[r]_h \ar[d] & U \ar[d] \\
% P05-ERR-0008: frozen English diagram uses X although the ambient target is T.
\Spec(K) \ar[r] & \Spec(V) \ar[r]^g & X
}
$$
由于 $\Spec(W)$ 有唯一闭点，可见对某个 $j$ 有 $\Im(h)$ 包含于 $U_j$。
因此，$h : \Spec(W) \to U_j$ 是所需的提升，从而
$\{U_j \to T\}$ 是标准 V 覆盖。
\end{proof}

\begin{lemma}
\label{topologies-lemma-base-change-standard-V}
设 $\{T_j \to T\}_{j = 1, \ldots, m}$ 为标准 V 覆盖，并设
$T' \to T$ 为仿射概形态射。则
$\{T_j \times_T T' \to T'\}_{j = 1, \ldots, m}$ 是标准 V 覆盖。
\end{lemma}

\begin{proof}
设 $\Spec(V) \to T'$ 为态射，其中 $V$ 是赋值环。由假设，可找到赋值环扩张
$V \subset W$、一个 $i$ 以及交换图
$$
\xymatrix{
\Spec(W) \ar[r] \ar[d] & T_i \ar[d] \\
\Spec(V) \ar[r] & T
}
$$
由纤维积的泛性质，得到所需的态射
$\Spec(W) \to T' \times_T T_i$。
\end{proof}

\begin{lemma}
\label{topologies-lemma-composition-standard-V}
设 $T$ 为仿射概形，$\{T_j \to T\}_{j = 1, \ldots, m}$ 为标准 V 覆盖，
并设 $\{T_{ji} \to T_j\}_{i = 1, \ldots n_j}$ 为标准 V 覆盖。
则 $\{T_{ji} \to T\}_{i, j}$ 是标准 V 覆盖。
\end{lemma}

\begin{proof}
这形式地源于以下观察：若 $V \subset W$ 与 $W \subset \Omega$
是赋值环扩张，则 $V \subset \Omega$ 也是赋值环扩张。
\end{proof}

\begin{lemma}
\label{topologies-lemma-refine-standard-V}
设 $T$ 为仿射概形，$\{T_j \to T\}_{j = 1, \ldots, m}$ 为一族态射，
其中对所有 $j$，$T_j$ 均仿射。下列条件等价：
\begin{enumerate}
\item $\{T_j \to T\}_{j = 1, \ldots, m}$ 是标准 V 覆盖；
\item 存在加细 $\{T_j \to T\}_{j = 1, \ldots, m}$ 的标准 V 覆盖；
\item $\{\coprod_{j = 1, \ldots, m} T_j \to T\}$ 是标准 V 覆盖。
\end{enumerate}
\end{lemma}

\begin{proof}
证明从略。提示：这几乎立即由定义得出。唯一略有意思之处是，从局部环的谱
到 $\coprod_{j = 1, \ldots, m} T_j$ 的态射必经由某个 $T_j$ 分解。
\end{proof}

\begin{definition}
\label{topologies-definition-V-covering}
设 $T$ 为概形。$T$ 的一个{\it V 覆盖}，是概形态射族
$\{T_i \to T\}_{i \in I}$，使得对每个仿射开集 $U \subset T$，
都存在标准 V 覆盖 $\{U_j \to U\}_{j = 1, \ldots, m}$，它加细态射族
$\{T_i \times_T U \to U\}_{i \in I}$。
\end{definition}

\noindent
V 拓扑具有与 fpqc 拓扑相同的集合论问题。因此，我们不定义 V 位点，
也不考虑关于 V 拓扑的上同调。另一方面，给定
$F : \Sch^{opp} \to \textit{Sets}$，考察 $F$ 是否满足 V 拓扑的层性质
是有意义的，见下文。此外，还可以研究 V 拓扑中对象的下降等问题。

\begin{lemma}
\label{topologies-lemma-refine-by-V}
设 $T$ 为概形，$\{f_i : T_i \to T\}_{i \in I}$ 为一族态射。下列条件等价：
\begin{enumerate}
\item $\{T_i \to T\}_{i \in I}$ 是 V 覆盖；
\item 存在加细 $\{T_i \to T\}_{i \in I}$ 的 V 覆盖；
\item $\{\coprod_{i \in I} T_i \to T\}$ 是 V 覆盖。
\end{enumerate}
\end{lemma}

\begin{proof}
证明从略。提示：与 Lemma \ref{topologies-lemma-refine-by-ph} 的证明比较。
\end{proof}

\begin{lemma}
\label{topologies-lemma-V}
设 $T$ 为概形。
\begin{enumerate}
\item 若 $T' \to T$ 是同构，则 $\{T' \to T\}$ 是 $T$ 的 V 覆盖。
\item 若 $\{T_i \to T\}_{i\in I}$ 是 V 覆盖，并且对每个 $i$ 都有
V 覆盖 $\{T_{ij} \to T_i\}_{j\in J_i}$，则
$\{T_{ij} \to T\}_{i \in I, j\in J_i}$ 是 V 覆盖。
\item 若 $\{T_i \to T\}_{i\in I}$ 是 V 覆盖，且 $T' \to T$
是概形态射，则 $\{T' \times_T T_i \to T'\}_{i\in I}$ 是 V 覆盖。
\end{enumerate}
\end{lemma}

\begin{proof}
断言 (1) 是显然的。

\medskip\noindent
第 (3) 项的证明。设 $U' \subset T'$ 为仿射开子概形。由于 $U'$ 拟紧，
% P05-ERR-0009: frozen finite-cover expression ends with the whole U' rather than an indexed member.
可找到有限仿射开覆盖 $U' = U'_1 \cup \ldots \cup U'$，使得
$U'_j \to T$ 映入某个仿射开集 $U_j \subset T$。选取加细
$\{T_i \times_T U_j \to U_j\}$ 的标准 V 覆盖
$\{U_{jl} \to U_j\}_{l = 1, \ldots, n_j}$。由 Lemma
\ref{topologies-lemma-base-change-standard-V}，基变换
$\{U_{jl} \times_{U_j} U'_j \to U'_j\}$ 是标准 V 覆盖。注意，
$\{U'_j \to U'\}$ 是标准 V 覆盖（例如由 Lemma
\ref{topologies-lemma-standard-fpqc-standard-V}）。由 Lemma
\ref{topologies-lemma-composition-standard-V}，态射族
$\{U_{jl} \times_{U_j} U'_j \to U'\}$ 是标准 V 覆盖。由于
$\{U_{jl} \times_{U_j} U'_j \to U'\}$ 加细
$\{T_i \times_T U' \to U'\}$，结论得证。

\medskip\noindent
第 (2) 项的证明。设 $U \subset T$ 为仿射开集。先选取加细
$\{T_i \times_T U \to U\}$ 的标准 V 覆盖
$\{U_k \to U\}_{k = 1, \ldots, m}$。设该加细由 $T$ 上的态射
$U_k \to T_{i_k}$ 给出。则
$$
\{T_{i_kj} \times_{T_{i_k}} U_k \to U_k\}_{j \in J_{i_k}}
$$
由第 (3) 项是 V 覆盖。由于 $U_k$ 仿射，可找到加细此态射族的标准 V 覆盖
$\{U_{ka} \to U_k\}_{a = 1, \ldots, b_k}$。再应用 Lemma
\ref{topologies-lemma-composition-standard-V} 可知，$\{U_{ka} \to U\}$ 是加细
$\{T_{ij} \times_T U \to U\}$ 的标准 V 覆盖。证明完毕。
\end{proof}

\begin{lemma}
\label{topologies-lemma-zariski-etale-smooth-syntomic-fppf-fpqc-ph-V}
任意 fpqc 覆盖都是 V 覆盖。从而尤其，任意 fppf、syntomic、光滑、
\'etale 或 Zariski 覆盖都是 V 覆盖。此外，ph 覆盖也是 V 覆盖。
\end{lemma}

\begin{proof}
fpqc 覆盖在仿射局部可由标准 fpqc 覆盖加细；见 Lemmas
\ref{topologies-lemma-fpqc-affine}。标准 fpqc 覆盖是标准 V 覆盖；见 Lemma
\ref{topologies-lemma-standard-fpqc-standard-V}。因此，第一个陈述由我们用标准 V 覆盖
给出的 V 覆盖定义得出。关于 fppf、syntomic、光滑、\'etale 或 Zariski
覆盖的结论，则源于它们都是 fpqc 覆盖；见 Lemma
\ref{topologies-lemma-zariski-etale-smooth-syntomic-fppf-fpqc}。

\medskip\noindent
关于 ph 覆盖的陈述以同样方式由 Lemma
\ref{topologies-lemma-standard-ph-standard-V} 得出。
\end{proof}

\begin{definition}
\label{topologies-definition-sheaf-property-V}
设 $F$ 是从概形范畴到集合范畴的反变函子。若 $F$ 对任意 V 覆盖都满足
层性质，则称它{\it 满足 V 拓扑的层性质}
（见 Definition \ref{topologies-definition-sheaf-property-fpqc}）。
\end{definition}

\noindent
在这种情形下，我们尽量避免使用“$F$ 是一个层”这一术语，因为如上所述，
我们并未定义 V 层的范畴。

\begin{lemma}
\label{topologies-lemma-sheaf-property-V}
设 $F$ 是从概形范畴到集合范畴的反变函子。则 $F$ 满足 V 拓扑的层性质，
当且仅当它满足：
\begin{enumerate}
\item 对每个 Zariski 覆盖的层性质；以及
\item 对任意标准 V 覆盖的层性质。
\end{enumerate}
此外，在 (1) 成立的前提下，性质 (2) 等价于下列性质：
\begin{enumerate}
\item[(2')] 对形如 $\{V \to U\}$（即只由一个箭头组成）的标准 V 覆盖
满足层性质。
\end{enumerate}
\end{lemma}

\begin{proof}
假设 (1) 和 (2) 成立。设 $\{f_i : T_i \to T\}_{i \in I}$ 为 V 覆盖，
并设 $s_i \in F(T_i)$ 为一族元素，使得 $s_i$ 与 $s_j$ 映到
$F(T_i \times_T T_j)$ 中的同一元素。令 $W \subset T$ 为满足下列性质的
最大开子集：存在唯一的 $s \in F(W)$，使得对所有 $i$ 都有
$s|_{f_i^{-1}(W)} = s_i|_{f_i^{-1}(W)}$。这样的最大开集存在，因为 $F$
对 Zariski 覆盖满足层性质；事实上，$W$ 是所有具有此性质的开集之并。
设 $t \in T$。我们将证明 $t \in W$。为此，选取仿射开集
$t \in U \subset T$，并证明存在唯一的 $s \in F(U)$，使得对所有 $i$
都有 $s|_{f_i^{-1}(U)} = s_i|_{f_i^{-1}(U)}$。

\medskip\noindent
可找到加细 $\{U \times_T T_i \to U\}$ 的标准 V 覆盖
$\{U_j \to U\}_{j = 1, \ldots, n}$；设给出加细的态射为
$h_j : U_j \to T_{i_j}$。由 (2)，得到唯一元素 $s \in F(U)$，满足
$s|_{U_j} = F(h_j)(s_{i_j})$。注意，对于任意 $U$ 上的概形 $V \to U$，
都存在唯一截面 $s_V \in F(V)$，使其在 $V \times_U U_j$ 上的限制为
$F(h_j \circ \text{pr}_2)(s_{i_j})$，其中 $j = 1, \ldots, n$。
确实，若 $V$ 仿射，则因 $\{V \times_U U_j \to V\}$ 是标准 V 覆盖
（Lemma \ref{topologies-lemma-base-change-standard-V}），这由 (2) 成立；一般情形则
通过选取 $V$ 的仿射开覆盖，由 (1) 和仿射情形得出。特别地，
$s_V = s|_V$。现在取 $V = U \times_T T_i$，并利用
$s_{i_j}|_{T_{i_j} \times_T T_i} = s_i|_{T_{i_j} \times_T T_i}$，可得
$s|_{U \times_T T_i} = s_V = s_i|_{U \times_T T_i}$，这正是所需结论。

\medskip\noindent
在 (1) 成立的前提下证明 (2) 与 (2') 等价。设
$\{T_i \to T\}_{i = 1, \ldots, n}$ 是标准 V 覆盖，则
$\coprod_{i = 1, \ldots, n} T_i \to T$ 是仿射概形态射，且显然也是标准
V 覆盖。在 (1) 成立的前提下，有 $F(\coprod T_i) = \prod F(T_i)$，
类似地还有
$F((\coprod T_i) \times_T (\coprod T_i)) = \prod F(T_i \times_T T_{i'})$。
因此，$\{T_i \to T\}$ 与 $\{\coprod T_i \to T\}$ 的层条件相同。
\end{proof}

\noindent
下列引理表明，成为 V 覆盖与能否提升特化有关。

\begin{lemma}
\label{topologies-lemma-refine-qcqs-V}
设 $X \to Y$ 为概形的拟紧态射。下列条件等价：
\begin{enumerate}
\item $\{X \to Y\}$ 是 V 覆盖；
\item 对任意赋值环 $V$ 及态射 $g : \Spec(V) \to Y$，都存在赋值环扩张
$V \subset W$ 以及交换图
$$
\xymatrix{
\Spec(W) \ar[r] \ar[d] & X \ar[d] \\
\Spec(V) \ar[r] & Y
}
$$
\item 对任意态射 $Z \to Y$ 及 $Z$ 中点的特化 $z' \leadsto z$，
都存在 $Z \times_Y X$ 中点的特化 $w' \leadsto w$，它映到
$z' \leadsto z$。
\end{enumerate}
\end{lemma}

\begin{proof}
假设 (1)，并设 $g : \Spec(V) \to Y$ 如 (2) 中所述。由于 $V$ 是局部环，
存在仿射开集 $U \subset Y$，使 $g$ 经由 $U$ 分解。由 Definition
\ref{topologies-definition-V-covering}，可找到加细 $\{X \times_Y U \to U\}$ 的
标准 V 覆盖 $\{U_j \to U\}$。由 Definition
\ref{topologies-definition-standard-V-covering}，可找到一个 $j$、赋值环扩张
$V \subset W$ 以及交换图
$$
\xymatrix{
\Spec(W) \ar[r] \ar[d] & U_j \ar[d] \ar@{..>}[r] & X \ar[ld] \\
\Spec(V) \ar[r] & Y
}
$$
由该覆盖的加细性质，存在使图交换的虚线箭头，从而 (2) 成立。

\medskip\noindent
假设 (2)，并设 $Z \to Y$ 与 $z' \leadsto z$ 如 (3) 中所述。由 Schemes,
Lemma \ref{schemes-lemma-points-specialize}，可找到赋值环 $V$ 以及态射
$\Spec(V) \to Z$，使 $\Spec(V)$ 的闭点映到 $z$，而 $\Spec(V)$ 的泛点映到 $z'$。
由 (2)，可找到赋值环扩张 $V \subset W$ 以及交换图
$$
\xymatrix{
\Spec(W) \ar[rr] \ar[d] & & X \ar[d] \\
\Spec(V) \ar[r] & Z \ar[r] & Y
}
$$
$\Spec(W)$ 的泛点和闭点，经由诱导态射
$\Spec(W) \to Z \times_Y X$ 映到 $Z \times_Y X$ 中的点
$w' \leadsto w$。这表明 (3) 成立。

\medskip\noindent
假设 (3) 成立，并设 $U \subset Y$ 为仿射开集。选取有限仿射开覆盖
$U \times_Y X = \bigcup_{j = 1, \ldots, m} U_j$。这是可能的，因为
$X \to Y$ 拟紧。我们断言 $\{U_j \to U\}$ 是标准 V 覆盖。
该断言蕴含 (1)，并将完成引理的证明。为证明此断言，设 $V$ 为赋值环，
并设 $g : \Spec(V) \to U$ 为态射。由 (3)，可找到下列概形中点的特化
$w' \leadsto w$：
% P05-ERR-0010: frozen fibre products reverse X and Y relative to X -> Y and g : Spec(V) -> U -> Y.
$$
T = \Spec(V) \times_X Y = \Spec(V) \times_U (U \times_X Y)
$$
其中 $w'$ 映到 $\Spec(V)$ 的泛点，而 $w$ 映到 $\Spec(V)$ 的闭点。
由 Schemes, Lemma \ref{schemes-lemma-points-specialize}，可找到赋值环 $W$
以及态射 $\Spec(W) \to T$，使 $\Spec(W)$ 的泛点映到 $w'$，而
$\Spec(W)$ 的闭点
映到 $w$。复合 $\Spec(W) \to T \to \Spec(V)$ 对应于一个嵌入
$V \subset W$，它把 $W$ 表为赋值环 $V$ 的一个扩张。由于
$T = \bigcup \Spec(V) \times_U U_j$ 是开覆盖，可见对某个 $j$，
$\Spec(W) \to T$ 经由 $\Spec(V) \times_U U_j$ 分解。因此得到交换图
$$
\xymatrix{
\Spec(W) \ar[d] \ar[r] & U_j \ar[d] \\
\Spec(V) \ar[r] & U
}
$$
断言得证。
\end{proof}

\noindent
V 覆盖给出一族普遍商映射。下列引理的逆命题不成立；见 Examples,
Section \ref{examples-section-universally-submersive-not-V}。

\begin{lemma}
\label{topologies-lemma-V-covering-universally-submersive}
设 $\{f_i : X_i \to X\}_{i \in I}$ 为 V 覆盖。则
$$
\coprod\nolimits_{i \in I} f_i :
\coprod\nolimits_{i \in I} X_i
\longrightarrow
X
$$
是概形的普遍商态射（Morphisms, Definition
\ref{morphisms-definition-submersive}）。
\end{lemma}

\begin{proof}
以下不再另行说明地使用 V 覆盖的基变换仍是 V 覆盖这一事实
（Lemma \ref{topologies-lemma-V}）。特别地，只需证明该态射是商态射。
商态射这一性质显然在基底上是 Zariski 局部的。因此可假设 $X$ 仿射。
此时 $\{X_i \to X\}$ 可由标准 V 覆盖 $\{Y_j \to X\}$ 加细。
若能证明 $\coprod Y_j \to X$ 是商态射，则由分解
$\coprod Y_j \to \coprod X_i \to X$ 可知 $\coprod X_i \to X$ 是商态射。
令 $Y = \coprod Y_j$，并考虑仿射概形态射 $f : Y \to X$。
由 Lemma \ref{topologies-lemma-refine-qcqs-V}，$X$ 中任意特化
$x' \leadsto x$ 均可提升为 $Y$ 中的某个特化 $y' \leadsto y$。
因此，若 $T \subset X$ 为使 $f^{-1}(T)$ 在 $Y$ 中闭的子集，则
$T \subset X$ 在特化下封闭。赋予 $f^{-1}(T) \subset Y$ 诱导的既约
闭子概形结构后，它是仿射概形，故由 Algebra, Lemma
\ref{algebra-lemma-image-stable-specialization-closed}，$T \subset X$ 是闭集。
因此 $f$ 是商态射。
\end{proof}





















\section{拓扑的变换}
\label{topologies-section-change-topologies}

\noindent
设 $f : X \to Y$ 是基概形 $S$ 上的概形态射。此时有下列位点态射
\footnote{这里没有列入 ph 拓扑与其他拓扑之间的比较；关于这一点，见
More on Morphisms, Remark \ref{more-morphisms-remark-change-topologies}。}
（按下面 Remark \ref{topologies-remark-choice-sites} 中的方式适当选取位点）：
\begin{enumerate}
\item $(\Sch/X)_{fppf} \longrightarrow (\Sch/Y)_{fppf}$,
\item $(\Sch/X)_{fppf} \longrightarrow (\Sch/Y)_{syntomic}$,
\item $(\Sch/X)_{fppf} \longrightarrow (\Sch/Y)_{smooth}$,
\item $(\Sch/X)_{fppf} \longrightarrow
(\Sch/Y)_\etale$,
\item $(\Sch/X)_{fppf} \longrightarrow (\Sch/Y)_{Zar}$,
\item $(\Sch/X)_{syntomic} \longrightarrow (\Sch/Y)_{syntomic}$,
\item $(\Sch/X)_{syntomic} \longrightarrow (\Sch/Y)_{smooth}$,
\item $(\Sch/X)_{syntomic} \longrightarrow
(\Sch/Y)_\etale$,
\item $(\Sch/X)_{syntomic} \longrightarrow (\Sch/Y)_{Zar}$,
\item $(\Sch/X)_{smooth} \longrightarrow (\Sch/Y)_{smooth}$,
\item $(\Sch/X)_{smooth} \longrightarrow
(\Sch/Y)_\etale$,
\item $(\Sch/X)_{smooth} \longrightarrow (\Sch/Y)_{Zar}$,
\item $(\Sch/X)_\etale \longrightarrow
(\Sch/Y)_\etale$,
\item $(\Sch/X)_\etale \longrightarrow (\Sch/Y)_{Zar}$,
\item $(\Sch/X)_{Zar} \longrightarrow (\Sch/Y)_{Zar}$,
\item $(\Sch/X)_{fppf} \longrightarrow Y_\etale$,
\item $(\Sch/X)_{syntomic} \longrightarrow Y_\etale$,
\item $(\Sch/X)_{smooth} \longrightarrow Y_\etale$,
\item $(\Sch/X)_\etale \longrightarrow Y_\etale$,
\item $(\Sch/X)_{fppf} \longrightarrow Y_{Zar}$,
\item $(\Sch/X)_{syntomic} \longrightarrow Y_{Zar}$,
\item $(\Sch/X)_{smooth} \longrightarrow Y_{Zar}$,
\item $(\Sch/X)_\etale \longrightarrow Y_{Zar}$,
\item $(\Sch/X)_{Zar} \longrightarrow Y_{Zar}$,
\item $X_\etale \longrightarrow Y_\etale$,
\item $X_\etale \longrightarrow Y_{Zar}$,
\item $X_{Zar} \longrightarrow Y_{Zar}$,
\end{enumerate}
在每种情形下，底层连续函子 $\Sch/Y \to \Sch/X$ 或
$Y_\tau \to \Sch/X$，都是函子 $Y'/Y \mapsto X \times_Y Y'/X$。
确实，在前面各节中，对上述 $\tau$ 已经见到了态射
$f_{big} : (\Sch/X)_\tau \to (\Sch/Y)_\tau$ 以及
$f_{small} : X_\tau \to Y_\tau$。我们还见到了位点态射
$\pi_Y : (\Sch/Y)_\tau \to Y_\tau$，其中
$\tau \in \{\etale, Zariski\}$。另一方面，当 $\tau$ 是强于 $\tau'$ 的
拓扑时，恒等函子 $(\Sch/X)_\tau \to (\Sch/X)_{\tau'}$ 显然定义一个
位点态射。因此，将这些态射复合便得到上面的所有可能态射。

\medskip\noindent
由底层函子的简单描述可立即看出，给定概形态射 $X \to Y \to Z$，
上面两个位点态射的复合，例如
$$
(\Sch/X)_{\tau_0} \longrightarrow
(\Sch/Y)_{\tau_1} \longrightarrow
(\Sch/Z)_{\tau_2}
$$
就是与概形态射 $X \to Z$ 相应的位点态射。

\begin{remark}
\label{topologies-remark-choice-sites}
从概形集合 $\{X, Y, S\}$ 出发，任取一个按 Sets, Lemma
\ref{sets-lemma-construct-category} 构造的范畴 $\Sch_\alpha$。
再从范畴 $\Sch_\alpha$ 和 fppf 覆盖类出发，按 Sets, Lemma
\ref{sets-lemma-coverings-site} 任取 $\Sch_\alpha$ 上的覆盖集合
$\text{Cov}_{fppf}$。以 $\Sch_{fppf}$ 表示由此得到的大 fppf 位点。
接着，对 $\tau \in \{Zariski, \etale, smooth, syntomic\}$，令
$\Sch_\tau$ 与 $\Sch_{fppf}$ 具有相同的底层范畴，并令其覆盖
$\text{Cov}_\tau \subset \text{Cov}_{fppf}$ 恰为其中的 $\tau$-覆盖。
容易验证，这给出一个大位点 $\Sch_\tau$。
\end{remark}









\section{大位点的变换}
\label{topologies-section-change-alpha}

\noindent
本节说明变换大 Zariski/fppf/\'etale 位点时会发生什么。

\medskip\noindent
设 $\tau, \tau' \in \{Zariski, \etale, smooth, syntomic, fppf\}$。
给定两个大位点 $\Sch_\tau$ 和 $\Sch'_{\tau'}$，若
$\Ob(\Sch_\tau) \subset \Ob(\Sch'_{\tau'})$ 且
$\text{Cov}(\Sch_\tau) \subset \text{Cov}(\Sch'_{\tau'})$，则称
{\it $\Sch_\tau$ 包含于 $\Sch'_{\tau'}$}。在这种情形下，$\tau$
强于 $\tau'$；例如，任何 fppf 位点都不可能包含于一个 \'etale 位点。

\begin{lemma}
\label{topologies-lemma-contained-in}
任意一组大 Zariski 位点都包含于某个共同的大 Zariski 位点。
作必要变更后，同样的结论对大 fppf 位点和大 \'etale 位点也成立。
\end{lemma}

\begin{proof}
这是因为一组集合的并仍是集合，而且 Sets, Lemmas
\ref{sets-lemma-construct-category} and
\ref{sets-lemma-coverings-site} 中的构造允许从任意预先给定的概形和覆盖集合
出发。
\end{proof}

\begin{lemma}
\label{topologies-lemma-change-alpha}
设 $\tau \in \{Zariski, \etale, smooth, syntomic, fppf\}$。给定大位点
$\Sch_\tau$ 和 $\Sch'_\tau$，并假设 $\Sch_\tau$ 包含于
$\Sch'_\tau$。包含函子 $\Sch_\tau \to \Sch'_\tau$ 满足 Sites, Lemma
\ref{sites-lemma-bigger-site} 的假设。存在拓扑斯态射
\begin{eqnarray*}
g : \Sh(\Sch_\tau) &
\longrightarrow &
\Sh(\Sch'_\tau) \\
f : \Sh(\Sch'_\tau) &
\longrightarrow &
\Sh(\Sch_\tau)
\end{eqnarray*}
使得 $f \circ g \cong \text{id}$。对 $\Sch_\tau$ 的任意对象 $S$，
包含函子 $(\Sch/S)_\tau \to (\Sch'/S)_\tau$ 也满足 Sites, Lemma
\ref{sites-lemma-bigger-site} 的假设。因此类似地得到态射
\begin{eqnarray*}
g : \Sh((\Sch/S)_\tau) &
\longrightarrow &
\Sh((\Sch'/S)_\tau) \\
f : \Sh((\Sch'/S)_\tau) &
\longrightarrow &
\Sh((\Sch/S)_\tau)
\end{eqnarray*}
满足 $f \circ g \cong \text{id}$。
\end{lemma}

\begin{proof}
对函子 $\Sch_\tau \to \Sch'_\tau$ 和
$(\Sch/S)_\tau \to (\Sch'/S)_\tau$，Sites, Lemma
\ref{sites-lemma-bigger-site} 的假设 (b)、(c)、(e) 是立即成立的。
性质 (a) 由 Lemma \ref{topologies-lemma-zariski-induced}、
\ref{topologies-lemma-etale-induced}、\ref{topologies-lemma-smooth-induced}、
\ref{topologies-lemma-syntomic-induced} 或 \ref{topologies-lemma-fppf-induced} 成立。
性质 (d) 成立，是因为范畴 $\Sch_\tau$、$\Sch'_\tau$ 中的纤维积存在，
并且与概形范畴中的纤维积相容。
\end{proof}

\noindent
讨论：函子 $g^{-1} = f_*$ 就是限制函子，它把 $\Sch'_\tau$ 上的层
$\mathcal{G}$ 送到其限制 $\mathcal{G}|_{\Sch_\tau}$。因此，该引理就是说，
给定 $\Sch_\tau$ 上的任意集合层 $\mathcal{F}$，都存在 $\Sch'_\tau$
上的一个典范层 $\mathcal{F}'$，使得
% P05-ERR-0011: frozen restriction formula has the smaller-site sheaf restricted to the larger site.
$\mathcal{F}|_{\Sch'_\tau} = \mathcal{F}'$。事实上，层 $\mathcal{F}'$
有如下描述：它是下列预层的层化：
$$
\Sch'_\tau \longrightarrow \textit{Sets}, \quad
V \longmapsto \colim_{V \to U} \mathcal{F}(U)
$$
其中 $U$ 是 $\Sch_\tau$ 的对象。确实，由 Sites, Lemmas
\ref{sites-lemma-when-shriek} and \ref{sites-lemma-bigger-site}，有
$\mathcal{F}' = f^{-1}\mathcal{F} = (u_p\mathcal{F})^\#$。






\section{函子的延拓}
\label{topologies-section-extending-functors}

\noindent
先从一个说明我们所做之事的简单例子开始。设 $R$ 为环，并设 $F$ 是定义在
如下形状的 $R$-代数所成范畴 $\mathcal{C}$ 上的函子：
$$
A = R[x_1, \ldots, x_n]/(f_1, \ldots, f_m)
$$
其中 $n, m \geq 0$ 为整数，而
$f_1, \ldots, f_m \in R[x_1, \ldots, x_n]$ 为元素。于是对任意
$R$-代数 $B$，可定义
$$
F'(B) = \colim_{A \to B,\ A \in \mathcal{C}} F(A)
$$
事实表明，$F'$ 是所有 $R$-代数所成范畴上唯一既延拓 $F$ 又与滤过余极限
交换的函子。若要求所给函子是 Zariski 层，同样的步骤也适用于概形范畴。

\begin{lemma}
\label{topologies-lemma-extend}
设 $S$ 为概形，$\mathcal{C}$ 为 $S$ 上所有概形所成范畴 $\Sch/S$ 的充满
子范畴。假设：
\begin{enumerate}
\item 若 $X \to S$ 是 $\mathcal{C}$ 的对象，且 $U \subset X$ 为仿射开集，
则 $U \to S$ 同构于 $\mathcal{C}$ 的一个对象；
\item 若 $V$ 是位于仿射开集 $U \subset S$ 上的仿射概形，且
$V \to U$ 是有限表示的，则 $V \to S$ 同构于 $\mathcal{C}$ 的一个对象。
\end{enumerate}
设 $F : \mathcal{C}^{opp} \to \textit{Sets}$ 为函子。假设：
\begin{enumerate}
\item[(a)] 对任意 Zariski 覆盖 $\{f_i : X_i \to X\}_{i \in I}$，若
$X, X_i$ 均为 $\mathcal{C}$ 的对象，则 $F$ 对此态射族满足层条件
\footnote{由于我们不知道 $X_i \times_X X_j$ 是否属于 $\mathcal{C}$，
这必须解释如下：由性质 (1)，存在 Zariski 覆盖
$\{U_{ijk} \to X_i \times_X X_j\}_{k \in K_{ij}}$，其中 $U_{ijk}$
为 $\mathcal{C}$ 的对象。此时层条件是说，$F(X)$ 是从
$\prod F(X_i)$ 到 $\prod F(U_{ijk})$ 的两个映射的等化子。}；
\item[(b)] 若 $X = \lim X_i$ 是 $S$ 上仿射概形的有向极限，且
$X, X_i$ 均为 $\mathcal{C}$ 的对象，则 $F(X) = \colim F(X_i)$。
\end{enumerate}
则存在唯一方式把 $F$ 延拓为函子
$F' : (\Sch/S)^{opp} \to \textit{Sets}$，使其满足 (a)、(b) 的类似条件；
也就是说，$F'$ 对任意 Zariski 覆盖满足层条件，并且每当
$X = \lim X_i$ 是 $S$ 上仿射概形的有向极限时，都有
$F'(X) = \colim F'(X_i)$。
\end{lemma}

\begin{proof}
思路是先把 $F$ 延拓到 $S$ 上足够大的一类仿射概形，再利用 Zariski 层性质
延拓到所有概形。

\medskip\noindent
设 $V$ 是 $S$ 上的仿射概形，其结构态射 $V \to S$ 经由某个仿射开集
$U \subset S$ 分解。在这种情形下，可写成
$$
V = \lim V_i
$$
这一余滤过极限，其中 $V_i \to U$ 为有限表示的，且 $V_i$ 仿射。见
Algebra, Lemma \ref{algebra-lemma-ring-colimit-fp}。由条件 (1) 和 (2)，
可将各 $V_i$ 替换为 $\mathcal{C}$ 的对象。注意，$V_i \to S$ 是局部有限表示
的（若 $S$ 拟分离，则这些态射事实上是有限表示的）。于是令
$$
F'(V) = \colim F(V_i)
$$
事实上，可以给出一个更典范的表达式：
$$
F'(V) = \colim_{V \to V'} F(V')
$$
其中余极限取遍 $S$ 上的态射 $V \to V'$ 所成范畴；这里 $V'$ 是
$\mathcal{C}$ 的对象，且其结构态射 $V' \to S$ 局部有限表示。这个表达式
与第一个公式相同，是因为由 Limits, Proposition
\ref{limits-proposition-characterize-locally-finite-presentation}，
逆系统 $V_i$ 在此范畴中共尾！最后注意，若 $V$ 是 $\mathcal{C}$ 的对象，
则由假设 (b) 有 $F'(V) = F(V)$。

\medskip\noindent
第二个公式把 $F'$ 变成如下范畴上的反变函子：其对象为 $S$ 上的仿射概形
$V$，且其结构态射经由 $S$ 的某个仿射开集分解。设 $V$ 是 $S$ 上这样的
仿射概形，并设 $V = \bigcup_{k = 1, \ldots, n} V_k$ 是由仿射开集组成的
有限开覆盖。此时可以问 $F'$ 对这个开覆盖是否满足层条件。答案是肯定的，
而且容易证明：如前一段写成 $V = \lim V_i$。由 Limits, Lemma
\ref{limits-lemma-descend-opens}，对所有充分大的 $i$，可找到与转移映射相容的
仿射开集 $V_{i, k} \subset V_i$，它们拉回为 $V$ 中的 $V_k$。因此
$$
F'(V_k) = \colim F(V_{i, k})
\quad\text{且}\quad
F'(V_k \cap V_l) = \colim F(V_{i, k} \cap V_{i, l})
$$
严格地说，在这些公式中，应用函子 $F$ 之前，必须把 $V_{i, k}$ 和
$V_{i, k} \cap V_{i, l}$ 替换为与其同构的 $\mathcal{C}$ 中仿射对象。
由于 $I$ 有向，余极限可穿过等化子。因此，
% P05-ERR-0012: frozen proof cites sheaf condition (b), while the lemma assigns sheaf descent to (a).
$F$ 对 Zariski 覆盖 $\{V_{i, k} \to V_i\}$ 的层条件 (b) 蕴含 $F'$
对所给覆盖的层条件。

\medskip\noindent
设 $X$ 为 $S$ 上的一般概形。以 $\mathcal{B}_X$ 表示 $X$ 的如下仿射开集
所成的集合：其到 $S$ 的结构态射映入 $S$ 的某个仿射开集。显然，
$\mathcal{B}_X$ 是 $X$ 的拓扑基。由前一段的结果及 Sheaves, Lemma
\ref{sheaves-lemma-cofinal-systems-coverings-standard-case}，可知 $F'$
是 $\mathcal{B}_X$ 上的层。因此，$F'$ 在 $\mathcal{B}_X$ 上的限制唯一地
延拓为 $X$ 上的层 $F'_X$；见 Sheaves, Lemma
\ref{sheaves-lemma-extend-off-basis}。若 $X$ 是 $\mathcal{C}$ 的对象，
则由于 $F'$ 与 $F$ 在二者均有定义的对象上一致，并且 $F$ 对 Zariski
覆盖满足层条件，得到典范识别 $F'_X(X) = F(X)$。

\medskip\noindent
设 $f : X \to Y$ 为 $S$ 上的概形态射。得到从 $F'_Y$ 到 $F'_X$ 的唯一
$f$-映射，它与所有映射 $F'(V) \to F'(U)$ 相容；这里
$U \in \mathcal{B}_X$、$V \in \mathcal{B}_Y$ 且 $f(U) \subset V$。
见 Sheaves, Lemma
\ref{sheaves-lemma-f-map-basis-above-and-below-structures}。关于给定 $S$
上概形态射 $X \to Y \to Z$ 时这些映射正确复合的验证从略。我们也略去
如下验证：若 $f$ 是 $\mathcal{C}$ 的态射，则经由上述识别
$F'_X(X) = F(X)$ 和 $F'_Y(Y) = F(Y)$，诱导映射
$F'_Y(Y) \to F'_X(X)$ 与映射 $F(Y) \to F(X)$ 相同。由此可见，
$F$ 的所需延拓就是把 $X/S$ 送到 $F'_X(X)$ 的函子。

\medskip\noindent
函子 $X \mapsto F'_X(X)$ 的性质 (a) 几乎立即由构造得到；细节从略。
设 $X = \lim_{i \in I} X_i$ 是 $S$ 上仿射概形的有向极限。必须证明
$$
F'_X(X) = \colim_{i \in I} F'_{X_i}(X_i)
$$
先假设存在某个 $i \in I$，使 $X_i \to S$ 经由仿射开集
$U \subset S$ 分解。则 $F'$ 在 $X$ 以及 $i' \geq i$ 时的 $X_{i'}$ 上
均有定义，并且当 $i' \geq i$ 时有
$F'_{X_{i'}}(X_{i'}) = F'(X_{i'})$，同时 $F'_X(X) = F'(X)$。
在这种情形下，每个态射 $X \to V$，其中 $V$ 在 $S$ 上局部有限表示，
都对某个 $i' \geq i$ 分解为 $X \to X_{i'} \to V$；见 Limits,
Proposition \ref{limits-proposition-characterize-locally-finite-presentation}。
因此有
\begin{align*}
F'_X(X)
& =
F'(X)  \\
& = \colim_{X \to V} F(V) \\
& =
\colim_{i' \geq i} \colim_{X_{i'} \to V} F(V) \\
& =
\colim_{i' \geq i} F'(X_{i'}) \\
& =
\colim_{i' \geq i} F'_{X_{i'}}(X_{i'}) \\
& =
\colim_{i' \in I} F'_{X_{i'}}(X_{i'})
\end{align*}
这正是所需结论。最后，在一般情形下，任取 $i \in I$，并选取有限仿射开覆盖
% P05-ERR-0013: frozen proof switches from X_i/X to undefined V_i/V.
$V_i = V_{i, 1} \cup \ldots \cup V_{i, n}$，使 $V_{i, k} \to S$
经由 $S$ 的某个仿射开集分解。令 $V_k \subset V$ 以及 $i' \geq i$ 时的
$V_{i', k}$ 为 $V_{i, k}$ 的逆像。由前一种情形可知
$$
F'_{V_k}(V_k) = \colim_{i' \geq i} F'_{V_{i', k}}(V_{i', k})
$$
and
$$
F'_{V_k \cap V_l}(V_k \cap V_l) =
\colim_{i' \geq i}
F'_{V_{i', k} \cap V_{i', l}}(V_{i', k} \cap V_{i', l})
$$
由层性质和滤过余极限的正合性，在这种情形下也得到
$F'_X(X) = \colim_{i \in I} F'_{X_i}(X_i)$。这完成了性质 (b) 的证明，
从而也完成了引理的证明。
\end{proof}

\begin{lemma}
\label{topologies-lemma-limit-fppf-topology}
设 $\tau \in \{Zariski, \etale, smooth, syntomic, fppf\}$。设 $T$ 为仿射概形，
它写成仿射概形的有向逆系统的极限 $T = \lim_{i \in I} T_i$。
\begin{enumerate}
\item 设 $\mathcal{V} = \{V_j \to T\}_{j = 1, \ldots, m}$ 为 $T$ 的标准
$\tau$-覆盖；见 Definitions
\ref{topologies-definition-standard-Zariski},
\ref{topologies-definition-standard-etale},
\ref{topologies-definition-standard-smooth},
\ref{topologies-definition-standard-syntomic}, and
\ref{topologies-definition-standard-fppf}。则存在指标 $i$ 以及标准 $\tau$-覆盖
$\mathcal{V}_i = \{V_{i, j} \to T_i\}_{j = 1, \ldots, m}$，使其到 $T$
的基变换 $T \times_{T_i} \mathcal{V}_i$ 同构于 $\mathcal{V}$。
\item 设 $\mathcal{V}_i$、$\mathcal{V}'_i$ 为 $T_i$ 的一对标准
$\tau$-覆盖。若
$f : T \times_{T_i} \mathcal{V}_i \to T \times_{T_i} \mathcal{V}'_i$
是 $T$ 的覆盖之间的态射，则存在指标 $i' \geq i$ 以及态射
% P05-ERR-0014: frozen domain drops the index i from the covering family.
$f_{i'} : T_{i'} \times_{T_i} \mathcal{V} \to
T_{i'} \times_{T_i} \mathcal{V}'_i$
其到 $T$ 的基变换为 $f$。
\item 若
% P05-ERR-0015: frozen domain again uses unindexed mathcal V although both coverings are over T_i.
$f, g : \mathcal{V} \to \mathcal{V}'_i$
是 $T_i$ 的标准 $\tau$-覆盖之间的态射，并且它们到 $T$ 的基变换
$f_T, g_T$ 相等，则存在指标 $i' \geq i$，使得
$f_{T_{i'}} = g_{T_{i'}}$。
\end{enumerate}
换言之，$T$ 的标准 $\tau$-覆盖所成范畴，是 $T_i$ 的标准 $\tau$-覆盖所成
范畴在 $I$ 上的余极限。
\end{lemma}

\begin{proof}
先对 $\tau = fppf$ 证明。由 Limits, Lemma
\ref{limits-lemma-descend-finite-presentation}，$T$ 上有限表示概形所成范畴，
是 $T_i$ 上有限表示范畴在 $I$ 上的余极限。由 Limits, Lemmas
\ref{limits-lemma-descend-affine-finite-presentation} and
\ref{limits-lemma-descend-flat-finite-presentation}，对 $T$ 上既仿射、平坦又
有限表示的概形所成范畴，同样的结论成立。为完成引理的证明，只需证明：若
$\{V_{j, i} \to T_i\}_{j = 1, \ldots, m}$ 是有限族平坦有限表示态射，
其中 $V_{j, i}$ 仿射，且基变换
$\coprod_j T \times_{T_i} V_{j, i} \to T$ 满射，则对某个 $i' \geq i$，态射
% P05-ERR-0016: frozen coproduct omits the j index used in the preceding base change.
$\coprod T_{i'} \times_{T_i} V_{j, i} \to T_{i'}$ 是满射。
是满射。分别以 $W_{i'} \subset T_{i'}$ 和 $W \subset T$ 表示其像。
由假设当然有 $W = T$。由于这些态射平坦且有限表示，可见 $W_i$ 是
$T_i$ 的拟紧开集；见 Morphisms, Lemma \ref{morphisms-lemma-fppf-open}。
此外，$W = T \times_{T_i} W_i$（取像与基变换交换）。因此，由 Limits,
Lemma \ref{limits-lemma-descend-opens}，对某个充分大的 $i'$ 有
$W_{i'} = T_{i'}$，结论得证。

\medskip\noindent
当 $\tau \in \{Zariski, \etale, smooth, syntomic\}$ 时，标准 $\tau$-覆盖
是标准 fppf 覆盖。因此，函子的全忠实性成立。唯一的问题是证明：给定某个
$i$ 上的标准 fppf 覆盖 $\mathcal{V}_i$，若
$\mathcal{V}_i \times_{T_i} T$ 是标准 $\tau$-覆盖，则对所有
$i' \gg i$，$\mathcal{V}_i \times_{T_i} T_{i'}$ 都是标准
$\tau$-覆盖。这立即由 Limits, Lemmas
\ref{limits-lemma-descend-open-immersion},
\ref{limits-lemma-descend-etale},
\ref{limits-lemma-descend-smooth}, and
\ref{limits-lemma-descend-syntomic}.
\end{proof}

\begin{lemma}
\label{topologies-lemma-extend-sheaf-general}
设 $S$、$\mathcal{C}$、$F$ 满足 Lemma \ref{topologies-lemma-extend} 的条件 (1)、(2)、
(a)、(b)，并以 $F' : (\Sch/S)^{opp} \to \textit{Sets}$ 表示该引理中构造的
唯一延拓。设\linebreak
$\tau \in \{Zariski, \etale, smooth, syntomic, fppf\}$。
假设：
\begin{enumerate}
\item[(c)] 对 $\Sch/S$ 中仿射概形的任意标准 $\tau$-覆盖
$\{V_i \to V\}_{i = 1, \ldots, n}$，若 $V \to S$ 经由仿射开集
$U \subset S$ 分解，且 $V \to U$ 是有限表示的，则 $F$ 对
$\{V_i \to V\}_{i = 1, \ldots, n}$ 满足层条件\footnote{这是有意义的，
因为由 (2)，$V$、$V_i$ 及 $V_i \times_V V_j$ 都同构于 $\mathcal{C}$
的对象。}。
\end{enumerate}
则 $F'$ 对所有 $\tau$-覆盖满足层条件。
\end{lemma}

\begin{proof}
设 $X$ 为 $S$ 上的概形，$\{X_i \to X\}_{i \in I}$ 为 $\tau$-覆盖。
设 $s_i \in F'(X_i)$ 为一族元素，使得对所有 $i, j \in I$，$s_i$ 与
$s_j$ 映到 $F'(X_i \times_X X_j)$ 中的同一元素。必须证明存在唯一元素
$s \in F'(X)$，使其对所有 $i \in I$ 的限制为 $s_i \in F'(X_i)$。

\medskip\noindent
特殊情形：$X$ 仿射，其结构态射映入 $S$ 的仿射开集 $U$，并且覆盖
$\{X_i \to X\}_{i \in I}$ 是标准 $\tau$-覆盖。在这种情形下，可写成
$$
X = \lim V_k
$$
这一余滤过极限，其中 $V_k \to U$ 为有限表示的，且 $V_k$ 仿射。见
Algebra, Lemma \ref{algebra-lemma-ring-colimit-fp}。由 Lemma
\ref{topologies-lemma-limit-fppf-topology}，存在一个 $k$ 以及标准 $\tau$-覆盖
$\{V_{k, i} \to V_k\}_{i \in I}$，其到 $X$ 的基变换就是所给覆盖。
当 $k' \geq k$ 时，以 $\{V_{k', i} \to V_{k'}\}_{i \in I}$ 表示该覆盖到
$V_{k'}$ 的基变换。于是可见
% P05-ERR-0017: frozen first colimit term keeps V_k although the index is k' >= k.
\begin{align*}
F'(X)
& =
\colim_{k' \geq k} F(V_k) \\
& =
\colim_{k' \geq k}
\text{Equalizer}(
\xymatrix{
\prod F(V_{k', i})
\ar@<1ex>[r] \ar@<-1ex>[r] &
\prod F(V_{k', i} \times_{V_{k'}} V_{k', j})
}
\\
& =
\text{Equalizer}(
\xymatrix{
\colim_{k' \geq k}
\prod F(V_{k', i})
\ar@<1ex>[r] \ar@<-1ex>[r] &
\colim_{k' \geq k}
\prod F(V_{k', i} \times_{V_{k'}} V_{k', j})
}
\\
& =
\text{Equalizer}(
\xymatrix{
\prod F'(X_i)
\ar@<1ex>[r] \ar@<-1ex>[r] &
\prod F'(X_i \times_X X_j)
}
\end{align*}
第一个等式由 $F'$ 的构造成立。第二个等式由假设 (c) 成立。第三个等式成立，
是因为滤过余极限正合。第四个等式再次由 $F'$ 的构造成立。由此可见，
$F'$ 对 $\{X_i \to X\}_{i \in I}$ 满足层性质。

\medskip\noindent
一般情形。选取仿射开覆盖 $X = \bigcup U_k$，使每个 $U_k$ 都映入 $S$
的某个仿射开集。对每个 $k$，可选取加细
$\{X_i \times_X U_k \to U_k\}_{i \in I}$ 的标准 $\tau$-覆盖
$\{V_{k, j} \to U_k\}_{j = 1, \ldots, m_k}$。对每个
$j \in \{1, \ldots, m_k\}$，选取指标 $i_{k, j} \in I$ 以及 $X$ 上的态射
$g_{k, j} : V_{k, j} \to X_{i_{k, j}}$。令 $s_{k, j}$ 为通过
$g_{k, j}$ 限制 $s_{i_{k, j}}$ 所得的 $F'(V_{k, j})$ 中元素。注意，
对所有 $k$、$k'$ 以及所有 $j \in \{1, \ldots, m_k\}$、
$j' \in \{1, \ldots, m_{k'}\}$，$s_{k, j}$ 与 $s_{k', j'}$ 在
$F'(V_{k, j} \times_X V_{k', j'})$ 中的限制相同；验证从略。特别地，
由前一段的结果，存在唯一元素 $s_k \in F'(U_k)$，其对所有 $j$ 的限制
均为 $s_{k, j}$。采用这一记号，即可完成证明。

\medskip\noindent
$s$ 唯一性的证明：由于 $F'$ 对 Zariski 覆盖满足层性质，且
$s|_{U_k}$ 必须等于 $s_k$，因为二者对所有 $j$ 的限制都是 $s_{k, j}$，
故唯一性成立。该唯一性又表明，$s_k$ 与 $s_{k'}$ 在（非仿射概形）
$U_k \cap U_{k'}$ 上必限制为 $F'$ 的同一截面，因为这些截面在
$\tau$-覆盖
$\{V_{k, j} \times_X V_{k', j'} \to U_k \cap U_{k'}\}$ 上限制为同一截面。
因此，由对 Zariski 覆盖的层性质，存在唯一的 $F'$ 在 $X$ 上的截面 $s$，
其在 $U_k$ 上的限制为 $s_k$。关于 $s$ 在 $X_i$ 上限制为 $s_i$ 的验证
（与上述类似）从略。
\end{proof}

\begin{lemma}
\label{topologies-lemma-extend-sheaf}
设 $\tau \in \{Zariski, \etale, smooth, syntomic, fppf\}$。设 $S$ 为包含在
大位点 $\Sch_\tau$ 中的概形。设
$F : (\Sch/S)_\tau^{opp} \to \textit{Sets}$ 为一个 $\tau$-层，并且在
$\mathcal{C} = (\Sch/S)_\tau$ 时满足 Lemma \ref{topologies-lemma-extend} 的性质 (b)。
则 $F$ 到 $S$ 上所有概形所成范畴的延拓 $F'$ 对所有 $\tau$-覆盖满足层条件。
\end{lemma}

\begin{proof}
这由取 $\mathcal{C} = (\Sch/S)_\tau$ 应用 Lemma
\ref{topologies-lemma-extend-sheaf-general} 得出。Lemma \ref{topologies-lemma-extend} 的条件
(1)、(2)、(a)、(b) 成立；细节从略。因此得到到 $S$ 上所有概形所成范畴的
唯一延拓 $F'$。最后注意，任意标准 $\tau$-覆盖都同义等价于
$(\Sch/S)_\tau$ 中的某个覆盖；见 Sets, Lemma
\ref{sets-lemma-what-is-in-it} 以及 Lemmas \ref{topologies-lemma-zariski-induced}、
\ref{topologies-lemma-etale-induced},
\ref{topologies-lemma-smooth-induced},
\ref{topologies-lemma-syntomic-induced}, and
\ref{topologies-lemma-fppf-induced}。由 Sites, Lemma
\ref{sites-lemma-tautological-same-sheaf}，层性质可沿覆盖的同义等价传递。
因此，$F$ 是 $\tau$-层这一事实蕴含 Lemma
\ref{topologies-lemma-extend-sheaf-general} 的性质 (c) 成立，结论得证。
\end{proof}
% Generated mechanically from the frozen descent.tex target.
% Do not edit this generated file; edit the bound locale source instead.

% OK, start here.
%

\setcounter{chapter}{34}
\chapter{下降}


\phantomsection
\label{descent-section-phantom}


\section{引言}
\label{descent-section-introduction}

\noindent
在概形上的拓扑一章中（见 Topologies, Section
\ref{topologies-section-introduction}），我们引入了概形的 Zariski、
\'etale、fppf、光滑、syntomic 和 fpqc 覆盖。本章讨论何种概形上结构
可以沿这类覆盖下降。例如可参阅 \cite{Gr-I}、\cite{Gr-II}、
\cite{Gr-III}、\cite{Gr-IV}、\cite{Gr-V} 和 \cite{Gr-VI}。
本章还旨在拟凝聚层、概形等对象的下降问题这一较少形式化的背景中，
引入下降、下降数据和有效下降数据等概念。形式化的概念，即位点上的叠，
将在叠一章中讨论（见 Stacks, Section
\ref{stacks-section-introduction}）。

\section{拟凝聚层的下降数据}
\label{descent-section-equivalence}

\noindent
本章采用如下约定：投影态射
$\text{pr}_i : X \times \ldots \times X \to X$ 的编号从 $i = 0$ 开始。
因此有 $\text{pr}_0, \text{pr}_1 : X \times X  \to X$、
$\text{pr}_0, \text{pr}_1, \text{pr}_2 : X \times X \times X  \to X$
等等。


\begin{definition}
\label{descent-definition-descent-datum-quasi-coherent}
设 $S$ 为概形，并设 $\{f_i : S_i \to S\}_{i \in I}$ 是以 $S$ 为靶的
态射族。
\begin{enumerate}
\item 相对于给定族的一个{\it 拟凝聚层下降数据
$(\mathcal{F}_i, \varphi_{ij})$}由如下资料组成：对每个 $i \in I$，
$S_i$ 上的一个拟凝聚层 $\mathcal{F}_i$；以及对每一对
$(i, j) \in I^2$，拟凝聚 $\mathcal{O}_{S_i \times_S S_j}$-模
$\varphi_{ij} : \text{pr}_0^*\mathcal{F}_i \to \text{pr}_1^*\mathcal{F}_j$
的一个同构，并要求对每个指标三元组 $(i, j, k) \in I^3$，图
$$
\xymatrix{
\text{pr}_0^*\mathcal{F}_i \ar[rd]_{\text{pr}_{01}^*\varphi_{ij}}
\ar[rr]_{\text{pr}_{02}^*\varphi_{ik}} & &
\text{pr}_2^*\mathcal{F}_k \\
& \text{pr}_1^*\mathcal{F}_j \ar[ru]_{\text{pr}_{12}^*\varphi_{jk}} &
}
$$
作为 $\mathcal{O}_{S_i \times_S S_j \times_S S_k}$-模的图可交换。
这称为{\it 余圈条件}。
\item 一个{\it 下降数据态射
$\psi : (\mathcal{F}_i, \varphi_{ij}) \to
(\mathcal{F}'_i, \varphi'_{ij})$}由一族 $\mathcal{O}_{S_i}$-模态射
$\psi_i : \mathcal{F}_i \to \mathcal{F}'_i$，即
$\psi = (\psi_i)_{i\in I}$，给出，并要求所有图
$$
\xymatrix{
\text{pr}_0^*\mathcal{F}_i \ar[r]_{\varphi_{ij}} \ar[d]_{\text{pr}_0^*\psi_i}
& \text{pr}_1^*\mathcal{F}_j \ar[d]^{\text{pr}_1^*\psi_j} \\
\text{pr}_0^*\mathcal{F}'_i \ar[r]^{\varphi'_{ij}} &
\text{pr}_1^*\mathcal{F}'_j \\
}
$$
均可交换。
\end{enumerate}
\end{definition}

\noindent
应当牢记的一个典型例子如下。设 $S = \bigcup S_i$ 是开覆盖。
在这种情形下，我们已经在 Sheaves, Section
\ref{sheaves-section-glueing-sheaves} 中见过集合层的下降数据，
在那里称之为“相对于给定覆盖的集合层粘合数据”。此外，我们证明了
粘合数据范畴等价于 $S$ 上的层范畴。我们将证明：当
$\{S_i \to S\}_{i\in I}$ 是 fpqc 覆盖时，上述背景中也有相应结论。

\medskip\noindent
在极端情形中，覆盖 $\{S \to S\}$ 由 $\text{id}_S$ 给出，此时下降数据
必为 $(\mathcal{F}, \text{id}_\mathcal{F})$ 的形式。余圈条件保证，
在这种情形下，$\mathcal{F}$ 上唯一允许的态射就是恒等态射。特别地，
下述引理表明：对每个拟凝聚 $\mathcal{O}_S$-模层和任意给定的态射族，
都相应地存在唯一的下降数据。

\begin{lemma}
\label{descent-lemma-refine-descent-datum}
设 $\mathcal{U} = \{U_i \to U\}_{i \in I}$ 和
$\mathcal{V} = \{V_j \to V\}_{j \in J}$ 是具有固定靶的概形态射族。
设 $(g, \alpha : I \to J, (g_i)) : \mathcal{U} \to \mathcal{V}$
是具有固定靶的态射族之间的态射；见 Sites, Definition
\ref{sites-definition-morphism-coverings}。设
$(\mathcal{F}_j, \varphi_{jj'})$ 是拟凝聚层相对于态射族
$\{V_j \to V\}_{j \in J}$ 的一个下降数据。则
\begin{enumerate}
\item 系统
$$
\left(g_i^*\mathcal{F}_{\alpha(i)},
(g_i \times g_{i'})^*\varphi_{\alpha(i)\alpha(i')}\right)
$$
是相对于态射族 $\{U_i \to U\}_{i \in I}$ 的下降数据。
\item 此构造关于下降数据 $(\mathcal{F}_j, \varphi_{jj'})$ 具有函子性。
\item 给定具有固定靶的态射族之间的第二个态射
$(g', \alpha' : I \to J, (g'_i))$，且 $g = g'$，则存在下降数据之间的
函子性同构
$$
(g_i^*\mathcal{F}_{\alpha(i)},
(g_i \times g_{i'})^*\varphi_{\alpha(i)\alpha(i')})
\cong
((g'_i)^*\mathcal{F}_{\alpha'(i)},
(g'_i \times g'_{i'})^*\varphi_{\alpha'(i)\alpha'(i')}).
$$
\end{enumerate}
\end{lemma}

\begin{proof}
略。提示：给出第 (3) 项中下降数据同构的态射
$g_i^*\mathcal{F}_{\alpha(i)} \to (g'_i)^*\mathcal{F}_{\alpha'(i)}$
是态射 $\varphi_{\alpha(i)\alpha'(i)}$ 沿
$(g_i, g'_i) : U_i \to V_{\alpha(i)} \times_V V_{\alpha'(i)}$ 的拉回。
\end{proof}

\noindent
任意态射族 $\mathcal{U} = \{S_i \to S\}_{i \in I}$ 都以唯一方式成为
平凡覆盖 $\{S \to S\}$ 的加细。若 $\mathcal{F}$ 是 $S$ 上的拟凝聚层，
则将上述步骤所得的相对于 $\mathcal{U}$ 的下降数据简记为
$(\mathcal{F}|_{S_i}, can)$。

\begin{definition}
\label{descent-definition-descent-datum-effective-quasi-coherent}
设 $S$ 为概形，并设 $\{S_i \to S\}_{i \in I}$ 是以 $S$ 为靶的态射族。
\begin{enumerate}
\item 设 $\mathcal{F}$ 为拟凝聚 $\mathcal{O}_S$-模。称 $\mathcal{F}$
相对于覆盖 $\{S \to S\}$ 的唯一下降数据为{\it 平凡下降数据}。
\item 平凡下降数据在 $\{S_i \to S\}$ 上的拉回称为{\it 典范下降数据}，
记作 $(\mathcal{F}|_{S_i}, can)$。
\item 若相对于给定覆盖的拟凝聚层下降数据
$(\mathcal{F}_i, \varphi_{ij})$ 满足：存在 $S$ 上的拟凝聚层
$\mathcal{F}$，使得 $(\mathcal{F}_i, \varphi_{ij})$ 同构于
$(\mathcal{F}|_{S_i}, can)$，则称它是{\it 有效的}。
\end{enumerate}
\end{definition}

\begin{lemma}
\label{descent-lemma-zariski-descent-effective}
设 $S$ 为概形，并设 $S = \bigcup U_i$ 是开覆盖。拟凝聚层相对于态射族
$\mathcal{U} = \{U_i \to S\}$ 的任意下降数据都是有效的。此外，从拟凝聚
$\mathcal{O}_S$-模范畴到相对于 $\mathcal{U}$ 的下降数据范畴的函子是全忠实的。
\end{lemma}

\begin{proof}
这立即由 Sheaves, Section \ref{sheaves-section-glueing-sheaves}
以及拟凝聚性是局部性质这一事实得出；见 Modules, Definition
\ref{modules-definition-quasi-coherent}。
\end{proof}

\noindent
为进一步推进，我们首先需要研究环上模的情形。










\section{模的下降}
\label{descent-section-descent-modules}

\noindent
设 $R \to A$ 为环同态。由 Simplicial, Example
\ref{simplicial-example-push-outs-simplicial-object}，它给出余单纯 $R$-代数
$$
\xymatrix{
A
\ar@<1ex>[r]
\ar@<-1ex>[r]
&
A \otimes_R A
\ar@<0ex>[l]
\ar@<2ex>[r]
\ar@<0ex>[r]
\ar@<-2ex>[r]
&
A \otimes_R A \otimes_R A
\ar@<1ex>[l]
\ar@<-1ex>[l]
}
$$
将其记为 $(A/R)_\bullet$，于是 $(A/R)_n$ 是 $A$ 在 $R$ 上的
$(n + 1)$ 重张量积。给定映射 $\varphi : [n] \to [m]$，$R$-代数同态
$(A/R)_\bullet(\varphi)$ 为
$$
a_0 \otimes \ldots \otimes a_n
\longmapsto
\prod\nolimits_{\varphi(i) = 0} a_i
\otimes
\prod\nolimits_{\varphi(i) = 1} a_i
\otimes \ldots \otimes
\prod\nolimits_{\varphi(i) = m} a_i
$$
其中约定空乘积为 $1$。因此，采用 Simplicial, Section
\ref{simplicial-section-cosimplicial-object} 中的记号，最初几个态射为
$$
\begin{matrix}
\delta^1_0 & : & a_0 & \mapsto & 1 \otimes a_0 \\
\delta^1_1 & : & a_0 & \mapsto & a_0 \otimes 1 \\
\sigma^0_0 & : & a_0 \otimes a_1 & \mapsto & a_0a_1 \\
\delta^2_0 & : & a_0 \otimes a_1 & \mapsto & 1 \otimes a_0 \otimes a_1 \\
\delta^2_1 & : & a_0 \otimes a_1 & \mapsto & a_0 \otimes 1 \otimes a_1 \\
\delta^2_2 & : & a_0 \otimes a_1 & \mapsto & a_0 \otimes a_1 \otimes 1 \\
\sigma^1_0 & : & a_0 \otimes a_1 \otimes a_2 & \mapsto & a_0a_1 \otimes a_2 \\
\sigma^1_1 & : & a_0 \otimes a_1 \otimes a_2 & \mapsto & a_0 \otimes a_1a_2
\end{matrix}
$$
等等。

\medskip\noindent
一个 $R$-模 $M$ 给出余单纯 $(A/R)_\bullet$-模
$(A/R)_\bullet \otimes_R M$。换言之，$M_n = (A/R)_n \otimes_R M$，
并利用 $R$-代数同态 $(A/R)_n \to (A/R)_m$ 定义
$M \otimes_R (A/R)_\bullet$ 上相应的态射。

\medskip\noindent
拟凝聚层下降数据在模的背景中的对应概念如下。

\begin{definition}
\label{descent-definition-descent-datum-modules}
设 $R \to A$ 为环同态。
\begin{enumerate}
\item 一个{\it 模相对于 $R \to A$ 的下降数据 $(N, \varphi)$}
由一个 $A$-模 $N$ 和一个 $A \otimes_R A$-模同构
$$
\varphi : N \otimes_R A \to A \otimes_R N
$$
给出，并要求满足{\it 余圈条件}：$A \otimes_R A \otimes_R A$-模态射图
$$
\xymatrix{
N \otimes_R A \otimes_R A \ar[rr]_{\varphi_{02}}
\ar[rd]_{\varphi_{01}}
& &
A \otimes_R A \otimes_R N \\
& A \otimes_R N \otimes_R A \ar[ru]_{\varphi_{12}} &
}
$$
可交换（记号见下文）。
\item 一个{\it 下降数据态射 $(N, \varphi) \to (N', \varphi')$}
是一个 $A$-模态射 $\psi : N \to N'$，并要求图
$$
\xymatrix{
N \otimes_R A \ar[r]_\varphi \ar[d]_{\psi \otimes \text{id}_A} &
A \otimes_R N \ar[d]^{\text{id}_A \otimes \psi} \\
N' \otimes_R A \ar[r]^{\varphi'} &
A \otimes_R N'
}
$$
可交换。
\end{enumerate}
\end{definition}

\noindent
在此定义中采用如下记号：$\varphi_{01} = \varphi \otimes \text{id}_A$，
$\varphi_{12} = \text{id}_A \otimes \varphi$；若
$\varphi(n \otimes 1) = \sum a_i \otimes n_i$，则
$\varphi_{02}(n \otimes 1 \otimes 1) = \sum a_i \otimes 1 \otimes n_i$。
三者都是 $A \otimes_R A \otimes_R A$-模同态。等价地，
$$
\varphi_{ij}
=
\varphi \otimes_{(A/R)_1, \ (A/R)_\bullet(\tau^2_{ij})} (A/R)_2
$$
其中 $\tau^2_{ij} : [1] \to [2]$ 是由 $0 \mapsto i$、$1 \mapsto j$
给出的映射。具体而言，
$(A/R)_{\bullet}(\tau^2_{02})(a_0 \otimes a_1) =
a_0 \otimes 1 \otimes a_1$,
其余情形类似\footnote{注意，若
$\{i, j, k\} = [2] = \{0, 1, 2\}$，则
$\tau^2_{ij} = \delta^2_k$；见 Simplicial, Definition
\ref{simplicial-definition-face-degeneracy}。}。

\medskip\noindent
为陈述下一个引理，还需要一些记号。设 $(N, \varphi)$ 是相对于环同态
$R \to A$ 的下降数据。对 $n \geq 0$ 和 $i \in [n]$，令
$$
N_{n, i} =
A \otimes_R
\ldots
\otimes_R A \otimes_R N \otimes_R A \otimes_R
\ldots
\otimes_R A
$$
其中因子 $N$ 位于第 $i$ 个位置。这是一个 $(A/R)_n$-模。引入映射
$\tau^n_i : [0] \to [n]$，$0 \mapsto i$，便可看出
$$
N_{n, i} = N \otimes_{(A/R)_0, \ (A/R)_\bullet(\tau^n_i)} (A/R)_n
$$
对 $0 \leq i \leq j \leq n$，令 $\tau^n_{ij} : [1] \to [n]$ 为把
$0$ 映至 $i$、把 $1$ 映至 $j$ 的映射。与上文类似，同态 $\varphi$
诱导同构
$$
\varphi^n_{ij}
=
\varphi \otimes_{(A/R)_1, \ (A/R)_\bullet(\tau^n_{ij})} (A/R)_n :
N_{n, i} \longrightarrow N_{n, j}
$$
（当 $i < j$ 时为 $(A/R)_n$-模同构）。若 $i = j$，则令
$\varphi^n_{ij} = \text{id}$。由于这些都是同构，我们可以把因子 $N$
移至任意所需位置。余圈条件恰好说明，移动方式并不影响结果（例如，
复合其中两个同构或一次完成移动所得结果相同）。最后，对任意
$\beta : [n] \to [m]$，将态射
$$
N_{\beta, i} : N_{n, i} \to N_{m, \beta(i)}
$$
定义为唯一满足下式的 $(A/R)_\bullet(\beta)$-半线性映射：
$$
N_{\beta, i}(1 \otimes \ldots \otimes n \otimes \ldots \otimes 1)
=
1 \otimes \ldots \otimes n \otimes \ldots \otimes 1
$$
对所有 $n \in N$ 均成立。这提示了下述引理。

\begin{lemma}
\label{descent-lemma-descent-datum-cosimplicial}
设 $R \to A$ 为环同态。给定下降数据 $(N, \varphi)$，可以按如下规则
为其配备一个余单纯 $(A/R)_\bullet$-模
$N_\bullet$\footnote{严格说来，应写成 $(N, \varphi)_\bullet$。}：
令 $N_n = N_{n, n}$，并对给定的 $\beta : [n] \to [m]$ 定义
$$
N_\bullet(\beta) = (\varphi^m_{\beta(n)m}) \circ N_{\beta, n} :
N_{n, n} \longrightarrow N_{m, m}.
$$
此步骤关于下降数据具有函子性。
\end{lemma}

\begin{proof}
若 $\varphi(n \otimes 1) = \sum \alpha_i \otimes x_i$，则最初几个态射为
$$
\begin{matrix}
\delta^1_0 & : & N & \to & A \otimes N & n & \mapsto & 1 \otimes n \\
\delta^1_1 & : & N & \to & A \otimes N & n & \mapsto &
\sum \alpha_i \otimes x_i\\
\sigma^0_0 & : & A \otimes N & \to & N & a_0 \otimes n & \mapsto & a_0n \\
\delta^2_0 & : & A \otimes N & \to & A \otimes A \otimes N &
a_0 \otimes n & \mapsto & 1 \otimes a_0 \otimes n \\
\delta^2_1 & : & A \otimes N & \to & A \otimes A \otimes N &
a_0 \otimes n & \mapsto & a_0 \otimes 1 \otimes n \\
\delta^2_2 & : & A \otimes N & \to & A \otimes A \otimes N &
a_0 \otimes n & \mapsto & \sum a_0 \otimes \alpha_i \otimes x_i \\
\sigma^1_0 & : & A \otimes A \otimes N & \to & A \otimes N &
a_0 \otimes a_1 \otimes n & \mapsto & a_0a_1 \otimes n \\
\sigma^1_1 & : & A \otimes A \otimes N & \to & A \otimes N &
a_0 \otimes a_1 \otimes n & \mapsto & a_0 \otimes a_1n
\end{matrix}
$$
其中采用 Simplicial, Section
\ref{simplicial-section-cosimplicial-object} 的记号。我们先验证两个关系
$\sigma^0_0 \circ \delta^1_0 = \text{id}$ 和
$\sigma^0_0 \circ \delta^1_1 = \text{id}$。第一个关系
$\sigma^0_0 \circ \delta^1_0 = \text{id}$ 由上述态射的显式描述立即可见。
证明第二个关系时必须使用余圈条件（因为任意同构
$\varphi : N \otimes_R A \to A \otimes_R N$ 未必满足该关系）。记
$p = \sigma^0_0 \circ \delta^1_1 : N \to N$。由上述态射的描述可知，
$p$ 也等于
$$
p = \varphi \otimes \text{id} :
N = (N \otimes_R A) \otimes_{(A \otimes_R A)} A
\longrightarrow
(A \otimes_R N) \otimes_{(A \otimes_R A)} A = N
$$
由于 $\varphi$ 是同构，$p$ 也是同构。对某些 $\alpha_i \in A$ 和
$x_i \in N$，写作
$\varphi(n \otimes 1) = \sum \alpha_i \otimes x_i$。于是
$p(n) = \sum \alpha_ix_i$。再对某些 $\alpha_{ij} \in A$ 和
$y_j \in N$，写作
$\varphi(x_i \otimes 1) = \sum \alpha_{ij} \otimes y_j$。余圈条件给出
$$
\sum \alpha_i \otimes \alpha_{ij} \otimes y_j
=
\sum \alpha_i \otimes 1 \otimes x_i.
$$
这意味着 $p(n) = \sum \alpha_ix_i = \sum \alpha_i\alpha_{ij}y_j =
\sum \alpha_i p(x_i) = p(p(n))$。因此 $p$ 是幂等算子；又因它是同构，
故它必为恒等态射。

\medskip\noindent
要完整证明 $N_\bullet$ 是余单纯模，需要验证 Simplicial, Remark
\ref{simplicial-remark-relations-cosimplicial} 中的全部五类关系。
$\sigma$ 之间的复合关系是显然的；$\delta$ 之间的复合关系归结为
$\varphi$ 的余圈条件。完全仿照上文即可验证关系
$\sigma_j \circ \delta_j = \sigma_j \circ \delta_{j + 1} = \text{id}$。
最后，$\delta$ 与 $\sigma$ 的其余复合关系对任意 $\varphi$ 都成立。
\end{proof}

\noindent
注意，可以为 $R$-模 $M$ 配备一个典范下降数据，即
$(M \otimes_R A, can)$，其中
$can : (M \otimes_R A) \otimes_R A \to A \otimes_R (M \otimes_R A)$
是显然的态射：
$(m \otimes a) \otimes a' \mapsto a \otimes (m \otimes a')$.

\begin{lemma}
\label{descent-lemma-canonical-descent-datum-cosimplicial}
设 $R \to A$ 为环同态，$M$ 为 $R$-模。与该典范下降数据相联系的余单纯
$(A/R)_\bullet$-模同构于余单纯模 $(A/R)_\bullet \otimes_R M$。
\end{lemma}

\begin{proof}
略。
\end{proof}

\begin{definition}
\label{descent-definition-descent-datum-effective-module}
设 $R \to A$ 为环同态。若存在一个 $R$-模 $M$ 以及从
$(M \otimes_R A, can)$ 到 $(N, \varphi)$ 的下降数据同构，
则称下降数据 $(N, \varphi)$ 是{\it 有效的}。
\end{definition}

\noindent
设 $R \to A$ 为环同态，$(N, \varphi)$ 为下降数据。可以取与
$N_\bullet$ 相联系的上链复形 $s(N_\bullet)$（见 Simplicial, Section
\ref{simplicial-section-dold-kan-cosimplicial}）。它具有如下形状：
$$
N \to A \otimes_R N \to A \otimes_R A \otimes_R N \to \ldots
$$
这些态射可描述如下。第一个态射是
$$
n \longmapsto 1 \otimes n - \varphi(n \otimes 1).
$$
第二个态射在纯张量上的取值为
$$
a \otimes n \longmapsto 1 \otimes a \otimes n
- a \otimes 1 \otimes n + a \otimes \varphi(n \otimes 1).
$$
此后的规律是清楚的。

\medskip\noindent
在特殊情形 $N = A \otimes_R M$ 中，对任意 $m \in M$，元素
$1 \otimes m$ 都属于与余单纯模 $(A/R)_\bullet \otimes_R M$ 相联系的
上链复形的第一个态射之核。因此得到增广上链复形
\begin{equation}
\label{descent-equation-extended-complex}
0 \to M \to A \otimes_R M \to A \otimes_R A \otimes_R M \to \ldots
\end{equation}
这里把 $0$ 置于次数 $-2$，把模 $M$ 置于次数 $-1$，把模
$A \otimes_R M$ 置于次数 $0$，依此类推。注意，当 $M = R$ 时，
此复形的形状为
$$
0 \to R \to A \to A \otimes_R A \to A \otimes_R A \otimes_R A \to \ldots
$$

\begin{lemma}
\label{descent-lemma-with-section-exact}
设 $R \to A$ 有截面。则对任意 $R$-模 $M$，增广上链复形
(\ref{descent-equation-extended-complex}) 正合。
\end{lemma}

\begin{proof}
由 Simplicial, Lemma
\ref{simplicial-lemma-push-outs-simplicial-object-w-section}，态射
$R \to (A/R)_\bullet$ 是余单纯 $R$-代数之间的同伦等价（这里 $R$
表示常值余单纯 $R$-代数）。由于 $\otimes_R M$ 是从 $R$-代数范畴到
$R$-模范畴的函子，$M \to (A/R)_\bullet \otimes_R M$ 因而是余单纯
$R$-模范畴中的同伦等价；见 Simplicial, Lemma
\ref{simplicial-lemma-functorial-homotopy}。由此，相联系复形之间的诱导态射
也是同伦等价；见 Simplicial, Lemma
\ref{simplicial-lemma-homotopy-s-Q}。与常值余单纯 $R$-模 $M$ 相联系的
复形为
$$
\xymatrix{
M \ar[r]^0 & M \ar[r]^1 & M \ar[r]^0 & M \ar[r]^1 & M \ldots
}
$$
故结论成立（增广版本只是在开头再添一个 $M$）。
\end{proof}

\begin{lemma}
\label{descent-lemma-ff-exact}
设 $R \to A$ 忠实平坦；见 Algebra, Definition
\ref{algebra-definition-flat}。则对任意 $R$-模 $M$，增广上链复形
(\ref{descent-equation-extended-complex}) 正合。
\end{lemma}

\begin{proof}
假设能证明存在忠实平坦环同态 $R \to R'$，使结论对环同态
$R' \to A' = R' \otimes_R A$ 成立，则结论也对 $R \to A$ 成立。
事实上，对任意 $R$-模 $M$，余单纯模
$(M \otimes_R R') \otimes_{R'} (A'/R')_\bullet$ 正是余单纯模
$R' \otimes_R (M \otimes_R (A/R)_\bullet)$。因此，复形
$(M \otimes_R R') \otimes_{R'} (A'/R')_\bullet$ 的上同调消失，
再结合 $R \to R'$ 的忠实平坦性，即推出与
$M \otimes_R (A/R)_\bullet$ 相联系的复形之上同调消失。增广复形在次数
$-1$ 和 $0$ 的上同调群消失也可同理证明（证明从略）。

\medskip\noindent
这样的忠实平坦扩张确实存在。取 $R' = A$ 即可，因为环同态
$R' = A \to A' = A \otimes_R A$ 有截面 $a \otimes a' \mapsto aa'$，
从而可应用 Lemma \ref{descent-lemma-with-section-exact}。
\end{proof}

\noindent
下面说明此复形与有效性问题的关系。

\begin{lemma}
\label{descent-lemma-recognize-effective}
设 $R \to A$ 为忠实平坦环同态，$(N, \varphi)$ 为下降数据。则
$(N, \varphi)$ 有效，当且仅当典范态射
$$
A \otimes_R H^0(s(N_\bullet)) \longrightarrow N
$$
是同构。
\end{lemma}

\begin{proof}
若 $(N, \varphi)$ 有效，则可写成 $N = A \otimes_R M$ 且
$\varphi = can$。由 Lemmas
\ref{descent-lemma-canonical-descent-datum-cosimplicial} 和
\ref{descent-lemma-ff-exact}，可得 $H^0(s(N_\bullet)) = M$。反之，设引理中的
态射为同构。在此情形令 $M = H^0(s(N_\bullet))$。这是 $N$ 的一个
$R$-子模，即 $M = \{n \in N \mid 1 \otimes n = \varphi(n \otimes 1)\}$。
只需验证：通过同构 $A \otimes_R M \to N$，典范下降数据与
$\varphi$ 相符。验证从略。
\end{proof}

\begin{lemma}
\label{descent-lemma-descent-descends}
设 $R \to A$ 为忠实平坦环同态，且 $R \to R'$ 忠实平坦。令
$A' = R' \otimes_R A$。若相对于 $R' \to A'$ 的全部下降数据均有效，
则相对于 $R \to A$ 的全部下降数据也均有效。
\end{lemma}

\begin{proof}
设 $(N, \varphi)$ 是相对于 $R \to A$ 的下降数据。令
$N' = R' \otimes_R N = A' \otimes_A N$，并以
$\varphi' = \text{id}_{R'} \otimes \varphi$ 表示下降数据 $\varphi$
的基变换。于是 $(N', \varphi')$ 是相对于 $R' \to A'$ 的下降数据，且
$H^0(s(N'_\bullet)) = R' \otimes_R H^0(s(N_\bullet))$。此外，态射
$A' \otimes_{R'} H^0(s(N'_\bullet)) \to N'$ 可与 $A$-模态射
% P05-ERR-0018
$A \otimes_R H^0(s(N)) \to N$ 沿忠实平坦态射 $A \to A'$ 的基变换
等同。故由 Lemma \ref{descent-lemma-recognize-effective} 得出结论。
\end{proof}

\noindent
下面是本节的主要结果。它的证明可能显得有些笨拙；更高观点的处理方式见
下文 Remark \ref{descent-remark-homotopy-equivalent-cosimplicial-algebras}。

\begin{proposition}
\label{descent-proposition-descent-module}
\begin{slogan}
模沿忠实平坦环同态的有效下降。
\end{slogan}
设 $R \to A$ 为忠实平坦环同态。则
\begin{enumerate}
\item 模相对于 $R \to A$ 的任意下降数据都是有效的；
\item 从 $R$-模到下降数据范畴的函子
$M \mapsto (A \otimes_R M, can)$ 是等价；
\item 其逆函子由 $(N, \varphi) \mapsto H^0(s(N_\bullet))$ 给出。
\end{enumerate}
\end{proposition}

\begin{proof}
我们只证明 (1)，省略 (2) 和 (3) 的证明。由于 $R \to A$ 忠实平坦，
存在忠实平坦基变换 $R \to R'$，使
$R' \to A' = R' \otimes_R A$ 有截面（如 Lemma
\ref{descent-lemma-ff-exact} 的证明中那样取 $R' = A$ 即可）。因此，利用
Lemma \ref{descent-lemma-descent-descends}，可以假设 $R \to A$ 有一个截面，
记为 $\sigma : A \to R$。设 $(N, \varphi)$ 是相对于 $R \to A$ 的
下降数据。令
$$
M = H^0(s(N_\bullet)) = \{n \in N \mid 1 \otimes n = \varphi(n \otimes 1)\}
\subset
N
$$
由 Lemma \ref{descent-lemma-recognize-effective}，只需证明
$A \otimes_R M \to N$ 是同构。

\medskip\noindent
取元素 $n \in N$。对某些 $a_i \in A$ 和 $x_i \in N$，写作
$\varphi(n \otimes 1) = \sum a_i \otimes x_i$。由 Lemma
\ref{descent-lemma-descent-datum-cosimplicial}，在 $N$ 中有
$n = \sum a_i x_i$（因为任意余单纯对象中都有
$\sigma^0_0 \circ \delta^1_1 = \text{id}$）。接着，对某些
$a_{ij} \in A$ 和 $y_j \in N$，写作
$\varphi(x_i \otimes 1) = \sum a_{ij} \otimes y_j$。余圈条件意味着
$$
\sum a_i \otimes a_{ij} \otimes y_j = \sum a_i \otimes 1 \otimes x_i
$$
在 $A \otimes_R A \otimes_R N$ 中成立。由此得到两点：
\begin{enumerate}
\item 对第一个 $A$ 应用 $\sigma$，得到
$\sum \sigma(a_i) \varphi(x_i \otimes 1) = \sum \sigma(a_i) \otimes x_i$；
\item 对中间的 $A$ 应用 $\sigma$，得到
$\sum_i a_i \otimes \sum_j \sigma(a_{ij}) y_j = \sum a_i \otimes x_i$。
\end{enumerate}
由 (1) 可知 $\sum \sigma(a_i) x_i \in M$。将这一点应用于 $x_i$，
便知对每个 $i$ 都有
% P05-ERR-0019
$\sum \sigma(a_{ij})y_i \in M$。把 (2) 中的等式相乘，得到
$\sum_i a_i (\sum_j \sigma(a_{ij}) y_j) = \sum a_i x_i = n$。
故 $A \otimes_R M \to N$ 是满射。最后，设 $m_i \in M$ 且
$\sum a_i m_i = 0$。对 $\sum a_im_i \otimes 1$ 应用 $\varphi$，可见
$\sum a_i \otimes m_i = 0$。换言之，$A \otimes_R M \to N$ 是单射，
结论成立。
\end{proof}

\begin{remark}
\label{descent-remark-standard-covering}
设 $R$ 为环，且 $f_1, \ldots, f_n\in R$ 生成单位理想。环
$A = \prod_i R_{f_i}$ 是忠实平坦 $R$-代数。注意，余单纯环
$(A/R)_\bullet$ 在次数 $n$ 的环为
$$
\prod\nolimits_{i_0, \ldots, i_n} R_{f_{i_0}\ldots f_{i_n}}
$$
因此上述结果重新得到 Algebra, Lemmas
\ref{algebra-lemma-standard-covering}、\ref{algebra-lemma-cover-module}
和 \ref{algebra-lemma-glue-modules}。不过，由于高次数中的正合性，
上述结果实际给出了更强的结论：它蕴含仿射概形上拟凝聚层的
{\v C}ech 上同调平凡。由此得到 Cohomology of Schemes, Lemma
\ref{coherent-lemma-cech-cohomology-quasi-coherent-trivial} 的另一个证明。
\end{remark}

\begin{remark}
\label{descent-remark-homotopy-equivalent-cosimplicial-algebras}
设 $R$ 为环，$A_\bullet$ 为余单纯 $R$-代数。在此背景中，下降数据
对应于满足如下性质的余单纯 $A_\bullet$-模 $M_\bullet$：对每个
$n, m \geq 0$ 和每个 $\varphi : [n] \to [m]$，态射
$M(\varphi) : M_n \to M_m$ 诱导同构
$$
M_n \otimes_{A_n, A(\varphi)} A_m \longrightarrow M_m.
$$
称这样的余单纯模为{\it 笛卡尔模}。在此背景中，Proposition
\ref{descent-proposition-descent-module} 的证明可分为以下步骤：
\begin{enumerate}
\item 若 $R \to R'$ 和 $R \to A$ 均忠实平坦，则当相对于
$(R' \otimes_R A)/R'$ 的下降数据有效时，相对于 $A/R$ 的下降数据也有效。
\item 设 $A$ 为 $R$-代数。相对于 $A/R$ 的下降数据对应于笛卡尔
$(A/R)_\bullet$-模。
\item 若 $R \to A$ 有截面，则 $(A/R)_\bullet$ 与 $R$ 同伦等价；
这里后者是取值为 $R$ 的常值余单纯 $R$-代数。
\item 若 $A_\bullet \to B_\bullet$ 是余单纯 $R$-代数之间的同伦等价，
则函子 $M_\bullet \mapsto M_\bullet \otimes_{A_\bullet} B_\bullet$
诱导笛卡尔 $A_\bullet$-模范畴与笛卡尔 $B_\bullet$-模范畴之间的等价。
\end{enumerate}
关于 (1)，见 Lemma \ref{descent-lemma-descent-descends}；(2) 使用 Lemma
\ref{descent-lemma-descent-datum-cosimplicial}；(3) 已在 Lemma
\ref{descent-lemma-with-section-exact} 的证明中见过（它依赖 Simplicial,
Lemma \ref{simplicial-lemma-push-outs-simplicial-object-w-section}）。
此外，只要从正确角度思考，(4) 是显然的！
\end{remark}








\section{普遍单射态射的下降}
\label{descent-section-descent-universally-injective}

\noindent
代数几何中的许多构造都使用{\it 下降}技术。例如，为了在给定空间上
构造对象，可以先在投影到该空间的某个稍大空间上工作；也可以沿覆盖态射
拉回，以验证空间或态射的某种性质。这类技术是否有用，当然取决于能否
识别出足够广泛的一类{\it 有效下降态射}。现代代数几何的 Grothendieck
式发展初期，人们已经证明：{\it 拟紧}且{\it 忠实平坦}的态射类对于
对象、态射及其许多性质的下降都是有效的。

\medskip\noindent
与通常一样，这一陈述归结为环与模的性质。Grothendieck 证明：要使环同态
$f: R \to S$ 成为模的有效下降态射，$f$ 忠实平坦便足够。然而，此条件
排除了许多自然例子；例如，任意可裂环同态都是有效下降态射。甚至在
忠实平坦下降的证明中就自然出现了一个这样的例子：对任意环同态
$f: R \to S$，无论它是否平坦，$1_S \otimes f: S \to S \otimes_R S$
都由乘法态射给出分裂。

\medskip\noindent
于是可以追问：是否存在一个自然的环论条件，既涵盖忠实平坦态射，
又涵盖可裂环同态，并且蕴含模的有效下降？令人意外的是（至少本节作者
对此感到意外），这个问题的完整答案早在约 1970 年就已为人所知！
不难验证，$f: R \to S$ 成为模的有效下降态射的一个必要条件是：$f$
在 $R$-模范畴中{\it 普遍单射}；也就是说，对任意 $R$-模 $M$，态射
$1_M  \otimes f: M \to M \otimes_R S$ 必须为单射。事实证明，此条件
也充分。例如，若 $f$ 在 $R$-模范畴中可裂（但在环范畴中未必可裂），
则 $f$ 是模的有效下降态射。
 
\medskip\noindent
这一结果的历史略显曲折。Olivier \cite{olivier} 最先断言了它，并把
普遍单射态射称为{\it 纯}态射，但没有清楚说明证明。从 Joyal 与 Tierney
的工作 \cite{joyal-tierney} 中可以提取出此结果；不过据我们所知，
文献中最早出现的独立证明是 Mesablishvili 的证明
\cite{mesablishvili1}。本节的第一个目的就是阐述 Mesablishvili 的证明。
除改正若干笔误外，他的原始叙述几乎不需修改；唯一的例外是，我们将明确
说明代数几何中下降数据的惯常定义与 \cite{mesablishvili1} 所用定义之间
的关系。该证明完全是范畴论的，因而可以用单子的语言表述（并应用于其他
背景）；见 \cite{janelidze-tholen}。

\medskip\noindent
本节的第二个目的是汇集一些信息，说明模、代数和态射的哪些性质可以沿
普遍单射环同态下降。有限模与平坦模的情形由 Mesablishvili
\cite{mesablishvili2} 处理。


\subsection{范畴论预备知识}
\label{descent-subsection-category-prelims}

\noindent
先回顾几个可简化叙述的范畴论基本概念。本小节固定一个环境范畴。

\medskip\noindent
回忆：对两个态射 $g_1, g_2: B \to C$，$g_1$ 和 $g_2$ 的
{\it 等化子}是满足
$g_1 \circ f = g_2 \circ f$ 且对该性质具有泛性质的态射
$f: A \to B$。后一个陈述意味着，任何不含虚线箭头的可交换图
$$
\xymatrix{A' \ar[rd]^e \ar@/^1.5pc/[rrd] \ar@{-->}[d] & & \\
A \ar[r]^f & B \ar@<1ex>[r]^{g_1} \ar@<-1ex>[r]_{g_2} & 
C
}
$$
都能以唯一方式补全。在此情形下，也称图
\begin{equation}
\label{descent-equation-equalizer}
\xymatrix{
A \ar[r]^f  & B \ar@<1ex>[r]^{g_1} \ar@<-1ex>[r]_{g_2} & C
}
\end{equation}
为等化子。反转箭头即得到{\it 余等化子}的定义。见 Categories,
Sections \ref{categories-section-equalizers} 和
\ref{categories-section-coequalizers}。

\medskip\noindent
由于等化子涉及泛性质，应用协变函子通常不会保持等化子性质。与单态射和
满态射的情形一样，在某些情况下可以通过给出分裂来绕过这一问题。

\begin{definition}
\label{descent-definition-split-equalizer}
一个{\it 可裂等化子}是满足 $g_1 \circ f = g_2 \circ f$ 的图
(\ref{descent-equation-equalizer})，并且存在辅助态射 $h : B \to A$ 和
$i : C \to B$，使得
\begin{equation}
\label{descent-equation-split-equalizer-conditions}
h \circ f = 1_A, \quad f \circ h = i \circ g_1, \quad i \circ g_2 = 1_B.
\end{equation}
\end{definition}

\noindent
关键在于，箭头之间的这些等式迫使 (\ref{descent-equation-equalizer}) 成为等化子：
写成 $e = f \circ (h \circ e)$，态射 $e$ 便唯一地通过 $f$ 分解。因此，
对可裂等化子应用协变函子仍得到可裂等化子；应用反变函子则得到
{\it 可裂余等化子}，其定义是显然的。

\subsection{普遍单射态射}
\label{descent-subsection-universally-injective}

\noindent
回忆，$\textit{Rings}$ 表示含 $1$ 交换环的范畴。对 $\textit{Rings}$ 的对象
$R$，以 $\text{Mod}_R$ 表示 $R$-模范畴。

\begin{remark}
\label{descent-remark-reflects}
设 $F : \mathcal{A} \to \mathcal{B}$ 是 Abel 范畴之间的正合函子，
且把非零对象映为非零对象，则它反映单射和满射。事实上，正合性意味着
$F$ 保持核与余核（比较 Homology, Section
\ref{homology-section-functors}）。例如，若 $f : R \to S$ 是忠实平坦
环同态，则 $\bullet \otimes_R S: \text{Mod}_R \to \text{Mod}_S$
具有这些性质。
\end{remark}

\noindent
设 $R$ 为环。回忆，若对所有 $P \in \text{Mod}_R$，态射
$f \otimes 1_P: M \otimes_R P \to N \otimes_R P$ 都是单射，
则称 $\text{Mod}_R$ 中的态射 $f : M \to N$ 是{\it 普遍单射}的。
见 Algebra, Definition \ref{algebra-definition-universally-injective}。

\begin{definition}
\label{descent-definition-universally-injective}
若环同态 $f: R \to S$ 作为 $\text{Mod}_R$ 中的态射是普遍单射的，
则称它是{\it 普遍单射}的。
\end{definition}

\begin{example}
\label{descent-example-split-injection-universally-injective}
$\text{Mod}_R$ 中任意可裂单射都是普遍单射的。特别地，
$\textit{Rings}$ 中任意可裂单射都是普遍单射的。
\end{example}

\begin{example}
\label{descent-example-cover-universally-injective}
设 $R$ 为环，且 $f_1, \ldots, f_n \in R$ 生成单位理想，则态射
$R \to R_{f_1} \oplus \ldots \oplus R_{f_n}$ 是普遍单射的。尽管这立即
由 Lemma \ref{descent-lemma-faithfully-flat-universally-injective} 得出，直接验证
仍有启发性：立刻可化归到 $R$ 为局部环的情形。此时某个 $f_i$ 必为单位，
故态射 $R \to R_{f_i}$ 是同构。
\end{example}

\begin{lemma}
\label{descent-lemma-faithfully-flat-universally-injective}
任意忠实平坦环同态都是普遍单射的。
\end{lemma}

\begin{proof}
这只是 Algebra, Lemma
\ref{algebra-lemma-faithfully-flat-universally-injective} 的改述。
\end{proof}

\noindent
\cite{mesablishvili1} 中的关键观察是：利用 $\text{Mod}_R$ 中内射余生成元
的通常构造，可以借助分裂有效地改述普遍单射性。

\begin{definition}
\label{descent-definition-C}
设 $R$ 为环。定义反变函子
{\it $C$} $ : \text{Mod}_R \to \text{Mod}_R$ 如下：
$$
C(M) = \Hom_{\textit{Ab}}(M, \mathbf{Q}/\mathbf{Z}),
$$
其中 $C(M)$ 上的 $R$-作用由 $rf(s) = f(rs)$ 给出。
\end{definition}

\noindent
在 More on Algebra, Section
\ref{more-algebra-section-injectives-modules} 中，此函子记为
$M \mapsto M^\vee$。

\begin{lemma}
\label{descent-lemma-C-is-faithful}
对任意环 $R$，函子 $C : \text{Mod}_R \to \text{Mod}_R$ 正合，
并且反映单射和满射。
\end{lemma}

\begin{proof}
正合性见 More on Algebra, Lemma \ref{more-algebra-lemma-vee-exact}；
其余性质由此得出，见 Remark \ref{descent-remark-reflects}。
\end{proof}

\begin{remark}
\label{descent-remark-adjunction}
我们将经常使用 $\Hom$ 与张量积之间的标准伴随，其形式是反变函子的
自然同构
\begin{equation}
\label{descent-equation-adjunction}
C(\bullet_1 \otimes_R \bullet_2) \cong \Hom_R(\bullet_1, C(\bullet_2)): 
\text{Mod}_R \times \text{Mod}_R \to \text{Mod}_R
\end{equation}
它把 $f: M_1 \otimes_R M_2 \to \mathbf{Q}/\mathbf{Z}$ 映为
$m_1 \mapsto 
(m_2 \mapsto f(m_1 \otimes m_2))$。见 Algebra, Lemma
\ref{algebra-lemma-hom-from-tensor-product-variant}。由此可得如下推论：若
$$
\xymatrix@C=9pc{
C(M) \ar@<1ex>[r] \ar@<-1ex>[r] & C(N) \ar[r] & C(P)
}
$$
是 $\text{Mod}_R$ 中的可裂余等化子图，则对任意
$Q \in \text{Mod}_R$，图
$$
\xymatrix@C=9pc{
C(M \otimes_R Q) \ar@<1ex>[r] \ar@<-1ex>[r] & C(N \otimes_R Q) \ar[r] & C(P 
\otimes_R Q)
}
$$
也是可裂余等化子图。
\end{remark}

\begin{lemma}
\label{descent-lemma-split-surjection}
设 $R$ 为环。$\text{Mod}_R$ 中的态射 $f: M \to N$ 普遍单射，当且仅当
$C(f): C(N) \to C(M)$ 是可裂满射。
\end{lemma}

\begin{proof}
由 (\ref{descent-equation-adjunction})，对任意 $P \in \text{Mod}_R$ 都有可交换图
$$
\xymatrix@C=9pc{
\Hom_R( P, C(N)) \ar[r]_{\Hom_R(P,C(f))} \ar[d]^{\cong} &
\Hom_R(P,C(M)) \ar[d]^{\cong} \\
C(P \otimes_R N ) \ar[r]^{C(1_{P} \otimes f)} & C(P \otimes_R M ).
}
$$
若 $f$ 普遍单射，则
$1_{C(M)} \otimes f: C(M) \otimes_R M \to 
C(M) \otimes_R N$ 是单射，
故取 $P = C(M)$ 时，上图的两行都是满射。因此，可以把
$1_{C(M)} \in \Hom_R(C(M), C(M))$ 提升为某个使 $C(f)$ 分裂的
% P05-ERR-0020
$g \in \Hom_R(C(N), C(M))$。反之，若 $C(f)$ 是可裂满射，则上图的
两行都是满射，故由 Lemma \ref{descent-lemma-C-is-faithful}，
$1_{P} \otimes f$ 是单射。
\end{proof}

\begin{remark}
\label{descent-remark-functorial-splitting}
设 $f: M \to N$ 是 $\text{Mod}_R$ 中的普遍单射态射。选取 $C(f)$ 的一个
分裂 $g$ 后，可对每个 $P \in \text{Mod}_R$ 构造
$C(1_P \otimes f)$ 的函子性分裂。事实上，由
(\ref{descent-equation-adjunction})，这等价于关于 $P$ 函子性地分裂
$\Hom_R(P, 
C(f))$，而态射 $g \circ \bullet$ 即实现了这一点。
\end{remark}


\subsection{模及其态射的下降}
\label{descent-subsection-descent-modules-morphisms}

\noindent
本小节始终固定一个环同态 $f: R \to S$。如 Section
\ref{descent-section-descent-modules} 所见，可以用余单纯代数的语言讨论模的
下降数据；但本小节倾向于使用更具体的术语。

\medskip\noindent
对 $i = 1, 2, 3$，令 $S_i$ 为 $S$ 在 $R$ 上的 $i$ 重张量积。按下列公式
定义环同态 $\delta_0^1, \delta_1^1: S_1 \to S_2$、
$\delta_{01}^1, \delta_{02}^1, \delta_{12}^1: S_1 \to S_3$ 以及
$\delta_0^2, \delta_1^2, \delta_2^2: S_2 \to S_3$：
\begin{align*}
\delta^1_0  (a_0) & =  1 \otimes a_0 \\
\delta^1_1  (a_0) & = a_0 \otimes 1 \\
\delta^2_0  (a_0 \otimes a_1) & =  1 \otimes a_0 \otimes a_1 \\
\delta^2_1  (a_0 \otimes a_1) & =  a_0 \otimes 1 \otimes a_1 \\
\delta^2_2  (a_0 \otimes a_1) & =  a_0 \otimes a_1 \otimes 1 \\
\delta_{01}^1(a_0) & = 1 \otimes 1 \otimes a_0 \\
\delta_{02}^1(a_0) & = 1 \otimes a_0 \otimes 1 \\
\delta_{12}^1(a_0) & = a_0 \otimes 1 \otimes 1.
\end{align*}
换言之，上指标表示源环，下指标表示插入因子 1 的位置。（此记号与
Section \ref{descent-section-descent-modules} 中引入的记号相容。）

\medskip\noindent
由 Definition \ref{descent-definition-descent-datum-modules} 回忆\footnote{准确地说，
这里的 $\theta$ 是 Definition \ref{descent-definition-descent-datum-modules} 中
$\varphi$ 的逆。}：对 $M \in \text{Mod}_S$，$M$ 相对于 $f$ 的一个
{\it 下降数据}是一个 $S_2$-模同构
$$
\theta :
M \otimes_{S,\delta^1_0} S_2
\longrightarrow
M \otimes_{S,\delta^1_1} S_2
$$
并且满足{\it 余圈条件}
\begin{equation}
\label{descent-equation-cocycle-condition}
(\theta \otimes \delta_2^2) \circ (\theta \otimes \delta_2^0) = (\theta \otimes 
\delta_2^1):
M \otimes_{S, \delta^1_{01}} S_3 \to M \otimes_{S,\delta^1_{12}} S_3.
\end{equation}
以 $DD_{S/R}$ 表示配备相对于 $f$ 的下降数据的 $S$-模所成的范畴。

\medskip\noindent
例如，给定 $M_0 \in \text{Mod}_R$ 并选取同构
$M \cong M_0 \otimes_R S$，把 $M \otimes_{S,\delta^1_0} S_2$ 和
$M \otimes_{S,\delta^1_1} S_2$ 自然地等同于 $M_0 \otimes_R S_2$，
便得到一个下降数据。特别地，此构造定义函子
$f^*: \text{Mod}_R \to DD_{S/R}$。

\begin{definition}
\label{descent-definition-effective-descent}
函子 $f^*: \text{Mod}_R \to DD_{S/R}$ 称为{\it 沿 $f$ 的底扩张}。
若 $f^*$ 全忠实，则称 $f$ 是{\it 模的下降态射}；若 $f^*$ 是范畴等价，
则称 $f$ 是{\it 模的有效下降态射}。
\end{definition}

\noindent
我们的目标是证明：当 $f$ 普遍单射时，可以利用 $\theta$ 在 $M$ 中定位
$M_0$。这一过程关键地使用若干等化子图。

\begin{lemma}
\label{descent-lemma-equalizer-M}
对 $(M,\theta) \in DD_{S/R}$，图
\begin{equation}
\label{descent-equation-equalizer-M}
\xymatrix@C=8pc{
M \ar[r]^{\theta \circ (1_M \otimes \delta_0^1)} &
M \otimes_{S, \delta_1^1} S_2
\ar@<1ex>[r]^{(\theta \otimes \delta_2^2) \circ (1_M \otimes \delta^2_0)}
\ar@<-1ex>[r]_{1_{M \otimes S_2} \otimes \delta^2_1} & 
M \otimes_{S, \delta_{12}^1} S_3
}
\end{equation}
是可裂等化子。
\end{lemma}

\begin{proof}
按下列公式定义环同态 $\sigma^0_0: S_2 \to S_1$ 以及 $\sigma_0^1, 
\sigma_1^1: S_3 \to S_2$：
\begin{align*}
\sigma^0_0 (a_0 \otimes a_1) & = a_0a_1 \\
\sigma^1_0 (a_0 \otimes a_1 \otimes a_2) & = a_0a_1 \otimes a_2 \\
\sigma^1_1 (a_0 \otimes a_1 \otimes a_2) & = a_0 \otimes a_1a_2.
\end{align*}
取辅助态射为
$1_M \otimes \sigma_0^0: M \otimes_{S, \delta_1^1} S_2 \to M$
和 $1_M \otimes \sigma_0^1: M \otimes_{S,\delta_{12}^1} S_3 \to M \otimes_{S, 
\delta_1^1} S_2$。(\ref{descent-equation-split-equalizer-conditions}) 所要求的
相容性中，第一个由余圈条件 (\ref{descent-equation-cocycle-condition}) 与
$\sigma_1^1$ 作张量积得出，其余则是直接的。
\end{proof}

\begin{lemma}
\label{descent-lemma-equalizer-CM}
对 $(M, \theta) \in DD_{S/R}$，图
\begin{equation}
\label{descent-equation-coequalizer-CM}
\xymatrix@C=8pc{
C(M \otimes_{S, \delta_{12}^1} S_3)
\ar@<1ex>[r]^{C((\theta \otimes \delta_2^2) \circ (1_M \otimes \delta^2_0))}
\ar@<-1ex>[r]_{C(1_{M \otimes S_2} \otimes \delta^2_1)} &
C(M \otimes_{S, \delta_1^1} S_2 )
\ar[r]^{C(\theta \circ (1_M \otimes \delta_0^1))} & C(M).
}
\end{equation}
由对 (\ref{descent-equation-equalizer-M}) 应用 $C$ 得到，并且是可裂余等化子。
\end{lemma}

\begin{proof}
略。
\end{proof}

\begin{lemma}
\label{descent-lemma-equalizer-S}
图
\begin{equation}
\label{descent-equation-equalizer-S}
\xymatrix@C=8pc{
S_1 \ar[r]^{\delta^1_1} &
S_2 \ar@<1ex>[r]^{\delta^2_2} \ar@<-1ex>[r]_{\delta^2_1} & 
S_3
}
\end{equation}
是可裂等化子。
\end{lemma}

\begin{proof}
在 Lemma \ref{descent-lemma-equalizer-M} 中取 $(M, \theta) = f^*(S)$。
\end{proof}

\noindent
这提示了 $f^*$ 的一个候选拟逆函子的定义。

\begin{definition}
\label{descent-definition-pushforward}
定义函子 {\it $f_*$} $: DD_{S/R} \to \text{Mod}_R$：令
$f_*(M, \theta)$ 为 $M$ 的 $R$-子模，使得图
\begin{equation}
\label{descent-equation-equalizer-f}
\xymatrix@C=8pc{f_*(M,\theta) \ar[r] & M \ar@<1ex>^{\theta \circ (1_M \otimes 
\delta_0^1)}[r] \ar@<-1ex>_{1_M \otimes \delta_1^1}[r] & 
M \otimes_{S, \delta_1^1} S_2 
}
\end{equation}
是等化子。
\end{definition}

\noindent
利用 Lemma \ref{descent-lemma-equalizer-M}，并注意限制函子
$\text{Mod}_S \to \text{Mod}_R$ 是底扩张函子
$\bullet \otimes_R S: \text{Mod}_R \to \text{Mod}_S$ 的右伴随，
可知 $f_*$ 是 $f^*$ 的右伴随。

\medskip\noindent
现在可以证明关键引理。在忠实平坦的情形，结论是显然的（见 Remark
\ref{descent-remark-descent-lemma}）；一般情形则需要一些论证。

\begin{lemma}
\label{descent-lemma-descent-lemma}
若 $f$ 普遍单射，则图
\begin{equation}
\label{descent-equation-equalizer-f2}
\xymatrix@C=8pc{
f_*(M, \theta) \otimes_R S
\ar[r]^{\theta \circ (1_M \otimes \delta_0^1)} &
M \otimes_{S, \delta_1^1} S_2 
\ar@<1ex>[r]^{(\theta \otimes \delta_2^2) \circ (1_M \otimes \delta^2_0)}
\ar@<-1ex>[r]_{1_{M \otimes S_2} \otimes \delta^2_1} &
M \otimes_{S, \delta_{12}^1} S_3
}
\end{equation}
由 (\ref{descent-equation-equalizer-f}) 在 $R$ 上与 $S$ 作张量积得到，
并且是等化子。
\end{lemma}

\begin{proof}
由 Lemma \ref{descent-lemma-split-surjection} 和 Remark
\ref{descent-remark-functorial-splitting}，态射
$C(1_N \otimes f): C(N \otimes_R S) \to C(N)$ 可关于 $N$ 函子性地分裂。
这给出下列可交换图中上方的竖直箭头：
$$
\xymatrix@C=8pc{
C(M \otimes_{S, \delta_1^1} S_2)
\ar@<1ex>^{C(\theta \circ (1_M \otimes \delta_0^1))}[r]
\ar@<-1ex>_{C(1_M \otimes \delta_1^1)}[r] \ar[d] &
C(M) \ar[r]\ar[d] & C(f_*(M,\theta)) \ar@{-->}[d] \\
C(M \otimes_{S,\delta_{12}^1} S_3)
\ar@<1ex>^{C((\theta \otimes \delta_2^2) \circ (1_M \otimes \delta^2_0))}[r]
\ar@<-1ex>_{C(1_{M \otimes S_2} \otimes \delta^2_1)}[r] \ar[d] &
C(M \otimes_{S, \delta_1^1} S_2 )
\ar[r]^{C(\theta \circ (1_M \otimes \delta_0^1))}
\ar[d]^{C(1_M \otimes \delta_1^1)} &
C(M) \ar[d] \ar@{=}[dl] \\
C(M \otimes_{S, \delta_1^1} S_2)
\ar@<1ex>[r]^{C(\theta \circ (1_M \otimes \delta_0^1))}
\ar@<-1ex>[r]_{C(1_M \otimes \delta_1^1)} &
C(M) \ar[r] &
C(f_*(M,\theta))
}
$$
其中沿各列的复合都是恒等态射。第二行是余等化子图
(\ref{descent-equation-coequalizer-CM})，因而产生虚线箭头。由右上方的方块得到
辅助态射 $C(f_*(M,\theta)) \to 
C(M)$ 和 $C(M) \to C(M\otimes_{S,\delta_1^1} S_2)$，它们蕴含第一行
是可裂余等化子图。由 Remark \ref{descent-remark-adjunction}，可在 $C$ 内部与
$S$ 作张量积，得到可裂余等化子图
$$
\xymatrix@C=8pc{
C(M \otimes_{S,\delta_2^2 \circ \delta_1^1} S_3)
\ar@<1ex>^{C((\theta \otimes \delta_2^2) \circ (1_M \otimes \delta^2_0))}[r] 
\ar@<-1ex>_{C(1_{M \otimes S_2} \otimes \delta^2_1)}[r] &
C(M \otimes_{S, \delta_1^1} S_2 )
\ar[r]^{C(\theta \circ (1_M \otimes \delta_0^1))} &
C(f_*(M,\theta) \otimes_R S).
}
$$
由 Lemma \ref{descent-lemma-C-is-faithful}，可知
(\ref{descent-equation-equalizer-f2}) 也必为等化子。
\end{proof}

\begin{remark}
\label{descent-remark-descent-lemma}
若 $f$ 是 $\text{Mod}_R$ 中的可裂单射，则可直接分裂 $f$，无须使用
$C$，从而简化论证。若 $f$ 忠实平坦，情况更为简单：此时 Lemma
\ref{descent-lemma-descent-lemma} 的结论立即成立，因为在 $R$ 上与 $S$ 作张量积
保持所有等化子。
\end{remark}

\begin{theorem}
\label{descent-theorem-descent}
下列条件等价。
\begin{enumerate}
\item[(a)] 态射 $f$ 是模的下降态射。
\item[(b)] 态射 $f$ 是模的有效下降态射。
\item[(c)] 态射 $f$ 普遍单射。
\end{enumerate}
\end{theorem}

\begin{proof}
(b) 蕴含 (a) 是显然的。下面验证 (a) 蕴含 (c)。若 $f$ 不是普遍单射，
则可找到 $M \in \text{Mod}_R$，使态射
$1_M \otimes f: M \to M \otimes_R S$ 有非零的核 $N$。自然投影
$M \to M/N$ 不是同构，但它在 $DD_{S/R}$ 中的像是同构。因此 $f^*$
并非全忠实。

\medskip\noindent
最后验证 (c) 蕴含 (b)。由 Lemmas \ref{descent-lemma-equalizer-M} 和
\ref{descent-lemma-descent-lemma}，对 $(M, \theta) \in DD_{S/R}$，自然态射
$f^* f_*(M,\theta) \to M$ 是 $S$-模同构。对
$M_0 \in \text{Mod}_R$，需要证明 $M_0 \to f_*f^*M_0$ 是同构。
由伴随性，复合 $f^*M_0 \to f^*f_*f^*M_0 \to f^*M_0$ 是恒等态射，
而上文表明第二个箭头是同构。因此
$f^*M_0 \to f^*f_*f^*M_0$ 是同构。由 Algebra, Lemma
\ref{algebra-lemma-base-change-along-universally-injective-ring-map}，
可知 $M_0 \to f_* f^* M_0$ 是同构。因此，$f_*$ 与 $f^*$ 互为拟逆函子，
结论得证。
\end{proof}

\subsection{模的性质的下降}
\label{descent-subsection-descent-properties-modules}

\noindent
本小节始终固定一个普遍单射环同态 $f : R \to S$、一个对象
$M \in \text{Mod}_R$ 以及一个环同态 $R \to A$。现在研究 $M$ 或 $A$
的哪些性质可以在沿 $f$ 作底扩张后检验。先从 \cite{mesablishvili2}
中的一些结果开始。

\begin{lemma}
\label{descent-lemma-flat-to-injective}
若 $M \in \text{Mod}_R$ 平坦，则 $C(M)$ 是内射 $R$-模。
\end{lemma}

\begin{proof}
设 $0 \to N \to P \to Q \to 0$ 是 $\text{Mod}_R$ 中的正合列。
由于 $M$ 平坦，
$$
0 \to N \otimes_R M \to P \otimes_R M \to Q \otimes_R M \to 0
$$
正合。由 Lemma \ref{descent-lemma-C-is-faithful}，
$$
0 \to C(Q \otimes_R M) \to C(P \otimes_R M) \to C(N \otimes_R M) \to 0
$$
正合。由 (\ref{descent-equation-adjunction})，后一个序列可改写为
$$
0 \to \Hom_R(Q, C(M)) \to \Hom_R(P, C(M)) \to \Hom_R(N, C(M)) \to 0.
$$
因此 $C(M)$ 是 $\text{Mod}_R$ 中的内射对象。
\end{proof}

\begin{theorem}
\label{descent-theorem-descend-module-properties}
若 $S$-模 $M \otimes_R S$ 具有下列性质之一：
\begin{enumerate}
\item[(a)]
有限生成；
\item[(b)]
有限表现；
\item[(c)]
平坦；
\item[(d)]
忠实平坦；
\item[(e)]
有限投射；
\end{enumerate}
则 $R$-模 $M$ 也具有相应性质（反之亦然）。
\end{theorem}

\begin{proof}
为证明 (a)，在 $\text{Mod}_S$ 中选取 $M \otimes_R S$ 的有限生成集
$\{n_i\}$。将每个 $n_i$ 写成 $\sum_j m_{ij} \otimes s_{ij}$，其中
$m_{ij} \in M$ 且 $s_{ij} \in S$。设 $F$ 是以 $e_{ij}$ 为基的有限自由
$R$-模，并令 $R$-模态射 $F \to M$ 把 $e_{ij}$ 映至 $m_{ij}$。于是
$F \otimes_R S\to M \otimes_R S$ 是满射，故
$\Coker(F \to M) \otimes_R S$ 为零，进而 $\Coker(F \to M)$ 为零。
这证明了 (a)。

\medskip\noindent
为证明 (b)，设 $M \otimes_R S$ 有限表现。由 (a)，$M$ 有限生成。选取核为
$K$ 的满射 $R^{\oplus n} \to M$。于是
% P05-ERR-0021
$K \otimes_R S \to S^{\oplus r} \to M \otimes_R S \to 0$ 正合。
由 Algebra, Lemma \ref{algebra-lemma-extension}，
$S^{\oplus r} \to M \otimes_R S$ 的核是有限 $S$-模。因此可找到有限多个
元素 $k_1, \ldots, k_t \in K$，使 $k_i \otimes 1$ 在
$S^{\oplus r}$ 中的像生成 $S^{\oplus r} \to M \otimes_R S$ 的核。
设 $K' \subset K$ 是由 $k_1, \ldots, k_t$ 生成的子模。则
$M' = R^{\oplus r}/K'$ 是有限表现 $R$-模，并有态射 $M' \to M$，
使 $M' \otimes_R S \to M \otimes_R S$ 为同构。故如所需，
$M' \cong M$。

\medskip\noindent
为证明 (c)，设 $0 \to M' \to M'' \to M \to 0$ 是 $\text{Mod}_R$ 中的短正合列。
由于 $\bullet \otimes_R S$ 是右正合函子，
$M'' \otimes_R S \to M \otimes_R S$ 为满射。故由 Lemma
\ref{descent-lemma-C-is-faithful}，态射
$C(M \otimes_R S) \to C(M'' \otimes_R S)$ 为单射。若
$M \otimes_R S$ 平坦，则 Lemma \ref{descent-lemma-flat-to-injective} 表明
$C(M \otimes_R S)$ 是 $\text{Mod}_S$ 中的内射对象。因此单射
$C(M \otimes_R S) \to C(M'' \otimes_R S)$ 在 $\text{Mod}_S$ 中可裂，
从而在 $\text{Mod}_R$ 中也可裂。由 Lemma
\ref{descent-lemma-split-surjection}，$C(M \otimes_R S) \to C(M)$ 是可裂满射，
所以 $C(M) \to C(M'')$ 是 $\text{Mod}_R$ 中的可裂单射。也就是说，序列
$$
0 \to C(M) \to C(M'') \to C(M') \to 0
$$
可裂正合。对 $N \in \text{Mod}_R$，由 (\ref{descent-equation-adjunction}) 可知
$$
0 \to C(M \otimes_R N) \to C(M'' \otimes_R N) \to C(M' \otimes_R N) \to 0
$$
可裂正合。由 Lemma \ref{descent-lemma-C-is-faithful}，
$$
0 \to M' \otimes_R N \to M'' \otimes_R N \to M \otimes_R N \to 0
$$
正合。这蕴含 $M$ 在 $R$ 上平坦。事实上，
% P05-ERR-0022
取 $M'$ 为满射到 $M$ 的自由模，可知对所有模 $N$ 都有
$\text{Tor}_1^R(M, N) = 0$，于是可应用 Algebra, Lemma
\ref{algebra-lemma-characterize-flat}。这证明了 (c)。

\medskip\noindent
为由 (c) 推出 (d)，注意：若 $N \in \text{Mod}_R$ 且 $M \otimes_R N$
为零，则 $M \otimes_R S \otimes_S (N \otimes_R S) \cong (M \otimes_R N) \otimes_R 
S$ 为零。因此 $N \otimes_R S$ 为零，进而 $N$ 为零。

\medskip\noindent
此时，为推出 (e)，只需回忆：$M$ 有限生成且投射，当且仅当它有限表现且
平坦。见 Algebra, Lemma \ref{algebra-lemma-finite-projective}。
\end{proof}

\noindent
对 $R$-代数也有如下版本。

\begin{theorem}
\label{descent-theorem-descend-algebra-properties}
若 $S$-代数 $A \otimes_R S$ 具有下列性质之一：
\begin{enumerate}
\item[(a)]
有限型；
\item[(b)]
有限表现；
\item[(c)]
形式非分歧；
\item[(d)]
非分歧；
\item[(e)]
\'etale；
\end{enumerate}
则 $R$-代数 $A$ 也具有相应性质（反之当然也成立）。
\end{theorem}

\begin{proof}
为证明 (a)，选取 $A \otimes_R S$ 在 $S$ 上的有限生成集 $\{x_i\}$。
将每个 $x_i$ 写成 $\sum_j y_{ij} \otimes s_{ij}$，其中
$y_{ij} \in A$ 且 $s_{ij} \in S$。令 $F$ 为以变量 $e_{ij}$ 生成的多项式
$R$-代数，并令把 $e_{ij}$ 映至 $y_{ij}$ 的 $R$-代数同态为
% P05-ERR-0023
$F \to M$。于是 $F \otimes_R S\to A \otimes_R S$ 是满射，故
$\Coker(F \to A) \otimes_R S$ 为零，进而 $\Coker(F \to A)$ 为零。
这证明了 (a)。

\medskip\noindent
为证明 (b)，设 $A \otimes_R S$ 是有限表现 $S$-代数。由 (a)，$A$ 在
$R$ 上有限型。选取核为 $I$ 的满射 $R[x_1, \ldots, x_n] \to A$。
于是 $I \otimes_R S \to S[x_1, \ldots, x_n] \to A \otimes_R S \to 0$
正合。由 Algebra, Lemma
\ref{algebra-lemma-finite-presentation-independent}，
$S[x_1, \ldots, x_n] \to A \otimes_R S$ 的核是有限生成理想。因此可找到
有限多个元素 $y_1, \ldots, y_t \in I$，使 $y_i \otimes 1$ 在
$S[x_1, \ldots, x_n]$ 中的像生成
$S[x_1, \ldots, x_n] \to A \otimes_R S$ 的核。令 $I' \subset I$ 为
$y_1, \ldots, y_t$ 生成的理想。则
$A' = R[x_1, \ldots, x_n]/I'$ 是有限表现 $R$-代数，并有态射
$A' \to A$，使 $A' \otimes_R S \to A \otimes_R S$ 为同构。
故如所需，$A' \cong A$。

\medskip\noindent
为证明 (c)，回忆：$A$ 在 $R$ 上形式非分歧，当且仅当相对微分模
$\Omega_{A/R}$ 消失；见 Algebra, Lemma
\ref{algebra-lemma-characterize-formally-unramified} 或
\cite[Proposition~17.2.1]{EGA4}。由于
$\Omega_{(A \otimes_R S)/S} = \Omega_{A/R} \otimes_R S$，
此消失性由 Theorem \ref{descent-theorem-descent} 下降。

\medskip\noindent
为由前述情形推出 (d)，回忆：$A$ 在 $R$ 上非分歧，当且仅当 $A$
在 $R$ 上形式非分歧且有限型；见 Algebra, Lemma
\ref{algebra-lemma-formally-unramified-unramified}。

\medskip\noindent
为证明 (e)，回忆：由 Algebra, Lemma
\ref{algebra-lemma-etale-flat-unramified-finite-presentation} 或
\cite[Th\'eor\`eme~17.6.1]{EGA4}，代数 $A$ 在 $R$ 上 \'etale，
当且仅当 $A$ 在 $R$ 上平坦、非分歧且有限表现。
\end{proof}

\begin{remark}
\label{descent-remark-when-locally-split}
若能找到忠实平坦环同态 $g: R \to T$，使 $T \to S \otimes_R T$ 具有
某种额外结构，事情就会更容易。例如，若能保证
$T \to S \otimes_R T$ 在 $\textit{Rings}$ 中可裂，则可以推出：
凡是在底扩张下稳定且能沿忠实平坦态射下降的模或代数性质，也能沿
普遍单射态射下降。一个显然的猜想是寻找 $g$，使 $T$ 不仅忠实平坦，
而且是 $\text{Mod}_R$ 中的内射对象；但即使对 $R = \mathbf{Z}$，
这样的同态也不可能存在。
\end{remark}
















\section{拟凝聚层的 fpqc 下降}
\label{descent-section-fpqc-descent-quasi-coherent}

\noindent
模的平坦下降的主要应用，是拟凝聚层相对于 fpqc 覆盖的相应下降命题。

\begin{lemma}
\label{descent-lemma-standard-fpqc-covering}
设 $S$ 为仿射概形，并设
$\mathcal{U} = \{f_i : U_i \to S\}_{i = 1, \ldots, n}$ 是 $S$ 的标准
fpqc 覆盖；见 Topologies, Definition
\ref{topologies-definition-standard-fpqc}。拟凝聚层相对于
$\mathcal{U} = \{U_i \to S\}$ 的任意下降数据都是有效的。此外，
从拟凝聚 $\mathcal{O}_S$-模范畴到相对于 $\mathcal{U}$ 的下降数据范畴的
函子是全忠实的。
\end{lemma}

\begin{proof}
这只是用概形语言改述 Proposition \ref{descent-proposition-descent-module}。
首先注意，拟凝聚层相对于 $\mathcal{U}$ 的下降数据 $\xi$，恰好等同于
拟凝聚层相对于覆盖
$\mathcal{U}' = \{\coprod_{i = 1, \ldots, n} U_i \to S\}$ 的下降数据
$\xi'$。此外，$\xi$ 的有效性等同于 $\xi'$ 的有效性。因此可以假设
$n = 1$，即 $\mathcal{U} = \{U \to S\}$，其中 $U$ 与 $S$ 均为仿射概形。
在此情形，下降数据对应于模相对于环同态
$$
\Gamma(S, \mathcal{O})
\longrightarrow
\Gamma(U, \mathcal{O}).
$$
的下降数据。由于 $U \to S$ 满且平坦，此环同态忠实平坦。换言之，
Proposition \ref{descent-proposition-descent-module} 适用，结论成立。
\end{proof}

\begin{proposition}
\label{descent-proposition-fpqc-descent-quasi-coherent}
设 $S$ 为概形，并设 $\mathcal{U} = \{\varphi_i : U_i \to S\}$ 是 fpqc
覆盖；见 Topologies, Definition \ref{topologies-definition-fpqc-covering}。
拟凝聚层相对于 $\mathcal{U} = \{U_i \to S\}$ 的任意下降数据都是有效的。
此外，从拟凝聚 $\mathcal{O}_S$-模范畴到相对于 $\mathcal{U}$ 的下降数据
范畴的函子是全忠实的。
\end{proposition}

\begin{proof}
设 $S = \bigcup_{j \in J} V_j$ 是仿射开覆盖。对 $j, j' \in J$，
以 $V_{jj'} = V_j \cap V_{j'}$ 表示交（它未必仿射）。对开子集
$V \subset S$，记
$\mathcal{U}_V = \{V \times_S U_i \to V\}_{i \in I}$；这是 fpqc 覆盖
（Topologies, Lemma \ref{topologies-lemma-fpqc}）。由 fpqc 覆盖的定义，
对每个 $j \in J$，可找到有限集 $K_j$、映射
$\underline{i} : K_j \to I$，以及仿射开集
$U_{\underline{i}(k), k} \subset U_{\underline{i}(k)}$（$k \in K_j$），
使 $\mathcal{V}_j = \{U_{\underline{i}(k), k} \to V_j\}_{k \in K_j}$
是 $V_j$ 的标准 fpqc 覆盖。当然，$\mathcal{V}_j$ 是
$\mathcal{U}_{V_j}$ 的加细。图示为
$$
\xymatrix{
\mathcal{V}_j \ar[r] \ar@{~>}[d] &
\mathcal{U}_{V_j} \ar[r] \ar@{~>}[d] &
\mathcal{U} \ar@{~>}[d] \\
V_j \ar@{=}[r] & V_j \ar[r] & S
}
$$
其中上方的水平箭头是具有固定靶的态射族之间的态射（见 Sites,
Definition \ref{sites-definition-morphism-coverings}）。

\medskip\noindent
为证明命题，依次证明从拟凝聚层到下降数据的函子忠实、全且本质满。

\medskip\noindent
忠实性。设 $\mathcal{F}$、$\mathcal{G}$ 是 $S$ 上的拟凝聚层，且
$a, b : \mathcal{F} \to \mathcal{G}$ 是 $\mathcal{O}_S$-模同态。
假设对所有 $i$ 都有 $\varphi_i^*(a) = \varphi_i^*(b)$。取 $s \in S$。
则对某个 $i \in I$ 和 $u \in U_i$，有 $s = \varphi_i(u)$。由于
$\mathcal{O}_{S, s} \to \mathcal{O}_{U_i, u}$ 平坦，因而忠实平坦
（Algebra, Lemma \ref{algebra-lemma-local-flat-ff}），可知
$a_s = b_s : \mathcal{F}_s \to \mathcal{G}_s$。故 $a = b$。

\medskip\noindent
全性。设 $\mathcal{F}$、$\mathcal{G}$ 是 $S$ 上的拟凝聚层，并设
$a_i : \varphi_i^*\mathcal{F} \to \varphi_i^*\mathcal{G}$ 是
$\mathcal{O}_{U_i}$-模同态，满足在
% P05-ERR-0024
$U_i \times_U U_j$ 上 $\text{pr}_0^*a_i = \text{pr}_1^*a_j$。
拉回这些态射，得到态射
$$
% P05-ERR-0025
a_k :
\varphi_{i(k)}^*\mathcal{F}|_{U_{\underline{i}(k), k}}
\longrightarrow
\varphi_{i(k)}^*\mathcal{G}|_{U_{\underline{i}(k), k}}
$$
其中 $k \in K_j$，记号同上。此外，Lemma
\ref{descent-lemma-refine-descent-datum} 保证这些态射定义了
$\mathcal{V}_j$ 上两个（典范）下降数据之间的态射。因此，由 Lemma
\ref{descent-lemma-standard-fpqc-covering}，相应地得到唯一态射
$a_j : \mathcal{F}|_{V_j} \to \mathcal{G}|_{V_j}$。为证明
$a_j|_{V_{jj'}} = a_{j'}|_{V_{jj'}}$，利用已经证明的忠实性，并注意：
$a_j$ 与 $a_{j'}$ 都等同于（典范）下降数据态射 $(a_i)_{i \in I}$
在任一同时加细 $\mathcal{V}_{j, V_{jj'}}$ 和
$\mathcal{V}_{j', V_{jj'}}$ 的覆盖上的拉回。例如，覆盖
% P05-ERR-0026
$\mathcal{V}_{jj'} =
\{V_k \times_S V_{k'} \to V_{jj'}\}_{k \in K_j, k' \in K_{j'}}$
即可。

\medskip\noindent
本质满性。设 $\xi = (\mathcal{F}_i, \varphi_{ii'})$ 是拟凝聚层相对于覆盖
$\mathcal{U}$ 的下降数据。拉回此下降数据，得到拟凝聚层相对于 $V_j$
的覆盖 $\mathcal{V}_j$ 的下降数据 $\xi_j$。再次由 Lemma
\ref{descent-lemma-standard-fpqc-covering}，存在 $V_j$ 上的拟凝聚层
$\mathcal{F}_j$，使其相联系的典范下降数据同构于 $\xi_j$。由上文证明的
全性，存在同构
$$
\phi_{jj'} :
\mathcal{F}_j|_{V_{jj'}}
\longrightarrow
\mathcal{F}_{j'}|_{V_{jj'}}
$$
它对应于 $\xi_j$ 与 $\xi_{j'}$ 在 $\mathcal{V}_{jj'}$ 上的拉回之间的
下降数据同构。为证明这些态射 $\phi_{jj'}$ 满足余圈条件，在三重交
$V_{jj'j''}$ 上使用上文证明的忠实性。于是，由 Lemma
\ref{descent-lemma-zariski-descent-effective}，各层 $\mathcal{F}_j$ 如所需地粘合成
拟凝聚层 $\mathcal{F}$。还需验证：与 $\mathcal{F}$ 相联系的、相对于
$\mathcal{U}$ 的典范下降数据同构于最初的下降数据。此验证从略。
\end{proof}









\section{拟凝聚层的 Galois 下降}
\label{descent-section-galois-descent}

\noindent
拟凝聚层的 Galois 下降只是拟凝聚层 fpqc 下降的一个特殊情形。本节说明
如何把 Galois 下降转化为 fpqc 下降，再应用先前的结果得出结论。

\medskip\noindent
设 $k'/k$ 为域扩张。则 $\{\Spec(k') \to \Spec(k)\}$ 是 fpqc 覆盖。
设 $X$ 为 $k$ 上的概形。对 $k$-代数 $A$，令
$X_A = X \times_{\Spec(k)} \Spec(A)$。由 Topologies, Lemma
\ref{topologies-lemma-fpqc}，$\{X_{k'} \to X\}$ 是 fpqc 覆盖。注意
$$
X_{k'} \times_X X_{k'} = X_{k' \otimes_k k'}
\quad\text{且}\quad
X_{k'} \times_X X_{k'} \times_X X_{k'} = X_{k' \otimes_k k' \otimes_k k'}
$$
因此，拟凝聚层相对于 $\{X_{k'} \to X\}$ 的下降数据由如下资料给出：
$X_{k'}$ 上的一个拟凝聚层 $\mathcal{F}$，以及 $X_{k' \otimes_k k'}$
上的同构 $\varphi : \text{pr}_0^*\mathcal{F} \to \text{pr}_1^*\mathcal{F}$，
它在 $X_{k' \otimes_k k' \otimes_k k'}$ 上满足显然的余圈条件。
下面具体说明 Galois 扩张情形中这意味着什么。

\medskip\noindent
设 $k'/k$ 是 Galois 群为 $G = \text{Gal}(k'/k)$ 的有限 Galois 扩张。
则有 $k$-代数同构
$$
k' \otimes_k k' \longrightarrow \prod\nolimits_{\sigma \in G} k',\quad
a \otimes b \longrightarrow \prod a\sigma(b)
$$
并且
$$
k' \otimes_k k' \otimes_k k' \longrightarrow
\prod\nolimits_{(\sigma, \tau) \in G \times G} k',\quad
a \otimes b \otimes c \longrightarrow \prod a\sigma(b)\sigma(\tau(c))
$$
这里选取 $a\sigma(b)\sigma(\tau(c))$ 而非 $a\sigma(b)\tau(c)$，是因为
前者使下述公式更简洁，但这并非必需。给定 $\sigma \in G$，记
$$
f_\sigma = \text{id}_X \times \Spec(\sigma) :
X_{k'} \longrightarrow X_{k'}
$$
请注意，由于 $\Spec(-)$ 是反变函子，有
$f_{\sigma \tau} = f_\tau \circ f_\sigma$，而不是相反的次序。利用上面的
第一个同构，得到等同
$$
X_{k' \otimes_k k'} = \coprod\nolimits_{\sigma \in G} X_{k'}
$$
使得 $\text{pr}_0$ 对应于态射
$$
\coprod\nolimits_{\sigma \in G} X_{k'}
\xrightarrow{\coprod \text{id}}
X_{k'}
$$
而 $\text{pr}_1$ 对应于态射
$$
\coprod\nolimits_{\sigma \in G} X_{k'}
\xrightarrow{\coprod f_\sigma}
X_{k'}
$$
由此可见，$X_{k'}$ 上 $\mathcal{F}$ 的下降数据 $\varphi$ 对应于一族同构
$\varphi_\sigma : \mathcal{F} \to f_\sigma^*\mathcal{F}$。为写出余圈条件，
使用等同
$$
X_{k' \otimes_k k' \otimes_k k'} =
\coprod\nolimits_{(\sigma, \tau) \in G \times G} X_{k'}.
$$
它由上述代数同构得到。在此等同下，态射 $\text{pr}_{01}$ 对应于
$$
\coprod\nolimits_{(\sigma, \tau) \in G \times G} X_{k'}
\longrightarrow
\coprod\nolimits_{\sigma \in G} X_{k'}
$$
该态射通过恒等态射把指标为 $(\sigma, \tau)$ 的分支映到指标为
$\sigma$ 的分支。态射 $\text{pr}_{12}$ 对应于
$$
\coprod\nolimits_{(\sigma, \tau) \in G \times G} X_{k'}
\longrightarrow
\coprod\nolimits_{\sigma \in G} X_{k'}
$$
该态射通过 $f_\sigma$ 把指标为 $(\sigma, \tau)$ 的分支映到指标为
$\tau$ 的分支。最后，态射 $\text{pr}_{02}$ 对应于
$$
\coprod\nolimits_{(\sigma, \tau) \in G \times G} X_{k'}
\longrightarrow
\coprod\nolimits_{\sigma \in G} X_{k'}
$$
该态射通过恒等态射把指标为 $(\sigma, \tau)$ 的分支映到指标为
$\sigma\tau$ 的分支。因此余圈条件
$$
\text{pr}_{02}^*\varphi = \text{pr}_{12}^*\varphi \circ \text{pr}_{01}^*\varphi
$$
转化为每一对 $(\sigma, \tau)$ 所对应的条件
$$
\varphi_{\sigma\tau} = f_\sigma^*\varphi_\tau \circ \varphi_\sigma
$$
这是从 $\mathcal{F} \to f_{\sigma\tau}^*\mathcal{F}$ 的态射等式。
（一切恰好相容；例如，$\varphi_\sigma$ 的靶是
$f_\sigma^*\mathcal{F}$，而 $f_\sigma^*\varphi_\tau$ 的源也正是
$f_\sigma^*\mathcal{F}$。）

\begin{lemma}
\label{descent-lemma-galois-descent}
设 $k'/k$ 是 Galois 群为 $G$ 的（有限）Galois 扩张，$X$ 为 $k$ 上的
概形。拟凝聚 $\mathcal{O}_X$-模范畴等价于系统
$(\mathcal{F}, (\varphi_\sigma)_{\sigma \in G})$ 所成的范畴，其中
\begin{enumerate}
\item $\mathcal{F}$ 是 $X_{k'}$ 上的拟凝聚模；
\item $\varphi_\sigma : \mathcal{F} \to f_\sigma^*\mathcal{F}$ 是模同构；
\item 对所有 $\sigma, \tau \in G$，均有
$\varphi_{\sigma\tau} = f_\sigma^*\varphi_\tau \circ \varphi_\sigma$。
\end{enumerate}
这里 $f_\sigma = \text{id}_X \times \Spec(\sigma) : X_{k'} \to X_{k'}$。
\end{lemma}

\begin{proof}
如上所见，引理中的资料 $(\mathcal{F}, (\varphi_\sigma)_{\sigma \in G})$
等同于 fpqc 覆盖 $\{X_{k'} \to X\}$ 的下降数据。因此引理由 Proposition
\ref{descent-proposition-fpqc-descent-quasi-coherent} 得出。
\end{proof}

\noindent
上述结论有如下稍一般的版本。设有满的有限 \'etale 态射 $X \to Y$、
有限群 $G$ 以及群同态
$G^{opp} \to \text{Aut}_Y(X), \sigma \mapsto f_\sigma$，使态射
$$
G \times X \longrightarrow X \times_Y X,\quad
(\sigma, x) \longmapsto (x, f_\sigma(x))
$$
为同构。则上述相同结论成立。

\begin{lemma}
\label{descent-lemma-galois-descent-more-general}
设 $X \to Y$、$G$ 和 $f_\sigma : X \to X$ 如上。拟凝聚
$\mathcal{O}_Y$-模范畴等价于系统
$(\mathcal{F}, (\varphi_\sigma)_{\sigma \in G})$ 所成的范畴，其中
\begin{enumerate}
\item $\mathcal{F}$ 是拟凝聚 $\mathcal{O}_X$-模；
\item $\varphi_\sigma : \mathcal{F} \to f_\sigma^*\mathcal{F}$ 是模同构；
\item 对所有 $\sigma, \tau \in G$，均有
$\varphi_{\sigma\tau} = f_\sigma^*\varphi_\tau \circ \varphi_\sigma$。
\end{enumerate}
\end{lemma}

\begin{proof}
由于 $X \to Y$ 是满的有限 \'etale 态射，$\{X \to Y\}$ 是 fpqc 覆盖。
又因 $G \times X \to X \times_Y X$，$(\sigma, x) \mapsto (x, f_\sigma(x))$
为同构，可知 $G \times G \times X \to X \times_Y X \times_Y X$，
$(\sigma, \tau, x) \mapsto (x, f_\sigma(x), f_{\sigma\tau}(x))$ 也是同构。
利用这些等同，引理中资料的范畴等同于拟凝聚层相对于覆盖
% P05-ERR-0027
$\{x \to Y\}$ 的下降数据范畴。因此引理由 Proposition
\ref{descent-proposition-fpqc-descent-quasi-coherent} 得出。
\end{proof}









\section{模的有限性性质的下降}
\label{descent-section-descent-finiteness}

\noindent
本节证明：要检验拟凝聚模是否满足某些有限性条件，可以在覆盖的各成员上
进行检验。

\begin{lemma}
\label{descent-lemma-finite-type-descends}
设 $X$ 为概形，$\mathcal{F}$ 为拟凝聚 $\mathcal{O}_X$-模，并设
$\{f_i : X_i \to X\}_{i \in I}$ 是 fpqc 覆盖，使每个
$f_i^*\mathcal{F}$ 都是有限型 $\mathcal{O}_{X_i}$-模。则
$\mathcal{F}$ 是有限型 $\mathcal{O}_X$-模。
\end{lemma}

\begin{proof}
略。仿射情形见 Algebra, Lemma
\ref{algebra-lemma-descend-properties-modules}。
\end{proof}

\begin{lemma}
\label{descent-lemma-finite-type-descends-fppf}
设 $f : (X, \mathcal{O}_X) \to (Y, \mathcal{O}_Y)$ 是局部赋环空间态射，
$\mathcal{F}$ 是 $\mathcal{O}_Y$-模层。若
\begin{enumerate}
\item $f$ 作为拓扑空间映射是开映射；
\item $f$ 满且平坦；
\item $f^*\mathcal{F}$ 有限型，
\end{enumerate}
则 $\mathcal{F}$ 有限型。
\end{lemma}

\begin{proof}
取一点 $y \in Y$，并选取映至 $y$ 的点 $x \in X$。选取开集
$x \in U \subset X$ 以及 $f^*\mathcal{F}(U)$ 中在 $U$ 上生成
$f^*\mathcal{F}$ 的元素 $s_1, \ldots, s_n$。由于
$f^*\mathcal{F} =
f^{-1}\mathcal{F} \otimes_{f^{-1}\mathcal{O}_Y} \mathcal{O}_X$，
缩小 $U$ 后可假设 $s_i = \sum t_{ij} \otimes a_{ij}$，其中
$t_{ij} \in f^{-1}\mathcal{F}(U)$ 且 $a_{ij} \in \mathcal{O}_X(U)$。
进一步缩小 $U$，可假设 $t_{ij}$ 来自某个截面
$s_{ij} \in \mathcal{F}(V)$，其中 $V \subset Y$ 为满足
$f(U) \subset V$ 的开集。令 $N$ 为截面 $s_{ij}$ 的个数，并考虑态射
$$
\sigma = (s_{ij}) : \mathcal{O}_V^{\oplus N} \to \mathcal{F}|_V
$$
由这些截面的选取可知 $f^*\sigma|_U$ 为满射。因此对每个 $u \in U$，
态射
$$
\sigma_{f(u)} \otimes_{\mathcal{O}_{Y, f(u)}} \mathcal{O}_{X, u} :
\mathcal{O}_{X, u}^{\oplus N}
\longrightarrow
\mathcal{F}_{f(u)} \otimes_{\mathcal{O}_{Y, f(u)}} \mathcal{O}_{X, u}
$$
为满射。由于 $f$ 平坦，局部环同态
$\mathcal{O}_{Y, f(u)} \to \mathcal{O}_{X, u}$ 平坦，因而忠实平坦
（Algebra, Lemma \ref{algebra-lemma-local-flat-ff}）。故
$\sigma_{f(u)}$ 为满射。又因 $f$ 是开映射，$f(U)$ 是 $y$ 的开邻域，
证明完成。
\end{proof}

\begin{lemma}
\label{descent-lemma-finite-presentation-descends}
设 $X$ 为概形，$\mathcal{F}$ 为拟凝聚 $\mathcal{O}_X$-模，并设
$\{f_i : X_i \to X\}_{i \in I}$ 是 fpqc 覆盖，使每个
$f_i^*\mathcal{F}$ 都是有限表现 $\mathcal{O}_{X_i}$-模。则
$\mathcal{F}$ 是有限表现 $\mathcal{O}_X$-模。
\end{lemma}

\begin{proof}
略。仿射情形见 Algebra, Lemma
\ref{algebra-lemma-descend-properties-modules}。
\end{proof}

\begin{lemma}
\label{descent-lemma-locally-generated-by-r-sections-descends}
设 $X$ 为概形，$\mathcal{F}$ 为拟凝聚 $\mathcal{O}_X$-模，并设
$\{f_i : X_i \to X\}_{i \in I}$ 是 fpqc 覆盖，使每个
$f_i^*\mathcal{F}$ 作为 $\mathcal{O}_{X_i}$-模都局部由 $r$ 个截面生成。
则 $\mathcal{F}$ 作为 $\mathcal{O}_X$-模局部由 $r$ 个截面生成。
\end{lemma}

\begin{proof}
由 Lemma \ref{descent-lemma-finite-type-descends}，$\mathcal{F}$ 有限型。因此
Nakayama 引理（Algebra, Lemma \ref{algebra-lemma-NAK}）表明：
$\mathcal{F}$ 在一点 $x \in X$ 的某个邻域上由 $r$ 个截面生成，当且仅当
$\dim_{\kappa(x)} \mathcal{F}_x \otimes \kappa(x) \leq r$。选取一个 $i$
以及映至 $x$ 的点 $x_i \in X_i$。则
$\dim_{\kappa(x)} \mathcal{F}_x \otimes \kappa(x) = 
\dim_{\kappa(x_i)} (f_i^*\mathcal{F})_{x_i} \otimes \kappa(x_i)$
由于 $f_i^*\mathcal{F}$ 局部由 $r$ 个截面生成，右端为 $\leq r$。
\end{proof}

\begin{lemma}
\label{descent-lemma-flat-descends}
设 $X$ 为概形，$\mathcal{F}$ 为拟凝聚 $\mathcal{O}_X$-模，并设
$\{f_i : X_i \to X\}_{i \in I}$ 是 fpqc 覆盖，使每个
$f_i^*\mathcal{F}$ 都是平坦 $\mathcal{O}_{X_i}$-模。则 $\mathcal{F}$
是平坦 $\mathcal{O}_X$-模。
\end{lemma}

\begin{proof}
略。仿射情形见 Algebra, Lemma
\ref{algebra-lemma-descend-properties-modules}。
\end{proof}

\begin{lemma}
\label{descent-lemma-finite-locally-free-descends}
设 $X$ 为概形，$\mathcal{F}$ 为拟凝聚 $\mathcal{O}_X$-模，并设
$\{f_i : X_i \to X\}_{i \in I}$ 是 fpqc 覆盖，使每个
$f_i^*\mathcal{F}$ 都是有限局部自由 $\mathcal{O}_{X_i}$-模。则
$\mathcal{F}$ 是有限局部自由 $\mathcal{O}_X$-模。
\end{lemma}

\begin{proof}
这由如下事实得出：拟凝聚层有限局部自由，当且仅当它有限表现且平坦；
见 Algebra, Lemma \ref{algebra-lemma-finite-projective}。具体而言，若每个
$f_i^*\mathcal{F}$ 都平坦且有限表现，则由 Lemmas
\ref{descent-lemma-flat-descends} 和 \ref{descent-lemma-finite-presentation-descends}，
$\mathcal{F}$ 也如此。
\end{proof}

\noindent
局部投射拟凝聚层的定义见 Properties, Section
\ref{properties-section-locally-projective}。

\begin{lemma}
\label{descent-lemma-locally-projective-descends}
设 $X$ 为概形，$\mathcal{F}$ 为拟凝聚 $\mathcal{O}_X$-模，并设
$\{f_i : X_i \to X\}_{i \in I}$ 是 fpqc 覆盖，使每个
$f_i^*\mathcal{F}$ 都是局部投射 $\mathcal{O}_{X_i}$-模。则
$\mathcal{F}$ 是局部投射 $\mathcal{O}_X$-模。
\end{lemma}

\begin{proof}
略。对 Zariski 覆盖，这是 Properties, Lemma
\ref{properties-lemma-locally-projective}；仿射情形则是 Algebra, Theorem
\ref{algebra-theorem-ffdescent-projectivity}。
\end{proof}

\begin{remark}
\label{descent-remark-locally-free-descends}
局部自由性是拟凝聚模的一个不在 fpqc 拓扑中下降的性质。事实上，设
$R$ 为环，$M$ 为投射 $R$-模，它是秩 1 局部自由模的可数直和
$M = \bigoplus L_n$，但自身并非局部自由；见 Examples, Lemma
\ref{examples-lemma-projective-not-locally-free}。作忠实平坦基变换
$$
R \longrightarrow
\bigoplus\nolimits_{m \geq 1}
\bigoplus\nolimits_{(i_1, \ldots, i_m) \in \mathbf{Z}^{\oplus m}}
L_1^{\otimes i_1} \otimes_R \ldots \otimes_R L_m^{\otimes i_m}
$$
后，$M$ 成为自由模。但我们不知道 fppf 覆盖的情形如何。换言之，
我们不知道下述问题的答案：设 $A \to B$ 是有限表现的忠实平坦环同态，
$M$ 是满足 $M \otimes_A B$ 自由的 $A$-模。$M$ 是否为局部自由
$A$-模？事实证明，若 $A$ 是 Noether 环，则答案肯定；这由
\cite{Bass} 的结果得出。但一般情形的答案仍未知。若你知道答案或有参考
文献，请发送邮件至
\href{mailto:stacks.project@gmail.com}{stacks.project@gmail.com}.
\end{remark}

\noindent
这里还补充两个与上述结果相关、但性质略有不同的结果。

\begin{lemma}
\label{descent-lemma-finite-over-finite-module}
设 $f : X \to Y$ 为概形态射，$\mathcal{F}$ 为拟凝聚
$\mathcal{O}_X$-模，并假设 $f$ 是有限态射。则 $\mathcal{F}$ 是有限型
$\mathcal{O}_X$-模，当且仅当 $f_*\mathcal{F}$ 是有限型
$\mathcal{O}_Y$-模。
\end{lemma}

\begin{proof}
由于 $f$ 有限，它是仿射的。这把问题化归到如下情形：$f$ 是由有限环同态
$A \to B$ 给出的态射 $\Spec(B) \to \Spec(A)$。此外，此时
$\mathcal{F} = \widetilde{M}$ 是与 $B$-模 $M$ 相联系的模层。注意，
$M$ 作为 $B$-模有限，当且仅当 $M$ 作为 $A$-模有限；见 Algebra,
Lemma \ref{algebra-lemma-finite-module-over-finite-extension}。结合
Properties, Lemma \ref{properties-lemma-finite-type-module} 即得引理。
\end{proof}

\begin{lemma}
\label{descent-lemma-finite-finitely-presented-module}
设 $f : X \to Y$ 为概形态射，$\mathcal{F}$ 为拟凝聚
$\mathcal{O}_X$-模，并假设 $f$ 有限且有限表现。则 $\mathcal{F}$ 是有限表现
$\mathcal{O}_X$-模，当且仅当 $f_*\mathcal{F}$ 是有限表现
$\mathcal{O}_Y$-模。
\end{lemma}

\begin{proof}
由于 $f$ 有限，它是仿射的。这把问题化归到如下情形：$f$ 是由有限且
有限表现的环同态 $A \to B$ 给出的态射 $\Spec(B) \to \Spec(A)$。
此外，此时 $\mathcal{F} = \widetilde{M}$ 是与 $B$-模 $M$ 相联系的模层。
注意，$M$ 作为 $B$-模有限表现，当且仅当 $M$ 作为 $A$-模有限表现；
见 Algebra, Lemma
\ref{algebra-lemma-finite-finitely-presented-extension}。结合 Properties,
Lemma \ref{properties-lemma-finite-presentation-module} 即得引理。
\end{proof}
















\section{拟凝聚层与拓扑（一）}
\label{descent-section-quasi-coherent-sheaves}

\noindent
本节的结果表明：概形 $S$ 上的拟凝聚模范畴，与拓扑一章中和 $S$
相联系的许多位点上的拟凝聚模范畴之间存在自然等价。

\medskip\noindent
设 $S$ 为概形，$\mathcal{F}$ 为拟凝聚 $\mathcal{O}_S$-模。考虑函子
\begin{equation}
\label{descent-equation-quasi-coherent-presheaf}
(\Sch/S)^{opp} \longrightarrow \textit{Ab},
\quad
(f : T \to S) \longmapsto \Gamma(T, f^*\mathcal{F}).
\end{equation}

\begin{lemma}
\label{descent-lemma-sheaf-condition-holds}
设 $S$ 为概形，$\mathcal{F}$ 为拟凝聚 $\mathcal{O}_S$-模，并设
$\tau \in \{Zariski, \linebreak[0] \etale, \linebreak[0] smooth,
\linebreak[0] syntomic, \linebreak[0] fppf, \linebreak[0] fpqc\}$。
(\ref{descent-equation-quasi-coherent-presheaf}) 中定义的函子，对 $S$ 上任意概形
$T$ 的任意 $\tau$-覆盖 $\{T_i \to T\}_{i \in I}$ 都满足层条件。
\end{lemma}

\begin{proof}
对 $\tau \in \{Zariski, \linebreak[0] \etale, \linebreak[0] smooth,
\linebreak[0] syntomic, \linebreak[0] fppf\}$，任意 $\tau$-覆盖也是 fpqc
覆盖；见 Topologies, Lemmas
\ref{topologies-lemma-zariski-etale},
\ref{topologies-lemma-zariski-etale-smooth},
\ref{topologies-lemma-zariski-etale-smooth-syntomic},
\ref{topologies-lemma-zariski-etale-smooth-syntomic-fppf}，以及
\ref{topologies-lemma-zariski-etale-smooth-syntomic-fppf-fpqc}.
因此只需对 fpqc 覆盖证明结论。设
$\{f_i : T_i \to T\}_{i \in I}$ 是 fpqc 覆盖，且给定
$f : T \to S$。假设有一族截面
$s_i \in \Gamma(T_i , f_i^*f^*\mathcal{F})$，满足
$s_i|_{T_i \times_T T_j} = s_j|_{T_i \times_T T_j}$。需要找出相应截面
$s \in \Gamma(T, f^*\mathcal{F})$。可以把 $s_i$ 重新解释为一族态射
$\varphi_i : f_i^*\mathcal{O}_T = \mathcal{O}_{T_i} \to f_i^*f^*\mathcal{F}$，
它与 $T$ 上拟凝聚层 $\mathcal{O}_T$ 和 $f^*\mathcal{F}$ 相联系的典范
下降数据相容。因此，由 Proposition
\ref{descent-proposition-fpqc-descent-quasi-coherent}，可将它们（唯一地）下降为
态射 $\mathcal{O}_T \to f^*\mathcal{F}$，后者给出所需截面 $s$。
\end{proof}

\noindent
特别地，可以作如下定义。

\begin{definition}
\label{descent-definition-structure-sheaf}
设 $\tau \in \{Zariski, \linebreak[0] \etale, \linebreak[0]
smooth, \linebreak[0] syntomic, \linebreak[0] fppf\}$，$S$ 为概形，
$\Sch_\tau$ 为包含 $S$ 的大位点，且 $\mathcal{F}$ 为拟凝聚
$\mathcal{O}_S$-模。
\begin{enumerate}
\item 大位点 $(\Sch/S)_\tau$ 的{\it 结构层}是环层
$T/S \mapsto \Gamma(T, \mathcal{O}_T)$，记为 $\mathcal{O}$ 或
$\mathcal{O}_S$。
\item 若 $\tau = Zariski$ 或 $\tau = \etale$，则小位点 $S_{Zar}$ 或
$S_\etale$ 的{\it 结构层}是环层 $T/S \mapsto \Gamma(T, \mathcal{O}_T)$，
记为 $\mathcal{O}$ 或 $\mathcal{O}_S$。
\item 大位点 $(\Sch/S)_\tau$ 上{\it 与 $\mathcal{F}$ 相联系的
$\mathcal{O}$-模层}是 $\mathcal{O}$-模层
$(f : T \to S) \mapsto \Gamma(T, f^*\mathcal{F})$，记为
$\mathcal{F}^a$（也常简记为 $\mathcal{F}$）。
\item 若 $\tau = Zariski$ 或 $\tau = \etale$，则小位点 $S_{Zar}$ 或
$S_\etale$ 上{\it 与 $\mathcal{F}$ 相联系的 $\mathcal{O}$-模层}是
$\mathcal{O}$-模层 $(f : T \to S) \mapsto \Gamma(T, f^*\mathcal{F})$，
记为 $\mathcal{F}^a$（也常简记为 $\mathcal{F}$）。
\end{enumerate}
\end{definition}

\noindent
注意，各种情形中都使用同一记号 $\mathcal{F}^a$。这不会真正造成混淆，
因为按定义，给出层 $\mathcal{F}^a$ 的规则不依赖于所考察的位点。

\begin{remark}
\label{descent-remark-Zariski-site-space}
在 Topologies, Lemma \ref{topologies-lemma-Zariski-usual} 中已经看到：
概形 $S$ 的小 Zariski 位点与作为拓扑空间的 $S$ 等价，其意义是相应的
层范畴自然等价。现在 $S_{Zar}$ 还配备了结构层 $\mathcal{O}$，由此可见，
赋环位点 $(S_{Zar}, \mathcal{O})$ 上的模层等同于赋环空间
$(S, \mathcal{O}_S)$ 上的模层。
\end{remark}

\begin{remark}
\label{descent-remark-change-topologies-ringed}
设 $f : T \to S$ 为概形态射。Topologies, Section
\ref{topologies-section-change-topologies} 所列的每个位点态射
$f_{sites}$ 都成为赋环位点态射。具体而言，每个此类位点态射
$f_{sites} : (\Sch/T)_\tau \to (\Sch/S)_{\tau'}$ 或
$f_{sites} : (\Sch/S)_\tau \to S_{\tau'}$ 均由连续函子
$S'/S \mapsto T \times_S S'/S$ 给出。因此，给定 $S'/S$，令
$$
f_{sites}^\sharp :
\mathcal{O}(S'/S)
\longrightarrow
f_{sites, *}\mathcal{O}(S'/S) =
\mathcal{O}(T \times_S S'/T)
$$
为通常的态射
$\text{pr}_{S'}^\sharp : \mathcal{O}(S') \to \mathcal{O}(T \times_S S')$。
类似地，态射
$i_f : \Sh(T_\tau) \to \Sh((\Sch/S)_\tau)$
（其中 $\tau \in \{Zar, \etale\}$；见 Topologies, Lemmas
\ref{topologies-lemma-put-in-T} 和 \ref{topologies-lemma-put-in-T-etale}）
成为赋环拓扑斯态射，因为 $i_f^{-1}\mathcal{O} = \mathcal{O}$。下面列出
若干特殊情形：
\begin{enumerate}
\item 大位点态射 $f_{big} : (\Sch/X)_{fppf} \to (\Sch/Y)_{fppf}$
成为赋环位点态射
$$
(f_{big}, f_{big}^\sharp) :
((\Sch/X)_{fppf}, \mathcal{O}_X)
\longrightarrow
((\Sch/Y)_{fppf}, \mathcal{O}_Y)
$$
如 Modules on Sites, Definition \ref{sites-modules-definition-ringed-site}
所述。大 syntomic、光滑、\'etale 和 Zariski 位点的情形类似。
\item 小位点态射 $f_{small} : X_\etale \to Y_\etale$ 成为赋环位点态射
$$
(f_{small}, f_{small}^\sharp) :
(X_\etale, \mathcal{O}_X)
\longrightarrow
(Y_\etale, \mathcal{O}_Y)
$$
如 Modules on Sites, Definition \ref{sites-modules-definition-ringed-site}
所述。小 Zariski 位点的情形类似。
\end{enumerate}
\end{remark}

\noindent
设 $S$ 为概形。显然，给定（例如）$(\Sch/S)_{Zar}$ 上的
$\mathcal{O}$-模，其到（例如）$(\Sch/S)_{fppf}$ 的拉回正是 fppf 层化。
为说明比较大位点与小位点时的情形，有如下引理。

\begin{lemma}
\label{descent-lemma-compare-sites}
设 $S$ 为概形。以
$$
\begin{matrix}
\text{id}_{\tau, Zar} & : & (\Sch/S)_\tau \to S_{Zar}, &
\tau \in \{Zar, \etale, smooth, syntomic, fppf\} \\
\text{id}_{\tau, \etale} & : &
(\Sch/S)_\tau \to S_\etale, &
\tau \in \{\etale, smooth, syntomic, fppf\} \\
\text{id}_{small, \etale, Zar} & : & S_\etale \to S_{Zar},
\end{matrix}
$$
表示 Remark \ref{descent-remark-change-topologies-ringed} 中的赋环位点态射。
设 $\mathcal{F}$ 为 $\mathcal{O}_S$-模层，并把它视为 $S_{Zar}$ 上的
$\mathcal{O}$-模层。则
\begin{enumerate}
\item $(\text{id}_{\tau, Zar})^*\mathcal{F}$ 是 $(\Sch/S)_\tau$ 上
Zariski 层
$$
(f : T \to S) \longmapsto \Gamma(T, f^*\mathcal{F})
$$
的 $\tau$-层化；
\item $(\text{id}_{small, \etale, Zar})^*\mathcal{F}$ 是 $S_\etale$ 上
Zariski 层
$$
(f : T \to S) \longmapsto \Gamma(T, f^*\mathcal{F})
$$
的 \'etale 层化。
\end{enumerate}
设 $\mathcal{G}$ 为 $S_\etale$ 上的 $\mathcal{O}$-模层。则
\begin{enumerate}
\item[(3)] $(\text{id}_{\tau, \etale})^*\mathcal{G}$ 是 \'etale 层
$$
(f : T \to S) \longmapsto \Gamma(T, f_{small}^*\mathcal{G})
$$
的 $\tau$-层化，其中 $f_{small} : T_\etale \to S_\etale$ 是 Remark
\ref{descent-remark-change-topologies-ringed} 中赋环小 \'etale 位点的态射。
\end{enumerate}
\end{lemma}

\begin{proof}
证明 (1)。首先注意，当 $\tau = Zar$ 时结论成立，因为此时有拓扑斯态射
$i_f : \Sh(T_{Zar}) \to \Sh((\Sch/S)_{Zar})$，并且作为
$T_{Zar} \to S_{Zar}$ 的态射，有
$\text{id}_{\tau, Zar} \circ i_f = f_{small}$；见 Topologies, Lemmas
\ref{topologies-lemma-put-in-T} 和
\ref{topologies-lemma-morphism-big-small}。由于拉回具有传递性（见
Modules on Sites, Lemma
\ref{sites-modules-lemma-push-pull-composition-modules}），可知如所需，
$i_f^*(\text{id}_{\tau, Zar})^*\mathcal{F} = f_{small}^*\mathcal{F}$。
因此，由本引理前的评注，$(\text{id}_{\tau, Zar})^*\mathcal{F}$ 是预层
$T \mapsto \Gamma(T, f^*\mathcal{F})$ 的 $\tau$-层化。

\medskip\noindent
(3) 的证明与 (1) 的证明完全相同，只是使用 Topologies, Lemmas
\ref{topologies-lemma-put-in-T-etale} 和
\ref{topologies-lemma-morphism-big-small-etale}。(2) 的证明从略。
\end{proof}

\begin{remark}
\label{descent-remark-change-topologies-ringed-sites}
Remark \ref{descent-remark-change-topologies-ringed} 和
Lemma \ref{descent-lemma-compare-sites} 有如下应用：
\begin{enumerate}
\item 设 $S$ 为概形。构造 $\mathcal{F} \mapsto \mathcal{F}^a$ 是沿赋环位点态射
$\text{id}_{\tau, Zar} : ((\Sch/S)_\tau, \mathcal{O})
\to (S_{Zar}, \mathcal{O})$
或态射
$\text{id}_{small, \etale, Zar} :
(S_\etale, \mathcal{O}) \to (S_{Zar}, \mathcal{O})$ 的拉回。
\item 设 $f : X \to Y$ 为概形态射。对 Remark
\ref{descent-remark-change-topologies-ringed} 中任一赋环位点态射 $f_{sites}$，都有
$$
(f^*\mathcal{F})^a = f_{sites}^*\mathcal{F}^a.
$$
这由 (1) 以及拉回与赋环位点态射的复合相容这一事实得出；见
Modules on Sites, Lemma
\ref{sites-modules-lemma-push-pull-composition-modules}。
\end{enumerate}
\end{remark}

\begin{lemma}
\label{descent-lemma-quasi-coherent-gives-quasi-coherent}
设 $S$ 为概形，$\mathcal{F}$ 为拟凝聚 $\mathcal{O}_S$-模，并设
$\tau \in \{Zariski, \linebreak[0] \etale, \linebreak[0]
smooth, \linebreak[0] syntomic, \linebreak[0] fppf\}$。
\begin{enumerate}
\item 层 $\mathcal{F}^a$ 是 $(\Sch/S)_\tau$ 上的拟凝聚 $\mathcal{O}$-模，
其定义见 Modules on Sites, Definition
\ref{sites-modules-definition-site-local}。
\item 若 $\tau = Zariski$ 或 $\tau = \etale$，则层 $\mathcal{F}^a$ 是
$S_{Zar}$ 或 $S_\etale$ 上的拟凝聚 $\mathcal{O}$-模，其定义见
Modules on Sites, Definition \ref{sites-modules-definition-site-local}。
\end{enumerate}
\end{lemma}

\begin{proof}
选取 Zariski 覆盖 $\{S_i \to S\}$，使得有正合列
$$
\bigoplus\nolimits_{k \in K_i} \mathcal{O}_{S_i} \longrightarrow
\bigoplus\nolimits_{j \in J_i} \mathcal{O}_{S_i} \longrightarrow
\mathcal{F}|_{S_i} \longrightarrow 0
$$
其中 $K_i$ 和 $J_i$ 是某些指标集。由赋环空间上拟凝聚层的定义，
这样的覆盖存在（见 Modules, Definition
\ref{modules-definition-quasi-coherent}）。

\medskip\noindent
证明 (1)。设 $\tau \in \{Zariski, \linebreak[0] fppf, \linebreak[0]
\etale, \linebreak[0] smooth, \linebreak[0] syntomic\}$。显然，
$\mathcal{F}^a|_{(\Sch/S_i)_\tau}$ 也位于正合列
$$
\bigoplus\nolimits_{k \in K_i} \mathcal{O}|_{(\Sch/S_i)_\tau}
\longrightarrow
\bigoplus\nolimits_{j \in J_i} \mathcal{O}|_{(\Sch/S_i)_\tau}
\longrightarrow
\mathcal{F}^a|_{(\Sch/S_i)_\tau} \longrightarrow 0
$$
中。因此，由 Modules on Sites, Lemma
\ref{sites-modules-lemma-local-final-object}，$\mathcal{F}^a$ 拟凝聚。

\medskip\noindent
证明 (2)。令 $\tau = \etale$。显然，$\mathcal{F}^a|_{(S_i)_\etale}$
也位于正合列
$$
\bigoplus\nolimits_{k \in K_i} \mathcal{O}|_{(S_i)_\etale}
\longrightarrow
\bigoplus\nolimits_{j \in J_i} \mathcal{O}|_{(S_i)_\etale}
\longrightarrow
\mathcal{F}^a|_{(S_i)_\etale} \longrightarrow 0
$$
中。因此，由 Modules on Sites, Lemma
\ref{sites-modules-lemma-local-final-object}，$\mathcal{F}^a$ 拟凝聚。
$\tau = Zariski$ 的情形类似（事实上，由于相应的赋环拓扑斯相同，
这只是同义反复）。
\end{proof}

\begin{lemma}
\label{descent-lemma-fully-faithful-associated}
设 $S$ 为概形，并设
$\tau \in \{Zariski, \linebreak[0] \etale, \linebreak[0]
smooth, \linebreak[0] syntomic, \linebreak[0] fppf\}$。Definition
\ref{descent-definition-structure-sheaf} 中每个函子
$\mathcal{F} \mapsto \mathcal{F}^a$，即
$$
\QCoh(\mathcal{O}_S) \to \QCoh((\Sch/S)_\tau, \mathcal{O})
\quad\text{或}\quad
\QCoh(\mathcal{O}_S) \to \QCoh(S_\tau, \mathcal{O})
$$
都是全忠实的。
\end{lemma}

\begin{proof}
（由 Lemma \ref{descent-lemma-quasi-coherent-gives-quasi-coherent}，确实得到所示的
函子。）我们把 $S$ 上的 $\mathcal{O}_S$-模与
$(S_{Zar}, \mathcal{O}_S)$ 上的模等同。对拟凝聚模 $\mathcal{F}$，函子
$\mathcal{F} \mapsto \mathcal{F}^a$ 由沿某个赋环位点态射 $f$ 的拉回给出；
见 Remark \ref{descent-remark-change-topologies-ringed-sites}。在每种情形中，
函子 $f_*$ 均由沿嵌入函子 $S_{Zar} \to S_\tau$ 或
$S_{Zar} \to (\Sch/S)_\tau$ 的限制给出（这些位点态射的定义方式见
Topologies, Section \ref{topologies-section-change-topologies}）。结合
$f^*\mathcal{F} = \mathcal{F}^a$ 的描述可知，只要 $\mathcal{F}$ 拟凝聚，
就有 $f_*f^*\mathcal{F} = \mathcal{F}$。于是
$$
\Hom_\mathcal{O}(\mathcal{F}^a, \mathcal{G}^a) =
\Hom_\mathcal{O}(f^*\mathcal{F}, f^*\mathcal{G}) =
\Hom_{\mathcal{O}_S}(\mathcal{F}, f_*f^*\mathcal{G}) =
\Hom_{\mathcal{O}_S}(\mathcal{F}, \mathcal{G})
$$
如所需成立。
\end{proof}

\begin{proposition}
\label{descent-proposition-equivalence-quasi-coherent}
设 $S$ 为概形，并设
$\tau \in \{Zariski, \linebreak[0] \etale, \linebreak[0]
smooth, \linebreak[0] syntomic, \linebreak[0] fppf\}$。
\begin{enumerate}
\item 函子 $\mathcal{F} \mapsto \mathcal{F}^a$ 定义范畴等价
$$
\QCoh(\mathcal{O}_S)
\longrightarrow
\QCoh((\Sch/S)_\tau, \mathcal{O})
$$
它连接 $S$ 上的拟凝聚层范畴与 $S$ 的大 $\tau$-位点上的拟凝聚
$\mathcal{O}$-模范畴。
\item 令 $\tau = Zariski$ 或 $\tau = \etale$。函子
$\mathcal{F} \mapsto \mathcal{F}^a$ 定义范畴等价
$$
\QCoh(\mathcal{O}_S)
\longrightarrow
\QCoh(S_\tau, \mathcal{O})
$$
它连接 $S$ 上的拟凝聚层范畴与 $S$ 的小 $\tau$-位点上的拟凝聚
$\mathcal{O}$-模范畴。
\end{enumerate}
\end{proposition}

\begin{proof}
Lemma \ref{descent-lemma-quasi-coherent-gives-quasi-coherent} 已说明该函子有良好定义；
Lemma \ref{descent-lemma-fully-faithful-associated} 已说明它全忠实。为完成证明，
将证明 $(\Sch/S)_\tau$ 上的拟凝聚 $\mathcal{O}$-模给出拟凝聚层相对于
$S$ 的某个 $\tau$-覆盖的下降数据。构造出此下降数据后，应用 Proposition
\ref{descent-proposition-fpqc-descent-quasi-coherent}，得到 $S$ 上相应的拟凝聚层。

\medskip\noindent
设 $\mathcal{G}$ 是 $S$ 的大 $\tau$-位点上的拟凝聚 $\mathcal{O}$-模。
由 Modules on Sites, Definition
\ref{sites-modules-definition-site-local}，存在 $S$ 的 $\tau$-覆盖
$\{S_i \to S\}_{i \in I}$，使每个限制
$\mathcal{G}|_{(\Sch/S_i)_\tau}$ 都有整体表现
$$
\bigoplus\nolimits_{k \in K_i} \mathcal{O}|_{(\Sch/S_i)_\tau}
\longrightarrow
\bigoplus\nolimits_{j \in J_i} \mathcal{O}|_{(\Sch/S_i)_\tau}
\longrightarrow
\mathcal{G}|_{(\Sch/S_i)_\tau} \longrightarrow 0
$$
其中 $J_i$ 和 $K_i$ 是某些指标集。我们断言，这蕴含
$\mathcal{G}|_{(\Sch/S_i)_\tau}$ 等于 $S_i$ 上某个拟凝聚层
$\mathcal{F}_i$ 所对应的 $\mathcal{F}_i^a$。事实上，这对直和
$\bigoplus\nolimits_{k \in K_i} \mathcal{O}|_{(\Sch/S_i)_\tau}$
以及
$\bigoplus\nolimits_{j \in J_i} \mathcal{O}|_{(\Sch/S_i)_\tau}$.
是显然的。因此，$\mathcal{G}|_{(\Sch/S_i)_\tau}$ 是某个态射
$\varphi : \mathcal{K}_i^a \to \mathcal{L}_i^a$ 的余核，其中
$\mathcal{K}_i$、$\mathcal{L}_i$ 是 $S_i$ 上的拟凝聚层。由
$(\ )^a$ 的全忠实性，$\varphi = \phi^a$，其中
$\phi : \mathcal{K}_i \to \mathcal{L}_i$ 是 $S_i$ 上的拟凝聚层态射。
于是显然
$\mathcal{G}|_{(\Sch/S_i)_\tau} \cong \Coker(\phi)^a$
如所断言。

\medskip\noindent
由于 $\mathcal{G}$ 定义在整个范畴 $(\Sch/S)_\tau$ 上，可知
$$
% P05-ERR-0028
(\text{pr}_0^*\mathcal{F}_i)^a
\cong
\mathcal{G}|_{(\Sch/(S_i \times_S S_j))_\tau}
\cong
(\text{pr}_1^*\mathcal{F})^a
$$
是 $(\Sch/(S_i \times_S S_j))_\tau$ 上的 $\mathcal{O}$-模同构。
因此，再次利用全忠实性，得到典范同构
$$
\phi_{ij} :
\text{pr}_0^*\mathcal{F}_i
\longrightarrow
\text{pr}_1^*\mathcal{F}_j
$$
它们是 $S_i \times_S S_j$ 上的拟凝聚模同构。省略它们满足余圈条件的验证。
既然它们确实满足该条件，由拟凝聚层相对于覆盖 $\{S_i \to S\}$ 的下降
有效性（Proposition \ref{descent-proposition-fpqc-descent-quasi-coherent}），
存在 $S$ 上的拟凝聚层 $\mathcal{F}$，使
$\mathcal{F}|_{S_i} \cong \mathcal{F}_i$ 与给定下降数据相容。换言之，
给定 $\mathcal{O}$-模同构
$$
\phi_i :
\mathcal{F}^a|_{(\Sch/S_i)_\tau}
\longrightarrow
\mathcal{G}|_{(\Sch/S_i)_\tau}
$$
它们在 $S_i \times_S S_j$ 上相符。由于
$\SheafHom_\mathcal{O}(\mathcal{F}^a, \mathcal{G})$ 是层（Modules on Sites,
Lemma \ref{sites-modules-lemma-internal-hom}），可知存在恢复上述同构
$\phi_i$ 的 $\mathcal{O}$-模态射 $\mathcal{F}^a \to \mathcal{G}$。
因此该态射是同构，结论成立。

\medskip\noindent
位点 $S_\etale$ 和 $S_{Zar}$ 的情形以完全相同的方式证明。
\end{proof}

\begin{lemma}
\label{descent-lemma-equivalence-quasi-coherent-properties}
设 $S$ 为概形，并设
$\tau \in \{Zariski, \linebreak[0] \etale, \linebreak[0]
smooth, \linebreak[0] syntomic, \linebreak[0] fppf\}$。令 $\mathcal{P}$
为 Modules on Sites, Definitions \ref{sites-modules-definition-global}、
\ref{sites-modules-definition-site-local} 和
\ref{sites-modules-definition-flat} 中定义的模的性质之一\footnote{列表如下：
自由、有限自由、由整体截面生成、由 $r$ 个整体截面生成、由有限多个整体
截面生成、具有整体表现、具有整体有限表现、局部自由、有限局部自由、
局部由截面生成、局部由 $r$ 个截面生成、有限型、有限表现、凝聚或平坦。}。
Proposition \ref{descent-proposition-equivalence-quasi-coherent} 中由规则
$\mathcal{F} \mapsto \mathcal{F}^a$ 定义的范畴等价
$$
\QCoh(\mathcal{O}_S)
\longrightarrow
\QCoh((\Sch/S)_\tau, \mathcal{O})
\quad\text{且}\quad
\QCoh(\mathcal{O}_S)
\longrightarrow
\QCoh(S_\tau, \mathcal{O})
$$
具有性质
$$
\mathcal{F}\text{ 具有 }\mathcal{P}
\Leftrightarrow
\mathcal{F}^a\text{ 具有 }\mathcal{P}\text{ 作为一个 }\mathcal{O}\text{-模}
$$
但 $\mathcal{P}$ 为“局部自由”或“凝聚”时可能例外。若
$\mathcal{P}=$“凝聚”，且 $S$ 局部 Noether、$\tau$ 为 Zariski 或
\'etale，则等价 $\QCoh(\mathcal{O}_S) \to \QCoh(S_\tau, \mathcal{O})$
满足上述性质。
\end{lemma}

\begin{proof}
对整体性质，即 Modules on Sites, Definition
\ref{sites-modules-definition-global} 中定义的性质，结论是直接的。
对局部性质，可以利用 Modules on Sites, Lemma
\ref{sites-modules-lemma-local-final-object}，把“$\mathcal{F}^a$ 具有
$\mathcal{P}$”转化为 $X$ 的某个覆盖的各成员上的性质。因此，结论由
Lemmas \ref{descent-lemma-finite-type-descends}、
\ref{descent-lemma-finite-presentation-descends}、
\ref{descent-lemma-locally-generated-by-r-sections-descends}、
\ref{descent-lemma-flat-descends} 和 \ref{descent-lemma-finite-locally-free-descends} 得出。
在局部 Noether 概形上，拟凝聚模凝聚等价于它有限型（见 Cohomology of
Schemes, Lemma \ref{coherent-lemma-coherent-Noetherian}），故这化归到
有限型模的情形（细节从略）。
\end{proof}










\section{拟凝聚模的上同调与拓扑}
\label{descent-section-quasi-coherent-cohomology}

\noindent
本节证明拟凝聚模的上同调不依赖于拓扑的选择。

\begin{lemma}
\label{descent-lemma-standard-covering-Cech}
设 $S$ 为概形。设
\begin{enumerate}
\item[(a)] $\tau \in \{Zariski, \linebreak[0] fppf, \linebreak[0]
\etale, \linebreak[0] smooth, \linebreak[0] syntomic\}$
且 $\mathcal{C} = (\Sch/S)_\tau$；或
\item[(b)] $\tau = \etale$ 且 $\mathcal{C} = S_\etale$；或
\item[(c)] $\tau = Zariski$ 且 $\mathcal{C} = S_{Zar}$。
\end{enumerate}
设 $\mathcal{F}$ 为 $\mathcal{C}$ 上的 Abel 层，
$U \in \Ob(\mathcal{C})$ 为仿射对象，并设
$\mathcal{U} = \{U_i \to U\}_{i = 1, \ldots, n}$ 是 $\mathcal{C}$ 中的
标准仿射 $\tau$-覆盖。则
\begin{enumerate}
\item $\mathcal{V} = \{\coprod_{i = 1, \ldots, n} U_i \to U\}$ 是 $U$ 的
$\tau$-覆盖；
\item $\mathcal{U}$ 是 $\mathcal{V}$ 的加细；
\item 诱导的 {\v C}ech 复形态射（Cohomology on Sites,
Equation (\ref{sites-cohomology-equation-map-cech-complexes})）
$$
\check{\mathcal{C}}^\bullet(\mathcal{V}, \mathcal{F})
\longrightarrow
\check{\mathcal{C}}^\bullet(\mathcal{U}, \mathcal{F})
$$
是复形同构。
\end{enumerate}
\end{lemma}

\begin{proof}
这是因为
$$
(\coprod\nolimits_{i_0 = 1, \ldots, n} U_{i_0}) \times_U
\ldots \times_U
(\coprod\nolimits_{i_p = 1, \ldots, n} U_{i_p})
=
\coprod\nolimits_{i_0, \ldots, i_p \in \{1, \ldots, n\}}
U_{i_0} \times_U \ldots \times_U U_{i_p}
$$
并且由于不交并是 $\tau$-覆盖，有
$\mathcal{F}(\coprod_a V_a) = \prod_a \mathcal{F}(V_a)$。
\end{proof}

\begin{lemma}
\label{descent-lemma-standard-covering-Cech-quasi-coherent}
设 $S$ 为概形，$\mathcal{F}$ 为 $S$ 上的拟凝聚层，并设
$\tau$、$\mathcal{C}$、$U$、$\mathcal{U}$ 如 Lemma
\ref{descent-lemma-standard-covering-Cech}。则有复形同构
$$
\check{\mathcal{C}}^\bullet(\mathcal{U}, \mathcal{F}^a)
\cong
s((A/R)_\bullet \otimes_R M)
$$
（见 Section \ref{descent-section-descent-modules}），其中
$R = \Gamma(U, \mathcal{O}_U)$、$M = \Gamma(U, \mathcal{F}^a)$，且
$R \to A$ 是忠实平坦环同态。特别地，
$$
\check{H}^p(\mathcal{U}, \mathcal{F}^a) = 0
$$
对所有 $p \geq 1$ 均成立。
\end{lemma}

\begin{proof}
由 Lemma \ref{descent-lemma-standard-covering-Cech}，
$\check{\mathcal{C}}^\bullet(\mathcal{U}, \mathcal{F}^a)$ 同构于
$\check{\mathcal{C}}^\bullet(\mathcal{V}, \mathcal{F}^a)$，其中
$\mathcal{V} = \{V \to U\}$ 且
$V = \coprod_{i = 1, \ldots n} U_i$ 也仿射。令
$A = \Gamma(V, \mathcal{O}_V)$。由于 $\{V \to U\}$ 是
$\tau$-覆盖，$R \to A$ 忠实平坦。另一方面，由 $\mathcal{F}^a$ 的定义，
次数 $p$ 项 $\check{\mathcal{C}}^p(\mathcal{V}, \mathcal{F}^a)$ 为
$$
\Gamma(V \times_U \ldots \times_U V, \mathcal{F}^a)
=
\Gamma(\Spec(A \otimes_R \ldots \otimes_R A), \mathcal{F}^a)
=
A \otimes_R \ldots \otimes_R A \otimes_R M
$$
省略验证 {\v C}ech 复形中的态射与复形
$s((A/R)_\bullet \otimes_R M)$ 中的态射一致。上同调的消失由
Lemma \ref{descent-lemma-ff-exact} 给出。
\end{proof}

\begin{proposition}
\label{descent-proposition-same-cohomology-quasi-coherent}
\begin{slogan}
无论使用哪种拓扑，拟凝聚层的上同调都相同。
\end{slogan}
设 $S$ 为概形，$\mathcal{F}$ 为 $S$ 上的拟凝聚层，并设
$\tau \in \{Zariski, \linebreak[0] \etale, \linebreak[0]
smooth, \linebreak[0] syntomic, \linebreak[0] fppf\}$。
\begin{enumerate}
\item 有典范同构
$$
H^q(S, \mathcal{F}) = H^q((\Sch/S)_\tau, \mathcal{F}^a).
$$
\item 有典范同构
$$
H^q(S, \mathcal{F}) =
H^q(S_{Zar}, \mathcal{F}^a) =
H^q(S_\etale, \mathcal{F}^a).
$$
\end{enumerate}
\end{proposition}

\begin{proof}
当 $q = 0$ 时，结论由 $\mathcal{F}^a$ 的定义立即得出。令
$\mathcal{C} = (\Sch/S)_\tau$，或 $\mathcal{C} = S_\etale$，
或 $\mathcal{C} = S_{Zar}$。

\medskip\noindent
应用 Cohomology on Sites, Lemma
\ref{sites-cohomology-lemma-cech-vanish-collection}，其中取
$\mathcal{F} = \mathcal{F}^a$，取
$\mathcal{B} \subset \Ob(\mathcal{C})$ 为 $\mathcal{C}$ 中仿射概形的集合，
并取 $\text{Cov} \subset \text{Cov}_\mathcal{C}$ 为标准仿射
$\tau$-覆盖的集合。该引理的假设 (3) 由 Lemma
\ref{descent-lemma-standard-covering-Cech-quasi-coherent} 满足。因此，对
$\mathcal{C}$ 的每个仿射对象 $U$，有 $H^p(U, \mathcal{F}^a) = 0$。

\medskip\noindent
接着，令 $U \in \Ob(\mathcal{C})$ 为任意分离对象，以
$f : U \to S$ 表示结构态射，并设 $U = \bigcup U_i$ 是仿射开覆盖。
也可把它视为 $\mathcal{C}$ 中 $U$ 的 $\tau$-覆盖
$\mathcal{U} = \{U_i \to U\}$。注意
$U_{i_0} \times_U \ldots \times_U U_{i_p} =
U_{i_0} \cap \ldots \cap U_{i_p}$ 是仿射的，因为假设了 $U$ 分离。
由 Cohomology on Sites, Lemma
\ref{sites-cohomology-lemma-cech-spectral-sequence-application} 及上述结果，
可知
$$
H^p(U, \mathcal{F}^a) = \check{H}^p(\mathcal{U}, \mathcal{F}^a)
= H^p(U, f^*\mathcal{F})
$$
其中最后一个等号由 Cohomology of Schemes, Lemma
\ref{coherent-lemma-cech-cohomology-quasi-coherent} 给出。特别地，若
$S$ 分离，可取 $U = S$ 且 $f = \text{id}_S$，命题即得证。建议读者
跳过证明的其余部分（或者重写它以给出更清楚的论述）。

\medskip\noindent
在 $S$ 上选取内射分解 $\mathcal{F} \to \mathcal{I}^\bullet$，并在
$\mathcal{C}$ 上选取内射分解 $\mathcal{F}^a \to \mathcal{J}^\bullet$。
记 $\mathcal{J}^n|_S$ 为 $\mathcal{J}^n$ 在 $S$ 的开集上的限制；由于开覆盖
也是 $\tau$-覆盖，它是拓扑空间 $S$ 上的一个层。由此得到复形
$$
0 \to \mathcal{F} \to \mathcal{J}^0|_S \to \mathcal{J}^1|_S \to \ldots
$$
它是正合的，因为它在任意仿射开集 $U \subset S$ 上的截面复形都是正合的
（这是由上面得到的 $H^p(U, \mathcal{F}^a)$ 在 $p > 0$ 时消失而知）。
因此，由 Derived Categories, Lemma
\ref{derived-lemma-morphisms-lift}，存在复形态射
$\mathcal{J}^\bullet|_S \to \mathcal{I}^\bullet$；特别地，它诱导态射
$$
R\Gamma(\mathcal{C}, \mathcal{F}^a)
=
\Gamma(S, \mathcal{J}^\bullet)
\longrightarrow
\Gamma(S, \mathcal{I}^\bullet)
=
R\Gamma(S, \mathcal{F}).
$$
取上同调便得到态射
$H^n(\mathcal{C}, \mathcal{F}^a) \to H^n(S, \mathcal{F})$，我们须证明
它是同构。令 $\mathcal{U} : S = \bigcup U_i$ 为一个仿射开覆盖；也可将其
视为一个 $\tau$-覆盖。由上述构造，得到双复形的态射
$$
\check{\mathcal{C}}^\bullet(\mathcal{U}, \mathcal{J})
=
\check{\mathcal{C}}^\bullet(\mathcal{U}, \mathcal{J}|_S)
\longrightarrow
\check{\mathcal{C}}^\bullet(\mathcal{U}, \mathcal{I}).
$$
它诱导谱序列之间的态射
$$
{}^\tau\! E_2^{p, q} = \check{H}^p(\mathcal{U}, \underline{H}^q(\mathcal{F}^a))
\longrightarrow
E_2^{p, q} = \check{H}^p(\mathcal{U}, \underline{H}^q(\mathcal{F}))
$$
% P05-ERR-0029
第一个谱序列收敛到 $H^{p + q}(\mathcal{C}, \mathcal{F})$，第二个收敛到
$H^{p + q}(S, \mathcal{F})$。另一方面，我们已经看到，所诱导的态射
${}^\tau\! E_2^{p, q} \to E_2^{p, q}$ 都是双射（因为所有交集都是仿射
概形中的开集，因而是分离的）。所以态射
$H^n(\mathcal{C}, \mathcal{F}^a) \to H^n(S, \mathcal{F})$ 也都是同构，
命题得证。
\end{proof}

\begin{proposition}
\label{descent-proposition-equivalence-quasi-coherent-functorial}
设 $f : T \to S$ 为概形态射。
\begin{enumerate}
\item Proposition \ref{descent-proposition-equivalence-quasi-coherent} 中的范畴等价
与拉回相容。更确切地说，对 $S$ 上的任意拟凝聚层 $\mathcal{G}$，有
$f^*(\mathcal{G}^a) = (f^*\mathcal{G})^a$。
\item Proposition \ref{descent-proposition-equivalence-quasi-coherent} 的第 (1) 部分
中的范畴等价一般{\bf 不}与前推相容。
\item 若 $f$ 拟紧且拟分离，并且 $\tau \in \{Zariski, \etale\}$，则
$f_*$ 和 $f_{small, *}$ 保持拟凝聚层，而且图表
$$
\xymatrix{
\QCoh(\mathcal{O}_T)
\ar[rr]_{f_*} \ar[d]_{\mathcal{F} \mapsto \mathcal{F}^a} & &
\QCoh(\mathcal{O}_S)
\ar[d]^{\mathcal{G} \mapsto \mathcal{G}^a} \\
\QCoh(T_\tau, \mathcal{O}) \ar[rr]^{f_{small, *}} & &
\QCoh(S_\tau, \mathcal{O})
}
$$
交换；也就是说，$f_{small, *}(\mathcal{F}^a) = (f_*\mathcal{F})^a$。
\end{enumerate}
\end{proposition}

\begin{proof}
第 (1) 部分由 Remark \ref{descent-remark-change-topologies-ringed-sites} 中的讨论
得出。第 (2) 部分只是一个警告，可以解释如下：首先，这个陈述无法作出
精确表述，因为一般而言，$f_*$ 不把拟凝聚层变为拟凝聚层。即使对 $f$
（以及 $f$ 的任意基变换）确实如此，大位点上的相容性也会意味着
$f_*\mathcal{F}$ 的形成与任意基变换交换，而这一般并不成立。一个具体
例子是拟紧开浸入
$j : X = \mathbf{A}^2_k \setminus \{0\} \to \mathbf{A}^2_k = Y$
，其中 $k$ 是一个域。我们有 $j_*\mathcal{O}_X = \mathcal{O}_Y$，但沿着
$0$ 态射基变换到 $\Spec(k)$ 后，可见此前推为零。

\medskip\noindent
下面在 $\tau = \etale$ 的情形证明 (3)。注意，$f$ 以及 $f$ 的任意基变换
都把拟凝聚层变为拟凝聚层；见 Schemes, Lemma
\ref{schemes-lemma-push-forward-quasi-coherent}。等式
$f_{small, *}(\mathcal{F}^a) = (f_*\mathcal{F})^a$ 意味着，对任意 \'etale
态射 $g : U \to S$，都有
$\Gamma(U, g^*f_*\mathcal{F}) = \Gamma(U \times_S T, (g')^*\mathcal{F})$
，其中 $g' : U \times_S T \to T$ 为投影。此式由 Cohomology of Schemes,
Lemma \ref{coherent-lemma-flat-base-change-cohomology} 成立。
\end{proof}

\begin{lemma}
\label{descent-lemma-higher-direct-images-small-etale}
设 $f : T \to S$ 为拟紧且拟分离的概形态射，并设 $\mathcal{F}$ 为 $T$ 上
的拟凝聚层。对于 \'etale 拓扑或 Zariski 拓扑，都有典范同构
$R^if_{small, *}(\mathcal{F}^a) = (R^if_*\mathcal{F})^a$.
\end{lemma}

\begin{proof}
我们对 \'etale 拓扑证明此结论；省略 Zariski 拓扑情形的证明。由
Cohomology of Schemes, Lemma
\ref{coherent-lemma-quasi-coherence-higher-direct-images}，层
$R^if_*\mathcal{F}$ 是拟凝聚的，故该断言有意义。层
$R^if_{small, *}\mathcal{F}^a$ 是与下列预层相关联的层：
$$
U \longmapsto H^i(U \times_S T, \mathcal{F}^a)
$$
其中 $g : U \to S$ 是 $S_\etale$ 的一个对象；见 Cohomology on Sites,
Lemma \ref{sites-cohomology-lemma-higher-direct-images}。按我们的约定，
右端是 $\mathcal{F}^a$ 在局部化 $T_\etale/U \times_S T$ 上的限制的
\'etale 上同调，而该局部化等于 $(U \times_S T)_\etale$。由
Proposition \ref{descent-proposition-same-cohomology-quasi-coherent}，这个预层与
下列预层相同：
$$
U \longmapsto
H^i(U \times_S T, (g')^*\mathcal{F}),
$$
其中 $g' : U \times_S T \to T$ 为投影。若 $U$ 仿射，则由 Cohomology of
Schemes, Lemma
\ref{coherent-lemma-quasi-coherence-higher-direct-images-application}，
它与 $H^0(U, R^if'_*(g')^*\mathcal{F})$ 相同。再由 Cohomology of
Schemes, Lemma \ref{coherent-lemma-flat-base-change-cohomology}，这等于
$H^0(U, g^*R^if_*\mathcal{F})$，即 $(R^if_*\mathcal{F})^a$ 在 $U$ 上的值。
因此，模层 $R^if_{small, *}(\mathcal{F}^a)$ 与
$(R^if_*\mathcal{F})^a$ 在 $S_\etale$ 的每个仿射对象上的值典范同构，
从而这两个层典范同构。
\end{proof}




\section{拟凝聚层与拓扑，II}
\label{descent-section-quasi-coherent-sheaves-bis}

\noindent
我们继续拓扑章中的讨论，比较概形 $S$ 上的拟凝聚模与关联于 $S$ 的任一
位点上的拟凝聚模。


\begin{lemma}
\label{descent-lemma-compare-etale-zariski-flat}
Lemma \ref{descent-lemma-compare-sites} 中的赋环位点态射
$\text{id}_{small, \etale, Zar} : S_\etale \to S_{Zar}$ 是平坦的。
\end{lemma}

\begin{proof}
记 $\epsilon = \text{id}_{small, \etale, Zar}$，并以
$\mathcal{O}_\etale$ 和 $\mathcal{O}_{Zar}$ 分别表示 $S_\etale$ 与
$S_{Zar}$ 上的结构层。我们须证明 $\mathcal{O}_\etale$ 是平坦的
$\epsilon^{-1}\mathcal{O}_{Zar}$-模。回忆 \'etale 态射是开映射；见
Morphisms, Lemma \ref{morphisms-lemma-etale-open}。由此（以及层的拉回
构造）可知，$\epsilon^{-1}\mathcal{O}_{Zar}$ 是 $S_\etale$ 上预层
$\mathcal{O}'$ 的层化；该预层把 \'etale 态射 $f : V \to S$ 映到
$\mathcal{O}_S(f(V))$。若 $V$ 与 $U = f(V) \subset S$ 都是仿射的，
则 $V \to U$ 是仿射概形之间的 \'etale 态射，因而对应于一个 \'etale
环同态。由于 \'etale 环同态平坦，可见
$\mathcal{O}_S(U) = \mathcal{O}'(V) \to
\mathcal{O}_\etale(V) = \mathcal{O}_V(V)$ 是平坦的。最后，对于每个
\'etale 态射 $f : V \to S$，即 $S_\etale$ 的每个对象，都存在仿射开覆盖
$V = \bigcup V_i$，使得对所有 $i$，$f(V_i)$ 都是 $S$ 中的仿射开集
\footnote{事实上，对 $y \in V$，选取仿射开集 $y \in V' \subset V$，
使 $f(V')$ 包含于某个仿射开集 $U \subset S$。再选取仿射开集
$f(y) \in U' \subset f(V')$。于是 $V'' = f^{-1}(U') \subset V'$ 是仿射
的，因为它等于 $U' \times_U V'$；而 $f(V'') = U'$ 也仿射。}。
结论遂由 Modules on Sites, Lemma
\ref{sites-modules-lemma-flatness-sheafification-refined} 得出。
\end{proof}

\begin{lemma}
\label{descent-lemma-equivalence-quasi-coherent-limits}
设 $S$ 为概形，并设 $\tau \in \{Zariski, \linebreak[0] \etale,
\linebreak[0] smooth, \linebreak[0] syntomic, \linebreak[0] fppf\}$。函子
$$
\QCoh(\mathcal{O}_S)
\longrightarrow
\textit{Mod}((\Sch/S)_\tau, \mathcal{O})
\quad\text{且}\quad
\QCoh(\mathcal{O}_S)
\longrightarrow
\textit{Mod}(S_\tau, \mathcal{O})
$$
由 Proposition \ref{descent-proposition-equivalence-quasi-coherent} 中的规则
$\mathcal{F} \mapsto \mathcal{F}^a$ 定义，并且具有下列性质：
\begin{enumerate}
\item 全忠实；
\item 与直和交换；
\item 与余极限交换；
\item 右正合；
\item 作为函子
$\QCoh(\mathcal{O}_S) \to \textit{Mod}(S_\etale, \mathcal{O})$,
是正合的；
\item 作为函子
$\QCoh(\mathcal{O}_S) \to
\textit{Mod}((\Sch/S)_\tau, \mathcal{O})$
一般{\bf 不}正合；
\item 给定两个拟凝聚 $\mathcal{O}_S$-模
$\mathcal{F}$、$\mathcal{G}$，有
$(\mathcal{F} \otimes_{\mathcal{O}_S} \mathcal{G})^a =
\mathcal{F}^a \otimes_\mathcal{O} \mathcal{G}^a$,
\item 若 $\tau = \etale$ 或 $\tau = Zariski$，并给定两个拟凝聚
$\mathcal{O}_S$-模 $\mathcal{F}$、$\mathcal{G}$，其中 $\mathcal{F}$
是有限表示的，则在 $\textit{Mod}(S_\tau, \mathcal{O})$ 中有
$(\SheafHom_{\mathcal{O}_S}(\mathcal{F}, \mathcal{G}))^a =
\SheafHom_\mathcal{O}(\mathcal{F}^a, \mathcal{G}^a)$；
\item 给定两个拟凝聚 $\mathcal{O}_S$-模 $\mathcal{F}$、$\mathcal{G}$，
一般在 $\textit{Mod}((\Sch/S)_\tau, \mathcal{O})$ 中{\bf 没有}
$(\SheafHom_{\mathcal{O}_S}(\mathcal{F}, \mathcal{G}))^a =
\SheafHom_\mathcal{O}(\mathcal{F}^a, \mathcal{G}^a)$
，即使 $\mathcal{F}$ 是有限表示的也如此；
\item 给定 $\mathcal{O}$-模的短正合列
$0 \to \mathcal{F}_1^a \to \mathcal{E} \to \mathcal{F}_2^a \to 0$
，则 $\mathcal{E}$ 是拟凝聚的\footnote{警告：这个说法容易引起误解。
见第 (6) 部分。}；也就是说，$\mathcal{E}$ 属于该函子的本质像。
\end{enumerate}
\end{lemma}

\begin{proof}
第 (1) 部分已见于 Proposition
\ref{descent-proposition-equivalence-quasi-coherent}。

\medskip\noindent
我们在 Schemes, Section \ref{schemes-section-quasi-coherent} 中已经看到，
概形上拟凝聚层的余极限仍为拟凝聚层。此外，Remark
\ref{descent-remark-change-topologies-ringed-sites} 表明
$\mathcal{F} \mapsto \mathcal{F}^a$ 是某个赋环位点态射的拉回函子，因而
与所有余极限交换；见 Modules on Sites, Lemma
\ref{sites-modules-lemma-exactness-pushforward-pullback}。所以 (3) 成立，
其特殊情形 (2) 也成立。

\medskip\noindent
这也表明该函子右正合（即与有限余极限交换），从而得到 (4)。

\medskip\noindent
函子 $\QCoh(\mathcal{O}_S) \to
\textit{Mod}(S_\etale, \mathcal{O})$，
$\mathcal{F} \mapsto \mathcal{F}^a$，是左正合的，因为 \'etale 态射
平坦；见 Morphisms, Lemma \ref{morphisms-lemma-etale-flat}。这证明了 (5)。

\medskip\noindent
为说明 (6)，设 $S = \Spec(\mathbf{Z})$。此时
$2 : \mathcal{O}_S \to \mathcal{O}_S$ 是单射，但在 $(\Sch/S)_\tau$ 上
所关联的 $\mathcal{O}$-模态射并不是单射，因为
$2 : \mathbf{F}_2 \to \mathbf{F}_2$ 不是单射，而
$\Spec(\mathbf{F}_2)$ 是 $(\Sch/S)_\tau$ 的一个对象。

\medskip\noindent
第 (7) 部分成立，因为如上所述，函子
$\mathcal{F} \mapsto \mathcal{F}^a$ 是赋环位点态射的拉回函子，而这类
函子由 Modules on Sites, Lemma
\ref{sites-modules-lemma-tensor-product-pullback} 与张量积交换。

\medskip\noindent
若 $\tau = Zariski$，则 (8) 显然成立，因为 $S_{Zar}$ 上
$\mathcal{O}$-模的范畴与拓扑空间 $S$ 上 $\mathcal{O}_S$-模的范畴相同。
若 $\tau = \etale$，则 (8) 成立，这是因为如上所述，函子
$\mathcal{F} \mapsto \mathcal{F}^a$ 是平坦赋环位点态射
$(S_\etale, \mathcal{O}) \to (S_{Zar}, \mathcal{O}_S)$ 的拉回函子；见
Lemma \ref{descent-lemma-compare-etale-zariski-flat}。由 Modules on Sites, Lemma
\ref{sites-modules-lemma-pullback-internal-hom}，沿平坦赋环位点态射的拉回
与从有限表示模出发取内部 Hom 交换。

\medskip\noindent
为说明 (9)，设 $S = \Spec(\mathbf{Z})$。令
$\mathcal{F} = \Coker(2 : \mathcal{O}_S \to \mathcal{O}_S)$
且 $\mathcal{G} = \mathcal{O}_S$。于是
$\mathcal{F}^a = \Coker(2 : \mathcal{O} \to \mathcal{O})$ 且
$\mathcal{G}^a = \mathcal{O}$。所以
$\SheafHom_\mathcal{O}(\mathcal{F}^a, \mathcal{G}^a) = \mathcal{O}[2]$
等于 $\mathcal{O}$ 中的 $2$-挠部分，它并非为零；见 (6) 的证明。另一
方面，模 $\SheafHom_{\mathcal{O}_S}(\mathcal{F}, \mathcal{G})$ 为零。

\medskip\noindent
证明 (10)。设
$0 \to \mathcal{F}_1^a \to \mathcal{E} \to \mathcal{F}_2^a \to 0$
为 $\mathcal{O}$-模的短正合列，其中 $\mathcal{F}_1$ 与
$\mathcal{F}_2$ 是 $S$ 上的拟凝聚层。考虑它在 $S_{Zar}$ 上的限制
$$
0 \to \mathcal{F}_1 \to \mathcal{E}|_{S_{Zar}} \to \mathcal{F}_2
$$
。由 Proposition \ref{descent-proposition-same-cohomology-quasi-coherent}，对任意
仿射开集 $U \subset S$，都有
$H^1(U, \mathcal{F}_1^a) = H^1(U, \mathcal{F}_1) = 0$。故上述序列在右端
也是正合的。由 Schemes, Section \ref{schemes-section-quasi-coherent}，
可知 $\mathcal{F} = \mathcal{E}|_{S_{Zar}}$ 是拟凝聚的。因此得到交换图
$$
\xymatrix{
& \mathcal{F}_1^a \ar[r] \ar[d] &
\mathcal{F}^a \ar[r] \ar[d] &
\mathcal{F}_2^a \ar[r] \ar[d] & 0 \\
0 \ar[r] &
\mathcal{F}_1^a \ar[r] &
\mathcal{E} \ar[r] &
\mathcal{F}_2^a \ar[r] & 0
}
$$
为完成证明，只须说明上行也右正合。为此，再次以
$U = \Spec(A) \subset S$ 表示 $S$ 的一个仿射开集。上面已经看到
$0 \to \mathcal{F}_1(U) \to \mathcal{E}(U) \to \mathcal{F}_2(U) \to 0$
是正合的。对任意仿射概形 $V/U$，若 $V = \Spec(B)$，则态射
$\mathcal{F}_1^a(V) \to \mathcal{E}(V)$ 为单射。按定义，有
$\mathcal{F}_1^a(V) = \mathcal{F}_1(U) \otimes_A B$。单射
$\mathcal{F}_1^a(V) \to \mathcal{E}(V)$ 分解为
$$
\mathcal{F}_1(U) \otimes_A B \to
\mathcal{E}(U) \otimes_A B \to \mathcal{E}(V)
$$
考虑形如 $B = A \oplus M$ 的 $A$-代数 $B$，可见
$\mathcal{F}_1(U) \to \mathcal{E}(U)$ 普遍单射（见 Algebra, Definition
\ref{algebra-definition-universally-injective}）。由于
$\mathcal{E}(U) = \mathcal{F}(U)$，可知 $\mathcal{F}_1 \to \mathcal{F}$
在任意基变换后仍为单射；等价地，
$\mathcal{F}_1^a \to \mathcal{F}^a$ 是单射。
\end{proof}

\begin{lemma}
\label{descent-lemma-properties-quasi-coherent}
设 $S$ 为概形。$S_\etale$ 上拟凝聚模的范畴
$\QCoh(S_\etale, \mathcal{O})$ 具有下列性质：
\begin{enumerate}
\item 拟凝聚层的任意直和都是拟凝聚的。
\item 拟凝聚层的任意余极限都是拟凝聚的。
\item 拟凝聚层态射的核与余核都是拟凝聚的。
\item 给定 $\mathcal{O}$-模的短正合列
$0 \to \mathcal{F}_1 \to \mathcal{F}_2 \to \mathcal{F}_3 \to 0$
，若三者中任意两个是拟凝聚的，则第三个也是拟凝聚的。
\item 给定两个拟凝聚 $\mathcal{O}$-模，其张量积是拟凝聚的。
\item 给定两个拟凝聚 $\mathcal{O}$-模 $\mathcal{F}$、$\mathcal{G}$，
其中 $\mathcal{F}$ 是有限表示的，则内部 Hom
$\SheafHom_\mathcal{O}(\mathcal{F}, \mathcal{G})$
是拟凝聚的。
\end{enumerate}
\end{lemma}

\begin{proof}
相应事实对于概形 $S$ 上的拟凝聚模成立；见 Schemes, Section
\ref{schemes-section-quasi-coherent}。我们利用 Lemma
\ref{descent-lemma-equivalence-quasi-coherent-limits} 把这些事实转移到
$S_\etale$ 上。

\medskip\noindent
证明 (1)。设 $\mathcal{F}_i$，$i \in I$，为
$\QCoh(S_\etale, \mathcal{O})$ 中的一族对象。对 $S$ 上某些拟凝聚模
$\mathcal{G}_i$，写成 $\mathcal{F}_i = \mathcal{G}_i^a$。由所引引理，
$\bigoplus \mathcal{F}_i = (\bigoplus \mathcal{G}_i)^a$，结论得证。

\medskip\noindent
证明 (2)。设 $\mathcal{I} \to \QCoh(S_\etale, \mathcal{O})$，
$i \mapsto \mathcal{F}_i$，为一个图。写成
$\mathcal{F}_i = \mathcal{G}_i^a$，便得到图
$\mathcal{I} \to \QCoh(\mathcal{O}_S)$。由所引引理，
$\colim \mathcal{F}_i = (\colim \mathcal{G}_i)^a$，结论得证。

\medskip\noindent
证明 (3)。设 $a : \mathcal{F} \to \mathcal{F}'$ 为
$\QCoh(S_\etale, \mathcal{O})$ 中的一个箭头。对 $S$ 上拟凝聚模之间的
某个态射 $b : \mathcal{G} \to \mathcal{G}'$，写成 $a = b^a$。由所引
引理，有 $\Ker(a) = \Ker(b)^a$ 且
$\Coker(a) = \Coker(b)^a$，结论得证。

\medskip\noindent
证明 (4)。除已知 $\mathcal{F}_1$ 与 $\mathcal{F}_3$ 拟凝聚的情形外，
结论都由 (3) 得出。在这一例外情形，写成
$\mathcal{F}_1 = \mathcal{G}_1^a$ 与
$\mathcal{F}_3 = \mathcal{G}_3^a$，其中 $\mathcal{G}_i$ 是 $S$ 上的
拟凝聚层。结论由 Lemma
\ref{descent-lemma-equivalence-quasi-coherent-limits} 的第 (10) 部分得出。

\medskip\noindent
证明 (5)。设 $\mathcal{F}$ 与 $\mathcal{F}'$ 属于
$\QCoh(S_\etale, \mathcal{O})$。写成
$\mathcal{F} = \mathcal{G}^a$ 与
$\mathcal{F}' = (\mathcal{G}')^a$，其中 $\mathcal{G}$ 与
$\mathcal{G}'$ 是 $S$ 上的拟凝聚层。由所引引理，有
$\mathcal{F} \otimes_\mathcal{O} \mathcal{F}' =
(\mathcal{G} \otimes_{\mathcal{O}_S} \mathcal{G}')^a$
，结论得证。

\medskip\noindent
证明 (6)。设 $\mathcal{F}$ 与 $\mathcal{G}$ 属于
$\QCoh(S_\etale, \mathcal{O})$，且 $\mathcal{F}$ 是有限表示的。写成
$\mathcal{F} = \mathcal{H}^a$ 与 $\mathcal{G} = (\mathcal{I})^a$，其中
$\mathcal{H}$ 与 $\mathcal{I}$ 是 $S$ 上的拟凝聚层。由 Lemma
\ref{descent-lemma-equivalence-quasi-coherent-properties}，可知 $\mathcal{H}$ 是
有限表示的。由 Lemma \ref{descent-lemma-equivalence-quasi-coherent-limits} 的
第 (8) 部分，有
$\SheafHom_\mathcal{O}(\mathcal{F}, \mathcal{G}) =
(\SheafHom_{\mathcal{O}_S}(\mathcal{H}, \mathcal{I}))^a$
，结论得证。
\end{proof}

\begin{lemma}
\label{descent-lemma-properties-quasi-coherent-on-big}
设 $S$ 为概形，并设 $\tau \in \{Zariski, \linebreak[0] \etale, \linebreak[0]
smooth, \linebreak[0] syntomic, \linebreak[0] fppf\}$。
$(\Sch/S)_\tau$ 上拟凝聚模的范畴
$\QCoh((\Sch/S)_\tau, \mathcal{O})$ 具有下列性质：
\begin{enumerate}
\item 拟凝聚层的任意直和都是拟凝聚的。
\item 拟凝聚层的任意余极限都是拟凝聚的。
\item 拟凝聚层态射的余核是拟凝聚的。
\item 给定 $\mathcal{O}$-模的短正合列
$0 \to \mathcal{F}_1 \to \mathcal{F}_2 \to \mathcal{F}_3 \to 0$
，若 $\mathcal{F}_1$ 与 $\mathcal{F}_3$ 拟凝聚，则
$\mathcal{F}_2$ 也拟凝聚。
\item 给定两个拟凝聚 $\mathcal{O}$-模，其张量积是拟凝聚的。
\item 给定两个拟凝聚 $\mathcal{O}$-模 $\mathcal{F}$、$\mathcal{G}$，
其中 $\mathcal{F}$ 有限局部自由，则内部 Hom
$\SheafHom_\mathcal{O}(\mathcal{F}, \mathcal{G})$
是拟凝聚的。
\end{enumerate}
\end{lemma}

\begin{proof}
相应事实对于概形 $S$ 上的拟凝聚模成立；见 Schemes, Section
\ref{schemes-section-quasi-coherent}。我们利用 Lemma
\ref{descent-lemma-equivalence-quasi-coherent-limits} 把这些事实转移到
$(\Sch/S)_\tau$ 上。

\medskip\noindent
证明 (1)。设 $\mathcal{F}_i$，$i \in I$，为
$\QCoh((\Sch/S)_\tau, \mathcal{O})$ 中的一族对象。对 $S$ 上某些拟凝聚模
$\mathcal{G}_i$，写成 $\mathcal{F}_i = \mathcal{G}_i^a$。由所引引理，
$\bigoplus \mathcal{F}_i = (\bigoplus \mathcal{G}_i)^a$，结论得证。

\medskip\noindent
证明 (2)。设 $\mathcal{I} \to \QCoh((\Sch/S)_\tau, \mathcal{O})$，
$i \mapsto \mathcal{F}_i$，为一个图。写成
$\mathcal{F}_i = \mathcal{G}_i^a$，便得到图
$\mathcal{I} \to \QCoh(\mathcal{O}_S)$。由所引引理，
$\colim \mathcal{F}_i = (\colim \mathcal{G}_i)^a$，结论得证。

\medskip\noindent
证明 (3)。设 $a : \mathcal{F} \to \mathcal{F}'$ 为
$\QCoh((\Sch/S)_\tau, \mathcal{O})$ 中的一个箭头。对 $S$ 上拟凝聚模之间
的某个态射 $b : \mathcal{G} \to \mathcal{G}'$，写成 $a = b^a$。由所引
引理，有 $\Coker(a) = \Coker(b)^a$（因为余核是余极限），结论得证。

\medskip\noindent
证明 (4)。写成 $\mathcal{F}_1 = \mathcal{G}_1^a$ 与
$\mathcal{F}_3 = \mathcal{G}_3^a$，其中 $\mathcal{G}_i$ 是 $S$ 上的
拟凝聚层。结论由 Lemma
\ref{descent-lemma-equivalence-quasi-coherent-limits} 的第 (10) 部分得出。

\medskip\noindent
证明 (5)。设 $\mathcal{F}$ 与 $\mathcal{F}'$ 属于
$\QCoh((\Sch/S)_\tau, \mathcal{O})$。写成
$\mathcal{F} = \mathcal{G}^a$ 与 $\mathcal{F}' = (\mathcal{G}')^a$，其中
$\mathcal{G}$ 与 $\mathcal{G}'$ 是 $S$ 上的拟凝聚层。由所引引理，有
$\mathcal{F} \otimes_\mathcal{O} \mathcal{F}' =
(\mathcal{G} \otimes_{\mathcal{O}_S} \mathcal{G}')^a$
，结论得证。

\medskip\noindent
证明 (6)。对某个拟凝聚 $\mathcal{O}_S$-模，写成
$\mathcal{F} = \mathcal{H}^a$。由 Lemma
\ref{descent-lemma-equivalence-quasi-coherent-properties}，可知 $\mathcal{H}$ 有限
局部自由。该问题在 $S$ 上是 Zariski 局部的（略去细节），所以可设
$\mathcal{H} = \mathcal{O}_S^{\oplus n}$ 有限自由。于是
$\mathcal{F} = \mathcal{O}^{\oplus n}$，而
$\SheafHom_\mathcal{O}(\mathcal{F}, \mathcal{G}) = \mathcal{G}^{\oplus n}$
是拟凝聚的。
\end{proof}

\begin{example}
\label{descent-example-internal-hom-not-qcoh}
设 $S$ 为概形。对 Lemma \ref{descent-lemma-properties-quasi-coherent-on-big} 中
考虑的某个拓扑 $\tau$，设 $\mathcal{F}$ 与 $\mathcal{G}$ 为
$(\Sch/S)_\tau$ 上的拟凝聚模。一般而言，
$\SheafHom_\mathcal{O}(\mathcal{F}, \mathcal{G})$
并非拟凝聚的，即使 $\mathcal{F}$ 是有限表示的也如此。事实上，取
$S = \Spec(\mathbf{Z})$，
$\mathcal{F} = \Coker(2 : \mathcal{O} \to \mathcal{O})$,
且 $\mathcal{G} = \mathcal{O}$。此时
$\SheafHom_\mathcal{O}(\mathcal{F}, \mathcal{G}) = \mathcal{O}[2]$
等于 $\mathcal{O}$ 中的 $2$-挠部分，而它不是拟凝聚的。
\end{example}

\begin{lemma}
\label{descent-lemma-qc-colimits}
设 $S$ 为概形。
\begin{enumerate}
\item 范畴 $\QCoh((\Sch/S)_{fppf}, \mathcal{O})$ 有余极限，并且这些
余极限与下列范畴中的余极限一致：
$\textit{Mod}((\Sch/S)_{Zar}, \mathcal{O})$,
$\textit{Mod}((\Sch/S)_\etale, \mathcal{O})$，以及
$\textit{Mod}((\Sch/S)_{fppf}, \mathcal{O})$.
\item 给定 $\QCoh((\Sch/S)_{fppf}, \mathcal{O})$ 中的
$\mathcal{F}, \mathcal{G}$，在 $\textit{Mod}((\Sch/S)_{Zar}, \mathcal{O})$、
$\textit{Mod}((\Sch/S)_\etale, \mathcal{O})$ 或
$\textit{Mod}((\Sch/S)_{fppf}, \mathcal{O})$ 中计算的张量积
$\mathcal{F} \otimes_\mathcal{O} \mathcal{G}$ 相互一致，其共同值是
$\QCoh((\Sch/S)_{fppf}, \mathcal{O})$ 的对象。
\item 给定 $\QCoh((\Sch/S)_{fppf}, \mathcal{O})$ 中的
$\mathcal{F}, \mathcal{G}$，且 $\mathcal{F}$ 有限局部自由（在 fppf 拓扑
中，等价地在 \'etale 拓扑中，或等价地在 Zariski 拓扑中），则在
$\textit{Mod}((\Sch/S)_{Zar}, \mathcal{O})$、
$\textit{Mod}((\Sch/S)_\etale, \mathcal{O})$ 或
$\textit{Mod}((\Sch/S)_{fppf}, \mathcal{O})$ 中计算的内部 Hom
$\SheafHom_\mathcal{O}(\mathcal{F}, \mathcal{G})$
相互一致，其共同值是 $\QCoh((\Sch/S)_{fppf}, \mathcal{O})$ 的对象。
\end{enumerate}
\end{lemma}

\begin{proof}
本引理以略有不同的方式汇集上面证明的结果。首先，由 Lemma
\ref{descent-lemma-properties-quasi-coherent-on-big}，我们已经知道 (1)、(2) 或
(3) 中构造的输出属于 $\QCoh((\Sch/S)_\tau, \mathcal{O})$。还须在每种
情形证明结果与所用拓扑无关。关键在于 Proposition
\ref{descent-proposition-equivalence-quasi-coherent} 中的等价
$\QCoh(\mathcal{O}_S) \to \QCoh((\Sch/S)_\tau, \mathcal{O})$,
$\mathcal{F} \mapsto \mathcal{F}^a$
由同一个公式给出，而与拓扑 $\tau \in \{Zariski, \etale, fppf\}$ 的选择
无关。

\medskip\noindent
证明 (1)。设 $\mathcal{I} \to \QCoh((\Sch/S)_{fppf}, \mathcal{O})$，
$i \mapsto \mathcal{F}_i$，为一个图。写成
$\mathcal{F}_i = \mathcal{G}_i^a$，便得到图
$\mathcal{I} \to \QCoh(\mathcal{O}_S)$。由 Lemma
\ref{descent-lemma-equivalence-quasi-coherent-limits}，对
$\tau \in \{Zariski, \etale, fppf\}$，在
$\textit{Mod}((\Sch/S)_\tau, \mathcal{O})$ 中有
$\colim \mathcal{F}_i = (\colim \mathcal{G}_i)^a$。这证明了 (1)。

\medskip\noindent
证明 (2)。写成 $\mathcal{F} = \mathcal{H}^a$ 与
$\mathcal{G} = (\mathcal{I})^a$，其中 $\mathcal{H}$ 与 $\mathcal{I}$
是 $S$ 上的拟凝聚层。于是
% P05-ERR-0030
$\mathcal{F} \otimes_\mathcal{O} \mathcal{G} =
(\mathcal{H} \otimes_\mathcal{O} \mathcal{I})^a$
，并且由 Lemma \ref{descent-lemma-equivalence-quasi-coherent-limits}，对
$\tau \in \{Zariski, \etale, fppf\}$，此等式在
$\textit{Mod}((\Sch/S)_\tau, \mathcal{O})$ 中成立。这证明了 (2)。

\medskip\noindent
证明 (3)。设 $\mathcal{F}$ 与 $\mathcal{G}$ 属于
$\QCoh((\Sch/S)_{fppf}, \mathcal{O})$。写成
$\mathcal{F} = \mathcal{H}^a$，其中 $\mathcal{H}$ 是 $S$ 上的拟凝聚层。
由 Lemma \ref{descent-lemma-equivalence-quasi-coherent-properties}，有
\begin{align*}
\mathcal{F}\text{ 在 fppf 拓扑中有限局部自由}
& \Leftrightarrow
\mathcal{H}\text{ 在下列对象上有限局部自由： }S \\
& \Leftrightarrow
\mathcal{F}\text{ 在平展拓扑中有限局部自由} \\
& \Leftrightarrow
\mathcal{H}\text{ 在下列对象上有限局部自由： }S \\
& \Leftrightarrow
\mathcal{F}\text{ 在 Zariski 拓扑中有限局部自由}
\end{align*}
这解释了第 (3) 部分的括号内陈述。现在，若这些等价条件成立，则
$\mathcal{H}$ 有限局部自由。Modules on Sites, Section
\ref{sites-modules-section-internal-hom} 中
$\SheafHom_\mathcal{O}(\mathcal{F}, \mathcal{G})$ in
的构造，仅依赖于作为模预层的 $\mathcal{F}$ 与 $\mathcal{G}$（只有
输出 $\SheafHom$ 是否为层，才依赖于 $\mathcal{F}$ 与 $\mathcal{G}$
是否为层）。
\end{proof}










\section{拟凝聚模与仿射概形}
\label{descent-section-alternative-quasi-coherent}

\noindent
设 $S$ 为概形\footnote{在本节中，正如 Topologies, Section
\ref{topologies-section-change-topologies} 中一样，我们选择各位点
$(\Sch/S)_\tau$，使得当
$\tau \in \{Zariski, \etale, smooth, syntomic, fppf\}$ 时，它们具有相同
的底范畴。于是各位点 $(\textit{Aff}/S)_\tau$ 也具有相同的底范畴。}。
设 $\tau \in \{Zariski, \etale, smooth, syntomic, fppf\}$。回忆，
$(\textit{Aff}/S)_\tau$ 是 $(\Sch/S)_\tau$ 中对象为仿射概形的全子范畴；
将覆盖规定为标准 $\tau$-覆盖，便把它变为一个位点。由 Topologies,
Lemmas
\ref{topologies-lemma-affine-big-site-Zariski},
\ref{topologies-lemma-affine-big-site-etale},
\ref{topologies-lemma-affine-big-site-smooth},
\ref{topologies-lemma-affine-big-site-syntomic}，以及
\ref{topologies-lemma-affine-big-site-fppf}
，有拓扑斯等价
$g : \Sh((\textit{Aff}/S)_\tau) \to \Sh((\Sch/S)_\tau)$
，其拉回函子由限制给出。回忆 $\mathcal{O}$ 表示 $(\Sch/S)_\tau$ 上的
结构层；为明确区分记号，暂以 $\mathcal{O}_{\textit{Aff}}$ 表示
$\mathcal{O}$ 在 $(\textit{Aff}/S)_\tau$ 上的限制。于是得到赋环拓扑斯的
等价
\begin{equation}
\label{descent-equation-alternative-ringed}
(\Sh((\textit{Aff}/S)_\tau), \mathcal{O}_{\textit{Aff}})
\longrightarrow
(\Sh((\Sch/S)_\tau), \mathcal{O})
\end{equation}
。由此也可以比较拟凝聚模。

\begin{lemma}
\label{descent-lemma-quasi-coherent-alternative}
设 $S$ 为概形，并设 $\mathcal{F}$ 为 $(\textit{Aff}/S)_{fppf}$ 上的
$\mathcal{O}_{\textit{Aff}}$-模预层。下列条件等价：
\begin{enumerate}
\item 对 $(\textit{Aff}/S)_{fppf}$ 的每个态射 $U \to U'$，态射
$\mathcal{F}(U') \otimes_{\mathcal{O}(U')} \mathcal{O}(U) \to \mathcal{F}(U)$
是同构；
\item $\mathcal{F}$ 是 $(\textit{Aff}/S)_{Zar}$ 上的层，并且按 Modules on
Sites, Definition \ref{sites-modules-definition-site-local} 的意义，是赋环
位点 $((\textit{Aff}/S)_{Zar}, \mathcal{O}_{\textit{Aff}})$ 上的拟凝聚模；
\item 对 \'etale 拓扑，条件与 (2) 相同；
\item 对 smooth 拓扑，条件与 (2) 相同；
\item 对 syntomic 拓扑，条件与 (2) 相同；
\item 对 fppf 拓扑，条件与 (2) 相同；
\item 经等价 (\ref{descent-equation-alternative-ringed})，$\mathcal{F}$ 对应于
$(\Sch/S)_{Zar}$,
$(\Sch/S)_\etale$,
$(\Sch/S)_{smooth}$,
$(\Sch/S)_{syntomic}$，或
$(\Sch/S)_{fppf}$
之一上的拟凝聚模；
\item $\mathcal{F}$ 由 Section \ref{descent-section-quasi-coherent-sheaves} 中
所述过程从唯一的拟凝聚 $\mathcal{O}_S$-模 $\mathcal{G}$ 得到。
\end{enumerate}
\end{lemma}

\begin{proof}
由于拟凝聚模的概念是内蕴的（Modules on Sites, Lemma
\ref{sites-modules-lemma-special-locally-free}），可知等价
(\ref{descent-equation-alternative-ringed}) 诱导拟凝聚模范畴之间的等价。
Proposition \ref{descent-proposition-equivalence-quasi-coherent} 说明，在
$\Sch/S$ 上研究拟凝聚模时使用哪种拓扑并不重要，并且还说明 (8) 与
(7) 相同。因此 (2) -- (8) 全都等价。

\medskip\noindent
假设等价条件 (2) -- (8) 成立，并令 $\mathcal{G}$ 如 (8) 中所述。设
$h : U \to U' \to S$ 为 $\textit{Aff}/S$ 的一个态射。以
$f : U \to S$ 与 $f' : U' \to S$ 表示结构态射，于是
$f = f' \circ h$。我们有
$\mathcal{F}(U') = \Gamma(U', (f')^*\mathcal{G})$ 以及
$\mathcal{F}(U) = \Gamma(U, f^*\mathcal{G}) = \Gamma(U, h^*(f')^*\mathcal{G})$。
故 (1) 由 Schemes, Lemma \ref{schemes-lemma-widetilde-pullback} 成立。

\medskip\noindent
假设 (1) 成立。为完成证明，只须证明 (2)。设 $U$ 为
$(\textit{Aff}/S)_{Zar}$ 的对象，记 $U = \Spec(R)$。标准开覆盖
$U = U_1 \cup \ldots \cup U_n$ 由 $U_i = D(f_i)$ 给出，其中某些元素
$f_1, \ldots, f_n \in R$ 生成 $R$ 的单位理想。由性质 (1)，有
$$
\mathcal{F}(U_i) =
\mathcal{F}(U) \otimes_R R_{f_i} =
\mathcal{F}(U)_{f_i}
$$
and
$$
\mathcal{F}(U_i \cap U_j) =
\mathcal{F}(U) \otimes_R R_{f_if_j} =
\mathcal{F}(U)_{f_if_j}
$$
所以由 Algebra, Lemma \ref{algebra-lemma-cover-module}，可知
$\mathcal{F}$ 是 $(\textit{Aff}/S)_{Zar}$ 上的层。选取自由 $R$-模表示
$$
\bigoplus\nolimits_{k \in K} R
\longrightarrow
\bigoplus\nolimits_{l \in L} R
\longrightarrow
\mathcal{F}(U)
\longrightarrow 0
$$
。由性质 (1) 以及张量积的右正合性，可知对
$(\textit{Aff}/S)_{Zar}$ 中的每个态射 $U' \to U$，都得到表示
$$
% P05-ERR-0031
\bigoplus\nolimits_{k \in K} \mathcal{O}_{Aff}(U')
\longrightarrow
\bigoplus\nolimits_{l \in L} \mathcal{O}_{Aff}(U')
\longrightarrow
\mathcal{F}(U')
\longrightarrow 0
$$
换言之，可见 $\mathcal{F}$ 在局部化范畴
$(\textit{Aff}/S)_{Zar}/U$ 上的限制有表示
$$
\bigoplus\nolimits_{k \in K} \mathcal{O}_{Aff}|_{(\textit{Aff}/S)_{Zar}/U}
\longrightarrow
\bigoplus\nolimits_{l \in L} \mathcal{O}_{Aff}|_{(\textit{Aff}/S)_{Zar}/U}
\longrightarrow
\mathcal{F}|_{(\textit{Aff}/S)_{Zar}/U}
\longrightarrow 0
$$
尽管记号十分繁复，这便完成了证明。
\end{proof}

\noindent
我们继续本节引言中开始的讨论。设 $\tau \in \{Zariski, \etale\}$。
回忆，$S_{affine, \tau}$ 是 $S_\tau$ 中对象为仿射概形的全子范畴；将覆盖
规定为标准 $\tau$-覆盖，便把它变为一个位点。见 Topologies,
Definitions \ref{topologies-definition-big-small-Zariski} 与
\ref{topologies-definition-big-small-etale}。分别由 Topologies, Lemmas
\ref{topologies-lemma-alternative-zariski} 与
\ref{topologies-lemma-alternative}，有拓扑斯等价
$g : \Sh(S_{affine, \tau}) \to \Sh(S_\tau)$，其拉回函子由限制给出。
回忆 $\mathcal{O}$ 表示 $S_\tau$ 上的结构层；为明确区分记号，暂以
$\mathcal{O}_{affine}$ 表示 $\mathcal{O}$ 在 $S_{affine, \tau}$ 上的限制。
于是得到赋环拓扑斯的等价
\begin{equation}
\label{descent-equation-alternative-small-ringed}
(\Sh(S_{affine, \tau}), \mathcal{O}_{affine})
\longrightarrow
(\Sh(S_\tau), \mathcal{O})
\end{equation}
。由此也可以比较拟凝聚模。

\begin{lemma}
\label{descent-lemma-quasi-coherent-alternative-small}
设 $S$ 为概形，设 $\tau \in \{Zariski, \etale\}$，并设 $\mathcal{F}$ 为
$S_{affine, \tau}$ 上的 $\mathcal{O}_{affine}$-模预层。下列条件等价：
\begin{enumerate}
\item 对 $S_{affine, \tau}$ 的每个态射 $U \to U'$，态射
$\mathcal{F}(U') \otimes_{\mathcal{O}(U')} \mathcal{O}(U) \to \mathcal{F}(U)$
是同构；
\item $\mathcal{F}$ 是 $S_{affine, \tau}$ 上的层，并且按 Modules on
Sites, Definition \ref{sites-modules-definition-site-local} 的意义，是赋环
位点 $(S_{affine, \tau}, \mathcal{O}_{affine})$ 上的拟凝聚模；
\item 经等价 (\ref{descent-equation-alternative-small-ringed})，$\mathcal{F}$ 对应
于 $S_\tau$ 上的拟凝聚模；
\item $\mathcal{F}$ 由 Section \ref{descent-section-quasi-coherent-sheaves} 中
所述过程从唯一的拟凝聚 $\mathcal{O}_S$-模 $\mathcal{G}$ 得到。
\end{enumerate}
\end{lemma}

\begin{proof}
我们对 \'etale 拓扑的情形证明此结论。

\medskip\noindent
假设 (1) 成立。为证明 $\mathcal{F}$ 是层，设
$\mathcal{U} = \{U_i \to U\}_{i = 1, \ldots, n}$ 为
$S_{affine, \etale}$ 的一个覆盖。由我们对 $\mathcal{F}$ 的假设，
% P05-ERR-0032
$\mathcal{F}$ 关于 $\mathcal{U}$ 的层条件归结为证明
$$
0 \to \mathcal{F}(U) \to
\prod \mathcal{F}(U) \otimes_{\mathcal{O}(U)} \mathcal{O}(U_i) \to
\prod \mathcal{F}(U) \otimes_{\mathcal{O}(U)} \mathcal{O}(U_i \times_U U_j)
$$
正合。这是因为 $\mathcal{O}(U) \to \prod \mathcal{O}(U_i)$ 忠实平坦
（由 Lemma \ref{descent-lemma-standard-covering-Cech} 以及
$S_{affine, \etale}$ 中的覆盖是标准 \'etale 覆盖这一事实），因而可以
应用 Lemma \ref{descent-lemma-ff-exact}。接着证明 $\mathcal{F}$ 在
$S_{affine, \etale}$ 上拟凝聚。具体而言，对 $S_{affine, \etale}$ 中的
$U$，置 $R = \mathcal{O}(U)$，并选取自由 $R$-模表示
$$
\bigoplus\nolimits_{k \in K} R
\longrightarrow
\bigoplus\nolimits_{l \in L} R
\longrightarrow
\mathcal{F}(U)
\longrightarrow 0
$$
。由性质 (1) 以及张量积的右正合性，可知对 $S_{affine, \etale}$ 中的
每个态射 $U' \to U$，都得到表示
$$
\bigoplus\nolimits_{k \in K} \mathcal{O}(U')
\longrightarrow
\bigoplus\nolimits_{l \in L} \mathcal{O}(U')
\longrightarrow
\mathcal{F}(U')
\longrightarrow 0
$$
换言之，可见 $\mathcal{F}$ 在局部化范畴 $S_{affine, etale}/U$ 上的限制
有表示
$$
% P05-ERR-0033
\bigoplus\nolimits_{k \in K} \mathcal{O}_{affine}|_{S_{affine, \etale}/U}
\longrightarrow
\bigoplus\nolimits_{l \in L} \mathcal{O}_{affine}|_{S_{affine, \etale}/U}
\longrightarrow
\mathcal{F}|_{S_{affine, \etale}/U}
\longrightarrow 0
$$
，这正是证明 $\mathcal{F}$ 拟凝聚所需的表示。尽管记号十分繁复，
这便完成了 (1) 蕴含 (2) 的证明。

\medskip\noindent
由于拟凝聚模的概念是内蕴的（Modules on Sites, Lemma
\ref{sites-modules-lemma-special-locally-free}），可知等价
(\ref{descent-equation-alternative-small-ringed}) 诱导拟凝聚模范畴之间的等价。
所以 (2) 与 (3) 等价。

\medskip\noindent
(3) 与 (4) 的等价性由 Proposition
\ref{descent-proposition-equivalence-quasi-coherent} 得出。

\medskip\noindent
假设 (4) 成立，下面证明 (1)。令 $\mathcal{G}$ 如 (4) 中所述。设
$h : U \to U' \to S$ 为 $S_{affine, \etale}$ 的一个态射。以
$f : U \to S$ 与 $f' : U' \to S$ 表示结构态射，于是
$f = f' \circ h$。我们有
$\mathcal{F}(U') = \Gamma(U', (f')^*\mathcal{G})$ 以及
$\mathcal{F}(U) = \Gamma(U, f^*\mathcal{G}) = \Gamma(U, h^*(f')^*\mathcal{G})$。
故 (1) 由 Schemes, Lemma \ref{schemes-lemma-widetilde-pullback} 成立。

\medskip\noindent
省略 Zariski 拓扑情形的证明。
\end{proof}





\section{寄生模}
\label{descent-section-parasitic}

\noindent
寄生模是限制到在基概形上平坦的概形后为零的模。其正式定义如下。

\begin{definition}
\label{descent-definition-parasitic}
设 $S$ 为概形，设 $\tau \in \{Zar, \etale,
smooth, syntomic, fppf\}$，并设 $\mathcal{F}$ 为 $(\Sch/S)_\tau$ 上的
$\mathcal{O}$-模预层。
\begin{enumerate}
\item 若对每个平坦态射 $U \to S$ 都有 $\mathcal{F}(U) = 0$，则称
$\mathcal{F}$ 是{\it 寄生的}\footnote{这可能不是标准术语。}。
\item 若对每个 $\tau$-覆盖 $\{U_i \to S\}_{i \in I}$ 以及所有 $i$，
都有 $\mathcal{F}(U_i) = 0$，则称 $\mathcal{F}$ 是
{\it 关于 $\tau$-拓扑寄生的}。
\end{enumerate}
\end{definition}

\noindent
若 $\tau = fppf$，这意味着只要 $U \to S$ 平坦且局部有限表示，就有
$\mathcal{F}|_{U_{Zar}} = 0$；其他情形类似。

\begin{lemma}
\label{descent-lemma-cohomology-parasitic}
设 $S$ 为概形，设 $\tau \in \{Zar, \etale, smooth,
syntomic, fppf\}$，并设 $\mathcal{G}$ 为 $(\Sch/S)_\tau$ 上的
$\mathcal{O}$-模预层。
\begin{enumerate}
\item 若 $\mathcal{G}$ 关于 $\tau$-拓扑寄生，则对每个在 $S$ 中开、
或在 $S$ 上 \'etale、或在 $S$ 上 smooth、或在 $S$ 上 syntomic、或在
$S$ 上平坦且局部有限表示的相应 $U$，都有
$H^p_\tau(U, \mathcal{G}) = 0$。
\item 若 $\mathcal{G}$ 寄生，则对每个在 $S$ 上平坦的 $U$，都有
$H^p_\tau(U, \mathcal{G}) = 0$。
\end{enumerate}
\end{lemma}

\begin{proof}
证明 $\tau = fppf$ 的情形；其他情形的证明完全相同。假设意味着，对任意
平坦且局部有限表示的 $U \to S$，都有 $\mathcal{G}(U) = 0$。把
Cohomology on Sites, Lemma
\ref{sites-cohomology-lemma-cech-vanish-collection} 应用于由平坦且局部
有限表示的 $U \to S$ 组成的子集
$\mathcal{B} \subset \Ob((\Sch/S)_{fppf})$，以及由 $\mathcal{B}$ 中对象的
所有 fppf 覆盖组成的集合 $\text{Cov}$。
\end{proof}

\begin{lemma}
\label{descent-lemma-direct-image-parasitic}
设 $f : T \to S$ 为概形态射。对 $(\Sch/T)_\tau$ 上的任意寄生
$\mathcal{O}$-模，前推 $f_*\mathcal{F}$ 与高阶直像
$R^if_*\mathcal{F}$ 都是 $(\Sch/S)_\tau$ 上的寄生 $\mathcal{O}$-模。
\end{lemma}

\begin{proof}
回忆，$R^if_*\mathcal{F}$ 是与下列预层相关联的层：
$$
U \mapsto H^i((\Sch/U \times_S T)_\tau, \mathcal{F})
$$
；见 Cohomology on Sites, Lemma
\ref{sites-cohomology-lemma-higher-direct-images}。若 $U \to S$ 平坦，则
$U \times_S T \to T$ 作为其基变换也是平坦的。故由 Lemma
\ref{descent-lemma-cohomology-parasitic}，上示群为零。若
$\{U_i \to U\}$ 是 $\tau$-覆盖，则 $U_i \times_S T \to T$ 也平坦。
因此很清楚，上示预层的层化在所有平坦于 $S$ 的概形 $U$ 上为零。
\end{proof}

\begin{lemma}
\label{descent-lemma-quasi-coherent-and-flat-base-change}
设 $S$ 为概形，设 $\tau \in \{Zar, \etale\}$，并设 $\mathcal{G}$ 为
$(\Sch/S)_{fppf}$ 上的 $\mathcal{O}$-模层，满足：
\begin{enumerate}
\item $\mathcal{G}|_{S_\tau}$ 拟凝聚；
\item 对每个平坦且局部有限表示的态射 $g : U \to S$，典范态射
$g_{\tau, small}^*(\mathcal{G}|_{S_\tau}) \to \mathcal{G}|_{U_\tau}$
是同构。
\end{enumerate}
那么，对每个在 $S$ 上平坦且局部有限表示的 $U$，都有
$H^p(U, \mathcal{G}) = H^p(U, \mathcal{G}|_{U_\tau})$。
\end{lemma}

\begin{proof}
令 $\mathcal{F}$ 为 $\mathcal{G}|_{S_\tau}$ 到大 fppf 位点
$(\Sch/S)_{fppf}$ 的拉回。注意 $\mathcal{F}$ 拟凝聚。存在典范比较态射
$\varphi : \mathcal{F} \to \mathcal{G}$；由假设 (1) 与 (2)，它对所有
平坦且局部有限表示的 $g : U \to S$ 诱导同构
$\mathcal{F}|_{U_\tau} \to \mathcal{G}|_{U_\tau}$
。所以在短正合列
$$
0 \to \Ker(\varphi) \to \mathcal{F} \to \Im(\varphi) \to 0
$$
and
$$
0 \to \Im(\varphi) \to \mathcal{G} \to \Coker(\varphi) \to 0
$$
中，层 $\Ker(\varphi)$ 与 $\Coker(\varphi)$ 关于 fppf 拓扑寄生。由
Lemma \ref{descent-lemma-cohomology-parasitic}，可知对平坦且局部有限表示的
$g : U \to S$，态射
$H^p(U, \mathcal{F}) \to H^p(U, \mathcal{G})$ 是同构。对于
$\mathcal{F}$，结论由 Proposition
\ref{descent-proposition-same-cohomology-quasi-coherent} 成立，故命题得证。
\end{proof}













\section{Fpqc 覆盖是普遍有效满态射}
\label{descent-section-fpqc-universal-effective-epimorphisms}

\noindent
我们应用上述材料证明一个有趣的结果，即 Lemma
\ref{descent-lemma-fpqc-universal-effective-epimorphisms}。由 Sites, Section
\ref{sites-section-representable-sheaves}，该引理蕴含：当
$\tau \in \{Zariski, fppf, \etale, smooth, syntomic\}$ 时，任一位点
$(\Sch/S)_\tau$ 上的可表预层都是层。先证明一个辅助引理。

\begin{lemma}
\label{descent-lemma-equiv-fibre-product}
对概形 $X$，以 $|X|$ 表示其底集合。设 $f : X \to S$ 为概形态射。则
$$
|X \times_S X| \to |X| \times_{|S|} |X|
$$
为满射。
\end{lemma}

\begin{proof}
这立即由 Schemes, Lemma \ref{schemes-lemma-points-fibre-product} 中对
纤维积上点的描述得出。
\end{proof}


\begin{lemma}
\label{descent-lemma-universal-effective-epimorphism-affine}
设 $\{f_i : X_i \to X\}_{i \in I}$ 为一族仿射概形态射。下列条件等价：
\begin{enumerate}
\item 对任意拟凝聚 $\mathcal{O}_X$-模 $\mathcal{F}$，有
$$
\Gamma(X, \mathcal{F}) =
\text{Equalizer}\left(
\xymatrix{
\prod\nolimits_{i \in I} \Gamma(X_i, f_i^*\mathcal{F})
\ar@<1ex>[r] \ar@<-1ex>[r] &
\prod\nolimits_{i, j \in I}
\Gamma(X_i \times_X X_j, (f_i \times f_j)^*\mathcal{F})
}
\right)
$$
\item $\{f_i : X_i \to X\}_{i \in I}$ 在仿射概形范畴中是普遍有效满态射
（Sites, Definition \ref{sites-definition-universal-effective-epimorphisms}）。
\end{enumerate}
\end{lemma}

\begin{proof}
假设 (2) 成立，并设 $\mathcal{F}$ 为拟凝聚 $\mathcal{O}_X$-模。考虑概形
（Constructions, Section \ref{constructions-section-spec}）
$$
X' = \underline{\Spec}_X(\mathcal{O}_X \oplus \mathcal{F})
$$
，其中 $\mathcal{O}_X \oplus \mathcal{F}$ 是 $\mathcal{O}_X$-代数，其乘法为
$(f, s)(f', s') = (ff', fs' + f's)$.
若 $s_i \in \Gamma(X_i, f_i^*\mathcal{F})$ 是截面，则 $s_i$ 确定下式中的
唯一元素：
$$
\Gamma(X' \times_X X_i, \mathcal{O}_{X' \times_X X_i}) =
\Gamma(X_i, \mathcal{O}_{X_i}) \oplus \Gamma(X_i, f_i^*\mathcal{F})
$$
省略该等式的证明。若 $(s_i)_{i \in I}$ 属于 (1) 中的等化子，则利用对
任意概形 $T$ 都成立的等式
$$
\Mor(T, \mathbf{A}^1_\mathbf{Z}) = \Gamma(T, \mathcal{O}_T)
$$
，可知这些截面定义一族态射
$h_i : X' \times_X X_i \to \mathbf{A}^1_\mathbf{Z}$，并且作为态射
$(X' \times_X X_i) \times_{X'} (X' \times_X X_j) \to \mathbf{A}^1_\mathbf{Z}$，
有 $h_i \circ \text{pr}_1 = h_j \circ \text{pr}_2$。由假设 (2)，得到与
各态射 $h_i$ 相容的态射 $h : X' \to \mathbf{A}^1_\mathbf{Z}$，继而确定
元素 $s \in \Gamma(X, \mathcal{F})$。省略验证 $s$ 在
$\Gamma(X_i, f_i^*\mathcal{F})$ 中映到 $s_i$。

\medskip\noindent
假设 (1)。设 $T$ 为仿射概形，并设 $h_i : X_i \to T$ 为一族态射，使得
对所有 $i, j \in I$，在 $X_i \times_X X_j$ 上有
$h_i \circ \text{pr}_1 = h_j \circ \text{pr}_2$。于是
$$
\prod h_i^\sharp :
\Gamma(T, \mathcal{O}_T)
\to
\prod \Gamma(X_i, \mathcal{O}_{X_i})
$$
映入等化子；对 $\mathcal{F} = \mathcal{O}_X$ 应用本引理的假设，便得到
环同态
$\Gamma(T, \mathcal{O}_T) \to \Gamma(X, \mathcal{O}_X)$
。该环同态对应于态射 $h : X \to T$，且 $h_i = h \circ f_i$。所以这族
态射是有效满态射。

\medskip\noindent
设 $p : Y \to X$ 为仿射概形态射。应用上一段的结果，我们来证明 $f_i$
的基变换 $g_i : Y_i \to Y$ 构成有效满态射。事实上，若 $\mathcal{G}$
为拟凝聚 $\mathcal{O}_Y$-模，则
$$
\Gamma(Y, \mathcal{G}) = \Gamma(X, p_*\mathcal{G}),\quad
\Gamma(Y_i, g_i^*\mathcal{G}) = \Gamma(X, f_i^*p_*\mathcal{G}),
$$
and
$$
\Gamma(Y_i \times_Y Y_j, (g_i \times g_j)^*\mathcal{G}) =
\Gamma(X, (f_i \times f_j)^*p_*\mathcal{G})
$$
，这是由平凡基变换公式（Cohomology of Schemes, Lemma
\ref{coherent-lemma-affine-base-change}）得到的。因此可见，这族 $g_i$
满足本引理的性质 (1)。
\end{proof}

\begin{lemma}
\label{descent-lemma-universal-effective-epimorphism-surjective}
设 $\{f_i : X_i \to X\}_{i \in I}$ 为一族概形态射。
\begin{enumerate}
\item 若这族态射在概形范畴中是普遍有效满态射，则 $\coprod f_i$ 为满射。
\item 若 $X$ 与 $X_i$ 仿射，且这族态射在仿射概形范畴中是普遍有效
满态射，则 $\coprod f_i$ 为满射。
\end{enumerate}
\end{lemma}

\begin{proof}
省略。提示：沿 $\Spec(\kappa(x)) \to X$ 作基变换，即可看出任意
$x \in X$ 都必须在像中。
\end{proof}

\begin{lemma}
\label{descent-lemma-check-universal-effective-epimorphism-affine}
设 $\{f_i : X_i \to X\}_{i \in I}$ 为一族概形态射。若对每个以 $Y$
为仿射概形的态射 $Y \to X$，基变换族 $g_i : Y_i \to Y$ 都构成有效
满态射，则 $f_i$ 这一族态射在概形范畴中构成普遍有效满态射。
\end{lemma}

\begin{proof}
设 $Y \to X$ 为概形态射。我们须证明基变换
$g_i : Y_i \to Y$ 构成有效满态射。为此，设给定概形 $T$ 以及态射
$h_i : Y_i \to T$，并且在 $Y_i \times_Y Y_j$ 上有
$h_i \circ \text{pr}_1 = h_j \circ \text{pr}_2$。选取仿射开覆盖
$Y = \bigcup V_\alpha$。令 $V_{\alpha, i}$ 为 $V_\alpha$ 在 $Y_i$ 中的
逆像。于是 $V_{\alpha, i}  \to V_\alpha$ 是 $f_i$ 沿
$V_\alpha \to X$ 的基变换。因此由假设，各限制
$h_i|_{V_{\alpha, i}}$ 来自概形态射 $h_\alpha : V_\alpha \to T$。
留给读者验证这些态射在交集上一致，并定义所需态射 $Y \to T$。见
Schemes, Section \ref{schemes-section-glueing-schemes} 中的讨论。
\end{proof}

\begin{lemma}
\label{descent-lemma-universal-effective-epimorphism}
设 $\{f_i : X_i \to X\}_{i \in I}$ 为一族仿射概形态射。假设 Lemma
\ref{descent-lemma-universal-effective-epimorphism-affine} 中的等价条件成立，并且
还假设对任意仿射概形态射 $Y \to X$，映射
$$
\coprod X_i \times_X Y \longrightarrow Y
$$
是拓扑空间的商映射（Topology, Definition
\ref{topology-definition-submersive}）。那么这族态射在概形范畴中是
普遍有效满态射。
\end{lemma}

\begin{proof}
由 Lemma \ref{descent-lemma-check-universal-effective-epimorphism-affine}，只须沿以
$Y$ 为仿射概形的 $Y \to X$ 对这族态射作基变换。置
$Y_i = X_i \times_X Y$。设 $T$ 为概形，并设 $h_i : Y_i \to T$ 为一族
态射，使得在 $Y_i \times_Y Y_j$ 上有
$h_i \circ \text{pr}_1 = h_j \circ \text{pr}_2$。注意，作为集合，$Y$ 是
从 $\coprod Y_i \times_Y Y_j$ 到 $\coprod Y_i$ 的两个映射的余等化子。
事实上，满射性由 Lemma
\ref{descent-lemma-universal-effective-epimorphism-surjective} 的仿射情形给出，
单射性由 Lemma \ref{descent-lemma-equiv-fibre-product} 给出。因此存在与各
$h_i$ 相容的底集合映射 $h : Y \to T$。由本引理的第二个条件，可知
$h$ 连续！所以，若 $y \in Y$，且 $U \subset T$ 是 $h(y)$ 的仿射开邻域，
则可找到仿射开集 $V \subset Y$，使得 $h(V) \subset U$。置
$V_i = Y_i \times_Y V = X_i \times_X V$，即可应用 Lemma
\ref{descent-lemma-universal-effective-epimorphism-affine} 中证明的结果，看出
$h|_V : V \to U \subset T$ 来自唯一的仿射概形态射
$h_V : V \to U$，且对所有 $i$，它作为概形态射与
$h_i|_{V_i}$ 一致。粘合这些 $h_V$（见 Schemes, Section
\ref{schemes-section-glueing-schemes}），便得到所需态射 $Y \to T$。
\end{proof}

\begin{lemma}
\label{descent-lemma-open-fpqc-covering}
设 $\{f_i : T_i \to T\}_{i \in I}$ 为 fpqc 覆盖。假设对每个 $i$ 都有开
子集 $W_i \subset T_i$，并且对所有 $i, j \in I$，作为
$T_i \times_T T_j$ 的开子集，有
$\text{pr}_0^{-1}(W_i) = \text{pr}_1^{-1}(W_j)$。则存在唯一的开子集
$W \subset T$，使得对每个 $i$ 都有 $W_i = f_i^{-1}(W)$。
\end{lemma}

\begin{proof}
把 Lemma \ref{descent-lemma-equiv-fibre-product} 应用于映射
$\coprod_{i \in I} T_i \to T$。它蕴含存在子集 $W \subset T$，使得对每个
$i$ 都有 $W_i = f_i^{-1}(W)$；具体地，$W = \bigcup f_i(W_i)$。为证明
$W$ 开，可以在 $T$ 上 Zariski 局部地工作。因此可设 $T$ 仿射。利用
Topologies, Definition \ref{topologies-definition-fpqc-covering}，可选取
一个加细 $\{T_i \to T\}_{i \in I}$ 的标准 fpqc 覆盖
$\{g_j : V_j \to T\}_{j \in J}$。令 $\alpha : J \to I$ 与
$h_j : V_j \to T_{\alpha(j)}$ 如 Sites, Definition
\ref{sites-definition-morphism-coverings} 中所述。于是
$g_j^{-1}(W) = h_j^{-1}(W_{\alpha(j)})$。所以可设
$\{f_i : T_i \to T\}$ 是标准 fpqc 覆盖。在这种情形，对态射
$\coprod T_i \to T$ 应用 Morphisms, Lemma
\ref{morphisms-lemma-fpqc-quotient-topology}，即可断定 $W$ 开。
\end{proof}

\begin{lemma}
\label{descent-lemma-fpqc-universal-effective-epimorphisms}
设 $\{T_i \to T\}$ 为 fpqc 覆盖；见 Topologies, Definition
\ref{topologies-definition-fpqc-covering}。则 $\{T_i \to T\}$ 在概形范畴
中是普遍有效满态射；见 Sites, Definition
\ref{sites-definition-universal-effective-epimorphisms}。换言之，概形范畴
上的每个可表函子都满足 fpqc 拓扑的层条件；见 Topologies, Definition
\ref{topologies-definition-sheaf-property-fpqc}。
\end{lemma}

\begin{proof}
设 $S$ 为概形。我们须证明如下陈述：给定态射
$\varphi_i : T_i \to S$，满足
$\varphi_i|_{T_i \times_T T_j} = \varphi_j|_{T_i \times_T T_j}$，则存在
唯一的态射 $T \to S$，其在每个 $T_i$ 上的限制为 $\varphi_i$。换言之，
我们须证明函子 $h_S = \Mor_{\Sch}( - , S)$ 满足 fpqc 拓扑的层性质。

\medskip\noindent
若 $\{T_i \to T\}$ 是 Zariski 覆盖，则结论由 Schemes, Lemma
\ref{schemes-lemma-glue} 得出。因此 Topologies, Lemma
\ref{topologies-lemma-sheaf-property-fpqc} 把问题归约到覆盖
$\{X \to Y\}$ 的情形，其中该覆盖由仿射概形之间的单个满平坦态射给出。

\medskip\noindent
第一种证明。由 Lemma \ref{descent-lemma-sheaf-condition-holds}，拟凝聚模关于
$\{X \to Y\}$ 满足层条件。由 Lemma \ref{descent-lemma-open-fpqc-covering}，态射
$X \to Y$ 是普遍商映射。因此可应用 Lemma
\ref{descent-lemma-universal-effective-epimorphism}，看出 $\{X \to Y\}$ 是普遍
有效满态射。

\medskip\noindent
第二种证明。令 $R \to A$ 为与满平坦态射 $\pi : X \to Y$ 对应的忠实
平坦环同态。设 $f : X \to S$ 为态射，并且作为态射
$X \times_Y X = \Spec(A \otimes_R A) \to S$，有
$f \circ \text{pr}_1 = f \circ \text{pr}_2$。由 Lemma
\ref{descent-lemma-equiv-fibre-product}，可知在底集合映射的意义下，对某个
（集合论）映射 $g : \Spec(R) \to S$，$f$ 具有形式
$f = g \circ \pi$。由 Morphisms, Lemma
\ref{morphisms-lemma-fpqc-quotient-topology} 以及 $f$ 连续这一事实，可知
$g$ 连续。

\medskip\noindent
取 $y \in Y = \Spec(R)$。选取包含 $g(y)$ 的仿射开集 $U \subset S$，
记 $U = \Spec(B)$。由上述结果，可选取 $r \in R$，使得
$y \in D(r) \subset g^{-1}(U)$。$f$ 在 $\pi^{-1}(D(r))$ 上取值于 $U$
的限制对应于环同态 $B \to A_r$。由对 $f$ 的假设，所诱导的两个环同态
$B \to A_r \otimes_{R_r} A_r = (A \otimes_R A)_r$ 相等。注意
$R_r \to A_r$ 忠实平坦。由 Lemma \ref{descent-lemma-ff-exact}，两个箭头
$A_r \to A_r \otimes_{R_r} A_r$ 的等化子是 $R_r$。所以
$B \to A_r$ 唯一地通过一个映射 $B \to R_r$ 分解。该映射继而给出概形
态射 $D(r) \to U \to S$；见 Schemes, Lemma
\ref{schemes-lemma-morphism-into-affine}。

\medskip\noindent
到目前为止我们证明了什么？我们已经证明，对任意素理想
$\mathfrak p \subset R$，都存在标准仿射开集 $D(r) \subset \Spec(R)$，
使得态射 $f|_{\pi^{-1}(D(r))} : \pi^{-1}(D(r)) \to S$ 唯一地通过某个
概形态射 $D(r) \to S$ 分解。省略验证这些态射粘合为所需态射
$\Spec(R) \to S$。
\end{proof}

\begin{lemma}
\label{descent-lemma-coequalizer-fpqc-local}
考虑概形 $X, Y, Z$、态射 $a, b : X \to Y$ 以及态射
$c : Y \to Z$，其中 $c \circ a = c \circ b$。置
$d = c \circ a = c \circ b$。若存在 fpqc 覆盖 $\{Z_i \to Z\}$，使得
\begin{enumerate}
\item 对所有 $i$，态射 $Y \times_{c, Z} Z_i \to Z_i$ 是
$(a, 1) : X \times_{d, Z} Z_i \to Y \times_{c, Z} Z_i$ 与
$(b, 1) : X \times_{d, Z} Z_i \to Y \times_{c, Z} Z_i$ 的余等化子；
\item 对所有 $i$ 与 $i'$，态射
$Y \times_{c, Z} (Z_i \times_Z Z_{i'}) \to (Z_i \times_Z Z_{i'})$
是
$(a, 1) : X \times_{d, Z} (Z_i \times_Z Z_{i'}) \to
Y \times_{c, Z} (Z_i \times_Z Z_{i'})$ 与
$(b, 1) : X \times_{d, Z} (Z_i \times_Z Z_{i'}) \to
Y \times_{c, Z} (Z_i \times_Z Z_{i'})$
的余等化子。
\end{enumerate}
则 $c$ 是 $a$ 与 $b$ 的余等化子。
\end{lemma}

\begin{proof}
事实上，由 Lemma \ref{descent-lemma-fpqc-universal-effective-epimorphisms}，
对概形 $T$，态射 $Z \to T$ 等同于一族在交叠上一致的态射
$Z_i \to T$。
\end{proof}















\section{态射有限性与光滑性性质的下降}
\label{descent-section-descent-finiteness-morphisms}

\noindent
本节证明态射的若干性质（光滑、局部有限表示等）沿忠实平坦态射下降。
先从代数版本开始。（关注 ``Noetherian'' 情形的读者应参看 Lemma
\ref{descent-lemma-finite-type-local-source-fppf-algebra}，而不是下一个引理。）

\begin{lemma}
\label{descent-lemma-flat-finitely-presented-permanence-algebra}
设 $R \to A \to B$ 为环同态。假设 $R \to B$ 有限表示，而
$A \to B$ 忠实平坦且有限表示。则 $R \to A$ 有限表示。
\end{lemma}

\begin{proof}
考虑代数 $C = B \otimes_A B$，以及由 $p(b) = b \otimes 1$ 与
$q(b) = 1 \otimes b$ 给出的两个映射 $p, q : B \to C$。两个复合
$A \to B \to C$ 当然相同。注意，$p : B \to C$ 是 $A \to B$ 的基变换，
因而平坦且有限表示；所以环同态 $R \to C$ 作为
$R \to B \to C$ 的复合也是有限表示的。

\medskip\noindent
我们将利用 Algebra, Lemma
\ref{algebra-lemma-characterize-finite-presentation} 的判别准则，证明
$R \to A$ 有限表示。设 $S$ 为任意 $R$-代数，并设
$S = \colim_{\lambda \in \Lambda} S_\lambda$ 写成 $R$-代数的有向余极限。
设 $A \to S$ 为 $R$-代数同态。我们须证明 $A \to S$ 通过某个
$S_\lambda$ 分解。考虑环 $B' = S \otimes_A B$ 与
$C' = S \otimes_A C = B' \otimes_S B'$。由于 $B$ 在 $A$ 上忠实平坦且
有限表示，$B'$ 在 $S$ 上也忠实平坦且有限表示。把 Algebra, Lemma
\ref{algebra-lemma-flat-finite-presentation-limit-flat} 的第 (2) 部分应用于
偶对 $(S \to B', B')$ 和系统 $(S_\lambda)$，可知存在
$\lambda_0 \in \Lambda$ 以及平坦有限表示的 $S_{\lambda_0}$-代数
$B_{\lambda_0}$，使得
$B' = S \otimes_{S_{\lambda_0}} B_{\lambda_0}$。对
$\lambda \geq \lambda_0$，置
$B_\lambda = S_\lambda \otimes_{S_{\lambda_0}} B_{\lambda_0}$ 且
$C_\lambda = B_\lambda \otimes_{S_\lambda} B_\lambda$。

\medskip\noindent
暂时中断论证，以证明当 $\lambda$ 足够大时，
$S_\lambda \to B_\lambda$ 忠实平坦。（这其实应该另立为其他地方的一个
引理，或许放在极限章中。）由于
$\Spec(B_{\lambda_0}) \to \Spec(S_{\lambda_0})$ 平坦且有限表示，它是开
映射（见 Morphisms, Lemma \ref{morphisms-lemma-fppf-open}）。令
$I \subset S_{\lambda_0}$ 为理想，使得 $V(I) \subset \Spec(S_{\lambda_0})$
是该像的补集。注意，像的形成与基变换交换。因此，由于
$\Spec(B') \to \Spec(S)$ 为满射，且
$B' = B_{\lambda_0} \otimes_{S_{\lambda_0}} S$，可知 $IS = S$。所以对
某个 $\lambda \geq \lambda_0$，有 $IS_{\lambda} = S_\lambda$。对这个指标
以及所有更大的 $\lambda$，态射
$\Spec(B_\lambda) \to \Spec(S_\lambda)$ 都是满射。

\medskip\noindent
仿照证明第一段的记号，以
$p_\lambda, q_\lambda : B_\lambda \to C_\lambda$ 表示两个典范映射。
于是 $B' = \colim_{\lambda \geq \lambda_0} B_\lambda$ 且
$C' = \colim_{\lambda \geq \lambda_0} C_\lambda$。由于 $B$ 与 $C$ 在
$R$ 上有限表示，反复应用 Algebra, Lemma
\ref{algebra-lemma-characterize-finite-presentation}，可知存在
$\lambda \geq \lambda_0$ 以及 $R$-代数同态
$B \to B_\lambda$、$C \to C_\lambda$，使得图表
$$
\xymatrix{
C \ar[rr] & &
C_\lambda \\
B \ar[rr]
\ar@<1ex>[u]^-p
\ar@<-1ex>[u]_-q
& &
B_\lambda
\ar@<1ex>[u]^-{p_\lambda}
\ar@<-1ex>[u]_-{q_\lambda}
}
$$
交换。这意味着 $A \to B \to B_\lambda$ 映入 $p_\lambda$ 与
$q_\lambda$ 的等化子。由 Lemma \ref{descent-lemma-ff-exact}，可知
$S_\lambda$ 是 $p_\lambda$ 与 $q_\lambda$ 的等化子。因此得到所需环同态
$A \to S_\lambda$，命题得证。
\end{proof}

\noindent
下面给出关于有限型性质的较容易版本。

\begin{lemma}
\label{descent-lemma-finite-type-local-source-fppf-algebra}
设 $R \to A \to B$ 为环同态。假设 $R \to B$ 有限型，而
$A \to B$ 忠实平坦且有限表示。则 $R \to A$ 有限型。
\end{lemma}

\begin{proof}
由 Algebra, Lemma
\ref{algebra-lemma-descend-faithfully-flat-finite-presentation}，存在交换图
$$
\xymatrix{
R \ar[r] \ar@{=}[d] &
A_0 \ar[d] \ar[r] &
B_0 \ar[d] \\
R \ar[r] & A \ar[r] & B
}
$$
，其中 $R \to A_0$ 有限表示，$A_0 \to B_0$ 忠实平坦且有限表示，并且
$B = A \otimes_{A_0} B_0$。由假设，$R \to B$ 有限型，所以可以向
$A_0$ 添加一些元素，并假设映射 $B_0 \to B$ 为满射！此时，由于
$A_0 \to B_0$ 忠实平坦，而
$$
(A_0 \to A) \otimes_{A_0} B_0 \cong (B_0 \to B)
$$
为满射，所以 $A_0 \to A$ 也为满射。命题得证。
\end{proof}

\begin{lemma}
\label{descent-lemma-flat-finitely-presented-permanence}
\begin{reference}
\cite[IV, 17.7.5 (i) and (ii)]{EGA}.
\end{reference}
设
$$
\xymatrix{
X \ar[rr]_f \ar[rd]_p & &
Y \ar[dl]^q \\
& S
}
$$
为概形态射的交换图。假设 $f$ 满、平坦且局部有限表示，并假设 $p$
局部有限表示（相应地，局部有限型）。则 $q$ 局部有限表示（相应地，
局部有限型）。
\end{lemma}

\begin{proof}
该问题在 $S$ 与 $Y$ 上是局部的。因此可设 $S$ 与 $Y$ 仿射。由于 $f$
平坦且局部有限表示，$f$ 是开映射（Morphisms, Lemma
\ref{morphisms-lemma-fppf-open}）。又因为 $Y$ 拟紧，存在有限多个仿射开集
$X_i \subset X$，使得 $Y = \bigcup f(X_i)$。显然可以用
$\coprod X_i$ 替换 $X$，从而也可设 $X$ 仿射。在这种情形，本引理等价
于上面的 Lemma \ref{descent-lemma-flat-finitely-presented-permanence-algebra}
（相应地，Lemma \ref{descent-lemma-finite-type-local-source-fppf-algebra}）。
\end{proof}

\noindent
我们用它改进前面得到的一些关于态射的结果。

\begin{lemma}
\label{descent-lemma-syntomic-smooth-etale-permanence}
设
$$
\xymatrix{
X \ar[rr]_f \ar[rd]_p & &
Y \ar[dl]^q \\
& S
}
$$
为概形态射的交换图。假设
\begin{enumerate}
\item $f$ 满且 syntomic（相应地，smooth；相应地，\'etale）；
\item $p$ syntomic（相应地，smooth；相应地，\'etale）。
\end{enumerate}
则 $q$ syntomic（相应地，smooth；相应地，\'etale）。
\end{lemma}

\begin{proof}
把 Morphisms, Lemmas
\ref{morphisms-lemma-syntomic-permanence},
\ref{morphisms-lemma-smooth-permanence}，以及
\ref{morphisms-lemma-etale-permanence-two}
与上面的 Lemma \ref{descent-lemma-flat-finitely-presented-permanence} 结合即可。
\end{proof}

\noindent
事实上，可以把这一结果加强如下。

\begin{lemma}
\label{descent-lemma-smooth-permanence}
设
$$
\xymatrix{
X \ar[rr]_f \ar[rd]_p & &
Y \ar[dl]^q \\
& S
}
$$
为概形态射的交换图。假设
\begin{enumerate}
\item $f$ 满、平坦且局部有限表示；
\item $p$ smooth（相应地，\'etale）。
\end{enumerate}
则 $q$ smooth（相应地，\'etale）。
\end{lemma}

\begin{proof}
假设 (1) 成立且 $p$ smooth。由 Lemma
\ref{descent-lemma-flat-finitely-presented-permanence}，可知 $q$ 局部有限表示。
由 Morphisms, Lemma \ref{morphisms-lemma-flat-permanence}，可知 $q$ 平坦。
所以现在只须证明 $q$ 的纤维光滑；见 Morphisms, Lemma
\ref{morphisms-lemma-smooth-flat-smooth-fibres}。对 $s \in S$，把 Varieties,
Lemma \ref{varieties-lemma-flat-under-smooth} 应用于平坦满态射
$X_s \to Y_s$，即得结论。省略 \'etale 情形的证明。
\end{proof}

\begin{remark}
\label{descent-remark-smooth-permanence}
在 Lemma \ref{descent-lemma-smooth-permanence} 的假设 (1) 以及 $p$ smooth 的条件
下，$X \to Y$ 不一定 smooth。一个反例是 $S = \Spec(k)$、
$X = \Spec(k[s])$、$Y = \Spec(k[t])$，且 $f$ 由 $t \mapsto s^2$ 给出。
不过，关于 $f$ 的结构的一些信息，另见 Lemma
\ref{descent-lemma-syntomic-permanence}。
\end{remark}

\begin{lemma}
\label{descent-lemma-syntomic-permanence}
设
$$
\xymatrix{
X \ar[rr]_f \ar[rd]_p & &
Y \ar[dl]^q \\
& S
}
$$
为概形态射的交换图。假设
\begin{enumerate}
\item $f$ 满、平坦且局部有限表示；
\item $p$ syntomic。
\end{enumerate}
则 $q$ 与 $f$ 都 syntomic。
\end{lemma}

\begin{proof}
% P05-ERR-0034
由 Lemma \ref{descent-lemma-flat-finitely-presented-permanence}，可知 $q$ 有限表示。
由 Morphisms, Lemma \ref{morphisms-lemma-flat-permanence}，可知 $q$ 平坦。
由 Morphisms, Lemma
\ref{morphisms-lemma-syntomic-locally-standard-syntomic}，现在只须证明
$Y \to S$ 的纤维以及 $X \to Y$ 的纤维的局部环都是局部完全交环。为此，
可取 $X \to Y \to S$ 在一点 $s \in S$ 上的纤维；也就是说，可以假设
$S$ 是某个域的谱。选取一点 $x \in X$，其像为 $y \in Y$，并考虑环同态
$$
\mathcal{O}_{Y, y} \longrightarrow \mathcal{O}_{X, x}
$$
。这是局部 Noether 环之间的平坦局部同态。局部环
$\mathcal{O}_{X, x}$ 是完全交。于是可应用 Avramov 的结果；见 Divided
Power Algebra, Lemma \ref{dpa-lemma-avramov}；从而断定
$\mathcal{O}_{Y, y}$ 与
$\mathcal{O}_{X, x}/\mathfrak m_y\mathcal{O}_{X, x}$ 都是完全交环。
\end{proof}

\noindent
下面这类引理有时很有用。

\begin{lemma}
\label{descent-lemma-curiosity}
设 $X \to Y \to Z$ 为概形态射。令 $P$ 为概形态射的下列性质之一：
平坦、局部有限型、局部有限表示。假设 $X \to Z$ 具有 $P$，并且
$\{X \to Y\}$ 可由 $Y$ 的一个 fppf 覆盖加细。则 $Y \to Z$ 具有 $P$。
\end{lemma}

\begin{proof}
设 $\Spec(C) \subset Z$ 为仿射开集，并设 $\Spec(B) \subset Y$ 为映入
$\Spec(C)$ 的仿射开集。关于 $X \to Y$ 的假设蕴含：可以找到标准仿射
fppf 覆盖 $\{\Spec(B_j) \to \Spec(B)\}$ 以及提升
$x_j : \Spec(B_j) \to X$。由于 $\Spec(B_j)$ 拟紧，可以找到位于
$\Spec(B)$ 上方的有限多个仿射开集 $\Spec(A_i) \subset X$，使得每个
$x_j$ 的像都包含在并集 $\bigcup \Spec(A_i)$ 中。因此，用各
$\Spec(B_j)$ 自身的标准仿射 Zariski 覆盖替换它之后，可以假设有标准
仿射 fppf 覆盖 $\{\Spec(B_i) \to \Spec(B)\}$，使得每个
$\Spec(B_i) \to Y$ 都通过位于 $\Spec(B)$ 上方的一个仿射开集
$\Spec(A_i) \subset X$ 分解。换言之，对每个 $i$，都有环同态
$C \to B \to A_i \to B_i$。注意，还可以考虑
$$
C \to B \to A = \prod A_i \to B' = \prod B_i
$$
，并且环同态 $B \to \prod B_i$ 忠实平坦且有限表示。

\medskip\noindent
$P = flat$ 的情形。此时已知 $C \to A$ 平坦，须证明 $C \to B$ 平坦。
设 $N \to N' \to N''$ 为 $C$-模的正合列。我们要证明
$N \otimes_C B \to N' \otimes_C B \to N'' \otimes_C B$ 正合。令 $H$
为其上同调，并令 $H'$ 为
$N \otimes_C B' \to N' \otimes_C B' \to N'' \otimes_C B'$ 的上同调。
由于 $B \to B'$ 平坦，有 $H' = H \otimes_B B'$。另一方面，
$N \otimes_C A \to N' \otimes_C A \to N'' \otimes_C A$ 正合，因而上同调
为零。所以态射 $H \to H'$ 为零（因为它通过零模分解）。故
$H' = 0$。由于 $B \to B'$ 忠实平坦，可知 $H = 0$，正如所需。

\medskip\noindent
$P = locally\ finite\ type$ 的情形。此时已知 $C \to A$ 有限型，须证明
$C \to B$ 有限型。由于 $B \to B'$ 有限表示（因而有限型），可知
$A \to B'$ 有限型；见 Algebra, Lemma
\ref{algebra-lemma-compose-finite-type}。所以 $C \to B'$ 有限型，结论由
Lemma \ref{descent-lemma-finite-type-local-source-fppf-algebra} 得出。

\medskip\noindent
$P = locally\ finite\ presentation$ 的情形。此时已知 $C \to A$ 有限表示，
须证明 $C \to B$ 有限表示。由于 $B \to B'$ 有限表示，而 $B \to A$
有限型，可知 $A \to B'$ 有限表示；见 Algebra, Lemma
\ref{algebra-lemma-compose-finite-type}。所以 $C \to B'$ 有限表示，结论
由 Lemma \ref{descent-lemma-flat-finitely-presented-permanence-algebra} 得出。
\end{proof}










\section{概形的局部性质}
\label{descent-section-descending-properties}

\noindent
常常可以证明概形的一个覆盖的各成员具有某种性质。在许多情形下，这蕴含
该概形也具有此性质。例如，若 $S$ 是概形，$f : S' \to S$ 是满平坦态射，
且 $S'$ 是既约概形，则 $S$ 既约。考察局部环并使用 Algebra, Lemma
\ref{algebra-lemma-descent-reduced} 即可证明这一点。我们称既约这一性质
{\it 沿平坦满态射下降}。这类结果的一部分汇集于 Algebra, Section
\ref{algebra-section-descending-properties}；关于概形的结果见 Section
\ref{descent-section-variants}。关于态射性质下降的一些类似结果见 Section
\ref{descent-section-descent-finiteness-morphisms}。

\medskip\noindent
另一方面，也存在满平坦态射 $f : S' \to S$，其中 $S$ 既约而 $S'$ 不
既约；例如态射 $\Spec(k[x]/(x^2)) \to \Spec(k)$。因此既约这一性质不
{\it 沿平坦态射上升}。剩余域无限这一性质确实沿平坦态射上升（当然不
沿满平坦态射下降）。这类结果的一部分汇集于 Algebra, Section
\ref{algebra-section-ascending-properties}。

\medskip\noindent
最后，若一个性质沿平坦态射上升，并沿平坦满态射下降，则称它
{\it 关于平坦拓扑为局部的}。一个略显琐碎的例子是剩余域具有给定特征
这一性质。为给出更精确的表述，并把它与概形上的各种拓扑联系起来，作
如下正式定义。

\begin{definition}
\label{descent-definition-property-local}
设 $\mathcal{P}$ 为概形的一项性质。设
$\tau \in \{fpqc, \linebreak[0] fppf, \linebreak[0] syntomic, \linebreak[0]
smooth, \linebreak[0] \etale, \linebreak[0] Zariski\}$.
若对任意 $\tau$-覆盖 $\{S_i \to S\}_{i \in I}$（见 Topologies, Section
\ref{topologies-section-procedure}）都有
$$
S \text{ 具有 }\mathcal{P}
\Leftrightarrow
\text{每个 }S_i \text{ 具有 }\mathcal{P}.
$$
则称 $\mathcal{P}$ {\it 在 $\tau$-拓扑中为局部的}。
\end{definition}

\noindent
确切地说，由于同构总是覆盖，可知（或要求）$S$ 具有性质
$\mathcal{P}$，当且仅当任意与 $S$ 同构的概形 $S'$ 具有该性质。事实上，
若 $\tau = fpqc, \linebreak[0] fppf, \linebreak[0] syntomic,
\linebreak[0] smooth, \linebreak[0] \etale$ 或 $Zariski$，且 $S$ 具有
$\mathcal{P}$，而 $S' \to S$ 相应地为平坦、平坦且局部有限表示、
syntomic、smooth、\'etale 或开浸入，则 $S'$ 具有 $\mathcal{P}$。这是
因为总能把 $\{S' \to S\}$ 扩充为一个 $\tau$-覆盖。

\medskip\noindent
有如下蕴含关系：
$\mathcal{P}$ 在 fpqc 拓扑中为局部的
$\Rightarrow$
$\mathcal{P}$ 在 fppf 拓扑中为局部的
$\Rightarrow$
$\mathcal{P}$ 在 syntomic 拓扑中为局部的
$\Rightarrow$
$\mathcal{P}$ 在 smooth 拓扑中为局部的
$\Rightarrow$
$\mathcal{P}$ 在 \'etale 拓扑中为局部的
$\Rightarrow$
$\mathcal{P}$ 在 Zariski 拓扑中为局部的。这由 Topologies, Lemmas
\ref{topologies-lemma-zariski-etale},
\ref{topologies-lemma-zariski-etale-smooth},
\ref{topologies-lemma-zariski-etale-smooth-syntomic},
\ref{topologies-lemma-zariski-etale-smooth-syntomic-fppf}, and
\ref{topologies-lemma-zariski-etale-smooth-syntomic-fppf-fpqc}.
得出。

\begin{lemma}
\label{descent-lemma-descending-properties}
设 $\mathcal{P}$ 为概形的一项性质。设
$\tau \in \{fpqc, \linebreak[0] fppf, \linebreak[0]
\etale, \linebreak[0] smooth, \linebreak[0] syntomic\}$。假设
\begin{enumerate}
\item 该性质在 Zariski 拓扑中为局部的；
\item 对任意仿射概形态射 $S' \to S$，若按照 $\tau$ 为 fpqc、fppf、
\'etale、smooth 或 syntomic，该态射相应地为平坦、平坦且有限表示、
\'etale、smooth 或 syntomic，则只要 $S$ 具有性质 $\mathcal{P}$，$S'$
也具有性质 $\mathcal{P}$；
\item 对任意满的仿射概形态射 $S' \to S$，若按照 $\tau$ 为 fpqc、fppf、
\'etale、smooth 或 syntomic，该态射相应地为平坦、平坦且有限表示、
\'etale、smooth 或 syntomic，则只要 $S'$ 具有性质 $\mathcal{P}$，$S$
也具有性质 $\mathcal{P}$。
\end{enumerate}
则 $\mathcal{P}$ 在基底上是 $\tau$-局部的。
\end{lemma}

\begin{proof}
这几乎立即由 $\tau$-覆盖的定义得出；见 Topologies, Definition
\ref{topologies-definition-fpqc-covering}
\ref{topologies-definition-fppf-covering}
\ref{topologies-definition-etale-covering}
\ref{topologies-definition-smooth-covering}，或
\ref{topologies-definition-syntomic-covering}
以及 Topologies, Lemma
\ref{topologies-lemma-fpqc-affine},
\ref{topologies-lemma-fppf-affine},
\ref{topologies-lemma-etale-affine},
\ref{topologies-lemma-smooth-affine}，或
\ref{topologies-lemma-syntomic-affine}。省略细节。
\end{proof}

\begin{remark}
\label{descent-remark-descending-properties-standard}
在上面的 Lemma \ref{descent-lemma-descending-properties} 中，若
$\tau = smooth$，则在条件 (3) 中可以假设该态射是（满的）标准 smooth
态射。$\tau = syntomic$ 或 $\tau = \etale$ 时也类似。
\end{remark}





\section{在 fppf 拓扑中为局部的概形性质}
\label{descent-section-descending-properties-fppf}

\noindent
本节给出一些在 fppf 拓扑中于基底上为局部的概形性质。

\begin{lemma}
\label{descent-lemma-Noetherian-local-fppf}
性质 $\mathcal{P}(S) =$``$S$ 局部 Noether'' 在 fppf 拓扑中为局部的。
\end{lemma}

\begin{proof}
我们使用 Lemma \ref{descent-lemma-descending-properties}。首先注意，``局部
Noether'' 在 Zariski 拓扑中为局部的。这由定义即得；见 Properties,
Definition \ref{properties-definition-noetherian}。接着，若
$S' \to S$ 是仿射概形之间的平坦有限表示态射，而 $S$ 局部 Noether，
则 $S'$ 局部 Noether；这是 Morphisms, Lemma
\ref{morphisms-lemma-finite-type-noetherian}。最后，若 $S' \to S$ 是仿射
概形之间的满、平坦、有限表示态射，而 $S'$ 局部 Noether，则须证明
$S$ 局部 Noether。这由 Algebra, Lemma
\ref{algebra-lemma-descent-Noetherian} 得出。因此 Lemma
\ref{descent-lemma-descending-properties} 的 (1)、(2)、(3) 成立，命题得证。
\end{proof}

\begin{lemma}
\label{descent-lemma-Jacobson-local-fppf}
性质 $\mathcal{P}(S) =$``$S$ 是 Jacobson 概形'' 在 fppf 拓扑中为局部的。
\end{lemma}

\begin{proof}
我们使用 Lemma \ref{descent-lemma-descending-properties}。首先注意，``Jacobson''
在 Zariski 拓扑中为局部的；这是 Properties, Lemma
\ref{properties-lemma-locally-jacobson}。接着，若 $S' \to S$ 是仿射概形
之间的平坦有限表示态射，而 $S$ 是 Jacobson 概形，则 $S'$ 是 Jacobson
概形；这是 Morphisms, Lemma
\ref{morphisms-lemma-Jacobson-universally-Jacobson}。最后，设
$f : S' \to S$ 是仿射概形之间的满、平坦、有限表示态射，且 $S'$ 是
Jacobson 概形；须证明 $S$ 是 Jacobson 概形。记 $S = \Spec(A)$、
$S' = \Spec(B)$，且 $S' \to S$ 由 $A \to B$ 给出。于是 $A \to B$
有限表示且忠实平坦。此外，环 $B$ 是 Jacobson 环；见 Properties,
Lemma \ref{properties-lemma-locally-jacobson}。

\medskip\noindent
由 Algebra, Lemma \ref{algebra-lemma-fppf-fpqf}，存在图表
$$
\xymatrix{
B \ar[rr] & & B' \\
& A \ar[ru] \ar[lu] &
}
$$
，其中 $A \to B'$ 有限表示、忠实平坦且拟有限。特别地，$B \to B'$
有限型，而由 Algebra, Proposition
\ref{algebra-proposition-Jacobson-permanence}，可知 $B'$ 是 Jacobson 环。
所以可以假设 $A \to B$ 不仅忠实平坦且有限表示，而且拟有限。

\medskip\noindent
反设 $A$ 不是 Jacobson 环以导出矛盾。由 Algebra, Lemma
\ref{algebra-lemma-characterize-jacobson}，存在非极大素理想
$\mathfrak p \subset A$ 以及元素 $f \in A$，$f \not \in \mathfrak p$，
使得 $V(\mathfrak p) \cap D(f) = \{\mathfrak p\}$。

\medskip\noindent
如下导出矛盾。先取 $A$ 的极大理想，使得
$\mathfrak p \subset \mathfrak m$。选取位于 $\mathfrak m$ 上方的素理想
$\mathfrak m' \subset B$（它存在，因为 $A \to B$ 忠实平坦；见 Algebra,
Lemma \ref{algebra-lemma-ff-rings}）。由于 $A \to B$ 平坦，由下降定理
（见 Algebra, Lemma \ref{algebra-lemma-flat-going-down}），可找到位于
$\mathfrak p$ 上方的素理想 $\mathfrak q \subset \mathfrak m'$。特别地，
$\mathfrak q$ 不是极大理想。因此再次由 Algebra, Lemma
\ref{algebra-lemma-characterize-jacobson}，集合
$V(\mathfrak q) \cap D(f)$ 是无限的（此处终于用到 $B$ 是 Jacobson 环）。
$V(\mathfrak q) \cap D(f)$ 的所有点都映到
$V(\mathfrak p) \cap D(f) = \{\mathfrak p\}$。因此
$\mathfrak p$ 上的纤维是无限的。这与 $A \to B$ 拟有限矛盾（见 Algebra,
Lemma \ref{algebra-lemma-quasi-finite}，或更明确地见 Morphisms, Lemma
\ref{morphisms-lemma-quasi-finite}）。引理得证。
\end{proof}

\begin{lemma}
\label{descent-lemma-locally-finite-nr-irred-local-fppf}
性质 $\mathcal{P}(S) =$``$S$ 的每个拟紧开集都只有有限多个不可约分支''
在 fppf 拓扑中为局部的。
\end{lemma}

\begin{proof}
我们使用 Lemma \ref{descent-lemma-descending-properties}。首先注意，
$\mathcal{P}$ 在 Zariski 拓扑中为局部的。接着，若 $T \to S$ 是仿射概形
之间的平坦有限表示态射，而 $S$ 只有有限多个不可约分支，则须证明
$T$ 也如此。事实上，由于 $T \to S$ 平坦，$T$ 的一般点映到 $S$ 的
一般点；见 Morphisms, Lemma
\ref{morphisms-lemma-generalizations-lift-flat}。所以只须证明，对
$s \in S$，纤维 $T_s$ 只有有限多个一般点。注意，$T_s$ 是
$\kappa(s)$ 上的有限型仿射概形；见 Morphisms, Lemma
\ref{morphisms-lemma-base-change-finite-type}。因此 $T_s$ 是 Noether
概形，并且只有有限多个不可约分支（Morphisms, Lemma
\ref{morphisms-lemma-finite-type-noetherian} 及 Properties, Lemma
\ref{properties-lemma-Noetherian-irreducible-components}）。最后，若
$T \to S$ 是仿射概形之间的满、平坦、有限表示态射，而 $T$ 只有有限
多个不可约分支，则须证明 $S$ 也如此。这由 Topology, Lemma
\ref{topology-lemma-surjective-continuous-irreducible-components} 得出。
因此 Lemma \ref{descent-lemma-descending-properties} 的 (1)、(2)、(3) 成立，
命题得证。
\end{proof}




\section{在 syntomic 拓扑中为局部的概形性质}
\label{descent-section-descending-properties-syntomic}

\noindent
本节给出一些在 syntomic 拓扑中于基底上为局部的概形性质。

\begin{lemma}
\label{descent-lemma-Sk-local-syntomic}
性质 $\mathcal{P}(S) =$``$S$ 局部 Noether 且满足 $(S_k)$'' 在 syntomic
拓扑中为局部的。
\end{lemma}

\begin{proof}
我们验证 Lemma \ref{descent-lemma-descending-properties} 的 (1)、(2)、(3)。由于
syntomic 态射平坦且有限表示（Morphisms, Lemmas
\ref{morphisms-lemma-syntomic-flat} 与
\ref{morphisms-lemma-syntomic-locally-finite-presentation}），在 Lemma
\ref{descent-lemma-Noetherian-local-fppf} 的证明中已经对 ``局部 Noether'' 验证
了这些条件；下文不再提及而直接使用。首先注意，$\mathcal{P}$ 在 Zariski
拓扑中为局部的。这由定义即得；见 Cohomology of Schemes, Definition
\ref{coherent-definition-depth}。接着，若 $S' \to S$ 是仿射概形之间的
syntomic 态射，而 $S$ 具有 $\mathcal{P}$，则 $S'$ 具有
$\mathcal{P}$。这是 Algebra, Lemma \ref{algebra-lemma-Sk-goes-up}
（使用 Morphisms, Lemma \ref{morphisms-lemma-syntomic-characterize}、
Algebra, Definition \ref{algebra-definition-lci} 以及 Lemma
\ref{algebra-lemma-lci-CM}）。最后，若 $S' \to S$ 是仿射概形之间的满
syntomic 态射，而 $S'$ 具有 $\mathcal{P}$，则 $S$ 具有
$\mathcal{P}$。这是 Algebra, Lemma \ref{algebra-lemma-descent-Sk}。
因此 Lemma \ref{descent-lemma-descending-properties} 的 (1)、(2)、(3) 成立，
命题得证。
\end{proof}

\begin{lemma}
\label{descent-lemma-CM-local-syntomic}
性质 $\mathcal{P}(S) =$``$S$ 是 Cohen--Macaulay 概形'' 在 syntomic 拓扑
中为局部的。
\end{lemma}

\begin{proof}
这由上面的 Lemma \ref{descent-lemma-Sk-local-syntomic} 立即得出，因为一个概形
是 Cohen--Macaulay 概形，当且仅当它局部 Noether，且对所有
$k \geq 0$ 都满足 $(S_k)$；见 Properties, Lemma
\ref{properties-lemma-scheme-CM-iff-all-Sk}。
\end{proof}







\section{在 smooth 拓扑中为局部的概形性质}
\label{descent-section-descending-properties-smooth}

\noindent
本节给出一些在 smooth 拓扑中于基底上为局部的概形性质。

\begin{lemma}
\label{descent-lemma-reduced-local-smooth}
性质 $\mathcal{P}(S) =$``$S$ 既约'' 在 smooth 拓扑中为局部的。
\end{lemma}

\begin{proof}
我们使用 Lemma \ref{descent-lemma-descending-properties}。首先注意，``既约'' 在
Zariski 拓扑中为局部的。这由定义即得；见 Schemes, Definition
\ref{schemes-definition-reduced}。接着，若 $S' \to S$ 是仿射概形之间的
smooth 态射，而 $S$ 既约，则 $S'$ 既约；这是 Algebra, Lemma
\ref{algebra-lemma-reduced-goes-up}。最后，若 $S' \to S$ 是仿射概形之间
的满 smooth 态射，而 $S'$ 既约，则 $S$ 既约；这是 Algebra, Lemma
\ref{algebra-lemma-descent-reduced}。因此 Lemma
\ref{descent-lemma-descending-properties} 的 (1)、(2)、(3) 成立，命题得证。
\end{proof}

\begin{lemma}
\label{descent-lemma-normal-local-smooth}
\begin{slogan}
正规性在 smooth 拓扑中为局部的。
\end{slogan}
性质 $\mathcal{P}(S) =$``$S$ 正规'' 在 smooth 拓扑中为局部的。
\end{lemma}

\begin{proof}
我们使用 Lemma \ref{descent-lemma-descending-properties}。首先证明 ``正规'' 在
Zariski 拓扑中为局部的。这由定义即得；见 Properties, Definition
\ref{properties-definition-normal}。接着，若 $S' \to S$ 是仿射概形之间的
smooth 态射，而 $S$ 正规，则 $S'$ 正规；这是 Algebra, Lemma
\ref{algebra-lemma-normal-goes-up}。最后，若 $S' \to S$ 是仿射概形之间的
满 smooth 态射，而 $S'$ 正规，则 $S$ 正规；这是 Algebra, Lemma
\ref{algebra-lemma-descent-normal}。因此 Lemma
\ref{descent-lemma-descending-properties} 的 (1)、(2)、(3) 成立，命题得证。
\end{proof}

\begin{lemma}
\label{descent-lemma-Rk-local-smooth}
性质 $\mathcal{P}(S) =$``$S$ 局部 Noether 且满足 $(R_k)$'' 在 smooth
拓扑中为局部的。
\end{lemma}

\begin{proof}
我们验证 Lemma \ref{descent-lemma-descending-properties} 的 (1)、(2)、(3)。由于
smooth 态射平坦且有限表示（Morphisms, Lemmas
\ref{morphisms-lemma-smooth-flat} 与
\ref{morphisms-lemma-smooth-locally-finite-presentation}），在 Lemma
\ref{descent-lemma-Noetherian-local-fppf} 的证明中已经对 ``局部 Noether'' 验证
了这些条件；下文不再提及而直接使用。首先注意，$\mathcal{P}$ 在 Zariski
拓扑中为局部的。这由定义即得；见 Properties, Definition
\ref{properties-definition-Rk}。接着，若 $S' \to S$ 是仿射概形之间的
smooth 态射，而 $S$ 具有 $\mathcal{P}$，则 $S'$ 具有
$\mathcal{P}$。这是 Algebra, Lemmas \ref{algebra-lemma-Rk-goes-up}
（使用 Morphisms, Lemma \ref{morphisms-lemma-smooth-characterize} 以及
Algebra, Lemmas \ref{algebra-lemma-base-change-smooth} 与
\ref{algebra-lemma-characterize-smooth-over-field}）。最后，若
$S' \to S$ 是仿射概形之间的满 smooth 态射，而 $S'$ 具有
$\mathcal{P}$，则 $S$ 具有 $\mathcal{P}$。这是 Algebra, Lemma
\ref{algebra-lemma-descent-Rk}。因此 Lemma
\ref{descent-lemma-descending-properties} 的 (1)、(2)、(3) 成立，命题得证。
\end{proof}

\begin{lemma}
\label{descent-lemma-regular-local-smooth}
性质 $\mathcal{P}(S) =$``$S$ 正则'' 在 smooth 拓扑中为局部的。
\end{lemma}

\begin{proof}
这由上面的 Lemma \ref{descent-lemma-Rk-local-smooth} 立即得出，因为局部 Noether
概形正则，当且仅当它局部 Noether，且对所有 $k \geq 0$ 都满足
$(R_k)$。
\end{proof}

\begin{lemma}
\label{descent-lemma-Nagata-local-smooth}
性质 $\mathcal{P}(S) =$``$S$ 是 Nagata 概形'' 在 smooth 拓扑中为局部的。
\end{lemma}

\begin{proof}
我们验证 Lemma \ref{descent-lemma-descending-properties} 的 (1)、(2)、(3)。首先
注意，Nagata 性在 Zariski 拓扑中为局部的；这是 Properties, Lemma
\ref{properties-lemma-locally-nagata}。接着，若 $S' \to S$ 是仿射概形
之间的 smooth 态射，而 $S$ 是 Nagata 概形，则 $S'$ 是 Nagata 概形；
这是 Morphisms, Lemma \ref{morphisms-lemma-finite-type-nagata}。最后，
若 $S' \to S$ 是仿射概形之间的满 smooth 态射，而 $S'$ 是 Nagata
概形，则 $S$ 是 Nagata 概形；这是 Algebra, Lemma
\ref{algebra-lemma-descent-nagata}。因此 Lemma
\ref{descent-lemma-descending-properties} 的 (1)、(2)、(3) 成立，命题得证。
\end{proof}





\section{下降性质的变体}
\label{descent-section-variants}

\noindent
有时可以下降并非局部的性质。本节汇集这类结果。关于光滑性等态射性质
的下降，另见 Section \ref{descent-section-descent-finiteness-morphisms}。

\begin{lemma}
\label{descent-lemma-descend-reduced}
若 $f : X \to Y$ 是概形的平坦满态射，且 $X$ 既约，则 $Y$ 既约。
\end{lemma}

\begin{proof}
考察局部环（Schemes, Definition \ref{schemes-definition-reduced}），并应用
Algebra, Lemma \ref{algebra-lemma-descent-reduced}，即得结论。
\end{proof}

\begin{lemma}
\label{descent-lemma-descend-regular}
设 $f : X \to Y$ 为代数空间态射。若 $f$ 局部有限表示、平坦且满，并且
$X$ 正则，则 $Y$ 正则。
\end{lemma}

\begin{proof}
本引理归约为下列代数陈述：若 $A \to B$ 是忠实平坦的有限表示环同态，
且 $B$ Noether 并正则，则 $A$ Noether 并正则。$A$ Noether 由 Algebra,
Lemma \ref{algebra-lemma-descent-Noetherian} 得出，而其正则性由 Algebra,
Lemma \ref{algebra-lemma-flat-under-regular} 得出。
\end{proof}









\section{概形芽}
\label{descent-section-germs}

\begin{definition}
\label{descent-definition-germs}
概形芽。
\begin{enumerate}
\item 由概形 $X$ 与点 $x \in X$ 组成的偶对 $(X, x)$ 称为
{\it $X$ 在 $x$ 处的芽}。
\item {\it 芽态射} $f : (X, x) \to (S, s)$ 是满足 $f(x) = s$ 的概形
态射 $f : U \to S$ 的等价类，其中 $U \subset X$ 是 $x$ 的开邻域。
两个这样的 $f$、$f'$ 等价，当且仅当 $f$ 与 $f'$ 在 $x$ 的某个开邻域
上一致。
\item 通过复合代表元定义{\it 芽态射的复合}（这是良定义的）。
\end{enumerate}
\end{definition}

\noindent
继续之前还需要一个定义。

\begin{definition}
\label{descent-definition-etale-morphism-germs}
设 $f : (X, x) \to (S, s)$ 为芽态射。若 $f$ 存在一个代表元
$f : U \to S$，它是概形的 \'etale 态射（相应地，smooth 态射），则称
$f$ {\it \'etale}（相应地，{\it smooth}）。
\end{definition}








\section{芽的局部性质}
\label{descent-section-properties-germs-local}

\begin{definition}
\label{descent-definition-local-at-point}
设 $\mathcal{P}$ 为概形芽的一项性质。若对任意 \'etale（相应地，smooth）
芽态射 $(U', u') \to (U, u)$，都有
$\mathcal{P}(U, u) \Leftrightarrow \mathcal{P}(U', u')$，则称
$\mathcal{P}$ {\it \'etale 局部}（相应地，{\it smooth 局部}）。
\end{definition}

\noindent
设 $(X, x)$ 为概形芽。$X$ 在 $x$ 处的维数，是 $x$ 在 $X$ 中各开邻域
维数的最小值；任意充分小的开邻域都具有这一维数。因此它是该芽同构类
的不变量，简记为 $\dim_x(X)$。下一个引理说明，断言
$\dim_x(X) = d$ 是芽的 \'etale 局部性质。

\begin{lemma}
\label{descent-lemma-dimension-at-point-local}
设 $f : U \to V$ 为概形的 \'etale 态射。设 $u \in U$ 且
$v = f(u)$。则 $\dim_u(U) = \dim_v(V)$。
\end{lemma}

\begin{proof}
陈述中的 $\dim_u(U)$ 是 Topology, Definition
\ref{topology-definition-Krull} 所定义的 $U$ 在 $u$ 处的维数，即 $u$
在 $U$ 中各开邻域的 Krull 维数的最小值。$\dim_v(V)$ 也类似。

\medskip\noindent
先证明 $\dim_v(V) \geq \dim_u(U)$。设 $V'$ 为 $v$ 在 $V$ 中的开邻域。
则存在 $u$ 在 $U$ 中的开邻域 $U'$，它包含于 $f^{-1}(V')$，且
$\dim_u(U) = \dim(U')$。设
$Z_0 \subset Z_1 \subset \ldots \subset Z_n$ 为 $U'$ 中不可约闭子概形的
一条链。若 $\xi_i \in Z_i$ 是一般点，则有特化
$\xi_n \leadsto \xi_{n - 1} \leadsto \ldots \leadsto \xi_0$.
这在 $V'$ 中给出特化
$f(\xi_n) \leadsto f(\xi_{n - 1}) \leadsto \ldots \leadsto f(\xi_0)$
。当 $i \not = j$ 时，$f(\xi_j) \not = f(\xi_i)$，因为 $f$ 的纤维离散
（见 Morphisms, Lemma \ref{morphisms-lemma-etale-over-field}）。所以
$\dim(V') \geq n$。形式地即可推出不等式
$\dim_v(V) \geq \dim_u(U)$。

\medskip\noindent
再证明 $\dim_u(U) \geq \dim_v(V)$。设 $U'$ 为 $u$ 在 $U$ 中的开邻域。
由 Morphisms, Lemma \ref{morphisms-lemma-fppf-open}，
$V' = f(U')$ 是 $v$ 的开邻域。因此 $\dim(V') \geq \dim_v(V)$。选取
$V'$ 中不可约闭子概形的一条链
$Z_0 \subset Z_1 \subset \ldots \subset Z_n$。令
$\xi_i \in Z_i$ 为一般点，于是有特化
$\xi_n \leadsto \xi_{n - 1} \leadsto \ldots \leadsto \xi_0$.
由于 $\xi_0 \in f(U')$，可找到一点 $\eta_0 \in U'$，满足
$f(\eta_0) = \xi_0$。考虑局部环同态
$$
\mathcal{O}_{V', \xi_0} \longrightarrow \mathcal{O}_{U', \eta_0}
$$
，它由 Morphisms, Lemma \ref{morphisms-lemma-etale-flat} 是平坦局部环同态。
注意，由 Schemes, Lemma \ref{schemes-lemma-specialize-points}，各点
$\xi_i$ 对应于左侧环的素理想。因此，对上示环同态应用下降定理（见
Algebra, Section \ref{algebra-section-going-up}），可以在 $U'$ 中找到一列
特化
$\eta_n \leadsto \eta_{n - 1} \leadsto \ldots \leadsto \eta_0$
，它们在 $f$ 下映到序列
$\xi_n \leadsto \xi_{n - 1} \leadsto \ldots \leadsto \xi_0$
。这蕴含 $\dim_u(U) \geq \dim_v(V)$。
\end{proof}

\noindent
设 $(X, x)$ 为概形芽。局部环 $\mathcal{O}_{X, x}$ 的同构类是该芽的
不变量。下一个引理说明，性质 $\dim(\mathcal{O}_{X, x}) = d$ 是芽的
\'etale 局部性质。

\begin{lemma}
\label{descent-lemma-dimension-local-ring-local}
设 $f : U \to V$ 为概形的 \'etale 态射。设 $u \in U$ 且
$v = f(u)$。则
$\dim(\mathcal{O}_{U, u}) = \dim(\mathcal{O}_{V, v})$.
\end{lemma}

\begin{proof}
须证明的代数陈述如下：若 $A \to B$ 是 \'etale 环同态，且
$\mathfrak q$ 是 $B$ 的一个素理想，位于 $\mathfrak p \subset A$ 上方，
则 $\dim(A_{\mathfrak p}) = \dim(B_{\mathfrak q})$。这是 More on Algebra,
Lemma \ref{more-algebra-lemma-dimension-etale-extension}。
\end{proof}

\noindent
设 $(X, x)$ 为概形芽。局部环 $\mathcal{O}_{X, x}$ 的同构类是该芽的
不变量。下一个引理说明，性质 ``$\mathcal{O}_{X, x}$ 正则'' 是芽的
\'etale 局部性质。

\begin{lemma}
\label{descent-lemma-regular-local-ring-local}
设 $f : U \to V$ 为概形的 \'etale 态射。设 $u \in U$ 且
$v = f(u)$。则 $\mathcal{O}_{U, u}$ 是正则局部环，当且仅当
$\mathcal{O}_{V, v}$ 是正则局部环。
\end{lemma}

\begin{proof}
须证明的代数陈述如下：若 $A \to B$ 是 \'etale 环同态，且
$\mathfrak q$ 是 $B$ 的一个素理想，位于 $\mathfrak p \subset A$ 上方，
则 $A_{\mathfrak p}$ 正则，当且仅当 $B_{\mathfrak q}$ 正则。这是 More on
Algebra, Lemma \ref{more-algebra-lemma-regular-etale-extension}。
\end{proof}







\section{在目标上为局部的态射性质}
\label{descent-section-descending-properties-morphisms}

\noindent
设 $f : X \to Y$ 为概形态射，设 $g : Y' \to Y$ 为概形态射，并设
$f' : X' \to Y'$ 为 $f$ 沿 $g$ 的基变换：
$$
\xymatrix{
X' \ar[d]_{f'} \ar[r]_{g'} & X \ar[d]^f \\
Y' \ar[r]^g & Y
}
$$
设 $\mathcal{P}$ 为概形态射的一项性质。可以问 (a)
$\mathcal{P}(f) \Rightarrow \mathcal{P}(f')$ 是否成立，也可以问其逆命题
(b) $\mathcal{P}(f') \Rightarrow \mathcal{P}(f)$ 是否成立。若每当 $g$
平坦时 (a) 都成立，则称 $\mathcal{P}$ 在平坦基变换下保持。若每当 $g$
满且平坦时 (b) 都成立，则称 $\mathcal{P}$ 沿满平坦基变换下降。若
$\mathcal{P}$ 在平坦基变换下保持，并沿满平坦基变换下降，则称
$\mathcal{P}$ 在目标上为平坦局部的。比较 Section
\ref{descent-section-descending-properties} 中的讨论。这是一个非常重要的概念，
正式定义如下。

\begin{definition}
\label{descent-definition-property-morphisms-local}
设 $\mathcal{P}$ 为基底上概形态射的一项性质，并设\linebreak[4]
$\tau \in \{fpqc, fppf, syntomic, smooth, \etale, Zariski\}$。若对任意
$\tau$-覆盖 $\{Y_i \to Y\}_{i \in I}$（见 Topologies, Section
\ref{topologies-section-procedure}）以及 $S$ 上的任意概形态射
$f : X \to Y$，都有
$$
f \text{ 具有 }\mathcal{P}
\Leftrightarrow
\text{每个 }Y_i \times_Y X \to Y_i\text{ 具有 }\mathcal{P}.
$$
则称 $\mathcal{P}$ {\it 在基底上为 $\tau$-局部的}，或
{\it 在目标上为 $\tau$-局部的}，或
{\it 关于 $\tau$-拓扑在基底上为局部的}。
\end{definition}

\noindent
确切地说，由于同构总是覆盖，可知（或要求）$X \to Y$ 具有性质
$\mathcal{P}$，当且仅当任意与 $X \to Y$ 同构的箭头 $X' \to Y'$ 具有
该性质。若一个性质在目标上为 $\tau$-局部的，则沿 $\tau$-覆盖中出现的
态射作基变换时，该性质保持。正式陈述如下。

\begin{lemma}
\label{descent-lemma-pullback-property-local-target}
设 $\tau \in \{fpqc, fppf, syntomic, smooth, \etale, Zariski\}$。设
$\mathcal{P}$ 为在目标上 $\tau$-局部的态射性质，并设 $f : X \to Y$
具有性质 $\mathcal{P}$。对任意态射 $Y' \to Y$，若它相应地为平坦、平坦
且局部有限表示、syntomic、\'etale、开浸入，则 $f$ 的基变换
% P05-ERR-0035
$f' : Y' \times_Y X \to Y'$ 具有性质 $\mathcal{P}$。
\end{lemma}

\begin{proof}
这是因为可以把 $Y' \to Y$ 加入一族构成 $\tau$-覆盖的态射中。
\end{proof}

\noindent
上述结果有一个经常使用的简单推论：若 $f : X \to Y$ 具有在目标上
$\tau$-局部的性质 $\mathcal{P}$，并且对某个开子概形 $V \subset Y$ 有
$f(X) \subset V$，则所诱导的态射 $X \to V$ 也具有 $\mathcal{P}$。
证明：沿 $V \to Y$ 对 $f$ 作基变换即得 $X \to V$。

\begin{lemma}
\label{descent-lemma-largest-open-of-the-base}
设 $\tau \in \{fppf, syntomic, smooth, \etale\}$，并设 $\mathcal{P}$ 为在
目标上 $\tau$-局部的态射性质。对任意概形态射 $f : X \to Y$，存在最大
开集 $W(f) \subset Y$，使得限制 $X_{W(f)} \to W(f)$ 具有
$\mathcal{P}$。此外，
\begin{enumerate}
\item 若 $g : Y' \to Y$ 平坦且局部有限表示、syntomic、smooth 或
\'etale，并且基变换 $f' : X_{Y'} \to Y'$ 具有 $\mathcal{P}$，则
$g(Y') \subset W(f)$；
\item 若 $g : Y' \to Y$ 平坦且局部有限表示、syntomic、smooth 或
\'etale，则 $W(f') = g^{-1}(W(f))$；
\item 若 $\{g_i : Y_i \to Y\}$ 是 $\tau$-覆盖，则
$g_i^{-1}(W(f)) = W(f_i)$，其中 $f_i$ 是 $f$ 沿 $Y_i \to Y$ 的基变换。
\end{enumerate}
\end{lemma}

\begin{proof}
考虑所有满足下列性质的态射 $g : Y' \to Y$ 的像 $g(Y') \subset Y$ 之并
$W$：
\begin{enumerate}
\item $g$ 平坦且局部有限表示、syntomic、smooth 或 \'etale；
\item 基变换 $Y' \times_{g, Y} X \to Y'$ 具有性质 $\mathcal{P}$。
\end{enumerate}
这类态射 $g$ 是开映射（见 Morphisms, Lemma
\ref{morphisms-lemma-fppf-open}），所以 $W \subset Y$ 是 $Y$ 的开子集。
由于 $\mathcal{P}$ 在 $\tau$-拓扑中为局部的，且给定了 $W$ 的覆盖
$\{Y' \to W\}$，使得各拉回具有 $\mathcal{P}$，故限制 $X_W \to W$
具有性质 $\mathcal{P}$。这证明了存在性，也证明了 $W(f)$ 具有性质 (1)。
为证明性质 (2)，注意 $W(f') \supset g^{-1}(W(f))$，因为在沿平坦且局部
有限表示、syntomic、smooth 或 \'etale 态射作基变换时，
$\mathcal{P}$ 保持；见 Lemma \ref{descent-lemma-pullback-property-local-target}。
另一方面，若 $Y'' \subset Y'$ 为开集，且 $X_{Y''} \to Y''$ 具有性质
$\mathcal{P}$，则由构造，$Y'' \to Y$ 通过 $W$ 分解；也就是说，
$Y'' \subset g^{-1}(W(f))$。这证明了 (2)。断言 (3) 由 (2) 得出，因为按
$\tau$-覆盖的定义，每个态射 $Y_i \to Y$ 都平坦且局部有限表示、
syntomic、smooth 或 \'etale。
\end{proof}

\begin{lemma}
\label{descent-lemma-descending-properties-morphisms}
设 $\mathcal{P}$ 为基底上概形态射的一项性质，并设
$\tau \in \{fpqc, fppf, \etale, smooth, syntomic\}$。假设
\begin{enumerate}
\item 按照 $\tau$ 为 fpqc、fppf、\'etale、smooth 或 syntomic，该性质
相应地在平坦、平坦且局部有限表示、\'etale、smooth 或 syntomic 基变换
下保持（比较 Schemes, Definition
\ref{schemes-definition-preserved-by-base-change}）；
\item 该性质在基底上为 Zariski 局部的；
\item 对任意满的仿射概形态射 $S' \to S$，若按照 $\tau$ 为 fpqc、fppf、
\'etale、smooth 或 syntomic，该态射相应地为平坦、平坦且有限表示、
\'etale、smooth 或 syntomic，并且对任意概形态射 $f : X \to S$，只要
基变换 $f' : X' = S' \times_S X \to S'$ 具有性质 $\mathcal{P}$，
$f$ 就具有性质 $\mathcal{P}$。
\end{enumerate}
则 $\mathcal{P}$ 在基底上为 $\tau$-局部的。
\end{lemma}

\begin{proof}
这几乎立即由 $\tau$-覆盖的定义得出；见 Topologies, Definition
\ref{topologies-definition-fpqc-covering}
\ref{topologies-definition-fppf-covering}
\ref{topologies-definition-etale-covering}
\ref{topologies-definition-smooth-covering}，或
\ref{topologies-definition-syntomic-covering}
以及 Topologies, Lemma
\ref{topologies-lemma-fpqc-affine},
\ref{topologies-lemma-fppf-affine},
\ref{topologies-lemma-etale-affine},
\ref{topologies-lemma-smooth-affine}，或
\ref{topologies-lemma-syntomic-affine}。省略细节。
\end{proof}

\begin{remark}
\label{descent-remark-descending-properties-morphisms-standard}
（这是对上面 Remark \ref{descent-remark-descending-properties-standard} 的重复。）
在上面的 Lemma \ref{descent-lemma-descending-properties-morphisms} 中，若
$\tau = smooth$，则在条件 (3) 中可以假设该态射是（满的）标准 smooth
态射。$\tau = syntomic$ 或 $\tau = \etale$ 时也类似。
\end{remark}




\section{在目标的 fpqc 拓扑中为局部的态射性质}
\label{descent-section-descending-properties-morphisms-fpqc}

\noindent
本节给出大量在 fpqc 拓扑中于基底上为局部的概形态射性质。与此相反，
我们将在 Examples, Section
\ref{examples-section-non-descending-property-projective} 中证明，即使在
Zariski 拓扑中，``射影'' 与 ``拟射影'' 也不在基底上为局部的。

\begin{lemma}
\label{descent-lemma-descending-property-quasi-compact}
性质 $\mathcal{P}(f) =$``$f$ 拟紧'' 在基底上为 fpqc 局部的。
\end{lemma}

\begin{proof}
拟紧态射的基变换拟紧；见 Schemes, Lemma
\ref{schemes-lemma-quasi-compact-preserved-base-change}。拟紧性在基底上为
Zariski 局部的；见 Schemes, Lemma
\ref{schemes-lemma-quasi-compact-affine}。最后，设 $S' \to S$ 为仿射概形
之间的平坦满态射，并设 $f : X \to S$ 为态射。假设基变换
$f' : X' \to S'$ 拟紧。则 $X'$ 拟紧，而 $X' \to X$ 满，所以 $X$
拟紧。这蕴含 $f$ 拟紧。因此应用 Lemma
\ref{descent-lemma-descending-properties-morphisms} 即得结论。
\end{proof}

\begin{lemma}
\label{descent-lemma-descending-property-quasi-separated}
性质 $\mathcal{P}(f) =$``$f$ 拟分离'' 在基底上为 fpqc 局部的。
\end{lemma}

\begin{proof}
拟分离态射的任意基变换拟分离；见 Schemes, Lemma
\ref{schemes-lemma-separated-permanence}。拟分离性在基底上为 Zariski
局部的（由定义，或由 Schemes, Lemma
\ref{schemes-lemma-characterize-quasi-separated}）。最后，设
$S' \to S$ 为仿射概形之间的平坦满态射，并设 $f : X \to S$ 为态射。
假设基变换 $f' : X' \to S'$ 拟分离。这意味着
$\Delta' : X' \to X'\times_{S'} X'$ 拟紧。注意，$\Delta'$ 是
$\Delta : X \to X \times_S X$ 沿 $S' \to S$ 的基变换。由 Lemma
\ref{descent-lemma-descending-property-quasi-compact}，这蕴含 $\Delta$ 拟紧，
所以 $f$ 拟分离。因此应用 Lemma
\ref{descent-lemma-descending-properties-morphisms} 即得结论。
\end{proof}

\begin{lemma}
\label{descent-lemma-descending-property-universally-closed}
性质 $\mathcal{P}(f) =$``$f$ 普遍闭'' 在基底上为 fpqc 局部的。
\end{lemma}

\begin{proof}
按定义，普遍闭态射的基变换普遍闭。普遍闭性在基底上为 Zariski 局部的
（由定义，或由 Morphisms, Lemma
\ref{morphisms-lemma-universally-closed-local-on-the-base}）。最后，设
$S' \to S$ 为仿射概形之间的平坦满态射，并设 $f : X \to S$ 为态射。
假设基变换 $f' : X' \to S'$ 普遍闭。设 $T \to S$ 为任意态射。考虑图表
$$
\xymatrix{
X' \ar[d] &
S' \times_S T \times_S X \ar[d] \ar[r] \ar[l] &
T \times_S X \ar[d] \\
S' &
S' \times_S T \ar[r] \ar[l] &
T
}
$$
，其中两个方块都是 Cartesian 的。因此假设蕴含中间的竖直箭头为闭
映射。右侧的水平箭头平坦、拟紧且满（因为它们是 $S' \to S$ 的基变换）。
所以 $T$ 的一个子集为闭集，当且仅当它在 $S' \times_S T$ 中的逆像为
闭集；见 Morphisms, Lemma
\ref{morphisms-lemma-fpqc-quotient-topology}。简单的图追踪表明右侧竖直
箭头也为闭映射，所以 $X \to S$ 普遍闭。因此应用 Lemma
\ref{descent-lemma-descending-properties-morphisms} 即得结论。
\end{proof}

\begin{lemma}
\label{descent-lemma-descending-property-universally-open}
性质 $\mathcal{P}(f) =$``$f$ 普遍开'' 在基底上为 fpqc 局部的。
\end{lemma}

\begin{proof}
证明与 Lemma \ref{descent-lemma-descending-property-universally-closed} 的证明相同。
\end{proof}

\begin{lemma}
\label{descent-lemma-descending-property-universally-submersive}
性质 $\mathcal{P}(f) =$``$f$ 普遍为商映射'' 在基底上为 fpqc 局部的。
\end{lemma}

\begin{proof}
证明与 Lemma \ref{descent-lemma-descending-property-universally-closed} 的证明相同；
其中使用 Morphisms, Lemma
\ref{morphisms-lemma-fpqc-quotient-topology} 所给出的事实：拟紧平坦满态射
普遍为商映射。
\end{proof}

\begin{lemma}
\label{descent-lemma-descending-property-separated}
性质 $\mathcal{P}(f) =$``$f$ 分离'' 在基底上为 fpqc 局部的。
\end{lemma}

\begin{proof}
分离态射的基变换分离；见 Schemes, Lemma
\ref{schemes-lemma-separated-permanence}。分离性在基底上为 Zariski 局部
的（由定义，或由 Schemes, Lemma
\ref{schemes-lemma-characterize-separated}）。最后，设 $S' \to S$ 为仿射
概形之间的平坦满态射，并设 $f : X \to S$ 为态射。假设基变换
$f' : X' \to S'$ 分离。这意味着
$\Delta' : X' \to X'\times_{S'} X'$ 是闭浸入，因而普遍闭。注意，
$\Delta'$ 是 $\Delta : X \to X \times_S X$ 沿 $S' \to S$ 的基变换。
由 Lemma \ref{descent-lemma-descending-property-universally-closed}，这蕴含
$\Delta$ 普遍闭。由于它是浸入（Schemes, Lemma
\ref{schemes-lemma-diagonal-immersion}），可知 $\Delta$ 是闭浸入。所以
$f$ 分离。因此应用 Lemma \ref{descent-lemma-descending-properties-morphisms}
即得结论。
\end{proof}

\begin{lemma}
\label{descent-lemma-descending-property-surjective}
性质 $\mathcal{P}(f) =$``$f$ 满'' 在基底上为 fpqc 局部的。
\end{lemma}

\begin{proof}
这是显然的。
\end{proof}

\begin{lemma}
\label{descent-lemma-descending-property-dominant}
性质 $\mathcal{P}(f) =$``$f$ 拟紧且支配'' 在基底上为 fpqc 局部的。
\end{lemma}

\begin{proof}
由 Morphisms, Lemma \ref{morphisms-lemma-base-change-dominant}，
拟紧支配态射在平坦拉回下得以保持。反方向较为容易。事实上，由 Lemma
\ref{descent-lemma-descending-property-quasi-compact}，拟紧性在基底上为 fpqc
局部的。支配性显然在基底上为 Zariski 局部的。最后，设 $S' \to S$
为概形的满态射（不必平坦），并设 $f : X \to S$ 为态射。假设基变换
$f' : X' \to S'$ 支配。于是 $X' \to S'\to S$ 是两个支配态射的复合，
因而支配。由于这也是复合 $X' \to X \to S$，可知 $X\to S$ 支配。
\end{proof}

\noindent
一般而言，支配态射在平坦拉回下未必保持，但下面这个特殊情形值得一提。
另见 Lemma \ref{descent-lemma-descending-fppf-property-dominant}。

\begin{lemma}
\label{descent-lemma-descending-property-dominant-over-field}
设 $E/k$ 为域扩张。则 $k$ 上的态射 $X \to Y$ 支配，当且仅当拉回
$X_E \to Y_E$ 支配。
\end{lemma}

\begin{proof}
由 Morphisms, Lemma
\ref{morphisms-lemma-scheme-over-field-universally-open}，态射
$\Spec(E) \to \Spec(k)$ 普遍开。因此 $Y_E \to Y$ 开。所以，若
$X \to Y$ 支配，则 Morphisms, Lemma
\ref{morphisms-lemma-open-base-change-dominant} 给出 $X_E \to Y_E$
支配。反之，假设 $X_E \to Y_E$ 支配。由 Morphisms, Lemma
\ref{morphisms-lemma-base-change-surjective}，态射 $Y_E \to Y$ 满，故复合
$X_E \to Y_E \to Y$ 支配。这也是复合 $X_E \to X \to Y$，所以
$X \to Y$ 支配。
\end{proof}

\begin{lemma}
\label{descent-lemma-descending-property-universally-injective}
性质 $\mathcal{P}(f) =$``$f$ 普遍单射'' 在基底上为 fpqc 局部的。
\end{lemma}

\begin{proof}
普遍单射态射的基变换普遍单射（这是形式的）。普遍单射性在基底上为
Zariski 局部的；这由定义显然可见。最后，设 $S' \to S$ 为仿射概形之间
的平坦满态射，并设 $f : X \to S$ 为态射。假设基变换
$f' : X' \to S'$ 普遍单射。设 $K$ 为域，并设
$a, b : \Spec(K) \to X$ 为两个满足 $f \circ a = f \circ b$ 的态射。
由于 $S' \to S$ 满，并由 Schemes, Section
\ref{schemes-section-points} 中的讨论，存在域扩张 $K'/K$ 和态射 $\Spec(K')
\to S'$，使下列实线图交换：
$$
\xymatrix{
\Spec(K') \ar[rrd] \ar@{-->}[rd]_{a', b'} \ar[dd] \\
 &
X' \ar[r] \ar[d] &
S' \ar[d] \\
\Spec(K) \ar[r]^{a, b} &
X \ar[r] &
S
}
$$
由于方块为 Cartesian 的，我们得到使图交换的两条虚线箭头 $a'$、$b'$。
因为 $X' \to S'$ 普遍单射，由 Morphisms, Lemma
\ref{morphisms-lemma-universally-injective} 可得 $a' = b'$。显然这迫使
$a = b$（由 Schemes, Section \ref{schemes-section-points} 中的讨论）。
因此应用 Lemma \ref{descent-lemma-descending-properties-morphisms} 即得结论。

\medskip\noindent
另一个证明可使用如下刻画：一个态射普遍单射，当且仅当其对角态射满；见
Morphisms, Lemma \ref{morphisms-lemma-universally-injective}。于是本引理由
满性在基底上为 fpqc 局部这一事实推出；见 Lemma
\ref{descent-lemma-descending-property-surjective}。（提示：使用对角态射的基变换
就是基变换的对角态射这一事实。）
\end{proof}

\begin{lemma}
\label{descent-lemma-descending-property-universal-homeomorphism}
性质 $\mathcal{P}(f) =$``$f$ 为普遍同胚'' 在基底上为 fpqc 局部的。
\end{lemma}

\begin{proof}
这可用与 Lemma \ref{descent-lemma-descending-property-universally-closed} 完全相同的
方式证明。或者，可以使用如下事实：拓扑空间的映射为同胚，当且仅当它
单射、满且开。因此，普遍同胚恰为满、普遍单射且普遍开的态射。于是本引理
由 Lemmas \ref{descent-lemma-descending-property-surjective}、
\ref{descent-lemma-descending-property-universally-injective} 和
\ref{descent-lemma-descending-property-universally-open} 推出。
\end{proof}

\begin{lemma}
\label{descent-lemma-descending-property-locally-finite-type}
性质 $\mathcal{P}(f) =$``$f$ 局部有限型'' 在基底上为 fpqc 局部的。
\end{lemma}

\begin{proof}
局部有限型性在基变换下保持；见 Morphisms, Lemma
\ref{morphisms-lemma-base-change-finite-type}。局部有限型性在基底上为
Zariski 局部的；见 Morphisms, Lemma
\ref{morphisms-lemma-locally-finite-type-characterize}。最后，设
$S' \to S$ 为仿射概形之间的平坦满态射，并设 $f : X \to S$ 为态射。
假设基变换 $f' : X' \to S'$ 局部有限型。设 $U \subset X$ 为仿射开集。
则 $U' = S' \times_S U$ 是 $S'$ 上的有限型仿射概形。写成
$S = \Spec(R)$,
$S' = \Spec(R')$,
$U = \Spec(A)$, and
$U' = \Spec(A')$.
我们知道 $R \to R'$ 忠实平坦，$A' = R' \otimes_R A$，并且
$R' \to A'$ 有限型。我们必须证明 $R \to A$ 有限型。这正是 Algebra,
Lemma \ref{algebra-lemma-finite-type-descends} 的结论。因此 $f$ 局部有限型。
所以应用 Lemma \ref{descent-lemma-descending-properties-morphisms} 即得结论。
\end{proof}

\begin{lemma}
\label{descent-lemma-descending-property-locally-finite-presentation}
性质 $\mathcal{P}(f) =$``$f$ 局部有限表示'' 在基底上为 fpqc 局部的。
\end{lemma}

\begin{proof}
局部有限表示性在基变换下保持；见 Morphisms, Lemma
\ref{morphisms-lemma-base-change-finite-presentation}。
% P05-ERR-0036
局部有限型性在基底上为 Zariski 局部的；见 Morphisms, Lemma
\ref{morphisms-lemma-locally-finite-presentation-characterize}。最后，设
$S' \to S$ 为仿射概形之间的平坦满态射，并设 $f : X \to S$ 为态射。
假设基变换 $f' : X' \to S'$ 局部有限表示。设 $U \subset X$ 为仿射开集。
则 $U' = S' \times_S U$ 是 $S'$ 上的有限型仿射概形。写成
$S = \Spec(R)$,
$S' = \Spec(R')$,
$U = \Spec(A)$, and
$U' = \Spec(A')$.
我们知道 $R \to R'$ 忠实平坦，$A' = R' \otimes_R A$，并且
$R' \to A'$ 有限表示。我们必须证明 $R \to A$ 有限表示。这正是
Algebra, Lemma \ref{algebra-lemma-finite-presentation-descends} 的结论。
因此 $f$ 局部有限表示。所以应用 Lemma
\ref{descent-lemma-descending-properties-morphisms} 即得结论。
\end{proof}

\begin{lemma}
\label{descent-lemma-descending-property-finite-type}
性质 $\mathcal{P}(f) =$``$f$ 有限型'' 在基底上为 fpqc 局部的。
\end{lemma}

\begin{proof}
合并 Lemmas \ref{descent-lemma-descending-property-quasi-compact} 和
\ref{descent-lemma-descending-property-locally-finite-type} 即得。
\end{proof}

\begin{lemma}
\label{descent-lemma-descending-property-finite-presentation}
性质 $\mathcal{P}(f) =$``$f$ 有限表示'' 在基底上为 fpqc 局部的。
\end{lemma}

\begin{proof}
合并 Lemmas \ref{descent-lemma-descending-property-quasi-compact}、
\ref{descent-lemma-descending-property-quasi-separated} 和
\ref{descent-lemma-descending-property-locally-finite-presentation} 即得。
\end{proof}

\begin{lemma}
\label{descent-lemma-descending-property-proper}
性质 $\mathcal{P}(f) =$``$f$ 固有'' 在基底上为 fpqc 局部的。
\end{lemma}

\begin{proof}
合并 Lemmas \ref{descent-lemma-descending-property-universally-closed}、
\ref{descent-lemma-descending-property-separated} 和
\ref{descent-lemma-descending-property-finite-type} 即得本引理。
\end{proof}

\begin{lemma}
\label{descent-lemma-descending-property-flat}
性质 $\mathcal{P}(f) =$``$f$ 平坦'' 在基底上为 fpqc 局部的。
\end{lemma}

\begin{proof}
平坦性在任意基变换下保持；见 Morphisms, Lemma
\ref{morphisms-lemma-base-change-flat}。由定义，平坦性在基底上为 Zariski
局部的。最后，设 $S' \to S$ 为仿射概形之间的平坦满态射，并设
$f : X \to S$ 为态射。假设基变换 $f' : X' \to S'$ 平坦。设
$U \subset X$ 为仿射开集，则 $U' = S' \times_S U$ 为仿射概形。写成
$S = \Spec(R)$,
$S' = \Spec(R')$,
$U = \Spec(A)$, and
$U' = \Spec(A')$.
我们知道 $R \to R'$ 忠实平坦，$A' = R' \otimes_R A$，并且
$R' \to A'$ 平坦。目标：证明 $R \to A$ 平坦。这立即由 Algebra,
Lemma \ref{algebra-lemma-flatness-descends} 推出。因此 $f$ 平坦。所以应用
Lemma \ref{descent-lemma-descending-properties-morphisms} 即得结论。
\end{proof}

\begin{lemma}
\label{descent-lemma-descending-property-open-immersion}
性质 $\mathcal{P}(f) =$``$f$ 为开浸入'' 在基底上为 fpqc 局部的。
\end{lemma}

\begin{proof}
开浸入性在基变换下稳定；见 Schemes, Lemma
\ref{schemes-lemma-base-change-immersion}。开浸入性在基底上为 Zariski
局部的（这是显然的）。

\medskip\noindent
设 $S' \to S$ 为仿射概形之间的平坦满态射，并设 $f : X \to S$ 为态射。
假设基变换 $f' : X' \to S'$ 为开浸入。我们断言 $f$ 为开浸入。此时
$f'$ 普遍开且普遍单射。因此由 Lemma
\ref{descent-lemma-descending-property-universally-open}，$f$ 普遍开；由 Lemma
\ref{descent-lemma-descending-property-universally-injective}，它普遍单射。特别地，
$f(X) \subset S$ 为开集。若对每个仿射开集 $U \subset f(X)$，我们都能证明
$f^{-1}(U) \to U$ 为同构，则 $f$ 为开浸入，结论即得。若以
$U' \subset S'$ 表示 $U$ 的逆像，则 $U' \to U$ 是仿射概形之间的忠实平坦
态射，且 $(f')^{-1}(U') \to U'$ 为同构（由于我们对 $U$ 的选取，
$f'(X')$ 包含 $U'$）。这样便归约到下一段讨论的情形。

\medskip\noindent
设 $S' \to S$ 为仿射概形之间的平坦满态射，设 $f : X \to S$ 为态射，并
假设基变换 $f' : X' \to S'$ 为同构。我们还必须证明 $f$ 为同构。显然，
$f$ 满、普遍单射且普遍开（后两者见上面的论证）。所以 $f$ 为双射，也即
$f$ 为同胚。因此由 Morphisms, Lemma
\ref{morphisms-lemma-homeomorphism-affine}，$f$ 为仿射态射。由于
$$
\mathcal{O}(S') \to
\mathcal{O}(X') =
\mathcal{O}(S') \otimes_{\mathcal{O}(S)} \mathcal{O}(X)
$$
为同构，并且 $\mathcal{O}(S) \to \mathcal{O}(S')$ 忠实平坦，这蕴含
$\mathcal{O}(S) \to \mathcal{O}(X)$ 为同构。因此 $f$ 为同构。这就完成了
上面断言的证明。所以应用 Lemma
\ref{descent-lemma-descending-properties-morphisms} 即得结论。
\end{proof}

\begin{lemma}
\label{descent-lemma-descending-property-isomorphism}
性质 $\mathcal{P}(f) =$``$f$ 为同构'' 在基底上为 fpqc 局部的。
\end{lemma}

\begin{proof}
合并 Lemmas \ref{descent-lemma-descending-property-surjective} 和
\ref{descent-lemma-descending-property-open-immersion} 即得。
\end{proof}

\begin{lemma}
\label{descent-lemma-descending-property-affine}
性质 $\mathcal{P}(f) =$``$f$ 为仿射态射'' 在基底上为 fpqc 局部的。
\end{lemma}

\begin{proof}
仿射态射的基变换为仿射态射；见 Morphisms, Lemma
\ref{morphisms-lemma-base-change-affine}。仿射性在基底上为 Zariski 局部的；
见 Morphisms, Lemma \ref{morphisms-lemma-characterize-affine}。最后，设
$g : S' \to S$ 为仿射概形之间的平坦满态射，并设 $f : X \to S$ 为态射。
假设基变换 $f' : X' \to S'$ 为仿射态射。换言之，$X'$ 为仿射概形，比如
$X' = \Spec(A')$。再写成 $S = \Spec(R)$ 和 $S' = \Spec(R')$。我们必须
证明 $X$ 为仿射概形。

\medskip\noindent
由 Lemmas \ref{descent-lemma-descending-property-quasi-compact} 和
\ref{descent-lemma-descending-property-separated}，$X \to S$ 分离且拟紧。因此
$f_*\mathcal{O}_X$ 是 $\mathcal{O}_S$-代数的拟凝聚层；见 Schemes,
Lemma \ref{schemes-lemma-push-forward-quasi-coherent}。所以对某个 $R$-代数
$A$，有 $f_*\mathcal{O}_X = \widetilde{A}$。当然，事实上
$A = \Gamma(X, \mathcal{O}_X)$。此外，由平坦基变换（例如见 Cohomology
of Schemes, Lemma \ref{coherent-lemma-flat-base-change-cohomology}），有
$g^*f_*\mathcal{O}_X = f'_*\mathcal{O}_{X'}$。换言之，
$A' = R' \otimes_R A$。考虑 Schemes, Lemma
\ref{schemes-lemma-morphism-into-affine} 中 $S$ 上的典范态射
$$
X \longrightarrow \Spec(A)
$$
。由上述结论，该态射到 $S'$ 的基变换为同构。因此由 Lemma
\ref{descent-lemma-descending-property-isomorphism}，它本身为同构。所以应用
Lemma \ref{descent-lemma-descending-properties-morphisms} 即得结论。
\end{proof}

\begin{lemma}
\label{descent-lemma-descending-property-closed-immersion}
性质 $\mathcal{P}(f) =$``$f$ 为闭浸入'' 在基底上为 fpqc 局部的。
\end{lemma}

\begin{proof}
设 $f : X \to Y$ 为概形态射，设 $\{Y_i \to Y\}$ 为 fpqc 覆盖。假设每个
$f_i : Y_i \times_Y X \to Y_i$ 都是闭浸入。这蕴含每个 $f_i$ 都是仿射
态射；见 Morphisms, Lemma
\ref{morphisms-lemma-closed-immersion-affine}。由 Lemma
\ref{descent-lemma-descending-property-affine}，可知 $f$ 为仿射态射。还需证明
$\mathcal{O}_Y \to f_*\mathcal{O}_X$ 满。对每个 $y \in Y$，存在 $i$ 和
映到 $y$ 的点 $y_i \in Y_i$。由 Cohomology of Schemes, Lemma
\ref{coherent-lemma-flat-base-change-cohomology}，层
$f_{i, *}(\mathcal{O}_{Y_i \times_Y X})$ 是 $f_*\mathcal{O}_X$ 的拉回。
由假设，它是 $\mathcal{O}_{Y_i}$ 的商。因此可见
$$
\Big(
\mathcal{O}_{Y, y} \longrightarrow (f_*\mathcal{O}_X)_y
\Big)
\otimes_{\mathcal{O}_{Y, y}} \mathcal{O}_{Y_i, y_i}
$$
满。由于 $\mathcal{O}_{Y_i, y_i}$ 在 $\mathcal{O}_{Y, y}$ 上忠实平坦，
这蕴含 $\mathcal{O}_{Y, y} \longrightarrow (f_*\mathcal{O}_X)_y$ 满，
正如所需。
\end{proof}

\begin{lemma}
\label{descent-lemma-descending-property-quasi-affine}
性质 $\mathcal{P}(f) =$``$f$ 为拟仿射态射'' 在基底上为 fpqc 局部的。
\end{lemma}

\begin{proof}
设 $f : X \to Y$ 为概形态射，设 $\{g_i : Y_i \to Y\}$ 为 fpqc 覆盖。
假设每个 $f_i : Y_i \times_Y X \to Y_i$ 都是拟仿射态射。这蕴含每个
$f_i$ 都拟紧且分离。由 Lemmas
\ref{descent-lemma-descending-property-quasi-compact} 和
\ref{descent-lemma-descending-property-separated}，这蕴含 $f$ 拟紧且分离。考虑
$\mathcal{O}_Y$-代数层 $\mathcal{A} = f_*\mathcal{O}_X$。由 Schemes,
Lemma \ref{schemes-lemma-push-forward-quasi-coherent}，它是拟凝聚
$\mathcal{O}_Y$-代数。考虑典范态射
$$
j : X \longrightarrow \underline{\Spec}_Y(\mathcal{A})
$$
；见 Constructions, Lemma \ref{constructions-lemma-canonical-morphism}。
由平坦基变换（例如见 Cohomology of Schemes, Lemma
\ref{coherent-lemma-flat-base-change-cohomology}），有
$g_i^*f_*\mathcal{O}_X = f_{i, *}\mathcal{O}_{X'}$，其中
$g_i : Y_i \to Y$ 是给定的平坦映射。因此，$j$ 沿 $g_i$ 的基变换 $j_i$
就是 Constructions, Lemma \ref{constructions-lemma-canonical-morphism}
中对应于态射 $f_i$ 的典范态射。由假设及 Morphisms, Lemma
\ref{morphisms-lemma-characterize-quasi-affine}，所有这些态射 $j_i$ 都是
拟紧开浸入。因此由 Lemmas \ref{descent-lemma-descending-property-quasi-compact}
和 \ref{descent-lemma-descending-property-open-immersion}，$j$ 为拟紧开浸入。再次
由 Morphisms, Lemma \ref{morphisms-lemma-characterize-quasi-affine}，可知
$f$ 为拟仿射态射。
\end{proof}

\begin{lemma}
\label{descent-lemma-descending-property-quasi-compact-immersion}
性质 $\mathcal{P}(f) =$``$f$ 为拟紧浸入'' 在基底上为 fpqc 局部的。
\end{lemma}

\begin{proof}
设 $f : X \to Y$ 为概形态射，设 $\{Y_i \to Y\}$ 为 fpqc 覆盖。写
$X_i = Y_i \times_Y X$，并以 $f_i : X_i \to Y_i$ 表示 $f$ 的基变换；
还以 $q_i : Y_i \to Y$ 表示给定的平坦态射。假设每个 $f_i$ 都是拟紧浸入。
由 Schemes, Lemma \ref{schemes-lemma-immersions-monomorphisms}，每个 $f_i$
都分离。由 Lemmas \ref{descent-lemma-descending-property-quasi-compact} 和
\ref{descent-lemma-descending-property-separated}，这蕴含 $f$ 拟紧且分离。令
$X \to Z \to Y$ 为 $f$ 经过其概形论像的分解。由于 $f$ 拟紧，由
Morphisms, Lemma
\ref{morphisms-lemma-quasi-compact-scheme-theoretic-image}，闭子概形
$Z \subset Y$ 由拟凝聚理想层
$\mathcal{I} = \Ker(\mathcal{O}_Y \to f_*\mathcal{O}_X)$ 截出。由平坦
基变换（例如见 Cohomology of Schemes, Lemma
\ref{coherent-lemma-flat-base-change-cohomology}；这里使用 $f$ 分离），
$f_{i, *}\mathcal{O}_{X_i}$ 是拉回 $q_i^*f_*\mathcal{O}_X$。因此，
$Y_i \times_Y Z$ 由拟凝聚理想层
$q_i^*\mathcal{I} =
\Ker(\mathcal{O}_{Y_i} \to f_{i, *}\mathcal{O}_{X_i})$ 截出。由
Morphisms, Lemma \ref{morphisms-lemma-quasi-compact-immersion}，态射
$X_i \to Y_i \times_Y Z$ 为开浸入。所以由 Lemma
\ref{descent-lemma-descending-property-open-immersion}，$X \to Z$ 为开浸入，
因而 $f$ 如所需为浸入（我们已知它拟紧）。
\end{proof}

\begin{lemma}
\label{descent-lemma-descending-property-integral}
性质 $\mathcal{P}(f) =$``$f$ 为整态射'' 在基底上为 fpqc 局部的。
\end{lemma}

\begin{proof}
整态射恰为仿射且普遍闭的态射；见 Morphisms, Lemma
\ref{morphisms-lemma-integral-universally-closed}。因此合并 Lemmas
\ref{descent-lemma-descending-property-universally-closed} 和
\ref{descent-lemma-descending-property-affine} 即得本引理。
\end{proof}

\begin{lemma}
\label{descent-lemma-descending-property-finite}
性质 $\mathcal{P}(f) =$``$f$ 为有限态射'' 在基底上为 fpqc 局部的。
\end{lemma}

\begin{proof}
有限态射恰为局部有限型的整态射；见 Morphisms, Lemma
\ref{morphisms-lemma-finite-integral}。因此合并 Lemmas
\ref{descent-lemma-descending-property-locally-finite-type} 和
\ref{descent-lemma-descending-property-integral} 即得本引理。
\end{proof}

\begin{lemma}
\label{descent-lemma-descending-property-quasi-finite}
性质 $\mathcal{P}(f) =$``$f$ 局部拟有限'' 和
$\mathcal{P}(f) =$``$f$ 拟有限'' 在基底上为 fpqc 局部的。
\end{lemma}

\begin{proof}
设 $f : X \to S$ 为概形态射，并设 $\{S_i \to S\}$ 为 fpqc 覆盖，使得每个
基变换 $f_i : X_i \to S_i$ 都局部拟有限。我们已经知道（Lemma
\ref{descent-lemma-descending-property-locally-finite-type}）``局部有限型'' 在基底上
为 fpqc 局部的，因而 $f$ 局部有限型。于是由 Morphisms, Lemma
\ref{morphisms-lemma-base-change-quasi-finite} 可知 $f$ 局部拟有限。拟有限的
情形也随之成立，因为我们已经知道``拟紧''在基底上为 fpqc 局部的（Lemma
\ref{descent-lemma-descending-property-quasi-compact}）。
\end{proof}

\begin{lemma}
\label{descent-lemma-descending-property-relative-dimension-d}
性质 $\mathcal{P}(f) =$``$f$ 局部有限型且相对维数为 $d$'' 在基底上为
fpqc 局部的。
\end{lemma}

\begin{proof}
这立即由局部有限型性在基底上为 fpqc 局部这一事实及 Morphisms, Lemma
\ref{morphisms-lemma-dimension-fibre-after-base-change} 推出。
\end{proof}

\begin{lemma}
\label{descent-lemma-descending-property-syntomic}
性质 $\mathcal{P}(f) =$``$f$ 为 syntomic 态射'' 在基底上为 fpqc 局部的。
\end{lemma}

\begin{proof}
一个态射为 syntomic 态射，当且仅当它局部有限表示、平坦，并且其纤维
为局部完全交。我们已经知道平坦性和局部有限表示性在基底上为 fpqc 局部
的（Lemmas \ref{descent-lemma-descending-property-flat} 和
\ref{descent-lemma-descending-property-locally-finite-presentation}）。因此 syntomic
的结论由 Morphisms, Lemma
\ref{morphisms-lemma-set-points-where-fibres-lci} 推出。
\end{proof}

\begin{lemma}
\label{descent-lemma-descending-property-smooth}
性质 $\mathcal{P}(f) =$``$f$ 光滑'' 在基底上为 fpqc 局部的。
\end{lemma}

\begin{proof}
一个态射光滑，当且仅当它局部有限表示、平坦，并且具有光滑纤维。我们已经
知道平坦性和局部有限表示性在基底上为 fpqc 局部的（Lemmas
\ref{descent-lemma-descending-property-flat} 和
\ref{descent-lemma-descending-property-locally-finite-presentation}）。因此光滑的
结论由 Morphisms, Lemma
\ref{morphisms-lemma-set-points-where-fibres-smooth} 推出。
\end{proof}

\begin{lemma}
\label{descent-lemma-descending-property-unramified}
性质 $\mathcal{P}(f) =$``$f$ 非分歧'' 在基底上为 fpqc 局部的。
性质 $\mathcal{P}(f) =$``$f$ 为 G-非分歧态射'' 在基底上为 fpqc 局部的。
\end{lemma}

\begin{proof}
一个态射非分歧（相应地，为 G-非分歧态射），当且仅当它局部有限型
（相应地，有限表示），并且其对角态射为开浸入（见 Morphisms, Lemma
\ref{morphisms-lemma-diagonal-unramified-morphism}）。我们已经知道，局部
有限型性（相应地，局部有限表示性）和开浸入性在基底上为 fpqc 局部的
（Lemmas \ref{descent-lemma-descending-property-locally-finite-presentation}、
\ref{descent-lemma-descending-property-locally-finite-type} 和
\ref{descent-lemma-descending-property-open-immersion}）。因此结论形式地成立。
\end{proof}

\begin{lemma}
\label{descent-lemma-descending-property-etale}
性质 $\mathcal{P}(f) =$``$f$ 为 \'etale 态射'' 在基底上为 fpqc 局部的。
\end{lemma}

\begin{proof}
一个态射为 \'etale 态射，当且仅当它平坦且为 G-非分歧态射；见
Morphisms, Lemma \ref{morphisms-lemma-flat-unramified-etale}。我们已经知道
平坦性和 G-非分歧性在基底上为 fpqc 局部的（Lemmas
\ref{descent-lemma-descending-property-flat} 和
\ref{descent-lemma-descending-property-unramified}）。故结论成立。
\end{proof}

\begin{lemma}
\label{descent-lemma-descending-property-finite-locally-free}
性质 $\mathcal{P}(f) =$``$f$ 有限局部自由'' 在基底上为 fpqc 局部的。
设 $d \geq 0$。性质 $\mathcal{P}(f) =$``$f$ 为次数 $d$ 的有限局部自由
态射'' 在基底上为 fpqc 局部的。
\end{lemma}

\begin{proof}
有限局部自由等价于有限、平坦且局部有限表示（Morphisms, Lemma
\ref{morphisms-lemma-finite-flat}）。因此这由 Lemmas
\ref{descent-lemma-descending-property-finite}、
\ref{descent-lemma-descending-property-flat} 和
\ref{descent-lemma-descending-property-locally-finite-presentation} 推出。若
$f : Z \to U$ 有限局部自由，且 $\{U_i \to U\}$ 是一族满的态射，使得每个
拉回 $Z \times_U U_i \to U_i$ 的次数均为 $d$，则 $Z \to U$ 的次数为
$d$；例如，这是因为在点 $u \in U$ 处，可以从纤维
$(f_*\mathcal{O}_Z)_u \otimes_{\mathcal{O}_{U, u}} \kappa(u)$ 读出次数。
\end{proof}

\begin{lemma}
\label{descent-lemma-descending-property-monomorphism}
性质 $\mathcal{P}(f) =$``$f$ 为单态射'' 在基底上为 fpqc 局部的。
\end{lemma}

\begin{proof}
设 $f : X \to S$ 为概形态射。设 $\{S_i \to S\}$ 为 fpqc 覆盖，并假设
$f$ 的每个基变换 $f_i : X_i \to S_i$ 都是单态射。设
$a, b : T \to X$ 为两个满足 $f \circ a = f \circ b$ 的态射。我们必须证明
$a = b$。由于 $f_i$ 为单态射，可知 $a_i = b_i$，其中
$a_i, b_i : S_i \times_S T \to X_i$ 为相应的基变换。特别地，复合
$S_i \times_S T \to T \to X$ 相等。由于
$\coprod S_i \times_S T \to T$ 为上态射（例如见 Lemma
\ref{descent-lemma-fpqc-universal-effective-epimorphisms}），故 $a = b$。
\end{proof}

\begin{lemma}
\label{descent-lemma-descending-property-regular-immersion}
下列性质
\begin{enumerate}
\item[] $\mathcal{P}(f) =$``$f$ 为 Koszul-正则浸入''，
\item[] $\mathcal{P}(f) =$``$f$ 为 $H_1$-正则浸入''，以及
\item[] $\mathcal{P}(f) =$``$f$ 为拟正则浸入''
\end{enumerate}
在基底上为 fpqc 局部的。
\end{lemma}

\begin{proof}
我们使用 Lemma \ref{descent-lemma-descending-properties-morphisms} 的判据来证明。
由 Divisors, Definition \ref{divisors-definition-regular-immersion}，
Koszul-正则（相应地，$H_1$-正则、拟正则）浸入性在基底上为 Zariski 局部
的。由 Divisors, Lemma
\ref{divisors-lemma-flat-base-change-regular-immersion}，Koszul-正则
（相应地，$H_1$-正则、拟正则）浸入性在平坦基变换下保持。Lemma
\ref{descent-lemma-descending-properties-morphisms} 的最后一个假设 (3) 转化为如下
代数陈述：设 $A \to B$ 为忠实平坦环映射，设 $I \subset A$ 为理想。若
$IB$ 在 $\Spec(B)$ 上局部地由 $B$ 中的 Koszul-正则（相应地，
$H_1$-正则、拟正则）序列生成，则 $I \subset A$ 在 $\Spec(A)$ 上局部地
由 $A$ 中的 Koszul-正则（相应地，$H_1$-正则、拟正则）序列生成。这正是
More on Algebra, Lemma
\ref{more-algebra-lemma-flat-descent-regular-ideal}。
\end{proof}





\section{在目标上的 fppf 拓扑中局部的态射性质}
\label{descent-section-descending-properties-morphisms-fppf}

\noindent
本节给出概形态射的一些性质：我们（尚）不能证明它们在 fpqc 拓扑中于基底
上为局部的，但它们确实在 fppf 拓扑中于基底上为局部的。

\begin{lemma}
\label{descent-lemma-descending-fppf-property-immersion}
性质 $\mathcal{P}(f) =$``$f$ 为浸入'' 在基底上为 fppf 局部的。
\end{lemma}

\begin{proof}
浸入性在基变换下稳定；见 Schemes, Lemma
\ref{schemes-lemma-base-change-immersion}。浸入性在基底上为 Zariski 局部
的。最后，设 $\pi : S' \to S$ 为仿射概形之间平坦且局部有限表示的满态射。
注意，由 Morphisms, Lemma \ref{morphisms-lemma-fppf-open}，
$\pi : S' \to S$ 为开映射。设 $f : X \to S$ 为态射。假设基变换
$f' : X' \to S'$ 为浸入。特别地，$f'(X') = \pi^{-1}(f(X))$ 局部闭。
因此由 Topology, Lemma \ref{topology-lemma-open-morphism-quotient-topology}，
$f(X) \subset S$ 局部闭。设 $Z \subset S$ 为闭子集
$Z = \overline{f(X)} \setminus f(X)$。再次由 Topology, Lemma
\ref{topology-lemma-open-morphism-quotient-topology}，$f'(X')$ 在
$S' \setminus Z'$ 中为闭集。因此可将 Lemma
\ref{descent-lemma-descending-property-closed-immersion} 应用于 fpqc 覆盖
$\{S' \setminus Z' \to S \setminus Z\}$，从而得到
$f : X \to S \setminus Z$ 为闭浸入。换言之，$f$ 为浸入。因此应用
Lemma \ref{descent-lemma-descending-properties-morphisms} 即得结论。
\end{proof}

\begin{lemma}
\label{descent-lemma-descending-fppf-property-dominant}
性质 $\mathcal{P}(f) =$``$f$ 支配'' 在基底上为 fppf 局部的。
\end{lemma}

\begin{proof}
由 Morphisms, Lemma \ref{morphisms-lemma-open-base-change-dominant}，支配态射
在沿开态射拉回时保持，因而在沿局部有限表示的平坦态射拉回时保持
（Morphisms, Lemma \ref{morphisms-lemma-fppf-open}）。反方向较为容易。
事实上，支配性显然在基底上为 Zariski 局部的。再设 $S' \to S$ 为概形的
满态射（不必平坦或局部有限表示），并设 $f : X \to S$ 为态射。假设基变换
$f' : X' \to S'$ 支配。则 $X' \to S'\to S$ 是两个支配态射的复合，故而
支配。由于这也是复合 $X' \to X \to S$，可知 $X\to S$ 支配。
\end{proof}






\section{态射性质的 fpqc 下降的应用}
\label{descent-section-application-descending-properties-morphisms}

\noindent
下面的引理看似有些轻巧，却是研究 \'etale 态射和非分歧态射的有用工具。

\begin{lemma}
\label{descent-lemma-flat-surjective-quasi-compact-monomorphism-isomorphism}
设 $f : X \to Y$ 为平坦、拟紧且满的单态射。则 f 为同构。
\end{lemma}

\begin{proof}
由于 $f$ 是平坦、拟紧且满的态射，$\{X \to Y\}$ 是 $Y$ 的 fpqc 覆盖。
对角态射 $\Delta : X \to X \times_Y X$ 为同构（Schemes, Lemma
\ref{schemes-lemma-monomorphism}）。这蕴含 $f$ 沿 $f$ 的基变换为同构。
因此由 Lemma \ref{descent-lemma-descending-property-isomorphism}，$f$ 为同构。
\end{proof}

\noindent
可以用这个引理证明下面的重要结果；我们还会给出一个不使用 fpqc 下降的
证明。我们将在 \'Etale Morphisms, Section
\ref{etale-section-topological-etale} 中更详细地讨论这个结果及相关结果。

\begin{lemma}
\label{descent-lemma-universally-injective-etale-open-immersion}
普遍单射的 \'etale 态射为开浸入。
\end{lemma}

\begin{proof}[第一种证明]
设 $f : X \to Y$ 为普遍单射的 \'etale 态射。则 $f$ 为开映射
（Morphisms, Lemma \ref{morphisms-lemma-etale-open}），所以可以用 $f(X)$
替换 $Y$，并假设 $f$ 满。于是 $f$ 双射且开，因而为同胚。所以 $f$ 拟紧。
因此由 Lemma
\ref{descent-lemma-flat-surjective-quasi-compact-monomorphism-isomorphism}，只需证明
$f$ 为单态射。由于 $X \to Y$ 为 \'etale 态射，由 Morphisms, Lemma
\ref{morphisms-lemma-diagonal-unramified-morphism}（以及 Morphisms, Lemma
\ref{morphisms-lemma-flat-unramified-etale}），态射
$\Delta_{X/Y} : X \to X \times_Y X$ 为开浸入。由于 $f$ 普遍单射，
$\Delta_{X/Y}$ 还满；见 Morphisms, Lemma
\ref{morphisms-lemma-universally-injective}。因此 $\Delta_{X/Y}$ 为同构，
也就是说，$X \to Y$ 为单态射。
\end{proof}

\begin{proof}[第二种证明]
设 $f : X \to Y$ 为普遍单射的 \'etale 态射。则 $f$ 为开映射（Morphisms,
Lemma \ref{morphisms-lemma-etale-open}），所以可以用 $f(X)$ 替换 $Y$，并
假设 $f$ 满。由于这些假设在任意基变换后仍然成立，可知 $f$ 为普遍同胚。
因此 $f$ 为整态射；见 Morphisms, Lemma
\ref{morphisms-lemma-universal-homeomorphism}。进而由 Morphisms, Lemma
\ref{morphisms-lemma-finite-integral}，$f$ 为有限态射；再由 Morphisms,
Lemma \ref{morphisms-lemma-finite-flat}，$f$ 有限局部自由。为完成证明，只需
证明 $f$ 是次数 $1$ 的有限局部自由态射（次数 $1$ 的有限局部自由态射为
同构）。存在 $Y$ 到开闭子概形 $V_d$ 的分解，使得
$f^{-1}(V_d) \to V_d$ 是次数 $d$ 的有限局部自由态射；见 Morphisms,
Lemma \ref{morphisms-lemma-finite-locally-free}。若 $V_d$ 非空，则可以选取
态射 $\Spec(k) \to V_d \subset Y$，其中 $k$ 为代数闭域（只需取 $V_d$
某点的剩余域的代数闭包）。于是由 Morphisms, Lemma
\ref{morphisms-lemma-etale-over-field} 及 $k$ 代数闭这一事实，
$\Spec(k) \times_Y X \to \Spec(k)$ 是若干份 $\Spec(k)$ 的不交并。然而，
由于 $f$ 普遍单射，只能有一份，因此如所需有 $d = 1$。
\end{proof}

\noindent
利用下面对平坦普遍单射态射的刻画，可以稍加改写上一个引理中的假设。

\begin{lemma}
\label{descent-lemma-flat-universally-injective}
设 $f : X \to Y$ 为概形态射。以 $X^0$ 表示 $X$ 的不可约分支的泛点集合。
若
\begin{enumerate}
\item $f$ 平坦且分离；
\item 对 $\xi \in X^0$，有 $\kappa(f(\xi)) = \kappa(\xi)$；并且
\item 若 $\xi, \xi' \in X^0$ 且 $\xi \not = \xi'$，则
$f(\xi) \not = f(\xi')$，
\end{enumerate}
则 $f$ 普遍单射。
\end{lemma}

\begin{proof}
我们必须证明 $\Delta : X \to X \times_Y X$ 满；见 Morphisms, Lemma
\ref{morphisms-lemma-universally-injective}。由于 $X \to Y$ 分离，$\Delta$
的像为闭集。所以若 $\Delta$ 不满，就可以找到 $X \times_S X$ 的某个不可约
分支的泛点
% P05-ERR-0037
$\eta \in X \times_S X$，它不在 $\Delta$ 的像中。投影
$\text{pr}_1 : X \times_Y X \to X$ 是平坦态射 $X \to Y$ 的基变换，因而
平坦；见 Morphisms, Lemma \ref{morphisms-lemma-base-change-flat}。因此沿
$\text{pr}_1$ 可以提升泛化；见 Morphisms, Lemma
\ref{morphisms-lemma-generalizations-lift-flat}。可知
$\xi = \text{pr}_1(\eta) \in X^0$。然而，假设 (2) 和 (3) 保证，对每个
$\xi \in X^0$，概形 $(X \times_Y X)_{f(\xi)}$ 至多有一个点。换言之，
$\Delta(\xi) = \eta$，矛盾。
\end{proof}

\noindent
因此可将 Lemma \ref{descent-lemma-universally-injective-etale-open-immersion} 改述如下。

\begin{lemma}
\label{descent-lemma-characterize-open-immersion}
设 $f : X \to Y$ 为概形态射。以 $X^0$ 表示 $X$ 的不可约分支的泛点集合。
若
\begin{enumerate}
\item $f$ 为 \'etale 态射且分离；
\item 对 $\xi \in X^0$，有 $\kappa(f(\xi)) = \kappa(\xi)$；并且
\item 若 $\xi, \xi' \in X^0$ 且 $\xi \not = \xi'$，则
$f(\xi) \not = f(\xi')$，
\end{enumerate}
则 $f$ 为开浸入。
\end{lemma}

\begin{proof}
立即由 Lemmas \ref{descent-lemma-flat-universally-injective} 和
\ref{descent-lemma-universally-injective-etale-open-immersion} 推出。
\end{proof}

\begin{lemma}
\label{descent-lemma-descending-property-proper-over-base}
设 $f : X \to Y$ 为局部有限型的概形态射，设 $Z$ 为 $X$ 的闭子集。若存在
fpqc 覆盖 $\{Y_i \to Y\}$，使得逆像 $Z_i \subset Y_i \times_Y X$ 在
$Y_i$ 上固有（Cohomology of Schemes, Definition
\ref{coherent-definition-proper-over-base}），则 $Z$ 在 $Y$ 上固有。
\end{lemma}

\begin{proof}
赋予 $Z$ 诱导的约化闭子概形结构；见 Schemes, Definition
\ref{schemes-definition-reduced-induced-scheme}。对每个 $i$，基变换
$Y_i \times_Y Z$ 是 $Y_i \times_Y X$ 的闭子概形，其底层闭子集为 $Z_i$。
由定义（借助 Cohomology of Schemes, Lemma
\ref{coherent-lemma-closed-proper-over-base}），投影
$Y_i \times_Y Z \to Y_i$ 为固有态射。因此由 Lemma
\ref{descent-lemma-descending-property-proper}，$Z \to Y$ 为固有态射。故由定义，
$Z$ 在 $Y$ 上固有。
\end{proof}

\begin{lemma}
\label{descent-lemma-descending-property-ample}
设 $f : X \to S$ 为概形态射，设 $\mathcal{L}$ 为可逆
$\mathcal{O}_X$-模，设 $\{g_i : S_i \to S\}_{i \in I}$ 为 fpqc 覆盖。
以 $f_i : X_i \to S_i$ 表示 $f$ 的基变换，并以 $\mathcal{L}_i$ 表示
$\mathcal{L}$ 到 $X_i$ 的拉回。下列条件等价：
\begin{enumerate}
\item $\mathcal{L}$ 在 $X/S$ 上丰沛；
\item 对每个 $i \in I$，$\mathcal{L}_i$ 在 $X_i/S_i$ 上丰沛。
\end{enumerate}
\end{lemma}

\begin{proof}
蕴含 (1) $\Rightarrow$ (2) 由 Morphisms, Lemma
\ref{morphisms-lemma-ample-base-change} 给出。假设对每个 $i \in I$，
$\mathcal{L}_i$ 在 $X_i/S_i$ 上丰沛。由 Morphisms, Definition
\ref{morphisms-definition-relatively-ample}，这蕴含 $X_i \to S_i$ 拟紧；
再由 Morphisms, Lemma \ref{morphisms-lemma-relatively-ample-separated}，
% P05-ERR-0038
这蕴含 $X_i \to S$ 分离。因此由 Lemmas
\ref{descent-lemma-descending-property-quasi-compact} 和
\ref{descent-lemma-descending-property-separated}，$f$ 拟紧且分离。

\medskip\noindent
这意味着
$\mathcal{A} = \bigoplus_{d \geq 0} f_*\mathcal{L}^{\otimes d}$
是拟凝聚分次 $\mathcal{O}_S$-代数（Schemes, Lemma
\ref{schemes-lemma-push-forward-quasi-coherent}）。此外，由 Cohomology of
Schemes, Lemma \ref{coherent-lemma-flat-base-change-cohomology}，
$\mathcal{A}$ 的构造与平坦基变换交换。特别地，若置
$\mathcal{A}_i = \bigoplus_{d \geq 0} f_{i, *}\mathcal{L}_i^{\otimes d}$
，则有 $\mathcal{A}_i = g_i^*\mathcal{A}$。由此，自然映射
$\psi_d : f^*\mathcal{A}_d \to \mathcal{L}^{\otimes d}$
（作为 $\mathcal{O}_X$-模的映射）拉回为自然映射
$\psi_{i, d} : f_i^*(\mathcal{A}_i)_d \to \mathcal{L}_i^{\otimes d}$
（作为 $\mathcal{O}_{X_i}$-模的映射）。由于 $\mathcal{L}_i$ 在 $X_i/S_i$
上丰沛，可知对任意点 $x_i \in X_i$，存在 $d \geq 1$，使得
$f_i^*(\mathcal{A}_i)_d \to \mathcal{L}_i^{\otimes d}$ 在 $x_i$ 处的茎上
为满射。这或者直接由相对丰沛模的定义推出，或者由 Morphisms, Lemma
\ref{morphisms-lemma-characterize-relatively-ample} 推出。若 $x \in X$，
则可以选取 $i$ 和映到 $x$ 的 $x_i \in X_i$。由于
$\mathcal{O}_{X, x} \to \mathcal{O}_{X_i, x_i}$ 平坦，因而忠实平坦，
可知对每个 $x \in X$，存在 $d \geq 1$，使得
$f^*\mathcal{A}_d \to \mathcal{L}^{\otimes d}$
在 $x$ 处的茎上为满射。这蕴含 Constructions, Lemma
\ref{constructions-lemma-invertible-map-into-relative-proj}
中对应于分次 $\mathcal{O}_X$-代数映射
$\psi : f^*\mathcal{A} \to \bigoplus_{d \geq 0} \mathcal{L}^{\otimes d}$
的开子集 $U(\psi) \subset X$ 等于 $X$。考虑对应的态射
$$
r_{\mathcal{L}, \psi} : X \longrightarrow \underline{\text{Proj}}_S(\mathcal{A})
$$
。由上述内容显然可知，$r_{\mathcal{L}, \psi}$ 到 $S_i$ 的基变换是态射
$r_{\mathcal{L}_i, \psi_i}$；由 Morphisms, Lemma
\ref{morphisms-lemma-characterize-relatively-ample}，后者为开浸入。因此由
Lemma \ref{descent-lemma-descending-property-open-immersion}，
$r_{\mathcal{L}, \psi}$ 为开浸入；再由 Morphisms, Lemma
\ref{morphisms-lemma-characterize-relatively-ample}，可知 $\mathcal{L}$ 在
$X/S$ 上丰沛。
\end{proof}















\section{在源上局部的态射性质}
\label{descent-section-properties-morphisms-local-source}

\noindent
经常会出现这样的情形：在与另一个态射预复合之后，可以证明某个态射具有
某一性质。在许多情况下，这也蕴含原态射具有该性质。我们用下面的定义将其
形式化。

\begin{definition}
\label{descent-definition-property-morphisms-local-source}
设 $\mathcal{P}$ 为概形态射的一个性质。设
$\tau \in \{Zariski, \linebreak[0] fpqc, \linebreak[0] fppf, \linebreak[0]
\etale, \linebreak[0] smooth, \linebreak[0] syntomic\}$。若对 $S$ 上任意概形
态射 $f : X \to Y$ 及任意 $\tau$-覆盖 $\{X_i \to X\}_{i \in I}$，都有
$$
f \text{ 具有 }\mathcal{P}
\Leftrightarrow
\text{每个 }X_i \to Y\text{ 具有 }\mathcal{P}.
$$
，则称 $\mathcal{P}$ 是{\it 在源上 $\tau$ 局部的}，或称它{\it 对
$\tau$-拓扑在源上局部}。
\end{definition}

\noindent
需要明确的是，由于同构总是覆盖，我们可见（或者要求）性质
$\mathcal{P}$ 对 $X \to Y$ 成立，当且仅当它对任意与 $X \to Y$ 同构的
箭头 $X' \to Y'$ 成立。若一个性质在源上为 $\tau$-局部的，则与出现在
$\tau$-覆盖中的态射预复合会保持该性质。下面是形式陈述。

\begin{lemma}
\label{descent-lemma-precompose-property-local-source}
设 $\tau \in \{fpqc, fppf, syntomic, smooth, \etale, Zariski\}$。设
$\mathcal{P}$ 为在源上 $\tau$-局部的态射性质。设 $f : X \to Y$ 具有性质
$\mathcal{P}$。对任意态射 $a : X' \to X$，若它分别为平坦态射、平坦且
局部有限表示的态射、syntomic 态射、
% P05-ERR-0039
\'etale 态射、开浸入，则复合 $f \circ a : X' \to Y$ 具有性质
$\mathcal{P}$。
\end{lemma}

\begin{proof}
这是因为可以把 $X' \to X$ 放入构成 $\tau$-覆盖的一族态射中。
\end{proof}

\begin{lemma}
\label{descent-lemma-largest-open-of-the-source}
设 $\tau \in \{fppf, syntomic, smooth, \etale\}$。设 $\mathcal{P}$ 为在源上
$\tau$-局部的态射性质。对任意概形态射 $f : X \to Y$，存在最大的开集
$W(f) \subset X$，使限制 $f|_{W(f)} : W(f) \to Y$ 具有 $\mathcal{P}$。
此外，若 $g : X' \to X$ 平坦且局部有限表示、为 syntomic 态射、光滑或为
\'etale 态射，并且 $f' = f \circ g : X' \to Y$，则
$g^{-1}(W(f)) = W(f')$。
\end{lemma}

\begin{proof}
考虑所有具有下列性质的态射 $g : X' \to X$ 的像 $g(X') \subset X$ 之并
$W$：
\begin{enumerate}
\item $g$ 平坦且局部有限表示、为 syntomic 态射、光滑或为 \'etale 态射；
\item 复合 $X' \to X \to Y$ 具有性质 $\mathcal{P}$。
\end{enumerate}
由于这样的态射 $g$ 为开映射（见 Morphisms, Lemma
\ref{morphisms-lemma-fppf-open}），$W \subset X$ 是 $X$ 的开子集。由于
$\mathcal{P}$ 在 $\tau$ 拓扑中为局部的，限制 $f|_W : W \to Y$ 具有性质
$\mathcal{P}$：事实上，我们得到 $W$ 的一个 $\tau$ 覆盖
$\{X' \to W\}$，其各拉回均具有 $\mathcal{P}$。这证明了 $W(f)$ 的存在性。
最后一句所述的相容性立即由 $W(f)$ 的构造推出。
\end{proof}

\begin{lemma}
\label{descent-lemma-properties-morphisms-local-source}
设 $\mathcal{P}$ 为概形态射的一个性质。设
$\tau \in \{fpqc, \linebreak[0] fppf, \linebreak[0]
\etale, \linebreak[0] smooth, \linebreak[0] syntomic\}$。假设：
\begin{enumerate}
\item 依照 $\tau$ 分别为 fpqc、fppf、\'etale、smooth 或 syntomic，该性质
在与平坦态射、平坦且局部有限表示的态射、\'etale 态射、光滑态射或
syntomic 态射预复合时保持；
\item 该性质在源上为 Zariski 局部的；
\item 该性质在目标上为 Zariski 局部的；
\item 对任意仿射概形态射 $f : X \to Y$ 及仿射概形的任意满态射
$X' \to X$，若依照 $\tau$ 分别为 fpqc、fppf、\'etale、smooth 或
syntomic，该满态射相应地为平坦、平坦有限表示、\'etale、光滑或
syntomic 态射，并且性质 $\mathcal{P}$ 对复合 $f' : X' \to Y$ 成立，
则性质 $\mathcal{P}$ 对 $f$ 成立。
\end{enumerate}
则 $\mathcal{P}$ 在源上为 $\tau$ 局部的。
\end{lemma}

\begin{proof}
这几乎立即由 $\tau$-覆盖的定义推出；见 Topologies, Definition
\ref{topologies-definition-fpqc-covering}
\ref{topologies-definition-fppf-covering}
\ref{topologies-definition-etale-covering}
\ref{topologies-definition-smooth-covering}, or
\ref{topologies-definition-syntomic-covering}
以及 Topologies, Lemma
\ref{topologies-lemma-fpqc-affine},
\ref{topologies-lemma-fppf-affine},
\ref{topologies-lemma-etale-affine},
\ref{topologies-lemma-smooth-affine}, or
\ref{topologies-lemma-syntomic-affine}.
省略细节。（提示：使用在源和目标上的局部性，把性质 $\mathcal{P}$ 的验证
归约到仿射概形之间的态射。然后应用 (1) 和 (4)。）
\end{proof}

\begin{remark}
\label{descent-remark-properties-morphisms-local-source-standard}
（这是对上面 Remarks \ref{descent-remark-descending-properties-standard} 和
\ref{descent-remark-descending-properties-morphisms-standard} 的重复。）在上述 Lemma
\ref{descent-lemma-properties-morphisms-local-source} 中，若 $\tau = smooth$，则在
条件 (4) 中可以假设该态射为（满的）标准光滑态射。$\tau = syntomic$ 或
$\tau = \etale$ 时也类似。
\end{remark}



\section{在源上的 fpqc 拓扑中局部的态射性质}
\label{descent-section-fpqc-local-source}

\noindent
下面给出一些在源上为 fpqc 局部的态射性质。

\begin{lemma}
\label{descent-lemma-flat-fpqc-local-source}
性质 $\mathcal{P}(f)=$``$f$ 平坦'' 在源上为 fpqc 局部的。
\end{lemma}

\begin{proof}
由于平坦性以局部环的映射来定义（Morphisms, Definition
\ref{morphisms-definition-flat}），必须证明如下代数事实：设
$A \to B \to C$ 为局部环的局部同态，并假设 $B \to C$ 平坦。则
$A \to B$ 平坦，当且仅当 $A \to C$ 平坦。若 $A \to B$ 平坦，则由
Algebra, Lemma \ref{algebra-lemma-composition-flat}，$A \to C$ 平坦。
反之，假设 $A \to C$ 平坦。注意 $B \to C$ 忠实平坦；见 Algebra,
Lemma \ref{algebra-lemma-local-flat-ff}。因此由 Algebra, Lemma
\ref{algebra-lemma-flat-permanence}，$A \to B$ 平坦。（另见 Morphisms,
Lemma \ref{morphisms-lemma-flat-permanence} 的直接证明。）
\end{proof}

\begin{lemma}
\label{descent-lemma-injective-local-rings-fpqc-local-source}
性质 $\mathcal{P}(f : X \to Y)=$``对每个 $x \in X$，局部环映射
$\mathcal{O}_{Y, f(x)} \to \mathcal{O}_{X, x}$ 单射'' 在源上为 fpqc 局部的。
\end{lemma}

\begin{proof}
省略。这只是一次（或许不太恰当的）幽默尝试。
\end{proof}





\section{在源上的 fppf 拓扑中局部的态射性质}
\label{descent-section-fppf-local-source}

\noindent
下面给出一些在源上为 fppf 局部的态射性质。

\begin{lemma}
\label{descent-lemma-locally-finite-presentation-fppf-local-source}
性质 $\mathcal{P}(f)=$``$f$ 局部有限表示'' 在源上为 fppf 局部的。
\end{lemma}

\begin{proof}
局部有限表示性在源和目标上都为 Zariski 局部的；见 Morphisms, Lemma
\ref{morphisms-lemma-locally-finite-presentation-characterize}。该性质在复合
下保持；见 Morphisms, Lemma
\ref{morphisms-lemma-composition-finite-presentation}。这证明了 Lemma
\ref{descent-lemma-properties-morphisms-local-source} 的 (1)、(2) 和 (3)。最后的
条件 (4) 是 Lemma \ref{descent-lemma-flat-finitely-presented-permanence-algebra}。
故结论成立。
\end{proof}

\begin{lemma}
\label{descent-lemma-locally-finite-type-fppf-local-source}
性质 $\mathcal{P}(f)=$``$f$ 局部有限型'' 在源上为 fppf 局部的。
\end{lemma}

\begin{proof}
局部有限型性在源和目标上都为 Zariski 局部的；见 Morphisms, Lemma
\ref{morphisms-lemma-locally-finite-type-characterize}。该性质在复合下保持；
见 Morphisms, Lemma \ref{morphisms-lemma-composition-finite-type}；而局部有限
表示的平坦态射局部有限型；见 Morphisms, Lemma
\ref{morphisms-lemma-finite-presentation-finite-type}。这证明了 Lemma
\ref{descent-lemma-properties-morphisms-local-source} 的 (1)、(2) 和 (3)。最后的
条件 (4) 是 Lemma \ref{descent-lemma-finite-type-local-source-fppf-algebra}。
故结论成立。
\end{proof}

\begin{lemma}
\label{descent-lemma-open-fppf-local-source}
性质 $\mathcal{P}(f)=$``$f$ 为开映射'' 在源上为 fppf 局部的。
\end{lemma}

\begin{proof}
开态射性显然在源和目标上都为 Zariski 局部的。该性质在复合下保持；见
Morphisms, Lemma \ref{morphisms-lemma-composition-open}；而有限表示的平坦
态射为开映射；见 Morphisms, Lemma \ref{morphisms-lemma-fppf-open}。这证明
了 Lemma \ref{descent-lemma-properties-morphisms-local-source} 的 (1)、(2) 和 (3)。
最后的条件 (4) 由 Morphisms, Lemma
\ref{morphisms-lemma-fpqc-quotient-topology} 推出。故结论成立。
\end{proof}

\begin{lemma}
\label{descent-lemma-universally-open-fppf-local-source}
性质 $\mathcal{P}(f)=$``$f$ 普遍开'' 在源上为 fppf 局部的。
\end{lemma}

\begin{proof}
设 $f : X \to Y$ 为概形态射，设 $\{X_i \to X\}_{i \in I}$ 为 fppf 覆盖。
% P05-ERR-0040
以 $f_i : X_i \to X$ 表示这些复合。我们必须证明 $f$ 普遍开，当且仅当每个
$f_i$ 都普遍开。若 $f$ 普遍开，则每个 $f_i$ 也普遍开，因为映射
$X_i \to X$ 普遍开，且普遍开态射的复合普遍开（Morphisms, Lemmas
\ref{morphisms-lemma-fppf-open} 和
\ref{morphisms-lemma-composition-open}）。反之，假设每个 $f_i$ 都普遍开。
设 $Y' \to Y$ 为概形态射，记 $X' = Y' \times_Y X$ 及
$X'_i = Y' \times_Y X_i$。注意，$\{X_i' \to X'\}_{i \in I}$ 也是 fppf
覆盖。由假设，态射 $f'_i : X_i' \to Y'$ 为开映射。因此由上面的 Lemma
\ref{descent-lemma-open-fppf-local-source}，$f' : X' \to Y'$ 为开映射，正如所需。
\end{proof}



\section{在源上的 syntomic 拓扑中局部的态射性质}
\label{descent-section-syntomic-local-source}

\noindent
下面给出一些在源上为 syntomic 局部的态射性质。

\begin{lemma}
\label{descent-lemma-syntomic-syntomic-local-source}
性质 $\mathcal{P}(f)=$``$f$ 为 syntomic 态射'' 在源上为 syntomic 局部的。
\end{lemma}

\begin{proof}
合并 Lemma \ref{descent-lemma-properties-morphisms-local-source}、Morphisms, Lemma
\ref{morphisms-lemma-syntomic-characterize}（在源和目标上为 Zariski 局部
的）、Morphisms, Lemma \ref{morphisms-lemma-composition-syntomic}
（预复合）以及 Lemma \ref{descent-lemma-syntomic-smooth-etale-permanence}
（第 (4) 部分）即得。
\end{proof}




\section{在源上的光滑拓扑中局部的态射性质}
\label{descent-section-smooth-local-source}

\noindent
下面给出一些在源上为光滑局部的态射性质。还请注意经由平坦态射下降光滑性
的结果 Lemma \ref{descent-lemma-smooth-permanence}；它在某些方面更强。

\begin{lemma}
\label{descent-lemma-smooth-smooth-local-source}
性质 $\mathcal{P}(f)=$``$f$ 光滑'' 在源上为光滑局部的。
\end{lemma}

\begin{proof}
合并 Lemma \ref{descent-lemma-properties-morphisms-local-source}、Morphisms, Lemma
\ref{morphisms-lemma-smooth-characterize}（在源和目标上为 Zariski 局部
的）、Morphisms, Lemma \ref{morphisms-lemma-composition-smooth}（预复合）
以及 Lemma \ref{descent-lemma-syntomic-smooth-etale-permanence}（第 (4) 部分）即得。
\end{proof}



\section{在源上的 \'etale 拓扑中局部的态射性质}
\label{descent-section-etale-local-source}

\noindent
下面给出一些在源上为 \'etale 局部的态射性质。

\begin{lemma}
\label{descent-lemma-etale-etale-local-source}
性质 $\mathcal{P}(f)=$``$f$ 为 \'etale 态射'' 在源上为 \'etale 局部的。
\end{lemma}

\begin{proof}
合并 Lemma \ref{descent-lemma-properties-morphisms-local-source}、Morphisms, Lemma
\ref{morphisms-lemma-etale-characterize}（在源和目标上为 Zariski 局部的）、
Morphisms, Lemma \ref{morphisms-lemma-composition-etale}（预复合）以及
Lemma \ref{descent-lemma-syntomic-smooth-etale-permanence}（第 (4) 部分）即得。
\end{proof}

\begin{lemma}
\label{descent-lemma-locally-quasi-finite-etale-local-source}
性质 $\mathcal{P}(f)=$``$f$ 局部拟有限'' 在源上为 \'etale 局部的。
\end{lemma}

\begin{proof}
我们使用 Lemma \ref{descent-lemma-properties-morphisms-local-source}。由 Morphisms,
Lemma \ref{morphisms-lemma-locally-quasi-finite-characterize}，局部拟有限性
在源和目标上为 Zariski 局部的。由 Morphisms, Lemmas
\ref{morphisms-lemma-composition-quasi-finite} 和
\ref{morphisms-lemma-etale-locally-quasi-finite}，局部拟有限态射与 \'etale
态射预复合后仍局部拟有限。最后，假设 $X \to Y$ 为仿射概形态射，且
$X' \to X$ 为仿射概形的满 \'etale 态射，并且 $X' \to Y$ 局部拟有限。
则 $X' \to Y$ 有限型；由 Lemma
\ref{descent-lemma-finite-type-local-source-fppf-algebra}，$X \to Y$ 也有限型。此外，
由假设 $X' \to Y$ 的纤维有限，因而 $X \to Y$ 的纤维也有限。由
Morphisms, Lemma \ref{morphisms-lemma-quasi-finite}，可知 $X \to Y$
拟有限。这证明了 Lemma \ref{descent-lemma-properties-morphisms-local-source} 的
最后一个假设，从而完成证明。
\end{proof}

\begin{lemma}
\label{descent-lemma-unramified-etale-local-source}
性质 $\mathcal{P}(f)=$``$f$ 非分歧'' 在源上为 \'etale 局部的。
性质 $\mathcal{P}(f)=$``$f$ 为 G-非分歧态射'' 在源上为 \'etale 局部的。
\end{lemma}

\begin{proof}
我们使用 Lemma \ref{descent-lemma-properties-morphisms-local-source}。由 Morphisms,
Lemma \ref{morphisms-lemma-unramified-characterize}，非分歧性（相应地，
G-非分歧性）在源和目标上为 Zariski 局部的。由 Morphisms, Lemmas
\ref{morphisms-lemma-composition-unramified} 和
\ref{morphisms-lemma-etale-smooth-unramified}，非分歧（相应地，G-非分歧）
态射与 \'etale 态射预复合后仍非分歧（相应地，G-非分歧）。最后，假设
$X \to Y$ 为仿射概形态射，且 $f : X' \to X$ 为仿射概形的满 \'etale
态射，并且 $X' \to Y$ 非分歧（相应地，G-非分歧）。则 $X' \to Y$
有限型（相应地，有限表示）；由 Lemma
\ref{descent-lemma-finite-type-local-source-fppf-algebra}（相应地，Lemma
\ref{descent-lemma-flat-finitely-presented-permanence-algebra}），$X \to Y$ 也有限型
（相应地，有限表示）。由 Morphisms, Lemma
\ref{morphisms-lemma-triangle-differentials-smooth}，有短正合列
$$
0 \to f^*\Omega_{X/Y} \to \Omega_{X'/Y} \to \Omega_{X'/X} \to 0.
$$
。由于 $X' \to Y$ 非分歧，中间项为零。又因 $f$ 忠实平坦，可知
$\Omega_{X/Y} = 0$。因此 $X \to Y$ 非分歧（相应地，G-非分歧）；见
Morphisms, Lemma \ref{morphisms-lemma-unramified-omega-zero}。这证明了
Lemma \ref{descent-lemma-properties-morphisms-local-source} 的最后一个假设，从而
完成证明。
\end{proof}




\section{在源与目标上 \'etale 局部的态射性质}
\label{descent-section-properties-etale-local-source-target}

\noindent
设 $\mathcal{P}$ 为概形态射的一个性质。``$\mathcal{P}$ 在源和目标上为
\'etale 局部的''这句话有一种直观含义。然而，这个概念并不等同于要求
$\mathcal{P}$ 既在源上为 \'etale 局部、又在目标上为 \'etale 局部。在
进一步讨论之前，先给出两个简单的例子。

\begin{example}
\label{descent-example-silly-one}
考虑按规则 $\mathcal{P}(X \to Y) = $``$Y$ 局部 Noether'' 定义的概形态射
性质 $\mathcal{P}$。读者可以验证，它在源上为 \'etale 局部的，也在目标上
为 \'etale 局部的（省略；见 Lemma \ref{descent-lemma-Noetherian-local-fppf}）。
但是，下列陈述{\bf 不}成立：若 $f : X \to Y$ 具有 $\mathcal{P}$，并且
$g : Y \to Z$ 为 \'etale 态射，则 $g \circ f$ 具有 $\mathcal{P}$。事实上，
$f$ 可以是 $Y$ 上的恒等态射，而 $g$ 可以是局部 Noether 概形 $Y$ 到非局部
Noether 概形 $Z$ 的开浸入。
\end{example}

\noindent
下面的例子在某种意义上更糟。

\begin{example}
\label{descent-example-silly-two}
考虑按下列规则定义的概形态射性质 $\mathcal{P}$：
$\mathcal{P}(f : X \to Y) = $``对每个 $y \in Y$，若它是某个 $f(x)$ 的特化，
其中 $x \in X$，则局部环 $\mathcal{O}_{Y, y}$ 为 Noether 环''。我们来验证
它在源上和目标上均为 \'etale 局部的。下面将自由使用 Schemes, Lemma
\ref{schemes-lemma-specialize-points}。

\medskip\noindent
在目标上局部：设 $\{g_i : Y_i \to Y\}$ 为 \'etale 覆盖。以
$f_i : X_i \to Y_i$ 表示 $f$ 的基变换，并以 $h_i : X_i \to X$ 表示投影。
假设 $\mathcal{P}(f)$。设
% P05-ERR-0041
$f(x_i) \leadsto y_i$ 为一个特化。则 $f(h_i(x_i)) \leadsto g_i(y_i)$，所以
$\mathcal{P}(f)$ 蕴含 $\mathcal{O}_{Y, g_i(y_i)}$ 为 Noether 环。此外，
$\mathcal{O}_{Y, g_i(y_i)} \to \mathcal{O}_{Y_i, y_i}$ 是某个 \'etale 环映射
的局部化。因此由 Algebra, Lemma \ref{algebra-lemma-Noetherian-permanence}，
$\mathcal{O}_{Y_i, y_i}$ 为 Noether 环。反之，假设对所有 $i$，
$\mathcal{P}(f_i)$ 均成立。设 $f(x) \leadsto y$ 为一个特化。选取 $i$ 和映到
$y$ 的 $y_i \in Y_i$。由于 $x$ 可视为
$\Spec(\mathcal{O}_{Y, y}) \times_Y X$ 的一个点，并且
$\mathcal{O}_{Y, y} \to \mathcal{O}_{Y_i, y_i}$ 忠实平坦，存在映到 $x$ 的点
$x_i \in \Spec(\mathcal{O}_{Y_i, y_i}) \times_Y X$。于是 $x_i \in X_i$，
并且 $f_i(x_i)$ 特化到 $y_i$。由 $\mathcal{P}(f_i)$ 可知
$\mathcal{O}_{Y_i, y_i}$ 为 Noether 环；再由 Algebra, Lemma
\ref{algebra-lemma-descent-Noetherian}，这蕴含 $\mathcal{O}_{Y, y}$ 为
Noether 环。

\medskip\noindent
在源上局部：设 $\{h_i : X_i \to X\}$ 为 \'etale 覆盖。以
$f_i : X_i \to Y$ 表示复合 $f \circ h_i$。假设 $\mathcal{P}(f)$。设
% P05-ERR-0042
$f(x_i) \leadsto y$ 为一个特化。则 $f(h_i(x_i)) \leadsto y$，所以
$\mathcal{P}(f)$ 蕴含 $\mathcal{O}_{Y, y}$ 为 Noether 环。因此
$\mathcal{P}(f_i)$ 成立。反之，假设对所有 $i$，$\mathcal{P}(f_i)$ 均成立。
设 $f(x) \leadsto y$ 为一个特化。选取 $i$ 和映到 $x$ 的 $x_i \in X_i$。
则 $y$ 是 $f_i(x_i) = f(x)$ 的特化。因此 $\mathcal{P}(f_i)$ 蕴含
$\mathcal{O}_{Y, y}$ 为 Noether 环，正如所需。

\medskip\noindent
我们断言，存在交换图
$$
\xymatrix{
U \ar[d]_a \ar[r]_h & V \ar[d]^b \\
X \ar[r]^f & Y
}
$$
，其中竖直箭头是满的 \'etale 态射，$h$ 具有 $\mathcal{P}$，而 $f$ 不具有
$\mathcal{P}$。事实上，令
$$
Y =
\Spec\Big(
\mathbf{C}[x_n; n \in \mathbf{Z}]/(x_n x_m; n \not = m)
\Big)
$$
，并设 $X \subset Y$ 为如下开子概形：它是所有坐标都满足 $x_n = 0$ 的点
的补集。令 $U = X$，令 $V = X \amalg Y$，令 $a, b$ 为显然的映射，并令
$h : U \to V$ 为 $U = X$ 到 $V$ 的第一个直和项的包含。上述断言成立，
因为 $U$ 局部 Noether，而 $Y$ 不是。
\end{example}

\noindent
那么，在源与目标上为 \'etale 局部的性质，其正确概念应是什么？我们认为，
类比 Morphisms, Definition \ref{morphisms-definition-property-local}，应取
下面的定义。

\begin{definition}
\label{descent-definition-local-source-target}
设 $\mathcal{P}$ 为概形态射的一个性质。若满足下列条件，则称
$\mathcal{P}$ {\it 在源与目标上为 \'etale 局部的}：
\begin{enumerate}
\item （在与 \'etale 映射预复合时稳定）若 $f : X \to Y$ 为 \'etale 态射，
且 $g : Y \to Z$ 具有 $\mathcal{P}$，则 $g \circ f$ 具有 $\mathcal{P}$；
\item （在 \'etale 基变换下稳定）若 $f : X \to Y$ 具有 $\mathcal{P}$，且
$Y' \to Y$ 为 \'etale 态射，则基变换 $f' : Y' \times_Y X \to Y'$ 具有
$\mathcal{P}$；
\item （局部性）给定态射 $f : X \to Y$，下列条件等价：
\begin{enumerate}
\item $f$ 具有 $\mathcal{P}$；
\item 对每个 $x \in X$，存在交换图
$$
\xymatrix{
U \ar[d]_a \ar[r]_h & V \ar[d]^b \\
X \ar[r]^f & Y
}
$$
，其中竖直箭头为 \'etale 态射，且存在满足 $a(u) = x$ 的 $u \in U$，使
$h$ 具有 $\mathcal{P}$。
\end{enumerate}
\end{enumerate}
\end{definition}

\noindent
事实证明，这个定义排除了 Examples \ref{descent-example-silly-one} 和
\ref{descent-example-silly-two} 中所见的行为。我们将在 Remark
\ref{descent-remark-compare-definitions} 中把它与 Deligne 和 Mumford 的论文
\cite{DM} 中的定义作比较。此外，一个在源与目标上为 \'etale 局部的性质，
在源上和目标上分别都是 \'etale 局部的。最后，正如 Lemma
\ref{descent-lemma-etale-local-source-target} 将表明的，其逆命题几乎成立。

\begin{lemma}
\label{descent-lemma-local-source-target-implies}
设 $\mathcal{P}$ 为在源与目标上 \'etale 局部的概形态射性质。则
\begin{enumerate}
\item $\mathcal{P}$ 在源上为 \'etale 局部的；
\item $\mathcal{P}$ 在目标上为 \'etale 局部的；
\item $\mathcal{P}$ 在与 \'etale 态射后复合时稳定：若 $f : X \to Y$
具有 $\mathcal{P}$，且 $g : Y \to Z$ 为 \'etale 态射，则 $g \circ f$
具有 $\mathcal{P}$；
\item $\mathcal{P}$ 具有保持性：给定 $f : X \to Y$ 及 \'etale 态射
$g : Y \to Z$，若 $g \circ f$ 具有 $\mathcal{P}$，则 $f$ 具有
$\mathcal{P}$。
\end{enumerate}
\end{lemma}

\begin{proof}
我们完整写出所有细节。

\medskip\noindent
(1) 的证明。设 $f : X \to Y$ 为概形态射，设
$\{X_i \to X\}_{i \in I}$ 为 $X$ 的 \'etale 覆盖。若每个复合
$h_i : X_i \to Y$ 都具有 $\mathcal{P}$，则对每个 $x \in X$，可以找到
$i \in I$ 及映到 $x$ 的点 $x_i \in X_i$。于是
$(X_i, x_i) \to (X, x)$ 是芽的 \'etale 态射，
$\text{id}_Y : Y \to Y$ 为 \'etale 态射，而 $h_i$ 正如 Definition
\ref{descent-definition-local-source-target} 第 (3) 部分所述。因此 $f$ 具有
$\mathcal{P}$。反之，若 $f$ 具有 $\mathcal{P}$，则由 Definition
\ref{descent-definition-local-source-target} 第 (1) 部分，每个 $X_i \to Y$ 都具有
$\mathcal{P}$。

\medskip\noindent
(2) 的证明。设 $f : X \to Y$ 为概形态射，设
$\{Y_i \to Y\}_{i \in I}$ 为 $Y$ 的 \'etale 覆盖。写
$X_i = Y_i \times_Y X$，并以 $h_i : X_i \to Y_i$ 表示 $f$ 的基变换。若
每个 $h_i : X_i \to Y_i$ 都具有 $\mathcal{P}$，则对每个 $x \in X$，选取
$i \in I$ 及映到 $x$ 的点 $x_i \in X_i$。于是
$(X_i, x_i) \to (X, x)$ 是芽的 \'etale 态射，$Y_i \to Y$ 为 \'etale
态射，而 $h_i$ 正如 Definition \ref{descent-definition-local-source-target}
第 (3) 部分所述。因此 $f$ 具有 $\mathcal{P}$。反之，若 $f$ 具有
$\mathcal{P}$，则由 Definition \ref{descent-definition-local-source-target}
第 (2) 部分，每个 $X_i \to Y_i$ 都具有 $\mathcal{P}$。

\medskip\noindent
(3) 的证明。假设 $f : X \to Y$ 具有 $\mathcal{P}$，且 $g : Y \to Z$ 为
\'etale 态射。对每个 $x \in X$，可将 $(X, x) \to (X, x)$ 视为芽的
\'etale 态射；$Y \to Z$ 为 \'etale 态射；而 $h = f$ 正如 Definition
\ref{descent-definition-local-source-target} 第 (3) 部分所述。因此 $g \circ f$
具有 $\mathcal{P}$。

\medskip\noindent
(4) 的证明。设 $f : X \to Y$ 为态射，$g : Y \to Z$ 为 \'etale 态射，
并且 $g \circ f$ 具有 $\mathcal{P}$。由 Definition
\ref{descent-definition-local-source-target} 第 (2) 部分，
$\text{pr}_Y : Y \times_Z X \to Y$ 具有 $\mathcal{P}$。另一方面，态射
$(f, 1) : X \to Y \times_Z X$ 是 \'etale 投影
$\text{pr}_X : Y \times_Z X \to X$ 的截面，因而为 \'etale 态射；见
Morphisms, Lemma \ref{morphisms-lemma-etale-permanence}。所以由 Definition
\ref{descent-definition-local-source-target} 第 (1) 部分，
$f = \text{pr}_Y \circ (f, 1)$ 具有 $\mathcal{P}$。
\end{proof}

\noindent
下面的引理是 Morphisms, Lemma
\ref{morphisms-lemma-locally-P-characterize} 的类似结果。

\begin{lemma}
\label{descent-lemma-local-source-target-characterize}
设 $\mathcal{P}$ 为在源与目标上 \'etale 局部的概形态射性质。设
$f : X \to Y$ 为概形态射。下列条件等价：
\begin{enumerate}
\item[(a)] $f$ 具有性质 $\mathcal{P}$；
\item[(b)] 对每个 $x \in X$，存在芽的 \'etale 态射
$a : (U, u) \to (X, x)$、\'etale 态射 $b : V \to Y$ 和态射
$h : U \to V$，使得 $f \circ a = b \circ h$，且 $h$ 具有
$\mathcal{P}$；
\item[(c)] 对任意交换图
$$
\xymatrix{
U \ar[d]_a \ar[r]_h & V \ar[d]^b \\
X \ar[r]^f & Y
}
$$
，若 $a$、$b$ 为 \'etale 态射，则态射 $h$ 具有 $\mathcal{P}$；
\item[(d)] 对某个如 (c) 的图，若 $a : U \to X$ 满，则 $h$ 具有
$\mathcal{P}$；
\item[(e)] 存在 \'etale 覆盖 $\{Y_i \to Y\}_{i \in I}$，使得每个基变换
$Y_i \times_Y X \to Y_i$ 都具有 $\mathcal{P}$；
\item[(f)] 存在 \'etale 覆盖 $\{X_i \to X\}_{i \in I}$，使得每个复合
$X_i \to Y$ 都具有 $\mathcal{P}$；
\item[(g)] 存在 \'etale 覆盖 $\{Y_i \to Y\}_{i \in I}$，并且对每个
$i \in I$，存在 \'etale 覆盖
$\{X_{ij} \to Y_i \times_Y X\}_{j \in J_i}$，使得每个态射
$X_{ij} \to Y_i$ 都具有 $\mathcal{P}$。
\end{enumerate}
\end{lemma}

\begin{proof}
(a) 与 (b) 的等价性是 Definition \ref{descent-definition-local-source-target} 的
一部分。(a) 与 (e) 的等价性是 Lemma
\ref{descent-lemma-local-source-target-implies} 的第 (2) 部分。(a) 与 (f) 的等价性
是 Lemma \ref{descent-lemma-local-source-target-implies} 的第 (1) 部分。由于现在
(a) 与 (e)、(f) 等价，可知 (a) 与 (g) 等价。

\medskip\noindent
显然 (c) 蕴含 (a)。若 (a) 成立，则对任意如 (c) 的图，由 Definition
\ref{descent-definition-local-source-target} 第 (1) 部分，态射 $f \circ a$ 具有
$\mathcal{P}$；进而由 Lemma \ref{descent-lemma-local-source-target-implies}
第 (4) 部分，$h$ 具有 $\mathcal{P}$。所以 (a) 与 (c) 等价。显然 (c)
蕴含 (d)。为证明 (d) 蕴含 (a)，假设有如 (c) 的图，其中
$a : U \to X$ 满，且 $h$ 具有 $\mathcal{P}$。由 Lemma
\ref{descent-lemma-local-source-target-implies} 第 (3) 部分，$b \circ h$ 具有
$\mathcal{P}$。由于 $\{a : U \to X\}$ 是 \'etale 覆盖，由 Lemma
\ref{descent-lemma-local-source-target-implies} 第 (1) 部分，可知 $f$ 具有
$\mathcal{P}$。
\end{proof}

\noindent
下面这个引理的结论似乎并非形式的；也就是说，它确实用到了 \'etale 态射
几何中的某些内容。

\begin{lemma}
\label{descent-lemma-etale-local-source-target}
设 $\mathcal{P}$ 为概形态射的一个性质。假设
\begin{enumerate}
\item $\mathcal{P}$ 在源上为 \'etale 局部的；
\item $\mathcal{P}$ 在目标上为 \'etale 局部的；
\item $\mathcal{P}$ 在与开浸入后复合时稳定：若 $f : X \to Y$ 具有
$\mathcal{P}$，且 $Y \subset Z$ 为开子概形，则 $X \to Z$ 具有
$\mathcal{P}$。
\end{enumerate}
则 $\mathcal{P}$ 在源与目标上为 \'etale 局部的。
\end{lemma}

\begin{proof}
设 $\mathcal{P}$ 为满足本引理条件 (1)、(2)、(3) 的概形态射性质。由 Lemma
\ref{descent-lemma-precompose-property-local-source}，$\mathcal{P}$ 在与 \'etale
态射预复合时稳定。由 Lemma \ref{descent-lemma-pullback-property-local-target}，
$\mathcal{P}$ 在 \'etale 基变换下稳定。因此只需证明 Definition
\ref{descent-definition-local-source-target} 的第 (3) 部分成立。

\medskip\noindent
更准确地，假设概形态射 $f : X \to Y$ 满足 Definition
\ref{descent-definition-local-source-target} 第 (3)(b) 部分。换言之，对每个
$x \in X$，存在 \'etale 态射 $a_x : U_x \to X$、映到 $x$ 的点
$u_x \in U_x$、\'etale 态射 $b_x : V_x \to Y$ 及态射
$h_x : U_x \to V_x$，使得 $f \circ a_x = b_x \circ h_x$，且 $h_x$ 具有
$\mathcal{P}$。一旦证明 $f$ 具有 $\mathcal{P}$，本引理的证明就完成了。
置 $U = \coprod U_x$、$a = \coprod a_x$、$V = \coprod V_x$、
$b = \coprod b_x$ 及 $h = \coprod h_x$。得到交换图
$$
\xymatrix{
U \ar[d]_a \ar[r]_h & V \ar[d]^b \\
X \ar[r]^f & Y
}
$$
，其中 $a$、$b$ 为 \'etale 态射，且 $a$ 满。注意，$h$ 具有
$\mathcal{P}$，因为每个 $h_x$ 都如此，且 $\mathcal{P}$ 在目标上为
\'etale 局部的。由于 $a$ 满，而 $\mathcal{P}$ 在源上为 \'etale 局部的，
只需证明 $b \circ h$ 具有 $\mathcal{P}$。这样，本引理归约为证明
$\mathcal{P}$ 在与 \'etale 态射后复合时稳定。

\medskip\noindent
在余下的证明中，设 $f : X \to Y$ 为具有性质 $\mathcal{P}$ 的态射，且
$g : Y \to Z$ 为 \'etale 态射。考虑下列陈述：
\begin{enumerate}
\item[(-)] 不加额外假设时，$g \circ f$ 具有性质 $\mathcal{P}$。
\item[(A)] 只要 $Z$ 为仿射概形，$g \circ f$ 就具有性质 $\mathcal{P}$。
\item[(AA)] 只要 $X$ 和 $Z$ 为仿射概形，$g \circ f$ 就具有性质
$\mathcal{P}$。
\item[(AAA)] 只要 $X$、$Y$ 和 $Z$ 均为仿射概形，$g \circ f$ 就具有性质
$\mathcal{P}$。
\end{enumerate}
一旦证明 (-)，本引理的证明就完成了。

\medskip\noindent
断言 1：(AAA) $\Rightarrow$ (AA)。事实上，设 $f : X \to Y$、
$g : Y \to Z$ 如上，且 $X$、$Z$ 为仿射概形。由于 $X$ 仿射，因而拟紧，
可以找到有限多个仿射开集 $Y_i \subset Y$，$i = 1, \ldots, n$，使得
$X = \bigcup_{i = 1, \ldots, n} f^{-1}(Y_i)$。置 $X_i = f^{-1}(Y_i)$。
由 Lemma \ref{descent-lemma-pullback-property-local-target}，每个态射
$X_i \to Y_i$ 都具有 $\mathcal{P}$。由于 $\mathcal{P}$ 在目标上为
\'etale 局部的，
$\coprod_{i = 1, \ldots, n} X_i \to \coprod_{i = 1, \ldots, n} Y_i$
具有 $\mathcal{P}$。将 (AAA) 应用于
$\coprod_{i = 1, \ldots, n} X_i \to \coprod_{i = 1, \ldots, n} Y_i$
及 \'etale 态射
$\coprod_{i = 1, \ldots, n} Y_i \to Z$，可知
$\coprod_{i = 1, \ldots, n} X_i \to Z$ 具有 $\mathcal{P}$。现在
$\{\coprod_{i = 1, \ldots, n} X_i \to X\}$ 是 \'etale 覆盖；由于
$\mathcal{P}$ 在源上为 \'etale 局部的，可知 $X \to Z$ 具有
$\mathcal{P}$，正如所需。

\medskip\noindent
断言 2：(AAA) $\Rightarrow$ (A)。事实上，设 $f : X \to Y$、
$g : Y \to Z$ 如上，且 $Z$ 为仿射概形。选取仿射开覆盖
$X = \bigcup X_i$。由于 $\mathcal{P}$ 在源上为 \'etale 局部的，每个
$f|_{X_i} : X_i \to Y$ 都具有 $\mathcal{P}$。由断言 1 从 (AAA) 推出的
(AA)，对每个 $i$，$X_i \to Z$ 都具有 $\mathcal{P}$。由于
$\{X_i \to X\}$ 是 \'etale 覆盖，而 $\mathcal{P}$ 在源上为 \'etale
局部的，可知 $X \to Z$ 具有 $\mathcal{P}$。

\medskip\noindent
断言 3：(AAA) $\Rightarrow$ (-)。事实上，设 $f : X \to Y$、
$g : Y \to Z$ 如上。选取仿射开覆盖 $Z = \bigcup Z_i$。置
$Y_i = g^{-1}(Z_i)$ 及 $X_i = f^{-1}(Y_i)$。由 Lemma
\ref{descent-lemma-pullback-property-local-target}，每个态射 $X_i \to Y_i$ 都具有
$\mathcal{P}$。由断言 2 从 (AAA) 推出的 (A)，对每个 $i$，
$X_i \to Z_i$ 都具有 $\mathcal{P}$。由于 $\mathcal{P}$ 在目标上为局部的，
且 $X_i = (g \circ f)^{-1}(Z_i)$，可知 $X \to Z$ 具有 $\mathcal{P}$。

\medskip\noindent
因此为证明本引理，只需证明 (AAA)。设 $f : X \to Y$ 和 $g : Y \to Z$
如上，且 $X, Y, Z$ 均为仿射概形。注意，仿射概形的 \'etale 态射具有
普遍有界的纤维；见 Morphisms, Lemma
\ref{morphisms-lemma-etale-locally-quasi-finite} 和 Lemma
\ref{morphisms-lemma-locally-quasi-finite-qc-source-universally-bounded}。
因此可以对界定 $Y \to Z$ 的纤维度数的整数 $n$ 作归纳。在 \'etale 态射
的情形，该整数的描述见 Morphisms, Lemma
\ref{morphisms-lemma-etale-universally-bounded}。若 $n = 1$，则
$Y \to Z$ 为开浸入；见 Lemma
\ref{descent-lemma-universally-injective-etale-open-immersion}；结论由本引理的假设
(3) 推出。假设 $n > 1$。

\medskip\noindent
考虑下列交换图：
$$
\xymatrix{
X \times_Z Y \ar[d] \ar[r]_{f_Y} &
Y \times_Z Y \ar[d] \ar[r]_-{\text{pr}} &
Y \ar[d] \\
X \ar[r]^f &
Y \ar[r]^g &
Z
}
$$
注意，我们有开闭子概形分解
$Y \times_Z Y = \Delta_{Y/Z}(Y) \amalg Y'$；见 Morphisms, Lemma
\ref{morphisms-lemma-diagonal-unramified-morphism}。作为基变换，第二投影
$\text{pr} : Y \times_Z Y \to Y$ 的纤维度数以 $n$ 为界；见 Morphisms,
Lemma \ref{morphisms-lemma-base-change-universally-bounded}。另一方面，
$\text{pr}|_{\Delta(Y)} : \Delta(Y) \to Y$ 为同构，且每个纤维恰有一个点。
因此应用 Morphisms, Lemma
\ref{morphisms-lemma-etale-universally-bounded}，可知限制
$\text{pr}|_{Y'} : Y' \to Y$ 的纤维度数以 $n - 1$ 为界。置
$X' = f_Y^{-1}(Y')$。图示如下：
$$
\xymatrix{
X \amalg X' \ar@{=}[d] \ar[r]_-{f \amalg f'} &
\Delta(Y) \amalg Y' \ar@{=}[d] \ar[r] &
Y \ar@{=}[d] \\
X \times_Z Y \ar[r]^{f_Y} &
Y \times_Z Y \ar[r]^-{\text{pr}} &
Y
}
$$
由于 $\mathcal{P}$ 在目标上为 \'etale 局部的，因而在 \'etale 基变换下
稳定（见 Lemma \ref{descent-lemma-pullback-property-local-target}），可知 $f_Y$
具有 $\mathcal{P}$。又因 $\mathcal{P}$ 在源上为 \'etale 局部的，
$f' = f_Y|_{X'}$ 具有 $\mathcal{P}$。由归纳假设，$X' \to Y$ 具有
$\mathcal{P}$。由于 $\mathcal{P}$ 在源上为局部的，并且
% P05-ERR-0043
$\{X \to X \times_Z Y, X' \to X \times_Y Z\}$ 是 \'etale 覆盖，可知
$\text{pr} \circ f_Y$ 具有 $\mathcal{P}$。注意，可将 $g \circ f$ 视为态射
$g \circ f : X \to g(Y)$。由于 $\text{pr} \circ f_Y$ 是态射
$g \circ f : X \to g(Y)$ 沿 \'etale 覆盖 $\{Y \to g(Y)\}$ 的拉回，且
$\mathcal{P}$ 在目标上为 \'etale 局部的，可知
$g \circ f : X \to g(Y)$ 具有性质 $\mathcal{P}$。最后，再次应用本引理的
假设 (3)，可知 $g \circ f : X \to Z$ 具有性质 $\mathcal{P}$。
\end{proof}

\begin{remark}
\label{descent-remark-list-local-source-target}
利用 Lemma \ref{descent-lemma-etale-local-source-target} 及本章前面各节完成的工作，
很容易列出在源与目标上为 \'etale 局部的态射类型。在每种情形中，我们
分别列出蕴含该性质在源上为 \'etale 局部的引理，以及蕴含它在目标上为
\'etale 局部的引理。在每种情形中，Lemma
\ref{descent-lemma-etale-local-source-target} 的第三个假设都显然成立，故予省略。
列表如下：
\begin{enumerate}
\item 平坦；见 Lemmas \ref{descent-lemma-flat-fpqc-local-source} 和
\ref{descent-lemma-descending-property-flat}；
\item 局部有限表示；见 Lemmas
\ref{descent-lemma-locally-finite-presentation-fppf-local-source} 和
\ref{descent-lemma-descending-property-locally-finite-presentation}；
\item 局部有限型；见 Lemmas \ref{descent-lemma-locally-finite-type-fppf-local-source}
和 \ref{descent-lemma-descending-property-locally-finite-type}；
\item 普遍开；见 Lemmas \ref{descent-lemma-universally-open-fppf-local-source} 和
\ref{descent-lemma-descending-property-universally-open}；
\item syntomic；见 Lemmas \ref{descent-lemma-syntomic-syntomic-local-source} 和
\ref{descent-lemma-descending-property-syntomic}；
\item 光滑；见 Lemmas \ref{descent-lemma-smooth-smooth-local-source} 和
\ref{descent-lemma-descending-property-smooth}；
\item \'etale；见 Lemmas \ref{descent-lemma-etale-etale-local-source} 和
\ref{descent-lemma-descending-property-etale}；
\item 局部拟有限；见 Lemmas
\ref{descent-lemma-locally-quasi-finite-etale-local-source} 和
\ref{descent-lemma-descending-property-quasi-finite}；
\item 非分歧；见 Lemmas \ref{descent-lemma-unramified-etale-local-source} 和
\ref{descent-lemma-descending-property-unramified}；
\item G-非分歧；见 Lemmas \ref{descent-lemma-unramified-etale-local-source} 和
\ref{descent-lemma-descending-property-unramified}；
\item 可按需在此添加更多性质。
\end{enumerate}
\end{remark}

\begin{remark}
\label{descent-remark-compare-definitions}
至此，对于态射性质 $\mathcal{P}$ ``在源和目标上为 \'etale 局部''的含义，
我们有三种可能的定义：
\begin{enumerate}
\item[(ST)] $\mathcal{P}$ 在源上为 \'etale 局部的，且 $\mathcal{P}$ 在
目标上为 \'etale 局部的；
\item[(DM)] （Deligne 和 Mumford 的论文 \cite[Page 100]{DM} 中的定义）
对每个图
$$
\xymatrix{
U \ar[d]_a \ar[r]_h & V \ar[d]^b \\
X \ar[r]^f & Y
}
$$
，若竖直箭头为满的 \'etale 态射，则
$\mathcal{P}(h) \Leftrightarrow \mathcal{P}(f)$；
\item[(SP)] $\mathcal{P}$ 在源与目标上为 \'etale 局部的。
\end{enumerate}
本节已经看到 (SP) $\Rightarrow$ (DM) $\Rightarrow$ (ST)。Examples
\ref{descent-example-silly-one} 和 \ref{descent-example-silly-two} 表明，两个蕴含均不能
反转。最后，Lemma \ref{descent-lemma-etale-local-source-target} 表明，对在与开浸入
后复合时稳定的态射性质，这一区别会消失；实践中总是这种情形。
\end{remark}

\begin{lemma}
\label{descent-lemma-etale-etale-local-source-target}
设 $\mathcal{P}$ 为在源与目标上 \'etale 局部的概形态射性质。给定概形的交换图
$$
\vcenter{
\xymatrix{
X' \ar[d]_{g'} \ar[r]_{f'} & Y' \ar[d]^g \\
X \ar[r]^f & Y
}
}
\quad\text{带标记点}\quad
\vcenter{
\xymatrix{
x' \ar[d] \ar[r] & y' \ar[d] \\
x \ar[r] & y
}
}
$$
，若 $g'$ 在 $x'$ 处为 \'etale 态射，且 $g$ 在 $y'$ 处为 \'etale 态射，
则
$x \in W(f) \Leftrightarrow x' \in W(f')$
，其中 $W(-)$ 如 Lemma \ref{descent-lemma-largest-open-of-the-source} 中所定义。
\end{lemma}

\begin{proof}
Lemma \ref{descent-lemma-largest-open-of-the-source} 适用，因为由 Lemma
\ref{descent-lemma-local-source-target-implies}，$\mathcal{P}$ 在源上为 \'etale
局部的。

\medskip\noindent
假设 $x \in W(f)$。设 $U' \subset X'$ 和 $V' \subset Y'$ 分别为 $x'$ 和
$y'$ 的开邻域，使得 $f'(U') \subset V'$、$g'(U') \subset W(f)$，且
$g'|_{U'}$ 与 $g|_{V'}$ 为 \'etale 态射。由 Definition
\ref{descent-definition-local-source-target} 的性质 (1)，
$f \circ g'|_{U'} = g \circ f'|_{U'}$ 具有 $\mathcal{P}$。再由 Lemma
\ref{descent-lemma-local-source-target-implies} 的 (4)，
$f'|_{U'} : U' \to V'$ 具有性质 $\mathcal{P}$。进而由 Lemma
\ref{descent-lemma-local-source-target-implies} 的 (3)，
% P05-ERR-0044
$f'_{U'} : U' \to Y'$ 具有 $\mathcal{P}$。所以由定义
$U' \subset W(f')$，进而 $x' \in W(f')$。

\medskip\noindent
假设 $x' \in W(f')$。设 $U' \subset X'$ 和 $V' \subset Y'$ 分别为 $x'$ 和
$y'$ 的开邻域，使得 $f'(U') \subset V'$、$U' \subset W(f')$，且
$g'|_{U'}$ 与 $g|_{V'}$ 为 \'etale 态射。由 $W(f')$ 的定义，
$U' \to Y'$ 具有 $\mathcal{P}$。由 Lemma
\ref{descent-lemma-local-source-target-implies} 的 (4)，$U' \to V'$ 具有
$\mathcal{P}$。再由 Lemma \ref{descent-lemma-local-source-target-implies} 的 (3)，
$U' \to Y$ 具有 $\mathcal{P}$。令
$U \subset X$ 为 \'etale（因而开）态射
% P05-ERR-0045
$g'|_U' : U' \to X$ 的像。则 $\{U' \to U\}$ 为 \'etale 覆盖；由 Lemma
\ref{descent-lemma-local-source-target-implies} 的 (1)，可知 $U \to Y$ 具有
$\mathcal{P}$。所以由定义 $U \subset W(f)$，进而 $x \in W(f)$。
\end{proof}

\begin{lemma}
\label{descent-lemma-orbits}
设 $k$ 为域，设 $n \geq 2$。对满足 $1 \leq i, j \leq n$、$i \not = j$ 和
$d \geq 0$ 的 i、j、d，以 $T_{i, j, d}$ 表示 $\mathbf{A}^n_k$ 的自同构；
它在坐标中由下式给出：
$$
(x_1, \ldots, x_n) \longmapsto
(x_1, \ldots, x_{i - 1}, x_i + x_j^d, x_{i + 1}, \ldots, x_n)
$$
设 $W \subset \mathbf{A}^n_k$ 为非空开子概形，并且对上述所有 $i, j, d$，
都有 $T_{i, j, d}(W) = W$。则或者 $W = \mathbf{A}^n_k$，或者 $k$ 的特征
为 $p > 0$，且 $\mathbf{A}^n_k \setminus W$ 是由闭点组成的有限集，这些点
的坐标在 $\mathbf{F}_p$ 上代数。
\end{lemma}

\begin{proof}
为证明本结论，可以用任意扩域替换 $k$。设 $Z$ 为
$\mathbf{A}^n_k \setminus W$ 的一个不可约分支。为得到矛盾，假设
$\dim(Z) \geq 1$。则存在域扩张 $k'/k$ 及
$Z_{k'} \subset \mathbf{A}^n_{k'}$ 的 $k'$-值点
$\xi = (\xi_1, \ldots, \xi_n) \in (k')^n$，使得 $x_1, \ldots, x_n$
中至少一个在素域上超越。断言：$\xi$ 在由变换 $T_{i, j, d}$ 生成的群作用
下的轨道在 $\mathbf{A}^n_{k'}$ 中 Zariski 稠密。该断言将给出所需的矛盾。

\medskip\noindent
若 $k'$ 的特征为零，则仅使用算子 $T_{i, j, 0}$ 就已经足够，因为这些算子
将 $\xi$ 变换为诸点
$$
(\xi_1 + a_1, \ldots, \xi_n + a_n)
$$
，其中 $(a_1, \ldots, a_n) \in \mathbf{Z}_{\geq 0}^n$ 任意。若特征为
$p > 0$，重新编号后可以假设 $\xi_n$ 在 $\mathbf{F}_p$ 上超越。对
$i < n$ 依次应用算子 $T_{i, n, d}$，可见 $\xi$ 的轨道包含诸元素
$$
(\xi_1 + P_1(\xi_n), \ldots, \xi_{n - 1} + P_{n - 1}(\xi_n), \xi_n)
$$
，其中 $(P_1, \ldots, P_{n - 1}) \in \mathbf{F}_p[t]$ 任意。因此轨道的
Zariski 闭包包含坐标超平面 $x_n = \xi_n$。对另一个坐标重复论证，可知
对任意 $P \in \mathbf{F}_p[t]$，只要 $\xi_i + P(\xi_n)$ 在
$\mathbf{F}_p$ 上超越，该 Zariski 闭包就包含
$x_i = \xi_i + P(\xi_n)$。这样的 $P$ 有无穷多个，故断言成立。

\medskip\noindent
当然，若 $Z = \{z\}$ 的维数为 $0$，且 $z$ 在 $\kappa(z)$ 中的坐标并非
都在 $\mathbf{F}_p$ 上代数，上一段的论证也同样适用。引理得证。
\end{proof}

\begin{lemma}
\label{descent-lemma-etale-tau-local-source-target}
设 $\mathcal{P}$ 为概形态射的一个性质。假设
\begin{enumerate}
\item $\mathcal{P}$ 在源上为 \'etale 局部的；
\item $\mathcal{P}$ 在目标上为光滑局部的；
\item $\mathcal{P}$ 在与开浸入后复合时稳定：若 $f : X \to Y$ 具有
$\mathcal{P}$，且 $Y \subset Z$ 为开子概形，则 $X \to Z$ 具有
$\mathcal{P}$。
\end{enumerate}
给定概形的交换图
$$
\vcenter{
\xymatrix{
X' \ar[d]_{g'} \ar[r]_{f'} & Y' \ar[d]^g \\
X \ar[r]^f & Y
}
}
\quad\text{带标记点}\quad
\vcenter{
\xymatrix{
x' \ar[d] \ar[r] & y' \ar[d] \\
x \ar[r] & y
}
}
$$
，使得
% P05-ERR-0046
$g$ 在 $y'$ 处光滑，并且 $X' \to X \times_Y Y'$ 在 $x'$ 处为 \'etale
态射，则 $x \in W(f) \Leftrightarrow x' \in W(f')$，其中 $W(-)$ 如
Lemma \ref{descent-lemma-largest-open-of-the-source} 中所定义。
\end{lemma}

\begin{proof}
由于 $\mathcal{P}$ 在源上为 \'etale 局部的，可知 $x \in W(f)$ 当且仅当
$x$ 在 $X \times_Y Y'$ 中的像属于 $W(X \times_Y Y' \to Y')$。因此可以
假设本引理中的图是 Cartesian 的。

\medskip\noindent
假设 $x \in W(f)$。由于 $\mathcal{P}$ 在目标上为光滑局部的，
$(g')^{-1}W(f) = W(f) \times_Y Y' \to Y'$ 具有 $\mathcal{P}$。所以
$(g')^{-1}W(f) \subset W(f')$，从而 $x' \in W(f')$。

\medskip\noindent
假设 $x' \in W(f')$。对 $y'$ 的任意开邻域 $V' \subset Y'$，可以用 $V'$
替换 $Y'$，并用 $U' = (f')^{-1}V'$ 替换 $X'$：这是因为
$V' \to Y'$ 光滑，因而 $W(f') \to Y'$ 的基变换
$W(f') \cap U' \to V'$ 具有性质 $\mathcal{P}$。所以可以假设存在 $Y$ 上的
\'etale 态射 $Y' \to \mathbf{A}^n_Y$；见 Morphisms, Lemma
\ref{morphisms-lemma-smooth-etale-over-affine-space}。图示如下：
$$
\xymatrix{
X' \ar[r] \ar[d] & Y' \ar[d] \\
\mathbf{A}^n_X \ar[r]_{f_n} \ar[d] & \mathbf{A}^n_Y \ar[d] \\
X \ar[r]^f & Y
}
$$
由 Lemma \ref{descent-lemma-etale-local-source-target}（并利用 \'etale 覆盖也是光滑
覆盖），$\mathcal{P}$ 在源与目标上为 \'etale 局部的。由 Lemma
\ref{descent-lemma-etale-etale-local-source-target}，$W(f')$ 是开集
$W(f_n) \subset \mathbf{A}^n_X$ 的逆像。特别地，$W(f_n)$ 包含一个位于
$x$ 上方的点。用 $W(f_n)$ 的像（它是开集）替换 $X$ 后，可以假设
$W(f_n) \to X$ 满。断言：$W(f_n) = \mathbf{A}^n_X$。该断言蕴含 $f$ 具有
$\mathcal{P}$，因为 $\mathcal{P}$ 在光滑拓扑中为局部的，并且
$\{\mathbf{A}^n_Y \to Y\}$ 是光滑覆盖。

\medskip\noindent
本质上，该断言源于以下事实：$W(f_n) \subset \mathbf{A}^n_X$ 是一个与每个
$\mathbf{A}^n_X \to X$ 的纤维相交的``平移不变''开集。不过，若要按这一路径
给出论证，就必须在 $Y$ 上作 \'etale 局部化以产生足够多的平移，这会有些
麻烦。我们改用 Lemma \ref{descent-lemma-orbits} 的自同构及仿射空间的 \'etale
态射。可以假设 $n \geq 2$。事实上，若 $n = 0$，结论已经成立。若
$n = 1$，考虑图
$$
\xymatrix{
\mathbf{A}^2_X \ar[r]_{f_2} \ar[d]_p & \mathbf{A}^2_Y \ar[d] \\
\mathbf{A}^1_X \ar[r]^{f_1} & \mathbf{A}^1_Y
}
$$
由本证明第一段，有 $p^{-1}(W(f_1)) \subset W(f_2)$。所以
$W(f_2) \to X$ 仍然满，我们可以改用 $f_2$。假设 $n \geq 2$。

\medskip\noindent
对满足 $1 \leq i, j \leq n$、$i \not = j$ 和 $d \geq 0$ 的任意
i、j、d，以 $T_{i, j, d}$ 表示 Lemma \ref{descent-lemma-orbits} 中定义的
$\mathbf{A}^n$ 的自同构。于是得到交换图
$$
\xymatrix{
\mathbf{A}^n_X \ar[r]_{f_n} \ar[d]_{T_{i, j, d}} &
\mathbf{A}^n_Y \ar[d]^{T_{i, j, d}} \\
\mathbf{A}^n_X \ar[r]^{f_n} & \mathbf{A}^n_Y
}
$$
，其竖直箭头为同构。因此
$T_{i, j, d}(W(f_n)) = W(f_n)$。应用 Lemma \ref{descent-lemma-orbits}，可知对任意
$x \in X$，纤维 $W(f_n)_x \subset \mathbf{A}^n_x$ 或者等于
$\mathbf{A}^n_x$（这正是我们要证明的），或者 $\kappa(x)$ 的特征为
$p > 0$，且 $W(f_n)_x$ 是由闭点组成的有限集
$Z_x \subset \mathbf{A}^n_x$ 的补集。第二种可能不成立。事实上，考虑
由下式给出的态射 $T_p : \mathbf{A}^n \to \mathbf{A}^n$：
$$
(x_1, \ldots, x_n) \mapsto (x_1 - x_1^p, \ldots, x_n - x_n^p)
$$
如上，我们得到交换图
$$
\xymatrix{
\mathbf{A}^n_X \ar[r]_{f_n} \ar[d]_{T_p} &
\mathbf{A}^n_Y \ar[d]^{T_p} \\
\mathbf{A}^n_X \ar[r]^{f_n} & \mathbf{A}^n_Y
}
$$
态射 $T_p : \mathbf{A}^n_X \to \mathbf{A}^n_X$ 在位于 $x$ 上方的每个点
处为 \'etale 态射，而态射 $T_p : \mathbf{A}^n_Y \to \mathbf{A}^n_Y$
在位于 $x$ 在 $Y$ 中的像上方的每个点处为 \'etale 态射。（省略细节；
提示：计算导数。）因此
% P05-ERR-0047
$$
T_p^{-1}(W) \cap \mathbf{A}^n_x = W \cap \mathbf{A}^n_x
$$
，这是由 Lemma \ref{descent-lemma-etale-etale-local-source-target} 得到的（我们已经
知道 $\mathcal{P}$ 在源与目标上为 \'etale 局部的）。由于
$T_p : \mathbf{A}^n_x \to \mathbf{A}^n_x$ 是次数 $p^n > 1$ 的有限
\'etale 态射，若 $Z_x$ 非空，则它包含更大的集合 $T_p^{-1}(Z_x)$。
这个矛盾完成了证明。
\end{proof}





\section{在源与目标上局部的芽态射性质}
\label{descent-section-local-source-target-at-point}

\noindent
本节讨论 Section \ref{descent-section-properties-etale-local-source-target} 的内容
对于概形芽态射的类似版本。

\begin{definition}
\label{descent-definition-local-source-target-at-point}
设 $\mathcal{Q}$ 为概形芽态射的一个性质。若对任意交换图
$$
\xymatrix{
(U', u') \ar[d]_a \ar[r]_{h'} & (V', v') \ar[d]^b \\
(U, u) \ar[r]^h & (V, v)
}
$$
，其中竖直箭头为芽的 \'etale 态射，都有
$\mathcal{Q}(h) \Leftrightarrow \mathcal{Q}(h')$，则称 $\mathcal{Q}$
{\it 在源与目标上为 \'etale 局部的}。
\end{definition}

\begin{lemma}
\label{descent-lemma-local-source-target-global-implies-local}
设 $\mathcal{P}$ 为在源与目标上 \'etale 局部的概形态射性质。考虑按下列规则
定义的芽态射性质 $\mathcal{Q}$：
$$
\mathcal{Q}((X, x) \to (S, s))
\Leftrightarrow
\text{存在一个代表元 }U \to S
\text{ 其具有 }\mathcal{P}
$$
则 $\mathcal{Q}$ 在源与目标上为 \'etale 局部的，含义如 Definition
\ref{descent-definition-local-source-target-at-point}。
\end{lemma}

\begin{proof}
若芽态射 $(X, x) \to (S, s)$ 具有 $\mathcal{Q}$，则存在任意小的 $x$ 的邻域
$U \subset X$ 和 $s$ 的邻域 $V \subset S$，使得
$(X, x) \to (S, s)$ 的代表 $U \to V$ 具有 $\mathcal{P}$。这由 Lemma
\ref{descent-lemma-local-source-target-implies} 推出。设
$$
\xymatrix{
(U', u') \ar[r]_{h'} \ar[d]_a & (V', v') \ar[d]^b \\
(U, u) \ar[r]^h & (V, v)
}
$$
如 Definition \ref{descent-definition-local-source-target-at-point} 所述。选取
$U_1 \subset U$ 及 $h$ 的代表 $h_1 : U_1 \to V$。选取
$V'_1 \subset V'$ 及 $b$ 的 \'etale 代表 $b_1 : V'_1 \to V$（Definition
\ref{descent-definition-etale-morphism-germs}）。选取 $U'_1 \subset U'$ 及 $a$、
$h'$ 的代表 $a_1 : U'_1 \to U_1$、$h'_1 : U'_1 \to V'_1$，其中 $a_1$
为 \'etale 态射。缩小 $U'_1$ 后，可以假设
$h_1 \circ a_1 = b_1 \circ h'_1$。由证明开头的说明，我们要证明
$u' \in W(h'_1) \Leftrightarrow u \in W(h_1)$，其中 $W(-)$ 如 Lemma
\ref{descent-lemma-largest-open-of-the-source} 中所定义。因此本引理由 Lemma
\ref{descent-lemma-etale-etale-local-source-target} 推出。
\end{proof}

\begin{lemma}
\label{descent-lemma-local-source-target-local-implies-global}
设 $\mathcal{P}$ 为在源与目标上 \'etale 局部的概形态射性质。令
% P05-ERR-0048
$Q$ 为相伴的芽态射性质；见 Lemma
\ref{descent-lemma-local-source-target-global-implies-local}。设 $f : X \to Y$ 为
概形态射。下列条件等价：
\begin{enumerate}
\item $f$ 具有性质 $\mathcal{P}$；
\item 对每个 $x \in X$，芽态射 $(X, x) \to (Y, f(x))$ 具有性质
$\mathcal{Q}$。
\end{enumerate}
\end{lemma}

\begin{proof}
蕴含 (1) $\Rightarrow$ (2) 直接由定义给出。蕴含 (2) $\Rightarrow$ (1)
也由 Definition \ref{descent-definition-local-source-target} 的第 (3) 部分推出。
\end{proof}

\noindent
芽态射 $(X, x) \to (S, s)$ 决定一个良定义的局部环映射。因此下面的引理
是有意义的。

\begin{lemma}
\label{descent-lemma-flat-at-point}
芽态射的性质
$$
\mathcal{P}((X, x) \to (S, s)) =
\mathcal{O}_{S, s} \to \mathcal{O}_{X, x}\text{ 平坦}
$$
在源与目标上为 \'etale 局部的。
\end{lemma}

\begin{proof}
给定如 Definition \ref{descent-definition-local-source-target-at-point} 中的图，得到
如下局部环的局部同态图：
$$
\xymatrix{
\mathcal{O}_{U', u'} & \mathcal{O}_{V', v'} \ar[l] \\
\mathcal{O}_{U, u} \ar[u] & \mathcal{O}_{V, v} \ar[l] \ar[u]
}
$$
注意，竖直箭头为 \'etale 环映射的局部化；特别地，它们本质有限表示、
平坦且非分歧（见 Algebra, Section \ref{algebra-section-etale}）。特别地，
竖直映射忠实平坦；见 Algebra, Lemma
\ref{algebra-lemma-local-flat-ff}。现在，若上方水平箭头平坦，则将 Algebra,
Lemma \ref{algebra-lemma-flat-permanence} 应用于
$R = \mathcal{O}_{V, v}$、$S = \mathcal{O}_{U, u}$ 和
$M = \mathcal{O}_{U', u'}$，可知下方水平箭头平坦。若下方水平箭头平坦，
则环映射
$$
\mathcal{O}_{V', v'} \otimes_{\mathcal{O}_{V, v}} \mathcal{O}_{U, u}
\longleftarrow
\mathcal{O}_{V', v'}
$$
由 Algebra, Lemma \ref{algebra-lemma-flat-base-change} 可知平坦。而环映射
$$
\mathcal{O}_{U', u'}
\longleftarrow
\mathcal{O}_{V', v'} \otimes_{\mathcal{O}_{V, v}} \mathcal{O}_{U, u}
$$
是 $\mathcal{O}_{U, u}$ 的两个 \'etale 环扩张之间映射的局部化，因此由
Algebra, Lemma \ref{algebra-lemma-map-between-etale} 可知平坦。
\end{proof}

\begin{lemma}
\label{descent-lemma-etale-on-fiber}
考虑概形态射的交换图
$$
\xymatrix{
U' \ar[r] \ar[d] & V' \ar[d] \\
U \ar[r] & V
}
$$
，其中竖直箭头为 \'etale 态射，且点 $v' \in V'$ 映到 $v \in V$。则纤维
的态射 $U'_{v'} \to U_v$ 为 \'etale 态射。
\end{lemma}

\begin{proof}
注意，$U'_v \to U_v$ 是 \'etale 态射 $U' \to U$ 的基变换，因而为
\'etale 态射。概形 $U'_v$ 是 $V'_v$ 上的概形。由 Morphisms, Lemma
\ref{morphisms-lemma-etale-over-field}，概形 $V'_v$ 是 $\kappa(v)$ 的有限
可分域扩张的谱的不交并。其中一个是 $v' = \Spec(\kappa(v'))$。所以
$U'_{v'}$ 是 $U'_v$ 的开闭子概形，从而 $U'_{v'} \to U'_v \to U_v$
为 \'etale 态射（它是开浸入与 \'etale 态射的复合；见 Morphisms,
Section \ref{morphisms-section-etale}）。
\end{proof}

\noindent
给定概形的芽态射 $(X, x) \to (S, s)$，可以把{\it 纤维}定义为芽
$(U_s, x)$ 的同构类，其中 $U \to S$ 为任意代表。我们常常略微滥用记号，
直接写成 $(X_s, x)$。

\begin{lemma}
\label{descent-lemma-dimension-local-ring-fibre}
设 $d \in \{0, 1, 2, \ldots, \infty\}$。芽态射的性质
$$
\mathcal{P}_d((X, x) \to (S, s)) =
\text{局部环 }
\mathcal{O}_{X_s, x}
\text{ 的纤维维数为 }d
$$
在源与目标上为 \'etale 局部的。
\end{lemma}

\begin{proof}
给定如 Definition \ref{descent-definition-local-source-target-at-point} 中的图，得到
把 $u'$ 映到 $u$ 的纤维的 \'etale 态射 $U'_{v'} \to U_v$；见 Lemma
\ref{descent-lemma-etale-on-fiber}。因此结论由 Lemma
\ref{descent-lemma-dimension-local-ring-local} 推出。
\end{proof}

\begin{lemma}
\label{descent-lemma-transcendence-degree-at-point}
设 $r \in \{0, 1, 2, \ldots, \infty\}$。芽态射的性质
$$
\mathcal{P}_r((X, x) \to (S, s))
\Leftrightarrow
\text{trdeg}_{\kappa(s)} \kappa(x) = r
$$
在源与目标上为 \'etale 局部的。
\end{lemma}

\begin{proof}
给定如 Definition \ref{descent-definition-local-source-target-at-point} 中的图，得到
如下局部环的局部同态图：
$$
\xymatrix{
\mathcal{O}_{U', u'} & \mathcal{O}_{V', v'} \ar[l] \\
\mathcal{O}_{U, u} \ar[u] & \mathcal{O}_{V, v} \ar[l] \ar[u]
}
$$
注意，竖直箭头为 \'etale 环映射的局部化；特别地，它们非分歧（见
Algebra, Section \ref{algebra-section-etale}）。所以
$\kappa(u')/\kappa(u)$ 和 $\kappa(v')/\kappa(v)$ 是有限可分域扩张。因此
% P05-ERR-0049
$\text{trdeg}_{\kappa(v)} \kappa(u) = \text{trdeg}_{\kappa(v')} \kappa(u)$
，这证明了本引理。
\end{proof}

\noindent
设 $(X, x)$ 为概形的芽。$X$ 在 $x$ 处的维数，是 $x$ 在 $X$ 中各开邻域的
维数之最小值；任意足够小的开邻域都具有此维数。因此，这是芽的同构类的
不变量。我们把它简记为 $\dim_x(X)$。

\begin{lemma}
\label{descent-lemma-dimension-at-point}
设 $d \in \{0, 1, 2, \ldots, \infty\}$。芽态射的性质
$$
\mathcal{P}_d((X, x) \to (S, s))
\Leftrightarrow
\dim_x (X_s) = d
$$
在源与目标上为 \'etale 局部的。
\end{lemma}

\begin{proof}
给定如 Definition \ref{descent-definition-local-source-target-at-point} 中的图，得到
把 $u'$ 映到 $u$ 的纤维的 \'etale 态射 $U'_{v'} \to U_v$；见 Lemma
\ref{descent-lemma-etale-on-fiber}。于是等式
$\dim_u(U_v) = \dim_{u'}(U'_{v'})$ 由 Lemma
\ref{descent-lemma-dimension-at-point-local} 推出。
\end{proof}







\section{概形上概形的下降数据}
\label{descent-section-descent-datum}

\noindent
本节的大部分论证都是形式的，仅依赖于下降数据的定义。在 Simplicial
Spaces, Section \ref{spaces-simplicial-section-simplicial-descent} 中，我们
将考察它与单纯概形的关系，这会在一定程度上澄清情形。

\begin{definition}
\label{descent-definition-descent-datum}
设 $f : X \to S$ 为概形态射。
\begin{enumerate}
\item 设 $V \to X$ 为 $X$ 上的概形。一个{\it $V/X/S$ 的下降数据}是
$X \times_S X$ 上的概形同构
$\varphi : V \times_S X \to X \times_S V$，它满足{\it 上循环条件}：图
$$
\xymatrix{
V \times_S X \times_S X \ar[rd]^{\varphi_{01}} \ar[rr]_{\varphi_{02}} &
&
X \times_S X \times_S V\\
&
X \times_S V \times_S X \ar[ru]^{\varphi_{12}}
}
$$
交换（记号含义显然）。
\item 也称对 $(V/X, \varphi)$ 是一个{\it 相对于 $X \to S$ 的下降数据}。
\item 相对于 $X \to S$ 的下降数据之间的一个{\it 态射
$g : (V/X, \varphi) \to (V'/X, \varphi')$}，是使下图交换的 $X$ 上概形
态射 $g : V \to V'$：
$$
\xymatrix{
V \times_S X \ar[r]_{\varphi} \ar[d]_{g \times \text{id}_X} &
X \times_S V \ar[d]^{\text{id}_X \times g} \\
V' \times_S X \ar[r]^{\varphi'} & X \times_S V'
}
$$
。
\end{enumerate}
\end{definition}

\noindent
上述定义会产生各种``奇妙的''恒等式。例如，$\varphi$ 沿对角态射
$\Delta : X \to X \times_S X$ 的拉回可视为态射
$\Delta^*\varphi : V \to V$。这是因为
$X \times_{\Delta, X \times_S X} (V \times_S X) = V$，并且
$X \times_{\Delta, X \times_S X} (X \times_S V) = V$。事实上，
$\Delta^*\varphi$ 等于恒等态射。若你不熟悉这些材料，这是一个很好的练习。

\begin{remark}
\label{descent-remark-easier}
设 $X \to S$ 为概形态射，设 $(V/X, \varphi)$ 为相对于 $X \to S$ 的下降
数据。可以把同构 $\varphi$ 看作同构
$$
(X \times_S X) \times_{\text{pr}_0, X} V
\longrightarrow
(X \times_S X) \times_{\text{pr}_1, X} V
$$
（作为 $X \times_S X$ 上的概形）。因此粗略地说，可以把 $\varphi$ 看作映射
$\varphi : \text{pr}_0^*V \to \text{pr}_1^*V$\footnote{遗憾的是，我们在这里
为箭头选择了``错误的''方向。根据``函数''与``空间''对偶的一般原则，
Definitions \ref{descent-definition-descent-datum} 和
\ref{descent-definition-descent-datum-for-family-of-morphisms} 中的方向，应当与
Definition \ref{descent-definition-descent-datum-quasi-coherent} 中所取的方向相反。}。
此时上循环条件断言
$\text{pr}_{02}^*\varphi =
\text{pr}_{12}^*\varphi \circ \text{pr}_{01}^*\varphi$.
从这个角度看，它与拟凝聚层上的下降数据情形十分相似。
\end{remark}

\noindent
下面给出具有固定目标的一族态射的情形下的定义。

\begin{definition}
\label{descent-definition-descent-datum-for-family-of-morphisms}
设 $S$ 为概形，设 $\{X_i \to S\}_{i \in I}$ 为以 $S$ 为目标的一族态射。
\begin{enumerate}
\item 一个{\it 相对于态射族 $\{X_i \to S\}$ 的下降数据
$(V_i, \varphi_{ij})$}由下列内容给出：对每个 $i \in I$，一个 $X_i$ 上的
概形 $V_i$；对每对 $(i, j) \in I^2$，一个 $X_i \times_S X_j$ 上的概形同构
$\varphi_{ij} : V_i \times_S X_j \to X_i \times_S V_j$
，使得对每个三元指标组 $(i, j, k) \in I^3$，图
$$
\xymatrix{
V_i \times_S X_j \times_S X_k
\ar[rd]^{\text{pr}_{01}^*\varphi_{ij}}
\ar[rr]_{\text{pr}_{02}^*\varphi_{ik}} &
&
X_i \times_S X_j \times_S V_k\\
&
X_i \times_S V_j \times_S X_k
\ar[ru]^{\text{pr}_{12}^*\varphi_{jk}}
}
$$
（作为 $X_i \times_S X_j \times_S X_k$ 上的概形）交换（记号含义显然）。
\item 下降数据之间的一个{\it 态射
$\psi : (V_i, \varphi_{ij}) \to (V'_i, \varphi'_{ij})$}由一族 $X_i$-概形
态射 $\psi_i : V_i \to V'_i$ 给出，记为
$\psi = (\psi_i)_{i \in I}$，并要求所有图
$$
\xymatrix{
V_i \times_S X_j \ar[r]_{\varphi_{ij}} \ar[d]_{\psi_i \times \text{id}} &
X_i \times_S V_j \ar[d]^{\text{id} \times \psi_j} \\
V'_i \times_S X_j \ar[r]^{\varphi'_{ij}} & X_i \times_S V'_j
}
$$
交换。
\end{enumerate}
\end{definition}

\noindent
例如，当询问相对曲线的纤维化范畴在 fpqc 拓扑中是否为叠时，这个概念会
自然出现（它不是——至少在坚持只使用概形时不是）。

\begin{remark}
\label{descent-remark-easier-family}
设 $S$ 为概形，设 $\{X_i \to S\}_{i \in I}$ 为以 $S$ 为目标的一族态射。
设 $(V_i, \varphi_{ij})$ 为相对于 $\{X_i \to S\}$ 的下降数据。可以把同构
$\varphi_{ij}$ 看作同构
$$
(X_i \times_S X_j) \times_{\text{pr}_0, X_i} V_i
\longrightarrow
(X_i \times_S X_j) \times_{\text{pr}_1, X_j} V_j
$$
（作为 $X_i \times_S X_j$ 上的概形）。因此粗略地说，可以把
$\varphi_{ij}$ 看作 $X_i \times_S X_j$ 上的同构
$\text{pr}_0^*V_i \to \text{pr}_1^*V_j$。此时上循环条件断言
$\text{pr}_{02}^*\varphi_{ik} =
\text{pr}_{12}^*\varphi_{jk} \circ \text{pr}_{01}^*\varphi_{ij}$.
从这个角度看，它与拟凝聚层上的下降数据情形十分相似。
\end{remark}

\noindent
我们通常使用仅由单个态射组成的族这一版本，理由是下面的引理。

\begin{lemma}
\label{descent-lemma-family-is-one}
设 $S$ 为概形，设 $\{X_i \to S\}_{i \in I}$ 为以 $S$ 为目标的一族态射。
置 $X = \coprod_{i \in I} X_i$，并将其视为 $S$-概形。存在范畴的典范等价
$$
\begin{matrix}
\text{下降数据范畴 } \\
\text{相对于该族 } \{X_i \to S\}_{i \in I}
\end{matrix}
\longrightarrow
\begin{matrix}
\text{ 下降数据范畴} \\
\text{ 相对于 } X/S
\end{matrix}
$$
，它把 $(V_i, \varphi_{ij})$ 映到 $(V, \varphi)$，其中
$V = \coprod_{i\in I} V_i$ 且 $\varphi = \coprod \varphi_{ij}$。
\end{lemma}

\begin{proof}
注意，$X \times_S X = \coprod_{ij} X_i \times_S X_j$，更高纤维积的情形
也类似。给出态射 $V \to X$ 恰好等同于给出一族态射 $V_i \to X_i$。
而给出下降数据 $\varphi$ 恰好等同于给出一族 $\varphi_{ij}$。
\end{proof}

\begin{lemma}
\label{descent-lemma-pullback}
概形上概形下降数据的拉回如下。
\begin{enumerate}
\item 设
$$
\xymatrix{
X' \ar[r]_f \ar[d]_{a'} & X \ar[d]^a \\
S' \ar[r]^h & S
}
$$
为概形态射的交换图。构造
$$
(V \to X, \varphi) \longmapsto f^*(V \to X, \varphi) = (V' \to X', \varphi')
$$
，其中 $V' = X' \times_X V$，且 $\varphi'$ 定义为复合
$$
\xymatrix{
V' \times_{S'} X' \ar@{=}[r] &
(X' \times_X V) \times_{S'} X' \ar@{=}[r] &
(X' \times_{S'} X') \times_{X \times_S X} (V \times_S X)
\ar[d]^{\text{id} \times \varphi} \\
X' \times_{S'} V' \ar@{=}[r] &
X' \times_{S'} (X' \times_X V) &
(X' \times_{S'} X') \times_{X \times_S X} (X \times_S V) \ar@{=}[l]
}
$$
，它定义了从相对于 $X \to S$ 的下降数据范畴到相对于 $X' \to S'$ 的
下降数据范畴的函子。
\item 给定两个使该图交换的态射 $f_i : X' \to X$，$i = 0, 1$，函子
$f_0^*$ 与 $f_1^*$ 典范同构。
\end{enumerate}
\end{lemma}

\begin{proof}
省略 (1) 的证明，但指出在 Remark \ref{descent-remark-easier} 引入的记号下，态射
$\varphi'$ 就是态射 $(f \times f)^*\varphi$。对 (2)，我们说明给出该函子
同构的是哪个态射 $f_0^*V \to f_1^*V$。事实上，由于 $f_0$ 和 $f_1$ 都
嵌入上述交换图，存在唯一的态射 $r : X' \to X \times_S X$，使得
$f_i = \text{pr}_i \circ r$。于是取
\begin{eqnarray*}
f_0^*V & = &
X' \times_{f_0, X} V \\
& = &
X' \times_{\text{pr}_0 \circ r, X} V \\
& = &
X' \times_{r, X \times_S X} (X \times_S X) \times_{\text{pr}_0, X} V \\
& \xrightarrow{\varphi} &
X' \times_{r, X \times_S X} (X \times_S X) \times_{\text{pr}_1, X} V \\
& = &
X' \times_{\text{pr}_1 \circ r, X} V \\
& = &
X' \times_{f_1, X} V \\
& = & f_1^*V
\end{eqnarray*}
省略对此构造确实有效的验证。
\end{proof}

\begin{definition}
\label{descent-definition-pullback-functor}
设 $S, S', X, X', f, a, a', h$ 如 Lemma \ref{descent-lemma-pullback} 中所述。该引理
构造的函子
$$
(V, \varphi) \longmapsto f^*(V, \varphi)
$$
称为下降数据上的{\it 拉回函子}。
\end{definition}

\begin{lemma}[概形族上概形下降数据的拉回]
\label{descent-lemma-pullback-family}
设 $\mathcal{U} = \{U_i \to S'\}_{i \in I}$ 和
$\mathcal{V} = \{V_j \to S\}_{j \in J}$ 为具有固定目标的态射族。设
$\alpha : I \to J$、$h : S' \to S$ 和 $g_i : U_i \to V_{\alpha(i)}$
构成具有固定目标的映射族之间的态射；见 Sites, Definition
\ref{sites-definition-morphism-coverings}。
\begin{enumerate}
\item 设 $(Y_j, \varphi_{jj'})$ 为相对于态射族
% P05-ERR-0050
$\{V_j \to S'\}$ 的下降数据。系统
$$
\left(
g_i^*Y_{\alpha(i)},
(g_i \times g_{i'})^*\varphi_{\alpha(i)\alpha(i')}
\right)
$$
（记号如 Remark \ref{descent-remark-easier-family}）
% P05-ERR-0051
是相对于 $\mathcal{V}$ 的下降数据。
\item 此构造定义了相对于 $\mathcal{U}$ 的下降数据与相对于
$\mathcal{V}$ 的下降数据之间的函子。
\item 再给定 $\alpha' : I \to J$、$h' : S' \to S$ 及具有固定目标的
映射族之间的态射 $g'_i : U_i \to V_{\alpha'(i)}$；若 $h = h'$，则所得到的
两个下降数据函子典范同构。
\item 经由 Lemma \ref{descent-lemma-family-is-one}，这些函子与 Lemma
\ref{descent-lemma-pullback} 中构造的拉回函子一致。
\end{enumerate}
\end{lemma}

\begin{proof}
这由 Lemma \ref{descent-lemma-pullback} 及 Lemma
\ref{descent-lemma-family-is-one} 的对应关系推出。
\end{proof}

\begin{definition}
\label{descent-definition-pullback-functor-family}
设 $\mathcal{U} = \{U_i \to S'\}_{i \in I}$、
$\mathcal{V} = \{V_j \to S\}_{j \in J}$、$\alpha : I \to J$、
$h : S' \to S$ 和 $g_i : U_i \to V_{\alpha(i)}$ 如 Lemma
\ref{descent-lemma-pullback-family} 中所述。该引理构造的函子
$$
(Y_j, \varphi_{jj'}) \longmapsto
(g_i^*Y_{\alpha(i)}, (g_i \times g_{i'})^*\varphi_{\alpha(i)\alpha(i')})
$$
称为下降数据上的{\it 拉回函子}。
\end{definition}

\noindent
若 $\mathcal{U}$ 和 $\mathcal{V}$ 具有同一目标 $S$，且 $\mathcal{U}$ 加细
$\mathcal{V}$（见 Sites, Definition
\ref{sites-definition-morphism-coverings}），但没有给出明确的偶
$(\alpha, g_i)$，我们仍可谈论拉回函子，因为 Lemma
\ref{descent-lemma-pullback-family} 已表明该偶的选择无关紧要（相差典范同构）。


\begin{definition}
\label{descent-definition-effective}
设 $S$ 为概形，设 $f : X \to S$ 为概形态射。
\begin{enumerate}
\item 给定 $S$ 上的概形 $U$，有 $U$ 相对于 $\text{id} : S \to S$ 的
{\it 平凡下降数据}，即 $U$ 上的恒等态射。
\item 由 Lemma \ref{descent-lemma-pullback}，沿 $f$ 拉回平凡下降数据，可在
$X \times_S U$ 上得到一个相对于 $X \to S$ 的{\it 典范下降数据}。常把
这个下降数据记为 $(X \times_S U, can)$。
\item 若相对于 $X/S$ 的下降数据 $(V, \varphi)$ 与某个 $S$ 上概形 $U$
所对应的典范下降数据 $(X \times_S U, can)$ 同构，则称 $(V, \varphi)$
{\it 有效}。
\end{enumerate}
\end{definition}

\noindent
因此，有效意味着存在 $S$ 上的概形 $U$ 及 $X$-概形同构
$\psi : V \to X \times_S U$，使得 $\varphi$ 等于复合
$$
V \times_S X \xrightarrow{\psi \times \text{id}_X}
X \times_S U \times_S X =
X \times_S X \times_S U
\xrightarrow{\text{id}_X \times \psi^{-1}}
X \times_S V
$$

\begin{definition}
\label{descent-definition-effective-family}
设 $S$ 为概形，设 $\{X_i \to S\}$ 为以 $S$ 为目标的一族态射。
\begin{enumerate}
\item 给定 $S$ 上的概形 $U$，拉回 $U$ 相对于
$\{\text{id} : S \to S\}$ 的平凡下降数据，即在概形族
$X_i \times_S U$ 上得到一个{\it 典范下降数据}。将其记为
$(X_i \times_S U, can)$。
\item 若存在 $S$ 上的概形 $U$，使相对于 $\{X_i \to S\}$ 的下降数据
$(V_i, \varphi_{ij})$ 与 $(X_i \times_S U, can)$ 同构，则称
$(V_i, \varphi_{ij})$ 为{\it 有效}。
\end{enumerate}
\end{definition}












\section{拉回函子的全忠实性}
\label{descent-section-fully-faithful}

\noindent
事实证明，fpqc 覆盖的下降数据之间的拉回函子是全忠实的。换言之，概形
态射满足 fpqc 下降。本节的目标是证明这一点；不过我们鼓励读者自行证明。
关键是使用 Lemma \ref{descent-lemma-fpqc-universal-effective-epimorphisms}。

\begin{lemma}
\label{descent-lemma-surjective-flat-epi}
满且平坦的态射在概形范畴中是上态射。
\end{lemma}

\begin{proof}
假设 $h : X' \to X$ 满且平坦，并有态射 $a, b : X \to Y$，满足
$a \circ h = b \circ h$。由于 $h$ 满，$a$ 与 $b$ 在底层拓扑空间上相同。
选取 $x' \in X'$，置 $x = h(x')$ 和 $y = a(x) = b(x)$。考虑局部环映射
$$
a^\sharp_x, b^\sharp_x : \mathcal{O}_{Y, y} \to \mathcal{O}_{X, x}
$$
。它们与平坦局部同态
$h^\sharp_{x'} : \mathcal{O}_{X, x} \to \mathcal{O}_{X', x'}$ 复合后相等。
平坦局部同态忠实平坦（Algebra, Lemma
\ref{algebra-lemma-local-flat-ff}），所以 $h^\sharp_{x'}$ 单射。因而
$a^\sharp_x = b^\sharp_x$，这蕴含 $a = b$，正如所需。
\end{proof}

\begin{lemma}
\label{descent-lemma-ff-base-change-faithful}
设 $h : S' \to S$ 为满平坦概形态射。基变换函子
$$
\Sch/S \longrightarrow \Sch/S', \quad
X \longmapsto S' \times_S X
$$
是忠实的。
\end{lemma}

\begin{proof}
设 $X_1$、$X_2$ 为 $S$ 上的概形，设
$\alpha, \beta : X_2 \to X_1$ 为 $S$ 上的态射。若 $\alpha$、$\beta$
基变换为同一态射，则得到如下交换图：
$$
\xymatrix{
X_2 \ar[d]^\alpha &
S' \times_S X_2 \ar[l] \ar[d] \ar[r] &
X_2 \ar[d]^\beta \\
X_1 &
S' \times_S X_1 \ar[l] \ar[r] &
X_1
}
$$
所以只需证明 $S' \times_S X_2 \to X_2$ 为上态射。作为满平坦态射的
基变换，它满且平坦（见 Morphisms, Lemmas
\ref{morphisms-lemma-base-change-surjective} 和
\ref{morphisms-lemma-base-change-flat}）。因此本引理由 Lemma
\ref{descent-lemma-surjective-flat-epi} 推出。
\end{proof}

\begin{lemma}
\label{descent-lemma-faithful}
在 Lemma \ref{descent-lemma-pullback} 的情形下，假设 $f : X' \to X$ 满且平坦。
则拉回函子忠实。
\end{lemma}

\begin{proof}
设 $(V_i, \varphi_i)$，$i = 1, 2$，为 $X \to S$ 的下降数据。设
$\alpha, \beta : V_1 \to V_2$ 为下降数据的态射。假设
$f^*\alpha = f^*\beta$。我们的任务是证明 $\alpha = \beta$。注意，
$\alpha$、$\beta$ 是 $X$ 上概形的态射，而 $f^*\alpha$、$f^*\beta$ 只是
$\alpha$、$\beta$ 到 $X'$ 上态射的基变换。因此本引理由 Lemma
\ref{descent-lemma-ff-base-change-faithful} 推出。
\end{proof}

\noindent
下面是本节的关键引理。

\begin{lemma}
\label{descent-lemma-fully-faithful}
在 Lemma \ref{descent-lemma-pullback} 的情形下，假设
\begin{enumerate}
\item $\{f : X' \to X\}$ 是 fpqc 覆盖（例如，若 $f$ 满、平坦且拟紧）；
\item $S = S'$.
\end{enumerate}
则拉回函子全忠实。
\end{lemma}

\begin{proof}
假设 (1) 蕴含 $f$ 满且平坦。因此由 Lemma \ref{descent-lemma-faithful}，拉回函子
忠实。设 $(V, \varphi)$ 和 $(W, \psi)$ 为相对于 $X \to S$ 的两个下降数据。
置 $(V', \varphi') = f^*(V, \varphi)$ 及
$(W', \psi') = f^*(W, \psi)$。设 $\alpha' : V' \to W'$ 为 $S$ 上 $X'$
的下降数据态射。我们必须证明存在 $S$ 上 $X$ 的下降数据态射
$\alpha : V \to W$，其拉回为 $\alpha'$。

\medskip\noindent
回忆 $V'$ 是 $V$ 沿 $f$ 的基变换，而 $\varphi'$ 是 $\varphi$ 沿
$f \times f$ 的基变换（见 Remark \ref{descent-remark-easier}）。由假设，图
$$
\xymatrix{
V' \times_S X' \ar[r]_{\varphi'} \ar[d]_{\alpha' \times \text{id}} &
X' \times_S V' \ar[d]^{\text{id} \times \alpha'} \\
W' \times_S X' \ar[r]^{\psi'} &
X' \times_S W'
}
$$
交换。我们断言两个复合
$$
\xymatrix{
V' \times_V V' \ar[r]^-{\text{pr}_i} &
V' \ar[r]^{\alpha'} &
W' \ar[r] &
W
}
, \quad i = 0, 1
$$
相等。建议读者自行证明，而不要阅读本段余下部分。（若你找到漂亮简洁的
论证，请发邮件。）设 $v_0, v_1$ 为 $V'$ 中映到同一点 $v \in V$ 的点。
设 $x_i \in X'$ 为 $v_i$ 的像，并设 $x$ 为 $X$ 中的点，即 $v$ 在 $X$ 中的像。换言之，
在 $V' = X' \times_X V$ 中，$v_i = (x_i, v)$。对 $V$ 的某个点 $v'$，写
$\varphi(v, x) = (x, v')$。这是可能的，因为 $\varphi$ 是
$X \times_S X$ 上的态射。记 $v_i' = (x_i, v')$，它是 $V'$ 的点。计算
（使用 $\varphi'$ 的定义）表明 $\varphi'(v_i, x_j) = (x_i, v'_j)$。记
$w_i = \alpha'(v_i)$ 及 $w'_i = \alpha'(v_i')$。现在可以把
$w_i = (x_i, u_i)$ 写成这种形式，其中 $u_i$ 是 $W$ 的某个点；类似地，
$w_i' = (x_i, u'_i)$，其中 $u_i'$ 是 $W$ 的某个点。该断言等价于
$u_0 = u_1$。使用 $\psi'$ 的定义作形式计算（见 Lemma
\ref{descent-lemma-pullback}），可知上述图的交换性意味着
$$
((x_i, x_j), \psi(u_i, x)) = ((x_i, x_j), (x, u'_j))
$$
（作为
$(X' \times_S X') \times_{X \times_S X} (X \times_S W)$
的点）对所有 $i, j \in \{0, 1\}$ 成立。这表明
$\psi(u_0, x) = \psi(u_1, x)$，取 $\psi^{-1}$ 即得 $u_0 = u_1$。上述论证
是形式的，并且可以取概形点（换言之，可以取
$(v_0, v_1) = \text{id}_{V' \times_V V'}$），故这证明了断言。

\medskip\noindent
此时可以使用 Lemma \ref{descent-lemma-fpqc-universal-effective-epimorphisms}。
事实上，$\{V' \to V\}$ 是态射 $f : X' \to X$ 的基变换，因而为 fpqc
覆盖。所以由 Lemma \ref{descent-lemma-fpqc-universal-effective-epimorphisms}，态射
$\alpha' : V' \to W' \to W$ 唯一地经过一个态射 $\alpha : V \to W$
分解，而后者的基变换必为 $\alpha'$。最后，图
$$
\xymatrix{
V \times_S X \ar[r]_{\varphi} \ar[d]_{\alpha \times \text{id}} &
X \times_S V \ar[d]^{\text{id} \times \alpha} \\
W \times_S X \ar[r]^{\psi} & X \times_S W
}
$$
交换，因为它到 $X' \times_S X'$ 的基变换交换，并且态射
$X' \times_S X' \to X \times_S X$ 满且平坦（使用 Lemma
\ref{descent-lemma-ff-base-change-faithful}）。所以 $\alpha$ 是下降数据态射
$(V, \varphi) \to (W, \psi)$，正如所需。
\end{proof}

\noindent
前面材料经过改进后的叙述已经使下面两个引理过时。但它们仍然正确！

\begin{lemma}
\label{descent-lemma-pullback-selfmap}
设 $X \to S$ 为概形态射，设 $f : X \to X$ 为 $X$ 在 $S$ 上的自映射。
此时，沿 $f$ 的拉回同构于相对于 $X \to S$ 的下降数据范畴上的恒等函子。
\end{lemma}

\begin{proof}
这由 Lemma \ref{descent-lemma-pullback} 显然可见，因为它给出
$f^* \cong \text{id}^*$。
\end{proof}

\begin{lemma}
\label{descent-lemma-morphism-with-section-equivalence}
设 $f : X' \to X$ 为基概形 $S$ 上的概形态射。假设存在 $S$ 上的态射
$g : X \to X'$；例如，$f$ 有截面时即如此。则 Lemma
\ref{descent-lemma-pullback} 的拉回函子定义了相对于 $X/S$ 与相对于 $X'/S$ 的
下降数据范畴之间的等价。
\end{lemma}

\begin{proof}
设 $g : X \to X'$ 为 $S$ 上的态射。上述 Lemma
\ref{descent-lemma-pullback-selfmap} 表明，函子
$f^* \circ g^* = (g \circ f)^*$ 和 $g^* \circ f^* = (f \circ g)^*$
分别同构于相应的恒等函子，正如所需。
\end{proof}

\begin{lemma}
\label{descent-lemma-morphism-source-faithfully-flat}
设 $f : X \to X'$ 为基概形 $S$ 上的概形态射。假设 $X \to S$ 满且平坦。
则 Lemma \ref{descent-lemma-pullback} 的拉回函子，是从相对于 $X'/S$ 的下降数据
范畴到相对于 $X/S$ 的下降数据范畴的忠实函子。
\end{lemma}

\begin{proof}
可以把 $X \to X'$ 分解为 $X \to X \times_S X' \to X'$。第一个态射有
截面，所以由 Lemma \ref{descent-lemma-morphism-with-section-equivalence}，它诱导
下降数据范畴的等价。第二个态射满且平坦，所以由 Lemma
\ref{descent-lemma-faithful}，它诱导忠实函子。
\end{proof}

\begin{lemma}
\label{descent-lemma-morphism-source-fpqc-covering}
设 $f : X \to X'$ 为基概形 $S$ 上的概形态射。假设 $\{X \to S\}$ 为 fpqc
覆盖（例如，若 $f$ 满、平坦且拟紧）。则 Lemma \ref{descent-lemma-pullback} 的
拉回函子，是从相对于 $X'/S$ 的下降数据范畴到相对于 $X/S$ 的下降数据
范畴的全忠实函子。
\end{lemma}

\begin{proof}
可以把 $X \to X'$ 分解为 $X \to X \times_S X' \to X'$。第一个态射有
截面，所以由 Lemma \ref{descent-lemma-morphism-with-section-equivalence}，它诱导
下降数据范畴的等价。第二个态射是 fpqc 覆盖，所以由 Lemma
\ref{descent-lemma-fully-faithful}，它诱导全忠实函子。
\end{proof}

\begin{lemma}
\label{descent-lemma-fpqc-refinement-coverings-fully-faithful}
设 $S$ 为概形。设 $\mathcal{U} = \{U_i \to S\}_{i \in I}$ 和
$\mathcal{V} = \{V_j \to S\}_{j \in J}$ 为以 $S$ 为目标的态射族。设
$\alpha : I \to J$、$\text{id} : S \to S$ 和
$g_i : U_i \to V_{\alpha(i)}$ 构成具有固定目标的映射族之间的态射；见
Sites, Definition \ref{sites-definition-morphism-coverings}。假设对每个
$j \in J$，态射族 $\{g_i : U_i \to V_j\}_{\alpha(i) = j}$ 是 $V_j$ 的
fpqc 覆盖。则 Lemma \ref{descent-lemma-pullback-family} 的拉回函子
$$
\text{相对下降数据 }
\mathcal{V}
\longrightarrow
\text{相对下降数据 }
\mathcal{U}
$$
是全忠实的。
\end{lemma}

\begin{proof}
考虑 $S$ 上的概形态射
$$
g :
X = \coprod\nolimits_{i \in I} U_i
\longrightarrow
Y = \coprod\nolimits_{j \in J} V_j
$$
；它在第 $i$ 个分支上经由态射
% P05-ERR-0053
$g_{\alpha(i)}$ 映入第 $\alpha(i)$ 个分支。我们断言
$\{g : X \to Y\}$ 是概形的 fpqc 覆盖。事实上，由 Topologies, Lemma
\ref{topologies-lemma-disjoint-union-is-fpqc-covering}，对每个 $j$，态射
$\{\coprod_{\alpha(i) = j} U_i \to V_j\}$ 是 fpqc 覆盖。因此，对每个仿射
开集 $V \subset V_j$（可把它视为 $Y$ 的仿射开集），可以找到有限多个
仿射开集 $W_1, \ldots, W_n \subset \coprod_{\alpha(i) = j} U_i$（可把它们
视为 $X$ 的仿射开集），使得
$V = \bigcup_{i = 1, \ldots, n} g(W_i)$。这给出了足够多的 $Y$ 的仿射开集，
它们可由有限多个 $X$ 的仿射开集覆盖，因而 Topologies, Lemma
\ref{topologies-lemma-recognize-fpqc-covering} 的第 (3) 部分适用，断言成立。
分别以 $DD(X/S)$ 和 $DD(\mathcal{U})$ 表示相对于 $X/S$ 和相对于
$\mathcal{U}$ 的下降数据范畴；对 $Y/S$ 和 $\mathcal{V}$ 也采用类似记号。
考虑图
$$
\xymatrix{
DD(Y/S) \ar[r] & DD(X/S) \\
DD(\mathcal{V}) \ar[u]^{\text{引理 }\ref{descent-lemma-family-is-one}} \ar[r] &
DD(\mathcal{U}) \ar[u]_{\text{引理 }\ref{descent-lemma-family-is-one}}
}
$$
该图交换；见 Lemma \ref{descent-lemma-pullback-family} 的证明。竖直箭头为等价。
因此本引理由 Lemma \ref{descent-lemma-fully-faithful} 推出，后者表明图中上方水平
箭头全忠实。
\end{proof}

\noindent
下一个引理表明，为了检验有效性，总可以对给定的以 $S$ 为目标的态射族作
Zariski 加细。

\begin{lemma}
\label{descent-lemma-Zariski-refinement-coverings-equivalence}
设 $S$ 为概形。设 $\mathcal{U} = \{U_i \to S\}_{i \in I}$ 和
$\mathcal{V} = \{V_j \to S\}_{j \in J}$ 为以 $S$ 为目标的态射族。设
$\alpha : I \to J$、$\text{id} : S \to S$ 和
$g_i : U_i \to V_{\alpha(i)}$ 构成具有固定目标的映射族之间的态射；见
Sites, Definition \ref{sites-definition-morphism-coverings}。假设对每个
$j \in J$，态射族 $\{g_i : U_i \to V_j\}_{\alpha(i) = j}$ 是 $V_j$ 的
Zariski 覆盖（见 Topologies, Definition
\ref{topologies-definition-zariski-covering}）。则 Lemma
\ref{descent-lemma-pullback-family} 的拉回函子
$$
\text{相对下降数据 }
\mathcal{V}
\longrightarrow
\text{相对下降数据 }
\mathcal{U}
$$
是范畴等价。特别地，$S$ 上概形的范畴等价于相对于 $S$ 的任意 Zariski
覆盖的下降数据范畴。
\end{lemma}

\begin{proof}
由 Lemma \ref{descent-lemma-fpqc-refinement-coverings-fully-faithful}，该函子忠实且
全忠实。下面说明如何证明它本质满。设 $(X_i, \varphi_{ii'})$ 为相对于
$\mathcal{U}$ 的下降数据。固定 $j \in J$，并置
$I_j = \{i \in I \mid \alpha(i) = j\}$。对 $i, i' \in I_j$，注意存在
典范态射
$$
c_{ii'} : U_i \times_{g_i, V_j, g_{i'}} U_{i'} \to U_i \times_S U_{i'}.
$$
。因此可沿此态射拉回 $\varphi_{ii'}$；对 $i, i' \in I_j$，置
$\psi_{ii'} = c_{ii'}^*\varphi_{ii'}$。这样得到相对于 Zariski 覆盖
$\{g_i : U_i \to V_j\}_{i \in I_j}$ 的下降数据
$(X_i, \psi_{ii'})$。注意，$\psi_{ii'}$ 是从 $X_i$ 的开集
$X_{i, U_i \times_{V_j} U_{i'}}$ 到 $X_{i'}$ 的相应开集的同构。由 Schemes,
Section \ref{schemes-section-glueing-schemes}，可以把
$(X_i, \psi_{ii'})$ 黏合成 $V_j$ 上的概形 $Y_j$。此外，对 $i \in I_j$ 和
$i' \in I_{j'}$，态射 $\varphi_{ii'}$ 黏合成满足上循环条件的态射
$\varphi_{jj'} : Y_j \times_S V_{j'} \to V_j \times_S Y_{j'}$（省略细节）。
因此得到所需的相对于 $\mathcal{V}$ 的下降数据
$(Y_j, \varphi_{jj'})$。
\end{proof}

\begin{lemma}
\label{descent-lemma-refine-coverings-fully-faithful}
设 $S$ 为概形。设 $\mathcal{U} = \{U_i \to S\}_{i \in I}$ 和
$\mathcal{V} = \{V_j \to S\}_{j \in J}$ 为 $S$ 的 fpqc 覆盖。若
$\mathcal{U}$ 是 $\mathcal{V}$ 的加细，则拉回函子
$$
\text{相对下降数据 }
\mathcal{V}
\longrightarrow
\text{相对下降数据 }
\mathcal{U}
$$
全忠实。特别地，$S$ 上概形的范畴可认同为相对于 $S$ 的任意 fpqc 覆盖的
下降数据范畴的一个全子范畴。
\end{lemma}

\begin{proof}
考虑 $S$ 的 fpqc 覆盖
$\mathcal{W} = \{U_i \times_S V_j \to S\}_{(i, j) \in I \times J}$。它既是
$\mathcal{U}$ 的加细，也是 $\mathcal{V}$ 的加细。因此有函子与范畴的
$2$-交换图
$$
\xymatrix{
DD(\mathcal{V}) \ar[rd] \ar[rr] & & DD(\mathcal{U}) \ar[ld] \\
& DD(\mathcal{W}) &
}
$$
记号如 Lemma \ref{descent-lemma-fpqc-refinement-coverings-fully-faithful} 的证明；交换
性由 Lemma \ref{descent-lemma-pullback-family} 第 (3) 部分给出。因此显然只需证明
函子 $DD(\mathcal{V}) \to DD(\mathcal{W})$ 和
$DD(\mathcal{U}) \to DD(\mathcal{W})$ 全忠实。这由 Lemma
\ref{descent-lemma-fpqc-refinement-coverings-fully-faithful} 推出，正如所需。
\end{proof}

\begin{remark}
\label{descent-remark-morphisms-of-schemes-satisfy-fpqc-descent}
Lemma \ref{descent-lemma-refine-coverings-fully-faithful} 断言概形态射满足 fpqc
下降。换言之，给定概形 $S$ 及 $S$ 上的概形 $X$、$Y$，函子
$$
(\Sch/S)^{opp} \longrightarrow \textit{Sets},
\quad
T \longmapsto \Mor_T(X_T, Y_T)
$$
满足 fpqc 拓扑的层条件。最简单的情形如下。假设 $T \to S$ 是仿射概形
之间的满平坦态射。设 $\psi_0 : X_T \to Y_T$ 为 $T$ 上与典范下降数据相容
的概形态射。则存在唯一的态射 $\psi : X \to Y$，其到 $T$ 的基变换为
$\psi_0$。事实上，这个特殊情形直接由 Lemma
\ref{descent-lemma-fully-faithful} 推出。而该引理又是下列两个事实的形式推论：
(a) 沿忠实平坦态射的基变换函子忠实，见 Lemma
\ref{descent-lemma-ff-base-change-faithful}；(b) 概形满足 fpqc 拓扑的层条件，见
Lemma \ref{descent-lemma-fpqc-universal-effective-epimorphisms}。
\end{remark}

\begin{lemma}
\label{descent-lemma-effective-for-fpqc-is-local-upstairs}
设 $X \to S$ 为满、拟紧、平坦的概形态射。设 $(V, \varphi)$ 为相对于
$X/S$ 的下降数据。假设对所有 $v \in V$，都存在开子概形
$v \in W \subset V$，使得 $\varphi(W \times_S X) \subset X \times_S W$，
且下降数据 $(W, \varphi|_{W \times_S X})$ 有效。则 $(V, \varphi)$ 有效。
\end{lemma}

\begin{proof}
设 $V = \bigcup W_i$ 为开覆盖，满足
$\varphi(W_i \times_S X) \subset X \times_S W_i$，且下降数据
$(W_i, \varphi|_{W_i \times_S X})$ 有效。设 $U_i \to S$ 为概形，并设
$\alpha_i : (X \times_S U_i, can) \to (W_i, \varphi|_{W_i \times_S X})$
为下降数据同构。对每对指标 $(i, j)$，考虑开集
$\alpha_i^{-1}(W_i \cap W_j) \subset X \times_S U_i$。由于所有内容都与
下降数据相容，且 $\{X \to S\}$ 为 fpqc 覆盖，可以应用 Lemma
\ref{descent-lemma-open-fpqc-covering}，找到开集 $U_{ij} \subset U_i$，使得
$\alpha_i^{-1}(W_i \cap W_j) = X \times_S U_{ij}$。现在，
$W_i \cap W_j$ 上的恒等态射与下降数据相容，因而来自唯一的 $S$ 上态射
$\varphi_{ij} : U_{ij} \to U_{ji}$（见 Remark
\ref{descent-remark-morphisms-of-schemes-satisfy-fpqc-descent}）。于是
$(U_i, U_{ij}, \varphi_{ij})$ 是 Schemes, Section
\ref{schemes-section-glueing-schemes} 中的黏合数据（省略证明）。因此可以
假设存在 $S$ 上的概形 $U$，使得 $U_i \subset U$ 为开集、
$U_{ij} = U_i \cap U_j$ 且
$\varphi_{ij} = \text{id}_{U_i \cap U_j}$；见 Schemes, Lemma
\ref{schemes-lemma-glue}。拉回到 $X$ 后，可以利用 $\alpha_i$ 得到所需
同构 $\alpha : X \times_S U \to V$。
\end{proof}







\section{态射类型的下降}
\label{descent-section-descending-types-morphisms}

\noindent
下面研究概形相对于 fpqc 覆盖的下降数据是否有效。首先要指出，这并非总是
成立；我们将在 Algebraic Spaces, Example
\ref{spaces-example-non-representable-descent} 中看到这一点。甚至投射态射也
不总是对 fpqc 覆盖满足下降，见 Examples, Lemma
\ref{examples-lemma-non-effective-descent-projective}。

\medskip\noindent
另一方面，如果试图下降的概形特别简单，那么对整个一类概形，下降数据有时
确实有效。这里引入抽象描述这一现象的术语，尽管以后若使用不当，它可能造成
混淆。

\begin{definition}
\label{descent-definition-descending-types-morphisms}
设 $\mathcal{P}$ 为基底上概形态射的一条性质，并设\linebreak[4]
$\tau \in \{Zariski, fpqc, fppf, \etale, smooth, syntomic\}$。如果对任意
$\tau$-覆盖 $\mathcal{U} : \{U_i \to S\}_{i \in I}$（见 Topologies,
Section \ref{topologies-section-procedure}），每个相对于 $\mathcal{U}$ 的
下降数据 $(X_i, \varphi_{ij})$，只要每个态射 $X_i \to U_i$ 具有性质
$\mathcal{P}$，就都是有效的，则称{\it $\mathcal{P}$ 型态射对
$\tau$-覆盖满足下降}。
\end{definition}

\noindent
注意，在每种情形下我们已经看到，从 $S$ 上概形到 $\mathcal{U}$ 上下降
数据的函子都是全忠实的（由 Lemma
\ref{descent-lemma-refine-coverings-fully-faithful} 以及 Topologies 中任意
$\tau$-覆盖也是 fpqc 覆盖的结果可知）。我们还看到，相对于 Zariski 覆盖，
下降数据总是有效的（Lemma
\ref{descent-lemma-Zariski-refinement-coverings-equivalence}）。为避免错误论证
（即在下降进行到一半时依赖 $\mathcal{P}$），谨慎的做法是：仅当
$\mathcal{P}$ 在任意基变换下稳定，或至少在 $\tau$-拓扑中对基底为局部性质
时，才研究刚引入的概念（见 Definition
\ref{descent-definition-property-morphisms-local}）。

\medskip\noindent
下面这个必备引理把问题约化到由仿射概形之间单个态射给出的覆盖。

\begin{lemma}
\label{descent-lemma-descending-types-morphisms}
设 $\mathcal{P}$ 为基底上概形态射的一条性质，并设
$\tau \in \{fpqc, fppf, \etale, smooth, syntomic\}$。假设：
\begin{enumerate}
\item $\mathcal{P}$ 在任意基变换下稳定（见 Schemes, Definition
\ref{schemes-definition-preserved-by-base-change}）；
\item 如果 $Y_j \to V_j$（$j = 1, \ldots, m$）具有 $\mathcal{P}$，则
$\coprod Y_j \to \coprod V_j$ 也具有该性质；并且
\item 对任意仿射概形之间的满射态射 $X \to S$，按照 $\tau$ 分别为 fpqc、
fppf、\'etale、smooth 或 syntomic，假设该态射分别是平坦、有限表现平坦、
\'etale、光滑或 syntomic 的；则相对于 $S$ 上 $X$ 的任意下降数据
$(V, \varphi)$，只要 $V \to X$ 具有 $\mathcal{P}$，就都是有效的。
\end{enumerate}
那么 $\mathcal{P}$ 型态射对 $\tau$-覆盖满足下降。
\end{lemma}

\begin{proof}
设 $S$ 为概形，设
$\mathcal{U} = \{\varphi_i : U_i \to S\}_{i \in I}$ 为 $S$ 的
$\tau$-覆盖。设 $(X_i, \varphi_{ii'})$ 为相对于 $\mathcal{U}$ 的下降
数据，并假设每个态射 $X_i \to U_i$ 具有性质 $\mathcal{P}$。我们必须证明
存在概形 $X \to S$，使得
$(X_i, \varphi_{ii'}) \cong (U_i \times_S X, can)$。

\medskip\noindent
在正式开始证明前先指出：对任意态射族
$\mathcal{V} : \{V_j \to S\}$ 及任意族态射
$\mathcal{V} \to \mathcal{U}$，如果把下降数据
$(X_i, \varphi_{ii'})$ 拉回成 $\mathcal{V}$ 上的下降数据
$(Y_j, \varphi_{jj'})$，那么每个态射 $Y_j \to V_j$ 也具有性质
$\mathcal{P}$。这是因为假设 (1) 说 $\mathcal{P}$ 在任意基变换下稳定，
并结合拉回的定义（见 Definition
\ref{descent-definition-pullback-functor-family}）。下文将不再说明地使用这一点。

\medskip\noindent
先在 $S$ 仿射时证明本引理。由 Topologies, Lemma
\ref{topologies-lemma-fpqc-affine},
\ref{topologies-lemma-fppf-affine},
\ref{topologies-lemma-etale-affine},
\ref{topologies-lemma-smooth-affine}, or
\ref{topologies-lemma-syntomic-affine}
，存在加细 $\mathcal{U}$ 的标准 $\tau$-覆盖
$\mathcal{V} : \{V_j \to S\}_{j = 1, \ldots, m}$。下降数据范畴之间的拉回
函子
$DD(\mathcal{U}) \to DD(\mathcal{V})$
由 Lemma \ref{descent-lemma-refine-coverings-fully-faithful} 可知是全忠实的。因此只需
证明标准 $\tau$-覆盖 $\mathcal{V}$ 上的下降数据有效。由假设 (2)，
$\coprod Y_j \to \coprod V_j$ 具有性质 $\mathcal{P}$。由 Lemma
\ref{descent-lemma-family-is-one}，问题约化到覆盖
$\{\coprod_{j = 1, \ldots, m} V_j \to S\}$；对此，引理的假设 (3) 已经给出
所需结果。因此，当 $S$ 仿射时引理成立。

\medskip\noindent
现设 $S$ 为一般概形，并设 $V \subset S$ 为仿射开集。由位点的性质，族
$\mathcal{U}_V = \{V \times_S U_i \to V\}_{i \in I}$ 是 $V$ 的
$\tau$-覆盖。用 $(X_i, \varphi_{ii'})_V$ 表示给定下降数据在
$\mathcal{U}_V$ 上的限制（或拉回）。由刚才的结论，得到 $V$ 上的概形
$X_V$，其相对于 $\mathcal{U}_V$ 的典范下降数据与
$(X_i, \varphi_{ii'})_V$ 同构。假设 $V' \subset V$ 是 $V$ 的仿射开集。
那么 $X_{V'}$ 与 $V' \times_V X_V$ 的典范下降数据都同构于
$(X_i, \varphi_{ii'})_{V'}$。再次由 Lemma
\ref{descent-lemma-refine-coverings-fully-faithful}，得到 $S$ 上的典范态射
$\rho^V_{V'} : X_{V'} \to X_V$，它把 $X_{V'}$ 等同于 $V'$ 在 $X_V$
中的逆像。略去如下验证：对 $S$ 的仿射开集
$V'' \subset V' \subset V$，有
$\rho^V_{V''} = \rho^V_{V'} \circ \rho^{V'}_{V''}$。

\medskip\noindent
由 Constructions, Lemma \ref{constructions-lemma-relative-glueing}，数据
$(X_V, \rho^V_{V'})$ 拼接成概形 $X \to S$。此外，我们有恢复态射
$\rho^V_{V'}$ 的同构 $V \times_S X \to X_V$。展开概形 $X_V$ 的构造，
得到同构
$$
V \times_S U_i \times_S X
\longrightarrow
V \times_S X_i
$$
，它们与态射 $\varphi_{ii'}$ 相容，也与限制到 $X$ 中更小的仿射开集相容。
这说明 $U_i \times_S X$ 上的典范下降数据同构于给定下降数据，证毕。
\end{proof}










\section{仿射态射的下降}
\label{descent-section-affine}

\noindent
本节证明“仿射态射对 fpqc 覆盖满足下降”。正式陈述如下。

\begin{lemma}
\label{descent-lemma-affine}
设 $S$ 为概形。设 $\{X_i \to S\}_{i\in I}$ 为 fpqc 覆盖，见
Topologies, Definition \ref{topologies-definition-fpqc-covering}。设
$(V_i/X_i, \varphi_{ij})$ 为相对于 $\{X_i \to S\}$ 的下降数据。如果每个
态射 $V_i \to X_i$ 都是仿射的，那么该下降数据有效。
\end{lemma}

\begin{proof}
仿射性是概形态射的一条性质，它在基底上为局部性质并在任意基变换下保持，
见 Morphisms, Lemmas \ref{morphisms-lemma-characterize-affine} 和
\ref{morphisms-lemma-base-change-affine}。因此可应用 Lemma
\ref{descent-lemma-descending-types-morphisms}，只需在 fpqc 覆盖由仿射概形之间的
单个平坦满射 $\{X \to S\}$ 给出时证明本引理。记
$X = \Spec(A)$ 及 $S = \Spec(R)$，于是 $R \to A$ 是忠实平坦环映射。设
$(V, \varphi)$ 为相对于 $S$ 上 $X$ 的下降数据，并假设 $V \to X$ 仿射。
由于 $V \to X$ 仿射，有 $V = \Spec(B)$，其中 $B$ 是某个 $A$-代数
（见 Morphisms, Definition \ref{morphisms-definition-affine}）。同构
$\varphi$ 对应于环同构
$$
\varphi^\sharp :
B \otimes_R A \longleftarrow A \otimes_R B
$$
，它是 $A \otimes_R A$-代数同构。$\varphi$ 的余圈条件说明
$$
\xymatrix{
B \otimes_R A \otimes_R A & &
A \otimes_R A \otimes_R B \ar[ll] \ar[ld]\\
& A \otimes_R B \otimes_R A \ar[lu] &
}
$$
交换。将这些箭头反向，便得到 Definition
\ref{descent-definition-descent-datum-modules} 意义下相对于 $R \to A$ 的模下降
数据。因此可应用 Proposition \ref{descent-proposition-descent-module}，得到
$R$-模
$C = \Ker(B \to A \otimes_R B)$
以及保持下降数据的同构 $A \otimes_R C \cong B$。对任意
$c, c' \in C$，因为态射 $\varphi$ 是代数同态，$B$ 中的乘积 $cc'$ 属于
$C$。因此 $C$ 是 $R$-代数，其到 $A$ 的基变换以与下降数据相容的方式
同构于 $B$。应用 $\Spec$，得到 $S$ 上的概形 $U$，使
$(V, \varphi) \cong (X \times_S U, can)$，正合所需。
\end{proof}

\begin{lemma}
\label{descent-lemma-closed-immersion}
设 $S$ 为概形。设 $\{X_i \to S\}_{i\in I}$ 为 fpqc 覆盖，见
Topologies, Definition \ref{topologies-definition-fpqc-covering}。设
$(V_i/X_i, \varphi_{ij})$ 为相对于 $\{X_i \to S\}$ 的下降数据。如果每个
态射 $V_i \to X_i$ 都是闭浸入，那么该下降数据有效。
\end{lemma}

\begin{proof}
这是因为闭浸入是仿射态射（Morphisms, Lemma
\ref{morphisms-lemma-closed-immersion-affine}），因而可应用 Lemma
\ref{descent-lemma-affine}。
\end{proof}


\section{拟仿射态射的下降}
\label{descent-section-quasi-affine}

\noindent
本节证明“拟仿射态射对 fpqc 覆盖满足下降”。正式陈述如下。

\begin{lemma}
\label{descent-lemma-quasi-affine}
设 $S$ 为概形。设 $\{X_i \to S\}_{i\in I}$ 为 fpqc 覆盖，见
Topologies, Definition \ref{topologies-definition-fpqc-covering}。设
$(V_i/X_i, \varphi_{ij})$ 为相对于 $\{X_i \to S\}$ 的下降数据。如果每个
态射 $V_i \to X_i$ 都是拟仿射的，那么该下降数据有效。
\end{lemma}

\begin{proof}
拟仿射性是概形态射的一条性质，它在任意基变换下保持，见 Morphisms,
Lemmas \ref{morphisms-lemma-characterize-quasi-affine} 和
\ref{morphisms-lemma-base-change-quasi-affine}。因此可应用 Lemma
\ref{descent-lemma-descending-types-morphisms}，只需在 fpqc 覆盖由仿射概形之间的
单个平坦满射 $\{X \to S\}$ 给出时证明本引理。记 $X = \Spec(A)$ 及
$S = \Spec(R)$，于是 $R \to A$ 是忠实平坦环映射。设 $(V, \varphi)$ 为
相对于 $S$ 上 $X$ 的下降数据，并假设 $\pi : V \to X$ 拟仿射。

\medskip\noindent
由 Morphisms, Lemma \ref{morphisms-lemma-characterize-quasi-affine}，这表示
$$
V \longrightarrow \underline{\Spec}_X(\pi_*\mathcal{O}_V) = W
$$
是 $X$ 上概形的拟紧开浸入。投影
$\text{pr}_i : X \times_S X \to X$ 是平坦的，因而由平坦基变换
（Cohomology of Schemes, Lemma
\ref{coherent-lemma-flat-base-change-cohomology}），有
$$
\text{pr}_0^*\pi_*\mathcal{O}_V =
(\pi \times \text{id}_X)_*\mathcal{O}_{V \times_S X}, \quad
\text{pr}_1^*\pi_*\mathcal{O}_V =
(\text{id}_X \times \pi)_*\mathcal{O}_{X \times_S V}
$$
。因此同构 $\varphi : V \times_S X \to X \times_S V$（它是
$X \times_S X$ 上的同构）诱导 $X \times_S X$ 上拟凝聚代数层的同构
$$
\varphi^\sharp :
\text{pr}_0^*\pi_*\mathcal{O}_V
\longrightarrow
\text{pr}_1^*\pi_*\mathcal{O}_V
$$
。$\varphi$ 的余圈条件蕴含 $\varphi^\sharp$ 的余圈条件。换言之，它在仿射
概形 $W$ 上产生相对于 $S$ 上 $X$ 的下降数据 $\varphi'$，并且态射
$V \to W$ 是下降数据的态射。因此由 Lemma \ref{descent-lemma-affine}（或由拟凝聚
代数下降的有效性），得到概形 $U' \to S$ 以及下降数据同构
$(W, \varphi') \cong (X \times_S U', can)$。顺便指出，由 Lemma
\ref{descent-lemma-descending-property-affine}，$U'$ 是仿射的。

\medskip\noindent
现在可把 $V$ 看成一个（拟紧）开集 $V \subset X \times_S U'$，它在下降
数据
$$
can : X \times_S U' \times_S X \to X \times_S X \times_S U',
(x_0, u', x_1) \mapsto (x_0, x_1, u').
$$
下稳定。换言之，对映到 $S$ 中同一点的任意 $x_0, x_1, u'$，有
$(x_0, u') \in V \Rightarrow (x_1, u') \in V$。因为 $X \to S$ 满射，
立即可知 $V$ 是某个子集 $U \subset U'$ 在态射
$X \times_S U' \to U'$ 下的逆像。因为 $X \to S$ 拟紧、平坦且满射，
$X \times_S U' \to U'$ 也拟紧、平坦且满射。因此由 Morphisms, Lemma
\ref{morphisms-lemma-fpqc-quotient-topology}，子集 $U \subset U'$ 是开集，
证毕。
\end{proof}












\section{以层表述下降数据}
\label{descent-section-descent-data-sheaves}


\noindent
对位点上的覆盖，下面给出理解下降数据的另一种方式。

\begin{lemma}
\label{descent-lemma-descent-data-sheaves}
设 $\tau \in \{Zariski, fppf, \etale, smooth, syntomic\}$\footnote{这里没有
fpqc 并非笔误。见 Topologies, Section
\ref{topologies-section-fpqc} 中的讨论。}。设 $\Sch_\tau$ 为大
$\tau$-位点，设 $S \in \Ob(\Sch_\tau)$。设
$\{S_i \to S\}_{i \in I}$ 为位点 $(\Sch/S)_\tau$ 中的覆盖。存在范畴等价
$$
\left\{
\begin{matrix}
\text{下降数据 }(X_i, \varphi_{ii'})\text{ 使得}\\
\text{每个 }X_i \in \Ob((\Sch/S)_\tau)
\end{matrix}
\right\}
\leftrightarrow
\left\{
\begin{matrix}
\text{层 }F\text{ 定义于 }(\Sch/S)_\tau\text{ 使得}\\
\text{每个 }h_{S_i} \times F\text{ 可表}
\end{matrix}
\right\}.
$$
此外：
\begin{enumerate}
\item 右侧表示 $h_{S_i} \times F$ 的对象对应于左侧的概形 $X_i$；并且
\item 层 $F$ 可表示，当且仅当相应的下降数据
$(X_i, \varphi_{ii'})$ 有效。
\end{enumerate}
\end{lemma}

\begin{proof}
我们已在 Section \ref{descent-section-fpqc-universal-effective-epimorphisms} 中看到，
可表示预层是位点 $(\Sch/S)_\tau$ 上的层。此外，米田引理（Categories,
Lemma \ref{categories-lemma-yoneda}）保证可表示层之间的映射与表示对象
之间的映射一一对应。在证明中将不再说明地使用这些事实。

\medskip\noindent
先构造从右到左的函子。设 $F$ 为 $(\Sch/S)_\tau$ 上的层，并且每个
$h_{S_i} \times F$ 都可表示。在这种情形下，令 $X_i$ 为
$(\Sch/S)_\tau$ 中的一个表示对象。它带有态射 $X_i \to S_i$。于是
$X_i \times_S S_{i'}$ 与 $S_i \times_S X_{i'}$ 都表示层
$h_{S_i} \times F \times h_{S_{i'}}$，因而得到同构
$$
\varphi_{ii'} : X_i \times_S S_{i'} \to S_i \times_S X_{i'}
$$
容易看出，态射 $\varphi_{ii'}$ 是 $S_i \times_S S_{i'}$ 上的态射并满足
余圈条件。从右到左的函子由构造
$F \mapsto (X_i, \varphi_{ii'})$ 给出。

\medskip\noindent
再构造从左到右的函子。对每个 $i$，用 $F_i$ 表示层 $h_{X_i}$。同构
$\varphi_{ii'}$ 给出同构
$$
\varphi_{ii'} :
F_i \times h_{S_{i'}}
\longrightarrow
h_{S_i} \times F_{i'}
$$
，它位于 $h_{S_i} \times h_{S_{i'}}$ 上。令 $F$ 等于下图中的余等化子
$$
\xymatrix{
\coprod_{i, i'} F_i \times h_{S_{i'}}
\ar@<1ex>[rr]^-{\text{pr}_0}
\ar@<-1ex>[rr]_-{\text{pr}_1 \circ \varphi_{ii'}}
& &
\coprod_i F_i \ar[r]
&
F
}
$$
余圈条件保证 $h_{S_i} \times F$ 同构于 $F_i$，因而可表示。从左到右的
函子由构造 $(X_i, \varphi_{ii'}) \mapsto F$ 给出。

\medskip\noindent
略去这些构造互为拟逆函子的验证。最后的陈述 (1) 与 (2) 由这些构造推出。
\end{proof}

\begin{remark}
\label{descent-remark-what-product-means}
在 Lemma \ref{descent-lemma-descent-data-sheaves} 的陈述中，
$h_{S_i} \times F$ 可表示这一条件，等价于 $F$ 在
$(\Sch/S_i)_\tau$ 上的限制可表示。
\end{remark}
% Generated mechanically from the frozen perfect.tex target.
% Do not edit this generated file; edit the bound locale source instead.

% OK, start here.
%

\setcounter{chapter}{35}
\chapter{概形的导出范畴}



\phantomsection
\label{perfect-section-phantom}


\section{引言}
\label{perfect-section-introduction}

\noindent
本章讨论概形上模的导出范畴。这里讨论的大部分材料可见于
\cite{TT}、\cite{Bokstedt-Neeman}、\cite{BvdB} 和 \cite{LN}。
当然还有许多其他参考文献。


\section{约定}
\label{perfect-section-conventions}

\noindent
若 $\mathcal{A}$ 为 Abel 范畴，且 $M$ 为 $\mathcal{A}$ 的对象，则也用
$M$ 表示 $K(\mathcal{A})$ 和/或 $D(\mathcal{A})$ 中对应于如下复形的
对象：它在次数 $0$ 上为 $M$，在所有其他次数上为零。

\medskip\noindent
若有环 $A$，则 $K(A)$ 表示 $A$-模复形的同伦范畴，$D(A)$ 表示相应的
导出范畴。类似地，若有赋环空间 $(X, \mathcal{O}_X)$，则
$K(\mathcal{O}_X)$ 表示 $\mathcal{O}_X$-模复形的同伦范畴，
$D(\mathcal{O}_X)$ 表示相应的导出范畴。










\section{拟凝聚模的导出范畴}
\label{perfect-section-derived-quasi-coherent}

\noindent
本节讨论概形 $X$ 上拟凝聚模与所有模之间的关系。可参考
\cite[Appendix B]{TT}。由 Schemes, Section
\ref{schemes-section-quasi-coherent} 中的讨论，嵌入
$\QCoh(\mathcal{O}_X) \subset \textit{Mod}(\mathcal{O}_X)$
把 $\QCoh(\mathcal{O}_X)$ 展示为 $\mathcal{O}_X$-模范畴的弱 Serre
子范畴。用
$$
D_\QCoh(\mathcal{O}_X) \subset D(\mathcal{O}_X)
$$
表示上同调层为拟凝聚层的复形所成的子范畴，见 Derived Categories,
Section \ref{derived-section-triangulated-sub}。于是得到典范函子
\begin{equation}
\label{perfect-equation-compare}
D(\QCoh(\mathcal{O}_X))
\longrightarrow
D_\QCoh(\mathcal{O}_X)
\end{equation}
，见 Derived Categories, Equation (\ref{derived-equation-compare})。

\begin{lemma}
\label{perfect-lemma-quasi-coherence-direct-sums}
设 $X$ 为概形。那么 $D_\QCoh(\mathcal{O}_X)$ 有直和。
\end{lemma}

\begin{proof}
由 Injectives, Lemma \ref{injectives-lemma-derived-products}，导出范畴
$D(\mathcal{O}_X)$ 有直和，并且可通过对任意代表逐项取直和来计算。由于
取直和（在任意 Grothendieck Abel 范畴中）是正合函子，直和的上同调层
显然是各上同调层的直和。拟凝聚层的直和仍为拟凝聚层，故结论成立，见
Schemes, Section \ref{schemes-section-quasi-coherent}。
\end{proof}

\noindent
我们需要一些关于导出极限的知识。提醒读者：在下面的引理中，导出极限通常
不是 $D_\QCoh$ 的对象。

\begin{lemma}
\label{perfect-lemma-Rlim-quasi-coherent}
设 $X$ 为概形。设 $(K_n)$ 为 $D_\QCoh(\mathcal{O}_X)$ 中的逆系统，其在
$D(\mathcal{O}_X)$ 中的导出极限为 $K = R\lim K_n$。假设对所有
$q \in \mathbf{Z}$ 及 $n \geq 1$，态射
$H^q(K_{n + 1}) \to H^q(K_n)$ 都是满射。那么：
\begin{enumerate}
\item $H^q(K) = \lim H^q(K_n)$；
\item $R\lim H^q(K_n) = \lim H^q(K_n)$；并且
\item 对每个仿射开集 $U \subset X$，当 $p > 0$ 时有
$H^p(U, \lim H^q(K_n)) = 0$。
\end{enumerate}
\end{lemma}

\begin{proof}
令 $\mathcal{B}$ 为 $X$ 的仿射开集全体。因为 $H^q(K_n)$ 拟凝聚，由
Cohomology of Schemes, Lemma
\ref{coherent-lemma-quasi-coherent-affine-cohomology-zero}，对
$U \in \mathcal{B}$ 有 $H^p(U, H^q(K_n)) = 0$。此外，由 Schemes,
Lemma \ref{schemes-lemma-equivalence-quasi-coherent}，对
$U \in \mathcal{B}$，映射
$H^0(U, H^q(K_{n + 1})) \to H^0(U, H^q(K_n))$ 是满射。我们刚验证了
Cohomology, Lemma \ref{cohomology-lemma-derived-limit-suitable-system} 的
条件，故 (1) 由其推出。(2) 与 (3) 由 Cohomology, Lemma
\ref{cohomology-lemma-inverse-limit-is-derived-limit} 推出。
\end{proof}

\noindent
下面的引理将帮助我们在 $D_\QCoh(\mathcal{O}_X)$ 的对象上“计算”右导出
函子。

\begin{lemma}
\label{perfect-lemma-nice-K-injective}
设 $X$ 为概形，设 $E$ 为 $D_\QCoh(\mathcal{O}_X)$ 的对象。那么 Derived
Categories, Remark
\ref{derived-remark-map-into-derived-limit-truncations}
中的态射 $E \to R\lim \tau_{\geq -n}E$ 是同构\footnote{特别地，由
Derived Categories, Lemma \ref{derived-lemma-difficulty-K-injectives}，
$E$ 有 K-内射代表。}。
\end{lemma}

\begin{proof}
用 $\mathcal{H}^i = H^i(E)$ 表示 $E$ 的第 $i$ 个上同调层。令
$\mathcal{B}$ 为 $X$ 的仿射开子集全体。由 Cohomology of Schemes,
Lemma \ref{coherent-lemma-quasi-coherent-affine-cohomology-zero}，对所有
$p > 0$、所有 $i \in \mathbf{Z}$ 以及所有 $U \in \mathcal{B}$，有
$H^p(U, \mathcal{H}^i) = 0$。因此结论由 Cohomology, Lemma
\ref{cohomology-lemma-is-limit-dimension} 推出。
\end{proof}

\begin{lemma}
\label{perfect-lemma-application-nice-K-injective}
设 $X$ 为概形。设 $F : \textit{Mod}(\mathcal{O}_X) \to \textit{Ab}$ 为
加性函子，$N \geq 0$ 为整数。假设：
\begin{enumerate}
\item $F$ 与可数直积交换；
\item 对所有 $p \geq N$ 及拟凝聚 $\mathcal{F}$，有
$R^pF(\mathcal{F}) = 0$。
\end{enumerate}
那么对 $E \in D_\QCoh(\mathcal{O}_X)$：
\begin{enumerate}
\item 当 $i \leq a$ 时，
$H^i(RF(\tau_{\leq a}E)) \to H^i(RF(E))$ 是同构；
\item 当 $i \geq b$ 时，
$H^i(RF(E)) \to H^i(RF(\tau_{\geq b - N + 1}E))$ 是同构；
\item 若对某个 $-\infty \leq a \leq b \leq \infty$，当
$i \not \in [a, b]$ 时有 $H^i(E) = 0$，则当
$i \not \in [a, b + N - 1]$ 时有 $H^i(RF(E)) = 0$。
\end{enumerate}
\end{lemma}

\begin{proof}
陈述 (1) 即 Derived Categories, Lemma
\ref{derived-lemma-negative-vanishing}。

\medskip\noindent
证明陈述 (2)。记 $E_n = \tau_{\geq -n}E$。由 Lemma
\ref{perfect-lemma-nice-K-injective}，有 $E = R\lim E_n$。于是由 Injectives,
Lemma \ref{injectives-lemma-RF-commutes-with-Rlim}，在
$D(\textit{Ab})$ 中有 $RF(E) = R\lim RF(E_n)$。因此对每个
$i \in \mathbf{Z}$，有短正合列
$$
0 \to R^1\lim H^{i - 1}(RF(E_n)) \to H^i(RF(E)) \to \lim H^i(RF(E_n)) \to 0
$$
，见 More on Algebra, Remark
\ref{more-algebra-remark-compare-derived-limit}。为证明 (2)，我们将说明左端
项为零，并且对任意满足 $i \geq b$ 的 $b$，右端项等于
$H^i(RF(E_{-b + N - 1}))$。

\medskip\noindent
对每个 $n$，有可区别三角
$$
H^{-n}(E)[n] \to E_n \to E_{n - 1} \to H^{-n}(E)[n + 1]
$$
，它位于 $D(\mathcal{O}_X)$ 中（Derived Categories, Remark
\ref{derived-remark-truncation-distinguished-triangle}）。因为
$H^{-n}(E)$ 拟凝聚，有
$$
H^i(RF(H^{-n}(E)[n])) = R^{i + n}F(H^{-n}(E)) = 0
$$
（当 $i + n \geq N$ 时），以及
$$
H^i(RF(H^{-n}(E)[n + 1])) = R^{i + n + 1}F(H^{-n}(E)) = 0
$$
（当 $i + n + 1 \geq N$ 时）。由此可知
$$
H^i(RF(E_n)) \to H^i(RF(E_{n - 1}))
$$
在 $n \geq N - i$ 时是同构。因此系统 $H^i(RF(E_n))$ 都满足 ML 条件，
短正合列中的 $R^1\lim$ 项为零（见 More on Algebra, Section
\ref{more-algebra-section-Rlim} 中的讨论）。此外，系统
$H^i(RF(E_n))$ 从 $n = N - i - 1$ 起为常值，正合所需。

\medskip\noindent
证明 (3)。在关于 $E$ 的假设下，有 $\tau_{\leq a - 1}E = 0$，故由 (1)
得到：当 $i \leq a - 1$ 时，$H^i(RF(E))$ 为零。类似地，
$\tau_{\geq b + 1}E = 0$，故由 (2) 得到：当 $i \geq b + N$ 时，
$H^i(RF(E))$ 为零。
\end{proof}

\noindent
下面的引理是本章许多结果的关键要素。

\begin{lemma}
\label{perfect-lemma-affine-compare-bounded}
设 $X = \Spec(A)$ 为仿射概形。图中的所有函子
$$
\xymatrix{
D(\QCoh(\mathcal{O}_X)) \ar[rr]_{(\ref{perfect-equation-compare})}
& &
D_\QCoh(\mathcal{O}_X) \ar[ld]^{R\Gamma(X, -)} \\
& D(A) \ar[lu]^{\widetilde{\ \ }}
}
$$
都是三角范畴的等价。此外，对 $D_\QCoh(\mathcal{O}_X)$ 中的 $E$，有
$H^0(X, E) = H^0(X, H^0(E))$。
\end{lemma}

\begin{proof}
函子 $R\Gamma(X, -)$ 给出函子 $D(\mathcal{O}_X) \to D(A)$，因而限制后
给出函子
\begin{equation}
\label{perfect-equation-back}
R\Gamma(X, -) : D_\QCoh(\mathcal{O}_X) \longrightarrow D(A).
\end{equation}
利用 $X$ 上拟凝聚模与 $A$-模范畴之间的等价，我们将说明此函子是
(\ref{perfect-equation-compare}) 的拟逆。

\medskip\noindent
说明。用 $(Y, \mathcal{O}_Y)$ 表示其环层由 $A$ 给出的单点空间，用
$\pi : (X, \mathcal{O}_X) \to (Y, \mathcal{O}_Y)$
表示显然的赋环空间态射。于是 $R\Gamma(X, -)$ 可等同于 $R\pi_*$，并且
通过等价
$\textit{Mod}(\mathcal{O}_Y) = \text{Mod}_A = \QCoh(\mathcal{O}_X)$
，函子 (\ref{perfect-equation-compare}) 可等同于
$L\pi^* = \pi^* = \widetilde{\ }$（见 Modules, Lemma
\ref{modules-lemma-construct-quasi-coherent-sheaves} 以及 Schemes, Lemmas
\ref{schemes-lemma-compare-constructions} 和
\ref{schemes-lemma-equivalence-quasi-coherent}）。因此函子
$$
\xymatrix{
D(A) \ar@<1ex>[r] & D(\mathcal{O}_X) \ar@<1ex>[l]
}
$$
互为伴随（由 Cohomology, Lemma \ref{cohomology-lemma-adjoint}）。特别地，
得到典范伴随映射
$$
a : \widetilde{R\Gamma(X, E)} \longrightarrow E
$$
（对 $D(\mathcal{O}_X)$ 中的 $E$），以及
$$
b : M^\bullet \longrightarrow R\Gamma(X, \widetilde{M^\bullet})
$$
（对 $A$-模复形 $M^\bullet$）。

\medskip\noindent
设 $E$ 为 $D_\QCoh(\mathcal{O}_X)$ 的对象。可以把 Lemma
\ref{perfect-lemma-application-nice-K-injective} 应用于函子
$F(-) = \Gamma(X, -)$；由 Cohomology of Schemes, Lemma
\ref{coherent-lemma-quasi-coherent-affine-cohomology-zero}，可取
$N = 1$。因此
$$
H^0(R\Gamma(X, E)) = H^0(R\Gamma(X, \tau_{\geq 0}E)) = \Gamma(X, H^0(E))
$$
（最后一个等号来自典范截断的定义）。利用这一点，我们将说明伴随映射
$a$ 与 $b$ 诱导同构 $H^0(a)$ 与 $H^0(b)$。于是 $a$ 与 $b$ 都是拟同构
（因为该陈述在平移下不变），从而引理得证。

\medskip\noindent
在两种情形下，都使用 $\widetilde{\ }$ 为正合函子这一事实（Schemes,
Lemma \ref{schemes-lemma-spec-sheaves}）。也就是说，这说明
$$
H^0\left(\widetilde{R\Gamma(X, E)}\right) =
\widetilde{H^0(R\Gamma(X, E))} =
\widetilde{\Gamma(X, H^0(E))}
$$
等于 $H^0(E)$，因为 $H^0(E)$ 拟凝聚。因此 $H^0(a)$ 是同构。另一方面，
有
$$
H^0(R\Gamma(X, \widetilde{M^\bullet})) =
\Gamma(X, H^0(\widetilde{M^\bullet})) =
\Gamma(X, \widetilde{H^0(M^\bullet)}) =
H^0(M^\bullet)
$$
，这证明 $H^0(b)$ 是同构。
\end{proof}

\begin{lemma}
\label{perfect-lemma-affine-K-flat}
设 $X = \Spec(A)$ 为仿射概形。若 $K^\bullet$ 为 $A$-模的 K-平坦复形，
则 $\widetilde{K^\bullet}$ 为 $\mathcal{O}_X$-模的 K-平坦复形。
\end{lemma}

\begin{proof}
由 More on Algebra, Lemma \ref{more-algebra-lemma-base-change-K-flat}，对每个
$\mathfrak p \in \Spec(A)$，$K^\bullet \otimes_A A_\mathfrak p$ 是
$A_\mathfrak p$-模的 K-平坦复形。因此由 Cohomology, Lemma
\ref{cohomology-lemma-check-K-flat-stalks}（以及 Schemes, Lemma
\ref{schemes-lemma-spec-sheaves}），可知 $\widetilde{K^\bullet}$ K-平坦。
\end{proof}

\begin{lemma}
\label{perfect-lemma-quasi-coherence-pushforward}
若 $f : X \to Y$ 为由环映射 $A \to B$ 给出的仿射概形态射，则图
$$
\xymatrix{
D(B) \ar[d] \ar[r] & D_\QCoh(\mathcal{O}_X) \ar[d]^{Rf_*} \\
D(A) \ar[r] & D_\QCoh(\mathcal{O}_Y)
}
$$
交换。
\end{lemma}

\begin{proof}
由 Lemma \ref{perfect-lemma-affine-compare-bounded} 以及
$R\Gamma(Y, Rf_*K) = R\Gamma(X, K)$ 得证；后一个等式见 Cohomology,
Lemma \ref{cohomology-lemma-Leray-unbounded}。
\end{proof}

\begin{lemma}
\label{perfect-lemma-quasi-coherence-pullback}
设 $f : Y \to X$ 为概形态射。
\begin{enumerate}
\item 函子 $Lf^*$ 把 $D_\QCoh(\mathcal{O}_X)$ 映入
$D_\QCoh(\mathcal{O}_Y)$。
\item 若 $X$ 与 $Y$ 都仿射，且 $f$ 由环映射 $A \to B$ 给出，则图
$$
\xymatrix{
D(B) \ar[r] & D_\QCoh(\mathcal{O}_Y) \\
D(A) \ar[r] \ar[u]^{- \otimes_A^\mathbf{L} B} &
D_\QCoh(\mathcal{O}_X) \ar[u]_{Lf^*}
}
$$
交换。
\end{enumerate}
\end{lemma}

\begin{proof}
先证明图
$$
\xymatrix{
D(B) \ar[r] & D(\mathcal{O}_Y) \\
D(A) \ar[r] \ar[u]^{- \otimes_A^\mathbf{L} B} &
D(\mathcal{O}_X) \ar[u]_{Lf^*}
}
$$
交换。这由 Lemma \ref{perfect-lemma-affine-K-flat} 及有关函子的构造立即可知。为证
(1)，设 $E$ 为 $D_\QCoh(\mathcal{O}_X)$ 的对象。要说明 $Lf^*E$ 的
上同调层拟凝聚，可以在 $X$ 上局部地工作。注意 $Lf^*$ 与限制到开子概形
相容。因此可假设 $f$ 是 (2) 中那样的仿射概形态射。由 Lemma
\ref{perfect-lemma-affine-compare-bounded}，$E$ 来自一个 $A$-模复形。由证明中
第一个图的交换性，$Lf^*E$ 也有同样性质，从而 (1) 成立。
\end{proof}

\begin{lemma}
\label{perfect-lemma-quasi-coherence-tensor-product}
设 $X$ 为概形。
\begin{enumerate}
\item 对 $D_\QCoh(\mathcal{O}_X)$ 的对象 $K, L$，导出张量积
$K \otimes^\mathbf{L}_{\mathcal{O}_X} L$ 属于
$D_\QCoh(\mathcal{O}_X)$。
\item 若 $X = \Spec(A)$ 仿射，则
$$
\widetilde{M^\bullet} \otimes_{\mathcal{O}_X}^\mathbf{L} \widetilde{K^\bullet}
=
\widetilde{M^\bullet \otimes_A^\mathbf{L} K^\bullet}
$$
对任意一对 $A$-模复形 $K^\bullet$、$M^\bullet$ 成立。
\end{enumerate}
\end{lemma}

\begin{proof}
(2) 中的等式由 Lemma \ref{perfect-lemma-affine-K-flat} 和导出张量积的构造立即
推出。为证 (1)，设 $K, L$ 为 $D_\QCoh(\mathcal{O}_X)$ 的对象。要检验
$K \otimes^\mathbf{L} L$ 属于 $D_\QCoh(\mathcal{O}_X)$，可以在 $X$ 上
局部地工作，故可假设 $X = \Spec(A)$ 仿射。由 Lemma
\ref{perfect-lemma-affine-compare-bounded}，可用 $A$-模复形表示 $K$ 与 $L$。
于是结论由 (2) 推出。
\end{proof}





\section{全直像}
\label{perfect-section-total-direct-image}

\noindent
下面的引理是 Cohomology of Schemes, Lemma
\ref{coherent-lemma-quasi-coherence-higher-direct-images} 的类似物。

\begin{lemma}
\label{perfect-lemma-quasi-coherence-direct-image}
设 $f : X \to S$ 为概形态射。假设 $f$ 拟分离且拟紧。
\begin{enumerate}
\item 函子 $Rf_*$ 把 $D_\QCoh(\mathcal{O}_X)$ 映入
$D_\QCoh(\mathcal{O}_S)$。
\item 若 $S$ 拟紧，则存在整数 $N = N(X, S, f)$，使得对
$D_\QCoh(\mathcal{O}_X)$ 的对象 $E$，只要对 $m > 0$ 有
$H^m(E) = 0$，就对 $m \geq N$ 有 $H^m(Rf_*E) = 0$。
\item 事实上，若 $S$ 拟紧，则可以找到 $N = N(X, S, f)$，使得对每个
概形态射 $S' \to S$，函子 $R(f')_*$ 都有同样结论，其中
$f' : X' \to S'$ 是 $f$ 的基变换。
\end{enumerate}
\end{lemma}

\begin{proof}
设 $E$ 为 $D_\QCoh(\mathcal{O}_X)$ 的对象。为证 (1)，必须说明 $Rf_*E$
有拟凝聚上同调层。该问题在 $S$ 上为局部问题，故可假设 $S$ 拟紧。按照
Cohomology of Schemes, Lemma
\ref{coherent-lemma-quasi-coherence-higher-direct-images} 选取
$N = N(X, S, f)$。于是对所有拟凝聚 $\mathcal{O}_X$-模 $\mathcal{F}$
及所有 $p \geq N$，有 $R^pf_*\mathcal{F} = 0$，并且基变换后仍然如此。

\medskip\noindent
先假设 $E$ 下有界。我们将说明，对这样的 $E$，以上选取的 $N$ 使 (1)、
(2) 与 (3) 成立。在这种情形下，例如可使用谱序列
$$
R^pf_*H^q(E) \Rightarrow R^{p + q}f_*E
$$
（Derived Categories, Lemma
\ref{derived-lemma-two-ss-complex-functor}），结合
$R^pf_*H^q(E)$ 的拟凝聚性以及 $R^pf_*H^q(E)$ 在 $p \geq N$ 时的消没，
得到此情形下
(1)、(2) 与 (3) 成立。

\medskip\noindent
下面证明 (2) 与 (3)。假设当 $m > 0$ 时 $H^m(E) = 0$。设
$U \subset S$ 为仿射开集。由 Cohomology of Schemes, Lemma
\ref{coherent-lemma-quasi-coherence-higher-direct-images-application}
以及 $N$ 的选取，对 $p \geq N$ 及任意拟凝聚 $\mathcal{O}_X$-模
$\mathcal{F}$，有 $H^p(f^{-1}(U), \mathcal{F}) = 0$。因此可把 Lemma
\ref{perfect-lemma-application-nice-K-injective} 应用于函子
$\Gamma(f^{-1}(U), -)$，得到
$$
R\Gamma(U, Rf_*E) = R\Gamma(f^{-1}(U), E)
$$
在次数 $\geq N$ 上的上同调消没。因为这对所有仿射开集 $U \subset S$
都成立，所以当 $m \geq N$ 时 $H^m(Rf_*E) = 0$。

\medskip\noindent
最后在一般情形下证明 (1)。回忆存在可区别三角
$$
\tau_{\leq -n - 1}E \to E \to \tau_{\geq -n}E \to
(\tau_{\leq -n - 1}E)[1]
$$
，它位于 $D(\mathcal{O}_X)$ 中，见 Derived Categories, Remark
\ref{derived-remark-truncation-distinguished-triangle}。由 (2)，
$Rf_*\tau_{\leq -n - 1}E$ 的上同调层在次数 $\geq -n + N$ 上消没。因此，
给定整数 $q$，可知对某个 $n$，$R^qf_*E$ 等于
$R^qf_*\tau_{\geq -n}E$，于是可应用上面的结果。
\end{proof}

\begin{lemma}
\label{perfect-lemma-acyclicity-lemma}
设 $f : X \to S$ 为拟分离且拟紧的概形态射。设
$\mathcal{F}^\bullet$ 为拟凝聚 $\mathcal{O}_X$-模复形，其每一项关于
$f_*$ 都右无环。那么 $f_*\mathcal{F}^\bullet$ 在 $D(\mathcal{O}_S)$ 中
表示 $Rf_*\mathcal{F}^\bullet$。
\end{lemma}

\begin{proof}
总有典范映射 $f_*\mathcal{F}^\bullet \to Rf_*\mathcal{F}^\bullet$。我们的
任务是说明它在上同调层上为同构。因为该陈述在平移下不变，只需证明
$H^0(f_*(\mathcal{F}^\bullet)) \to R^0f_*\mathcal{F}^\bullet$
是同构。该陈述在 $S$ 上为局部问题，故可假设 $S$ 仿射。由 Lemma
\ref{perfect-lemma-quasi-coherence-direct-image}，对所有充分大的 $n$，有
$R^0f_*\mathcal{F}^\bullet = R^0f_*\tau_{\geq -n}\mathcal{F}^\bullet$。
因此可假设 $\mathcal{F}^\bullet$ 下有界。按假设，每个
$\mathcal{F}^n$ 关于 $f_*$ 都右无环，故由 Leray 无环性引理（Derived
Categories, Lemma \ref{derived-lemma-leray-acyclicity}），
$f_*\mathcal{F}^\bullet \to Rf_*\mathcal{F}^\bullet$ 是拟同构。
\end{proof}

\begin{lemma}
\label{perfect-lemma-acyclicity-lemma-global}
设 $X$ 为拟分离拟紧概形。设 $\mathcal{F}^\bullet$ 为拟凝聚
$\mathcal{O}_X$-模复形，其每一项关于 $\Gamma(X, -)$ 都右无环。那么
$\Gamma(X, \mathcal{F}^\bullet)$ 在 $D(\Gamma(X, \mathcal{O}_X)$ 中表示
$R\Gamma(X, \mathcal{F}^\bullet)$。
\end{lemma}

\begin{proof}
把 Lemma \ref{perfect-lemma-acyclicity-lemma} 应用于典范态射
$X \to \Spec(\Gamma(X, \mathcal{O}_X))$。略去一些细节。
\end{proof}

\begin{lemma}
\label{perfect-lemma-spectral-sequence}
设 $X$ 为拟分离拟紧概形。对 $D_\QCoh(\mathcal{O}_X)$ 的任意对象 $K$，
谱序列
$$
E_2^{i, j} = H^i(X, H^j(K)) \Rightarrow H^{i + j}(X, K)
$$
（Cohomology, Example \ref{cohomology-example-spectral-sequence}）有界且收敛。
\end{lemma}

\begin{proof}
由 Cohomology, Lemma
\ref{cohomology-lemma-spectral-sequence-filtered-object} 中使用
$\tau_{\leq -p}K$ 所给滤过构造谱序列的方法可知，只需说明：给定
$n \in \mathbf{Z}$，有
$$
H^n(X, \tau_{\leq -p}K) = 0 \text{ 对 } p \gg 0
$$
以及
$$
H^n(X, K) = H^n(X, \tau_{\leq -p}K) \text{ 对 } p \ll 0
$$
。第一点来自 Lemma \ref{perfect-lemma-application-nice-K-injective}（取
$F = \Gamma(X, -)$）及 Cohomology of Schemes, Lemma
\ref{coherent-lemma-quasi-coherence-higher-direct-images} 中的界。对任意
赋环空间 $(X, \mathcal{O}_X)$ 及任意 $K \in D(\mathcal{O}_X)$，只要
$-p \leq n$，第二点就成立。
\end{proof}

\begin{lemma}
\label{perfect-lemma-quasi-coherence-pushforward-direct-sums}
设 $f : X \to S$ 为拟分离且拟紧的概形态射。那么
$Rf_* : D_\QCoh(\mathcal{O}_X) \to D_\QCoh(\mathcal{O}_S)$
与直和交换。
\end{lemma}

\begin{proof}
设 $E_i$ 为 $D_\QCoh(\mathcal{O}_X)$ 的一族对象，并令
$E = \bigoplus E_i$。要证明映射
$$
\bigoplus Rf_*E_i \longrightarrow Rf_*E
$$
是同构。我们将说明它在次数 $0$ 的上同调层上诱导同构，这就推出本引理。
为此可以在 $S$ 上局部地工作，故可以并且确实假设 $S$ 拟紧。按照 Lemma
\ref{perfect-lemma-quasi-coherence-direct-image} 选取整数 $N$。由所引引理，
$R^0f_*E = R^0f_*\tau_{\geq -N}E$ 且
$R^0f_*E_i = R^0f_*\tau_{\geq -N}E_i$。注意
$\tau_{\geq -N}E = \bigoplus \tau_{\geq -N}E_i$。因此可假设所有 $E_i$ 的
上同调层在次数 $< -N$ 上消没。接着使用谱序列
$$
R^pf_*H^q(E) \Rightarrow R^{p + q}f_*E
\quad\text{且}\quad
R^pf_*H^q(E_i) \Rightarrow R^{p + q}f_*E_i
$$
（Derived Categories, Lemma
\ref{derived-lemma-two-ss-complex-functor}），把问题约化到拟凝聚层直和的
情形。该情形由 Cohomology of Schemes, Lemma
\ref{coherent-lemma-colimit-cohomology} 处理。
\end{proof}









\section{仿射态射}
\label{perfect-section-affine-morphisms}

\noindent
本节汇集沿概形的仿射态射取直像的一些信息。

\begin{lemma}
\label{perfect-lemma-pushforward-affine-morphism}
设 $f : X \to S$ 为概形的仿射态射。设 $\mathcal{F}^\bullet$ 为拟凝聚
$\mathcal{O}_X$-模复形。那么
$f_*\mathcal{F}^\bullet = Rf_*\mathcal{F}^\bullet$.
\end{lemma}

\begin{proof}
合并 Lemma \ref{perfect-lemma-acyclicity-lemma} 与 Cohomology of Schemes, Lemma
\ref{coherent-lemma-relative-affine-vanishing}。另一种证明是在 $S$ 上仿射
局部地工作，并使用 Lemma \ref{perfect-lemma-quasi-coherence-pushforward}。
\end{proof}

\begin{lemma}
\label{perfect-lemma-affine-morphism}
设 $f : X \to S$ 为概形的仿射态射。那么
$Rf_* : D_\QCoh(\mathcal{O}_X) \to D_\QCoh(\mathcal{O}_S)$
反映同构。
\end{lemma}

\begin{proof}
该陈述是指：若 $Rf_*\alpha$ 是同构，则 $D_\QCoh(\mathcal{O}_X)$ 的态射
$\alpha : E \to F$ 是同构。可以在上同调层上检验这一点。特别地，该问题
在 $S$ 上为局部问题。因此可假设 $S$ 仿射，进而 $X$ 也仿射。在这种
情形下，由 Lemma \ref{perfect-lemma-affine-compare-bounded} 给出的导出范畴
$D_\QCoh(\mathcal{O}_X)$ 与 $D_\QCoh(\mathcal{O}_S)$ 的描述，结论显然。
略去一些细节。
\end{proof}

\begin{lemma}
\label{perfect-lemma-affine-morphism-pull-push}
设 $f : X \to S$ 为概形的仿射态射。对 $D_\QCoh(\mathcal{O}_S)$ 中的
$E$，有
$Rf_* Lf^* E = E \otimes^\mathbf{L}_{\mathcal{O}_S} f_*\mathcal{O}_X$.
\end{lemma}

\begin{proof}
因为 $f$ 仿射，映射 $f_*\mathcal{O}_X \to Rf_*\mathcal{O}_X$ 是同构
（Cohomology of Schemes, Lemma
\ref{coherent-lemma-relative-affine-vanishing}）。存在典范映射
$E \otimes^\mathbf{L} f_*\mathcal{O}_X =
E \otimes^\mathbf{L} Rf_*\mathcal{O}_X \to Rf_* Lf^* E$，它伴随于映射
$$
Lf^*(E \otimes^\mathbf{L} Rf_*\mathcal{O}_X) =
Lf^*E \otimes^\mathbf{L} Lf^*Rf_*\mathcal{O}_X \longrightarrow
Lf^* E \otimes^\mathbf{L} \mathcal{O}_X = Lf^* E
$$
；后者来自 $1 : Lf^*E \to Lf^*E$ 以及典范映射
$Lf^*Rf_*\mathcal{O}_X \to \mathcal{O}_X$。要检验这样构造的映射为同构，
可以在 $S$ 上局部地工作。因此可假设 $S$ 仿射，进而 $X$ 也仿射。在这种
情形下，由 Lemmas \ref{perfect-lemma-affine-compare-bounded} 和
\ref{perfect-lemma-quasi-coherence-pullback} 给出的导出范畴
$D_\QCoh(\mathcal{O}_X)$、$D_\QCoh(\mathcal{O}_S)$ 及函子 $Lf^*$ 的
描述，结论显然。略去一些细节。
\end{proof}

\noindent
设 $Y$ 为概形。设 $\mathcal{A}$ 为 $\mathcal{O}_Y$-代数层。用
$D_\QCoh(\mathcal{A})$ 表示 $D_\QCoh(\mathcal{O}_X)$ 在限制函子
$D(\mathcal{A}) \to D(\mathcal{O}_X)$ 下的逆像。换言之，
$K \in D(\mathcal{A})$ 属于 $D_\QCoh(\mathcal{A})$，当且仅当它的上同调层
作为 $\mathcal{O}_X$-模是拟凝聚的。若 $\mathcal{A}$ 本身拟凝聚，则这
等价于要求上同调层作为 $\mathcal{A}$-模拟凝聚，见 Morphisms, Lemma
\ref{morphisms-lemma-affine-equivalence-modules}。

\begin{lemma}
\label{perfect-lemma-affine-morphism-equivalence}
设 $f : X \to Y$ 为概形的仿射态射。那么 $f_*$ 诱导等价
$$
\Phi : D_\QCoh(\mathcal{O}_X) \longrightarrow D_\QCoh(f_*\mathcal{O}_X)
$$
，它与 $D_\QCoh(f_*\mathcal{O}_X) \to D_\QCoh(\mathcal{O}_Y)$ 的复合为
$Rf_* : D_\QCoh(\mathcal{O}_X) \to D_\QCoh(\mathcal{O}_Y)$。
\end{lemma}

\begin{proof}
回忆：对对象 $K \in D_\QCoh(\mathcal{O}_X)$，计算 $Rf_*$ 的方法是选取表示
$K$ 的 $\mathcal{O}_X$-模 K-内射复形 $\mathcal{I}^\bullet$，再取
$f_*\mathcal{I}^\bullet$。因此令 $\Phi(K)$ 为复形
$f_*\mathcal{I}^\bullet$，把它视作 $f_*\mathcal{O}_X$-模复形。用
$g : (X, \mathcal{O}_X) \to (Y, f_*\mathcal{O}_X)$ 表示显然的赋环空间
态射。那么 $g$ 是赋环空间的平坦态射（茎的描述见下文），且 $\Phi$ 是
$Rg_*$ 在 $D_\QCoh(\mathcal{O}_X)$ 上的限制。我们断言 $Lg^*$ 是拟逆。
首先注意，$Lg^*$ 把 $D_\QCoh(f_*\mathcal{O}_X)$ 映入
$D_\QCoh(\mathcal{O}_X)$，因为 $g^*$ 把拟凝聚模变成拟凝聚模（Modules,
Lemma \ref{modules-lemma-pullback-quasi-coherent}）。为完成证明，只需说明
伴随映射
$$
Lg^*\Phi(K) = Lg^*Rg_*K \to K
\quad\text{且}\quad
M \to Rg_*Lg^*M = \Phi(Lg^*M)
$$
对 $K \in D_\QCoh(\mathcal{O}_X)$ 及
$M \in D_\QCoh(f_*\mathcal{O}_X)$ 都是同构。这是局部问题，故可假设
$Y$ 仿射，进而 $X$ 也仿射。

\medskip\noindent
假设 $Y = \Spec(B)$ 且 $X = \Spec(A)$。设
$\mathfrak p = x \in \Spec(A) = X$ 为映到
$\mathfrak q = y \in \Spec(B) = Y$ 的点。那么
$(f_*\mathcal{O}_X)_y = A_\mathfrak q$ 且
$\mathcal{O}_{X, x} = A_\mathfrak p$，故 $g$ 平坦。因此 $g^*$ 正合，并且
对 $D(f_*\mathcal{O}_X)$ 中任意 $M$，有
$H^i(Lg^*M) = g^*H^i(M)$。对 $K \in D_\QCoh(\mathcal{O}_X)$，由高阶直像
的消没（Cohomology of Schemes, Lemma
\ref{coherent-lemma-relative-affine-vanishing}）以及 Lemma
\ref{perfect-lemma-application-nice-K-injective}（略去一个小细节），可知
$$
H^i(\Phi(K)) = H^i(Rf_*K) = f_*H^i(K)
$$
。因此只需说明
$$
g^*g_*\mathcal{F} \to \mathcal{F}
\quad\text{且}\quad
\mathcal{G} \to g_*g^*\mathcal{G}
$$
都是同构，其中 $\mathcal{F}$ 是拟凝聚 $\mathcal{O}_X$-模，而
$\mathcal{G}$ 是拟凝聚 $f_*\mathcal{O}_X$-模。这由 Morphisms, Lemma
\ref{morphisms-lemma-affine-equivalence-modules} 推出。
\end{proof}





\section{以闭子集为支集的上同调}
\label{perfect-section-cohomology-support}

\noindent
我们针对概形和拟凝聚模，详细阐述 Cohomology, Sections
\ref{cohomology-section-cohomology-support} 与
\ref{cohomology-section-cohomology-support-bis} 中的材料。

\begin{definition}
\label{perfect-definition-supported-on}
设 $X$ 为概形。设 $E$ 为 $D(\mathcal{O}_X)$ 的对象。设
$T \subset X$ 为闭子集。如果上同调层 $H^i(E)$ 的支集都在 $T$ 上，则称
$E$ {\it 支集在 $T$ 上}。
\end{definition}

\noindent
在该定义的情形下，我们复述 Cohomology, Section
\ref{cohomology-section-cohomology-support-bis} 中的部分讨论。设 $X$ 为
概形，设 $T \subset X$ 为闭子集。支集包含在 $T$ 中的
$\mathcal{O}_X$-模所成范畴，是所有 $\mathcal{O}_X$-模所成范畴的 Serre
子范畴，见 Homology, Definition
\ref{homology-definition-serre-subcategory} 和 Modules, Lemma
\ref{modules-lemma-support-section-closed}。因此可应用 Derived Categories,
Section \ref{derived-section-triangulated-sub} 中的讨论，得到
$D(\mathcal{O}_X)$ 的一个严格满、饱和三角子范畴，其对象都支集在 $T$
上；下文把该范畴记作 $D_T(\mathcal{O}_X)$。

\medskip\noindent
在 Definition \ref{perfect-definition-supported-on} 的情形下，用 $i : T \to X$
表示包含映射。回忆 Cohomology, Section
\ref{cohomology-section-cohomology-support-bis}：在这种情形下有函子
$R\mathcal{H}_T : D(\mathcal{O}_X) \to D(i^{-1}\mathcal{O}_X)$
，它是 $i_* : D(i^{-1}\mathcal{O}_X) \to D(\mathcal{O}_X)$ 的右伴随。

\begin{lemma}
\label{perfect-lemma-support-quasi-coherent}
设 $X$ 为概形。设 $T \subset X$ 为闭子集，使 $X \setminus T$ 是 $X$ 的
逆紧开子集。设 $i : T \to X$ 为包含映射。
\begin{enumerate}
\item 对 $D_\QCoh(\mathcal{O}_X)$ 中的 $E$，有
$i_*R\mathcal{H}_T(E)$ 属于 $D_{\QCoh, T}(\mathcal{O}_X)$。
\item 函子
$i_* \circ R\mathcal{H}_T : D_\QCoh(\mathcal{O}_X) \to
D_{\QCoh, T}(\mathcal{O}_X)$ 是包含函子
$D_{\QCoh, T}(\mathcal{O}_X) \to D_\QCoh(\mathcal{O}_X)$ 的右伴随。
\end{enumerate}
\end{lemma}

\begin{proof}
令 $U = X \setminus T$，并用 $j : U \to X$ 表示包含映射。由 Cohomology,
Lemma \ref{cohomology-lemma-triangle-sections-with-support-sheaves}，存在
可区别三角
$$
i_*R\mathcal{H}_T(E) \to E \to Rj_*(E|_U) \to i_*R\mathcal{H}_Z(E)[1]
$$
，它位于 $D(\mathcal{O}_X)$ 中。由 Lemma
\ref{perfect-lemma-quasi-coherence-direct-image}，复形 $Rj_*(E|_U)$ 有拟凝聚
上同调层（此处使用 $U$ 在 $X$ 中逆紧）。因此 (1) 成立。(2) 由此以及
Cohomology, Lemma
\ref{cohomology-lemma-complexes-with-support-on-closed} 中函子的伴随性推出。
\end{proof}

\begin{lemma}
\label{perfect-lemma-support-direct-sums}
设 $X$ 为概形。设 $T \subset X$ 为闭子集，使 $X \setminus T$ 是 $X$ 的
逆紧开子集。那么对 $D_\QCoh(\mathcal{O}_X)$ 的一族对象 $E_i$（
$i \in I$），有
$R\mathcal{H}_T(\bigoplus E_i) = \bigoplus R\mathcal{H}_T(E_i)$.
\end{lemma}

\begin{proof}
令 $U = X \setminus T$，并用 $j : U \to X$ 表示包含映射。由 Cohomology,
Lemma \ref{cohomology-lemma-triangle-sections-with-support-sheaves}，存在
可区别三角
$$
i_*R\mathcal{H}_T(E) \to E \to Rj_*(E|_U) \to i_*R\mathcal{H}_Z(E)[1]
$$
，它对 $D(\mathcal{O}_X)$ 中任意 $E$ 都位于 $D(\mathcal{O}_X)$ 中。由
Lemma \ref{perfect-lemma-quasi-coherence-pushforward-direct-sums}，函子
$E \mapsto Rj_*(E|_U)$ 在 $D_\QCoh(\mathcal{O}_X)$ 上与直和交换。由此，
函子 $i_* \circ R\mathcal{H}_T$ 也如此（略去细节）。由于
$i_* : D(i^{-1}\mathcal{O}_X) \to D_T(\mathcal{O}_X)$ 是等价
（Cohomology, Lemma
\ref{cohomology-lemma-complexes-with-support-on-closed}），结论成立。
\end{proof}

\begin{remark}
\label{perfect-remark-support-c-equations}
设 $X$ 为概形。设 $f_1, \ldots, f_c \in \Gamma(X, \mathcal{O}_X)$。用
$Z \subset X$ 表示由 $f_1, \ldots, f_c$ 截出的闭子概形。对
$0 \leq p < c$ 及 $1 \leq i_0 < \ldots < i_p \leq c$，用
$U_{i_0 \ldots i_p} \subset X$ 表示 $f_{i_0} \ldots f_{i_p}$ 可逆处的
开子概形。对任意 $\mathcal{O}_X$-模 $\mathcal{F}$，令
$$
\mathcal{F}_{i_0 \ldots i_p} =
(U_{i_0 \ldots i_p} \to X)_*(\mathcal{F}|_{U_{i_0 \ldots i_p}})
$$
在这种情形下，{\it 扩充交错 {\v C}ech 复形}是 $\mathcal{O}_X$-模复形
\begin{equation}
\label{perfect-equation-extended-alternating}
0 \to \mathcal{F} \to
\bigoplus\nolimits_{i_0} \mathcal{F}_{i_0} \to
\ldots \to
\bigoplus\nolimits_{i_0 < \ldots < i_p} \mathcal{F}_{i_0 \ldots i_p} \to
\ldots \to \mathcal{F}_{1 \ldots c} \to 0
\end{equation}
，其中 $\mathcal{F}$ 放在次数 $0$ 上。映射构造如下。给定
$1 \leq i_0 < \ldots < i_{p + 1} \leq c$ 及 $0 \leq j \leq p + 1$，
有典范映射
$$
\mathcal{F}_{i_0 \ldots \hat i_j \ldots i_{p + 1}} \to
\mathcal{F}_{i_0 \ldots i_p}
$$
，它来自包含
$U_{i_0 \ldots i_p} \subset U_{i_0 \ldots \hat i_j \ldots i_{p + 1}}$.
扩充交错复形的微分使用这些典范映射，并带符号 $(-1)^j$。
\end{remark}

\begin{lemma}
\label{perfect-lemma-extended-alternating-zero}
设 $X$、$f_1, \ldots, f_c \in \Gamma(X, \mathcal{O}_X)$ 及
$\mathcal{F}$ 如 Remark \ref{perfect-remark-support-c-equations}。复形
(\ref{perfect-equation-extended-alternating}) 在 $X \setminus Z$ 上的限制无环。
\end{lemma}

\noindent
指出：更一般地，对从有限开覆盖
$X \setminus Z = U_1 \cup \ldots \cup U_c$ 出发、按 Remark
\ref{perfect-remark-support-c-equations} 定义的任意扩充交错 {\v C}ech 复形，本
引理都成立。

\begin{proof}
设 $W \subset X \setminus Z$ 为开子集。在 $W$ 上取层复形
(\ref{perfect-equation-extended-alternating}) 的截面，得到复形
$$
\mathcal{F}(W) \to \bigoplus\nolimits_{i_0} \mathcal{F}(U_{i_0} \cap W) \to
\bigoplus\nolimits_{i_0 < i_1} \mathcal{F}(U_{i_0i_1} \cap W) \to \ldots
$$
换言之，得到覆盖 $W = \bigcup U_i \cap W$ 及 $\{1, \ldots, c\}$ 上标准
次序的扩充有序 {\v C}ech 复形，见 Cohomology, Section
\ref{cohomology-section-alternating-cech}。由 Cohomology, Lemma
\ref{cohomology-lemma-alternating-cech-trivial}，只要对某个
$1 \leq i \leq c$，$W$ 包含在 $V(f_i)$ 中，该复形就同伦于零。证毕。
\end{proof}

\begin{remark}
\label{perfect-remark-extended-alternating-map-to-support}
设 $X$、$f_1, \ldots, f_c \in \Gamma(X, \mathcal{O}_X)$ 及
$\mathcal{F}$ 如 Remark \ref{perfect-remark-support-c-equations}。用
$\mathcal{F}^\bullet$ 表示复形
(\ref{perfect-equation-extended-alternating})。由 Lemma
\ref{perfect-lemma-extended-alternating-zero}，$\mathcal{F}^\bullet$ 的上同调层
支集在 $Z$ 上，因而 $\mathcal{F}^\bullet$ 是 $D_Z(\mathcal{O}_X)$ 的
对象。另一方面，等式 $\mathcal{F}^0 = \mathcal{F}$ 决定
$D(\mathcal{O}_X)$ 中的典范映射
$\mathcal{F}^\bullet \to \mathcal{F}$。由于
$i_* \circ R\mathcal{H}_Z$ 是包含函子
$D_Z(\mathcal{O}_X) \to D(\mathcal{O}_X)$ 的右伴随（见 Cohomology,
Lemma \ref{cohomology-lemma-complexes-with-support-on-closed}），得到典范
交换图
$$
\xymatrix{
\mathcal{F}^\bullet \ar[rd] \ar[rr] & & \mathcal{F} \\
& i_*R\mathcal{H}_Z(\mathcal{F}) \ar[ru]
}
$$
，它位于 $D(\mathcal{O}_X)$ 中，并且关于 $\mathcal{O}_X$-模
$\mathcal{F}$ 函子性成立。
\end{remark}

\begin{lemma}
\label{perfect-lemma-extended-alternating-represented}
设 $X$、$f_1, \ldots, f_c \in \Gamma(X, \mathcal{O}_X)$ 及
$\mathcal{F}$ 如 Remark \ref{perfect-remark-support-c-equations}。若
$\mathcal{F}$ 拟凝聚，则复形 (\ref{perfect-equation-extended-alternating}) 在
$D_Z(\mathcal{O}_X)$ 中表示 $i_* R\mathcal{H}_Z(\mathcal{F})$。
\end{lemma}

\begin{proof}
用 $\mathcal{F}^\bullet$ 表示复形
(\ref{perfect-equation-extended-alternating})。本引理的陈述是指 Remark
\ref{perfect-remark-extended-alternating-map-to-support} 中的映射
$\mathcal{F}^\bullet \to i_*R\mathcal{H}_Z(\mathcal{F})$ 为同构。因为
$\mathcal{F}^\bullet$ 属于 $D_Z(\mathcal{O}_X)$（见所引评注），由
Cohomology, Lemma
\ref{cohomology-lemma-complexes-with-support-on-closed}，有
$i_*R\mathcal{H}_Z(\mathcal{F}^\bullet) = \mathcal{F}^\bullet$。态射
$U_{i_0 \ldots i_p} \to X$ 仿射，因为在 $X$ 的仿射开集上，它由把函数
$f_{i_0} \ldots f_{i_p}$ 取逆给出。因此
$$
\mathcal{F}_{i_0 \ldots i_p} =
(U_{i_0 \ldots i_p} \to X)_*\mathcal{F}|_{U_{i_0 \ldots i_p}} =
R(U_{i_0 \ldots i_p} \to X)_*\mathcal{F}|_{U_{i_0 \ldots i_p}}
$$
，这由 Cohomology of Schemes, Lemma
\ref{coherent-lemma-relative-affine-vanishing} 及 $\mathcal{F}$ 拟凝聚的假设
可知。由 Cohomology, Lemma
\ref{cohomology-lemma-sections-support-in-closed-disjoint-open}，可知
$R\mathcal{H}_Z(\mathcal{F}_{i_0 \ldots i_p}) = 0$。因此对 $p > 0$，有
$i_*R\mathcal{H}_Z(\mathcal{F}^p) = 0$。综合以上各点，得到
$$
\mathcal{F}^\bullet = i_*R\mathcal{H}_Z(\mathcal{F}^\bullet) =
i_*R\mathcal{H}_Z(\mathcal{F})
$$
，正合所需。
\end{proof}

\begin{lemma}
\label{perfect-lemma-supported-trivial-vanishing}
设 $X$ 为概形。设 $T \subset X$ 为闭子集，它局部可由结构层的至多 $c$
个元素截出。那么对 $i > c$ 及任意拟凝聚 $\mathcal{O}_X$-模
$\mathcal{F}$，有 $\mathcal{H}^i_Z(\mathcal{F}) = 0$。
\end{lemma}

\begin{proof}
这立即来自 Lemma \ref{perfect-lemma-extended-alternating-represented} 给出的
$R\mathcal{H}_T(\mathcal{F})$ 的局部描述。
\end{proof}

\begin{lemma}
\label{perfect-lemma-supported-vanishing}
设 $X$ 为概形。设 $Z \subset X$ 为闭子集，它局部可由含 $c$ 个元素的
Koszul 正则序列截出。那么对每个平坦拟凝聚 $\mathcal{O}_X$-模
$\mathcal{F}$，当 $i \not = c$ 时有
$\mathcal{H}^i_Z(\mathcal{F}) = 0$。
\end{lemma}

\begin{proof}
由 Lemma \ref{perfect-lemma-extended-alternating-represented} 给出的
$R\mathcal{H}_Z(\mathcal{F})$ 的描述，问题归结为如下代数陈述：给定环
$R$、Koszul 正则序列 $f_1, \ldots, f_c \in R$ 以及平坦 $R$-模 $M$，
扩充交错 {\v C}ech 复形
$M \to \bigoplus\nolimits_{i_0} M_{f_{i_0}} \to
\bigoplus\nolimits_{i_0 < i_1} M_{f_{i_0}f_{i_1}} \to
\ldots \to M_{f_1 \ldots f_c}$
（来自 More on Algebra, Section
\ref{more-algebra-section-alternating-cech}）仅在次数 $c$ 上有上同调。由
More on Algebra, Lemma
\ref{more-algebra-lemma-vanishing-extended-alternating-koszul}，得到 $R$ 的
扩充交错 {\v C}ech 复形所需的消没。由于 $M$ 的复形是把它与平坦
$R$-模 $M$ 作张量所得（More on Algebra, Lemma
\ref{more-algebra-lemma-extended-alternating-form-module}），结论成立。
\end{proof}

\begin{remark}
\label{perfect-remark-supported-map-c-equations}
设 $X$、$f_1, \ldots, f_c \in \Gamma(X, \mathcal{O}_X)$ 及
$\mathcal{F}$ 如 Remark \ref{perfect-remark-support-c-equations}。存在典范的
$\mathcal{O}_X|_Z$-线性映射
$$
c_{f_1, \ldots, f_c} :
i^*\mathcal{F}
\longrightarrow
\mathcal{H}^c_Z(\mathcal{F})
$$
，它关于 $\mathcal{F}$ 函子性成立。事实上，用 $\mathcal{F}^\bullet$
表示扩充交错 {\v C}ech 复形 (\ref{perfect-equation-extended-alternating})，则有
Remark \ref{perfect-remark-extended-alternating-map-to-support} 中的典范映射
$\mathcal{F}^\bullet \to i_*R\mathcal{H}_Z(\mathcal{F})$
。这在次数 $c$ 的上同调层上决定典范映射
$$
\Coker\left(\bigoplus \mathcal{F}_{1 \ldots \hat i \ldots c} \to
\mathcal{F}_{1 \ldots c}\right)
\longrightarrow
i_*\mathcal{H}^c_Z(\mathcal{F})
$$
。给定 $\mathcal{F}$ 的局部截面 $s$，可以考虑
$\mathcal{F}_{1 \ldots c}$ 的局部截面
$$
\frac{s}{f_1 \ldots f_c}
$$
。该截面在上述余核中的类，只依赖于 $s$ 模去映射
$(f_1, \ldots, f_c) : \mathcal{F}^{\oplus c} \to \mathcal{F}$.
的像。由于 $i_*i^*\mathcal{F}$ 等于
$(f_1, \ldots, f_c) : \mathcal{F}^{\oplus c} \to \mathcal{F}$
的余核，故得到 $\mathcal{O}_X$-模映射
$i_*i^*\mathcal{F} \to i_*\mathcal{H}_Z^c(\mathcal{F})$.
因为 $i_*$ 全忠实，得到映射 $c_{f_1, \ldots, f_c}$。
\end{remark}

\begin{example}
\label{perfect-example-affine-supported-map-c-equations}
设 $X = \Spec(A)$ 仿射，$f_1, \ldots, f_c \in A$，并设
$\mathcal{F} = \widetilde{M}$，其中 $M$ 为某个 $A$-模。Remark
\ref{perfect-remark-supported-map-c-equations} 中的映射
$c_{f_1, \ldots, f_c}$ 可描述为映射
$$
M/(f_1, \ldots, f_c)M
\longrightarrow
\Coker\left(
\bigoplus M_{f_1 \ldots \hat f_i \ldots f_c} \to M_{f_1 \ldots f_c}
\right)
$$
，它把 $s \in M$ 的类映到余核中 $s/f_1 \ldots f_c$ 的类。
\end{example}

\begin{lemma}
\label{perfect-lemma-supported-map-determinant}
设 $X$、$f_1, \ldots, f_c \in \Gamma(X, \mathcal{O}_X)$ 及
$\mathcal{F}$ 如 Remark \ref{perfect-remark-support-c-equations}。对
$1 \leq i, j \leq c$，设
$a_{ji} \in \Gamma(X, \mathcal{O}_X)$，并令
$g_j = \sum_{i = 1, \ldots, c} a_{ji}f_i$。假设 $g_1, \ldots, g_c$ 在概形
意义下截出 $Z$。若 $\mathcal{F}$ 拟凝聚，则
$$
c_{f_1, \ldots, f_c} = \det(a_{ji}) c_{g_1, \ldots, g_c}
$$
，其中 $c_{f_1, \ldots, f_c}$ 与 $c_{g_1, \ldots, g_c}$ 如 Remark
\ref{perfect-remark-supported-map-c-equations} 所述。
\end{lemma}

\begin{proof}
我们将证明：对任意 $s \in \mathcal{F}(X)$，作为
$\mathcal{H}_Z(\mathcal{F})$ 的整体截面，有
$c_{f_1, \ldots, f_c}(s) =
\det(a_{ij}) c_{g_1, \ldots, g_c}(s)$。这已经足够，因为随后对 $X$ 的任意
开子概形上的截面也得到同样结果。为此，对
$1 \leq i_0 < \ldots < i_p \leq c$ 及
$1 \leq j_0 < \ldots < j_q \leq c$，用
$U_{i_0 \ldots i_p} \subset X$,
$V_{j_0 \ldots j_q} \subset X$, and
$W_{i_0 \ldots i_p, j_0 \ldots j_q} \subset X$
分别表示 $f_{i_0} \ldots f_{i_p}$ 可逆处、
$g_{j_0} \ldots g_{j_q}$ 可逆处以及
$f_{i_0} \ldots f_{i_p}g_{j_0} \ldots g_{j_q}$ 可逆处的开子概形。用
$\mathcal{F}_{i_0 \ldots i_p}$、$\mathcal{F}'_{j_0 \ldots j_q}$、
$\mathcal{F}''_{i_0 \ldots i_p, j_0 \ldots j_q}$ 分别表示 $\mathcal{F}$
在 $U_{i_0 \ldots i_p}$、$V_{j_0 \ldots j_q}$、
$W_{i_0 \ldots i_p, j_0 \ldots j_q}$ 上的限制到 $X$ 的直像。于是得到三个
扩充交错 {\v C}ech 复形
$$
\mathcal{F}^\bullet :
\mathcal{F} \to \bigoplus\nolimits_{i_0} \mathcal{F}_{i_0} \to
\bigoplus\nolimits_{i_0 < i_1} \mathcal{F}_{i_0i_1} \to \ldots
$$
以及
$$
(\mathcal{F}')^\bullet :
\mathcal{F} \to \bigoplus\nolimits_{j_0} \mathcal{F}'_{j_0} \to
\bigoplus\nolimits_{j_0 < j_1} \mathcal{F}'_{j_0j_1} \to \ldots
$$
以及
$$
(\mathcal{F}'')^\bullet :
\mathcal{F} \to
\bigoplus\nolimits_{i_0} \mathcal{F}_{i_0} \oplus
\bigoplus\nolimits_{j_0} \mathcal{F}'_{j_0} \to
\bigoplus\nolimits_{i_0 < i_1} \mathcal{F}_{i_0i_1} \oplus
\bigoplus\nolimits_{i_0, j_0} \mathcal{F}''_{i_0, j_0} \oplus
\bigoplus\nolimits_{j_0 < j_1} \mathcal{F}'_{j_0j_1} \to
\ldots
$$
，其微分就是定义 (\ref{perfect-equation-extended-alternating}) 时所用的微分。存在
复形映射
$$
(\mathcal{F}'')^\bullet \to \mathcal{F}^\bullet
\quad\text{且}\quad
(\mathcal{F}'')^\bullet \to (\mathcal{F}')^\bullet
$$
，它们由各项上的投影映射给出（因而在次数 $0$ 上诱导恒等映射）。注意，由
Lemma \ref{perfect-lemma-extended-alternating-represented}，这些复形中的每一个都
表示 $i_*R\mathcal{H}_Z(\mathcal{F})$，而这些映射表示该对象上的恒等。
因此只需找到元素
$$
\sigma \in H^c((\mathcal{F}'')^\bullet(X))
$$
，使它在这两个映射下分别映到 $c_{f_1, \ldots, f_c}(s)$ 和
$\det(a_{ji})c_{g_1, \ldots, g_c}(s)$。事实上，可以显式给出表示
$\sigma$ 的余圈。具体地，取
$$
\sigma_{1 \ldots c} = \frac{s}{f_1 \ldots f_c} \in \mathcal{F}_{1 \ldots c}(X)
\quad\text{且}\quad
\sigma'_{1 \ldots c} = \frac{\det(a_{ji})s}{g_1 \ldots g_c} \in
\mathcal{F}'_{1 \ldots c}(X)
$$
，并取
$$
\sigma_{i_0 \ldots i_p, j_0 \ldots j_{c - p - 2}} =
\frac{\lambda(i_0 \ldots i_p, j_0 \ldots j_{c - p - 2})s}%
{f_{i_0} \ldots f_{i_p}g_{j_0} \ldots g_{j_{c - p - 2}}}
\in \mathcal{F}''_{i_0 \ldots i_p, j_0 \ldots j_{c - p - 2}}(X)
$$
，其中 $\lambda(i_0 \ldots i_p, j_0 \ldots j_{c - p - 2})$ 是形式表达式
$$
e_{i_0} \wedge \ldots \wedge e_{i_p} \wedge
(a_{j_01} e_1 + \ldots + a_{j_0c}e_c) \wedge \ldots \wedge
(a_{j_{c - p - 2}1} e_1 + \ldots + a_{j_{c - p - 2}c}e_c)
$$
中 $e_1 \wedge \ldots \wedge e_c$ 的系数。为验证 $\sigma$ 是余圈，必须
证明对 $1 \leq i_0 < \ldots < i_p \leq c$ 及
$1 \leq j_0 < \ldots < j_{c - p - 1} \leq c$，有
\begin{align*}
0 & =
\sum\nolimits_{a = 0, \ldots, p} (-1)^a
f_{i_a} \lambda(i_0 \ldots \hat i_a \ldots i_p, j_0 \ldots j_{c - p - 1}) \\
& +
\sum\nolimits_{b = 0, \ldots, c - p - 1} (-1)^{p + b + 1}g_{j_b}
\lambda(i_0 \ldots i_p, j_0 \ldots \hat j_b \ldots j_{c - p - 1})
\end{align*}
最容易的验证方法或许是指出，形式表达式
$$
\xi = e_{i_0} \wedge \ldots \wedge e_{i_p} \wedge
(a_{j_01} e_1 + \ldots + a_{j_0c}e_c) \wedge \ldots \wedge
(a_{j_{c - p - 1}1} e_1 + \ldots + a_{j_{c - p - 1}c}e_c)
$$
为 $0$，因为它是以 $e_1, \ldots, e_c$ 为基的自由模的第 $(c + 1)$ 个
外幂中的元素；并且上面的表达式是 $\xi$ 在把 $e_i \to f_i$ 的 Koszul
微分下的像。略去一些细节。
\end{proof}

\begin{lemma}
\label{perfect-lemma-supported-map-global}
设 $X$ 为概形。设 $Z \to X$ 为有限表示闭浸入，并且其余法层
$\mathcal{C}_{Z/X}$ 是秩为 $c$ 的局部自由层。则有典范映射
$$
c :
\wedge^c(\mathcal{C}_{Z/X})^\vee \otimes_{\mathcal{O}_Z} i^*\mathcal{F}
\longrightarrow
\mathcal{H}_Z^c(\mathcal{F})
$$
它关于拟凝聚模 $\mathcal{F}$ 函子性地成立。
\end{lemma}

\begin{proof}
这由注 \ref{perfect-remark-supported-map-c-equations} 中的构造，以及引理
\ref{perfect-lemma-supported-map-determinant} 所证明的、对理想层生成元选取的
无关性得到。略去一些细节。
\end{proof}

\begin{remark}
\label{perfect-remark-supported-functorial}
设 $g : X' \to X$ 为概形态射，并设
$f_1, \ldots, f_c \in \Gamma(X, \mathcal{O}_X)$。令
$f'_i = g^\sharp(f_i) \in \Gamma(X', \mathcal{O}_{X'})$。分别以
$Z \subset X$ 和 $Z' \subset X'$ 表示由 $f_1, \ldots, f_c$ 和
$f'_1, \ldots, f'_c$ 截出的闭子概形。于是
$Z' = Z \times_X X'$。以 $h : Z' \to Z$ 表示所诱导的概形态射。
设 $\mathcal{F}$ 为 $\mathcal{O}_X$-模，并令
$\mathcal{F}' = g^*\mathcal{F}$。在此情形下，若 $\mathcal{F}$ 拟凝聚，
则图表
$$
\xymatrix{
(i')^{-1}\mathcal{O}_{X'} \otimes_{h^{-1}i^{-1}\mathcal{O}_X}
h^{-1}\mathcal{H}^c_Z(\mathcal{F}) \ar[r] &
\mathcal{H}_{Z'}^c(\mathcal{F}') \\
h^*i^*\mathcal{F} \ar[r] \ar[u]_-{c_{f_1, \ldots, f_c}} &
(i')^*\mathcal{F}' \ar[u]^-{c_{f'_1, \ldots, f'_c}}
}
$$
交换，其中上方水平箭头是在 $c$ 次上同调层上取 Cohomology, Remark
\ref{cohomology-remark-support-functorial} 中的映射。具体而言，分别以
$\mathcal{F}^\bullet$ 和 $(\mathcal{F}')^\bullet$ 表示在注
\ref{perfect-remark-support-c-equations} 中用
$\mathcal{F}, f_1, \ldots, f_c$ 和
$\mathcal{F}', f'_1, \ldots, f'_c$ 构造的扩充交错 {\v C}ech 复形。
注意 $(\mathcal{F}')^\bullet = g^*\mathcal{F}^\bullet$。那么，即使不假设
$\mathcal{F}$ 拟凝聚，图表
$$
\xymatrix{
i'_* L(g|_{Z'})^* R\mathcal{H}_Z(\mathcal{F}) \ar[r] \ar@{=}[d] &
i'_*R\mathcal{H}_{Z'}(Lg^*\mathcal{F}) \ar[d] \\
Lg^*i_*R\mathcal{H}_Z(\mathcal{F}) &
i'_*R\mathcal{H}_{Z'}(\mathcal{F}') \\
Lg^*(\mathcal{F}^\bullet) \ar[u] \ar[r] &
(\mathcal{F}')^\bullet \ar[u]
}
$$
也交换，其中 $g|_{Z'} : (Z', (i')^{-1}\mathcal{O}_{X'}) \to
(Z, i^{-1}\mathcal{O}_X)$ 是所诱导的带环空间态射。这里上方水平箭头
由 Cohomology, Remark \ref{cohomology-remark-support-functorial} 给出，
等号也在那里得到解释。向上的箭头来自注
\ref{perfect-remark-extended-alternating-map-to-support}。下方水平箭头是映射
$Lg^*\mathcal{F}^\bullet
\to g^*\mathcal{F}^\bullet = (\mathcal{F}')^\bullet$，而向下的箭头由
$Lg^*\mathcal{F} \to g^*\mathcal{F} = \mathcal{F}'$ 诱导。此图表交换，
因为沿两条路径绕行图表所得的两个箭头
$Lg^*\mathcal{F}^\bullet \to
i'_*R\mathcal{H}_{Z'}(\mathcal{F}')$，在与
$i'_*R\mathcal{H}_{Z'}(\mathcal{F}') \to \mathcal{F}'$ 复合后，均为
典范映射 $Lg^*\mathcal{F}^\bullet \to \mathcal{F}'$。略去一些细节。
现在考察此图表在 $c$ 次上同调层上的情形，并利用注
\ref{perfect-remark-supported-map-c-equations} 中构造的映射
$i^*\mathcal{F} \to
\Coker(\bigoplus \mathcal{F}_{1 \ldots \hat i \ldots c} \to
\mathcal{F}_{1 \ldots c})$
与拉回相容，即得第一个图表的交换性。
\end{remark}








\section{拟凝聚化子}
\label{perfect-section-coherator}

\noindent
设 $X$ 为概形。所谓 {\it coherator}（拟凝聚化子）是函子
$$
Q_X :
\textit{Mod}(\mathcal{O}_X)
\longrightarrow
\QCoh(\mathcal{O}_X)
$$
它是包含函子
$\QCoh(\mathcal{O}_X) \to \textit{Mod}(\mathcal{O}_X)$ 的右伴随。
它对任意概形 $X$ 都存在；此外，对每个拟凝聚模 $\mathcal{F}$，伴随映射
$Q_X(\mathcal{F}) \to \mathcal{F}$ 都是同构，见 Properties,
Proposition \ref{properties-proposition-coherator}。由于 $Q_X$ 左正合（因为
它是右伴随），可以考虑其右导出延拓
$$
RQ_X :
D(\mathcal{O}_X)
\longrightarrow
D(\QCoh(\mathcal{O}_X)).
$$
由于 $Q_X$ 是包含函子
$\QCoh(\mathcal{O}_X) \to \textit{Mod}(\mathcal{O}_X)$ 的右伴随，
由 Derived Categories, Lemma
\ref{derived-lemma-derived-adjoint-functors} 可知，$RQ_X$ 是典范函子
$D(\QCoh(\mathcal{O}_X)) \to D(\mathcal{O}_X)$ 的右伴随。

\medskip\noindent
本节研究函子 $RQ_X$。在第 \ref{perfect-section-better-coherator} 节中，我们将研究
包含函子 $D_\QCoh(\mathcal{O}_X) \to D(\mathcal{O}_X)$ 的右伴随（当它
存在时）；二者密切相关。

\begin{lemma}
\label{perfect-lemma-affine-pushforward}
设 $f : X \to Y$ 为概形的仿射态射。则 $f_*$ 定义导出函子
$f_* : D(\QCoh(\mathcal{O}_X)) \to D(\QCoh(\mathcal{O}_Y))$，并且图表
$$
\xymatrix{
D(\QCoh(\mathcal{O}_X)) \ar[d]_{f_*} \ar[r] &
D_\QCoh(\mathcal{O}_X) \ar[d]^{Rf_*} \\
D(\QCoh(\mathcal{O}_Y)) \ar[r] &
D_\QCoh(\mathcal{O}_Y)
}
$$
交换。
\end{lemma}

\begin{proof}
函子
$f_* : \QCoh(\mathcal{O}_X) \to \QCoh(\mathcal{O}_Y)$
正合，见 Cohomology of Schemes, Lemma
\ref{coherent-lemma-relative-affine-vanishing}。因此 $f_*$ 定义如下导出函子；
具体只需把 $f_*$ 作用于任意代表复形：
$f_* : D(\QCoh(\mathcal{O}_X)) \to D(\QCoh(\mathcal{O}_Y))$
，见 Derived Categories, Lemma
\ref{derived-lemma-right-derived-exact-functor}。图表由引理
\ref{perfect-lemma-pushforward-affine-morphism} 交换。
\end{proof}

\begin{lemma}
\label{perfect-lemma-flat-pushforward-coherator}
设 $f : X \to Y$ 为概形态射，并假设 $f$ 拟紧、拟分离且平坦。以
$$
\Phi : D(\QCoh(\mathcal{O}_X)) \to D(\QCoh(\mathcal{O}_Y))
$$
表示 $f_* : \QCoh(\mathcal{O}_X) \to \QCoh(\mathcal{O}_Y)$ 的右导出函子，
则有 $RQ_Y \circ Rf_* = \Phi \circ RQ_X$。
\end{lemma}

\begin{proof}
我们将证明 $RQ_Y \circ Rf_*$ 和 $\Phi \circ RQ_X$ 都是同一函子
$D(\QCoh(\mathcal{O}_Y)) \to D(\mathcal{O}_X)$ 的右伴随。

\medskip\noindent
由于 $f$ 拟紧且拟分离，$f_*$ 保持拟凝聚性，见 Schemes, Lemma
\ref{schemes-lemma-push-forward-quasi-coherent}。回忆
$\QCoh(\mathcal{O}_X)$ 是 Grothendieck 阿贝尔范畴（Properties,
Proposition \ref{properties-proposition-coherator}）。因此
$D(\QCoh(\mathcal{O}_X))$ 中任意 $K$ 都可由
$\QCoh(\mathcal{O}_X)$ 中的 K-内射复形 $\mathcal{I}^\bullet$ 表示，见
Injectives, Theorem
\ref{injectives-theorem-K-injective-embedding-grothendieck}。于是可定义
$\Phi(K) = f_*\mathcal{I}^\bullet$。

\medskip\noindent
由于 $f$ 平坦，函子 $f^*$ 正合。因此 $f^*$ 同时定义
$f^* : D(\mathcal{O}_Y) \to D(\mathcal{O}_X)$ 和
$f^* : D(\QCoh(\mathcal{O}_Y)) \to D(\QCoh(\mathcal{O}_X))$。函子
$f^* = Lf^* : D(\mathcal{O}_Y) \to D(\mathcal{O}_X)$ 是
$Rf_* : D(\mathcal{O}_X) \to D(\mathcal{O}_Y)$ 的左伴随，见
Cohomology, Lemma \ref{cohomology-lemma-adjoint}。同理，函子
$f^* : D(\QCoh(\mathcal{O}_Y)) \to D(\QCoh(\mathcal{O}_X))$
是 $\Phi : D(\QCoh(\mathcal{O}_X)) \to D(\QCoh(\mathcal{O}_Y))$ 的左伴随，
这由 Derived Categories, Lemma
\ref{derived-lemma-derived-adjoint-functors} 得到。

\medskip\noindent
设 $A$ 为 $D(\QCoh(\mathcal{O}_Y))$ 的对象，而 $E$ 为
$D(\mathcal{O}_X)$ 的对象。于是
\begin{align*}
\Hom_{D(\QCoh(\mathcal{O}_Y))}(A, RQ_Y(Rf_*E))
& =
\Hom_{D(\mathcal{O}_Y)}(A, Rf_*E) \\
& =
\Hom_{D(\mathcal{O}_X)}(f^*A, E) \\
& =
\Hom_{D(\QCoh(\mathcal{O}_X))}(f^*A, RQ_X(E)) \\
& =
\Hom_{D(\QCoh(\mathcal{O}_Y))}(A, \Phi(RQ_X(E)))
\end{align*}
这便给出所需结论。
\end{proof}

\begin{lemma}
\label{perfect-lemma-affine-coherator}
设 $X = \Spec(A)$ 为仿射概形。则
\begin{enumerate}
\item $Q_X : \textit{Mod}(\mathcal{O}_X) \to \QCoh(\mathcal{O}_X)$
是把 $\mathcal{F}$ 送到与 $A$-模 $\Gamma(X, \mathcal{F})$ 相伴的拟凝聚
$\mathcal{O}_X$-模的函子；
\item $RQ_X : D(\mathcal{O}_X) \to D(\QCoh(\mathcal{O}_X))$
是把 $E$ 送到与 $D(A)$ 的对象 $R\Gamma(X, E)$ 相伴的拟凝聚
$\mathcal{O}_X$-模复形的函子；
\item 函子 $RQ_X$ 限制在 $D_\QCoh(\mathcal{O}_X)$ 上时，给出
(\ref{perfect-equation-compare}) 的拟逆。
\end{enumerate}
\end{lemma}

\begin{proof}
由 Schemes, Lemma \ref{schemes-lemma-compare-constructions}，函子 $Q_X$ 是
$$
\mathcal{F} \mapsto \widetilde{\Gamma(X, \mathcal{F})}
$$
这立即推出 (1) 和 (2)。第三个断言由引理
\ref{perfect-lemma-affine-compare-bounded}（的证明）得到。
\end{proof}

\noindent
现在可以证明函子
$D(\QCoh(\mathcal{O}_X)) \to D_\QCoh(\mathcal{O}_X)$ 何时为等价的一个判据。

\begin{lemma}
\label{perfect-lemma-argument-proves}
设 $X$ 为拟紧拟分离概形。假设对每个仿射开集 $U \subset X$，左正合函子
$j_* : \QCoh(\mathcal{O}_U) \to \QCoh(\mathcal{O}_X)$ 的右导出函子
$$
\Phi : D(\QCoh(\mathcal{O}_U)) \to D(\QCoh(\mathcal{O}_X))
$$
可置于交换图表
$$
\xymatrix{
D(\QCoh(\mathcal{O}_U)) \ar[d]_\Phi \ar[r]_{i_U} &
D_\QCoh(\mathcal{O}_U) \ar[d]^{Rj_*} \\
D(\QCoh(\mathcal{O}_X)) \ar[r]^{i_X} &
D_\QCoh(\mathcal{O}_X)
}
$$
中。则函子 (\ref{perfect-equation-compare})
$$
D(\QCoh(\mathcal{O}_X))
\longrightarrow
D_\QCoh(\mathcal{O}_X)
$$
是等价，且其拟逆由 $RQ_X$ 给出。
\end{lemma}

\begin{proof}
设 $E$ 为 $D_\QCoh(\mathcal{O}_X)$ 的对象，$A$ 为
$D(\QCoh(\mathcal{O}_X))$ 的对象。我们必须证明伴随映射
$$
A \to RQ_X(i_X(A))
\quad\text{且}\quad
i_X(RQ_X(E)) \to E
$$
是同构。考虑命题 $H_n$：只要 $E$ 和 $i_X(A)$ 支持在 $X$ 的某个闭子集
上（定义 \ref{perfect-definition-supported-on}），且该闭子集包含在 $X$ 的 $n$ 个
仿射开集之并中，上述伴随映射就是同构。我们对 $n$ 归纳证明 $H_n$。

\medskip\noindent
初始情形：$n = 0$。此时 $E = 0$，故映射
$i_X(RQ_X(E)) \to E$ 是同构。同样，$i_X(A) = 0$。因此 $i_X(A)$ 的
上同调层为零。由于包含函子
$\QCoh(\mathcal{O}_X) \to \textit{Mod}(\mathcal{O}_X)$ 完全忠实且正合，
可知 $A$ 的上同调对象为零，即 $A = 0$，从而
$A \to RQ_X(i_X(A))$ 也是同构。

\medskip\noindent
归纳步骤。假设 $E$ 和 $i_X(A)$ 支持在 $X$ 的闭子集 $T$ 上，而该闭子集
包含在 $U_1 \cup \ldots \cup U_n$ 中，其中 $U_i \subset X$ 为
仿射开集。令 $U = U_n$。考虑可区别三角
$$
A \to \Phi(A|_U) \to A' \to A[1]
\quad\text{且}\quad
E \to Rj_*(E|_U) \to E' \to E[1]
$$
其中 $\Phi$ 如引理陈述所定义。注意，$E \to Rj_*(E|_U)$ 在
$U = U_n$ 上是拟同构。由假设 $i_X \circ \Phi = Rj_* \circ i_U$，且
$i_X(A)|_U = i_U(A|_U)$，可知 $i_X(A) \to i_X(\Phi(A|_U))$ 在 $U$
上是拟同构。因此 $i_X(A')$ 和 $E'$ 支持在 $X$ 的闭子集
$T \setminus U$ 上，而该闭子集包含在
$U_1 \cup \ldots \cup U_{n - 1}$ 中。由归纳假设，断言对 $A'$ 和 $E'$
成立。根据 Derived Categories, Lemma
\ref{derived-lemma-third-isomorphism-triangle}，只需证明映射
$$
\Phi(A|_U) \to RQ_X(i_X(\Phi(A|_U)))
\quad\text{且}\quad
i_X(RQ_X(Rj_*E|_U)) \to Rj_*(E|_U)
$$
是同构。由假设和引理 \ref{perfect-lemma-flat-pushforward-coherator}（包含态射
$j : U \to X$ 平坦、拟紧且拟分离），有
$$
RQ_X(i_X(\Phi(A|_U))) = RQ_X(Rj_*(i_U(A|_U))) = \Phi(RQ_U(i_U(A|_U)))
$$
and
$$
i_X(RQ_X(Rj_*(E|_U))) = i_X(\Phi(RQ_U(E|_U))) = Rj_*(i_U(RQ_U(E|_U)))
$$
最后，映射
$$
A|_U \to RQ_U(i_U(A|_U))
\quad\text{且}\quad
i_U(RQ_U(E|_U)) \to E|_U
$$
由引理 \ref{perfect-lemma-affine-coherator} 是同构。结论得证。
\end{proof}

\begin{proposition}
\label{perfect-proposition-quasi-compact-affine-diagonal}
设 $X$ 为对角态射仿射的拟紧概形。则函子 (\ref{perfect-equation-compare})
$$
D(\QCoh(\mathcal{O}_X))
\longrightarrow
D_\QCoh(\mathcal{O}_X)
$$
是等价，且其拟逆由 $RQ_X$ 给出。
\end{proposition}

\begin{proof}
设 $U \subset X$ 为仿射开集。由 Morphisms, Lemma
\ref{morphisms-lemma-affine-permanence}，态射 $U \to X$ 是仿射的。因此引理
\ref{perfect-lemma-argument-proves} 的假设由引理 \ref{perfect-lemma-affine-pushforward} 保证，
结论随即得到。
\end{proof}

\begin{lemma}
\label{perfect-lemma-direct-image-coherator}
设 $f : X \to Y$ 为概形态射。假设 $X$ 和 $Y$ 拟紧且具有仿射对角态射。
以
$$
\Phi : D(\QCoh(\mathcal{O}_X)) \to D(\QCoh(\mathcal{O}_Y))
$$
表示 $f_* : \QCoh(\mathcal{O}_X) \to \QCoh(\mathcal{O}_Y)$ 的右导出函子，
则图表
$$
\xymatrix{
D(\QCoh(\mathcal{O}_X)) \ar[d]_\Phi \ar[r] &
D_\QCoh(\mathcal{O}_X) \ar[d]^{Rf_*} \\
D(\QCoh(\mathcal{O}_Y)) \ar[r] &
D_\QCoh(\mathcal{O}_Y)
}
$$
交换。
\end{lemma}

\begin{proof}
注意，由命题 \ref{perfect-proposition-quasi-compact-affine-diagonal}，图表中的水平
箭头均为范畴等价。因此可把这些范畴等同（对其他具有仿射对角态射的拟紧
概形也同样处理）。引理的断言是：对 $D(\QCoh(\mathcal{O}_X))$ 中每个
$K$，典范映射 $\Phi(K) \to Rf_*(K)$ 都是同构。注意，若
$K_1 \to K_2 \to K_3 \to K_1[1]$ 是 $D(\QCoh(\mathcal{O}_X))$ 中的
可区别三角，且断言对其中两个对象成立，则对第三个也成立。

\medskip\noindent
设 $U \subset X$ 为仿射开集。由于 $X$ 的对角态射仿射，包含态射
$j : U \to X$ 也是仿射的（Morphisms, Lemma
\ref{morphisms-lemma-affine-permanence}）。同理，复合
$g = f \circ j : U \to Y$ 仿射。设 $\mathcal{I}^\bullet$ 为
$\QCoh(\mathcal{O}_U)$ 中的 K-内射复形。由于
$j_* : \QCoh(\mathcal{O}_U) \to \QCoh(\mathcal{O}_X)$ 有正合左伴随
$j^* : \QCoh(\mathcal{O}_X) \to \QCoh(\mathcal{O}_U)$，可知
$j_*\mathcal{I}^\bullet$ 是 $\QCoh(\mathcal{O}_X)$ 中的 K-内射复形，见
Derived Categories, Lemma
\ref{derived-lemma-adjoint-preserve-K-injectives}。于是
$$
\Phi(j_*\mathcal{I}^\bullet) =
f_*j_*\mathcal{I}^\bullet =
g_*\mathcal{I}^\bullet
$$
由引理 \ref{perfect-lemma-affine-pushforward}，$j_*\mathcal{I}^\bullet$ 表示
$Rj_*\mathcal{I}^\bullet$，而 $g_*\mathcal{I}^\bullet$ 表示
$Rg_*\mathcal{I}^\bullet$。另一方面，有 $Rf_* \circ Rj_* = Rg_*$。
因此 $f_*j_*\mathcal{I}^\bullet$ 表示 $Rf_*(j_*\mathcal{I}^\bullet)$。
我们断定，引理对任意形如 $j_*\mathcal{G}^\bullet$ 的复形成立，其中
$\mathcal{G}^\bullet$ 是 $U$ 上拟凝聚模的复形。（注意，若
$\mathcal{G}^\bullet \to \mathcal{I}^\bullet$ 是拟同构，则
$j_*\mathcal{G}^\bullet \to j_*\mathcal{I}^\bullet$ 也是拟同构，因为
$j_*$ 在拟凝聚模上是正合函子。）

\medskip\noindent
设 $\mathcal{F}^\bullet$ 为拟凝聚 $\mathcal{O}_X$-模的复形。设
$T \subset X$ 为闭子集，使得每个 $p$ 对应的 $\mathcal{F}^p$ 的支集均
包含在 $T$ 中。令 $n$ 为满足
$T \subset U_1 \cup \ldots \cup U_n$ 的仿射开集 $U_1, \ldots, U_n$
的最少个数，我们对此数作归纳。初始情形 $n = 0$ 显然成立。若
$n \geq 1$，令 $U = U_1$，并以 $j : U \to X$ 表示如上的开浸入。考虑
复形映射 $c : \mathcal{F}^\bullet \to j_*j^*\mathcal{F}^\bullet$，得到
两个复形短正合列：
$$
0 \to \Ker(c) \to \mathcal{F}^\bullet \to \Im(c) \to 0
$$
and
$$
0 \to \Im(c) \to j_*j^*\mathcal{F}^\bullet \to \Coker(c) \to 0
$$
复形 $\Ker(c)$ 和 $\Coker(c)$ 支持在
$T \setminus U \subset U_2 \cup \ldots \cup U_n$ 上，故由归纳假设，
结论对它们成立。由上一段的讨论，结论对
$j_*j^*\mathcal{F}^\bullet$ 成立。考察与这些短正合列相伴的可区别三角，
并使用证明开头的观察，即得结论。
\end{proof}

\begin{remark}[警告]
\label{perfect-remark-warning-coherator}
设 $X$ 为对角态射仿射的拟紧概形。尽管由命题
\ref{perfect-proposition-quasi-compact-affine-diagonal} 已知
$D(\QCoh(\mathcal{O}_X)) = D_\QCoh(\mathcal{O}_X)$，这里仍可能出现反常
现象，也很容易出错。一个陷阱是轻率地认为这个等式意味着导出函子也相同。
例如，假设有拟紧开集 $U \subset X$。可以考虑左正合函子
$\Gamma(U, -)$ 的高阶右导出函子
$$
R^i(\QCoh)\Gamma(U, -) : \QCoh(\mathcal{O}_X) \to \textit{Ab}
$$
这是一个泛 $\delta$-函子；另一方面，对 $X$ 上所有阿贝尔层定义的函子
$H^i(U, -)$ 限制在 $\QCoh(\mathcal{O}_X)$ 上后构成一个 $\delta$-函子。
因而得到典范变换
$$
t^i : R^i(\QCoh)\Gamma(U, -) \to H^i(U, -).
$$
即使 $X = \Spec(A)$ 仿射，这些变换一般也不是同构！事实上，由右导出函子
的构造以及 $\QCoh(\mathcal{O}_X)$ 与 $\text{Mod}_A$ 的等价，若 $I$ 为内射
$A$-模，就有 $R^1(\QCoh)\Gamma(U, \widetilde{I}) = 0$。但 Examples,
Lemma \ref{examples-lemma-nonvanishing} 表明，存在 $A$、$I$ 和 $U$ 使得
$H^1(U, \widetilde{I}) \not = 0$。
\end{remark}




\section{诺特概形的拟凝聚化子}
\label{perfect-section-coherator-Noetherian}

\noindent
对于诺特概形，可以使用下列引理。

\begin{lemma}
\label{perfect-lemma-injective-quasi-coherent-sheaf-Noetherian}
设 $X$ 为诺特概形，$\mathcal{J}$ 为 $\QCoh(\mathcal{O}_X)$ 的内射对象。
则 $\mathcal{J}$ 是 $\mathcal{O}_X$-模的软层。
\end{lemma}

\begin{proof}
设 $U \subset X$ 为开子集，$s \in \mathcal{J}(U)$ 为截面。设
$\mathcal{I} \subset \mathcal{O}_X$ 为定义 $X \setminus U$ 上约化诱导概形
结构的拟凝聚理想层（见 Schemes, Definition
\ref{schemes-definition-reduced-induced-scheme}）。由 Cohomology of Schemes,
Lemma \ref{coherent-lemma-homs-over-open}，截面 $s$ 对应于某个 $n$ 下的映射
$\sigma : \mathcal{I}^n \to \mathcal{J}$。由于 $\mathcal{J}$ 是
$\QCoh(\mathcal{O}_X)$ 的内射对象，可把 $\sigma$ 延拓为映射
$\tilde s : \mathcal{O}_X \to \mathcal{J}$。于是 $\tilde s$ 对应于
$\mathcal{J}$ 的一个整体截面，其在该开集上的限制为 $s$。
\end{proof}

\begin{lemma}
\label{perfect-lemma-Noetherian-pushforward}
设 $f : X \to Y$ 为诺特概形间的态射。则拟凝聚层上的 $f_*$ 有右导出延拓
$\Phi : D(\QCoh(\mathcal{O}_X)) \to D(\QCoh(\mathcal{O}_Y))$
，使图表
$$
\xymatrix{
D(\QCoh(\mathcal{O}_X)) \ar[d]_{\Phi} \ar[r] &
D_\QCoh(\mathcal{O}_X) \ar[d]^{Rf_*} \\
D(\QCoh(\mathcal{O}_Y)) \ar[r] &
D_\QCoh(\mathcal{O}_Y)
}
$$
交换。
\end{lemma}

\begin{proof}
由于 $X$ 和 $Y$ 是诺特概形，该态射拟紧且拟分离（见 Properties, Lemma
\ref{properties-lemma-locally-Noetherian-quasi-separated} 及 Schemes, Remark
\ref{schemes-remark-quasi-compact-and-quasi-separated}）。因此 $f_*$ 保持拟凝聚
性，见 Schemes, Lemma \ref{schemes-lemma-push-forward-quasi-coherent}。接着，
设 $K$ 为 $D(\QCoh(\mathcal{O}_X))$ 的对象。由于
$\QCoh(\mathcal{O}_X)$ 是 Grothendieck 阿贝尔范畴（Properties,
Proposition \ref{properties-proposition-coherator}），可以用 K-内射复形
$\mathcal{I}^\bullet$ 表示 $K$，使每个 $\mathcal{I}^n$ 都是
$\QCoh(\mathcal{O}_X)$ 的内射对象，见 Injectives, Theorem
\ref{injectives-theorem-K-injective-embedding-grothendieck}。因此函子 $\Phi$
由下式定义：
$$
\Phi(K) = f_*\mathcal{I}^\bullet
$$
其中右端视为 $D(\QCoh(\mathcal{O}_Y))$ 的对象。为完成引理的证明，只需证明
典范映射
$$
f_*\mathcal{I}^\bullet \longrightarrow Rf_*\mathcal{I}^\bullet
$$
在 $D(\mathcal{O}_Y)$ 中是同构。由引理 \ref{perfect-lemma-acyclicity-lemma}，只需证明
对所有 $n \in \mathbf{Z}$，$\mathcal{I}^n$ 都是右 $f_*$-无环的。这成立，
因为由引理 \ref{perfect-lemma-injective-quasi-coherent-sheaf-Noetherian}，
$\mathcal{I}^n$ 是软层；而由 Cohomology, Lemma
\ref{cohomology-lemma-flasque-acyclic-pushforward}，软模是右 $f_*$-无环的。
\end{proof}

\begin{proposition}
\label{perfect-proposition-Noetherian}
设 $X$ 为诺特概形。则函子 (\ref{perfect-equation-compare})
$$
D(\QCoh(\mathcal{O}_X))
\longrightarrow
D_\QCoh(\mathcal{O}_X)
$$
是等价，且其拟逆由 $RQ_X$ 给出。
\end{proposition}

\begin{proof}
这由引理 \ref{perfect-lemma-argument-proves} 和引理
\ref{perfect-lemma-Noetherian-pushforward} 得到。
\end{proof}






\section{Koszul 复形}
\label{perfect-section-koszul}

\noindent
设 $A$ 为环，$f_1, \ldots, f_r$ 为 $A$ 中的一列元素。我们已在 More on
Algebra, Definition \ref{more-algebra-definition-koszul-complex} 中定义 Koszul
复形 $K_\bullet(f_1, \ldots, f_r)$。它是位于次数 $r, \ldots, 0$ 的链复形。
令 $K^{-n}(f_1, \ldots, f_r) = K_n(f_1, \ldots, f_r)$ 并采用相同的微分，
即可将其化为上链复形 $K^\bullet(f_1, \ldots, f_r)$。本节余下部分中的所有
复形都是上链复形。

\medskip\noindent
定义复形 $I^\bullet(f_1, \ldots, f_r)$，使得在 $K(A)$ 中有可区别三角
$$
I^\bullet(f_1, \ldots, f_r) \to
A \to
K^\bullet(f_1, \ldots, f_r) \to
I^\bullet(f_1, \ldots, f_r)[1]
$$
换言之，令
$$
I^i(f_1, \ldots, f_r) =
\left\{
\begin{matrix}
K^{i - 1}(f_1, \ldots, f_r) & \text{若 } i \leq 0 \\
0 & \text{否则}
\end{matrix}
\right.
$$
并采用 $K^\bullet(f_1, \ldots, f_r)$ 上微分的负值。可区别三角中的映射是
显然的映射。注意，$I^0(f_1, \ldots, f_r) = A^{\oplus r} \to A$ 在第
$i$ 个因子上由乘以 $f_i$ 给出。因此
$I^\bullet(f_1, \ldots, f_r) \to A$ 分解为
$$
I^\bullet(f_1, \ldots, f_r) \to I \to A
$$
其中 $I = (f_1, \ldots, f_r)$。事实上，有短正合列
$$
0 \to H^{-1}(K^\bullet(f_1, \ldots, f_s)) \to
H^0(I^\bullet(f_1, \ldots, f_s)) \to I \to 0
$$
并且对每个 $i < 0$，有
$H^i(I^\bullet(f_1, \ldots, f_r)) = H^{i - 1}(K^\bullet(f_1, \ldots, f_r))$.
注意，给定 $A$ 中的另一列元素 $g_1, \ldots, g_r$，有典范映射
$$
I^\bullet(f_1g_1, \ldots, f_rg_r) \to I^\bullet(f_1, \ldots, f_r)
\quad\text{且}\quad
K^\bullet(f_1g_1, \ldots, f_rg_r) \to K^\bullet(f_1, \ldots, f_r)
$$
它们与上述映射相容。其中第一个映射在
$I^0(f_1g_1, \ldots, f_rg_r) = A^{\oplus r}$ 的第 $i$ 个直和项上由乘以
$g_i$ 给出。特别地，给定 $f_1, \ldots, f_r$，得到复形的逆向系统
\begin{equation}
\label{perfect-equation-system}
I^\bullet(f_1, \ldots, f_r) \leftarrow
I^\bullet(f_1^2, \ldots, f_r^2) \leftarrow
I^\bullet(f_1^3, \ldots, f_r^3) \leftarrow \ldots
\end{equation}
它将在后文中起重要作用。为便于表述以下引理，我们固定一些记号。

\begin{situation}
\label{perfect-situation-complex}
这里 $A$ 为环，$f_1, \ldots, f_r$ 为 $A$ 中的一列元素。令
$X = \Spec(A)$ 及 $U = D(f_1) \cup \ldots \cup D(f_r) \subset X$。以
$\mathcal{U} : U = \bigcup_{i = 1, \ldots, r} D(f_i)$ 表示 $U$ 的这个给定
开覆盖。
\end{situation}

\noindent
第一个引理说明，上述复形可用于计算 $U$ 上拟凝聚层的上同调。

\begin{lemma}
\label{perfect-lemma-alternating-cech-complex}
在情形 \ref{perfect-situation-complex} 中，设 $M$ 为 $A$-模，并以 $\mathcal{F}$
表示相伴的 $\mathcal{O}_X$-模。则有复形的典范同构
$$
\Psi : \colim_e \Hom_A(I^\bullet(f_1^e, \ldots, f_r^e), M)
\longrightarrow
\check{\mathcal{C}}_{alt}^\bullet(\mathcal{U}, \mathcal{F})
$$
它关于 $M$ 函子性地成立，其中 $\Hom$-复形上的微分是
$I^\bullet(f_1^e, \ldots, f_r^e)$ 上微分的逆步变换。
\end{lemma}

\begin{proof}
回忆，交错 {\v C}ech 复形是通常 {\v C}ech 复形中由交错上链组成的子复形，
见 Cohomology, Section \ref{cohomology-section-alternating-cech}。通常把
$\check{\mathcal{C}}_{alt}^\bullet(\mathcal{U}, \mathcal{F})$ 中的 $p$-上链
视为 $\{1, \ldots, r\}^{p + 1}$ 上的交错函数 $s$，其在
$(i_0, \ldots, i_p)$ 处的值 $s_{i_0\ldots i_p}$ 属于
$M_{f_{i_0}\ldots f_{i_p}} = \mathcal{F}(U_{i_0\ldots i_p})$。另一方面，
$\Hom^\bullet(I^\bullet(f_1^e, \ldots, f_r^e), M)$ 中的 $p$-上链 $t$ 是映射
$t : \wedge^{p + 1}(A^{\oplus r}) \to M$。以 $[i] \in A^{\oplus r}$ 表示
第 $i$ 个基向量，并记
$$
[i_0, \ldots, i_p] = [i_0] \wedge \ldots \wedge [i_p]
\in \wedge^{p + 1}(A^{\oplus r})
$$
对上述 $t$，令
$$
\Psi(t)_{i_0 \ldots i_p} =
(-1)^p \frac{t([i_0, \ldots, i_p])}{f_{i_0}^e\ldots f_{i_p}^e}
$$
显然 $\Psi(t)$ 是交错上链。上述规则与系统的转移映射相容，因为
(\ref{perfect-equation-system}) 的转移映射
$$
I^\bullet(f_1^e, \ldots, f_r^e) \leftarrow
I^\bullet(f_1^{e + 1}, \ldots, f_r^{e + 1}),
$$
把 $[i_0, \ldots, i_p]$ 送到
$f_{i_0}\ldots f_{i_p}[i_0, \ldots, i_p]$。根据 Algebra, Lemma
\ref{algebra-lemma-localization-colimit} 中对局部化
$M_{f_{i_0} \ldots f_{i_p}}$ 的描述，显然规则 $\Psi$ 在余极限的次数 $p$
部分上定义上链模的同构。为完成证明，还需证明此映射与微分相容。对此，
对上述 $t$ 计算：
\begin{align*}
d(\Psi(t))_{i_0 \ldots i_{p + 1}}
& =
\sum\nolimits_{j = 0}^{p + 1} (-1)^j
\Psi(t)_{i_0\ldots \hat i_j \ldots i_{p + 1}} \\
& =
(-1)^p
\sum\nolimits_{j = 0}^{p + 1} (-1)^j t([i_0 \ldots \hat i_j \ldots i_{p + 1}])
(f_{i_0} \ldots \hat f_{i_j} \ldots f_{i_p})^{-e}
\end{align*}
回忆，$I^\bullet(f_1^e, \ldots, f_r^e)$ 上的微分是
$K^\bullet(f_1, \ldots, f_r)$ 上微分的负值。因此
\begin{align*}
\Psi(d(t))_{i_0 \ldots i_{p + 1}} & =
(-1)^{p + 1}
d(t)([i_0, \ldots, i_{p + 1}]) (f_{i_0} \ldots f_{i_{p + 1}})^{-e} \\
& =
(-1)^{p + 1}
t(d([i_0, \ldots, i_{p + 1}])) (f_{i_0} \ldots f_{i_{p + 1}})^{-e} \\
& =
(-1)^{p + 1}
t(-\sum\nolimits_{j = 0}^{p + 1}
(-1)^j f_{i_j}^e [i_0, \ldots, \hat i_j, \ldots i_{p + 1}])
(f_{i_0} \ldots f_{i_{p + 1}})^{-e} \\
& =
-(-1)^{p + 1}
\sum\nolimits_{j = 0}^{p + 1}
(-1)^j t([i_0, \ldots, \hat i_j, \ldots i_{p + 1}])
(f_{i_0} \ldots \hat f_{i_j} \ldots f_{i_p})^{-e}
\end{align*}
两个公式相同，证明完成。
\end{proof}

\noindent
设给定 $A$-模的有限复形 $I^\bullet$ 和 $A$-模复形 $M^\bullet$，则有相应的
Hom 复形 $\Hom^\bullet(I^\bullet, M^\bullet)$。此复形的 $n$ 次项为
$$
\bigoplus\nolimits_{p + q = n} \Hom_A(I^{-q}, M^p)
$$
其微分如 More on Algebra, Section
\ref{more-algebra-section-hom-complexes} 所述。由于复形 $I^\bullet$ 只有有限个
非零项，上述直和是有限直和。由一个复形的 {\v C}ech 复形取相伴总复形的
约定如 Cohomology, Section
\ref{cohomology-section-cech-cohomology-of-complexes} 所述。

\begin{lemma}
\label{perfect-lemma-alternating-cech-complex-complex}
在情形 \ref{perfect-situation-complex} 中，设 $M^\bullet$ 为 $A$-模复形，并以
$\mathcal{F}^\bullet$ 表示相伴的 $\mathcal{O}_X$-模复形。则有复形的典范同构
$$
\colim_e \Hom^\bullet(I^\bullet(f_1^e, \ldots, f_r^e), M^\bullet)
\longrightarrow
\text{Tot}(\check{\mathcal{C}}_{alt}^\bullet(\mathcal{U}, \mathcal{F}^\bullet))
$$
它关于 $M^\bullet$ 函子性地成立。
\end{lemma}

\begin{proof}
考虑 Cohomology, Section
\ref{cohomology-section-cech-cohomology-of-complexes} 中讨论的双复形
$F^{\bullet, \bullet}$，其项为
$F^{p, q} = \mathcal{C}_{alt}^p(\mathcal{U}, \mathcal{F}^q)$。再考虑双复形
$G^{\bullet, \bullet}$，其项为
$G^{p, q} = \colim_e \Hom_A(I^{-p}(f_1^e, \ldots, f_r^e), M^q)$，微分由
函子性给出（不插入符号）。引理 \ref{perfect-lemma-alternating-cech-complex} 的证明中
所构造的映射 $\psi^{p, q} : G^{p, q} \to F^{p, q}$ 是同构，并与微分
$d_1$（由该引理）和 $d_2$（显然）相容。然而，引理中箭头左右两端复形的
微分 $d$ 具有不同符号。具体地，对 $g \in G^{p, q}$，微分为
$$
d(g) =  d_2(g) - (-1)^{p + q} d_1(g) 
$$
（见 More on Algebra, Section \ref{more-algebra-section-hom-complexes}），而对
$f \in F^{p, q}$，微分为
$$
d(f) = d_1(f) + (-1)^p d_2(f)
$$
因此，把 $\psi^{p, q}$ 乘以 $(-1)^{pq + p(p - 1)/2}$ 即可校正符号。
\end{proof}

\begin{lemma}
\label{perfect-lemma-alternating-cech-complex-complex-computes-cohomology}
在情形 \ref{perfect-situation-complex} 中，设 $\mathcal{F}^\bullet$ 为拟凝聚
$\mathcal{O}_X$-模的复形。则在 $D(A)$ 中有典范同构
$$
\text{Tot}(\check{\mathcal{C}}_{alt}^\bullet(\mathcal{U}, \mathcal{F}^\bullet))
\longrightarrow
R\Gamma(U, \mathcal{F}^\bullet)
$$
它关于 $\mathcal{F}^\bullet$ 函子性地成立。
\end{lemma}

\begin{proof}
设 $\mathcal{B}$ 为 $U$ 的仿射开集之集。由于仿射概形上拟凝聚模的高阶
上同调群为零（Cohomology of Schemes, Lemma
\ref{coherent-lemma-quasi-coherent-affine-cohomology-zero}），这是 Cohomology,
Lemma \ref{cohomology-lemma-alternating-cech-complex-complex-ss} 的一个特例。
\end{proof}

\noindent
在情形 \ref{perfect-situation-complex} 中，以 $I_e$ 表示在引理
\ref{perfect-lemma-affine-compare-bounded} 的等价下，与 $A$-模复形
$I^\bullet(f_1^e, \ldots, f_r^e)$ 对应的 $D(\mathcal{O}_X)$ 的对象。映射
(\ref{perfect-equation-system}) 给出系统
$$
I_1 \leftarrow
I_2 \leftarrow
I_3 \leftarrow \ldots
$$
此外，有相容的映射系统 $I_e \to \mathcal{O}_X$，它们限制到 $U$ 上后成为
同构。因此对 $D(\mathcal{O}_X)$ 的每个对象 $E$，都有典范映射
\begin{equation}
\label{perfect-equation-comparison}
\colim_e \Hom_{D(\mathcal{O}_X)}(I_e, E) \longrightarrow H^0(U, E)
\end{equation}
其构造是把映射 $I_e \to E$ 送到它在 $U$ 上的限制，并使用等式
$\Hom_{D(\mathcal{O}_U)}(\mathcal{O}_U, E|_U) = H^0(U, E)$.

\begin{proposition}
\label{perfect-proposition-represent-cohomology-class-on-open}
在情形 \ref{perfect-situation-complex} 中，对 $D_\QCoh(\mathcal{O}_X)$ 的每个对象
$E$，映射 (\ref{perfect-equation-comparison}) 都是同构。
\end{proposition}

\begin{proof}
由引理 \ref{perfect-lemma-affine-compare-bounded}，可假设 $E$ 由拟凝聚层的复形
$\mathcal{F}^\bullet$ 给出。令
$M^\bullet = \Gamma(X, \mathcal{F}^\bullet)$ 为相应的 $A$-模复形。由引理
\ref{perfect-lemma-alternating-cech-complex-complex} 和
\ref{perfect-lemma-alternating-cech-complex-complex-computes-cohomology}，有拟同构
$$
\colim_e \Hom^\bullet(I^\bullet(f_1^e, \ldots, f_r^e), M^\bullet)
\longrightarrow
\text{Tot}(\check{\mathcal{C}}_{alt}^\bullet(\mathcal{U}, \mathcal{F}^\bullet))
\longrightarrow
R\Gamma(U, \mathcal{F}^\bullet)
$$
由 More on Algebra, Lemma \ref{more-algebra-lemma-RHom-out-of-projective} 和
Equation (\ref{more-algebra-equation-h0-RHom})，对复形
$\Hom^\bullet(I^\bullet(f_1^e, \ldots, f_r^e), M^\bullet)$
取 $H^0$ 即计算 $D(A)$ 中的 $\Hom$。因此两端取 $H^0$，得到
$$
\colim_e \Hom_{D(A)}(I^\bullet(f_1^e, \ldots, f_r^e), M^\bullet)
=
H^0(U, E)
$$
由引理 \ref{perfect-lemma-affine-compare-bounded}，
$\Hom_{D(A)}(I^\bullet(f_1^e, \ldots, f_r^e), M^\bullet) =
\Hom_{D(\mathcal{O}_X)}(I_e, E)$，结论得证。
\end{proof}

\noindent
在情形 \ref{perfect-situation-complex} 中，以 $K_e$ 表示在引理
\ref{perfect-lemma-affine-compare-bounded} 的等价下，与 $A$-模复形
$K^\bullet(f_1^e, \ldots, f_r^e)$ 对应的 $D(\mathcal{O}_X)$ 的对象。于是有
可区别三角
$$
I_e \to \mathcal{O}_X \to K_e \to I_e[1]
$$
及系统
$$
K_1 \leftarrow
K_2 \leftarrow
K_3 \leftarrow \ldots
$$
它与系统 $(I_e)$ 相容。此外，还有相容的映射系统
$$
K_e \to H^0(K_e) = \mathcal{O}_X/(f_1^e, \ldots, f_r^e)
$$

\begin{lemma}
\label{perfect-lemma-represent-cohomology-class-on-closed}
在情形 \ref{perfect-situation-complex} 中，设 $E$ 为
$D_\QCoh(\mathcal{O}_X)$ 的对象。假设对
$i = - r + 1, \ldots, 0$ 有 $H^i(E)|_U = 0$。那么给定
$s \in H^0(X, E)$，存在 $e \geq 0$ 及态射 $K_e \to E$，使得 $s$ 属于映射
$H^0(X, K_e) \to H^0(X, E)$ 的像。
\end{lemma}

\begin{proof}
由于 $U$ 被 $r$ 个仿射开集覆盖，对任意拟凝聚模，在 $j \geq r$ 时有
$H^j(U, \mathcal{F}) = 0$（Cohomology of Schemes, Lemma
\ref{coherent-lemma-vanishing-nr-affines}）。由引理
\ref{perfect-lemma-application-nice-K-injective}，$H^0(U, E)$ 等于
$H^0(U, \tau_{\geq -r + 1}E)$。有谱序列
$$
H^j(U, H^i(\tau_{\geq -r + 1}E)) \Rightarrow H^{i + j}(U, \tau_{\geq -N}E)
$$
见 Derived Categories, Lemma \ref{derived-lemma-two-ss-complex-functor}。因此由
对 $E$ 的上同调层所作的消失假设，$H^0(U, E) = 0$，从而 $s|_U = 0$。
把 $s$ 看作 $D(\mathcal{O}_X)$ 中的态射 $\mathcal{O}_X \to E$。由命题
\ref{perfect-proposition-represent-cohomology-class-on-open}，对某个 $e$，复合
$I_e \to \mathcal{O}_X \to E$ 为零。由可区别三角
$I_e \to \mathcal{O}_X \to K_e \to I_e[1]$，得到态射 $K_e \to E$，使
$s$ 是复合 $\mathcal{O}_X \to K_e \to E$。
\end{proof}


\section{伪凝聚复形与完美复形}
\label{perfect-section-spell-out}

\noindent
本节把 Cohomology, Sections \ref{cohomology-section-strictly-perfect},
\ref{cohomology-section-pseudo-coherent},
\ref{cohomology-section-tor}, and
\ref{cohomology-section-perfect} 中定义的一般概念，与 More on Algebra,
Sections
\ref{more-algebra-section-pseudo-coherent},
\ref{more-algebra-section-tor}, and
\ref{more-algebra-section-perfect} 中模复形的相应概念联系起来。

\begin{lemma}
\label{perfect-lemma-pseudo-coherent}
设 $X$ 为概形。若 $E$ 是 $D(\mathcal{O}_X)$ 的 $m$-伪凝聚对象，则对
$i > m$，$H^i(E)$ 是拟凝聚 $\mathcal{O}_X$-模，而 $H^m(E)$ 是某个拟凝聚
$\mathcal{O}_X$-模的商。若 $E$ 伪凝聚，则 $E$ 是
$D_\QCoh(\mathcal{O}_X)$ 的对象。
\end{lemma}

\begin{proof}
在 $X$ 上局部地，存在严格完美复形 $\mathcal{E}^\bullet$，使得对
$i > m$，$H^i(E)$ 同构于 $H^i(\mathcal{E}^\bullet)$，且 $H^m(E)$ 是
$H^m(\mathcal{E}^\bullet)$ 的商。层 $\mathcal{E}^i$ 是有限自由模的直和项，
故拟凝聚。结论得证。
\end{proof}

\begin{lemma}
\label{perfect-lemma-pseudo-coherent-affine}
设 $X = \Spec(A)$ 为仿射概形，$M^\bullet$ 为 $A$-模复形，$E$ 为
$D(\mathcal{O}_X)$ 中相应的对象。那么，$E$ 作为 $D(\mathcal{O}_X)$ 的对象
是 $m$-伪凝聚的（相应地，伪凝聚的），当且仅当 $M^\bullet$ 作为
$A$-模复形是 $m$-伪凝聚的（相应地，伪凝聚的）。
\end{lemma}

\begin{proof}
由定义立即可知，若 $M^\bullet$ 是 $m$-伪凝聚的，则 $E$ 也是。为证明反向，
假设 $E$ 是 $m$-伪凝聚的。由于 $X = \Spec(A)$ 拟紧，且其拓扑有由标准开集
组成的基，可以找到标准开覆盖 $X = D(f_1) \cup \ldots \cup D(f_n)$、
$D(f_i)$ 上的严格完美复形 $\mathcal{E}_i^\bullet$，以及映射
$\alpha_i : \mathcal{E}_i^\bullet \to E|_{U_i}$；这些映射在 $j > m$ 时于
$H^j$ 上诱导同构，并在 $H^m$ 上诱导满射。由 Cohomology, Lemma
\ref{cohomology-lemma-local-actual}，细分开覆盖后，可假设每个 $i$ 对应的
$\alpha_i$ 都由复形映射
$\mathcal{E}_i^\bullet \to \widetilde{M^\bullet}|_{U_i}$
给出。由 Modules, Lemma
\ref{modules-lemma-direct-summand-of-locally-free-is-locally-free}，各项
$\mathcal{E}_i^n$ 都是有限局部自由模。因此再次细分开覆盖后，可假设每个
$\mathcal{E}_i^n$ 都是有限自由 $\mathcal{O}_{U_i}$-模。由定义可知，
$M^\bullet_{f_i}$ 是 $A_{f_i}$-模的 $m$-伪凝聚复形。应用 More on Algebra,
Lemma \ref{more-algebra-lemma-glue-pseudo-coherent} 即得结论。

\medskip\noindent
“伪凝聚”的情形来自如下事实：按定义，$E$ 伪凝聚当且仅当对所有 $m$，
$E$ 都是 $m$-伪凝聚的；而由 More on Algebra, Lemma
\ref{more-algebra-lemma-pseudo-coherent}，$M^\bullet$ 也有相同的性质。
\end{proof}

\begin{lemma}
\label{perfect-lemma-identify-pseudo-coherent-noetherian}
设 $X$ 为诺特概形，$E$ 为 $D_\QCoh(\mathcal{O}_X)$ 的对象。对
$m \in \mathbf{Z}$，以下条件等价：
\begin{enumerate}
\item 对 $i \geq m$，$H^i(E)$ 凝聚，并且对 $i \gg 0$，它为零；
\item $E$ 是 $m$-伪凝聚的。
\end{enumerate}
特别地，$E$ 伪凝聚当且仅当 $E$ 是
$D^-_{\textit{Coh}}(\mathcal{O}_X)$ 的对象。
\end{lemma}

\begin{proof}
由于 $X$ 拟紧，可知在 (1) 和 (2) 两种情形下，对象 $E$ 都有上界。因此问题
关于 $X$ 是局部的，可假设 $X$ 仿射。设 $X = \Spec(A)$，其中 $A$ 为某个
诺特环。此时由引理 \ref{perfect-lemma-affine-compare-bounded}，$E$ 对应于
$A$-模复形 $M^\bullet$。由引理 \ref{perfect-lemma-pseudo-coherent-affine}，$E$ 是
$m$-伪凝聚的，当且仅当 $M^\bullet$ 是 $m$-伪凝聚的。另一方面，$H^i(E)$
凝聚当且仅当 $H^i(M^\bullet)$ 是有限 $A$-模（Properties, Lemma
\ref{properties-lemma-finite-type-module}）。因此结论由 More on Algebra,
Lemma \ref{more-algebra-lemma-Noetherian-pseudo-coherent} 得到。
\end{proof}

\begin{lemma}
\label{perfect-lemma-tor-dimension-affine}
设 $X = \Spec(A)$ 为仿射概形，$M^\bullet$ 为 $A$-模复形，$E$ 为
$D(\mathcal{O}_X)$ 中相应的对象。则
\begin{enumerate}
\item $E$ 的 Tor 振幅在 $[a, b]$ 内，当且仅当 $M^\bullet$ 的 Tor 振幅在
$[a, b]$ 内。
\item $E$ 具有有限 Tor 维数，当且仅当 $M^\bullet$ 具有有限 Tor 维数。
\end{enumerate}
\end{lemma}

\begin{proof}
(2) 显然由 (1) 得到。在 (1) 的证明中，我们将不再特别说明地使用引理
\ref{perfect-lemma-affine-compare-bounded} 的等价 $D(A) = D_\QCoh(X)$。假设
$M^\bullet$ 的 Tor 振幅在 $[a, b]$ 内。则 $K^\bullet$ 在 $D(A)$ 中同构于
由平坦 $A$-模构成的复形 $K^\bullet$，且对 $i \not \in [a, b]$ 有
$K^i = 0$，见 More on Algebra, Lemma
\ref{more-algebra-lemma-tor-amplitude}。于是 $E$ 同构于
$\widetilde{K^\bullet}$。由于每个 $\widetilde{K^i}$ 都是平坦
$\mathcal{O}_X$-模，由 Cohomology, Lemma
\ref{cohomology-lemma-tor-amplitude} 可知 $E$ 的 Tor 振幅在 $[a, b]$ 内。

\medskip\noindent
假设 $E$ 的 Tor 振幅在 $[a, b]$ 内。则 $E$ 有界，从而 $M^\bullet$ 属于
$K^-(A)$。因此可以把 $M^\bullet$ 替换为有上界的 $A$-模复形。甚至可以
选取投射消解，并假设 $M^\bullet$ 是由自由 $A$-模构成的有上界复形。于是
对任意 $A$-模 $N$，有
$$
E \otimes_{\mathcal{O}_X}^\mathbf{L} \widetilde{N}
\cong
\widetilde{M^\bullet} \otimes_{\mathcal{O}_X}^\mathbf{L} \widetilde{N}
\cong
\widetilde{M^\bullet \otimes_A N}
$$
于 $D(\mathcal{O}_X)$ 中。因此左端上同调层的消失性蕴含 $M^\bullet$ 的
Tor 振幅在 $[a, b]$ 内。
\end{proof}

\begin{lemma}
\label{perfect-lemma-tor-dimension-rel-affine}
设 $f : X \to S$ 为对应于环同态 $R \to A$ 的仿射概形间态射。设
$M^\bullet$ 为 $A$-模复形，$E$ 为 $D(\mathcal{O}_X)$ 中相应的对象。则
\begin{enumerate}
\item $E$ 作为 $D(f^{-1}\mathcal{O}_S)$ 的对象，其 Tor 振幅在 $[a, b]$
内，当且仅当 $M^\bullet$ 作为 $D(R)$ 的对象，其 Tor 振幅在 $[a, b]$ 内。
\item $E$ 作为 $D(f^{-1}\mathcal{O}_S)$ 的对象局部具有有限 Tor 维数，
当且仅当 $M^\bullet$ 作为 $D(R)$ 的对象具有有限 Tor 维数。
\end{enumerate}
\end{lemma}

\begin{proof}
考虑位于 $\mathfrak p \subset R$ 之上的素理想 $\mathfrak q \subset A$。
设 $x \in X$ 和 $s = f(x) \in S$ 为相应的点。于是
$(f^{-1}\mathcal{O}_S)_x = \mathcal{O}_{S, s} = R_\mathfrak p$ 且
$E_x = M^\bullet_\mathfrak q$。由此可如下看出等价性。

\medskip\noindent
若 $M^\bullet$ 作为 $R$-模复形的 Tor 振幅在 $[a, b]$ 内，则
$M^\bullet$ 在 $A$ 的任意素理想处局部化后仍有此性质。由 Cohomology,
Lemma \ref{cohomology-lemma-tor-amplitude-stalk} 可知，$E$ 作为
$f^{-1}\mathcal{O}_S$-模层复形的 Tor 振幅在 $[a, b]$ 内。反之，假设 $E$
作为 $D(f^{-1}\mathcal{O}_S)$ 的对象，其 Tor 振幅在 $[a, b]$ 内。再次使用
刚才引用的引理，可知对每个位于 $\mathfrak p \subset R$ 之上的素理想
$\mathfrak q \subset A$，$M^\bullet_\mathfrak q$ 作为
$R_\mathfrak p$-模复形的 Tor 振幅在 $[a, b]$ 内。由 More on Algebra,
Lemma \ref{more-algebra-lemma-tor-amplitude-localization}，$M^\bullet$ 作为
$R$-模复形的 Tor 振幅在 $[a, b]$ 内。这就完成了 (1) 的证明。

\medskip\noindent
由于 $X$ 拟紧，若 $E$ 作为 $f^{-1}\mathcal{O}_S$-模复形局部具有有限 Tor
维数，则事实上存在某些 $a, b$，使 $E$ 作为
$f^{-1}\mathcal{O}_S$-模复形的 Tor 振幅在 $[a, b]$ 内。因此 (2) 由 (1)
得到。
\end{proof}

\begin{lemma}
\label{perfect-lemma-tor-qc-qs}
设 $X$ 为拟分离概形，$E$ 为 $D_\QCoh(\mathcal{O}_X)$ 的对象，并设
$a \leq b$。以下条件等价：
\begin{enumerate}
\item $E$ 的 Tor 振幅在 $[a, b]$ 内；
\item 对 $\QCoh(\mathcal{O}_X)$ 中所有 $\mathcal{F}$ 及
$i \not \in [a, b]$，都有
$H^i(E \otimes_{\mathcal{O}_X}^\mathbf{L} \mathcal{F}) = 0$。
\end{enumerate}
\end{lemma}

\begin{proof}
显然 (1) 蕴含 (2)。假设 (2)。设 $U \subset X$ 为仿射开集。由于 $X$
拟分离，态射 $j : U \to X$ 拟紧且分离，因此 $j_*$ 把拟凝聚模变为拟凝聚
模（Schemes, Lemma \ref{schemes-lemma-push-forward-quasi-coherent}）。故函子
$\QCoh(\mathcal{O}_X) \to \QCoh(\mathcal{O}_U)$ 本质满。由此，条件 (2)
蕴含对所有拟凝聚 $\mathcal{O}_U$-模 $\mathcal{G}$ 及
$i \not \in [a, b]$，
$H^i(E|_U \otimes_{\mathcal{O}_U}^\mathbf{L} \mathcal{G})$
都消失。写成 $U = \Spec(A)$，并设 $M^\bullet$ 为由引理
\ref{perfect-lemma-affine-compare-bounded} 对应于 $E|_U$ 的 $A$-模复形。我们刚刚证明
$M^\bullet \otimes_A^\mathbf{L} N$ 在区间 $[a, b]$ 外的上同调群消失；换言之，
$M^\bullet$ 的 Tor 振幅在 $[a, b]$ 内。由引理
\ref{perfect-lemma-tor-dimension-affine}，$E|_U$ 的 Tor 振幅在 $[a, b]$ 内。引理得证。
\end{proof}

\begin{lemma}
\label{perfect-lemma-perfect-affine}
设 $X = \Spec(A)$ 为仿射概形，$M^\bullet$ 为 $A$-模复形，$E$ 为
$D(\mathcal{O}_X)$ 中相应的对象。则 $E$ 是 $D(\mathcal{O}_X)$ 的完美对象，
当且仅当 $M^\bullet$ 是 $D(A)$ 的完美对象。
\end{lemma}

\begin{proof}
这是引理 \ref{perfect-lemma-pseudo-coherent-affine} 和
\ref{perfect-lemma-tor-dimension-affine}、Cohomology, Lemma
\ref{cohomology-lemma-perfect} 以及 More on Algebra, Lemma
\ref{more-algebra-lemma-perfect} 的逻辑推论。
\end{proof}

\noindent
作为概形上伪凝聚复形之描述的推论，我们可以证明某些内 Hom 是拟凝聚的。

\begin{lemma}
\label{perfect-lemma-quasi-coherence-internal-hom}
设 $X$ 为概形。
\begin{enumerate}
\item 若 $L$ 属于 $D^+_\QCoh(\mathcal{O}_X)$，且 $D(\mathcal{O}_X)$ 中的
$K$ 伪凝聚，则 $R\SheafHom(K, L)$ 属于 $D_\QCoh(\mathcal{O}_X)$，并且
局部有下界。
\item 若 $L$ 属于 $D_\QCoh(\mathcal{O}_X)$，且 $D(\mathcal{O}_X)$ 中的
$K$ 完美，则 $R\SheafHom(K, L)$ 属于 $D_\QCoh(\mathcal{O}_X)$。
\item 若 $X = \Spec(A)$ 仿射且 $K, L \in D(A)$，则
$$
R\SheafHom(\widetilde{K}, \widetilde{L}) = \widetilde{R\Hom_A(K, L)}
$$
在以下两种情形成立：
\begin{enumerate}
\item $K$ 伪凝聚且 $L$ 有下界；
\item $K$ 完美而 $L$ 任意。
\end{enumerate}
\item 若 $X = \Spec(A)$ 且 $K, L$ 属于 $D(A)$，则
$R\SheafHom(\widetilde{K}, \widetilde{L})$ 的第 $n$ 个上同调层是与预层
$$
X \supset D(f) \longmapsto \Ext^n_{A_f}(K \otimes_A A_f, L \otimes_A A_f)
$$
相伴的层，其中 $f \in A$。
\end{enumerate}
\end{lemma}

\begin{proof}
在 $\mathcal{O}_X$ 的导出范畴中构造内 Hom 与局部化交换（见 Cohomology,
Section \ref{cohomology-section-internal-hom}）。因此为证明 (1) 和 (2)，可以
用仿射开集替换 $X$。由引理 \ref{perfect-lemma-affine-compare-bounded}、
\ref{perfect-lemma-pseudo-coherent-affine} 和
\ref{perfect-lemma-perfect-affine}
，为证明 (1) 和 (2)，只需证明 (3)。

\medskip\noindent
把 $K$ 表示为由有限投射 $A$-模构成的有上界（相应地，有限）复形，并把
$L$ 表示为 $A$-模的有下界（相应地，任意）复形，则 (3) 由 Cohomology,
Lemma
\ref{cohomology-lemma-Rhom-complex-of-direct-summands-finite-free}
中对内 Hom 的计算得到。

\medskip\noindent
为证明 (4)，回忆在任意带环空间上，$R\SheafHom(A, B)$ 的第 $n$ 个上同调层
是与预层
$$
U \mapsto \Hom_{D(U)}(A|_U, B|_U[n]) =
\Ext^n_{D(\mathcal{O}_U)}(A|_U, B|_U)
$$
相伴的层，见 Cohomology, Section \ref{cohomology-section-internal-hom}。另一
方面，$\widetilde{K}$ 在主开集 $D(f)$ 上的限制是 $K \otimes_A A_f$ 的像，
对 $L$ 亦然。因此 (4) 由引理 \ref{perfect-lemma-affine-compare-bounded} 的范畴等价
得到。
\end{proof}

\begin{lemma}
\label{perfect-lemma-internal-hom-evaluate-tensor-isomorphism}
设 $X$ 为概形，$K, L, M$ 为 $D_\QCoh(\mathcal{O}_X)$ 的对象。Cohomology,
Lemma \ref{cohomology-lemma-internal-hom-diagonal-better} 中的映射
$$
K \otimes_{\mathcal{O}_X}^\mathbf{L} R\SheafHom(M, L)
\longrightarrow
R\SheafHom(M, K \otimes_{\mathcal{O}_X}^\mathbf{L} L)
$$
在以下情形是同构：
\begin{enumerate}
\item $M$ 完美；
\item $K$ 完美；
\item $M$ 伪凝聚，$L \in D^+(\mathcal{O}_X)$，且 $K$ 具有有限 Tor 维数。
\end{enumerate}
\end{lemma}

\begin{proof}
引理 \ref{perfect-lemma-quasi-coherence-internal-hom} 把情形 (1) 和 (3) 化为仿射
情形，后者已在 More on Algebra, Lemma
\ref{more-algebra-lemma-internal-hom-evaluate-tensor-isomorphism} 中处理。（把
问题转写成代数语言时，还须使用引理 \ref{perfect-lemma-pseudo-coherent-affine}、
\ref{perfect-lemma-perfect-affine} 和 \ref{perfect-lemma-tor-dimension-affine}。）若 $K$ 完美而
不作其他假设，我们并不知道箭头的任何一端属于
$D_\QCoh(\mathcal{O}_X)$；但结论仍成立，因为可以局部工作，把问题化到
$K$ 是有限自由模的有限复形这一显然情形。
\end{proof}





\section{凝聚模的导出范畴}
\label{perfect-section-derived-coherent}

\noindent
设 $X$ 为局部诺特概形。此时，凝聚 $\mathcal{O}_X$-模的范畴
$\textit{Coh}(\mathcal{O}_X) \subset \textit{Mod}(\mathcal{O}_X)$ 是弱 Serre
子范畴，见 Homology, Section \ref{homology-section-serre-subcategories} 及
Cohomology of Schemes, Lemma
\ref{coherent-lemma-coherent-abelian-Noetherian}。以
$$
D_{\textit{Coh}}(\mathcal{O}_X) \subset D(\mathcal{O}_X)
$$
表示上同调层凝聚的复形所成的子范畴，见 Derived Categories, Section
\ref{derived-section-triangulated-sub}。于是得到典范函子
\begin{equation}
\label{perfect-equation-compare-coherent}
D(\textit{Coh}(\mathcal{O}_X))
\longrightarrow
D_{\textit{Coh}}(\mathcal{O}_X)
\end{equation}
见 Derived Categories, Equation (\ref{derived-equation-compare})。

\begin{lemma}
\label{perfect-lemma-coh-to-qcoh}
设 $X$ 为诺特概形。则函子
$$
D^-(\textit{Coh}(\mathcal{O}_X))
\longrightarrow
D^-_{\textit{Coh}(\mathcal{O}_X)}(\QCoh(\mathcal{O}_X))
$$
是等价。
\end{lemma}

\begin{proof}
注意，$\textit{Coh}(\mathcal{O}_X) \subset \QCoh(\mathcal{O}_X)$ 是 Serre
子范畴，见 Homology, Definition \ref{homology-definition-serre-subcategory}、
Lemma \ref{homology-lemma-characterize-serre-subcategory} 及 Cohomology of
Schemes, Lemmas \ref{coherent-lemma-coherent-abelian-Noetherian} 和
\ref{coherent-lemma-coherent-Noetherian-quasi-coherent-sub-quotient}。另一方面，
若 $\mathcal{G} \to \mathcal{F}$ 是从拟凝聚 $\mathcal{O}_X$-模到凝聚
$\mathcal{O}_X$-模的满射，则存在凝聚子模
$\mathcal{G}' \subset \mathcal{G}$ 满射到 $\mathcal{F}$。事实上，由
Properties, Lemma \ref{properties-lemma-quasi-coherent-colimit-finite-type}，
可以把 $\mathcal{G}$ 写成其凝聚子模的滤过并，其中必有一个满足要求。因此
引理由 Derived Categories, Lemma
\ref{derived-lemma-fully-faithful-embedding} 得到。
\end{proof}

\begin{proposition}
\label{perfect-proposition-DCoh}
设 $X$ 为诺特概形。则函子
$$
D^-(\textit{Coh}(\mathcal{O}_X))
\longrightarrow
D^-_{\textit{Coh}}(\mathcal{O}_X)
\quad\text{且}\quad
D^b(\textit{Coh}(\mathcal{O}_X))
\longrightarrow
D^b_{\textit{Coh}}(\mathcal{O}_X)
$$
都是等价。
\end{proposition}

\begin{proof}
考虑交换图表
$$
\xymatrix{
D^-(\textit{Coh}(\mathcal{O}_X)) \ar[r] \ar[d] &
D^-_{\textit{Coh}}(\mathcal{O}_X) \ar[d] \\
D^-(\QCoh(\mathcal{O}_X)) \ar[r] &
D^-_\QCoh(\mathcal{O}_X)
}
$$
由引理 \ref{perfect-lemma-coh-to-qcoh}，左侧竖直箭头完全忠实。由命题
\ref{perfect-proposition-Noetherian}，下方箭头是等价。按构造，右侧竖直箭头完全
忠实。因此上方水平箭头完全忠实。若 $K$ 是
$D^-_{\textit{Coh}}(\mathcal{O}_X)$ 的对象，则由命题
\ref{perfect-proposition-Noetherian} 与之对应的 $D^-(\QCoh(\mathcal{O}_X))$ 的对象
$K'$ 具有凝聚上同调层。因而由引理 \ref{perfect-lemma-coh-to-qcoh}，$K'$ 属于左侧
竖直箭头的本质像，从而上方水平箭头本质满。这完成了有上界情形的证明；
有界情形立即由有上界情形得到。
\end{proof}

\begin{lemma}
\label{perfect-lemma-direct-image-coherent}
设 $S$ 为诺特概形，$f : X \to S$ 为局部有限型概形态射。设 $E$ 为
$D^b_{\textit{Coh}}(\mathcal{O}_X)$ 的对象，使得对所有 $i$，$H^i(E)$ 的
支集在 $S$ 上固有。则 $Rf_*E$ 是
$D^b_{\textit{Coh}}(\mathcal{O}_S)$ 的对象。
\end{lemma}

\begin{proof}
考虑谱序列
$$
R^pf_*H^q(E) \Rightarrow R^{p + q}f_*E
$$
见 Derived Categories, Lemma \ref{derived-lemma-two-ss-complex-functor}。由假设
及 Cohomology of Schemes, Lemma
\ref{coherent-lemma-support-proper-over-base-pushforward}
，层 $R^pf_*H^q(E)$ 凝聚。因此 $R^{p + q}f_*E$ 凝聚，即
$Rf_*E \in D_{\textit{Coh}}(\mathcal{O}_S)$。有下界性显然；有上界性由
Cohomology of Schemes, Lemma
\ref{coherent-lemma-quasi-coherence-higher-direct-images}
或引理 \ref{perfect-lemma-quasi-coherence-direct-image} 得到。
\end{proof}

\begin{lemma}
\label{perfect-lemma-direct-image-coherent-bdd-below}
设 $S$ 为诺特概形，$f : X \to S$ 为局部有限型概形态射。设 $E$ 为
$D^+_{\textit{Coh}}(\mathcal{O}_X)$ 的对象，使得对所有 $i$，$H^i(E)$ 的
支集在 $S$ 上固有。则 $Rf_*E$ 是
$D^+_{\textit{Coh}}(\mathcal{O}_S)$ 的对象。
\end{lemma}

\begin{proof}
证明与引理 \ref{perfect-lemma-direct-image-coherent} 的证明相同。也可以考察正合函子
$Rf_*$ 对可区别三角
$\tau_{\leq a}E \to E \to \tau_{\geq a + 1}E \to \tau_{\leq a}E[1]$.
的作用，由引理 \ref{perfect-lemma-direct-image-coherent} 推出结论。
\end{proof}

\begin{lemma}
\label{perfect-lemma-coherent-internal-hom}
设 $X$ 为局部诺特概形。若 $L$ 属于
$D^+_{\textit{Coh}}(\mathcal{O}_X)$ 且 $K$ 属于
$D^-_{\textit{Coh}}(\mathcal{O}_X)$，则 $R\SheafHom(K, L)$ 属于
$D^+_{\textit{Coh}}(\mathcal{O}_X)$。
\end{lemma}

\begin{proof}
只需在 $X$ 是诺特环 $A$ 的谱时证明。由引理
\ref{perfect-lemma-identify-pseudo-coherent-noetherian}，$K$ 伪凝聚。再由引理
\ref{perfect-lemma-quasi-coherence-internal-hom}，问题转化为如下代数问题：若
$L \in D^+_{\textit{Coh}}(A)$ 且 $K$ 属于
$D^-_{\textit{Coh}}(A)$，则 $R\Hom_A(K, L)$ 属于
$D^+_{\textit{Coh}}(A)$。由于 $L$ 有下界而 $K$ 有上界，存在收敛谱序列
$$
\Ext^p_A(K, H^q(L)) \Rightarrow \text{Ext}^{p + q}_A(K, L)
$$
以及收敛谱序列
$$
\Ext^i_A(H^{-j}(K), H^q(L)) \Rightarrow \text{Ext}^{i + j}_A(K, H^q(L))
$$
见 Injectives, Remarks \ref{injectives-remark-spectral-sequences-ext} 和
\ref{injectives-remark-spectral-sequences-ext-variant}。由 Algebra, Lemma
\ref{algebra-lemma-ext-noetherian}，对有限 $A$-模 $M$、$N$，模
$\Ext^p_A(M, N)$ 有限，故证明完成。
\end{proof}

\begin{lemma}
\label{perfect-lemma-perfect-on-noetherian}
设 $X$ 为诺特概形，$D(\mathcal{O}_X)$ 中的 $E$ 完美。则
\begin{enumerate}
\item $E$ 属于 $D^b_{\textit{Coh}}(\mathcal{O}_X)$；
\item 若 $L$ 属于 $D_{\textit{Coh}}(\mathcal{O}_X)$，则
$E \otimes_{\mathcal{O}_X}^\mathbf{L} L$ 和
$R\SheafHom_{\mathcal{O}_X}(E, L)$ 都属于
$D_{\textit{Coh}}(\mathcal{O}_X)$；
\item 若 $L$ 属于 $D^b_{\textit{Coh}}(\mathcal{O}_X)$，则
$E \otimes_{\mathcal{O}_X}^\mathbf{L} L$ 和
$R\SheafHom_{\mathcal{O}_X}(E, L)$ 都属于
$D^b_{\textit{Coh}}(\mathcal{O}_X)$；
\item 若 $L$ 属于 $D^+_{\textit{Coh}}(\mathcal{O}_X)$，则
$E \otimes_{\mathcal{O}_X}^\mathbf{L} L$ 和
$R\SheafHom_{\mathcal{O}_X}(E, L)$ 都属于
$D^+_{\textit{Coh}}(\mathcal{O}_X)$；
\item 若 $L$ 属于 $D^-_{\textit{Coh}}(\mathcal{O}_X)$，则
$E \otimes_{\mathcal{O}_X}^\mathbf{L} L$ 和
$R\SheafHom_{\mathcal{O}_X}(E, L)$ 都属于
$D^-_{\textit{Coh}}(\mathcal{O}_X)$.
\end{enumerate}
\end{lemma}

\begin{proof}
由于 $X$ 拟紧，每个断言都可在 $X$ 任意开覆盖的各成员上检验。因此可假设
$E$ 由有限自由模的有界复形 $\mathcal{E}^\bullet$ 表示，见 Cohomology,
Lemma \ref{cohomology-lemma-perfect-on-locally-ringed}。此时各断言都很清楚，
因为 $R\SheafHom_{\mathcal{O}_X}(E, L)$ 和
$E \otimes_{\mathcal{O}_X}^\mathbf{L} L$ 都可用 $\mathcal{E}^\bullet$ 在复形
层面计算，见 Cohomology, Lemmas
\ref{cohomology-lemma-Rhom-strictly-perfect} 和
\ref{cohomology-lemma-bounded-flat-K-flat}。略去一些细节。
\end{proof}

\begin{lemma}
\label{perfect-lemma-ext-finite}
设 $A$ 为诺特环，$X$ 为 $A$ 上的固有概形。对属于
$D^+_{\textit{Coh}}(\mathcal{O}_X)$ 的 $L$ 和属于
$D^-_{\textit{Coh}}(\mathcal{O}_X)$ 的 $K$，$A$-模
$\Ext_{\mathcal{O}_X}^n(K, L)$ 是有限的。
\end{lemma}

\begin{proof}
回忆
$$
\Ext_{\mathcal{O}_X}^n(K, L) =
H^n(X, R\SheafHom_{\mathcal{O}_X}(K, L)) =
H^n(\Spec(A), Rf_*R\SheafHom_{\mathcal{O}_X}(K, L))
$$
见 Cohomology, Lemma \ref{cohomology-lemma-section-RHom-over-U} 和 Cohomology,
Section \ref{cohomology-section-Leray}。因此结论由引理
\ref{perfect-lemma-coherent-internal-hom} 和
\ref{perfect-lemma-direct-image-coherent-bdd-below} 得到。
\end{proof}

\begin{lemma}
\label{perfect-lemma-perfect-on-regular}
设 $X$ 为正则概形。则 $D^b_{\textit{Coh}}(\mathcal{O}_X)$ 的每个对象都
完美。若 $X$ 拟紧，即诺特且正则，则反过来，$D(\mathcal{O}_X)$ 的每个完美
对象都属于 $D^b_{\textit{Coh}}(\mathcal{O}_X)$。
\end{lemma}

\begin{proof}
设 $K$ 为 $D^b_{\textit{Coh}}(\mathcal{O}_X)$ 的对象。为检验 $K$ 完美，可以
在 $X$ 上仿射局部地工作（见 Cohomology, Section
\ref{cohomology-section-perfect}）。于是由引理 \ref{perfect-lemma-perfect-affine} 和
More on Algebra, Lemma \ref{more-algebra-lemma-regular-perfect}，$K$ 完美。
反向即引理 \ref{perfect-lemma-perfect-on-noetherian}。
\end{proof}




\section{复形有限性性质的下降}
\label{perfect-section-descent-finiteness}

\noindent
本节是 Descent, Section \ref{descent-section-descent-finiteness} 对概形导出
范畴对象的类似版本。其中最容易的结果或许如下。

\begin{lemma}
\label{perfect-lemma-tor-amplitude-descends}
设 $f : X \to Y$ 为概形（或更一般地，局部带环空间）的满射平坦态射。设
$E \in D(\mathcal{O}_Y)$ 且 $a, b \in \mathbf{Z}$。则 $E$ 的 Tor 振幅在
$[a, b]$ 内，当且仅当 $Lf^*E$ 的 Tor 振幅在 $[a, b]$ 内。
\end{lemma}

\begin{proof}
拉回总保持 Tor 振幅，见 Cohomology, Lemma
\ref{cohomology-lemma-tor-amplitude-pullback}。Tor 振幅在 $[a, b]$ 内这一性质
可在茎上检验，见 Cohomology, Lemma
\ref{cohomology-lemma-tor-amplitude-stalk}。由 Algebra, Lemma
\ref{algebra-lemma-local-flat-ff}，平坦局部环同态忠实平坦。因此结论由 More
on Algebra, Lemma \ref{more-algebra-lemma-flat-descent-tor-amplitude} 得到。
\end{proof}

\begin{lemma}
\label{perfect-lemma-pseudo-coherent-descends-fpqc}
设 $\{f_i : X_i \to X\}$ 为概形的 fpqc 覆盖，
$E \in D_\QCoh(\mathcal{O}_X)$，且 $m \in \mathbf{Z}$。则 $E$ 是
$m$-伪凝聚的，当且仅当每个 $Lf_i^*E$ 都是 $m$-伪凝聚的。
\end{lemma}

\begin{proof}
拉回总保持 $m$-伪凝聚性，见 Cohomology, Lemma
\ref{cohomology-lemma-pseudo-coherent-pullback}。反之，假设对所有 $i$，
$Lf_i^*E$ 都是 $m$-伪凝聚的。设 $U \subset X$ 为仿射开集。只需证明
$E|_U$ 是 $m$-伪凝聚的。由于 $\{f_i : X_i \to X\}$ 是 fpqc 覆盖，可以
找到有限多个仿射开集 $V_j \subset X_{a(j)}$，使得
$f_{a(j)}(V_j) \subset U$ 且 $U = \bigcup f_{a(j)}(V_j)$。令
$V = \coprod V_i$。因此可以用 $U$ 替换 $X$，用 $\{V \to U\}$ 替换
$\{f_i : X_i \to X\}$，并假设 $X$ 仿射且覆盖由仿射概形间的单个满射平坦
态射 $\{f : Y \to X\}$ 给出。此时结论由 More on Algebra, Lemma
\ref{more-algebra-lemma-flat-descent-pseudo-coherent} 以及引理
\ref{perfect-lemma-affine-compare-bounded} 和 \ref{perfect-lemma-pseudo-coherent-affine} 得到。
\end{proof}

\begin{lemma}
\label{perfect-lemma-pseudo-coherent-descends-fppf}
设 $\{f_i : X_i \to X\}$ 为概形的 fppf 覆盖，$E \in D(\mathcal{O}_X)$，
且 $m \in \mathbf{Z}$。则 $E$ 是 $m$-伪凝聚的，当且仅当每个 $Lf_i^*E$
都是 $m$-伪凝聚的。
\end{lemma}

\begin{proof}
拉回总保持 $m$-伪凝聚性，见 Cohomology, Lemma
\ref{cohomology-lemma-pseudo-coherent-pullback}。反之，假设对所有 $i$，
$Lf_i^*E$ 都是 $m$-伪凝聚的。设 $U \subset X$ 为仿射开集。只需证明
$E|_U$ 是 $m$-伪凝聚的。由于 $\{f_i : X_i \to X\}$ 是 fppf 覆盖，可以
找到有限多个仿射开集 $V_j \subset X_{a(j)}$，使得
$f_{a(j)}(V_j) \subset U$ 且 $U = \bigcup f_{a(j)}(V_j)$。令
$V = \coprod V_i$。因此可以用 $U$ 替换 $X$，用 $\{V \to U\}$ 替换
$\{f_i : X_i \to X\}$，并假设 $X$ 仿射且覆盖由单个有限表示满射平坦态射
$\{f : Y \to X\}$ 给出。

\medskip\noindent
由于 $f$ 平坦，导出函子 $Lf^*$ 就由 $f^*$ 给出，且 $f^*$ 正合。因此
$H^i(Lf^*E) = f^*H^i(E)$。由于 $Lf^*E$ 是 $m$-伪凝聚的，有
$Lf^*E \in D^-(\mathcal{O}_Y)$。由于 $f$ 满射且平坦，有
$E \in D^-(\mathcal{O}_X)$。设 $i \in \mathbf{Z}$ 为使 $H^i(E)$ 非零的最大
整数。若 $i < m$，则已经完成。否则，由 Cohomology, Lemma
\ref{cohomology-lemma-finite-cohomology}，$f^*H^i(E)$ 是有限型
$\mathcal{O}_Y$-模。再由 Descent, Lemma
\ref{descent-lemma-finite-type-descends-fppf}，$\mathcal{O}_X$-模 $H^i(E)$
是有限型的。因此用一个有限仿射开覆盖的各成员替换 $X$ 后，可假设存在映射
$$
\alpha : \mathcal{O}_X^{\oplus n}[-i] \longrightarrow E
$$
使 $H^i(\alpha)$ 为满射。设 $C$ 为 $D(\mathcal{O}_X)$ 中 $\alpha$ 的锥。
拉回到 $Y$ 并使用 Cohomology, Lemma
\ref{cohomology-lemma-cone-pseudo-coherent}，可知 $Lf^*C$ 是 $m$-伪凝聚的。
此外，对 $j \geq i$ 有 $H^j(C) = 0$。因此对 $i$ 归纳可知 $C$ 是
$m$-伪凝聚的。再次使用 Cohomology, Lemma
\ref{cohomology-lemma-cone-pseudo-coherent} 即得结论。
\end{proof}

\begin{lemma}
\label{perfect-lemma-perfect-descends-fpqc}
设 $\{f_i : X_i \to X\}$ 为概形的 fpqc 覆盖，$E \in D(\mathcal{O}_X)$。
则 $E$ 完美，当且仅当每个 $Lf_i^*E$ 都完美。
\end{lemma}

\begin{proof}
拉回总保持完美复形，见 Cohomology, Lemma
\ref{cohomology-lemma-perfect-pullback}。反之，假设对所有 $i$，$Lf_i^*E$
都完美。于是由引理 \ref{perfect-lemma-pseudo-coherent} 和 Cohomology, Lemma
\ref{cohomology-lemma-perfect}，每个 $Lf_i^*E$ 的上同调层都拟凝聚。由于态射
$f_i$ 平坦，有 $H^p(Lf_i^*E) = f_i^*H^p(E)$。因此由 Descent,
Proposition \ref{descent-proposition-fpqc-descent-quasi-coherent}，$E$ 的上同调
层拟凝聚。至此，引理由 Cohomology, Lemma
\ref{cohomology-lemma-perfect} 以及引理 \ref{perfect-lemma-tor-amplitude-descends} 和
\ref{perfect-lemma-pseudo-coherent-descends-fpqc} 形式地得到。
\end{proof}

\begin{lemma}
\label{perfect-lemma-closed-push-pseudo-coherent}
设 $i : Z \to X$ 为带环空间态射，使得 $i$ 在底拓扑空间上是闭浸入，且
$i_*\mathcal{O}_Z$ 作为 $\mathcal{O}_X$-模伪凝聚。设
$E \in D(\mathcal{O}_Z)$。则 $E$ 是 $m$-伪凝聚的，当且仅当 $Ri_*E$ 是
$m$-伪凝聚的。
\end{lemma}

\begin{proof}
整个证明中都将使用 $i_*$ 是正合函子、因而 $Ri_* = i_*$ 这一事实，见
Modules, Lemma \ref{modules-lemma-i-star-exact}。

\medskip\noindent
假设 $E$ 是 $m$-伪凝聚的。设 $x \in X$。我们要找 $x$ 的一个邻域，使
$i_*E$ 在其上是 $m$-伪凝聚的。若 $x \not \in Z$，则显然。故可假设
$x \in Z$。我们将使用如下事实：当 $x \in U \subset X$ 且该子集开时，
$U \cap Z$ 构成 $x$ 在 $Z$ 中的一个基本邻域系。缩小 $X$ 后，可假设 $E$
有上界。我们对使 $H^p(E)$ 非零的最大整数 $p$ 作归纳。若 $p < m$，则
无须证明。若 $p \geq m$，则由 Cohomology, Lemma
\ref{cohomology-lemma-finite-cohomology}，$H^p(E)$ 是有限型
$\mathcal{O}_Z$-模。因此缩小 $X$ 后，可选择映射
$\mathcal{O}_Z^{\oplus n}[-p] \to E$，它诱导满射
$\mathcal{O}_Z^{\oplus n} \to H^p(E)$。选择可区别三角
$$
\mathcal{O}_Z^{\oplus n}[-p] \to E \to C \to \mathcal{O}_Z^{\oplus n}[-p + 1]
$$
由 Cohomology, Lemma \ref{cohomology-lemma-cone-pseudo-coherent} 可知，对
$j \geq p$ 有 $H^j(C) = 0$，且 $C$ 是 $m$-伪凝聚的。由归纳可知 $i_*C$
在 $X$ 上是 $m$-伪凝聚的。由于 $i_*\mathcal{O}_Z$ 在 $X$ 上也是
$m$-伪凝聚的，由可区别三角
$$
i_*\mathcal{O}_Z^{\oplus n}[-p] \to i_*E \to i_*C \to
i_*\mathcal{O}_Z^{\oplus n}[-p + 1]
$$
及 Cohomology, Lemma \ref{cohomology-lemma-cone-pseudo-coherent} 可知，
$i_*E$ 是 $m$-伪凝聚的。

\medskip\noindent
反之，假设 $i_*E$ 是 $m$-伪凝聚的。设 $z \in Z$。我们要找 $z$ 的一个
邻域，使 $E$ 在其上是 $m$-伪凝聚的。我们将使用如下事实：当
$z \in U \subset X$ 且该子集开时，$U \cap Z$ 构成 $z$ 在 $Z$ 中的一个基本
邻域系。缩小 $X$ 后，可假设 $i_*E$、从而 $E$ 有上界。我们对使
$H^p(E)$ 非零的最大整数 $p$ 作归纳。若 $p < m$，则无须证明。若
$p \geq m$，则由 Cohomology, Lemma
\ref{cohomology-lemma-finite-cohomology}，
$H^p(i_*E) = i_*H^p(E)$ 是有限型 $\mathcal{O}_X$-模。选择一个表示 $E$ 的
$\mathcal{O}_Z$-模复形 $\mathcal{E}^\bullet$。缩小 $X$ 后，可选择映射
$\alpha : \mathcal{O}_X^{\oplus n}[-p] \to i_*\mathcal{E}^\bullet$，它诱导
满射 $\mathcal{O}_X^{\oplus n} \to i_*H^p(\mathcal{E}^\bullet)$。由伴随得到
映射 $\alpha : \mathcal{O}_Z^{\oplus n}[-p] \to \mathcal{E}^\bullet$，它诱导
满射 $\mathcal{O}_Z^{\oplus n} \to H^p(\mathcal{E}^\bullet)$。选择可区别三角
$$
\mathcal{O}_Z^{\oplus n}[-p] \to E \to C \to \mathcal{O}_Z^{\oplus n}[-p + 1]
$$
可知对 $j \geq p$ 有 $H^j(C) = 0$。由可区别三角
$$
i_*\mathcal{O}_Z^{\oplus n}[-p] \to i_*E \to i_*C \to
i_*\mathcal{O}_Z^{\oplus n}[-p + 1]
$$
、$i_*\mathcal{O}_Z$ 伪凝聚这一事实及 Cohomology, Lemma
\ref{cohomology-lemma-cone-pseudo-coherent}，可知 $i_*C$ 是 $m$-伪凝聚的。
由归纳，$C$ 是 $m$-伪凝聚的。再次应用 Cohomology, Lemma
\ref{cohomology-lemma-cone-pseudo-coherent}，可知 $E$ 是 $m$-伪凝聚的。
\end{proof}

\begin{lemma}
\label{perfect-lemma-finite-push-pseudo-coherent}
设 $f : X \to Y$ 为概形的有限态射，使得 $f_*\mathcal{O}_X$ 作为
$\mathcal{O}_Y$-模伪凝聚\footnote{这意味着 $f$ 是伪凝聚态射，见 More on
Morphisms, Lemma
\ref{more-morphisms-lemma-finite-pseudo-coherent}.}.
设 $E \in D_\QCoh(\mathcal{O}_X)$。则 $E$ 是 $m$-伪凝聚的，当且仅当
$Rf_*E$ 是 $m$-伪凝聚的。
\end{lemma}

\begin{proof}
这是把 More on Algebra, Lemma
\ref{more-algebra-lemma-finite-push-pseudo-coherent} 转写为概形语言。转写时使用
引理 \ref{perfect-lemma-affine-compare-bounded} 和
\ref{perfect-lemma-pseudo-coherent-affine}。
\end{proof}


\section{复形的提升}
\label{perfect-section-lift}

\noindent
设 $U \subset X$ 为带环空间的开子空间，并以 $j : U \to X$ 表示包含态射。
由于 $Rj_*$ 是限制函子的右逆，函子
$D(\mathcal{O}_X) \to D(\mathcal{O}_U)$ 本质满。本节把这一结论推广到具有
拟凝聚上同调层的复形等情形。

\begin{lemma}
\label{perfect-lemma-lift-quasi-coherent}
设 $X$ 为概形，$j : U \to X$ 为拟紧开浸入。函子
$$
D_\QCoh(\mathcal{O}_X) \to D_\QCoh(\mathcal{O}_U)
\quad\text{且}\quad
D^+_\QCoh(\mathcal{O}_X) \to D^+_\QCoh(\mathcal{O}_U)
$$
本质满。若 $X$ 拟紧，则函子
$$
D^-_\QCoh(\mathcal{O}_X) \to D^-_\QCoh(\mathcal{O}_U)
\quad\text{且}\quad
D^b_\QCoh(\mathcal{O}_X) \to D^b_\QCoh(\mathcal{O}_U)
$$
也本质满。
\end{lemma}

\begin{proof}
由引理 \ref{perfect-lemma-quasi-coherence-direct-image}，$Rj_*$ 把
$D_\QCoh(\mathcal{O}_U)$ 映到 $D_\QCoh(\mathcal{O}_X)$，故引理前的论证适用于
第一种情形。显然，$Rj_*$ 把 $D^+_\QCoh(\mathcal{O}_U)$ 映到
$D^+_\QCoh(\mathcal{O}_X)$，这给出关于有下界复形的断言。最后，引理
\ref{perfect-lemma-quasi-coherence-direct-image} 保证当 $X$ 拟紧时，$Rj_*$ 把
$D^-_\QCoh(\mathcal{O}_U)$ 映到 $D^-_\QCoh(\mathcal{O}_X)$。结合二者即得
最后一个断言。
\end{proof}

\begin{lemma}
\label{perfect-lemma-lift-coherent}
设 $X$ 为诺特概形，$j : U \to X$ 为开浸入。函子
$D^b_{\textit{Coh}}(\mathcal{O}_X) \to D^b_{\textit{Coh}}(\mathcal{O}_U)$
本质满。
\end{lemma}

\begin{proof}
设 $K$ 为 $D^b_{\textit{Coh}}(\mathcal{O}_U)$ 的对象。由命题
\ref{perfect-proposition-DCoh}，可用凝聚 $\mathcal{O}_U$-模的有界复形
$\mathcal{F}^\bullet$ 表示 $K$。设对某些 $a \leq b$，当
$i \not \in [a, b]$ 时 $\mathcal{F}^i = 0$。由于 $j$ 拟紧且分离，有界复形
$j_*\mathcal{F}^\bullet$ 的各项是 $X$ 上的拟凝聚模，见 Schemes, Lemma
\ref{schemes-lemma-push-forward-quasi-coherent}。如下归纳选取凝聚子模
$\mathcal{G}^i \subset j_*\mathcal{F}^i$。当 $i = a$ 时，任选一个在 $U$
上的限制为 $\mathcal{F}^a$ 的凝聚子模
$\mathcal{G}^a \subset j_*\mathcal{F}^a$；这由 Properties, Lemma
\ref{properties-lemma-extend} 可行。当 $i > a$ 时，先任选一个在 $U$ 上的
限制为 $\mathcal{F}^i$ 的凝聚子模
$\mathcal{H}^i \subset j_*\mathcal{F}^i$，然后令
$\mathcal{G}^i = \Im(\mathcal{H}^i \oplus \mathcal{G}^{i - 1}
\to j_*\mathcal{F}^i)$。显然，
$\mathcal{G}^\bullet \subset j_*\mathcal{F}^\bullet$ 是凝聚
$\mathcal{O}_X$-模的有界复形，其在 $U$ 上的限制为
$\mathcal{F}^\bullet$，正如所需。
\end{proof}

\begin{lemma}
\label{perfect-lemma-lift-pseudo-coherent}
设 $X$ 为仿射概形，$U \subset X$ 为拟紧开子概形。对
$D(\mathcal{O}_U)$ 中任意伪凝聚对象 $E$，存在有限自由
$\mathcal{O}_X$-模的有上界复形，其在 $U$ 上的限制同构于 $E$。
\end{lemma}

\begin{proof}
由引理 \ref{perfect-lemma-pseudo-coherent}，$E$ 是 $D_\QCoh(\mathcal{O}_U)$ 的对象。
由引理 \ref{perfect-lemma-lift-quasi-coherent}，可假设对
$D_\QCoh(\mathcal{O}_X)$ 的某个对象 $E'$ 有 $E = E'|U$。写成
$X = \Spec(A)$。由引理 \ref{perfect-lemma-affine-compare-bounded}，可以找到
$A$-模复形 $M^\bullet$，其相伴的 $\mathcal{O}_X$-模复形是 $E'$ 的一个
代表。

\medskip\noindent
选择 $f_1, \ldots, f_r \in A$，使得
$U = D(f_1) \cup \ldots \cup D(f_r)$。由引理
\ref{perfect-lemma-pseudo-coherent-affine}，复形 $M^\bullet_{f_j}$ 是
$A_{f_j}$-模的伪凝聚复形。设 $n$ 为整数。假设已有复形映射
$\alpha : F^\bullet \to M^\bullet$，其中 $F^\bullet$ 有上界，对 $i < n$
有 $F^i = 0$，且每个 $F^i$ 都是有限自由 $R$-模，并使得
$$
H^i(\alpha_{f_j}) : H^i(F^\bullet_{f_j}) \to H^i(M^\bullet_{f_j})
$$
在 $i > n$ 时为同构，在 $i = n$ 时为满射。图示如下：
$$
\xymatrix{
& F^n \ar[r] \ar[d]^\alpha & F^{n + 1} \ar[d]^\alpha \ar[r] & \ldots \\
M^{n-1} \ar[r] & M^n \ar[r] & M^{n + 1} \ar[r] & \ldots
}
$$
由于每个 $M^\bullet_{f_j}$ 的上同调在高次数消失，对 $n \gg 0$ 可以找到
这样的映射。我们将对 $n$ 归纳，把它延拓为复形映射
$F^\bullet \to M^\bullet$，使得对所有 $i$，$H^i(\alpha_{f_j})$ 都是同构。
取 $\widetilde{F^\bullet}$ 即得引理。

\medskip\noindent
归纳步骤是在上图中加入 $F^{n - 1}$。设 $C^\bullet$ 为 $\alpha$ 的锥
（Derived Categories, Definition \ref{derived-definition-cone}）。上同调长正合
列说明，对 $i \geq n$ 有 $H^i(C^\bullet_{f_j}) = 0$。由 More on Algebra,
Lemma \ref{more-algebra-lemma-cone-pseudo-coherent}，$C^\bullet_{f_j}$ 是
$(n - 1)$-伪凝聚的。由 More on Algebra, Lemma
\ref{more-algebra-lemma-finite-cohomology}，
$H^{n - 1}(C^\bullet_{f_j})$ 是有限 $A_{f_j}$-模。选择有限自由 $A$-模
$F^{n - 1}$ 和 $A$-模同态 $\beta : F^{n - 1} \to C^{n - 1}$，使得复合
$F^{n - 1} \to C^{n - 1} \to C^n$ 为零，且 $F^{n - 1}_{f_j}$ 满射到
$H^{n - 1}(C^\bullet_{f_j})$。（略去一些细节；提示：清除分母。）由于
$C^{n - 1} = M^{n - 1} \oplus F^n$，可以写成
$\beta = (\alpha^{n - 1}, -d^{n - 1})$。复合
$F^{n - 1} \to C^{n - 1} \to C^n$ 的消失性说明这些映射组成复形态射
$$
\xymatrix{
& F^{n - 1} \ar[d]^{\alpha^{n - 1}} \ar[r]_{d^{n - 1}} &
F^n \ar[r] \ar[d]^\alpha &
F^{n + 1} \ar[d]^\alpha \ar[r] & \ldots \\
\ldots \ar[r] &
M^{n - 1} \ar[r] & M^n \ar[r] & M^{n + 1} \ar[r] & \ldots
}
$$
此外，这些映射定义可区别三角的态射
$$
\xymatrix{
(F^n \to \ldots) \ar[r] \ar[d] &
(F^{n-1} \to \ldots) \ar[r] \ar[d] &
F^{n-1} \ar[r] \ar[d]_\beta &
(F^n \to \ldots)[1] \ar[d] \\
(F^n \to \ldots) \ar[r] &
M^\bullet \ar[r] &
C^\bullet \ar[r] &
(F^n \to \ldots)[1]
}
$$
因此，对 $\beta$ 的选取蕴含复形映射
$(F^{-1} \to \ldots) \to M^\bullet$ 在 $f_j$ 处局部化后的上同调上，于次数
$\geq n$ 诱导同构，并在次数 $n - 1$ 诱导满射。这完成了引理的证明。
\end{proof}

\noindent
以下两个引理或许应放在别处。

\begin{lemma}
\label{perfect-lemma-vanishing-ext}
设 $X$ 为拟紧拟分离概形，$E \in D^b_\QCoh(\mathcal{O}_X)$。存在整数
$n_0 > 0$，使得对每个有限局部自由 $\mathcal{O}_X$-模 $\mathcal{E}$ 和
每个 $n \geq n_0$，都有
$\Ext^n_{D(\mathcal{O}_X)}(\mathcal{E}, E) = 0$
\end{lemma}

\begin{proof}
回忆 $\Ext^n_{D(\mathcal{O}_X)}(\mathcal{E}, E) =
\Hom_{D(\mathcal{O}_X)}(\mathcal{E}, E[n])$。导出范畴中的态射满足
Mayer--Vietoris，见 Cohomology, Lemma
\ref{cohomology-lemma-mayer-vietoris-hom}。因此，若 $X = U \cup V$，且对某个
界 $n_0$，引理的结论对 $E|_U$、$E|_V$ 和 $E|_{U \cap V}$ 成立，则结论对
$E$ 以 $n_0 + 1$ 为界成立。因此只需在 $X$ 仿射时证明，见 Cohomology of
Schemes, Lemma \ref{coherent-lemma-induction-principle}。

\medskip\noindent
假设 $X = \Spec(A)$ 仿射。选择一个 $A$-模复形 $M^\bullet$，使其相伴的
拟凝聚模复形表示 $E$，见引理 \ref{perfect-lemma-affine-compare-bounded}。把
$\mathcal{E} = \widetilde{P}$ 写成某个 $A$-模 $P$ 的相伴层。由于
$\mathcal{E}$ 有限局部自由，$P$ 是有限投射 $A$-模。于是
\begin{align*}
\Hom_{D(\mathcal{O}_X)}(\mathcal{E}, E[n])
& = 
\Hom_{D(A)}(P, M^\bullet[n]) \\
& =
\Hom_{K(A)}(P, M^\bullet[n]) \\
& =
\Hom_A(P, H^n(M^\bullet))
\end{align*}
第一个等号由引理 \ref{perfect-lemma-affine-compare-bounded} 得到，第二个等号由
Derived Categories, Lemma
\ref{derived-lemma-morphisms-from-projective-complex} 得到，最后一个等号是因为
$\Hom_A(P, -)$ 是正合函子。由于 $E$、从而 $M^\bullet$ 有界，对所有充分大的
$n$，结果为零。
\end{proof}

\noindent
以下引理还可以加强（对所有仅在固定范围内具有非零上同调层的 $L$，其消失
性具有一致性）。

\begin{lemma}
\label{perfect-lemma-ext-from-perfect-into-bounded-QCoh}
设 $X$ 为拟紧拟分离概形，$K$ 为 $D(\mathcal{O}_X)$ 的完美对象。则
\begin{enumerate}
\item 存在整数 $a \leq b$，使得对任意 $L \in D_\QCoh(\mathcal{O}_X)$，
若对 $i \in [a, b]$ 有 $H^i(L) = 0$，则
$\Hom_{D(\mathcal{O}_X)}(K, L) = 0$；
\item 若 $L$ 有界，则除有限多个 $n$ 外，
$\Ext^n_{D(\mathcal{O}_X)}(K, L)$ 均为零。
\end{enumerate}
\end{lemma}

\begin{proof}
由于 $\Ext^n_{D(\mathcal{O}_X)}(K, L) =
\Hom_{D(\mathcal{O}_X)}(K, L[n])$，(2) 由 (1) 得到。下面证明 (1)。由于
$K$ 完美，有
$$
\Hom_{D(\mathcal{O}_X)}(K, L) =
H^0(X, K^\vee \otimes_{\mathcal{O}_X}^\mathbf{L} L)
$$
其中 $K^\vee$ 是 $K$ 的“对偶”完美复形，见 Cohomology, Lemma
\ref{cohomology-lemma-dual-perfect-complex}。注意，由引理
\ref{perfect-lemma-quasi-coherence-tensor-product} 和 \ref{perfect-lemma-pseudo-coherent}
（后者用于说明完美复形具有拟凝聚上同调层），
$K^\vee \otimes_{\mathcal{O}_X}^\mathbf{L} L$ 属于 $D_\QCoh(X)$。设
$K^\vee$ 的 Tor 振幅在 $[a, b]$ 内。则谱序列
$$
E_1^{p, q} = H^p(K^\vee \otimes_{\mathcal{O}_X}^\mathbf{L} H^q(L))
\Rightarrow
H^{p + q}(K^\vee \otimes_{\mathcal{O}_X}^\mathbf{L} L)
$$
说明，若对 $q \in [j - b, j - a]$ 有 $H^q(L) = 0$，则
$H^j(K^\vee \otimes_{\mathcal{O}_X}^\mathbf{L} L)$ 为零。令 $N$ 为
Cohomology of Schemes, Lemma
\ref{coherent-lemma-vanishing-nr-affines-quasi-separated} 中的整数 $d$。如果
上同调层
$$
H^{-N}(K^\vee \otimes_{\mathcal{O}_X}^\mathbf{L} L),
\ H^{-N + 1}(K^\vee \otimes_{\mathcal{O}_X}^\mathbf{L} L),
\ \ldots,
\ H^0(K^\vee \otimes_{\mathcal{O}_X}^\mathbf{L} L)
$$
均为零，则 $H^0(X, K^\vee \otimes_{\mathcal{O}_X}^\mathbf{L} L)$ 消失。事实上，
由刚才引用的引理和引理 \ref{perfect-lemma-application-nice-K-injective}，有
$$
H^0(X, K^\vee \otimes_{\mathcal{O}_X}^\mathbf{L} L) =
H^0(X, \tau_{\geq -N}(K^\vee \otimes_{\mathcal{O}_X}^\mathbf{L} L))
$$
并且由上同调层的消失性，它等于
$H^0(X, \tau_{\geq 1}(K^\vee \otimes_{\mathcal{O}_X}^\mathbf{L} L))$
；由 Derived Categories, Lemma \ref{derived-lemma-negative-vanishing}，后者为
零。因此，若对 $i \in [-b - N, -a]$ 有 $H^i(L) = 0$，则
$\Hom_{D(\mathcal{O}_X)}(K, L)$ 为零。
\end{proof}

\begin{lemma}
\label{perfect-lemma-lift-perfect-complex-plus-locally-free}
设 $X$ 为仿射概形，$U \subset X$ 为拟紧开集。对 $D(\mathcal{O}_U)$ 的每个
完美对象 $E$，存在整数 $r$ 和 $U$ 上的有限局部自由层 $\mathcal{F}$，使得
$\mathcal{F}[-r] \oplus E$ 是 $D(\mathcal{O}_X)$ 的某个完美对象的限制。
\end{lemma}

\begin{proof}
设 $X = \Spec(A)$。回忆完美复形是伪凝聚的，见 Cohomology, Lemma
\ref{cohomology-lemma-perfect}。由引理 \ref{perfect-lemma-lift-pseudo-coherent}，可以
找到有限自由 $A$-模的有上界复形 $\mathcal{F}^\bullet$，使 $E$ 在
$D(\mathcal{O}_U)$ 中同构于 $\mathcal{F}^\bullet|_U$。由 Cohomology, Lemma
\ref{cohomology-lemma-perfect} 以及 $U$ 的拟紧性，$E$ 具有有限 Tor 维数；
设 $E$ 的 Tor 振幅在 $[a, b]$ 内。选取 $r < a$，并令
$$
\mathcal{K} = \Ker(\mathcal{F}^{r} \to \mathcal{F}^{r + 1})
= \Im(\mathcal{F}^{r - 1} \to \mathcal{F}^r).
$$
由于 $E$ 的 Tor 振幅在 $[a, b]$ 内，由 Cohomology, Lemma
\ref{cohomology-lemma-last-one-flat}，$\mathcal{F} = \mathcal{K}|_U$ 平坦。因此
$\mathcal{F}$ 平坦且有限表示，从而有限局部自由（Properties, Lemma
\ref{properties-lemma-finite-locally-free}）。于是
$$
\mathcal{F} \to \mathcal{F}^r|_U \to \mathcal{F}^{r + 1}|_U \to \ldots
$$
是 $U$ 上表示 $E$ 的严格完美复形。另一方面，复形
$P = (\mathcal{F}^r \to \mathcal{F}^{r + 1} \to \ldots )$ 是 $X$ 上的完美
复形。使用愚截断，得到可区别三角
$$
P|_U \to E \to \mathcal{F}[-r - 1] \to (P|_U)[1]
$$
若映射 $E \to \mathcal{F}[-r - 1]$ 在 $D(\mathcal{O}_U)$ 中为零，则
$P|_U = \mathcal{F}[-r - 2] \oplus E$，见 Derived Categories, Lemma
\ref{derived-lemma-split}。例如由引理
\ref{perfect-lemma-ext-from-perfect-into-bounded-QCoh}，当 $r \ll 0$ 时这必定成立。
\end{proof}

\begin{lemma}
\label{perfect-lemma-lift-map}
设 $X$ 为仿射概形，$U \subset X$ 为拟紧开集。设 $E, E'$ 为
$D_\QCoh(\mathcal{O}_X)$ 的对象，其中 $E$ 完美。对每个映射
$\alpha : E|_U \to E'|_U$，存在 $X$ 上的复形映射
$$
E \xleftarrow{\beta} E_1 \xrightarrow{\gamma} E'
$$
，其中 $E_1$ 完美，$\beta : E_1 \to E$ 在 $U$ 上限制为同构，且
$\alpha = \gamma|_U \circ \beta|_U^{-1}$。此外，可以假设对 $X$ 上的某个
完美复形 $I$，有 $E_1 = E \otimes_{\mathcal{O}_X}^\mathbf{L} I$。
\end{lemma}

\begin{proof}
写成 $X = \Spec(A)$ 及 $U = D(f_1) \cup \ldots \cup D(f_r)$。选择表示 $E$
的有限投射 $A$-模的有限复形 $M^\bullet$（引理
\ref{perfect-lemma-perfect-affine}），并选择表示 $E'$ 的 $A$-模复形
$(M')^\bullet$（引理 \ref{perfect-lemma-affine-compare-bounded}）。此时复形
$H^\bullet = \Hom_A(M^\bullet, (M')^\bullet)$ 是 $A$-模复形，其相伴的拟凝聚
$\mathcal{O}_X$-模复形表示 $R\SheafHom(E, E')$，见 Cohomology, Lemma
\ref{cohomology-lemma-Rhom-strictly-perfect}。于是 $\alpha$ 确定
$H^0(U, R\SheafHom(E, E'))$ 的元素 $s$，见 Cohomology, Lemma
\ref{cohomology-lemma-section-RHom-over-U}。由命题
\ref{perfect-proposition-represent-cohomology-class-on-open}，存在 $e$ 及对应于 $s$
的映射
$$
\xi : I^\bullet(f_1^e, \ldots, f_r^e) \to \Hom_A(M^\bullet, (M')^\bullet)
$$
。令 $E_1$ 为与如下复形相伴的拟凝聚 $\mathcal{O}_X$-模复形所对应的对象：
$$
\text{Tot}(I^\bullet(f_1^e, \ldots, f_r^e) \otimes_A M^\bullet)
$$
利用典范映射 $I^\bullet(f_1^e, \ldots, f_r^e) \to A$ 得到
$E_1 \to E$；利用 $\xi$ 和 Cohomology, Lemma
\ref{cohomology-lemma-section-RHom-over-U} 得到 $E_1 \to E'$。
\end{proof}

\begin{lemma}
\label{perfect-lemma-lift-perfect-complex-plus-shift}
设 $X$ 为仿射概形，$U \subset X$ 为拟紧开集。对
$D(\mathcal{O}_U)$ 的每个完美对象 $F$，对象 $F \oplus F[1]$ 是
$D(\mathcal{O}_X)$ 的某个完美对象的限制。
\end{lemma}

\begin{proof}
由引理 \ref{perfect-lemma-lift-perfect-complex-plus-locally-free}，可找到
$D(\mathcal{O}_X)$ 的完美对象 $E$，使得对某个有限局部自由
$\mathcal{O}_U$-模 $\mathcal{F}$ 有 $E|_U = \mathcal{F}[r] \oplus F$。由引理
\ref{perfect-lemma-lift-map}，可找到完美复形的态射 $\alpha : E_1 \to E$，使
$(E_1)|_U \cong E|_U$，且 $\alpha|_U$ 是映射
$$
\left(
\begin{matrix}
\text{id}_{\mathcal{F}[r]} & 0 \\
0 & 0
\end{matrix}
\right)
:
\mathcal{F}[r] \oplus F \to \mathcal{F}[r] \oplus F
$$
于是 $\alpha$ 的锥就是一个所求对象。
\end{proof}

\begin{lemma}
\label{perfect-lemma-perfect-into-support-on-T}
设 $X$ 为拟紧拟分离概形，$f \in \Gamma(X, \mathcal{O}_X)$。对
$D_\QCoh(\mathcal{O}_X)$ 中任意态射 $\alpha : E \to E'$，若
\begin{enumerate}
\item $E$ 完美；
\item $E'$ 支持在 $T = V(f)$ 上，
\end{enumerate}
则存在 $n \geq 0$ 使得 $f^n \alpha  = 0$。
\end{lemma}

\begin{proof}
导出范畴中的态射满足 Mayer--Vietoris，见 Cohomology, Lemma
\ref{cohomology-lemma-mayer-vietoris-hom}。因此，若 $X = U \cup V$ 且引理的
结论对 $f|_U$、$f|_V$ 和 $f|_{U \cap V}$ 成立，则结论也对 $f$ 成立。
所以只需在 $X$ 仿射时证明，见 Cohomology of Schemes, Lemma
\ref{coherent-lemma-induction-principle}。

\medskip\noindent
设 $X = \Spec(A)$。则 $f \in A$。以下不再特别说明地使用引理
\ref{perfect-lemma-affine-compare-bounded} 的等价 $D(A) = D_\QCoh(X)$。用有限投射
$A$-模的有限复形 $P^\bullet$ 表示 $E$；这由引理
\ref{perfect-lemma-perfect-affine} 可行。设 $t$ 为使 $P^t$ 非零的最大整数。可区别
三角
$$
P^t[-t] \to P^\bullet \to \sigma_{\leq t - 1}P^\bullet \to P^t[-t + 1]
$$
说明，对复形 $P^\bullet$ 的长度归纳，可化到 $P^\bullet$ 只有一个非零项的
情形。再使用平移函子，可化到 $P^\bullet$ 仅由位于次数 $0$ 的单个有限投射
$A$-模 $P$ 组成的情形。用 $A$-模复形 $M^\bullet$ 表示 $E'$。于是
$\alpha$ 对应于映射 $P \to H^0(M^\bullet)$。由假设 (2)，模
$H^0(M^\bullet)$ 支持在 $V(f)$ 上，故 $H^0(M^\bullet)$ 的每个元素都被 $f$
的某个幂零化。
由于 $P$ 是有限 $A$-模，对某个 $n$，映射
$f^n\alpha : P \to H^0(M^\bullet)$ 为零，正如所需。
\end{proof}

\begin{lemma}
\label{perfect-lemma-lift-perfect-complex-plus-shift-support}
设 $X$ 为仿射概形，$T \subset X$ 为闭子集，使得 $X \setminus T$ 拟紧，并
设 $U \subset X$ 为拟紧开集。对 $D(\mathcal{O}_U)$ 中支持在 $T \cap U$ 上的
每个完美对象 $F$，对象 $F \oplus F[1]$ 是 $D(\mathcal{O}_X)$ 中某个支持在
$T$ 上的完美对象 $E$ 的限制。
\end{lemma}

\begin{proof}
设 $T = V(g_1, \ldots, g_s)$。用 $g_j$ 的某个幂替换它后，可假设乘以
$g_j$ 在 $F$ 上为零，见引理 \ref{perfect-lemma-perfect-into-support-on-T}。按引理
\ref{perfect-lemma-lift-perfect-complex-plus-shift} 选取 $E$。注意，
$g_j : E \to E$ 在 $U$ 上限制为零。选择可区别三角
$$
E \xrightarrow{g_1} E \to C_1 \to E[1]
$$
由 Derived Categories, Lemma \ref{derived-lemma-split}，对象 $C_1$ 在 $U$ 上
限制为 $F \oplus F[1] \oplus F[1] \oplus F[2]$。此外，由 Derived
Categories, Lemma \ref{derived-lemma-third-map-square-zero}，
$g_1 : C_1 \to C_1$ 的平方为零。事实上，图表
$$
\xymatrix{
E \ar[r] \ar[d]_0 & C_1 \ar[d]_{g_1} \ar[r] & E[1] \ar[d]_0 \\
E \ar[r] & C_1 \ar[r] & E[1]
}
$$
交换，因为复合 $E \xrightarrow{g_1} E \to C_1$ 和
$C_1 \to E[1] \xrightarrow{g_1} E[1]$ 均为零。继续进行，令 $C_{i + 1}$
等于映射 $g_i : C_i \to C_i$ 的锥，得到 $X$ 上支持在 $T$ 上的完美复形
$C_s$，它在 $U$ 上的限制为
$$
F \oplus F[1]^{\oplus s} \oplus F[2]^{\oplus {s \choose 2}}
\oplus \ldots \oplus F[s]
$$
按引理 \ref{perfect-lemma-lift-map} 选择完美复形的态射 $\beta : C' \to C_s$ 和
$\gamma : C' \to C_s$，使 $\beta|_U$ 为同构，且
$\gamma|_U \circ \beta|_U^{-1}$ 是态射
$$
F \oplus F[1]^{\oplus s} \oplus F[2]^{\oplus {s \choose 2}}
\oplus \ldots \oplus F[s]
\to
F \oplus F[1]^{\oplus s} \oplus F[2]^{\oplus {s \choose 2}}
\oplus \ldots \oplus F[s]
$$
，它在除 $F$ 外的所有直和项上都是恒等映射，而在该直和项上为零。由引理
\ref{perfect-lemma-lift-map}，对 $X$ 上的某个完美复形 $I$，还有
$C' = C_s \otimes^\mathbf{L} I$。因此 $g_j^2\text{id}_{C_s}$ 为零蕴含在
$C'$ 上也有同样的结论。故 $C'$ 也支持在 $T$ 上。于是
$\text{Cone}(\gamma)$ 就是所求对象。
\end{proof}

\noindent
下列引理的一个特例见 \cite{Neeman-Grothendieck}。

\begin{lemma}
\label{perfect-lemma-lift-map-from-perfect-complex-with-support}
设 $X$ 为拟紧拟分离概形，$U \subset X$ 为拟紧开集，$T \subset X$ 为闭
子集，且 $X \setminus T$ 在 $X$ 中逆紧。设 $E$ 为
$D_\QCoh(\mathcal{O}_X)$ 的对象，$\alpha : P \to E|_U$ 为映射，其中 $P$
是 $D(\mathcal{O}_U)$ 中支持在 $T \cap U$ 上的完美对象。则存在映射
$\beta : R \to E$，其中 $R$ 是 $D(\mathcal{O}_X)$ 中支持在 $T$ 上的完美
对象，使得 $P$ 是 $D(\mathcal{O}_U)$ 中 $R|_U$ 的直和项，并与 $\alpha$
和 $\beta|_U$ 相容。
\end{lemma}

\begin{proof}
由于 $X$ 拟紧，存在整数 $m$ 及 $X$ 的仿射开集 $V_j$，使得
$X = U \cup V_1 \cup \ldots \cup V_m$。对 $m$ 归纳，可假设 $m = 1$。
换言之，可假设 $X = U \cup V$，其中 $V$ 仿射。由引理
\ref{perfect-lemma-lift-perfect-complex-plus-shift-support}，可以在
$D(\mathcal{O}_V)$ 中选择支持在 $T \cap V$ 上的完美对象 $Q$，以及同构
$Q|_{U \cap V} \to (P \oplus P[1])|_{U \cap V}$。由引理
\ref{perfect-lemma-lift-map}，可用 $Q \otimes^\mathbf{L} I$ 替换 $Q$（它仍支持在
$T \cap V$ 上），并假设映射
$$
Q|_{U \cap V} \to (P \oplus P[1])|_{U \cap V}
\longrightarrow P|_{U \cap V}
\longrightarrow
E|_{U \cap V}
$$
可提升为 $Q \to E|_V$。由 Cohomology, Lemma
\ref{cohomology-lemma-glue}，得到 $D(\mathcal{O}_X)$ 的态射 $a : R \to E$，
使得 $a|_U$ 同构于 $P \oplus P[1] \to E|_U$，且 $a|_V$ 同构于
$Q \to E|_V$。于是 $R$ 完美且支持在 $T$ 上，正如所需。
\end{proof}

\begin{remark}
\label{perfect-remark-addendum}
引理 \ref{perfect-lemma-lift-map-from-perfect-complex-with-support} 的证明说明
$$
R|_U = P \oplus P^{\oplus n_1}[1] \oplus \ldots \oplus P^{\oplus n_m}[m]
$$
对某些 $m \geq 0$ 和 $n_j \geq 0$ 成立。因此 $R|_U$ 的最高次上同调层等于
$P$ 的最高次上同调层。对映射
$P^{\oplus n_1}[1] \oplus \ldots \oplus P^{\oplus n_m}[m] \to R|_U$ 重复该
构造、取锥并使用归纳，可以使 $R|_U$ 与 $P$ 在任意给定次数以上的上同调层
相等。
\end{remark}



\section{用完美复形逼近}
\label{perfect-section-approximation}

\noindent
本节讨论 Neeman 与 Lipman 的一个观察：伪凝聚复形可由完美复形“逼近”。

\begin{definition}
\label{perfect-definition-approximation-holds}
设 $X$ 为概形。考虑三元组 $(T, E, m)$，其中
\begin{enumerate}
\item $T \subset X$ 是闭子集；
\item $E$ 是 $D_\QCoh(\mathcal{O}_X)$ 的对象；
\item $m \in \mathbf{Z}$。
\end{enumerate}
若存在 $D(\mathcal{O}_X)$ 中支持在 $T$ 上的完美对象 $P$ 及映射
$\alpha : P \to E$，使它对 $i > m$ 诱导同构 $H^i(P) \to H^i(E)$，并诱导
满射 $H^m(P) \to H^m(E)$，则称{\it 三元组满足逼近性质} $(T, E, m)$。
\end{definition}

\noindent
逼近性质不可能对每个三元组都成立。事实上，若三元组 $(T, E, m)$ 满足逼近
性质，则显然
\begin{enumerate}
\item $E$ 是 $m$-伪凝聚的，见 Cohomology, Definition
\ref{cohomology-definition-pseudo-coherent}；
\item 对 $i \geq m$，上同调层 $H^i(E)$ 支持在 $T$ 上。
\end{enumerate}
此外，完美复形的“支集”是闭子概形，其补集在 $X$ 中逆紧（略去细节）。
因此若不对 $T$ 作此假设，就不能期待逼近性质成立。这部分地解释了下一个
定义中的条件。

\begin{definition}
\label{perfect-definition-approximation}
设 $X$ 为概形。若对任意满足 $X \setminus T$ 在 $X$ 中逆紧的闭子集
$T \subset X$，都存在整数 $r$，使得对定义
\ref{perfect-definition-approximation-holds} 中每个满足下列条件的三元组
$(T, E, m)$：
\begin{enumerate}
\item $E$ 是 $(m - r)$-伪凝聚的；
\item 对 $i \geq m - r$，$H^i(E)$ 支持在 $T$ 上，
\end{enumerate}
逼近性质都成立，则称 $X$ 上{\it 用完美复形逼近的性质成立}。
\end{definition}

\noindent
我们将证明，用完美复形逼近的性质对拟紧拟分离概形成立。第二个条件看来是
我们的证明方法所必需的。仔细分析以下论证，或许可以把第一个条件弱化为
“$E$ 是 $m$-伪凝聚的”。

\begin{lemma}
\label{perfect-lemma-open}
设 $X$ 为概形，$U \subset X$ 为开子概形，$(T, E, m)$ 为定义
\ref{perfect-definition-approximation-holds} 中的三元组。若
\begin{enumerate}
\item $T \subset U$；
\item $(T, E|_U, m)$ 满足逼近性质；
\item 对 $i \geq m$，层 $H^i(E)$ 支持在 $T$ 上，
\end{enumerate}
则 $(T, E, m)$ 满足逼近性质。
\end{lemma}

\begin{proof}
设 $j : U \to X$ 为包含态射。若 $P \to E|_U$ 是 $U$ 上三元组
$(T, E|_U, m)$ 的逼近，则 $j_!P = Rj_*P \to j_!(E|_U) \to E$ 是 $X$ 上
$(T, E, m)$ 的逼近。见 Cohomology, Lemmas
\ref{cohomology-lemma-pushforward-restriction} 和
\ref{cohomology-lemma-pushforward-perfect}。
\end{proof}

\begin{lemma}
\label{perfect-lemma-approximation-affine}
设 $X$ 为仿射概形。对定义 \ref{perfect-definition-approximation-holds} 中每个三元组
$(T, E, m)$，若存在整数 $r \geq 0$ 满足
\begin{enumerate}
\item $E$ 是 $m$-伪凝聚的；
\item 对 $i \geq m - r + 1$，$H^i(E)$ 支持在 $T$ 上；
\item $X \setminus T$ 是 $r$ 个仿射开集之并，
\end{enumerate}
则逼近性质成立。特别地，仿射概形上用完美复形逼近的性质成立。
\end{lemma}

\begin{proof}
设 $X = \Spec(A)$，并写成 $T = V(f_1, \ldots, f_r)$。（当 $r = 0$ 时，
即 $T = X$，结论立即由引理 \ref{perfect-lemma-pseudo-coherent-affine} 和定义得到。）
设 $(T, E, m)$ 为引理中的三元组。令 $t$ 为使 $H^t(E)$ 非零的最大整数，
并对 $t$ 归纳。初始情形为 $t < m$，此时结论显然。现假设 $t \geq m$。
由 Cohomology, Lemma \ref{cohomology-lemma-finite-cohomology}，层 $H^t(E)$
为有限型。由于它拟凝聚，可由有限多个截面生成（Properties, Lemma
\ref{properties-lemma-finite-type-module}）。对每个
$s \in \Gamma(X, H^t(E)) = H^t(X, E)$（见引理
\ref{perfect-lemma-affine-compare-bounded} 的证明），由引理
\ref{perfect-lemma-represent-cohomology-class-on-closed}，可以找到 $e > 0$ 及态射
$K_e[-t] \to E$，使 $s$ 属于
$H^0(K_e) = H^t(K_e[-t]) \to H^t(E)$ 的像。取这些映射的有限直和，得到
映射 $P \to E$，其中 $P$ 是支持在 $T$ 上的完美复形，对 $i > t$ 有
$H^i(P) = 0$，且 $H^t(P) \to E$ 为满射。选择可区别三角
$$
P \to E \to E' \to P[1]
$$
于是 $E'$ 是 $m$-伪凝聚的（Cohomology, Lemma
\ref{cohomology-lemma-cone-pseudo-coherent}），对 $i \geq t$ 有
$H^i(E') = 0$，并且对 $i \geq m - r + 1$，$H^i(E')$ 支持在 $T$ 上。由
归纳，找到 $(T, E', m)$ 的逼近 $P' \to E'$。把复合
$P' \to E' \to P[1]$ 置于可区别三角 $P \to P'' \to P' \to P[1]$ 中，
并使用 TR3，把态射 $P' \to E'$ 和 $P[1] \to P[1]$ 延拓为可区别三角的态射
$$
\xymatrix{
P \ar[r] \ar[d] & P'' \ar[d] \ar[r] & P' \ar[d] \ar[r] & P[1] \ar[d] \\
P \ar[r] &  E \ar[r] & E' \ar[r] & P[1]
}
$$
。于是 $P''$ 是支持在 $T$ 上的完美复形（Cohomology, Lemma
\ref{cohomology-lemma-two-out-of-three-perfect}）。简单的图追逐说明，
$P'' \to E$ 就是所求逼近。
\end{proof}

\begin{lemma}
\label{perfect-lemma-induction-step}
设 $X$ 为概形，$X = U \cup V$ 为开覆盖，其中 $U$ 拟紧，$V$ 仿射，且
$U \cap V$ 拟紧。若 $U$ 上用完美复形逼近的性质成立，则 $X$ 上也成立。
\end{lemma}

\begin{proof}
设 $T \subset X$ 为闭子集，且 $X \setminus T$ 在 $X$ 中逆紧。设 $r_U$ 为
定义 \ref{perfect-definition-approximation} 中适用于二元组 $(U, T \cap U)$ 的整数。
令 $T' = T \setminus U$。注意，$T' \subset V$；并且由对 $T$ 的假设，
$V \setminus T' = (X \setminus T) \cap U \cap V$ 拟紧。设 $r'$ 为覆盖
$V \setminus T'$ 所需的仿射开集数。我们断言
$r = \max(r_U, r')$ 对二元组 $(X, T)$ 适用。

\medskip\noindent
为此，选择三元组 $(T, E, m)$，使 $E$ 是 $(m - r)$-伪凝聚的，并且对
$i \geq m - r$，$H^i(E)$ 支持在 $T$ 上。令 $t$ 为使 $H^t(E)|_U$ 非零的
最大整数。（这样的整数存在，因为 $U$ 拟紧且 $E|_U$ 是
$(m - r)$-伪凝聚的。）我们对 $t$ 归纳证明 $E$ 可被逼近。

\medskip\noindent
初始情形：$t \leq m - r'$。这意味着对 $i \geq m - r'$，$H^i(E)$ 支持在
$T'$ 上。因此引理 \ref{perfect-lemma-approximation-affine} 保证 $V$ 上存在
$(T', E|_V, m)$ 的逼近 $P \to E|_V$。应用引理 \ref{perfect-lemma-open} 可知，
$(T', E, m)$ 可被逼近。这样的逼近也是 $(T, E, m)$ 的逼近。

\medskip\noindent
归纳步骤。选择 $(T \cap U, E|_U, m)$ 的逼近 $P \to E|_U$。这特别给出满射
$H^t(P) \to H^t(E|_U)$。由引理
\ref{perfect-lemma-lift-perfect-complex-plus-shift-support}，可在
$D(\mathcal{O}_V)$ 中选择支持在 $T \cap V$ 上的完美对象 $Q$，以及同构
$Q|_{U \cap V} \to (P \oplus P[1])|_{U \cap V}$。由引理
\ref{perfect-lemma-lift-map}，可用 $Q \otimes^\mathbf{L} I$ 替换 $Q$，并假设映射
$$
Q|_{U \cap V} \to (P \oplus P[1])|_{U \cap V}
\longrightarrow P|_{U \cap V}
\longrightarrow
E|_{U \cap V}
$$
可提升为 $Q \to E|_V$。由 Cohomology, Lemma
\ref{cohomology-lemma-glue}，得到 $D(\mathcal{O}_X)$ 中的态射 $a : R \to E$，
使得 $a|_U$ 同构于 $P \oplus P[1] \to E|_U$，而 $a|_V$ 同构于
$Q \to E|_V$。因此 $R$ 完美且支持在 $T$ 上，并且映射
$H^t(R) \to H^t(E)$ 限制到 $U$ 上后为满射。选择可区别三角
$$
R \to E \to E' \to R[1]
$$
于是 $E'$ 是 $(m - r)$-伪凝聚的（Cohomology, Lemma
\ref{cohomology-lemma-cone-pseudo-coherent}），对 $i \geq t$ 有
$H^i(E')|_U = 0$，并且对 $i \geq m - r$，$H^i(E')$ 支持在 $T$ 上。由
归纳，找到 $(T, E', m)$ 的逼近 $R' \to E'$。把复合
$R' \to E' \to R[1]$ 置于可区别三角 $R \to R'' \to R' \to R[1]$ 中，
并把态射 $R' \to E'$ 和 $R[1] \to R[1]$ 延拓为可区别三角的态射
$$
\xymatrix{
R \ar[r] \ar[d] & R'' \ar[d] \ar[r] & R' \ar[d] \ar[r] & R[1] \ar[d] \\
R \ar[r] &  E \ar[r] & E' \ar[r] & R[1]
}
$$
，这里使用 TR3。于是 $R''$ 是支持在 $T$ 上的完美复形（Cohomology,
Lemma \ref{cohomology-lemma-two-out-of-three-perfect}）。简单的图追逐说明，
$R'' \to E$ 是所求逼近。
\end{proof}

\begin{theorem}
\label{perfect-theorem-approximation}
设 $X$ 为拟紧拟分离概形。则 $X$ 上用完美复形逼近的性质成立。
\end{theorem}

\begin{proof}
这由 Cohomology of Schemes, Lemma
\ref{coherent-lemma-induction-principle} 的归纳原理，以及引理
\ref{perfect-lemma-induction-step} 和 \ref{perfect-lemma-approximation-affine} 得到。
\end{proof}






\section{生成导出范畴}
\label{perfect-section-generating}

\noindent
本节证明，拟紧拟分离概形的导出范畴 $D_\QCoh(\mathcal{O}_X)$ 可由单个完美
对象生成。我们强烈建议读者阅读 Bondal 与 van den Bergh 那篇精彩论文中对此
结果的证明，见 \cite{BvdB}。

\begin{lemma}
\label{perfect-lemma-direct-summand-of-a-restriction}
设 $X$ 为拟紧拟分离概形，$U$ 为拟紧开子概形，$P$ 为
$D(\mathcal{O}_U)$ 的完美对象。则 $P$ 是 $D(\mathcal{O}_X)$ 某个完美对象之
限制的直和项。
\end{lemma}

\begin{proof}
这是引理 \ref{perfect-lemma-lift-map-from-perfect-complex-with-support} 的特例。
\end{proof}

\begin{lemma}
\label{perfect-lemma-orthogonal-koszul-complex}
\begin{reference}
\cite[Proposition 6.1]{Bokstedt-Neeman}
\end{reference}
在情形 \ref{perfect-situation-complex} 中，以 $j : U \to X$ 表示开浸入，并设 $K$
为 $D(\mathcal{O}_X)$ 中对应于 $A$ 上 $f_1, \ldots, f_r$ 的 Koszul 复形的
完美对象。对 $E \in D_\QCoh(\mathcal{O}_X)$，以下条件等价：
\begin{enumerate}
\item $E = Rj_*(E|_U)$；
\item 对所有 $n \in \mathbf{Z}$，
$\Hom_{D(\mathcal{O}_X)}(K[n], E) = 0$。
\end{enumerate}
\end{lemma}

\begin{proof}
选择可区别三角 $E \to Rj_*(E|_U) \to N \to E[1]$。注意
$$
\Hom_{D(\mathcal{O}_X)}(K[n], Rj_*(E|_U)) =
\Hom_{D(\mathcal{O}_U)}(K|_U[n], E) = 0
$$
对所有 $n$ 都成立，因为 $K|_U = 0$。因此只需对 $N$ 证明。换言之，可假设
$E$ 在 $U$ 上的限制为零。注意，有 Koszul 复形的可区别三角
$$
K^\bullet(f_1^{e_1}, \ldots, f_i^{e'_i}, \ldots, f_r^{e_r}) \to
K^\bullet(f_1^{e_1}, \ldots, f_i^{e'_i + e''_i}, \ldots, f_r^{e_r}) \to
K^\bullet(f_1^{e_1}, \ldots, f_i^{e''_i}, \ldots, f_r^{e_r}) \to \ldots
$$
，见 More on Algebra, Lemma \ref{more-algebra-lemma-koszul-mult}。因此，若对
所有 $n \in \mathbf{Z}$ 有 $\Hom_{D(\mathcal{O}_X)}(K[n], E) = 0$，则用引理
\ref{perfect-lemma-represent-cohomology-class-on-closed} 中的 $K_e$ 替换 $K$ 后，同一
结论仍成立。于是本引理由该引理以及如下事实立即得到：$E$ 由 $A$-模复形
$R\Gamma(X, E)$ 确定，见引理 \ref{perfect-lemma-affine-compare-bounded}。
\end{proof}

\begin{theorem}
\label{perfect-theorem-bondal-van-den-Bergh}
设 $X$ 为拟紧拟分离概形。范畴 $D_\QCoh(\mathcal{O}_X)$ 可由单个完美对象
生成。更确切地说，存在 $D(\mathcal{O}_X)$ 的完美对象 $P$，使得对
$E \in D_\QCoh(\mathcal{O}_X)$，以下条件等价：
\begin{enumerate}
\item $E = 0$；
\item 对所有 $n \in \mathbf{Z}$，
$\Hom_{D(\mathcal{O}_X)}(P[n], E) = 0$。
\end{enumerate}
\end{theorem}

\begin{proof}
我们使用 Cohomology of Schemes, Lemma
\ref{coherent-lemma-induction-principle} 的归纳原理证明。

\medskip\noindent
若 $X$ 仿射，则 $\mathcal{O}_X$ 是完美生成元。这由引理
\ref{perfect-lemma-affine-compare-bounded} 得到。

\medskip\noindent
假设 $X = U \cup V$ 为开覆盖，其中 $U$ 拟紧，定理对 $U$ 成立，且 $V$ 为
仿射开集。设 $P$ 为 $D(\mathcal{O}_U)$ 的完美对象，且是
$D_\QCoh(\mathcal{O}_U)$ 的生成元。由引理
\ref{perfect-lemma-direct-summand-of-a-restriction}，可选择 $D(\mathcal{O}_X)$ 的完美
对象 $Q$，使其在 $U$ 上的限制是一个直和，而 $P$ 是其中一个直和项。设
$V = \Spec(A)$。令 $Z = X \setminus U$。这是 $V$ 的闭子集，且
$V \setminus Z$ 拟紧。选择 $f_1, \ldots, f_r \in A$，使
$Z = V(f_1, \ldots, f_r)$。设 $K \in D(\mathcal{O}_V)$ 为对应于 $A$ 上
$f_1, \ldots, f_r$ 的 Koszul 复形的完美对象。注意，由于 $K$ 支持在闭子集
$Z \subset V$ 上，推前 $K' = R(V \to X)_*K$ 是 $D(\mathcal{O}_X)$ 的完美
对象，其在 $V$ 上的限制为 $K$（见 Cohomology, Lemma
\ref{cohomology-lemma-pushforward-perfect}）。我们断言 $Q \oplus K'$ 是
$D_\QCoh(\mathcal{O}_X)$ 的生成元。

\medskip\noindent
设 $E$ 为 $D_\QCoh(\mathcal{O}_X)$ 的对象，使得从 $Q \oplus K'$ 的任意
平移到 $E$ 都没有非零映射。由 Cohomology, Lemma
\ref{cohomology-lemma-pushforward-restriction}，有
$K' =  R(V \to X)_! K$，因而
$$
\Hom_{D(\mathcal{O}_X)}(K'[n], E) = \Hom_{D(\mathcal{O}_V)}(K[n], E|_V)
$$
于是由引理 \ref{perfect-lemma-orthogonal-koszul-complex}，这些群的消失性蕴含
$E|_V$ 同构于 $R(U \cap V \to V)_*E|_{U \cap V}$。这又蕴含
$E = R(U \to X)_*E|_U$（略去一个小细节）。在此情形下，
$$
\Hom_{D(\mathcal{O}_X)}(Q[n], E) = \Hom_{D(\mathcal{O}_U)}(Q|_U[n], E|_U)
$$
其直和项之一是 $\Hom_{D(\mathcal{O}_U)}(P[n], E|_U)$。因此由对 $P$ 的
选取，这些群的消失性蕴含 $E|_U$ 为零，从而 $E$ 为零。
\end{proof}

\noindent
以下结果是定理 \ref{perfect-theorem-bondal-van-den-Bergh} 的加强，且用完全相同的
方法证明。回忆，对概形 $X$ 的闭子集 $T$，以 $D_T(\mathcal{O}_X)$ 表示
$D(\mathcal{O}_X)$ 中由支持在 $T$ 上的对象组成的严格满、饱和三角子范畴
（定义 \ref{perfect-definition-supported-on}）。类似地，以
$D_{\QCoh, T}(\mathcal{O}_X)$ 表示 $D(\mathcal{O}_X)$ 中由上同调层拟凝聚且
支持在 $T$ 上的复形组成的严格满、饱和三角子范畴。

\begin{lemma}
\label{perfect-lemma-generator-with-support}
\begin{reference}
\cite[Theorem 6.8]{Rouquier-dimensions}
\end{reference}
设 $X$ 为拟紧拟分离概形，$T \subset X$ 为闭子集，且 $X \setminus T$ 拟紧。
采用上述记号，范畴 $D_{\QCoh, T}(\mathcal{O}_X)$ 由单个完美对象生成。
\end{lemma}

\begin{proof}
我们使用 Cohomology of Schemes, Lemma
\ref{coherent-lemma-induction-principle} 的归纳原理证明。

\medskip\noindent
假设 $X = \Spec(A)$ 仿射。此时存在
$f_1, \ldots, f_r \in A$，使 $T = V(f_1, \ldots, f_r)$。设 $K$ 为引理
\ref{perfect-lemma-orthogonal-koszul-complex} 中 $f_1, \ldots, f_r$ 上的 Koszul 复形。
于是 $K$ 是上同调支持在 $T$ 上的完美对象，因而是
$D_{\QCoh, T}(\mathcal{O}_X)$ 的完美对象。另一方面，若
$E \in D_{\QCoh, T}(\mathcal{O}_X)$ 且对所有 $n$ 有
$\Hom(K, E[n]) = 0$，则引理 \ref{perfect-lemma-orthogonal-koszul-complex} 告诉我们
$E = Rj_*(E|_{X \setminus T}) = 0$。因此 $K$ 生成
$D_{\QCoh, T}(\mathcal{O}_X)$（按 Derived Categories, Definition
\ref{derived-definition-generators} 中三角范畴生成元的定义）。

\medskip\noindent
假设 $X = U \cup V$ 为开覆盖，其中 $V$ 仿射、$U$ 拟紧，且引理对 $U$
成立。设 $P$ 为 $D(\mathcal{O}_U)$ 中支持在 $T \cap U$ 上的完美对象，且是
$D_{\QCoh, T \cap U}(\mathcal{O}_U)$ 的生成元。由引理
\ref{perfect-lemma-lift-map-from-perfect-complex-with-support}，可选择
$D(\mathcal{O}_X)$ 中支持在 $T$ 上的完美对象 $Q$，使其在 $U$ 上的限制是
一个直和，而 $P$ 是其中一个直和项。写成 $V = \Spec(B)$。令
$Z = X \setminus U$。则 $Z$ 是 $V$ 的闭子集，且 $V \setminus Z$ 拟紧。
由于 $X$ 拟分离，$Z \cap T$ 是 $V$ 的闭子集，且
$W = V \setminus (Z \cap T)$ 拟紧。因此可以选择
$g_1, \ldots, g_s \in B$，使 $Z \cap T = V(g_1, \ldots, g_r)$。设
$K \in D(\mathcal{O}_V)$ 为对应于 $B$ 上 $g_1, \ldots, g_s$ 的 Koszul 复形
的完美对象。注意，由于 $K$ 支持在闭子集 $(Z \cap T) \subset V$ 上，推前
$K' = R(V \to X)_*K$ 是 $D(\mathcal{O}_X)$ 的完美对象，其在 $V$ 上的限制
为 $K$（见 Cohomology, Lemma
\ref{cohomology-lemma-pushforward-perfect}）。我们断言 $Q \oplus K'$ 是
$D_{\QCoh, T}(\mathcal{O}_X)$ 的生成元。

\medskip\noindent
设 $E$ 为 $D_{\QCoh, T}(\mathcal{O}_X)$ 的对象，使得从 $Q \oplus K'$ 的任意
平移到 $E$ 都没有非零映射。由 Cohomology, Lemma
\ref{cohomology-lemma-pushforward-restriction}，有
$K' =  R(V \to X)_! K$，因而
$$
\Hom_{D(\mathcal{O}_X)}(K'[n], E) = \Hom_{D(\mathcal{O}_V)}(K[n], E|_V)
$$
因此由引理 \ref{perfect-lemma-orthogonal-koszul-complex}，有
$E|_V = Rj_*E|_W$，其中 $j : W \to V$ 是包含态射。图示如下：
$$
\xymatrix{
W \ar[r]_j & V & Z \cap T \ar[l] \ar[d] \\
U \cap V \ar[u]^{j'} \ar[ru]_{j''} & & Z \ar[lu]
}
$$
由于 $E$ 支持在 $T$ 上，$E|_W$ 支持在
$T \cap W = T \cap U \cap V$ 上，而后者在 $W$ 中闭。因此
$$
E|_V = Rj_*(E|_W) = Rj_*(Rj'_*(E|_{U \cap V})) = Rj''_*(E|_{U \cap V})
$$
其中第二个等号是 Cohomology, Lemma
\ref{cohomology-lemma-pushforward-restriction} 的 (1)。这蕴含
$E = R(U \to X)_*E|_U$（略去一个小细节）。在此情形下，
$$
\Hom_{D(\mathcal{O}_X)}(Q[n], E) = \Hom_{D(\mathcal{O}_U)}(Q|_U[n], E|_U)
$$
其直和项之一是 $\Hom_{D(\mathcal{O}_U)}(P[n], E|_U)$。因此由对 $P$ 的
选取，这些群的消失性蕴含 $E|_U$ 为零，从而 $E$ 为零。
\end{proof}



\section{一个生成元例子}
\label{perfect-section-example-generator}

\noindent
本节证明，环上射影空间的导出范畴由一个向量丛生成；事实上，它由结构层的
若干平移之直和生成。

\medskip\noindent
下列引理说明，若 $\mathcal{L}$ 丰沛，则
$\bigoplus_{n \geq 0} \mathcal{L}^{\otimes -n}$ 是生成元。

\begin{lemma}
\label{perfect-lemma-nonzero-some-cohomology}
设 $X$ 为概形，$\mathcal{L}$ 为丰沛可逆 $\mathcal{O}_X$-模。若 $K$ 是
$D_\QCoh(\mathcal{O}_X)$ 的非零对象，则对某些 $n \geq 0$ 和
$p \in \mathbf{Z}$，上同调群
$H^p(X, K \otimes_{\mathcal{O}_X}^\mathbf{L} \mathcal{L}^{\otimes n})$
非零。
\end{lemma}

\begin{proof}
回忆，由于 $X$ 具有丰沛可逆层，它拟紧且分离（Properties, Definition
\ref{properties-definition-ample} 和 Lemma
\ref{properties-lemma-affine-s-opens-cover-quasi-separated}）。因此可应用命题
\ref{perfect-proposition-quasi-compact-affine-diagonal}，用拟凝聚模复形
$\mathcal{F}^\bullet$ 表示 $K$。选取任意 $p$，使
$\mathcal{H}^p = \Ker(\mathcal{F}^p \to \mathcal{F}^{p + 1})/
\Im(\mathcal{F}^{p - 1} \to \mathcal{F}^p)$ 非零。选择点 $x \in X$，使茎
$\mathcal{H}^p_x$ 非零。选择 $n \geq 0$ 及
$s \in \Gamma(X, \mathcal{L}^{\otimes n})$，使 $X_s$ 是 $x$ 的仿射开邻域。
选择 $\tau \in \mathcal{H}^p(X_s)$，使它映到茎 $\mathcal{H}^p_x$ 的非零
元素；这是可行的，因为 $\mathcal{H}^p$ 拟凝聚且 $X_s$ 仿射。由于在
$X_s$ 上取截面是拟凝聚模上的正合函子，可以找到截面
$\tau' \in \mathcal{F}^p(X_s)$，它在 $\mathcal{F}^{p + 1}(X_s)$ 中映为零，
并在 $\mathcal{H}^p(X_s)$ 中映到 $\tau$。由 Properties, Lemma
\ref{properties-lemma-invert-s-sections}，存在 $m$，使
$\tau' \otimes s^{\otimes m}$ 是某个截面
$\tau'' \in \Gamma(X, \mathcal{F}^p \otimes \mathcal{L}^{\otimes mn})$ 的像。
再次应用同一引理，得到 $l \geq 0$，使
$\tau'' \otimes s^{\otimes l}$ 在
$\mathcal{F}^{p + 1} \otimes \mathcal{L}^{\otimes (m + l)n}$ 中映为零。于是
$\tau''$ 给出非零类
$H^p(X, K \otimes^\mathbf{L}_{\mathcal{O}_X} \mathcal{L}^{(m + l)n})$
，正如所需。
\end{proof}

\begin{lemma}
\label{perfect-lemma-construct-the-next-one}
设 $A$ 为环，$X = \mathbf{P}^n_A$。对每个 $a \in \mathbf{Z}$，存在 $X$
上向量丛的正合复形
$$
0 \to \mathcal{O}_X(a) \to \ldots
\to \mathcal{O}_X(a + i)^{\oplus {n + 1 \choose i}} \to
\ldots \to \mathcal{O}_X(a + n + 1) \to 0
$$
\end{lemma}

\begin{proof}
回忆 $\mathbf{P}^n_A$ 是 $\text{Proj}(A[X_0, \ldots, X_n])$，见
Constructions, Definition \ref{constructions-definition-projective-space}。
考虑 $S = A[X_0, \ldots, X_n]$ 上 $X_0, \ldots, X_n$ 的 Koszul 复形
$$
K_\bullet = K_\bullet(A[X_0, \ldots, X_n], X_0, \ldots, X_n)
$$
。由于 $X_0, \ldots, X_n$ 显然是多项式环 $S$ 中的正则序列，由 More on
Algebra, Lemma \ref{more-algebra-lemma-regular-koszul-regular}，Koszul 复形
$K_\bullet$ 除次数 $0$ 外均正合，而此次数的上同调为
$S/(X_0, \ldots, X_n)$。若把 $K_i$ 的生成元置于次数 $+i$，则
$K_\bullet$ 成为分次模复形。换言之，有正合复形
$$
0 \to S(-n - 1) \to \ldots \to S(-n - 1 + i)^{\oplus {n \choose i}} \to \ldots
\to S \to S/(X_0, \ldots, X_n) \to 0
$$
应用 Constructions, Lemma \ref{constructions-lemma-proj-sheaves} 的正合函子
$\tilde{\ }$，并利用最后一项属于此函子的核，得到正合复形
$$
0 \to \mathcal{O}_X(-n - 1) \to \ldots
\to \mathcal{O}_X(-n - 1 + i)^{\oplus {n + 1 \choose i}} \to
\ldots \to \mathcal{O}_X \to 0
$$
再用可逆层 $\mathcal{O}_X(n + a + 1)$ 扭曲，即得引理中的正合复形。
\end{proof}

\begin{lemma}
\label{perfect-lemma-generator-P1}
设 $A$ 为环，$X = \mathbf{P}^n_A$。则
$$
E =
\mathcal{O}_X \oplus \mathcal{O}_X(-1) \oplus \ldots \oplus \mathcal{O}_X(-n)
$$
是 $D_\QCoh(X)$ 的一个生成元
（Derived Categories, Definition \ref{derived-definition-generators}）。
\end{lemma}

\begin{proof}
设 $K \in D_\QCoh(\mathcal{O}_X)$。假定对所有 $p \in \mathbf{Z}$ 都有
$\Hom(E, K[p]) = 0$。我们须证明 $K = 0$。
由 Derived Categories, Lemma
\ref{derived-lemma-right-orthogonal}
可知，对所有 $E' \in \langle E \rangle$ 和 $p \in \mathbf{Z}$，
$\Hom(E', K[p])$ 均为零。将 Lemma \ref{perfect-lemma-construct-the-next-one}
用于 $a = -n - 1$，可知
$\mathcal{O}_X(-n - 1) \in \langle E \rangle$，因为它拟同构于一个有限复形，
其各项都是 $E$ 的直和项的有限直和。对 $a = -n - 2$ 重复这一论证，得到
$\mathcal{O}_X(-n - 2) \in \langle E \rangle$。归纳可得，对所有
$m \geq 0$ 都有 $\mathcal{O}_X(-m) \in \langle E \rangle$。由于
$$
\Hom(\mathcal{O}_X(-m), K[p]) =
H^p(X, K \otimes_{\mathcal{O}_X}^\mathbf{L} \mathcal{O}_X(m)) =
H^p(X, K \otimes_{\mathcal{O}_X}^\mathbf{L} \mathcal{O}_X(1)^{\otimes m})
$$
由 Lemma \ref{perfect-lemma-nonzero-some-cohomology} 可知 $K = 0$。（这里还用到了
$\mathcal{O}_X(1)$ 是 $X$ 上的充足可逆层；这由
Properties, Lemma \ref{properties-lemma-open-in-proj-ample} 推出。）
\end{proof}

\begin{remark}
\label{perfect-remark-pullback-generator}
设 $f : X \to Y$ 为拟紧拟分离概形之间的态射，并设
$E \in D_\QCoh(\mathcal{O}_Y)$ 为一个生成元
（见 Theorem \ref{perfect-theorem-bondal-van-den-Bergh}）。则下列条件等价：
\begin{enumerate}
\item 对 $K \in D_\QCoh(\mathcal{O}_X)$，当且仅当 $K = 0$ 时
$Rf_*K = 0$；
\item $Rf_* : D_\QCoh(\mathcal{O}_X) \to D_\QCoh(\mathcal{O}_Y)$
反映同构；
\item $Lf^*E$ 是 $D_\QCoh(\mathcal{O}_X)$ 的一个生成元。
\end{enumerate}
(1) 与 (2) 的等价性是如下事实的形式推论：
$Rf_* : D_\QCoh(\mathcal{O}_X) \to D_\QCoh(\mathcal{O}_Y)$ 是三角范畴的正合函子。
类似地，(1) 与 (3) 的等价性由 $Lf^*$ 是 $Rf_*$ 的左伴随这一事实形式地推出。
当 $f$ 为仿射态射（Lemma \ref{perfect-lemma-affine-morphism}），或 $f$ 为开浸入，
或 $f$ 为这类态射的复合时，
这些条件成立。因此：
\begin{enumerate}
\item 若 $X$ 为拟仿射概形，则 $\mathcal{O}_X$ 是
$D_\QCoh(\mathcal{O}_X)$ 的一个生成元；
\item 若 $X \subset \mathbf{P}^n_A$ 为拟紧局部闭子概形，则由
Lemma \ref{perfect-lemma-generator-P1}，
$\mathcal{O}_X \oplus \mathcal{O}_X(-1) \oplus \ldots \oplus \mathcal{O}_X(-n)$
是 $D_\QCoh(\mathcal{O}_X)$ 的一个生成元。
\end{enumerate}
\end{remark}




\section{紧对象与完美对象}
\label{perfect-section-compact}

\noindent
设 $X$ 为有限维 Noether 概形。由
Cohomology, Proposition \ref{cohomology-proposition-vanishing-Noetherian}
及
Cohomology on Sites, Lemma \ref{sites-cohomology-lemma-when-jshriek-compact}
可知，对所有开集 $U \subset X$，模层 $j_!\mathcal{O}_U$ 都是
$D(\mathcal{O}_X)$ 的紧对象。这些层通常不是拟凝聚的，因而并不给出导出范畴
$D(\mathcal{O}_X)$ 的完美对象。不过，若限制于具有拟凝聚上同调层的复形，
这种情况便不会发生。精确陈述如下。

\begin{proposition}
\label{perfect-proposition-compact-is-perfect}
设 $X$ 为拟紧拟分离概形。$D_\QCoh(\mathcal{O}_X)$ 的一个对象为紧对象，
当且仅当它是完美的。
\end{proposition}

\begin{proof}
若 $K$ 是 $D(\mathcal{O}_X)$ 的完美对象，其对偶为
$K^\vee$（Cohomology, Lemma \ref{cohomology-lemma-dual-perfect-complex}），则
$$
\Hom_{D(\mathcal{O}_X)}(K, M) =
H^0(X, K^\vee \otimes_{\mathcal{O}_X}^\mathbf{L} M)
$$
对 $M$ 函子性地成立。由于 $K^\vee \otimes_{\mathcal{O}_X}^\mathbf{L} -$
与直和可交换，且由 Lemma \ref{perfect-lemma-quasi-coherence-pushforward-direct-sums}，
$H^0(X, -)$ 在 $D_\QCoh(\mathcal{O}_X)$ 上与直和可交换，故 $K$ 是
$D_\QCoh(\mathcal{O}_X)$ 中的紧对象。

\medskip\noindent
反之，设 $K$ 为 $D_\QCoh(\mathcal{O}_X)$ 的紧对象。为证明 $K$ 完美，
只需证明对每个仿射开集 $U \subset X$，$K|_U$ 都完美，见
Cohomology, Lemma \ref{cohomology-lemma-perfect-independent-representative}。
注意 $j : U \to X$ 是拟紧分离态射。因此
$Rj_* : D_\QCoh(\mathcal{O}_U) \to D_\QCoh(\mathcal{O}_X)$
与直和可交换，见 Lemma \ref{perfect-lemma-quasi-coherence-pushforward-direct-sums}。
所以限制到 $U$ 的函子与 $Rj_*$ 的伴随性蕴含 $K|_U$ 是
$D_\QCoh(\mathcal{O}_U)$ 的紧对象。于是归结为 $X$ 仿射的情形。

\medskip\noindent
设 $X = \Spec(A)$ 为仿射概形。由 Lemma \ref{perfect-lemma-affine-compare-bounded}，
问题转化为 $D(A)$ 中的同一问题。对 $D(A)$，结论即
More on Algebra, Proposition \ref{more-algebra-proposition-perfect-is-compact}。
\end{proof}

\begin{remark}
\label{perfect-remark-classical-generator}
设 $X$ 为拟紧拟分离概形，且设 $G$ 为 $D(\mathcal{O}_X)$ 的一个完美对象，
并为 $D_\QCoh(\mathcal{O}_X)$ 的生成元。由
Theorem \ref{perfect-theorem-bondal-van-den-Bergh}，这样的对象至少存在一个。
结合 Lemma \ref{perfect-lemma-quasi-coherence-direct-sums}、
Proposition \ref{perfect-proposition-compact-is-perfect} 与
Derived Categories, Proposition
\ref{derived-proposition-generator-versus-classical-generator}
可知 $G$ 是 $D_{perf}(\mathcal{O}_X)$ 的经典生成元。
\end{remark}

\noindent
下述结果加强了 Proposition \ref{perfect-proposition-compact-is-perfect}。
设 $T \subset X$ 为概形 $X$ 的闭子集。仍以 $D_T(\mathcal{O}_X)$ 表示
$D(\mathcal{O}_X)$ 中由支集于 $T$ 上的对象组成的严格满、饱和三角子范畴
（Definition \ref{perfect-definition-supported-on}）。由于取直和与取上同调层可交换，
$D_T(\mathcal{O}_X)$ 具有直和，且这些直和等于 $D(\mathcal{O}_X)$ 中的直和。

\begin{lemma}
\label{perfect-lemma-compact-is-perfect-with-support}
设 $X$ 为拟紧拟分离概形，$T \subset X$ 为闭子集，且 $X \setminus T$
拟紧。$D_{\QCoh, T}(\mathcal{O}_X)$ 的一个对象为紧对象，当且仅当它作为
$D(\mathcal{O}_X)$ 的对象是完美的。
\end{lemma}

\begin{proof}
由引理前的评注，$D_{\QCoh, T}(\mathcal{O}_X)$ 是具有直和的三角范畴。
由 Proposition \ref{perfect-proposition-compact-is-perfect}，完美对象给出
$D(\mathcal{O}_X)$ 的紧对象，因而更是任一在取直和下保持的子范畴中的紧对象。
反之，我们将用到存在一个生成元 $E \in D_{\QCoh, T}(\mathcal{O}_X)$，
它是 $\mathcal{O}_X$-模的完美复形，见
Lemma \ref{perfect-lemma-generator-with-support}。因此由上所述，$E$ 是紧对象。再由
Derived Categories, Proposition
\ref{derived-proposition-generator-versus-classical-generator}
可知，$E$ 是 $D_{\QCoh, T}(\mathcal{O}_X)$ 的紧对象所成满子范畴的经典生成元。
所以任一紧对象都可由 $E$ 经过有限次如下运算构造：
(a) 取平移，(b) 取有限直和，(c) 取锥，以及 (d) 取直和项。
每种运算都保持完美性，结论随即成立。
\end{proof}

\begin{remark}
\label{perfect-remark-classical-generator-with-support}
设 $X$ 为拟紧拟分离概形，$T \subset X$ 为闭子集，且 $X \setminus T$ 拟紧。
设 $G$ 为 $D_{\QCoh, T}(\mathcal{O}_X)$ 的完美对象，同时也是
$D_{\QCoh, T}(\mathcal{O}_X)$ 的生成元。由
Lemma \ref{perfect-lemma-generator-with-support}，这样的对象至少存在一个。
结合 $D_{\QCoh, T}(\mathcal{O}_X)$ 具有直和这一事实、
Lemma \ref{perfect-lemma-compact-is-perfect-with-support} 以及
Derived Categories, Proposition
\ref{derived-proposition-generator-versus-classical-generator}
可知 $G$ 是 $D_{perf, T}(\mathcal{O}_X)$ 的经典生成元。
\end{remark}

\noindent
下述引理应用了上一个引理证明中的思想。

\begin{lemma}
\label{perfect-lemma-map-from-pseudo-coherent-to-complex-with-support}
设 $X$ 为拟紧拟分离概形，$T \subset X$ 为闭子集，使得
$U = X \setminus T$ 拟紧。设 $\alpha : P \to E$ 为
$D_\QCoh(\mathcal{O}_X)$ 的一个态射，且满足下列条件之一：
\begin{enumerate}
\item $P$ 完美且 $E$ 支集于 $T$ 上；或
\item $P$ 伪凝聚，$E$ 支集于 $T$ 上，且 $E$ 下有界。
\end{enumerate}
则存在一个 $\mathcal{O}_X$-模的完美复形 $I$ 及一个映射
$I \to \mathcal{O}_X[0]$，使得 $I \otimes^\mathbf{L} P \to E$ 为零，
且 $I|_U \to \mathcal{O}_U[0]$ 为同构。
\end{lemma}

\begin{proof}
令 $\mathcal{D} = D_{\QCoh, T}(\mathcal{O}_X)$。在两种情形中，复形
$K = R\SheafHom(P, E)$ 都是 $\mathcal{D}$ 的对象。关于拟凝聚性，见
Lemma \ref{perfect-lemma-quasi-coherence-internal-hom}。由于 $R\SheafHom$ 的构造与
限制到开集可交换，显然 $K$ 支集于 $T$ 上。映射 $\alpha$ 定义了
$H^0(K) = \Hom_{D(\mathcal{O}_X)}(\mathcal{O}_X[0], K)$ 的一个元素。
因此只需对映射 $\alpha : \mathcal{O}_X[0] \to K$ 证明结论。

\medskip\noindent
取完美生成元 $E \in \mathcal{D}$，见
Lemma \ref{perfect-lemma-generator-with-support}。按
Derived Categories, Lemma \ref{derived-lemma-write-as-colimit}，利用生成元 $E$ 写成
$$
K = \text{hocolim} K_n
$$
。由于函子 $\mathcal{D} \to D(\mathcal{O}_X)$ 与直和可交换，
$K = \text{hocolim} K_n$ 在 $D(\mathcal{O}_X)$ 中也成立。由于
$\mathcal{O}_X$ 是 $D(\mathcal{O}_X)$ 的紧对象，存在 $n$ 及一个诱导出
$\alpha$ 的态射 $\alpha_n : \mathcal{O}_X \to K_n$，见
Derived Categories, Lemma \ref{derived-lemma-commutes-with-countable-sums}。
将 Derived Categories, Lemma \ref{derived-lemma-factor-through}
用于环境范畴 $D(\mathcal{O}_X)$ 中的态射 $\mathcal{O}_X[0] \to K_n$，
可知 $\alpha_n$ 分解为 $\mathcal{O}_X[0] \to Q \to K_n$，其中 $Q$ 是
$\langle E \rangle$ 的对象。因此 $Q$ 是支集于 $T$ 上的完美复形。

\medskip\noindent
选取可区别三角
$$
I \to \mathcal{O}_X[0] \to Q \to I[1]
$$
按构造，$I$ 完美，映射 $I \to \mathcal{O}_X[0]$ 限制到 $U$ 后为同构；
又因 $\alpha$ 经 $Q$ 分解，复合 $I \to K$ 为零。这就证明了引理。
\end{proof}






\section{作为模范畴的导出范畴}
\label{perfect-section-derived-is-dga}

\noindent
本节将从前面的结果推出若干结论。在此之前还需要几个引理。

\begin{lemma}
\label{perfect-lemma-tensor-with-QCoh-complex}
设 $X$ 为概形。设 $K^\bullet$ 为一个 $\mathcal{O}_X$-模复形，
其上同调层拟凝聚。令
$(E, d) = \Hom_{\text{Comp}^{dg}(\mathcal{O}_X)}(K^\bullet, K^\bullet)$
为自同态微分分次代数。则
$$
- \otimes_E^\mathbf{L} K^\bullet :
D(E, \text{d}) \longrightarrow D(\mathcal{O}_X)
$$
这一由
Differential Graded Algebra, Lemma
\ref{dga-lemma-tensor-with-complex-derived}
给出的函子之像包含于 $D_\QCoh(\mathcal{O}_X)$。
\end{lemma}

\begin{proof}
设 $P$ 为具有性质 (P) 的微分分次 $E$-模，并设 $F_\bullet$ 为
Differential Graded Algebra, Section \ref{dga-section-P-resolutions}
中所述的 $P$ 上滤过。则
$$
P \otimes_E K^\bullet = \text{hocolim}\ F_iP \otimes_E K^\bullet
$$
每个 $F_iP$ 都有一个有限滤过，其分次片是 $E[k]$ 的直和。结论容易推出。
\end{proof}

\noindent
下述结果取自 \cite{BvdB}。

\begin{theorem}
\label{perfect-theorem-DQCoh-is-Ddga}
设 $X$ 为拟紧拟分离概形。则存在一个微分分次代数 $(E, \text{d})$，
其非零上同调群 $H^i(E)$ 仅有有限多个，并且
$D_\QCoh(\mathcal{O}_X)$ 等价于 $D(E, \text{d})$。
\end{theorem}

\begin{proof}
设 $K^\bullet$ 为一个 $\mathcal{O}$-模的 K-内射复形，它是完美的并生成
$D_\QCoh(\mathcal{O}_X)$。由 Theorem \ref{perfect-theorem-bondal-van-den-Bergh}
及 K-内射解的存在性，这样的复形确实存在。我们将证明，取
$$
(E, \text{d}) = \Hom_{\text{Comp}^{dg}(\mathcal{O}_X)}(K^\bullet, K^\bullet)
$$
时定理成立；其中 $\text{Comp}^{dg}(\mathcal{O}_X)$ 是 $\mathcal{O}$-模复形的
微分分次范畴。参见
Differential Graded Algebra, Section \ref{dga-section-variant-base-change}。
由于 $K^\bullet$ 是 K-内射的，对所有 $n \in \mathbf{Z}$ 都有
\begin{equation}
\label{perfect-equation-E-is-OK}
H^n(E) = \Ext^n_{D(\mathcal{O}_X)}(K^\bullet, K^\bullet)
\end{equation}
。由 Lemma \ref{perfect-lemma-ext-from-perfect-into-bounded-QCoh}，这些 Ext 中只有
有限多个非零。考虑
$$
- \otimes_E^\mathbf{L} K^\bullet :
D(E, \text{d}) \longrightarrow D(\mathcal{O}_X)
$$
这一由
Differential Graded Algebra, Lemma
\ref{dga-lemma-tensor-with-complex-derived} 给出的函子。由于 $K^\bullet$ 完美，
它定义了 $D(\mathcal{O}_X)$ 的一个紧对象，见
Proposition \ref{perfect-proposition-compact-is-perfect}。结合
(\ref{perfect-equation-E-is-OK})，由
Differential Graded Algebra, Lemmas
\ref{dga-lemma-fully-faithful-in-compact-case} 可知上述函子全忠实。由
Differential Graded Algebra, Lemmas
\ref{dga-lemma-tensor-with-complex-hom-adjoint}，它有右伴随
$$
R\Hom(K^\bullet, - ) : D(\mathcal{O}_X) \longrightarrow D(E, \text{d})
$$
；根据伴随函子的一般理论，此函子是左拟逆。另一方面，由
Lemma \ref{perfect-lemma-tensor-with-QCoh-complex} 得到
$$
- \otimes_E^\mathbf{L} K^\bullet :
D(E, \text{d}) \longrightarrow D_\QCoh(\mathcal{O}_X)
$$
，而由于选取 $K^\bullet$ 为 $D_\QCoh(\mathcal{O}_X)$ 的生成元，
该伴随函子限制到 $D_\QCoh(\mathcal{O}_X)$ 后的核为零。
形式论证即给出所需等价，见
Derived Categories, Lemma
\ref{derived-lemma-fully-faithful-adjoint-kernel-zero}.
\end{proof}

\begin{remark}[带支集的变体]
\label{perfect-remark-DQCoh-is-Ddga-with-support}
设 $X$ 为拟紧拟分离概形。设 $T \subset X$ 为闭子集，且 $X \setminus T$
拟紧。Theorem \ref{perfect-theorem-DQCoh-is-Ddga} 对
$D_{\QCoh, T}(\mathcal{O}_X)$ 的类似结论也成立。这由定理证明中完全相同的论证推出，
其中使用 Lemmas \ref{perfect-lemma-generator-with-support}、
\ref{perfect-lemma-compact-is-perfect-with-support} 以及
Lemma \ref{perfect-lemma-tensor-with-QCoh-complex} 的带支集变体。
若以后需要这一结果，我们将在此给出精确陈述和详细证明。
\end{remark}

\begin{remark}[dga 的唯一性]
\label{perfect-remark-independence-choice}
设 $X$ 为环 $R$ 上的拟紧拟分离概形。由
Theorem \ref{perfect-theorem-DQCoh-is-Ddga} 证明中的构造，存在 $R$ 上的微分分次代数
$(A, \text{d})$，使得三角范畴 $D_\QCoh(X)$ 与 $D(A, \text{d})$
$R$-线性等价。可以问：$(A, \text{d})$ 有多唯一？答案（仅仅）略强于说
$(A, \text{d})$ 在导出等价意义下良定义。具体地，设 $(B, \text{d})$
是第二个这样的对。于是
$$
(A, \text{d}) = \Hom_{\text{Comp}^{dg}(\mathcal{O}_X)}(K^\bullet, K^\bullet)
$$
且
$$
(B, \text{d}) = \Hom_{\text{Comp}^{dg}(\mathcal{O}_X)}(L^\bullet, L^\bullet)
$$
，其中 $K^\bullet$ 与 $L^\bullet$ 是某些 $\mathcal{O}_X$-模的 K-内射复形，
分别对应于 $D_\QCoh(\mathcal{O}_X)$ 的完美生成元。令
$$
\Omega = \Hom_{\text{Comp}^{dg}(\mathcal{O}_X)}(K^\bullet, L^\bullet)
\quad
\Omega' = \Hom_{\text{Comp}^{dg}(\mathcal{O}_X)}(L^\bullet, K^\bullet)
$$
则 $\Omega$ 是微分分次 $B^{opp} \otimes_R A$-模，$\Omega'$ 是微分分次
$A^{opp} \otimes_R B$-模。此外，等价
$$
D(A, \text{d}) \to D_\QCoh(\mathcal{O}_X) \to
D(B, \text{d})
$$
由函子 $- \otimes_A^\mathbf{L} \Omega'$ 给出，拟逆也类似。因此正处于
Differential Graded Algebra, Remark \ref{dga-remark-hochschild-cohomology}
的情形。若以后需要此评注，我们将在这里给出精确陈述和详细证明。
\end{remark}






\section{伪凝聚复形的刻画，I}
\label{perfect-section-pseudo-coherent-hocolim}

\noindent
可以用上述方法把伪凝聚对象刻画为完美对象所给逼近的导出同伦极限。

\begin{lemma}
\label{perfect-lemma-pseudo-coherent-hocolim}
设 $X$ 为拟紧拟分离概形，且设 $K \in D(\mathcal{O}_X)$。下列条件等价：
\begin{enumerate}
\item $K$ 是伪凝聚的；
\item $K = \text{hocolim} K_n$，其中 $K_n$ 完美，且对所有 $n$，
$\tau_{\geq -n}K_n \to \tau_{\geq -n}K$ 均为同构。
\end{enumerate}
\end{lemma}

\begin{proof}
在任意环化空间上，(2) $\Rightarrow$ (1) 都成立。事实上，假定 (2) 成立。
回忆导出范畴的完美对象是伪凝聚的，见
Cohomology, Lemma \ref{cohomology-lemma-perfect}。于是由定义，
$\tau_{\geq -n}K_n$ 是 $(-n + 1)$-伪凝聚的，故
$\tau_{\geq -n}K$ 是 $(-n + 1)$-伪凝聚的，进而 $K$ 是
$(-n + 1)$-伪凝聚的。这对所有 $n$ 都成立，所以 $K$ 伪凝聚，见
Cohomology, Definition \ref{cohomology-definition-pseudo-coherent}。

\medskip\noindent
假定 (1)。首先选取由完美复形 $K_1$ 给出的 $(X, K, -2)$ 的一个逼近
$K_1 \to K$，见
Definitions \ref{perfect-definition-approximation-holds} and
\ref{perfect-definition-approximation} and
Theorem \ref{perfect-theorem-approximation}。归纳假设已经有
$$
K_1 \to K_2 \to \ldots \to K_n \to K
$$
，其中 $K_i$ 完美，且对所有 $1 \leq i \leq n$，
$\tau_{\geq -i}K_i \to \tau_{\geq -i}K$ 都是同构。对完美对象 $K_n$，
按 Lemma \ref{perfect-lemma-ext-from-perfect-into-bounded-QCoh} 选取 $a \leq b$。
选取 $(X, K, \min(a - 1, -n - 1))$ 的一个逼近 $K_{n + 1} \to K$。
选取可区别三角
$$
K_{n + 1} \to K \to C \to K_{n + 1}[1]
$$
。于是 $C \in D_\QCoh(\mathcal{O}_X)$ 且当 $i \geq a$ 时
$H^i(C) = 0$。故由 $a, b$ 的选取，
$\Hom_{D(\mathcal{O}_X)}(K_n, C) = 0$。所以复合
$K_n \to K \to C$ 为零。由
Derived Categories, Lemma \ref{derived-lemma-representable-homological}，
$K_n \to K$ 可经 $K_{n + 1}$ 分解，从而完成归纳步骤。

\medskip\noindent
还须证明 $K = \text{hocolim} K_n$。将
Derived Categories, Lemma \ref{derived-lemma-cohomology-of-hocolim}
用于函子
$H^i( - ) : D(\mathcal{O}_X) \to \textit{Mod}(\mathcal{O}_X)$
和我们所选的 $K_n$，即得此结论。
\end{proof}

\begin{lemma}
\label{perfect-lemma-pseudo-coherent-hocolim-with-support}
设 $X$ 为拟紧拟分离概形。设 $T \subset X$ 为闭子集，且 $X \setminus T$
拟紧。设 $K \in D(\mathcal{O}_X)$ 支集于 $T$ 上。下列条件等价：
\begin{enumerate}
\item $K$ 是伪凝聚的；
\item $K = \text{hocolim} K_n$，其中 $K_n$ 完美且支集于 $T$ 上，并且对所有
$n$，$\tau_{\geq -n}K_n \to \tau_{\geq -n}K$ 均为同构。
\end{enumerate}
\end{lemma}

\begin{proof}
本引理的证明与 Lemma \ref{perfect-lemma-pseudo-coherent-hocolim} 的证明完全相同，
唯一差别是在选取逼近时使用三元组 $(T, K, m)$。
\end{proof}









\section{一个等价的例子}
\label{perfect-section-example-equivalence}

\noindent
在 Section \ref{perfect-section-example-generator} 中，我们证明了环 $A$ 上射影空间
$\mathbf{P}^n_A$ 的导出范畴由一个向量丛生成；事实上，该向量丛是结构层若干平移的直和。
本节证明，这确定了 $D_\QCoh(\mathcal{O}_{\mathbf{P}^n_A})$ 与某个
$A$-代数的导出范畴之间的等价。

\medskip\noindent
在陈述结果之前先引入一些记号。设 $A$ 为环，
$X = \mathbf{P}^n_A = \text{Proj}(S)$，其中 $S = A[X_0, \ldots, X_n]$。
由 Lemma \ref{perfect-lemma-generator-P1} 可知
\begin{equation}
\label{perfect-equation-generator-Pn}
P =
\mathcal{O}_X \oplus \mathcal{O}_X(-1) \oplus \ldots \oplus \mathcal{O}_X(-n)
\end{equation}
是 $D_\QCoh(\mathcal{O}_X)$ 的完美生成元。考虑（非交换）$A$-代数
\begin{equation}
\label{perfect-equation-algebra-for-Pn}
R = \Hom_{\mathcal{O}_X}(P, P) =
\left(
\begin{matrix}
S_0 & S_1 & S_2 & \ldots & \ldots \\
0 & S_0 & S_1 & \ldots & \ldots\\
0 & 0 & S_0 & \ldots & \ldots \\
\ldots & \ldots & \ldots & \ldots & \ldots \\
0 & \ldots & \ldots & \ldots & S_0
\end{matrix}
\right)
\end{equation}
，其乘法和加法是显然的。若按通常方式将 $P$ 视为 $\mathcal{O}_X$-模复形
（即 $P$ 位于次数 $0$，其他各次数均为零），则
$$
R = \Hom_{\text{Comp}^{dg}(\mathcal{O}_X)}(P, P)
$$
，其中右端把 $R$ 视为 $A$ 上微分为零的微分分次代数
（即 $R$ 位于次数 $0$，其他各次数均为零）。根据
Differential Graded Algebra, Section \ref{dga-section-variant-base-change}
中的讨论，得到导出函子
$$
- \otimes_R^\mathbf{L} P : D(R) \longrightarrow D(\mathcal{O}_X),
$$
特别见 Differential Graded Algebra, Lemma
\ref{dga-lemma-tensor-with-complex-derived}。由
Lemma \ref{perfect-lemma-tensor-with-QCoh-complex}，此函子的本质像包含于
$D_\QCoh(\mathcal{O}_X)$。

\begin{lemma}
\label{perfect-lemma-Pn-module-category}
\begin{reference}
\cite{Beilinson}
\end{reference}
设 $A$ 为环，$X = \mathbf{P}^n_A = \text{Proj}(S)$，其中
$S = A[X_0, \ldots, X_n]$。取 (\ref{perfect-equation-generator-Pn}) 中的
$P$ 和 (\ref{perfect-equation-algebra-for-Pn}) 中的 $R$，则函子
$$
- \otimes_R^\mathbf{L} P : D(R) \longrightarrow D_\QCoh(\mathcal{O}_X)
$$
是三角范畴的一个 $A$-线性等价，并把 $R$ 送到 $P$。
\end{lemma}

\noindent
换言之，射影空间上拟凝聚模的导出范畴等价于某个（非交换）代数上模的导出范畴。
在 $A$ 上所有射影概形中，射影空间的这一性质似乎相当特殊。

\begin{proof}
为证明该函子全忠实，只需证明当 $i \not = 0$ 时 $\Ext^i_X(P, P)$ 为零，
当 $i = 0$ 时它等于 $R$，见
Differential Graded Algebra, Lemma
\ref{dga-lemma-fully-faithful-in-compact-case}。与
Lemma \ref{perfect-lemma-ext-from-perfect-into-bounded-QCoh} 的证明相同，有
$$
\Ext^i_X(P, P) = H^i(X, P^\wedge \otimes P) =
\bigoplus\nolimits_{0 \leq a, b \leq n} H^i(X, \mathcal{O}_X(a - b))
$$
。由射影空间上同调的计算（Cohomology of Schemes, Lemma
\ref{coherent-lemma-cohomology-projective-space-over-ring}），可知这些
$\Ext$-群除非 $i = 0$ 否则均为零。当 $i = 0$ 时得到 $R$，因为
(\ref{perfect-equation-algebra-for-Pn}) 正是如此定义 $R$ 的。由
Differential Graded Algebra, Lemma
\ref{dga-lemma-tensor-with-complex-hom-adjoint}
，该函子有右伴随 $R\Hom(P, -) : D_\QCoh(\mathcal{O}_X) \to D(R)$。
由 Lemma \ref{perfect-lemma-generator-P1}，$P$ 是 $D_\QCoh(\mathcal{O}_X)$ 的生成元，
故 $R\Hom(P, -)$ 的核为零。因此由 Derived Categories, Lemma
\ref{derived-lemma-fully-faithful-adjoint-kernel-zero}，该函子是三角范畴的等价。
\end{proof}






\section{再论拟凝聚化子}
\label{perfect-section-better-coherator}

\noindent
在 Section \ref{perfect-section-coherator} 中，我们构造并研究了典范函子
$D(\QCoh(\mathcal{O}_X)) \to D(\mathcal{O}_X)$ 的右伴随 $RQ_X$。
它被构造为拟凝聚化子
$Q_X : \textit{Mod}(\mathcal{O}_X) \to \QCoh(\mathcal{O}_X)$
的右导出扩张。本节研究包含函子
$$
D_\QCoh(\mathcal{O}_X) \longrightarrow D(\mathcal{O}_X)
$$
何时具有右伴随。若此右伴随存在，我们将把它记作
\footnote{这可能不是标准记号。不过，我们已经用 $Q_X$ 表示拟凝聚化子，
用 $RQ_X$ 表示其导出扩张。}
$$
DQ_X :
D(\mathcal{O}_X) \longrightarrow D_\QCoh(\mathcal{O}_X)
$$
。事实表明，拟紧拟分离概形具有这样的右伴随。

\begin{lemma}
\label{perfect-lemma-better-coherator}
设 $X$ 为拟紧拟分离概形。包含函子
$D_\QCoh(\mathcal{O}_X) \to D(\mathcal{O}_X)$ 具有右伴随 $DQ_X$。
\end{lemma}

\begin{proof}[第一种证明]
我们使用
Cohomology of Schemes, Lemma \ref{coherent-lemma-induction-principle}
中的归纳原理来证明。若
$D(\QCoh(\mathcal{O}_X)) \to D_\QCoh(\mathcal{O}_X)$ 是等价，则本引理成立，
因为 Section \ref{perfect-section-coherator} 的函子 $RQ_X$ 是函子
$D(\QCoh(\mathcal{O}_X)) \to D(\mathcal{O}_X)$ 的右伴随。特别地，本引理对仿射概形成立，
见 Lemma \ref{perfect-lemma-affine-coherator}。因此只需证明：若
$X = U \cup V$ 是两个拟紧开集之并，且引理对 $U$、$V$ 和 $U \cap V$
成立，则它对 $X$ 成立。

\medskip\noindent
该伴随存在，当且仅当对 $D(\mathcal{O}_X)$ 的每个对象 $K$，都能找到可区别三角
$$
E' \to E \to K \to E'[1]
$$
位于 $D(\mathcal{O}_X)$ 中，使得 $E'$ 属于 $D_\QCoh(\mathcal{O}_X)$，
且对 $D_\QCoh(\mathcal{O}_X)$ 中所有 $M$ 都有 $\Hom(M, K) = 0$。见
Derived Categories, Lemma \ref{derived-lemma-right-adjoint}。考虑
Cohomology, Lemma \ref{cohomology-lemma-exact-sequence-j-star}
在 $D(\mathcal{O}_X)$ 中给出的可区别三角
$$
E \to Rj_{U, *}E|_U \oplus Rj_{V, *}E|_V \to
Rj_{U \cap V, *}E|_{U \cap V} \to E[1]
$$
。由 Derived Categories, Lemma \ref{derived-lemma-prepare-adjoint}，
只需对 $Rj_{U, *}E|_U$、$Rj_{V, *}E|_V$ 和
$Rj_{U \cap V, *}E|_{U \cap V}$ 构造所需的可区别三角。
这把问题化为下一段所讨论的陈述。

\medskip\noindent
设 $j : U \to X$ 是开浸入，其中 $U$ 是本引理成立的拟紧开集。
设 $L$ 为 $D(\mathcal{O}_U)$ 的对象。则存在 $D(\mathcal{O}_X)$ 中的可区别三角
$$
E' \to Rj_*L \to K \to E'[1]
$$
，使得 $E'$ 属于 $D_\QCoh(\mathcal{O}_X)$，且对
$D_\QCoh(\mathcal{O}_X)$ 中所有 $M$ 都有 $\Hom(M, K) = 0$。
为此，选取 $D(\mathcal{O}_U)$ 中的可区别三角
$$
L' \to L \to Q \to L'[1]
$$
，使得 $L'$ 属于 $D_\QCoh(\mathcal{O}_U)$，且对
$D_\QCoh(\mathcal{O}_U)$ 中所有 $N$ 都有 $\Hom(N, Q) = 0$。
这是可能的，因为 Derived Categories, Lemma \ref{derived-lemma-right-adjoint}
中的陈述是充要条件。我们得到 $D(\mathcal{O}_X)$ 中的可区别三角
$$
Rj_*L' \to Rj_*L \to Rj_*Q \to Rj_*L'[1]
$$
。由 Lemma \ref{perfect-lemma-quasi-coherence-direct-image}，
$Rj_*L'$ 属于 $D_\QCoh(\mathcal{O}_X)$。另一方面，若 $M$ 属于
$D_\QCoh(\mathcal{O}_X)$，则
$$
\Hom(M, Rj_*Q) = \Hom(Lj^*M, Q) = 0
$$
，因为由 Lemma \ref{perfect-lemma-quasi-coherence-pullback}，
$Lj^*M$ 属于 $D_\QCoh(\mathcal{O}_U)$。证明完毕。
\end{proof}

\begin{proof}[第二种证明]
由 Derived Categories, Proposition \ref{derived-proposition-brown}，
该伴随存在。其假设均满足：首先，$D_\QCoh(\mathcal{O}_X)$ 具有直和，
且直和与包含函子可交换（Lemma \ref{perfect-lemma-quasi-coherence-direct-sums}）。
另一方面，$D_\QCoh(\mathcal{O}_X)$ 是紧生成的，因为由
Theorem \ref{perfect-theorem-bondal-van-den-Bergh} 它具有完美生成元，
而由 Proposition \ref{perfect-proposition-compact-is-perfect}，完美对象是紧的。
\end{proof}

\begin{lemma}
\label{perfect-lemma-pushforward-better-coherator}
设 $f : X \to Y$ 为概形的拟紧拟分离态射。若对 $X$ 和 $Y$，包含函子
$D_\QCoh \to D$ 的右伴随 $DQ_X$ 与 $DQ_Y$ 存在，则
$$
Rf_* \circ DQ_X = DQ_Y \circ Rf_*
$$
\end{lemma}

\begin{proof}
由 Lemma \ref{perfect-lemma-quasi-coherence-direct-image}，$Rf_*$ 把
$D_\QCoh(\mathcal{O}_X)$ 送入 $D_\QCoh(\mathcal{O}_Y)$，故该陈述有意义。
它之所以成立，是因为 $Lf^*$ 同样把 $D_\QCoh(\mathcal{O}_Y)$ 映入
$D_\QCoh(\mathcal{O}_X)$（Lemma \ref{perfect-lemma-quasi-coherence-pullback}），
因而 $Rf_* \circ DQ_X$ 与 $DQ_Y \circ Rf_*$ 都是
$Lf^* : D_\QCoh(\mathcal{O}_Y) \to D(\mathcal{O}_X)$ 的右伴随。
\end{proof}

\begin{remark}
\label{perfect-remark-explain-consequence}
设 $X$ 为拟紧拟分离概形，并设 $X = U \cup V$，其中 $U$ 与 $V$ 为拟紧开集。
由 Lemma \ref{perfect-lemma-better-coherator}，函子
$DQ_X$、$DQ_U$、$DQ_V$、$DQ_{U \cap V}$ 均存在。此外，对任意
$K \in D(\mathcal{O}_X)$，存在典范可区别三角
$$
DQ_X(K) \to Rj_{U, *}DQ_U(K|_U) \oplus Rj_{V, *}DQ_V(K|_V)
\to Rj_{U \cap V, *}DQ_{U \cap V}(K|_{U \cap V}) \to
$$
。这是将正合函子 $DQ_X$ 作用于
Cohomology, Lemma \ref{cohomology-lemma-exact-sequence-j-star}
的可区别三角，并三次使用 Lemma \ref{perfect-lemma-pushforward-better-coherator}
所得的结果。
\end{remark}

\begin{lemma}
\label{perfect-lemma-boundedness-better-coherator}
设 $X$ 为拟紧拟分离概形。Lemma \ref{perfect-lemma-better-coherator} 的函子
$DQ_X$ 具有如下有界性：存在整数 $N = N(X)$，使得若
$K$ 属于 $D(\mathcal{O}_X)$，并且对 $X$ 中的仿射开集 $U$ 及
$i \not \in [a, b]$ 有 $H^i(U, K) = 0$，则当
$i \not \in [a, b + N]$ 时，上同调层 $H^i(DQ_X(K))$ 为零。
\end{lemma}

\begin{proof}
我们用 Cohomology of Schemes, Lemma
\ref{coherent-lemma-induction-principle} 的归纳原理证明。

\medskip\noindent
若 $X$ 仿射，则取 $N = 0$ 时引理成立，因为此时 $RQ_X = DQ_X$ 是取
$R\Gamma(X, K)$ 所对应的拟凝聚层复形。见
Lemmas \ref{perfect-lemma-affine-compare-bounded} and \ref{perfect-lemma-affine-coherator}。

\medskip\noindent
假设 $U, V$ 是 $X$ 中的拟紧开集，且引理对 $U$、$V$ 及 $U \cap V$ 成立，
相应整数分别记为 $N(U)$、$N(V)$ 与 $N(U \cap V)$。现设 $K$ 属于
$D(\mathcal{O}_X)$，且对所有仿射开集 $W \subset X$ 及所有
$i \not \in [a, b]$ 都有 $H^i(W, K) = 0$。则
$K|_U$、$K|_V$、$K|_{U \cap V}$ 具有同一性质。因此
$RQ_U(K|_U)$、$RQ_V(K|_V)$ 与 $RQ_{U \cap V}(K|_{U \cap V})$
在区间 $[a, b + \max(N(U), N(V), N(U \cap V))$ 外的上同调层均为零。
由 Lemma \ref{perfect-lemma-quasi-coherence-direct-image}，函子
$Rj_{U, *}$、$Rj_{V, *}$、$Rj_{U \cap V, *}$ 在 $D_\QCoh$ 上具有有限上同调维数，
故存在 $N$，使得 $Rj_{U, *}DQ_U(K|_U)$、$Rj_{V, *}DQ_V(K|_V)$ 和
$Rj_{U \cap V, *}DQ_{U \cap V}(K|_{U \cap V})$ 在区间
$[a, b + N]$ 外的上同调层均为零。最后由
Remark \ref{perfect-remark-explain-consequence} 的可区别三角得出结论。
\end{proof}

\begin{example}
\label{perfect-example-inverse-limit-quasi-coherent}
设 $X$ 为拟紧拟分离概形，并设 $(\mathcal{F}_n)$ 为拟凝聚层的逆系统。
由于 $DQ_X$ 是右伴随，它与积可交换，因而也与导出极限可交换。因此
$$
DQ_X(R\lim \mathcal{F}_n) =
(R\lim\text{ 于 }D_\QCoh(\mathcal{O}_X))(\mathcal{F}_n)
$$
，其中第一个 $R\lim$ 在 $D(\mathcal{O}_X)$ 中取。事实上，把它记作
$K = R\lim \mathcal{F}_n$。对任意仿射开集 $U \subset X$，有
$$
H^i(U, K) =
H^i(R\Gamma(U, R\lim \mathcal{F}_n)) =
H^i(R\lim R\Gamma(U, \mathcal{F}_n)) =
H^i(R\lim \Gamma(U, \mathcal{F}_n))
$$
，因为上同调与导出极限可交换，且拟凝聚层 $\mathcal{F}_n$ 在仿射概形上没有高阶上同调。
由阿贝尔群范畴中 $R\lim$ 的计算，除非 $i \in [0, 1]$，否则
$H^i(U, K) = 0$。最后由 Lemma \ref{perfect-lemma-boundedness-better-coherator}
可知，$D_\QCoh(\mathcal{O}_X)$ 中的 $R\lim$，即上面的 $DQ_X(K)$，属于
$D^b_\QCoh(\mathcal{O}_X)$。
\end{example}













\section{上同调与基变换，IV}
\label{perfect-section-cohomology-and-base-change-perfect}

\noindent
本节继续
Cohomology of Schemes, Section
\ref{coherent-section-cohomology-and-base-change-perfect} 中的讨论。
首先，对概形的拟紧拟分离态射以及上同调层拟凝聚的复形，
我们有一个非常一般的投影公式版本。

\begin{lemma}
\label{perfect-lemma-cohomology-base-change}
设 $f : X \to Y$ 为概形的拟紧拟分离态射。对
$D_\QCoh(\mathcal{O}_X)$ 中的 $E$ 和 $D_\QCoh(\mathcal{O}_Y)$ 中的 $K$，
在 Cohomology, Equation
(\ref{cohomology-equation-projection-formula-map}) 中定义的映射
$$
Rf_*(E) \otimes_{\mathcal{O}_Y}^\mathbf{L} K
\longrightarrow
Rf_*(E \otimes_{\mathcal{O}_X}^\mathbf{L} Lf^*K)
$$
是同构。
\end{lemma}

\begin{proof}
为检验该映射是同构，可以在 $Y$ 上局部地工作。因此归结为 $Y$ 仿射的情形。

\medskip\noindent
设 $K = \bigoplus K_i$ 是若干复形 $K_i \in D_\QCoh(\mathcal{O}_Y)$ 的直和。
若陈述对每个 $K_i$ 成立，则也对 $K$ 成立。事实上，函子 $Lf^*$ 与
$\otimes^\mathbf{L}$ 按构造保持直和，而由
Lemma \ref{perfect-lemma-quasi-coherence-pushforward-direct-sums}，$Rf_*$
（对上同调层拟凝聚的复形）与直和可交换。此外，设
$K \to L \to M \to K[1]$ 是 $D_\QCoh(Y)$ 中的可区别三角。
若引理的陈述对 $K, L, M$ 中两个成立，则对第三个也成立
（因为所涉及的函子是三角范畴的正合函子）。

\medskip\noindent
假定 $Y$ 仿射，记 $Y = \Spec(A)$。函子
$\widetilde{\ } : D(A) \to D_\QCoh(\mathcal{O}_Y)$ 是等价
（Lemma \ref{perfect-lemma-affine-compare-bounded}）。令 $T$ 表示 $K \in D(A)$
所满足的如下性质：引理的陈述对 $\widetilde{K}$ 成立。由上述讨论及
More on Algebra, Remark \ref{more-algebra-remark-P-resolution}，
只需证明 $T$ 对 $A[k]$ 成立。由于引理的陈述对结构层的平移显然成立，证明完成。
\end{proof}

\begin{definition}
\label{perfect-definition-tor-independent}
设 $S$ 为概形，$X$、$Y$ 为 $S$ 上概形。若对每个映到同一点 $s \in S$ 的
$x \in X$ 与 $y \in Y$，环 $\mathcal{O}_{X, x}$ 和 $\mathcal{O}_{Y, y}$
在 $\mathcal{O}_{S, s}$ 上 Tor 独立（见 More on Algebra, Definition
\ref{more-algebra-definition-tor-independent}），则称 $X$ 与 $Y$
{\it 在 $S$ 上 Tor 独立}。
\end{definition}

\begin{lemma}
\label{perfect-lemma-tor-independent}
设 $f : X \to S$ 与 $g : Y \to S$ 为概形态射。下列条件等价：
\begin{enumerate}
\item $X$ 与 $Y$ 在 $S$ 上 Tor 独立；
\item 对每组仿射开集 $U \subset X$、$V \subset Y$、$W \subset S$，若
$f(U) \subset W$ 且 $g(V) \subset W$，则环 $\mathcal{O}_X(U)$ 与
$\mathcal{O}_Y(V)$ 在 $\mathcal{O}_S(W)$ 上 Tor 独立；
\item 存在仿射开覆盖 $S = \bigcup W_i$，且对每个 $i$ 存在仿射开覆盖
$f^{-1}(W_i) = \bigcup U_{ij}$ 与 $g^{-1}(W_i) = \bigcup V_{ik}$，
使得对所有 $i, j, k$，环 $\mathcal{O}_X(U_{ij})$ 与
$\mathcal{O}_Y(V_{ik})$ 在 $\mathcal{O}_S(W_i)$ 上 Tor 独立。
\end{enumerate}
\end{lemma}

\begin{proof}
略。提示：使用 More on Algebra, Lemma
\ref{more-algebra-lemma-tor-independent}。
\end{proof}

\begin{lemma}
\label{perfect-lemma-flat-base-change-tor-independent}
设 $X \to S$ 与 $Y \to S$ 为概形态射。设 $S' \to S$ 为概形态射，并记
$X' = X \times_S S'$、$Y' = Y \times_S S'$。若 $X$ 与 $Y$
在 $S$ 上 Tor 独立，且 $S' \to S$ 平坦，则 $X'$ 与 $Y'$ 在 $S'$ 上 Tor 独立。
\end{lemma}

\begin{proof}
略。提示：使用 Lemma \ref{perfect-lemma-tor-independent}，并在仿射开集上使用
More on Algebra, Lemma
\ref{more-algebra-lemma-flat-base-change-tor-independent}。
\end{proof}

\begin{lemma}
\label{perfect-lemma-compare-base-change}
设 $g : S' \to S$ 为概形态射，且设 $f : X \to S$ 拟紧拟分离。
考虑基变换图
$$
\xymatrix{
X' \ar[r]_{g'} \ar[d]_{f'} &
X \ar[d]^f \\
S' \ar[r]^g &
S
}
$$
若 $X$ 与 $S'$ 在 $S$ 上 Tor 独立，则对所有
$E \in D_\QCoh(\mathcal{O}_X)$，典范箭头
$Lg^*Rf_*E \to Rf'_*L(g')^*E$ 均为同构。
\end{lemma}

\begin{proof}
对 $D(\mathcal{O}_X)$ 的任意对象 $E$，由
Cohomology, Remark \ref{cohomology-remark-base-change} 得到典范基变换映射
$Lg^*Rf_*E \to Rf'_*L(g')^*E$。为检验它是同构，可以在 $S'$ 上局部地工作。
因此可假设 $g : S' \to S$ 是仿射概形之间的态射。特别地，$g$ 仿射，
只需证明
$$
Rg_*Lg^*Rf_*E \to Rg_*Rf'_*L(g')^*E = Rf_*(Rg'_* L(g')^* E)
$$
是同构，见 Lemma \ref{perfect-lemma-affine-morphism}（并使用 Lemmas
\ref{perfect-lemma-quasi-coherence-pullback}、
\ref{perfect-lemma-quasi-coherence-tensor-product}, and
\ref{perfect-lemma-quasi-coherence-direct-image}
可知对象 $Rf'_*L(g')^*E$ 与 $Lg^*Rf_*E$ 的上同调层拟凝聚）。注意 $g'$
也仿射（Morphisms, Lemma \ref{morphisms-lemma-base-change-affine}）。
由 Lemma \ref{perfect-lemma-affine-morphism-pull-push}，该映射成为
$$
Rf_*E \otimes_{\mathcal{O}_S}^\mathbf{L} g_*\mathcal{O}_{S'}
\longrightarrow
Rf_*(E \otimes_{\mathcal{O}_X}^\mathbf{L} g'_*\mathcal{O}_{X'})
$$
。注意 $g'_*\mathcal{O}_{X'} = f^*g_*\mathcal{O}_{S'}$（由仿射基变换，见
Cohomology of Schemes, Lemma \ref{coherent-lemma-affine-base-change}）。
因此由 Lemma \ref{perfect-lemma-cohomology-base-change}，只需证明
$Lf^*g_*\mathcal{O}_{S'} = f^*g_*\mathcal{O}_{S'}$。这由 $X$ 与 $S'$
在 $S$ 上 Tor 独立的假设推出。事实上，为检验这一点，可以在 $X$ 上局部地工作，
所以还可假设 $X$ 仿射。记 $X = \Spec(A)$、$S = \Spec(R)$ 且
$S' = \Spec(R')$。我们的假设蕴含 $A$ 与 $R'$ 在 $R$ 上 Tor 独立
（More on Algebra, Lemma \ref{more-algebra-lemma-tor-independent}），即当
$i > 0$ 时 $\text{Tor}_i^R(A, R') = 0$。换言之，
$A \otimes_R^\mathbf{L} R' = A \otimes_R R'$，而这恰好意味着
$Lf^*g_*\mathcal{O}_{S'} = f^*g_*\mathcal{O}_{S'}$
（使用 Lemma \ref{perfect-lemma-quasi-coherence-pullback}）。
\end{proof}

\noindent
下述引理将在对偶化复形一章中使用。

\begin{lemma}
\label{perfect-lemma-affine-morphism-and-hom-out-of-perfect}
考虑拟紧拟分离概形的笛卡尔方块
$$
\xymatrix{
X' \ar[r]_{g'} \ar[d]_{f'} & X \ar[d]^f \\
S' \ar[r]^g & S
}
$$
。假设 $g$ 与 $f$ Tor 独立，且 $S = \Spec(R)$、$S' = \Spec(R')$ 仿射。
对 $M, K \in D(\mathcal{O}_X)$，典范映射
$$
R\Hom_X(M, K) \otimes^\mathbf{L}_R R'
\longrightarrow
R\Hom_{X'}(L(g')^*M, L(g')^*K)
$$
在下列两种情形中是 $D(R')$ 中的同构：
\begin{enumerate}
\item $M \in D(\mathcal{O}_X)$ 完美，且 $K \in D_\QCoh(X)$；或
\item $M \in D(\mathcal{O}_X)$ 伪凝聚，$K \in D_\QCoh^+(X)$，且
$R'$ 在 $R$ 上具有有限 Tor 维数。
\end{enumerate}
\end{lemma}

\begin{proof}
由 Cohomology, Section \ref{cohomology-section-global-RHom}，
在整体 Hom 复形的 $D(\Gamma(X, \mathcal{O}_X))$ 中有典范映射
$R\Hom_X(M, K) \to R\Hom_{X'}(L(g')^*M, L(g')^*K)$
。限制标量后可将其视为 $D(R)$ 中的映射。再利用限制标量与
$- \otimes_R^\mathbf{L} R'$ 的伴随性，便得到引理中所显示的映射。
定义该映射后，只需证明它在阿贝尔群导出范畴中是同构。

\medskip\noindent
由 Lemma \ref{perfect-lemma-affine-morphism-pull-push}，右端等于
$$
R\Hom_X(M, R(g')_*L(g')^*K) =
R\Hom_X(M, K \otimes_{\mathcal{O}_X}^\mathbf{L} g'_*\mathcal{O}_{X'})
$$
。在两种情形中，由 Lemma \ref{perfect-lemma-quasi-coherence-internal-hom}，复形
$R\SheafHom(M, K)$ 都是 $D_\QCoh(\mathcal{O}_X)$ 的对象。有自然映射
$$
R\SheafHom(M, K) \otimes_{\mathcal{O}_X}^\mathbf{L} g'_*\mathcal{O}_{X'}
\longrightarrow
R\SheafHom(M, K \otimes_{\mathcal{O}_X}^\mathbf{L} g'_*\mathcal{O}_{X'})
$$
，由 Lemma \ref{perfect-lemma-internal-hom-evaluate-tensor-isomorphism}，
它在两种情形中都是同构。为说明该引理适用于情形 (2)，注意
$g'_*\mathcal{O}_{X'} = Rg'_*\mathcal{O}_{X'} =
Lf^*g_*\mathcal{O}_{S'}$，其中第二个等式来自
Lemma \ref{perfect-lemma-compare-base-change}。使用
Lemma \ref{perfect-lemma-tor-dimension-affine} 及
Cohomology, Lemma \ref{cohomology-lemma-tor-amplitude-pullback}
可知 $g'_*\mathcal{O}_{X'}$ 具有有限 Tor 维数。因此在两种情形中，
以 $R\SheafHom(M, K)$ 替换 $K$，问题都归结为证明
$$
R\Gamma(X, K) \otimes^\mathbf{L}_A A' \longrightarrow
R\Gamma(X, K \otimes^\mathbf{L}_{\mathcal{O}_X} g'_*\mathcal{O}_{X'})
$$
是同构。注意由 Lemma \ref{perfect-lemma-affine-morphism-pull-push}，左端等于
$R\Gamma(X', L(g')^*K)$。因此结论由 Lemma \ref{perfect-lemma-compare-base-change} 推出。
\end{proof}

\begin{remark}
\label{perfect-remark-multiplication-map}
采用 Lemma \ref{perfect-lemma-affine-morphism-and-hom-out-of-perfect} 的记号。图
$$
\xymatrix{
R\Hom_X(M, Rg'_*L) \otimes_R^\mathbf{L} R' \ar[r] \ar[d]_\mu &
R\Hom_{X'}(L(g')^*M, L(g')^*Rg'_*L) \ar[d]^a \\
R\Hom_X(M, R(g')_*L) \ar@{=}[r] &
R\Hom_{X'}(L(g')^*M, L)
}
$$
可交换，其中上方水平箭头是引理中的映射，$\mu$ 是乘法映射，而 $a$ 来自伴随映射
$L(g')^*Rg'_*L \to L$。乘法映射是任意 $K' \in D(R')$ 的伴随映射
$K' \otimes_R^\mathbf{L} R' \to K'$。
\end{remark}

\begin{lemma}
\label{perfect-lemma-tor-independence-and-tor-amplitude}
考虑概形的笛卡尔方块
$$
\xymatrix{
X' \ar[r]_{g'} \ar[d]_{f'} & X \ar[d]^f \\
S' \ar[r]^g & S
}
$$
。假设 $g$ 与 $f$ Tor 独立。
\begin{enumerate}
\item 若 $E \in D(\mathcal{O}_X)$ 作为 $f^{-1}\mathcal{O}_S$-模复形具有
$[a, b]$ 中的 Tor 振幅，则 $L(g')^*E$ 作为
$f^{-1}\mathcal{O}_{S'}$-模复形具有 $[a, b]$ 中的 Tor 振幅。
\item 若 $\mathcal{G}$ 是在 $S$ 上平坦的 $\mathcal{O}_X$-模，则
$L(g')^*\mathcal{G} = (g')^*\mathcal{G}$。
\end{enumerate}
\end{lemma}

\begin{proof}
可以在茎上计算 Tor 维数，见
Cohomology, Lemma \ref{cohomology-lemma-tor-amplitude-stalk}。
若 $x' \in X'$ 的像为 $x \in X$，则
$$
(L(g')^*E)_{x'} =
E_x \otimes_{\mathcal{O}_{X, x}}^\mathbf{L} \mathcal{O}_{X', x'}
$$
。设 $s' \in S'$ 与 $s \in S$ 分别为 $x'$ 与 $x$ 的像。由于 $X$ 与
$S'$ 在 $S$ 上 Tor 独立，可应用
More on Algebra, Lemma \ref{more-algebra-lemma-base-change-comparison}
，可知所显示公式的右端在 $D(\mathcal{O}_{S', s'})$ 中等于
$E_x \otimes_{\mathcal{O}_{S, s}}^\mathbf{L} \mathcal{O}_{S', s'}$。
于是 (1) 由 More on Algebra, Lemma
\ref{more-algebra-lemma-pull-tor-amplitude} 推出。为证明 (2)，注意
$\mathcal{G}$ 平坦等价于 $\mathcal{G}[0]$ 具有 $[0, 0]$ 中的 Tor 振幅。
应用 (1) 即得结论。
\end{proof}

\begin{lemma}
\label{perfect-lemma-compare-base-change-closed-immersion}
考虑概形的笛卡尔图
$$
\xymatrix{
Z' \ar[r]_{i'} \ar[d]_g & X' \ar[d]^f \\
Z \ar[r]^i & X
}
$$
，其中 $i$ 是闭浸入。若 $Z$ 与 $X'$ 在 $X$ 上 Tor 独立，则作为函子
$D(\mathcal{O}_Z) \to D(\mathcal{O}_{X'})$，有
$Ri'_* \circ Lg^* = Lf^* \circ Ri_*$。
\end{lemma}

\begin{proof}
注意本引理须对所有 $K \in D(\mathcal{O}_Z)$ 成立。函子 $i_*$ 与 $i'_*$
都是正合的，故 $Ri_*$ 与 $Ri'_*$ 可通过把 $i_*$ 与 $i'_*$ 作用于任意代表来计算。
因此，基变换映射
$$
Lf^*(Ri_*(K)) \longrightarrow Ri'_*(Lg^*(K))
$$
在点 $z' \in Z'$ 的茎上由下述映射给出，其中 $z \in Z$ 为其像：
$$
K_z \otimes_{\mathcal{O}_{X, z}}^\mathbf{L} \mathcal{O}_{X', z'}
\longrightarrow
K_z \otimes_{\mathcal{O}_{Z, z}}^\mathbf{L} \mathcal{O}_{Z', z'}
$$
由 More on Algebra, Lemma \ref{more-algebra-lemma-base-change-comparison}
及所假定的 Tor 独立性，该映射是同构。
\end{proof}








\section{K\"unneth 公式，II}
\label{perfect-section-kunneth}

\noindent
关于基底为域的情形，见 Varieties, Section
\ref{varieties-section-kunneth}。考虑概形的笛卡尔图
$$
\xymatrix{
& X \times_S Y \ar[ld]^p \ar[rd]_q \ar[dd]^f \\
X \ar[rd]_a & & Y \ar[ld]^b \\
& S
}
$$
设 $K \in D(\mathcal{O}_X)$ 且 $M \in D(\mathcal{O}_Y)$。存在典范映射
\begin{equation}
\label{perfect-equation-kunneth}
Ra_*K \otimes_{\mathcal{O}_S}^\mathbf{L} Rb_*M
\longrightarrow
Rf_*(Lp^*K \otimes_{\mathcal{O}_{X \times_S Y}}^\mathbf{L} Lq^*M)
\end{equation}
。事实上，可以使用映射
$Ra_*K \to Ra_*Rp_* Lp^*K = Rf_*Lp^*K$ 与
$Rb_*M \to Rb_*Rq_* Lq^*M = Rf_*Lq^*M$，然后使用相对杯积
（Cohomology, Remark \ref{cohomology-remark-cup-product}）。

\medskip\noindent
令 $A = \Gamma(S, \mathcal{O}_S)$。存在整体 K\"unneth 映射
\begin{equation}
\label{perfect-equation-kunneth-global}
R\Gamma(X, K) \otimes_A^\mathbf{L} R\Gamma(Y, M)
\longrightarrow
R\Gamma(X \times_S Y,
Lp^*K \otimes_{\mathcal{O}_{X \times_S Y}}^\mathbf{L} Lq^*M)
\end{equation}
位于 $D(A)$ 中。该映射由拉回映射
$R\Gamma(X, K) \to R\Gamma(X \times_S Y, Lp^*K)$ 与
$R\Gamma(Y, M) \to R\Gamma(X \times_S Y, Lq^*M)$，以及
Cohomology, Section \ref{cohomology-section-cup-product}
中构造的杯积得到。

\begin{lemma}
\label{perfect-lemma-kunneth}
在上述情形中，若 $a$ 与 $b$ 拟紧拟分离，且 $X$ 与 $Y$ 在 $S$ 上 Tor 独立，
则对 $K \in D_\QCoh(\mathcal{O}_X)$ 与 $M \in D_\QCoh(\mathcal{O}_Y)$，
(\ref{perfect-equation-kunneth}) 是同构。若此外 $S = \Spec(A)$ 仿射，则映射
(\ref{perfect-equation-kunneth-global}) 是同构。
\end{lemma}

\begin{proof}[第一种证明]
结论由下列一串同构推出：
\begin{align*}
Rf_*(Lp^*K \otimes_{\mathcal{O}_{X \times_S Y}}^\mathbf{L} Lq^*M)
& =
Ra_*Rp_*(Lp^*K \otimes_{\mathcal{O}_{X \times_S Y}}^\mathbf{L} Lq^*M) \\
& =
Ra_*(K \otimes_{\mathcal{O}_X}^\mathbf{L} Rp_*Lq^*M) \\
& =
Ra_*(K \otimes_{\mathcal{O}_X}^\mathbf{L} La^*Rb_*M) \\
& =
Ra_*K \otimes_{\mathcal{O}_S}^\mathbf{L} Rb_*M
\end{align*}
第一个等式成立是因为 $f = a \circ p$；第二个等式来自
Lemma \ref{perfect-lemma-cohomology-base-change}；第三个等式来自
Lemma \ref{perfect-lemma-compare-base-change}；第四个等式来自
Lemma \ref{perfect-lemma-cohomology-base-change}。略去验证这些同构的复合正是映射
(\ref{perfect-equation-kunneth})。若 $S$ 仿射，则箭头
(\ref{perfect-equation-kunneth-global}) 的源和靶分别是把 $R\Gamma(S, -)$
作用于 (\ref{perfect-equation-kunneth}) 的源和靶所得，故得到最后的陈述；细节从略。
\end{proof}

\begin{proof}[第二种证明]
箭头 (\ref{perfect-equation-kunneth}) 的构造与限制到 $S$ 的开子概形相容，
这由相对杯积的构造立即可见。因此只需在 $S$ 仿射时证明
(\ref{perfect-equation-kunneth}) 是同构。

\medskip\noindent
设 $S = \Spec(A)$ 仿射。由 Leray，
$R\Gamma(S, Rf_*K) = R\Gamma(X, K)$，其他情形也类似。由
Cohomology, Lemma \ref{cohomology-lemma-compose-cup-product}，
对映射 (\ref{perfect-equation-kunneth}) 取 $R\Gamma(S, -)$ 后诱导出映射
(\ref{perfect-equation-kunneth-global})。另一方面，由 Lemmas
\ref{perfect-lemma-quasi-coherence-direct-image} 与
\ref{perfect-lemma-quasi-coherence-tensor-product}
，映射 (\ref{perfect-equation-kunneth}) 的源与靶都属于 $D_\QCoh(\mathcal{O}_S)$。
所以由 Lemma \ref{perfect-lemma-affine-compare-bounded}，只需证明
(\ref{perfect-equation-kunneth-global}) 是同构。

\medskip\noindent
设 $S = \Spec(A)$、$X = \Spec(B)$、$Y = \Spec(C)$ 均仿射。
下文不再特别说明地使用 Lemma \ref{perfect-lemma-affine-compare-bounded}。
此时可选取各项平坦的 $B$-模 K-平坦复形 $K^\bullet$，使得
$K$ 由 $\widetilde{K}^\bullet$ 表示。类似地，可选取各项平坦的
$C$-模 K-平坦复形 $M^\bullet$，使得 $M$ 由 $\widetilde{M}^\bullet$ 表示。
见 More on Algebra, Lemma \ref{more-algebra-lemma-K-flat-resolution}。
于是 $\widetilde{K}^\bullet$ 是 $\mathcal{O}_X$-模的 K-平坦复形，
$\widetilde{M}^\bullet$ 也类似，见 Lemma \ref{perfect-lemma-affine-K-flat}。因此
$La^*K$ 由
$$
a^*\widetilde{K}^\bullet = \widetilde{K^\bullet \otimes_A C}
$$
表示，$Lb^*M$ 也类似。由上述引理及 More on Algebra, Lemma
\ref{more-algebra-lemma-base-change-K-flat}，这又是
$\mathcal{O}_{X \times_S Y}$-模的 K-平坦复形。最终可见，与
$$
\text{Tot}((K^\bullet \otimes_A C) \otimes_{B \otimes_A C}
(B \otimes_A M^\bullet)) =
\text{Tot}(K^\bullet \otimes_A M^\bullet)
$$
相伴的 $\mathcal{O}_{X \times_S Y}$-模复形在 $X \times_S Y$ 的导出范畴中表示
$La^*K \otimes_{\mathcal{O}_{X \times_S Y}}^\mathbf{L} Lb^*M$。
取整体截面得到 $\text{Tot}(K^\bullet \otimes_A M^\bullet)$；当然它也表示
$R\Gamma(X, K) \otimes_A^\mathbf{L} R\Gamma(Y, M)$。该同构由杯积给出这一事实，
可由 Cohomology, Section \ref{cohomology-section-cup-product} 中真正杯积与朴素杯积的关系推出
（特别见 Cohomology, Lemma \ref{cohomology-lemma-cup-compatible-with-naive}
及其后的讨论）。

\medskip\noindent
假设 $S = \Spec(A)$ 与 $Y$ 仿射。我们用
Cohomology of Schemes, Lemma \ref{coherent-lemma-induction-principle}
中的归纳原理证明该陈述。为此，只需证明：若 $X = U \cup V$ 是开覆盖，
且 $U$ 与 $V$ 拟紧，并且 $U$ 与 $Y$ 在 $S$ 上的映射
(\ref{perfect-equation-kunneth-global})
$$
R\Gamma(U, K) \otimes_A^\mathbf{L} R\Gamma(Y, M)
\longrightarrow
R\Gamma(U \times_S Y,
Lp^*K \otimes_{\mathcal{O}_{X \times_S Y}}^\mathbf{L} Lq^*M)
$$
，$V$ 与 $Y$ 在 $S$ 上的映射 (\ref{perfect-equation-kunneth-global})
$$
R\Gamma(V, K) \otimes_A^\mathbf{L} R\Gamma(Y, M)
\longrightarrow
R\Gamma(V \times_S Y,
Lp^*K \otimes_{\mathcal{O}_{X \times_S Y}}^\mathbf{L} Lq^*M)
$$
，以及 $U \cap V$ 与 $Y$ 在 $S$ 上的映射 (\ref{perfect-equation-kunneth-global})
$$
R\Gamma(U \cap V, K) \otimes_A^\mathbf{L} R\Gamma(Y, M)
\longrightarrow
R\Gamma((U \cap V) \times_S Y,
Lp^*K \otimes_{\mathcal{O}_{X \times_S Y}}^\mathbf{L} Lq^*M)
$$
均为同构，则 $X$ 与 $Y$ 在 $S$ 上的映射
(\ref{perfect-equation-kunneth-global}) 也是同构。然而，由 Cohomology, Lemma
\ref{cohomology-lemma-mayer-vietoris-cup}，这些映射组成可区别三角之间的一个映射，
而 (\ref{perfect-equation-kunneth-global}) 是最后一边。因此由
Derived Categories, Lemma \ref{derived-lemma-third-isomorphism-triangle}
得到结论。

\medskip\noindent
设 $S = \Spec(A)$ 仿射。为完成证明，可在 $Y$ 上使用
Cohomology of Schemes, Lemma \ref{coherent-lemma-induction-principle}
的归纳原理。事实上，由上已知当 $Y$ 仿射时该映射是同构。其余论证与上一段完全相同，
只是交换 $X$ 与 $Y$ 的角色。
\end{proof}

\begin{lemma}
\label{perfect-lemma-cohomology-de-rham-base-change}
设 $a : X \to S$ 为概形的拟紧拟分离态射。设 $\mathcal{F}^\bullet$
为 $a^{-1}\mathcal{O}_S$-模的局部有界复形。假设对所有 $n \in \mathbf{Z}$，
层 $\mathcal{F}^n$ 都是平坦 $a^{-1}\mathcal{O}_S$-模，且 $\mathcal{F}^n$ 具有与给定的
$a^{-1}\mathcal{O}_S$-模结构相容的拟凝聚 $\mathcal{O}_X$-模结构
（但复形 $\mathcal{F}^\bullet$ 中的微分不必是 $\mathcal{O}_X$-线性的）。则：
\begin{enumerate}
\item $Ra_*\mathcal{F}^\bullet$ 局部有界；
\item $Ra_*\mathcal{F}^\bullet$ 属于 $D_\QCoh(\mathcal{O}_S)$；
\item $Ra_*\mathcal{F}^\bullet$ 局部具有有限 Tor 维数；
\item $\mathcal{G} \otimes_{\mathcal{O}_S}^\mathbf{L} Ra_*\mathcal{F}^\bullet =
Ra_*(a^{-1}\mathcal{G} \otimes_{a^{-1}\mathcal{O}_S} \mathcal{F}^\bullet)$
对 $\mathcal{G} \in \QCoh(\mathcal{O}_S)$ 成立；
\item $K \otimes_{\mathcal{O}_S}^\mathbf{L} Ra_*\mathcal{F}^\bullet =
Ra_*(a^{-1}K \otimes_{a^{-1}\mathcal{O}_S}^\mathbf{L} \mathcal{F}^\bullet)$
对 $K \in D_\QCoh(\mathcal{O}_S)$ 成立。
\end{enumerate}
\end{lemma}

\begin{proof}
(1)、(2)、(3) 在 $S$ 上是局部的，故可且确实假设 $S$ 仿射。由于 $a$ 拟紧，
$X$ 拟紧。又因 $\mathcal{F}^\bullet$ 局部有界，故 $\mathcal{F}^\bullet$ 有界。

\medskip\noindent
对 (1) 与 (2)，可以使用 Derived Categories, Lemma
\ref{derived-lemma-two-ss-complex-functor} 的第一个谱序列
$R^pa_*\mathcal{F}^q \Rightarrow R^{p + q}a_*\mathcal{F}^\bullet$
。结合 Cohomology of Schemes, Lemma
\ref{coherent-lemma-quasi-coherence-higher-direct-images}
与 Homology, Lemma \ref{homology-lemma-biregular-ss-converges} 即得结论。

\medskip\noindent
用 Cohomology of Schemes, Lemma
\ref{coherent-lemma-induction-principle} 的归纳原理证明 (3)。具体地，对 $X$
的拟紧开集 $U$，考虑 $R(a|_U)_*(\mathcal{F}^\bullet|_U)$ 具有有限 Tor 维数
这一条件。若 $U, V$ 是 $X$ 中的拟紧开集，则由
Cohomology, Lemma \ref{cohomology-lemma-unbounded-relative-mayer-vietoris}
有相对 Mayer--Vietoris 可区别三角
$$
R(a|_{U \cup V})_*\mathcal{F}^\bullet|_{U \cup V} \to
R(a|_U)_*\mathcal{F}^\bullet|_U \oplus
R(a|_V)_*\mathcal{F}^\bullet|_V \to
R(a|_{U \cap V})_*\mathcal{F}^\bullet|_{U \cap V} \to
$$
。由 Tor 振幅在可区别三角中的行为
（见 Cohomology, Lemma \ref{cohomology-lemma-cone-tor-amplitude}），
若已知结论对 $U$、$V$、$U \cap V$ 成立，则也对 $U \cup V$ 成立。
这把问题归结为 $X$ 仿射的情形。此时
$$
Ra_*\mathcal{F}^\bullet = a_*\mathcal{F}^\bullet
$$
，这是由 Leray 无环引理
（Derived Categories, Lemma \ref{derived-lemma-leray-acyclicity}）
及仿射态射下拟凝聚模的高阶直像消失
（Cohomology of Schemes, Lemma \ref{coherent-lemma-relative-affine-vanishing}）
得到的。按假设 $\mathcal{F}^n$ 在 $S$ 上平坦，且 $X$ 仿射，所以对所有 $n$，
模 $a_*\mathcal{F}^n$ 都平坦。因此 $a_*\mathcal{F}^\bullet$ 是平坦
$\mathcal{O}_S$-模的有界复形，从而具有有限 Tor 维数。

\medskip\noindent
证明 (5)。记
$a' : (X, a^{-1}\mathcal{O}_S) \to (S, \mathcal{O}_S)$
为显然的环化空间平坦态射。(5) 断言
$$
K \otimes_{\mathcal{O}_S}^\mathbf{L} Ra'_*\mathcal{F}^\bullet =
Ra'_*(L(a')^*K \otimes_{a^{-1}\mathcal{O}_S}^\mathbf{L}
\mathcal{F}^\bullet)
$$
。因此 Cohomology, Equation
(\ref{cohomology-equation-projection-formula-map}) 给出从左端到右端的函子性映射，
我们要证明它是同构。这个问题在 $S$ 上是局部的，故可且确实假设 $S$ 仿射。
余下证明与 Lemma \ref{perfect-lemma-cohomology-base-change} 的证明{\it 完全}相同，
但须证明函子
$K \mapsto Ra'_*(L(a')^*K \otimes_{a^{-1}\mathcal{O}_S}^\mathbf{L}
\mathcal{F}^\bullet)$ 与直和可交换。此处要用到 $\mathcal{F}^n$
具有拟凝聚 $\mathcal{O}_X$-模结构。事实上，注意函子
$K \mapsto L(a')^*K
\otimes_{a^{-1}\mathcal{O}_S}^\mathbf{L} \mathcal{F}^\bullet$
与任意直和可交换。其次，若 $\mathcal{F}^\bullet$ 仅由单个拟凝聚
$\mathcal{O}_X$-模 $\mathcal{F}^\bullet = \mathcal{F}^n[-n]$ 构成，则由
Cohomology, Lemma \ref{cohomology-lemma-variant-derived-pullback} 有
$L(a')^*G
\otimes_{a^{-1}\mathcal{O}_S}^\mathbf{L} \mathcal{F}^\bullet =
La^*K \otimes_{\mathcal{O}_X}^\mathbf{L} \mathcal{F}^n[-n]$。
故在此情形中，与直和可交换由
Lemma \ref{perfect-lemma-quasi-coherence-pushforward-direct-sums} 推出。一般地，
由于 $S$ 仿射（因而 $X$ 拟紧），且 $\mathcal{F}^\bullet$ 局部有界，可知
$$
\mathcal{F}^\bullet = (\mathcal{F}^a \to \ldots \to \mathcal{F}^b)
$$
有界。对 $b - a$ 归纳并考虑可区别三角
$$
\mathcal{F}^b[-b] \to (\mathcal{F}^a \to \ldots \to \mathcal{F}^b)
\to (\mathcal{F}^a \to \ldots \to \mathcal{F}^{b - 1}) \to
\mathcal{F}^b[-b + 1]
$$
即可完成此部分的证明。略去若干细节。

\medskip\noindent
证明 (4)。令 $a' : (X, a^{-1}\mathcal{O}_S) \to (S, \mathcal{O}_S)$
如上。由于 $\mathcal{F}^\bullet$ 是平坦 $a^{-1}\mathcal{O}_S$-模的局部有界复形，
复形
$a^{-1}\mathcal{G} \otimes_{a^{-1}\mathcal{O}_S} \mathcal{F}^\bullet$
在 $D(a^{-1}\mathcal{O}_S)$ 中表示
$L(a')^*\mathcal{G}
\otimes_{a^{-1}\mathcal{O}_S}^\mathbf{L}
\mathcal{F}^\bullet$。因此 (4) 由 (5) 推出。
\end{proof}

\begin{lemma}
\label{perfect-lemma-K-flat}
设 $f : X \to Y$ 为概形态射，且 $Y = \Spec(A)$ 仿射。设
$\mathcal{U} : X = \bigcup_{i \in I} U_i$ 为有限仿射开覆盖，使得所有有限交
$U_{i_0 \ldots i_p} = U_{i_0} \cap \ldots \cap U_{i_p}$ 都仿射。
设 $\mathcal{F}^\bullet$ 为 $f^{-1}\mathcal{O}_Y$-模的有界复形。假设对所有
$n \in \mathbf{Z}$，层 $\mathcal{F}^n$ 都是平坦 $f^{-1}\mathcal{O}_Y$-模，
且 $\mathcal{F}^n$ 具有与给定的 $p^{-1}\mathcal{O}_Y$-模结构相容的拟凝聚
$\mathcal{O}_X$-模结构（但复形 $\mathcal{F}^\bullet$ 中的微分不必是
$\mathcal{O}_X$-线性的）。则复形
$\text{Tot}(\check{\mathcal{C}}^\bullet(\mathcal{U}, \mathcal{F}^\bullet))$
作为 $A$-模复形是 K-平坦的。
\end{lemma}

\begin{proof}
可以写成
$$
\mathcal{F}^\bullet = (\mathcal{F}^a \to \ldots \to \mathcal{F}^b)
$$
。对 $b - a$ 归纳，考虑可区别三角
$$
\mathcal{F}^b[-b] \to (\mathcal{F}^a \to \ldots \to \mathcal{F}^b)
\to (\mathcal{F}^a \to \ldots \to \mathcal{F}^{b - 1}) \to
\mathcal{F}^b[-b + 1]
$$
并使用
More on Algebra, Lemma \ref{more-algebra-lemma-K-flat-two-out-of-three}
，问题归结为 $\mathcal{F}^\bullet$ 仅由置于次数 $0$ 的单个拟凝聚
$\mathcal{O}_X$-模 $\mathcal{F}$ 构成的情形。此时，$\mathcal{F}$ 与
$\mathcal{U}$ 的 {\v C}ech 复形同伦等价于交错 {\v C}ech 复形，见
Cohomology, Lemma \ref{cohomology-lemma-alternating-usual}。由于
$U_{i_0 \ldots i_p}$ 总是仿射的，$\mathcal{F}(U_{i_0 \ldots i_p})$
在 $A$ 上平坦。因此
$\check{\mathcal{C}}_{alt}^\bullet(\mathcal{U}, \mathcal{F})$
是平坦 $A$-模的有界复形，从而由 More on Algebra, Lemma
\ref{more-algebra-lemma-derived-tor-quasi-isomorphism} 可知它是 K-平坦的。
\end{proof}

\noindent
令 $X, Y, S, a, b, p, q, f$ 如本节引言所述。设 $\mathcal{F}$ 为
$\mathcal{O}_X$-模，$\mathcal{G}$ 为 $\mathcal{O}_Y$-模，并令
$A = \Gamma(S, \mathcal{O}_S)$。考虑 $D(A)$ 中的映射
\begin{equation}
\label{perfect-equation-kunneth-single-sheaves}
R\Gamma(X, \mathcal{F})
\otimes_A^\mathbf{L}
R\Gamma(Y, \mathcal{G})
\longrightarrow
R\Gamma(X \times_S Y,
p^*\mathcal{F}
\otimes_{\mathcal{O}_{X \times_S Y}}
q^*\mathcal{G})
\end{equation}
。该映射由拉回映射
$R\Gamma(X, \mathcal{F}) \to R\Gamma(X \times_S Y, p^*\mathcal{F})$
与 $R\Gamma(Y, \mathcal{G}) \to R\Gamma(X \times_S Y, q^*\mathcal{G})$、
Cohomology, Section \ref{cohomology-section-cup-product} 中构造的杯积，
以及典范映射
$p^*\mathcal{F}
\otimes_{\mathcal{O}_{X \times_S Y}}^\mathbf{L}
q^*\mathcal{G}
\to
p^*\mathcal{F}
\otimes_{\mathcal{O}_{X \times_S Y}}
q^*\mathcal{G}$ 构造。

\begin{lemma}
\label{perfect-lemma-kunneth-single-sheaf}
在上述情形中，若 $S$ 仿射，$\mathcal{F}$ 与 $\mathcal{G}$ 在 $S$ 上平坦且拟凝聚，
并且 $X$ 与 $Y$ 拟紧且具有仿射对角态射，则映射
(\ref{perfect-equation-kunneth-single-sheaves}) 是同构。
\end{lemma}

\begin{proof}
强烈建议读者先阅读 Varieties, Lemma \ref{varieties-lemma-kunneth} 的证明。
选取有限仿射开覆盖 $\mathcal{U} : X = \bigcup_{i \in I} U_i$ 与
$\mathcal{V} : Y = \bigcup_{j \in J} V_j$。这确定了仿射开覆盖
$\mathcal{W} : X \times_S Y = \bigcup_{(i, j) \in I \times J} U_i \times_S V_j$。
注意 $\mathcal{W}$ 是 $\text{pr}_1^{-1}\mathcal{U}$ 与
$\text{pr}_2^{-1}\mathcal{V}$ 的加细。因此由 Cohomology, Section
\ref{cohomology-section-cech-cohomology-of-complexes} 中的讨论，得到映射
$$
\check{\mathcal{C}}^\bullet(\mathcal{U}, \mathcal{F})
\to
\check{\mathcal{C}}^\bullet(\mathcal{W}, p^*\mathcal{F})
\quad\text{且}\quad
\check{\mathcal{C}}^\bullet(\mathcal{V}, \mathcal{G})
\to
\check{\mathcal{C}}^\bullet(\mathcal{W}, q^*\mathcal{G})
$$
；它们在同伦意义下良定义，并与上同调上的拉回映射相容。在
Cohomology, Equation (\ref{cohomology-equation-needs-signs}) 中，
我们构造了复形映射
$$
\text{Tot}(
\check{\mathcal{C}}^\bullet(\mathcal{W}, p^*\mathcal{F})
\otimes_A
\check{\mathcal{C}}^\bullet(\mathcal{W}, q^*\mathcal{G}))
\longrightarrow
\check{\mathcal{C}}^\bullet(\mathcal{W},
p^*\mathcal{F} \otimes_{\mathcal{O}_{X \times_S Y}}
q^*\mathcal{G})
$$
；由 Cohomology, Lemma \ref{cohomology-lemma-diagrams-commute}，
它与上同调上的杯积相容。结合上述映射，得到复形映射
\begin{equation}
\label{perfect-equation-kunneth-on-cech}
\text{Tot}(
\check{\mathcal{C}}^\bullet(\mathcal{U}, \mathcal{F})
\otimes_A
\check{\mathcal{C}}^\bullet(\mathcal{V}, \mathcal{G}))
\to
\check{\mathcal{C}}^\bullet(\mathcal{W},
p^*\mathcal{F}
\otimes_{\mathcal{O}_{X \times_S Y}}
q^*\mathcal{G})
\end{equation}
我们断言，这就是引理陈述中的映射；即此箭头的源和靶分别与
(\ref{perfect-equation-kunneth-single-sheaves}) 的源和靶相同。事实上，由
Cohomology of Schemes, Lemma
\ref{coherent-lemma-quasi-coherent-affine-cohomology-zero}
及
Cohomology, Lemma \ref{cohomology-lemma-cech-complex-complex-computes}
，典范映射
$$
\check{\mathcal{C}}^\bullet(\mathcal{U}, \mathcal{F})
\to
R\Gamma(X, \mathcal{F}),
\quad
\check{\mathcal{C}}^\bullet(\mathcal{V}, \mathcal{G})
\to
R\Gamma(Y, \mathcal{G})
$$
and
$$
\check{\mathcal{C}}^\bullet(\mathcal{W},
p^*\mathcal{F} \otimes_{\mathcal{O}_{X \times_S Y}} q^*\mathcal{G})
\to
R\Gamma(X \times_S Y,
p^*\mathcal{F} \otimes_{\mathcal{O}_{X \times_S Y}}
q^*\mathcal{G})
$$
均为同构。另一方面，复形
$\check{\mathcal{C}}^\bullet(\mathcal{U}, \mathcal{F})$
由 Lemma \ref{perfect-lemma-K-flat} 是 K-平坦的，故
$\text{Tot}(
\check{\mathcal{C}}^\bullet(\mathcal{U}, \mathcal{F})
\otimes_A
\check{\mathcal{C}}^\bullet(\mathcal{V}, \mathcal{G}))$
表示导出张量积
$R\Gamma(X, \mathcal{F}) \otimes_A^\mathbf{L} R\Gamma(Y, \mathcal{G})$
，正如所断言。

\medskip\noindent
还须证明 (\ref{perfect-equation-kunneth-on-cech}) 是拟同构。我们用维数移位来证明。
令 $d(\mathcal{F}) = \max \{d \mid H^d(X, \mathcal{F}) \not = 0\}$。
假设 $d(\mathcal{F}) > 0$，并令 $U = \coprod\nolimits_{i \in I} U_i$。
由于 $I$ 有限，这是仿射概形。记 $j : U \to X$ 为在每个 $U_i$ 上等于包含
$U_i \to X$ 的态射。由于 $X$ 的对角态射仿射，态射 $j$ 仿射，见
Morphisms, Lemma \ref{morphisms-lemma-affine-permanence}。于是
$\mathcal{F}' = j_*j^*\mathcal{F}$ 在 $S$ 上平坦，见
Morphisms, Lemma \ref{morphisms-lemma-pushforward-flat-affine}。结合
Cohomology of Schemes, Lemmas
\ref{coherent-lemma-relative-affine-cohomology} and
\ref{coherent-lemma-quasi-coherent-affine-cohomology-zero}，还可得
$d(\mathcal{F}') = 0$。对所有 $x \in X$，映射
$\mathcal{F}_x  \to \mathcal{F}'_x$ 都是直和项的包含：若 $x \in U_i$，
则 $\mathcal{F}' \to (U_i \to X)_*\mathcal{F}|_{U_i}$ 给出分裂。
故 $\mathcal{F} \to \mathcal{F}'$ 为单射，且
$\mathcal{F}'' = \mathcal{F}'/\mathcal{F}$ 也在 $S$ 上平坦。平坦拟凝聚
$\mathcal{O}_X$-模的短正合列
$0 \to \mathcal{F} \to \mathcal{F}' \to \mathcal{F}'' \to 0$
给出复形的短正合列
$$
0 \to
\text{Tot}(
\check{\mathcal{C}}^\bullet(\mathcal{U}, \mathcal{F})
\otimes_A
\check{\mathcal{C}}^\bullet(\mathcal{V}, \mathcal{G})) \to
\text{Tot}(
\check{\mathcal{C}}^\bullet(\mathcal{U}, \mathcal{F}')
\otimes_A
\check{\mathcal{C}}^\bullet(\mathcal{V}, \mathcal{G})) \to
\text{Tot}(
\check{\mathcal{C}}^\bullet(\mathcal{U}, \mathcal{F}'')
\otimes_A
\check{\mathcal{C}}^\bullet(\mathcal{V}, \mathcal{G})) \to 0
$$
以及复形的短正合列
$$
0 \to
\check{\mathcal{C}}^\bullet(\mathcal{W},
p^*\mathcal{F}
\otimes_{\mathcal{O}_{X \times_S Y}}
q^*\mathcal{G}) \to
\check{\mathcal{C}}^\bullet(\mathcal{W},
p^*\mathcal{F}'
\otimes_{\mathcal{O}_{X \times_S Y}}
q^*\mathcal{G}) \to
\check{\mathcal{C}}^\bullet(\mathcal{W},
p^*\mathcal{F}''
\otimes_{\mathcal{O}_{X \times_S Y}}
q^*\mathcal{G}) \to 0
$$
此外，它们之间的映射 (\ref{perfect-equation-kunneth-on-cech}) 与这些短正合列相容。
所以只需对 $\mathcal{F}'$ 与 $\mathcal{F}''$ 证明
(\ref{perfect-equation-kunneth-on-cech}) 是同构。最后，
$d(\mathcal{F}'') < d(\mathcal{F})$。这样便归结为
$d(\mathcal{F}) = 0$ 的情形。

\medskip\noindent
对 $\mathcal{G}$ 作同样论证，可假设 $\mathcal{F}$ 与 $\mathcal{G}$
都只在次数 $0$ 上具有非零上同调。注意这意味着 $\Gamma(X, \mathcal{F})$
拟同构于位于次数 $\geq 0$ 的 $K$-平坦 $A$-模复形
$\check{\mathcal{C}}^\bullet(\mathcal{U}, \mathcal{F})$。
因此 $\Gamma(X, \mathcal{F})$ 是平坦 $A$-模
（因为与此模的高阶 Tor 可通过同 {\v C}ech 复形作张量来计算）。
设 $V \subset Y$ 为仿射开集。考虑仿射开覆盖
$\mathcal{U}_V : X \times_S V = \bigcup_{i \in I} U_i \times_S V$。
立即有
$$
\check{\mathcal{C}}^\bullet(\mathcal{U}, \mathcal{F})
\otimes_A \mathcal{G}(V) =
\check{\mathcal{C}}^\bullet(\mathcal{U}_V,
p^*\mathcal{F} \otimes_{\mathcal{O}_{X \times Y}}
q^*\mathcal{G})
$$
（复形相等）。由 $\mathcal{G}(V)$ 在 $A$ 上平坦，可知
$\Gamma(X, \mathcal{F}) \otimes_A \mathcal{G}(V) \to
\check{\mathcal{C}}^\bullet(\mathcal{U}, \mathcal{F})
\otimes_A \mathcal{G}(V)$ 是拟同构。由于预层
$V \mapsto \check{\mathcal{C}}^\bullet(\mathcal{U}_V,
p^*\mathcal{F} \otimes_{\mathcal{O}_{X \times Y}}
q^*\mathcal{G})$ 的层化表示
$Rq_*(p^*\mathcal{F} \otimes_{\mathcal{O}_{X \times Y}} q^*\mathcal{G})$
，见 Cohomology of Schemes, Lemma
\ref{coherent-lemma-separated-case-relative-cech}
，所以在 $Y$ 上有
$$
Rq_*(p^*\mathcal{F} \otimes_{\mathcal{O}_{X \times Y}} q^*\mathcal{G})
\cong
\Gamma(X, \mathcal{F}) \otimes_A \mathcal{G}
$$
，其中右端记号表示模
$$
b^*\widetilde{\Gamma(X, \mathcal{F})} \otimes_{\mathcal{O}_Y} \mathcal{G}
$$
。使用 $q$ 的 Leray 谱序列，得到
$$
H^n(X \times_S Y, p^*\mathcal{F} \otimes_{\mathcal{O}_{X \times Y}}
q^*\mathcal{G}) =
H^n(Y,
b^*\widetilde{\Gamma(X, \mathcal{F})} \otimes_{\mathcal{O}_Y} \mathcal{G})
$$
。对态射 $b : Y \to S = \Spec(A)$ 使用
Lemma \ref{perfect-lemma-cohomology-base-change}，并利用 $\Gamma(X, \mathcal{F})$
在 $A$ 上平坦，可知
$H^n(X \times_S Y, p^*\mathcal{F} \otimes_{\mathcal{O}_{X \times Y}}
q^*\mathcal{G})$ 在 $n > 0$ 时为零，在 $n = 0$ 时同构于
$H^0(X, \mathcal{F}) \otimes_A H^0(Y, \mathcal{G})$。当然，这里还用到
$\mathcal{G}$ 只在次数 $0$ 上有上同调。证明至此完成
（只差检验该同构在次数 $0$ 上确实由杯积给出；我们略去此检验）。
\end{proof}

\begin{remark}
\label{perfect-remark-annoying-compatibility}
设 $S = \Spec(A)$ 为仿射概形。设 $a : X \to S$ 与 $b : Y \to S$
为概形态射。设 $\mathcal{F}$、$\mathcal{G}$ 为拟凝聚 $\mathcal{O}_X$-模，
且 $\mathcal{E}$ 为拟凝聚 $\mathcal{O}_Y$-模。设
$\xi \in H^i(X, \mathcal{G})$，其拉回为
$p^*\xi \in H^i(X \times_S Y, p^*\mathcal{G})$。则下图可交换：
$$
\xymatrix{
R\Gamma(X, \mathcal{F})[-i] \otimes_A^\mathbf{L} R\Gamma(Y, \mathcal{E})
\ar[d] \ar[rr]_-{\xi \otimes \text{id}} & &
R\Gamma(X, \mathcal{G} \otimes_{\mathcal{O}_X} \mathcal{F})
\otimes_A^\mathbf{L} R\Gamma(Y, \mathcal{E}) \ar[d] \\
R\Gamma(X \times_S Y, p^*\mathcal{F} \otimes q^*\mathcal{E})[-i]
\ar[rr]^-{p^*\xi} & &
R\Gamma(X \times_S Y,
p^*(\mathcal{G} \otimes_{\mathcal{O}_X} \mathcal{F}) \otimes q^*\mathcal{E})
}
$$
，其中未带下标的张量积均在 $\mathcal{O}_{X \times_S Y}$ 上取。
水平箭头来自 Cohomology, Remark
\ref{cohomology-remark-cup-with-element-map-total-cohomology}
，竖直箭头是 (\ref{perfect-equation-kunneth-global})，因而由先拉回、再在
$X \times_S Y$ 上作杯积给出。该图可交换，是因为整体杯积
（在 $X \times_S Y$ 上对层 $p^*\mathcal{G}$、$p^*\mathcal{F}$ 与
$q^*\mathcal{E}$）具有结合性，见
Cohomology, Lemma \ref{cohomology-lemma-cup-product-associative}。
\end{remark}







\section{K\"unneth 公式，III}
\label{perfect-section-kunneth-complexes}

\noindent
令 $X, Y, S, a, b, p, q, f$ 如 Section \ref{perfect-section-kunneth} 的引言所述。
本节中，给定 $\mathcal{O}_X$-模 $\mathcal{F}$ 与
$\mathcal{O}_Y$-模 $\mathcal{G}$，令
$$
\mathcal{F} \boxtimes \mathcal{G} =
p^*\mathcal{F} \otimes_{\mathcal{O}_{X \times_S Y}} q^*\mathcal{G}
$$
。注意，与后面某节中的情形不同，这里取的是非导出张量积。

\medskip\noindent
在 $X$ 上，设 $\mathcal{F}^\bullet$ 为阿贝尔群层复形，其各项是拟凝聚
$\mathcal{O}_X$-模，且微分
$d^i_\mathcal{F} : \mathcal{F}^i \to \mathcal{F}^{i + 1}$
是 $X/S$ 上有限阶微分算子，见 Morphisms, Section
\ref{morphisms-section-differential-operators}。类似地，在 $Y$ 上，
设 $\mathcal{G}^\bullet$ 为阿贝尔群层复形，其各项是拟凝聚 $\mathcal{O}_Y$-模，
且微分
$d^j_\mathcal{G} : \mathcal{G}^j \to \mathcal{G}^{j + 1}$
是 $Y/S$ 上有限阶微分算子。应用
Morphisms, Lemma \ref{morphisms-lemma-base-change-differential-operators}
中的构造，得到双复形
$$
\xymatrix{
\ldots &
\ldots &
\ldots &
\ldots \\
\ldots \ar[r] &
\mathcal{F}^i \boxtimes \mathcal{G}^{j + 1}
\ar[r]^{d_1^{i, j + 1}} \ar[u] &
\mathcal{F}^{i + 1} \boxtimes \mathcal{G}^{j + 1} \ar[r] \ar[u] &
\ldots \\
\ldots \ar[r] &
\mathcal{F}^i \boxtimes \mathcal{G}^j
\ar[r]^{d_1^{i, j}} \ar[u]^{d_2^{i, j}} &
\mathcal{F}^{i + 1} \boxtimes \mathcal{G}^j \ar[r] \ar[u]_{d_2^{i + 1, j}} &
\ldots \\
\ldots &
\ldots \ar[u] &
\ldots \ar[u] &
\ldots
}
$$
，其各项是拟凝聚模，映射是 $X \times_S Y / S$ 上有限阶微分算子。参见
Morphisms, Remark \ref{morphisms-remark-base-change-differential-operators}
及
Homology, Example \ref{homology-example-double-complex-as-tensor-product-of}.
明确地说，我们令
$$
d_1^{i, j} = d^i_\mathcal{F} \boxtimes 1
\quad\text{且}\quad
d_2^{i, j} = 1 \boxtimes d^j_\mathcal{G}
$$
。下文记号
$$
\text{Tot}(\mathcal{F}^\bullet \boxtimes \mathcal{G}^\bullet)
$$
指与此双复形相伴的全复形。该复形的各项为拟凝聚
$\mathcal{O}_{X \times_S Y}$-模，其微分为 $X \times_S Y / S$ 上有限阶微分算子。

\medskip\noindent
在上述情形中，存在“相对杯积”映射
\begin{equation}
\label{perfect-equation-relative-de-rham-kunneth}
Ra_*(\mathcal{F}^\bullet)
\otimes_{\mathcal{O}_S}^\mathbf{L}
Rb_*(\mathcal{G}^\bullet)
\longrightarrow
Rf_*\left(\text{Tot}(\mathcal{F}^\bullet \boxtimes \mathcal{G}^\bullet)\right)
\end{equation}
。具体地，可组合下列映射来构造它：
\begin{enumerate}
\item $Ra_*(\mathcal{F}^\bullet) \to Rf_*(p^{-1}\mathcal{F}^\bullet)$,
\item $Rb_*(\mathcal{G}^\bullet) \to Rf_*(q^{-1}\mathcal{G}^\bullet)$,
\item
$Rf_*(p^{-1}\mathcal{F}^\bullet)
\otimes_{\mathcal{O}_S}^\mathbf{L}
Rf_*(q^{-1}\mathcal{G}^\bullet)
\to
Rf_*(p^{-1}\mathcal{F}^\bullet \otimes_{f^{-1}\mathcal{O}_S}^\mathbf{L}
q^{-1}\mathcal{G}^\bullet)$,
\item
$p^{-1}\mathcal{F}^\bullet \otimes_{f^{-1}\mathcal{O}_S}^\mathbf{L}
q^{-1}\mathcal{G}^\bullet \to
\text{Tot}(p^{-1}\mathcal{F}^\bullet \otimes_{f^{-1}\mathcal{O}_S}
q^{-1}\mathcal{G}^\bullet)$
\item
$\text{Tot}(p^{-1}\mathcal{F}^\bullet \otimes_{f^{-1}\mathcal{O}_S}
q^{-1}\mathcal{G}^\bullet) \to \text{Tot}(\mathcal{F}^\bullet \boxtimes
\mathcal{G}^\bullet)$.
\end{enumerate}
(1) 与 (2) 是拉回映射；(3) 是相对杯积，见
Cohomology, Remark \ref{cohomology-remark-cup-product}；
(4) 比较导出与非导出张量积；(5) 由底层双复形上的显然映射
$p^{-1}\mathcal{F}^i \otimes_{f^{-1}\mathcal{O}_S} q^{-1}\mathcal{G}^j
\to \mathcal{F}^i \boxtimes \mathcal{G}^j$ 给出。

\medskip\noindent
令 $A = \Gamma(S, \mathcal{O}_S)$。存在“整体杯积”映射
\begin{equation}
\label{perfect-equation-de-rham-kunneth}
R\Gamma(X, \mathcal{F}^\bullet)
\otimes_A^\mathbf{L}
R\Gamma(Y, \mathcal{G}^\bullet)
\longrightarrow
R\Gamma(X \times_S Y,
\text{Tot}(\mathcal{F}^\bullet \boxtimes \mathcal{G}^\bullet))
\end{equation}
位于 $D(A)$ 中。它与上述相对杯积类似，由下列映射构造：
\begin{enumerate}
\item $R\Gamma(X, \mathcal{F}^\bullet) \to
R\Gamma(X \times_S Y, p^{-1}\mathcal{F}^\bullet)$
\item $R\Gamma(Y, \mathcal{G}^\bullet) \to
R\Gamma(X \times_S Y, q^{-1}\mathcal{G}^\bullet)$,
\item $R\Gamma(X \times_S Y, p^{-1}\mathcal{F}^\bullet)
\otimes_A^\mathbf{L} R\Gamma(X \times_S Y, q^{-1}\mathcal{G}^\bullet) \to
R\Gamma(X \times_S Y,
p^{-1}\mathcal{F}^\bullet \otimes_{f^{-1}\mathcal{O}_S}^\mathbf{L}
q^{-1}\mathcal{G}^\bullet)$,
\item
$p^{-1}\mathcal{F}^\bullet \otimes_{f^{-1}\mathcal{O}_S}^\mathbf{L}
q^{-1}\mathcal{G}^\bullet \to
\text{Tot}(p^{-1}\mathcal{F}^\bullet \otimes_{f^{-1}\mathcal{O}_S}
q^{-1}\mathcal{G}^\bullet)$
\item
$\text{Tot}(p^{-1}\mathcal{F}^\bullet \otimes_{f^{-1}\mathcal{O}_S}
q^{-1}\mathcal{G}^\bullet) \to \text{Tot}(\mathcal{F}^\bullet \boxtimes
\mathcal{G}^\bullet)$.
\end{enumerate}
这里 (1) 与 (2) 是拉回映射，(3) 是 Cohomology, Section
\ref{cohomology-section-cup-product} 中构造的杯积。(4) 与 (5) 如上一段所述。

\begin{lemma}
\label{perfect-lemma-kunneth-special}
在上述情形中，若下列假设成立，则杯积 (\ref{perfect-equation-de-rham-kunneth})
是 $D(A)$ 中的同构：
\begin{enumerate}
\item $S = \Spec(A)$ 仿射；
\item $X$ 与 $Y$ 拟紧且具有仿射对角态射；
\item $\mathcal{F}^\bullet$ 有界；
\item $\mathcal{G}^\bullet$ 下有界；
\item $\mathcal{F}^n$ 在 $S$ 上平坦；
\item $\mathcal{G}^m$ 在 $S$ 上平坦。
\end{enumerate}
\end{lemma}

\begin{proof}
我们使用 Morphisms, Remark
\ref{morphisms-remark-base-change-differential-operators} 中引入的记号
$\mathcal{A}_{X/S}$ 与 $\mathcal{A}_{Y/S}$。假设在范畴 $\mathcal{A}_{X/S}$
中有复形映射
$$
\mathcal{F}_1^\bullet \to
\mathcal{F}_2^\bullet \to
\mathcal{F}_3^\bullet \to
\mathcal{F}_1^\bullet[1]
$$
。由杯积的函子性，得到可交换图
$$
\xymatrix{
R\Gamma(X, \mathcal{F}_1^\bullet)
\otimes_A^\mathbf{L}
R\Gamma(Y, \mathcal{G}^\bullet)
\ar[r] \ar[d] &
R\Gamma(X \times_S Y,
\text{Tot}(\mathcal{F}_1^\bullet \boxtimes \mathcal{G}^\bullet)) \ar[d] \\
R\Gamma(X, \mathcal{F}_2^\bullet)
\otimes_A^\mathbf{L}
R\Gamma(Y, \mathcal{G}^\bullet)
\ar[r] \ar[d] &
R\Gamma(X \times_S Y,
\text{Tot}(\mathcal{F}_2^\bullet \boxtimes \mathcal{G}^\bullet)) \ar[d] \\
R\Gamma(X, \mathcal{F}_3^\bullet)
\otimes_A^\mathbf{L}
R\Gamma(Y, \mathcal{G}^\bullet)
\ar[r] \ar[d] &
R\Gamma(X \times_S Y,
\text{Tot}(\mathcal{F}_3^\bullet \boxtimes \mathcal{G}^\bullet)) \ar[d] \\
R\Gamma(X, \mathcal{F}_1^\bullet[1])
\otimes_A^\mathbf{L}
R\Gamma(Y, \mathcal{G}^\bullet)
\ar[r] &
R\Gamma(X \times_S Y,
\text{Tot}(\mathcal{F}_1^\bullet[1] \boxtimes \mathcal{G}^\bullet))
}
$$
若原映射在 $\mathcal{A}_{X/S}$ 的同伦范畴中组成可区别三角，
则此图的各列在 $D(A)$ 中组成可区别三角。

\medskip\noindent
在本引理的情形中，假设当 $n < i$ 时 $\mathcal{F}^n = 0$。可以考虑逐项分裂的
复形短正合列
$$
0 \to \sigma_{\geq i + 1}\mathcal{F}^\bullet \to
\mathcal{F}^\bullet \to \mathcal{F}^i[-i] \to 0
$$
，其中截断如 Homology, Section \ref{homology-section-truncations} 所述。
这在 $\mathcal{A}_{X/S}$ 的同伦范畴中给出可区别三角
$$
\sigma_{\geq i + 1}\mathcal{F}^\bullet \to
\mathcal{F}^\bullet \to
\mathcal{F}^i[-i] \to
(\sigma_{\geq i + 1}\mathcal{F}^\bullet)[1]
$$
，其中最后一个箭头由边界映射 $\mathcal{F}^i \to \mathcal{F}^{i + 1}$ 给出。
由上述讨论，只需对 $\mathcal{F}^i[-i]$ 与
$\sigma_{\geq i + 1}\mathcal{F}^\bullet$ 证明本引理。由于
$\sigma_{\geq i + 1}\mathcal{F}^\bullet$ 的非零项更少，归纳可知，
只要能在 $\mathcal{F}^\bullet$ 仅于单个次数非零时证明本引理，结论就成立。
因此可假设 $\mathcal{F}^\bullet$ 只在一个次数上非零。

\medskip\noindent
假设 $\mathcal{F}^\bullet$ 是如下复形：拟凝聚 $\mathcal{O}_X$-模
$\mathcal{F}$ 在 $S$ 上平坦并位于次数 $0$，其他次数均为零。例如由
Lemma \ref{perfect-lemma-cohomology-de-rham-base-change}，
$R\Gamma(X, \mathcal{F})$ 具有有限 Tor 维数。设其 Tor 振幅在 $[i, j]$ 中。
取 $N \gg 0$，并在 $\mathcal{A}_{Y/S}$ 的同伦范畴中考虑可区别三角
$$
\sigma_{\geq N + 1}\mathcal{G}^\bullet \to
\mathcal{G}^\bullet \to
\sigma_{\leq N}\mathcal{G}^\bullet \to
(\sigma_{\geq N + 1}\mathcal{G}^\bullet)[1]
$$
。注意下列二者
$$
R\Gamma(X, \mathcal{F})
\otimes_A^\mathbf{L}
R\Gamma(Y, \sigma_{\geq N + 1}\mathcal{G}^\bullet)
\quad\text{且}\quad
R\Gamma(X \times_S Y,
\text{Tot}(\mathcal{F} \boxtimes \sigma_{\geq N + 1}\mathcal{G}^\bullet))
$$
在次数 $\leq N + i$ 上的上同调均消失。因此利用上述论证，若要在某个给定次数上
证明陈述，可以假设 $\mathcal{G}^\bullet$ 有界。再重复一次上述论证，还可假设
$\mathcal{G}^\bullet$ 只在一个次数上非零。此情形由
Lemma \ref{perfect-lemma-kunneth-single-sheaf} 处理。
\end{proof}







\section{Ext 的 K\"unneth 公式}
\label{perfect-section-kunneth-Ext}

\noindent
考虑概形的笛卡尔图
$$
\xymatrix{
& X \times_S Y \ar[ld]^p \ar[rd]_q \ar[dd]^f \\
X \ar[rd]_a & & Y \ar[ld]^b \\
& S
}
$$
本节中，对 $K \in D(\mathcal{O}_X)$ 与 $M \in D(\mathcal{O}_Y)$ 定义
$$
K \boxtimes M =
Lp^*K \otimes_{\mathcal{O}_{X \times_S Y}}^\mathbf{L} Lq^*M
$$
。我们断言，对 $K, K' \in D(\mathcal{O}_X)$ 与
$M, M' \in D(\mathcal{O}_Y)$，存在典范映射
\begin{equation}
\label{perfect-equation-kunneth-ext}
Ra_*R\SheafHom(K, K')
\otimes_{\mathcal{O}_S}^\mathbf{L}
Rb_*R\SheafHom(M, M')
\longrightarrow
Rf_*(R\SheafHom(K \boxtimes M, K' \boxtimes M'))
\end{equation}
。事实上，可取下述映射的伴随映射：
$$
\begin{matrix}
Lf^*\left(Ra_*R\SheafHom(K, K')
\otimes_{\mathcal{O}_S}^\mathbf{L}
Rb_* R\SheafHom(M, M')\right) = \\
Lf^* Ra_* R\SheafHom(K, K')
\otimes_{\mathcal{O}_{X \times_S Y}}^\mathbf{L}
Lf^* Rb_* R\SheafHom(M, M') = \\
Lp^* La^* Ra_* R\SheafHom(K, K')
\otimes_{\mathcal{O}_{X \times_S Y}}^\mathbf{L}
Lq^* Lb^* Rb_* R\SheafHom(M, M') \to \\
Lp^* R\SheafHom(K, K')
\otimes_{\mathcal{O}_{X \times_S Y}}^\mathbf{L}
Lq^* R\SheafHom(M, M') \to \\
R\SheafHom(Lp^*K, Lp^*K')
\otimes_{\mathcal{O}_{X \times_S Y}}^\mathbf{L}
R\SheafHom(Lq^*M, Lq^*M') \to \\
R\SheafHom(K \boxtimes M, K' \boxtimes M')
\end{matrix}
$$
这里第一个等式是拉回与张量积的相容性，见
Cohomology, Lemma \ref{cohomology-lemma-pullback-tensor-product}。
第二个等式使用 $f = a \circ p = b \circ q$ 以及拉回的复合，见
Cohomology, Lemma \ref{cohomology-lemma-derived-pullback-composition}。
第一个箭头由伴随映射
$La^* Ra_* \to \text{id}$
与 $Lb^* Rb_* \to \text{id}$ 给出，因为推前与拉回互为伴随，见
Cohomology, Lemma \ref{cohomology-lemma-adjoint}。第二个箭头由
Cohomology, Remark \ref{cohomology-remark-prepare-fancy-base-change} 给出。
第三个亦即最后一个箭头由
Cohomology, Remark \ref{cohomology-remark-tensor-internal-hom} 给出。
一个简单的特殊情形是下述结果。

\begin{lemma}
\label{perfect-lemma-kunneth-Ext}
在上述情形中，假设 $a$ 与 $b$ 拟紧拟分离，且 $X$ 与 $Y$ 在 $S$ 上 Tor 独立。
若 $K$ 完美、$K' \in D_\QCoh(\mathcal{O}_X)$、$M$ 完美且
$M' \in D_\QCoh(\mathcal{O}_Y)$，则 (\ref{perfect-equation-kunneth-ext}) 是同构。
\end{lemma}

\begin{proof}
在此情形中，有 $R\SheafHom(K, K') = K' \otimes^\mathbf{L} K^\vee$、
$R\SheafHom(M, M') = M' \otimes^\mathbf{L} M^\vee$，且
$$
R\SheafHom(K \boxtimes M, K' \boxtimes M') =
(K' \otimes^\mathbf{L} K^\vee) \boxtimes
(M' \otimes^\mathbf{L} M^\vee)
$$
。见 Cohomology, Lemma \ref{cohomology-lemma-dual-perfect-complex}；
这里还用到拉回与张量积保持完美性。因此该情形由
Lemma \ref{perfect-lemma-kunneth} 推出。（略去验证在这些识别下得到的是同一映射。）
\end{proof}








\section{上同调与基变换，V}
\label{perfect-section-cohomology-base-change}

\noindent
在 Section \ref{perfect-section-cohomology-and-base-change-perfect} 中，
我们看到，当态射 Tor 独立时，一个基变换定理成立。即使在仿射情形中，
没有这样的条件也不可能有基变换定理，见 More on Algebra, Section
\ref{more-algebra-section-tor-independence}。本节分析何时可以“一次对一个复形”
得到基变换结果。

\medskip\noindent
为实现这一点，假设有概形的可交换图
$$
\xymatrix{
X' \ar[r]_{g'} \ar[d]_{f'} &
X \ar[d]^f \\
S' \ar[r]^g &
S
}
$$
（通常将假设它是笛卡尔图）。设 $K \in D_\QCoh(\mathcal{O}_X)$，且设
$L(g')^*K \to K'$ 为 $D_\QCoh(\mathcal{O}_{X'})$ 中的映射。
对点 $x' \in X'$，令 $x = g'(x') \in X$、$s' = f'(x') \in S'$ 及
$s = f(x) = g(s')$。于是可考虑映射
$$
K_x \otimes_{\mathcal{O}_{S, s}}^\mathbf{L} \mathcal{O}_{S', s'} \to
K_x \otimes_{\mathcal{O}_{X, x}}^\mathbf{L} \mathcal{O}_{X', x'} \to
K'_{x'}
$$
，其中第一个箭头是 More on Algebra, Equation
(\ref{more-algebra-equation-comparison-map})，第二个来自
$(L(g')^*K)_{x'} =
K_x \otimes_{\mathcal{O}_{X, x}}^\mathbf{L} \mathcal{O}_{X', x'}$
及给定映射 $L(g')^*K \to K'$。对每个 $i \in \mathbf{Z}$，
$H^i(K_x \otimes_{\mathcal{O}_{S, s}}^\mathbf{L} \mathcal{O}_{S', s'})$.
上得到一个
$\mathcal{O}_{X, x} \otimes_{\mathcal{O}_{S, s}} \mathcal{O}_{S', s'}$-模结构。
综合这些数据，得到典范映射
\begin{equation}
\label{perfect-equation-bc}
H^i(K_x \otimes_{\mathcal{O}_{S, s}}^\mathbf{L} \mathcal{O}_{S', s'})
\otimes_{(\mathcal{O}_{X, x} \otimes_{\mathcal{O}_{S, s}} \mathcal{O}_{S', s'})}
\mathcal{O}_{X', x'}
\longrightarrow
H^i(K'_{x'})
\end{equation}
，它是 $\mathcal{O}_{X', x'}$-模映射。

\begin{lemma}
\label{perfect-lemma-single-complex-base-change-condition}
设
$$
\xymatrix{
X' \ar[r]_{g'} \ar[d]_{f'} &
X \ar[d]^f \\
S' \ar[r]^g &
S
}
$$
是概形的笛卡尔图。设 $K \in D_\QCoh(\mathcal{O}_X)$，且设
$L(g')^*K \to K'$ 为 $D_\QCoh(\mathcal{O}_{X'})$ 中的映射。下列条件等价：
\begin{enumerate}
\item 对任意 $x' \in X'$ 与 $i \in \mathbf{Z}$，映射
(\ref{perfect-equation-bc}) 是同构；
\item 对仿射开集 $U \subset X$ 与 $V' \subset S'$，若二者都映入仿射开集
$V \subset S$，且 $U' = V' \times_V U$，则复合
$$
R\Gamma(U, K) \otimes_{\mathcal{O}_S(U)}^\mathbf{L} \mathcal{O}_{S'}(V')
\to
R\Gamma(U, K) \otimes_{\mathcal{O}_X(U)}^\mathbf{L} \mathcal{O}_{X'}(U')
\to
R\Gamma(U', K')
$$
是 $D(\mathcal{O}_{S'}(V'))$ 中的同构；
\item 存在集合 $I$，以及如 (2) 中的四元组 $U_i, V_i', V_i, U_i'$，$i \in I$，
使得 $X' = \bigcup U'_i$。
\end{enumerate}
\end{lemma}

\begin{proof}
(2) 中第二个箭头来自
$$
R\Gamma(U, K) \otimes_{\mathcal{O}_X(U)}^\mathbf{L} \mathcal{O}_{X'}(U') =
R\Gamma(U', L(g')^*K)
$$
这一 Lemma \ref{perfect-lemma-quasi-coherence-pullback} 中的等式及给定箭头
$L(g')^*K \to K'$。(2) 中第一个箭头是 More on Algebra, Equation
(\ref{more-algebra-equation-comparison-map})。显然 (2) 蕴含 (3)。注意 (1)
在 $X'$ 上是局部的。因此只需证明：若 $X$、$S$、$S'$、$X'$ 均仿射，
则 (1) 等价于下述条件：
$$
R\Gamma(X, K) \otimes_{\mathcal{O}_S(S)}^\mathbf{L} \mathcal{O}_{S'}(S')
\to
R\Gamma(X, K) \otimes_{\mathcal{O}_X(X)}^\mathbf{L} \mathcal{O}_{X'}(X')
\to
R\Gamma(X', K')
$$
是 $D(\mathcal{O}_{S'}(S'))$ 中的同构。记
$S = \Spec(R)$、$X = \Spec(A)$、$S' = \Spec(R')$、$X' = \Spec(A')$；
$K$ 对应于 $A$-模复形 $M^\bullet$，$K'$ 对应于 $A'$-模复形
$N^\bullet$。注意 $A' = A \otimes_R R'$。上述条件即复合
$$
M^\bullet \otimes_R^\mathbf{L} R' \to
M^\bullet \otimes_A^\mathbf{L} A' \to
N^\bullet
$$
是 $D(R')$ 中的同构。等价地，对所有 $i \in \mathbf{Z}$，映射
$$
H^i(M^\bullet \otimes_R^\mathbf{L} R') \to
H^i(M^\bullet \otimes_A^\mathbf{L} A') \to
H^i(N^\bullet)
$$
是同构。注意这是 $A \otimes_R R'$-模映射，即 $A'$-模映射。
另一方面，(1) 要求对相容素理想
$\mathfrak q' \subset A'$, $\mathfrak q \subset A$,
$\mathfrak p' \subset R'$, $\mathfrak p \subset R$
，复合
$$
H^i(M^\bullet_\mathfrak q \otimes_{R_\mathfrak p}^\mathbf{L} R'_{\mathfrak p'})
\otimes_{(A_\mathfrak q \otimes_{R_\mathfrak p} R'_{\mathfrak p'})}
A'_{\mathfrak q'} \to
H^i(M^\bullet_{\mathfrak q}
\otimes_{A_\mathfrak q}^\mathbf{L} A'_{\mathfrak q'})
\to H^i(N^\bullet_{\mathfrak q'})
$$
是同构。由于
$$
H^i(M^\bullet_\mathfrak q \otimes_{R_\mathfrak p}^\mathbf{L} R'_{\mathfrak p'})
\otimes_{(A_\mathfrak q \otimes_{R_\mathfrak p} R'_{\mathfrak p'})}
A'_{\mathfrak q'} =
H^i(M^\bullet \otimes_R^\mathbf{L} R') \otimes_{A'} A'_{\mathfrak q'}
$$
是在 $\mathfrak q'$ 处的局部化，由 Algebra, Lemma
\ref{algebra-lemma-characterize-zero-local} 可知这两个条件等价。
\end{proof}

\begin{lemma}
\label{perfect-lemma-single-complex-base-change}
设
$$
\xymatrix{
X' \ar[r]_{g'} \ar[d]_{f'} &
X \ar[d]^f \\
S' \ar[r]^g &
S
}
$$
是概形的笛卡尔图。设 $K \in D_\QCoh(\mathcal{O}_X)$，且设
$L(g')^*K \to K'$ 为 $D_\QCoh(\mathcal{O}_{X'})$ 中的映射。若
\begin{enumerate}
\item Lemma \ref{perfect-lemma-single-complex-base-change-condition} 的等价条件成立；且
\item $f$ 拟紧拟分离，
\end{enumerate}
则复合 $Lg^*Rf_*K \to Rf'_*L(g')^*K \to Rf'_*K'$ 是同构。
\end{lemma}

\begin{proof}
可以使用 Lemma \ref{perfect-lemma-compare-base-change} 证明中的同一方法来证明，
但我们改用归纳原理与相对 Mayer--Vietoris 来证明。

\medskip\noindent
为检验该映射是同构，可以在 $S'$ 上局部地工作。因此可假设
$g : S' \to S$ 是仿射概形之间的态射。特别地，$X$ 是拟紧拟分离概形。
我们使用
Cohomology of Schemes, Lemma \ref{coherent-lemma-induction-principle}
的归纳原理证明：对任意拟紧开集 $U \subset X$，类似构造的映射
$Lg^*R(U \to S)_*K|_U \to R(U' \to S')_*K'|_{U'}$ 是同构；其中
$U' = (g')^{-1}(U)$。

\medskip\noindent
若 $U \subset X$ 为仿射开集，则由假设可知结论成立；见
Lemma \ref{perfect-lemma-single-complex-base-change-condition} 的 (2)，以及
Lemmas \ref{perfect-lemma-affine-compare-bounded} 与
\ref{perfect-lemma-quasi-coherence-pullback} 所给出的代数转译。

\medskip\noindent
归纳步骤。设 $X = U \cup V$ 是开覆盖，$U$、$V$、$U \cap V$ 均拟紧，
且结论对 $U$、$V$ 与 $U \cap V$ 成立。记 $a = f|_U$、$b = f|_V$、
$c = f|_{U \cap V}$。设 $a' : U' \to S'$、$b' : V' \to S'$ 与
$c' : U' \cap V' \to S'$ 分别为 $a$、$b$ 与 $c$ 的基变换。
使用相对 Mayer--Vietoris 的可区别三角
（Cohomology, Lemma \ref{cohomology-lemma-unbounded-relative-mayer-vietoris}），
得到可交换图
$$
\xymatrix{
Lg^*Rf_*K \ar[r] \ar[d] &
Rf'_* K' \ar[d] \\
Lg^*Ra_* K|_U \oplus
Lg^*Rb_* K|_V \ar[r] \ar[d] &
Ra'_* K'|_{U'} \oplus
Rb'_* K'|_{V'} \ar[d] \\
Lg^*Rc_* K|_{U \cap V} \ar[r] \ar[d] &
Rc'_* K'|_{U' \cap V'} \ar[d] \\
Lg^*Rf_* K[1] \ar[r] &
Rf'_* K'[1]
}
$$
由于第二个与第三个水平箭头是同构，第一个也是同构
（Derived Categories, Lemma \ref{derived-lemma-third-isomorphism-triangle}）。
引理得证。
\end{proof}

\begin{lemma}
\label{perfect-lemma-single-complex-base-change-condition-inherited}
设
$$
\xymatrix{
X' \ar[r]_{g'} \ar[d]_{f'} &
X \ar[d]^f \\
S' \ar[r]^g &
S
}
$$
是概形的笛卡尔图。设 $K \in D_\QCoh(\mathcal{O}_X)$，且设
$L(g')^*K \to K'$ 为 $D_\QCoh(\mathcal{O}_{X'})$ 中的映射。若
Lemma \ref{perfect-lemma-single-complex-base-change-condition} 的等价条件成立，则：
\begin{enumerate}
\item 对 $E \in D_\QCoh(\mathcal{O}_X)$，映射
$L(g')^*(E \otimes^\mathbf{L} K) \to L(g')^*E \otimes^\mathbf{L} K'$
满足 Lemma \ref{perfect-lemma-single-complex-base-change-condition} 的等价条件；
\item 若 $D(\mathcal{O}_X)$ 中的 $E$ 完美，则映射
$L(g')^*R\SheafHom(E, K) \to R\SheafHom(L(g')^*E, K')$
满足 Lemma \ref{perfect-lemma-single-complex-base-change-condition} 的等价条件；
\item 若 $K$ 下有界，且 $D(\mathcal{O}_X)$ 中的 $E$ 伪凝聚，则映射
$L(g')^*R\SheafHom(E, K) \to R\SheafHom(L(g')^*E, K')$
满足 Lemma \ref{perfect-lemma-single-complex-base-change-condition} 的等价条件。
\end{enumerate}
\end{lemma}

\begin{proof}
由 Lemmas \ref{perfect-lemma-quasi-coherence-pullback}、
\ref{perfect-lemma-quasi-coherence-tensor-product}、
\ref{perfect-lemma-quasi-coherence-internal-hom} 及 Cohomology, Lemmas
\ref{cohomology-lemma-pseudo-coherent-pullback}、
\ref{cohomology-lemma-perfect-pullback}，所涉及复形的上同调层拟凝聚，
所以该陈述有意义。在情形 (1) 中，利用张量积的传递性计算
(\ref{perfect-equation-bc}) 的源和靶，即可检验映射 (\ref{perfect-equation-bc}) 是同构，见
More on Algebra, Lemma \ref{more-algebra-lemma-triple-tensor-product}。
(2) 与 (3) 中的映射是复合
$$
L(g')^*R\SheafHom(E, K) \to R\SheafHom(L(g')^*E, L(g')^*K)
\to R\SheafHom(L(g')^*E, K')
$$
，其中第一个箭头是 Cohomology, Remark
\ref{cohomology-remark-prepare-fancy-base-change}，第二个箭头来自给定映射
$L(g')^*K \to K'$。为证明映射 (\ref{perfect-equation-bc}) 是同构，在情形 (2) 中
用有限投射 $\mathcal{O}_{X. x}$-模的有界复形表示 $E_x$，在情形 (3) 中
用有限自由模的有上界复形表示它，然后计算箭头的源和靶。略去若干细节。
\end{proof}

\begin{lemma}
\label{perfect-lemma-base-change-tensor}
设 $f : X \to S$ 为概形的拟紧拟分离态射。设
$E \in D_\QCoh(\mathcal{O}_X)$。设 $\mathcal{G}^\bullet$ 为在 $S$ 上平坦的
拟凝聚 $\mathcal{O}_X$-模的有上界复形。则构造
$$
Rf_*(E \otimes^\mathbf{L}_{\mathcal{O}_X} \mathcal{G}^\bullet)
$$
与任意基变换可交换（精确陈述见证明）。
\end{lemma}

\begin{proof}
该陈述意指如下。设 $g : S' \to S$ 为概形态射，并考虑基变换图
$$
\xymatrix{
X' \ar[r]_{g'} \ar[d]_{f'} &
X \ar[d]^f \\
S' \ar[r]^g &
S
}
$$
，换言之，$X' = S' \times_S X$。本引理断言
$$
Lg^*Rf_*(E \otimes^\mathbf{L}_{\mathcal{O}_X} \mathcal{G}^\bullet)
\longrightarrow
Rf'_*\left(
L(g')^*E \otimes^\mathbf{L}_{\mathcal{O}_{X'}} (g')^*\mathcal{G}^\bullet
\right)
$$
是同构。注意右端对 $\mathcal{G}^\bullet$ 并{\bf 未}使用导出拉回。
为证明这一点，应用 Lemmas \ref{perfect-lemma-single-complex-base-change} 与
\ref{perfect-lemma-single-complex-base-change-condition-inherited}，可知只需证明典范映射
$$
L(g')^*\mathcal{G}^\bullet \to (g')^*\mathcal{G}^\bullet
$$
满足 Lemma \ref{perfect-lemma-single-complex-base-change-condition} 的等价条件。
在茎上检验该条件即可；此时它立即由如下事实推出：按我们对复形
$\mathcal{G}^\bullet$ 的假设，
$\mathcal{G}^\bullet_x \otimes_{\mathcal{O}_{S, s}} \mathcal{O}_{S', s'}$
计算导出张量积。
\end{proof}

\begin{lemma}
\label{perfect-lemma-base-change-RHom}
设 $f : X \to S$ 为概形的拟紧拟分离态射。设 $E$ 为
$D(\mathcal{O}_X)$ 的对象，并设 $\mathcal{G}^\bullet$ 为拟凝聚
$\mathcal{O}_X$-模的复形。若
\begin{enumerate}
\item $E$ 完美，$\mathcal{G}^\bullet$ 有上界，且 $\mathcal{G}^n$ 在 $S$ 上平坦；或
\item $E$ 伪凝聚，$\mathcal{G}^\bullet$ 有界，且 $\mathcal{G}^n$ 在 $S$ 上平坦，
\end{enumerate}
则构造
$$
Rf_*R\SheafHom(E, \mathcal{G}^\bullet)
$$
与任意基变换可交换（精确陈述见证明）。
\end{lemma}

\begin{proof}
该陈述意指如下。设 $g : S' \to S$ 为概形态射，并考虑基变换图
$$
\xymatrix{
X' \ar[r]_{g'} \ar[d]_{f'} &
X \ar[d]^f \\
S' \ar[r]^g &
S
}
$$
，换言之，$X' = S' \times_S X$。本引理断言
$$
Lg^*Rf_*R\SheafHom(E, \mathcal{G}^\bullet)
\longrightarrow
R(f')_*R\SheafHom(L(g')^*E, (g')^*\mathcal{G}^\bullet)
$$
是同构。注意右端对 $\mathcal{G}^\bullet$ 并{\bf 未}使用导出拉回。为证明这一点，
应用 Lemmas \ref{perfect-lemma-single-complex-base-change} 与
\ref{perfect-lemma-single-complex-base-change-condition-inherited}，可知只需证明典范映射
$$
L(g')^*\mathcal{G}^\bullet \to (g')^*\mathcal{G}^\bullet
$$
满足 Lemma \ref{perfect-lemma-single-complex-base-change-condition} 的等价条件。
这已在 Lemma \ref{perfect-lemma-base-change-tensor} 的证明中说明。
\end{proof}





\section{构造完美复形}
\label{perfect-section-producing-perfect}

\noindent
下述引理是构造完美复形的主要技术工具。此结果的后续版本将通过 Noether 逼近
归结到这里，见 Section \ref{perfect-section-cohomology-and-base-change-final}。

\begin{lemma}
\label{perfect-lemma-perfect-direct-image}
设 $S$ 为 Noether 概形。设 $f : X \to S$ 为局部有限型概形态射。
设 $E \in D(\mathcal{O}_X)$，并满足：
\begin{enumerate}
\item $E \in D^b_{\textit{Coh}}(\mathcal{O}_X)$；
\item 对所有 $i$，$H^i(E)$ 的支集在 $S$ 上固有；
\item $E$ 作为 $D(f^{-1}\mathcal{O}_S)$ 的对象具有有限 Tor 维数。
\end{enumerate}
则 $Rf_*E$ 是 $D(\mathcal{O}_S)$ 的完美对象。
\end{lemma}

\begin{proof}
由 Lemma \ref{perfect-lemma-direct-image-coherent}，$Rf_*E$ 是
$D^b_{\textit{Coh}}(\mathcal{O}_S)$ 的对象。因此 $Rf_*E$ 伪凝聚
（Lemma \ref{perfect-lemma-identify-pseudo-coherent-noetherian}）。所以只需证明
$Rf_*E$ 具有有限 Tor 维数，见 Cohomology, Lemma
\ref{cohomology-lemma-perfect}。由 Lemma \ref{perfect-lemma-tor-qc-qs}，
只需检验对 $S$ 上所有拟凝聚层 $\mathcal{F}$，
$Rf_*(E) \otimes_{\mathcal{O}_S}^\mathbf{L} \mathcal{F}$
都具有一致有界的上同调。上有界性显然，因为 $Rf_*(E)$ 上有界。
令 $T \subset X$ 为所有 $i$ 的 $H^i(E)$ 的支集之并。由假设 (1) 与 (2)，
$T$ 在 $S$ 上固有，见 Cohomology of Schemes, Lemma
\ref{coherent-lemma-union-closed-proper-over-base}。特别地，存在包含 $T$ 的拟紧开集
$X' \subset X$。令 $f' = f|_{X'}$；由于 $E$ 限制到 $X \setminus T$ 上为零，
有 $Rf_*(E) = Rf'_*(E|_{X'})$。所以可用 $X'$ 替换 $X$ 并假设 $f$ 拟紧。
此外，由 Morphisms, Lemma
\ref{morphisms-lemma-finite-type-Noetherian-quasi-separated}，$f$ 拟分离。现在
$$
Rf_*(E) \otimes_{\mathcal{O}_S}^\mathbf{L} \mathcal{F} =
Rf_*\left(E \otimes_{\mathcal{O}_X}^\mathbf{L} Lf^*\mathcal{F}\right) =
Rf_*\left(E \otimes_{f^{-1}\mathcal{O}_S}^\mathbf{L} f^{-1}\mathcal{F}\right)
$$
，这是由 Lemma \ref{perfect-lemma-cohomology-base-change} 及
Cohomology, Lemma \ref{cohomology-lemma-variant-derived-pullback} 得到的。
由假设 (3)，复形
$E \otimes_{f^{-1}\mathcal{O}_S}^\mathbf{L} f^{-1}\mathcal{F}$
的上同调层位于某个给定有限区间内，记为 $[a, b]$。于是对它取 $Rf_*$ 后，
上同调位于区间 $[a, \infty)$ 内，结论得证。
\end{proof}

\begin{lemma}
\label{perfect-lemma-tensor-perfect}
设 $S$ 为 Noether 概形。设 $f : X \to S$ 为局部有限型概形态射。
设 $E \in D(\mathcal{O}_X)$ 完美。设 $\mathcal{G}^\bullet$ 为凝聚
$\mathcal{O}_X$-模的有界复形，它在 $S$ 上平坦且支集在 $S$ 上固有。
则 $K = Rf_*(E \otimes_{\mathcal{O}_X}^\mathbf{L} \mathcal{G}^\bullet)$
是 $D(\mathcal{O}_S)$ 的完美对象。
\end{lemma}

\begin{proof}
由 Lemma \ref{perfect-lemma-perfect-direct-image}，对象 $K$ 完美。我们检验该引理适用：
局部地，$E$ 同构于有限自由 $\mathcal{O}_X$-模的有限复形。因此局部地，
$E \otimes^\mathbf{L}_{\mathcal{O}_X} \mathcal{G}^\bullet$ 同构于一个有限复形，
其各项形如
$$
\bigoplus\nolimits_{i = a, \ldots, b} (\mathcal{G}^i)^{\oplus r_i}
$$
，其中 $a, b, r_a, \ldots, r_b$ 是某些整数。这立即蕴含上同调层
$H^i(E \otimes^\mathbf{L}_{\mathcal{O}_X} \mathcal{G})$ 凝聚。
Tor 维数的假设也由 $\mathcal{G}^i$ 在 $f^{-1}\mathcal{O}_S$ 上平坦推出。
\end{proof}

\begin{lemma}
\label{perfect-lemma-ext-perfect}
设 $S$ 为 Noether 概形。设 $f : X \to S$ 为局部有限型概形态射。
设 $E \in D(\mathcal{O}_X)$ 完美。设 $\mathcal{G}^\bullet$ 为凝聚
$\mathcal{O}_X$-模的有界复形，它在 $S$ 上平坦且支集在 $S$ 上固有。
则 $K = Rf_*R\SheafHom(E, \mathcal{G}^\bullet)$ 是
$D(\mathcal{O}_S)$ 的完美对象。
\end{lemma}

\begin{proof}
由于 $E$ 是完美复形，存在对偶完美复形 $E^\vee$，见
Cohomology, Lemma \ref{cohomology-lemma-dual-perfect-complex}。注意
$R\SheafHom(E, \mathcal{G}^\bullet) =
E^\vee \otimes^\mathbf{L}_{\mathcal{O}_X} \mathcal{G}^\bullet$。
因此 $K$ 的完美性由 Lemma \ref{perfect-lemma-tensor-perfect} 推出。
\end{proof}

\noindent
我们将在 Lemma \ref{perfect-lemma-flat-proper-perfect-direct-image-general}
中把下述引理推广到一般基底上的平坦固有态射，并在
More on Morphisms, Lemma
\ref{more-morphisms-lemma-perfect-proper-perfect-direct-image}
中推广到完美固有态射。

\begin{lemma}
\label{perfect-lemma-flat-proper-perfect-direct-image}
设 $S$ 为 Noether 概形。设 $f : X \to S$ 为平坦固有概形态射。
设 $E \in D(\mathcal{O}_X)$ 完美。则 $Rf_*E$ 是
$D(\mathcal{O}_S)$ 的完美对象。
\end{lemma}

\begin{proof}
我们断言 Lemma \ref{perfect-lemma-perfect-direct-image} 适用。条件 (1) 与 (2) 是立即的。
条件 (3) 在 $X$ 上是局部的。因此可假设 $X$ 与 $S$ 仿射，且 $E$
由 $\mathcal{O}_X$-模的严格完美复形表示。由于 $\mathcal{O}_X$ 作为
$f^{-1}\mathcal{O}_S$-模层是平坦的，条件 (3) 成立。
\end{proof}






\section{Ext 的投影公式}
\label{perfect-section-ext}

\noindent
Lemma \ref{perfect-lemma-compute-ext}（或文献中的类似结果）有时用于验证 Quot 函子、
Hilbert 概形及其他模问题的某个 Artin 判据。假设 $f : X \to S$
是固有、平坦、有限表示的概形态射，且 $E \in D(\mathcal{O}_X)$ 完美。
此时该引理断言
$$
\Ext^i_X(E, f^*\mathcal{F}) =
\Ext^i_S((Rf_*E^\vee)^\vee, \mathcal{F})
$$
对 $S$ 上拟凝聚的 $\mathcal{F}$ 成立。这样写使它看起来像 Ext 的投影公式；
事实上，该结果确实相当容易地由 Lemma \ref{perfect-lemma-cohomology-base-change} 推出。

\begin{lemma}
\label{perfect-lemma-compute-tensor-perfect}
采用 Lemma \ref{perfect-lemma-tensor-perfect} 的假设与记号。则对 $S$ 上拟凝聚的
$\mathcal{F}$，存在与边界映射相容的函子性同构（见证明）
$$
H^i(S, K \otimes^\mathbf{L}_{\mathcal{O}_S} \mathcal{F})
\longrightarrow
H^i(X, E \otimes_{\mathcal{O}_X}^\mathbf{L}
(\mathcal{G}^\bullet \otimes_{\mathcal{O}_X} f^*\mathcal{F}))
$$
。
\end{lemma}

\begin{proof}
有
$$
\mathcal{G}^\bullet \otimes_{\mathcal{O}_X}^\mathbf{L} Lf^*\mathcal{F} =
\mathcal{G}^\bullet \otimes_{f^{-1}\mathcal{O}_S}^\mathbf{L} f^{-1}\mathcal{F} =
\mathcal{G}^\bullet \otimes_{f^{-1}\mathcal{O}_S} f^{-1}\mathcal{F} =
\mathcal{G}^\bullet \otimes_{\mathcal{O}_X} f^*\mathcal{F}
$$
；第一个等式由 Cohomology, Lemma
\ref{cohomology-lemma-variant-derived-pullback} 给出；第二个等式是因为
$\mathcal{G}^n$ 是平坦 $f^{-1}\mathcal{O}_S$-模；第三个等式来自拉回的定义。
因此得到
\begin{align*}
H^i(X, E \otimes^\mathbf{L}_{\mathcal{O}_X}
(\mathcal{G}^\bullet \otimes_{\mathcal{O}_X} f^*\mathcal{F}))
& =
H^i(X, E \otimes^\mathbf{L}_{\mathcal{O}_X} \mathcal{G}^\bullet
\otimes_{\mathcal{O}_X}^\mathbf{L} Lf^*\mathcal{F}) \\
& =
H^i(S,
Rf_*(E \otimes^\mathbf{L}_{\mathcal{O}_X} \mathcal{G}^\bullet
\otimes^\mathbf{L}_{\mathcal{O}_X} Lf^*\mathcal{F})) \\
& =
H^i(S, Rf_*(E \otimes^\mathbf{L}_{\mathcal{O}_X} \mathcal{G}^\bullet)
\otimes^\mathbf{L}_{\mathcal{O}_S} \mathcal{F}) \\
& =
H^i(S, K \otimes^\mathbf{L}_{\mathcal{O}_S} \mathcal{F}) 
\end{align*}
。第一个等式由上文给出；第二个等式由 Leray 给出
（Cohomology, Lemma \ref{cohomology-lemma-before-Leray}）；第三个等式由
Lemma \ref{perfect-lemma-cohomology-base-change} 给出。关于边界映射的陈述意指如下：
给定拟凝聚 $\mathcal{O}_S$-模的短正合列
$0 \to \mathcal{F}_1 \to \mathcal{F}_2 \to \mathcal{F}_3 \to 0$，
这些同构组成可交换图
$$
\xymatrix{
H^i(S, K \otimes^\mathbf{L}_{\mathcal{O}_S} \mathcal{F}_3)
\ar[r] \ar[d]_\delta &
H^i(X, E \otimes^\mathbf{L}_{\mathcal{O}_X}
(\mathcal{G}^\bullet \otimes_{\mathcal{O}_X} f^*\mathcal{F}_3))
\ar[d]^\delta \\
H^{i + 1}(S, K \otimes^\mathbf{L}_{\mathcal{O}_S} \mathcal{F}_1)
\ar[r] &
H^{i + 1}(X, E \otimes^\mathbf{L}_{\mathcal{O}_X}
(\mathcal{G}^\bullet \otimes_{\mathcal{O}_X} f^*\mathcal{F}_1))
}
$$
，其中边界映射来自可区别三角
$$
K \otimes^\mathbf{L}_{\mathcal{O}_S} \mathcal{F}_1 \to
K \otimes^\mathbf{L}_{\mathcal{O}_S} \mathcal{F}_2 \to
K \otimes^\mathbf{L}_{\mathcal{O}_S} \mathcal{F}_3 \to
K \otimes^\mathbf{L}_{\mathcal{O}_S} \mathcal{F}_1[1]
$$
以及与 $D(\mathcal{O}_X)$ 中下述短正合列相伴的可区别三角：
$$
0 \to
\mathcal{G}^\bullet \otimes_{\mathcal{O}_X} f^*\mathcal{F}_1 \to
\mathcal{G}^\bullet \otimes_{\mathcal{O}_X} f^*\mathcal{F}_2 \to
\mathcal{G}^\bullet \otimes_{\mathcal{O}_X} f^*\mathcal{F}_3 \to 0
$$
它是 $\mathcal{O}_X$-模复形的短正合列。该列正合是因为 $\mathcal{G}^n$
在 $S$ 上平坦。略去验证所显示图的可交换性。
\end{proof}

\begin{lemma}
\label{perfect-lemma-compute-ext-perfect}
采用 Lemma \ref{perfect-lemma-ext-perfect} 的假设与记号。则对 $S$ 上拟凝聚的
$\mathcal{F}$，存在与边界映射相容的函子性同构（见证明）
$$
H^i(S, K \otimes^\mathbf{L}_{\mathcal{O}_S} \mathcal{F})
\longrightarrow
\Ext^i_{\mathcal{O}_X}(E,
\mathcal{G}^\bullet \otimes_{\mathcal{O}_X} f^*\mathcal{F})
$$
。
\end{lemma}

\begin{proof}
如 Lemma \ref{perfect-lemma-ext-perfect} 的证明，令 $E^\vee$ 为对偶完美复形，并回忆
$K = Rf_*(E^\vee \otimes_{\mathcal{O}_X}^\mathbf{L} \mathcal{G}^\bullet)$.
又有
$$
\Ext^i_{\mathcal{O}_X}(E,
\mathcal{G}^\bullet \otimes_{\mathcal{O}_X} f^*\mathcal{F})
=
H^i(X, E^\vee \otimes^\mathbf{L}_{\mathcal{O}_X}
(\mathcal{G}^\bullet \otimes_{\mathcal{O}_X} f^*\mathcal{F}))
$$
，这是由 $E^\vee$ 的构造得到的，所以对 $E^\vee$ 与 $\mathcal{G}^\bullet$
应用 Lemma \ref{perfect-lemma-compute-tensor-perfect} 即得这些同构的存在性。
关于边界映射的陈述意指如下：给定短正合列
$0 \to \mathcal{F}_1 \to \mathcal{F}_2 \to \mathcal{F}_3 \to 0$，
这些同构组成可交换图
$$
\xymatrix{
H^i(S, K \otimes^\mathbf{L}_{\mathcal{O}_S} \mathcal{F}_3)
\ar[r] \ar[d]_\delta &
\Ext^i_{\mathcal{O}_X}(E,
\mathcal{G}^\bullet \otimes_{\mathcal{O}_X} f^*\mathcal{F}_3) \ar[d]^\delta \\
H^{i + 1}(S, K \otimes^\mathbf{L}_{\mathcal{O}_S} \mathcal{F}_1)
\ar[r] &
\Ext^{i + 1}_{\mathcal{O}_X}(E,
\mathcal{G}^\bullet \otimes_{\mathcal{O}_X} f^*\mathcal{F}_1)
}
$$
，其中边界映射来自可区别三角
$$
K \otimes^\mathbf{L}_{\mathcal{O}_S} \mathcal{F}_1 \to
K \otimes^\mathbf{L}_{\mathcal{O}_S} \mathcal{F}_2 \to
K \otimes^\mathbf{L}_{\mathcal{O}_S} \mathcal{F}_3 \to
K \otimes^\mathbf{L}_{\mathcal{O}_S} \mathcal{F}_1[1]
$$
以及与 $D(\mathcal{O}_X)$ 中下述短正合列相伴的可区别三角：
$$
0 \to
\mathcal{G}^\bullet \otimes_{\mathcal{O}_X} f^*\mathcal{F}_1 \to
\mathcal{G}^\bullet \otimes_{\mathcal{O}_X} f^*\mathcal{F}_2 \to
\mathcal{G}^\bullet \otimes_{\mathcal{O}_X} f^*\mathcal{F}_3 \to 0
$$
。这是复形的短正合列。该列正合是因为 $\mathcal{G}$ 在 $S$ 上平坦。
略去验证所显示图的可交换性。
\end{proof}

\begin{lemma}
\label{perfect-lemma-compute-ext}
设 $f : X \to S$ 为概形态射，$E \in D(\mathcal{O}_X)$，且
$\mathcal{G}^\bullet$ 为 $\mathcal{O}_X$-模复形。假设：
\begin{enumerate}
\item $S$ 是 Noether 概形；
\item $f$ 局部有限型；
\item $E \in D^-_{\textit{Coh}}(\mathcal{O}_X)$；
\item $\mathcal{G}^\bullet$ 是凝聚 $\mathcal{O}_X$-模的有界复形，
它在 $S$ 上平坦且支集在 $S$ 上固有。
\end{enumerate}
则下列两个陈述成立：
\begin{enumerate}
\item[(A)] 对每个 $m \in \mathbf{Z}$，存在 $D(\mathcal{O}_S)$ 的完美对象 $K$
及函子性映射
$$
\alpha^i_\mathcal{F} :
\Ext^i_{\mathcal{O}_X}(E,
\mathcal{G}^\bullet \otimes_{\mathcal{O}_X} f^*\mathcal{F})
\longrightarrow
H^i(S, K \otimes^\mathbf{L}_{\mathcal{O}_S} \mathcal{F})
$$
，对 $S$ 上拟凝聚的 $\mathcal{F}$ 与边界映射相容（见证明），
并且当 $i \leq m$ 时 $\alpha^i_\mathcal{F}$ 是同构；
\item[(B)] 存在伪凝聚对象 $L \in D(\mathcal{O}_S)$ 及函子性同构
$$
\Ext^i_{\mathcal{O}_S}(L, \mathcal{F}) \longrightarrow
\Ext^i_{\mathcal{O}_X}(E,
\mathcal{G}^\bullet \otimes_{\mathcal{O}_X} f^*\mathcal{F})
$$
，对 $S$ 上拟凝聚的 $\mathcal{F}$ 与边界映射相容。
\end{enumerate}
\end{lemma}

\begin{proof}
先证明 (A)。
设 $\mathcal{G}^i$ 仅在 $i \in [a, b]$ 时非零。
我们可以把 $X$ 替换为
$\mathcal{G}^i$ 的支集之并的一个拟紧开邻域。
因此可以假设 $X$ 是诺特的。
此时 $X$ 与 $f$ 都是拟紧且拟分离的。
选取由完美复形 $P$ 给出的
$(X, E, -m - 1 + a)$
的一个逼近 $P \to E$
（其存在性由定理 \ref{perfect-theorem-approximation} 保证）。
于是诱导映射
$$
\Ext^i_{\mathcal{O}_X}(E,
\mathcal{G}^\bullet \otimes_{\mathcal{O}_X} f^*\mathcal{F})
\longrightarrow
\Ext^i_{\mathcal{O}_X}(P,
\mathcal{G}^\bullet \otimes_{\mathcal{O}_X} f^*\mathcal{F})
$$
在 $i \leq m$ 时是同构。事实上，此映射的核（相应地，余核）
是下列群的一个商（相应地，子模）：
$$
\Ext^i_{\mathcal{O}_X}(C,
\mathcal{G}^\bullet \otimes_{\mathcal{O}_X} f^*\mathcal{F})
\quad\text{分别}\quad
\Ext^{i + 1}_{\mathcal{O}_X}(C,
\mathcal{G}^\bullet \otimes_{\mathcal{O}_X} f^*\mathcal{F})
$$
其中 $C$ 是 $P \to E$ 的锥。由于 $C$ 在次数
$\geq -m - 1 + a$ 的上同调层均为零，导出范畴，
引理 \ref{derived-lemma-negative-exts} 表明，当
$i \leq m + 1$ 时这些 $\Ext$ 群为零。
这就把问题归结到 $E$ 是完美复形的情形，
而这正是引理 \ref{perfect-lemma-compute-ext-perfect}。
关于边界映射的断言已在
引理 \ref{perfect-lemma-compute-ext-perfect} 的证明中说明。

\medskip\noindent
证明 (B)。与 (A) 的证明相同，可以假设 $X$ 是诺特的。
由引理 \ref{perfect-lemma-identify-pseudo-coherent-noetherian}，
$E$ 是伪凝聚的。
根据引理 \ref{perfect-lemma-pseudo-coherent-hocolim}，可写成
$E = \text{hocolim} E_n$，其中 $E_n$ 完美，并且 $E_n \to E$
在截断 $\tau_{\geq -n}$ 上诱导同构。令 $E_n^\vee$
为其对偶完美复形
（上同调，引理 \ref{cohomology-lemma-dual-perfect-complex}）。
由此得到完美对象的逆系统
$\ldots \to E_3^\vee \to E_2^\vee \to E_1^\vee$。
它又给出逆系统
$$
\ldots \to K_3 \to K_2 \to K_1\quad\text{其中}\quad
K_n = Rf_*(E_n^\vee \otimes_{\mathcal{O}_X}^\mathbf{L} \mathcal{G}^\bullet)
$$
其各项在 $S$ 上都是完美的，见引理 \ref{perfect-lemma-tensor-perfect}。
由引理 \ref{perfect-lemma-compute-ext-perfect} 及其证明，
再结合上一段取 $P = E_n$ 时的论证，
对 $S$ 上任意拟凝聚层 $\mathcal{F}$，都有函子性的典范映射
$$
\xymatrix{
& \Ext^i_{\mathcal{O}_X}(E,
\mathcal{G}^\bullet \otimes_{\mathcal{O}_X} f^*\mathcal{F})
\ar[ld] \ar[rd] \\
H^i(S, K_{n + 1} \otimes_{\mathcal{O}_S}^\mathbf{L} \mathcal{F})
\ar[rr] & &
H^i(S, K_n \otimes_{\mathcal{O}_S}^\mathbf{L} \mathcal{F})
}
$$
它们在 $i \leq n + a$ 时是同构。
令 $L_n = K_n^\vee$ 为对偶完美复形。
于是 $L_1 \to L_2 \to L_3 \to \ldots$
是 $D(\mathcal{O}_S)$ 中的完美对象系统，
并且对 $S$ 上任意拟凝聚层 $\mathcal{F}$，映射
$$
\Ext^i_{\mathcal{O}_S}(L_{n + 1}, \mathcal{F})
\longrightarrow
\Ext^i_{\mathcal{O}_S}(L_n, \mathcal{F})
$$
在 $i \leq n + a - 1$ 时均为同构。这说明
$L_n \to L_{n + 1}$ 在截断
$\tau_{\geq -n - a + 2}$ 上诱导同构（提示：取
$L_n \to L_{n + 1}$ 的锥，并考察它最后一个非零上同调层）。
因此 $L = \text{hocolim} L_n$ 是伪凝聚的，见
引理 \ref{perfect-lemma-pseudo-coherent-hocolim}。同伦余极限的映射性质给出
$\Ext^i_{\mathcal{O}_S}(L, \mathcal{F}) =
\Ext^i_{\mathcal{O}_S}(L_n, \mathcal{F})$
在 $i \leq n + a - 3$ 时成立，证明完毕。
\end{proof}

\begin{remark}
\label{perfect-remark-base-change-of-L}
引理 \ref{perfect-lemma-compute-ext} (B) 中的伪凝聚复形 $L$
由该情形典范地确定。例如，按 (B) 构造 $L$
与基变换相容。换言之，给定笛卡尔图
$$
\xymatrix{
X' \ar[r]_{g'} \ar[d]_{f'} &
X \ar[d]^f \\
S' \ar[r]^g &
S
}
$$
以及 $S'$ 上的拟凝聚层 $\mathcal{F}'$，存在函子性的典范同构
$$
\Ext^i_{\mathcal{O}_{S'}}(Lg^*L, \mathcal{F}') \longrightarrow
\Ext^i_{\mathcal{O}_X}(L(g')^*E,
(g')^*\mathcal{G}^\bullet \otimes_{\mathcal{O}_{X'}} (f')^*\mathcal{F}')
$$
注意，在右侧我们对 $\mathcal{G}^\bullet$ {\bf 不}使用导出拉回。
若以后需要这种形式，我们将在此表述一个精确结果并给出详细证明。
\end{remark}






\section{极限与导出范畴}
\label{perfect-section-limits}

\noindent
本节汇集若干关于一个概形的导出范畴的结果，
其中该概形是某个概形逆系统的极限。
更确切地说，我们将在下述设定中工作。

\begin{situation}
\label{perfect-situation-descent}
设 $S = \lim_{i \in I} S_i$ 是一个概形有向系统的极限，
其转移态射 $f_{i'i} : S_{i'} \to S_i$ 均为仿射态射。
假设对所有 $i \in I$，$S_i$ 都是拟紧且拟分离的。
以 $f_i : S \to S_i$ 表示投影，并固定一个元素 $0 \in I$。
\end{situation}

\begin{lemma}
\label{perfect-lemma-descend-homomorphisms}
在情形 \ref{perfect-situation-descent} 中，令 $E_0$ 与 $K_0$ 为
$D(\mathcal{O}_{S_0})$ 的对象。
当 $i \geq 0$ 时置
$E_i = Lf_{i0}^*E_0$、$K_i = Lf_{i0}^*K_0$，
并置 $E = Lf_0^*E_0$、$K = Lf_0^*K_0$。则映射
$$
\colim_{i \geq 0} \Hom_{D(\mathcal{O}_{S_i})}(E_i, K_i)
\longrightarrow
\Hom_{D(\mathcal{O}_S)}(E, K)
$$
在下列任一情形下是同构：
\begin{enumerate}
\item $E_0$ 是完美的，且 $K_0 \in D_\QCoh(\mathcal{O}_{S_0})$；或者
\item $E_0$ 是伪凝聚的，且
$K_0 \in D_\QCoh(\mathcal{O}_{S_0})$ 具有有限 Tor 维数。
\end{enumerate}
\end{lemma}

\begin{proof}
对每个开集 $U_0 \subset S_0$，考虑条件 $P$：
典范映射
$$
\colim_{i \geq 0} \Hom_{D(\mathcal{O}_{U_i})}(E_i|_{U_i}, K_i|_{U_i})
\longrightarrow
\Hom_{D(\mathcal{O}_U)}(E|_U, K|_U)
$$
是同构，其中 $U = f_0^{-1}(U_0)$ 且 $U_i = f_{i0}^{-1}(U_0)$。
我们将利用概形的上同调，引理
\ref{coherent-lemma-induction-principle} 的归纳原理，
证明 $P$ 对所有拟紧开集 $U_0$ 成立。
该引理的条件 (2) 立即由导出范畴中的 Hom 的
Mayer--Vietoris 性质得到，见上同调，引理
\ref{cohomology-lemma-mayer-vietoris-hom}。
所以只需在 $S_0$ 仿射时证明本引理。

\medskip\noindent
假设 $S_0$ 仿射。记 $S_0 = \Spec(A_0)$、$S_i = \Spec(A_i)$，
以及 $S = \Spec(A)$。下文将不再说明地使用
引理 \ref{perfect-lemma-affine-compare-bounded}。

\medskip\noindent
在情形 (1) 中，对象 $E_0^\bullet$ 对应于有限生成投射
$A_0$-模的一个有限复形，见引理 \ref{perfect-lemma-perfect-affine}。
可以用 $A_0$-模的一个 K-平坦复形 $K_0^\bullet$
表示对象 $K_0$。在此情形中，需要证明
$$
\colim_{i \geq 0} \Hom_{D(A_i)}(E_0^\bullet \otimes_{A_0} A_i,
K_0^\bullet \otimes_{A_0} A_i)
\longrightarrow
\Hom_{D(A)}(E_0^\bullet \otimes_{A_0} A, K_0^\bullet \otimes_{A_0} A)
$$
由于 $E_0^\bullet$ 是投射模的上有界复形，
可将其改写为
$$
\colim_{i \geq 0} \Hom_{K(A_0)}(E_0^\bullet,
K_0^\bullet \otimes_{A_0} A_i)
\longrightarrow
\Hom_{K(A_0)}(E_0^\bullet, K_0^\bullet \otimes_{A_0} A)
$$
由于非零模 $E_0^n$ 只有有限多个，且它们都是有限表示模，
此映射是同构。

\medskip\noindent
在情形 (2) 中，对象 $E_0$ 对应于有限自由 $A_0$-模的一个
上有界复形 $E_0^\bullet$，见引理
\ref{perfect-lemma-pseudo-coherent-affine}。
可以用平坦 $A_0$-模的一个有限复形 $K_0^\bullet$
表示 $K_0$，见引理 \ref{perfect-lemma-tor-dimension-affine}
与代数续篇，引理 \ref{more-algebra-lemma-tor-amplitude}。
特别地，$K_0^\bullet$ 是 K-平坦的，因而可以像前面一样
归结到映射
$$
\colim_{i \geq 0} \Hom_{K(A_0)}(E_0^\bullet,
K_0^\bullet \otimes_{A_0} A_i)
\longrightarrow
\Hom_{K(A_0)}(E_0^\bullet, K_0^\bullet \otimes_{A_0} A)
$$
显然，此映射是同构（因为 $K_0^\bullet$ 有界，
其中只涉及有限多个项）。
\end{proof}

\begin{lemma}
\label{perfect-lemma-descend-perfect}
在情形 \ref{perfect-situation-descent} 中，$D(\mathcal{O}_S)$
的完美对象范畴是各 $D(\mathcal{O}_{S_i})$
的完美对象范畴的余极限。
\end{lemma}

\begin{proof}
对每个开集 $U_0 \subset S_0$，考虑条件 $P$：
函子
$$
\colim_{i \geq 0} D_{perf}(\mathcal{O}_{U_i})
\longrightarrow
D_{perf}(\mathcal{O}_U)
$$
是等价，其中 ${}_{perf}$ 表示完美对象构成的全子范畴，
且 $U = f_0^{-1}(U_0)$、$U_i = f_{i0}^{-1}(U_0)$。
我们将利用概形的上同调，引理
\ref{coherent-lemma-induction-principle} 的归纳原理，
证明 $P$ 对所有拟紧开集 $U_0$ 成立。
首先，由引理 \ref{perfect-lemma-descend-homomorphisms}，
已经知道该函子是全忠实的。因此只需证明本质满性。

\medskip\noindent
先验证归纳原理的条件 (2)。设
$S_0 = U_0 \cup V_0$，且 $P$ 对
$U_0$、$V_0$ 与 $U_0 \cap V_0$ 成立。令 $E$ 为
$D(\mathcal{O}_S)$ 的完美对象。可以找到 $i \geq 0$，
以及 $U_i$ 上的完美对象 $E_{U, i}$ 和 $V_i$ 上的完美对象
$E_{V, i}$，使它们到 $U$ 与 $V$ 的拉回分别同构于
$E|_U$ 与 $E|_V$。记
$$
a : E_{U, i} \to (Rf_{i, *}E)|_{U_i}
\quad\text{且}\quad
b : E_{V, i} \to (Rf_{i, *}E)|_{V_i}
$$
为同构 $Lf_i^*E_{U, i} \to E|_U$
与 $Lf_i^*E_{V, i} \to E|_V$ 的伴随映射。
由全忠实性，增大 $i$ 后，可以找到同构
$c : E_{U, i}|_{U_i \cap V_i} \to E_{V, i}|_{U_i \cap V_i}$
，其拉回为复合识别
$$
Lf_i^*E_{U, i}|_{U \cap V} \to E|_{U \cap V} \to Lf_i^*E_{V, i}|_{U \cap V}.
$$
应用上同调，引理 \ref{cohomology-lemma-glue}，
得到 $S_i$ 上的对象 $E_i$ 与映射 $d : E_i \to Rf_{i, *}E$，
其在 $U_i$ 与 $V_i$ 上分别限制为 $a$ 与 $b$。
于是显然 $E_i$ 是完美的，并且
$d$ 是某个同构 $Lf_i^*E_i \to E$ 的伴随映射。

\medskip\noindent
最后验证归纳原理的条件 (1)，即验证
$S_0$ 仿射时本引理成立。
记 $S_0 = \Spec(A_0)$、$S_i = \Spec(A_i)$，以及
$S = \Spec(A)$。利用引理 \ref{perfect-lemma-affine-compare-bounded}
与 \ref{perfect-lemma-perfect-affine}，只需证明
$$
D_{perf}(A) = \colim D_{perf}(A_i)
$$
这由以下事实立即可得：环上的完美复形由有限生成投射模
（因而是有限表示模）的有限复形给出。细节见代数续篇，
引理 \ref{more-algebra-lemma-colimit-perfect-complexes}。
\end{proof}





\section{上同调与基变换 VI}
\label{perfect-section-cohomology-and-base-change-final}

\noindent
这是关于上同调与基变换的最后一节，继续各节
\ref{perfect-section-cohomology-and-base-change-perfect},
\ref{perfect-section-cohomology-base-change} 与
\ref{perfect-section-producing-perfect} 的讨论。
评注 \ref{perfect-remark-explain-perfect-direct-image}
给出了一个容易理解的特殊情形。

\begin{lemma}
\label{perfect-lemma-base-change-tensor-perfect}
设 $f : X \to S$ 是有限表示态射，
$E \in D(\mathcal{O}_X)$ 是完美对象。设 $\mathcal{G}^\bullet$
是有限表示 $\mathcal{O}_X$-模的有界复形，
它在 $S$ 上平坦且支集在 $S$ 上固有。则
$$
K = Rf_*(E \otimes_{\mathcal{O}_X}^\mathbf{L} \mathcal{G}^\bullet)
$$
是 $D(\mathcal{O}_S)$ 的完美对象，且其构造与任意基变换相容。
\end{lemma}

\begin{proof}
关于基变换的断言就是引理 \ref{perfect-lemma-base-change-tensor}。
所以只需证明 $K$ 是完美对象。当 $S$ 是诺特概形时，
这由引理 \ref{perfect-lemma-tensor-perfect} 得到。
我们将利用诺特逼近归结到这种情形。
建议读者跳过本证明的余下部分。

\medskip\noindent
问题在 $S$ 上是局部的，所以可以假设 $S$ 仿射。
记 $S = \Spec(R)$。把 $R = \colim R_i$
写成诺特环 $R_i$ 的滤过余极限。由极限，引理
\ref{limits-lemma-descend-finite-presentation}
，存在某个 $i$ 及一个在 $R_i$ 上有限表示的概形 $X_i$，
其到 $R$ 的基变换为 $X$。由极限，引理
\ref{limits-lemma-descend-modules-finite-presentation}，
增大 $i$ 后，可以假设存在有限表示 $\mathcal{O}_{X_i}$-模的
一个有界复形 $\mathcal{G}_i^\bullet$，
其到 $X$ 的拉回为 $\mathcal{G}^\bullet$。
再次增大 $i$，可以假设 $\mathcal{G}_i^n$ 在 $R_i$ 上平坦，
见极限，引理 \ref{limits-lemma-descend-module-flat-finite-presentation}。
再增大 $i$，可以假设 $\mathcal{G}_i^n$ 的支集在 $R_i$ 上固有，
见极限，引理 \ref{limits-lemma-eventually-proper-support}
与概形的上同调，引理
\ref{coherent-lemma-module-support-proper-over-base}。
最后，由引理 \ref{perfect-lemma-descend-perfect}，增大 $i$ 后，
可以假设存在 $D(\mathcal{O}_{X_i})$ 的完美对象 $E_i$，
其到 $X$ 的拉回为 $E$。把引理 \ref{perfect-lemma-tensor-perfect}
应用于 $X_i \to \Spec(R_i)$、$E_i$、$\mathcal{G}_i^\bullet$，
并利用已经证明的基变换性质，即得结论。
\end{proof}

\begin{remark}
\label{perfect-remark-explain-perfect-direct-image}
设 $R$ 是环，$X$ 是在 $R$ 上有限表示的概形。
设 $\mathcal{G}$ 是有限表示 $\mathcal{O}_X$-模，
它在 $R$ 上平坦且支集在 $R$ 上固有。由
引理 \ref{perfect-lemma-base-change-tensor-perfect}，
存在有限生成投射 $R$-模的有限复形 $M^\bullet$，使得
$$
R\Gamma(X_{R'}, \mathcal{G}_{R'}) = M^\bullet \otimes_R R'
$$
对 $R$-代数 $R'$ 函子性地成立。
\end{remark}

\begin{lemma}
\label{perfect-lemma-base-change-tensor-pseudo-coherent}
设 $f : X \to S$ 是有限表示态射，
$E \in D(\mathcal{O}_X)$ 是伪凝聚对象。
设 $\mathcal{G}^\bullet$ 是有限表示 $\mathcal{O}_X$-模的上有界复形，
它在 $S$ 上平坦且支集在 $S$ 上固有。则
$$
K = Rf_*(E \otimes_{\mathcal{O}_X}^\mathbf{L} \mathcal{G}^\bullet)
$$
是 $D(\mathcal{O}_S)$ 的伪凝聚对象，且其构造与任意基变换相容。
\end{lemma}

\begin{proof}
关于基变换的断言就是引理 \ref{perfect-lemma-base-change-tensor}。
所以只需证明 $K$ 是伪凝聚对象。
这将由完美复形逼近及引理
\ref{perfect-lemma-base-change-tensor-perfect} 得到。
建议读者跳过本证明的余下部分。

\medskip\noindent
问题在 $S$ 上是局部的，所以可以假设 $S$ 仿射。
于是 $X$ 拟紧且拟分离。此外，存在整数 $N$，使全直像
$Rf_* : D_\QCoh(\mathcal{O}_X) \to D_\QCoh(\mathcal{O}_S)$
具有上同调维数 $N$，见引理
\ref{perfect-lemma-quasi-coherence-direct-image} 中的说明。
选取整数 $b$，使得当 $i > b$ 时 $\mathcal{G}^i = 0$。
只需证明对每个 $m$，$K$ 都是 $m$-伪凝聚的。
选取由完美复形 $P$ 给出的 $(X, E, m - N - 1 - b)$
的一个逼近 $P \to E$。其存在性由
定理 \ref{perfect-theorem-approximation} 保证。
在 $D_\QCoh(\mathcal{O}_X)$ 中选取可区别三角
$$
P \to E \to C \to P[1]
$$
。$C$ 的上同调层在次数 $\geq m - N - 1 - b$ 均为零。
因此 $C \otimes^\mathbf{L} \mathcal{G}^\bullet$
的上同调层在次数 $\geq m - N - 1$ 均为零。于是
$Rf_*(C \otimes^\mathbf{L} \mathcal{G}^\bullet)$
的上同调层在次数 $\geq m - 1$ 均为零。因此
$$
Rf_*(P \otimes^\mathbf{L} \mathcal{G}^\bullet) \to
Rf_*(E \otimes^\mathbf{L} \mathcal{G}^\bullet)
$$
在次数 $\geq m$ 的上同调层上是同构。
其次，设当 $i > a$ 时 $H^i(P) = 0$。则
$
P \otimes^\mathbf{L} \sigma_{\geq m - N - 1 - a}\mathcal{G}^\bullet
\longrightarrow
P \otimes^\mathbf{L} \mathcal{G}^\bullet
$
在次数 $\geq m - N - 1$ 的上同调层上是同构。
因而再次得到
$$
Rf_*(P \otimes^\mathbf{L} \sigma_{\geq m - N - 1 - a}\mathcal{G}^\bullet) \to
Rf_*(P \otimes^\mathbf{L} \mathcal{G}^\bullet)
$$
在次数 $\geq m$ 的上同调层上是同构。
由引理 \ref{perfect-lemma-base-change-tensor-perfect}，其源是完美复形。
因此 $K$ 如所需是 $m$-伪凝聚的。
\end{proof}

\begin{lemma}
\label{perfect-lemma-flat-proper-perfect-direct-image-general}
设 $S$ 是概形，$f : X \to S$ 是有限表示的固有态射。
\begin{enumerate}
\item 若 $E \in D(\mathcal{O}_X)$ 完美且 $f$ 平坦，则
$Rf_*E$ 是 $D(\mathcal{O}_S)$ 的完美对象，且其构造与任意基变换相容。
\item 若 $\mathcal{G}$ 是有限表示的 $\mathcal{O}_X$-模且在 $S$ 上平坦，
则 $Rf_*\mathcal{G}$ 是 $D(\mathcal{O}_S)$ 的完美对象，
且其构造与任意基变换相容。
\end{enumerate}
\end{lemma}

\begin{proof}
这是引理 \ref{perfect-lemma-base-change-tensor-perfect} 的特殊情形：
(1) 取 $\mathcal{G}^\bullet$ 为置于次数 $0$ 的 $\mathcal{O}_X$；
(2) 取 $E = \mathcal{O}_X$，并令 $\mathcal{G}^\bullet$
为置于次数 $0$ 的 $\mathcal{G}$。
\end{proof}

\begin{lemma}
\label{perfect-lemma-flat-proper-pseudo-coherent-direct-image-general}
设 $S$ 是概形，$f : X \to S$ 是有限表示的平坦固有态射。
设 $E \in D(\mathcal{O}_X)$ 是伪凝聚的。则 $Rf_*E$
是 $D(\mathcal{O}_S)$ 的伪凝聚对象，且其构造与任意基变换相容。
\end{lemma}

\noindent
更一般地，若 $f : X \to S$ 固有，且 $X$ 上的 $E$
相对于 $S$ 伪凝聚（态射续篇，定义
\ref{more-morphisms-definition-relative-pseudo-coherence}),
则 $Rf_*E$ 是伪凝聚的
（但在这种一般性下，其构造不与基变换相容）。
见 \cite{Kiehl}。

\begin{proof}
这是引理 \ref{perfect-lemma-base-change-tensor-pseudo-coherent} 中
取 $\mathcal{G}^\bullet$ 为置于次数 $0$ 的 $\mathcal{O}_X$
所得的特殊情形。
\end{proof}

\begin{lemma}
\label{perfect-lemma-pullback-and-limits}
设 $R$ 是环，$X$ 是概形，并设
$f : X \to \Spec(R)$ 是有限表示的平坦固有态射。
设 $(M_n)$ 为转移映射均满射的 $R$-模逆系统。
则典范映射
$$
\mathcal{O}_X \otimes_R (\lim M_n)
\longrightarrow
\lim \mathcal{O}_X \otimes_R M_n
$$
诱导从其源到对其目标应用 $DQ_X$ 后所得对象的同构。
\end{lemma}

\begin{proof}
该断言意指，对 $D_\QCoh(\mathcal{O}_X)$ 的任意对象 $E$，
诱导映射
$$
\Hom(E, \mathcal{O}_X \otimes_R (\lim M_n))
\longrightarrow
\Hom(E, \lim \mathcal{O}_X \otimes_R M_n)
$$
是同构。由于 $D_\QCoh(\mathcal{O}_X)$ 有完美生成元
（定理 \ref{perfect-theorem-bondal-van-den-Bergh}），
只需对完美的 $E$ 验证此事。
由引理 \ref{perfect-lemma-Rlim-quasi-coherent}，
$\lim \mathcal{O}_X \otimes_R M_n = R\lim \mathcal{O}_X \otimes_R M_n$.
正合函子
$R\Hom_X(E, -) : D_\QCoh(\mathcal{O}_X) \to D(R)$
（见上同调，第 \ref{cohomology-section-global-RHom} 节）
与积相容，因而也与导出极限相容，所以
$$
R\Hom_X(E, \lim \mathcal{O}_X \otimes_R M_n) =
R\lim R\Hom_X(E, \mathcal{O}_X \otimes_R M_n)
$$
令 $E^\vee$ 为对偶完美复形，见上同调，引理
\ref{cohomology-lemma-dual-perfect-complex}。有
$$
R\Hom_X(E, \mathcal{O}_X \otimes_R M_n) =
R\Gamma(X, E^\vee \otimes_{\mathcal{O}_X}^\mathbf{L} Lf^*M_n) =
R\Gamma(X, E^\vee) \otimes_R^\mathbf{L} M_n
$$
，这里使用了引理 \ref{perfect-lemma-cohomology-base-change}。
由引理 \ref{perfect-lemma-flat-proper-perfect-direct-image-general}，
$R\Gamma(X, E^\vee)$ 是 $R$-模的完美复形。
特别地，它是伪凝聚复形，因而由代数续篇，引理
\ref{more-algebra-lemma-pseudo-coherent-tensor-limit} 得到
$$
R\lim R\Gamma(X, E^\vee) \otimes_R^\mathbf{L} M_n =
R\Gamma(X, E^\vee) \otimes_R^\mathbf{L} \lim M_n
$$
，正如所需。
\end{proof}

\begin{lemma}
\label{perfect-lemma-base-change-RHom-perfect}
设 $f : X \to S$ 是有限表示态射，
$E \in D(\mathcal{O}_X)$ 是完美对象。设 $\mathcal{G}^\bullet$
是有限表示 $\mathcal{O}_X$-模的有界复形，
它在 $S$ 上平坦且支集在 $S$ 上固有。则
$$
K = Rf_*R\SheafHom(E, \mathcal{G}^\bullet)
$$
是 $D(\mathcal{O}_S)$ 的完美对象，且其构造与任意基变换相容。
\end{lemma}

\begin{proof}
关于基变换的断言就是引理 \ref{perfect-lemma-base-change-RHom}。
所以只需证明 $K$ 是完美对象。当 $S$ 是诺特概形时，
这由引理 \ref{perfect-lemma-ext-perfect} 得到。
我们将利用诺特逼近归结到这种情形。
建议读者跳过本证明的余下部分。

\medskip\noindent
问题在 $S$ 上是局部的，所以可以假设 $S$ 仿射。
记 $S = \Spec(R)$。把 $R = \colim R_i$
写成诺特环 $R_i$ 的滤过余极限。由极限，引理
\ref{limits-lemma-descend-finite-presentation}
，存在某个 $i$ 及一个在 $R_i$ 上有限表示的概形 $X_i$，
其到 $R$ 的基变换为 $X$。由极限，引理
\ref{limits-lemma-descend-modules-finite-presentation}，
增大 $i$ 后，可以假设存在有限表示 $\mathcal{O}_{X_i}$-模的
一个有界复形 $\mathcal{G}_i^\bullet$，
其到 $X$ 的拉回为 $\mathcal{G}^\bullet$。
再次增大 $i$，可以假设 $\mathcal{G}_i^n$ 在 $R_i$ 上平坦，
见极限，引理 \ref{limits-lemma-descend-module-flat-finite-presentation}。
再增大 $i$，可以假设 $\mathcal{G}_i^n$ 的支集在 $R_i$ 上固有，
见极限，引理 \ref{limits-lemma-eventually-proper-support}
与概形的上同调，引理
\ref{coherent-lemma-module-support-proper-over-base}。
最后，由引理 \ref{perfect-lemma-descend-perfect}，增大 $i$ 后，
可以假设存在 $D(\mathcal{O}_{X_i})$ 的完美对象 $E_i$，
其到 $X$ 的拉回为 $E$。把引理 \ref{perfect-lemma-ext-perfect}
应用于 $X_i \to \Spec(R_i)$、$E_i$、$\mathcal{G}_i^\bullet$，
并利用已经证明的基变换性质，即得结论。
\end{proof}









\section{完美复形}
\label{perfect-section-perfect-complexes}

\noindent
先讨论完美复形 Betti 数的跳跃轨迹。
给定概形 $X$ 上的复形 $E$ 与 $X$ 的点 $x$，
我们常以 $E \otimes_{\mathcal{O}_X}^\mathbf{L} \kappa(x)$
代替更准确的写法 $Li_x^*E$，其中
$i_x : x \to X$ 是典范态射。

\begin{lemma}
\label{perfect-lemma-jump-loci}
设 $X$ 是概形，$E \in D(\mathcal{O}_X)$ 伪凝聚
（例如它是完美的）。对任意 $i \in \mathbf{Z}$，考虑函数
$$
\beta_i : X \longrightarrow \{0, 1, 2, \ldots\},\quad
x \longmapsto
\dim_{\kappa(x)}
H^i(E \otimes_{\mathcal{O}_X}^\mathbf{L} \kappa(x))
$$
则有：
\begin{enumerate}
\item $\beta_i$ 的构造与任意基变换相容；
\item 函数 $\beta_i$ 上半连续；
\item $\beta_i$ 的等值集在 $X$ 中局部可构造。
\end{enumerate}
\end{lemma}

\begin{proof}
考虑概形态射 $f : Y \to X$ 与点 $y \in Y$。
令 $x$ 为 $y$ 的像，并考虑交换图
$$
\xymatrix{
y \ar[r]_j \ar[d]_g & Y \ar[d]^f \\
x \ar[r]^i & X
}
$$
于是 $Lg^* \circ Li^* = Lj^* \circ Lf^*$。
这说明，与伪凝聚复形 $Lf^*E$ 相伴的函数 $\beta'_i$
是函数 $\beta_i$ 的拉回，即
$\beta'_i = \beta_i \circ f$。这正是 (1) 的含义。

\medskip\noindent
固定 $i$ 并取 $x \in X$。只需在 $x$ 的某个开邻域中
证明 (2) 与 (3)，故可假设 $X$ 仿射。
于是可用有限自由模的一个上有界复形 $\mathcal{F}^\bullet$
表示 $E$（引理 \ref{perfect-lemma-lift-pseudo-coherent}）。
此时 $P = \sigma_{\geq i - 1}\mathcal{F}^\bullet$ 是完美对象，
且 $P \to E$ 对所有 $x' \in X$ 都诱导同构
$$
H^i(P \otimes_{\mathcal{O}_X}^\mathbf{L} \kappa(x')) \to
H^i(E \otimes_{\mathcal{O}_X}^\mathbf{L} \kappa(x'))
$$
。因此可以假设 $E$ 完美。在这种情形中，由代数续篇，引理
\ref{more-algebra-lemma-lift-perfect-from-residue-field}
，存在 $x$ 的仿射开邻域 $U$ 以及 $a \leq b$，
使 $E|_U$ 可由复形
$$
\ldots \to 0 \to \mathcal{O}_U^{\oplus \beta_a(x)}
\to \mathcal{O}_U^{\oplus \beta_{a + 1}(x)} \to
\ldots \to
\mathcal{O}_U^{\oplus \beta_{b - 1}(x)} \to
\mathcal{O}_U^{\oplus \beta_b(x)} \to 0
\to \ldots
$$
表示。（这里还使用了把问题化为代数问题的较早结果，例如
引理 \ref{perfect-lemma-affine-compare-bounded} 与
\ref{perfect-lemma-perfect-affine}。）
立即可见，对所有 $x' \in U$，
$\beta_i(x') \leq \beta_i(x)$。这证明了 $\beta_i$ 上半连续。

\medskip\noindent
为证明 (3)，可以假设 $X$ 仿射，并且
$E$ 由有限自由 $\mathcal{O}_X$-模的一个复形给出
（例如沿用上一段的论证，或使用上同调，引理
\ref{cohomology-lemma-perfect-on-locally-ringed}）。
因此需要证明：给定复形
$$
\mathcal{O}_X^{\oplus a} \to
\mathcal{O}_X^{\oplus b} \to
\mathcal{O}_X^{\oplus c}
$$
，把点 $x \in X$ 映到
$\kappa_x^{\oplus a} \to \kappa_x^{\oplus b} \to \kappa_x^{\oplus c}$
在中间项处上同调维数的函数具有可构造等值集。
令 $A \in \text{Mat}(a \times b, \Gamma(X, \mathcal{O}_X))$
为第一个箭头的矩阵。$A$ 在
$\text{Mat}(a \times b, \kappa(x))$ 中的像秩等于 $r$，
当且仅当 $A$ 的所有 $(r + 1) \times (r + 1)$ 子式在 $x$ 处为零，
并且 $A$ 的某个 $r \times r$ 子式在 $x$ 处非零。
所以秩等于 $r$ 的点集是可构造的局部闭集。
对第二个箭头作同样论证，再把这些结果合在一起，即得所需结论。
\end{proof}

\begin{lemma}
\label{perfect-lemma-chi-locally-constant}
设 $X$ 是概形，$E \in D(\mathcal{O}_X)$ 完美。函数
$$
\chi_E : X \longrightarrow \mathbf{Z},\quad
x \longmapsto \sum (-1)^i
\dim_{\kappa(x)} H^i(E \otimes_{\mathcal{O}_X}^\mathbf{L} \kappa(x))
$$
在 $X$ 上局部常值。
\end{lemma}

\begin{proof}
由上同调，引理
\ref{cohomology-lemma-perfect-on-locally-ringed}
可知，在 $X$ 上局部地，可以用有限自由
$\mathcal{O}_X$-模的有限复形 $\mathcal{E}^\bullet$ 表示 $E$。
在这样的开集上，函数 $\chi_E$ 为常值，其值是
$\sum (-1)^i \text{rank}(\mathcal{E}^i)$.
\end{proof}

\begin{lemma}
\label{perfect-lemma-open-where-cohomology-in-degree-i-rank-r}
设 $X$ 是概形，$E \in D(\mathcal{O}_X)$ 完美。
给定 $i, r \in \mathbf{Z}$，存在开子概形 $U \subset X$，
由下列性质刻画：
\begin{enumerate}
\item $E|_U \cong H^i(E|_U)[-i]$，且 $H^i(E|_U)$
是秩为 $r$ 的局部自由 $\mathcal{O}_U$-模；
\item 态射 $f : Y \to X$ 经过 $U$ 分解，当且仅当
$Lf^*E$ 同构于置于次数 $i$ 的一个秩为 $r$ 的局部自由模。
\end{enumerate}
\end{lemma}

\begin{proof}
对 $j \in \mathbf{Z}$，令
$\beta_j : X \to \{0, 1, 2, \ldots\}$
为引理 \ref{perfect-lemma-jump-loci} 中的函数。则集合
$$
W = \{x \in X \mid \beta_j(x) \leq 0\text{ 对所有 }j \not = i\}
$$
在 $X$ 中是开集，且其构造与到任意 $X$-概形 $Y$ 的拉回相容。
这是该引理的推论，因为先验地，在任一点的某个邻域中
只有有限多个 $\beta_j$ 非零。因此可以用 $W$ 替换 $X$，
并假设对所有 $x \in X$ 与所有 $j \not = i$，
$\beta_j(x) = 0$。在这种情形中，$H^i(E)$ 是有限局部自由模，
且 $E \cong H^i(E)[-i]$；例如见代数续篇，引理
\ref{more-algebra-lemma-lift-perfect-from-residue-field}。
于是 $X$ 是若干开子概形的不交并，并且 $H^i(E)$
在每个这样的开子概形上具有固定的秩，结论得证。
\end{proof}

\begin{lemma}
\label{perfect-lemma-locally-closed-where-H0-locally-free}
设 $X$ 是概形。设 $E \in D(\mathcal{O}_X)$ 完美，
并且对某些 $a, b \in \mathbf{Z}$，其 Tor 振幅包含于 $[a, b]$。
令 $r \geq 0$。则存在局部闭子概形 $j : Z \to X$，
由下列性质刻画：
\begin{enumerate}
\item $H^a(Lj^*E)$ 是秩为 $r$ 的局部自由 $\mathcal{O}_Z$-模；
\item 态射 $f : Y \to X$ 经过 $Z$ 分解，当且仅当
对所有态射 $g : Y' \to Y$，$\mathcal{O}_{Y'}$-模
$H^a(L(f \circ g)^*E)$ 都是秩为 $r$ 的局部自由模。
\end{enumerate}
此外，$j : Z \to X$ 是有限表示的，并且：
\begin{enumerate}
\item[(3)] 若 $f : Y \to X$ 分解为 $Y \xrightarrow{g} Z \to X$，则
$H^a(Lf^*E) = g^*H^a(Lj^*E)$,
\item[(4)] 若对所有 $x \in X$ 都有 $\beta_a(x) \leq r$，
则 $j$ 是闭浸入，并且对给定的 $f : Y \to X$，下列条件等价：
\begin{enumerate}
\item $f : Y \to X$ 经过 $Z$ 分解；
\item $H^a(Lf^*E)$ 是秩为 $r$ 的局部自由 $\mathcal{O}_Y$-模；
\end{enumerate}
若 $r = 1$，它们还等价于
\begin{enumerate}
\item[(c)] $\mathcal{O}_Y \to \SheafHom_{\mathcal{O}_Y}(H^a(Lf^*E), H^a(Lf^*E))$
是单射。
\end{enumerate}
\end{enumerate}
\end{lemma}

\begin{proof}
首先，令 $U \subset X$ 为局部可构造开子概形，
在其上引理 \ref{perfect-lemma-jump-loci} 的函数
$\beta_a$ 取值 $\leq r$。令 $f : Y \to X$ 如 (2) 所述。
则对任意 $y \in Y$，都有 $\beta_a(Lf^*E) = r$，
故由引理 \ref{perfect-lemma-jump-loci}，$y$ 映入 $U$。
所以 (2) 中的 $f$ 经过 $U$ 分解。
因此可以用 $U$ 替换 $X$，并假设对所有 $x \in X$，
$\beta_a(x) \in \{0, 1, \ldots, r\}$。
我们将证明，在这种情形中，存在由有限型拟凝聚理想截出的
闭子概形 $Z \subset X$，它由 (4) 中 (a)、(b)
（以及当 $r = 1$ 时的 (4)(c)）的等价性刻画，并且 (3) 成立。
这将完成证明，因为它尤其说明 (2) 中的态射都经过 $Z$ 分解。

\medskip\noindent
若 $x \in X$ 且 $\beta_a(x) < r$，则存在 $x$ 的开邻域，
在其上 $\beta_a < r$（引理 \ref{perfect-lemma-jump-loci}）。
由此可见，至少在集合论意义下，$Z$ 是闭子集。

\medskip\noindent
为赋予它概形论结构，考虑满足 $\beta_a(x) = r$ 的点
$x \in X$，并置 $\beta = \beta_{a + 1}(x)$。
由代数续篇，引理
\ref{more-algebra-lemma-lift-perfect-from-residue-field}
，存在 $x$ 的仿射开邻域 $U$，使得 $K|_U$ 可由复形
$$
\ldots \to 0 \to \mathcal{O}_U^{\oplus r}
\xrightarrow{(f_{ij})} \mathcal{O}_U^{\oplus \beta} \to
\ldots \to
\mathcal{O}_U^{\oplus \beta_{b - 1}(x)} \to
\mathcal{O}_U^{\oplus \beta_b(x)} \to 0
\to \ldots
$$
表示。（这里还使用了把问题化为代数问题的较早结果，例如
引理 \ref{perfect-lemma-affine-compare-bounded} 与
\ref{perfect-lemma-perfect-affine}。）现在，若 $g : Y \to U$ 是概形态射，
并且对某一对 $i, j$，$g^\sharp(f_{ij})$ 非零，则
$H^a(Lg^*E)$ 不是秩为 $r$ 的局部自由 $\mathcal{O}_Y$-模。
见代数续篇，引理 \ref{more-algebra-lemma-coker-injective-free}。
若对所有 $i, j$ 都有 $g^\sharp(f_{ij}) = 0$，
则 $H^a(Lg^*E)$ 显然是局部自由 $\mathcal{O}_Y$-模。
因此在 $U$ 上，由所有 $f_{ij}$ 截出的闭子概形满足 (3)，
并且 (4)(a) 与 (b) 等价。
$Z$ 的刻画说明这些局部构造的部分可以粘合（略去细节）。
最后，若 $r = 1$，则 (4)(c) 与 (4)(b) 等价，
因为此时局部地，$H^a(Lg^*E) \subset \mathcal{O}_Y$
是由各元素 $g^\sharp(f_{ij})$ 生成的理想的零化子。
\end{proof}





\section{应用}
\label{perfect-section-applications}

\noindent
这里大多是上同调与基变换的应用。将来可以把这些结果
推广到引理 \ref{perfect-lemma-base-change-tensor-perfect} 所讨论的情形。

\begin{lemma}
\label{perfect-lemma-jump-loci-geometric}
设 $f : X \to S$ 是有限表示的固有态射。
设 $\mathcal{F}$ 是有限表示的 $\mathcal{O}_X$-模，
并且在 $S$ 上平坦。对固定的 $i \in \mathbf{Z}$，考虑函数
$$
\beta_i : S \to \{0, 1, 2, \ldots\},\quad
s \longmapsto \dim_{\kappa(s)} H^i(X_s, \mathcal{F}_s)
$$
则有：
\begin{enumerate}
\item $\beta_i$ 的构造与任意基变换相容；
\item 函数 $\beta_i$ 上半连续；
\item $\beta_i$ 的等值集在 $S$ 中局部可构造。
\end{enumerate}
\end{lemma}

\begin{proof}
由上同调与基变换（更确切地说，由引理
\ref{perfect-lemma-flat-proper-perfect-direct-image-general}），
对象 $K = Rf_*\mathcal{F}$ 是 $S$ 的导出范畴中的完美对象，
且其构造与任意基变换相容。特别地，有
$$
H^i(X_s, \mathcal{F}_s) = H^i(K \otimes_{\mathcal{O}_S}^\mathbf{L} \kappa(s))
$$
因此本引理由引理 \ref{perfect-lemma-jump-loci} 得到。
\end{proof}

\begin{lemma}
\label{perfect-lemma-chi-locally-constant-geometric}
设 $f : X \to S$ 是有限表示的固有态射。
设 $\mathcal{F}$ 是有限表示的 $\mathcal{O}_X$-模，
并且在 $S$ 上平坦。函数
$$
s \longmapsto \chi(X_s, \mathcal{F}_s)
$$
在 $S$ 上局部常值，且此函数的构造与基变换相容。
\end{lemma}

\begin{proof}
由上同调与基变换（更确切地说，由引理
\ref{perfect-lemma-flat-proper-perfect-direct-image-general}），
对象 $K = Rf_*\mathcal{F}$ 是 $S$ 的导出范畴中的完美对象，
且其构造与任意基变换相容。因此需要证明映射
$$
s \longmapsto \sum (-1)^i \dim_{\kappa(s)}
H^i(K \otimes^\mathbf{L}_{\mathcal{O}_S} \kappa(s))
$$
在 $S$ 上局部常值。这正是引理 \ref{perfect-lemma-chi-locally-constant}。
\end{proof}

\begin{lemma}
\label{perfect-lemma-open-where-cohomology-in-degree-i-rank-r-geometric}
设 $f : X \to S$ 是有限表示的固有态射。
设 $\mathcal{F}$ 是有限表示的 $\mathcal{O}_X$-模，
并且在 $S$ 上平坦。固定 $i, r \in \mathbf{Z}$。
则存在具有下述性质的开子概形 $U \subset S$：
态射 $T \to S$ 经过 $U$ 分解，当且仅当
$Rf_{T, *}\mathcal{F}_T$ 同构于置于次数 $i$
的秩为 $r$ 的有限局部自由模。
\end{lemma}

\begin{proof}
由上同调与基变换（更确切地说，由引理
\ref{perfect-lemma-flat-proper-perfect-direct-image-general}），
对象 $K = Rf_*\mathcal{F}$ 是 $S$ 的导出范畴中的完美对象，
且其构造与任意基变换相容。因此本引理立即由
引理 \ref{perfect-lemma-open-where-cohomology-in-degree-i-rank-r} 得到。
\end{proof}

\begin{lemma}
\label{perfect-lemma-vanishing-implies-locally-free}
设 $f : X \to S$ 是有限表示态射。
设 $\mathcal{F}$ 是有限表示的 $\mathcal{O}_X$-模，
它在 $S$ 上平坦且支集在 $S$ 上固有。若对 $i > 0$，
$R^if_*\mathcal{F} = 0$，则 $f_*\mathcal{F}$ 局部自由，
且其构造与任意基变换相容（说明见证明）。
\end{lemma}

\begin{proof}
由引理 \ref{perfect-lemma-base-change-tensor-perfect}，
$D(\mathcal{O}_S)$ 的对象 $E = Rf_*\mathcal{F}$ 完美，
且其构造与任意基变换相容；这意味着，对任意笛卡尔图
$$
\xymatrix{
X' \ar[r]_{g'} \ar[d]_{f'} &
X \ar[d]^f \\
S' \ar[r]^g &
S
}
$$
，都有 $Rf'_*(g')^*\mathcal{F} = Lg^*E$。
由于次数 $< 0$ 绝不会出现上同调，可知对某个 $b$，
$E$（局部地）具有包含于 $[0, b]$ 的 Tor 振幅。
若对 $i > 0$，$H^i(E) = R^if_*\mathcal{F} = 0$，
则 $E$ 的 Tor 振幅包含于 $[0, 0]$。所以
$E = H^0(E)[0]$。由代数续篇，引理
\ref{more-algebra-lemma-perfect}
（以及代数，引理 \ref{algebra-lemma-finite-projective}
中对有限投射模的刻画），
$H^0(E) = f_*\mathcal{F}$ 是有限局部自由的。
与基变换相容意味着，对上述图有
$g^*f_*\mathcal{F} = f'_*(g')^*\mathcal{F}$，
而这由已经确立的 $E$ 与基变换相容性得到。
\end{proof}

\begin{lemma}
\label{perfect-lemma-proper-flat-h0}
设 $f : X \to S$ 是概形态射。假设
\begin{enumerate}
\item $f$ 固有、平坦且有限表示；
\item 对所有 $s \in S$，都有
$\kappa(s) = H^0(X_s, \mathcal{O}_{X_s})$。
\end{enumerate}
则有：
\begin{enumerate}
\item[(a)] $f_*\mathcal{O}_X = \mathcal{O}_S$，且在任意基变换后仍成立；
\item[(b)] 在 $S$ 上局部地，有
$$
Rf_*\mathcal{O}_X = \mathcal{O}_S \oplus P
$$
作为 $D(\mathcal{O}_S)$ 中的等式，
其中 $P$ 完美且 Tor 振幅包含于 $[1, \infty)$。
\end{enumerate}
\end{lemma}

\begin{proof}
由上同调与基变换
（引理 \ref{perfect-lemma-flat-proper-perfect-direct-image-general}），
复形 $E = Rf_*\mathcal{O}_X$ 完美，
且其构造与任意基变换相容。这首先说明 $E$
的 Tor 振幅包含于 $[0, \infty)$。其次，它说明对 $s \in S$ 有
$H^0(E \otimes^\mathbf{L} \kappa(s)) =
H^0(X_s, \mathcal{O}_{X_s}) = \kappa(s)$.
所以映射 $\mathcal{O}_S \to Rf_*\mathcal{O}_X = E$
诱导同构
$\mathcal{O}_S \otimes \kappa(s) \to H^0(E \otimes^\mathbf{L} \kappa(s))$.
因此
$H^0(E) \otimes \kappa(s) \to H^0(E \otimes^\mathbf{L} \kappa(s))$
是满射。应用代数续篇，引理
\ref{more-algebra-lemma-better-cut-complex-in-two} 可知，
把 $S$ 替换为 $s$ 的一个仿射开邻域后，有分解
$E = H^0(E) \oplus \tau_{\geq 1}E$，
其中 $\tau_{\geq 1}E$ 完美且 Tor 振幅包含于 $[1, \infty)$。
由于 $E$ 的 Tor 振幅包含于 $[0, \infty)$，
$H^0(E)$ 是平坦 $\mathcal{O}_S$-模。
因此 $H^0(E)$ 是平坦且完美的 $\mathcal{O}_S$-模，
从而有限局部自由，见代数续篇，引理
\ref{more-algebra-lemma-perfect}
（以及代数，引理 \ref{algebra-lemma-finite-projective}
所述的有限投射模为有限局部自由模这一事实）。
于是映射 $\mathcal{O}_S \to H^0(E)$ 是同构，
因为这可在与剩余域作张量后验证
（代数，引理 \ref{algebra-lemma-cokernel-flat}）。
\end{proof}

\begin{lemma}
\label{perfect-lemma-proper-flat-geom-red-connected}
设 $f : X \to S$ 是概形态射。假设
\begin{enumerate}
\item $f$ 固有、平坦且有限表示；
\item $f$ 的几何纤维既约且连通。
\end{enumerate}
则 $f_*\mathcal{O}_X = \mathcal{O}_S$，且在任意基变换后仍成立。
\end{lemma}

\begin{proof}
由引理 \ref{perfect-lemma-proper-flat-h0}，
只需证明对所有 $s \in S$ 都有
$\kappa(s) = H^0(X_s, \mathcal{O}_{X_s})$。
这由簇，引理
\ref{varieties-lemma-proper-geometrically-reduced-global-sections}
以及 $X_s$ 几何连通且几何既约这一事实得到。
\end{proof}

\begin{lemma}
\label{perfect-lemma-proper-idempotent-on-fibre}
设 $f : X \to S$ 是概形的固有态射。令 $s \in S$，
并令 $e \in H^0(X_s, \mathcal{O}_{X_s})$ 为幂等元。
则 $e$ 属于映射
$(f_*\mathcal{O}_X)_s \to H^0(X_s, \mathcal{O}_{X_s})$
的像。
\end{lemma}

\begin{proof}
令 $X_s = T_1 \amalg T_2$ 为与 $e$ 对应的不交并分解，
其中 $T_1$ 与 $T_2$ 非空，且都是 $X_s$ 中的开闭子集；
也就是说，$e$ 在 $T_1$ 上恒等于 $1$，
在 $T_2$ 上恒等于 $0$。

\medskip\noindent
假设 $S$ 是诺特概形。我们将使用概形的上同调，引理
\ref{coherent-lemma-formal-functions-stalk}
所给形式的形式函数定理。它告诉我们
$$
(f_*\mathcal{O}_X)_s^\wedge = \lim_n H^0(X_n, \mathcal{O}_{X_n})
$$
，其中 $X_n$ 是 $X_s$ 的第 $n$ 个无穷小邻域。
由于 $X_n$ 的底拓扑空间等于 $X_s$ 的底拓扑空间，
对所有 $n $ 都得到概形的不交并分解
$X_n = T_{1, n} \amalg T_{2, n}$，
其中当 $i = 1, 2$ 时，$T_{i, n}$ 的底拓扑空间是 $T_i$。
这意味着 $H^0(X_n, \mathcal{O}_{X_n})$
含有非平凡幂等元 $e_n$，即在 $T_{1, n}$ 上恒等于 $1$、
在 $T_{2, n}$ 上恒等于 $0$ 的函数。显然，
$e_{n + 1}$ 在 $X_n$ 上限制为 $e_n$。
所以 $e_\infty = \lim e_n$ 是该极限的非平凡幂等元。
于是 $e_\infty$ 是 $(f_*\mathcal{O}_X)_s$ 的完备化中的元素，
并映到 $H^0(X_s, \mathcal{O}_{X_s})$ 中的 $e$。
由于映射
$(f_*\mathcal{O}_X)_s^\wedge \to H^0(X_s, \mathcal{O}_{X_s})$
经过
$(f_*\mathcal{O}_X)^\wedge_s / \mathfrak m_s (f_*\mathcal{O}_X)_s^\wedge =
(f_*\mathcal{O}_X)_s / \mathfrak m_s (f_*\mathcal{O}_X)_s$
分解（代数，引理 \ref{algebra-lemma-hathat-finitely-generated}），
可知 $e$ 属于映射
$(f_*\mathcal{O}_X)_s \to H^0(X_s, \mathcal{O}_{X_s})$
的像，正如所需。

\medskip\noindent
一般情形：利用极限论证把一般情形归结到诺特情形。
建议读者跳过这段证明。
可以把 $S$ 替换为 $s$ 的一个仿射开邻域。
所以可以且确实假设 $S$ 仿射。由极限，引理
\ref{limits-lemma-proper-limit-of-proper-finite-presentation-noetherian}
，可以写成
$(f : X \to S) = \lim (f_i : X_i \to S_i)$，
其中 $f_i$ 固有且 $S_i$ 诺特。以 $s_i \in S_i$ 表示 $s$ 的像。
则 $s = \lim s_i$，见极限，引理
\ref{limits-lemma-inverse-limit-irreducibles}。
又因为极限与极限可交换（范畴，引理
\ref{categories-lemma-colimits-commute}），有
$X_s = X \times_S s = \lim X_i \times_{S_i} s_i = \lim X_{i, s_i}$。
因此由极限，引理 \ref{limits-lemma-descend-section}，
$e$ 是某个幂等元
$e_i \in H^0(X_{i, s_i}, \mathcal{O}_{X_{i, s_i}})$
的像。由诺特情形，茎
$(f_{i, *}\mathcal{O}_{X_i})_{s_i}$ 中存在元素 $\tilde e_i$
映到 $e_i$。拉回 $\tilde e_i$，得到
$(f_*\mathcal{O}_X)_s$ 中的元素 $\tilde e$，
它映到 $e$，证明完毕。
\end{proof}

\begin{lemma}
\label{perfect-lemma-proper-flat-geom-red}
设 $f : X \to S$ 是概形态射，$s \in S$。假设
\begin{enumerate}
\item $f$ 固有、平坦且有限表示；
\item 纤维 $X_s$ 几何既约。
\end{enumerate}
则把 $S$ 替换为 $s$ 的某个开邻域后，存在直和分解
$Rf_*\mathcal{O}_X = f_*\mathcal{O}_X \oplus P$
作为 $D(\mathcal{O}_S)$ 中的等式，其中 $f_*\mathcal{O}_X$
是有限平展 $\mathcal{O}_S$-代数，而
$P$ 是 Tor 振幅包含于 $[1, \infty)$ 的完美复形。
\end{lemma}

\begin{proof}
本引理的证明与引理 \ref{perfect-lemma-proper-flat-h0} 的证明相似，
建议读者先阅读后者。
由上同调与基变换
（引理 \ref{perfect-lemma-flat-proper-perfect-direct-image-general}），
复形 $E = Rf_*\mathcal{O}_X$ 完美，
且其构造与任意基变换相容。这首先说明
$E$ 的 Tor 振幅包含于 $[0, \infty)$。

\medskip\noindent
我们声称，把 $S$ 替换为 $s$ 的某个开邻域后，
可以在 $D(\mathcal{O}_S)$ 中找到直和分解
$E = H^0(E) \oplus \tau_{\geq 1}E$，
其中 $\tau_{\geq 1}E$ 的 Tor 振幅包含于 $[1, \infty)$。
暂且假设该断言成立，并假设已经作了这一替换，
因而具有上述直和分解。由于 $E$ 的 Tor 振幅包含于
$[0, \infty)$，$H^0(E)$ 是平坦 $\mathcal{O}_S$-模。
所以 $H^0(E)$ 是平坦且完美的 $\mathcal{O}_S$-模，
从而有限局部自由，见代数续篇，引理
\ref{more-algebra-lemma-perfect}
（以及代数，引理 \ref{algebra-lemma-finite-projective}
所述的有限投射模为有限局部自由模这一事实）。
当然，$H^0(E) = f_*\mathcal{O}_X$ 是 $\mathcal{O}_S$-代数。
由上同调与基变换，得到
$H^0(E) \otimes \kappa(s) = H^0(X_s, \mathcal{O}_{X_s})$.
由簇，引理
\ref{varieties-lemma-proper-geometrically-reduced-global-sections}
与 $X_s$ 几何既约这一假设可知，
$\kappa(s) \to H^0(E) \otimes \kappa(s)$ 是有限平展的。
对有限局部自由态射
$\underline{\Spec}_S(H^0(E)) \to S$
应用态射，引理
\ref{morphisms-lemma-set-points-where-fibres-etale}
可知，缩小 $S$ 后，$\mathcal{O}_S$-代数
$H^0(E)$ 是有限平展的。

\medskip\noindent
还需证明该断言。为此，只需证明映射
$$
(f_*\mathcal{O}_X)_s
\longrightarrow
H^0(X_s, \mathcal{O}_{X_s}) = H^0(E \otimes^\mathbf{L} \kappa(s))
$$
是满射，见代数续篇，引理
\ref{more-algebra-lemma-better-cut-complex-in-two}。
选取平坦局部环同态 $\mathcal{O}_{S, s} \to A$，
使 $A$ 的剩余域 $k$ 代数闭，见代数，引理
\ref{algebra-lemma-flat-local-given-residue-field}。
由平坦基变换（概形的上同调，引理
\ref{coherent-lemma-flat-base-change-cohomology})
，得到 $H^0(X_A, \mathcal{O}_{X_A}) =
(f_*\mathcal{O}_X)_s \otimes_{\mathcal{O}_{S, s}} A$
与 $H^0(X_k, \mathcal{O}_{X_k}) =
H^0(X_s, \mathcal{O}_{X_s}) \otimes_{\kappa(s)} k$.
所以只需证明
$H^0(X_A, \mathcal{O}_{X_A}) \to H^0(X_k, \mathcal{O}_{X_k})$
是满射。由于 $X_k$ 是 $k$ 上的既约固有概形，
且 $k$ 代数闭，由上面已经使用过的簇，引理
\ref{varieties-lemma-proper-geometrically-reduced-global-sections}，
$H^0(X_k, \mathcal{O}_{X_k})$ 是有限多个 $k$ 的乘积。
又由引理 \ref{perfect-lemma-proper-idempotent-on-fibre}，
这个 $k$-代数的幂等元都属于映射
$H^0(X_A, \mathcal{O}_{X_A}) \to H^0(X_k, \mathcal{O}_{X_k})$
的像，故结论成立。
\end{proof}





\section{其他应用}
\label{perfect-section-other-applications}

\noindent
本节陈述并证明若干可由上面建立的理论推出的结果。

\begin{lemma}
\label{perfect-lemma-cohomology-over-coherent-ring}
设 $R$ 是凝聚环，$X$ 是在 $R$ 上有限表示的概形。
设 $\mathcal{G}$ 是有限表示的 $\mathcal{O}_X$-模，
它在 $R$ 上平坦且支集在 $R$ 上固有。则
$H^i(X, \mathcal{G})$ 是凝聚 $R$-模。
\end{lemma}

\begin{proof}
结合引理 \ref{perfect-lemma-base-change-tensor-perfect}
与代数续篇，引理 \ref{more-algebra-lemma-coherent-pseudo-coherent}
及 \ref{more-algebra-lemma-perfect}。
\end{proof}

\begin{lemma}
\label{perfect-lemma-countable-cohomology}
设 $X$ 是拟紧拟分离概形。设 $K$ 为
$D_\QCoh(\mathcal{O}_X)$ 的对象，并且上同调层 $H^i(K)$
在仿射开集上的截面集均为可数集。则对任意拟紧开集
$U \subset X$ 以及 $D(\mathcal{O}_X)$ 中任意完美对象 $E$，集合
$$
H^i(U, K \otimes^\mathbf{L} E),\quad \Ext^i(E|_U, K|_U)
$$
都是可数的。
\end{lemma}

\begin{proof}
利用上同调，引理 \ref{cohomology-lemma-dual-perfect-complex}，
只需对群 $H^i(U, K \otimes^\mathbf{L} E)$ 证明结论。
我们将使用归纳原理证明本引理，见概形的上同调，引理
\ref{coherent-lemma-induction-principle}。

\medskip\noindent
先证明 $U = \Spec(A)$ 仿射时结论成立。可以用
$A$-模复形 $K^\bullet$ 表示 $K$，并用有限生成投射
$A$-模的有限复形 $P^\bullet$ 表示 $E$。见引理
\ref{perfect-lemma-affine-compare-bounded}、
\ref{perfect-lemma-perfect-affine}
以及我们对完美 $A$-模复形的定义
（代数续篇，定义 \ref{more-algebra-definition-perfect}）。
于是 $(E \otimes^\mathbf{L} K)|_U$ 由双复形
$P^\bullet \otimes_A K^\bullet$
相伴的全复形表示
（引理 \ref{perfect-lemma-quasi-coherence-tensor-product}）。
对复形 $P^\bullet$ 的长度作归纳
（或使用适当的谱序列），可知只需证明
对每个 $a$，$H^i(P^a \otimes_A K^\bullet)$ 都可数。
由于对某个 $n$，$P^a$ 是 $A^{\oplus n}$ 的直和因子，
这由上同调群 $H^i(K^\bullet)$ 可数这一假设得到。

\medskip\noindent
为完成证明，只需说明：若 $U = V \cup W$，
且结论对 $V$、$W$ 与 $V \cap W$ 成立，则它也对 $U$ 成立。
这是 Mayer--Vietoris 序列的直接推论，见上同调，引理
\ref{cohomology-lemma-unbounded-mayer-vietoris}。
\end{proof}

\begin{lemma}
\label{perfect-lemma-countable}
设 $X$ 是拟紧拟分离概形，并且 $\mathcal{O}_X$
在仿射开集上的截面集均为可数集。
设 $K$ 为 $D_\QCoh(\mathcal{O}_X)$ 的对象。下列条件等价：
\begin{enumerate}
\item $K = \text{hocolim} E_n$，其中 $E_n$ 是
$D(\mathcal{O}_X)$ 的完美对象；
\item 上同调层 $H^i(K)$ 在仿射开集上的截面集均为可数集。
\end{enumerate}
\end{lemma}

\begin{proof}
若 (1) 成立，则 (2) 成立，因为取同伦余极限与取上同调层相容
（由导出范畴，引理 \ref{derived-lemma-cohomology-of-hocolim}），
并且完美复形局部同构于有限自由 $\mathcal{O}_X$-模的有限复形，
故由对 $X$ 的假设满足 (2)。

\medskip\noindent
假设 (2)。选取表示 $K$ 的 K-内射复形 $\mathcal{K}^\bullet$。
选取 $D_\QCoh(\mathcal{O}_X)$ 的完美生成元 $E$，
并用 K-内射复形 $\mathcal{I}^\bullet$ 表示它。
根据定理 \ref{perfect-theorem-DQCoh-is-Ddga} 及其证明，
存在三角范畴的等价
$F : D_\QCoh(\mathcal{O}_X) \to D(A, \text{d})$，
其中 $(A, \text{d})$ 是微分分次代数
$$
(A, \text{d}) =
\Hom_{\text{Comp}^{dg}(\mathcal{O}_X)}
(\mathcal{I}^\bullet, \mathcal{I}^\bullet)
$$
，而该等价把 $K$ 送到微分分次模
$$
M = \Hom_{\text{Comp}^{dg}(\mathcal{O}_X)}
(\mathcal{I}^\bullet, \mathcal{K}^\bullet)
$$
注意，$H^i(A) = \Ext^i(E, E)$，并且
$H^i(M) = \Ext^i(E, K)$。
此外，由于 $F$ 是等价，它及其拟逆都与同伦余极限相容。
因此只需把 $M$ 写成 $D(A, \text{d})$ 的紧对象的同伦余极限。
由微分分次代数，引理 \ref{dga-lemma-countable}，
只需证明对每个 $i$，$\Ext^i(E, E)$ 与
$\Ext^i(E, K)$ 都可数。这由引理
\ref{perfect-lemma-countable-cohomology} 得到。
\end{proof}

\begin{lemma}
\label{perfect-lemma-computing-sections-as-colim}
设 $A$ 是环，$X$ 是在 $A$ 上有限表示的概形。
设 $f : U \to X$ 是有限表示的平坦态射。则：
\begin{enumerate}
\item 存在 $D(\mathcal{O}_X)$ 的完美对象 $L_n$ 的一个逆系统，使得
$$
R\Gamma(U, Lf^*K) = \text{hocolim}\ R\Hom_X(L_n, K)
$$
作为 $D(A)$ 中的等式，对 $D_\QCoh(\mathcal{O}_X)$ 中的 $K$
函子性地成立；
\item 存在 $D(\mathcal{O}_X)$ 的完美对象 $E_n$ 的一个系统，使得
$$
R\Gamma(U, Lf^*K) = \text{hocolim}\ R\Gamma(X, E_n \otimes^\mathbf{L} K)
$$
作为 $D(A)$ 中的等式，对 $D_\QCoh(\mathcal{O}_X)$ 中的 $K$
函子性地成立。
\end{enumerate}
\end{lemma}

\begin{proof}
由引理 \ref{perfect-lemma-cohomology-base-change}，有
$$
R\Gamma(U, Lf^*K) = R\Gamma(X, Rf_*\mathcal{O}_U \otimes^\mathbf{L} K)
$$
，并且对 $K$ 函子性地成立。注意，$R\Gamma(X, -)$
与同伦余极限相容，因为由引理
\ref{perfect-lemma-quasi-coherence-pushforward-direct-sums}，
它与直和相容。类似地，$- \otimes^\mathbf{L} K$
与导出余极限相容，因为 $- \otimes^\mathbf{L} K$ 与直和相容
（$D(\mathcal{O}_X)$ 中的直和由表示复形的直和给出）。
所以为证明 (2)，只需对 $D(\mathcal{O}_X)$ 的某个完美对象系统
$E_n$ 写出 $Rf_*\mathcal{O}_U = \text{hocolim} E_n$。
一旦做到这一点，置 $L_n = E_n^\vee$ 即得到 (1)，见上同调，
引理 \ref{cohomology-lemma-dual-perfect-complex}。

\medskip\noindent
写成 $A = \colim A_i$，其中 $A_i$ 在 $\mathbf{Z}$ 上有限型。
由极限，引理 \ref{limits-lemma-descend-finite-presentation}，
可以找到某个 $i$ 以及有限表示态射
$U_i \to X_i \to \Spec(A_i)$，其到 $\Spec(A)$ 的基变换恢复
$U \to X \to \Spec(A)$。
增大 $i$ 后，可以假设 $f_i : U_i \to X_i$ 平坦，见极限，
引理 \ref{limits-lemma-descend-flat-finite-presentation}。
由引理 \ref{perfect-lemma-compare-base-change}，
$Rf_{i, *}\mathcal{O}_{U_i}$ 沿 $g : X \to X_i$ 的导出拉回
等于 $Rf_*\mathcal{O}_U$。由于 $Lg^*$ 与导出余极限相容，
只需对 $f_i$ 证明所需结论。因此可以假设
$U$ 与 $X$ 都在 $\mathbf{Z}$ 上有限型。

\medskip\noindent
假设 $f : U \to X$ 是 $\mathbf{Z}$ 上有限型概形之间的态射。
为完成证明，我们将说明 $Rf_*\mathcal{O}_U$
是完美复形的同伦余极限。为此应用引理 \ref{perfect-lemma-countable}。
所以只需证明 $R^if_*\mathcal{O}_U$
在仿射开集上的截面集均为可数集。
这由把引理 \ref{perfect-lemma-countable-cohomology} 应用于结构层得到。
\end{proof}




\section{伪凝聚复形的刻画 II}
\label{perfect-section-pseudo-coherent}

\noindent
本节接续第 \ref{perfect-section-pseudo-coherent-hocolim} 节。
我们讨论用上同调刻画伪凝聚复形的方法。
更多这类结果见态射续篇，第
\ref{more-morphisms-section-characterize-pseudo-coherent} 节。

\begin{lemma}
\label{perfect-lemma-pseudo-coherent-over-algebra}
设 $A$ 是环，$R$ 是一个（可以非交换的）$A$-代数，
并且作为 $A$-模有限自由。则 $D(R)$ 中任何在 $D(A)$
里伪凝聚的对象 $M$，都可由有限自由（右）$R$-模的上有界复形表示。
\end{lemma}

\begin{proof}
选取表示 $M$ 的右 $R$-模复形 $M^\bullet$。
由于 $M$ 伪凝聚，对充分大的 $i$，有 $H^i(M) = 0$。
令 $m$ 为使 $H^m(M)$ 非零的最小指标。
由代数续篇，引理 \ref{more-algebra-lemma-finite-cohomology}，
$H^m(M)$ 是有限 $A$-模。
因此可以选取有限自由 $R$-模 $F^m$ 与映射
$F^m \to M^m$，使复合 $F^m \to M^m \to M^{m + 1}$ 为零，
并且 $F^m \to H^m(M)$ 满射。图示如下：
$$
\xymatrix{
& F^m \ar[d]^\alpha \ar[r] & 0 \ar[d] \ar[r] & \ldots \\
M^{m - 1} \ar[r] & M^m \ar[r] & M^{m + 1} \ar[r] & \ldots
}
$$
我们将对 $n \leq m$ 作降序归纳，构造：
当 $i \geq n$ 时的有限自由 $R$-模 $F^i$；
当 $i \geq n$ 时的微分 $d^i : F^i \to F^{i + 1}$；
以及与微分相容的映射 $\alpha : F^i \to K^i$，使得
(1) 当 $i > n$ 时 $H^i(\alpha)$ 是同构，而当 $i = n$ 时满射；
(2) 当 $i > m$ 时 $F^i = 0$。图示为
$$
\xymatrix{
& F^n \ar[r] \ar[d]^\alpha & F^{n + 1} \ar[d]^\alpha \ar[r] & \ldots \ar[r]
& F^i \ar[d]^\alpha \ar[r] & 0 \ar[d] \ar[r] & \ldots \\
M^{n - 1} \ar[r] & M^n \ar[r] & M^{n + 1} \ar[r] & \ldots \ar[r] &
M^i \ar[r] & M^{i + 1} \ar[r] & \ldots
}
$$
基本情形 $n = m$ 已在上面完成。
归纳步骤：令 $C^\bullet$ 为 $\alpha$ 的锥
（导出范畴，定义 \ref{derived-definition-cone}）。
上同调长正合列说明，当 $i \geq n$ 时，
$H^i(C^\bullet) = 0$。注意，因为 $R$ 作为 $A$-模有限自由，
$F^\bullet$ 作为 $A$-模复形是伪凝聚的。
所以由代数续篇，引理
\ref{more-algebra-lemma-cone-pseudo-coherent}，
$C^\bullet$ 作为 $A$-模复形是 $(n - 1)$-伪凝聚的。
由代数续篇，引理 \ref{more-algebra-lemma-finite-cohomology}，
$H^{n - 1}(C^\bullet)$ 是有限 $A$-模。
选取有限自由 $R$-模 $F^{n - 1}$ 与映射
$\beta : F^{n - 1} \to C^{n - 1}$，使复合
$F^{n - 1} \to C^{n - 1} \to C^n$ 为零，
并且 $F^{n - 1}$ 满射到 $H^{n - 1}(C^\bullet)$。
由于 $C^{n - 1} = M^{n - 1} \oplus F^n$，
可写成 $\beta = (\alpha^{n - 1}, -d^{n - 1})$。
复合 $F^{n - 1} \to C^{n - 1} \to C^n$ 为零说明
这些映射组成复形态射
$$
\xymatrix{
& F^{n - 1} \ar[d]^{\alpha^{n - 1}} \ar[r]_{d^{n - 1}} &
F^n \ar[r] \ar[d]^\alpha &
F^{n + 1} \ar[d]^\alpha \ar[r] & \ldots \\
\ldots \ar[r] &
M^{n - 1} \ar[r] & M^n \ar[r] & M^{n + 1} \ar[r] & \ldots
}
$$
此外，这些映射定义可区别三角之间的态射
$$
\xymatrix{
(F^n \to \ldots) \ar[r] \ar[d] &
(F^{n - 1} \to \ldots) \ar[r] \ar[d] &
F^{n - 1} \ar[r] \ar[d]_\beta &
(F^n \to \ldots)[1] \ar[d] \\
(F^n \to \ldots) \ar[r] &
M^\bullet \ar[r] &
C^\bullet \ar[r] &
(F^n \to \ldots)[1]
}
$$
因此，对 $\beta$ 的选择说明复形映射
$(F^{n - 1} \to \ldots) \to M^\bullet$
在次数 $\geq n$ 的上同调上诱导同构，并在次数 $n - 1$ 上诱导满射。
这完成了本引理的证明。
\end{proof}

\begin{lemma}
\label{perfect-lemma-pseudo-coherent-on-projective-space}
设 $A$ 是环，$n \geq 0$，并设
$K \in D_\QCoh(\mathcal{O}_{\mathbf{P}^n_A})$。
下列条件等价：
\begin{enumerate}
\item $K$ 伪凝聚；
\item 对 $D(\mathcal{O}_{\mathbf{P}^n_A})$ 的每个伪凝聚对象 $E$，
$R\Gamma(\mathbf{P}^n_A, E \otimes^\mathbf{L} K)$
都是 $D(A)$ 的伪凝聚对象；
\item 对 $D(\mathcal{O}_{\mathbf{P}^n_A})$ 的每个完美对象 $E$，
$R\Gamma(\mathbf{P}^n_A, E \otimes^\mathbf{L} K)$
都是 $D(A)$ 的伪凝聚对象；
\item 对 $D(\mathcal{O}_{\mathbf{P}^n_A})$ 的每个完美对象 $E$，
$R\Hom_{\mathbf{P}^n_A}(E, K)$ 都是 $D(A)$ 的伪凝聚对象；
\item $R\Gamma(\mathbf{P}^n_A,
K \otimes^\mathbf{L} \mathcal{O}_{\mathbf{P}^n_A}(d))$
对 $d = 0, 1, \ldots, n$ 都是 $D(A)$ 的伪凝聚对象。
\end{enumerate}
\end{lemma}

\begin{proof}
回忆
$$
R\Hom_{\mathbf{P}^n_A}(E, K) =
R\Gamma(\mathbf{P}^n_A, R\SheafHom_{\mathcal{O}_{\mathbf{P}^n_A}}(E, K))
$$
按定义成立，见上同调，第 \ref{cohomology-section-global-RHom} 节。
因此由上同调，引理 \ref{cohomology-lemma-dual-perfect-complex}，
(4) 与 (3) 等价。

\medskip\noindent
由于每个完美复形都是伪凝聚的，显然 (2) 蕴含 (3)。

\medskip\noindent
假设 (1) 成立。则对每个伪凝聚的 $E$，
$E \otimes^\mathbf{L} K$ 伪凝聚，见上同调，引理
\ref{cohomology-lemma-tensor-pseudo-coherent}。
由引理 \ref{perfect-lemma-flat-proper-pseudo-coherent-direct-image-general}，
这种伪凝聚复形的直像仍伪凝聚，故 (2) 成立。

\medskip\noindent
(3) 蕴含 (5)，因为可以对 $d = 0, 1, \ldots, n$
取 $E = \mathcal{O}_{\mathbf{P}^n_A}(d)$。

\medskip\noindent
为完成证明，还需说明 (5) 蕴含 (1)。
令 $P$ 如 (\ref{perfect-equation-generator-Pn})，并令
$R$ 如 (\ref{perfect-equation-algebra-for-Pn})。
由引理 \ref{perfect-lemma-Pn-module-category}，有等价
$$
- \otimes^\mathbf{L}_R P :
D(R) \longrightarrow D_\QCoh(\mathcal{O}_{\mathbf{P}^n_A})
$$
令 $M \in D(R)$ 为满足 $M \otimes^\mathbf{L} P = K$ 的对象。
由微分分次代数，引理
\ref{dga-lemma-upgrade-tensor-with-complex-derived}
，有同构
$$
R\Hom(R, M) = R\Hom_{\mathbf{P}^n_A}(P, K)
$$
作为 $D(A)$ 中的同构。与上面一样论证，得到
$$
R\Hom_{\mathbf{P}^n_A}(P, K) = R\Gamma(\mathbf{P}^n_A,
R\SheafHom_{\mathcal{O}_{\mathbf{P}^n_A}}(E, K)) =
R\Gamma(\mathbf{P}^n_A,
P^\vee \otimes^\mathbf{L}_{\mathcal{O}_{\mathbf{P}^n_A}} K).
$$
利用 $P^\vee$ 是各
$\mathcal{O}_{\mathbf{P}^n_A}(d)$（$d = 0, 1, \ldots, n$）
的直和以及 (5)，可知 $R\Hom(R, M)$ 作为 $A$-模复形伪凝聚。
当然，在 $D(A)$ 中 $M = R\Hom(R, M)$。
所以 $M$ 作为 $A$-模复形伪凝聚。
由引理 \ref{perfect-lemma-pseudo-coherent-over-algebra}，
可以用有限自由 $R$-模的上有界复形 $F^\bullet$ 表示 $M$。于是
$F^\bullet = \bigcup_{p \geq 0} \sigma_{\geq p}F^\bullet$
是一个滤过，它说明 $F^\bullet$ 是具有性质 (P)
的微分分次 $R$-模，见微分分次代数，第
\ref{dga-section-P-resolutions} 节。
因此 $K = M \otimes^\mathbf{L}_R P$ 由
$F^\bullet \otimes_R P$ 表示
（这由导出张量函子的构造得到；例如见微分分次代数，
引理 \ref{dga-lemma-tensor-with-complex-derived} 的证明）。
由于 $F^\bullet \otimes_R P$ 是上有界复形，
其各项均为若干个 $P$ 的直和，本引理得证。
\end{proof}

\begin{lemma}
\label{perfect-lemma-perfect-enough}
设 $A$ 是环，$X$ 是拟紧拟分离的 $A$-概形。
设 $K \in D^-_\QCoh(\mathcal{O}_X)$。
若对 $D(\mathcal{O}_X)$ 中每个完美的 $E$，
$R\Gamma(X, E \otimes^\mathbf{L} K)$ 在 $D(A)$ 中伪凝聚，
则对 $D(\mathcal{O}_X)$ 中每个伪凝聚的 $E$，
$R\Gamma(X, E \otimes^\mathbf{L} K)$ 在 $D(A)$ 中伪凝聚。
\end{lemma}

\begin{proof}
存在整数 $N$，使得
$R\Gamma(X, -) : D_\QCoh(\mathcal{O}_X) \to D(A)$
具有上同调维数 $N$，见引理
\ref{perfect-lemma-quasi-coherence-direct-image} 中的说明。
取 $b \in \mathbf{Z}$，使得当 $i > b$ 时 $H^i(K) = 0$。
令 $E$ 为 $X$ 上的伪凝聚对象。
只需证明对每个 $m$，
$R\Gamma(X, E \otimes^\mathbf{L} K)$ 都是 $m$-伪凝聚的。
选取由完美复形 $P$ 给出的 $(X, E, m - N - 1 - b)$
的一个逼近 $P \to E$。其存在性由
定理 \ref{perfect-theorem-approximation} 保证。
在 $D_\QCoh(\mathcal{O}_X)$ 中选取可区别三角
$$
P \to E \to C \to P[1]
$$
。$C$ 的上同调层在次数 $\geq m - N - 1 - b$ 均为零。
因此 $C \otimes^\mathbf{L} K$ 的上同调层
在次数 $\geq m - N - 1$ 均为零。
所以 $R\Gamma(X, C \otimes^\mathbf{L} K)$ 的上同调
在次数 $\geq m - 1$ 均为零。于是
$$
R\Gamma(X, P \otimes^\mathbf{L} K) \to R\Gamma(X, E \otimes^\mathbf{L} K)
$$
在次数 $\geq m$ 的上同调上是同构。
由假设，其源伪凝聚。
因此 $R\Gamma(X, E \otimes^\mathbf{L} K)$
如所需是 $m$-伪凝聚的。
\end{proof}









\section{相对完美对象}
\label{perfect-section-relatively-perfect}

\noindent
本节引入 \cite{lieblich-complexes} 中的一个概念。

\begin{definition}
\label{perfect-definition-relatively-perfect}
设 $f : X \to S$ 是平坦且局部有限表示的概形态射。
若 $D(\mathcal{O}_X)$ 的对象 $E$ 伪凝聚
（上同调，定义 \ref{cohomology-definition-pseudo-coherent}），
并且 $E$ 作为
$D(f^{-1}\mathcal{O}_S)$
的对象局部具有有限 Tor 维数
（上同调，定义 \ref{cohomology-definition-tor-amplitude}），
则称 $E$ {\it 相对于 $S$ 完美}，或称其为
{\it $S$-完美的}。
\end{definition}

\noindent
相关讨论见评注 \ref{perfect-remark-discuss-rel-perfect}。

\begin{example}
\label{perfect-example-relatively-perfect-field}
设 $k$ 是域，$X$ 是在 $k$ 上有限表示的概形
（特别地，$X$ 拟紧）。则 $D(\mathcal{O}_X)$ 的对象 $E$
是 $k$-完美的，当且仅当它有界且伪凝聚（由定义）；
也就是说，当且仅当它属于 $D^b_{\textit{Coh}}(X)$
（由引理 \ref{perfect-lemma-identify-pseudo-coherent-noetherian}）。
所以，相对完美{\bf 不}意味着“在纤维上完美”。
\end{example}

\noindent
相应的代数概念见代数续篇，第
\ref{more-algebra-section-relatively-perfect} 节。
概形上的概念与代数概念可如下联系。

\begin{lemma}
\label{perfect-lemma-affine-locally-rel-perfect}
设 $f : X \to S$ 是平坦且局部有限表示的概形态射。
设 $E$ 为 $D_\QCoh(\mathcal{O}_X)$ 的对象。下列条件等价：
\begin{enumerate}
\item $E$ 是 $S$-完美的；
\item 对任意映入仿射开集 $V \subset S$ 的仿射开集
$U \subset X$，复形 $R\Gamma(U, E)$ 是
$\mathcal{O}_S(V)$-完美的；
\item 存在仿射开覆盖 $S = \bigcup V_i$，并且对每个 $i$
存在仿射开覆盖 $f^{-1}(V_i) = \bigcup U_{ij}$，
使复形 $R\Gamma(U_{ij}, E)$ 是 $\mathcal{O}_S(V_i)$-完美的。
\end{enumerate}
\end{lemma}

\begin{proof}
伪凝聚性是局部性质，“局部具有有限 Tor 维数”也是局部性质。
因此本引理立即归结为下述断言：
若 $X$ 与 $S$ 仿射，则 $E$ 是 $S$-完美的，
当且仅当 $K = R\Gamma(X, E)$ 是 $\mathcal{O}_S(S)$-完美的。
记 $X = \Spec(A)$、$S = \Spec(R)$，并设 $E$ 对应于
$K \in D(A)$；也就是说，$K = R\Gamma(X, E)$，见引理
\ref{perfect-lemma-affine-compare-bounded}。

\medskip\noindent
注意，$K$ 是 $R$-完美的，当且仅当 $K$
作为 $R$-模复形伪凝聚且具有有限 Tor 维数
（代数续篇，定义 \ref{more-algebra-definition-relatively-perfect}）。
由引理 \ref{perfect-lemma-pseudo-coherent-affine}，
$E$ 伪凝聚当且仅当 $K$ 伪凝聚。
由引理 \ref{perfect-lemma-tor-dimension-rel-affine}，
$E$ 在 $f^{-1}\mathcal{O}_S$ 上具有有限 Tor 维数，
当且仅当 $K$ 作为 $R$-模复形具有有限 Tor 维数。
\end{proof}

\begin{lemma}
\label{perfect-lemma-triangulated}
设 $f : X \to S$ 是平坦且局部有限表示的概形态射。
$D(\mathcal{O}_X)$ 中由 $S$-完美对象组成的全子范畴
是饱和的\footnote{导出范畴，定义
\ref{derived-definition-saturated}。}三角子范畴。
\end{lemma}

\begin{proof}
这由上同调，诸引理
\ref{cohomology-lemma-cone-pseudo-coherent},
\ref{cohomology-lemma-summands-pseudo-coherent},
\ref{cohomology-lemma-cone-tor-amplitude} 与
\ref{cohomology-lemma-summands-tor-amplitude} 得到。
\end{proof}

\begin{lemma}
\label{perfect-lemma-perfect-relatively-perfect}
设 $f : X \to S$ 是平坦且局部有限表示的概形态射。
$D(\mathcal{O}_X)$ 的完美对象是 $S$-完美的。
若 $K, M \in D(\mathcal{O}_X)$，且 $K$ 完美而 $M$ 是
$S$-完美的，则
$K \otimes_{\mathcal{O}_X}^\mathbf{L} M$ 是 $S$-完美的。
\end{lemma}

\begin{proof}
第一种证明：利用引理 \ref{perfect-lemma-affine-locally-rel-perfect}
归结到仿射情形，再应用代数续篇，引理
\ref{more-algebra-lemma-perfect-relatively-perfect}。
\end{proof}

\begin{lemma}
\label{perfect-lemma-base-change-relatively-perfect}
设 $f : X \to S$ 是平坦且局部有限表示的概形态射。
设 $g : S' \to S$ 是概形态射。置 $X' = S' \times_S X$，
并以 $g' : X' \to X$ 表示投影。
若 $K \in D(\mathcal{O}_X)$ 是 $S$-完美的，则
$L(g')^*K$ 是 $S'$-完美的。
\end{lemma}

\begin{proof}
第一种证明：利用引理 \ref{perfect-lemma-affine-locally-rel-perfect}
归结到仿射情形，再应用代数续篇，引理
\ref{more-algebra-lemma-base-change-relatively-perfect}。

\medskip\noindent
第二种证明：由上同调，引理
\ref{cohomology-lemma-pseudo-coherent-pullback}，
$L(g')^*K$ 伪凝聚；有限 Tor 维数性质则由
引理 \ref{perfect-lemma-tor-independence-and-tor-amplitude} 得到。
\end{proof}

\begin{situation}
\label{perfect-situation-relative-descent}
设 $S = \lim_{i \in I} S_i$ 是概形有向系统的极限，
其转移态射 $g_{i'i} : S_{i'} \to S_i$ 均为仿射态射。
假设对所有 $i \in I$，$S_i$ 都拟紧且拟分离。
以 $g_i : S \to S_i$ 表示投影。固定元素 $0 \in I$
以及有限表示的平坦态射 $X_0 \to S_0$。
置 $X_i = S_i \times_{S_0} X_0$ 与 $X = S \times_{S_0} X_0$，
并以 $f_{i'i} : X_{i'} \to X_i$ 表示转移态射，
以 $f_i : X \to X_i$ 表示投影。
\end{situation}

\begin{lemma}
\label{perfect-lemma-relative-descend-homomorphisms}
在情形 \ref{perfect-situation-relative-descent} 中，
设 $K_0$ 与 $L_0$ 为 $D(\mathcal{O}_{X_0})$ 的对象。
当 $i \geq 0$ 时置
$K_i = Lf_{i0}^*K_0$、$L_i = Lf_{i0}^*L_0$，
并置 $K = Lf_0^*K_0$、$L = Lf_0^*L_0$。则映射
$$
\colim_{i \geq 0} \Hom_{D(\mathcal{O}_{X_i})}(K_i, L_i)
\longrightarrow
\Hom_{D(\mathcal{O}_X)}(K, L)
$$
在下列条件下是同构：$K_0$ 伪凝聚，并且
$L_0 \in D_\QCoh(\mathcal{O}_{X_0})$ 作为
$D((X_0 \to S_0)^{-1}\mathcal{O}_{S_0})$
的对象（局部）具有有限 Tor 维数。
\end{lemma}

\begin{proof}
对每个拟紧开集 $U_0 \subset X_0$，考虑条件 $P$：
$$
\colim_{i \geq 0} \Hom_{D(\mathcal{O}_{U_i})}(K_i|_{U_i}, L_i|_{U_i})
\longrightarrow
\Hom_{D(\mathcal{O}_U)}(K|_U, L|_U)
$$
是同构，其中 $U = f_0^{-1}(U_0)$ 且 $U_i = f_{i0}^{-1}(U_0)$。
若 $P$ 对 $U_0$、$V_0$ 与 $U_0 \cap V_0$ 成立，
则由导出范畴中的 Hom 的 Mayer--Vietoris 性质，
它也对 $U_0 \cup V_0$ 成立；见上同调，引理
\ref{cohomology-lemma-mayer-vietoris-hom}。

\medskip\noindent
以 $\pi_0 : X_0 \to S_0$ 表示给定态射。
先考虑 $U_0 = \pi_0^{-1}(W_0)$，其中
$W_0 \subset S_0$ 是拟紧开集。把概形的上同调，引理
\ref{coherent-lemma-induction-principle} 的归纳原理
应用于 $S_0$ 的拟紧开集，并结合上面的观察，
可知只需对 $U_0 = \pi_0^{-1}(W_0)$ 且 $W_0$ 仿射证明 $P$。
换言之，已经归结到 $S_0$ 仿射的情形。
接着再次应用归纳原理，这一次作用于 $X_0$
的所有拟紧拟分离开集，从而进一步归结到 $X_0$ 也仿射的情形。

\medskip\noindent
若 $X_0$ 与 $S_0$ 仿射，结论由代数续篇，引理
\ref{more-algebra-lemma-colimit-relatively-perfect} 得到。
事实上，由引理 \ref{perfect-lemma-pseudo-coherent} 与
\ref{perfect-lemma-affine-compare-bounded}
，该断言转化为模的导出范畴中的 Hom 计算。
引理 \ref{perfect-lemma-pseudo-coherent-affine} 说明
与 $K_0$ 对应的模复形伪凝聚，而引理
\ref{perfect-lemma-tor-dimension-rel-affine} 说明
与 $L_0$ 对应的模复形在 $\mathcal{O}_{S_0}(S_0)$
上具有有限 Tor 维数。因此代数续篇，引理
\ref{more-algebra-lemma-colimit-relatively-perfect}
的假设均满足，结论得证。
\end{proof}

\begin{lemma}
\label{perfect-lemma-descend-relatively-perfect}
在情形 \ref{perfect-situation-relative-descent} 中，
$D(\mathcal{O}_X)$ 的 $S$-完美对象范畴，是各
$D(\mathcal{O}_{X_i})$ 的 $S_i$-完美对象范畴的余极限。
\end{lemma}

\begin{proof}
对每个拟紧开集 $U_0 \subset X_0$，考虑条件 $P$：
函子
$$
\colim_{i \geq 0} D_{S_i\text{-完美}}(\mathcal{O}_{U_i})
\longrightarrow
D_{S\text{-完美}}(\mathcal{O}_U)
$$
是等价，其中 $U = f_0^{-1}(U_0)$ 且 $U_i = f_{i0}^{-1}(U_0)$。
由引理 \ref{perfect-lemma-relative-descend-homomorphisms}，
已经知道该函子全忠实。因此只需证明本质满性。

\medskip\noindent
设 $P$ 对 $X_0$ 的拟紧开集 $U_0$、$V_0$ 成立。
我们声称 $P$ 对 $U_0 \cup V_0$ 也成立。使用记号
$U_i = f_{i0}^{-1}U_0$、$U = f_0^{-1}U_0$、
$V_i = f_{i0}^{-1}V_0$ 与 $V = f_0^{-1}V_0$；
并且略微滥用记号，以 $f_i$ 统一表示态射
$U \to U_i$、$V \to V_i$、
$U \cap V \to U_i \cap V_i$ 与
$U \cup V \to U_i \cup V_i$。
设 $E$ 是 $D(\mathcal{O}_{U \cup V})$ 的 $S$-完美对象。
目标是证明 $E$ 属于该函子的本质像。
由假设，可以找到 $i \geq 0$、$U_i$ 上的 $S_i$-完美对象
$E_{U, i}$、$V_i$ 上的 $S_i$-完美对象 $E_{V, i}$，
以及同构 $Lf_i^*E_{U, i} \to E|_U$ 与
$Lf_i^*E_{V, i} \to E|_V$。令
$$
a : E_{U, i} \to (Rf_{i, *}E)|_{U_i}
\quad\text{且}\quad
b : E_{V, i} \to (Rf_{i, *}E)|_{V_i}
$$
为上述同构 $Lf_i^*E_{U, i} \to E|_U$
与 $Lf_i^*E_{V, i} \to E|_V$ 的伴随映射。
由全忠实性，增大 $i$ 后，可以找到同构
$c : E_{U, i}|_{U_i \cap V_i} \to E_{V, i}|_{U_i \cap V_i}$
，其拉回为复合识别
$$
Lf_i^*E_{U, i}|_{U \cap V} \to E|_{U \cap V} \to Lf_i^*E_{V, i}|_{U \cap V}.
$$
应用上同调，引理 \ref{cohomology-lemma-glue}，
得到 $U_i \cup V_i$ 上的对象 $E_i$
以及映射 $d : E_i \to Rf_{i, *}E$，
其在 $U_i$ 与 $V_i$ 上分别限制为 $a$ 与 $b$。
显然 $E_i$ 是 $S_i$-完美的
（因为相对完美性是局部性质），且
$d$ 是某个同构 $Lf_i^*E_i \to E$ 的伴随映射。

\medskip\noindent
采用引理 \ref{perfect-lemma-relative-descend-homomorphisms}
证明中完全相同的论证，并使用归纳原理
（概形的上同调，引理 \ref{coherent-lemma-induction-principle}），
可归结到 $X_0$ 与 $S_0$ 都仿射的情形。
（先处理 $S_0$ 的开集以归结到 $S_0$ 仿射，
再处理 $X_0$ 的开集以归结到 $X_0$ 仿射。）
在仿射情形中，结论由代数续篇，引理
\ref{more-algebra-lemma-colimit-relatively-perfect} 得到。
到代数问题的转化由引理
\ref{perfect-lemma-affine-locally-rel-perfect} 完成。
\end{proof}

\begin{lemma}
\label{perfect-lemma-derived-pushforward-rel-perfect}
设 $f : X \to S$ 是有限表示的平坦固有概形态射。
设 $E \in D(\mathcal{O}_X)$ 是 $S$-完美的。
则 $Rf_*E$ 是 $D(\mathcal{O}_S)$ 的完美对象，
且其构造与任意基变换相容。
\end{lemma}

\begin{proof}
关于基变换的断言就是引理 \ref{perfect-lemma-compare-base-change}。
所以只需证明 $Rf_*E$ 是完美对象。
我们将用极限论证归结到 $S$ 是诺特仿射概形的情形。

\medskip\noindent
问题在 $S$ 上是局部的，所以可以假设 $S$ 仿射。
记 $S = \Spec(R)$。把 $R = \colim R_i$
写成诺特环 $R_i$ 的滤过余极限。由极限，引理
\ref{limits-lemma-descend-finite-presentation}
，存在某个 $i$ 及一个在 $R_i$ 上有限表示的概形 $X_i$，
其到 $R$ 的基变换为 $X$。由极限，诸引理
\ref{limits-lemma-eventually-proper} 与
\ref{limits-lemma-descend-flat-finite-presentation}
，可以假设 $X_i$ 在 $R_i$ 上固有且平坦。
由引理 \ref{perfect-lemma-descend-relatively-perfect}，
可以假设存在 $D(\mathcal{O}_{X_i})$ 的 $R_i$-完美对象 $E_i$，
其到 $X$ 的拉回为 $E$。
把引理 \ref{perfect-lemma-perfect-direct-image}
应用于 $X_i \to \Spec(R_i)$ 与 $E_i$，
并利用已经证明的基变换性质，即得结论。
\end{proof}

\begin{lemma}
\label{perfect-lemma-compute-ext-rel-perfect}
设 $f : X \to S$ 是概形态射，并设
$E, K \in D(\mathcal{O}_X)$。假设
\begin{enumerate}
\item $S$ 拟紧且拟分离；
\item $f$ 固有、平坦且有限表示；
\item $E$ 是 $S$-完美的；
\item $K$ 伪凝聚。
\end{enumerate}
则存在伪凝聚对象 $L \in D(\mathcal{O}_S)$，使得
$$
Rf_*R\SheafHom(K, E) = R\SheafHom(L, \mathcal{O}_S)
$$
，且在任意基变换后仍成立：给定
$$
\vcenter{
\xymatrix{
X' \ar[r]_{g'} \ar[d]_{f'} &
X \ar[d]^f \\
S' \ar[r]^g &
S
}
}
\quad\quad
\begin{matrix}
\text{为 Cartesian，此时有 } \\
Rf'_*R\SheafHom(L(g')^*K, L(g')^*E) \\
= R\SheafHom(Lg^*L, \mathcal{O}_{S'})
\end{matrix}
$$
\end{lemma}

\begin{proof}
由于 $S$ 拟紧且拟分离，$X$ 也具有这些性质。
由引理 \ref{perfect-lemma-pseudo-coherent-hocolim}，可写成
$K = \text{hocolim} K_n$，其中 $K_n$ 完美，并且 $K_n \to K$
在截断 $\tau_{\geq -n}$ 上诱导同构。令 $K_n^\vee$
为对偶完美复形
（上同调，引理 \ref{cohomology-lemma-dual-perfect-complex}）。
由此得到完美对象的逆系统
$\ldots \to K_3^\vee \to K_2^\vee \to K_1^\vee$。
由引理 \ref{perfect-lemma-perfect-relatively-perfect}，
$K_n^\vee \otimes_{\mathcal{O}_X} E$ 是 $S$-完美的。
所以可以把引理 \ref{perfect-lemma-derived-pushforward-rel-perfect}
应用于 $K_n^\vee \otimes_{\mathcal{O}_X} E$，从而得到
$S$ 上完美复形的逆系统
$$
\ldots \to M_3 \to M_2 \to M_1
$$
，其中
$$
M_n = Rf_*(K_n^\vee \otimes_{\mathcal{O}_X}^\mathbf{L} E) =
Rf_*R\SheafHom(K_n, E)
$$
此外，这些复形的构造与任意基变换相容，即 $Lg^*M_n =
Rf'_*((L(g')^*K_n)^\vee \otimes_{\mathcal{O}_{X'}}^\mathbf{L} L(g')^*E) =
Rf'_*R\SheafHom(L(g')^*K_n, L(g')^*E)$.

\medskip\noindent
由于 $K_n \to K$ 在 $\tau_{\geq -n}$ 上诱导同构，
$K_n \to K_{n + 1}$ 也在 $\tau_{\geq -n}$ 上诱导同构。
又因为 $K_n^\vee = R\SheafHom(K_n, \mathcal{O}_X)$，
$K_{n + 1}^\vee \to K_n^\vee$
在 $\tau_{\leq n}$ 上诱导同构。
设 $E$ 作为 $f^{-1}\mathcal{O}_Y$-模复形
具有包含于 $[a, b]$ 的 Tor 振幅。
由引理 \ref{perfect-lemma-tor-independence-and-tor-amplitude}，
任意基变换后仍然如此。因此
$K_{n + 1}^\vee \otimes_{\mathcal{O}_X} E \to
K_n^\vee \otimes_{\mathcal{O}_X} E$
在 $\tau_{\leq n + a}$ 上诱导同构，
且在任意基变换后仍然如此。
应用右导出函子 $Rf_*$，可知映射
$M_{n + 1} \to M_n$ 在 $\tau_{\leq n + a}$
上诱导同构，且任意基变换后仍然如此。
选取可区别三角
$$
M_{n + 1} \to M_n \to C_n \to M_{n + 1}[1]
$$
取 $S'$ 为某点 $s \in S$ 的剩余域的谱，并作拉回，可见
$C_n \otimes_{\mathcal{O}_S}^\mathbf{L} \kappa(s)$
仅在次数 $\geq n + a$ 可能有非零上同调。由代数续篇，引理
\ref{more-algebra-lemma-lift-perfect-from-residue-field}
，对某个整数 $m_n$，完美复形 $C_n$
的 Tor 振幅包含于 $[n + a, m_n]$。
特别地，对偶完美复形 $C_n^\vee$
的 Tor 振幅包含于 $[-m_n, -n - a]$。

\medskip\noindent
令 $L_n = M_n^\vee$ 为对偶完美复形。上一段的讨论给出：
$L_n \to L_{n + 1}$ 在 $\tau_{\geq -n - a}$ 上诱导同构。
所以 $L = \text{hocolim} L_n$ 伪凝聚，见引理
\ref{perfect-lemma-pseudo-coherent-hocolim}。
由于
$$
R\SheafHom(K, E) = R\SheafHom(\text{hocolim} K_n, E) =
R\lim R\SheafHom(K_n, E) = R\lim K_n^\vee \otimes_{\mathcal{O}_X} E
$$
（上同调，引理 \ref{cohomology-lemma-colim-and-lim-of-duals}），
并且 $R\lim$ 与 $Rf_*$ 相容，得到
$$
Rf_*R\SheafHom(K, E) = R\lim M_n = R\lim R\SheafHom(L_n, \mathcal{O}_S) =
R\SheafHom(L, \mathcal{O}_S)
$$
这证明了 $S$ 上的公式。由于 $M_n$ 的构造与基变换相容，
任意基变换后该公式仍成立。
\end{proof}

\begin{remark}
\label{perfect-remark-compare-L}
读者或许已注意到引理 \ref{perfect-lemma-compute-ext-rel-perfect}
与引理 \ref{perfect-lemma-compute-ext} 之间的相似性。
事实上，引理 \ref{perfect-lemma-compute-ext-rel-perfect}
中的伪凝聚复形 $L$，可刻画为 $S$ 上唯一的伪凝聚复形，
使得存在函子性同构
$$
\Ext^i_{\mathcal{O}_S}(L, \mathcal{F}) \longrightarrow
\Ext^i_{\mathcal{O}_X}(K,
E \otimes_{\mathcal{O}_X}^\mathbf{L} Lf^*\mathcal{F})
$$
；当 $\mathcal{F}$ 取遍 $\QCoh(\mathcal{O}_S)$ 时，
这些同构与边界映射相容。若以后需要这一形式，
我们将在此表述精确结果并给出详细证明。
\end{remark}

\begin{lemma}
\label{perfect-lemma-bounded-on-fibres}
设 $f : X \to S$ 是平坦且局部有限表示的概形态射。
设 $E$ 为 $D(\mathcal{O}_X)$ 的伪凝聚对象。下列条件等价：
\begin{enumerate}
\item $E$ 是 $S$-完美的；
\item $E$ 局部下有界，并且对每个点 $s \in S$，
$D(\mathcal{O}_{X_s})$ 的对象 $L(X_s \to X)^*E$
局部下有界。
\end{enumerate}
\end{lemma}

\begin{proof}
由于一切都是局部的，立即归结到 $X$ 与 $S$ 都仿射的情形，
见引理 \ref{perfect-lemma-affine-locally-rel-perfect}。
设 $X \to S$ 对应于 $\Spec(A) \to \Spec(R)$，
而 $E$ 对应于 $D(A)$ 中的 $K$。若 $s$ 对应于素理想
$\mathfrak p \subset R$，则由于 $R \to A$ 平坦，
$L(X_s \to X)^*E$ 对应于
$K \otimes_R^\mathbf{L} \kappa(\mathfrak p)$；
例如见引理 \ref{perfect-lemma-compare-base-change}。
因此本引理由相应的代数结果得到，见代数续篇，引理
\ref{more-algebra-lemma-bounded-on-fibres-relatively-perfect}。
\end{proof}

\begin{remark}
\label{perfect-remark-discuss-rel-perfect}
在我们的定义适用时，定义 \ref{perfect-definition-relatively-perfect}
给出的相对完美复形概念，与 \cite{lieblich-complexes}
中的概念等价\footnote{为说明这一点，使用引理
\ref{perfect-lemma-affine-locally-rel-perfect} 与代数续篇，引理
\ref{more-algebra-lemma-structure-relatively-perfect}。}。
接着，假设只要求 $f : X \to S$ 局部有限型
（不必平坦，也不必局部有限表示）。
上述论文中的定义是：$E \in D(\mathcal{O}_X)$ 相对完美，如果
\begin{enumerate}
\item[(A)] 在 $X$ 上局部地，对象 $E$ 与一个有限复形拟同构，
该复形由在 $S$ 上平坦的有限表示 $\mathcal{O}_X$-模组成。
\end{enumerate}
另一方面，我们的定义 \ref{perfect-definition-relatively-perfect}
的自然推广是：
\begin{enumerate}
\item[(B)] $E$ 相对于 $S$ 伪凝聚
（态射续篇，定义
\ref{more-morphisms-definition-relative-pseudo-coherence})
，并且 $E$ 作为
$D(f^{-1}\mathcal{O}_S)$
的对象局部具有有限 Tor 维数
（上同调，定义 \ref{cohomology-definition-tor-amplitude}）。
\end{enumerate}
条件 (B) 的优点是，它显然定义了 $D(\mathcal{O}_X)$
的一个三角子范畴；而我们怀疑条件 (A) 并非如此。
条件 (A) 的优点是更易于操作，尤其是在涉及极限时。
\end{remark}









\section{消解性质}
\label{perfect-section-resolution-property}

\noindent
论文 \cite{totaro_resolution} 讨论了这一概念；
\cite{Gross-thesis}、\cite{Gross-surface} 与 \cite{Gross-stack}
继续了相关讨论。目前尚不知道域上的固有概形是否总具有消解性质，
还是这一断言为假。若你知道该问题的答案，请发送电子邮件至
\href{mailto:stacks.project@gmail.com}{stacks.project@gmail.com}.

\medskip\noindent
可以作出下述定义，尽管对一般概形考虑它几乎没有意义。

\begin{definition}
\label{perfect-definition-resolution-property}
设 $X$ 是概形。若每个有限型拟凝聚 $\mathcal{O}_X$-模
都是某个有限局部自由 $\mathcal{O}_X$-模的商，
则称 $X$ 具有{\it 消解性质}。
\end{definition}

\noindent
若 $X$ 是拟紧拟分离概形，则只需验证每个有限表示
$\mathcal{O}_X$-模（自动拟凝聚）都是有限局部自由
$\mathcal{O}_X$-模的商，见性质，引理
\ref{properties-lemma-finite-directed-colimit-surjective-maps}.
若 $X$ 是诺特概形，则有限型拟凝聚模恰好就是凝聚
$\mathcal{O}_X$-模，见概形的上同调，引理
\ref{coherent-lemma-coherent-Noetherian}。

\begin{lemma}
\label{perfect-lemma-resolution-property-ample}
设 $X$ 是概形。若 $X$ 有充足可逆 $\mathcal{O}_X$-模，
则 $X$ 具有消解性质。
\end{lemma}

\begin{proof}
这是性质，命题 \ref{properties-proposition-characterize-ample}
的直接推论。
\end{proof}

\begin{lemma}
\label{perfect-lemma-resolution-property-ample-relative}
设 $f : X \to Y$ 是概形态射。假设
\begin{enumerate}
\item $Y$ 拟紧、拟分离并具有消解性质；
\item $X$ 上存在一个 $f$-充足可逆模。
\end{enumerate}
则 $X$ 具有消解性质。
\end{lemma}

\begin{proof}
令 $\mathcal{F}$ 为有限型拟凝聚 $\mathcal{O}_X$-模，
并令 $\mathcal{L}$ 为 $f$-充足可逆模。
选取仿射开覆盖 $Y = V_1 \cup \ldots \cup V_m$。
置 $U_j = f^{-1}(V_j)$。由性质，命题
\ref{properties-proposition-characterize-ample}
，对每个 $j$ 都存在有限多个映射
$s_{j, i} : \mathcal{L}^{\otimes n_{j, i}}|_{U_j} \to \mathcal{F}|_{U_j}$
，它们联合满射。考虑拟凝聚 $\mathcal{O}_Y$-模
$$
\mathcal{H}_n =
f_*(\mathcal{F} \otimes_{\mathcal{O}_X} \mathcal{L}^{\otimes n})
$$
可以把 $s_{j, i}$ 视为层 $\mathcal{H}_{-n_{j, i}}$
在 $V_j$ 上的截面。假设能找到有限局部自由
$\mathcal{O}_Y$-模 $\mathcal{E}_{i, j}$ 以及映射
$\mathcal{E}_{i, j} \to \mathcal{H}_{-n_{j, i}}$，
使 $s_{j, i}$ 属于其像。则相应映射
$$
f^*\mathcal{E}_{i, j}
\otimes_{\mathcal{O}_X}
\mathcal{L}^{\otimes n_{i, j}} \longrightarrow \mathcal{F}
$$
将联合满射，从而本引理得证。由性质，引理
\ref{properties-lemma-quasi-coherent-colimit-finite-type}，
对每个 $i, j$ 都可以找到有限型拟凝聚子模
$\mathcal{H}'_{i, j} \subset  \mathcal{H}_{-n_{j, i}}$
，它包含 $V_j$ 上的截面 $s_{i, j}$。
所以 $Y$ 的消解性质给出满射
$\mathcal{E}_{i, j} \to \mathcal{H}'_{j, i}$，结论成立。
\end{proof}

\begin{lemma}
\label{perfect-lemma-resolution-property-goes-up-affine}
设 $f : X \to Y$ 是仿射或拟仿射概形态射，
且 $Y$ 拟紧拟分离。若 $Y$ 具有消解性质，则 $X$ 也具有。
\end{lemma}

\begin{proof}
由态射，引理 \ref{morphisms-lemma-quasi-affine-O-ample}，
这是引理 \ref{perfect-lemma-resolution-property-ample-relative} 的特殊情形。
\end{proof}

\noindent
下面是一种可证明消解性质下降的情形。

\begin{lemma}
\label{perfect-lemma-resolution-property-goes-down-finite-flat}
设 $f : X \to Y$ 是满射的有限局部自由概形态射。
若 $X$ 具有消解性质，则 $Y$ 也具有。
\end{lemma}

\begin{proof}
该条件意味着 $f$ 仿射，并且 $f_*\mathcal{O}_X$
是正秩有限局部自由 $\mathcal{O}_Y$-模。
令 $\mathcal{G}$ 为有限型拟凝聚 $\mathcal{O}_Y$-模。
由假设，对某个有限局部自由 $\mathcal{O}_X$-模 $\mathcal{E}$，
存在满射 $\mathcal{E} \to f^*\mathcal{G}$。
由于 $f_*$ 在拟凝聚模上正合
（概形的上同调，引理 \ref{coherent-lemma-relative-affine-vanishing}），
得到满射
$$
f_*\mathcal{E}
\longrightarrow
f_*f^*\mathcal{G} = \mathcal{G} \otimes_{\mathcal{O}_Y} f_*\mathcal{O}_X
$$
取对偶，得到满射
$$
f_*\mathcal{E}
\otimes_{\mathcal{O}_Y}
\SheafHom_{\mathcal{O}_Y}(f_*\mathcal{O}_X, \mathcal{O}_Y)
\longrightarrow
\mathcal{G}
$$
由于 $f_*\mathcal{E}$ 有限局部自由\footnote{具体地，
若 $A \to B$ 是有限局部自由环映射，且 $N$
是有限局部自由 $B$-模，则 $N$ 是有限局部自由 $A$-模。
首先，$N$ 在 $B$ 上有限局部自由蕴含 $N$ 作为 $B$-模
平坦且有限表示，见代数，引理
\ref{algebra-lemma-finite-projective}。其次，由代数，引理
\ref{algebra-lemma-finite-finitely-presented-extension}，
$N$ 是有限表示 $A$-模；由代数，引理
\ref{algebra-lemma-composition-flat}，它是平坦 $A$-模。
最后对 $A$ 应用代数，引理
\ref{algebra-lemma-finite-projective} 即得结论。}，结论成立。
\end{proof}

\begin{lemma}
\label{perfect-lemma-resolution-property-finite-number}
设 $X$ 是概形。假设给定：
\begin{enumerate}
\item 有限仿射开覆盖 $X = U_1 \cup \ldots \cup U_m$；
\item 有限型拟凝聚理想 $\mathcal{I}_j$，满足
$V(\mathcal{I}_j) = X \setminus U_j$。
\end{enumerate}
则 $X$ 具有消解性质，当且仅当对
$j = 1, \ldots, m$，$\mathcal{I}_j$
都是有限局部自由 $\mathcal{O}_X$-模的商。
\end{lemma}

\begin{proof}
一个方向显然。反过来，设
$\mathcal{E}_j \to \mathcal{I}_j$ 是满射，
且 $\mathcal{E}_j$ 有限局部自由。下一段将归结到诺特情形；
建议读者跳过。

\medskip\noindent
首先注意，$U_j \cap U_{j'}$ 拟紧，因为它是仿射概形
$U_j$ 上有限型拟凝聚理想
$\mathcal{I}_{j'}|{U_j}$ 的零点概形之补，见性质，引理
\ref{properties-lemma-quasi-coherent-finite-type-ideals}。
所以 $X$ 拟紧且拟分离，见概形，引理
\ref{schemes-lemma-characterize-quasi-separated}。
由极限，命题 \ref{limits-proposition-approximate}，
可把 $X = \lim X_i$ 写成诺特概形有向系统的极限，
其转移态射均仿射。对每个 $j$，可以找到某个 $i$
以及有限局部自由 $\mathcal{O}_{X_i}$-模
$\mathcal{E}_{i, j}$，其拉回为 $\mathcal{E}_j$，见极限，
引理 \ref{limits-lemma-descend-invertible-modules}。
增大 $i$ 后，可以假设复合
$\mathcal{E}_j \to \mathcal{I}_j \to \mathcal{O}_X$
是某映射 $\mathcal{E}_{i, j} \to \mathcal{O}_{X_i}$ 的拉回，
见极限，引理 \ref{limits-lemma-descend-modules-finite-presentation}。
以 $\mathcal{I}_{i, j} \subset \mathcal{O}_{X_i}$
表示该映射的像；这是诺特概形 $X_i$ 上的拟凝聚理想层，
其到 $X$ 的拉回为 $\mathcal{I}_j$。
以 $U_{i, j} \subset X_i$ 表示相应补开集；
由极限，引理 \ref{limits-lemma-limit-affine}，
可以假设对所有 $i, j$，这些开集都仿射。
若能对 $X_i$ 上的开集 $U_{i, j}$ 与理想层
$\mathcal{I}_{i, j}$ 证明本引理，则由于 $X$ 在 $X_i$ 上仿射，
引理 \ref{perfect-lemma-resolution-property-goes-up-affine}
说明它具有消解性质。这样便归结到诺特概形的情形。

\medskip\noindent
假设 $X$ 诺特。对每个凝聚模 $\mathcal{F}$，
可以选取有限列截面 $s_{jk} \in \mathcal{F}(U_j)$，
$k = 1, \ldots, e_j$，它们生成 $\mathcal{F}$ 到 $U_j$ 的限制。
由概形的上同调，引理 \ref{coherent-lemma-homs-over-open}，
对某个 $n_{jk} \geq 1$，可以把 $s_{jk}$ 延拓为映射
$s'_{jk} : \mathcal{I}_i^{n_{jk}} \to \mathcal{F}$。
于是考虑复合
$$
\mathcal{E}_j^{\otimes n_{jk}} \to \mathcal{I}_j^{n_{jk}} \to \mathcal{F}
$$
即可得结论。
\end{proof}

\begin{lemma}
\label{perfect-lemma-resolution-property-ample-family}
设 $X$ 是概形。若 $X$ 有可逆模的充足族
（态射，定义
\ref{morphisms-definition-family-ample-invertible-modules}),
则 $X$ 具有消解性质。
\end{lemma}

\begin{proof}
由于 $X$ 拟紧，存在 $n$ 以及各对
$(\mathcal{L}_i, s_i)$，$i = 1, \ldots, n$，
其中 $\mathcal{L}_i$ 是可逆 $\mathcal{O}_X$-模，
$s_i \in \Gamma(X, \mathcal{L}_i)$ 是截面；
并且 $s_i$ 非零的点集 $U_i \subset X$ 仿射，
且 $X = U_1 \cup \ldots \cup U_n$。
令 $\mathcal{I}_i \subset \mathcal{O}_X$ 为
$s_i : \mathcal{L}_i^{\otimes -1} \to \mathcal{O}_X$ 的像。
应用引理 \ref{perfect-lemma-resolution-property-finite-number}，
可知 $X$ 具有消解性质。
\end{proof}

\begin{lemma}
\label{perfect-lemma-regular-resolution-property}
设 $X$ 是具有仿射对角的拟紧正则概形。
则 $X$ 具有消解性质。
\end{lemma}

\begin{proof}
结合除子，引理 \ref{divisors-lemma-regular-ample-family}
与上面的引理 \ref{perfect-lemma-resolution-property-ample-family}。
\end{proof}

\begin{lemma}
\label{perfect-lemma-resolution-property-descends}
设 $X = \lim X_i$ 是拟紧拟分离概形有向系统的极限，
其转移态射均仿射。则 $X$ 具有消解性质，
当且仅当对某个 $i$，$X_i$ 具有消解性质。
\end{lemma}

\begin{proof}
若 $X_i$ 具有消解性质，则由引理
\ref{perfect-lemma-resolution-property-goes-up-affine}，$X$ 也具有。
反过来，假设 $X$ 具有消解性质。选取 $i \in I$，
并选取有限仿射开覆盖
$X_i = U_{i, 1} \cup \ldots \cup U_{i, m}$.
对每个 $j$，选取有限型拟凝聚理想层
$\mathcal{I}_{i, j} \subset \mathcal{O}_{X_i}$，使
$X_i \setminus V(\mathcal{I}_{i, j}) = U_{i, j}$，见性质，引理
\ref{properties-lemma-quasi-coherent-finite-type-ideals}。
以 $U_j \subset X$ 表示 $U_{i, j}$ 的逆像，
并以 $\mathcal{I}_j \subset \mathcal{O}_X$
表示 $\mathcal{I}_{i, j}$ 的拉回。
由于 $X$ 具有消解性质，可以选取有限局部自由
$\mathcal{O}_X$-模 $\mathcal{E}_j$
以及满射 $\mathcal{E}_j \to \mathcal{I}_j$。
由极限，诸引理 \ref{limits-lemma-descend-invertible-modules} 与
\ref{limits-lemma-descend-modules-finite-presentation}
，增大 $i$ 后，可以找到有限局部自由
$\mathcal{O}_{X_i}$-模 $\mathcal{E}_{i, j}$
与映射 $\mathcal{E}_{i, j} \to \mathcal{O}_{X_i}$，
它们到 $X$ 的基变换恢复各复合
$\mathcal{E}_j \to \mathcal{I}_j \to \mathcal{O}_X$, $j = 1, \ldots, m$.
有限表示 $\mathcal{O}_{X_i}$-模
$\Coker(\mathcal{E}_{i, j} \to \mathcal{O}_{X_i})$
与 $\mathcal{O}_{X_i}/\mathcal{I}_{i, j}$
到 $X$ 的拉回，作为 $\mathcal{O}_X$ 的商是一致的。
所以由极限，引理
\ref{limits-lemma-descend-modules-finite-presentation}
，可以假设这两个模相同；换言之，
$\mathcal{E}_{i, j} \to \mathcal{O}_{X_i}$ 的像等于
$\mathcal{I}_{i, j}$。于是由引理
\ref{perfect-lemma-resolution-property-finite-number}，
$X_i$ 具有消解性质。
\end{proof}

\begin{lemma}
\label{perfect-lemma-resolution-property-affine-diagonal}
\begin{reference}
\cite[Proposition 1.3]{totaro_resolution} 的特殊情形。
\end{reference}
设 $X$ 是具有消解性质的拟紧拟分离概形。
则 $X$ 具有仿射对角。
\end{lemma}

\begin{proof}
结合极限，命题 \ref{limits-proposition-approximate}
与引理 \ref{perfect-lemma-resolution-property-descends}，
可归结到 $X$ 诺特的情形（略去一个小细节）。
假设 $X$ 诺特。回忆，$X \times X$ 被仿射开集
$U \times V$ 覆盖，其中 $U$、$V$ 为 $X$ 的仿射开集；
见概形，第 \ref{schemes-section-fibre-products} 节。
所以为证明对角 $\Delta : X \to X \times X$ 仿射，
只需证明对 $X$ 的所有仿射开集 $U$、$V$，
$U \cap V = \Delta^{-1}(U \times V)$ 仿射；见态射，引理
\ref{morphisms-lemma-characterize-affine}。
特别地，只需证明：若 $U$ 是 $X$ 的仿射开集，
则包含态射 $j : U \to X$ 仿射。
由概形的上同调，引理
\ref{coherent-lemma-criterion-affine-morphism}，
只需证明对任意拟凝聚 $\mathcal{O}_U$-模 $\mathcal{G}$，
$R^1j_*\mathcal{G} = 0$。
由命题 \ref{perfect-proposition-Noetherian}
（这里使用了已归结到诺特情形这一点），
可以用拟凝聚 $\mathcal{O}_X$-模复形
$\mathcal{H}^\bullet$ 表示 $Rj_*\mathcal{G}$。
为得矛盾，假设
$$
H^1(\mathcal{H}^\bullet) =
\Ker(\mathcal{H}^1 \to \mathcal{H}^2)/\Im(\mathcal{H}^0 \to \mathcal{H}^1)
$$
非零。于是可以找到凝聚 $\mathcal{O}_X$-模
$\mathcal{F}$ 与映射
$$
\mathcal{F} \longrightarrow
\Ker(\mathcal{H}^1 \to \mathcal{H}^2)
$$
，使其与到 $H^1(\mathcal{H}^\bullet)$ 的投影之复合非零。
事实上，由性质，引理
\ref{properties-lemma-quasi-coherent-colimit-finite-type}，
可以把 $\Ker(\mathcal{H}^1 \to \mathcal{H}^2)$
写成其凝聚子模的滤过并，而其中某一个会满足要求。
接着，利用 $X$ 的消解性质，选取有限局部自由
$\mathcal{O}_X$-模 $\mathcal{E}$
以及满射 $\mathcal{E} \to \mathcal{F}$。
这在导出范畴中产生映射
$$
\mathcal{E}[-1] \longrightarrow Rj_*\mathcal{G}
$$
，它在上同调层上非零，因而在 $D(\mathcal{O}_X)$ 中非零。
由伴随，这等同于映射
$$
j^*\mathcal{E}[-1] \to \mathcal{G}
$$
在 $D(\mathcal{O}_U)$ 中非零。由于 $\mathcal{E}$ 有限局部自由，
这又等同于下列群中的非零元素：
$$
H^1(U, j^*\mathcal{E}^\vee \otimes_{\mathcal{O}_U} \mathcal{G})
$$
其中
$\mathcal{E}^\vee = \SheafHom_{\mathcal{O}_X}(\mathcal{E}, \mathcal{O}_X)$
是对偶有限局部自由模。但由概形的上同调，引理
\ref{coherent-lemma-quasi-coherent-affine-cohomology-zero}
，该群为零，这就是所需矛盾。
（若对使用对偶的步骤有疑问，请见更一般的上同调，引理
\ref{cohomology-lemma-dual-perfect-complex}。）
\end{proof}






\section{消解性质与完美复形}
\label{perfect-section-resolution-property-perfect}

\noindent
本节讨论具有消解性质的概形上，完美复形与严格完美复形之间的关系。

\begin{lemma}
\label{perfect-lemma-construct-strictly-perfect}
设 $X$ 是具有消解性质的拟紧拟分离概形。
设 $\mathcal{F}^\bullet$ 是拟凝聚 $\mathcal{O}_X$-模的下有界复形，
它表示 $D(\mathcal{O}_X)$ 的一个完美对象。
则存在有限局部自由 $\mathcal{O}_X$-模的有界复形
$\mathcal{E}^\bullet$，以及拟同构
$\mathcal{E}^\bullet \to \mathcal{F}^\bullet$。
\end{lemma}

\begin{proof}
取整数 $a, b \in \mathbf{Z}$，使得 $\mathcal{F}^\bullet$
的 Tor 振幅包含于 $[a, b]$，并且当 $n < a$ 时
$\mathcal{F}^n = 0$。这种整数对的存在性由上同调，引理
\ref{cohomology-lemma-perfect} 与 $X$ 拟紧这一事实得到。
若 $b < a$，则 $\mathcal{F}^\bullet$ 在导出范畴中为零，
本引理成立。我们将对 $b - a \geq 0$ 作归纳，证明存在复形
$\mathcal{E}^a \to \ldots \to \mathcal{E}^b$
，其中 $\mathcal{E}^i$ 有限局部自由，并且存在拟同构
$\mathcal{E}^\bullet \to \mathcal{F}^\bullet$。

\medskip\noindent
基本情形为 $b - a = 0$。此时
$H^b(\mathcal{F}^\bullet) = H^a(\mathcal{F}^\bullet) =
\Ker(\mathcal{F}^a \to \mathcal{F}^{a + 1})$
有限局部自由。事实上，它是 Tor 维数为 $0$
的有限表示 $\mathcal{O}_X$-模，因而有限局部自由。
见上同调，诸引理 \ref{cohomology-lemma-perfect}、
\ref{cohomology-lemma-finite-cohomology} 以及性质，引理
\ref{properties-lemma-finite-locally-free}。
所以可以取 $\mathcal{E}^\bullet$ 为置于次数 $b$ 的
$H^b(\mathcal{F}^\bullet)$。证明的余下部分处理归纳步骤。

\medskip\noindent
假设 $b > a$。注意
$$
H^b(\mathcal{F}^\bullet) =
\Ker(\mathcal{F}^b \to \mathcal{F}^{b + 1})/
\Im(\mathcal{F}^{b - 1} \to \mathcal{F}^b)
$$
是有限型拟凝聚 $\mathcal{O}_X$-模，见上同调，诸引理
\ref{cohomology-lemma-perfect} 与
\ref{cohomology-lemma-finite-cohomology}。于是可以找到
有限型拟凝聚 $\mathcal{O}_X$-模 $\mathcal{F}$ 以及映射
$$
\mathcal{F} \longrightarrow
\Ker(\mathcal{F}^b \to \mathcal{F}^{b + 1})
$$
，使它与到 $H^b(\mathcal{F}^\bullet)$ 的投影之复合满射。
事实上，由性质，引理
\ref{properties-lemma-quasi-coherent-colimit-finite-type}，
可以把 $\Ker(\mathcal{F}^b \to \mathcal{F}^{b + 1})$
写成其有限型拟凝聚子模的滤过并，而其中某一个会满足要求。
接着，利用 $X$ 的消解性质，选取有限局部自由
$\mathcal{O}_X$-模 $\mathcal{E}^b$
以及满射 $\mathcal{E}^b \to \mathcal{F}$。考虑复形映射
$$
\alpha : \mathcal{E}^b[-b] \to \mathcal{F}^\bullet
$$
及其锥 $C(\alpha)^\bullet$，见导出范畴，定义
\ref{derived-definition-cone}。
注意，$C(\alpha)^\bullet$ 仅在次数 $\geq a$ 时可能非零；
由上同调，引理 \ref{cohomology-lemma-cone-tor-amplitude}，
其 Tor 振幅包含于 $[a, b]$；并且由构造，
$H^b(C(\alpha)^\bullet) = 0$。
所以实际上 $C(\alpha)^\bullet$ 的 Tor 振幅包含于
$[a, b - 1]$。因此可对 $C(\alpha)^\bullet$
应用归纳假设，得到具有归纳假设中所述性质的复形映射
$$
(\mathcal{E}^a \to \ldots \to \mathcal{E}^{b - 1})
\longrightarrow
C(\alpha)^\bullet
$$
。展开锥的定义，得到交换图
$$
\xymatrix{
\ldots \ar[r] &
\mathcal{E}^{b - 2} \ar[r] \ar[d] &
\mathcal{E}^{b - 1} \ar[r] \ar[d] &
0 \ar[r] \ar[d] &
\ldots \\
\ldots \ar[r] &
\mathcal{F}^{b - 2} \ar[r] &
\mathcal{F}^{b - 1} \oplus \mathcal{E}^b \ar[r] &
\mathcal{F}^b \ar[r] &
\ldots
}
$$
显然，由此得到复形映射
$(\mathcal{E}^a \to \ldots \to \mathcal{E}^b) \to \mathcal{F}^\bullet$.
略去验证此映射是拟同构。
\end{proof}

\begin{lemma}
\label{perfect-lemma-resolution-property-perfect-complex}
设 $X$ 是具有消解性质的拟紧拟分离概形。
则 $D(\mathcal{O}_X)$ 的每个完美对象，
都可由有限局部自由 $\mathcal{O}_X$-模的有界复形表示。
\end{lemma}

\begin{proof}
令 $E$ 为 $D(\mathcal{O}_X)$ 的完美对象。
由引理 \ref{perfect-lemma-resolution-property-affine-diagonal}，
$X$ 具有仿射对角。所以由命题
\ref{perfect-proposition-quasi-compact-affine-diagonal}，
可用拟凝聚 $\mathcal{O}_X$-模复形
$\mathcal{F}^\bullet$ 表示 $E$。
注意，因为 $X$ 拟紧，$E$ 属于 $D^b(\mathcal{O}_X)$。
所以当 $n \ll 0$ 时，
$\tau_{\geq n}\mathcal{F}^\bullet$
是表示 $E$ 的拟凝聚 $\mathcal{O}_X$-模下有界复形。
于是应用引理 \ref{perfect-lemma-construct-strictly-perfect} 即得结论。
\end{proof}

\begin{lemma}
\label{perfect-lemma-resolution-property-map-perfect-complex}
设 $X$ 是具有消解性质的拟紧拟分离概形。
设 $\mathcal{E}^\bullet$ 与 $\mathcal{F}^\bullet$
为有限局部自由 $\mathcal{O}_X$-模的有限复形。
则任意
$\alpha \in \Hom_{D(\mathcal{O}_X)}(\mathcal{E}^\bullet, \mathcal{F}^\bullet)$
都可由图
$$
\mathcal{E}^\bullet \leftarrow \mathcal{G}^\bullet \to \mathcal{F}^\bullet
$$
表示，其中 $\mathcal{G}^\bullet$
是有限局部自由 $\mathcal{O}_X$-模的有界复形，
且 $\mathcal{G}^\bullet \to \mathcal{E}^\bullet$ 是拟同构。
\end{lemma}

\begin{proof}
由引理 \ref{perfect-lemma-resolution-property-affine-diagonal}，
$X$ 具有仿射对角。所以由命题
\ref{perfect-proposition-quasi-compact-affine-diagonal}，
可用图
$$
\mathcal{E}^\bullet \leftarrow \mathcal{H}^\bullet \to \mathcal{F}^\bullet
$$
表示 $\alpha$，其中 $\mathcal{H}^\bullet$
是拟凝聚 $\mathcal{O}_X$-模复形，且
$\mathcal{H}^\bullet \to \mathcal{E}^\bullet$
是拟同构。当 $n \ll 0$ 时，映射
$\mathcal{H}^\bullet \to \mathcal{E}^\bullet$ 与
$\mathcal{H}^\bullet \to \mathcal{F}^\bullet$
都经过拟同构
$\mathcal{H}^\bullet \to \tau_{\geq n}\mathcal{H}^\bullet$
分解，这仅仅因为 $\mathcal{E}^\bullet$ 与
$\mathcal{F}^\bullet$ 都是有界复形。
因此可以用 $\tau_{\geq n}\mathcal{H}^\bullet$
替换 $\mathcal{H}^\bullet$，并假设 $\mathcal{H}^\bullet$ 下有界。
随后应用引理 \ref{perfect-lemma-construct-strictly-perfect} 即得结论。
\end{proof}

\begin{lemma}
\label{perfect-lemma-resolution-property-homotopy-map-perfect-complex}
设 $X$ 是具有消解性质的拟紧拟分离概形。
设 $\mathcal{E}^\bullet$ 与 $\mathcal{F}^\bullet$
为有限局部自由 $\mathcal{O}_X$-模的有限复形。
设 $\alpha^\bullet, \beta^\bullet :\mathcal{E}^\bullet \to \mathcal{F}^\bullet$
为两个复形映射，它们在 $D(\mathcal{O}_X)$ 中定义同一映射。
则存在拟同构
$\gamma^\bullet : \mathcal{G}^\bullet \to \mathcal{E}^\bullet$
，其中 $\mathcal{G}^\bullet$
是有限局部自由 $\mathcal{O}_X$-模的有界复形，
并且 $\alpha^\bullet \circ \gamma^\bullet$ 与
$\beta^\bullet \circ \gamma^\bullet$ 是同伦的复形映射。
\end{lemma}

\begin{proof}
由引理 \ref{perfect-lemma-resolution-property-affine-diagonal}，
$X$ 具有仿射对角。所以由命题
\ref{perfect-proposition-quasi-compact-affine-diagonal}
（以及导出范畴的定义），存在拟同构
$\gamma^\bullet : \mathcal{G}^\bullet \to \mathcal{E}^\bullet$
，其中 $\mathcal{G}^\bullet$ 是拟凝聚
$\mathcal{O}_X$-模复形，并且
$\alpha^\bullet \circ \gamma^\bullet$ 与
$\beta^\bullet \circ \gamma^\bullet$ 是同伦的复形映射。
选取证明此事的同伦
$h^i : \mathcal{G}^i \to \mathcal{F}^{i - 1}$。
选取 $n \ll 0$。由于 $\mathcal{E}^\bullet$ 下有界，
映射 $\gamma^\bullet$ 典范地经过商映射
$\mathcal{G}^\bullet \to \tau_{\geq n}\mathcal{G}^\bullet$
分解。出于完全相同的原因，各映射 $h^i$ 经过满射
$\mathcal{G}^i \to (\tau_{\geq n}\mathcal{G})^i$.
分解。因此可以用 $\tau_{\geq n}\mathcal{G}^\bullet$
替换 $\mathcal{G}^\bullet$。
随后应用引理 \ref{perfect-lemma-construct-strictly-perfect} 即得结论。
\end{proof}

\begin{proposition}
\label{perfect-proposition-perfect-resolution-property}
设 $X$ 是具有消解性质的拟紧拟分离概形。记
\begin{enumerate}
\item $\mathcal{A}$ 为有限局部自由 $\mathcal{O}_X$-模的加性范畴；
\item $K^b(\mathcal{A})$ 为 $\mathcal{A}$
中有界复形的同伦范畴，见导出范畴，第
\ref{derived-section-homotopy} 节；
\item $D_{perf}(\mathcal{O}_X)$ 为 $D(\mathcal{O}_X)$
中由完美对象组成的严格全、饱和三角子范畴。
\end{enumerate}
在此记号下，显然的函子
$$
K^b(\mathcal{A}) \longrightarrow D_{perf}(\mathcal{O}_X)
$$
是三角范畴的正合函子，并经过三角范畴的等价
$S^{-1}K^b(\mathcal{A}) \to D_{perf}(\mathcal{O}_X)$ 分解，
其中 $S$ 是 $K^b(\mathcal{A})$ 中由拟同构组成的饱和乘法系。
\end{proposition}

\begin{proof}
若能读懂本命题的陈述，请跳过第一段。关于其中使用的一些定义，见
导出范畴，定义
\ref{derived-definition-triangulated-subcategory}
（三角子范畴），导出范畴，定义
\ref{derived-definition-saturated}（饱和三角子范畴），
导出范畴，定义 \ref{derived-definition-localization}
（与三角结构相容的乘法系），以及范畴，定义
\ref{categories-definition-saturated-multiplicative-system}
（饱和乘法系）。
由上同调，诸引理
\ref{cohomology-lemma-two-out-of-three-perfect} 与
\ref{cohomology-lemma-summands-perfect}，
$D_{perf}(\mathcal{O}_X)$ 是 $D(\mathcal{O}_X)$
的饱和三角子范畴。还要注意，$K^b(\mathcal{A})$
是三角范畴，见导出范畴，引理
\ref{derived-lemma-bounded-triangulated-subcategories}.

\medskip\noindent
显然，该函子把可区别三角送到可区别三角，即它是正合的。
于是由导出范畴，引理
\ref{derived-lemma-triangle-functor-localize}，
$S$ 是与 $K^b(\mathcal{A})$ 上三角结构相容的饱和乘法系。
所以由导出范畴，命题
\ref{derived-proposition-construct-localization}.
，局部化 $S^{-1}K^b(\mathcal{A})$ 存在且为三角范畴。
由导出范畴，引理
\ref{derived-lemma-universal-property-localization}.
，得到正合分解
$S^{-1}K^b(\mathcal{A}) \to D_{perf}(\mathcal{O}_X)$。
由诸引理 \ref{perfect-lemma-resolution-property-perfect-complex}、
\ref{perfect-lemma-resolution-property-map-perfect-complex} 与
\ref{perfect-lemma-resolution-property-homotopy-map-perfect-complex}
，该函子是等价。最后，函子
$S^{-1}K^b(\mathcal{A}) \to D_{perf}(\mathcal{O}_X)$
由导出范畴，引理 \ref{derived-lemma-exact-equivalence}
是三角范畴的等价（即可区别三角相互对应）。
\end{proof}






\section{K-群}
\label{perfect-section-K-groups}

\noindent
简要讨论与概形相伴的各类范畴的 $K_0$。
先前材料见代数，第 \ref{algebra-section-K-groups} 节，
同调，第 \ref{homology-section-K-groups} 节，
导出范畴，第 \ref{derived-section-K-groups} 节，以及
代数续篇，引理 \ref{more-algebra-lemma-perfect-to-K-group-universal}。

\medskip\noindent
类似于代数，第 \ref{algebra-section-K-groups} 节，
对任意诺特概形 $X$，我们将定义两个 $K$-群
$K'_0(X)$ 与 $K_0(X)$。前者使用凝聚 $\mathcal{O}_X$-模，
后者使用有限局部自由 $\mathcal{O}_X$-模。

\begin{lemma}
\label{perfect-lemma-Noetherian-Kprime}
设 $X$ 是诺特概形。则
$$
K_0(\textit{Coh}(\mathcal{O}_X)) =
K_0(D^b(\textit{Coh}(\mathcal{O}_X)) =
K_0(D^b_{\textit{Coh}}(\mathcal{O}_X))
$$
\end{lemma}

\begin{proof}
第一个等式由导出范畴，引理
\ref{derived-lemma-K-bounded-derived} 给出。
由导出范畴，引理 \ref{derived-lemma-DBA-map-K}，有
$K_0(\textit{Coh}(\mathcal{O}_X)) =
K_0(D^b_{\textit{Coh}}(\mathcal{O}_X))$
。这证明了本引理。
（也可以利用命题 \ref{perfect-proposition-DCoh} 给出的
$D^b(\textit{Coh}(\mathcal{O}_X)) = D^b_{\textit{Coh}}(\mathcal{O}_X)$
来得到第二个等式。）
\end{proof}

\noindent
定义如下。

\begin{definition}
\label{perfect-definition-K-group}
设 $X$ 是概形。
\begin{enumerate}
\item 以 $K_0(X)$ 表示 {\it $X$ 的 Grothendieck 群}。
它是 $D(\mathcal{O}_X)$ 中由完美对象组成的严格全、饱和三角子范畴
$D_{perf}(\mathcal{O}_X)$ 的第零 K-群。公式为
$$
K_0(X) = K_0(D_{perf}(\mathcal{O}_X))
$$
\item 若 $X$ 局部诺特，则以 $K'_0(X)$ 表示
{\it $X$ 上凝聚层的 Grothendieck 群}。
它是凝聚 $\mathcal{O}_X$-模的阿贝尔范畴的第零 $K$-群。公式为
$$
K'_0(X) = K_0(\textit{Coh}(\mathcal{O}_X))
$$
\end{enumerate}
\end{definition}

\noindent
我们将证明，若 $X$ 具有消解性质
（例如，若 $X$ 有充足可逆模），则我们对 $K_0(X)$ 的定义
与常用的有限局部自由模定义一致。见引理 \ref{perfect-lemma-K-is-old-K}。

\begin{lemma}
\label{perfect-lemma-K-agrees-affine}
设 $X = \Spec(R)$ 是仿射概形。则 $K_0(X) = K_0(R)$；
若 $R$ 诺特，则 $K'_0(X) = K'_0(R)$。
\end{lemma}

\begin{proof}
回忆，$K'_0(R)$ 与 $K_0(R)$ 已在代数，
第 \ref{algebra-section-K-groups} 节中定义。

\medskip\noindent
由代数续篇，引理
\ref{more-algebra-lemma-perfect-to-K-group-universal}，
有 $K_0(R) = K_0(D_{perf}(R))$。
由诸引理 \ref{perfect-lemma-perfect-affine} 与
\ref{perfect-lemma-affine-compare-bounded}，
有 $D_{perf}(R) = D_{perf}(\mathcal{O}_X)$。
这证明了等式 $K_0(R) = K_0(X)$。

\medskip\noindent
等式 $K'_0(R) = K'_0(X)$ 成立，因为由概形的上同调，引理
\ref{coherent-lemma-coherent-Noetherian}，
$\textit{Coh}(\mathcal{O}_X)$ 等价于有限 $R$-模的范畴。
此外，由定义显然可知，$K'_0(R)$ 正是有限 $R$-模范畴的第零 K-群。
\end{proof}

\noindent
设 $X$ 是诺特概形。此时 $K'_0(X)$ 与 $K_0(X)$ 都有定义，
并且存在典范映射
$$
K_0(X) = K_0(D_{perf}(\mathcal{O}_X))
\longrightarrow
K_0(D^b_{\textit{Coh}}(\mathcal{O}_X)) = K'_0(X)
$$
事实上，完美复形属于 $D^b_{\textit{Coh}}(\mathcal{O}_X)$
（由引理 \ref{perfect-lemma-identify-pseudo-coherent-noetherian}），
包含函子
$D_{perf}(\mathcal{O}_X) \to D^b_{\textit{Coh}}(\mathcal{O}_X)$
在第零 $K$-群上诱导映射
（导出范畴，引理 \ref{derived-lemma-map-K}），
而右侧的等式由引理 \ref{perfect-lemma-Noetherian-Kprime} 给出。

\begin{lemma}
\label{perfect-lemma-Kprime-K}
设 $X$ 是诺特正则概形。则映射
$K_0(X) \to K'_0(X)$ 是同构。
\end{lemma}

\begin{proof}
这立即由引理 \ref{perfect-lemma-perfect-on-regular}
以及上面对映射 $K_0(X) \to K'_0(X)$ 的构造得到。
\end{proof}

\noindent
设 $X$ 是概形。以 $\textit{Vect}(X)$ 表示有限局部自由
$\mathcal{O}_X$-模的范畴。尽管 $\textit{Vect}(X)$
一般不是阿贝尔范畴，但 $\textit{Vect}(X)$ 中短正合列的含义是清楚的。
以 $K_0(\textit{Vect}(X))$ 表示具有下列性质的唯一阿贝尔群
\footnote{这里正确的一般性应当是对任意正合范畴定义 $K_0$，
见内射对象，评注 \ref{injectives-remark-embed-exact-category}。}：
\begin{enumerate}
\item 对每个有限局部自由 $\mathcal{O}_X$-模 $\mathcal{E}$，
给定 $K_0(\textit{Vect}(X))$ 中的元素 $[\mathcal{E}]$；
\item 对有限局部自由 $\mathcal{O}_X$-模的每个短正合列
$0 \to \mathcal{E}' \to \mathcal{E} \to \mathcal{E}'' \to 0$
，在 $K_0(\textit{Vect}(X))$ 中有关系
$[\mathcal{E}] = [\mathcal{E}'] + [\mathcal{E}'']$；
\item 群 $K_0(\textit{Vect}(X))$ 由各元素 $[\mathcal{E}]$ 生成；
\item $K_0(\textit{Vect}(X))$ 中生成元
$[\mathcal{E}]$ 之间的所有关系，都是上述正合列所给关系的
$\mathbf{Z}$-线性组合。
\end{enumerate}
略去 $K_0(\textit{Vect}(X))$ 的详细构造。存在自然映射
$$
K_0(\textit{Vect}(X)) \longrightarrow K_0(X)
$$
具体地，对有限局部自由 $\mathcal{O}_X$-模 $\mathcal{E}$，
以 $\mathcal{E}[0]$ 表示 $X$ 上的完美复形，
它在次数 $0$ 为 $\mathcal{E}$，在其他次数为零。
给定有限局部自由 $\mathcal{O}_X$-模的短正合列
$0 \to \mathcal{E} \to \mathcal{E}' \to \mathcal{E}'' \to 0$
，得到可区别三角
$\mathcal{E}[0] \to \mathcal{E}'[0] \to \mathcal{E}''[0] \to \mathcal{E}[1]$，
见导出范畴，第 \ref{derived-section-canonical-delta-functor} 节。
这说明，把 $[\mathcal{E}]$ 送到 $[\mathcal{E}[0]]$，
就得到映射
$K_0(\textit{Vect}(X)) \to K_0(D_{perf}(\mathcal{O}_X)) = K_0(X)$
；对这一可怕记号深表歉意。

\begin{lemma}
\label{perfect-lemma-K-is-old-K}
设 $X$ 是具有消解性质的拟紧拟分离概形。
则映射 $K_0(\textit{Vect}(X)) \to K_0(X)$ 是同构。
\end{lemma}

\begin{proof}
本引理将直接由诸引理
\ref{perfect-lemma-resolution-property-perfect-complex}、
\ref{perfect-lemma-resolution-property-map-perfect-complex} 与
\ref{perfect-lemma-resolution-property-homotopy-map-perfect-complex}
得到；下文将不再说明地使用这些引理的结果。
构造逆映射
$$
c : K_0(X) = K_0(D_{perf}(\mathcal{O}_X)) \longrightarrow
K_0(\textit{Vect}(X))
$$
具体地，$D_{perf}(\mathcal{O}_X)$ 的任意对象，
都可由有限局部自由 $\mathcal{O}_X$-模的有界复形
$\mathcal{E}^\bullet$ 表示。置
$$
c([\mathcal{E}^\bullet]) = \sum (-1)^i[\mathcal{E}^i]
$$
当然，还需证明这一定义良好。暂且把 $c$
看成定义在有限局部自由 $\mathcal{O}_X$-模有界复形上的映射。

\medskip\noindent
假设 $\mathcal{E}^\bullet \to \mathcal{F}^\bullet$
是有限局部自由 $\mathcal{O}_X$-模有界复形之间的满射。
令 $\mathcal{K}^\bullet$ 为其核。于是得到
$\mathcal{O}_X$-模的短正合列
$$
0 \to \mathcal{K}^n \to \mathcal{E}^n \to \mathcal{F}^n \to 0
$$
。由于 $\mathcal{F}^n$ 有限局部自由，这些正合列局部分裂。
所以 $\mathcal{K}^\bullet$ 也是有限局部自由
$\mathcal{O}_X$-模的有界复形，并且
$c(\mathcal{E}^\bullet) = c(\mathcal{K}^\bullet) + c(\mathcal{F}^\bullet)$
在 $K_0(\textit{Vect}(X))$ 中成立。

\medskip\noindent
设给定有限局部自由 $\mathcal{O}_X$-模的无环有界复形
$\mathcal{E}^\bullet$。设当 $n \not \in [a, b]$ 时
$\mathcal{E}^n = 0$。可以把 $\mathcal{E}^\bullet$
分解成短正合列
$$
\begin{matrix}
0 \to \mathcal{E}^a \to \mathcal{E}^{a + 1} \to \mathcal{F}^{a + 1} \to 0,\\
0 \to \mathcal{F}^{a + 1} \to \mathcal{E}^{a + 2} \to
\mathcal{F}^{a + 3} \to 0, \\
\ldots \\
0 \to \mathcal{F}^{b - 3} \to \mathcal{E}^{b - 2} \to
\mathcal{F}^{b - 2} \to 0, \\
0 \to \mathcal{F}^{b - 2} \to \mathcal{E}^{b - 1} \to \mathcal{E}^b \to 0
\end{matrix}
$$
作降序归纳可知
$\mathcal{F}^{b - 2}, \ldots, \mathcal{F}^{a + 1}$
都是有限局部自由 $\mathcal{O}_X$-模，并且
$$
c(\mathcal{E}^\bullet) = \sum (-1)[\mathcal{E}^n] =
\sum (-1)^n([\mathcal{F}^{n - 1}] + [\mathcal{F}^n]) = 0
$$
所以我们的构造在无环复形上取零。

\medskip\noindent
由前两段的结果可知，$c$ 定义良好。具体地，设有限局部自由
$\mathcal{O}_X$-模的有界复形
$\mathcal{E}^\bullet$ 与 $\mathcal{F}^\bullet$
表示 $D(\mathcal{O}_X)$ 中的同一对象。则可找到拟同构
$a : \mathcal{G}^\bullet \to \mathcal{E}^\bullet$ 与
$b : \mathcal{G}^\bullet \to \mathcal{F}^\bullet$，
其中 $\mathcal{G}^\bullet$ 是有限局部自由
$\mathcal{O}_X$-模的有界复形。得到复形的短正合列
$$
0 \to \mathcal{E}^\bullet \to C(a)^\bullet \to \mathcal{G}^\bullet[1] \to 0
$$
，见导出范畴，定义 \ref{derived-definition-cone}。
由于 $a$ 是拟同构，其锥 $C(a)^\bullet$ 无环
（例如这可由导出范畴，第
\ref{derived-section-canonical-delta-functor} 节中的讨论推出）。因此
$$
0 = c(C(f)^\bullet) = c(\mathcal{E}^\bullet) + c(\mathcal{G}^\bullet[1]) =
c(\mathcal{E}^\bullet) - c(\mathcal{G}^\bullet)
$$
，正如所需。对 $b$ 使用同一论证可知
$0 = c(\mathcal{F}^\bullet) - c(\mathcal{G}^\bullet)$.
所以 $c(\mathcal{E}^\bullet) = c(\mathcal{F}^\bullet)$，
而 $c$ 定义良好。

\medskip\noindent
类似地，对有限局部自由 $\mathcal{O}_X$-模有界复形之间的映射
$\mathcal{E}^\bullet \to \mathcal{F}^\bullet$ 的锥作同样论证，
可知若 $X \to Y \to Z$ 是 $D_{perf}(\mathcal{O}_X)$
中的可区别三角，则 $c(Y) = c(X) + c(Z)$。略去细节。
于是得到所需的阿贝尔群同态
$c : K_0(X) \to K_0(\textit{Vect}(X))$。

\medskip\noindent
显然，复合
$K_0(\textit{Vect}(X)) \to K_0(X) \to K_0(\textit{Vect}(X))$
是恒等映射。另一方面，设 $\mathcal{E}^\bullet$
是有限局部自由 $\mathcal{O}_X$-模的有界复形。
朴素截断的可区别三角
（见同调，第 \ref{homology-section-truncations} 节）
$$
\sigma_{\geq n}\mathcal{E}^\bullet \to
\sigma_{\geq n - 1}\mathcal{E}^\bullet \to
\mathcal{E}^{n - 1}[-n + 1] \to
(\sigma_{\geq n}\mathcal{E}^\bullet)[1]
$$
的存在性与归纳法表明
$$
[\mathcal{E}^\bullet] = \sum (-1)^i[\mathcal{E}^i[0]]
$$
在 $K_0(X) = K_0(D_{perf}(\mathcal{O}_X))$ 中成立；
仍请原谅这里的记号。因此映射
$K_0(\textit{Vect}(X)) \to K_0(D_{perf}(\mathcal{O}_X)) = K_0(X)$
是满射，证明完毕。
\end{proof}

\begin{remark}
\label{perfect-remark-K-ring}
设 $X$ 是概形。K 群 $K_0(X)$ 典范地是交换环。
具体地，利用导出张量积
$$
\otimes = \otimes^\mathbf{L}_{\mathcal{O}_X} :
D_{perf}(\mathcal{O}_X) \times D_{perf}(\mathcal{O}_X)
\longrightarrow
D_{perf}(\mathcal{O}_X)
$$
以及导出范畴，引理 \ref{derived-lemma-bilinear-map-K}，
得到双线性乘法。由于 $K \otimes L \cong L \otimes K$，
这个乘法是交换的。由于
$(K \otimes L) \otimes M = K \otimes (L \otimes M)$
，这个乘法是结合的。最后，$K_0(X)$ 的单位元是
$1 = [\mathcal{O}_X]$。

\medskip\noindent
若 $\textit{Vect}(X)$ 与 $K_0(\textit{Vect}(X))$ 如上，
则显然 $K_0(\textit{Vect}(X))$ 也有环结构：若
$\mathcal{E}$ 与 $\mathcal{F}$ 是有限局部自由
$\mathcal{O}_X$-模，则置
$$
[\mathcal{E}] \cdot [\mathcal{F}] =
[\mathcal{E} \otimes_{\mathcal{O}_X} \mathcal{F}]
$$
读者容易验证，这确实定义了双线性、交换且结合的乘法。
略去细节。上面构造的映射
$$
K_0(\textit{Vect}(X)) \longrightarrow K_0(X)
$$
对于这些定义是环映射。

\medskip\noindent
现假设 $X$ 是 Noether 概形。导出张量积还给出映射
$$
\otimes = \otimes^\mathbf{L}_{\mathcal{O}_X} :
D_{perf}(\mathcal{O}_X) \times D^b_{\textit{Coh}}(\mathcal{O}_X)
\longrightarrow
D^b_{\textit{Coh}}(\mathcal{O}_X)
$$
再次利用导出范畴，引理 \ref{derived-lemma-bilinear-map-K}，
得到双线性乘法 $K_0(X) \times K'_0(X) \to K'_0(X)$；
这是因为由引理 \ref{perfect-lemma-Noetherian-Kprime}，
$K'_0(X) = K_0(D^b_{\textit{Coh}}(\mathcal{O}_X))$。
读者容易证明，这赋予 $K'_0(X)$ 以环 $K_0(X)$ 上模的结构。
\end{remark}

\begin{remark}
\label{perfect-remark-pushforward-K}
设 $f : X \to Y$ 是局部 Noether 概形之间的固有态射。
存在映射
$$
f_* : K'_0(X) \longrightarrow K'_0(Y)
$$
，它把 $[\mathcal{F}]$ 送到
$$
[\bigoplus\nolimits_{i \geq 0} R^{2i}f_*\mathcal{F}] -
[\bigoplus\nolimits_{i \geq 0} R^{2i + 1}f_*\mathcal{F}]
$$
这是定义良好的，因为层 $R^if_*\mathcal{F}$ 是凝聚的
（概形上同调，命题
\ref{coherent-proposition-proper-pushforward-coherent}），
因为局部只有有限多个这样的层非零，也因为 $X$
上凝聚层的短正合列会给出 $Y$ 上 $R^if_*$ 的长正合列。
若 $Y$ 拟紧（这是实践中最常使用的情形），则上式可改写为
$$
f_*[\mathcal{F}] = \sum (-1)^i[R^if_*\mathcal{F}] = [Rf_*\mathcal{F}]
$$
，这里使用了引理 \ref{perfect-lemma-Noetherian-Kprime} 中的等式
$K'_0(Y) = K_0(D^b_{\textit{Coh}}(Y))$。
\end{remark}

\begin{lemma}
\label{perfect-lemma-projection-formula}
设 $f : X \to Y$ 是局部 Noether 概形之间的固有态射。
则对 $\alpha \in K'_0(X)$ 与 $\beta \in K_0(Y)$，有
$f_*(\alpha \cdot f^*\beta) = f_*\alpha \cdot \beta$。
\end{lemma}

\begin{proof}
这由引理 \ref{perfect-lemma-cohomology-base-change}、评注
\ref{perfect-remark-pushforward-K} 中的讨论，以及评注
\ref{perfect-remark-K-ring} 中乘法
$K'_0(X) \times K_0(X) \to K'_0(X)$ 的定义得出。
\end{proof}

\begin{remark}
\label{perfect-remark-perf-Z}
设 $X$ 是概形，$Z \subset X$ 是闭子概形。考虑严格满、饱和的三角子范畴
$$
D_{Z, perf}(\mathcal{O}_X) \subset D(\mathcal{O}_X)
$$
，其对象是这样的 $\mathcal{O}_X$-模完美复形：
它们的同调层在集合论意义下支集于 $Z$。
这个三角范畴的第零 $K$ 群
$K_0(D_{Z, perf}(\mathcal{O}_X))$
有时记作 $K_Z(X)$ 或 $K_{0, Z}(X)$。
完全如评注 \ref{perfect-remark-K-ring} 中那样利用导出张量积，
可知 $K_0(D_{Z, perf}(\mathcal{O}_X))$
具有结合且交换的乘法，但一般
$K_0(D_{Z, perf}(\mathcal{O}_X))$ 没有单位元。
\end{remark}













\section{复形的行列式}
\label{perfect-section-det-two-terms}

\noindent
本节续接代数续篇，第
\ref{more-algebra-section-determinants-complexes} 节。
对任意赋环空间 $(X, \mathcal{O}_X)$，存在函子
$$
\det :
\left\{
\begin{matrix}
\text{完美复形范畴} \\
\text{态射为同构}
\end{matrix}
\right\}
\longrightarrow
\left\{
\begin{matrix}
\text{可逆模范畴} \\
\text{态射为同构}
\end{matrix}
\right\}
$$
此外，给定滤过导出范畴 $DF(\mathcal{O}_X)$ 的对象 $(L, F)$，
若其滤过有限且各分次部分都是完美复形，则存在典范同构
$\det(\text{gr}L) \to \det(L)$。原始论述见
\cite{determinant}。我们稍后会加入这一材料（插入将来的引用）。

\medskip\noindent
目前先给出如下特设构造：$X$ 是概形，而我们考虑
$D(\mathcal{O}_X)$ 中 Tor 振幅位于 $[-1, 0]$ 的完美对象 $L$。

\begin{lemma}
\label{perfect-lemma-determinant-two-term-complexes}
设 $X$ 是概形。存在函子
$$
\det :
\left\{
\begin{matrix}
\text{Tor 振幅位于 }[-1, 0]\text{ 的} \\
\text{完美复形范畴} \\
\text{态射为同构}
\end{matrix}
\right\}
\longrightarrow
\left\{
\begin{matrix}
\text{可逆模范畴} \\
\text{态射为同构}
\end{matrix}
\right\}
$$
此外，给定 $D(\mathcal{O}_X)$ 中秩为 $0$、Tor 振幅位于
$[-1, 0]$ 的完美对象 $L$，存在典范元素
$\delta(L) \in \Gamma(X, \det(L))$，使得对
$D(\mathcal{O}_X)$ 中任意同构 $a : L \to K$，都有
$\det(a)(\delta(L)) = \delta(K)$。
而且此构造在仿射局部由代数续篇，第
\ref{more-algebra-section-determinants-complexes} 节中的构造给出。
\end{lemma}

\begin{proof}
设 $L$ 是左侧范畴的对象。若 $\Spec(A) = U \subset X$
是仿射开集，则 $L|_U$ 对应于 Tor 振幅位于 $[-1, 0]$
的 $A$-模完美复形 $L^\bullet$，见诸引理
\ref{perfect-lemma-affine-compare-bounded}、
\ref{perfect-lemma-tor-dimension-affine} 与
\ref{perfect-lemma-perfect-affine}。
于是可以考虑代数续篇，引理
\ref{more-algebra-lemma-determinant-two-term-complexes}
中构造的可逆 $A$-模 $\det(L^\bullet)$。
若 $\Spec(B) = V \subset U$ 是包含于 $U$ 的另一仿射开集，则
$\det(L^\bullet) \otimes_A B = \det(L^\bullet \otimes_A B)$，
所以此构造与限制映射相容
（见引理 \ref{perfect-lemma-quasi-coherence-pullback}，并注意 $A \to B$ 平坦）。
因此可将这些可逆模黏合，得到 $X$ 上的可逆模 $\det(L)$。
函子性与典范截面完全以同样方式构造。略去细节。
\end{proof}

\begin{remark}
\label{perfect-remark-functorial-det}
引理 \ref{perfect-lemma-determinant-two-term-complexes} 的构造与拉回相容。
更确切地，给定概形态射 $f : X \to Y$ 以及
$D(\mathcal{O}_Y)$ 中 Tor 振幅位于 $[-1, 0]$ 的完美对象 $K$，
则 $Lf^*K$ 是 $D(\mathcal{O}_X)$ 中 Tor 振幅位于
$[-1, 0]$ 的完美对象 $K$，并且有典范等同
$$
f^*\det(K) \longrightarrow \det(Lf^*K)
$$
此外，若 $K$ 的秩为 $0$，则在此映射下，
$\delta(K)$ 拉回为 $\delta(Lf^*K)$。
这由行列式的仿射局部构造显然可见。
\end{remark}






\section{判定有界性}
\label{perfect-section-detecting-boundedness}

\noindent
本节证明，第 \ref{perfect-section-generating} 节所构造的拟紧拟分离概形
之 $D_\QCoh$ 的紧生成元具有一项特殊性质。
由于该节与本节十分相似，建议先阅读该节。

\begin{lemma}
\label{perfect-lemma-orthogonal-koszul-first-variant}
在环境 \ref{perfect-situation-complex} 中，将开浸入记作
$j : U \to X$，并令 $K$ 为 $D(\mathcal{O}_X)$ 中对应于
$A$ 上由 $f_1, \ldots, f_r$ 所给 Koszul 复形的完美对象。
设 $E \in D_\QCoh(\mathcal{O}_X)$ 且 $a \in \mathbf{Z}$。
考虑下列条件：
\begin{enumerate}
\item 典范映射 $\tau_{\geq a}E \to \tau_{\geq a} Rj_*(E|_U)$
是同构。
\item 对所有 $n \geq a$，有
$\Hom_{D(\mathcal{O}_X)}(K[-n], E) = 0$。
\end{enumerate}
则 (2) 蕴含 (1)，而 (1) 蕴含把 $a$ 替换为 $a + 1$ 后的 (2)。
\end{lemma}

\begin{proof}
选取可区别三角 $N \to E \to Rj_*(E|_U) \to N[1]$。
于是 (1) 蕴含 $\tau_{\geq a + 1} N = 0$，而
$\tau_{\geq a}N = 0$ 蕴含 (1)。注意
$$
\Hom_{D(\mathcal{O}_X)}(K[-n], Rj_*(E|_U)) =
\Hom_{D(\mathcal{O}_U)}(K|_U[-n], E) = 0
$$
对所有 $n$ 都成立，因为 $K|_U = 0$。所以 (2) 等价于：
对所有 $n \geq a$，有
$\Hom_{D(\mathcal{O}_X)}(K[-n], N) = 0$。
注意，存在 Koszul 复形的可区别三角
$$
K^\bullet(f_1^{e_1}, \ldots, f_i^{e'_i}, \ldots, f_r^{e_r}) \to
K^\bullet(f_1^{e_1}, \ldots, f_i^{e'_i + e''_i}, \ldots, f_r^{e_r}) \to
K^\bullet(f_1^{e_1}, \ldots, f_i^{e''_i}, \ldots, f_r^{e_r}) \to \ldots
$$
，见代数续篇，引理 \ref{more-algebra-lemma-koszul-mult}。
因此，对所有 $n \geq a$ 有
$\Hom_{D(\mathcal{O}_X)}(K[-n], N) = 0$，等价于：
对所有 $n \geq a$ 及所有 $e \geq 1$，有
$\Hom_{D(\mathcal{O}_X)}(K_e[-n], N) = 0$，
其中 $K_e$ 如引理
\ref{perfect-lemma-represent-cohomology-class-on-closed} 所述。
由于 $N|_U = 0$，该引理又表明，这等价于当 $n \geq a$ 时
$H^n(X, N) = 0$。故 (2) 等价于 $\tau_{\geq a}N = 0$，
因为 $N$ 由 $A$-模复形 $R\Gamma(X, N)$ 决定，见引理
\ref{perfect-lemma-affine-compare-bounded}。引理得证。
\end{proof}

\begin{lemma}
\label{perfect-lemma-orthogonal-koszul-second-variant}
在环境 \ref{perfect-situation-complex} 中，将开浸入记作
$j : U \to X$，并令 $K$ 为 $D(\mathcal{O}_X)$ 中对应于
$A$ 上由 $f_1, \ldots, f_r$ 所给 Koszul 复形的完美对象。
设 $E \in D_\QCoh(\mathcal{O}_X)$ 且 $a \in \mathbf{Z}$。
考虑下列条件：
\begin{enumerate}
\item 典范映射 $\tau_{\leq a}E \to \tau_{\leq a} Rj_*(E|_U)$
是同构；并且
\item $\Hom_{D(\mathcal{O}_X)}(K[-n], E) = 0$ for all $n \leq a$.
\end{enumerate}
则 (2) 蕴含 (1)，而 (1) 蕴含把 $a$ 替换为 $a - 1$ 后的 (2)。
\end{lemma}

\begin{proof}
选取可区别三角 $E \to Rj_*(E|_U) \to N \to E[1]$。
于是 (1) 蕴含 $\tau_{\leq a - 1}N = 0$，而
$\tau_{\leq a}N = 0$ 蕴含 (1)。注意
$$
\Hom_{D(\mathcal{O}_X)}(K[-n], Rj_*(E|_U)) =
\Hom_{D(\mathcal{O}_U)}(K|_U[-n], E) = 0
$$
对所有 $n$ 都成立，因为 $K|_U = 0$。所以 (2) 等价于：
对所有 $n \leq a$，有
$\Hom_{D(\mathcal{O}_X)}(K[-n], N) = 0$。
注意，存在 Koszul 复形的可区别三角
$$
K^\bullet(f_1^{e_1}, \ldots, f_i^{e'_i}, \ldots, f_r^{e_r}) \to
K^\bullet(f_1^{e_1}, \ldots, f_i^{e'_i + e''_i}, \ldots, f_r^{e_r}) \to
K^\bullet(f_1^{e_1}, \ldots, f_i^{e''_i}, \ldots, f_r^{e_r}) \to \ldots
$$
，见代数续篇，引理 \ref{more-algebra-lemma-koszul-mult}。
因此，对所有 $n \leq a$ 有
$\Hom_{D(\mathcal{O}_X)}(K[-n], N) = 0$，等价于：
对所有 $n \leq a$ 及所有 $e \geq 1$，有
$\Hom_{D(\mathcal{O}_X)}(K_e[-n], N) = 0$，
其中 $K_e$ 如引理
\ref{perfect-lemma-represent-cohomology-class-on-closed} 所述。
由于 $N|_U = 0$，该引理又表明，这等价于当 $n \leq a$ 时
$H^n(X, N) = 0$。故 (2) 等价于 $\tau_{\leq a}N = 0$，
因为 $N$ 由 $A$-模复形 $R\Gamma(X, N)$ 决定，见引理
\ref{perfect-lemma-affine-compare-bounded}。引理得证。
\end{proof}

\begin{lemma}
\label{perfect-lemma-bounded-truncation}
设 $X$ 是拟紧拟分离概形。
设 $P \in D_{perf}(\mathcal{O}_X)$、
$E \in D_{\QCoh}(\mathcal{O}_X)$ 且 $a \in \mathbf{Z}$。
下列条件等价：
\begin{enumerate}
\item $\Hom_{D(\mathcal{O}_X)}(P[-i], E) = 0$ for $i \gg 0$, and
\item $\Hom_{D(\mathcal{O}_X)}(P[-i], \tau_{\geq a} E) = 0$ for $i \gg 0$.
\end{enumerate}
\end{lemma}

\begin{proof}
利用三角 $ \tau_{< a} E \to E \to \tau_{\geq a} E \to$，
可知只需证明
$$
\Hom_{D(\mathcal{O}_X)}(P[-i], \tau_{< a} E) =
\Hom_{D(\mathcal{O}_X)}(P, (\tau_{< a} E)[i]) = 0 
$$
当 $i \gg 0$ 时成立。由于 $P$ 完美，这由引理
\ref{perfect-lemma-ext-from-perfect-into-bounded-QCoh} 得出。
\end{proof}

\begin{lemma}
\label{perfect-lemma-bounded-below-truncation}
设 $X$ 是拟紧拟分离概形。设
$P \in D_{perf}(\mathcal{O}_X)$、
$E \in D_{\QCoh}(\mathcal{O}_X)$ 且 $a \in \mathbf{Z}$。
下列条件等价：
\begin{enumerate}
\item $\Hom_{D(\mathcal{O}_X)}(P[-i], E) = 0$ for $i \ll 0$, and
\item $\Hom_{D(\mathcal{O}_X)}(P[-i], \tau_{\leq a} E) = 0$ for $i \ll 0$.
\end{enumerate}
\end{lemma}

\begin{proof}
利用三角 $ \tau_{\leq a} E \to E \to \tau_{> a} E \to$，
可知只需证明
$$
\Hom_{D(\mathcal{O}_X)}(P[-i], \tau_{> a} E) =
\Hom_{D(\mathcal{O}_X)}(P, (\tau_{> a} E)[i]) = 0
$$
当 $i \ll 0$ 时成立。由于 $P$ 完美，这由引理
\ref{perfect-lemma-ext-from-perfect-into-bounded-QCoh} 得出。
\end{proof}

\begin{proposition}
\label{perfect-proposition-detecting-bounded-above}
设 $X$ 是拟紧拟分离概形。设
$G \in D_{perf}(\mathcal{O}_X)$ 是生成
$D_\QCoh (\mathcal{O}_X)$ 的完美复形，并设
$E \in D_\QCoh (\mathcal{O}_X)$。下列条件等价：
\begin{enumerate}
\item $E \in D^-_\QCoh (\mathcal{O}_X)$,
\item $\Hom_{D(\mathcal{O}_X)}(G[-i], E) = 0$ for $i \gg 0$,
\item $\Ext^i_X(G, E) = 0$ for $i \gg 0$,
\item $R\Hom_X(G, E)$ 属于 $D^-(\mathbf{Z})$；
\item $H^i(X, G^\vee \otimes_{\mathcal{O}_X}^\mathbf{L} E) = 0$
，其中 $i \gg 0$；
\item $R\Gamma(X, G^\vee \otimes_{\mathcal{O}_X}^\mathbf{L} E)$
属于 $D^-(\mathbf{Z})$；
\item 对 $D(\mathcal{O}_X)$ 的每个完美对象 $P$，
\begin{enumerate}
\item 把 $G$ 替换为 $P$ 后，断言 (2)、(3)、(4) 成立；
\item $H^i(X, P \otimes_{\mathcal{O}_X}^\mathbf{L} E) = 0$ for $i \gg 0$,
\item $R\Gamma(X, P \otimes_{\mathcal{O}_X}^\mathbf{L} E)$
属于 $D^-(\mathbf{Z})$。
\end{enumerate}
\end{enumerate}
\end{proposition}

\begin{proof}
假设 (1) 成立。由于
$\Hom_{D(\mathcal{O}_X)}(G[-i], E) = \Hom_{D(\mathcal{O}_X)}(G, E[i])$
，由引理 \ref{perfect-lemma-ext-from-perfect-into-bounded-QCoh}，
当 $i \gg 0$ 时该群为零。这证明了 (1) 蕴含 (2)。

\medskip\noindent
由同调，第 \ref{cohomology-section-global-RHom} 节中的讨论，
(2)、(3)、(4) 等价。按定义
$H^i(X, -) = H^i(R\Gamma(X, -))$，故 (5) 与 (6) 等价。
由同调，引理 \ref{cohomology-lemma-dual-perfect-complex}
以及 $(G[-i])^\vee = G^\vee[i]$ 这一事实，
彼此等价的条件 (2)、(3)、(4) 又等价于彼此等价的条件 (5)、(6)。

\medskip\noindent
显然 (7) 蕴含 (2)。反过来，证明彼此等价的条件
(2) -- (6) 蕴含 (7)。由评注 \ref{perfect-remark-classical-generator}，
$G$ 是 $D_{perf}(\mathcal{O}_X)$ 的经典生成元。
对 $P \in D_{perf}(\mathcal{O}_X)$，令 $T(P)$ 表示断言：
$R\Hom_X(P, E)$ 属于 $D^-(\mathbf{Z})$。
显然，$T$ 对直和遗传，对可区别三角满足三取二性质，
对直和项遗传，并在平移下保持。因此由导出范畴，评注
\ref{derived-remark-check-on-generator}，(4) 蕴含 $T$
在整个 $D_{perf}(\mathcal{O}_X)$ 上成立。
同一论证适用于其余所有性质；不过对性质 (7)(b) 与 (7)(c)，
还要使用 $P \mapsto P^\vee$ 是 $D_{perf}(\mathcal{O}_X)$
的自等价。略去这一小细节。

\medskip\noindent
我们将利用概形上同调，引理
\ref{coherent-lemma-induction-principle} 的归纳原理，
证明彼此等价的条件 (2) -- (7) 蕴含 (1)。

\medskip\noindent
首先在 $X$ 仿射时证明 (2) -- (7) $\Rightarrow$ (1)。
置 $P = \mathcal{O}_X[0]$。由 (7)，当 $i \gg 0$ 时
$H^i (X, E) = 0$。由于 $E$ 由 $R\Gamma (X, E)$ 决定，
(1) 随即成立，见引理 \ref{perfect-lemma-affine-compare-bounded}。

\medskip\noindent
现设 $X = U \cup V$，其中 $U$ 是 $X$ 的拟紧开集，
$V$ 是仿射开集；并假设对概形 $U$、$V$ 与 $U \cap V$，
已知蕴含关系 (2) -- (7) $\Rightarrow$ (1)。
设 $E \in D_\QCoh(\mathcal{O}_X)$ 满足 (2) -- (7)。
由引理 \ref{perfect-lemma-direct-summand-of-a-restriction} 与定理
\ref{perfect-theorem-bondal-van-den-Bergh}，存在 $X$ 上完美复形 $Q$，
使得 $Q|_U$ 生成 $D_\QCoh (\mathcal{O}_U)$。
取 $f_1, \dots , f_r \in \Gamma (V, \mathcal{O}_V)$，
使得作为 $V$ 的子集有
$V \setminus U = V(f_1, \dots , f_r)$。
令 $K \in D_{perf}(\mathcal{O}_V)$ 为对应于
$f_1, \dots , f_r$ 上 Koszul 复形的对象，并令
$K' \in D_{perf}(\mathcal{O}_X)$ 为
\begin{equation}
\label{perfect-equation-detecting-bounded-above}
K' = R (V \to X)_* K = R (V \to X)_! K,
\end{equation}
，见同调，诸引理 \ref{cohomology-lemma-pushforward-restriction}
与 \ref{cohomology-lemma-pushforward-perfect}。
这是 $X$ 上支集于闭集 $X \setminus U \subset V$
的完美复形，并且在 $V$ 上同构于 $K$。由假设，
$R\Hom_{\mathcal{O}_X}(Q, E)$ 与
$R\Hom_{\mathcal{O}_X}(K', E)$ 都上有界。

\medskip\noindent
由 (\ref{perfect-equation-detecting-bounded-above}) 中对 $K'$ 的第二种描述，有
$$
\Hom_{D(\mathcal{O}_V)}(K[-i], E|_V) = \Hom_{D(\mathcal{O}_X)}(K'[-i], E) = 0
$$
，其中 $i \gg 0$。因此可对 $E|_V$ 应用引理
\ref{perfect-lemma-orthogonal-koszul-first-variant}，得到整数 $a$，使得
$\tau_{\geq a}(E|_V) = \tau_{\geq a} R (U \cap V \to V)_* (E|_{U \cap V})$.
于是 $\tau_{\geq a} E = \tau_{\geq a} R (U \to X)_* (E |_U)$
（限制到 $U$ 与 $V$ 后检查典范映射是同构）。
因而两次使用引理 \ref{perfect-lemma-bounded-truncation} 可知
$$
\Hom_{D(\mathcal{O}_U)}(Q|_U [-i], E|_U) =
\Hom_{D(\mathcal{O}_X)}(Q[-i], R (U \to X)_* (E|_U)) = 0
$$
，其中 $i \gg 0$。由于本命题对 $U$ 及生成元 $Q|_U$ 成立，
有 $E|_U \in D^-_\QCoh(\mathcal{O}_U)$。
又因函子 $R (U \to X)_*$ 保持 $D^-_\QCoh$
（由引理 \ref{perfect-lemma-quasi-coherence-direct-image}），得到
$\tau_{\geq a}E \in D^-_\QCoh(\mathcal{O}_X)$。
因此 $E \in D^-_\QCoh (\mathcal{O}_X)$。
\end{proof}

\begin{proposition}
\label{perfect-proposition-detecting-bounded-below}
设 $X$ 是拟紧拟分离概形。设
$G \in D_{perf}(\mathcal{O}_X)$
是生成 $D_\QCoh (\mathcal{O}_X)$ 的完美复形，并设
$E \in D_\QCoh (\mathcal{O}_X)$。下列条件等价：
\begin{enumerate}
\item $E \in D^+_\QCoh (\mathcal{O}_X)$,
\item $\Hom_{D(\mathcal{O}_X)}(G[-i], E) = 0$ for $i \ll 0$,
\item $\Ext^i_X(G, E) = 0$ for $i \ll 0$,
\item $R\Hom_X(G, E)$ 属于 $D^+(\mathbf{Z})$；
\item $H^i(X, G^\vee \otimes_{\mathcal{O}_X}^\mathbf{L} E) = 0$
，其中 $i \ll 0$；
\item $R\Gamma(X, G^\vee \otimes_{\mathcal{O}_X}^\mathbf{L} E)$
属于 $D^+(\mathbf{Z})$；
\item 对 $D(\mathcal{O}_X)$ 的每个完美对象 $P$，
\begin{enumerate}
\item 把 $G$ 替换为 $P$ 后，断言 (2)、(3)、(4) 成立；
\item $H^i(X, P \otimes_{\mathcal{O}_X}^\mathbf{L} E) = 0$ for $i \ll 0$,
\item $R\Gamma(X, P \otimes_{\mathcal{O}_X}^\mathbf{L} E)$
属于 $D^+(\mathbf{Z})$。
\end{enumerate}
\end{enumerate}
\end{proposition}

\begin{proof}
假设 (1) 成立。由于
$\Hom_{D(\mathcal{O}_X)}(G[-i], E) = \Hom_{D(\mathcal{O}_X)}(G, E[i])$
，由引理 \ref{perfect-lemma-ext-from-perfect-into-bounded-QCoh}，
当 $i \ll 0$ 时该群为零。这证明了 (1) 蕴含 (2)。

\medskip\noindent
由同调，第 \ref{cohomology-section-global-RHom} 节中的讨论，
(2)、(3)、(4) 等价。按定义
$H^i(X, -) = H^i(R\Gamma(X, -))$，故 (5) 与 (6) 等价。
由同调，引理 \ref{cohomology-lemma-dual-perfect-complex}
以及 $(G[-i])^\vee = G^\vee[i]$ 这一事实，
彼此等价的条件 (2)、(3)、(4) 又等价于彼此等价的条件 (5)、(6)。

\medskip\noindent
显然 (7) 蕴含 (2)。反过来，证明彼此等价的条件
(2) -- (6) 蕴含 (7)。由评注 \ref{perfect-remark-classical-generator}，
$G$ 是 $D_{perf}(\mathcal{O}_X)$ 的经典生成元。
对 $P \in D_{perf}(\mathcal{O}_X)$，令 $T(P)$ 表示断言：
$R\Hom_X(P, E)$ 属于 $D^+(\mathbf{Z})$。
显然，$T$ 对直和遗传，对可区别三角满足三取二性质，
对直和项遗传，并在平移下保持。因此由导出范畴，评注
\ref{derived-remark-check-on-generator}，(4) 蕴含 $T$
在整个 $D_{perf}(\mathcal{O}_X)$ 上成立。
同一论证适用于其余所有性质；不过对性质 (7)(b) 与 (7)(c)，
还要使用 $P \mapsto P^\vee$ 是 $D_{perf}(\mathcal{O}_X)$
的自等价。略去这一小细节。

\medskip\noindent
我们将利用概形上同调，引理
\ref{coherent-lemma-induction-principle} 的归纳原理，
证明彼此等价的条件 (2) -- (7) 蕴含 (1)。

\medskip\noindent
首先在 $X$ 仿射时证明 (2) -- (7) $\Rightarrow$ (1)。
取 $P = \mathcal{O}_X[0]$。由 (7)，当 $i \ll 0$ 时
$H^i (X, E) = 0$。由于 $E$ 由 $R\Gamma (X, E)$ 决定，
(1) 随即成立，见引理 \ref{perfect-lemma-affine-compare-bounded}。

\medskip\noindent
现设 $X = U \cup V$，其中 $U$ 是 $X$ 的拟紧开集，
$V$ 是仿射开集；并假设对概形 $U$、$V$ 与 $U \cap V$，
已知蕴含关系 (2) -- (7) $\Rightarrow$ (1)。
设 $E \in D_\QCoh(\mathcal{O}_X)$ 满足 (2) -- (7)。
由引理 \ref{perfect-lemma-direct-summand-of-a-restriction} 与定理
\ref{perfect-theorem-bondal-van-den-Bergh}，存在 $X$ 上完美复形 $Q$，
使得 $Q|_U$ 生成 $D_\QCoh (\mathcal{O}_U)$。
取 $f_1, \dots , f_r \in \Gamma (V, \mathcal{O}_V)$，
使得作为 $V$ 的子集有
$V \setminus U = V(f_1, \dots , f_r)$。
令 $K \in D_{perf}(\mathcal{O}_V)$ 为对应于
$f_1, \dots , f_r$ 上 Koszul 复形的对象，并令
$K' \in D_{perf}(\mathcal{O}_X)$ 为
\begin{equation}
\label{perfect-equation-detecting-bounded-below}
K' = R (V \to X)_* K = R (V \to X)_! K,
\end{equation}
，见同调，诸引理 \ref{cohomology-lemma-pushforward-restriction}
与 \ref{cohomology-lemma-pushforward-perfect}。
这是 $X$ 上支集于闭集 $X \setminus U \subset V$
的完美复形，并且在 $V$ 上同构于 $K$。由假设，
$R\Hom_{\mathcal{O}_X}(Q, E)$ 与
$R\Hom_{\mathcal{O}_X}(K', E)$ 都下有界。

\medskip\noindent
由 (\ref{perfect-equation-detecting-bounded-below}) 中对 $K'$ 的第二种描述，有
$$
\Hom_{D(\mathcal{O}_V)}(K[-i], E|_V) = \Hom_{D(\mathcal{O}_X)}(K'[-i], E) = 0
$$
，其中 $i \ll 0$。因此可对 $E|_V$ 应用引理
\ref{perfect-lemma-orthogonal-koszul-second-variant}，得到整数 $a$，使得
$\tau_{\leq a}(E|_V) = \tau_{\leq a} R(U \cap V \to V)_*(E|_{U \cap V})$.
于是 $\tau_{\leq a} E = \tau_{\leq a} R(U \to X)_*(E|_U)$
（限制到 $U$ 与 $V$ 后检查典范映射是同构）。
因而两次使用引理 \ref{perfect-lemma-bounded-below-truncation} 可知
$$
\Hom_{D(\mathcal{O}_U)}(Q|_U [-i], E|_U) =
\Hom_{D(\mathcal{O}_X)}(Q[-i], R (U \to X)_* (E|_U)) = 0
$$
，其中 $i \ll 0$。由于本命题对 $U$ 及生成元 $Q|_U$ 成立，
有 $E|_U \in D^+_\QCoh (\mathcal{O}_U)$。
又因函子 $R(U \to X)_*$ 保持下有界对象
（见同调，第 \ref{cohomology-section-derived-functors} 节），得到
$\tau_{\leq a} E \in D^+_\QCoh(\mathcal{O}_X)$。
因此 $E \in D^+_\QCoh (\mathcal{O}_X)$。
\end{proof}










\section{导出范畴中的拟凝聚对象}
\label{perfect-section-QC}

\noindent
设 $X$ 是概形。回忆 $X_{affine, Zar}$ 表示 $X$
的仿射开集所成的范畴，其拓扑由标准 Zariski 覆盖给出，
见拓扑，定义 \ref{topologies-definition-big-small-Zariski}。
提醒读者，$X_{affine, Zar}$ 的拓扑斯就是 $X$ 的小
Zariski 拓扑斯，见拓扑，引理
\ref{topologies-lemma-alternative-zariski}。
位点 $X_{affine, Zar}$ 带有结构层 $\mathcal{O}$，
并且存在赋环拓扑斯的等价
$$
(\Sh(X_{affine, Zar}), \mathcal{O})
\longrightarrow
(\Sh(X_{Zar}), \mathcal{O})
$$
见下降，等式 (\ref{descent-equation-alternative-small-ringed})
及第 \ref{descent-section-alternative-quasi-coherent} 节中围绕该式的讨论；
那里使用了略有不同的记号。

\medskip\noindent
本节以 $X_{affine}$ 表示 $X_{affine, Zar}$
的底层范畴，并赋予它混沌拓扑，即使层与预层相同的拓扑。
特别地，结构层 $\mathcal{O}$ 也成为 $X_{affine}$ 上的层。
由此得到赋环位点态射
$$
\epsilon :
(X_{affine, Zar}, \mathcal{O})
\longrightarrow
(X_{affine}, \mathcal{O})
$$
，如位点上同调，第 \ref{sites-cohomology-section-compare} 节所述。
本节将把 $D_\QCoh(\mathcal{O}_X)$ 与位点上同调，第
\ref{sites-cohomology-section-modules-cohomology} 节中引入的范畴
$\mathit{QC}(X_{affine}, \mathcal{O})$ 等同起来。

\begin{lemma}
\label{perfect-lemma-DQCoh-alternative-small}
在上述情形中，下列三角范畴之间存在典范正合等价：
\begin{enumerate}
\item $D_\QCoh(\mathcal{O}_X)$,
\item $D_\QCoh(X_{Zar}, \mathcal{O})$,
\item $D_\QCoh(X_{affine, Zar}, \mathcal{O})$,
\item $D_\QCoh(X_{affine}, \mathcal{O}_X)$, and
\item $\mathit{QC}(X_{affine}, \mathcal{O})$.
\end{enumerate}
\end{lemma}

\begin{proof}
若 $U \subset V \subset X$ 是仿射开集，则环映射
$\mathcal{O}(V) \to \mathcal{O}(U)$ 平坦。
因此 (4) 与 (5) 之间的等价是位点上同调，引理
\ref{sites-cohomology-lemma-cartesian-flat-quasi-coherent}
的特例（该证明也阐明了陈述的含义）。

\medskip\noindent
由下降，评注 \ref{descent-remark-Zariski-site-space}，
赋环位点 $(X_{Zar}, \mathcal{O})$ 与赋环空间
$(X, \mathcal{O}_X)$ 具有相同的模范畴。
由下降，命题 \ref{descent-proposition-equivalence-quasi-coherent}，
在此等价下拟凝聚模彼此对应。
所以得到 (1) 与 (2) 中三角范畴之间的典范正合等价。

\medskip\noindent
引理之前的讨论表明，存在赋环拓扑斯的等价
$(\Sh(X_{affine, Zar}), \mathcal{O}) \to (\Sh(X_{Zar}), \mathcal{O})$，
从而存在模的阿贝尔范畴之间的等价。
由于拟凝聚模的概念是内在的
（位点上的模，引理 \ref{sites-modules-lemma-special-locally-free}），
这个等价保持拟凝聚模子范畴。
因此得到 (2) 与 (3) 中三角范畴之间的典范正合等价。

\medskip\noindent
为得到 (3) 与 (4) 中三角范畴之间的正合等价，
我们把位点上同调，引理
\ref{sites-cohomology-lemma-compare-topologies-derived-adequate-modules}
应用于上述态射
$\epsilon : (X_{affine, Zar}, \mathcal{O}) \to
(X_{affine}, \mathcal{O})$。取
$\mathcal{B} = \Ob(X_{affine})$，并令
$\mathcal{A} \subset \textit{PMod}(X_{affine}, \mathcal{O})$
为这样的预层 $\mathcal{F}$ 所成的满子范畴：
$\mathcal{F}(V) \otimes_{\mathcal{O}(V)} \mathcal{O}(U)
\to \mathcal{F}(U)$ 是同构。注意，由下降，引理
\ref{descent-lemma-quasi-coherent-alternative-small}，
$\mathcal{A}$ 的对象恰好是 Zariski 拓扑下这样的层：
它们是 $(X_{affine, Zar}, \mathcal{O})$ 上的拟凝聚模。
另一方面，由位点上的模，引理
\ref{sites-modules-lemma-chaotic-quasi-coherent}，
$\mathcal{A}$ 的对象也恰好是 $(X_{affine}, \mathcal{O})$
上，即混沌拓扑下的拟凝聚模。因此，只要证明位点上同调，引理
\ref{sites-cohomology-lemma-compare-topologies-derived-adequate-modules}
可以应用，就确实可利用 $\epsilon^*$ 与 $R\epsilon_*$
得到 (3) 与 (4) 中范畴之间的典范等价。

\medskip\noindent
需要验证四个条件：
\begin{enumerate}
\item $\mathcal{A}$ 的每个对象都是 Zariski 拓扑下的层。
上面已经看到这一点。
\item $\mathcal{A}$ 是
$\textit{Mod}(X_{affine, Zar}, \mathcal{O})$ 的弱 Serre 子范畴。
上面已经看到
$\mathcal{A} = \QCoh(X_{affine, Zar}, \mathcal{O})$，
并且在等价
$\textit{Mod}(X_{affine, Zar}, \mathcal{O}) =
\textit{Mod}(X, \mathcal{O}_X)$
下，这些对象对应于 $X$ 上的拟凝聚模。因此结论由概形，第
\ref{schemes-section-quasi-coherent} 节中的讨论得出。
\item $X_{affine}$ 的每个对象在混沌拓扑下都有一个覆盖，
其成员属于 $\mathcal{B}$。这是因为 $\mathcal{B}$ 包含所有对象。
\item 对 $X_{affine}$ 的每个对象 $U$ 以及
$\mathcal{A}$ 中的每个 $\mathcal{F}$，当 $p > 0$ 时有
$H^p_{Zar}(U, \mathcal{F}) = 0$。
这由仿射概形上拟凝聚模的同调消灭性得出，见概形上同调，引理
\ref{coherent-lemma-quasi-coherent-affine-cohomology-zero}。
\end{enumerate}
证明完毕。
\end{proof}

\begin{remark}
\label{perfect-remark-QC-compare}
设 $S$ 是概形。稍后还将证明
$\mathit{QC}((\textit{Aff}/S), \mathcal{O})$
典范等价于 $D_\QCoh(\mathcal{O}_S)$。
见叠上的层，命题 \ref{stacks-sheaves-proposition-QC-compare}。
\end{remark}
% Generated mechanically from the frozen more-morphisms.tex target.
% Do not edit this generated file; edit the bound locale source instead.

% OK, start here.
%

\setcounter{chapter}{36}
\chapter{态射续篇}



\phantomsection
\label{more-morphisms-section-phantom}


\section{引言}
\label{more-morphisms-section-introduction}

\noindent
本章继续研究概形态射的性质。基本参考文献是 \cite{EGA}。






\section{加厚}
\label{more-morphisms-section-thickenings}

\noindent
下列术语未必完全标准，但使用起来很方便。

\begin{definition}
\label{more-morphisms-definition-thickening}
加厚。
\begin{enumerate}
\item 若 $X$ 是 $X'$ 的闭子概形，且二者的底层拓扑空间相同，
则称概形 $X'$ 是概形 $X$ 的一个{\it 加厚}。
\item 若 $X$ 是 $X'$ 的闭子概形，且定义 $X$ 的拟凝聚理想层
$\mathcal{I} \subset \mathcal{O}_{X'}$ 平方为零，
则称概形 $X'$ 是概形 $X$ 的一个{\it 一阶加厚}。
\item 若 $X$ 是 $X'$ 的闭子概形，且定义 $X$ 的拟凝聚理想层
$\mathcal{I} \subset \mathcal{O}_{X'}$ 幂零，即存在整数
$n \geq 0$ 使得 $\mathcal{I}^{n + 1} = 0$，
则称概形 $X'$ 是概形 $X$ 的一个{\it 有限阶加厚}。
\item 若 $X$ 是 $X$ 的闭子概形，且
$\mathcal{I}^{n + 1} = 0$，其中
$\mathcal{I} \subset \mathcal{O}_{X'}$ 是定义 $X$
的拟凝聚理想层，则称概形 $X'$ 是概形 $X$ 的一个
{\it $n$ 阶加厚}。
% P05-ERR-CH37-0001: frozen source says X is a closed subscheme of X.
\item 给定两个加厚 $X \subset X'$ 与 $Y \subset Y'$，
所谓{\it 加厚的态射}，是满足 $f'(X) \subset Y$
的态射 $f' : X' \to Y'$；换言之，$f'|_X$
经闭子概形 $Y$ 分解。在此情形中，置
$f = f'|_X : X \to Y$，并称
$(f, f') : (X \subset X') \to (Y \subset Y')$
为加厚的态射。
\item 设 $S$ 是概形。类似地定义{\it $S$ 上的加厚}及
{\it $S$ 上加厚的态射}。这意味着上述概形
$X, X', Y, Y'$ 都是 $S$-概形，而且态射
$X \to X'$、$Y \to Y'$ 与 $f' : X' \to Y'$ 都是 $S$-态射。
\end{enumerate}
\end{definition}

\noindent
设 $i_X : X \to X'$ 是加厚。核
$\mathcal{I} = \Ker(i_X^\sharp)$ 的任意局部截面都是局部幂零的。
这使我们能在一定程度上控制 $\mathcal{O}_{X'}$ 的结构，
但从技术上说，有限阶加厚 $X \subset X'$ 这一类要容易处理得多。
具体地，若 $X'$ 是 $n$ 阶加厚，则有滤过
% P05-ERR-CH37-0002: frozen source says “nth order thinking”.
$$
0 = \mathcal{I}^{n + 1} \subset
\mathcal{I}^n \subset
\mathcal{I}^{n - 1} \subset \ldots \subset
\mathcal{I} \subset \mathcal{O}_{X'}
$$
并且可见 $X'$ 具有闭子空间滤过
$$
X = X_1 \subset X_2 \subset \ldots \subset X_n \subset X_{n + 1} = X'
$$
，使每一对 $X_i \subset X_{i + 1}$ 都是 $S$ 上的一阶加厚。
利用简单的归纳论证，许多针对一阶加厚证明的结果，
都可改述为关于有限阶加厚的结果。

\medskip\noindent
一阶加厚可描述如下（见模，引理
\ref{modules-lemma-double-structure-gives-derivation}）。

\begin{lemma}
\label{more-morphisms-lemma-first-order-thickening}
设 $X$ 是基概形 $S$ 上的概形。考虑 $X$ 上层的短正合列
$$
0 \to \mathcal{I} \to \mathcal{A} \to \mathcal{O}_X \to 0
$$
，其中 $\mathcal{A}$ 是 $f^{-1}\mathcal{O}_S$-代数层，
$\mathcal{A} \to \mathcal{O}_X$ 是
$f^{-1}\mathcal{O}_S$-代数层的满射，而 $\mathcal{I}$ 是其核。
若
\begin{enumerate}
\item $\mathcal{I}$ 是 $\mathcal{A}$ 中平方为零的理想；并且
\item $\mathcal{I}$ 作为 $\mathcal{O}_X$-模是拟凝聚的，
\end{enumerate}
则 $X' = (X, \mathcal{A})$ 是概形，且 $X \to X'$ 是
$S$ 上的一阶加厚。此外，$S$ 上任意一阶加厚都具有这种形式。
\end{lemma}

\begin{proof}
显然 $X'$ 是局部赋环空间。设 $U = \Spec(B)$ 是 $X$
的仿射开集，并置 $A = \Gamma(U, \mathcal{A})$。
注意，由于 $H^1(U, \mathcal{I}) = 0$
（见概形上同调，引理
\ref{coherent-lemma-quasi-coherent-affine-cohomology-zero}），
映射 $A \to B$ 是满射。按假设，其核
$I = \mathcal{I}(U)$ 是环 $A$ 中平方为零的理想。
由概形，引理 \ref{schemes-lemma-morphism-into-affine}，
存在局部赋环空间的典范态射
$$
(U, \mathcal{A}|_U) \longrightarrow \Spec(A)
$$
，它来自映射 $B \to \Gamma(U, \mathcal{A})$。
% P05-ERR-CH37-0003: direction conflicts with the preceding A -> B surjection.
由于此态射嵌入交换图
$$
\xymatrix{
(U, \mathcal{O}_X|_U) \ar[d] \ar[r] & \Spec(B) \ar[d] \\
(U, \mathcal{A}|_U) \ar[r] & \Spec(A)
}
$$
，可见它在底层拓扑空间上是同胚。因此为证明它是同构，
只需检查它在局部环上诱导同构。
对与素理想 $\mathfrak p \subset A$ 对应的 $u \in U$，
得到短正合列的交换图
$$
\xymatrix{
0 \ar[r] &
I_{\mathfrak p} \ar[r] \ar[d] &
A_{\mathfrak p} \ar[r] \ar[d] &
B_{\mathfrak p} \ar[r] \ar[d] & 0 \\
0 \ar[r] &
\mathcal{I}_u \ar[r] &
\mathcal{A}_u \ar[r] &
\mathcal{O}_{X, u} \ar[r] & 0.
}
$$
由于 $\mathcal{I}$ 与 $\mathcal{O}_X$ 都是拟凝聚层，
左右两条竖直箭头均为同构，故中间的映射也是同构。
所以 $X' = (X, \mathcal{A})$ 的每一点都有仿射邻域，
而 $X'$ 如所需是概形。
\end{proof}

\begin{lemma}
\label{more-morphisms-lemma-thickening-affine-scheme}
\begin{slogan}
仿射性不受加厚影响
\end{slogan}
\begin{reference}
有限阶加厚的情形见 \cite[Proposition 5.1.9]{EGA1}。
\end{reference}
仿射概形的任意加厚都是仿射的。
\end{lemma}

\begin{proof}
这是极限，命题 \ref{limits-proposition-affine} 的特殊情形。
\end{proof}

\begin{proof}[有限阶加厚情形的证明]
设 $X \subset X'$ 是有限阶加厚，且 $X$ 仿射。
可以利用 Serre 判别法证明 $X'$ 仿射。更确切地，将使用概形上同调，引理
\ref{coherent-lemma-quasi-compact-h1-zero-covering}。
设 $\mathcal{F}$ 是拟凝聚 $\mathcal{O}_{X'}$-模。
只需证明 $H^1(X', \mathcal{F}) = 0$。
以 $i : X \to X'$ 表示给定的闭浸入，并记
$\mathcal{I} = \Ker(i^\sharp : \mathcal{O}_{X'} \to i_*\mathcal{O}_X)$.
由定义 \ref{more-morphisms-definition-thickening} 之后关于有限阶加厚的讨论，
存在 $n \geq 0$ 以及滤过
$$
0 = \mathcal{F}_{n + 1} \subset \mathcal{F}_n \subset
\mathcal{F}_{n - 1} \subset \ldots \subset
\mathcal{F}_0 = \mathcal{F}
$$
，其各项为拟凝聚子模，且
$\mathcal{F}_a/\mathcal{F}_{a + 1}$ 被 $\mathcal{I}$ 零化。
具体可取 $\mathcal{F}_a = \mathcal{I}^a\mathcal{F}$。
于是对某个拟凝聚 $\mathcal{O}_X$-模 $\mathcal{G}_a$，有
$\mathcal{F}_a/\mathcal{F}_{a + 1} = i_*\mathcal{G}_a$，
见态射，引理 \ref{morphisms-lemma-i-star-equivalence}。由此得到
$$
H^1(X', \mathcal{F}_a/\mathcal{F}_{a + 1}) =
H^1(X', i_*\mathcal{G}_a) = H^1(X, \mathcal{G}_a) = 0
$$
第二个等号来自概形上同调，引理
\ref{coherent-lemma-relative-affine-cohomology}，
最后一个等号来自概形上同调，引理
\ref{coherent-lemma-quasi-coherent-affine-cohomology-zero}。
因此 $\mathcal{F}$ 有一个有限滤过，其逐次商的第一同调均消失；
简单归纳即可推出 $H^1(X', \mathcal{F}) = 0$。
\end{proof}

\begin{lemma}
\label{more-morphisms-lemma-base-change-thickening}
设 $S \subset S'$ 是概形的加厚，给定态射 $X' \to S'$，
并置 $X = S \times_{S'} X'$。则
$(X \subset X') \to (S \subset S')$ 是加厚的态射。
若 $S \subset S'$ 是一阶（分别地，有限阶）加厚，
则 $X \subset X'$ 也是一阶（分别地，有限阶）加厚。
\end{lemma}

\begin{proof}
略。
\end{proof}

\begin{lemma}
\label{more-morphisms-lemma-composition-thickening}
\begin{slogan}
加厚的复合仍是加厚
\end{slogan}
若 $S \subset S'$ 与 $S' \subset S''$ 都是加厚，
则 $S \subset S''$ 也是加厚。
\end{lemma}

\begin{proof}
略。
\end{proof}

\begin{lemma}
\label{more-morphisms-lemma-descending-property-thickening}
成为加厚这一性质关于 fpqc 是局部的。
一阶加厚也同样如此。
\end{lemma}

\begin{proof}
该陈述的含义如下。设 $X \to X'$ 是概形态射，
$\{g_i : X'_i \to X'\}$ 是 fpqc 覆盖，
使得对所有 $i$，基变换 $X_i \to X'_i$ 都是加厚。
则 $X \to X'$ 是加厚。由于诸态射 $g_i$ 联合满射，
可知 $X \to X'$ 满射。由下降，引理
\ref{descent-lemma-descending-property-closed-immersion}，
$X \to X'$ 是闭浸入。因此 $X \to X'$ 是加厚。
略去一阶加厚情形的证明。
\end{proof}









\section{加厚的态射}
\label{more-morphisms-section-morphisms-thickenings}

\noindent
若 $(f, f') : (X \subset X') \to (Y \subset Y')$
是概形加厚的态射，则态射 $f$ 的性质往往由 $f'$ 继承。
这有若干种变体。

\begin{lemma}
\label{more-morphisms-lemma-thicken-property-morphisms}
设 $(f, f') : (X \subset X') \to (S \subset S')$
是加厚的态射。则
\begin{enumerate}
\item $f$ 是仿射态射，当且仅当 $f'$ 是仿射态射；
\item $f$ 是满射，当且仅当 $f'$ 是满射；
\item $f$ 拟紧，当且仅当 $f'$ 拟紧；
\item $f$ 泛闭，当且仅当 $f'$ 泛闭；
\item $f$ 整，当且仅当 $f'$ 整；
\item $f$（拟）分离，当且仅当 $f'$（拟）分离；
\item $f$ 泛单射，当且仅当 $f'$ 泛单射；
\item $f$ 泛开，当且仅当 $f'$ 泛开；
\item $f$ 拟仿射，当且仅当 $f'$ 拟仿射；以及
\item 此处添加更多内容。
\end{enumerate}
\end{lemma}

\begin{proof}
注意，$S \to S'$ 与 $X \to X'$ 都是泛同胚
（例如见态射，引理
\ref{morphisms-lemma-reduction-universal-homeomorphism}）。
这立即蕴含 (2)、(3)、(4)、(7)、(8)。
(1) 由引理 \ref{more-morphisms-lemma-thickening-affine-scheme} 得出；
该引理说明 $S$ 与 $S'$ 的仿射开集之间，以及
$X$ 与 $X'$ 的仿射开集之间，都存在一一对应。
(9) 由极限，引理 \ref{limits-lemma-thickening-quasi-affine}
以及刚才关于 $S$ 与 $S'$ 的仿射开集所作评注得出。
由态射，引理
\ref{morphisms-lemma-integral-universally-closed}，
(5) 由 (1) 与 (4) 得出。最后注意
$$
S \times_X S = S \times_{X'} S \to S \times_{X'} S' \to S' \times_{X'} S'
$$
是加厚（由引理 \ref{more-morphisms-lemma-base-change-thickening}，
两条箭头都是加厚）。
% P05-ERR-CH37-0004: frozen source fibre products reverse X/S roles.
所以把 (3) 与 (4) 应用于态射
$(S \subset S') \to (S \times_X S \to S' \times_{X'} S')$
，就得到 (6)。
\end{proof}

\begin{lemma}
\label{more-morphisms-lemma-thicken-property-relatively-ample}
设 $(f, f') : (X \subset X') \to (S \subset S')$
是加厚的态射。设 $\mathcal{L}'$ 是 $X'$ 上的可逆层，
并以 $\mathcal{L}$ 表示它到 $X$ 的限制。
则 $\mathcal{L}'$ 是 $f'$-丰沛的，当且仅当
$\mathcal{L}$ 是 $f$-丰沛的。
\end{lemma}

\begin{proof}
回忆，相对丰沛性是关于基底中每个仿射开集的条件，
见态射，定义 \ref{morphisms-definition-relatively-ample}。
由引理 \ref{more-morphisms-lemma-thickening-affine-scheme}，
$S$ 与 $S'$ 的仿射开集之间存在一一对应。
所以可假设 $S$ 与 $S'$ 都仿射，并归结为证明：
$\mathcal{L}'$ 丰沛，当且仅当 $\mathcal{L}$ 丰沛。
这就是极限，引理 \ref{limits-lemma-ample-on-reduction}。
\end{proof}

\begin{lemma}
\label{more-morphisms-lemma-thicken-property-morphisms-cartesian}
设 $(f, f') : (X \subset X') \to (S \subset S')$
是加厚的态射，并且 $X = S \times_{S'} X'$。
若 $S \subset S'$ 是有限阶加厚，则
\begin{enumerate}
\item $f$ 是闭浸入，当且仅当 $f'$ 是闭浸入；
\item $f$ 局部有限型，当且仅当 $f'$ 局部有限型；
\item $f$ 局部拟有限，当且仅当 $f'$ 局部拟有限；
\item $f$ 是相对维数为 $d$ 的局部有限型态射，当且仅当
$f'$ 是相对维数为 $d$ 的局部有限型态射；
\item $\Omega_{X/S} = 0$，当且仅当 $\Omega_{X'/S'} = 0$；
\item $f$ 非分歧，当且仅当 $f'$ 非分歧；
\item $f$ 固有，当且仅当 $f'$ 固有；
\item $f$ 有限，当且仅当 $f'$ 有限；
\item $f$ 是单态射，当且仅当 $f'$ 是单态射；
\item $f$ 是浸入，当且仅当 $f'$ 是浸入；以及
\item 此处添加更多内容。
\end{enumerate}
\end{lemma}

\begin{proof}
引理所列性质 $\mathcal{P}$ 在基变换下都稳定，
所以若 $f'$ 具有性质 $\mathcal{P}$，则 $f$ 也具有。
见概形，诸引理 \ref{schemes-lemma-base-change-immersion} 与
\ref{schemes-lemma-base-change-monomorphism}，以及态射，诸引理
\ref{morphisms-lemma-base-change-finite-type},
\ref{morphisms-lemma-base-change-quasi-finite},
\ref{morphisms-lemma-base-change-relative-dimension-d},
\ref{morphisms-lemma-base-change-differentials},
\ref{morphisms-lemma-base-change-unramified},
\ref{morphisms-lemma-base-change-proper}, 及
\ref{morphisms-lemma-base-change-finite}.

\medskip\noindent
因此每种情形中有内容的方向，都是从 $f$ 具有该性质出发，
推出 $f'$ 也具有。对加厚的阶数作归纳，可假设
$S \subset S'$ 是一阶加厚，见定义
\ref{more-morphisms-definition-thickening} 之后紧接着的讨论。

\medskip\noindent
大多数证明都要归结到仿射情形。设 $U' \subset S'$ 是仿射开集，
$V' \subset X'$ 是位于 $U'$ 上方的仿射开集。
令 $U' = \Spec(A')$，并以 $I \subset A'$ 表示定义闭子概形
$U' \cap S$ 的理想。设 $V' = \Spec(B')$。
于是 $V' \cap X = \Spec(B'/IB')$。置
$A = A'/I$、$B = B'/IB'$，得到交换图
$$
\xymatrix{
0 \ar[r] &
IB' \ar[r] &
B' \ar[r] &
B \ar[r] & 0 \\
0 \ar[r] &
IA' \ar[r] \ar[u] &
A' \ar[r] \ar[u] &
A \ar[r] \ar[u] & 0
}
$$
，其两行正合，且 $I^2 = 0$。

\medskip\noindent
把 (1) 译成代数：若 $A \to B$ 满射，则 $A' \to B'$ 满射。
这由 Nakayama 引理（代数，引理 \ref{algebra-lemma-NAK}）得出。

\medskip\noindent
把 (2) 译成代数：若 $A \to B$ 是有限型环映射，则
$A' \to B'$ 也是有限型环映射。这可将 Nakayama 引理
（代数，引理 \ref{algebra-lemma-NAK}）
应用于映射 $A'[x_1, \ldots, x_n] \to B'$ 得出，
其中 $A[x_1, \ldots, x_n] \to B$ 满射。

\medskip\noindent
(3) 的证明。它由 (2) 以及如下事实得出：局部有限型态射的拟有限性
可在纤维上检验，见态射，引理
\ref{morphisms-lemma-quasi-finite-at-point-characterize}。

\medskip\noindent
(4) 的证明。它由 (2) 以及如下事实得出：附加性质“相对维数为
$d$”可在纤维上检验（按定义如此，见态射，定义
\ref{morphisms-definition-relative-dimension-d}。
% P05-ERR-CH37-0005: frozen source is missing the closing parenthesis.

\medskip\noindent
把 (5) 译成代数：若 $\Omega_{B/A} = 0$，则
$\Omega_{B'/A'} = 0$。由代数，引理
\ref{algebra-lemma-differentials-base-change}，有
$0 = \Omega_{B/A} = \Omega_{B'/A'}/I\Omega_{B'/A'}$。
因此由 Nakayama 引理（代数，引理 \ref{algebra-lemma-NAK}），
$\Omega_{B'/A'} = 0$。

\medskip\noindent
把 (6) 译成代数：若 $A \to B$ 是非分歧映射，则
$A' \to B'$ 非分歧。由于 $A \to B$ 有限型，
由上面的 (2)，$A' \to B'$ 有限型。
又因 $A \to B$ 非分歧，有 $\Omega_{B/A} = 0$；
由 (5)，有 $\Omega_{B'/A'} = 0$。所以 $A' \to B'$ 非分歧。

\medskip\noindent
(7) 的证明。把 (2)、引理
\ref{more-morphisms-lemma-thicken-property-morphisms} 的结果，
以及“固有等价于拟紧 $+$ 分离 $+$ 局部有限型 $+$ 泛闭”这一事实结合即可。

\medskip\noindent
(8) 的证明。把 (2) 与引理
\ref{more-morphisms-lemma-thicken-property-morphisms} 的结果结合，
并使用“有限等价于整 $+$ 局部有限型”这一事实即可
（态射，引理 \ref{morphisms-lemma-finite-integral}）。

\medskip\noindent
(9) 的证明。由于 $f$ 是单态射，有 $X = X \times_S X$。
可把迄今证明的结果应用于加厚的态射
$(X \subset X') \to (X \times_S X \subset X' \times_{S'} X')$.
。由 (1)，得到 $X' \to X' \times_{S'} X'$ 是闭浸入。
事实上它是一阶加厚，因为定义闭浸入
$X' \to X' \times_{S'} X'$ 的理想，包含在理想
$\mathcal{I} \subset \mathcal{O}_{S'}$ 的拉回中；
该理想定义 $S'$ 中的 $S$。确实，
$X = X \times_S X = (X' \times_{S'} X') \times_{S'} S$
包含于 $X'$。所以由态射，引理
\ref{morphisms-lemma-differentials-diagonal}，
只需证明 $\Omega_{X'/S'} = 0$；
这由 (5) 以及 $X/S$ 的相应陈述得出。

\medskip\noindent
(10) 的证明。若 $f : X \to S$ 是浸入，则它分解为
$X \to U \to S$，其中 $U \to S$ 是开浸入，
$X \to U$ 是闭浸入。令 $U' \subset S'$ 为底层拓扑空间
与 $U$ 相同的开子概形。于是 $X' \to S'$ 经 $U'$ 分解，
由 (1)，$X' \to U'$ 是闭浸入。证明完毕。
\end{proof}

\noindent
下一个引理是前一引理的变体。这里不再假设所涉加厚为有限阶
（该假设使我们能把局部有限型这一性质从 $f$ 传递给 $f'$），
而是直接假定 $f$ 与 $f'$ 都局部有限型。

\begin{lemma}
\label{more-morphisms-lemma-properties-that-extend-over-thickenings}
设 $(f, f') : (X \subset X') \to (Y \to Y')$ 是加厚的态射。
% P05-ERR-CH37-0006: frozen source uses Y -> Y' instead of Y subset Y'.
假设 $f$ 与 $f'$ 局部有限型，且 $X = Y \times_{Y'} X'$。则
\begin{enumerate}
\item $f$ 局部拟有限，当且仅当 $f'$ 局部拟有限；
\item $f$ 有限，当且仅当 $f'$ 有限；
\item $f$ 是闭浸入，当且仅当 $f'$ 是闭浸入；
\item $\Omega_{X/Y} = 0$，当且仅当 $\Omega_{X'/Y'} = 0$；
\item $f$ 非分歧，当且仅当 $f'$ 非分歧；
\item $f$ 是单态射，当且仅当 $f'$ 是单态射；
\item $f$ 是浸入，当且仅当 $f'$ 是浸入；
\item $f$ 固有，当且仅当 $f'$ 固有；以及
\item 此处添加更多内容。
\end{enumerate}
\end{lemma}

\begin{proof}
引理所列性质 $\mathcal{P}$ 在基变换下都稳定，
所以若 $f'$ 具有性质 $\mathcal{P}$，则 $f$ 也具有。
见概形，诸引理 \ref{schemes-lemma-base-change-immersion} 与
\ref{schemes-lemma-base-change-monomorphism}，以及态射，诸引理
\ref{morphisms-lemma-base-change-quasi-finite},
\ref{morphisms-lemma-base-change-relative-dimension-d},
\ref{morphisms-lemma-base-change-differentials},
\ref{morphisms-lemma-base-change-unramified},
\ref{morphisms-lemma-base-change-proper}, 及
\ref{morphisms-lemma-base-change-finite}.
因此每种情形中只需证明：若 $f$ 具有所需性质，
则 $f'$ 也具有。

\medskip\noindent
一个态射局部拟有限，当且仅当它局部有限型且概形论纤维为离散空间，
见态射，引理 \ref{morphisms-lemma-locally-quasi-finite-fibres}。
由于传到加厚不会改变底层拓扑空间，可知若（且仅若）
$f$ 局部拟有限，则 $f'$ 局部拟有限。这证明了 (1)。

\medskip\noindent
(2) 由引理 \ref{more-morphisms-lemma-thicken-property-morphisms} 的 (5)
以及下列事实得出：有限态射恰好是局部有限型的整态射
（态射，引理 \ref{morphisms-lemma-finite-integral}）。

\medskip\noindent
(3)。这立即由态射，引理
\ref{morphisms-lemma-check-closed-infinitesimally} 得出。

\medskip\noindent
(4) 由下列代数陈述得出：设 $A$ 是环，
$I \subset A$ 是局部幂零理想，$B$ 是有限型 $A$-代数。
若 $\Omega_{(B/IB)/(A/I)} = 0$，则 $\Omega_{B/A} = 0$。
具体地，该假设意味着 $I\Omega_{B/A} = 0$，见代数，引理
\ref{algebra-lemma-differentials-base-change}。
另一方面，$\Omega_{B/A}$ 是有限 $B$-模，见代数，引理
\ref{algebra-lemma-differentials-finitely-generated}。
所以 $\Omega_{B/A}$ 的消失性由 Nakayama 引理
（代数，引理 \ref{algebra-lemma-NAK}）以及
$IB$ 包含于 $B$ 的 Jacobson 根这一事实得出。

\medskip\noindent
(5) 立即由 (4) 与态射，引理
\ref{morphisms-lemma-unramified-omega-zero} 得出。

\medskip\noindent
(6) 的证明。由于 $f$ 是单态射，有 $X = X \times_Y X$。
可把迄今证明的结果应用于加厚的态射
$(X \subset X') \to (X \times_Y X \subset X' \times_{Y'} X')$.
。由 (3)，得到
$\Delta_{X'/Y'} : X' \to X' \times_{Y'} X'$ 是闭浸入。
事实上，$\Delta_{X'/Y'}$ 在底层集合上是双射，
故 $\Delta_{X'/Y'}$ 是加厚。另一方面，由态射，引理
\ref{morphisms-lemma-diagonal-morphism-finite-type}，
$\Delta_{X'/Y'}$ 局部有限表示。换言之，
$\Delta_{X'/Y'}(X')$ 由有限型拟凝聚理想层
$\mathcal{J} \subset \mathcal{O}_{X' \times_{Y'} X'}$ 截出。
由 (5)，$\Omega_{X'/Y'} = 0$；再由态射，引理
\ref{morphisms-lemma-differentials-diagonal}，
$X' \to X' \times_{Y'} X'$ 的余法层为零。
换言之，$\mathcal{J}/\mathcal{J}^2 = 0$。
例如由代数，引理
\ref{algebra-lemma-ideal-is-squared-union-connected}，
这蕴含 $\Delta_{X'/Y'}$ 是同构。

\medskip\noindent
(7) 的证明。若 $f : X \to Y$ 是浸入，则它分解为
$X \to V \to Y$，其中 $V \to Y$ 是开浸入，
$X \to V$ 是闭浸入。令 $V' \subset Y'$ 为底层拓扑空间
与 $V$ 相同的开子概形。于是 $X' \to V'$ 经 $V'$ 分解，
% P05-ERR-CH37-0007: frozen source's factorization sentence is tautological.
并由 (3) 得到 $X' \to V'$ 是闭浸入。

\medskip\noindent
(8) 由引理 \ref{more-morphisms-lemma-thicken-property-morphisms}
以及固有态射的定义得出：它是拟紧、泛闭、分离且局部有限型的态射。
\end{proof}








\section{加厚的 Picard 群}
\label{more-morphisms-section-picard-group-thickening}

\noindent
这里给出一些关于加厚的 Picard 群的材料。

\begin{lemma}
\label{more-morphisms-lemma-picard-group-first-order-thickening}
设 $X \subset X'$ 是理想层为 $\mathcal{I}$ 的一阶加厚。则有阿贝尔群的典范正合列
$$
\xymatrix{
0 \ar[r] &
H^0(X, \mathcal{I}) \ar[r] &
H^0(X', \mathcal{O}_{X'}^*) \ar[r] &
H^0(X, \mathcal{O}^*_X) \ar `r[d] `d[l] `l[llld] `d[dll] [dll] \\
& H^1(X, \mathcal{I}) \ar[r] &
\Pic(X') \ar[r] &
\Pic(X) \ar `r[d] `d[l] `l[llld] `d[dll] [dll] \\
& H^2(X, \mathcal{I}) \ar[r] & \ldots \ar[r] & \ldots
}
$$
。
\end{lemma}

\begin{proof}
这是与阿贝尔群层的短正合列
$$
0 \to \mathcal{I} \to \mathcal{O}_{X'}^* \to \mathcal{O}_X^* \to 0
$$
相伴的长上同调正合列，其中第一个映射把 $\mathcal{I}$ 的局部截面 $f$
送到 $\mathcal{O}_{X'}$ 的可逆截面 $1 + f$。
这里还使用了环化空间的 Picard 群与其可逆函数层的第一上同调群之间的等同；
见 上同调章引理 \ref{cohomology-lemma-h1-invertible}。
\end{proof}

\begin{lemma}
\label{more-morphisms-lemma-torsion-pic-thickening}
设 $X \subset X'$ 是一个加厚，且整数 $n$ 在 $\mathcal{O}_X$ 中可逆。
则映射 $\Pic(X')[n] \to \Pic(X)[n]$ 是双射。
\end{lemma}

\begin{proof}[有限阶加厚情形的证明]
由定义 \ref{more-morphisms-definition-thickening} 后所述的一般原理，问题可归结为一阶加厚的情形。
此时由引理 \ref{more-morphisms-lemma-picard-group-first-order-thickening}，只需证明
$H^1(X, \mathcal{I})[n]$、$H^1(X, \mathcal{I})/n$ 和
$H^2(X, \mathcal{I})[n]$ 均为零。
这是因为 $\mathcal{I}$ 作为 $\mathcal{O}_X$-模，其上的乘 $n$ 映射是同构。
\end{proof}

\begin{proof}[一般情形的证明]
设 $\mathcal{I} \subset \mathcal{O}_{X'}$ 是截出 $X$ 的拟凝聚理想层。
则有阿贝尔群层的短正合列
$$
0 \to (1 + \mathcal{I})^* \to \mathcal{O}_{X'}^* \to \mathcal{O}_X^* \to 0
$$
从而得到如引理 \ref{more-morphisms-lemma-picard-group-first-order-thickening} 所述的长上同调正合列，
只是把 $H^i(X, \mathcal{I})$ 换成 $H^i(X, (1 + \mathcal{I})^*)$。
因此只需证明取 $n$ 次幂给出同构
$(1 + \mathcal{I})^* \to (1 + \mathcal{I})^*$。
在仿射开集上取截面后，这由
代数章引理 \ref{algebra-lemma-lift-nth-roots} 得出。
\end{proof}





\section{无穷小邻域}
\label{more-morphisms-section-first-order-infinitesimal-neighbourhood}

\noindent
有限阶加厚有如下自然构造。设 $i : Z \to X$ 是概形的浸入。
选择开子概形 $U \subset X$，使得 $i$ 将 $Z$ 等同于闭子概形
$Z \subset U$。设 $\mathcal{I} \subset \mathcal{O}_U$
是在 $U$ 中定义 $Z$ 的拟凝聚理想层。对 $n \geq 1$，
可以考虑由拟凝聚理想层 $\mathcal{I}^{n + 1}$ 定义的闭子概形
$Z_n \subset U$。

\begin{definition}
\label{more-morphisms-definition-first-order-infinitesimal-neighbourhood}
设 $i : Z \to X$ 是概形的浸入。
\begin{enumerate}
\item $Z$ 在 $X$ 中的{\it 一阶无穷小邻域}是上述 $X$ 上的一阶加厚
$Z \subset Z_1$。
\item $Z$ 在 $X$ 中的{\it $n$ 阶无穷小邻域}是上述 $X$ 上的
$n$ 阶加厚 $Z \subset Z_n$。
\end{enumerate}
\end{definition}

\noindent
这些加厚具有如下泛性质（它也消除了对上述构造依赖于开集 $U$ 之选取的疑虑）。

\begin{lemma}
\label{more-morphisms-lemma-first-order-infinitesimal-neighbourhood}
设 $i : Z \to X$ 是概形的浸入。
\begin{enumerate}
\item $Z$ 在 $X$ 中的一阶无穷小邻域 $Z'$ 具有如下泛性质：给定任意交换图
$$
\xymatrix{
Z \ar[d]_i & T \ar[l]^a \ar[d] \\
X & T' \ar[l]_b
}
$$
其中 $T \subset T'$ 是 $X$ 上的一阶加厚，则存在唯一的 $X$ 上加厚的态射
$(a', a) : (T \subset T') \to (Z \subset Z')$。
\item 对 $n \geq 1$，$Z$ 在 $X$ 中的 $n$ 阶无穷小邻域 $Z_n$
具有如下泛性质：给定任意交换图
$$
\xymatrix{
Z \ar[d]_i & T \ar[l]^a \ar[d] \\
X & T' \ar[l]_b
}
$$
其中 $T \subset T'$ 是 $X$ 上的 $n$ 阶加厚，则存在唯一的 $X$ 上加厚的态射
$(a', a) : (T \subset T') \to (Z \subset Z_n)$。
\end{enumerate}
\end{lemma}

\begin{proof}
只证明 (1)。设 $U \subset X$ 是构造 $Z'$ 时所用的开集；也就是说，
$Z$ 被等同于 $U$ 中由拟凝聚理想层 $\mathcal{I}$ 截出的闭子概形。
由于 $|T| = |T'|$，可知 $b(T') \subset U$，因而可把 $b$ 看作到 $U$ 的态射。
设 $\mathcal{J} \subset \mathcal{O}_{T'}$ 是截出 $T$ 的理想。
由上图中的 $b(T) \subset Z$，可知
$b^\sharp(b^{-1}\mathcal{I}) \subset \mathcal{J}$。
由于 $T'$ 是 $T$ 的一阶加厚，有 $\mathcal{J}^2 = 0$，故
$b^\sharp(b^{-1}(\mathcal{I}^2)) = 0$。由
概形章引理 \ref{schemes-lemma-characterize-closed-subspace}，
这表明 $b$ 经由 $Z'$ 分解。以 $a' : T' \to Z'$ 表示此分解，结论即明。
\end{proof}

\begin{lemma}
\label{more-morphisms-lemma-infinitesimal-neighbourhood-conormal}
设 $i : Z \to X$ 是概形的浸入，并设 $Z \subset Z'$ 是 $Z$ 在 $X$
中的一阶无穷小邻域。则图
$$
\xymatrix{
Z \ar[r] \ar[d] & Z' \ar[d] \\
Z \ar[r] & X
}
$$
由 态射章引理 \ref{morphisms-lemma-conormal-functorial}
诱导余法层的映射 $\mathcal{C}_{Z/X} \to \mathcal{C}_{Z/Z'}$。
此映射是同构。
\end{lemma}

\begin{proof}
这由上述 $Z'$ 的构造立即可见。
\end{proof}
















\section{形式非分歧态射}
\label{more-morphisms-section-formally-unramified}

\noindent
回顾：若对每个实线部分交换的图
$$
\xymatrix{
A \ar[r] \ar@{-->}[rd] & B/I \\
R \ar[r] \ar[u] & B \ar[u]
}
$$
（其中 $I \subset B$ 是平方为零的理想），至多存在一条使该图交换的虚线箭头，
则称环映射 $R \to A$ 是{\it 形式非分歧的}；见
代数章定义 \ref{algebra-definition-formally-unramified}。
这引出概形态射的如下类似概念。

\begin{definition}
\label{more-morphisms-definition-formally-unramified}
设 $f : X \to S$ 是概形态射。若给定任意实线部分交换的图
$$
\xymatrix{
X \ar[d]_f & T \ar[d]^i \ar[l] \\
S & T' \ar[l] \ar@{-->}[lu]
}
$$
其中 $T \subset T'$ 是 $S$ 上仿射概形的一阶加厚时，至多存在一条使该图交换的
虚线箭头，则称 $f$ 是{\it 形式非分歧的}。
\end{definition}

\noindent
先证明一些形式性引理，也就是只需画出相应交换图即可证明的引理。

\begin{lemma}
\label{more-morphisms-lemma-formally-unramified-not-affine}
若 $f : X \to S$ 是形式非分歧态射，则给定任意实线部分交换的图
$$
\xymatrix{
X \ar[d]_f & T \ar[d]^i \ar[l] \\
S & T' \ar[l] \ar@{-->}[lu]
}
$$
其中 $T \subset T'$ 是 $S$ 上概形的一阶加厚，至多存在一条使该图交换的虚线箭头。
换言之，在定义 \ref{more-morphisms-definition-formally-unramified} 中可以去掉 $T$ 为仿射的条件。
\end{lemma}

\begin{proof}
这是因为一个态射由它在各仿射开集上的限制所决定。
\end{proof}

\begin{lemma}
\label{more-morphisms-lemma-composition-formally-unramified}
形式非分歧态射的复合仍是形式非分歧的。
\end{lemma}

\begin{proof}
这是形式性的。
\end{proof}

\begin{lemma}
\label{more-morphisms-lemma-base-change-formally-unramified}
形式非分歧态射的基变换仍是形式非分歧的。
\end{lemma}

\begin{proof}
这是形式性的。
\end{proof}

\begin{lemma}
\label{more-morphisms-lemma-formally-unramified-on-opens}
设 $f : X \to S$ 是概形态射，并设 $U \subset X$ 与 $V \subset S$
是满足 $f(U) \subset V$ 的开集。若 $f$ 形式非分歧，则
$f|_U : U \to V$ 也形式非分歧。
\end{lemma}

\begin{proof}
考虑定义 \ref{more-morphisms-definition-formally-unramified} 中那样的实线部分交换图
$$
\xymatrix{
U \ar[d]_{f|_U} & T \ar[d]^i \ar[l]^a \\
V & T' \ar[l] \ar@{-->}[lu]
}
$$
。
% P05-ERR-CH37-0008: Frozen source says ``formally ramified'' here; retained literally pending canon review.
若 $f$ 形式分歧，则至多存在一个满足 $a'|_T = a$ 的 $S$-态射
$a' : T' \to X$。因此显然至多存在一个这样的到 $U$ 的态射。
\end{proof}

\begin{lemma}
\label{more-morphisms-lemma-affine-formally-unramified}
设 $f : X \to S$ 是概形态射，并假设 $X$ 与 $S$ 均为仿射。
则 $f$ 形式非分歧，当且仅当
$\mathcal{O}_S(S) \to \mathcal{O}_X(X)$ 是形式非分歧环映射。
\end{lemma}

\begin{proof}
由环范畴与仿射概形范畴之间的等价，这立即从定义
（定义 \ref{more-morphisms-definition-formally-unramified} 与
代数章定义 \ref{algebra-definition-formally-unramified}）得出；见
概形章引理 \ref{schemes-lemma-category-affine-schemes}。
\end{proof}

\noindent
下面用微分层给出一个刻画。

\begin{lemma}
\label{more-morphisms-lemma-formally-unramified-differentials}
设 $f : X \to S$ 是概形态射。则 $f$ 形式非分歧，当且仅当
$\Omega_{X/S} = 0$。
\end{lemma}

\begin{proof}
回顾 态射章引理 \ref{morphisms-lemma-differentials-affine}
证明中的部分论证。取开集 $W \subset X \times_S X$，使对角态射
$\Delta : X \to X \times_S X$ 诱导到 $W$ 的闭浸入。
设 $\mathcal{J} \subset \mathcal{O}_W$ 是此闭浸入的理想层，并设
$X' \subset W$ 是由拟凝聚理想层 $\mathcal{J}^2$ 定义的闭子概形。
考虑由两个投影 $X \times_S X \to X$ 诱导的两个态射
$p_1, p_2 : X' \to X$。注意，$p_1$ 与 $p_2$ 同
$\Delta : X \to X'$ 复合后相等，而且 $X \to X'$ 是由平方为零的理想
定义的闭浸入。此外有短正合列
$$
0 \to \mathcal{J}/\mathcal{J}^2 \to \mathcal{O}_{X'} \to \mathcal{O}_X \to 0
$$
且 $\Omega_{X/S} = \mathcal{J}/\mathcal{J}^2$。再者，
$\mathcal{J}/\mathcal{J}^2$ 由局部截面
$p_1^\sharp(f) - p_2^\sharp(f)$ 生成，其中 $f$ 遍历
$\mathcal{O}_X$ 的局部截面。

\medskip\noindent
假设 $f : X \to S$ 形式非分歧。按假设，这意味着 $p_1 = p_2$
限制到任意仿射开集 $T' \subset X'$ 后相等，故 $p_1 = p_2$。
由上可得 $\Omega_{X/S} = \mathcal{J}/\mathcal{J}^2 = 0$。

\medskip\noindent
反之，假设 $\Omega_{X/S} = 0$，于是 $X' = X$。任取定义
\ref{more-morphisms-definition-formally-unramified} 中图的两条虚线箭头
$f'_1, f'_2 : T' \to X$。它们给出态射
$(f'_1, f'_2) : T' \to X \times_S X$。
由于 $f'_1|_T = f'_2|_T$ 且 $|T| =|T'|$，$T'$ 在
$(f'_1, f'_2)$ 下的像包含于上面选取的开集 $W$ 中。
又因 $(f'_1, f'_2)(T) \subset \Delta(X)$，且 $T$ 在 $T'$ 中由平方为零的
理想定义，可知 $(f'_1, f'_2)$ 经由 $X'$ 分解。因 $X' = X$，
遂得所需的 $f_1' = f'_2$。
\end{proof}

\begin{lemma}
\label{more-morphisms-lemma-formally-unramified}
设 $f : X \to Y$ 是概形态射。以下条件等价：
\begin{enumerate}
\item $f$ 形式非分歧；
\item 对每个 $x \in X$，存在满足 $f(U) \subset V$ 的开集
$x \in U \subset X$ 与 $f(x) \in V \subset Y$，使得
$f|_U : U \to V$ 形式非分歧；
\item 对每一对满足 $f(U) \subset V$ 的仿射开集 $U \subset X$ 与
$V \subset Y$，环映射 $\mathcal{O}_Y(V) \to \mathcal{O}_X(U)$
形式非分歧；
\item 存在仿射开覆盖 $Y = \bigcup V_j$，且对每个 $j$ 存在仿射开覆盖
$f^{-1}(V_j) = \bigcup U_{ji}$，使得对所有 $j$ 与 $i$，环映射
$\mathcal{O}_Y(V) \to \mathcal{O}_X(U)$ 形式非分歧。
\end{enumerate}
\end{lemma}

\begin{proof}
结合引理 \ref{more-morphisms-lemma-formally-unramified-on-opens} 与
\ref{more-morphisms-lemma-affine-formally-unramified} 可得 (1) $\Rightarrow$ (3)。
(3) $\Rightarrow$ (4) 是立即的。由引理
\ref{more-morphisms-lemma-affine-formally-unramified} 得 (4) $\Rightarrow$ (2)。
若 (2) 成立，则由引理 \ref{more-morphisms-lemma-formally-unramified-differentials}
（这里还使用 态射章引理
\ref{morphisms-lemma-differentials-restrict-open}），
$\Omega_{X/S}$ 在每一点的某个邻域上消失；于是再由引理
\ref{more-morphisms-lemma-formally-unramified-differentials}，$f$ 形式非分歧。
\end{proof}

\begin{lemma}
\label{more-morphisms-lemma-unramified-formally-unramified}
\begin{slogan}
非分歧态射恰为局部有限型的形式非分歧态射。
\end{slogan}
设 $f : X \to S$ 是概形态射。以下条件等价：
\begin{enumerate}
\item 态射 $f$ 非分歧（分别地，G-非分歧）；
\item 态射 $f$ 局部有限型（分别地，局部有限表示）且形式非分歧。
\end{enumerate}
\end{lemma}

\begin{proof}
使用引理 \ref{more-morphisms-lemma-formally-unramified-differentials} 与
态射章引理 \ref{morphisms-lemma-unramified-omega-zero}。
\end{proof}









\section{泛一阶加厚}
\label{more-morphisms-section-universal-thickening}

\noindent
设 $h : Z \to X$ 是概形态射。$Z$ 在 $X$ 上的一个{\it 泛一阶加厚}
是 $X$ 上的一阶加厚 $Z \subset Z'$，使得给定 $X$ 上的任意一阶加厚
$T \subset T'$ 和实线部分交换图
$$
\xymatrix{
& Z \ar[ld] & & T \ar[rd] \ar[ll]^a \\
Z' \ar[rrd] & & & & T' \ar@{..>}[llll]_{a'} \ar[lld]^b \\
 & & X
}
$$
时，存在唯一的虚线箭头使该图交换。注意，在这种情形下
$(a, a') : (T \subset T') \to (Z \subset Z')$ 是 $X$ 上加厚的态射。
因此若泛一阶加厚存在，则它在唯一同构的意义下唯一。
一般而言泛一阶加厚未必存在，但当 $h$ 形式非分歧时它确实存在。

\begin{lemma}
\label{more-morphisms-lemma-universal-thickening}
设 $h : Z \to X$ 是形式非分歧的概形态射。则存在 $Z$ 在 $X$ 上的
泛一阶加厚 $Z \subset Z'$。
\end{lemma}

\begin{proof}
在本证明中，若 $Z \subset Z'$ 满足引理中的条件，就称它为 $Z$ 在 $X$
上的泛一阶加厚。我们将通过粘合构造泛一阶加厚
$Z \subset Z'$，先从仿射情形开始；该情形即
代数章引理 \ref{algebra-lemma-universal-thickening}。
首先作一些一般性说明。

\medskip\noindent
若 $Z$ 在 $X$ 上的泛一阶加厚存在，则它在唯一同构的意义下唯一。
再设 $V \subset Z$ 与 $U \subset X$ 是满足 $h(V) \subset U$ 的开子概形，
并设 $Z \subset Z'$ 是 $Z$ 在 $X$ 上的泛一阶加厚。令
$V' \subset Z'$ 为满足 $V = Z \cap V'$ 的开子概形。我们断言
$V \subset V'$ 是 $V$ 在 $U$ 上的泛一阶加厚。事实上，设给定任意图
$$
\xymatrix{
V \ar[d]_h & T \ar[l]^a \ar[d] \\
U & T' \ar[l]_b
}
$$
其中 $T \subset T'$ 是 $U$ 上的一阶加厚。由 $Z'$ 的泛性质，得到
$(a, a') : (T \subset T') \to (Z \subset Z')$。
但由于底拓扑空间相等 $|T| = |T'|$，可知 $a'(T') \subset V'$。
因此可把 $(a, a')$ 看作 $U$ 上加厚的态射
$(a, a') : (T \subset T') \to (V \subset V')$，唯一性也很清楚。
完全类似地可证：若 $h(Z) \subset U$，且 $Z \subset Z'$ 是 $U$ 上的
泛一阶加厚，则 $Z \subset Z'$ 也是 $X$ 上的泛一阶加厚。

\medskip\noindent
在粘合各仿射片之前，先证明 $Z$ 与 $X$ 均为仿射时引理成立。
写 $X = \Spec(R)$、$Z = \Spec(S)$。由
代数章引理 \ref{algebra-lemma-universal-thickening}，
存在 $X$ 上的一阶加厚 $Z \subset Z'$，它对如下图具有引理中的泛性质：
$$
\xymatrix{
Z \ar[d]_h & T \ar[l]^a \ar[d] \\
X & T' \ar[l]_b
}
$$
其中 $T, T'$ 为仿射。对一般的图，可选择仿射开覆盖
$T' = \bigcup T'_i$，并得到 $X$ 上的态射 $a'_i : T'_i \to Z'$，
满足 $a'_i|_{T_i} = a|_{T_i}$。由唯一性，$a'_i$ 与 $a'_j$
在 $T'_i \cap T'_j$ 的任意仿射开集上相等。因此诸态射 $a'_i$
粘合成所需的 $X$ 上整体态射 $a' : T' \to Z'$。
故当 $X$ 与 $Z$ 均为仿射时引理成立。

\medskip\noindent
选择仿射开覆盖 $Z = \bigcup Z_i$，使每个 $Z_i$ 都映入 $X$ 的某个仿射开集
$U_i$。由引理 \ref{more-morphisms-lemma-formally-unramified-on-opens}，态射
$Z_i \to U_i$ 形式非分歧。因此由仿射情形，得到 $U_i$ 上的泛一阶加厚
$Z_i \subset Z_i'$。由上述一般性说明，$Z_i \subset Z_i'$ 也是 $Z_i$
在 $X$ 上的泛一阶加厚。令 $Z'_{i, j} \subset Z'_i$ 为满足
$Z_i \cap Z_j = Z'_{i, j} \cap Z_i$ 的开子概形。由一般性说明，
$Z'_{i, j}$ 与 $Z'_{j, i}$ 都是 $Z_i \cap Z_j$ 在 $X$ 上的泛一阶加厚。
于是由第一条一般性说明，存在典范同构
$\varphi_{ij} : Z'_{i, j} \to Z'_{j, i}$，它在 $Z_i \cap Z_j$ 上诱导恒等映射。
我们断言这些态射满足
概形章第 \ref{schemes-section-glueing-schemes} 节 的上闭链条件。
（验证从略。提示：使用 $Z'_{i, j} \cap Z'_{i, k}$ 是
$Z_i \cap Z_j \cap Z_k$ 在 $X$ 上的泛一阶加厚，并由上文知它在唯一同构的
意义下被确定。）因此可以使用
概形章第 \ref{schemes-section-glueing-schemes} 节 的结果，
得到 $X$ 上的一阶加厚 $Z \subset Z'$，使满足
$Z_i = Z'_i \cap Z$ 的开子概形 $Z'_i \subset Z'$ 是 $Z_i$
在 $X$ 上的泛一阶加厚。

\medskip\noindent
由此可形式地推出 $Z'$ 是 $Z$ 在 $X$ 上的泛一阶加厚。事实上，
对于任意如下图，我们已经有泛性质：
$$
\xymatrix{
Z \ar[d]_h & T \ar[l]^a \ar[d] \\
X & T' \ar[l]_b
}
$$
其中 $a(T)$ 包含于某个 $Z_i$。给定一般的图，可选择开覆盖
$T' = \bigcup T'_i$，使 $a(T_i) \subset Z_i$。于是得到 $X$ 上的态射
$a'_i : T'_i \to Z'$，满足 $a'_i|_{T_i} = a|_{T_i}$。
$a'_i$ 与 $a'_j$ 在 $T'_i \cap T'_j$ 上必然相等，因为
$a'_i|_{T'_i \cap T'_j}$ 与 $a'_j|_{T'_i \cap T'_j}$ 都解出了映入
$Z_i \cap Z_j$ 在 $X$ 上的泛一阶加厚 $Z'_i \cap Z'_j$ 的问题。
因此诸态射 $a'_i$ 粘合成所需的 $X$ 上整体态射
$a' : T' \to Z'$。证明完毕。
\end{proof}

\begin{definition}
\label{more-morphisms-definition-universal-thickening}
设 $h : Z \to X$ 是形式非分歧的概形态射。
\begin{enumerate}
\item $Z$ 在 $X$ 上的{\it 泛一阶加厚}是引理
\ref{more-morphisms-lemma-universal-thickening} 中构造的加厚 $Z \subset Z'$。
\item $Z$ 在 $X$ 上的{\it 余法层}是 $Z$ 在其 $X$ 上泛一阶加厚
$Z'$ 中的余法层。
\end{enumerate}
在这种情形下，我们常将该余法层记为 $\mathcal{C}_{Z/X}$。
\end{definition}

\noindent
因此在 $Z$ 上有层的短正合列
$$
0 \to \mathcal{C}_{Z/X} \to \mathcal{O}_{Z'} \to \mathcal{O}_Z \to 0
$$
。下面的引理证明，此定义与 $Z \to X$ 为浸入时的定义并不冲突。

\begin{lemma}
\label{more-morphisms-lemma-immersion-universal-thickening}
设 $i : Z \to X$ 是概形的浸入。则
\begin{enumerate}
\item $i$ 形式非分歧；
\item $Z$ 在 $X$ 上的泛一阶加厚就是定义
\ref{more-morphisms-definition-first-order-infinitesimal-neighbourhood} 中
$Z$ 在 $X$ 中的一阶无穷小邻域；
\item 态射章定义 \ref{morphisms-definition-conormal-sheaf}
意义下 $i$ 的余法层与定义 \ref{more-morphisms-definition-universal-thickening}
意义下 $i$ 的余法层一致。
\end{enumerate}
\end{lemma}

\begin{proof}
由 态射章引理 \ref{morphisms-lemma-open-immersion-unramified} 与
\ref{morphisms-lemma-closed-immersion-unramified}，浸入是非分歧的，
故由引理 \ref{more-morphisms-lemma-unramified-formally-unramified} 它形式非分歧。
其余断言由引理 \ref{more-morphisms-lemma-first-order-infinitesimal-neighbourhood}、
\ref{more-morphisms-lemma-infinitesimal-neighbourhood-conormal} 与诸定义合并得出。
\end{proof}

\begin{lemma}
\label{more-morphisms-lemma-universal-thickening-unramified}
设 $Z \to X$ 是形式非分歧的概形态射。则泛一阶加厚 $Z'$
在 $X$ 上形式非分歧。
\end{lemma}

\begin{proof}
有两个证明。第一个是在仿射局部运算并应用
代数章引理 \ref{algebra-lemma-differentials-universal-thickening}
来证明 $\Omega_{Z'/X} = 0$；此后引理
\ref{more-morphisms-lemma-formally-unramified-differentials} 即给出所需结论。
第二个是如下直接论证。

\medskip\noindent
设 $T \subset T'$ 是一阶加厚，并设
$$
\xymatrix{
Z' \ar[d] & T \ar[l]^c \ar[d] \\
X & T' \ar[l] \ar[lu]^{a, b}
}
$$
为交换图。考虑图中的两个态射 $a, b : T' \to Z'$。令
$T_0 = c^{-1}(Z) \subset T$，并令
$T'_a = a^{-1}(Z)$（均取概形论意义）。
由于 $Z'$ 是 $Z$ 的一阶加厚，$T'$ 是 $T'_a$ 的一阶加厚。
又因 $c = a|_T$，有 $T_0 = T \cap T'_a$（概形论意义）。
由于 $T'$ 是 $T$ 的一阶加厚，遂知 $T'_a$ 是 $T_0$ 的一阶加厚。
现在 $a|_{T'_a}$ 与 $b|_{T'_a}$ 是 $X$ 上从 $T'_a$ 到 $Z'$ 的态射，
且它们限制到 $T_0$ 后作为到 $Z$ 的态射相等。因此由 $Z'$ 的泛性质，
$a|_{T'_a} = b|_{T'_a}$。于是 $a$ 与 $b$ 是从 $T'_a$ 的一阶加厚
$T'$ 出发的态射，它们在 $T'_a$ 上的限制作为到 $Z$ 的态射相等。
再次使用 $Z'$ 的泛性质，得到 $a = b$。换言之，形式非分歧态射的定义性质
对 $Z' \to X$ 成立，正是所求。
\end{proof}

\begin{lemma}
\label{more-morphisms-lemma-universal-thickening-functorial}
考虑概形的交换图
$$
\xymatrix{
Z \ar[r]_h \ar[d]_f & X \ar[d]^g \\
W \ar[r]^{h'} & Y
}
$$
其中 $h$ 与 $h'$ 形式非分歧。设 $Z \subset Z'$ 是 $Z$ 在 $X$ 上的泛一阶加厚，
$W \subset W'$ 是 $W$ 在 $Y$ 上的泛一阶加厚。存在 $Y$ 上加厚的典范态射
$(f, f') : (Z, Z') \to (W, W')$，它嵌入如下交换图：
$$
\xymatrix{
& & & Z' \ar[ld] \ar[d]^{f'} \\
Z \ar[rr] \ar[d]_f \ar[rrru] & & X \ar[d] & W' \ar[ld] \\
W \ar[rrru]|!{[rr];[rruu]}\hole \ar[rr] & & Y
}
$$
特别地，加厚的态射 $(f, f')$ 诱导余法层的态射
$f^*\mathcal{C}_{W/Y} \to \mathcal{C}_{Z/X}$。
\end{lemma}

\begin{proof}
第一个断言由 $W'$ 的泛性质立即可见。余法层上的诱导映射，是将
态射章引理 \ref{morphisms-lemma-conormal-functorial}
应用于 $(Z \subset Z') \to (W \subset W')$ 所得的映射。
\end{proof}

\begin{lemma}
\label{more-morphisms-lemma-universal-thickening-fibre-product}
设
$$
\xymatrix{
Z \ar[r]_h \ar[d]_f & X \ar[d]^g \\
W \ar[r]^{h'} & Y
}
$$
是概形范畴中的纤维积图，并设 $h'$ 形式非分歧。则 $h$ 形式非分歧；
若 $W \subset W'$ 是 $W$ 在 $Y$ 上的泛一阶加厚，则
$Z = X \times_Y W \subset X \times_Y W'$ 是 $Z$ 在 $X$ 上的泛一阶加厚。
特别地，引理 \ref{more-morphisms-lemma-universal-thickening-functorial} 的典范映射
$f^*\mathcal{C}_{W/Y} \to \mathcal{C}_{Z/X}$ 是满射。
\end{lemma}

\begin{proof}
由引理 \ref{more-morphisms-lemma-base-change-formally-unramified}，态射 $h$ 形式非分歧。
显然 $X \times_Y W'$ 是一阶加厚。因为 $W'$ 具有泛性质，
由纤维积的映射性质容易检验它也具有泛性质。
关于这为何推出余法层之间的映射为满射，见
态射章引理 \ref{morphisms-lemma-conormal-functorial-flat}。
\end{proof}

\begin{lemma}
\label{more-morphisms-lemma-universal-thickening-fibre-product-flat}
设
$$
\xymatrix{
Z \ar[r]_h \ar[d]_f & X \ar[d]^g \\
W \ar[r]^{h'} & Y
}
$$
是概形范畴中的纤维积图，其中 $h'$ 形式非分歧且 $g$ 平坦。
在这种情形下，相应的泛一阶加厚之间的映射 $Z' \to W'$ 平坦，并且
$f^*\mathcal{C}_{W/Y} \to \mathcal{C}_{Z/X}$ 是同构。
\end{lemma}

\begin{proof}
平坦性在基变换下保持；见
态射章引理 \ref{morphisms-lemma-base-change-flat}。
因此第一个断言由引理 \ref{more-morphisms-lemma-universal-thickening-fibre-product}
中对 $W'$ 的描述得出。显然 $X \times_Y W'$ 是一阶加厚。
因为 $W'$ 具有泛性质，由纤维积的映射性质容易检验
它也具有泛性质。关于这为何推出余法层之间的映射是同构，见
态射章引理 \ref{morphisms-lemma-conormal-functorial-flat}。
\end{proof}

\begin{lemma}
\label{more-morphisms-lemma-universal-thickening-localize}
取泛一阶加厚与取开集可交换。更精确地，设 $h : Z \to X$
是形式非分歧的概形态射，$V \subset Z$、$U \subset X$ 是满足
$h(V) \subset U$ 的开集，并设 $Z'$ 是 $Z$ 在 $X$ 上的泛一阶加厚。
则 $h|_V : V \to U$ 形式非分歧，而且 $V$ 在 $U$ 上的泛一阶加厚
是满足 $V = Z \cap V'$ 的开子概形 $V' \subset Z'$。特别地，
$\mathcal{C}_{Z/X}|_V = \mathcal{C}_{V/U}$。
\end{lemma}

\begin{proof}
第一个断言即引理 \ref{more-morphisms-lemma-formally-unramified-on-opens}。
泛加厚的相容性可以从引理 \ref{more-morphisms-lemma-universal-thickening} 的证明、
代数章引理 \ref{algebra-lemma-universal-thickening-localize}，
或引理 \ref{more-morphisms-lemma-universal-thickening-fibre-product-flat} 推出。
\end{proof}

\begin{lemma}
\label{more-morphisms-lemma-differentials-universally-unramified}
设 $h : Z \to X$ 是 $S$ 上形式非分歧的概形态射。设 $Z \subset Z'$
是 $Z$ 在 $X$ 上的泛一阶加厚，其结构态射为 $h' : Z' \to X$。典范映射
$$
c_{h'} : (h')^*\Omega_{X/S} \longrightarrow \Omega_{Z'/S}
$$
诱导同构 $h^*\Omega_{X/S} \to \Omega_{Z'/S} \otimes \mathcal{O}_Z$。
\end{lemma}

\begin{proof}
映射 $c_{h'}$ 是 态射章引理
\ref{morphisms-lemma-functoriality-differentials} 中定义的映射。
若 $i : Z \to Z'$ 是给定的闭浸入，则 $i^*c_{h'}$ 是映射
$h^*\Omega_{X/S} \to \Omega_{Z'/S} \otimes \mathcal{O}_Z$。
通过局部化，检验它是同构可归结为仿射情形；见引理
\ref{more-morphisms-lemma-universal-thickening-localize} 与
态射章引理 \ref{morphisms-lemma-differentials-restrict-open}。
在仿射情形下，结论即
代数章引理 \ref{algebra-lemma-differentials-universal-thickening}。
\end{proof}

\begin{lemma}
\label{more-morphisms-lemma-universally-unramified-differentials-sequence}
设 $h : Z \to X$ 是 $S$ 上形式非分歧的概形态射。则有典范正合列
$$
\mathcal{C}_{Z/X} \to h^*\Omega_{X/S} \to \Omega_{Z/S} \to 0.
$$
第一条箭头由 $\text{d}_{Z'/S}$ 诱导，其中
% P05-ERR-CH37-0009: Frozen source calls Z' a universal first order ``neighbourhood'' here; retained pending canon review.
$Z'$ 是 $Z$ 在 $X$ 上的泛一阶邻域。
\end{lemma}

\begin{proof}
我们知道有典范正合列
$$
\mathcal{C}_{Z/Z'} \to
\Omega_{Z'/S} \otimes \mathcal{O}_Z \to
\Omega_{Z/S} \to 0.
$$
；见 态射章引理
\ref{morphisms-lemma-differentials-relative-immersion}。
因此应用引理 \ref{more-morphisms-lemma-differentials-universally-unramified} 即得结论。
\end{proof}

\begin{lemma}
\label{more-morphisms-lemma-two-unramified-morphisms}
设
$$
\xymatrix{
Z \ar[r]_i \ar[rd]_j & X \ar[d] \\
& Y
}
$$
是概形的交换图，其中 $i$ 与 $j$ 形式非分歧。则有典范正合列
$$
\mathcal{C}_{Z/Y} \to
\mathcal{C}_{Z/X} \to
i^*\Omega_{X/Y} \to 0
$$
其中第一条箭头来自引理 \ref{more-morphisms-lemma-universal-thickening-functorial}，
第二条来自引理 \ref{more-morphisms-lemma-universally-unramified-differentials-sequence}。
\end{lemma}

\begin{proof}
以 $Z \to Z'$ 表示 $Z$ 在 $X$ 上的泛一阶加厚，以 $Z \to Z''$
表示 $Z$ 在 $Y$ 上的泛一阶加厚。
% P05-ERR-CH37-0010: Frozen source cites the differential-sequence lemma for this functorial morphism; retained pending canon review.
由引理 \ref{more-morphisms-lemma-universally-unramified-differentials-sequence}，
有典范态射 $Z' \to Z''$，从而有交换图
$$
\xymatrix{
Z \ar[r]_{i'} \ar[rd]_{j'} & Z' \ar[r] \ar[d] & X \ar[d] \\
& Z'' \ar[r] & Y
}
$$
将 态射章引理 \ref{morphisms-lemma-two-immersions}
应用于左边的三角形，得到正合列
$$
\mathcal{C}_{Z/Z''} \to
\mathcal{C}_{Z/Z'} \to
(i')^*\Omega_{Z'/Z''} \to 0
$$
由于 $Z''$ 在 $Y$ 上形式非分歧（见引理
\ref{more-morphisms-lemma-universal-thickening-unramified}），有
% P05-ERR-CH37-0011: Frozen source writes Omega_{Z/Y}; the next sentence uses Omega_{Z'/Y}. Formula retained pending canon review.
$\Omega_{Z'/Z''} = \Omega_{Z/Y}$（结合引理
\ref{more-morphisms-lemma-formally-unramified-differentials} 与
态射章引理 \ref{morphisms-lemma-triangle-differentials}）。
再由引理 \ref{more-morphisms-lemma-differentials-universally-unramified}，有
$(i')^*\Omega_{Z'/Y} = i^*\Omega_{X/Y}$。
\end{proof}

\begin{lemma}
\label{more-morphisms-lemma-transitivity-conormal}
设 $Z \to Y \to X$ 是形式非分歧的概形态射。
\begin{enumerate}
\item 若 $Z \subset Z'$ 是 $Z$ 在 $X$ 上的泛一阶加厚，且
$Y \subset Y'$ 是 $Y$ 在 $X$ 上的泛一阶加厚，则存在态射
$Z' \to Y'$，并且 $Y \times_{Y'} Z'$ 是 $Z$ 在 $Y$ 上的泛一阶加厚。
\item 有典范正合列
$$
i^*\mathcal{C}_{Y/X} \to
\mathcal{C}_{Z/X} \to
\mathcal{C}_{Z/Y} \to 0
$$
其中各映射来自引理 \ref{more-morphisms-lemma-universal-thickening-functorial}，
且 $i : Z \to Y$ 是第一个态射。
\end{enumerate}
\end{lemma}

\begin{proof}
(1) 中的映射 $h : Z' \to Y'$ 来自引理
\ref{more-morphisms-lemma-universal-thickening-functorial}。由 $Z'$ 与 $Y'$ 的泛性质，
$Y \times_{Y'} Z'$ 是 $Z$ 在 $Y$ 上的泛一阶加厚这一断言立即可见。
由 态射章引理 \ref{morphisms-lemma-transitivity-conormal}，有正合列
$$
(i')^*\mathcal{C}_{Y \times_{Y'} Z'/Z'} \to
\mathcal{C}_{Z/Z'} \to
\mathcal{C}_{Z/Y \times_{Y'} Z'} \to 0
$$
其中 $i' : Z \to Y \times_{Y'} Z'$ 是给定态射。由
态射章引理 \ref{morphisms-lemma-conormal-functorial-flat}，
存在满射
$h^*\mathcal{C}_{Y/Y'} \to \mathcal{C}_{Y \times_{Y'} Z'/Z'}$。
再结合等式 $\mathcal{C}_{Y/Y'} = \mathcal{C}_{Y/X}$、
$\mathcal{C}_{Z/Z'} = \mathcal{C}_{Z/X}$ 与
$\mathcal{C}_{Z/Y \times_{Y'} Z'} = \mathcal{C}_{Z/Y}$，即证引理。
\end{proof}











\section{形式平展态射}
\label{more-morphisms-section-formally-etale}

\noindent
回顾：若对每个实线部分交换的图
$$
\xymatrix{
A \ar[r] \ar@{-->}[rd] & B/I \\
R \ar[r] \ar[u] & B \ar[u]
}
$$
（其中 $I \subset B$ 是平方为零的理想），恰好存在一条使该图交换的虚线箭头，
则称环映射 $R \to A$ 是{\it 形式平展的}；见
代数章定义 \ref{algebra-definition-formally-etale}。
这引出概形态射的如下类似概念。

\begin{definition}
\label{more-morphisms-definition-formally-etale}
设 $f : X \to S$ 是概形态射。若给定任意实线部分交换的图
$$
\xymatrix{
X \ar[d]_f & T \ar[d]^i \ar[l] \\
S & T' \ar[l] \ar@{-->}[lu]
}
$$
其中 $T \subset T'$ 是 $S$ 上仿射概形的一阶加厚时，恰好存在一条使该图交换的
虚线箭头，则称 $f$ 是{\it 形式平展的}。
\end{definition}

\noindent
显然，形式平展态射是形式非分歧的。因此若 $f : X \to S$ 形式平展，
则 $\Omega_{X/S}$ 为零；见引理
\ref{more-morphisms-lemma-formally-unramified-differentials}。

\begin{lemma}
\label{more-morphisms-lemma-formally-etale-not-affine}
若 $f : X \to S$ 是形式平展态射，则给定任意实线部分交换的图
$$
\xymatrix{
X \ar[d]_f & T \ar[d]^i \ar[l] \\
S & T' \ar[l] \ar@{-->}[lu]
}
$$
其中 $T \subset T'$ 是 $S$ 上概形的一阶加厚，恰好存在一条使该图交换的虚线箭头。
换言之，在定义 \ref{more-morphisms-definition-formally-etale} 中可以去掉 $T$ 为仿射的条件。
\end{lemma}

\begin{proof}
设 $T' = \bigcup T'_i$ 是仿射开覆盖，并令 $T_i = T \cap T'_i$。
于是得到嵌入图中的态射 $a'_i : T'_i \to X$。由唯一性，
$a'_i$ 与 $a'_j$ 在 $T'_i \cap T'_j$ 的任意仿射开子概形上相等，
故 $a'_i$ 与 $a'_j$ 在 $T'_i \cap T'_j$ 上相等。因此诸态射 $a'_i$ 粘合成整体态射
$a' : T' \to X$。$a'$ 的唯一性已见于引理
\ref{more-morphisms-lemma-formally-unramified-not-affine}。
\end{proof}

\begin{lemma}
\label{more-morphisms-lemma-composition-formally-etale}
形式平展态射的复合仍是形式平展的。
\end{lemma}

\begin{proof}
这是形式性的。
\end{proof}

\begin{lemma}
\label{more-morphisms-lemma-base-change-formally-etale}
形式平展态射的基变换仍是形式平展的。
\end{lemma}

\begin{proof}
这是形式性的。
\end{proof}

\begin{lemma}
\label{more-morphisms-lemma-formally-etale-on-opens}
设 $f : X \to S$ 是概形态射，并设 $U \subset X$ 与 $V \subset S$
是满足 $f(U) \subset V$ 的开子概形。若 $f$ 形式平展，
则 $f|_U : U \to V$ 也形式平展。
\end{lemma}

\begin{proof}
考虑定义 \ref{more-morphisms-definition-formally-etale} 中那样的实线部分交换图
$$
\xymatrix{
U \ar[d]_{f|_U} & T \ar[d]^i \ar[l]^a \\
V & T' \ar[l] \ar@{-->}[lu]
}
$$
。
% P05-ERR-CH37-0012: Frozen source says ``formally ramified'' here; retained literally pending canon review.
若 $f$ 形式分歧，则恰好存在一个满足 $a'|_T = a$ 的 $S$-态射
$a' : T' \to X$。由于 $|T'| = |T|$，可知 $a'(T') \subset U$，
这就给出从 $T'$ 到 $U$ 的唯一态射。
\end{proof}

\begin{lemma}
\label{more-morphisms-lemma-characterize-formally-etale}
设 $f : X \to S$ 是概形态射。以下条件等价：
\begin{enumerate}
\item $f$ 形式平展；
\item $f$ 形式非分歧，且 $X$ 在 $S$ 上的泛一阶加厚等于 $X$；
\item $f$ 形式非分歧，且 $\mathcal{C}_{X/S} = 0$；
\item $\Omega_{X/S} = 0$ 且 $\mathcal{C}_{X/S} = 0$。
\end{enumerate}
\end{lemma}

\begin{proof}
事实上，最后一个断言之所以有意义，是因为 $\Omega_{X/S} = 0$
意味着可由引理 \ref{more-morphisms-lemma-formally-unramified-differentials} 与
定义 \ref{more-morphisms-definition-universal-thickening} 定义 $\mathcal{C}_{X/S}$。
这也说明 (3) 与 (4) 等价。

\medskip\noindent
假设 (1)、(2)、(3) 中任一个，都推出 $f$ 形式非分歧。因此可以假设
$f$ 形式非分歧。(1)、(2)、(3) 的等价性来自 $X$ 在 $S$ 上的泛一阶加厚
$X'$ 的泛性质以及如下事实：
$X = X' \Leftrightarrow \mathcal{C}_{X/S} = 0$；毕竟按定义，
$\mathcal{C}_{X/S} = \mathcal{C}_{X/X'}$ 是 $X$ 在 $X'$ 中的理想层。
\end{proof}

\begin{lemma}
\label{more-morphisms-lemma-unramified-flat-formally-etale}
非分歧平坦态射是形式平展的。
\end{lemma}

\begin{proof}
设 $X \to S$ 非分歧且平坦。则 $\Delta : X \to X \times_S X$
是开浸入；见 态射章引理
\ref{morphisms-lemma-diagonal-unramified-morphism}。
我们须证明 $\mathcal{C}_{X/S}$ 为零。考虑两个投影
$p, q : X \times_S X \to X$。由于 $f$ 形式非分歧（见引理
\ref{more-morphisms-lemma-unramified-formally-unramified}），$q$ 形式非分歧（见引理
\ref{more-morphisms-lemma-base-change-formally-unramified}）。由于 $f$ 平坦，$p$ 平坦；
见 态射章引理 \ref{morphisms-lemma-base-change-flat}。
因此由引理 \ref{more-morphisms-lemma-universal-thickening-fibre-product-flat}，
$p^*\mathcal{C}_{X/S} = \mathcal{C}_q$，其中 $\mathcal{C}_q$
表示形式非分歧态射 $q : X \times_S X \to X$ 的余法层。
但 $\Delta(X) \subset X \times_S X$ 是经 $q$ 同构地映到 $X$ 的开子概形。
故由引理 \ref{more-morphisms-lemma-universal-thickening-localize}，
$\mathcal{C}_q|_{\Delta(X)} = \mathcal{C}_{X/X} = 0$。
换言之，$\mathcal{C}_{X/S}$ 经恒等态射拉回到 $X$ 上为零，即
$\mathcal{C}_{X/S} = 0$。
\end{proof}

\begin{lemma}
\label{more-morphisms-lemma-affine-formally-etale}
设 $f : X \to S$ 是概形态射，并假设 $X$ 与 $S$ 均为仿射。
则 $f$ 形式平展，当且仅当
$\mathcal{O}_S(S) \to \mathcal{O}_X(X)$ 是形式平展环映射。
\end{lemma}

\begin{proof}
由环范畴与仿射概形范畴之间的等价，这立即从定义
（定义 \ref{more-morphisms-definition-formally-etale} 与
代数章定义 \ref{algebra-definition-formally-etale}）得出；见
概形章引理 \ref{schemes-lemma-category-affine-schemes}。
\end{proof}

\begin{lemma}
\label{more-morphisms-lemma-formally-etale}
设 $f : X \to Y$ 是概形态射。以下条件等价：
\begin{enumerate}
\item $f$ 形式平展；
\item 对每个 $x \in X$，存在满足 $f(U) \subset V$ 的开集
$x \in U \subset X$ 与 $f(x) \in V \subset Y$，使得
$f|_U : U \to V$ 形式平展；
\item 对每一对满足 $f(U) \subset V$ 的仿射开集 $U \subset X$ 与
$V \subset Y$，环映射 $\mathcal{O}_Y(V) \to \mathcal{O}_X(U)$ 形式平展；
\item 存在仿射开覆盖 $Y = \bigcup V_j$，且对每个 $j$ 存在仿射开覆盖
$f^{-1}(V_j) = \bigcup U_{ji}$，使得对所有 $j$ 与 $i$，环映射
$\mathcal{O}_Y(V) \to \mathcal{O}_X(U)$ 形式平展。
\end{enumerate}
\end{lemma}

\begin{proof}
结合引理 \ref{more-morphisms-lemma-formally-etale-on-opens} 与
\ref{more-morphisms-lemma-affine-formally-etale} 得 (1) $\Rightarrow$ (3)。
(3) $\Rightarrow$ (4) 是立即的。由引理 \ref{more-morphisms-lemma-affine-formally-etale}
得 (4) $\Rightarrow$ (2)。若 (2) 成立，则由引理
\ref{more-morphisms-lemma-characterize-formally-etale}（这里还使用 态射章引理
\ref{morphisms-lemma-differentials-restrict-open} 与
引理 \ref{more-morphisms-lemma-universal-thickening-localize}），
$\Omega_{X/S}$ 与 $\mathcal{C}_{X/S}$ 在每一点的某个邻域上消失；于是由引理
\ref{more-morphisms-lemma-characterize-formally-etale}，
% P05-ERR-CH37-0013: Frozen source concludes ``formally unramified'' rather than ``formally etale''; retained pending canon review.
$f$ 形式非分歧。
\end{proof}

\begin{lemma}
\label{more-morphisms-lemma-etale-formally-etale}
设 $f : X \to S$ 是概形态射。以下条件等价：
\begin{enumerate}
\item 态射 $f$ 平展；
\item 态射 $f$ 局部有限表示且形式平展。
\end{enumerate}
\end{lemma}

\begin{proof}
假设 $f$ 平展。平展态射局部有限表示、平坦且非分歧；见
态射章第 \ref{morphisms-section-etale} 节。
故 $f$ 局部有限表示且形式平展；见引理
\ref{more-morphisms-lemma-unramified-flat-formally-etale}。

\medskip\noindent
反之，假设 $f$ 局部有限表示且形式平展。平展性在 $X$ 与 $S$ 的 Zariski
拓扑上是局部的；见 态射章引理
\ref{morphisms-lemma-etale-characterize}。由引理
\ref{more-morphisms-lemma-formally-etale-on-opens}，可用映入仿射开集 $V$ 的仿射开集
$U$ 覆盖 $X$，使 $U \to V$ 形式平展（且有限表示；见 Morphisms,
引理 \ref{morphisms-lemma-locally-finite-presentation-characterize}）。
由引理 \ref{more-morphisms-lemma-affine-formally-etale}，环映射
$\mathcal{O}(V) \to \mathcal{O}(U)$ 形式平展（且有限表示）。
于是由 代数章引理 \ref{algebra-lemma-formally-etale-etale} 得证。
（讨论形式光滑态射时，我们还将给出这一蕴含的另一个证明。）
\end{proof}













\section{映射的无穷小形变}
\label{more-morphisms-section-action-by-derivations}

\noindent
本节说明如何用导子使映射作无穷小移动。贯穿本节，我们使用如下事实：
$X$ 的加厚 $X'$ 上的层可以视为 $X$ 上的层。

\begin{lemma}
\label{more-morphisms-lemma-difference-derivation}
设 $S$ 是概形，$X \subset X'$ 与 $Y \subset Y'$ 是 $S$ 上的两个一阶加厚，
且 $(a, a'), (b, b') : (X \subset X') \to (Y \subset Y')$
是 $S$ 上加厚的两个态射。假设
\begin{enumerate}
\item $a = b$；
\item 两个映射 $a^*\mathcal{C}_{Y/Y'} \to \mathcal{C}_{X/X'}$
（态射章引理 \ref{morphisms-lemma-conormal-functorial}）相等。
\end{enumerate}
则映射 $(a')^\sharp - (b')^\sharp$ 分解为
$$
\mathcal{O}_{Y'} \to \mathcal{O}_Y \xrightarrow{D}
a_*\mathcal{C}_{X/X'} \to a_*\mathcal{O}_{X'}
$$
其中 $D$ 是 $\mathcal{O}_S$-导子。
\end{lemma}

\begin{proof}
不在 $Y$ 上而在 $X$ 上运算；其优点是拉回函子 $a^{-1}$ 正合。
由 (1) 与 (2)，得到行正合的交换图
$$
\xymatrix{
0 \ar[r] &
\mathcal{C}_{X/X'} \ar[r] &
\mathcal{O}_{X'} \ar[r] &
\mathcal{O}_X \ar[r] & 0 \\
0 \ar[r] &
a^{-1}\mathcal{C}_{Y/Y'} \ar[r] \ar[u] &
a^{-1}\mathcal{O}_{Y'}
\ar[r] \ar@<1ex>[u]^{(a')^\sharp} \ar@<-1ex>[u]_{(b')^\sharp} &
a^{-1}\mathcal{O}_Y \ar[r] \ar[u] & 0
}
$$
一般而言，在这种情形下，$\mathcal{O}_S$-代数映射
$(a')^\sharp$ 与 $(b')^\sharp$ 之差是从 $a^{-1}\mathcal{O}_Y$ 到
$\mathcal{C}_{X/X'}$ 的 $\mathcal{O}_S$-导子。由函子 $a^{-1}$ 与 $a_*$
的伴随性，这等价于从 $\mathcal{O}_Y$ 到
$a_*\mathcal{C}_{X/X'}$ 的 $\mathcal{O}_S$-导子。略去一些细节。
\end{proof}

\noindent
注意，在上述引理的情形下，可以把 $D$ 写成
\begin{equation}
\label{more-morphisms-equation-D}
D = \text{d}_{Y/S} \circ \theta
\end{equation}
其中 $\theta$ 是 $\mathcal{O}_Y$-线性映射
$\theta : \Omega_{Y/S} \to a_*\mathcal{C}_{X/X'}$。
当然，再次由伴随性，也可把 $\theta$ 看作 $\mathcal{O}_X$-线性映射
$\theta : a^*\Omega_{Y/S} \to \mathcal{C}_{X/X'}$。

\begin{lemma}
\label{more-morphisms-lemma-action-by-derivations}
设 $S$ 是概形，且
$(a, a') : (X \subset X') \to (Y \subset Y')$
是 $S$ 上一阶加厚的态射。设
$$
\theta : a^*\Omega_{Y/S} \to \mathcal{C}_{X/X'}
$$
是 $\mathcal{O}_X$-线性映射。则存在唯一的对的态射
$(b, b') : (X \subset X') \to (Y \subset Y')$，使引理
\ref{more-morphisms-lemma-difference-derivation} 的 (1)、(2) 成立，并且导子 $D$
与 $\theta$ 由等式 (\ref{more-morphisms-equation-D}) 联系。
\end{lemma}

\begin{proof}
只需令 $b = a$，并将 $(b')^\sharp$ 定义为映射
$$
(a')^\sharp + D : a^{-1}\mathcal{O}_{Y'} \to \mathcal{O}_{X'}
$$
其中 $D$ 如等式 (\ref{more-morphisms-equation-D}) 所示。略去如下验证：
$(b')^\sharp$ 是 $\mathcal{O}_S$-代数层的映射，且引理
\ref{more-morphisms-lemma-difference-derivation} 的 (1)、(2) 成立。
等式 (\ref{more-morphisms-equation-D}) 由构造成立。
\end{proof}

\begin{remark}
\label{more-morphisms-remark-action-by-derivations}
假设与记号同引理 \ref{more-morphisms-lemma-action-by-derivations}。
局部截面 $\theta$ 在 $a'$ 上的作用有时记为 $\theta \cdot a'$。
注意，这无非表示 $(a')^\sharp$ 与 $(\theta \cdot a')^\sharp$ 之差为导子
$D$，它与 $\theta$ 由等式 (\ref{more-morphisms-equation-D}) 联系。
\end{remark}

\begin{lemma}
\label{more-morphisms-lemma-sheaf}
设 $S$ 是概形，$X \subset X'$ 与 $Y \subset Y'$ 是 $S$ 上的一阶加厚。
假设给定态射 $a : X \to Y$ 以及 $\mathcal{O}_X$-模的映射
$A : a^*\mathcal{C}_{Y/Y'} \to \mathcal{C}_{X/X'}$。
对开子概形 $U' \subset X'$，考虑满足下列条件的态射 $a' : U' \to Y'$：
\begin{enumerate}
\item $a'$ 是 $S$ 上的态射；
\item $a'|_U = a|_U$；
\item 诱导映射 $a^*\mathcal{C}_{Y/Y'}|_U \to \mathcal{C}_{X/X'}|_U$
是 $A$ 在 $U$ 上的限制。
\end{enumerate}
这里 $U = X \cap U'$。则规则
\begin{equation}
\label{more-morphisms-equation-sheaf}
U' \mapsto
\{a' : U' \to Y'\text{ such that (1), (2), (3) hold.}\}
\end{equation}
定义了 $X'$ 上的集合层。
\end{lemma}

\begin{proof}
用 $\mathcal{F}$ 表示引理中的规则。对
$V' \subset U' \subset X'$，$\mathcal{F}$ 的限制映射
$\mathcal{F}(U') \to \mathcal{F}(V')$ 就是限制映射
$a' \mapsto a'|_{V'}$。在此定义下，由于态射是局部定义的，
$\mathcal{F}$ 显然是层。
\end{proof}

\noindent
在下一个引理中，我们将 $X$ 上的层同 $X$ 的任意加厚上的层等同起来。

\begin{lemma}
\label{more-morphisms-lemma-action-sheaf}
记号与假设同引理 \ref{more-morphisms-lemma-sheaf}。层
$$
\SheafHom_{\mathcal{O}_X}(a^*\Omega_{Y/S}, \mathcal{C}_{X/X'})
$$
作用在层 (\ref{more-morphisms-equation-sheaf}) 上。此外，对层
(\ref{more-morphisms-equation-sheaf}) 在其上有截面的任意开集 $U' \subset X'$，
该作用是单传递的。
\end{lemma}

\begin{proof}
这是引理 \ref{more-morphisms-lemma-difference-derivation}、
\ref{more-morphisms-lemma-action-by-derivations} 与 \ref{more-morphisms-lemma-sheaf} 的合并。
\end{proof}

\begin{remark}
\label{more-morphisms-remark-special-case}
引理 \ref{more-morphisms-lemma-difference-derivation}、
\ref{more-morphisms-lemma-action-by-derivations}、\ref{more-morphisms-lemma-sheaf} 与
\ref{more-morphisms-lemma-action-sheaf} 的一个特殊情形是 $Y = Y'$。
此时映射 $A$ 总为零，引理 \ref{more-morphisms-lemma-sheaf} 的层就是由规则
$$
U' \mapsto
\{a' : U' \to Y\text{ 相对于 }S\text{ 其中 } a'|_U = a|_U\}
$$
给出的层，而层
$\SheafHom_{\mathcal{O}_X}(a^*\Omega_{Y/S}, \mathcal{C}_{X/X'})$
作用在其上。
\end{remark}

\begin{remark}
\label{more-morphisms-remark-another-special-case}
引理 \ref{more-morphisms-lemma-difference-derivation}、
\ref{more-morphisms-lemma-action-by-derivations}、\ref{more-morphisms-lemma-sheaf} 与
\ref{more-morphisms-lemma-action-sheaf} 的另一个特殊情形是：$S$ 本身为加厚
$Z \subset Z' = S$，且 $Y = Z \times_{Z'} Y'$。图示为
$$
\xymatrix{
(X \subset X') \ar@{..>}[rr]_{(a, ?)} \ar[rd]_{(g, g')} & &
(Y \subset Y') \ar[ld]^{(h, h')} \\
& (Z \subset Z')
}
$$
此时映射 $A : a^*\mathcal{C}_{Y/Y'} \to \mathcal{C}_{X/X'}$
由 $a$ 决定：映射 $h^*\mathcal{C}_{Z/Z'} \to \mathcal{C}_{Y/Y'}$
是满射（因为假设了 $Y = Z \times_{Z'} Y'$），所以拉回
$g^*\mathcal{C}_{Z/Z'} = a^*h^*\mathcal{C}_{Z/Z'} \to
a^*\mathcal{C}_{Y/Y'}$ 是满射；而复合
$g^*\mathcal{C}_{Z/Z'} \to a^*\mathcal{C}_{Y/Y'} \to \mathcal{C}_{X/X'}$
必须是由 $g'$ 诱导的典范映射。因此引理 \ref{more-morphisms-lemma-sheaf} 的层
就是由规则
$$
U' \mapsto
\{a' : U' \to Y'\text{ 相对于 }Z'\text{ 其中 } a'|_U = a|_U\}
$$
给出的层，而层
$\SheafHom_{\mathcal{O}_X}(a^*\Omega_{Y/Z}, \mathcal{C}_{X/X'})$
作用在其上。
\end{remark}

\begin{lemma}
\label{more-morphisms-lemma-omega-deformation}
设 $S$ 是概形，$X \subset X'$ 是 $S$ 上的一阶加厚，$Y$ 是 $S$ 上的概形。
设 $a', b' : X' \to Y$ 是 $S$ 上满足
$a = a'|_X = b'|_X$ 的两个态射。由此得到交换图
$$
\xymatrix{
X \ar[r] \ar[d]_a & X' \ar[d]^{(b', a')} \\
Y \ar[r]^-{\Delta_{Y/S}} & Y \times_S Y
}
$$
由于水平箭头都是浸入，其余法层分别为
$\mathcal{C}_{X/X'}$ 与 $\Omega_{Y/S}$，由
态射章引理 \ref{morphisms-lemma-conormal-functorial}，
得到映射 $\theta : a^*\Omega_{Y/S} \to \mathcal{C}_{X/X'}$。
于是此 $\theta$ 与引理 \ref{more-morphisms-lemma-difference-derivation} 的导子 $D$
由等式 (\ref{more-morphisms-equation-D}) 联系。
\end{lemma}

\begin{proof}
从略。提示：可在仿射开集上检验该等式，此时它来自如下计算。
若 $f$ 是 $\mathcal{O}_Y$ 的局部截面，则
$1 \otimes f - f \otimes 1$ 是 $\mathcal{C}_{Y/(Y \times_S Y)}$
中对应于 $\text{d}_{Y/S}(f)$ 的局部截面。它被映到
$\mathcal{C}_{X/X'}$ 的局部截面
$(a')^\sharp(f) - (b')^\sharp(f) = D(f)$。换言之，
$\theta(\text{d}_{Y/S}(f)) = D(f)$。
\end{proof}

\noindent
为后续用途，需要一个粗略地说如下的结果：引理
\ref{more-morphisms-lemma-action-by-derivations} 的构造与平展局部化相容。

\begin{lemma}
\label{more-morphisms-lemma-sheaf-differentials-etale-localization}
设
$$
\xymatrix{
X_1 \ar[d] & X_2 \ar[l]^f \ar[d] \\
S_1 & S_2 \ar[l]
}
$$
是概形的交换图，其中 $X_2 \to X_1$ 与 $S_2 \to S_1$ 平展。则
态射章引理 \ref{morphisms-lemma-functoriality-differentials}
的映射 $c_f : f^*\Omega_{X_1/S_1} \to \Omega_{X_2/S_2}$ 是同构。
\end{lemma}

\begin{proof}
回顾，平展态射 $U \to V$ 是满足 $\Omega_{U/V} = 0$ 的光滑态射。
由此可见，态射章引理
\ref{morphisms-lemma-triangle-differentials} 推出
$\Omega_{X_2/S_2} = \Omega_{X_2/S_1}$，而 态射章引理
\ref{morphisms-lemma-triangle-differentials-smooth} 推出映射
$f^*\Omega_{X_1/S_1} \to \Omega_{X_2/S_1}$（将态射 $f$ 看作
$S_1$ 上的态射）是同构。故引理成立。
\end{proof}

\begin{lemma}
\label{more-morphisms-lemma-action-by-derivations-etale-localization}
考虑一阶加厚的交换图
$$
\vcenter{
\xymatrix{
(T_2 \subset T_2') \ar[d]_{(h, h')} \ar[rr]_{(a_2, a_2')} & &
(X_2 \subset X_2') \ar[d]^{(f, f')} \\
(T_1 \subset T_1') \ar[rr]^{(a_1, a_1')} & &
(X_1 \subset X_1')
}
}
\quad
\begin{matrix}
\text{and a commutative} \\
\text{diagram of schemes}
\end{matrix}
\quad
\vcenter{
\xymatrix{
X_2' \ar[r] \ar[d] & S_2 \ar[d] \\
X_1' \ar[r] & S_1
}
}
$$
其中 $X_2 \to X_1$ 与 $S_2 \to S_1$ 平展。对任意
$\mathcal{O}_{T_1}$-线性映射
$\theta_1 : a_1^*\Omega_{X_1/S_1} \to \mathcal{C}_{T_1/T'_1}$，
令 $\theta_2$ 为复合
$$
\xymatrix{
a_2^*\Omega_{X_2/S_2} \ar@{=}[r] &
h^*a_1^*\Omega_{X_1/S_1} \ar[r]^-{h^*\theta_1} &
h^*\mathcal{C}_{T_1/T'_1} \ar[r] &
\mathcal{C}_{T_2/T'_2}
}
$$
（等号将在证明中说明）。则图
$$
\xymatrix{
T_2' \ar[rr]_{\theta_2 \cdot a_2'} \ar[d] & & X'_2 \ar[d] \\
T_1' \ar[rr]^{\theta_1 \cdot a_1'} & & X'_1
}
$$
交换，其中作用 $\theta_2 \cdot a_2'$ 与 $\theta_1 \cdot a_1'$
如注记 \ref{more-morphisms-remark-action-by-derivations} 所述。
\end{lemma}

\begin{proof}
该等号来自引理 \ref{more-morphisms-lemma-sheaf-differentials-etale-localization}
的等同 $f^*\Omega_{X_1/S_1} = \Omega_{X_2/S_2}$。事实上，由此以及
$f \circ a_2 = a_1 \circ h$，有
$a_2^*\Omega_{X_2/S_2} = a_2^*f^*\Omega_{X_1/S_1} =
h^*a_1^*\Omega_{X_1/S_1}$。说明这一点后，可以在仿射开集上检验该图的
交换性。因此可假设最初大图中的概形都是仿射的。于是得到交换图
$$
\vcenter{
\xymatrix{
(B'_2, I_2) & & (A'_2, J_2) \ar[ll]^{a_2'} \\
(B'_1, I_1) \ar[u]^{h'} & & (A'_1, J_1) \ar[ll]_{a_1'} \ar[u]_{f'}
}
}
\quad\text{且}\quad
\vcenter{
\xymatrix{
A'_2 & & R_2 \ar[ll] \\
A'_1 \ar[u] & & R_1 \ar[ll] \ar[u]
}
}
$$
该记号表示 $I_1, I_2, J_1, J_2$ 是平方为零的理想，而对之间的映射是把理想
送入理想的环映射。令 $A_1 = A'_1/J_1$、$A_2 = A'_2/J_2$、
$B_1 = B'_1/I_1$ 且 $B_2 = B'_2/I_2$。已知
$$
A_2 \otimes_{A_1} \Omega_{A_1/R_1} \longrightarrow \Omega_{A_2/R_2}
$$
是同构。映射
$\theta_1 : B_1 \otimes_{A_1} \Omega_{A_1/R_1} \to I_1$
是 $B_1$-线性的，因而给出 $R_1$-导子
$D_1 = \theta_1 \circ \text{d}_{A_1/R_1} : A_1 \to I_1$。
类似地，映射
$\theta_2 : B_2 \otimes_{A_2} \Omega_{A_2/R_2} \to I_2$
给出 $R_2$-导子
$D_2 = \theta_2 \circ \text{d}_{A_2/R_2} : A_2 \to I_2$。
$\theta_2$ 的构造蕴含 $\theta_1$ 与 $\theta_2$ 的如下相容性：
对每个 $x \in A_1$，有
$$
h'(D_1(x)) = D_2(f'(x))
$$
，作为 $I_2$ 的元素。通过
$A'_1 \to A_1 \xrightarrow{D_1} I_1 \to B_1$，可把 $D_1$
看作映射 $A'_1 \to B'_1$；类似地，可把 $D_2$ 看作映射
$A'_2 \to B'_2$。于是上述等式对 $x \in A'_1$ 成立。
由引理 \ref{more-morphisms-lemma-action-by-derivations} 与注记
\ref{more-morphisms-remark-action-by-derivations} 中作用的构造，
$\theta_1 \cdot a_1'$ 对应于环映射 $a_1' + D_1 : A'_1 \to B'_1$，
而 $\theta_2 \cdot a_2'$ 对应于环映射
$a_2' + D_2 : A'_2 \to B'_2$。由上述等式得到
$h' \circ (a_1' + D_1) = (a_2' + D_2) \circ f'$
，正是所需结论。
\end{proof}

\begin{remark}
\label{more-morphisms-remark-tiny-improvement}
引理 \ref{more-morphisms-lemma-action-by-derivations-etale-localization}
可作如下改进。假设有引理
\ref{more-morphisms-lemma-action-by-derivations-etale-localization} 中的交换图，但不假设
$X_2 \to X_1$ 与 $S_2 \to S_1$ 平展。再假设有
$\theta_1 : a_1^*\Omega_{X_1/S_1} \to \mathcal{C}_{T_1/T'_1}$
及
$\theta_2 : a_2^*\Omega_{X_2/S_2} \to \mathcal{C}_{T_2/T'_2}$
使得
$$
\xymatrix{
f_*\mathcal{O}_{X_2} \ar[rr]_{f_*D_2} & &
f_*a_{2, *}\mathcal{C}_{T_2/T_2'} \\
\mathcal{O}_{X_1} \ar[rr]^{D_1} \ar[u]^{f^\sharp} & &
a_{1, *}\mathcal{C}_{T_1/T_1'} \ar[u]_{\text{由下列对象诱导： }(h')^\sharp}
}
$$
交换，其中 $D_i$ 如等式 (\ref{more-morphisms-equation-D}) 那样对应于 $\theta_i$。
则引理 \ref{more-morphisms-lemma-action-by-derivations-etale-localization} 的结论成立。
$X_2 \to X_1$ 与 $S_2 \to S_1$ 均平展这一条件的重要性在于，
它使我们能够从 $\theta_1$ 构造 $\theta_2$。
\end{remark}








\section{概形的无穷小形变}
\label{more-morphisms-section-deform}

\noindent
下面的简单引理常是检验映射的无穷小形变是否平坦的便利工具。

\begin{lemma}
\label{more-morphisms-lemma-deform}
设 $(f, f') : (X \subset X') \to (S \subset S')$ 是一阶加厚的态射，
并假设 $f$ 平坦。则以下条件等价：
\begin{enumerate}
\item $f'$ 平坦，且 $X = S \times_{S'} X'$；
\item 典范映射 $f^*\mathcal{C}_{S/S'} \to \mathcal{C}_{X/X'}$ 是同构。
\end{enumerate}
\end{lemma}

\begin{proof}
由于问题在 $X'$ 上是局部的，可假设 $X, X', S, S'$ 都是仿射概形。
写 $S' = \Spec(A')$、$X' = \Spec(B')$、$S = \Spec(A)$、
$X = \Spec(B)$，其中对某些平方为零的理想有
$A = A'/I$ 且 $B = B'/J$。于是得到交换图
$$
\xymatrix{
0 \ar[r] &
J \ar[r] &
B' \ar[r] &
B \ar[r] & 0 \\
0 \ar[r] &
I \ar[r] \ar[u] &
A' \ar[r] \ar[u] &
A \ar[r] \ar[u] & 0
}
$$
，其行正合。引理中的典范映射是
$$
I \otimes_A B = I \otimes_{A'} B' \longrightarrow J.
$$
$f$ 平坦这一假设表示 $A \to B$ 平坦。

\medskip\noindent
假设 (1)。则 $A' \to B'$ 平坦且 $J = IB'$。平坦性推出
$\text{Tor}_1^{A'}(B', A) = 0$（见
代数章引理 \ref{algebra-lemma-characterize-flat}）。
这表示 $I \otimes_{A'} B' \to B'$ 为单射（见
代数章说明 \ref{algebra-remark-Tor-ring-mod-ideal}）。
故 $I \otimes_A B \to J$ 是同构。

\medskip\noindent
假设 (2)。由此 $J = IB'$，故 $X = S \times_{S'} X'$。
此外，反向使用上一段的诸蕴含，得到
$\text{Tor}_1^{A'}(B', A'/I) = 0$。因此由
代数章引理 \ref{algebra-lemma-what-does-it-mean}，
$B'$ 在 $A'$ 上平坦。
\end{proof}

\noindent
下面的引理是“逐纤维平坦性判据”（``crit\`ere de platitude par fibres''）
的“幂零”版本；见节 \ref{more-morphisms-section-criterion-flat-fibres}。

\begin{lemma}
\label{more-morphisms-lemma-flatness-morphism-thickenings}
考虑交换图
$$
\xymatrix{
(X \subset X') \ar[rr]_{(f, f')} \ar[rd] & & (Y \subset Y') \ar[ld] \\
& (S \subset S')
}
$$
，其中各对象为加厚。假设
\begin{enumerate}
\item $X'$ 在 $S'$ 上平坦；
\item $f$ 平坦；
\item $S \subset S'$ 是有限阶加厚；
\item $X = S \times_{S'} X'$ 且 $Y = S \times_{S'} Y'$。
\end{enumerate}
则 $f'$ 平坦，并且 $Y'$ 在 $f'$ 的像中的每一点处都在 $S'$ 上平坦。
\end{lemma}

\begin{proof}
这是 代数章引理
\ref{algebra-lemma-criterion-flatness-fibre-nilpotent} 的直接推论。
\end{proof}

\noindent
概形态射的许多性质在平坦形变下保持。

\begin{lemma}
\label{more-morphisms-lemma-deform-property}
考虑交换图
$$
\xymatrix{
(X \subset X') \ar[rr]_{(f, f')} \ar[rd] & & (Y \subset Y') \ar[ld] \\
& (S \subset S')
}
$$
，其中各对象为加厚。假设 $S \subset S'$ 是有限阶加厚，$X'$ 在 $S'$ 上平坦，
$X = S \times_{S'} X'$，且 $Y = S \times_{S'} Y'$。则
\begin{enumerate}
\item $f$ 平坦，当且仅当 $f'$ 平坦；
\label{more-morphisms-item-flat}
\item $f$ 是同构，当且仅当 $f'$ 是同构；
\label{more-morphisms-item-isomorphism}
\item $f$ 是开浸入，当且仅当 $f'$ 是开浸入；
\label{more-morphisms-item-open-immersion}
\item $f$ 拟紧，当且仅当 $f'$ 拟紧；
\label{more-morphisms-item-quasi-compact}
\item $f$ 泛闭，当且仅当 $f'$ 泛闭；
\label{more-morphisms-item-universally-closed}
\item $f$（拟）分离，当且仅当 $f'$（拟）分离；
\label{more-morphisms-item-separated}
\item $f$ 是单态射，当且仅当 $f'$ 是单态射；
\label{more-morphisms-item-monomorphism}
\item $f$ 是满射，当且仅当 $f'$ 是满射；
\label{more-morphisms-item-surjective}
\item $f$ 泛单射，当且仅当 $f'$ 泛单射；
\label{more-morphisms-item-universally-injective}
\item $f$ 是仿射态射，当且仅当 $f'$ 是仿射态射；
\label{more-morphisms-item-affine}
\item
\label{more-morphisms-item-finite-type}
$f$ 局部有限型，当且仅当 $f'$ 局部有限型；
\item $f$ 局部拟有限，当且仅当 $f'$ 局部拟有限；
\label{more-morphisms-item-quasi-finite}
\item
\label{more-morphisms-item-finite-presentation}
$f$ 局部有限表示，当且仅当 $f'$ 局部有限表示；
\item
\label{more-morphisms-item-relative-dimension-d}
$f$ 局部有限型且相对维数为 $d$，当且仅当 $f'$ 局部有限型且相对维数为 $d$；
\item $f$ 泛开，当且仅当 $f'$ 泛开；
\label{more-morphisms-item-universally-open}
\item $f$ 是 syntomic 态射，当且仅当 $f'$ 是 syntomic 态射；
\label{more-morphisms-item-syntomic}
\item $f$ 光滑，当且仅当 $f'$ 光滑；
\label{more-morphisms-item-smooth}
\item $f$ 非分歧，当且仅当 $f'$ 非分歧；
\label{more-morphisms-item-unramified}
\item $f$ 平展，当且仅当 $f'$ 平展；
\label{more-morphisms-item-etale}
\item $f$ 固有，当且仅当 $f'$ 固有；
\label{more-morphisms-item-proper}
\item $f$ 是整态射，当且仅当 $f'$ 是整态射；
\label{more-morphisms-item-integral}
\item $f$ 是有限态射，当且仅当 $f'$ 是有限态射；
\label{more-morphisms-item-finite}
\item
\label{more-morphisms-item-finite-locally-free}
$f$ 是有限局部自由态射（秩为 $d$），当且仅当 $f'$ 是有限局部自由态射
（秩为 $d$）；
\item 此处还应添加更多性质。
\end{enumerate}
\end{lemma}

\begin{proof}
关于 $X$ 与 $Y$ 的假设意味着，$f$ 是 $f'$ 经 $X \to X'$ 的基变换。
上述 (1)--(23) 中列出的性质 $\mathcal{P}$ 在基变换下都保持，
因此若 $f'$ 具有性质 $\mathcal{P}$，则 $f$ 也具有。见
概形章引理 \ref{schemes-lemma-base-change-immersion},
\ref{schemes-lemma-quasi-compact-preserved-base-change},
\ref{schemes-lemma-separated-permanence}, 及
\ref{schemes-lemma-base-change-monomorphism}
及
态射章引理
\ref{morphisms-lemma-base-change-surjective},
\ref{morphisms-lemma-base-change-universally-injective},
\ref{morphisms-lemma-base-change-affine},
\ref{morphisms-lemma-base-change-finite-type},
\ref{morphisms-lemma-base-change-quasi-finite},
\ref{morphisms-lemma-base-change-finite-presentation},
\ref{morphisms-lemma-base-change-relative-dimension-d},
\ref{morphisms-lemma-base-change-syntomic},
\ref{morphisms-lemma-base-change-smooth},
\ref{morphisms-lemma-base-change-unramified},
\ref{morphisms-lemma-base-change-etale},
\ref{morphisms-lemma-base-change-proper},
\ref{morphisms-lemma-base-change-finite}, 及
\ref{morphisms-lemma-base-change-finite-locally-free}.

\medskip\noindent
因此每种情形中有趣的方向都是：假设 $f$ 具有该性质，并推出 $f'$ 也具有。
对加厚的阶数作归纳，可假设 $S \subset S'$ 是一阶加厚；见紧接定义
\ref{more-morphisms-definition-thickening} 之后的讨论。先作几条一般性说明，
以下论证中将不再逐次提及。
(I) 设 $W' \subset S'$ 是仿射开集，并设 $U' \subset X'$ 与
$V' \subset Y'$ 是位于 $W'$ 上方、满足 $f'(U') \subset V'$ 的仿射开集。
令 $W' = \Spec(R')$，并以 $I \subset R'$ 表示定义闭子概形
$W' \cap S$ 的理想。写 $U' = \Spec(B')$ 且 $V' = \Spec(A')$。
于是得到交换图
$$
\xymatrix{
0 \ar[r] &
IB' \ar[r] &
B' \ar[r] &
B \ar[r] & 0 \\
0 \ar[r] &
IA' \ar[r] \ar[u] &
A' \ar[r] \ar[u] &
A \ar[r] \ar[u] & 0
}
$$
，其行正合。此外 $IB' \cong I \otimes_R B$；见引理
\ref{more-morphisms-lemma-deform} 的证明。
(II) 态射 $X \to X'$ 与 $Y \to Y'$ 是泛同胚。因此映射 $f$ 与 $f'$
（在任意基变换后）的拓扑相同。
(III) 若 $f$ 平坦，则 $f'$ 平坦，且 $Y' \to S'$ 在 $f'$ 的像中的每一点处
都平坦；见引理 \ref{more-morphisms-lemma-flatness-morphism-thickenings}。

\medskip\noindent
关于 (\ref{more-morphisms-item-flat})。这就是一般性说明 (III)。

\medskip\noindent
关于 (\ref{more-morphisms-item-isomorphism})。假设 $f$ 是同构。由 (III)，
$Y' \to S'$ 平坦。选择仿射开集 $V' \subset Y'$，并令
$U' = (f')^{-1}(V')$。于是 $V = Y \cap V'$ 是仿射的，从而
$V \cong f^{-1}(V) = U = Y \times_{Y'} U'$ 是仿射的。由引理
\ref{more-morphisms-lemma-thickening-affine-scheme}，$U'$ 是仿射的。
因此有一般性说明 (I) 中那样的图；又因 $R' \to A'$ 平坦，
$IA \cong I \otimes_R A$。于是
$IB' \cong I \otimes_R B \cong I \otimes_R A \cong IA'$，
且 $A \cong B$。由上图各行的正合性，$A' \cong B'$，即
$U' \cong V'$。故 $f'$ 是同构。

\medskip\noindent
关于 (\ref{more-morphisms-item-open-immersion})。假设 $f$ 是开浸入。
则 $f$ 给出 $X$ 与开子概形 $V \subset Y$ 之间的同构。
令 $V' \subset Y'$ 为底拓扑空间等于 $V$ 的开子概形。于是 $f'$
是从 $X'$ 到 $V'$ 的映射，由 (\ref{more-morphisms-item-isomorphism}) 它是同构。
故 $f'$ 是开浸入。

\medskip\noindent
关于 (\ref{more-morphisms-item-quasi-compact})。由说明 (II) 立即可得。更一般的陈述亦见
引理 \ref{more-morphisms-lemma-thicken-property-morphisms}。

\medskip\noindent
关于 (\ref{more-morphisms-item-universally-closed})。由说明 (II) 立即可得。更一般的陈述亦见
引理 \ref{more-morphisms-lemma-thicken-property-morphisms}。

\medskip\noindent
关于 (\ref{more-morphisms-item-separated})。注意
$X \times_Y X = Y \times_{Y'} (X' \times_{Y'} X')$，故
$X' \times_{Y'} X'$ 是 $X \times_Y X$ 的一个加厚。
映射 $\Delta_{X/Y}$ 与 $\Delta_{X'/Y'}$
的拓扑相符，故得证。更一般的陈述亦见引理
\ref{more-morphisms-lemma-thicken-property-morphisms}。

\medskip\noindent
关于 (\ref{more-morphisms-item-monomorphism})。假设 $f$ 是单态射。考虑对角态射
$\Delta_{X'/Y'} : X' \to X' \times_{Y'} X'$。
$\Delta_{X'/Y'}$ 经 $S \to S'$ 的基变换为 $\Delta_{X/Y}$，
后者按假设是同构。由 (\ref{more-morphisms-item-isomorphism})，
$\Delta_{X'/Y'}$ 是同构。

\medskip\noindent
关于 (\ref{more-morphisms-item-surjective})。这是显然的。更一般的陈述亦见引理
\ref{more-morphisms-lemma-thicken-property-morphisms}。

\medskip\noindent
关于 (\ref{more-morphisms-item-universally-injective})。由说明 (II) 立即可得。
更一般的陈述亦见引理 \ref{more-morphisms-lemma-thicken-property-morphisms}。

\medskip\noindent
关于 (\ref{more-morphisms-item-affine})。假设 $f$ 是仿射态射。选择仿射开集
$V' \subset Y'$ 并令 $U' = (f')^{-1}(V')$。于是
$V = Y \cap V'$ 是仿射的，从而 $U = Y \times_{Y'} U'$ 是仿射的。
由引理 \ref{more-morphisms-lemma-thickening-affine-scheme}，$U'$ 是仿射的。
故 $f'$ 是仿射态射。更一般的陈述亦见引理
\ref{more-morphisms-lemma-thicken-property-morphisms}。

\medskip\noindent
关于 (\ref{more-morphisms-item-finite-type})。由说明 (I)，问题化为证明：若
$A \to B$ 有限型，则 $A' \to B'$ 有限型。设
$x_1, \ldots, x_n \in B'$ 在 $B$ 中的像作为 $A$-代数生成 $B$。
% P05-ERR-CH37-0014: Frozen source repeats the surjectivity to B and appears to omit the intended conclusion for B'; retained pending canon review.
则 $A'[x_1, \ldots, x_n] \to B$ 是满射，因为
$A'[x_1, \ldots, x_n] \to B$ 与
$I \otimes_R A[x_1, \ldots, x_n] \to I \otimes_R B$ 都是满射。
更一般的陈述亦见引理 \ref{more-morphisms-lemma-thicken-property-morphisms-cartesian}。

\medskip\noindent
关于 (\ref{more-morphisms-item-quasi-finite})。由 (\ref{more-morphisms-item-finite-type}) 以及有限型态射的
拟有限性可在纤维上检验这一事实得出；见 态射章引理
\ref{morphisms-lemma-quasi-finite-at-point-characterize}。
更一般的陈述亦见引理 \ref{more-morphisms-lemma-thicken-property-morphisms-cartesian}。

\medskip\noindent
关于 (\ref{more-morphisms-item-finite-presentation})。由说明 (I)，问题化为证明：若
$A \to B$ 有限表示，则 $A' \to B'$ 有限表示。由
(\ref{more-morphisms-item-finite-type})，可假设对某个理想 $K'$ 有
$B' = A'[x_1, \ldots, x_n]/K'$。得到短正合列
$$
0 \to K' \to A'[x_1, \ldots, x_n] \to B' \to 0
$$
由于 $B'$ 在 $R'$ 上平坦，$K' \otimes_{R'} R$ 是满射
$A[x_1, \ldots, x_n] \to B$ 的核。由关于 $A \to B$ 的假设，
存在有限多个 $f'_1, \ldots, f'_m \in K'$，其在
$A[x_1, \ldots, x_n]$ 中的像生成此核。由于 $I$ 幂零，
由 Nakayama 引理可知 $f'_1, \ldots, f'_m$ 生成 $K'$；见
代数章引理 \ref{algebra-lemma-NAK}。

\medskip\noindent
关于 (\ref{more-morphisms-item-relative-dimension-d})。由 (\ref{more-morphisms-item-finite-type})
与一般性说明 (II) 得出。更一般的陈述亦见引理
\ref{more-morphisms-lemma-thicken-property-morphisms-cartesian}。

\medskip\noindent
关于 (\ref{more-morphisms-item-universally-open})。由一般性说明 (II) 立即可得。
更一般的陈述亦见引理 \ref{more-morphisms-lemma-thicken-property-morphisms}。

\medskip\noindent
关于 (\ref{more-morphisms-item-syntomic})。假设 $f$ 是 syntomic 态射。由
(\ref{more-morphisms-item-finite-presentation})，$f'$ 局部有限表示；由一般性说明 (III)，
$f'$ 平坦，且 $f'$ 的纤维就是 $f$ 的纤维。因此由
态射章引理 \ref{morphisms-lemma-syntomic-flat-fibres}，
$f'$ 是 syntomic 态射。

\medskip\noindent
关于 (\ref{more-morphisms-item-smooth})。假设 $f$ 光滑。由
(\ref{more-morphisms-item-finite-presentation})，$f'$ 局部有限表示；由一般性说明 (III)，
$f'$ 平坦，且 $f'$ 的纤维就是 $f$ 的纤维。因此由
态射章引理 \ref{morphisms-lemma-smooth-flat-smooth-fibres}，
$f'$ 光滑。

\medskip\noindent
关于 (\ref{more-morphisms-item-unramified})。假设 $f$ 非分歧。由
(\ref{more-morphisms-item-finite-type})，$f'$ 局部有限型，且 $f'$ 的纤维就是 $f$ 的纤维。
因此由 态射章引理
\ref{morphisms-lemma-unramified-etale-fibres}，$f'$ 非分歧。
更一般的陈述亦见引理 \ref{more-morphisms-lemma-thicken-property-morphisms-cartesian}。

\medskip\noindent
关于 (\ref{more-morphisms-item-etale})。假设 $f$ 平展。由
(\ref{more-morphisms-item-finite-presentation})，$f'$ 局部有限表示；由一般性说明 (III)，
$f'$ 平坦，且 $f'$ 的纤维就是 $f$ 的纤维。因此由
态射章引理 \ref{morphisms-lemma-etale-flat-etale-fibres}，
$f'$ 平展。

\medskip\noindent
关于 (\ref{more-morphisms-item-proper})。合并 (\ref{more-morphisms-item-separated})、
(\ref{more-morphisms-item-finite-type})、(\ref{more-morphisms-item-quasi-compact}) 与
(\ref{more-morphisms-item-universally-closed}) 即得。更一般的陈述亦见引理
\ref{more-morphisms-lemma-thicken-property-morphisms-cartesian}。

\medskip\noindent
关于 (\ref{more-morphisms-item-integral})。将 (\ref{more-morphisms-item-universally-closed})、
(\ref{more-morphisms-item-affine}) 与 态射章引理
\ref{morphisms-lemma-integral-universally-closed} 合并。更一般的陈述亦见引理
\ref{more-morphisms-lemma-thicken-property-morphisms}。

\medskip\noindent
关于 (\ref{more-morphisms-item-finite})。将 (\ref{more-morphisms-item-integral})、
(\ref{more-morphisms-item-finite-type}) 与 态射章引理
\ref{morphisms-lemma-finite-integral} 合并。更一般的陈述亦见引理
\ref{more-morphisms-lemma-thicken-property-morphisms-cartesian}。

\medskip\noindent
关于 (\ref{more-morphisms-item-finite-locally-free})。假设 $f$ 有限局部自由。
由 (\ref{more-morphisms-item-finite})，$f'$ 有限；由一般性说明 (III)，$f'$ 平坦；
由 (\ref{more-morphisms-item-finite-presentation})，$f'$ 局部有限表示。因此由
态射章引理 \ref{morphisms-lemma-finite-flat}，
$f'$ 有限局部自由。
\end{proof}

\noindent
下面的引理是“逐纤维平坦性判据”（``crit\`ere de platitude par fibres''）
的“局部幂零”版本；见节 \ref{more-morphisms-section-criterion-flat-fibres}。

\begin{lemma}
\label{more-morphisms-lemma-flatness-morphism-thickenings-fp-over-ft}
考虑交换图
$$
\xymatrix{
(X \subset X') \ar[rr]_{(f, f')} \ar[rd] & & (Y \subset Y') \ar[ld] \\
& (S \subset S')
}
$$
，其中各对象为加厚。假设
\begin{enumerate}
\item $Y' \to S'$ 局部有限型；
\item $X' \to S'$ 平坦且局部有限表示；
\item $f$ 平坦；
\item $X = S \times_{S'} X'$ 且 $Y = S \times_{S'} Y'$。
\end{enumerate}
则 $f'$ 平坦，并且对 $f'$ 的像中的每个 $y' \in Y'$，局部环
$\mathcal{O}_{Y', y'}$ 在 $\mathcal{O}_{S', s'}$ 上平坦且本质有限表示。
\end{lemma}

\begin{proof}
这是 代数章引理
\ref{algebra-lemma-criterion-flatness-fibre-locally-nilpotent} 的直接推论。
\end{proof}

\noindent
概形态射的许多性质在上述引理那样的平坦形变下保持。

\begin{lemma}
\label{more-morphisms-lemma-deform-property-fp-over-ft}
考虑交换图
$$
\xymatrix{
(X \subset X') \ar[rr]_{(f, f')} \ar[rd] & & (Y \subset Y') \ar[ld] \\
& (S \subset S')
}
$$
，其中各对象为加厚。假设 $Y' \to S'$ 局部有限型，
$X' \to S'$ 平坦且局部有限表示，$X = S \times_{S'} X'$，
且 $Y = S \times_{S'} Y'$。则
\begin{enumerate}
\item $f$ 平坦，当且仅当 $f'$ 平坦；
\label{more-morphisms-item-flat-fp-over-ft}
\item $f$ 是同构，当且仅当 $f'$ 是同构；
\label{more-morphisms-item-isomorphism-fp-over-ft}
\item $f$ 是开浸入，当且仅当 $f'$ 是开浸入；
\label{more-morphisms-item-open-immersion-fp-over-ft}
\item $f$ 拟紧，当且仅当 $f'$ 拟紧；
\label{more-morphisms-item-quasi-compact-fp-over-ft}
\item $f$ 泛闭，当且仅当 $f'$ 泛闭；
\label{more-morphisms-item-universally-closed-fp-over-ft}
\item $f$（拟）分离，当且仅当 $f'$（拟）分离；
\label{more-morphisms-item-separated-fp-over-ft}
\item $f$ 是单态射，当且仅当 $f'$ 是单态射；
\label{more-morphisms-item-monomorphism-fp-over-ft}
\item $f$ 是满射，当且仅当 $f'$ 是满射；
\label{more-morphisms-item-surjective-fp-over-ft}
\item $f$ 泛单射，当且仅当 $f'$ 泛单射；
\label{more-morphisms-item-universally-injective-fp-over-ft}
\item $f$ 是仿射态射，当且仅当 $f'$ 是仿射态射；
\label{more-morphisms-item-affine-fp-over-ft}
\item $f$ 局部拟有限，当且仅当 $f'$ 局部拟有限；
\label{more-morphisms-item-quasi-finite-fp-over-ft}
\item
\label{more-morphisms-item-relative-dimension-d-fp-over-ft}
$f$ 局部有限型且相对维数为 $d$，当且仅当 $f'$ 局部有限型且相对维数为 $d$；
\item $f$ 泛开，当且仅当 $f'$ 泛开；
\label{more-morphisms-item-universally-open-fp-over-ft}
\item $f$ 是 syntomic 态射，当且仅当 $f'$ 是 syntomic 态射；
\label{more-morphisms-item-syntomic-fp-over-ft}
\item $f$ 光滑，当且仅当 $f'$ 光滑；
\label{more-morphisms-item-smooth-fp-over-ft}
\item $f$ 非分歧，当且仅当 $f'$ 非分歧；
\label{more-morphisms-item-unramified-fp-over-ft}
\item $f$ 平展，当且仅当 $f'$ 平展；
\label{more-morphisms-item-etale-fp-over-ft}
\item $f$ 固有，当且仅当 $f'$ 固有；
\label{more-morphisms-item-proper-fp-over-ft}
\item $f$ 是有限态射，当且仅当 $f'$ 是有限态射；
\label{more-morphisms-item-finite-fp-over-ft}
\item
\label{more-morphisms-item-finite-locally-free-fp-over-ft}
$f$ 是有限局部自由态射（秩为 $d$），当且仅当 $f'$ 是有限局部自由态射
（秩为 $d$）；
\item 此处还应添加更多性质。
\end{enumerate}
\end{lemma}

\begin{proof}
关于 $X$ 与 $Y$ 的假设意味着，$f$ 是 $f'$ 经 $X \to X'$ 的基变换。
上述 (1)--(20) 中列出的性质 $\mathcal{P}$ 在基变换下都保持，
因此若 $f'$ 具有性质 $\mathcal{P}$，则 $f$ 也具有。见
概形章引理 \ref{schemes-lemma-base-change-immersion},
\ref{schemes-lemma-quasi-compact-preserved-base-change},
\ref{schemes-lemma-separated-permanence}, 及
\ref{schemes-lemma-base-change-monomorphism}
及
态射章引理
\ref{morphisms-lemma-base-change-surjective},
\ref{morphisms-lemma-base-change-universally-injective},
\ref{morphisms-lemma-base-change-affine},
\ref{morphisms-lemma-base-change-quasi-finite},
\ref{morphisms-lemma-base-change-relative-dimension-d},
\ref{morphisms-lemma-base-change-syntomic},
\ref{morphisms-lemma-base-change-smooth},
\ref{morphisms-lemma-base-change-unramified},
\ref{morphisms-lemma-base-change-etale},
\ref{morphisms-lemma-base-change-proper},
\ref{morphisms-lemma-base-change-finite}, 及
\ref{morphisms-lemma-base-change-finite-locally-free}.

\medskip\noindent
因此每种情形中有趣的方向都是：假设 $f$ 具有该性质，并推出 $f'$ 也具有。
先作几条一般性说明，以下论证中将不再逐次提及。
(I) 设 $W' \subset S'$ 是仿射开集，并设 $U' \subset X'$ 与
$V' \subset Y'$ 是位于 $W'$ 上方、满足 $f'(U') \subset V'$ 的仿射开集。
令 $W' = \Spec(R')$，并以 $I \subset R'$ 表示定义闭子概形
$W' \cap S$ 的理想。写 $U' = \Spec(B')$ 且 $V' = \Spec(A')$。
于是得到交换图
$$
\xymatrix{
0 \ar[r] &
IB' \ar[r] &
B' \ar[r] &
B \ar[r] & 0 \\
0 \ar[r] &
IA' \ar[r] \ar[u] &
A' \ar[r] \ar[u] &
A \ar[r] \ar[u] & 0
}
$$
，其行正合。
(II) 态射 $X \to X'$ 与 $Y \to Y'$ 是泛同胚。因此映射 $f$ 与 $f'$
（在任意基变换后）的拓扑相同。
(III) 若 $f$ 平坦，则 $f'$ 平坦，且 $Y' \to S'$ 在 $f'$ 的像中的每一点处
都平坦；见引理 \ref{more-morphisms-lemma-flatness-morphism-thickenings}。

\medskip\noindent
关于 (\ref{more-morphisms-item-flat-fp-over-ft})。这就是一般性说明 (III)。

\medskip\noindent
关于 (\ref{more-morphisms-item-isomorphism-fp-over-ft})。假设 $f$ 是同构。
选择仿射开集 $V' \subset Y'$，并令 $U' = (f')^{-1}(V')$。
于是 $V = Y \cap V'$ 是仿射的，从而
$V \cong f^{-1}(V) = U = Y \times_{Y'} U'$ 是仿射的。由引理
\ref{more-morphisms-lemma-thickening-affine-scheme}，$U'$ 是仿射的。
因此有一般性说明 (I) 中那样的图。由 代数章引理
\ref{algebra-lemma-isomorphism-modulo-locally-nilpotent}，
$A' \to B'$ 是同构，即 $U' \cong V'$。故 $f'$ 是同构。

\medskip\noindent
关于 (\ref{more-morphisms-item-open-immersion-fp-over-ft})。假设 $f$ 是开浸入。
则 $f$ 给出 $X$ 与开子概形 $V \subset Y$ 之间的同构。
令 $V' \subset Y'$ 为底拓扑空间等于 $V$ 的开子概形。于是 $f'$
是从 $X'$ 到 $V'$ 的映射，由 (\ref{more-morphisms-item-isomorphism-fp-over-ft})
它是同构。故 $f'$ 是开浸入。

\medskip\noindent
关于 (\ref{more-morphisms-item-quasi-compact-fp-over-ft})。由说明 (II) 立即可得。
更一般的陈述亦见引理 \ref{more-morphisms-lemma-thicken-property-morphisms}。

\medskip\noindent
关于 (\ref{more-morphisms-item-universally-closed-fp-over-ft})。由说明 (II) 立即可得。
更一般的陈述亦见引理 \ref{more-morphisms-lemma-thicken-property-morphisms}。

\medskip\noindent
关于 (\ref{more-morphisms-item-separated-fp-over-ft})。注意
$X \times_Y X = Y \times_{Y'} (X' \times_{Y'} X')$，故
$X' \times_{Y'} X'$ 是 $X \times_Y X$ 的一个加厚。
映射 $\Delta_{X/Y}$ 与 $\Delta_{X'/Y'}$
的拓扑相符，故得证。更一般的陈述亦见引理
\ref{more-morphisms-lemma-thicken-property-morphisms}。

\medskip\noindent
关于 (\ref{more-morphisms-item-monomorphism-fp-over-ft})。假设 $f$ 是单态射。
考虑对角态射 $\Delta_{X'/Y'} : X' \to X' \times_{Y'} X'$。
注意 $X' \times_{Y'} X' \to S'$ 局部有限型。
$\Delta_{X'/Y'}$ 经 $S \to S'$ 的基变换为 $\Delta_{X/Y}$，
后者按假设是同构。由 (\ref{more-morphisms-item-isomorphism-fp-over-ft})，
$\Delta_{X'/Y'}$ 是同构。

\medskip\noindent
关于 (\ref{more-morphisms-item-surjective-fp-over-ft})。这是显然的。更一般的陈述亦见引理
\ref{more-morphisms-lemma-thicken-property-morphisms}。

\medskip\noindent
关于 (\ref{more-morphisms-item-universally-injective-fp-over-ft})。由说明 (II) 立即可得。
更一般的陈述亦见引理 \ref{more-morphisms-lemma-thicken-property-morphisms}。

\medskip\noindent
关于 (\ref{more-morphisms-item-affine-fp-over-ft})。假设 $f$ 是仿射态射。
选择仿射开集 $V' \subset Y'$，并令 $U' = (f')^{-1}(V')$。
于是 $V = Y \cap V'$ 是仿射的，从而 $U = Y \times_{Y'} U'$ 是仿射的。
由引理 \ref{more-morphisms-lemma-thickening-affine-scheme}，$U'$ 是仿射的。
故 $f'$ 是仿射态射。更一般的陈述亦见引理
\ref{more-morphisms-lemma-thicken-property-morphisms}。

\medskip\noindent
关于 (\ref{more-morphisms-item-quasi-finite-fp-over-ft})。这是因为 $f'$ 局部有限型
（由 态射章引理 \ref{morphisms-lemma-permanence-finite-type}），
而有限型态射的拟有限性可在纤维上检验；见 态射章引理
\ref{morphisms-lemma-quasi-finite-at-point-characterize}。

\medskip\noindent
关于 (\ref{more-morphisms-item-relative-dimension-d-fp-over-ft})。由一般性说明 (II)
以及 $f'$ 局部有限型这一事实得出（态射章引理
\ref{morphisms-lemma-permanence-finite-type}）。

\medskip\noindent
关于 (\ref{more-morphisms-item-universally-open-fp-over-ft})。由一般性说明 (II) 立即可得。
更一般的陈述亦见引理 \ref{more-morphisms-lemma-thicken-property-morphisms}。

\medskip\noindent
关于 (\ref{more-morphisms-item-syntomic-fp-over-ft})。假设 $f$ 是 syntomic 态射。
由 态射章引理 \ref{morphisms-lemma-finite-presentation-permanence}，
$f'$ 局部有限表示。由一般性说明 (III)，$f'$ 平坦。
$f'$ 的纤维就是 $f$ 的纤维。因此由 态射章引理
\ref{morphisms-lemma-syntomic-flat-fibres}，$f'$ 是 syntomic 态射。

\medskip\noindent
关于 (\ref{more-morphisms-item-smooth-fp-over-ft})。假设 $f$ 光滑。由
态射章引理 \ref{morphisms-lemma-finite-presentation-permanence}，
$f'$ 局部有限表示。由一般性说明 (III)，$f'$ 平坦。
$f'$ 的纤维就是 $f$ 的纤维。因此由 态射章引理
\ref{morphisms-lemma-smooth-flat-smooth-fibres}，$f'$ 光滑。

\medskip\noindent
关于 (\ref{more-morphisms-item-unramified-fp-over-ft})。假设 $f$ 非分歧。由
态射章引理 \ref{morphisms-lemma-permanence-finite-type}，
$f'$ 局部有限型。$f'$ 的纤维就是 $f$ 的纤维。因此由
态射章引理 \ref{morphisms-lemma-unramified-etale-fibres}，
$f'$ 非分歧。

\medskip\noindent
关于 (\ref{more-morphisms-item-etale-fp-over-ft})。假设 $f$ 平展。由
态射章引理 \ref{morphisms-lemma-finite-presentation-permanence}，
$f'$ 局部有限表示。由一般性说明 (III)，$f'$ 平坦。
$f'$ 的纤维就是 $f$ 的纤维。因此由 态射章引理
\ref{morphisms-lemma-etale-flat-etale-fibres}，$f'$ 平展。

\medskip\noindent
关于 (\ref{more-morphisms-item-proper-fp-over-ft})。合并
(\ref{more-morphisms-item-separated-fp-over-ft})、$f$ 局部有限型这一事实
（态射章引理 \ref{morphisms-lemma-permanence-finite-type}）、
(\ref{more-morphisms-item-quasi-compact-fp-over-ft}) 与
(\ref{more-morphisms-item-universally-closed-fp-over-ft}) 即得。

\medskip\noindent
关于 (\ref{more-morphisms-item-finite-fp-over-ft})。合并
(\ref{more-morphisms-item-universally-closed-fp-over-ft})、
(\ref{more-morphisms-item-affine-fp-over-ft})、态射章引理
\ref{morphisms-lemma-integral-universally-closed}、$f$ 局部有限型这一事实
（态射章引理 \ref{morphisms-lemma-permanence-finite-type}）以及
态射章引理 \ref{morphisms-lemma-finite-integral}。

\medskip\noindent
关于 (\ref{more-morphisms-item-finite-locally-free-fp-over-ft})。
假设 $f$ 有限局部自由。由 (\ref{more-morphisms-item-finite-fp-over-ft})，$f'$ 有限。
由一般性说明 (III)，$f'$ 平坦。由 态射章引理
\ref{morphisms-lemma-finite-presentation-permanence}，$f'$ 局部有限表示。
因此由 态射章引理 \ref{morphisms-lemma-finite-flat}，
$f'$ 有限局部自由。
\end{proof}

\begin{lemma}[射影概形的形变]
\label{more-morphisms-lemma-deform-projective}
设 $f : X \to S$ 是固有、平坦且有限表示的概形态射，
$\mathcal{L}$ 是 $f$-丰沛的，并假设 $S$ 拟紧。则存在 $d_0 \geq 0$，
使得对每个笛卡尔图
$$
\vcenter{
\xymatrix{
X \ar[r]_{i'} \ar[d]_f & X' \ar[d]^{f'} \\
S \ar[r]^i & S'
}
}
\quad\text{且}\quad
\begin{matrix}
\text{invertible }\mathcal{O}_{X'}\text{-模}\\
\mathcal{L}'\text{ 其中 }\mathcal{L} \cong (i')^*\mathcal{L}'
\end{matrix}
$$
，其中 $S \subset S'$ 是加厚且 $f'$ 固有、平坦、有限表示，都有
\begin{enumerate}
\item $R^p(f')_*(\mathcal{L}')^{\otimes d} = 0$
对所有 $p > 0$ 与 $d \geq d_0$ 成立；
\item $\mathcal{A}'_d = (f')_*(\mathcal{L}')^{\otimes d}$
在 $d \geq d_0$ 时有限局部自由；
\item $\mathcal{A}' =
\mathcal{O}_{S'} \oplus \bigoplus_{d \geq d_0} \mathcal{A}'_d$
是有限表示的拟凝聚 $\mathcal{O}_{S'}$-代数；
\item 存在典范同构
$r' : X' \to \underline{\text{Proj}}_{S'}(\mathcal{A}')$；
\item 存在典范同构
$\theta' : (r')^*\mathcal{O}_{\underline{\text{Proj}}_{S'}(\mathcal{A}')}(1)
\to \mathcal{L}'$.
\end{enumerate}
$\mathcal{A}'$、$r'$、$\theta'$ 的构造关于数据
$(X', S', i, i', f', \mathcal{L}')$ 是函子性的。
\end{lemma}

\begin{proof}
先描述映射 $r'$ 与 $\theta'$。注意 $\mathcal{L}'$ 是 $f'$-丰沛的；
见引理 \ref{more-morphisms-lemma-thicken-property-relatively-ample}。
有拟凝聚分次 $\mathcal{O}_{S'}$-代数的典范映射
$\mathcal{A}' \to \bigoplus_{d \geq 0} (f')_*(\mathcal{L}')^{\otimes d}$
它在次数 $\geq d_0$ 上为同构。因此，它在相对 Proj 上诱导同构，
并与结构层的 Serre 扭曲相容；见 构造章引理
\ref{constructions-lemma-eventual-iso-graded-rings-map-relative-proj}。
故由 态射章引理
\ref{morphisms-lemma-characterize-relatively-ample} 得到态射 $r'$
（该引理又援用了 构造章引理
\ref{constructions-lemma-invertible-map-into-relative-proj} 的构造），
而由 态射章引理 \ref{morphisms-lemma-proper-ample-is-proj}，
它是同构。映射 $\theta'$ 来自 构造章引理
\ref{constructions-lemma-invertible-map-into-relative-proj}。
由 性质章引理 \ref{properties-lemma-ample-gcd-is-one}，
$\theta'$ 是同构（这里还使用了如下事实：在 $S'$ 的各仿射开集上，
态射 $r'$ 与 性质章引理 \ref{properties-lemma-map-into-proj}
中的态射相同；这已在 态射章引理
\ref{morphisms-lemma-proper-ample-is-proj} 的证明中说明）。

\medskip\noindent
假设 (1)、(2) 中所述的消失与局部自由性成立，则该构造的函子性可如下看出。
设 $h : T \to S'$ 是概形态射，以 $f_T : X'_T \to T$ 表示 $f'$ 的基变换，
并以 $\mathcal{L}_T$ 表示 $\mathcal{L}$ 到 $X'_T$ 的拉回。
% P05-ERR-CH37-0015: Frozen source says pullback of L although the deformation datum uses L'; retained pending canon review.
由上同调与基变换（例如 概形的导出范畴章引理
\ref{perfect-lemma-compare-base-change} 的表述），在 $T$ 上有相应的消失，并且
$h^*\mathcal{A}'_d = f_{T, *}\mathcal{L}_T^{\otimes d}$
（因而也有推出层的局部自由性以及相应分次
$\mathcal{O}_T$-代数 $\mathcal{A}_T$ 的有限生成性）。因此态射
$r_T : X_T \to
\underline{\text{Proj}}_T(\bigoplus f_{T, *}\mathcal{L}_T^{\otimes d})$
% P05-ERR-CH37-0016: Frozen source writes X_T after introducing X'_T; formula retained pending canon review.
就是 $r'$ 到 $T$ 的基变换，而 $\theta'$ 的拉回就是映射 $\theta_T$。

\medskip\noindent
综上，只需证明 (1)、(2)、(3)。选择 $d_0$，使对所有
$d \geq d_0$ 与 $p > 0$ 都有
$R^pf_*\mathcal{L}^{\otimes d} = 0$；见 概形上同调章引理
\ref{coherent-lemma-coherent-proper-ample}。我们断言此 $d_0$ 可用。

\medskip\noindent
由上同调与基变换（概形的导出范畴章引理
\ref{perfect-lemma-flat-proper-perfect-direct-image-general}），
$E'_d = Rf'_*(\mathcal{L}')^{\otimes d}$ 是
$D(\mathcal{O}_{S'})$ 的完美对象，且其构造与任意基变换可交换。
特别地，$E_d = Li^*E'_d = Rf_*\mathcal{L}^{\otimes d}$。
由 概形的导出范畴章引理
\ref{perfect-lemma-vanishing-implies-locally-free}，当 $d \geq d_0$ 时，
复形 $E_d$ 同构于置于上同调次数 $0$ 的有限局部自由
$\mathcal{O}_S$-模 $f_*\mathcal{L}^{\otimes d}$。再由
概形的导出范畴章引理
\ref{perfect-lemma-open-where-cohomology-in-degree-i-rank-r}，
$E'_d$ 同构于置于上同调次数 $0$ 的有限局部自由模。
这当然表示 $E'_d = \mathcal{A}'_d[0]$，对 $p > 0$ 有
$R^pf'_*(\mathcal{L}')^{\otimes d} = 0$，且 $\mathcal{A}'_d$
有限局部自由。这证明了 (1) 与 (2)。

\medskip\noindent
最后还须证明，$\mathcal{A}'$ 作为 $\mathcal{O}_{S'}$-代数层有限表示
（此概念在 Properties, Section
\ref{properties-section-extending-quasi-coherent-sheaves} 中引入）。
设 $U' = \Spec(R') \subset S'$ 是仿射开集。则
$A' = \mathcal{A}'(U')$ 是分次 $R'$-代数，其各分次部分都是有限投射
$R'$-模。须证明 $A'$ 是有限表示的 $R'$-代数。
我们将把它归结到 Noether 情形。具体而言，可以找到有限型
$\mathbf{Z}$-子代数 $R'_0 \subset R'$ 以及 $R'_0$ 上的一对
\footnote{它具有与 $X' \to S'$ 及 $\mathcal{L}'$ 相同的性质；即
$X'_0 \to \Spec(R'_0)$ 平坦且固有，而 $\mathcal{L}'_0$ 丰沛。}
$(X'_0, \mathcal{L}'_0)$，其基变换为
$(X'_{U'}, \mathcal{L}'|_{X'_{U'}})$；见
极限章引理
\ref{limits-lemma-descend-modules-finite-presentation},
\ref{limits-lemma-descend-invertible-modules},
\ref{limits-lemma-eventually-proper},
\ref{limits-lemma-descend-flat-finite-presentation}, 及
\ref{limits-lemma-limit-ample}.
概形上同调章引理 \ref{coherent-lemma-coherent-proper-ample}
推出
$A'_0 = \bigoplus_{d \geq 0} H^0(X'_0, (\mathcal{L}'_0)^{\otimes d})$
是有限生成的分次 $R'_0$-代数，并推出存在 $d'_0$，使得
$H^p(X'_0, (\mathcal{L}'_0)^{\otimes d}) = 0$, $p > 0$ for $d \geq d'_0$.
。将上述论证应用于 $X'_0 \to \Spec(R'_0)$ 与 $\mathcal{L}'_0$，
可知 $(A'_0)_d$ 是有限投射 $R'_0$-模，并且
$$
A'_d = \mathcal{A}'_d(U') =
H^0(X'_{U'}, (\mathcal{L}')^{\otimes d}|_{X'_{U'}}) =
H^0(X'_0, (\mathcal{L}'_0)^{\otimes d}) \otimes_{R'_0} R' =
(A'_0)_d \otimes_{R'_0} R'
$$
对 $d \geq d'_0$ 成立。论证中的一个小转折是：我们不知道能否选择
$d'_0$ 等于 $d_0$\footnote{事实上，通过更多极限论证可以归结到这一情形。}。
为绕开这一点，用以下一系列论证完成证明：
\begin{enumerate}
\item[(a)] 代数
$B = R'_0 \oplus \bigoplus_{d \geq \max(d_0, d'_0)} (A'_0)_d$
是有限型 $R'_0$-代数：对 $B \subset A'_0$ 应用 Artin--Tate 引理；见
代数章引理 \ref{algebra-lemma-Artin-Tate}。
\item[(b)] 由于 $R'_0$ 是 Noether 环，$B$ 是有限表示的 $R'_0$-代数。
\item[(c)] 由张量积的右正合性，$B \otimes_{R'_0} R'$
是有限表示的 $R'$-代数。
\item[(d)] 由上述等式，这恰好表示
$C = R' \oplus \bigoplus_{d \geq \max(d_0, d'_0)} A'_d$
是有限表示的 $R'$-代数。
\item[(e)] 商 $A'/C$ 是有限投射 $R'$-模 $A'_d$（其中
% P05-ERR-CH37-0017: Frozen source includes the upper endpoint although C already contains that degree; retained pending canon review.
$d_0 \leq d \leq \max(d_0, d'_0)$）的直和，故作为 $R'$-模有限表示。
\item[(f)] 由 代数章引理
\ref{algebra-lemma-finitely-presented-over-subring}，商 $A'/C$
作为 $C$-模有限表示。
\item[(g)] 因此由 代数章引理 \ref{algebra-lemma-extension}，
$A'$ 作为 $C$-模有限表示。
\item[(h)] 由 代数章引理 \ref{algebra-lemma-finite-finite-type}，
这推出 $A'$ 作为 $C$-代数有限表示。
\item[(i)] 最后，将 代数章引理
\ref{algebra-lemma-compose-finite-type} 应用于 $R' \to C \to A'$，
这推出 $A'$ 作为 $R'$-代数有限表示。
\end{enumerate}
证明完毕。
\end{proof}









\section{形式光滑态射}
\label{more-morphisms-section-formally-smooth}

\noindent
Michael Artin 对光滑性的微分判据（例如 态射章引理
\ref{morphisms-lemma-smooth-at-point}）的看法是，它们（在实践中）基本无用。
本节引入形式光滑态射 $X \to S$ 的概念。这种态射由如下性质刻画：
只要 $T$ 仿射，$X$ 的 $T$-值点就能提升到 $T$ 的无穷小加厚上。
主要结果是：形式光滑且局部有限表示的态射是光滑的；见引理
\ref{more-morphisms-lemma-smooth-formally-smooth}。事实证明，这一判据往往比上述微分判据
更容易使用。

\medskip\noindent
回顾：若对每个实线部分交换的图
$$
\xymatrix{
A \ar[r] \ar@{-->}[rd] & B/I \\
R \ar[r] \ar[u] & B \ar[u]
}
$$
（其中 $I \subset B$ 是平方为零的理想），存在一条使该图交换的虚线箭头，
则称环映射 $R \to A$ 是{\it 形式光滑的}；见
代数章定义 \ref{algebra-definition-formally-smooth}。
这引出概形态射的如下类似概念。

\begin{definition}
\label{more-morphisms-definition-formally-smooth}
设 $f : X \to S$ 是概形态射。若给定任意实线部分交换的图
$$
\xymatrix{
X \ar[d]_f & T \ar[d]^i \ar[l] \\
S & T' \ar[l] \ar@{-->}[lu]
}
$$
其中 $T \subset T'$ 是 $S$ 上仿射概形的一阶加厚时，存在一条使该图交换的
虚线箭头，则称 $f$ 是{\it 形式光滑的}。
\end{definition}

\noindent
对形式非分歧与形式平展态射，可以去掉 $T'$ 为仿射的条件；见引理
\ref{more-morphisms-lemma-formally-unramified-not-affine} 与
\ref{more-morphisms-lemma-formally-etale-not-affine}。对形式光滑态射则不再如此。
事实上，一个稍微更自然的条件是：应能在 $T'$ 上 Zariski 局部地补出虚线箭头。
分析引理 \ref{more-morphisms-lemma-formally-smooth} 的证明可知，这与现有定义等价。
特别地，形式光滑性在源上是 Zariski 局部的（事实上它在源上是光滑局部的，
% P05-ERR-CH37-0018: Frozen source contains the literal placeholder ``insert future reference here''.
此处插入将来的参考文献）。

\begin{lemma}
\label{more-morphisms-lemma-composition-formally-smooth}
形式光滑态射的复合仍是形式光滑的。
\end{lemma}

\begin{proof}
从略。
\end{proof}

\begin{lemma}
\label{more-morphisms-lemma-base-change-formally-smooth}
形式光滑态射的基变换仍是形式光滑的。
\end{lemma}

\begin{proof}
从略；其代数版本见 代数章引理 \ref{algebra-lemma-base-change-fs}。
\end{proof}

\begin{lemma}
\label{more-morphisms-lemma-formally-etale-unramified-smooth}
设 $f : X \to S$ 是概形态射。则 $f$ 形式平展，当且仅当
$f$ 形式光滑且形式非分歧。
\end{lemma}

\begin{proof}
从略。
\end{proof}

\begin{lemma}
\label{more-morphisms-lemma-formally-smooth-on-opens}
设 $f : X \to S$ 是概形态射，并设 $U \subset X$ 与 $V \subset S$
是满足 $f(U) \subset V$ 的开子概形。若 $f$ 形式光滑，
则 $f|_U : U \to V$ 也形式光滑。
\end{lemma}

\begin{proof}
考虑定义 \ref{more-morphisms-definition-formally-smooth} 中那样的实线部分交换图
$$
\xymatrix{
U \ar[d]_{f|_U} & T \ar[d]^i \ar[l]^a \\
V & T' \ar[l] \ar@{-->}[lu]
}
$$
。若 $f$ 形式光滑，则存在满足 $a'|_T = a$ 的 $S$-态射
$a' : T' \to X$。由于 $T$ 与 $T'$ 的底集合相同，$a'$ 是到 $U$ 的态射
（见 概形章第 \ref{schemes-section-open-immersion} 节），
而且它显然也是 $V$-态射。因此得到所需的上述虚线箭头。
\end{proof}

\begin{lemma}
\label{more-morphisms-lemma-affine-formally-smooth}
设 $f : X \to S$ 是概形态射，并假设 $X$ 与 $S$ 均为仿射。
则 $f$ 形式光滑，当且仅当
$\mathcal{O}_S(S) \to \mathcal{O}_X(X)$ 是形式光滑环映射。
\end{lemma}

\begin{proof}
由环范畴与仿射概形范畴之间的等价，这立即从定义
（定义 \ref{more-morphisms-definition-formally-smooth} 与
代数章定义 \ref{algebra-definition-formally-smooth}）得出；见
概形章引理 \ref{schemes-lemma-category-affine-schemes}。
\end{proof}

\noindent
下面的引理是本节的主要结果。它是函子观点的一次胜利：与 Limits,
命题 \ref{limits-proposition-characterize-locally-finite-presentation}
合并后，它表明可用函子 $h_X : \Sch/S \to \textit{Sets}$ 的“简单”性质
来识别态射 $f : X \to S$ 是否光滑。

\begin{lemma}[无穷小提升判据]
\label{more-morphisms-lemma-smooth-formally-smooth}
设 $f : X \to S$ 是概形态射。以下条件等价：
\begin{enumerate}
\item 态射 $f$ 光滑；
\item 态射 $f$ 局部有限表示且形式光滑。
\end{enumerate}
\end{lemma}

\begin{proof}
假设 $f : X \to S$ 局部有限表示且形式光滑。考虑满足
$f(U) \subset V$ 的一对仿射开集 $\Spec(A) = U \subset X$ 与
$\Spec(R) = V \subset S$。由引理 \ref{more-morphisms-lemma-formally-smooth-on-opens}，
$U \to V$ 形式光滑。由引理 \ref{more-morphisms-lemma-affine-formally-smooth}，
$R \to A$ 形式光滑。由 态射章引理
\ref{morphisms-lemma-locally-finite-presentation-characterize}，
$R \to A$ 有限表示。由 代数章命题
\ref{algebra-proposition-smooth-formally-smooth}，$R \to A$ 光滑。
故由光滑态射的定义，$X \to S$ 光滑。

\medskip\noindent
反之，假设 $f : X \to S$ 光滑。考虑实线部分交换图
$$
\xymatrix{
X \ar[d]_f & T \ar[d]^i \ar[l]^a \\
S & T' \ar[l] \ar@{-->}[lu]
}
$$
，如定义 \ref{more-morphisms-definition-formally-smooth} 所述。我们将证明虚线箭头存在，
从而证明 $f$ 形式光滑。

\medskip\noindent
令 $\mathcal{F}$ 为引理 \ref{more-morphisms-lemma-sheaf} 在注记
\ref{more-morphisms-remark-special-case} 所讨论的特殊情形下给出的 $T'$ 上集合层。令
$$
\mathcal{H} =
\SheafHom_{\mathcal{O}_T}(a^*\Omega_{X/S}, \mathcal{C}_{T/T'})
$$
为 $\mathcal{O}_T$-模层，其上有引理 \ref{more-morphisms-lemma-action-sheaf}
中的作用 $\mathcal{H} \times \mathcal{F} \to \mathcal{F}$。
我们的目标就是证明 $\mathcal{F}(T) \not = \emptyset$。换言之，
要证明 $\mathcal{F}$ 是 $T$ 上平凡的 $\mathcal{H}$-挠子（见 Cohomology,
第 \ref{cohomology-section-h1-torsors} 节）。分两步：
(I) 为证明 $\mathcal{F}$ 是挠子，须对所有 $t \in T$ 证明
$\mathcal{F}_t \not = \emptyset$（见 上同调章定义
\ref{cohomology-definition-torsor}）。
(II) 为证明 $\mathcal{F}$ 是平凡挠子，只需证明
$H^1(T, \mathcal{H}) = 0$（见 上同调章引理
\ref{cohomology-lemma-torsors-h1}——可把 $\mathcal{H}$ 看作阿贝尔层或
$\mathcal{O}_T$-模来计算其上同调；见 上同调章引理
\ref{cohomology-lemma-modules-abelian}）。

\medskip\noindent
先证明 (I)。对每个 $t \in T$，可选择 $t$ 的仿射开邻域
$U \subset T$，使 $a(U)$ 包含于仿射开集 $\Spec(A) = W \subset X$，
而后者映到仿射开集 $\Spec(R) = V \subset S$。
由 态射章引理 \ref{morphisms-lemma-smooth-characterize}，
环映射 $R \to A$ 光滑。因此由 代数章命题
\ref{algebra-proposition-smooth-formally-smooth}，环映射 $R \to A$
形式光滑。进而由引理 \ref{more-morphisms-lemma-affine-formally-smooth}，
$W \to V$ 形式光滑。故可将 $a|_U : U \to W$ 提升为 $V$-态射
$a' : U' \to W \subset X$，这表明 $\mathcal{F}(U) \not = \emptyset$。

\medskip\noindent
最后证明 (II)。由 态射章引理
\ref{morphisms-lemma-finite-presentation-differentials}，
$\Omega_{X/S}$ 有限表示（由 态射章引理
\ref{morphisms-lemma-smooth-omega-finite-locally-free}，
它甚至有限局部自由）。因此 $a^*\Omega_{X/S}$ 有限表示（见
模章引理 \ref{modules-lemma-pullback-finite-presentation}）。故层
$\mathcal{H} =
\SheafHom_{\mathcal{O}_T}(a^*\Omega_{X/S}, \mathcal{C}_{T/T'})$
由 概形章第 \ref{schemes-section-quasi-coherent} 节 中的讨论是拟凝聚的。
因此由 概形上同调章引理
\ref{coherent-lemma-quasi-coherent-affine-cohomology-zero}，
有所需的 $H^1(T, \mathcal{H}) = 0$。
\end{proof}

\noindent
局部投射拟凝聚模定义于 Properties, Section
\ref{properties-section-locally-projective}。

\begin{lemma}
\label{more-morphisms-lemma-formally-smooth-sheaf-differentials}
设 $f : X \to Y$ 是形式光滑的概形态射。则 $\Omega_{X/Y}$
在 $X$ 上局部投射。
\end{lemma}

\begin{proof}
选择满足 $f(U) \subset V$ 的仿射开集 $U \subset X$ 与 $V \subset Y$。
由引理 \ref{more-morphisms-lemma-formally-smooth-on-opens}，
$f|_U : U \to V$ 形式光滑。因此由引理
\ref{more-morphisms-lemma-affine-formally-smooth}，环映射
$\Gamma(V, \mathcal{O}_V) \to \Gamma(U, \mathcal{O}_U)$ 形式光滑。
故由 代数章引理 \ref{algebra-lemma-characterize-formally-smooth-again}，
$\Gamma(U, \mathcal{O}_U)$-模
$\Omega_{\Gamma(U, \mathcal{O}_U)/\Gamma(V, \mathcal{O}_V)}$
是投射的。因此 $\Omega_{U/V}$ 局部投射；见 Properties, Section
\ref{properties-section-locally-projective}。
\end{proof}

\begin{lemma}
\label{more-morphisms-lemma-h1-is-zero}
设 $T$ 是仿射概形，$\mathcal{F}$、$\mathcal{G}$ 是拟凝聚
$\mathcal{O}_T$-模。考虑
$\mathcal{H} = \SheafHom_{\mathcal{O}_T}(\mathcal{F}, \mathcal{G})$。
若 $\mathcal{F}$ 局部投射，则 $H^1(T, \mathcal{H}) = 0$。
\end{lemma}

\begin{proof}
由概形上局部投射层的定义（见 性质章定义
\ref{properties-definition-locally-projective}），$\mathcal{F}$ 是自由
$\mathcal{O}_T$-模的直和项。因此可假设
$\mathcal{F} = \bigoplus_{i \in I} \mathcal{O}_T$ 是自由模。
此时 $\mathcal{H} = \prod_{i \in I} \mathcal{G}$ 是拟凝聚模的积。
由 上同调章引理 \ref{cohomology-lemma-cohomology-products}，
$H^1 = 0$，因为仿射概形上拟凝聚层的上同调为零；见
概形上同调章引理
\ref{coherent-lemma-quasi-coherent-affine-cohomology-zero}。
\end{proof}

\begin{lemma}
\label{more-morphisms-lemma-formally-smooth}
设 $f : X \to Y$ 是概形态射。以下条件等价：
\begin{enumerate}
\item $f$ 形式光滑；
\item 对每个 $x \in X$，存在满足 $f(U) \subset V$ 的开集
$x \in U \subset X$ 与 $f(x) \in V \subset Y$，使得
$f|_U : U \to V$ 形式光滑；
\item 对每一对满足 $f(U) \subset V$ 的仿射开集 $U \subset X$ 与
$V \subset Y$，环映射 $\mathcal{O}_Y(V) \to \mathcal{O}_X(U)$ 形式光滑；
\item 存在仿射开覆盖 $Y = \bigcup V_j$，且对每个 $j$ 存在仿射开覆盖
$f^{-1}(V_j) = \bigcup U_{ji}$，使得对所有 $j$ 与 $i$，
% P05-ERR-CH37-0019: In this and the analogous formally unramified/etale item (4), frozen source uses undefined V,U instead of the indexed opens.
环映射 $\mathcal{O}_Y(V) \to \mathcal{O}_X(U)$ 形式光滑。
\end{enumerate}
\end{lemma}

\begin{proof}
蕴含 (1) $\Rightarrow$ (2)、(1) $\Rightarrow$ (3) 与
(2) $\Rightarrow$ (4) 来自引理 \ref{more-morphisms-lemma-formally-smooth-on-opens}。
蕴含 (3) $\Rightarrow$ (4) 是立即的。

\medskip\noindent
假设 (4)。证明 $f$ 形式光滑的论证同引理
\ref{more-morphisms-lemma-smooth-formally-smooth} 证明的第二部分。考虑实线部分交换图
$$
\xymatrix{
X \ar[d]_f & T \ar[d]^i \ar[l]^a \\
Y & T' \ar[l] \ar@{-->}[lu]
}
$$
，如定义 \ref{more-morphisms-definition-formally-smooth} 所述。我们将证明虚线箭头存在，
从而证明 $f$ 形式光滑。令 $\mathcal{F}$ 为引理 \ref{more-morphisms-lemma-sheaf}
在注记 \ref{more-morphisms-remark-special-case} 所讨论的特殊情形下给出的
$T'$ 上集合层。令
$$
\mathcal{H} =
\SheafHom_{\mathcal{O}_T}(a^*\Omega_{X/Y}, \mathcal{C}_{T/T'})
$$
为 $T$ 上的 $\mathcal{O}_T$-模层，其上有引理
\ref{more-morphisms-lemma-action-sheaf} 中的作用
$\mathcal{H} \times \mathcal{F} \to \mathcal{F}$。
此作用 $\mathcal{H} \times \mathcal{F} \to \mathcal{F}$
使 $\mathcal{F}$ 成为伪 $\mathcal{H}$-挠子；见 Cohomology,
定义 \ref{cohomology-definition-torsor}。目标是证明
$\mathcal{F}$ 是平凡的 $\mathcal{H}$-挠子。分两步：
(I) 为证明 $\mathcal{F}$ 是挠子，须证明 $\mathcal{F}$ 局部有截面。
(II) 为证明 $\mathcal{F}$ 是平凡挠子，只需证明
$H^1(T, \mathcal{H}) = 0$；见 上同调章引理
\ref{cohomology-lemma-torsors-h1}。

\medskip\noindent
先证明 (I)。对每个 $t \in T$，可选择 $t$ 的仿射开邻域
$W \subset T$，使对某些 $i, j$ 有 $a(W)$ 包含于 $U_{ji}$。
令 $W' \subset T'$ 为相应的开子概形。由假设 (4)，可将
$a|_W : W \to U_{ji}$ 提升为 $V_j$-态射
$a' : W' \to U_{ji}$，这表明 $\mathcal{F}(W')$ 非空。

\medskip\noindent
最后证明 (II)。由引理 \ref{more-morphisms-lemma-formally-smooth-sheaf-differentials}，
$\Omega_{U_{ji}/V_j}$ 局部投射。因此 $\Omega_{X/Y}$ 局部投射；见
性质章引理 \ref{properties-lemma-locally-projective}。
故 $a^*\Omega_{X/Y}$ 局部投射；见 性质章引理
\ref{properties-lemma-locally-projective-pullback}。因此
$$
H^1(T, \mathcal{H}) =
H^1(T, \SheafHom_{\mathcal{O}_T}(a^*\Omega_{X/Y}, \mathcal{C}_{T/T'})) = 0
$$
，由引理 \ref{more-morphisms-lemma-h1-is-zero} 得到所需结论。
\end{proof}

\begin{lemma}
\label{more-morphisms-lemma-triangle-differentials-formally-smooth}
设 $f : X \to Y$、$g : Y \to S$ 是概形态射，并假设 $f$ 形式光滑。则
$$
0 \to f^*\Omega_{Y/S} \to \Omega_{X/S} \to \Omega_{X/Y} \to 0
$$
（见 态射章引理
\ref{morphisms-lemma-triangle-differentials}）是短正合列。
\end{lemma}

\begin{proof}
此引理的代数版本如下：给定环映射 $A \to B \to C$，其中
$B \to C$ 形式光滑，则
$$
0 \to C \otimes_B \Omega_{B/A} \to \Omega_{C/A} \to \Omega_{C/B} \to 0
$$
这一 代数章引理 \ref{algebra-lemma-exact-sequence-differentials}
中的序列正合。这就是 代数章引理
\ref{algebra-lemma-ses-formally-smooth}。
\end{proof}

\begin{lemma}
\label{more-morphisms-lemma-differentials-formally-unramified-formally-smooth}
设 $h : Z \to X$ 是 $S$ 上形式非分歧的概形态射，并假设 $Z$
在 $S$ 上形式光滑。则典范正合列
$$
0 \to \mathcal{C}_{Z/X} \to h^*\Omega_{X/S} \to \Omega_{Z/S} \to 0
$$
（见引理 \ref{more-morphisms-lemma-universally-unramified-differentials-sequence}）
是短正合列。
\end{lemma}

\begin{proof}
设 $Z \to Z'$ 是 $Z$ 在 $X$ 上的泛一阶加厚。由引理
\ref{more-morphisms-lemma-universally-unramified-differentials-sequence} 的证明，
我们的序列与序列
$$
\mathcal{C}_{Z/Z'} \to \Omega_{Z'/S} \otimes \mathcal{O}_Z \to
\Omega_{Z/S} \to 0.
$$
等同。由于 $Z \to S$ 形式光滑，可在 $Z'$ 上局部地找到包含映射
$Z \to Z'$ 的一个 $S$ 上左逆 $Z' \to Z$。因此该序列局部分裂；见
态射章引理
\ref{morphisms-lemma-differentials-relative-immersion-section}。
\end{proof}

\begin{lemma}
\label{more-morphisms-lemma-two-unramified-morphisms-formally-smooth}
设
$$
\xymatrix{
Z \ar[r]_i \ar[rd]_j & X \ar[d]^f \\
& Y
}
$$
是概形的交换图，其中 $i$ 与 $j$ 形式非分歧，且 $f$ 形式光滑。
则典范正合列
$$
0 \to
\mathcal{C}_{Z/Y} \to
\mathcal{C}_{Z/X} \to
i^*\Omega_{X/Y} \to 0
$$
（见引理 \ref{more-morphisms-lemma-two-unramified-morphisms}）正合且局部分裂。
\end{lemma}

\begin{proof}
以 $Z \to Z'$ 表示 $Z$ 在 $X$ 上的泛一阶加厚，以 $Z \to Z''$
表示 $Z$ 在 $Y$ 上的泛一阶加厚。
% P05-ERR-CH37-0020: Frozen source again cites the differential-sequence lemma for the functorial map Z' -> Z''.
由引理 \ref{more-morphisms-lemma-universally-unramified-differentials-sequence}，
有典范态射 $Z' \to Z''$，从而有交换图
$$
\xymatrix{
Z \ar[r]_{i'} \ar[rd]_{j'} & Z' \ar[r]_a \ar[d]^k & X \ar[d]^f \\
& Z'' \ar[r]^b & Y
}
$$
在引理 \ref{more-morphisms-lemma-two-unramified-morphisms} 的证明中，
我们已将上述序列与序列
$$
\mathcal{C}_{Z/Z''} \to
\mathcal{C}_{Z/Z'} \to
(i')^*\Omega_{Z'/Z''} \to 0
$$
等同。设 $U'' \subset Z''$ 是仿射开集，以 $U \subset Z$ 与
$U' \subset Z'$ 表示相应的仿射开子概形。由于 $f$ 形式光滑，存在态射
$h : U'' \to X$，它在 $U$ 上与 $i$ 相同，且 $f \circ h$ 等于
$b|_{U''}$。由于 $Z'$ 是泛一阶加厚，得到唯一的态射
% P05-ERR-CH37-0021: Frozen source writes the ill-typed equality g = a \circ h; formula retained pending canon review.
$g : U'' \to Z'$，满足 $g = a \circ h$。$Z''$ 的泛性质推出
$k \circ g$ 是包含映射 $U'' \to Z''$。故 $g$ 是 $k$ 的左逆。图示为
$$
\xymatrix{
U \ar[d] \ar[r] & Z' \ar[d]^k \\
U'' \ar[r] \ar[ru]^g & Z''
}
$$
因此 $g$ 诱导映射
$\mathcal{C}_{Z/Z'}|_U \to \mathcal{C}_{Z/Z''}|_U$，
它是 $U$ 上映射 $\mathcal{C}_{Z/Z''} \to \mathcal{C}_{Z/Z'}$ 的左逆。
\end{proof}



































\section{Noether 基底上的光滑性}
\label{more-morphisms-section-smooth-Noetherian}

\noindent
事实证明，若基底是 Noether 的，则在形式光滑性的表述中可以少要求一些。
在某种意义上，下面这些引理是形变理论的开端。

\begin{lemma}
\label{more-morphisms-lemma-lifting-along-artinian-at-point}
设 $f : X \to S$ 是概形态射，$x \in X$。假设 $S$ 局部 Noether，
且 $f$ 局部有限型。以下条件等价：
\begin{enumerate}
\item $f$ 在 $x$ 处光滑；
\item 对每个实线部分交换的图
$$
\xymatrix{
X \ar[d]_f & \Spec(B) \ar[d]^i \ar[l]^-\alpha \\
S & \Spec(B') \ar[l]_-{\beta} \ar@{-->}[lu]
}
$$
，其中 $B' \to B$ 是局部环的满射，$\Ker(B' \to B)$ 平方为零，
且 $\alpha$ 把 $\Spec(B)$ 的闭点映到 $x$，存在一条使该图交换的虚线箭头；
\item 与 (2) 相同，但 $B' \to B$ 遍历小扩张（见 代数章定义
\ref{algebra-definition-small-extension}）；
\item 与 (2) 相同，但 $B' \to B$ 遍历如下小扩张：$\alpha$ 诱导同构
$\kappa(x) \to \kappa(\mathfrak m)$，其中
$\mathfrak m \subset B$ 是极大理想。
\end{enumerate}
\end{lemma}

\begin{proof}
选择 $f(x)$ 的仿射邻域 $V \subset S$，并选择 $x$ 的仿射邻域
$U \subset X$，使得 $f(U) \subset V$。对 (2) 中任意“测试”图，
态射 $\alpha$ 都把 $\Spec(B)$ 映入 $U$，而态射 $\beta$ 把
$\Spec(B')$ 映入 $V$（见 Schemes, Section
\ref{schemes-section-points}）。故引理归结为仿射概形之间的态射
$f|_U : U \to V$。（事实上，$V$ 是 Noether 的，且 $f|_U$ 有限型；见
性质章引理 \ref{properties-lemma-locally-Noetherian} 与
态射章引理 \ref{morphisms-lemma-locally-finite-type-characterize}。）
在此仿射情形下，引理与 代数章引理
\ref{algebra-lemma-smooth-test-artinian} 完全相同。
\end{proof}

\noindent
有时知道如下事实很有用：只需对以有限型点为“中心”的小扩张检验提升判据
（见 态射章第 \ref{morphisms-section-points-finite-type} 节）。

\begin{lemma}
\label{more-morphisms-lemma-lifting-along-artinian}
设 $f : X \to S$ 是概形态射。假设 $S$ 局部 Noether，且 $f$ 局部有限型。
以下条件等价：
\begin{enumerate}
\item $f$ 光滑；
\item 对每个实线部分交换的图
$$
\xymatrix{
X \ar[d]_f & \Spec(B) \ar[d]^i \ar[l]^-\alpha \\
S & \Spec(B') \ar[l]_-{\beta} \ar@{-->}[lu]
}
$$
，其中 $B' \to B$ 是 Artin 局部环的小扩张，且 $\beta$ 有限型（！），
存在一条使该图交换的虚线箭头。
\end{enumerate}
\end{lemma}

\begin{proof}
若 $f$ 光滑，则无穷小提升判据（引理
\ref{more-morphisms-lemma-smooth-formally-smooth}）表明 $f$ 形式光滑，故 (2) 成立。

\medskip\noindent
假设 (2)。$f$ 不光滑的点 $x \in X$ 构成 $X$ 的闭子集 $T$。
由 态射章第 \ref{morphisms-section-points-finite-type} 节 中的讨论，
若 $T \not = \emptyset$，则存在点 $x \in T \subset X$，使态射
$$
\Spec(\kappa(x)) \to X \to S
$$
有限型（具体地，在 $T$ 中任取一个在 $X$ 的某个仿射开集内为闭点的点 $x$）。
由 态射章引理 \ref{morphisms-lemma-artinian-finite-type}，
给定剩余域为 $\kappa(x)$ 的任意 Artin 局部环 $B'$，任意态射
$\beta : \Spec(B') \to S$ 都有限型。因此，引理
\ref{more-morphisms-lemma-lifting-along-artinian-at-point} (4) 中使用的所有图都对应于
本引理 (2) 中的图。故 $X \to S$ 在 $x$ 处光滑，矛盾。
% P05-ERR-CH37-0022: Frozen source has the typo ``smooth a x''; translated as intended while preserving a non-rendering erratum marker.
\end{proof}

\noindent
下面是一个有用的应用。

\begin{lemma}
\label{more-morphisms-lemma-check-smoothness-on-infinitesimal-nbhds}
设 $f : X \to S$ 是局部 Noether 概形之间的有限型态射。
设 $Z \subset S$ 是闭子概形，其 $n$ 阶无穷小邻域为
$Z_n \subset S$。令 $X_n = Z_n \times_S X$。
\begin{enumerate}
\item 若对所有 $n$，$X_n \to Z_n$ 都光滑，则 $f$ 在
$f^{-1}(Z)$ 的每一点处光滑。
\item 若对所有 $n$，$X_n \to Z_n$ 都平展，则 $f$ 在
$f^{-1}(Z)$ 的每一点处平展。
\end{enumerate}
\end{lemma}

\begin{proof}
假设对所有 $n$，$X_n \to Z_n$ 都光滑。设 $x \in X$ 位于 $Z$ 的某一点上方。
给定引理 \ref{more-morphisms-lemma-lifting-along-artinian-at-point} (3) 中那样的小扩张
$B' \to B$ 与态射 $\alpha$、$\beta$，$B'$ 的极大理想幂零（因为 $B'$
是 Artin 的），故对适当的 $n$，态射 $\beta$ 经由 $Z_n$ 分解，而
$\alpha$ 经由 $X_n$ 分解。于是应用 $X_n \to Z_n$ 的提升性质，得到图中
所需的虚线箭头。这证明了 (1)。(2) 由 (1) 以及如下事实得出：一个态射平展，
当且仅当它光滑且相对维数为 $0$。
\end{proof}

\begin{lemma}
\label{more-morphisms-lemma-check-flatness-on-infinitesimal-nbhds}
设 $f : X \to S$ 是局部 Noether 概形之间的态射。设 $Z \subset S$
是闭子概形，其 $n$ 阶无穷小邻域为 $Z_n \subset S$。令
$X_n = Z_n \times_S X$。若对所有 $n$，$X_n \to Z_n$ 都平坦，
则 $f$ 在 $f^{-1}(Z)$ 的每一点处平坦。
\end{lemma}

\begin{proof}
这就是将 代数章引理 \ref{algebra-lemma-flat-module-powers}
翻译成概形的语言。
\end{proof}








\section{朴素余切复形}
\label{more-morphisms-section-netherlander}

\noindent
本节延续 模章第 \ref{modules-section-netherlander} 节，
而后者又延续 代数章第 \ref{algebra-section-netherlander} 节 中的讨论。

\begin{definition}
\label{more-morphisms-definition-netherlander}
设 $f : X \to Y$ 是概形态射。$f$ 的{\it 朴素余切复形}是 Modules,
定义 \ref{modules-definition-cotangent-complex-morphism-ringed-topoi}
中定义的复形。记作 $\NL_f$ 或 $\NL_{X/Y}$。
\end{definition}

\begin{lemma}
\label{more-morphisms-lemma-NL-affine}
设 $f : X \to Y$ 是概形态射。设 $\Spec(A) = U \subset X$ 与
$\Spec(R) = V \subset S$ 是满足 $f(U) \subset V$ 的仿射开集。
存在复形的典范映射
$$
\widetilde{\NL_{A/R}} \longrightarrow \NL_{X/Y}|_U
$$
，它在 $D(\mathcal{O}_U)$ 中是同构。
% P05-ERR-CH37-0023: Frozen source names V as an open of S although the morphism target is Y.
\end{lemma}

\begin{proof}
由 模章第 \ref{modules-section-netherlander} 节 中
$\NL_{X/Y}$ 的构造，有 $A = \mathcal{O}_X(U)$-模复形的典范映射\linebreak[4]
$\NL_{\mathcal{O}_X(U)/f^{-1}\mathcal{O}_Y(U)} \to \NL_{X/Y}(U)$，
并且它与进一步的限制相容。使用典范映射
$R \to f^{-1}\mathcal{O}_Y(U)$，得到 $A$-模复形的典范映射
$\NL_{A/R} \to \NL_{\mathcal{O}_X(U)/f^{-1}\mathcal{O}_Y(U)}$。
使用 $\widetilde{\ }$ 函子的泛性质（见 概形章引理
\ref{schemes-lemma-compare-constructions}），得到引理所述的映射。
要检验此映射在上同调层上为同构，只需检验它在茎上诱导同构。
这由 代数章引理 \ref{algebra-lemma-NL-localize-bottom}、
\ref{algebra-lemma-localize-NL} 与 模章引理
\ref{modules-lemma-stalk-NL} 得出（这里还使用了如下描述：在点
$\mathfrak p \in \Spec(A)$ 处，$\mathcal{O}_X$ 与
$f^{-1}\mathcal{O}_Y$ 的茎分别是 $A_\mathfrak p$ 与 $R_\mathfrak q$，
其中 $\mathfrak q = R \cap \mathfrak p$；所用参考为 概形章引理
\ref{schemes-lemma-spec-sheaves} 与 层章引理
\ref{sheaves-lemma-stalk-pullback}）。
\end{proof}

\begin{lemma}
\label{more-morphisms-lemma-netherlander-quasi-coherent}
设 $f : X \to Y$ 是概形态射。复形 $\NL_{X/Y}$ 的上同调层拟凝聚，
在次数 $-1$、$0$ 之外为零，并且在次数 $0$ 等于 $\Omega_{X/Y}$。
\end{lemma}

\begin{proof}
由 模章第 \ref{modules-section-netherlander} 节 中朴素余切复形的构造，
$\NL_{X/Y}$ 位于次数 $-1$、$0$，且其次数 $0$ 上同调为 $\Omega_{X/Y}$。
微分层拟凝聚（由 态射章引理
\ref{morphisms-lemma-differentials-diagonal}）。因此只需证明
$H^{-1}(\NL_{X/Y})$ 拟凝聚。使用引理 \ref{more-morphisms-lemma-NL-affine}
在各仿射开集上检验即可。
\end{proof}

\begin{lemma}
\label{more-morphisms-lemma-netherlander-fp}
设 $f : X \to Y$ 是概形态射。若 $f$ 局部有限表示，则在 $X$ 上局部地，
$\NL_{X/Y}$ 拟同构于复形
$$
\ldots \to 0 \to \mathcal{F}^{-1} \to \mathcal{F}^0 \to 0 \to \ldots
$$
，其中各项为拟凝聚 $\mathcal{O}_X$-模，$\mathcal{F}^0$ 有限表示，
而 $\mathcal{F}^{-1}$ 有限型。
\end{lemma}

\begin{proof}
由引理 \ref{more-morphisms-lemma-NL-affine}，只需证明当 $R \to A$ 是有限表示环映射时，
$\NL_{A/R}$ 具有这种形状。写
$A = R[x_1, \ldots, x_n]/I$，其中 $I$ 有限生成。则 $I/I^2$
是有限 $A$-模，且 $\NL_{A/R}$ 拟同构于
$$
\ldots \to 0 \to I/I^2 \to
\bigoplus\nolimits_{i = 1, \ldots, n} A\text{d}x_i \to 0 \to \ldots
$$
；这由 代数章第 \ref{algebra-section-netherlander} 节，
特别是 代数章引理 \ref{algebra-lemma-NL-homotopy} 得出。
\end{proof}

\begin{lemma}
\label{more-morphisms-lemma-NL-formally-smooth}
设 $f : X \to Y$ 是概形态射。以下条件等价：
\begin{enumerate}
\item $f$ 形式光滑；
\item $H^{-1}(\NL_{X/Y}) = 0$ 及 $H^0(\NL_{X/Y}) = \Omega_{X/Y}$
局部投射。
\end{enumerate}
\end{lemma}

\begin{proof}
这由 代数章命题
\ref{algebra-proposition-characterize-formally-smooth} 与引理
\ref{more-morphisms-lemma-formally-smooth} 得出。
\end{proof}

\begin{lemma}
\label{more-morphisms-lemma-NL-formally-etale}
设 $f : X \to Y$ 是概形态射。以下条件等价：
\begin{enumerate}
\item $f$ 形式平展；
\item $H^{-1}(\NL_{X/Y}) = H^0(\NL_{X/Y}) = 0$.
\end{enumerate}
\end{lemma}

\begin{proof}
形式平展态射形式光滑，故由引理 \ref{more-morphisms-lemma-NL-formally-smooth}，
$H^{-1}(\NL_{X/Y}) = 0$。另一方面，由引理
\ref{more-morphisms-lemma-characterize-formally-etale}，$\Omega_{X/Y} = 0$。
反之，若 (2) 成立，则由引理 \ref{more-morphisms-lemma-NL-formally-smooth}，
$f$ 形式光滑；由引理 \ref{more-morphisms-lemma-formally-unramified-differentials}
又知其形式非分歧；从而由引理
\ref{more-morphisms-lemma-formally-etale-unramified-smooth}，它形式平展。
\end{proof}

\begin{lemma}
\label{more-morphisms-lemma-NL-smooth}
设 $f : X \to Y$ 是概形态射。以下条件等价：
\begin{enumerate}
\item $f$ 光滑；
\item $f$ 局部有限表示，
$H^{-1}(\NL_{X/Y}) = 0$, 及 $H^0(\NL_{X/Y}) = \Omega_{X/Y}$
有限局部自由。
\end{enumerate}
\end{lemma}

\begin{proof}
这由光滑环同态的定义（代数章定义
\ref{algebra-definition-smooth}）、引理 \ref{more-morphisms-lemma-NL-affine}
以及光滑概形态射的定义（态射章定义
\ref{morphisms-definition-smooth}）得出。这里还使用了如下事实：
对环上的模而言，有限局部自由等价于有限投射（代数章引理
\ref{algebra-lemma-finite-projective}）。
\end{proof}

\begin{lemma}
\label{more-morphisms-lemma-NL-etale}
设 $f : X \to Y$ 是概形态射。以下条件等价：
\begin{enumerate}
\item $f$ 平展；
\item $f$ 局部有限表示，且
$H^{-1}(\NL_{X/Y}) = H^0(\NL_{X/Y}) = 0$.
\end{enumerate}
\end{lemma}

\begin{proof}
这由平展环同态的定义（代数章定义
\ref{algebra-definition-etale}）、引理 \ref{more-morphisms-lemma-NL-affine}
以及平展概形态射的定义（态射章定义
\ref{morphisms-definition-etale}）得出。
\end{proof}


\begin{lemma}
\label{more-morphisms-lemma-NL-immersion}
设 $i : Z \to X$ 是概形的浸入。则在 $D(\mathcal{O}_Z)$ 中，
$\NL_{Z/X}$ 同构于 $\mathcal{C}_{Z/X}[1]$，其中
$\mathcal{C}_{Z/X}$ 是 $Z$ 在 $X$ 中的余法层。
\end{lemma}

\begin{proof}
这由 代数章引理 \ref{algebra-lemma-NL-surjection}、
态射章引理 \ref{morphisms-lemma-affine-conormal} 与引理
\ref{more-morphisms-lemma-NL-affine} 得出。
\end{proof}

\begin{lemma}
\label{more-morphisms-lemma-exact-sequence-NL}
设 $f : X \to Y$ 与 $g : Y \to Z$ 是概形态射。则有上同调层的典范六项正合列
$$
H^{-1}(f^*\NL_{Y/Z}) \to
H^{-1}(\NL_{X/Z}) \to
H^{-1}(\NL_{X/Y}) \to
f^*\Omega_{Y/Z} \to \Omega_{X/Z} \to \Omega_{X/Y} \to 0
$$
。
\end{lemma}

\begin{proof}
这是 模章引理
\ref{modules-lemma-exact-sequence-NL-ringed-topoi} 的特殊情形。
\end{proof}

\begin{lemma}
\label{more-morphisms-lemma-get-triangle-NL}
设 $f : X \to Y$ 与 $Y \to Z$ 是概形态射。假设 $X \to Y$
是完全交态射。则在 $D(\mathcal{O}_X)$ 中有典范可区别三角
$$
f^*\NL_{Y/Z} \to \NL_{X/Z} \to \NL_{X/Y} \to f^*\NL_{Y/Z}[1]
$$
，它恢复引理 \ref{more-morphisms-lemma-exact-sequence-NL} 的 $6$ 项正合列。
\end{lemma}

\begin{proof}
只需证明 模章引理
\ref{modules-lemma-exact-sequence-NL-ringed-topoi} 的典范映射
$$
f^*\NL_{Y/Z} \to \text{Cone}(\NL_{X/Z} \to \NL_{X/Y})[-1]
$$
在 $D(\mathcal{O}_X)$ 中是同构。为此，只需证明该 $6$ 项序列左端为零，
即 $H^{-1}(f^*\NL_{Y/Z}) \to H^{-1}(\NL_{X/Z})$ 是单射。
在仿射局部，这由 代数进阶章引理
\ref{more-algebra-lemma-transitive-lci-at-end} 中相应的代数结果得出。
使用引理 \ref{more-morphisms-lemma-NL-affine} 将其翻译成代数即可。
\end{proof}

\begin{lemma}
\label{more-morphisms-lemma-smooth-etale-permanence}
设 $X \to Y \to Z$ 是概形态射。假设 $X \to Z$ 光滑，
且 $Y \to Z$ 平展。则 $X \to Y$ 光滑。
\end{lemma}

\begin{proof}
由 态射章引理
\ref{morphisms-lemma-finite-presentation-permanence}，态射 $X \to Y$
局部有限表示。由引理 \ref{more-morphisms-lemma-NL-smooth}，
$H^{-1}(\NL_{X/Z}) = 0$，且模 $\Omega_{X/Z}$ 有限局部自由。
由引理 \ref{more-morphisms-lemma-NL-etale}，
$H^{-1}(\NL_{Y/Z}) = H^0(\NL_{Y/Z}) = 0$。
由引理 \ref{more-morphisms-lemma-exact-sequence-NL}，得到
$H^{-1}(\NL_{X/Y}) = 0$，且
$\Omega_{X/Y} \cong \Omega_{X/Z}$ 有限局部自由。
由引理 \ref{more-morphisms-lemma-NL-smooth}，态射 $X \to Y$ 光滑。
\end{proof}

\begin{lemma}
\label{more-morphisms-lemma-get-NL}
设概形态射 $f : X \to Y$ 分解为 $f = g \circ i$，其中 $i$ 是浸入，
而 $g : P \to Y$ 形式光滑（例如光滑）。则在 $D(\mathcal{O}_X)$ 中有典范同构
$$
\NL_{X/Y} \cong \left(\mathcal{C}_{X/P} \to i^*\Omega_{P/Y}\right)
$$
，其中余法层 $\mathcal{C}_{X/P}$ 置于次数 $-1$。
\end{lemma}

\begin{proof}
（括号中的断言见引理 \ref{more-morphisms-lemma-smooth-formally-smooth}。）
由引理 \ref{more-morphisms-lemma-NL-immersion} 与 \ref{more-morphisms-lemma-NL-formally-smooth}，有
$\NL_{X/P} = \mathcal{C}_{X/P}[1]$ 与 $\NL_{P/Y} = \Omega_{P/Y}$，
其中 $\Omega_{P/Y}$ 局部投射。这还推出
$i^*\NL_{P/Y} \to i^*\Omega_{P/Y}$ 是拟同构（略去一个小细节；
理由是 $i^*\NL_{P/Y}$ 与 $\tau_{\geq -1}Li^*\NL_{P/Y}$ 相同；见
代数进阶章引理 \ref{more-algebra-lemma-tensor-NL}）。
因此 模章引理
\ref{modules-lemma-exact-sequence-NL-ringed-topoi} 的典范映射
$$
i^*\NL_{P/Y} \to \text{Cone}(\NL_{X/Y} \to \NL_{X/P})[-1]
$$
在 $D(\mathcal{O}_X)$ 中是同构，因为由上可知上同调群
$H^{-1}(i^*\NL_{P/Y})$ 为零。换言之，有可区别三角
$$
i^*\NL_{P/Y} \to \NL_{X/Y} \to \NL_{X/P} \to i^*\NL_{P/Y}[1]
$$
显然，这表示 $\NL_{X/Y}$ 是映射
$\NL_{X/P}[-1] \to i^*\NL_{P/Y}$ 的锥。根据上面对这些导出范畴对象的
上同调层的计算，这等价于引理的断言。
\end{proof}

\begin{lemma}
\label{more-morphisms-lemma-base-change-NL}
考虑概形的笛卡尔图
$$
\xymatrix{
X' \ar[r]_{g'} \ar[d] & X \ar[d] \\
Y' \ar[r] & Y
}
$$
典范映射 $(g')^*\NL_{X/Y} \to \NL_{X'/Y'}$ 在 $H^0$ 上诱导同构，
在 $H^{-1}$ 上诱导满射。
\end{lemma}

\begin{proof}
翻译成代数后，这就是 代数进阶章引理
\ref{more-algebra-lemma-base-change-NL}。用引理 \ref{more-morphisms-lemma-NL-affine}
完成翻译。
\end{proof}

\begin{lemma}
\label{more-morphisms-lemma-flat-base-change-NL}
考虑概形的笛卡尔图
$$
\xymatrix{
X' \ar[d] \ar[r]_{g'} & X \ar[d] \\
Y' \ar[r] & Y
}
$$
若 $Y' \to Y$ 平坦，则典范映射
$(g')^*\NL_{X/Y} \to \NL_{X'/Y'}$ 是拟同构。
\end{lemma}

\begin{proof}
由引理 \ref{more-morphisms-lemma-NL-affine}，这从 代数章引理
\ref{algebra-lemma-change-base-NL} 得出。
\end{proof}

\begin{lemma}
\label{more-morphisms-lemma-base-change-NL-flat}
考虑概形的笛卡尔图
$$
\xymatrix{
X' \ar[r]_{g'} \ar[d] & X \ar[d] \\
Y' \ar[r] & Y
}
$$
若 $X \to Y$ 平坦，则典范映射
$(g')^*\NL_{X/Y} \to \NL_{X'/Y'}$ 是拟同构。
若此外 $\NL_{X/Y}$ 的 Tor 振幅包含于 $[-1, 0]$，则
$L(g')^*\NL_{X/Y} \to \NL_{X'/Y'}$ 也是拟同构。
\end{lemma}

\begin{proof}
翻译成代数后，这就是 代数进阶章引理
\ref{more-algebra-lemma-base-change-NL-flat}。使用引理
\ref{more-morphisms-lemma-NL-affine} 与 概形的导出范畴章引理
\ref{perfect-lemma-affine-compare-bounded}、
\ref{perfect-lemma-tor-dimension-affine} 完成翻译。
\end{proof}











\section{概形范畴中的推出，I}
\label{more-morphisms-section-pushouts}

\noindent
本节构造 $Y \leftarrow X \rightarrow X'$ 的推出，其中 $X \to Y$ 仿射，
而 $X \to X'$ 是加厚。事实上，这对我们是一个重要情形，值得详细讨论。
在节 \ref{more-morphisms-section-pushouts-II} 中将讨论一个更有趣也更困难的情形。
关于任意范畴中推出的一般讨论，见 Categories, Section
\ref{categories-section-pushouts}。

\begin{lemma}
\label{more-morphisms-lemma-basic-example-pushout}
设 $A' \to A$ 是环的满射，$B \to A$ 是环映射，并设
$B' = B \times_A A'$ 为环的纤维积。令
$S = \Spec(A)$、$S' = \Spec(A')$、$T = \Spec(B)$ 且
$T' = \Spec(B')$。则
$$
\vcenter{
\xymatrix{
S \ar[r]_i \ar[d]_f & S' \ar[d]^{f'} \\
T \ar[r]^{i'} & T'
}
}
\quad\text{corresponding to}\quad
\vcenter{
\xymatrix{
A & A' \ar[l] \\
B \ar[u] & B' \ar[l] \ar[u]
}
}
$$
是概形的推出。
\end{lemma}

\begin{proof}
由 代数进阶章引理
\ref{more-algebra-lemma-points-of-fibre-product}，作为拓扑空间有
$T' = T \amalg_S S'$；也就是说，该图在拓扑空间范畴中是推出。
接着考虑映射
$$
((i')^\sharp, (f')^\sharp) :
\mathcal{O}_{T'}
\longrightarrow
i'_*\mathcal{O}_T \times_{g_*\mathcal{O}_S} f'_*\mathcal{O}_{S'}
$$
，其中 $g = i' \circ f = f' \circ i$。我们断言此映射是环层的同构。
事实上，可把两边都看成拟凝聚 $\mathcal{O}_{T'}$-模（右边使用
概形章引理 \ref{schemes-lemma-push-forward-quasi-coherent}），
且该映射为 $\mathcal{O}_{T'}$-线性的。因此只需证明它在整体截面层面上
是同构（概形章引理
\ref{schemes-lemma-equivalence-quasi-coherent}）。在整体截面上，
恢复引理陈述中的等同 $B' \to B \times_A A'$（这正是选取 $B'$ 的方式）。

\medskip\noindent
设 $X$ 是概形，并假设给定满足 $m' \circ i = n \circ f$ 的概形态射
$m' : S' \to X$ 与 $n : T \to X$（将此复合记为 $m$）。由于
$T' = T \amalg_S S'$（见上文），得到与 $m'$、$n$ 相容的唯一拓扑空间映射
$n' : T' \to X$。由上一段对 $\mathcal{O}_{T'}$ 的描述，
得到环层的唯一同态
$$
(n')^\sharp :
\mathcal{O}_X
\longrightarrow
(n')_*\mathcal{O}_{T'} =
m'_*\mathcal{O}_T \times_{m_*\mathcal{O}_T} n_*\mathcal{O}_S
$$
，它由 $(m')^\sharp$ 与 $n^\sharp$ 给出。因此
$(n', (n')^\sharp)$ 是与 $m'$、$n$ 相容的唯一环化空间态射
$T' \to X$。为完成证明，只需证明 $n'$ 是概形态射，
即局部环化空间的态射。

\medskip\noindent
设 $t' \in T'$ 的像为 $x \in X$。须证明
$\mathcal{O}_{X, x} \to \mathcal{O}_{T', t'}$ 是局部同态。
若 $t' \not \in T$，则 $t'$ 是唯一一点 $s' \in S'$ 的像，且
$\mathcal{O}_{T', t'} = \mathcal{O}_{S', s'}$。事实上，
$S' \setminus S \to T' \setminus T$ 是概形的同构，因为
$B' \to A'$ 诱导同构
$\Ker(B' \to B) = \Ker(A' \to A)$。若 $t'$ 是 $t \in T$ 的像，
则已知复合
$\mathcal{O}_{X, x} \to \mathcal{O}_{T', t'} \to \mathcal{O}_{T, t}$
是局部同态，故同样得到结论。
\end{proof}

\begin{lemma}
\label{more-morphisms-lemma-pushout-fpqc-local}
设 $\mathcal{I} \to (\Sch/S)_{fppf}$、$i \mapsto X_i$ 是概形的图，
而 $(W, X_i \to W)$ 是该图在概形范畴中的余锥（范畴章说明
\ref{categories-remark-cones-and-cocones}）。若存在概形的 fpqc 覆盖
$\{W_a \to W\}_{a \in A}$，使得
\begin{enumerate}
\item 对所有 $a \in A$，在概形范畴中有
$W_a = \colim X_i \times_W W_a$
；
\item 对所有 $a, b \in A$，在概形范畴中有
$W_a \times_W W_b = \colim X_i \times_W W_a \times_W W_b$
；
\end{enumerate}
则在概形范畴中 $W = \colim X_i$。
\end{lemma}

\begin{proof}
事实上，对概形 $T$，一个态射 $W \to T$ 等价于一族态射
$W_a \to T$、$a \in A$，它们在交叠 $W_a \times_W W_b$ 上相等；见
下降章引理 \ref{descent-lemma-fpqc-universal-effective-epimorphisms}。
\end{proof}

\begin{lemma}
\label{more-morphisms-lemma-pushout-along-thickening}
设 $X \to X'$ 是概形的加厚，而 $X \to Y$ 是仿射概形态射。
则在概形范畴中存在推出
$$
\xymatrix{
X \ar[r] \ar[d]_f
&
X' \ar[d]^{f'}
\\
Y \ar[r]
&
Y'
}
$$
。此外，$Y \subset Y'$ 是加厚，$X = Y \times_{Y'} X'$，并且在
$|Y| = |Y'|$ 上有层的等式
$$
\mathcal{O}_{Y'} = \mathcal{O}_Y \times_{f_*\mathcal{O}_X} f'_*\mathcal{O}_{X'}
$$
。
\end{lemma}

\begin{proof}
先将 $Y'$ 构造为环化空间。具体地，作为拓扑空间取 $Y' = Y$。
以 $f' : X' \to Y'$ 表示等于 $f$ 的拓扑空间映射。取引理等式的右边作为
结构层 $\mathcal{O}_{Y'}$。为证明 $Y'$ 是概形，须证明每一点都有仿射邻域。
由于层的纤维积构造与限制到开集可交换，可假设 $Y$ 仿射。
于是 $X$ 仿射（因为 $f$ 仿射），且 $X'$ 也仿射（见引理
\ref{more-morphisms-lemma-thickening-affine-scheme}）。设
$Y \leftarrow X \rightarrow X'$ 对应于 $B \rightarrow A \leftarrow A'$。
令 $B' = B \times_A A'$；这就是 $\mathcal{O}_{Y'}$ 的整体截面。
由于 $A' \to A$ 是核局部幂零的满射，$B' \to B$ 也是核局部幂零的满射。
故作为拓扑空间，$\Spec(B') = \Spec(B)$。我们断言
$Y' = \Spec(B')$。为此，将证明：对 $g' \in B'$ 及其像 $g \in B$，有
$\mathcal{O}_{Y'}(D(g)) = B'_{g'}$。事实上，由 代数进阶章引理
\ref{more-algebra-lemma-diagram-localize} 有
$$
(B')_{g'} = B_g \times_{A_h} A'_{h'}
$$
，其中 $h \in A$、$h' \in A'$ 是 $g'$ 的像。由于 $B_g$、分别地
$A_h$、分别地 $A'_{h'}$ 等于 $\mathcal{O}_Y(D(g))$、分别地
$f_*\mathcal{O}_X(D(g))$、分别地 $f'_*\mathcal{O}_{X'}(D(g))$，
断言得证。

\medskip\noindent
还须证明 $Y'$ 是推出。上述讨论表明，概形 $Y'$ 有仿射开覆盖
$Y' = \bigcup W'_i$，使相应的开集 $U'_i \subset X'$、$W_i \subset Y$
与 $U_i \subset X$ 都仿射。此外，若 $A'_i$、$B_i$、$A_i$
分别是对应于 $U'_i$、$W_i$、$U_i$ 的环，则 $W'_i$ 对应于
$B_i \times_{A_i} A'_i$。因此可应用引理
\ref{more-morphisms-lemma-basic-example-pushout} 与 \ref{more-morphisms-lemma-pushout-fpqc-local}，
得出我们的构造在概形范畴中是推出。
\end{proof}

\noindent
在下面的引理中，使用 范畴章例
\ref{categories-example-2-fibre-product-categories} 中定义的范畴纤维积。

\begin{lemma}
\label{more-morphisms-lemma-equivalence-categories-schemes-over-pushout}
设 $X \to X'$ 是概形的加厚，$X \to Y$ 是仿射概形态射，并设
$Y' = Y \amalg_X X'$ 为推出（见引理
\ref{more-morphisms-lemma-pushout-along-thickening}）。基变换给出函子
$$
F : (\Sch/Y') \longrightarrow (\Sch/Y) \times_{(\Sch/X)} (\Sch/X')
$$
，它由 $V' \longmapsto (V' \times_{Y'} Y, V' \times_{Y'} X', 1)$ 给出，
并有左伴随函子
$$
G : (\Sch/Y) \times_{(\Sch/X)} (\Sch/X') \longrightarrow (\Sch/Y')
$$
，它把三元组 $(V, U', \varphi)$ 送到推出
$V \amalg_{(V \times_Y X)} U'$。最后，$F \circ G$ 同构于恒等函子。
\end{lemma}

\begin{proof}
设 $(V, U', \varphi)$ 是纤维积范畴的对象。令
$U = U' \times_{X'} X$。注意 $U \to U'$ 是加厚。由于
$\varphi : V \times_Y X \to U' \times_{X'} X = U$ 是同构，
有位于 $X \to Y$ 上方的态射 $U \to V$，它将 $U$ 等同于纤维积
$X \times_Y V$。特别地，$U \to V$ 仿射；见 态射章引理
\ref{morphisms-lemma-base-change-affine}。因此可应用引理
\ref{more-morphisms-lemma-pushout-along-thickening} 得到推出 $V' = V \amalg_U U'$。
由于 $V'$ 是推出，而且给定了在 $U$ 上作为到 $Y'$ 的态射相等的态射
$V \to Y$ 与 $U' \to X'$，遂得到态射 $V' \to Y'$。
令 $G(V, U', \varphi) = V'$，便得到函子 $G$。

\medskip\noindent
下面证明 $G$ 是 $F$ 的左伴随。设 $Z$ 是 $Y'$ 上的概形。须证明
$$
\Mor(V', Z) = \Mor((V, U', \varphi), F(Z))
$$
，其中态射集取自各自的范畴。设 $g' : V' \to Z$ 是态射。
以 $\tilde g$、分别地 $\tilde f'$ 表示 $g'$ 与态射 $V \to V'$、
分别地 $U' \to V'$ 的复合。分别经 $Y \to Y'$、$X' \to Y'$ 对
$\tilde g$、$\tilde f'$ 作基变换，得到态射
$g : V \to Z \times_{Y'} Y$、分别地
$f' : U' \to Z \times_{Y'} X'$。于是 $(g, f')$ 是上述等式右边的元素
（略去细节）。反之，设
$(g, f') : (V, U', \varphi) \to F(Z)$ 是右边的元素。
可考虑 $g$、分别地 $f$ 与 $Z \times_{Y'} X' \to Z$、分别地
$Z \times_{Y'} Y \to Z$ 的复合 $\tilde g : V \to Z$、分别地
$\tilde f' : U' \to Z$。则 $\tilde g$ 与 $\tilde f'$ 作为从 $U$
到 $Z$ 的态射相等。由推出的泛性质，得到态射 $g' : V' \to Z$，
即等式左边的元素。略去这些构造互为逆的验证。
% P05-ERR-CH37-0024: Frozen source swaps f/f' and the two projection targets in the converse composition sentence.

\medskip\noindent
为证明 $F \circ G$ 同构于恒等函子，须证明伴随映射
$(V, U', \varphi) \to F(G(V, U', \varphi))$ 是同构。
为此可在仿射局部运算。写 $X = \Spec(A)$、$X' = \Spec(A')$、
$Y = \Spec(B)$。则 $A' \to A$ 与 $B \to A$ 是 More on Algebra,
引理 \ref{more-algebra-lemma-module-over-fibre-product} 中那样的环映射，
且 $Y' = \Spec(B')$，其中 $B' = B \times_A A'$。再设
$V = \Spec(D)$、$U' = \Spec(C')$，而 $\varphi$ 由 $A$-代数同构
$D \otimes_B A \to C' \otimes_{A'} A = C'/IC'$ 给出。
令 $D' = D \times_{C'/IC'} C'$。在这种情形下，须证明
$D' \otimes_{B'} B \cong D$ 且 $D' \otimes_{B'} A' \cong C'$。
这是 代数进阶章引理
\ref{more-algebra-lemma-module-over-fibre-product} 的特殊情形。
\end{proof}

\begin{lemma}
\label{more-morphisms-lemma-scheme-over-pushout-flat-modules}
设 $X \to X'$ 是概形的加厚，$X \to Y$ 是仿射概形态射，并设
$Y' = Y \amalg_X X'$ 为推出（见引理
\ref{more-morphisms-lemma-pushout-along-thickening}）。设 $V' \to Y'$ 是概形态射。令
$V = Y \times_{Y'} V'$、$U' = X' \times_{Y'} V'$ 且
$U = X \times_{Y'} V'$。以下两者的范畴等价：
\begin{enumerate}
\item 在 $Y'$ 上平坦的拟凝聚 $\mathcal{O}_{V'}$-模；
\item 三元组 $(\mathcal{G}, \mathcal{F}', \varphi)$ 的范畴，其中
\begin{enumerate}
\item $\mathcal{G}$ 是在 $Y$ 上平坦的拟凝聚 $\mathcal{O}_V$-模；
\item $\mathcal{F}'$ 是在 $X'$ 上平坦的拟凝聚 $\mathcal{O}_{U'}$-模；
\item $\varphi : (U \to V)^*\mathcal{G} \to (U \to U')^*\mathcal{F}'$
是 $\mathcal{O}_U$-模的同构。
\end{enumerate}
\end{enumerate}
该等价把 $\mathcal{G}'$ 映到
$((V \to V')^*\mathcal{G}', (U' \to V')^*\mathcal{G}', can)$。
假设 $\mathcal{G}'$ 对应于三元组
$(\mathcal{G}, \mathcal{F}', \varphi)$。则
\begin{enumerate}
\item[(a)] $\mathcal{G}'$ 是有限型 $\mathcal{O}_{V'}$-模，当且仅当
$\mathcal{G}$ 与 $\mathcal{F}'$ 分别是有限型 $\mathcal{O}_Y$-模与
$\mathcal{O}_{U'}$-模。
\item[(b)] 若 $V' \to Y'$ 局部有限表示，则 $\mathcal{G}'$ 是有限表示的
$\mathcal{O}_{V'}$-模，当且仅当 $\mathcal{G}$ 与 $\mathcal{F}'$
分别是有限表示的 $\mathcal{O}_Y$-模与 $\mathcal{O}_{U'}$-模。
\end{enumerate}
\end{lemma}

\begin{proof}
一个拟逆函子把三元组 $(\mathcal{G}, \mathcal{F}', \varphi)$ 映到纤维积
$$
(V \to V')_*\mathcal{G}
\times_{(U \to V')_*\mathcal{F}}
(U' \to V')_*\mathcal{F}'
$$
，其中 $\mathcal{F} = (U \to U')^*\mathcal{F}'$。这是可行的，因为在仿射情形下，
恢复 代数进阶章引理
\ref{more-algebra-lemma-relative-flat-module-over-fibre-product} 的等价。
略去一些细节。

\medskip\noindent
(a)、(b) 由 代数进阶章引理
\ref{more-algebra-lemma-relative-finite-module-over-fibre-product} 与
\ref{more-algebra-lemma-relative-finitely-presented-module-over-fibre-product}
得出。
\end{proof}

\begin{lemma}
\label{more-morphisms-lemma-equivalence-categories-schemes-over-pushout-flat}
在引理 \ref{more-morphisms-lemma-equivalence-categories-schemes-over-pushout} 的情形下，
若对某个三元组 $(V, U', \varphi)$ 有 $V' = G(V, U', \varphi)$，则
\begin{enumerate}
\item $V' \to Y'$ 局部有限型，当且仅当 $V \to Y$ 与 $U' \to X'$
局部有限型；
\item $V' \to Y'$ 平坦，当且仅当 $V \to Y$ 与 $U' \to X'$ 平坦；
\item $V' \to Y'$ 平坦且局部有限表示，当且仅当 $V \to Y$ 与
$U' \to X'$ 平坦且局部有限表示；
\item $V' \to Y'$ 光滑，当且仅当 $V \to Y$ 与 $U' \to X'$ 光滑；
\item $V' \to Y'$ 平展，当且仅当 $V \to Y$ 与 $U' \to X'$ 平展；
\item 视需要在此添加更多性质。
\end{enumerate}
若 $W'$ 在 $Y'$ 上平坦，则伴随映射 $G(F(W')) \to W'$ 是同构。
因此 $F$ 与 $G$ 在以下两个范畴之间定义互为拟逆的函子：在 $Y'$ 上平坦的
概形范畴，以及满足 $V \to Y$ 与 $U' \to X'$ 平坦的三元组
$(V, U', \varphi)$ 的范畴。
\end{lemma}

\begin{proof}
逐个考察仿射片后，本引理的断言等价于 代数进阶章引理
\ref{more-algebra-lemma-properties-algebras-over-fibre-product} 中的相应断言。
\end{proof}







\section{平坦轨迹的开性}
\label{more-morphisms-section-open-flat}

\noindent
这个结果的证明需要一些工作，并且（或许）值得单列一节。结果如下。

\begin{theorem}
\label{more-morphisms-theorem-openness-flatness}
\begin{reference}
\cite[IV Theorem 11.3.1]{EGA}
\end{reference}
设 $S$ 是概形。设 $f : X \to S$ 是局部有限表示的态射。
设 $\mathcal{F}$ 是局部有限表示的拟凝聚 $\mathcal{O}_X$-模。则
$$
U = \{x \in X \mid \mathcal{F}\text{ is flat over }S\text{ 在 }x\}
$$
是 $X$ 的开子集。
\end{theorem}

\begin{proof}
开性可以在 $X$ 上局部检验，故可以假设 $f$ 是仿射概形之间的态射。
此时本定理恰为 代数章定理
\ref{algebra-theorem-openness-flatness}。
\end{proof}

\begin{lemma}
\label{more-morphisms-lemma-flat-locus-base-change}
设 $S$ 是概形。设
$$
\xymatrix{
X' \ar[r]_{g'} \ar[d]_{f'} & X \ar[d]^f \\
S' \ar[r]^g & S
}
$$
是概形的笛卡尔图。设 $\mathcal{F}$ 是拟凝聚 $\mathcal{O}_X$-模。
设 $x' \in X'$ 的像为 $x = g'(x')$ 与 $s' = f'(x')$。
\begin{enumerate}
\item 若 $\mathcal{F}$ 在 $x$ 处于 $S$ 上平坦，则
$(g')^*\mathcal{F}$ 在 $x'$ 处于 $S'$ 上平坦。
\item 若 $g$ 在 $s'$ 处平坦，且 $(g')^*\mathcal{F}$ 在 $x'$ 处于
$S'$ 上平坦，则 $\mathcal{F}$ 在 $x$ 处于 $S$ 上平坦。
\end{enumerate}
特别地，若 $g$ 平坦、$f$ 局部有限表示，且 $\mathcal{F}$ 局部有限表示，
则定理 \ref{more-morphisms-theorem-openness-flatness} 中开子集的形成与基变换交换。
\end{lemma}

\begin{proof}
考虑局部环的交换图
$$
\xymatrix{
\mathcal{O}_{X', x'} & \mathcal{O}_{X, x} \ar[l] \\
\mathcal{O}_{S', s'} \ar[u] & \mathcal{O}_{S, s} \ar[l] \ar[u]
}
$$
注意到 $\mathcal{O}_{X', x'}$ 是
$\mathcal{O}_{X, x} \otimes_{\mathcal{O}_{S, s}} \mathcal{O}_{S', s'}$
的一个局部化，并且 $((g')^*\mathcal{F})_{x'}$ 等于
$\mathcal{F}_x \otimes_{\mathcal{O}_{X, x}} \mathcal{O}_{X', x'}$。
故本引理由 代数章引理
\ref{algebra-lemma-base-change-flat-up-down} 得出。
\end{proof}






\section{按纤维判别平坦性}
\label{more-morphisms-section-criterion-flat-fibres}

\noindent
考虑概形的交换图（左图）
$$
\xymatrix{
X \ar[rr]_f \ar[dr] & & Y \ar[dl] \\
& S
}
\quad
\xymatrix{
X_s \ar[rr]_{f_s} \ar[rd] & & Y_s \ar[dl] \\
& \Spec(\kappa(s))
}
$$
以及拟凝聚 $\mathcal{O}_X$-模 $\mathcal{F}$。
给定落在 $s \in S$ 上方的点 $x \in X$，其像为 $y = f(x)$，
我们考虑如下问题：$\mathcal{F}$ 在 $x$ 处是否于 $Y$ 上平坦？
若 $\mathcal{F}$ 在 $x$ 处于 $S$ 上平坦，则定理表明，这一问题与
$\mathcal{F}$ 在纤维上的限制
$$
\mathcal{F}_s = (X_s \to X)^*\mathcal{F}
$$
在 $x$ 处是否于 $Y_s$ 上平坦密切相关。下面将给出一个“诺特”版本和
一个“有限表示”版本；此前我们还处理过一个“幂零”版本，见引理
\ref{more-morphisms-lemma-flatness-morphism-thickenings}。

\begin{theorem}
\label{more-morphisms-theorem-criterion-flatness-fibre-Noetherian}
设 $S$ 是概形。设 $f : X \to Y$ 是 $S$ 上的概形态射。
设 $\mathcal{F}$ 是拟凝聚 $\mathcal{O}_X$-模。设 $x \in X$，令
$y = f(x)$，并令 $s \in S$ 为 $x$ 在 $S$ 中的像。
假设 $S$、$X$、$Y$ 局部诺特，$\mathcal{F}$ 凝聚，且
$\mathcal{F}_x \not = 0$。则以下条件等价：
\begin{enumerate}
\item $\mathcal{F}$ 在 $x$ 处于 $S$ 上平坦，且
$\mathcal{F}_s$ 在 $x$ 处于 $Y_s$ 上平坦；
\item $Y$ 在 $y$ 处于 $S$ 上平坦，且 $\mathcal{F}$ 在 $x$ 处于
$Y$ 上平坦。
\end{enumerate}
\end{theorem}

\begin{proof}
考虑环映射
$$
\mathcal{O}_{S, s} \longrightarrow
\mathcal{O}_{Y, y} \longrightarrow \mathcal{O}_{X, x}
$$
以及模 $\mathcal{F}_x$。$\mathcal{F}_s$ 在 $x$ 处的茎为模
$\mathcal{F}_x/\mathfrak m_s \mathcal{F}_x$，而 $Y_s$ 在 $y$ 处的局部环为
$\mathcal{O}_{Y, y}/\mathfrak m_s \mathcal{O}_{Y, y}$。
因此蕴含 (1) $\Rightarrow$ (2) 即 代数章引理
\ref{algebra-lemma-criterion-flatness-fibre-Noetherian}。
若 (2) 成立，则第一个环映射忠实平坦，且 $\mathcal{F}_x$ 在
$\mathcal{O}_{Y, y}$ 上平坦，故由 代数章引理
\ref{algebra-lemma-composition-flat}，$\mathcal{F}_x$ 在
$\mathcal{O}_{S, s}$ 上平坦。此外，
$\mathcal{F}_x/\mathfrak m_s \mathcal{F}_x$ 是平坦模 $\mathcal{F}_x$
沿 $\mathcal{O}_{Y, y} \to \mathcal{O}_{Y, y}/\mathfrak m_s \mathcal{O}_{Y, y}$
的基变换，故由 代数章引理
\ref{algebra-lemma-flat-base-change}，它也是平坦的。
\end{proof}

\noindent
下面是非诺特版本。

\begin{theorem}
\label{more-morphisms-theorem-criterion-flatness-fibre}
设 $S$ 是概形。设 $f : X \to Y$ 是 $S$ 上的概形态射。
设 $\mathcal{F}$ 是拟凝聚 $\mathcal{O}_X$-模。假设
\begin{enumerate}
\item $X$ 在 $S$ 上局部有限表示；
\item $\mathcal{F}$ 是有限表示的 $\mathcal{O}_X$-模；
\item $Y$ 在 $S$ 上局部有限型。
\end{enumerate}
设 $x \in X$，令 $y = f(x)$，并令 $s \in S$ 为 $x$ 在 $S$ 中的像。
若 $\mathcal{F}_x \not = 0$，则以下条件等价：
\begin{enumerate}
\item $\mathcal{F}$ 在 $x$ 处于 $S$ 上平坦，且
$\mathcal{F}_s$ 在 $x$ 处于 $Y_s$ 上平坦；
\item $Y$ 在 $y$ 处于 $S$ 上平坦，且 $\mathcal{F}$ 在 $x$ 处于
$Y$ 上平坦。
\end{enumerate}
此外，使 (1) 与 (2) 成立的点 $x$ 所成的集合在
$\text{Supp}(\mathcal{F})$ 中是开的。
\end{theorem}

\begin{proof}
考虑环映射
$$
\mathcal{O}_{S, s} \longrightarrow
\mathcal{O}_{Y, y} \longrightarrow \mathcal{O}_{X, x}
$$
以及模 $\mathcal{F}_x$。$\mathcal{F}_s$ 在 $x$ 处的茎为模
$\mathcal{F}_x/\mathfrak m_s \mathcal{F}_x$，而 $Y_s$ 在 $y$ 处的局部环为
$\mathcal{O}_{Y, y}/\mathfrak m_s \mathcal{O}_{Y, y}$。
因此蕴含 (1) $\Rightarrow$ (2) 即 代数章引理
\ref{algebra-lemma-criterion-flatness-fibre-fp-over-ft}。
若 (2) 成立，则第一个环映射忠实平坦，且 $\mathcal{F}_x$ 在
$\mathcal{O}_{Y, y}$ 上平坦，故由 代数章引理
\ref{algebra-lemma-composition-flat}，$\mathcal{F}_x$ 在
$\mathcal{O}_{S, s}$ 上平坦。此外，
$\mathcal{F}_x/\mathfrak m_s \mathcal{F}_x$ 是平坦模 $\mathcal{F}_x$
沿 $\mathcal{O}_{Y, y} \to \mathcal{O}_{Y, y}/\mathfrak m_s \mathcal{O}_{Y, y}$
的基变换，故由 代数章引理
\ref{algebra-lemma-flat-base-change}，它也是平坦的。

\medskip\noindent
由 态射章引理
\ref{morphisms-lemma-finite-presentation-permanence}，态射 $f$
局部有限表示。考虑集合
\begin{equation}
\label{more-morphisms-equation-open}
U = \{x \in X \mid \mathcal{F} \text{ flat at }x
\text{ over both }Y\text{ 且 }S\}.
\end{equation}
由定理 \ref{more-morphisms-theorem-openness-flatness}，此集合在 $X$ 中是开的。
注意，若 $x \in U$，则 $\mathcal{F}_s$ 是平坦模沿态射
$Y_s \to Y$ 的基变换，因而在 $x$ 处于 $Y_s$ 上平坦；见
态射章引理 \ref{morphisms-lemma-base-change-module-flat}。
故 $U \cap \text{Supp}(\mathcal{F})$ 的每一点都满足条件 (1)。
另一方面，显然若 $x \in \text{Supp}(\mathcal{F})$ 满足 (1) 与 (2)，
则 $x \in U$。因此所求开集为 $U \cap \text{Supp}(\mathcal{F})$。
\end{proof}

\noindent
这些定理常以如下简化形式使用。这里只给出整体陈述——当然也有逐点版本。

\begin{lemma}
\label{more-morphisms-lemma-morphism-between-flat-Noetherian}
设 $S$ 是概形。设 $f : X \to Y$ 是 $S$ 上的概形态射。假设
\begin{enumerate}
\item $S$、$X$、$Y$ 局部诺特；
\item $X$ 在 $S$ 上平坦；
\item 对每个 $s \in S$，态射 $f_s : X_s \to Y_s$ 平坦。
\end{enumerate}
则 $f$ 平坦。若 $f$ 还满射，则 $Y$ 在 $S$ 上平坦。
\end{lemma}

\begin{proof}
这是定理 \ref{more-morphisms-theorem-criterion-flatness-fibre-Noetherian} 的特殊情形。
\end{proof}

\begin{lemma}
\label{more-morphisms-lemma-morphism-between-flat}
设 $S$ 是概形。设 $f : X \to Y$ 是 $S$ 上的概形态射。假设
\begin{enumerate}
\item $X$ 在 $S$ 上局部有限表示；
\item $X$ 在 $S$ 上平坦；
\item 对每个 $s \in S$，态射 $f_s : X_s \to Y_s$ 平坦；
\item $Y$ 在 $S$ 上局部有限型。
\end{enumerate}
则 $f$ 平坦。若 $f$ 还满射，则 $Y$ 在 $S$ 上平坦。
\end{lemma}

\begin{proof}
这是定理 \ref{more-morphisms-theorem-criterion-flatness-fibre} 的特殊情形。
\end{proof}

\begin{lemma}
\label{more-morphisms-lemma-base-change-criterion-flatness-fibre}
设 $S$ 是概形。设 $f : X \to Y$ 是 $S$ 上的概形态射。
设 $\mathcal{F}$ 是拟凝聚 $\mathcal{O}_X$-模。假设
\begin{enumerate}
\item $X$ 在 $S$ 上局部有限表示；
\item $\mathcal{F}$ 是有限表示的 $\mathcal{O}_X$-模；
\item $\mathcal{F}$ 在 $S$ 上平坦；
\item $Y$ 在 $S$ 上局部有限型。
\end{enumerate}
则集合
$$
U = \{x \in X \mid \mathcal{F} \text{ flat at }x \text{ 相对于 }Y\}.
$$
在 $X$ 中是开的，且其形成与任意基变换交换：若 $S' \to S$ 是概形态射，
而 $U'$ 是 $X' = X \times_S S'$ 中使
$\mathcal{F}' = \mathcal{F} \times_S S'$ 在 $Y' = Y \times_S S'$ 上平坦
的点集，则 $U' = U \times_S S'$。
\end{lemma}

\begin{proof}
由 态射章引理
\ref{morphisms-lemma-finite-presentation-permanence}，态射 $f$
局部有限表示。故由定理 \ref{more-morphisms-theorem-openness-flatness}，$U$ 是开的。
由于已假设 $\mathcal{F}$ 在 $S$ 上平坦，定理
\ref{more-morphisms-theorem-criterion-flatness-fibre} 给出
$$
U = \{x \in X \mid \mathcal{F}_s \text{ flat at }x \text{ 相对于 }Y_s\}.
$$
其中 $s$ 始终表示 $x$ 在 $S$ 中的像。（当 $\mathcal{F}_x = 0$ 时，
这一描述也平凡成立。）此外，本引理的假设对态射 $f' : X' \to Y'$
与层 $\mathcal{F}'$ 仍然成立，故 $U'$ 有类似的描述。换言之，只需证明：
给定映到 $s \in S$ 的 $s' \in S'$，集合
$$
\{x' \in X'_{s'} \mid
\mathcal{F}'_{s'} \text{ flat at }x' \text{ 相对于 }Y'_{s'}\}
$$
是 $X_s$ 中相应轨迹的逆像。由引理
\ref{more-morphisms-lemma-flat-locus-base-change}，这确实成立，因为在笛卡尔图
$$
\xymatrix{
X'_{s'} \ar[d] \ar[r] & X_s \ar[d] \\
Y'_{s'} \ar[r] & Y_s
}
$$
中，水平态射是平坦态射
$\Spec(\kappa(s')) \to \Spec(\kappa(s))$ 的基变换，因而平坦。
\end{proof}

\begin{lemma}
\label{more-morphisms-lemma-base-change-flatness-fibres}
设 $S$ 是概形。设 $f : X \to Y$ 是 $S$ 上的概形态射。假设
\begin{enumerate}
\item $X$ 在 $S$ 上局部有限表示；
\item $X$ 在 $S$ 上平坦；
\item $Y$ 在 $S$ 上局部有限型。
\end{enumerate}
则集合
$$
U = \{x \in X \mid X\text{ flat at }x \text{ 相对于 }Y\}.
$$
在 $X$ 中是开的，且其形成与任意基变换交换。
\end{lemma}

\begin{proof}
这是引理 \ref{more-morphisms-lemma-base-change-criterion-flatness-fibre} 的特殊情形。
\end{proof}

\noindent
下一个引理是 代数章引理
\ref{algebra-lemma-free-fibre-flat-free} 的一个变体。注意，假设
$(\mathcal{F}_s)_x$ 是平坦 $\mathcal{O}_{X_s, x}$-模，意味着
$(\mathcal{F}_s)_x$ 是自由 $\mathcal{O}_{X_s, x}$-模。若 $x \in X_s$
是 $X_s$ 某个不可约分支的泛点且 $X_s$ 约化，则这一点总是成立
（此时 $\mathcal{O}_{X_s, x}$ 是域；见 代数章引理
\ref{algebra-lemma-minimal-prime-reduced-ring}）。

\begin{lemma}
\label{more-morphisms-lemma-flat-and-free-at-point-fibre}
设 $f : X \to S$ 是有限表示的概形态射，$\mathcal{F}$ 是有限表示的
$\mathcal{O}_X$-模。设 $x \in X$ 的像为 $s \in S$。若 $\mathcal{F}$
在 $x$ 处于 $S$ 上平坦，且 $(\mathcal{F}_s)_x$ 是平坦
$\mathcal{O}_{X_s, x}$-模，则 $\mathcal{F}$ 在 $x$ 的某个邻域上有限自由。
\end{lemma}

\begin{proof}
若 $\mathcal{F}_x \otimes \kappa(x)$ 为零，则由 Nakayama 引理
（代数章引理 \ref{algebra-lemma-NAK}），$\mathcal{F}_x = 0$，
因而 $\mathcal{F}$ 在 $x$ 的某个邻域上为零（模章引理
\ref{modules-lemma-finite-type-stalk-zero}），此时结论成立。因此可假设
$\mathcal{F}_x \otimes \kappa(x)$ 非零，并将定理
\ref{more-morphisms-theorem-criterion-flatness-fibre} 应用于 $f = \text{id} : X \to X$。
由此得出 $\mathcal{F}_x$ 在 $\mathcal{O}_{X, x}$ 上平坦，故
$\mathcal{F}_x$ 自由；例如见 代数章引理
\ref{algebra-lemma-finite-flat-local}。选择开邻域 $x \in U \subset X$
以及截面 $s_1, \ldots, s_r \in \mathcal{F}(U)$，使它们在
$\mathcal{F}_x$ 中的像构成一组基。缩小 $U$ 后，相应映射
$\psi : \mathcal{O}_U^{\oplus r} \to \mathcal{F}|_U$ 是满射
（模章引理 \ref{modules-lemma-finite-type-stalk-zero}）。于是
$\Ker(\psi)$ 有限型（见 模章引理
\ref{modules-lemma-kernel-surjection-finite-free-onto-finite-presentation}），
且 $\Ker(\psi)_x = 0$。再次缩小 $U$ 后，$\psi$ 即为同构。
\end{proof}

\begin{lemma}
\label{more-morphisms-lemma-finite-free-open}
设 $f : X \to S$ 是局部有限表示的概形态射。设 $\mathcal{F}$ 是
在 $S$ 上平坦的有限表示 $\mathcal{O}_X$-模。则集合
$$
\{x \in X : \mathcal{F}\text{ free in a neighbourhood of }x\}
$$
在 $X$ 中是开的，且其形成与任意基变换 $S' \to S$ 交换。
\end{lemma}

\begin{proof}
开性是平凡的。设 $x \in X$ 映到 $s \in S$。由引理
\ref{more-morphisms-lemma-flat-and-free-at-point-fibre}，$x$ 属于上述集合，当且仅当
$\mathcal{F}|_{X_s}$ 在 $x$ 处于 $X_s$ 上平坦。显然这又等价于
$\mathcal{F}$ 在 $x$ 处于 $X$ 上平坦（因为 $\mathcal{F}_x$ 的自由性
蕴含这一陈述，而这一陈述又蕴含 $\mathcal{F}|_{X_s}$ 在 $x$ 处于
$X_s$ 上平坦）。因此，将引理
\ref{more-morphisms-lemma-base-change-criterion-flatness-fibre} 应用于 $S$ 上的
$\text{id} : X \to X$，即得基变换陈述。
\end{proof}






\section{光滑概形之间的闭浸入}
\label{more-morphisms-section-closed-immersions-smooth}

\noindent
这里收录一些不太适合放在其他地方的结果。

\begin{lemma}
\label{more-morphisms-lemma-etale-local-structure}
设 $S$ 是概形。设 $Y \to X$ 是在 $S$ 上光滑的概形之间的闭浸入。
对每个 $y \in Y$，存在整数 $0 \leq m, n$ 以及交换图
$$
\xymatrix{
Y \ar[d] &
V \ar[l] \ar[d] \ar[r] &
\mathbf{A}^m_S
\ar[d]^{(a_1, \ldots, a_m) \mapsto (a_1, \ldots, a_m, 0 \ldots, 0)} \\
X &
U \ar[l] \ar[r]^-\pi &
\mathbf{A}^{m + n}_S
}
$$
其中 $U \subset X$ 是开集，$V = Y \cap U$，$\pi$ 平展，
$V = \pi^{-1}(\mathbf{A}^m_S)$，且 $y \in V$。
\end{lemma}

\begin{proof}
这个问题在 $X$ 上是局部的，故可以用 $y$ 的一个开邻域替换 $X$。
由 除子章引理
\ref{divisors-lemma-immersion-smooth-into-smooth-regular-immersion}，
$Y \to X$ 是正则浸入，因而可以假设 $X = \Spec(A)$ 仿射，并存在正则序列
$f_1, \ldots, f_n \in A$，使得 $Y = V(f_1, \ldots, f_n)$。
进一步缩小 $X$（从而也缩小 $Y$）后，可以假设存在平展态射
$Y \to \mathbf{A}^m_S$；见 态射章引理
\ref{morphisms-lemma-smooth-etale-over-affine-space}。设
$\mathcal{O}_Y(Y)$ 中的 $\overline{g}_1, \ldots, \overline{g}_m$
是此平展态射的坐标函数。选取这些函数的提升
$g_1, \ldots, g_m \in A$，并考虑态射
$$
(g_1, \ldots, g_m, f_1, \ldots, f_n) :
X
\longrightarrow
\mathbf{A}^{m + n}_S
$$
于 $S$ 上。两边概形均在 $S$ 上局部有限表示，故此态射局部有限表示
（态射章引理
\ref{morphisms-lemma-finite-presentation-permanence}）。按构造，此态射在
$\mathbf{A}^m_S \subset \mathbf{A}^{m + n}_S$ 上的限制平展。因此，为证明
$X \to \mathbf{A}^{m + n}_S$ 在 $y$ 处平展，只需证明
$X \to \mathbf{A}^{m + n}_S$ 在 $y$ 处平坦；见 态射章引理
\ref{morphisms-lemma-etale-at-point}。设 $s \in S$ 为 $y$ 的像。由定理
\ref{more-morphisms-theorem-criterion-flatness-fibre}，只需检验
$X_s \to \mathbf{A}^{m + n}_s$ 在 $y$ 处平坦。设
$z \in \mathbf{A}^{m + n}_s$ 为 $y$ 的像。局部环映射
$$
\mathcal{O}_{\mathbf{A}^{m + n}_s, z}
\longrightarrow
\mathcal{O}_{X_s, y}
$$
由 代数章引理 \ref{algebra-lemma-CM-over-regular-flat} 是平坦的。
事实上，域上的光滑概形是正则的，而正则环是 Cohen--Macaulay 环；见
簇章引理 \ref{varieties-lemma-smooth-regular} 与 代数章引理
\ref{algebra-lemma-regular-ring-CM}。因此源与靶都是正则局部环
（从而都是 CM 环）。源与靶的维数相同：由 代数进阶章引理
\ref{more-algebra-lemma-dimension-etale-extension}，有
$\dim(\mathcal{O}_{Y_s, y}) = \dim(\mathcal{O}_{\mathbf{A}^m_s, z})$
；又有 $\dim(\mathcal{O}_{\mathbf{A}^{m + n}_s, z}) = n +
\dim(\mathcal{O}_{\mathbf{A}^m_s, z})$，并且
$\dim(\mathcal{O}_{X_s, y}) = n + \dim(\mathcal{O}_{Y_s, y})$
，因为 $\mathcal{O}_{Y_s, y}$ 是 $\mathcal{O}_{X_s, y}$ 对长度为 $n$
的正则序列 $f_1, \ldots, f_n$ 的商（见 除子章说明
\ref{divisors-remark-relative-regular-immersion-elements}）。最后，由于
$Y_s \to \mathbf{A}^m_s$ 在 $y$ 处平展，上述箭头的纤维环在
$\kappa(z)$ 上有限。证明完毕。
\end{proof}

\begin{remark}
\label{more-morphisms-remark-relative-blowup}
固定一个环 $R$，并令 $S = \Spec(R)$。固定整数 $0 \leq m$ 与
$1 \leq n$。考虑闭浸入
$$
Z = \mathbf{A}^m_S \longrightarrow \mathbf{A}^{m + n}_S = X,\quad
(a_1, \ldots, a_m) \mapsto (a_1, \ldots, a_m, 0, \ldots 0).
$$
我们考察 $X$ 沿闭子概形 $Z$ 的吹起 $X'$。写成
$$
X = \Spec(A)
\quad\text{其中}\quad
A = R[x_1, \ldots, x_m, y_1, \ldots, y_n]
$$
则 $X'$ 是 $A$ 关于理想 $(y_1, \ldots, y_n)$ 的 Rees 代数的 Proj。
此 Rees 代数等于 $B = A[T_1, \ldots, T_n]/(y_iT_j - y_jT_i)$；略去细节。
因此 $X' = \text{Proj}(B)$ 在 $S$ 上光滑，因为它由仿射开集
\begin{align*}
D_+(T_i)
& =
\Spec(B_{(T_i)}) \\
& =
\Spec(A[t_1, \ldots, \hat t_i, \ldots t_n]/(y_j - y_i t_j)) \\
& =
\Spec(R[x_1, \ldots, x_m, y_i, t_1, \ldots, \hat t_i, \ldots, t_n])
\end{align*}
覆盖，而这些开集同构于 $\mathbf{A}^{n + m}_S$。在此坐标图中，
例外除子由令 $y_i = 0$ 截出，故例外除子也在 $S$ 上光滑。
\end{remark}

\begin{lemma}
\label{more-morphisms-lemma-blowup}
设 $S$ 是概形。设 $Z \to X$ 是在 $S$ 上光滑的概形之间的闭浸入。
设 $b : X' \to X$ 是沿 $Z$ 的吹起，例外除子为 $E \subset X'$。
则 $X'$ 与 $E$ 在 $S$ 上光滑。态射 $p : E \to Z$ 典范同构于射影空间丛
$$
\mathbf{P}(\mathcal{I}/\mathcal{I}^2) \longrightarrow Z
$$
，其中 $\mathcal{I} \subset \mathcal{O}_X$ 是 $Z$ 的理想层。
来自射影空间丛结构的相对 $\mathcal{O}_E(1)$ 同构于
$\mathcal{O}_{X'}(-E)$ 在 $E$ 上的限制。
\end{lemma}

\begin{proof}
% P05-ERR-0025: frozen source has typographical "immmersion".
由 除子章引理
\ref{divisors-lemma-immersion-smooth-into-smooth-regular-immersion}，
浸入 $Z \to X$ 是正则浸入，故理想层 $\mathcal{I}$ 有限型，从而 $b$
是射影态射，并带有相对丰沛可逆层
$\mathcal{O}_{X'}(1) = \mathcal{O}_{X'}(-E)$；见 除子章引理
\ref{divisors-lemma-blowing-up-gives-effective-Cartier-divisor} 及
\ref{divisors-lemma-blowing-up-projective}.
典范映射 $\mathcal{I} \to b_*\mathcal{O}_{X'}(1)$ 按照吹起的构造给出闭浸入
$$
X' \longrightarrow
\mathbf{P}\left(\bigoplus\nolimits_{n \geq 0}
\text{Sym}^n_{\mathcal{O}_X}(\mathcal{I})\right)
$$
。此态射在 $E$ 上的限制给出 $Z$ 上的典范映射
$$
E \longrightarrow
\mathbf{P}\left(\bigoplus\nolimits_{n \geq 0}
\text{Sym}^n_{\mathcal{O}_Z}(\mathcal{I}/\mathcal{I}^2)\right)
$$
。由于 $\mathcal{I}/\mathcal{I}^2$ 有限局部自由，若此典范映射是同构，
则引理的最后一部分成立。说完这些之后，问题现在在 $X$ 上是平展局部的。
事实上，由 除子章引理
\ref{divisors-lemma-flat-base-change-blowing-up}，吹起与平坦基变换交换；
而且在预合成一个满平展态射之后可以检验光滑性。因此，利用引理
\ref{more-morphisms-lemma-etale-local-structure} 所给出的 $S$ 上概形闭浸入的平展局部结构，
问题归结为评注 \ref{more-morphisms-remark-relative-blowup} 中讨论的情形。
\end{proof}








\section{平坦模与相对伴随点}
\label{more-morphisms-section-flat-relative-assassin}

\noindent
本节将证明：在几何情形下，平坦模的支集在基底上（在某种意义下）
是等维的。诺特情形可参见 \cite[IV Proposition 12.1.1.5]{EGA}。
首先证明两个辅助引理。

\begin{lemma}
\label{more-morphisms-lemma-mod-injective-valuation-ring}
% P05-ERR-0026: frozen source has ungrammatical "Let A -> B is".
设 $A$ 是赋值环。设 $A \to B$ 是局部环之间的局部同态，且本质有限型。
设 $u : N \to M$ 是有限 $B$-模之间的映射。假设 $M$ 在 $A$ 上平坦，且
$\overline{u} : N/\mathfrak m_A N \to M/\mathfrak m_A M$ 是单射。
则 $u$ 是单射，且 $M/u(N)$ 在 $A$ 上平坦。
\end{lemma}

\begin{proof}
我们将从 代数章引理 \ref{algebra-lemma-mod-injective-general}
推出本引理（请注意，我们交换了 $M$ 与 $N$ 的角色）。为完成化约，
将使用 代数进阶章引理
\ref{more-algebra-lemma-valuation-ring-flat-essentially-finite-type}
把问题化到有限表示情形。

\medskip\noindent
由假设，可以把 $B$ 写成多项式代数 $P = A[x_1, \ldots, x_n]$
在素理想 $\mathfrak q$ 处的局部化的一个商。于是可以把
$u : N \to M$ 看成有限 $P_\mathfrak q$-模之间的映射。因此不妨并且确实
假设 $B$ 在 $A$ 上本质有限表示。

\medskip\noindent
其次，$B$-模 $N$ 有限，但未必有限表示。将 $N = \colim N_\lambda$
写成有限表示 $B$-模的滤过余极限，其转移映射均为满射。例如，选择表示
$0 \to K \to B^{\oplus r} \to N \to 0$，把 $K$ 写成其有限子模
$K_\lambda$ 的并，并令 $N_\lambda = \Coker(K_\lambda \to B^{\oplus r})$。
模 $N/\mathfrak m_A N$ 在诺特环 $B/\mathfrak m_A B$ 上有限表示。
故对充分大的 $\lambda$，有
$N_\lambda/\mathfrak m_A N_\lambda = N/\mathfrak m_A N$。
现在，若能对复合映射 $u_\lambda : N_\lambda \to M$ 证明本引理，
则可推出 $N_\lambda = N$，并且结论对 $u$ 成立。因此不妨并且确实假设
$N$ 在 $B$ 上有限表示。

\medskip\noindent
由 代数进阶章引理
\ref{more-algebra-lemma-valuation-ring-flat-essentially-finite-type}，模 $M$
在 $B$ 上有限表示。因此 代数章引理
\ref{algebra-lemma-mod-injective-general} 的所有假设都成立，结论随之得出。
\end{proof}

\begin{lemma}
\label{more-morphisms-lemma-make-base-change}
\begin{reference}
这可见于 \cite[IV Proposition 12.1.1.5]{EGA} 的证明
\end{reference}
设 $f : X \to S$ 是概形态射。设点 $y \in X$ 的像为 $t \in S$。
以 $Y \subset X$ 表示 $\{y\}$ 的闭包，并赋予它 $X$ 的整闭子概形结构。
设 $s \in S$，并设 $x \in Y_s$ 是 $Y_s$ 某个不可约分支的泛点。
存在笛卡尔图
$$
\xymatrix{
X' \ar[r]_{g'} \ar[d]_{f'} & X \ar[d]^f \\
S' \ar[r]^g & S
}
$$
满足以下性质：
\begin{enumerate}
\item $S'$ 是某个赋值环的谱，其泛点为 $t'$、闭点为 $s'$；
\item $g(t') = t$ 且 $g(s') = s$；
\item 存在点 $y' \in X'_{t'}$，它是
$(S' \times_S Y)_{t'} = Y_t \times_t t'$ 某个不可约分支的泛点，
并满足 $g'(y') = y$；
\item 以 $Y' \subset X'$ 表示 $\{y'\}$ 的闭包，并赋予它 $X'$
的整闭子概形结构；则存在点 $x' \in Y'_{s'}$，它是
$Y'_{s'}$ 某个不可约分支的泛点，并满足 $g'(x') = x$。
\end{enumerate}
\end{lemma}

\begin{proof}
选择赋值环 $R$，令 $S' = \Spec(R)$，其泛点为 $t'$、闭点为 $s'$；
再选择满足 $h(t') = y$ 与 $h(s') = x$ 的态射 $h : S' \to X$。
见 概形章引理 \ref{schemes-lemma-points-specialize}。令
$g = f \circ h$，于是 $g(t') = t$ 且 $g(s') = s$。考虑基变换
$$
\xymatrix{
X' \ar[r]_{g'} \ar[d] & X \ar[d] \\
S' \ar@/^1em/[u]^\sigma \ar[r]^-g & S
}
$$
得到此基变换的一个截面 $\sigma$，满足 $h = g' \circ \sigma$。

\medskip\noindent
由于 $h$ 经由 $Y$ 分解，$\sigma$ 当然经由 $Y$ 的基变换
$S' \times_S Y$ 分解。设 $y' \in X'_{t'} \subset X'$ 是纤维
$$
(S' \times_S Y)_{t'} = Y_t \times_t t'
$$
的某个不可约分支的泛点，该分支包含点 $\sigma(t')$；也就是说，
$y' \leadsto \sigma(t')$。
% P05-ERR-0027: frozen source applies g:S'->S to points of X'; likely g'.
由于 $g'(y') \in Y_t$ 且 $g(y') \leadsto g(\sigma(t')) = y$，而 $y$
是纤维 $Y_t$ 的泛点，故 $g'(y') = y$。以 $Y' \subset X'$ 表示
$\{y'\}$ 在 $X'$ 中的闭包，并赋予它整闭子概形结构。由于
$\sigma(t') \in Y'$，$\sigma$ 经由 $Y'$ 分解。选择 $Y'_{s'}$
某个包含 $\sigma(s')$ 的不可约分支的泛点 $x' \in Y'_{s'}$；也就是说，
$x' \leadsto \sigma(s')$，从而 $g'(x') \leadsto g'(\sigma(s')) = x$。
再次利用假设中 $x$ 是 $Y_s$ 某个不可约分支的泛点，以及
$g'(Y') \subset Y$，得出 $g'(x') = x$。
\end{proof}

\begin{lemma}
\label{more-morphisms-lemma-associated-point-specializes}
设 $f : X \to S$ 是局部有限型的概形态射，$\mathcal{F}$ 是有限型拟凝聚
$\mathcal{O}_X$-模。设 $y \in \text{Ass}_{X/S}(\mathcal{F})$ 的像为
$t \in S$。以 $Y \subset X$ 表示 $\{y\}$ 在 $X$ 中的闭包，并赋予它
整闭子概形结构。设 $s \in S$，并设 $x \in Y_s$ 是 $Y_s$
某个不可约分支的泛点。若 $\mathcal{F}$ 在 $x$ 处于 $S$ 上平坦，则
$x \in \text{Ass}_{X/S}(\mathcal{F})$，且
$\dim_x(Y_s) = \dim(Y_t)$。
\end{lemma}

\begin{proof}
选择引理 \ref{more-morphisms-lemma-make-base-change} 中的图，并令
$\mathcal{F}' = (g')^*\mathcal{F}$。除子章引理
\ref{divisors-lemma-base-change-relative-assassin} 推出
$y' \in \text{Ass}_{X'/S'}(\mathcal{F}')$。由 $y'$ 的选取还可得
$\dim(Y'_{t'}) = \dim(Y_t)$；例如见 代数章引理
\ref{algebra-lemma-inequalities-under-field-extension}。由 代数章引理
\ref{algebra-lemma-finite-type-domain-over-valuation-ring-dim-fibres}，
$Y'_{s'}$ 等维，且其维数等于 $\dim(Y_t)$。由于 $\mathcal{F}$ 在 $x$
处于 $S$ 上平坦，$\mathcal{F}'$ 在 $x'$ 处于 $S'$ 上平坦；见
态射章引理 \ref{morphisms-lemma-base-change-module-flat}。

\medskip\noindent
% P05-ERR-0028: frozen source uses X'/S here, while the surrounding base is S'.
假设能够证明 $x' \in \text{Ass}_{X'/S}(\mathcal{F}')$。则 除子章引理
\ref{divisors-lemma-base-change-relative-assassin} 推出
$x \in \text{Ass}_{X/S}(\mathcal{F})$，并且 $Y'_{s'}$ 中包含 $x'$
的不可约分支 $C'$ 是 $C \times_s s'$ 的一个不可约分支，其中
$C \subset Y_s$ 是包含 $x$ 的不可约分支。因此
$\dim(C) = \dim(C') = \dim(Y_t)$（见上文），证明完成。
这把问题化到下一段讨论的情形。

\medskip\noindent
假设 $S = \Spec(A)$，其中 $A$ 是赋值环，而 $t$ 与 $s$ 分别是 $S$
的泛点与闭点。为导出矛盾，假设
$x \not \in \text{Ass}_{X/S}(\mathcal{F})$。换言之，假设
$x \not \in \text{Ass}_{X_s}(\mathcal{F}_s)$，其中 $\mathcal{F}_s$
是 $\mathcal{F}$ 在 $X_s$ 上的拉回。考虑环映射
$$
A \longrightarrow \mathcal{O}_{X, x} = B
$$
以及 $B = \mathcal{O}_{X, x}$ 上的模 $N = \mathcal{F}_x$。于是
$B/\mathfrak m_A B = \mathcal{O}_{X_s, x}$，而 $N/\mathfrak m_A N$
是 $\mathcal{F}_s$ 在点 $x$ 处的茎。以 $\mathfrak q \subset B$
表示点 $y$ 对应的素理想；见 概形章引理
\ref{schemes-lemma-specialize-points}。由于 $x$ 是 $Y_s$ 的泛点，
$\mathfrak q + \mathfrak m_A B$ 的根为 $\mathfrak m_B$。集合
$\text{Ass}_{B/\mathfrak m_A B}(N/\mathfrak m_A N)$ 是有限素理想集
（代数章引理 \ref{algebra-lemma-finite-ass}），并且由于
$x \not \in \text{Ass}_{X/S}(\mathcal{F})$，它不含
$B/\mathfrak m_A B$ 的极大理想。
% P05-ERR-0029: frozen source has duplicated word "of of".
因此 $\mathfrak q$ 在 $B/\mathfrak m_A B$ 中的像不包含于这些素理想中的
任何一个。由素理想回避（代数章引理 \ref{algebra-lemma-silly}），
可以找到元素 $g \in \mathfrak q$，使其在 $B/\mathfrak m_A B$ 中的像是
$N/\mathfrak m_A N$ 上的非零因子（这里使用 代数章引理
\ref{algebra-lemma-ass-zero-divisors} 对零因子的描述）。由于
$N = \mathcal{F}_x$ 在 $A$ 上平坦，引理
\ref{more-morphisms-lemma-mod-injective-valuation-ring} 表明
$$
g : N \longrightarrow N
$$
是单射。特别地，若 $K = \text{Frac}(A)$ 是 $A$ 的分式域，则
$$
g : N \otimes_A K \longrightarrow N \otimes_A K
$$
是单射。注意，$\mathfrak q$ 对应于 $B \otimes_A K$ 的一个素理想。
以 $\mathcal{F}_t$ 表示 $\mathcal{F}$ 在泛纤维 $X_t$ 上的限制。我们有
$(B \otimes_A K)_{\mathfrak q} = \mathcal{O}_{X_t, y}$，且
$(N \otimes_A K)_\mathfrak q$ 是 $\mathcal{F}_t$ 在 $y$ 处的茎。
因此 $g \in \mathfrak m_y \subset \mathcal{O}_{X_t, y}$ 是茎
$(\mathcal{F}_t)_y$ 上的非零因子，这与假设
$y \in \text{Ass}_{X/S}(\mathcal{F})$ 矛盾。
\end{proof}

\begin{lemma}
\label{more-morphisms-lemma-relative-dimension-support-flat}
设 $f : X \to S$ 是局部有限型的概形态射。设 $\mathcal{F}$ 是在 $S$
上平坦的有限型拟凝聚 $\mathcal{O}_X$-模。假设 $S$ 不可约，泛点为
$\eta$。若 $\dim(\text{Supp}(\mathcal{F}_\eta)) \leq r$，则对所有
$s \in S$ 都有 $\dim(\text{Supp}(\mathcal{F}_s)) \leq r$。
\end{lemma}

\begin{proof}
设 $x \in \text{Supp}(\mathcal{F}_s)$ 是
$\text{Supp}(\mathcal{F}_s)$ 某个不可约分支的泛点。由 代数章引理
\ref{algebra-lemma-going-down-flat-module}，可以在
$\text{Supp}(\mathcal{F})$ 中找到特化 $y \leadsto x$，满足
$f(y) = \eta$。当然可以假设 $y$ 是
$\text{Supp}(\mathcal{F}_\eta)$ 某个不可约分支的泛点。由引理
\ref{more-morphisms-lemma-associated-point-specializes}，$\overline{\{x\}}$ 的维数至多为
$r$。
\end{proof}

\begin{lemma}
\label{more-morphisms-lemma-flat-associated-equidimensional}
设 $f : X \to S$ 是局部有限型的概形态射，$\mathcal{F}$ 是有限型拟凝聚
$\mathcal{O}_X$-模。设 $y \in \text{Ass}_{X/S}(\mathcal{F})$。
以 $Y \subset X$ 表示 $\{y\}$ 在 $X$ 中的闭包，并赋予它整闭子概形结构；
以 $T \subset S$ 表示 $\{f(y)\}$ 的闭包，并赋予它整闭子概形结构。
得到交换图
$$
\xymatrix{
Y \ar[r] \ar[d] & X \ar[d] \\
T \ar[r] & S
}
$$
，其中 $Y \to T$ 是支配态射。假设在 $Y \to T$ 的各纤维的所有不可约
分支的泛点处，$\mathcal{F}$ 都在 $S$ 上平坦（例如，若
$\mathcal{F}$ 在 $S$ 上平坦）。则
\begin{enumerate}
\item 若 $s \in S$，且 $x \in Y_s$ 是 $Y_s$ 某个不可约分支的泛点，
则 $x \in \text{Ass}_{X/S}(\mathcal{F})$；
\item 存在整数 $d \geq 0$，使得 $Y \to T$ 的相对维数为 $d$；见
态射章定义 \ref{morphisms-definition-relative-dimension-d}。
\end{enumerate}
\end{lemma}

\begin{proof}
这立即由逐点版本引理 \ref{more-morphisms-lemma-associated-point-specializes} 得出。
注意，为计算局部代数概形 $Y_s$ 的维数，只需考察泛点附近；见
簇章第 \ref{varieties-section-algebraic-schemes} 节。
\end{proof}

\begin{remark}
\label{more-morphisms-remark-flat-equidimensional}
下面列出一些情形，在这些情形中，上述材料，尤其是引理
\ref{more-morphisms-lemma-flat-associated-equidimensional}，可以用来断定概形态射
$f : X \to S$ 具有 态射章定义
\ref{morphisms-definition-relative-dimension-d} 所定义的相对维数 $d$。
例如，以下条件足够：
\begin{enumerate}
\item $X$ 是整概形，泛点为 $\xi$；
\item $\kappa(\xi)$ 在 $\kappa(f(\xi))$ 上的超越次数为 $d$；
\item $f$ 局部有限型；
\item 存在有限型拟凝聚 $\mathcal{O}_X$-模 $\mathcal{F}$，它在 $S$
上平坦，且 $\text{Supp}(\mathcal{F}) = X$。
\end{enumerate}
另一组足够的假设如下：
\begin{enumerate}
\item $S$ 不可约，泛点为 $\eta$；
\item $X_\eta$ 在 $X$ 中稠密；
\item $X_\eta$ 的每个不可约分支都具有维数 $d$；
\item $f$ 局部有限型；
\item 存在有限型拟凝聚 $\mathcal{O}_X$-模 $\mathcal{F}$，它在 $S$
上平坦，且 $\text{Supp}(\mathcal{F}) = X$。
\end{enumerate}
当然，可以放宽 $\mathcal{F}$ 的平坦性条件，只要求 $\mathcal{F}$
在 $S$ 上余维 $0$ 处平坦；也就是说，$\mathcal{F}$ 在每个纤维的每个泛点处都于
$S$ 上平坦。若以后需要这些结果，我们会在此仔细陈述并证明它们。
\end{remark}











\section{再论正规化}
\label{more-morphisms-section-normalization}

\noindent
正规化与光滑基变换交换。

\begin{lemma}
\label{more-morphisms-lemma-integral-closure-smooth-pullback}
设 $f : Y \to X$ 是概形的光滑态射，$\mathcal{A}$ 是拟凝聚
$\mathcal{O}_X$-代数层。$\mathcal{O}_Y$ 在 $f^*\mathcal{A}$ 中的整闭包
等于 $f^*\mathcal{A}'$，其中 $\mathcal{A}' \subset \mathcal{A}$ 是
$\mathcal{O}_X$ 在 $\mathcal{A}$ 中的整闭包。
\end{lemma}

\begin{proof}
这是用概形语言对 代数章引理
\ref{algebra-lemma-integral-closure-commutes-smooth} 的重述。略去细节。
\end{proof}

\begin{lemma}[正规化与光滑基变换交换]
\label{more-morphisms-lemma-normalization-smooth-localization}
设
$$
\xymatrix{
Y_2 \ar[r] \ar[d]_{f_2} & Y_1 \ar[d]^{f_1} \\
X_2 \ar[r]^\varphi & X_1
}
$$
是概形范畴中的纤维方块。假设 $f_1$ 拟紧且拟分离，并且 $\varphi$ 光滑。
设 $Y_i \to X_i' \to X_i$ 是 $X_i$ 在 $Y_i$ 中的正规化。则
$X_2' \cong X_2 \times_{X_1} X_1'$。
\end{lemma}

\begin{proof}
将分解 $Y_1 \to X_1' \to X_1$ 基变换到 $X_2$，得到分解
$Y_2 \to X_2 \times_{X_1} X_1' \to X_2$，并且
$X_2 \times_{X_1} X_1' \to X_2$ 是整态射（态射章引理
\ref{morphisms-lemma-base-change-finite}）。因此，由 态射章引理
\ref{morphisms-lemma-characterize-normalization} 的泛性质，得到态射
$h : X_2' \to X_2 \times_{X_1} X_1'$。注意，$X_2'$ 是
$\mathcal{O}_{X_2}$ 在 $f_{2, *}\mathcal{O}_{Y_2}$ 中的整闭包的相对谱。
若 $\mathcal{A}' \subset f_{1, *}\mathcal{O}_{Y_1}$ 表示
$\mathcal{O}_{X_1}$ 的整闭包，则 $X_2 \times_{X_1} X_1'$ 是
$\varphi^*\mathcal{A}'$ 的相对谱；见 构造章引理
\ref{constructions-lemma-spec-properties}。由 概形上同调章引理
\ref{coherent-lemma-flat-base-change-cohomology}，有
$f_{2, *}\mathcal{O}_{Y_2} = \varphi^*f_{1, *}\mathcal{O}_{Y_1}$。
故结论由引理 \ref{more-morphisms-lemma-integral-closure-smooth-pullback} 得出。
\end{proof}

\begin{lemma}[正规化与光滑态射]
\label{more-morphisms-lemma-normalization-and-smooth}
设 $X \to Y$ 是概形的光滑态射。假设 $Y$ 的每个拟紧开集都只有有限多个
不可约分支。则 $X$ 也具有此性质，并且在 $X$ 上存在唯一同构
$X^\nu = X \times_Y Y^\nu$，其中 $X^\nu$、$Y^\nu$ 分别是 $X$、$Y$
的正规化。
\end{lemma}

\begin{proof}
由 下降章引理
\ref{descent-lemma-locally-finite-nr-irred-local-fppf}，$X$ 的每个拟紧开集
都只有有限多个不可约分支。注意 $X_{red} = X \times_Y Y_{red}$，因为
约化概形上的光滑概形是约化的；见 下降章引理
\ref{descent-lemma-reduced-local-smooth}。故可假设 $X$ 与 $Y$ 都约化
（按定义，一个概形的正规化等于其约化的正规化）。其次，由 Descent,
引理 \ref{descent-lemma-normal-local-smooth}，
$X' = X \times_Y Y^\nu$ 是正规概形。态射 $X' \to Y^\nu$ 光滑
（从而平坦），因此 $X'$ 各不可约分支的泛点位于 $Y^\nu$ 各不可约
分支的泛点上方。由于 $Y^\nu \to Y$ 是双有理态射，可知
$X' \to X$ 也是双有理态射（因为 $X' \to Y^\nu$ 在 $Y$ 的泛点上方的
纤维上诱导同构）。于是存在分解 $X^\nu \to X' \to X$；见 Morphisms,
引理 \ref{morphisms-lemma-normalization-normal}。由于 $X'$ 正规且在
$X$ 上整，此分解是同构。
\end{proof}

\begin{lemma}[正规化与 Hensel 化]
\label{more-morphisms-lemma-normalization-and-henselization}
设 $X$ 是局部诺特概形，$\nu : X^\nu \to X$ 是正规化态射。则对任意点
$x \in X$，基变换
$$
X^\nu \times_X \Spec(\mathcal{O}_{X, x}^h) \to \Spec(\mathcal{O}_{X, x}^h),
\quad\text{分别}\quad
X^\nu \times_X \Spec(\mathcal{O}_{X, x}^{sh}) \to \Spec(\mathcal{O}_{X, x}^{sh})
$$
分别是 $\Spec(\mathcal{O}_{X, x}^h)$ 与
$\Spec(\mathcal{O}_{X, x}^{sh})$ 的正规化。
\end{lemma}

\begin{proof}
设 $\eta_1, \ldots, \eta_r$ 是 $X$ 中经过 $x$ 的不可约分支的泛点。
将正规化基变换到 $\Spec(\mathcal{O}_{X, x})$，得到
$\mathcal{O}_{X, x}$ 在 $\prod \kappa(\eta_i)$ 中的整闭包的谱。
这由 态射章定义 \ref{morphisms-definition-normalization}
中 $X$ 的正规化构造以及 态射章引理
\ref{morphisms-lemma-integral-closure} 得出；也可使用 态射章引理
\ref{morphisms-lemma-description-normalization} 对正规化的描述。
因此问题化为如下代数问题。设 $A$ 是诺特局部环；回忆这蕴含其 Hensel 化
$A^h$ 与严格 Hensel 化 $A^{sh}$ 也都是诺特的（代数进阶章引理
\ref{more-algebra-lemma-henselization-noetherian}）。设
$\mathfrak p_1, \ldots, \mathfrak p_r$ 是其极小素理想，并设 $A'$ 是
$A$ 在 $\prod \kappa(\mathfrak p_i)$ 中的整闭包。问题是证明：
$A' \otimes_A A^h$ 以及相应的 $A' \otimes_A A^{sh}$，分别可由诺特局部环
$A^h$ 与 $A^{sh}$ 以同样方式构造。

\medskip\noindent
% P05-ERR-0030: frozen source says minimal primes of A (likely A^h) and A^sh.
由于 $A^h$ 以及相应的 $A^{sh}$ 都是平展 $A$-代数的余极限，
$A$ 与 $A^{sh}$ 的极小素理想恰好分别是 $A^h$ 与 $A^{sh}$ 中位于
$A$ 的极小素理想上方的素理想（由下降性质；见 代数章引理
\ref{algebra-lemma-flat-going-down} 与
\ref{algebra-lemma-minimal-prime-image-minimal-prime}）。因此 More on Algebra,
引理 \ref{more-algebra-lemma-fibres-henselization} 告诉我们，
$A^h \otimes_A \prod \kappa(\mathfrak p_i)$,
分别
$A^{sh} \otimes_A \prod \kappa(\mathfrak p_i)$
分别是 $A^h$ 与 $A^{sh}$ 的极小素理想处的剩余域之积。已知在一个上环中
取整闭包与平展基变换交换；见 代数章引理
\ref{algebra-lemma-integral-closure-commutes-etale}。把 $A^h$ 与 $A^{sh}$
写成平展 $A$-代数的极限，可知对到 $A^h$ 与 $A^{sh}$ 的基变换也有同样
结论（也可使用更一般的 代数章引理
\ref{algebra-lemma-integral-closure-commutes-colim-smooth}）。
\end{proof}








\section{正规态射}
\label{more-morphisms-section-normal}

\noindent
Deligne 与 Mumford 的文章 \cite{DM} 提到了正规态射的概念。这只是可以用
类似方式定义的一系列态射类型之一\footnote{其他类型包括 coprof
$\leq k$、Cohen--Macaulay、$(S_k)$、正则、$(R_k)$ 与约化。见
\cite[IV Definition 6.8.1.]{EGA}。Gorenstein 态射将在 Duality for Schemes,
第 \ref{duality-section-gorenstein} 节 中定义。}随着需要，我们会逐步用
各自的章节加入这些类型。

\begin{definition}
\label{more-morphisms-definition-normal}
设 $f : X \to Y$ 是概形态射。假设所有纤维 $X_y$ 都是局部诺特概形。
\begin{enumerate}
\item 设 $x \in X$ 且 $y = f(x)$。若 $f$ 在 $x$ 处平坦，并且概形
$X_y$ 在 $x$ 处于 $\kappa(y)$ 上几何正规（见 簇章定义
\ref{varieties-definition-geometrically-normal}），则称 $f$ 在
$x$ 处{\it 正规}。
\item 若 $f$ 在 $X$ 的每一点处都正规，则称 $f$ 是{\it 正规态射}。
\end{enumerate}
\end{definition}

\noindent
因此，态射 $X \to Y$ 正规这一条件，比仅仅要求所有纤维都是正规局部诺特
概形更强。

\begin{lemma}
\label{more-morphisms-lemma-normal}
设 $f : X \to Y$ 是概形态射。假设 $f$ 的所有纤维都局部诺特。
以下条件等价：
\begin{enumerate}
\item $f$ 正规；
\item $f$ 平坦，且其纤维都是几何正规概形。
\end{enumerate}
\end{lemma}

\begin{proof}
这直接由定义得出。
\end{proof}

\begin{lemma}
\label{more-morphisms-lemma-smooth-normal}
光滑态射是正规的。
\end{lemma}

\begin{proof}
设 $f : X \to Y$ 是光滑态射。由 态射章引理
\ref{morphisms-lemma-smooth-locally-finite-presentation}，$f$ 局部有限表示，
故纤维 $X_y$ 在某个域上局部有限型，从而局部诺特。此外，$f$ 平坦；见
态射章引理 \ref{morphisms-lemma-smooth-flat}。最后，纤维 $X_y$
在域上光滑（由 态射章引理
\ref{morphisms-lemma-base-change-smooth}），故由 簇章引理
\ref{varieties-lemma-smooth-geometrically-normal} 几何正规。因此由引理
\ref{more-morphisms-lemma-normal}，$f$ 正规。
\end{proof}

\noindent
我们要证明，对光滑拓扑而言，这一概念在源与靶上都是局部的。先处理纤维
局部诺特这一性质。

\begin{lemma}
\label{more-morphisms-lemma-locally-Noetherian-fibres-fppf-local-source-and-target}
性质 $\mathcal{P}(f)=$“$f$ 的纤维局部诺特”在源与靶的 fppf 拓扑下都是
局部的。
\end{lemma}

\begin{proof}
设 $f : X \to Y$ 是概形态射。设
$\{\varphi_i : Y_i \to Y\}_{i \in I}$ 是 $Y$ 的 fppf 覆盖，并以
$f_i : X_i \to Y_i$ 表示 $f$ 沿 $\varphi_i$ 的基变换。设 $i \in I$，
且 $y_i \in Y_i$ 是一点。令 $y = \varphi_i(y_i)$。注意
$$
X_{i, y_i} = \Spec(\kappa(y_i)) \times_{\Spec(\kappa(y))} X_y.
$$
此外，由于 $\varphi_i$ 有限表示，域扩张 $\kappa(y_i)/\kappa(y)$ 有限生成。
因此在此情形中，$X_y$ 局部诺特，当且仅当 $X_{i, y_i}$ 局部诺特；见
簇章引理 \ref{varieties-lemma-locally-Noetherian-base-change}。
这就推出在靶上的局部性。

\medskip\noindent
设 $\{X_i \to X\}$ 是 $X$ 的 fppf 覆盖，并设 $y \in Y$。此时
$\{X_{i, y} \to X_y\}$ 是该纤维的 fppf 覆盖。因此源上的局部性由
下降章引理 \ref{descent-lemma-Noetherian-local-fppf} 得出。
\end{proof}

\begin{lemma}
\label{more-morphisms-lemma-normal-fppf-local-source-and-target}
性质 $\mathcal{P}(f)=$“$f$ 的纤维局部诺特且 $f$ 正规”在靶的 fppf
拓扑下是局部的，并在源的光滑拓扑下是局部的。
\end{lemma}

\begin{proof}
我们有
$\mathcal{P}(f) =
\mathcal{P}_1(f) \wedge \mathcal{P}_2(f) \wedge \mathcal{P}_3(f)$
，其中 $\mathcal{P}_1(f)=$“$f$ 的纤维局部诺特”，
$\mathcal{P}_2(f)=$“$f$ 平坦”，而
$\mathcal{P}_3(f)=$“$f$ 的纤维几何正规”。我们已经知道
$\mathcal{P}_1$ 与 $\mathcal{P}_2$ 在源与靶的 fppf 拓扑下都是局部的；见
引理 \ref{more-morphisms-lemma-locally-Noetherian-fibres-fppf-local-source-and-target}
以及 下降章引理 \ref{descent-lemma-descending-property-flat} 与
\ref{descent-lemma-flat-fpqc-local-source}。故只需处理 $\mathcal{P}_3$。

\medskip\noindent
设 $f : X \to Y$ 是概形态射。设
$\{\varphi_i : Y_i \to Y\}_{i \in I}$ 是 $Y$ 的 fpqc 覆盖，并以
$f_i : X_i \to Y_i$ 表示 $f$ 沿 $\varphi_i$ 的基变换。设 $i \in I$，
且 $y_i \in Y_i$ 是一点。令 $y = \varphi_i(y_i)$。注意
$$
X_{i, y_i} = \Spec(\kappa(y_i)) \times_{\Spec(\kappa(y))} X_y.
$$
因此在此情形中，$X_y$ 几何正规，当且仅当 $X_{i, y_i}$ 几何正规；见
簇章引理 \ref{varieties-lemma-geometrically-normal-upstairs}。
这表明 $\mathcal{P}_3$ 在靶的 fpqc 拓扑下是局部的。

\medskip\noindent
设 $\{X_i \to X\}$ 是 $X$ 的光滑覆盖，并设 $y \in Y$。此时
$\{X_{i, y} \to X_y\}$ 是该纤维的光滑覆盖。因此 $\mathcal{P}_3$
在源的光滑拓扑下的局部性由 下降章引理
\ref{descent-lemma-normal-local-smooth} 得出。综合以上结果即得本引理。
\end{proof}



\section{正则态射}
\label{more-morphisms-section-regular}

\noindent
与第 \ref{more-morphisms-section-normal} 节比较。这一概念的代数版本见 More on Algebra,
第 \ref{more-algebra-section-regular} 节。

\begin{definition}
\label{more-morphisms-definition-regular}
设 $f : X \to Y$ 是概形态射。假设所有纤维 $X_y$ 都是局部诺特概形。
\begin{enumerate}
\item 设 $x \in X$ 且 $y = f(x)$。若 $f$ 在 $x$ 处平坦，并且概形
$X_y$ 在 $x$ 处于 $\kappa(y)$ 上几何正则（见 簇章定义
\ref{varieties-definition-geometrically-regular}），则称 $f$ 在
$x$ 处{\it 正则}。
\item 若 $f$ 在 $X$ 的每一点处都正则，则称 $f$ 是{\it 正则态射}。
\end{enumerate}
\end{definition}

\noindent
态射 $X \to Y$ 正则这一条件，比仅仅要求所有纤维都是正则局部诺特概形
更强。

\begin{lemma}
\label{more-morphisms-lemma-regular}
设 $f : X \to Y$ 是概形态射。假设 $f$ 的所有纤维都局部诺特。
以下条件等价：
\begin{enumerate}
\item $f$ 正则；
\item $f$ 平坦，且其纤维都是几何正则概形；
\item 对每一对满足 $f(U) \subset V$ 的仿射开集 $U \subset X$、
$V \subset Y$，环映射 $\mathcal{O}_Y(V) \to \mathcal{O}_X(U)$ 正则；
\item 存在开覆盖 $Y = \bigcup_{j \in J} V_j$ 以及开覆盖
$f^{-1}(V_j) = \bigcup_{i \in I_j} U_i$，使每个态射
$U_i \to V_j$ 都正则；
\item 存在仿射开覆盖 $Y = \bigcup_{j \in J} V_j$ 以及仿射开覆盖
$f^{-1}(V_j) = \bigcup_{i \in I_j} U_i$，使环映射
$\mathcal{O}_Y(V_j) \to \mathcal{O}_X(U_i)$ 都正则。
\end{enumerate}
\end{lemma}

\begin{proof}
由定义立即得到 (1) 与 (2) 等价。设 $x \in X$ 且 $y = f(x)$。
按定义，$f$ 在 $x$ 处平坦，当且仅当
$\mathcal{O}_{Y, y} \to \mathcal{O}_{X, x}$ 是平坦环映射；而 $X_y$
在 $x$ 处于 $\kappa(y)$ 上几何正则，当且仅当
$\mathcal{O}_{X_y, x} = \mathcal{O}_{X, x}/\mathfrak m_y\mathcal{O}_{X, x}$
是 $\kappa(y)$ 上的几何正则代数。因此，
% P05-ERR-0032: frozen source unexpectedly capitalizes "Whether" mid-sentence.
$f$ 在 $x$ 处是否正则，仅取决于局部环之间的局部同态
$\mathcal{O}_{Y, y} \to \mathcal{O}_{X, x}$。故 (1) 与 (4) 显然等价。

\medskip\noindent
回忆 代数进阶章定义 \ref{more-algebra-definition-regular}：
环映射 $A \to B$ 正则，当且仅当它平坦，并且对所有素理想
$\mathfrak p \subset A$，纤维环 $B \otimes_A \kappa(\mathfrak p)$
都是诺特且几何正则的。由 簇章引理
\ref{varieties-lemma-geometrically-regular}，这等价于
$\Spec(B \otimes_A \kappa(\mathfrak p))$ 是 $\kappa(\mathfrak p)$ 上的
几何正则概形。因此 (2) 推出 (3)，而 (3) 显然推出 (5)。最后假设 (5)。
这推出 $f$ 平坦（见 态射章引理
\ref{morphisms-lemma-flat-characterize}）。此外，若 $y \in Y$，则对某个
$j$ 有 $y \in V_j$，并且 $X_y = \bigcup_{i \in I_j} U_{i, y}$；由
簇章引理 \ref{varieties-lemma-geometrically-regular}，每个
$U_{i, y}$ 都在 $\kappa(y)$ 上几何正则。再次应用 簇章引理
\ref{varieties-lemma-geometrically-regular}，可知 $X_y$ 几何正则。
故 (2) 成立，本引理得证。
\end{proof}

\begin{lemma}
\label{more-morphisms-lemma-smooth-regular}
光滑态射是正则的。
\end{lemma}

\begin{proof}
设 $f : X \to Y$ 是光滑态射。由 态射章引理
\ref{morphisms-lemma-smooth-locally-finite-presentation}，$f$ 局部有限表示，
故纤维 $X_y$ 在某个域上局部有限型，从而局部诺特。此外，$f$ 平坦；见
态射章引理 \ref{morphisms-lemma-smooth-flat}。最后，纤维 $X_y$
在域上光滑（由 态射章引理
\ref{morphisms-lemma-base-change-smooth}），并因此由
% P05-ERR-0031: frozen source cites smooth-geometrically-normal for regularity.
簇章引理 \ref{varieties-lemma-smooth-geometrically-normal}
几何正则。于是由引理
\ref{more-morphisms-lemma-regular}，$f$ 正则。
\end{proof}

\begin{lemma}
\label{more-morphisms-lemma-regular-fppf-local-source-and-target}
性质 $\mathcal{P}(f)=$“$f$ 的纤维局部诺特且 $f$ 正则”在靶的 fppf
拓扑下是局部的，并在源的光滑拓扑下是局部的。
\end{lemma}

\begin{proof}
我们有
$\mathcal{P}(f) =
\mathcal{P}_1(f) \wedge \mathcal{P}_2(f) \wedge \mathcal{P}_3(f)$
，其中 $\mathcal{P}_1(f)=$“$f$ 的纤维局部诺特”，
$\mathcal{P}_2(f)=$“$f$ 平坦”，而
$\mathcal{P}_3(f)=$“$f$ 的纤维几何正则”。我们已经知道
$\mathcal{P}_1$ 与 $\mathcal{P}_2$ 在源与靶的 fppf 拓扑下都是局部的；见
引理 \ref{more-morphisms-lemma-locally-Noetherian-fibres-fppf-local-source-and-target}
以及 下降章引理 \ref{descent-lemma-descending-property-flat} 与
\ref{descent-lemma-flat-fpqc-local-source}。故只需处理 $\mathcal{P}_3$。

\medskip\noindent
设 $f : X \to Y$ 是概形态射。设
$\{\varphi_i : Y_i \to Y\}_{i \in I}$ 是 $Y$ 的 fpqc 覆盖，并以
$f_i : X_i \to Y_i$ 表示 $f$ 沿 $\varphi_i$ 的基变换。设 $i \in I$，
且 $y_i \in Y_i$ 是一点。令 $y = \varphi_i(y_i)$。注意
$$
X_{i, y_i} = \Spec(\kappa(y_i)) \times_{\Spec(\kappa(y))} X_y.
$$
因此在此情形中，$X_y$ 几何正则，当且仅当 $X_{i, y_i}$ 几何正则；见
簇章引理 \ref{varieties-lemma-geometrically-regular-upstairs}。
这表明 $\mathcal{P}_3$ 在靶的 fpqc 拓扑下是局部的。

\medskip\noindent
设 $\{X_i \to X\}$ 是 $X$ 的光滑覆盖，并设 $y \in Y$。此时
$\{X_{i, y} \to X_y\}$ 是该纤维的光滑覆盖。因此 $\mathcal{P}_3$
在源的光滑拓扑下的局部性由 下降章引理
\ref{descent-lemma-regular-local-smooth} 得出。综合以上结果即得本引理。
\end{proof}




\section{Cohen--Macaulay 态射}
\label{more-morphisms-section-CM}

\noindent
与第 \ref{more-morphisms-section-normal} 节比较。注意，正如 Algebra, Section
\ref{algebra-section-geometrically-CM} 与 Varieties, Section
\ref{varieties-section-CM} 所指出的，“几何 Cohen--Macaulay”等同于通常的
Cohen--Macaulay。

\begin{definition}
\label{more-morphisms-definition-CM}
设 $f : X \to Y$ 是概形态射。假设所有纤维 $X_y$ 都是局部诺特概形。
\begin{enumerate}
\item 设 $x \in X$ 且 $y = f(x)$。若 $f$ 在 $x$ 处平坦，并且概形
$X_y$ 在 $x$ 处的局部环是 Cohen--Macaulay 的，则称 $f$ 在
$x$ 处{\it Cohen--Macaulay}。
\item 若 $f$ 在 $X$ 的每一点处都 Cohen--Macaulay，则称 $f$ 是
{\it Cohen--Macaulay 态射}。
\end{enumerate}
\end{definition}

\noindent
下面给出一种等价表述。

\begin{lemma}
\label{more-morphisms-lemma-CM}
设 $f : X \to Y$ 是概形态射。假设 $f$ 的所有纤维都局部诺特。
以下条件等价：
\begin{enumerate}
\item $f$ 是 Cohen--Macaulay 态射；
\item $f$ 平坦，且其纤维都是 Cohen--Macaulay 概形。
\end{enumerate}
\end{lemma}

\begin{proof}
这直接由定义得出。
\end{proof}

\begin{lemma}
\label{more-morphisms-lemma-CM-dimension}
设 $f : X \to Y$ 是局部诺特概形之间的局部有限型 Cohen--Macaulay
态射。对 $X$ 中每个在 $Y$ 中的像为 $y$ 的点 $x$，有
$$
\dim_x(X) = \dim_y(Y) + \dim_x(X_y),
$$
，其中 $X_y$ 表示 $y$ 上方的纤维。
\end{lemma}

\begin{proof}
用 $x$ 的一个开邻域替换 $X$ 后，存在自然数 $d$，使 $X \to Y$
的所有纤维在每一点处的维数均为 $d$；见 态射章引理
\ref{morphisms-lemma-flat-finite-presentation-CM-fibres-relative-dimension}。
于是 $f$ 平坦、局部有限型，且相对维数为 $d$。故结论由 Morphisms,
引理 \ref{morphisms-lemma-rel-dimension-dimension} 得出。
\end{proof}

\begin{lemma}
\label{more-morphisms-lemma-composition-CM}
设 $f : X \to Y$ 与 $g : Y \to Z$ 是概形态射。假设 $f$、$g$ 与
$g \circ f$ 的纤维都局部诺特。设 $x \in X$ 的像分别为
$y \in Y$ 与 $z \in Z$。
\begin{enumerate}
\item 若 $f$ 在 $x$ 处 Cohen--Macaulay，且 $g$ 在 $f(x)$ 处
Cohen--Macaulay，则 $g \circ f$ 在 $x$ 处 Cohen--Macaulay。
\item 若 $f$ 与 $g$ 都是 Cohen--Macaulay 态射，则 $g \circ f$
也是 Cohen--Macaulay 态射。
\item 若 $g \circ f$ 在 $x$ 处 Cohen--Macaulay，且 $f$ 在 $x$
处平坦，则 $f$ 在 $x$ 处 Cohen--Macaulay，且 $g$ 在 $f(x)$ 处
Cohen--Macaulay。
\item 若 $g \circ f$ 是 Cohen--Macaulay 态射且 $f$ 平坦，则 $f$
是 Cohen--Macaulay 态射，并且 $g$ 在 $f$ 的像中每一点处都
Cohen--Macaulay。
\end{enumerate}
\end{lemma}

\begin{proof}
考虑诺特局部环的映射
$$
\mathcal{O}_{Y_z, y} \to \mathcal{O}_{X_z, x}
$$
并注意到其纤维为
$$
\mathcal{O}_{X_z, x}/\mathfrak m_{Y_z, y}\mathcal{O}_{X_z, x} =
\mathcal{O}_{X_y, x}
$$
% P05-ERR-0033: frozen source has ungrammatical "Thus the lemma this follows".
因此本引理由 代数章引理 \ref{algebra-lemma-CM-goes-up} 得出。
\end{proof}

\begin{lemma}
\label{more-morphisms-lemma-flat-morphism-from-CM-scheme}
设 $f : X \to Y$ 是局部诺特概形之间的平坦态射。若 $X$ 是
Cohen--Macaulay 概形，则 $f$ 是 Cohen--Macaulay 态射，并且对所有
$x \in X$，$\mathcal{O}_{Y, f(x)}$ 都是 Cohen--Macaulay 环。
\end{lemma}

\begin{proof}
化为代数陈述后，结论由 代数章引理
\ref{algebra-lemma-CM-goes-up} 得出。
\end{proof}

\begin{lemma}
\label{more-morphisms-lemma-base-change-CM}
设 $f : X \to Y$ 是概形态射，并假设所有纤维 $X_y$ 都是局部诺特概形。
设 $Y' \to Y$ 局部有限型，并设 $f' : X' \to Y'$ 是 $f$ 的基变换。
设点 $x' \in X'$ 的像为 $x \in X$。
\begin{enumerate}
\item 若 $f$ 在 $x$ 处 Cohen--Macaulay，则 $f' : X' \to Y'$
在 $x'$ 处 Cohen--Macaulay。
\item 若 $f$ 在 $x$ 处平坦，且 $f'$ 在 $x'$ 处 Cohen--Macaulay，
则 $f$ 在 $x$ 处 Cohen--Macaulay。
\item 若 $Y' \to Y$ 在 $f'(x')$ 处平坦，且 $f'$ 在 $x'$ 处
Cohen--Macaulay，则 $f$ 在 $x$ 处 Cohen--Macaulay。
\end{enumerate}
\end{lemma}

\begin{proof}
注意，关于 $Y' \to Y$ 的假设蕴含：对映到 $y \in Y$ 的
$y' \in Y'$，域扩张 $\kappa(y')/\kappa(y)$ 有限生成。因此所有纤维
$X'_{y'} = (X_y)_{\kappa(y')}$ 也局部诺特；见 簇章引理
\ref{varieties-lemma-locally-Noetherian-base-change}。故本引理的陈述有意义。
令 $y' = f'(x')$ 且 $y = f(x)$，得到局部环的交换图
$$
\xymatrix{
\mathcal{O}_{X', x'} & \mathcal{O}_{X, x} \ar[l] \\
\mathcal{O}_{Y', y'} \ar[u] & \mathcal{O}_{Y, y} \ar[l] \ar[u]
}
$$
，其中左上角是右上角与左下角在右下角上的张量积的一个局部化。

\medskip\noindent
假设 $f$ 在 $x$ 处 Cohen--Macaulay。映射
$\mathcal{O}_{Y, y} \to \mathcal{O}_{X, x}$ 的平坦性蕴含
$\mathcal{O}_{Y', y'} \to \mathcal{O}_{X', x'}$ 的平坦性；见 Algebra,
引理 \ref{algebra-lemma-base-change-flat-up-down}。环
$\mathcal{O}_{X, x}/\mathfrak m_y\mathcal{O}_{X, x}$ 是
Cohen--Macaulay 的，蕴含
$\mathcal{O}_{X', x'}/\mathfrak m_{y'}\mathcal{O}_{X', x'}$ 是
Cohen--Macaulay 的；见 簇章引理
\ref{varieties-lemma-CM-base-change}。故 $f'$ 在 $x'$ 处 Cohen--Macaulay。

\medskip\noindent
假设 $f$ 在 $x$ 处平坦，且 $f'$ 在 $x'$ 处 Cohen--Macaulay。
$\mathcal{O}_{X', x'}/\mathfrak m_{y'}\mathcal{O}_{X', x'}$ 是
Cohen--Macaulay 的，蕴含
$\mathcal{O}_{X, x}/\mathfrak m_y\mathcal{O}_{X, x}$ 是
Cohen--Macaulay 的；见 簇章引理
\ref{varieties-lemma-CM-base-change}。故 $f$ 在 $x$ 处 Cohen--Macaulay。

\medskip\noindent
假设 $Y' \to Y$ 在 $y'$ 处平坦，且 $f'$ 在 $x'$ 处
Cohen--Macaulay。映射
$\mathcal{O}_{Y', y'} \to \mathcal{O}_{X', x'}$ 与
$\mathcal{O}_{Y, y} \to \mathcal{O}_{Y', y'}$ 的平坦性蕴含
$\mathcal{O}_{Y, y} \to \mathcal{O}_{X, x}$ 的平坦性；见 Algebra,
引理 \ref{algebra-lemma-base-change-flat-up-down}。环
$\mathcal{O}_{X', x'}/\mathfrak m_{y'}\mathcal{O}_{X', x'}$ 是
Cohen--Macaulay 的，蕴含
$\mathcal{O}_{X, x}/\mathfrak m_y\mathcal{O}_{X, x}$ 是
Cohen--Macaulay 的；见 簇章引理
\ref{varieties-lemma-CM-base-change}。故 $f$ 在 $x$ 处 Cohen--Macaulay。
\end{proof}

\begin{lemma}
\label{more-morphisms-lemma-flat-finite-presentation-CM-open}
\begin{reference}
\cite[IV Corollary 12.1.7(iii)]{EGA}
\end{reference}
设 $f : X \to S$ 是平坦且局部有限表示的概形态射。令
$$
W = \{x \in X \mid f\text{ is Cohen-Macaulay at }x\}
$$
则
\begin{enumerate}
\item $W = \{x \in X \mid \mathcal{O}_{X_{f(x)}, x}\text{ 为 Cohen--Macaulay}\}$；
\item $W$ 在 $X$ 中是开的；
\item $W$ 在 $X \to S$ 的每个纤维中稠密；
\item $W$ 的形成与 $f$ 的任意基变换交换：对任意态射
$g : S' \to S$，考虑 $f$ 的基变换 $f' : X' \to S'$ 以及投影
$g' : X' \to X$。则态射 $f'$ 对应的集合 $W'$ 等于
$W' = (g')^{-1}(W)$。
\end{enumerate}
\end{lemma}

\begin{proof}
由于 $f$ 平坦且具有局部诺特纤维，(1) 中的等式由定义成立。(2) 与 (3)
由 代数章引理
\ref{algebra-lemma-generic-CM-flat-finite-presentation} 得出。(4) 可由
代数章引理 \ref{algebra-lemma-CM-locus-commutes-base-change} 或
簇章引理 \ref{varieties-lemma-CM-base-change} 得出。
\end{proof}

\begin{lemma}
\label{more-morphisms-lemma-flat-finite-presentation-characterize-CM}
设 $f : X \to S$ 是平坦且局部有限表示的概形态射。设 $x \in X$
的像为 $s \in S$，并令 $d = \dim_x(X_s)$。以下条件等价：
\begin{enumerate}
\item $f$ 在 $x$ 处 Cohen--Macaulay；
\item 存在 $x$ 的开邻域 $U \subset X$ 以及 $S$ 上的局部拟有限态射
$U \to \mathbf{A}^d_S$，它在 $x$ 处平坦；
\item 存在 $x$ 的开邻域 $U \subset X$ 以及 $S$ 上的局部拟有限平坦态射
$U \to \mathbf{A}^d_S$；
\item 对 $x$ 的任意开邻域 $U \subset X$ 及任意 $S$-态射
$g : U \to \mathbf{A}^d_S$，都有：$g$ 在 $x$ 处拟有限
$\Rightarrow$ $g$ 在 $x$ 处平坦。
\end{enumerate}
\end{lemma}

\begin{proof}
平坦性的开性说明 (2) 与 (3) 等价；见定理
\ref{more-morphisms-theorem-openness-flatness}。

\medskip\noindent
选择满足 $x \in U$ 的仿射开集 $U = \Spec(A) \subset X$，以及满足
$f(U) \subset V$ 的仿射开集 $V = \Spec(R) \subset S$。则 $R \to A$
是有限表示平坦环映射。设 $\mathfrak p \subset A$ 是 $x$ 对应的素理想。
用一个主局部化替换 $A$ 后，可以假设存在拟有限映射
$R[x_1, \ldots, x_d]  \to A$；见 代数章引理
\ref{algebra-lemma-quasi-finite-over-polynomial-algebra}。因此至少存在一对
$(U, g)$，由 $x$ 的开邻域 $U \subset X$ 与局部\footnote{若 $S$
拟分离，则 $g$ 将是拟有限的。}拟有限态射
$g : U \to \mathbf{A}^d_S$ 组成。

\medskip\noindent
断言：给定有限表示平坦映射 $R \to A$、素理想 $\mathfrak p \subset A$
以及在 $\mathfrak p$ 处拟有限的
$\varphi : R[x_1, \ldots, x_d] \to A$，则 $\Spec(\varphi)$ 在
$\mathfrak p$ 处平坦，当且仅当 $\Spec(A) \to \Spec(R)$ 在
$\mathfrak p$ 处 Cohen--Macaulay。事实上，由定理
\ref{more-morphisms-theorem-criterion-flatness-fibre}，平坦性可以在纤维上检验。
Cohen--Macaulay 性也如此（因为已经假设 $A$ 在 $R$ 上平坦）。
故该断言由 代数章引理 \ref{algebra-lemma-where-CM} 得出。

\medskip\noindent
该断言说明 (1) 与 (4) 等价；再结合第二段中已经构造出合适的
$(U, g)$ 这一事实，该断言还说明 (1) 与 (2) 等价。
\end{proof}

\begin{lemma}
\label{more-morphisms-lemma-flat-finite-presentation-CM-pieces}
设 $f : X \to S$ 是平坦且局部有限表示的概形态射。对 $d \geq 0$，
存在开集 $U_d \subset X$，满足以下性质：
\begin{enumerate}
\item $W = \bigcup_{d \geq 0} U_d$ 在 $f$ 的每个纤维中稠密；
\item $U_d \to S$ 的相对维数为 $d$（见 态射章定义
\ref{morphisms-definition-relative-dimension-d}）。
\end{enumerate}
\end{lemma}

\begin{proof}
这由引理 \ref{more-morphisms-lemma-flat-finite-presentation-CM-open} 与 态射章引理
\ref{morphisms-lemma-flat-finite-presentation-CM-fibres-relative-dimension}
结合得出。
\end{proof}

\begin{lemma}
\label{more-morphisms-lemma-flat-finite-presentation-specialization-dimension}
设 $f : X \to S$ 是平坦且局部有限表示的概形态射。假设
$x' \leadsto x$ 是 $X$ 中点的特化，其在 $S$ 中的像为
$s' \leadsto s$。若 $x$ 是 $X_s$ 某个不可约分支的泛点，则
$\dim_{x'}(X_{s'}) = \dim_x(X_s)$。
\end{lemma}

\begin{proof}
对某个 $d$，点 $x$ 包含在 $U_d$ 中，其中 $U_d$ 如引理
\ref{more-morphisms-lemma-flat-finite-presentation-CM-pieces} 所述。
\end{proof}

\begin{lemma}
\label{more-morphisms-lemma-CM-local-source-and-target}
性质 $\mathcal{P}(f)=$“$f$ 的纤维局部诺特且 $f$ 是
Cohen--Macaulay 态射”在靶的 fppf 拓扑下是局部的，并在源的 syntomic
拓扑下是局部的。
\end{lemma}

\begin{proof}
我们有
$\mathcal{P}(f) =
\mathcal{P}_1(f) \wedge \mathcal{P}_2(f)$
，其中 $\mathcal{P}_1(f)=$“$f$ 平坦”，而
$\mathcal{P}_2(f)=$“$f$ 的纤维局部诺特且 Cohen--Macaulay”。已知
$\mathcal{P}_1$ 在源与靶的 fppf 拓扑下都是局部的；见 下降章引理
\ref{descent-lemma-descending-property-flat} 与
\ref{descent-lemma-flat-fpqc-local-source}。故只需处理 $\mathcal{P}_2$。

\medskip\noindent
设 $f : X \to Y$ 是概形态射。设
$\{\varphi_i : Y_i \to Y\}_{i \in I}$ 是 $Y$ 的 fppf 覆盖，并以
$f_i : X_i \to Y_i$ 表示 $f$ 沿 $\varphi_i$ 的基变换。设 $i \in I$，
且 $y_i \in Y_i$ 是一点。令 $y = \varphi_i(y_i)$。注意
$$
X_{i, y_i} = \Spec(\kappa(y_i)) \times_{\Spec(\kappa(y))} X_y.
$$
，并且 $\kappa(y_i)/\kappa(y)$ 是有限生成域扩张。因此若 $X_y$
局部诺特，则 $X_{i, y_i}$ 局部诺特；见 簇章引理
\ref{varieties-lemma-locally-Noetherian-base-change}。若此外 $X_y$
是 Cohen--Macaulay 概形，则 $X_{i, y_i}$ 也是 Cohen--Macaulay 概形；
见 簇章引理 \ref{varieties-lemma-CM-base-change}。因此
$\mathcal{P}_2$ 在靶的 fppf 拓扑下是局部的。

\medskip\noindent
设 $\{X_i \to X\}$ 是 $X$ 的 syntomic 覆盖，并设 $y \in Y$。此时
$\{X_{i, y} \to X_y\}$ 是该纤维的 syntomic 覆盖。因此
$\mathcal{P}_2$ 在源的 syntomic 拓扑下的局部性由 下降章引理
\ref{descent-lemma-CM-local-syntomic} 得出。综合以上结果即得本引理。
\end{proof}






\section{切割 Cohen--Macaulay 态射}
\label{more-morphisms-section-slicing-CM}

\noindent
本节的结果最终导向如下断言：fppf 拓扑与“有限表示、平坦、拟有限”拓扑
相同。下一个引理与 除子章引理
\ref{divisors-lemma-fibre-Cartier} 密切相关。

\begin{lemma}
\label{more-morphisms-lemma-slice-given-element}
设 $f : X \to S$ 是概形态射。设点 $x \in X$ 的像为 $s \in S$，并设
$h \in \mathfrak m_x \subset \mathcal{O}_{X, x}$。假设
\begin{enumerate}
\item $f$ 局部有限表示；
\item $f$ 在 $x$ 处平坦；
\item $h$ 在
$\mathcal{O}_{X_s, x} = \mathcal{O}_{X, x}/\mathfrak m_s\mathcal{O}_{X, x}$
中的像 $\overline{h}$ 是非零因子。
\end{enumerate}
则存在 $x$ 的仿射开邻域 $U \subset X$，使 $h$ 来自
$h \in \Gamma(U, \mathcal{O}_U)$，并且 $D = V(h)$ 是 $U$ 中包含
$x \in D$ 的有效 Cartier 除子，且 $D \to S$ 平坦并局部有限表示。
\end{lemma}

\begin{proof}
我们把问题化到诺特情形来证明。由平坦性的开性（见定理
\ref{more-morphisms-theorem-openness-flatness}），用 $x$ 的一个开邻域替换 $X$ 后，
可以假设 $X \to S$ 平坦。还可假设 $X$ 与 $S$ 仿射。稍微缩小 $X$ 后，
可以假设存在 $h \in \Gamma(X, \mathcal{O}_X)$，它映到给定的 $h$。

\medskip\noindent
可写 $S = \Spec(A)$，并把 $A = \colim_i A_i$ 写成有限型
$\mathbf{Z}$-代数的定向余极限。于是，由 代数章引理
\ref{algebra-lemma-flat-finite-presentation-limit-flat}，或 极限章引理
\ref{limits-lemma-descend-finite-presentation}、
\ref{limits-lemma-descend-affine-finite-presentation} 以及
% P05-ERR-0034: frozen source repeats limits-lemma-descend-finite-presentation.
\ref{limits-lemma-descend-finite-presentation}，可以找到笛卡尔图
$$
\xymatrix{
X \ar[r] \ar[d]_f & X_0 \ar[d]^{f_0} \\
S \ar[r] & S_0
}
$$
其中 $f_0$ 平坦且有限表示，$X_0$ 仿射，而 $S_0$ 仿射且诺特。令
$x_0 \in X_0$ 与相应的 $s_0 \in S_0$ 分别为 $x$ 与相应的 $s$ 的像。
还可假设存在元素 $h_0 \in \Gamma(X_0, \mathcal{O}_{X_0})$，其在 $X$
上的限制为 $h$。（若使用上面的代数参考，这是显然的；若使用极限一章
的参考，则把 $h$ 看作态射 $X \to \mathbf{A}^1_S$，结论由 Limits,
引理 \ref{limits-lemma-descend-finite-presentation} 得出。）注意，
$\mathcal{O}_{X_s, x}$ 是
$\mathcal{O}_{(X_0)_{s_0}, x_0} \otimes_{\kappa(s_0)} \kappa(s)$
的一个局部化，故
$\mathcal{O}_{(X_0)_{s_0}, x_0} \to \mathcal{O}_{X_s, x}$ 是平坦局部
环映射，特别是忠实平坦的。因此像
$\overline{h}_0 \in \mathcal{O}_{(X_0)_{s_0}, x_0}$ 属于
$\mathfrak m_{(X_0)_{s_0}, x_0}$，且是非零因子。我们断言：用 $x_0$
的一个主开邻域替换 $X_0$ 后，元素 $h_0$ 是
$B_0 = \Gamma(X_0, \mathcal{O}_{X_0})$ 中的非零因子，并且
$B_0/h_0B_0$ 在 $A_0 = \Gamma(S_0, \mathcal{O}_{S_0})$ 上平坦。
若断言成立，则
$$
0 \to B_0 \xrightarrow{h_0} B_0 \to B_0/h_0B_0 \to 0
$$
是平坦 $A_0$-模的短正合列。故与 $A$ 作张量后仍正合（由 Algebra,
引理 \ref{algebra-lemma-flat-tor-zero}），从而本引理得证。

\medskip\noindent
只剩证明上述断言。相应的代数陈述如下（这里略去下标 ${}_0$）：设
$A \to B$ 是诺特环之间的有限型平坦环映射。设素理想
$\mathfrak q \subset B$ 位于 $\mathfrak p \subset A$ 上方。假设
$h \in \mathfrak q$ 在 $B_{\mathfrak q}/\mathfrak p B_{\mathfrak q}$
中的像是非零因子。目标是证明：
% P05-ERR-0035: frozen source has ungrammatical "after possible replacing".
对某个 $g \in B$、$g \not \in \mathfrak q$，以 $B_g$ 替换 $B$ 后，
元素 $h$ 成为非零因子，而 $B/hB$ 在 $A$ 上平坦。由 代数章引理
\ref{algebra-lemma-grothendieck}，$h$ 是 $B_{\mathfrak q}$ 中的非零因子，
且 $B_{\mathfrak q}/hB_{\mathfrak q}$ 在 $A$ 上平坦。由平坦性的开性，
见 代数章定理 \ref{algebra-theorem-openness-flatness} 或定理
\ref{more-morphisms-theorem-openness-flatness}，对某个 $g \in B$、
$g \not \in \mathfrak q$，以 $B_g$ 替换 $B$ 后，$B/hB$ 在 $A$ 上平坦。
最后，令 $I = \{b \in B \mid hb = 0\}$ 为 $h$ 的零化子。由于 $h$
是 $B_{\mathfrak q}$ 中的非零因子，$IB_{\mathfrak q} = 0$。又由于
$B$ 诺特，$I$ 有限生成。因此存在 $g \in B$、$g \not \in \mathfrak q$，
使 $IB_g = 0$。以 $B_g$ 替换 $B$ 后，$h$ 是非零因子。
\end{proof}

\begin{lemma}
\label{more-morphisms-lemma-slice-given-elements}
设 $f : X \to S$ 是概形态射。设点 $x \in X$ 的像为 $s \in S$，并设
$h_1, \ldots, h_r \in \mathcal{O}_{X, x}$。假设
\begin{enumerate}
\item $f$ 局部有限表示；
\item $f$ 在 $x$ 处平坦；
\item $h_1, \ldots, h_r$ 在
$\mathcal{O}_{X_s, x} = \mathcal{O}_{X, x}/\mathfrak m_s\mathcal{O}_{X, x}$
中的像构成正则序列。
\end{enumerate}
则存在 $x$ 的仿射开邻域 $U \subset X$，使 $h_1, \ldots, h_r$ 来自
$h_1, \ldots, h_r \in \Gamma(U, \mathcal{O}_U)$，并且
$Z = V(h_1, \ldots, h_r) \to U$ 是满足 $x \in Z$ 的正则浸入，而
$Z \to S$ 平坦且局部有限表示。此外，对 $S$ 上的任意概形 $S'$，
基变换 $Z_{S'} \to U_{S'}$ 都是正则浸入。
\end{lemma}

\begin{proof}
（我们关于正则序列的约定蕴含：对每个 $i$，$h_i \in \mathfrak m_x$。）
当 $r = 1$ 时，结论由引理 \ref{more-morphisms-lemma-slice-given-element} 与 Divisors,
引理 \ref{divisors-lemma-relative-Cartier} 结合得出；后一个引理说明
$V(h_1)$ 在基变换后仍是有效 Cartier 除子。当 $r > 1$ 时，对 $r$
作直接归纳即可（恰好应用 $r$ 次 $r = 1$ 的结果；略去细节）。

\medskip\noindent
另一种证明方法使用 Divisors, Section
\ref{divisors-section-relative-regular-immersion} 的材料。首先，由平坦性的
开性（见定理 \ref{more-morphisms-theorem-openness-flatness}），用 $x$ 的一个开邻域替换
$X$ 后，可以假设 $X \to S$ 平坦。还可假设 $X$ 与 $S$ 仿射。稍微缩小
$X$ 后，可以假设 $h_1, \ldots, h_r \in \Gamma(X, \mathcal{O}_X)$。
令 $Z = V(h_1, \ldots, h_r)$。注意，$X_s$ 是诺特概形（因为它是代数
$\kappa(s)$-概形；见 Varieties, Section
\ref{varieties-section-algebraic-schemes}），且 $X_s$ 上的拓扑由 $X$
上的拓扑诱导（见 概形章引理
\ref{schemes-lemma-fibre-topological}）。故进一步缩小 $X$ 后，可以假设
$Z_s \subset X_s$ 是由 $r$ 个元素 $h_i|_{X_s}$ 截出的正则浸入；见
除子章引理 \ref{divisors-lemma-Noetherian-scheme-regular-ideal}
及其证明。还显然有 $r = \dim_x(X_s) - \dim_x(Z_s)$，因为
\begin{align*}
\dim_x(X_s) & = \dim(\mathcal{O}_{X_s, x}) +
\text{trdeg}_{\kappa(s)}(\kappa(x)), \\
\dim_x(Z_s) & = \dim(\mathcal{O}_{Z_s, x}) +
\text{trdeg}_{\kappa(s)}(\kappa(x)), \\
\dim(\mathcal{O}_{X_s, x}) & = \dim(\mathcal{O}_{Z_s, x}) + r
\end{align*}
前两个等式由 代数章引理
\ref{algebra-lemma-dimension-at-a-point-finite-type-field} 得出，
% P05-ERR-0036: frozen source says "the second" after already citing first two.
而第二个等式由把 代数章引理
\ref{algebra-lemma-one-equation} 应用 $r$ 次得出。因此 除子章引理
\ref{divisors-lemma-fibre-quasi-regular} 的 (3) 表明（在 Zariski 意义下
缩小 $X$ 后）态射 $Z \to X$ 是正则浸入，可以对它应用 除子章引理
\ref{divisors-lemma-relative-regular-immersion-flat-in-neighbourhood}
（该引理给出平坦性与基变换陈述）。
\end{proof}

\begin{lemma}
\label{more-morphisms-lemma-slice-once}
设 $f : X \to S$ 是概形态射。设点 $x \in X$ 的像为 $s \in S$。假设
\begin{enumerate}
\item $f$ 局部有限表示；
\item $f$ 在 $x$ 处平坦；
\item $\mathcal{O}_{X_s, x}$ 具有 $\text{depth} \geq 1$。
\end{enumerate}
则存在 $x$ 的仿射开邻域 $U \subset X$，以及包含 $x$ 的有效 Cartier
除子 $D \subset U$，使 $D \to S$ 平坦且有限表示。
\end{lemma}

\begin{proof}
选取任意 $h \in \mathfrak m_x \subset \mathcal{O}_{X, x}$，使它在
$\mathcal{O}_{X_s, x}$ 中的像是非零因子，然后应用引理
\ref{more-morphisms-lemma-slice-given-element}。
\end{proof}

\begin{lemma}
\label{more-morphisms-lemma-slice-CM}
\begin{reference}
\cite[IV Proposition 17.16.1]{EGA}
\end{reference}
设 $f : X \to S$ 是概形态射。设点 $x \in X$ 的像为 $s \in S$。假设
\begin{enumerate}
\item $f$ 局部有限表示；
\item $f$ 在 $x$ 处 Cohen--Macaulay；
\item $x$ 是 $X_s$ 的闭点。
\end{enumerate}
则存在包含 $x$ 的正则浸入 $Z \to X$，使得
\begin{enumerate}
\item[(a)] $Z \to S$ 平坦且局部有限表示；
\item[(b)] $Z \to S$ 局部拟有限；
\item[(c)] 集合论意义下 $Z_s = \{x\}$。
\end{enumerate}
\end{lemma}

\begin{proof}
不妨并且确实用 $s$ 的一个仿射开邻域替换 $S$。对仿射 $S$，我们将对
$d = \dim_x(X_s)$ 作归纳来证明本引理。

\medskip\noindent
情形 $d = 0$。此时证明可以取 $Z$ 为 $x$ 的一个开邻域。（注意，开浸入
是正则浸入。）事实上，若 $d = 0$，则 $X \to S$ 在 $x$ 处拟有限；见
态射章引理
\ref{morphisms-lemma-locally-quasi-finite-rel-dimension-0}。因此存在仿射开邻域
$U \subset X$，使 $U \to S$ 拟有限；见 态射章引理
\ref{morphisms-lemma-quasi-finite-points-open}。用 $U$ 替换 $X$ 后，纤维
$X_s$ 是有限离散集。
% P05-ERR-0037: frozen source says an affine open neighbourhood "of X", likely x.
于是进一步用 $X$ 的一个仿射开邻域替换 $X$ 后，有
$f^{-1}(\{s\}) = \{x\}$（因为 $X_s$ 上的拓扑由 $X$ 上的拓扑诱导；
见 概形章引理 \ref{schemes-lemma-fibre-topological}）。
这证明了该情形。

\medskip\noindent
下面假设 $d > 0$。注意，由于 $x$ 是其纤维的闭点，扩张
$\kappa(x)/\kappa(s)$ 有限（由 Hilbert 零点定理；见 态射章引理
\ref{morphisms-lemma-closed-point-fibre-locally-finite-type}）。因此
$$
\text{depth}(\mathcal{O}_{X_s, x}) = \dim(\mathcal{O}_{X_s, x}) = d > 0
$$
，第一个等式因为 $\mathcal{O}_{X_s, x}$ 是 Cohen--Macaulay 环而成立，
第二个等式由 态射章引理
\ref{morphisms-lemma-dimension-fibre-at-a-point} 得出。因此可应用引理
\ref{more-morphisms-lemma-slice-once} 找到图
$$
\xymatrix{
D \ar[r] \ar[rrd] & U \ar[r] \ar[rd] & X \ar[d] \\
& & S
}
$$
，其中 $x \in D$。注意，对某个非零因子 $\overline{h}$，有
$\mathcal{O}_{D_s, x} = \mathcal{O}_{X_s, x}/(\overline{h})$；见
除子章引理 \ref{divisors-lemma-relative-Cartier}。因此
$\mathcal{O}_{D_s, x}$ 是 Cohen--Macaulay 环，其维数比
$\mathcal{O}_{X_s, x}$ 的维数小一；例如见 代数章引理
\ref{algebra-lemma-reformulate-CM}。于是态射 $D \to S$ 平坦、局部有限表示，
并在 $x$ 处 Cohen--Macaulay，且
$\dim_x(D_s) = \dim_x(X_s) - 1 = d - 1$。由归纳假设，可以找到具有
(a)、(b)、(c) 性质的正则浸入 $Z \to D$。由于 $Z \to D \to U$
都是正则浸入，由 除子章引理
\ref{divisors-lemma-composition-regular-immersion}，$Z \to U$ 也是
正则浸入。证明完毕。
\end{proof}

\begin{lemma}
\label{more-morphisms-lemma-qf-fp-flat-neighbourhood-dominates-fppf}
设 $f : X \to S$ 是平坦且局部有限表示的概形态射。设 $s \in S$
是 $f$ 的像中的一点。则存在交换图
$$
\xymatrix{
S' \ar[rr] \ar[rd]_g & & X \ar[ld]^f \\
& S
}
$$
，其中 $g : S' \to S$ 平坦、局部有限表示、局部拟有限，且
$s \in g(S')$。
\end{lemma}

\begin{proof}
由假设，纤维 $X_s$ 非空。因此存在闭点 $x \in X_s$，使 $f$ 在该点处
Cohen--Macaulay；见引理 \ref{more-morphisms-lemma-flat-finite-presentation-CM-open}。
应用引理 \ref{more-morphisms-lemma-slice-CM}，并令 $S' = Z$。
\end{proof}

\noindent
下一个引理说明，fppf 拓扑的层与“拟有限、平坦、有限表示”拓扑的层是
同一回事。

\begin{lemma}
\label{more-morphisms-lemma-qf-fp-flat-dominates-fppf}
设 $S$ 是概形。设 $\mathcal{U} = \{S_i \to S\}_{i \in I}$ 是 $S$
的 fppf 覆盖；见 拓扑章定义
\ref{topologies-definition-fppf-covering}。则存在加细 $\mathcal{U}$
（见 位点章定义 \ref{sites-definition-morphism-coverings}）的 fppf
覆盖 $\mathcal{V} = \{T_j \to S\}_{j \in J}$，使每个
$T_j \to S$ 都局部拟有限。
\end{lemma}

\begin{proof}
对每个 $s \in S$，存在 $i \in I$，使 $s$ 属于 $S_i \to S$ 的像。
由引理 \ref{more-morphisms-lemma-qf-fp-flat-neighbourhood-dominates-fppf}，可以找到态射
$g_s : T_s \to S$，使 $s \in g_s(T_s)$；该态射平坦、局部有限表示、
局部拟有限，并且 $g_s$ 经由 $S_i \to S$ 分解。因此
$\{T_s \to S\}$ 即为所求加细 $\mathcal{U}$ 的 $S$ 的覆盖。
\end{proof}



\section{泛纤维}
\label{more-morphisms-section-generic}

\noindent
这里给出若干关于泛纤维与邻近纤维之间关系的结果。

\begin{lemma}
\label{more-morphisms-lemma-empty-generic-fibre}
设 $f : X \to Y$ 是有限型概形态射。假设 $Y$ 不可约，泛点为 $\eta$。
若 $X_\eta = \emptyset$，则存在非空开集 $V \subset Y$，使
$X_V = V \times_Y X = \emptyset$。
\end{lemma}

\begin{proof}
这立即由更一般的 态射章引理
\ref{morphisms-lemma-quasi-compact-generic-point-not-in-image} 得出。
\end{proof}

\begin{lemma}
\label{more-morphisms-lemma-nonempty-generic-fibre}
设 $f : X \to Y$ 是有限型概形态射。假设 $Y$ 不可约，泛点为 $\eta$。
若 $X_\eta \not = \emptyset$，则存在非空开集 $V \subset Y$，使
$X_V = V \times_Y X \to V$ 满射。
\end{lemma}

\begin{proof}
取仿射开集后，这由 代数章引理
\ref{algebra-lemma-characterize-image-finite-type} 得出。（当然，这也可由
泛平坦性得出。）
\end{proof}

\begin{lemma}
\label{more-morphisms-lemma-nowhere-dense-generic-fibre}
设 $f : X \to Y$ 是有限型概形态射。假设 $Y$ 不可约，泛点为 $\eta$。
若 $Z \subset X$ 是闭子集，且 $Z_\eta$ 在 $X_\eta$ 中处处不稠密，
则存在非空开集 $V \subset Y$，使对所有 $y \in V$，$Z_y$ 在 $X_y$
中处处不稠密。
\end{lemma}

\begin{proof}
设 $Y' \subset Y$ 是 $Y$ 的约化，并令 $X' = Y' \times_Y X$、
$Z' = Y' \times_Y Z$。由 态射章引理
\ref{morphisms-lemma-reduction-universal-homeomorphism}，$Y' \to Y$
是泛同胚，故只需对 $Z' \subset X' \to Y'$ 证明本引理。因此可假设
$Y$ 是整概形；见 性质章引理
\ref{properties-lemma-characterize-integral}。由 态射章命题
\ref{morphisms-proposition-generic-flatness}，存在非空仿射开集
$V \subset Y$，使 $X_V \to V$ 与 $Z_V \to V$ 都平坦且有限表示。
我们断言这个 $V$ 可用。取 $y \in V$。若 $Z_y$ 有非空内部，则 $Z_y$
包含 $X_y$ 某个不可约分支的泛点 $\xi$。注意
$\eta \leadsto f(\xi)$。由于 $Z_V \to V$ 平坦，可以选择特化
$\xi' \leadsto \xi$，其中 $\xi' \in Z$ 且 $f(\xi') = \eta$；见
态射章引理 \ref{morphisms-lemma-generalizations-lift-flat}。由引理
\ref{more-morphisms-lemma-flat-finite-presentation-specialization-dimension}，有
$$
\dim_{\xi'}(Z_\eta) = \dim_{\xi}(Z_y) = \dim_{\xi}(X_y) = \dim_{\xi'}(X_\eta).
$$
因此，$Z_\eta$ 中某个经过 $\xi'$ 的不可约分支具有维数
$\dim_{\xi'}(X_\eta)$，这与 $Z_\eta$ 在 $X_\eta$ 中处处不稠密的假设
矛盾，结论得证。
\end{proof}

\begin{lemma}
\label{more-morphisms-lemma-scheme-theoretically-dense-generic-fibre}
设 $f : X \to Y$ 是有限型概形态射。假设 $Y$ 不可约，泛点为 $\eta$。
设 $U \subset X$ 是开子概形，且 $U_\eta$ 在 $X_\eta$ 中概形论稠密。
则存在非空开集 $V \subset Y$，使对所有 $y \in V$，$U_y$ 在 $X_y$
中概形论稠密。
\end{lemma}

\begin{proof}
设 $Y' \subset Y$ 是 $Y$ 的约化，并令 $X' = Y' \times_Y X$、
$U' = Y' \times_Y U$。由于 $Y' \to Y$ 在点上诱导双射，而
$U' \to U$ 与 $X' \to X$ 在概形论纤维上诱导同构，可以用 $Y'$ 替换
$Y$、用 $X'$ 替换 $X$。因此可假设 $Y$ 是整概形；见 性质章引理
\ref{properties-lemma-characterize-integral}。还可以用非空仿射开集替换
$Y$。换言之，可假设 $Y = \Spec(A)$，其中 $A$ 是分式域为 $K$ 的整环。

\medskip\noindent
由于 $f$ 有限型，$X$ 拟紧。对某些仿射开集 $X_i$，写成
$X = X_1 \cup \ldots \cup X_n$。由 态射章定义
\ref{morphisms-definition-scheme-theoretically-dense}，
$U_i = X_i \cap U$ 是 $X_i$ 的开子概形，且 $U_{i, \eta}$ 在
$X_{i, \eta}$ 中概形论稠密。因此只需对各对 $(X_i, U_i)$ 证明结论；
换言之，可假设 $X$ 仿射。

\medskip\noindent
写 $X = \Spec(B)$。注意，$B_K$ 是有限型 $K$-代数，故是诺特的。
因此 $U_\eta$ 拟紧。于是可找到有限多个 $g_1, \ldots, g_m \in B$，
使 $D(g_j) \subset U$，且
$U_\eta = D(g_1)_\eta \cup \ldots \cup D(g_m)_\eta$。$U_\eta$ 在
$X_\eta$ 中概形论稠密，意味着
$B_K \to \bigoplus_j (B_K)_{g_j}$ 是单射；见 态射章例
\ref{morphisms-example-scheme-theoretic-closure}。由 代数章引理
\ref{algebra-lemma-when-injective-covering}，这等价于
$B_K \to \bigoplus\nolimits_{j = 1, \ldots, m} B_K$，
$b \mapsto (g_1b, \ldots, g_mb)$ 是单射。令 $M$ 为此映射在 $A$
上的余核；也就是说，有正合列
$$
0 \to I \to B \xrightarrow{(g_1, \ldots, g_m)}
\bigoplus\nolimits_{j = 1, \ldots, m} B \to M \to 0
$$
对某个非零 $h$，用 $A_h$ 替换 $A$ 后，可以假设 $B$ 是平坦有限表示
$A$-代数，且 $M$ 在 $A$ 上平坦；见 代数章引理
\ref{algebra-lemma-generic-flatness}。$B$ 在 $A$ 上平坦，蕴含 $B$
作为 $A$-模无挠；见 代数进阶章引理
\ref{more-algebra-lemma-flat-torsion-free}。故 $B \subset B_K$。
由假设 $I_K = 0$，这蕴含 $I = 0$（因为
$I \subset B \subset B_K$ 是 $I_K$ 的子集）。于是现在有短正合列
$$
0 \to B \xrightarrow{(g_1, \ldots, g_m)}
\bigoplus\nolimits_{j = 1, \ldots, m} B \to M \to 0
$$
，其中 $M$ 在 $A$ 上平坦。因此，对每个到域 $\kappa$ 的同态
$A \to \kappa$，得到短正合列
$$
0 \to B \otimes_A \kappa \xrightarrow{(g_1 \otimes 1, \ldots, g_m \otimes 1)}
\bigoplus\nolimits_{j = 1, \ldots, m} B \otimes_A \kappa \to
M \otimes_A \kappa \to 0
$$
；见 代数章引理 \ref{algebra-lemma-flat-tor-zero}。反向应用上述论证，
这意味着 $\bigcup D(g_j \otimes 1)$ 在
$\Spec(B \otimes_A \kappa)$ 中概形论稠密。由于
$\bigcup D(g_j \otimes 1) = \bigcup D(g_j)_\kappa \subset U_\kappa$，
可知 $U_\kappa$ 在 $X_\kappa$ 中概形论稠密，这正是所要证明的。
\end{proof}

\noindent
设有概形态射 $f : X \to Y$ 与点 $y \in Y$。回忆纤维 $X_y$ 同胚于
$X$ 的子集 $f^{-1}(\{y\})$，后者赋予诱导拓扑；见 概形章引理
\ref{schemes-lemma-fibre-topological}。设有闭子集 $T(y) \subset X_y$。
令 $T$ 为 $T(y)$ 在 $X$ 中的闭包，并赋予 $T$ 诱导的约化概形结构。
则 $T$ 是 $X$ 的闭子概形，并且集合论意义下 $T_y = T(y)$。事实上，
$T$ 是具有此性质的 $X$ 的最小闭子概形。因此，若愿意，把 $X_y$
的闭子集记作 $T_y$ 是“无害的”。在下一个引理中，我们把这一点应用于
$f$ 的泛纤维。

\begin{lemma}
\label{more-morphisms-lemma-cover-generic-fibre-neighbourhood}
设 $f : X \to Y$ 是有限型概形态射。假设 $Y$ 不可约，泛点为 $\eta$。
设 $X_\eta = Z_{1, \eta} \cup \ldots \cup Z_{n, \eta}$ 是泛纤维的覆盖，
其中各项都是 $X_\eta$ 的闭子集。令 $Z_i$ 为 $Z_{i, \eta}$ 在 $X$
中的闭包（见上面的讨论）。则存在非空开集 $V \subset Y$，使对所有
$y \in V$，都有 $X_y = Z_{1, y} \cup \ldots \cup Z_{n, y}$。
\end{lemma}

\begin{proof}
若 $Y$ 诺特，则 $U = X \setminus (Z_1 \cup \ldots \cup Z_n)$ 在 $Y$
上有限型，可以直接应用引理 \ref{more-morphisms-lemma-empty-generic-fibre}，得到对某个
非空开集 $V \subset Y$ 有 $U_V = \emptyset$。一般情形下论证如下。
由于问题是拓扑性的，可以用其约化替换 $Y$。于是 $Y$ 是整概形；
见 性质章引理 \ref{properties-lemma-characterize-integral}。缩小 $Y$
后，可以假设 $X \to Y$ 平坦；见 态射章命题
\ref{morphisms-proposition-generic-flatness}。此时，由 态射章引理
\ref{morphisms-lemma-generalizations-lift-flat}，$X_y$ 中每个点 $x$
都是某个点 $x' \in X_\eta$ 的特化。由于各 $Z_i$ 在 $X$ 中闭，并覆盖
泛纤维，这蕴含对 $y \in Y$ 有 $X_y = \bigcup Z_{i, y}$，正合所求。
\end{proof}

\noindent
下一个引理说明，源约化的态射具有约化泛纤维。

\begin{lemma}
\label{more-morphisms-lemma-reduction-generic-fibre}
设 $f : X \to Y$ 是概形态射。设 $\eta \in Y$ 是 $Y$ 某个不可约分支
的泛点。则
$(X_\eta)_{red} = (X_{red})_\eta$.
\end{lemma}

\begin{proof}
选择 $\eta$ 的仿射邻域 $\Spec(A) \subset Y$，并选择经态射 $f$ 映入
$\Spec(A)$ 的仿射开集 $\Spec(B) \subset X$。设
$\mathfrak p \subset A$ 是 $\eta$ 对应的极小素理想。设 $B_{red}$
为 $B$ 对幂零根 $\sqrt{(0)}$ 的商。本引理的代数内容是
$C = B_{red} \otimes_A \kappa(\mathfrak p)$ 约化。以 $I \subset A$
表示幂零根，故 $A_{red} = A/I$。令
$\mathfrak p_{red} = \mathfrak p/I$；这是 $A_{red}$ 的极小素理想，且
$\kappa(\mathfrak p) = \kappa(\mathfrak p_{red})$。由于
$A \to B_{red}$ 与 $A \to \kappa(\mathfrak p)$ 都经由
$A \to A_{red}$ 分解，有
$C = B_{red} \otimes_{A_{red}} \kappa(\mathfrak p_{red})$。由 Algebra,
引理 \ref{algebra-lemma-minimal-prime-reduced-ring}，
$\kappa(\mathfrak p_{red}) = (A_{red})_{\mathfrak p_{red}}$ 是一个局部化。
故 $C$ 是 $B_{red}$ 的一个局部化（代数章引理
\ref{algebra-lemma-tensor-localization}），从而约化。
\end{proof}

\begin{lemma}
\label{more-morphisms-lemma-make-generic-fibre-geometrically-reduced}
设 $f : X \to Y$ 是概形态射。假设 $Y$ 不可约且 $f$ 有限型。
存在图
$$
\xymatrix{
X' \ar[d]_{f'} \ar[r]_{g'} & X_V \ar[r] \ar[d] & X \ar[d]^f \\
Y' \ar[r]^g & V \ar[r] & Y
}
$$
，其中
\begin{enumerate}
\item $V$ 是 $Y$ 的非空开集；
\item $X_V = V \times_Y X$；
\item $g : Y' \to V$ 是有限泛同胚；
\item $X' = (Y' \times_Y X)_{red} = (Y' \times_V X_V)_{red}$；
\item $g'$ 是有限泛同胚；
\item $Y'$ 是整仿射概形；
\item $f'$ 平坦且有限表示；
\item $f'$ 的泛纤维几何约化。
\end{enumerate}
\end{lemma}

\begin{proof}
设 $V = \Spec(A)$ 是 $Y$ 的非空仿射开集。由假设，$A$ 的幂零根
$\mathfrak p = \sqrt{(0)}$ 是素理想。令 $K = \kappa(\mathfrak p)$。
令 $p$ 在 $K$ 的特征为正时等于其特征，在特征为零时等于 $1$。
由 簇章引理
\ref{varieties-lemma-finite-extension-geometrically-reduced}，存在有限纯不可分
域扩张 $K'/K$，使 $(X_{K'})_{red}$ 在 $K'$ 上几何约化。选择元素
$x_1, \ldots, x_n \in K'$，使它们在 $K$ 上生成 $K'$，并且每个 $x_i$
的某个 $p$ 次幂属于 $A/\mathfrak p$。令 $A' \subset K'$ 为由
$x_1, \ldots, x_n$ 生成的 $K'$ 的有限 $A$-子代数。注意，$A'$ 是分式域为
$K'$ 的整环。由 代数章引理 \ref{algebra-lemma-p-ring-map}，
$A \to A'$ 在谱上诱导泛同胚。令 $Y' = \Spec(A')$，并令
$X' = (Y' \times_Y X)_{red}$。由引理
\ref{more-morphisms-lemma-reduction-generic-fibre}，$X' \to Y'$ 的泛纤维为
$(X_{K'})_{red}$，按构造它几何约化。注意，$X' \to X_V$ 是有限泛同胚，
因为它是约化态射 $X' \to Y' \times_Y X$（见 态射章引理
\ref{morphisms-lemma-reduction-universal-homeomorphism}）与 $g$ 的基变换的
复合。此时除可能的 (7) 外，本引理所有性质都成立。缩小 $Y'$，从而缩小
$V$，即可得到 (7)；见 态射章命题
\ref{morphisms-proposition-generic-flatness}。
\end{proof}

\begin{lemma}
\label{more-morphisms-lemma-make-components-generic-fibre-geometrically-irreducible}
设 $f : X \to Y$ 是概形态射。假设 $Y$ 不可约且 $f$ 有限型。
存在图
$$
\xymatrix{
X' \ar[d]_{f'} \ar[r]_{g'} & X_V \ar[r] \ar[d] & X \ar[d]^f \\
Y' \ar[r]^g & V \ar[r] & Y
}
$$
，其中
\begin{enumerate}
\item $V$ 是 $Y$ 的非空开集；
\item $X_V = V \times_Y X$；
\item $g : Y' \to V$ 是满有限平展态射；
\item $X' = Y' \times_Y X = Y' \times_V X_V$；
\item $g'$ 是满有限平展态射；
\item $Y'$ 是不可约仿射概形；
\item $f'$ 的泛纤维的所有不可约分支都几何不可约。
\end{enumerate}
\end{lemma}

\begin{proof}
设 $V = \Spec(A)$ 是 $Y$ 的非空仿射开集。由假设，$A$ 的幂零根
$\mathfrak p = \sqrt{(0)}$ 是素理想。令 $K = \kappa(\mathfrak p)$。
由 簇章引理
\ref{varieties-lemma-finite-extension-geometrically-irreducible-components}，
存在有限可分域扩张 $K'/K$，使 $X_{K'}$ 的所有不可约分支都在 $K'$
上几何不可约。选择在 $K$ 上生成 $K'$ 的元素 $\alpha \in K'$；见
域章引理 \ref{fields-lemma-primitive-element}。设
$P(T) \in K[T]$ 是 $\alpha$ 在 $K$ 上的极小多项式。对某个
$f \in A$、$f \not \in \mathfrak p$，用 $f \alpha$ 替换 $\alpha$ 后，
可以假设存在首一多项式
$T^d + a_1T^{d - 1} + \ldots + a_d \in A[T]$，它在映射
$A[T] \to K[T]$ 下映到 $P(T) \in K[T]$。
% P05-ERR-0038: frozen source quotients A[T] by P, previously defined in K[T].
令 $A' = A[T]/(P)$。则 $A \to A'$ 是有限自由环映射，并存在唯一的
位于 $\mathfrak p$ 上方的素理想 $\mathfrak q$，使
$K = \kappa(\mathfrak p) \subset \kappa(\mathfrak q) = K'$ 有限可分，
且 $\mathfrak pA'_{\mathfrak q}$ 是 $A'_{\mathfrak q}$ 的极大理想。
因此 $g : Y' = \Spec(A') \to V = \Spec(A)$ 在 $\mathfrak q$ 处平展；
见 代数章引理 \ref{algebra-lemma-characterize-etale}。这意味着存在开集
$W \subset \Spec(A')$，使 $g|_W : W \to \Spec(A)$ 平展。由于 $g$ 有限，
且 $\mathfrak q$ 是位于 $\mathfrak p$ 上方的唯一点，
$Z = g(Y' \setminus W)$ 是 $V$ 的不含 $\mathfrak p$ 的闭子集。
因此，用不与 $Z$ 相交的 $V$ 的一个主仿射开集替换 $V$ 后，$g$
成为有限平展态射。
\end{proof}







\section{相对伴随点}
\label{more-morphisms-section-assassin}

\begin{lemma}
\label{more-morphisms-lemma-relative-assassin-in-neighbourhood}
设 $f : X \to S$ 是概形态射，$\mathcal{F}$ 是拟凝聚
$\mathcal{O}_X$-模。设 $\xi \in \text{Ass}_{X/S}(\mathcal{F})$，并令
$Z = \overline{\{\xi\}} \subset X$。若 $f$ 局部有限型，且
$\mathcal{F}$ 是有限型 $\mathcal{O}_X$-模，则存在非空开集
$V \subset Z$，使对每个 $s \in f(V)$，$V_s$ 的泛点都属于
$\text{Ass}_{X/S}(\mathcal{F})$。
\end{lemma}

\begin{proof}
可以用 $f(\xi)$ 的一个仿射开邻域替换 $S$，并用 $\xi$ 的一个仿射开
邻域替换 $X$。故可假设 $S = \Spec(A)$、$X = \Spec(B)$，而 $f$ 由
有限型环映射 $A \to B$ 给出；见 态射章引理
\ref{morphisms-lemma-locally-finite-type-characterize}。此外，对某个有限
$B$-模 $M$，可写 $\mathcal{F} = \widetilde{M}$；见 性质章引理
\ref{properties-lemma-finite-type-module}。设 $\mathfrak q \subset B$
是 $\xi$ 对应的素理想，并设 $\mathfrak p \subset A$ 是 $A$ 中相应的
素理想。由假设，
$\mathfrak q \in \text{Ass}_B(M \otimes_A \kappa(\mathfrak p))$；见
代数章说明 \ref{algebra-remark-bourbaki} 与 除子章引理
\ref{divisors-lemma-associated-affine-open}。在此记号下，
$Z = V(\mathfrak q) \subset \Spec(B)$。特别地，
$f(Z) \subset V(\mathfrak p)$。因此，显然只需在分别用
$A/\mathfrak pA$、$B/\mathfrak pB$ 与 $M/\mathfrak pM$ 替换
$A$、$B$ 与 $M$ 后证明本引理。换言之，可假设 $A$ 是分式域为 $K$
的整环，且 $\mathfrak q \subset B$ 是 $M \otimes_A K$ 的伴随素理想。

\medskip\noindent
此时可以使用泛平坦性。确切地说，由 代数章引理
\ref{algebra-lemma-generic-flatness}，存在非零 $g \in A$，使 $M_g$
作为 $A_g$-模平坦。用 $A_g$ 替换 $A$ 后，可以假设 $M$ 作为 $A$-模
平坦。

\medskip\noindent
此时，由 代数章引理 \ref{algebra-lemma-post-bourbaki}，
$\mathfrak q$ 也是 $M$ 的伴随素理想。因此得到单射 $B$-模映射
$B/\mathfrak q \to M$。令 $Q$ 为其余核，得到有限 $B$-模的短正合列
$$
0 \to B/\mathfrak q \to M \to Q \to 0
$$
。再次应用泛平坦性 代数章引理
\ref{algebra-lemma-generic-flatness}，这次应用于 $B$-模 $Q$，可以假设
$Q$ 是平坦 $A$-模。特别地，可以假设上述短正合列泛单射；见 Algebra,
引理 \ref{algebra-lemma-flat-tor-zero}。在此情形下，对 $A$ 的任意素理想
$\mathfrak p'$，有
$(B/\mathfrak q) \otimes_A \kappa(\mathfrak p')
\subset M \otimes_A \kappa(\mathfrak p')$
。本引理随即得出，因为由 除子章引理
\ref{divisors-lemma-minimal-support-in-ass}，
$(B/\mathfrak q) \otimes_A \kappa(\mathfrak p')$ 的支集的极小素理想
$\mathfrak q'$ 是
$(B/\mathfrak q) \otimes_A \kappa(\mathfrak p')$ 的伴随素理想。
\end{proof}

\begin{lemma}
\label{more-morphisms-lemma-bad-case}
设 $f : X \to Y$ 是概形态射，$\mathcal{F}$ 是拟凝聚
$\mathcal{O}_X$-模，$U \subset X$ 是开子概形。假设
\begin{enumerate}
\item $f$ 有限型；
\item $\mathcal{F}$ 有限型；
\item $Y$ 不可约，泛点为 $\eta$；
\item $\text{Ass}_{X_\eta}(\mathcal{F}_\eta)$ 不包含于 $U_\eta$。
\end{enumerate}
则存在非空开子概形 $V \subset Y$，使对所有 $y \in V$，集合
$\text{Ass}_{X_y}(\mathcal{F}_y)$ 不包含于 $U_y$。
\end{lemma}

\begin{proof}
设 $Z \subset X$ 是 $\mathcal{F}$ 的概形论支集；见 态射章定义
\ref{morphisms-definition-scheme-theoretic-support}。则 $Z_\eta$ 是
$\mathcal{F}_\eta$ 的概形论支集（态射章引理
\ref{morphisms-lemma-flat-pullback-support}）。因此，由 除子章引理
\ref{divisors-lemma-minimal-support-in-ass}，$Z_\eta$ 各不可约分支的泛点
包含于 $\text{Ass}_{X_\eta}(\mathcal{F}_\eta)$。
% P05-ERR-0039: frozen "bad case" proof then asserts disjointness/empty T_eta;
% this conflicts with hypothesis (4) and appears to prove the opposite case.
由此可见 $Z_\eta \cap U_\eta = \emptyset$。因此
$T = Z \setminus U$ 是 $Z$ 的闭子集，且 $T_\eta = \emptyset$。若赋予
$T$ 诱导的约化概形结构，则 $T \to Y$ 是有限型态射。由引理
\ref{more-morphisms-lemma-empty-generic-fibre}，存在非空开集 $V \subset Y$，使
$T_V = \emptyset$。这个 $V$ 即为所求。
\end{proof}

\begin{lemma}
\label{more-morphisms-lemma-good-case}
设 $f : X \to Y$ 是概形态射，$\mathcal{F}$ 是拟凝聚
$\mathcal{O}_X$-模，$U \subset X$ 是开子概形。假设
\begin{enumerate}
\item $f$ 有限型；
\item $\mathcal{F}$ 有限型；
\item $Y$ 不可约，泛点为 $\eta$；
\item $\text{Ass}_{X_\eta}(\mathcal{F}_\eta) \subset U_\eta$。
\end{enumerate}
则存在非空开子概形 $V \subset Y$，使对所有 $y \in V$，都有
$\text{Ass}_{X_y}(\mathcal{F}_y) \subset U_y$。
\end{lemma}

\begin{proof}
（这个证明与引理 \ref{more-morphisms-lemma-scheme-theoretically-dense-generic-fibre}
的证明相同。我们建议读者先阅读那个证明。）由于陈述只涉及纤维，显然
可以用其约化替换 $Y$。因此可假设 $Y$ 是整概形；见 Properties,
引理 \ref{properties-lemma-characterize-integral}。还可假设
$Y = \Spec(A)$ 仿射。于是 $A$ 是分式域为 $K$ 的整环。

\medskip\noindent
由于 $f$ 有限型，$X$ 拟紧。对某些仿射开集 $X_i$，写成
$X = X_1 \cup \ldots \cup X_n$，并令
$\mathcal{F}_i = \mathcal{F}|_{X_i}$。由假设，
$U_i = X_i \cap U$ 的泛纤维包含
$\text{Ass}_{X_{i, \eta}}(\mathcal{F}_{i, \eta})$。因此只需对各三元组
$(X_i, \mathcal{F}_i, U_i)$ 证明结果；换言之，可假设 $X$ 仿射。

\medskip\noindent
写 $X = \Spec(B)$。设 $N$ 是满足 $\mathcal{F} = \widetilde{N}$ 的有限
$B$-模。注意，$B_K$ 是有限型 $K$-代数，故诺特。因此 $U_\eta$ 拟紧。
于是可找到有限多个 $g_1, \ldots, g_m \in B$，使 $D(g_j) \subset U$，
且 $U_\eta = D(g_1)_\eta \cup \ldots \cup D(g_m)_\eta$。由于
$\text{Ass}_{X_\eta}(\mathcal{F}_\eta) \subset U_\eta$，映射
$N_K \to \bigoplus_j (N_K)_{g_j}$ 是单射。由 代数章引理
\ref{algebra-lemma-when-injective-covering}，这等价于
$N_K \to \bigoplus\nolimits_{j = 1, \ldots, m} N_K$，
$n \mapsto (g_1n, \ldots, g_mn)$ 是单射。令 $I$ 与 $M$ 分别为此映射
在 $A$ 上的核与余核；也就是说，有正合列
$$
0 \to I \to N \xrightarrow{(g_1, \ldots, g_m)}
\bigoplus\nolimits_{j = 1, \ldots, m} N \to M \to 0
$$
对某个非零 $h$，用 $A_h$ 替换 $A$ 后，可以假设 $B$ 是平坦有限表示
$A$-代数，并且 $M$ 与 $N$ 都在 $A$ 上平坦；见 代数章引理
\ref{algebra-lemma-generic-flatness}。$N$ 在 $A$ 上平坦，蕴含 $N$
作为 $A$-模无挠；见 代数进阶章引理
\ref{more-algebra-lemma-flat-torsion-free}。故 $N \subset N_K$。
按构造 $I_K = 0$，这蕴含 $I = 0$（因为
$I \subset N \subset N_K$ 是 $I_K$ 的子集）。于是现在有短正合列
$$
0 \to N \xrightarrow{(g_1, \ldots, g_m)}
\bigoplus\nolimits_{j = 1, \ldots, m} N \to M \to 0
$$
，其中 $M$ 在 $A$ 上平坦。因此，对每个到域 $\kappa$ 的同态
$A \to \kappa$，得到短正合列
$$
0 \to N \otimes_A \kappa \xrightarrow{(g_1 \otimes 1, \ldots, g_m \otimes 1)}
\bigoplus\nolimits_{j = 1, \ldots, m} N \otimes_A \kappa \to
M \otimes_A \kappa \to 0
$$
；见 代数章引理 \ref{algebra-lemma-flat-tor-zero}。反向应用上述论证，
这意味着 $\bigcup D(g_j \otimes 1)$ 包含
$\text{Ass}_{B \otimes_A \kappa}(N \otimes_A \kappa)$。由于
$\bigcup D(g_j \otimes 1) = \bigcup D(g_j)_\kappa \subset U_\kappa$，
可知 $U_\kappa$ 包含
$\text{Ass}_{X \otimes \kappa}(\mathcal{F} \otimes \kappa)$，
这正是所要证明的。
\end{proof}

\begin{lemma}
\label{more-morphisms-lemma-base-change-assassin-in-U}
设 $f : X \to S$ 是局部有限型态射，$\mathcal{F}$ 是有限型拟凝聚
$\mathcal{O}_X$-模，而 $U \subset X$ 是开子概形。设 $g : S' \to S$
是概形态射，设 $f' : X' = X_{S'} \to S'$ 是 $f$ 的基变换，设
$g' : X' \to X$ 是投影，并令 $\mathcal{F}' = (g')^*\mathcal{F}$、
$U' = (g')^{-1}(U)$。最后，设 $s' \in S'$ 的像为 $s = g(s')$。
此时
$$
\text{Ass}_{X_s}(\mathcal{F}_s) \subset U_s
\Leftrightarrow
\text{Ass}_{X'_{s'}}(\mathcal{F}'_{s'}) \subset U'_{s'}.
$$
\end{lemma}

\begin{proof}
这立即由 除子章引理
\ref{divisors-lemma-base-change-relative-assassin} 得出。另见 Divisors,
说明 \ref{divisors-remark-base-change-relative-assassin}。
\end{proof}

\begin{lemma}
\label{more-morphisms-lemma-relative-assassin-constructible}
设 $f : X \to Y$ 是有限表示态射，$\mathcal{F}$ 是有限表示拟凝聚
$\mathcal{O}_X$-模。设 $U \subset X$ 是开子概形，且 $U \to Y$ 拟紧。
则集合
$$
E = \{y \in Y \mid \text{Ass}_{X_y}(\mathcal{F}_y) \subset U_y\}
$$
在 $Y$ 中局部可构造。
\end{lemma}

\begin{proof}
设 $y \in Y$。须证明存在 $y$ 在 $Y$ 中的开邻域 $V$，使
$E \cap V$ 在 $V$ 中可构造。因此可假设 $Y$ 仿射。写
$Y = \Spec(A)$，并把 $A = \colim A_i$ 写成有限型
$\mathbf{Z}$-代数的定向极限。由 极限章引理
\ref{limits-lemma-descend-finite-presentation}，可以找到 $i$ 以及有限表示态射
$f_i : X_i \to \Spec(A_i)$，其到 $Y$ 的基变换恢复 $f$。必要时增大 $i$，
可以假设存在有限表示拟凝聚 $\mathcal{O}_{X_i}$-模 $\mathcal{F}_i$，
其到 $X$ 的拉回同构于 $\mathcal{F}$；见 极限章引理
\ref{limits-lemma-descend-modules-finite-presentation}。必要时再增大一次 $i$，
可以假设存在开子概形 $U_i \subset X_i$，其在 $X$ 中的逆像为 $U$；
见 极限章引理 \ref{limits-lemma-descend-opens}。由引理
\ref{more-morphisms-lemma-base-change-assassin-in-U}，只需对 $f_i$ 证明本引理。因此问题
化到 $Y$ 是诺特环的谱的情形。

\medskip\noindent
在 $Y$ 为诺特概形时，我们用 拓扑章引理
\ref{topology-lemma-characterize-constructible-Noetherian} 的判别法证明 $E$
可构造。为此，设 $Z \subset Y$ 是不可约闭子概形。须证明
$E \cap Z$ 要么包含非空开子集，要么在 $Z$ 中不稠密。把引理
\ref{more-morphisms-lemma-bad-case} 与 \ref{more-morphisms-lemma-good-case} 应用于 $Z$ 上的基变换
$(X, \mathcal{F}, U) \times_Y Z$，即得结论。
\end{proof}











\section{约化纤维}
\label{more-morphisms-section-reduced}

\begin{lemma}
\label{more-morphisms-lemma-nonreduced-in-neighbourhood}
设 $f : X \to Y$ 是概形态射。假设 $Y$ 不可约，泛点为 $\eta$，且
$f$ 有限型。若 $X_\eta$ 非约化，则存在非空开集 $V \subset Y$，
使对所有 $y \in V$，纤维 $X_y$ 都非约化。
\end{lemma}

\begin{proof}
设 $Y' \subset Y$ 是 $Y$ 的约化，并设 $X' \to Y'$ 是 $f$ 的基变换。
注意，$Y' \to Y$ 在点上诱导双射，而 $X' \to X$ 识别各纤维。
% P05-ERR-0040: frozen source says assume Y' reduced after replacement, likely Y.
故可假设 $Y'$ 约化，也就是整；见 性质章引理
\ref{properties-lemma-characterize-integral}。还可用仿射开集替换 $Y$，
故可假设 $Y = \Spec(A)$，其中 $A$ 是整环。以 $K$ 表示 $A$ 的分式域。
选择仿射开集 $\Spec(B) = U \subset X$ 以及非零幂零截面
$h_\eta \in \Gamma(U_\eta, \mathcal{O}_{U_\eta}) = B_K$。缩小 $Y$ 后，
可以假设 $h$ 来自 $h \in \Gamma(U, \mathcal{O}_U) = B$。再稍微缩小
$Y$ 后，可以假设 $h$ 幂零。令
$I = \{b \in B \mid hb = 0\}$ 为 $h$ 的零化子。则 $C = B/I$ 是有限型
$A$-代数，其泛纤维 $(B/I)_K$ 非零（因为 $h_\eta \not = 0$）。
把泛平坦性应用于 $A \to C$ 与 $A \to B/hB$；见 代数章引理
\ref{algebra-lemma-generic-flatness}。得到 $g \in A$、$g \not = 0$，使
$C_g$ 作为 $A_g$-模自由，而 $(B/hB)_g$ 作为 $A_g$-模平坦。用
$D(g) \subset Y$ 替换 $Y$。现在有短正合列
$$
0 \to C \to B \to B/hB \to 0.
$$
，其中 $B/hB$ 在 $A$ 上平坦，而 $C$ 是非零自由 $A$-模。由此可知，
对到域的任意同态 $A \to \kappa$，环 $C \otimes_A \kappa$ 非零，且序列
$$
0 \to C \otimes_A \kappa \to B \otimes_A \kappa \to
B/hB \otimes_A \kappa \to 0
$$
正合；见 代数章引理 \ref{algebra-lemma-flat-tor-zero}。注意，由张量积的
右正合性，
$B/hB \otimes_A \kappa = (B \otimes_A \kappa) / h(B \otimes_A \kappa)$。
因此乘以 $h$ 在 $B \otimes_A \kappa$ 上不是零映射。这显然意味着，
对任意点 $y \in Y$，元素 $h$ 在 $U_y$ 上限制为非零元素，故 $X_y$
非约化。
\end{proof}

\begin{lemma}
\label{more-morphisms-lemma-base-change-fibres-geometrically-reduced}
设 $f : X \to Y$ 是概形态射。设 $g : Y' \to Y$ 是任意态射，并以
$f' : X' \to Y'$ 表示 $f$ 的基变换。则
\begin{align*}
\{y' \in Y' \mid X'_{y'}\text{ is geometrically reduced}\} \\
= g^{-1}(\{y \in Y \mid X_y\text{ is geometrically reduced}\}).
\end{align*}
\end{lemma}

\begin{proof}
这归结为如下陈述：若 $y' \in Y'$ 的像为 $y \in Y$，则纤维
$X'_{y'} = X_y \times_y y'$ 在 $\kappa(y')$ 上几何约化，当且仅当
$X_y$ 在 $\kappa(y)$ 上几何约化。这由 簇章引理
\ref{varieties-lemma-geometrically-reduced-upstairs} 得出。
\end{proof}

\begin{lemma}
\label{more-morphisms-lemma-not-geometrically-reduced-in-neighbourhood}
设 $f : X \to Y$ 是概形态射。假设 $Y$ 不可约，泛点为 $\eta$，且
$f$ 有限型。若 $X_\eta$ 非几何约化，则存在非空开集 $V \subset Y$，
使对所有 $y \in V$，纤维 $X_y$ 都非几何约化。
\end{lemma}

\begin{proof}
应用引理 \ref{more-morphisms-lemma-make-generic-fibre-geometrically-reduced}，得到
$$
\xymatrix{
X' \ar[d]_{f'} \ar[r]_{g'} & X_V \ar[r] \ar[d] & X \ar[d]^f \\
Y' \ar[r]^g & V \ar[r] & Y
}
$$
并具有该引理中列出的全部性质。设 $\eta'$ 是 $Y'$ 的泛点。考虑态射
$X' \to X_{Y'}$（即约化态射）以及由此得到的泛纤维态射
$X'_{\eta'} \to X_{\eta'}$。由于 $X'_{\eta'}$ 几何约化，而 $X_\eta$
不是几何约化，此态射不可能是同构；见 簇章引理
\ref{varieties-lemma-geometrically-reduced-upstairs}。故 $X_{\eta'}$ 非约化。
因此由引理 \ref{more-morphisms-lemma-nonreduced-in-neighbourhood}，对某个非空开集
$V' \subset Y'$ 的所有点 $y' \in V'$，$X_{Y'} \to Y'$ 的纤维都非约化。
由于 $g : Y' \to V$ 是同胚，引理
\ref{more-morphisms-lemma-base-change-fibres-geometrically-reduced} 说明 $g(V')$
就是所求开集。
\end{proof}

\begin{lemma}
\label{more-morphisms-lemma-geometrically-reduced-generic-fibre}
设 $f : X \to Y$ 是概形态射。假设
\begin{enumerate}
\item $Y$ 不可约，泛点为 $\eta$；
\item $X_\eta$ 几何约化；
\item $f$ 有限型。
\end{enumerate}
则存在非空开子概形 $V \subset Y$，使 $X_V \to V$ 具有几何约化纤维。
\end{lemma}

\begin{proof}
设 $Y' \subset Y$ 是 $Y$ 的约化，并设 $X' \to Y'$ 是 $f$ 的基变换。
注意，$Y' \to Y$ 在点上诱导双射，而 $X' \to X$ 识别各纤维。
% P05-ERR-0041: frozen source again says assume Y' reduced, likely Y.
故可假设 $Y'$ 约化，也就是整；见 性质章引理
\ref{properties-lemma-characterize-integral}。还可用仿射开集替换 $Y$，
故可假设 $Y = \Spec(A)$，其中 $A$ 是整环。以 $K$ 表示 $A$ 的分式域。
稍微缩小 $Y$ 后，还可假设 $X \to Y$ 平坦且有限表示；见 Morphisms,
命题 \ref{morphisms-proposition-generic-flatness}。

\medskip\noindent
由于 $X_\eta$ 几何约化，存在开稠密子集 $V \subset X_\eta$，使
$V \to \Spec(K)$ 光滑；见 簇章引理
\ref{varieties-lemma-geometrically-reduced-dense-smooth-open}。令
$U \subset X$ 为 $f$ 光滑的点所成的集合。由 态射章引理
\ref{morphisms-lemma-set-points-where-fibres-smooth}，有
$V \subset U_\eta$。因此 $U$ 的泛纤维在 $X$ 的泛纤维中稠密。由于
$X_\eta$ 约化，$U_\eta$ 在 $X_\eta$ 中概形论稠密；见 Morphisms,
引理 \ref{morphisms-lemma-reduced-scheme-theoretically-dense}。注意，由于
$U \to Y$ 光滑，$U \to Y$ 的所有纤维都几何约化。因此只需证明：缩小
$Y$ 后，对所有 $y \in Y$，概形 $U_y$ 在 $X_y$ 中概形论稠密；见
态射章引理 \ref{morphisms-lemma-reduced-subscheme-closure}。
这由引理 \ref{more-morphisms-lemma-scheme-theoretically-dense-generic-fibre} 得出。
\end{proof}

\begin{lemma}
\label{more-morphisms-lemma-geometrically-reduced-constructible}
设 $f : X \to Y$ 是拟紧且局部有限表示的态射。则集合
$$
E = \{y \in Y \mid X_y\text{ is geometrically reduced}\}
$$
在 $Y$ 中局部可构造。
\end{lemma}

\begin{proof}
设 $y \in Y$。须证明存在 $y$ 在 $Y$ 中的开邻域 $V$，使
$E \cap V$ 在 $V$ 中可构造。因此可假设 $Y$ 仿射。于是 $X$ 拟紧。
选择有限仿射开覆盖 $X = U_1 \cup \ldots \cup U_n$。各 $U_i \to Y$
在 $y$ 处的纤维构成 $X \to Y$ 在 $y$ 处纤维的仿射开覆盖。因此还可
假设 $X$ 仿射。写 $Y = \Spec(A)$，并把 $A = \colim A_i$ 写成有限型
$\mathbf{Z}$-代数的定向极限。由 极限章引理
\ref{limits-lemma-descend-finite-presentation}，可以找到 $i$ 与有限表示态射
$f_i : X_i \to \Spec(A_i)$，其到 $Y$ 的基变换恢复 $f$。由引理
\ref{more-morphisms-lemma-base-change-fibres-geometrically-reduced}，只需对 $f_i$
证明本引理。因此问题化到 $Y$ 是诺特环的谱的情形。

\medskip\noindent
在 $Y$ 为诺特概形时，我们用 拓扑章引理
\ref{topology-lemma-characterize-constructible-Noetherian} 的判别法证明 $E$
可构造。为此，设 $Z \subset Y$ 是泛点为 $\xi$ 的不可约闭子概形。
须证明 $E \cap Z$ 要么包含非空开子集，要么在 $Z$ 中不稠密。
若 $X_\xi$ 几何约化，则把引理
\ref{more-morphisms-lemma-geometrically-reduced-generic-fibre} 应用于态射
$X_Z \to Z$，可知对某个非空开集 $V \subset Z$，所有纤维 $X_y$
都几何约化。若 $X_\xi$ 非几何约化，则把引理
\ref{more-morphisms-lemma-not-geometrically-reduced-in-neighbourhood} 应用于态射
$X_Z \to Z$，可知对某个非空开集 $V \subset Z$，所有纤维 $X_y$
都非几何约化。结论得证。
\end{proof}

\begin{lemma}
\label{more-morphisms-lemma-proper-flat-over-dvr-reduced-fibre}
设 $f : X \to \Spec(R)$ 是固有态射，其中 $R$ 是离散赋值环。假设
$X$ 的每个不可约分支都支配 $\Spec(R)$（例如，当 $f$ 平坦时），
并假设特殊纤维约化。则 $X$ 与泛纤维 $X_\eta$ 都约化。
\end{lemma}

\begin{proof}
（平坦态射满足关于不可约分支的假设，例如可由 除子章引理
\ref{divisors-lemma-bourbaki} 或平坦环映射的下降性质得出。）假设特殊纤维
$X_s$ 约化。设 $x \in X$ 是任意点。我们证明 $\mathcal{O}_{X, x}$
约化；这将证明 $X$ 与 $X_\eta$ 约化。选择特化 $x \leadsto x'$，其中
$x'$ 位于特殊纤维中；由于固有态射是闭的，这样的特化存在。考虑局部环
$A = \mathcal{O}_{X, x'}$。由于 $\mathcal{O}_{X, x}$ 是 $A$ 的局部化，
只需证明 $A$ 约化。设 $\pi \in R$ 是一个素元。若 $a \in A$ 非零，
则存在 $n \geq 0$ 与元素 $a' \in A$，使 $a = \pi^n a'$ 且
$a' \not \in \pi A$。这由 Krull 交定理（代数章引理
\ref{algebra-lemma-intersect-powers-ideal-module-zero}）得出。
% P05-ERR-0042: frozen source has typo "If a is nilpotent the a is".
若 $a$ 幂零，则 $a$ 属于 $A$ 的所有极小素理想；而 $a'$ 也如此，
因为由关于 $X$ 的不可约分支的假设，$\pi$ 不属于 $A$ 的任何极小素理想。
故 $a'$ 幂零。但 $a'$ 在约化环
$A/\pi A = \mathcal{O}_{X_s, x'}$ 中的像非零。这产生矛盾，除非 $A$
约化，而这正是所要证明的。
\end{proof}

\begin{lemma}
\label{more-morphisms-lemma-geometrically-reduced-open}
设 $f : X \to Y$ 是有限表示平坦固有态射。则集合\linebreak[4]
$\{y \in Y \mid X_y\text{ is geometrically reduced}\}$ 在 $Y$ 中是开的。
\end{lemma}

\begin{proof}
可假设 $Y$ 仿射。于是 $Y$ 是有限型仿射 $\mathbf{Z}$-概形的余滤过极限。
故可假设 $X \to Y$ 是 $X_0 \to Y_0$ 的基变换，其中 $Y_0$ 是有限型
$\mathbf{Z}$-代数的谱，而 $X_0 \to Y_0$ 平坦且固有。见 极限章引理
\ref{limits-lemma-descend-finite-presentation}、
\ref{limits-lemma-descend-flat-finite-presentation} 与
\ref{limits-lemma-eventually-proper}。由于纤维几何约化的点集的形成与
基变换交换（引理 \ref{more-morphisms-lemma-base-change-fibres-geometrically-reduced}），
可以假设基底诺特。

\medskip\noindent
假设 $Y$ 诺特。由引理
\ref{more-morphisms-lemma-geometrically-reduced-constructible}，该集合可构造。因此只需
证明它在一般化下稳定（拓扑章引理
\ref{topology-lemma-characterize-closed-Noetherian}）。由 性质章引理
\ref{properties-lemma-locally-Noetherian-specialization-dvr}，问题化到如下
情形：$Y = \Spec(R)$，$R$ 是离散赋值环，且闭纤维 $X_y$ 几何约化。
须证泛纤维 $X_\eta$ 几何约化。

\medskip\noindent
若不然，则存在 $R$ 的分式域的有限扩张 $L$，使 $X_L$ 非约化；见
簇章引理 \ref{varieties-lemma-geometrically-reduced}。存在支配 $R$
的离散赋值环 $R' \subset L$，其分式域为 $L$；见 代数章引理
\ref{algebra-lemma-integral-closure-Dedekind}。用 $R'$ 替换 $R$ 后，
问题归结为引理 \ref{more-morphisms-lemma-proper-flat-over-dvr-reduced-fibre}。
\end{proof}







\section{纤维的不可约分支}
\label{more-morphisms-section-irreducible}

\begin{lemma}
\label{more-morphisms-lemma-irreducible-components-in-neighbourhood}
设 $f : X \to Y$ 是概形态射。假设 $Y$ 不可约，泛点为 $\eta$，且
$f$ 有限型。若 $X_\eta$ 有 $n$ 个不可约分支，则存在非空开集
$V \subset Y$，使对所有 $y \in V$，纤维 $X_y$ 至少有 $n$ 个不可约
分支。
\end{lemma}

\begin{proof}
由于问题纯粹是拓扑性的，可以用 $X$ 与 $Y$ 的约化分别替换它们。
特别地，这蕴含 $Y$ 是整概形；见 性质章引理
\ref{properties-lemma-characterize-integral}。设
$X_\eta = X_{1, \eta} \cup \ldots \cup X_{n, \eta}$ 是 $X_\eta$
的不可约分支分解。设 $X_i \subset X$ 是泛纤维为 $X_{i, \eta}$
的约化闭子概形。注意，$Z_{i, j} = X_i \cap X_j$ 是 $X_i$ 的闭子集，
其泛纤维 $Z_{i, j, \eta}$ 在 $X_{i, \eta}$ 中处处不稠密。因此缩小
$Y$ 后，可以假设对每个 $y \in Y$，$Z_{i, j, y}$ 在 $X_{i, y}$ 中
处处不稠密；见引理 \ref{more-morphisms-lemma-nowhere-dense-generic-fibre}。进一步缩小
$Y$ 后，可以假设对 $y \in Y$ 有 $X_y = \bigcup X_{i, y}$；见引理
\ref{more-morphisms-lemma-cover-generic-fibre-neighbourhood}。此外，缩小 $Y$ 后，可以
假设每个 $X_i \to Y$ 都平坦且有限表示；见 态射章命题
\ref{morphisms-proposition-generic-flatness}。态射 $X_i \to Y$ 是开的；见
态射章引理 \ref{morphisms-lemma-fppf-open}。因此存在 $\eta$ 的开邻域
$V$，对每个 $i$ 该邻域都包含于 $f(X_i)$。对每个 $y \in V$，概形
$X_{i, y}$ 是 $X_y$ 的非空闭子集，有 $X_y = \bigcup X_{i, y}$，且交
$Z_{i, j, y} = X_{i, y} \cap X_{j, y}$ 在 $X_{i, y}$ 中不稠密。
这显然蕴含 $X_y$ 至少有 $n$ 个不可约分支。
\end{proof}

\begin{lemma}
\label{more-morphisms-lemma-base-change-fibres-geometrically-irreducible}
设 $f : X \to Y$ 是概形态射。设 $g : Y' \to Y$ 是任意态射，并以
$f' : X' \to Y'$ 表示 $f$ 的基变换。则
\begin{align*}
\{y' \in Y' \mid X'_{y'}\text{ is geometrically irreducible}\} \\
= g^{-1}(\{y \in Y \mid X_y\text{ is geometrically irreducible}\}).
\end{align*}
\end{lemma}

\begin{proof}
这归结为如下陈述：若 $y' \in Y'$ 的像为 $y \in Y$，则纤维
$X'_{y'} = X_y \times_y y'$ 在 $\kappa(y')$ 上几何不可约，当且仅当
$X_y$ 在 $\kappa(y)$ 上几何不可约。这由 簇章引理
\ref{varieties-lemma-geometrically-irreducible-check-after-extension} 得出。
\end{proof}

\begin{lemma}
\label{more-morphisms-lemma-base-change-fibres-nr-geometrically-irreducible-components}
设 $f : X \to Y$ 是概形态射。设
$$
n_{X/Y} : Y \to \{0, 1, 2, 3, \ldots, \infty\}
$$
是如下函数：它把 $y \in Y$ 映到 $(X_y)_K$ 的不可约分支数，其中 $K$
是 $\kappa(y)$ 的可分闭扩张。此函数良定义；若 $g : Y' \to Y$ 是态射，
则
$$
n_{X'/Y'} = n_{X/Y} \circ g
$$
，其中 $X' \to Y'$ 是 $f$ 的基变换。
\end{lemma}

\begin{proof}
假设 $y' \in Y'$ 的像为 $y \in Y$。假设
$K \supset \kappa(y)$ 与 $K' \supset \kappa(y')$ 是可分闭扩张。
则可以选择域的交换图
$$
\xymatrix{
K \ar[r] & K'' & K' \ar[l] \\
\kappa(y) \ar[u] \ar[rr] & & \kappa(y') \ar[u]
}
$$
结论随即得出，因为概形态射
$$
\xymatrix{
(X'_{y'})_{K'} &
(X'_{y'})_{K''} = (X_y)_{K''} \ar[l] \ar[r] &
(X_y)_K
}
$$
在不可约分支之间诱导双射；见 簇章引理
\ref{varieties-lemma-separably-closed-field-irreducible-components}。
\end{proof}

\begin{lemma}
\label{more-morphisms-lemma-irreducible-polynomial-over-domain}
设 $A$ 是分式域为 $K$ 的整环，$P \in A[x_1, \ldots, x_n]$。
以 $\overline{K}$ 表示 $K$ 的代数闭包。假设 $P$ 在
$\overline{K}[x_1, \ldots, x_n]$ 中不可约。则存在 $f \in A$，使对所有
到域的同态 $\varphi : A_f \to \kappa$，多项式
$P^\varphi \in \kappa[x_1, \ldots, x_n]$ 都不可约。
\end{lemma}

\begin{proof}
存在 $A$ 上的 $A[x_1, \ldots, x_n]$ 的自同构 $\Psi$，使
$\Psi(P) = ax_n^d +$ 关于 $x_n$ 的低次项，其中 $a \not = 0$；见
代数章引理 \ref{algebra-lemma-helper-polynomial}。可以用 $\Psi(P)$
替换 $P$，并用 $A_a$ 替换 $A$。因此可假设 $P$ 是关于 $x_n$ 的次数
为 $d > 0$ 的首一多项式。对 $i = 1, \ldots, n - 1$，令 $d_i$
为 $P$ 关于 $x_i$ 的次数。注意，这蕴含：对每个到域 $\kappa$ 的同态
$\varphi : A \to \kappa$，$P^\varphi$ 是关于 $x_n$ 的次数为 $d$
的首一多项式，并且关于 $x_i$ 的次数为 $\leq d_i$。因此若
$P^\varphi$ 可约，则可以写成
$$
P^\varphi = Q_1 Q_2
$$
% P05-ERR-0044: frozen proof allows e_1 or e_2 = 0, hence trivial factors.
，其中 $Q_1, Q_2$ 关于 $x_n$ 分别为次数 $e_1, e_2 \geq 0$ 的首一
多项式，满足 $e_1 + e_2 = d$，并且对 $i = 1, \ldots, n - 1$，
关于 $x_i$ 的次数均为 $\leq d_i$。换言之，可以写成
\begin{equation}
\label{more-morphisms-equation-factors}
Q_j = x_n^{e_j} + \sum\nolimits_{0 \leq l < e_j}
\left( \sum\nolimits_{L \in \mathcal{L}} a_{j, l, L} x^L \right) x_n^l
\end{equation}
，其中求和遍历形如 $L = (l_1, \ldots, l_{n - 1})$ 且满足
$0 \leq l_i \leq d_i$ 的多重指标 $L$ 所成的集合 $\mathcal{L}$。
对任意满足 $e_1 + e_2 = d$ 的 $e_1, e_2 \geq 0$，考虑 $A$-代数
$$
B_{e_1, e_2} =
A[\{a_{1, l, L}\}_{0 \leq l < e_1, L \in \mathcal{L}},
\{a_{2, l, L}\}_{0 \leq l < e_2, L \in \mathcal{L}}]/(\text{relations})
$$
% P05-ERR-0045: frozen source has subject agreement "relations is".
，其中 $(\text{relations})$ 是由多项式
$$
P - Q_1Q_2 \in
A[\{a_{1, l, L}\}_{0 \leq l < e_1, L \in \mathcal{L}},
\{a_{2, l, L}\}_{0 \leq l < e_2, L \in \mathcal{L}}][x_1, \ldots, x_n]
$$
的系数生成的理想，这里 $Q_1$ 与 $Q_2$ 如
(\ref{more-morphisms-equation-factors}) 中定义。
% P05-ERR-0046: frozen source contains the editorial fragment "OK, and".
好，并且 $P$ 在 $\overline{K}$ 上不可约这一假设蕴含不存在任何
$A$-代数同态 $B_{e_1, e_2} \to \overline{K}$。由 Hilbert 零点定理，
见 代数章定理 \ref{algebra-theorem-nullstellensatz}，这意味着
$B_{e_1, e_2} \otimes_A K = 0$。由于 $B_{e_1, e_2}$ 是有限生成
$A$-代数，这表示可以找到 $f_{e_1, e_2} \in A$，使
$(B_{e_1, e_2})_{f_{e_1, e_2}} = 0$。按构造，这意味着：若
$\varphi : A_{f_{e_1, e_2}} \to \kappa$ 是到域的同态，则
$P^\varphi$ 不存在分解 $P^\varphi = Q_1 Q_2$，其中 $Q_1$ 关于
$x_n$ 的次数为 $e_1$，而 $Q_2$ 关于 $x_n$ 的次数为 $e_2$。
因此取 $f = \prod_{e_1, e_2 \geq 0, e_1 + e_2 = d} f_{e_1, e_2}$ 即得结论。
\end{proof}

\begin{lemma}
\label{more-morphisms-lemma-geom-irreducible-generic-fibre}
设 $f : X \to Y$ 是概形态射。假设
\begin{enumerate}
\item $Y$ 不可约，泛点为 $\eta$；
\item $X_\eta$ 几何不可约；
\item $f$ 有限型。
\end{enumerate}
则存在非空开子概形 $V \subset Y$，使 $X_V \to V$ 具有几何不可约纤维。
\end{lemma}

\begin{proof}[引理 \ref{more-morphisms-lemma-geom-irreducible-generic-fibre} 的第一种证明]
我们给出本引理的两种证明。二者本质等价；第二种更自足，但稍长。
选择图
$$
\xymatrix{
X' \ar[d]_{f'} \ar[r]_{g'} & X_V \ar[r] \ar[d] & X \ar[d]^f \\
Y' \ar[r]^g & V \ar[r] & Y
}
$$
，如引理 \ref{more-morphisms-lemma-make-generic-fibre-geometrically-reduced} 中所述。
注意，$f'$ 的泛纤维是 $f$ 的泛纤维的约化（见引理
\ref{more-morphisms-lemma-reduction-generic-fibre}），故几何不可约。假设本引理对态射
$f'$ 成立。则缩小 $V$ 后，$f'$ 的所有纤维都几何不可约。由于
$X' = (Y' \times_V X_V)_{red}$，这蕴含 $Y' \times_V X_V$ 的所有纤维
都几何不可约。因此由引理
\ref{more-morphisms-lemma-base-change-fibres-geometrically-irreducible}，
$X_V \to V$ 的所有纤维都几何不可约，结论得证。这样便可假设泛纤维
既几何约化又几何不可约，并可假设 $Y = \Spec(A)$，其中 $A$ 是整环。

\medskip\noindent
设 $x \in X_\eta$ 是泛点。由于 $X_\eta$ 几何不可约且约化，
$L = \kappa(x)$ 是 $K = \kappa(\eta)$ 的有限生成扩张，并且几何约化、
几何不可约；见 簇章引理
\ref{varieties-lemma-geometrically-reduced-at-point} 与
\ref{varieties-lemma-geometrically-irreducible-function-field}。特别地，
域扩张 $L/K$ 可分；见 代数章引理
\ref{algebra-lemma-characterize-separable-field-extensions}。因此可以找到
$x_1, \ldots, x_{r + 1} \in L$，使它们在 $K$ 上生成 $L$，并且
$x_1, \ldots, x_r$ 是 $L$ 在 $K$ 上的一组超越基；见 代数章引理
\ref{algebra-lemma-generating-finitely-generated-separable-field-extensions}。
设 $P \in K(x_1, \ldots, x_r)[T]$ 是 $x_{r + 1}$ 的极小多项式。
清除分母后，可以假设 $P$ 的系数属于 $A[x_1, \ldots, x_r]$。
注意，由于 $L$ 在 $K$ 上几何约化且几何不可约，多项式 $P$ 在
$\overline{K}[x_1, \ldots, x_r, T]$ 中不可约，其中 $\overline{K}$
是 $K$ 的代数闭包。记
$$
B' = A[x_1, \ldots, x_{r + 1}]/(P(x_{r + 1}))
$$
，并令 $X' = \Spec(B')$。按构造，$B'$ 的分式域作为 $K$-扩张同构于
$L = \kappa(x)$。因此存在开集 $U \subset X$、开集 $U' \subset X'$
以及 $Y$-同构 $U \to U'$；见 态射章引理
\ref{morphisms-lemma-common-open}。图示如下：
$$
\xymatrix{
X \ar[rd] &
U \ar[l] \ar@{=}[r] \ar[d] &
U' \ar[r] \ar[d] &
X' \ar[ld] \ar@{=}[r] & \Spec(B') \\
& Y \ar@{=}[r] & Y &
}
$$
注意，$U_\eta \subset X_\eta$ 与 $U'_\eta \subset X'_\eta$ 都是稠密开集。
因此，应用引理 \ref{more-morphisms-lemma-nowhere-dense-generic-fibre} 缩小 $Y$ 后，
对所有 $y \in Y$，$U_y$ 在 $X_y$ 中稠密，且 $U'_y$ 在 $X'_y$
中稠密。因此只需对 $X' \to Y$ 证明本引理，而这正是引理
\ref{more-morphisms-lemma-irreducible-polynomial-over-domain} 的内容。
\end{proof}

\begin{proof}[引理 \ref{more-morphisms-lemma-geom-irreducible-generic-fibre} 的第二种证明]
设 $Y' \subset Y$ 是 $Y$ 的约化，并设 $X' \to X$ 是 $X$ 的约化。
注意，$X' \to X  \to Y$ 经由 $Y'$ 分解；见 概形章引理
\ref{schemes-lemma-map-into-reduction}。由 态射章引理
\ref{morphisms-lemma-reduction-universal-homeomorphism}，$Y' \to Y$
与 $X' \to X$ 都是泛同胚，故只需对 $X' \to Y'$ 证明本引理。因此可
假设 $X$ 与 $Y$ 都约化。特别地，$Y$ 是整概形；见 性质章引理
\ref{properties-lemma-characterize-integral}。于是由 态射章命题
\ref{morphisms-proposition-generic-flatness}，存在非空仿射开集
$V \subset Y$，使 $X_V \to V$ 平坦且有限表示。用 $V$ 替换 $Y$ 后，
除 (1)、(2)、(3) 外，还可假设 $Y$ 是整仿射概形、$X$ 约化，而 $f$
平坦且有限表示。特别地，$f$ 泛开；见 态射章引理
\ref{morphisms-lemma-fppf-open}。

\medskip\noindent
选择非空仿射开集 $U \subset X$。则 $U \to Y$ 平坦且有限表示，泛纤维
几何不可约。补集 $X_\eta \setminus U_\eta$ 处处不稠密。因此缩小 $Y$
后，可以假设对所有 $y \in Y$，$U_y \subset X_y$ 是开稠密子集；见
引理 \ref{more-morphisms-lemma-nowhere-dense-generic-fibre}。故可以用 $U$ 替换 $X$，
把问题化到如下情形：$Y$ 是整仿射概形，$X$ 是在 $Y$ 上平坦且有限表示
的约化仿射概形，其泛纤维 $X_\eta$ 几何不可约。

\medskip\noindent
写 $X = \Spec(B)$、$Y = \Spec(A)$。则 $A$ 是整环，$B$ 约化，
$A \to B$ 是有限表示平坦映射，而 $B_K$ 在 $A$ 的分式域 $K$ 上几何
不可约。特别地，$B_K$ 是整环。设 $L$ 为 $B_K$ 的分式域。注意，由于
$B$ 是有限表示 $A$-代数，$L$ 是 $K$ 的有限生成域扩张。设 $K'/K$
是有限纯不可分扩张，使 $(L \otimes_K K')_{red}$ 是可分生成域扩张；
见 代数章引理 \ref{algebra-lemma-make-separable}。选择
$x_1, \ldots, x_n \in K'$，使它们在 $K$ 上生成域扩张 $K'$，并且对某个
素数幂 $q_i$ 有 $x_i^{q_i} \in A$（略去 $x_i$ 存在性的证明）。
设 $A'$ 为由 $x_1, \ldots, x_n$ 生成的 $K'$ 的 $A$-子代数。则 $A'$
是有限 $A$-子代数 $A' \subset K'$，其分式域为 $K'$。注意，
$\Spec(A') \to \Spec(A)$ 是泛同胚；见 代数章引理
\ref{algebra-lemma-p-ring-map}。因此只需在基变换到 $\Spec(A')$ 后证明
结论。我们将用 $A'$ 替换 $A$，并用 $(B \otimes_A A')_{red}$ 替换
$B$，从而达到 $L$ 是 $K$ 的可分生成域扩张的情形。当然，
$(B \otimes_A A')_{red}$ 可能不再平坦，或不再在 $A'$ 上有限表示；
但对合适的 $f \in A'$，用 $A'_f$ 替换 $A'$ 即可补救这一点；见
代数章引理 \ref{algebra-lemma-generic-flatness}。

\medskip\noindent
此时已知：$A$ 是整环，$B$ 约化，$A \to B$ 平坦且有限表示，$B_K$
是整环，其分式域 $L$ 是 $A$ 的分式域 $K$ 的可分生成域扩张。由
代数章引理
\ref{algebra-lemma-generating-finitely-generated-separable-field-extensions}，
可以写 $L = K(x_1, \ldots, x_{r + 1})$，其中
$x_1, \ldots, x_r$ 在 $K$ 上代数无关，而 $x_{r + 1}$ 在
$K(x_1, \ldots, x_r)$ 上可分。清除分母后，可以假设 $x_{r + 1}$
在 $K(x_1, \ldots, x_r)$ 上的极小多项式
$P \in K(x_1, \ldots, x_r)[T]$ 的系数属于
$A[x_1, \ldots, x_r]$。注意，由于 $L/K$ 可分，且 $L$ 在 $K$
上几何不可约，多项式 $P$ 在 $K$ 的代数闭包 $\overline{K}$ 上不可约。
记
$$
B' = A[x_1, \ldots, x_{r + 1}]/(P(x_{r + 1})).
$$
按构造，$B$ 与 $B'$ 的分式域作为 $K$-扩张同构。因此，对合适的
$h \in B$ 与 $h' \in B'$，存在 $A$-代数同构
$B_h \cong B'_{h'}$；见 态射章引理
\ref{morphisms-lemma-common-open}。换言之，$X$ 与
$X' = \Spec(B')$ 有一个共同的仿射开集 $U$。图示如下：
$$
\xymatrix{
X = \Spec(B) \ar[rd] &
U \ar[l] \ar[r] \ar[d] &
\Spec(B') = X' \ar[ld] \\
& Y = \Spec(A) &
}
$$
再次缩小 $Y$（把引理 \ref{more-morphisms-lemma-nowhere-dense-generic-fibre} 分别应用于
$X$ 中的 $Z = X \setminus U$ 与 $X'$ 中的 $Z' = X' \setminus U$）后，
对所有 $y \in Y$，$U_y$ 在 $X_y$ 中稠密，且 $U_y$ 在 $X'_y$
中稠密。因此只需对 $X' \to Y$ 证明本引理，而这正是引理
\ref{more-morphisms-lemma-irreducible-polynomial-over-domain} 的内容。
\end{proof}

\begin{lemma}
\label{more-morphisms-lemma-nr-geom-irreducible-components-good}
设 $f : X \to Y$ 是概形态射。设 $n_{X/Y}$ 是引理
\ref{more-morphisms-lemma-base-change-fibres-nr-geometrically-irreducible-components}
中引入的 $Y$ 上的函数，它计数 $f$ 的纤维的几何不可约分支数。
假设 $f$ 有限型。设 $y \in Y$ 是一点。则存在非空开集
$V \subset \overline{\{y\}}$，使 $n_{X/Y}|_V$ 为常数。
\end{lemma}

\begin{proof}
在 $\overline{\{y\}}$ 上赋予诱导的约化概形结构，并记为 $Z$。
设 $f_Z : X_Z \to Z$ 是 $f$ 的基变换。显然，只需对 $f_Z$ 与 $Z$
的泛点证明本引理。因此可假设 $Y$ 是整概形；见 性质章引理
\ref{properties-lemma-characterize-integral}。此时的目标是构造非空开集
$V \subset Y$，使 $n_{X/Y}|_V$ 为常数。

\medskip\noindent
把引理
\ref{more-morphisms-lemma-make-components-generic-fibre-geometrically-irreducible}
应用于 $f : X \to Y$，得到 $g : Y' \to V \subset Y$。由于
$g : Y' \to V$ 是满有限平展态射，特别地是开态射（见 态射章引理
\ref{morphisms-lemma-etale-open}），只需证明存在开集 $V' \subset Y'$，
使 $n_{X'/Y'}|_{V'}$ 为常数；见引理
\ref{more-morphisms-lemma-base-change-fibres-nr-geometrically-irreducible-components}。
因此可假设泛纤维 $X_\eta$ 的所有不可约分支都在 $\kappa(\eta)$ 上
几何不可约。

\medskip\noindent
此时，假设
$X_\eta = X_{1, \eta} \bigcup \ldots \bigcup X_{n, \eta}$
是泛纤维的（几何）不可约分支分解。特别地，
$n_{X/Y}(\eta) = n$。设 $X_i$ 是 $X_{i, \eta}$ 在 $X$ 中的闭包。
缩小 $Y$ 后，可以假设 $X = \bigcup X_i$；见引理
\ref{more-morphisms-lemma-cover-generic-fibre-neighbourhood}。进一步缩小 $Y$ 后，
$f$ 的每个纤维至少有 $n$ 个不可约分支；见引理
\ref{more-morphisms-lemma-irreducible-components-in-neighbourhood}。故对所有 $y \in Y$，
$n_{X/Y}(y) \geq n$。再次缩小 $Y$ 后，对每个 $i$ 与所有 $y \in Y$，
$X_{i, y}$ 几何不可约；见引理
\ref{more-morphisms-lemma-geom-irreducible-generic-fibre}。由于
$X_y = \bigcup X_{i, y}$，这说明 $n_{X/Y}(y) \leq n$，证明完成。
\end{proof}

\begin{lemma}
\label{more-morphisms-lemma-nr-geom-irreducible-components-constructible}
设 $f : X \to Y$ 是概形态射。设 $n_{X/Y}$ 是引理
\ref{more-morphisms-lemma-base-change-fibres-nr-geometrically-irreducible-components}
中引入的 $Y$ 上的函数，它计数 $f$ 的纤维的几何不可约分支数。
假设 $f$ 有限表示。则 $n_{X/Y}$ 的水平集
$$
E_n = \{y \in Y \mid n_{X/Y}(y) = n\}
$$
在 $Y$ 中局部可构造。
\end{lemma}

\begin{proof}
固定 $n$，并设 $y \in Y$。须证明存在 $y$ 在 $Y$ 中的开邻域 $V$，
使 $E_n \cap V$ 在 $V$ 中可构造。因此可假设 $Y$ 仿射。写
$Y = \Spec(A)$，并把 $A = \colim A_i$ 写成有限型
$\mathbf{Z}$-代数的定向极限。由 极限章引理
\ref{limits-lemma-descend-finite-presentation}，可以找到 $i$ 与有限表示态射
$f_i : X_i \to \Spec(A_i)$，其到 $Y$ 的基变换恢复 $f$。由引理
\ref{more-morphisms-lemma-base-change-fibres-nr-geometrically-irreducible-components}，
只需对 $f_i$ 证明本引理。因此问题化到 $Y$ 是诺特环的谱的情形。

\medskip\noindent
在 $Y$ 为诺特概形时，我们用 拓扑章引理
\ref{topology-lemma-characterize-constructible-Noetherian} 的判别法证明
$E_n$ 可构造。为此，设 $Z \subset Y$ 是不可约闭子概形。须证明
$E_n \cap Z$ 要么包含非空开子集，要么在 $Z$ 中不稠密。设
$\xi \in Z$ 是泛点。引理
\ref{more-morphisms-lemma-nr-geom-irreducible-components-good} 表明，$n_{X/Y}$ 在 $Z$
中 $\xi$ 的一个邻域上为常数。这显然蕴含所需结论。
\end{proof}













\section{纤维的连通分支}
\label{more-morphisms-section-connected}

\begin{lemma}
\label{more-morphisms-lemma-connected-components-in-neighbourhood}
设 $f : X \to Y$ 是概形态射。假设 $Y$ 不可约、泛点为 $\eta$，且
$f$ 是有限型的。若 $X_\eta$ 有 $n$ 个连通分支，则存在非空开子集
$V \subset Y$，使得对所有 $y \in V$，纤维 $X_y$ 至少有 $n$ 个连通分支。
\end{lemma}

\begin{proof}
由于问题纯粹是拓扑性的，可将 $X$ 和 $Y$ 换成其约化。特别地，这说明
$Y$ 是整概形，见 性质章引理
\ref{properties-lemma-characterize-integral}。设
$X_\eta = X_{1, \eta} \cup \ldots \cup X_{n, \eta}$ 是 $X_\eta$
的连通分支分解。令 $X_i \subset X$ 为泛纤维是 $X_{i, \eta}$ 的约化闭子概形。
注意，$Z_{i, j} = X_i \cap X_j$ 是 $X$ 的闭子集，且其泛纤维
$Z_{i, j, \eta}$ 为空。因此缩小 $Y$ 后，可假设
$Z_{i, j} = \emptyset$，见引理 \ref{more-morphisms-lemma-empty-generic-fibre}。进一步缩小
$Y$ 后，可假设对 $y \in Y$ 有 $X_y = \bigcup X_{i, y}$，见引理
\ref{more-morphisms-lemma-cover-generic-fibre-neighbourhood}。此外，缩小 $Y$ 后，可假设每个
$X_i \to Y$ 都平坦且有限表现，见 态射章命题
\ref{morphisms-proposition-generic-flatness}。态射 $X_i \to Y$ 是开态射，见
态射章引理 \ref{morphisms-lemma-fppf-open}。因此存在 $\eta$ 的开邻域
$V$，它对每个 $i$ 都包含于 $f(X_i)$。对每个 $y \in V$，概形
$X_{i, y}$ 是 $X_y$ 的非空闭子集，且 $X_y = \bigcup X_{i, y}$，而交
$Z_{i, j, y} = X_{i, y} \cap X_{j, y}$ 为空！显然，这说明 $X_y$ 至少有
$n$ 个连通分支。
\end{proof}

\begin{lemma}
\label{more-morphisms-lemma-base-change-fibres-geometrically-connected}
设 $f : X \to Y$ 是概形态射。任取态射 $g : Y' \to Y$，并将 $f$ 的基变换记为
$f' : X' \to Y'$。则
\begin{align*}
\{y' \in Y' \mid X'_{y'}\text{ 几何连通}\} \\
= g^{-1}(\{y \in Y \mid X_y\text{ 几何连通}\}).
\end{align*}
\end{lemma}

\begin{proof}
这归结为如下陈述：若 $y' \in Y'$ 的像为 $y \in Y$，则纤维
$X'_{y'} = X_y \times_y y'$ 在 $\kappa(y')$ 上几何连通，当且仅当 $X_y$
在 $\kappa(y)$ 上几何连通。这由 簇章引理
\ref{varieties-lemma-geometrically-connected-check-after-extension} 得出。
\end{proof}

\begin{lemma}
\label{more-morphisms-lemma-base-change-fibres-nr-geometrically-connected-components}
设 $f : X \to Y$ 是概形态射。令
$$
n_{X/Y} : Y \to \{0, 1, 2, 3, \ldots, \infty\}
$$
为如下函数：对 $y \in Y$，它取值为 $(X_y)_K$ 的连通分支数，其中 $K$
是 $\kappa(y)$ 的一个可分闭扩张。此定义良定；若 $g : Y' \to Y$ 是态射，则
$$
n_{X'/Y'} = n_{X/Y} \circ g
$$
其中 $X' \to Y'$ 是 $f$ 的基变换。
\end{lemma}

\begin{proof}
设 $y' \in Y'$ 的像为 $y \in Y$。设 $K \supset \kappa(y)$ 和
$K' \supset \kappa(y')$ 是可分闭扩张。于是可以选取域的交换图
$$
\xymatrix{
K \ar[r] & K'' & K' \ar[l] \\
\kappa(y) \ar[u] \ar[rr] & & \kappa(y') \ar[u]
}
$$
下列概形态射诱导连通分支之间的双射，因而得到结论：
$$
\xymatrix{
(X'_{y'})_{K'} &
(X'_{y'})_{K''} = (X_y)_{K''} \ar[l] \ar[r] &
(X_y)_K
}
$$
见 簇章引理
\ref{varieties-lemma-separably-closed-field-connected-components}。
\end{proof}

\begin{lemma}
\label{more-morphisms-lemma-geometrically-connected-generic-fibre}
设 $f : X \to Y$ 是概形态射。假设
\begin{enumerate}
\item $Y$ 不可约，泛点为 $\eta$；
\item $X_\eta$ 几何连通；并且
\item $f$ 是有限型的。
\end{enumerate}
则存在非空开子概形 $V \subset Y$，使得 $X_V \to V$ 的纤维几何连通。
\end{lemma}

\begin{proof}
按引理 \ref{more-morphisms-lemma-make-components-generic-fibre-geometrically-irreducible}
选取图
$$
\xymatrix{
X' \ar[d]_{f'} \ar[r]_{g'} & X_V \ar[r] \ar[d] & X \ar[d]^f \\
Y' \ar[r]^g & V \ar[r] & Y
}
$$
。注意 $f'$ 的泛纤维几何连通（例如由引理
\ref{more-morphisms-lemma-base-change-fibres-nr-geometrically-connected-components}）。假设本引理对态射
$f'$ 成立。这意味着存在非空开子集 $W \subset Y'$，使得 $X' \to Y'$ 在
$W$ 上的每个纤维都几何连通。由 态射章引理
\ref{morphisms-lemma-etale-open}，$g$ 是开态射，故 $f$ 在非空开子集
$V = g(W)$ 的各点处的纤维全都几何连通，见引理
\ref{more-morphisms-lemma-base-change-fibres-nr-geometrically-connected-components}。由此可见，可以假设泛纤维
$X_\eta$ 的不可约分支均几何不可约。

\medskip\noindent
令 $Y'$ 为 $Y$ 的约化，并置 $X' = Y' \times_Y X$。只需对态射
$X' \to Y'$ 证明本引理（例如再次使用引理
\ref{more-morphisms-lemma-base-change-fibres-nr-geometrically-connected-components}）。由于 $X' \to Y'$
的泛纤维与 $X \to Y$ 的泛纤维相同，可假设 $Y$ 不可约且约化（即为整概形，见
性质章引理 \ref{properties-lemma-characterize-integral}），并且泛纤维
$X_\eta$ 的不可约分支均几何不可约。

\medskip\noindent
现在设
$X_\eta = X_{1, \eta} \bigcup \ldots \bigcup X_{n, \eta}$
是泛纤维的（几何）不可约分支分解。令 $X_i$ 为 $X_{i, \eta}$ 在 $X$ 中的闭包。
缩小 $Y$ 后，由引理 \ref{more-morphisms-lemma-cover-generic-fibre-neighbourhood} 可假设
$X = \bigcup X_i$。令 $Z_{i, j} = X_i \cap X_j$。作分解
$$
\{1, \ldots, n\} \times \{1, \ldots, n\} = I \amalg J
$$
，其中当 $Z_{i, j, \eta} = \emptyset$ 时 $(i, j) \in I$，而当
$Z_{i, j, \eta} \not = \emptyset$ 时 $(i, j) \in J$。缩小 $Y$ 后，由引理
\ref{more-morphisms-lemma-empty-generic-fibre} 可假设对所有 $(i, j) \in I$ 都有
$Z_{i, j} = \emptyset$。再缩小 $Y$ 后，由引理
\ref{more-morphisms-lemma-geom-irreducible-generic-fibre}，对每个 $i$ 以及所有 $y \in Y$，
$X_{i, y}$ 都几何不可约。进一步缩小 $Y$，由 态射章命题
\ref{morphisms-proposition-generic-flatness}，可使每个 $Z_{i, j} \to Y$ 对所有
$(i, j) \in J$ 都平坦且有限表现。这意味着 $f(Z_{i, j}) \subset Y$ 是开集，见
态射章引理 \ref{morphisms-lemma-fppf-open}。我们断言
$$
V  = \bigcap\nolimits_{(i, j) \in J} f(Z_{i, j})
$$
满足要求，即对每个 $y \in V$，$X_y$ 都几何连通。事实上，$X_\eta$ 连通说明
由 $J$ 中各对生成的等价关系只有一个等价类。现在若 $y \in V$，且
$K \supset \kappa(y)$ 是可分闭扩张，则 $(X_y)_K$ 的不可约分支就是各纤维
$(X_{i, y})_K$。此外，由构造及 $y \in V$ 可见，$(X_{i, y})_K$ 与
$(X_{j, y})_K$ 相交，当且仅当 $(i, j) \in J$。因此关于等价类的上述说明表明
$(X_y)_K$ 连通，结论成立。
\end{proof}

\begin{lemma}
\label{more-morphisms-lemma-nr-geom-connected-components-good}
设 $f : X \to Y$ 是概形态射。令 $n_{X/Y}$ 为引理
\ref{more-morphisms-lemma-base-change-fibres-nr-geometrically-connected-components} 中引入的、在 $Y$
上计算 $f$ 的纤维之几何连通分支数的函数。假设 $f$ 是有限型的。设
$y \in Y$ 是一点。则存在非空开子集 $V \subset \overline{\{y\}}$，使得
$n_{X/Y}|_V$ 为常数。
\end{lemma}

\begin{proof}
令 $Z$ 为 $\overline{\{y\}}$ 上的诱导约化概形结构。令 $f_Z : X_Z \to Z$
为 $f$ 的基变换。显然，只需对 $f_Z$ 和 $Z$ 的泛点证明本引理。因此可以假设
$Y$ 是整概形，见 性质章引理
\ref{properties-lemma-characterize-integral}。此时目标是构造非空开子集
$V \subset Y$，使得 $n_{X/Y}|_V$ 为常数。

\medskip\noindent
对 $f : X \to Y$ 应用引理
\ref{more-morphisms-lemma-make-components-generic-fibre-geometrically-irreducible}，得到
$g : Y' \to V \subset Y$。由于 $g : Y' \to V$ 是满的有限 \'etale 态射，特别地是开态射
（见 态射章引理 \ref{morphisms-lemma-etale-open}），只需证明存在开子集
$V' \subset Y'$，使得 $n_{X'/Y'}|_{V'}$ 为常数，见
% P05-ERR-CH37-0047: Frozen source cites the irreducible-component base-change lemma here although the argument concerns connected components.
引理 \ref{more-morphisms-lemma-base-change-fibres-nr-geometrically-irreducible-components}.
因此可假设泛纤维 $X_\eta$ 的所有不可约分支在 $\kappa(\eta)$ 上几何不可约。由
簇章引理
\ref{varieties-lemma-irreducible-components-geometrically-irreducible}
可知，$X_\eta$ 的连通分支也几何连通。

\medskip\noindent
现在设
$X_\eta = X_{1, \eta} \bigcup \ldots \bigcup X_{n, \eta}$
是泛纤维的（几何）连通分支分解。特别地，$n_{X/Y}(\eta) = n$。令 $X_i$
为 $X_{i, \eta}$ 在 $X$ 中的闭包。缩小 $Y$ 后，由引理
\ref{more-morphisms-lemma-cover-generic-fibre-neighbourhood} 可假设 $X = \bigcup X_i$。进一步缩小
$Y$ 后，由引理 \ref{more-morphisms-lemma-connected-components-in-neighbourhood} 可知 $f$ 的每个纤维至少有
$n$ 个连通分支。因此对所有 $y \in Y$，均有 $n_{X/Y}(y) \geq n$。再次缩小
$Y$ 后，由引理 \ref{more-morphisms-lemma-geometrically-connected-generic-fibre}，对每个 $i$ 和所有
$y \in Y$，$X_{i, y}$ 都几何连通。由于 $X_y = \bigcup X_{i, y}$，这说明
$n_{X/Y}(y) \leq n$，证明完成。
\end{proof}

\begin{lemma}
\label{more-morphisms-lemma-nr-geom-connected-components-constructible}
设 $f : X \to Y$ 是概形态射。令 $n_{X/Y}$ 为引理
\ref{more-morphisms-lemma-base-change-fibres-nr-geometrically-connected-components} 中引入的、在 $Y$
上计算 $f$ 的纤维之几何连通分支数的函数。假设 $f$ 是有限表现的。则
$n_{X/Y}$ 的水平集
$$
E_n = \{y \in Y \mid n_{X/Y}(y) = n\}
$$
在 $Y$ 中局部可构造。
\end{lemma}

\begin{proof}
固定 $n$。设 $y \in Y$。须证明在 $Y$ 中存在 $y$ 的开邻域 $V$，使得
$E_n \cap V$ 在 $V$ 中可构造。因此可假设 $Y$ 仿射。写成
$Y = \Spec(A)$，并将 $A = \colim A_i$ 写成有限型 $\mathbf{Z}$-代数的有向极限。
由 极限章引理 \ref{limits-lemma-descend-finite-presentation}，可以找到 $i$ 及有限表现态射
$f_i : X_i \to \Spec(A_i)$，使其到 $Y$ 的基变换恢复 $f$。由引理
\ref{more-morphisms-lemma-base-change-fibres-nr-geometrically-connected-components}，只需对 $f_i$
证明本引理。因此问题化到 $Y$ 是诺特环的谱的情形。

\medskip\noindent
在 $Y$ 为诺特概形时，我们用 拓扑章引理
\ref{topology-lemma-characterize-constructible-Noetherian} 的判别法证明 $E_n$ 可构造。
为此，设 $Z \subset Y$ 是不可约闭子概形。须证明 $E_n \cap Z$ 要么包含非空开子集，
要么在 $Z$ 中不稠密。令 $\xi \in Z$ 为泛点。引理
\ref{more-morphisms-lemma-nr-geom-connected-components-good} 表明，$n_{X/Y}$ 在 $Z$ 中 $\xi$
的一个邻域上为常数。这显然蕴含所需结论。
\end{proof}

\begin{lemma}
\label{more-morphisms-lemma-connected-flat-over-dvr}
\begin{slogan}
不连通概形的平坦退化要么不连通，要么非约化。
\end{slogan}
设 $f : X \to S$ 是概形态射。假设
\begin{enumerate}
\item $S$ 是离散赋值环的谱；
\item $f$ 平坦；
\item $X$ 连通；
\item 闭纤维 $X_s$ 约化。
\end{enumerate}
则泛纤维 $X_\eta$ 连通。
\end{lemma}

\begin{proof}
写成 $S = \Spec(R)$，并令 $\pi \in R$ 为一致化元。反设 $X_\eta$ 不连通以导出矛盾。
这意味着存在非平凡幂等元 $e \in \Gamma(X_\eta, \mathcal{O}_{X_\eta})$。任取
$X$ 中的仿射开集 $U = \Spec(A)$。由于 $A$ 在 $R$ 上平坦，$\pi$ 是 $A$
上的非零因子，例如见 代数进阶章引理
\ref{more-algebra-lemma-flat-torsion-free}。于是 $e|_{U_\eta}$ 对应于一个元素
$e \in A[1/\pi]$。取 $z \in A$，使得 $e = z/\pi^n$，其中 $n \geq 0$ 最小。
注意 $z^2 = \pi^nz$。若 $n > 0$，这说明 $z \bmod \pi A$ 是幂零元。按假设，
$A/\pi A$ 约化，故 $n$ 的最小性蕴含 $n = 0$。因此 $e \in A$！换言之，
$e \in \Gamma(X, \mathcal{O}_X)$。由于 $X$ 连通，$e$ 必为平凡幂等元，矛盾。
\end{proof}







\section{与截面相交的连通分支}
\label{more-morphisms-section-connected-components}

\noindent
本节结果特别适用于群概形 $G \to S$ 及其单位截面 $e : S \to G$。

\begin{situation}
\label{more-morphisms-situation-connected-along-section}
% P05-ERR-CH37-0048: Frozen source has the ungrammatical phrase ``Here f ... be a morphism''.
这里 $f : X \to Y$ 是概形态射，而 $s : Y \to X$ 是 $f$ 的一个截面。对每个
$y \in Y$，以 $X^0_y$ 表示 $X_y$ 中包含 $s(y)$ 的连通分支。最后，置
$X^0 = \bigcup_{y \in Y} X^0_y$。
\end{situation}

\begin{lemma}
\label{more-morphisms-lemma-base-change-connected-along-section}
设 $f : X \to Y$、$s : Y \to X$ 如情形
\ref{more-morphisms-situation-connected-along-section}。若 $g : Y' \to Y$ 是任意态射，考虑基变换图
$$
\xymatrix{
X' \ar[r]_{g'} \ar[d]^{f'} & X \ar[d]_f \\
Y' \ar@/^1pc/[u]^{s'} \ar[r]^g & Y \ar@/_1pc/[u]_s
}
$$
，由此得到 $(X')^0 \subset X'$。则 $(X')^0 = (g')^{-1}(X^0)$。
\end{lemma}

\begin{proof}
设 $y' \in Y'$ 的像为 $y \in Y$。可将 $X^0_y$ 视为 $X_y$ 的闭子概形，例如见
态射章定义
\ref{morphisms-definition-scheme-structure-connected-component}。由于
$s(y) \in X^0_y$，簇章引理
\ref{varieties-lemma-geometrically-connected-if-connected-and-point} 表明，$X_y^0$
是 $\kappa(y)$ 上的几何连通概形。因此 $X_y^0 \times_y y' \to X'_{y'}$
是包含 $s'(y')$ 的连通闭子概形，从而
$X_y^0 \times_y y' \subset (X'_{y'})^0$。反向包含
$X_y^0 \times_y y' \supset (X'_{y'})^0$ 是显然的，因为 $(X'_{y'})^0$
在 $X_y$ 中的像是 $X_y$ 中包含 $s(y)$ 的连通子集。
\end{proof}

\begin{lemma}
\label{more-morphisms-lemma-connected-along-section-good}
设 $f : X \to Y$、$s : Y \to X$ 如情形
\ref{more-morphisms-situation-connected-along-section}。假设 $f$ 是有限型的，并设 $y \in Y$
为一点。则存在非空开子集 $V \subset \overline{\{y\}}$，使得 $X^0$
在基变换 $X_V$ 中的逆像在 $X_V$ 中既开又闭。
\end{lemma}

\begin{proof}
令 $Z \subset Y$ 为 $\overline{\{y\}}$ 上的诱导约化闭子概形结构。令
$f_Z : X_Z \to Z$ 和 $s_Z : Z \to X_Z$ 为 $f$ 和 $s$ 的基变换。由引理
\ref{more-morphisms-lemma-base-change-connected-along-section}，有 $(X_Z)^0 = (X^0)_Z$。因此只需对态射
% P05-ERR-CH37-0049: Frozen source introduces an undefined point x in X_Z although the lemma's distinguished point lies in the base Z.
$X_Z \to Z$ 以及映到 $Z$ 的泛点的点 $x \in X_Z$ 证明本引理。换言之，问题已化到
$Y$ 是泛点为 $\eta$ 的整概形的情形（见 性质章引理
\ref{properties-lemma-characterize-integral}）。目标是证明，缩小 $Y$ 后，子集 $X^0$
成为 $X$ 的既开又闭子集。

\medskip\noindent
注意，概形 $X_\eta$ 在一个域上有限型，因而是诺特概形。因此其连通分支既开又闭。
故可写成 $X_\eta = X_\eta^0 \amalg T_\eta$，其中 $T_\eta$ 是 $X_\eta$
的既开又闭子集。接着，令 $T \subset X$ 为 $T_\eta$ 的闭包，并令
$X^{00} \subset X$ 为 $X_\eta^0$ 的闭包。注意，$T_\eta$（相应地 $X^0_\eta$）
是 $T$（相应地 $X^{00}$）的泛纤维，见引理
\ref{more-morphisms-lemma-cover-generic-fibre-neighbourhood} 之前的讨论。此外，该引理表明，缩小 $Y$
后，可假设 $X = X^{00} \cup T$（在集合论意义下）。注意
$(T \cap X^{00})_\eta = T_\eta \cap X^0_\eta = \emptyset$。因此缩小 $Y$ 后，由引理
\ref{more-morphisms-lemma-empty-generic-fibre} 可假设 $T \cap X^{00} = \emptyset$。特别地，
$X^{00}$ 在 $X$ 中是开集。注意 $X^0_\eta$ 连通且有有理点，即 $s(\eta)$，故它几何连通，见
簇章引理
\ref{varieties-lemma-geometrically-connected-if-connected-and-point}。于是缩小 $Y$ 后，
由引理 \ref{more-morphisms-lemma-geometrically-connected-generic-fibre} 可假设 $X^{00} \to Y$
的所有纤维都几何连通。此时可知各纤维 $X^{00}_y$ 是 $X_y$ 中包含 $s(y)$
的开、闭且连通的子集。因此 $X^0 = X^{00}$，结论成立。
\end{proof}

\begin{lemma}
\label{more-morphisms-lemma-connected-along-section-locally-constructible}
设 $f : X \to Y$、$s : Y \to X$ 如情形
\ref{more-morphisms-situation-connected-along-section}。若 $f$ 是有限表现的，则 $X^0$ 在 $X$
中局部可构造。
\end{lemma}

\begin{proof}
设 $x \in X$。须证明存在 $x$ 的开邻域 $U$，使得 $X^0 \cap U$ 在 $U$
中可构造。这将问题化到 $Y$ 为仿射的情形。写成 $Y = \Spec(A)$，并将
$A = \colim A_i$ 写成有限型 $\mathbf{Z}$-代数的有向极限。由 极限章引理
\ref{limits-lemma-descend-finite-presentation}，可以找到 $i$ 及有限表现态射
$f_i : X_i \to \Spec(A_i)$，并配备截面 $s_i : \Spec(A_i) \to X_i$，使得到
$Y$ 的基变换恢复 $f$ 和截面 $s$。由引理
\ref{more-morphisms-lemma-base-change-connected-along-section}，只需对 $f_i, s_i$ 证明本引理。因此问题化到
$Y$ 是诺特环的谱的情形。

\medskip\noindent
假设 $Y$ 是诺特仿射概形。由于 $f$ 是有限表现的，因而是有限型的，故 $X$
也是诺特概形，见 态射章引理
\ref{morphisms-lemma-finite-type-noetherian}。为在此情形证明本引理，只需证明：对每个不可约闭子集
$Z \subset X$，交 $Z \cap X^0$ 要么包含 $Z$ 的一个非空开集，要么在 $Z$
中不稠密；见 拓扑章引理
\ref{topology-lemma-characterize-constructible-Noetherian}。令 $x \in Z$ 为泛点，并令
$y = f(x)$。由引理 \ref{more-morphisms-lemma-connected-along-section-good}，存在非空开子集
$V \subset \overline{\{y\}}$，使得 $X^0 \cap X_V$ 在 $X_V$ 中既开又闭。由于
$f(Z) \subset \overline{\{y\}}$ 且 $f(x) = y \in V$，可见
$W = f^{-1}(V) \cap Z$ 是 $Z$ 的非空开子集。因此 $X^0 \cap W$ 在 $W$
中既开又闭。由于 $W$ 不可约，$X^0 \cap W$ 要么为空，要么等于 $W$。本引理得证。
\end{proof}

\begin{lemma}
\label{more-morphisms-lemma-connected-along-section-open-neighbourhood}
设 $f : X \to Y$、$s : Y \to X$ 如情形
\ref{more-morphisms-situation-connected-along-section}。设 $y \in Y$ 为一点。假设
\begin{enumerate}
\item $f$ 有限表现且平坦；并且
\item 纤维 $X_y$ 几何约化。
\end{enumerate}
则 $X^0$ 是 $X$ 中 $X^0_y$ 的一个邻域。
\end{lemma}

\begin{proof}
可将 $Y$ 换成 $y$ 的一个仿射开邻域。写成 $Y = \Spec(A)$，并将
$A = \colim A_i$ 写成有限型 $\mathbf{Z}$-代数的有向极限。由 极限章引理
\ref{limits-lemma-descend-finite-presentation}，可以找到 $i$ 及有限表现态射
$f_i : X_i \to \Spec(A_i)$，并配备截面 $s_i : \Spec(A_i) \to X_i$，使得到
$Y$ 的基变换恢复 $f$ 和截面 $s$。必要时增大 $i$ 后，还可假设 $f_i$ 平坦，见
极限章引理 \ref{limits-lemma-descend-flat-finite-presentation}。令 $y_i$ 为 $y$
在 $Y_i$ 中的像。注意 $X_y = (X_{i, y_i}) \times_{y_i} y$。因此
$X_{i, y_i}$ 几何约化，见 簇章引理
\ref{varieties-lemma-geometrically-reduced-upstairs}。由引理
\ref{more-morphisms-lemma-base-change-connected-along-section}，只需对系统 $f_i, s_i, y_i \in Y_i$
证明本引理。因此问题化到 $Y$ 是诺特环的谱的情形。

\medskip\noindent
假设 $Y$ 是诺特环的谱。由于 $f$ 是有限表现的，因而是有限型的，故 $X$
也是诺特概形，见 态射章引理
\ref{morphisms-lemma-finite-type-noetherian}。设 $x \in X^0$ 是位于 $y$ 上方的一点。
由 拓扑章引理
\ref{topology-lemma-constructible-neighbourhood-Noetherian}，只需证明：对任何经过 $x$
的不可约闭子集 $Z \subset X$，交 $X^0 \cap Z$ 在 $Z$ 中稠密。特别地，只需证明
泛点 $x' \in Z$ 属于 $X^0$。由 性质章引理
\ref{properties-lemma-locally-Noetherian-specialization-dvr}，可以找到离散赋值环 $R$
和态射 $\Spec(R) \to X$，它把闭点映到 $x$，把泛点映到 $x'$。通过复合
$\Spec(R) \to X \to Y$，将 $\Spec(R)$ 视为 $Y$ 上的概形。由引理
\ref{more-morphisms-lemma-base-change-connected-along-section}，$(X_R)^0$ 是 $X^0$ 的逆像。按构造，结构态射
$X_R \to \Spec(R)$ 除 $s$ 的基变换 $s_R$ 外还有第二个截面
$t : \Spec(R) \to X_R$，使得 $t(\eta_R)$ 是 $X_R$ 中映到 $x'$ 的点，而
$t(0_R)$ 是 $X_R$ 中映到 $x$ 的点。注意 $t(0_R)$ 属于 $(X_R)^0$，并且
$t(\eta_R) \leadsto t(0_R)$。因此只需证明这蕴含
$t(\eta_R) \in (X_R)^0$。故只需在 $Y$ 是离散赋值环的谱、$y$ 是其闭点的情形证明本引理。

\medskip\noindent
假设 $Y$ 是离散赋值环的谱，而 $y$ 是其闭点。目标是证明 $X^0$ 是 $X^0_y$
的一个邻域。由于 $X_y$ 只有有限多个不可约分支，$X^0_y$ 在 $X_y$ 中既开又闭。
因此补集 $C = X_y \setminus X_y^0$ 在 $X$ 中是闭集。于是
$U = X \setminus C$ 是 $X^0_y$ 的开邻域，并且 $U^0 = X^0$。故只需对态射
$U \to Y$ 证明结论。换言之，可假设 $X_y$ 连通。假设 $X$ 不连通，并设
$X = X_1 \amalg \ldots \amalg X_n$ 为其连通分支分解。则 $s(Y)$ 完全包含在某个
$X_i$ 中。不妨设 $s(Y) \subset X_1$，于是 $X^0 \subset X_1$。故可将 $X$
换成 $X_1$，并假设 $X$ 连通。此时引理 \ref{more-morphisms-lemma-connected-flat-over-dvr}
表明 $X_\eta$ 连通，即 $X^0 = X$，结论成立。
\end{proof}

\begin{lemma}
\label{more-morphisms-lemma-connected-along-section-open}
设 $f : X \to Y$、$s : Y \to X$ 如情形
\ref{more-morphisms-situation-connected-along-section}。假设
\begin{enumerate}
\item $f$ 有限表现且平坦；并且
\item $f$ 的所有纤维均几何约化。
\end{enumerate}
则 $X^0$ 在 $X$ 中是开集。
\end{lemma}

\begin{proof}
这立即由引理 \ref{more-morphisms-lemma-connected-along-section-open-neighbourhood} 得出。
\end{proof}






\section{纤维的维数}
\label{more-morphisms-section-dimension}

\begin{lemma}
\label{more-morphisms-lemma-dimension-in-neighbourhood}
设 $f : X \to Y$ 是概形态射。假设 $Y$ 不可约、泛点为 $\eta$，且 $f$
是有限型的。若 $X_\eta$ 的维数为 $n$，则存在非空开子集 $V \subset Y$，使得对所有
$y \in V$，纤维 $X_y$ 的维数均为 $n$。
\end{lemma}

\begin{proof}
令 $Z = \{x \in X \mid \dim_x(X_{f(x)}) > n \}$。由 态射章引理
\ref{morphisms-lemma-openness-bounded-dimension-fibres}，这是 $X$ 的闭子集。按假设
$Z_\eta = \emptyset$。因此由引理 \ref{more-morphisms-lemma-empty-generic-fibre}，可缩小 $Y$
并假设 $Z = \emptyset$。令
$Z' = \{x \in X \mid \dim_x(X_{f(x)}) > n - 1 \} =
\{x \in X \mid \dim_x(X_{f(x)}) = n\}$。与前面一样，这是 $X$ 的闭子集。
按假设有 $Z'_\eta \not = \emptyset$。因此缩小
$Y$ 后，由引理 \ref{more-morphisms-lemma-nonempty-generic-fibre} 可假设 $Z' \to Y$ 是满射。
结论成立。
\end{proof}

\begin{lemma}
\label{more-morphisms-lemma-base-change-dimension-fibres}
设 $f : X \to Y$ 是有限型态射。令
$$
n_{X/Y} : Y \to \{0, 1, 2, 3, \ldots, \infty\}
$$
为把 $y \in Y$ 映到 $X_y$ 的维数的函数。若 $g : Y' \to Y$ 是态射，则
$$
n_{X'/Y'} = n_{X/Y} \circ g
$$
其中 $X' \to Y'$ 是 $f$ 的基变换。
\end{lemma}

\begin{proof}
这由 态射章引理
\ref{morphisms-lemma-dimension-fibre-after-base-change} 得出。
\end{proof}

\begin{lemma}
\label{more-morphisms-lemma-dimension-fibres-constructible}
设 $f : X \to Y$ 是概形态射。令 $n_{X/Y}$ 为引理
\ref{more-morphisms-lemma-base-change-dimension-fibres} 中引入的、在 $Y$ 上给出 $f$
之纤维维数的函数。假设 $f$ 是有限表现的。则 $n_{X/Y}$ 的水平集
$$
E_n = \{y \in Y \mid n_{X/Y}(y) = n\}
$$
在 $Y$ 中局部可构造。
\end{lemma}

\begin{proof}
固定 $n$。设 $y \in Y$。须证明在 $Y$ 中存在 $y$ 的开邻域 $V$，使得
$E_n \cap V$ 在 $V$ 中可构造。因此可假设 $Y$ 仿射。写成
$Y = \Spec(A)$，并将 $A = \colim A_i$ 写成有限型 $\mathbf{Z}$-代数的有向极限。
由 极限章引理 \ref{limits-lemma-descend-finite-presentation}，可以找到 $i$
及有限表现态射 $f_i : X_i \to \Spec(A_i)$，使其到 $Y$ 的基变换恢复 $f$。
由引理 \ref{more-morphisms-lemma-base-change-dimension-fibres}，只需对 $f_i$ 证明本引理。因此问题化到
$Y$ 是诺特环的谱的情形。

\medskip\noindent
在 $Y$ 为诺特概形时，我们用 拓扑章引理
\ref{topology-lemma-characterize-constructible-Noetherian} 的判别法证明 $E_n$ 可构造。
为此，设 $Z \subset Y$ 是不可约闭子概形。须证明 $E_n \cap Z$ 要么包含非空开子集，
要么在 $Z$ 中不稠密。令 $\xi \in Z$ 为泛点。引理
\ref{more-morphisms-lemma-dimension-in-neighbourhood} 表明，$n_{X/Y}$ 在 $Z$ 中 $\xi$
的一个邻域上为常数。这蕴含所需结论。
\end{proof}

\begin{lemma}
\label{more-morphisms-lemma-dimension-fibres-flat}
设 $f : X \to Y$ 是有限表现的平坦概形态射。令 $n_{X/Y}$ 为引理
\ref{more-morphisms-lemma-base-change-dimension-fibres} 中引入的、在 $Y$ 上给出 $f$
之纤维维数的函数。则 $n_{X/Y}$ 下半连续。
\end{lemma}

\begin{proof}
令 $W \subset X$、$W = \coprod_{d \geq 0} U_d$ 为引理
\ref{more-morphisms-lemma-flat-finite-presentation-CM-open} 和
\ref{more-morphisms-lemma-flat-finite-presentation-CM-pieces} 中构造的开集。设 $y \in Y$ 为一点。若
$n_{X/Y}(y) = \dim(X_y) = n$，则 $y$ 属于 $U_n \to Y$ 的像。由 Morphisms,
引理 \ref{morphisms-lemma-fppf-open}，$f(U_n)$ 在 $Y$ 中是开集。因此存在
$y$ 的开邻域，使得其上 $n_{X/Y}$ 均为 $\geq n$。
\end{proof}

\begin{lemma}
\label{more-morphisms-lemma-dimension-fibres-proper}
设 $f : X \to Y$ 是概形的固有态射。令 $n_{X/Y}$ 为引理
\ref{more-morphisms-lemma-base-change-dimension-fibres} 中引入的、在 $Y$ 上给出 $f$
之纤维维数的函数。则 $n_{X/Y}$ 上半连续。
\end{lemma}

\begin{proof}
令 $Z_d = \{x \in X \mid \dim_x(X_{f(x)}) > d\}$。由 态射章引理
\ref{morphisms-lemma-openness-bounded-dimension-fibres}，$Z_d$ 是 $X$ 的闭子集。
由于 $f$ 固有，$f(Z_d)$ 是闭集。又因
$y \in f(Z_d) \Leftrightarrow n_{X/Y}(y) > d$，本引理成立。
\end{proof}

\begin{lemma}
\label{more-morphisms-lemma-dimension-fibres-proper-flat}
设 $f : X \to Y$ 是有限表现的固有平坦概形态射。令 $n_{X/Y}$ 为引理
\ref{more-morphisms-lemma-base-change-dimension-fibres} 中引入的、在 $Y$ 上给出 $f$
之纤维维数的函数。则 $n_{X/Y}$ 局部常值。
\end{lemma}

\begin{proof}
这立即由引理 \ref{more-morphisms-lemma-dimension-fibres-flat} 和
\ref{more-morphisms-lemma-dimension-fibres-proper} 得出。
\end{proof}

















\section{弱相对 Noether 正规化}
\label{more-morphisms-section-relative-noether-normalization}

\noindent
本节目标是证明引理 \ref{more-morphisms-lemma-weak-relative-noether-normalization}。

\begin{lemma}
\label{more-morphisms-lemma-silly}
设 $R$ 是环。设 $\mathfrak p_1, \ldots, \mathfrak p_r$ 是 $R$ 的素理想，且当
$i \not = j$ 时有 $\mathfrak p_i \not \subset \mathfrak p_j$。设
$k_i \subset \kappa(\mathfrak p_i)$ 是子域，使得扩张
$\kappa(\mathfrak p_i)/k_i$ 不是代数扩张。设 $J \subset R$ 是不包含于任何
$\mathfrak p_i$ 的理想。则存在元素 $x \in J$，使得对 $i = 1, \ldots, r$，
$x$ 在 $\kappa(\mathfrak p_i)$ 中的像在 $k_i$ 上超越。
\end{lemma}

\begin{proof}
理想 $J_i = J \mathfrak p_1 \ldots \hat{\mathfrak p}_i \ldots \mathfrak p_r$
不包含于 $\mathfrak p_i$，见 代数章引理
\ref{algebra-lemma-product-ideals-in-prime}。于是
% P05-ERR-CH37-0050: Frozen source uses the undefined ring B in Frac(B/p_i); the ambient ring is R.
$\kappa(\mathfrak p_i) = \text{Frac}(B/\mathfrak p_i)$ 的每个元素 $\xi$
都可写成 $\xi = a/b$，其中 $a, b \in J_i$ 且
$b \not \in \mathfrak p_i$。选取在 $k_i$ 上超越的 $\xi$，可见 $a$ 或 $b$
在 $\kappa(\mathfrak p_i)$ 中的像在 $k_i$ 上超越。因此，对每个
$i = 1, \ldots, r$，可以找到元素
$x_i \in J_i = J \mathfrak p_1 \ldots \hat{\mathfrak p}_i \ldots \mathfrak p_r$，
使其在 $\kappa(\mathfrak p_i)$ 中的像在 $k_i$ 上超越。于是
$x = x_1 + \ldots + x_r$ 满足要求。
\end{proof}

\begin{lemma}
\label{more-morphisms-lemma-sum-is-ok}
设 $R \to S$ 是有限型环同态。设 $d \geq 0$，并设 $a, b \in S$。假设由把
$x$ 映到 $a$ 的 $R$-代数同态 $R[x] \to S$ 给出的态射
$$
f_a : \Spec(S) \longrightarrow \mathbf{A}^1_R
$$
的纤维维数为 $\leq d$。则存在 $n_0$，使得对 $n \geq n_0$，由把 $x$
映到 $a^n + b$ 的 $R$-代数同态 $R[x] \to S$ 给出的态射
$$
f_{a^n + b} : \Spec(S) \longrightarrow \mathbf{A}^1_R
$$
的纤维维数为 $\leq d$。
\end{lemma}

\begin{proof}
本段将问题化到 $R \to S$ 有限表现的情形。具体地，将
$S = R[A, B, x_1, \ldots, x_n]/J$ 写成这样的形式，其中
% P05-ERR-CH37-0051: Frozen source places J in R[x_1,...,x_n] although the presentation also uses indeterminates A and B.
$J \subset R[x_1, \ldots, x_n]$ 是某个理想，而 $A$ 和 $B$ 在 $S$ 中分别映到
$a$ 和 $b$。于是 $J$ 是其有限生成理想 $J_\lambda \subset J$ 的并。置
$S_\lambda = R[A, B, x_1, \ldots, x_n]/J_\lambda$，并以
$a_\lambda, b_\lambda \in S_\lambda$ 表示 $A$ 和 $B$ 的像。于是对某个 $\lambda$，
态射
$$
f_{a_\lambda} : \Spec(S_\lambda) \longrightarrow \mathbf{A}^1_R
$$
的纤维维数为 $\leq d$，见 极限章引理
\ref{limits-lemma-limit-dimension}。固定这样的 $\lambda$。若能找到适用于
$R \to S_\lambda$、$a_\lambda$、$b_\lambda$ 的 $n_0$，则该 $n_0$ 也适用于
$R \to S$。事实上，$f_{a_\lambda^n + b_\lambda} : \Spec(S_\lambda) \to \mathbf{A}^1_R$
的纤维包含 $f_{a^n + b} : \Spec(S) \to \mathbf{A}^1_R$ 的纤维。这将问题化到下一段所讨论的情形。

\medskip\noindent
假设 $R \to S$ 有限表现。本段将问题化到 $R$ 在 $\mathbf{Z}$ 上有限型的情形。
由 代数章引理 \ref{algebra-lemma-limit-module-finite-presentation}，可以找到有向集
$\Lambda$ 及定义在 $\Lambda$ 上的环同态系统 $R_\lambda \to S_\lambda$，其余极限为 $R \to S$，
使得当 $\mu \geq \lambda$ 时 $S_\mu = S_\lambda \otimes_{R_\lambda} R_\mu$，
并且每个 $R_\lambda$ 和 $S_\lambda$ 都在 $\mathbf{Z}$ 上有限型。选取
$\lambda_0 \in \Lambda$ 以及映到 $a, b \in S$ 的元素
$a_{\lambda_0}, b_{\lambda_0} \in S_{\lambda_0}$。对
$\lambda \geq \lambda_0$，以 $a_\lambda, b_\lambda \in S_\lambda$ 表示
$a_{\lambda_0}, b_{\lambda_0}$ 的像。于是当 $\lambda \geq \lambda_0$ 足够大时，态射
$$
f_{a_\lambda} : 
\Spec(S_\lambda) \longrightarrow \mathbf{A}^1_{R_\lambda}
$$
的纤维维数为 $\leq d$，见 极限章引理
\ref{limits-lemma-descend-dimension-d}。固定这样的 $\lambda$。若能找到适用于
$R_\lambda \to S_\lambda$、$a_\lambda$、$b_\lambda$ 的 $n_0$，则该 $n_0$
也适用于 $R \to S$。事实上，由 态射章引理
\ref{morphisms-lemma-dimension-fibre-after-base-change}，
$f_{a^n + b} : \Spec(S) \to \mathbf{A}^1_R$ 的任一纤维都与
$f_{a_\lambda^n + b_\lambda} : \Spec(S_\lambda) \to \mathbf{A}^1_{R_\lambda}$
的某个纤维具有相同维数。
% P05-ERR-CH37-0052: Frozen source has the duplicated phrase ``the the'' in this reduction sentence.
这将问题化到下一段所讨论的情形。

\medskip\noindent
假设 $R$ 和 $S$ 在 $\mathbf{Z}$ 上有限型。特别地，$R$ 的维数有限，因而可对
$\dim(R)$ 使用归纳法。因此可假设结论对引理中所有如下的 $R' \to S'$、$a$、$b$
成立：$R'$ 和 $S'$ 在 $\mathbf{Z}$ 上有限型，但
$\dim(R') < \dim(R)$。

\medskip\noindent
由于此陈述只涉及 $S$ 的谱的拓扑，可假设 $S$ 约化。令 $S^\nu$ 为 $S$
的正规化。由于 $S$ 优秀，$S \subset S^\nu$ 是有限扩张，见 Algebra,
命题 \ref{algebra-proposition-ubiquity-nagata} 和 态射章引理
\ref{morphisms-lemma-nagata-normalization}。因此 $\Spec(S^\nu) \to \Spec(S)$
是满的有限态射（代数章引理
\ref{algebra-lemma-integral-overring-surjective}）。由此可知，若结论对
$R \to S^\nu$ 以及 $a$、$b$ 在 $S^\nu$ 中的像成立，则它对
$R \to S$、$a$、$b$ 也成立。（略去一个小细节。）这将问题化到下一段所讨论的情形。

\medskip\noindent
假设 $R$ 和 $S$ 在 $\mathbf{Z}$ 上有限型，并且 $S$ 正规。于是对某些正规整环
$S_i$，有 $S = S_1 \times \ldots \times S_r$。若结论对每个 $R \to S_i$
及 $a$、$b$ 在 $S_i$ 中的像成立，则它对 $R \to S$、$a$、$b$
也成立。（略去一个小细节。）这将问题化到下一段所讨论的情形。

\medskip\noindent
假设 $R$ 和 $S$ 在 $\mathbf{Z}$ 上有限型，并且 $S$ 是正规整环。可将 $R$
换成 $R$ 在 $S$ 中的像（这不会增加 $R$ 的维数）。这将问题化到下一段所讨论的情形。

\medskip\noindent
假设 $R \subset S$ 在 $\mathbf{Z}$ 上有限型，且 $S$ 是正规整环。考虑态射
$$
f_a : \Spec(S) \to \mathbf{A}^1_R
$$
假设表明 $f_a$ 的纤维维数为 $\leq d$。因此
$f : \Spec(S) \to \Spec(R)$ 的纤维维数为 $\leq d + 1$（Morphisms,
引理 \ref{morphisms-lemma-dimension-fibre-at-a-point-additive}）。考虑整概形间的态射
$$
\phi : \Spec(S) \to \mathbf{A}^2_R = \Spec(R[x, y])
$$
，它对应于把 $x$ 映到 $a$、把 $y$ 映到 $b$ 的 $R$-代数同态
$R[x, y] \to S$。须考虑两种情形：
\begin{enumerate}
\item $\phi$ 是支配态射；
\item $\phi$ 不是支配态射。
\end{enumerate}
我们断言，在两种情形下，都存在整数 $n_0$ 和非空开子集
$V \subset \Spec(R)$，使得当 $n \geq n_0$ 时，$f_{a^n + b}$ 在各点
$q \in \mathbf{A}^1_V$ 处的纤维维数为 $\leq d$。

\medskip\noindent
在情形 (1) 中证明该断言。有 $f_{a^n + b} = \pi_n \circ \phi$，其中
$$
\pi_n :  \mathbf{A}^2_R \to \mathbf{A}^1_R
$$
是对应于把 $x$ 映到 $x^n + y$ 的 $R$-代数同态 $R[x] \to R[x, y]$
的平坦态射。由于 $\phi$ 支配，存在稠密开子集 $U \subset \Spec(S)$，使得
$\phi|_U : U \to \mathbf{A}^2_R$ 平坦（例如由泛平坦性得出，见 Morphisms,
命题 \ref{morphisms-proposition-generic-flatness}）。于是复合
$$
f_{a^n + b}|_U :
U \xrightarrow{\phi|_U}
\mathbf{A}^2_R
\xrightarrow{\pi_n}
\mathbf{A}^1_R
$$
也平坦。因此由 态射章引理
\ref{morphisms-lemma-dimension-fibre-at-a-point-additive}，此态射的纤维在
$f|_U : U \to \Spec(R)$ 的纤维中至少有余维 $1$。换言之，
$f_{a^n + b}|_U$ 的纤维维数为 $\leq d$。另一方面，由于 $U$ 在
$\Spec(S)$ 中稠密，可以找到非空开子集 $V \subset \Spec(R)$，使得对所有
$p \in V$，$U \cap f^{-1}(p) \subset f^{-1}(p)$ 都稠密（例如见引理
\ref{more-morphisms-lemma-nowhere-dense-generic-fibre}）。因此
$\dim(f^{-1}(p) \setminus U \cap f^{-1}(p)) \leq d$，从而该断言成立（因为
$f_{a^n + b} : \Spec(S) \to \mathbf{A}^1_R$ 的任一纤维都包含于 $f$ 的一个纤维中）。

\medskip\noindent
情形 (2)。此时可在 $R[x, y]$ 中找到非零元 $g = \sum c_{ij} x^i y^j$，使得
$\Im(\phi) \subset V(g)$。事实上，可假设 $g$ 在 $\text{Frac}(R)$ 上不可约。
若 $g \in R[x]$，设其首项系数为 $c$，则在
$V = D(c) \subset \Spec(R)$ 上，$f$ 的纤维维数已经为 $\leq d$（因为 $f_a$
的像包含于 $V(g) \subset \mathbf{A}^1_R$，而后者在 $V$ 上的纤维有限）。因此可假设
$g$ 不属于 $R[x]$。令 $s \geq 1$ 为 $g$ 作为 $y$ 的多项式的次数，并令 $t$
为 $\sum c_{is} x^i$ 作为 $x$ 的多项式的次数。于是 $c_{ts}$ 非零，并且
$$
g(x, -x^n) = (-1)^s c_{ts} x^{t + sn} + l.o.t.
$$
，只要 $n$ 大于 $g$ 作为 $x$ 的多项式的次数（略去一个小细节）。对这样的
$n$，多项式 $g(x, -x^n)$ 是 $x$ 的非零多项式，而且对所有
% P05-ERR-CH37-0053: Frozen source switches c_{ts} to c_{st} in the nonvanishing condition.
$\mathfrak p \subset R$，只要 $c_{st} \not \in \mathfrak p$，它在
$\kappa(\mathfrak p)[x]$ 中的像也是非零多项式。因此当 $V$ 取为
$\Spec(R)$ 的主开集 $D(c_{ts})$ 时，我们的断言成立。

\medskip\noindent
现在可以对 $\dim(R)$ 使用归纳法。具体地，令 $I \subset R$ 为满足
$V(I) = \Spec(R) \setminus V$ 的理想。由于 $R$ 是整环，注意有
$\dim(R/I) < \dim(R)$。令 $n'_0$ 为对 $\dim(R)$ 归纳应用于
$R/I \to S/IS$ 以及 $a$、$b$ 在 $S/IS$ 中的像而得到的整数。于是
$\max(n_0, n'_0)$ 满足要求。
\end{proof}

\begin{lemma}
\label{more-morphisms-lemma-weak-relative-noether-normalization}
设 $R \to S$ 是有限型环同态。令 $d$ 为 $\Spec(S) \to \Spec(R)$
的各纤维维数的最大值。则存在拟有限环同态 $R[t_1, \ldots, t_d] \to S$。
\end{lemma}

\begin{proof}
本段将问题化到 $R \to S$ 有限表现的情形。具体地，将
$S = R[x_1, \ldots, x_n]/J$ 写成这样的形式，其中
$J \subset R[x_1, \ldots, x_n]$ 是某个理想。于是 $J$ 是其有限生成理想
$J_\lambda \subset J$ 的并。置 $S_\lambda = R[x_1, \ldots, x_n]/J_\lambda$。
于是对某个 $\lambda$，$\Spec(S_\lambda) \to \Spec(R)$ 的纤维维数为
$\leq d$，见 极限章引理 \ref{limits-lemma-limit-dimension}。固定这样的
$\lambda$。若能找到拟有限同态 $R[t_1, \ldots, t_d] \to S_\lambda$，则复合
$R[t_1, \ldots, t_d] \to S$ 当然拟有限。这将问题化到下一段所讨论的情形。

\medskip\noindent
假设 $R \to S$ 有限表现。本段将问题化到 $R$ 在 $\mathbf{Z}$ 上有限型的情形。
由 代数章引理 \ref{algebra-lemma-limit-module-finite-presentation}，可以找到有向集
$\Lambda$ 及定义在 $\Lambda$ 上的环同态系统 $R_\lambda \to S_\lambda$，其余极限为 $R \to S$，
使得当 $\mu \geq \lambda$ 时 $S_\mu = S_\lambda \otimes_{R_\lambda} R_\mu$，
并且每个 $R_\lambda$ 和 $S_\lambda$ 都在 $\mathbf{Z}$ 上有限型。于是当
$\lambda$ 足够大时，$\Spec(S_\lambda) \to \Spec(R_\lambda)$ 的纤维维数为
$\leq d$，见 极限章引理 \ref{limits-lemma-descend-dimension-d}。固定这样的
$\lambda$。若能找到拟有限环同态
$R_\lambda[t_1, \ldots, t_d] \to S_\lambda$，则其基变换
$R[t_1, \ldots, t_d] \to S$ 也拟有限（代数章引理
\ref{algebra-lemma-quasi-finite-base-change}）。
% P05-ERR-CH37-0054: Frozen source again has the duplicated phrase ``the the'' in this reduction sentence.
这将问题化到下一段所讨论的情形。

\medskip\noindent
假设 $R$ 和 $S$ 在 $\mathbf{Z}$ 上有限型。若 $d = 0$，则环同态拟有限，证明完成。
假设 $d > 0$。我们将找到元素 $a \in S$，使得 $R$-代数同态
$R[x] \to S$、$x \mapsto a$ 的纤维维数为 $< d$。这将通过归纳完成证明。

\medskip\noindent
我们对 $\dim(R)$ 归纳证明 $a$ 的存在性。

\medskip\noindent
在 $S$ 的极小素理想中，令 $\mathfrak q_1, \ldots, \mathfrak q_r \subset S$
为满足 $\dim_{\mathfrak q_i}(S/R) = d$ 的那些。记号见 代数章定义
\ref{algebra-definition-relative-dimension}。设 $\mathfrak q_i$ 位于素理想
$\mathfrak p_i \subset R$ 上方。由于 $\mathfrak q_i$ 是其纤维的泛点，有
$\text{trdeg}_{\kappa(\mathfrak p_i)}(\kappa(\mathfrak q_i)) = d$；例如可对
$S \otimes_R \kappa(\mathfrak p_i)$ 应用 代数章引理
\ref{algebra-lemma-dimension-at-a-point-finite-type-field}。因此由引理
\ref{more-morphisms-lemma-silly}，可以找到元素 $a \in S$，使得对 $i = 1, \ldots, r$，
$a$ 在 $\kappa(\mathfrak q_i)$ 中的像在 $\kappa(\mathfrak p_i)$ 上超越。考虑态射
$$
f_a : \Spec(S) \longrightarrow \mathbf{A}^1_R
$$
% P05-ERR-CH37-0055: Frozen source has the malformed phrase ``corresponding the ... to mapping''.
，它对应于把 $x$ 映到 $a$ 的 $R$-代数同态 $R[x] \to S$。令
$U \subset \Spec(S)$ 为纤维维数为 $\leq d - 1$ 的开子集，见 Morphisms,
引理 \ref{morphisms-lemma-openness-bounded-dimension-fibres}。按构造，$U$ 包含
$\Spec(S)$ 的所有泛点。特别地，$U$ 包含 $\Spec(S) \to \Spec(R)$
的所有泛纤维的全部泛点，因为这样的点必为 $\Spec(S)$ 的泛点。置
$T = \Spec(S) \setminus U$，并将其视为 $\Spec(S)$ 的约化闭子概形。由刚才所述及假设
$\dim(S/R) \leq d$，$T \to \Spec(R)$ 的泛纤维维数为 $\leq d - 1$。因此多次应用引理
\ref{more-morphisms-lemma-dimension-in-neighbourhood}，以构造 $\Spec(R)$ 各泛点的开邻域，可以找到稠密开子集
$V \subset \Spec(R)$，使得 $T_V \to V$ 的纤维维数为 $\leq d - 1$。由此可知，对
$q \in \mathbf{A}^1_V$，$f_a$ 在 $q$ 上的纤维维数为 $\leq d - 1$（因为已经分别界定了
$f_a|_U$ 与 $f_a|_T$ 的纤维维数）。

\medskip\noindent
由素理想回避，可假设对某个 $f \in R$ 有 $V = D(f)$。于是环同态
$R_f[x] \to S_f$、$x \mapsto a$ 的纤维维数为 $\leq d - 1$。可将 $a$
换成 $fa$，并假设 $a \in (f)$。对 $\dim(R)$ 使用归纳法，可以找到元素
$\overline{b} \in S/fS$，使得
% P05-ERR-CH37-0056: Frozen source writes Spec(R/fR[x]), an ambiguous/ill-typed polynomial-ring quotient; likely Spec((R/fR)[x]).
$\Spec(S/fS) \to \Spec(R/fR[x])$、$x \mapsto \overline{b}$ 的纤维维数为
$\leq d - 1$。令 $b \in S$ 为 $\overline{b}$ 的一个提升。由引理
\ref{more-morphisms-lemma-sum-is-ok}，存在 $n > 0$，使得 $a^n + b$ 对 $R_f \to S_f$
仍然满足要求。另一方面，$a^n + b$ 在 $S/fS$ 中的像是 $\overline{b}$，证明完成。
\end{proof}













\section{Bertini 定理}
\label{more-morphisms-section-bertini}

\noindent
我们继续 簇章第 \ref{varieties-section-bertini} 节 中开始的讨论。本节遵循
\cite{Jou} 给出的精彩论证，证明几何不可约簇的一般超平面截面几何不可约。

\begin{lemma}
\label{more-morphisms-lemma-amazing-bertini-lemma}
\begin{reference}
见 \cite{Jou} 第 71、72 页
\end{reference}
设 $K/k$ 是几何不可约的有限生成域扩张。设 $n \geq 1$。设
$g_1, \ldots, g_n \in K$ 是元素，并且存在 $c_1, \ldots, c_n \in k$，使得元素
$$
x_1, \ldots, x_n, \sum g_ix_i, \sum c_ig_i \in K(x_1, \ldots, x_n)
$$
在 $k$ 上代数无关。则 $K(x_1, \ldots, x_n)$ 在
$k(x_1, \ldots, x_n, \sum g_ix_i)$ 上几何不可约。
\end{lemma}

\begin{proof}
取引理陈述中的 $c_1, \ldots, c_n \in k$。记
$\xi = \sum g_ix_i$、$\delta = \sum c_ig_i$。对 $a \in k$，考虑
$K(x_1, \ldots, x_n)$ 的自同构 $\sigma_a$；它在 $K$ 上为恒等，并由规则
$$
\sigma_a(x_i) = x_i + a c_i
$$
给出。注意 $\sigma_a(\xi) = \xi + a \delta$ 且
$\sigma_a(\delta) = \delta$。考虑域塔
$$
K_0 = k(x_1, \ldots, x_n) \subset
K_1 = K_0(\xi) \subset
K_2 = K_0(\xi, \delta) \subset K(x_1, \ldots, x_n) = \Omega
$$
注意 $\sigma_a(K_0) = K_0$ 且 $\sigma_a(K_2) = K_2$。设
$\theta \in \Omega$ 在 $K_1$ 上可分代数。须证明 $\theta \in K_1$，见
代数章引理
\ref{algebra-lemma-geometrically-irreducible-separable-elements}。

\medskip\noindent
以 $K'_2$ 表示 $K_2$ 在 $\Omega$ 中的可分代数闭包。由于 $K'_2/K_2$
有限（代数章引理 \ref{algebra-lemma-make-geometrically-irreducible}）且可分，
$K'_2$ 与 $K_2$ 之间只有有限多个中间域（域章引理
\ref{fields-lemma-primitive-element}）。若 $k$ 无限\footnote{有限域情形将在证明的最后一段处理。}，
则可在 $k$ 中找到不同的元素 $a_1, a_2$，使得
$$
K_2(\sigma_{a_1}(\theta)) = K_2(\sigma_{a_2}(\theta))
$$
作为 $\Omega$ 的子域相等。记 $\theta_i = \sigma_{a_i}(\theta)$ 以及
$\xi_i = \sigma_{a_i}(\xi) = \xi + a_i \delta$。注意
$$
K_2 = K_0(\xi_1, \xi_2)
$$
，因为 $\xi_i = \xi + a_i \delta$、
$\xi = (a_2 \xi_1 - a_1 \xi_2)/(a_2 - a_1)$，且
$\delta = (\xi_1 - \xi_2)/(a_1 - a_2)$。由于 $K_2/K_0$ 是超越次数为
$2$ 的纯超越扩张，$\xi_1$ 和 $\xi_2$ 在 $K_0$ 上代数无关。由于
$\theta_1$ 在 $K_0(\xi_1)$ 上代数，$\xi_2$ 在
$K_0(\xi_1, \theta_1)$ 上超越。

\medskip\noindent
按假设，$K/k$ 几何不可约。这蕴含 $K(x_1, \ldots, x_n)/K_0$ 几何不可约
（代数章引理
\ref{algebra-lemma-geometrically-irreducible-base-change-transcendental}）。进而，作为子扩张，
$K_0(\xi_1, \theta_1)/K_0$ 几何不可约（代数章引理
\ref{algebra-lemma-subalgebra-geometrically-irreducible}）。由于 $\xi_2$ 在
$K_0(\xi_1, \theta_1)$ 上超越，$K_0(\xi_1, \xi_2, \theta_1)/K_0(\xi_2)$
几何不可约（代数章引理
\ref{algebra-lemma-geometrically-irreducible-add-transcendental}）。按上面对
$a_1, a_2$ 的选取，有
$$
K_0(\xi_1, \xi_2, \theta_1) =
K_2(\sigma_{a_1}(\theta)) =
K_2(\sigma_{a_2}(\theta)) =
K_0(\xi_1, \xi_2, \theta_2)
$$
由于 $\theta_2$ 在 $K_0(\xi_2)$ 上可分代数，再次由 代数章引理
\ref{algebra-lemma-geometrically-irreducible-separable-elements} 得出
$\theta_2 \in K_0(\xi_2)$。对这一关系施以 $\sigma_{a_2}^{-1}$，便得到
% P05-ERR-CH37-0057: Frozen source has the typo ``givens''.
$\theta \in K_0(\xi) = K_1$，正是所需。

\medskip\noindent
这完成了 $k$ 无限时的证明。若 $k$ 有限，则可选取变量 $t$，并考虑扩张
$K(t)/k(t)$；由 代数章引理
\ref{algebra-lemma-geometrically-irreducible-base-change-transcendental}，它几何不可约。
% P05-ERR-CH37-0058: Frozen source has the ungrammatical phrase ``it is still be true''.
仍有
$x_1, \ldots, x_n, \sum g_ix_i, \sum c_ig_i$
在 $K(t, x_1, \ldots, x_n)$ 中于 $k(t)$ 上代数无关，故由上述论证，
$K(t, x_1, \ldots, x_n)$ 在
$k(t, x_1, \ldots, x_n, \sum g_ix_i)$ 上几何不可约。最后再用一次 Algebra,
引理
\ref{algebra-lemma-geometrically-irreducible-base-change-transcendental}
即可完成证明。
\end{proof}

\begin{lemma}
\label{more-morphisms-lemma-algebraically-independent}
设 $A$ 是域 $k$ 上的有限型整环。设 $n \geq 2$。设
$g_1, \ldots, g_n \in A$ 是使得 $V(g_1, g_2)$ 有一个维数为
$\dim(A) - 2$ 的不可约分支的元素。则存在 $c_1, \ldots, c_n \in k$，使得元素
$$
x_1, \ldots, x_n, \sum g_ix_i, \sum c_ig_i \in
\text{Frac}(A)(x_1, \ldots, x_n)
$$
在 $k$ 上代数无关。
\end{lemma}

\begin{proof}
在 $k$ 上代数无关意味着由 $y = \sum g_ix_i$ 和 $z = \sum c_ig_i$ 给出的态射
$$
T = \Spec(A[x_1, \ldots, x_n])
\longrightarrow
\Spec(k[x_1, \ldots, x_n, y, z]) = S
$$
为支配态射。置 $d = \dim(A)$。若 $T \to S$ 不支配，则其像的维数
$< n + 2$，因而每个纤维的每个不可约分支的维数
$> d + n - (n + 2) = d - 2$，见 簇章引理
\ref{varieties-lemma-dimension-fibres-locally-algebraic}。选取闭点
$u \in V(g_1, g_2)$，使其包含在一个维数为 $d - 2$ 的不可约分支中，而不属于
$V(g_1, g_2)$ 的任何其他分支。考虑 $T$ 中位于 $u$ 上方的闭点
$t = (u, 1, 0, \ldots 0)$。置
$(c_1, \ldots, c_n) = (0, 1, 0, \ldots, 0)$。于是 $t$ 映到 $S$ 中的点
$s = (1, 0, \ldots, 0)$。$T \to S$ 在 $s$ 上的纤维由下列元素截出：
$$
x_1 - 1, x_2, \ldots, x_n, \sum x_ig_i, g_2
$$
，因而等价地由下列元素截出：
$$
x_1 - 1, x_2, \ldots, x_n, g_1, g_2
$$
。由对 $g_1, g_2$ 的条件，此子概形有一个维数为 $d - 2$ 的不可约分支。
\end{proof}

\begin{lemma}
\label{more-morphisms-lemma-bertini-irreducible}
\begin{reference}
\cite[Theorem 6.3 part 4)]{Jou}
\end{reference}
在 Varieties, Situation \ref{varieties-situation-family-divisors} 中，假设
\begin{enumerate}
\item $X$ 在 $k$ 上有限型；
\item $X$ 在 $k$ 上几何不可约；
\item 存在 $v_1, v_2, v_3 \in V$ 以及 $H_{v_2} \cap H_{v_3}$ 的一个不可约分支
$Z$，使得 $Z \not \subset H_{v_1}$ 且 $\text{codim}(Z, X) = 2$；并且
\item $\bigcap_{v \in V} H_v$ 的每个不可约分支 $Y$ 都满足
$\text{codim}(Y, X) \geq 2$。
\end{enumerate}
则对一般的 $v \in V \otimes_k k'$，概形 $H_v$ 在 $k'$ 上几何不可约。
\end{lemma}

\begin{proof}
要使假设 (3) 成立，元素 $v_1, v_2, v_3$ 必须在 $V$ 中于 $k$ 上线性无关
（略去一个小细节）。因此可选取 $V$ 的基 $v_1, \ldots, v_r$，以前 $3$
个基向量包含这些元素。回忆 $H_{univ} \subset \mathbf{A}^r_k \times_k X$
是“泛除子”。考虑投影 $q : H_{univ} \to \mathbf{A}^r_k$，其概形论纤维为除子
$H_v$。由引理 \ref{more-morphisms-lemma-geom-irreducible-generic-fibre}，只需证明 $q$
的泛纤维几何不可约。为此，可将 $X$ 换成其约化，因而可假设 $X$
是在 $k$ 上有限型的整概形。

\medskip\noindent
设 $U \subset X$ 是满足 $\mathcal{L}|_U \cong \mathcal{O}_U$ 的非空仿射开集。
写成 $U = \Spec(A)$。以 $f_i \in A$ 表示在同构
$\mathcal{L}|_U \cong \mathcal{O}_U$ 下对应于截面 $\psi(v_i)|_U$ 的元素。于是
$H_{univ} \cap (\mathbf{A}^r_k \times_k U)$ 由下式给出：
$$
H_U = \Spec(A[x_1, \ldots, x_r]/(x_1f_1 + \ldots + x_rf_r))
$$
由所选基可知 $f_1$ 不可能为零，因为否则 $v_1 = 0$，进而
$H_{v_1} = X$，与假设 (3) 矛盾。因此 $\sum x_if_i$ 是
$A[x_1, \ldots, x_r]$ 中的非零因子。由此，$H_U$ 的每个不可约分支的维数均为
$d + r - 1$，其中 $d = \dim(X) = \dim(A)$。若 $U' = U \cap D(f_1)$，则
$$
H_{U'} =
\Spec(A_{f_1}[x_1, \ldots, x_r]/(x_1f_1 + \ldots + x_rf_r)) \cong
% P05-ERR-CH37-0059: Frozen source omits the comma after \ldots in the polynomial-variable list below.
\Spec(A_{f_1}[x_2, \ldots x_r]) =
\mathbf{A}^{r - 1}_k \times_k U'
$$
不可约。另一方面，
$$
H_U \setminus H_{U'} =
\Spec(A/(f_1)[x_1, \ldots, x_r]/(x_2f_2 + \ldots + x_rf_r))
$$
的维数至多为 $d + r - 2$。事实上，对 $i \not = 1$，概形
$(H_U \setminus H_{U'}) \times_U D(f_i)$ 要么为空（若 $f_i = 0$），要么由与上面相同的论证，
同构于 $\Spec(A/(f_1))$ 的某个开集上的 $r - 1$ 维仿射空间，因而维数至多为
$d + r - 2$。另一方面，$(H_U \setminus H_{U'}) \times_U V(f_2, \ldots, f_r)$
是 $\Spec(A/(f_1, \ldots, f_r))$ 上的 $r$ 维仿射空间，故假设 (4) 表明其维数至多为
$d + r - 2$。于是对上述每个 $U$，$H_U$ 都不可约。因此 $H_{univ}$ 不可约。

\medskip\noindent
因此只需证明 $H_{univ}$ 的泛点在 $\mathbf{A}^r_k$ 的泛点上几何不可约，见
簇章引理
\ref{varieties-lemma-geometrically-irreducible-function-field}。在 $X$ 中选取包含于
$X \setminus H_{v_1}$ 的非空仿射开集 $U = \Spec(A)$，使其与假设 (3) 所断言存在的
$H_{v_2} \cap H_{v_3}$ 的不可约分支 $Z$ 相交。沿用上述记号，须证明域扩张
$$
\text{Frac}(A[x_1, \ldots, x_r]/(x_1f_1 + \ldots + x_rf_r))/
k(x_1, \ldots , x_r)
$$
几何不可约。由对 $U$ 的选取，注意 $f_1$ 在 $A$ 中可逆。置
$K = \text{Frac}(A)$ 为 $A$ 的分式域。像上面一样消去变量 $x_1$ 后，只须证明域扩张
$$
K(x_2, \ldots, x_r)/
k(x_2, \ldots, x_r, -\sum\nolimits_{i = 2, \ldots, r} f_1^{-1}f_i x_i)
$$
几何不可约。由引理 \ref{more-morphisms-lemma-amazing-bertini-lemma}，只需证明对某些
$c_2, \ldots, c_r \in k$，元素
$$
x_2, \ldots, x_r, \sum\nolimits_{i = 2, \ldots, r} f_1^{-1}f_i x_i,
\sum\nolimits_{i = 2, \ldots, r} c_if_1^{-1}f_i
$$
在 $A[x_2, \ldots, x_r]$ 的分式域中于 $k$ 上代数无关。这由引理
\ref{more-morphisms-lemma-algebraically-independent} 以及如下事实得出：$Z \cap U$ 是
$V(f_1^{-1}f_2, f_1^{-1}f_3) \subset U$ 的一个不可约分支。
\end{proof}

\begin{remark}
\label{more-morphisms-remark-geometric-proof-bertini-irreducible}
我们概述引理 \ref{more-morphisms-lemma-bertini-irreducible} 一个特殊情形的“几何”证明。设 $k$
是代数闭域，而 $X \subset \mathbf{P}^n_k$ 光滑、不可约且维数 $\geq 2$。我们断言存在超平面
$H \subset \mathbf{P}^n_k$，使得 $X \cap H$ 光滑且不可约。事实上，由 Varieties,
引理 \ref{varieties-lemma-bertini}，对一般的
$v \in V = kT_0 \oplus \ldots \oplus kT_n$，相应的超平面截面 $X \cap H_v$
是光滑的。另一方面，由 Enriques--Severi--Zariski，概形 $X \cap H_v$ 连通，见
簇章引理 \ref{varieties-lemma-connectedness-ample-divisor}。因此
$X \cap H_v$ 光滑且不可约。
\end{remark}













\section{立方体定理}
\label{more-morphisms-section-theorem-cube}

\noindent
下述引理说明，在引理所作假设下，Picard 函子的对角可由局部闭浸入表示。

\begin{lemma}
\label{more-morphisms-lemma-diagonal-picard-flat-proper}
设 $f : X \to S$ 是有限表现的平坦固有态射。设 $\mathcal{E}$ 是有限局部自由
$\mathcal{O}_X$-模。对态射 $g : T \to S$，考虑基变换图
$$
\xymatrix{
X_T \ar[d]_p \ar[r]_q & X \ar[d]^f \\
T \ar[r]^g & S
}
$$
。假设对所有 $g : T \to S$，$\mathcal{O}_T \to p_*\mathcal{O}_{X_T}$
均为同构。则存在有限表现浸入 $j : Z \to S$，使得态射 $g : T \to S$
经 $Z$ 分解，当且仅当存在有限局部自由 $\mathcal{O}_T$-模 $\mathcal{N}$，满足
$p^*\mathcal{N} \cong q^*\mathcal{E}$。
\end{lemma}

\begin{proof}
由假设 $H^0(X_s, \mathcal{O}_{X_s}) = \kappa(s)$，注意 $f$ 的纤维 $X_s$
连通。因此 $\mathcal{E}$ 的秩在纤维上为常数。由于 $f$ 是开态射（Morphisms,
引理 \ref{morphisms-lemma-fppf-open}）且是闭态射，$S$ 有分解
$S = \coprod S_r$ 为开闭子概形，使得 $\mathcal{E}$ 在 $S_r$
的逆像上具有常秩 $r$。因此可假设 $\mathcal{E}$ 具有常秩 $r$。以
$\mathcal{E}^\vee = \SheafHom(\mathcal{E}, \mathcal{O}_X)$ 表示其秩为 $r$ 的对偶模。

\medskip\noindent
由上同调与基变换（更确切地，由 概形的导出范畴章引理
\ref{perfect-lemma-flat-proper-perfect-direct-image-general}），
$E = Rf_*\mathcal{E}$ 是 $S$ 的导出范畴中的完美对象，且其构造与任意基变换交换。
$E' = Rf_*\mathcal{E}^\vee$ 同样如此。由于在次数 $< 0$ 中绝无上同调，可见对某个
$b$，$E$ 和 $E'$（局部）具有位于 $[0, b]$ 中的 tor-振幅。注意，对任意
$g : T \to S$，有 $p_*(q^*\mathcal{E}) = H^0(Lg^*E)$ 以及
$p_*(q^*\mathcal{E}^\vee) = H^0(Lg^*E')$。令 $j : Z \to S$ 和
$j' : Z' \to S$ 为 概形的导出范畴章引理
\ref{perfect-lemma-locally-closed-where-H0-locally-free} 对 $E$ 和 $E'$、取
$a = 0$ 和 $r = r$ 时构造的有限表现浸入；粗略地说，它们由如下性质刻画：
$H^0(Lj^*E)$ 和 $H^0((j')^*E')$ 是与拉回相容的有限局部自由模。

\medskip\noindent
设 $g : T \to S$ 是态射。若存在引理中的 $\mathcal{N}$，则使用投影公式
上同调章引理 \ref{cohomology-lemma-projection-formula}，可见各模
$$
p_*(q^*\mathcal{E}) \cong
p_*(p^*\mathcal{N}) \cong
\mathcal{N} \otimes_{\mathcal{O}_T} p_*\mathcal{O}_{X_T} \cong
\mathcal{N}\quad\text{and similarly }\quad
p_*(q^*\mathcal{E}^\vee) \cong \mathcal{N}^\vee
$$
是秩为 $r$ 的有限局部自由模，而且在任何进一步的基变换 $T' \to T$ 后仍然如此。
因此此时 $T \to S$ 经 $j$ 和 $j'$ 分解。故可将 $S$ 换成 $Z \times_S Z'$，
并假设 $f_*\mathcal{E}$ 和 $f_*\mathcal{E}^\vee$ 是秩为 $r$ 的有限局部自由
$\mathcal{O}_S$-模，其构造与任意基变换交换（略去一个小细节）。

\medskip\noindent
% P05-ERR-CH37-0060: Frozen source has the ungrammatical phrase ``if g ... be a morphism''.
在此情形，若 $g : T \to S$ 是态射并且存在引理中的 $\mathcal{N}$，则映射（次数
$0$ 中的杯积）
$$
p_*(q^*\mathcal{E})
\otimes_{\mathcal{O}_T}
p_*(q^*\mathcal{E}^\vee)
\longrightarrow \mathcal{O}_T
$$
是完美配对。反之，若此杯积映射是完美配对，则在 $T$ 上局部可以分别选取
$p_*(q^*\mathcal{E})$ 中的截面基 $\sigma_1, \ldots, \sigma_r$ 和
$p_*(q^*\mathcal{E}^\vee)$ 中的截面基 $\tau_1, \ldots, \tau_r$，使其乘积满足
$\sigma_i \tau_j = \delta_{ij}$。将 $\sigma_i$ 视为 $X_T$ 上 $q^*\mathcal{E}$
的截面，并将 $\tau_j$ 视为 $X_T$ 上 $q^*\mathcal{E}^\vee$ 的截面，可知
$$
\sigma_1, \ldots, \sigma_r :
\mathcal{O}_{X_T}^{\oplus r}
\longrightarrow
q^*\mathcal{E}
$$
是同构，其逆由下式给出：
$$
\tau_1, \ldots, \tau_r :
q^*\mathcal{E}
\longrightarrow
\mathcal{O}_{X_T}^{\oplus r}
$$
。换言之，有 $p^*p_*q^*\mathcal{E} \cong q^*\mathcal{E}$。而杯积非退化这一条件截出一个逆紧开子概形
（即某个适当行列式非零的轨迹），证明完成。
\end{proof}

\noindent
上述引理特别说明：若向量丛在一个固有空间的固有平坦族的各纤维上平凡，则它是某个向量丛的拉回。下面稍作展开。

\begin{lemma}
\label{more-morphisms-lemma-trivial-on-fibres}
设 $f : X \to S$ 是有限表现的平坦固有态射，满足
$f_*\mathcal{O}_X = \mathcal{O}_S$，并且此性质在任意基变换后仍成立。设
$\mathcal{E}$ 是有限局部自由 $\mathcal{O}_X$-模。假设
\begin{enumerate}
\item 对所有 $s \in S$，$\mathcal{E}|_{X_s}$ 同构于
$\mathcal{O}_{X_s}^{\oplus r_s}$；并且
\item $S$ 约化。
\end{enumerate}
则对某个有限局部自由 $\mathcal{O}_S$-模 $\mathcal{N}$，有
$\mathcal{E} = f^*\mathcal{N}$。
\end{lemma}

\begin{proof}
事实上，此时引理 \ref{more-morphisms-lemma-diagonal-picard-flat-proper} 中的局部闭浸入
$j : Z \to S$ 是双射，因而是闭浸入。但由于 $S$ 约化，$j$ 是同构。
\end{proof}

\begin{lemma}
\label{more-morphisms-lemma-triviality-generic-fibre-valuation-ring}
设 $f : X \to S$ 是有限表现的固有平坦态射。设 $\mathcal{L}$ 是可逆
$\mathcal{O}_X$-模。假设
\begin{enumerate}
\item $S$ 是赋值环的谱；
\item $\mathcal{L}$ 在 $f$ 的泛纤维 $X_\eta$ 上平凡；
\item $f$ 的闭纤维 $X_0$ 是整概形；
\item $H^0(X_\eta, \mathcal{O}_{X_\eta})$ 等于 $S$ 的函数域。
\end{enumerate}
则 $\mathcal{L}$ 平凡。
\end{lemma}

\begin{proof}
写成 $S = \Spec(A)$。先在 $A$ 为离散赋值环时证明本引理（这是实践中最常用的情形）。
令 $\pi \in A$ 为一致化元。取平凡化截面
$s \in \Gamma(X_\eta, \mathcal{L}_\eta)$。必要时以 $\pi^n s$ 代替 $s$ 后，
可假设 $s \in \Gamma(X, \mathcal{L})$。若 $s|_{X_0} = 0$，则 $s$ 可被
$\pi$ 整除（因为 $X_0$ 是概形论纤维且 $X$ 在 $A$ 上平坦）。因此可假设
$s|_{X_0}$ 非零。于是 $s$ 的零化轨迹 $Z(s)$ 包含于 $X_0$，但不包含 $X_0$
的泛点（因为 $X_0$ 是整概形）。这意味着
% P05-ERR-CH37-0061: Frozen source has the malformed phrase ``the Z(s)''.
$Z(s)$ 在 $X$ 中的余维 $\geq 2$；由 除子章引理
\ref{divisors-lemma-effective-Cartier-codimension-1}，这将导致矛盾，除非
$Z(s) = \emptyset$，正是所需。

\medskip\noindent
现在证明一般情形。赋值环 $A$ 是凝聚环（代数章例
\ref{algebra-example-valuation-ring-coherent}），故 $H^0(X, \mathcal{L})$ 是凝聚
$A$-模，见 概形的导出范畴章引理
\ref{perfect-lemma-cohomology-over-coherent-ring}。等价地，$H^0(X, \mathcal{L})$
是有限表现 $A$-模（代数章引理 \ref{algebra-lemma-coherent-ring}）。由于
$H^0(X, \mathcal{L})$ 无挠（由 $X$ 在 $A$ 上平坦），代数进阶章引理
\ref{more-algebra-lemma-generalized-valuation-ring-modules} 表明，对某个 $n$，有
$H^0(X, \mathcal{L}) = A^{\oplus n}$。由平坦基变换（Cohomology of Schemes,
引理 \ref{coherent-lemma-flat-base-change-cohomology}），有
$$
K = H^0(X_\eta, \mathcal{O}_{X_\eta}) \cong
H^0(X_\eta, \mathcal{L}_\eta) =
H^0(X, \mathcal{L}) \otimes_A K
$$
其中 $K$ 是 $A$ 的分式域。因此 $n = 1$。选取生成元
$s \in H^0(X, \mathcal{L})$。令 $\mathfrak m \subset A$ 为极大理想。则
$\kappa = A/\mathfrak m = \colim A/\pi$，这里取遍非零元
$\pi \in \mathfrak m$ 的滤过余极限（此处使用 $A$ 是赋值环）。因此
$X_0 = \lim X \times_S \Spec(A/\pi)$。若 $s|_{X_0}$ 为零，则对某个 $\pi$，
$s$ 在 $X \times_S \Spec(A/\pi)$ 上的限制为零，见 极限章引理
\ref{limits-lemma-descend-section}。但若如此，$\pi^{-1} s$ 就是 $\mathcal{L}$
的整体截面，这与 $s$ 是 $H^0(X, \mathcal{L})$ 的生成元矛盾。因此
$s|_{X_0}$ 非零。令 $Z(s) \subset X$ 为 $s$ 的零概形。由于 $s|_{X_0}$
非零且 $X_0$ 是整概形，$Z(s)_0 \subset X_0$ 是有效 Cartier 除子。由于 $f$
固有且 $S$ 局部，$Z(s)$ 的每一点都特殊化到 $Z(s)_0$ 中的一点。因此由 Divisors,
引理 \ref{divisors-lemma-fibre-Cartier} 的第 (3) 部分，$Z(s)$ 是相对有效
Cartier 除子；特别地，$Z(s) \to S$ 平坦。因此，若
% P05-ERR-CH37-0062: Frozen source has the typo ``nonemtpy''.
$Z(s)$ 非空，则 $Z(s)_\eta$ 也非空，这与 $s|_{X_\eta}$ 是
$\mathcal{L}_\eta$ 的平凡化矛盾。故 $Z(s) = \emptyset$，正是所需。
\end{proof}

\begin{lemma}
\label{more-morphisms-lemma-get-a-closed}
设 $f : X \to S$ 和 $\mathcal{E}$ 如引理
\ref{more-morphisms-lemma-diagonal-picard-flat-proper}，并且进一步假设 $\mathcal{E}$ 是可逆
$\mathcal{O}_X$-模。若 $f$ 的几何纤维还是整概形，则 $Z$ 在 $S$ 中闭。
\end{lemma}

\begin{proof}
由于 $j : Z \to S$ 有限表现，只需证明：对任意态射
$g : \Spec(A) \to S$，其中 $A$ 是分式域为 $K$ 的赋值环，若
$g(\Spec(K)) \in j(Z)$，则 $g(\Spec(A)) \subset j(Z)$。见 Morphisms,
引理 \ref{morphisms-lemma-reach-points-scheme-theoretic-image}。这由引理
\ref{more-morphisms-lemma-triviality-generic-fibre-valuation-ring} 以及引理
\ref{more-morphisms-lemma-diagonal-picard-flat-proper} 中对 $j : Z \to S$ 的刻画得出。
\end{proof}

\begin{lemma}
\label{more-morphisms-lemma-H1-O-picard-flat-proper}
考虑概形的交换图
$$
\xymatrix{
X' \ar[rr] \ar[dr]_{f'} & & X \ar[dl]^f \\
& S
}
$$
，其中 $f' : X' \to S$ 和 $f : X \to S$ 满足引理
\ref{more-morphisms-lemma-diagonal-picard-flat-proper} 的假设。设 $\mathcal{L}$ 是可逆
$\mathcal{O}_X$-模，并令 $\mathcal{L}'$ 为其在 $X'$ 上的拉回。分别令
$Z \subset S$ 和 $Z' \subset S$ 为引理
\ref{more-morphisms-lemma-diagonal-picard-flat-proper} 对 $(f, \mathcal{L})$ 和
$(f', \mathcal{L}')$ 构造的局部闭子概形，于是 $Z \subset Z'$。若 $s \in Z$ 且
$$
H^1(X_s, \mathcal{O}) \longrightarrow H^1(X'_s, \mathcal{O})
$$
为单射，则对 $s$ 的某个开邻域 $U$，有 $Z \cap U = Z' \cap U$。
\end{lemma}

\begin{proof}
可将 $S$ 换成 $Z'$。将 $S$ 缩小为 $s$ 的仿射开邻域后，可假设
$\mathcal{L}' = \mathcal{O}_{X'}$。令
$E = Rf_*\mathcal{L}$ 以及
$E' = Rf'_*\mathcal{L}' = Rf'_*\mathcal{O}_{X'}$。它们是完美复形，且其构造与任意基变换交换
（概形的导出范畴章引理
\ref{perfect-lemma-flat-proper-perfect-direct-image-general}）。特别地，
$$
E \otimes_{\mathcal{O}_S}^\mathbf{L} \kappa(s) =
R\Gamma(X_s, \mathcal{L}_s) = R\Gamma(X_s, \mathcal{O}_{X_s})
$$
，其中第二个等式源于 $s \in Z$。置
$h_i = \dim_{\kappa(s)} H^i(X_s, \mathcal{O}_{X_s})$。缩小 $S$ 后，可用复形
$$
\mathcal{O}_S \to \mathcal{O}_S^{\oplus h_1} \to
\mathcal{O}_S^{\oplus h_2} \to \ldots
$$
表示 $E$，见 代数进阶章引理
\ref{more-algebra-lemma-lift-perfect-from-residue-field}（严格地说，这还用到
概形的导出范畴章引理
\ref{perfect-lemma-affine-compare-bounded} 和
\ref{perfect-lemma-perfect-affine}）。类似地，可假设 $E'$ 由复形
$$
\mathcal{O}_S \to \mathcal{O}_S^{\oplus h'_1} \to
\mathcal{O}_S^{\oplus h'_2} \to \ldots
$$
表示，其中 $h'_i = \dim_{\kappa(s)} H^i(X'_s, \mathcal{O}_{X'_s})$。由上同调的函子性，有映射
$$
E \longrightarrow E'
$$
位于 $D(\mathcal{O}_S)$ 中，且其构造与基变换交换。由于表示 $E$ 的复形是有限自由模的有限复形，
而 $S$ 仿射，可以选取复形映射
$$
\xymatrix{
\mathcal{O}_S \ar[r]_d \ar[d]_a &
\mathcal{O}_S^{\oplus h_1} \ar[r] \ar[d]_b &
\mathcal{O}_S^{\oplus h_2} \ar[r] \ar[d]_c & \ldots \\
\mathcal{O}_S \ar[r]^{d'} &
\mathcal{O}_S^{\oplus h'_1} \ar[r] &
\mathcal{O}_S^{\oplus h'_2} \ar[r] & \ldots
}
$$
表示给定映射 $E \to E'$。由于 $s \in Z$，$\mathcal{L}_s$ 的平凡化截面拉回为
$\mathcal{L}'_s = \mathcal{O}_{X'_s}$ 的平凡化截面。因此
$a \otimes \kappa(s)$ 是同构，故缩小 $S$ 后，$a$ 是同构。最后，使用
$H^1(X_s, \mathcal{O}) \to H^1(X'_s, \mathcal{O})$ 为单射这一假设，可知定义
$b$ 的矩阵有一个 $h_1 \times h_1$ 子式映到 $\kappa(s)$ 中的非零元。因此缩小
$S$ 后，可假设 $b$ 为单射。然而，由于 $\mathcal{L}' = \mathcal{O}_{X'}$，有
$d' = 0$。于是 $d = 0$。由此可见，$\mathcal{L}_s$ 的平凡化截面提升为 $X$
上 $\mathcal{L}$ 的截面。一个直接的拓扑论证（略）表明，必要时进一步略微缩小
$S$ 后，$\mathcal{L}$ 平凡。
\end{proof}

\begin{lemma}
\label{more-morphisms-lemma-H1-O-multiple-picard-flat-proper}
考虑概形的 $n$ 个交换图
$$
\xymatrix{
X_i \ar[rr] \ar[dr]_{f_i} & & X \ar[dl]^f \\
& S
}
$$
，其中 $f_i : X_i \to S$ 和 $f : X \to S$ 满足引理
\ref{more-morphisms-lemma-diagonal-picard-flat-proper} 的假设。设 $\mathcal{L}$ 是可逆
$\mathcal{O}_X$-模，并令 $\mathcal{L}_i$ 为其在 $X_i$ 上的拉回。分别令
$Z \subset S$ 和 $Z_i \subset S$ 为引理
\ref{more-morphisms-lemma-diagonal-picard-flat-proper} 对 $(f, \mathcal{L})$ 和
$(f_i, \mathcal{L}_i)$ 构造的局部闭子概形，于是
$Z \subset \bigcap_{i = 1, \ldots, n} Z_i$。若 $s \in Z$ 且
$$
H^1(X_s, \mathcal{O}) \longrightarrow
\bigoplus\nolimits_{i = 1, \ldots, n} H^1(X_{i, s}, \mathcal{O})
$$
为单射，则对 $s$ 的某个开邻域 $U$，有
$Z \cap U = (\bigcap_{i = 1, \ldots, n} Z_i) \cap U$（概形论交）。
\end{lemma}

\begin{proof}
本引理是引理 \ref{more-morphisms-lemma-H1-O-picard-flat-proper} 的一个变体；我们强烈建议读者先阅读该证明。
这里的证明基本上是对该证明稍作修改后的复本。由 $Z$ 和各 $Z_i$ 的（概形值）点的描述，
有 $Z \subset \bigcap_{i = 1, \ldots, n} Z_i$，其中取概形论交。因此可将 $S$
换成概形论交 $\bigcap_{i = 1, \ldots, n} Z_i$。将 $S$ 缩小为 $s$ 的仿射开邻域后，
可假设对 $i = 1, \ldots, n$ 有 $\mathcal{L}_i = \mathcal{O}_{X_i}$。令
$E = Rf_*\mathcal{L}$ 以及
$E_i = Rf_{i, *}\mathcal{L}_i = Rf_{i, *}\mathcal{O}_{X_i}$。它们是完美复形，
且其构造与任意基变换交换（概形的导出范畴章引理
\ref{perfect-lemma-flat-proper-perfect-direct-image-general}）。特别地，
$$
E \otimes_{\mathcal{O}_S}^\mathbf{L} \kappa(s) =
R\Gamma(X_s, \mathcal{L}_s) = R\Gamma(X_s, \mathcal{O}_{X_s})
$$
，其中第二个等式源于 $s \in Z$。置
$h_j = \dim_{\kappa(s)} H^j(X_s, \mathcal{O}_{X_s})$。缩小 $S$ 后，可用复形
$$
\mathcal{O}_S \to \mathcal{O}_S^{\oplus h_1} \to
\mathcal{O}_S^{\oplus h_2} \to \ldots
$$
表示 $E$，见 代数进阶章引理
\ref{more-algebra-lemma-lift-perfect-from-residue-field}（严格地说，这还用到
概形的导出范畴章引理
\ref{perfect-lemma-affine-compare-bounded} 和
\ref{perfect-lemma-perfect-affine}）。类似地，可假设 $E_i$ 由复形
$$
\mathcal{O}_S \to \mathcal{O}_S^{\oplus h_{i, 1}} \to
\mathcal{O}_S^{\oplus h_{i, 2}} \to \ldots
$$
表示，其中 $h_{i, j} = \dim_{\kappa(s)} H^j(X_{i, s}, \mathcal{O}_{X_{i, s}})$。
由上同调的函子性，有映射
$$
E \longrightarrow E_i
$$
位于 $D(\mathcal{O}_S)$ 中，且其构造与基变换交换。由于表示 $E$ 的复形是有限自由模的有限复形，
而 $S$ 仿射，可以选取复形映射
$$
\xymatrix{
\mathcal{O}_S \ar[r]_d \ar[d]_{a_i} &
\mathcal{O}_S^{\oplus h_1} \ar[r] \ar[d]_{b_i} &
\mathcal{O}_S^{\oplus h_2} \ar[r] \ar[d]_{c_i} & \ldots \\
\mathcal{O}_S \ar[r]^{d_i} &
\mathcal{O}_S^{\oplus h_{i, 1}} \ar[r] &
\mathcal{O}_S^{\oplus h_{i, 2}} \ar[r] & \ldots
}
$$
表示给定映射 $E \to E_i$。由于 $s \in Z$，$\mathcal{L}_s$ 的平凡化截面拉回为
$\mathcal{L}_{i, s} = \mathcal{O}_{X_{i, s}}$ 的平凡化截面。因此
$a_i \otimes \kappa(s)$ 是同构，故缩小 $S$ 后，$a_i$ 是同构。最后，使用
$H^1(X_s, \mathcal{O}) \to
\bigoplus_{i = 1, \ldots, n} H^1(X_{i, s}, \mathcal{O})$ 为单射这一假设，
可知定义 $\oplus b_i$ 的矩阵有一个 $h_1 \times h_1$ 子式映到 $\kappa(s)$
中的非零元。因此缩小 $S$ 后，可假设
% P05-ERR-CH37-0063: Frozen source writes O_S^{h_1} here, unlike the direct-sum notation used for the surrounding free modules.
$(b_1, \ldots, b_n) : \mathcal{O}_S^{h_1}
\to \bigoplus_{i = 1, \ldots, n} \mathcal{O}_S^{h_{i, 1}}$
为单射。然而，由于 $\mathcal{L}_i = \mathcal{O}_{X_i}$，对
$i = 1, \ldots n$ 有 $d_i = 0$。又因
$(b_1, \ldots, b_n) \circ d = (\oplus d_i) \circ (a_1, \ldots, a_n)$，故
$d = 0$。由此可见，$\mathcal{L}_s$ 的平凡化截面提升为 $X$ 上 $\mathcal{L}$
的截面。一个直接的拓扑论证（略）表明，必要时进一步略微缩小 $S$ 后，
$\mathcal{L}$ 平凡。
\end{proof}

\begin{lemma}
\label{more-morphisms-lemma-pic-of-product}
设 $f : X \to S$ 和 $g : Y \to S$ 是满足引理
\ref{more-morphisms-lemma-diagonal-picard-flat-proper} 之假设的概形态射。设
$\sigma : S \to X$ 和 $\tau : S \to Y$ 分别为 $f$ 和 $g$ 的截面。设
$s \in S$。设 $\mathcal{L}$ 是 $X \times_S Y$ 上的可逆层。若 $X$ 上的
$(1 \times \tau)^*\mathcal{L}$、$Y$ 上的 $(\sigma \times 1)^*\mathcal{L}$ 以及
$\mathcal{L}|_{(X \times_S Y)_s}$ 均平凡，则存在 $s$ 的开邻域 $U$，使得
$\mathcal{L}$ 在 $(X \times_S Y)_U$ 上平凡。
\end{lemma}

\begin{proof}
由 K\"unneth 公式（簇章引理 \ref{varieties-lemma-kunneth}），映射
$$
H^1(X_s \times_{\Spec(\kappa(s))} Y_s, \mathcal{O}) \to
H^1(X_s, \mathcal{O}) \oplus H^1(Y_s, \mathcal{O})
$$
为单射。因此可将引理 \ref{more-morphisms-lemma-H1-O-multiple-picard-flat-proper} 应用于两个态射
$$
1 \times \tau : X \to X \times_S Y
\quad\text{且}\quad
\sigma \times 1 : Y \to X \times_S Y
$$
，从而得到结论。
\end{proof}

\begin{theorem}[立方体定理]
\label{more-morphisms-theorem-of-the-cube}
设 $S$ 是概形。设 $X$、$Y$ 和 $Z$ 是 $S$ 上的概形。设 $x : S \to X$ 和
$y : S \to Y$ 为结构态射的截面。设 $\mathcal{L}$ 是
$X \times_S Y \times_S Z$ 上的可逆模。若
\begin{enumerate}
\item $X \to S$ 和 $Y \to S$ 是具有几何整纤维的有限表现平坦固有态射；
\item $\mathcal{L}$ 经 $x \times \text{id}_Y \times \text{id}_Z$ 和
$\text{id}_X \times y \times \text{id}_Z$ 的拉回分别在 $Y \times_S Z$ 和
$X \times_S Z$ 上平凡；
\item 存在一点 $z \in Z$，使得 $\mathcal{L}$ 在
$X \times_S Y \times_S z$ 上的限制平凡；并且
\item $Z$ 连通；
\end{enumerate}
则 $\mathcal{L}$ 平凡。
\end{theorem}

\noindent
一个常用的特殊情形如下。设 $k$ 是域。设 $X, Y, Z$ 是带有 $k$-有理点
$x, y, z$ 的簇。设 $\mathcal{L}$ 是 $X \times Y \times Z$ 上的可逆模。若
\begin{enumerate}
\item $\mathcal{L}$ 在 $x \times Y \times Z$、$X \times y \times Z$ 和
$X \times Y \times z$ 上平凡；并且
\item $X$ 和 $Y$ 在 $k$ 上几何整且固有；
\end{enumerate}
则 $\mathcal{L}$ 平凡。

\begin{proof}
% P05-ERR-CH37-0064: Frozen source omits the verb in ``whose geometrically integral fibres''.
注意，态射 $X \times_S Y \to S$ 是具有几何整纤维的有限表现平坦固有态射（关于纤维的陈述，
见 簇章引理 \ref{varieties-lemma-geometrically-integral}、
\ref{varieties-lemma-bijection-irreducible-components} 和
\ref{varieties-lemma-geometrically-reduced-any-base-change}）。由 Derived Categories of Schemes,
引理 \ref{perfect-lemma-proper-flat-geom-red-connected}，结构层经 $X \to S$、
$Y \to S$ 或 $X \times_S Y \to S$ 的推前都是 $S$ 的结构层，并且这在任何基变换后仍成立。
因此可将引理 \ref{more-morphisms-lemma-diagonal-picard-flat-proper} 应用于态射
$$
p : X \times_S Y \times_S Z \longrightarrow Z
$$
和可逆模 $\mathcal{L}$，得到“泛”局部闭子概形 $Z' \subset Z$，使得
$\mathcal{L}|_{X \times_S Y \times_S Z'}$ 是 $Z'$ 上某个可逆模 $\mathcal{N}$ 的拉回。
$z$ 的存在表明 $Z'$ 非空。由引理 \ref{more-morphisms-lemma-get-a-closed}，$Z' \subset Z$
是闭子概形。设 $z' \in Z'$ 为一点。注意，可将 $p$ 写成积态射
$$
(X \times_S Z) \times_Z (Y \times_S Z) \longrightarrow Z
$$
因此可将引理 \ref{more-morphisms-lemma-pic-of-product} 应用于态射 $p$、点 $z'$ 以及由 $x$ 和 $y$
给出的截面 $\sigma : Z \to X \times_S Z$ 和 $\tau : Z \to Y \times_S Z$。
由此 $Z'$ 是开集。因此 $Z' = Z$，且对 $Z$ 上某个可逆模 $\mathcal{N}$，有
$\mathcal{L} = p^*\mathcal{N}$。经
$x \times y \times \text{id}_Z : Z  \to X \times_S Y \times_S Z$
拉回，一方面得到 $\mathcal{N}$，另一方面由假设 (2) 得到平凡可逆模。因此
$\mathcal{N} = \mathcal{O}_Z$，证明完成。
\end{proof}









\section{极限论证}
\label{more-morphisms-section-limits}

\noindent
本节给出一些涉及概形极限与 Noether 逼近的引理。我们主要限于仿射情形。其中一些引理是极限章中引理的特殊情形。

\begin{lemma}
\label{more-morphisms-lemma-Noetherian-approximation}
设 $f : X \to S$ 是有限表现的仿射概形态射。则存在 Cartesian 图
$$
\xymatrix{
X_0 \ar[d]_{f_0} & X \ar[l]^g \ar[d]^f \\
S_0 & S \ar[l]
}
$$
使得
\begin{enumerate}
\item $X_0$、$S_0$ 是仿射概形；
\item $S_0$ 在 $\mathbf{Z}$ 上有限型；
\item $f_0$ 是有限型的。
\end{enumerate}
\end{lemma}

\begin{proof}
写成 $S = \Spec(A)$ 和 $X = \Spec(B)$。由于 $f$ 有限表现，$B$
作为 $A$-代数有限表现，见 态射章引理
\ref{morphisms-lemma-locally-finite-presentation-characterize}。因此本引理由
代数章引理 \ref{algebra-lemma-limit-module-finite-presentation} 得出。
\end{proof}

\begin{lemma}
\label{more-morphisms-lemma-Noetherian-approximation-module}
设 $f : X \to S$ 是有限表现的仿射概形态射。设 $\mathcal{F}$ 是有限表现的拟凝聚
$\mathcal{O}_X$-模。则存在引理 \ref{more-morphisms-lemma-Noetherian-approximation} 中的图，并且存在凝聚
$\mathcal{O}_{X_0}$-模 $\mathcal{F}_0$，满足 $g^*\mathcal{F}_0 = \mathcal{F}$。
\end{lemma}

\begin{proof}
写成 $S = \Spec(A)$、$X = \Spec(B)$ 以及 $\mathcal{F} = \widetilde{M}$。
由于 $f$ 有限表现，$B$ 作为 $A$-代数有限表现，见 态射章引理
\ref{morphisms-lemma-locally-finite-presentation-characterize}。由于 $\mathcal{F}$
作为 $\mathcal{O}_X$-模有限表现，$M$ 作为 $B$-模有限表现，见 Properties,
引理 \ref{properties-lemma-finite-presentation-module}。因此本引理由 代数章引理
\ref{algebra-lemma-limit-module-finite-presentation} 得出。
\end{proof}

\begin{lemma}
\label{more-morphisms-lemma-Noetherian-approximation-flat-module}
设 $f : X \to S$ 是有限表现的仿射概形态射。设 $\mathcal{F}$ 是有限表现的拟凝聚
$\mathcal{O}_X$-模，并且在 $S$ 上平坦。则可选取引理
\ref{more-morphisms-lemma-Noetherian-approximation-module} 中的图和层 $\mathcal{F}_0$，使得此外
$\mathcal{F}_0$ 在 $S_0$ 上平坦。
\end{lemma}

\begin{proof}
写成 $S = \Spec(A)$、$X = \Spec(B)$ 以及 $\mathcal{F} = \widetilde{M}$。
由于 $f$ 有限表现，$B$ 作为 $A$-代数有限表现，见 态射章引理
\ref{morphisms-lemma-locally-finite-presentation-characterize}。由于 $\mathcal{F}$
作为 $\mathcal{O}_X$-模有限表现，$M$ 作为 $B$-模有限表现，见 Properties,
引理 \ref{properties-lemma-finite-presentation-module}。由于 $\mathcal{F}$ 在 $S$
上平坦，$M$ 在 $A$ 上平坦，见 态射章引理
\ref{morphisms-lemma-flat-module-characterize}。因此本引理由 代数章引理
\ref{algebra-lemma-flat-finite-presentation-limit-flat} 得出。
\end{proof}

\begin{lemma}
\label{more-morphisms-lemma-Noetherian-approximation-flat}
设 $f : X \to S$ 是有限表现且平坦的仿射概形态射。则存在引理
\ref{more-morphisms-lemma-Noetherian-approximation} 中的图，使得此外 $f_0$ 平坦。
\end{lemma}

\begin{proof}
这是引理 \ref{more-morphisms-lemma-Noetherian-approximation-flat-module} 的特殊情形。
\end{proof}

\begin{lemma}
\label{more-morphisms-lemma-Noetherian-approximation-smooth}
设 $f : X \to S$ 是仿射概形的光滑态射。则存在引理
\ref{more-morphisms-lemma-Noetherian-approximation} 中的图，使得此外 $f_0$ 光滑。
\end{lemma}

\begin{proof}
写成 $S = \Spec(A)$、$X = \Spec(B)$。由于 $f$ 光滑，$B$ 作为 $A$-代数光滑，
见 态射章引理 \ref{morphisms-lemma-smooth-characterize}。因此本引理由
代数章引理 \ref{algebra-lemma-finite-presentation-fs-Noetherian} 得出。
\end{proof}

\begin{lemma}
\label{more-morphisms-lemma-Noetherian-approximation-geometrically-reduced}
设 $f : X \to S$ 是具有几何约化纤维的有限表现仿射概形态射。则存在引理
\ref{more-morphisms-lemma-Noetherian-approximation} 中的图，使得此外 $f_0$ 具有几何约化纤维。
\end{lemma}

\begin{proof}
应用引理 \ref{more-morphisms-lemma-Noetherian-approximation} 得到仿射概形的 Cartesian 图
$$
\xymatrix{
X_0 \ar[d]_{f_0} & X \ar[l]^g \ar[d]^f \\
S_0 & S \ar[l]_h
}
$$
，其中 $X_0 \to S_0$ 是在 $\mathbf{Z}$ 上有限型的概形之间的有限型态射。由引理
\ref{more-morphisms-lemma-geometrically-reduced-constructible}，其中 $f_0$ 的纤维几何约化的点集
$E \subset S_0$ 是可构造子集。由引理
\ref{more-morphisms-lemma-base-change-fibres-geometrically-reduced}，有 $h(S) \subset E$。写成
$S_0 = \Spec(A_0)$ 和 $S = \Spec(A)$。将 $A = \colim_i A_i$ 写成有限型
$A_0$-代数的直余极限。由 极限章引理
\ref{limits-lemma-limit-contained-in-constructible}，对某个 $i$，
$\Spec(A_i) \to S_0$ 的像包含于 $E$。将 $S_0$ 换成 $\Spec(A_i)$，并将
$X_0$ 换成 $X_0 \times_{S_0} \Spec(A_i)$ 后，可见 $f_0$ 的所有纤维均几何约化。
\end{proof}

\begin{lemma}
\label{more-morphisms-lemma-Noetherian-approximation-geometrically-irreducible}
设 $f : X \to S$ 是具有几何不可约纤维的有限表现仿射概形态射。则存在引理
\ref{more-morphisms-lemma-Noetherian-approximation} 中的图，使得此外 $f_0$ 具有几何不可约纤维。
\end{lemma}

\begin{proof}
应用引理 \ref{more-morphisms-lemma-Noetherian-approximation} 得到仿射概形的 Cartesian 图
$$
\xymatrix{
X_0 \ar[d]_{f_0} & X \ar[l]^g \ar[d]^f \\
S_0 & S \ar[l]_h
}
$$
，其中 $X_0 \to S_0$ 是在 $\mathbf{Z}$ 上有限型的概形之间的有限型态射。由引理
\ref{more-morphisms-lemma-nr-geom-irreducible-components-constructible}，其中 $f_0$
的纤维几何不可约的点集 $E \subset S_0$ 是可构造子集。由引理
\ref{more-morphisms-lemma-base-change-fibres-geometrically-irreducible}，有 $h(S) \subset E$。写成
$S_0 = \Spec(A_0)$ 和 $S = \Spec(A)$。将 $A = \colim_i A_i$ 写成有限型
$A_0$-代数的直余极限。由 极限章引理
\ref{limits-lemma-limit-contained-in-constructible}，对某个 $i$，
$\Spec(A_i) \to S_0$ 的像包含于 $E$。将 $S_0$ 换成 $\Spec(A_i)$，并将
$X_0$ 换成 $X_0 \times_{S_0} \Spec(A_i)$ 后，可见 $f_0$ 的所有纤维均几何不可约。
\end{proof}

\begin{lemma}
\label{more-morphisms-lemma-Noetherian-approximation-geometrically-connected}
设 $f : X \to S$ 是具有几何连通纤维的有限表现仿射概形态射。则存在引理
\ref{more-morphisms-lemma-Noetherian-approximation} 中的图，使得此外 $f_0$ 具有几何连通纤维。
\end{lemma}

\begin{proof}
应用引理 \ref{more-morphisms-lemma-Noetherian-approximation} 得到仿射概形的 Cartesian 图
$$
\xymatrix{
X_0 \ar[d]_{f_0} & X \ar[l]^g \ar[d]^f \\
S_0 & S \ar[l]_h
}
$$
，其中 $X_0 \to S_0$ 是在 $\mathbf{Z}$ 上有限型的概形之间的有限型态射。由引理
\ref{more-morphisms-lemma-nr-geom-connected-components-constructible}，其中 $f_0$
的纤维几何连通的点集 $E \subset S_0$ 是可构造子集。由引理
\ref{more-morphisms-lemma-base-change-fibres-geometrically-connected}，有 $h(S) \subset E$。写成
$S_0 = \Spec(A_0)$ 和 $S = \Spec(A)$。将 $A = \colim_i A_i$ 写成有限型
$A_0$-代数的直余极限。由 极限章引理
\ref{limits-lemma-limit-contained-in-constructible}，对某个 $i$，
$\Spec(A_i) \to S_0$ 的像包含于 $E$。将 $S_0$ 换成 $\Spec(A_i)$，并将
$X_0$ 换成 $X_0 \times_{S_0} \Spec(A_i)$ 后，可见 $f_0$ 的所有纤维均几何连通。
\end{proof}

\begin{lemma}
\label{more-morphisms-lemma-Noetherian-approximation-dimension-d}
设 $d \geq 0$ 是整数。设 $f : X \to S$ 是有限表现的仿射概形态射，且其所有纤维的维数均为
$d$。则存在引理 \ref{more-morphisms-lemma-Noetherian-approximation} 中的图，使得此外 $f_0$
的所有纤维维数均为 $d$。
\end{lemma}

\begin{proof}
应用引理 \ref{more-morphisms-lemma-Noetherian-approximation} 得到仿射概形的 Cartesian 图
$$
\xymatrix{
X_0 \ar[d]_{f_0} & X \ar[l]^g \ar[d]^f \\
S_0 & S \ar[l]_h
}
$$
，其中 $X_0 \to S_0$ 是在 $\mathbf{Z}$ 上有限型的概形之间的有限型态射。由引理
\ref{more-morphisms-lemma-dimension-fibres-constructible}，其中 $f_0$ 的纤维维数为 $d$
的点集 $E \subset S_0$ 是可构造子集。由引理
\ref{more-morphisms-lemma-base-change-dimension-fibres}，有 $h(S) \subset E$。写成
$S_0 = \Spec(A_0)$ 和 $S = \Spec(A)$。将 $A = \colim_i A_i$ 写成有限型
$A_0$-代数的直余极限。由 极限章引理
\ref{limits-lemma-limit-contained-in-constructible}，对某个 $i$，
$\Spec(A_i) \to S_0$ 的像包含于 $E$。将 $S_0$ 换成 $\Spec(A_i)$，并将
$X_0$ 换成 $X_0 \times_{S_0} \Spec(A_i)$ 后，可见 $f_0$ 的所有纤维维数均为 $d$。
\end{proof}

\begin{lemma}
\label{more-morphisms-lemma-Noetherian-approximation-standard-syntomic}
设 $f : X \to S$ 是仿射概形的标准 syntomic 态射（见 态射章定义
\ref{morphisms-definition-syntomic}）。则存在引理
\ref{more-morphisms-lemma-Noetherian-approximation} 中的图，使得此外 $f_0$ 是标准 syntomic 态射。
\end{lemma}

\begin{proof}
本引理是 代数章引理
\ref{algebra-lemma-relative-global-complete-intersection-Noetherian} 的复本。
\end{proof}

\begin{lemma}
\label{more-morphisms-lemma-Noetherian-approximation-combine}
（Noether 逼近与性质合并。）设 $P$、$Q$ 是在基变换下稳定的概形态射性质。设
$f : X \to S$ 是有限表现的仿射概形态射。假设可以找到 Cartesian 图
$$
\vcenter{
\xymatrix{
X_1 \ar[d]_{f_1} & X \ar[l] \ar[d]^f \\
S_1 & S \ar[l]
}
}
\quad\text{且}\quad
\vcenter{
\xymatrix{
X_2 \ar[d]_{f_2} & X \ar[l] \ar[d]^f \\
S_2 & S \ar[l]
}
}
$$
，其中 $S_1$、$S_2$ 在 $\mathbf{Z}$ 上有限型，$f_1$、$f_2$ 是有限型的，且
$f_1$ 具有性质 $P$，$f_2$ 具有性质 $Q$。则可以找到 Cartesian 图
$$
\xymatrix{
X_0 \ar[d]_{f_0} & X \ar[l] \ar[d]^f \\
S_0 & S \ar[l]
}
$$
，其中 $S_0$ 在 $\mathbf{Z}$ 上有限型，$f_0$ 是有限型的，并且 $f_0$
同时具有性质 $P$ 与性质 $Q$。
\end{lemma}

\begin{proof}
给定的一对图对应于环的余 Cartesian 图
$$
\vcenter{
\xymatrix{
B_1 \ar[r] & B \\
A_1 \ar[u] \ar[r] & A \ar[u]
}
}
\quad\text{且}\quad
\vcenter{
\xymatrix{
B_2 \ar[r] & B \\
A_2 \ar[u] \ar[r] & A \ar[u]
}
}
$$
令 $A_0 \subset A$ 为 $A$ 的有限型 $\mathbf{Z}$-子代数，包含
$A_1 \to A$ 和 $A_2 \to A$ 的像。这样的子代数存在，因为按假设 $A_1$ 和
$A_2$ 均在 $\mathbf{Z}$ 上有限型。注意，各环
$B_{0, 1} = B_1 \otimes_{A_1} A_0$ 和 $B_{0, 2} = B_2 \otimes_{A_2} A_0$
是有限型 $A_0$-代数，并满足
$B_{0, 1} \otimes_{A_0} A \cong B \cong B_{0, 2} \otimes_{A_0} A$
作为 $A$-代数同构。由于 $A$ 是其有限型 $A_0$-子代数的有向余极限，由 Limits,
引理 \ref{limits-lemma-descend-finite-presentation}，增大 $A_0$ 后可假设存在
$A_0$-代数同构 $B_{0, 1} \cong B_{0, 2}$。由于假设性质 $P$ 和 $Q$
在基变换下稳定，置 $S_0 = \Spec(A_0)$ 以及
$$
X_0 = X_1 \times_{S_1} S_0 =
\Spec(B_{0, 1}) \cong \Spec(B_{0, 2}) = X_2 \times_{S_2} S_0
$$
即可满足要求。
\end{proof}












\section{\'Etale 邻域}
\label{more-morphisms-section-etale-neighbourhoods}

\noindent
事实表明，作 \'etale 基变换后，态射的某些性质更易研究。引入以下术语会很方便。

\begin{definition}
\label{more-morphisms-definition-etale-neighbourhood}
设 $S$ 是概形，$s \in S$ 是一点。
\begin{enumerate}
\item $(S, s)$ 的一个{\it \'etale 邻域}是一个二元组 $(U, u)$，连同概形的
\'etale 态射 $\varphi : U \to S$，满足 $\varphi(u) = s$。
\item $(S, s)$ 的 \'etale 邻域之间的一个{\it 态射}
$f : (V, v) \to (U, u)$，就是满足 $f(v) = u$ 的 $S$-概形态射 $f : V \to U$。
\item {\it 初等 \'etale 邻域}是满足 $\kappa(s) = \kappa(u)$ 的 \'etale 邻域
$\varphi : (U, u) \to (S, s)$。
\end{enumerate}
\end{definition}

\noindent
% P05-ERR-CH37-0065: Frozen source has mismatched parentheses around the two literature citations in this sentence.
初等 \'etale 邻域这一概念在文献中有许多不同名称，例如有时称为
“\'etale 邻域”（\cite[Page 36]{Milne}）或“强 \'etale”（\cite[Page 108]{KPR}）。
这里遵循论文 \cite{GruRay} 的约定，称其为初等 \'etale 邻域。

\medskip\noindent
若 $f : (V, v) \to (U, u)$ 是 \'etale 邻域之间的态射，则 $f$ 自动为
\'etale，见 态射章引理 \ref{morphisms-lemma-etale-permanence}。因此它使
$(V, v)$ 成为 $(U, u)$ 的 \'etale 邻域。当然，由于 \'etale 态射的复合仍为
\'etale（态射章引理 \ref{morphisms-lemma-composition-etale}），反过来，
$(U, u)$ 的任意 \'etale 邻域 $(V, v)$ 也为 $(S, s)$ 的 \'etale 邻域。还要指出，
若 $U \subset S$ 是 $s$ 的开邻域，则 $(U, s) \to (S, s)$ 是 \'etale 邻域。
这是因为开浸入为 \'etale 态射（态射章引理
\ref{morphisms-lemma-open-immersion-etale}）。本节将不再说明地使用这些评注。

\medskip\noindent
注意，若 $(U, u) \to (S, s)$ 是 \'etale 邻域，则
$\kappa(u)/\kappa(s)$ 是有限可分扩张，见 态射章引理
\ref{morphisms-lemma-etale-at-point}。

\begin{lemma}
\label{more-morphisms-lemma-realize-prescribed-residue-field-extension-etale}
设 $S$ 是概形，$s \in S$。设 $k/\kappa(s)$ 是有限可分域扩张。则存在 \'etale 邻域
$(U, u) \to (S, s)$，使得域扩张 $\kappa(u)/\kappa(s)$ 同构于 $k/\kappa(s)$。
\end{lemma}

\begin{proof}
可假设 $S$ 仿射。此时本引理由 代数章引理
\ref{algebra-lemma-make-etale-map-prescribed-residue-field} 得出。
\end{proof}

\begin{lemma}
\label{more-morphisms-lemma-etale-neighbourhoods-not-quite-filtered}
设 $S$ 是概形，并设 $s$ 是 $S$ 的一点。\'Etale 邻域范畴具有以下性质：
\begin{enumerate}
\item 设 $(U_i, u_i)_{i=1, 2}$ 是 $S$ 中 $s$ 的两个 \'etale 邻域。则存在第三个
\'etale 邻域 $(U, u)$ 以及态射
$(U, u) \to (U_i, u_i)$, $i = 1, 2$.
\item 设 $h_1, h_2: (U, u) \to (U', u')$ 是 $s$ 的 \'etale 邻域之间的两个态射。
假设 $h_1$、$h_2$ 诱导相同的剩余域映射 $\kappa(u') \to \kappa(u)$。则存在
\'etale 邻域 $(U'', u'')$ 以及态射 $h : (U'', u'') \to (U, u)$，使 $h_1$
与 $h_2$ 等化，即满足 $h_1 \circ h = h_2 \circ h$。
\end{enumerate}
\end{lemma}

\begin{proof}
对第 (1) 部分，考虑纤维积 $U = U_1 \times_S U_2$。由于 \'etale 态射在基变换下保持，
它在 $U_1$ 和 $U_2$ 上均为 \'etale 态射，见 态射章引理
\ref{morphisms-lemma-base-change-etale}。例如由 概形章引理
\ref{schemes-lemma-points-fibre-product} 对纤维积之点的描述，$U$ 中有一点同时映到
$u_1$ 和 $u_2$。对第 (2) 部分，将 $U''$ 定义为纤维积
$$
\xymatrix{
U'' \ar[r] \ar[d] & U \ar[d]^{(h_1, h_2)} \\
U' \ar[r]^-\Delta & U' \times_S U'.
}
$$
。由于 $h_1$ 和 $h_2$ 诱导相同的剩余域映射
$\kappa(u') \to \kappa(u)$，存在位于 $u'$ 上方的点 $u'' \in U''$，满足
$\kappa(u'') = \kappa(u')$。特别地，$U'' \not = \emptyset$。此外，由于
$U'$ 在 $S$ 上 \'etale，纤维积 $U'\times_S U'$ 也如此（见 态射章引理
\ref{morphisms-lemma-base-change-etale} 和
\ref{morphisms-lemma-composition-etale}）。因此由 态射章引理
\ref{morphisms-lemma-etale-permanence}，竖直箭头 $(h_1, h_2)$ 为 \'etale 态射。
故由基变换，$U''$ 在 $U'$ 上 \'etale，进而也在 $S$ 上 \'etale（因为 \'etale
态射的复合仍为 \'etale 态射）。因此 $(U'', u'')$ 满足要求。
\end{proof}

\begin{lemma}
\label{more-morphisms-lemma-elementary-etale-neighbourhoods}
设 $S$ 是概形，并设 $s$ 是 $S$ 的一点。$(S, s)$ 的初等 \'etale 邻域范畴是余滤过的
（见 范畴章定义 \ref{categories-definition-codirected}）。
\end{lemma}

\begin{proof}
这立即由定义和引理 \ref{more-morphisms-lemma-etale-neighbourhoods-not-quite-filtered} 得出。
\end{proof}

\begin{lemma}
\label{more-morphisms-lemma-describe-henselization}
设 $S$ 是概形，并设 $s \in S$。则
$$
\mathcal{O}_{S, s}^h =
\colim_{(U, u)} \mathcal{O}(U)
$$
，其中余极限取遍 $(S, s)$ 的初等 \'etale 邻域 $(U, u)$ 的范畴之反范畴；后者为滤过范畴。
\end{lemma}

\begin{proof}
设 $\Spec(A) \subset S$ 是 $s$ 的仿射邻域。令 $\mathfrak p \subset A$ 为对应于
$s$ 的素理想。作此选取后，有典范同构
$\mathcal{O}_{S, s} = A_{\mathfrak p}$ 和 $\kappa(s) = \kappa(\mathfrak p)$。
由满足如下条件的初等 \'etale 邻域 $(U, u)$ 给出一个余终系统：$U$ 仿射且
$U \to S$ 经 $\Spec(A)$ 分解。换言之，右端等于
$\colim_{(B, \mathfrak q)} B$，其中余极限取遍配备一个位于 $\mathfrak p$
上方且满足 $\kappa(\mathfrak p) = \kappa(\mathfrak q)$ 的素理想 $\mathfrak q$
的 \'etale $A$-代数 $B$。因此本引理由 代数章引理
\ref{algebra-lemma-henselization-different} 得出。
\end{proof}

\noindent
可以把纤维上各点的 \'etale 邻域提升到全空间。

\begin{lemma}
\label{more-morphisms-lemma-lift-etale-neighbourhood-fibre}
\begin{slogan}
把纤维上一点的 \'etale 邻域提升到全空间。
\end{slogan}
设 $X \to S$ 是概形态射。设 $x \in X$ 的像为 $s \in S$。设
$(V, v) \to (X_s, x)$ 是 \'etale 邻域。则存在 \'etale 邻域
$(U, u) \to (X, x)$，并且存在 $(X_s, x)$ 的 \'etale 邻域之间的态射
$(U_s, u) \to (V, v)$，它是开浸入。
\end{lemma}

\begin{proof}
可假设 $X$、$V$ 和 $S$ 仿射。设态射 $X \to S$ 由 $A \to B$ 给出，点
$x$ 由素理想 $\mathfrak q \subset B$ 给出，点 $s$ 由
$\mathfrak p = A \cap \mathfrak q$ 给出，而态射 $V \to X_s$ 由
$B \otimes_A \kappa(\mathfrak p) \to C$ 给出。由于 $\kappa(\mathfrak p)$
是 $A/\mathfrak p$ 的一个局部化，存在 $f \in A$、$f \not \in \mathfrak p$
以及 \'etale 环同态 $B \otimes_A (A/\mathfrak p)_f \to D$，使得
$$
C = (B \otimes_A \kappa(\mathfrak p))
\otimes_{B \otimes_A (A/\mathfrak p)_f} D
$$
，见 代数章引理 \ref{algebra-lemma-etale} 的第 (9) 部分。将 $A$ 换成
$A_f$、将 $B$ 换成 $B_f$ 后，可假设 $D$ 在
$B \otimes_A A/\mathfrak p = B/\mathfrak p B$ 上 \'etale。于是可以应用
代数章引理 \ref{algebra-lemma-lift-etale}。本引理得证。
\end{proof}














\section{\'Etale 邻域与分支}
\label{more-morphisms-section-etale-nbhds-branches}

\noindent
概形在一点处的（几何）分支数已在 Properties, Section
\ref{properties-section-unibranch} 中定义。在 Varieties, Section
\ref{varieties-section-number-of-branches} 中，我们把它与正规化态射的纤维联系起来。
本节讨论用 \'etale 邻域刻画这一数目。

\begin{lemma}
\label{more-morphisms-lemma-nr-minimal-primes}
设 $R = \colim R_i$ 是一个环的有向系统的余极限，且其过渡同态忠实平坦。
则将 $R$ 的极小素理想数视为 $\{0, 1, 2, \ldots, \infty\}$ 中的元素时，
它等于各 $R_i$ 的极小素理想数的上确界。
\end{lemma}

\begin{proof}
若 $A \to B$ 是平坦环同态，则由下降定理，$\Spec(B) \to \Spec(A)$
把极小素理想映到极小素理想（代数章引理
\ref{algebra-lemma-flat-going-down}）。若 $A \to B$ 忠实平坦，则每个极小素理想都是某个极小素理想的像
（由 代数章引理 \ref{algebra-lemma-ff-rings} 和
\ref{algebra-lemma-minimal-prime-image-minimal-prime}）。因此，若 $i \leq i'$，则
% P05-ERR-CH37-0066: Frozen source reverses the monotonicity inequality for minimal-prime counts under a faithfully flat transition R_i -> R_i'.
$R_i$ 的极小素理想数 $\geq$ $R_{i'}$ 的极小素理想数。由 代数章引理
\ref{algebra-lemma-colimit-faithfully-flat}，每个同态 $R_i \to R$ 都忠实平坦，且还可见
$R$ 的极小素理想数 $\geq$ $R_i$ 的极小素理想数。最后，设
$\mathfrak q_1, \ldots, \mathfrak q_n$ 是 $R$ 中两两不同的极小素理想。则可以找到
$i$，使得 $R_i \cap \mathfrak q_1, \ldots, R_i \cap \mathfrak q_n$
两两不同（作为集合，因而也作为素理想）。这蕴含本引理。
\end{proof}

\begin{lemma}
\label{more-morphisms-lemma-nr-branches}
设 $X$ 是概形，$x \in X$ 是一点。则
\begin{enumerate}
\item $X$ 在 $x$ 处的分支数，等于当初等 \'etale 邻域
$(U, u) \to (X, x)$ 变化时，$U$ 中经过 $u$ 的不可约分支数的上确界；
\item $X$ 在 $x$ 处的几何分支数，等于当 \'etale 邻域
$(U, u) \to (X, x)$ 变化时，$U$ 中经过 $u$ 的不可约分支数的上确界；
\item $X$ 在 $x$ 处单分支，当且仅当对每个初等 \'etale 邻域
$(U, u) \to (X, x)$，$U$ 中恰有一个经过 $u$ 的不可约分支；并且
\item $X$ 在 $x$ 处几何单分支，当且仅当对每个 \'etale 邻域
$(U, u) \to (X, x)$，$U$ 中恰有一个经过 $u$ 的不可约分支。
\end{enumerate}
\end{lemma}

\begin{proof}
第 (3)、(4) 部分由第 (1)、(2) 部分和 性质章引理
\ref{properties-lemma-number-of-branches-1} 得出。

\medskip\noindent
证明 (1)。设 $\Spec(A)$ 是 $x$ 的仿射开邻域，并令 $\mathfrak p \subset A$
为对应于 $x$ 的素理想。可将 $X$ 换成 $\Spec(A)$；由于仿射初等 \'etale 邻域
$(U, u)$ 构成余终子系统，在取上确界时只需考虑它们。回忆 Hensel 化
$A_\mathfrak p^h$ 是各环 $B_\mathfrak q$ 在二元组 $(B, \mathfrak q)$
范畴上的余极限，其中 $B$ 是 \'etale $A$-代数，而 $\mathfrak q$ 是位于
$\mathfrak p$ 上方且满足 $\kappa(\mathfrak q) = \kappa(\mathfrak p)$ 的素理想，见
代数章引理 \ref{algebra-lemma-henselization-different}。由环与仿射概形之间的对应，
这些二元组 $(B, \mathfrak q)$ 恰好对应于仿射初等 \'etale 邻域 $(U, u)$。
注意，$\Spec(B)$ 中经过 $\mathfrak q$ 的不可约分支，恰为
$B_\mathfrak q$ 的极小素理想。由 性质章定义
\ref{properties-definition-number-of-branches}，$A_\mathfrak p^h$ 的极小素理想数就是
$X$ 在 $x$ 处的分支数。注意，系统中的过渡同态
$B_\mathfrak q \to B'_{\mathfrak q'}$ 全都平坦。由于平坦局部环同态忠实平坦
（代数章引理 \ref{algebra-lemma-local-flat-ff}），本引理由引理
\ref{more-morphisms-lemma-nr-minimal-primes} 得出。

\medskip\noindent
证明 (2)。除使用 代数章引理
\ref{algebra-lemma-strict-henselization-different} 外，证明与 (1) 相同。有一个细微差别：
给定 $\kappa(x)$ 的可分代数闭包 $\kappa^{sep}$，对每个 \'etale 邻域
$(U, u)$，都可以选取 $\kappa(x)$-嵌入
$\phi : \kappa(u) \to \kappa^{sep}$，因为 $\kappa(u)/\kappa(x)$ 有限可分
（态射章引理 \ref{morphisms-lemma-etale-at-point}）。因此可在所有三元组
$(U, u, \phi)$ 上取上确界，其中 $(U, u) \to (X, x)$ 是仿射 \'etale 邻域，
而 $\phi : \kappa(u) \to \kappa^{sep}$ 是 $\kappa(x)$-嵌入。这些三元组恰好对应于
代数章引理 \ref{algebra-lemma-strict-henselization-different} 中的三元组，
其余证明完全相同。
\end{proof}

\noindent
我们还需要上述引理的相对版本。

\begin{lemma}
\label{more-morphisms-lemma-nr-branches-fibre}
设 $X \to S$ 是概形态射，$x \in X$ 是像为 $s$ 的点。则
\begin{enumerate}
\item 纤维 $X_s$ 在 $x$ 处的分支数，等于当初等 \'etale 邻域
$(U, u) \to (X, x)$ 变化时，纤维 $U_s$ 中经过 $u$ 的不可约分支数的上确界；
\item 纤维 $X_s$ 在 $x$ 处的几何分支数，等于当 \'etale 邻域
$(U, u) \to (X, x)$ 变化时，纤维 $U_s$ 中经过 $u$ 的不可约分支数的上确界；
\item 纤维 $X_s$ 在 $x$ 处单分支，当且仅当对每个初等 \'etale 邻域
$(U, u) \to (X, x)$，纤维 $U_s$ 中恰有一个经过 $u$ 的不可约分支；并且
% P05-ERR-CH37-0067: Frozen source switches from the fibre X_s to X in item (4), unlike the relative statement and the three preceding items.
\item $X$ 在 $x$ 处几何单分支，当且仅当对每个 \'etale 邻域
$(U, u) \to (X, x)$，$U_s$ 中恰有一个经过 $u$ 的不可约分支。
\end{enumerate}
\end{lemma}

\begin{proof}
合并引理 \ref{more-morphisms-lemma-nr-branches} 与
\ref{more-morphisms-lemma-lift-etale-neighbourhood-fibre}。
\end{proof}

\begin{lemma}
\label{more-morphisms-lemma-number-of-branches-and-smooth}
设 $X \to S$ 是概形的光滑态射。设 $x \in X$ 的像为 $s \in S$。则
\begin{enumerate}
\item $X$ 在 $x$ 处的几何分支数等于 $S$ 在 $s$ 处的几何分支数。
\item 若 $\kappa(x)/\kappa(s)$ 是纯不可分\footnote{事实上，只需
$\kappa(x)$ 在 $\kappa(s)$ 上几何不可约。若日后需要，我们会补充详细证明。}
域扩张，则 $X$ 在 $x$ 处的分支数等于 $S$ 在 $s$ 处的分支数。
\end{enumerate}
\end{lemma}

\begin{proof}
这立即由 代数进阶章引理
\ref{more-algebra-lemma-invariance-number-branches-smooth} 和定义得出。
\end{proof}




\section{非分歧态射与 \'etale 态射}
\label{more-morphisms-section-unramified-is-etale}

\noindent
有时非分歧态射自动为 \'etale 态射。

\begin{lemma}
\label{more-morphisms-lemma-unramified-dominant-unibranch-is-etale}
设 $f : X \to Y$ 是概形态射。设 $x \in X$ 的像为 $y \in Y$。假设
\begin{enumerate}
\item $Y$ 是整概形，且在 $y$ 处几何单分支；
\item $f$ 局部有限型；
\item 存在特殊化 $x' \leadsto x$，使得 $f(x')$ 是 $Y$ 的泛点；
\item $f$ 在 $x$ 处非分歧。
\end{enumerate}
则 $f$ 在 $x$ 处为 \'etale 态射。
\end{lemma}

\begin{proof}
可将 $X$ 和 $Y$ 换成 $x$ 和 $y$ 的适当仿射开邻域。于是 $Y$ 是整环 $A$
的谱，而 $X$ 是有限型 $A$-代数 $B$ 的谱。令 $\mathfrak q \subset B$
为对应于 $x$ 的素理想，$\mathfrak p \subset A$ 为对应于 $y$ 的素理想。局部环
$A_\mathfrak p = \mathcal{O}_{Y, y}$ 几何单分支。环同态 $A \to B$ 在
$\mathfrak q$ 处非分歧。此外，(3) 中的点 $x'$ 对应于满足
$A \cap \mathfrak q' = (0)$ 的素理想 $\mathfrak q' \subset \mathfrak q$。
因此 $A_\mathfrak p \to B_\mathfrak q$ 为单射。结论由 代数进阶章引理
\ref{more-algebra-lemma-local-unramified-extension-unibranch-domain-is-etale} 得出。
\end{proof}

\begin{lemma}
\label{more-morphisms-lemma-global-unramified-dominant-unibranch-is-etale}
\begin{reference}
\cite[Expose I, Corollary 9.11]{SGA1}
\end{reference}
设 $f : X \to Y$ 是概形态射。假设
\begin{enumerate}
\item $Y$ 是整的且几何单分支；
\item $X$ 至少有一个不可约分支支配 $Y$；
\item $f$ 非分歧；并且
\item $X$ 连通。
\end{enumerate}
则 $f$ 为 \'etale 态射，且 $X$ 不可约。
\end{lemma}

\begin{proof}
令 $X' \subset X$ 为支配 $Y$ 的不可约分支。这意味着 $X'$ 的泛点映到 $Y$
的泛点（例如见 态射章引理
\ref{morphisms-lemma-dominant-finite-number-irreducible-components}）。由引理
\ref{more-morphisms-lemma-unramified-dominant-unibranch-is-etale}，$f$ 在 $X'$ 的每一点处为
\'etale 态射。特别地，$f$ 为 \'etale 的开子概形 $U \subset X$ 包含 $X'$。
注意，$U$ 的每个拟紧开集只有有限多个不可约分支，见 下降章引理
\ref{descent-lemma-locally-finite-nr-irred-local-fppf}。另一方面，由于 $Y$
几何单分支且 $U$ 在 $Y$ 上 \'etale，概形 $U$ 几何单分支。特别地，$U$
的每一点至多经过一个不可约分支。一个简单的拓扑论证表明，$X' \subset U$
既开又闭。于是 $X'$ 在 $X$ 中当然也既开又闭；由连通性得到
$X = U = X'$，正是所需。
\end{proof}

\begin{lemma}
\label{more-morphisms-lemma-weird-permanence-etale}
设 $f : X \to Y$ 和 $g : Y \to Z$ 是概形态射。设 $x \in X$ 的像为
$y \in Y$。假设
\begin{enumerate}
\item $Y$ 是整概形，且在 $y$ 处几何单分支；
\item $g \circ f$ 在 $x$ 处为 \'etale 态射；
\item 存在特殊化 $x' \leadsto x$，使得 $f(x')$ 是 $Y$ 的泛点。
\end{enumerate}
则 $f$ 在 $x$ 处为 \'etale 态射，且 $g$ 在 $y$ 处为 \'etale 态射。
\end{lemma}

\begin{proof}
将 $X$ 换成 $x$ 的一个开邻域后，可假设 $g \circ f$ 为 \'etale 态射。于是由
态射章引理 \ref{morphisms-lemma-unramified-permanence} 和
\ref{morphisms-lemma-etale-smooth-unramified}，$f$ 非分歧。因此由引理
\ref{more-morphisms-lemma-unramified-dominant-unibranch-is-etale}，$f$ 在 $x$ 处为 \'etale 态射。
再由 \'etale 下降，$g$ 在 $y$ 处为 \'etale 态射，例如见 下降章引理
\ref{descent-lemma-syntomic-smooth-etale-permanence}。
\end{proof}

\begin{lemma}
\label{more-morphisms-lemma-global-weird-permanence-etale}
设 $f : X \to Y$ 和 $g : Y \to Z$ 是概形态射。假设
\begin{enumerate}
\item $Y$ 是整的且几何单分支；
\item $g \circ f$ 为 \'etale 态射；
\item $X$ 的每个不可约分支都支配 $Y$。
\end{enumerate}
则 $f$ 为 \'etale 态射，且 $g$ 在 $f$ 的像中的每一点处为 \'etale 态射。
\end{lemma}

\begin{proof}
这立即由逐点版本引理 \ref{more-morphisms-lemma-weird-permanence-etale} 得出。
\end{proof}










\section{光滑态射的切片}
\label{more-morphisms-section-etale-over-smooth}

\noindent
本节阐释一个大意如下的结果：概形 $S$ 的光滑覆盖可由 \'etale 覆盖加细。其证明技术依赖于切片论证。

\begin{lemma}
\label{more-morphisms-lemma-slice-smooth-given-element}
设 $f : X \to S$ 是概形态射。设 $x \in X$ 的像为 $s \in S$。设
$h \in \mathfrak m_x \subset \mathcal{O}_{X, x}$。假设
\begin{enumerate}
\item $f$ 在 $x$ 处光滑；并且
\item $\text{d}h$ 在
$$
\Omega_{X_s/s, x} \otimes_{\mathcal{O}_{X_s, x}} \kappa(x) =
\Omega_{X/S, x} \otimes_{\mathcal{O}_{X, x}} \kappa(x)
$$
中的像 $\text{d}\overline{h}$ 非零。
\end{enumerate}
则存在 $x$ 的仿射开邻域 $U \subset X$，使得 $h$ 来自
$h \in \Gamma(U, \mathcal{O}_U)$，并且 $D = V(h)$ 是 $U$ 中的有效 Cartier
除子，满足 $x \in D$，且 $D \to S$ 光滑。
\end{lemma}

\begin{proof}
由于 $f$ 在 $x$ 处光滑，将 $X$ 替换为 $x$ 的一个开邻域后，可假设 $f$
光滑。特别地，$f$ 平坦且局部有限表现。由
引理 \ref{more-morphisms-lemma-slice-given-element}
已知存在 $x$ 的开邻域 $U \subset X$，使得 $h$ 来自
$h \in \Gamma(U, \mathcal{O}_U)$，并且 $D = V(h)$ 是 $U$ 中满足
$x \in D$ 的有效 Cartier 除子，而 $D \to S$ 平坦且有限表现。由
态射章引理 \ref{morphisms-lemma-differentials-relative-immersion}
有短正合列
$$
\mathcal{C}_{D/U} \to i^*\Omega_{U/S} \to \Omega_{D/S} \to 0
$$
其中 $i : D \to U$ 是闭浸入，而 $\mathcal{C}_{D/U}$ 是 $D$ 在 $U$ 中的余法层。
由于 $D$ 是由 $h \in \Gamma(U, \mathcal{O}_U)$ 截出的有效 Cartier 除子，可知
% P05-ERR-CH37-0068: Frozen source identifies the conormal sheaf with h · O_S, although the sheaf is on D and the usual invertible conormal is I/I^2.
$\mathcal{C}_{D/U} = h \cdot \mathcal{O}_S$。由于 $U \to S$ 光滑，层
$\Omega_{U/S}$ 有限局部自由，故其拉回 $i^*\Omega_{U/S}$ 也有限局部自由。
该序列的第一个箭头把自由生成元 $h$ 映到 $i^*\Omega_{U/S}$ 的截面
$\text{d}h|_D$；按假设，它在纤维
$\Omega_{U/S, x} \otimes \kappa(x)$ 中的值非零。由 $\otimes \kappa(x)$ 的右正合性可得
$$
\dim_{\kappa(x)} \left( \Omega_{D/S, x} \otimes \kappa(x) \right)
=
\dim_{\kappa(x)} \left( \Omega_{U/S, x} \otimes \kappa(x) \right) - 1.
$$
由
态射章引理 \ref{morphisms-lemma-smooth-at-point}
可知，$\Omega_{U/S, x} \otimes \kappa(x)$ 至多由 $\dim_x(U_s)$ 个元素生成。
由上式，$\Omega_{D/S, x} \otimes \kappa(x)$ 至多由
$\dim_x(U_s) - 1$ 个元素生成。注意到
$\dim_x(D_s) = \dim_x(U_s) - 1$；例如，由
代数章引理 \ref{algebra-lemma-one-equation}
有 $\dim(\mathcal{O}_{D_s, x}) = \dim(\mathcal{O}_{U_s, x}) - 1$
（$D_s \subset U_s$ 也是有效 Cartier 除子，见
除子章引理 \ref{divisors-lemma-relative-Cartier})
，再使用
态射章引理 \ref{morphisms-lemma-dimension-fibre-at-a-point}.
因此，$\Omega_{D/S, x} \otimes \kappa(x)$ 至多由 $\dim_x(D_s)$ 个元素生成，
再由
态射章引理 \ref{morphisms-lemma-smooth-at-point}
可知 $D \to S$ 在 $x$ 处光滑。缩小 $U$ 后，$D \to S$ 光滑，结论得证。
\end{proof}

\begin{lemma}
\label{more-morphisms-lemma-slice-smooth-once}
设 $f : X \to S$ 是概形态射。设 $x \in X$ 的像为 $s \in S$。假设
\begin{enumerate}
\item $f$ 在 $x$ 处光滑；并且
\item 映射
$$
\Omega_{X_s/s, x} \otimes_{\mathcal{O}_{X_s, x}} \kappa(x)
\longrightarrow
\Omega_{\kappa(x)/\kappa(s)}
$$
的核非零。
\end{enumerate}
则存在 $x$ 的仿射开邻域 $U \subset X$，以及包含 $x$ 的有效 Cartier 除子
$D \subset U$，使得 $D \to S$ 光滑。
\end{lemma}

\begin{proof}
记 $k = \kappa(s)$ 且 $R = \mathcal{O}_{X_s, x}$。以 $\mathfrak m$ 表示 $R$
的极大理想，并令 $\kappa = R/\mathfrak m$，于是 $\kappa = \kappa(x)$。
微分模的形成与局部化交换（见
代数章引理 \ref{algebra-lemma-differentials-localize})
，故 $\Omega_{X_s/s, x} = \Omega_{R/k}$。由
代数章引理 \ref{algebra-lemma-differential-seq}
有正合列
$$
\mathfrak m/\mathfrak m^2 \xrightarrow{\text{d}}
\Omega_{R/k} \otimes_R \kappa \to
\Omega_{\kappa/k} \to 0.
$$
因此，若 (2) 成立，则存在元素 $\overline{h} \in \mathfrak m$，使得
$\text{d}\overline{h}$ 非零。选取 $\overline{h}$ 的提升
$h \in \mathcal{O}_{X, x}$，并应用
引理 \ref{more-morphisms-lemma-slice-smooth-given-element}.
\end{proof}

\begin{remark}
\label{more-morphisms-remark-necessary-condition-slice-smooth}
即使 $x$ 是正维纤维的闭点，
引理 \ref{more-morphisms-lemma-slice-smooth-once}
中的第二个条件也是必要的。下面给出一个例子：设 $k$ 是特征为 $p > 0$
的非完美域。取不是 $p$ 次幂的元素 $a \in k$。令
$\mathfrak m = (x, y^p - a) \subset k[x, y]$。它对应于
$X = \mathbf{A}^2_k$ 的闭点 $w$。令 $S = \mathbf{A}^1_k$，并令
$f : X \to S$ 为对应于 $k[x] \to k[x, y]$ 的态射。此时不存在交换图
$$
\xymatrix{
S' \ar[rr]_h \ar[rd]_g & & X \ar[ld]^f \\
& S
}
$$
使得 $g$ 为 \'etale 且 $w$ 位于 $h$ 的像中。这是显然的，因为剩余域扩张
$\kappa(w)/\kappa(f(w))$ 是纯不可分的，而对任何满足
$g(s') = f(w)$ 的 $s' \in S'$，扩张 $\kappa(s')/\kappa(f(w))$ 都将是可分的。
\end{remark}

\noindent
若假设剩余域扩张可分，则
说明 \ref{more-morphisms-remark-necessary-condition-slice-smooth}
中的现象不会发生。准确结果如下。

\begin{lemma}
\label{more-morphisms-lemma-slice-smooth-once-separable-residue-field-extension}
设 $f : X \to S$ 是概形态射。设 $x \in X$ 的像为 $s \in S$。假设
\begin{enumerate}
\item $f$ 在 $x$ 处光滑；
\item 剩余域扩张 $\kappa(x)/\kappa(s)$ 可分；并且
\item $x$ 不是 $X_s$ 的泛点。
\end{enumerate}
则存在 $x$ 的仿射开邻域 $U \subset X$，以及包含 $x$ 的有效 Cartier 除子
$D \subset U$，使得 $D \to S$ 光滑。
\end{lemma}

\begin{proof}
记 $k = \kappa(s)$ 且 $R = \mathcal{O}_{X_s, x}$。以 $\mathfrak m$ 表示 $R$
的极大理想，并令 $\kappa = R/\mathfrak m$，于是 $\kappa = \kappa(x)$。
微分模的形成与局部化交换（见
代数章引理 \ref{algebra-lemma-differentials-localize})
，故 $\Omega_{X_s/s, x} = \Omega_{R/k}$。由假设 (2) 及
代数章引理 \ref{algebra-lemma-computation-differential}
，映射
$$
\text{d} :
\mathfrak m/\mathfrak m^2
\longrightarrow
\Omega_{R/k} \otimes_R \kappa(\mathfrak m)
$$
是单射。假设 (3) 蕴含 $\mathfrak m/\mathfrak m^2 \not = 0$。
因此存在元素 $\overline{h} \in \mathfrak m$，使得
$\text{d}\overline{h}$ 非零。选取 $\overline{h}$ 的提升
$h \in \mathcal{O}_{X, x}$，并应用
引理 \ref{more-morphisms-lemma-slice-smooth-given-element}.
\end{proof}

\noindent
下一引理构造的子概形 $Z$ 实际上是 $x$ 的某个仿射开邻域中的完全交。若以后需要这一事实，
我们会另行明确表述一个引理。

\begin{lemma}
\label{more-morphisms-lemma-slice-smooth}
设 $f : X \to S$ 是概形态射。设 $x \in X$ 的像为 $s \in S$。假设
\begin{enumerate}
\item $f$ 在 $x$ 处光滑；并且
\item $x$ 是 $X_s$ 的闭点，且 $\kappa(s) \subset \kappa(x)$ 可分。
\end{enumerate}
则存在包含 $x$ 的浸入 $Z \to X$，使得
\begin{enumerate}
\item $Z \to S$ 为 \'etale；并且
\item 集合论意义下 $Z_s = \{x\}$。
\end{enumerate}
\end{lemma}

\begin{proof}
我们可以并且确实把 $S$ 替换为 $s$ 的一个仿射开邻域。我们也可以并且确实把
$X$ 替换为 $x$ 的一个仿射开邻域，使得 $X \to S$ 光滑。下文对仿射概形间的
光滑态射按 $d = \dim_x(X_s)$ 归纳证明此引理。

\medskip\noindent
情形 $d = 0$。此时证明可取 $Z$ 为 $x$ 的一个开邻域。事实上，若 $d = 0$，则
$X \to S$ 在 $x$ 处拟有限，见
态射章引理 \ref{morphisms-lemma-locally-quasi-finite-rel-dimension-0}.
因此存在仿射开邻域 $U \subset X$，使得 $U \to S$ 拟有限，见
态射章引理 \ref{morphisms-lemma-quasi-finite-points-open}.
于是把 $X$ 替换为 $U$ 后，$X$ 在 $S$ 上拟有限且光滑，因而在 $S$ 上相对维数为
$0$ 且光滑，故在 $S$ 上为 \'etale。此外，纤维 $X_s$ 是有限离散集。因此，
% P05-ERR-CH37-0069: Frozen source says an affine open neighbourhood of X; the intended distinguished point is x.
把 $X$ 替换为 $X$ 的另一个仿射开邻域后，可使 $f^{-1}(\{s\}) = \{x\}$
（因为 $X_s$ 的拓扑由 $X$ 的拓扑诱导，见
概形章引理 \ref{schemes-lemma-fibre-topological}).
这就证明了此情形。

\medskip\noindent
接着假设 $d > 0$。注意到，由于 $x$ 是其纤维的闭点，扩张
$\kappa(x)/\kappa(s)$ 有限（由 Hilbert 零点定理，见
态射章引理 \ref{morphisms-lemma-closed-point-fibre-locally-finite-type}).
由于代数可分域扩张具有此性质，故 $\Omega_{\kappa(x)/\kappa(s)} = 0$。因此可以应用
引理 \ref{more-morphisms-lemma-slice-smooth-once}
得到图
$$
\xymatrix{
D \ar[r] \ar[rrd] & U \ar[r] \ar[rd] & X \ar[d] \\
& & S
}
$$
其中 $x \in D$。注意到 $\dim_x(D_s) = \dim_x(X_s) - 1$；例如，由
代数章引理 \ref{algebra-lemma-one-equation}
有 $\dim(\mathcal{O}_{D_s, x}) = \dim(\mathcal{O}_{X_s, x}) - 1$
（$D_s \subset X_s$ 也是有效 Cartier 除子，见
除子章引理 \ref{divisors-lemma-relative-Cartier})
，再使用
态射章引理 \ref{morphisms-lemma-dimension-fibre-at-a-point}.
因此态射 $D \to S$ 光滑，且
$\dim_x(D_s) = \dim_x(X_s) - 1 = d - 1$。由归纳假设，可找到所需的浸入
$Z \to D$，证明完成。
\end{proof}

\begin{lemma}
\label{more-morphisms-lemma-etale-nbhd-dominates-smooth}
\begin{slogan}
光滑态射具有 \'etale 局部截面。
\end{slogan}
\begin{reference}
见 \cite[Corollaire 17.16.3 (ii)]{EGA4}。
\end{reference}
设 $f : X \to S$ 是概形的光滑态射。设 $s \in S$ 是 $f$ 像中的点。
则存在 \'etale 邻域 $(S', s') \to (S, s)$ 以及 $S$-态射 $S' \to X$。
\end{lemma}

\begin{proof}[引理 \ref{more-morphisms-lemma-etale-nbhd-dominates-smooth} 的第一种证明]
按假设 $X_s \not = \emptyset$。由
簇章引理 \ref{varieties-lemma-smooth-separable-closed-points-dense}
，存在闭点 $x \in X_s$，使得 $\kappa(x)$ 是 $\kappa(s)$ 的有限可分域扩张。因此由
引理 \ref{more-morphisms-lemma-slice-smooth}
，存在浸入 $Z \to X$，使得 $Z \to S$ 为 \'etale 且 $x \in Z$。取
$(S' , s') = (Z, x)$。
\end{proof}

\begin{proof}[引理 \ref{more-morphisms-lemma-etale-nbhd-dominates-smooth} 的第二种证明]
选取满足 $f(x) = s$ 的点 $x \in X$。选择图
$$
\xymatrix{
X \ar[d] & U \ar[l] \ar[d] \ar[r]_-\pi & \mathbf{A}^d_V \ar[ld] \\
S & V \ar[l]
}
$$
使得 $\pi$ 为 \'etale，$x \in U$，且 $V = \Spec(R)$ 仿射，见
态射章引理 \ref{morphisms-lemma-smooth-etale-over-affine-space}.
特别地，$s \in V$。态射 $\pi : U \to \mathbf{A}^d_V$ 是开映射，见
态射章引理 \ref{morphisms-lemma-etale-open}.
因此 $W = \pi(U) \cap \mathbf{A}^d_s$ 是 $\mathbf{A}^d_s$ 的非空开子集。
取点 $w \in W$，使得 $\kappa(s) \subset \kappa(w)$ 有限可分，见
簇章引理 \ref{varieties-lemma-affine-space-over-field}.
由
代数章引理 \ref{algebra-lemma-dim-affine-space}
，存在 $d$ 个元素
$\overline{f}_1, \ldots, \overline{f}_d \in \kappa(s)[x_1, \ldots, x_d]$，
它们生成 $\kappa(s)[x_1, \ldots, x_d]$ 中对应于 $w$ 的极大理想。
把 $R$ 替换为一个主局部化后，可假设存在
$f_1, \ldots, f_d \in R[x_1, \ldots, x_d]$，其像为
$\overline{f}_1, \ldots, \overline{f}_d \in \kappa(s)[x_1, \ldots, x_d]$。
考虑 $R$-代数
$$
R' = R[x_1, \ldots, x_d]/(f_1, \ldots, f_d)
$$
并令 $S' = \Spec(R')$。按构造，有 $V$ 上的闭浸入
$j : S' \to \mathbf{A}^d_V$。按构造，$S' \to V$ 在 $s$ 上的纤维是单点
$s'$，其剩余域是 $\kappa(s)$ 的有限可分扩张。设
$\mathfrak q' \subset R'$ 为相应素理想。由
代数章引理 \ref{algebra-lemma-localize-relative-complete-intersection}
可知，对某个 $g \in R'$，$(R')_g$ 是 $R$ 上的相对全局完全交，且
% P05-ERR-CH37-0070: Frozen source uses undefined q here after defining q'.
$g \not \in \mathfrak q$。因此 $S' \to V$ 在 $s'$ 的一个邻域上平坦且有限表现，见
代数章引理 \ref{algebra-lemma-relative-global-complete-intersection}.
按构造，$S' \to V$ 在 $s$ 上的概形论纤维是 $\Spec(\kappa(s'))$。故由
态射章引理 \ref{morphisms-lemma-etale-at-point}
可知，$S' \to S$ 在 $s'$ 处为 \'etale。令
$$
S'' = U \times_{\pi, \mathbf{A}^d_V, j} S'.
$$
按构造，存在点 $s'' \in S''$，它经投影 $p : S'' \to S'$ 映到 $s'$。
注意到 $p$ 是 \'etale 态射 $\pi$ 的基变换，故为 \'etale，见
态射章引理 \ref{morphisms-lemma-base-change-etale}.
选取 $s''$ 的一个小仿射邻域 $S''' \subset S''$，使它映入 $s' \in S'$
的非空开邻域，而态射 $S' \to S$ 在该邻域上为 \'etale。于是 \'etale 邻域
$(S''', s'') \to (S, s)$ 即为本引理所求。
\end{proof}

\noindent
下一引理表明，光滑拓扑上的层与 \'etale 拓扑上的层是同一回事。

\begin{lemma}
\label{more-morphisms-lemma-etale-dominates-smooth}
设 $S$ 是概形。设 $\mathcal{U} = \{S_i \to S\}_{i \in I}$ 是 $S$ 的光滑覆盖，见
拓扑章定义 \ref{topologies-definition-smooth-covering}.
则存在 \'etale 覆盖 $\mathcal{V} = \{T_j \to S\}_{j \in J}$（见
拓扑章定义 \ref{topologies-definition-etale-covering})
，它加细 $\mathcal{U}$（见
位点章定义 \ref{sites-definition-morphism-coverings})
。
\end{lemma}

\begin{proof}
对每个 $s \in S$，存在 $i \in I$，使得 $s$ 位于 $S_i \to S$ 的像中。由
引理 \ref{more-morphisms-lemma-etale-nbhd-dominates-smooth}
可找到 \'etale 态射 $g_s : T_s \to S$，使得 $s \in g_s(T_s)$，且 $g_s$
经 $S_i \to S$ 分解。因此 $\{T_s \to S\}$ 是加细 $\mathcal{U}$ 的 $S$ 的
\'etale 覆盖。
\end{proof}

\begin{lemma}
\label{more-morphisms-lemma-cover-smooth-by-special}
设 $f : X \to S$ 是概形的光滑态射。则存在 \'etale 覆盖
$\{U_i \to X\}_{i \in I}$，使得 $U_i \to S$ 分解为
$U_i \to V_i \to S$，其中 $V_i \to S$ 为 \'etale，而 $U_i \to V_i$
是仿射概形间的光滑态射，具有截面且纤维几何连通。
\end{lemma}

\begin{proof}
设 $s \in S$。由
簇章引理 \ref{varieties-lemma-smooth-separable-closed-points-dense}
，使得 $\kappa(x)/\kappa(s)$ 可分的闭点 $x \in X_s$ 在 $X_s$ 中稠密。
因此，只需构造像包含 $x$ 的 \'etale 态射 $U \to X$，使得 $U \to S$
按引理所述方式分解。为此，选取经过 $x$ 的浸入 $Z \to X$，使得
$Z \to S$ 为 \'etale（引理 \ref{more-morphisms-lemma-slice-smooth}）。把 $S$ 替换为 $Z$，
把 $X$ 替换为 $Z \times_S X$ 后，可假设 $X \to S$ 具有截面
$\sigma : S \to X$，且 $\sigma(s) = x$。接着先把 $S$ 替换为 $s$ 的仿射开邻域，
再把 $X$ 替换为 $x$ 的仿射开邻域。最后，考虑 Section
\ref{more-morphisms-section-connected-components} 的子集 $X^0 \subset X$。由 Lemmas
\ref{more-morphisms-lemma-connected-along-section-open} 和
\ref{more-morphisms-lemma-connected-along-section-locally-constructible}
可知，这是一个包含 $\sigma$ 的逆紧开子概形，并且 $X^0 \to S$ 的纤维几何连通。
若 $X^0$ 不仿射，则选取包含 $x$ 的仿射开集 $U \subset X^0$。由于
$X^0 \to S$ 光滑，$U$ 的像是开集。选取 $s$ 的仿射开邻域 $V \subset S$，
使它包含于 $\sigma^{-1}(U)$ 及 $U \to S$ 的像中。最后，读者可看出
$U \cap f^{-1}(V) \to V$ 具有全部所需性质。例如，
$U \cap f^{-1}(V)$ 等于 $U \times_S V$，因而作为仿射概形的纤维积而仿射。此外，
% P05-ERR-CH37-0071: Frozen source says the fibres of X^0 -> S are irreducible although only geometric connectedness was established.
$U \cap f^{-1}(V) \to V$ 的几何纤维是 $X^0 \to S$ 的不可约纤维的非空开子概形，
故为连通的。略去若干细节。
\end{proof}











\section{\'Etale 邻域与 Artin 逼近}
\label{more-morphisms-section-etale-nbhds-artin}

\noindent
本节证明如下形式的结果：若两个带点概形具有同构的完备局部环，则它们具有同构的
\'etale 邻域。我们将依赖 Popescu 定理，见
环同态光滑化章定理 \ref{smoothing-theorem-popescu}.

\begin{lemma}
\label{more-morphisms-lemma-map-approximation}
设 $S$ 是局部 Noether 概形。设 $X$、$Y$ 是 $S$ 上局部有限型的概形。设
$x \in X$ 与 $y \in Y$ 位于同一点 $s \in S$ 之上。假设
$\mathcal{O}_{S, s}$ 是 G-环。再假设给定局部 $\mathcal{O}_{S, s}$-代数映射
$$
\varphi : \mathcal{O}_{Y, y} \longrightarrow \mathcal{O}_{X, x}^\wedge
$$
。对每个 $N \geq 1$，存在初等 \'etale 邻域
$(U, u) \to (X, x)$ 以及把 $u$ 映到 $y$ 的 $S$-态射 $f : U \to Y$，使得图
$$
\xymatrix{
\mathcal{O}_{X, x}^\wedge \ar[r] &
\mathcal{O}_{U, u}^\wedge \\
\mathcal{O}_{Y, y} \ar[r]^{f^\sharp_u} \ar[u]^\varphi &
\mathcal{O}_{U, u} \ar[u]
}
$$
模 $\mathfrak m_u^N$ 交换。
\end{lemma}

\begin{proof}
该问题在 $X$ 上是局部的，故可假设 $X$、$Y$、$S$ 仿射。记
$S = \Spec(R)$、$X = \Spec(A)$、$Y = \Spec(B)$。写成
$B = R[x_1, \ldots, x_n]/(f_1, \ldots, f_m)$。设
$\mathfrak p \subset A$ 是对应于 $x$ 的素理想。由
代数进阶章命题
\ref{more-algebra-proposition-finite-type-over-G-ring}，局部环
$\mathcal{O}_{X, x} = A_\mathfrak p$ 是 G-环。映射 $\varphi$ 是映射
$$
B_\mathfrak q^\wedge \longrightarrow A_\mathfrak p^\wedge
$$
，其中 $\mathfrak q \subset B$ 是对应于 $y$ 的素理想。令
$a_1, \ldots, a_n \in A_\mathfrak p^\wedge$ 为 $x_1, \ldots, x_n$ 经
$R[x_1, \ldots, x_n] \to B \to B_\mathfrak q^\wedge \to A_\mathfrak p^\wedge$
所得的像。于是可以应用 环同态光滑化章引理
\ref{smoothing-lemma-approximation-property-variant}
，得到 \'etale 环映射 $A \to A'$、素理想 $\mathfrak p' \subset A'$ 以及
$b_1, \ldots, b_n \in A'$，使得
$\kappa(\mathfrak p) = \kappa(\mathfrak p')$,
$a_i - b_i \in (\mathfrak p')^N(A'_{\mathfrak p'})^\wedge$, 及
% P05-ERR-CH37-0072: Frozen source indexes f_j by j = 1,...,n although the presentation defines f_1,...,f_m.
$f_j(b_1, \ldots, b_n) = 0$ for $j = 1, \ldots, n$.
把 $x_i$ 的类映到 $b_i \in A'$，便确定一个 $R$-代数映射 $B \to A'$。
取 $U = \Spec(A') \to \Spec(B)$ 为态射 $f$，并取 $u = \mathfrak p'$，证明完成。
\end{proof}

\begin{lemma}
\label{more-morphisms-lemma-isomorphism-approximation}
设 $S$ 是局部 Noether 概形。设 $X$、$Y$ 是 $S$ 上局部有限型的概形。设
$x \in X$ 与 $y \in Y$ 位于同一点 $s \in S$ 之上。假设
$\mathcal{O}_{S, s}$ 是 G-环。假设有完备局部环之间的
$\mathcal{O}_{S, s}$-代数同构
$$
\varphi : \mathcal{O}_{Y, y}^\wedge \longrightarrow \mathcal{O}_{X, x}^\wedge
$$
。则对每个 $N \geq 1$，存在 $S$ 上带点概形的态射
$$
(X, x) \leftarrow (U, u) \rightarrow (Y, y)
$$
，使得两个箭头都定义初等 \'etale 邻域，且图
$$
\xymatrix{
& \mathcal{O}_{U, u}^\wedge \\
\mathcal{O}_{Y, y}^\wedge \ar[rr]^\varphi \ar[ru] & &
\mathcal{O}_{X, x}^\wedge \ar[lu]
}
$$
模 $\mathfrak m_u^N$ 交换。
\end{lemma}

\begin{proof}
可假设 $N \geq 2$。应用 引理 \ref{more-morphisms-lemma-map-approximation} 得到
$(U, u) \to (X, x)$ 与 $f : (U, u) \to (Y, y)$。我们断言 $f$ 在 $u$
处为 \'etale，这将完成证明。事实上，将证明诱导映射
$\mathcal{O}_{Y, y}^\wedge \to \mathcal{O}_{U, u}^\wedge$
是同构。证明这一点后，
引理 \ref{more-morphisms-lemma-lifting-along-artinian-at-point}
将表明 $f$ 在 $u$ 处光滑；当然 $f$ 在 $u$ 处也非分歧，故
态射章引理 \ref{morphisms-lemma-etale-smooth-unramified}
表明 $f$ 在 $u$ 处为 \'etale。对局部环 $(R, \mathfrak m)$，令
$\text{Gr}_\mathfrak m(R) =
\bigoplus_{n \geq 0} \mathfrak m^n/\mathfrak m^{n + 1}$.
为证明该断言，考察诱导的分次环图
$$
\xymatrix{
& \text{Gr}_{\mathfrak m_u}(\mathcal{O}_{U, u}) \\
\text{Gr}_{\mathfrak m_y}(\mathcal{O}_{Y, y}) \ar[rr]^\varphi \ar[ru] & &
\text{Gr}_{\mathfrak m_x}(\mathcal{O}_{X, x}) \ar[lu]
}
$$
由于 $N \geq 2$，且图中的分次代数由次数 $1$ 的部分生成，所以此图确实交换！
按假设，下方箭头是同构。由 代数进阶章引理
\ref{more-algebra-lemma-flat-unramified}（例如），映射
$\mathcal{O}_{X, x}^\wedge \to \mathcal{O}_{U, u}^\wedge$
是同构，因而图中西北方向的箭头是同构。于是 $f$ 诱导同构
% P05-ERR-CH37-0073: Frozen source gives f an X-based graded source and an O_U target graded by m_y; f has source Y and the target maximal ideal is m_u.
$\text{Gr}_{\mathfrak m_x}(\mathcal{O}_{X, x}) \to
\text{Gr}_{\mathfrak m_y}(\mathcal{O}_{U, u})$.
对两个局部环使用归纳法及短正合列
$$
0 \to \text{Gr}^n_{\mathfrak m}(R) \to R/\mathfrak m^{n + 1} \to
R/\mathfrak m^n \to 0
$$
，由蛇引理可得 $f$ 对所有 $n$ 均诱导同构
$\mathcal{O}_{Y, y}/\mathfrak m_y^n \to \mathcal{O}_{U, u}/\mathfrak m_u^n$
，这正是所要证明的。
\end{proof}

\begin{lemma}
\label{more-morphisms-lemma-relative-map-approximation-pre}
设 $X \to S$、$Y \to T$、$x$、$s$、$y$、$t$、$\sigma$、$y_\sigma$ 与
$\varphi$ 如下给定：有概形态射
$$
\vcenter{
\xymatrix{
X \ar[d] & Y \ar[d] \\
S & T
}
}
\quad\text{带标记点}\quad
\vcenter{
\xymatrix{
x \ar[d] & y \ar[d] \\
s & t
}
}
$$
其中 $S$ 局部 Noether，$T$ 在 $\mathbf{Z}$ 上有限型。态射 $X \to S$ 与
$Y \to T$ 局部有限型。局部环 $\mathcal{O}_{S, s}$ 是 G-环。映射
$$
\sigma : \mathcal{O}_{T, t} \longrightarrow \mathcal{O}_{S, s}^\wedge
$$
是局部同态。令
$Y_\sigma = Y \times_{T, \sigma} \Spec(\mathcal{O}_{S, s}^\wedge)$.
再设 $y_\sigma$ 是 $Y_\sigma$ 中的点，它映到 $y$ 以及
$\Spec(\mathcal{O}_{S, s}^\wedge)$ 的闭点。最后，
$$
\varphi :
\mathcal{O}_{X, x}^\wedge
\longrightarrow
\mathcal{O}_{Y_\sigma, y_\sigma}^\wedge
$$
是 $\mathcal{O}_{S, s}^\wedge$-代数的同构。此时存在交换图
$$
\xymatrix{
X \ar[d] &
W \ar[l] \ar[rd] \ar[rr] & &
Y \times_{T, \tau} V \ar[r] \ar[ld] & Y \ar[d] \\
S & &
V \ar[ll] \ar[rr]^\tau & &
T
}
$$
以及点 $w \in W$、$v \in V$，使得
\begin{enumerate}
\item $(V, v) \to (S, s)$ 是初等 \'etale 邻域；
\item $(W, w) \to (X, x)$ 是初等 \'etale 邻域；并且
\item $\tau(v) = t$。
\end{enumerate}
令 $y_\tau \in Y \times_T V$ 为在恒等式
$(Y_\sigma)_s = (Y \times_T V)_v$ 下对应于 $y_\sigma$ 的点。则
\begin{enumerate}
\item[(4)] $(W, w) \to (Y \times_{T, \tau} V, y_\tau)$ 是初等 \'etale 邻域。
\end{enumerate}
\end{lemma}

\begin{proof}
记 $X_\sigma = X \times_S \Spec(\mathcal{O}_{S, s}^\wedge)$，并以
$x_\sigma \in X_\sigma$ 表示位于 $x$ 上方的唯一点。由
代数进阶章命题
\ref{more-algebra-proposition-Noetherian-complete-G-ring}，
$\mathcal{O}_{S, s}^\wedge$ 是 G-环。由 Lemma
\ref{more-morphisms-lemma-isomorphism-approximation}，可以选择
$$
(X_\sigma, x_\sigma) \leftarrow (U, u) \rightarrow (Y_\sigma, y_\sigma)
$$
，其中两个箭头都是初等 \'etale 邻域。

\medskip\noindent
把 $S$ 替换为 $s$ 的一个开邻域后，可假设 $S = \Spec(R)$ 仿射。由于
$\mathcal{O}_{S, s}$ 是 G-环，由 环同态光滑化章定理
\ref{smoothing-theorem-popescu}，环 $\mathcal{O}_{S, s}^\wedge$ 是光滑
$R$-代数的滤过余极限。因此可写成
$$
\Spec(\mathcal{O}_{S, s}^\wedge) = \lim S_i
$$
，这是在 $S$ 上光滑的仿射概形 $S_i$ 的有向极限。以 $s_i \in S_i$ 表示
$\Spec(\mathcal{O}_{S, s}^\wedge)$ 的闭点的像。注意
$\kappa(s) = \kappa(s_i)$。令 $X_i = X \times_S S_i$，并以
$x_i \in X_i$ 表示映到 $x$ 的唯一点。注意 $\kappa(x) = \kappa(x_i)$。
由于 $T$ 在 $\mathbf{Z}$ 上有限型，由 极限章命题
\ref{limits-proposition-characterize-locally-finite-presentation}
，可选择一个 $i$ 以及带点概形的态射
$\sigma_i : (S_i, s_i) \to (T, t)$，使它与
$\Spec(\mathcal{O}_{S, s}^\wedge) \to S_i$ 的复合等于 $\sigma$。令
$Y_i = Y \times_T S_i$，并以 $y_i$ 表示 $y_\sigma$ 的像。注意
$\kappa(y_i) = \kappa(y_\sigma)$。由 极限章引理
\ref{limits-lemma-descend-finite-presentation}，可选择一个 $i$ 以及图
$$
\xymatrix{
X_i \ar[rd] &
U_i \ar[l] \ar[d] \ar[r] &
Y_i \ar[ld] \\
& S_i
}
$$
，它到 $\Spec(\mathcal{O}_{S, s}^\wedge)$ 的基变换恢复
$X_\sigma \leftarrow U \rightarrow Y_\sigma$。由 极限章引理
\ref{limits-lemma-descend-etale}，增大 $i$ 后可假设态射
$X_i \leftarrow U_i \rightarrow Y_i$ 为 \'etale。设 $u_i \in U_i$ 是 $u$ 的像。
由于 $u_i \mapsto x_i$，故
$\kappa(x) = \kappa(x_\sigma) = \kappa(u) \supset \kappa(u_i) \supset
\kappa(x_i) = \kappa(x)$，从而 $\kappa(u_i) = \kappa(x_i)$。因此
$(X_i, x_i) \leftarrow (U_i, u_i)$ 是初等 \'etale 邻域。又由于
$\kappa(y_i) = \kappa(y_\sigma) = \kappa(u)$，可知
$(U_i, u_i) \to (Y_i, y_i)$ 也是初等 \'etale 邻域。

\medskip\noindent
至此已构造图
$$
\xymatrix{
X \ar[d] &
X \times_S S_i \ar[l] \ar[rd] &
U_i \ar[l] \ar[r] \ar[d] &
Y \times_T S_i \ar[r] \ar[ld] &
Y \ar[d] \\
S & &
S_i \ar[ll] \ar[rr] & &
T
}
$$
，它与引理陈述中的图相同，但 $S_i \to S$ 是光滑的。由 Lemma
\ref{more-morphisms-lemma-slice-smooth}，并缩小 $S_i$ 后，可假设存在经过 $s_i$ 的闭子概形
$V \subset S_i$，使得 $V \to S$ 为 \'etale。令 $W$ 等于 $V$ 在 $U_i$
中的概形论逆像，即得结论。
\end{proof}

\noindent
我们强烈建议读者跳过本节余下部分。

\begin{lemma}
\label{more-morphisms-lemma-relative-map-approximation}
考虑图
$$
\vcenter{
\xymatrix{
X \ar[d] & Y \ar[d] \\
S & T \ar[l]
}
}
\quad\text{带标记点}\quad
\vcenter{
\xymatrix{
x \ar[d] & y \ar[d] \\
s & t \ar[l]
}
}
$$
% P05-ERR-CH37-0074: Frozen source has the ungrammatical phrase “where S be”.
，其中 $S$ 是局部 Noether 概形，且各态射局部有限型。假设
$\mathcal{O}_{S, s}$ 是 G-环。再假设给定局部 $\mathcal{O}_{S, s}$-代数映射
$$
\sigma : \mathcal{O}_{T, t} \longrightarrow \mathcal{O}_{S, s}^\wedge
$$
以及局部 $\mathcal{O}_{S, s}$-代数映射
$$
\varphi :
\mathcal{O}_{X, x}
\longrightarrow
\mathcal{O}_{Y_\sigma, y_\sigma}^\wedge
$$
，其中 $Y_\sigma = Y \times_{T, \sigma} \Spec(\mathcal{O}_{S, s}^\wedge)$，
而 $y_\sigma$ 是 $Y_\sigma$ 中位于 $y$ 上方的唯一点。对每个 $N \geq 1$，
存在 $S$ 上概形的交换图
$$
\xymatrix{
X \ar[d] & X \times_S V \ar[l] \ar[rd] &
W \ar[l]^-f \ar[r] \ar[d] &
Y \times_{T, \tau} V \ar[r] \ar[ld] & Y \ar[d] \\
S & & V \ar[ll] \ar[rr]^\tau & & T
}
$$
以及点 $w \in W$、$v \in V$，使得
\begin{enumerate}
\item $v \mapsto s$、$\tau(v) = t$、$f(w) = (x, v)$，且 $w \mapsto (y, v)$；
\item $(V, v) \to (S, s)$ 是初等 \'etale 邻域；
\item 图
$$
\xymatrix{
\mathcal{O}_{S, s}^\wedge \ar[r] & \mathcal{O}_{V, v}^\wedge \\
\mathcal{O}_{T, t} \ar[r]^{\tau^\sharp_v} \ar[u]_\sigma &
\mathcal{O}_{V, v} \ar[u]
}
$$
% P05-ERR-CH37-0075: Frozen source says “commutes module” rather than “commutes modulo”.
模 $\mathfrak m_v^N$ 交换；
\item $(W, w) \to (Y \times_{T, \tau} V, (y, v))$ 是初等 \'etale 邻域；
\item 图
$$
\xymatrix{
\mathcal{O}_{X, x} \ar[r]_\varphi &
\mathcal{O}_{Y_\sigma, y_\sigma}^\wedge \ar[r] &
\mathcal{O}_{Y_\sigma, y_\sigma}/\mathfrak m_{y_\sigma}^N \ar@{=}[r] &
\mathcal{O}_{Y \times_{T, \tau} V, (y, v)}/\mathfrak m_{(y, v)}^N
\ar[d]_{\cong} \\
\mathcal{O}_{X, x} \ar[r] \ar@{=}[u] &
\mathcal{O}_{X \times_S V, (x, v)} \ar[r]^{f^\sharp_w} &
\mathcal{O}_{W, w} \ar[r] &
\mathcal{O}_{W, w}/\mathfrak m_w^N
}
$$
交换。该等号源于以下事实：由 (2) 与 (3)，$Y_\sigma$ 与
$Y \times_{T, \tau} V$ 在
$\mathcal{O}_{V, v}/\mathfrak m_v^N = \mathcal{O}_{S, s}/\mathfrak m_s^N$
上典范同构。
\end{enumerate}
\end{lemma}

\begin{proof}
把 $X$、$S$、$T$、$Y$ 替换为仿射开子概形后，可假设引理陈述中的图来自对下图
施行 $\Spec$：
$$
\vcenter{
\xymatrix{
A & B \\
R \ar[u] \ar[r] & C \ar[u]
}
}
\quad\text{with primes}\quad
\vcenter{
\xymatrix{
\mathfrak p_A & \mathfrak p_B \\
\mathfrak p_R \ar@{-}[u] \ar@{-}[r] & \mathfrak p_C \ar@{-}[u]
}
}
$$
，其中各环均 Noether，各环映射均有限型。在本证明中，每个环 $E$ 都是 Noether
$R$-代数，并带有位于 $\mathfrak p_R$ 上方的素理想 $\mathfrak p_E$；所有环映射
都是与给定素理想相容的 $R$-代数映射。此外，写 $E^\wedge$ 时，是指 $E$ 在
$\mathfrak p_E$ 处局部化后的完备化。我们还将不再另作说明地使用如下事实：若
\'etale 环映射 $E_1 \to E_2$ 满足
$\kappa(\mathfrak p_{E_1}) = \kappa(\mathfrak p_{E_2})$，则由
代数进阶章引理 \ref{more-algebra-lemma-flat-unramified}，它诱导同构
$E_1^\wedge = E_2^\wedge$。

\medskip\noindent
在此记号下，$\sigma$ 与 $\varphi$ 对应于环映射
$$
\sigma : C \to R^\wedge
\quad\text{且}\quad
\varphi : A \longrightarrow (B \otimes_{C, \sigma} R^\wedge)^\wedge
$$
。图示如下：
$$
\xymatrix{
A \ar@/^1em/[rrr]^\varphi &
B \ar[r] &
B \otimes_{C, \sigma} R^\wedge \ar[r] &
(B \otimes_{C, \sigma} R^\wedge)^\wedge \\
R \ar[r] \ar[u] &
C \ar[r]^\sigma \ar[u] & R^\wedge \ar[u] \ar[ru]
}
$$
由
代数进阶章命题
\ref{more-algebra-proposition-Noetherian-complete-G-ring}，$R^\wedge$ 是 G-环。因此由
代数进阶章命题
\ref{more-algebra-proposition-finite-type-over-G-ring}，
$B \otimes_{C, \sigma} R^\wedge$ 是 G-环。由 Lemma
\ref{more-morphisms-lemma-map-approximation}（转写成代数形式），存在 \'etale 环映射
$B \otimes_{C, \sigma} R^\wedge \to B'$，它诱导同构
$\kappa(\mathfrak p_{B \otimes_{C, \sigma} R^\wedge})
\to \kappa(\mathfrak p_{B'})$
；还存在 $R$-代数映射 $A \to B'$，使得复合
$$
A \to B' \to (B')^\wedge = (B \otimes_{C, \sigma} R^\wedge)^\wedge
$$
模 $(\mathfrak p_{(B \otimes_{C, \sigma} R^\wedge)^\wedge})^N$ 与
$\varphi$ 相同。因此可以用这一复合替换 $\varphi$，因为 $\varphi$ 在结论中
只通过引理陈述第 (5) 项的交换性要求出现。图示为：
$$
\xymatrix{
& & B' \ar[r] & (B')^\wedge \ar@{=}[d] \\
A \ar[rru] &
B \ar[r] &
B \otimes_{C, \sigma} R^\wedge \ar[r] \ar[u] &
(B \otimes_{C, \sigma} R^\wedge)^\wedge \\
R \ar[r] \ar[u] &
C \ar[r]^\sigma \ar[u] & R^\wedge \ar[u] \ar[ru]
}
$$
接着使用如下事实：由于按假设 $R_{\mathfrak p_R}$ 是 G-环，$R^\wedge$ 是光滑
$R$-代数的滤过余极限（环同态光滑化章定理
\ref{smoothing-theorem-popescu}）。由于 $C$ 在 $R$ 上有限表现，得到分解
$$
C \to R' \to R^\wedge
$$
，其中某个 $R \to R'$ 光滑，见
代数章引理 \ref{algebra-lemma-characterize-finite-presentation}.
增大 $R'$ 后，可假设存在 \'etale 的 $B \otimes_C R'$-代数 $B''$，其到
$B \otimes_{C, \sigma} R^\wedge$ 的基变换是 $B'$，见
代数章引理 \ref{algebra-lemma-etale}.
于是 $B'$ 是这些 $B''$ 的滤过余极限，从而增大 $R'$ 后，可假设存在
$R$-代数映射 $A \to B''$，使得 $A \to B'' \to B'$ 是先前构造的映射
（仍用上述引理）。图示为
$$
\xymatrix{
& & B'' \ar[r] & B' \ar[r] & (B')^\wedge \ar@{=}[d] \\
A \ar[rru] &
B \ar[r] &
B \otimes_C R' \ar[r] \ar[u] &
B \otimes_{C, \sigma} R^\wedge \ar[r] \ar[u] &
(B \otimes_{C, \sigma} R^\wedge)^\wedge \\
R \ar[r] \ar[u] &
C \ar[r] \ar[u] &
R' \ar[r] \ar[u] &
R^\wedge \ar[u] \ar[ru]
}
$$
并且
$$
B' = B'' \otimes_{(B \otimes_C R')} (B \otimes_{C, \sigma} R^\wedge)
$$
。这意味着可用 $R'$ 替换 $C$，用 $R' \to R^\wedge$ 替换
$\sigma : C \to R^\wedge$，并用 $B''$ 替换 $B$，从而化简为图
$$
\xymatrix{
A \ar[r] &
B \ar[r] &
B \otimes_{C, \sigma} R^\wedge \\
R \ar[r] \ar[u] &
C \ar[r]^\sigma \ar[u] & R^\wedge \ar[u]
}
$$
，其中 $\varphi$ 等于各水平箭头的复合再接上从
$B \otimes_{C, \sigma} R^\wedge$ 到其完备化的典范映射。
证明的最后一步，是把 引理 \ref{more-morphisms-lemma-map-approximation}（或其证明）再应用一次于
$\Spec(R)$ 上的 $\Spec(C)$ 和 $\Spec(R)$，以及映射 $C \to R^\wedge$。
该引理给出环映射 $C \to D$，使得 $R \to D$ 为 \'etale，
$\kappa(\mathfrak p_R) = \kappa(\mathfrak p_D)$，并且
$$
C \to D \to D^\wedge = R^\wedge
$$
模 $(\mathfrak p_{R^\wedge})^N$ 等于 $\sigma : C \to R^\wedge$。于是可取
$$
V = \Spec(D)
\quad\text{且}\quad
W = \Spec(B \otimes_C D)
$$
作为引理所求。事实上，图
$$
\xymatrix{
A \ar[r] &
B \otimes_{C, \sigma} R^\wedge \ar[r] &
B \otimes_{C, \sigma} R^\wedge/(\mathfrak p_{R^\wedge})^N \ar@{=}[r] &
B \otimes_C D/(\mathfrak p_D)^N \\
A \ar@{=}[u] \ar[r] &
A \otimes_R D \ar[r] &
B \otimes_R D \ar[r] &
B \otimes_C D/(\mathfrak p_D)^N \ar@{=}[u]
}
$$
交换，因为 $C \to D \to D^\wedge = R^\wedge$ 模
$(\mathfrak p_{R^\wedge})^N$ 等于 $\sigma$。这证明了第 (5) 项，其余性质由构造立即得到。
\end{proof}

\begin{lemma}
\label{more-morphisms-lemma-control-agreement}
% P05-ERR-CH37-0076: Frozen source uses plural “morphisms” for the single map T -> S.
设 $T \to S$ 是 Noether 概形的有限型态射。设 $t \in T$ 映到 $s \in S$，并设
$\sigma : \mathcal{O}_{T, t} \to \mathcal{O}_{S, s}^\wedge$
是局部 $\mathcal{O}_{S, s}$-代数映射。对每个 $N \geq 1$，存在有限型态射
$(T', t') \to (T, t)$，使得 $\sigma$ 经
$\mathcal{O}_{T, t} \to \mathcal{O}_{T', t'}$
分解，并且对每个经 $\mathcal{O}_{T, t} \to \mathcal{O}_{T', t'}$ 分解的局部
$\mathcal{O}_{S, s}$-代数映射
$\sigma' : \mathcal{O}_{T, t} \to \mathcal{O}_{S, s}^\wedge$
，映射 $\sigma$ 与 $\sigma'$ 模 $\mathfrak m_s^N$ 相同。
\end{lemma}

\begin{proof}
可假设 $S$ 与 $T$ 仿射。记 $S = \Spec(R)$ 且 $T = \Spec(C)$。设
$c_1, \ldots, c_n \in C$ 是 $C$ 作为 $R$-代数的一组生成元。设
$\mathfrak p \subset R$ 是对应于 $s$ 的素理想。记
$\mathfrak p = (f_1, \ldots, f_m)$。把 $R$ 替换为一个主局部化（以消去
$R_\mathfrak p$ 中的分母）后，可假设存在
$r_1, \ldots, r_n \in R$ 与 $a_{i, I} \in \mathcal{O}_{S, s}^\wedge$，
其中 $I = (i_1, \ldots, i_m)$ 且 $\sum i_j = N$，使得
$$
\sigma(c_i) = r_i + \sum\nolimits_I a_{i, I} f_1^{i_1} \ldots f_m^{i_m}
$$
在 $\mathcal{O}_{S, s}^\wedge$ 中成立。于是考虑
$$
C' = C[t_{i, I}]/
\left(c_i - r_i - \sum\nolimits_I t_{i, I} f_1^{i_1} \ldots f_m^{i_m}\right)
$$
，取 $\mathfrak p' = \mathfrak pC' + (t_{i, I})$；把 $t_{i, I}$ 映到
$a_{i, I}$，便给出 $\sigma : C \to \mathcal{O}_{S, s}^\wedge$ 经 $C'$ 的分解。
取 $T' = \Spec(C')$ 即可，因为引理陈述中的任何 $\sigma'$ 都把 $c_i$ 模极大理想的
$N$ 次幂映到 $r_i$。
\end{proof}

\begin{lemma}
\label{more-morphisms-lemma-control-graded}
设 $Y \to T \to S$ 是 Noether 概形的有限型态射。设 $t \in T$ 映到
$s \in S$，并设
$\sigma : \mathcal{O}_{T, t} \to \mathcal{O}_{S, s}^\wedge$
是局部 $\mathcal{O}_{S, s}$-代数映射。存在有限型态射
$(T', t') \to (T, t)$，使得 $\sigma$ 经
$\mathcal{O}_{T, t} \to \mathcal{O}_{T', t'}$
分解，并且对每个经 $\mathcal{O}_{T, t} \to \mathcal{O}_{T', t'}$ 分解的局部
$\mathcal{O}_{S, s}$-代数映射
$\sigma' : \mathcal{O}_{T, t} \to \mathcal{O}_{S, s}^\wedge$
，闭浸入
$$
Y \times_{T, \sigma} \Spec(\mathcal{O}_{S, s}^\wedge) = Y_\sigma
\longleftarrow Y_t \longrightarrow
Y_{\sigma'} =
Y \times_{T, \sigma'} \Spec(\mathcal{O}_{S, s}^\wedge)
$$
具有同构的余法代数。
\end{lemma}

\begin{proof}
一个有用的观察是，$\sigma$ 的存在蕴含 $\kappa(s) = \kappa(t)$。陈述确有意义，
因为 $Y_\sigma$ 与 $Y_{\sigma'}$ 在
$s \in \Spec(\mathcal{O}_{S, s}^\wedge)$ 上的纤维都典范同构于 $Y_t$。
我们把“$\sigma'$ 经 $\mathcal{O}_{T, t} \to \mathcal{O}_{T', t'}$ 分解”这一性质
视为对 $\sigma'$ 的约束。若有若干这样的约束，例如由
$(T'_i, t'_i) \to (T, t)$、$i = 1, \ldots, n$ 给出，则可考虑
% P05-ERR-CH37-0077: Frozen source has the ungrammatical phrase “can combined”.
$(T'_1 \times_T \ldots \times_T T'_n, (t'_1, \ldots, t'_n)) \to
(T, t)$ 把它们合并。下文将不再另作说明地使用这一点。

\medskip\noindent
由 引理 \ref{more-morphisms-lemma-control-agreement}，可假设引理陈述中的任何 $\sigma'$ 模
$\mathfrak m_s^2$ 都与 $\sigma$ 相同。注意，$Y_t$ 在 $Y_\sigma$ 中的余法代数
正是拟凝聚分次 $\mathcal{O}_{Y_t}$-代数
$$
\bigoplus\nolimits_{n \geq 0}
\mathfrak m_s^n\mathcal{O}_{Y_\sigma}/
\mathfrak m_s^{n + 1}\mathcal{O}_{Y_\sigma}
$$
；$Y_{\sigma'}$ 的情形类似。由于 $\sigma$ 与 $\sigma'$ 模
$\mathfrak m_s^2$ 相同，这两个代数在次数 $0$ 与 $1$ 上相同。另一方面，这些余法代数
在次数 $0$ 上由次数 $1$ 的部分生成。因此，若存在延拓刚在次数 $0$ 与 $1$ 上构造的
同构的同构，则它是唯一的。

\medskip\noindent
可假设 $S$ 与 $T$ 仿射。设 $Y = Y_1 \cup \ldots \cup Y_n$ 是仿射开覆盖。
若能按引理构造 $(T_i', t'_i) \to (T, t)$，使得对 $Y_i \to T \to S$ 与
$\sigma$ 存在所需同构（见上一段），则由唯一性，这些同构可以黏合，从而证明
$Y \to T$ 的结果。因此可假设 $Y$ 仿射。

\medskip\noindent
写成 $S = \Spec(R)$、$T = \Spec(C)$ 与 $Y = \Spec(B)$。选择表现
$B = C[x_1, \ldots, x_n]/(f_1, \ldots, f_m)$。记
$R^\wedge = \mathcal{O}_{S, s}^\wedge$。取多项式
$a_{kj} \in R^\wedge[x_1, \ldots, x_n]$，使得
$$
\sum\nolimits_{j = 1, \ldots, m} a_{kj}\sigma(f_j) = 0,\quad
\text{对 }k = 1, \ldots, K
$$
给出 $\sigma(f_j) \in R^\wedge[x_1, \ldots, x_n]$ 之间的关系模的一组生成元。
于是有正合列
\begin{equation}
\label{more-morphisms-equation-resolution}
R^\wedge[x_1, \ldots, x_n]^{\oplus K} \to
R^\wedge[x_1, \ldots, x_n]^{\oplus m} \to
R^\wedge[x_1, \ldots, x_n] \to B \otimes_{C, \sigma} R^\wedge \to 0
\end{equation}
按
代数进阶章第 \ref{more-algebra-section-artin-rees} 节.
中的定义，取整数 $c$，使得 Artin--Rees 引理对该序列的第一、第二个映射以及理想
$\mathfrak m_{R^\wedge}R^\wedge[x_1, \ldots, x_n]$ 均适用。
用多重指标记号写成
$$
a_{kj} = \sum\nolimits_{I \in \Omega} a_{kj, I} x^I
\quad\text{且}\quad
f_j = \sum\nolimits_{I \in \Omega} f_{j, I} x^I
$$
，其中 $a_{kj, I} \in R^\wedge$、$f_{j, I} \in C$，而 $\Omega$ 是有限多重指标集。
% P05-ERR-CH37-0079: Frozen source omits “is” in “and Omega a finite set of multiindices”.
于是可知
$$
\sum\nolimits_{j = 1, \ldots, m,\ I, I' \in \Omega,\ I + I' = I''}
a_{kj, I} \sigma(f_{j, I'}) = 0,\quad
I''\text{ a multiindex}
$$
在 $R^\wedge$ 中成立。因此取
$$
% P05-ERR-CH37-0078: Frozen source defines C' with t_{jk,I}, uses t_{kj,I} in the relations, and then maps t_{kj,I} to a_{jk,I}; the k,j indices are inconsistent.
C' = C[t_{jk, I}]/
\left(
\sum\nolimits_{j = 1, \ldots, m,\ I, I' \in \Omega,\ I + I' = I''}
t_{kj, I} f_{j, I'},\ I''\text{ a multiindex}\right)
$$
。于是 $\sigma$ 经映射 $\tilde\sigma : C' \to R^\wedge$ 分解，该映射把
$t_{kj, I}$ 映到 $a_{jk, I}$。因此 $T' = \Spec(C')$ 带有点 $t' \in T'$，
使得 $\sigma$ 经 $\mathcal{O}_{T, t} \to \mathcal{O}_{T', t'}$ 分解。在
$C'[x_1, \ldots, x_n]$ 中令 $t_{kj} = \sum t_{kj, I} x^I$。于是有复形
\begin{equation}
\label{more-morphisms-equation-resolution-new}
C'[x_1, \ldots, x_n]^{\oplus K} \to
C'[x_1, \ldots, x_n]^{\oplus m} \to
C'[x_1, \ldots, x_n] \to B \otimes_C C' \to 0
\end{equation}
，它在 $C'[x_1, \ldots, x_n]$ 处正合，且经 $\tilde\sigma$ 的基变换给出
(\ref{more-morphisms-equation-resolution})。

\medskip\noindent
由 引理 \ref{more-morphisms-lemma-control-agreement}，可找到进一步的态射
$(T'', t'') \to (T', t')$，使得 $\tilde\sigma$ 经
$\mathcal{O}_{T', t'} \to \mathcal{O}_{T'', t''}$ 分解，并且若
$\sigma' : C \to R^\wedge$ 经 $\mathcal{O}_{T'', t''}$ 分解，则诱导映射
$\tilde \sigma' : C' \to R^\wedge$ 模 $\mathfrak m_s^{c + 1}$ 与
$\tilde \sigma$ 相同。因此，若 $\sigma'$ 是这样的映射，就得到复形
$$
R^\wedge[x_1, \ldots, x_n]^{\oplus K} \to
R^\wedge[x_1, \ldots, x_n]^{\oplus m} \to
R^\wedge[x_1, \ldots, x_n] \to B \otimes_{C, \sigma'} R^\wedge \to 0
$$
：这是把 $\tilde\sigma'$ 应用于多项式 $t_{kj}$ 与 $f_j$ 而得到的
$R^\wedge[x_1, \ldots, x_n]$ 上复形。换言之，它是复形
(\ref{more-morphisms-equation-resolution-new}) 经 $\tilde\sigma'$ 的基变换。由于
$\tilde \sigma$ 与 $\tilde \sigma'$ 模 $\mathfrak m_s^{c + 1}$ 同余，定义此复形的
矩阵与定义复形 (\ref{more-morphisms-equation-resolution}) 的矩阵模
$\mathfrak m_s^{c + 1}$ 同余。由于 (\ref{more-morphisms-equation-resolution}) 正合，可以应用
代数进阶章引理 \ref{more-algebra-lemma-approximate-complex-graded}
得出
$$
\text{Gr}_{\mathfrak m_s}(B \otimes_{C, \sigma'} R^\wedge)
\cong
\text{Gr}_{\mathfrak m_s}(B \otimes_{C, \sigma} R^\wedge)
$$
，如所求。
\end{proof}

\begin{lemma}
\label{more-morphisms-lemma-relative-isomorphism-approximation}
% P05-ERR-CH37-0080: Frozen source has “notation an assumptions” instead of “notation and assumptions”.
沿用 引理 \ref{more-morphisms-lemma-relative-map-approximation} 的记号与假设，并假设 $\varphi$
在完备化上诱导同构。则可以选择图，使得 $f$ 为 \'etale。
\end{lemma}

\begin{proof}
可假设 $N \geq 2$，并可按 引理 \ref{more-morphisms-lemma-control-graded} 用 $(T', t')$ 替换
$(T, t)$。由于 $(V, v) \to (S, s)$ 是初等 \'etale 邻域，
$(X \times_S V, (x, v)) \to (X, x)$ 也是如此。因此，由
代数进阶章引理 \ref{more-algebra-lemma-flat-unramified}.
，$\mathcal{O}_{X, x} \to \mathcal{O}_{X \times_S V, (x, v)}$ 在完备化上诱导同构。
我们断言 $\mathcal{O}_{X, x} \to \mathcal{O}_{W, w}$ 在完备化上诱导同构。
证明这一点后，
引理 \ref{more-morphisms-lemma-lifting-along-artinian-at-point}
将表明 $f$ 在 $w$ 处光滑；当然，
% P05-ERR-CH37-0081: Frozen source says f is unramified at u although the distinguished point in this lemma is w.
$f$ 在 $u$ 处也非分歧，故
态射章引理 \ref{morphisms-lemma-etale-smooth-unramified}
表明 $f$ 在 $w$ 处为 \'etale。

\medskip\noindent
首先使用
引理 \ref{more-morphisms-lemma-relative-map-approximation}
第 (5) 项的交换性，看到对 $i \leq N$ 有交换图
$$
\xymatrix{
\text{Gr}^i_{\mathfrak m_x}(\mathcal{O}_{X, x})
\ar[r]_-\varphi &
\text{Gr}^i_{\mathfrak m_{y_\sigma}}(\mathcal{O}_{Y_\sigma, y_\sigma}^\wedge)
\ar@{=}[r] &
\text{Gr}^i_{\mathfrak m_{(y, v)}}(\mathcal{O}_{Y \times_{T, \tau} V, (y, v)})
\ar[d]_{\cong} \\
\text{Gr}^i_{\mathfrak m_x}(\mathcal{O}_{X, x})
\ar[r]^-{\cong} \ar@{=}[u] &
\text{Gr}^i_{\mathfrak m_{(x, v)}}(\mathcal{O}_{X \times_S V, (x, v)})
\ar[r]^{f^\sharp_w} &
\text{Gr}^i_{\mathfrak m_w}(\mathcal{O}_{W, w})
}
$$
这蕴含 $f^\sharp_w$ 在剩余域上给出同构 $\kappa(x) \to \kappa(w)$，在余切空间上
给出同构 $\mathfrak m_x/\mathfrak m_x^2 \to \mathfrak m_w/\mathfrak m_w^2$。
因此 $f^\sharp_w$ 在完备局部环上给出满射
$\mathcal{O}_{X, x}^\wedge \to \mathcal{O}_{W, w}^\wedge$
。

\medskip\noindent
由 引理 \ref{more-morphisms-lemma-control-graded}，有
$\text{Gr}_{\mathfrak m_s}(\mathcal{O}_{(Y \times_{T, \tau} V, (y, v)})$
与
$\text{Gr}_{\mathfrak m_s}(\mathcal{O}_{Y_\sigma, y_\sigma})$.
之间的同构。这可通过在点 $y$ 处取余法层同构的茎得到。由于所涉局部环 Noether，
关于 $\mathfrak m_s$ 取相伴分次与完备化交换，因为关于理想的完备化是 Noether 环上
有限模范畴的正合函子。这给出下列分次环图中的右侧竖直同构：
$$
\xymatrix{
\text{Gr}_{\mathfrak m_s}(\mathcal{O}_{W, w}^\wedge) &
% P05-ERR-CH37-0082: Frozen source omits a closing parenthesis in the point subscript of the upper-right local ring.
\text{Gr}_{\mathfrak m_s}
(\mathcal{O}_{(Y \times_{T, \tau} V, (y, v)}^\wedge) \ar[l] \\
\text{Gr}_{\mathfrak m_s}(\mathcal{O}_{X, x}^\wedge)
\ar[r]^\varphi \ar[u] &
\text{Gr}_{\mathfrak m_s}(\mathcal{O}_{Y_\sigma, y_\sigma}^\wedge)
\ar[u]_{\cong}
}
$$
我们不声称此图交换。由上一段的结果，左侧箭头满射；另外三个箭头是同构。因此左侧
箭头是两个同构 Noether 环之间的满射。故由 代数章引理
\ref{algebra-lemma-surjective-endo-noetherian-ring-is-iso}，它是同构
（也可用 Hilbert 函数直接论证）。特别地，
$\mathcal{O}_{X, x}^\wedge \to \mathcal{O}_{W, w}^\wedge$ 既单射又满射，证明完成。
\end{proof}















\section{有限自由覆盖在局部上支配 \'etale 覆盖}
\label{more-morphisms-section-finite-free-over-etale}

\noindent
本节说明一个大意如下的结果：概形 $S$ 的 \'etale 覆盖可由 $S$ 的有限局部自由覆盖的
Zariski 覆盖加细。

\begin{lemma}
\label{more-morphisms-lemma-dominate-etale-neighbourhood-finite-flat}
设 $S$ 是概形，$s \in S$。设 $f : (U, u) \to (S, s)$ 是 \'etale 邻域。
则存在仿射开邻域 $s \in V \subset S$ 以及满的有限局部自由态射
$\pi : T \to V$，使得对每个 $t \in \pi^{-1}(s)$，都存在开邻域
$t \in W_t \subset T$ 及交换图
$$
\xymatrix{
T \ar[d]^\pi & W_t \ar[l] \ar[rr]_{h_t} \ar[rd] & & U \ar[dl] \\
V \ar[rr] & & S
}
$$
，其中 $h_t(t) = u$。
\end{lemma}

\begin{proof}
该问题在 $S$ 上是局部的，故可把 $S$ 替换为 $s$ 的任意开邻域。也可把 $U$ 替换为
$u$ 的一个开邻域。因此，由 态射章引理
\ref{morphisms-lemma-etale-locally-standard-etale}，可假设 $U \to S$ 是仿射概形间的
标准 \'etale 态射。在此情形，引理（取 $V = S$）由
代数章引理 \ref{algebra-lemma-standard-etale-finite-flat-Zariski}.
得到。
\end{proof}

\begin{lemma}
\label{more-morphisms-lemma-dominate-etale-affine-finite-flat}
设 $f : U \to S$ 是仿射概形间的满 \'etale 态射。则存在满的有限局部自由态射
$\pi : T \to S$ 以及有限开覆盖 $T = T_1 \cup \ldots \cup T_n$，使得每个
$T_i \to S$ 都经 $U \to S$ 分解。图示为：
$$
\xymatrix{
& \coprod T_i  \ar[rd] \ar[ld] & \\
T \ar[rd]^\pi & & U \ar[ld]_f \\
& S &
}
$$
其中西南方向的箭头是 Zariski 覆盖。
\end{lemma}

\begin{proof}
这正是
代数章引理 \ref{algebra-lemma-etale-finite-flat-zariski}.
的改述。
\end{proof}

\begin{remark}
\label{more-morphisms-remark-topologies}
用拓扑的语言，引理 \ref{more-morphisms-lemma-dominate-etale-neighbourhood-finite-flat} 与
\ref{more-morphisms-lemma-dominate-etale-affine-finite-flat} 表示如下结论。设 $S$ 是任意概形，
$\{f_i : U_i \to S\}$ 是 $S$ 的 \'etale 覆盖。存在 Zariski 开覆盖
$S = \bigcup V_j$；对每个 $j$，存在有限局部自由满态射 $W_j \to V_j$；并且对每个
$j$，存在 Zariski 开覆盖 $\{W_{j, k} \to W_j\}$，使得族
$\{W_{j, k} \to S\}$ 加细给定的 \'etale 覆盖
$\{f_i : U_i \to S\}$。这在实践中意味着什么？例如，假设有一个下降问题，我们知道
如何对 Zariski 覆盖以及形如 $\{\pi : T \to S\}$ 的 fppf 覆盖解决它，其中
$\pi$ 有限局部自由且满。那么，这个下降问题对 \'etale 覆盖也有肯定答案。Gabber
在证明仿射概形 $X$ 满足 $\text{Br}(X) = \text{Br}'(X)$ 时使用了这一技巧，见
\cite{Hoobler}。
\end{remark}




\section{拟有限态射的 \'etale 局部化}
\label{more-morphisms-section-etale-localization}

\noindent
现在给出围绕“拟有限态射经 \'etale 局部化后变为有限态射”这一主题的一系列引理。
一般思路如下。设给定概形态射 $f : X \to S$ 以及点 $s \in S$。设
$\varphi : (U, u) \to (S, s)$ 是 $S$ 中 $s$ 的 \'etale 邻域。考虑纤维积
$X_U = U \times_S X$ 与基本图
\begin{equation}
\label{more-morphisms-equation-basic-diagram}
\vcenter{
\xymatrix{
V \ar[r] \ar[dr] & X_U \ar[d] \ar[r] & X \ar[d]^f \\
& U \ar[r]^\varphi & S
}
}
\end{equation}
，其中 $V \subset X_U$ 是开集。对态射 $f_U : X_U \to U$，或者对适当开集 $V$
的态射 $V \to U$，是否存在某种标准模型？一般而言答案当然是否定的，但对拟有限态射
可以得到一些结论。

\begin{lemma}
\label{more-morphisms-lemma-etale-makes-quasi-finite-finite-at-point}
设 $f : X \to S$ 是概形态射。设 $x \in X$，令 $s = f(x)$。假设
\begin{enumerate}
\item $f$ 局部有限型；并且
\item $x \in X_s$ 是孤立点\footnote{在 (1) 成立时，这意味着 $f$ 在 $x$ 处拟有限，
见 态射章引理 \ref{morphisms-lemma-quasi-finite-at-point-characterize}。}。
\end{enumerate}
则存在
\begin{enumerate}
\item[(a)] 初等 \'etale 邻域 $(U, u) \to (S, s)$；
\item[(b)] 开子概形 $V \subset X_U$（见 \ref{more-morphisms-equation-basic-diagram}）
\end{enumerate}
，使得
\begin{enumerate}
\item[(\romannumeral1)] $V \to U$ 是有限态射；
\item[(\romannumeral2)] $V$ 中恰有一个映到 $U$ 中 $u$ 的点 $v$；并且
\item[(\romannumeral3)] 点 $v$ 在态射 $X_U \to X$ 下映到 $x$，并诱导
$\kappa(x) = \kappa(v)$。
\end{enumerate}
此外，对任意初等 \'etale 邻域 $(U', u') \to (U, u)$，令
$V' = U' \times_U V \subset X_{U'}$，则三元组 $(U', u', V')$ 仍满足性质
(\romannumeral1)、(\romannumeral2) 与 (\romannumeral3)。
\end{lemma}

\begin{proof}
取仿射开集 $Y \subset X$、$W \subset S$，使得 $f(Y) \subset W$ 且
$x \in Y$。注意，$x$ 也是态射 $f|_Y : Y \to W$ 的纤维的孤立点。若能对
$f|_Y : Y \to W$ 证明该定理，则可得到 $f$ 的结论。因此归约到 $f$ 是仿射概形间
态射的情形；这正是
代数章引理 \ref{algebra-lemma-etale-makes-quasi-finite-finite-one-prime}.
。
\end{proof}

\noindent
在上一引理和下一引理中，我们不假设态射 $f$ 分离。这意味着其中构造的开集
$V$、$V_i$ 未必在 $X_U$ 中闭。此外，若选择仿射邻域 $U$，则每个 $V_i$ 都仿射，
但在非分离情形，交 $V_i \cap V_j$ 未必仿射。

\begin{lemma}
\label{more-morphisms-lemma-etale-makes-quasi-finite-finite-multiple-points}
设 $f : X \to S$ 是概形态射。设 $x_1, \ldots, x_n \in X$ 是在 $S$ 中具有
同一像 $s$ 的点。假设
\begin{enumerate}
\item $f$ 局部有限型；并且
\item 对 $i = 1, \ldots, n$，$x_i \in X_s$ 是孤立点。
\end{enumerate}
则存在
\begin{enumerate}
\item[(a)] 初等 \'etale 邻域 $(U, u) \to (S, s)$；
\item[(b)] 对每个 $i$，开子概形 $V_i \subset X_U$，
\end{enumerate}
使得对每个 $i$ 均有
\begin{enumerate}
\item[(\romannumeral1)] $V_i \to U$ 是有限态射；
\item[(\romannumeral2)] $V_i$ 中恰有一个映到 $U$ 中 $u$ 的点 $v_i$；并且
\item[(\romannumeral3)] 点 $v_i$ 映到 $X$ 中的 $x_i$，且
$\kappa(x_i) = \kappa(v_i)$。
\end{enumerate}
\end{lemma}

\begin{proof}
对 $n$ 使用归纳法。具体而言，假设 $(U, u) \to (S, s)$ 与
$V_i \subset X_U$（$i = 1, \ldots, n - 1$）对 $x_1, \ldots, x_{n - 1}$
适用。由于 $\kappa(s) = \kappa(u)$，纤维 $(X_U)_u = X_s$。因此存在唯一的点
$x'_n \in X_u \subset X_U$ 对应于 $x_n \in X_s$。此外，$x'_n$ 在 $X_u$
中是孤立的。故由 引理 \ref{more-morphisms-lemma-etale-makes-quasi-finite-finite-at-point}，
存在初等 \'etale 邻域 $(U', u') \to (U, u)$ 以及对 $x'_n$、因而对 $x_n$
适用的开集 $V_n \subset X_{U'}$。由
引理 \ref{more-morphisms-lemma-etale-makes-quasi-finite-finite-at-point}
的最后一个断言，对 $i = 1, \ldots, n - 1$，开子概形
$V'_i = U'\times_U V_i$ 对 $x_1, \ldots, x_{n - 1}$ 仍适用。结论得证。
\end{proof}

\noindent
若允许非平凡域扩张 $\kappa(u)/\kappa(s)$，即允许一般的 \'etale 邻域，则可如下分裂这些点。

\begin{lemma}
\label{more-morphisms-lemma-etale-makes-quasi-finite-finite-multiple-points-var}
设 $f : X \to S$ 是概形态射。设 $x_1, \ldots, x_n \in X$ 是在 $S$ 中具有
同一像 $s$ 的点。假设
\begin{enumerate}
\item $f$ 局部有限型；并且
\item 对 $i = 1, \ldots, n$，$x_i \in X_s$ 是孤立点。
\end{enumerate}
则存在
\begin{enumerate}
\item[(a)] \'etale 邻域 $(U, u) \to (S, s)$；
\item[(b)] 对每个 $i$，整数 $m_i$ 及开子概形
$V_{i, j} \subset X_U$，$j = 1, \ldots, m_i$
\end{enumerate}
，使得
\begin{enumerate}
\item[(\romannumeral1)] 每个 $V_{i, j} \to U$ 都是有限态射；
\item[(\romannumeral2)] $V_{i, j}$ 中恰有一个映到 $U$ 中 $u$ 的点 $v_{i, j}$，且
$\kappa(u) \subset \kappa(v_{i, j})$ 是有限纯不可分扩张；
% P05-ERR-CH37-0083: Frozen source orders item (iv) before item (iii).
\item[(\romannumeral4)] 若 $v_{i, j} = v_{i', j'}$，则 $i = i'$ 且 $j = j'$；并且
\item[(\romannumeral3)] 各点 $v_{i, j}$ 映到 $X$ 中的 $x_i$，且
$(X_U)_u$ 中没有其他点映到 $x_i$。
\end{enumerate}
\end{lemma}

\begin{proof}
本证明是
代数章引理 \ref{algebra-lemma-etale-makes-quasi-finite-finite-variant}
证明的概形语言变体。由 态射章引理
\ref{morphisms-lemma-quasi-finite-at-point-characterize}，态射 $f$ 在每个点 $x_i$
处拟有限。因此对每个 $i$，$\kappa(s) \subset \kappa(x_i)$ 有限
（态射章引理 \ref{morphisms-lemma-residue-field-quasi-finite}）。
对每个 $i$，令 $\kappa(s) \subset L_i \subset \kappa(x_i)$ 为满足如下性质的子域：
$L_i/\kappa(s)$ 可分，而 $\kappa(x_i)/L_i$ 纯不可分。选取有限 Galois 扩张
$L/\kappa(s)$，使得对 $i = 1, \ldots, n$ 存在 $\kappa(s)$-嵌入
$L_i \to L$。选择 \'etale 邻域 $(U, u) \to (S, s)$，使得作为
$\kappa(s)$-扩张有 $L \cong \kappa(u)$（Lemma
\ref{more-morphisms-lemma-realize-prescribed-residue-field-extension-etale}）。

\medskip\noindent
设 $y_{i, j}$（$j = 1, \ldots, m_i$）为 $X_U$ 中位于 $x_i \in X$ 与
$u \in U$ 上方的点。由
概形章引理 \ref{schemes-lemma-points-fibre-product}
，这些点 $y_{i, j}$ 恰好对应于环
$\kappa(u) \otimes_{\kappa(s)} \kappa(x_i)$ 中的素理想。这也说明它们为何只有有限个；
事实上 $m_i = [L_i : \kappa(s)]$，但我们不需要这一点。由 $L$ 的选择
（以及初等域论），对每对 $i, j$，$\kappa(u) \subset \kappa(y_{i, j})$
都是有限纯不可分扩张。此外，例如由 态射章引理
\ref{morphisms-lemma-base-change-quasi-finite}，态射 $X_U \to U$ 在所有点
$y_{i, j}$（所有 $i, j$）处拟有限。

\medskip\noindent
把 引理 \ref{more-morphisms-lemma-etale-makes-quasi-finite-finite-multiple-points} 应用于态射
$X_U \to U$、点 $u \in U$ 及各点 $y_{i, j} \in (X_U)_u$。这给出 \'etale 邻域
$(U', u') \to (U, u)$，满足 $\kappa(u) = \kappa(u')$；还给出开集
$V_{i, j} \subset X_{U'}$，满足该引理的性质 (\romannumeral1)、
(\romannumeral2) 与 (\romannumeral3)。我们断言 \'etale 邻域
$(U', u') \to (S, s)$ 及开集 $V_{i, j} \subset X_{U'}$ 即为本引理所求。
验证从略。
\end{proof}

\begin{lemma}
\label{more-morphisms-lemma-etale-splits-off-quasi-finite-part-technical}
设 $f : X \to S$ 是概形态射。设 $s \in S$，并设
$x_1, \ldots, x_n \in X_s$。假设
\begin{enumerate}
\item $f$ 局部有限型；
\item $f$ 分离；并且
\item $x_1, \ldots, x_n$ 是 $X_s$ 中两两不同的孤立点。
\end{enumerate}
则存在初等 \'etale 邻域 $(U, u) \to (S, s)$ 以及分解
$$
U \times_S X = W \amalg V_1 \amalg \ldots \amalg V_n
$$
为开闭子概形，使得态射 $V_i \to U$ 有限，$V_i \to U$ 在 $u$ 上的纤维是单点集
$\{v_i\}$，每个 $v_i$ 映到 $x_i$ 且 $\kappa(x_i) = \kappa(v_i)$，而
$W \to U$ 在 $u$ 上的纤维不含映到任何 $x_i$ 的点。
\end{lemma}

\begin{proof}
按
引理 \ref{more-morphisms-lemma-etale-makes-quasi-finite-finite-multiple-points}.
选择 $(U, u) \to (S, s)$ 与 $V_i \subset X_U$。由于 $X_U \to U$ 分离
(概形章引理 \ref{schemes-lemma-separated-permanence})
，而 $V_i \to U$ 有限、因而固有
(态射章引理 \ref{morphisms-lemma-finite-proper})
，由
态射章引理 \ref{morphisms-lemma-image-proper-scheme-closed}.
可知 $V_i \subset X_U$ 闭。因此 $V_i \cap V_j$ 是 $V_i$ 中不含 $v_i$ 的闭子集。
故 $V_i \cap V_j$ 在 $U$ 中的像是不含 $u$ 的闭集（因为 $V_i \to U$ 固有）。
缩小 $U$ 后，因而可假设
% P05-ERR-CH37-0084: Frozen source says V_i cap V_j is empty for all i,j, which is impossible when i=j; pairwise disjointness requires i != j.
$V_i \cap V_j = \emptyset$ 对所有 $i, j$ 成立。这给出引理中的分解。
\end{proof}

\noindent
下面给出归约到纯不可分域扩张的变体。

\begin{lemma}
\label{more-morphisms-lemma-etale-splits-off-quasi-finite-part-technical-variant}
设 $f : X \to S$ 是概形态射。设 $s \in S$，并设
$x_1, \ldots, x_n \in X_s$。假设
\begin{enumerate}
\item $f$ 局部有限型；
\item $f$ 分离；并且
\item $x_1, \ldots, x_n$ 是 $X_s$ 中两两不同的孤立点。
\end{enumerate}
则存在 \'etale 邻域 $(U, u) \to (S, s)$ 以及分解
$$
U \times_S X =
W \amalg
\ \coprod\nolimits_{i = 1, \ldots, n}
\ \coprod\nolimits_{j = 1, \ldots, m_i}
V_{i, j}
$$
为开闭子概形，使得态射 $V_{i, j} \to U$ 有限，$V_{i, j} \to U$ 在 $u$ 上的纤维
是单点集 $\{v_{i, j}\}$，每个 $v_{i, j}$ 映到 $x_i$，
$\kappa(u) \subset \kappa(v_{i, j})$ 纯不可分，而 $W \to U$ 在 $u$ 上的纤维
不含映到任何 $x_i$ 的点。
\end{lemma}

\begin{proof}
证明与 引理 \ref{more-morphisms-lemma-etale-splits-off-quasi-finite-part-technical}
完全相同，只是使用 Lemma
\ref{more-morphisms-lemma-etale-makes-quasi-finite-finite-multiple-points-var} 代替 Lemma
\ref{more-morphisms-lemma-etale-makes-quasi-finite-finite-multiple-points}。
\end{proof}

\noindent
下面的版本可能稍易理解。

\begin{lemma}
\label{more-morphisms-lemma-etale-splits-off-quasi-finite-part}
设 $f : X \to S$ 是概形态射。设 $s \in S$。假设
\begin{enumerate}
\item $f$ 局部有限型；
\item $f$ 分离；并且
\item $X_s$ 至多有有限多个孤立点。
\end{enumerate}
则存在初等 \'etale 邻域 $(U, u) \to (S, s)$ 以及分解
$$
U \times_S X = W \amalg V
$$
为开闭子概形，使得态射 $V \to U$ 有限，并且态射 $W \to U$ 的纤维 $W_u$ 不含
孤立点。特别地，若 $f^{-1}(s)$ 是有限集，则 $W_u = \emptyset$。
\end{lemma}

\begin{proof}
这由
引理 \ref{more-morphisms-lemma-etale-splits-off-quasi-finite-part-technical}
立即可得：取 $x_1, \ldots, x_n$ 为 $X_s$ 的全部孤立点，并令
$V = \bigcup V_i$。
\end{proof}






\section{整态射的 \'etale 局部化}
\label{more-morphisms-section-etale-localization-integral}

\noindent
下面给出 第 \ref{more-morphisms-section-etale-localization} 节 中结果在整态射情形的若干变体。

\begin{lemma}
\label{more-morphisms-lemma-etale-makes-integral-split}
设 $R \to S$ 是整环映射。设 $\mathfrak p \subset R$ 是素理想。假设
\begin{enumerate}
\item 位于 $\mathfrak p$ 上方的素理想 $\mathfrak q_1, \ldots, \mathfrak q_n$ 只有有限个；并且
\item 对每个 $i$，极大可分子扩张
% P05-ERR-CH37-0085: Frozen source uses undefined q in the extension tower; the surrounding indexed field is kappa(q_i).
$\kappa(\mathfrak q)/\kappa(\mathfrak q_i)_{sep}/\kappa(\mathfrak p)$
(域章引理 \ref{fields-lemma-separable-first})
在 $\kappa(\mathfrak p)$ 上有限。
\end{enumerate}
则存在 \'etale 环映射 $R \to R'$ 以及位于 $\mathfrak p$ 上方的素理想
$\mathfrak p'$，使得
$$
S \otimes_R R' = A_1 \times \ldots \times A_m
$$
，其中整映射 $R' \to A_j$ 恰有一个位于 $\mathfrak p'$ 上方的素理想
$\mathfrak r_j$，且 $\kappa(\mathfrak r_j)/\kappa(\mathfrak p')$ 纯不可分。
\end{lemma}

\begin{proof}[第一种证明]
本证明使用
代数章引理 \ref{algebra-lemma-etale-makes-quasi-finite-finite-variant}.
具体而言，选取该域在 $\kappa(\mathfrak p)$ 上的生成元
$\theta_i \in \kappa(\mathfrak q_i)_{sep}$（域章引理
\ref{fields-lemma-primitive-element}）。纤维环 $S \otimes_R \kappa(\mathfrak p)$
的谱是有限离散空间，其点对应于 $\mathfrak q_1, \ldots, \mathfrak q_n$。由中国剩余定理
(代数章引理 \ref{algebra-lemma-chinese-remainder})
可知 $S \otimes_R \kappa(\mathfrak p) \to \prod \kappa(\mathfrak q_i)$ 满射。
因此，对某个 $g \in R$、$g \not \in \mathfrak p$，用 $R_g$ 替换 $R$ 后，
可假设
$(0, \ldots, 0, \theta_i, 0, \ldots, 0) \in \prod \kappa(\mathfrak q_i)$
是某个 $x_i \in S$ 的像。设 $S' \subset S$ 是由这些 $x_i$ 生成的 $R$-子代数。
由于 $\Spec(S) \to \Spec(S')$ 满射
(代数章引理 \ref{algebra-lemma-integral-overring-surjective})
，可知 $\mathfrak q_i' = S' \cap \mathfrak q_i$ 是 $S'$ 中位于
$\mathfrak p$ 上方的各素理想。由 $x_i$ 的选择可知这些素理想彼此不同，并且
% P05-ERR-CH37-0086: Frozen source has the malformed phrase “distinct that and”.
$\kappa(\mathfrak q'_i)_{sep} = \kappa(\mathfrak q_i)_{sep}$.
特别地，域扩张 $\kappa(\mathfrak q_i)/\kappa(\mathfrak q'_i)$ 纯不可分。
由于 $R \to S'$ 有限，可以应用
代数章引理 \ref{algebra-lemma-etale-makes-quasi-finite-finite-variant}.
% P05-ERR-CH37-0087: Frozen source places a period before “and we get”, breaking the sentence.
从而得到 $R \to R'$、$\mathfrak p'$ 以及分解
$$
S' \otimes_R R' = A'_1 \times \ldots \times A'_m \times B'
$$
，其中整映射 $R' \to A'_j$ 恰有一个位于 $\mathfrak p'$ 上方的素理想
$\mathfrak r'_j$，使得 $\kappa(\mathfrak r'_j)/\kappa(\mathfrak p')$ 纯不可分；
而 $B'$ 没有位于 $\mathfrak p'$ 上方的素理想。由于 $R' \to B'$ 有限
（因为 $R \to S'$ 有限），在某个 $g' \in R'$、$g' \not \in \mathfrak p'$
处局部化 $R'$ 后，可假设 $B' = 0$。经映射
$S' \otimes_R R' \to S \otimes_R R'$，得到 $S$ 的相应分解。
\end{proof}

\begin{proof}[第二种证明]
本证明使用严格 Hensel 化。首先，
% P05-ERR-CH37-0088: Frozen source says “R is strictly henselization”; the intended condition is that R is strictly henselian/a strict henselization.
假设 $R$ 是以 $\mathfrak p$ 为极大理想的严格 Hensel 化。则
$S/\mathfrak p S$ 有有限多个对应于 $\mathfrak q_1, \ldots, \mathfrak q_n$ 的素理想；
它们均为极大理想，且其剩余域在 $\kappa(\mathfrak p)$ 上纯不可分。因此
% P05-ERR-CH37-0089: Frozen source localizes at undefined p_i after naming the relevant primes q_i.
$S/\mathfrak p S$ 等于 $\prod (S/\mathfrak p S)_{\mathfrak p_i}$。由
代数进阶章引理 \ref{more-algebra-lemma-characterize-henselian-pair}，可把这一乘积分解
提升为 $S$ 的如下所述的
% P05-ERR-CH37-0090: Frozen source says “product composition” instead of “product decomposition”.
乘积复合。

\medskip\noindent
在一般情形，设 $R^{sh}$ 是 $R_\mathfrak p$ 的严格 Hensel 化。于是可把第一段的结果
应用于 $R^{sh} \to S \otimes_R R^{sh}$。考虑 $S \otimes_R R^{sh}$ 中对应于
该乘积分解的 $m$ 个两两正交幂等元。由
代数章引理 \ref{algebra-lemma-strict-henselization-different}
，$R^{sh}$ 是 \'etale 环映射 $(R, \mathfrak p) \to (R', \mathfrak p')$ 的滤过余极限，
故这些幂等元按要求下降到某个 $R'$。
\end{proof}







\section{Zariski 主定理}
\label{more-morphisms-section-application-etale-neighbourhoods}

\noindent
本节证明 Grothendieck 重述的 Zariski 主定理。在本文中说“Zariski 主定理”时，通常指
% P05-ERR-CH37-0091: Frozen source says “in this content”; the idiom requires “in this context”.
引理 \ref{more-morphisms-lemma-finite-type-separated},
引理 \ref{more-morphisms-lemma-quasi-finite-separated-quasi-affine} 或
引理 \ref{more-morphisms-lemma-quasi-finite-separated-pass-through-finite} 中的任一个。在多数文献中，
最后一个结果被称为 Zariski 主定理。

\medskip\noindent
我们已经在
代数章定理 \ref{algebra-theorem-main-theorem}
中证明了代数版本，也已经用概形语言重述了该版本，见
态射章定理 \ref{morphisms-theorem-main-theorem}.
本节的版本更为精细；为得到完整结果，我们使用 Section
\ref{more-morphisms-section-etale-localization} 的 \'etale 局部化技术归约到代数情形。

\begin{lemma}
\label{more-morphisms-lemma-finite-type-separated}
设 $f : X \to S$ 是概形态射。假设 $f$ 有限型且分离。设 $S'$ 是 $S$ 在 $X$
中的正规化，见
态射章定义 \ref{morphisms-definition-normalization-X-in-Y}.
图示为：
$$
\xymatrix{
X \ar[rd]_f \ar[rr]_{f'} & & S' \ar[ld]^\nu \\
& S &
}
$$
则存在开子概形 $U' \subset S'$，使得
\begin{enumerate}
\item $(f')^{-1}(U') \to U'$ 是同构；并且
\item $(f')^{-1}(U') \subset X$ 是 $f$ 在其处拟有限的点集。
\end{enumerate}
\end{lemma}

\begin{proof}
由 态射章引理 \ref{morphisms-lemma-quasi-finite-points-open}，$f$ 在其处拟有限的
点组成开子集 $U \subset X$。本引理等价于
\begin{enumerate}
\item[(a)] $U' = f'(U) \subset S'$ 是开集；
\item[(b)] $U = (f')^{-1}(U')$；并且
\item[(c)] $U \to U'$ 是同构。
\end{enumerate}
任取 $x \in U$。我们断言存在开邻域 $f'(x) \in V \subset S'$，使得
$(f')^{-1}V \to V$ 是同构。先证明该断言蕴含本引理。事实上，
$(f')^{-1}V \cong V$ 在 $S$ 上局部有限型（因为它是 $X$ 的开子概形）；而对
$v \in V$，剩余域扩张 $\kappa(v)/\kappa(\nu(v))$ 是代数的（因为
$V \subset S'$ 且 $S'$ 在 $S$ 上整）。因此 $V \to S$ 的纤维离散
（态射章引理
\ref{morphisms-lemma-algebraic-residue-field-extension-closed-point-fibre})
，而 $(f')^{-1}V \to S$ 局部拟有限
(态射章引理 \ref{morphisms-lemma-locally-quasi-finite-fibres}).
这蕴含 $(f')^{-1}V \subset U$ 与 $V \subset U'$。由于 $x$ 任意，(a)、(b)、(c) 成立。

\medskip\noindent
令 $s = f(x)$。设 $(T, t) \to (S, s)$ 是初等 \'etale 邻域。用下标 ${}_T$
表示到 $T$ 的基变换。设 $y = (x, t) \in X_T$ 是纤维 $X_t$ 中位于 $x$
上方的唯一点。注意，$U_T \subset X_T$ 是 $f_T$ 在其处拟有限的点集，见
态射章引理 \ref{morphisms-lemma-base-change-quasi-finite}.
还要注意
$$
X_T \xrightarrow{f'_T} S'_T \xrightarrow{\nu_T} T
$$
是 $T$ 在 $X_T$ 中的正规化，见
引理 \ref{more-morphisms-lemma-normalization-smooth-localization}.
假设该断言对 $y \in U_T \subset X_T \to S'_T \to T$ 成立，即假设可找到开邻域
$f'_T(y) \in V' \subset S'_T$，使得 $(f'_T)^{-1}V' \to V'$ 是同构。态射
$S'_T \to S'$ 为 \'etale，故 $V'$ 的像 $V \subset S'$ 是开集。由于
$f'_T(y) \in V'$，有 $f'(x) \in V$。注意图
$$
\xymatrix{
(f'_T)^{-1}V' \ar[r] \ar[d] & (f')^{-1}(V) \ar[d] \\
V' \ar[r] & V
}
$$
是纤维方块（因为 $S'_T \times_{S'} X = X_T$）。由于左侧竖直箭头是同构，且
$\{V' \to V\}$ 是 \'etale 覆盖，由
下降章引理 \ref{descent-lemma-descending-property-isomorphism}.
可知右侧竖直箭头是同构。换言之，该断言对 $x \in U \subset X \to S' \to S$ 成立。

\medskip\noindent
由上一段的结果，为证明该断言，可以把 $S$ 替换为 $s = f(x)$ 的初等 \'etale 邻域。
因此可假设有分解
$$
X = V \amalg W
$$
为开闭子概形，其中 $V \to S$ 有限且 $x \in V$，见 Lemma
\ref{more-morphisms-lemma-etale-splits-off-quasi-finite-part-technical}。由于 $X$ 是 $S$ 上 $V$ 与
$W$ 的不交并，且 $V \to S$ 有限，可知 $S$ 在 $X$ 中的正规化是态射
$$
X = V \amalg W \longrightarrow V \amalg W' \longrightarrow S
$$
，其中 $W'$ 是 $S$ 在 $W$ 中的正规化，见
态射章引理 \ref{morphisms-lemma-normalization-in-disjoint-union},
\ref{morphisms-lemma-finite-integral} 与
\ref{morphisms-lemma-normalization-in-integral}.
该断言随即成立，证明完成。
\end{proof}

\begin{lemma}
\label{more-morphisms-lemma-quasi-finite-separated-quasi-affine}
\begin{slogan}
拟有限分离态射是拟仿射态射
\end{slogan}
设 $f : X \to S$ 是概形态射。假设 $f$ 拟有限且分离。设 $S'$ 是 $S$ 在 $X$
中的正规化，见
态射章定义 \ref{morphisms-definition-normalization-X-in-Y}.
图示为：
$$
\xymatrix{
X \ar[rd]_f \ar[rr]_{f'} & & S' \ar[ld]^\nu \\
& S &
}
$$
则 $f'$ 是拟紧开浸入，而 $\nu$ 是整态射。特别地，$f$ 拟仿射。
\end{lemma}

\begin{proof}
这由 引理 \ref{more-morphisms-lemma-finite-type-separated} 得到。事实上，由该引理，存在开子概形
$U' \subset S'$，使得 $(f')^{-1}(U') = X$，且 $X \to U'$ 是同构。换言之，
$f'$ 是开浸入。注意，$f'$ 拟紧，因为 $f$ 拟紧且 $\nu : S' \to S$ 分离
（概形章引理 \ref{schemes-lemma-quasi-compact-permanence}）。于是由
态射章引理 \ref{morphisms-lemma-characterize-quasi-affine}.
可知 $f$ 拟仿射。
\end{proof}

\begin{lemma}[Zariski's Main Theorem]
\label{more-morphisms-lemma-quasi-finite-separated-pass-through-finite}
\begin{reference}
\cite[IV Corollary 18.12.13]{EGA}
\end{reference}
设 $f : X \to S$ 是概形态射。假设 $f$ 拟有限且分离，并假设 $S$ 拟紧且拟分离。
则存在分解
$$
\xymatrix{
X \ar[rd]_f \ar[rr]_j & & T \ar[ld]^\pi \\
& S &
}
$$
，其中 $j$ 是拟紧开浸入，而 $\pi$ 有限。
\end{lemma}

\begin{proof}
设 $X \to S' \to S$ 如 Lemma
\ref{more-morphisms-lemma-quasi-finite-separated-quasi-affine} 的结论。由
性质章引理
\ref{properties-lemma-integral-algebra-directed-colimit-finite}
，可写成 $\nu_*\mathcal{O}_{S'} = \colim_{i \in I} \mathcal{A}_i$，即有限拟凝聚
$\mathcal{O}_S$-代数
$\mathcal{A}_i \subset \nu_*\mathcal{O}_{S'}$ 的有向余极限。于是对每个 $i$，
$\pi_i : T_i = \underline{\Spec}_S(\mathcal{A}_i) \to S$
是有限态射。注意，转移态射 $T_{i'} \to T_i$ 仿射，且 $S' = \lim T_i$。

\medskip\noindent
由 极限章引理 \ref{limits-lemma-descend-opens}，存在一个 $i$ 及拟紧开集
$U_i \subset T_i$，其在 $S'$ 中的逆像等于 $f'(X)$。对 $i' \geq i$，令
$U_{i'}$ 为 $U_i$ 在 $T_{i'}$ 中的逆像。于是
$X \cong f'(X) = \lim_{i' \geq i} U_{i'}$，见
极限章引理 \ref{limits-lemma-directed-inverse-system-has-limit}.
由 极限章引理 \ref{limits-lemma-finite-type-eventually-closed} 可知，对某个
$i' \geq i$，$X \to U_{i'}$ 是闭浸入。（事实上，对充分大的 $i'$，
$X \cong U_{i'}$，但我们不需要这一点。）因此 $X \to T_{i'}$ 是浸入。由
态射章引理 \ref{morphisms-lemma-factor-quasi-compact-immersion}
可将其分解为 $X \to T \to T_{i'}$，其中第一个箭头是开浸入，第二个是闭浸入。
结论得证。
\end{proof}

\begin{lemma}
\label{more-morphisms-lemma-quasi-finite-separated-pass-through-finite-addendum}
沿用 引理 \ref{more-morphisms-lemma-quasi-finite-separated-pass-through-finite} 的记号与假设。
再假设 $f$ 局部有限表现。则可选择该分解，使得 $T$ 在 $S$ 上有限且有限表现。
\end{lemma}

\begin{proof}
由 极限章引理
\ref{limits-lemma-finite-in-finite-and-finite-presentation}，可写
$T = \lim T_i$，其中所有 $T_i$ 在 $S$ 上有限且有限表现，转移态射
$T_{i'} \to T_i$ 是闭浸入。由
极限章引理 \ref{limits-lemma-descend-opens}
，存在一个 $i$ 及开子概形 $U_i \subset T_i$，其在 $T$ 中的逆像是 $X$。由
极限章引理
\ref{limits-lemma-finite-type-eventually-closed}
可知对充分大的 $i$ 有 $X \cong U_i$。用 $T_i$ 替换 $T$ 即完成证明。
\end{proof}








\section{Zariski 主定理的应用，I}
\label{more-morphisms-section-applications-zmt}

\noindent
第一个应用是把有限态射刻画为具有有限纤维的固有态射。

\begin{lemma}
\label{more-morphisms-lemma-characterize-finite}
设 $f : X \to S$ 是概形态射。下列条件等价：
\begin{enumerate}
\item $f$ 有限；
\item $f$ 固有且纤维有限；
\item $f$ 固有且局部拟有限；
\item $f$ 普遍闭、分离、局部有限型且纤维有限。
\end{enumerate}
\end{lemma}

\begin{proof}
由 态射章引理 \ref{morphisms-lemma-finite-proper},
\ref{morphisms-lemma-quasi-finite},
以及 \ref{morphisms-lemma-finite-quasi-finite}，(1) 蕴含 (2)。由 态射章引理
\ref{morphisms-lemma-finite-fibre}，(2) 蕴含 (3)。由固有态射的定义以及 Morphisms,
引理 \ref{morphisms-lemma-quasi-finite-locally-quasi-compact} 和
\ref{morphisms-lemma-quasi-finite}，(3) 蕴含 (4)。

\medskip\noindent
假设 (4)。取 $s \in S$。由 态射章引理
\ref{morphisms-lemma-finite-fibre}，$X_s$ 的有限多个点在 $X_s$ 中全都是孤立点。
按 引理 \ref{more-morphisms-lemma-etale-splits-off-quasi-finite-part}，选择初等 \'etale 邻域
$(U, u) \to (S, s)$ 以及分解 $X_U = V \amalg W$。
注意，由于 $X_s$ 的所有点都是孤立点，$W_u = \emptyset$。由于 $f$ 普遍闭，$W$
在 $U$ 中的像是不含 $u$ 的闭集。缩小 $U$ 后，可假设 $W = \emptyset$。换言之，
$X_U = V$ 在 $U$ 上有限。由于 $s \in S$ 任意，这意味着存在像覆盖 $S$ 的
\'etale 态射族 $\{U_i \to S\}$，使得基变换 $X_{U_i} \to U_i$ 有限。注意，
$\{U_i \to S\}$ 是 \'etale 覆盖，见 拓扑章定义
\ref{topologies-definition-etale-covering}。因此它是 fpqc 覆盖，见
拓扑章
引理 \ref{topologies-lemma-zariski-etale-smooth-syntomic-fppf-fpqc}。故由
下降章引理 \ref{descent-lemma-descending-property-finite} 可知 $f$ 有限。
\end{proof}

\noindent
由此得到以下有用结果。

\begin{lemma}
\label{more-morphisms-lemma-proper-finite-fibre-finite-in-neighbourhood}
设 $f : X \to S$ 是概形态射。设 $s \in S$。假设 $f$ 固有且
$f^{-1}(\{s\})$ 是有限集。则存在 $s$ 的开邻域 $V \subset S$，使得
$f|_{f^{-1}(V)} : f^{-1}(V) \to V$ 有限。
\end{lemma}

\begin{proof}
由 态射章引理 \ref{morphisms-lemma-finite-fibre}，态射 $f$ 在
$f^{-1}(\{s\})$ 的所有点处拟有限。由 态射章引理
\ref{morphisms-lemma-quasi-finite-points-open}，$f$ 在其处拟有限的点集是开集
$U \subset X$。令 $Z = X \setminus U$。则 $s \not \in f(Z)$。由于 $f$ 固有，
集合 $f(Z) \subset S$ 闭。选取 $s$ 的任意开邻域 $V \subset S$，使得
$f(Z) \cap V = \emptyset$。则 $f^{-1}(V) \to V$ 局部拟有限且固有，故拟有限
（态射章引理 \ref{morphisms-lemma-quasi-finite-locally-quasi-compact}），
从而纤维有限（态射章引理 \ref{morphisms-lemma-quasi-finite}），进而由
引理 \ref{more-morphisms-lemma-characterize-finite} 可知它有限。
\end{proof}

\begin{lemma}
\label{more-morphisms-lemma-flat-proper-family-cannot-collapse-fibre}
考虑概形的交换图
$$
\xymatrix{
X \ar[rr]_h \ar[rd]_f & & Y \ar[ld]^g \\
& S
}
$$
设 $s \in S$。假设
\begin{enumerate}
\item $X \to S$ 是固有态射；
\item $Y \to S$ 分离且局部有限型；并且
\item $X_s \to Y_s$ 的像有限。
\end{enumerate}
则存在包含 $s$ 的开子空间 $U \subset S$，使得 $X_U \to Y_U$ 经在 $U$ 上有限的
闭子概形 $Z \subset Y_U$ 分解。
% P05-ERR-CH37-0092: Frozen source says “open subspace” although S is a scheme and the intended object is an open subscheme.
\end{lemma}

\begin{proof}
设 $Z \subset Y$ 是 $h$ 的概形论像，见
态射章第 \ref{morphisms-section-scheme-theoretic-image} 节.
由 态射章引理 \ref{morphisms-lemma-scheme-theoretic-image-is-proper}，态射
$X \to Z$ 满，而 $Z \to S$ 固有。因此 $X_s \to Z_s$ 满。可知
% P05-ERR-CH37-0093: Frozen source has the stray word “either” in “We see that either (3) implies...”.
(3) 蕴含 $Z_s$ 有限。故由
引理 \ref{more-morphisms-lemma-proper-finite-fibre-finite-in-neighbourhood}.
$Z \to S$ 在 $s$ 的一个开邻域上有限。
\end{proof}





\section{Zariski 主定理的应用，II}
\label{more-morphisms-section-applications-zmt-II}

\noindent
本节给出 Zariski 主定理关于拟有限态射结构的另外若干推论。

\begin{lemma}
\label{more-morphisms-lemma-separated-locally-quasi-finite-over-affine}
设 $f : X \to Y$ 是分离且局部拟有限的态射，并且 $Y$ 仿射。则 $X$ 的每个有限点集
都包含在 $X$ 的某个仿射开集内。
\end{lemma}

\begin{proof}
设 $x_1, \ldots, x_n \in X$。选择拟紧开集 $U \subset X$，使得 $x_i \in U$。
由 引理 \ref{more-morphisms-lemma-quasi-finite-separated-quasi-affine}，$U \to Y$ 拟仿射。
故由
性质章引理 \ref{properties-lemma-ample-finite-set-in-affine}.
，存在包含 $x_1, \ldots, x_n$ 的仿射开集 $V \subset U$。
\end{proof}

\begin{lemma}
\label{more-morphisms-lemma-quasi-finite-finite-over-dense-open}
设 $f : Y \to X$ 是拟有限态射。则存在稠密开集 $U \subset X$，使得
$f|_{f^{-1}(U)} : f^{-1}(U) \to U$ 有限。
\end{lemma}

\begin{proof}
若 $U_i \subset X$（$i \in I$）是一族开集，使得限制
$f|_{f^{-1}(U_i)} : f^{-1}(U_i) \to U_i$ 有限，则令 $U = \bigcup U_i$，限制
$f|_{f^{-1}(U)} : f^{-1}(U) \to U$ 也有限，见
态射章引理 \ref{morphisms-lemma-finite-local}.
因此该问题在 $X$ 上是局部的，可假设 $X$ 仿射。

\medskip\noindent
假设 $X$ 仿射。写成 $Y = \bigcup_{j = 1, \ldots, m} V_j$，其中 $V_j$ 仿射。
这是可能的，因为 $f$ 拟有限，特别地拟紧。每个 $V_j \to X$ 拟有限且分离。
设 $\eta \in X$ 是 $X$ 的某个不可约分支的泛点。由 态射章引理
\ref{morphisms-lemma-quasi-finite} 与 \ref{morphisms-lemma-generically-finite}
可知，存在开邻域 $\eta \in U_\eta$，使得 $f^{-1}(U_\eta) \cap V_j \to U_\eta$
有限。可以选择对每个 $j = 1, \ldots, m$ 都适用的 $U_\eta$。注意，$X$ 的泛点集
在 $X$ 中稠密。因此存在稠密开集 $W = \bigcup_\eta U_\eta$，使得每个
$f^{-1}(W) \cap V_j \to W$ 都有限。只需证明存在稠密开集 $U \subset W$，使得
$f|_{f^{-1}(U)} : f^{-1}(U) \to U$ 有限。因此可把 $X$ 替换为 $W$ 的一个仿射开子概形，
并假设每个 $V_j \to X$ 有限。

\medskip\noindent
假设 $X$ 仿射，$Y = \bigcup_{j = 1, \ldots, m} V_j$ 且 $V_j$ 仿射，并且限制
$f|_{V_j} : V_j \to X$ 有限。令
$$
\Delta_{ij} =
\Big(\overline{V_i \cap V_j} \setminus V_i \cap V_j\Big) \cap V_j.
$$
这是 $V_j$ 的无处稠密闭子集，因为它是开子集 $V_i \cap V_j$ 在 $V_j$ 中的边界。由
态射章引理 \ref{morphisms-lemma-image-nowhere-dense-finite}
，像 $f(\Delta_{ij})$ 是 $X$ 的无处稠密闭子集。由
拓扑章引理 \ref{topology-lemma-nowhere-dense}
，并集 $T = \bigcup f(\Delta_{ij})$ 是 $X$ 的无处稠密闭子集。因此
$U = X \setminus T$ 是 $X$ 的稠密开子集。我们断言
$f|_{f^{-1}(U)} : f^{-1}(U) \to U$ 有限。为此，设 $U' \subset U$ 是仿射开集。
令 $Y' = f^{-1}(U') = U' \times_X Y$、
$V_j' = Y' \cap V_j = U' \times_X V_j$。考虑 $f$ 的限制
$$
f' = f|_{Y'} : Y' \longrightarrow U'
$$
。此态射现在具有如下性质：$Y' = \bigcup_{j = 1, \ldots, m} V'_j$ 是仿射开覆盖，
每个 $V'_j \to U'$ 有限，而 $V_i' \cap V_j'$ 在 $V'_i$ 与 $V'_j$ 中都（开且）闭。
因此 $V_i' \cap V_j'$ 仿射，且映射
$$
\mathcal{O}(V'_i) \otimes_{\mathbf{Z}} \mathcal{O}(V'_j)
\longrightarrow
\mathcal{O}(V'_i \cap V'_j)
$$
满射。这蕴含 $Y'$ 分离，见
概形章引理 \ref{schemes-lemma-characterize-separated}.
最后考虑交换图
$$
\xymatrix{
\coprod_{j = 1, \ldots, m} V'_j \ar[rd] \ar[rr] & & Y' \ar[ld] \\
& U' &
}
$$
东南方向的箭头有限、因而固有，水平箭头满，而西南方向的箭头分离。因此由
态射章引理 \ref{morphisms-lemma-image-proper-is-proper}
可知 $Y' \to U'$ 固有。由于它也拟有限，由 Lemma
\ref{more-morphisms-lemma-characterize-finite} 可知它有限，结论得证。
\end{proof}

\begin{lemma}
\label{more-morphisms-lemma-stratify-flat-fp-lqf-universally-bounded}
设 $f : X \to S$ 平坦、局部有限表现、分离、局部拟有限，并且纤维普遍有界。
则存在闭子集
$$
\emptyset = Z_{-1} \subset Z_0 \subset Z_1 \subset Z_2 \subset
\ldots \subset Z_n = S
$$
，使得令 $S_r = Z_r \setminus Z_{r - 1}$ 后，分层
$S = \coprod_{r = 0, \ldots, n} S_r$ 由如下泛性质刻画：给定 $g : T \to S$，
投影 $X \times_S T \to T$ 是次数为 $r$ 的有限局部自由态射，当且仅当集合论意义下
$g(T) \subset S_r$。
\end{lemma}

\begin{proof}
设整数 $n$ 是 $X \to S$ 各纤维次数的上界。由 态射章引理
\ref{morphisms-lemma-base-change-universally-bounded}，任意基变换的纤维次数也以
$n$ 为上界。特别地，引理陈述中出现的所有整数 $r$ 都满足 $\leq n$。对 $n$
归纳证明本引理。基础情形 $n = 0$ 显然。

\medskip\noindent
我们断言，满足 $\deg_{\kappa(s)}(X_s) = n$ 的点 $s \in S$ 组成开子集
$S_n \subset S$，而 $X \times_S S_n \to S_n$ 是次数为 $n$ 的有限局部自由态射。
事实上，设 $s \in S$ 是这样的点。按
引理 \ref{more-morphisms-lemma-etale-splits-off-quasi-finite-part}.
选择初等 \'etale 态射 $(U, u) \to (S, s)$ 及分解
$U \times_S X = W \amalg V$。由于 $V \to U$ 有限、平坦且局部有限表现，
可知 $V \to U$ 有限局部自由，见
态射章引理 \ref{morphisms-lemma-finite-flat}.
把 $U$ 缩小为 $u$ 的更小邻域后，可假设 $V \to U$ 是某个次数 $d$ 的有限局部自由态射，见
态射章引理 \ref{morphisms-lemma-finite-locally-free}.
由于 $u \mapsto s$ 且 $W_u = \emptyset$，有 $d = n$。由于 $n$ 是纤维次数的最大值，
可知 $W = \emptyset$！因此 $U \times_S X \to U$ 是次数为 $n$ 的有限局部自由态射。
由 下降章引理 \ref{descent-lemma-descending-property-finite-locally-free}，
$X \to S$ 在 $s$ 的开邻域 $\Im(U \to S)$ 上是次数为 $n$ 的有限局部自由态射
（态射章引理 \ref{morphisms-lemma-etale-open}）。该断言得证。

\medskip\noindent
赋予 $S' = S \setminus S_n$ 约化诱导概形结构，并令 $X' = X \times_S S'$。
注意，$X' \to S'$ 的纤维次数普遍以 $n - 1$ 为上界。由归纳假设，得到适合于态射
$X' \to S'$ 的分层 $S' = S_0 \amalg \ldots \amalg S_{n - 1}$。我们断言
$S = \coprod_{r = 0, \ldots, n} S_r$ 对态射 $X \to S$ 适用。设
$g : T \to S$ 是概形态射，并假设 $X \times_S T \to T$ 是次数为 $r$ 的有限局部自由
态射。如上所述，这蕴含 $r \leq n$。若 $r = n$，则 $T \to S$ 显然经 $S_n$
分解。若 $r < n$，则
% P05-ERR-CH37-0094: Frozen source uses undefined S_d here; the complement fixed in the preceding sentence is S \ S_n.
$g(T) \subset S' = S \setminus S_d$ (集合论意义下) hence
集合论意义下成立，故 $T_{red} \to S$ 经 $S'$ 分解，见
概形章引理 \ref{schemes-lemma-map-into-reduction}.
注意，作为基变换，$X \times_S T_{red} \to T_{red}$ 也是次数为 $r$ 的有限局部自由态射。
由分层 $S' = \coprod_{r = 0, \ldots, n - 1} S_r$ 的泛性质，
$g(T) = g(T_{red})$ 包含于 $S_r$。反过来，假设有 $g : T \to S$，使得集合论意义下
$g(T) \subset S_r$。若 $r = n$，则 $g$ 经 $S_n$ 分解，而作为基变换，
$X \times_S T \to T$ 显然是次数为 $n$ 的有限局部自由态射。若 $r < n$，则
$X \times_S T \to T$ 分离、平坦且局部有限表现，并且它在 $T_{red}$ 上的限制是次数为
$r$ 的有限局部自由态射。由于 $T_{red} \to T$ 是普遍同胚，可知
$X \times_S T_{red} \to X \times_S T$ 也是普遍同胚，因而
$X \times_S T \to T$ 普遍闭（因为有限态射
$X \times_S T_{red} \to T_{red}$ 具有这一性质）。例如由 Lemma
\ref{more-morphisms-lemma-characterize-finite}，可知 $X \times_S T \to T$ 有限。再由 Morphisms,
引理 \ref{morphisms-lemma-finite-flat}，可知 $X \times_S T \to T$ 有限局部自由。最后，由于所有纤维的次数都是
$r$，其次数为 $r$。
\end{proof}

\begin{lemma}
\label{more-morphisms-lemma-stratify-flat-fp-qf}
设 $f : X \to S$ 是平坦、局部有限表现、分离且拟有限的概形态射。则存在闭子集
$$
\emptyset = Z_{-1} \subset Z_0 \subset Z_1 \subset Z_2 \subset
\ldots \subset S
$$
，使得令 $S_r = Z_r \setminus Z_{r - 1}$ 后，分层 $S = \coprod S_r$ 由如下泛性质
刻画：给定态射 $g : T \to S$，投影 $X \times_S T \to T$ 是次数为 $r$ 的有限局部自由
态射，当且仅当集合论意义下 $g(T) \subset S_r$。此外，包含映射 $S_r \to S$ 拟紧。
\end{lemma}

\begin{proof}
该问题在 $S$ 上是局部的，故可假设 $S$ 仿射。由 态射章引理
\ref{morphisms-lemma-locally-quasi-finite-qc-source-universally-bounded}
，在此情形 $f$ 的纤维普遍有界。因此该分层的存在性由
引理 \ref{more-morphisms-lemma-stratify-flat-fp-lqf-universally-bounded}.
得到。

\medskip\noindent
我们将证明对每个 $r \geq 0$，$U_r = S \setminus Z_r \to S$ 拟紧。由初等拓扑，
这将证明最后一个断言。由于拟紧映射的复合拟紧，只需证明 $U_r \to U_{r - 1}$ 拟紧。
选择仿射开集 $W \subset U_{r - 1}$。写成 $W = \Spec(A)$。则对某个理想
$I \subset A$ 有 $Z_r \cap W = V(I)$，且
$X \times_S \Spec(A/I) \to \Spec(A/I)$ 是次数为 $r$ 的有限局部自由态射。注意
$A/I = \colim A/I_i$，其中 $I_i \subset I$ 遍历有限生成理想。由
极限章引理 \ref{limits-lemma-descend-finite-locally-free}
可知，对某个 $i$，$X \times_S \Spec(A/I_i) \to \Spec(A/I_i)$ 是次数为 $r$
的有限局部自由态射。（这里使用了 $X \to S$ 有限表现；事实上，它局部有限表现、分离且
拟紧。）因此，集合论意义下，$\Spec(A/I_i) \to \Spec(A) = W$ 经
$Z_r \cap W$ 分解。于是 $Z_r \cap W = V(I_i)$ 是 $A$ 中某个有限元素集的零点集。
这意味着 $W \setminus Z_r$ 是有限多个标准开集之并，因而拟紧，如所求。
\end{proof}

\begin{lemma}
\label{more-morphisms-lemma-stratify-flat-fp-lqf}
设 $f : X \to S$ 是平坦、局部有限表现、分离且局部拟有限的概形态射。则存在开子概形
$$
S = U_0 \supset U_1 \supset U_2 \supset \ldots
$$
，使得当 $k$ 是域时，态射 $\Spec(k) \to S$ 经 $U_d$ 分解，当且仅当
$X \times_S \Spec(k)$ 在 $k$ 上的次数 $\geq d$。
\end{lemma}

\begin{proof}
该陈述仅仅是说，纤维次数 $\geq d$ 的点集是开集。因此可在 $S$ 上局部工作，并假设
$S$ 仿射。在此情形，对每个拟紧开集 $W \subset X$，由 Lemma
\ref{more-morphisms-lemma-stratify-flat-fp-qf}，$W \to S$ 的纤维次数 $\geq d$ 的点集 $U_d(W)$
是开集。由于 $U_d = \bigcup_W U_d(W)$，结论成立。
\end{proof}

\begin{lemma}
\label{more-morphisms-lemma-go-down-with-annihilators}
设 $f : X \to S$ 是平坦、局部有限表现且局部拟有限的概形态射。设
$g \in \Gamma(X, \mathcal{O}_X)$ 非零。则存在满足 $g|_V \not = 0$ 的开集
$V \subset X$、开集 $U \subset S$，它们构成交换图
$$
\xymatrix{
V \ar[r] \ar[d]_\pi & X \ar[d]^f \\
U \ar[r] & S,
}
$$
；还存在拟凝聚子层 $\mathcal{F} \subset \mathcal{O}_U$、整数 $r > 0$ 以及单射的
$\mathcal{O}_U$-模映射 $\mathcal{F}^{\oplus r} \to \pi_*\mathcal{O}_V$，其像包含
$g|_V$。
\end{lemma}

\begin{proof}
可假设 $X$ 与 $S$ 仿射。由 Lemmas
\ref{more-morphisms-lemma-stratify-flat-fp-lqf-universally-bounded} 与
\ref{more-morphisms-lemma-stratify-flat-fp-qf}，得到滤过
$\emptyset = Z_{-1} \subset Z_0 \subset Z_1 \subset Z_2 \subset \ldots
\subset Z_n = S$。设 $T \subset X$ 是有限 $\mathcal{O}_X$-模
$\Im(g : \mathcal{O}_X \to \mathcal{O}_X)$ 的概形论支撑。注意，$T$ 是 $g$ 作为
$\mathcal{O}_X$ 截面的支撑（模章定义 \ref{modules-definition-support}），
且对任意开集 $V \subset X$，$g|_V \not = 0$ 当且仅当
$V \cap T \not = \emptyset$。设 $r$ 是使得集合论意义下 $f(T) \subset Z_r$ 的最小整数。
取 $\xi \in T$ 为 $T$ 的某个不可约分支的泛点，使得
$f(\xi) \not \in Z_{r - 1}$（因而 $f(\xi) \in Z_r$）。可把 $S$ 替换为 $f(\xi)$
的一个包含于 $S \setminus Z_{r - 1}$ 的仿射邻域。写成 $S = \Spec(A)$，并设
$I = (a_1, \ldots, a_m) \subset A$ 是有限生成理想，使得集合论意义下
$V(I) = Z_r$（见
代数章引理 \ref{algebra-lemma-qc-open}).
由 $r$ 的选择，$g$ 的支撑包含于 $f^{-1}V(I)$，故存在整数 $N$，使得对
$j = 1, \ldots, m$ 有 $a_j^N g = 0$。
% P05-ERR-CH37-0095: Frozen source replaces a_j by a_j^r although the annihilating exponent just constructed is N.
用 $a_j^r$ 替换 $a_j$ 后，可假设 $Ig = 0$。对任意 $A$-模 $M$，以 $M[I]$
表示 $M$ 的 $I$-挠部分，即 $M[I] = \{m \in M \mid Im = 0\}$。写成
$X = \Spec(B)$，于是 $g \in B[I]$。由于 $A \to B$ 平坦，可知
$$
B[I] = A[I] \otimes_A B \cong A[I] \otimes_{A/I} B/IB
$$
由 $Z_r$ 的选择，$A/I$-模 $B/IB$ 是秩为 $r$ 的有限局部自由模。因此，把 $S$
替换为 $f(\xi)$ 的更小仿射开邻域后，可假设作为 $A/I$-模有
$B/IB \cong (A/IA)^{\oplus r}$。选择映射 $\psi : A^{\oplus r} \to B$，使其模
$I$ 约化为上一句的同构。于是诱导映射
$$
A[I]^{\oplus r} \longrightarrow B[I]
$$
是同构。取 $\mathcal{F}$ 为与 $A$-模 $A[I]$ 相伴的拟凝聚层，并取映射
$\mathcal{F}^{\oplus r} \to \pi_*\mathcal{O}_V$ 为对应于
$A[I]^{\oplus r} \subset A^{\oplus r} \to B$ 的映射，即得本引理。
\end{proof}

\begin{lemma}
\label{more-morphisms-lemma-there-is-a-scheme-integral-over}
设 $U \to X$ 是概形的满 \'etale 态射。假设 $X$
拟紧且拟分离。则存在满整态射 $Y \to X$，使得对
每个 $y \in Y$，都有一个开邻域 $V \subset Y$，
且 $V \to X$ 经由 $U$ 分解。事实上，可以假设
$Y \to X$ 是有限且有限表示的。
\end{lemma}

\begin{proof}
由于 $X$ 拟紧，存在有限多个仿射开集
$U_i \subset U$，使得 $U' = \coprod U_i \to X$ 是满射。
用 $U'$ 替换 $U$ 后，可假设 $U$ 是仿射的。
特别地，$U \to X$ 是分离的
（概形章引理 \ref{schemes-lemma-affine-separated}）。
于是存在一个整数 $d$，界定 $U \to X$ 的几何纤维的次数
（见 态射章引理
\ref{morphisms-lemma-locally-quasi-finite-qc-source-universally-bounded}）。
我们对 $d$ 归纳，为所有满 \'etale 映到 $X$ 的拟紧分离概形
$U$ 证明本引理。若 $d = 1$，则 $U = X$，取 $Y = U$
即可得到结论。以下假设 $d > 1$。

\medskip\noindent
应用引理 \ref{more-morphisms-lemma-quasi-finite-separated-quasi-affine}，
得到一个分解
$$
\xymatrix{
U \ar[rr]_j \ar[rd] & & Y \ar[ld]^\pi \\
& X
}
$$
其中 $\pi$ 是整态射，$j$ 是拟紧开浸入。我们可以且确实假设
$j(U)$ 在 $Y$ 中概形论稠密。注意
$$
U \times_X Y = U \amalg W
$$
其中第一个直和项是 $U \to U \times_X Y$ 的像
（由 概形章引理 \ref{schemes-lemma-semi-diagonal}，它是闭的；
又因为它作为两个在 $Y$ 上 \'etale 的概形之间的态射是 \'etale 的，
所以也是开的），第二个直和项是其开闭补集。
$W$ 的像 $V \subset Y$ 是包含 $Y \setminus U$ 的开子概形。

\medskip\noindent
\'Etale 态射 $W \to Y$ 的几何纤维的基数均为 $< d$。
事实上，对 $U \subset Y$ 的几何点直接考察即可看出这一点。
由于 $U \subset Y$ 稠密，例如由引理
\ref{more-morphisms-lemma-stratify-flat-fp-lqf-universally-bounded}，
它对 $Y$ 的所有几何点均成立
（拟紧分离 \'etale 态射的纤维次数在特化下不会增大）。因此，可以将
归纳假设应用于 $W \to V$，得到一个满整态射
$Z \to V$，其中 $Z$ 是概形，且该态射在 Zariski 局部经由 $W$ 分解。
选取分解 $Z \to Z' \to Y$，使得 $Z' \to Y$ 是整态射，
$Z \to Z'$ 是开浸入
（引理 \ref{more-morphisms-lemma-quasi-finite-separated-quasi-affine}）。
用 $Z$ 在 $Z'$ 中的概形论闭包替换 $Z'$ 后，
可假设 $Z$ 在 $Z'$ 中概形论稠密。这样便有
$Z' \times_Y V = Z$。最后，令 $T \subset Y$ 为赋予
$Y \setminus V$ 的诱导既约闭子概形结构。考虑态射
$$
Z' \amalg T \longrightarrow X
$$
按构造，这是满整态射。由于 $T \subset U$，态射 $T \to X$
显然经由 $U$ 分解。另一方面，取一点 $z \in Z'$。
若 $z \not \in Z$，则 $z$ 映到 $Y \setminus V \subset U$
中的一点，因而可找到 $z$ 的一个邻域，使该邻域上的态射经由 $U$ 分解。
若 $z \in Z$，则存在邻域 $\Omega \subset Z$，它经由
$W \subset U \times_X Y$ 分解，因而也经由 $U$ 分解。
这证明了所需对象的存在性。

\medskip\noindent
假设已经找到整且满的 $Y \to X$，并且它在 Zariski 局部经由
$U$ 分解。选取有限仿射开覆盖 $Y = \bigcup V_j$，使得
$V_j \to X$ 经由 $U$ 分解。可以写成 $Y = \lim Y_i$，
其中 $Y_i \to X$ 是有限且有限表示的；见 极限章引理
\ref{limits-lemma-integral-limit-finite-and-finite-presentation}。
当 $i$ 充分大时，可以找到仿射开集 $V_{i, j} \subset Y_i$，
其在 $Y$ 中的逆像恰为 $V_j$；见
极限章引理 \ref{limits-lemma-descend-opens}。
进一步增大 $i$，则 $X$ 上的态射 $V_j \to U$ 来自
$X$ 上的态射 $V_{i, j} \to U$；见 极限章命题
\ref{limits-proposition-characterize-locally-finite-presentation}。
证明完毕。
\end{proof}







\section{连通纤维态射的应用}
\label{more-morphisms-section-application-geometrically-connected}

\noindent
本节证明若干引理，用以构造所有纤维均几何连通或几何整的态射。
这些结果稍后在研究有限型态射的局部结构时会很有用。

% P05-EN-0096: the frozen source reads "an a point"; Chinese follows the evident intended syntax.
\begin{lemma}
\label{more-morphisms-lemma-descent-connected-fibres}
考虑概形态射的图表
$$
\xymatrix{
Z \ar[r]_{\sigma} \ar[rd] & X \ar[d] \\
& Y
}
$$
以及一点 $y \in Y$。假设
\begin{enumerate}
\item $X \to Y$ 是有限表示且平坦的，
\item $Z \to Y$ 是有限局部自由的，
\item $Z_y \not = \emptyset$,
\item $X \to Y$ 的所有纤维均几何既约，并且
\item $X_y$ 在 $\kappa(y)$ 上几何连通。
\end{enumerate}
则存在拟紧开集 $X^0 \subset X$，使得 $X^0_y = X_y$，
并且 $X^0 \to Y$ 的所有非空纤维均几何连通。
\end{lemma}

\begin{proof}
本证明将用到平坦、有限表示和有限局部自由这些性质在基变换与复合下保持，
还将用到有限局部自由态射既是开态射又是闭态射。这些事实见
态射章引理
\ref{morphisms-lemma-base-change-flat},
\ref{morphisms-lemma-base-change-finite-presentation},
\ref{morphisms-lemma-base-change-finite-locally-free},
\ref{morphisms-lemma-composition-flat},
\ref{morphisms-lemma-composition-finite-presentation},
\ref{morphisms-lemma-composition-finite-locally-free},
\ref{morphisms-lemma-fppf-open}, 及
\ref{morphisms-lemma-finite-proper}.

\medskip\noindent
注意，$X_Z \to Z$ 是有限表示的平坦态射，并具有由 $\sigma$
给出的截面 $s$。以 $X_Z^0$ 表示
Situation \ref{more-morphisms-situation-connected-along-section}
所定义的 $X_Z$ 的子集。由
引理 \ref{more-morphisms-lemma-connected-along-section-open}，
它是 $X_Z$ 的开子集。

\medskip\noindent
$X$ 到 $Z \times_Y Z$ 的拉回 $X_{Z \times_Y Z}$ 带有两个截面
$s_0, s_1$，即 $s$ 沿
$\text{pr}_0, \text{pr}_1 : Z \times_Y Z \to Z$ 的基变换。
Situation \ref{more-morphisms-situation-connected-along-section} 的构造给出两个子集
$(X_{Z \times_Y Z})_{s_0}^0$ 和
$(X_{Z \times_Y Z})_{s_1}^0$。由
引理 \ref{more-morphisms-lemma-base-change-connected-along-section}，
它们分别是 $X_Z^0$ 在态射
$1_X \times \text{pr}_0, 1_X \times \text{pr}_1 : X_{Z \times_Y Z} \to X_Z$
下的逆像。特别地，这些子集是开的。

\medskip\noindent
令 $(Z \times_Y Z)_y = \{z_1, \ldots, z_n\}$。
由于 $X_y$ 几何连通，可见
$(X_{Z \times_Y Z})_{s_0}^0$ 和 $(X_{Z \times_Y Z})_{s_1}^0$
在每个 $z_i$ 上的纤维相同（都等于整个纤维）。换言之，
$$
s_0(z_i) \in (X_{Z \times_Y Z})_{s_1}^0
\quad\text{且}\quad
s_1(z_i) \in (X_{Z \times_Y Z})_{s_0}^0.
$$
由于子集 $(X_{Z \times_Y Z})_{s_0}^0$ 和 $(X_{Z \times_Y Z})_{s_1}^0$
在 $X_{Z \times_Y Z}$ 中是开的，存在 $(Z \times_Y Z)_y$ 的一个开邻域
$W \subset Z \times_Y Z$，使得
$$
s_0(W) \subset (X_{Z \times_Y Z})_{s_1}^0
\quad\text{且}\quad
s_1(W) \subset (X_{Z \times_Y Z})_{s_0}^0.
$$
于是由 Situation \ref{more-morphisms-situation-connected-along-section}
中的构造直接可得
$$
p^{-1}(W) \cap (X_{Z \times_Y Z})_{s_0}^0
=
p^{-1}(W) \cap (X_{Z \times_Y Z})_{s_1}^0
$$
% P05-EN-0097: the frozen source gives the projection target as Z \times_W Z;
% the evident ambient base is Z \times_Y Z. Formula retained unchanged.
其中 $p : X_{Z \times_Y Z} \to Z \times_W Z$ 是投影。
由于 $Z \times_Y Z \to Y$ 有限局部自由，因而既开又闭，
存在 $y$ 的一个仿射开邻域 $V \subset Y$，使得
$q^{-1}(V) \subset W$，其中 $q : Z \times_Y Z \to Y$ 是结构态射。
为证明本引理，可以用 $V$ 替换 $Y$。这样便看出
% P05-EN-0098: the frozen source says X_Z^0 \subset Y_Z; expected ambient is X_Z.
$X_Z^0 \subset Y_Z$ 是开集，并且
$$
(1_X \times \text{pr}_0)^{-1}(X_Z^0) =
(1_X \times \text{pr}_1)^{-1}(X_Z^0).
$$
这意味着 $X_Z^0$ 的像 $X^0 \subset X$ 是开集，且
$(X_Z \to X)^{-1}(X^0) = X_Z^0$；见
下降章引理 \ref{descent-lemma-open-fpqc-covering}。
最后，由 引理 \ref{more-morphisms-lemma-connected-along-section-locally-constructible}，
$X_Z^0$ 是拟紧的，因而 $X^0$ 也是拟紧的
（这里用到此时 $Y$ 仿射，故 $X$ 拟紧且拟分离；于是局部可构造等同于
可构造，特别是拟紧的；略去细节）。由此可见，$X^0$ 具有全部所需性质。
\end{proof}

\begin{lemma}
\label{more-morphisms-lemma-fibre-geometrically-connected-reduced}
设 $h : Y \to S$ 是概形态射，$s \in S$ 是一点，
$T \subset Y_s$ 是开子概形。假设
\begin{enumerate}
\item $h$ 平坦且有限表示，
\item $h$ 的所有纤维均几何既约，并且
\item $T$ 在 $\kappa(s)$ 上几何连通。
\end{enumerate}
则可以找到一个仿射初等 \'etale 邻域
$(S', s') \to (S, s)$ 和一个拟紧开集 $V \subset Y_{S'}$，使得
\begin{enumerate}
\item[(a)] $V \to S'$ 的所有纤维均几何连通；
\item[(b)] $V_{s'} = T \times_s s'$.
\end{enumerate}
\end{lemma}

\begin{proof}
问题显然在 $S$ 上是局部的，故可以用 $s$ 的一个仿射开邻域替换 $S$。
$Y_s$ 上的拓扑由 $Y$ 上的拓扑诱导；见
概形章引理 \ref{schemes-lemma-fibre-topological}。
因此，可以找到拟紧开集 $V \subset Y$，使得 $V_s = T$。
$h$ 在 $V$ 上的限制是拟紧的（因为 $S$ 仿射而 $V$ 拟紧）、
拟分离的、局部有限表示的且平坦的，因而是有限表示的平坦态射。
所以用 $V$ 替换 $Y$ 后，除 (1) 和 (2) 外，还可假设
$Y_s = T$ 且 $S$ 仿射。

\medskip\noindent
选取闭点 $y \in Y_s$，使得 $h$ 在 $y$ 处是 Cohen--Macaulay 的；见
引理 \ref{more-morphisms-lemma-flat-finite-presentation-CM-open}。
由 引理 \ref{more-morphisms-lemma-slice-CM}，存在图表
$$
\xymatrix{
Z \ar[r] \ar[rd] & Y \ar[d] \\
& S
}
$$
使得 $Z \to S$ 平坦、局部有限表示且局部拟有限，并且
$Z_s = \{y\}$。应用
引理 \ref{more-morphisms-lemma-etale-makes-quasi-finite-finite-at-point}，
找到一个初等邻域 $(S', s') \to (S, s)$ 和一个开集
$Z' \subset Z_{S'} = S' \times_S Z$，使得 $Z' \to S'$ 有限，
并且其中恰有一点 $z' \in Z'$ 位于 $s$ 上方。注意，
$Z' \to S'$ 还是局部有限表示且平坦的（因为它是 $Z \to S$
的基变换中的开集），故 $Z' \to S'$ 有限局部自由；见
态射章引理 \ref{morphisms-lemma-finite-flat}。
注意，作为 $h$ 的基变换，$Y_{S'} \to S'$ 平坦且有限表示，
并具有几何既约纤维。此外，$Y_{s'} = Y_s$ 几何连通。
将 引理 \ref{more-morphisms-lemma-descent-connected-fibres}
应用于 $S'$ 上的 $Z' \to Y_{S'}$，得到拟紧开集
$V \subset Y_{S'}$，它满足 (2)，且在 $S'$ 上的纤维或为空，
或几何连通。由于 $V \to S'$ 是开态射
（态射章引理 \ref{morphisms-lemma-fppf-open}），用 $s'$
的一个仿射开邻域替换 $S'$ 后，可以假设 $V \to S'$ 是满射，
从而 (1) 成立。
\end{proof}

\begin{lemma}
\label{more-morphisms-lemma-cover-by-geometrically-connected}
设 $f : X \to S$ 是局部有限表示且平坦、并具有几何既约纤维的概形态射。
则存在一个 \'etale 覆盖 $\{X_i \to X\}_{i \in I}$，
使得 $X_i \to S$ 分解为 $X_i \to S_i \to S$，其中
$S_i \to S$ 是 \'etale 态射，而 $X_i \to S_i$ 是有限表示的平坦态射，
并具有几何连通且几何既约的纤维。
\end{lemma}

\begin{proof}
取一点 $x \in X$，其像为 $s \in S$。我们将构造图表
$$
\xymatrix{
X' \ar[r] \ar[rd] & S' \times_S X \ar[r] \ar[d] & X \ar[d] \\
& S' \ar[r] & S
}
$$
以及点 $s' \in S'$、$x' \in X'$、$y \in S' \times_S X$，使得
$x'$ 映到 $x$，$(S', s') \to (S, s)$ 是 \'etale 邻域，
$(X', x') \to (S' \times_S X, y)$ 是 \'etale 邻域\footnote{证明实际上给出
一个开集 $X' \subset S' \times_S X$。}，并且
$X' \to S'$ 具有几何连通纤维。如果对每个 $x \in X$ 都能如此构造，
则本引理成立（覆盖的成员由这样得到的诸 \'etale 态射
$X' \to X$ 给出）。
% P05-EN-0099: the frozen source reads "is the replace"; Chinese follows the intended infinitive.
第一步是用 $x$ 与 $s$ 的仿射开邻域分别替换 $X$ 与 $S$，
从而化归到 $X$ 和 $S$ 均仿射的情形（此时 $f$ 为有限表示）。

\medskip\noindent
选取 $\kappa(s)$ 的一个可分代数扩张 $\overline{k}$，并以
$X_{\overline{k}}$ 表示 $X_s$ 的相应基变换。选取
$X_{\overline{k}}$ 中映到 $x \in X_s$ 的一点 $\overline{x}$，
以及 $\overline{x}$ 的一个连通拟紧开邻域
$\overline{V} \subset X_{\overline{k}}$。
（这是可能的，因为域上局部有限型的任意概形作为局部 Noether 拓扑空间
是局部连通的。）由
簇章引理 \ref{varieties-lemma-Galois-action-quasi-compact-open}，
可以找到有限可分扩张 $k'/\kappa(s)$ 和拟紧开集
$V' \subset X_{k'}$，使其基变换为 $\overline{V}$。特别地，
$V'$ 在 $k'$ 上几何连通；见
簇章引理 \ref{varieties-lemma-characterize-geometrically-connected}。
由 引理 \ref{more-morphisms-lemma-realize-prescribed-residue-field-extension-etale}，
可以找到 \'etale 邻域 $(S', s') \to (S, s)$，使得
$\kappa(s')$ 作为 $\kappa(s)$ 的扩张同构于 $k'$。
以 $x' \in (S' \times_S X)_{s'} = X_{k'}$ 表示 $\overline{x}$ 的像。
于是，用 $(S', s')$ 替换 $(S, s)$，并用
$(S' \times_S X, x')$ 替换 $(X, x)$ 后，就化归到下一段处理的情形。

\medskip\noindent
假设存在包含 $x$ 且几何连通的拟紧开集 $V \subset X_s$。
于是可以应用 引理 \ref{more-morphisms-lemma-fibre-geometrically-connected-reduced}，
找到一个仿射 \'etale 邻域 $(S', s') \to (S, s)$ 和一个拟紧开集
$X' \subset S' \times_S X$，使得 $X' \to S'$ 具有几何连通纤维，
并且 $X'$ 含有一个映到 $x$ 的点。证明完毕。
\end{proof}

\begin{lemma}
\label{more-morphisms-lemma-normal-morphism-irreducible}
设 $h : Y \to S$ 是概形态射，$s \in S$ 是一点，
$T \subset Y_s$ 是开子概形。假设
\begin{enumerate}
\item $h$ 是有限表示的，
\item $h$ 是正规态射，并且
\item $T$ 在 $\kappa(s)$ 上几何不可约。
\end{enumerate}
则可以找到一个仿射初等 \'etale 邻域
$(S', s') \to (S, s)$ 和一个拟紧开集 $V \subset Y_{S'}$，使得
\begin{enumerate}
\item[(a)] $V \to S'$ 的所有纤维均几何整；
\item[(b)] $V_{s'} = T \times_s s'$.
\end{enumerate}
\end{lemma}

\begin{proof}
应用 引理 \ref{more-morphisms-lemma-fibre-geometrically-connected-reduced}，
找到一个仿射初等 \'etale 邻域 $(S', s') \to (S, s)$ 和一个拟紧开集
$V \subset Y_{S'}$，使得 $V \to S'$ 的所有纤维均几何连通，并且
$V_{s'} = T \times_s s'$。由于 $V$ 是 $h$ 的基变换中的开集，
$V \to S'$ 的所有纤维均几何正规；见 引理 \ref{more-morphisms-lemma-normal}。
特别地，它们几何既约。为完成证明，只需说明它们几何不可约。
若 $t \in S'$，则 $V_t$ 在 $\kappa(t)$ 上有限型，因此
$V_t \times_{\kappa(t)} \overline{\kappa(t)}$ 在
$\overline{\kappa(t)}$ 上有限型，从而是 Noether 的。
由 $S' \to S$ 的选取，概形
$V_t \times_{\kappa(t)} \overline{\kappa(t)}$ 是连通的。
故由 性质章引理 \ref{properties-lemma-normal-Noetherian}，
$V_t \times_{\kappa(t)} \overline{\kappa(t)}$ 不可约，结论得证。
\end{proof}








\section{有限型态射结构的应用}
\label{more-morphisms-section-application-finite-type}

\noindent
本节的结果可见于 \cite{GruRay}。粗略地说，它断言：在源与目标上
\'etale 局部地，一个有限型态射是一个有限态射与一个具有几何连通纤维的
光滑态射之复合，而该光滑态射的相对维数等于原态射的纤维维数。

\begin{lemma}
\label{more-morphisms-lemma-local-structure-finite-type}
设 $f : X \to S$ 是态射。取 $x \in X$，并令 $s = f(x)$。
假设 $f$ 局部有限型，且 $n = \dim_x(X_s)$。则存在交换图表
$$
\xymatrix{
X \ar[dd] & X' \ar[l]^g \ar[d]^\pi & x \ar@{|->}[dd] &
x' \ar@{|->}[l]  \ar@{|->}[d] \\
& Y \ar[d]^h & & y \ar@{|->}[d] \\
S \ar@{=}[r] & S & s & s \ar@{=}[l]
}
$$
以及满足 $g(x') = x$ 的一点 $x' \in X'$，使得令 $y = \pi(x')$ 后，
\begin{enumerate}
\item $h : Y \to S$ 是相对维数为 $n$ 的光滑态射；
\item $g : (X', x') \to (X, x)$ 是初等 \'etale 邻域；
\item $\pi$ 有限，且 $\pi^{-1}(\{y\}) = \{x'\}$；
\item $\kappa(y)$ 是 $\kappa(s)$ 的纯超越扩张。
\end{enumerate}
此外，若 $f$ 局部有限表示，则 $\pi$ 是有限表示的。
\end{lemma}

\begin{proof}
问题在 $X$ 和 $S$ 上是局部的，故可假设 $X$ 与 $S$ 均仿射。
由 代数章引理 \ref{algebra-lemma-refined-quasi-finite-over-polynomial-algebra}，
用 $x$ 在 $X$ 中的一个标准开邻域替换 $X$ 后，可假设存在分解
$$
\xymatrix{
X \ar[r]^\pi & \mathbf{A}^n_S \ar[r] & S
}
$$
使得 $\pi$ 拟有限，并且 $\kappa(\pi(x))$ 是 $\kappa(s)$ 的纯超越扩张。
由 引理 \ref{more-morphisms-lemma-etale-makes-quasi-finite-finite-at-point}，
存在初等 \'etale 邻域
$$
(Y, y) \to (\mathbf{A}^n_S, \pi(x))
$$
以及开集 $X' \subset X \times_{\mathbf{A}^n_S} Y$；它包含唯一一个
位于 $y$ 上方的点 $x'$，并且 $X' \to Y$ 有限。
这证明了 (1) -- (4)。最后一个断言由
态射章引理 \ref{morphisms-lemma-finite-presentation-permanence} 得到。
\end{proof}

\begin{lemma}
\label{more-morphisms-lemma-local-local-structure-finite-type}
\begin{slogan}
在 \'etale 邻域中，有限型态射是光滑态射上的有限态射。
\end{slogan}
设 $f : X \to S$ 是态射。取 $x \in X$，并令 $s = f(x)$。
假设 $f$ 局部有限型，且 $n = \dim_x(X_s)$。则存在交换图表
$$
\xymatrix{
X \ar[dd] & X' \ar[l]^g \ar[d]^\pi & x \ar@{|->}[dd] &
x' \ar@{|->}[l] \ar@{|->}[d] \\
& Y' \ar[d]^h & & y' \ar@{|->}[d] \\
S & S' \ar[l]_e & s & s' \ar@{|->}[l]
}
$$
以及满足 $g(x') = x$ 的一点 $x' \in X'$，使得令 $y' = \pi(x')$、
$s' = h(y')$ 后，
\begin{enumerate}
\item $h : Y' \to S'$ 是相对维数为 $n$ 的光滑态射；
\item $Y' \to S'$ 的所有纤维均几何整；
\item $g : (X', x') \to (X, x)$ 是初等 \'etale 邻域；
\item $\pi$ 有限，且 $\pi^{-1}(\{y'\}) = \{x'\}$；
\item $\kappa(y')$ 是 $\kappa(s')$ 的纯超越扩张；
\item $e : (S', s') \to (S, s)$ 是初等 \'etale 邻域。
\end{enumerate}
此外，若 $f$ 局部有限表示，则 $\pi$ 是有限表示的。
\end{lemma}

\begin{proof}
问题在 $S$ 上是局部的，故可用 $s$ 的一个仿射开邻域替换 $S$。
接着应用 引理 \ref{more-morphisms-lemma-local-structure-finite-type}，得到交换图表
$$
\xymatrix{
X \ar[dd] & X' \ar[l]^g \ar[d]^\pi & x \ar@{|->}[dd] &
x' \ar@{|->}[l] \ar@{|->}[d] \\
& Y \ar[d]^h & & y \ar@{|->}[d] \\
S \ar@{=}[r] & S & s & s \ar@{=}[l]
}
$$
其中 $h$ 是相对维数为 $n$ 的光滑态射，而 $\kappa(y)$ 是
$\kappa(s)$ 的纯超越扩张。由于问题在 $X$ 上也是局部的，
可以用 $y$ 的一个仿射邻域替换 $Y$（并用其在 $\pi$ 下的逆像替换
$X'$）。由于 $S$ 仿射，这保证 $Y \to S$ 拟紧、分离且光滑，
特别是有限表示的。令 $T$ 为 $Y_s$ 中含有 $y$ 的连通分支。
由于 $Y_s$ 是 Noether 的，可见 $T$ 是开的。由 Varieties,
引理 \ref{varieties-lemma-geometrically-connected-if-connected-and-point}，
$T$ 在 $\kappa(s)$ 上几何连通。又因为 $T$ 在 $\kappa(s)$ 上光滑，
所以它几何正规；见
簇章引理 \ref{varieties-lemma-smooth-geometrically-normal}。
由此断定 $T$ 在 $\kappa(s)$ 上几何不可约
（因为连通的 Noether 正规概形不可约；见
性质章引理 \ref{properties-lemma-normal-Noetherian}）。
最后，由 引理 \ref{more-morphisms-lemma-smooth-normal}，光滑态射 $h$ 是正规态射。
% P05-EN-0100: the frozen source has singular "all assumption"; Chinese follows the intended plural.
至此已经验证：对态射 $h : Y \to S$ 和开集 $T \subset Y_s$，
引理 \ref{more-morphisms-lemma-normal-morphism-irreducible} 的全部假设均成立。
应用 引理 \ref{more-morphisms-lemma-normal-morphism-irreducible} 后，得到
$e : S' \to S$、$s' \in S'$ 和 $Y'$，如交换图表
$$
\xymatrix{
X \ar[dd] & X' \ar[l]^g \ar[d]^\pi & X' \times_Y Y' \ar[l] \ar[d] &
x \ar@{|->}[dd] & x' \ar@{|->}[l] \ar@{|->}[d] &
(x', s') \ar@{|->}[l] \ar@{|->}[d] \\
& Y \ar[d]^h & Y' \ar[d] \ar[l] & & y \ar@{|->}[d] &
(y, s') \ar@{|->}[l] \ar@{|->}[d] \\
S \ar@{=}[r] & S & S' \ar[l]_e & s & s \ar@{=}[l] & s' \ar@{|->}[l]
}
$$
所示，其中 $e : (S', s') \to (S, s)$ 是初等 \'etale 邻域，
而 $Y' \subset Y_{S'}$ 是一个开邻域，其在 $S'$ 上的所有纤维均几何不可约，
并且在等同 $Y_s = Y_{S', s'}$ 下有 $Y'_{s'} = T$。
令 $(y, s') \in Y'$ 为对应于 $y \in T$ 的点；它也是
$Y \times_S S'$ 中位于 $y$ 上方、剩余域等于 $\kappa(y)$ 且映到
$S'$ 中 $s'$ 的唯一一点。类似地，令
$(x', s') \in X' \times_Y Y' \subset X' \times_S S'$ 为位于
$x'$ 上方、剩余域等于 $\kappa(x')$ 且位于 $s'$ 上方的唯一一点。
于是，该图表的外部部分给出了本引理所提问题的解。略去若干次要细节。
\end{proof}

\begin{lemma}
\label{more-morphisms-lemma-local-local-structure-finite-type-affine}
假设与记号同 引理 \ref{more-morphisms-lemma-local-local-structure-finite-type}。
除性质 (1) -- (6) 外，还可以安排使得
\begin{enumerate}
\item[(7)] $S'$、$Y'$、$X'$ 均仿射。
\end{enumerate}
\end{lemma}

\begin{proof}
注意，若 $Y'$ 仿射，则由于 $\pi$ 有限，$X'$ 也仿射。
选取 $s'$ 的仿射开邻域 $U' \subset S'$，再选取 $y'$ 的仿射开邻域
$V' \subset h^{-1}(U')$。令 $W' = h(V')$。由
态射章引理 \ref{morphisms-lemma-smooth-open}，
这是 $S'$ 中包含于 $U'$ 的 $s'$ 的开邻域。选取 $s'$ 的仿射开邻域
$U'' \subset W'$。于是 $h^{-1}(U'') \cap V'$ 仿射，因为它等于
$U'' \times_{U'} V'$。按构造，
$h^{-1}(U'') \cap V' \to U''$ 是满光滑态射，其纤维是
$Y' \to S'$ 的几何整纤维的（非空）开子概形，因而也是几何整的。
所以可以用 $U''$ 替换 $S'$，并用 $h^{-1}(U'') \cap V'$ 替换 $Y'$。
\end{proof}

\noindent
性质 $\pi^{-1}(\{y'\}) = \{x'\}$ 的意义由下述引理部分说明。

\begin{lemma}
\label{more-morphisms-lemma-finite-morphism-single-point-in-fibre}
设 $\pi : X \to Y$ 是有限态射。取 $x \in X$，令 $y = \pi(x)$，
并假设 $\pi^{-1}(\{y\}) = \{x\}$。则
\begin{enumerate}
\item 对 $x$ 在 $X$ 中的每个邻域 $U \subset X$，都存在 $y$
的邻域 $V \subset Y$，使得 $\pi^{-1}(V) \subset U$。
\item 环映射 $\mathcal{O}_{Y, y} \to \mathcal{O}_{X, x}$ 是有限的。
\item 若 $\pi$ 是有限表示的，则
$\mathcal{O}_{Y, y} \to \mathcal{O}_{X, x}$ 是有限表示的。
\item 对任意拟凝聚 $\mathcal{O}_X$-模 $\mathcal{F}$，作为
$\mathcal{O}_{Y, y}$-模，有 $\mathcal{F}_x = \pi_*\mathcal{F}_y$。
\end{enumerate}
\end{lemma}

\begin{proof}
第一个断言纯属拓扑性质；只需使用 $\pi$ 是连续闭映射并且
$\pi^{-1}(\{y\}) = \{x\}$。为证明第二、第三部分，可以假设
$X = \Spec(B)$ 且 $Y = \Spec(A)$。此时 $A \to B$ 是有限环映射，
而 $y$ 对应于 $A$ 的素理想 $\mathfrak p$，并且 $B$ 中恰有一个
位于 $\mathfrak p$ 上方的素理想 $\mathfrak q$。于是
$B_{\mathfrak q} = B_{\mathfrak p}$；见
代数章引理 \ref{algebra-lemma-unique-prime-over-localize-below}。
换言之，映射 $A_{\mathfrak p} \to B_{\mathfrak q}$ 等于将
$A \to B$ 在 $\mathfrak p$ 处局部化所得的映射
$A_{\mathfrak p} \to B_{\mathfrak p}$。因此，(2) 和 (3) 由局部化的
简单性质得到（略去若干细节）。对最后一个断言，设对某个 $B$-模 $M$
有 $\mathcal{F} = \widetilde M$。
% P05-EN-0101: frozen source has F = M_q here; the stalk F_x is evidently intended.
则 $\mathcal{F} = M_{\mathfrak q}$ 且
$\pi_*\mathcal{F}_y = M_{\mathfrak p}$。由上述讨论，这两个局部化相同。
也可以直接利用 (1) 和茎的定义看出
$\mathcal{F}_x = \pi_*\mathcal{F}_y$。
\end{proof}






\section{fppf 拓扑的应用}
\label{more-morphisms-section-application-fppf}

\noindent
可以利用上述 \'etale 局部化技术证明下述结果：fppf 拓扑等于由 Zariski
覆盖以及形如 $\{f : T \to S\}$ 的覆盖“生成”的拓扑，其中 $f$
是满的有限局部自由态射。

\begin{lemma}
\label{more-morphisms-lemma-dominate-fppf-etale-locally}
设 $S$ 是概形，$\{S_i \to S\}_{i \in I}$ 是 fppf 覆盖。则存在
\begin{enumerate}
\item 一个 \'etale 覆盖 $\{S'_a \to S\}$；
\item 满的有限局部自由态射 $V_a \to S'_a$；
\end{enumerate}
使得 fppf 覆盖 $\{V_a \to S\}$ 加细给定覆盖 $\{S_i \to S\}$。
\end{lemma}

\begin{proof}
可以假设每个 $S_i \to S$ 都局部拟有限；见
引理 \ref{more-morphisms-lemma-qf-fp-flat-dominates-fppf}。

\medskip\noindent
固定一点 $s \in S$。选取 $i \in I$ 以及映到 $s$ 的一点
$s_i \in S_i$。选取初等 \'etale 邻域
$(S', s) \to (S, s)$，使得存在开集
$$
S_i \times_S S'  \supset V
$$
它包含唯一一个映到 $s \in S'$ 的点 $v \in V$，并且
$V \to S'$ 有限；见
引理 \ref{more-morphisms-lemma-etale-makes-quasi-finite-finite-at-point}。
于是 $V \to S'$ 有限局部自由：它是有限的，而
$S_i \times_S S' \to S'$ 作为 $S_i \to S$ 的基变换平坦且局部有限表示；见
态射章引理 \ref{morphisms-lemma-base-change-finite-presentation},
\ref{morphisms-lemma-base-change-flat}, 及
\ref{morphisms-lemma-finite-flat}。
因此 $V \to S'$ 是开态射；缩小 $S'$ 后，可以假设 $V \to S'$
是满的有限局部自由态射。由于对 $S$ 的每一点都可以这样做，断定
$\{S_i \to S\}$ 可由形如 $\{V_a \to S\}_{a \in A}$ 的覆盖加细，
其中每个 $V_a \to S$ 分解为 $V_a \to S'_a \to S$，
$S'_a \to S$ 为 \'etale 态射，而 $V_a \to S'_a$ 为满的有限局部自由态射。
\end{proof}

\begin{lemma}
\label{more-morphisms-lemma-dominate-fppf}
设 $S$ 是概形，$\{S_i \to S\}_{i \in I}$ 是 fppf 覆盖。则存在
\begin{enumerate}
\item Zariski 开覆盖 $S = \bigcup U_j$；
\item 满的有限局部自由态射 $W_j \to U_j$；
\item Zariski 开覆盖 $W_j = \bigcup_k W_{j, k}$；
\item 满的有限局部自由态射 $T_{j, k} \to W_{j, k}$；
\end{enumerate}
使得 fppf 覆盖 $\{T_{j, k} \to S\}$ 加细给定覆盖
$\{S_i \to S\}$。
\end{lemma}

\begin{proof}
令 $\{V_a \to S\}_{a \in A}$ 为
引理 \ref{more-morphisms-lemma-dominate-fppf-etale-locally} 中得到的 fppf 覆盖。
换言之，该覆盖加细 $\{S_i \to S\}$，并且每个 $V_a \to S$
分解为 $V_a \to S'_a \to S$，其中 $S'_a \to S$ 是 \'etale 态射，
$V_a \to S'_a$ 是满的有限局部自由态射。

\medskip\noindent
由 说明 \ref{more-morphisms-remark-topologies}，存在 Zariski 开覆盖
$S = \bigcup U_j$；对每个 $j$，存在有限局部自由满态射
$W_j \to U_j$；并且对每个 $j$，存在 Zariski 开覆盖
$\{W_{j, k} \to W_j\}$，使得族 $\{W_{j, k} \to S\}$ 加细
\'etale 覆盖 $\{S'_a \to S\}$。也就是说，对每一对 $j, k$，
存在 $a(j, k)$ 以及 $W_{j, k} \to S$ 的分解
$W_{j, k} \to S'_a \to S$。令
$T_{j, k} = W_{j, k} \times_{S'_a} V_a$，一切即清楚。
\end{proof}

\begin{lemma}
\label{more-morphisms-lemma-extend-integral-surjective-morphisms}
设 $S$ 是概形。若 $U \subset S$ 是开集且 $V \to U$ 是满整态射，
则存在满整态射 $\overline{V} \to S$，使得作为 $U$ 上的概形，
$\overline{V} \times_S U$ 同构于 $V$。
\end{lemma}

\begin{proof}
% P05-EN-0102: the frozen source says normalization of S in U, but the claimed inverse image is V;
% the construction evidently requires normalization/integral closure of S in V. Rendered formula/variables retained.
令 $V' \to S$ 为 $S$ 在 $U$ 中的正规化；见
态射章第 \ref{morphisms-section-normalization-X-in-Y} 节。
按构造，$V' \to S$ 是整态射。由 态射章引理
\ref{morphisms-lemma-normalization-localization} 及
\ref{morphisms-lemma-normalization-in-integral}，可见 $U$ 在 $V'$ 中的逆像是
$V$。令 $Z$ 为赋予 $S \setminus U$ 的既约诱导概形结构。于是
$\overline{V} = V' \amalg Z$ 即满足要求。
\end{proof}

\begin{lemma}
\label{more-morphisms-lemma-extend-finite-surjective-morphisms}
设 $S$ 是拟紧拟分离概形。若 $U \subset S$ 是拟紧开集，且
$V \to U$ 是满有限态射，则存在满有限态射
$\overline{V} \to S$，使得作为 $U$ 上的概形，
$\overline{V} \times_S U$ 同构于 $V$。
\end{lemma}

\begin{proof}
由 Zariski 主定理
（引理 \ref{more-morphisms-lemma-quasi-finite-separated-pass-through-finite}），
可以假设 $V$ 是某个在 $S$ 上有限的概形 $V'$ 中的拟紧开集。
用 $V$ 的概形论像替换 $V'$ 后，可以假设 $V$ 在 $V'$ 中稠密。
于是 $V' \times_S U = V$，因为 $V$ 在 $U$ 上有限，故
$V \to V' \times_S U$ 是闭的。令 $Z$ 为赋予
$S \setminus U$ 的既约诱导概形结构。则
$\overline{V} = V' \amalg Z$ 即满足要求。
\end{proof}

\begin{lemma}
\label{more-morphisms-lemma-dominate-fppf-integral}
设 $S$ 是概形，$\{S_i \to S\}_{i \in I}$ 是 fppf 覆盖。
则存在满整态射 $S' \to S$ 和开覆盖 $S' = \bigcup U'_\alpha$，
使得对每个 $\alpha$，态射 $U'_\alpha \to S$ 经由某个 $i$ 所对应的
$S_i \to S$ 分解。
\end{lemma}

\begin{proof}
按 引理 \ref{more-morphisms-lemma-dominate-fppf} 选取
$S = \bigcup U_j$、$W_j \to U_j$、$W_j = \bigcup W_{j, k}$ 以及
$T_{j, k} \to W_{j, k}$。由
引理 \ref{more-morphisms-lemma-extend-integral-surjective-morphisms}，可以将
$W_j \to U_j$ 延拓为满整态射 $\overline{W}_j \to S$。
之后，可以将 $T_{j, k} \to W_{j, k}$ 延拓为满整态射
$\overline{T}_{j, k} \to \overline{W}_j$。
令 $\overline{T}_j$ 等于所有概形 $\overline{T}_{j, k}$ 在
$\overline{W}_j$ 上的积
（极限章引理 \ref{limits-lemma-infinite-product}）。
再令 $S'$ 等于所有概形 $\overline{T}_j$ 在 $S$ 上的积。
若 $x \in S'$，则存在 $j$，使得 $x$ 在 $S$ 中的像属于 $U_j$。
因此存在 $k$，使得 $x$ 在投影 $S' \to \overline{W}_j$ 下的像属于
$W_{j, k}$。从而在投影
$S' \to \overline{T}_j \to \overline{T}_{j, k}$ 下，点 $x$ 落在
$T_{j, k}$ 中，而 $T_{j, k} \to S$ 经由某个 $i$ 所对应的 $S_i$ 分解。
最后，由 极限章引理 \ref{limits-lemma-infinite-product-integral} 及
\ref{limits-lemma-infinite-product-surjective}，态射 $S' \to S$ 整且满。
\end{proof}

\begin{lemma}
\label{more-morphisms-lemma-dominate-fppf-finite}
设 $S$ 是拟紧拟分离概形，$\{S_i \to S\}_{i \in I}$ 是 fppf 覆盖。
则存在有限表示的满有限态射 $S' \to S$ 以及开覆盖
$S' = \bigcup U'_\alpha$，使得对每个 $\alpha$，态射
$U'_\alpha \to S$ 经由某个 $i$ 所对应的 $S_i \to S$ 分解。
\end{lemma}

\begin{proof}
% P05-EN-0103: the frozen proof switches from the statement's base S to undefined X throughout.
% Variables are retained; the evident intended base is S.
令 $Y \to X$ 为 引理 \ref{more-morphisms-lemma-dominate-fppf-integral} 中得到的整满态射。
选取有限仿射开覆盖 $Y = \bigcup V_j$，使得
$V_j \to X$ 经由 $S_{i(j)}$ 分解。可以写成 $Y = \lim Y_\lambda$，
其中 $Y_\lambda \to X$ 有限且有限表示；见
极限章引理 \ref{limits-lemma-integral-limit-finite-and-finite-presentation}。
当 $\lambda$ 充分大时，可以找到仿射开集
$V_{\lambda, j} \subset Y_\lambda$，其在 $Y$ 中的逆像恰为 $V_j$；见
极限章引理 \ref{limits-lemma-descend-opens}。
进一步增大 $\lambda$，则 $X$ 上的态射 $V_j \to S_{i(j)}$ 来自
$X$ 上的态射 $V_{\lambda, j} \to S_{i(j)}$；见 极限章命题
\ref{limits-proposition-characterize-locally-finite-presentation}。
对这个 $\lambda$ 令 $S' = Y_\lambda$，证明完毕。
\end{proof}

\begin{lemma}
\label{more-morphisms-lemma-fppf-ph}
概形的 fppf 覆盖是 ph 覆盖。
\end{lemma}

\begin{proof}
设 $\{T_i \to T\}$ 是概形的 fppf 覆盖；见
拓扑章定义 \ref{topologies-definition-fppf-covering}。
注意，$T_i \to T$ 局部有限型。设 $U \subset T$ 是仿射开集。
只需说明 $\{T_i \times_T U \to U\}$ 可由一个标准 ph 覆盖加细；见
拓扑章定义 \ref{topologies-definition-ph-covering}。
这立即由 引理 \ref{more-morphisms-lemma-dominate-fppf-finite} 以及有限态射为固有态射这一事实得到
（态射章引理 \ref{morphisms-lemma-finite-proper}）。
\end{proof}

\begin{remark}
\label{more-morphisms-remark-change-topologies}
作为 引理 \ref{more-morphisms-lemma-fppf-ph} 的推论，得到比较态射
$$
\epsilon : (\Sch/S)_{ph} \longrightarrow (\Sch/S)_{fppf}
$$
这是由底层范畴上的恒等函子给出的位点态射
（按 拓扑章说明 \ref{topologies-remark-choice-sites} 适当选取位点）。
函子 $\epsilon_*$ 在底层预层上是恒等函子，而函子
% P05-EN-0104: frozen source omits the verb in "epsilon^{-1} associated ... its ph sheafification".
$\epsilon^{-1}$ 将一个 fppf 层对应到其 ph 层化。
通过复合，还可以将 ph 拓扑与 syntomic、光滑、\'etale 及 Zariski
拓扑进行比较。
\end{remark}




\section{拟投射概形}
\label{more-morphisms-section-quasi-projective}

\noindent
术语“拟投射概形”尚未定义。一种可能的定义是：具有充足可逆层的概形。
然而，若 $X$ 是基概形 $S$ 上的概形，则当态射 $X \to S$ 拟投射时，
我们称{\it $X$ 在 $S$ 上拟投射}
（态射章定义 \ref{morphisms-definition-quasi-projective}）。
由于任意概形的恒等态射都是拟投射的，可见在 $S$ 上拟投射的概形未必具有
充足可逆层。因此，最好仍不定义术语“拟投射概形”。

\begin{lemma}
\label{more-morphisms-lemma-quasi-projective}
设 $S$ 是具有充足可逆层的概形，$f : X \to S$ 是概形态射。
下列条件等价：
\begin{enumerate}
\item $X \to S$ 拟投射；
\item $X \to S$ 是 H-拟投射的；
\item 存在 $S$ 上概形的拟紧开浸入 $X \to X'$，并且
$X' \to S$ 投射；
\item $X \to S$ 有限型，且 $X$ 具有充足可逆层；
\item $X \to S$ 有限型，且存在一个 $f$-极充足可逆层。
\end{enumerate}
\end{lemma}

\begin{proof}
蕴含 (2) $\Rightarrow$ (1) 是
态射章引理 \ref{morphisms-lemma-H-quasi-projective-quasi-projective}。
蕴含 (1) $\Rightarrow$ (2) 是 态射章引理
\ref{morphisms-lemma-projective-over-quasi-projective-is-H-projective}。
蕴含 (2) $\Rightarrow$ (3) 是
态射章引理 \ref{morphisms-lemma-H-quasi-projective-open-H-projective}

\medskip\noindent
假设 $X \subset X'$ 如 (3) 所述。特别地，$X \to S$ 有限型。
% P05-EN-0105: frozen source calls X -> S H-projective here, while the cited lemma and ensuing immersion
% characterize H-quasi-projective; rendered terminology follows the frozen source.
由 态射章引理 \ref{morphisms-lemma-H-quasi-projective-open-H-projective}，
态射 $X \to S$ 是 H-投射的。因此存在拟紧浸入
$i : X \to \mathbf{P}^n_S$。于是
$\mathcal{L} = i^*\mathcal{O}_{\mathbf{P}^n_S}(1)$ 是 $f$-极充足的。
由于 $X \to S$ 拟紧，由
态射章引理 \ref{morphisms-lemma-ample-very-ample}，
$\mathcal{L}$ 是 $f$-充足的。因此按定义 $X \to S$ 拟投射。

\medskip\noindent
蕴含 (4) $\Rightarrow$ (2) 是
态射章引理 \ref{morphisms-lemma-quasi-projective-finite-type-over-S}.

\medskip\noindent
蕴含 (5) $\Rightarrow$ (4) 由
态射章引理 \ref{morphisms-lemma-pullback-ample-tensor-relatively-ample}.

\medskip\noindent
条件 (1) 与 (5) 的等价性由
态射章引理 \ref{morphisms-lemma-ample-very-ample} 及
\ref{morphisms-lemma-finite-type-ample-very-ample}.
\end{proof}

\begin{lemma}
\label{more-morphisms-lemma-category-quasi-projective}
设 $S$ 是具有充足可逆层的概形。令 $\text{QP}_S$ 为 $S$ 上概形范畴的
全子范畴，其对象满足 引理 \ref{more-morphisms-lemma-quasi-projective} 中的等价条件。
\begin{enumerate}
\item 若 $S' \to S$ 是概形态射且 $S'$ 具有充足可逆层，则基变换确定函子
$\text{QP}_S \to \text{QP}_{S'}$；
\item 若 $X \in \text{QP}_S$ 且 $Y \in \text{QP}_X$，则 $Y \in \text{QP}_S$；
\item 范畴 $\text{QP}_S$ 对纤维积封闭；
\item 范畴 $\text{QP}_S$ 对有限不交并封闭；
\item 若 $X \to S$ 投射，则 $X \in \text{QP}_S$；
\item 若 $X \to S$ 是有限型拟仿射态射，则 $X$ 属于 $\text{QP}_S$；
\item 若 $X \to S$ 拟有限且分离，则 $X \in \text{QP}_S$；
\item 若 $X \to S$ 是拟紧浸入，则 $X \in \text{QP}_S$；
% P05-EN-0106: the frozen source contains this editorial placeholder as a numbered mathematical item.
\item 此处添加更多内容。
\end{enumerate}
\end{lemma}

\begin{proof}
第 (1) 部分由 态射章引理
\ref{morphisms-lemma-base-change-quasi-projective}.

\medskip\noindent
第 (2) 部分由 引理 \ref{more-morphisms-lemma-quasi-projective} 的第四种刻画得到。

\medskip\noindent
若 $X \to S$ 与 $Y \to S$ 拟投射，则由
态射章引理 \ref{morphisms-lemma-base-change-quasi-projective}，
$X \times_S Y \to Y$ 拟投射。因此 (3) 由 (2) 得到。

\medskip\noindent
若 $X = Y \amalg Z$ 是概形的不交并，而 $\mathcal{L}$ 是可逆
$\mathcal{O}_X$-模，并且 $\mathcal{L}|_Y$ 与 $\mathcal{L}|_Z$
均充足，则 $\mathcal{L}$ 充足（略去细节）。因此，第 (4) 部分由
引理 \ref{more-morphisms-lemma-quasi-projective} 的第四种刻画得到。

\medskip\noindent
第 (5) 部分由
态射章引理 \ref{morphisms-lemma-projective-quasi-projective}.

\medskip\noindent
第 (6) 部分由
态射章引理
\ref{morphisms-lemma-quasi-affine-finite-type-quasi-projective}.

\medskip\noindent
第 (7) 部分由第 (6) 部分及
引理 \ref{more-morphisms-lemma-quasi-finite-separated-quasi-affine}.

\medskip\noindent
第 (8) 部分由第 (7) 部分及
态射章引理 \ref{morphisms-lemma-immersion-locally-quasi-finite}.
\end{proof}

\noindent
下述引理其实并不属于本节，但在别处似乎也没有合适的位置。

\begin{lemma}
\label{more-morphisms-lemma-integral-over-quasi-affine}
设 $X$ 是拟仿射概形，$f : U \to X$ 是整态射。
则 $U$ 拟仿射，并且图表
$$
\xymatrix{
U \ar[r] \ar[d] & \Spec(\Gamma(U, \mathcal{O}_U)) \ar[d] \\
X \ar[r] & \Spec(\Gamma(X, \mathcal{O}_X))
}
$$
是笛卡尔的。
\end{lemma}

\begin{proof}
概形 $U$ 拟仿射，因为整态射是仿射的，仿射态射是拟仿射的；一个概形
拟仿射，当且仅当它到 $\Spec(\mathbf{Z})$ 的结构态射拟仿射；并且
拟仿射态射的复合拟仿射。前两个陈述立即由定义得到，第三个是
态射章引理 \ref{morphisms-lemma-composition-quasi-affine}。
令 $U' =
X \times_{\Spec(\Gamma(X, \mathcal{O}_X))} \Spec(\Gamma(U, \mathcal{O}_U))$
，并考虑扩展图表
$$
\xymatrix{
U \ar[r]_j \ar[rd] & U' \ar[d] \ar[r] &
\Spec(\Gamma(U, \mathcal{O}_U)) \ar[d] \\
& X \ar[r] & \Spec(\Gamma(X, \mathcal{O}_X))
}
$$
由 态射章引理 \ref{morphisms-lemma-image-proper-scheme-closed}，
结合整态射普遍闭这一事实
（态射章引理 \ref{morphisms-lemma-integral-universally-closed}）
% P05-EN-0107: frozen source reads "vertical arrows are in the diagram are separated".
以及图表中的竖直箭头分离这一事实，态射 $j$ 是闭的。
另一方面，由 性质章引理 \ref{properties-lemma-quasi-affine}，
引理图表中的水平箭头是开的，故 $j$ 是开的。因此，$j$ 将 $U$
等同于 $U'$ 的一个开闭子概形。若 $U \not = U'$，则 $U$ 在 $U'$
中不稠密，更不可能在 $\Gamma(U, \mathcal{O}_U)$ 的谱中稠密。
然而，$U$ 在 $\Spec(\Gamma(U, \mathcal{O}_U))$ 中的概形论像是
$\Spec(\Gamma(U, \mathcal{O}_U))$，因为
$\Gamma(U, \mathcal{O}_U)$ 中任何切出一个使 $U$ 的态射经其分解的
闭子概形的理想都必须为零。因此，例如由
态射章引理 \ref{morphisms-lemma-quasi-compact-scheme-theoretic-image}，
$U$ 在 $\Spec(\Gamma(U, \mathcal{O}_U))$ 中稠密。
所以 $U = U'$，结论得证。
\end{proof}









\section{投射概形}
\label{more-morphisms-section-projective}

\noindent
本节是 第 \ref{more-morphisms-section-quasi-projective} 节 针对投射态射的对应版本。

\begin{lemma}
\label{more-morphisms-lemma-projective}
设 $S$ 是具有充足可逆层的概形，$f : X \to S$ 是概形态射。
下列条件等价：
\begin{enumerate}
\item $X \to S$ 投射；
\item $X \to S$ 是 H-投射的；
\item $X \to S$ 拟投射且固有；
\item $X \to S$ 是 H-拟投射且固有的；
\item $X \to S$ 固有，且 $X$ 具有充足可逆层；
\item $X \to S$ 固有，且存在一个 $f$-充足可逆层；
\item $X \to S$ 固有，且存在一个 $f$-极充足可逆层；
\item 存在拟凝聚分次 $\mathcal{O}_S$-代数 $\mathcal{A}$，它由
$\mathcal{A}_1$ 在 $\mathcal{A}_0$ 上生成，其中 $\mathcal{A}_1$
是有限型 $\mathcal{O}_S$-模，并且
$X = \underline{\text{Proj}}_S(\mathcal{A})$. 
\end{enumerate}
\end{lemma}

\begin{proof}
首先注意，在每种情形下态射 $f$ 都是固有的；见
态射章引理 \ref{morphisms-lemma-H-projective} 及
\ref{morphisms-lemma-locally-projective-proper}。
因此，只需在 $f$ 为固有态射的情形证明这些概念等价。
以下将直接使用这一点而不再说明。

\medskip\noindent
等价关系 (1) $\Leftrightarrow$ (3) 和 (2) $\Leftrightarrow$ (4) 是
态射章引理 \ref{morphisms-lemma-projective-is-quasi-projective-proper}.

\medskip\noindent
蕴含 (2) $\Rightarrow$ (1) 是
态射章引理 \ref{morphisms-lemma-H-projective}.

\medskip\noindent
蕴含 (1) $\Rightarrow$ (2) 和 (3) $\Rightarrow$ (4) 是
态射章引理
\ref{morphisms-lemma-projective-over-quasi-projective-is-H-projective}.

\medskip\noindent
蕴含 (1) $\Rightarrow$ (7) 立即由
态射章定义 \ref{morphisms-definition-projective} 及
\ref{morphisms-definition-very-ample}.

\medskip\noindent
由 态射章定义 \ref{morphisms-definition-quasi-projective}，
条件 (3) 和 (6) 等价。

\medskip\noindent
因此，(1) -- (4)、(6) 等价且蕴含 (7)。由
引理 \ref{more-morphisms-lemma-quasi-projective}，条件 (3)、(5) 和 (7) 等价。
由此可见，(1) -- (7) 等价。

\medskip\noindent
由 除子章引理 \ref{divisors-lemma-relative-proj-projective}，
(8) 蕴含 (1)。反之，若 (2) 成立，则可选取闭浸入
$$
i :
X
\longrightarrow
\mathbf{P}^n_S = \underline{\text{Proj}}_S(\mathcal{O}_S[T_0, \ldots, T_n]).
$$
该等式见 构造章引理 \ref{constructions-lemma-projective-space-bundle}。
由 除子章引理 \ref{divisors-lemma-closed-subscheme-proj}，
$X$ 是 $\mathcal{O}_S[T_0, \ldots, T_n]$ 的某个拟凝聚分次商代数
$\mathcal{A}$ 的相对 Proj。于是 $\mathcal{A}$ 满足 (8) 的条件。
\end{proof}

\begin{lemma}
\label{more-morphisms-lemma-category-projective}
设 $S$ 是具有充足可逆层的概形。令 $\text{P}_S$ 为 $S$ 上概形范畴的
全子范畴，其对象满足 引理 \ref{more-morphisms-lemma-projective} 中的等价条件。
\begin{enumerate}
\item 若 $S' \to S$ 是概形态射且 $S'$ 具有充足可逆层，则基变换确定函子
$\text{P}_S \to \text{P}_{S'}$；
\item 若 $X \in \text{P}_S$ 且 $Y \in \text{P}_X$，则 $Y \in \text{P}_S$；
\item 范畴 $\text{P}_S$ 对纤维积封闭；
\item 范畴 $\text{P}_S$ 对有限不交并封闭；
\item 若 $X \to S$ 有限，则 $X$ 属于 $\text{P}_S$；
% P05-EN-0108: the frozen source contains this editorial placeholder as a numbered mathematical item.
\item 此处添加更多内容。
\end{enumerate}
\end{lemma}

\begin{proof}
第 (1) 部分由 态射章引理
\ref{morphisms-lemma-base-change-projective}.

\medskip\noindent
第 (2) 部分由 引理 \ref{more-morphisms-lemma-projective} 的第五种刻画，
以及固有态射的复合仍固有这一事实得到
（态射章引理 \ref{morphisms-lemma-composition-proper}）。

\medskip\noindent
若 $X \to S$ 与 $Y \to S$ 投射，则由
态射章引理 \ref{morphisms-lemma-base-change-projective}，
$X \times_S Y \to Y$ 投射。因此 (3) 由 (2) 得到。

\medskip\noindent
若 $X = Y \amalg Z$ 是概形的不交并，而 $\mathcal{L}$ 是可逆
$\mathcal{O}_X$-模，并且 $\mathcal{L}|_Y$ 与 $\mathcal{L}|_Z$
均充足，则 $\mathcal{L}$ 充足（略去细节）。因此，第 (4) 部分由
引理 \ref{more-morphisms-lemma-projective} 的第五种刻画得到。

\medskip\noindent
第 (5) 部分由
态射章引理 \ref{morphisms-lemma-finite-projective}.
\end{proof}

\noindent
下面给出一个类型稍有不同的结果。

\begin{lemma}
\label{more-morphisms-lemma-ample-in-neighbourhood}
\begin{reference}
\cite[IV Corollary 9.6.4]{EGA}
\end{reference}
设 $f : X \to Y$ 是概形的固有态射，$\mathcal{L}$ 是可逆
$\mathcal{O}_X$-模。设 $y \in Y$ 是一点，并且 $\mathcal{L}_y$
在 $X_y$ 上充足。则存在 $y$ 的开邻域 $V \subset Y$，使得
$\mathcal{L}|_{f^{-1}(V)}$ 在 $f^{-1}(V)/V$ 上充足。
\end{lemma}

\begin{proof}
可以假设 $Y$ 仿射。于是找到有向集 $I$ 以及概形态射的逆系统
$X_i \to Y_i$，其中 $Y_i$ 在 $\mathbf{Z}$ 上有限型，转移态射
$X_i \to X_{i'}$ 与 $Y_i \to Y_{i'}$ 仿射，$X_i \to Y_i$ 固有，
并且 $X \to Y = \lim (X_i \to Y_i)$。见 极限章引理
\ref{limits-lemma-proper-limit-of-proper-finite-presentation-noetherian}。
缩小 $I$ 后，可以假设有一族相容的可逆 $\mathcal{O}_{X_i}$-模
$\mathcal{L}_i$，其拉回为 $\mathcal{L}$；见
极限章引理 \ref{limits-lemma-descend-invertible-modules}。
令 $y_i \in Y_i$ 为 $y$ 的像。则
$\kappa(y) = \colim \kappa(y_i)$。因此，由
极限章引理 \ref{limits-lemma-limit-ample}，对某个 $i$，
$\mathcal{L}_{i, y_i}$ 在 $X_{i, y_i}$ 上充足。
由 概形上同调章引理 \ref{coherent-lemma-ample-in-neighbourhood}，
可找到 $y_i$ 的开邻域 $V_i \subset Y_i$，使得 $\mathcal{L}_i$
在 $f_i^{-1}(V_i)$ 上的限制相对于 $V_i$ 充足。
令 $V \subset Y$ 为 $V_i$ 的逆像，即完成证明
（提示：使用 态射章引理 \ref{morphisms-lemma-ample-base-change}、
$X \to Y \times_{Y_i} X_i$ 为仿射态射这一事实，以及由
态射章引理 \ref{morphisms-lemma-pullback-ample-tensor-relatively-ample}
可知充足可逆层沿仿射态射的拉回仍充足这一事实）。
\end{proof}







\section{Proj 与 Spec}
\label{more-morphisms-section-proj-spec}

\noindent
本节阐明分次环的 Proj 与谱之间的关系。

\medskip\noindent
设 $R$ 是环，$A$ 是分次 $R$-代数；见
代数章第 \ref{algebra-section-graded} 节。
对 $m \geq 0$，记 $A_{\geq m} = \bigoplus_{d \geq m} A_d$。
考虑分次环
$$
B = \bigoplus\nolimits_{d \geq 0} A_{\geq d}
$$
对 $d' \geq d$ 及 $a \in A_{d'}$，以 $a^{(d)} \in B$ 表示
$B_d$ 中与 $a$ 对应的元素。以 $\sigma : A \to B$ 和
$\psi : A \to B$ 表示两个显然的环映射：若 $a \in A_d$，则
$\sigma(a) = a^{(0)}$ 且 $\psi(a) = a^{(d)}$。
于是 $\psi$ 是分次环映射，而 $\sigma$ 使 $B$ 成为 $A$ 上的分次代数。
还有一个满分次环映射 $\tau : B \to A$：对 $d' \geq d$ 及
$a \in A_{d'}$，若 $d' > d$，它将 $a^{(d)}$ 映到 $0$；
若 $d' = d$，则映到 $a$。

\medskip\noindent
仿射概形与谱。令 $X = \Spec(A)$。无关理想 $A_+$ 切出闭子概形
$Z = V(A_+) = \Spec(A/A_+) = \Spec(A_0)$。令 $U = X \setminus Z$。
$$
U \longrightarrow X \longrightarrow Z
$$
投射概形与 Proj。令 $P = \text{Proj}(A)$。我们可以且确实将 $P$
视为 $\Spec(A_0) = Z$ 上的概形。令 $L = \text{Proj}(B)$。
我们可以且确实将 $L$ 视为 $\Spec(B_0) = \Spec(A) = X$ 上的概形；
注意，$B_0$ 与 $A$ 的等同由 $\sigma$ 给出。满射 $\tau$ 定义闭浸入
$0 : P \to L$。由于 $A \xrightarrow{\sigma} B \to A$ 等于映射
$A \to A_0 \to A$，断定
$$
\xymatrix{
P \ar[d] \ar[r]_0 & L \ar[d] \\
Z \ar[r] & X
}
$$
交换。

\medskip\noindent
我们断言 $\psi$ 定义态射 $L \to P$。由 构造章引理
\ref{constructions-lemma-morphism-proj}，只需对每个满足
$B_+ \not \subset \mathfrak p$ 的齐次素理想 $\mathfrak p \subset B$
验证 $\psi(A_+) \not \subset \mathfrak p$。首先，选取齐次元素
$g \in B_+$，使得 $g \not \in \mathfrak p$。可以将 $g$ 写成有限和
$g = \sum a_i^{(d)}$，其中对某些 $d_i \geq d$ 有
$a_i \in A_{d_i}$。从而存在 $d' \geq d$ 及 $a \in A_{d'}$，
使得 $a^{(d)} \not \in \mathfrak p$。于是
$$
(a^{(d)})^{d'} =
(a^{d'})^{(d'd)} =
a^{(d)} (a^{d' - 1})^{(d(d' - 1))} =
\psi(a) (a^{d' - 1})^{(d(d' - 1))}
$$
（这个记号并不十分理想）不属于 $\mathfrak p$。因此
$\psi(a) \not \in \mathfrak p$，断言得证。于是可以将上图扩展为交换图表
$$
\xymatrix{
P \ar[d] \ar[r]_0 & L \ar[d] \ar[r]_\pi & P \ar[d] \\
Z \ar[r] & X \ar[r] & Z
}
$$
其中 $X \to Z$ 由 $A_0 \to A$ 给出。由于
$\tau \circ \psi = \text{id}_A$，可见 $\pi \circ 0 = \text{id}_P$。

\medskip\noindent
注意，$\pi$ 是仿射态射。这由 构造章引理
\ref{constructions-lemma-morphism-proj} 中的构造显然可见。
事实上，若对某个 $d > 0$ 有 $f \in A_d$，则令 $g = \psi(f)$ 后，有
$\pi^{-1}(D_+(f)) = D_+(g)$。在此情形，齐次部分有如下等式
$$
(B[1/g])_{m'} = \bigoplus\nolimits_{m \geq m'} (A[1/f])_m
$$
该同构与进一步的局部化相容。取 $m' = 0$，可见
$\pi_*\mathcal{O}_L$ 是所有 $m \geq 0$ 的 $\mathcal{O}_P(m)$ 的直和\footnote{类似地可得
$\pi_*\mathcal{O}_L(i) = \bigoplus_{m \geq -i} \mathcal{O}_P(m)$。}。
于是 $L$ 等同于相对谱：
$$
L = \underline{\Spec}_P
\left(
\bigoplus\nolimits_{m \geq 0} \mathcal{O}_P(m)
\right)
$$
特别地，$L \to P$ 是锥\footnote{$L$ 往往是 $P$ 上的线丛；见下文。}；见
构造章第 \ref{constructions-section-cone} 节。
此外，$0 : P \to L$ 显然是该锥的顶点。

\medskip\noindent
如上一段，取某个 $d > 0$ 的 $f \in A_d$，并令
$g = \psi(f) \in B_d$。考察环映射的结构
$$
\xymatrix{
A_0 \ar[r] \ar[d] & A \ar[d]^\sigma \ar[r] & A_0 \ar[d] \\
(A[1/f])_0 \ar[r]^-\psi &
(B[1/g])_0 = \bigoplus\nolimits_{m \geq 0} (A[1/f])_m
\ar[r]^-\tau &
(A[1/f])_0
}
$$
% P05-EN-0109: the frozen source misspells "computations" as "compuations".
分次环中的一些计算\footnote{(1) 和 (2) 是清楚的。
为看出 (3)，注意若 $a \in A_d$，则
$\sigma(a) = \sigma(f) \psi(a/f)$。对 (4)，注意
$b/g^m$ 属于 $\tau$ 的核，当且仅当
$b \in A_{\geq md}$ 在 $A_{md}$ 中的像为零。
因此，只需说明：若 $m' > md$ 且 $a \in A_{m'}$，则
$a^{(md)}/g^m$ 的某个幂属于由 $\sigma(f)$ 生成的理想。
取 $e$ 使得 $em' - emd \geq d$。于是
$$
(a^{(md)}/g^m)^e = (a^e)^{(emd)}/g^{em} =
(fa^e)^{(emd + d)}/g^{em + 1} = \sigma(f) \cdot (a^e)^{(emd + d)}/g^{em + 1}
$$
正如所需（抱歉使用了如此糟糕的记号）。为看出 (5)，如前论证，并注意当
$d = 1$ 时，$a^{(md)}/g^m = \sigma(f) \cdot a^{(md + 1)}/g^{m + 1}$。}
表明
\begin{enumerate}
\item $\sigma(A_+)(B[1/g])_0 \subset \Ker(\tau : (B[1/g])_0 \to (A[1/f])_0)$,
\item $\sigma(f) \in (B[1/g])_0$ 是非零因子；
\item 作为理想，$\sigma(f) (B[1/g])_0 = \sigma(A_d) (B[1/g])_0$；
\item $\sigma(f) (B[1/g])_0$ 与
$\Ker(\tau : (B[1/g])_0 \to (A[1/f])_0)$ 具有相同的根；
\item 若 $d = 1$，则
$\sigma(f) (B[1/g])_0 = \Ker(\tau : (B[1/g])_0 \to (A[1/f])_0)$.
\end{enumerate}
特别地，可见
$$
0(D_+(f)) = V(\sigma(f)) \subset D_+(g) = \Spec((B[1/g])_0)
$$
在集合论意义下成立。换言之，由 $\sigma(A_d)$ 生成的理想在
$D_+(g)$ 上切出一个有效 Cartier 除子，它在集合论意义下等于闭浸入
$0 : P \to L$ 的像。

\medskip\noindent
我们断言 $L \to X$ 在 $U$ 上是同构。事实上，若对某个 $d > 0$
有 $f \in A_d$，则
$$
\Spec(A_f) \times_X L = \text{Proj}(A_f \otimes_A B) =
\text{Proj}(B_{\sigma(f)})
$$
对每个 $e$，有
$(B_{\sigma(f)})_e = A_f \otimes_B B_e = A_f \otimes_A A_{\geq e} = A_f$,
其中最后一个等式由单射 $A_{\geq e} \subset A$ 诱导。
因此 $B_{\sigma(f)} \cong A_f[T]$，其中 $T$ 的次数为 $1$。
由于 $\text{Proj}(A_f[T]) \to \Spec(A_f)$ 是同构，断言得证。
从现在起，将 $U$ 与 $L$ 中相应的开集等同。

\medskip\noindent
上一段所作的等同使我们可以考虑限制 $\pi|_U : U \to P$。
像上面多次所做的那样，选取某个 $d > 0$ 的 $f \in A_d$，并令
$g = \psi(f) \in B_d$。则
$$
U \cap \pi^{-1}(D_+(f)) = U \cap D_+(g)
$$
在将 $D_+(g)$ 与 $(B[1/g])_0$ 的谱等同后，它是
$\sigma(f) \in (B[1/g])_0$ 的零点轨迹的补集。这正是上面的断言 (4)。
因此 $U \cap D_+(g)$ 仿射，并且
$$
\mathcal{O}_L(U \cap D_+(g)) = (B[1/g])_0[1/\sigma(f)] =
\bigoplus\nolimits_{m \in \mathbf{Z}} (A[1/f])_m
$$
其中最后一个等号是上面等同
$(B[1/g])_0 = \bigoplus_{m \geq 0} (A[1/f])_m$ 的自然延拓。
完全如先前对 $\pi : L \to P$ 所作的那样，断定
$\pi|_U : U \to P$ 仿射，并且作为 $P$ 上的概形，
$$
U = \underline{\Spec}_P
\left(
\bigoplus\nolimits_{m \in \mathbf{Z}} \mathcal{O}_P(m)
\right)
$$
成立。

\medskip\noindent
综上，我们的构造给出概形的交换图表
\begin{equation}
\label{more-morphisms-equation-proj-and-spec}
\vcenter{
\xymatrix{
\underline{\Spec}_P \left(
\bigoplus\nolimits_{m \in \mathbf{Z}} \mathcal{O}_P(m)
\right)
\ar[r] \ar@{=}[d] &
L =
\underline{\Spec}_P \left(
\bigoplus\nolimits_{m \geq 0} \mathcal{O}_P(m)
\right) \ar[d]^\sigma \ar[r]_-\pi & P \ar[d] \\
U \ar[r] & X \ar[r] & Z
}
}
\end{equation}
，其中 $\pi$ 是一个锥，其零截面 $0 : P \to L$ 在集合论意义下映到
$Z$ 在 $L$ 中的逆像上。

\medskip\noindent
令 $W \subset P$ 为满足下列条件的最大开集：$\mathcal{O}_P(1)|_W$
可逆，并且对所有 $m \in \mathbf{Z}$，自然映射诱导同构
$\mathcal{O}_P(m)|_W \cong \mathcal{O}_P(1)^{\otimes m}|_W$；
也就是说，它是 构造章引理
\ref{constructions-lemma-where-invertible} 中 $d = 1$ 所对应的开集。
于是可见，$L|_W = \pi^{-1}(W) \to W$ 是秩为 $1$ 的向量丛
（构造章第 \ref{constructions-section-vector-bundle} 节），即
$$
L|_W = \mathbf{V}(\mathcal{O}_P(1)|_W)
$$
，这里采用 Grothendieck 式记号。上文已表明 $L|_W$ 等于
$\mathcal{O}_P(1)|_W$ 在 $\mathcal{O}_W$ 上的对称代数的相对谱，
故这一点立即成立。显然，态射 $0|_W : W \to L|_W$ 是该向量丛的
零截面。特别地，$0(W)$ 是 $L|_W$ 上的有效 Cartier 除子。
此外，开集 $U|_W = (\pi|_U)^{-1}(W)$ 是零截面的补集。

\medskip\noindent
若 $A$ 由 $f_1, \ldots, f_r \in A_1$ 在 $A_0$ 上生成，则对所有
$m \geq 0$，有 $(f_1, \ldots, f_r)^m = A_{\geq m}$；因此上述 $B$
是 $A_+ = (f_1, \ldots, f_r)$ 的 Rees 代数。在此情形，
$L \to X$ 是沿 $Z$ 的爆破，并且上一段中的 $W$ 满足 $W = P$。

\medskip\noindent
若 $P$ 拟紧，则当 $d$ 充分可除时，由
$\sigma(A_d)\mathcal{O}_L$ 切出的闭子概形 $D \subset L$ 是有效
Cartier 除子，$0 : P \to L$ 经由 $D$ 分解，并且在集合论意义下
$0(P) = D$。这由 构造章引理
\ref{constructions-lemma-proj-quasi-compact} 以及上面证明的 (1)、(2)、
(3)、(4) 得到。（取任意一个能被该引理中所得元素的次数之最小公倍数整除的
$d$。）

\medskip\noindent
继续假设 $P$ 拟紧。设 $\mathcal{F}$ 是拟凝聚 $\mathcal{O}_P$-模，
并令 $\mathcal{F}_U = \pi^*\mathcal{F}|_U$。则有
\begin{equation}
\label{more-morphisms-equation-cohomology-torsor}
R\Gamma(U, \mathcal{F}_U) =
\bigoplus\nolimits_{m \in \mathbf{Z}}
R\Gamma(P, \mathcal{F} \otimes_{\mathcal{O}_P} \mathcal{O}_P(m))
\end{equation}
此外，该直和分解关于 $\mathcal{F}$ 是函子性的；右侧由此诱导的
$A$-模结构与左侧由 $U \subset X$ 给出的 $A$-模结构相同。
为证明该公式，由于 $\pi|_U$ 仿射，并且
$(\pi|_U)_*\mathcal{O}_U = \bigoplus_{m \in \mathbf{Z}} \mathcal{O}_P(m)$，
得到
\begin{align*}
R(\pi|_U)_*\mathcal{F}_U
& =
(\pi|_U)_*\mathcal{F}_U \\
& =
(\pi|_U)_*(\pi|_U)^*\mathcal{F} \\
& =
\mathcal{F} \otimes_{\mathcal{O}_P} 
\bigoplus\nolimits_{m \in \mathbf{Z}} \mathcal{O}_P(m) \\
& =
\bigoplus\nolimits_{m \in \mathbf{Z}}
\mathcal{F} \otimes_{\mathcal{O}_P} \mathcal{O}_P(m)
\end{align*}
由 Leray，得到
$R\Gamma(U, \mathcal{F}_U) =
R\Gamma(P, R(\pi|_U)_*\mathcal{F}_U)$, see
上同调章引理 \ref{cohomology-lemma-apply-Leray}。
在此情形，取上同调与直和交换，故证明完成；见
概形的导出范畴章引理
\ref{perfect-lemma-quasi-coherence-pushforward-direct-sums}。
此处正用到了 $P$ 拟紧；由 构造章引理
\ref{constructions-lemma-proj-separated}，$P$ 是分离的。

\begin{lemma}
\label{more-morphisms-lemma-apply-proj-spec}
设 $R$ 是环，$P$ 是 $R$ 上的固有概形，$\mathcal{L}$ 是充足可逆
$\mathcal{O}_P$-模。令
$A = \bigoplus_{m \geq 0} \Gamma(P, \mathcal{L}^{\otimes m})$。
则 $P = \text{Proj}(A)$，且图表 (\ref{more-morphisms-equation-proj-and-spec}) 变为
$$
\xymatrix{
\underline{\Spec}_P \left(
\bigoplus\nolimits_{m \in \mathbf{Z}} \mathcal{L}^{\otimes m}
\right)
\ar[r] \ar@{=}[d] &
L =
\underline{\Spec}_P \left(
\bigoplus\nolimits_{m \geq 0} \mathcal{L}^{\otimes m}
\right) \ar[d]^\sigma \ar[r]_-\pi & P \ar[d] \\
U \ar[r] & X \ar[r] & Z
}
$$
并具有上述性质。
\end{lemma}

\begin{proof}
由 态射章引理 \ref{morphisms-lemma-proper-ample-is-proj}，
有 $P = \text{Proj}(A)$。此外，由
性质章引理 \ref{properties-lemma-ample-gcd-is-one}，在该等同下，
对所有 $m \in \mathbf{Z}$ 均有
$\mathcal{O}_P(m) = \mathcal{L}^{\otimes m}$。
\end{proof}









\section{纤维中的闭点}
\label{more-morphisms-section-closed-points-fibres}

\noindent
本节部分材料取自预印本 \cite{Osserman-Payne}。

\begin{lemma}
\label{more-morphisms-lemma-locally-principal-vertical}
设 $f : X \to S$ 是概形态射，$Z \subset X$ 是闭子概形，
$s \in S$。假设
\begin{enumerate}
\item $S$ 不可约，其泛点为 $\eta$；
\item $X$ 不可约；
\item $f$ 优势；
\item $f$ 局部有限型；
\item $\dim(X_s) \leq \dim(X_\eta)$,
\item $Z$ 在 $X$ 中局部主；
\item $Z_\eta = \emptyset$.
\end{enumerate}
则在集合论意义下，纤维 $Z_s$ 是 $X_s$ 的若干不可约分支之并。
\end{lemma}

\begin{proof}
以 $X_{red}$ 表示 $X$ 的既约化。则 $Z \cap X_{red}$ 是
$X_{red}$ 的局部主闭子概形；见
除子章引理 \ref{divisors-lemma-pullback-locally-principal}。
因此可以假设 $X$ 既约。换言之，$X$ 是整概形；见
性质章引理 \ref{properties-lemma-characterize-integral}。
在此情形，态射 $X \to S$ 经由 $S_{red}$ 分解；见
概形章引理 \ref{schemes-lemma-map-into-reduction}。
所以可以用 $S_{red}$ 替换 $S$，并假设 $S$ 也是整概形。

\medskip\noindent
$f$ 为优势态射这一断言意味着 $X$ 的泛点映到 $\eta$；见 Morphisms,
引理 \ref{morphisms-lemma-dominant-finite-number-irreducible-components}。
此外，概形 $X_\eta$ 是域 $\kappa(\eta)$ 上局部有限型的整概形。
因此，对 $X_\eta$ 的每一点 $\xi$，都有
$d = \dim(X_\eta) \geq 0$ 等于 $\dim_\xi(X_\eta)$；见
代数章引理 \ref{algebra-lemma-dimension-spell-it-out} 及
\ref{algebra-lemma-dimension-at-a-point-finite-type-over-field}。
结合 态射章引理 \ref{morphisms-lemma-openness-bounded-dimension-fibres}
与条件 (5)，断定对每个 $x \in X_s$ 都有 $\dim_x(X_s) = d$。

\medskip\noindent
在 Noether 情形，该断言可如下证明。
% P05-EN-0110: frozen source has grammatical "does not holds"; Chinese follows intended syntax.
若引理不成立，则存在 $x \in Z_s$，它是 $Z_s$ 的某个不可约分支的泛点，
却不是 $X_s$ 的任何不可约分支的泛点。于是
$\dim_x(Z_s) \leq d - 1$，因为 $\dim_x(X_s) = d$，并且在 $X_s$
中 $x$ 的一个邻域内，闭子概形 $Z_s$ 不包含 $X_s$ 的任何不可约分支。
所以用 $x$ 的一个开邻域替换 $X$ 后，可以假设对所有 $z \in Z$ 都有
$\dim_z(Z_{f(z)}) \leq d - 1$；见
态射章引理 \ref{morphisms-lemma-openness-bounded-dimension-fibres}。
令 $\xi' \in Z$ 为 $Z$ 的某个不可约分支的泛点，并令
% P05-EN-0111: frozen source sets s'=f(xi) before xi is introduced; expected f(xi').
$s' = f(\xi)$。由于 $Z \not = X$ 局部主，可见
% P05-EN-0112: frozen source uses O_{X,xi}; the codimension-one point is evidently xi'.
$\dim(\mathcal{O}_{X, \xi}) = 1$；见
代数章引理 \ref{algebra-lemma-minimal-over-1}
（此处用到了 $X$ 为 Noether 概形）。令 $\xi \in X$ 为 $X$ 的泛点，
并令 $\xi_1$ 为 $X_{s'}$ 的任意一个包含 $\xi'$ 的不可约分支的泛点。
于是有特化
$$
\xi \leadsto \xi_1 \leadsto \xi'.
$$
由于 $\dim(\mathcal{O}_{X, \xi}) = 1$，两个特化之一必为等式。
按假设 $s' \not = \eta$，故第一个特化不是等式。因此
$\xi' = \xi_1$ 是 $X_{s'}$ 的某个不可约分支的泛点。
再次应用 态射章引理
\ref{morphisms-lemma-openness-bounded-dimension-fibres}，可得
$\dim_{\xi'}(Z_{s'}) = \dim_{\xi'}(X_{s'}) \geq \dim(X_\eta) = d$
，这给出所需矛盾。

\medskip\noindent
在一般情形，按如下方式化归到 Noether 情形。若引理不成立，则存在位于
$s$ 上方的一点 $x \in X$，$x$ 是 $Z_s$ 的某个不可约分支的泛点，
却不是 $X_s$ 的任何不可约分支的泛点。令 $U \subset S$ 为 $s$
的仿射邻域，并令 $V \subset X$ 为满足 $f(V) \subset U$ 的 $x$
的仿射邻域。写成 $U = \Spec(A)$、$V = \Spec(B)$，使得 $f|_V$
由环映射 $A \to B$ 给出。令 $\mathfrak q \subset B$、相应地令
$\mathfrak p \subset A$ 为分别对应于 $x$、$s$ 的素理想。
必要时缩小 $V$ 后，可假设 $Z \cap V$ 由某个元素 $g \in B$ 切出。
以 $K$ 表示 $A$ 的分式域。此时已知如下事实：
\begin{enumerate}
\item $A \subset B$ 是整环的有限生成扩张；
\item 元素 $g \otimes 1$ 在 $B \otimes_A K$ 中可逆；
\item $d = \dim(B \otimes_A K) = \dim(B \otimes_A \kappa(\mathfrak p))$,
\item $g \otimes 1$ 不是 $B \otimes_A \kappa(\mathfrak p)$ 的单位；
\item $g \otimes 1$ 不属于 $B \otimes_A \kappa(\mathfrak p)$ 的任何极小素理想。
\end{enumerate}
我们要导出矛盾。

\medskip\noindent
选取在 $A$ 上生成 $B$ 的元素 $x_1, \ldots, x_n \in B$。
对有限生成的 $\mathbf{Z}$-代数 $A_0 \subset A$，令
$B_0 \subset B$ 为由 $x_1, \ldots, x_n$ 生成的 $A_0$-子代数，
以 $K_0$ 表示 $A_0$ 的分式域，并令
$\mathfrak p_0 = A_0 \cap \mathfrak p$。我们断言，当 $A_0$ 充分大时，
(1) -- (5) 对系统 $(A_0 \subset B_0, g, \mathfrak p_0)$ 也成立。

\medskip\noindent
逐一证明这些条件。第 (1) 部分按构造成立。对第 (2) 部分，
对某个 $h \otimes 1/a \in B \otimes_A K$ 写成
$(g \otimes 1) h = 1$。对某些系数 $a_I, a'_I \in A$，写成
$g = \sum a_I x^I$、$h = \sum a'_I x^I$（多重指标记号）。
只要 $A_0$ 包含 $a$ 及所有 $a_I, a'_I$，则 (2) 成立，因为
$B_0 \otimes_{A_0} K_0 \subset B \otimes_A K$
（它们是单射 $B_0 \to B$ 的局部化）。为得到 (3)，考虑正合列
$$
0 \to I \to A[X_1, \ldots, X_n] \to B \to 0
$$
它定义 $I$，其中第二个映射将 $X_i$ 映到 $x_i$。由于 $\otimes$
右正合，可见 $I \otimes_A K$、相应地
$I \otimes_A \kappa(\mathfrak p)$ 分别是满射
$K[X_1, \ldots, X_n] \to B \otimes_A K$、相应地
$\kappa(\mathfrak p)[X_1, \ldots, X_n] \to B \otimes_A \kappa(\mathfrak p)$
的核。由于域上的多项式环是 Noether 环，存在有限多个元素
$h_j \in I$，$j = 1, \ldots, m$，它们生成
$I \otimes_A K$ 和 $I \otimes_A \kappa(\mathfrak p)$。
写成 $h_j = \sum a_{j, I}X^I$。只要 $A_0$ 包含所有 $a_{j, I}$，
便得到
$$
B_0 \otimes_{A_0} K_0 \otimes_{K_0} K = B \otimes_A K
\quad\text{且}\quad
B_0 \otimes_{A_0} \kappa(\mathfrak p_0)
\otimes_{\kappa(\mathfrak p_0)} \kappa(\mathfrak p)
=
B \otimes_A \kappa(\mathfrak p).
$$
由 态射章引理 \ref{morphisms-lemma-dimension-fibre-after-base-change}
或 代数章引理 \ref{algebra-lemma-dimension-preserved-field-extension}，
可见 (3) 中的维数等式成立。第 (4) 部分立即成立。由于
$B_0 \otimes_{A_0} \kappa(\mathfrak p_0) \subset
B \otimes_A \kappa(\mathfrak p)$，
$B_0 \otimes_{A_0} \kappa(\mathfrak p_0)$ 的每个极小素理想，都位于
$B \otimes_A \kappa(\mathfrak p)$ 的某个极小素理想之下；见
代数章引理 \ref{algebra-lemma-image-dense-generic-points}。
这蕴含 (5) 成立。由此将问题化归到上面已经处理的 Noether 情形。
\end{proof}

\noindent
下面是上述引理的一个代数应用。若 $A \to B$ 平坦，则该引理的第四个假设
成立；见 引理 \ref{more-morphisms-lemma-equality-dimensions}。

\begin{lemma}
\label{more-morphisms-lemma-horizontal}
设 $A \to B$ 是局部环的局部同态，且 $g \in \mathfrak m_B$。假设
\begin{enumerate}
\item $A$ 与 $B$ 是整环，且 $A \subset B$；
\item $B$ 在 $A$ 上本质有限型；
\item $g$ 不属于 $\mathfrak m_AB$ 上方的任何极小素理想；
\item $\dim(B/\mathfrak m_AB) +
\text{trdeg}_{\kappa(\mathfrak m_A)}(\kappa(\mathfrak m_B)) =
\text{trdeg}_A(B)$.
\end{enumerate}
则 $A \subset B/gB$，即 $\Spec(A)$ 的泛点属于态射
$\Spec(B/gB) \to \Spec(A)$ 的像。
\end{lemma}

\begin{proof}
由 代数章引理 \ref{algebra-lemma-image-dense-generic-points}，
注意这两个断言等价。为开始证明，令 $C$ 为有限型 $A$-代数，
$\mathfrak q$ 为 $C$ 的素理想，使得 $B = C_{\mathfrak q}$。
当然可以假设 $C$ 是整环且 $g \in C$。用一个局部化替换 $C$ 后，有
$\dim(C/\mathfrak m_AC) = \dim(B/\mathfrak m_AB) +
\text{trdeg}_{\kappa(\mathfrak m_A)}(\kappa(\mathfrak m_B))$；见
态射章引理 \ref{morphisms-lemma-dimension-fibre-at-a-point}。
令 $K$ 等于 $A$ 的分式域，由同一引理可见
$\dim(C \otimes_A K) = \text{trdeg}_A(B)$。因此，假设 (4) 意味着态射
$\Spec(C) \to \Spec(A)$ 的泛纤维与闭纤维具有相同维数。

\medskip\noindent
假设引理不成立。则 $(B/gB) \otimes_A K = 0$。这意味着
$g \otimes 1$ 在 $B \otimes_A K
= C_{\mathfrak q} \otimes_A K$
中可逆。由于 $C_{\mathfrak q}$ 是主局部化的极限，断定对某个
$h \in C$、$h \not \in \mathfrak q$，$g \otimes 1$ 在
$C_h \otimes_A K$ 中可逆。因此，用 $C_h$ 替换 $C$ 后，可以假设
$(C/gC) \otimes_A K = 0$。再替换一次 $C$，以保证
$C/\mathfrak m_AC$ 的极小素理想与 $B/\mathfrak m_AB$ 的极小素理想
一一对应。此时，将 引理 \ref{more-morphisms-lemma-locally-principal-vertical}
应用于 $X = \Spec(C) \to \Spec(A) = S$ 以及局部闭子概形
$Z = \Spec(C/gC)$。由于 $Z_K = \emptyset$，可见
$Z \otimes \kappa(\mathfrak m_A)$ 必须包含
$X \otimes \kappa(\mathfrak m_A) = \Spec(C/\mathfrak m_AC)$ 的一个
不可约分支。但这与 $g$ 不属于 $\mathfrak m_AB$ 上方任何极小素理想的
假设矛盾。引理得证。
\end{proof}

\begin{lemma}
\label{more-morphisms-lemma-equality-dimensions}
设 $A \to B$ 是局部环的局部同态。假设
\begin{enumerate}
\item $A$ 与 $B$ 是整环，且 $A \subset B$；
\item $B$ 在 $A$ 上本质有限型；
\item $B$ 在 $A$ 上平坦。
\end{enumerate}
则有
$$
\dim(B/\mathfrak m_AB) +
\text{trdeg}_{\kappa(\mathfrak m_A)}(\kappa(\mathfrak m_B)) =
\text{trdeg}_A(B).
$$
\end{lemma}

\begin{proof}
令 $C$ 为有限型 $A$-代数，$\mathfrak q$ 为 $C$ 的素理想，使得
$B = C_{\mathfrak q}$。可以假设 $C$ 是整环。由 态射章引理
\ref{morphisms-lemma-dimension-fibre-at-a-point}，有
$\dim_{\mathfrak q}(C/\mathfrak m_AC) = \dim(B/\mathfrak m_AB) +
\text{trdeg}_{\kappa(\mathfrak m_A)}(\kappa(\mathfrak m_B))$。
令 $K$ 等于 $A$ 的分式域，由同一引理可见
$\dim(C \otimes_A K) = \text{trdeg}_A(B)$。
因此，实际要证明的是
$\dim_{\mathfrak q}(C/\mathfrak m_AC) = \dim(C \otimes_A K)$。
选取 $K$ 中支配 $A$ 的赋值环 $A'$；见
代数章引理 \ref{algebra-lemma-dominate}。令 $C' = C \otimes_A A'$。
选取 $C'$ 中位于 $\mathfrak q$ 上方的素理想 $\mathfrak q'$；
这样的素理想存在，因为
$$
C'/\mathfrak m_{A'}C' =
C/\mathfrak m_AC \otimes_{\kappa(\mathfrak m_A)} \kappa(\mathfrak m_{A'})
$$
这还证明 $C/\mathfrak m_AC \to C'/\mathfrak m_{A'}C'$ 忠实平坦。
此外，这证明
$\dim_{\mathfrak q}(C/\mathfrak m_AC) =
\dim_{\mathfrak q'}(C'/\mathfrak m_{A'}C')$；见
代数章引理 \ref{algebra-lemma-dimension-at-a-point-preserved-field-extension}。
注意，$B' = C'_{\mathfrak q'}$ 是 $B \otimes_A A'$ 的一个局部化。
因此 $B'$ 在 $A'$ 上平坦。泛纤维 $B' \otimes_{A'} K$ 是
$B \otimes_A K$ 的一个局部化，故 $B'$ 是整环。若对
$A' \subset B'$ 证明本引理，则得到等式
$\dim_{\mathfrak q'}(C'/\mathfrak m_{A'}C') = \dim(C' \otimes_{A'} K)$
，结合上述讨论，它蕴含所需等式
$\dim_{\mathfrak q}(C/\mathfrak m_AC) = \dim(C \otimes_A K)$
。由此将本引理化归到 $A$ 为赋值环的情形。

\medskip\noindent
设 $A \subset B$ 如引理所述，且 $A$ 是赋值环。像先前一样，
对某个有限型 $A$-整代数 $C$ 写成 $B = C_{\mathfrak q}$。
由 代数章引理
\ref{algebra-lemma-finite-type-domain-over-valuation-ring-dim-fibres}，
得到 $\dim(C/\mathfrak m_AC) = \dim(C \otimes_A K)$，结论得证。
\end{proof}

\begin{lemma}
\label{more-morphisms-lemma-closed-point-nearby-fibre}
设 $f : X \to S$ 是概形态射，$x \leadsto x'$ 是 $X$ 中点的特化。
令 $s = f(x)$ 且 $s' = f(x')$。假设
\begin{enumerate}
\item $x'$ 是 $X_{s'}$ 的闭点；
\item $f$ 局部有限型。
\end{enumerate}
则集合
$$
\{x_1 \in X
\text{ 使得 }
f(x_1) = s
\text{ 且 }
x_1\text{ is closed in }X_s
\text{ 且 }
x \leadsto x_1 \leadsto x'
\}
$$
在 $x$ 于 $X_s$ 中的闭包内稠密。
\end{lemma}

\begin{proof}
将 概形章引理 \ref{schemes-lemma-points-specialize}
应用于特化 $x \leadsto x'$。它给出态射
$\varphi : \Spec(B) \to X$，其中 $B$ 是赋值环，且 $\varphi$
将泛点映到 $x$、闭点映到 $x'$。还可以假设 $\kappa(x)$ 是 $B$
的分式域。令 $A = B \cap \kappa(s)$。注意，这是一个赋值环
（见 代数章引理 \ref{algebra-lemma-valuation-ring-cap-field}），
它支配 $\mathcal{O}_{S, s'} \to \kappa(s)$ 的像。考虑交换图表
$$
\xymatrix{
\Spec(B) \ar[rd] \ar[r] &
X_A \ar[d] \ar[r] & X \ar[d] \\
& \Spec(A) \ar[r] & S
}
$$
$B$ 的泛点（相应地，闭点）映到 $X_A$ 中位于 $\Spec(A)$ 泛点
（相应地，闭点）上方的一点 $x_A$（相应地，$x'_A$）。由 Morphisms,
引理 \ref{morphisms-lemma-base-change-closed-point-fibre-locally-finite-type}，
$x'_A$ 是 $X_A$ 的特殊纤维的闭点。注意，$X_A \to \Spec(A)$ 的泛纤维
同构于 $X_s$。因此，已经将引理化归到如下情形：$S$ 是一个赋值环的谱，
$s = \eta \in S$ 是泛点，而 $s' \in S$ 是闭点。

\medskip\noindent
对 $\dim_x(X_\eta)$ 归纳证明本引理。若 $\dim_x(X_\eta) = 0$，
则 $X_\eta$ 中没有其他点特化到 $x$，且 $x$ 在其纤维中闭；见
态射章引理 \ref{morphisms-lemma-quasi-finite-at-point-characterize}。
此时结论成立。以下假设 $\dim_x(X_\eta) > 0$。

\medskip\noindent
令 $X' \subset X$ 为 $X$ 的不可约闭子概形 $\overline{\{x\}}$
上赋予的既约诱导概形结构；见
概形章定义 \ref{schemes-definition-reduced-induced-scheme}。
为证明本引理，可以用 $X'$ 替换 $X$，因为这只会减小
$\dim_x(X_\eta)$。所以还可以假设 $X$ 是整概形，而 $x$ 是其泛点。
此外，可以用 $x'$ 的一个仿射邻域替换 $X$。于是
$X = \Spec(B)$，其中 $A \subset B$ 是整环的有限型扩张。
注意，在此情形
$\dim_x(X_\eta) = \dim(X_\eta) = \dim(X_{s'})$，并且事实上
$X_{s'}$ 等维；见 代数章引理
\ref{algebra-lemma-finite-type-domain-over-valuation-ring-dim-fibres}。

\medskip\noindent
令 $W \subset X_\eta$ 为真闭子集（这是我们想要“避开”的子集）。
由于 $X_s$ 在一个域上有限型，可见 $W$ 只有有限多个不可约分支，写成
$W = W_1 \cup \ldots \cup W_n$。令
% P05-EN-0113: frozen source indexes the n corresponding components by j=1,...,r.
$\mathfrak q_j \subset B$，$j = 1, \ldots, r$，为相应的素理想。
令 $\mathfrak q \subset B$ 为对应于点 $x'$ 的极大理想。
令 $\mathfrak p_1, \ldots, \mathfrak p_s \subset B$ 为位于
$\mathfrak m_AB$ 上方的极小素理想。这样的素理想只有有限多个，
因为它们对应于 Noether 概形 $X_{s'}$ 的不可约分支。此外，这些不可约分支
的维数都为 $> 0$（见上文），故对所有 $i$ 都有
$\mathfrak p_i \not = \mathfrak q$。现在选取元素
$g \in \mathfrak q$，使得对所有 $j$ 有 $g \not \in \mathfrak q_j$，
且对所有 $i$ 有 $g \not \in \mathfrak p_i$；见
代数章引理 \ref{algebra-lemma-silly}。
以 $Z \subset X$ 表示由 $g$ 定义的局部主闭子概形。令
$Z_\eta = Z_{1, \eta} \cup \ldots \cup Z_{n, \eta}$，$n \geq 0$，
为 $Z$ 的泛纤维的不可约分支分解（分支有限，因为泛纤维是 Noether 的）。
以 $Z_i \subset X$ 表示 $Z_{i, \eta}$ 的闭包。用一个更小的仿射邻域替换
$X$ 后，可以假设对每个 $i = 1, \ldots, n$ 都有 $x' \in Z_i$。
按构造，$Z \cap X_{s'}$ 不包含 $X_{s'}$ 的任何不可约分支。
因此，由 引理 \ref{more-morphisms-lemma-locally-principal-vertical}，断定
$Z_\eta \not = \emptyset$！换言之，$n \geq 1$。令
$x_1 \in Z_1$ 为泛点，则 $x_1 \leadsto x'$ 且 $f(x_1) = \eta$。
此外，按构造，$Z_{1, \eta} \cap W_j \subset W_j$ 是真闭子集。
所以 $Z_{1, \eta} \cap W_j$ 的每个不可约分支在 $X_\eta$ 中的余维
均为 $\geq 2$，而由 代数章引理 \ref{algebra-lemma-minimal-over-1}，
$\text{codim}(Z_{1, \eta}, X_\eta) = 1$。因此
$W \cap Z_{1, \eta}$ 是真闭子集。此时，归纳假设可应用于
$Z_1 \to S$ 及特化 $x_1 \leadsto x'$。由此得到
$Z_{1, \eta}$ 的一个不属于 $W$ 且特化到 $x'$ 的闭点 $x_2$。
于是得到 $x \leadsto x_2 \leadsto x'$；点 $x_2$ 在 $X_\eta$ 中闭，
且如所需 $x_2 \not \in W$。
\end{proof}

\begin{remark}
\label{more-morphisms-remark-full-specialization-sequence}
引理 \ref{more-morphisms-lemma-closed-point-nearby-fibre} 的证明实际上表明，
存在特化序列
$$
x \leadsto x_1 \leadsto x_2 \leadsto \ldots \leadsto x_d \leadsto x'
$$
其中所有 $x_i$ 都属于纤维 $X_s$，每个特化都是直接的，而
$x_d$ 是 $X_s$ 的闭点。整数
$d = \text{trdeg}_{\kappa(s)}(\kappa(x)) = \dim(\overline{\{x\}})$
，其中闭包取于 $X_s$ 中。此外，可以选取这些点 $x_i$，使其避开
$X_s$ 中任何不包含点 $x$ 的闭子集。
\end{remark}

\noindent
Examples, 第 \ref{examples-section-topology-finite-type} 节 表明，
若不假设 $A$ 为 Noether 环，则下述引理不成立。

\begin{lemma}
\label{more-morphisms-lemma-quasi-finite-quasi-section-meeting-nearby-open}
设 $\varphi : A \to B$ 是局部环的局部环映射。设
$V \subset \Spec(B)$ 是一个开子概形，且至少包含一个不位于
$\mathfrak m_A$ 上方的素理想。假设 $A$ 是 Noether 环，$\varphi$
本质有限型，并且 $A/\mathfrak m_A \subset B/\mathfrak m_B$ 有限。
则存在 $\mathfrak q \in V$，满足
$\mathfrak m_A \not = \mathfrak q \cap A$，使得
$A \to B/\mathfrak q$ 是某个拟有限环映射的局部化。
\end{lemma}

\begin{proof}
由于 $A$ 是 Noether 环，而 $A \to B$ 本质有限型，可知 $B$ 也是
Noether 环。由 性质章引理
\ref{properties-lemma-complement-closed-point-Jacobson}，拓扑空间
$\Spec(B) \setminus \{\mathfrak m_B\}$ 是 Jacobson 的。因此可以在非空开集
$$
V \setminus \{\mathfrak q \subset B \mid \mathfrak m_A = \mathfrak q \cap A\}.
$$
中选取闭点 $\mathfrak q$。（该集按假设非空；它是开的，因为
$\{\mathfrak m_A\}$ 是 $\Spec(A)$ 的闭子集。）
于是 $\Spec(B/\mathfrak q)$ 有两个点，即 $\mathfrak m_B$ 和
$\mathfrak q$，而 $\mathfrak q$ 不位于 $\mathfrak m_A$ 上方。
对某个有限型 $A$-代数 $C$ 及素理想 $\mathfrak m$，写成
$B/\mathfrak q = C_{\mathfrak m}$。由
代数章引理 \ref{algebra-lemma-isolated-point-fibre} (2)，
$A \to C$ 在 $\mathfrak m$ 处拟有限。因此，由
代数章引理 \ref{algebra-lemma-quasi-finite-open}，
用一个主局部化替换 $C$ 后，环映射 $A \to C$ 拟有限。
\end{proof}

\begin{lemma}
\label{more-morphisms-lemma-quasi-finite-quasi-section-meeting-nearby-open-X}
设 $f : X \to S$ 是概形态射。取 $x \in X$，其像为 $s \in S$。
设 $U \subset X$ 是开子概形。假设 $f$ 局部有限型，$S$ 局部
Noether，$x$ 是 $X_s$ 的闭点，并假设存在一点 $x' \in U$，满足
$x' \leadsto x$ 且 $f(x') \not = s$。则存在闭子概形 $Z \subset X$，
使得 (a) $x \in Z$；(b) $f|_Z : Z \to S$ 在 $x$ 处拟有限；
(c) 存在 $z \in Z$，$z \in U$、$z \leadsto x$ 且 $f(z) \not = s$。
\end{lemma}

\begin{proof}
这是 引理 \ref{more-morphisms-lemma-quasi-finite-quasi-section-meeting-nearby-open}
的改述。具体地，令 $A = \mathcal{O}_{S, s}$、
$B = \mathcal{O}_{X, x}$。以 $V \subset \Spec(B)$ 表示 $U$ 的逆像。
环映射 $f^\sharp : A \to B$ 本质有限型。按假设，$V$ 中至少有一点
不映到 $\Spec(A)$ 的闭点。因此，
引理 \ref{more-morphisms-lemma-quasi-finite-quasi-section-meeting-nearby-open}
的所有假设均成立，故得到素理想 $\mathfrak q \subset B$，它不位于
$\mathfrak m_A$ 上方，并且 $A \to B/\mathfrak q$ 是某个拟有限环映射的
局部化。令 $z \in X$ 为点 $\mathfrak q$ 在典范态射
$\Spec(B) \to X$ 下的像。令 $Z = \overline{\{z\}}$，并赋予诱导既约
概形结构。由于 $z \leadsto x$，可见 $x \in Z$ 且
$\mathcal{O}_{Z, x} = B/\mathfrak q$。按构造，$Z \to S$ 在 $x$ 处拟有限。
\end{proof}

\begin{remark}
\label{more-morphisms-remark-alternative-closed-point-nearby-fibre}
当 $S$ 局部 Noether 时，可以利用
引理 \ref{more-morphisms-lemma-quasi-finite-quasi-section-meeting-nearby-open}
或其变体
引理 \ref{more-morphisms-lemma-quasi-finite-quasi-section-meeting-nearby-open-X}，
给出 引理 \ref{more-morphisms-lemma-closed-point-nearby-fibre} 的另一证明。
下面是一个粗略纲要。首先，用 $s'$ 处局部环的谱替换 $S$，然后对
$\dim(S)$ 归纳。当 $\dim(S) = 0$ 时，由于此时 $s' = s$，结论平凡。
用 $\overline{\{x\}}$ 上的既约诱导概形结构替换 $X$。
将 引理 \ref{more-morphisms-lemma-quasi-finite-quasi-section-meeting-nearby-open-X}
应用于 $X \to S$、$x' \mapsto s'$ 以及任意含有 $x$ 的非空开集
$U \subset X$。由此得到闭子概形
% P05-EN-0114: frozen source omits conjunction/punctuation between x' in Z subset X and a point z in Z.
$x' \in Z \subset X$ 及一点 $z \in Z$，使得 $Z \to S$ 在 $x'$ 处
拟有限，并且 $f(z) \not = s'$。于是 $z$ 是 $X_{f(z)}$ 的闭点，
且 $z \leadsto x'$。由于 $f(z) \not = s'$，可见
$\dim(\mathcal{O}_{S, f(z)}) < \dim(S)$。由于 $x$ 是 $X$ 的泛点，
有 $x \leadsto z$，因而 $s = f(x) \leadsto f(z)$。
将归纳假设应用于 $s \leadsto f(z)$ 及 $z \mapsto f(z)$，即得结论。
\end{remark}

\begin{lemma}
\label{more-morphisms-lemma-change-hypotheses}
假设 $f : X \to S$ 局部有限型，$S$ 局部 Noether，$x \in X$ 是其纤维
$X_s$ 的闭点，而 $U \subset X$ 是满足 $U \cap X_s = \emptyset$
及 $x \in \overline{U}$ 的开子概形。则
引理 \ref{more-morphisms-lemma-quasi-finite-quasi-section-meeting-nearby-open-X}
的结论成立。
\end{lemma}

\begin{proof}
可以按如下方式化归到所引引理：首先分别用 $x$ 与 $s$ 的仿射邻域替换
$X$ 与 $S$。于是 $X$ 是 Noether 的，特别地，$U$ 拟紧（见
态射章引理 \ref{morphisms-lemma-finite-type-noetherian} 及
拓扑章引理 \ref{topology-lemma-Noetherian} 及
\ref{topology-lemma-Noetherian-quasi-compact}）。
因此存在满足 $x' \in U$ 的特化 $x' \leadsto x$（见
态射章引理 \ref{morphisms-lemma-reach-points-scheme-theoretic-image}）。
注意，$f(x') \not = s$。由此可见，该引理的全部假设均成立，结论得证。
\end{proof}



\section{Stein 分解}
\label{more-morphisms-section-stein-factorization}

\noindent
Stein 分解断言：若固有态射 $f : X \to S$ 满足
$f_*\mathcal{O}_X = \mathcal{O}_S$，则它具有连通纤维。

\begin{lemma}
\label{more-morphisms-lemma-stein-universally-closed}
设 $S$ 是概形，$f : X \to S$ 是普遍闭且拟分离的态射。
则存在分解
$$
\xymatrix{
X \ar[rr]_{f'} \ar[rd]_f & & S' \ar[dl]^\pi \\
& S &
}
$$
具有下列性质：
\begin{enumerate}
\item 态射 $f'$ 普遍闭、拟紧、拟分离且满；
\item 态射 $\pi : S' \to S$ 是整态射；
\item 有 $f'_*\mathcal{O}_X = \mathcal{O}_{S'}$；
\item 有 $S' = \underline{\Spec}_S(f_*\mathcal{O}_X)$；
\item $S'$ 是 $S$ 在 $X$ 中的正规化；见
态射章定义 \ref{morphisms-definition-normalization-X-in-Y}.
\end{enumerate}
分解 $f = \pi \circ f'$ 的形成与平坦基变换交换。
\end{lemma}

\begin{proof}
由 态射章引理 \ref{morphisms-lemma-universally-closed-quasi-compact}，
态射 $f$ 拟紧。因此，$S$ 在 $X$ 中的正规化 $S'$ 有定义
（态射章定义 \ref{morphisms-definition-normalization-X-in-Y}），
并且有分解 $X \to S' \to S$。由 态射章引理
\ref{morphisms-lemma-normalization-in-universally-closed}，得到 (2)、(4)、
(5)。由 态射章引理 \ref{morphisms-lemma-image-proper-scheme-closed}，
态射 $f'$ 普遍闭。由 概形章引理
\ref{schemes-lemma-quasi-compact-permanence}，它拟紧；由
概形章引理 \ref{schemes-lemma-compose-after-separated}，它拟分离。

\medskip\noindent
为证明其余陈述，可以假设基概形 $S$ 仿射，写成 $S = \Spec(R)$。
于是 $S' = \Spec(A)$，其中 $A = \Gamma(X, \mathcal{O}_X)$ 是整
$R$-代数。因而 $f'_*\mathcal{O}_X$ 显然就是 $\mathcal{O}_{S'}$
（由 概形章引理 \ref{schemes-lemma-push-forward-quasi-coherent}，
$f'_*\mathcal{O}_X$ 拟凝聚，故等于 $\widetilde{A}$）。这证明了 (3)。

\medskip\noindent
下面证明 $f'$ 是满射。由于 $f'$ 普遍闭（见上文），$f'$ 的像是闭子集
$V(I) \subset S' = \Spec(A)$。取 $h \in I$。则
$h|_X = f^\sharp(h)$ 是 $X$ 的结构层的一个整体截面，并在每一点消失。
由于 $X$ 拟紧，这意味着 $h|_X$ 是幂零截面，即对某个 $n > 0$ 有
% P05-EN-0115: frozen source writes h^n|X without the expected restriction underscore.
$h^n|X = 0$。但 $A = \Gamma(X, \mathcal{O}_X)$，故 $h^n = 0$。
由于 $I$ 的每个元素都幂零，断定 $V(I) = S'$，正如所需。

\medskip\noindent
由 概形上同调章引理
\ref{coherent-lemma-flat-base-change-cohomology}，$f_*\mathcal{O}_X$ 的形成
与平坦基变换交换。由 构造章引理
\ref{constructions-lemma-spec-properties}，相对谱的形成与任意基变换交换。
因此，该分解的形成与平坦基变换交换。
\end{proof}

\begin{lemma}
\label{more-morphisms-lemma-stein-universally-closed-residue-fields}
在 引理 \ref{more-morphisms-lemma-stein-universally-closed} 中进一步假设 $f$ 局部有限型。
则对 $s \in S$，纤维 $\pi^{-1}(\{s\}) = \{s_1, \ldots, s_n\}$
是有限的，并且域扩张 $\kappa(s_i)/\kappa(s)$ 有限。
\end{lemma}

\begin{proof}
回忆，$\pi^{-1}(\{s\})$ 的点之间不存在特化；见
代数章引理 \ref{algebra-lemma-integral-no-inclusion}。由于 $f'$ 满，
可知 $|X_s| \to \pi^{-1}(\{s\})$ 满。注意，$X_s$ 是域上的有限型
拟分离概形（拟紧性已在所引引理的证明中说明）。因此 $X_s$ 是
Noether 的（态射章引理 \ref{morphisms-lemma-finite-type-noetherian}）。
现在，一个拓扑论证（略去）表明 $\pi^{-1}(\{s\})$ 有限。
对每个 $i$，可以选取映到 $s_i$ 的有限型点 $x_i \in X_s$
（态射章引理 \ref{morphisms-lemma-enough-finite-type-points}）。
由此断定 $\kappa(s_i)/\kappa(s)$ 有限：$x_i$ 可由有限型态射
$\Spec(k_i) \to X_s$ 表示（按有限型点的定义），因此
$\Spec(k_i) \to s = \Spec(\kappa(s))$ 是有限型态射
（作为有限型态射的复合），故 $k_i/\kappa(s)$ 有限
（态射章引理 \ref{morphisms-lemma-point-finite-type}）。
\end{proof}

\begin{lemma}
\label{more-morphisms-lemma-characterize-geometrically-connected-fibres}
设 $f : X \to S$ 是概形态射，$s \in S$。则 $X_s$ 几何连通，
当且仅当对每个 \'etale 邻域 $(U, u) \to (S, s)$，基变换
$X_U \to U$ 的纤维 $X_u$ 连通。
\end{lemma}

\begin{proof}
若 $X_s$ 几何连通，则它的任意基变换都连通。反之，假设 $X_s$
不几何连通。由 簇章引理
\ref{varieties-lemma-characterize-geometrically-disconnected}，
对某个有限可分域扩张 $k/\kappa(s)$，
$X_s \times_{\Spec(\kappa(s))} \Spec(k)$ 不连通。
由 引理 \ref{more-morphisms-lemma-realize-prescribed-residue-field-extension-etale}，
存在仿射 \'etale 邻域 $(U, u) \to (S, s)$，使得
$\kappa(u)/\kappa(s)$ 与 $k/\kappa(s)$ 等同。在此情形，$X_u$ 不连通。
\end{proof}

\begin{theorem}[Stein 分解；Noether 情形]
\label{more-morphisms-theorem-stein-factorization-Noetherian}
设 $S$ 是局部 Noether 概形，$f : X \to S$ 是固有态射。
则存在分解
$$
\xymatrix{
X \ar[rr]_{f'} \ar[rd]_f & & S' \ar[dl]^\pi \\
& S &
}
$$
具有下列性质：
\begin{enumerate}
\item 态射 $f'$ 固有，并具有几何连通纤维；
\item 态射 $\pi : S' \to S$ 有限；
\item 有 $f'_*\mathcal{O}_X = \mathcal{O}_{S'}$；
\item 有 $S' = \underline{\Spec}_S(f_*\mathcal{O}_X)$；
\item $S'$ 是 $S$ 在 $X$ 中的正规化；见
态射章定义 \ref{morphisms-definition-normalization-X-in-Y}.
\end{enumerate}
\end{theorem}

\begin{proof}
令 $f = \pi \circ f'$ 为 引理 \ref{more-morphisms-lemma-stein-universally-closed}
中的分解。注意，除 引理 \ref{more-morphisms-lemma-stein-universally-closed} 的结论外，
$f'$ 还是分离的
（概形章引理 \ref{schemes-lemma-compose-after-separated}）且有限型
（态射章引理 \ref{morphisms-lemma-permanence-finite-type}）。
因此 $f'$ 固有。由 概形上同调章命题
\ref{coherent-proposition-proper-pushforward-coherent}，
$f_*\mathcal{O}_X$ 是凝聚 $\mathcal{O}_S$-模。因此 $\pi$ 有限，
即 (2) 成立。

\medskip\noindent
这证明了除最有趣的断言外的所有结论；剩下要证的是 $f'$ 的所有纤维均
几何连通。由上述讨论，显然可以用 $S'$ 替换 $S$，因而可假设 $S$
是 Noether 仿射概形，$f : X \to S$ 固有，并且
$f_*\mathcal{O}_X = \mathcal{O}_S$。取 $S$ 的一点 $s \in S$。
需要证明 $X_s$ 几何连通。由
引理 \ref{more-morphisms-lemma-characterize-geometrically-connected-fibres}，
只需对每个 \'etale 邻域 $(U, u) \to (S, s)$ 证明 $X_u$ 连通。
可以假设 $U$ 仿射。于是 $U$ 是 Noether 的
（态射章引理 \ref{morphisms-lemma-finite-type-noetherian}），
基变换 $f_U : X_U \to U$ 固有
（态射章引理 \ref{morphisms-lemma-base-change-proper}），并且
$(f_U)_*\mathcal{O}_{X_U} = \mathcal{O}_U$
（概形上同调章引理
\ref{coherent-lemma-flat-base-change-cohomology}）。
因此，用基变换 $(f_U : X_U \to U, u)$ 替换
$(f : X \to S, s)$ 后，只需在 $f_*\mathcal{O}_X = \mathcal{O}_S$
时证明纤维 $X_s$ 连通。可以从 概形的导出范畴章引理
\ref{perfect-lemma-proper-idempotent-on-fibre} 推出这一点
（考察 $X_s$ 结构层中的幂等元），不过下面还将给出直接论证。

\medskip\noindent
具体地，应用形式函数定理，更确切地说，应用 概形上同调章引理
\ref{coherent-lemma-formal-functions-stalk}。它给出
$$
\mathcal{O}^\wedge_{S, s} = (f_*\mathcal{O}_X)_s^\wedge =
\lim_n H^0(X_n, \mathcal{O}_{X_n})
$$
其中 $X_n$ 是 $X_s$ 的第 $n$ 个无穷小邻域。由于 $X_n$ 的底层拓扑空间
等于 $X_s$ 的底层拓扑空间，若 $X_s = T_1 \amalg T_2$ 是非空开闭子概形的
不交并，则对所有 $n$，同样有
$X_n = T_{1, n} \amalg T_{2, n}$。这又意味着
$H^0(X_n, \mathcal{O}_{X_n})$ 含有非平凡幂等元 $e_{1, n}$，即在
$T_{1, n}$ 上恒等于 $1$、在 $T_{2, n}$ 上恒等于 $0$ 的函数。
显然，$e_{1, n + 1}$ 在 $X_n$ 上限制为 $e_{1, n}$。所以
$e_1 = \lim e_{1, n}$ 是该极限的非平凡幂等元。这与
$\mathcal{O}^\wedge_{S, s}$ 为局部环矛盾。因此假设错误，即 $X_s$
连通，结论得证。
\end{proof}

\begin{theorem}[Stein 分解；一般情形]
\label{more-morphisms-theorem-stein-factorization-general}
设 $S$ 是概形，$f : X \to S$ 是固有态射。则存在分解
$$
\xymatrix{
X \ar[rr]_{f'} \ar[rd]_f & & S' \ar[dl]^\pi \\
& S &
}
$$
具有下列性质：
\begin{enumerate}
\item 态射 $f'$ 固有，并具有几何连通纤维；
\item 态射 $\pi : S' \to S$ 是整态射；
\item 有 $f'_*\mathcal{O}_X = \mathcal{O}_{S'}$；
\item 有 $S' = \underline{\Spec}_S(f_*\mathcal{O}_X)$；
\item $S'$ 是 $S$ 在 $X$ 中的正规化；见
态射章定义 \ref{morphisms-definition-normalization-X-in-Y}.
\end{enumerate}
\end{theorem}

\begin{proof}
应用 引理 \ref{more-morphisms-lemma-stein-universally-closed}，得到态射
$f' : X \to S'$。注意，除 引理 \ref{more-morphisms-lemma-stein-universally-closed}
的结论外，$f'$ 还是分离的
（概形章引理 \ref{schemes-lemma-compose-after-separated}）且有限型
（态射章引理 \ref{morphisms-lemma-permanence-finite-type}）。
因此 $f'$ 固有。至此，除 $f'$ 具有几何连通纤维这一陈述外，
已经证明了所有陈述。

\medskip\noindent
可以假设 $S = \Spec(R)$ 仿射。令
$R' = \Gamma(X, \mathcal{O}_X)$。则 $S' = \Spec(R')$。
所以可以用 $S'$ 替换 $S$，并假设
% P05-EN-0116: frozen source omits conjunction/punctuation between affine and R=Gamma(X,O_X).
$S = \Spec(R)$ 仿射且 $R = \Gamma(X, \mathcal{O}_X)$。
接着，取一点 $s \in S$。设 $U \to S$ 是仿射概形的 \'etale 态射，
$u \in U$ 是映到 $s$ 的一点。令 $X_U \to U$ 为 $X$ 的基变换。
由 引理 \ref{more-morphisms-lemma-characterize-geometrically-connected-fibres}，
只需证明 $X_U \to U$ 在 $u$ 上的纤维连通。由
概形上同调章引理 \ref{coherent-lemma-flat-base-change-cohomology}，
有 $\Gamma(X_U, \mathcal{O}_{X_U}) = \Gamma(U, \mathcal{O}_U)$。
因此需要证明：给定仿射概形 $S = \Spec(R)$、满足
$\Gamma(X, \mathcal{O}_X) = R$ 的固有态射 $X \to S$，以及一点
$s \in S$，纤维 $X_s$ 连通。

\medskip\noindent
为此，只需说明幂等元 $e \in H^0(X_s, \mathcal{O}_{X_s})$ 只有
$0$ 和 $1$（由 引理 \ref{more-morphisms-lemma-stein-universally-closed}，已经知道
$X_s$ 非空）。由 概形的导出范畴章引理
\ref{perfect-lemma-proper-idempotent-on-fibre}，用一个主局部化替换 $R$ 后，
可以假设 $e$ 是 $R$ 中某个元素的像。由于
$R \to H^0(X_s, \mathcal{O}_{X_s})$ 经由 $\kappa(s)$ 分解，结论即得。
\end{proof}

\noindent
下面给出一个应用。

\begin{lemma}
\label{more-morphisms-lemma-geometrically-connected-fibres-towards-normal}
设 $f : X \to S$ 是概形态射。假设
\begin{enumerate}
\item $f$ 固有；
\item $S$ 是泛点为 $\xi$ 的整概形；
\item $S$ 正规；
\item $X$ 既约；
\item $X$ 的每个不可约分支的泛点均映到 $\xi$；
\item 有 $H^0(X_\xi, \mathcal{O}) = \kappa(\xi)$。
\end{enumerate}
则 $f_*\mathcal{O}_X = \mathcal{O}_S$，并且 $f$ 具有几何连通纤维。
\end{lemma}

\begin{proof}
应用 定理 \ref{more-morphisms-theorem-stein-factorization-general}，得到分解
$X \to S' \to S$。只需证明 $S' = S$。这将由 态射章引理
\ref{morphisms-lemma-finite-birational-over-normal} 得到。
具体地，由于 $X$ 既约，$S'$ 既约
（态射章引理 \ref{morphisms-lemma-normalization-in-reduced}）。
由上述定理，态射 $S' \to S$ 是整态射。由 态射章引理
\ref{morphisms-lemma-normalization-generic} 及假设 (5)，$S'$ 的每个泛点
都位于 $\xi$ 上方。另一方面，由于 $S'$ 是 $f_*\mathcal{O}_X$ 的相对谱，
概形论纤维 $S'_\xi$ 是 $H^0(X_\xi, \mathcal{O})$ 的谱，而按假设后者
等于 $\kappa(\xi)$。因此 $S'$ 是整概形，其函数域等于 $S$ 的函数域。
证明完毕。
\end{proof}

\noindent
下面是另一个应用。

\begin{lemma}
\label{more-morphisms-lemma-proper-flat-nr-geom-conn-comps-lower-semicontinuous}
设 $X \to S$ 是有限表示的平坦固有态射。令 $n_{X/S}$ 为 $S$ 上计数
% P05-EN-0117: frozen source refers to fibres of undefined f and says "the numbers";
% the named morphism is X -> S and the function counts the number.
$f$ 的纤维的几何连通分支个数的函数，它引入于
引理 \ref{more-morphisms-lemma-base-change-fibres-nr-geometrically-connected-components}.
则 $n_{X/S}$ 下半连续。
\end{lemma}

\begin{proof}
取 $s \in S$，令 $n = n_{X/S}(s)$。注意 $n < \infty$：
$X \to S$ 在 $s$ 处的几何纤维是域上的固有概形，因而是 Noether 的，
故只有有限多个连通分支。需要找到 $s$ 的开邻域 $V$，使得
$n_{X/S}|_V \geq n$。令 $X \to S' \to S$ 为
定理 \ref{more-morphisms-theorem-stein-factorization-general} 中的 Stein 分解。
由 引理 \ref{more-morphisms-lemma-stein-universally-closed-residue-fields}，位于 $s$
上方的点只有有限多个 $s'_1, \ldots, s'_m \in S'$，并且扩张
$\kappa(s'_i)/\kappa(s)$ 有限。引理 \ref{more-morphisms-lemma-etale-makes-integral-split}
于是表明：用 $s$ 的一个 \'etale 邻域替换 $S$ 后，可以假设作为概形
$S' = V_1 \amalg \ldots \amalg V_m$，且 $s'_i \in V_i$、
$\kappa(s'_i)/\kappa(s)$ 纯不可分。于是概形 $X_{s_i'}$ 在
$\kappa(s)$ 上几何连通，故 $m = n$。概形
$X_i = (f')^{-1}(V_i)$，$i = 1, \ldots, n$，在 $S$ 上平坦且有限表示。
因此 $X_i \to S$ 的像是开的
（态射章引理 \ref{morphisms-lemma-fppf-open}）。
由此可见，在 $s$ 的一个邻域内，$n_{X/S}$ 至少为 $n$。
\end{proof}

\begin{lemma}
\label{more-morphisms-lemma-proper-flat-geom-red}
设 $f : X \to S$ 是概形态射。假设
\begin{enumerate}
\item $f$ 固有、平坦且有限表示；
\item $f$ 的几何纤维既约。
\end{enumerate}
则计数 $f$ 的纤维的几何连通分支个数的函数
$n_{X/S} : S \to \mathbf{Z}$ 局部常值。
\end{lemma}

\begin{proof}
由 引理 \ref{more-morphisms-lemma-proper-flat-nr-geom-conn-comps-lower-semicontinuous}，
% P05-EN-0118: frozen source misspells "semicontinuous" as "semincontinuous".
函数 $n_{X/S}$ 下半连续。对 $s \in S$，考虑 $\kappa(s)$-代数
$$
A_s = H^0(X_s, \mathcal{O}_{X_s})
$$
% P05-EN-0119: frozen source defines beta_0 on X but then sends s in S and compares functions on S;
% the evident domain is S. Formula retained.
定义函数 $\beta_0 : X \to \mathbf{Z}$，将 $s$ 映到
$\dim_{\kappa(s)} A_s$。由 簇章引理
\ref{varieties-lemma-proper-geometrically-reduced-global-sections} 以及
$X_s$ 几何既约这一事实，$A_s$ 是 $\kappa(s)$ 的有限可分扩张的有限积。
因此 $A_s \otimes_{\kappa(s)} \kappa(\overline{s})$ 是
$\beta_0(s)$ 个 $\kappa(\overline{s})$ 的积。故 $X_{\overline{s}}$
有 $\beta_0(s)$ 个连通分支。换言之，作为 $S$ 上的函数，有
$n_{X/S} = \beta_0$。所以由 概形的导出范畴章引理
\ref{perfect-lemma-jump-loci-geometric}，$n_{X/S}$ 上半连续。证明完毕。
\end{proof}

\noindent
最后一个应用。

\begin{lemma}
\label{more-morphisms-lemma-split-off-proper-part-henselian}
\begin{reference}
Adic Noether 基的情形可参见
\cite[III, Proposition 5.5.1]{EGA}。
\end{reference}
设 $(A, I)$ 是 Hensel 对。设 $X \to \Spec(A)$ 分离且有限型。
令 $X_0 = X \times_{\Spec(A)} \Spec(A/I)$。设 $Y \subset X_0$
是开闭子概形，并且 $Y \to \Spec(A/I)$ 固有。则存在在 $A$ 上固有的
开闭子概形 $W \subset X$，使得
$W \times_{\Spec(A)} \Spec(A/I) = Y$。
\end{lemma}

\begin{proof}
以 $T \mapsto T_0$ 表示沿 $\Spec(A/I) \to \Spec(A)$ 的基变换。
由 Chow 引理（采用 极限章引理 \ref{limits-lemma-chow-finite-type}
的形式），存在满固有态射 $\varphi : X' \to X$，使得 $X'$ 可浸入
$\mathbf{P}^n_A$。令 $Y' = \varphi^{-1}(Y)$。这是 $X'_0$ 的开闭子概形。
假设本引理对 $(X', Y')$ 成立。令 $W' \subset X'$ 为在 $A$ 上固有且满足
$Y' = W'_0$ 的开闭子概形。由 态射章引理
\ref{morphisms-lemma-image-proper-scheme-closed}，
$W = \varphi(W') \subset X$ 和 $Q = \varphi(X' \setminus W') \subset X$
是闭子集；由 态射章引理
\ref{morphisms-lemma-image-proper-is-proper}，$W$ 在 $A$ 上固有。
$W \cap Q$ 在 $\Spec(A)$ 中的像是闭的。由于 $(A, I)$ 是 Hensel 对，
若 $W \cap Q$ 非空，则 $W \cap Q$ 有一个位于 $\Spec(A/I)$ 上方的点。
但这是不可能的，因为 $W'_0 = Y' = \varphi^{-1}(Y)$。
所以 $W$ 是 $X$ 的开闭子概形，在 $A$ 上固有，并满足 $W_0 = Y$。
由此化归到下一段所述的情形。

\medskip\noindent
假设存在 $A$ 上的浸入 $j : X \to \mathbf{P}^n_A$。令
$\overline{X}$ 为 $j$ 的概形论像。由于 $j$ 是拟紧态射
（概形章引理 \ref{schemes-lemma-quasi-compact-permanence}），
可见 $j : X \to \overline{X}$ 是开浸入
（态射章引理 \ref{morphisms-lemma-quasi-compact-immersion}）。
因此，基变换 $j_0 : X_0 \to \overline{X}_0$ 也是开浸入。
所以 $j_0(Y) \subset \overline{X}_0$ 是开的。由 态射章引理
\ref{morphisms-lemma-image-proper-scheme-closed}，它也是闭的。
假设本引理对 $(\overline{X}, j_0(Y))$ 成立。令
$\overline{W} \subset \overline{X}$ 为相应的、在 $A$ 上固有的开闭子概形，
使得 $j_0(Y) = \overline{W}_0$。于是
$T = \overline{W} \setminus j(X)$ 在 $\overline{W}$ 中闭；由于
$\overline{W}$ 在 $A$ 上固有，该子集在 $\Spec(A)$ 中的像是闭的。
由于 $(A, I)$ 是 Hensel 对，若 $T$ 非空，则 $T$ 中有一点映到
$\Spec(A/I)$。这是不可能的，因为
$j_0(Y) = \overline{W}_0$ 包含于 $j(X)$。故 $\overline{W}$ 包含于
$j(X)$，并可令 $W \subset X$ 等于唯一一个经 $j$ 同构映到
$\overline{W}$ 的开闭子概形。由此化归到下一段所述的情形。

\medskip\noindent
假设 $X \subset \mathbf{P}^n_A$ 是闭子概形。则
$X \to \Spec(A)$ 是固有态射。令 $Z = X_0 \setminus Y$。
这是 $X_0$ 的开闭子概形，并且 $X_0 = Y \amalg Z$。
令 $X \to X' \to \Spec(A)$ 为
定理 \ref{more-morphisms-theorem-stein-factorization-general} 中的 Stein 分解。
令 $Y' \subset X'_0$ 与 $Z' \subset X'_0$ 分别为 $Y$ 与 $Z$ 的像。
% P05-EN-0120: frozen source says fibres of X -> Z; the Stein map with connected fibres is X -> X'.
由于 $X \to Z$ 的纤维几何连通，可见 $Y' \cap Z' = \emptyset$。
又因为 $X \to X'$ 满，所以 $X'_0 = Y' \amalg Z'$。
由于 $X' \to \Spec(A)$ 是整态射，可见 $X'$ 是某个在 $A$ 上整的
$A$-代数的谱。回忆，谱的开闭子集与相应环中的幂等元 $1$-到-$1$ 对应；见
代数章引理 \ref{algebra-lemma-disjoint-decomposition}。
因此，由 代数进阶章引理
\ref{more-algebra-lemma-characterize-henselian-pair}，可以写成
$X' = W' \amalg V'$，其中 $W'$ 与 $V'$ 开且闭，并且
$Y' = W'_0$、$Z' = V'_0$。令 $W$ 为其在 $X$ 中的逆像，即完成证明。
\end{proof}









\section{泛平坦性分层}
\label{more-morphisms-section-generic-flatness-stratification}

\noindent
可以利用泛平坦性构造基的一个分层，使给定模在各层上成为平坦模。

\begin{lemma}[泛平坦性分层]
\label{more-morphisms-lemma-generic-flatness-stratification}
设 $f : X \to S$ 是拟紧拟分离概形之间的有限表示态射，
$\mathcal{F}$ 是有限表示 $\mathcal{O}_X$-模。则存在 $t \geq 0$
及闭子概形
$$
S \supset S_0 \supset S_1 \supset \ldots \supset S_t = \emptyset
$$
使得 $S_i \to S$ 由有限型理想层定义，$S_0 \subset S$ 是增厚，
并且 $\mathcal{F}$ 到 $X \times_S (S_i \setminus S_{i + 1})$ 的拉回
在 $S_i \setminus S_{i + 1}$ 上平坦。
\end{lemma}

\begin{proof}
可以找到笛卡尔图表
$$
\xymatrix{
X \ar[d] \ar[r] & X_0 \ar[d] \\
S \ar[r] & S_0
}
$$
以及有限表示 $\mathcal{O}_{X_0}$-模 $\mathcal{F}_0$，其拉回为
$\mathcal{F}$，并且 $X_0$ 与 $S_0$ 在 $\mathbf{Z}$ 上有限型。见
极限章命题 \ref{limits-proposition-approximate} 及
引理 \ref{limits-lemma-descend-finite-presentation} 及
\ref{limits-lemma-descend-modules-finite-presentation}。
因此，可以假设 $X$ 与 $S$ 在 $\mathbf{Z}$ 上有限型，而
$\mathcal{F}$ 是凝聚 $\mathcal{O}_X$-模。

\medskip\noindent
假设 $X$ 与 $S$ 在 $\mathbf{Z}$ 上有限型，而 $\mathcal{F}$ 是凝聚
$\mathcal{O}_X$-模。在此情形，每个拟凝聚理想都是有限型的，故无需检查
$S_i$ 由有限型理想切出这一条件。令 $S_0 = S_{red}$ 为 $S$ 的既约化。
由 态射章命题
\ref{morphisms-proposition-generic-flatness-reduced} 所述的泛平坦性，
存在稠密开集 $U_0 \subset S_0$，使得 $\mathcal{F}$ 到
$X \times_S U_0$ 的拉回在 $U_0$ 上平坦。令 $S_1 \subset S_0$
为底层闭子集等于 $S \setminus U_0$ 的既约闭子概形。若
$S_1 \not = \emptyset$，便依此继续，得到
$S_0 \supset S_1 \supset \ldots$。由于 $S$ 是 Noether 的，
任意闭子集降链均稳定，所以对某个 $t \geq 0$ 有
$S_t = \emptyset$。
\end{proof}

\begin{lemma}
\label{more-morphisms-lemma-generic-flatness-stratification-scheme}
设 $f : X \to S$ 是拟紧拟分离概形之间的有限表示态射。
则存在 $t \geq 0$ 及闭子概形
$$
S \supset S_0 \supset S_1 \supset \ldots \supset S_t = \emptyset
$$
使得 $S_i \to S$ 由有限型理想层定义，$S_0 \subset S$ 是增厚，并且
$X \times_S (S_i \setminus S_{i + 1})$ 在
$S_i \setminus S_{i + 1}$ 上平坦。
\end{lemma}

\begin{proof}
在 $\mathcal{F} = \mathcal{O}_X$ 的情形应用
引理 \ref{more-morphisms-lemma-generic-flatness-stratification}。
\end{proof}

\begin{lemma}[Longke Tang]
\label{more-morphisms-lemma-flatness-criterion}
设 $f : X \to S$ 是有限表示的满态射。设 $M$ 是
$D_\QCoh(\mathcal{O}_S)$ 的对象，并且对 $i > 0$ 有
$H^i(M) = 0$。下列条件等价：
\begin{enumerate}
\item $M$ 同构于置于次数 $0$ 的平坦 $\mathcal{O}_S$-模；
\item $Lf^*M$ 同构于置于次数 $0$ 的平坦 $\mathcal{O}_X$-模。
\end{enumerate}
\end{lemma}

\begin{proof}
略去 (1) 蕴含 (2) 的证明。假设 (2)。证明 (1) 是局部问题，
故可假设 $S$ 仿射。按
引理 \ref{more-morphisms-lemma-generic-flatness-stratification-scheme} 选取分层
$S \supset S_0 \supset S_1 \supset \ldots \supset S_t = \emptyset$
。
% P05-EN-0121: frozen source defines T_i=S_i\S_{i-1}; the strata above are S_i\S_{i+1},
% and i=0 would require undefined S_{-1}. Formula retained.
令 $T_i = S_i \setminus S_{i - 1}$、$Y_i = X \times_S T_i$。
以 $g_i : T_i \to S$ 和 $h_i : Y_i \to X$ 表示包含态射。
% P05-EN-0122: frozen source uses undefined X_i as source of f_i after defining Y_i.
由于 $f$ 满，态射 $f_i : X_i \to T_i$ 满且平坦。由于
$Lh_i^* \circ Lf^* = Lf_i^* \circ Lg_i^*$，由 (2) 断定
$Lf_i^* Lg_i^*M$ 同构于置于次数 $0$ 的平坦 $\mathcal{O}_{Y_i}$-模。
由于 $f_i$ 平坦且满，可见 $Lg_i^*M$ 同构于置于次数 $0$ 的平坦
$\mathcal{O}_{T_i}$-模。

\medskip\noindent
我们对 $t$ 归纳证明这蕴含所需结论。采用如下记号：
$S = \Spec(A)$，并且 $M$ 是 $D_\QCoh(\mathcal{O}_S)$ 中与
$D(A)$ 中的 $N$ 对应的对象；见 概形的导出范畴章引理
\ref{perfect-lemma-affine-compare-bounded}。

\medskip\noindent
基础情形：$t = 1$。写成 $S_0 = \Spec(A/I)$，其中 $I$ 是有限生成理想。
由于在集合论意义下 $S_0 = S$，$I$ 幂零。假设表明
$N \otimes_A^\mathbf{L} A/I$ 的 Tor 振幅在 $[0, 0]$ 中。
由 代数进阶章引理
\ref{more-algebra-lemma-check-tor-dimension-modulo-nilpotent}，
$N$ 亦然。

\medskip\noindent
归纳步骤。假设 $t > 1$。写成 $S_{t - 1} = \Spec(A/I)$，其中
$I = (f_1, \ldots, f_r)$ 是有限生成理想。对 $r$ 归纳论证。
注意，由归纳假设，$D(A_{f_r})$ 中的
$N_{f_r} = N \otimes_A^\mathbf{L} A_{f_r}$ 的 Tor 振幅在
$[0, 0]$ 中（因为主开集 $D(f_r)$ 上的分层长度至多为 $t - 1$）。
另一方面，考虑 $D(A/f_rA)$ 中的
$N' = N \otimes_A^\mathbf{L} A/f_rA$。令
$S' = \Spec(A/f_rA)$，得到分层
$$
S' \supset S' \cap S_0 \supset S' \cap S_1 \supset \ldots \supset
S' \cap S_{t - 1} \supset
S' \cap S_t = \emptyset
$$
其中 $S' \cap S_{t - 1}$ 在 $S'$ 中由 $r - 1$ 个方程切出。
与先前类似，$M$ 到各部分
$S' \cap S_i \setminus S' \cap S_{i + 1}$ 的导出拉回的 Tor 振幅均在
$[0, 0]$ 中。因此，由对 $r$ 的归纳，$N'$ 的 Tor 振幅在
$[0, 0]$ 中。最后，由 代数进阶章引理
\ref{more-algebra-lemma-f-flat}，断定 $N$ 的 Tor 振幅在 $[0, 0]$ 中。
\end{proof}

\begin{lemma}
\label{more-morphisms-lemma-cokernel-flat}
设 $R$ 是 Noether 整环，$R \to A \to B$ 是有限型环映射。
设 $M$ 是有限 $A$-模，$N$ 是有限 $B$-模，$M \to N$ 是
$A$-线性映射。则存在非零的 $f \in R$，使得
$M_f \to N_f$ 的余核是平坦 $R_f$-模。
\end{lemma}

\begin{proof}
用 $M \to N$ 的像替换 $M$ 后，可以假设 $M \subset N$。
选取滤过 $0 = N_0 \subset N_1 \subset \ldots \subset N_t = N$，
使得对某个素理想 $\mathfrak q_i \subset B$，有
$N_i/N_{i - 1} = B/\mathfrak q_i$；见
代数章引理 \ref{algebra-lemma-filter-Noetherian-module}。
令 $M_i = M \cap N_i$。则 $Q = N/M$ 具有由子模
$Q_i = N_i/M_i$ 给出的滤过。只需证明在某个非零的 $f$ 的元素处局部化后，
% P05-EN-0123: frozen source says "a nonzero element of f"; the element should be of R.
$Q_i/Q_{i - 1}$ 成为平坦模（因为由 代数章引理
\ref{algebra-lemma-flat-ses}，平坦模的扩张平坦）。由于
$Q_i/Q_{i - 1}$ 同构于映射
$M_i/M_{i - 1} \to N_i/N_{i - 1}$ 的余核，故化归到下一段讨论的情形。

\medskip\noindent
假设 $B$ 是整环，且 $M \subset N = B$。用 $A$ 在 $B$ 中的像替换
$A$ 后，可以假设 $A \subset B$。由泛平坦性，可假设 $A$ 与 $B$
在 $R$ 上平坦
（代数章引理 \ref{algebra-lemma-generic-flatness-Noetherian}）。
现在，只需说明：用一个主局部化替换 $R$ 后，$M \to B$ 成为
$R$-普遍单射
（代数章引理 \ref{algebra-lemma-ui-flat-domain}）。
由泛自由性，可以找到非零的 $g \in A$，使得 $B_g$ 是自由
$A_g$-模（代数章引理
\ref{algebra-lemma-generic-flatness-Noetherian}）。
因此，可以选取作为 $A_g$-模的直和项 $M' \subset B_g$，使它是有限自由
$A_g$-模，并且 $M \to B \to B_g$ 经由 $M'$ 分解。
显然，只需说明：用一个主局部化替换 $R$ 后，$M \to M'$ 成为
$R$-普遍单射。

\medskip\noindent
设 $M' = A_g^{\oplus n}$。由于 $M \subset M'$ 是有限 $A$-模，
可见对某个 $m \geq 0$，$M$ 包含于 $(1/g^m)A^{\oplus n}$。
改变 $M'$ 的基后，可以假设 $M \subset A^{\oplus n}$。
于是只需说明：用一个主局部化替换 $R$ 后，$A^{\oplus n}/M$ 与
$A_g/A$ 成为 $R$-平坦模。事实上，此时由 代数章引理
\ref{algebra-lemma-flat-tor-zero}，$M \to A^{\oplus n}$ 与
$A^{\oplus n} \to A_g^{\oplus n}$ 普遍单射，因而它们的复合
$M \to M' = A_g^{\oplus n}$ 也普遍单射。

\medskip\noindent
由泛平坦性（见上述引文），可以假设模 $A^{\oplus n}/M$ 是
$R$-平坦的。对商 $A_g/A$，使用事实
$$
A_g/A = \colim (1/g^m)A/A \cong \colim A/g^mA
$$
以及模 $A/g^mA$ 具有长度为 $m$ 的滤过，其相继商均同构于 $A/gA$。
再次由泛平坦性，可以假设 $A/gA$ 是 $R$-平坦的，因而每个
$A/g^mA$ 都是 $R$-平坦的，从而 $A_g/A$ 亦然。
\end{proof}

\noindent
设 $f : X \to Y$ 是基概形 $S$ 上的概形态射。令 $Z \subset Y$
为 $f$ 的概形论像；见
态射章第 \ref{morphisms-section-scheme-theoretic-image} 节。
设 $g : S' \to S$ 是概形态射，并令
$f' : X \times_S S' \to Y \times_S S'$ 为 $f$ 沿 $g$ 的基变换。
$Z \times_S S' \subset Y \times_S S'$ 未必是 $f'$ 的概形论像。
若对每个上述 $g$，$f'$ 的概形论像都等于 $Z \times_S S'$，
则称{\it $f/S$ 的概形论像的形成与任意基变换交换}。

\begin{lemma}
\label{more-morphisms-lemma-base-change-scheme-theoretic-image}
设 $S$ 是拟紧拟分离概形。设 $f : X \to Y$ 是 $S$ 上的概形态射，
且 $X$ 与 $Y$ 均在 $S$ 上有限表示。则存在 $t \geq 0$ 及闭子概形
$$
S \supset S_0 \supset S_1 \supset \ldots \supset S_t = \emptyset
$$
具有下列性质：
\begin{enumerate}
\item $S_i \to S$ 由有限型理想层定义；
\item $S_0 \subset S$ 是增厚；
\item 令 $T_i = S_i \setminus S_{i + 1}$，并令 $f_i$ 为 $f$ 到
$T_i$ 的基变换，则 $f_i/T_i$ 的概形论像的形成与任意基变换交换
（见引理上方的讨论）。
\end{enumerate}
\end{lemma}

\begin{proof}
可以找到交换图表
$$
\xymatrix{
X \ar[d] \ar[r] & Y \ar[d] \ar[r] & S \ar[d] \\
U \ar[r] & V \ar[r] & W
}
$$
，其中两个方块均笛卡尔，且 $U$、$V$、$W$ 在 $\mathbf{Z}$ 上有限型。
具体地，首先利用 极限章命题 \ref{limits-proposition-approximate}，
将 $S$ 写成转移态射仿射的 $\mathbf{Z}$ 上有限型概形的余滤极限，
再利用 极限章引理 \ref{limits-lemma-descend-finite-presentation}
下降态射 $X \to Y$。由此化归到下一段讨论的情形。

\medskip\noindent
假设 $S$ 是 Noether 的。在此情形，每个拟凝聚理想都是有限型的，
故无需检查 $S_i$ 由有限型理想切出这一条件。令
$S_0 = S_{red}$ 为 $S$ 的既约化。令 $\eta \in S_0$ 为 $S_0$
的某个不可约分支的泛点。对 $S_0$ 的底层拓扑空间作 Noether 归纳，
可以假设结论对 $S_0$ 的任何不包含 $\eta$ 的闭子概形均成立。
所以只需说明存在开邻域 $U_0 \subset S_0$，使得 $f$ 到 $U_0$
的基变换 $f_0$ 具有性质 (3)。

\medskip\noindent
设 $R$ 是 Noether 整环，$f : X \to Y$ 是 $R$ 上有限型概形的态射。
由上一段的讨论，只需说明：对某个非零的 $g \in R$，用 $R_g$ 替换
$R$，并用到 $R_g$ 的基变换替换 $X$、$Y$ 后，$f/R$ 的概形论像的形成
与任意基变换交换。

\medskip\noindent
令 $Y = V_1 \cup \ldots V_n$ 为仿射开覆盖，并令
$U_i = f^{-1}(V_i)$。若该陈述对每个 $R$ 上的态射
$U_i \to V_i$ 均成立，则它对 $f$ 成立。事实上，由 态射章引理
\ref{morphisms-lemma-quasi-compact-scheme-theoretic-image}，
$U_i \to V_i$ 的概形论像是 $V_i$ 与 $f : X \to Y$ 的概形论像之交。
所以可以假设 $Y$ 仿射。

\medskip\noindent
令 $X = U_1 \cup \ldots U_n$ 为仿射开覆盖。则 $X \to Y$ 的概形论像
与 $\coprod U_i \to Y$ 的概形论像相同。
% P05-EN-0124: frozen source misspells "image" as "imge".
因此可以假设 $X$ 仿射。

\medskip\noindent
写成 $X = \Spec(A)$、$Y = \Spec(B)$，并设 $f$ 对应于 $R$-代数映射
% P05-EN-0125: frozen source reverses the affine contravariance: X=Spec(A)->Y=Spec(B)
% should correspond to B->A, and the image should be Spec(B/Ker). Formulas retained.
$\varphi : A \to B$。则 $f$ 的概形论像是
$\Spec(A/\Ker(\varphi))$；基变换后亦同（这里基变换采用仿射态射，
但只需对此检查）。因此，若对所有环映射 $R \to R'$ 都有
$\Ker(\varphi \otimes_R R') = \Ker(\varphi) \otimes_R R'$，
则概形论像的形成与基变换交换。

\medskip\noindent
对 $R$ 中适当的非零元素 $g$，分别用 $R_g$、$A_g$、$B_g$ 替换
$R$、$A$、$B$ 后，可以假设 $A$ 与 $B$ 在 $R$ 上平坦。
由 引理 \ref{more-morphisms-lemma-cokernel-flat}，还可假设 $B/A$ 是平坦 $R$-模。
于是 $0 \to \Ker(\varphi) \to A \to B \to B/A \to 0$
是平坦 $R$-模的正合列，这蕴含所需的基变换陈述。
\end{proof}






\section{对态射分层}
\label{more-morphisms-section-stratifying-morphisms}

\noindent
设 $f : X \to S$ 是拟紧拟分离概形之间的有限表现态射。
在 第 \ref{more-morphisms-section-generic-flatness-stratification} 节 中，
我们已经看到，可以对 $S$ 分层，使得 $X$ 在各层上平坦。
本节尝试同时对 $S$ 和 $X$ 分层，以得到光滑的各层；
这不能完全直接实现，我们还需要沿有限局部自由态射作基变换。

\begin{lemma}
\label{more-morphisms-lemma-smoothness-stratification-at-generic-point}
设 $f : X \to S$ 是有限表现的概形态射。
设 $\eta \in S$ 是 $S$ 的一个不可约分支的泛点。
假设 $S$ 既约。则存在
\begin{enumerate}
\item 包含 $\eta$ 的开子概形 $U \subset S$，
\item 满射、普遍单射的有限局部自由态射 $V \to U$，
\item 某个 $t \geq 0$ 以及闭子概形
$$
X \times_S V \supset Z_0 \supset Z_1 \supset \ldots \supset Z_t = \emptyset
$$
使得 $Z_i \to X \times_S V$ 由有限型理想层定义，
$Z_0 \subset X \times_S V$ 是一个增厚，并且态射
$Z_i \setminus Z_{i + 1} \to V$ 光滑。
\end{enumerate}
\end{lemma}

\begin{proof}
显然，我们可以用 $\eta$ 的一个开邻域替换 $S$，并用 $X$
在此开集上的限制替换之。因此可以假设 $S = \Spec(A)$，
其中 $A$ 是既约环，而 $\eta$ 对应于一个极小素理想 $\mathfrak p$。
回忆在这种情形下，局部环 $\mathcal{O}_{S, \eta} = A_\mathfrak p$
等于 $\kappa(\mathfrak p)$，见
代数章引理 \ref{algebra-lemma-minimal-prime-reduced-ring}.

\medskip\noindent
对 $k = \kappa(\eta)$ 上的概形 $X_\eta$ 应用
簇章引理 \ref{varieties-lemma-smooth-stratification}。
将由此得到的纯不可分域扩张记为 $k'/k$。
下一段把问题约化到 $k' = k$ 的情形。
（这一步对应于寻找引理陈述中的态射 $V \to U$；特别地，
若 $\kappa(\mathfrak p)$ 的特征为零，则可以取 $V = U$。）

\medskip\noindent
若 $k = \kappa(\mathfrak p)$ 的特征为零，则 $k' = k$。
若 $k = \kappa(\mathfrak p)$ 的特征为 $p > 0$，则 $p$
在 $A_\mathfrak p = \kappa(\mathfrak p)$ 中的像为零。
因此，用一个主局部化替换 $A$（即缩小 $S$）后，
可以假设在 $A$ 中有 $p = 0$。若 $k' \not = k$，则存在
$\beta \in k'$、$\beta \not \in k$，使得 $\beta^p \in k$。
用一个主局部化替换 $A$ 后，可以假设存在 $a \in A$，
使得 $\beta^p = a$。置 $A' = A[x]/(x^p - a)$。
于是 $S' = \Spec(A') \to \Spec(A) = S$ 是有限局部自由、满射且普遍单射的。
此外，若 $\mathfrak p' \subset A'$ 表示位于 $\mathfrak p$ 上方的唯一素理想，
则 $A'_{\mathfrak p'} = k(\beta)$，且 $k'/k(\beta)$ 的次数更小。
因此，用 $S'$ 替换 $S$，并用对应于 $\mathfrak p'$ 的点 $\eta'$
替换 $\eta$ 后，可见 $k'$ 在 $\eta$ 的剩余域上的次数已经减小。
如此继续并作归纳，就约化到
$k' = \kappa(\mathfrak p) = \kappa(\eta)$ 的情形。

\medskip\noindent
因此，可以假设 $S$ 仿射且既约，并且存在 $t \geq 0$ 及闭子概形
$$
X_\eta \supset Z_{\eta, 0} \supset Z_{\eta, 1}
\supset \ldots \supset Z_{\eta, t} = \emptyset
$$
使得 $Z_{\eta, 0} = (X_\eta)_{red}$，并且对所有 $i$，
$Z_{\eta, i} \setminus Z_{\eta, i + 1}$ 在 $\eta$ 上光滑。
回忆 $\kappa(\eta) = \kappa(\mathfrak p) = A_\mathfrak p$
是所有 $A_a$（其中 $a \in A$、$a \not \in \mathfrak p$）的滤过余极限。
见 代数章引理 \ref{algebra-lemma-localization-colimit}。
因此，对某个 $a \in A$、$a \not \in \mathfrak p$，
可以把上图下降为 $\Spec(A_a)$ 上的相应图。
更确切地，用 $\Spec(A_a)$ 替换 $S$ 后，可以假设存在
$t \geq 0$ 及闭子概形
$$
X \supset Z_0 \supset Z_1 \supset \ldots \supset Z_t = \emptyset
$$
使得 $Z_i \to X$ 是有限表现的闭浸入，$Z_0 \to X$ 是一个增厚，
并且 $Z_i \setminus Z_{i + 1}$ 在 $S$ 上光滑。换言之，引理成立。
更确切地，我们先用
极限章引理 \ref{limits-lemma-descend-finite-presentation}
得到态射
$$
Z_t \to Z_{t - 1} \to \ldots \to Z_0 \to X
$$
于 $S$ 上；每个态射都是有限表现的，并且它们到 $\eta$ 的基变换给出
上述给定闭子概形之间的包含。进一步缩小 $S$ 后，可以假设
每个态射都是闭浸入，见
极限章引理
\ref{limits-lemma-descend-closed-immersion-finite-presentation}.
缩小 $S$ 后，可以假设 $Z_0 \to X$ 满射，因而是一个增厚，见
极限章引理 \ref{limits-lemma-descend-surjective}.
再缩小一次 $S$，可以假设 $Z_i \setminus Z_{i + 1} \to S$ 光滑，见
极限章引理 \ref{limits-lemma-descend-smooth}.
证明完毕。
\end{proof}

\begin{lemma}
\label{more-morphisms-lemma-smoothness-stratification}
设 $f : X \to S$ 是拟紧拟分离概形之间的有限表现态射。
则存在 $t \geq 0$ 及闭子概形
$$
S \supset S_0 \supset S_1 \supset \ldots \supset S_t = \emptyset
$$
使得
\begin{enumerate}
\item $S_i \to S$ 由有限型理想层定义，
\item $S_0 \subset S$ 是一个增厚，
\item 对每个 $i$，存在满射的有限局部自由态射
$T_i \to S_i \setminus S_{i + 1}$，
\item 对每个 $i$，存在 $t_i \geq 0$ 及闭子概形
$$
X_i = X \times_S T_i \supset Z_{i, 0}
\supset Z_{i, 1} \supset \ldots \supset Z_{i, t_i} = \emptyset
$$
使得 $Z_{i, j} \to X_i$ 由有限型理想层定义，
$Z_{i, 0} \subset X_i$ 是一个增厚，并且态射
$Z_{i, j} \setminus Z_{i, j + 1} \to T_i$ 光滑。
\end{enumerate}
\end{lemma}

\begin{proof}
可以找到笛卡尔图
$$
\xymatrix{
X \ar[d] \ar[r] & X_0 \ar[d] \\
S \ar[r] & S_0
}
$$
使得 $X_0$ 和 $S_0$ 在 $\mathbf{Z}$ 上为有限型。见
极限章命题 \ref{limits-proposition-approximate} 及
引理 \ref{limits-lemma-descend-finite-presentation}.
因此，可以假设 $X$ 和 $S$ 在 $\mathbf{Z}$ 上为有限型。
事实上，引理对 $X_0 \to S_0$ 所提出问题的一个解，
经基变换后就给出 $S$ 上的一个解；细节略。

\medskip\noindent
假设 $X$ 和 $S$ 在 $\mathbf{Z}$ 上为有限型。
在这种情形下，每个拟凝聚理想都是有限型的，因而无须检验
$S_i$ 由有限型理想截出的条件。置 $S_0 = S_{red}$，即 $S$ 的既约化。
设 $\eta \in S_0$ 是一个不可约分支的泛点。
由 引理 \ref{more-morphisms-lemma-smoothness-stratification-at-generic-point}，
可以找到开子概形 $U \subset S_0$、满射且普遍单射的有限局部自由态射
$V \to U$、某个 $t_0 \geq 0$ 以及闭子概形
$$
X \times_S V \supset Z_{0, 0} \supset Z_{0, 1} \supset \ldots \supset
Z_{0, t_0} = \emptyset
$$
使得 $Z_{0, i} \to X \times_S V$ 由有限型理想层定义，
$Z_{0, 0} \subset X \times_S V$ 是一个增厚，并且态射
$Z_{0, i} \setminus Z_{0, i + 1} \to V$ 光滑。
然后，令 $S_1 \subset S_0$ 是 $S_0 \setminus U$ 上的既约诱导子概形结构。
对 $S$ 的底拓扑空间作诺特归纳，可以假设引理对
$X \times_S S_1 \to S_1$ 成立。这给出 $t \geq 1$ 以及
$$
S_1 = S_1 \supset S_2 \supset \ldots \supset S_t = \emptyset
$$
以及引理陈述中的 $t_i$ 和 $Z_{i, j}$。
这就证明了引理。
\end{proof}


















\section{改进相对维数为一的态射}
\label{more-morphisms-section-make-good-curves}

\noindent
通过扩张基域、紧化及正规化，可以使任意曲线成为光滑射影曲线。
这还蕴含关于泛纤维维数为 $1$ 的有限型态射的结果。

\begin{lemma}
\label{more-morphisms-lemma-make-good-curves}
设 $f : X \to S$ 是概形态射。设 $\eta \in S$ 是 $S$
的一个不可约分支的泛点。假设 $f$ 分离、有限表现，且
$\dim(X_\eta) \leq 1$。则存在交换图
$$
\xymatrix{
\overline{Y}_1 \amalg \ldots \amalg \overline{Y}_n \ar[rd] &
Y_1 \amalg \ldots \amalg Y_n \ar[r]_-\nu \ar[d] \ar[l]^j &
X_V \ar[r] \ar[d] &
X_U \ar[r] \ar[d] &
X \ar[d]^f \\
& T_1 \amalg \ldots \amalg T_n \ar[r] &
V \ar[r] &
U \ar[r] &
S
}
$$
满足下列性质：
\begin{enumerate}
\item $U \subset S$ 是 $\eta$ 的一个开邻域，
\item $V \to U$ 是有限、满射且普遍单射的态射，
\item $X_U = U \times_S X$ 和 $X_V = V \times_S X$ 是相应的基变换，
\item $\nu$ 有限且满射，并且存在开集 $W \subset X_V$，使得
\begin{enumerate}
\item $W$ 在 $X_V \to V$ 的每个纤维中稠密，
\item $\nu^{-1}(W) \cap Y_i$ 在 $Y_i \to T_i$ 的每个纤维中稠密，并且
\item $\nu^{-1}(W) \to W$ 是一个增厚，
\end{enumerate}
\item $j$ 是开浸入，
\item $T_i \to V$ 是有限 \'etale 态射，
\item $Y_i \to T_i$ 满射且光滑，
\item $\overline{Y}_i \to T_i$ 光滑且固有，其纤维几何连通且维数 $\leq 1$。
\end{enumerate}
\end{lemma}

\begin{proof}
显然，我们可以用 $\eta$ 的一个开邻域替换 $S$，并用 $X$
在此开集上的限制替换之。此外，可以用 $S$ 的既约化替换它，
并用 $X$ 到此既约化的基变换替换之。
因此可以假设 $S = \Spec(A)$，其中 $A$ 是既约环，
而 $\eta$ 对应于一个极小素理想 $\mathfrak p$。
回忆在这种情形下，局部环 $\mathcal{O}_{S, \eta} = A_\mathfrak p$
等于 $\kappa(\mathfrak p)$，见
代数章引理 \ref{algebra-lemma-minimal-prime-reduced-ring}.

\medskip\noindent
对 $k = \kappa(\eta)$ 上的概形 $X_\eta$ 应用 簇章引理
\ref{varieties-lemma-dim-1-projective-completion-after-insep}
。将由此得到的纯不可分域扩张记为 $k'/k$。
下一段把问题约化到 $k' = k$ 的情形。
（这一步对应于寻找引理陈述中的态射 $V \to U$；特别地，
若 $\kappa(\mathfrak p)$ 的特征为零，则可以取 $V = U$。）

\medskip\noindent
若 $k = \kappa(\mathfrak p)$ 的特征为零，则 $k' = k$。
若 $k = \kappa(\mathfrak p)$ 的特征为 $p > 0$，则 $p$
在 $A_\mathfrak p = \kappa(\mathfrak p)$ 中的像为零。
因此，用一个主局部化替换 $A$（即缩小 $S$）后，
可以假设在 $A$ 中有 $p = 0$。若 $k' \not = k$，则存在
$\beta \in k'$、$\beta \not \in k$，使得 $\beta^p \in k$。
用一个主局部化替换 $A$ 后，可以假设存在 $a \in A$，
使得 $\beta^p = a$。置 $A' = A[x]/(x^p - a)$。
于是 $S' = \Spec(A') \to \Spec(A) = S$ 有限、满射且普遍单射。
此外，若 $\mathfrak p' \subset A'$ 表示位于 $\mathfrak p$ 上方的唯一素理想，
则 $A'_{\mathfrak p'} = k(\beta)$，且 $k'/k(\beta)$ 的次数更小。
因此，用 $S'$ 替换 $S$，并用对应于 $\mathfrak p'$ 的点 $\eta'$
替换 $\eta$ 后，可见 $k'$ 在 $\eta$ 的剩余域上的次数已经减小。
如此继续并作归纳，就约化到
$k' = \kappa(\mathfrak p) = \kappa(\eta)$ 的情形。

\medskip\noindent
因此，可以假设 $S$ 仿射且既约，并且有图
$$
\xymatrix{
\overline{Y}_{1, \eta} \amalg \ldots \amalg \overline{Y}_{n, \eta} \ar[rd] &
Y_{1, \eta} \amalg \ldots \amalg Y_{n, \eta} \ar[r]_-\nu \ar[d] \ar[l]^j &
X_\eta \ar[d] \\
& \Spec(k_1) \amalg \ldots \amalg \Spec(k_n) \ar[r] &
\eta
}
$$
满足下列性质：
\begin{enumerate}
\item $\nu$ 是 $X_\eta$ 的正规化，
\item $j$ 是像稠密的开浸入，
\item 对 $i = 1, \ldots, n$，$k_i/\kappa(\eta)$ 是有限可分扩张，
\item $\overline{Y}_{i, \eta}$ 在 $k_i$ 上光滑、射影、几何不可约，
且维数 $\leq 1$。
\end{enumerate}
回忆 $\kappa(\eta) = \kappa(\mathfrak p) = A_\mathfrak p$
是所有 $A_a$（其中 $a \in A$、$a \not \in \mathfrak p$）的滤过余极限。
见 代数章引理 \ref{algebra-lemma-localization-colimit}。
因此，对某个 $a \in A$、$a \not \in \mathfrak p$，
可以把上图下降为 $\Spec(A_a)$ 上的相应图。
更确切地，用 $\Spec(A_a)$ 替换 $S$ 后，可以假设有交换图
$$
\xymatrix{
\overline{Y}_1 \amalg \ldots \amalg \overline{Y}_n \ar[rd] &
Y_1 \amalg \ldots \amalg Y_n \ar[r]_-\nu \ar[d] \ar[l]^j &
X \ar[d] \\
& T_1 \amalg \ldots \amalg T_n \ar[r] & S
}
$$
其到 $\eta$ 的基变换是上图，并满足下列性质：
\begin{enumerate}
\item $\nu$ 是有限满射，
\item $j$ 是开浸入，
\item 对 $i = 1, \ldots, n$，$T_i \to S$ 是有限 \'etale 态射，
\item $Y_i \to T_i$ 光滑且满射，
\item $\overline{Y}_i \to T_i$ 光滑且固有，并具有维数 $\leq 1$
的几何连通纤维。
\end{enumerate}
为此，先用
极限章引理 \ref{limits-lemma-descend-finite-presentation}
得到一个经基变换后为前图的图。然后使用 极限章引理
\ref{limits-lemma-descend-etale},
\ref{limits-lemma-descend-smooth},
\ref{limits-lemma-descend-finite-finite-presentation},
\ref{limits-lemma-limit-affine},
\ref{limits-lemma-descend-open-immersion},
\ref{limits-lemma-eventually-proper}, 及
\ref{limits-lemma-descend-surjective}
，得到 $\nu$ 有限且满射、$j$ 为开浸入、$T_i \to S$ 为有限
\'etale 态射、$Y_i \to T$ 光滑，以及 $\overline{Y}_i \to T_i$
固有且光滑。由于 $Y_i$ 不可能为空，光滑态射是开的，且
$T_i \to S$ 为有限 \'etale 态射，缩小 $S$ 后可以假设
$Y_i \to T_i$ 满射。最后，$\overline{Y}_i \to T_i$ 在 $T_i$
中位于 $\eta$ 上方的唯一点 $\eta_i = \Spec(k_i)$ 处的纤维几何连通。
因此，再次缩小后，可以假设所有纤维都具有同一性质，见
引理 \ref{more-morphisms-lemma-proper-flat-geom-red}.

\medskip\noindent
还需证明存在满足 (a)、(b)、(c) 的开集 $W \subset X$。
由于 $\nu_\eta : \coprod Y_{i, \eta} \to X_\eta$ 是正规化态射，
由
簇章引理 \ref{varieties-lemma-normalization-locally-algebraic}
可知，存在稠密开集 $W_\eta \subset X_\eta$，使得
$\nu^{-1}(W_\eta) \to W_\eta$ 等于 $W_\eta$ 的既约化到
$W_\eta$ 的包含。取拟紧开集 $W \subset X$，使其在 $\eta$
上的纤维就是刚找到的开集 $W_\eta$。
用另一个局部化替换 $A = \Gamma(S, \mathcal{O}_S)$ 后，
可以假设 $\nu^{-1}(W) \to W$ 是闭浸入，见 极限章引理
\ref{limits-lemma-descend-closed-immersion-finite-presentation}.
又因为 $\nu$ 满射，所以 $\nu^{-1}(W) \to W$ 是一个增厚。
置 $W_i = \nu^{-1}(W) \cap Y_i$。再次缩小 $S$，可以假设
对所有 $i$，$W_i \to T_i$ 满射（论证同上）。
由于 $Y_i \to T_i$ 具有几何不可约纤维，于是
$W_i \subset Y_i$ 在 $Y_i \to T_i$ 的每个纤维中稠密。
又因 $\nu$ 有限且满射，遂有 $W = \nu(\nu^{-1}(W))$
在 $X \to S$ 的每个纤维中也稠密。
\end{proof}











\section{分离且局部拟有限态射的下降}
\label{more-morphisms-section-separated-locally-quasi-finite}

\noindent
本节证明“分离且局部拟有限的态射对 fppf 覆盖满足下降”。
术语见 下降章定义
\ref{descent-definition-descending-types-morphisms}。
这个结果隐藏在 Raynaud 与 Gruson 那篇（在许多方面都）极精彩的论文
\cite[Lemma 5.7.1]{GruRay} 的证明中。
它也见于 \cite{Murre-representation} 以及
\cite[Expos\'e X, Lemma 5.4]{SGA3}
；后两者附加假设该态射局部有限表现。下面是正式陈述。

\begin{lemma}
\label{more-morphisms-lemma-separated-locally-quasi-finite-morphisms-fppf-descend}
设 $S$ 是概形。设 $\{X_i \to S\}_{i\in I}$ 是一个 fppf 覆盖，见
拓扑章定义 \ref{topologies-definition-fppf-covering}.
设 $(V_i/X_i, \varphi_{ij})$ 是相对于 $\{X_i \to S\}$ 的下降数据。
若每个态射 $V_i \to X_i$ 都分离且局部拟有限，则此下降数据有效。
\end{lemma}

\begin{proof}
分离性与局部拟有限性都是在任意基变换下保持的概形态射性质，见
概形章引理 \ref{schemes-lemma-separated-permanence} 及
态射章引理 \ref{morphisms-lemma-base-change-quasi-finite}.
因此可以应用 下降章引理 \ref{descent-lemma-descending-types-morphisms}，
只需证明 fppf 覆盖由仿射概形之间的单个平坦、满射、有限表现态射
$\{X \to S\}$ 给出的情形。写 $X = \Spec(A)$、$S = \Spec(R)$，
于是 $R \to A$ 是忠实平坦环映射。设 $(V, \varphi)$ 是相对于
$S$ 上 $X$ 的下降数据，并假设 $\pi : V \to X$ 分离且局部拟有限。

\medskip\noindent
任取仿射开集 $W^1 \subset V$。考虑
$W = \text{pr}_1(\varphi(W^1 \times_S X)) \subset V$。图示如下：
$$
\xymatrix{
W^1 \times_S X \ar[rrrrr] \ar[ddd] \ar[rd]
& & & & &
\varphi(W^1 \times_S X) \ar[ddd] \ar[ld] \\
& V \times_S X \ar[rrr]^\varphi \ar[rd] \ar[dd]
& & &
X \times_S V \ar[ld] \ar[dd] & \\
& &
X \times_S X \ar[r]^1 \ar[d]_{\text{pr}_0}
&
X \times_S X \ar[d]^{\text{pr}_1}
& & \\
W^1 \ar[r] &
V \ar[r] &
X &
X &
V \ar[l] &
W \ar[l]
}
$$
因为 $X \to S$ 平坦且有限表现，所以它普遍开
（态射章引理 \ref{morphisms-lemma-fppf-open}）。
因此 $W$ 是开集。显然，$W$ 还拟紧，并且 $W^1 \subset W$。
另外，由 $\varphi$ 的上闭链条件，
$\varphi(W \times_S X) = X \times_S W$。
因此，把 $\varphi$ 限制到 $W \times_S X$，得到新的下降数据
$(W, \varphi')$。注意态射 $W \to X$ 拟紧、分离且局部拟有限。
按定义，这蕴含它分离且拟有限。于是由
引理 \ref{more-morphisms-lemma-quasi-finite-separated-quasi-affine}
可知它是拟仿射的。故由
下降章引理 \ref{descent-lemma-quasi-affine}
可知下降数据 $(W, \varphi')$ 有效。

\medskip\noindent
换言之，存在由拟紧开集 $W_i$ 构成的开覆盖
$V = \bigcup W_i$，其中各开集在下降态射 $\varphi$ 下稳定。
此外，对每个这样的拟紧开集 $W \subset V$，相应的下降数据
$(W, \varphi')$ 都有效。把由各开集 $W_i$ 下降所得的概形粘合起来，
便知原下降数据有效，见
下降章引理 \ref{descent-lemma-effective-for-fpqc-is-local-upstairs}.
\end{proof}






\section{相对有限表现}
\label{more-morphisms-section-finite-type-finite-presentation}

\noindent
设 $R \to A$ 是有限型环映射，$M$ 是 $A$-模。在
代数进阶章
第 \ref{more-algebra-section-relative-finite-presentation} 节
中，我们定义了 $M$ 相对于 $R$ 有限表现的含义。
还证明了这一概念具有良好的局部化性质，并且可以粘合。
因此可以如下定义相应的整体概念。

\begin{definition}
\label{more-morphisms-definition-relatively-finitely-presented-sheaf}
设 $f : X \to S$ 是局部有限型的概形态射，$\mathcal{F}$ 是拟凝聚
$\mathcal{O}_X$-模。若存在仿射开覆盖 $S = \bigcup V_i$，并且对每个
$i$ 存在仿射开覆盖 $f^{-1}(V_i) = \bigcup_j U_{ij}$，使得
$\mathcal{F}(U_{ij})$ 作为 $\mathcal{O}_X(U_{ij})$-模相对于
$\mathcal{O}_S(V_i)$ 有限表现，则称 $\mathcal{F}$
{\it 相对于 $S$ 有限表现}，或称它{\it 在 $S$ 上相对有限表现}。
\end{definition}

\noindent
注意，这蕴含 $\mathcal{F}$ 是有限型 $\mathcal{O}_X$-模。
若 $X \to S$ 仅仅局部有限型，则即使 $X \to S$ 不是局部有限表现，
$\mathcal{F}$ 仍可能相对于 $S$ 有限表现。我们将看到，
$X \to S$ 局部有限表现，当且仅当 $\mathcal{O}_X$ 相对于 $S$ 有限表现。

\begin{lemma}
\label{more-morphisms-lemma-relative-finite-presentation-characterize}
设 $f : X \to S$ 是局部有限型的概形态射，$\mathcal{F}$ 是拟凝聚
$\mathcal{O}_X$-模。下列条件等价：
\begin{enumerate}
\item $\mathcal{F}$ 相对于 $S$ 有限表现；
\item 对每一对满足 $f(U) \subset V$ 的仿射开集 $U \subset X$、
$V \subset S$，$\mathcal{O}_X(U)$-模 $\mathcal{F}(U)$
相对于 $\mathcal{O}_S(V)$ 有限表现。
\end{enumerate}
此外，若这些条件成立，则对每一对满足 $f(U) \subset V$ 的开子概形
$U \subset X$、$V \subset S$，限制 $\mathcal{F}|_U$ 相对于 $V$ 有限表现。
\end{lemma}

\begin{proof}
最后一个陈述由 (1) 与 (2) 的等价性显然成立。(2) 蕴含 (1) 也显然。
假设 (1) 成立。取
$S = \bigcup V_i$ 和 $f^{-1}(V_i) = \bigcup U_{ij}$ 为
定义 \ref{more-morphisms-definition-relatively-finitely-presented-sheaf}.
中的仿射开覆盖。设 $U \subset X$ 和 $V \subset S$ 如 (2) 所述。
由 代数进阶章引理
\ref{more-algebra-lemma-glue-relative-finite-presentation}
，只需找到 $U$ 的标准开覆盖 $U = \bigcup U_k$，使得
$\mathcal{F}(U_k)$ 相对于 $\mathcal{O}_S(V)$ 有限表现。
换言之，对每个 $u \in U$，只需找到标准仿射开集
$u \in U' \subset U$，使得 $\mathcal{F}(U')$ 相对于
$\mathcal{O}_S(V)$ 有限表现。选取 $i$ 使 $f(u) \in V_i$，
再选取 $j$ 使 $u \in U_{ij}$。由
概形章引理 \ref{schemes-lemma-standard-open-two-affines}
，可以找到 $v \in V' \subset V \cap V_i$，使它在 $V'$ 与
$V_i$ 中都是标准仿射开集。于是 $f^{-1}V'  \cap U$ 与
$f^{-1}V' \cap U_{ij}$ 分别是 $U$ 与 $U_{ij}$ 的标准仿射开集。
再次应用该引理，可以找到
$u \in U' \subset f^{-1}V' \cap U \cap U_{ij}$，使它在
$f^{-1}V'  \cap U$ 与 $f^{-1}V' \cap U_{ij}$ 中都是标准仿射开集。
因此 $U'$ 也是 $U$ 与 $U_{ij}$ 的标准仿射开集。由 代数进阶章引理
\ref{more-algebra-lemma-localize-relative-finite-presentation}
，$\mathcal{F}(U_{ij})$ 相对于 $\mathcal{O}_S(V_i)$ 有限表现这一假设蕴含
$\mathcal{F}(U')$ 相对于 $\mathcal{O}_S(V_i)$ 有限表现。由于
$\mathcal{O}_X(U') =
\mathcal{O}_X(U') \otimes_{\mathcal{O}_S(V_i)} \mathcal{O}_S(V')$
，由 代数进阶章引理
\ref{more-algebra-lemma-base-change-relative-finite-presentation}
可知 $\mathcal{F}(U')$ 相对于 $\mathcal{O}_S(V')$ 有限表现。
再次应用 代数进阶章引理
\ref{more-algebra-lemma-localize-relative-finite-presentation}
，得到 $\mathcal{F}(U')$ 相对于 $\mathcal{O}_S(V)$ 有限表现。
证明完毕。
\end{proof}

\begin{lemma}
\label{more-morphisms-lemma-relative-finite-presentation}
设 $f : X \to S$ 是局部有限型的概形态射，$\mathcal{F}$ 是拟凝聚
$\mathcal{O}_X$-模。
\begin{enumerate}
\item 若 $f$ 局部有限表现，则 $\mathcal{F}$ 相对于 $S$ 有限表现，
当且仅当 $\mathcal{F}$ 有限表现。
\item 态射 $f$ 局部有限表现，当且仅当 $\mathcal{O}_X$
相对于 $S$ 有限表现。
\end{enumerate}
\end{lemma}

\begin{proof}
这直接由定义得到，见
代数进阶章定义
\ref{more-algebra-definition-relatively-finitely-presented} 后的讨论。
\end{proof}

\begin{lemma}
\label{more-morphisms-lemma-finite-morphism-relative-finite-presentation}
设 $\pi : X \to Y$ 是在基概形 $S$ 上局部有限型的概形之间的有限态射，
$\mathcal{F}$ 是拟凝聚 $\mathcal{O}_X$-模。则 $\mathcal{F}$
相对于 $S$ 有限表现，当且仅当 $\pi_*\mathcal{F}$ 相对于 $S$ 有限表现。
\end{lemma}

\begin{proof}
这是把
代数进阶章引理 \ref{more-algebra-lemma-finite-extension}
的结果翻译成概形语言。
\end{proof}

\begin{lemma}
\label{more-morphisms-lemma-base-change-relative-finite-presentation}
设 $f : X \to S$ 是局部有限型的概形态射，$\mathcal{F}$ 是拟凝聚
$\mathcal{O}_X$-模。设 $S' \to S$ 是概形态射，置 $X' = X \times_S S'$，
并以 $\mathcal{F}'$ 表示 $\mathcal{F}$ 到 $X'$ 的拉回。
若 $\mathcal{F}$ 相对于 $S$ 有限表现，则 $\mathcal{F}'$
相对于 $S'$ 有限表现。
\end{lemma}

\begin{proof}
这是把
代数进阶章引理
\ref{more-algebra-lemma-base-change-relative-finite-presentation}
的结果翻译成概形语言。
\end{proof}

\begin{lemma}
\label{more-morphisms-lemma-pull-relative-finite-presentation}
设 $X \to Y \to S$ 是局部有限型的概形态射，$\mathcal{G}$ 是拟凝聚
$\mathcal{O}_Y$-模。若 $f : X \to Y$ 局部有限表现，且
$\mathcal{G}$ 相对于 $S$ 有限表现，则 $f^*\mathcal{G}$
相对于 $S$ 有限表现。
\end{lemma}

\begin{proof}
这是把
代数进阶章引理
\ref{more-algebra-lemma-pull-relative-finite-presentation}
的结果翻译成概形语言。
\end{proof}

\begin{lemma}
\label{more-morphisms-lemma-composition-relative-finite-presentation}
设 $X \to Y \to S$ 是局部有限型的概形态射，$\mathcal{F}$ 是拟凝聚
$\mathcal{O}_X$-模。若 $Y \to S$ 局部有限表现，且 $\mathcal{F}$
相对于 $Y$ 有限表现，则 $\mathcal{F}$ 相对于 $S$ 有限表现。
\end{lemma}

\begin{proof}
这是把
代数进阶章引理
\ref{more-algebra-lemma-composition-relative-finite-presentation}
的结果翻译成概形语言。
\end{proof}

\begin{lemma}
\label{more-morphisms-lemma-ses-relatively-finite-presentation}
设 $X \to S$ 是局部有限型的概形态射。设
$0 \to \mathcal{F}' \to \mathcal{F} \to \mathcal{F}'' \to 0$
是拟凝聚 $\mathcal{O}_X$-模的短正合列。
\begin{enumerate}
\item 若 $\mathcal{F}', \mathcal{F}''$ 相对于 $S$ 有限表现，
则 $\mathcal{F}$ 也如此。
\item 若 $\mathcal{F}'$ 是有限型 $\mathcal{O}_X$-模，且
$\mathcal{F}$ 相对于 $S$ 有限表现，则 $\mathcal{F}''$
相对于 $S$ 有限表现。
\end{enumerate}
\end{lemma}

\begin{proof}
这是把
代数进阶章引理
\ref{more-algebra-lemma-ses-relatively-finite-presentation}
的结果翻译成概形语言。
\end{proof}

\begin{lemma}
\label{more-morphisms-lemma-sum-relatively-finite-presentation}
\begin{slogan}
模的直和项继承相对于基有限表现这一性质。
\end{slogan}
设 $X \to S$ 是局部有限型的概形态射，$\mathcal{F}, \mathcal{F}'$
是拟凝聚 $\mathcal{O}_X$-模。若 $\mathcal{F} \oplus \mathcal{F}'$
相对于 $S$ 有限表现，则 $\mathcal{F}$ 和 $\mathcal{F}'$ 也如此。
\end{lemma}

\begin{proof}
这是把
代数进阶章引理
\ref{more-algebra-lemma-sum-relatively-finite-presentation}
的结果翻译成概形语言。
\end{proof}










\section{相对伪凝聚性}
\label{more-morphisms-section-relative-pseudo-coherence}

\noindent
本节是
代数进阶章第 \ref{more-algebra-section-relative-pseudo-coherent} 节
在概形情形下的对应版本。我们强烈建议读者先阅读该节。
尽管本节的内容以及
上同调章第 \ref{cohomology-section-strictly-perfect} 节,
\ref{cohomology-section-pseudo-coherent},
\ref{cohomology-section-tor}, 及
\ref{cohomology-section-perfect}
中关于伪凝聚复形的内容，都是针对任意 $\mathcal{O}_X$-模复形建立的，
但若 $X$ 是概形，则只处理 $D_\QCoh(\mathcal{O}_X)$ 中的对象，
会极大简化许多引理和论证，常常能立即把手头的问题归结为其代数对应物。
此外，我们首先证明的结果之一是：相对伪凝聚性蕴含上同调层拟凝聚，
见 引理 \ref{more-morphisms-lemma-relative-pseudo-coherence}。
因此，初次阅读时，建议读者只处理 $D_\QCoh(\mathcal{O}_X)$ 中的对象。

\begin{lemma}
\label{more-morphisms-lemma-relatively-pseudo-coherent}
设 $X \to S$ 是仿射概形之间的有限型态射，$E$ 是
$D(\mathcal{O}_X)$ 的对象，并设 $m \in \mathbf{Z}$。
下列条件等价：
\begin{enumerate}
\item 对某个闭浸入 $i : X \to \mathbf{A}^n_S$，
$D(\mathcal{O}_{\mathbf{A}^n_S})$ 的对象 $Ri_*E$ 是 $m$-伪凝聚的；
\item 对所有闭浸入 $i : X \to \mathbf{A}^n_S$，
$D(\mathcal{O}_{\mathbf{A}^n_S})$ 的对象 $Ri_*E$ 都是 $m$-伪凝聚的。
\end{enumerate}
\end{lemma}

\begin{proof}
写 $S = \Spec(R)$、$X = \Spec(A)$。设 $i$ 对应于满射
$\alpha : R[x_1, \ldots, x_n] \to A$，而 $X \to \mathbf{A}^m_S$
对应于 $\beta : R[y_1, \ldots, y_m] \to A$。
选取满足 $\alpha(f_j) = \beta(y_j)$ 的
$f_j \in R[x_1, \ldots, x_n]$，以及满足
$\beta(g_i) = \alpha(x_i)$ 的 $g_i \in R[y_1, \ldots, y_m]$。
于是得到交换图
$$
\xymatrix{
R[x_1, \ldots, x_n, y_1, \ldots, y_m]
\ar[d]^{x_i \mapsto g_i} \ar[rr]_-{y_j \mapsto f_j} & &
R[x_1, \ldots, x_n] \ar[d] \\
R[y_1, \ldots, y_m] \ar[rr] & & A
}
$$
它对应于闭浸入的交换图
$$
\xymatrix{
\mathbf{A}^{n + m}_S & \mathbf{A}^n_S \ar[l] \\
\mathbf{A}^m_S \ar[u] & X \ar[u] \ar[l]
}
$$
因此，只需证明：对闭浸入
$$
f : \mathbf{A}^m_S \to \mathbf{A}^{n + m}_S
$$
，$D(\mathcal{O}_{\mathbf{A}^m_S})$ 的对象 $E$ 是 $m$-伪凝聚的，
当且仅当 $Rf_*E$ 是 $m$-伪凝聚的。这由
概形的导出范畴章引理
\ref{perfect-lemma-closed-push-pseudo-coherent}
以及 $f_*\mathcal{O}_{\mathbf{A}^m_S}$ 是伪凝聚
$\mathcal{O}_{\mathbf{A}^{n + m}_S}$-模这一事实得到。
$f_*\mathcal{O}_{\mathbf{A}^m_S}$ 的伪凝聚性很容易直接证明，
但它也可由
概形的导出范畴章引理 \ref{perfect-lemma-pseudo-coherent-affine}
和
代数进阶章引理 \ref{more-algebra-lemma-relatively-pseudo-coherent}.
得到。
\end{proof}

\noindent
回忆，若 $f : X \to S$ 是局部有限型的概形态射，则对每一对满足
$f(U) \subset V$ 的仿射开集 $U \subset X$、$V \subset S$，
环映射 $\mathcal{O}_S(V) \to \mathcal{O}_X(U)$ 都是有限型的
（态射章引理 \ref{morphisms-lemma-locally-finite-type-characterize}）。
因此总存在闭浸入 $U \to \mathbf{A}^n_V$，下面的定义是有意义的。

\begin{definition}
\label{more-morphisms-definition-relative-pseudo-coherence}
设 $f : X \to S$ 是局部有限型的概形态射。设 $E$ 是
$D(\mathcal{O}_X)$ 的对象，$\mathcal{F}$ 是 $\mathcal{O}_X$-模。
固定 $m \in \mathbf{Z}$。
\begin{enumerate}
\item 若存在仿射开覆盖 $S = \bigcup V_i$，且对每个 $i$
存在仿射开覆盖 $f^{-1}(V_i) = \bigcup U_{ij}$，使得对每一对
$(U_{ij} \to V_i, E|_{U_{ij}})$，
引理 \ref{more-morphisms-lemma-relatively-pseudo-coherent}
中的等价条件都成立，则称 $E${\it 相对于 $S$ 是 $m$-伪凝聚的}。
\item 若对所有 $m \in \mathbf{Z}$，$E$ 相对于 $S$ 都是
$m$-伪凝聚的，则称 $E${\it 相对于 $S$ 是伪凝聚的}。
\item 若把 $\mathcal{F}$ 看作 $D(\mathcal{O}_X)$ 的对象时，
它相对于 $S$ 是 $m$-伪凝聚的，则称 $\mathcal{F}$
{\it 相对于 $S$ 是 $m$-伪凝聚的}。
\item 若把 $\mathcal{F}$ 看作 $D(\mathcal{O}_X)$ 的对象时，
它相对于 $S$ 是伪凝聚的，则称 $\mathcal{F}$
{\it 相对于 $S$ 是伪凝聚的}。
\end{enumerate}
\end{definition}

\noindent
若 $X$ 拟紧，并且对某个 $m$，$E$ 相对于 $S$ 是 $m$-伪凝聚的，
则 $E$ 上有界。若 $E$ 相对于 $S$ 是伪凝聚的，
则 $E$ 具有拟凝聚上同调层。

\begin{lemma}
\label{more-morphisms-lemma-relative-pseudo-coherence}
设 $f : X \to S$ 是局部有限型的概形态射。若
$D(\mathcal{O}_X)$ 中的 $E$ 相对于 $S$ 是 $m$-伪凝聚的，
则对 $i > m$，$H^i(E)$ 是拟凝聚 $\mathcal{O}_X$-模。
若 $E$ 相对于 $S$ 是伪凝聚的，则 $E$ 是
$D_\QCoh(\mathcal{O}_X)$ 的对象。
\end{lemma}

\begin{proof}
选取仿射开覆盖 $S = \bigcup V_i$，并对每个 $i$ 选取仿射开覆盖
$f^{-1}(V_i) = \bigcup U_{ij}$，使得对每一对
$(U_{ij} \to V_i, E|_{U_{ij}})$，
引理 \ref{more-morphisms-lemma-relatively-pseudo-coherent}
中的等价条件都成立。由于拟凝聚性在 $X$ 上是局部的，
可以假设存在闭浸入 $i : X \to \mathbf{A}^n_S$，使得
$Ri_*E$ 在 $\mathbf{A}^n_S$ 上是 $m$-伪凝聚的。
由 概形的导出范畴章引理 \ref{perfect-lemma-pseudo-coherent}，
这意味着当 $q > m$ 时，$H^q(Ri_*E)$ 拟凝聚。
因为 $i_*$ 是正合函子，故
$i_*H^q(E) = H^q(Ri_*E)$ 在 $\mathbf{A}^n_S$ 上拟凝聚。
由 态射章引理 \ref{morphisms-lemma-i-star-equivalence}，
这蕴含 $H^q(E)$ 如所需是拟凝聚的（严格说来，它蕴含存在某个拟凝聚
$\mathcal{O}_X$-模 $\mathcal{F}$，使得
$i_*\mathcal{F} = i_*H^q(E)$；于是
模章引理 \ref{modules-lemma-i-star-equivalence}
告诉我们 $\mathcal{F} \cong H^q(E)$，从而得到结论）。
\end{proof}

\noindent
接下来证明，相对伪凝聚条件具有良好的局部化性质。

\begin{lemma}
\label{more-morphisms-lemma-localize-relative-pseudo-coherent}
设 $S$ 是仿射概形，$V \subset S$ 是标准开集，
$X \to V$ 是仿射概形之间的有限型态射，$U \subset X$ 是仿射开集，
$E$ 是 $D(\mathcal{O}_X)$ 的对象。若对 $(X \to V, E)$，
引理 \ref{more-morphisms-lemma-relatively-pseudo-coherent}
中的等价条件成立，则对 $(U \to S, E|_U)$，
引理 \ref{more-morphisms-lemma-relatively-pseudo-coherent}
中的等价条件也成立。
\end{lemma}

\begin{proof}
写 $S = \Spec(R)$、$V = D(f)$、$X = \Spec(A)$、$U = D(g)$。
假设对 $(X \to V, E)$，
引理 \ref{more-morphisms-lemma-relatively-pseudo-coherent}
中的等价条件成立。

\medskip\noindent
选取满射 $R_f[x_1, \ldots, x_n] \to A$。写
$R_f = R[x_0]/(fx_0 - 1)$。于是 $R[x_0, x_1, \ldots, x_n] \to A$
是满射，且 $R_f[x_1, \ldots, x_n]$ 作为
$R[x_0, \ldots, x_n]$-模是伪凝聚的。因此有
$$
X \to \mathbf{A}^n_V \to \mathbf{A}^{n + 1}_S
$$
，并可应用
概形的导出范畴章
引理 \ref{perfect-lemma-closed-push-pseudo-coherent}
断定 $E$ 到 $\mathbf{A}^{n + 1}_S$ 的推前 $E'$ 是 $m$-伪凝聚的。

\medskip\noindent
选取映到 $g \in A$ 的元素 $g' \in R[x_0, x_1, \ldots, x_n]$。
考虑满射
$R[x_0, \ldots, x_{n + 1}] \to R[x_0, \ldots, x_n, 1/g']$.
得到
$$
\xymatrix{
X \ar[d] & U \ar[d] \ar[l] \ar[dr] \\
\mathbf{A}^{n + 1}_S & D(g')\ar[l] \ar[r] & \mathbf{A}^{n + 2}_S
}
$$
其中左下箭头是开浸入，右下箭头是闭浸入。与前面一样可断定，
$E'|_{D(g')}$ 到 $\mathbf{A}^{n + 2}_S$ 的推前是 $m$-伪凝聚的。
由于这也正是 $E|_U$ 到 $\mathbf{A}^{n + 2}_S$ 的推前，引理得证。
\end{proof}

\begin{lemma}
\label{more-morphisms-lemma-glue-relative-pseudo-coherent}
设 $X \to S$ 是仿射概形之间的有限型态射，$E$ 是
$D(\mathcal{O}_X)$ 的对象，$m \in \mathbf{Z}$。
设 $X = \bigcup U_i$ 是标准仿射开覆盖。下列条件等价：
\begin{enumerate}
\item 对所有 $(U_i \to S, E|_{U_i})$，
引理 \ref{more-morphisms-lemma-relatively-pseudo-coherent}
中的等价条件成立；
\item 对 $(X \to S, E)$，
引理 \ref{more-morphisms-lemma-relatively-pseudo-coherent}
中的等价条件成立。
\end{enumerate}
\end{lemma}

\begin{proof}
蕴含关系 (2) $\Rightarrow$ (1) 就是
引理 \ref{more-morphisms-lemma-localize-relative-pseudo-coherent}.
假设 (1)。写 $S = \Spec(R)$、$X = \Spec(A)$、$U_i = D(f_i)$，
并在 $A$ 中写 $1 = \sum f_ig_i$。考虑满射
$$
R[x_i, y_i, z_i] \to R[x_i, y_i, z_i]/(\sum y_iz_i - 1) \to A.
$$
它把 $y_i$ 映到 $f_i$，把 $z_i$ 映到 $g_i$。注意
$R[x_i, y_i, z_i]/(\sum y_iz_i - 1)$ 作为
$R[x_i, y_i, z_i]$-模是伪凝聚的。因此，只需证明 $E$ 到
$T = \Spec(R[x_i, y_i, z_i]/(\sum y_iz_i - 1))$ 的推前是
$m$-伪凝聚的，见
概形的导出范畴章
引理 \ref{perfect-lemma-closed-push-pseudo-coherent}.
对每个 $i_0$，只需证明此推前到
$W_{i_0} = \Spec(R[x_i, y_i, z_i, 1/y_{i_0}]/(\sum y_iz_i - 1))$
的限制是 $m$-伪凝聚的。注意有交换图
$$
\xymatrix{
X \ar[d] & U_{i_0} \ar[l] \ar[d] \\
T & W_{i_0} \ar[l]
}
$$
它蕴含 $E$ 到 $T$ 的推前在 $W_{i_0}$ 上的限制，正是
$E|_{U_{i_0}}$ 到 $W_{i_0}$ 的推前。由于\linebreak[4]
$R[x_i, y_i, z_i, 1/y_{i_0}]/(\sum y_iz_i - 1)$
同构于 $R$ 上的多项式环，这就证明了所需结论。
\end{proof}

\begin{lemma}
\label{more-morphisms-lemma-relative-pseudo-coherence-characterize}
设 $f : X \to S$ 是局部有限型的概形态射，$E$ 是
$D(\mathcal{O}_X)$ 的对象。固定 $m \in \mathbf{Z}$。
下列条件等价：
\begin{enumerate}
\item $E$ 相对于 $S$ 是 $m$-伪凝聚的；
\item 对每一对满足 $f(U) \subset V$ 的仿射开集 $U \subset X$、
$V \subset S$，对 $(U \to V, E|_U)$，
引理 \ref{more-morphisms-lemma-relatively-pseudo-coherent}
中的等价条件成立。
\end{enumerate}
此外，若这些条件成立，则对每一对满足 $f(U) \subset V$
的开子概形 $U \subset X$、$V \subset S$，限制 $E|_U$
相对于 $V$ 是 $m$-伪凝聚的。
\end{lemma}

\begin{proof}
最后一个陈述由 (1) 与 (2) 的等价性显然成立。(2) 蕴含 (1) 也显然。
假设 (1) 成立。取 $S = \bigcup V_i$ 和
$f^{-1}(V_i) = \bigcup U_{ij}$ 为
定义 \ref{more-morphisms-definition-relative-pseudo-coherence}.
中的仿射开覆盖。设 $U \subset X$ 和 $V \subset S$ 如 (2) 所述。
由 引理 \ref{more-morphisms-lemma-glue-relative-pseudo-coherent}，
只需找到 $U$ 的标准开覆盖 $U = \bigcup U_k$，使得对每一对
$(U_k \to V, E|_{U_k})$，
引理 \ref{more-morphisms-lemma-relatively-pseudo-coherent}
中的等价条件成立。换言之，对每个 $u \in U$，只需找到标准仿射开集
$u \in U' \subset U$，使得对 $(U' \to V, E|_{U'})$，
引理 \ref{more-morphisms-lemma-relatively-pseudo-coherent}
中的等价条件成立。选取 $i$ 使 $f(u) \in V_i$，再选取 $j$
使 $u \in U_{ij}$。由
概形章引理 \ref{schemes-lemma-standard-open-two-affines}
，可以找到 $v \in V' \subset V \cap V_i$，使它在 $V'$ 与 $V_i$
中都是标准仿射开集。于是 $f^{-1}V'  \cap U$ 与
$f^{-1}V' \cap U_{ij}$ 分别是 $U$ 与 $U_{ij}$ 的标准仿射开集。
再次应用该引理，可以找到
$u \in U' \subset f^{-1}V' \cap U \cap U_{ij}$，使它在
$f^{-1}V'  \cap U$ 与 $f^{-1}V' \cap U_{ij}$ 中都是标准仿射开集。
因此 $U'$ 也是 $U$ 与 $U_{ij}$ 的标准仿射开集。
由 引理 \ref{more-morphisms-lemma-localize-relative-pseudo-coherent}，
对 $(U_{ij} \to V_i, E|_{U_{ij}})$，
引理 \ref{more-morphisms-lemma-relatively-pseudo-coherent}
中的等价条件成立这一假设，蕴含对 $(U' \to V, E|_{U'})$，
引理 \ref{more-morphisms-lemma-relatively-pseudo-coherent}
中的等价条件成立。
\end{proof}

\noindent
对于上同调层拟凝聚的导出范畴对象，可以把相对 $m$-伪凝聚性
与 代数进阶章定义
\ref{more-algebra-definition-relatively-pseudo-coherent}.
中定义的概念联系起来。我们将使用如下事实：对仿射概形
$U = \Spec(A)$，函子 $R\Gamma(U, -)$ 诱导
$D_\QCoh(\mathcal{O}_U)$ 与 $D(A)$ 之间的等价，见
概形的导出范畴章引理
\ref{perfect-lemma-affine-compare-bounded}.
此函子与拉回相容：若 $E$ 是 $D_\QCoh(\mathcal{O}_U)$ 的对象，
而 $A \to B$ 是对应于仿射概形态射
$g : V = \Spec(B) \to \Spec(A) = U$ 的环映射，则
$R\Gamma(V, Lg^*E) = R\Gamma(U, E) \otimes_A^\mathbf{L} B$.
见 概形的导出范畴章引理
\ref{perfect-lemma-quasi-coherence-pullback}.

\begin{lemma}
\label{more-morphisms-lemma-qcoh-relative-pseudo-coherence-characterize}
设 $f : X \to S$ 是局部有限型的概形态射，$E$ 是
$D_\QCoh(\mathcal{O}_X)$ 的对象。固定 $m \in \mathbf{Z}$。
下列条件等价：
\begin{enumerate}
\item $E$ 相对于 $S$ 是 $m$-伪凝聚的；
\item 存在仿射开覆盖 $S = \bigcup V_i$，且对每个 $i$
存在仿射开覆盖 $f^{-1}(V_i) = \bigcup U_{ij}$，使
$\mathcal{O}_X(U_{ij})$-模复形 $R\Gamma(U_{ij}, E)$
相对于 $\mathcal{O}_S(V_i)$ 是 $m$-伪凝聚的；
\item 对每一对满足 $f(U) \subset V$ 的仿射开集 $U \subset X$、
$V \subset S$，$\mathcal{O}_X(U)$-模复形 $R\Gamma(U, E)$
相对于 $\mathcal{O}_S(V)$ 是 $m$-伪凝聚的。
\end{enumerate}
\end{lemma}

\begin{proof}
设 $U$ 和 $V$ 如 (2) 所述，并选取闭浸入
$i : U \to \mathbf{A}^n_V$。利用
引理 \ref{more-morphisms-lemma-relative-pseudo-coherence-characterize} 作形式论证，
可见只需证明：$Ri_*(E|_U)$ 是 $m$-伪凝聚的，当且仅当
$R\Gamma(U, E)$ 相对于 $\mathcal{O}_S(V)$ 是 $m$-伪凝聚的。
写 $U = \Spec(A)$、$V = \Spec(R)$，并且
% P05-EN-0127: frozen source omits the closing parenthesis in this Spec expression.
$\mathbf{A}^n_V = \Spec(R[x_1, \ldots, x_n]$. 由上述说明
由引理前的说明，$E|_U$ 拟同构于 $U$ 上由 $D(A)$ 的对象
$R\Gamma(U, E)$ 所伴随的拟凝聚层复形。注意，因为 $i$ 是闭浸入
（因而 $i_*$ 正合），有
$R\Gamma(U, E) = R\Gamma(\mathbf{A}^n_V, Ri_*(E|_U))$。
于是 $Ri_*E$ 由视为 $D(R[x_1, \ldots, x_n])$ 对象的
$R\Gamma(U, E)$ 所伴随。由
概形的导出范畴章引理 \ref{perfect-lemma-pseudo-coherent-affine}
，$Ri_*(E|_U)$ 的 $m$-伪凝聚性等价于 $R\Gamma(U, E)$
在 $D(R[x_1, \ldots, x_n])$ 中的 $m$-伪凝聚性；按定义，
后者又等价于 $R\Gamma(U, E)$ 相对于 $R = \mathcal{O}_S(V)$
是 $m$-伪凝聚的。结论得证。
\end{proof}

\begin{lemma}
\label{more-morphisms-lemma-closed-morphism-relative-pseudo-coherence}
设 $i : X \to Y$ 是基概形 $S$ 上局部有限型的概形之间的态射。
假设 $i$ 诱导 $X$ 到 $Y$ 的一个闭子集的同胚。
设 $E$ 是 $D(\mathcal{O}_X)$ 的对象。则 $E$ 相对于 $S$
是 $m$-伪凝聚的，当且仅当 $Ri_*E$ 相对于 $S$ 是 $m$-伪凝聚的。
\end{lemma}

\begin{proof}
由 态射章引理 \ref{morphisms-lemma-homeomorphism-affine}，
态射 $i$ 是仿射的。因此可以假设 $S$、$Y$、$X$ 都仿射。
写 $S = \Spec(R)$、$Y = \Spec(A)$、$X = \Spec(B)$。
该条件意味着 $A/\text{rad}(A) \to B/\text{rad}(B)$ 满射；
% P05-EN-0128: frozen source calls rad(A), rad(B) Jacobson radicals although the nilpotence argument uses nilradicals.
这里 $\text{rad}(A)$ 和 $\text{rad}(B)$ 分别表示 $A$ 和 $B$
的 Jacobson 根。由于 $B$ 在 $A$ 上为有限型，可以找到
$b_1, \ldots, b_m \in \text{rad}(B)$，它们生成 $B$ 作为
$A$-代数。设对所有 $j$ 都有 $b_j^N = 0$。考虑环的图
$$
\xymatrix{
B & R[x_i, y_j]/(y_j^N) \ar[l] & R[x_i, y_j] \ar[l] \\
A \ar[u] & R[x_i] \ar[l] \ar[u] \ar[ru]
}
$$
它转化为仿射概形的图
$$
\xymatrix{
X \ar[d] \ar[r] & T \ar[d] \ar[r] & \mathbf{A}^{n + m}_S \ar[ld] \\
Y \ar[r] & \mathbf{A}^n_S
}
$$
。由 引理 \ref{more-morphisms-lemma-relative-pseudo-coherence-characterize}，
$E$ 相对于 $S$ 是 $m$-伪凝聚的，当且仅当它到
$\mathbf{A}^{n + m}_S$ 的推前是 $m$-伪凝聚的。由
概形的导出范畴章引理
\ref{perfect-lemma-closed-push-pseudo-coherent}
，这又等价于它到 $T$ 的推前是 $m$-伪凝聚的。同一引理表明，
后者等价于它到 $\mathbf{A}^n_S$ 的推前是 $m$-伪凝聚的。
再次由
引理 \ref{more-morphisms-lemma-relative-pseudo-coherence-characterize}
，这等价于 $Ri_*E$ 相对于 $S$ 是 $m$-伪凝聚的。
\end{proof}

\begin{lemma}
\label{more-morphisms-lemma-finite-morphism-relative-pseudo-coherence}
设 $\pi : X \to Y$ 是基概形 $S$ 上局部有限型的概形之间的有限态射。
设 $E$ 是 $D_\QCoh(\mathcal{O}_X)$ 的对象。则 $E$ 相对于 $S$
是 $m$-伪凝聚的，当且仅当 $R\pi_*E$ 相对于 $S$ 是 $m$-伪凝聚的。
\end{lemma}

\begin{proof}
这是把
代数进阶章引理
\ref{more-algebra-lemma-finite-extension-pseudo-coherent}
的结果翻译成概形语言。注意，由 概形的导出范畴章引理
\ref{perfect-lemma-quasi-coherence-direct-image}.
，$R\pi_*$ 确实把 $D_\QCoh(\mathcal{O}_X)$ 映入
$D_\QCoh(\mathcal{O}_Y)$。在这一翻译中使用
引理 \ref{more-morphisms-lemma-relative-pseudo-coherence-characterize}.
\end{proof}

\begin{lemma}
\label{more-morphisms-lemma-cone-relatively-pseudo-coherent}
设 $f : X \to S$ 是局部有限型的概形态射。
设 $(E, E', E'')$ 是 $D(\mathcal{O}_X)$ 的可区别三角，
并设 $m \in \mathbf{Z}$。
\begin{enumerate}
\item 若 $E$ 相对于 $S$ 是 $(m + 1)$-伪凝聚的，且 $E'$
相对于 $S$ 是 $m$-伪凝聚的，则 $E''$ 相对于 $S$ 是 $m$-伪凝聚的。
\item 若 $E, E''$ 相对于 $S$ 是 $m$-伪凝聚的，则 $E'$
相对于 $S$ 是 $m$-伪凝聚的。
\item 若 $E'$ 相对于 $S$ 是 $(m + 1)$-伪凝聚的，且 $E''$
相对于 $S$ 是 $m$-伪凝聚的，则 $E$ 相对于 $S$
是 $(m + 1)$-伪凝聚的。
\end{enumerate}
此外，若 $E, E', E''$ 中有两个相对于 $S$ 是伪凝聚的，
则第三个也如此。
\end{lemma}

\begin{proof}
这直接由 引理 \ref{more-morphisms-lemma-relative-pseudo-coherence-characterize} 和
上同调章引理 \ref{cohomology-lemma-cone-pseudo-coherent}.
得到。
\end{proof}

\begin{lemma}
\label{more-morphisms-lemma-rel-n-pseudo-module}
设 $X \to S$ 是局部有限型的概形态射，$\mathcal{F}$ 是
$\mathcal{O}_X$-模。则
\begin{enumerate}
\item 对所有 $m > 0$，$\mathcal{F}$ 相对于 $S$ 都是 $m$-伪凝聚的；
\item $\mathcal{F}$ 相对于 $S$ 是 $0$-伪凝聚的，当且仅当
$\mathcal{F}$ 是有限型 $\mathcal{O}_X$-模；
\item $\mathcal{F}$ 相对于 $S$ 是 $(-1)$-伪凝聚的，当且仅当
$\mathcal{F}$ 拟凝聚且相对于 $S$ 有限表现。
\end{enumerate}
\end{lemma}

\begin{proof}
第 (1) 部分直接由定义得到。为证明第 (3) 部分，可以在 $X$ 上局部工作
（两个性质都是局部的）。因此可以假设 $X$ 和 $S$ 仿射。
选取闭浸入 $i : X \to \mathbf{A}^n_S$。于是 $\mathcal{F}$
相对于 $S$ 是 $(-1)$-伪凝聚的，当且仅当 $i_*\mathcal{F}$
是 $(-1)$-伪凝聚的；而这又当且仅当 $i_*\mathcal{F}$
是有限表现的 $\mathcal{O}_{\mathbf{A}^n_S}$-模，见
上同调章引理 \ref{cohomology-lemma-finite-cohomology}.
有限表现的模拟凝聚，见
模章引理 \ref{modules-lemma-finite-presentation-quasi-coherent}.
由 态射章引理 \ref{morphisms-lemma-i-star-equivalence}，
$\mathcal{F}$ 拟凝聚，当且仅当 $i_*\mathcal{F}$ 拟凝聚。
至此第 (3) 部分得证。第 (2) 部分的证明类似但更简单。
\end{proof}

\begin{lemma}
\label{more-morphisms-lemma-summands-relative-pseudo-coherent}
设 $X \to S$ 是局部有限型的概形态射，$m \in \mathbf{Z}$，
$E, K$ 是 $D(\mathcal{O}_X)$ 的对象。若 $E \oplus K$
相对于 $S$ 是 $m$-伪凝聚的，则 $E$ 和 $K$ 也如此。
\end{lemma}

\begin{proof}
这由
上同调章引理 \ref{cohomology-lemma-summands-pseudo-coherent}
及定义得到。
\end{proof}

\begin{lemma}
\label{more-morphisms-lemma-complex-relative-pseudo-coherent-modules}
设 $X \to S$ 是局部有限型的概形态射，$m \in \mathbf{Z}$。
设 $\mathcal{F}^\bullet$ 是（局部）上有界的 $\mathcal{O}_X$-模复形，
使得对所有 $i$，$\mathcal{F}^i$ 相对于 $S$ 都是
$(m - i)$-伪凝聚的。则 $\mathcal{F}^\bullet$ 相对于 $S$
是 $m$-伪凝聚的。
\end{lemma}

\begin{proof}
这由
上同调章引理 \ref{cohomology-lemma-complex-pseudo-coherent-modules}
及定义得到。
\end{proof}

\begin{lemma}
\label{more-morphisms-lemma-cohomology-relative-pseudo-coherent}
设 $X \to S$ 是局部有限型的概形态射，$m \in \mathbf{Z}$，
$E$ 是 $D(\mathcal{O}_X)$ 的对象。若 $E$（局部）上有界，
并且对所有 $i$，$H^i(E)$ 相对于 $S$ 都是 $(m - i)$-伪凝聚的，
则 $E$ 相对于 $S$ 是 $m$-伪凝聚的。
\end{lemma}

\begin{proof}
这由
上同调章引理 \ref{cohomology-lemma-cohomology-pseudo-coherent}
及定义得到。
\end{proof}

\begin{lemma}
\label{more-morphisms-lemma-base-change-relative-pseudo-coherent}
设 $X \to S$ 是局部有限型的概形态射，$m \in \mathbf{Z}$。
设 $E$ 是 $D(\mathcal{O}_X)$ 的对象，且相对于 $S$ 是 $m$-伪凝聚的。
设 $S' \to S$ 是概形态射。置 $X' = X \times_S S'$，并以 $E'$
表示 $E$ 到 $X'$ 的导出拉回。若 $S'$ 与 $X$ 在 $S$ 上 Tor 独立，
则 $E'$ 相对于 $S'$ 是 $m$-伪凝聚的。
\end{lemma}

\begin{proof}
问题在 $X$ 与 $X'$ 上是局部的，故可假设 $X$、$S$、$S'$、$X'$
都是仿射的。选取闭浸入 $i : X \to \mathbf{A}^n_S$，并以
$i' : X' \to \mathbf{A}^n_{S'}$ 表示它到 $S'$ 的基变换。
以 $g : X' \to X$ 和 $g' : \mathbf{A}^n_{S'} \to \mathbf{A}^n_S$
表示投影，于是 $E' = Lg^*E$。由于 $X$ 与 $S'$ 在 $S$ 上 Tor 独立，
基变换映射
(上同调章说明 \ref{cohomology-remark-base-change})
诱导同构
$$
Ri'_*(Lg^*E) = L(g')^*Ri_*E
$$
。具体而言，对位于 $x \in X$ 上方的点 $x' \in X'$，
基变换映射在 $x'$ 处的茎上的映射为
$$
E_x \otimes_{\mathcal{O}_{\mathbf{A}^n_S, x}}^\mathbf{L}
\mathcal{O}_{\mathbf{A}^n_{S'}, x'}
\longrightarrow
E_x \otimes_{\mathcal{O}_{X, x}}^\mathbf{L}
\mathcal{O}_{X', x'}
$$
，它来自闭浸入 $i$ 与 $i'$。注意，源拟同构于
$E_x \otimes_{\mathcal{O}_{S, s}}^\mathbf{L} \mathcal{O}_{S', s'}$
的一个局部化；它同构于目标，因为根据假设，
$\mathcal{O}_{X', x'}$ 同构于
$\mathcal{O}_{X, x} \otimes_{\mathcal{O}_{S, s}}^\mathbf{L}
\mathcal{O}_{S', s'}$ 的（同一个）局部化。应用
上同调章引理 \ref{cohomology-lemma-pseudo-coherent-pullback}.
即得引理。
\end{proof}

\begin{lemma}
\label{more-morphisms-lemma-pull-relative-pseudo-coherent}
设 $f : X \to Y$ 是基 $S$ 上局部有限型的概形态射。
设 $m \in \mathbf{Z}$，$E$ 是 $D(\mathcal{O}_Y)$ 的对象。假设
\begin{enumerate}
\item $\mathcal{O}_X$ 相对于 $Y$ 是伪凝聚的\footnote{这意味着
$f$ 是伪凝聚态射，见 定义 \ref{more-morphisms-definition-pseudo-coherent}。}；
\item $E$ 相对于 $S$ 是 $m$-伪凝聚的。
\end{enumerate}
则 $Lf^*E$ 相对于 $S$ 是 $m$-伪凝聚的。
\end{lemma}

\begin{proof}
问题在 $X$ 上是局部的。因此可以假设 $X$、$Y$、$S$ 都仿射。
仿照 代数进阶章引理
\ref{more-algebra-lemma-pull-relative-pseudo-coherent}
的证明进行论证，可以找到交换图
$$
\xymatrix{
X \ar[r]_i \ar[d]_f &
\mathbf{A}^d_Y \ar[r]_j \ar[ld]^p &
\mathbf{A}^{n + d}_S \ar[ld] \\
Y \ar[r] &
\mathbf{A}^n_S
}
$$
注意
$$
Ri_* Lf^*E = Ri_* Li^* Lp^*E =
Lp^*E \otimes_{\mathcal{O}_{\mathbf{A}_Y^n}}^\mathbf{L} Ri_*\mathcal{O}_X
$$
，这是由 上同调章引理
\ref{cohomology-lemma-projection-formula-closed-immersion}.
得到的。由假设以及 $Y$ 仿射这一事实，可以用有限自由
% P05-EN-0129: frozen source switches the diagram's relative affine dimension d to n in coefficient sheaves.
$\mathcal{O}_{\mathbf{A}_Y^n}$-模复形 $\mathcal{F}^\bullet$
表示 $Ri_*\mathcal{O}_X = i_*\mathcal{O}_X$，并且
$\mathcal{F}^q = 0$ 对 $q > 0$ 成立
（细节略；使用 概形的导出范畴章引理
\ref{perfect-lemma-pseudo-coherent-affine}
和
代数进阶章引理
\ref{more-algebra-lemma-rel-n-pseudo-module}）。
由假设，$E$ 上有界；设对 $q > a$ 有 $H^q(E) = 0$。
用满足对 $q > a$ 有 $\mathcal{E}^q = 0$ 的 $\mathcal{O}_Y$-模复形
$\mathcal{E}^\bullet$ 表示 $E$。于是上述导出张量积由
$\text{Tot}(p^*\mathcal{E}^\bullet
\otimes_{\mathcal{O}_{\mathbf{A}_Y^n}} \mathcal{F}^\bullet)$ 表示。

\medskip\noindent
因为 $j$ 是闭浸入，函子 $j_*$ 正合，而且可以通过对任一代表层复形
应用 $j_*$ 来计算 $Rj_*$。因此必须证明
% P05-EN-0130: frozen source uses ambient dimension n+m here although the diagram has n+d.
$j_*\text{Tot}(p^*\mathcal{E}^\bullet \otimes_{\mathcal{O}_{\mathbf{A}_Y^n}}
\mathcal{F}^\bullet)$ 作为 $\mathcal{O}_{\mathbf{A}^{n + m}_S}$-模复形
是 $m$-伪凝聚的。注意
$\text{Tot}(p^*\mathcal{E}^\bullet \otimes_{\mathcal{O}_{\mathbf{A}_Y^n}}
\mathcal{F}^\bullet)$ 具有由子复形构成的滤过，其相继商为复形
$p^*\mathcal{E}^\bullet
\otimes_{\mathcal{O}_{\mathbf{A}_Y^n}} \mathcal{F}^q[-q]$。
注意，当 $q \ll 0$ 时，复形
$p^*\mathcal{E}^\bullet \otimes_{\mathcal{O}_{\mathbf{A}_Y^n}}
\mathcal{F}^q[-q]$
在次数 $\leq m$ 上的上同调为零，因而是 $m$-伪凝聚的。因此，应用
引理 \ref{more-morphisms-lemma-cone-relatively-pseudo-coherent}
并作归纳，只需证明
% P05-EN-0131: frozen source asks for full pseudo-coherence here but the next reduction proves only m-pseudo-coherence.
$p^*\mathcal{E}^\bullet \otimes_{\mathcal{O}_{\mathbf{A}_Y^n}}
\mathcal{F}^q[-q]$
对所有 $q$ 都相对于 $S$ 是伪凝聚的。注意，对 $q > 0$
有 $\mathcal{F}^q = 0$。又因为 $\mathcal{F}^q$ 有限自由，
这就归结为证明 $p^*\mathcal{E}^\bullet$ 相对于 $S$
是 $m$-伪凝聚的；例如，这由
引理 \ref{more-morphisms-lemma-base-change-relative-pseudo-coherent}
得到。
\end{proof}

\begin{lemma}
\label{more-morphisms-lemma-composition-relative-pseudo-coherent}
设 $f : X \to Y$ 是基 $S$ 上局部有限型的概形态射。
设 $m \in \mathbf{Z}$，$E$ 是 $D(\mathcal{O}_X)$ 的对象。
假设 $\mathcal{O}_Y$ 相对于 $S$ 是伪凝聚的\footnote{这意味着
$Y \to S$ 是伪凝聚态射，见 定义 \ref{more-morphisms-definition-pseudo-coherent}。}。
则下列条件等价：
\begin{enumerate}
\item $E$ 相对于 $Y$ 是 $m$-伪凝聚的；
\item $E$ 相对于 $S$ 是 $m$-伪凝聚的。
\end{enumerate}
\end{lemma}

\begin{proof}
问题在 $X$ 上是局部的，故可假设 $X$、$Y$、$S$ 都仿射。
仿照 代数进阶章引理
\ref{more-algebra-lemma-pull-relative-pseudo-coherent}
的证明进行论证，可以找到交换图
$$
\xymatrix{
X \ar[r]_i \ar[d]_f &
\mathbf{A}^m_Y \ar[r]_j \ar[ld]^p &
\mathbf{A}^{n + m}_S \ar[ld] \\
Y \ar[r] &
\mathbf{A}^n_S
}
$$
假设 $\mathcal{O}_Y$ 相对于 $S$ 是伪凝聚的，这蕴含
$\mathcal{O}_{\mathbf{A}^m_Y}$ 相对于 $\mathbf{A}^m_S$ 是伪凝聚的
（通过平坦基变换；例如可用
引理 \ref{more-morphisms-lemma-base-change-relative-pseudo-coherent} 看出）。
% P05-EN-0132: frozen source writes A^n_Y here although j has domain A^m_Y in the displayed diagram.
这又蕴含 $j_*\mathcal{O}_{\mathbf{A}^n_Y}$ 作为
$\mathcal{O}_{\mathbf{A}^{n + m}_S}$-模是伪凝聚的。于是引理中的等价性由
概形的导出范畴章引理
\ref{perfect-lemma-closed-push-pseudo-coherent}.
得到。
\end{proof}

\begin{lemma}
\label{more-morphisms-lemma-check-relative-pseudo-coherence-on-charts}
设
$$
\xymatrix{
X \ar[rd] \ar[rr]_i & & P \ar[ld] \\
& S
}
$$
是概形的交换图。假设 $i$ 是闭浸入，且 $P \to S$ 平坦并局部有限表现。
设 $E$ 是 $D(\mathcal{O}_X)$ 的对象。则下列条件等价：
\begin{enumerate}
\item $E$ 相对于 $S$ 是 $m$-伪凝聚的；
\item $Ri_*E$ 相对于 $S$ 是 $m$-伪凝聚的；
\item $Ri_*E$ 在 $P$ 上是 $m$-伪凝聚的。
\end{enumerate}
\end{lemma}

\begin{proof}
条件 (1) 与 (2) 的等价性就是
引理 \ref{more-morphisms-lemma-finite-morphism-relative-pseudo-coherence}.
条件 (2) 与 (3) 的等价性可由
引理
\ref{more-morphisms-lemma-composition-relative-pseudo-coherent}
应用于 $\text{id} : P \to P$ 得到，只须证明 $\mathcal{O}_P$
相对于 $S$ 是伪凝聚的。这由
代数进阶章引理
\ref{more-algebra-lemma-flat-finite-presentation-perfect}
及定义得到。
\end{proof}









\section{伪凝聚态射}
\label{more-morphisms-section-pseudo-coherent}

\noindent
务必避免阅读本节。如果确实需要其中的一些材料，请先阅读相应的代数章节，见
代数进阶章第 \ref{more-algebra-section-pseudo-coherent} 节,
\ref{more-algebra-section-relative-pseudo-coherent}, 及
\ref{more-algebra-section-pseudo-coherent-perfect-ring-map}.
目前只需知道：环映射 $A \to B$ 伪凝聚，当且仅当
$B = A[x_1, \ldots, x_n]/I$，并且 $B$ 作为
$A[x_1, \ldots, x_n]$-模具有由有限自由
$A[x_1, \ldots, x_n]$-模构成的分解。

\begin{lemma}
\label{more-morphisms-lemma-pseudo-coherent}
设 $f : X \to S$ 是概形态射。下列条件等价：
\begin{enumerate}
\item 存在仿射开覆盖 $S = \bigcup V_j$，且对每个 $j$
存在仿射开覆盖 $f^{-1}(V_j) = \bigcup U_{ji}$，使得
% P05-EN-0133: frozen source reverses the covering index U_ji to U_ij in the ring-map clause.
$\mathcal{O}_S(V_j) \to \mathcal{O}_X(U_{ij})$ 是伪凝聚环映射；
\item 对每一对满足 $f(U) \subset V$ 的仿射开集 $U \subset X$、
$V \subset S$，环映射 $\mathcal{O}_S(V) \to \mathcal{O}_X(U)$ 伪凝聚；
\item $f$ 局部有限型，且 $\mathcal{O}_X$ 相对于 $S$ 是伪凝聚的。
\end{enumerate}
\end{lemma}

\begin{proof}
为证明 (1) 与 (2) 等价，只需对伪凝聚环映射这一性质检验
态射章定义 \ref{morphisms-definition-property-local}
的条件 (1)(a)、(b)、(c)。这些性质由下列引理得到
（其中用到局部化是平坦的）：
代数进阶章引理
\ref{more-algebra-lemma-base-change-relative-pseudo-coherent},
\ref{more-algebra-lemma-localize-relative-pseudo-coherent}, 及
\ref{more-algebra-lemma-glue-relative-pseudo-coherent}.

\medskip\noindent
若 (1) 成立，则 $f$ 局部有限型，因为伪凝聚环映射按定义是有限型的。
此外，由 引理 \ref{more-morphisms-lemma-qcoh-relative-pseudo-coherence-characterize}
及定义，(1) 蕴含 $\mathcal{O}_X$ 相对于 $S$ 是伪凝聚的。
反之，若 (3) 成立，则对 (2) 中的每一对 $U$ 和 $V$，
环 $\mathcal{O}_X(U)$ 在 $\mathcal{O}_S(V)$ 上为有限型，
并且 $\mathcal{O}_X(U)$ 作为模相对于 $\mathcal{O}_S(V)$ 是伪凝聚的，见
引理 \ref{more-morphisms-lemma-relative-pseudo-coherence-characterize} 及
\ref{more-morphisms-lemma-qcoh-relative-pseudo-coherence-characterize}.
这正是伪凝聚环映射的定义，故 (2) 和 (1) 成立。
\end{proof}

\begin{definition}
\label{more-morphisms-definition-pseudo-coherent}
若概形态射 $f : X \to S$ 满足
引理 \ref{more-morphisms-lemma-pseudo-coherent}
中的等价条件，则称它是{\it 伪凝聚的}。在这种情形下，
也称 $X$ 在 $S$ 上伪凝聚。
\end{definition}

\noindent
注意，伪凝聚态射的基变换一般不伪凝聚。

\begin{lemma}
\label{more-morphisms-lemma-flat-base-change-pseudo-coherent}
伪凝聚态射的平坦基变换伪凝聚。
\end{lemma}

\begin{proof}
这转化为如下代数结果：设 $A \to B$ 是伪凝聚环映射，
$A \to A'$ 平坦。则 $A' \to B \otimes_A A'$ 伪凝聚。
这由更一般的
代数进阶章
引理 \ref{more-algebra-lemma-base-change-relative-pseudo-coherent}.
得到。
\end{proof}

\begin{lemma}
\label{more-morphisms-lemma-composition-pseudo-coherent}
伪凝聚概形态射的复合伪凝聚。
\end{lemma}

\begin{proof}
这转化为如下代数结果：若 $A \to B \to C$ 是可复合的伪凝聚环映射，
则 $A \to C$ 伪凝聚。这可由
代数进阶章
引理 \ref{more-algebra-lemma-pull-relative-pseudo-coherent}
或
代数进阶章
引理 \ref{more-algebra-lemma-composition-relative-pseudo-coherent}.
中的任一个得到。
\end{proof}

\begin{lemma}
\label{more-morphisms-lemma-pseudo-coherent-finite-presentation}
伪凝聚态射局部有限表现。
\end{lemma}

\begin{proof}
直接由定义得到。
\end{proof}

\begin{lemma}
\label{more-morphisms-lemma-flat-finite-presentation-pseudo-coherent}
平坦且局部有限表现的态射伪凝聚。
\end{lemma}

\begin{proof}
这是因为有限表现的平坦环映射伪凝聚（甚至完美），见
代数进阶章
引理 \ref{more-algebra-lemma-flat-finite-presentation-perfect}.
\end{proof}

\begin{lemma}
\label{more-morphisms-lemma-permanence-pseudo-coherent}
设 $f : X \to Y$ 是基概形 $S$ 上伪凝聚的概形之间的态射。
则 $f$ 伪凝聚。
\end{lemma}

\begin{proof}
这转化为如下代数结果：若 $R \to A \to B$ 是可复合的环映射，
$R \to A$、$R \to B$ 伪凝聚，则
% P05-EN-0134: frozen source concludes R->B, already assumed; the lemma requires the relative map A->B.
$R \to B$ 伪凝聚。这由
代数进阶章
引理 \ref{more-algebra-lemma-composition-relative-pseudo-coherent}.
得到。
\end{proof}

\begin{lemma}
\label{more-morphisms-lemma-finite-pseudo-coherent}
设 $f : X \to S$ 是有限概形态射。则 $f$ 伪凝聚，
当且仅当 $f_*\mathcal{O}_X$ 作为 $\mathcal{O}_S$-模伪凝聚。
\end{lemma}

\begin{proof}
把此引理翻译成代数语言，就是如下陈述：若 $R \to A$ 是有限环映射，
则 $R \to A$ 作为环映射伪凝聚（按定义，这意味着 $A$ 作为
$A$-模相对于 $R$ 伪凝聚），当且仅当 $A$ 作为 $R$-模伪凝聚。
这由更一般的
代数进阶章引理
\ref{more-algebra-lemma-finite-extension-pseudo-coherent}.
得到。
\end{proof}

\begin{lemma}
\label{more-morphisms-lemma-Noetherian-pseudo-coherent}
设 $f : X \to S$ 是概形态射。若 $S$ 局部诺特，
则 $f$ 伪凝聚，当且仅当 $f$ 局部有限型。
\end{lemma}

\begin{proof}
这转化为如下代数结果：若 $R \to A$ 是有限型环映射，且 $R$ 诺特，
则 $R \to A$ 伪凝聚，当且仅当 $R \to A$ 有限型。
为此，注意伪凝聚环映射按定义是有限型的。反之，若 $R \to A$
有限型，则可以写成 $A = R[x_1, \ldots, x_n]/I$；由
代数进阶章
引理 \ref{more-algebra-lemma-Noetherian-pseudo-coherent}
可知 $A$ 作为 $R[x_1, \ldots, x_n]$-模伪凝聚，
即 $R \to A$ 是伪凝聚环映射。
\end{proof}

\begin{lemma}
\label{more-morphisms-lemma-descending-property-pseudo-coherent}
性质 $\mathcal{P}(f) =$“$f$ 伪凝聚”在基上是 fpqc 局部的。
\end{lemma}

\begin{proof}
我们使用
下降章引理 \ref{descent-lemma-descending-properties-morphisms}
的判据来证明。由
定义 \ref{more-morphisms-definition-pseudo-coherent}
，伪凝聚性在基上是 Zariski 局部的。由
引理 \ref{more-morphisms-lemma-flat-base-change-pseudo-coherent}
，伪凝聚性在平坦基变换下保持。
下降章引理 \ref{descent-lemma-descending-properties-morphisms}
的最后一个假设 (3) 转化为如下代数陈述：设 $A \to B$ 是忠实平坦环映射，
$C = A[x_1, \ldots, x_n]/I$ 是 $A$-代数。若 $C \otimes_A B$
作为 $B[x_1, \ldots, x_n]$-模伪凝聚，则 $C$ 作为
$A[x_1, \ldots, x_n]$-模伪凝聚。这就是
代数进阶章引理 \ref{more-algebra-lemma-flat-descent-pseudo-coherent}.
\end{proof}

\begin{lemma}
\label{more-morphisms-lemma-quotient-of-flat-finitely-presented}
设 $A \to B$ 是有限表现的平坦环映射，$I \subset B$ 是理想。
则 $A \to B/I$ 伪凝聚，当且仅当 $I$ 作为 $B$-模伪凝聚。
\end{lemma}

\begin{proof}
选取表现 $B = A[x_1, \ldots, x_n]/J$。注意 $B$ 作为
$A[x_1, \ldots, x_n]$-模伪凝聚，因为由
引理 \ref{more-morphisms-lemma-flat-finite-presentation-pseudo-coherent}.
，$A \to B$ 是伪凝聚环映射。还要注意，$A \to B/I$ 伪凝聚，
当且仅当 $B/I$ 作为 $A[x_1, \ldots, x_n]$-模伪凝聚。由
代数进阶章引理 \ref{more-algebra-lemma-finite-push-pseudo-coherent}
，这等价于 $B/I$ 作为 $B$-模伪凝聚。短正合列
$0 \to I \to B \to B/I \to 0$ 表明，$I$ 伪凝聚当且仅当
$B/I$ 伪凝聚（见
代数进阶章
引理 \ref{more-algebra-lemma-two-out-of-three-pseudo-coherent}），
故引理得证。
\end{proof}

\noindent
下面的引理将被更强的
引理 \ref{more-morphisms-lemma-pseudo-coherent-fppf-local-source}.
取代。

\begin{lemma}
\label{more-morphisms-lemma-pseudo-coherent-syntomic-local-source}
性质 $\mathcal{P}(f) =$“$f$ 伪凝聚”在源上是 syntomic 局部的。
\end{lemma}

\begin{proof}
我们使用
下降章引理 \ref{descent-lemma-properties-morphisms-local-source}
的判据来证明。由
引理 \ref{more-morphisms-lemma-flat-finite-presentation-pseudo-coherent} 及
\ref{more-morphisms-lemma-composition-pseudo-coherent}
可知，伪凝聚性在与局部有限表现的平坦态射前复合时保持，
特别是在与 syntomic 态射前复合时保持（见
态射章引理 \ref{morphisms-lemma-syntomic-flat} 及
\ref{morphisms-lemma-syntomic-locally-finite-presentation}）。由
定义 \ref{more-morphisms-definition-pseudo-coherent}
显然可知，伪凝聚性在源和目标上是 Zariski 局部的。因此，按照上述
下降章引理 \ref{descent-lemma-properties-morphisms-local-source}
，只需证明如下陈述：设 $X' \to X \to Y$ 是仿射概形的态射，
$X' \to X$ syntomic，且 $X' \to Y$ 伪凝聚，则 $X \to Y$ 伪凝聚。
首先注意，无论如何，由
下降章引理 \ref{descent-lemma-flat-finitely-presented-permanence-algebra}.
$X \to Y$ 都有限表现。选取闭浸入 $X \to \mathbf{A}^n_Y$。由
代数章引理 \ref{algebra-lemma-lift-syntomic}
，可以找到仿射开覆盖 $X' = \bigcup_{i = 1, \ldots, n} X'_i$，
以及提升态射 $X'_i \to X$ 的 syntomic 态射
$W_i \to \mathbf{A}^n_Y$；也就是说，有纤维积图
$$
\xymatrix{
X'_i \ar[d] \ar[r] & W_i \ar[d] \\
X \ar[r] & \mathbf{A}^n_Y
}
$$
用 $\coprod X'_i$ 替换 $X'$ 并置 $W = \coprod W_i$，得到纤维积图
$$
\xymatrix{
X' \ar[d] \ar[r] & W \ar[d]^h \\
X \ar[r] & \mathbf{A}^n_Y
}
$$
，其中 $W \to \mathbf{A}^n_Y$ 平坦且有限表现，而 $X' \to Y$
仍伪凝聚。由于 $W \to \mathbf{A}^n_Y$ 是开的（见
态射章引理 \ref{morphisms-lemma-fppf-open})
），且 $X' \to X$ 满射，可以找到
$f \in \Gamma(\mathbf{A}^n_Y, \mathcal{O})$，使得
$X \subset D(f) \subset \Im(h)$。写
$Y = \Spec(R)$、$X = \Spec(A)$、$X' = \Spec(A')$、
$W = \Spec(B)$、$A = R[x_1, \ldots, x_n]/I$、$A' = B/IB$。
于是 $R \to A'$ 伪凝聚。图示为
$$
\xymatrix{
A' = B/IB & B \ar[l] \\
A = R[x_1, \ldots, x_n]/I \ar[u] & R[x_1, \ldots, x_n] \ar[l] \ar[u]
}
$$
由
引理 \ref{more-morphisms-lemma-quotient-of-flat-finitely-presented}
可知 $IB$ 作为 $B$-模伪凝聚。由上面对 $f$ 的选择，环映射
$R[x_1, \ldots, x_n]_f \to B_f$ 忠实平坦。这蕴含
$I_f \subset R[x_1, \ldots, x_n]_f$ 伪凝聚，见
代数进阶章引理 \ref{more-algebra-lemma-flat-descent-pseudo-coherent}.
再应用一次
引理 \ref{more-morphisms-lemma-quotient-of-flat-finitely-presented}
，可知 $R \to A$ 伪凝聚。
\end{proof}

\begin{lemma}
\label{more-morphisms-lemma-pseudo-coherent-fppf-local-source}
性质 $\mathcal{P}(f) =$“$f$ 伪凝聚”在源上是 fppf 局部的。
\end{lemma}

\begin{proof}
设 $f : X \to S$ 是概形态射。设 $\{g_i : X_i \to X\}$
是 fppf 覆盖，使得每个复合 $f \circ g_i$ 都伪凝聚。由
引理 \ref{more-morphisms-lemma-dominate-fppf}
，存在
\begin{enumerate}
\item Zariski 开覆盖 $X = \bigcup U_j$；
\item 满射的有限局部自由态射 $W_j \to U_j$；
\item Zariski 开覆盖 $W_j = \bigcup_k W_{j, k}$；
\item 满射的有限局部自由态射 $T_{j, k} \to W_{j, k}$，
\end{enumerate}
使得 fppf 覆盖 $\{h_{j, k} : T_{j, k} \to X\}$ 加细给定覆盖
$\{X_i \to X\}$。以 $\psi_{j, k} : T_{j, k} \to X_{\alpha(j, k)}$
表示见证 $\{T_{j, k} \to X\}$ 加细给定覆盖 $\{X_i \to X\}$ 的态射。
注意 $T_{j, k} \to X$ 平坦且局部有限表现，故由
引理 \ref{more-morphisms-lemma-flat-finite-presentation-pseudo-coherent}.
，$X_i$ 和 $T_{j, k}$ 都在 $X$ 上伪凝聚。因此
$\psi_{j, k} : T_{j, k} \to X_i$ 伪凝聚，见
引理 \ref{more-morphisms-lemma-permanence-pseudo-coherent}.
从而 $T_{j, k} \to S$ 作为 $\psi_{j, k}$ 与
$f \circ g_{\alpha(j, k)}$ 的复合是伪凝聚的，见
引理 \ref{more-morphisms-lemma-composition-pseudo-coherent}.
由此可见，已经把引理归结为 Zariski 开覆盖的情形（这没有问题），
以及覆盖由单个满射的有限局部自由态射给出的情形；后者在下一段处理。

\medskip\noindent
假设 $X' \to X \to S$ 是一列概形态射，其中 $X' \to X$
满射且有限局部自由，并且
% P05-EN-0135: frozen source says X'->Y although this paragraph has base S and no Y.
$X' \to Y$ 伪凝聚。目标是证明 $X \to S$ 伪凝聚。注意，由
下降章引理 \ref{descent-lemma-flat-finitely-presented-permanence}
，态射 $X \to S$ 局部有限表现。显然，问题可归结为 $X'$、$X$、$S$
都仿射，且 $X' \to X$ 为某个秩 $r > 0$ 的自由态射的情形。
相应的代数问题如下：设 $R \to A \to A'$ 是环映射，
$R \to A'$ 伪凝聚，$R \to A$ 有限表现，并且 $A' \cong A^{\oplus r}$
作为 $A$-模。目标是证明 $R \to A$ 伪凝聚。
假设 $R \to A'$ 伪凝聚，意味着 $A'$ 作为 $A'$-模相对于 $R$
伪凝聚。由
代数进阶章
引理 \ref{more-algebra-lemma-finite-extension-pseudo-coherent}
，这蕴含 $A'$ 作为 $A$-模相对于 $R$ 伪凝聚。
由于 $A' \cong A^{\oplus r}$ 作为 $A$-模，故 $A$ 作为 $A$-模
相对于 $R$ 伪凝聚，见
代数进阶章
引理 \ref{more-algebra-lemma-summands-relative-pseudo-coherent}.
按定义，这意味着 $R \to A$ 伪凝聚，证明完成。
\end{proof}








\section{完美态射}
\label{more-morphisms-section-perfect}

\noindent
要理解本节内容，必须稍微了解伪凝聚态射一节的内容。
目前只需知道：环映射 $A \to B$ 完美，当且仅当它伪凝聚，
且 $B$ 作为 $A$-模具有有限 Tor 维数。

\begin{lemma}
\label{more-morphisms-lemma-perfect}
设 $f : X \to S$ 是局部有限型的概形态射。下列条件等价：
\begin{enumerate}
\item 存在仿射开覆盖 $S = \bigcup V_j$，且对每个 $j$
存在仿射开覆盖 $f^{-1}(V_j) = \bigcup U_{ji}$，使得
% P05-EN-0136: frozen source reverses the covering index U_ji to U_ij in the ring-map clause.
$\mathcal{O}_S(V_j) \to \mathcal{O}_X(U_{ij})$ 是完美环映射；
\item 对每一对满足 $f(U) \subset V$ 的仿射开集 $U \subset X$、
$V \subset S$，环映射 $\mathcal{O}_S(V) \to \mathcal{O}_X(U)$ 完美。
\end{enumerate}
\end{lemma}

\begin{proof}
假设 (1)，并设 $U, V$ 如 (2) 所述。由
引理 \ref{more-morphisms-lemma-pseudo-coherent}
可知 $\mathcal{O}_S(V) \to \mathcal{O}_X(U)$ 伪凝聚。
因此，只需证明环映射“具有有限 Tor 维数”这一性质满足
态射章定义 \ref{morphisms-definition-property-local}.
的条件 (1)(a)、(b)、(c)。这些性质由
代数进阶章
引理 \ref{more-algebra-lemma-flat-push-tor-amplitude},
\ref{more-algebra-lemma-flat-base-change-finite-tor-dimension}, 及
\ref{more-algebra-lemma-glue-tor-amplitude}.
得到。略去一些细节。
\end{proof}

\begin{definition}
\label{more-morphisms-definition-perfect}
若概形态射 $f : X \to S$ 满足
引理 \ref{more-morphisms-lemma-perfect}
中的等价条件，则称它是{\it 完美的}。在这种情形下，
也称 $X$ 在 $S$ 上完美。
\end{definition}

\noindent
注意，完美态射特别地伪凝聚，因而局部有限表现。
还要注意，完美态射的基变换一般不完美。

\begin{lemma}
\label{more-morphisms-lemma-flat-base-change-perfect}
完美态射的平坦基变换完美。
\end{lemma}

\begin{proof}
这转化为如下代数结果：设 $A \to B$ 是完美环映射，
$A \to A'$ 平坦。则 $A' \to B \otimes_A A'$ 完美。
关于伪凝聚环映射的相应结果已在
引理 \ref{more-morphisms-lemma-flat-base-change-pseudo-coherent}.
中见过。有限 Tor 维数的相应事实由
代数进阶章
引理 \ref{more-algebra-lemma-flat-base-change-finite-tor-dimension}.
得到。
\end{proof}

\begin{lemma}
\label{more-morphisms-lemma-composition-perfect}
完美概形态射的复合完美。
\end{lemma}

\begin{proof}
这转化为如下代数结果：若 $A \to B \to C$ 是可复合的完美环映射，
则 $A \to C$ 完美。关于伪凝聚性的相应结论已见于
引理 \ref{more-morphisms-lemma-composition-pseudo-coherent}
及其证明。由假设，存在整数 $n$、$m$，使得 $B$ 在 $A$ 上的
Tor 维数 $\leq n$，而 $C$ 在 $B$ 上的 Tor 维数 $\leq m$。
于是对任意 $A$-模 $M$，有
$$
M \otimes_A^{\mathbf{L}} C =
(M \otimes_A^{\mathbf{L}} B) \otimes_B^{\mathbf{L}} C
$$
而
代数进阶章例 \ref{more-algebra-example-tor}
的谱序列表明，当 $p > n + m$ 时
$\text{Tor}^A_p(M, C) = 0$，正是所需结论。
\end{proof}

\begin{lemma}
\label{more-morphisms-lemma-flat-finite-presentation-perfect}
设 $f : X \to S$ 是概形态射。下列条件等价：
\begin{enumerate}
\item $f$ 平坦且完美；
\item $f$ 平坦且局部有限表现。
\end{enumerate}
\end{lemma}

\begin{proof}
蕴含关系 (2) $\Rightarrow$ (1) 就是
代数进阶章
引理 \ref{more-algebra-lemma-flat-finite-presentation-perfect}.
反向蕴含由伪凝聚态射局部有限表现这一事实得到，见
引理 \ref{more-morphisms-lemma-pseudo-coherent-finite-presentation}.
\end{proof}

\begin{lemma}
\label{more-morphisms-lemma-regular-target-perfect}
设 $f : X \to S$ 是概形态射。假设 $S$ 正则且 $f$ 局部有限型。
则 $f$ 完美。
\end{lemma}

\begin{proof}
见
代数进阶章引理 \ref{more-algebra-lemma-regular-perfect-ring-map}.
\end{proof}

\begin{lemma}
\label{more-morphisms-lemma-regular-immersion-perfect}
概形的正则浸入完美。概形的 Koszul-正则浸入完美。
\end{lemma}

\begin{proof}
由于正则浸入是 Koszul-正则浸入，见
除子章引理 \ref{divisors-lemma-regular-quasi-regular-immersion},
只需证明第二个陈述。这转化为如下代数陈述：假设 $I \subset A$
是由 $A$ 的 Koszul-正则列 $f_1, \ldots, f_r$ 生成的理想，
则 $A \to A/I$ 是完美环映射。由于 $A \to A/I$ 满射，
这给出 $A/I$ 由 $A$ 上多项式代数所得的一个表现。
因此，只需说明 $A/I$ 作为 $A$-模伪凝聚且具有有限 Tor 维数。
按 Koszul 列的定义，Koszul 复形 $K(A, f_1, \ldots, f_r)$
是 $A/I$ 的有限自由分解。因此 $A/I$ 是 $A$-模的完美复形，结论得证。
\end{proof}

\begin{lemma}
\label{more-morphisms-lemma-perfect-permanence}
设
$$
\xymatrix{
X \ar[rr]_f \ar[rd] & & Y \ar[ld] \\
& S
}
$$
是概形态射的交换图。假设 $Y \to S$ 光滑且 $X \to S$ 完美。
则 $f : X \to Y$ 完美。
\end{lemma}

\begin{proof}
可以把 $f$ 分解为复合
$$
X \longrightarrow X \times_S Y \longrightarrow Y
$$
，其中第一个态射是映射 $i = (1, f)$，第二个态射是投影。
由于 $Y \to S$ 平坦，见
态射章引理 \ref{morphisms-lemma-smooth-flat},
由 引理 \ref{more-morphisms-lemma-flat-base-change-perfect} 可知
$X \times_S Y \to Y$ 完美。
由于 $Y \to S$ 光滑，$X \times_S Y \to X$ 也光滑，见
态射章引理 \ref{morphisms-lemma-base-change-smooth}.
因此 $i$ 是光滑态射的截面，从而 $i$ 是正则浸入，见
除子章引理 \ref{divisors-lemma-section-smooth-regular-immersion}.
这蕴含 $i$ 完美，见
引理 \ref{more-morphisms-lemma-regular-immersion-perfect}.
因为完美态射的复合完美，故 $f$ 完美，见
引理 \ref{more-morphisms-lemma-composition-perfect}.
\end{proof}

\begin{remark}
\label{more-morphisms-remark-perfect-permanence}
基 $S$ 上完美概形 $X, Y$ 之间的态射未必完美。例如，取
$S = \Spec(k)$、$X = \Spec(k)$，
% P05-EN-0137: frozen source omits the closing parenthesis in this Spec expression.
$Y = \Spec(k[x]/(x^2)$，并取 $X \to Y$ 为唯一的 $S$-态射。
\end{remark}

\begin{lemma}
\label{more-morphisms-lemma-descending-property-perfect}
性质 $\mathcal{P}(f) =$“$f$ 完美”在基上是 fpqc 局部的。
\end{lemma}

\begin{proof}
我们使用
下降章引理 \ref{descent-lemma-descending-properties-morphisms}
的判据来证明。由
定义 \ref{more-morphisms-definition-perfect}
，完美性在基上是 Zariski 局部的。由
引理 \ref{more-morphisms-lemma-flat-base-change-perfect}
，完美性在平坦基变换下保持。
下降章引理 \ref{descent-lemma-descending-properties-morphisms}
的最后一个假设 (3) 转化为如下代数陈述：设 $A \to B$ 是忠实平坦环映射，
$C = A[x_1, \ldots, x_n]/I$ 是 $A$-代数。若 $C \otimes_A B$
作为 $B[x_1, \ldots, x_n]$-模完美，则 $C$ 作为
$A[x_1, \ldots, x_n]$-模完美。这就是
代数进阶章引理 \ref{more-algebra-lemma-flat-descent-perfect}.
\end{proof}

\begin{lemma}
\label{more-morphisms-lemma-check-perfect-stalks}
设 $f : X \to S$ 是伪凝聚概形态射。下列条件等价：
\begin{enumerate}
\item $f$ 完美；
\item $\mathcal{O}_X$ 作为 $f^{-1}\mathcal{O}_S$-模层局部具有有限 Tor 维数；
\item 对所有 $x \in X$，环 $\mathcal{O}_{X, x}$ 作为
$\mathcal{O}_{S, f(x)}$-模具有有限 Tor 维数。
\end{enumerate}
\end{lemma}

\begin{proof}
问题在 $X$ 和 $S$ 上是局部的。因此可以假设
$X = \Spec(B)$、$S = \Spec(A)$，且 $f$ 对应于伪凝聚环映射
$A \to B$。

\medskip\noindent
若 (1) 成立，则 $B$ 作为 $A$-模具有有限 Tor 维数 $d$。
于是对每个满足 $\mathfrak p = A \cap \mathfrak q$ 的素理想
$\mathfrak q \subset B$，$B_\mathfrak q$ 作为 $A_\mathfrak p$-模
具有 Tor 维数 $d$，见
代数进阶章引理 \ref{more-algebra-lemma-tor-amplitude-localization}.
于是由
上同调章引理 \ref{cohomology-lemma-tor-amplitude-stalk}.
，$\mathcal{O}_X$ 作为 $f^{-1}\mathcal{O}_S$-模层具有 Tor 维数 $d$。
故 (1) 蕴含 (2)。

\medskip\noindent
由 上同调章引理 \ref{cohomology-lemma-tor-amplitude-stalk}，
(2) 蕴含 (3)。

\medskip\noindent
假设 (3)。不能使用
代数进阶章引理 \ref{more-algebra-lemma-tor-amplitude-localization}
直接得出结论，因为并未给定 $B_\mathfrak q$ 在 $A_\mathfrak p$
上的 Tor 维数具有与 $\mathfrak q$ 无关的界。选取表现
$A[x_1, \ldots, x_n] \to B$。则 $B$ 作为
$A[x_1, \ldots, x_n]$-模伪凝聚。设
$\mathfrak q \subset A[x_1, \ldots, x_n]$ 是位于
$\mathfrak p \subset A$ 上方的素理想。则 $B_\mathfrak q$ 或者为零，
或者由假设，它作为 $A_\mathfrak p$-模具有有限 Tor 维数。
由于 $A \to A[x_1, \ldots, x_n]$ 的纤维具有有限整体维数，可以把
代数进阶章引理 \ref{more-algebra-lemma-perfect-over-polynomial-ring}
应用于 $A_\mathfrak p \to A[x_1, \ldots, x_n]_\mathfrak q$，
从而看出 $B_\mathfrak q$ 是完美
$A[x_1, \ldots, x_n]_\mathfrak q$-模。因此，由
代数进阶章引理 \ref{more-algebra-lemma-check-perfect-stalks}.
$B$ 是完美 $A[x_1, \ldots, x_n]$-模。故按定义，
$A \to B$ 是完美环映射。
\end{proof}

\begin{lemma}
\label{more-morphisms-lemma-perfect-closed-immersion-perfect-direct-image}
设 $i : Z \to X$ 是完美的概形闭浸入。则 $i_*\mathcal{O}_Z$
是完美 $\mathcal{O}_X$-模，即它是 $D(\mathcal{O}_X)$ 的完美对象。
\end{lemma}

\begin{proof}
这基本上直接由定义得到。具体地，设 $U = \Spec(A)$ 是 $X$
的仿射开集。则对某个理想 $I \subset A$，有
$i^{-1}(U) = \Spec(A/I)$；并且由
代数进阶章引理 \ref{more-algebra-lemma-perfect-ring-map}.
，$A/I$ 具有由有限投射 $A$-模构成的有限分解。因此
$i_*\mathcal{O}_Z|_U$ 可由有限局部自由 $\mathcal{O}_U$-模的
有限长复形表示。这正是所需结论，见
上同调章第 \ref{cohomology-section-perfect} 节.
\end{proof}

\begin{lemma}
\label{more-morphisms-lemma-perfect-proper-perfect-direct-image}
设 $S$ 是诺特概形，$f : X \to S$ 是完美固有概形态射。
设 $E \in D(\mathcal{O}_X)$ 完美。则 $Rf_*E$ 是
$D(\mathcal{O}_S)$ 的完美对象。
\end{lemma}

\begin{proof}
我们断言 概形的导出范畴章引理
\ref{perfect-lemma-perfect-direct-image} 适用。
适用。条件 (1) 和 (2) 是直接的。条件 (3) 在 $X$ 上是局部的。
因此可以假设 $X$ 和 $S$ 仿射，且 $E$ 由严格完美的
$\mathcal{O}_X$-模复形表示。于是只需证明 $\mathcal{O}_X$
作为 $f^{-1}\mathcal{O}_S$-模层具有有限 Tor 维数。
由 引理 \ref{more-morphisms-lemma-check-perfect-stalks}，这等价于完美性。
\end{proof}

\begin{lemma}
\label{more-morphisms-lemma-perfect-fppf-local-source}
性质 $\mathcal{P}(f) =$“$f$ 完美”在源上是 fppf 局部的。
\end{lemma}

\begin{proof}
设 $\{g_i : X_i \to X\}_{i \in I}$ 是概形的 fppf 覆盖，
$f : X \to S$ 是态射，使得每个 $f \circ g_i$ 都完美。由
引理 \ref{more-morphisms-lemma-pseudo-coherent-fppf-local-source}
可知 $f$ 伪凝聚。因此由
引理 \ref{more-morphisms-lemma-check-perfect-stalks}
，只需检验对所有 $x \in X$，$\mathcal{O}_{X, x}$ 作为
$\mathcal{O}_{S, f(x)}$-模具有有限 Tor 维数。
选取映到 $x$ 的 $i \in I$ 和 $x_i \in X_i$。于是
$\mathcal{O}_{X_i, x_i}$ 在 $\mathcal{O}_{S, f(x)}$ 上具有有限 Tor 维数，
且 $\mathcal{O}_{X, x} \to \mathcal{O}_{X_i, x_i}$ 忠实平坦。
所需结论由
代数进阶章引理 \ref{more-algebra-lemma-flat-descent-tor-amplitude}.
得到。
\end{proof}

\begin{lemma}
\label{more-morphisms-lemma-factor-regular-immersion}
设 $i : Z \to Y$ 和 $j : Y \to X$ 是概形浸入。假设
\begin{enumerate}
\item $X$ 局部诺特；
\item $j \circ i$ 是正则浸入；
\item $i$ 完美。
\end{enumerate}
则 $i$ 和 $j$ 都是正则浸入。
\end{lemma}

\begin{proof}
因为 $X$（从而 $Y$）局部诺特，四种正则浸入概念全部一致；
而且可以在局部环层面检验态射是否为正则浸入，见
除子章引理 \ref{divisors-lemma-Noetherian-scheme-regular-ideal}.
因此结论由
除幂代数章引理 \ref{dpa-lemma-regular-sequence}.
得到。
\end{proof}









\section{局部完全交态射}
\label{more-morphisms-section-lci}

\noindent
在
除子章第 \ref{divisors-section-regular-immersions} 节
中，我们定义了四种不同类型的正则浸入：正则、Koszul-正则、
$H_1$-正则及拟正则。本节考虑在 $X$ 上局部具有如下分解的态射
$f : X \to S$：
$$
\xymatrix{
X \ar[rr]_i \ar[rd] & & \mathbf{A}^n_S \ar[ld] \\
& S
}
$$
其中对 $* \in \{\emptyset, Koszul, H_1, quasi\}$，$i$ 是 $*$-正则浸入。
然而，当 $* = \emptyset$，即要求 $i$ 是正则浸入时，
我们不知道如何证明这一条件与分解的选择无关。另一方面，
我们希望局部完全交态射是完美的，而这只有在 $* = Koszul$
或 $* = \emptyset$ 时才成立。因此，我们把在 $X$ 上局部具有上述分解，
且其中 $i$ 是 Koszul-正则浸入的概形态射 $f : X \to S$，
定义为{\it 局部完全交态射}或{\it Koszul 态射}。
为说明这样定义可行，先证明此性质与所选分解无关。

\begin{lemma}
\label{more-morphisms-lemma-koszul-independence-factorization}
设 $S$ 是概形，$U$、$P$、$P'$ 是 $S$-概形，$u \in U$。
设 $i : U \to P$、$i' : U \to P'$ 是 $S$ 上的浸入。
假设 $P$ 和 $P'$ 在 $S$ 上光滑。则下列条件等价：
\begin{enumerate}
\item $i$ 在 $u$ 的一个邻域中是 Koszul-正则浸入；
\item $i'$ 在 $u$ 的一个邻域中是 Koszul-正则浸入。
\end{enumerate}
\end{lemma}

\begin{proof}
假设 $i$ 在 $u$ 的一个邻域中是 Koszul-正则浸入。
考虑态射 $j = (i, i') : U \to P \times_S P' = P''$。
因为 $P'' = P \times_S P' \to P$ 光滑，由
除子章引理 \ref{divisors-lemma-lift-regular-immersion-to-smooth}
可知 $j$ 是 Koszul-正则浸入；进而由
除子章引理 \ref{divisors-lemma-push-regular-immersion-thru-smooth}
可知 $i'$ 是 Koszul-正则浸入。
\end{proof}

\noindent
在给出定义之前，先作如下简单说明。设 $f : X \to S$
是局部有限型的概形态射，$x \in X$。则存在开邻域
$U \subset X$，以及 $f|_U$ 的一个分解：先作浸入
$i : U \to \mathbf{A}^n_S$，再作光滑投影
$\mathbf{A}^n_S \to S$。图示为
$$
\xymatrix{
X \ar[rd] & U \ar[l] \ar[d] \ar[r]_-i & \mathbf{A}^n_S = P \ar[ld]^\pi \\
& S
}
$$
事实上，对 $x$ 在 $X$ 中的任意仿射开邻域 $U$ 都可以这样做，见
态射章引理 \ref{morphisms-lemma-quasi-affine-finite-type-over-S}.

\begin{definition}
\label{more-morphisms-definition-lci}
设 $f : X \to S$ 是概形态射。
\begin{enumerate}
\item 设 $x \in X$。若 $f$ 在 $x$ 处为有限型，并且存在开邻域
及 $f|_U$ 的分解 $\pi \circ i$，其中 $i : U \to P$
是 Koszul-正则浸入，$\pi : P \to S$ 光滑，
则称 $f${\it 在 $x$ 处为 Koszul}。
\item 若 $f$ 在每一点处都为 Koszul，则称 $f$ 是
{\it Koszul 态射}，或称 $f$ 是{\it 局部完全交态射}。
\end{enumerate}
\end{definition}

\noindent
上面已经看到，分解 $f|_U = \pi \circ i$ 的选择无关紧要。
也就是说，给定 $f|_U$ 的分解——先作浸入 $i$，再作光滑态射 $\pi$——
$i$ 在 $x$ 的一个邻域中是否 Koszul-正则，是 $f$ 在 $x$ 处的内蕴性质。
为便于引用，把这一点明确记录为引理。

\begin{lemma}
\label{more-morphisms-lemma-lci}
设 $f : X \to S$ 是局部完全交态射，$P$ 是在 $S$ 上光滑的概形。
设 $U \subset X$ 是开子概形，$i : U \to P$ 是 $S$ 上的概形浸入。
则 $i$ 是 Koszul-正则浸入。
\end{lemma}

\begin{proof}
这就是局部完全交态射的定义性质。见上面的讨论。
\end{proof}

\noindent
下面汇集所有 Koszul 态射共有的一些性质。

\begin{lemma}
\label{more-morphisms-lemma-lci-properties}
设 $f : X \to S$ 是局部完全交态射。则
\begin{enumerate}
\item $f$ 局部有限表现；
\item $f$ 伪凝聚；
\item $f$ 完美。
\end{enumerate}
\end{lemma}

\begin{proof}
完美态射伪凝聚（因为完美环映射伪凝聚），而伪凝聚态射局部有限表现
（因为伪凝聚环映射有限表现），所以只需证明最后一个陈述。
完美性是局部性质，故可假设 $f$ 分解为 $\pi \circ i$，
其中 $\pi$ 光滑，而 $i$ 是 Koszul-正则浸入。
Koszul-正则浸入完美，见
引理 \ref{more-morphisms-lemma-regular-immersion-perfect}.
光滑态射平坦且局部有限表现，因而完美，见
引理 \ref{more-morphisms-lemma-flat-finite-presentation-perfect}.
最后，完美态射的复合完美，见
引理 \ref{more-morphisms-lemma-composition-perfect}.
\end{proof}

\begin{lemma}
\label{more-morphisms-lemma-affine-lci}
设 $f : X = \Spec(B) \to S = \Spec(A)$ 是仿射概形态射。
则 $f$ 是局部完全交态射，当且仅当 $A \to B$ 是局部完全交同态，见
代数进阶章定义
\ref{more-algebra-definition-local-complete-intersection}.
\end{lemma}

\begin{proof}
直接由定义得到。
\end{proof}

\noindent
注意，Koszul 态射的基变换一般不是 Koszul 态射。

\begin{lemma}
\label{more-morphisms-lemma-flat-base-change-lci}
局部完全交态射的平坦基变换仍是局部完全交态射。
\end{lemma}

\begin{proof}
略。提示：光滑态射的基变换光滑，而 Koszul-正则浸入的平坦基变换
仍是 Koszul-正则浸入，故结论成立，见
除子章引理 \ref{divisors-lemma-regular-immersion-noetherian}.
\end{proof}

\begin{lemma}
\label{more-morphisms-lemma-composition-lci}
局部完全交态射的复合仍是局部完全交态射。
\end{lemma}

\begin{proof}
设 $g : Y \to S$ 和 $f : X \to Y$ 是局部完全交态射。
设 $x \in X$，并置 $y = f(x)$。选取 $y$ 的开邻域 $V \subset Y$，
以及分解 $g|_V = \pi \circ i$，其中 $i : V \to P$
是 Koszul-正则浸入，$\pi : P \to S$ 光滑。
接着选取 $x \in X$ 的开邻域 $U$，以及分解
$f|_U = \pi' \circ i'$，其中 $i' : U \to P'$ 是 Koszul-正则浸入，
$\pi' : P' \to Y$ 光滑。事实上，可以假设
$P' = \mathbf{A}^n_V$，见 定义 \ref{more-morphisms-definition-lci}
前后的讨论。图示为：
$$
\xymatrix{
X \ar[d] & U \ar[l] \ar[r]_-{i'} & P' = \mathbf{A}^n_V \ar[d] \\
Y \ar[d] &  & V \ar[ll] \ar[r]_i & P \ar[d] \\
S & & & S \ar[lll]
}
$$
置 $P'' = \mathbf{A}^n_P$。则 $U \to P' \to P''$
是 Koszul-正则浸入的复合，即 $i'$ 与 $i$ 沿 $P'' \to P$
的平坦基变换之复合，因而是 Koszul-正则浸入，见
除子章
引理 \ref{divisors-lemma-regular-immersion-noetherian}
及
除子章引理 \ref{divisors-lemma-composition-regular-immersion}.
此外，$P'' \to P \to S$ 作为光滑态射的复合是光滑的，见
态射章引理 \ref{morphisms-lemma-composition-smooth}.
因此，$X \to S$ 如所需在 $x$ 处为 Koszul。
\end{proof}

\begin{lemma}
\label{more-morphisms-lemma-flat-lci}
\begin{slogan}
态射平坦且为局部完全交，当且仅当它是 syntomic 的。
\end{slogan}
设 $f : X \to S$ 是概形态射。下列条件等价：
\begin{enumerate}
\item $f$ 平坦且为局部完全交态射；
\item $f$ syntomic。
\end{enumerate}
\end{lemma}

\begin{proof}
在仿射局部工作时，这就是
代数进阶章引理 \ref{more-algebra-lemma-syntomic-lci}.
下面也给出一个更几何的证明。

\medskip\noindent
假设 (2)。由
态射章引理 \ref{morphisms-lemma-syntomic-locally-standard-syntomic}
，对 $X$ 的每个点 $x$，存在 $x$ 的仿射开邻域 $U$ 和 $f(x)$
的仿射开邻域 $V$，使得 $f|_U : U \to V$ 是标准 syntomic 态射。
这意味着 $U = \Spec(R[x_1, \ldots, x_n]/(f_1, \ldots, f_c))
\to V = \Spec(R)$，其中 $R[x_1, \ldots, x_n]/(f_1, \ldots, f_c)$
是 $R$ 上的相对整体完全交。由
代数章
引理 \ref{algebra-lemma-relative-global-complete-intersection-conormal}
，对每个素理想 $\mathfrak q \supset (f_1, \ldots, f_c)$，列
$f_1, \ldots, f_c$ 在局部环
$R[x_1, \ldots, x_n]_{\mathfrak q}$ 中都是正则列。
考虑 Koszul 复形
$K_\bullet = K_\bullet(R[x_1, \ldots, x_n], f_1, \ldots, f_c)$，
其同调群为 $H_i = H_i(K_\bullet)$。由
代数进阶章引理 \ref{more-algebra-lemma-regular-koszul-regular}
可知，对上述每个 $\mathfrak q$，当 $i > 0$ 时有
$(H_i)_{\mathfrak q} = 0$。另一方面，由
代数进阶章引理 \ref{more-algebra-lemma-homotopy-koszul}
可知 $(f_1, \ldots, f_c)$ 零化 $H_i$。因此，当 $i > 0$ 时
$H_i = 0$，并且 $f_1, \ldots, f_c$ 是 Koszul-正则列。
这证明 $U \to V$ 如所需分解为 Koszul-正则浸入
$U \to \mathbf{A}^n_V$ 后接一个光滑态射。

\medskip\noindent
假设 (1)。则 $f$ 平坦且局部有限表现
（引理 \ref{more-morphisms-lemma-lci-properties}）。因此，根据
态射章引理 \ref{morphisms-lemma-syntomic-locally-standard-syntomic}
，只需证明局部环 $\mathcal{O}_{X_s, x}$ 是局部完全交环。
在 $X$ 上局部选取分解 $f = \pi \circ i$，其中
$i : X \to P$ 是 Koszul-正则浸入，$\pi : P \to S$ 光滑。
注意 $X \to P$ 是 $S$ 上的相对拟正则浸入，见
除子章定义 \ref{divisors-definition-relative-H1-regular-immersion}.
因此，由
除子章
引理 \ref{divisors-lemma-relative-regular-immersion-flat-in-neighbourhood}
可知 $X \to P$ 是正则浸入，且这一性质在任意基变换后仍成立。
因此每个纤维都是正则浸入，从而 $X$ 的所有纤维的所有局部环
都是局部完全交。
\end{proof}

\begin{lemma}
\label{more-morphisms-lemma-regular-immersion-lci}
概形的正则浸入是局部完全交态射。
概形的 Koszul-正则浸入是局部完全交态射。
\end{lemma}

\begin{proof}
由于正则浸入是 Koszul-正则浸入，见
除子章引理 \ref{divisors-lemma-regular-quasi-regular-immersion},
只需证明第二个陈述，而它直接由定义得到。
\end{proof}

\begin{lemma}
\label{more-morphisms-lemma-lci-permanence}
设
$$
\xymatrix{
X \ar[rr]_f \ar[rd] & & Y \ar[ld] \\
& S
}
$$
是概形态射的交换图。假设 $Y \to S$ 光滑，且 $X \to S$
是局部完全交态射。则 $f : X \to Y$ 是局部完全交态射。
\end{lemma}

\begin{proof}
直接由定义得到。
\end{proof}

\begin{lemma}
\label{more-morphisms-lemma-morphism-regular-schemes-is-lci}
设 $f : X \to Y$ 是概形态射。若 $f$ 局部有限型，
且 $X$ 和 $Y$ 正则，则 $f$ 是局部完全交态射。
\end{lemma}

\begin{proof}
可以假设存在分解 $X \to \mathbf{A}^n_Y \to Y$，其中第一个箭头是浸入。
由于 $Y$ 正则，由
代数章引理 \ref{algebra-lemma-regular-goes-up}.
可知 $\mathbf{A}^n_Y$ 也正则。因此由
除子章引理 \ref{divisors-lemma-immersion-regular-regular-immersion}.
$X \to \mathbf{A}^n_Y$ 是正则浸入。
\end{proof}

\noindent
下面的引理具有不同的性质。

\begin{lemma}
\label{more-morphisms-lemma-lci-avramov}
设
$$
\xymatrix{
X \ar[rr]_f \ar[rd] & & Y \ar[ld] \\
& S
}
$$
是概形态射的交换图。假设
\begin{enumerate}
\item $S$ 局部诺特；
\item $Y \to S$ 局部有限型；
\item $f : X \to Y$ 完美；
\item $X \to S$ 是局部完全交态射。
\end{enumerate}
则 $X \to Y$ 是局部完全交态射，并且对所有 $x \in X$，
$Y \to S$ 在 $f(x)$ 处为 Koszul。
\end{lemma}

\begin{proof}
在此证明中，所有概形都局部诺特，所有环都诺特。我们将不再另行说明地使用：
在这种情形下，正则列与 Koszul-正则列一致，见 代数进阶章引理
\ref{more-algebra-lemma-noetherian-finite-all-equivalent}.
此外，理想（相应地，理想层）的正则性可以在局部环（相应地，茎）上检验，见
代数章引理 \ref{algebra-lemma-regular-sequence-in-neighbourhood}
（相应地，除子章引理 \ref{divisors-lemma-Noetherian-scheme-regular-ideal}）。

\medskip\noindent
问题是局部的。因此可以假设 $S$、$X$、$Y$ 都仿射。
在这种情形下，可以选取交换图
$$
\xymatrix{
\mathbf{A}^{n + m}_S \ar[d] & X \ar[l] \ar[d] \\
\mathbf{A}^n_S \ar[d] & Y \ar[l] \ar[ld] \\
S
}
$$
，其中水平箭头是闭浸入。设 $x \in X$ 是一点，并考虑相应的局部环交换图
$$
\xymatrix{
J \ar[r] &
\mathcal{O}_{\mathbf{A}^{n + m}_S, x} \ar[r] &
\mathcal{O}_{X, x} \\
I \ar[r] \ar[u] &
\mathcal{O}_{\mathbf{A}^n_S, f(x)} \ar[r] \ar[u] &
\mathcal{O}_{Y, f(x)} \ar[u]
}
$$
，其中 $J$ 和 $I$ 是水平箭头的核。由于 $X \to S$ 是局部完全交态射，
理想 $J$ 由正则列生成。由于 $X \to Y$ 完美，环
$\mathcal{O}_{X, x}$ 在 $\mathcal{O}_{Y, f(x)}$ 上具有有限 Tor 维数。
因此可以应用
除幂代数章引理 \ref{dpa-lemma-perfect-map-ci}
断定 $I$ 和 $J/I$ 都由正则列生成。由开头的说明，证明完毕。
\end{proof}

\begin{lemma}
\label{more-morphisms-lemma-lci-to-regular}
设
$$
\xymatrix{
X \ar[rr]_f \ar[rd] & & Y \ar[ld] \\
& S
}
$$
是概形态射的交换图。假设 $S$ 局部诺特，$Y \to S$ 局部有限型，
$Y$ 正则，且 $X \to S$ 是局部完全交态射。
则 $f : X \to Y$ 是局部完全交态射，并且对所有 $x \in X$，
$Y \to S$ 在 $f(x)$ 处为 Koszul。
\end{lemma}

\begin{proof}
这是 引理 \ref{more-morphisms-lemma-lci-avramov} 的特殊情形；这一点由
引理 \ref{more-morphisms-lemma-regular-target-perfect}
（以及 态射章引理 \ref{morphisms-lemma-permanence-finite-type}）得到。
\end{proof}

\begin{lemma}
\label{more-morphisms-lemma-perfect-conormal-free-lci}
设 $i : X \to Y$ 是浸入。若
\begin{enumerate}
\item $i$ 完美；
\item $Y$ 局部诺特；
\item 余法层 $\mathcal{C}_{X/Y}$ 有限局部自由，
\end{enumerate}
则 $i$ 是正则浸入。
\end{lemma}

\begin{proof}
翻译成代数语言，这就是
除幂代数章命题 \ref{dpa-proposition-regular-ideal}.
\end{proof}

\begin{lemma}
\label{more-morphisms-lemma-lci-NL}
% P05-EN-0138: frozen source calls the scheme morphism a homomorphism.
设 $f : X \to Y$ 是局部完全交同态。则朴素余切复形 $\NL_{X/Y}$
是 $D(\mathcal{O}_X)$ 中 Tor 振幅位于 $[-1, 0]$ 的完美对象。
\end{lemma}

\begin{proof}
翻译成代数语言，这就是
代数进阶章引理 \ref{more-algebra-lemma-lci-NL}。
在这一翻译中使用 引理 \ref{more-morphisms-lemma-affine-lci} 和
\ref{more-morphisms-lemma-NL-affine}，以及
概形的导出范畴章引理
\ref{perfect-lemma-affine-compare-bounded},
\ref{perfect-lemma-tor-dimension-affine} 及
\ref{perfect-lemma-perfect-affine}.
\end{proof}

\begin{lemma}
\label{more-morphisms-lemma-perfect-NL-lci}
设 $f : X \to Y$ 是局部诺特概形之间的完美态射。下列条件等价：
\begin{enumerate}
\item $f$ 是局部完全交态射；
\item $\NL_{X/Y}$ 的 Tor 振幅位于 $[-1, 0]$；
\item $\NL_{X/Y}$ 完美且 Tor 振幅位于 $[-1, 0]$。
\end{enumerate}
\end{lemma}

\begin{proof}
翻译成代数语言，这就是
除幂代数章引理 \ref{dpa-lemma-perfect-NL-lci}。
在这一翻译中使用 引理 \ref{more-morphisms-lemma-affine-lci} 和
\ref{more-morphisms-lemma-NL-affine}，以及
概形的导出范畴章引理
\ref{perfect-lemma-affine-compare-bounded},
\ref{perfect-lemma-tor-dimension-affine} 及
\ref{perfect-lemma-perfect-affine}.
\end{proof}

\begin{lemma}
\label{more-morphisms-lemma-flat-fp-NL-lci}
设 $f : X \to Y$ 是有限表现的平坦态射。下列条件等价：
\begin{enumerate}
\item $f$ 是局部完全交态射；
\item $f$ syntomic；
\item $\NL_{X/Y}$ 的 Tor 振幅位于 $[-1, 0]$；
\item $\NL_{X/Y}$ 完美且 Tor 振幅位于 $[-1, 0]$。
\end{enumerate}
\end{lemma}

\begin{proof}
翻译成代数语言，这就是
除幂代数章引理 \ref{dpa-lemma-flat-fp-NL-lci}。
在这一翻译中使用 引理 \ref{more-morphisms-lemma-affine-lci} 和
\ref{more-morphisms-lemma-NL-affine}，以及
概形的导出范畴章引理
\ref{perfect-lemma-affine-compare-bounded},
\ref{perfect-lemma-tor-dimension-affine} 及
\ref{perfect-lemma-perfect-affine}.
\end{proof}

\noindent
下面的引理把光滑态射刻画为对角态射完美的平坦态射。

\begin{lemma}
\label{more-morphisms-lemma-smooth-diagonal-perfect}
设 $f : X \to Y$ 是局部诺特概形之间的有限型态射。
以 $\Delta : X \to X \times_Y X$ 表示对角态射。下列条件等价：
\begin{enumerate}
\item $f$ 光滑；
\item $f$ 平坦，且 $\Delta : X \to X \times_Y X$ 是正则浸入；
\item $f$ 平坦，且 $\Delta : X \to X \times_Y X$ 是局部完全交态射；
\item $f$ 平坦，且 $\Delta : X \to X \times_Y X$ 完美。
\end{enumerate}
\end{lemma}

\begin{proof}
假设 (1)。由 态射章引理 \ref{morphisms-lemma-smooth-flat}，
$f$ 平坦。由 态射章引理 \ref{morphisms-lemma-base-change-smooth}，
投影 $X \times_Y X \to X$ 光滑。因此对角态射是光滑态射的截面，
从而是正则浸入，见
除子章引理 \ref{divisors-lemma-section-smooth-regular-immersion}.
故 (1) $\Rightarrow$ (2)。蕴含关系 (2) $\Rightarrow$ (3) 就是
引理 \ref{more-morphisms-lemma-regular-immersion-lci}.
蕴含关系 (3) $\Rightarrow$ (4) 就是
引理 \ref{more-morphisms-lemma-lci-properties}.
较有趣的蕴含关系 (4) $\Rightarrow$ (1) 直接由
除幂代数章引理 \ref{dpa-lemma-perfect-diagonal} 得到。
\end{proof}

\begin{lemma}
\label{more-morphisms-lemma-descending-property-lci}
性质 $\mathcal{P}(f) =$“$f$ 是局部完全交态射”在基上是 fpqc 局部的。
\end{lemma}

\begin{proof}
设 $f : X \to S$ 是概形态射，$\{S_i \to S\}$ 是 $S$ 的 fpqc 覆盖。
假设 $f$ 的每个基变换 $f_i : X_i \to S_i$ 都是局部完全交态射。
注意，这特别地蕴含 $f$ 局部有限型，见
引理 \ref{more-morphisms-lemma-lci-properties}
和
下降章引理 \ref{descent-lemma-descending-property-locally-finite-type}.
设 $x \in X$。选取 $x$ 的开邻域 $U$，以及 $S$ 上的浸入
$j : U \to \mathbf{A}^n_S$（见 定义 \ref{more-morphisms-definition-lci}
之前的讨论）。必须证明 $j$ 是 Koszul-正则浸入。
由于 $f_i$ 是局部完全交态射，基变换
$j_i : U \times_S S_i \to \mathbf{A}^n_{S_i}$
是 Koszul-正则浸入，见
引理 \ref{more-morphisms-lemma-lci}.
因为 $\{\mathbf{A}^n_{S_i} \to \mathbf{A}^n_S\}$ 是 fpqc 覆盖，由
下降章引理 \ref{descent-lemma-descending-property-regular-immersion}
可知 $j$ 如所需是 Koszul-正则浸入。
\end{proof}

\begin{lemma}
\label{more-morphisms-lemma-lci-syntomic-local-source}
性质 $\mathcal{P}(f) =$“$f$ 是局部完全交态射”在源上是 syntomic 局部的。
\end{lemma}

\begin{proof}
我们使用
下降章引理 \ref{descent-lemma-properties-morphisms-local-source}
的判据来证明。由
引理 \ref{more-morphisms-lemma-flat-lci} 及
\ref{more-morphisms-lemma-composition-lci}
可知，局部完全交性质在与 syntomic 态射前复合时保持。由
定义 \ref{more-morphisms-definition-lci}
显然可知，局部完全交性质在源和目标上是 Zariski 局部的。
因此，按照上述
下降章引理 \ref{descent-lemma-properties-morphisms-local-source}
，只需证明如下陈述：设 $X' \to X \to Y$ 是仿射概形的态射，
$X' \to X$ syntomic，且 $X' \to Y$ 是局部完全交态射，
则 $X \to Y$ 是局部完全交态射。首先注意，无论如何，由
下降章引理 \ref{descent-lemma-flat-finitely-presented-permanence-algebra}.
$X \to Y$ 都有限表现。选取闭浸入 $X \to \mathbf{A}^n_Y$。由
代数章引理 \ref{algebra-lemma-lift-syntomic}
，可以找到仿射开覆盖 $X' = \bigcup_{i = 1, \ldots, n} X'_i$，
以及提升态射 $X'_i \to X$ 的 syntomic 态射
$W_i \to \mathbf{A}^n_Y$；也就是说，有纤维积图
$$
\xymatrix{
X'_i \ar[d] \ar[r] & W_i \ar[d] \\
X \ar[r] & \mathbf{A}^n_Y
}
$$
用 $\coprod X'_i$ 替换 $X'$ 并置 $W = \coprod W_i$，
得到仿射概形的纤维积图
$$
\xymatrix{
X' \ar[d] \ar[r] & W \ar[d]^h \\
X \ar[r] & \mathbf{A}^n_Y
}
$$
，其中 $h : W \to \mathbf{A}^n_Y$ syntomic，而 $X' \to Y$
仍是局部完全交态射。由于 $W \to \mathbf{A}^n_Y$ 是开的（见
态射章引理 \ref{morphisms-lemma-fppf-open})
），且 $X' \to X$ 满射，可见 $X$ 包含于 $W \to \mathbf{A}^n_Y$ 的像中。
选取 $\mathbf{A}^n_Y$ 上的闭浸入
$W \to \mathbf{A}^{n + m}_Y$。现在图形如下：
$$
\xymatrix{
X' \ar[d] \ar[r] & W \ar[d]^h \ar[r] & \mathbf{A}^{n + m}_Y \ar[ld] \\
X \ar[r] & \mathbf{A}^n_Y
}
$$
因为 $h$ syntomic，因而是局部完全交态射（见上文），所以态射
$W \to \mathbf{A}^{n + m}_Y$ 是 Koszul-正则浸入。
因为 $X' \to Y$ 是局部完全交态射，所以态射
$X' \to \mathbf{A}^{n + m}_Y$ 是 Koszul-正则浸入。由
除子章引理 \ref{divisors-lemma-permanence-regular-immersion}
可知 $X' \to W$ 是 Koszul-正则浸入。又因为 Koszul-正则浸入
在目标上是 fpqc 局部的（见
下降章引理 \ref{descent-lemma-descending-property-regular-immersion})
），所以 $X \to \mathbf{A}^n_Y$ 是 Koszul-正则浸入，
这正是必须证明的。
\end{proof}

\begin{lemma}
\label{more-morphisms-lemma-base-change-lci-fibres}
设 $S$ 是概形，$f : X \to Y$ 是 $S$ 上的概形态射。
假设 $X$ 和 $Y$ 在 $S$ 上都平坦且局部有限表现。则集合
$$
\{x \in X \mid f\text{ Koszul at }x\}.
$$
在 $X$ 中是开的，且其构造与任意基变换 $S' \to S$ 可交换。
\end{lemma}

\begin{proof}
该集合按定义是开的（见 定义 \ref{more-morphisms-definition-lci}）。
设 $S' \to S$ 是概形态射。置 $X' = S' \times_S X$、
$Y' = S' \times_S Y$，并以 $f' : X' \to Y'$ 表示 $f$ 的基变换。
设 $x' \in X'$ 是一点，使得 $f'$ 在 $x'$ 处为 Koszul。
以 $s' \in S'$、$x \in X$、$y' \in Y'$、$y \in Y$、$s \in S$
表示 $x'$ 的像。注意 $f$ 局部有限表现，见
态射章引理 \ref{morphisms-lemma-finite-presentation-permanence}.
因此可以选取 $x$ 的仿射邻域 $U \subset X$ 及浸入
$i : U \to \mathbf{A}^n_Y$。以 $U' = S' \times_S U$ 和
$i' : U' \to \mathbf{A}^n_{Y'}$ 表示 $i$ 的基变换。
$f'$ 在 $x'$ 处为 Koszul 这一假设蕴含 $i'$ 在 $x'$ 的一个邻域中
是 Koszul-正则浸入，见
引理 \ref{more-morphisms-lemma-lci}.
概形 $X'$ 作为 $X$ 的基变换，在 $S'$ 上平坦且局部有限表现（见
态射章引理 \ref{morphisms-lemma-base-change-flat} 及
\ref{morphisms-lemma-base-change-finite-presentation}）。
因此 $i'$ 在 $x'$ 的一个邻域中是 $S'$ 上的相对
$H_1$-正则浸入（见
除子章定义 \ref{divisors-definition-relative-H1-regular-immersion}).
所以基变换 $i'_{s'} : U'_{s'} \to \mathbf{A}^n_{Y'_{s'}}$
在 $x'$ 的一个开邻域中是 $H_1$-正则浸入，见
除子章引理 \ref{divisors-lemma-relative-regular-immersion}
及 除子章定义 \ref{divisors-definition-relative-H1-regular-immersion}
之后的讨论。由于 $s' = \Spec(\kappa(s')) \to \Spec(\kappa(s)) = s$
是满射、平坦且普遍开的态射（见
态射章引理 \ref{morphisms-lemma-scheme-over-field-universally-open})
），可知基变换 $i_s : U_s \to \mathbf{A}^n_{Y_s}$ 在 $x$
的一个邻域中是 $H_1$-正则浸入，见
下降章引理 \ref{descent-lemma-descending-property-regular-immersion}.
最后，注意 $\mathbf{A}^n_Y$ 在 $S$ 上平坦且局部有限表现，故
除子章引理 \ref{divisors-lemma-fibre-quasi-regular}
蕴含 $i$ 在 $x$ 的一个邻域中如所需是（Koszul-）正则浸入。
\end{proof}

\begin{lemma}
\label{more-morphisms-lemma-unramified-lci}
设 $f : X \to Y$ 是局部完全交概形态射。则 $f$ 非分歧，
当且仅当 $f$ 形式非分歧；在这种情形下，余法层
$\mathcal{C}_{X/Y}$ 在 $X$ 上有限局部自由。
\end{lemma}

\begin{proof}
第一个断言直接由
引理 \ref{more-morphisms-lemma-unramified-formally-unramified}
及局部完全交态射局部有限型这一事实得到。为计算 $f$ 的余法层，
在 $X$ 上局部选取 $f$ 的分解 $f = p \circ i$，其中
$i : X \to V$ 是 Koszul-正则浸入，而 $V \to Y$ 光滑。由
引理 \ref{more-morphisms-lemma-two-unramified-morphisms-formally-smooth}
可知 $\mathcal{C}_{X/Y}$ 是 $\mathcal{C}_{X/V}$ 的局部直和项。
由于 $i$ 是 Koszul-正则（因而拟正则）浸入，后者有限局部自由，见
除子章引理 \ref{divisors-lemma-quasi-regular-immersion}.
\end{proof}

\begin{lemma}
\label{more-morphisms-lemma-transitivity-conormal-lci}
设 $Z \to Y \to X$ 是形式非分歧的概形态射。
假设 $Z \to Y$ 是局部完全交态射。则
$$
0 \to i^*\mathcal{C}_{Y/X} \to
\mathcal{C}_{Z/X} \to
\mathcal{C}_{Z/Y} \to 0
$$
这一来自
引理 \ref{more-morphisms-lemma-transitivity-conormal}
的正合列是短正合列。
\end{lemma}

\begin{proof}
问题在 $Z$ 上是局部的，故可假设态射 $Z \to Y$ 存在分解
$Z \to \mathbf{A}^n_Y \to Y$。于是得到交换图
$$
\xymatrix{
Z \ar[r]_{i'} \ar@{=}[d] &
\mathbf{A}^n_Y \ar[r] \ar[d] &
\mathbf{A}^n_X \ar[d] \\
Z \ar[r]^i & Y \ar[r] & X
}
$$
由于 $Z \to Y$ 是局部完全交态射，$Z \to \mathbf{A}^n_Y$
是 Koszul-正则浸入。因此由
除子章引理 \ref{divisors-lemma-transitivity-conormal-quasi-regular}
，正合列
$$
0 \to (i')^*\mathcal{C}_{\mathbf{A}^n_Y/\mathbf{A}^n_X} \to
\mathcal{C}_{Z/\mathbf{A}^n_X} \to
\mathcal{C}_{Z/\mathbf{A}^n_Y} \to 0
$$
是局部分裂的。注意，有
$i^*\mathcal{C}_{Y/X} = (i')^*\mathcal{C}_{\mathbf{A}^n_Y/\mathbf{A}^n_X}$
；这是由
引理 \ref{more-morphisms-lemma-universal-thickening-fibre-product-flat}
得到的。还要注意图
$$
\xymatrix{
(i')^*\mathcal{C}_{\mathbf{A}^n_Y/\mathbf{A}^n_X} \ar[r] &
\mathcal{C}_{Z/\mathbf{A}^n_X} \\
i^*\mathcal{C}_{Y/X} \ar[u]^{\cong} \ar[r] & \mathcal{C}_{Z/X} \ar[u]
}
$$
交换。因此下方的水平箭头是局部分裂的单射。引理得证。
\end{proof}













\section{微分层与余法层的正合列}
\label{more-morphisms-section-exact}

\noindent
本节汇集余法层与微分层的正合列的一些结果。在某种意义下，
它们都是与一对可复合概形态射相伴的余切复形三角的具体实现。

\medskip\noindent
设 $g : Z \to Y$ 和 $f : Y \to X$ 是概形态射。
\begin{enumerate}
\item 有典范正合列
$$
g^*\Omega_{Y/X} \to \Omega_{Z/X} \to \Omega_{Z/Y} \to 0,
$$
，见
态射章引理 \ref{morphisms-lemma-triangle-differentials}.
若 $g : Z \to Y$ 光滑，或更一般地形式光滑，则此列是短正合列，见
态射章引理 \ref{morphisms-lemma-triangle-differentials-smooth}
或 引理 \ref{more-morphisms-lemma-triangle-differentials-formally-smooth}。
\item 若 $g$ 是浸入，或更一般地形式非分歧，则有典范正合列
$$
\mathcal{C}_{Z/Y} \to g^*\Omega_{Y/X} \to \Omega_{Z/X} \to 0,
$$
，见 态射章引理 \ref{morphisms-lemma-differentials-relative-immersion}
或 引理 \ref{more-morphisms-lemma-universally-unramified-differentials-sequence}。
若 $f \circ g : Z \to X$ 光滑，或更一般地形式光滑，
则此列是短正合列，见
态射章引理 \ref{morphisms-lemma-differentials-relative-immersion-smooth}
或 引理 \ref{more-morphisms-lemma-differentials-formally-unramified-formally-smooth}。
\item 若 $g$ 与 $f \circ g$ 是浸入，或更一般地形式非分歧，
则有典范正合列
$$
\mathcal{C}_{Z/X} \to \mathcal{C}_{Z/Y} \to g^*\Omega_{Y/X} \to 0,
$$
，见 态射章引理 \ref{morphisms-lemma-two-immersions} 或
引理 \ref{more-morphisms-lemma-two-unramified-morphisms}。若 $f : Y \to X$ 光滑，
或更一般地形式光滑，则此列是短正合列，见
态射章引理 \ref{morphisms-lemma-two-immersions-smooth} 或
引理 \ref{more-morphisms-lemma-two-unramified-morphisms-formally-smooth}。
\item 若 $g$ 和 $f$ 是浸入，或更一般地形式非分歧，则有典范正合列
$$
g^*\mathcal{C}_{Y/X} \to \mathcal{C}_{Z/X} \to \mathcal{C}_{Z/Y} \to 0.
$$
见 态射章引理 \ref{morphisms-lemma-transitivity-conormal} 或
引理 \ref{more-morphisms-lemma-transitivity-conormal}。
若 $g : Z \to Y$ 是正则浸入\footnote{$g$ 是 $H_1$-正则浸入即已足够。
注意，作为局部完全交态射的浸入是 Koszul-正则的。}，
或更一般地是局部完全交态射，则此列是短正合列，见
除子章引理 \ref{divisors-lemma-transitivity-conormal-quasi-regular}
或 引理 \ref{more-morphisms-lemma-transitivity-conormal-lci}。
\end{enumerate}










\section{弱 \'etale 态射}
\label{more-morphisms-section-weakly-etale}

\noindent
若 $A \to B$ 和 $B \otimes_A B \to B$ 都平坦，
则称环同态 $A \to B$ 弱 \'etale，见
代数进阶章定义 \ref{more-algebra-definition-weakly-etale}.
概形态射的对应概念如下。

\begin{definition}
\label{more-morphisms-definition-weakly-etale}
若 $X \to Y$ 与对角态射 $X \to X \times_Y X$ 都平坦，
则称概形态射 $X \to Y$ 是{\it 弱 \'etale 的}或{\it 绝对平坦的}。
\end{definition}

\noindent
\'etale 态射是弱 \'etale 的；反过来，弱 \'etale 态射确实在某种程度上
类似 \'etale 态射。例如，若 $X \to Y$ 弱 \'etale，则
$L_{X/Y} = 0$，这是由
余切复形章引理 \ref{cotangent-lemma-when-zero}.
得到的。稍后将证明一个非常精确的结果，把弱 \'etale 态射与
\'etale 态射联系起来（见
Pro-\'etale 上同调章第 \ref{proetale-section-weakly-etale} 节).
本节只讨论基础内容。

\begin{lemma}
\label{more-morphisms-lemma-check-weakly-etale-stalks}
设 $f : X \to Y$ 是概形态射。下列条件等价：
\begin{enumerate}
\item $X \to Y$ 弱 \'etale；
\item 对每个 $x \in X$，环映射
$\mathcal{O}_{Y, f(x)} \to \mathcal{O}_{X, x}$ 弱 \'etale。
\end{enumerate}
\end{lemma}

\begin{proof}
注意，在假设 (1) 和 (2) 中，态射 $f$ 都平坦。因此可以假设 $f$ 平坦。
设 $x \in X$ 在 $Y$ 中的像为 $y = f(x)$。有典范环映射
$$
\mathcal{O}_{X, x} \otimes_{\mathcal{O}_{Y, y}} \mathcal{O}_{X, x}
\longrightarrow
\mathcal{O}_{X \times_Y X, \Delta_{X/Y}(x)}
\longrightarrow
\mathcal{O}_{X, x}
$$
，其中第一个映射是局部化（因而平坦），第二个映射是满射。
条件 (1) 意味着对所有 $x$，第二个箭头平坦；条件 (2) 意味着对所有
$x$，复合平坦。因此等价性由
代数章引理 \ref{algebra-lemma-flat-localization} 的第 (2) 部分得到。
\end{proof}

\begin{lemma}
\label{more-morphisms-lemma-key}
设 $X \to Y$ 是概形态射，使得 $X \to X \times_Y X$ 平坦。
设 $\mathcal{F}$ 是 $\mathcal{O}_X$-模。若 $\mathcal{F}$
在 $Y$ 上平坦，则 $\mathcal{F}$ 在 $X$ 上平坦。
\end{lemma}

\begin{proof}
% P05-EN-0139: frozen source uses f(x) although the morphism in the lemma is unnamed.
设 $x \in X$ 在 $Y$ 中的像为 $y = f(x)$。
由于 $X \to X \times_Y X$ 平坦，可知
$\mathcal{O}_{X, x} \otimes_{\mathcal{O}_{Y, y}} \mathcal{O}_{X, x} \to
\mathcal{O}_{X, x}$ 平坦。因此结论由
代数进阶章引理 \ref{more-algebra-lemma-key}
及定义得到。
\end{proof}

\begin{lemma}
\label{more-morphisms-lemma-weakly-etale-characterize}
设 $f : X \to S$ 是概形态射。下列条件等价：
\begin{enumerate}
\item 态射 $f$ 弱 \'etale。
\item 对每一对满足 $f(U) \subset V$ 的仿射开集 $U \subset X$、
$V \subset S$，环映射 $\mathcal{O}_S(V) \to \mathcal{O}_X(U)$ 弱 \'etale。
\item 存在开覆盖 $S = \bigcup_{j \in J} V_j$，以及开覆盖
$f^{-1}(V_j) = \bigcup_{i \in I_j} U_i$，使得对
$j\in J, i\in I_j$，每个态射 $U_i \to V_j$ 都弱 \'etale。
\item 存在仿射开覆盖 $S = \bigcup_{j \in J} V_j$，以及仿射开覆盖
$f^{-1}(V_j) = \bigcup_{i \in I_j} U_i$，使得对所有
$j\in J, i\in I_j$，环映射
$\mathcal{O}_S(V_j) \to \mathcal{O}_X(U_i)$ 都弱 \'etale。
\end{enumerate}
此外，若 $f$ 弱 \'etale，则对任意满足 $f(U) \subset V$
的开子概形 $U \subset X$、$V \subset S$，限制
$f|_U : U \to V$ 弱 \'etale。
\end{lemma}

\begin{proof}
设给定满足 $f(U) \subset V$ 的开子概形 $U \subset X$、$V \subset S$。
% P05-EN-0140: frozen source switches the base S to an undefined Y in the ambient fibre product and diagonal.
则 $U \times_V U \subset X \times_Y X$ 是开集
（概形章引理 \ref{schemes-lemma-open-fibre-product}），而
$f|_U : U \to V$ 的对角态射 $\Delta_{U/V}$ 是限制
$\Delta_{X/Y}|_U : U \to U \times_V U$。由于平坦性是概形态射的局部性质
（态射章引理 \ref{morphisms-lemma-flat-characterize}），
引理的最后一个陈述以及 (1) 与 (3) 的等价性都随之成立。
若 $X$ 和 $Y$ 仿射，则 $X \to Y$ 弱 \'etale，当且仅当
$\mathcal{O}_Y(Y) \to \mathcal{O}_X(X)$ 弱 \'etale（再次使用
态射章引理 \ref{morphisms-lemma-flat-characterize}).
因此，(1) 与 (3) 也等价于 (2) 与 (4)。
\end{proof}

\begin{lemma}
\label{more-morphisms-lemma-composition-weakly-etale}
设 $X \to Y \to Z$ 是概形态射。
\begin{enumerate}
\item 若 $X \to X \times_Y X$ 和 $Y \to Y \times_Z Y$ 平坦，
则 $X \to X \times_Z X$ 平坦。
\item 若 $X \to Y$ 和 $Y \to Z$ 弱 \'etale，则 $X \to Z$ 弱 \'etale。
\end{enumerate}
\end{lemma}

\begin{proof}
第 (1) 部分由 $X$ 在 $Z$ 上的对角态射的分解
$$
X \to X \times_Y X \to X \times_Z X
$$
、等式
$$
X \times_Y X = (X \times_Z X) \times_{(Y \times_Z Y)} Y,
$$
、平坦态射的基变换平坦这一事实，以及平坦态射的复合平坦这一事实得到
（态射章引理 \ref{morphisms-lemma-base-change-flat} 和
\ref{morphisms-lemma-composition-flat}）。
第 (2) 部分由第 (1) 部分以及刚刚使用的平坦态射复合仍平坦这一事实得到。
\end{proof}

\begin{lemma}
\label{more-morphisms-lemma-base-change-weakly-etale}
设 $X \to Y$ 和 $Y' \to Y$ 是概形态射，并设
$X' = Y' \times_Y X$ 是 $X$ 的基变换。
\begin{enumerate}
\item 若 $X \to X \times_Y X$ 平坦，则
$X' \to X' \times_{Y'} X'$ 平坦。
\item 若 $X \to Y$ 弱 \'etale，则 $X' \to Y'$ 弱 \'etale。
\end{enumerate}
\end{lemma}

\begin{proof}
假设 $X \to X \times_Y X$ 平坦。态射 $X' \to X' \times_{Y'} X'$
是 $X \to X \times_Y X$ 沿 $Y' \to Y$ 的基变换。因此由
态射章引理 \ref{morphisms-lemma-base-change-flat}，它平坦。
这证明了 (1)。第 (2) 部分由 (1) 及刚刚使用的平坦态射的基变换平坦得到。
\end{proof}

\begin{lemma}
\label{more-morphisms-lemma-go-down}
设 $X \to Y \to Z$ 是概形态射。假设 $X \to Y$ 平坦且满射，
并且 $X \to X \times_Z X$ 平坦。则 $Y \to Y \times_Z Y$ 平坦。
\end{lemma}

\begin{proof}
考虑交换图
$$
\xymatrix{
X \ar[r] \ar[d] & X \times_Z X \ar[d] \\
Y \ar[r] & Y \times_Z Y
}
$$
上方的水平箭头平坦，竖直箭头也平坦。因此 $X$ 在
$Y \times_Z Y$ 上平坦。由
态射章引理 \ref{morphisms-lemma-flat-permanence}
可知 $Y$ 在 $Y \times_Z Y$ 上平坦。
\end{proof}

\begin{lemma}
\label{more-morphisms-lemma-weakly-etale-formally-unramified}
设 $f : X \to Y$ 是弱 \'etale 概形态射。
则 $f$ 形式非分歧，即 $\Omega_{X/Y} = 0$。
\end{lemma}

\begin{proof}
回忆，由
引理 \ref{more-morphisms-lemma-formally-unramified-differentials}.
，$f$ 形式非分歧当且仅当 $\Omega_{X/Y} = 0$。通过
引理 \ref{more-morphisms-lemma-weakly-etale-characterize} 和
态射章引理 \ref{morphisms-lemma-differentials-affine}
，这归结为环的情形，即
代数进阶章引理 \ref{more-algebra-lemma-formally-unramified}.
\end{proof}

\begin{lemma}
\label{more-morphisms-lemma-when-weakly-etale}
设 $f : X \to Y$ 是概形态射。在下列每种情形中，$X \to Y$ 都弱 \'etale：
\begin{enumerate}
\item $X \to Y$ 是平坦单态射；
\item $X \to Y$ 是开浸入；
\item $X \to Y$ 平坦且非分歧；
\item $X \to Y$ 是 \'etale 态射。
\end{enumerate}
\end{lemma}

\begin{proof}
若 (1) 成立，则 $\Delta_{X/Y}$ 是同构
（概形章引理 \ref{schemes-lemma-monomorphism}），故 $f$ 当然弱 \'etale。
情形 (2) 是 (1) 的特殊情形。非分歧态射的对角态射是开浸入
（态射章引理 \ref{morphisms-lemma-diagonal-unramified-morphism}），
因而平坦。所以平坦非分歧态射弱 \'etale。
\'etale 态射平坦且非分歧
（态射章引理 \ref{morphisms-lemma-etale-smooth-unramified}），
故 (4) 由 (3) 得到。
\end{proof}

\begin{lemma}
\label{more-morphisms-lemma-reduced-goes-up}
设 $f : X \to Y$ 是概形态射。若 $Y$ 既约且 $f$ 弱 \'etale，
则 $X$ 既约。
\end{lemma}

\begin{proof}
通过 引理 \ref{more-morphisms-lemma-weakly-etale-characterize}，
这归结为环的情形，即
代数进阶章引理
\ref{more-algebra-lemma-absolutely-flat-over-absolutely-flat}.
\end{proof}

\noindent
下面的引理使用了关于弱 \'etale 环映射的一个非平凡结果。

\begin{lemma}
\label{more-morphisms-lemma-weakly-etale-strictly-henselian-local-rings}
设 $f : X \to Y$ 是概形态射。下列条件等价：
\begin{enumerate}
\item $f$ 弱 \'etale；
\item 对 $x \in X$，局部环映射
$\mathcal{O}_{Y, f(x)} \to \mathcal{O}_{X, x}$
在严格 Hensel 化上诱导同构。
\end{enumerate}
\end{lemma}

\begin{proof}
设 $x \in X$ 在 $Y$ 中的像为 $y = f(x)$。
选取 $\kappa(x)$ 的可分代数闭包 $\kappa^{sep}$。
设 $\mathcal{O}_{X, x}^{sh}$ 是对应于 $\kappa^{sep}$ 的严格 Hensel 化，
而 $\mathcal{O}_{Y, y}^{sh}$ 是相对于 $\kappa(y)$ 在
$\kappa^{sep}$ 中的可分代数闭包的严格 Hensel 化。考虑交换图
$$
\xymatrix{
\mathcal{O}_{X, x} \ar[r] & \mathcal{O}_{X, x}^{sh} \\
\mathcal{O}_{Y, y} \ar[u] \ar[r] & \mathcal{O}_{Y, y}^{sh} \ar[u]
}
$$
，其中各映射为局部环的局部同态，见
代数章引理 \ref{algebra-lemma-strictly-henselian-functorial}.
由于严格 Hensel 化是 \'etale 环映射的滤过余极限，
代数进阶章引理 \ref{more-algebra-lemma-when-weakly-etale}
表明水平映射弱 \'etale。此外，由
代数进阶章引理 \ref{more-algebra-lemma-dumb-properties-henselization}.
，水平映射忠实平坦。

\medskip\noindent
假设 $f$ 弱 \'etale。由 引理 \ref{more-morphisms-lemma-check-weakly-etale-stalks}，
左侧竖直箭头弱 \'etale。由
代数进阶章引理 \ref{more-algebra-lemma-composition-weakly-etale} 及
\ref{more-algebra-lemma-weakly-etale-permanence}
可知右侧竖直箭头弱 \'etale。由
代数进阶章定理 \ref{more-algebra-theorem-olivier}
可知右侧竖直映射是同构。

\medskip\noindent
假设 $\mathcal{O}_{Y, y}^{sh} \to \mathcal{O}_{X, x}^{sh}$ 是同构。
则 $\mathcal{O}_{Y, y} \to \mathcal{O}_{X, x}^{sh}$ 弱 \'etale。
由于 $\mathcal{O}_{X, x} \to \mathcal{O}_{X, x}^{sh}$ 忠实平坦，
由
代数进阶章引理 \ref{more-algebra-lemma-go-down}.
可知 $\mathcal{O}_{Y, y} \to \mathcal{O}_{X, x}$ 弱 \'etale。
因此由 引理 \ref{more-morphisms-lemma-check-weakly-etale-stalks}，(2) 蕴含 (1)。
\end{proof}

\begin{lemma}
\label{more-morphisms-lemma-normal-goes-up}
设 $f : X \to Y$ 是概形态射。若 $Y$ 是正规概形且 $f$ 弱 \'etale，
则 $X$ 是正规概形。
\end{lemma}

\begin{proof}
由 代数进阶章引理 \ref{more-algebra-lemma-henselization-normal}，
概形 $S$ 正规，当且仅当对所有 $s \in S$，
$\mathcal{O}_{S, s}$ 的严格 Hensel 化都是正规整环。因此引理由
引理 \ref{more-morphisms-lemma-weakly-etale-strictly-henselian-local-rings}.
得到。
\end{proof}

\begin{lemma}
\label{more-morphisms-lemma-weakly-etale-permanence}
设 $S$ 是概形，$f : X \to Y$ 是 $S$ 上的概形态射。
若 $X$、$Y$ 在 $S$ 上弱 \'etale，则 $f$ 弱 \'etale。
\end{lemma}

\begin{proof}
我们将不再另行说明地使用 态射章引理
\ref{morphisms-lemma-base-change-flat} 和
\ref{morphisms-lemma-composition-flat}。把 $X \to Y$ 写成复合
$X \to X \times_S Y \to Y$。第二个态射是平坦态射 $X \to S$
的基变换，因而平坦。第一个态射是平坦态射
$Y \to Y \times_S Y$ 沿态射 $X \times_S Y \to Y \times_S Y$
的基变换，因而也平坦。所以 $X \to Y$ 平坦。
态射 $X \times_Y X \to X \times_S X$ 是浸入。
因此 引理 \ref{more-morphisms-lemma-key} 蕴含：由于 $X$ 在 $X \times_S X$
上平坦，$X$ 在 $X \times_Y X$ 上也平坦。
\end{proof}

\noindent
下面是这样一个观察的概形论推广：同时可分且纯不可分的域扩张必为同构。

\begin{lemma}
\label{more-morphisms-lemma-weakly-etale-universal-homeomorphism}
设 $f : X \to Y$ 是概形态射。若 $f$ 弱 \'etale 且为普遍同胚，
则它是同构。
\end{lemma}

\begin{proof}
因为 $f$ 是普遍同胚，由
态射章引理 \ref{morphisms-lemma-homeomorphism-affine} 及
\ref{morphisms-lemma-universally-injective}，对角态射
$\Delta : X \to X \times_Y X$ 是满射闭浸入。由于 $\Delta$ 也平坦，
由
态射章引理 \ref{morphisms-lemma-characterize-flat-closed-immersions}.
可知 $\Delta$ 必为同构。换言之，$f$ 是单态射
（概形章引理 \ref{schemes-lemma-monomorphism}）。
由于 $f$ 是普遍同胚，它当然拟紧。因此由 下降章引理
\ref{descent-lemma-flat-surjective-quasi-compact-monomorphism-isomorphism}
可知 $f$ 是同构。
\end{proof}

\noindent
下面是
\'Etale 态射章引理 \ref{etale-lemma-relative-frobenius-etale}.
的弱 \'etale 推广。

\begin{lemma}
\label{more-morphisms-lemma-relative-frobenius-weakly-etale}
设 $U \to X$ 是弱 \'etale 概形态射，其中 $X$ 是特征 $p$ 的概形。
则相对 Frobenius 态射 $F_{U/X} : U \to U \times_{X, F_X} X$ 是同构。
\end{lemma}

\begin{proof}
由 簇章引理 \ref{varieties-lemma-relative-frobenius}，
$F_{U/X}$ 是普遍同胚。由于它是 $X$ 上弱 \'etale 概形之间的态射，
由 引理 \ref{more-morphisms-lemma-weakly-etale-permanence}
可知 $F_{U/X}$ 弱 \'etale。因此由
引理 \ref{more-morphisms-lemma-weakly-etale-universal-homeomorphism}，$F_{U/X}$ 是同构。
\end{proof}






\section{约化纤维定理}
\label{more-morphisms-section-reduced-fibre-theorem}

\noindent
本节讨论题名所预告的此类定理中最简单的一种。虽然该结果的证明
稍显繁复，但从本质上说，它直接由 Epp 关于消除分歧的结果推出，见
代数进阶，定理 \ref{more-algebra-theorem-epp}。

\medskip\noindent
设 $A$ 是分式域为 $K$ 的戴德金整环。
设 $X$ 是 $A$ 上平坦且有限型的概形。
设 $L$ 是 $K$ 的有限扩张，并设 $B$ 是 $A$ 在 $L$ 中的整闭包。
于是 $B$ 是戴德金整环
（代数，引理 \ref{algebra-lemma-integral-closure-Dedekind}）。
令 $X_B = X \times_{\Spec(A)} \Spec(B)$ 为相应的基变换。
则 $X_B \to \Spec(B)$ 是有限型的
（态射，引理 \ref{morphisms-lemma-base-change-finite-type}）。
因此 $X_B$ 是诺特概形
（态射，引理 \ref{morphisms-lemma-finite-type-noetherian}）。
从而正规化 $\nu : Y \to X_B$ 存在（见
态射，定义 \ref{morphisms-definition-normalization}
及其后的讨论）。图示为
\begin{equation}
\label{more-morphisms-equation-normalized-base-change}
\xymatrix{
Y \ar[rd] \ar[r]_\nu & X_B \ar[r] \ar[d] & X \ar[d] \\
 & \Spec(B) \ar[r] & \Spec(A)
}
\end{equation}
我们有时称 $Y$ 为 $X$ 的{\it 正规化基变换}。
一般而言，态射 $\nu$ 未必有限。但若 $A$ 是纳伽塔环（这是一项在实践中
几乎总能满足的条件），则 $\nu$ 有限，且 $Y$ 在 $B$ 上有限型，见
态射，引理 \ref{morphisms-lemma-nagata-normalization} 和
\ref{morphisms-lemma-finite-type-nagata}。

\medskip\noindent
取正规化基变换与复合相容。更确切地说，若 $M/L/K$ 是域的有限扩张，
相应的整闭包为 $A \subset B \subset C$，则 $Y \to \Spec(B)$
关于 $M/L$ 的正规化基变换 $Z$，等于 $X \to \Spec(A)$
关于 $M/K$ 的正规化基变换。

\begin{theorem}
\label{more-morphisms-theorem-normalized-base-change-with-reduced-fibre}
设 $A$ 是分式域为 $K$ 的戴德金环。
设 $X$ 是 $A$ 上平坦且有限型的概形，并假设 $A$ 是纳伽塔环。
存在有限扩张 $L/K$，使得正规化基变换 $Y$ 在所有纤维的全部泛点处
都在 $\Spec(B)$ 上光滑。
\end{theorem}

\begin{proof}
在证明中，我们将反复使用如下事实：对于 $X \to \Spec(A)$ 这样的
（平坦、有限表示）态射，其光滑点集的形成与基变换相容，见
态射，引理 \ref{morphisms-lemma-set-points-where-fibres-smooth}。

\medskip\noindent
先选择有限扩张 $L/K$，使得 $(X_L)_{red}$ 在 $L$ 上几何约化，见
簇，引理 \ref{varieties-lemma-finite-extension-geometrically-reduced}。
由于 $Y \to (X_B)_{red}$ 是双有理的，应用
簇，引理 \ref{varieties-lemma-generic-points-geometrically-reduced}
可知 $Y_L$ 也在 $L$ 上几何约化。因此由
簇，引理 \ref{varieties-lemma-geometrically-reduced-dense-smooth-open}，
$Y_L \to \Spec(L)$ 在某个稠密开子集 $V \subset Y_L$ 上光滑。
于是态射 $Y \to \Spec(B)$ 的光滑轨迹 $U \subset Y$ 是开集
（由态射，定义 \ref{morphisms-definition-smooth}），
并且在泛纤维中稠密。以 $B$ 代替 $A$、以 $Y$ 代替 $X$，
便归约到下一段所处理的情形。

\medskip\noindent
假设 $X$ 正规，且 $X \to \Spec(A)$ 的光滑轨迹 $U \subset X$
在泛纤维中稠密。这蕴含 $U$ 在除有限多个之外的所有纤维中都稠密，
见 引理 \ref{more-morphisms-lemma-nowhere-dense-generic-fibre}。
令 $x_1, \ldots, x_r \in X \setminus U$ 为如下有限多个点：它们既是
$X \setminus U$ 的不可约分支的泛点，又是 $X \to \Spec(A)$ 的纤维的
不可约分支的泛点。置 $\mathcal{O}_i = \mathcal{O}_{X, x_i}$。
令 $A_i$ 为 $A$ 在如下极大理想处的局部化：该极大理想对应于 $x_i$
在 $\Spec(A)$ 中的像。由
代数进阶，Proposition
\ref{more-algebra-proposition-epp-essentially-finite-type}，
存在有限扩张 $K_i/K$，它们是离散赋值环扩张
$A_i \to \mathcal{O}_i$ 的解。取一个支配所有扩张 $K_i/K$ 的有限扩张
$L/K$。由代数进阶，引理 \ref{more-algebra-lemma-solution-goes-up}，
$L/K$ 仍是 $A_i \to \mathcal{O}_i$ 的解。

\medskip\noindent
考虑刚构造的扩张 $L/K$ 所对应的图
(\ref{more-morphisms-equation-normalized-base-change})。注意 $U_B \subset X_B$ 在 $B$
上光滑，因而正规（例如使用代数，Lemma
\ref{algebra-lemma-normal-goes-up}）。故 $Y \to X_B$ 在 $U_B$ 上是同构。
设 $y \in Y$ 是 $Y \to \Spec(B)$ 的某个纤维的一个不可约分支的泛点，
该纤维位于极大理想 $\mathfrak m \subset B$ 之上。假设
$y \not \in U_B$。则 $y$ 映到某个点 $x_i$。由此可知
$\mathcal{O}_{Y, y}$ 是 $\mathcal{O}_i$ 在 $R(X) \otimes_K L$
中的整闭包的一个局部环（略去细节）。由于 $L/K$ 是
$A_i \to \mathcal{O}_i$ 的解，故
$B_\mathfrak m \to \mathcal{O}_{Y, y}$ 对 $\mathfrak m_y$-进拓扑形式光滑
（这正是成为“解”的定义）。换言之，
$\mathfrak m\mathcal{O}_{Y, y} = \mathfrak m_y$，并且剩余域扩张可分，见
代数进阶，引理 \ref{more-algebra-lemma-extension-dvrs-formally-smooth}。
因此纤维在 $y$ 处的局部环是 $\kappa(y)$。例如由代数，Lemma
\ref{algebra-lemma-separable-smooth}，这蕴含该纤维在 $y$ 处于
$\kappa(\mathfrak m)$ 上光滑。证明完成。
\end{proof}

\begin{lemma}[曲线上的变体]
\label{more-morphisms-lemma-normalized-base-change-with-reduced-fibre-over-curve}
设 $f : X \to S$ 是平坦有限型的概形态射。假设 $S$ 是纳伽塔整概形，
其函数域为 $K$，并且是维数为 $1$ 的正则概形。
则存在有限扩张 $L/K$，使得在图
$$
\xymatrix{
Y \ar[rd]_g \ar[r]_-\nu & X \times_S T \ar[d] \ar[r] & X \ar[d]_f \\
& T \ar[r] & S
}
$$
中，态射 $g$ 在纤维的所有泛点处光滑。这里，$T$ 是 $S$ 在
$\Spec(L)$ 中的正规化，而 $\nu : Y \to X \times_S T$ 是正规化。
\end{lemma}

\begin{proof}
选择有限仿射开覆盖 $S = \bigcup \Spec(A_i)$。于是对每个 $i$，
$K$ 都等于 $A_i$ 的分式域。令 $X_i = X \times_S \Spec(A_i)$。
对态射 $X_i \to \Spec(A_i)$，按照
定理 \ref{more-morphisms-theorem-normalized-base-change-with-reduced-fibre}
选择 $L_i/K$。令 $B_i \subset L_i$ 为 $A_i$ 的整闭包，并令 $Y_i$
为 $X$ 到 $B_i$ 的正规化基变换。取支配每个 $L_i$ 的有限扩张
$L/K$。令 $T_i \subset T$ 为 $\Spec(A_i)$ 的原像。对每个 $i$，
得到交换图
$$
\xymatrix{
g^{-1}(T_i) \ar[r] \ar[d] & Y_i \ar[r] \ar[d] & X \times_S \Spec(A_i) \ar[d] \\
T_i \ar[r] & \Spec(B_i) \ar[r] & \Spec(A_i)
}
$$
而且左侧方块正是本节开头所讨论的正规化基变换。在
定理 \ref{more-morphisms-theorem-normalized-base-change-with-reduced-fibre}
的证明中，我们已经看到，$Y \to T$ 的光滑轨迹包含如下点集在
$g^{-1}(T_i)$ 中的原像：$Y_i$ 在这些点处于 $B_i$ 上光滑。
这就证明了引理。
\end{proof}

\begin{lemma}[可分扩张情形的变体]
\label{more-morphisms-lemma-normalized-base-change-with-reduced-fibre-separable}
设 $A$ 是分式域为 $K$ 的戴德金环。设 $X$ 是 $A$ 上平坦且有限型的概形。
假设 $A$ 是纳伽塔环，并且对 $X$ 的每个不可约分支的每个泛点
$\eta$，域扩张 $\kappa(\eta)/K$ 都可分。则存在有限可分扩张 $L/K$，
使得正规化基变换 $Y$ 在所有纤维的全部泛点处都在 $\Spec(B)$ 上光滑。
\end{lemma}

\begin{proof}
证明方法与 定理 \ref{more-morphisms-theorem-normalized-base-change-with-reduced-fibre}
完全相同，只需作几处小修改。最重要的改动是使用代数进阶，Lemma
\ref{more-algebra-lemma-epp-essentially-finite-type-separable}，
而不是代数进阶，Proposition
\ref{more-algebra-proposition-epp-essentially-finite-type}。
在证明中，我们将反复使用如下事实：对于 $X \to \Spec(A)$ 这样的
（平坦、有限表示）态射，其光滑点集的形成与基变换相容，见
态射，引理 \ref{morphisms-lemma-set-points-where-fibres-smooth}。

\medskip\noindent
由于 $X$ 在 $A$ 上平坦，$X$ 的每个泛点 $\eta$ 都映到 $\Spec(A)$
的泛点。以 $X$ 的约化概形代替它后，可以假设 $X$ 约化。
在这种情形下，由簇，Lemma
\ref{varieties-lemma-generic-points-geometrically-reduced}，$X_K$ 在 $K$
上几何约化。因此由簇，Lemma
\ref{varieties-lemma-geometrically-reduced-dense-smooth-open}，
$X_K \to \Spec(K)$ 在某个稠密开集上光滑。于是态射
$X \to \Spec(A)$ 的光滑轨迹 $U \subset X$ 是开集
（由态射，定义 \ref{morphisms-definition-smooth}），
并且在泛纤维中稠密。这就把问题归约到下一段的情形。

\medskip\noindent
假设 $X$ 正规，且 $X \to \Spec(A)$ 的光滑轨迹 $U \subset X$
在泛纤维中稠密。这蕴含 $U$ 在除有限多个之外的所有纤维中都稠密，
见 引理 \ref{more-morphisms-lemma-nowhere-dense-generic-fibre}。
令 $x_1, \ldots, x_r \in X \setminus U$ 为如下有限多个点：它们既是
$X \setminus U$ 的不可约分支的泛点，又是 $X \to \Spec(A)$ 的纤维的
不可约分支的泛点。置 $\mathcal{O}_i = \mathcal{O}_{X, x_i}$。
注意，$\mathcal{O}_i$ 的分式域是 $X$ 的某个泛点的剩余域。
令 $A_i$ 为 $A$ 在如下极大理想处的局部化：该极大理想对应于 $x_i$
在 $\Spec(A)$ 中的像。我们可以应用代数进阶，Lemma
\ref{more-algebra-lemma-epp-essentially-finite-type-separable}，
得到有限可分扩张 $K_i/K$，它们是 $A_i \to \mathcal{O}_i$ 的解。
取一个支配所有扩张 $K_i/K$ 的有限可分扩张 $L/K$。由代数进阶，
引理 \ref{more-algebra-lemma-solution-goes-up}，$L/K$ 仍是
$A_i \to \mathcal{O}_i$ 的解。

\medskip\noindent
考虑刚构造的扩张 $L/K$ 所对应的图
(\ref{more-morphisms-equation-normalized-base-change})。注意 $U_B \subset X_B$ 在 $B$
上光滑，因而正规（例如使用代数，Lemma
\ref{algebra-lemma-normal-goes-up}）。故 $Y \to X_B$ 在 $U_B$ 上是同构。
设 $y \in Y$ 是 $Y \to \Spec(B)$ 的某个纤维的一个不可约分支的泛点，
该纤维位于极大理想 $\mathfrak m \subset B$ 之上。假设
$y \not \in U_B$。则 $y$ 映到某个点 $x_i$。由此可知
$\mathcal{O}_{Y, y}$ 是 $\mathcal{O}_i$ 在 $R(X) \otimes_K L$
中的整闭包的一个局部环（略去细节）。由于 $L/K$ 是
$A_i \to \mathcal{O}_i$ 的解，故
$B_\mathfrak m \to \mathcal{O}_{Y, y}$ 形式光滑
（这正是成为“解”的定义）。换言之，
$\mathfrak m\mathcal{O}_{Y, y} = \mathfrak m_y$，并且剩余域扩张可分。
因此纤维在 $y$ 处的局部环是 $\kappa(y)$。例如由代数，Lemma
\ref{algebra-lemma-separable-smooth}，这蕴含该纤维在 $y$ 处于
$\kappa(\mathfrak m)$ 上光滑。证明完成。
\end{proof}

\begin{lemma}[曲线上的可分扩张变体]
\label{more-morphisms-lemma-normalized-base-change-with-reduced-fibre-over-curve-separable}
设 $f : X \to S$ 是平坦有限型的概形态射。假设 $S$ 是纳伽塔整概形，
其函数域为 $K$，并且是维数为 $1$ 的正则概形。假设对 $X$
的每个不可约分支的每个泛点 $\eta$，域扩张 $\kappa(\eta)/K$
都可分。则存在有限可分扩张 $L/K$，使得在图
$$
\xymatrix{
Y \ar[rd]_g \ar[r]_-\nu & X \times_S T \ar[d] \ar[r] & X \ar[d]_f \\
& T \ar[r] & S
}
$$
中，态射 $g$ 在纤维的所有泛点处光滑。这里，$T$ 是 $S$ 在
$\Spec(L)$ 中的正规化，而 $\nu : Y \to X \times_S T$ 是正规化。
\end{lemma}

\begin{proof}
这由 Lemma
\ref{more-morphisms-lemma-normalized-base-change-with-reduced-fibre-separable}
推出，其方式与 Lemma
\ref{more-morphisms-lemma-normalized-base-change-with-reduced-fibre-over-curve}
由 定理 \ref{more-morphisms-theorem-normalized-base-change-with-reduced-fibre}
推出的方式完全相同。
\end{proof}








\section{Ind-拟仿射态射}
\label{more-morphisms-section-ind-quasi-affine}

\noindent
下面给出稍后要用到的一点理论。

\begin{definition}
\label{more-morphisms-definition-ind-quasi-affine}
如果 $X$ 的每个拟紧开子概形都是拟仿射的，则称概形 $X$ 是
{\it ind-拟仿射的}。类似地，如果对 $Y$ 中的每个仿射开集 $V$，
$f^{-1}(V)$ 都是 ind-拟仿射的，则称概形态射 $X \to Y$ 是
{\it ind-拟仿射的}。
\end{definition}

\noindent
仿射概形的任一开子概形都是 ind-拟仿射概形。若
$X = \bigcup_{i \in I} U_i$ 是拟仿射开集的并，且任意两个 $U_i$
都包含在第三个这样的开集中，则 $X$ 是 ind-拟仿射的。
Ind-拟仿射概形 $X$ 是分离的，因为任意两个仿射开集 $U, V$
都包含在 $X$ 的一个分离开子概形中，即 $U \cup V$ 中。
类似地，ind-拟仿射态射是分离的。

\begin{lemma}
\label{more-morphisms-lemma-ind-quasi-affine-alternative-definition}
对概形态射 $f : X \to Y$，下列条件等价：
\begin{enumerate}
\item $f$ 是 ind-拟仿射的；
\item 对每个仿射开子概形 $V \subset Y$ 以及每个拟紧开子概形
$U \subset f^{-1}(V)$，所诱导的态射 $U \to V$ 是拟仿射的；
\item
存在 $Y$ 的一个由拟紧拟分离开子概形 $V_j \subset Y$ 组成的覆盖
$\{ V_j \}_{j \in J}$，使得对每个 $j \in J$ 以及每个拟紧开子概形
$U \subset f^{-1}(V_j)$，所诱导的态射 $U \to V_j$ 是拟仿射的；
\item 对每个拟紧拟分离开子概形 $V \subset Y$ 以及每个拟紧开子概形
$U \subset f^{-1}(V)$，所诱导的态射 $U \to V$ 是拟仿射的。
\end{enumerate}
特别地，成为 ind-拟仿射态射这一性质在基底上是 Zariski 局部的。
\end{lemma}

\begin{proof}
(1) $\Leftrightarrow$ (2) 的等价性由定义和
态射，引理 \ref{morphisms-lemma-characterize-quasi-affine} 得出。
为证 (2) $\Rightarrow$ (4)，令 $U$ 和 $V$ 如 (4) 所述。由
概形，引理 \ref{schemes-lemma-quasi-compact-permanence}，所诱导的态射
$U \to V$ 是拟紧的。因此对每个仿射开集 $V' \subset V$，纤维积
$V' \times_V U$ 是拟紧的，所以由 (2)，所诱导的映射
$V' \times_V U \to V'$ 是拟仿射的。于是由
态射，引理 \ref{morphisms-lemma-characterize-quasi-affine}，
$U \to V$ 也是拟仿射的。
这一论证也给出 (3) $\Rightarrow$ (4)：沿用相同记号，包含在某个
$V_j$ 中的那些仿射开集 $V' \subset V$ 覆盖 $V$，因此只需证明拟紧映射
$V' \times_V U \to V'$ 是拟仿射的。然而由 (3)，复合
$V' \times_V U \to V' \to V_j$ 是拟仿射的；而由
概形，引理 \ref{schemes-lemma-compose-after-separated}，映射
$V' \to V_j$ 是拟分离的。因此由
态射，引理 \ref{morphisms-lemma-quasi-affine-permanence}，
$V' \times_V U \to V'$ 是拟仿射的。最后的蕴含
(4) $\Rightarrow$ (2) 和 (4) $\Rightarrow$ (3) 是显然的。
\end{proof}

\begin{lemma}
\label{more-morphisms-lemma-ind-quasi-affine-composition}
成为 ind-拟仿射态射这一性质在复合下稳定。
\end{lemma}

\begin{proof}
设 $f : X \to Y$ 和 $g : Y \to Z$ 是 ind-拟仿射态射。设
$V \subset Z$ 和 $U \subset f^{-1}(g^{-1}(V))$ 是拟紧开集，并且
$V$ 还是拟分离的。像 $f(U)$ 是 $g^{-1}(V)$ 的拟紧子集，故包含在某个
拟紧开集 $W \subset g^{-1}(V)$ 中（它是有限多个仿射开集的并）。
我们得到分解 $U \to W \to V$。由
引理 \ref{more-morphisms-lemma-ind-quasi-affine-alternative-definition}，映射
$W \to V$ 是拟仿射的，特别地，$W$ 是拟分离的。再次应用
引理 \ref{more-morphisms-lemma-ind-quasi-affine-alternative-definition}，可知
$U \to W$ 也是拟仿射的。因此由态射，Lemma
\ref{morphisms-lemma-composition-quasi-affine}，复合 $U \to V$
也是拟仿射的；最后再应用一次
引理 \ref{more-morphisms-lemma-ind-quasi-affine-alternative-definition} 即可。
\end{proof}

\begin{lemma}
\label{more-morphisms-lemma-ind-quasi-affine-examples}
任一拟仿射态射都是 ind-拟仿射的。
任一浸入都是 ind-拟仿射的。
\end{lemma}

\begin{proof}
第一项断言直接由定义得出。特别地，仿射态射（例如闭浸入）
是 ind-拟仿射的。因此由 Lemma
\ref{more-morphisms-lemma-ind-quasi-affine-composition}，只需证明开浸入是
ind-拟仿射的，而这也直接由定义得出。
\end{proof}

\begin{lemma}
\label{more-morphisms-lemma-ind-quasi-affine-permanence}
设 $f : X \to Y$ 和 $g : Y \to Z$ 是概形态射。若 $g \circ f$
是 ind-拟仿射的，则 $f$ 是 ind-拟仿射的。
\end{lemma}

\begin{proof}
由 引理 \ref{more-morphisms-lemma-ind-quasi-affine-alternative-definition}，可以先在
$Z$ 上、再在 $Y$ 上 Zariski 局部地工作，故不妨假设 $Z$ 仿射，
进而也假设 $Y$ 仿射。此时 $X$ 的任一拟紧开集都是拟仿射的，
故结论由 引理 \ref{more-morphisms-lemma-ind-quasi-affine-alternative-definition} 得出。
\end{proof}

\begin{lemma}
\label{more-morphisms-lemma-base-change-ind-quasi-affine}
成为 ind-拟仿射态射这一性质在基变换下稳定。
\end{lemma}

\begin{proof}
设 $f : X \to Y$ 是 ind-拟仿射态射。为检验 $f$ 的每个基变换都是
ind-拟仿射的，由 Lemma
\ref{more-morphisms-lemma-ind-quasi-affine-alternative-definition}，可以在 $Y$ 上
Zariski 局部地工作，故假设 $Y$ 仿射。进一步，也可以假设基变换态射
$Z \to Y$ 中的概形 $Z$ 仿射。基变换 $X \times_Y Z \to X$ 是仿射态射，
故由 引理 \ref{more-morphisms-lemma-ind-quasi-affine-composition} 和
\ref{more-morphisms-lemma-ind-quasi-affine-examples}，映射 $X \times_Y Z \to Y$
是 ind-拟仿射的。再由 引理 \ref{more-morphisms-lemma-ind-quasi-affine-permanence}，
基变换 $X \times_Y Z \to Z$ 是 ind-拟仿射的，正合所需。
\end{proof}

\begin{lemma}
\label{more-morphisms-lemma-descending-property-ind-quasi-affine}
成为 ind-拟仿射态射这一性质在基底上是 fpqc 局部的。
\end{lemma}

\begin{proof}
引理 \ref{more-morphisms-lemma-base-change-ind-quasi-affine} 给出的 ind-拟仿射性在
基变换下的稳定性证明了一个方向。反过来，设 $f : X \to Y$ 是概形态射，
$\{g_i : Y_i \to Y\}$ 是一个 fpqc 覆盖，并且对每个 $i$，基变换
$f_i : X_i \to Y_i$ 都是 ind-拟仿射的。我们要证明 $f$ 是
ind-拟仿射的。

\medskip\noindent
由 引理 \ref{more-morphisms-lemma-ind-quasi-affine-alternative-definition}，可以在 $Y$
上 Zariski 局部地工作，故假设 $Y$ 仿射。再利用
引理 \ref{more-morphisms-lemma-base-change-ind-quasi-affine} 所保证的基变换稳定性细化覆盖，
并假设该覆盖由单个仿射忠实平坦态射 $g : Y' \to Y$ 给出。
对任一拟紧开集 $U \subset X$，其 $Y'$-基变换
$U \times_Y Y' \subset X \times_Y Y'$ 也是拟紧的。最后注意，由
下降，引理 \ref{descent-lemma-descending-property-quasi-affine}，
映射 $U \to Y$ 是拟仿射的，当且仅当 $U \times_Y Y' \to Y'$ 是拟仿射的。
\end{proof}

\begin{lemma}
\label{more-morphisms-lemma-etale-separated-ind-quasi-affine}
分离局部拟有限的概形态射是 ind-拟仿射的。
\end{lemma}

\begin{proof}
设 $f : X \to Y$ 是分离局部拟有限的概形态射。设 $V \subset Y$
仿射，而 $U \subset f^{-1}(V)$ 是拟紧开集。我们要证明 $U$ 拟仿射。
由于 $U \to V$ 是分离拟有限的概形态射，结论由 Zariski 主定理得出。
见 引理 \ref{more-morphisms-lemma-quasi-finite-separated-quasi-affine}。
\end{proof}




\section{概形范畴中的推出，II}
\label{more-morphisms-section-pushouts-II}

\noindent
本节是 第 \ref{more-morphisms-section-pushouts} 节 的延续。我们将在 $Z \to X$
为闭浸入、$Z \to Y$ 为整态射且满足一项附加条件时，构造
$Y \leftarrow Z \rightarrow X$ 的推出。详细讨论见
\cite{Ferrand-Conducteur}。

\begin{situation}
\label{more-morphisms-situation-pushout-along-closed-immersion-and-integral}
这里，$S$ 是概形，$i : Z \to X$ 和 $j : Z \to Y$ 是 $S$-概形的态射。
我们假设
\begin{enumerate}
\item $i$ 是闭浸入；
\item $j$ 是整的概形态射；
\item 对 $y \in Y$，存在仿射开集 $U \subset X$，使得
$j^{-1}(\{y\}) \subset i^{-1}(U)$。
\end{enumerate}
\end{situation}

\begin{lemma}
\label{more-morphisms-lemma-prepare-pushout-along-closed-immersion-and-integral}
在 Situation \ref{more-morphisms-situation-pushout-along-closed-immersion-and-integral}
中，对 $y \in Y$，存在仿射开集 $U \subset X$ 和 $V \subset Y$，
使得 $i^{-1}(U) = j^{-1}(V)$ 且 $y \in V$。
\end{lemma}

\begin{proof}
令 $y \in Y$。选择仿射开集 $U \subset X$，使得
$j^{-1}(\{y\}) \subset i^{-1}(U)$（由假设可以做到）。选择 $y$ 的一个
仿射开邻域 $V \subset Y$，使得 $j^{-1}(V) \subset i^{-1}(U)$。
这是可能的，因为 $j : Z \to Y$ 是闭态射
（态射，引理 \ref{morphisms-lemma-integral-universally-closed}），
且 $i^{-1}(U)$ 包含 $y$ 上的纤维。由于 $j$ 是整的，概形论纤维 $Z_y$
是某个域上的整代数的谱。由极限，Lemma
\ref{limits-lemma-ample-profinite-set-in-principal-affine}，可以找到
$\overline{f} \in \Gamma(i^{-1}(U), \mathcal{O}_{i^{-1}(U)})$，使得
$Z_y \subset D(\overline{f}) \subset j^{-1}(V)$。由于
$i|_{i^{-1}(U)} : i^{-1}(U) \to U$ 是仿射概形之间的闭浸入，可以选择
$f \in \Gamma(U, \mathcal{O}_U)$，其在 $i^{-1}(U)$ 上的限制是
$\overline{f}$。以主开集 $D(f) \subset U$ 代替 $U$ 后，得到仿射开集
$y \in V \subset Y$ 和 $U \subset X$，满足
$$
j^{-1}(\{y\}) \subset i^{-1}(U) \subset j^{-1}(V)
$$
现在我们（在某种意义下）重复这一论证。具体来说，选择
$g \in \Gamma(V, \mathcal{O}_V)$，使得 $y \in D(g)$ 且
$j^{-1}(D(g)) \subset i^{-1}(U)$（与上面相同的论证保证这可以做到）。
然后可以取 $f \in \Gamma(U, \mathcal{O}_U)$，使其在 $i^{-1}(U)$
上的限制是 $g$ 沿 $i^{-1}(U) \to V$ 的拉回（仍由上述同一理由）。
于是最终得到仿射开集 $y \in V' = D(g) \subset V \subset Y$ 和
% P05-EN-0141: frozen source writes i^{-1}(V') although V' is an open of Y; the defined open of X is U'.
$U' = D(f) \subset U \subset X$，满足 $j^{-1}(V') = i^{-1}(V')$。
\end{proof}

\begin{proposition}
\label{more-morphisms-proposition-pushout-along-closed-immersion-and-integral}
\begin{reference}
\cite[Theorem 7.1 part iii]{Ferrand-Conducteur}
\end{reference}
在 Situation \ref{more-morphisms-situation-pushout-along-closed-immersion-and-integral}
中，推出 $Y \amalg_Z X$ 在概形范畴中存在。图示为
$$
\xymatrix{
Z \ar[r]_i \ar[d]_j & X \ar[d]^a \\
Y \ar[r]^-b & Y \amalg_Z X
}
$$
该图是纤维方块，态射 $a$ 是整的，态射 $b$ 是闭浸入，并且
$$
\mathcal{O}_{Y \amalg_Z X} =
b_*\mathcal{O}_Y \times_{c_*\mathcal{O}_Z} a_*\mathcal{O}_X
$$
作为环层成立，其中 $c = a \circ i = b \circ j$。
\end{proposition}

\begin{proof}
作为拓扑空间，我们把 $Y \amalg_Z X$ 定义为该图在拓扑空间范畴中的推出
（拓扑，第 \ref{topology-section-colimits} 节）。它就是底层集合的推出
（拓扑，引理 \ref{topology-lemma-colimits}）赋予商拓扑。
在 $Y \amalg_Z X$ 上，有环层的映射
$$
b_*\mathcal{O}_Y \longrightarrow c_*\mathcal{O}_Z
\longleftarrow a_*\mathcal{O}_X
$$
并且可以把
$$
\mathcal{O}_{Y \amalg_Z X} =
b_*\mathcal{O}_Y \times_{c_*\mathcal{O}_Z} a_*\mathcal{O}_X
$$
定义为环层范畴中的纤维积。为证明所得对象是概形，需要证明每个点都有
仿射开邻域。对不在 $c$ 的像中的点，这是清楚的，因为 $c$ 的像是闭子集，
其补集作为环化空间同构于
$(Y \setminus j(Z)) \amalg (X \setminus i(Z))$。

\medskip\noindent
$c$ 的像中的一个点对应于 $j$ 的像中的唯一一点 $y \in Y$。由
引理 \ref{more-morphisms-lemma-prepare-pushout-along-closed-immersion-and-integral}，
可以找到仿射开集 $U \subset X$ 和 $V \subset Y$，使得
$y \in V$ 且 $i^{-1}(U) = j^{-1}(V)$。第一段的构造显然与限制到
相容的开子概形相容，所以为证明它给出一个概形，可以假设
$X$、$Y$ 和 $Z$ 都仿射。

\medskip\noindent
若 $X = \Spec(A)$、$Y = \Spec(B)$ 和 $Z = \Spec(C)$ 都仿射，则
代数进阶，引理 \ref{more-algebra-lemma-points-of-fibre-product} 表明，
作为拓扑空间有 $Y \amalg_Z X = \Spec(B \times_C A)$。
% P05-EN-0142: frozen source switches the pushout to the undefined fibre product Y times_Z X.
为完成 $Y \times_Z X$ 是概形的证明，只需说明在
$\Spec(B \times_C A)$ 上，结构层是各直像的纤维积。这可由把
代数进阶，引理 \ref{more-algebra-lemma-diagram-localize}
应用于 $\Spec(B \times_C A)$ 的主仿射开集得到。

\medskip\noindent
上述讨论表明，概形 $Y \amalg_Z X$ 有仿射开覆盖
$Y \amalg_Z X = \bigcup W_i$，使得 $U_i = a^{-1}(W_i)$、
$V_i = b^{-1}(W_i)$ 和 $\Omega_i = c^{-1}(W_i)$ 分别是 $X$、$Y$ 和 $Z$
中的仿射开集。因此 $a$ 和 $b$ 都是仿射的。进一步，若 $A_i$、$B_i$、
$C_i$ 分别是对应于 $U_i$、$V_i$、$\Omega_i$ 的环，则
$A_i \to C_i$ 满射，而 $W_i$ 对应于 $A_i \times_{C_i} B_i$，
后者满射到 $B_i$。故 $b$ 是闭浸入。由代数进阶，Lemma
\ref{more-algebra-lemma-fibre-product-integral}，环映射
$A_i \times_{C_i} B_i \to A_i$ 是整的，故 $a$ 是整的。
该图是笛卡尔的，因为
$$
C_i \cong B_i \otimes_{B_i \times_{C_i} A_i} A_i
$$
这是因为 $B_i \times_{C_i} A_i \to B_i$ 和 $A_i \to C_i$
是具有相同核的满射。

\medskip\noindent
最后，应用 引理 \ref{more-morphisms-lemma-basic-example-pushout} 和
\ref{more-morphisms-lemma-pushout-fpqc-local}，可知我们的构造是概形范畴中的推出。
\end{proof}

\begin{lemma}
\label{more-morphisms-lemma-pushout-separated}
在 Situation \ref{more-morphisms-situation-pushout-along-closed-immersion-and-integral}
中，若 $X$ 和 $Y$ 分离，则推出 $Y \amalg_Z X$
（命题 \ref{more-morphisms-proposition-pushout-along-closed-immersion-and-integral}）
分离。把“分离”替换为“在 $S$ 上分离”、“拟分离”或“在 $S$ 上拟分离”，
结论同样成立。
\end{lemma}

\begin{proof}
态射 $Y \amalg X \to Y \amalg_Z X$ 满射且普遍闭。因此可以应用
态射，引理 \ref{morphisms-lemma-image-universally-closed-separated}。
\end{proof}

\begin{lemma}
\label{more-morphisms-lemma-pushout-finite-type}
在 Situation \ref{more-morphisms-situation-pushout-along-closed-immersion-and-integral}
中，假设 $S$ 是局部诺特概形，而 $X$、$Y$ 和 $Z$ 都在 $S$ 上局部有限型。
则推出 $Y \amalg_Z X$
（命题 \ref{more-morphisms-proposition-pushout-along-closed-immersion-and-integral}）
在 $S$ 上局部有限型。
\end{lemma}

\begin{proof}
在仿射开集上考察，结论化为代数进阶，Lemma
\ref{more-algebra-lemma-fibre-product-finite-type} 的结果。
\end{proof}

\begin{lemma}
\label{more-morphisms-lemma-pushout-functor}
在 Situation \ref{more-morphisms-situation-pushout-along-closed-immersion-and-integral}
中，假设给定交换图
$$
\xymatrix{
Y' \ar[d]^g & Z' \ar[l]^{j'} \ar[r]_{i'} \ar[d]^h & X' \ar[d]^f \\
Y & Z \ar[l] \ar[r] & X
}
$$
其中两个方块均为笛卡尔方块，且 $f, g, h$ 分离并局部拟有限。则
\begin{enumerate}
\item 推出 $Y \amalg_Z X$ 和 $Y' \amalg_{Z'} X'$ 存在；
\item $Y' \amalg_{Z'} X' \to Y \amalg_Z X$ 分离且局部拟有限；
\item 图中的两个方块
$$
\xymatrix{
Y' \ar[r] \ar[d] & Y' \amalg_{Z'} X' \ar[d] & X' \ar[l] \ar[d] \\
Y \ar[r] & Y \amalg_Z X & X \ar[l]
}
$$
均为笛卡尔方块。
\end{enumerate}
\end{lemma}

\begin{proof}
由 命题 \ref{more-morphisms-proposition-pushout-along-closed-immersion-and-integral}，
推出 $Y \amalg_Z X$ 存在。为证明推出 $Y' \amalg_{Z'} X'$ 存在，
只需检验 Situation \ref{more-morphisms-situation-pushout-along-closed-immersion-and-integral}
的条件 (3) 对 $(X', Y', Z', i', j')$ 成立。具体地，令
$y' \in Y'$，并把它的像记作 $y \in Y$。选择仿射开集 $U \subset X$，
使得 $i(j^{-1}(y)) \subset U$。选择包含在 $f^{-1}(U)$ 中的拟紧开集
$U' \subset X'$，使其包含拟紧子集 $i'((j')^{-1}(\{y'\}))$。由
引理 \ref{more-morphisms-lemma-etale-separated-ind-quasi-affine}，$U'$ 是拟仿射的。
由于 $Z'_{y'}$ 是某个域上的整代数的谱，可以应用极限，Lemma
\ref{limits-lemma-ample-profinite-set-in-principal-affine}，从而找到
$U'$ 的一个包含 $i'((j')^{-1}(\{y'\}))$ 的仿射开子概形，正合所需。

\medskip\noindent
验证存在性后，我们检验其余断言。在仿射局部，这恰好是代数进阶，Lemma
\ref{more-algebra-lemma-properties-algebras-over-fibre-product}
中 $B \to D$ 和 $A' \to C'$ 局部拟有限的情形\footnote{确切地说，
$X, Y, Z, Y \amalg_Z X, X', Y', Z', Y' \amalg_{Z'} X'$
分别对应于 $A', B, A, B', C', D, C, D'$。}。特别地，态射
$Y' \amalg_{Z'} X' \to Y \amalg_Z X$ 局部有限型。由代数进阶，Lemma
\ref{more-algebra-lemma-module-over-fibre-product}，图中的方块是笛卡尔的。
由于局部拟有限性可以在纤维上检验
（态射，引理 \ref{morphisms-lemma-quasi-finite-at-point-characterize}），
可知 $Y' \amalg_{Z'} X' \to Y \amalg_Z X$ 局部拟有限。

\medskip\noindent
还需检验 $Y' \amalg_{Z'} X' \to Y \amalg_Z X$ 分离。由
命题 \ref{more-morphisms-proposition-pushout-along-closed-immersion-and-integral}，
$Y' \amalg X' \to Y' \amalg_{Z'} X'$ 普遍闭且满射。又因为态射
$Y' \amalg X' \to Y \amalg_Z X$ 分离（它分解为
$Y' \amalg X' \to Y \amalg X \to Y \amalg_Z X$），结论由
态射，引理 \ref{morphisms-lemma-image-universally-closed-separated} 得出。
\end{proof}

\begin{lemma}
\label{more-morphisms-lemma-pushout-functor-equivalence-flat}
在 Situation \ref{more-morphisms-situation-pushout-along-closed-immersion-and-integral}
中，在推出 $Y \amalg_Z X$ 上平坦、分离且局部拟有限的概形所成的范畴，
等价于 引理 \ref{more-morphisms-lemma-pushout-functor} 中满足 $f, g, h$ 平坦的
$(X', Y', Z', i', j', f, g, h)$ 所成的范畴。把“平坦”替换为
“\'etale”，同样成立。
\end{lemma}

\begin{proof}
若从 引理 \ref{more-morphisms-lemma-pushout-functor} 中的
$(X', Y', Z', i', j', f, g, h)$ 出发，并且 $f, g, h$ 平坦或 \'etale，
则由代数进阶，Lemma
\ref{more-algebra-lemma-properties-algebras-over-fibre-product}，
$Y' \amalg_{Z'} X' \to Y \amalg_Z X$ 平坦或 \'etale。

\medskip\noindent
反过来，设 $W \to Y \amalg_Z X$ 是分离且局部拟有限的态射。置
$X' = W \times_{Y \amalg_Z X} X$、$Y' = W \times_{Y \amalg_Z X} Y$ 和
$Z' = W \times_{Y \amalg_Z X} Z$，并赋予它们显然的态射
$i', j', f, g, h$。形成推出 $Y' \amalg_{Z'} X'$，便得到态射
$$
Y' \amalg_{Z'} X' \longrightarrow W
$$
% P05-EN-0143: frozen source gives the base as Y amalg_X Z; the constructed pushout throughout is Y amalg_Z X.
它由推出的泛性质成为 $Y \amalg_X Z$ 上的概形态射。若不假设
$W \to Y \amalg_Z X$ 平坦，则一般而言该态射不是同构。
（事实上，代数进阶，Lemma
\ref{more-algebra-lemma-module-over-fibre-product-bis}
表明所显示的箭头是闭浸入，但一般不是同构。）然而，如果
% P05-EN-0144: frozen source switches the pushout to the undefined fibre product Y times_Z X in the flatness condition.
$W \to Y \times_Z X$ 平坦，则由代数进阶，Lemma
\ref{more-algebra-lemma-properties-algebras-over-fibre-product}，它是同构。
\end{proof}

\noindent
下面讨论两个态射均为闭浸入时的存在性。

\begin{lemma}
\label{more-morphisms-lemma-pushout-along-closed-immersions}
设 $i : Z \to X$ 和 $j : Z \to Y$ 是概形的闭浸入。则推出
$Y \amalg_Z X$ 在概形范畴中存在。图示为
$$
\xymatrix{
Z \ar[r]_i \ar[d]_j & X \ar[d]^a \\
Y \ar[r]^-b & Y \amalg_Z X
}
$$
该图是纤维方块，态射 $a$ 和 $b$ 都是闭浸入，并且有短正合列
$$
0 \to \mathcal{O}_{Y \amalg_Z X} \to
a_*\mathcal{O}_X \oplus b_*\mathcal{O}_Y \to
c_*\mathcal{O}_Z \to 0
$$
其中 $c = a \circ i = b \circ j$。
\end{lemma}

\begin{proof}
这是 Proposition
\ref{more-morphisms-proposition-pushout-along-closed-immersion-and-integral} 的特殊情形。
注意，Situation \ref{more-morphisms-situation-pushout-along-closed-immersion-and-integral}
中的假设 (3) 直接成立，因为 $j$ 的各纤维都是单点集。最后交换两个箭头的
角色，即可得出 $a$ 和 $b$ 都是闭浸入。
\end{proof}

\begin{lemma}
\label{more-morphisms-lemma-pushout-along-closed-immersions-properties-above}
设 $i : Z \to X$ 和 $j : Z \to Y$ 是概形的闭浸入。设
$f : X' \to X$ 和 $g : Y' \to Y$ 是概形态射，并设
$\varphi : X' \times_{X, i} Z \to Y' \times_{Y, j} Z$ 是 $Z$-概形的同构。
考虑态射
$$
h :
X' \amalg_{X' \times_{X, i} Z, \varphi} Y'
\longrightarrow
X \amalg_Z Y
$$
则有
\begin{enumerate}
\item $h$ 局部有限型，当且仅当 $f$ 和 $g$ 局部有限型；
\item $h$ 平坦，当且仅当 $f$ 和 $g$ 平坦；
\item $h$ 平坦且局部有限表示，当且仅当 $f$ 和 $g$ 平坦且局部有限表示；
\item $h$ 光滑，当且仅当 $f$ 和 $g$ 光滑；
\item $h$ 是 \'etale 的，当且仅当 $f$ 和 $g$ 是 \'etale 的；
% P05-EN-0145: frozen source contains the editorial placeholder "add more here as needed" as a numbered mathematical assertion.
\item 可按需要在此添加更多性质。
\end{enumerate}
\end{lemma}

\begin{proof}
由 引理 \ref{more-morphisms-lemma-pushout-along-closed-immersions}，我们知道这些推出存在，
特别地，得到态射 $h$。因此可以把出现的所有概形都替换为仿射概形。
在这种情形下，本引理的各项断言等价于代数进阶，Lemma
\ref{more-algebra-lemma-properties-algebras-over-fibre-product}
的相应断言。
\end{proof}





\section{相对态射}
\label{more-morphisms-section-relative-morphisms}

\noindent
本节证明一个可表性结果；我们将在基本群，Section
\ref{pione-section-finite-etale} 中用它证明关于一个概形的有限 \'etale
覆盖范畴的结果。本节材料在可表性判据，Section
\ref{criteria-section-relative-morphisms} 中以适当的一般性加以讨论。

\medskip\noindent
设 $S$ 是概形，$Z$ 和 $X$ 是 $S$ 上的概形。给定 $S$ 上的概形 $T$，
可以考虑 $S$ 上的态射 $b : T \times_S Z \to T \times_S X$。图示为
\begin{equation}
\label{more-morphisms-equation-hom}
\vcenter{
\xymatrix{
T \times_S Z \ar[rd] \ar[rr]_b & &
T \times_S X \ar[ld] & Z \ar[rd] & & X \ar[ld] \\
& T \ar[rrr] & & & S
}
}
\end{equation}
当然，也可以把 $b$ 看作态射 $b : T \times_S Z \to X$，使得
$$
\xymatrix{
T \times_S Z \ar[r] \ar[d] \ar@/^1pc/[rrr]_-b &
Z \ar[rd] & & X \ar[ld] \\
T \ar[rr] & & S
}
$$
交换。在这种情形下，可以定义函子
\begin{equation}
\label{more-morphisms-equation-hom-functor}
\mathit{Mor}_S(Z, X) : (\Sch/S)^{opp} \longrightarrow \textit{Sets},
\quad
T \longmapsto \{b\text{ 如上}\}
\end{equation}
下面是一个基本的可表性结果。

\begin{lemma}
\label{more-morphisms-lemma-hom-from-finite-free-into-affine}
设 $Z \to S$ 和 $X \to S$ 是仿射概形的态射。假设
$\Gamma(Z, \mathcal{O}_Z)$ 是有限自由的
$\Gamma(S, \mathcal{O}_S)$-模。则 $\mathit{Mor}_S(Z, X)$
可由 $S$ 上的一个仿射概形表示。
\end{lemma}

\begin{proof}
写成 $S = \Spec(R)$。为 $\Gamma(Z, \mathcal{O}_Z)$ 选择一个 $R$-基
$\{e_1, \ldots, e_m\}$。选择一个表示
$$
\Gamma(X, \mathcal{O}_X) = R[\{x_i\}_{i \in I}]/(\{f_k\}_{k \in K}).
$$
用 $\overline{x}_i$ 表示 $x_i$ 在该商中的像。写成
$$
P = R[\{a_{ij}\}_{i \in I, 1 \leq j \leq m}].
$$
考虑 $R$-代数映射
$$
\Psi :
R[\{x_i\}_{i \in I}]
\longrightarrow
P \otimes_R \Gamma(Z, \mathcal{O}_Z), \quad
x_i \longmapsto \sum\nolimits_j a_{ij} \otimes e_j.
$$
写成 $\Psi(f_k) = \sum c_{kj} \otimes e_j$，其中 $c_{kj} \in P$。
最后，用 $J \subset P$ 表示由元素 $c_{kj}$（$k \in K$，
$1 \leq j \leq m$）生成的理想。我们断言 $W = \Spec(P/J)$
表示函子 $\mathit{Mor}_S(Z, X)$。

\medskip\noindent
首先注意，按构造 $P/J$ 是 $R$-代数，因而给出态射 $W \to S$。
其次，按构造，映射 $\Psi$ 分解经过 $\Gamma(X, \mathcal{O}_X)$，
所以得到 $P/J$-代数同态
$$
P/J \otimes_R \Gamma(X, \mathcal{O}_X)
\longrightarrow
P/J \otimes_R \Gamma(Z, \mathcal{O}_Z)
$$
它确定一个态射 $b_{univ} : W \times_S Z \to W \times_S X$。
由 Yoneda 引理，$b_{univ}$ 确定函子变换
$W \to \mathit{Mor}_S(Z, X)$；我们断言它是同构。为证明这一点，
只需证明对任一仿射概形 $T$，它诱导集合的双射
$W(T) \to \mathit{Mor}_S(Z, X)(T)$。

\medskip\noindent
假设 $T = \Spec(R')$ 是 $S$ 上的仿射概形，并且
$b \in \mathit{Mor}_S(Z, X)(T)$。结构态射 $T \to S$ 在 $R'$ 上定义
$R$-代数结构，而 $b$ 定义一个 $R'$-代数映射
$$
b^\sharp :
R' \otimes_R \Gamma(X, \mathcal{O}_X)
\longrightarrow
R' \otimes_R \Gamma(Z, \mathcal{O}_Z).
$$
特别地，可以写成
$b^\sharp(1 \otimes \overline{x}_i) = \sum \alpha_{ij} \otimes e_j$，
其中 $\alpha_{ij} \in R'$。这对应于由规则
$a_{ij} \mapsto \alpha_{ij}$ 确定的 $R$-代数映射 $P \to R'$。
由理想 $J$ 的构造，该映射分解经过商 $P/J$，给出映射
$P/J \to R'$。这又对应于态射 $T \to W$，使得 $b$ 是 $b_{univ}$
的拉回。略去一些细节。
\end{proof}

\begin{lemma}
\label{more-morphisms-lemma-hom-from-finite-locally-free-into-affine}
设 $Z \to S$ 和 $X \to S$ 是概形态射。若 $Z \to S$ 有限局部自由，
且 $X \to S$ 仿射，则 $\mathit{Mor}_S(Z, X)$ 可由一个在 $S$
上仿射的概形表示。
\end{lemma}

\begin{proof}
选择仿射开覆盖 $S = \bigcup U_i$，使得
$\Gamma(Z \times_S U_i, \mathcal{O}_{Z \times_S U_i})$ 是有限自由的
$\mathcal{O}_S(U_i)$-模。令 $F_i \subset \mathit{Mor}_S(Z, X)$
为如下子函子：若 $T \to S$ 不分解经过 $U_i$，则它把 $T/S$
映到空集；否则映到 $\mathit{Mor}_S(Z, X)(T)$。这些子函子的集合满足
概形，引理 \ref{schemes-lemma-glue-functors} 的条件
(2)(a)、(2)(b)、(2)(c)，从而证明引理。条件 (2)(a) 由
引理 \ref{more-morphisms-lemma-hom-from-finite-free-into-affine} 得出，
其余两个条件由直接的论证得出。
\end{proof}

\noindent
下一个引理中关于态射 $f : X \to S$ 的条件，对证明这类断言非常有用。
在下列任一情形中，该条件成立：$X$ 拟仿射；$f$ 拟仿射；$f$ 拟射影；
$f$ 局部射影；$X$ 上存在充足可逆层；$X$ 上存在 $f$-充足可逆层；
或者 $X$ 上存在 $f$-很充足可逆层。

\begin{lemma}
\label{more-morphisms-lemma-hom-from-finite-locally-free-representable}
设 $Z \to S$ 和 $X \to S$ 是概形态射。假设
\begin{enumerate}
\item $Z \to S$ 有限局部自由；
\item 对所有满足 $s \in S$ 且 $x_1, \ldots, x_d \in X_s$ 的
$(s, x_1, \ldots, x_d)$，存在仿射开集 $U \subset X$，使得
$x_1, \ldots, x_d \in U$。
\end{enumerate}
则 $\mathit{Mor}_S(Z, X)$ 可由概形表示。
\end{lemma}

\begin{proof}
考虑由偶对 $(U, V)$ 组成的集合 $I$，其中 $U \subset X$ 和 $V \subset S$
是仿射开集，且 $U \to S$ 分解经过 $V$。对 $i \in I$，用
$(U_i, V_i)$ 表示相应的偶对。置
$F_i = \mathit{Mor}_{V_i}(Z_{V_i}, U_i)$。显然，$F_i$ 是
$\mathit{Mor}_S(Z, X)$ 的子函子。于是我们断言概形，Lemma
% P05-EN-0146: frozen source omits the predicate "hold" after claiming the three gluing conditions.
\ref{schemes-lemma-glue-functors} 的条件 (2)(a)、(2)(b)、(2)(c) 成立，
这就证明了引理。

\medskip\noindent
条件 (2)(a) 由
引理 \ref{more-morphisms-lemma-hom-from-finite-locally-free-into-affine} 得出。

\medskip\noindent
为检验条件 (2)(b)，考虑 $T/S$ 和 $b \in \mathit{Mor}_S(Z, X)(T)$。
把 $b$ 看作态射 $T \times_S Z \to X$，便得到开集
$b^{-1}(U_i) \subset T \times_S Z$。显然，$b \in F_i(T)$ 当且仅当
$b^{-1}(U_i) = T \times_S Z$。由于投影 $p : T \times_S Z \to T$
有限，因而是闭态射，满足 $p^{-1}(\{t\}) \subset b^{-1}(U_i)$
的点 $t \in T$ 所成的集合 $U_{i, b} \subset T$ 是开集。
于是 $f : T' \to T$ 分解经过 $U_{i, b}$，当且仅当
% P05-EN-0147: frozen source writes b composed with f although b has domain T times_S Z and f has codomain T; the base-changed map is required.
$b \circ f \in F_i(T')$；条件 (2)(b) 得证。

\medskip\noindent
最后检验条件 (2)(c)，此处要使用我们对 $X \to S$ 所加的条件。
具体地，考虑 $T/S$ 和 $b \in \mathit{Mor}_S(Z, X)(T)$。
只需证明每个 $t \in T$ 都包含在上一段定义的某个开集 $U_{i, b}$ 中。
这等价于存在某个 $i$，使得 $b(p^{-1}(\{t\})) \subset U_i$，
其中 $p : T \times_S Z \to T$ 是投影，而 $b : T \times_S Z \to X$
是给定的态射。由于 $p$ 有限，集合 $b(p^{-1}(\{t\})) \subset X$
是有限集，并且包含在 $X \to S$ 的如下纤维中：该纤维位于 $t$
在 $S$ 中的像 $s$ 之上。因此，我们对 $X \to S$ 的条件恰好表明
存在合适的偶对。
\end{proof}

\begin{lemma}
\label{more-morphisms-lemma-hom-from-finite-locally-free-separated-lqf}
设 $Z \to S$ 和 $X \to S$ 是概形态射。假设 $Z \to S$ 有限局部自由，
而 $X \to S$ 分离且局部拟有限。则 $\mathit{Mor}_S(Z, X)$
可由概形表示。
\end{lemma}

\begin{proof}
这由 引理 \ref{more-morphisms-lemma-hom-from-finite-locally-free-representable} 和
\ref{more-morphisms-lemma-separated-locally-quasi-finite-over-affine} 得出。
\end{proof}









\section{伪凝聚复形的刻画，III}
\label{more-morphisms-section-characterize-pseudo-coherent}

\noindent
本节讨论以其上同调刻画伪凝聚复形的方法。这是概形的导出范畴，Section
\ref{perfect-section-pseudo-coherent} 的延续。基本工具是使用 Chow 引理的
一个导出版本，把问题归约到射影空间的情形，见
引理 \ref{more-morphisms-lemma-derived-chow}。

\begin{lemma}
\label{more-morphisms-lemma-case-of-tor-independence}
考虑概形的交换图
$$
\xymatrix{
Z' \ar[d] \ar[r] & Y' \ar[d] \\
X' \ar[r] & S'
}
$$
设 $S \to S'$ 是态射。用 $X$ 和 $Y$ 分别表示 $X'$ 和 $Y'$ 到 $S$
的基变换。假设 $Y' \to S'$ 和 $Z' \to X'$ 平坦。则
$X \times_S Y$ 和 $Z'$ 在 $X' \times_{S'} Y'$ 上 Tor 独立。
\end{lemma}

\begin{proof}
问题是局部的，故可以假设所有概形都仿射（略去一些细节）。注意
$$
\xymatrix{
X \times_S Y \ar[r] \ar[d] & X' \times_{S'} Y' \ar[d] \\
X \ar[r] & X'
}
$$
是竖直箭头平坦的笛卡尔方块。写成 $X = \Spec(A)$、
$X' = \Spec(A')$、$X' \times_{S'} Y' = \Spec(B')$。则
$X \times_S Y = \Spec(A \otimes_{A'} B')$。再写成
$Z' = \Spec(C')$。我们要证明
$$
\text{Tor}_p^{B'}(A \otimes_{A'} B', C') = 0,
\quad\text{对 } p > 0
$$
由于 $A' \to B'$ 平坦，有
$A \otimes_{A'} B' = A \otimes_{A'}^\mathbf{L} B'$。因此
$$
(A \otimes_{A'} B') \otimes_{B'}^\mathbf{L} C' =
(A \otimes_{A'}^\mathbf{L} B') \otimes_{B'}^\mathbf{L} C' =
A \otimes_{A'}^\mathbf{L} C' =
A \otimes_{A'} C'
$$
第二个等号由代数进阶，Lemma
\ref{more-algebra-lemma-double-base-change} 得出；最后一个等号则因为
$A' \to C'$ 平坦。这就证明了引理。
\end{proof}

\begin{lemma}[导出 Chow 引理]
\label{more-morphisms-lemma-derived-chow}
设 $A$ 是环，$X$ 是 $A$ 上分离且有限表示的概形，并设 $x \in X$。
则存在 $x$ 的开邻域 $U \subset X$、整数 $n \geq 0$、
开集 $V \subset \mathbf{P}^n_A$、闭子概形
$Z \subset X \times_A \mathbf{P}^n_A$、点 $z \in Z$ 以及
$D(\mathcal{O}_{X \times_A \mathbf{P}^n_A})$ 中的对象 $E$，使得
\begin{enumerate}
\item $Z \to X \times_A \mathbf{P}^n_A$ 是有限表示的；
\item $b : Z \to X$ 在 $U$ 上是同构，且 $b(z) = x$；
\item $c : Z \to \mathbf{P}^n_A$ 在 $V$ 上是闭浸入；
\item $b^{-1}(U) = c^{-1}(V)$，特别地，$c(z) \in V$；
\item $E|_{X \times_A V} \cong
(b, c)_*\mathcal{O}_Z|_{X \times_A V}$；
\item $E$ 伪凝聚，且支撑在 $Z$ 上。
\end{enumerate}
\end{lemma}

\begin{proof}
可以找到有限型 $\mathbf{Z}$-子代数 $A' \subset A$，以及 $A'$ 上
分离且有限表示的概形 $X'$，使其到 $A$ 的基变换是 $X$。见极限，
引理 \ref{limits-lemma-descend-finite-presentation} 和
\ref{limits-lemma-descend-separated-finite-presentation}。
令 $x' \in X'$ 为 $x$ 的像。若能对 $x' \in X'/A'$ 证明本引理，
则对 $x \in X/A$ 的引理随即成立。具体地，若
$U', n', V', Z', z', E'$ 给出了 $x' \in X'/A'$ 的解，则可以令
$U \subset X$ 为 $U'$ 的原像，令 $n = n'$，令
$V \subset \mathbf{P}^n_A$ 为 $V'$ 的原像，令
$Z \subset X \times \mathbf{P}^n$ 为 $Z'$ 的概形论原像，
令 $z \in Z$ 为映到 $x$ 的唯一一点，并令 $E$ 为 $E'$ 的导出拉回。
由上同调，引理 \ref{cohomology-lemma-pseudo-coherent-pullback}，
$E$ 是伪凝聚的。只需再检验 (5)。为此，置
$W = b^{-1}(U) = c^{-1}(V)$ 和
% P05-EN-0148: frozen source uses U instead of U' in the definition of W' over X'.
$W' = (b')^{-1}(U) = (c')^{-1}(V')$，并考虑笛卡尔方块
$$
\xymatrix{
W \ar[d]_{(b, c)} \ar[r] & W' \ar[d]^{(b', c')} \\
X \times_A V \ar[r] & X' \times_{A'} V'
}
$$
由 引理 \ref{more-morphisms-lemma-case-of-tor-independence}，概形 $X \times_A V$ 和
$W'$ 在 $X' \times_{A'} V'$ 上 Tor 独立。因此由概形的导出范畴，
引理 \ref{perfect-lemma-compare-base-change}，
$(b', c')_*\mathcal{O}_{W'}$ 到 $X \times_A V$ 的导出拉回是
$(b, c)_*\mathcal{O}_W$。这里还用到
$R(b', c')_*\mathcal{O}_{Z'} = (b', c')_*\mathcal{O}_{Z'}$，
因为 $(b', c')$ 是闭浸入；对 $(b, c)_*\mathcal{O}_Z$ 也类似。
由于 $E'|_{U' \times_{A'} V'} =
(b', c')_*\mathcal{O}_{W'}$，我们得到
$E|_{U \times_A V} = (b, c)_*\mathcal{O}_W$，故 (5) 成立。
这把我们归约到下一段所描述的情形。

\medskip\noindent
假设 $A$ 在 $\mathbf{Z}$ 上有限型。选择 $x$ 的仿射开邻域
$U \subset X$。则 $U$ 在 $A$ 上有限型。选择闭浸入
$U \to \mathbf{A}^n_A$，并用 $j : U \to \mathbf{P}^n_A$ 表示它与开浸入
$\mathbf{A}^n_A \to \mathbf{P}^n_A$ 复合所得的浸入。令 $Z$ 为
$$
(\text{id}_U, j) : U \longrightarrow X \times_A \mathbf{P}^n_A
$$
的概形论闭包。由于投影 $X \times \mathbf{P}^n \to X$ 分离，
由态射，Lemma
\ref{morphisms-lemma-scheme-theoretic-image-of-partial-section} 可知
$b : Z \to X$ 在 $U$ 上是同构。令 $z \in Z$ 为位于 $x$ 上的唯一一点。

\medskip\noindent
令 $Y \subset \mathbf{P}^n_A$ 为 $j$ 的概形论闭包。于是显然，
$Z \subset X \times_A Y$ 是
$(\text{id}_U, j) : U \to X \times_A Y$ 的概形论闭包。
由于 $X$ 分离，态射 $X \times_A Y \to Y$ 也分离。因此由上面使用的
同一引理，$Z \to Y$ 在开子概形 $j(U) \subset Y$ 上是同构。
选择开集 $V \subset \mathbf{P}^n_A$，使得 $V \cap Y = j(U)$。
于是 (3) 和 (4) 成立。

\medskip\noindent
由于 $A$ 是诺特环，$X$ 和 $X \times_A \mathbf{P}^n_A$ 都是诺特概形。
因此在这种情形下可以取 $E = (b, c)_*\mathcal{O}_Z$，见概形的导出范畴，
引理 \ref{perfect-lemma-identify-pseudo-coherent-noetherian}。
证明完成。
\end{proof}

\begin{lemma}
\label{more-morphisms-lemma-compute-Fourier-Mukai-for-derived-chow}
设 $A$、$x \in X$ 以及 $U, n, V, Z, z, E$ 如
引理 \ref{more-morphisms-lemma-derived-chow} 所述。对任意
$K \in D_\QCoh(\mathcal{O}_X)$，有
$$
Rq_*(Lp^*K \otimes^\mathbf{L} E)|_V = R(U \to V)_*K|_U
$$
其中 $p : X \times_A \mathbf{P}^n_A \to X$ 和
$q : X \times_A \mathbf{P}^n_A \to \mathbf{P}^n_A$ 是投影，
而态射 $U \to V$ 是有限表示的闭浸入 $c \circ (b|_U)^{-1}$。
\end{lemma}

\begin{proof}
由于 $b^{-1}(U) = c^{-1}(V)$，且 $c$ 在 $V$ 上是闭浸入，
$c \circ (b|_U)^{-1}$ 是闭浸入。它是有限表示的，因为 $U$ 和 $V$
都在 $A$ 上有限表示，见态射，Lemma
\ref{morphisms-lemma-finite-presentation-permanence}。首先有
$$
Rq_*(Lp^*K \otimes^\mathbf{L} E)|_V =
Rq'_*\left((Lp^*K \otimes^\mathbf{L} E)|_{X \times_A V}\right)
$$
其中 $q' : X \times_A V \to V$ 是投影；这里使用了总直像的形成与局部化相容。
置 $W = b^{-1}(U) = c^{-1}(V)$，并用 $i : W \to X \times_A V$
表示闭浸入 $i = (b, c)|_W$。则
$$
Rq'_*\left((Lp^*K \otimes^\mathbf{L} E)|_{X \times_A V}\right) =
Rq'_*(Lp^*K|_{X \times_A V} \otimes^\mathbf{L} i_*\mathcal{O}_W)
$$
由性质 (5) 得到该等式。由于 $i$ 是闭浸入，有
$i_*\mathcal{O}_W = Ri_*\mathcal{O}_W$。使用概形的导出范畴，Lemma
\ref{perfect-lemma-cohomology-base-change}，可以把它改写为
$$
Rq'_* Ri_* Li^* Lp^*K|_{X \times_A V} =
R(q' \circ i)_* Lb^*K|_W =
R(U \to V)_* K|_U
$$
这正是所需结论。
\end{proof}

\begin{lemma}
\label{more-morphisms-lemma-characterize-pseudo-coherent}
设 $A$ 是环，$X$ 是 $A$ 上分离且有限表示的概形，并设
$K \in D_\QCoh(\mathcal{O}_X)$。若对 $D(\mathcal{O}_X)$ 中的每个
伪凝聚对象 $E$，$R\Gamma(X, E \otimes^\mathbf{L} K)$ 在 $D(A)$
中都是伪凝聚的，则 $K$ 相对于 $A$ 伪凝聚。
\end{lemma}

\begin{proof}
假设 $K \in D_\QCoh(\mathcal{O}_X)$，并且对 $D(\mathcal{O}_X)$
中的每个伪凝聚对象 $E$，$R\Gamma(X, E \otimes^\mathbf{L} K)$
在 $D(A)$ 中都是伪凝聚的。令 $x \in X$。我们将证明 $K$
在 $x$ 的某个邻域中相对于 $A$ 伪凝聚；这便证明了引理。

\medskip\noindent
选择 引理 \ref{more-morphisms-lemma-derived-chow} 中的 $U, n, V, Z, z, E$。
用 $p : X \times \mathbf{P}^n \to X$ 和
$q : X \times \mathbf{P}^n \to \mathbf{P}^n_A$ 表示投影。
则对任意 $i \in \mathbf{Z}$，有
\begin{align*}
& R\Gamma(\mathbf{P}^n_A,
Rq_*(Lp^*K \otimes^\mathbf{L} E)
\otimes^\mathbf{L}
\mathcal{O}_{\mathbf{P}^n_A}(i)) \\
& =
R\Gamma(X \times \mathbf{P}^n,
Lp^*K \otimes^\mathbf{L} E \otimes^\mathbf{L}
Lq^*\mathcal{O}_{\mathbf{P}^n_A}(i)) \\
& =
R\Gamma(X,
K \otimes^\mathbf{L} Rp_*(E \otimes^\mathbf{L}
Lq^*\mathcal{O}_{\mathbf{P}^n_A}(i)))
\end{align*}
这是由概形的导出范畴，Lemma
\ref{perfect-lemma-cohomology-base-change} 得到的。由概形的导出范畴，
引理 \ref{perfect-lemma-flat-proper-pseudo-coherent-direct-image-general}，
复形 $Rp_*(E \otimes^\mathbf{L} Lq^*\mathcal{O}_{\mathbf{P}^n_A}(i))$
在 $X$ 上伪凝聚。因此假设表明，所显示公式中的表达式是 $D(A)$
中的伪凝聚对象。由概形的导出范畴，Lemma
\ref{perfect-lemma-pseudo-coherent-on-projective-space}，可知
$Rq_*(Lp^*K \otimes^\mathbf{L} E)$ 在 $\mathbf{P}^n_A$ 上伪凝聚。
由 引理 \ref{more-morphisms-lemma-compute-Fourier-Mukai-for-derived-chow}，有
$$
% P05-EN-0149: frozen source restricts a complex on projective space to X times_A V; Lemma compute-Fourier-Mukai states the restriction is to V.
Rq_*(Lp^*K \otimes^\mathbf{L} E)|_{X \times_A V} =
R(U \to V)_*K|_U
$$
由于 $U \to V$ 是到 $\mathbf{P}^n_A$ 的一个开子概形中的闭浸入，
由 引理 \ref{more-morphisms-lemma-check-relative-pseudo-coherence-on-charts}，
这表明 $K|_U$ 相对于 $A$ 伪凝聚。
\end{proof}

\begin{lemma}
\label{more-morphisms-lemma-characterize-pseudo-coh-improved}
设 $A$ 是环，$X$ 是 $A$ 上分离且有限表示的概形。设
$K \in D_\QCoh(\mathcal{O}_X).$ 若对每个完美对象
$E \in D(\mathcal{O}_X)$，$R \Gamma (X, E \otimes ^{\mathbf{L}} K)$
在 $D(A)$ 中都伪凝聚，则 $K$ 相对于 $A$ 伪凝聚。
\end{lemma}

\begin{proof}
鉴于 引理 \ref{more-morphisms-lemma-characterize-pseudo-coherent}，只需证明对每个
伪凝聚对象 $E \in D(\mathcal{O}_X)$，
$R \Gamma (X, E \otimes ^{\mathbf{L}} K)$ 在 $D(A)$ 中都伪凝聚。
由概形的导出范畴，Proposition
\ref{perfect-proposition-detecting-bounded-above}，可知
$K \in D^-_\QCoh (\mathcal{O}_X)$。现在结论由概形的导出范畴，Lemma
\ref{perfect-lemma-perfect-enough} 得出。
\end{proof}

\begin{lemma}
\label{more-morphisms-lemma-characterize-relatively-perfect}
设 $A$ 是环，$X$ 是 $A$ 上分离、有限表示且平坦的概形。设
$K \in D_\QCoh(\mathcal{O}_X).$ 若对每个完美对象
$E \in D(\mathcal{O}_X)$，$R \Gamma (X, E \otimes^\mathbf{L} K)$
在 $D(A)$ 中都完美，则 $K$ 是 $\Spec(A)$-完美的。
\end{lemma}

\begin{proof}
由 引理 \ref{more-morphisms-lemma-characterize-pseudo-coh-improved}，$K$ 相对于 $A$
伪凝聚。由 引理 \ref{more-morphisms-lemma-check-relative-pseudo-coherence-on-charts}，
$K$ 在 $D( \mathcal{O}_X)$ 中伪凝聚。由概形的导出范畴，Proposition
\ref{perfect-proposition-detecting-bounded-below}，可知 $K$ 属于
$D^-(\mathcal{O}_X)$。设 $\mathfrak{p}$ 是 $A$ 的素理想，并用
$i : Y \to X$ 表示 $\mathfrak{p}$ 上的概形论纤维的嵌入，即 $Y$
是 $\kappa(\mathfrak p)$ 上的概形。由概形的导出范畴，Lemma
\ref{perfect-lemma-bounded-on-fibres}，只要证明 $Li^*(K)$ 下有界即可。
令 $G \in D_{perf} (\mathcal{O}_X)$ 为生成 $D_\QCoh (\mathcal{O}_X)$
的完美复形，见概形的导出范畴，Theorem
\ref{perfect-theorem-bondal-van-den-Bergh}。我们有
\begin{align*}
R\Hom _{\mathcal{O}_Y}(Li^*(G), Li^*(K))
& =
R\Gamma(Y, Li^*(G ^\vee \otimes ^\mathbf{L} K)) \\
& =
R\Gamma(X, G^\vee \otimes ^{\mathbf{L}} K)
\otimes^\mathbf{L}_A \kappa(\mathfrak{p})
\end{align*}
第一个等号使用了 $Li^*$ 保持完美对象及其对偶这一事实，以及上同调，
引理 \ref{cohomology-lemma-dual-perfect-complex}；略去一些细节。
由于 $X$ 在 $A$ 上平坦，第二个等号由概形的导出范畴，Lemma
\ref{perfect-lemma-compare-base-change} 得出。由假设，它是
$D(\kappa(\mathfrak{p}))$ 中的完美对象。由概形的导出范畴，Remark
\ref{perfect-remark-pullback-generator}，对象
$Li^*(G) \in D_{perf}(\mathcal{O}_Y)$ 生成 $D_\QCoh(\mathcal{O}_Y)$。
因此概形的导出范畴，Proposition
\ref{perfect-proposition-detecting-bounded-below} 现在表明 $Li^*(K)$
下有界，证明完成。
\end{proof}




\section{复形有限性性质的下降}
\label{more-morphisms-section-descent-finiteness}

\noindent
本节是概形的导出范畴，Section
\ref{perfect-section-descent-finiteness} 的延续。

\begin{lemma}
\label{more-morphisms-lemma-relative-pseudo-coherent-descends-fppf}
设 $X \to S$ 局部有限型，$\{f_i : X_i \to X\}$ 是概形的 fppf 覆盖。
设 $E \in D_\QCoh(\mathcal{O}_X)$，并设 $m \in \mathbf{Z}$。
则 $E$ 相对于 $S$ 是 $m$-伪凝聚的，当且仅当每个 $Lf_i^*E$
相对于 $S$ 都是 $m$-伪凝聚的。
\end{lemma}

\begin{proof}
假设 $E$ 相对于 $S$ 是 $m$-伪凝聚的。由
引理 \ref{more-morphisms-lemma-flat-finite-presentation-pseudo-coherent}，态射 $f_i$
是伪凝聚的。因此由 引理 \ref{more-morphisms-lemma-pull-relative-pseudo-coherent}，
$Lf_i^*E$ 相对于 $S$ 是 $m$-伪凝聚的。

\medskip\noindent
反过来，假设对每个 $i$，$Lf_i^*E$ 相对于 $S$ 都是 $m$-伪凝聚的。
按照 引理 \ref{more-morphisms-lemma-dominate-fppf} 选择 $S = \bigcup U_j$、
$W_j \to U_j$、$W_j = \bigcup W_{j, k}$、$T_{j, k} \to W_{j, k}$，
以及 $S$ 上的态射
$\alpha_{j, k} : T_{j, k} \to X_{i(j, k)}$。由于态射
% P05-EN-0150: frozen source switches the lower index k to an uppercase K in T_{j,K}.
$T_{j, K} \to S$ 平坦且有限表示，由
引理 \ref{more-morphisms-lemma-permanence-pseudo-coherent} 可知 $\alpha_{j, k}$
是伪凝聚的。因此
$$
% P05-EN-0151: frozen source switches T_{j,k} to T_{i,k} on the right side of the pullback identity.
L\alpha_{j, k}^*Lf_{i(j, k)}^*E = L(T_{i, k} \to S)^*E
$$
由 引理 \ref{more-morphisms-lemma-pull-relative-pseudo-coherent}，它相对于 $S$
是 $m$-伪凝聚的。现在要沿覆盖
$\{T_{j, k} \to W_{j, k}\}$、$W_j = \bigcup W_{j, k}$、
$\{W_j \to U_j\}$ 和 $S = \bigcup U_j$ 下降这一性质。
对 Zariski 覆盖，结论由相对于 $S$ 的 $m$-伪凝聚性的定义成立，
故可以假设有单个满射有限局部自由态射 $\pi : Y \to X$，使得
% P05-EN-0152: frozen source drops the qualifier m- at the decisive descent step and in the following direct-image claim.
$L\pi^*E$ 相对于 $S$ 伪凝聚。此时由
引理 \ref{more-morphisms-lemma-finite-morphism-relative-pseudo-coherence}，
$R\pi_*L\pi^*E$ 相对于 $S$ 伪凝聚（这是我们第一次使用 $E$
具有拟凝聚上同调层）。我们有
$R\pi_*L\pi^*E = E \otimes^\mathbf{L}_{\mathcal{O}_X} \pi_*\mathcal{O}_Y$
例如由概形的导出范畴，Lemma
\ref{perfect-lemma-cohomology-base-change} 得到该等式；并且在 $X$
上局部地，映射 $\mathcal{O}_X \to \pi_*\mathcal{O}_Y$ 是一个直和项的嵌入。
因此结论由 引理 \ref{more-morphisms-lemma-summands-relative-pseudo-coherent} 得出。
\end{proof}

\begin{lemma}
\label{more-morphisms-lemma-relative-pseudo-coherent-post-compose}
设 $X \to T \to S$ 是概形态射。假设 $T \to S$ 平坦且局部有限表示，
而 $X \to T$ 局部有限型。设 $E \in D(\mathcal{O}_X)$，并设
$m \in \mathbf{Z}$。则 $E$ 相对于 $S$ 是 $m$-伪凝聚的，当且仅当
$E$ 相对于 $T$ 是 $m$-伪凝聚的。
\end{lemma}

\begin{proof}
在 $X$ 上局部地，可以选择闭浸入 $i : X \to \mathbf{A}^n_T$。
于是 $\mathbf{A}^n_T \to S$ 平坦且局部有限表示。因此应用
引理 \ref{more-morphisms-lemma-composition-relative-pseudo-coherent} 即得所述等价性。
\end{proof}

\begin{lemma}
\label{more-morphisms-lemma-relative-pseudo-coherent-descends-fppf-base}
设 $f : X \to S$ 局部有限型，$\{S_i \to S\}$ 是概形的 fppf 覆盖。
用 $f_i : X_i \to S_i$ 表示 $f$ 的基变换，并用 $g_i : X_i \to X$
表示投影。设 $E \in D_\QCoh(\mathcal{O}_X)$，并设 $m \in \mathbf{Z}$。
则 $E$ 相对于 $S$ 是 $m$-伪凝聚的，当且仅当每个 $Lg_i^*E$
相对于 $S_i$ 都是 $m$-伪凝聚的。
\end{lemma}

\begin{proof}
这形式地由 引理 \ref{more-morphisms-lemma-relative-pseudo-coherent-descends-fppf} 和
\ref{more-morphisms-lemma-relative-pseudo-coherent-post-compose} 得出。具体地，若
$E$ 相对于 $S$ 是 $m$-伪凝聚的，则 $Lg_i^*E$ 相对于 $S$
是 $m$-伪凝聚的（由第一个引理），因而 $Lg_i^*E$ 相对于 $S_i$ 是
$m$-伪凝聚的（由第二个引理）。反过来，若 $Lg_i^*E$ 相对于 $S_i$
是 $m$-伪凝聚的，则 $Lg_i^*E$ 相对于 $S$ 是 $m$-伪凝聚的（由第二个引理），
因而 $E$ 相对于 $S$ 是 $m$-伪凝聚的（由第一个引理）。
\end{proof}






\section{相对完美对象}
\label{more-morphisms-section-relatively-perfect}

\noindent
本节是概形的导出范畴，Section
\ref{perfect-section-relatively-perfect} 中讨论的延续。

\begin{lemma}
\label{more-morphisms-lemma-thickening-pseudo-coherent}
设 $i : X \to X'$ 是概形的有限阶增厚。设
$K' \in D(\mathcal{O}_{X'})$ 是对象，并且 $K = Li^*K'$ 伪凝聚。
则 $K'$ 伪凝聚。
\end{lemma}

\begin{proof}
先证明 $K'$ 具有拟凝聚上同调层。为此，可以归约到一阶增厚的情形，见
第 \ref{more-morphisms-section-thickenings} 节。令
$\mathcal{I} \subset \mathcal{O}_{X'}$ 为定义 $X$ 的拟凝聚理想层。
把短正合列
$$
0 \to \mathcal{I} \to \mathcal{O}_{X'} \to i_*\mathcal{O}_X \to 0
$$
与 $K'$ 作张量积，得到优三角
$$
K' \otimes_{\mathcal{O}_{X'}}^\mathbf{L} \mathcal{I}
\to K' \to
K' \otimes_{\mathcal{O}_{X'}}^\mathbf{L} i_*\mathcal{O}_X
\to
(K' \otimes_{\mathcal{O}_{X'}}^\mathbf{L} \mathcal{I})[1]
$$
由于 $i_* = Ri_*$，并且可以把 $\mathcal{I}$ 看成拟凝聚
$\mathcal{O}_X$-模（因为这里是一阶增厚），可以把它改写为
$$
i_*(K \otimes_{\mathcal{O}_X}^\mathbf{L} \mathcal{I})
\to K' \to
i_*K \to
i_*(K \otimes_{\mathcal{O}_X}^\mathbf{L} \mathcal{I})[1]
$$
使用上同调，Lemma
\ref{cohomology-lemma-projection-formula-closed-immersion} 来识别各项。
由于 $K$ 属于 $D_\QCoh(\mathcal{O}_X)$，可知 $K'$ 属于
$D_\QCoh(\mathcal{O}_{X'})$；这里使用了概形的导出范畴，Lemmas
\ref{perfect-lemma-pseudo-coherent}、
\ref{perfect-lemma-quasi-coherence-tensor-product} 和
\ref{perfect-lemma-quasi-coherence-direct-image}。

\medskip\noindent
假设 $K'$ 属于 $D_\QCoh(\mathcal{O}_{X'})$。问题在 $X'$ 上是局部的，
故可以假设 $X'$ 仿射。设 $X' = \Spec(A')$、$X = \Spec(A)$，其中
$A = A'/I$ 且 $I$ 幂零。于是 $K'$ 来自某个对象 $M' \in D(A')$，
见概形的导出范畴，引理 \ref{perfect-lemma-affine-compare-bounded}。
因此由概形的导出范畴，Lemma
\ref{perfect-lemma-pseudo-coherent-affine} 以及对 $K$ 的假设，
$M = M' \otimes_{A'}^\mathbf{L} A$ 是 $D(A)$ 中的伪凝聚对象。
进而由代数进阶，Lemma
\ref{more-algebra-lemma-check-pseudo-coherent-modulo-nilpotent}，
$M'$ 伪凝聚。证明完成。
\end{proof}

\begin{lemma}
\label{more-morphisms-lemma-thickening-relatively-perfect}
考虑概形的笛卡尔图
$$
\xymatrix{
X \ar[r]_i \ar[d]_f & X' \ar[d]^{f'} \\
Y \ar[r]^j & Y'
}
$$
假设 $X' \to Y'$ 平坦且局部有限表示，而 $Y \to Y'$ 是有限阶增厚。
设 $E' \in D(\mathcal{O}_{X'})$。若 $E = Li^*(E')$ 是 $Y$-完美的，
则 $E'$ 是 $Y'$-完美的。
\end{lemma}

\begin{proof}
回忆，$E$ 是 $Y$-完美的，意味着 $E$ 伪凝聚，并且作为
$f^{-1}\mathcal{O}_Y$-模复形局部具有有限 Tor 维数
（概形的导出范畴，Definition
\ref{perfect-definition-relatively-perfect}）。由
引理 \ref{more-morphisms-lemma-thickening-pseudo-coherent}，$E'$ 伪凝聚。特别地，
$E'$ 属于 $D_\QCoh(\mathcal{O}_{X'})$，见概形的导出范畴，Lemma
\ref{perfect-lemma-pseudo-coherent}。为证明 $E'$ 局部具有有限 Tor 维数，
可以在 $X'$ 上局部地工作。因此可以假设
% P05-EN-0153: frozen source switches the diagram's bases Y',Y to undefined S',S in this affine reduction.
$X'$、$S'$、$X$、$S$ 都仿射，设它们分别由环 $A'$、$R'$、$A$、$R$
给出。于是由概形的导出范畴，Lemma
\ref{perfect-lemma-affine-locally-rel-perfect}，问题归约到交换代数版本；
该版本见代数进阶，Lemma
\ref{more-algebra-lemma-thickening-relatively-perfect}。
\end{proof}

\begin{lemma}
\label{more-morphisms-lemma-henselian-relatively-perfect}
设 $(R, I)$ 是由环和包含在 Jacobson 根中的理想 $I$ 组成的二元组。
置 $S = \Spec(R)$ 和 $S_0 = \Spec(R/I)$。设 $f : X \to S$
固有、平坦且有限表示。记 $X_0 = S_0 \times_S X$。设
$E \in D(\mathcal{O}_X)$ 伪凝聚。若 $E$ 到 $X_0$ 的导出限制 $E_0$
是 $S_0$-完美的，则 $E$ 是 $S$-完美的。
\end{lemma}

\begin{proof}
选择有限仿射开覆盖 $X = U_1 \cup \ldots \cup U_n$。对每个 $i$，
可以选择闭浸入 $U_i \to \mathbf{A}^{d_i}_S$。置
$U_{i, 0} = S_0 \times_S U_i$。对每个 $i$，复形
$E_0|_{U_{i, 0}}$ 的 Tor 振幅包含在某个 $[a_i, b_i]$ 中，其中
$a_i, b_i \in \mathbf{Z}$。令 $x \in X$。我们将证明对某个 $i$，
$E_x$ 在 $R$ 上的 Tor 振幅包含在 $[a_i - d_i, b_i]$ 中。
这将完成证明，因为由上同调，Lemma
\ref{cohomology-lemma-tor-amplitude-stalk}，Tor 振幅可以从茎读出。

\medskip\noindent
由于 $f$ 固有，$f(\overline{\{x\}})$ 是 $S$ 的闭子集。由于 $I$
包含在 Jacobson 根中，$f(\overline{\{x\}})$ 与闭子集
% P05-EN-0154: frozen source has the malformed finite verb "meeting" where the proof requires "meets".
$S_0 \subset S$ 相交。因此存在特化 $x \leadsto x_0$，其中
$x_0 \in X_0$。选择满足 $x_0 \in U_i$ 的 $i$，于是
$x_0 \in U_{i, 0}$。在余下的证明中固定该 $i$。
写成 $U_i = \Spec(A)$。则 $A$ 是平坦有限表示的 $R$-代数，
并且是一个含 $d_i$ 个变量的多项式 $R$-代数的商。限制
$E|_{U_i}$ 对应于 $D(A)$ 中的某个伪凝聚对象 $K$
（由概形的导出范畴，Lemmas
\ref{perfect-lemma-affine-compare-bounded} 和
\ref{perfect-lemma-pseudo-coherent-affine}）。注意，$E_0$ 对应于
$K \otimes_A^\mathbf{L} A/IA$。令
$\mathfrak q \subset \mathfrak q_0 \subset A$ 为对应于
$x \leadsto x_0$ 的素理想。则 $E_x = K_{\mathfrak q}$，而
$K_{\mathfrak q}$ 是 $K_{\mathfrak q_0}$ 的局部化。因此只需证明
$K_{\mathfrak q_0}$ 作为 $R$-模复形的 Tor 振幅包含在
$[a_i - d_i, b_i]$ 中。令 $I \subset \mathfrak p_0 \subset R$
为对应于 $f(x_0)$ 的素理想。于是有
\begin{align*}
K \otimes_R^\mathbf{L} \kappa(\mathfrak p_0)
& =
(K \otimes_R^\mathbf{L} R/I) \otimes_{R/I}^\mathbf{L}
\kappa(\mathfrak p_0) \\
& =
(K \otimes_A^\mathbf{L} A/IA) \otimes_{R/I}^\mathbf{L} \kappa(\mathfrak p_0)
\end{align*}
第二个等号成立是因为 $R \to A$ 平坦。由对 $a_i, b_i$ 的选择，
该复形仅在区间 $[a_i, b_i]$ 中的次数上有上同调。因此最后把
代数进阶，Lemma
\ref{more-algebra-lemma-lift-from-fibre-relatively-perfect}
应用于 $R \to A$、$\mathfrak q_0$、$\mathfrak p_0$ 和 $K$ 即得结论。
\end{proof}










\section{收缩有理曲线}
\label{more-morphisms-section-contracting}

\noindent
本节研究固有态射 $f : X \to Y$，其纤维维数 $\leq 1$，并且满足
$R^1f_*\mathcal{O}_X = 0$。为理解本节题名，请参阅代数曲线，Sections
\ref{curves-section-contracting-rational-tails}、
\ref{curves-section-contracting-rational-bridges} 和
\ref{curves-section-contracting-to-stable}。

\begin{lemma}
\label{more-morphisms-lemma-check-h1-fibre-zero}
设 $f : X \to Y$ 是概形的固有态射，$y \in Y$ 是满足
$\dim(X_y) \leq 1$ 的点。若
\begin{enumerate}
\item $R^1f_*\mathcal{O}_X = 0$；或者更一般地，
\item 存在态射 $g : Y' \to Y$，使得 $y$ 位于 $g$ 的像中，并且
$R^1f'_*\mathcal{O}_{X'} = 0$，其中 $f' : X' \to Y'$ 是 $f$
沿 $g$ 的基变换，
\end{enumerate}
则 $H^1(X_y, \mathcal{O}_{X_y}) = 0$。
\end{lemma}

\begin{proof}
为证明引理，可以用 $y$ 的开邻域代替 $Y$。因此可以假设 $Y$ 仿射，
并且 $f$ 的所有纤维维数均 $\leq 1$，见态射，Lemma
\ref{morphisms-lemma-openness-bounded-dimension-fibres}。在这种情形下，
$R^1f_*\mathcal{O}_X$ 是有限型拟凝聚 $\mathcal{O}_Y$-模，且其形成与
任意基变换相容，见极限，Lemmas
\ref{limits-lemma-proper-top-cohomology-finite-type} 和
\ref{limits-lemma-higher-direct-images-zero-above-dimension-fibre}。
引理立即得证。
\end{proof}

\begin{lemma}
\label{more-morphisms-lemma-h1-fibre-zero}
设 $f : X \to Y$ 是概形的固有态射，$y \in Y$ 是满足
$\dim(X_y) \leq 1$ 且 $H^1(X_y, \mathcal{O}_{X_y}) = 0$ 的点。
则存在 $y$ 的开邻域 $V \subset Y$，使得
$R^1f_*\mathcal{O}_X|_V = 0$，并且沿任意 $Y' \to V$ 基变换后仍成立。
\end{lemma}

\begin{proof}
为证明引理，可以用 $y$ 的开邻域代替 $Y$。因此可以假设 $Y$ 仿射，
并且 $f$ 的所有纤维维数均 $\leq 1$，见态射，Lemma
\ref{morphisms-lemma-openness-bounded-dimension-fibres}。在这种情形下，
$R^1f_*\mathcal{O}_X$ 是有限型拟凝聚 $\mathcal{O}_Y$-模，且其形成与
任意基变换相容，见极限，Lemmas
\ref{limits-lemma-proper-top-cohomology-finite-type} 和
\ref{limits-lemma-higher-direct-images-zero-above-dimension-fibre}。
设 $Y = \Spec(A)$，$y$ 对应于素理想 $\mathfrak p \subset A$，而
$R^1f_*\mathcal{O}_X$ 对应于有限 $A$-模 $M$。由上述基变换断言，
$H^1(X_y, \mathcal{O}_{X_y}) = 0$ 意味着
$\mathfrak pM_\mathfrak p = M_\mathfrak p$。由 Nakayama 引理，
$M_\mathfrak p = 0$。由于 $M$ 有限，可以找到 $f \in A$、
$f \not \in \mathfrak p$，使得 $M_f = 0$。因此取主开集
$V$ 为 $D(f)$，便得到所需结论。
\end{proof}

\begin{lemma}
\label{more-morphisms-lemma-globally-generated-vanishing}
设 $f : X \to Y$ 是概形的固有态射，并且对所有 $y \in Y$ 均有
$\dim(X_y) \leq 1$ 和 $H^1(X_y, \mathcal{O}_{X_y}) = 0$。
设 $\mathcal{F}$ 是 $X$ 上的拟凝聚层。则
\begin{enumerate}
\item 对 $p > 1$，有 $R^pf_*\mathcal{F} = 0$；
\item 若存在满射 $f^*\mathcal{G} \to \mathcal{F}$，其中
$\mathcal{G}$ 是 $Y$ 上的拟凝聚层，则 $R^1f_*\mathcal{F} = 0$。
\end{enumerate}
若 $Y$ 仿射，则还有
\begin{enumerate}
\item[(3)] 对 $p \not \in \{0, 1\}$，有 $H^p(X, \mathcal{F}) = 0$；
\item[(4)] 若 $\mathcal{F}$ 全局生成，则 $H^1(X, \mathcal{F}) = 0$。
\end{enumerate}
\end{lemma}

\begin{proof}
(1) 中的消失就是极限，Lemma
\ref{limits-lemma-higher-direct-images-zero-above-dimension-fibre}。
为证明 (2)，可以在 $Y$ 上局部地工作，并假设 $Y$ 仿射。此时
$R^1f_*\mathcal{F}$ 是与模 $H^1(X, \mathcal{F})$ 相伴的 $Y$ 上拟凝聚模。
这里使用了 $Y$ 仿射、高阶直像的拟凝聚性（概形的上同调，Lemma
\ref{coherent-lemma-quasi-coherence-higher-direct-images}），以及概形的上同调，
引理 \ref{coherent-lemma-quasi-coherence-higher-direct-images-application}。
由于 $Y$ 仿射，拟凝聚模 $\mathcal{G}$ 全局生成，因而
$f^*\mathcal{G}$ 和 $\mathcal{F}$ 也全局生成。这样便看出 (4) 蕴含 (2)。
(3) 由 (1) 以及刚才关于直像拟凝聚性的说明得出。因此只需证明 (4)。
若 $\mathcal{F}$ 全局生成，则存在满射
$\bigoplus_{i \in I} \mathcal{O}_X \to \mathcal{F}$。由 (1) 和上同调长正合列，
它诱导 $H^1$ 上的满射。由 引理 \ref{more-morphisms-lemma-h1-fibre-zero}，
$R^1f_*\mathcal{O}_X = 0$，故 $H^1(X, \mathcal{O}_X) = 0$；并且
$H^1(X, -)$ 与直和相容（上同调，Lemma
\ref{cohomology-lemma-quasi-separated-cohomology-colimit}）。结论随即得出。
\end{proof}

\begin{lemma}
\label{more-morphisms-lemma-h1-fibre-zero-check-h0-kappa}
设 $f : X \to Y$ 是概形的固有态射。假设
\begin{enumerate}
\item 对所有 $y \in Y$，均有 $\dim(X_y) \leq 1$ 和
$H^1(X_y, \mathcal{O}_{X_y}) = 0$；
\item $\mathcal{O}_Y \to f_*\mathcal{O}_X$ 满射。
\end{enumerate}
则对 $f$ 的任一基变换 $f' : X' \to Y'$，态射
$\mathcal{O}_{Y'} \to f'_*\mathcal{O}_{X'}$ 满射。
\end{lemma}

\begin{proof}
可以假设 $Y$ 和 $Y'$ 仿射。于是可以选择闭浸入 $Y' \to Y''$，
其中 $Y'' \to Y$ 是仿射概形之间的平坦态射。由平坦基变换
（概形的上同调，引理 \ref{coherent-lemma-flat-base-change-cohomology}），
结论对 $X'' \to Y''$ 成立。因此可以假设 $Y'$ 是 $Y$ 的闭子概形。
令 $\mathcal{I} \subset \mathcal{O}_Y$ 为定义 $Y'$ 的理想。于是有短正合列
$$
0 \to \mathcal{I}\mathcal{O}_X \to
\mathcal{O}_X \to \mathcal{O}_{X'} \to 0
$$
其中把 $\mathcal{O}_{X'}$ 看成 $X$ 上的拟凝聚模。由
引理 \ref{more-morphisms-lemma-globally-generated-vanishing}，有
$H^1(X, \mathcal{I}\mathcal{O}_X) = 0$。因此
$$
H^0(Y, \mathcal{O}_Y) \to
H^0(Y, f_*\mathcal{O}_X) = H^0(X, \mathcal{O}_X) \to
H^0(X, \mathcal{O}_{X'})
$$
如所需是满射。第一个箭头满射，是因为 $Y$ 仿射，并且我们假设
$\mathcal{O}_Y \to f_*\mathcal{O}_X$ 满射；第二个箭头满射，则由上述
短正合列所伴随的上同调长正合列以及刚证明的消失得到。
\end{proof}

\begin{lemma}
\label{more-morphisms-lemma-h1-fibre-zero-h0-kappa}
考虑交换图表
$$
\xymatrix{
X \ar[rr]_f \ar[rd] & & Y \ar[ld] \\
& S
}
$$
为概形态射的交换图。设 $s \in S$ 是一点。假设
\begin{enumerate}
\item $X \to S$ 局部有限表示，并且在 $X_s$ 的各点处平坦；
\item $f$ 固有；
\item $f_s : X_s \to Y_s$ 的各纤维维数 $\leq 1$，并且
$R^1f_{s, *}\mathcal{O}_{X_s} = 0$；
\item $\mathcal{O}_{Y_s} \to f_{s, *}\mathcal{O}_{X_s}$ 满射。
\end{enumerate}
则存在开集 $Y_s \subset V \subset Y$，使得
(a) $f^{-1}(V)$ 在 $S$ 上平坦，
(b) 对 $y \in V$ 有 $\dim(X_y) \leq 1$，
(c) $R^1f_*\mathcal{O}_X|_V = 0$，
(d) $\mathcal{O}_V \to f_*\mathcal{O}_X|_V$ 满射，
并且沿任意 $Y' \to V$ 基变换后，(b)、(c)、(d) 仍成立。
\end{lemma}

\begin{proof}
令 $y \in Y$ 为位于 $s$ 上的一点。只需找到 $y$ 的一个具有所需性质的
开邻域。第一步，用 引理 \ref{more-morphisms-lemma-h1-fibre-zero} 中找到的开集 $V$
代替 $Y$，使得 $R^1f_*\mathcal{O}_X$ 普遍为零（该引理的假设由
引理 \ref{more-morphisms-lemma-check-h1-fibre-zero} 成立）。还可缩小 $Y$，使得 $f$
的所有纤维维数均 $\leq 1$（使用态射，Lemma
\ref{morphisms-lemma-openness-bounded-dimension-fibres} 以及 $f$ 的固有性）。
因此可以假设在 $V = Y$ 时 (b)、(c) 成立，并且沿任意
$Y' \to Y$ 基变换后仍成立。于是由
引理 \ref{more-morphisms-lemma-h1-fibre-zero-check-h0-kappa}，只需在 $Y$ 上证明 (d)。
仍可进一步缩小 $Y$；例如，可以并且确实假设 $Y$ 和 $S$ 都仿射。

\medskip\noindent
由 定理 \ref{more-morphisms-theorem-openness-flatness}，存在开子集 $U \subset X$，
$X \to S$ 在其上平坦；按假设，该开集包含 $X_s$。于是
$f(X \setminus U)$ 是不含 $y$ 的闭子集。因此缩小 $Y$ 后，
可以假设 $X$ 在 $S$ 上平坦。

\medskip\noindent
设 $S = \Spec(R)$。选择闭浸入 $Y \to Y'$，其中 $Y'$ 是集合 $E$
上的多项式环 $R[x_e; e \in E]$ 的谱。用 $f' : X \to Y'$ 表示 $f$
与 $Y \to Y'$ 的复合。于是对 $f'$ 和 $s$，假设 (1)--(4) 以及
(b)、(c) 均成立。若能证明 $\mathcal{O}_{Y'} \to f'_*\mathcal{O}_X$
在 $y$ 的某个开邻域中满射，则 $\mathcal{O}_Y \to f_*\mathcal{O}_X$
也如此。因此可以假设 $Y$ 是 $R[x_e; e \in E]$ 的谱。

\medskip\noindent
此时 $X$ 和 $Y$ 都在 $S$ 上平坦。于是 $Y_s$ 和 $X$ 在 $Y$ 上
Tor 独立。我们鼓励读者自己寻找证明；它也可由
引理 \ref{more-morphisms-lemma-case-of-tor-independence} 推出，只需把该引理应用于
顶点为 $X, Y, S, S$ 的方块及其沿 $s \to S$ 的基变换。因此
$$
Rf_{s, *}\mathcal{O}_{X_s} = L(Y_s \to Y)^*Rf_*\mathcal{O}_X
$$
这是由概形的导出范畴，Lemma
\ref{perfect-lemma-compare-base-change} 得到的。由已经建立的消失，它蕴含
$f_{s, *}\mathcal{O}_{X_s} = (Y_s \to Y)^*f_*\mathcal{O}_X$。
因此 $\mathcal{O}_Y \to f_*\mathcal{O}_X$ 是拟凝聚
$\mathcal{O}_Y$-模的映射，其到 $Y_s$ 的拉回满射。我们断言
$f_*\mathcal{O}_X$ 是有限型 $\mathcal{O}_Y$-模。若此断言成立，则
$\mathcal{O}_Y \to f_*\mathcal{O}_X$ 的余核 $\mathcal{F}$ 是有限型拟凝聚
$\mathcal{O}_Y$-模，并且 $\mathcal{F}_y \otimes \kappa(y) = 0$。
由 Nakayama 引理（代数，引理 \ref{algebra-lemma-NAK}），
$\mathcal{F}_y = 0$。因此 $\mathcal{F}$ 在 $y$ 的某个开邻域中为零
（模，引理 \ref{modules-lemma-finite-type-stalk-zero}），证明完成。

\medskip\noindent
证明该断言。对有限子集 $E' \subset E$，置
$Y' = \Spec(R[x_e; e \in E'])$。当 $E'$ 足够大时，态射
$f' : X \to Y \to Y'$ 固有，见极限，Lemma
\ref{limits-lemma-finite-type-eventually-proper}。以下固定 $E'$ 和 $Y'$。
把 $R = \colim R_i$ 写成其有限型 $\mathbf{Z}$-子代数的余极限。置
$S_i = \Spec(R_i)$ 和 $Y'_i = \Spec(R_i[x_e; e \in E'])$。
当 $i$ 足够大时，可以找到图
$$
\xymatrix{
X \ar[d] \ar[r]_{f'} & Y' \ar[d] \ar[r] & S \ar[d] \\
X_i \ar[r]^{f'_i} & Y'_i \ar[r] & S_i
}
$$
其中两个方块均为笛卡尔方块，$X_i$ 在 $S_i$ 上平坦，且
$X_i \to Y'_i$ 固有。见极限，Lemmas
\ref{limits-lemma-descend-finite-presentation}、
\ref{limits-lemma-descend-flat-finite-presentation} 和
\ref{limits-lemma-eventually-proper}。与上面相同的论证表明，$Y'$ 和
$X_i$ 在 $Y'_i$ 上 Tor 独立，因此
$$
R\Gamma(X, \mathcal{O}_X) =
R\Gamma(X_i, \mathcal{O}_{X_i})
\otimes^\mathbf{L}_{R_i[x_e; e \in E']}
R[x_e; e \in E']
$$
这里使用与上面相同的参考结果。由概形的上同调，Lemma
\ref{coherent-lemma-proper-over-affine-cohomology-finite}，复形
$R\Gamma(X_i, \mathcal{O}_{X_i})$ 在诺特环 $R_i[x_e; e \in E']$
的导出范畴中伪凝聚（见代数进阶，Lemma
\ref{more-algebra-lemma-Noetherian-pseudo-coherent}）。因此
$R\Gamma(X, \mathcal{O}_X)$ 在 $R[x_e; e \in E']$ 的导出范畴中伪凝聚，
见代数进阶，引理 \ref{more-algebra-lemma-pull-pseudo-coherent}。
由于唯一不消失的上同调模是 $H^0(X, \mathcal{O}_X)$，可知它是有限
$R[x_e; e \in E']$-模，见代数进阶，Lemma
\ref{more-algebra-lemma-n-pseudo-module}。证明完成。
\end{proof}

\begin{lemma}
\label{more-morphisms-lemma-h1-fibre-zero-h0-flat}
考虑交换图表
$$
\xymatrix{
X \ar[rr]_f \ar[rd] & & Y \ar[ld] \\
& S
}
$$
为概形态射的交换图。假设 $X \to S$ 平坦，$f$ 固有，对 $y \in Y$
有 $\dim(X_y) \leq 1$，并且 $R^1f_*\mathcal{O}_X = 0$。
则 $f_*\mathcal{O}_X$ 在 $S$ 上平坦，且 $f_*\mathcal{O}_X$
的形成与任意基变换 $S' \to S$ 相容。
\end{lemma}

\begin{proof}
可以假设 $Y$ 和 $S$ 仿射，设 $S = \Spec(A)$。为证明拟凝聚
$\mathcal{O}_Y$-模 $f_*\mathcal{O}_X$ 相对于 $S$ 平坦，只需证明
$H^0(X, \mathcal{O}_X)$ 在 $A$ 上平坦（略去一些细节）。由
引理 \ref{more-morphisms-lemma-globally-generated-vanishing}，对每个 $A$-模 $M$ 均有
$H^1(X, \mathcal{O}_X \otimes_A M) = 0$。又因 $\mathcal{O}_X$ 在 $A$
上平坦，可知函子 $M \mapsto H^0(X, \mathcal{O}_X \otimes_A M)$ 正合。
此外，由上同调，Lemma
\ref{cohomology-lemma-quasi-separated-cohomology-colimit}，该函子与直和相容。
于是容易验证，关于 $M$ 函子性地有
$H^0(X, \mathcal{O}_X \otimes_A M) = M \otimes_A H^0(X, \mathcal{O}_X)$；
这就给出所需的平坦性。最后，若 $S' \to S$ 是由环映射
$A \to A'$ 给出的仿射概形态射，则在刚讨论的仿射情形中有
$$
H^0(X \times_S S', \mathcal{O}_{X \times_S S'}) =
H^0(X, \mathcal{O}_X \otimes_A A') = H^0(X, \mathcal{O}_X) \otimes_A A'
$$
这表明 $f_*\mathcal{O}_X$ 的形成与任意基变换 $S' \to S$ 相容。
略去一些细节。
\end{proof}

\begin{lemma}
\label{more-morphisms-lemma-h1-fibre-zero-isom}
考虑交换图表
$$
\xymatrix{
X \ar[rr]_f \ar[rd] & & Y \ar[ld] \\
& S
}
$$
为概形态射的交换图。设 $s \in S$ 是一点。假设
\begin{enumerate}
\item $X \to S$ 局部有限表示，并且在 $X_s$ 的各点处平坦；
\item $Y \to S$ 局部有限表示；
\item $f$ 固有；
\item $f_s : X_s \to Y_s$ 的各纤维维数 $\leq 1$，并且
$R^1f_{s, *}\mathcal{O}_{X_s} = 0$；
\item $\mathcal{O}_{Y_s} \to f_{s, *}\mathcal{O}_{X_s}$ 是同构。
\end{enumerate}
则存在开集 $Y_s \subset V \subset Y$，使得
(a) $V$ 在 $S$ 上平坦，
(b) $f^{-1}(V)$ 在 $S$ 上平坦，
(c) 对 $y \in V$ 有 $\dim(X_y) \leq 1$，
(d) $R^1f_*\mathcal{O}_X|_V = 0$，
(e) $\mathcal{O}_V \to f_*\mathcal{O}_X|_V$ 是同构，
并且沿任意 $S' \to S$ 对 $f^{-1}(V) \to V$ 作基变换后，
(a)--(e) 仍成立。
\end{lemma}

\begin{proof}
令 $y \in Y_s$。总可以用 $y$ 的开邻域代替 $Y$。因此可以假设
$Y$ 和 $S$ 仿射。还可以假设 $X$ 在 $S$ 上平坦，对 $y \in Y$ 有
$\dim(X_y) \leq 1$，$R^1f_*\mathcal{O}_X = 0$ 普遍成立，并且
$\mathcal{O}_Y \to f_*\mathcal{O}_X$ 满射，见
引理 \ref{more-morphisms-lemma-h1-fibre-zero-h0-kappa}。（我们不会用到所有这些条件。）

\medskip\noindent
假设 $S$ 和 $Y$ 仿射。把 $S = \lim S_i$ 写成仿射诺特概形 $S_i$
的余滤逆极限。由极限，Lemma
\ref{limits-lemma-descend-finite-presentation}，存在元素 $0 \in I$ 以及图
$$
\xymatrix{
X_0 \ar[rr]_{f_0} \ar[rd] & & Y_0 \ar[ld] \\
& S_0
}
$$
它由概形的有限型态射组成，并且到 $S$ 的基变换是本引理中的图。
增大 $0$ 后，可以假设 $Y_0$ 仿射，且 $X_0 \to S_0$ 固有，见极限，
引理 \ref{limits-lemma-eventually-proper} 和
\ref{limits-lemma-limit-affine}。令 $s_0 \in S_0$ 为 $s$ 的像。
由于 $Y_s$ 仿射，$R^1f_{s, *}\mathcal{O}_{X_s} = 0$ 等价于
$H^1(X_s, \mathcal{O}_{X_s}) = 0$。由于 $X_s$ 是 $X_{0, s_0}$
沿忠实平坦映射 $\kappa(s_0) \to \kappa(s)$ 的基变换，可知
$H^1(X_{0, s_0}, \mathcal{O}_{X_{0, s_0}}) = 0$，因而
% P05-EN-0155: frozen source applies f_0 directly to the fibre X_{0,s_0}; the restricted morphism should be indexed f_{0,s_0}.
$R^1f_{0, *}\mathcal{O}_{X_{0, s_0}} = 0$。类似地，由于
$\mathcal{O}_{Y_s} \to f_{s, *}\mathcal{O}_{X_s}$ 是同构，
$\mathcal{O}_{Y_{0, s_0}} \to f_{0, *}\mathcal{O}_{X_{0, s_0}}$ 也是同构。
由于 $X_s \to Y_s$ 的各纤维维数至多为 $1$，对态射
$X_{0, s_0} \to Y_{0, s_0}$ 也如此。最后，由于 $X \to S$ 平坦，
增大 $0$ 后，可以假设 $X_0$ 在 $S_0$ 上平坦，见极限，Lemma
\ref{limits-lemma-descend-flat-finite-presentation}。因此只需对
$X_0 \to Y_0 \to S_0$ 和点 $s_0$ 证明本引理。

\medskip\noindent
合并上述归约论证，我们归约到如下情形：$S$ 和 $Y$ 仿射，$S$ 诺特，
$f$ 的各纤维维数 $\leq 1$，并且 $R^1f_*\mathcal{O}_X = 0$ 普遍成立。
设 $y \in Y_s$ 是一点。断言：
$$
\mathcal{O}_{Y, y} \longrightarrow (f_*\mathcal{O}_X)_y
$$
是同构。该断言蕴含本引理。事实上，$f_*\mathcal{O}_X$ 是凝聚的
（概形的上同调，Proposition
\ref{coherent-proposition-proper-pushforward-coherent}），故该断言表明，
可以用 $y$ 的开邻域代替 $Y$，并得到同构
$\mathcal{O}_Y \to f_*\mathcal{O}_X$。于是由
引理 \ref{more-morphisms-lemma-h1-fibre-zero-h0-flat}，$Y$ 在 $S$ 上平坦。
最后，由 引理 \ref{more-morphisms-lemma-h1-fibre-zero-h0-flat} 的最后一个断言，
同构 $\mathcal{O}_Y \to f_*\mathcal{O}_X$ 沿任意基变换 $S' \to S$
后仍为同构。

\medskip\noindent
证明该断言。我们已经知道
$\mathcal{O}_{Y, y} \longrightarrow (f_*\mathcal{O}_X)_y$ 满射
（引理 \ref{more-morphisms-lemma-h1-fibre-zero-h0-kappa}），并且
$(f_*\mathcal{O}_X)_y$ 是 $\mathcal{O}_{S, s}$-平坦的
（引理 \ref{more-morphisms-lemma-h1-fibre-zero-h0-flat}），而且所诱导的映射
$$
\mathcal{O}_{Y_s, y} =  \mathcal{O}_{Y, y}/\mathfrak m_s\mathcal{O}_{Y, y}
\longrightarrow
(f_*\mathcal{O}_X)_y/\mathfrak m_s (f_*\mathcal{O}_X)_y
\to
(f_{s, *}\mathcal{O}_{X_s})_y
$$
由本引理的假设是单射。于是由代数，Lemma
\ref{algebra-lemma-mod-injective}，
$\mathcal{O}_{Y, y} \longrightarrow (f_*\mathcal{O}_X)_y$ 是单射，
正合所需。
\end{proof}

\begin{lemma}
\label{more-morphisms-lemma-bijection-on-Pic}
设 $f : X \to Y$ 是诺特概形的固有态射，满足
$f_*\mathcal{O}_X = \mathcal{O}_Y$，$f$ 的各纤维维数 $\leq 1$，
并且对 $y \in Y$ 有 $H^1(X_y, \mathcal{O}_{X_y}) = 0$。
则 $f^* : \Pic(Y) \to \Pic(X)$ 双射到由如下元素组成的子群：
$\mathcal{L} \in \Pic(X)$，并且对所有 $y \in Y$ 都有
$\mathcal{L}|_{X_y} \cong \mathcal{O}_{X_y}$。
\end{lemma}

\begin{proof}
由投影公式（上同调，引理 \ref{cohomology-lemma-projection-formula}），
对 $\mathcal{N} \in \Pic(Y)$ 有 $f_*f^*\mathcal{N} \cong \mathcal{N}$。
我们断言，若 $\mathcal{L} \in \Pic(X)$ 满足对所有 $y \in Y$ 都有
$\mathcal{L}|_{X_y} \cong \mathcal{O}_{X_y}$，则
$\mathcal{N} = f_*\mathcal{L}$ 可逆，并且
$\mathcal{L} \cong f^*\mathcal{N}$。这将完成证明。

\medskip\noindent
由概形的上同调，Proposition
\ref{coherent-proposition-proper-pushforward-coherent}，
$\mathcal{O}_Y$-模 $\mathcal{N} = f_*\mathcal{L}$ 是凝聚的。
因此为证明它是可逆 $\mathcal{O}_Y$-模，只需在茎上检验（代数，Lemma
\ref{algebra-lemma-finite-projective}）。由于诺特局部环到其完备化的映射
忠实平坦，只需检验完备化 $(f_*\mathcal{L})_y^\wedge$ 自由（见代数，
第 \ref{algebra-section-completion-noetherian} 节 和
引理 \ref{algebra-lemma-finite-projective-descends}）。为此，将使用概形的上同调，
引理 \ref{coherent-lemma-formal-functions-stalk} 所表述的形式函数定理。
由于 $f_*\mathcal{O}_X = \mathcal{O}_Y$，从而
$(f_*\mathcal{O}_X)_y^\wedge \cong \mathcal{O}_{Y, y}^\wedge$，
只需证明对每个 $n$ 都有 $\mathcal{L}|_{X_n} \cong \mathcal{O}_{X_n}$
% P05-EN-0156: frozen source opens but does not close the parenthetical "compatibly for varying n".
（并且关于变化的 $n$ 相容。
由 引理 \ref{more-morphisms-lemma-picard-group-first-order-thickening}，有正合列
$$
H^1(X_y, \mathfrak m_y^n\mathcal{O}_X/\mathfrak m_y^{n + 1}\mathcal{O}_X)
\to \Pic(X_{n + 1}) \to \Pic(X_n)
$$
记号与形式函数定理中相同。注意有满射
$$
\mathcal{O}_{X_y}^{\oplus r_n} \cong
\mathfrak m_y^n/\mathfrak m_y^{n + 1} \otimes_{\kappa(y)} \mathcal{O}_{X_y}
\longrightarrow
\mathfrak m_y^n\mathcal{O}_X/\mathfrak m_y^{n + 1}\mathcal{O}_X
$$
其中 $r_n \geq 0$ 是某个整数。由于 $\dim(X_y) \leq 1$，该满射诱导
第一上同调群之间的满射（由上同调，Proposition
\ref{cohomology-proposition-vanishing-Noetherian}，次数 $\geq 2$
的上同调消失）。因此正合列中的 $H^1$ 为零，而转移映射
$\Pic(X_{n + 1}) \to \Pic(X_n)$ 如所需是单射。

\medskip\noindent
还需证明 $f^*\mathcal{N} \cong \mathcal{L}$。这可用相同方法证明，
略去细节。
\end{proof}









\section{仿射分层}
\label{more-morphisms-section-affine-stratifications}

\noindent
本节材料取自 \cite{RV}。阅读本节之前，请先阅读拓扑，Section
\ref{topology-section-stratifications} 中关于分层的一些内容。

\medskip\noindent
若 $X$ 是概形，则 $X$ 的分层通常是指 $X$ 的底层拓扑空间的分层。
各层是局部闭子集。使用概形，Remark
\ref{schemes-remark-reduced-induced-locally-closed}，我们把这些层看作
$X$ 的约化局部闭子概形。

\begin{definition}
\label{more-morphisms-definition-affine-stratification}
设 $X$ 是概形。{\it 仿射分层}是局部有限分层
$X = \coprod_{i \in I} X_i$，其中各层 $X_i$ 仿射，并且嵌入态射
$X_i \to X$ 是仿射的。
\end{definition}

\noindent
在嵌入态射 $X_i \to X$ 拟紧的条件下，分层 $X = \coprod X_i$
局部有限，等价于各层是 $X$ 的局部可构造子集，见性质，Lemma
\ref{properties-lemma-stratification-locally-finite-constructible}。

\medskip\noindent
$X_i \to X$ 是仿射态射这一条件，与我们赋予局部闭子集 $X_i$
的概形结构无关，见 引理 \ref{more-morphisms-lemma-thicken-property-morphisms}。
此外，若 $X$ 分离（或更一般地，其对角态射仿射），且
$X = \coprod X_i$ 是具有仿射层的局部有限分层，则态射
$X_i \to X$ 仿射。见态射，Lemma
\ref{morphisms-lemma-affine-permanence}。因此在大多数关注的情形中，
可以不再单独考虑嵌入态射 $X_i \to X$ 的仿射性条件。

\medskip\noindent
我们通常关注分层的偏序指标集 $I$ 有限的情形。回忆，偏序集 $I$
的{\it 长度}是 $I$ 的元素链 $i_0 < i_1 < \ldots < i_p$ 的长度 $p$
的上确界。

\begin{lemma}
\label{more-morphisms-lemma-improve-stratification}
设 $X$ 是概形，$X = \coprod_{i \in I} X_i$ 是有限仿射分层。
则存在指标集为 $\{0, \ldots, n\}$ 的仿射分层，其中 $n$ 是 $I$ 的长度。
\end{lemma}

\begin{proof}
回忆，$I$ 上有偏序，使得对所有 $i \in I$，$X_i$ 的闭包包含在
$\bigcup_{j \leq i} X_j$ 中。令 $I' \subset I$ 为 $I$ 的极大指标所成的集合。
若 $i \in I'$，则 $X_i$ 在 $X$ 中开，因为其余各层的闭包之并正是
$X_i$ 的补集。把 $U = \bigcup_{i \in I'} X_i$ 看作 $X$ 的开子概形，
于是作为概形有 $U_{red} = \coprod_{i \in I'} X_i$。由概形，Lemma
\ref{schemes-lemma-disjoint-union-affines} 和
引理 \ref{more-morphisms-lemma-thickening-affine-scheme}，$U$ 是仿射概形。态射
$U \to X$ 是仿射的，因为使用 Lemma
\ref{more-morphisms-lemma-thicken-property-morphisms} 作同样的论证可知，每个
$X_i \to X$（$i \in I'$）都是仿射的。给补集 $Z = X \setminus U$
赋予约化诱导概形结构，则它有仿射分层
$Z = \bigcup_{i \in I \setminus I'} X_i$。这里使用了如下事实：
概形态射 $T \to Z$ 仿射，当且仅当复合 $T \to X$ 仿射；这由态射，
引理 \ref{morphisms-lemma-closed-immersion-affine}、
\ref{morphisms-lemma-composition-affine} 和
\ref{morphisms-lemma-affine-permanence} 得出。注意，偏序集
$I \setminus I'$ 的长度恰好比 $I$ 的长度小一。因此由归纳可知，$Z$
有仿射分层 $Z = Z_0 \amalg \ldots \amalg Z_{n - 1}$，其指标集为
% P05-EN-0157: frozen source labels strata Z_0 through Z_{n-1} but states the index set is {1,...,n}; it should be {0,...,n-1}.
$\{1, \ldots, n\}$。置 $Z_n = U$，即得 $X$ 的所需分层。
\end{proof}

\noindent
若概形 $X$ 有有限仿射分层，则 $X$ 当然拟紧。稍不明显的是，
这还迫使 $X$ 拟分离。

\begin{lemma}
\label{more-morphisms-lemma-qc-affine-stratification}
设 $X$ 是概形。下列条件等价：
\begin{enumerate}
\item $X$ 有有限仿射分层；
\item $X$ 拟紧且拟分离。
\end{enumerate}
\end{lemma}

\begin{proof}
设 $X = \bigcup X_i$ 是有限仿射分层。由于每个 $X_i$ 仿射，因而拟紧，
可知 $X$ 拟紧。设 $U, V \subset X$ 是仿射开集。由于 $X_i \to X$
是仿射态射，$U \cap X_i$ 和 $V \cap X_i$ 是 $X_i$ 中的仿射开集。
因此 $U \cap V \cap X_i$ 是仿射概形 $X_i$ 的仿射开集（例如见概形，
引理 \ref{schemes-lemma-characterize-separated}）。于是
$U \cap V = \coprod U \cap V \cap X_i$ 作为有限多个仿射层的并是拟紧的。
由概形，引理 \ref{schemes-lemma-characterize-quasi-separated}，
$X$ 拟分离。

\medskip\noindent
假设 $X$ 拟紧且拟分离。可以使用概形的上同调，Lemma
\ref{coherent-lemma-induction-principle} 的归纳原理，证明 $X$
有有限仿射分层。若 $X$ 为空，则它有空的仿射分层。若 $X$
是非空仿射概形，则它有仅含一层的仿射分层。接下来假设
$X = U \cup V$，其中 $U$ 是拟紧开集，$V$ 是仿射开集，并且已有有限
仿射分层 $U = \bigcup_{i \in I} U_i$ 和
$U \cap V = \coprod_{j \in J} W_j$。记 $Z = X \setminus V$ 和
$Z' = X \setminus U$。注意，$Z$ 在 $U$ 中闭，而 $Z'$ 在 $V$ 中闭。
注意，$U_i \cap Z$ 和 $U_i \cap W_j = U_i \times_U W_j$ 是在 $U$
上仿射的仿射概形。（提示：$U_i \times_U W_j \to W_j$ 是
$U_i \to U$ 的基变换，故仿射；从而 $U_i \cap W_j$ 仿射，进而
$U_i \cap W_j \to U_i$ 仿射，最后 $U_i \cap W_j \to U$ 仿射。）
因此
$$
U = \coprod\nolimits_{i \in I} (U_i \cap Z) \amalg
\coprod\nolimits_{(i, j) \in I \times J} (U_i \cap W_j)
$$
是有限仿射分层；$I \amalg I \times J$ 上的偏序由
$i' \leq (i, j) \Leftrightarrow i' \leq i$ 和
$(i', j') \leq (i, j) \Leftrightarrow i' \leq i$ 且 $j' \leq j$ 给出。
注意，$(U_i \cap Z) \times_X V = \emptyset$ 和
$(U_i \cap W_j) \times_X V = U_i \cap W_j$ 都仿射。因此态射
$U_i \cap Z \to X$ 和 $U_i \cap W_j \to X$ 仿射，因为态射的仿射性
可以在靶上局部地检验（态射，Lemma
\ref{morphisms-lemma-characterize-affine}），而我们已经有了在 $U$ 和 $V$
上的仿射性。为完成证明，在上述分层中再加入一层，即 $Z'$，
并把它的指标作为该偏序集的一个极小元。
\end{proof}

\begin{definition}
\label{more-morphisms-definition-affine-stratification-number}
设 $X$ 是非空拟紧拟分离概形。{\it 仿射分层数}是使下列等价条件成立的
最小整数 $n \geq 0$：
\begin{enumerate}
\item 存在有限仿射分层 $X = \coprod_{i \in I} X_i$，其中 $I$ 的长度为 $n$；
\item 存在仿射分层 $X = X_0 \amalg X_1 \amalg \ldots \amalg X_n$，
其指标集为 $\{0, \ldots, n\}$。
\end{enumerate}
\end{definition}

\noindent
这些条件的等价性由 引理 \ref{more-morphisms-lemma-improve-stratification} 成立。
有限仿射分层的存在性已在 Lemma
\ref{more-morphisms-lemma-qc-affine-stratification} 中证明。

\begin{lemma}
\label{more-morphisms-lemma-affine-stratification-number-bound}
设 $X$ 是分离概形，且有一个由 $n + 1$ 个仿射开集组成的开覆盖。
则 $X$ 的仿射分层数至多为 $n$。
\end{lemma}

\begin{proof}
设 $X = U_0 \cup \ldots \cup U_n$ 是仿射开覆盖。置
$$
X_i = (U_i \cup \ldots \cup U_n) \setminus
(U_{i + 1} \cup \ldots \cup U_n)
$$
于是 $X_i$ 作为 $U_i$ 的闭子概形是仿射的。由态射，Lemma
\ref{morphisms-lemma-affine-permanence}，态射 $X_i \to X$ 仿射。
最后有 $\overline{X_i} \subset X_i \cup X_{i - 1} \cup \ldots X_0$。
\end{proof}

\begin{lemma}
\label{more-morphisms-lemma-affine-stratification-number-bound-Noetherian}
设 $X$ 是维数为 $\infty > d \geq 0$ 的诺特概形。
则 $X$ 的仿射分层数至多为 $d$。
\end{lemma}

\begin{proof}
对 $d$ 归纳。若 $d = 0$，则 $X$ 仿射，见性质，Lemma
\ref{properties-lemma-locally-Noetherian-dimension-0}。假设 $d > 0$。
令 $\eta_1, \ldots, \eta_n$ 为 $X$ 的不可约分支的泛点（性质，Lemma
\ref{properties-lemma-Noetherian-irreducible-components}）。可以用包含
$\eta_1, \ldots, \eta_n$ 的仿射开集覆盖 $X$，见性质，Lemma
\ref{properties-lemma-point-and-maximal-points-affine}。由于 $X$ 拟紧，
可以找到有限仿射开覆盖 $X = \bigcup_{j = 1, \ldots, m} U_j$，
使得对所有 $j = 1, \ldots, m$ 都有
$\eta_1, \ldots, \eta_n \in U_j$。选择包含
$\eta_1, \ldots, \eta_n$ 的仿射开集
$U \subset U_1 \cap \ldots \cap U_m$（由刚引用的引理可以做到）。
于是态射 $U \to X$ 仿射，因为对所有 $j$，$U \to U_j$ 都仿射，
见态射，引理 \ref{morphisms-lemma-characterize-affine}。令
$Z = X \setminus U$。按构造，$\dim(Z) < \dim(X)$。由归纳假设，
可以找到 $Z$ 的仿射分层
$Z = \bigcup_{i \in \{0, \ldots, n\}} Z_i$，其中 $n \leq \dim(Z)$。
% P05-EN-0158: frozen source reverses the assignment and writes U = X_{n+1}; the new top stratum should be X_{n+1} = U.
置 $U = X_{n + 1}$，并对 $i \leq n$ 置 $X_i = Z_i$，即得结论。
\end{proof}

\begin{proposition}
\label{more-morphisms-proposition-vanishing-affine-stratification-number}
设 $X$ 是仿射分层数为 $n$ 的非空拟紧拟分离概形。则对每个拟凝聚
$\mathcal{O}_X$-模 $\mathcal{F}$，当 $p > n$ 时有
$H^p(X, \mathcal{F}) = 0$。
\end{proposition}

\begin{proof}
我们对仿射分层数 $n$ 归纳。若 $n = 0$，则 $X$ 仿射，结论就是
概形的上同调，Lemma
\ref{coherent-lemma-quasi-coherent-affine-cohomology-zero}。假设 $n > 0$。
由 定义 \ref{more-morphisms-definition-affine-stratification-number}，存在仿射概形
$U$ 和仿射开浸入 $j : U \to X$，使得补集 $Z$ 的仿射分层数为 $n - 1$。
由于 $U$ 和 $j$ 仿射，对 $p > 0$ 有
$H^p(X, j_*(\mathcal{F}|_U)) = 0$，见概形的上同调，Lemmas
\ref{coherent-lemma-relative-affine-cohomology} 和
\ref{coherent-lemma-relative-affine-vanishing}。用 $\mathcal{K}$ 和
$\mathcal{Q}$ 表示映射 $\mathcal{F} \to j_*(\mathcal{F}|_U)$ 的核和余核。
于是得到正合列
$$
0 \to \mathcal{K} \to \mathcal{F} \to j_*(\mathcal{F}|_U)
\to \mathcal{Q} \to 0
$$
它由拟凝聚 $\mathcal{O}_X$-模组成（见概形，Section
\ref{schemes-section-quasi-coherent}）。把该正合列分成短正合列并使用上同调
长正合列的标准论证表明，只需证明当 $p \geq n$ 时，
$H^p(X, \mathcal{K}) = 0$ 且 $H^p(X, \mathcal{Q}) = 0$。
由于 $\mathcal{F} \to j_*(\mathcal{F}|_U)$ 在 $U$ 上限制为同构，
$\mathcal{K}$ 和 $\mathcal{Q}$ 都支撑在 $Z$ 上。由性质，Lemma
\ref{properties-lemma-quasi-coherent-colimit-finite-type}，可以把这些模
写成其有限型拟凝聚子模的滤过余极限。利用 $X$ 上层的上同调与滤过余极限
相容这一事实，见上同调，Lemma
\ref{cohomology-lemma-quasi-separated-cohomology-colimit}，只需证明如下断言：
若 $\mathcal{G}$ 是支撑包含在 $Z$ 中的有限型拟凝聚模，则对
$p \geq n$ 有 $H^p(X, \mathcal{G}) = 0$。令 $Z' \subset X$ 为
$\mathcal{G} \oplus \mathcal{O}_Z$ 的概形论支撑；可以并且确实把
$\mathcal{G}$ 看成 $Z'$ 上的拟凝聚模，见态射，Section
\ref{morphisms-section-support}。于是 $Z'$ 和 $Z$ 有相同的底层拓扑空间，
从而有相同的仿射分层数，即 $n - 1$。因此由归纳假设，
$H^p(X, \mathcal{G}) = H^p(Z', \mathcal{G})$（该等式由概形的上同调，
引理 \ref{coherent-lemma-relative-affine-cohomology} 得出）在
$p \geq n$ 时为零。
\end{proof}

\begin{example}
\label{more-morphisms-example-affine-stratification-number}
设 $k$ 是域，$X = \mathbf{P}^n_k$ 是 $k$ 上的 $n$ 维射影空间。
引理 \ref{more-morphisms-lemma-affine-stratification-number-bound} 由构造，Lemma
\ref{constructions-lemma-standard-covering-projective-space} 可用于它。因此
$\mathbf{P}^n_k$ 的仿射分层数至多为 $n$。另一方面，
$\mathbf{P}^n_k$ 上的某些拟凝聚模在次数 $n$ 上有非零上同调，
见概形的上同调，Lemma
\ref{coherent-lemma-cohomology-projective-space-over-ring}。使用
命题 \ref{more-morphisms-proposition-vanishing-affine-stratification-number}，
可知 $\mathbf{P}^n_k$ 的仿射分层数等于 $n$。
\end{example}








\section{普遍开态射}
\label{more-morphisms-section-universally-open}

\noindent
下面给出关于普遍开态射的一些材料。

\begin{lemma}
\label{more-morphisms-lemma-test-universally-open}
设 $f : X \to S$ 是概形态射。下列条件等价：
\begin{enumerate}
\item $f$ 普遍开；
\item 对每个局部有限表示的态射 $S' \to S$，基变换
$X_{S'} \to S'$ 是开态射；
\item 对每个 $n$，态射 $\mathbf{A}^n \times X \to \mathbf{A}^n \times S$
都是开态射。
\end{enumerate}
\end{lemma}

\begin{proof}
显然，(1) 蕴含 (2)，而 (2) 蕴含 (3)。下面证明 (3) 蕴含 (1)。
假设对某个概形态射 $g : T \to S$，基变换 $X_T \to T$ 不是开态射。
则可以找到仿射开集 $V \subset S$、$U \subset X$、$W \subset T$，
满足 $f(U) \subset V$ 和 $g(W) \subset V$，并且
$U \times_V W \to W$ 不是开态射。若能证明这蕴含
$\mathbf{A}^n \times U \to \mathbf{A}^n \times V$ 不是开态射，
则 $\mathbf{A}^n \times X \to \mathbf{A}^n \times S$ 不是开态射，
证明完成。这把问题归约到下一段证明的结果。

\medskip\noindent
设 $A \to B$ 是环映射，并且对某个 $A$-代数 $A'$，
$A' \to B' = A' \otimes_A B$ 所诱导的谱映射不是开映射。
由于主开集构成 $\Spec(B')$ 的拓扑基，可知对某个 $g \in B'$，
$D(g)$ 在 $\Spec(A')$ 中的像不是开集。写成
$g = \sum_{i = 1, \ldots, n} a'_i \otimes b_i$，其中 $n$、
$a'_i \in A'$ 和 $b_i \in B$ 如此选取。考虑
$B[x_1, \ldots, x_n]$ 中的元素
$h = \sum_{i = 1, \ldots, n} x_i b_i$。反设 $D(h)$ 在态射
$$
\Spec(B[x_1, \ldots, x_n]) \longrightarrow \Spec(A[x_1, \ldots, x_n])
$$
下的像为开集，以期导出矛盾。此时 $D(h)$ 将满射到某个拟紧开集
$U \subset \Spec(A[x_1, \ldots, x_n])$。令
$A[x_1, \ldots, x_n] \to A'$ 为把 $x_i$ 映到 $a'_i$ 的 $A$-代数同态。
它还诱导把 $h$ 映到 $g$ 的 $B$-代数同态
$B[x_1, \ldots, x_n] \to B'$。由于
$$
\xymatrix{
\Spec(B[x_1, \ldots, x_n]) \ar[d] &
\Spec(B') \ar[l] \ar[d] \\
\Spec(A[x_1, \ldots, x_n]) &
\Spec(A') \ar[l]
}
$$
是笛卡尔图，$D(g)$ 在 $\Spec(A')$ 中的像等于 $U$ 在 $\Spec(A')$
中的原像，因而是开集。这就是所需的矛盾。
\end{proof}

\begin{lemma}
\label{more-morphisms-lemma-quasi-finite-Noetherian-universally-open}
设 $f : X \to Y$ 是概形态射。若
\begin{enumerate}
\item $f$ 局部拟有限；
\item $Y$ 几何单分支且局部诺特；
\item $X$ 的每个不可约分支都支配 $Y$ 的某个不可约分支，
\end{enumerate}
则 $f$ 普遍开。
\end{lemma}

\begin{proof}
对任意 $n$，由 引理 \ref{more-morphisms-lemma-number-of-branches-and-smooth} 和性质，
引理 \ref{properties-lemma-number-of-branches-1}，概形
$\mathbf{A}^n \times Y$ 几何单分支。因此对所有 $n$，态射
$\mathbf{A}^n \times X \to \mathbf{A}^n \times Y$ 满足本引理的假设。
由 引理 \ref{more-morphisms-lemma-test-universally-open}，只需证明 $f$ 是开态射。
由态射，Lemma
\ref{morphisms-lemma-locally-finite-presentation-universally-open}，
只需证明一般化可以沿 $f$ 提升。假设 $y' \leadsto y$ 是 $Y$
中点的特化，而 $x \in X$ 映到 $y$。按照
引理 \ref{more-morphisms-lemma-etale-makes-quasi-finite-finite-at-point} 选择图
$$
\xymatrix{
u \ar[d] & U \ar[d] \ar[r] & X \ar[d] \\
v & V \ar[r] & Y
}
$$
其中 $(V, v) \to (Y, y)$ 是初等 \'etale 邻域，$U \to V$ 有限，
$u$ 是 $U$ 中唯一映到 $v$ 的点，$U \subset V \times_Y X$ 是开集，
并且 $v \mapsto y$、$u \mapsto x$。令 $E$ 为 $U$ 的一个经过 $u$
的不可约分支（至少存在一个）。由于 $U \to X$ 是 \'etale 的，
$E$ 映到 $X$ 的某个不可约分支，而后者又支配 $Y$ 的某个不可约分支
（由假设）。由于 $U \to V$ 有限，因而闭，$E$ 的像 $E' \subset V$
是经过 $v$ 的不可约闭子集，并支配 $Y$ 的某个不可约分支。由于
$V \to Y$ 是 \'etale 的，$E'$ 必为 $V$ 的一个经过 $v$ 的不可约分支。
由于 $Y$ 几何单分支，$E'$ 是 $V$ 的唯一一个经过 $v$ 的不可约分支
（引理 \ref{more-morphisms-lemma-nr-branches}）。由于 $V$ 局部诺特，缩小 $V$ 后，
可以假设 $E' = V$（集合相等）。

\medskip\noindent
由于 $V \to Y$ 是 \'etale 的，可以找到特化 $v' \leadsto v$，
其像为 $y' \leadsto y$。由上可找到映到 $v'$ 的 $u' \in U$。
于是 $u' \leadsto u$，因为 $u$ 是 $U$ 中唯一映到 $v$ 的点，
且 $U \to V$ 是闭态射。最后，$u'$ 的像 $x' \in X$ 特化到 $x$，
并映到 $y'$。证明完成。
\end{proof}

\begin{lemma}
\label{more-morphisms-lemma-characterize-universally-open-finite}
设 $A \to B$ 是环映射，并设 $B$ 作为 $A$-模由
$b_1, \ldots, b_d \in B$ 生成。置
$h = \sum x_ib_i \in B[x_1, \ldots, x_d]$。则
$\Spec(B) \to \Spec(A)$ 普遍开，当且仅当 $D(h)$ 在
$\Spec(A[x_1, \ldots, x_d])$ 中的像为开集。
\end{lemma}

\begin{proof}
若 $\Spec(B) \to \Spec(A)$ 普遍开，则 $D(h)$ 的像当然为开集。
反过来，假设 $D(h)$ 的像 $U$ 为开集。设 $A \to A'$ 是环映射。
只需证明任一主开集 $D(g) \subset \Spec(A' \otimes_A B)$ 在
$\Spec(A')$ 中的像为开集。可以写成
$g = \sum_{i = 1, \ldots, d} a'_i \otimes b_i$，其中
$a'_i \in A'$。令
% P05-EN-0159: frozen source switches the generator count d to an undefined n in both polynomial rings of this proof.
$A[x_1, \ldots, x_n] \to A'$ 为把 $x_i$ 映到 $a'_i$ 的 $A$-代数同态。
它还诱导把 $h$ 映到 $g$ 的 $B$-代数同态
$B[x_1, \ldots, x_n] \to A' \otimes_A B$。由于
$$
\xymatrix{
\Spec(B[x_1, \ldots, x_n]) \ar[d] &
\Spec(B') \ar[l] \ar[d] \\
\Spec(A[x_1, \ldots, x_n]) &
\Spec(A') \ar[l]
}
$$
% P05-EN-0160: frozen source uses Spec(B') in this diagram without defining B' in the lemma; the base-changed ring is A' tensor_A B.
是笛卡尔图，$D(g)$ 在 $\Spec(A')$ 中的像等于 $U$ 在 $\Spec(A')$
中的原像，因而为开集。
\end{proof}

\begin{lemma}
\label{more-morphisms-lemma-descend-quasi-finite-universally-open}
设 $S = \lim S_i$ 是转移态射均仿射的概形有向系的极限。设
$0 \in I$，并设 $f_0 : X_0 \to Y_0$ 是 $S_0$ 上的概形态射。
假设 $S_0$、$X_0$、$Y_0$ 拟紧且拟分离。令
$f_i : X_i \to Y_i$ 为 $f_0$ 到 $S_i$ 的基变换，并令
$f : X \to Y$ 为 $f_0$ 到 $S$ 的基变换。若
\begin{enumerate}
\item $f$ 局部拟有限且普遍开；
\item $f_0$ 局部有限表示，
\end{enumerate}
则存在 $i \geq 0$，使得 $f_i$ 局部拟有限且普遍开。
\end{lemma}

\begin{proof}
由极限，引理 \ref{limits-lemma-descend-quasi-finite}，增大 $0$ 后，
可以假设 $f_0$ 局部拟有限。令 $x \in X$。通过拟有限态射的 \'etale
局部化，可以找到图
$$
\xymatrix{
X \ar[d] & U \ar[l] \ar[d] \\
Y & V \ar[l]
}
$$
其中 $V \to Y$ 是 \'etale 的，$U \subset X_V$ 是开集，$U \to V$ 有限，
且 $x$ 位于 $U \to X$ 的像中，见
引理 \ref{more-morphisms-lemma-etale-makes-quasi-finite-finite-at-point}。缩小 $V$ 后，
可以假设 $V$ 和 $U$ 仿射。由于 $X$ 拟紧，取有限多个这样的 $V$ 和 $U$
的不交并，便可使上述图中的 $U \to X$ 满射。由极限，Lemmas
\ref{limits-lemma-descend-finite-presentation}、
\ref{limits-lemma-descend-opens}、
\ref{limits-lemma-descend-surjective}、
\ref{limits-lemma-descend-finite-finite-presentation}、
\ref{limits-lemma-descend-etale} 和
\ref{limits-lemma-limit-affine}，必要时增大 $0$ 后，可以假设有图
$$
\xymatrix{
X_0 \ar[d] & U_0 \ar[l] \ar[d] \\
Y_0 & V_0 \ar[l]
}
$$
其中 $V_0$ 仿射，$V_0 \to Y_0$ 是 \'etale 的，
$U_0 \subset (X_0)_{V_0}$ 是开集，$U_0 \to V_0$ 有限，且
$U_0 \to X_0$ 满射。由于 $V_i \to Y_i$ 是 \'etale 的，因而普遍开，
只需证明当 $i$ 足够大时 $U_i \to V_i$ 普遍开。这把我们归约到
下一段讨论的情形。

\medskip\noindent
设 $A = \colim A_i$ 是环的滤过余极限，$A_0 \to B_0$ 是环映射。
置 $B = A \otimes_{A_0} B_0$ 和 $B_i = A_i \otimes_{A_0} B_0$。
假设 $A_0 \to B_0$ 有限且有限表示，并且 $A \to B$ 普遍开。
我们要证明当 $i$ 足够大时 $A_i \to B_i$ 普遍开。选择生成 $B_0$
这一 $A_0$-模的元素 $b_{0, 1}, \ldots, b_{0, d} \in B_0$。
在 $B_0[x_1, \ldots, x_d]$ 中置
$h_0 = \sum_{j = 1, \ldots, d} x_jb_{0, j}$。分别用 $h$ 和 $h_i$
表示 $h_0$ 在 $B[x_1, \ldots, x_d]$ 和
$B_i[x_1, \ldots, x_d]$ 中的像。由于 $A \to B$ 普遍开，$D(h)$
在 $\Spec(A[x_1, \ldots, x_d])$ 中的像 $U$ 是开集。由于它是仿射概形的像，
$U$ 当然拟紧。当 $i$ 足够大时，存在拟紧开集
$U_i \subset \Spec(A_i[x_1, \ldots, x_d])$，其在
$\Spec(A[x_1, \ldots, x_d])$ 中的原像为 $U$，见极限，Lemma
\ref{limits-lemma-descend-opens}。增大 $i$ 后，可以假设 $D(h_i)$ 映入
$U_i$；这可把同一引理应用于 $U_i$ 在 $D(h_i)$ 中的拉回得到。
最后，当 $i$ 更大时，概形态射 $D(h_i) \to U_i$ 将满射，
这是再次应用极限，引理 \ref{limits-lemma-descend-surjective} 的结果。
由 引理 \ref{more-morphisms-lemma-characterize-universally-open-finite}，可知
$A_i \to B_i$ 普遍开。
\end{proof}

\begin{lemma}
\label{more-morphisms-lemma-count-geometric-fibres}
设 $f : X \to Y$ 是局部拟有限态射。则
\begin{enumerate}
\item Lemmas
\ref{more-morphisms-lemma-base-change-fibres-nr-geometrically-irreducible-components} 和
\ref{more-morphisms-lemma-base-change-fibres-nr-geometrically-connected-components}
中的函数 $n_{X/Y}$ 相同；
\item 若 $X$ 拟紧，则 $n_{X/Y}$ 取得最大值 $d < \infty$。
\end{enumerate}
\end{lemma}

\begin{proof}
局部拟有限态射的（几何）纤维是离散的，故两个函数相同，见态射，Lemma
\ref{morphisms-lemma-locally-quasi-finite-fibres}。有界性由态射，Lemmas
\ref{morphisms-lemma-characterize-universally-bounded} 和
\ref{morphisms-lemma-locally-quasi-finite-qc-source-universally-bounded} 得出。
\end{proof}

\begin{lemma}
\label{more-morphisms-lemma-count-geometric-fibres-universally-open}
设 $f : X \to Y$ 是分离、局部拟有限且普遍开的概形态射。设
$n_{X/Y}$ 如 引理 \ref{more-morphisms-lemma-count-geometric-fibres} 中所述。
若对某个 $y \in Y$ 和 $d \geq 0$ 有 $n_{X/Y}(y) \geq d$，
则在 $y$ 的某个开邻域中有 $n_{X/Y} \geq d$。
\end{lemma}

\begin{proof}
问题在 $Y$ 上是局部的，故可以假设 $Y$ 仿射。令 $K$ 为剩余域
$\kappa(y)$ 的代数闭包。假设表示 $(X_y)_K$ 有 $\geq d$ 个连通分支。
于是对某个合适的拟紧开集 $X' \subset X$，概形 $(X'_y)_K$
有 $\geq d$ 个连通分支；略去细节。以 $X'$ 代替 $X$ 后，可以假设
$X$ 拟紧。于是 $f$ 拟有限。令 $x_1, \ldots, x_n$ 为 $X$ 中位于
$y$ 上的点。应用
引理 \ref{more-morphisms-lemma-etale-splits-off-quasi-finite-part-technical-variant}，
得到 \'etale 邻域 $(U, u) \to (Y, y)$ 以及分解
$$
U \times_Y X =
W \amalg
\ \coprod\nolimits_{i = 1, \ldots, n}
\ \coprod\nolimits_{j = 1, \ldots, m_i}
V_{i, j}
$$
如所引之处所述。注意，在这种情形下
$n_{X/Y}(y) = \sum_i m_i$；略去一些细节。由于 $f$ 普遍开，对所有
$i, j$，态射 $V_{i, j} \to U$ 都是开态射。因此缩小 $U$ 后，
可以假设对所有 $i, j$，$V_{i, j} \to U$ 都满射。这证明
$n_{U \times_Y X/U} \geq \sum_i m_i = n_{X/Y}(y) \geq d$.
由于 $n_{X/Y}$ 的构造与基变换相容，证明完成。
\end{proof}

\begin{lemma}
\label{more-morphisms-lemma-large-open}
设 $f : X \to Y$ 是分离、局部拟有限且普遍开的概形态射。设
$n_{X/Y}$ 如 引理 \ref{more-morphisms-lemma-count-geometric-fibres} 中所述。
若 $n_{X/Y}$ 取得最大值 $d < \infty$，则集合
$$
Y_d = \{y \in Y \mid n_{X/Y}(y) = d\}
$$
在 $Y$ 中开，并且态射 $f^{-1}(Y_d) \to Y_d$ 有限。
\end{lemma}

\begin{proof}
$Y_d$ 的开性直接由
引理 \ref{more-morphisms-lemma-count-geometric-fibres-universally-open} 得出。
为证明在 $Y_d$ 上有限，重复该引理证明中的论证。具体地，令
$y \in Y_d$。则 $X$ 中至多有 $d$ 个点位于 $y$ 上。设
$x_1, \ldots, x_n$ 为 $X$ 中位于 $y$ 上的点。应用
引理 \ref{more-morphisms-lemma-etale-splits-off-quasi-finite-part-technical-variant}，
得到 \'etale 邻域 $(U, u) \to (Y, y)$ 以及分解
$$
U \times_Y X =
W \amalg
\ \coprod\nolimits_{i = 1, \ldots, n}
\ \coprod\nolimits_{j = 1, \ldots, m_i}
V_{i, j}
$$
如所引之处所述。注意，在这种情形下
$d = n_{X/Y}(y) = \sum_i m_i$；略去一些细节。由于 $f$ 普遍开，
对所有 $i, j$，态射 $V_{i, j} \to U$ 都是开态射。因此缩小 $U$ 后，
可以假设对所有 $i, j$，$V_{i, j} \to U$ 都满射，并且可以假设
% P05-EN-0161: frozen source says U maps into W, although W is a summand of U times_Y X and no such map is defined; likely Y_d was intended.
$U$ 映入 $W$。这证明
$n_{U \times_Y X/U} \geq \sum_i m_i = d$。由于 $n_{X/Y}$ 的构造
与基变换相容，有 $n_{U \times_Y X/U} = d$。这意味着 $W$ 必为空，
从而 $U \times_Y X \to U$ 有限。由下降，Lemma
\ref{descent-lemma-descending-property-finite}，这蕴含 $X \to Y$
在开态射 $U \to Y$ 的像上有限。换言之，$f$ 在 $y$ 的某个开邻域上有限，
正合所需。
\end{proof}






\section{赋权}
\label{more-morphisms-section-weightings}

\noindent
本节材料取自 \cite[Exposee XVII, 6.2.4]{SGA4}。

\medskip\noindent
设 $\pi : U \to V$ 是具有有限纤维的局部拟有限概形态射。给定函数
$w : U \to \mathbf{Z}$，定义函数
$$
\textstyle{\int}_\pi w : V \longrightarrow \mathbf{Z},\quad
v \longmapsto
\sum\nolimits_{u \in U,\ \pi(u) = v} w(u) [\kappa(u) : \kappa(v)]_s
$$
注意，这些域扩张有限（态射，Lemma
\ref{morphisms-lemma-residue-field-quasi-finite}），
$[\kappa' : \kappa]_s$ 表示可分次数（域，Definition
\ref{fields-definition-insep-degree}），并且因为假设 $\pi$ 的纤维有限，
该和为有限和。还可以按如下方式计算 $\int_\pi w$ 在点 $v \in V$
处的值。选择代数闭域 $k$ 以及像为 $v$ 的态射
$\overline{v} : \Spec(k) \to V$。则有
$$
(\textstyle{\int}_\pi w)(v) =
\sum\nolimits_{\overline{u} \in U_{\overline{v}}} w(\overline{u})
$$
这里，$w(\overline{u})$ 当然表示 $w$ 在点 $\overline{u}$ 沿态射
$U_{\overline{v}} \to U$ 的像 $u$ 处的值。注意，可以把
$\overline{u} \in U_{\overline{v}}$ 看成满足
$\pi \circ \overline{u} = \overline{v}$ 的态射
$\overline{u} : \Spec(k) \to U$。事实上，由于 $U \to V$
局部拟有限且纤维有限，概形 $U_{\overline{v}}$ 是 $k$ 上有限维代数的谱，
其所有素理想都是剩余域为 $k$ 的极大理想。为验证该等式，注意由域，
引理 \ref{fields-lemma-separable-degree}，位于给定 $u$ 上的态射
$\overline{u}$ 的数目等于 $[\kappa(u) : \kappa(v)]_s$。

\begin{lemma}
\label{more-morphisms-lemma-weighting-check-after-etale-base-change}
给定笛卡尔方块
$$
\xymatrix{
U \ar[d]_\pi & U' \ar[l]^h \ar[d]^{\pi'} \\
V & V' \ar[l]_g
}
$$
其中 $\pi$ 局部拟有限且纤维有限，并给定函数
$w : U \to \mathbf{Z}$，则有
$(\int_\pi w) \circ g = \int_{\pi'} (w \circ h)$。
\end{lemma}

\begin{proof}
这直接由上面对 $\int_\pi w$ 的第二种描述得出。若从定义证明，
需使用如下事实：若 $E/F$ 是域的有限扩张，$F'/F$ 是另一个域扩张，
并把 $(E \otimes_F F')_{red} = \prod E'_i$ 写成 $F'$ 上有限域的乘积，
则有
$$
[E : F]_s = \sum [E'_i : F']_s
$$
为证明该等式，选择代数闭域扩张 $\Omega/F'$，并注意
\begin{align*}
[E : F]_s
& =
|\Mor_F(E, \Omega)| \\
& =
|\Mor_{F'}(E \otimes_F F', \Omega)| \\
& =
|\Mor_{F'}((E \otimes_F F')_{red}, \Omega)| \\
& =
\sum |\Mor_{F'}(E'_i, \Omega)| \\
& =
\sum [E'_i : F']_s
\end{align*}
这里使用了域，引理 \ref{fields-lemma-separable-degree}。
\end{proof}

\begin{definition}
\label{more-morphisms-definition-weighting}
设 $f : X \to Y$ 是局部拟有限态射。$f$ 的一个{\it 赋权}，
或一个 {\it pond\'eration}，是映射 $w : X \to \mathbf{Z}$，
使得对任意图
$$
\xymatrix{
X \ar[d]_f & U \ar[l]^h \ar[d]^\pi \\
Y & V \ar[l]_g
}
$$
若 $V \to Y$ 是 \'etale 的，$U \subset X_V$ 是开集，且 $U \to V$ 有限，
则函数 $\int_\pi (w \circ h)$ 局部常值。
\end{definition}

\noindent
当然，取 $w = 0$，便得到任意局部拟有限态射 $f$ 的一个赋权，
尽管它并不十分有趣。事实将表明，对各种用途而言，最有趣的是
{\it 正}赋权，即 $w : X \to \mathbf{Z}_{> 0}$。

\begin{lemma}
\label{more-morphisms-lemma-weighting-base-change}
设 $f : X \to Y$ 是局部拟有限态射，$w : X \to \mathbf{Z}$ 是一个赋权。
设 $f' : X' \to Y'$ 是 $f$ 沿态射 $Y' \to Y$ 的基变换。
则 $w$ 与投影 $X' \to X$ 的复合 $w' : X' \to \mathbf{Z}$ 是 $f'$ 的赋权。
\end{lemma}

\begin{proof}
考虑图表
$$
\xymatrix{
X' \ar[d]_{f'} & U' \ar[l]^{h'} \ar[d]^{\pi'} \\
Y' & V' \ar[l]_{g'}
}
$$
它是 定义 \ref{more-morphisms-definition-weighting} 中针对态射 $f'$ 的图。
对任意 $v' \in V'$，需证明 $\int_{\pi'} (w' \circ h')$ 在 $v'$
的某个开邻域中常值。由
引理 \ref{more-morphisms-lemma-weighting-check-after-etale-base-change}
（以及 \'etale 态射是开态射这一事实），可以用 $v'$ 的任意 \'etale
邻域代替 $V'$。以 $v'$ 的某个 \'etale 邻域代替 $V'$ 后，
可以假设 $U' = U'_1 \amalg \ldots \amalg U'_n$，其中每个 $U'_i$
都恰有一点 $u'_i$ 位于 $v'$ 上，并且 $\kappa(u'_i)/\kappa(v')$
纯不可分，见
引理 \ref{more-morphisms-lemma-etale-splits-off-quasi-finite-part-technical-variant}。
显然，只需证明 $\int_{U'_i \to V'} w'|_{U'_i}$ 在 $v'$
的某个邻域中常值。这把我们归约到下一段讨论的情形。

\medskip\noindent
我们有 $v' \in V'$，且 $U'$ 中恰有一点 $u'$ 位于 $v'$ 上，
并满足 $\kappa(u')/\kappa(v')$ 纯不可分。用 $x \in X$ 和 $y \in Y$
分别表示 $u'$ 和 $v'$ 的像。可以找到 \'etale 邻域
$(V, v) \to (Y, y)$ 以及开集 $U \subset X_V$，使得
$\pi : U \to V$ 有限，并且其中恰有一点 $u \in U$ 位于 $v$ 上；
该点沿投影 $h : U \to X$ 映到 $x \in X$，而且
$\kappa(u)/\kappa(v)$ 纯不可分。由上面使用的引理，这可以做到。
考虑态射
$$
U'' = U \times_X U' \longrightarrow V \times_Y V' = V''
$$
由于 $u$ 和 $u'$ 都映到 $x \in X$，存在映到 $(u, u')$ 的点
$u'' \in U''$。用 $v'' \in V''$ 表示 $u''$ 的像。以 $V'', v''$
代替 $V', v'$ 后，可以假设复合 $V' \to Y' \to Y$ 分解经过 \'etale
邻域的映射 $(V', v') \to (V, v)$，并且所诱导的态射
$X'_{V'} = X_{V'} \to X_V$ 把 $u'$ 映到 $u$。在基变换
$X'_{V'} = X_{V'}$ 中有两个开子概形，即 $U'$ 以及
$U \subset X_V$ 的原像 $U_{V'}$。按构造，二者都恰有一个点位于
$v'$ 上，而且对二者该点都是 $u'$。因此缩小 $V'$ 后，可以假设这两个
开子概形相同；事实上，$U' \setminus (U' \cap U_{V'})$ 和
$U_{V'} \setminus (U' \cap U_{V'})$ 在 $V'$ 中有闭像，而这些像不含
$v'$。因此 $U' = U_{V'}$，从而得到
引理 \ref{more-morphisms-lemma-weighting-check-after-etale-base-change} 中的笛卡尔图。
由假设，$\int_\pi (w \circ h)$ 局部常值，故结论成立。
\end{proof}

\begin{lemma}
\label{more-morphisms-lemma-weighting-open}
设 $f : X \to Y$ 是局部拟有限态射，$w : X \to \mathbf{Z}$ 是 $f$
的赋权。若 $X' \subset X$ 是开集，则 $w|_{X'}$ 是
$f|_{X'} : X' \to Y$ 的赋权。
\end{lemma}

\begin{proof}
直接由定义得出。
\end{proof}

\begin{lemma}
\label{more-morphisms-lemma-weighting-composition}
设 $f : X \to Y$ 和 $g : Y \to Z$ 是局部拟有限态射。设
$w_f : X \to \mathbf{Z}$ 是 $f$ 的赋权，而
$w_g : Y \to \mathbf{Z}$ 是 $g$ 的赋权。则函数
$$
X \longrightarrow \mathbf{Z},\quad
x \longmapsto w_f(x) w_g(f(x))
$$
是 $g \circ f$ 的赋权。
\end{lemma}

\begin{proof}
对 $x \in X$，置 $w_{g \circ f}(x) = w_f(x) w_g(f(x))$。考虑图
$$
\xymatrix{
X \ar[d]_{g \circ f} & U \ar[l] \ar[d]^\pi \\
Z & W \ar[l]
}
$$
其中 $W \to Z$ 是 \'etale 的，$U \subset X_W$ 是开集，且 $U \to W$ 有限。
我们要证明 $\int_\pi w_{g \circ f}|_U$ 局部常值。选择一点 $w \in W$。
由 引理 \ref{more-morphisms-lemma-weighting-check-after-etale-base-change}
（以及 \'etale 态射是开态射这一事实），只需证明用一个 \'etale
邻域代替 $(W, w)$ 后，$\int_\pi w_{g \circ f}|_U$ 常值。
以一个 \'etale 邻域代替 $(W, w)$ 后，可以假设
$U = U_1 \amalg \ldots \amalg U_n$，其中每个 $U_i$ 都恰有一点 $u_i$
位于 $w$ 上，并且 $\kappa(u_i)/\kappa(w)$ 纯不可分，见
引理 \ref{more-morphisms-lemma-etale-splits-off-quasi-finite-part-technical-variant}。
显然，只需证明 $\int_{U_i \to W} w_{g \circ f}|_{U_i}$ 在 $w$
的某个 \'etale 邻域中常值。这把我们归约到下一段讨论的情形。

\medskip\noindent
我们有 $w \in W$，且该开集中恰有一点 $u \in U$ 位于 $w$ 上，
并满足 $\kappa(u)/\kappa(w)$ 纯不可分。考虑点 $v = f(u) \in Y$。
用一个初等 \'etale 邻域代替 $(W, w)$ 后，可以假设存在 $v$
的开邻域 $V \subset Y_W$，使得 $V \to W$ 有限，见
引理 \ref{more-morphisms-lemma-etale-makes-quasi-finite-finite-at-point}。于是
$f_W^{-1}(V) \cap U$ 是 $u$ 的开邻域，其中
$f_W : X_W \to Y_W$ 是 $f$ 到 $W$ 的基变换。因此在 Zariski 意义下
缩小 $W$ 后，可以假设 $f_W(U) \subset V$。于是得到态射
$$
U \xrightarrow{a} V \xrightarrow{b} W
$$
由于 $V \to W$ 分离（因为它有限），$U \to V$ 有限。由于 $w_f$ 和
$w_g$ 分别是 $f$ 和 $g$ 的赋权，$\int_a w_f|_U$ 在 $V$ 上局部常值，
而 $\int_b w_g|_V$ 在 $W$ 上局部常值。因此再次缩小 $W$ 后，
可以假设这两个函数常值，设其值分别为 $n$ 和 $m$。于是立即可知
$\int_\pi w_{g \circ f}|_U = \int_{b \circ a} w_{g \circ f}|_U$
常值为 $nm$，正合所需。
\end{proof}

\begin{lemma}
\label{more-morphisms-lemma-weighting-universally-open}
设 $f : X \to Y$ 是局部拟有限态射，$w : X \to \mathbf{Z}$ 是一个赋权。
若对所有 $x \in X$ 都有 $w(x) > 0$，则 $f$ 普遍开。
\end{lemma}

\begin{proof}
该性质在基变换下保持，见 引理 \ref{more-morphisms-lemma-weighting-base-change}，
故只需证明 $f$ 是开态射。又因为可以用 $X$ 的任一开集代替 $X$，
只需证明 $f(X)$ 为开集。令 $y \in f(X)$。选择 $x \in X$，使得
$f(x) = y$。只需证明 $f(X)$ 包含 $y$ 的一个开邻域，并且只需在用
$y$ 的一个 \'etale 邻域代替 $Y$ 后证明这一点。通过拟有限态射的
\'etale 局部化，见 第 \ref{more-morphisms-section-etale-localization} 节，可以假设存在
$x$ 的开邻域 $U \subset X$，使得 $\pi = f|_U : U \to Y$ 有限。
于是 $\int_\pi w|_U$ 局部常值，并且在 $y$ 处取正值。因此
$\pi(U)$ 包含 $y$ 的一个开邻域，证明完成。
\end{proof}

\begin{lemma}
\label{more-morphisms-lemma-weighting-flat-quasi-finite}
设 $f : X \to Y$ 是概形态射。假设 $f$ 局部拟有限、局部有限表示且平坦。
则 $f$ 有一个正赋权 $w : X \to \mathbf{Z}_{> 0}$，它把位于
$y \in Y$ 上的 $x \in X$ 映到
$$
w(x) =
\text{length}_{\mathcal{O}_{X, x}}
(\mathcal{O}_{X, x}/\mathfrak m_y \mathcal{O}_{X, x})
[\kappa(x) : \kappa(y)]_i
$$
其中 $[\kappa' : \kappa]_i$ 是不可分次数
（域，定义 \ref{fields-definition-insep-degree}）。
\end{lemma}

\begin{proof}
考虑 定义 \ref{more-morphisms-definition-weighting} 中的图。设 $u \in U$
在 $X, Y, V$ 中的像分别为 $x, y, v$。我们断言
$$
\text{length}_{\mathcal{O}_{X, x}}
(\mathcal{O}_{X, x}/\mathfrak m_y \mathcal{O}_{X, x}) =
\text{length}_{\mathcal{O}_{U, u}}
(\mathcal{O}_{U, u}/\mathfrak m_v \mathcal{O}_{U, u})
$$
及
$$
[\kappa(x) : \kappa(y)]_i =
[\kappa(u) : \kappa(v)]_i
$$
第一个等式成立，是因为
$\mathcal{O}_{X, x} \to \mathcal{O}_{U, u}$ 是平坦局部同态，满足
$\mathfrak m_y \mathcal{O}_{U, u} = \mathfrak m_v \mathcal{O}_{U, u}$ 和
$\mathfrak m_x \mathcal{O}_{U, u} = \mathfrak m_u$（因为
$\mathcal{O}_{Y, y} \to \mathcal{O}_{V, v}$ 和
$\mathcal{O}_{X, x} \to \mathcal{O}_{U, u}$ 都非分歧），再应用代数，
引理 \ref{algebra-lemma-pullback-module} 即得该等式。第二个等式成立，
是因为 $\kappa(v)/\kappa(y)$ 是有限可分扩张，而 $\kappa(u)$ 是
$\kappa(x) \otimes_{\kappa(y)} \kappa(v)$ 的一个因子，故不可分次数不变。
由此可见，引理中的函数之形成与 \'etale 基变换相容。这把问题归约到
下一段的讨论。

\medskip\noindent
假设 $f$ 是有限、平坦且有限表示的态射。我们要证明 $\int_f w$
在 $Y$ 上局部常值。事实上，$f$ 有限局部自由（态射，Lemma
\ref{morphisms-lemma-finite-flat}），而我们将证明 $\int_f w$ 等于 $f$
的次数（它是 $Y$ 上的局部常值函数）。具体地，对 $y \in Y$ 有
\begin{align*}
(\textstyle{\int}_f w)(y)
& =
\sum\nolimits_{f(x) = y}
\text{length}_{\mathcal{O}_{X, x}}
(\mathcal{O}_{X, x}/\mathfrak m_y \mathcal{O}_{X, x})
[\kappa(x) : \kappa(y)]_i
[\kappa(x) : \kappa(y)]_s \\
& =
\sum\nolimits_{f(x) = y}
\text{length}_{\mathcal{O}_{X, x}}
(\mathcal{O}_{X, x}/\mathfrak m_y \mathcal{O}_{X, x})
[\kappa(x) : \kappa(y)] \\
& =
\text{length}_{\mathcal{O}_{Y, y}}((f_*\mathcal{O}_X)_y/
\mathfrak m_y (f_*\mathcal{O}_X)_y)
\end{align*}
最后一个等号由代数，引理 \ref{algebra-lemma-pushdown-module} 得出。
最后所得数正是 $f_*\mathcal{O}_X$ 在 $y$ 处的秩，正合所需。
\end{proof}

\begin{lemma}
\label{more-morphisms-lemma-weighting-quasi-finite-Noetherian}
设 $f : X \to Y$ 是概形态射。假设
\begin{enumerate}
\item $f$ 局部拟有限；
\item $Y$ 几何单分支且局部诺特。
\end{enumerate}
则存在赋权 $w : X \to \mathbf{Z}_{\geq 0}$，它把位于 $y \in Y$
上的 $x \in X$ 映到 $\mathcal{O}_{X, x}^{sh}$ 在
$\mathcal{O}_{Y, y}^{sh}$ 上的“泛可分次数”。
\end{lemma}

\begin{proof}
由代数，引理 \ref{algebra-lemma-quasi-finite-strict-henselization}，
$\mathcal{O}_{Y, y}^{sh} \to \mathcal{O}_{X, x}^{sh}$ 有限。由于 $Y$
几何单分支，$\mathcal{O}_{Y, y}^{sh}$ 中有唯一的极小素理想
$\mathfrak p$，见代数进阶，Lemma
\ref{more-algebra-lemma-geometrically-unibranch}。写成
$$
(\kappa(\mathfrak p) \otimes_{\mathcal{O}_{Y, y}^{sh}}
\mathcal{O}_{X, x}^{sh})_{red} =
\prod K_i
$$
其中右侧是域的有限乘积。置
$w(x) = \sum [K_i : \kappa(\mathfrak p)]_s$。

\medskip\noindent
该定义显然不受 \'etale 局部化影响。因此为证明 $w$ 是赋权，
只需证明当 $f$ 是有限态射时，$\int_f w$ 局部常值。注意，
$\int_f w$ 在 $Y$ 的泛点 $\eta$ 处的值，就是 $X \to Y$ 在 $\eta$
上的几何纤维 $X_{\overline{\eta}}$ 的点数。此外，由于 $Y$ 单分支，
$Y$ 的一点 $y$ 是唯一一个泛点 $\eta$ 的特化。因此只需证明
$(\int_f w)(y)$ 等于 $X_{\overline{\eta}}$ 的点数。转到 $y$
的仿射邻域后，可以假设 $X \to Y$ 由有限环映射 $A \to B$ 给出。
假设 $\mathcal{O}_{Y, y}^{sh}$ 是利用到代数闭域 $k$ 的映射
$\kappa(y) \to k$ 构造的。则
$$
\mathcal{O}_{Y, y}^{sh} \otimes_A B =
\prod\nolimits_{f(x) = y}
\prod\nolimits_{\varphi \in \Mor_{\kappa(y)}(\kappa(x), k)}
\mathcal{O}_{X, x}^{sh}
$$
这是由代数，引理 \ref{algebra-lemma-finite-over-henselian} 和上面使用的
引理得到的。注意，$\mathcal{O}_{Y, y}^{sh}$ 的极小素理想 $\mathfrak p$
映到 $A$ 中对应于 $\eta$ 的素理想。因此所需等式成立，因为几何纤维的
点数在域扩张下不变。
\end{proof}

\begin{lemma}
\label{more-morphisms-lemma-weighting-specialization}
设 $f : X \to Y$ 是局部拟有限概形态射，$w : X \to \mathbf{Z}$ 是 $f$
的赋权。设 $y \leadsto y'$ 是 $Y$ 中点的特化。假设对所有满足
$x \in X$ 且 $f(x) = y$ 的点都有 $w(x) = 0$。则对所有满足
$x' \in X$ 且 $f(x') = y'$ 的点都有 $w(x') = 0$。
\end{lemma}

\begin{proof}
令 $x' \in X$ 为映到 $y'$ 的点。由
引理 \ref{more-morphisms-lemma-etale-makes-quasi-finite-finite-at-point}，可以选择图
$$
\xymatrix{
U \ar[r] \ar[dr]_\pi & X_V \ar[d] \ar[r] & X \ar[d]^f \\
& V \ar[r]^\varphi & Y
}
$$
其中 $V \to Y$ 是 \'etale 的，$v' \in V$ 映到 $y'$，
% P05-EN-0162: frozen source says U is open in Y_V, but the diagram and quasi-finite localization require U open in X_V.
$U \subset Y_V$ 是开集，$U \to V$ 有限，并且恰有一点 $u' \in U$
映到 $v'$，且该点在 $X$ 中映到 $x'$。注意，由于 $V \to Y$ 是开态射，
存在点 $v \in V$ 特化到 $v'$ 并映到 $y$。由 Definition
\ref{more-morphisms-definition-weighting}，函数 $\int_\pi (w \circ h)$ 局部常值，
其中 $h : U \to X$ 是上方两个水平箭头的复合。由对 $w$ 的假设，
$\int_\pi (w \circ h)$ 在 $v$ 处为零，故在 $v'$ 处也为零。然而，
它在 $v'$ 处的值是 $w(x')[\kappa(u'):\kappa(v')]_s$，结论随即得出。
\end{proof}




\section{赋权的进一步性质}
\label{more-morphisms-section-more-weightings}

\noindent
我们再证明赋权的若干基本性质。虽然乍看之下赋权可能表现得十分不规则，
但事实上定义 \ref{more-morphisms-definition-weighting} 中所施加的条件相当强。

\begin{lemma}
\label{more-morphisms-lemma-jumps-w}
设 $f : X \to Y$ 为局部拟有限态射，
$w : X \to \mathbf{Z}$ 为 $f$ 的一个赋权。则函数 $w$ 的各水平集在 $X$ 中局部可构造。
\end{lemma}

\begin{proof}
在以下证明中，我们将不再另行说明地使用引理 \ref{more-morphisms-lemma-weighting-open} 和
\ref{more-morphisms-lemma-weighting-base-change}。我们还将使用概形及拓扑空间的可构造子集的
初等性质，见拓扑章第 \ref{topology-section-constructible} 节和
性质章第 \ref{properties-section-constructible} 节。由此读者可见，该问题关于
$X$ 和 $Y$ 都是局部的；细节从略。因此可设 $X$ 和 $Y$ 都是仿射的。
如果能够找到一个有限呈现满态射 $Y' \to Y$，使得 $w$ 的各水平集在
$X' = Y' \times_Y X$ 中的拉回均为局部可构造子集，那么结论即由
态射章定理 \ref{morphisms-theorem-chevalley} 得出。

\medskip\noindent
设 $X$ 和 $Y$ 为仿射概形。由引理
\ref{more-morphisms-lemma-quasi-finite-separated-pass-through-finite}，可选取一个浸入
$X \to T$，其中 $T \to Y$ 为有限态射。根据态射章引理
\ref{morphisms-lemma-massage-finite}，将 $Y$ 替换为一个在 $Y$ 上有限局部自由的
满态射 $Y'$，将 $X$ 替换为 $Y' \times_Y X$，并将 $T$ 替换为一个在 $Y'$ 上
有限局部自由且以 $Y' \times_Y T$ 为闭子概形的概形之后，可设 $T$ 在 $Y$ 上
有限局部自由，并且包含与 $Y$ 同构地映射的闭子概形 $T_i$，使得
$T = \bigcup_{i = 1, \ldots, n} T_i$（集合论意义下）。由于 $T_i \subset T$
是可构造闭子集（它是概形的一个有限呈现态射 $Y \to T$ 的像），故对于
$I \subset \{1, \ldots, n\}$，交集 $\bigcap_{i \in I} T_i$ 是 $T$ 的可构造闭子集，
因而映为 $Y$ 的可构造闭子集。

\medskip\noindent
对于一个各部分均非空的不交并分解
$\{1, \ldots, n\} = I_1 \amalg \ldots \amalg I_r$，考虑子集
$Y_{I_1, \ldots, I_r} \subset Y$，其元素为满足如下条件的点 $y \in Y$：
$T_y = \{x_1, \ldots, x_r\}$ 恰由 $r$ 个点组成，且
$x_j \in T_i \Leftrightarrow i \in I_j$。由以上说明，这些子集构成 $Y$ 的一个
可构造分拆。由代数章引理 \ref{algebra-lemma-constructible-is-image}，存在一个
在 $Y$ 上有限呈现的仿射概形 $Y'$，使得 $Y' \to Y$ 的像恰为
$Y_{I_1, \ldots, I_r}$。因此可设对于某个不交并分解
$\{1, \ldots, n\} = I_1 \amalg \ldots \amalg I_r$，有
$Y = Y_{I_1, \ldots, I_r}$。此时，令 $T(j) = \bigcap_{i \in I_j} T_i$，则
$T = T(1) \amalg \ldots \amalg T(r)$ 是 $T$ 分解为两两不交闭子集（因而也是
开子集）的分解。与局部闭子概形 $X$ 相交，得到类似的开闭分解
$X = X(1) \amalg \ldots \amalg X(r)$。态射 $X(j) \to Y$ 为浸入。
由于 $w$ 是赋权，$w|_{X(j)}$ 局部常值\footnote{事实上，若
$f : X \to Y$ 为浸入且 $w$ 是 $f$ 的一个赋权，则在 $w$ 非零的轨迹上，
$f$ 的限制为开映射。}，结论随即成立。
\end{proof}

\begin{lemma}
\label{more-morphisms-lemma-weighting-on-dense-set}
设 $f : X \to Y$ 为概形的局部拟有限态射，
$w : X \to \mathbf{Z}$ 为 $f$ 的一个赋权，$E \subset X$ 为稠密子集。
若对所有 $x \in E$ 均有 $w(x) = 0$，则 $w = 0$。
\end{lemma}

\begin{proof}
建议先阅读引理 \ref{more-morphisms-lemma-weighting-specialization}。
取 $x \in X$，其像为 $y \in Y$。我们将证明 $w(x) = 0$。由引理
\ref{more-morphisms-lemma-etale-makes-quasi-finite-finite-at-point}，可选取图表
$$
\xymatrix{
U \ar[r] \ar[dr]_\pi & X_V \ar[d] \ar[r] & X \ar[d]^f \\
& V \ar[r]^\varphi & Y
}
$$
其中 $V \to Y$ 为 \'etale 态射，点 $v \in V$ 映到 $y$，
开子集 $U \subset Y_V$ 满足 $U \to V$ 为有限态射，并且存在唯一一点
$u \in U$ 映到 $v$，且它在 $X$ 中进一步映到 $x$。我们可以并且确实设
$V$（从而 $U$）为仿射的。由定义 \ref{more-morphisms-definition-weighting}，函数
$\int_\pi (w \circ h)$ 局部常值，其中 $h : U \to X$ 是上方横向箭头的复合。
由引理 \ref{more-morphisms-lemma-jumps-w}，$w \circ h$ 为 $0$ 的点集 $U_0$ 是 $U$ 的
可构造子集；又因它包含 $h^{-1}(E)$（使用 $h$ 为开映射），故在 $U$ 中稠密。
由拓扑章引理 \ref{topology-lemma-dense-in-constructible}，$U_0$ 包含一个稠密
开子集；换言之，$w \circ h$ 非零的点集包含在某个无处稠密闭子集
$Z \subset U$ 中。于是由态射章引理
\ref{morphisms-lemma-image-nowhere-dense-finite}，$\pi(Z)$ 是 $V$ 的无处稠密
闭子集。因此 $\int_\pi (w \circ h)$ 在一个稠密开子集上为零，从而处处为零。
然而它在 $v$ 处的值为 $w(x)[\kappa(u):\kappa(v)]_s$，结论得证。
\end{proof}

\begin{lemma}
\label{more-morphisms-lemma-jumps-int-w}
设 $f : X \to Y$ 为有限呈现的局部拟有限态射，
$w : X \to \mathbf{Z}$ 为 $f$ 的一个赋权。则函数 $\int_f w$ 的各水平集在
$Y$ 中局部可构造。
\end{lemma}

\begin{proof}
由引理 \ref{more-morphisms-lemma-weighting-check-after-etale-base-change}，函数 $\int_f w$ 的构造
与任意基变换交换；由引理 \ref{more-morphisms-lemma-weighting-base-change}，基变换后仍得到赋权。
这意味着，若能找到一个有限呈现满态射 $Y' \to Y$，则只须在基变换到 $Y'$
之后证明结论，见态射章定理 \ref{morphisms-theorem-chevalley}。

\medskip\noindent
该问题在 $Y$ 上是局部的，故可设 $Y$ 为仿射概形。于是 $X$ 拟紧且拟分离
（因为 $f$ 为有限呈现）。设 $X = U \cup V$，其中两个开子集均拟紧。
则有
$$
\textstyle{\int}_f w =
\textstyle{\int}_{f|_U} w|_U +
\textstyle{\int}_{f|_V} w|_V -
\textstyle{\int}_{f|_{U \cap V}} w|_{U \cap V}
$$
因此，若结论对 $w|_U$、$w|_V$ 和 $w|_{U \cap V}$ 成立，则它也对 $w$ 成立。
由归纳原理（概形上同调章引理 \ref{coherent-lemma-induction-principle}），
只须在 $X$ 为仿射概形时证明本引理。

\medskip\noindent
设 $X$ 和 $Y$ 都是仿射概形。由引理
\ref{more-morphisms-lemma-quasi-finite-separated-pass-through-finite}，可选取一个开浸入
$X \to T$，其中 $T \to Y$ 为有限态射。因为仍可沿适当的 $Y' \to Y$ 作基变换，
所以可使用态射章引理 \ref{morphisms-lemma-massage-finite}，将问题约化到如下情形：
态射 $T \to Y$ 所诱导的所有剩余域扩张（因而当然也包括 $X \to Y$ 所诱导的）
均为平凡扩张。在此情形下，$\int_f w$ 就是对每条纤维中的 $w$ 值求和。
由引理 \ref{more-morphisms-lemma-jumps-w}，$w$ 的各水平集在 $X$ 中局部可构造；由性质章引理
\ref{properties-lemma-stratification-locally-finite-constructible}，函数 $w$ 仅取
有限多个值。因此结论由态射章定理 \ref{morphisms-theorem-chevalley} 以及关于整数
之和的若干初等论证得出。
\end{proof}

\begin{lemma}
\label{more-morphisms-lemma-semicontinuous-w}
设 $f : X \to Y$ 为局部拟有限态射，
$w : X \to \mathbf{Z}_{> 0}$ 为 $f$ 的一个正赋权。则 $w$ 上半连续。
\end{lemma}

\begin{proof}
取 $x \in X$，其像为 $y \in Y$。选取一个 \'etale 邻域
$(V, v) \to (Y, y)$ 以及开子集 $U \subset X_V$，使得
$\pi : U \to V$ 为有限态射，并且存在唯一一点 $u \in U$ 映到 $v$，且
$\kappa(u)/\kappa(v)$ 为纯不可分扩张。见引理
\ref{more-morphisms-lemma-etale-makes-quasi-finite-finite-multiple-points-var}。于是
$(\int_\pi w|_U)(v) = w(u)$。由定义 \ref{more-morphisms-definition-weighting}，将 $V$ 替换为
$v$ 的一个邻域后，对所有 $u' \in U$ 均有
$w|_U(u') \leq w|_U(u) = w(x)$。事实上，$w|_U(u')$ 是表达式
$(\int_\pi w|_U)(\pi(u'))$ 中的一个加项。由于 \'etale 态射 $U \to X$ 为开映射，
这就证明了本引理。
\end{proof}

\begin{lemma}
\label{more-morphisms-lemma-semicontinuous-int-w}
设 $f : X \to Y$ 为纤维有限的分离局部拟有限态射，
$w : X \to \mathbf{Z}_{> 0}$ 为 $f$ 的一个正赋权。则 $\int_f w$ 下半连续。
\end{lemma}

\begin{proof}
取 $y \in Y$，设 $x_1, \ldots, x_r \in X$ 为位于 $y$ 上方的点。应用引理
\ref{more-morphisms-lemma-etale-splits-off-quasi-finite-part-technical-variant}，得到一个 \'etale 邻域
$(U, u) \to (Y, y)$ 以及分解
$$
U \times_Y X =
W \amalg
\ \coprod\nolimits_{i = 1, \ldots, n}
\ \coprod\nolimits_{j = 1, \ldots, m_i}
V_{i, j}
$$
如所引之处。注意 $(\int_f w)(y) = \sum w(v_{i, j})$，其中
$w(v_{i, j}) = w(x_i)$。依定义，$\int_{V_{i, j} \to U} w|_{V_{i, j}}$ 局部常值，
故缩小 $U$ 后可设这些函数均为常值函数，其值为 $w(v_{i, j})$。于是
$$
\textstyle{\int}_{U \times_Y X \to U} w|_{U \times_Y X} =
\textstyle{\int}_{W \to U} w|_W +
\sum \textstyle{\int}_{V_{i, j} \to U} w|_{V_{i, j}} =
\textstyle{\int}_{W \to U} w|_W + (\int_f w)(y)
$$
该值 $\geq (\int_f w)(y)$。又因 $U \to Y$ 为开映射，且积分的构造与基变换交换
（引理 \ref{more-morphisms-lemma-weighting-check-after-etale-base-change}），结论得证。
\end{proof}

\begin{lemma}
\label{more-morphisms-lemma-max-int-w}
设 $f : X \to Y$ 为局部拟有限态射，其中 $X$ 拟紧；设
$w : X \to \mathbf{Z}$ 为 $f$ 的一个赋权。则 $\int_f w$ 取得其最大值。
\end{lemma}

\begin{proof}
由引理 \ref{more-morphisms-lemma-jumps-w} 和性质章引理
\ref{properties-lemma-stratification-locally-finite-constructible}，$w$ 在 $X$ 上仅取
有限多个值。由态射章引理
\ref{morphisms-lemma-locally-quasi-finite-qc-source-universally-bounded}，
$X \to Y$ 的几何纤维一致有界。这表明 $\int_f w$ 有界。
\end{proof}

\begin{lemma}
\label{more-morphisms-lemma-max-int-finite}
设 $f : X \to Y$ 为分离局部拟有限态射，
$w : X \to \mathbf{Z}_{> 0}$ 为 $f$ 的一个正赋权。假设 $\int_w f$ 取得最大值
$d$，并令 $Y_d \subset Y$ 为满足 $(\int_f w)(y) = d$ 的点 $y$ 所组成的开集。
则态射 $f^{-1}(Y_d) \to Y_d$ 为有限态射。
\end{lemma}

\begin{proof}
由引理 \ref{more-morphisms-lemma-semicontinuous-int-w}，$Y_d$ 为开集。取 $y \in Y_d$，设
$x_1, \ldots, x_n$ 为 $X$ 中位于 $y$ 上方的点。应用引理
\ref{more-morphisms-lemma-etale-splits-off-quasi-finite-part-technical-variant}，得到一个 \'etale 邻域
$(U, u) \to (Y, y)$ 以及分解
$$
U \times_Y X =
W \amalg
\ \coprod\nolimits_{i = 1, \ldots, n}
\ \coprod\nolimits_{j = 1, \ldots, m_i}
V_{i, j}
$$
如所引之处。注意 $d = \sum w(v_{i, j})$，其中 $w(v_{i, j}) = w(x_i)$。
依定义，$\int_{V_{i, j} \to U} w|_{V_{i, j}}$ 局部常值，故缩小 $U$ 后可设这些
函数均为常值函数，其值为 $w(v_{i, j})$。于是
$$
\textstyle{\int}_{U \times_Y X \to U} w|_{U \times_Y X} =
\textstyle{\int}_{W \to U} w|_W +
\sum \textstyle{\int}_{V_{i, j} \to U} w|_{V_{i, j}} =
\textstyle{\int}_{W \to U} w|_W + (\int_f w)(y)
$$
该值 $\geq (\int_f w)(y) = d$，故 $W$ 必为空集。因此
$U \times_Y X \to U$ 为有限态射。由下降章引理
\ref{descent-lemma-descending-property-finite}，这说明 $X \to Y$ 在开态射
$U \to Y$ 的像上为有限态射。换言之，$f$ 在 $y$ 的某个开邻域上为有限态射，
正如所需。
\end{proof}

\begin{lemma}
\label{more-morphisms-lemma-open-and-closed-in-finite}
设 $A \to B$ 为有限且有限呈现的环同态。则存在一个有限呈现环同态
$A \to A_{univ}$ 以及一个幂等元 $e_{univ} \in B \otimes_A A_{univ}$，
使得对于任意环同态 $A \to A'$ 和幂等元 $e \in B \otimes_A A'$，
都存在一个将 $e_{univ}$ 映到 $e$ 的环同态 $A_{univ} \to A'$。
\end{lemma}

\begin{proof}
选取 $b_1, \ldots, b_n \in B$，使其作为 $A$-模生成 $B$。对每个 $i$，选取首一
多项式 $P_i \in A[x]$，使得 $P_i(b_i) = 0$ 在 $B$ 中成立，见代数章引理
\ref{algebra-lemma-finite-is-integral}。于是 $B$ 是有限自由 $A$-代数
$B' = A[x_1, \ldots, x_n]/(P_1(x_1), \ldots, P_n(x_n))$ 的商。令
$J \subset B'$ 为满射 $B' \to B$ 的核。由于 $B$ 是有限生成 $A$-代数，
$J =(f_1, \ldots, f_m)$ 有限生成，见代数章引理
\ref{algebra-lemma-compose-finite-type}。选取 $B'$ 的一个 $A$-基
$b'_1, \ldots, b'_N$。考虑代数
$$
A_{univ} = A[z_1, \ldots, z_N, y_1, \ldots, y_m]/I
$$
其中 $I$ 是由下列表达式关于基元素 $b'_1, \ldots, b'_N$ 的系数在
$A[z_1, \ldots, z_n, y_1, \ldots, y_m]$ 中所生成的理想：
$$
(\sum z_j b'_j)^2 - \sum z_j b'_j + \sum y_k f_k
$$
该表达式位于 $B'[z_1, \ldots, z_N, y_1, \ldots, y_m]$ 中。依构造，元素
$\sum z_j b'_j$ 在代数 $B \otimes_A A_{univ}$ 中映为幂等元 $e_{univ}$。
此外，若 $e \in B \otimes_A A'$ 为幂等元，则可将 $e$ 提升为
$B' \otimes_A A'$ 中形如 $\sum b'_j \otimes a'_j$ 的元素，并且可以找到
$a''_k \in A'$，使得
$$
(\sum b'_j \otimes a'_j)^2 - \sum b'_j \otimes a'_j + \sum f_k \otimes a''_k
$$
在 $B' \otimes_A A'$ 中为零。因此得到一个 $A$-代数同态
$A_{univ} \to A$，它将 $z_j$ 映到 $a'_j$、将 $y_k$ 映到 $a''_k$，并将
$e_{univ}$ 映到 $e$。证明完毕。
\end{proof}

\begin{lemma}
\label{more-morphisms-lemma-open-and-closed-in-quasi-finite}
设 $X \to Y$ 为仿射概形间拟有限且有限呈现的态射。则存在一个有限呈现态射
$Y_{univ} \to Y$ 以及一个开子概形 $U_{univ} \subset Y_{univ} \times_Y X$，
使得 $U_{univ} \to Y_{univ}$ 为有限态射，并具有如下性质：给定任意仿射概形间的
态射 $Y' \to Y$ 以及开子概形 $U' \subset Y' \times_Y X$，若
$U' \to Y'$ 为有限态射，则存在态射 $Y' \to Y_{univ}$，使得 $U_{univ}$ 的逆像
为 $U'$。
\end{lemma}

\begin{proof}
回忆有限型态射为拟有限态射，当且仅当其相对维数为 $0$，见态射章引理
\ref{morphisms-lemma-locally-quasi-finite-rel-dimension-0}。在 $d = 0$ 时应用引理
\ref{more-morphisms-lemma-Noetherian-approximation-dimension-d}，可约化到 $X$ 和 $Y$ 均为
Noether 概形的情形。由代数章引理
\ref{algebra-lemma-quasi-finite-open-integral-closure}，可选取一个开浸入
$X \to X'$，使得 $X' \to Y$ 为有限态射。注意，若有如 (2) 中的 $Y' \to Y$
和 $U'$，则
$$
U' \to Y' \times_Y X \to Y' \times_Y X'
$$
是两个在 $Y'$ 上有限的概形之间的开浸入，因而也是闭浸入。由此可知，$U'$ 对应于
下列代数中的一个幂等元：
$$
\Gamma(Y', \mathcal{O}_{Y'})
\otimes_{\Gamma(Y, \mathcal{O}_Y)}
\Gamma(X', \mathcal{O}_{X'})
$$
该幂等元所对应的开闭子集包含在开集 $Y' \times_Y X$ 中。令
$Y'_{univ} \to Y$ 及幂等元
$$
e'_{univ} \in
\Gamma(Y_{univ}, \mathcal{O}_{Y_{univ}})
\otimes_{\Gamma(Y, \mathcal{O}_Y)}
\Gamma(X', \mathcal{O}_{X'})
$$
为引理 \ref{more-morphisms-lemma-open-and-closed-in-finite} 对环同态
$\Gamma(Y, \mathcal{O}_Y) \to \Gamma(X', \mathcal{O}_{X'})$ 所构造的二元组
（这里使用 $Y$ 为 Noether 概形，以知 $X'$ 在 $Y$ 上有限呈现）。令
$U'_{univ} \subset Y'_{univ} \times_Y X'$ 为相应的开闭子概形。于是
$$
U'_{univ} \setminus Y'_{univ} \times_Y X
$$
是 $U'_{univ}$ 的闭子集，故其像 $T \subset Y'_{univ}$ 为闭集。若令
$Y_{univ} = Y'_{univ} \setminus T$，并令 $U_{univ}$ 为 $U'_{univ}$ 在
$Y_{univ} \times_Y X$ 上的限制，则本引理成立。
\end{proof}

\begin{lemma}
\label{more-morphisms-lemma-descend-weighting}
设 $Y = \lim Y_i$ 为仿射概形的有向极限。取 $0 \in I$，并设
$f_0 : X_0 \to Y_0$ 为仿射概形间拟有限且有限呈现的态射。令
$f : X  \to Y$ 以及对 $i \geq 0$ 的 $f_i : X_i \to Y_i$ 为 $f_0$ 的基变换。
若 $w : X \to \mathbf{Z}$ 为 $f$ 的一个赋权，则当 $i$ 充分大时，存在 $f_i$
的一个赋权 $w_i : X_i \to \mathbf{Z}$，其到 $X$ 的拉回为 $w$。
\end{lemma}

\begin{proof}
由引理 \ref{more-morphisms-lemma-jumps-w}，$w$ 的各水平集是 $X$ 的可构造子集 $E_k$。
由性质章引理 \ref{properties-lemma-stratification-locally-finite-constructible}，
这说明函数 $w$ 仅取有限多个值。因此存在 $i$，使得对所有 $k$，$E_k$ 均下降为
$X_i$ 中的一个可构造子集 $E_{i, k}$；此外，可设
$X_i = \coprod E_{i, k}$。这是因为 $X$ 的拓扑空间是谱空间 $X_i$ 沿具有谱过渡
映射的有向系在拓扑空间范畴中的极限。见极限章第 \ref{limits-section-descent} 节和
拓扑章第 \ref{topology-section-limits-spectral} 节。定义
$w_i : X_i \to \mathbf{Z}$，使其各水平集为可构造集 $E_{i, k}$。

\medskip\noindent
按引理 \ref{more-morphisms-lemma-open-and-closed-in-quasi-finite} 选取 $Y_{i, univ} \to Y_i$ 和
$U_{i, univ} \subset Y_{i, univ} \times_{Y_i} X_i$。由该构造的泛性质，为了证明
$w_i$ 是赋权，只须证明
$$
\tau_i = \textstyle{\int}_{U_{i, univ} \to Y_{i, univ}} w_i|_{U_{i, univ}}
$$
在 $Y_{i, univ}$ 上局部常值。由引理 \ref{more-morphisms-lemma-jumps-int-w}，该函数具有可构造
水平集，但它（尚）不一定局部常值。令 $Y_{univ} = Y_{i, univ} \times_{Y_i} Y$，
并令 $U_{univ} \subset Y_{univ} \times_Y X$ 为 $U_{i, univ}$ 的逆像。由于 $w$
到 $Y_{univ} \times_Y X$ 的拉回是 $Y_{univ} \times_Y X \to Y_{univ}$ 的一个赋权
（引理 \ref{more-morphisms-lemma-weighting-base-change}），故
$$
\tau = \textstyle{\int}_{U_{univ} \to Y_{univ}} w_i|_{U_{univ}}
$$
在 $Y_{univ}$ 上局部常值。因此 $\tau$ 的各水平集均开且闭。最后，
$Y_{univ} = \lim_{i' \geq i} Y_{i', univ}$，并且 $\tau$ 的各水平集是相应定义的
$\tau_{i'}$ 的各水平集的逆极限。因此，上述引文蕴含：当 $i'$ 充分大时，
$\tau_{i'}$ 的各水平集也为开集。对这样的指标 $i'$，可知 $w_{i'}$ 是 $f_{i'}$
的一个赋权，正如所需。
\end{proof}







\section{赋权与仿射分层数}
\label{more-morphisms-section-bounds-asn}

\noindent
本节对一类由仿射概形覆盖的概形之仿射分层数给出一个界。

\begin{lemma}
\label{more-morphisms-lemma-affineness-of-large-open}
设 $f : X \to Y$ 为仿射概形间拟有限且有限呈现的态射，
$w : X \to \mathbf{Z}_{> 0}$ 为 $f$ 的一个正赋权。设 $d < \infty$ 为
$\int_f w$ 的最大值。则 $Y$ 的开子集
$$
Y_d = \{y \in Y \mid (\textstyle{\int}_f w)(y) = d \}
$$
为仿射概形。
\end{lemma}

\begin{proof}
由引理 \ref{more-morphisms-lemma-max-int-w}，$\int_f w$ 取得其最大值；由引理
\ref{more-morphisms-lemma-semicontinuous-int-w}，集合 $Y_d$ 为开集。因此本引理的陈述确有意义。

\medskip\noindent
约化到 Noether 情形；请跳过本段。回忆有限型态射为拟有限态射，当且仅当其相对维数
为 $0$，见态射章引理
\ref{morphisms-lemma-locally-quasi-finite-rel-dimension-0}。在 $d = 0$ 时应用引理
\ref{more-morphisms-lemma-Noetherian-approximation-dimension-d}，可找到仿射 Noether 概形间的
拟有限态射 $f_0 : X_0 \to Y_0$ 以及态射 $Y \to Y_0$，使得 $f$ 是 $f_0$ 的
基变换。于是可将 $Y = \lim Y_i$ 写成在 $Y_0$ 上有限型的仿射概形的有向极限，
见代数章引理 \ref{algebra-lemma-ring-colimit-fp}。由引理
\ref{more-morphisms-lemma-descend-weighting}，可找到 $i$，使我们的赋权 $w$ 下降为 $f_0$ 的基变换
$f_i : X_i \to Y_i$ 的一个赋权 $w_i$。若本引理对 $f_i, w_i$ 成立，则它也对
$f$ 成立，因为 $\int_f w$ 的构造与基变换交换，见引理
\ref{more-morphisms-lemma-weighting-check-after-etale-base-change}。

\medskip\noindent
设 $X$ 和 $Y$ 为 Noether 概形。令 $X' \to Y'$ 为 $f$ 沿态射
$g : Y' \to Y$ 的基变换。$\int_f w$ 的构造以及由此得到的开集 $Y_d$ 都与基变换
交换。若 $g$ 有限且满，则 $Y'_d \to Y_d$ 有限且满。在此情形下，证明 $Y_d$
为仿射概形等价于证明 $Y'_d$ 为仿射概形，见概形上同调章引理
\ref{coherent-lemma-image-affine-finite-morphism-affine-Noetherian}。

\medskip\noindent
由引理 \ref{more-morphisms-lemma-quasi-finite-separated-pass-through-finite}，可选取一个浸入
$X \to T$，其中 $T$ 在 $Y$ 上有限。我们对有限态射 $T \to Y$ 应用态射章引理
\ref{morphisms-lemma-massage-finite}。该引理说明，存在有限满态射 $Y' \to Y$，
使得 $Y' \times_Y T$ 是某个具有特殊形式且在 $Y'$ 上有限的概形 $T'$ 的闭子概形。
由第一段的讨论，可将 $Y$ 替换为 $Y'$，将 $T$ 替换为 $T'$，并将 $X$ 替换为
$Y' \times_Y X$。因此可设存在一个浸入 $X \to T$（未必开或闭）以及闭子概形
$T_i \subset T$，$i = 1, \ldots, n$，其中
\begin{enumerate}
\item $T \to Y$ 有限（且局部自由）；
\item $T_i \to Y$ 为同构；并且
\item $T = \bigcup_{i = 1, \ldots, n} T_i$ 在集合论意义下成立。
\end{enumerate}
令 $Y' = \coprod Y_k$ 为 $Y$ 的各不可约分支（视为 $Y$ 的整闭子概形）的不交并。
此处利用 $Y$ 为 Noether 概形，我们可以再次沿 $Y' \to Y$ 作基变换。因此还可设
$Y$ 为整 Noether 概形。

\medskip\noindent
去除重复项后，我们还可并且确实设当 $i \not = j$ 时 $T_i \not = T_j$。由于
$Y$ 以及所有 $T_i$ 都是整概形，这意味着若 $T_i$ 与 $T_j$ 相交，则它们的交是
映到 $Y$ 的真闭子集中的闭子集。

\medskip\noindent
注意 $V_i = X \cap T_i$ 是局部闭子集，而且是 $X$ 的闭子概形，故为仿射概形。
令 $\eta \in Y$ 和 $\eta_i \in T_i$ 为泛点。若 $\eta \not \in Y_d$，则
$Y_d = \emptyset$，结论成立。设 $\eta \in Y_d$。以
$I \in \{1, \ldots, n\}$ 表示满足 $\eta_i \in V_i$ 的指标 $i$ 所成的子集。
对 $i \in I$，局部闭子集 $V_i \subset T_i$ 包含不可约空间 $T_i$ 的泛点，
因而为开集。另一方面，由于 $f$ 为开映射（引理
\ref{more-morphisms-lemma-weighting-universally-open}），对任意 $x \in X$，均可找到 $i \in I$
以及特化 $\eta_i \leadsto x$。于是 $x \in T_i$，从而 $x \in V_i$。换言之，
在集合论意义下 $X = \bigcup_{i \in I} V_i$。我们断言
$Y_d = \bigcap_{i \in I} \Im(V_i \to Y)$；这将完成证明，因为 $Y$ 的各仿射开集
$\Im(V_i \to Y)$ 的交仍为仿射概形。

\medskip\noindent
对 $y \in Y$，设 $X$ 中的 $f^{-1}(\{y\}) = \{x_1, \ldots, x_r\}$。对每个
$i \in I$，至多存在一个 $j(i) \in \{1, \ldots, x_r\}$，使得
$\eta_i \leadsto x_{j(i)}$。事实上，$j(i)$ 存在且等于 $j$，当且仅当
$x_j \in V_i$。若 $i \in I$ 使得 $j = j(i)$ 存在，则 $V_i \to Y$ 在
$x_j \mapsto y$ 的一个邻域上为同构。因此，将源和靶分别替换为
$x_j \mapsto y$ 的邻域后，$\bigcup_{i \in I,\ j(i) = j} V_i \to Y$ 为有限态射。
于是赋权的定义给出 $w(x_j) = \sum_{i \in I,\ j(i) = j} w(\eta_i)$。故
$$
(\textstyle{\int}_f w)(\eta) =
\sum\nolimits_{i \in I} w(\eta_i) \geq
\sum\nolimits_{j(i)\text{ exists}} w(\eta_i) =
\sum\nolimits_j w(x_j) = (\textstyle{\int}_f w)(y)
$$
等号成立，当且仅当 $y$ 属于 $\bigcap_{i \in I} \Im(V_i \to Y)$，这正是所要证明的。
\end{proof}

\begin{proposition}
\label{more-morphisms-proposition-asn-weighting}
设 $f : X \to Y$ 为概形的拟有限满态射，
$w : X \to \mathbf{Z}_{> 0}$ 为 $f$ 的一个正赋权。假设 $X$ 为仿射概形，而
$Y$ 分离\footnote{$Y$ 的对角态射为仿射态射即已足够。}且非空。则 $Y$ 的仿射分层数
至多为 $\int_f w$ 的不同取值个数减 $1$。
\end{proposition}

\begin{proof}
注意，由于 $Y$ 分离，态射 $X \to Y$ 为仿射态射（态射章引理
\ref{morphisms-lemma-affine-permanence}）。由引理 \ref{more-morphisms-lemma-max-int-w}，函数
$\int_f w$ 取得其最大值 $d$。我们对 $d$ 作归纳。考虑 $Y$ 的开子概形
$Y_d = \{y \in Y \mid (\int_f w)(y) = d\}$；回忆
$f^{-1}(Y_d) \to Y_d$ 为有限态射，见引理 \ref{more-morphisms-lemma-max-int-finite}。
由引理 \ref{more-morphisms-lemma-affineness-of-large-open}，对每个仿射开集 $W \subset Y$，
$Y_d \cap W$ 均为仿射概形（这里使用 $W \times_Y X$ 在 $X$ 上仿射，因而为仿射
概形）。故 $Y_d \to Y$ 为仿射态射。于是
$f^{-1}(Y_d) = Y_d \times_Y X$ 在 $X$ 上仿射，从而是仿射概形。又
$f^{-1}(Y_d) \to Y_d$ 为满态射，故由极限章引理 \ref{limits-lemma-affine}，
$Y_d$ 为仿射概形。若 $Y_d = Y$，则 $Y$ 的仿射分层数为 $0$，结论成立。
否则，令 $Y' = Y \setminus Y_d$，视为 $Y$ 的既约闭子概形，并令
$X' = Y' \times_Y X$。由于 $X'$ 在 $X$ 中闭，故为仿射概形；由于 $Y'$ 在 $Y$
中闭，故为分离概形。态射 $f' : X' \to Y'$ 为满态射，而 $w$ 诱导 $f'$ 的一个
赋权 $w'$（见引理 \ref{more-morphisms-lemma-weighting-base-change}），且 $\int_{f'} w'$ 是
$\int_f w$ 在 $Y'$ 上的限制。由归纳假设，$Y'$ 具有长度 $\leq$
$\int_{f'} w'$ 的不同取值个数减 $1$ 的仿射分层，证明完成。
\end{proof}









\section{完全分解态射}
\label{more-morphisms-section-completely-decomposed}

\noindent
Nishnevich 为了定义如今所谓的 Nishnevich 拓扑，研究了 \'etale 态射的完全分解族
这一概念；例如见 \cite{Nishnevich}。

\begin{definition}
\label{more-morphisms-definition-cd-morphism}
若概形态射 $f : X \to Y$ 满足：对每个点 $y \in Y$，都存在点 $x \in X$，
使得 $f(x) = y$ 且域扩张 $\kappa(x)/\kappa(y)$ 为平凡扩张，则称它为
{\it 完全分解的}\footnote{这一术语可能并非标准术语。}。若具有固定靶的概形态射族
$\{f_i : X_i \to Y\}_{i \in I}$ 满足 $\coprod f_i : \coprod Y_i \to X$
完全分解，则称该族{\it 完全分解}。
\end{definition}

\noindent
我们先给出若干基本引理。

\begin{lemma}
\label{more-morphisms-lemma-composition-cd}
两个完全分解概形态射的复合仍完全分解。若
$\{f_i : X_i \to Y\}_{i \in I}$ 完全分解，并且对每个 $i$，族
$\{X_{ij} \to X_i\}_{j \in J_i}$ 完全分解，则族
$\{X_{ij} \to Y\}_{i \in I, j \in J_i}$ 完全分解。
\end{lemma}

\begin{proof}
从略。
\end{proof}

\begin{lemma}
\label{more-morphisms-lemma-base-change-cd}
完全分解概形态射的基变换仍完全分解。若
$\{f_i : X_i \to Y\}_{i \in I}$ 完全分解，且 $Y' \to Y$ 为概形态射，则
$\{X_i \times_Y Y' \to Y'\}_{i \in I}$ 完全分解。
\end{lemma}

\begin{proof}
设 $f : X \to Y$ 和 $g : Y' \to Y$ 为概形态射。取点 $y' \in Y'$，其在 $Y$
中的像为 $y = g(y')$。若点 $x \in X$ 满足 $f(x) = y$ 且
$\kappa(x) = \kappa(y)$，则存在唯一一点 $x' \in X' = X \times_Y Y'$，
它在 $Y'$ 中映到 $y'$、在 $X$ 中映到 $x$，并且
$\kappa(x') = \kappa(y')$，见概形章引理
\ref{schemes-lemma-points-fibre-product}。本引理由此事实容易得出；细节从略。
\end{proof}

\begin{lemma}
\label{more-morphisms-lemma-decompose}
\begin{reference}
\cite[Lemma 2.1.2]{EHIK}
\end{reference}
设 $f : X \to Y$ 为概形态射。假设 $f$ 完全分解、$f$ 局部有限呈现，并且 $Y$ 拟紧且
拟分离。则存在 $n \geq 0$ 以及态射 $Z_i \to Y$，$i = 1, \ldots, n$，
具有如下性质：
\begin{enumerate}
\item $\coprod Z_i \to Y$ 为满态射；
\item 对所有 $i$，$Z_i \to Y$ 均为浸入；
\item 对所有 $i$，$Z_i \to Y$ 均为有限呈现；并且
\item 对所有 $i$，基变换 $X \times_Y Z_i \to Z_i$ 均具有截面。
\end{enumerate}
\end{lemma}

\begin{proof}
取 $y \in Y$。由假设，存在 $Y$ 上的态射
$\sigma : \Spec(\kappa(y)) \to X$。可将 $\Spec(\kappa(y))$ 写成 $Y$ 上仿射概形
$Z$ 的有向极限，其中 $Z \to Y$ 为有限呈现浸入。事实上，选取仿射开集
$y \in \Spec(A) \subset Y$，并设 $y$ 对应于 $A$ 的素理想 $\mathfrak p$。则
$\kappa(\mathfrak p)$ 是环 $(A/I)_f$ 的滤过余极限，其中
$I \subset \mathfrak p$ 为有限生成理想，且 $f \in A$、$f \not \in \mathfrak p$。
态射 $Z = \Spec((A/I)_f) \to Y$ 是有限呈现浸入；由于 $Y$ 拟分离，
$Z \to Y$ 拟紧，见概形章引理 \ref{schemes-lemma-quasi-compact-permanence}。
由极限章命题 \ref{limits-proposition-characterize-locally-finite-presentation}，
对某个这样的 $Z$，存在 $Y$ 上的态射 $\sigma' : Z \to X$，它在
$\kappa(\mathfrak p)$ 的谱上与 $\sigma$ 一致。由于 $\sigma'$ 是 $Y$ 上的态射，
我们得到投影 $X \times_Y Z \to Z$ 的一个截面。

\medskip\noindent
由此可知，$Y$ 是若干有限呈现浸入 $Z \to Y$ 的像之并，其中
$X \times_Y Z \to Z$ 具有截面。由于 $Z \to Y$ 的像可构造（态射章引理
\ref{morphisms-lemma-chevalley}），且 $Y$ 在可构造拓扑下为紧空间（性质章引理
\ref{properties-lemma-quasi-compact-quasi-separated-spectral} 及拓扑章引理
\ref{topology-lemma-constructible-hausdorff-quasi-compact}），可知有限多个这样的像
即已足够。
\end{proof}

\begin{lemma}
\label{more-morphisms-lemma-descend-cd}
设 $S = \lim_{\lambda \in \Lambda} S_\lambda$ 为具有仿射过渡态射的概形有向系的
极限。取 $0 \in \Lambda$，并设 $f_0 : X_0 \to Y_0$ 为 $S_0$ 上的概形态射。
对 $\lambda \geq 0$，令 $f_\lambda : X_\lambda \to Y_\lambda$ 为 $f_0$ 到
$S_\lambda$ 的基变换，并令 $f : X \to Y$ 为 $f_0$ 到 $S$ 的基变换。若
\begin{enumerate}
\item $f$ 完全分解；
\item $Y_0$ 拟紧且拟分离；并且
\item $f_0$ 局部有限呈现，
\end{enumerate}
则存在 $\lambda \geq 0$，使得 $f_\lambda$ 完全分解。
\end{lemma}

\begin{proof}
由于 $Y_0$ 拟紧且拟分离，而概形 $Y$ 在 $Y_0$ 上仿射，故其拟紧且拟分离。
按引理 \ref{more-morphisms-lemma-decompose} 选取 $n \geq 0$ 以及 $Z_i \to Y$，
$i = 1, \ldots, n$。以 $\sigma_i : Z_i \to X$ 表示由 $Z_i$ 的选取而存在的
$Y$ 上态射。增大 $0 \in \Lambda$ 后，可设存在有限呈现态射
$Z_{i, 0} \to Y_0$，其到 $S$ 的基变换为 $Z_i \to Y$，见极限章引理
\ref{limits-lemma-descend-finite-presentation}。由极限章引理
\ref{limits-lemma-descend-immersion}，必要时再次增大 $0$ 后，可设
$Z_{i, 0} \to Y_0$ 为浸入。由于 $\coprod Z_i \to Y$ 为满态射，必要时增大 $0$
后，可设 $\coprod Z_{i, 0} \to Y_0$ 为满态射，见极限章引理
\ref{limits-lemma-descend-surjective}。注意
$Z_i = \lim_{\lambda \geq 0} Z_{i, \lambda}$，其中
$Z_{i, \lambda} = Y_\lambda \times_{Y_0} Z_{i, 0}$。将复合
$$
Z_i \xrightarrow{\sigma_i} X \to X_0
$$
视为 $Y_0$ 上的态射。由于 $f_0$ 局部有限呈现，由极限章命题
\ref{limits-proposition-characterize-locally-finite-presentation}，可找到
$\lambda \geq 0$，使得存在 $Y_0$ 上的态射
$\sigma'_{i, \lambda} : Z_{i, \lambda} \to X_0$，其与
$Z_i \to Z_{i, \lambda}$ 的预复合就是所显示的箭头。当然，此时
$\sigma'_{i, \lambda}$ 确定一个 $Y_\lambda$ 上的态射
$\sigma_{i, \lambda} : Z_{i, \lambda} \to
X_\lambda = X_0 \times_{Y_0} Y_\lambda$。由于
$\coprod Z_{i, \lambda} \to Y_\lambda$ 为满态射，可知
$X_\lambda \to Y_\lambda$ 完全分解。
\end{proof}




\section{丰富可逆模族}
\label{more-morphisms-section-families-ample-invertible-modules}

\noindent
我们继续态射章第 \ref{morphisms-section-families-ample-invertible-modules} 节的讨论。

\begin{lemma}
\label{more-morphisms-lemma-ample-family-ample-relative}
设 $f : X \to Y$ 为概形态射。假设
\begin{enumerate}
\item $Y$ 具有一个丰富可逆模族；
\item $X$ 上存在一个 $f$-丰富可逆模。
\end{enumerate}
则 $X$ 具有一个丰富可逆模族。
\end{lemma}

\begin{proof}
令 $\mathcal{L}$ 为 $X$ 上的一个 $f$-丰富可逆模。特别地，这蕴含 $f$ 拟紧，见
态射章定义 \ref{morphisms-definition-relatively-ample}。由态射章定义
\ref{morphisms-definition-family-ample-invertible-modules}，$Y$ 拟紧，故 $X$ 拟紧
（因而 $X$ 本身满足态射章定义
\ref{morphisms-definition-family-ample-invertible-modules} 的第一个条件）。
取 $x \in X$，其像为 $y \in Y$。由假设 (2)，可找到一个可逆
$\mathcal{O}_Y$-模 $\mathcal{N}$ 以及截面 $t \in \Gamma(Y, \mathcal{N})$，
使得 $t$ 不消失的轨迹 $Y_t$ 为仿射概形。于是 $\mathcal{L}$ 在
$f^{-1}(Y_t) = X_{f^*t}$ 上丰富，故可找到截面
$s \in \Gamma(X_{f^*t}, \mathcal{L})$，使得 $(X_{f^*t})_s$ 为仿射概形且包含
$x$。由性质章引理 \ref{properties-lemma-invert-s-sections}，对某个
$n \geq 0$，乘积 $(f^*t)^n s$ 延拓为截面
$s' \in \Gamma(X, f^*\mathcal{N}^{\otimes n} \otimes \mathcal{L})$。
最后，$f^*\mathcal{N}^{\otimes n + 1} \otimes \mathcal{L}$ 的截面
$s'' = f^* ts'$ 在 $X \setminus X_{f^*t}$ 的每一点均消失，故
$X_{s''} = (X_{f^*t})_s$ 为仿射概形，正如所需。
\end{proof}

\begin{lemma}
\label{more-morphisms-lemma-resolution-property-goes-up-affine}
设 $f : X \to Y$ 为概形的仿射或拟仿射态射。若 $Y$ 具有一个丰富可逆模族，
则 $X$ 也具有一个丰富可逆模族。
\end{lemma}

\begin{proof}
由态射章引理 \ref{morphisms-lemma-quasi-affine-O-ample}，这是引理
\ref{more-morphisms-lemma-ample-family-ample-relative} 的特殊情形。
\end{proof}






\section{吹起与丰富可逆模族}
\label{more-morphisms-section-ample-family-by-blowing-up}

\noindent
我们证明 \cite{Gross-thesis} 中的一个结果。

\begin{lemma}
\label{more-morphisms-lemma-get-ample-family}
设 $X$ 为概形。给定 $X$ 上的有效 Cartier 除子
$D_1, \ldots, D_m$ 以及可逆模 $\mathcal{L}_1, \ldots, \mathcal{L}_m$，满足
$\bigcap D_i = \emptyset$ 且 $\mathcal{L}_i|_{X \setminus D_i}$ 丰富。
则 $X$ 具有一个丰富可逆模族。
\end{lemma}

\begin{proof}
取 $x \in X$。选取指标 $i \in \{1, \ldots, m\}$，使得 $x \not \in D_i$。
令 $U_i = X \setminus D_i$。由 $\mathcal{L}_i|_{U_i}$，可找到 $n \geq 1$ 以及
截面 $s \in \Gamma(U_i, \mathcal{L}_i^{\otimes n})$，使得 $s$ 不消失的轨迹
$(U_i)_s$ 为仿射概形（性质章定义 \ref{properties-definition-ample}）。
由于 $U_i$ 是典范截面 $1 \in \mathcal{O}_X(D_i)$ 不消失的轨迹，故由性质章引理
\ref{properties-lemma-invert-s-sections}，存在 $N \geq 0$，使得 $s$ 延拓为截面
$$
s' \in \Gamma(X, \mathcal{L}_i^{\otimes n} \otimes_{\mathcal{O}_X}
\mathcal{O}_X(N D_i))
$$
将 $N$ 替换为 $N + 1$ 后，可见 $s'$ 在 $D_i$ 的每一点都消失，因而
$X_{s'} = (U_i)_s$ 为仿射概形。这证明 $X$ 具有一个丰富可逆模族，见态射章定义
\ref{morphisms-definition-family-ample-invertible-modules}。
\end{proof}

\begin{lemma}
\label{more-morphisms-lemma-blow-up-ample-family}
\begin{reference}
\cite[Proposition 1.3.1]{Gross-thesis}
\end{reference}
设 $X$ 为仅有有限多个不可约分支的拟紧拟分离概形。则存在一个拟紧稠密开子集
$U \subset X$ 以及一个 $U$-容许吹起 $X' \to X$，使得概形 $X'$ 具有一个丰富
可逆模族。
\end{lemma}

\begin{proof}
令 $\eta_1, \ldots, \eta_n \in X$ 为 $X$ 的各不可约分支的泛点。由性质章引理
\ref{properties-lemma-point-and-maximal-points-affine} 以及 $X$ 拟紧这一事实，
可找到有限仿射开覆盖 $X = U_1 \cup \ldots \cup U_m$，使得每个 $U_i$ 都包含
$\eta_1, \ldots, \eta_n$。特别地，拟紧开子集
$U = U_1 \cap \ldots \cap U_m$ 在 $X$ 中稠密。令
$\mathcal{I}_i \subset \mathcal{O}_X$ 为有限型拟凝聚理想层，使得
$U_i = X \setminus Z_i$，其中 $Z_i = V(\mathcal{I}_i)$，见性质章引理
\ref{properties-lemma-quasi-coherent-finite-type-ideals}。令
$$
f : X' \longrightarrow X
$$
为 $X$ 沿理想层 $\mathcal{I} = \mathcal{I}_1 \cdots \mathcal{I}_m$ 的吹起。
注意，$f$ 是 $U$-容许吹起，因为 $V(\mathcal{I})$ 在集合论意义下是与 $U$
不交的各 $Z_i$ 之并。此外，$f$ 为射影态射，且 $\mathcal{O}_{X'}(1)$ 是
$f$-相对丰富的，见除子章引理 \ref{divisors-lemma-blowing-up-projective}。
由除子章引理 \ref{divisors-lemma-blowing-up-two-ideals}，对每个 $i$，态射
$f'$ 分解为 $X' \to X'_i \to X$，其中 $X'_i \to X$ 是沿 $\mathcal{I}_i$
的吹起，而 $X' \to X'_i$ 是另一个吹起（即沿略去 $\mathcal{I}_i$ 后诸理想
$\mathcal{I}_j$ 之积的拉回所作的吹起）。由此可知
$D_i = f^{-1}(Z_i) \subset X'$ 是有效 Cartier 除子，见除子章引理
\ref{divisors-lemma-blow-up-pullback-effective-Cartier} 和
\ref{divisors-lemma-blowing-up-gives-effective-Cartier-divisor}。我们有
$X' \setminus D_i = f^{-1}(U_i)$。由于 $\mathcal{O}_{X'}(1)$ 为 $f$-丰富，
$\mathcal{O}_{X'}(1)$ 在 $X' \setminus D_i$ 上的限制丰富。由引理 \ref{more-morphisms-lemma-get-ample-family}，
$X'$ 具有一个丰富可逆模族。
\end{proof}

\begin{proposition}
\label{more-morphisms-proposition-envelope-with-resolution-property}
设 $X$ 为拟紧拟分离概形。则存在一个有限呈现、固有且完全分解（定义
\ref{more-morphisms-definition-cd-morphism}）的态射 $f : Y \to X$，使得概形 $Y$ 具有一个
丰富可逆模族。
\end{proposition}

\begin{proof}
由极限章命题 \ref{limits-proposition-approximate}，存在仿射态射 $X \to X_0$，
其中 $X_0$ 为 $\mathbf{Z}$ 上有限型概形。下面构造一个具有所有所需性质的态射
$Y_0 \to X_0$。令 $Y = X \times_{X_0} Y_0$，并令 $f$ 等于投影
$f : Y \to X$，即可得出结论。事实上，$f$ 为有限呈现（态射章引理
\ref{morphisms-lemma-base-change-finite-presentation}），$f$ 为固有态射
（态射章引理 \ref{morphisms-lemma-base-change-proper}），且 $f$ 完全分解
（引理 \ref{more-morphisms-lemma-base-change-cd}）。最后，由于 $Y \to Y_0$ 为仿射态射
（它是 $X \to X_0$ 的基变换），由引理
\ref{more-morphisms-lemma-resolution-property-goes-up-affine}，$Y$ 具有一个丰富可逆模族。
这将问题约化到下一段讨论的情形。

\medskip\noindent
设 $X$ 在 $\mathbf{Z}$ 上有限型。特别地，$\dim(X) < \infty$。我们对
$\dim(X)$ 作归纳。若 $\dim(X) = 0$，则 $X$ 为仿射概形并具有消解性质。
一般地，由引理 \ref{more-morphisms-lemma-blow-up-ample-family}，存在稠密开子集 $U \subset X$
以及一个 $U$-容许吹起 $X' \to X$，使得 $X'$ 具有一个丰富可逆模族。由于
$f : X' \to X$ 在 $U$ 上为同构，$U$ 的每一点都可提升为 $X'$ 中具有相同剩余域
的一点。令 $Z = X \setminus U$，赋予既约诱导概形结构。由于 $U$ 在 $X$ 中稠密
（见上文），有 $\dim(Z) < \dim(X)$。由归纳假设，可找到一个固有、完全分解的态射
$W \to Z$，使得 $W$ 具有一个丰富可逆模族。于是
$Y = X' \amalg W \to X$ 即为所需态射。
\end{proof}




\section{闭浸入的广延判据}
\label{more-morphisms-section-extensive-criterion-closed-immersions}

\noindent
本节给出概形态射成为闭浸入的一个判据。

\begin{lemma}
\label{more-morphisms-lemma-affine-extensive-criterion}
仿射概形态射 $f : X \to Y$ 为闭浸入，当且仅当对每个单射环同态 $A \to B$
以及交换方块
$$
\xymatrix{
\Spec(B) \ar[d] \ar[r] & X \ar[d]^f \\
\Spec(A) \ar[r] \ar@{..>}[ur] & Y
}
$$
都存在提升 $\Spec(A) \to X$，使两个三角形交换。
\end{lemma}

\begin{proof}
设态射 $f$ 由环同态 $\phi : R \to S$ 给出。则 $f$ 为闭浸入，当且仅当
$\phi$ 为满射。

\medskip\noindent
首先假设 $\phi$ 为满射。设 $\psi : A \to B$ 为单射环同态，并且给定交换图表
$$
\xymatrix{
R \ar[r]^\alpha \ar[d]^\phi & A \ar[d]^\psi \\
S \ar[r]^\beta \ar@{..>}[ur] & B
}
$$
定义提升 $S \to A$，令 $s \mapsto \alpha(r)$，其中 $r \in R$ 满足
$\phi(r) = s$。由于 $\psi$ 为单射且方块交换，这一定义适定。取环的谱给出交换环
范畴与仿射概形范畴之间的反等价，故 $f$ 具有所需的提升性质。

\medskip\noindent
其次，假设 $\phi$ 对所有单射环同态 $\psi: A \to B$ 均具有提升。注意
$\phi(R)$ 是 $S$ 的子环，故得到交换方块
$$
\xymatrix{
R \ar[r] \ar[d]^\phi  & \phi(R) \ar[d] \\
S \ar@{=}[r] \ar@{..>}[ur] & S
}
$$
其中存在提升 $S \to \phi(R)$。因此包含映射 $\phi(R) \to S$ 必为同构，
这说明 $\phi$ 为满射，证明完成。
\end{proof}

\begin{lemma}
\label{more-morphisms-lemma-scheme-affine-if-aff-is-split-mono}
设 $X$ 为概形。若概形章引理 \ref{schemes-lemma-morphism-into-affine} 的典范态射
$X \to \Spec(\Gamma(X, \mathcal{O}_X))$ 具有收缩，则 $X$ 为仿射概形。
\end{lemma}

\begin{proof}
记 $S = \Spec(\Gamma(X, \mathcal{O}_X))$，并以 $f : X \to S$ 表示引理中给出的
态射。设 $s : S \to X$ 为一个收缩；于是 $\text{id}_X = sf$。从而
$f s f = \text{id}_S f$。由于 $f$ 诱导同构
$\Gamma(S, \mathcal{O}_S) \to \Gamma(X, \mathcal{O}_X)$，这说明 $fs$ 与
$\text{id}_S$ 在 $\Gamma(S, \mathcal{O}_S)$ 上诱导相同的映射。由于 $S$
为仿射概形，遂有 $f s = \text{id}_S$。因此 $f$ 为同构，$X$ 为仿射概形，
正如所要证明。
\end{proof}

\begin{lemma}
\label{more-morphisms-lemma-aff-is-injective-on-sheaves}
设 $X$ 为概形，并令 $f : X \to S = \Spec(\Gamma(X, \mathcal{O}_X))$ 为概形章
引理 \ref{schemes-lemma-morphism-into-affine} 的典范态射。包含在
$f^\sharp : \mathcal{O}_S \to f_*\mathcal{O}_X$ 的核中的最大拟凝聚
$\mathcal{O}_S$-模为零。若 $X$ 拟紧，则 $f^\sharp$ 为单射。特别地，若 $X$
拟紧，则 $f$ 为支配态射。
\end{lemma}

\begin{proof}
令 $M \subset \Gamma(S, \mathcal{O}_S)$ 为与包含在 $f^\sharp$ 的核中的最大
拟凝聚 $\mathcal{O}_S$-模相对应的子模。则任意元素 $a \in M$ 均被
$f^\sharp$ 映到零。然而，$f^\sharp(a)$ 是下列模中的元素
$$
\Gamma(S, f_*\mathcal{O}_X) = \Gamma(X, \mathcal{O}_X) =
\Gamma(S, \mathcal{O}_S)
$$
且它对应于 $a$ 本身！故 $a = 0$。于是 $M = 0$，这证明了第一个断言。
注意，这等价于态射 $f : X \to S$ 在概形论意义下为满态射。

\medskip\noindent
若 $X$ 拟紧，则由态射章引理
\ref{morphisms-lemma-quasi-compact-scheme-theoretic-image}，
$\Ker(f^\sharp)$ 拟凝聚。因此 $\Ker(f^\sharp) = 0$，且 $f^\sharp$ 为单射。
在此情形下，由态射章引理
\ref{morphisms-lemma-quasi-compact-scheme-theoretic-image} 的第 (4) 部分，
$f$ 为支配态射。
\end{proof}

\begin{lemma}
\label{more-morphisms-lemma-extensive-criterion}
设 $f: X \to Y$ 为概形的拟紧态射。则 $f$ 为闭浸入，当且仅当对每个单射环同态
$A \to B$ 以及交换方块
$$
\xymatrix{
\Spec(B) \ar[d] \ar[r] & X \ar[d]^f \\
\Spec(A) \ar[r] \ar@{..>}[ur] & Y
}
$$
都存在提升 $\Spec A \to X$，使图表交换。
\end{lemma}

\begin{proof}
假设 $f$ 为闭浸入。设 $A \to B$ 为单射环同态，并考虑交换方块
$$
\xymatrix{
\Spec(B) \ar[d] \ar[r]  & X \ar[d]^f \\
\Spec(A) \ar[r] \ar@{..>}[ur] & Y
}
$$
则 $\Spec(A) \times_Y X \to \Spec(A)$ 为闭浸入，故得到理想 $I \subset A$
以及交换图表
$$
\xymatrix{
\Spec(B) \ar[d] \ar[r] & \Spec(A/I) \ar[r] \ar[d] & X \ar[d]^f \\
\Spec(A) \ar[r] \ar@{..>}[ur] & \Spec(A) \ar[r] & Y
}
$$
由引理 \ref{more-morphisms-lemma-affine-extensive-criterion}，得到一个提升。

\medskip\noindent
假设 $f$ 具有引理中所述的提升性质。证明 $f$ 为闭浸入这一问题在 $Y$ 上是局部的，
故可并且确实设 $Y$ 为仿射概形。特别地，$Y$ 拟紧，从而 $X$ 拟紧。因此存在有限
仿射开覆盖 $X = U_1 \cup \ldots \cup U_n$。态射
$$
\pi : U = \coprod U_i \longrightarrow X
$$
的源为仿射概形，并且所诱导的环同态
$\Gamma(X, \mathcal{O}_X) \to \Gamma(U, \mathcal{O}_U)$ 为单射。由假设，
下图中存在提升：
$$
\xymatrix{
U \ar[r]^\pi \ar[d] & X \ar[d]^f \\
\Spec(\Gamma(X, \mathcal{O}_X)) \ar[r]^-{f'} \ar@{..>}[ur]^h & Y
}
$$
其中 $f'$ 是与环同态
$\Gamma(Y, \mathcal{O}_Y) \to \Gamma(X, \mathcal{O}_X)$ 相对应的仿射概形态射。
由 $\pi$ 为满射可知，态射 $h$ 是典范态射
$X \to \Spec(\Gamma(X, \mathcal{O}_X))$ 的一个收缩；细节从略。因此，由引理
\ref{more-morphisms-lemma-scheme-affine-if-aff-is-split-mono}，$X$ 为仿射概形。再由引理
\ref{more-morphisms-lemma-affine-extensive-criterion}，可知 $f$ 为闭浸入。
\end{proof}
% Generated mechanically from the frozen flat.tex target.
% Do not edit this generated file; edit the bound locale source instead.

% OK, start here.
%

\setcounter{chapter}{37}
\chapter{平坦性续篇}


\phantomsection
\label{flat-section-phantom}




\section{引言}
\label{flat-section-introduction}

\noindent
本章讨论平坦模与概形平坦态射的一些高阶结果及其应用。关于平坦性的多数结果
可见 Raynaud 与 Gruson 的论文 \cite{GruRay}。

\medskip\noindent
在阅读本章之前，建议读者先查看下列结果（此表也可作为先前结果的索引）：
\begin{enumerate}
\item 关于平坦模的一般讨论，见《代数》第 \ref{algebra-section-flat} 节。
\item $\text{Tor}$ 群与平坦性之间的关系，见《代数》第
\ref{algebra-section-tor} 节。
\item 平坦性判据，见《代数》第 \ref{algebra-section-criteria-flatness} 节
（诺特情形）、《代数》第 \ref{algebra-section-flatness-artinian} 节
（阿廷情形）、《代数》第 \ref{algebra-section-more-flatness-criteria} 节
（非诺特情形），最后见《态射进阶》第
\ref{more-morphisms-section-criterion-flat-fibres} 节。
\item 一般平坦性，见《代数》第 \ref{algebra-section-generic-flatness} 节与
《态射》第 \ref{morphisms-section-generic-flatness} 节。
\item 平坦轨迹的开性，见《代数》第 \ref{algebra-section-open-flat} 节与
《态射进阶》第 \ref{more-morphisms-section-open-flat} 节。
\item 平坦化，见《代数进阶》诸节
\ref{more-algebra-section-flattening},
\ref{more-algebra-section-flattening-artinian},
\ref{more-algebra-section-flattening-local-base},
\ref{more-algebra-section-flattening-local-source-base} 以及
\ref{more-algebra-section-flattening-Noetherian-complete-local}.
\item 其他结果见《代数进阶》诸节
\ref{more-algebra-section-descent-flatness-integral}、
\ref{more-algebra-section-torsion-flat}、
\ref{more-algebra-section-flat-finite} 以及
\ref{more-algebra-section-blowup-flat}.
\end{enumerate}
作为平坦性材料的应用，我们讨论以下主题：Grothendieck 存在定理的非诺特版本、
爆破与平坦性、Nagata 的紧化定理、h 拓扑、爆破方块与下降、弱正规化、
正特征下向量丛的下降，以及 h 拓扑中完美复形的局部结构。






\section{关于 \'etale 局部化的引理}
\label{flat-section-etale-localization}

\noindent
本节列出一些关于 \'etale 局部化的引理，它们将在本章后文中使用。
初次阅读时可以跳过本节。

\begin{lemma}
\label{flat-lemma-lift-etale}
设 $i : Z \to X$ 是仿射概形的闭浸入，$Z' \to Z$ 是 \'etale 态射，
且 $Z'$ 仿射。则存在 \'etale 态射 $X' \to X$，其中 $X'$ 仿射，
使得作为 $Z$ 上概形有 $Z' \cong Z \times_X X'$。
\end{lemma}

\begin{proof}
见《代数》引理 \ref{algebra-lemma-lift-etale}。
\end{proof}

\begin{lemma}
\label{flat-lemma-etale-at-point}
设
$$
\xymatrix{
X \ar[d] & X' \ar[l] \ar[d] \\
S & S' \ar[l]
}
$$
是概形的交换图，其中 $X' \to X$ 与 $S' \to S$ 均为 \'etale 态射。
设 $s' \in S'$ 为一点。则
$$
X' \times_{S'} \Spec(\mathcal{O}_{S', s'})
\longrightarrow
X \times_S \Spec(\mathcal{O}_{S', s'})
$$
是 \'etale 态射。
\end{lemma}

\begin{proof}
这是因为，作为 $X$ 上 \'etale 概形之间的态射，$X' \to X_{S'}$
是 \'etale 的；见《态射》引理 \ref{morphisms-lemma-etale-permanence}。
而且 \'etale 态射的基变换仍为 \'etale；见《态射》引理
\ref{morphisms-lemma-base-change-etale}。
\end{proof}

\begin{lemma}
\label{flat-lemma-etale-flat-up-down}
设 $X \to T \to S$ 是概形态射，且 $T \to S$ 为 \'etale 态射。
设 $\mathcal{F}$ 是拟凝聚 $\mathcal{O}_X$-模，$x \in X$ 为一点。则
$$
\mathcal{F}\text{ 在 }S\text{ 上于 }x\text{ 处平坦}
\Leftrightarrow
\mathcal{F}\text{ 在 }T\text{ 上于 }x\text{ 处平坦}
$$
特别地，$\mathcal{F}$ 在 $S$ 上平坦，当且仅当 $\mathcal{F}$ 在 $T$ 上平坦。
\end{lemma}

\begin{proof}
\'Etale 态射是平坦态射（见《态射》引理
\ref{morphisms-lemma-etale-flat}），故蕴含 ``$\Leftarrow$'' 由《代数》引理
\ref{algebra-lemma-composition-flat} 得出。反之，假设 $\mathcal{F}$
在点 $x$ 处于 $S$ 上平坦。记 $\tilde x \in X \times_S T$ 为这样一点：
它在 $X$ 中位于 $x$ 上方，在 $T$ 中位于 $x$ 在 $T$ 中的像上方。于是经由
$\text{pr}_2 : X \times_S T \to T$，
$(X \times_S T \to X)^*\mathcal{F}$ 在点 $\tilde x$ 处于 $T$ 上平坦；
见《态射》引理 \ref{morphisms-lemma-base-change-module-flat}。
对角态射 $\Delta_{T/S} : T \to T \times_S T$ 是开浸入；合并《态射》引理
\ref{morphisms-lemma-diagonal-unramified-morphism} 与
\ref{morphisms-lemma-etale-smooth-unramified} 即得。因此可将 $X$
等同于 $X \times_S T$ 的开子概形；$\text{pr}_2$ 在此开子概形上的限制
就是给定态射 $X \to T$，点 $\tilde x$ 对应于其中的点 $x$，并且
$(X \times_S T \to X)^*\mathcal{F}$ 在该开子概形上的限制就是
$\mathcal{F}$。由此可见，$\mathcal{F}$ 在点 $x$ 处于 $T$ 上平坦。
\end{proof}

\begin{lemma}
\label{flat-lemma-etale-flat-up-down-local-ring}
设 $T \to S$ 是 \'etale 态射，$t \in T$ 的像为 $s \in S$。
设 $M$ 是 $\mathcal{O}_{T, t}$-模。则
$$
M\text{ 在 }\mathcal{O}_{S, s}\text{ 上平坦}
\Leftrightarrow
M\text{ 在 }\mathcal{O}_{T, t}\text{ 上平坦}.
$$
\end{lemma}

\begin{proof}
可先用 $s$ 的仿射邻域替换 $S$，再用 $t$ 的仿射邻域替换 $T$。
置 $\mathcal{F} = (\Spec(\mathcal{O}_{T, t}) \to T)_*\widetilde M$。
这是 $T$ 上的拟凝聚层（见概形章引理
\ref{schemes-lemma-push-forward-quasi-coherent}，或直接论证），
其在 $t$ 处的茎为 $M$（细节从略）。应用引理
\ref{flat-lemma-etale-flat-up-down} 即得。
\end{proof}

\begin{lemma}
\label{flat-lemma-flat-up-down-henselization}
设 $S$ 是概形，$s \in S$ 为一点。记 $\mathcal{O}_{S, s}^h$
（分别地，$\mathcal{O}_{S, s}^{sh}$）为亨泽尔化（分别地，严格亨泽尔化）；
见《代数》定义 \ref{algebra-definition-henselization}。
设 $M^{sh}$ 是 $\mathcal{O}_{S, s}^{sh}$-模。下列条件等价：
\begin{enumerate}
\item $M^{sh}$ 在 $\mathcal{O}_{S, s}$ 上平坦；
\item $M^{sh}$ 在 $\mathcal{O}_{S, s}^h$ 上平坦；
\item $M^{sh}$ 在 $\mathcal{O}_{S, s}^{sh}$ 上平坦。
\end{enumerate}
若 $M^{sh} = M^h \otimes_{\mathcal{O}_{S, s}^h} \mathcal{O}_{S, s}^{sh}$，
则上述条件还等价于：
\begin{enumerate}
\item[(4)] $M^h$ 在 $\mathcal{O}_{S, s}$ 上平坦；
\item[(5)] $M^h$ 在 $\mathcal{O}_{S, s}^h$ 上平坦。
\end{enumerate}
若 $M^h = M \otimes_{\mathcal{O}_{S, s}} \mathcal{O}_{S, s}^h$，
则上述条件还等价于：
\begin{enumerate}
\item[(6)] $M$ 在 $\mathcal{O}_{S, s}$ 上平坦。
\end{enumerate}
\end{lemma}

\begin{proof}
由《代数进阶》引理 \ref{more-algebra-lemma-dumb-properties-henselization}，局部环同态
$\mathcal{O}_{S, s} \to \mathcal{O}_{S, s}^h \to \mathcal{O}_{S, s}^{sh}$
是忠实平坦的。因此 (3) $\Rightarrow$ (2) $\Rightarrow$ (1) 与
(5) $\Rightarrow$ (4) 由《代数》引理
\ref{algebra-lemma-composition-flat} 得出。由忠实平坦性，等价关系
(6) $\Leftrightarrow$ (5) 与 (5) $\Leftrightarrow$ (3) 由《代数》引理
\ref{algebra-lemma-flatness-descends} 得出。因此只需证明
(1) $\Rightarrow$ (2) $\Rightarrow$ (3) 与 (4) $\Rightarrow$ (5)。
为此可假设 $S$ 是仿射概形。

\medskip\noindent
假设 (1) 成立。由引理 \ref{flat-lemma-etale-flat-up-down-local-ring}，
对任意 \'etale 邻域 $(T, t) \to (S, s)$，$M^{sh}$ 在
$\mathcal{O}_{T, t}$ 上平坦。由于 $\mathcal{O}_{S, s}^h$ 与
$\mathcal{O}_{S, s}^{sh}$ 是形如 $\mathcal{O}_{T, t}$ 的局部环的有向余极限
（见《代数》引理 \ref{algebra-lemma-henselization-different} 与
\ref{algebra-lemma-strict-henselization-different}），由《代数》引理
\ref{algebra-lemma-colimit-rings-flat}，$M^{sh}$ 在
$\mathcal{O}_{S, s}^h$ 与 $\mathcal{O}_{S, s}^{sh}$ 上平坦。
故 (1) 蕴含 (2) 与 (3)。将 $\mathcal{O}_{S, s}$ 替换为
$\mathcal{O}_{S, s}^h$，当然也得到 (2) $\Rightarrow$ (3)。
同样的论证证明 (4) $\Rightarrow$ (5)。
\end{proof}

\begin{lemma}
\label{flat-lemma-tor-amplitude-up-down-henselization}
设 $S$ 是概形，$s \in S$ 为一点。记 $\mathcal{O}_{S, s}^h$
（分别地，$\mathcal{O}_{S, s}^{sh}$）为亨泽尔化（分别地，严格亨泽尔化）；
见《代数》定义 \ref{algebra-definition-henselization}。
设 $M^{sh}$ 是 $D(\mathcal{O}_{S, s}^{sh})$ 的对象，
$a, b \in \mathbf{Z}$。下列条件等价：
\begin{enumerate}
\item $M^{sh}$ 在 $\mathcal{O}_{S, s}$ 上的 Tor 振幅包含于 $[a, b]$；
\item $M^{sh}$ 在 $\mathcal{O}_{S, s}^h$ 上的 Tor 振幅包含于 $[a, b]$；
\item $M^{sh}$ 在 $\mathcal{O}_{S, s}^{sh}$ 上的 Tor 振幅包含于 $[a, b]$。
\end{enumerate}
若 $M^{sh} =
M^h \otimes_{\mathcal{O}_{S, s}^h}^\mathbf{L} \mathcal{O}_{S, s}^{sh}$
其中 $M^h \in D(\mathcal{O}_{S, s}^h)$，则上述条件还等价于：
\begin{enumerate}
\item[(4)] $M^h$ 在 $\mathcal{O}_{S, s}$ 上的 Tor 振幅包含于 $[a, b]$；
\item[(5)] $M^h$ 在 $\mathcal{O}_{S, s}^h$ 上的 Tor 振幅包含于 $[a, b]$。
\end{enumerate}
若 $M^h = M \otimes_{\mathcal{O}_{S, s}}^\mathbf{L} \mathcal{O}_{S, s}^h$
其中 $M \in D(\mathcal{O}_{S, s})$，则上述条件还等价于：
\begin{enumerate}
\item[(6)] $M$ 在 $\mathcal{O}_{S, s}$ 上的 Tor 振幅包含于 $[a, b]$。
\end{enumerate}
\end{lemma}

\begin{proof}
由《代数进阶》引理 \ref{more-algebra-lemma-dumb-properties-henselization}，局部环同态
$\mathcal{O}_{S, s} \to \mathcal{O}_{S, s}^h \to \mathcal{O}_{S, s}^{sh}$
是忠实平坦的。因此 (3) $\Rightarrow$ (2) $\Rightarrow$ (1) 与
(5) $\Rightarrow$ (4) 由《代数进阶》引理
\ref{more-algebra-lemma-flat-push-tor-amplitude} 得出。由忠实平坦性，
等价关系 (6) $\Leftrightarrow$ (5) 与 (5) $\Leftrightarrow$ (3) 由
《代数进阶》引理 \ref{more-algebra-lemma-flat-descent-tor-amplitude} 得出。
因此只需证明 (1) $\Rightarrow$ (3)、(2) $\Rightarrow$ (3) 与
(4) $\Rightarrow$ (5)。

\medskip\noindent
假设 (1) 成立。特别地，$M^{sh}$ 在次数 $< a$ 与 $> b$ 上的上同调为零。
故可用自由 $\mathcal{O}_{X, x}^{sh}$-模复形 $P^\bullet$ 表示
$M^{sh}$，且当 $i > b$ 时 $P^i = 0$（例如见一般性很强的导出范畴章引理
\ref{derived-lemma-subcategory-left-resolution}）。注意，对所有 $n$，
$P^n$ 在 $\mathcal{O}_{S, s}$ 上平坦。考虑 $\Coker(d_P^{a - 1})$。
由《代数进阶》引理 \ref{more-algebra-lemma-last-one-flat}，这是平坦
$\mathcal{O}_{S, s}$-模。故由引理 \ref{flat-lemma-flat-up-down-henselization}，
它也是平坦 $\mathcal{O}_{S, s}^{sh}$-模。因此
$\tau_{\geq a}P^\bullet$ 是平坦 $\mathcal{O}_{S, s}^{sh}$-模复形，
它在 $D(\mathcal{O}_{S, s}^{sh}$ 中表示 $M^{sh}$，于是可知
$M^{sh}$ 的 Tor 振幅包含于 $[a, b]$；见《代数进阶》引理
\ref{more-algebra-lemma-tor-amplitude}。故 (1) 蕴含 (3)。将
$\mathcal{O}_{S, s}$ 替换为 $\mathcal{O}_{S, s}^h$，当然也得到
(2) $\Rightarrow$ (3)。同样的论证证明 (4) $\Rightarrow$ (5)。
\end{proof}

\begin{lemma}
\label{flat-lemma-finite-flat-weak-assassin-up-down}
设 $g : T \to S$ 是概形的有限平坦态射，$\mathcal{G}$ 是拟凝聚
$\mathcal{O}_S$-模。设 $t \in T$ 的像为 $s \in S$。则
$$
t \in \text{WeakAss}(g^*\mathcal{G})
\Leftrightarrow
s \in \text{WeakAss}(\mathcal{G})
$$
\end{lemma}

\begin{proof}
蕴含 ``$\Leftarrow$'' 立即由《因子》引理
\ref{divisors-lemma-weakly-ass-pullback} 得出。假设
$t \in \text{WeakAss}(g^*\mathcal{G})$。设
$\Spec(A) \subset S$ 是 $s$ 的仿射开邻域，且 $\mathcal{G}$ 是与
$A$-模 $M$ 相伴的拟凝聚层。设 $\mathfrak p \subset A$ 是与 $s$
对应的素理想。由于 $g$ 有限平坦，对某个有限平坦 $A$-代数 $B$，有
$g^{-1}(\Spec(A)) = \Spec(B)$。注意，$g^*\mathcal{G}$ 是与
$B$-模 $M \otimes_A B$ 相伴的拟凝聚 $\mathcal{O}_{\Spec(B)}$-模，
而 $g_*g^*\mathcal{G}$ 是与 $A$-模 $M \otimes_A B$ 相伴的拟凝聚
$\mathcal{O}_{\Spec(A)}$-模。由《代数》引理
\ref{algebra-lemma-finite-flat-local}，对某个整数 $n \geq 0$，有
$B_{\mathfrak p} \cong A_{\mathfrak p}^{\oplus n}$。由于假设至少存在
一个位于 $s$ 上方的 $T$ 中点，故 $n \geq 1$。于是考察茎可知
$$
s \in \text{WeakAss}(\mathcal{G})
\Leftrightarrow
s \in \text{WeakAss}(g_*g^*\mathcal{G})
$$
现在，由《因子》引理 \ref{divisors-lemma-weakly-associated-finite}，
假设 $t \in \text{WeakAss}(g^*\mathcal{G})$ 蕴含
$s \in \text{WeakAss}(g_*g^*\mathcal{G})$，故由上式有
$s \in  \text{WeakAss}(\mathcal{G})$。
\end{proof}

\begin{lemma}
\label{flat-lemma-etale-weak-assassin-up-down}
设 $h : U \to S$ 是概形的 \'etale 态射，$\mathcal{G}$ 是拟凝聚
$\mathcal{O}_S$-模。设 $u \in U$ 的像为 $s \in S$。则
$$
u \in \text{WeakAss}(h^*\mathcal{G})
\Leftrightarrow
s \in \text{WeakAss}(\mathcal{G})
$$
\end{lemma}

\begin{proof}
用 $s$ 与 $u$ 的仿射邻域分别替换 $S$ 与 $U$ 后，可假设 $g$
是仿射概形的标准 \'etale 态射；见《态射》引理
\ref{morphisms-lemma-etale-locally-standard-etale}。
% P05-EN-CH38-0003: frozen proof names g and X although the lemma's morphism and source are h and U.
因此可假设 $S = \Spec(A)$ 且 $X = \Spec(A[x, 1/g]/(f))$，其中 $f$
首一，且 $f'$ 在 $A[x, 1/g]$ 中可逆。注意，
$A[x, 1/g]/(f) = (A[x]/(f))_g$ 也是有限自由 $A$-代数
$A[x]/(f)$ 的局部化。因此可将 $U$ 看作概形
$T = \Spec(A[x]/(f))$ 的开子概形，而后者在 $S$ 上有限局部自由。
这就归结到上面的引理 \ref{flat-lemma-finite-flat-weak-assassin-up-down}。
\end{proof}

\begin{lemma}
\label{flat-lemma-weakly-associated-henselization}
设 $S$ 是概形，$s \in S$ 为一点。记 $\mathcal{O}_{S, s}^h$
（分别地，$\mathcal{O}_{S, s}^{sh}$）为亨泽尔化（分别地，严格亨泽尔化）；
见《代数》定义 \ref{algebra-definition-henselization}。
设 $\mathcal{F}$ 是拟凝聚 $\mathcal{O}_S$-模。下列条件等价：
\begin{enumerate}
\item $s$ 是 $\mathcal{F}$ 的弱伴随点；
\item $\mathfrak m_s$ 是 $\mathcal{F}_s$ 的弱伴随素理想；
\item $\mathfrak m_s^h$ 是
$\mathcal{F}_s \otimes_{\mathcal{O}_{S, s}} \mathcal{O}_{S, s}^h$
的弱伴随素理想；
\item $\mathfrak m_s^{sh}$ 是
$\mathcal{F}_s \otimes_{\mathcal{O}_{S, s}} \mathcal{O}_{S, s}^{sh}$
的弱伴随素理想。
\end{enumerate}
\end{lemma}

\begin{proof}
(1) 与 (2) 的等价性就是定义；见《因子》定义
\ref{divisors-definition-weakly-associated}。将《因子》引理
\ref{divisors-lemma-weakly-ass-pullback} 应用于下列平坦态射
（见《代数进阶》引理 \ref{more-algebra-lemma-dumb-properties-henselization}）：
$$
\Spec(\mathcal{O}_{S, s}^{sh}) \to
\Spec(\mathcal{O}_{S, s}^h) \to
\Spec(\mathcal{O}_{S, s})
$$
以及闭点，即得 (2) $\Rightarrow$ (3) $\Rightarrow$ (4)。
为证明 (4) $\Rightarrow$ (2)，可用一个仿射邻域替换 $S$。假设
$x \in \mathcal{F}_s \otimes_{\mathcal{O}_{S, s}} \mathcal{O}_{S, s}^{sh}$
是这样一个元素：其零化子之根基等于 $\mathfrak m_s^{sh}$
（见《代数》引理 \ref{algebra-lemma-weakly-ass-local}）。由《代数》引理
\ref{algebra-lemma-strict-henselization-different}，
$\mathcal{O}_{S, s}^{sh}$ 等于当 \'etale 邻域
$f : (U, u) \to (S, s)$ 变化时 $\mathcal{O}_{U, u}$ 的极限，
故可假设 $x$ 是某个
$x' \in \mathcal{F}_s \otimes_{\mathcal{O}_{S, s}} \mathcal{O}_{U, u}$
的像。局部环同态 $\mathcal{O}_{U, u} \to \mathcal{O}_{S, s}^{sh}$
忠实平坦（因为这是严格亨泽尔化），因而普遍单射（《代数》引理
\ref{algebra-lemma-faithfully-flat-universally-injective}）。于是 $x'$
的零化子是 $x$ 的零化子的逆像，故该零化子之根基等于
$\mathfrak m_u$。因此 $u$ 是 $f^*\mathcal{F}$ 的弱伴随点。
由引理 \ref{flat-lemma-etale-weak-assassin-up-down}，$s$ 是
$\mathcal{F}$ 的弱伴随点。
\end{proof}





\section{有限型模的局部结构}
\label{flat-section-local-structure-module}

\noindent
使本章许多论证得以成立的关键技术结果，是几何引理
\ref{flat-lemma-elementary-devissage}。

\begin{lemma}
\label{flat-lemma-sheaf-lives-on-subscheme}
设 $f : X \to S$ 是仿射概形的有限型态射，$\mathcal{F}$ 是有限型
拟凝聚 $\mathcal{O}_X$-模。设 $x \in X$ 在 $S$ 中的像为
$s = f(x)$，并置 $\mathcal{F}_s = \mathcal{F}|_{X_s}$。
则存在有限表示闭浸入 $i : Z \to X$ 与有限型拟凝聚
$\mathcal{O}_Z$-模 $\mathcal{G}$，使得 $i_*\mathcal{G} = \mathcal{F}$ 且
$Z_s = \text{Supp}(\mathcal{F}_s)$。
\end{lemma}

\begin{proof}
设态射 $f : X \to S$ 由环同态 $A \to B$ 给出，并设 $\mathcal{F}$
是与 $B$-模 $M$ 相伴的拟凝聚层。由《态射》引理
\ref{morphisms-lemma-locally-finite-type-characterize}，$A \to B$
是有限型环同态；由《性质》引理 \ref{properties-lemma-finite-type-module}，
$M$ 是有限 $B$-模。特别地，$\mathcal{F}$ 的支集是 $\Spec(B)$ 中由
$M$ 的零化子 $I = \{x \in B \mid xm = 0\ \forall m \in M\}$
截出的闭子概形；见《代数》引理 \ref{algebra-lemma-support-closed}。
设 $\mathfrak q \subset B$ 是与 $x$ 对应的素理想，
$\mathfrak p \subset A$ 是与 $s$ 对应的素理想。注意，
$X_s = \Spec(B \otimes_A \kappa(\mathfrak p))$，而 $\mathcal{F}_s$
是与 $B \otimes_A \kappa(\mathfrak p)$-模
$M \otimes_A \kappa(\mathfrak p)$ 相伴的拟凝聚层。由《态射》引理
\ref{morphisms-lemma-support-finite-type}，$\mathcal{F}_s$ 的支集等于
$V(I(B \otimes_A \kappa(\mathfrak p)))$。由于
$B \otimes_A \kappa(\mathfrak p)$ 在 $\kappa(\mathfrak p)$ 上有限型，
存在有限多个元素 $f_1, \ldots, f_m \in I$，使得
$$
I(B \otimes_A \kappa(\mathfrak p)) =
(f_1, \ldots, f_n)(B \otimes_A \kappa(\mathfrak p)).
$$
% P05-EN-CH38-0004: frozen display uses f_1,...,f_n after choosing f_1,...,f_m.
记 $i : Z \to X$ 为由 $(f_1, \ldots, f_m)$ 截出的闭子概形，即
$Z = \Spec(B/(f_1, \ldots, f_m))$。由于 $I$ 零化 $M$，可将 $M$
视为 $B/(f_1, \ldots, f_m)$-模。换言之，$\mathcal{F}$ 是 $Z$
上某个有限型模的推前。按构造有
$Z_s = \text{Supp}(\mathcal{F}_s)$，故引理得证。
\end{proof}

\begin{lemma}
\label{flat-lemma-elementary-devissage}
设 $f : X \to S$ 是局部有限型概形态射，$\mathcal{F}$ 是有限型
拟凝聚 $\mathcal{O}_X$-模。设 $x \in X$ 在 $S$ 中的像为
$s = f(x)$。置 $\mathcal{F}_s = \mathcal{F}|_{X_s}$ 且
$n = \dim_x(\text{Supp}(\mathcal{F}_s))$。则可构造：
\begin{enumerate}
\item 初等 \'etale 邻域 $g : (X', x') \to (X, x)$ 与
$e : (S', s') \to (S, s)$；
\item 交换图
$$
\xymatrix{
X \ar[dd]_f & X' \ar[dd] \ar[l]^g & Z' \ar[l]^i \ar[d]^\pi \\
& & Y' \ar[d]^h \\
S & S' \ar[l]_e & S' \ar@{=}[l]
}
$$
\item 点 $z' \in Z'$，满足 $i(z') = x'$、$y' = \pi(z')$、
$h(y') = s'$；
\item 有限型拟凝聚 $\mathcal{O}_{Z'}$-模 $\mathcal{G}$；
\end{enumerate}
使得下列性质成立：
\begin{enumerate}
\item $X'$、$Z'$、$Y'$、$S'$ 都是仿射概形；
\item $i$ 是有限表示闭浸入；
\item $i_*(\mathcal{G}) \cong g^*\mathcal{F}$；
\item $\pi$ 有限，且 $\pi^{-1}(\{y'\}) = \{z'\}$；
\item 扩张 $\kappa(y')/\kappa(s')$ 是纯超越的；
\item $h$ 是相对维数为 $n$ 且纤维几何整的光滑态射。
\end{enumerate}
\end{lemma}

\begin{proof}
设 $V \subset S$ 是 $s$ 的仿射邻域，$U \subset f^{-1}(V)$ 是 $x$
的仿射邻域。只需对 $f|_U : U \to V$ 与 $\mathcal{F}|_U$ 证明引理。
因此在余下证明中假设 $X$ 与 $S$ 均仿射。

\medskip\noindent
首先假设 $X_s = \text{Supp}(\mathcal{F}_s)$；特别地，
$n = \dim_x(X_s)$。应用《态射进阶》引理
\ref{more-morphisms-lemma-local-local-structure-finite-type} 与
\ref{more-morphisms-lemma-local-local-structure-finite-type-affine}，
得到交换图
$$
\xymatrix{
X \ar[dd] & X' \ar[l]^g \ar[d]^\pi \\
& Y' \ar[d]^h  \\
S & S' \ar[l]_e
}
$$
以及点 $x' \in X'$。置 $Z' = X'$、$i = \text{id}$ 且
$\mathcal{G} = g^*\mathcal{F}$，便得到此情形的解。

\medskip\noindent
一般地，依引理 \ref{flat-lemma-sheaf-lives-on-subscheme} 选取闭浸入
$Z \to X$ 与 $Z$ 上的层 $\mathcal{G}$。将上一段的结果应用于
$Z \to S$ 与 $\mathcal{G}$，得到图表
$$
\xymatrix{
X \ar[dd]_f & Z \ar[l] \ar[dd]_{f|_Z} & Z' \ar[l]^g \ar[d]^\pi \\
& & Y' \ar[d]^h \\
S & S \ar@{=}[l] & S' \ar[l]_e
}
$$
以及满足全部所需性质的点 $z' \in Z'$。我们将用引理
\ref{flat-lemma-lift-etale} 把 $Z'$ 嵌入一个在 $X$ 上 \'etale 的概形中。
由于要求 $X'$ 是 $S'$ 上的概形，不能直接应用该引理。改为考虑态射
$$
\xymatrix{
Z' \ar[r] & Z \times_S S' \ar[r] & X \times_S S'
}
$$
由《态射》引理 \ref{morphisms-lemma-etale-permanence}，第一个态射为
\'etale；第二个态射是闭浸入的基变换，因而为闭浸入。最后，由于
$X$、$S$、$S'$、$Z$、$Z'$ 均仿射，可应用引理
\ref{flat-lemma-lift-etale}，得到仿射概形的 \'etale 态射
$X' \to X \times_S S'$，使得
$$
Z' = (Z \times_S S') \times_{(X \times_S S')} X' = Z \times_X X'.
$$
由于 $Z \to X$ 是有限表示闭浸入，$Z' \to X'$ 亦然。设
$x' \in X'$ 为与 $z' \in Z'$ 对应的点。于是补全后的图表
$$
\xymatrix{
X \ar[dd] & X' \ar[dd] \ar[l] & Z' \ar[l]^i \ar[d]^\pi \\
& & Y' \ar[d]^h \\
S & S' \ar[l]_e & S' \ar@{=}[l]
}
$$
就是原问题的一个解。
\end{proof}

\begin{lemma}
\label{flat-lemma-devissage-finite-presentation}
假设与记号同引理 \ref{flat-lemma-elementary-devissage}。若 $f$ 局部有限表示，
则 $\pi$ 有限表示。在此情形，下列条件等价：
\begin{enumerate}
\item $\mathcal{F}$ 在 $x$ 的某个邻域上是有限表示
$\mathcal{O}_X$-模；
\item $\mathcal{G}$ 在 $z'$ 的某个邻域上是有限表示
$\mathcal{O}_{Z'}$-模；
\item $\pi_*\mathcal{G}$ 在 $y'$ 的某个邻域上是有限表示
$\mathcal{O}_{Y'}$-模。
\end{enumerate}
仍假设 $f$ 局部有限表示，则下列条件也彼此等价：
\begin{enumerate}
\item[(a)] $\mathcal{F}_x$ 是有限表示 $\mathcal{O}_{X, x}$-模；
\item[(b)] $\mathcal{G}_{z'}$ 是有限表示 $\mathcal{O}_{Z', z'}$-模；
\item[(c)] $(\pi_*\mathcal{G})_{y'}$ 是有限表示
$\mathcal{O}_{Y', y'}$-模。
\end{enumerate}
\end{lemma}

\begin{proof}
假设 $f$ 局部有限表示。由于 $Z' \to S$ 是这类态射的复合，它也局部
有限表示；见《态射》引理
\ref{morphisms-lemma-composition-finite-presentation}。注意，
$Y' \to S$ 是光滑态射与 \'etale 态射的复合，因而也局部有限表示。
故《态射》引理 \ref{morphisms-lemma-finite-presentation-permanence}
蕴含 $\pi$ 局部有限表示。由于 $\pi$ 有限，它也分离且拟紧，
因而 $\pi$ 实际上有限表示。

\medskip\noindent
为证明 (1)、(2)、(3) 等价，还考虑：
(4) $g^*\mathcal{F}$ 在 $x'$ 的某个邻域上是有限表示
$\mathcal{O}_{X'}$-模。有限表示模的拉回仍有限表示；见模章引理
\ref{modules-lemma-pullback-finite-presentation}。故
(1) $\Rightarrow$ (4)。\'Etale 态射 $g$ 是开的；见《态射》引理
\ref{morphisms-lemma-etale-open}。因此，对 $x'$ 的任意开邻域
$U' \subset X'$，像 $g(U')$ 是 $x$ 的开邻域，且
$\{U' \to g(U')\}$ 是 \'etale 覆盖。于是由下降章引理
\ref{descent-lemma-finite-presentation-descends}，有
(4) $\Rightarrow$ (1)。利用下降章引理
\ref{descent-lemma-finite-finitely-presented-module} 与一些简单拓扑论证
（见《态射进阶》引理
\ref{more-morphisms-lemma-finite-morphism-single-point-in-fibre}），可知
(4) $\Leftrightarrow$ (2) $\Leftrightarrow$ (3)。

\medskip\noindent
为证明 (a)、(b)、(c) 等价，考虑环同态
$$
\mathcal{O}_{X, x} \to
\mathcal{O}_{X', x'} \to
\mathcal{O}_{Z', z'} \leftarrow
\mathcal{O}_{Y', y'}
$$
第一个环同态忠实平坦。因此 $\mathcal{F}_x$ 在
$\mathcal{O}_{X, x}$ 上有限表示，当且仅当 $g^*\mathcal{F}_{x'}$
在 $\mathcal{O}_{X', x'}$ 上有限表示；见《代数》引理
\ref{algebra-lemma-descend-properties-modules}。第二个环同态满射
（故有限），且按假设有限表示，因此 $g^*\mathcal{F}_{x'}$
在 $\mathcal{O}_{X', x'}$ 上有限表示，当且仅当 $\mathcal{G}_{z'}$
在 $\mathcal{O}_{Z', z'}$ 上有限表示；见《代数》引理
\ref{algebra-lemma-finite-finitely-presented-extension}。由于 $\pi$ 有限、
有限表示且 $\pi^{-1}(\{y'\}) = \{x'\}$，环同态
$\mathcal{O}_{Y', y'} \leftarrow \mathcal{O}_{Z', z'}$ 有限且有限表示；
见《态射进阶》引理
\ref{more-morphisms-lemma-finite-morphism-single-point-in-fibre}。
% P05-EN-CH38-0005: frozen proof writes fibre {x'} although pi has domain Z' and the distinguished point is z'.
因此 $\mathcal{G}_{z'}$ 在 $\mathcal{O}_{Z', z'}$ 上有限表示，
当且仅当 $\pi_*\mathcal{G}_{y'}$ 在 $\mathcal{O}_{Y', y'}$ 上有限表示；
见《代数》引理 \ref{algebra-lemma-finite-finitely-presented-extension}。
\end{proof}

\begin{lemma}
\label{flat-lemma-devissage-flat}
假设与记号同引理 \ref{flat-lemma-elementary-devissage}。下列条件等价：
\begin{enumerate}
\item $\mathcal{F}$ 在 $x$ 的某个邻域上于 $S$ 上平坦；
\item $\mathcal{G}$ 在 $z'$ 的某个邻域上于 $S'$ 上平坦；
\item $\pi_*\mathcal{G}$ 在 $y'$ 的某个邻域上于 $S'$ 上平坦。
\end{enumerate}
下列条件也等价：
\begin{enumerate}
\item[(a)] $\mathcal{F}_x$ 在 $\mathcal{O}_{S, s}$ 上平坦；
\item[(b)] $\mathcal{G}_{z'}$ 在 $\mathcal{O}_{S', s'}$ 上平坦；
\item[(c)] $(\pi_*\mathcal{G})_{y'}$ 在 $\mathcal{O}_{S', s'}$ 上平坦。
\end{enumerate}
\end{lemma}

\begin{proof}
为证明 (1)、(2)、(3) 等价，还考虑：
(4) $g^*\mathcal{F}$ 在 $x'$ 的某个邻域上于 $S$ 上平坦。
我们将使用引理 \ref{flat-lemma-etale-flat-up-down}，不再说明地把在 $S$
上平坦与在 $S'$ 上平坦等同起来。\'Etale 态射 $g$ 平坦且为开映射；
见《态射》引理 \ref{morphisms-lemma-etale-open}。因此，对 $x'$ 的任意
开邻域 $U' \subset X'$，像 $g(U')$ 是 $x$ 的开邻域，且态射
$U' \to g(U')$ 满且平坦。故由《态射》引理
\ref{morphisms-lemma-flat-permanence}，有 (4) $\Leftrightarrow$ (1)。注意
$$
\Gamma(X', g^*\mathcal{F}) =
\Gamma(Z', \mathcal{G}) =
\Gamma(Y', \pi_*\mathcal{G})
$$
因此，$g^*\mathcal{F}$、$\mathcal{G}$ 与 $\pi_*\mathcal{G}$ 在
$S'$ 上平坦三者彼此等价（这里使用了 $X'$、$Z'$、$Y'$、$S'$
均仿射）。再用关于仿射邻域的若干略去的拓扑论证（比较《态射进阶》引理
\ref{more-morphisms-lemma-finite-morphism-single-point-in-fibre}），便得到
(4) $\Leftrightarrow$ (2) $\Leftrightarrow$ (3)。

\medskip\noindent
为证明 (a)、(b)、(c) 等价，考虑局部环同态的交换图
$$
\xymatrix{
\mathcal{O}_{X', x'} \ar[r]_\iota &
\mathcal{O}_{Z', z'} &
\mathcal{O}_{Y', y'} \ar[l]^\alpha &
\mathcal{O}_{S', s'} \ar[l]^\beta \\
\mathcal{O}_{X, x} \ar[u]^\gamma & & &
\mathcal{O}_{S, s} \ar[lll]_\varphi \ar[u]_\epsilon
}
$$
我们将使用引理 \ref{flat-lemma-etale-flat-up-down-local-ring}，不再说明地把
在 $\mathcal{O}_{S, s}$ 上平坦与在 $\mathcal{O}_{S', s'}$ 上平坦
等同起来。同态 $\gamma$ 忠实平坦。因此 $\mathcal{F}_x$ 在
$\mathcal{O}_{S, s}$ 上平坦，当且仅当 $g^*\mathcal{F}_{x'}$
在 $\mathcal{O}_{S', s'}$ 上平坦；见《代数》引理
\ref{algebra-lemma-flatness-descends-more-general}。作为
$\mathcal{O}_{S', s'}$-模，$g^*\mathcal{F}_{x'}$、
$\mathcal{G}_{z'}$ 与 $\pi_*\mathcal{G}_{y'}$ 全都同构；见《态射进阶》引理
\ref{more-morphisms-lemma-finite-morphism-single-point-in-fibre}。
证明完毕。
\end{proof}









\section{一步 d\'evissage}
\label{flat-section-one-step-devissage}

\noindent
本节解释模的一步 d\'evissage 是什么。一步 d\'evissage 在底与靶上
\'etale 局部存在。我们讨论一步 d\'evissage 的基变换、Zariski 缩小与
\'etale 局部化。

\begin{definition}
\label{flat-definition-one-step-devissage}
设 $S$ 是概形，$X$ 在 $S$ 上局部有限型，$\mathcal{F}$ 是有限型
拟凝聚 $\mathcal{O}_X$-模，且 $s \in S$ 为一点。所谓
{\it $s$ 上 $\mathcal{F}/X/S$ 的一步 d\'evissage}，是指给定
$S$ 上概形的态射
$$
\xymatrix{
X & Z \ar[l]_i \ar[r]^\pi & Y
}
$$
以及有限型拟凝聚 $\mathcal{O}_Z$-模 $\mathcal{G}$，满足：
\begin{enumerate}
\item $X$、$S$、$Z$、$Y$ 均仿射；
\item $i$ 是有限表示闭浸入；
\item $\mathcal{F} \cong i_*\mathcal{G}$；
\item $\pi$ 有限；
\item 结构态射 $Y \to S$ 光滑，且其纤维几何不可约、维数为
$\dim(\text{Supp}(\mathcal{F}_s))$。
\end{enumerate}
此时称 $(Z, Y, i, \pi, \mathcal{G})$ 为 $s$ 上
$\mathcal{F}/X/S$ 的一步 d\'evissage。
\end{definition}

\noindent
注意，这样的一步 d\'evissage 只有在 $X$ 与 $S$ 仿射时才可能存在。
上述定义只要求 $X$ 在 $S$ 上（局部）有限型，下文仍在这一设定中工作。
在 \cite{GruRay} 中，作者始终采用如下设定：$X \to S$ 局部有限表示，
而 $\mathcal{F}$ 是有限型拟凝聚 $\mathcal{O}_X$-模。这样选择的好处是，
询问 $\mathcal{F}$ 作为 $\mathcal{O}_X$-模是否有限表示是“有意义”的，
而在我们的设定中则“没有意义”。有关讨论见《态射进阶》第
\ref{more-morphisms-section-finite-type-finite-presentation} 节；
其中的观察表明，在我们的设定下，可考虑 $\mathcal{F}$ “相对于 $S$
局部有限表示”这一条件，并可始终使用该概念。不过，我们将依靠引理
\ref{flat-lemma-devissage-finite-presentation} 的结果与注记
\ref{flat-remark-finite-presentation} 中的观察，在问题出现时逐案处理。

\begin{definition}
\label{flat-definition-one-step-devissage-at-x}
设 $S$ 是概形，$X$ 在 $S$ 上局部有限型，$\mathcal{F}$ 是有限型
拟凝聚 $\mathcal{O}_X$-模。设 $x \in X$ 在 $S$ 中的像为 $s$。
所谓 {\it $x$ 处 $\mathcal{F}/X/S$ 的一步 d\'evissage}，是指系统
$(Z, Y, i, \pi, \mathcal{G}, z, y)$，其中
$(Z, Y, i, \pi, \mathcal{G})$ 是 $s$ 上 $\mathcal{F}/X/S$
的一步 d\'evissage，并且
\begin{enumerate}
\item $\dim_x(\text{Supp}(\mathcal{F}_s)) = \dim(\text{Supp}(\mathcal{F}_s))$,
\item $z \in Z$ 是满足 $i(z) = x$ 与 $\pi(z) = y$ 的点；
\item $\pi^{-1}(\{y\}) = \{z\}$；
\item 扩张 $\kappa(y)/\kappa(s)$ 是纯超越的。
\end{enumerate}
\end{definition}

\noindent
$x$ 处 $\mathcal{F}/X/S$ 的一步 d\'evissage 只有在 $X$ 与 $S$
仿射时才可能存在。结合定义 \ref{flat-definition-one-step-devissage} 的条件
(5)，条件 (1) 保证 $Y \to S$ 的相对维数等于
$\dim_x(\text{Supp}(\mathcal{F}_s))$。

\begin{lemma}
\label{flat-lemma-elementary-devissage-variant}
设 $f : X \to S$ 是局部有限型概形态射，$\mathcal{F}$ 是有限型
拟凝聚 $\mathcal{O}_X$-模，且 $x \in X$ 在 $S$ 中的像为
$s = f(x)$。则存在带点概形的交换图
$$
\xymatrix{
(X, x) \ar[d]_f & (X', x') \ar[l]^g \ar[d] \\
(S, s) & (S', s') \ar[l] \\
}
$$
使得 $(S', s') \to (S, s)$ 与 $(X', x') \to (X, x)$ 是初等
\'etale 邻域，且 $g^*\mathcal{F}/X'/S'$ 在 $x'$ 处有一步
d\'evissage。
\end{lemma}

\begin{proof}
这立即由定义 \ref{flat-definition-one-step-devissage-at-x} 与引理
\ref{flat-lemma-elementary-devissage} 得出。
\end{proof}

\begin{lemma}
\label{flat-lemma-base-change-one-step}
设 $S$、$X$、$\mathcal{F}$、$s$ 如定义
\ref{flat-definition-one-step-devissage}，且
$(Z, Y, i, \pi, \mathcal{G})$ 是 $s$ 上 $\mathcal{F}/X/S$
的一步 d\'evissage。设 $(S', s') \to (S, s)$ 是任意带点概形态射。
由这些数据，令 $X', Z', Y', i', \pi'$ 分别为
$X, Z, Y, i, \pi$ 沿 $S' \to S$ 的基变换；令 $\mathcal{F}'$
为 $\mathcal{F}$ 到 $X'$ 的拉回，$\mathcal{G}'$ 为
$\mathcal{G}$ 到 $Z'$ 的拉回。若 $S'$ 仿射，则
$(Z', Y', i', \pi', \mathcal{G}')$ 是 $s'$ 上
$\mathcal{F}'/X'/S'$ 的一步 d\'evissage。
\end{lemma}

\begin{proof}
仿射概形的纤维积仿射；见概形章引理
\ref{schemes-lemma-fibre-product-affines}。基变换保持闭浸入、有限表示态射、
有限态射、光滑态射、纤维几何不可约的态射与相对维数为 $n$ 的态射；
见《态射》诸引理 \ref{morphisms-lemma-base-change-closed-immersion}、
\ref{morphisms-lemma-base-change-finite-presentation},
\ref{morphisms-lemma-base-change-finite},
\ref{morphisms-lemma-base-change-smooth},
\ref{morphisms-lemma-base-change-relative-dimension-d} 以及《态射进阶》引理
\ref{more-morphisms-lemma-base-change-fibres-geometrically-irreducible}.
由于沿有限态射 $i$ 的推前与基变换交换，有
$i'_*\mathcal{G}' \cong \mathcal{F}'$；见概形上同调章引理
\ref{coherent-lemma-affine-base-change}。又有
$\dim(\text{Supp}(\mathcal{F}_s)) = \dim(\text{Supp}(\mathcal{F}'_{s'}))$
这是因为下式成立，并可应用《态射》引理
\ref{morphisms-lemma-dimension-fibre-after-base-change}：
$$
\text{Supp}(\mathcal{F}_s) \times_s s' = \text{Supp}(\mathcal{F}'_{s'}).
$$
引理得证。
\end{proof}

\begin{lemma}
\label{flat-lemma-base-change-one-step-at-x}
设 $S$、$X$、$\mathcal{F}$、$x$、$s$ 如定义
\ref{flat-definition-one-step-devissage-at-x}，且
$(Z, Y, i, \pi, \mathcal{G}, z, y)$ 是 $x$ 处
$\mathcal{F}/X/S$ 的一步 d\'evissage。设带点概形态射
$(S', s') \to (S, s)$ 诱导同构 $\kappa(s) = \kappa(s')$。
令 $(Z', Y', i', \pi', \mathcal{G}')$ 如引理
\ref{flat-lemma-base-change-one-step} 中所构造，并令 $x' \in X'$
（分别地，$z' \in Z'$、$y' \in Y'$）为同时映到
$x \in X$（分别地，$z \in Z$、$y \in Y$）与 $s' \in S'$
的唯一一点。若 $S'$ 仿射，则
$(Z', Y', i', \pi', \mathcal{G}', z', y')$ 是 $x'$ 处
$\mathcal{F}'/X'/S'$ 的一步 d\'evissage。
\end{lemma}

\begin{proof}
由引理 \ref{flat-lemma-base-change-one-step}，
$(Z', Y', i', \pi', \mathcal{G}')$ 是 $s'$ 上
$\mathcal{F}'/X'/S'$ 的一步 d\'evissage。由于假设
$\kappa(s) = \kappa(s')$ 保证纤维 $X'_{s'}$、$Z'_{s'}$、$Y'_{s'}$
分别同构于 $X_s$、$Z_s$、$Y_s$，定义
\ref{flat-definition-one-step-devissage-at-x} 的性质 (1) -- (4) 对
$(Z', Y', i', \pi', \mathcal{G}', z', y')$ 成立。
\end{proof}

\begin{definition}
\label{flat-definition-shrink}
设 $S$、$X$、$\mathcal{F}$、$x$、$s$ 如定义
\ref{flat-definition-one-step-devissage-at-x}，且
$(Z, Y, i, \pi, \mathcal{G}, z, y)$ 是 $x$ 处
$\mathcal{F}/X/S$ 的一步 d\'evissage。此情形的一个
{\it 标准缩小}，是指给定标准开集 $S' \subset S$、$X' \subset X$、
$Z' \subset Z$、$Y' \subset Y$，满足 $s \in S'$、$x \in X'$、
$z \in Z'$、$y \in Y'$，并且
$$
(Z', Y', i|_{Z'}, \pi|_{Z'}, \mathcal{G}|_{Z'}, z, y)
$$
是 $x$ 处 $\mathcal{F}|_{X'}/X'/S'$ 的一步 d\'evissage。
\end{definition}

\begin{lemma}
\label{flat-lemma-shrink}
假设与记号同定义 \ref{flat-definition-shrink}。则：
\begin{enumerate}
\item
\label{flat-item-shrink-base}
若 $S' \subset S$ 是 $s$ 的标准开邻域，则置
$X' = X_{S'}$、$Z' = Z_{S'}$、$Y' = Y_{S'}$，得到一个标准缩小。
\item
\label{flat-item-shrink-on-Y}
设 $W \subset Y$ 是 $y$ 的标准开邻域。则存在满足
$Y' = W \times_S S'$ 的标准缩小。
\item
\label{flat-item-shrink-on-X}
设 $U \subset X$ 是 $x$ 的开邻域。则存在满足 $X' \subset U$
的标准缩小。
\end{enumerate}
\end{lemma}

\begin{proof}
(1) 立即由引理 \ref{flat-lemma-base-change-one-step-at-x} 与如下事实得出：
仿射概形态射下标准开集的逆像仍为标准开集；见《代数》引理
\ref{algebra-lemma-spec-functorial}。

\medskip\noindent
设 $W \subset Y$ 如 (2)。由于 $Y \to S$ 光滑，它是开映射；见《态射》引理
\ref{morphisms-lemma-smooth-open}。故可找到包含于 $W$ 之像中的
$s$ 的标准开邻域 $S'$。于是 $W_{S'} \to S'$ 的纤维是 $Y \to S$
在 $S'$ 上各纤维的非空开子概形，因而也几何不可约。置
$Y' = W_{S'}$ 与 $Z' = \pi^{-1}(Y')$，可见 $Z' \subset Z$
是 $z$ 的标准开邻域。取函数 $\overline{h} \in \Gamma(Z, \mathcal{O}_Z)$
使 $Z' = D(\overline{h})$。由于 $i : Z \to X$ 是闭浸入，可找到
$h \in \Gamma(X, \mathcal{O}_X)$ 使得
$i^\sharp(h) = \overline{h}$。取 $X' = D(h) \subset X$，
便得到 (2) 所述的标准缩小。

\medskip\noindent
设 $U \subset X$ 如 (3)。缩小 $U$ 后，可假设 $U$ 是标准开集。
由《态射进阶》引理
\ref{more-morphisms-lemma-finite-morphism-single-point-in-fibre}，
存在 $y$ 的标准开邻域 $W \subset Y$，使得
$\pi^{-1}(W) \subset i^{-1}(U)$。应用 (2)，得到标准缩小
$X', S', Z', Y'$，其中 $Y' = W_{S'}$。由于
$Z' \subset \pi^{-1}(W) \subset i^{-1}(U)$，可用 $X' \cap U$
替换 $X'$（因 $U$ 也是标准开集，交集仍为标准开集），而不破坏
标准缩小定义中的任何条件。结论得证。
\end{proof}

\begin{lemma}
\label{flat-lemma-elementary-etale-neighbourhood}
设 $S$、$X$、$\mathcal{F}$、$x$、$s$ 如定义
\ref{flat-definition-one-step-devissage-at-x}，且
$(Z, Y, i, \pi, \mathcal{G}, z, y)$ 是 $x$ 处
$\mathcal{F}/X/S$ 的一步 d\'evissage。设
$$
\xymatrix{
(Y, y) \ar[d] & (Y', y') \ar[l] \ar[d] \\
(S, s) & (S', s') \ar[l]
}
$$
是带点概形的交换图，且水平箭头为初等 \'etale 邻域。则存在交换图
$$
\xymatrix{
& & (X'', x'') \ar[lld] \ar[d] & (Z'', z'') \ar[l] \ar[lld] \ar[d] \\
(X, x) \ar[d] & (Z, z) \ar[l] \ar[d] &
(S'', s'') \ar[lld] & (Y'', y'') \ar[lld] \ar[l] \\
(S, s) & (Y, y) \ar[l]
}
$$
满足下列性质：
\begin{enumerate}
\item $(S'', s'') \to (S', s')$ 是初等 \'etale 邻域，且态射
$S'' \to S$ 是复合 $S'' \to S' \to S$；
\item $Y''$ 是 $Y' \times_{S'} S''$ 的开子概形；
\item $Z'' = Z \times_Y Y''$,
\item $(X'', x'') \to (X, x)$ 是初等 \'etale 邻域；
\item $(Z'', Y'', i'', \pi'', \mathcal{G}'', z'', y'')$ 是层
$\mathcal{F}''$ 在 $x''$ 处的一步 d\'evissage。
\end{enumerate}
这里 $\mathcal{F}''$（分别地，$\mathcal{G}''$）是 $\mathcal{F}$
（分别地，$\mathcal{G}$）沿态射 $X'' \to X$（分别地，$Z'' \to Z$）
的拉回，而 $i'' : Z'' \to X''$ 与 $\pi'' : Z'' \to Y''$ 如图所示。
\end{lemma}

\begin{proof}
任取初等 \'etale 邻域 $(S'', s'') \to (S', s')$，其中 $S''$ 仿射。
任取包含点 $y'' = (y', s'')$ 的仿射开邻域
$Y'' \subset Y' \times_{S'} S''$。由 (3) 得到仿射的
$(Z'', z'')$。此外，$Z_{S''} \to X_{S''}$ 是闭浸入，而
$Z'' \to Z_{S''}$ 是 \'etale 态射。因此可应用引理
\ref{flat-lemma-lift-etale}，找到仿射概形的 \'etale 态射
$X'' \to X_{S'}$，使 $Z'' \cong X'' \times_{X_{S'}} Z_{S'}$。
% P05-EN-CH38-0006: frozen proof constructs over S'' but then writes X_{S'} and Z_{S'}.
记 $i'' : Z'' \to X''$ 为相应的闭浸入。置 $x'' = i''(z'')$，
得到引理所述的交换图。性质 (1)、(2)、(3)、(4) 按构造成立。
故只需为适当选择的 $(S'', s'') \to (S', s')$ 与 $Y''$ 证明 (5)。

\medskip\noindent
先列出对第一段中任意选择的 $(S'', s'') \to (S', s')$ 与 $Y''$
都成立的性质。按构造有 $Z'' = X'' \times_X Z$，故
$i''_*\mathcal{G}'' = \mathcal{F}''$（记号同引理陈述）；见概形上同调章
引理 \ref{coherent-lemma-affine-base-change}。置
$n = \dim(\text{Supp}(\mathcal{F}_s)) =
\dim_x(\text{Supp}(\mathcal{F}_s))$。态射 $Y'' \to S''$ 光滑且相对
维数为 $n$：事实上，$Y' \to S'$ 是相对维数为 $n$ 的 \'etale 态射与
相对维数为 $n$ 的光滑态射之复合 $Y' \to Y_{S'} \to S'$，且基变换保持相对维数为
$n$ 的光滑态射。又有 $\kappa(y'') = \kappa(y)$ 与
$\kappa(s) = \kappa(s'')$，故 $\kappa(y'')$ 是
$\kappa(s'')$ 的纯超越扩张。纤维态射 $X''_{s''} \to X_s$ 是
$\kappa(s) = \kappa(s'')$ 上仿射概形的 \'etale 态射，把点 $x''$
映到点 $x$，并把 $\mathcal{F}_s$ 拉回为 $\mathcal{F}''_{s''}$。因此
$$
\dim(\text{Supp}(\mathcal{F}''_{s''})) =
\dim(\text{Supp}(\mathcal{F}_s)) = n =
\dim_x(\text{Supp}(\mathcal{F}_s)) =
\dim_{x''}(\text{Supp}(\mathcal{F}''_{s''}))
$$
这是因为维数在 \'etale 局部化下不变；见下降章引理
\ref{descent-lemma-dimension-at-point-local}。由于
$\pi'' : Z'' \to Y''$ 是 $\pi$ 的基变换，$\pi''$ 有限；又因
$\kappa(y) = \kappa(y'')$，有 $\pi^{-1}(\{y''\}) = \{z''\}$。
% P05-EN-CH38-0007: frozen proof uses pi^{-1} on a point of Y''; the base-changed map is pi''.

\medskip\noindent
至此，定义 \ref{flat-definition-one-step-devissage} 的所有条件均已验证，
唯有尚未验证 $Y'' \to S''$ 的纤维几何不可约。一般而言这当然未必成立，
此处将使用 $\kappa(s) \subset \kappa(y)$ 为纯超越扩张。具体地，
设 $T \subset Y'_{s'}$ 为包含
$y' = (y, s')$ 的不可约分支。由于 $Y'_{s'}$ 在 $\kappa(s')$ 上光滑，
$T$ 是 $Y'_{s'}$ 的开子概形。又因 $\kappa(s') = \kappa(s)$
在 $\kappa(y') = \kappa(y)$ 中代数闭，由簇章引理
\ref{varieties-lemma-geometrically-connected-if-connected-and-point}，
$T$ 几何连通。由于 $T$ 光滑，$T$ 几何不可约。于是可应用《态射进阶》引理
\ref{more-morphisms-lemma-normal-morphism-irreducible}，找到初等
\'etale 态射 $(S'', s'') \to (S', s')$ 与仿射开集
$Y'' \subset Y'_{S''}$，使得 $Y'' \to S''$ 的全部纤维几何不可约，
且 $T = Y''_{s''}$。缩小（先缩小 $Y''$，再缩小 $S''$）后，可假设
$Y''$ 与 $S''$ 均仿射。引理得证。
\end{proof}

\begin{lemma}
\label{flat-lemma-existence-alpha}
设 $S$、$X$、$\mathcal{F}$、$s$ 如定义
\ref{flat-definition-one-step-devissage}，且
$(Z, Y, i, \pi, \mathcal{G})$ 是 $s$ 上 $\mathcal{F}/X/S$
的一步 d\'evissage。设 $\xi \in Y_s$ 为（唯一的）一般点。
则存在整数 $r > 0$ 与 $\mathcal{O}_Y$-模同态
$$
\alpha : \mathcal{O}_Y^{\oplus r} \longrightarrow \pi_*\mathcal{G}
$$
使得
$$
\alpha :
\kappa(\xi)^{\oplus r}
\longrightarrow
(\pi_*\mathcal{G})_\xi \otimes_{\mathcal{O}_{Y, \xi}} \kappa(\xi)
$$
为同构。此外，此时有
$$
\dim(\text{Supp}(\Coker(\alpha)_s)) < \dim(\text{Supp}(\mathcal{F}_s)).
$$
\end{lemma}

\begin{proof}
按假设，概形 $S$ 与 $Y$ 仿射。写作 $S = \Spec(A)$ 与
$Y = \Spec(B)$。由于 $\pi$ 有限，$\mathcal{O}_Y$-模
$\pi_*\mathcal{G}$ 是有限型拟凝聚 $\mathcal{O}_Y$-模。因此对某个有限
$B$-模 $N$，有 $\pi_*\mathcal{G} = \widetilde{N}$。设
$\mathfrak p \subset B$ 为与 $\xi$ 对应的素理想。为得到 $\alpha$，
置 $r = \dim_{\kappa(\mathfrak p)} N \otimes_B \kappa(\mathfrak p)$，
并选取 $x_1, \ldots, x_r \in N$，使其构成
$N \otimes_B \kappa(\mathfrak p)$ 的一组基。取
$\alpha : B^{\oplus r} \to N$ 为公式
$\alpha(b_1, \ldots, b_r) = \sum b_ix_i$ 给出的同态。显然，
$\alpha : \kappa(\mathfrak p)^{\oplus r} \to
N \otimes_B \kappa(\mathfrak p)$ 如所需为同构。最后，设 $\alpha$
是任意具有此性质的同态。则 $N' = \Coker(\alpha)$ 是有限 $B$-模，
且 $N' \otimes \kappa(\mathfrak p) = 0$。由 Nakayama 引理
（《代数》引理 \ref{algebra-lemma-NAK}），有
$N'_{\mathfrak p} = 0$。纤维 $Y_s$ 几何不可约、维数为 $n$，其一般点为
$\xi$；而刚刚已经看到 $\xi$ 不在 $\Coker(\alpha)$ 的支集中，
故引理最后一个断言成立。
\end{proof}


\section{完全 d\'evissage}
\label{flat-section-complete-devissage}

\noindent
本节解释模的完全 d\'evissage 是什么，并证明它存在。本节内容主要是记账工作。

\begin{definition}
\label{flat-definition-complete-devissage}
设 $S$ 是概形，$X$ 在 $S$ 上局部有限型，$\mathcal{F}$ 是有限型
拟凝聚 $\mathcal{O}_X$-模，且 $s \in S$ 为一点。所谓
{\it $s$ 上 $\mathcal{F}/X/S$ 的完全 d\'evissage}，是指给定图表
$$
\xymatrix{
X & Z_1 \ar[l]^{i_1} \ar[d]^{\pi_1} \\
& Y_1 & Z_2 \ar[l]^{i_2} \ar[d]^{\pi_2} \\
& & Y_2 & Z_3 \ar[l] \ar[d] \\
& & & ... & ... \ar[l] \ar[d] \\
& & & & Y_n
}
$$
其中各概形均在 $S$ 上，并给定有限型拟凝聚
$\mathcal{O}_{Z_k}$-模 $\mathcal{G}_k$ 与 $\mathcal{O}_{Y_k}$-模同态
$$
\alpha_k :
\mathcal{O}_{Y_k}^{\oplus r_k}
\longrightarrow
\pi_{k, *}\mathcal{G}_k,
\quad
k = 1, \ldots, n
$$
满足下列性质：
\begin{enumerate}
\item $(Z_1, Y_1, i_1, \pi_1, \mathcal{G}_1)$ 是 $s$ 上
$\mathcal{F}/X/S$ 的一步 d\'evissage；
\item 同态 $\alpha_k$ 诱导同构
$$
\kappa(\xi_k)^{\oplus r_k} \longrightarrow
(\pi_{k, *}\mathcal{G}_k)_{\xi_k}
\otimes_{\mathcal{O}_{Y_k, \xi_k}} \kappa(\xi_k)
$$
其中 $\xi_k \in (Y_k)_s$ 是唯一的一般点；
\item 对 $k = 2, \ldots, n$，系统
$(Z_k, Y_k, i_k, \pi_k, \mathcal{G}_k)$
是 $s$ 上 $\Coker(\alpha_{k - 1})/Y_{k - 1}/S$ 的一步 d\'evissage；
\item $\Coker(\alpha_n) = 0$.
\end{enumerate}
此时称
$(Z_k, Y_k, i_k, \pi_k, \mathcal{G}_k, \alpha_k)_{k = 1, \ldots, n}$
为 $s$ 上 $\mathcal{F}/X/S$ 的完全 d\'evissage。
\end{definition}

\begin{definition}
\label{flat-definition-complete-devissage-at-x}
设 $S$ 是概形，$X$ 在 $S$ 上局部有限型，$\mathcal{F}$ 是有限型
拟凝聚 $\mathcal{O}_X$-模。设 $x \in X$ 的像为 $s \in S$。
所谓 {\it $x$ 处 $\mathcal{F}/X/S$ 的完全 d\'evissage}，是指给定系统
$$
(Z_k, Y_k, i_k, \pi_k, \mathcal{G}_k, \alpha_k, z_k, y_k)_{k = 1, \ldots, n}
$$
使得 $(Z_k, Y_k, i_k, \pi_k, \mathcal{G}_k, \alpha_k)$ 是 $s$ 上
$\mathcal{F}/X/S$ 的完全 d\'evissage，并且
\begin{enumerate}
\item $(Z_1, Y_1, i_1, \pi_1, \mathcal{G}_1, z_1, y_1)$ 是 $x$ 处
$\mathcal{F}/X/S$ 的一步 d\'evissage；
\item 对 $k = 2, \ldots, n$，系统
$(Z_k, Y_k, i_k, \pi_k, \mathcal{G}_k, z_k, y_k)$
是 $y_{k - 1}$ 处 $\Coker(\alpha_{k - 1})/Y_{k - 1}/S$
的一步 d\'evissage。
\end{enumerate}
\end{definition}

\noindent
再次注意，完全 d\'evissage 只有在 $X$ 与 $S$ 仿射时才可能存在。

\begin{lemma}
\label{flat-lemma-base-change-complete}
设 $S$、$X$、$\mathcal{F}$、$s$ 如定义
\ref{flat-definition-complete-devissage}，且
$(S', s') \to (S, s)$ 是任意带点概形态射。设
$(Z_k, Y_k, i_k, \pi_k, \mathcal{G}_k, \alpha_k)_{k = 1, \ldots, n}$
是 $s$ 上 $\mathcal{F}/X/S$ 的完全 d\'evissage。由这些数据，
令 $X', Z'_k, Y'_k, i'_k, \pi'_k$ 分别为
$X, Z_k, Y_k, i_k, \pi_k$ 沿 $S' \to S$ 的基变换；令
$\mathcal{F}'$ 为 $\mathcal{F}$ 到 $X'$ 的拉回，$\mathcal{G}'_k$
为 $\mathcal{G}_k$ 到 $Z'_k$ 的拉回，$\alpha'_k$ 为
$\alpha_k$ 到 $Y'_k$ 的拉回。若 $S'$ 仿射，则
$(Z'_k, Y'_k, i'_k, \pi'_k, \mathcal{G}'_k, \alpha'_k)_{k = 1, \ldots, n}$
是 $s'$ 上 $\mathcal{F}'/X'/S'$ 的完全 d\'evissage。
\end{lemma}

\begin{proof}
由引理 \ref{flat-lemma-base-change-one-step}，一步 d\'evissage 的基变换仍为
一步 d\'evissage。因此只需证明 $\Coker(\alpha_k)$ 的形成与基变换交换，
且定义 \ref{flat-definition-complete-devissage} 的条件 (2) 在基变换下保持。
第一个断言成立，因为 $\pi'_{k, *}\mathcal{G}'_k$ 是
$\pi_{k, *}\mathcal{G}_k$ 的拉回（由概形上同调章引理
\ref{coherent-lemma-affine-base-change}），且 $\otimes$ 右正合。
第二个断言成立，因为由同一理由有
$$
(\pi_{k, *}\mathcal{G}_k)_{\xi_k}
\otimes_{\mathcal{O}_{Y_k, \xi_k}} \kappa(\xi_k)
\otimes_{\kappa(\xi_k)} \kappa(\xi'_k)
\cong
(\pi'_{k, *}\mathcal{G}'_k)_{\xi'_k}
\otimes_{\mathcal{O}_{Y'_k, \xi'_k}} \kappa(\xi'_k)
$$
其中记号含义显然。
\end{proof}

\begin{lemma}
\label{flat-lemma-base-change-complete-at-x}
设 $S$、$X$、$\mathcal{F}$、$x$、$s$ 如定义
\ref{flat-definition-complete-devissage-at-x}。设带点概形态射
$(S', s') \to (S, s)$ 诱导同构 $\kappa(s) = \kappa(s')$。设
$(Z_k, Y_k, i_k, \pi_k, \mathcal{G}_k, \alpha_k, z_k, y_k)_{k = 1, \ldots, n}$
是 $x$ 处 $\mathcal{F}/X/S$ 的完全 d\'evissage。令
$(Z'_k, Y'_k, i'_k, \pi'_k, \mathcal{G}'_k, \alpha'_k)_{k = 1, \ldots, n}$
如引理 \ref{flat-lemma-base-change-complete} 中所构造；令 $x' \in X'$
（分别地，$z'_k \in Z'$、$y'_k \in Y'$）为同时映到
$x \in X$（分别地，$z_k \in Z_k$、$y_k \in Y_k$）与
$s' \in S'$ 的唯一一点。若 $S'$ 仿射，则
$(Z'_k, Y'_k, i'_k, \pi'_k, \mathcal{G}'_k, \alpha'_k,
z'_k, y'_k)_{k = 1, \ldots, n}$
是 $x'$ 处 $\mathcal{F}'/X'/S'$ 的完全 d\'evissage。
\end{lemma}

\begin{proof}
合并引理 \ref{flat-lemma-base-change-complete} 与
\ref{flat-lemma-base-change-one-step-at-x} 即得。
\end{proof}

\begin{definition}
\label{flat-definition-shrink-complete}
设 $S$、$X$、$\mathcal{F}$、$x$、$s$ 如定义
\ref{flat-definition-complete-devissage-at-x}。考虑 $x$ 处
$(Z_k, Y_k, i_k, \pi_k, \mathcal{G}_k, \alpha_k, z_k, y_k)_{k = 1, \ldots, n}$
的 $\mathcal{F}/X/S$ 完全 d\'evissage。此情形的一个
{\it 标准缩小}，是指给定标准开集 $S' \subset S$、$X' \subset X$、
$Z'_k \subset Z_k$、$Y'_k \subset Y_k$，使得 $s_k \in S'$、
$x_k \in X'$、$z_k \in Z'$、$y_k \in Y'$，并且
% P05-EN-CH38-0008: frozen definition uses undefined s_k,x_k and unindexed Z',Y' in its membership conditions.
$$
(Z'_k, Y'_k, i'_k, \pi'_k,
\mathcal{G}'_k, \alpha'_k, z_k, y_k)_{k = 1, \ldots, n}
$$
是 $x$ 处 $\mathcal{F}'/X'/S'$ 的一步 d\'evissage，其中
% P05-EN-CH38-0009: frozen definition calls the whole indexed system a one step rather than complete devissage.
$\mathcal{G}'_k = \mathcal{G}_k|_{Z'_k}$，且
$\mathcal{F}' = \mathcal{F}|_{X'}$.
\end{definition}

\begin{lemma}
\label{flat-lemma-shrink-complete}
假设与记号同定义 \ref{flat-definition-shrink-complete}。则：
\begin{enumerate}
\item
\label{flat-item-shrink-base-complete}
若 $S' \subset S$ 是 $s$ 的标准开邻域，则置
$X' = X_{S'}$、$Z'_k = Z_{S'}$、$Y'_k = Y_{S'}$，得到标准缩小。
% P05-EN-CH38-0010: frozen formula drops the k indices from Z_k and Y_k on the right.
\item
\label{flat-item-shrink-on-Y-complete}
设 $W \subset Y_n$ 是 $y$ 的标准开邻域。则存在满足
$Y'_n = W \times_S S'$ 的标准缩小。
\item
\label{flat-item-shrink-on-X-complete}
设 $U \subset X$ 是 $x$ 的开邻域。则存在满足 $X' \subset U$
的标准缩小。
\end{enumerate}
\end{lemma}

\begin{proof}
(1) 立即由引理 \ref{flat-lemma-base-change-complete-at-x} 与
\ref{flat-lemma-shrink} 得出。

\medskip\noindent
证明 (2)。为方便起见，记 $X = Y_0$。应用引理
\ref{flat-lemma-shrink} (\ref{flat-item-shrink-on-Y})，找到如下标准缩小：
$S', Y'_{n - 1}, Z'_n, Y'_n$
它是 $y_{n - 1}$ 处 $\Coker(\alpha_{n - 1})/Y_{n - 1}/S$
的一步 d\'evissage 的标准缩小，并满足 $Y'_n = W \times_S S'$。
重复此过程，找到标准缩小
$S'', Y''_{n - 2}, Z''_{n - 1}, Y''_{n - 1}$
它是 $y_{n - 2}$ 处 $\Coker(\alpha_{n - 2})/Y_{n - 2}/S$
的一步 d\'evissage 的标准缩小，并满足
$Y''_{n - 1} = Y'_{n - 1} \times_S S''$。如此继续，直至得到
$S^{(n)}, Y^{(n)}_0, Z^{(n)}_1, Y^{(n)}_1$.
此时显然，取 $S^{(n)}$、$X^{(n)}$、
$Z_k^{(n - k)} \times_S S^{(n)}$ 与
$Y_k^{(n - k)} \times_S S^{(n)}$，便得到所需性质的标准缩小。

\medskip\noindent
证明 (3)。对完全 d\'evissage 的长度作归纳。首先应用引理
\ref{flat-lemma-shrink} (\ref{flat-item-shrink-on-X})，找到标准缩小
$S', X', Z'_1, Y'_1$
它来自 $x$ 处 $\mathcal{F}/X/S$ 的一步 d\'evissage，且
$X' \subset U$。若 $n = 1$，结论已得。若 $n > 1$，则由归纳假设，
可找到完全 d\'evissage 的标准缩小 $S''$、$Y''_1$、$Z''_k$、$Y''_k$：
$(Z_k, Y_k, i_k, \pi_k, \mathcal{G}_k, \alpha_k, z_k, y_k)_{k = 2, \ldots, n}$
它是 $x$ 处 $\Coker(\alpha_1)/Y_1/S$ 的完全 d\'evissage，且
$Y''_1 \subset Y'_1$。利用引理 \ref{flat-lemma-shrink}
(\ref{flat-item-shrink-on-Y})，可找到 $S''' \subset S'$、
$X''' \subset X'$、$Z'''_1$ 与
$Y'''_1 = Y''_1 \times_S S'''$，它们构成标准缩小。问题的解取为
$$
S''', X''', Z'''_1, Y'''_1, Z''_2 \times_S S''',
Y''_2 \times_S S''', \ldots, Z''_n \times_S S''', Y''_n \times_S S'''
$$
引理证明完毕。
\end{proof}

\begin{proposition}
\label{flat-proposition-existence-complete-at-x}
设 $S$ 是概形，$X$ 在 $S$ 上局部有限型，$x \in X$ 的像为
$s \in S$。则存在交换图
$$
\xymatrix{
(X, x) \ar[d] & (X', x') \ar[l]^g \ar[d] \\
(S, s) & (S', s') \ar[l]
}
$$
其中水平箭头为初等 \'etale 邻域，且 $g^*\mathcal{F}/X'/S'$ 在 $x$
处有完全 d\'evissage。
% P05-EN-CH38-0011: frozen conclusion locates the pulled-back sheaf's devissage at x rather than x'.
\end{proposition}

\begin{proof}
对整数 $d = \dim_x(\text{Supp}(\mathcal{F}_s))$ 作归纳。由引理
\ref{flat-lemma-elementary-devissage-variant}，存在图表
$$
\xymatrix{
(X, x) \ar[d] & (X', x') \ar[l]^g \ar[d] \\
(S, s) & (S', s') \ar[l]
}
$$
其中水平箭头为初等 \'etale 邻域，且 $g^*\mathcal{F}/X'/S'$ 在
$x'$ 处有一步 d\'evissage。由问题的局部性，可用
$(X', x') \to (S', s')$ 替换 $(X, x) \to (S, s)$。替换后可假设
存在 $\mathcal{F}/X/S$ 在 $x$ 处的一步 d\'evissage
$(Z_1, Y_1, i_1, \pi_1, \mathcal{G}_1)$。

\medskip\noindent
应用引理 \ref{flat-lemma-existence-alpha}，找到同态
$$
\alpha_1 :
\mathcal{O}_{Y_1}^{\oplus r_1}
\longrightarrow
\pi_{1, *}\mathcal{G}_1
$$
它诱导 $\kappa(\xi_1)$ 上向量空间的同构，其中
$\xi_1 \in Y_1$ 是 $Y_1$ 在 $s$ 上纤维的唯一一般点。此外，
$\dim_{y_1}(\text{Supp}(\Coker(\alpha_1)_s)) < d$。
$\Coker(\alpha_1)_s$ 在 $y_1$ 处的茎可能为零。在此情形，由引理
\ref{flat-lemma-shrink} (\ref{flat-item-shrink-on-Y}) 缩小 $Y_1$ 后，
可假设 $\Coker(\alpha_1) = 0$，从而得到长度为零的完全 d\'evissage。

\medskip\noindent
现在假设 $\Coker(\alpha_1)_s$ 在 $y_1$ 处的茎非零。在此情形，
由归纳假设存在交换图
\begin{equation}
\label{flat-equation-overcome-this}
\vcenter{
\xymatrix{
(Y_1, y_1) \ar[d] & (Y'_1, y'_1) \ar[l]^h \ar[d] \\
(S, s) & (S', s') \ar[l]
}
}
\end{equation}
其中水平箭头为初等 \'etale 邻域，且
$h^*\Coker(\alpha_1)/Y'_1/S'$ 有完全 d\'evissage
$$
(Z_k, Y_k, i_k, \pi_k, \mathcal{G}_k, \alpha_k, z_k, y_k)_{k = 2, \ldots, n}
$$
于 $y'_1$。（特别地，$i_2 : Z_2 \to Y'_1$ 是到 $Y'_2$ 的闭浸入。）
% P05-EN-CH38-0012: frozen parenthetical gives i_2 target Y'_1, then calls it an immersion into Y'_2.
此时，把引理 \ref{flat-lemma-elementary-etale-neighbourhood} 应用于
$S, X, \mathcal{F}, x, s$、系统
$(Z_1, Y_1, i_1, \pi_1, \mathcal{G}_1)$ 与图
(\ref{flat-equation-overcome-this})，得到图表
$$
\xymatrix{
& & (X'', x'') \ar[lld] \ar[d] & (Z''_1, z''_1) \ar[l] \ar[lld] \ar[d] \\
(X, x) \ar[d] & (Z_1, z_1) \ar[l] \ar[d] &
(S'', s'') \ar[lld] & (Y''_1, y''_1) \ar[lld] \ar[l] \\
(S, s) & (Y_1, y_1) \ar[l]
}
$$
它具有所引引理列出的全部性质。特别地，
$Y''_1 \subset Y'_1 \times_{S'} S''$。置
$X_1 = Y'_1 \times_{S'} S''$，并记 $\mathcal{F}_1$ 为
$\Coker(\alpha_1)$ 的拉回。由引理
\ref{flat-lemma-base-change-complete-at-x}，系统
\begin{equation}
\label{flat-equation-shrink-this}
(Z_k \times_{S'} S'',
Y_k \times_{S'} S'', i''_k, \pi''_k, \mathcal{G}''_k,
\alpha''_k, z''_k, y''_k)_{k = 2, \ldots, n}
\end{equation}
是 $\mathcal{F}_1$ 到 $X_1$ 的完全 d\'evissage。再次由问题的性质，
可用 $(X'', x'') \to (S'', s'')$ 替换
$(X, x) \to (S, s)$。由此可假设：
\begin{enumerate}
\item[(a)] 存在 $\mathcal{F}/X/S$ 在 $x$ 处的一步 d\'evissage
$(Z_1, Y_1, i_1, \pi_1, \mathcal{G}_1)$；
\item[(b)] 存在 $\alpha_1 : \mathcal{O}_{Y_1}^{\oplus r_1}
\to \pi_{1, *}\mathcal{G}_1$，使得
$\alpha \otimes \kappa(\xi_1)$ 为同构；
% P05-EN-CH38-0013: frozen condition switches alpha_1 to alpha.
\item[(c)] $Y_1 \subset X_1$ 是开集，$y_1 = x_1$，且
$\mathcal{F}_1|_{Y_1} \cong \Coker(\alpha_1)$；
\item[(d)] 存在完全 d\'evissage
$(Z_k, Y_k, i_k, \pi_k, \mathcal{G}_k, \alpha_k, z_k, y_k)_{k = 2, \ldots, n}$
它是 $\mathcal{F}_1/X_1/S$ 在 $x_1$ 处的完全 d\'evissage。
\end{enumerate}
为完成证明，只需缩小这一步 d\'evissage 与完全 d\'evissage，使二者拼接成
一个完全 d\'evissage。（建议读者用引理 \ref{flat-lemma-shrink} 与
\ref{flat-lemma-shrink-complete} 自行完成，而不必阅读下述证明。）
由于 $Y_1 \subset X_1$ 是 $x_1$ 的开邻域，可应用引理
\ref{flat-lemma-shrink-complete} (\ref{flat-item-shrink-on-X-complete})，
找到数据 (d) 的标准缩小 $S', X'_1, Z'_2, Y'_2, \ldots, Y'_n$，
使 $X'_1 \subset Y_1$。注意，$X'_1$ 也是仿射概形 $Y_1$ 的标准开集。
接着如下缩小数据 (a)：先由引理 \ref{flat-lemma-shrink}
(\ref{flat-item-shrink-base}) 把底 $S$ 缩小到 $S'$，再用引理
\ref{flat-lemma-shrink} (\ref{flat-item-shrink-on-Y}) 把所得结果缩小为
$S''$、$X''$、$Z''_1$、$Y''_1$，最终满足
$Y''_1 = X'_1 \times_S S''$ 与 $S'' \subset S'$。于是
$$
Z''_1, Y''_1, Z'_2 \times_{S'} S'', Y'_2 \times_{S'} S'', \ldots,
Y'_n \times_{S'} S''
$$
给出所求的完全 d\'evissage。
\end{proof}

\noindent
再做一些记账工作，得到下列推论。

\begin{lemma}
\label{flat-lemma-existence-complete}
设 $X \to S$ 是有限型概形态射，$\mathcal{F}$ 是有限型拟凝聚
$\mathcal{O}_X$-模，且 $s \in S$ 为一点。则存在初等 \'etale 邻域
$(S', s') \to (S, s)$ 与 \'etale 态射
$h_i : Y_i \to X_{S'}$，$i = 1, \ldots, n$，使得对每个 $i$，
在 $s'$ 上存在 $\mathcal{F}_i/Y_i/S'$ 的完全 d\'evissage，其中
$\mathcal{F}_i$ 是 $\mathcal{F}$ 到 $Y_i$ 的拉回，并且
$X_s = (X_{S'})_{s'} \subset \bigcup h_i(Y_i)$。
\end{lemma}

\begin{proof}
对每一点 $x \in X_s$，可找到图表
$$
\xymatrix{
(X, x) \ar[d] & (X', x') \ar[l]^g \ar[d] \\
(S, s) & (S', s') \ar[l]
}
$$
其中水平箭头为初等 \'etale 邻域，且 $g^*\mathcal{F}/X'/S'$ 在
$x'$ 处有完全 d\'evissage。由于 $X \to S$ 有限型，纤维 $X_s$
拟紧；又由于上述每个 $g : X' \to X$ 都是开映射，可用有限多个
$g(X'_{s'})$ 覆盖 $X_s$。因此可找到这类图表的有限族
$$
\vcenter{
\xymatrix{
(X, x) \ar[d] & (X'_i, x'_i) \ar[l]^{g_i} \ar[d] \\
(S, s) & (S'_i, s'_i) \ar[l]
}
}
\quad i = 1, \ldots, n
$$
使得 $X_s = \bigcup g_i(X'_i)$。置
$S' = S'_1 \times_S \ldots \times_S S'_n$，并令
$Y_i = X_i \times_{S'_i} S'$ 为 $X'_i$ 到 $S'$ 的基变换。
由引理 \ref{flat-lemma-base-change-complete}，$\mathcal{F}$ 到 $Y_i$
的拉回在 $s$ 上有完全 d\'evissage，结论得证。
% P05-EN-CH38-0014: frozen proof defines Y_i from X_i although only X'_i was introduced, and concludes over s rather than s'.
\end{proof}



\section{代数表述}
\label{flat-section-translation}

\noindent
把具有完全 d\'evissage 的含义用代数语言写清楚或许有用。为此引入下列概念
（它并不十分有用，所以我们给它一个长得不可思议的名称）。

\begin{definition}
\label{flat-definition-elementary-etale-neighbourhood}
设 $R \to S$ 是环同态，$\mathfrak q$ 是 $S$ 中位于 $R$ 的素理想
$\mathfrak p$ 上方的素理想。所谓{\it 环同态 $R \to S$ 在
$\mathfrak q$ 处的初等 \'etale 局部化}，是指给定环及相伴素理想的交换图
$$
\xymatrix{
S \ar[r] & S' \\
R \ar[u] \ar[r] & R' \ar[u]
}
\quad\quad
\xymatrix{
\mathfrak q \ar@{-}[r] & \mathfrak q' \\
\mathfrak p \ar@{-}[u] \ar@{-}[r] & \mathfrak p' \ar@{-}[u]
}
$$
使得 $R \to R'$ 与 $S \to S'$ 是 \'etale 环同态，且
$\kappa(\mathfrak p) = \kappa(\mathfrak p')$ 与
$\kappa(\mathfrak q) = \kappa(\mathfrak q')$。
\end{definition}

\begin{definition}
\label{flat-definition-complete-devissage-algebra}
设 $R \to S$ 是有限型环同态，$\mathfrak r$ 是 $R$ 的素理想，
$N$ 是有限 $S$-模。所谓 {\it $\mathfrak r$ 上 $N/S/R$ 的完全
d\'evissage}，是指给定 $R$-代数同态
$$
\xymatrix{
& A_1 & & A_2 & & ... & & A_n \\
S \ar[ru] & & B_1 \ar[lu] \ar[ru] & & ... \ar[lu] \ar[ru] & &
... \ar[lu] \ar[ru] & & B_n \ar[lu]
}
$$
有限 $A_i$-模 $M_i$ 与 $B_i$-模同态
$\alpha_i : B_i^{\oplus r_i} \to M_i$，满足：
\begin{enumerate}
\item $S \to A_1$ 满且有限表示；
\item $B_i \to A_{i + 1}$ 满且有限表示；
\item $B_i \to A_i$ 有限；
\item $R \to B_i$ 光滑且纤维几何不可约；
\item 作为 $S$-模有 $N \cong M_1$；
\item 作为 $B_i$-模有 $\Coker(\alpha_i) \cong M_{i + 1}$；
\item 同态 $\alpha_i : \kappa(\mathfrak p_i)^{\oplus r_i}
\to M_i \otimes_{B_i} \kappa(\mathfrak p_i)$ 为同构，
其中 $\mathfrak p_i = \mathfrak rB_i$；
\item $\Coker(\alpha_n) = 0$.
\end{enumerate}
% P05-EN-CH38-0015: frozen conditions (2) and (6) lack an i<n range, and condition (7) treats rB_i as a prime without qualification.
在此情形，称
$(A_i, B_i, M_i, \alpha_i)_{i = 1, \ldots, n}$
为 $\mathfrak r$ 上 $N/S/R$ 的完全 d\'evissage。
\end{definition}

\begin{remark}
\label{flat-remark-finite-presentation}
注意，对所有 $i$，$R$-代数 $B_i$ 在 $R$ 上有限表示；当 $i \geq 2$
时，$A_i$ 亦然。若 $S$ 在 $R$ 上有限表示，则 $A_1$ 也在 $R$
上有限表示。在此情形，完全 d\'evissage 中的全部环同态均有限表示。
见《代数》引理 \ref{algebra-lemma-compose-finite-type}。仍假设 $S$
在 $R$ 上有限表示，则下列条件等价：
\begin{enumerate}
\item $M$ 在 $S$ 上有限表示；
\item $M_1$ 在 $A_1$ 上有限表示；
\item $M_1$ 在 $B_1$ 上有限表示；
\item 每个 $M_i$ 作为 $A_i$-模与作为 $B_i$-模均有限表示。
\end{enumerate}
% P05-EN-CH38-0016: frozen remark switches the finite module from N to undefined M.
等价关系 (1) $\Leftrightarrow$ (2) 与 (2) $\Leftrightarrow$ (3)
由《代数》引理 \ref{algebra-lemma-finite-finitely-presented-extension} 得出。
若 $M_1$ 有限表示，则 $\Coker(\alpha_1)$ 也有限表示（见《代数》引理
\ref{algebra-lemma-extension}），因而 $M_2$ 有限表示，依此类推。
\end{remark}

\begin{definition}
\label{flat-definition-complete-devissage-at-x-algebra}
设 $R \to S$ 是有限型环同态，$\mathfrak q$ 是 $S$ 中位于 $R$
的素理想 $\mathfrak r$ 上方的素理想，且 $N$ 是有限 $S$-模。
所谓 {\it $N/S/R$ 在 $\mathfrak q$ 处的完全 d\'evissage}，是指给定
$\mathfrak r$ 上 $N/S/R$ 的完全 d\'evissage
$(A_i, B_i, M_i, \alpha_i)_{i = 1, \ldots, n}$ 与位于
$\mathfrak r$ 上方的素理想 $\mathfrak q_i \subset B_i$，使得
\begin{enumerate}
\item $\kappa(\mathfrak r) \subset \kappa(\mathfrak q_i)$ 是纯超越扩张；
\item 存在唯一的素理想 $\mathfrak q'_i \subset A_i$ 位于
$\mathfrak q_i \subset B_i$ 上方；
\item $\mathfrak q = \mathfrak q'_1 \cap S$，且
$\mathfrak q_i = \mathfrak q'_{i + 1} \cap A_i$,
\item $R \to B_i$ 的相对维数为
$\dim_{\mathfrak q_i}(\text{Supp}(M_i \otimes_R \kappa(\mathfrak r)))$.
\end{enumerate}
\end{definition}

\begin{remark}
\label{flat-remark-same-notion}
设 $A \to B$ 是有限型环同态，$N$ 是有限 $B$-模，$\mathfrak q$
是 $B$ 中位于 $A$ 的素理想 $\mathfrak r$ 上方的素理想。置
$X = \Spec(B)$、$S = \Spec(A)$，并在 $X$ 上置
$\mathcal{F} = \widetilde{N}$。设 $x$ 为与 $\mathfrak q$ 对应的点，
$s \in S$ 为与 $\mathfrak p$ 对应的点。则
% P05-EN-CH38-0017: frozen remark introduces r but later calls the corresponding base prime p.
\begin{enumerate}
\item 若存在 $s$ 上 $\mathcal{F}/X/S$ 的完全 d\'evissage，则存在
$\mathfrak p$ 上 $N/B/A$ 的完全 d\'evissage；
\item 存在 $\mathcal{F}/X/S$ 在 $x$ 处的完全 d\'evissage，当且仅当存在
$N/B/A$ 在 $\mathfrak q$ 处的完全 d\'evissage。
\end{enumerate}
这里有一点细微差别：在“$\mathfrak p$ 上 $N/B/A$ 的完全
d\'evissage”的表述中，我们省略了相对维数条件，所以 (1) 中的蕴含
只有一个方向。在 $\mathfrak q$ 处完全 d\'evissage 的概念中已经包含
此条件。无论如何，我们只会使用如下事实：$\mathcal{F}/X/S$
存在完全 d\'evissage 蕴含 $N/B/A$ 存在完全 d\'evissage。
\end{remark}

\begin{lemma}
\label{flat-lemma-existence-algebra}
设 $R \to S$ 是有限型环同态，$M$ 是有限 $S$-模，
$\mathfrak q$ 是 $S$ 的素理想。则存在环同态 $R \to S$ 在
$\mathfrak q$ 处的初等 \'etale 局部化
$R' \to S', \mathfrak q', \mathfrak p'$，使得
$(M \otimes_S S')/S'/R'$ 在 $\mathfrak q'$ 处存在完全 d\'evissage。
\end{lemma}

\begin{proof}
这是命题 \ref{flat-proposition-existence-complete-at-x} 经由注记
\ref{flat-remark-same-notion} 得到的重述。
\end{proof}




\section{局部化与普遍单射同态}
\label{flat-section-localize-universally-injective}


\begin{lemma}
\label{flat-lemma-homothety-spectrum}
设 $R \to S$ 是环同态，$N$ 是 $S$-模。假设
\begin{enumerate}
\item $R$ 是极大理想为 $\mathfrak m$ 的局部环；
\item $\overline{S} = S/\mathfrak m S$ 是诺特环；并且
% P05-EN-CH38-0018: frozen source writes m_R although only the maximal ideal m is defined.
\item $\overline{N} = N/\mathfrak m_R N$ 是有限 $\overline{S}$-模。
\end{enumerate}
令 $\Sigma \subset S$ 为由在 $\overline{N}$ 上不是零因子的元素组成的乘法子集。
则 $\Sigma^{-1}S$ 是半局部环，其谱由包含于
$\text{Ass}_S(\overline{N})$ 某一元素的素理想 $\mathfrak q \subset S$ 组成。
此外，$\Sigma^{-1}S$ 的任一极大理想都对应于 $\overline{N}$
作为 $\overline{S}$-模的一个伴随素理想。
\end{lemma}

\begin{proof}
注意
$\text{Ass}_S(\overline{N}) = \text{Ass}_{\overline{S}}(\overline{N})$，参见
代数，引理 \ref{algebra-lemma-ass-quotient-ring}。
由代数，引理 \ref{algebra-lemma-finite-ass}，这是有限集。
记 $\{\mathfrak q_1, \ldots, \mathfrak q_r\} = \text{Ass}_S(\overline{N})$。
由代数，引理 \ref{algebra-lemma-ass-zero-divisors}，有
$\Sigma = S \setminus (\bigcup \mathfrak q_i)$。
结合代数，引理 \ref{algebra-lemma-spec-localization} 中对
$\Spec(\Sigma^{-1}S)$ 的描述和代数，引理 \ref{algebra-lemma-silly}，
可知 $\Sigma^{-1}S$ 的素理想对应于 $S$ 中包含于某个
$\mathfrak q_i$ 的素理想。因此，$\Sigma^{-1}S$ 的极大理想与集合
$\{\mathfrak q_1, \ldots, \mathfrak q_r\}$ 中关于包含关系的极大元一一对应。
引理得证。
\end{proof}

\begin{lemma}
\label{flat-lemma-homothety-universally-injective}
采用引理 \ref{flat-lemma-homothety-spectrum} 的假设与记号。再假设
\begin{enumerate}
\item $S$ 是局部环，且 $R \to S$ 是局部同态；
\item $S$ 在 $R$ 上本质有限表示；
\item $N$ 是有限表示 $S$-模；并且
\item $N$ 在 $R$ 上平坦。
\end{enumerate}
则每个 $s \in \Sigma$ 都定义一个普遍单射的 $R$-模同态
$s : N \to N$，并且同态 $N \to \Sigma^{-1}N$ 是 $R$-普遍单射的。
\end{lemma}

\begin{proof}
由代数，引理 \ref{algebra-lemma-mod-injective-general}，序列
$0 \to N \to N \to N/sN \to 0$ 正合，且 $N/sN$ 在 $R$ 上平坦。
这说明 $s : N \to N$ 普遍单射；参见代数，引理
\ref{algebra-lemma-flat-tor-zero}。同态 $N \to \Sigma^{-1}N$
是同态 $s : N \to N$ 的有向余极限，因而普遍单射。
\end{proof}

\begin{lemma}
\label{flat-lemma-base-change-universally-flat-local}
设 $R \to S$ 是环同态，$N$ 是 $S$-模，$S \to S'$ 是环同态。假设
\begin{enumerate}
\item $R \to S$ 是局部环之间的局部同态；
\item $S$ 在 $R$ 上本质有限表示；
\item $N$ 是有限表示 $S$-模；
\item $N$ 在 $R$ 上平坦；
\item $S \to S'$ 平坦；并且
\item $\Spec(S') \to \Spec(S)$ 的像包含 $S$ 中位于
$\mathfrak m_R$ 上方的所有素理想 $\mathfrak q$，其中
$\mathfrak q$ 是 $N/\mathfrak m_R N$ 的伴随素理想。
\end{enumerate}
则 $N \to N \otimes_S S'$ 是 $R$-普遍单射的。
\end{lemma}

\begin{proof}
% P05-EN-CH38-0019: frozen source defines N' using tensor over R although S' is introduced as an S-algebra and the proof later uses tensor over S.
置 $N' = N \otimes_R S'$。考虑交换图
$$
\xymatrix{
N \ar[d] \ar[r] & N' \ar[d] \\
\Sigma^{-1}N \ar[r] & \Sigma^{-1}N'
}
$$
其中 $\Sigma \subset S$ 是由在 $N/\mathfrak m_R N$ 上不是零因子的元素组成的集合。
若能证明同态 $N \to \Sigma^{-1}N'$ 普遍单射，则 $N \to N'$ 也普遍单射
（参见代数，引理 \ref{algebra-lemma-universally-injective-permanence}）。

\medskip\noindent
由引理 \ref{flat-lemma-homothety-spectrum}，环 $\Sigma^{-1}S$ 是半局部环，
其极大理想对应于 $N/\mathfrak m_R N$ 的伴随素理想。因此，由假设，
$\Spec(\Sigma^{-1}S') \to \Spec(\Sigma^{-1}S)$ 的像包含所有这些极大理想。
由代数，引理 \ref{algebra-lemma-ff-rings}，环同态
$\Sigma^{-1}S \to \Sigma^{-1}S'$ 忠实平坦。因而
$\Sigma^{-1}N \to \Sigma^{-1}N'$，即同态
$$
N \otimes_S \Sigma^{-1}S \longrightarrow N \otimes_S \Sigma^{-1}S'
$$
普遍单射；参见代数，引理
\ref{algebra-lemma-faithfully-flat-universally-injective} 和
\ref{algebra-lemma-universally-injective-tensor}。最后应用引理
\ref{flat-lemma-homothety-universally-injective}，可知
$N \to \Sigma^{-1}N$ 普遍单射。普遍单射模同态的复合仍普遍单射
（参见代数，引理 \ref{algebra-lemma-composition-universally-injective}），
故 $N \to \Sigma^{-1}N'$ 普遍单射，得证。
\end{proof}

\begin{lemma}
\label{flat-lemma-base-change-universally-flat}
设 $R \to S$ 是环同态，$N$ 是 $S$-模，$S \to S'$ 是环同态。假设
\begin{enumerate}
\item $R \to S$ 有限表示，且 $N$ 是有限表示 $S$-模；
\item $N$ 在 $R$ 上平坦；
\item $S \to S'$ 平坦；并且
\item $\Spec(S') \to \Spec(S)$ 的像包含所有满足下述条件的素理想
$\mathfrak q$：若 $\mathfrak p$ 是 $\mathfrak q$ 在 $R$ 中的逆像，则
$\mathfrak q$ 是 $N \otimes_R \kappa(\mathfrak p)$ 的伴随素理想。
\end{enumerate}
则 $N \to N \otimes_S S'$ 是 $R$-普遍单射的。
\end{lemma}

\begin{proof}
由代数，引理 \ref{algebra-lemma-universally-injective-check-stalks}，只须证明：
对于 $S$ 中位于 $R$ 的 $\mathfrak p$ 上方的任一素理想 $\mathfrak q$，
% P05-EN-CH38-0020: frozen source tensors over R and localizes S' at q although the map is S to S' and no prime above q is specified.
$N_{\mathfrak q} \to (N \otimes_R S')_{\mathfrak q}$ 是
$R_{\mathfrak p}$-普遍单射的。因此可将引理
\ref{flat-lemma-base-change-universally-flat-local} 应用于环同态
$R_{\mathfrak p} \to S_{\mathfrak q} \to S'_{\mathfrak q}$
以及模 $N_{\mathfrak q}$。
\end{proof}

\noindent
读者不妨将下述引理与代数，引理 \ref{algebra-lemma-mod-injective}、
\ref{algebra-lemma-mod-injective-general} 以及
第 \ref{flat-section-variants-mod-injective} 节的结果比较。
在每种情形中，结论都是同态 $u : M \to N$ 普遍单射且余核平坦。

\begin{lemma}
\label{flat-lemma-universally-injective-local}
设 $(R, \mathfrak m)$ 是局部环，$u : M \to N$ 是 $R$-模同态。
若 $M$ 是投射 $R$-模，$N$ 是平坦 $R$-模，并且
$\overline{u} : M/\mathfrak mM \to N/\mathfrak mN$ 单射，
则 $u$ 普遍单射。
\end{lemma}

\begin{proof}
由代数，定理 \ref{algebra-theorem-projective-free-over-local-ring}，模 $M$ 自由。
只要对 $M$ 的每个有限生成直和因子证明结论，引理便随之成立。因此可设
$M$ 有限自由。将 $N = \colim_i N_i$ 写成有限自由模的有向余极限；
参见代数，定理 \ref{algebra-theorem-lazard}。由于 $M$ 有限自由，
同态 $u : M \to N$ 对某个 $i$ 通过 $N_i$ 分解。将相应的 $R$-模同态记为
$u_i : M \to N_i$。因为 $\overline{u}$ 单射，所以
$\overline{u_i} : M/\mathfrak mM \to N_i/\mathfrak mN_i$ 单射，
并且对每个 $i' \geq i$，它与同态
$N_i/\mathfrak mN_i \to N_{i'}/\mathfrak mN_{i'}$ 的复合仍单射。
由于 $M$ 和 $N_{i'}$ 都是局部环 $R$ 上的有限自由模，这说明对每个
$i' \geq i$，$M \to N_{i'}$ 都是分裂单射。因此，对任一 $R$-模 $Q$，
以及所有 $i' \geq i$，同态 $M \otimes_R Q \to N_{i'} \otimes_R Q$
均单射。又因为 $- \otimes_R Q$ 与余极限交换，故可得
% P05-EN-CH38-0021: frozen conclusion retains N_{i'} after passing to the colimit; the target should be N.
$M \otimes_R Q \to N_{i'} \otimes_R Q$ 单射，如所要求。
\end{proof}

\begin{lemma}
\label{flat-lemma-invert-universally-injective}
采用引理 \ref{flat-lemma-homothety-spectrum} 的假设与记号。
再假设 $N$ 作为 $R$-模是投射的。则每个 $s \in \Sigma$ 都定义一个
普遍单射的 $R$-模同态 $s : N \to N$，并且同态
$N \to \Sigma^{-1}N$ 是 $R$-普遍单射的。
\end{lemma}

\begin{proof}
取 $s \in \Sigma$。由引理 \ref{flat-lemma-universally-injective-local}，
同态 $s : N \to N$ 普遍单射。同态 $N \to \Sigma^{-1}N$
是同态 $s : N \to N$ 的有向余极限，因而普遍单射。
\end{proof}






\section{完备化与 Mittag-Leffler 模}
\label{flat-section-completion-ML}


\begin{lemma}
\label{flat-lemma-completed-direct-sum-ML}
设 $R$ 是环，$I \subset R$ 是理想，$A$ 是集合。
假设 $R$ 是诺特环且关于 $I$ 完备。则完备化
$(\bigoplus\nolimits_{\alpha \in A} R)^\wedge$ 平坦且为 Mittag-Leffler 模。
\end{lemma}

\begin{proof}
由《代数进阶》，引理
\ref{more-algebra-lemma-ui-completion-direct-sum-into-product}，同态
$(\bigoplus\nolimits_{\alpha \in A} R)^\wedge
\to \prod_{\alpha \in A} R$ 普遍单射。因此，由代数，引理
\ref{algebra-lemma-ui-flat-domain} 和 \ref{algebra-lemma-pure-submodule-ML}，
只须证明 $\prod_{\alpha \in A} R$ 平坦且为 Mittag-Leffler 模。
由代数，命题 \ref{algebra-proposition-characterize-coherent}
（以及代数，引理 \ref{algebra-lemma-Noetherian-coherent}），可知
$\prod_{\alpha \in A} R$ 平坦。又因为 $R$ 的诸副本之积是
Mittag-Leffler 模，结论成立；参见代数，引理
\ref{algebra-lemma-product-over-Noetherian-ring}。
\end{proof}

\begin{lemma}
\label{flat-lemma-lift-ML}
设 $R$ 是环，$I \subset R$ 是理想，$M$ 是 $R$-模。假设
\begin{enumerate}
\item $R$ 是诺特环且 $I$-进完备；
\item $M$ 在 $R$ 上平坦；并且
\item $M/IM$ 是投射 $R/I$-模。
\end{enumerate}
则 $I$-进完备化 $M^\wedge$ 是平坦 Mittag-Leffler $R$-模。
\end{lemma}

\begin{proof}
选择满射 $F \to M$，其中 $F$ 是自由 $R$-模。由代数，引理
\ref{algebra-lemma-split-completed-sequence}，模 $M^\wedge$ 是模
$F^\wedge$ 的直和因子。因此只须对 $F$ 证明引理；此时结论由引理
\ref{flat-lemma-completed-direct-sum-ML} 得出。
\end{proof}

\noindent
在引理 \ref{flat-lemma-universally-injective-to-completion} 和
\ref{flat-lemma-universally-injective-to-completion-flat} 中，若 $R \to S$
有限型，则 $S$ 为诺特环这一假设成立；参见代数，引理
\ref{algebra-lemma-Noetherian-permanence}。

\begin{lemma}
\label{flat-lemma-universally-injective-to-completion}
设 $R$ 是环，$I \subset R$ 是理想，$R \to S$ 是环同态，
$N$ 是 $S$-模。假设
\begin{enumerate}
\item $R$ 是诺特环；
\item $S$ 是诺特环；
\item $N$ 是有限 $S$-模；并且
\item 对任一有限 $R$-模 $Q$，每个
$\mathfrak q \in \text{Ass}_S(Q \otimes_R N)$ 均满足
$IS + \mathfrak q \not = S$。
\end{enumerate}
则从 $N$ 到 $N$ 的 $I$-进完备化的同态 $N \to N^\wedge$，
作为 $R$-模同态是普遍单射的。
\end{lemma}

\begin{proof}
须证明对任一有限 $R$-模 $Q$，同态
$Q \otimes_R N \to Q \otimes_R N^\wedge$ 单射；参见代数，定理
\ref{algebra-theorem-universally-exact-criteria}。由于存在典范同态
$Q \otimes_R N^\wedge \to (Q \otimes_R N)^\wedge$，只须证明典范同态
$Q \otimes_R N \to (Q \otimes_R N)^\wedge$ 单射。因此可用
$Q \otimes_R N$ 替换 $N$，并且只须证明同态 $N \to N^\wedge$ 单射。

\medskip\noindent
令 $K = \Ker(N \to N^\wedge)$。由于 $N$ 是
$\prod_{\mathfrak q \in \text{Ass}(N)} N_{\mathfrak q}$ 的子模，
只须对每个 $\mathfrak q \in \text{Ass}(N)$ 证明
$K_{\mathfrak q} = 0$；参见代数，引理
\ref{algebra-lemma-zero-at-ass-zero}。取
$\mathfrak q \in \text{Ass}(N)$。由最后一个假设，存在素理想
$\mathfrak q' \supset IS + \mathfrak q$。由于 $K_{\mathfrak q}$
是 $K_{\mathfrak q'}$ 的局部化，只须证明 $K_{\mathfrak q'}$ 为零。
注意 $K = \bigcap I^nN$，故
$K_{\mathfrak q'} \subset \bigcap I^nN_{\mathfrak q'}$。
由代数，引理 \ref{algebra-lemma-intersect-powers-ideal-module-zero}，
$K_{\mathfrak q'} = 0$。
\end{proof}

\begin{lemma}
\label{flat-lemma-universally-injective-to-completion-flat}
设 $R$ 是环，$I \subset R$ 是理想，$R \to S$ 是环同态，
$N$ 是 $S$-模。假设
\begin{enumerate}
\item $R$ 是诺特环；
\item $S$ 是诺特环；
\item $N$ 是有限 $S$-模；
\item $N$ 在 $R$ 上平坦；并且
\item 对任一素理想 $\mathfrak q \subset S$，若它是
$N \otimes_R \kappa(\mathfrak p)$ 的伴随素理想，其中
$\mathfrak p = R \cap \mathfrak q$，则有 $IS + \mathfrak q \not = S$。
\end{enumerate}
则从 $N$ 到 $N$ 的 $I$-进完备化的同态 $N \to N^\wedge$，
作为 $R$-模同态是普遍单射的。
\end{lemma}

\begin{proof}
结论由引理 \ref{flat-lemma-universally-injective-to-completion} 得出，
因为代数，引理 \ref{algebra-lemma-bourbaki-fibres} 和注记
\ref{algebra-remark-bourbaki} 保证张量积 $N \otimes_R Q$
的伴随素理想集合包含于诸模
$N \otimes_R \kappa(\mathfrak p)$ 的伴随素理想集合之中。
\end{proof}




\section{投射模}
\label{flat-section-projective}

\noindent
下述引理可用于在基底上作诺特归纳以证明投射性；参见引理
\ref{flat-lemma-fibres-irreducible-flat-projective}。

\begin{lemma}
\label{flat-lemma-flat-pure-over-complete-projective}
设 $R$ 是环，$I \subset R$ 是理想，$R \to S$ 是环同态，
$N$ 是 $S$-模。假设
\begin{enumerate}
\item $R$ 是诺特环且 $I$-进完备；
\item $R \to S$ 有限型；
\item $N$ 是有限 $S$-模；
\item $N$ 在 $R$ 上平坦；
\item $N/IN$ 作为 $R/I$-模是投射的；并且
\item 对任一素理想 $\mathfrak q \subset S$，若它是
$N \otimes_R \kappa(\mathfrak p)$ 的伴随素理想，其中
$\mathfrak p = R \cap \mathfrak q$，则有 $IS + \mathfrak q \not = S$。
\end{enumerate}
则 $N$ 作为 $R$-模是投射的。
\end{lemma}

\begin{proof}
由引理 \ref{flat-lemma-universally-injective-to-completion-flat}，同态
$N \to N^\wedge$ 普遍单射。由引理 \ref{flat-lemma-lift-ML}，模
$N^\wedge$ 是 Mittag-Leffler 模。再由代数，引理
\ref{algebra-lemma-pure-submodule-ML}，$N$ 是 Mittag-Leffler 模。
因此 $N$ 作为 $R$-模可数生成、平坦且为 Mittag-Leffler 模，
故由代数，引理 \ref{algebra-lemma-countgen-projective}，它是投射模。
\end{proof}

\begin{lemma}
\label{flat-lemma-fibres-irreducible-flat-projective}
设 $R$ 是环，$R \to S$ 是环同态。假设
\begin{enumerate}
\item $R$ 是诺特环；
\item $R \to S$ 有限型且平坦；并且
\item 每个纤维环 $S \otimes_R \kappa(\mathfrak p)$ 在
$\kappa(\mathfrak p)$ 上几何整。
\end{enumerate}
则 $S$ 作为 $R$-模是投射的。
\end{lemma}

\begin{proof}
考虑集合
$$
\{I \subset R \mid S/IS\text{ 作为 }R/I\text{-模不投射}\}
$$
须证明此集合为空。反设它非空，则它含有极大元 $I$。
令 $J = \sqrt{I}$ 为其根理想。若 $I \not = J$，则 $S/JS$
作为 $R/J$-模是投射的，$S/IS$ 在 $R/I$ 上平坦，并且 $J/I$
是 $R/I$ 中的幂零理想。应用代数，引理
\ref{algebra-lemma-lift-projective}，可知 $S/IS$ 是投射
$R/I$-模，矛盾。因此可设 $I$ 是根理想。换言之，问题归结为在下述情形
证明引理：$R$ 是既约环，并且对 $R$ 的每个非零理想 $I$，
$S/IS$ 都是投射 $R/I$-模。

\medskip\noindent
假设 $R$ 是既约环，并且对 $R$ 的每个非零理想 $I$，$S/IS$
都是投射 $R/I$-模。由一般平坦性，即代数，引理
\ref{algebra-lemma-generic-flatness-Noetherian}
（应用于为整环的局部化 $R_g$），或更一般的代数，引理
\ref{algebra-lemma-generic-flatness-reduced}，存在非零元 $f \in R$，
使得 $S_f$ 是自由 $R_f$-模。将 $R$ 的 $(f)$-进完备化记为
$R^\wedge = \lim R/(f^n)$。注意环同态
$$
R \longrightarrow R_f \times R^\wedge
$$
忠实平坦；参见代数，引理 \ref{algebra-lemma-completion-flat}。
因此，由投射性的忠实平坦下降（参见代数，定理
\ref{algebra-theorem-ffdescent-projectivity}），只须证明
$S \otimes_R R^\wedge$ 是投射 $R^\wedge$-模。为此将使用引理
\ref{flat-lemma-flat-pure-over-complete-projective} 的判据。首先注意，
$S/fS = (S \otimes_R R^\wedge)/f(S \otimes_R R^\wedge)$
是投射 $R/(f)$-模，并且 $S \otimes_R R^\wedge$ 作为相应性质的基变换，
在 $R^\wedge$ 上平坦且有限型。其次，设 $\mathfrak p^\wedge$ 是
$R^\wedge$ 的素理想，令 $\mathfrak p \subset R$ 为 $R$ 中相应的素理想。
由于 $R \to S$ 的纤维环几何整，任一基变换的纤维环也如此。因此
% P05-EN-CH38-0022: frozen source has the number-agreement typo “is a prime ideals lying”.
$\mathfrak q^\wedge = \mathfrak p^\wedge(S \otimes_R R^\wedge)$
是位于 $\mathfrak p^\wedge$ 上方的素理想，并且是
$S \otimes_R \kappa(\mathfrak p^\wedge)$ 的唯一伴随素理想。因此，只须有
$f(S \otimes_R R^\wedge) + \mathfrak q^\wedge \not = S \otimes_R R^\wedge$
便得证。这确实成立：$f$ 位于 $f$-进完备环 $R^\wedge$ 的雅可比根中，
所以 $\mathfrak p^\wedge + fR^\wedge \not = R^\wedge$；又因为各纤维非空
（不可约空间非空），$R^\wedge \to S \otimes_R R^\wedge$ 在谱上满射。
\end{proof}

\begin{lemma}
\label{flat-lemma-fibres-irreducible-flat-projective-nonnoetherian}
设 $R$ 是环，$R \to S$ 是环同态。假设
\begin{enumerate}
\item $R \to S$ 有限表示且平坦；并且
\item 每个纤维环 $S \otimes_R \kappa(\mathfrak p)$ 在
$\kappa(\mathfrak p)$ 上几何整。
\end{enumerate}
则 $S$ 作为 $R$-模是投射的。
\end{lemma}

\begin{proof}
可以找到环的余笛卡儿图
$$
\xymatrix{
S_0 \ar[r] & S \\
R_0 \ar[u] \ar[r] & R \ar[u]
}
$$
其中 $R_0$ 在 $\mathbf{Z}$ 上有限型，且同态 $R_0 \to S_0$
有限型、平坦并具有几何整纤维；参见《态射进阶》，引理
\ref{more-morphisms-lemma-Noetherian-approximation-flat}、
\ref{more-morphisms-lemma-Noetherian-approximation-geometrically-reduced},
\ref{more-morphisms-lemma-Noetherian-approximation-geometrically-irreducible},
以及 \ref{more-morphisms-lemma-Noetherian-approximation-combine}。
由引理 \ref{flat-lemma-fibres-irreducible-flat-projective}，$S_0$ 是投射
$R_0$-模。因此 $S = S_0 \otimes_{R_0} R$ 是投射 $R$-模；
参见代数，引理 \ref{algebra-lemma-ascend-properties-modules}。
\end{proof}

\begin{remark}
\label{flat-remark-how-in-RG}
引理 \ref{flat-lemma-fibres-irreducible-flat-projective-nonnoetherian}
是本章后续结果发展中的关键一步。该引理在 \cite{GruRay} 中的对应结果是
\cite[I 命题 3.3.1]{GruRay}：若 $R \to S$ 光滑且纤维几何整，
则 $S$ 作为 $R$-模是投射的。这是引理
\ref{flat-lemma-fibres-irreducible-flat-projective-nonnoetherian} 的特殊情形；
不过我们稍后终究会改进该引理，所以此时拥有一个更强结果并无太大收益。
下面简述 \cite{GruRay} 中给出的证明。
\begin{enumerate}
\item 先如上归结到 $R$ 为诺特环的情形。
\item 投射性沿忠实平坦环同态下降；参见代数，定理
\ref{algebra-theorem-ffdescent-projectivity}。因此可在 $R$ 的 fppf
拓扑下局部工作，从而可设 $R \to S$ 有截面 $\sigma : S \to R$
（只须使用基变换到 $S$ 的通常技巧）。置 $I = \Ker(S \to R)$。
\item 在 $R$ 上再作一些局部化，可设 $I/I^2$ 是自由 $R$-模，且 $S$
关于 $I$ 的完备化 $S^\wedge$ 同构于 $R[[t_1, \ldots, t_n]]$；
参见《态射》，引理 \ref{morphisms-lemma-section-smooth-morphism}。
这里使用了 $R \to S$ 光滑这一事实。
\item 为证明 $S$ 作为 $R$-模是投射的，只须证明 $S$ 作为 $R$-模
平坦、可数生成且为 Mittag-Leffler 模；参见代数，引理
\ref{algebra-lemma-countgen-projective}。前两个性质显然，故只须证明
$S$ 作为 $R$-模是 Mittag-Leffler 模。由代数，引理
\ref{algebra-lemma-power-series-ML}，模 $R[[t_1, \ldots, t_n]]$
在 $R$ 上为 Mittag-Leffler 模。因此，代数，引理
\ref{algebra-lemma-pure-submodule-ML} 表明，只须证明
% P05-EN-CH38-0023: frozen source has the grammatical surplus article in “show that the S to S^wedge”.
$S \to S^\wedge$ 作为 $R$-模同态普遍单射。
\item 应用引理 \ref{flat-lemma-base-change-universally-flat} 可知
$S \to S^\wedge$ 是 $R$-普遍单射的。事实上，由于 $R \to S$
具有几何整纤维，任一纤维环的每个伴随点都只是该纤维环的一般点，
而这个一般点属于 $\Spec(S^\wedge) \to \Spec(S)$ 的像。
\end{enumerate}
刚才概述的证明与导向引理
\ref{flat-lemma-fibres-irreducible-flat-projective-nonnoetherian} 之证明的论证发展
彼此类似。在二者中，完备化都起到关键作用；在两处，具有几何整纤维这一
假设都保证从 $S$ 到 $S$ 的完备化的同态是 $R$-普遍单射的。
\end{remark}













\section{有限型平坦模，第一部分}
\label{flat-section-finite-type-flat-I}

\noindent
在某些情形下，给定有限表示环同态 $R \to S$ 与有限 $S$-模 $N$，
$N$ 在 $R$ 上平坦蕴含 $N$ 有限表示。本节将证明这一结论“逐点”成立。
需要指出，命题 \ref{flat-proposition-finite-type-flat-at-point} 的第一个证明使用
第 \ref{flat-section-local-structure-module} 节的几何结果，但不使用完全
d\'evissage 的存在性。

\begin{lemma}
\label{flat-lemma-induction-step}
设 $(R, \mathfrak m)$ 是局部环，$R \to S$ 是具有几何整纤维的有限表示
平坦环同态。记 $\mathfrak p = \mathfrak mS$。设
$\mathfrak q \subset S$ 是位于 $\mathfrak m$ 上方的素理想，$N$ 是有限
$S$-模。存在 $r \geq 0$ 以及 $S$-模同态
$$
\alpha : S^{\oplus r} \longrightarrow N
$$
使得
$\alpha : \kappa(\mathfrak p)^{\oplus r} \to N \otimes_S \kappa(\mathfrak p)$
是同构。对任一这样的 $\alpha$，下列条件等价：
\begin{enumerate}
\item $N_{\mathfrak q}$ 在 $R$ 上平坦；
\item $\alpha$ 是 $R$-普遍单射的，且
$\Coker(\alpha)_{\mathfrak q}$ 在 $R$ 上平坦；
\item $\alpha$ 单射，且 $\Coker(\alpha)_{\mathfrak q}$ 在 $R$ 上平坦；
\item $\alpha_{\mathfrak p}$ 是同构，且
$\Coker(\alpha)_{\mathfrak q}$ 在 $R$ 上平坦；并且
\item $\alpha_{\mathfrak q}$ 单射，且
$\Coker(\alpha)_{\mathfrak q}$ 在 $R$ 上平坦。
\end{enumerate}
\end{lemma}

\begin{proof}
为得到 $\alpha$，置
$r = \dim_{\kappa(\mathfrak p)} N \otimes_S \kappa(\mathfrak p)$，并选取
$x_1, \ldots, x_r \in N$，使其构成
$N \otimes_S \kappa(\mathfrak p)$ 的一组基。定义
$\alpha(s_1, \ldots, s_r) = \sum s_i x_i$，便证明了存在性。

\medskip\noindent
固定一个 $\alpha$。最有意义的蕴含是 (1) $\Rightarrow$ (2)，先来证明它。
假设 (1)。因为 $S/\mathfrak mS$ 是分式域为
$\kappa(\mathfrak p)$ 的整环，可知
$(S/\mathfrak mS)^{\oplus r} \to
N_{\mathfrak p}/\mathfrak mN_{\mathfrak p} = N \otimes_S \kappa(\mathfrak p)$
单射。因此，由引理 \ref{flat-lemma-universally-injective-local} 和
\ref{flat-lemma-fibres-irreducible-flat-projective-nonnoetherian}.
同态 $S^{\oplus r} \to N_{\mathfrak p}$ 是 $R$-普遍单射的。
于是 $S^{\oplus r} \to N$ 是 $R$-普遍单射的；参见代数，引理
\ref{algebra-lemma-universally-injective-permanence}。其局部化
$\alpha_{\mathfrak q}$ 也为 $R$-普遍单射；参见代数，引理
\ref{algebra-lemma-universally-injective-localize}。由代数，引理
\ref{algebra-lemma-ui-flat-domain}，可知
$\Coker(\alpha)_{\mathfrak q}$ 在 $R$ 上平坦。

\medskip\noindent
蕴含 (2) $\Rightarrow$ (3) 是立即的。若 (3) 成立，则
$\alpha_{\mathfrak p}$ 作为单射模同态的局部化是单射的。由 Nakayama
引理（代数，引理 \ref{algebra-lemma-NAK}），$\alpha_{\mathfrak p}$
也满射。因此 (3) $\Rightarrow$ (4)。若 (4) 成立，则
$\alpha_{\mathfrak p}$ 是同构，而
$S_{\mathfrak q} \to S_{\mathfrak p}$ 单射，故
% P05-EN-CH38-0024: frozen proof says alpha is injective although this argument establishes the localized map alpha_q required by (5).
$\alpha$ 单射。事实上，结合引理 \ref{flat-lemma-invert-universally-injective}
与 \ref{flat-lemma-fibres-irreducible-flat-projective-nonnoetherian}，可知
$S \setminus \mathfrak p$ 中的元素在 $S$ 上都是非零因子。
因此 (4) $\Rightarrow$ (5)。最后，若 (5) 成立，则
$N_{\mathfrak q}$ 作为平坦模的扩张在 $R$ 上平坦；参见代数，引理
\ref{algebra-lemma-flat-ses}。故 (5) $\Rightarrow$ (1)，证明完成。
\end{proof}

\begin{lemma}
\label{flat-lemma-complete-devissage-flat-finite-type-module}
设 $(R, \mathfrak m)$ 是局部环，$R \to S$ 是有限表示环同态，
$N$ 是有限 $S$-模，$\mathfrak q$ 是 $S$ 中位于 $\mathfrak m$
上方的素理想。假设 $N_{\mathfrak q}$ 在 $R$ 上平坦，并且在
$\mathfrak q$ 处存在 $N/S/R$ 的完全 d\'evissage。则 $N$ 是有限表示
$S$-模，作为 $R$-模自由，并且存在 $R$-模同构
$$
N \cong B_1^{\oplus r_1} \oplus \ldots \oplus B_n^{\oplus r_n}
$$
其中每个 $B_i$ 都是具有几何不可约纤维的光滑 $R$-代数。
\end{lemma}

\begin{proof}
设 $(A_i, B_i, M_i, \alpha_i, \mathfrak q_i)_{i = 1, \ldots, n}$
为给定的完全 d\'evissage。对 $n$ 作归纳以证明引理。注意，$N$ 作为
$S$-模有限表示，当且仅当 $M_1$ 作为 $B_1$-模有限表示；参见注记
\ref{flat-remark-finite-presentation}。
% P05-EN-CH38-0025: frozen source has the article error “an B_1-module”.
还要注意，作为 $R$-模有 $N_{\mathfrak q} \cong (M_1)_{\mathfrak q_1}$，
因为：(a) $N_{\mathfrak q} \cong (M_1)_{\mathfrak q'_1}$，其中
$\mathfrak q'_1$ 是 $A_1$ 中位于 $\mathfrak q_1$ 上方的唯一素理想；
(b) 由代数，引理 \ref{algebra-lemma-unique-prime-over-localize-below}，
$(A_1)_{\mathfrak q'_1} = (A_1)_{\mathfrak q_1}$；从而
(c) $(M_1)_{\mathfrak q'_1} \cong (M_1)_{\mathfrak q_1}$。
因此 $(M_1)_{\mathfrak q_1}$ 是平坦 $R$-模。为证明引理，故可用
$(B_1, M_1)$ 替换 $(S, N)$。由引理 \ref{flat-lemma-induction-step}，同态
$\alpha_1 : B_1^{\oplus r_1} \to M_1$ 是 $R$-普遍单射的，并且
% P05-EN-CH38-0026: frozen proof retains q after replacing the active prime by q_1.
$\Coker(\alpha_1)_{\mathfrak q}$ 在 $R$ 上平坦。注意
$(A_i, B_i, M_i, \alpha_i, \mathfrak q_i)_{i = 2, \ldots, n}$ 是
$\mathfrak q_1$ 处 $\Coker(\alpha_1)/B_1/R$ 的完全 d\'evissage。
因此，归纳假设说明 $\Coker(\alpha_1)$ 是有限表示 $B_1$-模，
作为 $R$-模自由，并且具有引理中的分解。这又说明 $M_1$ 是有限表示
$B_1$-模；参见代数，引理 \ref{algebra-lemma-extension}。进而，作为
$R$-模有
$M_1 \cong B_1^{\oplus r_1} \oplus \Coker(\alpha_1)$，
从而得到引理中的分解。最后，由引理
\ref{flat-lemma-fibres-irreducible-flat-projective-nonnoetherian}，$B_1$
作为 $R$-模是投射的；再由代数，定理
\ref{algebra-theorem-projective-free-over-local-ring}，它作为 $R$-模自由。
证明完成。
\end{proof}

\begin{proposition}
\label{flat-proposition-finite-type-flat-at-point}
设 $f : X \to S$ 是概形态射，$\mathcal{F}$ 是 $X$ 上的拟凝聚层，
$x \in X$ 的像为 $s \in S$。假设
\begin{enumerate}
\item $f$ 局部有限表示；
\item $\mathcal{F}$ 有限型；并且
\item $\mathcal{F}$ 在点 $x$ 处于 $S$ 上平坦。
\end{enumerate}
则存在初等 \'etale 邻域 $(S', s') \to (S, s)$ 以及开子概形
$$
V \subset X \times_S \Spec(\mathcal{O}_{S', s'})
$$
它包含 $X \times_S \Spec(\mathcal{O}_{S', s'})$ 中映到 $x$ 的唯一点，
并且 $\mathcal{F}$ 在 $V$ 上的拉回是有限表示 $\mathcal{O}_V$-模，
且在 $\mathcal{O}_{S', s'}$ 上平坦。
\end{proposition}

\begin{proof}[第一个证明]
这个证明较长，但不使用完全 d\'evissage 的存在性。问题在 $x$ 与 $s$
附近是局部的，故可设 $X$ 和 $S$ 仿射。在证明过程中，我们会有限次地
% P05-EN-CH38-0027: frozen source is missing “a finite number of” in “we will finitely many times replace”.
用 $(S, s)$ 的初等 \'etale 邻域替换 $S$。此时目标是在这样的替换之后，
找到包含 $x$ 的开集
$V \subset X \times_S \Spec(\mathcal{O}_{S, s})$，使得
$\mathcal{F}|_V$ 在 $S$ 上平坦且有限表示。当然，在证明的任意阶段，也可
用 $\Spec(\mathcal{O}_{S, s})$ 替换 $S$；也就是说，可设 $S$ 为局部概形。
对整数 $n = \dim_x(\text{Supp}(\mathcal{F}_s))$ 作归纳以证明命题。

\medskip\noindent
可以选择
\begin{enumerate}
\item 初等 \'etale 邻域 $g : (X', x') \to (X, x)$ 与
$e : (S', s') \to (S, s)$；
\item 交换图
$$
\xymatrix{
X \ar[dd]_f & X' \ar[dd] \ar[l]^g & Z' \ar[l]^i \ar[d]^\pi \\
& & Y' \ar[d]^h \\
S & S' \ar[l]_e & S' \ar@{=}[l]
}
$$
\item 点 $z' \in Z'$，满足 $i(z') = x'$、$y' = \pi(z')$、
$h(y') = s'$；
\item 有限型拟凝聚 $\mathcal{O}_{Z'}$-模 $\mathcal{G}$；
\end{enumerate}
如引理 \ref{flat-lemma-elementary-devissage} 中所述。按照证明第一段的说明，
将用 $\Spec(\mathcal{O}_{S', s'})$ 替换 $S$。考虑图
$$
\xymatrix{
X_{\mathcal{O}_{S', s'}} \ar[ddr]_f &
X'_{\mathcal{O}_{S', s'}} \ar[dd] \ar[l]^g &
Z'_{\mathcal{O}_{S', s'}} \ar[l]^i \ar[d]^\pi \\
& & Y'_{\mathcal{O}_{S', s'}} \ar[dl]^h \\
& \Spec(\mathcal{O}_{S', s'})
}
$$
这里通过 $\Spec(\mathcal{O}_{S', s'}) \to S'$ 对 $S'$ 上的概形
$X', Z', Y'$ 作了基变换，并通过
$\Spec(\mathcal{O}_{S', s'}) \to S$ 对 $S$ 上的概形 $X$ 作了基变换。
$g$ 仍为 \'etale；参见引理 \ref{flat-lemma-etale-at-point}。分别用
$X_{\mathcal{O}_{S', s'}}$、$X'_{\mathcal{O}_{S', s'}}$、
$Z'_{\mathcal{O}_{S', s'}}$、$Y'_{\mathcal{O}_{S', s'}}$
替换 $X$、$X'$、$Z'$、$Y'$ 后，可设已有引理
\ref{flat-lemma-elementary-devissage} 中的图，并且 $S = S'$ 是闭点为 $s$
的局部概形。由引理 \ref{flat-lemma-devissage-finite-presentation} 和
\ref{flat-lemma-devissage-flat}，对 $Y' \to S$、层 $\pi_*\mathcal{G}$ 与点
$y'$ 的结论蕴含对 $X \to S$、$\mathcal{F}$ 与 $x$ 的结论。因此可设
$S$ 局部，且 $X \to S$ 是仿射概形之间的光滑态射，其纤维几何不可约、
维数为 $n$。

\medskip\noindent
归纳的初始情形：$n = 0$。因为 $X \to S$ 光滑，且纤维几何不可约、
维数为 $0$，所以 $X \to S$ 是开浸入；参见下降章，引理
\ref{descent-lemma-universally-injective-etale-open-immersion}。又因为 $S$
局部且闭点属于 $X \to S$ 的像，故 $X = S$。于是 $\mathcal{F}$
对应于有限平坦 $\mathcal{O}_{S, s}$ 模。在此情形，结论由代数，引理
\ref{algebra-lemma-finite-flat-local} 得出；该引理说明 $\mathcal{F}$
事实上有限自由。

\medskip\noindent
归纳步骤。假设只要闭纤维中支撑的维数 $< n$，结论就成立。记
$S = \Spec(A)$、$X = \Spec(B)$，并对某个 $B$-模 $N$ 记
$\mathcal{F} = \widetilde{N}$。注意 $A$ 是局部环，将其极大理想记为
$\mathfrak m$。由于 $X \to S$ 的纤维几何不可约，
$\mathfrak p = \mathfrak mB$ 是位于 $\mathfrak m$ 上方的唯一极小素理想。
最后，令 $\mathfrak q \subset B$ 为与 $x$ 对应的素理想。由引理
\ref{flat-lemma-induction-step}
可选择同态
$$
\alpha : B^{\oplus r} \to N
$$
使得 $\kappa(\mathfrak p)^{\oplus r} \to N \otimes_B \kappa(\mathfrak p)$
是同构。此外，因为 $N_{\mathfrak q}$ 在 $A$ 上平坦，该引理还说明
$\alpha$ 单射，且 $\Coker(\alpha)_{\mathfrak q}$ 在 $A$ 上平坦。
置 $Q = \Coker(\alpha)$。注意 $Q/\mathfrak mQ$ 的支撑不包含
$\mathfrak p$，所以必有
$\dim_{\mathfrak q}(\text{Supp}(Q/\mathfrak mQ)) < n$。综合关于 $Q$
的全部已知信息，可将归纳假设应用于 $Q$。因此存在初等 \'etale 态射
% P05-EN-CH38-0028: frozen source writes (S', s) although the local ring on the next line is at s'.
$(S', s) \to (S, s)$，使得结论对 $B \otimes_A A'$ 上的
$Q \otimes_A A'$ 成立，其中 $A' = \mathcal{O}_{S', s'}$。
用 $A'$ 替换 $A$ 后，有正合序列
$$
0 \to B^{\oplus r} \to N \to Q \to 0
$$
（这里使用了上面提到的 $\alpha$ 单射），其各项均为有限 $B$-模；
还得到元素 $g \in B$、$g \not \in \mathfrak q$，使得 $Q_g$
在 $B_g$ 上有限表示且在 $A$ 上平坦。由于局部化正合，序列
$$
0 \to B_g^{\oplus r} \to N_g \to Q_g \to 0
$$
仍正合。因为 $B_g$ 和 $Q_g$ 在 $A$ 上平坦，所以 $N_g$ 在 $A$
上平坦；参见代数，引理 \ref{algebra-lemma-flat-ses}。又因为 $B_g$
和 $Q_g$ 在 $B_g$ 上有限表示，$N_g$ 也如此；参见代数，引理
\ref{algebra-lemma-extension}。
\end{proof}

\begin{proof}[第二个证明]
应用命题 \ref{flat-proposition-existence-complete-at-x}，得到带点概形的交换图
$$
\xymatrix{
(X, x) \ar[d] & (X', x') \ar[l]^g \ar[d] \\
(S, s) & (S', s') \ar[l]
}
$$
其中水平箭头都是初等 \'etale 邻域，并且
% P05-EN-CH38-0029: frozen source locates the complete devissage of g^*F on X' at x rather than x'.
$g^*\mathcal{F}/X'/S'$ 在 $x$ 处有完全 d\'evissage（特别地，$S'$ 和
$X'$ 仿射）。由《态射》，引理 \ref{morphisms-lemma-flat-permanence}，
$g^*\mathcal{F}$ 在点 $x'$ 处于 $S$ 上平坦；再由引理
\ref{flat-lemma-etale-flat-up-down}，它在点 $x'$ 处于 $S'$ 上平坦。
由注记 \ref{flat-remark-same-notion}，可得
$$
\Gamma(X', g^*\mathcal{F})/
\Gamma(X', \mathcal{O}_{X'})/
\Gamma(S', \mathcal{O}_{S'})
$$
在 $\Gamma(X', \mathcal{O}_{X'})$ 中与 $x'$ 对应的素理想处有完全
d\'evissage。可将这个完全 d\'evissage 基变换到
$\Gamma(S', \mathcal{O}_{S'})$ 在与 $s'$ 对应的素理想处的局部环
$\mathcal{O}_{S', s'}$。因此，引理
\ref{flat-lemma-complete-devissage-flat-finite-type-module} 说明
$$
% P05-EN-CH38-0030: frozen display switches from g^*F to an undefined F'.
\Gamma(X', \mathcal{F}')
\otimes_{\Gamma(S', \mathcal{O}_{S'})}
\mathcal{O}_{S', s'}
$$
在 $\mathcal{O}_{S', s'}$ 上平坦，并且在
$\Gamma(X', \mathcal{O}_{X'})
\otimes_{\Gamma(S', \mathcal{O}_{S'})} \mathcal{O}_{S', s'}$ 上有限表示。
换言之，$\mathcal{F}$ 在
$X' \times_{S'} \Spec(\mathcal{O}_{S', s'})$ 上的限制有限表示，且在
$\mathcal{O}_{S', s'}$ 上平坦。
% P05-EN-CH38-0031: frozen prose names F on X' rather than the previously defined pullback g^*F.
由于态射
$X' \times_{S'} \Spec(\mathcal{O}_{S', s'})
\to X \times_S \Spec(\mathcal{O}_{S', s'})$
是 \'etale 的（引理 \ref{flat-lemma-etale-at-point}），其像
$V \subset X \times_S \Spec(\mathcal{O}_{S', s'})$ 是开子概形；由
\'etale 下降，$\mathcal{F}$ 在 $V$ 上的限制有限表示，且在
$\mathcal{O}_{S', s'}$ 上平坦。这里使用了《态射》，引理
\ref{morphisms-lemma-etale-open}，下降章，引理
\ref{descent-lemma-finite-presentation-descends}，以及《态射》，引理
\ref{morphisms-lemma-flat-permanence}。
\end{proof}

\begin{lemma}
\label{flat-lemma-open-in-fibre-where-flat}
设 $f : X \to S$ 是局部有限型概形态射，$\mathcal{F}$ 是有限型拟凝聚
$\mathcal{O}_X$-模，且 $s \in S$。则集合
$$
\{x \in X_s \mid \mathcal{F}\text{ 在 }x\text{ 处于 }S\text{ 上平坦}\}
$$
在纤维 $X_s$ 中是开集。
\end{lemma}

\begin{proof}
% P05-EN-CH38-0032: frozen proof refers to U although the displayed flat locus is not named U in the statement.
设 $x \in U$。按命题 \ref{flat-proposition-finite-type-flat-at-point}，选择初等
\'etale 邻域 $(S', s') \to (S, s)$ 与开集
$V \subset X \times_S \Spec(\mathcal{O}_{S', s'})$。注意，由于
$\kappa(s) = \kappa(s')$，有 $X_{s'} = X_s$。若
$x' \in V \cap X_{s'}$，则 $\mathcal{F}$ 在 $X \times_S S'$ 上的拉回
在点 $x'$ 处于 $S'$ 上平坦。因此，$\mathcal{F}$ 在点 $x'$ 处于
$S$ 上平坦；参见《态射》，引理 \ref{morphisms-lemma-flat-permanence}。
换言之，$X_s \cap V \subset U$ 是 $U$ 中点 $x$ 的开邻域。
\end{proof}

\begin{lemma}
\label{flat-lemma-finite-type-flat-at-point}
设 $f : X \to S$ 是概形态射，$\mathcal{F}$ 是 $X$ 上的拟凝聚层，
$x \in X$ 的像为 $s \in S$。假设
\begin{enumerate}
\item $f$ 局部有限型；
\item $\mathcal{F}$ 有限型；并且
\item $\mathcal{F}$ 在点 $x$ 处于 $S$ 上平坦。
\end{enumerate}
则存在初等 \'etale 邻域 $(S', s') \to (S, s)$ 以及开子概形
$$
V \subset X \times_S \Spec(\mathcal{O}_{S', s'})
$$
它包含 $X \times_S \Spec(\mathcal{O}_{S', s'})$ 中映到 $x$ 的唯一点，
并且 $\mathcal{F}$ 在 $V$ 上的拉回在 $\mathcal{O}_{S', s'}$ 上平坦。
\end{lemma}

\begin{proof}
（本引理与命题 \ref{flat-proposition-finite-type-flat-at-point} 的唯一区别是，
这里不假设 $f$ 有限表示。）问题在 $X$ 和 $S$ 上是局部的，故可设
$X$、$S$ 仿射。记 $X = \Spec(B)$、$S = \Spec(A)$，并写成
$B = A[x_1, \ldots, x_n]/I$。换言之，得到闭浸入
$i : X \to \mathbf{A}^n_S$。记 $t = i(x) \in \mathbf{A}^n_S$。
将命题 \ref{flat-proposition-finite-type-flat-at-point} 应用于
$\mathbf{A}^n_S \to S$、层 $i_*\mathcal{F}$ 与点 $t$，得到初等
\'etale 邻域 $(S', s') \to (S, s)$ 以及开子概形
$$
W \subset \mathbf{A}^n_{\mathcal{O}_{S', s'}}
$$
使得 $i_*\mathcal{F}$ 在 $W$ 上的拉回在 $\mathcal{O}_{S', s'}$ 上平坦。
这意味着
$V := W \cap \big(X \times_S \Spec(\mathcal{O}_{S', s'})\big)$
就是所需的开子概形。
\end{proof}

\begin{lemma}
\label{flat-lemma-finite-type-flat-along-fibre}
设 $f : X \to S$ 是概形态射，$\mathcal{F}$ 是 $X$ 上的拟凝聚层，
$s \in S$。假设
\begin{enumerate}
\item $f$ 有限表示；
\item $\mathcal{F}$ 有限型；并且
\item $\mathcal{F}$ 在纤维 $X_s$ 的每一点处于 $S$ 上平坦。
\end{enumerate}
则存在初等 \'etale 邻域 $(S', s') \to (S, s)$ 以及开子概形
$$
V \subset X \times_S \Spec(\mathcal{O}_{S', s'})
$$
它包含纤维 $X_s = X \times_S s'$，并且 $\mathcal{F}$ 在 $V$ 上的拉回
是有限表示 $\mathcal{O}_V$-模，且在 $\mathcal{O}_{S', s'}$ 上平坦。
\end{lemma}

\begin{proof}
对每个点 $x \in X_s$，可应用命题
\ref{flat-proposition-finite-type-flat-at-point}，找到初等 \'etale 邻域
$(S_x, s_x) \to (S, s)$ 与开集
$V_x \subset X \times_S \Spec(\mathcal{O}_{S_x, s_x})$，使得
$x \in X_s = X \times_S s_x$ 包含于 $V_x$，并且 $\mathcal{F}$ 在
$V_x$ 上的拉回是有限表示 $\mathcal{O}_{V_x}$-模，且在
$\mathcal{O}_{S_x, s_x}$ 上平坦。特别地，可将纤维 $(V_x)_{s_x}$
视为 $X_s$ 中点 $x$ 的开邻域。因为 $X_s$ 拟紧，可找到有限个点
$x_1, \ldots, x_n \in X_s$，使得 $X_s$ 是各
$(V_{x_i})_{s_{x_i}}$ 的并。选择支配每个邻域
$(S_{x_i}, s_{x_i})$ 的初等 \'etale 邻域
$(S' , s') \to (S, s)$；参见《态射进阶》，引理
\ref{more-morphisms-lemma-elementary-etale-neighbourhoods}。
置 $V = \bigcup V_i$，其中 $V_i$ 是开集 $V_{x_i}$ 在态射
% P05-EN-CH38-0033: frozen source has the number-agreement error “V_i is the inverse images”.
$$
X \times_S \Spec(\mathcal{O}_{S', s'})
\longrightarrow
X \times_S \Spec(\mathcal{O}_{S_{x_i}, s_{x_i}})
$$
下的逆像。按构造，$V$ 包含 $X_s$；同样按构造，$\mathcal{F}$ 在
$V$ 上的拉回是有限表示 $\mathcal{O}_V$-模，且在
$\mathcal{O}_{S', s'}$ 上平坦。
\end{proof}

\begin{lemma}
\label{flat-lemma-finite-type-flat-along-fibre-variant}
设 $f : X \to S$ 是概形态射，$\mathcal{F}$ 是 $X$ 上的拟凝聚层，
$s \in S$。假设
\begin{enumerate}
\item $f$ 有限型；
\item $\mathcal{F}$ 有限型；并且
\item $\mathcal{F}$ 在纤维 $X_s$ 的每一点处于 $S$ 上平坦。
\end{enumerate}
则存在初等 \'etale 邻域 $(S', s') \to (S, s)$ 以及开子概形
$$
V \subset X \times_S \Spec(\mathcal{O}_{S', s'})
$$
它包含纤维 $X_s = X \times_S s'$，并且 $\mathcal{F}$ 在 $V$ 上的拉回
在 $\mathcal{O}_{S', s'}$ 上平坦。
\end{lemma}

\begin{proof}
（本引理与引理 \ref{flat-lemma-finite-type-flat-along-fibre} 的唯一区别是，
这里不假设 $f$ 有限表示。）对每个点 $x \in X_s$，可应用引理
\ref{flat-lemma-finite-type-flat-at-point}，找到初等 \'etale 邻域
$(S_x, s_x) \to (S, s)$ 与开集
$V_x \subset X \times_S \Spec(\mathcal{O}_{S_x, s_x})$，使得
$x \in X_s = X \times_S s_x$ 包含于 $V_x$，并且 $\mathcal{F}$ 在
$V_x$ 上的拉回在 $\mathcal{O}_{S_x, s_x}$ 上平坦。特别地，可将纤维
$(V_x)_{s_x}$ 视为 $X_s$ 中点 $x$ 的开邻域。因为 $X_s$ 拟紧，
可找到有限个点 $x_1, \ldots, x_n \in X_s$，使得 $X_s$ 是各
$(V_{x_i})_{s_{x_i}}$ 的并。选择支配每个邻域
$(S_{x_i}, s_{x_i})$ 的初等 \'etale 邻域
$(S' , s') \to (S, s)$；参见《态射进阶》，引理
\ref{more-morphisms-lemma-elementary-etale-neighbourhoods}。
置 $V = \bigcup V_i$，其中 $V_i$ 是开集 $V_{x_i}$ 在态射
% P05-EN-CH38-0034: frozen source repeats the number-agreement error “V_i is the inverse images”.
$$
X \times_S \Spec(\mathcal{O}_{S', s'})
\longrightarrow
X \times_S \Spec(\mathcal{O}_{S_{x_i}, s_{x_i}})
$$
下的逆像。按构造，$V$ 包含 $X_s$；同样按构造，$\mathcal{F}$ 在
$V$ 上的拉回在 $\mathcal{O}_{S', s'}$ 上平坦。
\end{proof}

\begin{lemma}
\label{flat-lemma-finite-type-flat-at-point-X}
设 $S$ 是概形，$X$ 在 $S$ 上局部有限型，$x \in X$ 的像为
$s \in S$。若 $X$ 在点 $x$ 处于 $S$ 上平坦，则存在初等
\'etale 邻域 $(S', s') \to (S, s)$ 以及开子概形
$$
V \subset X \times_S \Spec(\mathcal{O}_{S', s'})
$$
它包含 $X \times_S \Spec(\mathcal{O}_{S', s'})$ 中映到 $x$ 的唯一点，
并且 $V \to \Spec(\mathcal{O}_{S', s'})$ 平坦且有限表示。
\end{lemma}

\begin{proof}
问题在 $X$ 和 $S$ 上是局部的，故可设 $X$、$S$ 仿射。记
$X = \Spec(B)$、$S = \Spec(A)$，并写成
$B = A[x_1, \ldots, x_n]/I$。换言之，得到闭浸入
$i : X \to \mathbf{A}^n_S$。记 $t = i(x) \in \mathbf{A}^n_S$。
将命题 \ref{flat-proposition-finite-type-flat-at-point} 应用于
$\mathbf{A}^n_S \to S$、层 $\mathcal{F} = i_*\mathcal{O}_X$ 与点 $t$，
得到初等 \'etale 邻域 $(S', s') \to (S, s)$ 以及开子概形
$$
W \subset \mathbf{A}^n_{\mathcal{O}_{S', s'}}
$$
使得 $i_*\mathcal{O}_X$ 的拉回平坦且有限表示。这意味着
$V := W \cap \big(X \times_S \Spec(\mathcal{O}_{S', s'})\big)$
就是所需的开子概形。
\end{proof}

\begin{lemma}
\label{flat-lemma-finite-type-flat-at-point-local}
设 $f : X \to S$ 是局部有限表示态射，$\mathcal{F}$ 是有限型拟凝聚
$\mathcal{O}_X$-模。若 $x \in X$，且 $\mathcal{F}$ 在点 $x$ 处于
$S$ 上平坦，则 $\mathcal{F}_x$ 是有限表示 $\mathcal{O}_{X, x}$-模。
\end{lemma}

\begin{proof}
令 $s = f(x)$。由命题 \ref{flat-proposition-finite-type-flat-at-point}，
存在初等 \'etale 邻域 $(S', s') \to (S, s)$，使得 $\mathcal{F}$ 在
$X \times_S \Spec(\mathcal{O}_{S', s'})$ 上的拉回，在与 $x$ 对应的点
$x' \in X_{s'} = X_s$ 的某个邻域中有限表示。环同态
$$
\mathcal{O}_{X, x} \longrightarrow
\mathcal{O}_{X \times_S \Spec(\mathcal{O}_{S', s'}), x'}
=
\mathcal{O}_{X \times_S S', x'}
$$
作为 \'etale 环同态的局部化，是平坦局部同态。因此由下降，
$\mathcal{F}_x$ 在 $\mathcal{O}_{X, x}$ 上有限表示；参见代数，引理
\ref{algebra-lemma-descend-properties-modules}（还要使用平坦局部环同态
忠实平坦这一事实；参见代数，引理 \ref{algebra-lemma-local-flat-ff}）。
\end{proof}

\begin{lemma}
\label{flat-lemma-finite-type-flat-at-point-local-X}
设 $f : X \to S$ 是局部有限型态射，$x \in X$ 的像为 $s \in S$。
若 $f$ 在点 $x$ 处于 $S$ 上平坦，则 $\mathcal{O}_{X, x}$ 在
$\mathcal{O}_{S, s}$ 上本质有限表示。
\end{lemma}

\begin{proof}
可设 $X$ 和 $S$ 仿射。记 $X = \Spec(B)$、$S = \Spec(A)$，并写成
$B = A[x_1, \ldots, x_n]/I$。换言之，得到闭浸入
$i : X \to \mathbf{A}^n_S$。记 $t = i(x) \in \mathbf{A}^n_S$。
将引理 \ref{flat-lemma-finite-type-flat-at-point-local} 应用于
$\mathbf{A}^n_S \to S$、层 $\mathcal{F} = i_*\mathcal{O}_X$ 与点 $t$，
可得 $\mathcal{O}_{X, x}$ 在 $\mathcal{O}_{\mathbf{A}^n_S, t}$
上有限表示，这蕴含所求结论。
\end{proof}







\section{从开集延拓性质}
\label{flat-section-extending-properties}

\noindent
本节汇集若干下述形式的结果：若 $f : X \to S$ 是平坦概形态射，且
$f$ 在 $S$ 的某个稠密开集上满足某一性质，则 $f$ 在整个 $S$ 上满足同一性质。

\begin{lemma}
\label{flat-lemma-flat-finite-type-finitely-presented-over-dense-open}
\begin{slogan}
在一个（良好）开集上延拓有限表示模所得的 $S$-平坦有限型模，
也是 $X$-有限表示的。
\end{slogan}
设 $f : X \to S$ 是概形态射，$\mathcal{F}$ 是拟凝聚
$\mathcal{O}_X$-模，$U \subset S$ 是开集。假设
\begin{enumerate}
\item $f$ 局部有限表示；
\item $\mathcal{F}$ 有限型且在 $S$ 上平坦；
\item $U \subset S$ 逆紧且概形论稠密；
\item $\mathcal{F}|_{f^{-1}U}$ 有限表示。
\end{enumerate}
则 $\mathcal{F}$ 有限表示。
\end{lemma}

\begin{proof}
问题在 $X$ 和 $S$ 上是局部的，故可设 $X$、$S$ 仿射。记
$S = \Spec(A)$、$X = \Spec(B)$。设 $N$ 是有限 $B$-模，且
$\mathcal{F}$ 是与 $N$ 相伴的拟凝聚层。对某些 $f_i \in A$，有
$U = D(f_1) \cup \ldots \cup D(f_n)$；参见代数，引理
\ref{algebra-lemma-qc-open}。因为 $U$ 概形论稠密，同态
$A \to A_{f_1} \times \ldots \times A_{f_n}$ 单射。取位于
$\mathfrak p \subset A$ 上方的素理想 $\mathfrak q \subset B$；它对应于
映到 $s \in S$ 的点 $x \in X$。由引理
\ref{flat-lemma-finite-type-flat-at-point-local}，模 $N_\mathfrak q$ 在
$B_\mathfrak q$ 上有限表示。选择 $B$-模满射
$\varphi : B^{\oplus m} \to N$。选择
$k_1, \ldots, k_t \in \Ker(\varphi)$，并置
$N' = B^{\oplus m}/\sum Bk_j$。存在典范满射 $N' \to N$，且 $N$
是以这种方式构造的 $B$-模 $N'$ 的滤过余极限。因此可选择
$k_1, \ldots, k_t$，使得 (a) $N'_{f_i} \cong N_{f_i}$，
$i = 1, \ldots, n$；以及 (b) $N'_\mathfrak q \cong N_\mathfrak q$。
特别地，这说明 $N'_\mathfrak q$ 在 $A$ 上平坦。由平坦性的开性
（参见代数，定理 \ref{algebra-theorem-openness-flatness}），存在
$g \in B$、$g \not \in \mathfrak q$，使得 $N'_g$ 在 $A$ 上平坦。
考虑交换图
$$
\xymatrix{
N'_g \ar[r] \ar[d] & N_g \ar[d] \\
\prod N'_{gf_i} \ar[r] & \prod N_{gf_i}
}
$$
由 $k_1, \ldots, k_t$ 的选择，底部箭头是同构。因为
$A \to \prod A_{f_i}$ 单射且 $N'_g$ 在 $A$ 上平坦，左侧竖直箭头单射。
因此顶部水平箭头单射，从而是同构。这证明 $N_g$ 在 $B_g$ 上有限表示。
最后应用代数，引理 \ref{algebra-lemma-cover} 即得结论。
\end{proof}

\begin{lemma}
\label{flat-lemma-flat-finite-type-finitely-presented-over-dense-open-X}
设 $f : X \to S$ 是概形态射，$U \subset S$ 是开集。假设
\begin{enumerate}
\item $f$ 局部有限型且平坦；
\item $U \subset S$ 逆紧且概形论稠密；
\item $f|_{f^{-1}U} : f^{-1}U \to U$ 局部有限表示。
\end{enumerate}
% P05-EN-CH38-0035: frozen conclusion has the duplicated preposition “is of locally of finite presentation”.
则 $f$ 局部有限表示。
\end{lemma}

\begin{proof}
问题在 $X$ 和 $S$ 上是局部的，故可设 $X$、$S$ 仿射。选择闭浸入
$i : X \to \mathbf{A}^n_S$，并将引理
\ref{flat-lemma-flat-finite-type-finitely-presented-over-dense-open} 应用于
$i_*\mathcal{O}_X$。省略若干细节。
\end{proof}

\begin{lemma}
\label{flat-lemma-flat-finite-presentation-dimension-over-dense-open}
设 $f : X \to S$ 是平坦且局部有限型的概形态射，$U \subset S$ 是稠密
开集，并且 $X_U \to U$ 的相对维数 $\leq e$；参见《态射》，定义
\ref{morphisms-definition-relative-dimension-d}。若还满足下列条件之一：
\begin{enumerate}
\item $f$ 局部有限表示；或
\item $U \subset S$ 逆紧；
\end{enumerate}
则 $f$ 的相对维数 $\leq e$。
\end{lemma}

\begin{proof}
先证明情形 (1)。设 $W \subset X$ 是《态射进阶》下列引理中构造并研究的开子概形：
\ref{more-morphisms-lemma-flat-finite-presentation-CM-open} 和
\ref{more-morphisms-lemma-flat-finite-presentation-CM-pieces}.
注意，每个纤维的每个一般点都包含于 $W$，所以只须对 $W$ 证明结论。
由于 $W = \bigcup_{d \geq 0} U_d$，只须证明当 $d > e$ 时
$U_d = \emptyset$。因为 $f$ 平坦且局部有限表示，所以它是开态射，
因而 $f(U_d)$ 是开集（《态射》，引理 \ref{morphisms-lemma-fppf-open}）。
因此，若 $U_d$ 非空，则 $f(U_d) \cap U \not = \emptyset$，如所需。

\medskip\noindent
再证明情形 (2)。可用其既约化替换 $S$，此时 $U$ 概形论稠密。
于是由引理 \ref{flat-lemma-flat-finite-type-finitely-presented-over-dense-open-X}，
$f$ 局部有限表示。这样便归结到情形 (1)。
\end{proof}

\begin{lemma}
\label{flat-lemma-proper-flat-finite-over-dense-open}
设 $f : X \to S$ 是平坦固有概形态射，$U \subset S$ 是稠密开集，
且 $X_U \to U$ 有限。若 $f$ 局部有限表示或 $U \subset S$ 逆紧，
则 $f$ 有限。
\end{lemma}

\begin{proof}
由引理 \ref{flat-lemma-flat-finite-presentation-dimension-over-dense-open}，
$f$ 的纤维维数为零。因此 $f$ 拟有限（《态射》，引理
\ref{morphisms-lemma-locally-quasi-finite-rel-dimension-0}），从而具有有限
纤维（《态射》，引理 \ref{morphisms-lemma-quasi-finite}）。故由《态射进阶》，
引理 \ref{more-morphisms-lemma-characterize-finite}，$f$ 有限。
\end{proof}

\begin{lemma}
\label{flat-lemma-zariski}
设 $f : X \to S$ 是概形态射，$U \subset S$ 是开集。若
\begin{enumerate}
\item $f$ 分离、局部有限型且平坦；
\item $f^{-1}(U) \to U$ 是同构；并且
\item $U \subset S$ 逆紧且概形论稠密；
\end{enumerate}
则 $f$ 是开浸入。
\end{lemma}

\begin{proof}
由引理 \ref{flat-lemma-flat-finite-type-finitely-presented-over-dense-open-X}，
态射 $f$ 局部有限表示。像 $f(X) \subset S$ 是开集（《态射》，引理
\ref{morphisms-lemma-fppf-open}），故可用 $f(X)$ 替换 $S$。于是须证明
$f$ 是同构。可设 $S$ 仿射。还可归结到 $X$ 拟紧的情形，因为只须证明：
像为 $S$ 的任一拟紧开集 $X' \subset X$ 都同构地映到 $S$。因此可设
$f$ 拟紧。由引理
\ref{flat-lemma-flat-finite-presentation-dimension-over-dense-open}，$f$ 的全部纤维
维数都是 $0$。故 $f$ 拟有限；参见《态射》，引理
\ref{morphisms-lemma-locally-quasi-finite-rel-dimension-0}。
取 $s \in S$。选择初等 \'etale 邻域
$g : (T, t) \to (S, s)$，使得 $X \times_S T = V \amalg W$，其中
$V \to T$ 有限且 $W_t = \emptyset$；参见《态射进阶》，引理
\ref{more-morphisms-lemma-etale-splits-off-quasi-finite-part}。将给定态射记为
$\pi : V \amalg W \to T$。因为 $\pi$ 平坦且局部有限表示，
$\pi(V)$ 在 $T$ 中是开集（《态射》，引理 \ref{morphisms-lemma-fppf-open}）。
缩小 $T$ 后，可设 $T = \pi(V)$。由于 $f$ 在 $U$ 上是同构，
$\pi$ 在 $g^{-1}U$ 上是同构。结合 $\pi(V) = T$，这说明
$\pi^{-1}g^{-1}U$ 包含于 $V$。由《态射》，引理
\ref{morphisms-lemma-flat-morphism-scheme-theoretically-dense-open}，
$\pi^{-1}g^{-1}U \subset V \amalg W$ 概形论稠密。因此
$W = \emptyset$。于是 $X \times_S T = V$ 在 $T$ 上有限。这说明在用
$s$ 的某个开邻域替换 $S$ 后，$f$ 有限；例如可用下降章，引理
\ref{descent-lemma-descending-property-finite}。此时 $f$ 有限局部自由
（《态射》，引理 \ref{morphisms-lemma-finite-flat}）；再将 $S$ 缩小为
$s$ 的更小开邻域，可设 $f$ 是某个次数 $d$ 的有限局部自由态射
（《态射》，引理 \ref{morphisms-lemma-finite-locally-free}）。但由于在稠密
开集 $U$ 上次数为 $1$，显然有 $d = 1$。故 $f$ 是同构。
\end{proof}





\section{有限表示平坦模}
\label{flat-section-finitely-presented-flat}

\noindent
在某些情形下，给定一个有限表示的环映射 $R \to S$ 以及一个有限表示的
$S$-模 $N$，$N$ 在 $R$ 上的平坦性蕴含 $N$ 作为 $R$-模是投射的，
至少在用一个 \'etale 扩张替换 $S$ 之后如此。
% P05-EN-CH38-0036: Frozen English has the grammatical typo "collect a some results".
本节汇集若干这类结果。

\begin{lemma}
\label{flat-lemma-induction-step-fp}
设 $R$ 是一个环。设 $R \to S$ 是一个有限表示的平坦环映射，
其纤维几何整。设 $\mathfrak q \subset S$ 是位于素理想
$\mathfrak r \subset R$ 之上的一个素理想。置
$\mathfrak p = \mathfrak r S$。设 $N$ 是一个有限表示的 $S$-模。
则存在 $r \geq 0$ 以及一个 $S$-模映射
$$
\alpha : S^{\oplus r} \longrightarrow N
$$
使得
$\alpha : \kappa(\mathfrak p)^{\oplus r} \to N \otimes_S \kappa(\mathfrak p)$
是同构。对于任意这样的 $\alpha$，以下条件等价：
\begin{enumerate}
\item $N_{\mathfrak q}$ 是 $R$-平坦的；
\item 存在 $f \in R$，$f \not \in \mathfrak r$，使得
$\alpha_f : S_f^{\oplus r} \to N_f$ 是 $R_f$-普遍单射的，并且存在
$g \in S$，$g \not \in \mathfrak q$，使得 $\Coker(\alpha)_g$
是 $R$-平坦的；
\item $\alpha_{\mathfrak r}$ 是 $R_{\mathfrak r}$-普遍单射的，并且
$\Coker(\alpha)_{\mathfrak q}$ 是 $R$-平坦的；
\item $\alpha_{\mathfrak r}$ 是单射，并且
$\Coker(\alpha)_{\mathfrak q}$ 是 $R$-平坦的；
\item $\alpha_{\mathfrak p}$ 是同构，并且
$\Coker(\alpha)_{\mathfrak q}$ 是 $R$-平坦的；
\item $\alpha_{\mathfrak q}$ 是单射，并且
$\Coker(\alpha)_{\mathfrak q}$ 是 $R$-平坦的。
\end{enumerate}
\end{lemma}

\begin{proof}
为得到 $\alpha$，置
$r = \dim_{\kappa(\mathfrak p)} N \otimes_S \kappa(\mathfrak p)$，并选取
$x_1, \ldots, x_r \in N$，使它们构成
$N \otimes_S \kappa(\mathfrak p)$ 的一组基。定义
$\alpha(s_1, \ldots, s_r) = \sum s_i x_i$。这就证明了存在性。

\medskip\noindent
固定 $\alpha$ 的一个选择。我们可以将
引理 \ref{flat-lemma-induction-step}
应用于映射
$\alpha_{\mathfrak r} : S_{\mathfrak r}^{\oplus r} \to N_{\mathfrak r}$.
于是可见 (1)、(3)、(4)、(5) 和 (6) 彼此等价。
又因 (2) 蕴含 (3) 也是显然的，我们只需证明 (1) 蕴含 (2)。

\medskip\noindent
假设 (1) 成立。由平坦性的开放性，参见
《代数》，定理 \ref{algebra-theorem-openness-flatness}，集合
$$
U_1 = \{\mathfrak q' \subset S \mid N_{\mathfrak q'}\text{ 在 }R\text{ 上平坦}\}
$$
在 $\Spec(S)$ 中是开的。由假设，它包含 $\mathfrak q$，因而也包含
$\mathfrak p$。由于 $S^{\oplus r}$ 和 $N$ 都是有限表示的 $S$-模，集合
$$
U_2 = \{\mathfrak q' \subset S \mid
\alpha_{\mathfrak q'}\text{ 是同构}\}
$$
在 $\Spec(S)$ 中是开的，参见
《代数》，引理 \ref{algebra-lemma-map-between-finitely-presented}。
由 (5)，它包含 $\mathfrak p$。因为 $R \to S$ 是有限表示且平坦的，映射
$\Phi : \Spec(S) \to \Spec(R)$ 是开映射，参见
《代数》，命题 \ref{algebra-proposition-fppf-open}。
对于任意素理想 $\mathfrak r' \in \Phi(U_1 \cap U_2)$，存在一个位于
$\mathfrak r'$ 之上的素理想 $\mathfrak q'$，使得 $N_{\mathfrak q'}$
是平坦的且 $\alpha_{\mathfrak q'}$ 是同构。这蕴含
$\alpha \otimes \kappa(\mathfrak p')$ 是同构，其中
$\mathfrak p' = \mathfrak r' S$。因此，由蕴含关系
(1) $\Rightarrow$ (3)，$\alpha_{\mathfrak r'}$ 是
$R_{\mathfrak r'}$-普遍单射的。
故若选取 $f \in R$，$f \not \in \mathfrak r$，使得
$D(f) \subset \Phi(U_1 \cap U_2)$，则可推出 $\alpha_f$ 是
$R_f$-普遍单射的，参见
《代数》，引理 \ref{algebra-lemma-universally-injective-check-stalks}。
同样的论证还表明，对于任意
$\mathfrak q' \in U_1 \cap \Phi^{-1}(\Phi(U_1 \cap U_2))$
模 $\Coker(\alpha)_{\mathfrak q'}$ 都是 $R$-平坦的。注意
$\mathfrak q \in U_1 \cap \Phi^{-1}(\Phi(U_1 \cap U_2))$。
因此可以找到 $g \in S$，$g \not \in \mathfrak q$，使得
$D(g) \subset U_1 \cap \Phi^{-1}(\Phi(U_1 \cap U_2))$，结论得证。
\end{proof}

\begin{lemma}
\label{flat-lemma-complete-devissage-flat-finitely-presented-module}
设 $R \to S$ 是一个有限表示的环映射。设 $N$ 是一个有限表示的
$S$-模，并且在 $R$ 上平坦。设 $\mathfrak r \subset R$ 是一个素理想。
假设在 $\mathfrak r$ 上存在 $N/S/R$ 的一个完全 d\'evissage。
则存在 $f \in R$，$f \not \in \mathfrak r$，使得
$$
N_f \cong B_1^{\oplus r_1} \oplus \ldots \oplus B_n^{\oplus r_n}
$$
% P05-EN-CH38-0037: Frozen English says "as R-modules" after localization;
% the displayed decomposition naturally concerns R_f-modules. Translation follows source.
作为 $R$-模成立，其中每个 $B_i$ 都是一个纤维几何不可约的光滑
$R_f$-代数。此外，$N_f$ 作为 $R_f$-模是投射的。
\end{lemma}

\begin{proof}
设 $(A_i, B_i, M_i, \alpha_i)_{i = 1, \ldots, n}$ 是给定的完全
d\'evissage。我们对 $n$ 作归纳来证明该引理。注意，该引理的断言完全
关乎 $N$ 作为 $R$-模的结构。因此可以用 $M_1$ 替换 $N$，并将 $M_1$
视为一个 $B_1$-模。关于 $M_1$ 为何是有限表示的 $B_1$-模，参见
注记 \ref{flat-remark-finite-presentation}。由
引理 \ref{flat-lemma-induction-step-fp}，在对某个
$f \in R$，$f \not \in \mathfrak r$ 用 $R_f$ 替换 $R$ 之后，可以假设
映射 $\alpha_1 : B_1^{\oplus r_1} \to M_1$ 是 $R$-普遍单射的。
由于 $M_1$ 和 $B_1^{\oplus r_1}$ 都是 $R$-平坦的，并且作为
$B_1$-模有限表示，可见 $\Coker(\alpha_1)$ 是 $R$-平坦的
（《代数》，引理 \ref{algebra-lemma-ui-flat-domain}），
且作为 $B_1$-模有限表示。注意，
$(A_i, B_i, M_i, \alpha_i)_{i = 2, \ldots, n}$ 是
$\Coker(\alpha_1)$ 的一个完全 d\'evissage。因此，归纳假设蕴含：
在对某个 $f \in R$，$f \not \in \mathfrak r$ 用 $R_f$ 替换 $R$ 后，
可以假设 $\Coker(\alpha_1)$ 具有该引理所述的分解，并且是投射的。
特别地，
% P05-EN-CH38-0038: Frozen English uses equality for a splitting obtained up to isomorphism.
$M_1 = B_1^{\oplus r_1} \oplus \Coker(\alpha_1)$.
这证明了关于分解的断言。关于投射性的断言来自以下事实：由
引理 \ref{flat-lemma-fibres-irreducible-flat-projective-nonnoetherian}，
$B_1$ 作为 $R$-模是投射的。
\end{proof}

\begin{remark}
\label{flat-remark-complete-devissage-flat-finitely-presented-module}
引理 \ref{flat-lemma-complete-devissage-flat-finitely-presented-module}
有如下变体：我们弱化平坦性条件，只假设 $N$ 在位于 $\mathfrak r$
之上的某个给定素理想 $\mathfrak q$ 处平坦；但同时加强 d\'evissage
条件，假设存在一个{\it 在 $\mathfrak q$ 处}的完全 d\'evissage。
与引理 \ref{flat-lemma-complete-devissage-flat-finite-type-module} 比较。
\end{remark}

\noindent
以下是本节的主要结果。

\begin{proposition}
\label{flat-proposition-finite-presentation-flat-at-point}
设 $f : X \to S$ 是一个概形态射。设 $\mathcal{F}$ 是 $X$ 上的一个
拟凝聚层。设 $x \in X$ 的像为 $s \in S$。假设
\begin{enumerate}
\item $f$ 局部有限表示；
\item $\mathcal{F}$ 有限表示；
\item $\mathcal{F}$ 在 $x$ 处于 $S$ 上平坦。
\end{enumerate}
则存在带基点概形的一个交换图
$$
\xymatrix{
(X, x) \ar[d] & (X', x') \ar[l]^g \ar[d] \\
(S, s) & (S', s') \ar[l]
}
$$
其中水平箭头都是初等 \'etale 邻域，$X'$、$S'$ 都是仿射的，并且
$\Gamma(X', g^*\mathcal{F})$ 是一个投射的
$\Gamma(S', \mathcal{O}_{S'})$-模。
\end{proposition}

\begin{proof}
由平坦性的开放性，参见
《态射进阶》，定理 \ref{more-morphisms-theorem-openness-flatness}，
可以用 $x$ 的一个开邻域替换 $X$，并假设 $\mathcal{F}$ 在 $S$ 上平坦。
接着应用命题 \ref{flat-proposition-existence-complete-at-x}，得到命题所述的
一个图，使得 $g^*\mathcal{F}/X'/S'$ 在 $s'$ 上有一个完全
d\'evissage。（特别地，$S'$ 和 $X'$ 都是仿射的。）由
《态射》，引理 \ref{morphisms-lemma-flat-permanence}，可见
$g^*\mathcal{F}$ 在 $S$ 上平坦；再由
引理 \ref{flat-lemma-etale-flat-up-down}，可见它在 $S'$ 上平坦。
通过注记 \ref{flat-remark-same-notion}，我们推出
$$
\Gamma(X', g^*\mathcal{F})/
\Gamma(X', \mathcal{O}_{X'})/
\Gamma(S', \mathcal{O}_{S'})
$$
在 $\Gamma(S', \mathcal{O}_{S'})$ 中对应于 $s'$ 的素理想上有一个完全
d\'evissage。因此，引理
\ref{flat-lemma-complete-devissage-flat-finitely-presented-module} 蕴含：
用 $s'$ 的一个标准开邻域替换 $S'$ 后，命题的结论成立。
\end{proof}

\noindent
在本节余下部分，我们证明该结果的若干变体。第一个是一个“整体”版本。

\begin{lemma}
\label{flat-lemma-finite-presentation-flat-along-fibre}
设 $f : X \to S$ 是一个概形态射。设 $\mathcal{F}$ 是 $X$ 上的一个
拟凝聚层。设 $s \in S$。假设
\begin{enumerate}
\item $f$ 有限表示；
\item $\mathcal{F}$ 有限表示；
\item $\mathcal{F}$ 在纤维 $X_s$ 的每一点处于 $S$ 上平坦。
\end{enumerate}
则存在一个初等 \'etale 邻域 $(S', s') \to (S, s)$ 以及概形的一个交换图
$$
\xymatrix{
X \ar[d] & X' \ar[l]^g \ar[d] \\
S & S' \ar[l]
}
$$
使得 $g$ 是 \'etale 的，$X_s \subset g(X')$，概形 $X'$、$S'$ 都是
仿射的，并且 $\Gamma(X', g^*\mathcal{F})$ 是一个投射的
$\Gamma(S', \mathcal{O}_{S'})$-模。
\end{lemma}

\begin{proof}
对于每一点 $x \in X_s$，可以应用
命题 \ref{flat-proposition-finite-presentation-flat-at-point} 找到一个交换图
$$
\xymatrix{
(X, x) \ar[d] & (Y_x, y_x) \ar[l]^{g_x} \ar[d] \\
(S, s) & (S_x, s_x) \ar[l]
}
$$
其水平箭头都是初等 \'etale 邻域，$Y_x$、$S_x$ 都是仿射的，并且
$\Gamma(Y_x, g_x^*\mathcal{F})$ 是一个投射的
$\Gamma(S_x, \mathcal{O}_{S_x})$-模。特别地，
$g_x(Y_x) \cap X_s$ 是 $X_s$ 中 $x$ 的一个开邻域。
由于 $X_s$ 拟紧，可以找到有限多个点 $x_1, \ldots, x_n \in X_s$，
使 $X_s$ 是各 $g_{x_i}(Y_{x_i}) \cap X_s$ 的并。选取一个支配每个邻域
$(S_{x_i}, s_{x_i})$ 的初等 \'etale 邻域
$(S' , s') \to (S, s)$，参见《态射进阶》，
引理 \ref{more-morphisms-lemma-elementary-etale-neighbourhoods}。
还可以假设 $S'$ 是仿射的。置
$X' = \coprod Y_{x_i} \times_{S_{x_i}} S'$，并赋予它显然的态射
$g : X' \to X$。由构造，$g(X')$ 包含 $X_s$，并且
$$
\Gamma(X', g^*\mathcal{F})
=
\bigoplus \Gamma(Y_{x_i}, g_{x_i}^*\mathcal{F})
\otimes_{\Gamma(S_{x_i}, \mathcal{O}_{S_{x_i}})}
\Gamma(S', \mathcal{O}_{S'}).
$$
这是一个投射的 $\Gamma(S', \mathcal{O}_{S'})$-模，参见
《代数》，引理 \ref{algebra-lemma-ascend-properties-modules}。
\end{proof}

\noindent
以下两个引理是上述结果在 $\mathcal{F} = \mathcal{O}_X$ 情形下的改述。

\begin{lemma}
\label{flat-lemma-finite-presentation-flat-at-point-X}
设 $f : X \to S$ 局部有限表示。设 $x \in X$ 的像为 $s \in S$。
若 $f$ 在 $x$ 处于 $S$ 上平坦，则存在带基点概形的一个交换图
$$
\xymatrix{
(X, x) \ar[d] & (X', x') \ar[l]^g \ar[d] \\
(S, s) & (S', s') \ar[l]
}
$$
其水平箭头都是初等 \'etale 邻域，$X'$、$S'$ 都是仿射的，并且
$\Gamma(X', \mathcal{O}_{X'})$ 是一个投射的
$\Gamma(S', \mathcal{O}_{S'})$-模。
\end{lemma}

\begin{proof}
这是命题 \ref{flat-proposition-finite-presentation-flat-at-point} 的一个特例。
\end{proof}

\begin{lemma}
\label{flat-lemma-finite-presentation-flat-along-fibre-X}
设 $f : X \to S$ 有限表示。设 $s \in S$。若 $X$ 在 $X_s$ 的所有点处
于 $S$ 上平坦，则存在一个初等 \'etale 邻域
$(S', s') \to (S, s)$ 以及概形的一个交换图
$$
\xymatrix{
X \ar[d] & X' \ar[l]^g \ar[d] \\
S & S' \ar[l]
}
$$
其中 $g$ 是 \'etale 的，$X_s \subset g(X')$，$X'$、$S'$ 都是仿射的，
并且 $\Gamma(X', \mathcal{O}_{X'})$ 是一个投射的
$\Gamma(S', \mathcal{O}_{S'})$-模。
\end{lemma}

\begin{proof}
这是引理 \ref{flat-lemma-finite-presentation-flat-along-fibre} 的一个特例。
\end{proof}

\noindent
以下各引理说明命题 \ref{flat-proposition-finite-presentation-flat-at-point}
在只假设态射与层为有限型（而不必有限表示）时的推论。

\begin{lemma}
\label{flat-lemma-finite-type-flat-at-point-free}
设 $f : X \to S$ 是一个概形态射。设 $\mathcal{F}$ 是 $X$ 上的一个
拟凝聚层。设 $x \in X$ 的像为 $s \in S$。假设
\begin{enumerate}
\item $f$ 局部有限表示；
\item $\mathcal{F}$ 为有限型；
\item $\mathcal{F}$ 在 $x$ 处于 $S$ 上平坦。
\end{enumerate}
则存在一个初等 \'etale 邻域 $(S', s') \to (S, s)$ 以及带基点概形的
一个交换图
$$
\xymatrix{
(X, x) \ar[d] & (X', x') \ar[l]^g \ar[d] \\
(S, s) & (\Spec(\mathcal{O}_{S', s'}), s') \ar[l]
}
$$
使得 $X' \to X \times_S \Spec(\mathcal{O}_{S', s'})$ 是 \'etale 的，
$\kappa(x) = \kappa(x')$，概形 $X'$ 在 $\mathcal{O}_{S', s'}$ 上仿射且
有限表示，层 $g^*\mathcal{F}$ 在 $\mathcal{O}_{X'}$ 上有限表示，并且
$\Gamma(X', g^*\mathcal{F})$ 是一个自由的
$\mathcal{O}_{S', s'}$-模。
\end{lemma}

\begin{proof}
为证明该引理，可以用任意初等 \'etale 邻域替换 $(S, s)$，也可以用
$\Spec(\mathcal{O}_{S, s})$ 替换 $S$。因此，由
命题 \ref{flat-proposition-finite-type-flat-at-point}，可以假设
$\mathcal{F}$ 在 $x$ 的一个邻域中有限表示且在 $S$ 上平坦。
在此情形下，结论来自
命题 \ref{flat-proposition-finite-presentation-flat-at-point}，因为
《代数》，定理 \ref{algebra-theorem-projective-free-over-local-ring}
保证在局部环上投射 $=$ 自由。
\end{proof}

\begin{lemma}
\label{flat-lemma-finite-type-flat-at-point-free-variant}
设 $f : X \to S$ 是一个概形态射。设 $\mathcal{F}$ 是 $X$ 上的一个
拟凝聚层。设 $x \in X$ 的像为 $s \in S$。假设
\begin{enumerate}
\item $f$ 局部有限型；
\item $\mathcal{F}$ 为有限型；
\item $\mathcal{F}$ 在 $x$ 处于 $S$ 上平坦。
\end{enumerate}
则存在一个初等 \'etale 邻域 $(S', s') \to (S, s)$ 以及带基点概形的
一个交换图
$$
\xymatrix{
(X, x) \ar[d] & (X', x') \ar[l]^g \ar[d] \\
(S, s) & (\Spec(\mathcal{O}_{S', s'}), s') \ar[l]
}
$$
使得 $X' \to X \times_S \Spec(\mathcal{O}_{S', s'})$ 是 \'etale 的，
$\kappa(x) = \kappa(x')$，概形 $X'$ 是仿射的，并且
$\Gamma(X', g^*\mathcal{F})$ 是一个自由的
$\mathcal{O}_{S', s'}$-模。
\end{lemma}

\begin{proof}
（与引理 \ref{flat-lemma-finite-type-flat-at-point-free} 唯一的区别在于，
这里不假设 $f$ 有限表示。）该问题在 $X$ 与 $S$ 上都是局部的。
因此可以假设 $X$ 与 $S$ 都是仿射的，例如
$X = \Spec(B)$ 且 $S = \Spec(A)$。由于 $B$ 是有限型 $A$-代数，
可以找到一个满射 $A[x_1, \ldots, x_n] \to B$。换言之，可以选取
一个闭浸入 $i : X \to \mathbf{A}^n_S$。置 $t = i(x)$ 且
$\mathcal{G} = i_*\mathcal{F}$。注意，作为 $\mathcal{O}_{S, s}$-模，
$\mathcal{G}_t \cong \mathcal{F}_x$。故 $\mathcal{G}$ 在 $t$ 处于
$S$ 上平坦。将引理 \ref{flat-lemma-finite-type-flat-at-point-free}
应用于态射 $\mathbf{A}^n_S \to S$、点 $t$ 以及层 $\mathcal{G}$。
于是可以找到一个初等 \'etale 邻域 $(S', s') \to (S, s)$ 以及带基点
概形的一个交换图
$$
\xymatrix{
(\mathbf{A}^n_S, t) \ar[d] & (Y, y) \ar[l]^h \ar[d] \\
(S, s) & (\Spec(\mathcal{O}_{S', s'}), s') \ar[l]
}
$$
使得 $Y \to \mathbf{A}^n_{\mathcal{O}_{S', s'}}$ 是 \'etale 的，
$\kappa(t) = \kappa(y)$，概形 $Y$ 是仿射的，并且
$\Gamma(Y, h^*\mathcal{G})$ 是一个投射的
$\mathcal{O}_{S', s'}$-模。于是，$Y$ 的闭子概形
$X' = Y \times_{\mathbf{A}^n_S} X$ 给出原问题的一个解。
\end{proof}

\begin{lemma}
\label{flat-lemma-finite-type-flat-along-fibre-free}
设 $f : X \to S$ 是一个概形态射。设 $\mathcal{F}$ 是 $X$ 上的一个
拟凝聚层。设 $s \in S$。假设
\begin{enumerate}
\item $f$ 有限表示；
\item $\mathcal{F}$ 为有限型；
\item $\mathcal{F}$ 在 $X_s$ 的所有点处于 $S$ 上平坦。
\end{enumerate}
则存在一个初等 \'etale 邻域 $(S', s') \to (S, s)$ 以及概形的一个交换图
$$
\xymatrix{
X \ar[d] & X' \ar[l]^g \ar[d] \\
S & \Spec(\mathcal{O}_{S', s'}) \ar[l]
}
$$
使得 $X' \to X \times_S \Spec(\mathcal{O}_{S', s'})$ 是 \'etale 的，
$X_s = g((X')_{s'})$，概形 $X'$ 在 $\mathcal{O}_{S', s'}$ 上仿射且
有限表示，层 $g^*\mathcal{F}$ 在 $\mathcal{O}_{X'}$ 上有限表示，并且
$\Gamma(X', g^*\mathcal{F})$ 是一个自由的
$\mathcal{O}_{S', s'}$-模。
\end{lemma}

\begin{proof}
对于每一点 $x \in X_s$，可以应用
引理 \ref{flat-lemma-finite-type-flat-at-point-free} 找到一个初等 \'etale 邻域
$(S_x , s_x) \to (S, s)$ 以及一个交换图
$$
\xymatrix{
(X, x) \ar[d] & (Y_x, y_x) \ar[l]^{g_x} \ar[d] \\
(S, s) & (\Spec(\mathcal{O}_{S_x, s_x}), s_x) \ar[l]
}
$$
使得 $Y_x \to X \times_S \Spec(\mathcal{O}_{S_x, s_x})$ 是 \'etale 的，
$\kappa(x) = \kappa(y_x)$，概形 $Y_x$ 在 $\mathcal{O}_{S_x, s_x}$ 上
仿射且有限表示，层 $g_x^*\mathcal{F}$ 在 $\mathcal{O}_{Y_x}$ 上有限表示，
并且 $\Gamma(Y_x, g_x^*\mathcal{F})$ 是一个自由的
$\mathcal{O}_{S_x, s_x}$-模。特别地，$g_x((Y_x)_{s_x})$ 是 $X_s$ 中
$x$ 的一个开邻域。由于 $X_s$ 拟紧，可以找到有限多个点
$x_1, \ldots, x_n \in X_s$，使 $X_s$ 是各
$g_{x_i}((Y_{x_i})_{s_{x_i}})$ 的并。选取一个支配每个邻域
$(S_{x_i}, s_{x_i})$ 的初等 \'etale 邻域
$(S' , s') \to (S, s)$，参见《态射进阶》，
引理 \ref{more-morphisms-lemma-elementary-etale-neighbourhoods}。置
$$
X' = \coprod Y_{x_i} \times_{\Spec(\mathcal{O}_{S_{x_i}, s_{x_i}})}
\Spec(\mathcal{O}_{S', s'})
$$
并赋予它显然的态射 $g : X' \to X$。由构造，
$X_s = g(X'_{s'})$，并且
$$
\Gamma(X', g^*\mathcal{F})
=
\bigoplus \Gamma(Y_{x_i}, g_{x_i}^*\mathcal{F})
\otimes_{\mathcal{O}_{S_{x_i}, s_{x_i}}}
\mathcal{O}_{S', s'}.
$$
这是自由模经基变换所得模的直和，因而是一个自由的
$\mathcal{O}_{S', s'}$-模。略去若干次要细节。
\end{proof}

\begin{lemma}
\label{flat-lemma-finite-type-flat-along-fibre-free-variant}
设 $f : X \to S$ 是一个概形态射。设 $\mathcal{F}$ 是 $X$ 上的一个
拟凝聚层。设 $s \in S$。假设
\begin{enumerate}
\item $f$ 有限型；
\item $\mathcal{F}$ 为有限型；
\item $\mathcal{F}$ 在 $X_s$ 的所有点处于 $S$ 上平坦。
\end{enumerate}
则存在一个初等 \'etale 邻域 $(S', s') \to (S, s)$ 以及概形的一个交换图
$$
\xymatrix{
X \ar[d] & X' \ar[l]^g \ar[d] \\
S & \Spec(\mathcal{O}_{S', s'}) \ar[l]
}
$$
使得 $X' \to X \times_S \Spec(\mathcal{O}_{S', s'})$ 是 \'etale 的，
$X_s = g((X')_{s'})$，概形 $X'$ 是仿射的，并且
$\Gamma(X', g^*\mathcal{F})$ 是一个自由的
$\mathcal{O}_{S', s'}$-模。
\end{lemma}

\begin{proof}
（与引理 \ref{flat-lemma-finite-type-flat-along-fibre-free} 唯一的区别在于，
这里不假设 $f$ 有限表示。）对于每一点 $x \in X_s$，可以应用
引理 \ref{flat-lemma-finite-type-flat-at-point-free-variant}
找到一个初等 \'etale 邻域 $(S_x , s_x) \to (S, s)$ 以及一个交换图
$$
\xymatrix{
(X, x) \ar[d] & (Y_x, y_x) \ar[l]^{g_x} \ar[d] \\
(S, s) & (\Spec(\mathcal{O}_{S_x, s_x}), s_x) \ar[l]
}
$$
使得 $Y_x \to X \times_S \Spec(\mathcal{O}_{S_x, s_x})$ 是 \'etale 的，
$\kappa(x) = \kappa(y_x)$，概形 $Y_x$ 是仿射的，并且
$\Gamma(Y_x, g_x^*\mathcal{F})$ 是一个自由的
$\mathcal{O}_{S_x, s_x}$-模。特别地，$g_x((Y_x)_{s_x})$ 是 $X_s$ 中
$x$ 的一个开邻域。由于 $X_s$ 拟紧，可以找到有限多个点
$x_1, \ldots, x_n \in X_s$，使 $X_s$ 是各
$g_{x_i}((Y_{x_i})_{s_{x_i}})$ 的并。选取一个支配每个邻域
$(S_{x_i}, s_{x_i})$ 的初等 \'etale 邻域
$(S' , s') \to (S, s)$，参见《态射进阶》，
引理 \ref{more-morphisms-lemma-elementary-etale-neighbourhoods}。置
$$
X' = \coprod Y_{x_i} \times_{\Spec(\mathcal{O}_{S_{x_i}, s_{x_i}})}
\Spec(\mathcal{O}_{S', s'})
$$
并赋予它显然的态射 $g : X' \to X$。由构造，
$X_s = g(X'_{s'})$，并且
$$
\Gamma(X', g^*\mathcal{F})
=
\bigoplus \Gamma(Y_{x_i}, g_{x_i}^*\mathcal{F})
\otimes_{\mathcal{O}_{S_{x_i}, s_{x_i}}}
\mathcal{O}_{S', s'}.
$$
这是自由模经基变换所得模的直和，因而是一个自由的
$\mathcal{O}_{S', s'}$-模。
\end{proof}




\section{有限型平坦模，第二部分}
\label{flat-section-finite-type-flat-II}

\noindent
我们将需要以下引理。

\begin{lemma}
\label{flat-lemma-weak-bourbaki-pre-pre}
设 $R \to S$ 是一个有限表示的环映射。设 $N$ 是一个有限表示的
$S$-模。设 $\mathfrak q \subset S$ 是一个位于
$\mathfrak p \subset R$ 之上的素理想。置
$\overline{S} = S \otimes_R \kappa(\mathfrak p)$，
$\overline{\mathfrak q} = \mathfrak q \overline{S}$，以及
$\overline{N} = N \otimes_R \kappa(\mathfrak p)$。则可以找到
$g \in S$，$g \not \in \mathfrak q$，使得
$\overline{g} \in \mathfrak r$ 对所有满足
$\mathfrak r \in \text{Ass}_{\overline{S}}(\overline{N})$
且 $\mathfrak r \not \subset \overline{\mathfrak q}$ 的素理想成立。
\end{lemma}

\begin{proof}
事实上，若 $\text{Ass}_{\overline{S}}(\overline{N}) =
\{\mathfrak r_1, \ldots, \mathfrak r_n\}$（有限性来自《代数》，
引理 \ref{algebra-lemma-finite-ass}），则重新编号后可以假设
$$
\mathfrak r_1 \subset \overline{\mathfrak q},
\ldots,
\mathfrak r_r \subset \overline{\mathfrak q}, \quad
\mathfrak r_{r + 1} \not \subset \overline{\mathfrak q},
\ldots,
\mathfrak r_n \not \subset \overline{\mathfrak q}
$$
由于 $\overline{\mathfrak q}$ 是素理想，乘积
$\mathfrak r_{r + 1} \ldots \mathfrak r_n$ 不包含于
$\overline{\mathfrak q}$。因此，可以选取 $\overline{S}$ 中的一个元素
$a$，它属于 $\mathfrak r_{r + 1}, \ldots, \mathfrak r_n$，但不属于
$\overline{\mathfrak q}$。若存在映到 $a$ 的 $g \in S$，则这个 $g$
可用。一般而言，可以找到一个非零元素
$\lambda \in \kappa(\mathfrak p)$，使得 $\lambda a$ 是某个
$g \in S$ 的像。
\end{proof}

\noindent
以下引理在下文有一个稍强的变体，即引理 \ref{flat-lemma-weak-bourbaki}。

\begin{lemma}
\label{flat-lemma-weak-bourbaki-pre}
设 $R \to S$ 是一个有限表示的环映射。设 $N$ 是一个有限表示的
$S$-模，并且作为 $R$-模平坦。设 $M$ 是一个 $R$-模。
设 $\mathfrak q$ 是 $S$ 中位于 $\mathfrak p \subset R$ 之上的一个
素理想。则
$$
\mathfrak q \in \text{WeakAss}_S(M \otimes_R N)
\Leftrightarrow
\Big(
\mathfrak p \in \text{WeakAss}_R(M)
\text{ 且 }
\overline{\mathfrak q} \in \text{Ass}_{\overline{S}}(\overline{N})
\Big)
$$
这里 $\overline{S} = S \otimes_R \kappa(\mathfrak p)$，
$\overline{\mathfrak q} = \mathfrak q \overline{S}$，并且
$\overline{N} = N \otimes_R \kappa(\mathfrak p)$.
\end{lemma}

\begin{proof}
按引理 \ref{flat-lemma-weak-bourbaki-pre-pre} 选取 $g \in S$。
将命题 \ref{flat-proposition-finite-presentation-flat-at-point}
应用于概形态射 $\Spec(S_g) \to \Spec(R)$、与 $N_g$ 对应的拟凝聚模，
以及分别对应于素理想 $\mathfrak qS_g$ 和 $\mathfrak p$ 的点。
翻译成代数语言，我们得到环的一个交换图
$$
\xymatrix{
S \ar[r] & S_g \ar[r] & S' \\
& R \ar[lu] \ar[u] \ar[r] & R' \ar[u]
}
\quad\quad
\xymatrix{
\mathfrak q \ar@{-}[r] \ar@{-}[rd] &
\mathfrak qS_g \ar@{-}[d] \ar@{-}[r] & \mathfrak q' \ar@{-}[d] \\
& \mathfrak p \ar@{-}[r] & \mathfrak p'
}
$$
并带有图中所示的素理想；水平箭头都是 \'etale 的，并且
$N \otimes_S S'$ 作为 $R'$-模是投射的。置
$N' = N \otimes_S S'$, $M' = M \otimes_R R'$,
% P05-EN-CH38-0039: Frozen English tensors S' over R' with kappa(q'),
% although q' is an S'-prime and the next displayed identity uses kappa(p').
$\overline{S}' = S' \otimes_{R'} \kappa(\mathfrak q')$,
$\overline{\mathfrak q}' = \mathfrak q' \overline{S}'$,
以及
$$
\overline{N}' = N' \otimes_{R'} \kappa(\mathfrak p') =
\overline{N} \otimes_{\overline{S}} \overline{S}'
$$
由引理 \ref{flat-lemma-etale-weak-assassin-up-down}，有
\begin{align*}
\text{WeakAss}_{S'}(M' \otimes_{R'} N') & =
(\Spec(S') \to \Spec(S))^{-1}\text{WeakAss}_S(M \otimes_R N) \\
\text{WeakAss}_{R'}(M') & =
(\Spec(R') \to \Spec(R))^{-1}\text{WeakAss}_R(M) \\
\text{Ass}_{\overline{S}'}(\overline{N}') & =
(\Spec(\overline{S}') \to \Spec(\overline{S}))^{-1}
\text{Ass}_{\overline{S}}(\overline{N})
\end{align*}
对 $\overline{N}$ 和 $\overline{N}'$ 应用《代数》，
引理 \ref{algebra-lemma-ass-weakly-ass}。特别地，有
\begin{align*}
\mathfrak q \in \text{WeakAss}_S(M \otimes_R N)
& \Leftrightarrow
\mathfrak q' \in \text{WeakAss}_{S'}(M' \otimes_{R'} N') \\
\mathfrak p \in \text{WeakAss}_R(M)
& \Leftrightarrow
\mathfrak p' \in \text{WeakAss}_{R'}(M') \\
\overline{\mathfrak q} \in \text{Ass}_{\overline{S}}(\overline{N})
& \Leftrightarrow
\overline{\mathfrak q}' \in \text{WeakAss}_{\overline{S}'}(\overline{N}')
\end{align*}
对 $g$ 的审慎选择以及上面关于
$\text{Ass}_{\overline{S}'}(\overline{N}')$ 的公式表明
\begin{equation}
\label{flat-equation-key-observation}
\text{若 }\mathfrak r' \in \text{Ass}_{\overline{S}'}(\overline{N}')
\text{ 位于 }\mathfrak r \subset \overline{S}\text{ 上方，则 }
\mathfrak r \subset \overline{\mathfrak q}
\end{equation}
这将是证明后面部分的一个关键观察。下文将不加说明地使用
《代数》，引理 \ref{algebra-lemma-weakly-ass-local} 给出的弱伴随素理想
刻画。

\medskip\noindent
假设
$\overline{\mathfrak q} \not \in \text{Ass}_{\overline{S}}(\overline{N})$。
那么
$\overline{\mathfrak q}' \not \in \text{Ass}_{\overline{S}'}(\overline{N}')$。
由《代数》，引理 \ref{algebra-lemma-ass-zero-divisors}、
\ref{algebra-lemma-finite-ass} 和
\ref{algebra-lemma-silly}
，存在一个元素 $\overline{a}' \in \overline{\mathfrak q}'$，它不是
$\overline{N}'$ 上的零因子。对某个非零
$\lambda \in \kappa(\mathfrak p)$ 用 $\lambda \overline{a}'$ 替换
$\overline{a}'$ 后，可以找到映到 $\overline{a}'$ 的
$a' \in \mathfrak q'$。由
引理 \ref{flat-lemma-invert-universally-injective}，映射
$a' : N'_{\mathfrak p'} \to N'_{\mathfrak p'}$ 是
$R'_{\mathfrak p'}$-普遍单射的。特别地，映射
$a' : M' \otimes_{R'} N' \to M' \otimes_{R'} N'$ 在
$\mathfrak p'$ 处局部化后是单射，因而在 $\mathfrak q'$ 处局部化后
也是单射。显然，这蕴含
$\mathfrak q' \not \in \text{WeakAss}_{S'}(M' \otimes_{R'} N')$。
我们得出：$\mathfrak q \in \text{WeakAss}_S(M \otimes_R N)$ 蕴含
$\overline{\mathfrak q} \in \text{Ass}_{\overline{S}}(\overline{N})$。

\medskip\noindent
假设 $\mathfrak q \in \text{WeakAss}_S(M \otimes_R N)$。我们要证明
% P05-EN-CH38-0040: Frozen English writes WeakAss_S(M), although p is an R-prime.
$\mathfrak p \in \text{WeakAss}_S(M)$。取元素
$z \in M \otimes_R N$，使得 $\mathfrak q$ 是包含
$J = \text{Ann}_S(z)$ 的极小素理想。设 $f_i \in \mathfrak p$，
$i \in I$ 是理想 $\mathfrak p$ 的一组生成元。由于 $\mathfrak q$
位于 $\mathfrak p$ 之上，对于每个 $i$，可以选取 $n_i \geq 1$ 以及
$g_i \in S$，$g_i \not \in \mathfrak q$，使得
$g_i f_i^{n_i} \in J$，即 $g_i f_i^{n_i} z = 0$。
设 $z' \in (M' \otimes_{R'} N')_{\mathfrak p'}$ 为 $z$ 的像。
注意，$z'$ 非零，因为 $z$ 在 $(M \otimes_R N)_\mathfrak q$ 中的像非零，
并且 $S_\mathfrak q \to S'_{\mathfrak q'}$ 忠实平坦。我们断言
$f_i^{n_i} z' = 0$。

\medskip\noindent
断言的证明：设 $g'_i \in S'$ 为 $g_i$ 的像。由关键观察
(\ref{flat-equation-key-observation})，可见像
$\overline{g}'_i \in \overline{S}'$ 不属于 $\mathfrak r'$，其中
% P05-EN-CH38-0041: Frozen English uses overline N instead of overline N' here.
$\mathfrak r' \in \text{Ass}_{\overline{S}'}(\overline{N})$。
于是由引理 \ref{flat-lemma-invert-universally-injective}，可见
$g'_i : N'_{\mathfrak p'} \to N'_{\mathfrak p'}$ 是
$R'_{\mathfrak p'}$-普遍单射的。特别地，映射
$g'_i : M' \otimes_{R'} N' \to M' \otimes_{R'} N'$ 在
$\mathfrak p'$ 处局部化后是单射。由于 $g_i f_i^{n_i} z' = 0$，
断言得证。

\medskip\noindent
我们的断言表明，$z'$ 在 $R'_{\mathfrak p'}$ 中的零化理想包含各元素
$f_i^{n_i}$。由于 $R \to R'$ 是 \'etale 的，由《代数》，
引理 \ref{algebra-lemma-etale-at-prime}，有
$\mathfrak p'R'_{\mathfrak p'} = \mathfrak pR'_{\mathfrak p'}$
。因此，$z'$ 在 $R'_{\mathfrak p'}$ 中的零化理想的根等于
% P05-EN-CH38-0042: Frozen English localizes R at the R'-prime p'; likely R' was intended.
$\mathfrak p' R_{\mathfrak p'}$（这里使用了 $z'$ 非零）。另一方面，
$$
z' \in (M' \otimes_{R'} N')_{\mathfrak p'} =
M'_{\mathfrak p'} \otimes_{R'_{\mathfrak p'}} N'_{\mathfrak p'}
$$
模 $N'_{\mathfrak p'}$ 在局部环 $R'_{\mathfrak p'}$ 上投射，因而自由
（《代数》，定理 \ref{algebra-theorem-projective-free-over-local-ring}）。
因此，可以找到一个有限自由直和项 $F' \subset N'_{\mathfrak p'}$，
使得 $z' \in M'_{\mathfrak p'} \otimes_{R'_{\mathfrak p'}} F'$。
若 $F'$ 的秩为 $n$，则可推出
$\mathfrak p' R'_{\mathfrak p'} \in
\text{WeakAss}_{R'_{\mathfrak p'}}({M'_{\mathfrak p'}}^{\oplus n})$.
这蕴含
$\mathfrak p'R'_{\mathfrak p'} \in \text{WeakAss}(M'_{\mathfrak p'})$
，例如可由《代数》，引理 \ref{algebra-lemma-weakly-ass} 得到。
于是 $\mathfrak p' \in \text{WeakAss}_{R'}(M')$，进而得到
$\mathfrak p \in \text{WeakAss}_R(M)$。这就完成了引理中等价关系的
蕴含“$\Rightarrow$”的证明。

\medskip\noindent
假设 $\mathfrak p \in \text{WeakAss}_R(M)$ 且
$\overline{\mathfrak q} \in \text{Ass}_{\overline{S}}(\overline{N})$。
我们要证明 $\mathfrak q$ 是 $M \otimes_R N$ 的一个弱伴随素理想。
注意，$\overline{\mathfrak q}'$ 是
$\text{Ass}_{\overline{S}'}(\overline{N}')$ 的一个极大元。
这来自 (\ref{flat-equation-key-observation})，以及如下事实：
$\overline{S}'$ 中位于 $\overline{\mathfrak q}$ 之上的各素理想之间
不存在包含关系（因为 \'etale 态射的纤维是离散的，参见《态射》，
引理 \ref{morphisms-lemma-etale-over-field}）。因此，用
$R', S', \mathfrak p', \mathfrak q', M', N'$ 分别替换
$R, S, \mathfrak p, \mathfrak q, M, N$ 后，除引理的假设外，还可假设
\begin{enumerate}
\item $\mathfrak p \in \text{WeakAss}_R(M)$；
\item $\overline{\mathfrak q} \in \text{Ass}_{\overline{S}}(\overline{N})$；
\item $N$ 作为 $R$-模是投射的；
\item $\overline{\mathfrak q}$ 在
$\text{Ass}_{\overline{S}}(\overline{N})$ 中是极大的。
\end{enumerate}
还有一步约化：可以用 $R, S, M, N$ 在 $\mathfrak p$ 处的局部化分别
替换它们。这又给出一个条件，即
\begin{enumerate}
\item[(5)] $R$ 是极大理想为 $\mathfrak p$ 的局部环。
\end{enumerate}
最后我们将证明 (1) -- (5) 蕴含
$\mathfrak q \in \text{WeakAss}(M \otimes_R N)$.

\medskip\noindent
由于 $R$ 是局部的且 $\mathfrak p \in \text{WeakAss}_R(M)$，可以选取
$y \in M$，使其零化理想 $I$ 的根等于 $\mathfrak p$。
对某些 $\overline{g}_i \in \overline{S}$，写成
$\overline{\mathfrak q} = (\overline{g}_1, \ldots, \overline{g}_n)$。
选取映到 $\overline{g}_i$ 的 $g_i \in S$。那么
$\mathfrak q = \mathfrak pS + g_1S + \ldots + g_nS$。考虑映射
$$
\Psi : N/IN \longrightarrow (N/IN)^{\oplus n}, \quad
z \longmapsto (g_1z, \ldots, g_nz).
$$
这是投射 $R/I$-模之间的一个同态。由于 $\mathfrak p/I$ 局部幂零，
局部环 $R/I$ 是自伴随的（《代数进阶》，
定义 \ref{more-algebra-definition-auto-ass}）。映射
$\Psi \otimes \kappa(\mathfrak p)$ 不是单射，因为
$\overline{\mathfrak q} \in \text{Ass}_{\overline{S}}(\overline{N})$。
因此，《代数进阶》，引理 \ref{more-algebra-lemma-P-fPD-zero}
蕴含 $\Psi$ 不是单射。在 $\Psi$ 的核中选取一个非零元
$z \in N/IN$。由构造，零化理想 $J = \text{Ann}_S(z)$ 包含 $IS$
以及各 $g_i$。因此，$\sqrt{J} \subset S$ 包含 $\mathfrak q$。
设 $\mathfrak s \subset S$ 是包含 $J$ 的一个极小素理想。那么
$\mathfrak q \subset \mathfrak s$，$\mathfrak s$ 位于 $\mathfrak p$
之上，并且 $\mathfrak s \in \text{WeakAss}_S(N/IN)$。
最后一个事实来自弱伴随素理想的定义。将该引理已证明的
“$\Rightarrow$”部分应用于环映射 $R \to S$ 以及模 $R/I$ 和 $N$，
得到
$\overline{\mathfrak s} \in \text{Ass}_{\overline{S}}(\overline{N})$.
由于 $\overline{\mathfrak q} \subset \overline{\mathfrak s}$，
$\overline{\mathfrak q}$ 的极大性（见上面的条件 (4)）蕴含
$\overline{\mathfrak q} = \overline{\mathfrak s}$。这表明
$\mathfrak q = \mathfrak s$，从而得到所需结论。
\end{proof}

\begin{lemma}
\label{flat-lemma-bourbaki-finite-type-general-base-at-point}
设 $S$ 是一个概形。设 $f : X \to S$ 局部有限型。设 $x \in X$
的像为 $s \in S$。设 $\mathcal{F}$ 是 $X$ 上的一个有限型拟凝聚层，
$\mathcal{G}$ 是 $S$ 上的一个拟凝聚层。若 $\mathcal{F}$ 在 $x$ 处于
$S$ 上平坦，则
$$
x \in \text{WeakAss}_X(\mathcal{F} \otimes_{\mathcal{O}_X} f^*\mathcal{G})
\Leftrightarrow
s \in \text{WeakAss}_S(\mathcal{G})
\text{ 且 }
x \in \text{Ass}_{X_s}(\mathcal{F}_s).
$$
\end{lemma}

\begin{proof}
本段将问题约化到 $f$ 有限表示的情形。该问题在 $X$ 与 $S$ 上都是
局部的，故可假设 $X$ 与 $S$ 都是仿射的。写成
$X = \Spec(B)$、$S = \Spec(A)$，并写
$B = A[x_1, \ldots, x_n]/I$。换言之，我们得到 $S$ 上的一个闭浸入
$i : X \to \mathbf{A}^n_S$。记 $t = i(x) \in \mathbf{A}^n_S$。
注意，$i_*\mathcal{F}$ 是 $\mathbf{A}^n_S$ 上的一个有限型拟凝聚层，
它在 $t$ 处于 $S$ 上平坦；并且
$$
i_*(\mathcal{F} \otimes_{\mathcal{O}_X} f^*\mathcal{G}) =
i_*\mathcal{F} \otimes_{\mathcal{O}_{\mathbf{A}^n_S}} p^*\mathcal{G}
$$
其中 $p : \mathbf{A}^n_S \to S$ 是投影。注意，$t$ 是
$i_*(\mathcal{F} \otimes_{\mathcal{O}_X} f^*\mathcal{G})$ 的弱伴随点，
当且仅当 $x$ 是
$\mathcal{F} \otimes_{\mathcal{O}_X} f^*\mathcal{G}$ 的弱伴随点，参见
《除子》，引理 \ref{divisors-lemma-weakly-associated-finite}。
类似地，$x \in \text{Ass}_{X_s}(\mathcal{F}_s)$ 当且仅当
$t \in \text{Ass}_{\mathbf{A}^n_s}((i_*\mathcal{F})_s)$（参见
《代数》，引理 \ref{algebra-lemma-ass-quotient-ring}）。
因此，只需在 $X = \mathbf{A}^n_S$ 的情形下证明该引理。
于是可以假设 $X \to S$ 有限表示。

\medskip\noindent
本段将问题约化到 $\mathcal{F}$ 有限表示且在 $S$ 上平坦的情形。
按命题 \ref{flat-proposition-finite-type-flat-at-point} 选取一个初等
\'etale 邻域 $e : (S', s') \to (S, s)$ 以及一个开集
$V \subset X \times_S \Spec(\mathcal{O}_{S', s'})$。
设 $x' \in X' = X \times_S S'$ 是映到 $x$ 与 $s'$ 的唯一一点。
那么只需对 $X' \to S'$、$x'$、$s'$、
$(X' \to X)^*\mathcal{F}$ 以及 $e^*\mathcal{G}$ 证明该断言，参见
引理 \ref{flat-lemma-etale-weak-assassin-up-down}。
% P05-EN-CH38-0043: Frozen English omits "be" in "Let v in V the unique point".
设 $v \in V$ 是映到 $x'$ 的唯一一点，并设
$s' \in \Spec(\mathcal{O}_{S', s'})$ 为闭点。那么
$\mathcal{O}_{V, v} = \mathcal{O}_{X', x'}$，且
$\mathcal{O}_{\Spec(\mathcal{O}_{S', s'}), s'} =
\mathcal{O}_{S', s'}$；$\mathcal{F}$ 与 $\mathcal{G}$ 的拉回的茎也同样。
此外，$V_{s'} \subset X'_{s'}$ 是一个开子概形。由于一个点为弱伴随点
这一条件只依赖于层的茎，可以将
$X' \to S'$、$x'$、$s'$、$(X' \to X)^*\mathcal{F}$ 以及 $e^*\mathcal{G}$
替换为
$V \to \Spec(\mathcal{O}_{S', s'})$, $v$, $s'$, $(V \to X)^*\mathcal{F}$,
以及 $(\Spec(\mathcal{O}_{S', s'}) \to S)^*\mathcal{G}$。
因此可以假设 $f$ 有限表示，并且 $\mathcal{F}$ 有限表示且在 $S$ 上平坦。

\medskip\noindent
假设 $f$ 有限表示且 $\mathcal{F}$ 有限表示并在 $S$ 上平坦。
将 $X$ 与 $S$ 分别缩小为 $x$ 与 $s$ 的仿射邻域后，此情形由
引理 \ref{flat-lemma-weak-bourbaki-pre} 处理。
\end{proof}

\begin{lemma}
\label{flat-lemma-weak-bourbaki}
设 $R \to S$ 是一个本质有限型的环映射。设 $N$ 是某个有限 $S$-模的
一个局部化，并且在 $R$ 上平坦。设 $M$ 是一个 $R$-模。则
$$
\text{WeakAss}_S(M \otimes_R N)
=
\bigcup\nolimits_{\mathfrak p \in \text{WeakAss}_R(M)}
\text{Ass}_{S \otimes_R \kappa(\mathfrak p)}(N \otimes_R \kappa(\mathfrak p))
$$
\end{lemma}

\begin{proof}
该引理是引理 \ref{flat-lemma-bourbaki-finite-type-general-base-at-point}
用代数语言作出的翻译。略去翻译的细节。
\end{proof}

\begin{lemma}
\label{flat-lemma-bourbaki-finite-type-general-base}
设 $f : X \to S$ 是一个局部有限型的态射。设 $\mathcal{F}$ 是 $X$
上的一个有限型拟凝聚层，并且在 $S$ 上平坦。设 $\mathcal{G}$ 是 $S$
上的一个拟凝聚层。则有
$$
\text{WeakAss}_X(\mathcal{F} \otimes_{\mathcal{O}_X} f^*\mathcal{G}) =
\bigcup\nolimits_{s \in \text{WeakAss}_S(\mathcal{G})}
\text{Ass}_{X_s}(\mathcal{F}_s)
$$
\end{lemma}

\begin{proof}
这是引理 \ref{flat-lemma-bourbaki-finite-type-general-base-at-point}
的直接推论。
\end{proof}

\begin{theorem}
\label{flat-theorem-finite-type-flat}
\begin{slogan}
平坦轨迹是开的（非 Noether 情形）。
\end{slogan}
设 $f : X \to S$ 是一个概形态射。设 $\mathcal{F}$ 是一个拟凝聚
$\mathcal{O}_X$-模。假设
\begin{enumerate}
\item $X \to S$ 局部有限表示；
\item $\mathcal{F}$ 是有限型 $\mathcal{O}_X$-模；
\item $S$ 的弱伴随点集在 $S$ 中局部有限。
\end{enumerate}
则 $U = \{x \in X \mid \mathcal{F}\text{ 在 }x\text{ 处于 }S\text{ 上平坦}\}$
在 $X$ 中是开的，并且 $\mathcal{F}|_U$ 是有限表示的
$\mathcal{O}_U$-模且在 $S$ 上平坦。
\end{theorem}

\begin{proof}
设 $x \in X$，使得 $\mathcal{F}$ 在 $x$ 处于 $S$ 上平坦。
我们必须找到 $x$ 的一个开邻域，使 $\mathcal{F}$ 在该邻域上的限制
是一个有限表示的 $S$-平坦模。该问题在 $X$ 与 $S$ 上都是局部的，
故可假设 $X$ 与 $S$ 都是仿射的。由
引理 \ref{flat-lemma-finite-type-flat-at-point-local}，$\mathcal{F}_x$ 是
一个有限表示的 $\mathcal{O}_{X, x}$-模。因此，由
《代数》，引理 \ref{algebra-lemma-construct-fp-module-from-stalk}，
存在一个有限表示的 $\mathcal{O}_X$-模 $\mathcal{F}'$ 以及一个映射
$\varphi : \mathcal{F}' \to \mathcal{F}$，它诱导同构
$\varphi_x : \mathcal{F}'_x \to \mathcal{F}_x$。特别地，可见
$\mathcal{F}'$ 在 $x$ 处于 $S$ 上平坦。因此，由平坦性的开放性
（《态射进阶》，定理 \ref{more-morphisms-theorem-openness-flatness}），
缩小 $X$ 后，可以假设 $\mathcal{F}'$ 在 $S$ 上平坦。
由于 $\mathcal{F}$ 为有限型，再次缩小 $X$ 后可以假设 $\varphi$ 是
满射，参见《模》，引理 \ref{modules-lemma-finite-type-surjective-on-stalk}；
也可以使用 Nakayama 引理（《代数》，引理 \ref{algebra-lemma-NAK}）。
由引理 \ref{flat-lemma-bourbaki-finite-type-general-base}，有
$$
\text{WeakAss}_X(\mathcal{F}') \subset
\bigcup\nolimits_{s \in \text{WeakAss}(S)} \text{Ass}_{X_s}(\mathcal{F}'_s)
$$
由假设，$\text{WeakAss}(S)$ 是有限的；又由《除子》，
引理 \ref{divisors-lemma-finite-ass}，
$\text{Ass}_{X_s}(\mathcal{F}'_s)$ 是有限的，故可推出
$\text{WeakAss}_X(\mathcal{F}')$ 是有限的。利用《代数》，
引理 \ref{algebra-lemma-silly}，再次缩小 $X$ 后，可以假设
$\text{WeakAss}_X(\mathcal{F}')$ 包含于 $x$ 的一般化点集。
现在考虑 $\mathcal{K} = \Ker(\varphi)$。我们有
$\text{WeakAss}_X(\mathcal{K}) \subset \text{WeakAss}_X(\mathcal{F}')$
（由《除子》，引理 \ref{divisors-lemma-ses-weakly-ass}）。
另一方面，$\varphi_x$ 是同构，而且对于所有 $x' \leadsto x$，
$\varphi_{x'}$ 也是同构。我们推出
$\text{WeakAss}_X(\mathcal{K}) = \emptyset$，从而由《除子》，
引理 \ref{divisors-lemma-weakly-ass-zero}，有 $\mathcal{K} = 0$。
\end{proof}

\begin{lemma}
\label{flat-lemma-finite-type-flat-algebra}
设 $R \to S$ 是一个有限表示的环映射。设 $M$ 是一个有限 $S$-模。
假设 $\text{WeakAss}_R(R)$ 是有限的。则
$$
U = \{\mathfrak q \subset S \mid M_{\mathfrak q}\text{ 在 }R\text{ 上平坦}\}
$$
在 $\Spec(S)$ 中是开的；并且对于每个满足 $D(g) \subset U$ 的
$g \in S$，局部化 $M_g$ 是一个有限表示的 $S_g$-模，且在 $R$ 上平坦。
\end{lemma}

\begin{proof}
这立即来自定理 \ref{flat-theorem-finite-type-flat}。
\end{proof}

\begin{lemma}
\label{flat-lemma-finite-type-flat-X}
设 $f : X \to S$ 是一个局部有限型的概形态射。假设 $S$ 的弱伴随点集
在 $S$ 中局部有限。则 $f$ 平坦的各点 $x \in X$ 构成一个开子概形
$U \subset X$，并且 $U \to S$ 平坦且局部有限表示。
\end{lemma}

\begin{proof}
该问题在 $X$ 与 $S$ 上都是局部的，故可假设 $X$ 与 $S$ 都是仿射的。
那么 $X \to S$ 对应于一个有限型环映射 $A \to B$。选取一个满射
$A[x_1, \ldots, x_n] \to B$，并将 $B$ 视为一个
$A[x_1, \ldots, x_n]$-模。应用
引理 \ref{flat-lemma-finite-type-flat-algebra} 即完成证明。
\end{proof}

\begin{lemma}
\label{flat-lemma-finite-type-flat-over-integral}
设 $f : X \to S$ 是一个局部有限型且平坦的概形态射。若 $S$ 是整的，
则 $f$ 局部有限表示。
\end{lemma}

\begin{proof}
这是引理 \ref{flat-lemma-finite-type-flat-X} 的特例。
\end{proof}

\begin{proposition}
\label{flat-proposition-flat-finite-type-finite-presentation-domain}
设 $R$ 是一个整环。设 $R \to S$ 是一个有限型环映射。
设 $M$ 是一个有限 $S$-模。
\begin{enumerate}
\item 若 $S$ 在 $R$ 上平坦，则 $S$ 是一个有限表示的 $R$-代数。
\item 若 $M$ 作为 $R$-模平坦，则 $M$ 作为 $S$-模有限表示。
\end{enumerate}
\end{proposition}

\begin{proof}
(1) 是引理 \ref{flat-lemma-finite-type-flat-over-integral} 的一个特例。
对于 (2)，选取一个满射 $R[x_1, \ldots, x_n] \to S$。
由引理 \ref{flat-lemma-finite-type-flat-algebra}，可见 $M$ 作为
$R[x_1, \ldots, x_n]$-模有限表示。最后应用《代数》，
引理 \ref{algebra-lemma-finitely-presented-over-subring}。
\end{proof}

\begin{lemma}[定理 \ref{flat-theorem-finite-type-flat} 的有限型版本]
\label{flat-lemma-finite-type-flat}
设 $f : X \to S$ 是一个概形态射。设 $\mathcal{F}$ 是一个拟凝聚
$\mathcal{O}_X$-模。假设
\begin{enumerate}
\item $X \to S$ 局部有限型；
\item $\mathcal{F}$ 是有限型 $\mathcal{O}_X$-模；
\item $S$ 的弱伴随点集在 $S$ 中局部有限。
\end{enumerate}
则 $U = \{x \in X \mid \mathcal{F}\text{ 在 }x\text{ 处于 }S\text{ 上平坦}\}$
在 $X$ 中是开的，并且 $\mathcal{F}|_U$ 在 $S$ 上平坦且局部相对于
$S$ 有限表示（参见《态射进阶》，定义
\ref{more-morphisms-definition-relatively-finitely-presented-sheaf}).
\end{lemma}

\begin{proof}
该问题在 $X$ 与 $S$ 上都是局部的。因此可以假设 $X$ 与 $S$ 都是
仿射的。那么可以选取一个闭浸入 $i : X \to \mathbf{A}^n_S$。
将定理 \ref{flat-theorem-finite-type-flat} 应用于
$X' = \mathbf{A}^n_S \to S$ 以及有限型拟凝聚模
$\mathcal{F}' = i_*\mathcal{F}$，可得
$$
U' = \{x' \in X' \mid \mathcal{F}'\text{ 在 }x'\text{ 处于 }S\text{ 上平坦}\}
$$
在 $X'$ 中是开的，并且 $\mathcal{F}'|_{U'}$ 有限表示。
由于 $\mathcal{F}'$ 在 $X' \setminus i(X)$ 上的限制为零，且对于所有
$x \in X$ 都有 $\mathcal{F}'_{i(x)} \cong \mathcal{F}_x$，可见
$$
U' = i(U) \amalg (X' \setminus i(X))
$$
因此 $U = i^{-1}(U')$ 是开的。此外，显然
$\mathcal{F}'|_{U'} = (i|_U)_*(\mathcal{F}|_U)$.
故由《态射进阶》，引理
\ref{more-morphisms-lemma-relative-finite-presentation} 和
\ref{more-morphisms-lemma-finite-morphism-relative-finite-presentation}.
可推出 $\mathcal{F}|_U$ 相对于 $S$ 有限表示。
\end{proof}

\begin{lemma}
\label{flat-lemma-finite-type-flat-algebra-bis}
设 $R \to S$ 是一个有限型环映射。设 $M$ 是一个有限 $S$-模。
假设 $\text{WeakAss}_R(R)$ 是有限的。则
$$
U = \{\mathfrak q \subset S \mid M_{\mathfrak q}\text{ 在 }R\text{ 上平坦}\}
$$
在 $\Spec(S)$ 中是开的；并且对于每个满足 $D(g) \subset U$ 的
$g \in S$，局部化 $M_g$ 在 $R$ 上平坦，且作为 $S_g$-模相对于
$R$ 有限表示（参见《代数进阶》，定义
\ref{more-algebra-definition-relatively-finitely-presented}).
\end{lemma}

\begin{proof}
这是引理 \ref{flat-lemma-finite-type-flat} 用代数语言作出的翻译。
\end{proof}








\section{相对纯模的例子}
\label{flat-section-examples-pure-modules}

\noindent
在本节中，我们讨论一些结果的例子；它们将为稍后引入的
{\it 相对纯模}概念和 {\it 非纯性}概念提供动机。每个例子都表述为一个
引理。请注意，关于伴随素理想的条件与下列引理中出现的条件相似：
引理 \ref{flat-lemma-base-change-universally-flat}、
\ref{flat-lemma-universally-injective-to-completion},
\ref{flat-lemma-universally-injective-to-completion-flat} 和
\ref{flat-lemma-flat-pure-over-complete-projective}。
另请参见《代数》，引理 \ref{algebra-lemma-compare-relative-assassins}，
其中有相关讨论。

\begin{lemma}
\label{flat-lemma-explain-why-pure}
设 $R$ 是极大理想为 $\mathfrak m$ 的局部环。设 $R \to S$ 是一个环映射。
设 $N$ 是一个 $S$-模。假设
\begin{enumerate}
\item $N$ 作为 $R$-模是投射的；
\item $S/\mathfrak mS$ 是 Noether 环，并且 $N/\mathfrak mN$ 是有限的
$S/\mathfrak mS$-模。
\end{enumerate}
则对于任意素理想 $\mathfrak q \subset S$，若它是
$N \otimes_R \kappa(\mathfrak p)$ 的伴随素理想，其中
$\mathfrak p = R \cap \mathfrak q$，则有
$\mathfrak q + \mathfrak m S \not = S$。
\end{lemma}

\begin{proof}
注意，引理 \ref{flat-lemma-homothety-spectrum} 与
\ref{flat-lemma-invert-universally-injective} 的假设都满足。以下不再另行说明，
直接使用这两个引理的结论。令 $\Sigma \subset S$ 为在 $N/\mathfrak mN$
上不是零因子的元素所成的乘法集。映射
$N \to \Sigma^{-1}N$ 是 $R$-普遍单射的。因此，任何作为
$N \otimes_R \kappa(\mathfrak p)$ 伴随素理想的
$\mathfrak q \subset S$，也是
$\Sigma^{-1}N \otimes_R \kappa(\mathfrak p)$ 的伴随素理想。
显然，这蕴含 $\mathfrak q$ 对应于 $\Sigma^{-1}S$ 的一个素理想。
于是 $\mathfrak q \subset \mathfrak q'$，其中 $\mathfrak q'$ 对应于
$N/\mathfrak mN$ 的一个伴随素理想，结论得证。
\end{proof}

\noindent
以下引理给出这一现象的另一个（稍显无聊的）例子。

\begin{lemma}
\label{flat-lemma-explain-why-pure-complete}
设 $R$ 是一个环。设 $I \subset R$ 是一个理想。设 $R \to S$ 是一个环映射。
设 $N$ 是一个 $S$-模。若 $N$ 是 $I$-进完备的，则对于任意 $R$-模 $M$，
以及任意素理想 $\mathfrak q \subset S$，若它是
$N \otimes_R M$ 的伴随素理想，则有 $\mathfrak q + I S \not = S$。
\end{lemma}

\begin{proof}
令 $S^\wedge$ 表示 $S$ 的 $I$-进完备化。注意，$N$ 是一个
$S^\wedge$-模，因而 $N \otimes_R M$ 也是一个 $S^\wedge$-模。
取 $z \in N \otimes_R M$，使得 $\mathfrak q = \text{Ann}_S(z)$。
由于 $z \not = 0$，可见 $\text{Ann}_{S^\wedge}(z) \not = S^\wedge$，
从而 $\mathfrak q S^\wedge \not = S^\wedge$。因此存在一个极大理想
$\mathfrak m \subset S^\wedge$，使得 $\mathfrak q S^\wedge \subset \mathfrak m$。
由《代数》，引理 \ref{algebra-lemma-radical-completion}，
$IS^\wedge \subset \mathfrak m$，结论得证。
\end{proof}

\noindent
注意，以下引理给出引理 \ref{flat-lemma-explain-why-pure} 的另一种证明，
因为局部环上的投射模是自由模，参见《代数》，定理
\ref{algebra-theorem-projective-free-over-local-ring}。

\begin{lemma}
\label{flat-lemma-explain-why-pure-direct-sum-finite-modules}
设 $R$ 是极大理想为 $\mathfrak m$ 的局部环。设 $R \to S$ 是一个环映射。
设 $N$ 是一个 $S$-模，并假设 $N$ 作为 $R$-模同构于有限 $R$-模的一个直和。
则对于任意 $R$-模 $M$，
以及任意素理想 $\mathfrak q \subset S$，若它是 $N \otimes_R M$ 的伴随素理想，
则有 $\mathfrak q + \mathfrak m S \not = S$。
\end{lemma}

\begin{proof}
写 $N = \bigoplus_{i \in I} M_i$，其中每个 $M_i$ 都是有限的 $R$-模。
设 $M$ 是一个 $R$-模，设 $\mathfrak q \subset S$ 是
$N \otimes_R M$ 的伴随素理想，并且 $\mathfrak q + \mathfrak m S = S$。
取 $z \in N \otimes_R M$，使得 $\mathfrak q = \text{Ann}_S(z)$。
稍微调整直和分解后，可以假设对于某个 $1 \in I$ 有
$z \in M_1 \otimes_R M$。写成
$1 = f + \sum x_j g_j$，其中 $f \in \mathfrak q$、$x_j \in \mathfrak m$、
$g_j \in S$。对于任意 $g \in S$，记 $g'$ 为以下 $R$-线性映射：
$$
M_1 \to N \xrightarrow{g} N \to M_1
$$
其中第一个箭头是包含映射，第二个箭头是乘以 $g$，第三个箭头是投影映射。
由于每个 $x_j \in R$，我们得到等式
$$
f' + \sum x_j g'_j = \text{id}_{M_1} \in \text{End}_R(M_1)
$$
由 Nakayama 引理（《代数》，引理 \ref{algebra-lemma-NAK}），可见 $f'$
是满射；于是由《代数》，引理 \ref{algebra-lemma-fun}，可见 $f'$ 是同构。
特别地，以下映射
$$
M_1 \otimes_R M \to N \otimes_R M \xrightarrow{f} N \otimes_R M
\to M_1 \otimes_R M
$$
是同构。这与 $fz = 0$ 的假设矛盾。
\end{proof}

\begin{lemma}
\label{flat-lemma-explain-why-pure-ML}
设 $R$ 是极大理想为 $\mathfrak m$ 的 Hensel 局部环。设 $R \to S$ 是一个环映射。
设 $N$ 是一个 $S$-模。假设 $N$ 作为 $R$-模是可数生成的且为
Mittag-Leffler 模。则对于任意 $R$-模 $M$ 以及任意素理想
$\mathfrak q \subset S$，若它是 $N \otimes_R M$ 的伴随素理想，则有
$\mathfrak q + \mathfrak m S \not = S$.
\end{lemma}

\begin{proof}
由《代数》，引理 \ref{algebra-lemma-split-ML-henselian}，该引理约化为
引理 \ref{flat-lemma-explain-why-pure-direct-sum-finite-modules}。
\end{proof}

\noindent
设 $f : X \to S$ 是一个概形态射，$\mathcal{F}$ 是 $X$ 上的一个
拟凝聚模。设 $\xi \in \text{Ass}_{X/S}(\mathcal{F})$，并令
$Z = \overline{\{\xi\}}$。图示如下：
$$
\xymatrix{
\xi \ar@{|->}[d] & Z \ar[r] \ar[rd] & X \ar[d]^f \\
f(\xi) & & S
}
$$
注意，$f(Z) \subset \overline{\{f(\xi)\}}$，并且 $f(Z)$ 是闭集
当且仅当等式成立，即 $f(Z) = \overline{\{f(\xi)\}}$。
由引理 \ref{flat-lemma-explain-why-pure} 可知，若 $S$、$X$ 都是仿射的，
纤维 $X_s$ 是 Noether 概形，$\mathcal{F}$ 是有限型的，并且
$\Gamma(X, \mathcal{F})$ 是投射的 $\Gamma(S, \mathcal{O}_S)$-模，
则 $f(Z) = \overline{\{f(\xi)\}}$ 是一个闭子集。
对于下列情形，也有略有不同但类似的陈述：
引理 \ref{flat-lemma-explain-why-pure-complete}、
\ref{flat-lemma-explain-why-pure-direct-sum-finite-modules} 和
\ref{flat-lemma-explain-why-pure-ML}.




\section{非纯性}
\label{flat-section-impure}

\noindent
我们想用 $\text{Ass}_{X/S}(\mathcal{F})$ 中点的特化来形式化
第 \ref{flat-section-examples-pure-modules} 节中举例说明的现象。
我们还希望在 $s \in S$ 的附近局部地工作。为此作出以下定义。

\begin{situation}
\label{flat-situation-pre-pure}
这里 $S$、$X$ 是概形，且 $f : X \to S$ 是有限型态射。
此外，$\mathcal{F}$ 是有限型拟凝聚 $\mathcal{O}_X$-模。
最后，$s$ 是 $S$ 的一个点。
\end{situation}

\noindent
在此情形下，考虑一个态射 $g : T \to S$、一个满足 $g(t) = s$ 的
点 $t \in T$、一个特化 $t' \leadsto t$，以及 $X$ 的基变换中位于
$t'$ 之上的一点 $\xi \in X_T$。图示如下：
\begin{equation}
\label{flat-equation-impurity}
\vcenter{
\xymatrix{
\xi \ar@{|->}[d] & \\
t' \ar@{~>}[r] & t \ar@{|->}[r] & s
}
}
\quad\quad
\vcenter{
\xymatrix{
X_T \ar[d] \ar[r] & X \ar[d] \\
T \ar[r]^g & S
}
}
\end{equation}
此外，记 $\mathcal{F}_T$ 为 $\mathcal{F}$ 到 $X_T$ 的拉回。

\begin{definition}
\label{flat-definition-impurity}
在情形 \ref{flat-situation-pre-pure} 中，如果
图 (\ref{flat-equation-impurity}) 满足
$\xi \in \text{Ass}_{X_T/T}(\mathcal{F}_T)$ 且
$\overline{\{\xi\}} \cap X_t = \emptyset$，则称它定义了
$\mathcal{F}$ 在 $s$ 之上的一个 {\it 非纯性}。
我们将用如下说法表示这一点：“设
$(g : T \to S, t' \leadsto t, \xi)$ 是 $\mathcal{F}$ 在 $s$ 之上的
一个非纯性”。
\end{definition}

\begin{lemma}
\label{flat-lemma-impure-finite-presentation}
在情形 \ref{flat-situation-pre-pure} 中，若存在 $\mathcal{F}$ 在 $s$ 之上的
一个非纯性，则存在一个 $\mathcal{F}$ 在 $s$ 之上的非纯性
$(g : T \to S, t' \leadsto t, \xi)$，使得 $g$ 局部有限表示，且 $t$
是 $g$ 在 $s$ 之上的纤维中的闭点。
\end{lemma}

\begin{proof}
设 $(g : T \to S, t' \leadsto t, \xi)$ 是 $\mathcal{F}$ 在 $s$ 之上的
任意非纯性。将《极限》，引理 \ref{limits-lemma-separate} 应用于
$t \in T$ 与 $Z = \overline{\{\xi\}}$，得到 $t$ 的一个开邻域
$V \subset T$、一个交换图
$$
\xymatrix{
V \ar[d] \ar[r]_a & T' \ar[d]^b \\
T \ar[r]^g & S,
}
$$
以及一个闭子概形 $Z' \subset X_{T'}$，满足
\begin{enumerate}
\item 态射 $b : T' \to S$ 局部有限表示；
\item $Z' \cap X_{a(t)} = \emptyset$；
\item $Z \cap X_V$ 通过态射 $X_V \to X_{T'}$ 映入 $Z'$。
\end{enumerate}
由于 $t'$ 特化到 $t$，可以用 $t$ 的开邻域 $V$ 替换 $T$。
于是有交换图
$$
\xymatrix{
X_T \ar[d] \ar[r] &
X_{T'} \ar[d] \ar[r] &
X \ar[d] \\
T \ar[r]^a & T' \ar[r]^b & S
}
$$
其中 $b \circ a = g$。令 $\xi' \in X_{T'}$ 表示 $\xi$ 的像。由
《除子》，引理 \ref{divisors-lemma-base-change-relative-assassin}，
可见 $\xi' \in \text{Ass}_{X_{T'}/T'}(\mathcal{F}_{T'})$。
此外，按构造，$\overline{\{\xi'\}}$ 的闭包包含于闭子集 Z' 中，
而 $Z'$ 避开纤维 $X_{a(t)}$。因此
$(T' \to S, a(t') \leadsto a(t), \xi')$ 是 $\mathcal{F}$ 在 $s$ 之上的
一个非纯性。

\medskip\noindent
因此可以假设 $g : T \to S$ 局部有限表示。令
$Z = \overline{\{\xi\}}$。由假设 $Z_t = \emptyset$。由《态射进阶》，
引理 \ref{more-morphisms-lemma-empty-generic-fibre}，这意味着对于
$\overline{\{t\}}$ 的某个开子集中的 $t''$，有 $Z_{t''} = \emptyset$。
由于 $T \to S$ 在 $s$ 上的纤维是 Jacobson 概形，参见《态射》，
引理 \ref{morphisms-lemma-ubiquity-Jacobson-schemes}，可以找到
$t'' \in \overline{\{t\}}$ 中的一个闭点，使 $Z_{t''} = \emptyset$。
于是 $(g : T \to S, t' \leadsto t'', \xi)$ 就是所需的非纯性。
\end{proof}

\begin{lemma}
\label{flat-lemma-impure-limit}
在情形 \ref{flat-situation-pre-pure} 中，设
$(g : T \to S, t' \leadsto t, \xi)$ 是 $\mathcal{F}$ 在 $s$ 之上的一个
非纯性。假设 $T = \lim_{i \in I} T_i$ 是 $S$ 上仿射概形的一个有向极限。
则对于某个 $i$，三元组
$(T_i \to S, t'_i \leadsto t_i, \xi_i)$ 是 $\mathcal{F}$ 在 $s$ 之上的
一个非纯性。
\end{lemma}

\begin{proof}
命题中的记号含义如下：令 $p_i : T \to T_i$ 为投影态射，令
$t_i = p_i(t)$ 且 $t'_i = p_i(t')$。最后，$\xi_i \in X_{T_i}$ 是
$\xi$ 的像。由《除子》，引理 \ref{divisors-lemma-base-change-relative-assassin}，
$\xi_i$ 确实是 $\mathcal{F}_{T_i}$ 关于 $T_i$ 的相对伴随点。
因此唯一需要证明的是：对某个 $i$，有
$\overline{\{\xi_i\}} \cap X_{t_i} = \emptyset$。

\medskip\noindent
第一种证明。令 $Z_i = \overline{\{\xi_i\}} \subset X_{T_i}$，
$Z = \overline{\{\xi\}} \subset X_T$，并赋予它们约化的诱导概形结构。
由《极限》，引理 \ref{limits-lemma-inverse-limit-irreducibles}，
$Z = \lim Z_i$。取一个域 $k$ 以及一个像为 $t$ 的态射
$\Spec(k) \to T$。则
$$
\emptyset =
Z \times_T \Spec(k) = (\lim Z_i) \times_{(\lim T_i)} \Spec(k)
= \lim Z_i \times_{T_i} \Spec(k)
$$
这是因为极限与纤维积交换（极限与极限交换）。每个
$Z_i \times_{T_i} \Spec(k)$ 都是拟紧的，因为 $X_{T_i} \to T_i$
有限型，从而 $Z_i \to T_i$ 有限型。由《极限》，引理
\ref{limits-lemma-limit-nonempty}，因此对于某个 $i$，
$Z_i \times_{T_i} \Spec(k)$ 为空。由于复合态射
$\Spec(k) \to T \to T_i$ 的像是 $t_i$，所需结论成立。

\medskip\noindent
第二种证明。置 $Z = \overline{\{\xi\}}$。将《极限》，引理
\ref{limits-lemma-separate} 应用于此情形，得到 $t$ 的一个开邻域
$V \subset T$、一个交换图
$$
\xymatrix{
V \ar[d] \ar[r]_a & T' \ar[d]^b \\
T \ar[r]^g & S,
}
$$
以及一个闭子概形 $Z' \subset X_{T'}$，满足
\begin{enumerate}
\item 态射 $b : T' \to S$ 局部有限表示；
\item $Z' \cap X_{a(t)} = \emptyset$；
\item $Z \cap X_V$ 通过态射 $X_V \to X_{T'}$ 映入 $Z'$。
\end{enumerate}
可以假设 $V$ 是 $T$ 的仿射开子集，因此由《极限》，引理
\ref{limits-lemma-descend-opens} 和 \ref{limits-lemma-limit-affine}，
可以找到某个 $i$ 和一个仿射开子集 $V_i \subset T_i$，使得
$V = f_i^{-1}(V_i)$。由《极限》，命题
\ref{limits-proposition-characterize-locally-finite-presentation}，
在必要时将 $i$ 增大后，可以找到态射 $a_i : V_i \to T'$，使得
$a = a_i \circ f_i|_V$。诱导的态射 $X_{V_i} \to X_{T'}$ 将 $\xi_i$
映入 $Z'$。由于 $Z' \cap X_{a(t)} = \emptyset$，可得
$(T_i \to S, t'_i \leadsto t_i, \xi_i)$ 是 $\mathcal{F}$ 在 $s$ 之上的
一个非纯性。
\end{proof}

\begin{lemma}
\label{flat-lemma-quasi-finite-impurity-elementary}
在情形 \ref{flat-situation-pre-pure} 中，若存在一个 $\mathcal{F}$ 在 $s$ 之上的
非纯性 $(g : T \to S, t' \leadsto t, \xi)$，且 $g$ 在 $t$ 处拟有限，
则存在一个非纯性 $(g : T \to S, t' \leadsto t, \xi)$，使
$(T, t) \to (S, s)$ 是初等 \'etale 邻域。
\end{lemma}

\begin{proof}
设 $(g : T \to S, t' \leadsto t, \xi)$ 是 $\mathcal{F}$ 在 $s$ 之上的一个
非纯性，且 $g$ 在 $t$ 处拟有限。缩小 $T$ 后，可以假设 $g$ 局部有限型。
将《态射进阶》，引理
\ref{more-morphisms-lemma-etale-makes-quasi-finite-finite-at-point}
应用于 $T \to S$ 与 $t \mapsto s$，得到图
$$
\xymatrix{
T \ar[d] & T \times_S U \ar[l] \ar[d] & V \ar[l] \ar[ld] \\
S & U \ar[l]
}
$$
其中 $(U, u) \to (S, s)$ 是一个初等 \'etale 邻域，且
$V \subset T \times_S U$ 是 $v = (t, u)$ 的一个开邻域，满足
$V \to U$ 有限，并且 $v$ 是 $V$ 中位于 $u$ 之上的唯一点。
由于态射 $V \to T$ 是 \'etale 的，因而平坦，可见存在一个特化
$v' \leadsto v$，使得 $v' \mapsto t'$。注意
$\kappa(t') \subset \kappa(v')$ 是有限可分扩张。任取
任取点 $\zeta \in X_{v'}$，它映到 $\xi \in X_{t'}$。由《除子》，引理
\ref{divisors-lemma-base-change-relative-assassin}，可见
$\zeta \in \text{Ass}_{X_V/V}(\mathcal{F}_V)$。
此外，$\overline{\{\zeta\}}$ 不与纤维 $X_v$ 相交，因为按假设
$\overline{\{\xi\}}$ 不与 $X_t$ 相交。换言之，
$(V \to S, v' \leadsto v, \zeta)$ 是 $\mathcal{F}$ 在 $S$ 之上的一个
非纯性。

\medskip\noindent
接着，令 $u' \in U'$ 为 $v'$ 的像，令 $\theta \in X_U$ 为 $\zeta$ 的像。
则 $\theta \mapsto u'$ 且 $u' \leadsto u$。由《除子》，引理
\ref{divisors-lemma-base-change-relative-assassin}，可见
$\theta \in \text{Ass}_{X_U/U}(\mathcal{F})$。
此外，由于 $\pi : X_V \to X_U$ 是有限的，有
$\pi\big(\overline{\{\zeta\}}\big) = \overline{\{\pi(\zeta)\}}$。
由于 $v$ 是 $V$ 中位于 $u$ 之上的唯一点，可见
$X_u \cap \overline{\{\pi(\zeta)\}} = \emptyset$，因为
$X_v \cap \overline{\{\zeta\}} = \emptyset$。因此
$(U \to S, u' \leadsto u, \theta)$ 是 $\mathcal{F}$ 在 $s$ 之上的一个
非纯性，结论得证。
\end{proof}

\begin{lemma}
\label{flat-lemma-Noetherian-impurity-quasi-finite}
在情形 \ref{flat-situation-pre-pure} 中，假设 $S$ 局部 Noether。
若存在 $\mathcal{F}$ 在 $s$ 之上的一个非纯性，则存在一个
$\mathcal{F}$ 在 $s$ 之上的非纯性 $(g : T \to S, t' \leadsto t, \xi)$，
使得 $g$ 在 $t$ 处拟有限。
\end{lemma}

\begin{proof}
可以用 $s$ 的一个仿射邻域替换 $S$。由引理
\ref{flat-lemma-impure-finite-presentation}，可以假设存在一个非纯性
$(g : T \to S, t' \leadsto t, \xi)$，其中 $g$ 局部有限型，且 $t$ 是
$g$ 在 $s$ 之上的纤维中的闭点。可以用
$\overline{\{t'\}}$ 上的约化诱导概形结构替换 $T$。令
$Z = \overline{\{\xi\}} \subset X_T$。由假设 $Z_t = \emptyset$，且
$Z \to T$ 的像包含 $t'$。由《态射进阶》，引理
\ref{more-morphisms-lemma-relative-assassin-in-neighbourhood}，存在
Z 的一个非空开子集 $V \subset Z$，使得对于任意 $w \in f(V)$，
$V_w$ 的任意一般点 $\xi'$ 都属于
$\text{Ass}_{X_T/T}(\mathcal{F}_T)$。由《态射进阶》，引理
\ref{more-morphisms-lemma-nonempty-generic-fibre}，存在一个非空开集
$W \subset T$，满足 $W \subset f(V)$。由《态射进阶》，引理
\ref{more-morphisms-lemma-quasi-finite-quasi-section-meeting-nearby-open-X}，
存在一个闭子概形 $T' \subset T$，使得 $t \in T'$，$T' \to S$ 在 $t$
处拟有限，并且存在 $z \in T' \cap W$，满足 $z \leadsto t$ 且不映到 $s$。
取非空概形 $V_z$ 的任意一般点 $\xi'$。则
$(T' \to S, z \leadsto t, \xi')$ 是所需的非纯性。
\end{proof}

\noindent
下文将使用 $S$ 在 $s$ 处的 Hensel 化
$S^h = \Spec(\mathcal{O}_{S, s}^h)$，参见《\'Etale 上同调》，定义
\ref{etale-cohomology-definition-etale-local-rings}。
由于 $S^h \to S$ 将 $S^h$ 的闭点映到 $s$，并诱导剩余域的同构，
我们也用 $s \in S^h$ 表示这个闭点。因此
$(S^h, s) \to (S, s)$ 是带基点概形的一个态射。

\begin{lemma}
\label{flat-lemma-impurity-on-henselization}
在情形 \ref{flat-situation-pre-pure} 中，若存在 $\mathcal{F}$ 在 $s$ 之上的
非纯性 $(S^h \to S, s' \leadsto s, \xi)$，则存在一个
$\mathcal{F}$ 在 $s$ 之上的非纯性 $(T \to S, t' \leadsto t, \xi)$，
其中 $(T, t) \to (S, s)$ 是初等 \'etale 邻域。
\end{lemma}

\begin{proof}
可以用 $s$ 的一个仿射邻域替换 $S$。设 $S = \Spec(A)$，且 $s$ 对应于
$\mathfrak p \subset A$。则
$\mathcal{O}_{S, s}^h = \colim_{(T, t)} \Gamma(T, \mathcal{O}_T)$，
其中极限遍历 $(S, s)$ 的仿射初等 \'etale 邻域 $(T, t)$ 所成的
余滤范畴的相反范畴，参见《态射进阶》，引理
\ref{more-morphisms-lemma-describe-henselization} 及其证明。因此
$S^h = \lim_i T_i$，由引理 \ref{flat-lemma-impure-limit} 即得结论。
\end{proof}

\begin{lemma}
\label{flat-lemma-pure-along-X-s}
在情形 \ref{flat-situation-pre-pure} 中，以下各条件等价：
\begin{enumerate}
\item 存在 $\mathcal{F}$ 在 $s$ 之上的非纯性
$(S^h \to S, s' \leadsto s, \xi)$，其中 $S^h$ 是 $S$ 在 $s$ 处的
Hensel 化；
\item 存在 $\mathcal{F}$ 在 $s$ 之上的非纯性
$(T \to S, t' \leadsto t, \xi)$，使 $(T, t) \to (S, s)$ 是初等
\'etale 邻域；
\item 存在 $\mathcal{F}$ 在 $s$ 之上的非纯性
$(T \to S, t' \leadsto t, \xi)$，使 $T \to S$ 在 $t$ 处拟有限。
\end{enumerate}
\end{lemma}

\begin{proof}
由于 \'etale 态射局部拟有限，(2) 蕴含 (3) 是显然的。在引理
\ref{flat-lemma-quasi-finite-impurity-elementary} 中已经证明 (3) 蕴含 (2)。
在引理 \ref{flat-lemma-impurity-on-henselization} 中已经证明 (1) 蕴含 (2)。
最后，若 $(T \to S, t' \leadsto t, \xi)$ 是 $\mathcal{F}$ 在 $s$ 之上的
一个非纯性，且 $(T, t) \to (S, s)$ 是初等 \'etale 邻域，则可以选取
结构态射 $S^h \to T \to S$ 的分解 $S^h \to S$。
选取任意点 $s' \in S^h$，它映到 $t'$，并选取任意
$\xi' \in X_{s'}$，它映到 $\xi \in X_{t'}$。则
$(S^h \to S, s' \leadsto s, \xi')$ 是 $\mathcal{F}$ 在 $s$ 之上的
一个非纯性。细节略去。
\end{proof}





\section{相对纯模}
\label{flat-section-pure}

\noindent
相对于底空间纯的模这一概念由 \cite{GruRay} 引入。

\begin{definition}
\label{flat-definition-pure}
设 $f : X \to S$ 是有限型概形态射。
设 $\mathcal{F}$ 是有限型拟凝聚 $\mathcal{O}_X$-模。
\begin{enumerate}
\item 设 $s \in S$。若不存在 $\mathcal{F}$ 在 $s$ 之上的非纯性
$(g : T \to S, t' \leadsto t, \xi)$，且 $(T, t) \to (S, s)$ 是一个
初等 \'etale 邻域，则称 $\mathcal{F}$ {\it 沿 $X_s$ 纯}。
\item 若不存在 $\mathcal{F}$ 在 $s$ 之上的任何非纯性，则称
$\mathcal{F}$ {\it 沿 $X_s$ 普遍纯}。
\item 若 $\mathcal{O}_X$ 沿 $X_s$ 纯，则称 $X$ {\it 沿 $X_s$ 纯}。
\item 若对每个 $s \in S$，$\mathcal{F}$ 都沿 $X_s$ 普遍纯，则称
$\mathcal{F}$ {\it 普遍 $S$-纯}，或称其 {\it 相对于 $S$ 普遍纯}。
\item 若对每个 $s \in S$，$\mathcal{F}$ 都沿 $X_s$ 纯，则称
$\mathcal{F}$ {\it $S$-纯}，或称其 {\it 相对于 $S$ 纯}。
\item 若 $\mathcal{O}_X$ 相对于 $S$ 纯，则称 $X$ {\it $S$-纯}，或称其
{\it 相对于 $S$ 纯}。
\end{enumerate}
\end{definition}

\noindent
这里我们有意将范围限制为有限型态射，而不只是局部有限型态射；讨论见
注记 \ref{flat-remark-discuss-finite-type}。在该定义的情形下，引理
\ref{flat-lemma-pure-along-X-s} 告诉我们，下列条件等价：
\begin{enumerate}
\item $\mathcal{F}$ 沿 $X_s$ 纯；
\item 不存在非纯性 $(g : T \to S, t' \leadsto t, \xi)$，且 $g$
在 $t$ 处拟有限；
\item 不存在形如 $(S^h \to S, s' \leadsto s, \xi)$ 的非纯性，其中
$S^h$ 是 $S$ 在 $s$ 处的 Hensel 化。
\end{enumerate}
记 $X^h = X \times_S S^h$，并令 $\mathcal{F}^h$ 为 $\mathcal{F}$
到 $X^h$ 的拉回，则最后一个条件可以用下面更积极的方式表述：
\begin{enumerate}
\item[(4)] $\text{Ass}_{X^h/S^h}(\mathcal{F}^h)$ 的所有点都专化到
$X_s$ 中的点。
\end{enumerate}
特别地，显然 $\mathcal{F}$ 沿 $X_s$ 纯，当且仅当 $\mathcal{F}$ 拉回到
$X \times_S \Spec(\mathcal{O}_{S, s})$ 后沿 $X_s$ 纯。

\begin{remark}
\label{flat-remark-discuss-finite-type}
设 $f : X \to S$ 是局部有限型态射，且 $\mathcal{F}$ 是拟凝聚有限型
$\mathcal{O}_X$-模。在此情形，(1) 与 (2) 仍然等价，因为引理
\ref{flat-lemma-quasi-finite-impurity-elementary} 的证明并未使用 $f$ 是拟紧的
事实。(3) 与 (4) 等价也很清楚。不过，我们不知道 (1) 与 (3) 是否等价。
此时有时使用等价条件 (3) 和 (4) 来定义纯性会更方便，正如
\cite{GruRay} 中所做的那样。另一方面，对于许多应用，正确的概念似乎
确实应当是普遍纯。
\end{remark}

\noindent
一个自然问题是，相对于底空间纯这一性质是否在基变换下保持；也就是说，
纯是否等同于普遍纯。事实证明，在 Noether 环概形上（见引理
\ref{flat-lemma-Noetherian-base-change}），或者层是平坦的情况下（见引理
\ref{flat-lemma-finite-type-flat-pure-along-fibre-is-universal} 与
\ref{flat-lemma-finite-type-flat-pure-is-universal}），答案是肯定的。
一般而言这并不成立，即使态射和层都是有限表示的也不成立；反例见
例子一节 \ref{examples-section-pure-not-universally}。首先，我们将
这里对“普遍”的用法与通常的概念对应起来。

\begin{lemma}
\label{flat-lemma-base-change-universally}
设 $f : X \to S$ 是有限型概形态射。设 $\mathcal{F}$ 是有限型拟凝聚
$\mathcal{O}_X$-模。设 $s \in S$。下列条件等价：
\begin{enumerate}
\item $\mathcal{F}$ 沿 $X_s$ 普遍纯；
\item 对每个带点概形态射 $(S', s') \to (S, s)$，拉回
$\mathcal{F}_{S'}$ 沿 $X_{s'}$ 纯。
\end{enumerate}
特别地，$\mathcal{F}$ 相对于 $S$ 普遍纯，当且仅当 $\mathcal{F}$ 的每个
基变换 $\mathcal{F}_{S'}$ 都相对于 $S'$ 纯。
\end{lemma}

\begin{proof}
这是形式性的。
\end{proof}

\begin{lemma}
\label{flat-lemma-quasi-finite-base-change}
设 $f : X \to S$ 是有限型概形态射。设 $\mathcal{F}$ 是有限型拟凝聚
$\mathcal{O}_X$-模。设 $s \in S$。设 $(S', s') \to (S, s)$ 是带点概形态射。
若 $S' \to S$ 在 $s'$ 处拟有限，且 $\mathcal{F}$ 沿 $X_s$ 纯，则
$\mathcal{F}_{S'}$ 沿 $X_{s'}$ 纯。
\end{lemma}

\begin{proof}
% P05-EN-CH38-0044: Frozen English line 4917 reads "It" where the conditional sentence requires "If"; the canonical target renders the intended conditional and records the source defect.
若 $(T \to S', t' \leadsto t, \xi)$ 是 $\mathcal{F}_{S'}$ 在 $s'$ 之上的
一个非纯性，且 $T \to S'$ 在 $t$ 处拟有限，则
$(T \to S, t' \to t, \xi)$ 是 $\mathcal{F}$ 在 $s$ 之上的非纯性，且
$T \to S$ 在 $t$ 处拟有限；见《态射》，引理
\ref{morphisms-lemma-composition-quasi-finite}。因此，立即由
\ref{flat-definition-pure} 之后给出的纯性刻画 (2) 得到本引理。
\end{proof}

\begin{lemma}
\label{flat-lemma-Noetherian-base-change}
设 $f : X \to S$ 是有限型概形态射。设 $\mathcal{F}$ 是有限型拟凝聚
$\mathcal{O}_X$-模。设 $s \in S$。若 $\mathcal{O}_{S, s}$ 是 Noether 环，
则 $\mathcal{F}$ 沿 $X_s$ 纯，当且仅当 $\mathcal{F}$ 沿 $X_s$ 普遍纯。
\end{lemma}

\begin{proof}
首先可以将 $S$ 替换为 $\Spec(\mathcal{O}_{S, s})$，即可以假设 $S$ 是
Noether 环。然后使用引理 \ref{flat-lemma-Noetherian-impurity-quasi-finite} 和
\ref{flat-definition-pure} 之后讨论中给出的纯性刻画 (2)，即可得出结论。
\end{proof}

\noindent
纯性满足平坦下降。

\begin{lemma}
\label{flat-lemma-flat-descend-pure}
设 $f : X \to S$ 是有限型概形态射。设 $\mathcal{F}$ 是有限型拟凝聚
$\mathcal{O}_X$-模。设 $s \in S$。设 $(S', s') \to (S, s)$ 是带点概形态射。
假设 $S' \to S$ 在 $s'$ 处平坦。
\begin{enumerate}
\item 若 $\mathcal{F}_{S'}$ 沿 $X_{s'}$ 纯，则 $\mathcal{F}$ 沿 $X_s$ 纯。
\item 若 $\mathcal{F}_{S'}$ 沿 $X_{s'}$ 普遍纯，则 $\mathcal{F}$ 沿 $X_s$
普遍纯。
\end{enumerate}
\end{lemma}

\begin{proof}
设 $(T \to S, t' \leadsto t, \xi)$ 是 $\mathcal{F}$ 在 $s$ 之上的非纯性。
置 $T_1 = T \times_S S'$，并令 $t_1$ 为 $T_1$ 中映到 $t$ 和 $s'$ 的唯一点。
由于 $T_1 \to T$ 在 $t_1$ 处平坦，见《态射》，引理
\ref{morphisms-lemma-base-change-flat}，存在位于 $t' \leadsto t$ 之上的
专化 $t'_1 \leadsto t_1$，见《代数》，一节
\ref{algebra-section-going-up}。选取 $\xi_1 \in X_{t'_1}$，使其对应于
$\Spec(\kappa(t'_1) \otimes_{\kappa(t')} \kappa(\xi))$ 的一个泛点，见
《概形》，引理 \ref{schemes-lemma-points-fibre-product}。由《除子》，引理
\ref{divisors-lemma-base-change-relative-assassin}，有
$\xi_1 \in \text{Ass}_{X_{T_1}/T_1}(\mathcal{F}_{T_1})$。
由于 $\{\xi_1\}$ 在 $X_{T_1}$ 中的 Zariski 闭包映入 $\{\xi\}$ 在 $X_T$
中的 Zariski 闭包，我们得到该闭包与 $X_{t_1}$ 不相交。因此
$(T_1 \to S', t'_1 \leadsto t_1, \xi_1)$ 是 $\mathcal{F}_{S'}$ 在 $s'$
之上的非纯性。换言之，我们证明了本引理第 (2) 部分的逆否命题。最后，若
$(T, t) \to (S, s)$ 是初等 \'etale 邻域，则 $(T_1, t_1) \to (S', s')$
也是初等 \'etale 邻域，由此可见 (1) 成立。
\end{proof}

\begin{lemma}
\label{flat-lemma-supported-on-closed}
设 $i : Z \to X$ 是在概形 $S$ 上有限型概形的闭浸入。设 $s \in S$。
设 $\mathcal{F}$ 是 $Z$ 上的有限型拟凝聚层。则 $\mathcal{F}$ 沿 $Z_s$
（普遍）纯，当且仅当 $i_*\mathcal{F}$ 沿 $X_s$（普遍）纯。
\end{lemma}

\begin{proof}
这由《除子》，引理
\ref{divisors-lemma-relative-weak-assassin-finite} 得出。
\end{proof}








\section{相对纯层的例子}
\label{flat-section-examples-pure-sheaves}

\noindent
下面给出一些可以看出纯性含义的例子情形。

\begin{lemma}
\label{flat-lemma-proper-pure}
设 $f : X \to S$ 是有限型概形态射。
设 $\mathcal{F}$ 是有限型拟凝聚 $\mathcal{O}_X$-模。
\begin{enumerate}
\item 若 $\mathcal{F}$ 的支集在 $S$ 上适当，则 $\mathcal{F}$ 相对于 $S$
普遍纯。
\item 若 $f$ 适当，则 $\mathcal{F}$ 相对于 $S$ 普遍纯。
\item 若 $f$ 适当，则 $X$ 相对于 $S$ 普遍纯。
\end{enumerate}
\end{lemma}

\begin{proof}
首先将 (1) 归约为 (2)。具体地，令 $Z \subset X$ 为 $\mathcal{F}$ 的
概形论支集。令 $i : Z \to X$ 为相应的闭浸入，并写成
$\mathcal{F} = i_*\mathcal{G}$，其中 $\mathcal{G}$ 是有限型拟凝聚
$\mathcal{O}_Z$-模，见《态射》，一节 \ref{morphisms-section-support}。
在情形 (1) 中，按假设 $Z \to S$ 适当。因此由引理
\ref{flat-lemma-supported-on-closed}，情形 (1) 归约为情形 (2)。

\medskip\noindent
假设 $f$ 适当。设 $(g : T \to S, t' \leadsto t, \xi)$ 是 $\mathcal{F}$
在 $s \in S$ 之上的非纯性。由于 $f$ 适当，它是普遍闭的。因此
$f_T : X_T \to T$ 是闭态射。由于 $f_T(\xi) = t'$，这蕴含
$t \in f(\overline{\{\xi\}})$，矛盾。
\end{proof}

\begin{lemma}
\label{flat-lemma-quasi-finite-pure}
设 $f : X \to S$ 是概形的分离有限型态射。设 $\mathcal{F}$ 是有限型
拟凝聚 $\mathcal{O}_X$-模。假设对每个 $s \in S$，$\text{Supp}(\mathcal{F}_s)$
都是有限的。则下列条件等价：
\begin{enumerate}
\item $\mathcal{F}$ 相对于 $S$ 纯；
\item $\mathcal{F}$ 的概形论支集在 $S$ 上有限；
\item $\mathcal{F}$ 相对于 $S$ 普遍纯。
\end{enumerate}
特别地，给定拟有限分离态射 $X \to S$，可知 $X$ 相对于 $S$ 纯，当且仅当
$X \to S$ 有限。
\end{lemma}

\begin{proof}
令 $Z \subset X$ 为 $\mathcal{F}$ 的概形论支集，见《态射》，定义
\ref{morphisms-definition-scheme-theoretic-support}。则 $Z \to S$ 是带有
有限纤维的分离有限型概形态射。因此它是分离且拟有限的，见《态射》，引理
\ref{morphisms-lemma-quasi-finite}。由引理 \ref{flat-lemma-supported-on-closed}，
只需对 $Z \to S$ 以及视为 $Z$ 上有限型拟凝聚模的层 $\mathcal{F}$ 证明本引理。
于是可以假设 $X \to S$ 分离且拟有限，并且 $\text{Supp}(\mathcal{F}) = X$。

\medskip\noindent
由引理 \ref{flat-lemma-proper-pure} 以及《态射》，引理
\ref{morphisms-lemma-finite-proper}，(2) 蕴含 (3)。显然 (3) 蕴含 (1)。
假设 (1) 成立。我们证明 (2) 成立。显然可以假设 $S$ 仿射。由《态射进阶》，
引理 \ref{more-morphisms-lemma-quasi-finite-separated-pass-through-finite}，
可以找到图表
$$
\xymatrix{
X \ar[rd]_f \ar[rr]_j & & T \ar[ld]^\pi \\
& S &
}
$$
其中 $\pi$ 有限，$j$ 是拟紧开浸入。若证明 $j$ 是闭的，则 $j$ 是闭浸入，
从而 $f = \pi \circ j$ 有限。要证明 $j$ 闭，只需证明专化沿 $j$ 可提升，见
《概形》，引理 \ref{schemes-lemma-quasi-compact-closed}。令 $x \in X$，置
$t' = j(x)$，并设 $t' \leadsto t$ 是一个专化。我们要证明 $t \in j(X)$。
置 $s' = f(x)$、$s = \pi(t)$，于是 $s' \leadsto s$。由《态射进阶》，引理
\ref{more-morphisms-lemma-etale-splits-off-quasi-finite-part-technical}，
可以找到初等 \'etale 邻域 $(U, u) \to (S, s)$ 以及分解
$$
T_U = T \times_S U = V \amalg W
$$
为开闭子概形，使得 $V \to U$ 有限，并且 $V$ 中存在唯一映到 $u$ 的点 $v$，
且 $v$ 在 $T$ 中映到 $t$。由于 $V \to T$ 是 \'etale，可以提升一般化，见
《态射》，引理 \ref{morphisms-lemma-generalizations-lift-flat} 与
\ref{morphisms-lemma-etale-flat}。因此存在专化 $v' \leadsto v$，使得 $v'$
映到 $t' \in T$。特别地，$v' \in X_U \subset T_U$。记 $u' \in U$ 为 $t'$ 的像。
注意，由于 $X_{u'}$ 是有限离散集且
$X_{u'} = \text{Supp}(\mathcal{F}_{u'})$，有
$v' \in \text{Ass}_{X_U/U}(\mathcal{F})$。由于 $\mathcal{F}$ 相对于 $S$ 纯，
可知 $v'$ 必须专化到 $X_u$ 中的一个点。由于 $v$ 是 $V$ 中位于 $u$ 之上的唯一点
（且 $W$ 中没有点能成为 $v'$ 的专化），可知 $v \in X_u$。故 $t \in X$。
\end{proof}

\begin{lemma}
\label{flat-lemma-flat-geometrically-integral-fibres-pure}
设 $f : X \to S$ 是有限型平坦概形态射，且纤维几何整。则 $X$ 在 $S$ 上
普遍纯。
\end{lemma}

\begin{proof}
令 $\xi \in X$，其中 $s' = f(\xi)$，并设 $s' \leadsto s$ 是 $S$ 中的专化。
若 $\xi$ 是 $X_{s'}$ 的伴随点，则 $\xi$ 是唯一的泛点，因为 $X_{s'}$ 是整概形。
令 $\xi_0$ 为 $X_s$ 的唯一泛点。由于 $X \to S$ 平坦，可以在 $X$ 中将
$s' \leadsto s$ 提升为专化 $\xi' \leadsto \xi_0$，见《态射》，引理
\ref{morphisms-lemma-generalizations-lift-flat}。由于 $\xi$ 是 $X_{s'}$ 的泛点，
% P05-EN-CH38-0045: Frozen English line 5138 begins "The $\xi \leadsto \xi'$" and omits the copula; the canonical target preserves the mathematical relation and records the grammatical omission.
有 $\xi \leadsto \xi'$，从而 $\xi \leadsto \xi_0$。这意味着
$(\text{id}_S, s' \to s, \xi)$ 不是 $\mathcal{O}_X$ 在 $s$ 之上的非纯性。
由于“$f$ 有限型、平坦且纤维几何整”这一假设在基变换下保持，可知任意基变换后
都不存在非纯性。因此 $X$ 是普遍 $S$-纯。
\end{proof}

\begin{lemma}
\label{flat-lemma-affine-locally-projective-pure}
设 $f : X \to S$ 是有限型仿射概形态射。设 $\mathcal{F}$ 是有限型拟凝聚
$\mathcal{O}_X$-模，且 $f_*\mathcal{F}$ 在 $S$ 上局部投射，见《性质》，定义
\ref{properties-definition-locally-projective}。则 $\mathcal{F}$ 在 $S$ 上普遍纯。
\end{lemma}

\begin{proof}
将情形归约到 $S$ 是 Hensel 局部环的谱后，这由引理
\ref{flat-lemma-explain-why-pure} 得出。
\end{proof}




\section{纯性的一个判据}
\label{flat-section-criterion-purity}

\noindent
我们首先证明，对于有限型拟凝聚层的一个平坦族，相对伴随点会专化到邻近
纤维的相对伴随点（如果它们确实发生专化）。

\begin{lemma}
\label{flat-lemma-associated-point-specializes}
设 $f : X \to S$ 是有限型概形态射。设 $\mathcal{F}$ 是有限型拟凝聚
$\mathcal{O}_X$-模。设 $s \in S$。假设 $\mathcal{F}$ 在 $X_s$ 的所有点处
都在 $S$ 上平坦。设 $x' \in \text{Ass}_{X/S}(\mathcal{F})$，满足
$f(x') = s'$，且 $s' \leadsto s$ 是 $S$ 中的专化。若 $x'$ 专化到 $X_s$
中的一个点，则 $x' \leadsto x$，其中
$x \in \text{Ass}_{X_s}(\mathcal{F}_s)$。
\end{lemma}

\begin{proof}
设 $x' \leadsto t$，其中 $t \in X_s$。则可以找到专化
$x' \leadsto x \leadsto t$，使得 $x$ 对应于
$\overline{\{x'\}} \cap f^{-1}(\{s\})$ 的某个不可约分支的泛点。按假设，
$\mathcal{F}$ 在 $x$ 处在 $S$ 上平坦。由《态射进阶》，引理
\ref{more-morphisms-lemma-associated-point-specializes}，可知
$x \in \text{Ass}_{X/S}(\mathcal{F})$，即得所需结论。
\end{proof}

\begin{lemma}
\label{flat-lemma-criterion}
设 $f : X \to S$ 是有限型概形态射。设 $\mathcal{F}$ 是有限型拟凝聚
$\mathcal{O}_X$-模。设 $s \in S$。设 $(S', s') \to (S, s)$ 是一个
初等 \'etale 邻域，并设
$$
\xymatrix{
X \ar[d] & X' \ar[l]^g \ar[d] \\
S & S' \ar[l]
}
$$
是概形态射的交换图表。假设
\begin{enumerate}
\item $\mathcal{F}$ 在 $X_s$ 的所有点处都在 $S$ 上平坦；
\item $X' \to S'$ 是有限型的；
\item $g^*\mathcal{F}$ 沿 $X'_{s'}$ 纯；
\item $g : X' \to X$ 是 \'etale；
\item $g(X')$ 包含 $\text{Ass}_{X_s}(\mathcal{F}_s)$。
\end{enumerate}
在此情形，$\mathcal{F}$ 沿 $X_s$ 纯，当且仅当 $X' \to X \times_S S'$ 的像
包含 $\text{Ass}_{X \times_S S'/S'}(\mathcal{F} \times_S S')$ 中位于
$S'$ 内专化到 $s'$ 的点之上的那些点。
\end{lemma}

\begin{proof}
由于态射 $S' \to S$ 是 \'etale，若 $\mathcal{F}$ 沿 $X_s$ 纯，则
$\mathcal{F} \times_S S'$ 沿 $X_s$ 纯，见引理
\ref{flat-lemma-quasi-finite-base-change}。由于纯性满足平坦下降，见引理
\ref{flat-lemma-flat-descend-pure}，若 $\mathcal{F} \times_S S'$ 沿 $X_{s'}$ 纯，
则 $\mathcal{F}$ 沿 $X_s$ 纯。因此可以用 $S'$ 替换 $S$，并假设 $S = S'$，
于是 $g : X' \to X$ 是 $S$ 上有限型概形之间的 \'etale 态射。此外，可以用
$\Spec(\mathcal{O}_{S, s})$ 替换 $S$，并假设 $S$ 是局部的。

\medskip\noindent
首先假设 $\mathcal{F}$ 沿 $X_s$ 纯。在此情形，由纯性
$\text{Ass}_{X/S}(\mathcal{F})$ 的每个点都专化到 $X_s$ 中的某个点。因此由
引理 \ref{flat-lemma-associated-point-specializes}，可知
$\text{Ass}_{X/S}(\mathcal{F})$ 的每个点都专化到
$\text{Ass}_{X_s}(\mathcal{F}_s)$ 中的某个点。于是
$\text{Ass}_{X/S}(\mathcal{F})$ 的每个点都在 $g$ 的像中（因为该像是开集，
并包含 $\text{Ass}_{X_s}(\mathcal{F}_s)$）。

\medskip\noindent
反之，假设 $g(X')$ 包含 $\text{Ass}_{X/S}(\mathcal{F})$。令
$S^h = \Spec(\mathcal{O}_{S, s}^h)$ 为 $S$ 在 $s$ 处的 Hensel 化。记
$g^h : (X')^h \to X^h$ 为沿 $S^h \to S$ 对 $g$ 作的基变换，并记
$\mathcal{F}^h$ 为 $\mathcal{F}$ 到 $X^h$ 的拉回。由《除子》，引理
\ref{divisors-lemma-base-change-relative-assassin} 和注记
\ref{divisors-remark-base-change-relative-assassin}，相对伴随点集
$\text{Ass}_{X^h/S^h}(\mathcal{F}^h)$ 是经投影 $X^h \to X$ 的
$\text{Ass}_{X/S}(\mathcal{F})$ 的逆像。由于假设 $g(X')$ 包含
$\text{Ass}_{X/S}(\mathcal{F})$，可知基变换
$g^h((X')^h) = g(X') \times_S S^h$ 包含
$\text{Ass}_{X^h/S^h}(\mathcal{F}^h)$。因此可以归约到 $S$ 是 Hensel 局部环的谱的情形。
令 $x \in \text{Ass}_{X/S}(\mathcal{F})$。要完成引理的证明，需要证明 $x$ 专化到
$X_s$ 中的一个点；这正是 \ref{flat-definition-pure} 之后关于纯性的讨论中的
判据 (4)。按假设，存在 $x' \in X'$ 使得 $g(x') = x$。由于
$g : X' \to X$ 是 \'etale，可知 $x' \in \text{Ass}_{X'/S}(g^*\mathcal{F})$，见
引理 \ref{flat-lemma-etale-weak-assassin-up-down}（应用于纤维态射
$X'_w \to X_w$，其中 $w \in S$ 是 $x'$ 的像）。由于 $g^*\mathcal{F}$ 沿
$X'_s$ 纯，可知对某个 $y \in X'_s$ 有 $x' \leadsto y$。于是
$x = g(x') \leadsto g(y)$，且 $g(y) \in X_s$，即得所需结论。
\end{proof}

\begin{lemma}
\label{flat-lemma-finite-type-flat-pure-along-fibre-is-universal}
设 $f : X \to S$ 是概形态射。设 $\mathcal{F}$ 是拟凝聚
$\mathcal{O}_X$-模。设 $s \in S$。假设
\begin{enumerate}
\item $f$ 是有限型的；
\item $\mathcal{F}$ 是有限型的；
\item $\mathcal{F}$ 在 $X_s$ 的所有点处都在 $S$ 上平坦；
\item $\mathcal{F}$ 沿 $X_s$ 纯。
\end{enumerate}
则 $\mathcal{F}$ 沿 $X_s$ 普遍纯。
\end{lemma}

\begin{proof}
先作一个预备说明。设 $(S', s') \to (S, s)$ 是初等 \'etale 邻域。记
$\mathcal{F}'$ 为 $\mathcal{F}$ 到 $X' = X \times_S S'$ 的拉回。
由 \ref{flat-definition-pure} 之后的讨论可知，$\mathcal{F}'$ 沿 $X'_{s'}$ 纯。
此外，$\mathcal{F}'$ 沿 $X'_{s'}$ 在 $S'$ 上平坦。因此只需证明
$\mathcal{F}'$ 沿 $X'_{s'}$ 普遍纯。事实上，给定带点概形的任意态射
$(T, t) \to (S, s)$，纤维积 $(T', t') = (T \times_S S', (t, s'))$ 在
$(T, t)$ 上平坦；因此若 $\mathcal{F}_{T'}$ 沿 $X_{t'}$ 纯，则由引理
\ref{flat-lemma-flat-descend-pure}，$\mathcal{F}_T$ 沿 $X_t$ 纯。
所以在证明过程中总可以用初等 \'etale 邻域替换 $(s, S)$。
由于问题具有局部性，也可以用 $\Spec(\mathcal{O}_{S, s})$ 替换 $S$。

\medskip\noindent
选取初等 \'etale 邻域 $(S', s') \to (S, s)$ 以及交换图表
$$
\xymatrix{
X \ar[d] & X' \ar[l]^g \ar[d] \\
S & \Spec(\mathcal{O}_{S', s'}) \ar[l]
}
$$
使得 $X' \to X \times_S \Spec(\mathcal{O}_{S', s'})$ 是 \'etale，
$X_s = g((X')_{s'})$，概形 $X'$ 仿射，并且
$\Gamma(X', g^*\mathcal{F})$ 是自由的 $\mathcal{O}_{S', s'}$-模；见引理
\ref{flat-lemma-finite-type-flat-along-fibre-free-variant}。注意，
$X' \to \Spec(\mathcal{O}_{S', s'})$ 是有限型的（因为它是一个拟紧态射，
由一个 \'etale 态射与一个有限型态射的基变换复合而成）。根据证明第一段的
预备说明，可以用 $\Spec(\mathcal{O}_{S', s'})$ 替换 $S$。因此可以假设存在
交换图表
$$
\xymatrix{
X \ar[dr] & & X' \ar[ll]^g \ar[ld] \\
& S &
}
$$
其中的概形都是 $S$ 上有限型的，$g$ 是 \'etale，$X_s \subset g(X')$，
$S$ 是闭点为 $s$ 的局部概形，$X'$ 仿射，并且
$\Gamma(X', g^*\mathcal{F})$ 是自由的 $\Gamma(S, \mathcal{O}_S)$-模。
注意此时 $g^*\mathcal{F}$ 在 $S$ 上普遍纯，见引理
\ref{flat-lemma-affine-locally-projective-pure}。

\medskip\noindent
在此情形，应用引理 \ref{flat-lemma-criterion}，由 $\mathcal{F}$ 沿 $X_s$ 纯且
沿 $X_s$ 在 $S$ 上平坦的假设，可推出
$\text{Ass}_{X/S}(\mathcal{F}) \subset g(X')$。由《除子》，引理
\ref{divisors-lemma-base-change-relative-assassin} 及注记
\ref{divisors-remark-base-change-relative-assassin}，对任意带点概形态射
$(T, t) \to (S, s)$，有
$$
\text{Ass}_{X_T/T}(\mathcal{F}_T) \subset
(X_T \to X)^{-1}(\text{Ass}_{X/S}(\mathcal{F})) \subset
g(X') \times_S T = g_T(X'_T).
$$
因此，将所显示图表基变换到 $(T, t)$ 后应用引理 \ref{flat-lemma-criterion}，
可得 $\mathcal{F}_T$ 沿 $X_t$ 纯，这就是所需结论。
\end{proof}

\begin{lemma}
\label{flat-lemma-finite-type-flat-pure-is-universal}
设 $f : X \to S$ 是有限型概形态射。设 $\mathcal{F}$ 是有限型拟凝聚
$\mathcal{O}_X$-模。假设 $\mathcal{F}$ 在 $S$ 上平坦。在此情形，
$\mathcal{F}$ 相对于 $S$ 纯，当且仅当 $\mathcal{F}$ 相对于 $S$ 普遍纯。
\end{lemma}

\begin{proof}
这是引理 \ref{flat-lemma-finite-type-flat-pure-along-fibre-is-universal} 与定义的
直接推论。
\end{proof}

\begin{lemma}
\label{flat-lemma-limit-purity}
设 $I$ 是有向集。设 $(S_i, g_{ii'})$ 是 $I$ 上的仿射概形逆系统。
置 $S = \lim_i S_i$，并令 $s \in S$。记 $g_i : S \to S_i$ 为投影，并置
$s_i = g_i(s)$。假设 $f : X \to S$ 是有限表示态射，$\mathcal{F}$ 是有限表示的
拟凝聚 $\mathcal{O}_X$-模，沿 $X_s$ 纯，并且在 $X_s$ 的所有点处都在 $S$ 上平坦。
则存在 $i \in I$、有限表示态射 $X_i \to S_i$ 以及有限表示的拟凝聚
$\mathcal{O}_{X_i}$-模 $\mathcal{F}_i$，使得沿 $(X_i)_{s_i}$ 纯，
并在 $(X_i)_{s_i}$ 的所有点处在 $S_i$ 上平坦，且
$X \cong X_i \times_{S_i} S$，并且 $\mathcal{F}_i$ 到 $X$ 的拉回同构于
$\mathcal{F}$。
\end{lemma}

\begin{proof}
令 $U \subset X$ 为 $\mathcal{F}$ 在 $S$ 上平坦的点集。由《态射进阶》，定理
\ref{more-morphisms-theorem-openness-flatness}，这是 $X$ 的开子概形。按假设
$X_s \subset U$。由于 $X_s$ 拟紧，可以找到一个拟紧开集 $U' \subset U$，满足
$X_s \subset U'$。由《极限》，引理 \ref{limits-lemma-descend-finite-presentation}，
可以找到 $i \in I$ 以及有限表示态射 $f_i : X_i \to S_i$，其到 $S$ 的基变换
同构于 $f_i$。固定这样的选择，并置 $X_{i'} = X_i \times_{S_i} S_{i'}$。
于是 $X = \lim_{i'} X_{i'}$，过渡态射是仿射的。由《极限》，引理
\ref{limits-lemma-descend-modules-finite-presentation}，增大 $i$ 后可以假设存在
有限表示的拟凝聚 $\mathcal{O}_{X_i}$-模 $\mathcal{F}_i$，其到 $S$ 的基变换同构于
$\mathcal{F}$。由《极限》，引理 \ref{limits-lemma-descend-opens}，再增大 $i$ 后可以
假设存在开集 $U'_i \subset X_i$，其在 $X$ 中的逆像为 $U'$。特别注意
$(X_i)_{s_i} \subset U'_i$。由《极限》，引理
\ref{limits-lemma-descend-module-flat-finite-presentation}（再一次增大 $i$ 后），
可以假设 $\mathcal{F}_i$ 在 $U'_i$ 上平坦。特别地，$\mathcal{F}_i$ 沿
$(X_i)_{s_i}$ 平坦。

\medskip\noindent
接下来，使用引理 \ref{flat-lemma-finite-presentation-flat-along-fibre}，选取初等
\'etale 邻域 $(S_i', s_i') \to (S_i, s_i)$ 以及概形的交换图表
$$
\xymatrix{
X_i \ar[d] & X_i' \ar[l]^{g_i} \ar[d] \\
S_i & S_i' \ar[l]
}
$$
使得 $g_i$ 是 \'etale，$(X_i)_{s_i} \subset g_i(X_i')$，概形 $X_i'$、$S_i'$ 都
仿射，并且 $\Gamma(X_i', g_i^*\mathcal{F}_i)$ 是
$\Gamma(S_i', \mathcal{O}_{S_i'})$-投射模。注意 $g_i^*\mathcal{F}_i$ 在 $S'_i$
上普遍纯，见引理 \ref{flat-lemma-affine-locally-projective-pure}。对任意 $i' \geq i$，
可以将上述图表基变换为 $S_{i'}$ 上带有态射
$(S'_{i'}, s'_{i'}) \to (S_{i'}, s_{i'})$ 和
$g_{i'} : X'_{i'} \to X_{i'}$ 的图表；也可以将图表基变换为 $S$ 上带有态射
$(S', s') \to (S, s)$ 和 $g : X' \to X$ 的图表。

\medskip\noindent
此时可以使用纯性判据。令 $W'_i \subset X_i \times_{S_i} S'_i$ 等于
\'etale 态射 $X'_i \to X_i \times_{S_i} S'_i$ 的像。对每个 $i' \geq i$，
类似地有像 $W'_{i'} \subset X_{i'} \times_{S_{i'}} S'_{i'}$，并有像
$W' \subset X \times_S S'$。取像与基变换交换，故
$W'_{i'} = W'_i \times_{S'_i} S'_{i'}$ 且 $W' = W_i \times_{S'_i} S'$。
由于 $\mathcal{F}$ 沿 $X_s$ 纯，引理 \ref{flat-lemma-criterion} 蕴含
\begin{equation}
\label{flat-equation-inclusion}
f^{-1}(\Spec(\mathcal{O}_{S', s'})) \cap
\text{Ass}_{X \times_S S'/S'}(\mathcal{F} \times_S S') \subset W'
\end{equation}
由《态射进阶》，引理
\ref{more-morphisms-lemma-relative-assassin-constructible}，可知
$$
E = \{t \in S' \mid \text{Ass}_{X_t}(\mathcal{F}_t) \subset W' \}
\quad\text{且}\quad
E_{i'} = \{t \in S'_{i'} \mid
\text{Ass}_{X_t}(\mathcal{F}_{i', t}) \subset W'_{i'} \}
$$
是 $S'$ 与 $S'_{i'}$ 的局部可构造子集。由《态射进阶》，引理
\ref{more-morphisms-lemma-base-change-assassin-in-U}，可知 $E_{i'}$ 是关于态射
$S'_{i'} \to S'_i$ 的 $E_i$ 逆像，而 $E$ 是关于态射 $S' \to S'_i$ 的 $E_i$
逆像。因此，等式（\ref{flat-equation-inclusion}）等价于
$\Spec(\mathcal{O}_{S', s'})$ 映入 $E_i$。由于
$\mathcal{O}_{S', s'} = \colim_{i' \geq i} \mathcal{O}_{S'_{i'}, s'_{i'}}$，
可知对某个 $i' \geq i$，$\Spec(\mathcal{O}_{S'_{i'}, s'_{i'}})$ 映入 $E_i$；见
《极限》，引理 \ref{limits-lemma-limit-contained-in-constructible}。然后将引理
\ref{flat-lemma-criterion} 应用于 $S_{i'}$ 上的情形，得到
$\mathcal{F}_{i'}$ 沿 $(X_{i'})_{s_{i'}}$ 纯。
\end{proof}

\begin{lemma}
\label{flat-lemma-flat-finite-presentation-purity-open}
设 $f : X \to S$ 是有限表示态射。设 $\mathcal{F}$ 是在 $S$ 上平坦的有限表示
拟凝聚 $\mathcal{O}_X$-模。则集合
$$
U = \{s \in S \mid \mathcal{F}\text{ 沿 }X_s\text{ 纯}\}
$$
在 $S$ 中是开的。
\end{lemma}

\begin{proof}
设 $s \in U$。使用引理 \ref{flat-lemma-finite-presentation-flat-along-fibre}，可以
找到初等 \'etale 邻域 $(S', s') \to (S, s)$ 以及交换图表
$$
\xymatrix{
X \ar[d] & X' \ar[l]^g \ar[d] \\
S & S' \ar[l]
}
$$
使得 $g$ 是 \'etale，$X_s \subset g(X')$，概形 $X'$、$S'$ 都仿射，且
$\Gamma(X', g^*\mathcal{F})$ 是 $\Gamma(S', \mathcal{O}_{S'})$-投射模。
注意 $g^*\mathcal{F}$ 在 $S'$ 上普遍纯，见引理
\ref{flat-lemma-affine-locally-projective-pure}。令 $W' \subset X \times_S S'$ 等于
\'etale 态射 $X' \to X \times_S S'$ 的像。% P05-EN-CH38-0046: Frozen English line 5517 names $W$ although the defined open is $W'$; the canonical target preserves the frozen symbol and records the source inconsistency.
注意 $W$ 在 $S'$ 上是开且拟紧的。
置
$$
E = \{t \in S' \mid \text{Ass}_{X_t}(\mathcal{F}_t) \subset W' \}.
$$
由《态射进阶》，引理
\ref{more-morphisms-lemma-relative-assassin-constructible}，$E$ 是 $S'$ 中的
可构造子集。由引理 \ref{flat-lemma-criterion}，可知
$\Spec(\mathcal{O}_{S', s'}) \subset E$。由《态射》，引理
\ref{morphisms-lemma-constructible-containing-open}，可知 $E$ 包含 $s'$ 的
一个开邻域 $V'$。再次应用引理 \ref{flat-lemma-criterion}，可知对 $V'$ 在 $S$ 中像内
的任意点 $s_1$，层 $\mathcal{F}$ 沿 $X_{s_1}$ 纯。由于 $S' \to S$ 是 \'etale，
$V'$ 在 $S$ 中的像是开的，结论得证。
\end{proof}


\section{纯性的应用}
\label{flat-section-applications-purity}

\noindent
下面给出一些使用纯性的例子。第一个引理实际上使用了稍弱形式的纯性。

\begin{lemma}
\label{flat-lemma-injectivity-map-source-flat-pure}
设 $f : X \to S$ 是有限型概形态射。设 $\mathcal{F}$ 是 $X$ 上有限型拟凝聚层。
假设 $S$ 是闭点为 $s$ 的局部概形。假设 $\mathcal{F}$ 沿 $X_s$ 纯，且
$\mathcal{F}$ 在 $S$ 上平坦。设 $\varphi : \mathcal{F} \to \mathcal{G}$ 是
拟凝聚 $\mathcal{O}_X$-模的态射。下列条件等价：
\begin{enumerate}
\item 对所有 $x \in \text{Ass}_{X_s}(\mathcal{F}_s)$，茎上的映射
$\varphi_x$ 都是单射；
\item $\varphi$ 是单射。
\end{enumerate}
\end{lemma}

\begin{proof}
令 $\mathcal{K} = \Ker(\varphi)$。我们的目标是证明 $\mathcal{K} = 0$。为此，
只需证明 $\text{WeakAss}_X(\mathcal{K}) = \emptyset$，见《除子》，引理
\ref{divisors-lemma-weakly-ass-zero}。有
$\text{WeakAss}_X(\mathcal{K}) \subset \text{WeakAss}_X(\mathcal{F})$，见
《除子》，引理 \ref{divisors-lemma-ses-weakly-ass}。由于 $\mathcal{F}$ 平坦，
由引理 \ref{flat-lemma-bourbaki-finite-type-general-base} 可知
$\text{WeakAss}_X(\mathcal{F}) \subset \text{Ass}_{X/S}(\mathcal{F})$。
由纯性，$\text{Ass}_{X/S}(\mathcal{F})$ 的任意点 $x'$ 都是 $X_s$ 中某点的
一般化，因此由引理 \ref{flat-lemma-associated-point-specializes}，它是某个
$x \in \text{Ass}_{X_s}(\mathcal{F}_s)$ 的专化。于是 $\varphi_x$ 的单射性蕴含
$\varphi_{x'}$ 的单射性，从而 $\mathcal{K}_{x'} = 0$。
\end{proof}

\begin{proposition}
\label{flat-proposition-finite-presentation-flat-pure-is-projective}
设 $f : X \to S$ 是概形的仿射有限表示态射。设 $\mathcal{F}$ 是有限表示、
在 $S$ 上平坦的拟凝聚 $\mathcal{O}_X$-模。下列条件等价：
\begin{enumerate}
\item $f_*\mathcal{F}$ 在 $S$ 上局部投射；
\item $\mathcal{F}$ 相对于 $S$ 纯。
\end{enumerate}
特别地，给定有限表示环映射 $A \to B$ 以及在 $A$ 上平坦的有限表示
$B$-模 $N$，有：$N$ 作为 $A$-模投射，当且仅当
$\widetilde{N}$ 在 $\Spec(B)$ 上相对于 $\Spec(A)$ 纯。
\end{proposition}

\begin{proof}
蕴含 (1) $\Rightarrow$ (2) 正是引理 \ref{flat-lemma-affine-locally-projective-pure}。
假设 $\mathcal{F}$ 相对于 $S$ 纯。注意，由引理
\ref{flat-lemma-finite-type-flat-pure-along-fibre-is-universal}，这蕴含 $\mathcal{F}$
在任意基变换后仍然纯。由《下降》，引理
\ref{descent-lemma-locally-projective-descends}，只需证明 $f_*\mathcal{F}$ 在
$S$ 上 fpqc 局部投射。取 $s \in S$。我们证明 $f_*\mathcal{F}$ 限制到 $s$ 的
某个 \'etale 邻域后局部投射。具体地，由引理
\ref{flat-lemma-finite-presentation-flat-along-fibre}，用 $s$ 的一个仿射初等 \'etale
邻域替换 $S$ 后，可以假设存在图表
$$
\xymatrix{
X \ar[dr] & & X' \ar[ll]^g \ar[ld] \\
& S &
}
$$
其中概形都是 $S$ 上仿射且有限表示的，$g$ 是 \'etale，$X_s \subset g(X')$，且
$\Gamma(X', g^*\mathcal{F})$ 是 $\Gamma(S, \mathcal{O}_S)$-投射模。
注意此时 $g^*\mathcal{F}$ 在 $S$ 上普遍纯，见引理
\ref{flat-lemma-affine-locally-projective-pure}。因此由引理 \ref{flat-lemma-criterion}，
可知开集 $g(X')$ 包含 $\text{Ass}_{X/S}(\mathcal{F})$ 中位于
$\Spec(\mathcal{O}_{S, s})$ 之上的点。置
$$
E = \{t \in S \mid \text{Ass}_{X_t}(\mathcal{F}_t) \subset g(X') \}.
$$
由《态射进阶》，引理
\ref{more-morphisms-lemma-relative-assassin-constructible}，$E$ 是 $S$ 中的
可构造子集。我们已经看到 $\Spec(\mathcal{O}_{S, s}) \subset E$。由《态射》，
引理 \ref{morphisms-lemma-constructible-containing-open}，可知 $E$ 包含 $s$ 的
一个开邻域。因此用 $s$ 的仿射邻域替换 $S$ 后，可以假设
$\text{Ass}_{X/S}(\mathcal{F}) \subset g(X')$。由引理
\ref{flat-lemma-base-change-universally-flat}，这意味着
$$
\Gamma(X, \mathcal{F}) \longrightarrow \Gamma(X', g^*\mathcal{F})
$$
是 $\Gamma(S, \mathcal{O}_S)$-普遍单射。由《代数》，引理
\ref{algebra-lemma-pure-submodule-ML}，可知 $\Gamma(X, \mathcal{F})$ 作为
$\Gamma(S, \mathcal{O}_S)$-模是 Mittag-Leffler 的。由于
$\Gamma(X, \mathcal{F})$ 是可数生成且作为 $\Gamma(S, \mathcal{O}_S)$-模平坦的，
由《代数》，引理 \ref{algebra-lemma-countgen-projective}，可知它是投射模。
\end{proof}

\noindent
我们可以利用该命题改进之前的一些结果。下面的引理改进了命题
\ref{flat-proposition-finite-presentation-flat-at-point}。

\begin{lemma}
\label{flat-lemma-flat-finite-presentation-affine-neighbourhood-projective}
设 $f : X \to S$ 是局部有限表示态射。设 $\mathcal{F}$ 是有限表示的
拟凝聚 $\mathcal{O}_X$-模。设 $x \in X$，且 $s = f(x) \in S$。
若 $\mathcal{F}$ 在 $x$ 处在 $S$ 上平坦，则存在仿射初等 \'etale 邻域
$(S', s') \to (S, s)$ 以及包含 $x' = (x, s')$ 的仿射开集
$U' \subset X \times_S S'$，使得 $\Gamma(U', \mathcal{F}|_{U'})$ 是
$\Gamma(S', \mathcal{O}_{S'})$-投射模。
\end{lemma}

\begin{proof}
在证明过程中，可以用 $x$ 的一个开邻域替换 $X$，并用 $s$ 的初等 \'etale
邻域替换 $S$。因此，由平坦性的开性（见《态射进阶》，定理
\ref{more-morphisms-theorem-openness-flatness}），可以假设 $\mathcal{F}$ 在
$S$ 上平坦。还可以假设 $S$ 和 $X$ 都仿射。再缩小 $X$ 后，可以假设
$\text{Ass}_{X_s}(\mathcal{F}_s)$ 的每个点都是 $x$ 的一般化。将 $(S, s)$
替换为初等 \'etale 邻域时这一性质保持。因此可以应用引理
\ref{flat-lemma-finite-presentation-flat-along-fibre}，归约到存在图表的情形
$$
\xymatrix{
X \ar[dr] & & X' \ar[ll]^g \ar[ld] \\
& S &
}
$$
其中概形都是 $S$ 上仿射且有限表示的，$g$ 是 \'etale，$X_s \subset g(X')$，且
$\Gamma(X', g^*\mathcal{F})$ 是 $\Gamma(S, \mathcal{O}_S)$-投射模。
注意此时 $g^*\mathcal{F}$ 在 $S$ 上普遍纯，见引理
\ref{flat-lemma-affine-locally-projective-pure}。

\medskip\noindent
令 $U \subset g(X')$ 为包含 $x$ 的仿射开邻域。我们断言 $\mathcal{F}|_U$ 沿
$U_s$ 纯。若证得此点，则由引理
\ref{flat-lemma-flat-finite-presentation-purity-open}，缩小 $S$ 后
$\mathcal{F}|_U$ 相对于 $S$ 纯，从而由命题
\ref{flat-proposition-finite-presentation-flat-pure-is-projective} 得到投射性。
为证明断言，需要证明：将 $(S, s)$ 替换为任意初等 \'etale 邻域后，
$\text{Ass}_{U/S}(\mathcal{F}|_U)$ 中位于某个 $s' \in S$ 之上的任意点
$\xi$（其中 $s' \leadsto s$）都专化到 $U_s$ 中的点。由于 $U \subset g(X')$，
可找到 $\xi' \in X'$ 使 $g(\xi') = \xi$。由于 $g^*\mathcal{F}$ 在 $S$ 上纯，
应用引理 \ref{flat-lemma-associated-point-specializes}，可知存在专化
$\xi' \leadsto x'$，其中 $x' \in \text{Ass}_{X'_s}(g^*\mathcal{F}_s)$。
于是 $g(x') \in \text{Ass}_{X_s}(\mathcal{F}_s)$（例如见引理
\ref{flat-lemma-etale-weak-assassin-up-down}，应用于 Noether 概形的 \'etale 态射
$X'_s \to X_s$），因而按上面对 $X$ 的选择有 $g(x') \leadsto x$！由于 $x \in U$，
可知 $g(x') \in U$。于是 $\xi = g(\xi') \leadsto g(x') \in U_s$，即得所需结论。
\end{proof}

\noindent
下面的引理改进了引理 \ref{flat-lemma-finite-type-flat-at-point-free-variant}。

\begin{lemma}
\label{flat-lemma-flat-finite-type-affine-neighbourhood-projective}
设 $f : X \to S$ 是局部有限型态射。设 $\mathcal{F}$ 是有限型拟凝聚
$\mathcal{O}_X$-模。设 $x \in X$，且 $s = f(x) \in S$。若 $\mathcal{F}$ 在 $x$
处在 $S$ 上平坦，则存在仿射初等 \'etale 邻域 $(S', s') \to (S, s)$ 以及
包含 $x' = (x, s')$ 的仿射开集
$U' \subset X \times_S \Spec(\mathcal{O}_{S', s'})$，使得
$\Gamma(U', \mathcal{F}|_{U'})$ 是自由的 $\mathcal{O}_{S', s'}$-模。
\end{lemma}

\begin{proof}
问题在 $X$ 和 $S$ 上都是 Zariski 局部的。因此可以假设 $X$ 和 $S$ 仿射。
于是可找到 $S$ 上的闭浸入 $i : X \to \mathbf{A}^n_S$。显然，只需对
$\mathbf{A}^n_S$ 上的层 $i_*\mathcal{F}$ 和点 $i(x)$ 证明引理即可。这样问题
归约到 $X \to S$ 有限表示的情形。用 $\Spec(\mathcal{O}_{S', s'})$ 替换 $S$，
并用 $X \times_S \Spec(\mathcal{O}_{S', s'})$ 的一个开子集替换 $X$ 后，可以假设
$\mathcal{F}$ 有限表示，见命题 \ref{flat-proposition-finite-type-flat-at-point}。
此时应用引理 \ref{flat-lemma-flat-finite-presentation-affine-neighbourhood-projective}
和《代数》，定理 \ref{algebra-theorem-projective-free-over-local-ring}，即可得出结论。
\end{proof}

\begin{lemma}
\label{flat-lemma-flat-finite-type-local-colimit-free}
设 $A \to B$ 是本质有限型局部环之间的局部环映射。设 $N$ 是有限 $B$-模，
作为 $A$-模平坦。若 $A$ 是 Hensel 环，则 $N$ 是自由 $A$-模 $F_i$ 的滤过余极限
$$
N = \colim_i F_i
$$
其中系统的所有过渡映射 $u_i : F_i \to F_{i'}$ 都诱导单射
$\overline{u}_i : F_i/\mathfrak m_AF_i \to F_{i'}/\mathfrak m_AF_{i'}$。
此外，$N$ 是 Mittag-Leffler $A$-模。
\end{lemma}

\begin{proof}
可以找到有限型态射 $X \to S = \Spec(A)$、位于 $S$ 闭点 $s$ 之上的点
$x \in X$，以及有限型拟凝聚 $\mathcal{O}_X$-模 $\mathcal{F}$，使得作为 $A$-模
$\mathcal{F}_x \cong N$。缩小 $X$ 后，可以假设
$\text{Ass}_{X_s}(\mathcal{F}_s)$ 的每个点都专化到 $x$。由引理
\ref{flat-lemma-flat-finite-type-affine-neighbourhood-projective}，存在 $x$ 的仿射开邻域
基本系 $U_i \subset X$，使得 $\Gamma(U_i, \mathcal{F})$ 是自由 $A$-模 $F_i$。
注意，若 $U_{i'} \subset U_i$，则
$$
F_i/\mathfrak m_AF_i = \Gamma(U_{i, s}, \mathcal{F}_s)
\longrightarrow
\Gamma(U_{i', s}, \mathcal{F}_s) = F_{i'}/\mathfrak m_AF_{i'}
$$
是单射，因为核中的一个截面会支撑在 $X_s$ 的某个不含 $x$ 的闭子集上，
这与上面对 $X$ 的选择矛盾。由于映射 $F_i \to F_{i'}$ 是 $A$-普遍单射
（引理 \ref{flat-lemma-universally-injective-local}），由《代数》，引理
\ref{algebra-lemma-colimit-universally-injective-ML}，可知 $N$ 是
Mittag-Leffler 模。
\end{proof}

\noindent
如果是第一次通读，可以跳过下面的引理。

\begin{lemma}
\label{flat-lemma-flat-finite-type-local-valuation-ring-has-content}
设 $A \to B$ 是本质有限型局部环之间的局部环映射。设 $N$ 是有限 $B$-模，
作为 $A$-模平坦。若 $A$ 是赋值环，则 $N$ 的任意元素都有内容理想
$I \subset A$（见《代数进阶》，定义
\ref{more-algebra-definition-content-ideal}）。此外，$I$ 是主理想。
\end{lemma}

\begin{proof}
最后一个断言由 $I$ 是有限生成理想这一事实得出；见《代数进阶》，引理
\ref{more-algebra-lemma-content-finitely-generated} 与《代数》，引理
\ref{algebra-lemma-characterize-valuation-ring}。

\medskip\noindent
$I$ 存在性的证明。令 $A \subset A^h$ 为 Hensel 化。令 $B'$ 为
$B \otimes_A A^h$ 在极大理想
$\mathfrak m_B \otimes A^h + B \otimes \mathfrak m_{A^h}$ 处的局部化。
则 $B \to B'$ 平坦，因而忠实平坦。令 $N' = N \otimes_B B'$。令
$x \in N$，并令 $x' \in N'$ 为其像。我们断言，对理想 $I \subset A$，有
$x \in IN \Leftrightarrow x' \in IN'$。事实上，$N/IN \to N'/IN'$ 是将
$B \to B'$ 与 $N/IN$ 张量所得的映射，而由《代数》，引理
\ref{algebra-lemma-faithfully-flat-universally-injective}，$B \to B'$ 普遍单射。
由《代数进阶》，引理 \ref{more-algebra-lemma-henselization-valuation-ring}，
以及《代数》，引理 \ref{algebra-lemma-ideals-valuation-ring}，映射 $A \to A^h$
在理想集上给出保持包含关系的双射 $I \mapsto IA^h$。因此，$x$ 在 $A$ 中有内容
理想，当且仅当 $x'$ 在 $A^h$ 中有内容理想。对于 $x' \in N'$ 的断言由引理
\ref{flat-lemma-flat-finite-type-local-colimit-free} 和《代数》，引理
\ref{algebra-lemma-minimal-contains} 得出。
\end{proof}

\noindent
应用如下。

\begin{lemma}
\label{flat-lemma-proper-flat-over-dvr-reduced-fibre}
设 $X \to \Spec(R)$ 是适当平坦态射，其中 $R$ 是赋值环。若特殊纤维约化，
则 $X$ 以及 $X \to \Spec(R)$ 的每个纤维都约化。
\end{lemma}

\begin{proof}
假设特殊纤维 $X_s$ 约化。令 $x \in X$ 为任意点，我们证明
$\mathcal{O}_{X, x}$ 约化；这将证明 $X$ 约化。令 $x \leadsto x'$ 是一个专化，
其中 $x'$ 位于特殊纤维中；这样的专化存在，因为适当态射是闭的。考虑局部环
$A = \mathcal{O}_{X, x'}$。则 $\mathcal{O}_{X, x}$ 是 $A$ 的局部化，所以只需证明
$A$ 约化。令 $a \in A$ 非零，并令 $I = (\pi) \subset R$ 为其内容理想，见引理
\ref{flat-lemma-flat-finite-type-local-valuation-ring-has-content}。于是 $a = \pi a'$，
且 $a'$ 映到 $A/\mathfrak mA$ 中的非零元素，其中 $\mathfrak m \subset R$ 是极大理想。
若 $a$ 幂零，则 $a'$ 也幂零，因为由 $A$ 在 $R$ 上平坦，$\pi$ 不是零因子。
但是 $a'$ 映到约化环 $A/\mathfrak m A = \mathcal{O}_{X_s, x'}$ 中的非零元素。
除非 $A$ 约化，否则这矛盾；而这正是要证明的结论。

\medskip\noindent
当然，若 $X$ 约化，则 $X$ 在 $R$ 上的泛纤维也约化。若 $\mathfrak p \subset R$
是素理想，则由《代数》，引理 \ref{algebra-lemma-make-valuation-rings}，
$R/\mathfrak p$ 是赋值环。因此将 $X$ 基变换到 $R/\mathfrak p$ 后重复上述论证，
即可证明位于 $\mathfrak p$ 之上的纤维约化。
\end{proof}






\section{扁平化函子}
\label{flat-section-flattening-functors}

\noindent
设 $S$ 是概形。回顾：若对每个以仿射概形组成、极限为 $T$ 的有向逆系统
$\{T_i\}_{i \in I}$，都有 $F(T) = \colim_i F(T_i)$，则称函子
$F : (\Sch/S)^{opp} \to \textit{Sets}$ 保持极限。

\begin{situation}
\label{flat-situation-iso}
设 $f : X \to S$ 是概形态射。设 $u : \mathcal{F} \to \mathcal{G}$ 是拟凝聚
$\mathcal{O}_X$-模的同态。对于任意在 $S$ 上的概形 $T$，记
$u_T : \mathcal{F}_T \to \mathcal{G}_T$ 为 $u$ 到 $T$ 的基变换；换言之，
$u_T$ 是沿投影态射 $X_T = X \times_S T \to X$ 对 $u$ 的拉回。
在此情形可以考虑函子
\begin{equation}
\label{flat-equation-iso}
F_{iso} : (\Sch/S)^{opp} \longrightarrow \textit{Sets}, \quad
T \longrightarrow \left\{
\begin{matrix}
\{*\} & \text{若} & u_T \text{ 是同构}, \\
\emptyset & \text{否则。} &
\end{matrix}
\right.
\end{equation}
还有变体 $F_{inj}$、$F_{surj}$、$F_{zero}$，分别要求 $u_T$ 是单射、满射或零映射。
\end{situation}

\begin{lemma}
\label{flat-lemma-iso-sheaf}
在情形 \ref{flat-situation-iso} 中。
\begin{enumerate}
\item 函子 $F_{iso}$、$F_{inj}$、$F_{surj}$、$F_{zero}$ 都满足 fpqc 拓扑的层性质。
\item 若 $f$ 拟紧且 $\mathcal{G}$ 有限型，则 $F_{surj}$ 保持极限。
\item 若 $f$ 拟紧且 $\mathcal{F}$ 有限型，则 $F_{zero}$ 保持极限。
\item 若 $f$ 拟紧、$\mathcal{F}$ 有限型且 $\mathcal{G}$ 有限表示，则
$F_{iso}$ 保持极限。
\end{enumerate}
\end{lemma}

\begin{proof}
设 $\{T_i \to T\}_{i \in I}$ 是 $S$ 上概形的 fpqc 覆盖。置
$X_i = X_{T_i} = X \times_S T_i$，并置 $u_i = u_{T_i}$。注意
$\{X_i \to X_T\}_{i \in I}$ 是 $X_T$ 的 fpqc 覆盖，见《拓扑》，引理
\ref{topologies-lemma-fpqc}。特别地，对每个 $x \in X_T$，存在 $i \in I$ 以及
映到 $x$ 的 $x_i \in X_i$。由于
$\mathcal{O}_{X_T, x} \to \mathcal{O}_{X_i, x_i}$ 平坦，因而忠实平坦（见《代数》，
引理 \ref{algebra-lemma-local-flat-ff}），可知 $(u_i)_{x_i}$ 是单射、满射、双射
或零映射，当且仅当 $(u_T)_x$ 具有相应性质。这就证明了引理的第 (1) 部分。

\medskip\noindent
(2) 的证明。假设 $f$ 拟紧且 $\mathcal{G}$ 有限型。令
$T = \lim_{i \in I} T_i$ 为仿射 $S$-概形的有向极限，并假设 $u_T$ 是满射。
置 $X_i = X_{T_i} = X \times_S T_i$，以及
$u_i = u_{T_i} : \mathcal{F}_i = \mathcal{F}_{T_i}
\to \mathcal{G}_i = \mathcal{G}_{T_i}$。为证明 (2)，需证明对某个 $i$，$u_i$ 是满射。
取 $i_0 \in I$，并用 $\{i \mid i \geq i_0\}$ 替换 $I$。由于 $f$ 拟紧，概形
$X_{i_0}$ 拟紧。因此可以选择仿射开集 $W_1, \ldots, W_m \subset X$ 以及仿射开覆盖
$X_{i_0} = U_{1, i_0} \cup \ldots \cup U_{m, i_0}$，使得 $U_{j, i_0}$
在投影态射 $X_{i_0} \to X$ 下映入 $W_j$。对任意 $i \in I$，令 $U_{j, i}$
为 $U_{j, i_0}$ 的逆像。置 $U_j = \lim_i U_{j, i}$，则
$X_T = U_1 \cup \ldots \cup U_m$ 是 $X_T$ 的仿射开覆盖。现在只需证明：对给定
$j \in \{1, \ldots, m\}$，存在 $i = i(j) \in I$ 使 $u_i|_{U_{j, i}}$
为满射。由《性质》，引理 \ref{properties-lemma-finite-type-module}，这转化为下列代数问题：
设 $A$ 是环，$u : M \to N$ 是 $A$-模映射。假设
$R = \colim_{i \in I} R_i$ 是 $A$-代数的有向余极限。若 $N$ 是有限 $A$-模，且
$u \otimes 1 : M \otimes_A R \to N \otimes_A R$ 是满射，则对某个 $i$，映射
$u \otimes 1 : M \otimes_A R_i \to N \otimes_A R_i$ 是满射。这就是《代数》，
引理 \ref{algebra-lemma-module-map-property-in-colimit} 的第 (2) 部分。

\medskip\noindent
(3) 的证明。与 (2) 的证明完全相同的论证将其归约为下列代数问题：
设 $A$ 是环，$u : M \to N$ 是 $A$-模映射。假设
$R = \colim_{i \in I} R_i$ 是 $A$-代数的有向余极限。若 $M$ 是有限 $A$-模，且
$u \otimes 1 : M \otimes_A R \to N \otimes_A R$ 为零，则对某个 $i$，映射
$u \otimes 1 : M \otimes_A R_i \to N \otimes_A R_i$ 为零。这就是《代数》，
引理 \ref{algebra-lemma-module-map-property-in-colimit} 的第 (1) 部分。

\medskip\noindent
(4) 的证明。假设 $f$ 拟紧且 $\mathcal{F}, \mathcal{G}$ 有限表示。以与上一段
完全相同的方式论证（另外使用《性质》，引理
\ref{properties-lemma-finite-presentation-module}），(3) 转化为下列代数断言：
设 $A$ 是环，$u : M \to N$ 是 $A$-模映射。假设
$R = \colim_{i \in I} R_i$ 是 $A$-代数的有向余极限。假设 $M$ 是有限 $A$-模，
$N$ 是有限表示 $A$-模，且
$u \otimes 1 : M \otimes_A R \to N \otimes_A R$ 是同构。则对某个 $i$，映射
$u \otimes 1 : M \otimes_A R_i \to N \otimes_A R_i$ 是同构。这就是《代数》，
引理 \ref{algebra-lemma-module-map-property-in-colimit} 的第 (3) 部分。
\end{proof}

\begin{situation}
\label{flat-situation-flat-at-point}
设 $(A, \mathfrak m_A)$ 为局部环。
记 $\mathcal{C}$ 为如下范畴：其对象是 $A$-代数 $A'$，
其中这些环是局部环，且代数结构
$A \to A'$ 是局部环之间的局部同态。
$\mathcal{C}$ 中对象 $A', A''$ 之间的态射是 $A$-代数的
局部同态 $A' \to A''$。
设 $A \to B$ 是局部环之间的局部环映射，且 $M$ 是 $B$-模。
若 $A'$ 是 $\mathcal{C}$ 的对象，则置 $B' = B \otimes_A A'$，
并将 $M' = M \otimes_A A'$ 视为 $B'$-模。
给定 $A' \in \Ob(\mathcal{C})$，考虑条件
\begin{equation}
\label{flat-equation-flat-at-primes}
\forall \mathfrak q \in V(\mathfrak m_{A'}B' + \mathfrak m_B B')
\subset \Spec(B') :
M'_{\mathfrak q}\text{ 在 }A'\text{ 上平坦}.
\end{equation}
注意这与《代数进阶》，方程
(\ref{more-algebra-equation-flat-at-primes}) 相似。
特别地，若 $A' \to A''$ 是 $\mathcal{C}$ 中的态射，且
(\ref{flat-equation-flat-at-primes}) 对 $A'$ 成立，则它对 $A''$ 也成立，见《代数进阶》，
引理 \ref{more-algebra-lemma-base-change-flat-at-primes}。
因此得到函子
\begin{equation}
\label{flat-equation-flat-at-point}
F_{lf} : \mathcal{C} \longrightarrow \textit{Sets}, \quad
A' \longrightarrow \left\{
\begin{matrix}
\{*\} & \text{若 }(\ref{flat-equation-flat-at-primes})\text{ 成立}, \\
\emptyset & \text{否则。} &
\end{matrix}
\right.
\end{equation}
\end{situation}

\begin{lemma}
\label{flat-lemma-flat-at-point}
在情形 \ref{flat-situation-flat-at-point} 中。
\begin{enumerate}
\item 若 $A' \to A''$ 是 $\mathcal{C}$ 中的平坦态射，
则 $F_{lf}(A') = F_{lf}(A'')$。
\item 若 $A \to B$ 本质有限表示，且
$M$ 是有限表示的 $B$-模，则 $F_{lf}$ 保持极限：
若 $\{A_i\}_{i \in I}$ 是 $\mathcal{C}$ 中对象的有向系统，则
$F_{lf}(\colim_i A_i) = \colim_i F_{lf}(A_i)$.
\end{enumerate}
\end{lemma}

\begin{proof}
第 (1) 部分是《代数进阶》引理
\ref{more-algebra-lemma-flat-descent-flat-at-primes} 的特例。
第 (2) 部分是《代数进阶》引理
\ref{more-algebra-lemma-limit-preserving-flat-at-primes} 的特例。
\end{proof}

\begin{lemma}
\label{flat-lemma-flat-at-point-finite}
在情形 \ref{flat-situation-flat-at-point} 中。设 $B \to C$ 是局部
$A$-代数之间的局部映射，且 $N$ 是 $C$-模。记
$F'_{lf} : \mathcal{C} \to \textit{Sets}$ 为与对 $(C, N)$ 关联的函子。
若 $M \cong N$ 作为 $B$-模且 $B \to C$ 有限，则
$F_{lf} = F'_{lf}$.
\end{lemma}

\begin{proof}
设 $A'$ 是 $\mathcal{C}$ 的对象。置 $C' = C \otimes_A A'$，
并仿照情形 \ref{flat-situation-flat-at-point} 中 $B'$、$M'$ 的定义
置 $N' = N \otimes_A A'$。
注意 $M' \cong N'$ 作为 $B'$-模。
假设 $B \to C$ 有限，则有两个结论：
(a) $\mathfrak m_C = \sqrt{\mathfrak m_B C}$，且 (b)
$B' \to C'$ 有限。结论 (a) 蕴含
$$
V(\mathfrak m_{A'}C' + \mathfrak m_C C')
=
\left(
\Spec(C') \to \Spec(B')
\right)^{-1}V(\mathfrak m_{A'}B' + \mathfrak m_B B').
$$
设 $\mathfrak q \subset V(\mathfrak m_{A'}B' + \mathfrak m_B B')$。
则 $M'_{\mathfrak q}$ 在 $A'$ 上平坦，当且仅当
$C'_{\mathfrak q}$-模 $N'_{\mathfrak q}$ 在 $A'$ 上平坦
（因为它们作为 $A'$-模同构），当且仅当
$C'_{\mathfrak q}$ 的每个极大理想 $\mathfrak r$ 都满足
$N'_{\mathfrak r}$ 在 $A'$ 上平坦（见《代数》，引理
\ref{algebra-lemma-flat-localization}）。
由 (b)，$B'_{\mathfrak q} \to C'_{\mathfrak q}$ 有限，
$C'_{\mathfrak q}$ 的极大理想恰好对应于
位于 $\mathfrak q$ 之上的 $C'$ 的素理想（见《代数》，引理
\ref{algebra-lemma-integral-going-up}），
而根据上面的显示方程，这些素理想全都包含在
$V(\mathfrak m_{A'}C' + \mathfrak m_C C')$ 中。因此引理得证。
\end{proof}

\begin{lemma}
\label{flat-lemma-flat-at-point-go-up}
在情形
\ref{flat-situation-flat-at-point} 中，设 $B \to C$ 是局部
环之间的平坦局部同态。置 $N = M \otimes_B C$。记
$F'_{lf} : \mathcal{C} \to \textit{Sets}$ 为与对 $(C, N)$ 关联
的函子。则 $F_{lf} = F'_{lf}$。
\end{lemma}

\begin{proof}
设 $A'$ 是 $\mathcal{C}$ 的对象。置 $C' = C \otimes_A A'$，
并仿照情形 \ref{flat-situation-flat-at-point} 中 $B'$、$M'$ 的定义
置 $N' = N \otimes_A A' = M' \otimes_{B'} C'$。注意
$$
V(\mathfrak m_{A'}B' + \mathfrak m_B B')
=
\Spec( \kappa(\mathfrak m_B) \otimes_A \kappa(\mathfrak m_{A'}) )
$$
并且对 $V(\mathfrak m_{A'}C' + \mathfrak m_C C')$ 也有类似等式。
环映射
$$
\kappa(\mathfrak m_B) \otimes_A \kappa(\mathfrak m_{A'})
\longrightarrow
\kappa(\mathfrak m_C) \otimes_A \kappa(\mathfrak m_{A'})
$$
是忠实平坦的，因此
$V(\mathfrak m_{A'}C' + \mathfrak m_C C') \to
V(\mathfrak m_{A'}B' + \mathfrak m_B B')$ 是满射。
最后，若 $\mathfrak r \in V(\mathfrak m_{A'}C' + \mathfrak m_C C')$
映到 $\mathfrak q \in V(\mathfrak m_{A'}B' + \mathfrak m_B B')$，则
$M'_{\mathfrak q}$ 在 $A'$ 上平坦，当且仅当
$N'_{\mathfrak r}$ 在 $A'$ 上平坦，因为 $B' \to C'$ 平坦，见
《代数》，引理 \ref{algebra-lemma-flatness-descends-more-general}。
引理由这些观察形式地推出。
\end{proof}

\begin{situation}
\label{flat-situation-free-at-generic-points}
设 $f : X \to S$ 是具有几何不可约纤维的光滑态射。
设 $\mathcal{F}$ 是有限型拟凝聚 $\mathcal{O}_X$-模。对任意
$S$ 上的概形 $T$，记 $\mathcal{F}_T$ 为 $\mathcal{F}$ 到 $T$ 的基变换，
换言之，$\mathcal{F}_T$ 是沿投影态射
$X_T = X \times_S T \to X$ 的 $\mathcal{F}$ 的拉回。注意 $X_T \to T$
是具有几何不可约纤维的光滑态射，见《态射》引理
\ref{morphisms-lemma-base-change-smooth} 及《态射进阶》引理
\ref{more-morphisms-lemma-base-change-fibres-geometrically-irreducible}。
设 $p \geq 0$ 为整数。给定点 $t \in T$，考虑条件
\begin{equation}
\label{flat-equation-free-at-generic-point-fibre}
\mathcal{F}_T\text{ 在 }\xi_t\text{ 的一个邻域内自由且秩为 }p
\end{equation}
其中 $\xi_t$ 是纤维 $X_t$ 的泛点。对所有 $t \in T$，此条件在基变换下稳定，
因此得到函子
\begin{equation}
\label{flat-equation-free-at-generic-points}
H_p : (\Sch/S)^{opp} \longrightarrow \textit{Sets}, \quad
T \longrightarrow \left\{
\begin{matrix}
\{*\} & \text{若 }\mathcal{F}_T\text{ 满足
(\ref{flat-equation-free-at-generic-point-fibre}) }\forall t\in T, \\
\emptyset & \text{否则。}
\end{matrix}
\right.
\end{equation}
\end{situation}

\begin{lemma}
\label{flat-lemma-free-at-generic-points}
在情形 \ref{flat-situation-free-at-generic-points} 中。
\begin{enumerate}
\item 函子 $H_p$ 对 fpqc 拓扑满足层性质。
\item 若 $\mathcal{F}$ 有限表示，则函子 $H_p$ 保持极限。
\end{enumerate}
\end{lemma}

\begin{proof}
设 $\{T_i \to T\}_{i \in I}$ 是概形在 $S$ 上的 fpqc\footnote{利用有限表示的平坦态射是开映射，很容易证明 $H_p$ 是 fppf 拓扑下的层。这正是我们稍后真正需要的。不过，直接证明它也满足 fpqc 拓扑的层条件也颇有趣。} 覆盖。
置 $X_i = X_{T_i} = X \times_S T_i$，记 $\mathcal{F}_i$ 为 $\mathcal{F}$ 到 $X_i$ 的拉回。
假设对所有 $i$，$\mathcal{F}_i$ 都满足
(\ref{flat-equation-free-at-generic-point-fibre})。
取 $t \in T$，令 $\xi_t \in X_T$ 表示 $X_t$ 的泛点。
我们需要证明 $\mathcal{F}$ 在 $\xi_t$ 的某个邻域内自由。
对某个 $i \in I$，可找到映到 $t$ 的 $t_i \in T_i$。
令 $\xi_i \in X_i$ 表示 $X_{t_i}$ 的泛点，于是 $\xi_i$ 映到 $\xi_t$。
由于 $\mathcal{F}_i$ 在 $\xi_i$ 的某个邻域内秩为 $p$ 的自由性，得到
$(\mathcal{F}_i)_{x_i} \cong \mathcal{O}_{X_i, x_i}^{\oplus p}$
这蕴含
$\mathcal{F}_{T, \xi_t} \cong \mathcal{O}_{X_T, \xi_t}^{\oplus p}$
因为 $\mathcal{O}_{X_T, \xi_t} \to \mathcal{O}_{X_i, x_i}$ 平坦，见例如
《代数》，引理 \ref{algebra-lemma-finite-projective-descends}。
因此在 $X_T$ 中存在 $\xi_t$ 的仿射邻域 $U$ 以及满射
$\mathcal{O}_U^{\oplus p} \to \mathcal{F}_U = \mathcal{F}_T|_U$，见
《模》，引理 \ref{modules-lemma-finite-type-surjective-on-stalk}。
缩小 $T$ 后可假设 $U \to T$ 是满射。因此 $U \to T$ 是具有几何不可约
纤维的仿射光滑态射。此外，对每个 $t' \in T$，诱导映射
$$
\alpha :
\mathcal{O}_{U, \xi_{t'}}^{\oplus p}
\longrightarrow
\mathcal{F}_{U, \xi_{t'}}
$$
是同构（因为由同样的论证，右侧的模是秩为 $p$ 的自由模）。由引理
\ref{flat-lemma-induction-step} 得
$$
\Gamma(U, \mathcal{O}_U^{\oplus p})
\otimes_{\Gamma(T, \mathcal{O}_T)} \mathcal{O}_{T, t'}
\longrightarrow
\Gamma(U, \mathcal{F}_U)
\otimes_{\Gamma(T, \mathcal{O}_T)} \mathcal{O}_{T, t'}
$$
对每个 $t' \in T$ 都是单射。因此满射 $\alpha$ 是同构。这就完成了 (1) 的证明。

\medskip\noindent
假设 $\mathcal{F}$ 有限表示。
设 $T = \lim_{i \in I} T_i$ 是仿射 $S$-概形的有向极限，
并假设 $\mathcal{F}_T$ 满足
(\ref{flat-equation-free-at-generic-point-fibre}).
置 $X_i = X_{T_i} = X \times_S T_i$，记 $\mathcal{F}_i$ 为
$\mathcal{F}$ 到 $X_i$ 的拉回。
记 $U \subset X_T$ 为 $\mathcal{F}_T$ 在 $T$ 上平坦的点所成的开子概形，见
《态射进阶》，定理 \ref{more-morphisms-theorem-openness-flatness}。
由假设，每个纤维的泛点都是 $U$ 中的点，即
$U \to T$ 是具有几何不可约纤维的光滑满射。可稍微缩小 $U$，并假设 $U$ 是拟紧的。
利用《极限》引理 \ref{limits-lemma-descend-opens}，
可找到 $i \in I$ 及拟紧开集 $U_i \subset X_i$，
使其在 $X_T$ 中的逆像为 $U$。增大 $i$ 后可假设
$\mathcal{F}_i|_{U_i}$ 在 $T_i$ 上平坦，见《极限》引理
\ref{limits-lemma-descend-module-flat-finite-presentation}。
特别地，$\mathcal{F}_i|_{U_i}$ 有限局部自由，
因此定义局部常值的
秩函数 $\rho : U_i \to \{0, 1, 2, \ldots \}$。
令 $(U_i)_p \subset U_i$ 表示 $\rho$ 取值为 $p$ 的开闭子集。令
$V_i \subset T_i$ 为 $(U_i)_p$ 的像；注意 $V_i$ 是开且拟紧的。
由假设，$T \to T_i$ 的像包含于 $V_i$ 中。
因此由《极限》引理 \ref{limits-lemma-descend-opens}，
存在 $i' \geq i$ 使 $T_{i'} \to T_i$ 经由 $V_i$ 分解。
于是 $\mathcal{F}_{i'}$ 满足
(\ref{flat-equation-free-at-generic-point-fibre})，如所需。细节略去。
\end{proof}

\begin{lemma}
\label{flat-lemma-pre-flat-dimension-n}
设 $f : X \to S$ 是局部有限型的概形态射。
设 $\mathcal{F}$ 是有限型拟凝聚 $\mathcal{O}_X$-模。设 $n \geq 0$。
以下条件等价：
\begin{enumerate}
\item 对于 $s \in S$，令 $Z \subset X_s$ 为 $\mathcal{F}$ 在 $S$ 上不平坦的点所成的闭子集（见引理
\ref{flat-lemma-open-in-fibre-where-flat}），则 $\dim(Z) < n$，并且
\item 对于 $x \in X$，若 $\mathcal{F}$ 在 $x$ 处不是在 $S$ 上平坦的，则
$\text{trdeg}_{\kappa(f(x))}(\kappa(x)) < n$。
\end{enumerate}
若这些条件成立，则在任意基变换后仍成立。
\end{lemma}

\begin{proof}
设 $x \in X$ 是位于 $s \in S$ 之上的点。
由《簇》引理 \ref{varieties-lemma-dimension-locally-algebraic}，
$\{x\}$ 在 $X_s$ 中的闭包的维数为
$\text{trdeg}_{\kappa(s)}(\kappa(x))$。
反之，若 $Z \subset X_s$ 是维数为 $d$ 的闭子集，则存在 $x \in Z$，
使 $\text{trdeg}_{\kappa(s)}(\kappa(x)) = d$（同一引理）。
因此 (1) 与 (2) 等价（甚至逐纤维成立）。
基变换的断言由《态射》引理
\ref{morphisms-lemma-base-change-module-flat} 和
\ref{morphisms-lemma-dimension-fibre-after-base-change} 得出。
\end{proof}

\begin{definition}
\label{flat-definition-flat-dimension-n}
设 $f : X \to S$ 是局部有限型的概形态射。
设 $\mathcal{F}$ 是有限型拟凝聚 $\mathcal{O}_X$-模。设 $n \geq 0$。
若引理 \ref{flat-lemma-pre-flat-dimension-n} 的等价条件成立，
则称 {\it $\mathcal{F}$ 在维数 $\geq n$ 上关于 $S$ 平坦}。
\end{definition}

\begin{situation}
\label{flat-situation-flat-dimension-n}
设 $f : X \to S$ 是局部有限型的概形态射。
设 $\mathcal{F}$ 是有限型拟凝聚 $\mathcal{O}_X$-模。对任意
$S$ 上的概形 $T$，记 $\mathcal{F}_T$ 为 $\mathcal{F}$ 到 $T$ 的基变换，
换言之，$\mathcal{F}_T$ 是沿投影态射
$X_T = X \times_S T \to X$ 的 $\mathcal{F}$ 的拉回。注意 $X_T \to T$ 是有限型的，
且 $\mathcal{F}_T$ 是有限型 $\mathcal{O}_{X_T}$-模（《态射》引理
\ref{morphisms-lemma-base-change-finite-type} 及《模》引理
\ref{modules-lemma-pullback-finite-type}）。
设 $n \geq 0$。由定义 \ref{flat-definition-flat-dimension-n} 与引理
\ref{flat-lemma-pre-flat-dimension-n} 得到函子
\begin{equation}
\label{flat-equation-flat-dimension-n}
F_n : (\Sch/S)^{opp} \longrightarrow \textit{Sets}, \quad
T \longrightarrow \left\{
\begin{matrix}
\{*\} & \text{若 }\mathcal{F}_T\text{ 在 }T\text{ 上于 }\dim \geq n\text{ 的部分平坦}, \\
\emptyset & \text{否则。}
\end{matrix}
\right.
\end{equation}
\end{situation}

\begin{lemma}
\label{flat-lemma-flat-dimension-n}
在情形 \ref{flat-situation-flat-dimension-n} 中。
\begin{enumerate}
\item 函子 $F_n$ 对 fpqc 拓扑满足层性质。
\item 若 $f$ 拟紧且局部有限表示，并且 $\mathcal{F}$ 有限表示，
则函子 $F_n$ 保持极限。
\end{enumerate}
\end{lemma}

\begin{proof}
设 $\{T_i \to T\}_{i \in I}$ 是 $S$ 上概形的 fpqc 覆盖。
置 $X_i = X_{T_i} = X \times_S T_i$，记 $\mathcal{F}_i$ 为
$\mathcal{F}$ 到 $X_i$ 的拉回。假设对所有 $i$，$\mathcal{F}_i$
在维数 $\geq n$ 上关于 $T_i$ 平坦。取 $t \in T$。
选取指标 $i$ 及映到 $t$ 的点 $t_i \in T_i$。
考虑笛卡尔图
$$
\xymatrix{
X_{\Spec(\mathcal{O}_{T, t})} \ar[d] &
X_{\Spec(\mathcal{O}_{T_i, t_i})} \ar[d] \ar[l] \\
\Spec(\mathcal{O}_{T, t}) &
\Spec(\mathcal{O}_{T_i, t_i}) \ar[l]
}
$$
由于下方水平态射平坦，由《态射进阶》引理
\ref{more-morphisms-lemma-flat-locus-base-change}，$\mathcal{F}_i$ 在
$T_i$ 上不平坦的集合 $Z_i \subset X_{t_i}$ 与 $\mathcal{F}_T$ 在 $T$ 上不平坦的
集合 $Z \subset X_t$ 满足关系 $Z_i = Z_{\kappa(t_i)}$。
于是由《态射》引理 \ref{morphisms-lemma-dimension-fibre-after-base-change}，
得到 $\mathcal{F}_T$ 在维数 $\geq n$ 上关于 $T$ 平坦。

\medskip\noindent
假设 $f$ 拟紧且局部有限表示，并且 $\mathcal{F}$ 有限表示。
本段先将 (2) 的证明约化到 $f$ 有限表示的情形。
设 $T = \lim_{i \in I} T_i$ 是仿射 $S$-概形的有向极限，
并假设 $\mathcal{F}_T$ 在维数 $\geq n$ 上平坦。置
$X_i = X_{T_i} = X \times_S T_i$，记 $\mathcal{F}_i$ 为 $\mathcal{F}$ 到 $X_i$ 的拉回。
我们需要证明对某个 $i$，$\mathcal{F}_i$ 在维数 $\geq n$ 上平坦。
取 $i_0 \in I$，并将 $I$ 换为 $\{i \mid i \geq i_0\}$。
由于 $T_{i_0}$ 仿射（因而拟紧），存在有限个仿射开集
$W_j \subset S$（$j = 1, \ldots, m$）以及一个仿射开覆盖
$T_{i_0} = \bigcup_{j = 1, \ldots, m} V_{j, i_0}$，
使得 $T_{i_0} \to S$ 将 $V_{j, i_0}$ 映入 $W_j$。
对 $i \geq i_0$，记 $V_{j, i}$ 为 $V_{j, i_0}$ 在 $T_i$ 中的逆像。
若能对每个 $j$ 证明存在某个 $i$ 使 $\mathcal{F}_{V_{j, i_0}}$
在维数 $\geq n$ 上平坦，则结论成立。这样便可约化到 $S$ 仿射的情形。
此时 $X$ 拟紧，可取有限仿射开覆盖 $X = W_1 \cup \ldots \cup W_m$。
此时 $(X \to S, \mathcal{F})$ 的结论等价于
$(\coprod W_j, \coprod \mathcal{F}|_{W_j})$ 的结论。因此可假设
$f$ 有限表示。

\medskip\noindent
假设 $f$ 有限表示且 $\mathcal{F}$ 有限表示。记 $U \subset X_T$ 为
$\mathcal{F}_T$ 在 $T$ 上平坦的点所成的开子概形，见《态射进阶》定理
\ref{more-morphisms-theorem-openness-flatness}。
由假设，$Z = X_T \setminus U$ 在 $T$ 上每个纤维的维数都为 $< n$。
由《极限》引理 \ref{limits-lemma-approximate-given-relative-dimension}，
可找到闭子概形 $Z \subset Z' \subset X_T$，使对所有 $t \in T$ 都有
$\dim(Z'_t) < n$，且 $Z' \to X_T$ 有限表示。由《极限》引理
\ref{limits-lemma-descend-finite-presentation} 和
\ref{limits-lemma-descend-closed-immersion-finite-presentation}，
存在 $i \in I$ 及有限表示闭子概形 $Z'_i \subset X_i$，
其到 $T$ 的基变换为 $Z'$。由《极限》引理
\ref{limits-lemma-limit-dimension}，可假设 $Z'_i \to T_i$ 的所有纤维维数都为 $< n$。
由《极限》引理 \ref{limits-lemma-descend-module-flat-finite-presentation}，
可假设 $\mathcal{F}_i|_{X_i \setminus T'_i}$ 在 $T_i$ 上平坦。
这蕴含 $\mathcal{F}_i$ 在维数 $\geq n$ 上平坦；这里使用
$Z' \to X_T$ 有限表示，因而补集 $X_T \setminus Z'$ 拟紧！
于是 (2) 得证，引理的证明完成。
\end{proof}

\begin{situation}
\label{flat-situation-flat}
设 $f : X \to S$ 是概形态射。
设 $\mathcal{F}$ 是拟凝聚 $\mathcal{O}_X$-模。
对任意 $S$ 上的概形 $T$，记 $\mathcal{F}_T$ 为 $\mathcal{F}$ 到 $T$ 的
基变换，换言之，$\mathcal{F}_T$ 是沿投影态射
$X_T = X \times_S T \to X$ 的 $\mathcal{F}$ 的拉回。由于平坦模的基变换
仍平坦，得到函子
\begin{equation}
\label{flat-equation-flat}
F_{flat} : (\Sch/S)^{opp} \longrightarrow \textit{Sets}, \quad
T \longrightarrow \left\{
\begin{matrix}
\{*\} & \text{若 } \mathcal{F}_T\text{ 在 }T\text{ 上平坦}, \\
\emptyset & \text{否则。}
\end{matrix}
\right.
\end{equation}
\end{situation}

\begin{lemma}
\label{flat-lemma-flat}
在情形 \ref{flat-situation-flat} 中。
\begin{enumerate}
\item 函子 $F_{flat}$ 对 fpqc 拓扑满足层性质。
\item 若 $f$ 拟紧且局部有限表示，并且 $\mathcal{F}$ 有限表示，
则函子 $F_{flat}$ 保持极限。
\end{enumerate}
\end{lemma}

\begin{proof}
第 (1) 部分由以下断言推出：若 $T' \to T$ 是 $S$ 上概形的满射平坦态射，
则 $\mathcal{F}_{T'}$ 在 $T'$ 上平坦，当且仅当 $\mathcal{F}_T$ 在 $T$ 上平坦，见
《态射进阶》引理 \ref{more-morphisms-lemma-flat-locus-base-change}。
第 (2) 部分在将 $X$、$S$ 约化到仿射情形后，由《极限》引理
\ref{limits-lemma-descend-module-flat-finite-presentation} 得出（比较引理
\ref{flat-lemma-flat-dimension-n} 的证明）。
\end{proof}




\section{扁平化分层}
\label{flat-section-flattening}

\noindent
这里只给出定义。寻找“泛平坦性分层”的读者应参阅
《态射进阶》第
\ref{more-morphisms-section-generic-flatness-stratification} 节。

\begin{definition}
\label{flat-definition-flattening}
设 $X \to S$ 是概形态射。
设 $\mathcal{F}$ 是拟凝聚 $\mathcal{O}_X$-模。
若情形 \ref{flat-situation-flat} 中定义的函子 $F_{flat}$
可由 $S$ 上的概形 $S'$ 表示，则称 {\it $\mathcal{F}$ 的泛扁平化存在}。
若 $\mathcal{O}_X$ 的泛扁平化存在，则称 {\it $X$ 的泛扁平化存在}。
\end{definition}

\noindent
注意，若 $\mathcal{F}$ 的泛扁平化 $S'$\footnote{概形 $S'$ 有时称为 {\it universal flatificator}；在 \cite{GruRay} 中称为 {\it platificateur universel}。泛扁平化的存在性不应与《代数进阶》第 \ref{more-algebra-section-blowup-flat} 节讨论的结果类型混淆。} 存在，
则态射 $S' \to S$ 是概形的满射单射，$\mathcal{F}_{S'}$ 在 $S'$ 上平坦，
并且态射 $T \to S$ 穿过 $S'$ 分解，当且仅当 $\mathcal{F}_T$ 在 $T$ 上平坦。

\begin{example}
\label{flat-example-no-universal-flattening}
设 $X = S = \Spec(k[x, y])$，其中 $k$ 是域。令
$\mathcal{F} = \widetilde{M}$，其中 $M = k[x, x^{-1}, y]/(y)$。
对 $k[x, y]$-代数 $A$，置 $F_{flat}(A) = F_{flat}(\Spec(A))$。
则对所有 $n$ 有 $F_{flat}(k[x, y]/(x, y)^n) = \{*\}$，但
$F_{flat}(k[[x, y]]) = \emptyset$。这意味着 $F_{flat}$ 不可表示
（甚至不能由代数空间表示，见《形式概形》引理
\ref{formal-spaces-lemma-map-into-algebraic-space}）。
因此在此情形泛扁平化不存在。
\end{example}

\noindent
我们定义（比较《拓扑》注记
\ref{topology-remark-locally-finite-stratification}）：概形 $S$ 的（局部有限的、
概形论的）{\it 分层} 由偏序集 $I$ 所指标的一族闭子概形 $Z_i \subset S$ 给出，
满足 $S = \bigcup Z_i$（集合论意义下），$S$ 的每个点都有一个邻域只与有限多个
$Z_i$ 相交，并且满足
$$
Z_i \cap Z_j = \bigcup\nolimits_{k \leq i, j} Z_k.
$$
置 $S_i = Z_i \setminus \bigcup_{j < i} Z_j$，实际的分层是将
该概形分解为局部闭子概形的 $S = \coprod S_i$。我们常只标出分层
$S_i$，把闭子概形 $Z_i$ 的构造留给读者。给定一个分层，可得到单射
$$
S' = \coprod\nolimits_{i \in I} S_i \longrightarrow S.
$$
称之为与该分层关联的 {\it 单射}。用此术语可定义扁平化分层的含义。

\begin{definition}
\label{flat-definition-flattening-stratification}
设 $X \to S$ 是概形态射。
设 $\mathcal{F}$ 是拟凝聚 $\mathcal{O}_X$-模。
若情形 \ref{flat-situation-flat} 中定义的函子 $F_{flat}$
可由单射 $S' \to S$ 表示，且该单射与由局部闭子概形给出的
$S$ 的一个分层关联，则称 $\mathcal{F}$ 有 {\it 扁平化分层}。
若 $\mathcal{O}_X$ 有扁平化分层，则称 $X$ 有 {\it 扁平化分层}。
\end{definition}

\noindent
当扁平化分层存在时，理解给分层编号的指标集及其偏序通常很重要。
这往往与模的秩有关。例如，当 $X = S$ 且 $\mathcal{F}$ 是有限表示的
$\mathcal{O}_S$-模时，扁平化分层存在，并由 $\mathcal{F}$ 的 Fitting 理想给出，见
《因子》引理 \ref{divisors-lemma-finite-presentation-module}。




\section{Artin 环上的扁平化分层}
\label{flat-section-flattening-artinian}

\noindent
当底概形是 Artin 环的谱时，扁平化分层存在。

\begin{lemma}
\label{flat-lemma-flattening-stratification-artinian}
设 $S$ 是 Artin 环的谱。
对 $S$ 上任意概形 $X$ 及任意拟凝聚 $\mathcal{O}_X$-模，
泛扁平化存在。事实上，泛扁平化由闭浸入 $S' \to S$ 给出，
因而也是 $\mathcal{F}$ 的扁平化分层。
\end{lemma}

\begin{proof}
取仿射开覆盖 $X = \bigcup U_i$。
则 $F_{flat}$ 是与各对 $(U_i, \mathcal{F}|_{U_i})$ 关联的函子的乘积。
因此只需对每个 $(U_i, \mathcal{F}|_{U_i})$ 证明结论。
在仿射情形，引理由《代数进阶》引理
\ref{more-algebra-lemma-flattening-artinian-universal-property} 立即推出。
\end{proof}






\section{态射的扁平化}
\label{flat-section-flattening-map}

\noindent
定理 \ref{flat-theorem-flattening-map}
是后续扁平化论断的关键。

\begin{lemma}
\label{flat-lemma-universally-separating}
设 $S$ 是概形。
设 $g : X' \to X$ 是 $S$ 上概形的平坦态射，且 $X$ 在 $S$ 上局部有限型。
设 $\mathcal{F}$ 是有限型拟凝聚 $\mathcal{O}_X$-模，并且在 $S$ 上平坦。
若 $\text{Ass}_{X/S}(\mathcal{F}) \subset g(X')$，则典范映射
$$
\mathcal{F} \longrightarrow g_*g^*\mathcal{F}
$$
是单射，并且在任意基变换后仍为单射。
\end{lemma}

\begin{proof}
最后的断言意味着，对任意态射 $T \to S$，映射
$\mathcal{F}_T \to (g_T)_*g_T^*\mathcal{F}_T$ 是单射。条件
$\text{Ass}_{X/S}(\mathcal{F}) \subset g(X')$ 在基变换下保持，见《因子》引理
\ref{divisors-lemma-base-change-relative-assassin} 及注记
\ref{divisors-remark-base-change-relative-assassin}。
平坦性与有限型的假设也具有同样性质。
因此只需证明所显示箭头的单射性。
令 $\mathcal{K} = \Ker(\mathcal{F} \to g_*g^*\mathcal{F})$。
目标是证明 $\mathcal{K} = 0$。
为此只需证明 $\text{WeakAss}_X(\mathcal{K}) = \emptyset$，见《因子》引理
\ref{divisors-lemma-weakly-ass-zero}。
我们有
$\text{WeakAss}_X(\mathcal{K}) \subset \text{WeakAss}_X(\mathcal{F})$，见《因子》引理
\ref{divisors-lemma-ses-weakly-ass}。
由于 $\mathcal{F}$ 平坦，由引理 \ref{flat-lemma-bourbaki-finite-type-general-base} 得
$\text{WeakAss}_X(\mathcal{F}) \subset \text{Ass}_{X/S}(\mathcal{F})$。
由假设，$\text{Ass}_{X/S}(\mathcal{F})$ 的任意点 $x$ 都是某个 $x' \in X'$ 的像。
由于 $g$ 平坦，局部环映射 $\mathcal{O}_{X, x} \to \mathcal{O}_{X', x'}$
忠实平坦，故映射
$$
\mathcal{F}_x
\longrightarrow
g^*\mathcal{F}_{x'} =
\mathcal{F}_x \otimes_{\mathcal{O}_{X, x}} \mathcal{O}_{X', x'}
$$
是单射（见《代数》引理
\ref{algebra-lemma-faithfully-flat-universally-injective}）。
这就蕴含 $\mathcal{K}_x = 0$，即得所需结论。
\end{proof}

\begin{lemma}
\label{flat-lemma-flattening-module-map}
设 $A$ 是环，$u : M \to N$ 是 $A$-模的满射。
若 $M$ 是投射 $A$-模，则存在理想 $I \subset A$，使对任意环映射
$\varphi : A \to B$，以下条件等价：
\begin{enumerate}
\item $u \otimes 1 : M \otimes_A B \to N \otimes_A B$ 是同构；
\item $\varphi(I) = 0$.
\end{enumerate}
\end{lemma}

\begin{proof}
由于 $M$ 投射，可取投射 $A$-模 $C$，使 $F = M \oplus C$ 是自由 $A$-模。
将 $u$ 替换为 $u \oplus 1 : F = M \oplus C \to N \oplus C$ 后，
可假设 $M$ 自由。在此情形，令 $I$ 为相对于 $M$ 的某个（固定）基，
$\Ker(u)$ 中所有元素的系数在 $A$ 中生成的理想。
这是因为 $u$ 满射且 $\otimes_A B$ 右正合，所以
$\Ker(u \otimes 1)$ 是 $\Ker(u) \otimes_A B$ 在 $M \otimes_A B$ 中的像。
\end{proof}

\begin{theorem}
\label{flat-theorem-flattening-map}
在情形
\ref{flat-situation-iso} 中，假设
\begin{enumerate}
\item $f$ 有限表示；
\item $\mathcal{F}$ 有限表示、在 $S$ 上平坦且相对于 $S$ 纯；
\item $u$ 是满射。
\end{enumerate}
则 $F_{iso}$ 可由闭浸入 $Z \to S$ 表示。
此外，若 $\mathcal{G}$ 有限表示，则 $Z \to S$ 有限表示。
\end{theorem}

\begin{proof}
下文不再特别说明，我们使用 $\mathcal{F}$ 在 $S$ 上普遍纯这一事实，见引理
\ref{flat-lemma-finite-type-flat-pure-along-fibre-is-universal}。
由引理 \ref{flat-lemma-iso-sheaf} 以及《下降》引理
\ref{descent-lemma-closed-immersion} 和
\ref{descent-lemma-descent-data-sheaves}，该问题对 $S$ 上的 \'etale 拓扑是局部的。
因此只需证明：给定 $s \in S$，存在 $(S, s)$ 的一个 \'etale 邻域，使定理成立。

\medskip\noindent
利用引理 \ref{flat-lemma-finite-presentation-flat-along-fibre}，
并将 $S$ 替换为 $s$ 的一个初等 \'etale 邻域后，可假设存在交换图
$$
\xymatrix{
X \ar[dr] & & X' \ar[ll]^g \ar[ld] \\
& S &
}
$$
其中概形均在 $S$ 上有限表示，且 $g$ 是 \'etale，$X_s \subset g(X')$，
概形 $X'$ 与 $S$ 仿射，并且
$\Gamma(X', g^*\mathcal{F})$ 是 $\Gamma(S, \mathcal{O}_S)$-投射模。
注意 $g^*\mathcal{F}$ 在 $S$ 上普遍纯，见引理
\ref{flat-lemma-affine-locally-projective-pure}。
因此由引理 \ref{flat-lemma-criterion}，开集 $g(X')$ 包含
位于 $\Spec(\mathcal{O}_{S, s})$ 之上的
$\text{Ass}_{X/S}(\mathcal{F})$ 中的点。置
$$
E = \{t \in S \mid \text{Ass}_{X_t}(\mathcal{F}_t) \subset g(X') \}.
$$
由《态射进阶》引理
\ref{more-morphisms-lemma-relative-assassin-constructible}，$E$ 是 $S$ 的可构造子集。
我们已知 $\Spec(\mathcal{O}_{S, s}) \subset E$。由《态射》引理
\ref{morphisms-lemma-constructible-containing-open}，$E$ 包含 $s$ 的一个开邻域。
因此将 $S$ 替换为 $s$ 的较小仿射邻域后，可假设
$\text{Ass}_{X/S}(\mathcal{F}) \subset g(X')$。

\medskip\noindent
由于假设 $u$ 是满射，有 $F_{iso} = F_{inj}$。由引理
\ref{flat-lemma-universally-separating}，$u : \mathcal{F} \to \mathcal{G}$ 是单射，
当且仅当 $g^*u : g^*\mathcal{F} \to g^*\mathcal{G}$ 是单射；在任意基变换后
同样成立。因此可约化到如下情形：除定理假设外，$X \to S$ 是仿射概形的态射，
且 $\Gamma(X, \mathcal{F})$ 是 $\Gamma(S, \mathcal{O}_S)$-投射模。
此情形立即由引理 \ref{flat-lemma-flattening-module-map} 得出。

\medskip\noindent
若 $\mathcal{G}$ 有限表示，则将引理
\ref{flat-lemma-iso-sheaf} 第 (4) 部分与《极限》注记
\ref{limits-remark-limit-preserving} 合并，即得 $Z$ 有限表示。
\end{proof}

\begin{lemma}
\label{flat-lemma-Weil-restriction-closed-subschemes}
设 $f:X\to S$ 是有限表示、平坦且纯的概形态射。设 $Y$ 是 $X$ 的闭子概形。
令 $F=f_*Y$ 为沿 $f$ 的 $Y$ 的 Weil 限制函子，定义为
$$
F : (\Sch/S)^{opp} \to \textit{Sets}, \quad
T \mapsto
\left\{
\begin{matrix}
\{*\} & \text{若} & Y_T\to X_T \text{ 是同构，}\\
\emptyset & \text{否则。} &
\end{matrix}
\right.
$$
则 $F$ 可由闭浸入 $Z\to S$ 表示。此外，若 $Y\to S$ 有限表示，
则 $Z\to S$ 有限表示。
\end{lemma}

\begin{proof}
令 $\mathcal{I}$ 是 $X$ 中定义 $Y$ 的理想层，并令
$u:\mathcal{O}_X\to\mathcal{O}_X/\mathcal{I}$ 为满射。
对 $S$-概形 $T$，闭浸入 $Y_T\to X_T$ 是同构，当且仅当
$u_T$ 是同构。因此结论由定理 \ref{flat-theorem-flattening-map} 得出。
\end{proof}



\section{局部情形下的扁平化}
\label{flat-section-flattening-local}

\noindent
本节开始应用前面的材料，得到扁平化分层的一个影子。

\begin{theorem}
\label{flat-theorem-flattening-local}
在情形 \ref{flat-situation-flat-at-point} 中，假设 $A$ 是 Hensel 环，
$B$ 是 $A$ 上本质有限型的，且 $M$ 是有限 $B$-模。
则存在理想 $I \subset A$，使 $A/I$ 在范畴 $\mathcal{C}$ 上共表示函子 $F_{lf}$。
换言之，给定局部环的局部同态
$\varphi : A \to A'$，其中 $B' = B \otimes_A A'$ 且 $M' = M \otimes_A A'$，
以下条件等价：
\begin{enumerate}
\item $\forall \mathfrak q \in V(\mathfrak m_{A'}B' + \mathfrak m_B B')
\subset \Spec(B') :
M'_{\mathfrak q}\text{ 在 }A'\text{ 上平坦}$，且
\item $\varphi(I) = 0$。
\end{enumerate}
若 $B$ 在 $A$ 上本质有限表示且 $M$ 在 $B$ 上有限表示，
则 $I$ 是有限生成理想。
\end{theorem}

\begin{proof}
取有限型环映射 $A \to C$、有限 $C$-模 $N$ 及 $C$ 的素理想
$\mathfrak q$，使 $B = C_{\mathfrak q}$ 且 $M = N_{\mathfrak q}$。
下文中，当我们说
“对 $(N/C/A, \mathfrak q)$ 定理成立”是指
是指它对 $(A \to B, M)$ 成立，其中 $B = C_{\mathfrak q}$ 且
$M = N_{\mathfrak q}$。由引理 \ref{flat-lemma-flat-at-point-go-up}，
若将 $B$ 替换为在 $B$ 上平坦的局部环，函子 $F_{lf}$ 不变。
因此由于 $A$ 是 Hensel 环，可应用引理 \ref{flat-lemma-existence-algebra}，
并假设 $N/C/A$ 在 $\mathfrak q$ 处存在完全 d\'evissage。

\medskip\noindent
令 $(A_i, B_i, M_i, \alpha_i, \mathfrak q_i)_{i = 1, \ldots, n}$
为 $N/C/A$ 在 $\mathfrak q$ 处的一个完全 d\'evissage。令
$\mathfrak q'_i \subset A_i$ 为位于 $\mathfrak q_i \subset B_i$ 之上的唯一素理想，
如定义 \ref{flat-definition-complete-devissage-at-x-algebra} 所述。
由于 $C \to A_1$ 满射且 $N \cong M_1$ 作为 $C$-模，由引理
\ref{flat-lemma-flat-at-point-finite}，只需证明定理对
$(M_1/A_1/A, \mathfrak q'_1)$ 成立。
由于 $B_1 \to A_1$ 有限，且 $\mathfrak q_1$ 是位于 $\mathfrak q'_1$ 之上的
$B_1$ 的唯一素理想，故 $(A_1)_{\mathfrak q'_1} \to (B_1)_{\mathfrak q_1}$ 有限
（见《代数》引理 \ref{algebra-lemma-unique-prime-over-localize-below} 或《态射进阶》引理
\ref{more-morphisms-lemma-finite-morphism-single-point-in-fibre}）。
因此由引理 \ref{flat-lemma-flat-at-point-finite}，只需证明定理对
$(M_1/B_1/A, \mathfrak q_1)$ 成立。

\medskip\noindent
此时可对 d\'evissage 的长度 $n$ 作归纳，假设定理对
$(M_2/B_2/A, \mathfrak q_2)$ 成立。（若 $n = 1$，则 $M_2 = 0$，
它在 $A$ 上平坦。）反向执行上一段的最后几个步骤，并使用
$M_2 \cong \Coker(\alpha_2)$ 作为 $B_1$-模这一事实，可知定理对
$(\Coker(\alpha_1)/B_1/A, \mathfrak q_1)$ 成立。

\medskip\noindent
设 $A'$ 是 $\mathcal{C}$ 的对象。此时应用引理
\ref{flat-lemma-induction-step}：若对 $B_1 \otimes_A A'$ 中位于
$\mathfrak m_{A'}$ 之上的素理想 $\mathfrak q'$，有
$(M_1 \otimes_A A')_{\mathfrak q'}$ 在 $A'$ 上平坦，则
$(\Coker(\alpha_1) \otimes_A A')_{\mathfrak q'}$ 在 $A'$ 上平坦。
因此 $F_{lf}$ 是与 $(B_1)_{\mathfrak q_1}$ 上的模
$\Coker(\alpha_1)_{\mathfrak q_1}$ 关联的函子 $F'_{lf}$ 的子函子。
由上一段，$F'_{lf}$ 由某个理想 $J \subset A$ 的商环 $A/J$ 共表示。
故可将 $A$ 替换为 $A/J$，并假设
$\Coker(\alpha_1)_{\mathfrak q_1}$ 在 $A$ 上平坦。

\medskip\noindent
由于 $\Coker(\alpha_1)$ 是 $B_1$-模，且 $N_1/B_1/A$ 在 $\mathfrak q_1$ 处
存在完全 d\'evissage，又由引理
\ref{flat-lemma-complete-devissage-flat-finite-type-module}，
$\Coker(\alpha_1)_{\mathfrak q_1}$ 在 $A$ 上平坦，
可知 $\Coker(\alpha_1)$ 是自由 $A$-模，特别地在 $A$ 上平坦。因此由引理
\ref{flat-lemma-induction-step}，$F_{lf}(A')$ 非空，当且仅当
$\alpha \otimes 1_{A'}$ 是单射。令 $N_1 = \Im(\alpha_1) \subset M_1$，于是有正合列
$$
0 \to N_1 \to M_1 \to \Coker(\alpha_1) \to 0
\quad\text{且}\quad
B_1^{\oplus r_1} \to N_1 \to 0
$$
由于 $\Coker(\alpha_1)$ 平坦，第一列在普遍意义下正合（见《代数》引理
\ref{algebra-lemma-flat-universally-injective}）。
因此 $\alpha \otimes 1_{A'}$ 是单射，当且仅当
$B_1^{\oplus r_1} \otimes_A A' \to N_1 \otimes_A A'$ 是同构。
最后，应用定理 \ref{flat-theorem-flattening-map}，可知该函子由某个理想
$I$ 的商环 $A/I$ 共表示，从而 $F_{lf}$ 也由 $A/I$ 共表示。

\medskip\noindent
为证明最后的断言，假设 $A \to B$ 本质有限表示，且 $M$ 在 $B$ 上有限表示。
令 $I \subset A$ 为使 $F_{lf}$ 由 $A/I$ 共表示的理想。
写成 $I = \bigcup I_\lambda$，其中 $I_\lambda$ 遍历包含于 $I$ 中的有限生成理想。
由于 $F_{lf}(A/I) = \{*\}$，由引理 \ref{flat-lemma-flat-at-point} 第 (2) 部分可知，
对某个 $\lambda$ 有 $F_{lf}(A/I_\lambda) = \{*\}$。
显然这蕴含 $I = I_\lambda$。
\end{proof}

\begin{remark}
\label{flat-remark-flattening-local-scheme-theoretic}
这是定理 \ref{flat-theorem-flattening-local} 的一个概形论重述。
设 $(X, x) \to (S, s)$ 是带点概形的态射，且该态射局部有限型。
设 $\mathcal{F}$ 是有限型拟凝聚 $\mathcal{O}_X$-模。
假设 $S$ 是闭点为 $s$ 的 Hensel 局部概形。
存在闭子概形 $Z \subset S$，具有如下性质：对任意带点概形态射
$(T, t) \to (S, s)$，以下两个条件等价：
\begin{enumerate}
\item 在纤维 $X_t$ 中映到 $x \in X_s$ 的所有点处，$\mathcal{F}_T$ 都在 $T$ 上平坦；
\item $\Spec(\mathcal{O}_{T, t}) \to S$ 经由 $Z$ 分解。
\end{enumerate}
此外，如果 $X \to S$ 在 $x$ 处有限表示，且 $\mathcal{F}_x$
作为 $\mathcal{O}_{X, x}$-模有限表示，那么 $Z \to S$ 有限表示。
\end{remark}

\noindent
现在可以完全免费地从上述结果得到一些非常一般的结果。注意，
最有趣的情形可能是 $E = X_s$！

\begin{lemma}
\label{flat-lemma-freebie}
设 $S$ 是闭点为 $s$ 的 Hensel 局部环的谱。
设 $X \to S$ 是局部有限型的概形态射。
设 $\mathcal{F}$ 是有限型拟凝聚 $\mathcal{O}_X$-模。
设 $E \subset X_s$ 是一个子集。存在闭子概形 $Z \subset S$，具有如下性质：
对任意带点概形态射 $(T, t) \to (S, s)$，以下两个条件等价：
\begin{enumerate}
\item 在纤维 $X_t$ 中映到 $E \subset X_s$ 中某点的所有点处，
$\mathcal{F}_T$ 都在 $T$ 上平坦；
\item $\Spec(\mathcal{O}_{T, t}) \to S$ 经由 $Z$ 分解。
\end{enumerate}
此外，如果 $X \to S$ 局部有限表示，$\mathcal{F}$ 有限表示，且
$E \subset X_s$ 是闭且拟紧的，那么 $Z \to S$ 有限表示。
\end{lemma}

\begin{proof}
对 $x \in X_s$，记在注记
\ref{flat-remark-flattening-local-scheme-theoretic} 中得到的闭子概形为 $Z_x \subset S$。
显然，$Z = \bigcap_{x \in E} Z_x$ 即可！

\medskip\noindent
为证明最后一个断言，假设 $X$ 局部有限表示，$\mathcal{F}$ 有限表示，
且 $Z$ 闭且拟紧。首先选取有限多个仿射开集 $W_j \subset X$，使得
$E \subset \bigcup W_j$。显然，只需对每个态射 $W_j \to S$、层
$\mathcal{F}|_{X_j}$ 以及闭子集 $E \cap W_j$ 证明该结果即可。因此不妨设 $X$ 仿射。
此时，《代数进阶》引理
\ref{more-algebra-lemma-limit-preserving-flat-at-primes} 表明，由 (1) 定义的函子
“保持极限”。于是，完全按照定理 \ref{flat-theorem-flattening-local} 证明的最后一部分，
可以证明 $Z \to S$ 有限表示。
\end{proof}

\begin{remark}
\label{flat-remark-flattening-complete-noetherian}
追溯引理 \ref{flat-lemma-freebie} 的证明直至其根源，会发现这是一条漫长而曲折的道路。
但如果假设
\begin{enumerate}
\item $f$ 有限型；
\item $\mathcal{F}$ 是有限型 $\mathcal{O}_X$-模；
\item $E = X_s$；以及
\item $S$ 是 Noether 完备局部环的谱。
\end{enumerate}
那么可以完全依靠更初等的代数给出如下证明：首先取有限仿射开覆盖，
将情形约化到 $X$ 仿射。在此情形，$Z$ 的存在性由《代数进阶》引理
\ref{more-algebra-lemma-flattening-complete-local-universal-property} 给出。
证明的关键步骤是在截断
$\Spec(\mathcal{O}_{S, s}/\mathfrak m_s^n)$ 中逐步构造闭子概形 $Z$。
这依赖于底为阿廷概形时扁平化分层总是存在这一事实，以及
$\mathcal{O}_{S, s} = \lim \mathcal{O}_{S, s}/\mathfrak m_s^n$ 这一事实。
\end{remark}




\section{一个引理的变体}
\label{flat-section-variants-mod-injective}

\noindent
本节讨论《代数》引理 \ref{algebra-lemma-mod-injective-general} 与
\ref{algebra-lemma-mod-injective} 的变体。
最一般的版本是命题
\ref{flat-proposition-finite-type-injective-into-flat-mod-m}；该命题曾作为
\cite[引理 4.2.2]{GruRay} 陈述，但其中的证明只给出了
引理 \ref{flat-lemma-upstairs-finite-type-injective-into-flat-mod-m} 所陈述的较弱结果。
命题 \ref{flat-proposition-finite-type-injective-into-flat-mod-m} 的精巧证明归功于 Ofer Gabber。
由于目前没有命题的应用，我们建议读者读完引理
\ref{flat-lemma-upstairs-finite-type-injective-into-flat-mod-m} 的证明后跳到下一节；
下一节将用此引理证明定理
\ref{flat-theorem-check-flatness-at-associated-points}。

\begin{situation}
\label{flat-situation-mod-injective}
设 $\varphi : A \to B$ 是局部环之间的局部环同态，且本质有限型。
设 $M$ 是平坦 $A$-模，$N$ 是有限 $B$-模，且
$u : N \to M$ 是一个 $A$-模映射，满足
$\overline{u} : N/\mathfrak m_AN \to M/\mathfrak m_AM$ 是单射。
\end{situation}

\noindent
在此情形，我们的目标是证明 $u$ 是 $A$-普遍单射，
$N$ 作为 $B$-模有限表示，并且 $N$ 作为 $A$-模平坦。
若这些结论成立，就称该情形中{\it 引理成立}。

\begin{lemma}
\label{flat-lemma-Noetherian-finite-type-injective-into-flat-mod-m}
若在情形 \ref{flat-situation-mod-injective} 中环 $A$ 是 Noether 环，则引理成立。
\end{lemma}

\begin{proof}
应用《代数》引理 \ref{algebra-lemma-mod-injective}，可知 $u$ 是单射且
$N/u(M)$ 在 $A$ 上平坦。于是 $u$ 是 $A$-普遍单射（《代数》引理
\ref{algebra-lemma-flat-tor-zero}），并且 $N$ 在 $A$ 上平坦（《代数》引理
\ref{algebra-lemma-flat-ses}）。此时 $B$ 是 Noether 环，故 $N$ 有限表示。
\end{proof}

\begin{lemma}
\label{flat-lemma-reduce-finite-type-injective-into-flat-mod-m}
设 $A_0$ 是局部环。若对每个满足下列条件的情形
\ref{flat-situation-mod-injective}，引理均成立：$A = A_0$，$B$ 是 $A$ 上多项式代数的
一个局部化，且 $N$ 作为 $B$-模有限表示；那么对每个满足 $A = A_0$ 的情形
\ref{flat-situation-mod-injective}，引理都成立。
\end{lemma}

\begin{proof}
设 $A \to B$、$u : N \to M$ 如情形 \ref{flat-situation-mod-injective}。
写成 $B = C/I$，其中 $C$ 是在某个素理想处对 $A$ 上多项式代数作的局部化。
如果能证明 $N$ 作为 $C$-模有限表示，那么当然也能证明 $N$ 作为 $B$-模有限表示（见
《代数》引理 \ref{algebra-lemma-finitely-presented-over-subring}）。
因此不妨设 $B$ 是多项式代数的局部化。
接着写 $N = B^{\oplus n}/K$，其中 $K \subset B^{\oplus n}$ 是某个子模。
由于 $B/\mathfrak m_AB$ 是 Noether 环（它本质有限型于一个域），存在有限多个元素
$k_1, \ldots, k_s \in K$，使得令 $K' = \sum Bk_i$、$N' = B^{\oplus n}/K'$ 时，
典范满射 $N' \to N$ 在商模上诱导同构
$N'/\mathfrak m_AN' \cong N/\mathfrak m_AN$。
现在，若引理对复合映射 $u' : N' \to M$ 成立，则 $u'$ 是单射，
从而 $N' = N$ 且 $u' = u$。因此原情形中的引理也成立。
\end{proof}

\begin{lemma}
\label{flat-lemma-henselian-finite-type-injective-into-flat-mod-m}
若在情形 \ref{flat-situation-mod-injective} 中环 $A$ 是 Hensel 的，
则引理成立。
\end{lemma}

\begin{proof}
只需在 $B$ 本质有限表示于 $A$ 且 $N$ 作为 $B$-模有限表示时证明即可，见
引理 \ref{flat-lemma-reduce-finite-type-injective-into-flat-mod-m}。
暂时再假设 $N$ 在 $A$ 上平坦。于是 $N$ 是自由 $A$-模 $F_i$ 的滤余极限
$N = \colim_i F_i$，其中转移映射
$u_{ii'} : F_i \to F_{i'}$ 模去 $\mathfrak m_A$ 后为单射，见
引理 \ref{flat-lemma-flat-finite-type-local-colimit-free}。
由引理 \ref{flat-lemma-universally-injective-local}，每个复合映射
$u_i : F_i \to M$ 都是 $A$-普遍单射，
因此 $u = \colim u_i$ 如所需是 $A$-普遍单射。

\medskip\noindent
假设 $A$ 是 Hensel 局部环，$B$ 本质有限表示于 $A$，且 $N$ 作为 $B$-模有限表示。
由定理 \ref{flat-theorem-flattening-local}，存在有限生成理想 $I \subset A$，使得
$N/IN$ 在 $A/I$ 上平坦，并且除非 $I = 0$，否则 $N/I^2N$ 不在 $A/I^2$ 上平坦。
上一段的结果表明，引理对 $A/I$ 上的
$u \bmod I : N/IN \to M/IM$ 成立。
考虑交换图
$$
\xymatrix{
0 \ar[r] &
M \otimes_A I/I^2 \ar[r] &
M/I^2M \ar[r] &
M/IM \ar[r] & 0 \\
&
N \otimes_A I/I^2 \ar[r] \ar[u]^u &
N/I^2N \ar[r] \ar[u]^u &
N/IN \ar[r] \ar[u]^u & 0
}
$$
由于 $\otimes$ 的右正合性以及 $M$ 在 $A$ 上平坦，图中的两行正合。
注意左侧竖直箭头是映射
$N/IN \otimes_{A/I} I/I^2 \to M/IM \otimes_{A/I} I/I^2$，因而是单射。
作图追逐可知左下方箭头是单射，即
$\text{Tor}^1_{A/I^2}(I/I^2, M/I^2) = 0$；见
《代数》注记 \ref{algebra-remark-Tor-ring-mod-ideal}。
于是由《代数》引理 \ref{algebra-lemma-what-does-it-mean}，$N/I^2N$ 在 $A/I^2$ 上平坦，
除非 $I = 0$，这便矛盾。
\end{proof}

\noindent
下面的引理讨论情形 \ref{flat-situation-mod-injective} 的特殊情形：$M$ 带有
$B$-模结构且 $u$ 是 $B$-线性的。这是实践中最常用的情形，
证明也比一般情形容易得多。

\begin{lemma}
\label{flat-lemma-upstairs-finite-type-injective-into-flat-mod-m}
设 $A \to B$ 是局部环之间本质有限型的局部环同态，且 $u : N \to M$ 是
$B$-模映射。若 $N$ 是有限 $B$-模，$M$ 在 $A$ 上平坦，并且
$\overline{u} : N/\mathfrak m_A N \to M/\mathfrak m_A M$ 是单射，
那么 $u$ 是 $A$-普遍单射，$N$ 作为 $B$-模有限表示，且 $N$ 在 $A$ 上平坦。
\end{lemma}

\begin{proof}
设 $A \to A^h$ 是 $A$ 的 Hensel 化。令 $B'$ 为
$B \otimes_A A^h$ 在极大理想
$\mathfrak m_B \otimes A^h + B \otimes \mathfrak m_{A^h}$ 处的局部化。
由于 $B \to B'$ 平坦（因而忠实平坦，见《代数》引理
\ref{algebra-lemma-local-flat-ff}），我们可以将 $A \to B$ 替换为
$A^h \to B'$，将模 $M$ 替换为 $M \otimes_B B'$，将模 $N$ 替换为
$N \otimes_B B'$，并将 $u$ 替换为 $u \otimes \text{id}_{B'}$；见《代数》引理
\ref{algebra-lemma-descend-properties-modules} 与
\ref{algebra-lemma-flatness-descends-more-general}。
因此不妨设 $A$ 是 Hensel 局部环。在此情形，我们的引理由更一般的
引理 \ref{flat-lemma-henselian-finite-type-injective-into-flat-mod-m} 推出。
\end{proof}

\begin{lemma}
\label{flat-lemma-valuation-ring-finite-type-injective-into-flat-mod-m}
若在情形 \ref{flat-situation-mod-injective} 中环 $A$ 是赋值环，则引理成立。
\end{lemma}

\begin{proof}
回忆一下，$A$-模平坦当且仅当无挠，见《代数进阶》引理
\ref{more-algebra-lemma-valuation-ring-torsion-free-flat}。
令 $T \subset N$ 为 $A$-挠。则 $u(T) = 0$，且 $N/T$ 在 $A$ 上平坦。
因此 $N/T$ 作为 $B$-模有限表示，见《代数进阶》引理
\ref{more-algebra-lemma-flat-finite-type-valuation-ring-finite-presentation}。
于是 $T$ 是有限 $B$-模，见《代数》引理 \ref{algebra-lemma-extension}。
由于 $N/T$ 在 $A$ 上平坦，我们有
$T/\mathfrak m_A T \subset N/\mathfrak m_A N$，见《代数》引理
\ref{algebra-lemma-flat-tor-zero}。
因为 $\overline{u}$ 是单射而 $u(T) = 0$，可得 $T/\mathfrak m_A T = 0$。
由 Nakayama 引理，$T = 0$，见《代数》引理 \ref{algebra-lemma-NAK}。
至此三个断言中的两个已证得（$N$ 在 $A$ 上平坦且作为 $B$-模有限表示），
剩下的是证明 $u$ 普遍单射。

\medskip\noindent
由《代数》定理 \ref{algebra-theorem-universally-exact-criteria}，只需证明对每个有限表示
$A$-模 $Q$，映射 $N \otimes_A Q \to M \otimes_A Q$ 是单射。由《代数进阶》引理
\ref{more-algebra-lemma-generalized-valuation-ring-modules}，不妨设
$Q = A/(a)$，其中 $a \in \mathfrak m_A$ 非零。因此只需证明
$N/aN \to M/aM$ 是单射。
设 $x \in N$ 且 $u(x) \in aM$。由引理
\ref{flat-lemma-flat-finite-type-local-valuation-ring-has-content}，$x$ 有一个内容理想
$I \subset A$。由于 $I$ 有限生成（《代数进阶》引理
\ref{more-algebra-lemma-content-finitely-generated}），且 $A$ 是赋值环，
有 $I = (b)$，其中某个 $b$（由《代数》引理
\ref{algebra-lemma-characterize-valuation-ring}）。由《代数进阶》引理
\ref{more-algebra-lemma-equal-content}，元素 $u(x)$ 也有内容理想 $I$。
由于 $u(x) \in aM$，由《代数进阶》定义 \ref{more-algebra-definition-content-ideal} 可知
$(b) \subset (a)$。又因 $x \in bN$，遂有 $x \in aN$，即得所需。
\end{proof}

\noindent
考虑以下情形
\begin{equation}
\label{flat-equation-star}
\begin{matrix}
A \to B\text{ 为有限表示，}S \subset B
\text{ 是乘法闭子集，并且}\\
N\text{ 是有限表示的 }S^{-1}B\text{-模}
\end{matrix}
\end{equation}
在此情形，{\it 纯展开}是指一个仿射开集
$U \subset \Spec(B)$，满足 $\Spec(S^{-1}B) \subset U$，以及一个延拓 $N$ 的
有限表示 $\mathcal{O}(U)$-模 $N'$，使得 $N'$ 是 $A$-投射的，且
$N' \to N = S^{-1}N'$ 是 $A$-普遍单射。

\medskip\noindent
在 (
\ref{flat-equation-star}) 中，若 $A \to A_1$ 是环映射，则可以作基变换：取
$B_1 = B \otimes_A A_1$，令 $S_1 \subset B_1$ 为 $S$ 的像，并令
$N_1 = N \otimes_A A_1$。这是可行的，因为
$S_1^{-1}B_1 = S^{-1}B \otimes_A A_1$。下文将不再特别说明而使用这一点。

\begin{lemma}
\label{flat-lemma-properties-pure-spreadout}
在 (\ref{flat-equation-star}) 中，若存在纯展开，则
\begin{enumerate}
\item $N$ 的元素在 $A$ 中都有内容理想；以及
\item 若 $u : N \to M$ 是到平坦 $A$-模 $M$ 的态射，且对 $A$ 的每个极大理想
$\mathfrak m$，映射 $N/\mathfrak m N \to M/\mathfrak m M$ 都是单射，
则 $u$ 是 $A$-普遍单射。
\end{enumerate}
\end{lemma}

\begin{proof}
取纯展开定义中的 $U$、$N'$。任意 $x' \in N'$ 在 $A$ 中都有内容理想，
因为 $N'$ 是 $A$-投射的（这可以直接看出，也可由《代数进阶》引理
\ref{more-algebra-lemma-content-exists-flat-Mittag-Leffler} 与《代数》例
\ref{algebra-example-ML} 推出）。由于 $N' \to N$ 是 $A$-普遍单射，
任意 $x' \in N'$ 的像 $x \in N$ 也在 $A$ 中有内容理想
（它与 $x'$ 的内容理想相同）。对一般的 $x \in N$，选取 $s \in S$ 使得
$s x$ 属于 $N' \to N$ 的像，并利用 $x$ 与 $sx$ 具有相同内容理想这一点。

\medskip\noindent
令 $u : N \to M$ 如 (2) 中所述。为证明 $u$ 是 $A$-普遍单射，
可将 $A$ 替换为某个极大理想处的局部化（略去一个小细节）。
假设 $A$ 是极大理想为 $\mathfrak m$ 的局部环。取 $s \in S$，考虑复合映射
$$
N' \to N \xrightarrow{1/s} N \xrightarrow{u} M
$$
这些映射中的每一个模去 $\mathfrak m$ 后都是单射，因此由引理
\ref{flat-lemma-universally-injective-local}，其复合是 $A$-普遍单射。
由于 $N = \colim_{s \in S} (1/s)N'$，作为普遍单射映射的余极限，$u$
是 $A$-“逆向单射”的。
% P05-EN-CH38-0047: frozen source line 7266 says “inversally injective” after universally-injective component maps; candidate only, do not silently repair.
\end{proof}

\begin{lemma}
\label{flat-lemma-find-pure-spreadout}
在 (\ref{flat-equation-star}) 中，对每个 $\mathfrak p \in \Spec(A)$，
存在有限生成理想 $I \subset \mathfrak pA_\mathfrak p$，使得在 $A_\mathfrak p/I$ 上有纯展开。
\end{lemma}

\begin{proof}
可将 $A$ 替换为 $A_\mathfrak p$。因此不妨设 $A$ 是局部环，且
$\mathfrak p$ 是 $A$ 的极大理想 $\mathfrak m$。通过清除 $S^{-1}B$ 上 $N$ 的表示中的分母，
可写成 $N = S^{-1}N'$，其中 $N'$ 是有限表示 $B$-模。
由于 $B/\mathfrak m B$ 是 Noether 环，映射
$N'/\mathfrak m N' \to N/\mathfrak m N$ 的核 $K$ 有限生成。
于是可取 $s \in S$ 使 $K$ 被 $s$ 消灭。将 $B$ 替换为 $B_s$ 后（这只相当于
经过 $\Spec(B)$ 的一个仿射开子概形，因而是允许的），可知 $S$ 的元素在
$N'/\mathfrak m N'$ 上作用为单射。此时选取局部子环 $A_0 \subset A$，使其本质有限型于
$\mathbf{Z}$；选取有限型环映射 $A_0 \to B_0$，使得
$B = A \otimes_{A_0} B_0$；再选取有限 $B_0$-模 $N'_0$，使得
$N' = B \otimes_{B_0} N'_0 = A \otimes_{A_0} N'_0$。
断言 $I = \mathfrak m_{A_0} A$ 可行。确实有
$$
N'/IN' = N'_0/\mathfrak m_{A_0} N'_0 \otimes_{\kappa_{A_0}} A/I
$$
该模在 $A/I$ 上自由。$S$ 的元素模去极大理想后作用为单射，
因此例如由引理 \ref{flat-lemma-invert-universally-injective}，
$N'/IN' \to N/IN$ 是普遍单射。
\end{proof}

\begin{lemma}
\label{flat-lemma-universally-injective-if-flat}
在 (\ref{flat-equation-star}) 中，假设 $N$ 在 $A$ 上平坦，$M$ 是平坦 $A$-模，
且 $u : N \to M$ 是 $A$-模映射，满足对所有
$\mathfrak p \in \Spec(A)$，$u \otimes \text{id}_{\kappa(\mathfrak p)}$ 都是单射。
则 $u$ 是 $A$-普遍单射。
\end{lemma}

\begin{proof}
由《代数》引理 \ref{algebra-lemma-check-universally-injective-into-flat}，
只需检查对每个理想 $I \subset A$，映射 $N/IN \to M/IM$ 是单射。
将 $A$ 替换为 $A/I$ 后，可见只需证明 $u$ 是单射。

\medskip\noindent
证明 $u$ 为单射。设 $x \in N$ 是 $u$ 核中的非零元素。则模 $Ax$ 有一个弱伴随素理想
$\mathfrak p$，见《代数》引理 \ref{algebra-lemma-weakly-ass-zero}。
将 $A$ 替换为 $A_\mathfrak p$，不妨设 $A$ 局部，并可找到非零元素
$y \in Ax$，使其湮灭子的根基等于 $\mathfrak m_A$，见
《代数》引理 \ref{algebra-lemma-weakly-ass-local}。
因此 $\text{Supp}(y) \subset \Spec(S^{-1}B)$ 非空，且包含于
$\Spec(S^{-1}B) \to \Spec(A)$ 的闭纤维中。
取 $I \subset \mathfrak m_A$ 为有限生成理想，使得在 $A/I$ 上有纯展开，见
引理 \ref{flat-lemma-find-pure-spreadout}。于是对某个 $n$ 有 $I^n y = 0$。由平坦性，
$y \in \text{Ann}_M(I^n) = \text{Ann}_A(I^n) \otimes_R N$。
因此，为得到所需矛盾，只需证明
$$
\text{Ann}_A(I^n) \otimes_R N
\longrightarrow
\text{Ann}_A(I^n) \otimes_R M
$$
是单射。由于 $N$、$M$ 都平坦，且 $\text{Ann}_A(I^n)$ 被 $I^n$ 消灭，
只需证明对每个被 $I$ 消灭的 $A$-模 $Q$，映射
$Q \otimes_A N \to Q \otimes_A M$ 是单射。这由 $I$ 的选取以及
引理 \ref{flat-lemma-properties-pure-spreadout} 的 (2) 得到。
\end{proof}

\begin{lemma}
\label{flat-lemma-big-intersection-is-zero}
设 $A$ 是非域的局部整环。设 $S$ 是 $A$ 的有限生成理想组成的集合。
假设 $S$ 对乘积封闭，且
$\bigcup_{I \in S} V(I)$ 是 $\Spec(A)$ 的泛点的补集。
则 $\bigcap_{I \in S} I = (0)$。
\end{lemma}

\begin{proof}
由于 $\mathfrak m_A \subset A$ 不是 $\Spec(A)$ 的泛点，可知至少有一个 $I \in S$
满足 $I \subset \mathfrak m_A$。故 $\bigcap_{I \in S} I \subset \mathfrak m_A$。
设 $f \in \mathfrak m_A$ 非零。则
$V(f) \subset \bigcup_{I \in S} V(I)$。
由于 $V(f)$ 上的可构造拓扑拟紧（拓扑学引理
\ref{topology-lemma-constructible-hausdorff-quasi-compact} 以及
《代数》引理 \ref{algebra-lemma-spec-spectral}），可得
$V(f) \subset V(I_1) \cup \ldots \cup V(I_n)$，其中某些 $I_j \in S$。
因为 $I_1 \ldots I_n \in S$，可知对某个 $I$ 有 $V(f) \subset V(I)$。
由于 $I$ 有限生成，这意味着对某个 $m$ 有 $I^m \subset (f)$；又因 $S$ 对乘积封闭，
可知对某个 $I \in S$ 有 $I \subset (f^2)$。于是不能有 $f \in I$。
\end{proof}

\begin{lemma}
\label{flat-lemma-closed-points-complement}
设 $A$ 是局部环，设 $I, J \subset A$ 是理想。
若 $J$ 有限生成且对所有 $n \geq 1$ 都有 $I \subset J^n$，
则 $V(I)$ 包含 $\Spec(A) \setminus V(J)$ 的闭点。
\end{lemma}

\begin{proof}
设 $\mathfrak p \subset A$ 是 $\Spec(A) \setminus V(J)$ 的闭点。
我们要证明 $I \subset \mathfrak p$。否则，某个 $f \in I$ 在 $A/\mathfrak p$ 中的像非零。
注意 $V(J) \cap \Spec(A/\mathfrak p)$ 是所有非泛点组成的集合。
因此将引理 \ref{flat-lemma-big-intersection-is-zero} 应用于理想族
$J^nA/\mathfrak p$，可知 $f$ 在 $A/\mathfrak p$ 中的像为零。
\end{proof}

\begin{lemma}
\label{flat-lemma-make-smaller-flatness-ideal}
设 $A$ 是局部环，设 $I \subset A$ 是理想。
设 $U \subset \Spec(A)$ 是拟紧开集。
设 $M$ 是 $A$-模。假设
\begin{enumerate}
\item $M/IM$ 在 $A/I$ 上平坦；
\item $M$ 在 $U$ 上平坦；
\end{enumerate}
则 $M/I_2M$ 在 $A/I_2$ 上平坦，其中
$I_2 = \Ker(I \to \Gamma(U, I/I^2))$。
\end{lemma}

\begin{proof}
只需证明
$M \otimes_A I/I_2 \to IM/I_2M$ 是单射，见《代数》引理
\ref{algebra-lemma-what-does-it-mean-again}。
由假设 (2)，该映射在 $U$ 上是单射。因此只需证明
$M \otimes_A I/I_2$ 到其在 $U$ 上的截面的映射是单射。
我们有 $M \otimes_A I/I_2 = M/IM \otimes_A I/I_2$，且
$M/IM$ 是有限自由 $A/I$-模的滤余极限
（《代数》定理 \ref{algebra-theorem-lazard}）。
所以只需证明 $I/I_2$ 到其在 $U$ 上截面的映射是单射，
这由 $I_2$ 的构造立即得到。
\end{proof}

\begin{proposition}
\label{flat-proposition-finite-type-injective-into-flat-mod-m}
设 $A \to B$ 是局部环之间本质有限型的局部环同态。设 $M$ 是平坦 $A$-模，
$N$ 是有限 $B$-模，且 $u : N \to M$ 是 $A$-模映射，满足
$\overline{u} : N/\mathfrak m_AN \to M/\mathfrak m_AM$ 是单射。
则 $u$ 是 $A$-普遍单射，$N$ 作为 $B$-模有限表示，且 $N$ 在 $A$ 上平坦。
\end{proposition}

\begin{proof}
不妨设 $B$ 是有限表示 $A$-代数 $B_0$ 的局部化，且 $N$ 是有限表示
$B_0$-模 $M_0$ 的局部化，见引理
\ref{flat-lemma-reduce-finite-type-injective-into-flat-mod-m}。
由《态射进阶》引理
\ref{more-morphisms-lemma-generic-flatness-stratification}，在
$\Spec(B_0)$ 上相对于 $\Spec(A)$ 存在 $\widetilde{M_0}$ 的“泛平坦性分层”。
翻译回 $N$，得到一列闭子概形
$$
S = \Spec(A) \supset S_0 \supset S_1 \supset \ldots \supset S_t = \emptyset
$$
其中 $S_i \subset S$ 由 $A$ 的有限生成理想截出，并且
$\widetilde{N}$ 拉回到
$\Spec(B) \times_S (S_i \setminus S_{i + 1})$ 后，在
$S_i \setminus S_{i + 1}$ 上平坦。我们对 $t$ 作归纳来证明命题
（基例 $t = 1$ 将与其他步骤并行证明）。令 $\Spec(A/J_i)$ 为
$S_i \setminus S_{i + 1}$ 的概形论闭包。

\medskip\noindent
{\bf 断言 1。} $N/J_iN$ 在 $A/J_i$ 上平坦。对 $i = t - 1$ 这是直接的，
而对 $i > 0$ 则由归纳假设得到。因此不妨设 $t > 1$、$S_{t - 1} \not = \emptyset$，
且 $J_0 = 0$；我们需要证明 $N$ 平坦。令 $J \subset A$ 为定义 $S_1$ 的理想。
再次对 $t$ 归纳，我们也知道模去 $J$ 的幂后平坦。令 $A^h$ 为 $A$ 的 Hensel 化，
令 $B'$ 为 $B \otimes_A A^h$ 在极大理想
$\mathfrak m_B \otimes A^h + B \otimes \mathfrak m_{A^h}$ 处的局部化。
则 $B \to B'$ 忠实平坦。置 $N' = N \otimes_B B'$。注意，$N'$ 在 $A^h$ 上平坦
当且仅当 $N$ 在 $A$ 上平坦。由定理 \ref{flat-theorem-flattening-local}，存在最小理想
$I \subset A^h$，使得 $N'/IN'$ 在 $A^h/I$ 上平坦，且 $I$ 有限生成。
由上面的结论，对所有 $n \geq 1$ 都有 $I \subset J^nA^h$。令
$S_i^h \subset \Spec(A^h)$ 为 $S_i \subset \Spec(A)$ 的逆像。由引理
\ref{flat-lemma-closed-points-complement}，$V(I)$ 包含
$U = \Spec(A^h) - S_1^h$ 的闭点。按构造，$N'$ 在 $U$ 上对 $A^h$ 平坦。
由引理 \ref{flat-lemma-make-smaller-flatness-ideal}，$N'/I_2N'$ 在 $A/I_2$ 上平坦，其中
$I_2 = \Ker(I \to \Gamma(U, I/I^2))$。由 $I$ 的极小性，$I = I_2$。
这意味着 $I = I^2$ 在 $U$ 上局部成立，即对所有 $u \in U$，有
$I\mathcal{O}_{U, u} = (0)$ 或 $I\mathcal{O}_{U, u} = (1)$。
由于 $V(I)$ 包含 $U$ 的闭点，可知 $I = 0$ 在 $U$ 上成立。
又因 $U \subset \Spec(A^h)$ 概形论稠密（因为本段开头已将 $A$ 替换为 $A/J_0$），
可知 $I = 0$。于是 $N'$ 在 $A^h$ 上平坦，断言 1 得证。

\medskip\noindent
回到断言 1 之前所述的情形。令 $A^h$ 为 $A$ 的 Hensel 化，令 $B'$ 为
$B \otimes_A A^h$ 在极大理想
$\mathfrak m_B \otimes A^h + B \otimes \mathfrak m_{A^h}$ 处的局部化，并令
$N' = N \otimes_B B'$。现在可知，定理 \ref{flat-theorem-flattening-local} 的扁平化理想
$I \subset A^h$ 是幂零的。若 $nil(A^h)$ 表示幂零元素理想，则
$nil(A^h) = nil(A) A^h$
（《代数进阶》引理 \ref{more-algebra-lemma-henselization-nil}）。因此存在有限生成的
幂零理想 $I_0 \subset A$，使得 $N/I_0N$ 在 $A/I_0$ 上平坦。

\medskip\noindent
{\bf 断言 2。} 对每个素理想 $\mathfrak p \subset A$，映射
$\kappa(\mathfrak p) \otimes_A N \to \kappa(\mathfrak p) \otimes_A M$
是单射。若该映射不是单射，则称 $\mathfrak p$ 为坏素理想。
% P05-EN-CH38-0048: frozen source line 7490 says “We say p is bad it this is false”; candidate only, do not silently repair.
设 $C$ 是非空坏素理想链，并令
$\mathfrak p^* = \bigcup_{\mathfrak p \in C} \mathfrak p$.
由引理 \ref{flat-lemma-find-pure-spreadout}，存在有限生成理想
$\mathfrak a \subset \mathfrak p^*A_{\mathfrak p^*}$
使得在 $V(\mathfrak a)$ 上存在纯展开。
若 $\mathfrak p^*$ 是好的，则由
引理 \ref{flat-lemma-properties-pure-spreadout}
可知 $V(\mathfrak a)$ 的点都是好的。
然而，由于 $\mathfrak a$ 有限生成，且
$\mathfrak p^*A_{\mathfrak p^*} = \bigcup_{\mathfrak p \in C}A_{\mathfrak p^*}$
可知 $V(\mathfrak a)$ 包含某个 $\mathfrak p \in C$，矛盾。
所以 $\mathfrak p^*$ 是坏的。由 Zorn 引理，若存在坏素理想，则存在一个极大的坏素理想，
记为 $\mathfrak p$。换言之，不妨设每个
$\mathfrak p' \supset \mathfrak p$ 且 $\mathfrak p' \not = \mathfrak p$ 都是好的。
在此情形，对每个 $f \in A$ 且 $f \not \in \mathfrak p$，映射
$u \otimes \text{id}_{A/(\mathfrak p + f)}$ 都是普遍单射，见引理
\ref{flat-lemma-universally-injective-if-flat}。
因此只需证明 $N/\mathfrak p N$ 关于由子模 $f(N/\mathfrak pN)$ 定义的拓扑是分离的。
由于 $B \to B'$ 忠实平坦，只需对模 $N'/\mathfrak p N'$ 证明同样结论。
由引理 \ref{flat-lemma-flat-finite-type-local-colimit-free} 以及《代数进阶》引理
\ref{more-algebra-lemma-content-exists-flat-Mittag-Leffler}，$N'/\mathfrak pN'$ 的元素
在 $A^h/\mathfrak pA^h$ 中都有内容理想。因此只需证明
$\bigcap_{f \in A, f \not \in \mathfrak p} f(A^h/\mathfrak p A^h) = 0$.
然后对 $A^h/\mathfrak q A^h$ 也证明同样结论，其中每个素理想
$\mathfrak q \subset A^h$ 都在 $\mathfrak p A^h$ 上极小。由于 $A \to A^h$ 是 Hensel 化，
每个 $\mathfrak q$ 都收缩到 $\mathfrak p$，而每个
$\mathfrak q' \supset \mathfrak q$ 且 $\mathfrak q' \not = \mathfrak q$ 都收缩到严格包含
$\mathfrak p$ 的素理想 $\mathfrak p'$。因此由引理
\ref{flat-lemma-big-intersection-is-zero} 得到交的消失。

\medskip\noindent
现在可以把一切合并起来。具体地，利用断言 1 和断言 2，可知
$N/I_0 N \to M/I_0M$ 由引理
\ref{flat-lemma-universally-injective-if-flat} 是 $A/I_0$-普遍单射。
于是下图
$$
\xymatrix{
N \otimes_A (I_0^n/I_0^{n + 1}) \ar[r] \ar[d] &
M \otimes_A (I_0^n/I_0^{n + 1}) \ar@{=}[d] \\
I_0^n N /I_0^{n + 1} N \ar[r] &
I_0^n M /I_0^{n + 1} M
}
$$
表明左侧竖直箭头是单射。因此由《代数》引理
\ref{algebra-lemma-what-does-it-mean-again}，可知 $N$ 平坦。类似地，
$u$ 的普遍单射性可以约化到模去 $I_0$ 的情形（甚至无需先证明 $N$ 平坦）。证明完毕。
\end{proof}


\section{有限型平坦模，第三部分}
\label{flat-section-finite-type-flat-III}

\noindent
下面的结果是本章的主要结果之一。

\begin{theorem}
\label{flat-theorem-check-flatness-at-associated-points}
设 $f : X \to S$ 局部有限型。
设 $\mathcal{F}$ 是有限型拟凝聚 $\mathcal{O}_X$-模。
设 $x \in X$ 的像为 $s \in S$。以下条件等价：
\begin{enumerate}
\item $\mathcal{F}$ 在 $x$ 处关于 $S$ 平坦；
\item 对每个专化到 $x$ 的 $x' \in \text{Ass}_{X_s}(\mathcal{F}_s)$，
$\mathcal{F}$ 在 $x'$ 处关于 $S$ 平坦。
\end{enumerate}
\end{theorem}

\begin{proof}
显然 (1) 蕴含 (2)，因为对每个专化到 $x$ 的点，$\mathcal{F}_{x'}$ 都是
$\mathcal{F}_x$ 的局部化。置 $A = \mathcal{O}_{S, s}$、$B = \mathcal{O}_{X, x}$，以及
$N = \mathcal{F}_x$。令 $\Sigma \subset B$ 为 $B$ 中在 $N/\mathfrak m_AN$ 上
作为非零因子作用的元素组成的乘法子集。由
引理 \ref{flat-lemma-homothety-spectrum} 对 $\Spec(\Sigma^{-1}N)$ 的描述，
假设 (2) 蕴含 $\Sigma^{-1}N$ 在 $A$ 上平坦。
另一方面，按构造映射 $N \to \Sigma^{-1}N$ 模去 $\mathfrak m_A$ 后是单射。
因此应用引理 \ref{flat-lemma-upstairs-finite-type-injective-into-flat-mod-m} 即得结论。
\end{proof}

\noindent
现在直接应用这一点，得到以下有用结果。

\begin{lemma}
\label{flat-lemma-check-along-closed-fibre}
设 $S$ 是闭点为 $s$ 的局部概形。
设 $f : X \to S$ 局部有限型。
设 $\mathcal{F}$ 是有限型拟凝聚 $\mathcal{O}_X$-模。
假设
\begin{enumerate}
\item $\text{Ass}_{X/S}(\mathcal{F})$ 的每个点都专化到闭纤维 $X_s$ 中的某点
\footnote{例如，当 $f$ 有限型且 $\mathcal{F}$ 沿 $X_s$ 纯时，或当 $f$ 适当时，这成立。}；
\item $\mathcal{F}$ 在 $X_s$ 的每个点处关于 $S$ 平坦。
\end{enumerate}
则 $\mathcal{F}$ 在 $S$ 上平坦。
\end{lemma}

\begin{proof}
这是直接的，因为由定理
\ref{flat-theorem-check-flatness-at-associated-points}，只需在 $\mathcal{F}$ 关于 $S$ 的相对伴随点处
检查平坦性即可。
\end{proof}











\section{泛扁平化}
\label{flat-section-flattening-final}

\noindent
若 $f : X \to S$ 是概形的适当有限表示态射，则可以找到 $f$ 的泛扁平化。
本节讨论这一点及其若干变体。

\begin{lemma}
\label{flat-lemma-free-at-generic-points-representable}
在情形 \ref{flat-situation-free-at-generic-points} 中，对每个 $p \geq 0$，函子 $H_p$
（\ref{flat-equation-free-at-generic-points}）由局部闭浸入 $S_p \to S$ 表示。
若 $\mathcal{F}$ 有限表示，则 $S_p \to S$ 有限表示。
\end{lemma}

\begin{proof}
对每个 $S$，我们将同时对所有 $p \geq 0$ 证明该断言。
由引理 \ref{flat-lemma-free-at-generic-points}，函子 $H_p$ 是 fppf 拓扑下的层。
因此结合
《下降》，引理 \ref{descent-lemma-descent-data-sheaves},
《态射进阶》，引理
\ref{more-morphisms-lemma-separated-locally-quasi-finite-morphisms-fppf-descend}
以及
下降理论引理 \ref{descent-lemma-descending-fppf-property-immersion}，可知该问题
关于 $S$ 上的 \'etale 拓扑是局部的。特别地，该问题关于 $S$ 的 Zariski 拓扑是局部的。

\medskip\noindent
对 $s \in S$，记 $\xi_s$ 为纤维 $X_s$ 的唯一泛点。
注意，对每个 $s \in S$，$\mathcal{F}$ 的限制 $\mathcal{F}_s$
在 $\xi_s$ 的某个邻域内局部自由，秩为某个 $p(s) \geq 0$。
（由于 $X_s$ 既约且光滑，这由 $\mathcal{F}_s$ 关于 $X_s$ 的泛平坦性得到，见
《代数》引理 \ref{algebra-lemma-generic-flatness-Noetherian}，尽管这是过强的条件。）
为今后引用，注意
$$
p(s) =
\dim_{\kappa(\xi_s)}(
\mathcal{F}_{\xi_s} \otimes_{\mathcal{O}_{X, \xi_s}} \kappa(\xi_s)
).
$$
特别地，$H_{p(s)}(s)$ 非空，而当 $q \not = p(s)$ 时 $H_q(s)$ 为空。

\medskip\noindent
设 $U \subset X$ 是开子概形。由于 $f : X \to S$ 光滑，故为开态射。
由 (\ref{flat-equation-free-at-generic-points}) 立即可知，对于成对对象
$(f|_U : U \to f(U), \mathcal{F}|_U)$ 的函子 $H_p$ 与对于成对对象
$(f|_{f^{-1}(f(U))}, \mathcal{F}|_{f^{-1}(f(U))})$ 的函子 $H_p$ 相同。
因此，为证明 $S_p$ 在 $f(U)$ 上的存在性，总可以将 $X$ 替换为 $U$。

\medskip\noindent
取 $s \in S$。存在 $\xi_s$ 的仿射开邻域 $U$，使得 $\mathcal{F}|_U$
至多由 $p(s)$ 个元素生成。由上述论证，为了在 $s$ 的一个邻域内证明
$H_{p(s)}$ 的断言，不妨设 $\mathcal{F}$ 由 $p(s)$ 个元素生成，即存在满射
$$
u : \mathcal{O}_X^{\oplus p(s)} \longrightarrow \mathcal{F}
$$
此时显然 $H_{p(s)}$ 等于映射 $u$ 的 $F_{iso}$
（\ref{flat-equation-iso}）（这立即由引理
\ref{flat-lemma-injectivity-map-source-flat-pure} 得到；也可在再缩小一些、使 $S$ 与 $X$ 都仿射后，
由引理 \ref{flat-lemma-induction-step-fp} 得到）。
因此可应用定理 \ref{flat-theorem-flattening-map}，得到 $H_{p(s)}$
在 $s$ 的某个邻域内由闭浸入表示。

\medskip\noindent
结果形式上由上述内容推出。具体地，上述论证表明在 $S$ 的局部，函数
$s \mapsto p(s)$ 有界。因此可对 $p = \max_{s \in S} p(s)$ 作归纳。
由上文，函子 $H_p$ 由闭浸入 $S_p \to S$ 表示。将 $S$ 替换为
$S \setminus S_p$，最大值至少降低一，遂由归纳假设得证。

\medskip\noindent
假设 $\mathcal{F}$ 有限表示。由引理
\ref{flat-lemma-free-at-generic-points} 的 (2) 与极限论注记
\ref{limits-remark-limit-preserving}，$S_p \to S$ 局部有限表示。
然后重做本段中的归纳论证，可知当 $S$ 仿射时每个 $S_p$ 都拟紧：
首先，若 $p = \max_{s \in S} p(s)$，则 $S_p \subset S$ 闭（见上文），因而拟紧。
由于 $S_p \to S$ 是有限表示的闭浸入，$U = S \setminus S_p$ 是 $S$ 中拟紧开集
（例如见态射论第 \ref{morphisms-section-constructible} 节的讨论）。
接着 $S_{p - 1} \to U$ 是有限表示的闭浸入，故 $S_{p - 1}$ 拟紧，且
$U' = S \setminus (S_p \cup S_{p - 1})$ 拟紧。依此类推。
\end{proof}

\begin{lemma}
\label{flat-lemma-localize-flat-dimension-n}
在情形 \ref{flat-situation-flat-dimension-n} 中，设 $h : X' \to X$ 是 \'etale 态射。
置 $\mathcal{F}' = h^*\mathcal{F}$ 及 $f' = f \circ h$。
令 $F_n'$ 为与 $(f' : X' \to S, \mathcal{F}')$ 相关的
(\ref{flat-equation-flat-dimension-n})。则 $F_n$ 是 $F_n'$ 的子函子；若
$h(X') \supset \text{Ass}_{X/S}(\mathcal{F})$，则 $F_n = F'_n$。
\end{lemma}

\begin{proof}
设 $T \to S$ 是任意态射。则 $h_T : X'_T \to X_T$ 是 \'etale 的，因为它是 \'etale 态射 $g$ 的基变换。对 $t \in T$，记
% P05-EN-CH38-0049: frozen source line 7739 names g although this lemma defines h; candidate only, do not silently repair.
$Z \subset X_t$ 为 $\mathcal{F}_T$ 在 $T$ 上不平坦的点集；类似地，记
$Z' \subset X'_t$ 为 $\mathcal{F}'_T$ 在 $T$ 上不平坦的点集。由于
$\mathcal{F}'_T = h_T^*\mathcal{F}_T$，有
$Z' = h_t^{-1}(Z)$，见态射论引理
\ref{morphisms-lemma-flat-permanence}。
因此 $Z' \to Z$ 是 \'etale 态射，从而 $\dim(Z') \leq \dim(Z)$
（例如由下降理论引理
\ref{descent-lemma-dimension-at-point-local}，或直接因为 \'etale 态射是相对维数为 $0$ 的光滑态射）。
这就得到 $F_n \subset F_n'$。


\medskip\noindent
最后，假设 $h(X') \supset \text{Ass}_{X/S}(\mathcal{F})$，且 $T \to S$ 是满足
$F_n'(T)$ 非空的态射，即 $\mathcal{F}'_T$ 在 $T$ 上维数 $\geq n$ 处平坦。
取 $t \in T$，令 $Z \subset X_t$ 和 $Z' \subset X'_t$ 如上。为得到矛盾，假设
$\dim(Z) \geq n$。取泛点 $\xi \in Z$，它对应于一个维数 $\geq n$ 的不可约分支。
令 $x \in \text{Ass}_{X_t}(\mathcal{F}_t)$ 是 $\xi$ 的一般化点。由
除子论引理 \ref{divisors-lemma-base-change-relative-assassin} 及
注记 \ref{divisors-remark-base-change-relative-assassin}，$x$ 映到
$\text{Ass}_{X/S}(\mathcal{F})$ 中的一个点。因此 $x$ 在 $h_T$ 的像中，
即存在 $x' \in X'_T$ 使 $x = h_T(x')$。但因 $x \leadsto \xi$ 且 $\dim(Z') < n$，
$x' \not \in Z'$。于是 $\mathcal{F}'_T$ 在 $x'$ 处关于 $T$ 平坦，
从而 $\mathcal{F}_T$ 在 $x$ 处关于 $T$ 平坦（由态射论引理
\ref{morphisms-lemma-flat-permanence}）。由于这对每个这样的 $x$ 都成立，
由定理 \ref{flat-theorem-check-flatness-at-associated-points}，$\mathcal{F}_T$ 在 $\xi$ 处
关于 $T$ 平坦，得到所需矛盾。
\end{proof}

\begin{lemma}
\label{flat-lemma-compare-H-F}
假设 $X \to S$ 是仿射概形之间的光滑态射，纤维几何不可约且维数为 $d$，并且
$\mathcal{F}$ 是有限表示的拟凝聚 $\mathcal{O}_X$-模。则存在某个 $c \geq 0$，使得
$F_d = \coprod_{p = 0, \ldots, c} H_p$，其中 $F_d$ 如
(\ref{flat-equation-flat-dimension-n})，而 $H_p$ 如
(\ref{flat-equation-free-at-generic-points})。
\end{lemma}

\begin{proof}
由于 $X$ 仿射且 $\mathcal{F}$ 是有限表示的拟凝聚模，$\mathcal{F}$ 可由
$c \geq 0$ 个元素生成。因此对任意 $x \in X$，都有
$\dim_{\kappa(x)}(\mathcal{F}_x \otimes \kappa(x))$ 不超过 $c$。
特别地，当 $p > c$ 时 $H_p = \emptyset$。此外，显然有包含
$\coprod H_p \to F_d$。引理的内容是：若基变换 $\mathcal{F}_T$ 在 $T$ 上维数
$\geq d$ 处平坦，且 $t \in T$，那么 $\mathcal{F}_T$ 在 $X_t$ 的唯一泛点
$\xi$ 的某个开邻域 $U \subset X_T$ 上自由，秩为某个 $r$。也就是说，$H_r$
包含 $U$ 的像，而该像是 $t$ 的一个开邻域。$U$ 的存在性由《态射进阶》引理
\ref{more-morphisms-lemma-flat-and-free-at-point-fibre} 得到。
\end{proof}

\begin{lemma}
\label{flat-lemma-flat-dimension-n-representable}
在情形 \ref{flat-situation-flat-dimension-n} 中，设 $s \in S$，$d \geq 0$。假设
\begin{enumerate}
\item 在某个点 $s \in S$ 上存在 $\mathcal{F}/X/S$ 的完全 d\'evissage；
\item $X$ 在 $S$ 上有限表示；
\item $\mathcal{F}$ 是有限表示的 $\mathcal{O}_X$-模；以及
\item $\mathcal{F}$ 在 $S$ 上维数 $\geq d + 1$ 处平坦。
\end{enumerate}
那么经过可能将 $S$ 替换为 $s$ 的一个开邻域后，函子
$F_d$（\ref{flat-equation-flat-dimension-n}）由有限表示的单态射
$Z_d \to S$ 表示。
\end{lemma}

\begin{proof}
先作一个预备说明：$X$、$S$ 是仿射概形，而且只需在 $S$ 上仿射概形范畴中证明
$F_d$ 由有限表示的单态射 $Z_d \to S$ 表示。（当然不要求 $Z_d$ 仿射。）
因此在本引理的整个证明中，我们都在 $S$ 上的仿射概形范畴中工作。

\medskip\noindent
令 $(Z_k, Y_k, i_k, \pi_k, \mathcal{G}_k, \alpha_k)_{k = 1, \ldots, n}$
为 $s$ 上 $\mathcal{F}/X/S$ 的一个完全 d\'evissage，见定义
\ref{flat-definition-complete-devissage}。我们对 d\'evissage 的长度 $n$ 作归纳。
回忆 $Y_k \to S$ 光滑且纤维几何不可约，见定义
\ref{flat-definition-one-step-devissage}。令 $d_k$ 为 $Y_k$ 关于 $S$ 的相对维数。
回忆 $i_{k, *}\mathcal{G}_k = \Coker(\alpha_k)$，且 $i_k$ 是闭浸入。
由上述定义有
$d_1 = \dim(\text{Supp}(\mathcal{F}_s))$，且
$$
d_k = \dim(\text{Supp}(\Coker(\alpha_{k - 1})_s))
= \dim(\text{Supp}(\mathcal{G}_{k, s}))
$$
对 $k = 2, \ldots, n$ 成立。由于 $\alpha_k$ 在 $(Y_k)_s$ 的泛点处是同构，故
$d_1 > d_2 > \ldots > d_n \geq 0$。

\medskip\noindent
注意 $i_1$ 是闭浸入且 $\mathcal{F} = i_{1, *}\mathcal{G}_1$。
因此对任意 $T \to S$（其中 $T$ 仿射），有
$\mathcal{F}_T = i_{1, T, *}\mathcal{G}_{1, T}$，且 $i_{1, T}$ 仍是
在 $T$ 上的概形之间的闭浸入。于是 $\mathcal{F}_T$ 在 $T$ 上维数 $\geq d$ 处平坦
当且仅当 $\mathcal{G}_{1, T}$ 在 $T$ 上维数 $\geq d$ 处平坦。
由于 $\pi_1 : Z_1 \to Y_1$ 有限，同理可知 $\mathcal{G}_{1, T}$ 在 $T$ 上维数
$\geq d$ 处平坦，当且仅当 $\pi_{1, T, *}\mathcal{G}_{1, T}$ 在 $T$ 上维数
$\geq d$ 处平坦。对
“在维数 $\geq d + 1$ 处平坦”也有同样论证，
由 $\mathcal{F}$ 的假设，特别得到 $\pi_{1, *}\mathcal{G}_1$ 在 $S$ 上维数
$\geq d + 1$ 处平坦。

\medskip\noindent
假设 $d_1 > d$。由上面的讨论，特别可知 $\pi_{1, *}\mathcal{G}_1$ 在 $S$ 上
于 $(Y_1)_s$ 的泛点处平坦。由引理
\ref{flat-lemma-induction-step-fp}，可将 $S$ 替换为 $s$ 的仿射邻域，并假设
$\alpha_1$ 是 $S$-普遍单射。由于 $\alpha_1$ 是 $S$-普遍单射，对任意仿射 $T$ 上的
态射 $T \to S$，有短正合列
$$
0 \to \mathcal{O}_{Y_{1, T}}^{\oplus r_1}
\to \pi_{1, T, *}\mathcal{G}_{1, T} \to \Coker(\alpha_1)_T \to 0
$$
且第一个箭头仍是 $T$-普遍单射。因此 $(Y_1)_T$ 中
$\pi_{1, T, *}\mathcal{G}_{1, T}$ 在 $T$ 上平坦的点集，与 $(Y_1)_T$ 中
$\Coker(\alpha_1)_T$ 在 $S$ 上平坦的点集相同。这样问题约化为层
$\Coker(\alpha_1)$；它有长度为 $n - 1$ 的完全 d\'evissage，故由归纳得证。

\medskip\noindent
若 $d_1 < d$，则 $F_d$ 由 $S$ 表示，结论成立。

\medskip\noindent
最后一种情形是 $d_1 = d$。此情形由引理
\ref{flat-lemma-compare-H-F} 与引理
\ref{flat-lemma-free-at-generic-points-representable} 的组合得到。
\end{proof}

\begin{theorem}
\label{flat-theorem-flat-dimension-n-representable}
在情形 \ref{flat-situation-flat-dimension-n} 中，进一步假设 $f$ 有限表示，
$\mathcal{F}$ 是有限表示的 $\mathcal{O}_X$-模，且 $\mathcal{F}$ 关于 $S$ 纯。
则 $F_n$ 由有限表示的单态射 $Z_n \to S$ 表示。
\end{theorem}

\begin{proof}
由引理 \ref{flat-lemma-flat-dimension-n}，函子 $F_n$ 是 fppf 拓扑下的层。
注意，有限表示的单态射是分离且拟有限的（态射论引理
\ref{morphisms-lemma-monomorphism-loc-finite-type-loc-quasi-finite}）。
因此结合
《下降》，引理 \ref{descent-lemma-descent-data-sheaves},
《态射进阶》，引理
\ref{more-morphisms-lemma-separated-locally-quasi-finite-morphisms-fppf-descend}
以及下降理论引理 \ref{descent-lemma-descending-property-monomorphism} 和
\ref{descent-lemma-descending-property-finite-presentation}，可知该问题关于 $S$ 的
\'etale 拓扑是局部的。

\medskip\noindent
特别地，该情形关于 $S$ 上的 Zariski 拓扑是局部的，因此可设 $S$ 为仿射。
在此情形，$f$ 的纤维维数有上界，所以当 $n$ 足够大时 $F_n$ 可表示。
因此我们可以对 $n$ 使用递降归纳。
设我们已知 $F_{n + 1}$ 由有限表示的单态射
$Z_{n + 1} \to S$ 表示。考虑基变换
$X_{n + 1} = Z_{n + 1} \times_S X$，以及 $\mathcal{F}$ 到
$X_{n + 1}$ 的拉回 $\mathcal{F}_{n + 1}$。
由于 $Z_{n + 1} \to S$ 是有限表示的单态射，故它拟有限；于是由引理
\ref{flat-lemma-quasi-finite-base-change} 可知，$\mathcal{F}_{n + 1}$ 关于
$Z_{n + 1}$ 是纯的。
由于 $F_n$ 是 $F_{n + 1}$ 的子函子，我们得出：要证明 $F_n$ 的结论，
只需证明情形 $\mathcal{F}_{n + 1}/X_{n + 1}/Z_{n + 1}$ 所对应的函子
具有该结论。
这样就把问题化为在 $S_{n + 1} = S$ 时证明 $F_n$ 的结论；也就是说，
我们可以假设 $\mathcal{F}$ 在 $S$ 上的维数 $\geq n + 1$ 的部分平坦。

\medskip\noindent
固定 $n$，并假设 $\mathcal{F}$ 在 $S$ 上的维数 $\geq n + 1$ 的部分平坦。
为了完成证明，我们必须说明 $F_n$ 由有限表示的单态射
$Z_n \to S$ 表示。
由于问题关于 $S$ 的 \'etale 拓扑是局部的，只需证明：对每个 $s \in S$，
存在一个初等 \'etale 邻域 $(S', s') \to (S, s)$，使得基变换到 $S'$ 后
结论成立。
因此由引理 \ref{flat-lemma-existence-complete}，我们可以假设存在
\'etale 态射 $h_j : Y_j \to X$，其中 $j = 1, \ldots, m$，并且对每个
$j$，在 $s$ 上存在 $\mathcal{F}_j/Y_j/S$ 的完全 d\'evissage；
这里 $\mathcal{F}_j$ 是 $\mathcal{F}$ 到 $Y_j$ 的拉回，并且
$X_s \subset \bigcup h_j(Y_j)$。注意，由引理
\ref{flat-lemma-localize-flat-dimension-n}，层 $\mathcal{F}_j$ 在 $S$ 上的
维数 $\geq n + 1$ 的部分仍然平坦。
置 $W = \bigcup h_j(Y_j)$，这是 $X$ 的拟紧开子集。
由于 $\mathcal{F}$ 沿着 $X_s$ 是纯的，我们得到
$$
E = \{t \in S \mid \text{Ass}_{X_t}(\mathcal{F}_t) \subset W \}.
$$
它包含 $s$ 的所有泛化点。由态射论进一步的引理
\ref{more-morphisms-lemma-relative-assassin-constructible}，
$E$ 是 $S$ 的可构造子集。我们已经看到
$\Spec(\mathcal{O}_{S, s}) \subset E$。由态射论引理
\ref{morphisms-lemma-constructible-containing-open}，可知 $E$ 包含 $s$ 的
一个开邻域。因此缩小 $S$ 后可设 $E = S$。
由引理 \ref{flat-lemma-localize-flat-dimension-n} 可知，只需对
$X = \coprod Y_j$ 且 $\mathcal{F} = \coprod \mathcal{F}_j$ 所对应的函子
$F_n$ 证明该引理。
% P05-EN-CH38-0050: source line 7974 uses \mathcal{F}_i although the indexed sheaf is \mathcal{F}_j; preserved literally.
若 $F_{j, n}$ 表示对应于 $Y_j \to S$ 及层 $\mathcal{F}_i$ 的函子，
则 $F_n = \prod F_{j, n}$。因此，只需证明每个 $F_{j, n}$ 都由某个有限表示的
单态射 $Z_{j, n} \to S$ 表示，因为此时
$$
Z_n = Z_{1, n} \times_S \ldots \times_S Z_{m, n}
$$
于是定理化归为引理 \ref{flat-lemma-flat-dimension-n-representable} 所处理的
特殊情形。
\end{proof}

\noindent
下面的引理用泛扁平化明确说明该定理的含义。

\begin{lemma}
\label{flat-lemma-when-universal-flattening}
设 $f : X \to S$ 是概形态射。
设 $\mathcal{F}$ 是拟凝聚的 $\mathcal{O}_X$-模。
\begin{enumerate}
\item 若 $f$ 是有限表示的，$\mathcal{F}$ 是有限表示的
$\mathcal{O}_X$-模，且 $\mathcal{F}$ 关于 $S$ 纯，则存在 $\mathcal{F}$ 的泛扁平化
$S' \to S$。此外，$S' \to S$ 是有限表示的单态射。
\item 若 $f$ 是有限表示的且 $X$ 关于 $S$ 纯，则存在 $X$ 的泛扁平化
$S' \to S$。此外，$S' \to S$ 是有限表示的单态射。
\item 若 $f$ 适当且有限表示，且 $\mathcal{F}$ 是有限表示的
$\mathcal{O}_X$-模，则存在 $\mathcal{F}$ 的泛扁平化 $S' \to S$。
此外，$S' \to S$ 是有限表示的单态射。
\item 若 $f$ 适当且有限表示，则存在 $X$ 的泛扁平化 $S' \to S$。
\end{enumerate}
\end{lemma}

\begin{proof}
这些陈述立即由定理
\ref{flat-theorem-flat-dimension-n-representable} 应用于 $F_0 = F_{flat}$ 得出；
另外，若 $f$ 适当，则 $\mathcal{F}$ 自动相对于基纯，见引理
\ref{flat-lemma-proper-pure}。
\end{proof}







\section{Grothendieck 存在定理 IV}
\label{flat-section-existence}

\noindent
本节继续讨论概形上同调中第
\ref{coherent-section-existence}、\ref{coherent-section-existence-proper} 和
\ref{coherent-section-existence-proper-support} 节的内容。
我们将在以下情形中工作。

\begin{situation}
\label{flat-situation-existence}
这里给定环的逆系统 $(A_n)$，其过渡映射是满射且核局部幂零。
置 $A = \lim A_n$。给定一个在 $A$ 上分离且有限表示的概形 $X$。
置 $X_n = X \times_{\Spec(A)} \Spec(A_n)$，并将其视为 $X$ 的闭子概形。
进一步给定一个系统 $(\mathcal{F}_n, \varphi_n)$，其中 $\mathcal{F}_n$ 是
有限表示的 $\mathcal{O}_{X_n}$-模，在 $A_n$ 上平坦，且其支集在 $A_n$ 上适当，
并且
$$
\varphi_n :
\mathcal{F}_n \otimes_{\mathcal{O}_{X_n}} \mathcal{O}_{X_{n - 1}}
\longrightarrow
\mathcal{F}_{n - 1}
$$
是同构（记号使用态射论引理
\ref{morphisms-lemma-i-star-equivalence} 中的等价）。
\end{situation}

 \noindent
我们的目标是考察能否在 $X$ 上找到拟凝聚层 $\mathcal{F}$，使得
$\mathcal{F}_n = \mathcal{F} \otimes_{\mathcal{O}_X} \mathcal{O}_{X_n}$
对所有 $n$ 都有。

\begin{lemma}
\label{flat-lemma-compute-what-it-should-be}
在情形 \ref{flat-situation-existence} 中，考虑
$$
K = R\lim_{D_\QCoh(\mathcal{O}_X)}(\mathcal{F}_n) =
DQ_X(R\lim_{D(\mathcal{O}_X)}\mathcal{F}_n)
$$
则 $K$ 属于 $D^b_{\QCoh}(\mathcal{O}_X)$；事实上，$K$ 的非零上同调层
只出现在次数 $\geq 0$ 中。
\end{lemma}

\begin{proof}
这是概形导出范畴中例
\ref{perfect-example-inverse-limit-quasi-coherent}.
\end{proof}

\begin{lemma}
\label{flat-lemma-compute-against-perfect}
在情形 \ref{flat-situation-existence} 中，令 $K$ 如引理
\ref{flat-lemma-compute-what-it-should-be} 所述。对 $D(\mathcal{O}_X)$ 的任意完全对象
$E$，有
\begin{enumerate}
\item $M = R\Gamma(X, K \otimes^\mathbf{L} E)$ 是 $D(A)$ 的完美对象，
并且存在规范同构
$R\Gamma(X_n, \mathcal{F}_n \otimes^\mathbf{L} E|_{X_n}) =
M \otimes_A^\mathbf{L} A_n$
在 $D(A_n)$ 中成立，
\item $N = R\Hom_X(E, K)$ 是 $D(A)$ 的完美对象，
并且存在规范同构
$R\Hom_{X_n}(E|_{X_n}, \mathcal{F}_n) = N \otimes_A^\mathbf{L} A_n$
在 $D(A_n)$ 中成立。
\end{enumerate}
在两个陈述中，$E|_{X_n}$ 表示 $E$ 到 $X_n$ 的导出拉回。
\end{lemma}

\begin{proof}
证明 (2)。记 $E_n = E|_{X_n}$，并令
$N_n = R\Hom_{X_n}(E_n, \mathcal{F}_n)$.
回忆 $R\Hom_{X_n}(-, -)$ 等于
$R\Gamma(X_n, R\SheafHom(-, -))$，见概形上同调第
\ref{cohomology-section-global-RHom} 节。因此由概形导出范畴引理
\ref{perfect-lemma-base-change-RHom-perfect}
可知 $N_n$ 是 $D(A_n)$ 的完全对象，且其构造与基变换相容。
因此，由 $\varphi_n$ 给出的映射
$N_n \otimes_{A_n}^\mathbf{L} A_{n - 1} \to N_{n - 1}$
是同构。由《代数进阶》引理
\ref{more-algebra-lemma-Rlim-perfect-gives-perfect}，可知 $R\lim N_n$ 是完全的，
且其基变换回 $A_n$ 恢复 $N_n$。
另一方面，正合函子
$R\Hom_X(E, -) : D_\QCoh(\mathcal{O}_X) \to D(A)$
作为三角范畴之间的函子与乘积交换，因而也与导出极限交换，于是
$$
R\Hom_X(E, K) =
R\lim R\Hom_X(E, \mathcal{F}_n) =
R\lim R\Hom_X(E_n, \mathcal{F}_n) =
R\lim N_n
$$
这证明了 (2)。要看 (1) 成立，使用概形上同调引理
\ref{cohomology-lemma-dual-perfect-complex} 将其转化为 (2)。
\end{proof}

\begin{lemma}
\label{flat-lemma-relative-pseudo-coherence}
在情形 \ref{flat-situation-existence} 中，令 $K$ 如引理
\ref{flat-lemma-compute-what-it-should-be} 所述。则 $K$ 关于 $A$ 拟相干。
\end{lemma}

\begin{proof}
% P05-EN-CH38-0051: source line 8135 spells Combining as “Combinging”; frozen source wording is recorded without a rendered correction.
结合引理 \ref{flat-lemma-compute-against-perfect} 和概形导出范畴引理
\ref{perfect-lemma-perfect-enough}，可知对 $D(\mathcal{O}_X)$ 中的所有拟相干
$E$，$R\Gamma(X, K \otimes^\mathbf{L} E)$ 在 $D(A)$ 中拟相干。
因此由态射论进一步的引理
\ref{more-morphisms-lemma-characterize-pseudo-coherent}.
\end{proof}

\begin{lemma}
\label{flat-lemma-compute-over-affine}
在情形 \ref{flat-situation-existence} 中，令 $K$ 如引理
\ref{flat-lemma-compute-what-it-should-be} 所述。对任意拟紧开集 $U \subset X$，有
$$
R\Gamma(U, K) \otimes_A^\mathbf{L} A_n =
R\Gamma(U_n, \mathcal{F}_n)
$$
在 $D(A_n)$ 中成立，其中 $U_n = U \cap X_n$。
\end{lemma}

\begin{proof}
固定 $n$。由概形导出范畴引理
\ref{perfect-lemma-computing-sections-as-colim}
存在 $X$ 上完全复形的系统 $E_m$，使得
$R\Gamma(U, K) = \text{hocolim} R\Gamma(X, K \otimes^\mathbf{L} E_m)$.
事实上，该公式不仅对 $K$ 成立，而且对
$D_\QCoh(\mathcal{O}_X)$ 的每个对象都成立。将其应用于 $\mathcal{F}_n$，得到
\begin{align*}
R\Gamma(U_n, \mathcal{F}_n)
& =
R\Gamma(U, \mathcal{F}_n) \\
& =
\text{hocolim}_m R\Gamma(X, \mathcal{F}_n \otimes^\mathbf{L} E_m) \\
& =
\text{hocolim}_m R\Gamma(X_n, \mathcal{F}_n \otimes^\mathbf{L} E_m|_{X_n})
\end{align*}
使用引理 \ref{flat-lemma-compute-against-perfect} 以及
$- \otimes_A^\mathbf{L} A_n$ 与同伦余极限交换这一事实，即得结论。
\end{proof}

\begin{lemma}
\label{flat-lemma-finitely-presented}
在情形 \ref{flat-situation-existence} 中，令 $K$ 如引理
\ref{flat-lemma-compute-what-it-should-be} 所述。
记 $X_0 \subset X$ 为由位于 $\Spec(A)$ 的闭子集
$\Spec(A_1) = \Spec(A_2) = \ldots$ 之上的点组成的闭子集。
存在包含 $X_0$ 的开子集 $W \subset X$，使得
\begin{enumerate}
\item 除非 $i = 0$，否则 $H^i(K)|_W$ 为零；
\item $\mathcal{F} = H^0(K)|_W$ 是有限表示的；
\item $\mathcal{F}_n = \mathcal{F} \otimes_{\mathcal{O}_X} \mathcal{O}_{X_n}$。
\end{enumerate}
\end{lemma}

\begin{proof}
固定 $n \geq 1$。按构造，在 $D_\QCoh(\mathcal{O}_X)$ 中存在规范映射
$K \to \mathcal{F}_n$，从而存在拟凝聚层之间的规范映射
$H^0(K) \to \mathcal{F}_n$。这解释了 (3) 的含义。

\medskip\noindent
设 $x \in X_0$ 为一点。我们将找到 $x$ 的一个开邻域 $W$，
使得 (1)、(2) 和 (3) 成立。由于 $X_0$ 拟紧，这将证明引理。
取 $x$ 的一个仿射开邻域 $U \subset X$，记 $U = \Spec(B)$。
取一个满射 $P \to B$，其中 $P$ 在 $A$ 上光滑。
由引理 \ref{flat-lemma-relative-pseudo-coherence} 及相对拟相干的定义，
存在一个由有限自由 $P$-模组成的有界上复形 $F^\bullet$ 表示 $Ri_*K$，
其中 $i : U \to \Spec(P)$ 是该表示诱导的闭浸入。
令 $M_n$ 为与 $\mathcal{F}_n|_U$ 对应的 $B$-模。
由引理 \ref{flat-lemma-compute-over-affine}
$$
H^i(F^\bullet \otimes_A A_n) =
\left\{
\begin{matrix}
0 & \text{若} & i \not = 0 \\
M_n & \text{若} & i = 0
\end{matrix}
\right.
$$
令 $i$ 为使 $F^i$ 非零的最大指标。若 $i \leq 0$，则 (1)、(2) 和 (3) 成立。
否则 $i > 0$，且我们看到映射
$$
F^{i - 1} \to F^i
$$
在点 $x$ 处的秩最大。因此在 $\Spec(P)$ 中 $x$ 的某个开邻域上秩最大。
于是将 $P$ 替换为主元局部化后，可设上面显示的映射是满射。
由于 $F^i$ 有限自由，我们可以选择分裂 $F^{i - 1} = F' \oplus F^i$。
然后可以将 $F^\bullet$ 替换为复形
$$
\ldots \to F^{i - 2} \to F' \to 0 \to \ldots
$$
并对 $i$ 作归纳即可。
\end{proof}

\begin{lemma}
\label{flat-lemma-proper-support}
在情形 \ref{flat-situation-existence} 中，令 $K$ 如引理
\ref{flat-lemma-compute-what-it-should-be} 所述。令 $W \subset X$
如引理 \ref{flat-lemma-finitely-presented} 所述。置
$\mathcal{F} = H^0(K)|_W$。则在必要时缩小开集 $W$ 后，
$\mathcal{F}$ 的支集在 $A$ 上适当。
\end{lemma}

\begin{proof}
固定 $n \geq 1$。令 $I_n = \Ker(A \to A_n)$。
由《代数进阶》引理 \ref{more-algebra-lemma-limit-henselian}，
$(A, I_n)$ 是 Hensel 对。
令 $Z \subset W$ 为 $\mathcal{F}$ 的支集。由于 $\mathcal{F}$ 有限表示，
这是一个闭子集。由引理 \ref{flat-lemma-finitely-presented} 的 (3)，可知
$Z \times_{\Spec(A)} \Spec(A_n)$ 等于 $\mathcal{F}_n$ 的支集，因而在
% P05-EN-CH38-0052: source line 8262 uses A/I although the construction is indexed by I_n; frozen token is retained literally.
$\Spec(A/I)$ 上适当。
由态射论进一步的引理
\ref{more-morphisms-lemma-split-off-proper-part-henselian}
可写成 $Z = Z_1 \amalg Z_2$，其中 $Z_1, Z_2$ 在 $Z$ 中既开又闭，
$Z_1$ 在 $A$ 上适当，并且
$Z_1 \times_{\Spec(A)} \Spec(A/I_n)$ 等于 $\mathcal{F}_n$ 的支集。
换言之，$Z_2$ 不与 $X_0$ 相交。因此将 $W$ 替换为
$W \setminus Z_2$ 后即得引理。
\end{proof}

\begin{lemma}
\label{flat-lemma-monomorphism-isomorphism}
令 $A = \lim A_n$ 为环系统的极限，其过渡映射为满射且核局部幂零。
令 $S = \Spec(A)$。令 $T \to S$ 为局部有限型的单态射。
若对所有 $n$，$\Spec(A_n) \to S$ 都经由 $T$ 分解，则 $T = S$。
\end{lemma}

\begin{proof}
置 $S_n = \Spec(A_n)$。令 $T_0 \subset T$ 为分解 $S_n \to T$ 的公共像。
则 $T_0$ 拟紧。取包含 $T_0$ 的拟紧开集 $T' \subset T$。
于是 $S_n \to T$ 经由 $T'$ 分解。若能证明 $T' = S$，则
$T' = T = S$。因此可设 $T$ 拟紧。

\medskip\noindent
假设 $T$ 拟紧。在此情形，$T \to S$ 分离且拟有限（态射论引理
\ref{morphisms-lemma-monomorphism-loc-finite-type-loc-quasi-finite}).
利用 Zariski 主定理（采用态射论进一步引理
\ref{more-morphisms-lemma-quasi-finite-separated-pass-through-finite})
可选取分解 $T \to W \to S$，其中 $W \to S$ 有限且
$T \to W$ 为开浸入。写作 $W = \Spec(B)$。
分解 $S_n \to T$ 是唯一的，也可视为到 $W$ 的态射，从而得到
$$
A \longrightarrow B \longrightarrow \lim A_n = A
$$
考虑由右侧箭头得到的态射
$h : S = \Spec(A) \to \Spec(B) = W$。则
$$
T \times_{W, h} S
$$
是 $S$ 的开子概形，并包含所有 $n$ 的 $S_n \to S$ 的像。
完成证明只需说明：若开集 $U \subset S$ 包含某个 $n \geq 1$ 的
$S_n \to S$ 的像，则该开集等于 $S$。这是因为
$(A, \Ker(A \to A_n))$ 是 Hensel 对（《代数进阶》引理
\ref{more-algebra-lemma-limit-henselian}），从而 $S$ 的每个闭点都包含在
$S_n \to S$ 的像中。
\end{proof}

\begin{theorem}[Grothendieck 存在定理]
\label{flat-theorem-existence}
在情形 \ref{flat-situation-existence} 中，存在有限表示的
$\mathcal{O}_X$-模 $\mathcal{F}$，其在 $A$ 上平坦且支集在 $A$ 上适当，并满足
$\mathcal{F}_n = \mathcal{F} \otimes_{\mathcal{O}_X} \mathcal{O}_{X_n}$
对所有 $n$ 均成立，且与映射 $\varphi_n$ 相容。
\end{theorem}

\begin{proof}
应用引理 \ref{flat-lemma-compute-what-it-should-be}、
\ref{flat-lemma-compute-against-perfect},
\ref{flat-lemma-relative-pseudo-coherence},
\ref{flat-lemma-compute-over-affine},
\ref{flat-lemma-finitely-presented} 和
\ref{flat-lemma-proper-support}
，得到一个开子概形 $W \subset X$，包含所有位于 $\Spec(A_n)$
之上的点，以及一个有限表示的 $\mathcal{O}_W$-模 $\mathcal{F}$，
其支集在 $A$ 上适当，并且
$\mathcal{F}_n = \mathcal{F} \otimes_{\mathcal{O}_W} \mathcal{O}_{X_n}$
对所有 $n \geq 1$ 均成立。（这是有意义的，因为 $X_n \subset W$。）
由引理 \ref{flat-lemma-proper-pure}，$\mathcal{F}$ 关于 $\Spec(A)$ 普遍纯。
由定理 \ref{flat-theorem-flat-dimension-n-representable}（解释见引理
\ref{flat-lemma-when-universal-flattening}），存在 $\mathcal{F}$ 的泛扁平化
$S' \to \Spec(A)$，且态射 $S' \to \Spec(A)$ 是有限表示的单态射。
由于 $\mathcal{F}$ 到 $\Spec(A_n)$ 的基变换是 $\mathcal{F}_n$，可知对每个
$n$，$\Spec(A_n) \to \Spec(A)$ 唯一地经由 $S'$ 分解。
由引理 \ref{flat-lemma-monomorphism-isomorphism}，可知 $S' = \Spec(A)$。
这意味着 $\mathcal{F}$ 在 $A$ 上平坦。
最后，由于 $\mathcal{F}$ 的概形论支集 $Z$ 在 $\Spec(A)$ 上适当，
态射 $Z \to X$ 是闭的。因此推前 $(W \to X)_*\mathcal{F}$ 支持在 $W$ 上，
并具有所需的全部性质。
\end{proof}






\section{Grothendieck 存在定理 V}
\label{flat-section-existence-derived}

\noindent
本节在导出范畴中证明 Grothendieck 存在定理的一个类比，沿用第
\ref{flat-section-existence} 节处理拟凝聚模的方法。
经典情形在概形上同调的
\ref{coherent-section-existence}、\ref{coherent-section-existence-proper} 和
\ref{coherent-section-existence-proper-support} 节中讨论。
我们将在以下情形中工作。

\begin{situation}
\label{flat-situation-existence-derived}
这里给定环的逆系统 $(A_n)$，其过渡映射为满射且核局部幂零。
置 $A = \lim A_n$。给定一个在 $A$ 上适当、平坦且有限表示的概形 $X$。
置 $X_n = X \times_{\Spec(A)} \Spec(A_n)$，并将其视为 $X$ 的闭子概形。
进一步给定系统 $(K_n, \varphi_n)$，其中 $K_n$ 是
$D(\mathcal{O}_{X_n})$ 的拟相干对象，并且
$$
\varphi_n : K_n \longrightarrow K_{n - 1}
$$
是 $D(\mathcal{O}_{X_n})$ 中的映射，并诱导 $D(\mathcal{O}_{X_{n - 1}})$ 中的同构
$K_n \otimes_{\mathcal{O}_{X_n}}^\mathbf{L} \mathcal{O}_{X_{n - 1}}
\to K_{n - 1}$。
\end{situation}

 \noindent
更准确地说，应写成
$\varphi_n : K_n \to Ri_{n - 1, *}K_{n - 1}$
其中 $i_{n - 1} : X_{n - 1} \to X_n$ 是包含态射；在此记号下，条件是伴随映射
$Li_{n - 1}^*K_n \to K_{n - 1}$ 为同构。
我们的目标是找到拟相干对象 $K \in D(\mathcal{O}_X)$，使得
$K_n = K \otimes_{\mathcal{O}_X}^\mathbf{L} \mathcal{O}_{X_n}$
对所有 $n$ 成立（仍以同样方式滥用记号）。

\begin{lemma}
\label{flat-lemma-compute-what-it-should-be-derived}
在情形 \ref{flat-situation-existence-derived} 中，考虑
$$
K = R\lim_{D_\QCoh(\mathcal{O}_X)}(K_n) =
DQ_X(R\lim_{D(\mathcal{O}_X)} K_n)
$$
则 $K$ 属于 $D^-_{\QCoh}(\mathcal{O}_X)$。
\end{lemma}

\begin{proof}
函子 $DQ_X$ 存在，因为 $X$ 拟紧且拟分离，见概形导出范畴引理
\ref{perfect-lemma-better-coherator}.
由于 $DQ_X$ 是右伴随，它与乘积交换，因而也与导出极限交换。
这就得到引理陈述中的等式。

\medskip\noindent
由概形导出范畴引理 \ref{perfect-lemma-boundedness-better-coherator}，
函子 $DQ_X$ 的上同调维数有界。因此只需证明
$R\lim K_n \in D^-(\mathcal{O}_X)$。
为此，取 $U \subset X$ 为仿射开集。则存在规范正合序列
$$
0 \to
R^1\lim H^{m - 1}(U, K_n) \to H^m(U, R\lim K_n) \to
\lim H^m(U, K_n) \to 0
$$
由概形上同调引理 \ref{cohomology-lemma-RGamma-commutes-with-Rlim} 得到。
由于 $U$ 仿射且 $K_n$ 拟相干（因而由概形导出范畴引理
\ref{perfect-lemma-pseudo-coherent} 具有拟凝聚上同调层），由概形导出范畴引理
\ref{perfect-lemma-affine-compare-bounded} 可知
$H^m(U, K_n) = H^m(K_n)(U)$。
因此只需证明 $K_n$ 的上界与 $n$ 无关。

\medskip\noindent
由于 $K_n$ 拟相干，有 $K_n \in D^-(\mathcal{O}_{X_n})$。
设 $a_n$ 是使 $H^{a_n}(K_n)$ 非零的最大指标。当然
$a_1 \leq a_2 \leq a_3 \leq \ldots$。
注意，$H^{a_n}(K_n)$ 是一个
$\mathcal{O}_{X_n}$-有限表示模
(概形上同调引理 \ref{cohomology-lemma-finite-cohomology})。
我们有 $H^{a_n}(K_{n - 1}) =
H^{a_n}(K_n) \otimes_{\mathcal{O}_{X_n}} \mathcal{O}_{X_{n - 1}}$.
由于 $X_{n - 1} \to X_n$ 是加厚，由 Nakayama 引理（代数引理
\ref{algebra-lemma-NAK}）可知，若
$H^{a_n}(K_n) \otimes_{\mathcal{O}_{X_n}} \mathcal{O}_{X_{n - 1}}$
为零，则 $H^{a_n}(K_n)$ 也为零。因此对所有 $n$ 有
$a_n = a_{n - 1}$，结论得证。
\end{proof}

\begin{lemma}
\label{flat-lemma-compute-against-perfect-derived}
在情形 \ref{flat-situation-existence-derived} 中，令 $K$ 如引理
\ref{flat-lemma-compute-what-it-should-be-derived} 所述。对
$D(\mathcal{O}_X)$ 的任意完全对象 $E$，上同调
$$
M = R\Gamma(X, K \otimes^\mathbf{L} E)
$$
是 $D(A)$ 的拟相干对象，且存在规范同构
$$
R\Gamma(X_n, K_n \otimes^\mathbf{L} E|_{X_n}) = M \otimes_A^\mathbf{L} A_n
$$
在 $D(A_n)$ 中成立。这里 $E|_{X_n}$ 表示 $E$ 到 $X_n$ 的导出拉回。
\end{lemma}

\begin{proof}
记 $E_n = E|_{X_n}$，并令
$M_n = R\Gamma(X_n, K_n \otimes^\mathbf{L} E|_{X_n})$.
由概形导出范畴引理
\ref{perfect-lemma-flat-proper-pseudo-coherent-direct-image-general}
可知 $M_n$ 是 $D(A_n)$ 的拟相干对象，且其构造与基变换相容。
因此，由 $\varphi_n$ 给出的映射
$M_n \otimes_{A_n}^\mathbf{L} A_{n - 1} \to M_{n - 1}$
是同构。由《代数进阶》引理
\ref{more-algebra-lemma-Rlim-pseudo-coherent-gives-pseudo-coherent}
可知 $R\lim M_n$ 拟相干，且其基变换回 $A_n$ 恢复 $M_n$。
另一方面，正合函子
$R\Gamma(X, -) : D_\QCoh(\mathcal{O}_X) \to D(A)$
作为三角范畴之间的函子与乘积交换，因而也与导出极限交换，于是
$$
R\Gamma(X, E \otimes^\mathbf{L} K) =
R\lim R\Gamma(X,  E \otimes^\mathbf{L} K_n) =
R\lim R\Gamma(X_n, E_n \otimes^\mathbf{L} K_n) =
R\lim M_n
$$
这正是所需结论。
\end{proof}

\begin{lemma}
\label{flat-lemma-relative-pseudo-coherence-derived}
在情形 \ref{flat-situation-existence-derived} 中，令 $K$ 如引理
\ref{flat-lemma-compute-what-it-should-be-derived} 所述。则 $K$ 在 $X$ 上拟相干。
\end{lemma}

\begin{proof}
% P05-EN-CH38-0053: source line 8510 spells Combining as “Combinging”; translated by intended meaning.
结合引理 \ref{flat-lemma-compute-against-perfect-derived} 和概形导出范畴引理
\ref{perfect-lemma-perfect-enough}
可知，对 $D(\mathcal{O}_X)$ 中所有拟相干的 $E$，
$R\Gamma(X, K \otimes^\mathbf{L} E)$ 在 $D(A)$ 中拟相干。
因此由态射论进一步的引理
\ref{more-morphisms-lemma-characterize-pseudo-coherent}
可知 $K$ 关于 $A$ 拟相干。
% P05-EN-CH38-0054: source line 8519 has “is of flat”; translated by intended “is flat”.
由于 $X$ 在 $A$ 上平坦且有限表示，这等价于在 $X$ 上拟相干，见
态射论进一步的引理
\ref{more-morphisms-lemma-check-relative-pseudo-coherence-on-charts}.
\end{proof}

\begin{lemma}
\label{flat-lemma-compute-over-affine-derived}
在情形 \ref{flat-situation-existence-derived} 中，令 $K$ 如引理
\ref{flat-lemma-compute-what-it-should-be-derived} 所述。对任意拟紧开集
$U \subset X$，有
$$
R\Gamma(U, K) \otimes_A^\mathbf{L} A_n =
R\Gamma(U_n, K_n)
$$
在 $D(A_n)$ 中成立，其中 $U_n = U \cap X_n$。
\end{lemma}

\begin{proof}
固定 $n$。由概形导出范畴引理
\ref{perfect-lemma-computing-sections-as-colim}
存在 $X$ 上完全复形的系统 $E_m$，使得
$R\Gamma(U, K) = \text{hocolim} R\Gamma(X, K \otimes^\mathbf{L} E_m)$.
事实上，该公式不仅对 $K$ 成立，而且对
$D_\QCoh(\mathcal{O}_X)$ 的每个对象都成立。将其应用于 $K_n$，得到
\begin{align*}
R\Gamma(U_n, K_n)
& =
R\Gamma(U, K_n) \\
& =
\text{hocolim}_m R\Gamma(X, K_n \otimes^\mathbf{L} E_m) \\
& =
\text{hocolim}_m R\Gamma(X_n, K_n \otimes^\mathbf{L} E_m|_{X_n})
\end{align*}
利用引理 \ref{flat-lemma-compute-against-perfect-derived} 以及
$- \otimes_A^\mathbf{L} A_n$ 与同伦余极限交换这一事实，即得结论。
\end{proof}

\begin{theorem}[导出 Grothendieck 存在定理]
\label{flat-theorem-existence-derived}
在情形 \ref{flat-situation-existence-derived} 中，存在
$D(\mathcal{O}_X)$ 中的拟相干对象 $K$，使得
$K_n = K \otimes_{\mathcal{O}_X}^\mathbf{L} \mathcal{O}_{X_n}$
对所有 $n$ 均成立，且与映射 $\varphi_n$ 相容。
\end{theorem}

\begin{proof}
应用引理 \ref{flat-lemma-compute-what-it-should-be-derived}、
\ref{flat-lemma-compute-against-perfect-derived},
\ref{flat-lemma-relative-pseudo-coherence-derived}
，得到 $D(\mathcal{O}_X)$ 的拟相干对象 $K$。
在引理 \ref{flat-lemma-compute-over-affine-derived} 中取仿射开集，
立即可知 $K$ 在 $X_n$ 上的限制为 $K_n$。
\end{proof}

\begin{remark}
\label{flat-remark-correct-generality}
本节的结果可以推广。若只假设 $X \to \Spec(A)$ 分离且有限表示，
$K_n$ 关于 $A_n$ 拟相干，并支持在 $X_n$ 的一个在 $A_n$ 上适当的闭子集上，
结论很可能仍然成立。所得 $K$ 将关于 $A$ 拟相干，并支持在一个
在 $A$ 上适当的闭子集上。若将来需要这一结果，我们会在此给出精确陈述并证明。
\end{remark}






\section{爆破与平坦性}
\label{flat-section-blowup-flat}

\noindent
本节继续讨论如下形式的结果：“爆破后严格变换变得平坦”，见《代数进阶》第
\ref{more-algebra-section-blowup-flat} 节和除子论第
\ref{divisors-section-blowup-flat} 节。
本节采用以下（多少算是标准的）记号。若 $X \to S$ 是概形态射，
$\mathcal{F}$ 是 $X$ 上的拟凝聚模，且 $T \to S$ 是概形态射，
则用 $\mathcal{F}_T$ 表示 $\mathcal{F}$ 到基变换
$X_T = X \times_S T$ 的拉回。

\begin{remark}
\label{flat-remark-successive-blowups}
设 $S$ 为拟紧拟分离概形。设 $f : X \to S$ 为概形态射。
设 $\mathcal{F}$ 为 $X$ 上的拟凝聚模。设 $U \subset S$ 为拟紧开子概形。
给定一个 $U$-容许爆破 $S' \to S$，用 $X'$ 表示 $X$ 的严格变换，
用 $\mathcal{F}'$ 表示 $\mathcal{F}$ 的严格变换，并通过除子论引理
\ref{divisors-lemma-strict-transform} 将其视为 $X'$ 上的拟凝聚模。
设 $P$ 是 $\mathcal{F}/X/S$ 的一个性质，它在 $U$-容许爆破下作上述严格变换时
保持不变。本节的一般问题是：在 $\mathcal{F}/X/S$ 满足辅助条件时，证明存在
$U$-容许爆破 $S' \to S$，使严格变换 $\mathcal{F}'/X'/S'$ 具有性质 $P$。

\medskip\noindent
一般策略是使用如下事实：$U$-容许爆破的复合仍是 $U$-容许爆破，见除子论引理
\ref{divisors-lemma-composition-admissible-blowups}。
实际上，我们将使用更精确的除子论引理
\ref{divisors-lemma-composition-finite-type-blowups}，并结合除子论引理
\ref{divisors-lemma-strict-transform-composition-blowups}。
因此，只需找到一列 $U$-容许爆破
$$
S = S_0 \leftarrow S_1 \leftarrow \ldots \leftarrow S_n
$$
使得置 $\mathcal{F}_0 = \mathcal{F}$、$X_0 = X$，并令
$\mathcal{F}_i/X_i$ 等于 $\mathcal{F}_{i - 1}/X_{i - 1}$ 的严格变换后，
最终得到具有性质 $P$ 的 $\mathcal{F}_n/X_n/S_n$。

\medskip\noindent
特别地，选择有限型拟凝聚理想层
$\mathcal{I} \subset \mathcal{O}_S$，使得 $V(\mathcal{I}) = S \setminus U$，
见性质论引理 \ref{properties-lemma-quasi-coherent-finite-type-ideals}。
令 $S' \to S$ 为沿 $\mathcal{I}$ 的爆破，令 $E \subset S'$ 为例外除子
（除子论引理
\ref{divisors-lemma-blowing-up-gives-effective-Cartier-divisor}）。
于是问题化为存在有效 Cartier 除子 $D \subset S$ 且其支集为
% P05-EN-CH38-0055: source line 8651 uses X \setminus U although D and U lie in S; frozen token retained.
$X \setminus U$ 的情形。特别地，可设 $U$ 在 $S$ 中概形论稠密
（除子论引理 \ref{divisors-lemma-complement-effective-Cartier-divisor}）。

\medskip\noindent
设 $P$ 关于 $S$ 是局部的：若 $S = \bigcup S_i$ 是由拟紧开集组成的有限开覆盖，
且 $P$ 对 $\mathcal{F}_{S_i}/X_{S_i}/S_i$ 成立，则 $P$ 对
$\mathcal{F}/X/S$ 成立。
在此情形，上述一般问题也关于 $S$ 局部；也就是说，若给定 $s \in S$，
可以找到 $s$ 的拟紧开邻域 $W$，使 $\mathcal{F}_W/X_W/W$ 的问题可解，
则 $\mathcal{F}/X/S$ 的问题可解。这由除子论引理
\ref{divisors-lemma-extend-admissible-blowups} 和
\ref{divisors-lemma-dominate-admissible-blowups} 得出。
\end{remark}

\begin{lemma}
\label{flat-lemma-helper-blowup-affine-space}
设 $R$ 为环，$f \in R$。设 $r\geq 0$ 为整数。
设 $R \to S$ 为环同态，$M$ 为 $S$-模。假设
\begin{enumerate}
\item $R \to S$ 有限表示且平坦；
% P05-EN-CH38-0056: source line 8674 says a fibre over kappa(p) is geometrically integral over R; frozen base token retained.
\item 每个纤维环 $S \otimes_R \kappa(\mathfrak p)$ 在 $R$ 上几何整；
\item $M$ 是有限 $S$-模；
\item $M_f$ 是有限表示的 $S_f$-模；
\item 对所有 $\mathfrak p \in R$，若 $f \not \in \mathfrak p$ 且
$\mathfrak q = \mathfrak pS$，则模 $M_{\mathfrak q}$ 是秩为 $r$ 的
自由 $S_\mathfrak q$-模。
\end{enumerate}
则存在有限生成理想 $I \subset R$，满足 $V(f) = V(I)$，并且对所有
$a \in I$，若 $R' = R[\frac{I}{a}]$，则商
$$
M' = (M \otimes_R R')/a\text{-幂挠部分}
$$
作为 $S' = S \otimes_R R'$-模满足如下性质：对每个素理想
$\mathfrak p' \subset R'$，存在 $g \in S'$，使
$g \not \in \mathfrak p'S'$，且 $M'_g$ 是秩为 $r$ 的自由
$S'_g$-模。
\end{lemma}

\begin{proof}
本引理推广了《代数进阶》引理 \ref{more-algebra-lemma-blowup-module}；
我们建议读者先阅读该证明。
% P05-EN-CH38-0057: source line 8697 attributes finite generation of M to (1), while hypothesis (3) supplies it; frozen citation retained.
选择满射 $S^{\oplus n} \to M$，由 (1) 这是可行的。
选择有限子模 $K \subset \Ker(S^{\oplus n} \to M)$，使得
$S^{\oplus n}/K \to M$ 在逆转 $f$ 后成为同构。由 (4) 这是可行的。
置 $M_1 = S^{\oplus n}/K$，并假设可以对 $M_1$ 证明本引理。
设 $I \subset R$ 为相应理想。则对 $a \in I$，映射
$$
M_1' = (M_1 \otimes_R R')/a\text{-幂挠部分}
\longrightarrow
M' = (M \otimes_R R')/a\text{-幂挠部分}
$$
是满射。由于 $R'_a = R_f$，它在 $R'$ 中逆转 $a$ 后也是同构，见代数引理
\ref{algebra-lemma-blowup-in-principal}。
但 $a$ 在 $M'_1$ 上是非零因子，故上面显示的映射是同构。
因此只需在 $M$ 为有限表示的 $S$-模时证明本引理。

\medskip\noindent
假设 $M$ 是满足 (3) 的有限表示 $S$-模。
则 $J = \text{Fit}_r(M) \subset S$ 是有限生成理想。
由引理 \ref{flat-lemma-fibres-irreducible-flat-projective-nonnoetherian}，
可以将 $S$ 写成自由 $R$-模的直和项：
$\bigoplus_{\alpha \in A} R = S \oplus C$。
对任意元素 $h \in S$，在上述分解中写成 $h = \sum a_\alpha$，
称 $a_\alpha$ 为 $h$ 的系数。
令 $I' \subset R$ 为由 $J$ 中元素的系数组成的理想。
乘以 $S$ 的元素定义 $R$-线性映射 $S \to S$，故 $I'$ 由 $J$ 的生成元的
系数生成；也就是说，$I'$ 是有限生成理想。我们断言 $I = fI'$ 可行。

\medskip\noindent
先验证 $V(f) = V(I)$。包含关系 $V(f) \subset V(I)$ 是显然的。
反过来，若 $f \not \in \mathfrak p$，则由性质 (5) 和《代数进阶》引理
\ref{more-algebra-lemma-fitting-ideal-generate-locally}，
$\mathfrak q =  \mathfrak p S$ 不属于 $V(J)$。
因此，$J$ 中存在一个元素，其在 $S \otimes_R \kappa(\mathfrak p)$ 中的像不为零。
所以 $I'$ 中存在不包含于 $\mathfrak p$ 的元素，故
$\mathfrak p \not \in V(fI') = V(I)$。

\medskip\noindent
令 $a \in I$，并置 $R' = R[\frac{I}{a}]$。
对某个 $a' \in I'$，可写成 $a = fa'$。
由代数引理 \ref{algebra-lemma-affine-blowup} 和
\ref{algebra-lemma-blowup-add-principal}，可知 $I' R' = a'R'$，
且 $a'$ 在 $R'$ 中是非零因子。
% P05-EN-CH38-0058: source line 8743 bases S over itself rather than R; frozen tensor base is retained.
置 $S' = S \otimes_S R'$。
$JS' = \text{Fit}_r(M \otimes_S S')$ 的每个元素 $g$ 都可写成
$g = \sum_\alpha c_\alpha$，其中某些 $c_\alpha \in I'R'$。
由于 $I'R' = a'R'$，可对某些 $c'_\alpha \in R'$ 写成
$c_\alpha = a'c'_\alpha$，并且在 $S'$ 中有
$g = (\sum c'_\alpha)a' = g' a'$。
此外，存在 $g_0 \in J$，使某个 $\alpha$ 满足 $a' = c_\alpha$。
对该元素，在 $S'$ 中有 $g_0 = g'_0 a'$，其中 $g'_0$ 是 $S'$ 中的单位。
令 $\mathfrak p' \subset R'$ 为素理想，且
$\mathfrak q' = \mathfrak p'S'$。
由上可知，$JS'_{\mathfrak q'}$ 是由非零因子 $a'$ 生成的主理想。
由《代数进阶》引理
\ref{more-algebra-lemma-principal-fitting-ideal}
可知 $M'_{\mathfrak q'}$ 可由 $r$ 个元素生成。
由于 $M'$ 有限，存在 $m_1, \ldots, m_r \in M'$ 和
$g \in S'$，且 $g \not \in \mathfrak q'$，使得相应映射
$(S')^{\oplus r} \to M'$ 在逆转 $g$ 后成为满射。

\medskip\noindent
最后，考虑理想 $J' = \text{Fit}_{r - 1}(M')$。
注意 $J'S'_g$ 由 $m_1, \ldots, m_r$ 之间关系的系数生成
（Fitting 理想与基变换相容）。因此只需证明 $J' = 0$，见《代数进阶》引理
\ref{more-algebra-lemma-fitting-ideal-finite-locally-free}.
由于 $R'_a = R_f$（代数引理 \ref{algebra-lemma-blowup-in-principal}）
且 $M'_a = M_f$，由 (5) 可知，对任意形如
$\mathfrak q'' = \mathfrak p''S'$ 的素理想 $\mathfrak q'' \subset S'$，
其中 $\mathfrak p'' \subset R'_a$，$J'_a$ 在
$S_{\mathfrak q''}$ 中的像为零。由于
$S'_a \subset \prod_{\mathfrak q''\text{ 如上}} S'_{\mathfrak q''}$
（由引理 \ref{flat-lemma-base-change-universally-flat}，
$(S'_a)_{\mathfrak p''} \subset S'_{\mathfrak q''}$），
% P05-EN-CH38-0059: source line 8777 multiplies the S'-ideal J' by R'_a; frozen expression retained.
可知 $J'R'_a = 0$。由于 $a$ 在 $R'$ 中是非零因子，遂得
$J' = 0$，证明完成。
\end{proof}

\begin{lemma}
\label{flat-lemma-flatten-module-pre}
设 $S$ 为拟紧拟分离概形。设 $X \to S$ 为概形态射。
设 $\mathcal{F}$ 为 $X$ 上的拟凝聚模。设 $U \subset S$ 为拟紧开集。假设
\begin{enumerate}
\item $X \to S$ 仿射、有限表示、平坦，且纤维几何整；
\item $\mathcal{F}$ 是有限型模；
\item $\mathcal{F}_U$ 有限表示；
\item $\mathcal{F}$ 在位于 $U$ 中各点之上的纤维的所有一般点处在 $S$ 上平坦。
\end{enumerate}
则存在 $U$-容许爆破 $S' \to S$ 和开子概形 $V \subset X_{S'}$，使得：
(a) $\mathcal{F}$ 的严格变换 $\mathcal{F}'$ 限制为有限局部自由
$\mathcal{O}_V$-模；(b) $V \to S'$ 为满射。
\end{lemma}

\begin{proof}
给定满足引理假设的 $\mathcal{F}/X/S$ 和 $U \subset S$，以 $P$ 表示性质：
“$\mathcal{F}$ 在纤维的所有一般点处在 $S$ 上平坦”。
显然，$P$ 在严格变换下保持不变，见除子论引理
\ref{divisors-lemma-strict-transform-flat} 和态射论引理
\ref{morphisms-lemma-base-change-module-flat}。同样显然，$P$ 关于 $S$ 是局部的。
因此注 \ref{flat-remark-successive-blowups} 的全部观察均适用于本引理提出的问题。

\medskip\noindent
考虑函数 $r : U \to \mathbf{Z}_{\geq 0}$，它把 $u \in U$ 映为整数
$$
r(u) = \dim_{\kappa(\xi_u)}(\mathcal{F}_{\xi_u} \otimes \kappa(\xi_u))
$$
其中 $\xi_u$ 是纤维 $X_u$ 的一般点。
由态射论进一步的引理
\ref{more-morphisms-lemma-flat-and-free-at-point-fibre}
以及 $X_S$ 中开集在 $S$ 中的像为开集这一事实，可知 $r(u)$ 局部常值。
因此 $U = U_0 \amalg U_1 \amalg \ldots \amalg U_c$ 是开闭子概形的有限不交并，
且 $r$ 在 $U_i$ 上恒等于 $i$。
由除子论引理 \ref{divisors-lemma-separate-disjoint-opens-by-blowing-up}，
可以找到一个 $U$-容许爆破，将 $S$ 分解为两个概形的不交并，其中第一个包含
$U_0$，第二个包含 $U_1 \cup \ldots \cup U_c$。
再重复此过程 $c - 1$ 次，可设 $S$ 是不交并
$S = S_0 \amalg S_1 \amalg \ldots \amalg S_c$，且 $U_i \subset S_i$。
因此可设上述函数 $r$ 为常值，记其值为 $r$。

\medskip\noindent
由注 \ref{flat-remark-successive-blowups}，可设存在有效 Cartier 除子
$D \subset S$，其支集为 $S \setminus U$。
再次应用注 \ref{flat-remark-successive-blowups}，并结合除子论引理
\ref{divisors-lemma-characterize-effective-Cartier-divisor}，可设
$S = \Spec(R)$ 且 $D = \Spec(R/(f))$，其中 $f \in R$ 为非零因子。
此情形由引理 \ref{flat-lemma-helper-blowup-affine-space} 处理。
\end{proof}

\begin{lemma}
\label{flat-lemma-trick-fitting-ideal}
设 $A \to C$ 为秩为 $d$ 的有限局部自由环同态。
设 $h \in C$，使 $C_h$ 在 $A$ 上 \'etale。
设 $J \subset C$ 为理想。置 $I = \text{Fit}_0(C/J)$，其中将
$C/J$ 视为有限 $A$-模。则对某个理想 $J' \subset C_h$，有
$IC_h = JJ'$。若 $J$ 有限生成，则 $I$ 和 $J'$ 也有限生成。
\end{lemma}

\begin{proof}
我们将使用 Fitting 理想的基本性质，见《代数进阶》引理
\ref{more-algebra-lemma-fitting-ideal-basics}。
于是 $IC$ 是 $C/J \otimes_A C$ 的 Fitting 理想。
注意，$C \to C \otimes_A C$，$c \mapsto 1 \otimes c$ 有一个截面
（乘法映射）。按假设，$C \to C \otimes_A C$ 在 $\Spec(C_h)$ 经该截面所得像中
的每个素理想处都是 \'etale 的。
因此乘法映射 $C \otimes_A C_h \to C_h$ 是 \'etale 的，特别地是平坦的，
见代数引理 \ref{algebra-lemma-map-between-etale}。
因此存在一个这样的 $C_h$-代数，使得作为 $C_h$-代数有
$C \otimes_A C_h \cong C_h \oplus C'$，见代数引理
\ref{algebra-lemma-surjective-flat-finitely-presented}。
于是对某个理想 $I' \subset C'$，作为 $C_h$-模有
$(C/J) \otimes_A C_h \cong (C_h/J_h) \oplus C'/I'$。
故 $IC_h = JJ'$，其中 $J' = \text{Fit}_0(C'/I')$，并且
% P05-EN-CH38-0060: source line 8872 switches from C'/I' to C'/J'; frozen quotient retained.
将 $C'/J'$ 视为 $C_h$-模。
\end{proof}

\begin{lemma}
\label{flat-lemma-push-ideal}
设 $A \to B$ 为 \'etale 环同态。设 $a \in A$ 为非零因子。
设 $J \subset B$ 为有限型理想，满足 $V(J) \subset V(aB)$。
对每个 $\mathfrak q \subset B$，存在有限型理想 $I \subset A$，
满足 $V(I) \subset V(a)$，并存在 $g \in B$，$g \not \in \mathfrak q$，
使得对某个有限型理想 $J' \subset B_g$ 有 $IB_g = JJ'$。
\end{lemma}

\begin{proof}
可将 $B$ 替换为在某个元素 $g \in B$ 处的主元局部化，其中
$g \not \in \mathfrak q$。因此可设 $B$ 为标准 \'etale，见代数命题
\ref{algebra-proposition-etale-locally-standard}。
于是可设 $B$ 是 $C = A[x]/(f)$ 的一个局部化，其中
$f \in A[x]$ 为某个次数 $d$ 的首一多项式。
设对某个 $h \in C$ 有 $B = C_h$。选择在 $B$ 上生成 $J$ 的元素
$h_1, \ldots, h_n \in C$。
条件 $V(J) \subset V(aB)$ 意味着对某个充分大的 $m$，在 $B$ 中有
$a^m = \sum b_i h_i$。置 $h_{n + 1} = a^m$。
如引理 \ref{flat-lemma-trick-fitting-ideal} 中一样，取
% P05-EN-CH38-0061: source line 8897 uses r+1 after defining n generators; frozen indices retained.
$I = \text{Fit}_0(C/(h_1, \ldots, h_{r + 1}))$。
由于模 $C/(h_1, \ldots, h_{r + 1})$ 被 $a^m$ 零化，
可知 $a^{dm} \in I$，从而 $V(I) \subset V(a)$。
\end{proof}

\begin{lemma}
\label{flat-lemma-flatten-module-etale-localize}
设 $S$ 为拟紧拟分离概形。设 $X \to S$ 为概形态射。
设 $\mathcal{F}$ 为 $X$ 上的拟凝聚模。设 $U \subset S$ 为拟紧开集。
假设存在有限多个交换图
$$
\xymatrix{
& X_i \ar[r]_{j_i} \ar[d] & X \ar[d] \\
S_i^* \ar[r] & S_i \ar[r]^{e_i} & S
}
$$
其中
\begin{enumerate}
\item $e_i : S_i \to S$ 是拟紧 \'etale 态射，且
$S = \bigcup e_i(S_i)$；
\item $j_i : X_i \to X$ 是 \'etale 态射，且
$X = \bigcup j_i(X_i)$；
\item $S^*_i \to S_i$ 是 $e_i^{-1}(U)$-容许爆破，使
$j_i^*\mathcal{F}$ 的严格变换 $\mathcal{F}_i^*$ 在 $S^*_i$ 上平坦。
\end{enumerate}
则存在 $U$-容许爆破 $S' \to S$，使 $\mathcal{F}$ 的严格变换在 $S'$ 上平坦。
\end{lemma}

\begin{proof}
我们断言，本引理的假设在 $U$-容许爆破下保持不变。
具体地，设 $b : S' \to S$ 是沿拟凝聚理想层 $\mathcal{I}$ 的
$U$-容许爆破。此外，设 $S^*_i \to S_i$ 是沿拟凝聚理想层
$\mathcal{J}_i$ 的爆破。
则对 $\mathcal{F}$ 关于爆破 $S' \to S$ 的严格变换 $\mathcal{F}'$，
态射族 $e'_i : S'_i = S_i \times_S S' \to S'$ 和
$j'_i : X_i' = X_i \times_S S' \to X \times_S S'$ 满足条件 (1)、(2)、(3)。
首先，注意 $S_i'$ 是 $S_i$ 沿 $\mathcal{I}$ 的拉回的爆破，见除子论引理
\ref{divisors-lemma-flat-base-change-blowing-up}。
其次，考虑 $S_i'$ 沿理想 $\mathcal{J}_i$ 的拉回的爆破
$S_i^{\prime *} \to S_i'$。
由除子论引理 \ref{divisors-lemma-blowing-up-two-ideals}，得到交换图
$$
\xymatrix{
S_i^{\prime *} \ar[r] \ar[rd] \ar[d] & S'_i \ar[d] \\
S_i^* \ar[r] & S_i
}
$$
且上图中的所有态射都是爆破。因此由除子论引理
\ref{divisors-lemma-strict-transform-flat} 和
\ref{divisors-lemma-strict-transform-composition-blowups}，可知
% P05-EN-CH38-0062: source lines 8963-8967 use an undefined F_i' and mismatched blowup base; formula retained.
\begin{align*}
& (j'_i)^*\mathcal{F}'\text{ 关于 }
S_i^{\prime *} \to S_i'\text{ 的严格变换} \\
= &
j_i^*\mathcal{F}\text{ 关于 }
S_i^{\prime *} \to S_i\text{ 的严格变换} \\
= &
\mathcal{F}_i'\text{ 关于 }
S_i^{\prime *} \to S_i'\text{ 的严格变换} \\
= &
\mathcal{F}_i^*\text{ 经由 }
X_i \times_{S_i} S_i^{\prime *} \to X_i\text{ 的拉回}
\end{align*}
因此它在 $S_i^{\prime *}$ 上平坦
（态射论引理 \ref{morphisms-lemma-base-change-module-flat}）。
由此可见，注 \ref{flat-remark-successive-blowups} 的全部观察均适用于如下问题：
在本引理的假设下，寻找一个 $U$-容许爆破，使 $\mathcal{F}$ 的严格变换
在基上变得平坦。特别地，可设 $S \setminus U$ 是某个有效 Cartier 除子
$D \subset S$ 的支集。再次应用注 \ref{flat-remark-successive-blowups}，并结合
除子论引理 \ref{divisors-lemma-characterize-effective-Cartier-divisor}，可设
$S = \Spec(A)$ 且 $D = \Spec(A/(a))$，其中 $a \in A$ 为非零因子。

\medskip\noindent
选取一个 $i$ 和 $s \in S_i$。由引理 \ref{flat-lemma-push-ideal}，可以找到开邻域
$s \in W_i \subset S_i$ 和有限型拟凝聚理想
$\mathcal{I} \subset \mathcal{O}_S$，使得对某个有限型拟凝聚理想
$\mathcal{J}'_i \subset \mathcal{O}_{W_i}$ 有
$\mathcal{I} \cdot \mathcal{O}_{W_i} = \mathcal{J}_i \mathcal{J}'_i$，
并且 $V(\mathcal{I}) \subset V(a) = S \setminus U$。
由于 $S_i$ 拟紧，可将 $S_i$ 替换为由这类开集组成的有限族
$W_1, \ldots, W_n$，并假设对每个 $i$ 存在拟凝聚理想层
$\mathcal{I}_i \subset \mathcal{O}_S$，使得对某个有限型拟凝聚理想
$\mathcal{J}'_i \subset \mathcal{O}_{S_i}$ 有
$\mathcal{I}_i \cdot \mathcal{O}_{S_i} = \mathcal{J}_i \mathcal{J}'_i$。
如证明第一段的讨论，考虑 $S$ 沿乘积
$\mathcal{I}_1 \ldots \mathcal{I}_n$ 的爆破 $S'$
（按构造，该爆破是 $U$-容许的）。$S' \to S$ 到 $S_i$ 的基变换是沿
$$
\mathcal{J}_i \cdot
\mathcal{J}'_i \mathcal{I}_1 \ldots \hat{\mathcal{I}_i} \ldots \mathcal{I}_n
$$
的爆破；它经由给定爆破 $S_i^* \to S_i$ 分解
（除子论引理 \ref{divisors-lemma-blowing-up-two-ideals}）。
在上图的记号中，这意味着 $S_i^{\prime *} = S_i'$。
因此以 $S'$ 替换 $S$ 后，达到 $j_i^*\mathcal{F}$ 在 $S_i$ 上平坦的情形。
于是 $j_i^*\mathcal{F}$ 在 $S$ 上平坦，见引理
\ref{flat-lemma-etale-flat-up-down}。由态射论引理
\ref{morphisms-lemma-flat-permanence}，可知 $\mathcal{F}$ 在 $S$ 上平坦。
\end{proof}

\begin{theorem}
\label{flat-theorem-flatten-module}
设 $S$ 为拟紧拟分离概形。设 $X$ 为 $S$ 上的概形。
设 $\mathcal{F}$ 为 $X$ 上的拟凝聚模。设 $U \subset S$ 为拟紧开集。假设
\begin{enumerate}
\item $X$ 拟紧；
\item $X$ 在 $S$ 上局部有限表示；
\item $\mathcal{F}$ 是有限型模；
\item $\mathcal{F}_U$ 有限表示；
\item $\mathcal{F}_U$ 在 $U$ 上平坦。
\end{enumerate}
则存在 $U$-容许爆破 $S' \to S$，使 $\mathcal{F}$ 的严格变换
$\mathcal{F}'$ 是有限表示的 $\mathcal{O}_{X \times_S S'}$-模，
且在 $S'$ 上平坦。
\end{theorem}

\begin{proof}
先证明可以找到一个 $U$-容许爆破，使严格变换平坦。
该问题在源和目标上关于 \'etale 拓扑局部；精确陈述见引理
\ref{flat-lemma-flatten-module-etale-localize}。
特别地，可设 $S = \Spec(R)$ 和 $X = \Spec(A)$ 为仿射。
对 $s \in S$，记 $\mathcal{F}_s = \mathcal{F}|_{X_s}$
（即 $\mathcal{F}$ 到纤维的拉回）。由于 $X \to S$ 有限型，
$d = \max_{s \in S} \dim(\text{Supp}(\mathcal{F}_s))$ 是整数。
我们对 $d$ 作归纳。

\medskip\noindent
设 $x \in X$ 是位于 $s \in S$ 之上的 $X$ 中一点，且
$\dim_x(\text{Supp}(\mathcal{F}_s)) = d$。
应用引理 \ref{flat-lemma-elementary-devissage}，得到
$g : X' \to X$、$e : S' \to S$、$i : Z' \to X'$ 和 $\pi : Z' \to Y'$。
注意 $Y' \to S'$ 是仿射概形之间的光滑态射，其纤维几何不可约且维数为 $d$。
由于问题关于 \'etale 拓扑局部，只需对 $g^*\mathcal{F}/X'/S'$ 证明定理。
由于 $i : Z' \to X'$ 是有限表示的闭浸入（且严格变换与仿射推前交换，见除子论引理
\ref{divisors-lemma-strict-transform-affine}），只需对 $\mathcal{G}$ 证明扁平化结论。
由于 $\pi$ 有限（因而也仿射），只需对
$\pi_*\mathcal{G}/Y'/S'$ 证明扁平化结论。因此可设 $X \to S$
是仿射概形之间的光滑态射，其纤维几何不可约且维数为 $d$。

\medskip\noindent
接着，应用引理 \ref{flat-lemma-flatten-module-pre} 中的爆破。
这样达到如下情形：存在满射到 $S$ 的开集 $V \subset X$，使
$\mathcal{F}|_V$ 有限局部自由。设 $\xi \in X$ 为 $X_s$ 的一般点。令
$r = \dim_{\kappa(\xi)} \mathcal{F}_\xi \otimes \kappa(\xi)$。
选择映射 $\alpha : \mathcal{O}_X^{\oplus r} \to \mathcal{F}$，它诱导同构
$\kappa(\xi)^{\oplus r} \to \mathcal{F}_\xi \otimes \kappa(\xi)$。
由于 $\mathcal{F}$ 在 $V$ 上局部自由，可找到 $\xi$ 的一个开邻域 $W$，
使 $\alpha$ 在其上为同构。将 $S$ 缩小为 $s$ 的一个仿射开邻域，使
$W \to S$ 为满射。设 $\mathcal{F}$ 是由 $A$-模 $N$ 所伴随的拟凝聚模。
由于 $\mathcal{F}$ 在纤维的所有一般点处在 $S$ 上平坦
（事实上在 $W$ 的所有点处如此），可知
$$
\alpha_\mathfrak p : A_\mathfrak p^{\oplus r} \to N_\mathfrak p
$$
对 $R$ 的所有素理想 $\mathfrak p$ 都普遍单射，见引理
\ref{flat-lemma-induction-step}。因此 $\alpha$ 普遍单射，见代数引理
\ref{algebra-lemma-universally-injective-check-stalks}。
置 $\mathcal{H} = \Coker(\alpha)$。
由除子论引理 \ref{divisors-lemma-strict-transform-universally-injective}，
给定 $U$-容许爆破 $S' \to S$ 后，严格变换
$\mathcal{F}'$ 和 $\mathcal{H}'$ 嵌入正合序列
$$
0 \to \mathcal{O}_{X \times_S S'}^{\oplus r} \to \mathcal{F}'
\to \mathcal{H}' \to 0
$$
因此引理 \ref{flat-lemma-induction-step} 还说明：$\mathcal{F}'$ 在一点
$x'$ 处平坦，当且仅当 $\mathcal{H}'$ 在该点处平坦。
特别地，$\mathcal{H}_U$ 在 $U$ 上平坦，且 $\mathcal{H}_U$ 是有限表示的模。
可将归纳假设应用于 $\mathcal{H}$，从而存在一个 $U$-容许爆破，
使严格变换 $\mathcal{H}'$ 如所需般平坦。

\medskip\noindent
为完成定理证明，还必须说明 $\mathcal{F}'$ 是有限表示的模
（可能还要再作一个 $U$-容许爆破）。由于可设 $U \subset S$ 概形论稠密
（见注 \ref{flat-remark-successive-blowups} 第三段），这由引理
\ref{flat-lemma-flat-finite-type-finitely-presented-over-dense-open} 得出。
定理证毕。
\end{proof}





\section{应用}
\label{flat-section-applications}

\noindent
本节应用上面的一些结果。

\begin{lemma}
\label{flat-lemma-flat-after-blowing-up}
设 $S$ 为拟紧拟分离概形。设 $X$ 为 $S$ 上的概形。
设 $U \subset S$ 为拟紧开集。假设
\begin{enumerate}
\item $X \to S$ 有限型且拟分离；
\item $X_U \to U$ 平坦且局部有限表示。
\end{enumerate}
则存在 $U$-容许爆破 $S' \to S$，使 $X$ 的严格变换在 $S'$ 上平坦且有限表示。
\end{lemma}

\begin{proof}
按假设，$X \to S$ 拟紧且拟分离，所以 $X$ 关于爆破 $S' \to S$ 的严格变换
也拟紧且拟分离。因此，为证明引理，只需找到一个 $U$-容许爆破，
使严格变换平坦且局部有限表示。
令 $X = W_1 \cup \ldots \cup W_n$ 为有限仿射开覆盖。
若能找到 $U$-容许爆破 $S_i \to S$，使 $W_i$ 的严格变换平坦且局部有限表示，
则存在支配所有 $S_i \to S$ 的 $U$-容许爆破 $S' \to S$，并满足要求
（见除子论引理 \ref{divisors-lemma-dominate-admissible-blowups}；
另见注 \ref{flat-remark-successive-blowups}）。因此可设 $X$ 仿射。

\medskip\noindent
假设 $X$ 仿射。由态射论引理
\ref{morphisms-lemma-quasi-affine-finite-type-over-S}，可选取 $S$ 上的浸入
$j : X \to \mathbf{A}^n_S$。
取拟紧开子概形 $V \subset \mathbf{A}^n_S$，使 $j$ 诱导 $S$ 上的闭浸入
$i : X \to V$。将定理 \ref{flat-theorem-flatten-module} 应用于
$V \to S$ 和拟凝聚模 $i_*\mathcal{O}_X$，得到 $U$-容许爆破
$S' \to S$，使 $i_*\mathcal{O}_X$ 的严格变换在 $S'$ 上平坦，且作为
$\mathcal{O}_{V \times_S S'}$-模有限表示。
令 $X'$ 为 $X$ 关于 $S' \to S$ 的严格变换，令
$i' : X' \to V \times_S S'$ 为诱导态射。
由于严格变换与沿仿射态射的推前交换（除子论引理
\ref{divisors-lemma-strict-transform-affine}），可知
% P05-EN-CH38-0063: source line 9167 says flat over S after a base change to S'; frozen base retained.
$i'_*\mathcal{O}_{X'}$ 在 $S$ 上平坦，且作为
$\mathcal{O}_{V \times_S S'}$-模有限表示。这推出引理。
\end{proof}

\begin{lemma}
\label{flat-lemma-finite-after-blowing-up}
设 $S$ 为拟紧拟分离概形。设 $X$ 为 $S$ 上的概形。
设 $U \subset S$ 为拟紧开集。假设
\begin{enumerate}
\item $X \to S$ 适当；
\item $X_U \to U$ 有限局部自由。
\end{enumerate}
则存在 $U$-容许爆破 $S' \to S$，使 $X$ 的严格变换在 $S'$ 上有限局部自由。
\end{lemma}

\begin{proof}
由引理 \ref{flat-lemma-flat-after-blowing-up}，可设 $X \to S$ 平坦且有限表示。
必要时以一个 $U$-容许爆破替换 $S$ 后，可设 $U \subset S$ 概形论稠密。
% P05-EN-CH38-0064: source lines 9190-9193 use an undefined f for the morphism X -> S; symbol retained.
于是由引理 \ref{flat-lemma-proper-flat-finite-over-dense-open}，$f$ 有限。
进而由态射论引理 \ref{morphisms-lemma-finite-flat}，$f$ 有限局部自由。
\end{proof}

\begin{lemma}
\label{flat-lemma-zariski-after-blowup}
设 $\varphi : X \to S$ 为有限型分离态射，其中 $S$ 拟紧且拟分离。
设 $U \subset S$ 为拟紧开集，使 $\varphi^{-1}U \to U$ 为同构。
则存在 $U$-容许爆破 $S' \to S$，使 $X$ 的严格变换 $X'$
同构于 $S'$ 的一个开子概形。
\end{lemma}

\begin{proof}
注 \ref{flat-remark-successive-blowups} 的讨论适用。
因此可以先作一次 $U$-容许爆破，并假设补集 $S \setminus U$ 是某个
有效 Cartier 除子 $D$ 的支集。特别地，$U$ 在 $S$ 中概形论稠密。
接着再作一次 $U$-容许爆破，达到 $X \to S$ 平坦且有限表示的情形，
见引理 \ref{flat-lemma-flat-after-blowing-up}。此时结论由引理
\ref{flat-lemma-zariski} 得出。
\end{proof}

\noindent
下一个引理说明，适当修改可由一个爆破支配。

\begin{lemma}
\label{flat-lemma-dominate-modification-by-blowup}
设 $\varphi : X \to S$ 为适当态射，其中 $S$ 拟紧且拟分离。
设 $U \subset S$ 为拟紧开集，使 $\varphi^{-1}U \to U$ 为同构。
则存在支配 $X$ 的 $U$-容许爆破 $S' \to S$；也就是说，爆破态射存在分解
$S' \to X \to S$。
\end{lemma}

\begin{proof}
注 \ref{flat-remark-successive-blowups} 的讨论适用。
因此可以先作一次 $U$-容许爆破，并假设补集 $S \setminus U$ 是某个
有效 Cartier 除子 $D$ 的支集。特别地，$U$ 在 $S$ 中概形论稠密。
选择另一个 $U$-容许爆破 $S' \to S$，使 $X$ 的严格变换 $X'$
是 $S'$ 的开子概形，见引理 \ref{flat-lemma-zariski-after-blowup}。
由于 $X' \to S'$ 适当，且 $U \subset S'$ 稠密，可知 $X' = S'$。
略去一些细节。
\end{proof}

\begin{lemma}
\label{flat-lemma-get-section-after-blowup}
设 $S$ 为概形。设 $U \subset W \subset S$ 为开子概形。
设 $f : X \to W$ 和 $s : U \to X$ 为态射，且
$f \circ s = \text{id}_U$。假设
\begin{enumerate}
\item $f$ 适当；
\item $S$ 拟紧且拟分离；
\item $U$ 和 $W$ 拟紧。
\end{enumerate}
则存在 $U$-容许爆破 $b : S' \to S$，以及延拓 $s$ 的态射
$s' : b^{-1}(W) \to X$，满足 $f \circ s' = b|_{b^{-1}(W)}$。
\end{lemma}

\begin{proof}
我们可以并且确实用 $s$ 的概形论像替换 $X$。
于是 $X \to W$ 在 $U$ 上是同构，见态射论引理
\ref{morphisms-lemma-scheme-theoretic-image-of-partial-section}.
由引理 \ref{flat-lemma-dominate-modification-by-blowup}，存在 $U$-容许爆破
$W' \to W$ 和 $s$ 的延拓 $W' \to X$。
最后应用除子论引理 \ref{divisors-lemma-extend-admissible-blowups}，
将 $W' \to W$ 延拓为 $S$ 的 $U$-容许爆破，证明完成。
\end{proof}









\section{紧化}
\label{flat-section-compactify}

\noindent
设 $S$ 为拟紧拟分离概形。若存在到一个在 $S$ 上适当的概形
$\overline{X}$ 的拟紧开浸入 $X \to \overline{X}$，则称 $S$ 上的概形
$X$ {\it 在 $S$ 上有紧化}，或称 {\it 在 $S$ 上可紧化}。
若 $X$ 在 $S$ 上有紧化，则 $X \to S$ 分离且有限型。
Nagata 的一个定理说明逆命题也成立，见 \cite{Lutkebohmert}、
\cite{Conrad-Nagata}、\cite{Nagata-1}、\cite{Nagata-2}、
\cite{Nagata-3} 和 \cite{Nagata-4}。
我们将在下一节证明该定理，见定理 \ref{flat-theorem-nagata}。

\medskip\noindent
设 $S$ 为拟紧拟分离概形。设 $X \to S$ 为有限型分离态射。
{\it $X$ 在 $S$ 上的紧化范畴} 定义如下：
\begin{enumerate}
\item 对象是 $S$ 上的开浸入 $j : X \to \overline{X}$，其中
$\overline{X} \to S$ 适当。
\item 态射 $(j' : X \to \overline{X}') \to (j : X \to \overline{X})$
是 $S$ 上的概形态射 $f : \overline{X}' \to \overline{X}$，满足
$f \circ j' = j$。
\end{enumerate}
若 $j : X \to \overline{X}$ 是紧化，则 $j$ 是拟紧开浸入，见概形论注
\ref{schemes-remark-quasi-compact-and-quasi-separated}。

\medskip\noindent
{\bf 警告。} 我们{\it 不}假设紧化 $j : X \to \overline{X}$ 的像稠密。
因此，若 $f : \overline{X}' \to \overline{X}$ 是紧化之间的态射，
未必有 $f^{-1}(j(X)) = j'(X)$。

\begin{lemma}
\label{flat-lemma-compactifications-cofiltered}
设 $S$ 为拟紧拟分离概形。设 $X$ 为 $S$ 上可紧化的概形。
\begin{enumerate}
\item[(a)] $X$ 在 $S$ 上的紧化范畴是余滤的。
\item[(b)] 由满足下述条件的紧化 $j : X \to \overline{X}$ 组成的
全子范畴是初始的：$j(X)$ 在 $\overline{X}$ 中稠密且概形论稠密
（范畴论定义 \ref{categories-definition-initial}）。
\item[(c)] 若 $f : \overline{X}' \to \overline{X}$ 是 $X$ 的紧化之间的态射，
且 $j'(X)$ 在 $\overline{X}'$ 中稠密，则
$f^{-1}(j(X)) = j'(X)$。
\end{enumerate}
\end{lemma}

\begin{proof}
为证明 (a)，必须验证范畴论定义
\ref{categories-definition-codirected} 的条件 (1)、(2)、(3)。
条件 (1) 成立，正是因为我们假设 $X$ 可紧化。
设 $j_i : X \to \overline{X}_i$，$i = 1, 2$ 为两个紧化。
考虑 $(j_1, j_2) : X \to \overline{X}_1 \times_S \overline{X}_2$
的概形论像 $\overline{X}$。这确定了第三个紧化
$j : X \to \overline{X}$，它支配两个 $j_i$：
$$
\xymatrix{
(X, \overline{X}_1) & (X, \overline{X}) \ar[l] \ar[r] & (X, \overline{X}_2)
}
$$
因此 (2) 成立。设 $f_1, f_2 : \overline{X}_1 \to \overline{X}_2$
为紧化 $j_i : X \to \overline{X}_i$，$i = 1, 2$ 之间的两个态射。
令 $\overline{X} \subset \overline{X}_1$ 为 $f_1$ 和 $f_2$ 的等化子。
由于 $\overline{X}_2 \to S$ 分离，$\overline{X}$ 是
$\overline{X}_1$ 的闭子概形，因而在 $S$ 上适当。
此外，由于 $f_1|_X = f_2|_X = \text{id}_X$，得到开浸入
$X \to \overline{X}$。由闭浸入
$\overline{X} \to \overline{X}_1$ 给出的态射
$(X \to \overline{X}) \to (j_1 : X \to \overline{X}_1)$
等化 $f_1$ 和 $f_2$，这证明条件 (3)。

\medskip\noindent
证明 (b)。设 $j : X \to \overline{X}$ 为紧化。
若 $\overline{X}'$ 表示 $X$ 在 $\overline{X}$ 中的概形论闭包，
则由态射论引理 \ref{morphisms-lemma-quasi-compact-immersion}，
$X$ 在 $\overline{X}'$ 中稠密且概形论稠密。
这证明了范畴论定义 \ref{categories-definition-initial} 的第一个条件。
由于已证明 $X$ 的紧化范畴余滤，范畴论定义
\ref{categories-definition-initial} 的第二个条件由第一个条件推出
（我们略去这一范畴论练习的解答）。

\medskip\noindent
证明 (c)。用 $j'(X)$ 的概形论闭包替换 $\overline{X}'$
（这不改变底层拓扑空间）后，结论由态射论引理
\ref{morphisms-lemma-scheme-theoretic-image-of-partial-section} 得出。
\end{proof}

\noindent
我们还可以考虑所有紧化的范畴（其中 $X$ 可变）。
事实表明，将此范畴对在内部诱导同构的态射集合做局部化，
所得范畴等价于 $S$ 上可紧化概形的范畴。

\begin{lemma}
\label{flat-lemma-compactifyable}
设 $S$ 为拟紧拟分离概形。设 $f : X \to Y$ 为 $S$ 上的概形态射，
其中 $Y$ 在 $S$ 上分离且有限型，而 $X$ 在 $S$ 上可紧化。
则 $X$ 在 $Y$ 上有紧化。
\end{lemma}

\begin{proof}
设 $j : X \to \overline{X}$ 为 $X$ 在 $S$ 上的紧化。
令 $\overline{X}'$ 为
$(j, f) : X \to \overline{X} \times_S Y$ 的概形论像。
态射 $\overline{X}' \to Y$ 适当，因为
$\overline{X} \times_S Y \to Y$ 是 $\overline{X} \to S$ 的基变换。
另一方面，由于 $Y$ 在 $S$ 上分离，态射
$(1, f) : X \to X \times_S Y$ 是闭浸入
（概形论引理 \ref{schemes-lemma-semi-diagonal}）。
因此，将态射论引理
\ref{morphisms-lemma-scheme-theoretic-image-of-partial-section}
应用于投影 $\overline{X} \times_S Y \to \overline{X}$ 的“部分截面”
$s = (j, f)$，可知 $X \to \overline{X}'$ 是开浸入。
\end{proof}

\noindent
设 $S$ 为拟紧拟分离概形。定义{\it 紧化范畴}如下：
其对象是偶 $(X, \overline{X})$，其中 $\overline{X}$ 是在 $S$ 上适当的概形，
而 $X \subset \overline{X}$ 为拟紧开集；其态射是 $S$ 上概形态射的交换图
$$
\xymatrix{
X \ar[d] \ar[r]_f & Y \ar[d] \\
\overline{X} \ar[r]^{\overline{f}} & \overline{Y}
}
$$

\begin{lemma}
\label{flat-lemma-right-multiplicative-system}
设 $S$ 为拟紧拟分离概形。满足下述条件的态射族
$(u, \overline{u}) : (X', \overline{X}') \to (X, \overline{X})$
构成紧化范畴中箭头的一个右乘法系
（范畴论定义 \ref{categories-definition-multiplicative-system}）：
$u$ 是同构。
\end{lemma}

\begin{proof}
公理 RMS1 的验证是平凡的。下面验证 RMS2。设给定图
$$
\xymatrix{
& (X', \overline{X}') \ar[d]_{(u, \overline{u})} \\
(Y, \overline{Y}) \ar[r]^{(f, \overline{f})} & (X, \overline{X})
}
$$
其中 $u : X' \to X$ 是同构。令 $Y' = Y \times_X X'$，
其投影映射为 $v : Y' \to Y$（这是同构）。
再置 $\overline{Y}' = \overline{Y} \times_{\overline{X}} \overline{X}'$，
其投影映射为 $\overline{v} : \overline{Y}' \to \overline{Y}$。
显然，$Y' \to \overline{Y}'$ 是开浸入。图
$$
\xymatrix{
(Y', \overline{Y}') \ar[r]_{(g, \overline{g})} \ar[d]_{(v, \overline{v})} &
(X', \overline{X}') \ar[d]_{(u, \overline{u})} \\
(Y, \overline{Y}) \ar[r]^{(f, \overline{f})} & (X, \overline{X})
}
$$
说明公理 RMS2 成立。

\medskip\noindent
下面验证 RMS3。设给定紧化之间的一对态射
$(f, \overline{f}), (g, \overline{g}) :
(X, \overline{X}) \to (Y, \overline{Y})$
以及态射
$(v, \overline{v}) : (Y, \overline{Y}) \to (Y', \overline{Y}')$
使得 $v$ 是同构，并且
$(v, \overline{v}) \circ (f, \overline{f}) =
(v, \overline{v}) \circ (g, \overline{g})$。则 $f = g$。
因此，若令 $\overline{X}' \subset \overline{X}$ 为
$\overline{f}$ 和 $\overline{g}$ 的等化子，则
$(u, \overline{u}) : (X, \overline{X}') \to (X, \overline{X})$
是紧化范畴中的态射，并且如所需地满足
$(f, \overline{f}) \circ (u, \overline{u}) =
(g, \overline{g}) \circ (u, \overline{u})$。
\end{proof}

\begin{lemma}
\label{flat-lemma-invert-right-multiplicative-system}
设 $S$ 为拟紧拟分离概形。函子
$(X, \overline{X}) \mapsto X$ 定义了一个等价：从紧化范畴在引理
\ref{flat-lemma-right-multiplicative-system} 的右乘法系处的局部化
（范畴论引理 \ref{categories-lemma-right-localization}），
到 $S$ 上可紧化概形的范畴。
\end{lemma}

\begin{proof}
记 $\mathcal{C}$ 为紧化范畴，记
$Q : \mathcal{C} \to \mathcal{C}'$ 为范畴论引理
\ref{categories-lemma-properties-right-localization} 的局部化函子。
记 $\mathcal{D}$ 为 $S$ 上可紧化概形的范畴。
由刚才引用的引理及我们选择的乘法系，显然得到函子
$\mathcal{C}' \to \mathcal{D}$。该函子显然本质满。
若 $f : X \to Y$ 是可紧化概形之间的态射，则选择到一个在 $S$ 上适当的概形
中的开浸入 $Y \to \overline{Y}$，再选择到一个在 $\overline{Y}$ 上适当的概形
$\overline{X}$ 中的嵌入 $X \to \overline{X}$
（将引理 \ref{flat-lemma-compactifyable} 应用于 $X \to \overline{Y}$ 即可）。
这给出紧化之间的态射
$(X, \overline{X}) \to (Y, \overline{Y})$，它产生给定态射 $X \to Y$。
最后，设局部化范畴中给定一对源和目标相同的态射，记为
$$
a = ((f, \overline{f}) : (X', \overline{X}') \to (Y, \overline{Y}),
(u, \overline{u}) : (X', \overline{X}') \to (X, \overline{X}))
$$
以及
$$
b = ((g, \overline{g}) : (X'', \overline{X}'') \to (Y, \overline{Y}),
(v, \overline{v}) : (X'', \overline{X}'') \to (X, \overline{X}))
$$
它们在 $S$ 上产生同一个态射 $X \to Y$；换言之，
$f \circ u^{-1} = g \circ v^{-1}$。
由范畴论引理 \ref{categories-lemma-morphisms-right-localization}，可设
$(X', \overline{X}') = (X'', \overline{X}'')$ 且
$(u, \overline{u}) = (v, \overline{v})$。
此时可考虑 $\overline{f}$ 和 $\overline{g}$ 的等化子
$\overline{X}''' \subset \overline{X}'$。
态射 $(w, \overline{w}) : (X', \overline{X}''') \to (X', \overline{X}')$
属于该乘法子集；通过预复合 $(w, \overline{w})$，可知局部化范畴中
$a = b$。
\end{proof}






\section{Nagata 紧化}
\label{flat-section-compactifications}

\noindent
本节证明在第 \ref{flat-section-compactify} 节中宣布的定理。

\begin{lemma}
\label{flat-lemma-check-separated}
设 $X \to S$ 是概形态射。若 $X = U \cup V$ 是一个开覆盖，
其中 $U \to S$ 和 $V \to S$ 均为分离态射，且
$U \cap V \to U \times_S V$ 为闭态射，则
$X \to S$ 为分离态射。
\end{lemma}

\begin{proof}
证明从略。提示：利用 $X \times_S X$ 的开覆盖
$U \times_S U$、$U \times_S V$、$V \times_S U$ 和
$V \times_S V$，验证 $\Delta : X \to X \times_S X$ 是闭态射。
\end{proof}

\begin{lemma}
\label{flat-lemma-separate-disjoint-locally-closed-by-blowing-up}
设 $X$ 是拟紧且拟分离概形，$U \subset X$ 是拟紧开子集。
\begin{enumerate}
\item 若 $Z_1, Z_2 \subset X$ 是有限表示的闭子概形，且
$Z_1 \cap Z_2 \cap U = \emptyset$，则存在一个
$U$-容许爆破 $X' \to X$，使 $Z_1$ 与 $Z_2$ 的严格变换不相交。
\item 若 $T_1, T_2 \subset U$ 是互不相交的构造性闭子集，
则存在一个 $U$-容许爆破 $X' \to X$，使 $T_1$ 与 $T_2$
的闭包不相交。
\end{enumerate}
\end{lemma}

\begin{proof}
(1) 的证明。假设 $Z_i \to X$ 为有限表示，意味着 $Z_i$ 的
拟相干理想层 $\mathcal{I}_i$ 为有限型；参见
《态射》，引理 \ref{morphisms-lemma-closed-immersion-finite-presentation}。
令 $Z \subset X$ 为由乘积 $\mathcal{I}_1 \mathcal{I}_2$
切出的闭子概形。注意，$Z \cap U$ 是 $Z_1 \cap U$ 与
$Z_2 \cap U$ 的不交并。由《因子》，引理
\ref{divisors-lemma-separate-disjoint-opens-by-blowing-up}
，存在一个 $U \cap Z$-容许爆破 $Z' \to Z$，使 $Z_1$ 与
$Z_2$ 的严格变换不相交。记 $Y \subset Z$ 为此爆破的中心。
作为 $Y \to Z$ 与 $Z \to X$ 的复合，$Y \to X$
是有限表示的闭浸入（《因子》，定义
\ref{divisors-definition-admissible-blowup}
和 《态射》，引理
\ref{morphisms-lemma-composition-finite-presentation}）。
因此，以 $Y$ 为中心的爆破 $X' \to X$ 是 $U$-容许爆破。
由严格变换的一般性质，$Z_1, Z_2$ 关于 $X' \to X$
的严格变换与 $Z_1, Z_2$ 关于 $Z' \to Z$ 的严格变换相同；参见
《因子》，引理 \ref{divisors-lemma-strict-transform}。
这就证明了 (1)。

\medskip\noindent
(2) 的证明。由《性质》，引理
\ref{properties-lemma-quasi-coherent-finite-type-ideals}
，存在有限型拟相干理想层
$\mathcal{J}_i \subset \mathcal{O}_U$，使
$T_i = V(\mathcal{J}_i)$（作为集合）。由《性质》，引理
\ref{properties-lemma-extend}，存在有限型拟相干理想层
$\mathcal{I}_i \subset \mathcal{O}_X$，其在 $U$ 上的限制为
$\mathcal{J}_i$。将 (1) 的结论应用于闭子概形
$Z_i = V(\mathcal{I}_i)$ 即得结论。
\end{proof}

\begin{lemma}
\label{flat-lemma-blowup-iso-along}
设 $f : X \to Y$ 是拟紧且拟分离概形之间的适当态射。
设 $V \subset Y$ 是拟紧开子集，并令 $U = f^{-1}(V)$。
设 $T \subset V$ 是闭子集，并且 $f|_U : U \to V$
在 $V$ 中 $T$ 的某个开邻域上是同构。则存在一个
$V$-容许爆破 $Y' \to Y$，使 $f$ 的严格变换
$f' : X' \to Y'$ 在 $T$ 于 $Y'$ 中的闭包的某个开邻域上是同构。
\end{lemma}

\begin{proof}
令 $T' \subset V$ 为使 $f|_U$ 成为同构的最大开集之补。
于是 $T', T$ 在 $V$ 中均为闭集，且 $T \cap T' = \emptyset$。
由于 $V$ 是谱拓扑空间，可以找到构造性闭子集 $T_c, T'_c$，
使 $T \subset T_c$、$T' \subset T'_c$ 且
$T_c \cap T'_c = \emptyset$（先取 $V$ 中一个包含 $T'$ 而不与
$T$ 相交的拟紧开集 $W$，并令 $T_c = V \setminus W$；再取
$V$ 中一个包含 $T_c$ 而不与 $T'$ 相交的拟紧开集 $W'$，
并令 $T'_c = V \setminus W'$）。
由引理 \ref{flat-lemma-separate-disjoint-locally-closed-by-blowing-up}，
在用一个 $V$-容许爆破替换 $Y$ 后，可以假设 $T_c$ 与 $T'_c$
在 $Y$ 中的闭包不相交。令
$Y_0 = Y \setminus \overline{T}'_c$、$V_0 = V \setminus T'_c$、
$U_0 = U \times_V V_0$ 以及 $X_0 = X \times_Y Y_0$。
由于 $U_0 \to V_0$ 是同构，由引理
\ref{flat-lemma-zariski-after-blowup}，可以找到一个
$V_0$-容许爆破 $Y'_0 \to Y_0$，使 $X_0$ 的严格变换
$X'_0$ 同构地映到 $Y'_0$。由《因子》，引理
\ref{divisors-lemma-extend-admissible-blowups}，存在一个
$V$-容许爆破 $Y' \to Y$，其在 $Y_0$ 上的限制是
$Y'_0 \to Y_0$。若以 $f' : X' \to Y'$ 表示 $f$ 的严格变换，
则所需结论成立，因为 $f'$ 在 $Y'_0$ 上的限制是同构。
\end{proof}

\begin{lemma}
\label{flat-lemma-find-common-blowups}
设 $S$ 是拟紧且拟分离概形。设 $U \to X_1$ 和
$U \to X_2$ 是 $S$ 上概形的开浸入，并假设
$U$、$X_1$、$X_2$ 在 $S$ 上均为有限型且分离。
则存在 $S$ 上概形的交换图
$$
\xymatrix{
X_1' \ar[d] \ar[r] & X & X_2' \ar[l] \ar[d] \\
X_1 & U \ar[l] \ar[lu] \ar[u] \ar[ru] \ar[r] & X_2
}
$$
其中 $X_i' \to X_i$ 是 $U$-容许爆破，$X_i' \to X$
是开浸入，且 $X$ 在 $S$ 上分离并为有限型。
\end{lemma}

\begin{proof}
在整个证明中，所有概形都在 $S$ 上分离且为有限型。
特别地，这些概形拟紧且拟分离，而它们之间的态射拟紧且分离。
参见《概形》，第 \ref{schemes-section-quasi-compact} 节和第 \ref{schemes-section-separation-axioms} 节。
我们将使用如下事实：若 $U \to W$ 是这类 $S$ 上概形的浸入，
则 $U$ 在 $W$ 中的概形论像 $Z$ 是 $W$ 的闭子概形，
$U \to Z$ 是开浸入，$U \subset Z$ 在概形论意义下稠密，
并且 $U \subset Z$ 在拓扑意义下稠密。参见《态射》，引理
\ref{morphisms-lemma-quasi-compact-immersion}。

\medskip\noindent
令 $X_{12} \subset X_1 \times_S X_2$ 为
$U \to X_1 \times_S X_2$ 的概形论像。由《态射》，引理
\ref{morphisms-lemma-scheme-theoretic-image-of-partial-section}，
投影 $p_i : X_{12} \to X_i$ 诱导同构
$p_i^{-1}(U) \to U$。选取 $U$-容许爆破
$X_i^i \to X_i$，使 $X_{12}$ 的严格变换 $X_{12}^i$
同构于 $X_i^i$ 的一个开子概形；参见引理
\ref{flat-lemma-zariski-after-blowup}。令
$\mathcal{I}_i \subset \mathcal{O}_{X_i}$ 为相应的有限型拟相干理想层。
回顾 $X_{12}^i \to X_{12}$ 是以
$p_i^{-1}\mathcal{I}_i \mathcal{O}_{X_{12}}$ 为中心的爆破；
参见《因子》，引理 \ref{divisors-lemma-strict-transform}。
令 $X_{12}'$ 为 $X_{12}$ 以
$p_1^{-1}\mathcal{I}_1 p_2^{-1}\mathcal{I}_2 \mathcal{O}_{X_{12}}$
为中心的爆破；其含义参见《因子》，引理
\ref{divisors-lemma-blowing-up-two-ideals}。特别地，我们得到交换图
$$
\xymatrix{
X_{12}' \ar[d] \ar[r] & X_{12}^2 \ar[d] \\
X_{12}^1 \ar[r] & X_{12}
}
$$
其中所有态射均为 $U$-容许爆破。由于
$X_{12}^i \subset X_i^i$ 是开子概形，由《因子》，引理
\ref{divisors-lemma-extend-admissible-blowups}，可以选取
$U$-容许爆破 $X_i' \to X_i^i$，其限制为
$X_{12}' \to X_{12}^i$。于是 $X_{12}' \subset X_i'$
是开子概形，且交换图
$$
\xymatrix{
X_{12}' \ar[d] \ar[r] & X_i' \ar[d] \\
X_{12}^i \ar[r] & X_i^i
}
$$
的竖直箭头是爆破，水平箭头是开浸入。注意，
$X'_{12} \to X_1' \times_S X_2'$ 是浸入且适当
（使用 $X'_{12} \to X_{12}$ 为适当态射、
$X_{12} \to X_1 \times_S X_2$ 为闭态射，并且
$X_1' \times_S X_2' \to X_1 \times_S X_2$ 为分离态射，
再应用 《态射》，引理
\ref{morphisms-lemma-image-proper-scheme-closed}）。
因此 $X'_{12} \to  X_1' \times_S X_2'$ 是闭浸入。
沿公共开子概形 $X_{12}'$ 黏合 $X_1'$ 和 $X_2'$ 来定义 $X$，
则 $X \to S$ 为有限型且分离（引理 \ref{flat-lemma-check-separated}）。
由于 $U$-容许爆破的复合仍是 $U$-容许爆破
（《因子》，引理
\ref{divisors-lemma-composition-admissible-blowups}），引理得证。
\end{proof}

\begin{lemma}
\label{flat-lemma-replaced-by-strict-transform}
设 $X \to S$ 和 $Y \to S$ 是概形态射。设 $U \subset X$
是开子概形。设 $V \to X \times_S Y$ 是拟紧态射，且它与第一投影
的复合映入 $U$。设 $Z \subset X \times_S Y$ 是
$V \to X \times_S Y$ 的概形论像，并设
$X' \to X$ 是 $U$-容许爆破。则 $V \to X' \times_S Y$
的概形论像是 $Z$ 关于该爆破的严格变换。
\end{lemma}

\begin{proof}
以 $Z' \to Z$ 表示严格变换。态射 $Z' \to X'$ 诱导态射
$Z' \to X' \times_S Y$，后者是闭浸入
（因为依定义，$Z'$ 是 $X' \times_X Z$ 的闭子概形）。
因此只需证明 $V \to Z'$ 的概形论像 $Z''$ 就是 $Z'$。
注意，$Z'' \subset Z'$ 是闭子概形，且 $V \to Z'$
经 $Z''$ 分解。态射 $V \to X \times_S Y$ 与
$V \to X' \times_S Y$ 均拟紧；对后者，这由《概形》，引理
\ref{schemes-lemma-quasi-compact-permanence}
以及 $X' \times_S Y \to X \times_S Y$ 作为适当态射的基变换为分离态射而得。
故由《态射》，引理
\ref{morphisms-lemma-quasi-compact-scheme-theoretic-image}
可知 $Z \cap (U \times_S Y) = Z'' \cap (U \times_S Y)$。
因此，浸入态射 $Z'' \to Z'$ 在 $Z' \to Z$ 的例外除子
$E$ 之外是同构。然而，$Z'$ 的结构层没有支集在 $E$ 上的非零截面
（这是严格变换的定义），故满射
$\mathcal{O}_{Z'} \to \mathcal{O}_{Z''}$ 必为同构。
\end{proof}

\begin{lemma}
\label{flat-lemma-compactification-dominates}
设 $S$ 是拟紧且拟分离概形。设 $U$ 是在 $S$ 上有限型且分离的概形，
$V \subset U$ 是拟紧开子集。若 $V$ 在 $S$ 上有紧化
$V \subset Y$，则存在一个 $V$-容许爆破 $Y' \to Y$
和开子集 $V \subset V' \subset Y'$，使 $V \to U$
延拓为适当态射 $V' \to U$。
\end{lemma}

\begin{proof}
考虑“对角”态射 $V \to Y \times_S U$ 的概形论像
$Z \subset Y \times_S U$。若用一个 $V$-容许爆破替换 $Y$，
则 $Z$ 被替换为它关于该爆破的严格变换；参见引理
\ref{flat-lemma-replaced-by-strict-transform}。因此由引理
\ref{flat-lemma-zariski-after-blowup}，可以假设 $Z \to Y$ 是开浸入。
若 $V' \subset Y$ 表示其像，则诱导态射 $V' \to U$ 是适当态射：
投影 $Y \times_S U \to U$ 是适当态射，而
$V' \cong Z$ 是 $Y \times_S U$ 的闭子概形。
\end{proof}

\noindent
下述引理仅在 Noether 情形下表述。拟紧且拟分离概形的版本也成立，
但它将直接由本节的主定理推出。

\begin{lemma}
\label{flat-lemma-two-compactifications}
设 $S$ 是 Noether 概形，$U$ 是在 $S$ 上有限型且分离的概形。
设 $U = U_1 \cup U_2$ 是开覆盖，$U_1$ 与 $U_2$ 在 $S$ 上均有紧化，
并且 $U_1 \cap U_2$ 在 $U$ 中稠密。则 $U$ 在 $S$ 上有紧化。
\end{lemma}

\begin{proof}
对 $i = 1, 2$，选取紧化 $U_i \subset X_i$。可以假设
$U_i$ 在 $X_i$ 中概形论稠密。由引理
\ref{flat-lemma-compactification-dominates}，还可以假设存在开子集
$V_i \subset X_i$ 和适当态射 $\psi_i : V_i \to U$，
它延拓 $\text{id} : U_i \to U_i$。图示为
$$
\xymatrix{
U_i \ar[r] \ar[d] & V_i \ar[r] \ar[dl]^{\psi_i} & X_i \\
U
}
$$
若 $\{i, j\} = \{1, 2\}$，记
$Z_i = U \setminus U_j = U_i \setminus (U_1 \cap U_2)$
以及
$Z_j = U \setminus U_i = U_j \setminus (U_1 \cap U_2)$.
于是
$$
U = U_1 \amalg Z_2 = Z_1 \amalg U_2 = Z_1 \amalg (U_1 \cap U_2) \amalg Z_2
$$
记 $Z_{i, i} \subset V_i$ 为 $Z_i$ 在 $\psi_i$ 下的逆像。
注意，$\psi_i$ 在 $Z_i$ 的某个开邻域上是同构。
记 $Z_{i, j} \subset V_i$ 为 $Z_j$ 在 $\psi_i$ 下的逆像。
注意，$\psi_i : Z_{i, j} \to Z_j$ 是适当态射。
由于 $Z_i$ 与 $Z_j$ 是 $U$ 中互不相交的闭子集，
$Z_{i, i}$ 与 $Z_{i, j}$ 是 $V_i$ 中互不相交的闭子集。

\medskip\noindent
以 $\overline{Z}_{i, i}$ 和 $\overline{Z}_{i, j}$ 表示
$Z_{i, i}$ 和 $Z_{i, j}$ 在 $X_i$ 中的闭包。
用一个 $V_i$-容许爆破替换 $X_i$ 后，可以假设
$\overline{Z}_{i, i}$ 与 $\overline{Z}_{i, j}$ 不相交；
参见引理 \ref{flat-lemma-separate-disjoint-locally-closed-by-blowing-up}。
我们假设这对 $X_1$ 和 $X_2$ 都成立。注意，用进一步的
$V_i$-容许爆破替换 $X_i$ 时，此性质仍得保持。

\medskip\noindent
令 $V_{12} = V_1 \times_U V_2$。有一个浸入
$V_{12} \to X_1 \times_S X_2$，它是闭浸入
$V_{12} = V_1 \times_U V_2 \to V_1 \times_S V_2$
（《概形》，引理 \ref{schemes-lemma-fibre-product-after-map}）
与开浸入 $V_1 \times_S V_2 \to X_1 \times_S X_2$ 的复合。
令 $X_{12} \subset X_1 \times_S X_2$ 为
$V_{12} \to X_1 \times_S X_2$ 的概形论像。投影态射
$$
p_1 : X_{12} \to X_1
\quad\text{且}\quad
p_2 : X_{12} \to X_2
$$
是适当态射，因为 $X_1$ 和 $X_2$ 在 $S$ 上适当。
若用一个 $V_1$-容许爆破替换 $X_1$，则 $X_{12}$
被替换为它关于该爆破的严格变换；参见引理
\ref{flat-lemma-replaced-by-strict-transform}。

\medskip\noindent
以 $\psi : V_{12} \to U$ 表示复合
$\psi = \psi_1 \circ p_1|_{V_{12}} = \psi_2 \circ p_2|_{V_{12}}$.
考虑闭子概形
$$
Z_{12, 2} =
(p_1|_{V_{12}})^{-1}(Z_{1, 2}) =
(p_2|_{V_{12}})^{-1}(Z_{2, 2}) =
\psi^{-1}(Z_2) \subset V_{12}
$$
由于 $\psi_2 : V_2 \to U$ 在 $Z_2$ 的某个开邻域上是同构，
且 $V_{12} = V_1 \times_U V_2$，态射
$p_1|_{V_{12}} : V_{12} \to V_1$ 在 $Z_{1, 2}$
的某个开邻域上是同构。由引理 \ref{flat-lemma-blowup-iso-along}，
存在一个 $V_1$-容许爆破 $X_1' \to X_1$，使 $p_1$
的严格变换 $p'_1 : X'_{12} \to X'_1$ 在
$Z_{1, 2}$ 于 $X'_1$ 中的闭包的某个开邻域上是同构。
用 $X'_1$ 替换 $X_1$，并用 $X'_{12}$ 替换 $X_{12}$ 后，
可以假设 $p_1$ 在 $\overline{Z}_{1, 2}$ 的某个开邻域上是同构。

\medskip\noindent
上一段的化简告诉我们
$$
X_{12} \cap (\overline{Z}_{1, 2} \times_S \overline{Z}_{2, 1}) = \emptyset
$$
% P05-EN-CH38-0065: frozen source omits “is” in “where the intersection taken”; candidate “where the intersection is taken”.
，其中交集取于 $X_1 \times_S X_2$ 中。事实上，
$X_{12}$ 中的逆像 $p_1^{-1}(\overline{Z}_{1, 2})$
同构地映到 $\overline{Z}_{1, 2}$。特别地，$Z_{12, 2}$
在 $p_1^{-1}(\overline{Z}_{1, 2})$ 中稠密。因此 $p_2$
将 $p_1^{-1}(\overline{Z}_{1, 2})$ 映入
$\overline{Z}_{2, 2}$。由于
$\overline{Z}_{2, 2} \cap \overline{Z}_{2, 1} = \emptyset$，结论成立。

\medskip\noindent
考虑由黏合得到的概形
$$
W_i = U \coprod\nolimits_{U_i} (X_i \setminus \overline{Z}_{i, j}),
\quad i = 1, 2
$$
应用引理 \ref{flat-lemma-check-separated} 可见 $W_i \to S$ 是分离态射。
首先，$U \to S$ 与 $X_i \to S$ 均为分离态射。浸入
$U_i \to U \times_S (X_i \setminus \overline{Z}_{i, j})$
是闭浸入：任何满足 $u_i \in U_i$ 且
$u \in U \setminus U_i$ 的特化 $u_i \leadsto u$，都能沿适当态射
$\psi_i : V_i \to U$ 唯一提升为 $V_i$ 中的特化
$u_i \leadsto v_i$，而 $v_i$ 必属于 $Z_{i, j}$。
因此该浸入的像为闭集，从而该浸入是闭浸入。

\medskip\noindent
另一方面，设 $A$ 是 $S$ 上的任意赋值环，分式域为 $K$，
并设 $\gamma : \Spec(K) \to (U_1 \cap U_2)$ 是 $S$ 上态射。
则存在某个 $i$，使 $\gamma$ 可延拓为态射
$h_i : \Spec(A) \to W_i$。事实上，对 $i = 1, 2$，
由 $X_i$ 在 $S$ 上适当的赋值判别法，存在延拓 $\gamma$ 的态射
$g_i : \Spec(A) \to X_i$；参见《态射》，引理
\ref{morphisms-lemma-characterize-proper}。
因此，唯一可能出现的问题是对 $i = 1, 2$ 均有
$g_i(\mathfrak m_A) \in \overline{Z}_{i, j}$。
% P05-EN-CH38-0066: frozen source misspells “emptiness” as “emptyness”.
但这与上面已证明的 $X_{12}$ 和
$\overline{Z}_{1, 2} \times_S \overline{Z}_{2, 1}$ 交为空相矛盾。

\medskip\noindent
考虑如引理 \ref{flat-lemma-find-common-blowups} 中的图
$$
\xymatrix{
W_1' \ar[d] \ar[r] & W & W_2' \ar[l] \ar[d] \\
W_1 & U \ar[l] \ar[lu] \ar[u] \ar[ru] \ar[r] & W_2
}
$$
由上一段，对每个实线图
$$
\xymatrix{
\Spec(K) \ar[r]_\gamma  \ar[d] & W \ar[d] \\
\Spec(A) \ar@{..>}[ru] \ar[r] & S
}
$$
若 $\Im(\gamma) \subset U_1 \cap U_2$，则存在某个 $i$
以及 $\gamma$ 的延拓 $h_i : \Spec(A) \to W_i$。
利用 $W'_i \to W_i$ 为适当态射的赋值判别法，可以进一步将
$h_i$ 提升为 $h'_i : \Spec(A) \to W'_i$。故图中的虚线箭头存在。
由于 $W$ 在 $S$ 上分离，该箭头也唯一。由《态射》，引理
\ref{morphisms-lemma-refined-valuative-criterion-universally-closed}，
这说明 $W \to S$ 是普遍闭态射。又因 $W \to S$
已经有限型且分离，故它是适当态射，结论得证。
\end{proof}

\begin{theorem}
\label{flat-theorem-nagata}
\begin{reference}
参见 \cite{Lutkebohmert}、\cite{Conrad-Nagata}、\cite{Nagata-1}、
\cite{Nagata-2}、\cite{Nagata-3} 以及 \cite{Nagata-4}
\end{reference}
设 $S$ 是拟紧且拟分离概形，$X \to S$ 是分离的有限型态射。
则 $X$ 在 $S$ 上有紧化。
\end{theorem}

\begin{proof}
首先化归到 Noether 情形。我们强烈建议读者跳过本段。
存在闭浸入 $X \to X'$，其中 $X' \to S$ 分离且为有限表示；
参见《极限》，命题
\ref{limits-proposition-separated-closed-in-finite-presentation}。
若能找到 $X'$ 在 $S$ 上的一个紧化，则取 $X$ 在其中的概形论像，
便得到 $X$ 在 $S$ 上的紧化。因此可以假设
$X \to S$ 分离且为有限表示。可以将
$S = \lim S_i$ 写成一族 Noether 概形的有向极限，
其转移态射均为仿射态射；参见《极限》，命题
\ref{limits-proposition-approximate}。可以选取某个 $i$
和有限表示态射 $X_i \to S_i$，使它到 $S$ 的基变换为
$X \to S$；参见《极限》，引理
\ref{limits-lemma-descend-finite-presentation}。增大 $i$ 后，
可以假设 $X_i \to S_i$ 分离；参见《极限》，引理
\ref{limits-lemma-descend-separated-finite-presentation}。
若能找到 $X_i$ 在 $S_i$ 上的紧化，则它到 $S$ 的基变换
就是 $X$ 在 $S$ 上的紧化。这就化归到下一段讨论的情形。

\medskip\noindent
现设 $S$ 为 Noether 概形。可以选取有限仿射开覆盖
$X = \bigcup_{i = 1, \ldots, n} U_i$，使
$U_1 \cap \ldots \cap U_n$ 在 $X$ 中稠密。这由
《性质》，引理 \ref{properties-lemma-point-and-maximal-points-affine}
以及 $X$ 拟紧且只有有限多个不可约分支这一事实推出。
对每个 $i$，由《态射》，引理
\ref{morphisms-lemma-quasi-affine-finite-type-over-S}，
可以选取 $n_i \geq 0$ 和浸入 $U_i \to \mathbf{A}^{n_i}_S$。
在 $\mathbf{P}^{n_i}_S$ 中取概形论像，可知对
$i = 1, \ldots, n$，$U_i$ 在 $S$ 上有紧化。
应用引理 \ref{flat-lemma-two-compactifications} 共
$(n - 1)$ 次，即得定理。
\end{proof}






\section{h 拓扑}
\label{flat-section-h}

\noindent
粗略地说，在本文中，h 层是 Zariski 拓扑下的层，并且对满射的
有限表示适当态射满足层条件；参见引理
\ref{flat-lemma-characterize-sheaf-h}。不过值得指出的是，
概形范畴上 h 拓扑的定义随参考文献而异。

\medskip\noindent
Voevodsky 最初将 h 覆盖定义为有限族有限型态射，并要求它们联合地
普遍为商态射（《态射》，定义
\ref{morphisms-definition-submersive}）；参见
\cite[定义 3.1.2]{Voevodsky}。当所用概形范畴限制为某个固定
Noether 基概形上的全体有限型概形时，这一定义最为适用。
在此情形中，Voevodsky 将 h 覆盖与 ph 覆盖联系起来。
ph 拓扑由 Zariski 覆盖和适当满射生成。更多内容参见
《拓扑》，第 \ref{topologies-section-ph} 节。

\medskip\noindent
我们在《拓扑》，第 \ref{topologies-section-V} 节中研究 V 拓扑。
若 $Y$ 中点的任意特化都是 $X$ 中点的某个特化之像，
且任意基变换后仍如此，则拟紧态射 $X \to Y$ 定义一个 V 覆盖
（《拓扑》，引理 \ref{topologies-lemma-refine-qcqs-V}）。
此时 $X \to Y$ 是普遍商态射
（《拓扑》，引理
\ref{topologies-lemma-V-covering-universally-submersive}）。
处理非 Noether 概形时，V 覆盖这一概念可以很好地代替
“具有固定靶且联合地普遍为商态射的态射族”。

\medskip\noindent
我们的方法是先对局部有限表示的态射族（允许无限族）证明
ph 覆盖与 V 覆盖等价。随后在 Stacks Project 中以这些态射族作为
h 覆盖的定义。对 Noether 概形和有限族而言，这些覆盖与
Voevodsky 的定义一致；参见引理 \ref{flat-lemma-Noetherian-h-covering}。
在固定拟紧且拟分离概形 $S$ 上有限表示概形的范畴中，
这些覆盖所确定的拓扑与 \cite[定义 2.7]{Witt-Grass}
中的拓扑相同。

\begin{lemma}
\label{flat-lemma-equivalence-h-v-locally-finite-presentation}
设 $\{f_i : X_i \to X\}_{i \in I}$ 是具有固定靶的概形态射族，
并且对每个 $i$，$f_i$ 都局部为有限表示。下列条件等价：
\begin{enumerate}
\item $\{X_i \to X\}$ 是 ph 覆盖；
\item $\{X_i \to X\}$ 是 V 覆盖。
\end{enumerate}
\end{lemma}

\begin{proof}
设 $U \subset X$ 为仿射开子集。由《拓扑》，定义
\ref{topologies-definition-ph-covering} 和
\ref{topologies-definition-V-covering}，只需证明基变换
$\{X_i \times_X U \to U\}$ 可由标准 ph 覆盖加细，当且仅当
它可由标准 V 覆盖加细。因此可以假设 $X$ 仿射，而需证明
$\{X_i \to X\}$ 可由标准 ph 覆盖加细，当且仅当它可由
标准 V 覆盖加细。标准 ph 覆盖是标准 V 覆盖；参见
《拓扑》，引理 \ref{topologies-lemma-standard-ph-standard-V}。
故只需证明反向蕴含。

\medskip\noindent
假设 $X$ 仿射，且 $\{f_i : X_i \to X\}_{i \in I}$
可由标准 V 覆盖
$\{g_j : Y_j \to X\}_{j = 1, \ldots, m}$ 加细。
对每个 $j$，选取 $i_j$ 和态射 $h_j : Y_j \to X_{i_j}$，
使 $g_j = f_{i_j} \circ h_j$。由于 $Y_j$ 仿射，因而拟紧，
对每个 $j$ 可以找到有限多个仿射开子集
$U_{j, k} \subset X_{i_j}$，使
$\Im(h_j) \subset \bigcup U_{j, k}$。于是
$\{U_{j, k} \to X\}_{j, k}$ 加细 $\{X_i \to X\}$，
并且是标准 V 覆盖（它是仿射概形之间态射的有限族，
而且从 $\{Y_j \to X\}$ 的相应性质继承赋值环提升性质）。
这就化归到下一段讨论的情形。

\medskip\noindent
假设 $\{f_i : X_i \to X\}_{i = 1, \ldots, n}$
是标准 V 覆盖，且 $f_i$ 为有限表示。需证明
$\{X_i \to X\}$ 可由标准 ph 覆盖加细。选取一般平坦性分层
$$
X = S \supset S_0 \supset S_1 \supset \ldots \supset S_t = \emptyset
$$
它是 《态射进阶》，引理
\ref{more-morphisms-lemma-generic-flatness-stratification-scheme}，
对仿射概形之间的有限表示态射
$$
\coprod\nolimits_{i = 1, \ldots, n} f_i :
\coprod\nolimits_{i = 1, \ldots, n} X_i
\longrightarrow
X
$$
所给出的分层。下文将不再说明地使用该分层的全部性质。
依构造，每个 $f_i$ 到 $U_k = S_k \setminus S_{k + 1}$
的基变换都是平坦态射。以 $Y_k$ 表示 $U_k$ 在 $S_k$
中的概形论闭包。由于 $U_k \to S_k$ 是拟紧开浸入
（参见《性质》，引理
\ref{properties-lemma-quasi-coherent-finite-type-ideals}），
$U_k \subset Y_k$ 是拟紧稠密（且概形论稠密）的开浸入；
参见《态射》，引理
\ref{morphisms-lemma-quasi-compact-scheme-theoretic-image}。
态射 $\coprod_{k = 0, \ldots, t - 1} Y_k \to X$ 有限且满，
故 $\{Y_k \to X\}$ 是标准 ph 覆盖，因而也是标准 V 覆盖
（见上文）。由标准 V 覆盖的传递性质
（《拓扑》，引理 \ref{topologies-lemma-composition-standard-V}），
只需证明覆盖 $\{X_i \to X\}$ 到每个 $Y_k$ 的拉回
可由标准 V 覆盖加细。这就化归到下一段的情形。

\medskip\noindent
假设 $\{f_i : X_i \to X\}_{i = 1, \ldots, n}$ 是标准 V 覆盖，
$f_i$ 为有限表示，并且存在稠密拟紧开子集 $U \subset X$，
使 $X_i \times_X U \to U$ 平坦。由定理
\ref{flat-theorem-flatten-module}，存在 $U$-容许爆破
$X' \to X$，使 $f_i$ 的严格变换
$f'_i : X'_i \to X'$ 平坦。注意，投射（因而闭）态射
$X' \to X$ 是满射，因为 $U \subset X$ 稠密且 $U$
被视为 $X'$ 的开子集。必要时用进一步的 $U$-容许爆破替换
$X'$，还可假设 $U \subset X'$ 概形论稠密
（参见注 \ref{flat-remark-successive-blowups}）。
因此，对每个点 $x \in X'$，存在赋值环 $V$ 和态射
$g : \Spec(V) \to X'$，使 $\Spec(V)$ 的一般点映入 $U$，
而 $\Spec(V)$ 的闭点映到 $x$；参见《态射》，引理
\ref{morphisms-lemma-reach-points-scheme-theoretic-image}。
由于 $\{X_i \to X\}$ 是标准 V 覆盖，可以选取赋值环扩张
$V \subset W$、指标 $i$ 和态射 $\Spec(W) \to X_i$，使图
$$
\xymatrix{
\Spec(W) \ar[d] \ar[rr] & & X_i \ar[d] \\
\Spec(V) \ar[r] & X' \ar[r] & X
}
$$
交换。由于 $X'_i \subset X' \times_X X_i$ 是包含开子概形
$U \times_X X_i$ 的闭子概形，$\Spec(W)$ 是整概形，
而诱导态射 $h : \Spec(W) \to X' \times_X X_i$
将 $\Spec(W)$ 的一般点映入 $U \times_X X_i$，可知 $h$
经闭子概形 $X'_i \subset X' \times_X X_i$ 分解。
因此 $\{f'_i : X'_i \to X'\}$ 是 V 覆盖。特别地，
$\coprod f'_i$ 是满射，故 $\{X'_i \to X'\}$ 是 fppf 覆盖。
由于 fppf 覆盖是 ph 覆盖
（《态射进阶》，引理 \ref{more-morphisms-lemma-fppf-ph}），
可以找到加细 $\{X'_i \to X\}$ 的标准 ph 覆盖
$\{Y_j \to X'\}$。设此覆盖由适当满射 $Y \to X'$
和有限仿射开覆盖 $Y = \bigcup Y_j$ 给出。于是复合
$Y \to X$ 是适当满射，从而 $\{Y_j \to X\}$ 是标准 ph 覆盖。
证明完毕。
\end{proof}

\noindent
现给出我们的定义。

\begin{definition}
\label{flat-definition-h-covering}
设 $T$ 是概形。$T$ 的一个{\it h 覆盖}是态射族
$\{f_i : T_i \to T\}_{i \in I}$，其中每个 $f_i$
都局部为有限表示，并满足引理
\ref{flat-lemma-equivalence-h-v-locally-finite-presentation}
中的一个（从而全部）等价条件。
\end{definition}

\noindent
对 Noether 概形，这与 ph 覆盖相同
（我们将在下述引理 \ref{flat-lemma-Noetherian-h-ph} 中记录此事），
因而恢复了 Voevodsky 的概念。

\begin{lemma}
\label{flat-lemma-Noetherian-h-covering}
设 $X$ 是 Noether 概形，并设 $\{X_i \to X\}_{i \in I}$
是有限族有限型态射。下列条件等价：
\begin{enumerate}
\item $\coprod_{i \in I} X_i \to X$ 是普遍商态射
（《态射》，定义 \ref{morphisms-definition-submersive}）；
\item $\{X_i \to X\}_{i \in I}$ 是 h 覆盖。
\end{enumerate}
\end{lemma}

\begin{proof}
由更一般的 《拓扑》，引理
\ref{topologies-lemma-V-covering-universally-submersive}
和 h 覆盖的定义，(2) $\Rightarrow$ (1) 立即成立。
假设 $\coprod X_i \to X$ 是普遍商态射。我们将证明
$\{X_i \to X\}$ 可由 ph 覆盖加细；由《拓扑》，引理
\ref{topologies-lemma-refine-by-ph} 和 h 覆盖的定义，这已经足够。
论证与引理 \ref{flat-lemma-equivalence-h-v-locally-finite-presentation}
的证明中所用者相同。

\medskip\noindent
选取一般平坦性分层
$$
X = S \supset S_0 \supset S_1 \supset \ldots \supset S_t = \emptyset
$$
它是 《态射进阶》，引理
\ref{more-morphisms-lemma-generic-flatness-stratification-scheme}
中对有限表示态射
$$
\coprod\nolimits_{i = 1, \ldots, n} f_i :
\coprod\nolimits_{i = 1, \ldots, n} X_i
\longrightarrow
X
$$
所给出的分层。下文将不再说明地使用该分层的全部性质。
依构造，每个 $f_i$ 到 $U_k = S_k \setminus S_{k + 1}$
的基变换都是平坦态射。以 $Y_k$ 表示 $U_k$ 在 $S_k$
中的概形论闭包。由于 $U_k \to S_k$ 是拟紧开浸入
（本段所有概形均为 Noether 概形），$U_k \subset Y_k$
是拟紧稠密（且概形论稠密）的开浸入；参见《态射》，引理
\ref{morphisms-lemma-quasi-compact-scheme-theoretic-image}。
态射 $\coprod_{k = 0, \ldots, t - 1} Y_k \to X$ 有限且满，
故 $\{Y_k \to X\}$ 是 ph 覆盖。由 ph 覆盖的传递性质
（《拓扑》，引理 \ref{topologies-lemma-ph}），
只需证明覆盖 $\{X_i \to X\}$ 到每个 $Y_k$ 的拉回
可由 ph 覆盖加细。这就化归到下一段的情形。

\medskip\noindent
假设 $\coprod X_i \to X$ 是普遍商态射，并且存在稠密开子集
$U \subset X$，使对每个 $i$，$X_i \times_X U \to U$
均为平坦态射。由定理 \ref{flat-theorem-flatten-module}，
存在 $U$-容许爆破 $X' \to X$，使对每个 $i$，
$f_i$ 的严格变换 $f'_i : X'_i \to X'$ 均为平坦态射。
注意，投射（因而闭）态射 $X' \to X$ 是满射，因为
$U \subset X$ 稠密且 $U$ 被视为 $X'$ 的开子集。必要时用进一步的
$U$-容许爆破替换 $X'$，还可假设 $U \subset X'$ 稠密
（参见注 \ref{flat-remark-successive-blowups}）。
因此，对每个点 $x \in X'$，存在离散赋值环 $A$ 和态射
$g : \Spec(A) \to X'$，使 $\Spec(A)$ 的一般点映入 $U$，
而 $\Spec(A)$ 的闭点映到 $x$；参见《极限》，引理
\ref{limits-lemma-reach-point-closure-Noetherian}。令
$$
W = \Spec(A) \times_X \coprod X_i = \coprod \Spec(A) \times_X X_i
$$
由于 $\coprod X_i \to X$ 是普遍商态射，$W$ 中存在特化
$w' \leadsto w$，使 $w'$ 映到 $\Spec(A)$ 的一般点，
而 $w$ 映到 $\Spec(A)$ 的闭点。
（否则，$W \to \Spec(A)$ 的闭纤维在一般化下稳定，因而为开集，
这与 $W \to \Spec(A)$ 是商态射相矛盾。）
设 $w' \in \Spec(A) \times_X X_i$，则当然也有
$w \in \Spec(A) \times_X X_i$。令
$x'_i \leadsto x_i$ 为 $w' \leadsto w$ 在
$X' \times_X X_i$ 中的像。由于 $x'_i \in X'_i$，且
$X'_i \subset X' \times_X X_i$ 是闭子概形，可知
$x_i \in X'_i$。又因为 $x_i$ 映到 $x \in X'$，所以
$\coprod X'_i \to X'$ 是满射！特别地，
$\{X'_i \to X'\}$ 是 fppf 覆盖，而 fppf 覆盖是 ph 覆盖
（《态射进阶》，引理 \ref{more-morphisms-lemma-fppf-ph}）。
由于 $X' \to X$ 是适当满射，$\{X'_i \to X\}$ 是 ph 覆盖，
证明完毕。
\end{proof}

\begin{lemma}
\label{flat-lemma-Noetherian-h-ph}
设 $X$ 是局部 Noether 概形。以 $X$ 为靶的态射族
$\{f_i : X_i \to X\}_{i \in I}$ 是 h 覆盖，当且仅当它是 ph 覆盖。
\end{lemma}

\begin{proof}
由定义 \ref{flat-definition-h-covering}，h 覆盖是 ph 覆盖。
反之，若 $\{f_i : X_i \to X\}$ 是 ph 覆盖，则态射 $f_i$
局部为有限型（《拓扑》，定义
\ref{topologies-definition-ph-covering}）。由于 $X$ 局部 Noether，
每个 $f_i$ 都局部为有限表示，故依定义得到 h 覆盖。
\end{proof}

\noindent
% P05-EN-CH38-0067: frozen source has compound subject “lemma and ... theorem” with singular “shows”; candidate “show”.
下述引理与 \cite[定理 8.4]{rydh_descent} 表明，
我们的定义与 David Rydh 的论文 \cite{rydh_descent} 中的定义一致
（或至少密切相关）。为简便起见，我们限制在仿射基情形。

\begin{lemma}
\label{flat-lemma-approximate-h-cover}
设 $X$ 是仿射概形，$\{X_i \to X\}_{i \in I}$ 是 h 覆盖。
则存在满的适当态射
$$
Y \longrightarrow X
$$
它还是有限表示的（！），并且存在有限仿射开覆盖
$Y = \bigcup_{j = 1, \ldots, m} Y_j$，使
$\{Y_j \to X\}_{j = 1, \ldots, m}$ 加细
$\{X_i \to X\}_{i \in I}$。
\end{lemma}

\begin{proof}
依假设，存在适当满射 $Y \to X$ 和有限仿射开覆盖
$Y = \bigcup_{j = 1, \ldots, m} Y_j$，使
$\{Y_j \to X\}_{j = 1, \ldots, m}$ 加细
$\{X_i \to X\}_{i \in I}$。这意味着对每个 $j$，
存在指标 $i_j \in I$ 及 $X$ 上态射
$h_j : Y_j \to X_{i_j}$。参见定义 \ref{flat-definition-h-covering}
和 《拓扑》，定义 \ref{topologies-definition-ph-covering}。
问题在于尚不知道 $Y \to X$ 是否为有限表示。
由《极限》，引理
\ref{limits-lemma-proper-limit-of-proper-finite-presentation}，可将
$$
Y = \lim Y_\lambda
$$
写成 $X$ 上概形 $Y_\lambda$ 的有向极限，其中这些概形
适当且为有限表示，态射 $Y \to Y_\lambda$ 及转移态射均为闭浸入。
% P05-EN-CH38-0068: frozen source duplicates “the” before “transition morphisms”.
注意，每个 $Y_\lambda \to X$ 都是满射。由《极限》，引理
\ref{limits-lemma-descend-opens}，可以找到某个 $\lambda$
和覆盖 $Y_\lambda$ 的拟紧开子集
$Y_{\lambda, j} \subset Y_\lambda$，$j = 1, \ldots, m$，
它们在 $Y$ 上限制为 $Y_j$。于是
$Y_j = \lim Y_{\lambda, j}$。增大 $\lambda$ 后，可假设对所有
$j$，$Y_{\lambda, j}$ 均仿射；参见《极限》，引理
\ref{limits-lemma-limit-affine}。最后，由于 $X_i \to X$
局部为有限表示，可以使用局部有限表示态射的函子性刻画
（《极限》，命题
\ref{limits-proposition-characterize-locally-finite-presentation}）
找到某个 $\lambda$，使对每个 $j$ 都存在态射
$h_{\lambda, j} : Y_{\lambda, j} \to X_{i_j}$，
它在 $Y_j$ 上的限制就是上面选取的态射 $h_j$。
因此 $\{Y_{\lambda, j} \to X\}$ 加细
$\{X_i \to X\}$，证明完毕。
\end{proof}

\noindent
现在回到 h 覆盖一般理论的建立。

\begin{lemma}
\label{flat-lemma-zariski-h}
fppf 覆盖是 h 覆盖。因此 syntomic、光滑、\'etale 以及
Zariski 覆盖也都是 h 覆盖。
\end{lemma}

\begin{proof}
这是因为 fppf 覆盖中的态射按要求局部为有限表示，
且 fppf 覆盖是 ph 覆盖；参见《态射进阶》，引理
\ref{more-morphisms-lemma-fppf-ph}。第二个断言由第一个断言和
《拓扑》，引理
\ref{topologies-lemma-zariski-etale-smooth-syntomic-fppf} 推出。
\end{proof}

\begin{lemma}
\label{flat-lemma-surjective-proper-finite-presentation-h}
设 $f : Y \to X$ 是有限表示的适当满概形态射。
则 $\{Y \to X\}$ 是 h 覆盖。
\end{lemma}

\begin{proof}
合并应用 《拓扑》，引理
\ref{topologies-lemma-zariski-etale-smooth-syntomic-fppf-fpqc-ph-V} 和
\ref{topologies-lemma-surjective-proper-ph}.
\end{proof}

\begin{lemma}
\label{flat-lemma-refine-by-h}
设 $T$ 是概形，并设 $\{f_i : T_i \to T\}_{i \in I}$ 是态射族，
使对所有 $i$，$f_i$ 均局部为有限表示。下列条件等价：
\begin{enumerate}
\item $\{T_i \to T\}_{i \in I}$ 是 h 覆盖；
\item 存在加细 $\{T_i \to T\}_{i \in I}$ 的 h 覆盖；
\item $\{\coprod_{i \in I} T_i \to T\}$ 是 h 覆盖。
\end{enumerate}
\end{lemma}

\begin{proof}
这由 ph 覆盖的相应断言
（《拓扑》，引理 \ref{topologies-lemma-refine-by-ph}）
或 V 覆盖的相应断言
（《拓扑》，引理 \ref{topologies-lemma-refine-by-V}）推出。
\end{proof}

\noindent
下面证明我们的 h 覆盖概念满足 《站点》，定义
\ref{sites-definition-site} 的条件。

\begin{lemma}
\label{flat-lemma-h}
设 $T$ 是概形。
\begin{enumerate}
\item 若 $T' \to T$ 是同构，则 $\{T' \to T\}$ 是 $T$ 的 h 覆盖。
\item 若 $\{T_i \to T\}_{i\in I}$ 是 h 覆盖，且对每个 $i$
都有 h 覆盖 $\{T_{ij} \to T_i\}_{j\in J_i}$，则
$\{T_{ij} \to T\}_{i \in I, j\in J_i}$ 是 h 覆盖。
\item 若 $\{T_i \to T\}_{i\in I}$ 是 h 覆盖，而
$T' \to T$ 是概形态射，则
$\{T' \times_T T_i \to T'\}_{i\in I}$ 是 h 覆盖。
\end{enumerate}
\end{lemma}

\begin{proof}
这立即由 ph 覆盖或 V 覆盖的相应断言
（《拓扑》，引理 \ref{topologies-lemma-ph} 或
\ref{topologies-lemma-V}），以及局部有限表示态射类
在基变换和复合下保持这一事实推出。
\end{proof}

\noindent
下面定义我们将在 Stacks Project 中使用的大 h 站点。
继续之前，宜先阅读 《拓扑》，第 \ref{topologies-section-procedure} 节中的一般讨论。

\begin{definition}
\label{flat-definition-big-h-site}
一个{\it 大 h 站点}是按如下方式构造的、《站点》，定义
\ref{sites-definition-site} 意义下的任意站点 $\Sch_h$：
\begin{enumerate}
\item 任取一组概形 $S_0$，并任取这些概形之间的一组 h 覆盖
$\text{Cov}_0$。
\item 以任意范畴 $\Sch_\alpha$ 作为底层范畴；该范畴从集合
$S_0$ 出发，按《集合》，引理
\ref{sets-lemma-construct-category} 构造。
\item 从范畴 $\Sch_\alpha$、h 覆盖类和上面选取的集合
$\text{Cov}_0$ 出发，按《集合》，引理
\ref{sets-lemma-coverings-site} 任取一组覆盖。
\end{enumerate}
\end{definition}

\noindent
关于大站点定义的动机与说明，参见《拓扑》，定义
\ref{topologies-definition-big-zariski-site} 之后的注释。

\begin{definition}
\label{flat-definition-standard-h}
设 $T$ 是仿射概形。$T$ 的一个{\it 标准 h 覆盖}是态射族
$\{f_i : T_i \to T\}_{i = 1, \ldots, n}$，其中每个 $T_i$
都仿射、每个 $f_i$ 都为有限表示，并满足下列等价条件之一：
(1) $\{T_i \to T\}$ 可由标准 ph 覆盖加细；或
(2) $\{T_i \to T\}$ 是 V 覆盖。
\end{definition}

\noindent
这些条件的等价性由引理
\ref{flat-lemma-equivalence-h-v-locally-finite-presentation}、
《拓扑》，定义 \ref{topologies-definition-ph-covering}
和引理 \ref{topologies-lemma-refine-by-ph} 推出。

\medskip\noindent
在继续引入概形 $S$ 的大 h 站点之前，先指出：
大 h 站点 $\Sch_h$ 上的拓扑在某种意义下由全体概形范畴上的
h 拓扑诱导而来。

\begin{lemma}
\label{flat-lemma-h-induced}
设 $\Sch_h$ 是定义 \ref{flat-definition-big-h-site} 中的大 h 站点，
$T \in \Ob(\Sch_h)$，并设
$\{T_i \to T\}_{i \in I}$ 是 $T$ 的任意 h 覆盖。
\begin{enumerate}
\item 站点 $\Sch_h$ 中存在 $T$ 的覆盖
$\{U_j \to T\}_{j \in J}$，它加细 $\{T_i \to T\}_{i \in I}$。
\item 若 $\{T_i \to T\}_{i \in I}$ 是标准 h 覆盖，
则它按定义自然等价于 $\Sch_h$ 的一个覆盖。
\item 若 $\{T_i \to T\}_{i \in I}$ 是 Zariski 覆盖，
则它按定义自然等价于 $\Sch_h$ 的一个覆盖。
\end{enumerate}
\end{lemma}

\begin{proof}
证明从略。提示：这与 《拓扑》，引理
\ref{topologies-lemma-ph-induced} 的证明完全相同。
\end{proof}

\begin{definition}
\label{flat-definition-big-small-h}
设 $S$ 是概形，$\Sch_h$ 是包含 $S$ 的大 h 站点。
\begin{enumerate}
\item $S$ 的{\it 大 h 站点}记为 $(\Sch/S)_h$，
即 《站点》，第 \ref{sites-section-localize} 节中引入的站点
$\Sch_h/S$。
\item $S$ 的{\it 大仿射 h 站点}记为
$(\textit{Aff}/S)_h$，它是 $(\Sch/S)_h$ 的满子范畴，
其对象为仿射的 $U/S$。$(\textit{Aff}/S)_h$ 的覆盖是
$(\Sch/S)_h$ 的任意覆盖 $\{U_i \to U\}$，并要求它是标准 h 覆盖。
\end{enumerate}
\end{definition}

\noindent
我们明确陈述大仿射 h 站点确实是一个站点。

\begin{lemma}
\label{flat-lemma-verify-site-h}
设 $S$ 是概形，$\Sch_h$ 是包含 $S$ 的大 h 站点。
则 $(\textit{Aff}/S)_h$ 是一个站点。
\end{lemma}

\begin{proof}
仿照 《拓扑》，引理
\ref{topologies-lemma-verify-site-etale} 的证明，只需说明
标准 h 覆盖族满足 《站点》，定义
\ref{sites-definition-site} 的性质 (1)、(2)、(3)。
这很清楚。例如，给定标准 h 覆盖
$\{T_i \to T\}_{i\in I}$，并对每个 $i$ 给定标准 h 覆盖
$\{T_{ij} \to T_i\}_{j \in J_i}$，则
$\{T_{ij} \to T\}_{i \in I, j\in J_i}$ 是 h 覆盖
（引理 \ref{flat-lemma-h}），$\bigcup_{i\in I} J_i$ 是有限集，
而每个 $T_{ij}$ 都仿射。因此
$\{T_{ij} \to T\}_{i \in I, j\in J_i}$ 是标准 h 覆盖。
\end{proof}

\begin{lemma}
\label{flat-lemma-fibre-products-h}
设 $S$ 是概形，$\Sch_h$ 是包含 $S$ 的大 h 站点。
站点 $\Sch_h$、$(\Sch/S)_h$ 和 $(\textit{Aff}/S)_h$
的底层范畴都有纤维积。在每种情形中，到全体概形范畴
$\Sch$ 的显然函子均与取纤维积交换。范畴
$(\Sch/S)_h$ 有终对象，即 $S/S$。
\end{lemma}

\begin{proof}
对 $\Sch_h$，结论依构造成立；参见《集合》，引理
\ref{sets-lemma-what-is-in-it}。设有概形态射
$U \to S$、$V \to U$、$W \to U$，其中
$U, V, W \in \Ob(\Sch_h)$。$\Sch_h$ 中的纤维积
$V \times_U W$ 也是 $\Sch$ 中的纤维积；在全体 $S$ 上概形的范畴中，
它是 $V/S$ 与 $W/S$ 在 $U/S$ 上的纤维积，因而也是
$(\Sch/S)_h$ 中的纤维积。这证明了 $(\Sch/S)_h$ 的结论。
若 $U, V, W$ 均仿射，则 $V \times_U W$ 也仿射，
故 $(\textit{Aff}/S)_h$ 的结论也成立。
\end{proof}

\noindent
下面验证大仿射站点与大站点定义同一个拓扑斯。

\begin{lemma}
\label{flat-lemma-affine-big-site-h}
设 $S$ 是概形，$\Sch_h$ 是包含 $S$ 的大 h 站点。
函子 $(\textit{Aff}/S)_h \to (\Sch/S)_h$ 是余连续的，
并诱导从 $\Sh((\textit{Aff}/S)_h)$ 到
$\Sh((\Sch/S)_h)$ 的拓扑斯等价。
\end{lemma}

\begin{proof}
特殊余连续函子的概念见《站点》，定义
\ref{sites-definition-special-cocontinuous-functor}。
因此需验证 《站点》，引理 \ref{sites-lemma-equivalence}
的假设 (1) -- (5)。记包含函子为
$u : (\textit{Aff}/S)_h \to (\Sch/S)_h$。
若 $T$ 仿射，则 $T/S$ 的任意 h 覆盖都可由标准 h 覆盖加细，
例如可用引理 \ref{flat-lemma-approximate-h-cover}；故余连续性成立，
从而 (1) 成立。由于标准 h 覆盖本身就是 h 覆盖，$u$ 显然连续，
故 (2) 成立。(3) 和 (4) 立即由 $u$ 全忠实推出。
最后，每个概形都有仿射开覆盖（这是 h 覆盖），故 (5) 成立。
\end{proof}

\begin{lemma}
\label{flat-lemma-characterize-sheaf-h}
设 $\mathcal{F}$ 是 $(\Sch/S)_h$ 上的预层。
则 $\mathcal{F}$ 是层，当且仅当
\begin{enumerate}
\item $\mathcal{F}$ 对 Zariski 覆盖满足层条件；
\item 若 $f : V \to U$ 适当、满且为有限表示，则
$\mathcal{F}(U)$ 双射地映到两个映射
$\mathcal{F}(V) \to \mathcal{F}(V \times_U V)$ 的等化子。
\end{enumerate}
此外，在 (1) 成立时，性质 (2) 等价于性质
\begin{enumerate}
\item[(2')] 对 (2) 中 $U$ 仿射的 $\{V \to U\}$ 满足层性质。
\end{enumerate}
\end{lemma}

\begin{proof}
先证明若 (1) 和 (2) 成立，则 $\mathcal{F}$ 是层。
设 $\{T_i \to T\}$ 是 $(\Sch/S)_h$ 中的覆盖。
我们验证该覆盖的层条件。设 $s_i \in \mathcal{F}(T_i)$ 是一族截面，
它们在 $T_i \times_T T_{i'}$ 上限制为同一截面。
我们将证明存在唯一截面 $s \in \mathcal{F}(T)$，
它在 $T_i$ 上限制为 $s_i$。设
$T = \bigcup U_j$ 是仿射开覆盖。由性质 (1)，为构造 $s$，
只需构造在 $U_j \cap U_{j'}$ 上相符的截面
$s_j \in \mathcal{F}(U_j)$。考虑覆盖
$\{T_i \times_T U_j \to U_j\}$。则截面
$s_{ji} = s_i|_{T_i \times_T U_j}$ 在
$(T_i \times_T U_j) \times_{U_j} (T_{i'} \times_T U_j)$ 上相符。
选取有限表示的适当满射 $V_j \to U_j$ 和有限仿射开覆盖
$V_j = \bigcup V_{jk}$，使 $\{V_{jk} \to U_j\}$
加细 $\{T_i \times_T U_j \to U_j\}$；参见引理
\ref{flat-lemma-approximate-h-cover}。若
$s_{jk} \in \mathcal{F}(V_{jk})$ 表示通过隐含态射将 $s_{ji}$
拉回到 $V_{jk}$ 所得的截面，则 $s_{jk}$ 黏合为截面
$s'_j \in \mathcal{F}(V_j)$。再次使用交叠处的相符性，可见
$s'_j$ 属于两个映射
$\mathcal{F}(V_j) \to \mathcal{F}(V_j \times_{U_j} V_j)$ 的等化子。
故由 (2)，$s'_j$ 来自唯一截面
$s_j \in \mathcal{F}(U_j)$。这些截面 $s_j$
具有全部所需性质的验证从略。

\medskip\noindent
现证明在 (1) 成立时 (2) 与 (2') 等价。设
$V \to U$ 是 $(\Sch/S)_h$ 中适当、满且为有限表示的态射。
选取仿射开覆盖 $U = \bigcup U_i$，并令
$V_i = V \times_U U_i$。由 (2')，映射
$\mathcal{F}(U_i) \to \mathcal{F}(V_i)$ 为单射；
由 (1)，映射 $\mathcal{F}(U) \to \prod \mathcal{F}(U_i)$
为单射。因此 $\mathcal{F}(U) \to \mathcal{F}(V)$ 为单射。
最后，设给定 $t \in \mathcal{F}(V)$，它属于两个映射
$\mathcal{F}(V) \to \mathcal{F}(V \times_U V)$ 的等化子。
则对所有 $i$，$t|_{V_i}$ 属于两个映射
$\mathcal{F}(V_i) \to \mathcal{F}(V_i \times_{U_i} V_i)$
的等化子。故由 (2')，对所有 $i$，得到唯一截面
$s_i \in \mathcal{F}(U_i)$，它映到 $t|_{V_i}$。
对所有 $i, j$ 验证
$s_i|_{U_i \cap U_j} = s_j|_{U_i \cap U_j}$ 从略；
这里使用刚证明的唯一性。由覆盖 $U = \bigcup U_i$ 的层性质，
得到截面 $s \in \mathcal{F}(U)$。$s$ 在
$\mathcal{F}(V)$ 中映到 $t$ 的证明从略。
\end{proof}

\noindent
下面建立与这些站点相关的拓扑斯之间的一些关系。

\begin{lemma}
\label{flat-lemma-morphism-big-h}
设 $\Sch_h$ 是大 h 站点，$f : T \to S$ 是 $\Sch_h$ 中的态射。
函子
$$
u : (\Sch/T)_h \longrightarrow (\Sch/S)_h,
\quad
V/T \longmapsto V/S
$$
是余连续的，并有连续右伴随
$$
v : (\Sch/S)_h \longrightarrow (\Sch/T)_h,
\quad
(U \to S) \longmapsto (U \times_S T \to T).
$$
它们诱导同一个拓扑斯态射
$$
f_{big} :
\Sh((\Sch/T)_h)
\longrightarrow
\Sh((\Sch/S)_h)
$$
此外，
$f_{big}^{-1}(\mathcal{G})(U/T) = \mathcal{G}(U/S)$，
$f_{big, *}(\mathcal{F})(U/S) = \mathcal{F}(U \times_S T/T)$。
此外，$f_{big}^{-1}$ 有左伴随 $f_{big!}$，后者与纤维积和等化子交换。
\end{lemma}

\begin{proof}
函子 $u$ 余连续、连续，并与纤维积和等化子交换。因此可应用
《站点》，引理 \ref{sites-lemma-when-shriek} 和
\ref{sites-lemma-preserve-equalizers}，从而得到
$f_{big}^{-1}$ 的公式及 $f_{big!}$ 的存在性。此外，函子 $v$
是右伴随，因为给定 $U/T$ 与 $V/S$，有预期的等式
$\Mor_S(u(U), V) = \Mor_T(U, V \times_S T)$。
故可应用 《站点》，引理
\ref{sites-lemma-have-functor-other-way} 和
\ref{sites-lemma-have-functor-other-way-morphism}，
得到 $f_{big, *}$ 的公式。
\end{proof}

\begin{lemma}
\label{flat-lemma-composition-h}
% P05-EN-CH38-0069: frozen source lists X, Y, Y although the following morphism uses Z; candidate third scheme is Z.
给定 $(\Sch/S)_h$ 中的概形 $X$、$Y$、$Y$ 以及态射
$f : X \to Y$、$g : Y \to Z$，有
$g_{big} \circ f_{big} = (g \circ f)_{big}$.
\end{lemma}

\begin{proof}
这由引理 \ref{flat-lemma-morphism-big-h} 中大站点上函子的
推前与拉回的简单描述推出。
\end{proof}










\section{h 拓扑的进一步结果}
\label{flat-section-h-more}

\noindent
本节证明关于 h 拓扑的若干进一步结果。先看一些反例。

\begin{example}
\label{flat-example-structure-sheaf-not-h-sheaf}
“结构层”$\mathcal{O}$ 不是 h 拓扑中的层。例如，考虑有限表示的
满闭浸入 $X \to Y$。由引理
\ref{flat-lemma-surjective-proper-finite-presentation-h}，
$\{X \to Y\}$ 是 h 覆盖。还要注意 $X \times_Y X = X$。
% P05-EN-CH38-0070: frozen source has “where a sheaf”; candidate “were a sheaf”.
因此，若 $\mathcal{O}$ 是 h 拓扑中的层，则
$\mathcal{O}_Y(Y) \to \mathcal{O}_X(X)$ 应为双射。
一旦 $X$, $Y$ 仿射而态射 $X \to Y$ 不是同构，事实便非如此。
\end{example}

\begin{example}
\label{flat-example-representable-sheaf-not-h-sheaf}
在任意站点 $(\Sch/S)_h$ 上，该拓扑都不是次典范的；
换言之，可表预层未必是层。事实上，“结构层”$\mathcal{O}$
是可表的，因为在 $(\Sch/S)_h$ 中
$\mathcal{O}(X) = \Mor_S(X, \mathbf{A}^1_S)$；
而例 \ref{flat-example-structure-sheaf-not-h-sheaf} 已说明
$\mathcal{O}$ 不是层。
\end{example}

\begin{lemma}
\label{flat-lemma-limit-h-topology}
设 $T$ 是仿射概形，并被写成一族仿射概形所成有向逆系统的极限
$T = \lim_{i \in I} T_i$。
\begin{enumerate}
\item 设 $\mathcal{V} = \{V_j \to T\}_{j = 1, \ldots, m}$
是 $T$ 的标准 h 覆盖；参见定义 \ref{flat-definition-standard-h}。
则存在指标 $i$ 和标准 h 覆盖
$\mathcal{V}_i = \{V_{i, j} \to T_i\}_{j = 1, \ldots, m}$，
使它到 $T$ 的基变换 $T \times_{T_i} \mathcal{V}_i$
同构于 $\mathcal{V}$。
\item 设 $\mathcal{V}_i$, $\mathcal{V}'_i$ 是 $T_i$
的两个标准 h 覆盖。若
$f : T \times_{T_i} \mathcal{V}_i \to T \times_{T_i} \mathcal{V}'_i$
是 $T$ 的覆盖之间的态射，则存在指标 $i' \geq i$ 和态射
% P05-EN-CH38-0071: frozen source uses unindexed \mathcal{V} in the domain below although this item starts from \mathcal{V}_i.
$f_{i'} : T_{i'} \times_{T_i} \mathcal{V} \to
T_{i'} \times_{T_i} \mathcal{V}'_i$
，其到 $T$ 的基变换为 $f$。
\item 若
% P05-EN-CH38-0072: frozen source again uses \mathcal{V} rather than \mathcal{V}_i for a covering of T_i.
$f, g : \mathcal{V} \to \mathcal{V}'_i$
是 $T_i$ 的标准 h 覆盖之间的态射，并且它们到 $T$
的基变换 $f_T, g_T$ 相等，则存在指标 $i' \geq i$，
使 $f_{T_{i'}} = g_{T_{i'}}$。
\end{enumerate}
换言之，$T$ 的标准 h 覆盖范畴是 $T_i$ 的标准 h 覆盖范畴
在 $I$ 上的余极限。
\end{lemma}

\begin{proof}
由《极限》，引理 \ref{limits-lemma-descend-finite-presentation}，
$T$ 上有限表示概形的范畴是 $T_i$ 上有限表示概形范畴
在 $I$ 上的余极限。
% P05-EN-CH38-0073: frozen source says “categories of finite presentation over T_i”; candidate inserts “schemes”.
由《极限》，引理
\ref{limits-lemma-descend-affine-finite-presentation}，
对 $T$ 上仿射且有限表示的概形范畴也有相同结论。
为完成引理的证明，只需证明如下断言：若
$\{V_{j, i} \to T_i\}_{j = 1, \ldots, m}$ 是有限表示态射的有限族，
每个 $V_{j, i}$ 仿射，且基变换族
$\{T \times_{T_i} V_{j, i} \to T\}$ 是 h 覆盖，
则对某个 $i' \geq i$，态射族
$\{T_{i'} \times_{T_i} V_{j, i} \to T_{i'}\}$ 是 h 覆盖。
为此，利用引理 \ref{flat-lemma-approximate-h-cover} 选取有限表示的
适当满射 $Y \to T$ 和有限仿射开覆盖
$Y = \bigcup_{k = 1, \ldots, n} Y_k$，使
$\{Y_k \to T\}_{k = 1, \ldots, n}$ 加细
$\{T \times_{T_i} V_{j, i} \to T\}$。
利用上面的论证及 《极限》，引理
\ref{limits-lemma-eventually-proper}、
\ref{limits-lemma-descend-surjective} 和
\ref{limits-lemma-descend-opens}，可找到 $i' \geq i$、
有限表示的适当满射 $Y_{i'} \to T_{i'}$ 以及仿射开覆盖
$Y_{i'} = \bigcup_{k = 1, \ldots, n} Y_{i', k}$，
并使 $\{Y_{i', k} \to Y_{i'}\}$ 加细
$\{T_{i'} \times_{T_i} V_{j, i} \to T_{i'}\}$。
因此最后这个态射族是 h 覆盖，证明完毕。
\end{proof}

\begin{lemma}
\label{flat-lemma-extend-sheaf-h}
设 $S$ 是包含在大站点 $\Sch_h$ 中的概形。设
$F : (\Sch/S)_h^{opp} \to \textit{Sets}$ 是 h 层，
并在 $\mathcal{C} = (\Sch/S)_h$ 时满足 《拓扑》，引理
\ref{topologies-lemma-extend} 的性质 (b)。
则 $F$ 到全体 $S$ 上概形范畴的延拓 $F'$ 对所有 h 覆盖
满足层条件，并保持极限（《极限》，注记
\ref{limits-remark-limit-preserving}）。
\end{lemma}

\begin{proof}
使用引理 \ref{flat-lemma-limit-h-topology} 和 \ref{flat-lemma-h-induced}，
并采用 《拓扑》，引理
\ref{topologies-lemma-extend-sheaf-general} 和
\ref{topologies-lemma-extend-sheaf} 证明中的论证即可。细节从略。
\end{proof}








\section{爆破方块与 ph 拓扑}
\label{flat-section-blow-up-ph}

\noindent
设 $X$ 是概形，$Z \subset X$ 是闭子概形，且包含态射为有限表示；
换言之，与 $Z$ 对应的拟相干理想层为有限型。
设 $b : X' \to X$ 是 $X$ 沿 $Z$ 的爆破，并令
$E = b^{-1}(Z)$ 为例外除子。参见《因子》，第 \ref{divisors-section-blowing-up} 节。在此情形下，并在本节中，
称
\begin{equation}
\label{flat-equation-blow-up-square}
\vcenter{
\xymatrix{
E \ar[d] \ar[r] & X' \ar[d]^b \\
Z \ar[r] & X
}
}
\end{equation}
为一个{\it 爆破方块}。

\begin{lemma}
\label{flat-lemma-blow-up-square-ph}
设 $\mathcal{F}$ 是站点 $(\Sch/S)_{ph}$ 上的层；参见
《拓扑》，定义 \ref{topologies-definition-big-small-ph}。
则对范畴 $(\Sch/S)_{ph}$ 中的任意爆破方块
(\ref{flat-equation-blow-up-square})，图
$$
\xymatrix{
\mathcal{F}(E) & \mathcal{F}(X') \ar[l] \\
\mathcal{F}(Z) \ar[u] & \mathcal{F}(X) \ar[u] \ar[l]
}
$$
在集合范畴中为笛卡儿图。
\end{lemma}

\begin{proof}
由于 $Z \amalg X' \to X$ 是适当满射，
$\{Z \amalg X' \to X\}$ 是 ph 覆盖
（《拓扑》，引理 \ref{topologies-lemma-surjective-proper-ph}）。
并且
$$
(Z \amalg X') \times_X (Z \amalg X') =
Z \amalg E \amalg E \amalg X' \times_X X'
$$
由于 $\mathcal{F}$ 是 Zariski 层，$\mathcal{F}$ 将不交并映到乘积。
因此，覆盖 $\{Z \amalg X' \to X\}$ 的层条件说明
$\mathcal{F}(X) \to \mathcal{F}(Z) \times \mathcal{F}(X')$
为单射，其像是满足下列条件的对 $(t, s')$ 所成集合：
(a) $t|_E = s'|_E$；(b) $s'$ 属于两个映射
$\mathcal{F}(X') \to \mathcal{F}(X' \times_X X')$ 的等化子。
接着注意，显然态射
$$
E \times_Z E \amalg X' \longrightarrow X' \times_X X'
$$
是适当满射，因为 $b$ 诱导同构
$X' \setminus E \to X \setminus Z$。故
$\mathcal{F}(X' \times_X X') \to
\mathcal{F}(E \times_Z E) \times \mathcal{F}(X')$ 为单射。
于是 (a) $\Rightarrow$ (b)，引理得证。
\end{proof}

\begin{lemma}
\label{flat-lemma-thickening-ph}
设 $\mathcal{F}$ 是 《拓扑》，定义
\ref{topologies-definition-big-small-ph} 中站点
$(\Sch/S)_{ph}$ 上的层。设 $X \to X'$ 是
$(\Sch/S)_{ph}$ 中的态射，并且是增厚。则
$\mathcal{F}(X') \to \mathcal{F}(X)$ 为双射。
\end{lemma}

\begin{proof}
% P05-EN-CH38-0074: frozen source has stray “of” in “a proper surjective morphism of and”.
注意，$X \to X'$ 是适当满射，且 $X \times_{X'} X = X$。
对 ph 覆盖 $\{X \to X'\}$ 应用层性质
（《拓扑》，引理
\ref{topologies-lemma-surjective-proper-ph}）即得结论。
\end{proof}








\section{几乎爆破方块与 h 拓扑}
\label{flat-section-blow-up-h}

\noindent
考虑爆破方块 (\ref{flat-equation-blow-up-square})。虽然态射
$b : X' \to X$ 是投射态射
（《因子》，引理 \ref{divisors-lemma-blowing-up-projective}），
一般却没有简单方法保证 $b$ 为有限表示。由于 h 覆盖使用有限表示态射
构造，我们需要一个变体。具体地，称概形的交换图
\begin{equation}
\label{flat-equation-almost-blow-up-square}
\vcenter{
\xymatrix{
E \ar[d] \ar[r] & X' \ar[d]^b \\
Z \ar[r] & X
}
}
\end{equation}
为一个{\it 几乎爆破方块}，如果满足下列条件：
\begin{enumerate}
\item $Z \to X$ 是有限表示的闭浸入；
\item $E = b^{-1}(Z)$ 是 $X'$ 的局部主闭子概形；
\item $b$ 适当且为有限表示；
\item 由 $\mathcal{O}_{X'}$ 中支集在 $E$ 上的截面所成拟相干理想
切出的闭子概形 $X'' \subset X'$
（《性质》，引理
\ref{properties-lemma-sections-supported-on-closed-subset}）
是 $X$ 沿 $Z$ 的爆破。
\end{enumerate}
由此可知，态射 $b$ 诱导同构
$X' \setminus E \to X \setminus Z$。
几乎爆破方块的一些很简单的例子参见例
\ref{flat-example-one-generator} 和 \ref{flat-example-two-generators}。

\medskip\noindent
爆破的基变换通常不是爆破，但几乎爆破与基变换相容。

\begin{lemma}
\label{flat-lemma-base-change-almost-blow-up}
考虑几乎爆破方块 (\ref{flat-equation-almost-blow-up-square})，
并设 $Y \to X$ 为任意态射。则基变换
$$
\xymatrix{
Y \times_X E \ar[d] \ar[r] & Y \times_X X' \ar[d] \\
Y \times_X Z \ar[r] & Y
}
$$
也是几乎爆破方块。
\end{lemma}

\begin{proof}
由《态射》，引理 \ref{morphisms-lemma-base-change-proper} 和
\ref{morphisms-lemma-base-change-finite-presentation}，
态射 $Y \times_X X' \to Y$ 适当且为有限表示。态射
$Y \times_X Z \to Y$ 是有限表示的闭浸入
（《态射》，引理
\ref{morphisms-lemma-base-change-closed-immersion}）。
$Y \times_X Z$ 在 $Y \times_X X'$ 中的逆像等于 $E$
在 $Y \times_X X'$ 中的逆像，因而局部为主
（《因子》，引理
\ref{divisors-lemma-pullback-locally-principal}）。
分别令 $X'' \subset X'$ 和 $Y'' \subset Y \times_X X'$
为与 $\mathcal{O}_{X'}$ 和
% P05-EN-CH38-0075: frozen source has Y \times_Y X' where the ambient space is Y \times_X X'.
$\mathcal{O}_{Y \times_Y X'}$ 中分别支集在 $E$ 和
$Y \times_X E$ 上的截面所成拟相干理想对应的闭子概形。
显然，$Y'' \subset Y \times_X X''$ 是与
% P05-EN-CH38-0076: frozen source has Y \times_Y X'' where the ambient base change is Y \times_X X''.
$\mathcal{O}_{Y \times_Y X''}$ 中支集在
$Y \times_X (E \cap X'')$ 上的截面所成拟相干理想对应的闭子概形。
因此 $Y''$ 是 $Y$ 关于爆破 $X'' \to X$ 的严格变换；
参见《因子》，定义
\ref{divisors-definition-strict-transform}。故由《因子》，引理
\ref{divisors-lemma-strict-transform}，$Y''$
是 $Y$ 沿 $Y \times_X Z$ 的爆破。
\end{proof}

\noindent
可以缩小几乎爆破方块。

\begin{lemma}
\label{flat-lemma-shrink-almost-blow-up}
考虑几乎爆破方块 (\ref{flat-equation-almost-blow-up-square})，
并设 $W \to X'$ 是有限表示的闭浸入。下列条件等价：
\begin{enumerate}
\item $X' \setminus E$ 在概形论意义下包含于 $W$；
\item $X$ 沿 $Z$ 的爆破 $X''$ 在概形论意义下包含于 $W$；
\item 图
$$
\xymatrix{
E \cap W \ar[d] \ar[r] & W \ar[d] \\
Z \ar[r] & X
}
$$
是几乎爆破方块。这里 $E \cap W$ 是概形论交。
\end{enumerate}
\end{lemma}

\begin{proof}
假设 (1)。则满射 $\mathcal{O}_{X'} \to \mathcal{O}_W$
在开集
% P05-EN-CH38-0077: frozen source writes X' \subset E; condition (1) and the argument require X' \setminus E.
$X' \subset E$ 上是同构。由于 $X'' \subset X'$ 的理想层
由 $\mathcal{O}_{X'}$ 中支集在 $E$ 上的截面组成
（依我们对几乎爆破方块的定义），可知 (2) 成立。
若 (2) 成立，则 (3) 成立。若 (3) 成立，则 (1) 成立，
因为 $X'' \cap (X' \setminus E)$ 同构于 $X \setminus Z$，
而后者又同构于 $X' \setminus E$。
\end{proof}

\noindent
实际爆破是任意给定几乎爆破的各个缩小之极限。

\begin{lemma}
\label{flat-lemma-blow-up-limit-almost-blow-up}
考虑几乎爆破方块 (\ref{flat-equation-almost-blow-up-square})，
其中 $X$ 拟紧且拟分离。则 $X$ 沿 $Z$ 的爆破 $X''$ 可写成
$$
X'' = \lim X'_i
$$
这里的极限取遍由有限表示闭子概形 $X'_i \subset X'$
组成的有向系统；这些闭子概形满足引理
\ref{flat-lemma-shrink-almost-blow-up} 的等价条件。
\end{lemma}

\begin{proof}
令 $\mathcal{I} \subset \mathcal{O}_{X'}$ 为与 $X''$
对应的拟相干理想层。由《性质》，引理
\ref{properties-lemma-quasi-coherent-colimit-finite-type}，
可将 $\mathcal{I}$ 写成其有限型拟相干子模的滤过余极限
$\mathcal{I} = \colim \mathcal{I}_i$。这些模与闭子概形
$X'_i$ 按 $1$-对-$1$ 对应，故证明完毕。
\end{proof}

\noindent
几乎爆破方块确实存在。

\begin{lemma}
\label{flat-lemma-almost-blow-up-square}
设 $X$ 是拟紧且拟分离概形，$Z \subset X$ 是由有限型拟相干理想层
切出的闭子概形。则存在如
(\ref{flat-equation-almost-blow-up-square}) 所示的几乎爆破方块。
\end{lemma}

\begin{proof}
可以将 $X = \lim X_i$ 写成一族 Noether 概形所成逆系统的有向极限，
其转移态射均为仿射态射；参见《极限》，命题
\ref{limits-proposition-approximate}。可以找到指标 $i$
和闭浸入 $Z_i \to X_i$，其到 $X$ 的基变换是闭浸入
$Z \to X$。参见《极限》，引理
\ref{limits-lemma-descend-finite-presentation} 和
\ref{limits-lemma-descend-closed-immersion-finite-presentation}。
令 $b_i : X'_i \to X_i$ 为以 $Z_i$ 为中心的爆破。
这给出爆破方块
$$
\xymatrix{
E_i \ar[r] \ar[d] & X'_i \ar[d]^{b_i} \\
Z_i \ar[r] & X_i
}
$$
其中所有态射都是 Noether 概形之间的有限型态射，因而为有限表示。
所以这是几乎爆破方块。由引理
\ref{flat-lemma-base-change-almost-blow-up}，将此图基变换到 $X$
即可得到所需的几乎爆破方块。
\end{proof}

\noindent
几乎爆破方块在引理 \ref{flat-lemma-shrink-almost-blow-up}
所述的缩小意义下是唯一的。

\begin{lemma}
\label{flat-lemma-almost-blow-up-unique}
设 $X$ 是拟紧且拟分离概形，$Z \subset X$
是由有限型拟相干理想层切出的闭子概形。设给定几乎爆破方块
(\ref{flat-equation-almost-blow-up-square})
$$
\xymatrix{
E_k \ar[r] \ar[d] & X_k' \ar[d] \\
Z \ar[r] & X
}
$$
其中 $k = 1, 2$。则存在几乎爆破方块
$$
\xymatrix{
E \ar[r] \ar[d] & X' \ar[d] \\
Z \ar[r] & X
}
$$
以及 $X$ 上闭浸入 $i_k : X' \to X'_k$，使
$E = i_k^{-1}(E_k)$。
\end{lemma}

\begin{proof}
以 $X'' \to X$ 表示 $X$ 沿 $Z$ 的爆破。将 $X''$
同时视为 $X'_1$ 和 $X'_2$ 的闭子概形。如引理
\ref{flat-lemma-blow-up-limit-almost-blow-up} 中那样写
$X'' = \lim X'_{1, i}$。由《极限》，命题
\ref{limits-proposition-characterize-locally-finite-presentation}，
存在某个 $i$ 和态射 $h : X'_{1, i} \to X'_2$，
它与包含 $X'' \subset X'_{1, i}$ 及
$X'' \subset X'_2$ 相容。由《极限》，引理
\ref{limits-lemma-eventually-closed-immersion}，对某个
$i' \geq i$，$h$ 在 $X'_{1, i'}$ 上的限制是闭浸入。
证明完毕。
\end{proof}

\noindent
在适当情形下，我们关于爆破的扁平化技巧也适用于几乎爆破。

\begin{lemma}
\label{flat-lemma-flat-after-almost-blowing-up}
设 $Y$ 是拟紧且拟分离概形，$X$ 是 $Y$ 上有限表示概形。
设 $V \subset Y$ 是拟紧开子集，且 $X_V \to V$ 平坦。
则存在交换图
$$
\xymatrix{
E \ar[ddd] \ar[rd] & & & D \ar[lll] \ar[ddd] \ar[ld] \\
& Y' \ar[d] & X' \ar[l] \ar[d] \\
& Y & X \ar[l] \\
Z \ar[ru] & & & T \ar[lll] \ar[lu]
}
$$
其左右两个方块为几乎爆破方块，上下两个方块为笛卡儿方块，
并且 $Z \cap V = \emptyset$，而 $X' \to Y'$ 平坦
（且为有限表示）。
\end{lemma}

\begin{proof}
若 $Y$ 是 Noether 概形，则引理立即由引理
\ref{flat-lemma-flat-after-blowing-up} 推出，因为在此情形下，
爆破方块就是几乎爆破方块（还要用到严格变换是爆破）。
一般情形通过绝对 Noether 逼近化归到 Noether 情形。

\medskip\noindent
可以将 $Y = \lim Y_i$ 写成一族 Noether 概形所成逆系统的有向极限，
其转移态射均仿射；参见《极限》，命题
\ref{limits-proposition-approximate}。可以找到指标 $i$
和有限表示态射 $X_i \to Y_i$，其到 $Y$ 的基变换是
$X \to Y$；参见《极限》，引理
\ref{limits-lemma-descend-finite-presentation}。增大 $i$ 后，
可以假设 $V$ 是某个开子概形 $V_i \subset Y_i$ 的逆像；
参见《极限》，引理 \ref{limits-lemma-descend-opens}。
最后，再增大 $i$，可以假设 $X_{i, V_i} \to V_i$ 平坦；
参见《极限》，引理
\ref{limits-lemma-descend-flat-finite-presentation}。
由 Noether 情形，可对 $X_i \to Y_i \supset V_i$
构造引理中的图。沿 $Y \to Y_i$ 对此图作基变换即得所求。
这里使用了基变换保持相关态射性质；参见《态射》，引理
\ref{morphisms-lemma-base-change-proper},
\ref{morphisms-lemma-base-change-finite-presentation},
\ref{morphisms-lemma-base-change-closed-immersion} 和
\ref{morphisms-lemma-base-change-flat}
；还使用了几乎爆破方块的基变换仍为几乎爆破方块，
参见引理 \ref{flat-lemma-base-change-almost-blow-up}。
\end{proof}

\begin{lemma}
\label{flat-lemma-blow-up-square-h}
设 $\mathcal{F}$ 是定义 \ref{flat-definition-big-small-h}
中构造的某个站点 $(\Sch/S)_h$ 上的层。
则对范畴 $(\Sch/S)_h$ 中的任意几乎爆破方块
(\ref{flat-equation-almost-blow-up-square})，图
$$
\xymatrix{
\mathcal{F}(E) & \mathcal{F}(X') \ar[l] \\
\mathcal{F}(Z) \ar[u] & \mathcal{F}(X) \ar[u] \ar[l]
}
$$
在集合范畴中为笛卡儿图。
\end{lemma}

\begin{proof}
由于 $Z \amalg X' \to X$ 是有限表示的适当满射，
$\{Z \amalg X' \to X\}$ 是 h 覆盖
（引理 \ref{flat-lemma-surjective-proper-finite-presentation-h}）。
并且
$$
(Z \amalg X') \times_X (Z \amalg X') =
Z \amalg E \amalg E \amalg X' \times_X X'
$$
由于 $\mathcal{F}$ 是 Zariski 层，$\mathcal{F}$ 将不交并映到乘积。
因此覆盖 $\{Z \amalg X' \to X\}$ 的层条件说明
$\mathcal{F}(X) \to \mathcal{F}(Z) \times \mathcal{F}(X')$
为单射，其像是满足下列条件的对 $(t, s')$ 所成集合：
(a) $t|_E = s'|_E$；(b) $s'$ 属于两个映射
$\mathcal{F}(X') \to \mathcal{F}(X' \times_X X')$ 的等化子。
接着注意，显然态射
$$
E \times_Z E \amalg X' \longrightarrow X' \times_X X'
$$
是有限表示的适当满射，因为 $b$ 诱导同构
$X' \setminus E \to X \setminus Z$。故
$\mathcal{F}(X' \times_X X') \to
\mathcal{F}(E \times_Z E) \times \mathcal{F}(X')$ 为单射。
于是 (a) $\Rightarrow$ (b)，引理得证。
\end{proof}

\begin{lemma}
\label{flat-lemma-thickening-h}
设 $\mathcal{F}$ 是定义 \ref{flat-definition-big-small-h}
中构造的某个站点 $(\Sch/S)_h$ 上的层。设
$X \to X'$ 是 $(\Sch/S)_h$ 中有限表示的增厚。则
$\mathcal{F}(X') \to \mathcal{F}(X)$ 为双射。
\end{lemma}

\begin{proof}
第一种证明。注意，$X \to X'$ 是有限表示的适当满射，且
$X \times_{X'} X = X$。对 h 覆盖 $\{X \to X'\}$
应用层性质（引理
\ref{flat-lemma-surjective-proper-finite-presentation-h}）即得结论。

\medskip\noindent
第二种证明（有些可笑）。$X'$ 沿 $X$ 的爆破是空概形。
理由是，若 $a$ 是 $A$ 的幂零元素，则仿射爆破代数
$A[\frac{I}{a}]$（《代数》，第 \ref{algebra-section-blow-up} 节）为零。细节从略。
因此得到如下形式的几乎爆破方块：
$$
\xymatrix{
\emptyset \ar[r] \ar[d] & \emptyset \ar[d] \\
X \ar[r] & X'
}
$$
由于 $\mathcal{F}$ 是层，$\mathcal{F}(\emptyset)$ 是单点集。
应用引理 \ref{flat-lemma-blow-up-square-h} 即得结论。
\end{proof}

\begin{proposition}
\label{flat-proposition-check-h}
设 $\mathcal{F}$ 是定义 \ref{flat-definition-big-small-h}
中构造的某个站点 $(\Sch/S)_h$ 上的预层。
则 $\mathcal{F}$ 是层，当且仅当满足下列条件：
\begin{enumerate}
\item $\mathcal{F}$ 是 Zariski 拓扑下的层；
\item 给定 $(\Sch/S)_h$ 中的态射 $f : X \to Y$，
其中 $Y$ 仿射，而 $f$ 满、平坦、适当且为有限表示，
则 $\mathcal{F}(Y)$ 是两个映射
$\mathcal{F}(X) \to \mathcal{F}(X \times_Y X)$ 的等化子；
\item 给定范畴 $(\Sch/S)_h$ 中的几乎爆破方块
(\ref{flat-equation-almost-blow-up-square})，其中 $X$ 仿射，则图
$$
\xymatrix{
\mathcal{F}(E) & \mathcal{F}(X') \ar[l] \\
\mathcal{F}(Z) \ar[u] & \mathcal{F}(X) \ar[u] \ar[l]
}
$$
在集合范畴中为笛卡儿图。
\end{enumerate}
\end{proposition}

\begin{proof}
假设 $\mathcal{F}$ 是层。由于 Zariski 覆盖是 h 覆盖，
条件 (1) 成立；参见引理 \ref{flat-lemma-zariski-h}。
对 (2) 中的 $f$，$\{X \to Y\}$ 显然是 fppf 覆盖，
因而是 h 覆盖；参见引理 \ref{flat-lemma-zariski-h}。
故条件 (2) 成立。条件 (3) 由引理
\ref{flat-lemma-blow-up-square-h} 成立。

\medskip\noindent
反之，假设 $\mathcal{F}$ 满足 (1)、(2)、(3)。
我们将应用引理 \ref{flat-lemma-characterize-sheaf-h}
证明 $\mathcal{F}$ 是层。考虑 $(\Sch/S)_h$ 中满、有限表示且适当的态射
$f : X \to Y$，其中 $Y$ 仿射。只需证明
$\mathcal{F}(Y)$ 是两个映射
$\mathcal{F}(X) \to \mathcal{F}(X \times_Y X)$ 的等化子。

\medskip\noindent
先进一步假设 $f : X \to Y$ 是闭浸入
（换言之，$f$ 是增厚）。则 $Y$ 沿 $X$ 的爆破为空概形，
从而产生一个顶点为 $\emptyset, \emptyset, X, Y$ 的几乎爆破方块
% P05-EN-CH38-0078: frozen source says “consisting with”; candidate “consisting of”.
（比较引理 \ref{flat-lemma-thickening-h} 的第二种证明）。
因此条件 (3) 告诉我们图
$$
\xymatrix{
\mathcal{F}(\emptyset) & \mathcal{F}(\emptyset) \ar[l] \\
\mathcal{F}(X) \ar[u] & \mathcal{F}(Y) \ar[u] \ar[l]
}
$$
在集合范畴中为笛卡儿图。由于 $\mathcal{F}$ 是 Zariski 拓扑下的层，
$\mathcal{F}(\emptyset)$ 是单点集。因此
$\mathcal{F}(X) = \mathcal{F}(Y)$。

\medskip\noindent
插曲 A：设 $T \to T'$ 是 $(\Sch/S)_h$ 中有限表示的增厚。
则 $\mathcal{F}(T') \to \mathcal{F}(T)$ 为双射。
事实上，选取仿射开覆盖 $T' = \bigcup T'_i$，并令
$T_i = T \times_{T'} T'_i$ 为 $T$ 中相应的仿射开集。
由上一段的结论，对所有 $i$，映射
$\mathcal{F}(T'_i) \to \mathcal{F}(T_i)$ 都为双射。
利用 Zariski 层性质，可见
$\mathcal{F}(T') \to \mathcal{F}(T)$ 为单射。
重复该论证可知它是双射。略去少量细节。

\medskip\noindent
插曲 B：考虑范畴 $(\Sch/S)_h$ 中的几乎爆破方块
(\ref{flat-equation-almost-blow-up-square})。我们断言图
$$
\xymatrix{
\mathcal{F}(E) & \mathcal{F}(X') \ar[l] \\
\mathcal{F}(Z) \ar[u] & \mathcal{F}(X) \ar[u] \ar[l]
}
$$
在集合范畴中为笛卡儿图。选取 $X$ 的仿射开覆盖，
并仿照插曲 A 论证，即可由条件 (3) 推出此事。细节从略。

\medskip\noindent
接下来，设 $f : X \to Y$ 是 $(\Sch/S)_h$ 中满、有限表示且适当的态射，
其中 $Y$ 仿射。选取一般平坦性分层
$$
Y \supset Y_0 \supset Y_1 \supset \ldots \supset Y_t = \emptyset
$$
它是 《态射进阶》，引理
\ref{more-morphisms-lemma-generic-flatness-stratification-scheme}
中为 $f : X \to Y$ 给出的分层。下文将不再说明地使用该分层的全部性质。
令 $X_0 = X \times_Y Y_0$。由插曲 B，有
$\mathcal{F}(Y_0) = \mathcal{F}(Y)$、
$\mathcal{F}(X_0) = \mathcal{F}(X)$，且
$\mathcal{F}(X_0 \times_{Y_0} X_0) = \mathcal{F}(X \times_Y X)$.

\medskip\noindent
我们对 $t$ 归纳证明结论。若 $t = 1$，则
$X_0 \to Y_0$ 满、适当、平坦且为有限表示，故由性质 (2) 得到结论。
若 $t > 1$，则可用 $Y_0$ 替换 $Y$，并用 $X_0$ 替换 $X$
（见上文），从而假设 $Y = Y_0$。

\medskip\noindent
考虑拟紧开子概形
$V = Y \setminus Y_1 = Y_0 \setminus Y_1$.
为 $f : X \to Y \supset V$ 选取如引理
\ref{flat-lemma-flat-after-almost-blowing-up} 中的图
$$
\xymatrix{
E \ar[ddd] \ar[rd] & & & D \ar[lll] \ar[ddd] \ar[ld] \\
& Y' \ar[d] & X' \ar[l] \ar[d] \\
& Y & X \ar[l] \\
Z \ar[ru] & & & T \ar[lll] \ar[lu]
}
$$
于是 $f' : X' \to Y'$ 平坦且为有限表示。$f'$ 也适当
（使用《态射》，引理
\ref{morphisms-lemma-composition-proper} 和
\ref{morphisms-lemma-image-proper-scheme-closed}）。
因此，像 $W = f'(X') \subset Y'$ 是 $Y'$ 的开闭子概形
（对开性使用《态射》，引理
\ref{morphisms-lemma-fppf-open}）。注意，
$Y' \setminus E$ 包含于 $W$。由引理
\ref{flat-lemma-shrink-almost-blow-up}，这意味着可在上图中用
$W$ 替换 $Y'$。换言之，可以并且确实假设 $f'$ 为满射。
此时由插曲 B 可知
$$
\vcenter{
\xymatrix{
\mathcal{F}(E) & \mathcal{F}(Y') \ar[l] \\
\mathcal{F}(Z) \ar[u] & \mathcal{F}(Y) \ar[u] \ar[l]
}
}
\quad\text{且}\quad
\vcenter{
\xymatrix{
\mathcal{F}(D) & \mathcal{F}(X') \ar[l] \\
\mathcal{F}(T) \ar[u] & \mathcal{F}(X) \ar[u] \ar[l]
}
}
$$
中的两个方块均为笛卡儿图。注意，
$Z \cap Y_1 \to Z$ 是有限表示的增厚
（因为 $Z$ 作为 $Y$ 中与 $V$ 不相交的闭子概形，
在集合论意义下包含于 $Y_1$）。因此得到滤过
$$
% P05-EN-CH38-0079: frozen source reverses the last inclusions; a decreasing stratification requires \supset throughout.
Z \supset Z \cap Y_1 \supset Z \cap Y_2 \subset \ldots \subset
Z \cap Y_t = \emptyset
$$
该滤过用于 $f$ 在 $T$ 上的限制
$T = Z \times_Y X \to Z$。故由归纳假设，
$\mathcal{F}(Z) \to \mathcal{F}(T)$
是集合的单射，其像为两个映射
$\mathcal{F}(T) \to \mathcal{F}(T \times_Z T)$ 的等化子。

\medskip\noindent
设 $s \in \mathcal{F}(X)$ 属于两个映射
$\mathcal{F}(X) \to \mathcal{F}(X \times_Y X)$ 的等化子。
由上文，限制 $s|_T$ 来自唯一元素
$t \in \mathcal{F}(Z)$；同理，限制 $s|_{X'}$ 来自唯一元素
$t' \in \mathcal{F}(Y')$。沿上面巨大交换图中箭头所对应的
$\mathcal{F}$ 限制映射追踪截面，可知 $t$ 与 $t'$
在 $\mathcal{F}(E)$ 中限制为同一元素：它们在
$\mathcal{F}(D)$ 中限制为同一元素，并且有 (2)；
这里使用了 $D \to E$ 作为 $X' \to Y'$ 的限制，是满、平坦、适当
且有限表示的态射。因此，由上面两个笛卡儿方块中的第一个，
得到唯一截面
$u \in \mathcal{F}(Y)$
，它在 $Z$ 和 $Y'$ 上分别限制为 $t$ 和 $t'$。
% P05-EN-CH38-0080: frozen source has “u restrict”; candidate “u restricts”.
要验证 $u$ 在 $X$ 上限制为 $s$，使用第二个图即可。
\end{proof}

\begin{example}
\label{flat-example-one-generator}
设 $A$ 是环，$f \in A$ 是元素，$J \subset A$
是被 $f$ 的某个幂零化的有限生成理想。则
$$
\xymatrix{
E = \Spec(A/fA + J) \ar[r] \ar[d] & \Spec(A/J) = X' \ar[d] \\
Z = \Spec(A/fA) \ar[r] & \Spec(A) = X
}
$$
是几乎爆破方块。
\end{example}

\begin{example}
\label{flat-example-two-generators}
设 $A$ 是环，$f_1, f_2 \in A$ 是元素。
$$
\xymatrix{
E = \text{Proj}(A/(f_1, f_2)[T_0, T_1]) \ar[r] \ar[d] &
% P05-EN-CH38-0081: frozen source appears to omit the closing parenthesis of \text{Proj}(...) before = X'.
\text{Proj}(A[T_0, T_1]/(f_2 T_0 - f_1 T_1) = X' \ar[d] \\
Z = \Spec(A/(f_1, f_2)) \ar[r] & \Spec(A) = X
}
$$
是几乎爆破方块。
\end{example}

\begin{lemma}
\label{flat-lemma-refine-check-h}
设 $\mathcal{F}$ 是定义 \ref{flat-definition-big-small-h}
中构造的某个站点 $(\Sch/S)_h$ 上的预层。
则 $\mathcal{F}$ 是层，当且仅当满足下列条件：
\begin{enumerate}
\item $\mathcal{F}$ 是 Zariski 拓扑下的层；
\item 给定 $(\Sch/S)_h$ 中的态射 $f : X \to Y$，
其中 $Y$ 仿射，而 $f$ 满、平坦、适当且为有限表示，
则 $\mathcal{F}(Y)$ 是两个映射
$\mathcal{F}(X) \to \mathcal{F}(X \times_Y X)$ 的等化子；
\item $\mathcal{F}$ 将范畴 $(\Sch/S)_h$ 中如例
\ref{flat-example-one-generator} 的几乎爆破方块变为集合的笛卡儿图；
\item $\mathcal{F}$ 将范畴 $(\Sch/S)_h$ 中如例
\ref{flat-example-two-generators} 的几乎爆破方块变为集合的笛卡儿图。
\end{enumerate}
\end{lemma}

\begin{proof}
由命题 \ref{flat-proposition-check-h}，只需证明：给定范畴
$(\Sch/S)_h$ 中的几乎爆破方块
(\ref{flat-equation-almost-blow-up-square})，其中 $X$ 仿射，则图
$$
\xymatrix{
\mathcal{F}(E) & \mathcal{F}(X') \ar[l] \\
\mathcal{F}(Z) \ar[u] & \mathcal{F}(X) \ar[u] \ar[l]
}
$$
在集合范畴中为笛卡儿图。证明的大致思路是用其他几乎爆破方块
支配该态射，以便局部应用假设 (3) 和 (4)。

\medskip\noindent
假设有范畴 $(\Sch/S)_h$ 中的几乎爆破方块
(\ref{flat-equation-almost-blow-up-square})、开覆盖
$X = \bigcup U_i$，以及开覆盖
$U_i \cap U_j = \bigcup U_{ijk}$，使图
$$
\vcenter{
\xymatrix{
\mathcal{F}(E \cap b^{-1}(U_i)) & \mathcal{F}(b^{-1}(U_i)) \ar[l] \\
\mathcal{F}(Z \cap U_i) \ar[u] & \mathcal{F}(U_i) \ar[u] \ar[l]
}
}
\quad\text{且}\quad
\vcenter{
\xymatrix{
\mathcal{F}(E \cap b^{-1}(U_{ijk})) &
\mathcal{F}(b^{-1}(U_{ijk})) \ar[l] \\
\mathcal{F}(Z \cap U_{ijk}) \ar[u] &
\mathcal{F}(U_{ijk}) \ar[u] \ar[l]
}
}
$$
均为笛卡儿图。则图
$$
\xymatrix{
\mathcal{F}(E) & \mathcal{F}(X') \ar[l] \\
\mathcal{F}(Z) \ar[u] & \mathcal{F}(X) \ar[u] \ar[l]
}
$$
也为笛卡儿图。这是因为 $\mathcal{F}$ 是 Zariski 拓扑下的层。

\medskip\noindent
特别地，若有爆破方块
% P05-EN-CH38-0082: frozen source calls equation-almost-blow-up-square a “blow up square”; candidate “almost blow up square”.
(\ref{flat-equation-almost-blow-up-square})，其中 $b : X' \to X$
为闭浸入，而 $Z$ 是局部主闭子概形，则
$\mathcal{F}(X) = \mathcal{F}(X') \times_{\mathcal{F}(E)} \mathcal{F}(Z)$.
事实上，在 $X$ 上仿射局部地，可得到如 (3) 中的几乎爆破方块。

\medskip\noindent
令 $Z \subset X$、$E_k \subset X'_k \to X$、
$E \subset X' \to X$ 以及 $i_k : X' \to X'_k$
如引理 \ref{flat-lemma-almost-blow-up-unique} 的陈述中所述。则
$$
\xymatrix{
E \ar[d] \ar[r] & X' \ar[d] \\
E_k \ar[r] & X'_k
}
$$
是上一段所讨论类型的几乎爆破方块。因此
$$
\mathcal{F}(X'_k) = \mathcal{F}(X') \times_{\mathcal{F}(E)} \mathcal{F}(E_k)
$$
对 $k = 1, 2$ 成立。由此可知
$$
\mathcal{F}(X) \longrightarrow
\mathcal{F}(X'_k) \times_{\mathcal{F}(E_k)} \mathcal{F}(Z)
$$
对 $k = 1$ 为双射，当且仅当它对 $k = 2$ 为双射。
因此，给定有限表示的闭浸入 $Z \to X$，其中 $X$ 拟紧且拟分离，
等式
$\mathcal{F}(X) = \mathcal{F}(X') \times_{\mathcal{F}(E)} \mathcal{F}(Z)$
是否成立，与所选几乎爆破方块
(\ref{flat-equation-almost-blow-up-square}) 无关。
（而且由引理 \ref{flat-lemma-almost-blow-up-square}，
$Z \subset X$ 的几乎爆破方块确实存在。）

\medskip\noindent
最后，考虑 $(\Sch/S)_h$ 的仿射对象 $X$ 和有限表示闭浸入
$Z \to X$。对定义 $X$ 中 $Z$ 的理想的生成元数目 $r$ 归纳，
证明概形对 $(X, Z)$ 具有所需性质。若生成元数目为 $\leq 2$，
则可选取例 \ref{flat-example-two-generators} 中的几乎爆破方块，
并由假设 (4) 得到结论。

\medskip\noindent
归纳步骤。设 $X = \Spec(A)$ 且
$Z = \Spec(A/(f_1, \ldots, f_r))$，其中 $r > 2$。
为概形对 $(X, Z)$ 选取爆破方块
(\ref{flat-equation-almost-blow-up-square})。令
$Z_1 = \Spec(A/(f_1, f_2))$，并令
$$
\xymatrix{
E_1 \ar[d] \ar[r] & Y \ar[d] \\
Z_1 \ar[r] & X
}
$$
为例 \ref{flat-example-two-generators} 中构造的几乎爆破方块。
由引理 \ref{flat-lemma-base-change-almost-blow-up}，基变换
$$
(I)
\vcenter{
\xymatrix{
Y \times_X E \ar[r] \ar[d] & Y \times_X X' \ar[d] \\
Y \times_X Z \ar[r] & Y
}
}
\quad\text{且}\quad
(II)
\vcenter{
\xymatrix{
E \ar[r] \ar[d] & Z_1 \times_X X' \ar[d] \\
Z \ar[r] & Z_1
}
}
$$
都是几乎爆破方块。$Z_1$ 中 $Z$ 的理想由 $r - 2$ 个元素生成。
$Y \times_X Z$ 的理想由 $f_1, \ldots, f_r$ 到 $Y$ 的拉回生成。
在 $Y$ 上局部地，由 $f_1, f_2$ 生成的理想可由一个元素生成；
故 $Y \times_X Z$ 在 $Y$ 上仿射局部地由至多
$r - 1$ 个元素切出。由归纳假设和上面的讨论，
$$
\mathcal{F}(Y) =
\mathcal{F}(Y \times_X X') \times_{\mathcal{F}(Y \times_X E)}
\mathcal{F}(Y \times_X Z)
$$
并且
$$
\mathcal{F}(Z_1) =
\mathcal{F}(Z_1 \times_X X') \times_{\mathcal{F}(E)}
\mathcal{F}(Z)
$$
由假设 (4)，有
$$
\mathcal{F}(X) =
\mathcal{F}(Y) \times_{\mathcal{F}(E_1)} \mathcal{F}(Z_1)
$$
现设给定一对 $(s', t)$，其中 $s' \in \mathcal{F}(X')$、
$t \in \mathcal{F}(Z)$，且二者在 $\mathcal{F}(E)$ 中的限制相同。
则
% P05-EN-CH38-0083: frozen source writes s'|{...} rather than restriction syntax s'|_{...}.
$(s'|{Z_1 \times_X X'}, t)$ 是某个唯一元素
$t_1 \in \mathcal{F}(Z_1)$ 的像。同理，
$(s'|_{Y \times_X X'}, t|_{Y \times_X Z})$ 是某个唯一元素
$s_Y \in \mathcal{F}(Y)$ 的像。我们断言 $s_Y$ 与 $t_1$
在 $\mathcal{F}(E_1)$ 中限制为同一元素。这是因为几乎爆破方块
$$
\xymatrix{
E_1 \times_X E \ar[r] \ar[d] & E_1 \times_X X' \ar[d] \\
E_1 \times_X Z \ar[r] & E_1
}
$$
既是 (I) 沿 $E_1 \to Y$ 的基变换，又是 (II) 沿
$E_1 \to Z_1$ 的基变换；并且用于构造 $s_Y$ 与 $t_1$
的截面对彼此相符。因此，由第三个纤维积等式，存在唯一的
$s \in \mathcal{F}(X)$，它在 $\mathcal{F}(Y)$ 中映到 $s_Y$，
并在
% P05-EN-CH38-0084: frozen source says t_1 is reached in F(Z), but t_1 was constructed in F(Z_1) and the fibre product uses F(Z_1).
$\mathcal{F}(Z)$ 中映到 $t_1$。验证 $s$ 在
$\mathcal{F}(X')$ 中映到 $s'$，并在 $\mathcal{F}(Z)$ 中映到
$t$ 的过程从略；提示：使用刚构造的 $s$ 的唯一性，并在仿射局部工作。
\end{proof}

\begin{lemma}
\label{flat-lemma-refine-check-h-stack}
设 $p : \mathcal{S} \to (\Sch/S)_h$ 是群胚纤维化范畴。
则 $\mathcal{S}$ 是群胚叠，当且仅当满足下列条件：
\begin{enumerate}
\item $\mathcal{S}$ 是 Zariski 拓扑下的群胚叠；
\item 给定 $(\Sch/S)_h$ 中的态射 $f : X \to Y$，
其中 $Y$ 仿射，而 $f$ 满、平坦、适当且为有限表示，则
$$
\mathcal{S}_Y \longrightarrow
\mathcal{S}_X \times_{\mathcal{S}_{X \times_Y X}} \mathcal{S}_X
$$
是范畴等价；
\item 对范畴 $(\Sch/S)_h$ 中如例 \ref{flat-example-one-generator}
或 \ref{flat-example-two-generators} 的几乎爆破方块，函子
$$
\mathcal{S}_X \longrightarrow
\mathcal{S}_Z \times_{\mathcal{S}_E} \mathcal{S}_{X'}
$$
是范畴等价。
\end{enumerate}
\end{lemma}

\begin{proof}
本引理形式地由引理 \ref{flat-lemma-refine-check-h}
和我们对群胚叠的定义推出。例如，假设 (1)、(2)、(3)。
为证明 $\mathcal{S}$ 是叠，需证明态射和对象的下降；
参见《叠》，定义
\ref{stacks-definition-stack-in-groupoids}。

\medskip\noindent
若 $x, y$ 是 $\mathcal{S}$ 中位于 $(\Sch/S)_h$ 的对象 $U$
之上的对象，则我们的假设蕴含 $\mathit{Isom}(x, y)$
是 $(\Sch/U)_h$ 上满足引理 \ref{flat-lemma-refine-check-h}
中 (1)、(2)、(3)、(4) 的预层，因而是层。略去一些细节。

\medskip\noindent
设 $\{U_i \to U\}_{i \in I}$ 是 $(\Sch/S)_h$ 的覆盖，
$(x_i, \varphi_{ij})$ 是 $\mathcal{S}$ 中相对于态射族
$\{U_i \to U\}_{i \in I}$ 的下降资料；参见《叠》，定义
\ref{stacks-definition-descent-data}。考虑规则 $F$：
它把 $(\Sch/U)_h$ 中的 $V/U$ 映到由对 $(y, \psi_i)$ 组成的集合，
其中 $y$ 是 $\mathcal{S}_V$ 的对象，而
$\psi_i : y|_{U_i \times_U V} \to x_i|_{U_i \times_U V}$
是 $\mathcal{S}$ 中位于 $U_i \times_U V$ 之上的态射，并满足
$$
\varphi_{ij}|_{U_i \times_U U_j \times_U V}
\circ
\psi_i|_{U_i \times_U U_j \times_U V}
=
\psi_j|_{U_i \times_U U_j \times_U V}
$$
在同构意义下成立。由于已经有态射的下降，显然
$F(V/U)$ 或为空集，或为单点集。另一方面，
$F(U_{i_0}/U)$ 非空，因为它包含
$(x_{i_0}, \varphi_{i_0i})$。我们的目标是证明
$F(U/U)$ 非空，因此只需说明 $F$ 是 $(\Sch/U)_h$ 上的层。
可使用引理 \ref{flat-lemma-refine-check-h} 的判据。
而我们的假设 (1)、(2)、(3) 蕴含（画出一些略去的交换图即可）
引理 \ref{flat-lemma-refine-check-h} 的性质 (1)、(2)、(3)、(4)
对 $F$ 成立。

\medskip\noindent
若 $\mathcal{S}$ 是群胚叠，则 (1)、(2)、(3) 成立的验证从略。
\end{proof}





\section{绝对弱正规化与 h 覆盖}
\label{flat-section-h-and-weak-normalization}

\noindent
本节使用第 \ref{flat-section-blow-up-h} 节中的判据给出若干 h 层，
并将结构层的 h 层化与绝对弱正规化联系起来。
为此需要下述初等引理。

\begin{lemma}
\label{flat-lemma-funny-blow-up}
设 $Z, X, X', E$ 构成如例 \ref{flat-example-two-generators}
中的几乎爆破方块。则对 $p > 0$，
$H^p(X', \mathcal{O}_{X'}) = 0$，并且
$\Gamma(X, \mathcal{O}_X) \to \Gamma(X', \mathcal{O}_{X'})$
是环的满射，其核为平方零理想。
\end{lemma}

\begin{proof}
先假设 $A = \mathbf{Z}[f_1, f_2]$ 为多项式环。
此时，我们的几乎爆破方块就是 $X = \Spec(A)$ 沿闭子概形
$Z$ 的爆破；事实上，$X' \subset \mathbf{P}^1_X$
是由 $\mathcal{O}_{\mathbf{P}^1_X}(1)$ 的整体截面
$f_2T_0 - f_1 T_1$ 切出的有效 Cartier 除子。因此有消解
$$
0 \to
\mathcal{O}_{\mathbf{P}^1_X}(-1) \to
\mathcal{O}_{\mathbf{P}^1_X} \to
\mathcal{O}_{X'} \to
0
$$
使用《概形上同调》，第 \ref{coherent-section-cohomology-projective-space} 节
给出的上同调描述，可知在此情形下
$\Gamma(X, \mathcal{O}_X) \to \Gamma(X', \mathcal{O}_{X'})$
为同构，且 $H^1(X', \mathcal{O}_{X'}) = 0$。

\medskip\noindent
接着注意，如例 \ref{flat-example-two-generators} 中的任意图，
都是上一段之图沿环映射 $\mathbf{Z}[f_1, f_2] \to A$ 的基变换。
故由《态射进阶》，引理
\ref{more-morphisms-lemma-check-h1-fibre-zero},
\ref{more-morphisms-lemma-h1-fibre-zero} 和
\ref{more-morphisms-lemma-h1-fibre-zero-check-h0-kappa}
，可知一般地 $H^1(X', \mathcal{O}_{X'})$ 为零，并且映射
$H^0(X, \mathcal{O}_X) \to H^0(X', \mathcal{O}_{X'})$ 为满射。
% P05-EN-CH38-0085: frozen source ends this inference with the fragment “and the surjectivity ... in general”.

\medskip\noindent
接下来在一般情形中研究其核。若 $a \in A$ 映到零，
则在仿射坐标片上考察可得
$$
a = (f_1x - f_2)(a_0 + a_1x + \ldots + a_rx^r)\text{ 于 }A[x]\text{ 中}
$$
其中 $r \geq 0$ 且 $a_0, \ldots, a_r \in A$；同理
$$
a = (f_1 - f_2y)(b_0 + b_1y + \ldots + b_s y^s)\text{ 于 }A[y]\text{ 中}
$$
其中 $s \geq 0$ 且 $b_0, \ldots, b_s \in A$。这意味着
$$
a = f_2 a_0,\ f_1 a_0 = f_2 a_1,\ \ldots,\ f_1 a_r = 0,
\ a = f_1 b_0,\ f_2 b_0 = f_1 b_1,\ \ldots,\ f_2 b_s = 0
$$
若 $(a', r', a'_i, s', b'_j)$ 是另一组这样的数据，则
$$
aa' = f_1f_2a_0b'_0 = f_1f_2a_1b'_1 = f_1f_2a_2b'_2 = \ldots = 0
$$
这正合所需。
\end{proof}

\noindent
对 $\mathbf{F}_p$-代数 $A$，定义 $\colim_F A$ 为系统
$$
A \xrightarrow{F} A \xrightarrow{F} A \xrightarrow{F} \ldots
$$
的余极限，其中 $F : A \to A$、$a \mapsto a^p$
是 Frobenius 自同态。

\begin{lemma}
\label{flat-lemma-h-sheaf-colim-F}
设 $p$ 是素数，$S$ 是 $\mathbf{F}_p$ 上概形，
$(\Sch/S)_h$ 是定义 \ref{flat-definition-big-small-h} 中的站点。
则 $(\Sch/S)_h$ 上存在唯一的层 $\mathcal{F}$，使对
$(\Sch/S)_h$ 的任意拟紧且拟分离对象 $X$，
$$
\mathcal{F}(X) = \colim_F \Gamma(X, \mathcal{O}_X)
$$
\end{lemma}

\begin{proof}
将函子
$$
X \longrightarrow \colim_F \Gamma(X, \mathcal{O}_X)
$$
的 Zariski 层化记作 $\mathcal{F}$。
对于拟紧且拟分离概形 $X$，由《层》，引理
\ref{sheaves-lemma-directed-colimits-sections}
以及 $\mathcal{O}$ 是 Zariski 拓扑的层这一事实，有
$\mathcal{F}(X) = \colim_F \Gamma(X, \mathcal{O}_X)$。
% P05-EN-CH38-0086: source 11871--11874 has a full stop before lowercase ``by''.
因此，只需证明 $\mathcal{F}$ 是 h 层。为此，我们核验
引理 \ref{flat-lemma-refine-check-h} 的条件 (1)、(2)、(3) 和 (4)。
条件 (1) 成立，因为我们作了（几乎不必要的）Zariski 层化。
条件 (2) 成立，因为 $\mathcal{O}$ 是 fppf 层（《下降》，引理
\ref{descent-lemma-sheaf-condition-holds}），而且若 $A$ 是两个
$\mathbf{F}_p$-代数同态 $B \to C$ 的等化子，则 $\colim_F A$
是两个映射 $\colim_F B \to \colim_F C$ 的等化子。

\medskip\noindent
我们核验条件 (3)。令 $A, f, J$ 如
例 \ref{flat-example-one-generator} 中所述。我们必须证明
$$
\colim_F A = \colim_F A/J \times_{\colim_F A/fA + J} \colim_F A/fA
$$
这归结为如下代数问题：设 $a', a'' \in A$ 满足
$F^n(a' - a'') \in fA + J$。求出 $a \in A$ 和 $m \geq 0$，使得
$a - F^m(a') \in J$ 且 $a - F^m(a'') \in fA$，并证明序对
$(a, m)$ 在形如 $(a, m) \mapsto (F(a), m + 1)$ 的替换之下唯一确定。
为此，只须将 $F^n(a' - a'') = f h + g$ 写成 $h \in A$、$g \in J$
的形式，并令 $a = F^n(a') - g = F^n(a'') + fh$、$m = n$。
为证明唯一性，设 $(a_1, m_1)$ 是第二个解。作上述形式的替换后，
可以假设 $m = m_1$。于是 $a - a_1 \in J$ 且 $a - a_1 \in fA$。
由于 $J$ 被 $f$ 的某个幂零化，$a - a_1$ 是幂零元。因此对充分大的
$k$，$F^k(a - a_1)$ 为零。再作若干次替换后，便得到 $a = a_1$。

\medskip\noindent
我们核验条件 (4)。令 $X, X', Z, E$ 如
例 \ref{flat-example-two-generators} 中所述。由
引理 \ref{flat-lemma-funny-blow-up} 可知
$$
\mathcal{F}(X) = \colim_F \Gamma(X, \mathcal{O}_X)
\longrightarrow
\colim_F \Gamma(X', \mathcal{O}_{X'}) = \mathcal{F}(X')
$$
是双射。在此情形下 $E = \mathbf{P}^1_Z$，故
$\mathcal{F}(Z) \to \mathcal{F}(E)$ 也是双射。因此结论在此情形下同样成立。
\end{proof}

\noindent
设 $p$ 是素数。对于 $\mathbf{F}_p$-代数 $A$，令 $\lim_F A$
等于逆系统
$$
\ldots \xrightarrow{F} A \xrightarrow{F} A \xrightarrow{F} A
$$
的极限，其中 $F : A \to A$、$a \mapsto a^p$ 是 Frobenius 自同态。

\begin{lemma}
\label{flat-lemma-h-sheaf-lim-F}
设 $p$ 是素数，$S$ 是 $\mathbf{F}_p$ 上概形，
$(\Sch/S)_h$ 是定义 \ref{flat-definition-big-small-h} 中的站点。规则
$$
\mathcal{F}(X) = \lim_F \Gamma(X, \mathcal{O}_X)
$$
定义了 $(\Sch/S)_h$ 上的一个层。
\end{lemma}

\begin{proof}
为证明 $\mathcal{F}$ 是层，我们核验 引理 \ref{flat-lemma-refine-check-h}
的条件 (1)、(2)、(3) 和 (4)。条件 (1) 成立，因为层的极限仍是层，
且 $\mathcal{O}$ 是 Zariski 层。条件 (2) 成立，因为 $\mathcal{O}$
是 fppf 层（《下降》，引理
\ref{descent-lemma-sheaf-condition-holds}），而且若 $A$ 是两个
$\mathbf{F}_p$-代数同态 $B \to C$ 的等化子，则 $\lim_F A$
是两个映射 $\lim_F B \to \lim_F C$ 的等化子。

\medskip\noindent
我们核验条件 (3)。令 $A, f, J$ 如
例 \ref{flat-example-one-generator} 中所述。我们必须证明
\begin{align*}
\lim_F A
& \to
\lim_F A/J \times_{\lim_F A/fA + J} \lim_F A/fA \\
& =
\lim_F (A/J \times_{A/fA + J} A/fA) \\
& =
\lim_F A/(fA \cap J)
\end{align*}
是双射。由于 $J$ 被 $f$ 的某个幂零化，
$\mathfrak a = fA \cap J$ 是幂零理想，即存在 $n$ 使得
$\mathfrak a^n = 0$。不难验证，在此情形下
$\lim_F A \to \lim_F A/\mathfrak a$ 是双射。

\medskip\noindent
我们核验条件 (4)。令 $X, X', Z, E$ 如
例 \ref{flat-example-two-generators} 中所述。由
引理 \ref{flat-lemma-funny-blow-up} 和上述同一论证可知
$$
\mathcal{F}(X) = \lim_F \Gamma(X, \mathcal{O}_X)
\longrightarrow
\lim_F \Gamma(X', \mathcal{O}_{X'}) = \mathcal{F}(X')
$$
是双射。在此情形下 $E = \mathbf{P}^1_Z$，故
$\mathcal{F}(Z) \to \mathcal{F}(E)$ 也是双射。因此结论在此情形下同样成立。
\end{proof}

\noindent
在下一个引理中，我们使用概形 $X$ 的绝对弱正规化 $X^{awn}$；见
《态射》，第 \ref{morphisms-section-seminormalization} 节。

\begin{lemma}
\label{flat-lemma-weak-normalization-ph-sheaf}
设 $(\Sch/S)_{ph}$ 是 《拓扑》，定义
\ref{topologies-definition-big-small-ph} 中的站点。规则
$$
X \longmapsto \Gamma(X^{awn}, \mathcal{O}_{X^{awn}})
$$
是 $(\Sch/S)_{ph}$ 上的层。
\end{lemma}

\begin{proof}
为证明 $\mathcal{F}$ 是层，我们核验 《拓扑》，引理
\ref{topologies-lemma-characterize-sheaf} 的条件 (1) 和 (2)。
条件 (1) 成立，因为 $X^{awn}$ 的形成与开覆盖可交换；见《态射》，
引理 \ref{morphisms-lemma-seminormalization} 及其证明。

\medskip\noindent
设 $\pi : Y \to X$ 是满的固有态射。我们必须证明，两个映射
$$
\Gamma(Y^{awn}, \mathcal{O}_{Y^{awn}}) \to
\Gamma((Y \times_X Y)^{awn}, \mathcal{O}_{(Y \times_X Y)^{awn}})
$$
的等化子等于 $\Gamma(X^{awn}, \mathcal{O}_{X^{awn}})$。令 $f$
是该等化子的元素。考虑态射
$$
f : Y^{awn} \longrightarrow \mathbf{A}^1_X
$$
由于 $Y^{awn} \to X$ 是泛闭的，$f$ 的概形论像 $Z$ 是
$\mathbf{A}^1_X$ 的一个在 $X$ 上固有的闭子概形，并且
$f : Y^{awn} \to Z$ 是满态射；见《态射》，引理
\ref{morphisms-lemma-scheme-theoretic-image-is-proper}。因此 $Z \to X$
是有限满态射（《态射》，引理 \ref{morphisms-lemma-finite-proper}）。

\medskip\noindent
设 $k$ 是域，$z_1, z_2 : \Spec(k) \to Z$ 是两个经 $Z \to X$
等化的态射。我们断言 $z_1 = z_2$。只须证明像
$\lambda_i = z_i^*f \in k$ 相等（因为 $Z$ 的结构层在 $X$ 的结构层上
由 $f$ 生成）。为此，选取域扩张 $K/k$ 和态射
$y_1, y_2 : \Spec(K) \to Y^{awn}$，使得
$z_i \circ (\Spec(K) \to \Spec(k)) = f \circ y_i$。由于
$Y^{awn} \to Z$ 是满态射，这是可以做到的。选取一个基数非常大的
代数闭扩张 $\Omega/k$。对任意 $k$-代数同态
$\sigma_i : K \to \Omega$，得到
$$
\Spec(\Omega)
\xrightarrow{\sigma_1, \sigma_2}
\Spec(K \otimes_k K)
\xrightarrow{y_1, y_2}
Y^{awn} \times_X Y^{awn}
$$
由于典范态射
$(Y \times_X Y)^{awn} \to Y^{awn} \times_X Y^{awn}$ 是泛同胚，
且 $\Omega$ 代数闭，上述复合唯一提升为态射
$\Spec(\Omega) \to (Y \times_X Y)^{awn}$。由于 $f$ 属于上述等化子，
这证明 $\sigma_1(\lambda_1) = \sigma_2(\lambda_2)$。关于域扩张的一个
简单引理表明，这蕴含 $\lambda_1 = \lambda_2$；细节从略。

\medskip\noindent
由此可知 $Z \to X$ 是泛单射，即 $Z \to X$ 在点上是单射，并诱导
纯不可分的剩余域扩张（《态射》，引理
\ref{morphisms-lemma-universally-injective}）。
% P05-EN-CH38-0087: source 12060 says ``purely inseparated residue field extensions''.
综上，$Z \to X$ 是泛同胚；见《态射》，引理
\ref{morphisms-lemma-universal-homeomorphism}。

\medskip\noindent
令 $g : X^{awn} \to Z$ 是由 $X^{awn}$ 的泛性质得到的映射。于是
$Y^{awn} \to X^{awn} \to Z$ 与 $f : Y^{awn} \to Z$ 是两个 $X$-态射。
由 $Y^{awn} \to Y$ 的泛性质，相应的两个 $Y$-态射
$Y^{awn} \to Y \times_X Z$ 必须相等。这蕴含作为到
$\mathbf{A}^1_X$ 的态射有 $g \circ \pi^{wan} = f$，从而
$g \in \Gamma(X^{awn}, \mathcal{O}_{X^{awn}})$ 正是所求元素。
% P05-EN-CH38-0088: source 12071 uses \pi^{wan}, unlike the surrounding awn notation.
\end{proof}

\begin{lemma}
\label{flat-lemma-weak-normalization-h-sheaf}
设 $S$ 是概形。选取定义 \ref{flat-definition-big-small-h} 中的站点
$(\Sch/S)_h$。规则
$$
X \longmapsto \Gamma(X^{awn}, \mathcal{O}_{X^{awn}})
$$
是 $(\Sch/S)_h$ 上“结构层”$\mathcal{O}$ 的层化。对 ph 拓扑亦然。
\end{lemma}

\begin{proof}
在 引理 \ref{flat-lemma-weak-normalization-ph-sheaf} 中，我们已经看到
本引理的规则 $\mathcal{F}$ 定义了 ph 拓扑的层，因而更是 h 拓扑的层。
显然，存在典范的环预层态射 $\mathcal{O} \to \mathcal{F}$。
为完成证明，只须证明
\begin{enumerate}
\item 若 $f \in \mathcal{O}(X)$ 在 $\mathcal{F}(X)$ 中的像为零，
则存在 h 覆盖 $\{X_i \to X\}$，使得 $f|_{X_i} = 0$；
\item 给定 $f \in \mathcal{F}(X)$，存在 h 覆盖 $\{X_i \to X\}$，
使得 $f|_{X_i}$ 是某个 $f_i \in \mathcal{O}(X_i)$ 的像。
\end{enumerate}
令 $f$ 如 (1) 中所述。则 $f|_{X^{awn}} = 0$，这意味着 $f$ 局部幂零。
因此，若 $X' \subset X$ 是由 $f$ 截出的闭子概形，则 $X' \to X$
是有限表现的满闭浸入。故 $\{X' \to X\}$ 是所需的 h 覆盖。
令 $f$ 如 (2) 中所述。用一个仿射开覆盖的成员替换 $X$ 后，可以假设
$X = \Spec(A)$ 是仿射的。于是 $f \in A^{awn}$；见
《态射》，引理 \ref{morphisms-lemma-seminormalization-ring}。
由《态射》，引理 \ref{morphisms-lemma-universal-homeo-limit}，
可找到有限表现的环同态 $A \to B$，使得
$\Spec(B) \to \Spec(A)$ 是泛同胚，并且 $f$ 是某个 $b \in B$
在典范映射 $B \to A^{awn}$ 下的像。于是
$\{\Spec(B) \to \Spec(A)\}$ 是 h 覆盖，结论得证。关于 ph 拓扑的
断言可由相同方法推出（也可由 h 拓扑的断言推出）。
\end{proof}

\noindent
设 $p$ 是素数。若映射 $F : A \to A$、$x \mapsto x^p$ 是
$A$ 的自同构，则称 $\mathbf{F}_p$-代数 $A$ 是{\it 完美的}。

\begin{lemma}
\label{flat-lemma-perfect-weankly-normal}
设 $p$ 是素数。$\mathbf{F}_p$-代数 $A$ 绝对弱正规，当且仅当它是完美的。
\end{lemma}

\begin{proof}
由《态射》，定义 \ref{morphisms-definition-seminormal-ring}
中的条件 (2)(b) 立即可知，若 $A$ 绝对弱正规，则它是完美的。

\medskip\noindent
假设 $A$ 是完美的。设 $x, y \in A$ 且 $x^3 = y^2$。
若 $p > 3$，则可对某些 $n, m > 0$ 写成 $p = 2n + 3m$。
选取 $a, b \in A$，使得 $a^p = x$ 且 $b^p = y$。令
$c = a^n b^m$，则
$$
c^{2p} = x^{2n} y^{2m} = x^{2n + 3m} = x^p
$$
从而 $c^2 = x$。类似地，$c^3 = y$。若 $p = 2$，则将
$x = a^2$ 代入得到 $a^6 = y^2$，因而 $a^3 = y$。
若 $p = 3$，则将 $y = a^3$ 代入得到 $x^3 = a^6$，因而 $x = a^2$。

\medskip\noindent
设对某个素数 $\ell$，$x, y \in A$ 满足
$\ell^\ell x = y^\ell$。若 $\ell \not = p$，则 $a = y/\ell$
满足 $a^\ell = x$ 且 $\ell a = y$。若 $\ell = p$，则
$y = 0$，而对某个 $a$ 有 $x = a^p$。
\end{proof}

\begin{lemma}
\label{flat-lemma-char-p}
设 $p$ 是素数。
\begin{enumerate}
\item 若 $A$ 是 $\mathbf{F}_p$-代数，则 $\colim_F A = A^{awn}$。
\item 若 $S$ 是 $\mathbf{F}_p$ 上概形，则 $\mathcal{O}$ 的 h 层化
将拟紧且拟分离的 $X$ 映到 $\colim_F \Gamma(X, \mathcal{O}_X)$。
\end{enumerate}
\end{lemma}

\begin{proof}
证明 (1)。由《代数》，引理 \ref{algebra-lemma-p-ring-map}，
$A \to \colim_F A$ 在谱上诱导泛同胚。因此只须证明
$B = \colim_F A$ 绝对弱正规；见《态射》，引理
\ref{morphisms-lemma-seminormalization-ring}。注意，环同态
$F : B \to B$ 是自同构；换言之，$B$ 是完美环。因此可应用
引理 \ref{flat-lemma-perfect-weankly-normal}。

\medskip\noindent
证明 (2)。在仿射局部考察，由 (1) 以及 引理
\ref{flat-lemma-h-sheaf-colim-F} 和
\ref{flat-lemma-weak-normalization-h-sheaf} 即得。
\end{proof}






\section{正特征下向量丛的下降}
\label{flat-section-descent-vectorbundles-h}

\noindent
本节的参考文献是 \cite{Witt-Grass}。

\medskip\noindent
对概形 $S$，以 $\textit{Vect}(S)$ 表示有限局部自由
$\mathcal{O}_S$-模的范畴。设 $p$ 是素数，$S$ 是
$\mathbf{F}_p$ 上拟紧且拟分离的概形。本节将使用范畴
$$
\colim_F \textit{Vect}(S) =
\colim
\left(
\textit{Vect}(S) \xrightarrow{F^*}
\textit{Vect}(S) \xrightarrow{F^*}
\textit{Vect}(S) \xrightarrow{F^*} \ldots
\right)
$$
其中 $F : S \to S$ 是绝对 Frobenius 态射。具体而言，该范畴的对象是
序对 $(\mathcal{E}, n)$，其中 $\mathcal{E}$ 是有限局部自由
$\mathcal{O}_S$-模，$n \geq 0$ 是整数。其态射取为
% P05-EN-CH38-0089: source 12210--12211 uses unhyphenated “down to earth” attributively.
$$
\Hom_{\colim_F \textit{Vect}(S)}((\mathcal{E}, n), (\mathcal{G}, m)) =
\colim_N \Hom_S(F^{N - n, *}\mathcal{E}, F^{N - m, *}\mathcal{G})
$$
其中 $F : S \to S$ 是 $S$ 的绝对 Frobenius 态射。因此对象
$(\mathcal{E}, n)$ 同构于对象 $(F^*\mathcal{E}, n + 1)$。

\begin{lemma}
\label{flat-lemma-colim-F-Vect}
设 $p$ 是素数，$S$ 是 $\mathbf{F}_p$ 上拟紧且拟分离的概形。
范畴 $\colim_F \textit{Vect}(S)$ 等价于 $S$ 上环层
$\colim_F \mathcal{O}_S$ 的有限局部自由模范畴。
\end{lemma}

\begin{proof}
从略。
\end{proof}

\begin{lemma}
\label{flat-lemma-vector-bundle-I}
设 $p$ 是素数。考虑特征 $p$ 中如
例 \ref{flat-example-one-generator} 的几乎爆破方块 $X, X', Z, E$。
则函子
$$
\colim_F \textit{Vect}(X)
\longrightarrow
\colim_F \textit{Vect}(Z)
\times_{\colim_F \textit{Vect}(E)}
\colim_F \textit{Vect}(X')
$$
是等价。
\end{lemma}

\begin{proof}
令 $A, f, J$ 如 例 \ref{flat-example-one-generator} 中所述。
由于这里的所有概形都是仿射的，并且向量丛范畴中有内 Hom，
只要证明对有限投射 $A$-模 $P$ 有
% P05-EN-CH38-0090: source 12252 says “the fully faithfulness”; standard noun is “full faithfulness”.
$$
\colim P \otimes_{A, F^N} A =
\colim P \otimes_{A, F^N} A/J
\times_{\colim P \otimes_{A, F^N} A/fA + J}
\colim P \otimes_{A, F^N} A/fA
$$
由此即可推出该函子的全忠实性。将 $P$ 写成有限自由模的直和项后，
此式可由 $P$ 有限自由的情形推出，而该情形立即归结为 $P = A$。
$P = A$ 的情形由 引理 \ref{flat-lemma-h-sheaf-colim-F} 得到
（事实上，我们在该引理的证明中直接证明了这一情形）。

\medskip\noindent
本质满性。这里得到如下代数问题。设 $P_1$ 是有限投射 $A/J$-模，
$P_2$ 是有限投射 $A/fA$-模，并且
$$
\varphi :
P_1 \otimes_{A/J} A/fA + J
\longrightarrow
P_2 \otimes_{A/fA} A/fA + J
$$
是同构。目标是证明：存在 $N$、有限投射 $A$-模 $P$、同构
$\varphi_1 : P \otimes_A A/J \to P_1 \otimes_{A/J, F^N} A/J$,
以及同构
$\varphi_2 : P \otimes_A A/fA \to P_2 \otimes_{A/fA, F^N} A/fA$
，并且它们以显然的方式与 $\varphi$ 相容。证明如下。首先注意
$$
A/(J \cap fA) = A/J \times_{A/fA + J} A/fA
$$
因此由《代数进阶》，引理
\ref{more-algebra-lemma-finitely-presented-module-over-fibre-product}，
存在 $A/(J \cap fA)$ 上有限投射模 $P'$，连同同构
$\varphi'_1 : P' \otimes_A A/J \to P_1$ 和
$\varphi_2 : P' \otimes_A A/fA \to P_2$，它们与 $\varphi$ 相容。
% P05-EN-CH38-0091: source 12290--12291 primes only \varphi'_1 although both maps arise from P'.
由于 $J$ 是有限生成的 $f$-幂挠理想，$J \cap fA$ 是幂零理想。
因此对某个 $N$，$F^N$ 有分解
$$
A \xrightarrow{\alpha} A/(J \cap fA) \xrightarrow{\beta} A
$$
令 $P = P' \otimes_{A/(J \cap fA), \beta} A$，结论即得。
\end{proof}

\begin{lemma}
\label{flat-lemma-vector-bundle-II}
设 $p$ 是素数。考虑特征 $p$ 中如
例 \ref{flat-example-two-generators} 的几乎爆破方块 $X, X', Z, E$。
则函子
$$
G :
\colim_F \textit{Vect}(X)
\longrightarrow
\colim_F \textit{Vect}(Z)
\times_{\colim_F \textit{Vect}(E)}
\colim_F \textit{Vect}(X')
$$
是等价。
\end{lemma}

\begin{proof}
全忠实性。设 $(\mathcal{E}, n)$ 和 $(\mathcal{F}, m)$ 是
$\colim_F \textit{Vect}(X)$ 的对象。令
$(a, b) : G(\mathcal{E}, n) \to G(\mathcal{F}, m)$ 是右端中的态射。
可以选取 $N \gg 0$，并将 $a$ 视为映射
$a : F^{N - n, *}\mathcal{E}|_Z \to F^{N - m, *}\mathcal{F}|_Z$
，将 $b$ 视为映射
$b : F^{N - n, *}\mathcal{E}|_{X'} \to F^{N - m, *}\mathcal{F}|_{X'}$
，二者在 $E$ 上相同。选取有限仿射开覆盖
$X = X_1 \cup \ldots \cup X_n$，使得 $\mathcal{E}|_{X_i}$ 和
$\mathcal{F}|_{X_i}$ 都是有限自由 $\mathcal{O}_{X_i}$-模。
对每个 $i$，基变换
$$
\xymatrix{
E_i \ar[r] \ar[d] & X'_i \ar[d] \\
Z_i \ar[r] & X_i
}
$$
是另一个如 例 \ref{flat-example-two-generators} 的几乎爆破方块。
对于这些方块，由 引理 \ref{flat-lemma-h-sheaf-colim-F}
（见该引理的证明）可知
$$
\colim_F H^0(X_i, \mathcal{O}_{X_i}) =
\colim_F H^0(Z_i, \mathcal{O}_{Z_i})
\times_{\colim_F H^0(E_i, \mathcal{O}_{E_i})}
\colim_F H^0(X'_i, \mathcal{O}_{X'_i})
$$
因此增大 $N$ 后，可以假设映射 $a|_{Z_i}$ 和 $b|_{X'_i}$ 来自映射
$c_i : F^{N - n, *}\mathcal{E}|_{X_i} \to F^{N - m, *}\mathcal{F}|_{X_i}$。
必要时再增大 $N$，可以假设 $c_i$ 与 $c_j$ 在 $X_i \cap X_j$ 上相同。
故这些映射黏合成左端中所需的态射
$(\mathcal{E}, n) \to (\mathcal{F}, m)$。

\medskip\noindent
本质满性。令 $(\mathcal{F}, \mathcal{G}, \varphi)$ 是由有限局部自由
$\mathcal{O}_Z$-模 $\mathcal{F}$、有限局部自由
$\mathcal{O}_{X'}$-模 $\mathcal{G}$ 以及同构
$\varphi : \mathcal{F}|_E \to \mathcal{G}|_E$ 构成的三元组。
我们必须证明，用某个 Frobenius 幂拉回替换该三元组后，它来自一个
有限局部自由 $\mathcal{O}_X$-模。

\medskip\noindent
Noether 情形的归约；建议读者跳过本段。回顾
$X = \Spec(A)$、$Z = \Spec(A/(f_1, f_2))$、
$X' = \text{Proj}(A[T_0, T_1]/(f_2T_0 - f_1T_1))$、
$E = \mathbf{P}^1_Z$。由《极限》，引理
\ref{limits-lemma-descend-invertible-modules}，可找到包含 $f_1$ 和 $f_2$
的有限生成 $\mathbf{F}_p$-子代数 $A_0 \subset A$，使三元组
$(\mathcal{F}, \mathcal{G}, \varphi)$ 下降到
$X_0 = \Spec(A_0)$、$Z_0 = \Spec(A_0/(f_1, f_2))$、
$X_0' = \text{Proj}(A_0[T_0, T_1]/(f_2T_0 - f_1T_1))$、
$E_0 = \mathbf{P}^1_{Z_0}$。因此可以假设这些概形都是 Noether 的。

\medskip\noindent
假设 $X$ 是 Noether 的。可以选取有限仿射开覆盖
$X = X_1 \cup \ldots \cup X_n$，使 $\mathcal{F}|_{Z \cap X_i}$ 自由。
由于可以在 Zariski 拓扑中黏合 $\colim_F \textit{Vect}(X)$ 的对象
（引理 \ref{flat-lemma-colim-F-Vect}），并且我们已经知道在 $X_i$ 和
$X_i \cap X_j$ 上的全忠实性（见证明第一段），只须在每个 $X_i$
上证明存在性。这就归结到下一段讨论的情形。

\medskip\noindent
假设 $X$ 是 Noether 的，且 $\mathcal{F} = \mathcal{O}_Z^{\oplus r}$。
利用 $\varphi$ 得到同构 $\mathcal{O}_E^{\oplus r} \to \mathcal{G}|_E$。
令 $I = (f_1, f_2) \subset A$。令
$\mathcal{I} \subset \mathcal{O}_{X'}$ 是 $E$ 的理想层；它由
$f_1$ 和 $f_2$ 整体生成。对任意 $n$，存在满射
$$
(\mathcal{I}^n/\mathcal{I}^{n + 1})^{\oplus r} =
\mathcal{I}^n/\mathcal{I}^{n + 1} \otimes_{\mathcal{O}_E}
\mathcal{G}|_E \longrightarrow
\mathcal{I}^n\mathcal{G}/\mathcal{I}^{n + 1}\mathcal{G}
$$
因此该模的一阶上同调群为零。这里用到 $E = \mathbf{P}^1_Z$，故其结构层，
事实上任意整体生成的拟凝聚模，都有消失的 $H^1$。比较 《态射进阶》，
引理 \ref{more-morphisms-lemma-globally-generated-vanishing}。
再利用短正合列
$$
0 \to  \mathcal{I}^n\mathcal{G}/\mathcal{I}^{n + 1}\mathcal{G} \to
\mathcal{G}/\mathcal{I}^{n + 1}\mathcal{G} \to
\mathcal{G}/\mathcal{I}^n\mathcal{G} \to 0
$$
和归纳法，可知
$$
\lim H^0(X', \mathcal{G}/\mathcal{I}^n\mathcal{G})
\to
H^0(E, \mathcal{G}|_E) = H^0(E, \mathcal{O}_E^{\oplus r}) =
A/I^{\oplus r}
$$
是满射。由形式函数定理（《概形上同调》，定理
\ref{coherent-theorem-formal-functions}），这蕴含
$$
H^0(X', \mathcal{G}) \to
H^0(E, \mathcal{G}|_E) = H^0(E, \mathcal{O}_E^{\oplus r}) =
A/I^{\oplus r}
$$
是满射。因此可选取与给定的 $\mathcal{G}|_E$ 平凡化相容的映射
$\alpha : \mathcal{O}_{X'}^{\oplus r} \to \mathcal{G}$。
故 $\alpha$ 在 $X'$ 中 $E$ 的某个开邻域上是同构。于是 $Z$ 的每一点都有
一个可在其上解决问题的仿射开邻域。由于
$X' \setminus E \to X \setminus Z$ 是同构，对 $X$ 中不在 $Z$ 内的点亦然。
再作一次 Zariski 黏合即完成证明。
% P05-EN-CH38-0092: source 12418 and 12426 typesets A/I^{\oplus r}; the direct-sum exponent should apply to A/I.
\end{proof}

\begin{proposition}
\label{flat-proposition-h-descent-vector-bundles-p}
设 $p$ 是素数，$S$ 是特征 $p$ 的概形。则群胚纤维化范畴
$$
p : \mathcal{S} \longrightarrow (\Sch/S)_h
$$
是群胚叠，其中 $U$ 上的纤维范畴是 $U$ 上有限局部自由
$\colim_F \mathcal{O}_U$-模的范畴。此外，若 $U$ 拟紧且拟分离，
则 $\mathcal{S}_U$ 是 $\colim_F \textit{Vect}(U)$。
\end{proposition}

\begin{proof}
最后一个断言正是 引理 \ref{flat-lemma-colim-F-Vect} 的内容。
为证明该命题，我们核验 引理 \ref{flat-lemma-refine-check-h-stack}
的条件 (1)、(2) 和 (3)。

\medskip\noindent
条件 (1) 成立，因为按照定义，我们在 Zariski 拓扑中可以作黏合。

\medskip\noindent
为证明条件 (2)，设 $f : X \to Y$ 是 $S$ 上有限表现、满、平坦且固有的态射，
并且 $Y$ 仿射。由于 $Y, X, X \times_Y X$ 都拟紧且拟分离，
可以使用命题陈述中给出的纤维范畴描述。显然只须证明
$$
\textit{Vect}(Y)
\longrightarrow
\textit{Vect}(X)
\times_{\textit{Vect}(X \times_Y X)}
\textit{Vect}(X)
$$
是等价（因为这将蕴含余极限的相同结论）。这立即由有限局部自由模的
fppf 下降得到；见《下降》，命题
\ref{descent-proposition-fpqc-descent-quasi-coherent} 和
引理 \ref{descent-lemma-finite-locally-free-descends}。

\medskip\noindent
条件 (3) 正是 引理 \ref{flat-lemma-vector-bundle-I} 和
\ref{flat-lemma-vector-bundle-II} 的内容。
\end{proof}

\begin{lemma}
\label{flat-lemma-trivial-fibres-dvr}
设 $f : X \to S$ 是具有几何连通纤维的固有态射，其中 $S$ 是离散赋值环的谱。
以 $\eta \in S$ 表示一般点，并以 $X_n \subset X$ 表示由 $S$ 上某个
素元的第 $n$ 次幂截出的闭子概形。则存在整数 $n$，使得如下断言成立：
任意有限局部自由 $\mathcal{O}_X$-模 $\mathcal{E}$，只要
$\mathcal{E}|_{X_\eta}$ 和 $\mathcal{E}|_{X_n}$ 都自由，它就是自由的。
% P05-EN-CH38-0093: source 12490 writes the phrasal verb “cutout” as one word.
\end{lemma}

\begin{proof}
首先归约到 $X \to S$ 有截面的情形。记 $S = \Spec(A)$。
选取 $X_\eta$ 的闭点 $\xi$。选取离散赋值环扩张 $A \subset B$，
使得 $B$ 的分式域是 $\kappa(\xi)$。由 Krull--Akizuki
（《代数》，引理 \ref{algebra-lemma-integral-closure-Dedekind}）
以及 $\kappa(\xi)$ 是 $A$ 的分式域的有限扩张，这是可以做到的。
由固有性的赋值判据（《态射》，引理
\ref{morphisms-lemma-characterize-proper}），得到 $B$-值点
$\tau : \Spec(B) \to X$，它诱导截面 $\sigma : \Spec(B) \to X_B$。
对有限局部自由 $\mathcal{O}_X$-模 $\mathcal{E}$，令 $\mathcal{E}_B$
为它在基变换 $X_B$ 上的拉回。由平坦基变换（《概形上同调》，
引理 \ref{coherent-lemma-flat-base-change-cohomology}），有
$H^0(X_B, \mathcal{E}_B) = H^0(X, \mathcal{E}) \otimes_A B$。
因此，若 $\mathcal{E}_B$ 是秩 $r$ 的自由模，则
$H^0(X, \mathcal{E})$ 中的截面生成自由 $B$-模
$\tau^*\mathcal{E} = \sigma^*\mathcal{E}_B$。特别地，可以找到
$\mathcal{E}$ 的 $r$ 个整体截面 $s_1, \ldots, s_r$，它们生成
$\tau^*\mathcal{E}$。于是
$$
s_1, \ldots, s_r :
\mathcal{O}_X^{\oplus r}
\longrightarrow
\mathcal{E}
$$
是秩 $r$ 的有限局部自由 $\mathcal{O}_X$-模之间的映射；它在 $X_B$
上的拉回是自由 $\mathcal{O}_{X_B}$-模之间的映射，并且在每个纤维的
某一点处限制为同构。取行列式，得到函数
$g \in \Gamma(X_\eta, \mathcal{O}_{X_B})$，它在每个纤维的某一点处可逆。
由于纤维固有且连通，$g$ 必须可逆（细节从略；提示：使用《簇》，
引理 \ref{varieties-lemma-proper-geometrically-reduced-global-sections}）。
因此只须对基变换 $X_B \to \Spec(B)$ 证明本引理。
% P05-EN-CH38-0094: source 12530 places an \mathcal{O}_{X_B}-section in \Gamma(X_\eta, -), although the determinant is taken on X_B.

\medskip\noindent
假设已有截面 $\sigma : S \to X$。令 $\mathcal{E}$ 是有限局部自由
$\mathcal{O}_X$-模，并假设它在一般纤维和 $X_n$ 上自由
（稍后选取 $n$）。选取同构
$\sigma^*\mathcal{E} = \mathcal{O}_S^{\oplus r}$。考虑 $D(A)$ 中的映射
$$
K = R\Gamma(X, \mathcal{E}) \longrightarrow
R\Gamma(S, \sigma^*\mathcal{E}) = A^{\oplus r}
$$
同上论证可知，$\mathcal{E}$ 自由，当且仅当所诱导的映射
$H^0(K) = H^0(X, \mathcal{E}) \to A^{\oplus r}$ 是满射。

\medskip\noindent
令 $L = R\Gamma(X, \mathcal{O}_X^{\oplus r})$，并注意相应的映射
$L \to A^{\oplus r}$ 具有所需性质。由平坦基变换以及
$\mathcal{E}$ 在一般纤维上自由的假设，有
$K \otimes_A Q(A) \cong L \otimes_A Q(A)$。令 $\pi \in A$ 是素元。
注意
$$
K \otimes_A^\mathbf{L} A/\pi^m A =
R\Gamma(X, \mathcal{E} \xrightarrow{\pi^m} \mathcal{E})
$$
对 $L$ 亦然。以 $\mathcal{E}_{tors} \subset \mathcal{E}$ 表示由支集在
特殊纤维上的截面组成的凝聚子层，对其他 $\mathcal{O}_X$-模也作同样记号。
选取 $k > 0$，使得
$(\mathcal{O}_X)_{tors} \to \mathcal{O}_X/\pi^k \mathcal{O}_X$
是单射（《概形上同调》，引理
\ref{coherent-lemma-Artin-Rees}）。由于 $\mathcal{E}$ 局部自由，可知
$\mathcal{E}_{tors} \subset \mathcal{E}/\pi^k\mathcal{E}$。
于是对 $n \geq m + k$，在 $D(\mathcal{O}_X)$ 中有同构
\begin{align*}
(\mathcal{E} \xrightarrow{\pi^m} \mathcal{E})
& \cong
(\mathcal{E}/\pi^k\mathcal{E} \xrightarrow{\pi^m}
\mathcal{E}/\pi^{k + m}\mathcal{E}) \\
& \cong
(\mathcal{O}_X^{\oplus r}/\pi^k\mathcal{O}_X^{\oplus r} \xrightarrow{\pi^m}
\mathcal{O}_X^{\oplus r}/\pi^{k + m}\mathcal{O}_X^{\oplus r}) \\
& \cong
(\mathcal{O}_X^{\oplus r} \xrightarrow{\pi^m} \mathcal{O}_X^{\oplus r})
\end{align*}
这在 $D(A)$ 中确定一个同构
$$
K \otimes_A^\mathbf{L} A/\pi^m A \cong L \otimes_A^\mathbf{L} A/\pi^m A
$$
（当 $n \geq m + k$ 时成立）。注意，这些同构与沿 $\sigma$ 拉回相容，
因而特别可知
$K \otimes_A^\mathbf{L} A/\pi^m A \to (A/\pi^m A)^{\oplus r}$
在 $0$ 次上同调模上诱导满射（因为对 $L$ 而言这是成立的）。
% P05-EN-CH38-0095: source 12588 has the article error “an surjection”.
由于 $A$ 是离散赋值环，在 $D(A)$ 中有
$$
K \cong \bigoplus H^i(K)[-i]
\quad\text{且}\quad
L \cong
\bigoplus H^i(L)[-i]
$$
参见《代数进阶》，例
\ref{more-algebra-example-finite-injective-finite-global-dimension}。
上同调群 $H^i(K) = H^i(X, \mathcal{E})$ 和
$H^i(L) = H^i(X, \mathcal{O}_X)^{\oplus r}$
由《概形上同调》，引理
\ref{coherent-lemma-proper-over-affine-cohomology-finite} 是有限 $A$-模。
由《代数进阶》，引理
\ref{more-algebra-lemma-generalized-valuation-ring-modules}，这些模都是
循环模的直和。我们已经看到，$H^i(K)$ 与 $H^i(L)$ 的自由部分具有
相同的秩 $\beta_i$。接着注意
$$
H^i(L \otimes_A^\mathbf{L} A/\pi^m A) =
H^i(L)/\pi^m H^i(L) \oplus H^{i + 1}(L)[\pi^m]
$$
对 $K$ 亦然。令 $e$ 是满足如下条件的最大整数：对某个 $i$，
$A/\pi^eA$ 是 $H^i(X, \mathcal{O}_X)$（等价地，$H^i(L)$）的直和项。
取 $m = e + 1$，则 $H^i(L \otimes_A^\mathbf{L} A/\pi^m A)$ 是
$\beta_i$ 个 $A/\pi^m A$ 与另外若干循环模的直和，其中每个循环模都被
$\pi^e$ 零化。对 $K$ 作同样论证，并利用同构
% P05-EN-CH38-0096: source 12613--12615 does not define e when all relevant cohomology modules are torsion-free.
$K \otimes_A^\mathbf{L} A/\pi^m A \cong L \otimes_A^\mathbf{L} A/\pi^m A$
可知，$H^i(K)$ 不可能有形如 $A/\pi^l A$、其中 $l > e$ 的循环直和项。
（这还蕴含 $K$ 与 $L$ 作为 $D(A)$ 的对象同构，但我们不需要这一点。）
于是映射
$$
H^0(K \otimes^\mathbf{L}_A A/\pi^{e + 1} A) =
H^0(K)/\pi^{e + 1}H^0(K) \oplus H^1(K)[\pi^{e + 1}]
\longrightarrow
(A/\pi^{e + 1} A)^{\oplus r}
$$
为满射的唯一可能，是它在第一个直和项上为满射。这正是所要证明的。
（更准确地说，若有截面 $\sigma$，则引理陈述中的整数 $n$ 应等于
$k + e + 1$，其中 $k$ 与 $e$ 如上，并且只依赖于 $X$。）
\end{proof}

\begin{lemma}
\label{flat-lemma-trivial-over-dvrs}
设 $f : X \to S$ 是概形态射，$\mathcal{E}$ 是有限局部自由
$\mathcal{O}_X$-模。假设
\begin{enumerate}
\item $f$ 平坦且固有，并且 $\mathcal{O}_S = f_*\mathcal{O}_X$；
\item $S$ 是正规 Noether 概形；
\item 对每个余维 $1$ 的点 $s \in S$，$\mathcal{E}$ 在
$X \times_S \Spec(\mathcal{O}_{S, s})$ 上的拉回是自由的。
\end{enumerate}
则 $\mathcal{E}$ 同构于某个有限局部自由 $\mathcal{O}_S$-模的拉回。
\end{lemma}

\begin{proof}
我们证明典范映射
$$
\Phi : f^*f_*\mathcal{E} \longrightarrow \mathcal{E}
$$
是同构。由平坦基变换（《概形上同调》，引理
\ref{coherent-lemma-flat-base-change-cohomology}）以及假设 (1) 和 (3)，
对每个余维 $1$ 的点 $s \in S$，该映射在
$X \times_S \Spec(\mathcal{O}_{S, s})$ 上的拉回是同构。
由《因子》，引理 \ref{divisors-lemma-check-isomorphism-via-depth-and-ass}，
只须证明：对任意映到余维 $\geq 2$ 的点 $s \in S$ 的点 $x \in X$，有
$\text{depth}((f^*f_*\mathcal{E})_x) \geq 2$。由于 $f$ 平坦且
$(f^*f_*\mathcal{E})_x = (f_*\mathcal{E})_s \otimes_{\mathcal{O}_{S, s}}
\mathcal{O}_{X, x}$，只须证明 $\text{depth}((f_*\mathcal{E})_s) \geq 2$；见
《代数》，引理 \ref{algebra-lemma-apply-grothendieck}。
由于 $S$ 是正规 Noether 概形且 $\dim(\mathcal{O}_{S, s}) \geq 2$，
有 $\text{depth}(\mathcal{O}_{S, s}) \geq 2$；见《性质》，引理
\ref{properties-lemma-criterion-normal}。于是所需结论由《因子》，
引理 \ref{divisors-lemma-depth-pushforward} 得到。
\end{proof}

\noindent
利用上述结果，可以证明如下惊人的断言。

\begin{theorem}
\label{flat-theorem-pullback-trivial-fibres}
设 $p$ 是素数，$Y$ 是 $\mathbf{F}_p$ 上拟紧且拟分离的概形。
设 $f : X \to Y$ 是具有几何连通纤维的有限表现满固有态射。则函子
$$
\colim_F \textit{Vect}(Y) \longrightarrow \colim_F \textit{Vect}(X)
$$
全忠实，其本质像描述如下。令 $\mathcal{E}$ 是有限局部自由
$\mathcal{O}_X$-模。假设对所有 $y \in Y$，存在整数
$n_y, r_y \geq 0$，使得
% P05-EN-CH38-0097: source 12697 has singular “there exists integers”.
$$
F^{n_y, *}\mathcal{E}|_{X_{y, red}}
\cong
\mathcal{O}_{X_{y, red}}^{\oplus r_y}
$$
则对某个 $n \geq 0$，第 $n$ 次 Frobenius 幂拉回
$F^{n, *}\mathcal{E}$ 是某个有限局部自由 $\mathcal{O}_Y$-模的拉回。
\end{theorem}

\begin{proof}
证明全忠实性。由于 $Y$ 上的向量丛局部平凡，这归结为证明
% P05-EN-CH38-0098: source 12710 closes “vector bundles” into “vectorbundles”.
$$
\colim_F \Gamma(Y, \mathcal{O}_Y)
\longrightarrow
\colim_F \Gamma(X, \mathcal{O}_X)
$$
是双射。由于 $\{X \to Y\}$ 是 h 覆盖，只要证明两个映射
$$
\colim_F \Gamma(X, \mathcal{O}_X)
\longrightarrow
\colim_F \Gamma(X \times_Y X, \mathcal{O}_{X \times_Y X})
$$
相等，即由 引理 \ref{flat-lemma-h-sheaf-colim-F} 得到该结论。
令 $g \in \Gamma(X, \mathcal{O}_X)$，以 $g_1$ 和 $g_2$ 表示 $g$
在 $X \times_Y X$ 上的两个拉回。由于 $X_{y, red}$ 几何连通，
$H^0(X_{y, red}, \mathcal{O}_{X_{y, red}})$ 是 $\kappa(y)$ 的纯不可分扩张；
见《簇》，引理
\ref{varieties-lemma-proper-geometrically-reduced-global-sections}。
因此对某个 $p$ 的幂 $q$（它可以依赖于 $y$），
$g^q|_{X_{y, red}}$ 来自 $\kappa(y)$ 中的元素。于是对
$(X \times_Y X)_y = X_y \times_y X_y$ 的任意点，$g_1^q$ 和 $g_2^q$
在其剩余域中的像相同。因此 $g_1 - g_2$ 在
$(X \times_Y X)_{red}$ 上限制为零。故对某个 $n$ 有
$(g_1 - g_2)^n = 0$；按需要，可将它取为 $p$ 的幂。

\medskip\noindent
本质像的描述。令 $\mathcal{E}$ 如该命题陈述中所述。我们首先归约到
Noether 情形。
% P05-EN-CH38-0099: source 12741--12742 calls the enclosing theorem a proposition.

\medskip\noindent
令 $y \in Y$ 是一点，并将其视为从剩余域的谱到 $Y$ 的态射
$y \to Y$。可以将 $y \to Y$ 写成有限表现态射 $Y_i \to Y$ 的滤过极限，
其中 $Y_i$ 仿射。（最好自行证明这一点；它也可形式地由《极限》，
引理 \ref{limits-lemma-relative-approximation} 和
\ref{limits-lemma-limit-affine} 推出。）对每个 $i$，令
$Z_i = Y_i \times_Y X$。则 $X_y = \lim Z_i$ 且
$X_{y, red} = \lim Z_{i, red}$。由《极限》，引理
\ref{limits-lemma-descend-modules-finite-presentation}，可找到 $i$，使得
$F^{n_y, *}\mathcal{E}|_{Z_{i, red}} \cong
\mathcal{O}_{Z_{i, red}}^{\oplus r_y}$.
固定该 $i$。有 $Z_{i, red} = \lim Z_{i, j}$，其中
$Z_{i, j} \to Z_i$ 是有限表现的增厚（《极限》，引理
\ref{limits-lemma-closed-is-limit-closed-and-finite-presentation}）。
使用与此前相同的引理，可找到 $j$，使得
$F^{n_y, *}\mathcal{E}|_{Z_{i, j}} \cong \mathcal{O}_{Z_{i, j}}^{\oplus r_y}$.
由此可知，对每个 $y \in Y$，存在像包含 $y$ 的有限表现态射
$Y_y \to Y$ 以及增厚 $Z_y \to Y_y \times_Y X$，使得
$F^{n_y, *}\mathcal{E}|_{Z_y} \cong \mathcal{O}_{Z_y}^{\oplus r_y}$.
注意，$Y_y \to Y$ 的像是可构造的（《态射》，引理
\ref{morphisms-lemma-chevalley}）。由于 $Y$ 在可构造拓扑中拟紧
（《拓扑》，引理 \ref{topology-lemma-constructible-hausdorff-quasi-compact}
和 《性质》，引理
\ref{properties-lemma-quasi-compact-quasi-separated-spectral}），
存在有限多个有限表现态射
$$
Y_1 \to Y,\ Y_2 \to Y,\ \ldots,\ Y_N \to Y
$$
使得集合论意义下 $Y = \bigcup \Im(Y_a \to Y)$，并且对每个
$a \in \{1, \ldots, N\}$，存在整数 $n_a, r_a \geq 0$ 和有限表现的增厚
$Z_a \subset Y_a \times_Y X$，使得
$F^{n_a, *}\mathcal{E}|_{Z_a} \cong \mathcal{O}_{Z_a}^{\oplus r_a}$.

\medskip\noindent
按此方式表述后，该条件可下降到一个绝对 Noether 逼近。强烈建议读者跳过
本段。首先，将 $Y = \lim_{i \in I} Y_i$ 写成 $\mathbf{F}_p$ 上有限型概形
的余滤过极限，其中转移态射仿射（《极限》，引理
\ref{limits-lemma-relative-approximation}）。其次，可以假设已有固有态射
$f_i : X_i \to Y_i$，其到 $Y$ 的基变换恢复 $f : X \to Y$；见《极限》，
引理 \ref{limits-lemma-descend-finite-presentation}。增大 $i$ 后，可以假设
存在有限局部自由 $\mathcal{O}_{X_i}$-模 $\mathcal{E}_i$，它在 $X$ 上的
拉回同构于 $\mathcal{E}$；见《极限》，引理
\ref{limits-lemma-descend-invertible-modules}。
选取 $0 \in I$，以 $E \subset Y_0$ 表示 $f_0$ 的几何纤维连通的可构造子集；
见《态射进阶》，引理
\ref{more-morphisms-lemma-nr-geom-connected-components-constructible}。
于是 $Y \to Y_0$ 映入 $E$；见《态射进阶》，引理
\ref{more-morphisms-lemma-base-change-fibres-geometrically-connected}。
故当 $i \gg 0$ 时，$Y_i \to Y_0$ 映入 $E$；见《极限》，引理
\ref{limits-lemma-limit-contained-in-constructible}。因此当 $i \gg 0$ 时，
$f_i$ 的纤维几何连通。由《极限》，引理
\ref{limits-lemma-descend-finite-presentation}，对充分大的 $i$，可以找到
有限型态射 $Y_{i, a} \to Y_i$，其到 $Y$ 的基变换恢复
$Y_a \to Y$，其中 $a  \in \{1, \ldots, N\}$。必要时增大 $i$ 后，
可以找到增厚 $Z_{i, a} \subset Y_{i, a} \times_{Y_i} X_i$，其到
$Y_a \times_Y X$ 的基变换恢复 $Z_a$（使用与此前相同的参考结果，并结合
《极限》，引理 \ref{limits-lemma-descend-closed-immersion-finite-presentation}
和 \ref{limits-lemma-descend-surjective}）。由于
$Z_a = \lim Z_{i, a}$，再增大 $i$ 后，可以假设
$F^{n_a, *}\mathcal{E}_i|_{Z_{i, a}} \cong
\mathcal{O}_{Z_{i, a}}^{\oplus r_a}$；见《极限》，引理
\ref{limits-lemma-descend-modules-finite-presentation}。最后，再增大一次 $i$，
由《极限》，引理 \ref{limits-lemma-descend-surjective}，可以假设
$\coprod Y_{i, a} \to Y_i$ 是满态射。此时所有假设对
$X_i \to Y_i$ 和 $\mathcal{E}_i$ 都成立，因此只须对
$X_i \to Y_i$ 和 $\mathcal{E}_i$ 证明该结果。
% P05-EN-CH38-0100: source 12825 omits “the” in “prove result”.

\medskip\noindent
假设 $Y$ 在 $\mathbf{F}_p$ 上有限型。为证明该结果，我们对
$\dim(Y)$ 作归纳。我们要找到 $\colim_F \textit{Vect}(Y)$ 的一个对象，
其拉回是由 $\mathcal{E}$ 确定的 $\colim_F \textit{Vect}(X)$ 对象。
由已经证明的全忠实性以及 命题
\ref{flat-proposition-h-descent-vector-bundles-p}，用一个 h 覆盖的成员替换
$Y$ 并对 $X$ 作相应基变换后，只须构造 $\mathcal{E}$ 的下降。
这意味着，只要不增加 $Y$ 的维数，就可以用一个在 $Y$ 上固有且满的
概形替换 $Y$。若 $T \subset T'$ 是 $\mathbf{F}_p$ 上有限型概形的增厚，则
$\colim_F \textit{Vect}(T) = \colim_F \textit{Vect}(T')$
，因为 $\{T \to T'\}$ 是满足 $T \times_{T'} T = T$ 的 h 覆盖。
若 $T' \to T$ 是 $\mathbf{F}_p$ 上有限型概形之间的泛同胚，则
$\colim_F \textit{Vect}(T) = \colim_F \textit{Vect}(T')$
，因为 $\{T \to T'\}$ 是 h 覆盖，且对角态射
$T \subset T \times_{T'} T$ 是增厚。
% P05-EN-CH38-0101: source 12845--12849 reverses T and T' in both the h-cover arrow and the diagonal fibre product for T' -> T.

\medskip\noindent
利用上述一般说明，我们用其约化替换 $X$，并假设 $X$ 既约。
考虑 Stein 分解 $X \to Y' \to Y$；见《态射进阶》，定理
\ref{more-morphisms-theorem-stein-factorization-Noetherian}。
则 $Y' \to Y$ 是 $\mathbf{F}_p$ 上有限型概形之间的泛同胚。
由上述说明，可以用 $Y'$ 替换 $Y$。因此可以假设
$f_*\mathcal{O}_X = \mathcal{O}_Y$，且 $Y$ 既约。这就归结到下一段
讨论的情形。

\medskip\noindent
假设 $Y$ 既约，并且在 $Y$ 的某个稠密开子概形上有
$f_*\mathcal{O}_X = \mathcal{O}_Y$。则 $X \to Y$ 在某个稠密开子概形
$V \subset Y$ 上平坦；见《态射》，命题
\ref{morphisms-proposition-generic-flatness-reduced}。
由 引理 \ref{flat-lemma-flat-after-blowing-up}，存在 $V$-容许爆破
$Y' \to Y$，使 $X$ 的严格变换 $X'$ 在 $Y'$ 上平坦。由于 $Y$ 与 $Y'$
有公共稠密开子概形，$\dim(Y') = \dim(Y)$。由《态射进阶》，引理
\ref{more-morphisms-lemma-proper-flat-nr-geom-conn-comps-lower-semicontinuous}
以及 $V \subset Y'$ 稠密这一事实，$f' : X' \to Y'$ 的所有纤维都几何连通。
仍有 $(f'_*\mathcal{O}_{X'})|_V = \mathcal{O}_V$。写成
$$
Y' \times_Y X = X' \cup E \times_Y X
$$
其中 $E \subset Y'$ 是该爆破的例外除子。由上述一般说明，只须对
$Y' \times_Y X \to Y'$ 以及 $\mathcal{E}$ 在 $Y' \times_Y X$ 上的限制
证明存在性。设我们找到了 $\colim_F \textit{Vect}(Y')$ 中的对象
$\xi'$，它拉回到 $\mathcal{E}$ 在 $X'$ 上的限制（视为该余极限范畴的对象）。
对 $\dim(Y)$ 作归纳，可找到 $\colim_F \textit{Vect}(E)$ 中的对象
$\xi''$，它拉回到 $\mathcal{E}$ 在 $E \times_Y X$ 上的限制。
于是全忠实性确定唯一同构 $\xi'|_E \to \xi''$，它与给定的、到
$\mathcal{E}$ 在 $E \times_{Y'} X'$ 上限制的等同相容。由于
% P05-EN-CH38-0102: source 12892 misspells “full faithfulness” as “fully faithfullness”.
$$
\{E \times_Y X \to Y' \times_Y X, X' \to Y' \times_Y X\}
$$
是由一对闭浸入给出的 h 覆盖，且
$$
(E \times_Y X) \times_{(Y' \times_Y X)} X' = E \times_{Y'} X'
$$
由此可知 $\xi'$ 拉回到 $\mathcal{E}$ 在 $Y' \times_Y X$ 上的限制。
因此只须找到 $\xi'$，这就归结到下一段讨论的情形。

\medskip\noindent
假设 $Y$ 既约，$f$ 平坦，并且在 $Y$ 的某个稠密开子概形上有
$f_*\mathcal{O}_X = \mathcal{O}_Y$。在此情形下，考虑正规化
$Y^\nu \to Y$（《态射》，第 \ref{morphisms-section-normalization} 节）。这是在某个稠密开集上为同构的
有限满态射（《态射》，引理
\ref{morphisms-lemma-nagata-normalization} 和
\ref{morphisms-lemma-ubiquity-nagata}）。因此由上述一般说明，可以用
$Y^\nu$ 替换 $Y$，并用 $Y^\nu \times_Y X$ 替换 $X$。替换后有
$\mathcal{O}_Y = f_*\mathcal{O}_X$（因为在此情形下 Stein 分解必为同构；
略去一个小细节）。

\medskip\noindent
假设 $Y$ 是正规 Noether 概形，$f$ 平坦，且
$f_*\mathcal{O}_X = \mathcal{O}_Y$。用适当的 Frobenius 幂拉回替换
$\mathcal{E}$ 后，可以假设 $\mathcal{E}$ 在 $f$ 位于 $Y$ 各不可约分支
一般点处的概形论纤维上平凡（因为 $\colim_F \textit{Vect}(-)$ 在
泛同胚上是等价；见上文）。与上述论证（归约到 Noether 情形时）类似，
可知存在稠密开子概形 $V \subset Y$，使得
$\mathcal{E}|_{f^{-1}(V)}$ 自由。令 $Z \subset Y$ 是一个闭子概形，
使集合论意义下 $Y = V \amalg Z$。令 $z_1, \ldots, z_t \in Z$ 是 $Z$
中余维 $1$ 的不可约分支的一般点。则 $A_i = \mathcal{O}_{Y, z_i}$
是离散赋值环。令 $n_i$ 是 引理 \ref{flat-lemma-trivial-fibres-dvr}
对 $A_i$ 上概形 $X_{A_i}$ 所给出的整数。用适当的 Frobenius 幂拉回替换
$\mathcal{E}$ 后，可以假设 $\mathcal{E}$ 在
$X_{A_i/\mathfrak m_i^{n_i}}$ 上自由（同样因为余极限范畴在泛同胚下不变；
见上文）。于是 引理 \ref{flat-lemma-trivial-fibres-dvr} 表明
$\mathcal{E}$ 在 $X_{A_i}$ 上自由。最后应用
引理 \ref{flat-lemma-trivial-over-dvrs} 即得结论。
\end{proof}










\section{复形的爆破}
\label{flat-section-blowup-complexes}

\noindent
本节寻找导出范畴的完美对象在爆破后的正规形。

\begin{lemma}
\label{flat-lemma-fitting-ideals-complex}
设 $X$ 是概形，$E \in D(\mathcal{O}_X)$ 是伪凝聚的。对每个
$p, k \in \mathbf{Z}$，存在有限型拟凝聚理想层
$\text{Fit}_{p, k}(E) \subset \mathcal{O}_X$，具有如下性质：
若 $U \subset X$ 是开子集，且 $E|_U$ 同构于
% P05-EN-CH38-0103: source 12965 has the article error “an finite type”.
$$
\ldots \to
\mathcal{O}_U^{\oplus n_{b - 2}}
\xrightarrow{d_{b - 2}}
\mathcal{O}_U^{\oplus n_{b - 1}}
\xrightarrow{d_{b - 1}}
\mathcal{O}_U^{\oplus n_b} \to 0 \to \ldots
$$
则限制 $\text{Fit}_{p, k}(E)|_U$ 由 $d_p$ 的矩阵中如下阶数的
子式生成：
$$
- k + n_{p + 1} - n_{p + 2} + \ldots + (-1)^{b - p + 1} n_b
$$
约定：若 $r \leq 0$，则由 $r \times r$ 子式生成的理想是
$\mathcal{O}_U$；若 $r > \min(n_p, n_{p + 1})$，则由
$r \times r$ 子式生成的理想为零。
\end{lemma}

\begin{proof}
注意，$E$ 在 $X$ 上局部具有引理中所述的形式；见《代数进阶》，
第 \ref{more-algebra-section-pseudo-coherent} 节，《上同调》，
第 \ref{cohomology-section-pseudo-coherent} 节，以及
《概形的导出范畴》，第 \ref{perfect-section-spell-out} 节。
因此只须证明该子式理想与所选代表无关。为此，只须在局部环中核验。
在局部环 $(R, \mathfrak m, \kappa)$ 上，考虑上有界复形
$$
F^\bullet :
\ldots \to
R^{\oplus n_{b - 2}}
\xrightarrow{d_{b - 2}}
R^{\oplus n_{b - 1}}
\xrightarrow{d_{b - 1}}
R^{\oplus n_b} \to 0 \to \ldots
$$
以 $\text{Fit}_{k, p}(F^\bullet) \subset R$ 表示由 $d_p$ 的矩阵中
阶数为 $k - n_{p + 1} + n_{p + 2} - \ldots + (-1)^{b - p} n_b$ 的
子式生成的理想。假设 $F^\bullet$ 某个微分的某个矩阵系数可逆。
选取最大的整数 $i$，使 $d_i$ 有可逆的矩阵系数。由《代数》，引理
\ref{algebra-lemma-add-trivial-complex}，复形 $F^\bullet$ 同构于如下直和：
一个仅在次数 $i$ 和 $i + 1$ 有非零项的平凡复形
$\ldots \to 0 \to R \to R \to 0 \to \ldots$，以及一个复形
$(F')^\bullet$。留给读者验证
% P05-EN-CH38-0104: source 13005--13006 swaps the Fit indices and reverses the stated minor-size formula.
$\text{Fit}_{p, k}(F^\bullet) = \text{Fit}_{p, k}((F')^\bullet)$;
；这里用到了子式阶数的公式。若 $(F')^\bullet$ 还有某个微分具有可逆的
矩阵系数，就再次作同样处理，依此类推。
继续这一过程，最终得到复形 $(F^\infty)^\bullet$，其所有微分矩阵的系数
都在 $\mathfrak m$ 中。这里可能需要作无穷多步，但对任意截断，
只有有限多步会影响次数 $\geq $ 该截断的复形。因此，“极限”
$(F^\infty)^\bullet$ 是有限自由模的良定义上有界复形，并配有到起始复形的
拟同构 $(F^\infty)^\bullet \to F^\bullet$，且
% P05-EN-CH38-0105: source 13023 leaves “the cutoff” outside a malformed math fragment “degrees $\geq $”.
$\text{Fit}_{p, k}(F^\bullet) = \text{Fit}_{p, k}((F^\infty)^\bullet)$.
由《代数进阶》，引理
\ref{more-algebra-lemma-lift-pseudo-coherent-from-residue-field}，
复形 $(F^\infty)^\bullet$ 在同构意义下唯一，证明完成。
\end{proof}

\begin{lemma}
\label{flat-lemma-blowup-complex}
设 $X$ 是概形，$E \in D(\mathcal{O}_X)$ 是完美的。设
$U \subset X$ 是概形论稠密的开子概形，使得对所有
$i \in \mathbf{Z}$，$H^i(E|_U)$ 都是常秩 $r_i$ 的有限局部自由模。
则存在 $U$-容许爆破 $b : X' \to X$，使得对所有
$i \in \mathbf{Z}$，$H^i(Lb^*E)$ 都是 Tor 维数 $\leq 1$ 的完美
$\mathcal{O}_{X'}$-模。
% P05-EN-CH38-0106: source 13038, 13084, and 13129 use “scheme theoretically” without the compound-adverb hyphen.
\end{lemma}

\begin{proof}
我们在仿射局部构造并研究该爆破。具体而言，设 $V \subset X$ 是仿射开子概形，
使得 $E|_V$ 可由复形
$$
\mathcal{O}_V^{\oplus n_a} \xrightarrow{d_a}
\ldots \xrightarrow{d_{b - 1}} \mathcal{O}_V^{\oplus n_b}
$$
表示。令 $k_ i = r_{i + 1} - r_{i + 2} + \ldots + (-1)^{b - i + 1}r_b$。
略去的一个计算表明，$d_i$ 在 $U \cap V$ 上的秩是
% P05-EN-CH38-0107: source 13054 has a stray space in the subscript token k_ i.
% P05-EN-CH38-0108: source 13055 has singular “A computation” with bare “show”.
$$
\rho_i = - k_i + n_{i + 1} - n_{i + 2} + \ldots + (-1)^{b - i + 1}n_b
$$
这里的含义是 $d_i$ 的余核是秩 $n_{i + 1} - \rho_i$ 的有限局部自由模。
令 $\mathcal{I}_i \subset \mathcal{O}_V$ 是由 $d_i$ 的矩阵中
$\rho_i \times \rho_i$ 子式生成的理想。

\medskip\noindent
一方面，与 引理 \ref{flat-lemma-fitting-ideals-complex} 比较可知，
理想 $\mathcal{I}_i$ 对应于整体理想 $\text{Fit}_{i, k_i}(E)$；已经证明
后者与表示 $E|_V$ 的复形的选择无关。另一方面，$\mathcal{I}_i$ 是
$\Coker(d_i)$ 的第 $(n_{i + 1} - \rho_i)$ 个 Fitting 理想。请记住这一点。

\medskip\noindent
令 $b : X' \to X$ 是沿理想 $\text{Fit}_{i, k_i}(E)$ 之积的爆破；
这是有意义的，因为在 $X$ 上局部，这些理想中几乎全都等于单位理想
（见上文）。该爆破控制沿理想 $\text{Fit}_{i, k_i}(E)$ 的各个爆破
$b_i : X'_i \to X$；见《因子》，引理
\ref{divisors-lemma-blowing-up-two-ideals}。由《因子》，引理
\ref{divisors-lemma-blowup-fitting-ideal}，每个 $b_i$ 都是 $U$-容许爆破。
因此 $b$ 是 $U$-容许爆破（略去一个小细节；比较 《因子》，引理
\ref{divisors-lemma-dominate-admissible-blowups} 的证明）。最后，$U$
仍是 $X'$ 的概形论稠密开子概形。因此用 $X'$ 替换 $X$ 后，便处于下一段
讨论的情形。

\medskip\noindent
假设对所有 $i$，$\text{Fit}_{i, k_i}(E)$ 都是可逆理想。选取仿射开集
$V$ 以及如上表示 $E|_V$ 的有限自由模复形。由《因子》，引理
\ref{divisors-lemma-blowup-fitting-ideal}，$\Coker(d_i)$ 的 Tor 维数
$\leq 1$。因此，作为从有限局部自由模到 Tor 维数 $\leq 1$ 的有限表现模
之映射的核，$\Im(d_i)$ 是有限局部自由的。于是 $\Ker(d_i)$ 也有限局部自由
（同一论证）。故短正合列
$$
0 \to \Im(d_{i - 1}) \to \Ker(d_i) \to H^i(E)|_V \to 0
$$
给出所需结论，证明完成。
\end{proof}

\begin{lemma}
\label{flat-lemma-blowup-complex-integral}
设 $X$ 是整概形，$E \in D(\mathcal{O}_X)$ 是完美的。则存在非空开集
$U \subset X$，使得对所有 $i \in \mathbf{Z}$，$H^i(E|_U)$ 都是
常秩 $r_i$ 的有限局部自由模；并且存在 $U$-容许爆破
$b : X' \to X$，使得对所有 $i \in \mathbf{Z}$，$H^i(Lb^*E)$ 都是
Tor 维数 $\leq 1$ 的完美 $\mathcal{O}_{X'}$-模。
\end{lemma}

\begin{proof}
强烈建议读者自行证明 $U$ 的存在性。令 $\eta \in X$ 是一般点。
$E$ 在 $\eta$ 上的限制在 $D(\kappa(\eta))$ 中同构于一个微分全为零的
有限维向量空间有限复形 $V^\bullet$。令
$r_i = \dim_{\kappa(\eta)} V^i$。由各项为
$\mathcal{O}_X^{\oplus r_i}$、微分全为零的复形表示的
$D(\mathcal{O}_X)$ 中完美对象 $E'$，拉回到 $\eta$ 后与 $E$ 同构。
因此由《概形的导出范畴》，引理
\ref{perfect-lemma-descend-relatively-perfect}，存在 $\eta$ 的开邻域 $U$，
使 $E|_U$ 与 $E'|_U$ 同构。这证明第一个断言。第二个断言由第一个断言和
引理 \ref{flat-lemma-blowup-complex} 推出，因为整概形 $X$ 中任意非空开集都
概形论稠密。
\end{proof}

\begin{remark}
\label{flat-remark-when-you-have-a-complex}
设 $X$ 是概形，$E \in D(\mathcal{O}_X)$ 是完美对象，并且对所有
$i \in \mathbf{Z}$，$H^i(E)$ 都是 Tor 维数 $\leq 1$ 的完美
$\mathcal{O}_X$-模。该性质有时可将关于 $E$ 的问题归结为关于
$H^i(E)$ 的问题。例如，设
$$
\mathcal{E}^a \xrightarrow{d^a} \ldots
\xrightarrow{d^{b - 2}} \mathcal{E}^{b - 1}
\xrightarrow{d^{b - 1}} \mathcal{E}^b
$$
是表示 $E$ 的有限局部自由 $\mathcal{O}_X$-模有界复形。则对所有 $i$，
$\Im(d^i)$ 和 $\Ker(d^i)$ 都是有限局部自由 $\mathcal{O}_X$-模。
事实上，归纳假设对所有大于 $i$ 的指标均已知该结论。首先利用短正合列
$$
0 \to \Im(d^i) \to \Ker(d^{i + 1}) \to H^{i + 1}(E) \to 0
$$
以及 $H^{i + 1}(E)$ 是 Tor 维数 $\leq 1$ 的完美模这一假设，可知
$\Im(d^i)$ 有限局部自由。对短正合列
$$
0 \to \Ker(d^i) \to \mathcal{E}^i \to \Im(d^i) \to 0
$$
再次使用相同论证，得到 $\Ker(d^i)$ 有限局部自由。因此，特异三角
$$
\tau_{\leq k - 1}E \to \tau_{\leq k}E \to H^k(E)[-k] \to
(\tau_{\leq k - 1}E)[1]
$$
由如下有限局部自由模有界复形的短正合列表示：
$$
\begin{matrix}
& &
& &
& &
0 \\
& &
& &
& &
\downarrow \\
\mathcal{E}^a & \to &
\ldots & \to &
\mathcal{E}^{k - 2} & \to &
\Ker(d^{k - 1}) \\
\downarrow & &
& &
\downarrow & &
\downarrow \\
\mathcal{E}^a & \to &
\ldots & \to &
\mathcal{E}^{k - 2} & \to &
\mathcal{E}^{k - 1} & \to &
\Ker(d^k) \\
& &
& &
& &
\downarrow & &
\downarrow \\
& &
& &
& &
\Im(d^{k - 1}) & \to &
\Ker(d^k) \\
& &
& &
& &
\downarrow \\
& &
& &
& &
0
\end{matrix}
$$
这里各行是复形，图中省略了“显然”的零项。
\end{remark}





\section{完美模的爆破}
\label{flat-section-blowup-perfect-modules}

\noindent
本节试图寻找 Tor 维数 $\leq 1$ 的完美模在爆破后的正规形。
我们只取得了部分成功。

\begin{lemma}
\label{flat-lemma-blowup-pd1-derived}
设 $X$ 是概形，$\mathcal{F}$ 是 Tor 维数 $\leq 1$ 的完美
$\mathcal{O}_X$-模。对任意爆破 $b : X' \to X$，有
$Lb^*\mathcal{F} = b^*\mathcal{F}$，并且 $b^*\mathcal{F}$ 是
Tor 维数 $\leq 1$ 的完美 $\mathcal{O}_X$-模。
% P05-EN-CH38-0109: source 13228 and 13256 retain \mathcal{O}_X although the modules live on X'.
\end{lemma}

\begin{proof}
可以假设 $X = \Spec(A)$ 仿射，并且对应于 $\mathcal{F}$ 的 $A$-模 $M$
有一个表现
$$
0 \to A^{\oplus m} \to A^{\oplus n} \to M \to 0
$$
设 $I \subset A$ 是理想且 $a \in I$。回顾仿射爆破代数
$A[\frac{I}{a}]$ 是 $A_a$ 的子环。由于局部化正合，
$A_a^{\oplus m} \to A_a^{\oplus n}$ 是单射。因此
$A[\frac{I}{a}]^{\oplus m} \to A[\frac{I}{a}]^{\oplus n}$
也是单射。这证明了引理。
\end{proof}

\begin{lemma}
\label{flat-lemma-blowup-pd1}
设 $X$ 是概形，$\mathcal{F}$ 是 Tor 维数 $\leq 1$ 的完美
$\mathcal{O}_X$-模。设 $U \subset X$ 是概形论稠密的开集，使得
$\mathcal{F}|_U$ 是常秩 $r$ 的有限局部自由模。则存在 $U$-容许爆破
$b : X' \to X$，以及典范短正合列
% P05-EN-CH38-0110: source 13248 and 13370 add two further unhyphenated “scheme theoretically dense” loci.
$$
0 \to \mathcal{K} \to b^*\mathcal{F} \to \mathcal{Q} \to 0
$$
其中 $\mathcal{Q}$ 是秩 $r$ 的有限局部自由模，而 $\mathcal{K}$ 是
Tor 维数 $\leq 1$ 的完美 $\mathcal{O}_X$-模，且其在 $U$ 上的限制为零。
\end{lemma}

\begin{proof}
结合 《因子》，引理 \ref{divisors-lemma-blowup-fitting-ideal} 和
引理 \ref{flat-lemma-blowup-pd1-derived} 即得。
\end{proof}

\begin{lemma}
\label{flat-lemma-canonical-blowup-torsion-pd1}
设 $X$ 是概形，$\mathcal{F}$ 是 Tor 维数 $\leq 1$ 的完美
$\mathcal{O}_X$-模。设 $U \subset X$ 是满足 $\mathcal{F}|_U = 0$
的开集。则存在 $U$-容许爆破
$$
b : X' \to X
$$
使得 $\mathcal{F}' = b^*\mathcal{F}$ 配有两个典范的局部有限滤过
$$
0 = F^0 \subset F^1 \subset F^2 \subset \ldots \subset \mathcal{F}'
\quad\text{且}\quad
\mathcal{F}' = F_1 \supset F_2 \supset F_3 \supset \ldots \supset 0
$$
并且对每个 $n \geq 1$，存在有效 Cartier 除子 $D_n \subset X'$，使得
$$
F^i/F^{i - 1}
\quad\text{且}\quad
F_i/F_{i + 1}
$$
在 $D_i$ 上是秩 $i$ 的有限局部自由模。
\end{lemma}

\begin{proof}
选取仿射开集 $V \subset X$，使得存在某个 $n$ 和某个矩阵 $A$，以及表现
$$
0 \to \mathcal{O}_V^{\oplus n} \xrightarrow{A} \mathcal{O}_V^{\oplus n} \to
\mathcal{F} \to 0
$$
我们将沿所有 $k \geq 0$ 的 Fitting 理想
$\text{Fit}_k(\mathcal{F})$ 之积作爆破。这是有意义的，因为在上述仿射
情形中，对 $k > n$ 有
$\text{Fit}_k(\mathcal{F})|_V = \mathcal{O}_V$。显然，这是
$U$-容许爆破。由《因子》，引理
\ref{divisors-lemma-blowing-up-two-ideals}，在 $X'$ 上各理想
$\text{Fit}_k(\mathcal{F})$ 都可逆。因此归结到下一段讨论的情形。

\medskip\noindent
假设对 $k \geq 0$，$\text{Fit}_k(\mathcal{F})$ 都是可逆理想。
若 $E_k \subset X$ 是由 $\text{Fit}_k(\mathcal{F})$ 定义的有效
Cartier 除子（$k \geq 0$），则引理陈述中的有效 Cartier 除子 $D_k$ 将满足
$$
E_k = D_{k + 1} + 2 D_{k + 2} + 3 D_{k + 3} + \ldots
$$
这是有意义的，因为族 $D_k$ 将是局部有限的。此外，该等式唯一确定
有效 Cartier 除子 $D_k$，因此只须在局部构造 $D_k$。

\medskip\noindent
选取仿射开集 $V \subset X$ 以及如上的 $\mathcal{F}|_V$ 表现。
我们对表现中的整数 $n$ 作归纳，构造这些除子和滤过。对 $k > n$，令
$D_k|_V = \emptyset$，并令 $D_n|V = E_{n - 1}|_V$。
% P05-EN-CH38-0111: source 13323 omits the restriction subscript in D_n|V.
缩小 $V$ 后，可以假设 $\text{Fit}_{n - 1}(\mathcal{F})|_V$
由单个非零因子 $f \in \Gamma(V, \mathcal{O}_V)$ 生成。
由于 $\text{Fit}_{n - 1}(\mathcal{F})|_V$ 是由 $A$ 的各项生成的理想，
存在 $\Gamma(V, \mathcal{O}_V)$ 中的矩阵 $A'$，使得 $A = fA'$。
由短正合列在 $V$ 上定义 $\mathcal{F}'$：
$$
0 \to \mathcal{O}_V^{\oplus n} \xrightarrow{A'} \mathcal{O}_V^{\oplus n} \to
\mathcal{F}' \to 0
$$
由于 $A'$ 的各项在 $\Gamma(V, \mathcal{O}_V)$ 中生成单位理想，
$\mathcal{F}'$ 在 $V$ 上局部有一个将 $n$ 减少 $1$ 的表现；见《代数》，
引理 \ref{algebra-lemma-add-trivial-complex}。还要注意
$f^{n - k}\text{Fit}_k(\mathcal{F}') = \text{Fit}_k(\mathcal{F})|_V$
对 $k = 0, \ldots, n$ 成立。因此对所有 $k$，
$\text{Fit}_k(\mathcal{F}')$ 都是可逆理想。由归纳可知，存在有效
Cartier 除子 $D'_k \subset V$，使得 $\mathcal{F}'$ 具有引理陈述中的
两个典范滤过。然后对 $k = 1, \ldots, n - 1$，令
$D_k|_V = D'_k$。注意，在这些选择下，上述有效 Cartier 除子间的等式成立。
最后构造滤过。事实上，有短正合列
$$
0 \to
\mathcal{O}_{D_n \cap V}^{\oplus n} \to
\mathcal{F} \to \mathcal{F}' \to 0
\quad\text{且}\quad
0 \to
\mathcal{F}' \to \mathcal{F} \to
\mathcal{O}_{D_n \cap V}^{\oplus n} \to 0
$$
它们来自 $A$ 的两个分解 $A = A'f = f A'$。这些序列是典范的，
因为第一个序列中的子模是
$\Ker(f : \mathcal{F} \to \mathcal{F})$，第二个序列中的商模是
$\Coker(f : \mathcal{F} \to \mathcal{F})$。
\end{proof}

\begin{lemma}
\label{flat-lemma-blowup-map-pd1}
设 $X$ 是概形。设 $\varphi : \mathcal{F} \to \mathcal{G}$ 是
Tor 维数 $\leq 1$ 的完美 $\mathcal{O}_X$-模之间的同态。
设 $U \subset X$ 是概形论稠密的开集，使得
$\mathcal{F}|_U = 0$ 且 $\mathcal{G}|_U = 0$。则存在 $U$-容许爆破
$b : X' \to X$，使得 $b^*\varphi$ 的核、像和余核都是
Tor 维数 $\leq 1$ 的完美 $\mathcal{O}_{X'}$-模。
% P05-EN-CH38-0112: source 13369 misspells “homomorphism” as “homorphism”.
\end{lemma}

\begin{proof}
这些假设表明 $D(\mathcal{O}_X)$ 的对象
$(\mathcal{F} \to \mathcal{G})$ 是完美的。因此由 引理
\ref{flat-lemma-blowup-complex} 和 \ref{flat-lemma-blowup-pd1-derived}
（后者用来考察拉回后的复形形状），得到一个对余核和核都适用的
$U$-容许爆破。像是到余核之映射的核，因而也是 Tor 维数
$\leq 1$ 的完美模。
\end{proof}











\section{Berthelot 与 Ogus 引入的算子}
\label{flat-section-eta}

\noindent
请先阅读 《上同调》，第 \ref{cohomology-section-eta} 节。

\medskip\noindent
设 $X$ 是概形，$D \subset X$ 是有效 Cartier 除子。令
$\mathcal{I} = \mathcal{I}_D \subset \mathcal{O}_X$ 是 $D$ 的理想层；
见《因子》，第 \ref{divisors-section-effective-Cartier-invertible} 节。显然，可以将
《上同调》，第 \ref{cohomology-section-eta} 节的讨论应用于
$X$ 和 $\mathcal{I}$。

\begin{lemma}
\label{flat-lemma-eta-stalks}
设 $X$ 是概形，$D \subset X$ 是理想层为
$\mathcal{I} \subset \mathcal{O}_X$ 的有效 Cartier 除子。设
$\mathcal{F}^\bullet$ 是拟凝聚 $\mathcal{O}_X$-模复形，并且对所有 $i$，
$\mathcal{F}^i$ 都无 $\mathcal{I}$-挠。则
$\eta_\mathcal{I}\mathcal{F}^\bullet$ 是拟凝聚 $\mathcal{O}_X$-模复形。
此外，若 $U = \Spec(A) \subset X$ 是仿射开集且 $D \cap U = V(f)$，
则 $\eta_f(\mathcal{F}^\bullet(U))$ 典范同构于
$(\eta_\mathcal{I}\mathcal{F}^\bullet)(U)$。
\end{lemma}

\begin{proof}
从略。
\end{proof}

\begin{lemma}
\label{flat-lemma-Leta}
设 $X$ 是概形，$D \subset X$ 是理想层为
$\mathcal{I} \subset \mathcal{O}_X$ 的有效 Cartier 除子。《上同调》，
引理 \ref{cohomology-lemma-Leta} 的函子
$L\eta_\mathcal{I} : D(\mathcal{O}_X) \to D(\mathcal{O}_X)$
将 $D_\QCoh(\mathcal{O}_X)$ 映入自身。此外，若
$X = \Spec(A)$ 仿射且 $D = V(f)$，则 《代数进阶》，引理
\ref{more-algebra-lemma-Leta} 所定义的 $D(A)$ 上函子 $L\eta_f$，
与 $D_\QCoh(\mathcal{O}_X)$ 上函子 $L\eta_\mathcal{I}$，经由
《概形的导出范畴》，引理
\ref{perfect-lemma-affine-compare-bounded} 的等价相对应。
\end{lemma}

\begin{proof}
从略。
\end{proof}





\section{复形的爆破，II}
\label{flat-section-blowup-complexes-II}

\noindent
本节的材料将用于在 第 \ref{flat-section-blowup-complexes-III} 节
中构造 Macpherson 图构造的一个版本。
 
\begin{situation}
\label{flat-situation-complex-and-divisor}
这里 $X$ 是概形，$D \subset X$ 是理想层为
$\mathcal{I} \subset \mathcal{O}_X$ 的有效 Cartier 除子，而
$M$ 是 $D(\mathcal{O}_X)$ 的完美对象。
\end{situation}

\noindent
令 $(X, D, M)$ 是 情形 \ref{flat-situation-complex-and-divisor}
中的三元组。考虑满足如下条件的仿射开集
$U = \Spec(A) \subset X$：
\begin{enumerate}
\item 对某个非零因子 $f \in A$，有 $D \cap U = V(f)$；
\item 存在表示 $M|_U$ 的有限自由 $A$-模有界复形 $M^\bullet$
（经由《概形的导出范畴》，引理
\ref{perfect-lemma-affine-compare-bounded} 的等价）。
\end{enumerate}
称 $(U, A, f, M^\bullet)$ 是 $(X, D, M)$ 的一个{\it 仿射图}。
考虑《代数进阶》，第 \ref{more-algebra-section-perfect-eta} 节中定义的理想
$I_i(M^\bullet, f) \subset A$。若对每个 $x \in D$，存在满足
$x \in U$ 的仿射图 $(U, A, f, M^\bullet)$，使得对所有 $i \in \mathbf{Z}$，
$I_i(M^\bullet, f)$ 都是主理想，则称 $(X, S, M)$ 是{\it 良三元组}。
% P05-EN-CH38-0113: source 13482 changes D to the undefined S in the good-triple notation.

\begin{lemma}
\label{flat-lemma-pullback-triple-ideals-good}
在 情形 \ref{flat-situation-complex-and-divisor} 中，设 $h : Y \to X$
是概形态射，使得 $D$ 的拉回 $E = h^{-1}D$ 有定义（《因子》，
定义 \ref{divisors-definition-pullback-effective-Cartier-divisor}）。
设 $(U, A, f, M^\bullet)$ 是 $(X, D, M)$ 的仿射图，
$V = \Spec(B) \subset Y$ 是满足 $h(V) \subset U$ 的仿射开集。
以 $g \in B$ 表示 $f \in A$ 的像。则
% P05-EN-CH38-0114: source 13493--13494 uses “Let ... is” twice instead of “Let ... be”.
\begin{enumerate}
\item $(V, B, g, M^\bullet \otimes_A B)$ 是 $(Y, E, Lh^*M)$ 的仿射图；
\item 在 $B$ 中有 $I_i(M^\bullet, f)B = I_i(M^\bullet \otimes_A B, g)$；
\item 若 $(X, D, M)$ 是良三元组，则 $(Y, E, Lh^*M)$ 是良三元组。
\end{enumerate}
\end{lemma}

\begin{proof}
第一个断言由如下观察得到：$g$ 是 $B$ 中定义
$E \cap V \subset V$ 的非零因子；而 $M^\bullet \otimes_A B$ 表示
$M^\bullet \otimes_A^\mathbf{L} B$，因而由《概形的导出范畴》，
引理 \ref{perfect-lemma-quasi-coherence-pullback}，它表示 $M$ 在 $V$ 上的
拉回。(2) 由 (1) 和 《代数进阶》，引理
\ref{more-algebra-lemma-eta-base-change-pre} 推出。再结合 《代数进阶》，
引理 \ref{more-algebra-lemma-eta-base-change-pre}，可知引理的第二个断言成立。
% P05-EN-CH38-0115: source 13512--13516 cites the same lemma twice, concludes (2) twice, and never proves item (3).
\end{proof}

\begin{lemma}
\label{flat-lemma-good-triple}
设 $X, D, \mathcal{I}, M$ 如 情形
\ref{flat-situation-complex-and-divisor} 中所述。若 $(X, D, M)$ 是良三元组，
则 $L\eta_\mathcal{I}M$ 是 $D(\mathcal{O}_X)$ 的完美对象。
\end{lemma}

\begin{proof}
这是 《代数进阶》，引理
\ref{more-algebra-lemma-eta-locally-free} 的翻译。使用
引理 \ref{flat-lemma-Leta} 完成该翻译。
\end{proof}

\begin{lemma}
\label{flat-lemma-good-triple-bdd-loc-free}
设 $X, D, \mathcal{I}, M$ 如 情形
\ref{flat-situation-complex-and-divisor} 中所述。假设 $(X, D, M)$ 是良三元组。
若存在表示 $M$ 的有限局部自由 $\mathcal{O}_X$-模局部有界复形
$\mathcal{M}^\bullet$，则存在表示 $L\eta_\mathcal{I}M$ 的有限局部自由
$\mathcal{O}_{X'}$-模局部有界复形 $\mathcal{Q}^\bullet$。
% P05-EN-CH38-0116: source 13542 uses the undefined X'; the eta object still lives on X.
\end{lemma}

\begin{proof}
由《上同调》，引理 \ref{cohomology-lemma-Leta}，复形
$\mathcal{Q}^\bullet = \eta_\mathcal{I}\mathcal{M}^\bullet$
表示 $L\eta_\mathcal{I}M$。为核验该复形局部有界且由有限局部自由模组成，
可以在仿射局部工作，此时有界性显然。选取 $(X, D, M)$ 的仿射图
$(U, A, f, M^\bullet)$，使各理想 $I_i(M^\bullet, f)$ 都是主理想，
并使对每个 $i$，$\mathcal{M}^i|_U$ 都有限自由。由 $(X, D, M)$ 是良三元组
的假设，可以作此选择。写成
$N^i = \Gamma(U, \mathcal{M}^i|_U)$，得到有限自由 $A$-模有界复形
$N^\bullet$；它与复形 $M^\bullet$ 表示 $D(A)$ 中的同一对象
（由《概形的导出范畴》，引理
\ref{perfect-lemma-affine-compare-bounded}）。由《代数进阶》，引理
\ref{more-algebra-lemma-ideal-well-defined}，对所有 $i$，
$I_i(N^\bullet, f)$ 都是主理想。因此复形 $\eta_fN^\bullet$ 是有限局部自由
$A$-模有界复形。由 引理 \ref{flat-lemma-eta-stalks}，
$\mathcal{Q}^i|_U$ 是对应于 $\eta_fN^i$ 的拟凝聚
$\mathcal{O}_U$-模，结论得证。
\end{proof}

\begin{lemma}
\label{flat-lemma-complex-and-divisor-eta-pull}
在 情形 \ref{flat-situation-complex-and-divisor} 中，设 $h : Y \to X$
是概形态射，使拉回 $E = h^{-1}D$ 有定义。若 $(X, D, M)$ 是良三元组，则
$$
Lh^*(L\eta_\mathcal{I}M) = L\eta_\mathcal{J}(Lh^*M)
$$
在 $D(\mathcal{O}_Y)$ 中成立，其中 $\mathcal{J}$ 是 $E$ 的理想层。
\end{lemma}

\begin{proof}
这是 《代数进阶》，引理
\ref{more-algebra-lemma-eta-base-change} 的翻译。使用 引理
\ref{flat-lemma-eta-stalks} 和 \ref{flat-lemma-Leta} 完成该翻译。
\end{proof}

\begin{lemma}
\label{flat-lemma-complex-and-divisor-blowup-pre}
在情形 \ref{flat-situation-complex-and-divisor} 中，
存在唯一态射 $b : X' \to X$，使得
\begin{enumerate}
\item 拉回 $D' = b^{-1}D$ 有定义，且 $(X', D', M')$ 是良三元组，
其中 $M' = Lb^*M$；
\item 对任意概形态射 $h : Y \to X$，若拉回 $E = h^{-1}D$ 有定义且
$(Y, E, Lh^*M)$ 是良三元组，则 $h$ 唯一地通过 $b$ 分解。
\end{enumerate}
此外，对任意仿射图 $(U, A, f, M^\bullet)$，限制
$b^{-1}(U) \to U$ 是沿理想 $I_i(M^\bullet, f)$ 之积的爆破；而对任意
拟紧开集 $W \subset X$，限制
$b|_{b^{-1}(W)} : b^{-1}(W) \to W$ 是 $W \setminus D$-容许爆破。
\end{lemma}

\begin{proof}
证明的思路只是：对任意仿射图 $(U, A, f, M^\bullet)$，在局部沿积理想
$I_i(M^\bullet, f)$ 爆破 $X$。《代数进阶》，第 \ref{more-algebra-section-perfect-eta} 节 开头的几个引理表明，这是良定义的。
泛性质 (2) 随即由爆破的泛性质得到。以下给出细节。

\medskip\noindent
令 $U, A, f, M^\bullet$ 是 $(X, D, M)$ 的仿射图。除有限多个以外，
理想 $I_i(M^\bullet, f)$ 都等于 $A$，故可以考虑
$$
I = \prod\nolimits_i I_i(M^\bullet, f)
$$
这是 $A$ 的有限生成理想。以
$$
b_U : U' \to U
$$
表示 $U$ 沿 $I$ 的爆破。由《因子》，引理
\ref{divisors-lemma-blow-up-pullback-effective-Cartier}，
$b_U^{-1}(U \cap D)$ 有定义。回顾 $f^{r_i} \in I_i(M^\bullet, f)$，
故 $b_U$ 是 $(U \setminus D)$-容许爆破。由《因子》，引理
\ref{divisors-lemma-blowing-up-two-ideals}，对每个 $i$，态射 $b_U$
分解为 $U' \to U'_i \to U$，其中 $U'_i \to U$ 是沿
$I_i(M^\bullet, f)$ 的爆破，而 $U' \to U'_i$ 是另一个爆破。
因此 $I_i(M^\bullet, f)$ 在 $U'$ 上的拉回
$I_i(M^\bullet, f)\mathcal{O}_{U'}$ 是可逆理想层；见《因子》，引理
\ref{divisors-lemma-blow-up-pullback-effective-Cartier} 和
\ref{divisors-lemma-blowing-up-gives-effective-Cartier-divisor}。
故 $(U', b^{-1}D, Lb^*M|_U)$ 是良三元组；关于理想 $I_i(-,-)$
在拉回下的行为，见 引理 \ref{flat-lemma-pullback-triple-ideals-good}。
% P05-EN-CH38-0117: source 13638 uses global b inside the local b_U construction before b exists.
最后，我们断言 $b_U : U' \to U$ 具有引理陈述 (2) 中的泛性质。
事实上，设 $h : Y \to U$ 是概形态射，使得拉回
$E = h^{-1}(D \cap U)$ 有定义且 $(Y, E, Lh^*M)$ 是良三元组。
则 $Y$ 由仿射图 $(V, B, g, N^\bullet)$ 覆盖，并且对每个 $i$，
$I_i(N^\bullet, g)$ 都是可逆理想。由于 $g$ 与 $f$ 在 $B$ 中的像
都截出有效 Cartier 除子 $E \cap V$，二者相差一个单位。因此由
《代数进阶》，引理 \ref{more-algebra-lemma-eta-change-unit}，
可以假设 $g$ 是 $f$ 的像。再由《代数进阶》，引理
\ref{more-algebra-lemma-ideal-well-defined}，$I_i(N^\bullet, g)$
作为 $B$-模同构于 $I_i(M^\bullet \otimes_A B, g)$。因此
$I_i(M^\bullet \otimes_A B, g) = I_i(M^\bullet, f)B$
（引理 \ref{flat-lemma-pullback-triple-ideals-good}）是可逆 $B$-模。
故理想 $IB$ 可逆，从而 $I\mathcal{O}_Y$ 可逆。于是由《因子》，引理
\ref{divisors-lemma-universal-property-blowing-up}，得到 $h$ 通过
$b_U$ 的唯一分解。

\medskip\noindent
令 $\mathcal{B}$ 是如下仿射开集 $U \subset X$ 的集合：存在
$(X, D, M)$ 的仿射图 $(U, A, f, M^\bullet)$。则 $\mathcal{B}$
是 $X$ 的拓扑的一组基；细节从略。对 $U \in \mathcal{B}$，有上述构造的
态射 $b_U : U' \to U$，它在 $U$ 上满足该泛性质。若
$U_1 \subset U_2 \subset X$ 都属于 $\mathcal{B}$，则由泛性质，
$b_{U_1} : U'_1 \to U_1$ 典范同构于
$$
b_{U_2}|_{b_{U_2}^{-1}(U_1)} : b_{U_2}^{-1}(U_1) \longrightarrow U_1
$$
换言之，得到 $U_1$ 上同构
$U'_1 \to b_{U_2}^{-1}(U_1)$。这些同构满足上链条件（仍由泛性质），
故由《构造》，引理 \ref{constructions-lemma-relative-glueing}，
得到态射 $b : X' \to X$，其在 $\mathcal{B}$ 中每个 $U$ 上的限制同构于
$U' \to U$。于是态射 $b : X' \to X$ 满足引理陈述中的性质 (1) 和 (2)，
因为这些性质可在局部核验（细节从略）。

\medskip\noindent
还须证明引理的最后一个断言。设 $W \subset X$ 是拟紧开集。选取有限覆盖
$W = U_1 \cup \ldots \cup U_T$，使得对每个 $1 \leq t \leq T$，
存在仿射图 $(U_t, A_t, f_t, M_t^\bullet)$。下文将使用如下事实：
对任意仿射开集 $V = \Spec(B) \subset U_t \cap U_{t'}$，
(a) $f_t$ 与 $f_{t'}$ 在 $B$ 中的像相差一个单位；且
(b) 复形 $M_t^\bullet \otimes_A B$ 与 $M_{t'} \otimes_A B$
定义 $D(B)$ 中同构的对象。对 $i \in \mathbf{Z}$，令
% P05-EN-CH38-0119: source 13691 tensors both chart complexes over undefined A instead of A_t and A_{t'}.
$$
N_i = \max\nolimits_{t = 1, \ldots, T}
\left(\sum\nolimits_{j \geq i} (-1)^{j - i} rk(M_t^j)\right)
$$
则 $N_t - \sum\nolimits_{j \geq i} (-1)^{j - i} rk(M_t^j) \geq 0$，
因而可以考虑理想
% P05-EN-CH38-0118: source 13697 uses N_t although the maximum was defined as N_i.
$$
I_{t, i} =
f_t^{N_i - \sum\nolimits_{j \geq i} (-1)^{j - i} rk(M_t^j)}
I_i(M_t^\bullet, f_t)
\subset A_t
$$
由《代数进阶》，引理 \ref{more-algebra-lemma-eta-change-unit} 和
\ref{more-algebra-lemma-ideal-well-defined}，理想 $I_{t, i}$ 黏合成
拟凝聚有限型理想 $\mathcal{I}_i \subset \mathcal{O}_W$。此外，除有限多个
以外，这些理想都等于 $\mathcal{O}_W$。显然，上述构造的态射
$X' \to X$ 在 $W$ 上限制为沿理想 $\mathcal{I}_i$ 之积的爆破。
证明完成。
\end{proof}

\begin{lemma}
\label{flat-lemma-complex-and-divisor-blowup}
在 情形 \ref{flat-situation-complex-and-divisor} 中，令 $b : X' \to X$
是 引理 \ref{flat-lemma-complex-and-divisor-blowup-pre} 的态射。考虑理想层为
$\mathcal{I}' \subset \mathcal{O}_{X'}$ 的有效 Cartier 除子
$D' = b^{-1}D$。则 $Q = L\eta_{\mathcal{I}'}Lb^*M$ 是
$D(\mathcal{O}_{X'})$ 的完美对象。
\end{lemma}

\begin{proof}
由 引理 \ref{flat-lemma-complex-and-divisor-blowup-pre} 和
\ref{flat-lemma-good-triple} 得到。
\end{proof}

\begin{lemma}
\label{flat-lemma-complex-and-divisor-blowup-base-change}
在 情形 \ref{flat-situation-complex-and-divisor} 中，设 $h : Y \to X$
是概形态射，使拉回 $E = h^{-1}D$ 有定义。令 $b : X' \to X$ 和
$c : Y' \to Y$ 分别是 引理
\ref{flat-lemma-complex-and-divisor-blowup-pre} 对 $D \subset X$ 与 $M$、
以及对 $E \subset Y$ 与 $Lh^*M$ 所构造的态射。则 $Y'$ 是 $Y$ 关于
$b : X' \to X$ 的严格变换（精确表述见证明），并且
$$
L\eta_{\mathcal{J}'}L(h \circ c)^*M = L(Y' \to X')^*Q
$$
其中 $Q = L\eta_{\mathcal{I}'}Lb^*M$ 如
引理 \ref{flat-lemma-complex-and-divisor-blowup} 中所述。特别地，若
$(Y, E, Lh^*M)$ 是良三元组，且 $k : Y \to X'$ 是满足
$h = b \circ k$ 的唯一态射，则
$L\eta_\mathcal{J}Lh^*M = Lk^*Q$。
\end{lemma}

\begin{proof}
记 $E' = c^{-1}E$。则 $(Y', E', L(h \circ c)^*M)$ 是良三元组。
故由 引理 \ref{flat-lemma-complex-and-divisor-blowup-pre} 的泛性质，
存在唯一态射
$$
h' : Y' \longrightarrow X'
$$
使得 $b \circ h' = h \circ c$。特别地，存在态射
$(h', c) : Y' \to X' \times_X Y$。我们断言：给定拟紧开集
$W \subset X$，使 $b^{-1}(W) \to W$ 是爆破，该态射将 $Y'|_W$
等同于 $Y_W$ 关于 $b^{-1}(W) \to W$ 的严格变换。该断言是在 $W$
上的局部问题，因此可以在仿射图上证明。下一段即作此事。

\medskip\noindent
令 $(U, A, f, M^\bullet)$ 是 $(X, D, M)$ 的仿射图。回顾
引理 \ref{flat-lemma-complex-and-divisor-blowup} 的证明，$b : X' \to X$
在 $U$ 上的限制是 $U = \Spec(A)$ 沿理想 $I_i(M^\bullet, f)$ 之积的爆破。
% P05-EN-CH38-0120: source 13767 cites the later consequence lemma, while the local blowup construction is in lemma-complex-and-divisor-blowup-pre.
若 $V = \Spec(B) \subset Y$ 是满足 $h(V) \subset U$ 的任意仿射开集，
则 $(V, B, g, M^\bullet \otimes_A B)$ 是 $(Y, E, Lh^*M)$ 的仿射图，
其中 $g \in B$ 是 $f$ 的像；见
引理 \ref{flat-lemma-pullback-triple-ideals-good}。因此 $c : Y' \to Y$
在 $V$ 上的限制是沿理想 $I_i(M^\bullet, f)B$ 之积的爆破；换言之，
$c : Y' \to Y$ 在 $h^{-1}(U)$ 上是 $h^{-1}(U)$ 沿理想
$\prod I_i(M^\bullet, f) \mathcal{O}_{h^{-1}(U)}$ 的爆破。由于对严格变换
也有相同描述，关于严格变换的断言成立。

\medskip\noindent
由此，等式
$L\eta_{\mathcal{J}'}L(h \circ c)^*M = L(Y' \to X')^*Q$
由 引理 \ref{flat-lemma-complex-and-divisor-eta-pull} 得到。
最后一个断言是其特例（即 $c = \text{id}_Y$ 且 $k = h'$ 的情形）。
\end{proof}

\begin{lemma}
\label{flat-lemma-complex-and-divisor-blowup-good}
在 情形 \ref{flat-situation-complex-and-divisor} 中，令 $W \subset X$
是 $M$ 的上同调层在其上局部自由的最大开子概形。则 引理
\ref{flat-lemma-complex-and-divisor-blowup-pre} 的态射 $b : X' \to X$
在 $W$ 上是同构。
\end{lemma}

\begin{proof}
这是因为，对任意满足 $U \subset W$ 的仿射图
$(U, A, f, M^\bullet)$，由《代数进阶》，引理
\ref{more-algebra-lemma-eta-cohomology-free}，各
$I_i(M^\bullet, f)$ 在局部由 $f$ 的某个幂生成。由于 $f$ 是非零因子，
爆破 $b^{-1}(U) \to U$ 是同构。
\end{proof}

\begin{lemma}
\label{flat-lemma-complex-and-divisor-blowup-good-T}
设 $X, D, \mathcal{I}, M$ 如 情形
\ref{flat-situation-complex-and-divisor} 中所述。若 $(X, D, M)$ 是良三元组，
则存在有限表现的闭浸入
$$
i : T \longrightarrow D
$$
该态射具有如下性质：
\begin{enumerate}
\item $T$ 在概形论意义下包含 $D \cap W$，其中 $W \subset X$
是 $M$ 的上同调层在其上局部自由的最大开集；
\item $Li^*L\eta_\mathcal{I}M$ 的上同调层局部自由；
\item 对任意像满足 $x = i(t) \in W$ 的点 $t \in T$，
$H^i(M)_x$ 作为 $\mathcal{O}_{X, x}$-模的秩，等于
$H^i(Li^*L\eta_\mathcal{I}M)_t$ 作为 $\mathcal{O}_{T, t}$-模的秩。
\end{enumerate}
% P05-EN-CH38-0121: source 13820, 13870, 13898, 13949, and 14056 add five further unhyphenated “scheme theoretically” loci.
\end{lemma}

\begin{proof}
令 $(U, A, f, M^\bullet)$ 是 $(X, D, M)$ 的仿射图，使得对所有
$i \in \mathbf{Z}$，$I_i(M^\bullet, f)$ 都是主理想。将
$T \cap U \subset D \cap U$ 定义为由理想
$$
J(M^\bullet, f) = \sum J_i(M^\bullet, f) \subset A/fA
$$
定义的闭子概形；该理想在《代数进阶》，引理
\ref{more-algebra-lemma-eta-vanishing-beta-plus-pre} 和
\ref{more-algebra-lemma-eta-vanishing-beta-plus} 中研究。按第二个引理的
表述，$T \cap U \to D \cap U$ 由其中研究的环同态
$A/fA \to C$ 给出。由于 $(X, D, M)$ 是良三元组，可以用这种形式的
仿射图覆盖 $X$；由两个引理中的第一个，这些构造可以黏合。因此得到闭子概形
$T \subset D$，它在上述良仿射图上由理想 $J(M^\bullet, f)$ 给出。
性质 (1) 和 (2) 随即由第二个引理得到；细节从略。给读者一个有用的小观察：
由《代数进阶》，引理 \ref{more-algebra-lemma-eta-locally-free}，
$\eta_fM^\bullet$ 是局部自由模复形，故
$Li^*L\eta_\mathcal{I}M|_{T \cap U}$ 由 $C$-模复形
$\eta_fM^\bullet \otimes_A C$ 表示。关于秩的断言 (3) 由《上同调》，
引理 \ref{cohomology-lemma-eta-cohomology-locally-free} 得到。
\end{proof}

\begin{lemma}
\label{flat-lemma-complex-and-divisor-blowup-T}
在 情形 \ref{flat-situation-complex-and-divisor} 中，令 $b : X' \to X$
和 $D'$ 如 引理 \ref{flat-lemma-complex-and-divisor-blowup-pre} 中所述，
并令 $Q = L\eta_{\mathcal{I}'}Lb^*M$ 如
引理 \ref{flat-lemma-complex-and-divisor-blowup} 中所述。令 $W \subset X$
是 $M$ 的上同调模在其上局部自由的最大开集。则存在有限表现的闭浸入
$i : T \to D'$，使得
\begin{enumerate}
\item 在概形论意义下 $D' \cap b^{-1}(W) \subset T$；
\item $Li^*Q$ 具有局部自由上同调层；
\item 对映到 $w \in W$ 的 $t \in T$，$H^i(Li^*Q)_t$ 作为
$\mathcal{O}_{T, t}$-模的秩等于 $H^i(M)_x$ 作为
$\mathcal{O}_{X, x}$-模的秩。
\end{enumerate}
% P05-EN-CH38-0122: source 13872 introduces w but item (3) uses undefined x in both H^i(M)_x and O_{X,x}.
\end{lemma}

\begin{proof}
引理 \ref{flat-lemma-complex-and-divisor-blowup-good} 表明 $b$ 在 $W$ 上
是同构。因此 $b^{-1}(W) \subset X'$ 包含于 $Lb^*M$ 具有局部自由
上同调层的最大开集 $W' \subset X'$。于是本引理的各个断言立即由
引理 \ref{flat-lemma-complex-and-divisor-blowup-good-T} 对
$(X', D', Lb^*M)$ 的应用以及陈述中提到的其他引理得到。
\end{proof}

\begin{lemma}
\label{flat-lemma-complex-and-divisor-blowup-T-ranks}
在 情形 \ref{flat-situation-complex-and-divisor} 中，令 $b : X' \to X$、
$D' \subset X'$ 和 $Q$ 如 引理
\ref{flat-lemma-complex-and-divisor-blowup} 中所述。令
$\rho = (\rho_i)_{i \in \mathbf{Z}}$ 是整数族。令
$W(\rho) \subset X$ 是满足如下性质的最大开子概形：对所有 $i$，
$H^i(M)$ 都是秩 $\rho_i$ 的局部自由模。令 $i : T \to D'$ 如
引理 \ref{flat-lemma-complex-and-divisor-blowup-T} 中所述。则存在既开又闭的
子概形 $T(\rho) \subset T$，它在概形论意义下包含
$D' \cap b^{-1}(W(\rho))$，并且对所有 $i$，
$H^i(Li^*Q|_{T(\rho)})$ 都是秩 $\rho_i$ 的局部自由模。
\end{lemma}

\begin{proof}
令 $T(\rho) \subset T$ 是如下既开又闭的子概形：对所有 $i$，
$H^i(Li^*Q)$ 的秩都是 $\rho_i$。则该断言立即由 引理
\ref{flat-lemma-complex-and-divisor-blowup-T} 中关于上同调模秩的断言得到。
\end{proof}

\begin{lemma}
\label{flat-lemma-complex-and-divisor-blowup-T-represented-bdd-loc-free}
在 情形 \ref{flat-situation-complex-and-divisor} 中，令 $b : X' \to X$、
$D' \subset X'$ 和 $Q$ 如 引理
\ref{flat-lemma-complex-and-divisor-blowup} 中所述。若存在表示 $M$ 的
有限局部自由 $\mathcal{O}_X$-模局部有界复形 $\mathcal{M}^\bullet$，
则存在表示 $Q$ 的有限局部自由 $\mathcal{O}_{X'}$-模局部有界复形
$\mathcal{Q}^\bullet$。
\end{lemma}

\begin{proof}
回顾 $Q = L\eta_{\mathcal{I}'}Lb^*M$，其中 $\mathcal{I}'$ 是有效
Cartier 除子 $D'$ 的理想层。有限局部自由 $\mathcal{O}_{X'}$-模
局部有界复形 $(\mathcal{M}')^\bullet = b^*\mathcal{M}^\bullet$
表示 $Lb^*M$。因此本引理由
引理 \ref{flat-lemma-good-triple-bdd-loc-free} 得到。
\end{proof}

\begin{lemma}
\label{flat-lemma-complex-and-divisor-derived}
设 $X$ 是概形，$D \subset X$ 是有效 Cartier 除子，且
$M \in D(\mathcal{O}_X)$ 是完美对象。令 $W \subset X$ 是上同调层
$H^i(M)$ 在其上局部自由的最大开集。存在固有态射
$b : X' \longrightarrow X$ 和 $D(\mathcal{O}_{X'})$ 中的对象 $Q$，
具有如下性质：
\begin{enumerate}
\item $b : X' \to X$ 在 $X \setminus D$ 上是同构；
\item $b : X' \to X$ 在 $W$ 上是同构；
\item $D' = b^{-1}D$ 是有效 Cartier 除子；
\item $Q = L\eta_{\mathcal{I}'}Lb^*M$，其中 $\mathcal{I}'$ 是 $D'$ 的理想层；
\item $Q$ 是 $D(\mathcal{O}_{X'})$ 的完美对象；
\item 存在有限表现的闭浸入 $i : T \to D'$，使得
\begin{enumerate}
\item 在概形论意义下 $D' \cap b^{-1}(W) \subset T$；
\item $Li^*Q$ 具有有限局部自由上同调层；
\item 对像为 $w \in W$ 的 $t \in T$，$H^i(Li^*Q)_t$ 作为
$\mathcal{O}_{T, t}$-模的秩等于 $H^i(M)_x$ 作为
$\mathcal{O}_{X, x}$-模的秩；
\end{enumerate}
\item 对任意 $(X, D, M)$ 的仿射图 $(U, A, f, M^\bullet)$，
$b$ 在 $U$ 上的限制是 $U = \Spec(A)$ 沿理想
$I = \prod I_i(M^\bullet, f)$ 的爆破；
\item 对任意 $(X', D', Lb^*N)$ 的仿射图 $(V, B, g, N^\bullet)$，
若 $I_i(N^\bullet, g)$ 是主理想，则
% P05-EN-CH38-0123: source 13958 uses Lb^*N in the triple although N^\bullet is the chart complex and the pulled-back object is Lb^*M.
\begin{enumerate}
\item $Q|_V$ 对应于 $\eta_gN^\bullet$；
\item $T \subset V \cap D'$ 对应于 《代数进阶》，引理
\ref{more-algebra-lemma-eta-vanishing-beta-plus} 中研究的理想
$J(N^\bullet, g) = \sum J_i(N^\bullet, g) \subset B/gB$
。
% P05-EN-CH38-0128: source 13962 places all of T inside an arbitrary chart V; the local closed subscheme should be T \cap V.
\end{enumerate}
\item 若 $M$ 可由有限局部自由 $\mathcal{O}_X$-模局部有界复形表示，
则 $Q$ 可由有限局部自由 $\mathcal{O}_{X'}$-模有界复形表示。
% P05-EN-CH38-0124: source 13967--13970 strengthens “locally bounded” to “bounded”, unlike the lemma it collects.
\end{enumerate}
\end{lemma}

\begin{proof}
该断言汇集了 引理 \ref{flat-lemma-pullback-triple-ideals-good}、
\ref{flat-lemma-good-triple},
\ref{flat-lemma-complex-and-divisor-eta-pull},
\ref{flat-lemma-complex-and-divisor-blowup-pre},
\ref{flat-lemma-complex-and-divisor-blowup},
\ref{flat-lemma-complex-and-divisor-blowup-base-change},
\ref{flat-lemma-complex-and-divisor-blowup-good},
\ref{flat-lemma-complex-and-divisor-blowup-good-T},
\ref{flat-lemma-complex-and-divisor-blowup-T} 以及
\ref{flat-lemma-complex-and-divisor-blowup-T-represented-bdd-loc-free}
中得到的信息。
\end{proof}










\section{复形的爆破，III}
\label{flat-section-blowup-complexes-III}

\noindent
本节给出 \cite[Section 18.1]{F} 中 Macpherson 图构造版本的一个
“代数版本”。
% P05-EN-CH38-0127: source 14001 repeats “version” in “an algebra version of the version”.

\medskip\noindent
设 $X$ 是概形，$E$ 是 $D(\mathcal{O}_X)$ 的完美对象。令
$U \subset X$ 是使得 $E|_U$ 具有局部自由上同调层的最大开子概形。

\medskip\noindent
考虑交换图
$$
\xymatrix{
\mathbf{A}^1_X \ar[r] \ar[rd] &
\mathbf{P}^1_X \ar[d]^p &
(\mathbf{P}^1_X)_\infty \ar[l] \ar[ld] \\
& X \ar@/_1em/[ur]_\infty
}
$$
回顾 $\mathbf{A}^1 = D_+(T_0)$ 是 $\mathbf{P}^1$ 的第一个标准仿射开集，
而 $\infty = V_+(T_0)$ 是其补有效 Cartier 除子；上图是这些概形到 $X$
的拉回。注意，$\infty : X \to (\mathbf{P}^1_X)_\infty$ 是同构。于是
$$
(\mathbf{P}^1_X, (\mathbf{P}^1_X)_\infty, Lp^*E)
$$
是第 \ref{flat-section-blowup-complexes-II} 节情形
\ref{flat-situation-complex-and-divisor} 中的三元组。令
$$
b : W \longrightarrow \mathbf{P}^1_X\quad\text{且}\quad
W_\infty = b^{-1}((\mathbf{P}^1_X)_\infty)
$$
是引理 \ref{flat-lemma-complex-and-divisor-blowup-pre} 从该三元组出发构造的
爆破和有效 Cartier 除子。另记
$$
Q = L\eta_{\mathcal{I}} Lb^*M = L\eta_\mathcal{I} L(p \circ b)^*E
$$
为 引理 \ref{flat-lemma-complex-and-divisor-blowup} 中考虑的
$D(\mathcal{O}_W)$ 的完美对象。这里
$\mathcal{I} \subset \mathcal{O}_W$ 是 $W_\infty$ 的理想层。
% P05-EN-CH38-0125: source 14040 uses undefined M in Lb^*M; the triple contains Lp^*E.

\begin{lemma}
\label{flat-lemma-graph-construction}
上述构造具有如下性质：
\begin{enumerate}
\item $b$ 在 $\mathbf{P}^1_U \cup \mathbf{A}^1_X$ 上是同构；
\item $Q$ 在 $\mathbf{A}^1_X$ 上的限制等于 $E$ 的拉回；
\item 存在有限表现的闭浸入 $i : T \to W_\infty$，使得在概形论意义下
$(W_\infty \to X)^{-1}U \subset T$，且 $Li^*Q$ 具有局部自由上同调层；
\item 对像为 $u \in U$ 的 $t \in T$，$H^i(Li^*Q)_t$ 作为
$\mathcal{O}_{T, t}$-模的秩等于 $H^i(M)_u$ 作为
$\mathcal{O}_{U, u}$-模的秩；
\item 若 $E$ 可由有限局部自由 $\mathcal{O}_X$-模局部有界复形表示，
则 $Q$ 可由有限局部自由 $\mathcal{O}_W$-模局部有界复形表示。
\end{enumerate}
% P05-EN-CH38-0126: source 14060 uses undefined M in the rank comparison; the section's object is E.
\end{lemma}

\begin{proof}
这立即由第 \ref{flat-section-blowup-complexes-II} 节中的结果得到；
汇集全部所需内容的断言见
引理 \ref{flat-lemma-complex-and-divisor-derived}。
\end{proof}
% Generated mechanically from the frozen groupoids.tex target.
% Do not edit this generated file; edit the bound locale source instead.

% OK, start here.
%

\setcounter{chapter}{38}
\chapter{群胚概形}



\phantomsection
\label{groupoids-section-phantom}


\section{引言}
\label{groupoids-section-introduction}

\noindent
本章讨论群胚概形的一般理论。例如，可参阅 Keel 和 Mori
的优美论文 \cite{K-M}。





\section{记号}
\label{groupoids-section-notation}

\noindent
设 $S$ 为概形。若 $U$、$T$ 是 $S$ 上的概形，则以
$U(T)$ 表示 $U$ {\it 在} $S$ 上的 $T$-值点集合。换言之，
$U(T) = \Mor_S(T, U)$。在下文的讨论中，我们尽量保留字母 $T$
来表示 $S$ 上的一个“测试概形”。设 $X$、$Y$ 是
$S$ 上的概形，并设 $f : X \to Y$ 是 $S$ 上的概形态射。
对任意 $S$ 上的概形 $T$，都得到一个诱导的集合映射
$$
f : X(T) \longrightarrow Y(T)
$$
如上所示，我们也以 $f$ 表示此映射。事实上，该构造关于概形
$T/S$ 具有函子性。米田引理（见《范畴》引理 \ref{categories-lemma-yoneda}）说明，$f$ 决定函子变换
$f : h_X \to h_Y$，且也由这一变换唯一决定。更一般地，
对于纤维积之间的映射，我们采用相同的记号。例如，若
$X$、$Y$、$Z$ 是 $S$ 上的概形，且
$m : X \times_S Y \to Z \times_S Z$
是 $S$ 上的概形态射，则可将 $m$ 视为对应于一族 $T$-值点之间的映射
$$
X(T) \times Y(T) \longrightarrow Z(T) \times Z(T).
$$
其余情形依此类推。

\medskip\noindent
我们继续沿用从下标 $0$ 开始标记投影映射的约定；因此有
$\text{pr}_0 : X \times_S Y \to X$ 以及
$\text{pr}_1 : X \times_S Y \to Y$。






\section{等价关系}
\label{groupoids-section-equivalence-relations}

\noindent
回顾集合 $A$ 上的一个{\it 关系} $R$，它就是满足
$R \subset A \times A$ 的子集。通常用 $a R b$ 表示
$(a, b) \in R$。若
$a R b, b R c \Rightarrow a R c$，则称该关系具有{\it 传递性}；
若对所有 $a \in A$ 都有 $a R a$，则称其具有{\it 自反性}；
若 $a R b \Rightarrow b R a$，则称其具有{\it 对称性}。
同时具有传递性、自反性和对称性的关系称为{\it 等价关系}。

\medskip\noindent
在概形的情形中，我们略微放宽关系的概念，只要求
$R \to A \times A$ 是一个映射。定义如下。

\begin{definition}
\label{groupoids-definition-equivalence-relation}
设 $S$ 为概形，$U$ 为 $S$ 上的概形。
\begin{enumerate}
\item $U$ 在 $S$ 上的一个{\it 预关系}，是任意概形态射
$j : R \to U \times_S U$。在此情形下，令
$t = \text{pr}_0 \circ j$、$s = \text{pr}_1 \circ j$，从而
$j = (t, s)$。
\item $U$ 在 $S$ 上的一个{\it 关系}，是概形的单态射
$j : R \to U \times_S U$。
\item 一个{\it 预等价关系}是预关系
$j : R \to U \times_S U$，并且对所有 $T/S$，
$j : R(T) \to U(T) \times U(T)$ 的像都是等价关系。
\item 当且仅当对每个 $S$ 上的概形 $T$，$R$ 的 $T$-值点
在 $U$ 的 $T$-值点集合上给出一个等价关系时，我们称概形态射
$R \to U \times_S U$ 是
{\it $U$ 在 $S$ 上的一个等价关系}。
\end{enumerate}
\end{definition}

\noindent
换言之，等价关系就是一个预等价关系，并且 $j$ 是关系。

\begin{lemma}
\label{groupoids-lemma-restrict-relation}
设 $S$ 为概形，
$U$ 为 $S$ 上的概形，
$j : R \to U \times_S U$ 为预关系，
$g : U' \to U$ 为概形态射。最后令
$$
R' = (U' \times_S U')\times_{U \times_S U} R
\xrightarrow{j'}
U' \times_S U'
$$
则 $j'$ 是 $U'$ 在 $S$ 上的预关系。
若 $j$ 是关系，则 $j'$ 是关系；
若 $j$ 是预等价关系，则 $j'$ 是预等价关系；
若 $j$ 是等价关系，则 $j'$ 是等价关系。
\end{lemma}

\begin{proof}
证明略。
\end{proof}

\begin{definition}
\label{groupoids-definition-restrict-relation}
设 $S$ 为概形，$U$ 为 $S$ 上的概形，
$j : R \to U \times_S U$ 为预关系，
$g : U' \to U$ 为概形态射。
预关系 $j' : R' \to U' \times_S U'$ 称为预关系 $j$
在 $U'$ 上的{\it 限制}或{\it 拉回}。
在这种情形下，有时记 $R' = R|_{U'}$。
\end{definition}

\begin{lemma}
\label{groupoids-lemma-pre-equivalence-equivalence-relation-points}
设 $j : R \to U \times_S U$ 为预关系。
考虑由下述规则在概形 $U$ 的点上定义的关系：
$$
x \sim y
\Leftrightarrow
\exists\ r \in R :
t(r) = x,
s(r) = y.
$$
若 $j$ 是预等价关系，则这是一个等价关系。
\end{lemma}

\begin{proof}
假设 $x \sim y$ 且 $y \sim z$。
取 $r \in R$，使得 $t(r) = x$、$s(r) = y$；再取
$r' \in R$，使得 $t(r') = y$、$s(r') = z$。
取域 $K$，使其出现在如下交换图中：
$$
\xymatrix{
\kappa(r) \ar[r] & K \\
\kappa(y) \ar[u] \ar[r] & \kappa(r') \ar[u]
}
$$
用 $x_K, y_K, z_K : \Spec(K) \to U$
表示态射
$$
\begin{matrix}
\Spec(K) \to \Spec(\kappa(r))
\to
\Spec(\kappa(x)) \to U \\
\Spec(K) \to \Spec(\kappa(r))
\to
\Spec(\kappa(y)) \to U \\
\Spec(K) \to \Spec(\kappa(r'))
\to
\Spec(\kappa(z)) \to U
\end{matrix}
$$
由构造可知 $(x_K, y_K) \in j(R(K))$ 且
$(y_K, z_K) \in j(R(K))$。由于 $j$ 是预等价关系，
还可得到 $(x_K, z_K) \in j(R(K))$。
这显然蕴含 $x \sim z$。

\medskip\noindent
$\sim$ 的自反性和对称性的证明从略。
\end{proof}

\begin{lemma}
\label{groupoids-lemma-etale-equivalence-relation}
设 $j : R \to U \times_S U$ 为预关系。假设
\begin{enumerate}
\item $s, t$ 都是非分歧的；
\item 对 $S$ 上的任意代数闭域 $k$，映射
$R(k) \to U(k) \times U(k)$ 是等价关系；
\item 存在态射 $e : U \to R$、$i : R \to R$、
$c : R \times_{s, U, t} R \to R$，使得
$$
\xymatrix{
U \ar[r]_e \ar[d]_\Delta &
R \ar[d]_j &
R \ar[d]^j \ar[r]_i &
R \ar[d]^j &
R \times_{s, U, t} R \ar[d]^{j \times j} \ar[r]_c &
R \ar[d]^j \\
U \times_S U \ar[r] &
U \times_S U &
U \times_S U \ar[r]^{flip} &
U \times_S U &
U \times_S U \times_S U \ar[r]^{\text{pr}_{02}} &
U \times_S U
}
$$
均为交换图。
\end{enumerate}
则 $j$ 是等价关系。
\end{lemma}

\begin{proof}
由条件 (1) 以及
《概形态射》引理 \ref{morphisms-lemma-unramified-permanence}
可知 $j$ 是非分歧的。因此根据
《概形态射》引理 \ref{morphisms-lemma-diagonal-unramified-morphism}，
$\Delta_j : R \to R \times_{U \times_S U} R$ 是开浸入。
另一方面，条件 (2) 表明 $\Delta_j$ 在 $k$-值点上是双射，
因而 $\Delta_j$ 是同构，所以 $j$ 是单态射。再由上述交换图，
容易推出对 $S$ 上的所有概形 $T$，
$R(T) \subset U(T) \times U(T)$ 都是等价关系。
\end{proof}













\section{群概形}
\label{groupoids-section-group-schemes}

\noindent
回顾一个{\it 群}是二元组 $(G, m)$，其中 $G$ 是集合，
$m : G \times G \to G$ 是满足下列性质的集合映射：
\begin{enumerate}
\item （结合律）$m(g, m(g', g'')) = m(m(g, g'), g'')$
对所有 $g, g', g'' \in G$ 均成立；
\item （单位元）存在唯一的元素 $e \in G$
（称为 $G$ 的{\it 恒等元}、{\it 单位元}或 $1$），使得
$m(g, e) = m(e, g) = g$ 对所有 $g \in G$ 均成立；
\item （逆元）对所有 $g \in G$，存在 $i(g) \in G$，
使得 $m(g, i(g)) = m(i(g), g) = e$，其中 $e$ 是单位元。
\end{enumerate}
于是得到映射 $e : \{*\} \to G$ 和映射 $i : G \to G$，
并且四元组 $(G, m, e, i)$ 满足上列公理。

\medskip\noindent
一个{\it 群同态} $\psi : (G, m) \to (G', m')$ 是集合映射
$\psi : G \to G'$，且满足
$m'(\psi(g), \psi(g')) = \psi(m(g, g'))$。这自动保证
$\psi(e) = e'$ 以及 $i'(\psi(g)) = \psi(i(g))$。
（记号含义显然。）下文将使用这一点。

\begin{definition}
\label{groupoids-definition-group-scheme}
设 $S$ 为概形。
\begin{enumerate}
\item 一个{\it $S$ 上的群概形}是二元组 $(G, m)$，其中
$G$ 是 $S$ 上的概形，而 $m : G \times_S G \to G$
是 $S$ 上的概形态射，且满足如下性质：
对每个 $S$ 上的概形 $T$，二元组 $(G(T), m)$ 都是群。
\item 一个{\it $S$ 上群概形的态射
$\psi : (G, m) \to (G', m')$}
是 $S$ 上的概形态射 $\psi : G \to G'$，并且对每个 $T/S$，
诱导映射 $\psi : G(T) \to G'(T)$ 都是群同态。
\end{enumerate}
\end{definition}

\noindent
设 $(G, m)$ 是概形 $S$ 上的群概形。
由上述讨论（以及第 \ref{groupoids-section-notation} 节中的讨论），
可得到 $S$ 上的概形态射：
（单位元）$e : S \to G$ 以及（取逆）$i : G \to G$；
对每个 $T$，四元组 $(G(T), m, e, i)$ 都满足上列群公理。

\medskip\noindent
设 $(G, m)$、$(G', m')$ 是 $S$ 上的群概形，
$f : G \to G'$ 是 $S$ 上的概形态射。
由定义可知，$f$ 是 $S$ 上群概形的态射，当且仅当下图交换：
$$
\xymatrix{
G \times_S G \ar[r]_-{f \times f} \ar[d]_m &
G' \times_S G' \ar[d]^{m'} \\
G \ar[r]^f & G'
}
$$

\begin{lemma}
\label{groupoids-lemma-base-change-group-scheme}
设 $(G, m)$ 是 $S$ 上的群概形，
$S' \to S$ 是概形态射。
则拉回 $(G_{S'}, m_{S'})$ 是 $S'$ 上的群概形。
\end{lemma}

\begin{proof}
证明略。
\end{proof}

\begin{definition}
\label{groupoids-definition-closed-subgroup-scheme}
设 $S$ 为概形，$(G, m)$ 为 $S$ 上的群概形。
\begin{enumerate}
\item $G$ 的一个{\it 闭子群概形}是闭子概形
$H \subset G$，使得 $m|_{H \times_S H}$ 分解通过 $H$，
并在 $H$ 上诱导出 $S$ 上的群概形结构。
\item $G$ 的一个{\it 开子群概形}是开子概形
$G' \subset G$，使得 $m|_{G' \times_S G'}$ 分解通过 $G'$，
并在 $G'$ 上诱导出 $S$ 上的群概形结构。
\end{enumerate}
\end{definition}

\noindent
等价地，可以说 $H$ 是 $G$ 的闭子群概形，如果它是
$S$ 上的群概形，并带有一个 $S$ 上的群概形态射
$i : H \to G$，且该态射将 $H$ 识别为 $G$ 的闭子概形。

\begin{lemma}
\label{groupoids-lemma-closed-subgroup-scheme}
设 $S$ 为概形，$(G, m, e, i)$ 为 $S$ 上的群概形。
\begin{enumerate}
\item 闭子概形 $H \subset G$ 是闭子群概形，当且仅当
$e : S \to G$、$m|_{H \times_S H} : H \times_S H \to G$
及 $i|_H : H \to G$ 均分解通过 $H$。
\item 开子概形 $H \subset G$ 是开子群概形，当且仅当
$e : S \to G$、$m|_{H \times_S H} : H \times_S H \to G$
及 $i|_H : H \to G$ 均分解通过 $H$。
\end{enumerate}
\end{lemma}

\begin{proof}
考察 $T$-值点后，这恰好转化为刻画一个群的子集何时为子群的熟知条件。
\end{proof}

\begin{definition}
\label{groupoids-definition-smooth-group-scheme}
设 $S$ 为概形，$(G, m)$ 为 $S$ 上的群概形。
\begin{enumerate}
\item 若结构态射 $G \to S$ 是光滑的，则称 $G$ 为{\it 光滑群概形}。
\item 若结构态射 $G \to S$ 是平坦的，则称 $G$ 为{\it 平坦群概形}。
\item 若结构态射 $G \to S$ 是分离的，则称 $G$ 为{\it 分离群概形}。
\end{enumerate}
可按需增补。
\end{definition}






\section{群概形的例子}
\label{groupoids-section-examples-group-schemes}

\begin{example}[乘法群概形]
\label{groupoids-example-multiplicative-group}
考虑这样一个函子：它把任意概形 $T$ 映到结构层整体截面环的单位群
$\Gamma(T, \mathcal{O}_T^*)$。该函子由概形
$$
\mathbf{G}_m = \Spec(\mathbf{Z}[x, x^{-1}])
$$
表示。给出群结构的态射为
\begin{eqnarray*}
\mathbf{G}_m \times \mathbf{G}_m & \to & \mathbf{G}_m \\
\Spec(\mathbf{Z}[x, x^{-1}] \otimes_{\mathbf{Z}} \mathbf{Z}[x, x^{-1}])
& \to &
\Spec(\mathbf{Z}[x, x^{-1}]) \\
\mathbf{Z}[x, x^{-1}] \otimes_{\mathbf{Z}} \mathbf{Z}[x, x^{-1}]
& \leftarrow &
\mathbf{Z}[x, x^{-1}] \\
x \otimes x & \leftarrow & x
\end{eqnarray*}
由此可见，$\mathbf{G}_m$ 是 $\mathbf{Z}$ 上的群概形。
对任意概形 $S$，基变换 $\mathbf{G}_{m, S}$ 是 $S$ 上的群概形，
其点函子仍为
$$
T/S
\longmapsto
\mathbf{G}_{m, S}(T) = \mathbf{G}_m(T) = \Gamma(T, \mathcal{O}_T^*)
$$
。
\end{example}

\begin{example}[单位根]
\label{groupoids-example-roots-of-unity}
设 $n \in \mathbf{N}$。
考虑这样一个函子：它把任意概形 $T$ 映到
$\Gamma(T, \mathcal{O}_T^*)$ 中由 $n$ 次单位根组成的子群。
该函子由概形
$$
\mu_n = \Spec(\mathbf{Z}[x]/(x^n - 1)).
$$
表示。给出群结构的态射为
\begin{eqnarray*}
\mu_n \times \mu_n & \to & \mu_n \\
\Spec(
\mathbf{Z}[x]/(x^n - 1)
\otimes_{\mathbf{Z}}
\mathbf{Z}[x]/(x^n - 1))
& \to &
\Spec(\mathbf{Z}[x]/(x^n - 1)) \\
\mathbf{Z}[x]/(x^n - 1) \otimes_{\mathbf{Z}} \mathbf{Z}[x]/(x^n - 1)
& \leftarrow &
\mathbf{Z}[x]/(x^n - 1) \\
x \otimes x & \leftarrow & x
\end{eqnarray*}
由此可见，$\mu_n$ 是 $\mathbf{Z}$ 上的群概形。
对任意概形 $S$，基变换 $\mu_{n, S}$ 是 $S$ 上的群概形，
其点函子仍为
$$
T/S
\longmapsto
\mu_{n, S}(T) = \mu_n(T) = \{f \in \Gamma(T, \mathcal{O}_T^*) \mid f^n = 1\}
$$
。
\end{example}


\begin{example}[加法群概形]
\label{groupoids-example-additive-group}
考虑这样一个函子：它把任意概形 $T$ 映到结构层整体截面所成的群
$\Gamma(T, \mathcal{O}_T)$。该函子由概形
$$
\mathbf{G}_a = \Spec(\mathbf{Z}[x])
$$
表示。给出群结构的态射为
\begin{eqnarray*}
\mathbf{G}_a \times \mathbf{G}_a & \to & \mathbf{G}_a \\
\Spec(\mathbf{Z}[x] \otimes_{\mathbf{Z}} \mathbf{Z}[x])
& \to &
\Spec(\mathbf{Z}[x]) \\
\mathbf{Z}[x] \otimes_{\mathbf{Z}} \mathbf{Z}[x]
& \leftarrow &
\mathbf{Z}[x] \\
x \otimes 1 + 1 \otimes x & \leftarrow & x
\end{eqnarray*}
由此可见，$\mathbf{G}_a$ 是 $\mathbf{Z}$ 上的群概形。
对任意概形 $S$，基变换 $\mathbf{G}_{a, S}$ 是 $S$ 上的群概形，
其点函子仍为
$$
T/S
\longmapsto
\mathbf{G}_{a, S}(T) = \mathbf{G}_a(T) = \Gamma(T, \mathcal{O}_T)
$$
。
\end{example}

\begin{example}[一般线性群概形]
\label{groupoids-example-general-linear-group}
设 $n \geq 1$。
考虑这样一个函子：它把任意概形 $T$ 映到群
$$
\text{GL}_n(\Gamma(T, \mathcal{O}_T))
$$
即以结构层整体截面环为系数的可逆 $n \times n$ 矩阵所成的群。
该函子由概形
$$
\text{GL}_n = \Spec(\mathbf{Z}[\{x_{ij}\}_{1 \leq i, j \leq n}][1/d])
$$
表示，其中 $d = \det((x_{ij}))$，而 $(x_{ij})$ 是在
$(i, j)$ 位置取值 $x_{ij}$ 的 $n \times n$ 矩阵。
给出群结构的态射为
\begin{eqnarray*}
\text{GL}_n \times \text{GL}_n & \to & \text{GL}_n \\
\Spec(\mathbf{Z}[x_{ij}, 1/d] \otimes_{\mathbf{Z}}
\mathbf{Z}[x_{ij}, 1/d])
& \to &
\Spec(\mathbf{Z}[x_{ij}, 1/d]) \\
\mathbf{Z}[x_{ij}, 1/d] \otimes_{\mathbf{Z}} \mathbf{Z}[x_{ij}, 1/d]
& \leftarrow &
\mathbf{Z}[x_{ij}, 1/d] \\
\sum x_{ik} \otimes x_{kj} & \leftarrow & x_{ij}
\end{eqnarray*}
由此可见，$\text{GL}_n$ 是 $\mathbf{Z}$ 上的群概形。
对任意概形 $S$，基变换 $\text{GL}_{n, S}$ 是 $S$ 上的群概形，
其点函子仍为
$$
T/S
\longmapsto
\text{GL}_{n, S}(T) = \text{GL}_n(T) =\text{GL}_n(\Gamma(T, \mathcal{O}_T))
$$
。
\end{example}

\begin{example}
\label{groupoids-example-determinant}
行列式定义了 $\mathbf{Z}$ 上的群概形态射
$$
\det : \text{GL}_n \longrightarrow \mathbf{G}_m
$$
经基变换，对任意基概形 $S$，它给出其上的群概形态射
$\text{GL}_{n, S} \to \mathbf{G}_{m, S}$。
\end{example}

\begin{example}[常值群]
\label{groupoids-example-constant-group}
设 $G$ 为抽象群。考虑这样一个函子：它把任意概形 $T$
映到局部常值映射 $T \to G$ 所成的群（这里 $T$ 带扎里斯基拓扑，
$G$ 带离散拓扑）。该函子由概形
$$
G_{\Spec(\mathbf{Z})} =
\coprod\nolimits_{g \in G} \Spec(\mathbf{Z}).
$$
表示。给出群结构的态射为
$$
G_{\Spec(\mathbf{Z})}
\times_{\Spec(\mathbf{Z})}
G_{\Spec(\mathbf{Z})}
\longrightarrow
G_{\Spec(\mathbf{Z})}
$$
它把对应于二元组 $(g, g')$ 的分支映到对应于 $gg'$ 的分支。
对任意概形 $S$，基变换 $G_S$ 是 $S$ 上的群概形，
其点函子仍为
$$
T/S
\longmapsto
G_S(T) = \{f : T \to G \text{ 局部常值}\}
$$
。
\end{example}





\section{群概形的性质}
\label{groupoids-section-properties-group-schemes}

\noindent
本节汇集若干在任意基上都成立的群概形的简单性质。

\begin{lemma}
\label{groupoids-lemma-group-scheme-separated}
设 $S$ 为概形，$G$ 为 $S$ 上的群概形。
则 $G \to S$ 是分离的（相应地，\ 拟分离的），当且仅当单位元态射
$e : S \to G$ 是闭浸入（相应地，\ 拟紧的）。
\end{lemma}

\begin{proof}
回顾由《概形》引理 \ref{schemes-lemma-section-immersion}，
$e$ 是一个浸入；而当 $G \to S$ 分离（相应地，\ 拟分离）时，
它是闭浸入（相应地，\ 拟紧的）。
反过来，考虑下图：
$$
\xymatrix{
G \ar[r]_-{\Delta_{G/S}} \ar[d] &
G \times_S G \ar[d]^{(g, g') \mapsto m(i(g), g')} \\
S \ar[r]^e & G
}
$$
从代数几何的函子观点出发，验证此图为笛卡尔图是一个练习。
换言之，$\Delta_{G/S}$ 是 $e$ 的一个基变换。因此，若 $e$
是闭浸入（相应地，\ 拟紧的），则 $\Delta_{G/S}$ 也具有相同性质；
参见《概形》引理 \ref{schemes-lemma-base-change-immersion}
（相应地，\ 《概形》引理 \ref{schemes-lemma-quasi-compact-preserved-base-change}）。
\end{proof}

\begin{lemma}
\label{groupoids-lemma-flat-action-on-group-scheme}
设 $S$ 为概形，$G$ 为 $S$ 上的群概形。
设 $T$ 为 $S$ 上的概形，并设 $\psi : T \to G$ 是 $S$ 上的态射。
若 $T$ 在 $S$ 上平坦，则态射
$$
T \times_S G \longrightarrow G, \quad
(t, g) \longmapsto m(\psi(t), g)
$$
是平坦的。特别地，若 $G$ 在 $S$ 上平坦，则
$m : G \times_S G \to G$ 是平坦的。
\end{lemma}

\begin{proof}
考虑下图：
$$
\xymatrix{
T \times_S G \ar[rrr]_{(t, g) \mapsto (t, m(\psi(t), g))} & & &
T \times_S G \ar[r]_{\text{pr}} \ar[d] &
G \ar[d] \\
& & &
T \ar[r] &
S
}
$$
左上方的水平箭头是同构，方块为笛卡尔图。
因此结论由《概形态射》引理 \ref{morphisms-lemma-base-change-flat} 得出。
\end{proof}

\begin{lemma}
\label{groupoids-lemma-group-scheme-module-differentials}
\begin{reference}
\cite[Proposition 3.15]{BookAV}
\end{reference}
设 $(G, m, e, i)$ 为概形 $S$ 上的群概形，
$f : G \to S$ 为结构态射。则存在典范同构
$$
\Omega_{G/S} \cong f^*\mathcal{C}_{S/G} \cong f^*e^*\Omega_{G/S}
$$
其中 $\mathcal{C}_{S/G}$ 表示浸入 $e$ 的余法层。
特别地，若 $S$ 是一个域的谱，则 $\Omega_{G/S}$
是自由 $\mathcal{O}_G$-模。
\end{lemma}

\begin{proof}
由《概形态射》引理 \ref{morphisms-lemma-base-change-differentials} 有
$$
\Omega_{G \times_S G/G} = \text{pr}_0^*\Omega_{G/S}
$$
其中左端借助 $\text{pr}_1$ 将 $G \times_S G$ 视为 $G$ 上的概形。
令 $\tau : G \times_S G \to G \times_S G$ 为“剪切映射”，
它在点上由 $(g, h) \mapsto (m(g, h), h)$ 给出。
通过投影 $\text{pr}_1$ 将 $G \times_S G$ 视为 $G$ 上的概形时，
此映射是其自同构。结合以上两点，得到同构
$$
\tau^*\text{pr}_0^*\Omega_{G/S} \to \text{pr}_0^*\Omega_{G/S}
$$
由于 $\text{pr}_0 \circ \tau = m$，可将它改写为同构
$$
m^*\Omega_{G/S} \to \text{pr}_0^*\Omega_{G/S}
$$
再沿
$(e \circ f, \text{id}_G) : G \to G \times_S G$
拉回此同构，并利用 $m \circ (e \circ f, \text{id}_G) = \text{id}_G$
以及 $\text{pr}_0 \circ (e \circ f, \text{id}_G) = e \circ f$，
便得到所需的同构
$$
\Omega_{G/S} \to f^*e^*\Omega_{G/S}
$$
。由
《概形态射》引理 \ref{morphisms-lemma-differentials-relative-immersion-section}
还有 $\mathcal{C}_{S/G} \cong e^*\Omega_{G/S}$。
若 $S$ 是一个域的谱，则 $S$ 上的任意 $\mathcal{O}_S$-模都是自由的，
故最后一个断言成立。
\end{proof}

\begin{lemma}
\label{groupoids-lemma-group-scheme-addition-tangent-vectors}
设 $S$ 为概形，$G$ 为 $S$ 上的群概形，并设 $s \in S$。
则复合映射
$$
T_{G/S, e(s)} \oplus T_{G/S, e(s)} = T_{G \times_S G/S, (e(s), e(s))}
\rightarrow T_{G/S, e(s)}
$$
就是切向量的加法。这里的 $=$ 来自
《代数簇》引理 \ref{varieties-lemma-tangent-space-product}，
而右侧箭头由 $m : G \times_S G \to G$ 通过
《代数簇》引理 \ref{varieties-lemma-map-tangent-spaces} 诱导。
\end{lemma}

\begin{proof}
我们将使用 Varieties, Equation (\ref{varieties-equation-tangent-space-fibre})，
并在纤维中处理切向量。第一因子 $T_{G_s/s, e(s)}$ 中的元素
$\theta$，经由 $(1, e) : G_s \to G_s \times G_s$ 所诱导的映射
$T_{G_s/s, e(s)} \to T_{G_s \times G_s/s, (e(s), e(s))}$，
其像就是 $\theta$。由于 $m \circ (1, e) = 1$，
由函子性可知 $\theta$ 映到 $\theta$。
又因该映射是线性的，$(\theta_1, \theta_2)$ 映到
$\theta_1 + \theta_2$。
\end{proof}





\section{域上群概形的性质}
\label{groupoids-section-properties-group-schemes-field}

\noindent
本节汇集群概形在域上的若干性质。对于域上（局部）代数的群概形，
还可以得到更多结论；参见第 \ref{groupoids-section-algebraic-group-schemes} 节。

\begin{lemma}
\label{groupoids-lemma-group-scheme-over-field-open-multiplication}
若 $(G, m)$ 是域 $k$ 上的群概形，则乘法映射
$m : G \times_k G \to G$ 是开映射。
\end{lemma}

\begin{proof}
乘法映射同构于投影映射
$\text{pr}_0 : G \times_k G \to G$，因为交换图
$$
\xymatrix{
G \times_k G \ar[d]^m \ar[rrr]_{(g, g') \mapsto (m(g, g'), g')} & & &
G \times_k G \ar[d]^{(g, g') \mapsto g} \\
G \ar[rrr]^{\text{id}} & & & G
}
$$
的水平箭头都是同构。根据
《概形态射》引理 \ref{morphisms-lemma-scheme-over-field-universally-open}，
该投影是开映射。
\end{proof}

\begin{lemma}
\label{groupoids-lemma-group-scheme-over-field-translate-open}
设 $(G, m)$ 是域 $k$ 上的群概形，$U \subset G$ 是开集，
$T \to G$ 是概形态射。则复合映射
$T \times_k U \to G \times_k G \to G$ 的像是开集。
\end{lemma}

\begin{proof}
对任意域扩张 $K/k$，态射 $G_K \to G$ 是开映射
（《概形态射》引理 \ref{morphisms-lemma-scheme-over-field-universally-open}）。
$T \times_k U$ 的每个点 $\xi$ 都是某个 $K$ 下态射
$(t, u) : \Spec(K) \to T \times_k U$ 的像。此时
$T_K \times_K U_K = (T \times_k U)_K \to G_K$ 的像包含开平移
$t \cdot U_K$。综合这些事实可知，$T \times_k U \to G$ 的像
包含 $\xi$ 之像的一个开邻域。由于 $\xi$ 任意，结论成立。
\end{proof}

\begin{lemma}
\label{groupoids-lemma-group-scheme-over-field-separated}
设 $G$ 为域上的群概形。则 $G$ 是分离概形。
\end{lemma}

\begin{proof}
记 $S = \Spec(k)$，其中 $k$ 为域，并设 $G$ 为 $S$ 上的群概形。
由引理 \ref{groupoids-lemma-group-scheme-separated}，只需证明
$e : S \to G$ 是闭浸入。根据《概形态射》引理 \ref{morphisms-lemma-algebraic-residue-field-extension-closed-point-fibre}，
$e : S \to G$ 的像是 $G$ 的闭点。
显然，$\mathcal{O}_G \to e_*\mathcal{O}_S$ 是满射，因为
$e_*\mathcal{O}_S$ 是支集在 $G$ 的单位元处、取值为 $k$ 的摩天层。
由《概形》引理 \ref{schemes-lemma-characterize-closed-immersions}，
$e$ 是闭浸入。
\end{proof}

\begin{lemma}
\label{groupoids-lemma-group-scheme-field-geometrically-irreducible}
设 $G$ 是域 $k$ 上的群概形。则
\begin{enumerate}
\item $G$ 的每个局部环 $\mathcal{O}_{G, g}$ 都有唯一的极小素理想；
\item $G$ 中恰有一个经过 $e$ 的不可约分支 $Z$；且
\item $Z$ 在 $k$ 上几何不可约。
\end{enumerate}
\end{lemma}

\begin{proof}
对任意点 $g \in G$，存在域扩张 $K/k$ 以及映到 $g$ 的
$K$-值点 $g' \in G(K)$。把 $g'$ 看成群概形 $G_K$ 的
$K$-有理点，则
$\mathcal{O}_{G, g} \to \mathcal{O}_{G_K, g'}$
是忠实平坦的局部环同态（因为 $G_K \to G$ 平坦，且平坦的局部环同态
是忠实平坦的；见《交换代数》引理 \ref{algebra-lemma-local-flat-ff}）。
$\mathcal{O}_{G_K, g'}$ 的结论蕴含 $\mathcal{O}_{G, g}$ 的结论；
参见《交换代数》引理 \ref{algebra-lemma-injective-minimal-primes-in-image}。
因此，为证明 (1)，只需对 $G$ 的 $k$-有理点 $g$ 证明。
此时由 $g$ 作平移定义了自同构 $G \to G$，它把 $e$ 映到 $g$，
故 $\mathcal{O}_{G, g} \cong \mathcal{O}_{G, e}$。
这样便看出 (2) 蕴含 (1)，因为经过 $e$ 的不可约分支与
$\mathcal{O}_{G, e}$ 的极小素理想一一对应。

\medskip\noindent
为证明 (2) 和 (3)，只需在 $k$ 代数闭时证明 (2)。
在这种情形下，设 $Z_1$、$Z_2$ 是 $G$ 的两个经过 $e$ 的不可约分支。
由于 $k$ 代数闭，闭子概形
$Z_1 \times_k Z_2 \subset G \times_k G$ 也不可约；参见
《代数簇》引理 \ref{varieties-lemma-bijection-irreducible-components}。
因此 $m(Z_1 \times_k Z_2)$ 包含在 $G$ 的某个不可约分支中。
另一方面，由于 $m|_{e \times G} = \text{id}_G$ 且
$m|_{G \times e} = \text{id}_G$，它同时包含 $Z_1$ 和 $Z_2$。
故如所需，$Z_1 = Z_2$。
\end{proof}

\begin{remark}
\label{groupoids-remark-warning-group-scheme-geometrically-irreducible}
警告：引理 \ref{groupoids-lemma-group-scheme-field-geometrically-irreducible}
并不意味着 $G/k$ 的每个不可约分支都是几何不可约的。
例如，$\mathbf{Q}$ 上的群概形
$\mu_{3, \mathbf{Q}} = \Spec(\mathbf{Q}[x]/(x^3 - 1))$
有两个不可约分支，它们对应于分解
$x^3 - 1 = (x - 1)(x^2 + x + 1)$。
第一个因子对应经过单位元的不可约分支；第二个不可约分支在
$\Spec(\mathbf{Q})$ 上并非几何不可约。
\end{remark}

\begin{lemma}
\label{groupoids-lemma-reduced-subgroup-scheme-perfect}
设 $G$ 是完美域 $k$ 上的群概形。则 $G$ 的约化 $G_{red}$
是 $G$ 的闭子群概形。
\end{lemma}

\begin{proof}
证明略。提示：利用《代数簇》引理 \ref{varieties-lemma-perfect-reduced} 和
\ref{varieties-lemma-geometrically-reduced-any-base-change}
可知 $G_{red} \times_k G_{red}$ 是约化的。
\end{proof}

\begin{lemma}
\label{groupoids-lemma-open-subgroup-closed-over-field}
设 $k$ 为域，$\psi : G' \to G$ 是 $k$ 上的群概形态射。
若 $\psi(G')$ 在 $G$ 中是开的，则 $\psi(G')$ 在 $G$ 中是闭的。
\end{lemma}

\begin{proof}
令 $U = \psi(G') \subset G$，并令
$Z = G \setminus \psi(G') = G \setminus U$ 带有诱导的约化闭子概形结构。
由引理 \ref{groupoids-lemma-group-scheme-over-field-translate-open}，映射
$$
Z \times_k G' \longrightarrow
Z \times_k U \longrightarrow G
$$
的像是开的（第一个箭头是满射）。另一方面，由于 $\psi$
是群概形同态，$Z \times_k G' \to G$ 的像包含在 $Z$ 中
（因为对 $G'$ 的所有点 $g'$，由 $\psi(g')$ 作平移都保持
$U$；略去一个小细节）。因此 $Z \subset G$ 是开子集
（尽管未必是开子概形）。所以 $U = \psi(G')$ 是闭的。
\end{proof}

\begin{lemma}
\label{groupoids-lemma-immersion-group-schemes-closed-over-field}
设 $i : G' \to G$ 是域 $k$ 上群概形的浸入。
则 $i$ 是闭浸入；也就是说，$i(G')$ 是 $G$ 的闭子群概形。
\end{lemma}

\begin{proof}
为证明 $i$ 是闭浸入，只需证明 $i(G')$ 是 $G$ 的闭子集。
设 $k \subset k'$ 是 $k$ 的一个完美扩张。若
$i(G'_{k'}) \subset G_{k'}$ 是闭的，则由
《概形态射》引理 \ref{morphisms-lemma-fpqc-quotient-topology}，
$i(G') \subset G$ 也是闭的（因为 $G_{k'} \to G$ 平坦、拟紧且满射）。
因此可以并且确实假设 $k$ 完美。下文将直接使用这样一个事实：
$k$ 上约化概形的乘积仍是约化的。由
引理 \ref{groupoids-lemma-reduced-subgroup-scheme-perfect}，
可将 $G'$ 和 $G$ 换成各自的约化。
令 $\overline{G'} \subset G$ 为 $i(G')$ 的闭包，并赋予它约化闭子概形结构。
由《代数簇》引理 \ref{varieties-lemma-closure-of-product}，
$\overline{G'} \times_k \overline{G'}$ 是
$G' \times_k G' \to G \times_k G$ 之像的闭包。因此，由 $m$ 的连续性，
$$
m\Big(\overline{G'} \times_k \overline{G'}\Big)
\subset \overline{G'}
$$
成立。于是 $\overline{G'} \subset G$ 是一个约化闭子群概形。
由引理 \ref{groupoids-lemma-open-subgroup-closed-over-field}，
$i(G') \subset \overline{G'}$ 也是闭的，故如所需，
$i(G') = \overline{G'}$。
\end{proof}

\begin{lemma}
\label{groupoids-lemma-irreducible-group-scheme-over-field-qc}
设 $G$ 是域 $k$ 上的群概形。若 $G$ 不可约，则 $G$ 拟紧。
\end{lemma}

\begin{proof}
设 $K/k$ 为域扩张。若 $G_K$ 拟紧，则因 $G_K \to G$ 满射，
$G$ 也拟紧。由
引理 \ref{groupoids-lemma-group-scheme-field-geometrically-irreducible}，
$G_K$ 不可约。因此，只需把 $k$ 换成某个扩张后证明本引理。
取 $K$ 为基数非常大的代数闭域扩张。由
《代数簇》引理 \ref{varieties-lemma-make-Jacobson}，
$G_K$ 是 Jacobson 概形，并且其所有闭点的剩余域都等于 $K$。
换言之，可假设 $G$ 是 Jacobson 概形，且其所有闭点的剩余域都是 $k$。

\medskip\noindent
设 $U \subset G$ 是非空仿射开集，$g \in G(k)$。于是
$gU \cap U \not = \emptyset$，所以 $g$ 位于态射
$$
U \times_{\Spec(k)} U \longrightarrow G, \quad
(u_1, u_2) \longmapsto u_1u_2^{-1}
$$
的像中。根据引理 \ref{groupoids-lemma-group-scheme-over-field-open-multiplication}，该像是开的。
又因为 $G$ 是 Jacobson 概形且闭点都是 $k$-有理的，所以该像等于整个
$G$。由于 $U$ 仿射，$U \times_{\Spec(k)} U$ 也仿射。
因此 $G$ 是一个拟紧概形的像，从而拟紧。
\end{proof}

\begin{lemma}
\label{groupoids-lemma-connected-group-scheme-over-field-irreducible}
设 $G$ 是域 $k$ 上的群概形。若 $G$ 连通，则 $G$ 不可约。
\end{lemma}

\begin{proof}
由《代数簇》引理 \ref{varieties-lemma-geometrically-connected-if-connected-and-point}，
$G$ 几何连通。若对某个域扩张 $K/k$ 能证明 $G_K$ 不可约，
则本引理成立。因此可应用
《代数簇》引理 \ref{varieties-lemma-make-Jacobson}，
将问题约化到如下情形：$k$ 代数闭，$G$ 是 Jacobson 概形，
且所有闭点都是 $k$-有理的。

\medskip\noindent
设 $Z \subset G$ 是 $G$ 中唯一经过单位元的不可约分支；见
引理 \ref{groupoids-lemma-group-scheme-field-geometrically-irreducible}。
赋予 $Z$ 诱导的约化闭子概形结构，则 $Z \times_k Z$
约化且不可约（《代数簇》引理 \ref{varieties-lemma-geometrically-reduced-any-base-change} 和
\ref{varieties-lemma-bijection-irreducible-components}）。
因此 $m|_{Z \times_k Z} : Z \times_k Z \to G$ 分解通过 $Z$，
从而 $Z$ 成为 $G$ 的闭子群概形。

\medskip\noindent
为导出矛盾，假设还存在另一个不可约分支 $Z' \subset G$。
由引理 \ref{groupoids-lemma-group-scheme-field-geometrically-irreducible}，
$Z \cap Z' = \emptyset$。由
引理 \ref{groupoids-lemma-irreducible-group-scheme-over-field-qc}，$Z$ 拟紧。
于是可以选取拟紧开集 $U \subset G$，使得
$Z \subset U$ 且 $U \cap Z' = \emptyset$。
由引理 \ref{groupoids-lemma-group-scheme-over-field-translate-open}，
$Z \times_k U \to G$ 的像 $W$ 在 $G$ 中是开的。
另一方面，$W$ 是拟紧空间的像，因而拟紧。我们断言 $W$ 是闭的。

\medskip\noindent
证明该断言。由于 $W$ 拟紧而 $G$ 分离
（引理 \ref{groupoids-lemma-group-scheme-over-field-separated}），根据
《概形》引理 \ref{schemes-lemma-quasi-compact-permanence}，
包含映射 $W \to G$ 是拟紧的。因此，$W$ 闭包中的点都是
$W$ 中某些点的特化（《概形态射》引理 \ref{morphisms-lemma-reach-points-scheme-theoretic-image}）。
所以只需证明：$G$ 中任何与 $W$ 相交的不可约分支
$Z'' \subset G$ 都包含在 $W$ 中。
由于 $G$ 是 Jacobson 概形且闭点均为有理点，
$Z'' \cap W$ 有一个有理点 $g \in Z''(k) \cap W(k)$，
从而 $Z'' = Zg$。但由构造 $W = m(Z \times_k W)$，所以
$Z'' \cap W \not = \emptyset$ 蕴含 $Z'' \subset W$。

\medskip\noindent
由此断言，$W \subset G$ 是 $G$ 的既开又闭子集。
此外 $W \cap Z' = \emptyset$，否则由上一段的论证，
对某个 $g \in W(k)$ 会有 $Z' = Zg$。由于 $Z$ 是子群，
甚至可以取 $g \in U(k)$，这与 $Z' \cap U = \emptyset$ 矛盾。
因此 $W \subset G$ 是真既开又闭子集，与 $G$ 连通的假设矛盾。
\end{proof}

\begin{proposition}
\label{groupoids-proposition-connected-component}
设 $G$ 为域 $k$ 上的群概形。存在一个典范闭子群概形
$G^0 \subset G$，它具有下列性质：
\begin{enumerate}
\item $G^0 \to G$ 是平坦闭浸入；
\item $G^0 \subset G$ 是单位元的连通分支；
\item $G^0$ 几何不可约；且
\item $G^0$ 拟紧。
\end{enumerate}
\end{proposition}

\begin{proof}
令 $G^0$ 为单位元的连通分支，并赋予它典范概形结构
（《概形态射》定义 \ref{morphisms-definition-scheme-structure-connected-component}）。
为证明 $G^0$ 是闭子群概形，使用
引理 \ref{groupoids-lemma-closed-subgroup-scheme} 的判据。
态射 $e : \Spec(k) \to G$ 分解通过 $G^0$，因为我们所取的
$G^0$ 正是 $G$ 中包含 $e$ 的连通分支。
由于 $i : G \to G$ 是固定 $e$ 的自同构，$i$ 把 $G^0$ 映入自身。
根据《代数簇》引理 \ref{varieties-lemma-geometrically-connected-criterion}，
概形 $G^0$ 在 $k$ 上几何连通。因此 $G^0 \times_k G^0$ 连通
（《代数簇》引理 \ref{varieties-lemma-bijection-connected-components}）。
故在集合论意义下 $m(G^0 \times_k G^0) \subset G^0$。
由《概形态射》引理 \ref{morphisms-lemma-characterize-flat-closed-immersions}，
$m|_{G^0 \times_k G^0} : G^0 \times_k G^0 \to G$
分解通过 $G^0$。所以 $G^0$ 是 $G$ 的闭子群概形。
由引理 \ref{groupoids-lemma-connected-group-scheme-over-field-irreducible}，
$G^0$ 不可约；由
引理 \ref{groupoids-lemma-group-scheme-field-geometrically-irreducible}，
$G^0$ 几何不可约；再由
引理 \ref{groupoids-lemma-irreducible-group-scheme-over-field-qc}，
$G^0$ 拟紧。
\end{proof}

\begin{lemma}
\label{groupoids-lemma-profinite-product-over-field}
设 $k$ 为域，$T = \Spec(A)$，其中 $A$ 是一族代数的有向余极限，
而每个代数都是若干份 $k$ 的有限乘积。对任意 $k$ 上的概形 $X$，
作为拓扑空间有 $|T \times_k X| = |T| \times |X|$。
\end{lemma}

\begin{proof}
取仿射开覆盖后，可约化到 $X$ 仿射的情形。记
$X = \Spec(B)$。写成 $A = \colim A_i$，其中
$A_i = \prod_{t \in T_i} k$ 且 $T_i$ 有限。
于是 $T_i = |\Spec(A_i)|$ 带离散拓扑，而转移态射
$A_i \to A_{i'}$ 由集合映射 $T_{i'} \to T_i$ 给出。
因此，作为拓扑空间有 $|T| = \lim T_i$；见
《概形的极限》引理 \ref{limits-lemma-topology-limit}。类似地，
\begin{align*}
|T \times_k X| & =
|\Spec(A \otimes_k B)| \\
& =
|\Spec(\colim A_i \otimes_k B)| \\
& =
\lim |\Spec(A_i \otimes_k B)| \\
& =
\lim |\Spec(\prod\nolimits_{t \in T_i} B)| \\
& =
\lim T_i \times |X| \\
& =
(\lim T_i) \times |X| \\
& =
|T| \times |X|
\end{align*}
这里使用了上面的引理以及极限与极限可交换这一事实。
\end{proof}

\noindent
下一个引理说明，事实上可以把一个“点的代数性预有限族”
置于某个仿射开集中。强烈建议读者先阅读
引理 \ref{groupoids-lemma-points-in-affine}。

\begin{lemma}
\label{groupoids-lemma-compact-set-in-affine}
设 $k$ 为代数闭域，$G$ 为 $k$ 上的群概形。
假设 $G$ 是 Jacobson 概形，且所有闭点均为 $k$-有理点。
设 $T = \Spec(A)$，其中 $A$ 是一族代数的有向余极限，
而每个代数都是若干份 $k$ 的有限乘积。对任意态射 $f : T \to G$，
存在包含 $f(T)$ 的仿射开集 $U \subset G$。
\end{lemma}

\begin{proof}
令 $G^0 \subset G$ 为命题 \ref{groupoids-proposition-connected-component} 中得到的闭子群概形。
前两段将问题约化到 $G = G^0$ 的情形。

\medskip\noindent
注意，$T$ 是有限拓扑空间的有向逆极限
（《概形的极限》引理 \ref{limits-lemma-topology-limit}），
因而作为拓扑空间是预有限的
（《拓扑学》定义 \ref{topology-definition-profinite}）。
令 $W \subset G$ 为包含 $T \to G$ 之像的拟紧开集。
把 $W$ 替换为 $G^0 \times W \to G \times G \to G$ 的像后，
可假设 $W$ 在 $G^0$ 的左平移作用下不变；见
引理 \ref{groupoids-lemma-group-scheme-over-field-translate-open}。
考虑复合映射
$$
\psi = \pi \circ f : T \xrightarrow{f} W \xrightarrow{\pi} \pi_0(W)
$$
空间 $\pi_0(W)$ 是预有限的
（《拓扑学》引理 \ref{topology-lemma-spectral-pi0} 以及
《概形的性质》引理 \ref{properties-lemma-quasi-compact-quasi-separated-spectral}）。
对 $\xi \in \pi_0(W)$，令 $F_\xi \subset T$ 为
$T \to \pi_0(W)$ 在该点上的纤维。
假设对每个 $\xi \in \pi_0(W)$，都已找到仿射开集
$U_\xi \subset W$，使得 $f(F_\xi) \subset U_\xi$。
由《拓扑学》引理 \ref{topology-lemma-closed-map}，$\psi$ 是闭映射，
所以 $\psi(T \setminus f^{-1}(U_\xi))$ 是 $\pi_0(W)$
中不含 $\xi$ 的闭子集。其补集 $V_\xi$ 因而是 $\xi$ 的开邻域，
且 $\psi^{-1}(V_\xi) \subseteq f^{-1}(U_\xi)$。
由 $V_\xi$ 组成的 $\pi_0(W)$ 的开覆盖，可以细化为两两不交的
既开又闭子集所成的有限覆盖
$\pi_0(W) = V_1 \amalg \ldots \amalg V_n$
（《拓扑学》引理 \ref{topology-lemma-profinite-refine-open-covering}）。
对每个 $i = 1, \ldots, n$，取 $\xi_i \in \pi_0(W)$，
使得 $V_i \subset V_{\xi_i}$。由于 $V_i$ 闭而 $U_{\xi_i}$ 仿射，
$U_{\xi_i} \cap \pi^{-1}(V_i)$ 也仿射。
把 $U_{\xi_i}$ 替换为 $U_{\xi_i} \cap \pi^{-1}(V_i)$ 后，有
$\psi^{-1}(V_i) = f^{-1}(U_{\xi_i})$。因为这些 $V_i$ 覆盖
$\pi_0(W)$，所以这些 $f^{-1}(U_{\xi_i})$ 覆盖 $T$，从而
$f(T) \subseteq U = \bigcup_i U_{\xi_i}$。此时 $U$ 是有限多个
仿射开集的不交并，因此如所需为仿射的。

\medskip\noindent
设 $Z$ 为一个与 $f(T)$ 相交的 $G$ 的连通分支。则 $Z$
有一个 $k$-有理点 $z$（因为概形 $T$ 的所有剩余域都同构于 $k$）。
因此 $Z = G^0 z$。由 $W$ 的选取，$Z \subset W$。
上一段的论证把问题约化为：在 $W$ 中寻找 $f(T) \cap Z$
的仿射开邻域。经一个有理点平移后，可假设 $Z = G^0$
（细节略）。注意，概形论逆像
$T' = f^{-1}(G^0) \subset T$ 是同类型的闭子概形。
把 $T$ 替换为 $T'$ 后，可假设 $f(T) \subset G^0$。
选取 $e \in G$ 的一个仿射开邻域 $U \subset G$，
从而特别有 $U \cap G^0$ 非空。我们将证明，存在
$g \in G^0(k)$，使得 $f(T) \subset g^{-1}U$。
由于 $W$ 在 $G^0$ 的左平移下不变，$g^{-1}U \subset W$，
这将完成证明。

\medskip\noindent
前两段的论证允许我们转到 $G^0$，并将问题约化如下：
假设 $G$ 不可约，$U \subset G$ 是 $e$ 的仿射开邻域。
证明对某个 $g \in G(k)$ 有 $f(T) \subset g^{-1}U$。
考虑态射
$$
U \times_k T \longrightarrow G \times_k T,\quad
(t, u) \longrightarrow (uf(t)^{-1}, t)
$$
这是开浸入，因为它向 $G \times_k T \to G \times_k T$
的延拓是同构。由对 $T$ 的假设及
引理 \ref{groupoids-lemma-profinite-product-over-field}，有
$|U \times_k T| = |U| \times |T|$，对 $G \times_k T$ 也类似。
因此，上述开浸入的像是有限多个盒
$\bigcup_{i = 1, \ldots, n} U_i \times V_i$ 的并，其中
$V_i \subset T$，而 $U_i \subset G$ 是拟紧开集。
这意味着可能出现的开集 $Uf(t)^{-1}$（$t \in T$）只有有限多个，
记为 $Uf(t_1)^{-1}, \ldots, Uf(t_r)^{-1}$。
由于 $G$ 不可约，交集
$$
Uf(t_1)^{-1} \cap \ldots \cap Uf(t_r)^{-1}
$$
非空。又由于 $G$ 是 Jacobson 概形且闭点均为 $k$-有理点，
可在此交集中选取一个 $k$-值点 $g \in G(k)$。
于是对所有 $t \in T$ 都有 $g \in Uf(t)^{-1}$，
即如所需，$f(t) \in g^{-1}U$。
\end{proof}

\begin{remark}
\label{groupoids-remark-easy}
若 $G$ 是域上的群概形，是否总存在一个拟紧的既开又闭子群概形？
由命题 \ref{groupoids-proposition-connected-component}，
只有当 $G$ 在几何意义下有无穷多个连通分支时，这个问题才有意义。
\end{remark}

\begin{lemma}
\label{groupoids-lemma-group-scheme-field-countable-affine}
设 $G$ 为域上的群概形。存在既开又闭的子概形
$G' \subset G$，它是仿射概形的可数并。
\end{lemma}

\begin{proof}
设 $e \in U(k)$ 为单位元的拟紧开邻域。
把 $U$ 替换为 $U \cap i(U)$ 后，可假设 $U$ 在取逆映射下不变。
由于 $G$ 分离，这仍是拟紧集。令
$$
G' = \bigcup\nolimits_{n \geq 1} m_n(U \times_k \ldots \times_k U)
$$
其中 $m_n : G \times_k \ldots \times_k G \to G$ 是 $n$ 元乘法映射
$(g_1, \ldots, g_n) \mapsto m(m(\ldots (m(g_1, g_2), g_3), \ldots ), g_n)$。
这些映射都是开的（见
引理 \ref{groupoids-lemma-group-scheme-over-field-open-multiplication}），
故 $G'$ 是开子群概形。由
引理 \ref{groupoids-lemma-open-subgroup-closed-over-field}，它也是闭子群概形。
\end{proof}






\section{代数群概形的性质}
\label{groupoids-section-algebraic-group-schemes}

\noindent
回顾域 $k$ 上的概形称为（局部）代数的，是指它在
$\Spec(k)$ 上（局部）为有限型；见
《代数簇》定义 \ref{varieties-definition-algebraic-scheme}。
本节标题中的“代数”正是这个含义。

\begin{lemma}
\label{groupoids-lemma-group-scheme-finite-type-field}
设 $k$ 为域，$G$ 为 $k$ 上的局部代数群概形。
则 $G$ 等维，并且对所有 $g \in G$ 都有
$\dim(G) = \dim_g(G)$。对任意闭点 $g \in G$，有
$\dim(G) = \dim(\mathcal{O}_{G, g})$。
\end{lemma}

\begin{proof}
先证明，对任意一对点 $g, g' \in G$，都有
$\dim_g(G) = \dim_{g'}(G)$。由
《概形态射》引理 \ref{morphisms-lemma-dimension-fibre-after-base-change}，
可以任意扩张基域。因此可假设 $g$ 和 $g'$ 都在 $k$ 上有定义。
于是存在 $G$ 的一个自同构，把 $g$ 映到 $g'$，故二者相等。
根据《概形态射》引理 \ref{morphisms-lemma-dimension-fibre-at-a-point}，有
$\dim_g(G) = \dim(\mathcal{O}_{G, g}) +
\text{trdeg}_k(\kappa(g))$。
另一方面，$G$（或 $G$ 的任意开子集）的维数是 $G$ 的局部环维数的上确界；
见《概形的性质》引理 \ref{properties-lemma-codimension-local-ring}。
显然，该值在闭点 $g$ 处达到最大；此时由 Hilbert 零点定理，
$\text{trdeg}_k(\kappa(g)) = 0$（见
《概形态射》第 \ref{morphisms-section-points-finite-type} 节）。
引理得证。
\end{proof}

\noindent
下述结果有时称为 Cartier 定理。

\begin{lemma}
\label{groupoids-lemma-group-scheme-characteristic-zero-smooth}
设 $k$ 为特征 $0$ 的域，$G$ 为 $k$ 上的局部代数群概形。
则结构态射 $G \to \Spec(k)$ 是光滑的；也就是说，
$G$ 是光滑群概形。
\end{lemma}

\begin{proof}
由引理 \ref{groupoids-lemma-group-scheme-module-differentials}，
$G$ 在 $k$ 上的微分模是自由的。于是光滑性由
《代数簇》引理 \ref{varieties-lemma-char-zero-differentials-free-smooth}
推出。
\end{proof}

\begin{remark}
\label{groupoids-remark-when-reduced}
特征 $0$ 的域上的任意群概形都是约化的；参见
\cite[I, Theorem 1.1 and I, Corollary 3.9, and II, Theorem 2.4]{Perrin-thesis}
以及 \cite[Proposition 4.2.8]{Perrin}。
这一问题由 \cite[page 80]{Oort} 提出。
我们已在引理 \ref{groupoids-lemma-group-scheme-characteristic-zero-smooth}
中看到，当群概形局部为有限型时，该结论成立。
\end{remark}

\begin{lemma}
\label{groupoids-lemma-reduced-group-scheme-prefect-field-characteristic-p-smooth}
设 $k$ 为特征 $p > 0$ 的完美域（特征零情形见
引理 \ref{groupoids-lemma-group-scheme-characteristic-zero-smooth}）。
设 $G$ 为 $k$ 上的局部代数群概形。若 $G$ 约化，则结构态射
$G \to \Spec(k)$ 是光滑的；也就是说，$G$ 是光滑群概形。
\end{lemma}

\begin{proof}
由引理 \ref{groupoids-lemma-group-scheme-module-differentials}，
层 $\Omega_{G/k}$ 是自由的。因此结论由
《代数簇》引理 \ref{varieties-lemma-char-p-differentials-free-smooth}
推出。
\end{proof}

\begin{remark}
\label{groupoids-remark-reduced-smooth-not-true-general}
设 $k$ 为特征 $p > 0$ 的域，$\alpha \in k$ 不是 $p$ 次幂。
闭子群概形
$$
G = V(x^p + \alpha y^p) \subset \mathbf{G}_{a, k}^2
$$
约化且不可约，但不光滑，甚至不正规。
\end{remark}

\noindent
下一个引理是引理 \ref{groupoids-lemma-compact-set-in-affine} 的特殊情形，
其证明稍为简单。

\begin{lemma}
\label{groupoids-lemma-points-in-affine}
设 $k$ 为代数闭域，$G$ 为 $k$ 上的局部代数群概形，
$g_1, \ldots, g_n \in G(k)$ 为 $k$-有理点。
则存在包含 $g_1, \ldots, g_n$ 的仿射开集 $U \subset G$。
\end{lemma}

\begin{proof}
首先对 $n$ 归纳，说明可假设所有 $g_i$ 都位于 $G$ 的同一连通分支上。
否则，可以找到分解 $G = W_1 \amalg W_2$，其中 $W_i$ 在 $G$ 中开；
适当重编号后，对某个 $0 < r < n$ 有
$g_1, \ldots, g_r \in W_1$ 以及
$g_{r + 1}, \ldots, g_n \in W_2$。
由归纳假设，可找到 $G$ 的仿射开集 $U_1$ 和 $U_2$，使得
$g_1, \ldots, g_r \in U_1$ 且
$g_{r + 1}, \ldots, g_n \in U_2$。于是
$$
g_1, \ldots, g_n \in (U_1 \cap W_1) \cup (U_2 \cap W_2)
$$
给出所求解。因此可假设 $g_1, \ldots, g_n$ 都位于 $G$
的同一连通分支。由 $g_1^{-1}$ 作平移后，可假设
$g_1, \ldots, g_n \in G^0$，其中 $G^0 \subset G$ 如
命题 \ref{groupoids-proposition-connected-component} 所定义。
选取 $e$ 的仿射开邻域 $U$；特别地，$U \cap G^0$ 非空。
由于 $G^0$ 不可约，
$$
G^0 \cap (Ug_1^{-1} \cap \ldots \cap Ug_n^{-1})
$$
非空。由于 $G \to \Spec(k)$ 局部为有限型，
$G^0 \to \Spec(k)$ 也局部为有限型，故每个非空开集都有
$k$-有理点。因此可取 $g \in G^0(k)$，使得对所有 $i$ 都有
$g \in Ug_i^{-1}$。于是对所有 $i$ 都有 $g_i \in g^{-1}U$，
而 $g^{-1}U$ 就是所求仿射开集。
\end{proof}

\begin{lemma}
\label{groupoids-lemma-algebraic-quasi-projective}
设 $k$ 为域，$G$ 为 $k$ 上的代数群概形。则 $G$ 在 $k$ 上拟射影。
\end{lemma}

\begin{proof}
由《代数簇》引理 \ref{varieties-lemma-ample-after-field-extension}，
可假设 $k$ 代数闭。令 $G^0 \subset G$ 为命题 \ref{groupoids-proposition-connected-component} 中的 $G$ 的连通分支。
$G$ 的其他每个连通分支都有 $k$-有理点，因而作为概形同构于 $G^0$。
由于 $G$ 拟紧且 Noether，这样的连通分支只有有限多个。
于是问题约化到下一段所讨论的情形。

\medskip\noindent
设 $G$ 是代数闭域 $k$ 上的连通代数群概形。
若 $k$ 的特征为零，则由
引理 \ref{groupoids-lemma-group-scheme-characteristic-zero-smooth}，
$G$ 在 $k$ 上光滑。若 $k$ 的特征为 $p > 0$，令
$H = G_{red}$ 为 $G$ 的约化。根据
《除子》命题 \ref{divisors-proposition-push-down-ample}，
只需证明 $H$ 有充足可逆层。（对 $k$ 上的代数概形，
有充足可逆层等价于在 $k$ 上拟射影；例如见相当一般的
《态射进阶》引理 \ref{more-morphisms-lemma-quasi-projective}。）
由引理 \ref{groupoids-lemma-reduced-subgroup-scheme-perfect}，
$H$ 是 $k$ 上的群概形；由
引理 \ref{groupoids-lemma-reduced-group-scheme-prefect-field-characteristic-p-smooth}，
$H$ 在 $k$ 上光滑。这将问题约化到下一段所讨论的情形。

\medskip\noindent
设 $G$ 是代数闭域 $k$ 上拟紧、不可约且光滑的群概形。
注意，$G$ 的局部环都是正则的，因而是唯一分解整环
（《代数簇》引理 \ref{varieties-lemma-smooth-regular} 以及
《代数进阶》引理 \ref{more-algebra-lemma-regular-local-UFD}）。
$G$ 的一个非空仿射开集的补集，是某个有效 Cartier 除子 $D$ 的支集；
这由《除子》引理 \ref{divisors-lemma-complement-open-affine-effective-cartier-divisor} 得到。
（由引理 \ref{groupoids-lemma-group-scheme-over-field-separated}，$G$ 分离。）
因此存在有效 Cartier 除子 $D \subset G$，使得
$G \setminus D$ 仿射。下文还将用到：对任意 $n \geq 1$ 以及
$g_1, \ldots, g_n \in G(k)$，补集
$G \setminus \bigcup D g_i$ 仿射。事实上，它是分离概形 $G$ 中
诸仿射开集 $G \setminus Dg_i \cong G \setminus D$ 的交。

\medskip\noindent
可以选取下图的第一行
$$
\xymatrix{
G & U \ar[l]_j \ar[r]^\pi & \mathbf{A}^d_k \\
& W \ar[r]^{\pi'} \ar[u] & V \ar[u]
}
$$
使得 $U \not = \emptyset$，$j : U \to G$ 是开浸入，而
$\pi$ 是平展的（\'etale）；参见
《概形态射》引理 \ref{morphisms-lemma-smooth-etale-over-affine-space}。
存在非空仿射开集 $V \subset \mathbf{A}^d_k$，使得取
$W = \pi^{-1}(V)$ 时，态射 $\pi' = \pi|_W : W \to V$
是有限平展的（finite \'etale）。特别地，$\pi'$ 有限局部自由，
设其次数为 $n$。考虑有效 Cartier 除子
$$
\mathcal{D} = \{(g, w) \mid m(g, j(w)) \in D\} \subset G \times W
$$
（这是把 $D \subset G$ 沿平坦态射 $m : G \times G \to G$
拉回后，再限制到 $G \times W$ 所得的除子。）
考虑闭子集\footnote{利用《除子》第 \ref{divisors-section-norms} 节的材料，可以把有效 Cartier 除子
$\mathcal{D}$ 沿有限局部自由态射 $1 \times \pi'$ 的范数取作
有效 Cartier 除子 $E$，从而省去部分论证。}
$T = (1 \times \pi')(\mathcal{D}) \subset G \times V$。
由于 $\pi'$ 有限局部自由，$T$ 的每个不可约分支在
$G \times V$ 中的余维都是 $1$。又因 $G \times V$ 在 $k$ 上光滑，
这些分支都是有效 Cartier 除子（《除子》引理 \ref{divisors-lemma-weil-divisor-is-cartier-UFD} 以及上引诸引理），
所以 $T$ 是 $G \times V$ 中某个有效 Cartier 除子 $E$ 的支集。
若 $v \in V(k)$，则
$(\pi')^{-1}(v) = \{w_1, \ldots, w_n\} \subset W(k)$，并且在集合论意义下，
$$
E_v = \bigcup\nolimits_{i = 1, \ldots, n} D j(w_i)^{-1}
$$
在 $G$ 中成立。特别地，$G \setminus E_v$ 是仿射开集（见上文）。
此外，对任意 $g \in G(k)$，存在 $v \in V$，使得
$g \not \in E_v$。事实上，满足 $g \not \in Dj(w)^{-1}$ 的
$w \in W$ 所成的集合 $W'$ 是非空开集；只需选取 $v$，
使得 $W' \to V$ 在 $v$ 上的纤维有 $n$ 个元素。

\medskip\noindent
考虑 $G \times V$ 上的可逆层
$\mathcal{M} = \mathcal{O}_{G \times V}(E)$。
由《代数簇》引理 \ref{varieties-lemma-rational-equivalence-for-Pic}，
限制 $\mathcal{M}_v = \mathcal{O}_G(E_v)$ 的同构类
$\mathcal{L}$ 与 $v \in V(k)$ 无关。另一方面，对每个
$g \in G(k)$，都可找到 $v$，使得 $g \not \in E_v$ 且
$G \setminus E_v$ 仿射。因此，$\mathcal{O}_G(E_v)$ 的典范截面
（《除子》定义 \ref{divisors-definition-invertible-sheaf-effective-Cartier-divisor}）
对应于 $\mathcal{L}$ 的一个截面 $s_v$；它在 $g$ 处不消失，
并且 $G_{s_v}$ 仿射。按照定义，这说明 $\mathcal{L}$ 充足
（《概形的性质》定义 \ref{properties-definition-ample}）。
\end{proof}

\begin{lemma}
\label{groupoids-lemma-algebraic-center}
设 $k$ 为域，$G$ 为 $k$ 上的局部代数群概形。
则 $G$ 的中心是 $G$ 的闭子群概形。
\end{lemma}

\begin{proof}
以 $\text{Aut}(G)$ 表示 $k$ 上概形范畴中的反变函子：
它把 $S/k$ 映到基变换 $G_S$ 作为 $S$ 上群概形的自同构集合。
有自然变换
$$
G \longrightarrow \text{Aut}(G),\quad
g \longmapsto \text{inn}_g
$$
它把 $G$ 的一个 $S$-值点 $g$ 映到由 $g$ 决定的 $G$ 的内自同构。
按照定义，$G$ 的中心 $C$ 是该变换的核；也就是说，它是把
$S$ 映到如下 $g \in G(S)$ 所成集合的函子：其对应的内自同构
是恒等的。本引理的断言是，该函子可由 $G$ 的一个闭子群概形表示。

\medskip\noindent
选取整数 $n \geq 1$。令 $G_n \subset G$ 为 $G$ 的单位元 $e$
的第 $n$ 阶无穷小邻域。对每个概形 $S/k$，基变换 $G_{n, S}$
都是 $e_S : S \to G_S$ 的第 $n$ 阶无穷小邻域。
因此有自然变换
$\text{Aut}(G) \to \text{Aut}(G_n)$，右端是 $G_n$
作为概形的自同构函子（一般而言，$G_n$ 不是群概形）。
注意，$G_n$ 是某个 Artin 局部环 $A_n$ 的谱；其剩余域为 $k$，
并且作为 $k$-向量空间是有限维的
（《代数簇》引理 \ref{varieties-lemma-algebraic-scheme-dim-0}）。
由于 $G_n$ 的每个自同构都特别诱导一个可逆线性映射
$A_n \to A_n$，得到函子变换
$$
G \to \text{Aut}(G) \to \text{Aut}(G_n) \to \text{GL}(A_n)
$$
最后一个群值函子是可表示的；见
例 \ref{groupoids-example-general-linear-group}。并且最后一个箭头显然为单射。
因此，对每个 $n$ 都得到闭子群概形
$$
H_n = \Ker(G \to \text{Aut}(G_n)) = \Ker(G \to \text{GL}(A_n)).
$$
作为第一步近似，令 $H = \bigcap_{n \geq 1} H_n$
（概形论交）。这是一个包含中心 $C$ 的闭子群概形。

\medskip\noindent
设 $h$ 是 $H$ 的一个 $S$-值点，其中 $S$ 局部 Noether。
于是自同构 $\text{inn}_h$ 在所有闭子概形 $G_{n, S}$ 上诱导恒等映射。
考虑核
$K = \Ker(\text{inn}_h : G_S \to G_S)$。
它是 $S$ 上 $G_S$ 的闭子群概形，并包含所有 $n \geq 1$ 的
闭子概形 $G_{n, S}$。这蕴含 $K$ 包含
$e(S) \subset G_S$ 的一个开邻域；见
《交换代数》注 \ref{algebra-remark-intersection-powers-ideal}。
令 $G^0 \subset G$ 如命题 \ref{groupoids-proposition-connected-component} 所定义。
由于 $G^0$ 几何不可约，可断定 $K$ 包含 $G^0_S$
（对任意非空开集 $U \subset G^0_{k'}$ 和任意域扩张 $k'/k$，
都有 $U \cdot U^{-1} = G^0_{k'}$；见
引理 \ref{groupoids-lemma-irreducible-group-scheme-over-field-qc} 的证明）。
取 $S = H$ 应用上述结论，得到 $G^0$ 和 $H$ 是 $G$
中其点两两交换的子群概形：对任意概形 $S$ 以及任意 $S$-值点
$g \in G^0(S)$、$h \in H(S)$，在 $G(S)$ 中都有 $gh = hg$。

\medskip\noindent
假设 $k$ 代数闭。在 $G$ 的每个不可约分支 $G_i$ 中选取一个
$k$-值点 $g_i$。注意，在此情形下，$G$ 的连通分支就是 $G$ 的不可约分支，
而它们都是 $G^0$ 经各 $g_i$ 平移所得。我们断言
$$
C = H \cap \bigcap\nolimits_i \Ker(\text{inn}_{g_i} : G \to G)
\quad (\text{概形论交})
$$
事实上，$C$ 包含于右端。反过来，右端的每个 $S$-值点 $h$
都与 $G^0$ 以及各 $g_i$ 交换，因而与
$G = \bigcup G^0g_i$ 中的一切元素交换。

\medskip\noindent
一般基域 $k$ 的情形由代数闭包 $\overline{k}$ 上的结果经下降得到。
具体地，令 $A \subset G_{\overline{k}}$ 为表示
$G_{\overline{k}}$ 之中心的闭子群概形。由中心构造的函子性，
作为 $G_{\overline{k} \otimes_k \overline{k}}$ 的闭子概形有
$$
A \times_{\Spec(k)} \Spec(\overline{k}) =
\Spec(\overline{k}) \times_{\Spec(k)} A
$$
所以由《下降》引理 \ref{descent-lemma-closed-immersion}
（以及《下降》引理 \ref{descent-lemma-descending-property-closed-immersion}），
$A$ 下降为闭子群概形 $Z \subset G$。
于是 $Z$ 表示 $C$（略去一个小论证），证明完毕。
\end{proof}







\section{阿贝尔簇}
\label{groupoids-section-abelian-varieties}

\noindent
Mumford 关于阿贝尔簇的著作是这一主题的极佳参考文献；
见 \cite{AVar}，建议读者参阅。存在许多等价定义，下面给出其中一个。

\begin{definition}
\label{groupoids-definition-abelian-variety}
设 $k$ 为域。一个{\it 阿贝尔簇}是 $k$ 上的群概形，同时也是
$k$ 上固有且几何整的簇\footnote{等价定义见
注 \ref{groupoids-remark-16-equivalent}。}。
\end{definition}

\noindent
下面先证明关于这一概念的几个引理，再在
命题 \ref{groupoids-proposition-review-abelian-varieties}
中汇总所有结论。

\begin{lemma}
\label{groupoids-lemma-abelian-variety-projective}
设 $k$ 为域，$A$ 为 $k$ 上的阿贝尔簇。则 $A$ 是射影的。
\end{lemma}

\begin{proof}
这由引理 \ref{groupoids-lemma-algebraic-quasi-projective} 以及
《态射进阶》引理 \ref{more-morphisms-lemma-projective} 得出。
\end{proof}

\begin{lemma}
\label{groupoids-lemma-abelian-variety-change-field}
设 $k$ 为域，$A$ 为 $k$ 上的阿贝尔簇。
对任意域扩张 $K/k$，基变换 $A_K$ 是 $K$ 上的阿贝尔簇。
\end{lemma}

\begin{proof}
证明略。注意，这正是我们坚持要求 $A$ 几何整的原因；
若无此条件，本引理以及下文许多结论都会错误。
\end{proof}

\begin{lemma}
\label{groupoids-lemma-abelian-variety-smooth}
设 $k$ 为域，$A$ 为 $k$ 上的阿贝尔簇。则 $A$ 在 $k$ 上光滑。
\end{lemma}

\begin{proof}
若 $k$ 完美，则该结论由
引理 \ref{groupoids-lemma-group-scheme-characteristic-zero-smooth}（特征零）
以及引理 \ref{groupoids-lemma-reduced-group-scheme-prefect-field-characteristic-p-smooth}
（正特征）得出。一般情形可利用光滑性的下降
（《下降》引理 \ref{descent-lemma-descending-property-smooth}）
以及引理 \ref{groupoids-lemma-abelian-variety-change-field} 转到完美闭包，
从而约化到上述情形。
\end{proof}

\begin{lemma}
\label{groupoids-lemma-abelian-variety-abelian}
阿贝尔簇是交换群概形；也就是说，群律可交换。
\end{lemma}

\begin{proof}
设 $k$ 为域，$A$ 为 $k$ 上的阿贝尔簇。
由引理 \ref{groupoids-lemma-abelian-variety-change-field}，可把 $k$
替换为其代数闭包。考虑态射
$$
h : A \times_k A \longrightarrow A \times_k A,\quad
(x, y) \longmapsto (x, xyx^{-1}y^{-1})
$$
两边都取第一投影时，这是 $A$ 上的态射。设 $e \in A(k)$ 为单位元。
于是 $h|_{e \times A}$ 是取常值 $(e, e)$ 的映射。根据
《态射进阶》引理 \ref{more-morphisms-lemma-flat-proper-family-cannot-collapse-fibre}，
存在 $e$ 的开邻域 $U \subset A$，使得 $h|_{U \times A}$
分解通过某个在 $U$ 上有限的 $Z \subset U \times A$。
这意味着，对 $x \in U(k)$，态射
$A \to A$，$y \mapsto xyx^{-1}y^{-1}$ 只取有限多个值。
当然，这意味着它恒取值 $e$。所以
$(x, y) \mapsto xyx^{-1}y^{-1}$ 在 $U \times A$ 上恒取值 $e$，
从而 $A$ 上的群律可交换。
\end{proof}

\begin{lemma}
\label{groupoids-lemma-apply-cube}
设 $k$ 为域，$A$ 为 $k$ 上的阿贝尔簇，并设 $\mathcal{L}$
为可逆 $\mathcal{O}_A$-模。则在 $A \times_k A \times_k A$ 上
存在可逆模的同构
$$
m_{1, 2, 3}^*\mathcal{L} \otimes
m_1^*\mathcal{L} \otimes
m_2^*\mathcal{L} \otimes
m_3^*\mathcal{L} \cong
m_{1, 2}^*\mathcal{L} \otimes
m_{1, 3}^*\mathcal{L} \otimes
m_{2, 3}^*\mathcal{L}
$$
其中 $m_{i_1, \ldots, i_t} : A \times_k A \times_k A \to A$
是态射 $(x_1, x_2, x_3) \mapsto \sum x_{i_j}$。
\end{lemma}

\begin{proof}
把立方定理（《态射进阶》定理 \ref{more-morphisms-theorem-of-the-cube}）应用于差
$$
\mathcal{M} =
m_{1, 2, 3}^*\mathcal{L} \otimes
m_1^*\mathcal{L} \otimes
m_2^*\mathcal{L} \otimes
m_3^*\mathcal{L} \otimes
m_{1, 2}^*\mathcal{L}^{\otimes -1} \otimes
m_{1, 3}^*\mathcal{L}^{\otimes -1} \otimes
m_{2, 3}^*\mathcal{L}^{\otimes -1}
$$
即可。其理由是，$\mathcal{M}$ 在
$A \times A \times e = A \times A$ 上的限制等于
$$
n_{1, 2}^*\mathcal{L} \otimes
n_1^*\mathcal{L} \otimes
n_2^*\mathcal{L} \otimes
n_{1, 2}^*\mathcal{L}^{\otimes -1} \otimes
n_1^*\mathcal{L}^{\otimes -1} \otimes
n_2^*\mathcal{L}^{\otimes -1} \cong \mathcal{O}_{A \times_k A}
$$
其中 $n_{i_1, \ldots, i_t} : A \times_k A \to A$
为态射 $(x_1, x_2) \mapsto \sum x_{i_j}$。
对 $A \times e \times A$ 和 $e \times A \times A$ 亦然。
\end{proof}

\begin{lemma}
\label{groupoids-lemma-pullbacks-by-n}
设 $k$ 为域，$A$ 为 $k$ 上的阿贝尔簇，并设 $\mathcal{L}$
为可逆 $\mathcal{O}_A$-模。则
$$
[n]^*\mathcal{L} \cong
\mathcal{L}^{\otimes n(n + 1)/2} \otimes
([-1]^*\mathcal{L})^{\otimes n(n - 1)/2}
$$
其中 $[n] : A \to A$ 把 $x$ 映到有 $n$ 个加项的
$x + x + \ldots + x$，而 $[-1] : A \to A$ 是 $A$ 的取逆映射。
\end{lemma}

\begin{proof}
考虑态射 $A \to A \times_k A \times_k A$，
$x \mapsto (x, x, -x)$，其中 $-x = [-1](x)$。
拉回引理 \ref{groupoids-lemma-apply-cube} 中的关系，得到
$$
\mathcal{L} \otimes
\mathcal{L} \otimes
\mathcal{L} \otimes
[-1]^*\mathcal{L} \cong
[2]^*\mathcal{L}
$$
这证明了 $n = 2$ 的情形。归纳假设结论对
$1, 2, \ldots, n$ 成立。再考虑态射
$A \to A \times_k A \times_k A$，$x \mapsto (x, x, [n - 1]x)$。
拉回引理 \ref{groupoids-lemma-apply-cube} 中的关系，得到
$$
[n + 1]^*\mathcal{L} \otimes
\mathcal{L} \otimes
\mathcal{L} \otimes
[n - 1]^*\mathcal{L} \cong
[2]^*\mathcal{L} \otimes
[n]^*\mathcal{L} \otimes
[n]^*\mathcal{L}
$$
再作初等算术运算即得结论。
\end{proof}

\begin{lemma}
\label{groupoids-lemma-degree-multiplication-by-d}
设 $k$ 为域，$A$ 为 $k$ 上的阿贝尔簇，并设
$[d] : A \to A$ 为乘以 $d$ 的映射。则 $[d]$ 有限局部自由，
次数为 $d^{2\dim(A)}$。
\end{lemma}

\begin{proof}
由引理 \ref{groupoids-lemma-abelian-variety-projective}
（以及《态射进阶》引理 \ref{more-morphisms-lemma-projective}），$A$ 有一个充足可逆模
$\mathcal{L}$。因为 $[-1] : A \to A$ 是自同构，
$[-1]^*\mathcal{L}$ 也是充足可逆 $\mathcal{O}_A$-模。
因此 $\mathcal{N} = \mathcal{L} \otimes [-1]^*\mathcal{L}$ 充足；
见《概形的性质》引理 \ref{properties-lemma-ample-tensor-globally-generated}。
由于 $\mathcal{N} \cong [-1]^*\mathcal{N}$，由
引理 \ref{groupoids-lemma-pullbacks-by-n} 有
$[d]^*\mathcal{N} \cong \mathcal{N}^{\otimes d^2}$。

\medskip\noindent
为导出矛盾，设 $C$ 是包含在 $[d]$ 的某一纤维中的固有曲线。
此时 $\mathcal{N}^{\otimes d^2}|_C \cong \mathcal{O}_C$
是次数为 $0$ 的充足可逆 $\mathcal{O}_C$-模，这例如与
《代数簇》引理 \ref{varieties-lemma-ample-curve} 矛盾。
（也可使用《代数簇》引理 \ref{varieties-lemma-ample-positive}。）
因此 $[d]$ 的每个纤维维数都是 $0$，故例如由
《概形上同调》引理 \ref{coherent-lemma-characterize-finite}，
$[d]$ 有限。此外，由引理 \ref{groupoids-lemma-abelian-variety-smooth}，
$A$ 在 $k$ 上光滑，因而根据
《交换代数》引理 \ref{algebra-lemma-CM-over-regular-flat}，
$[d] : A \to A$ 平坦（这里还使用了域上光滑的概形正则，以及
正则环是 Cohen--Macaulay 环；见
《代数簇》引理 \ref{varieties-lemma-smooth-regular} 和
《交换代数》引理 \ref{algebra-lemma-regular-ring-CM}）。
所以由《概形态射》引理 \ref{morphisms-lemma-finite-flat}，
$[d]$ 有限局部自由。

\medskip\noindent
最后计算次数。由《代数簇》引理 \ref{varieties-lemma-degree-finite-morphism-in-terms-degrees}，
$$
\deg_{\mathcal{N}^{\otimes d^2}}(A) = \deg([d]) \deg_\mathcal{N}(A)
$$
由于 $A$ 关于 $\mathcal{N}^{\otimes d^2}$（相应地，
$\mathcal{N}$）的次数，是下列多项式中 $n^{\dim(A)}$ 的系数：
$$
n \longmapsto \chi(A, \mathcal{N}^{\otimes nd^2}),\quad
\text{分别}\quad n \longmapsto \chi(A, \mathcal{N}^{\otimes n})
$$
所以 $\deg([d]) = d^{2 \dim(A)}$。
\end{proof}

\begin{lemma}
\label{groupoids-lemma-abelian-variety-multiplication-by-d-etale}
\begin{slogan}
阿贝尔簇上乘以一个整数的态射为平展态射，当且仅当该整数在基域中可逆。
\end{slogan}
设 $k$ 为域，$A$ 为 $k$ 上的非零阿贝尔簇。
则 $[d] : A \to A$ 为平展（\'etale）态射，当且仅当 $d$ 在 $k$ 中可逆。
\end{lemma}

\begin{proof}
注意 $[d](x + y) = [d](x) + [d](y)$。由于由一点作平移是 $A$
的自同构，$[d] : A \to A$ 为平展（\'etale）的点集要么为空，
要么等于 $A$（略去一些细节）。因此，只需检验 $[d]$
在单位元 $e \in A(k)$ 处是否平展（\'etale）。
已知 $[d]$ 有限局部自由
（引理 \ref{groupoids-lemma-degree-multiplication-by-d}），故它在 $e$ 处
平展（\'etale），等价于证明
$\text{d}[d] : T_{A/k, e} \to T_{A/k, e}$ 是单射；见
《代数簇》引理 \ref{varieties-lemma-injective-tangent-spaces-unramified}
以及《概形态射》引理 \ref{morphisms-lemma-flat-unramified-etale}。
由引理 \ref{groupoids-lemma-group-scheme-addition-tangent-vectors}，
$\text{d}[d]$ 是 $T_{A/k, e}$ 上乘以 $d$ 的映射。
\end{proof}

\begin{lemma}
\label{groupoids-lemma-abelian-variety-multiplication-by-p}
设 $k$ 为特征 $p > 0$ 的域，$A$ 为 $k$ 上维数为 $g$ 的阿贝尔簇。
则 $[p] : A \to A$ 在 $0$ 上的纤维至多有 $p^g$ 个不同的点。
\end{lemma}

\begin{proof}
可将并且确实将 $k$ 替换为其代数闭包。由
引理 \ref{groupoids-lemma-group-scheme-addition-tangent-vectors}，
$[p]$ 的导数作为映射 $T_{A/k, e} \to T_{A/k, e}$
是乘以 $p$，因而为零（比较引理 \ref{groupoids-lemma-abelian-variety-multiplication-by-d-etale} 的证明）。
由于 $[p]$ 与平移可交换，$[p]$ 的导数处处为零；也就是说，诱导映射
$[p]^*\Omega_{A/k} \to \Omega_{A/k}$ 为零。
考察泛点可知，相应的函数域映射
$[p]^* : k(A) \to k(A)$ 在 $\Omega_{k(A)/k}$ 上诱导零映射。
设 $t_1, \ldots, t_g$ 为 $k(A)$ 在 $k$ 上的一个 p-基
（《代数进阶》定义 \ref{more-algebra-definition-p-basis}
以及引理 \ref{more-algebra-lemma-p-basis}）。
由《交换代数》引理 \ref{algebra-lemma-derivative-zero-pth-power}，
$[p]^*(t_i)$ 有一个 $p$ 次根。因此，
$k(A)[x_1, \ldots, x_g]/(x_1^p - t_1, \ldots, x_g^p - t_g)$
是 $[p]^* : k(A) \to k(A)$ 的一个子扩张。
于是可找到仿射开集 $U \subset A$，使得
$t_i \in \mathcal{O}_A(U)$ 且
$x_i \in \mathcal{O}_A([p]^{-1}(U))$。得到 $[p]$ 在 $U$ 上的分解
$$
[p]^{-1}(U)
\xrightarrow{\pi_1}
\Spec(\mathcal{O}(U)[x_1, \ldots, x_g]/(x_1^p - t_1, \ldots, x_g^p - t_g))
\xrightarrow{\pi_2}
U
$$
缩小 $U$ 后，可假设 $\pi_1$ 有限局部自由
（例如使用一般平坦性；实际上在本情形中它已经有限局部自由）。
由引理 \ref{groupoids-lemma-degree-multiplication-by-d}，$[p]$ 的次数为
$p^{2g}$。由于 $\pi_2$ 的次数为 $p^g$，$\pi_1$ 的次数也为 $p^g$。
态射 $\pi_2$ 是万有同胚，因而其纤维都是单点集。
所以 $[p]^{-1}(U) \to U$ 的集合论纤维就是 $\pi_1$ 的纤维，
至多有 $p^g$ 个元素。由于 $[p]$ 是 $k$ 上的群概形同态，
对每个 $a \in A(k)$，$[p] : A(k) \to A(k)$ 的纤维基数都相同，
证明完毕。
\end{proof}

\begin{proposition}
\label{groupoids-proposition-review-abelian-varieties}
\begin{reference}
\cite{AVar} 对此作了精彩阐释。
\end{reference}
设 $A$ 为域 $k$ 上的阿贝尔簇。则
\begin{enumerate}
\item $A$ 在 $k$ 上射影；
\item $A$ 是交换群概形；
\item 对所有 $n \geq 1$，态射 $[n] : A \to A$ 都是满射；
\item 若 $k$ 代数闭，则 $A(k)$ 是可除阿贝尔群；
\item $A[n] = \Ker([n] : A \to A)$ 是 $k$ 上次数为
$n^{2\dim A}$ 的有限群概形；
\item $A[n]$ 在 $k$ 上平展（\'etale），当且仅当 $n \in k^*$；
\item 若 $n \in k^*$ 且 $k$ 代数闭，则
$A(k)[n] \cong (\mathbf{Z}/n\mathbf{Z})^{\oplus 2\dim(A)}$；
\item 若 $k$ 是特征 $p > 0$ 的代数闭域，则存在整数
$0 \leq f \leq \dim(A)$，使得对所有 $m \geq 1$ 都有
$A(k)[p^m] \cong (\mathbf{Z}/p^m\mathbf{Z})^{\oplus f}$。
\end{enumerate}
\end{proposition}

\begin{proof}
(1) 由引理 \ref{groupoids-lemma-abelian-variety-projective} 得出。
(2) 由引理 \ref{groupoids-lemma-abelian-variety-abelian} 得出。
(3) 由引理 \ref{groupoids-lemma-degree-multiplication-by-d} 得出。
若 $k$ 代数闭，则 $k$ 上簇的满态射在 $k$-有理点上诱导满射，
故由 (3) 得 (4)。
(5) 由引理 \ref{groupoids-lemma-degree-multiplication-by-d} 以及如下事实得出：
次数为 $N$ 的有限局部自由态射经基变换后，仍是次数为 $N$
的有限局部自由态射。(6) 由
引理 \ref{groupoids-lemma-abelian-variety-multiplication-by-d-etale} 得出。
具体地，若 $n$ 在 $k$ 中可逆，则 $[n]$ 平展（\'etale），
因而 $A[n]$ 在 $k$ 上平展（\'etale）。
另一方面，若 $n$ 在 $k$ 中不可逆，则 $[n]$ 在 $e$ 处不平展
（\'etale），从而 $A[n]$ 在 $k$ 上的 $e$ 处不平展（\'etale）
（使用《概形态射》引理 \ref{morphisms-lemma-flat-unramified-etale} 和
\ref{morphisms-lemma-set-points-where-fibres-unramified}）。

\medskip\noindent
假设 $k$ 代数闭，令 $g = \dim(A)$。证明 (7)。
设 $\ell$ 是在 $k$ 中可逆的素数。于是
$$
A[\ell](k) = A(k)[\ell]
$$
是被 $\ell$ 零化、阶为 $\ell^{2g}$ 的有限阿贝尔群。
由有限阿贝尔群结构定理，它同构于
$(\mathbf{Z}/\ell\mathbf{Z})^{2g}$。接着考虑短正合列
$$
0 \to A(k)[\ell] \to A(k)[\ell^2] \xrightarrow{\ell} A(k)[\ell] \to 0
$$
用与上面相同的论证可得
$A(k)[\ell^2] \cong (\mathbf{Z}/\ell^2\mathbf{Z})^{2g}$。
对指数归纳，得到
$A(k)[\ell^m] \cong (\mathbf{Z}/\ell^m\mathbf{Z})^{2g}$。
对与 $k$ 的特征互素的合数 $n$，取各主分量，即得到
$A(k)$ 中 $n$-挠子群的正确形式。(8) 的证明完全相同；
使用引理 \ref{groupoids-lemma-abelian-variety-multiplication-by-p} 所给出的
$A(k)[p] \cong (\mathbf{Z}/p\mathbf{Z})^{\oplus f}$，
其中某个 $0 \leq f \leq g$。
\end{proof}

\begin{remark}
\label{groupoids-remark-16-equivalent}
设 $k$ 为域。阿贝尔簇有 $2 \times 4 \times 2 = 16$ 个等价定义。
考虑以下三组性质：
\begin{itemize}
\item 射影、固有；
\item 几何不可约、不可约、几何连通、连通；
\item 光滑、几何约化。
\end{itemize}
从每组各选一个性质，并设 $A$ 为 $k$ 上的群概形，
且这些所选性质都在 $k$ 上成立，
则 $A$ 是阿贝尔簇。若选取“固有、几何不可约、几何约化”，
便恢复定义 \ref{groupoids-definition-abelian-variety}
（使用《代数簇》引理 \ref{varieties-lemma-geometrically-integral}）。
最弱的选择是“固有、连通、几何约化”；例如参见
《概形态射》引理 \ref{morphisms-lemma-locally-projective-proper} 和
《代数簇》引理 \ref{varieties-lemma-smooth-geometrically-normal}。
因此，设 $A$ 是 $k$ 上固有、连通且几何约化的群概形。
由引理 \ref{groupoids-lemma-connected-group-scheme-over-field-irreducible} 和
\ref{groupoids-lemma-group-scheme-field-geometrically-irreducible}，
$A$ 几何不可约，因而是阿贝尔簇。
最后，若 $A/k$ 是阿贝尔簇，则由命题 \ref{groupoids-proposition-review-abelian-varieties}，它在 $k$ 上射影且光滑，
因而满足最强的选择“射影、几何不可约、光滑”。
\end{remark}







\section{群概形的作用}
\label{groupoids-section-action-group-scheme}

\noindent
设 $(G, m)$ 为群，$V$ 为集合。回顾 $G$ 在 $V$ 上的一个
{\it （左）作用}由映射 $a : G \times V \to V$ 给出，并满足
\begin{enumerate}
\item （结合律）对所有 $g, g' \in G$ 和 $v \in V$，
$a(m(g, g'), v) = a(g, a(g', v))$；
\item （单位元）对所有 $v \in V$，$a(e, v) = v$。
\end{enumerate}
也称 $V$ 是一个{\it $G$-集}（这通常意味着从记号中省略 $a$，
是一种记号滥用）。{\it $G$-集的映射}
$\psi : V \to V'$ 是满足如下条件的任意集合映射：
对所有 $v \in V$ 都有 $\psi(a(g, v)) = a(g, \psi(v))$。

\begin{definition}
\label{groupoids-definition-action-group-scheme}
设 $S$ 为概形，$(G, m)$ 为 $S$ 上的群概形。
\begin{enumerate}
\item {\it $G$ 在概形 $X/S$ 上的作用}是 $S$ 上的态射
$a : G \times_S X \to X$，使得对每个 $T/S$，映射
$a : G(T) \times X(T) \to X(T)$ 都在 $X(T)$ 上定义
$G(T)$-集结构。
\item 设 $X$、$Y$ 是 $S$ 上各自带有 $G$-作用的概形。
一个{\it 等变}态射，更准确地说，一个{\it $G$-等变态射}
$\psi : X \to Y$，是 $S$ 上的概形态射，并使得对每个 $T/S$，
$\psi : X(T) \to Y(T)$ 都是 $G(T)$-集的态射。
\end{enumerate}
\end{definition}

\noindent
在情形 (1) 中，这表示下列诸图交换：
\begin{equation}
\label{groupoids-equation-action}
\vcenter{
\xymatrix{
G \times_S G \times_S X \ar[r]_-{1_G \times a} \ar[d]_{m \times 1_X} &
G \times_S X \ar[d]^a \\
G \times_S X \ar[r]^a & X
}
}
\quad\quad
\vcenter{
\xymatrix{
G \times_S X \ar[r]_-a & X \\
X\ar[u]^{e \times 1_X} \ar[ru]_{1_X}
}
}
\end{equation}
在情形 (2) 中，这仅表示下图交换：
$$
\xymatrix{
G \times_S X \ar[r]_-{\text{id} \times \psi} \ar[d]_a &
G \times_S Y \ar[d]^a \\
X \ar[r]^\psi & Y
}
$$

\begin{definition}
\label{groupoids-definition-free-action}
设 $S$、$G \to S$、$X \to S$ 如
定义 \ref{groupoids-definition-action-group-scheme} 所述。
设 $a : G \times_S X \to X$ 是 $G$ 在 $X/S$ 上的作用。
若对每个 $S$ 上的概形 $T$，作用
$a : G(T) \times X(T) \to X(T)$ 都是群 $G(T)$
在集合 $X(T)$ 上的自由作用，则称该作用是{\it 自由的}。
\end{definition}

\begin{lemma}
\label{groupoids-lemma-free-action}
情形如定义 \ref{groupoids-definition-free-action}。
作用 $a$ 自由，当且仅当
$$
G \times_S X \to X \times_S X, \quad (g, x) \mapsto (a(g, x), x)
$$
是单态射。
\end{lemma}

\begin{proof}
由定义立即可得。
\end{proof}






\section{主齐性空间}
\label{groupoids-section-principal-homogeneous}

\noindent
在《位点上同调》定义 \ref{sites-cohomology-definition-torsor} 中，我们定义了一个位点上
群层的扭子。设 $\tau \in \{Zariski, \linebreak[0] \etale,
\linebreak[0] smooth, \linebreak[0] syntomic, \linebreak[0] fppf\}$
为一种拓扑，而 $(G, m)$ 是 $S$ 上的群概形。
由于 $\tau$ 强于典范拓扑（见《下降》引理 \ref{descent-lemma-fpqc-universal-effective-epimorphisms}），
$\underline{G}$（见《位点与层》定义 \ref{sites-definition-representable-sheaf}）是
$(\Sch/S)_\tau$ 上的群层。因此，我们已经知道
$(\Sch/S)_\tau$ 上 $\underline{G}$-扭子的含义。
若该层可表示，便出现一种特殊情形。下述定义直接说明：
表示它的概形何时为 $G$-扭子。

\begin{definition}
\label{groupoids-definition-pseudo-torsor}
设 $S$ 为概形，$(G, m)$ 为 $S$ 上的群概形，$X$ 为 $S$ 上的概形，
并设 $a : G \times_S X \to X$ 是 $G$ 在 $X$ 上的作用。
\begin{enumerate}
\item 若诱导的 $S$ 上概形态射
$G \times_S X \to X \times_S X$，
$(g, x) \mapsto (a(g, x), x)$ 是同构，则称 $X$ 是
{\it 伪 $G$-扭子}，或称 $X$ {\it 在 $G$ 下形式主齐}。
\item 若存在 $S$ 上的 $G$-等变同构 $G \to X$，其中 $G$
通过左乘作用于 $G$，则称伪 $G$-扭子 $X$ 是{\it 平凡的}。
\end{enumerate}
\end{definition}

\noindent
显然，若 $S' \to S$ 是概形态射，则 $S$ 上伪 $G$-扭子的拉回
$X_{S'}$ 是 $S'$ 上的伪 $G_{S'}$-扭子。

\begin{lemma}
\label{groupoids-lemma-characterize-trivial-pseudo-torsors}
情形如定义 \ref{groupoids-definition-pseudo-torsor}。
\begin{enumerate}
\item 概形 $X$ 是伪 $G$-扭子，当且仅当对 $S$ 上的每个概形 $T$，
集合 $X(T)$ 要么为空，要么群 $G(T)$ 在 $X(T)$ 上的作用是单传递的。
\item 伪 $G$-扭子 $X$ 平凡，当且仅当态射 $X \to S$ 有截面。
\end{enumerate}
\end{lemma}

\begin{proof}
证明略。
\end{proof}

\begin{definition}
\label{groupoids-definition-principal-homogeneous-space}
设 $S$ 为概形，$(G, m)$ 为 $S$ 上的群概形，$X$ 为
$S$ 上的伪 $G$-扭子。
\begin{enumerate}
\item 若存在 fpqc 覆盖\footnote{这表示默认类型的扭子是在某个
fpqc 覆盖上平凡的伪扭子。这是
\cite[Expos\'e IV, 6.5]{SGA3} 中的定义。由于我们最常在 fppf
拓扑中工作，这一定义对我们略有不便。}
$\{S_i \to S\}_{i \in I}$，使得每个
$X_{S_i} \to S_i$ 都有截面（即为平凡伪 $G_{S_i}$-扭子），
则称 $X$ 是{\it 主齐性空间}或{\it $G$-扭子}。
\item 设 $\tau \in \{Zariski, \etale, smooth, syntomic, fppf\}$。
若存在 $\tau$-覆盖 $\{S_i \to S\}_{i \in I}$，使得每个
$X_{S_i} \to S_i$ 都有截面，则称 $X$ 是
{\it $\tau$ 拓扑中的 $G$-扭子}、{\it $\tau$ $G$-扭子}，
或简称{\it $\tau$ 扭子}。
\item 若 $X$ 是 $G$-扭子，并且它是 \'etale 拓扑的扭子，
则称其{\it 拟等平凡}。
\item 若 $X$ 是 $G$-扭子，并且它是 Zariski 拓扑的扭子，
则称其{\it 局部平凡}。
\end{enumerate}
\end{definition}

\noindent
我们有时说“设 $X$ 是 $S$ 上的 $G$-扭子”，意指 $X$ 是
$S$ 上带有 $G$-作用的概形，并且该作用使它成为 $S$ 上的主齐性空间。
下面说明，当两种说法都适用时，这与先前引入的记号一致。

\begin{lemma}
\label{groupoids-lemma-torsor}
设 $S$ 为概形，$(G, m)$ 为 $S$ 上的群概形，$X$ 为 $S$ 上的概形，
并设 $a : G \times_S X \to X$ 是 $G$ 在 $X$ 上的作用。
设 $\tau \in \{Zariski, \etale, smooth, syntomic, fppf\}$。
则 $X$ 是 $\tau$-拓扑中的 $G$-扭子，当且仅当
$\underline{X}$ 是 $(\Sch/S)_\tau$ 上的 $\underline{G}$-扭子。
\end{lemma}

\begin{proof}
证明略。
\end{proof}

\begin{remark}
\label{groupoids-remark-fun-with-torsors}
设 $(G, m)$ 为概形 $S$ 上的群概形。在这种情形下，自然会问：
\begin{enumerate}
\item 若 $X \to S$ 是伪 $G$-扭子且 $X \to S$ 满射，
则 $X$ 是否必为 $G$-扭子？
\item $(\Sch/S)_{fppf}$ 上的每个 $\underline{G}$-扭子是否可表示？
换言之，每个 $\underline{G}$-扭子是否都来自某个 fppf $G$-扭子？
\item 每个 $G$-扭子是否都是
fppf（相应地，\ smooth；相应地，\ \'etale；相应地，\ Zariski）扭子？
\end{enumerate}
一般而言，这些问题的答案都是否定的。要得到肯定答案，需要对
$G \to S$ 施加附加条件。例如：若 $S$ 是一个域的谱，则 (1)
的答案为肯定，因为此时 $\{X \to S\}$ 是使 $X$ 平凡化的 fpqc 覆盖。
若 $G \to S$ 仿射，则 (2) 的答案为肯定
（由《下降》引理 \ref{descent-lemma-affine} 得出）。
若 $G = \text{GL}_{n, S}$，则 (3) 的答案为肯定；
事实上，任意 $\text{GL}_{n, S}$-扭子都局部平凡
（由《下降》引理 \ref{descent-lemma-finite-locally-free-descends} 得出）。
\end{remark}



\section{等变拟凝聚层}
\label{groupoids-section-equivariant}

\noindent
我们把“函数”看作与“空间”对偶。因此，空间的态射所诱导的函数映射
方向相反。此外，也把模层的截面看作“函数”。这自然导向下述定义中
所选的箭头方向。

\begin{definition}
\label{groupoids-definition-equivariant-module}
设 $S$ 为概形，$(G, m)$ 为 $S$ 上的群概形，并设
$a : G \times_S X \to X$ 是群概形 $G$ 在 $X/S$ 上的作用。
一个{\it $G$-等变拟凝聚 $\mathcal{O}_X$-模}，或简称
{\it 等变拟凝聚 $\mathcal{O}_X$-模}，是二元组
$(\mathcal{F}, \alpha)$，其中 $\mathcal{F}$ 是拟凝聚
$\mathcal{O}_X$-模，而 $\alpha$ 是
$\mathcal{O}_{G \times_S X}$-模同态
$$
\alpha : a^*\mathcal{F} \longrightarrow \text{pr}_1^*\mathcal{F}
$$
这里 $\text{pr}_1 : G \times_S X \to X$ 是投影，并且满足：
\begin{enumerate}
\item 下图在 $\mathcal{O}_{G \times_S G \times_S X}$-模范畴中交换：
$$
\xymatrix{
(1_G \times a)^*\text{pr}_1^*\mathcal{F} \ar[r]_-{\text{pr}_{12}^*\alpha} &
\text{pr}_2^*\mathcal{F} \\
(1_G \times a)^*a^*\mathcal{F} \ar[u]^{(1_G \times a)^*\alpha} \ar@{=}[r] &
(m \times 1_X)^*a^*\mathcal{F} \ar[u]_{(m \times 1_X)^*\alpha}
}
$$
\item 拉回
$$
(e \times 1_X)^*\alpha : \mathcal{F} \longrightarrow \mathcal{F}
$$
是恒等映射。
\end{enumerate}
有关说明，可与 Equation (\ref{groupoids-equation-action}) 中的相应图比较。
\end{definition}

\noindent
注意，第一个图的交换性保证 $(e \times 1_X)^*\alpha$
是 $\mathcal{F}$ 上的幂等算子；因此条件 (2) 仅仅要求它是同构。

\begin{lemma}
\label{groupoids-lemma-pullback-equivariant}
设 $S$ 为概形，$G$ 为 $S$ 上的群概形。
设 $f : Y \to X$ 是两个带 $G$-作用的 $S$-概形之间的
$G$-等变态射。规则
$(\mathcal{F}, \alpha) \mapsto (f^*\mathcal{F}, (1_G \times f)^*\alpha)$
定义了从 $G$-等变拟凝聚 $\mathcal{O}_X$-模范畴到
$G$-等变拟凝聚 $\mathcal{O}_Y$-模范畴的函子。
\end{lemma}

\begin{proof}
证明略。
\end{proof}

\noindent
下面给出一个例子。

\begin{example}
\label{groupoids-example-Gm-on-affine}
设 $A$ 为 $\mathbf{Z}$-分次环；也就是说，$A$ 带有直和分解
$A = \bigoplus_{n \in \mathbf{Z}} A_n$，且
$A_n \cdot A_m \subset A_{n + m}$。
令 $X = \Spec(A)$。环同态
$\mu : A \to A \otimes \mathbf{Z}[x, x^{-1}]$，
$f \mapsto f \otimes x^{\deg(f)}$，给出一个 $\mathbf{G}_m$-作用
$$
a : \mathbf{G}_m \times X \longrightarrow X
$$
。确切地说，只需验证
$$
\xymatrix{
A \ar[r]_\mu \ar[d]_\mu &
A \otimes \mathbf{Z}[x, x^{-1}] \ar[d]^{\mu \otimes 1} \\
A \otimes \mathbf{Z}[x, x^{-1}] \ar[r]^-{1 \otimes m} &
A \otimes \mathbf{Z}[x, x^{-1}] \otimes \mathbf{Z}[x, x^{-1}]
}
$$
其中 $m(x) = x \otimes x$；见例 \ref{groupoids-example-multiplicative-group}。在齐次元上求值便立即可见。
设 $M$ 是分次 $A$-模。使用定义 \ref{groupoids-definition-equivariant-module} 中的 $\alpha$，即可得到
$\mathbf{G}_m$-等变拟凝聚 $\mathcal{O}_X$-模
$\mathcal{F} = \widetilde{M}$；该映射对应于
$A \otimes \mathbf{Z}[x, x^{-1}]$-模同态
$$
M \otimes_{A, \mu} (A \otimes \mathbf{Z}[x, x^{-1}])
\longrightarrow
M \otimes_{A, \text{id}_A \otimes 1} (A \otimes \mathbf{Z}[x, x^{-1}])
$$
它对齐次的 $m \in M$，把 $m \otimes 1 \otimes 1$
映到 $m \otimes 1 \otimes x^{\deg(m)}$。
\end{example}

\begin{lemma}
\label{groupoids-lemma-complete-reducibility-Gm}
设 $a : \mathbf{G}_m \times X \to X$ 是仿射概形上的作用。
则 $X$ 是某个 $\mathbf{Z}$-分次环的谱，且该作用如
例 \ref{groupoids-example-Gm-on-affine} 所述。
\end{lemma}

\begin{proof}
设 $f \in A = \Gamma(X, \mathcal{O}_X)$。可以唯一地写成有限和
$$
a^\sharp(f) = \sum\nolimits_{n \in \mathbf{Z}} f_n \otimes x^n
\quad\text{于}\quad
A \otimes \mathbf{Z}[x, x^{-1}] =
\Gamma(\mathbf{G}_m \times X, \mathcal{O}_{\mathbf{G}_m \times X})
$$
其中 $f_n$ 属于 $A$。于是得到映射 $A \to A$，$f \mapsto f_n$。
由于 $a$ 是作用，在 $x = 1$ 处求值可得 $f = \sum f_n$。
再由 $a$ 是作用，有
$$
\sum (f_n)_m \otimes x^m \otimes x^n = \sum f_n x^n \otimes x^n
$$
（比较例 \ref{groupoids-example-Gm-on-affine} 中的计算）。
所以当 $n \not = m$ 时 $(f_n)_m = 0$，且 $(f_n)_n = f_n$。
令
$$
A_n = \{f \in A \mid f_n = f\}
$$
则 $A = \sum A_n$。另一方面，在上述情形中，$f = 0$
蕴含 $f_n = 0$，所以该和必为直和。
\end{proof}

\begin{lemma}
\label{groupoids-lemma-Gm-equivariant-module}
设 $A$ 为分次环，$X = \Spec(A)$ 带有例 \ref{groupoids-example-Gm-on-affine} 中的作用
$a : \mathbf{G}_m \times X \to X$。
设 $\mathcal{F}$ 为 $\mathbf{G}_m$-等变拟凝聚
$\mathcal{O}_X$-模。则 $M = \Gamma(X, \mathcal{F})$
具有典范分次，使其成为分次 $A$-模，并使同构
$\widetilde{M} \to \mathcal{F}$
（《概形》引理 \ref{schemes-lemma-quasi-coherent-affine}）
成为 $\mathbf{G}_m$-等变模的同构；其中 $\widetilde{M}$
上的 $\mathbf{G}_m$-等变结构来自
例 \ref{groupoids-example-Gm-on-affine}。
\end{lemma}

\begin{proof}
可以对模 $M$ 重复引理 \ref{groupoids-lemma-complete-reducibility-Gm}
的论证。也可以考虑概形
$(X', \mathcal{O}_{X'}) = (X, \mathcal{O}_X \oplus \mathcal{F})$，
其中把 $\mathcal{F}$ 看成平方为零的理想。
利用 $X$ 上以及 $\mathcal{F}$ 上的作用，可定义自然作用
$a' : \mathbf{G}_m \times X' \to X'$。然后对 $X'$ 应用
引理 \ref{groupoids-lemma-complete-reducibility-Gm} 即得结论。
（这一论证的优点是，它立即说明 $A$ 和 $M$ 上的分次相容；
也就是说，$M$ 是分次 $A$-模。）
细节略。
\end{proof}



\section{群胚}
\label{groupoids-section-groupoids}

\noindent
回顾群胚是每个态射均为同构的范畴；见
《范畴》定义 \ref{categories-definition-groupoid}。
因此，群胚有对象集 $\text{Ob}$、箭头集 $\text{Arrows}$、
{\it 源}与{\it 靶}映射
$s, t : \text{Arrows} \to \text{Ob}$，以及{\it 合成律}
$c : \text{Arrows} \times_{s, \text{Ob}, t} \text{Arrows}
\to \text{Arrows}$。这些映射恰好满足下列公理：
\begin{enumerate}
\item （结合律）作为映射
$\text{Arrows} \times_{s, \text{Ob}, t}
\text{Arrows} \times_{s, \text{Ob}, t}
\text{Arrows} \to \text{Arrows}$，有
$c \circ (1, c) = c \circ (c, 1)$；
\item （恒等）存在映射 $e : \text{Ob} \to \text{Arrows}$，使得
\begin{enumerate}
\item 作为映射 $\text{Ob} \to \text{Ob}$，有
$s \circ e = t \circ e = \text{id}$；
\item 作为映射 $\text{Arrows} \to \text{Arrows}$，有
$c \circ (1, e \circ s) = c \circ (e \circ t, 1) = 1$；
\end{enumerate}
\item （取逆）存在映射 $i : \text{Arrows} \to \text{Arrows}$，使得
\begin{enumerate}
\item 作为映射 $\text{Arrows} \to \text{Ob}$，有
$s \circ i = t$、$t \circ i = s$；且
\item 作为映射 $\text{Arrows} \to \text{Arrows}$，有
$c \circ (1, i) = e \circ t$ 以及
$c \circ (i, 1) = e \circ s$。
\end{enumerate}
\end{enumerate}
若这些条件成立，则映射 $e$ 和 $i$ 唯一确定，且 $i$ 是双射。
注意，若 $(\text{Ob}', \text{Arrows}', s', t', c')$
是另一个群胚范畴，则函子
$f : (\text{Ob}, \text{Arrows}, s, t, c) \to
(\text{Ob}', \text{Arrows}', s', t', c')$
由一对集合映射 $f : \text{Ob} \to \text{Ob}'$ 与
$f : \text{Arrows} \to \text{Arrows}'$ 给出，它们满足
$s' \circ f = f \circ s$、$t' \circ f = f \circ t$ 以及
$c' \circ (f, f) = f \circ c$。与恒等和取逆的相容性自动成立。
下文将使用这一点。（警告：对一般范畴，必须另行要求与恒等的相容性。）

\begin{definition}
\label{groupoids-definition-groupoid}
设 $S$ 为概形。
\begin{enumerate}
\item 一个{\it $S$ 上的群胚概形}，或简称
{\it $S$ 上的群胚}，是五元组 $(U, R, s, t, c)$，其中
$U$ 和 $R$ 是 $S$ 上的概形，而
$s, t : R \to U$ 以及 $c : R \times_{s, U, t} R \to R$
是 $S$ 上的概形态射，并满足如下性质：对任意 $S$ 上的概形
$T$，五元组
$$
(U(T), R(T), s, t, c)
$$
都是上述意义下的群胚范畴。
\item 一个{\it $S$ 上群胚概形的态射
$f : (U, R, s, t, c) \to (U', R', s', t', c')$}
由概形态射 $f : U \to U'$ 与 $f : R \to R'$ 给出，
并满足如下性质：对任意 $S$ 上的概形 $T$，映射 $f$
定义了从群胚范畴 $(U(T), R(T), s, t, c)$ 到群胚范畴
$(U'(T), R'(T), s', t', c')$ 的函子。
\end{enumerate}
\end{definition}

\noindent
设 $(U, R, s, t, c)$ 为 $S$ 上的群胚。
由定义前的讨论以及米田引理，存在唯一的 $S$ 上概形态射
$e : U \to R$ 和 $i : R \to R$，使得对每个 $S$ 上的概形 $T$，
诱导映射 $e : U(T) \to R(T)$ 是恒等，而
$i : R(T) \to R(T)$ 是群胚范畴中的取逆。
七元组 $(U, R, s, t, c, e, i)$ 满足对应于上述公理
(1)、(2)(a)、(2)(b)、(3)(a)、(3)(b) 的各交换图。
反之，具有这一性质的七元组，其五元组 $(U, R, s, t, c)$
是群胚概形。注意，$i$ 是同构，而 $e$ 同时是 $s$ 和 $t$ 的截面。
此外，对一个给定的 $S$ 上群胚概形，记
$$
j = (t, s) : R \longrightarrow U \times_S U
$$
这与上文第 \ref{groupoids-section-equivalence-relations} 节中的约定相容。
我们有时说“设 $(U, R, s, t, c, e, i)$ 为 $S$ 上的群胚”，
以强调恒等与取逆的存在。

\begin{lemma}
\label{groupoids-lemma-groupoid-pre-equivalence}
给定 $S$ 上的群胚概形 $(U, R, s, t, c)$，态射
$j : R \to U \times_S U$ 是预等价关系。
\end{lemma}

\begin{proof}
证明略。这是一个很好的定义练习。
\end{proof}

\begin{lemma}
\label{groupoids-lemma-equivalence-groupoid}
给定 $S$ 上的等价关系 $j : R \to U \times_S U$，
存在唯一的方式把它扩张为 $S$ 上的群胚 $(U, R, s, t, c)$。
\end{lemma}

\begin{proof}
证明略。这是一个很好的定义练习。
\end{proof}

\begin{lemma}
\label{groupoids-lemma-diagram}
设 $S$ 为概形，$(U, R, s, t, c)$ 为 $S$ 上的群胚。
在交换图
$$
\xymatrix{
& U & \\
R \ar[d]_s \ar[ru]^t &
R \times_{s, U, t} R
\ar[l]^-{\text{pr}_0} \ar[d]^{\text{pr}_1} \ar[r]_-c &
R \ar[d]^s \ar[lu]_t \\
U & R \ar[l]_t \ar[r]^s & U
}
$$
中，下面两个方块都是纤维积方块。此外，上面的三角形
（实际上也是方块）同样是笛卡尔的。
\end{lemma}

\begin{proof}
证明略。这是关于定义以及代数几何函子观点的练习。
\end{proof}

\begin{lemma}
\label{groupoids-lemma-diagram-pull}
设 $S$ 为概形，$(U, R, s, t, c, e, i)$ 为 $S$ 上的群胚。图
\begin{equation}
\label{groupoids-equation-pull}
\xymatrix{
R \times_{t, U, t} R
\ar@<1ex>[r]^-{\text{pr}_1} \ar@<-1ex>[r]_-{\text{pr}_0}
\ar[d]_{(\text{pr}_0, c \circ (i, 1))} &
R \ar[r]^t \ar[d]^{\text{id}_R} &
U \ar[d]^{\text{id}_U} \\
R \times_{s, U, t} R
\ar@<1ex>[r]^-c \ar@<-1ex>[r]_-{\text{pr}_0} \ar[d]_{\text{pr}_1} &
R \ar[r]^t \ar[d]^s &
U \\
R \ar@<1ex>[r]^s \ar@<-1ex>[r]_t &
U
}
\end{equation}
交换。给出的竖直映射使上面两行彼此同构。
左下方的两个方块都是笛卡尔的。
\end{lemma}

\begin{proof}
该图的交换性由群胚公理得出。用群胚的语言说，左上方的竖直箭头
把具有相同靶的一对态射 $(\alpha, \beta)$ 映到一对态射
$(\alpha, \alpha^{-1} \circ \beta)$。在任意群胚中，这定义了
$\text{Arrows} \times_{t, \text{Ob}, t} \text{Arrows}$
与
$\text{Arrows} \times_{s, \text{Ob}, t} \text{Arrows}$
之间的双射，故引理的第二个断言成立。最后一个断言由
引理 \ref{groupoids-lemma-diagram} 得出。
\end{proof}

\begin{lemma}
\label{groupoids-lemma-base-change-groupoid}
设 $(U, R, s, t, c)$ 为概形 $S$ 上的群胚，$S' \to S$ 为态射。
令 $U' = S' \times_S U$、$R' = S' \times_S R$，并赋予它们
态射 $s, t, c$ 的基变换 $s'$、$t'$、$c'$。
则 $(U', R', s', t', c')$ 是 $S'$ 上的群胚概形，而诸投影决定
$S$ 上群胚概形的态射
$(U', R', s', t', c') \to (U, R, s, t, c)$。
\end{lemma}

\begin{proof}
证明略。提示：
$R' \times_{s', U', t'} R' = S' \times_S (R \times_{s, U, t} R)$。
\end{proof}









\section{群胚上的拟凝聚层}
\label{groupoids-section-groupoids-quasi-coherent}

\noindent
关于箭头方向的约定，见第 \ref{groupoids-section-equivariant} 节的引言。

\begin{definition}
\label{groupoids-definition-groupoid-module}
设 $S$ 为概形，$(U, R, s, t, c)$ 为 $S$ 上的群胚概形。
所谓 $(U, R, s, t, c)$ 上的一个{\it 拟凝聚模}，是指一对
$(\mathcal{F}, \alpha)$，其中 $\mathcal{F}$ 是拟凝聚
$\mathcal{O}_U$-模，而 $\alpha$ 是 $\mathcal{O}_R$-模态射
$$
\alpha : t^*\mathcal{F} \longrightarrow s^*\mathcal{F}
$$
满足：
\begin{enumerate}
\item 图
$$
\xymatrix{
& \text{pr}_1^*t^*\mathcal{F} \ar[r]_-{\text{pr}_1^*\alpha} &
\text{pr}_1^*s^*\mathcal{F} \ar@{=}[rd] & \\
\text{pr}_0^*s^*\mathcal{F} \ar@{=}[ru] & & & c^*s^*\mathcal{F} \\
& \text{pr}_0^*t^*\mathcal{F} \ar[lu]^{\text{pr}_0^*\alpha} \ar@{=}[r] &
c^*t^*\mathcal{F} \ar[ru]_{c^*\alpha}
}
$$
在 $\mathcal{O}_{R \times_{s, U, t} R}$-模范畴中交换；
\item 拉回
$$
e^*\alpha : \mathcal{F} \longrightarrow \mathcal{F}
$$
是恒等态射。
\end{enumerate}
可与引理 \ref{groupoids-lemma-diagram} 中的交换图比较。
\end{definition}

\noindent
第一个图的交换性迫使算子 $e^*\alpha$ 幂等。因此，第二个条件也可表述为
$e^*\alpha$ 是同构。事实上，该条件蕴含 $\alpha$ 是同构。

\begin{lemma}
\label{groupoids-lemma-isomorphism}
设 $S$ 为概形，$(U, R, s, t, c)$ 为 $S$ 上的群胚概形。
若 $(\mathcal{F}, \alpha)$ 是 $(U, R, s, t, c)$ 上的拟凝聚模，
则 $\alpha$ 是同构。
\end{lemma}

\begin{proof}
沿态射 $(i, 1) : R \to R \times_{s, U, t} R$ 拉回
定义 \ref{groupoids-definition-groupoid-module} 中的交换图。
由此得到 $i^*\alpha \circ \alpha = s^*e^*\alpha$。
再沿态射 $(1, i)$ 拉回，得到关系
$\alpha \circ i^*\alpha = t^*e^*\alpha$。由第二个假设，
这些态射都是恒等态射。因此 $i^*\alpha$ 是 $\alpha$ 的逆。
\end{proof}

\begin{lemma}
\label{groupoids-lemma-pullback}
设 $S$ 为概形，并考虑 $S$ 上群胚概形的态射
$f : (U, R, s, t, c) \to (U', R', s', t', c')$。
则由
$$
(\mathcal{F}, \alpha) \mapsto (f^*\mathcal{F}, f^*\alpha)
$$
给出的拉回 $f^*$，定义了从 $(U', R', s', t', c')$ 上的拟凝聚层范畴
到 $(U, R, s, t, c)$ 上的拟凝聚层范畴的函子。
\end{lemma}

\begin{proof}
证明略。
\end{proof}

\begin{lemma}
\label{groupoids-lemma-pushforward}
设 $S$ 为概形，并考虑 $S$ 上群胚概形的态射
$f : (U, R, s, t, c) \to (U', R', s', t', c')$。
假设：
\begin{enumerate}
\item $f : U \to U'$ 拟紧且拟分离；
\item 方块
$$
\xymatrix{
R \ar[d]_t \ar[r]_f & R' \ar[d]^{t'} \\
U \ar[r]^f & U'
}
$$
是笛卡尔的；
\item $s'$ 和 $t'$ 平坦。
\end{enumerate}
则由
$$
(\mathcal{F}, \alpha) \mapsto (f_*\mathcal{F}, f_*\alpha)
$$
给出的推前 $f_*$，定义了从 $(U, R, s, t, c)$ 上的拟凝聚层范畴
到 $(U', R', s', t', c')$ 上的拟凝聚层范畴的函子；它是
引理 \ref{groupoids-lemma-pullback} 所定义的拉回的右伴随。
\end{lemma}

\begin{proof}
由于 $U \to U'$ 拟紧且拟分离，$f_*$ 把拟凝聚层变为拟凝聚层
（《概形》引理 \ref{schemes-lemma-push-forward-quasi-coherent}）。
又因方块
$$
\vcenter{
\xymatrix{
R \ar[d]_t \ar[r]_f & R' \ar[d]^{t'} \\
U \ar[r]^f & U'
}
}
\quad\text{且}\quad
\vcenter{
\xymatrix{
R \ar[d]_s \ar[r]_f & R' \ar[d]^{s'} \\
U \ar[r]^f & U'
}
}
$$
均为笛卡尔的，故
$(t')^*f_*\mathcal{F} = f_*t^*\mathcal{F}$ 以及
$(s')^*f_*\mathcal{F} = f_*s^*\mathcal{F}$，见
《概形上同调》引理 \ref{coherent-lemma-flat-base-change-cohomology}。
所以可把 $f_*\alpha$ 视为态射
$(t')^*f_*\mathcal{F} \to (s')^*f_*\mathcal{F}$。
类似论证表明 $f_*\alpha$ 满足上循环条件。
该函子与拉回函子伴随，因为环空间上模的拉回与推前彼此伴随。
其余细节略去。
\end{proof}

\begin{lemma}
\label{groupoids-lemma-colimits}
设 $S$ 为概形，$(U, R, s, t, c)$ 为 $S$ 上的群胚概形。
则 $(U, R, s, t, c)$ 上的拟凝聚模范畴有余极限。
\end{lemma}

\begin{proof}
设 $i \mapsto (\mathcal{F}_i, \alpha_i)$ 是以 $\mathcal{I}$ 为指标范畴的图。
可以形成余极限
$\mathcal{F} = \colim \mathcal{F}_i$；
它是 $U$ 上的拟凝聚层，见
《概形》第 \ref{schemes-section-quasi-coherent} 节。
由于余极限与拉回交换，有
$s^*\mathcal{F} = \colim s^*\mathcal{F}_i$，类似地还有
$t^*\mathcal{F} = \colim t^*\mathcal{F}_i$。因此可以令
$\alpha = \colim \alpha_i$。我们略去证明：在
$(U, R, s, t, c)$ 上的拟凝聚模范畴中，$(\mathcal{F}, \alpha)$
确为该图的余极限。
\end{proof}

\begin{lemma}
\label{groupoids-lemma-abelian}
设 $S$ 为概形，$(U, R, s, t, c)$ 为 $S$ 上的群胚概形。
若 $s$、$t$ 平坦，则 $(U, R, s, t, c)$ 上的拟凝聚模范畴
是阿贝尔范畴。
\end{lemma}

\begin{proof}
设
$\varphi : (\mathcal{F}, \alpha) \to (\mathcal{G}, \beta)$
是 $(U, R, s, t, c)$ 上拟凝聚模的同态。由于 $s$ 平坦，
$$
0 \to s^*\Ker(\varphi)
\to s^*\mathcal{F} \to s^*\mathcal{G} \to s^*\Coker(\varphi) \to 0
$$
正合；沿 $t$ 拉回亦然。因此 $\alpha$ 和 $\beta$ 诱导同构
$\kappa : t^*\Ker(\varphi) \to s^*\Ker(\varphi)$ 以及
$\lambda : t^*\Coker(\varphi) \to s^*\Coker(\varphi)$，
且二者满足上循环条件。直接验证可知，
$(\Ker(\varphi), \kappa)$ 和
$(\Coker(\varphi), \lambda)$ 分别是
$(U, R, s, t, c)$ 上拟凝聚模范畴中的核与余核。
此外，条件 $\Coim(\varphi) = \Im(\varphi)$ 成立，
因为它在 $U$ 上成立。
\end{proof}













\section{拟凝聚模的余极限}
\label{groupoids-section-colimits}

\noindent
本节证明若干技术性结果：在适当假设下，群胚上的每个拟凝聚模
都是“较小”拟凝聚模的滤过余极限。

\begin{lemma}
\label{groupoids-lemma-construct-quasi-coherent}
设 $(U, R, s, t, c)$ 为 $S$ 上的群胚概形。
假设 $s, t$ 平坦、拟紧且拟分离。对于 $U$ 上任意拟凝聚模
$\mathcal{G}$，存在典范同构
$\alpha : t^*s_*t^*\mathcal{G} \to s^*s_*t^*\mathcal{G}$，
使 $(s_*t^*\mathcal{G}, \alpha)$ 成为 $(U, R, s, t, c)$ 上的拟凝聚模。
这一构造定义函子
$$
\QCoh(\mathcal{O}_U) \longrightarrow \QCoh(U, R, s, t, c)
$$
且该函子是遗忘函子
$(\mathcal{F}, \beta) \mapsto \mathcal{F}$ 的右伴随。
\end{lemma}

\begin{proof}
拟凝聚模沿拟紧拟分离态射的推前仍是拟凝聚的，见
《概形》引理 \ref{schemes-lemma-push-forward-quasi-coherent}。
因此 $s_*t^*\mathcal{G}$ 拟凝聚。采用引理 \ref{groupoids-lemma-diagram}
中的记号，有
$$
t^*s_*t^*\mathcal{G} =
\text{pr}_{1, *}\text{pr}_0^*t^*\mathcal{G} =
\text{pr}_{1, *}c^*t^*\mathcal{G} =
s^*s_*t^*\mathcal{G}
$$
中间的等式源于态射
$R \times_{s, U, t} R \to U$ 满足
$t \circ c = t \circ \text{pr}_0$；首尾两个等式则源于这些步骤中
基变换与推前交换，见《概形上同调》引理 \ref{coherent-lemma-flat-base-change-cohomology}。

\medskip\noindent
为了验证 $\alpha$ 满足定义 \ref{groupoids-definition-groupoid-module}
中的上循环条件以及伴随性质，我们换一种方式描述构造
$\mathcal{G} \mapsto (s_*t^*\mathcal{G}, \alpha)$。
考虑由 $R$ 上等价关系 $R \times_{t, U, t} R$ 所关联的群胚概形
$(R, R \times_{t, U, t} R, \text{pr}_0, \text{pr}_1, \text{pr}_{02})$，
见引理 \ref{groupoids-lemma-equivalence-groupoid}。存在群胚概形态射
$$
f :
(R, R \times_{t, U, t} R, \text{pr}_1, \text{pr}_0, \text{pr}_{02})
\longrightarrow
(U, R, s, t, c)
$$
它在对象上由 $s : R \to U$ 给出，在箭头上由
$R \times_{t, U, t} R \to R$，
$(r_0, r_1) \mapsto r_0^{-1} \circ r_1$ 给出；
所需诸图交换性的验证从略。由于
$t, s : R \to U$ 拟紧、拟分离且平坦，并且由
引理 \ref{groupoids-lemma-diagram-pull} 有笛卡尔方块
$$
\xymatrix{
R \times_{t, U, t} R \ar[d]_{\text{pr}_0}
\ar[rr]_-{(r_0, r_1) \mapsto r_0^{-1} \circ r_1} & & R \ar[d]^t \\
R \ar[rr]^s & & U
}
$$
故引理 \ref{groupoids-lemma-pushforward} 适用于 $f$。
于是拟凝聚模沿 $f$ 的推前与拉回互为伴随函子。
为完成证明，我们把这些函子与上面描述的函子识别起来。注意
$$
t^* :
\QCoh(\mathcal{O}_U)
\longrightarrow
\QCoh(R, R \times_{t, U, t} R, \text{pr}_1, \text{pr}_0, \text{pr}_{02})
$$
是范畴等价；这是拟凝聚层的下降理论，因为
$\{t : R \to U\}$ 是 fpqc 覆盖，见
《下降》命题 \ref{descent-proposition-fpqc-descent-quasi-coherent}。

\medskip\noindent
沿 $f$ 推前并预复合等价 $t^*$，会把
$\mathcal{G}$ 送到 $(s_*t^*\mathcal{G}, \alpha)$。
我们略去验证以此方式得到的同构 $\alpha$ 与上面构造的同构相同。

\medskip\noindent
沿 $f$ 拉回并后复合等价 $t^*$ 的逆，会把
$(\mathcal{F}, \beta)$ 送到模 $s^*\mathcal{F}$
相对于 $\{t : R \to U\}$ 的下降；这里该模配备
$R \times_{t, U, t} R$ 上的下降数据
$\gamma$，它是 $\beta$ 沿
$(r_0, r_1) \mapsto r_0^{-1} \circ r_1$ 的拉回。
考虑同构 $\beta : t^*\mathcal{F} \to s^*\mathcal{F}$。
$t^*\mathcal{F}$ 相对于 $\{t : R \to U\}$ 的典范下降数据
（《下降》定义 \ref{descent-definition-descent-datum-effective-quasi-coherent}）
经 $\beta$ 转化为态射
$$
\text{pr}_0^*s^*\mathcal{F}
\xrightarrow{\text{pr}_0^*\beta^{-1}}
\text{pr}_0^*t^*\mathcal{F}
\xrightarrow{can}
\text{pr}_1^*t^*\mathcal{F}
\xrightarrow{\text{pr}_1^*\beta}
\text{pr}_1^*s^*\mathcal{F}
$$
由于 $\beta$ 满足上循环条件，此态射等于 $\beta$ 沿
$(r_0, r_1) \mapsto r_0^{-1} \circ r_1$ 的拉回。
为此，只需取定义 \ref{groupoids-definition-groupoid-module}
中实际的上循环关系，并沿态射
$(\text{pr}_0, c \circ (i, 1)) : R \times_{t, U, t} R \to R \times_{s, U, t} R$
将其拉回；该态射也出现在
引理 \ref{groupoids-lemma-diagram-pull} 的交换图中。因此
$(s^*\mathcal{F}, \gamma)$ 同构于 $(t^*\mathcal{F}, can)$。
综上，沿 $f$ 拉回再后复合等价 $t^*$ 的逆所得函子，
同构于遗忘函子
$(\mathcal{F}, \beta) \mapsto \mathcal{F}$。
\end{proof}

\begin{remark}
\label{groupoids-remark-adjunction-map}
在引理 \ref{groupoids-lemma-construct-quasi-coherent} 的情形下，记遗忘函子为
$$
F : \QCoh(U, R, s, t, c) \to \QCoh(\mathcal{O}_U),\quad
(\mathcal{F}, \beta) \mapsto \mathcal{F}
$$
并记该引理中构造的右伴随为
$$
G : \QCoh(\mathcal{O}_U) \to \QCoh(U, R, s, t, c),\quad
\mathcal{G} \mapsto (s_*t^*\mathcal{G}, \alpha)
$$
则伴随的单位 $\eta : \text{id} \to G \circ F$
在 $(\mathcal{F}, \beta)$ 上取值为态射
$$
\mathcal{F} \to s_*s^*\mathcal{F} \xrightarrow{\beta^{-1}} s_*t^*\mathcal{F}
$$
验证从略。
\end{remark}

\begin{lemma}
\label{groupoids-lemma-push-pull}
设 $f : Y \to X$ 为概形态射，$\mathcal{F}$ 为拟凝聚
$\mathcal{O}_X$-模，$\mathcal{G}$ 为拟凝聚
$\mathcal{O}_Y$-模，并设
$\varphi : \mathcal{G} \to f^*\mathcal{F}$ 为模态射。假设：
\begin{enumerate}
\item $\varphi$ 是单射；
\item $f$ 拟紧、拟分离、平坦且满射；
\item $X$, $Y$ 局部 Noether；
\item $\mathcal{G}$ 是凝聚 $\mathcal{O}_Y$-模。
\end{enumerate}
则由拉回
$$
\xymatrix{
\mathcal{F} \ar[r] & f_*f^*\mathcal{F} \\
\mathcal{F} \cap f_*\mathcal{G} \ar[u] \ar[r] &
f_*\mathcal{G} \ar[u]
}
$$
定义的 $\mathcal{F} \cap f_*\mathcal{G}$ 是凝聚
$\mathcal{O}_X$-模。
\end{lemma}

\begin{proof}
我们将直接使用《概形上同调》引理 \ref{coherent-lemma-coherent-Noetherian} 对凝聚模的刻画，
以及凝聚模构成 $\QCoh(\mathcal{O}_X)$ 的 Serre 子范畴这一事实；见
《概形上同调》引理 \ref{coherent-lemma-coherent-Noetherian-quasi-coherent-sub-quotient}。
若 $f$ 有截面 $\sigma$，则
$\mathcal{F} \cap f_*\mathcal{G}$ 包含于态射
$\sigma^*\mathcal{G} \to \sigma^*f^*\mathcal{F} = \mathcal{F}$
的像中，因而是凝聚的。一般地，为证明
$\mathcal{F} \cap f_*\mathcal{G}$ 凝聚，只需证明
$f^*(\mathcal{F} \cap f_*\mathcal{G})$ 凝聚（见
《下降》引理 \ref{descent-lemma-finite-type-descends}）。
由于 $f$ 平坦，它等于
$f^*\mathcal{F} \cap f^*f_*\mathcal{G}$。
又因 $f$ 平坦、拟紧且拟分离，故
$f^*f_*\mathcal{G} = p_*q^*\mathcal{G}$，其中
$p, q : Y \times_X Y \to Y$ 是投影；见
《概形上同调》引理 \ref{coherent-lemma-flat-base-change-cohomology}。
由于 $p$ 有截面，结论成立。
\end{proof}

\noindent
设 $S$ 为概形，$(U, R, s, t, c)$ 为 $S$ 上的概形群胚。
假设 $U$ 局部 Noether。在下面的引理中，如果
$(U, R, s, t, c)$ 上拟凝聚层 $(\mathcal{F}, \alpha)$ 的
$\mathcal{O}_U$-模 $\mathcal{F}$ 是凝聚的，我们就称它是
{\it 凝聚的}。

\begin{lemma}
\label{groupoids-lemma-colimit-coherent}
设 $(U, R, s, t, c)$ 为 $S$ 上的群胚概形。假设：
\begin{enumerate}
\item $U$, $R$ 是 Noether 概形；
\item $s, t$ 平坦、拟紧且拟分离。
\end{enumerate}
则 $(U, R, s, t, c)$ 上的每个拟凝聚模
$(\mathcal{F}, \beta)$ 都是凝聚模的滤过余极限。
\end{lemma}

\begin{proof}
我们将不再特别说明而使用《概形上同调》引理 \ref{coherent-lemma-coherent-Noetherian} 对局部 Noether 概形上凝聚模的刻画。
由《概形上同调》引理 \ref{coherent-lemma-directed-colimit-coherent}，可将
$\mathcal{F} = \colim \mathcal{H}_i$ 写成凝聚子模
$\mathcal{H}_i \subset \mathcal{F}$ 的滤过余极限。
给定 $U$ 上拟凝聚层 $\mathcal{H}$，以
$(s_*t^*\mathcal{H}, \alpha)$ 表示
引理 \ref{groupoids-lemma-construct-quasi-coherent} 给出的
$(U, R, s, t, c)$ 上拟凝聚层。考虑
$\QCoh(U, R, s, t, c)$ 中的伴随态射
$(\mathcal{F}, \beta) \to (s_*t^*\mathcal{F}, \alpha)$，
见注 \ref{groupoids-remark-adjunction-map}。令
$$
(\mathcal{F}_i, \beta_i) =
(\mathcal{F}, \beta)
\times_{(s_*t^*\mathcal{F}, \alpha)}
(s_*t^*\mathcal{H}_i, \alpha)
$$
这是在 $\QCoh(U, R, s, t, c)$ 中取的纤维积。根据
引理 \ref{groupoids-lemma-abelian} 的证明，限制到 $U$ 是
$\QCoh(U, R, s, t, c)$ 上的正合函子，故得到拉回图
$$
\xymatrix{
\mathcal{F} \ar[r] & s_*t^*\mathcal{F} \\
\mathcal{F}_i \ar[r] \ar[u] & s_*t^*\mathcal{H}_i \ar[u]
}
$$
换言之，$\mathcal{F}_i = \mathcal{F} \cap s_*t^*\mathcal{H}_i$。
由注 \ref{groupoids-remark-adjunction-map} 中对伴随态射的描述，
该图同构于
$$
\xymatrix{
\mathcal{F} \ar[r] & s_*s^*\mathcal{F} \\
\mathcal{F}_i \ar[r] \ar[u] & s_*t^*\mathcal{H}_i \ar[u]
}
$$
其中右侧竖直箭头由对态射
$$
t^*\mathcal{H}_i \to t^*\mathcal{F} \xrightarrow{\beta} s^*\mathcal{F}
$$
施以 $s_*$ 得到。由于 $t$ 平坦，该箭头是单射。因此由
引理 \ref{groupoids-lemma-push-pull}，$\mathcal{F}_i$ 凝聚。
最后，因为 $s$ 拟紧且拟分离，$s_*$ 与余极限交换
（见《概形上同调》引理 \ref{coherent-lemma-colimit-cohomology}）。故
$s_*t^*\mathcal{F} = \colim s_*t^*\mathcal{H}_i$，进而如所需有
$(\mathcal{F}, \beta) = \colim (\mathcal{F}_i, \beta_i)$。
\end{proof}

\noindent
下面这个颇为有趣的引理在处理域上的群胚时很有用。事实上，
它正是证明代数群的任意表示都是有限维表示之余极限的标准论证。

\begin{lemma}
\label{groupoids-lemma-colimit-finite-type}
设 $(U, R, s, t, c)$ 为 $S$ 上的群胚概形。假设：
\begin{enumerate}
\item $U$, $R$ 仿射；
\item 存在 $e_i \in \mathcal{O}_R(R)$，使得每个元素
$g \in \mathcal{O}_R(R)$ 都能唯一写成
$\sum s^*(f_i)e_i$，其中 $f_i \in \mathcal{O}_U(U)$。
\end{enumerate}
则 $(U, R, s, t, c)$ 上每个拟凝聚模
$(\mathcal{F}, \alpha)$ 都是有限型拟凝聚模的滤过余极限。
\end{lemma}

\begin{proof}
该假设意指 $\mathcal{O}_R(R)$ 经由 $s$ 成为以 $e_i$ 为基的自由
$\mathcal{O}_U(U)$-模。因此对任意拟凝聚
$\mathcal{O}_U$-模 $\mathcal{G}$，有
$s^*\mathcal{G}(R) = \bigoplus_i \mathcal{G}(U)e_i$。
我们用 $s(-)$ 表示截面沿 $s$ 的拉回，其他态射亦采用类似记号。
设 $(\mathcal{F}, \alpha)$ 是 $(U, R, s, t, c)$ 上的拟凝聚模，
并取 $\sigma \in \mathcal{F}(U)$。由上可唯一写成
$$
\alpha(t(\sigma)) = \sum s(\sigma_i) e_i
$$
其中 $\sigma_i \in \mathcal{F}(U)$（当然，除有限多个外均为零）。
还可把 $R \times_{s, U, t}R$ 上的函数写成
$$
c(e_i) = \sum \text{pr}_1(f_{ij}) \text{pr}_0(e_j)
$$
定义 \ref{groupoids-definition-groupoid-module} 中图的交换性意味着
$$
\sum \text{pr}_1(\alpha(t(\sigma_i))) \text{pr}_0(e_i)
=
\sum \text{pr}_1(s(\sigma_i)f_{ij}) \text{pr}_0(e_j)
$$
（计算略）。取出 $\text{pr}_0(e_l)$ 的系数，得到
$\alpha(t(\sigma_l)) = \sum s(\sigma_i)f_{il}$。
于是，由诸元素 $\sigma_i$ 生成的子模
$\mathcal{G} \subset \mathcal{F}$ 定义一个由 $\alpha$ 保持的
有限型拟凝聚模，因而它是 $\QCoh(U, R, s, t, c)$ 中
$\mathcal{F}$ 的子对象。这个子模包含 $\sigma$
（将第一个关系沿 $e$ 拉回即可看出），故结论成立。
\end{proof}

\noindent
我们建议读者跳过本节的其余部分。设 $S$ 为概形，
$(U, R, s, t, c)$ 为 $S$ 上的概形群胚，$\kappa$ 为一个基数。
下文称 $(U, R, s, t, c)$ 上的拟凝聚层
$(\mathcal{F}, \alpha)$ 是 $\kappa$-生成的，是指
$\mathcal{F}$ 是 $\kappa$-生成的 $\mathcal{O}_U$-模；见
《概形的性质》定义 \ref{properties-definition-kappa-generated}。

\begin{lemma}
\label{groupoids-lemma-set-of-iso-classes}
设 $(U, R, s, t, c)$ 为 $S$ 上的群胚概形，$\kappa$ 为基数。
存在集合 $T$ 以及 $(U, R, s, t, c)$ 上一族
$\kappa$-生成拟凝聚模
$(\mathcal{F}_t, \alpha_t)_{t \in T}$，使得
$(U, R, s, t, c)$ 上每个 $\kappa$-生成拟凝聚模
都同构于某个 $(\mathcal{F}_t, \alpha_t)$。
\end{lemma}

\begin{proof}
对 $U$ 上每个拟凝聚模 $\mathcal{F}$，使
$(\mathcal{F}, \alpha)$ 成为 $(U, R, s, t, c)$ 上拟凝聚模的态射
$\alpha : t^*\mathcal{F} \to s^*\mathcal{F}$ 构成一个
（可能为空的）集合。由《概形的性质》引理 \ref{properties-lemma-set-of-iso-classes}，
$\kappa$-生成拟凝聚 $\mathcal{O}_U$-模的同构类构成一个集合。
\end{proof}

\begin{lemma}
\label{groupoids-lemma-colimit-kappa}
设 $(U, R, s, t, c)$ 为 $S$ 上的群胚概形，并假设 $s, t$ 平坦。
则存在基数 $\kappa$，使得 $(U, R, s, t, c)$ 上每个拟凝聚模
$(\mathcal{F}, \alpha)$ 都是其 $\kappa$-生成拟凝聚子模的有向余极限。
\end{lemma}

\begin{proof}
在引理的陈述及本证明中，拟凝聚模 $(\mathcal{F}, \alpha)$ 的
{\it 子模}，是指拟凝聚子模
$\mathcal{G} \subset \mathcal{F}$，满足作为
$s^*\mathcal{F}$ 的子层有
$\alpha(t^*\mathcal{G}) = s^*\mathcal{G}$。
这是有意义的，因为 $s, t$ 平坦，所以拉回 $s^*$ 与 $t^*$
正合，也就是说它们保持子层。本证明将重复
《概形的性质》引理 \ref{properties-lemma-colimit-kappa} 的证明。
我们建议读者先阅读该证明。

\medskip\noindent
选择仿射开覆盖 $U = \bigcup_{i \in I} U_i$。
对每一对 $i, j$，选择仿射开覆盖
$$
U_i \cap U_j = \bigcup\nolimits_{k \in I_{ij}} U_{ijk}
\quad\text{且}\quad
s^{-1}(U_i) \cap t^{-1}(U_j) = \bigcup\nolimits_{k \in J_{ij}} W_{ijk}.
$$
记 $U_i = \Spec(A_i)$、$U_{ijk} = \Spec(A_{ijk})$、
$W_{ijk} = \Spec(B_{ijk})$。令 $\kappa$ 为任意无限基数，
它 $\geq$ 任一集合 $I$, $I_{ij}$, $J_{ij}$ 的基数。

\medskip\noindent
设 $(\mathcal{F}, \alpha)$ 为 $(U, R, s, t, c)$ 上的拟凝聚模。
令 $M_i = \mathcal{F}(U_i)$、$M_{ijk} = \mathcal{F}(U_{ijk})$。
注意
$$
M_i \otimes_{A_i} A_{ijk} = M_{ijk} = M_j \otimes_{A_j} A_{ijk}
$$
并且 $\alpha$ 给出同构
$$
\alpha|_{W_{ijk}} :
M_i \otimes_{A_i, t} B_{ijk}
\longrightarrow
M_j \otimes_{A_j, s} B_{ijk}
$$
见《概形》引理 \ref{schemes-lemma-widetilde-pullback}。
利用选择公理，我们选择映射
$$
(i, j, k, m) \mapsto S(i, j, k, m)
$$
使它对每个 $i, j \in I$、$k \in I_{ij}$ 或 $k \in J_{ij}$
以及 $m \in M_i$，指定有限子集
$S(i, j, k, m) \subset M_j$，从而有
$$
m \otimes 1 = \sum\nolimits_{m' \in S(i, j, k, m)} m' \otimes a_{m'}
\quad\text{或}\quad
\alpha(m \otimes 1) = \sum\nolimits_{m' \in S(i, j, k, m)} m' \otimes b_{m'}
$$
在 $M_{ijk}$ 中成立，其中
$a_{m'} \in A_{ijk}$ 或 $b_{m'} \in B_{ijk}$。
此外，约定当 $k \in I_{ij}$ 时，对所有
$i, j = i, k, m$ 都有 $S(i, i, k, m) = \{m\}$。
固定这样的一族 $S(i, j, k, m)$。

\medskip\noindent
给定一族子集 $\mathcal{S} = (S_i)_{i \in I}$，
其中 $S_i \subset M_i$ 的基数至多为 $\kappa$，定义
$\mathcal{S}' = (S'_i)$，其中
$$
S'_j = \bigcup\nolimits_{(i, j, k, m)\text{ 使得 }m \in S_i}
S(i, j, k, m)
$$
注意 $S_i \subset S'_i$。由于 $S'_i$ 是在一个基数至多为
$\kappa$ 的集合上对有限集取并，其基数至多为 $\kappa$。
令 $\mathcal{S}^{(0)} = \mathcal{S}$、
$\mathcal{S}^{(1)} = \mathcal{S}'$，并归纳定义
$\mathcal{S}^{(n + 1)} = (\mathcal{S}^{(n)})'$。再令
$\mathcal{S}^{(\infty)} = \bigcup_{n \geq 0} \mathcal{S}^{(n)}$。
写作 $\mathcal{S}^{(\infty)} = (S^{(\infty)}_i)$，则对任意元素
$m \in S^{(\infty)}_i$，$m$ 在 $M_{ijk}$ 中的像可写成有限和
$\sum m' \otimes a_{m'}$，其中 $m' \in S_j^{(\infty)}$。
因此，令
$$
N_i = A_i\text{-子模，属于 }M_i\text{ 由下列对象生成： }S^{(\infty)}_i
$$
便有
$$
N_i \otimes_{A_i} A_{ijk} = N_j \otimes_{A_j} A_{ijk}
\quad\text{且}\quad
\alpha(N_i \otimes_{A_i, t} B_{ijk}) = N_j \otimes_{A_j, s} B_{ijk}
$$
这些等式分别作为 $M_{ijk}$ 或
$M_j \otimes_{A_j, s} B_{ijk}$ 的子模等式成立。
故存在拟凝聚子模
$\mathcal{G} \subset \mathcal{F}$，满足
$\mathcal{G}(U_i) = N_i$，且作为 $s^*\mathcal{F}$ 的子模有
$\alpha(t^*\mathcal{G}) = s^*\mathcal{G}$。
换言之，$(\mathcal{G}, \alpha|_{t^*\mathcal{G}})$ 是
$(\mathcal{F}, \alpha)$ 的子模。此外，依构造
$\mathcal{G}$ 是 $\kappa$-生成的。

\medskip\noindent
设 $\{(\mathcal{G}_t, \alpha_t)\}_{t \in T}$ 为
$(\mathcal{F}, \alpha)$ 的全体 $\kappa$-生成拟凝聚子模之集合。
若 $t, t' \in T$，则 $\mathcal{G}_t + \mathcal{G}_{t'}$
也是 $\kappa$-生成拟凝聚子模，因为它是态射
$\mathcal{G}_t \oplus \mathcal{G}_{t'} \to \mathcal{F}$ 的像。
因此，该系统按包含关系排序后是有向的。上述论证表明，
$\mathcal{F}$ 在 $U_i$ 上的每个截面都属于某个
$\mathcal{G}_t$（因为可从使该截面属于 $S_i$ 的
$\mathcal{S}$ 开始）。故
$\colim_t \mathcal{G}_t \to \mathcal{F}$ 既单又满，如所需。
\end{proof}





\section{群胚与群概形}
\label{groupoids-section-groupoids-group-schemes}

\noindent
由群 $G$ 在集合 $V$ 上的作用 $a$ 构造群胚有许多方式。
我们采用如下方式：把元素 $g \in G$ 看作一条以 $v$ 为源、
以 $a(g, v)$ 为靶的箭头。这就导出概形之群作用的下述构造。

\begin{lemma}
\label{groupoids-lemma-groupoid-from-action}
设 $S$ 为概形，$Y$ 为 $S$ 上的概形，$(G, m)$ 为 $Y$ 上的群概形，
其单位元为 $e_G$、逆为 $i_G$。设 $X/Y$ 为 $Y$ 上的概形，
$a : G \times_Y X \to X$ 是 $G$ 在 $X/Y$ 上的作用。
则可按如下方式得到 $S$ 上的群胚概形
$(U, R, s, t, c, e, i)$：
\begin{enumerate}
\item 令 $U = X$，$R = G \times_Y X$。
\item 令 $s : R \to U$ 为 $(g, x) \mapsto x$。
\item 令 $t : R \to U$ 为 $(g, x) \mapsto a(g, x)$。
\item 令 $c : R \times_{s, U, t} R \to R$ 为
$((g, x), (g', x')) \mapsto (m(g, g'), x')$。
\item 令 $e : U \to R$ 为 $x \mapsto (e_G(x), x)$。
\item 令 $i : R \to R$ 为 $(g, x) \mapsto (i_G(g), a(g, x))$。
\end{enumerate}
\end{lemma}

\begin{proof}
证明略。提示：只需在集合层面验证这一构造即可。
为此，使用引理上方的描述，把 $g$ 看成从 $v$ 到
$a(g, v)$ 的箭头。
\end{proof}

\begin{lemma}
\label{groupoids-lemma-action-groupoid-modules}
设 $S$ 为概形，$Y$ 为 $S$ 上的概形，$(G, m)$ 为 $Y$ 上的群概形。
设 $X$ 为 $Y$ 上的概形，$a : G \times_Y X \to X$
是 $G$ 在 $X$ 上且相对于 $Y$ 的作用。令 $(U, R, s, t, c)$ 为
引理 \ref{groupoids-lemma-groupoid-from-action} 构造的群胚概形。
规则 $(\mathcal{F}, \alpha) \mapsto (\mathcal{F}, \alpha)$
定义了 $G$-等变 $\mathcal{O}_X$-模范畴与
$(U, R, s, t, c)$ 上拟凝聚模范畴之间的范畴等价。
\end{lemma}

\begin{proof}
该断言是有意义的，因为作为态射
$R = G \times_Y X \to X$，有
$t = a$ 以及 $s = \text{pr}_1$；见
定义 \ref{groupoids-definition-equivariant-module} 和
\ref{groupoids-definition-groupoid-module}。
采用引理 \ref{groupoids-lemma-groupoid-from-action} 中的转译，
两个定义对交换性的要求完全一致。
\end{proof}





\section{稳定子群概形}
\label{groupoids-section-stabilizer}

\noindent
给定群胚概形，可如下得到一个群概形。

\begin{lemma}
\label{groupoids-lemma-groupoid-stabilizer}
设 $S$ 为概形，$(U, R, s, t, c)$ 为 $S$ 上的群胚。
由笛卡尔方块
$$
\xymatrix{
G \ar[r] \ar[d] & R \ar[d]^{j = (t, s)} \\
U \ar[r]^-{\Delta} & U \times_S U
}
$$
定义的概形 $G$ 是 $U$ 上的群概形，其乘法
$m$ 由复合律 $c$ 诱导。
\end{lemma}

\begin{proof}
这是因为在群胚范畴中，任一对象的自态射集构成一个群。
\end{proof}

\noindent
由于 $\Delta$ 是浸入，$G = j^{-1}(\Delta_{U/S})$
是 $R$ 的局部闭子概形。从这个角度看，结构态射
$j^{-1}(\Delta_{U/S}) \to U$ 可由 $s$ 或 $t$ 诱导
（二者相同），而 $m$ 由 $c$ 诱导。

\begin{definition}
\label{groupoids-definition-stabilizer-groupoid}
设 $S$ 为概形，$(U, R, s, t, c)$ 为 $S$ 上的群胚。
群概形 $j^{-1}(\Delta_{U/S})\to U$ 称为
群胚概形 $(U, R, s, t, c)$ 的{\it 稳定子}。
\end{definition}

\noindent
文献中常把稳定子群概形记作 $S$
（大概是因为 stabilizer 一词以字母 s 开头）；
这里不能这样做，因为我们已经用 $S$ 表示基概形。

\begin{lemma}
\label{groupoids-lemma-groupoid-action-stabilizer}
设 $S$ 为概形，$(U, R, s, t, c)$ 为 $S$ 上的群胚，
$G/U$ 为其稳定子。把 $R$ 经态射 $t : R \to U$
视为 $U$ 上的概形，记作 $R_t/U$。
存在由复合律 $c$ 诱导的典范左作用
$$
a : G \times_U R_t \longrightarrow R_t
$$
\end{lemma}

\begin{proof}
就 $T/S$ 上的点而言，定义 $a(g, r) = c(g, r)$。
\end{proof}

\begin{lemma}
\label{groupoids-lemma-groupoid-action-stabilizer-pseudo-torsor}
设 $S$ 为概形，$(U, R, s, t, c)$ 为 $S$ 上的群胚概形，
$G$ 为 $R$ 的稳定子群概形。令
$$
G_0 = G \times_{U, \text{pr}_0} (U \times_S U) = G \times_S U
$$
并把它视为 $U \times_S U$ 上的群概形。引理 \ref{groupoids-lemma-groupoid-action-stabilizer} 中 $G$ 在 $R$ 上的作用
诱导 $G_0$ 在 $R$ 上相对于 $U \times_S U$ 的作用，使 $R$
成为 $U \times_S U$ 上的伪 $G_0$-扭子。
\end{lemma}

\begin{proof}
这是因为在群胚范畴 $\mathcal{C}$ 中，集合
$\Mor_\mathcal{C}(x, y)$ 是群
$\Mor_\mathcal{C}(y, y)$ 下的主齐性集。
\end{proof}

\begin{lemma}
\label{groupoids-lemma-fibres-j}
设 $S$ 为概形，$(U, R, s, t, c)$ 为 $S$ 上的群胚概形。
设 $p \in U \times_S U$ 为一点，并以 $R_p$ 表示
$j = (t, s) : R \to U \times_S U$ 的概形论纤维。
若 $R_p \not = \emptyset$，则作用
$$
G_{0, \kappa(p)} \times_{\kappa(p)} R_p \longrightarrow R_p
$$
（见引理 \ref{groupoids-lemma-groupoid-action-stabilizer-pseudo-torsor}）
使 $R_p$ 成为 $\kappa(p)$ 上的 $G_{\kappa(p)}$-扭子。
\end{lemma}

\begin{proof}
由陈述中引用的引理，该作用是伪扭子。若 $R_p$ 不是空概形，
则 $\{R_p \to p\}$ 是使该伪扭子平凡化的 fpqc 覆盖。
\end{proof}







\section{群胚的限制}
\label{groupoids-section-restrict-groupoid}

\noindent
考虑一个（通常意义下的）群胚
$\mathcal{C} = (\text{Ob}, \text{Arrows}, s, t, c)$，
并假设给定集合映射 $g : \text{Ob}' \to \text{Ob}$。
可以构造群胚
$\mathcal{C}' = (\text{Ob}', \text{Arrows}', s', t', c')$：
把 $\text{Ob}'$ 中元素 $x', y'$ 之间的态射理解为
$\mathcal{C}$ 中 $g(x'), g(y')$ 之间的态射。换言之，令
$$
\text{Arrows}' =
\text{Ob}'
\times_{g, \text{Ob}, t}
\text{Arrows}
\times_{s, \text{Ob}, g}
\text{Ob}'.
$$
并以显然的方式选择 $s'$、$t'$ 和 $c'$。存在典范函子
$\mathcal{C}' \to \mathcal{C}$，它全忠实，但未必本质满。
配备函子 $\mathcal{C}' \to \mathcal{C}$ 的群胚
$\mathcal{C}'$ 称为群胚 $\mathcal{C}$ 在 $\text{Ob}'$ 上的
{\it 限制}。

\begin{lemma}
\label{groupoids-lemma-restrict-groupoid}
设 $S$ 为概形，$(U, R, s, t, c)$ 为 $S$ 上的群胚概形，
$g : U' \to U$ 为概形态射。考虑下图
$$
\xymatrix{
R' \ar[d] \ar[r] \ar@/_3pc/[dd]_{t'} \ar@/^1pc/[rr]^{s'}&
R \times_{s, U} U' \ar[r] \ar[d] &
U' \ar[d]^g \\
U' \times_{U, t} R \ar[d] \ar[r] &
R \ar[r]^s \ar[d]_t &
U \\
U' \ar[r]^g &
U
}
$$
其中所有方块都是纤维积方块。则存在典范复合律
$c' : R' \times_{s', U', t'} R' \to R'$，使
$(U', R', s', t', c')$ 成为 $S$ 上的群胚概形，并使
$U' \to U$, $R' \to R$ 定义 $S$ 上群胚概形的态射
$(U', R', s', t', c') \to (U, R, s, t, c)$。
此外，对 $S$ 上任意概形 $T$，群胚函子
$$
(U'(T), R'(T), s', t', c') \to (U(T), R(T), s, t, c)
$$
是 $(U(T), R(T), s, t, c)$ 经映射
$U'(T) \to U(T)$ 的限制（如上所述）。
\end{lemma}

\begin{proof}
证明略。
\end{proof}

\begin{definition}
\label{groupoids-definition-restrict-groupoid}
设 $S$ 为概形，$(U, R, s, t, c)$ 为 $S$ 上的群胚概形，
$g : U' \to U$ 为概形态射。引理 \ref{groupoids-lemma-restrict-groupoid} 中构造的群胚态射
$(U', R', s', t', c') \to (U, R, s, t, c)$
称为 $(U, R, s, t, c)$ 在 $U'$ 上的{\it 限制}。
在这种情形下，有时使用记号 $R' = R|_{U'}$。
\end{definition}

\begin{lemma}
\label{groupoids-lemma-restrict-groupoid-relation}
定义 \ref{groupoids-definition-restrict-groupoid} 与
\ref{groupoids-definition-restrict-relation} 所定义的群胚限制和
（预）等价关系限制，经引理 \ref{groupoids-lemma-groupoid-pre-equivalence} 与
\ref{groupoids-lemma-equivalence-groupoid} 的构造彼此一致。
\end{lemma}

\begin{proof}
这里所断言的是，引理 \ref{groupoids-lemma-restrict-groupoid} 中的
$R'$ 也等于
$$
R' = (U' \times_S U')\times_{U \times_S U} R
\longrightarrow
U' \times_S U'
$$
事实上，这也许是陈述该引理更清楚的方式。
\end{proof}

\begin{lemma}
\label{groupoids-lemma-restrict-stabilizer}
设 $S$ 为概形，$(U, R, s, t, c)$ 为 $S$ 上的群胚概形，
$g : U' \to U$ 为概形态射。令
$(U', R', s', t', c')$ 为 $(U, R, s, t, c)$ 经 $g$ 的限制。
设 $G$ 为 $(U, R, s, t, c)$ 的稳定子，$G'$ 为
$(U', R', s', t', c')$ 的稳定子。则 $G'$ 是 $G$ 沿 $g$ 的基变换；
换言之，存在典范识别 $G' = U' \times_{g, U} G$。
\end{lemma}

\begin{proof}
证明略。
\end{proof}






\section{不变子概形}
\label{groupoids-section-invariant}

\noindent
本节简要讨论不变子概形的概念。

\begin{definition}
\label{groupoids-definition-invariant-open}
设 $(U, R, s, t, c)$ 为基概形 $S$ 上的群胚概形。
\begin{enumerate}
\item 若 $t(s^{-1}(W)) \subset W$，则称子集
$W \subset U$ 在{\it 集合论意义下是 $R$-不变的}。
\item 若 $t(s^{-1}(W)) \subset W$，则称开子集
$W \subset U$ 是{\it $R$-不变的}。
\item 若 $t^{-1}(Z) = s^{-1}(Z)$，则称闭子概形
$Z \subset U$ 是{\it $R$-不变的}。这里采用概形论逆像；见
《概形》定义 \ref{schemes-definition-inverse-image-closed-subscheme}。
\item 若作为 $R$ 上的概形有
$T \times_{U, t} R = R \times_{s, U} T$，则称概形单态射
$T \to U$ 是{\it $R$-不变的}。
\end{enumerate}
\end{definition}

\noindent
对于子集与开子概形 $W \subset U$，$R$-不变性也等价于要求
作为 $R$ 的子集有 $s^{-1}(W) = t^{-1}(W)$。
若 $W \subset U$ 是 $R$-等变开子概形，则 $R$ 在 $W$ 上的限制就是
$R_W = s^{-1}(W) = t^{-1}(W)$。类似地，若
$Z \subset U$ 是 $R$-不变闭子概形，则 $R$ 在 $Z$ 上的限制就是
$R_Z = s^{-1}(Z) = t^{-1}(Z)$。

\begin{lemma}
\label{groupoids-lemma-constructing-invariant-opens}
设 $S$ 为概形，$(U, R, s, t, c)$ 为 $S$ 上的群胚概形。
\begin{enumerate}
\item 对任意子集 $W \subset U$，子集 $t(s^{-1}(W))$
在集合论意义下是 $R$-不变的。
\item 若 $s$ 与 $t$ 是开态射，则对每个开子集 $W \subset U$，
开集 $t(s^{-1}(W))$ 是 $R$-不变开子概形。
\item 若 $s$ 与 $t$ 开且拟紧，则 $U$ 有一个由
$R$-不变拟紧开子概形组成的开覆盖。
\end{enumerate}
\end{lemma}

\begin{proof}
第 (1) 点由引理 \ref{groupoids-lemma-pre-equivalence-equivalence-relation-points} 和
\ref{groupoids-lemma-groupoid-pre-equivalence} 得出：$t(s^{-1}(W))$
正是 $U$ 中与 $W$ 的某一点等价的点集。
接着假设 $s$ 和 $t$ 开，且 $W \subset U$ 开。
由于 $s$ 开，集合 $W' = t(s^{-1}(W))$ 是 $U$ 的开子集。
最后，假设 $s$, $t$ 都开且拟紧。若 $W \subset U$ 拟紧开，
则 $W' = t(s^{-1}(W))$ 也拟紧开，并由上面的讨论可知它不变。
让 $W$ 遍历 $U$ 的全体仿射开集，即得 (3)。
\end{proof}

\begin{lemma}
\label{groupoids-lemma-first-observation}
设 $S$ 为概形，$(U, R, s, t, c)$ 为 $S$ 上的群胚概形。
假设 $s$ 与 $t$ 拟紧且平坦，而 $U$ 拟分离。若
$W \subset U$ 拟紧开，则 $t(s^{-1}(W))$ 是
$U$ 中某个非空族拟紧开子集的交。
\end{lemma}

\begin{proof}
注意 $s^{-1}(W)$ 是 $R$ 中的拟紧开集。作为连续映射，
$t$ 把拟紧子集 $s^{-1}(W)$ 映到拟紧子集 $t(s^{-1}(W))$。
由于 $t$ 平坦且 $s^{-1}(W)$ 对泛化封闭，$t(s^{-1}(W))$
也对泛化封闭；见
（《概形态射》引理 \ref{morphisms-lemma-generalizations-lift-flat} 及
《拓扑学》引理 \ref{topology-lemma-lift-specializations-images}）。
选取包含 $t(s^{-1}(W))$ 的拟紧开集 $W' \subset U$。
由《概形的性质》引理 \ref{properties-lemma-quasi-compact-quasi-separated-spectral}，
$W'$ 是谱空间（这里用到了 $U$ 拟分离）。对
$t(s^{-1}(W)) \subset W'$ 应用《拓扑学》引理 \ref{topology-lemma-make-spectral-space}，即得结论。
\end{proof}

\begin{lemma}
\label{groupoids-lemma-second-observation}
假设与记号同引理 \ref{groupoids-lemma-first-observation}。
存在 $R$-不变开集 $V \subset U$ 以及拟紧开集 $W'$，
使 $W \subset V \subset W' \subset U$。
\end{lemma}

\begin{proof}
令 $E = t(s^{-1}(W))$。回忆 $E$ 在集合论意义下是 $R$-不变的
（引理 \ref{groupoids-lemma-constructing-invariant-opens}）。
由引理 \ref{groupoids-lemma-first-observation}，存在包含 $E$ 的拟紧开集
$W'$。令 $Z = U \setminus W'$，并考虑
$T = t(s^{-1}(Z))$。注意 $Z \subset T$，且
$E \cap T = \emptyset$，因为
$s^{-1}(E) = t^{-1}(E)$ 与 $s^{-1}(Z)$ 不交。
由于 $T$ 是闭子集 $s^{-1}(Z) \subset R$ 在拟紧态射
$t : R \to U$ 下的像，闭包 $\overline{T}$ 中任意点
$\xi$ 都是 $T$ 中某点的特化；见
《概形态射》引理 \ref{morphisms-lemma-reach-points-scheme-theoretic-image}
（并见《概形态射》引理 \ref{morphisms-lemma-quasi-compact-scheme-theoretic-image}，
它说明概形论像就是像的闭包）。
设 $\xi' \leadsto \xi$，其中 $\xi' \in T$。若 $r \in R$ 且
$s(r) = \xi$，则由 $s$ 平坦，可在 $R$ 中找到一个特化
$r' \leadsto r$，使 $s(r') = \xi'$
（《概形态射》引理 \ref{morphisms-lemma-generalizations-lift-flat}）。
于是 $t(r') \leadsto t(r)$。由于 $T$ 由
引理 \ref{groupoids-lemma-constructing-invariant-opens}
在集合论意义下不变，故 $t(r') \in T$。
因此 $\overline{T}$ 是集合论意义下 $R$-不变的闭子集，
而 $V = U \setminus \overline{T}$ 正是所求开集。
它包含于 $W'$，证明完成。
\end{proof}




\section{商层}
\label{groupoids-section-quotient-sheaves}

\noindent
设 $\tau \in \{Zariski, \etale, fppf, smooth, syntomic\}$，
$S$ 为概形，$j : R \to U \times_S U$ 为 $S$ 上的预关系。
假设 $U, R, S$ 是某个 $\tau$-位点 $\Sch_\tau$ 的对象
（见《概形上的拓扑》第 \ref{topologies-section-procedure} 节）。
于是可考虑函子
$$
h_U, h_R :
(\Sch/S)_\tau^{opp}
\longrightarrow
\textit{Sets}.
$$
它们都是层，见《下降》引理 \ref{descent-lemma-fpqc-universal-effective-epimorphisms}。
态射 $j$ 诱导映射 $j : h_R \to h_U \times h_U$。
对每个对象 $T \in \Ob((\Sch/S)_\tau)$，
可取由 $j(T) : R(T) \to U(T) \times U(T)$ 生成的等价关系
$\sim_T$，并取其商。这样得到预层
\begin{equation}
\label{groupoids-equation-quotient-presheaf}
(\Sch/S)_\tau^{opp}
\longrightarrow
\textit{Sets}, \quad
T \longmapsto U(T)/\sim_T
\end{equation}

\begin{definition}
\label{groupoids-definition-quotient-sheaf}
令 $\tau$, $S$ 及预关系 $j : R \to U \times_S U$ 如上。
在此情形下，与 $j$ 关联的{\it 商层 $U/R$}，是预层
(\ref{groupoids-equation-quotient-presheaf}) 在 $\tau$-拓扑下的层化。
若 $j : R \to U \times_S U$ 来自群概形 $G/S$
在 $U$ 上的作用，如引理 \ref{groupoids-lemma-groupoid-from-action}，
则有时把该商层记作 $U/G$。
\end{definition}

\noindent
这恰好意味着图
$$
\xymatrix{
h_R \ar@<1ex>[r] \ar@<-1ex>[r] &
h_U \ar[r] &
U/R
}
$$
是 $(\Sch/S)_\tau$ 上集合层范畴中的余等化子图。
利用 Yoneda 嵌入，可以把 $(\Sch/S)_\tau$
看成 $(\Sch/S)_\tau$ 上层范畴的全子范畴，从而把概形与可表函子相识别。
沿用这一记号滥用，通常把上图简写为
$$
\xymatrix{
R \ar@<1ex>[r]^s \ar@<-1ex>[r]_t &
U \ar[r] &
U/R
}
$$
在考虑群胚或等价关系的商层时，我们主要使用 fppf 拓扑。

\begin{definition}
\label{groupoids-definition-representable-quotient}
在定义 \ref{groupoids-definition-quotient-sheaf} 的情形下，
若层 $U/R$ 可表，则称预关系 $j$ 有{\it 可表商}。
若取 $j = (t, s)$ 时商 $U/R$ 可表，则称群胚
$(U, R, s, t, c)$ 有{\it 可表商}。
\end{definition}

\noindent
下述引理刻画表示该商的概形 $M$。例如，当
$\tau = fppf$、$U \to M$ 平坦、有限表现且满，并且
$R \cong U \times_M U$ 时，该引理适用。

\begin{lemma}
\label{groupoids-lemma-criterion-quotient-representable}
在定义 \ref{groupoids-definition-quotient-sheaf} 的情形下，
假设存在概形 $M$ 以及态射 $U \to M$，满足：
\begin{enumerate}
\item 态射 $U \to M$ 等化 $s, t$；
\item 态射 $U \to M$ 在 $\tau$-拓扑下诱导层的满射
$h_U \to h_M$；
\item 诱导映射 $(t, s) : R \to U \times_M U$
在 $\tau$-拓扑下诱导层的满射
$h_R \to h_{U \times_M U}$。
\end{enumerate}
此时 $M$ 表示商层 $U/R$。
\end{lemma}

\begin{proof}
条件 (1) 说明 $h_U \to h_M$ 通过 $U/R$ 分解；
条件 (2) 说明层映射 $U/R \to h_M$ 是满射；
条件 (3) 说明层映射 $U/R \to h_M$ 是单射。
故结论成立。
\end{proof}

\noindent
若不要求 $j$ 为预等价关系（例如仅为预关系），则下述引理不成立。

\begin{lemma}
\label{groupoids-lemma-quotient-pre-equivalence}
设 $\tau \in \{Zariski, \etale, fppf, smooth, syntomic\}$，
$S$ 为概形，$j : R \to U \times_S U$ 为 $S$ 上的预等价关系，
并假设 $U, R, S$ 是某个 $\tau$-位点 $\Sch_\tau$ 的对象。
对 $T \in \Ob((\Sch/S)_\tau)$ 及
$a, b \in U(T)$，下列条件等价：
\begin{enumerate}
\item $a$ 与 $b$ 映到 $(U/R)(T)$ 中的同一元素；
\item 存在 $T$ 的一个 $\tau$-覆盖
$\{f_i : T_i \to T\}$ 以及态射 $r_i : T_i \to R$，使
$a \circ f_i = s \circ r_i$ 且 $b \circ f_i = t \circ r_i$。
\end{enumerate}
换言之，此时 $\tau$-层的映射
$$
h_R \longrightarrow h_U \times_{U/R} h_U
$$
是满射。
\end{lemma}

\begin{proof}
证明略。提示：该论证之所以成立，是因为此时预层
(\ref{groupoids-equation-quotient-presheaf}) 实际由
$T \mapsto U(T)/j(R(T))$ 给出，而
$j(R(T)) \subset U(T) \times U(T)$ 是等价关系；见
定义 \ref{groupoids-definition-equivalence-relation}。
\end{proof}

\begin{lemma}
\label{groupoids-lemma-quotient-pre-equivalence-relation-restrict}
设 $\tau \in \{Zariski, \etale, fppf, smooth, syntomic\}$，
$S$ 为概形，$j : R \to U \times_S U$ 为 $S$ 上的预等价关系，
$g : U' \to U$ 为 $S$ 上的概形态射。令
$j' : R' \to U' \times_S U'$ 为 $j$ 在 $U'$ 上的限制。
假设 $U, U', R, S$ 是某个 $\tau$-位点 $\Sch_\tau$ 的对象。
则商层的映射
$$
U'/R' \longrightarrow U/R
$$
是单射。若 $g$ 定义层的满射 $h_{U'} \to h_U$，
其拓扑为 $\tau$（例如，若 $\{g : U' \to U\}$
是 $\tau$-覆盖），则 $U'/R' \to U/R$ 是同构。
\end{lemma}

\begin{proof}
假设截面 $\xi, \xi' \in (U'/R')(T)$ 映到 $U/R$ 的同一截面。
可以找到 $T$ 的一个 $\tau$-覆盖
$\mathcal{T} = \{T_i \to T\}$，使
$\xi|_{T_i}, \xi'|_{T_i}$ 分别由
$a_i, a_i' \in U'(T_i)$ 给出。由
引理 \ref{groupoids-lemma-quotient-pre-equivalence} 及位点公理，
加细 $\mathcal{T}$ 后可假设存在态射 $r_i : T_i \to R$，使
$g \circ a_i = s \circ r_i$、$g \circ a_i' = t \circ r_i$。
按照构造，
$R' = R \times_{U \times_S U} (U' \times_S U')$，
故 $(r_i, (a_i, a_i')) \in R'(T_i)$。这说明 $a_i$ 与
$a_i'$ 在 $T_i$ 上定义 $U'/R'$ 的同一截面。
由层条件可知 $\xi = \xi'$。

\medskip\noindent
若 $h_{U'} \to h_U$ 是层的满射，则
$U'/R' \to U/R$ 当然也是满射。若
$\{g : U' \to U\}$ 是 $\tau$-覆盖，则层映射
$h_{U'} \to h_U$ 是满射，见《位点与层》引理 \ref{sites-lemma-covering-surjective-after-sheafification}。
因此在此情形下 $U'/R' \to U/R$ 也是满射。
\end{proof}

\begin{lemma}
\label{groupoids-lemma-quotient-groupoid-restrict}
设 $\tau \in \{Zariski, \etale, fppf, smooth, syntomic\}$，
$S$ 为概形，$(U, R, s, t, c)$ 为 $S$ 上的群胚概形，
$g : U' \to U$ 为 $S$ 上的概形态射。令
$(U', R', s', t', c')$ 为 $(U, R, s, t, c)$ 在 $U'$ 上的限制。
假设 $U, U', R, S$ 是某个 $\tau$-位点 $\Sch_\tau$ 的对象。
则商层的映射
$$
U'/R' \longrightarrow U/R
$$
是单射。若复合
$$
\xymatrix{
U' \times_{g, U, t} R \ar[r]_-{\text{pr}_1} \ar@/^3ex/[rr]^h
& R \ar[r]_s & U
}
$$
在 $\tau$-拓扑下定义层的满射，则该映射是双射。
例如，下列任一条件均足够：
$\{h : U' \times_{g, U, t} R \to U\}$ 是 $\tau$-覆盖；
$U' \to U$ 在 $\tau$-拓扑下定义层的满射；或
$\{g : U' \to U\}$ 是 $\tau$-拓扑下的覆盖。
\end{lemma}

\begin{proof}
由引理 \ref{groupoids-lemma-groupoid-pre-equivalence} 和
\ref{groupoids-lemma-quotient-pre-equivalence-relation-restrict} 即得单射性。
为证明满射性（层满射的刻画见
《位点与层》第 \ref{sites-section-sheaves-injective} 节），作如下论证。
假设 $T$ 是概形，$\sigma \in U/R(T)$。存在覆盖
$\{T_i \to T\}$，使 $\sigma|_{T_i}$ 是某个元素
$f_i \in U(T_i)$ 的像。因此可假设 $\sigma$ 是
$f \in U(T)$ 的像。由 $h$ 是层满射这一假设，可以找到
$\tau$-覆盖 $\{\varphi_i : T_i \to T\}$ 以及态射
$f_i : T_i \to U' \times_{g, U, t} R$，使
$f \circ \varphi_i = h \circ f_i$。记
$f'_i = \text{pr}_0 \circ f_i : T_i \to U'$。
于是 $f'_i \in U'(T_i)$ 映到
$g \circ f'_i \in U(T_i)$，并且
$g \circ f'_i \sim_{T_i} h \circ f_i = f \circ \varphi_i$；
记号同 (\ref{groupoids-equation-quotient-presheaf})。具体地，
给出该关系的 $R(T_i)$ 中元素是 $\text{pr}_1 \circ f_i$。
这说明 $\sigma$ 在 $T_i$ 上的限制属于
$U'/R'(T_i) \to U/R(T_i)$ 的像，如所需。

\medskip\noindent
若 $\{h\}$ 是 $\tau$-覆盖，则它诱导层的满射，见
《位点与层》引理 \ref{sites-lemma-covering-surjective-after-sheafification}。
若 $U' \to U$ 满，则 $h$ 也满，因为 $s$ 有截面
（即群胚概形的单位元 $e$）。
\end{proof}

\begin{lemma}
\label{groupoids-lemma-criterion-fibre-product}
设 $S$ 为概形，$f : (U, R, j) \to (U', R', j')$
为 $S$ 上等价关系之间的态射。假设
$$
\xymatrix{
R \ar[d]_s \ar[r]_f & R' \ar[d]^{s'} \\
U \ar[r]^f & U'
}
$$
是笛卡尔的。对任意
$\tau \in  \{Zariski, \etale, fppf, smooth, syntomic\}$，图
$$
\xymatrix{
U \ar[d] \ar[r] & U/R \ar[d]^f \\
U' \ar[r] & U'/R'
}
$$
是 $\tau$-层的纤维积方块。
\end{lemma}

\begin{proof}
由引理 \ref{groupoids-lemma-quotient-pre-equivalence}，商层有一个简单描述；
下文将不再说明而直接使用它。先证明
$$
U \longrightarrow U' \times_{U'/R'} U/R
$$
是单射。假设 $a, b \in U(T)$ 映到右侧同一元素，则
$f(a) = f(b)$。把 $T$ 替换为某个 $\tau$-覆盖的各成员后，
可假设存在 $r \in R(T)$，使 $a = s(r)$ 且 $b = t(r)$。
于是 $r' = f(r)$ 是 $R'$ 的一个 $T$-值点，满足
$s'(r') = t'(r')$。因此 $r' = e'(f(a))$（其中 $e'$
是与 $j'$ 关联的群胚概形的单位，见
引理 \ref{groupoids-lemma-equivalence-groupoid}）。
由于引理中的第一个图是笛卡尔的，这蕴含 $r$ 必须等于
$e(a)$，故 $a = b$。

\medskip\noindent
最后证明所示箭头满。设 $T$ 为 $S$ 上的概形，
$(a', \overline{b})$ 为层
$U' \times_{U'/R'} U/R$ 在 $T$ 上的截面。
把 $T$ 替换为某个 $\tau$-覆盖的各成员后，可假设
$\overline{b}$ 是元素 $b \in U(T)$ 的类。
再次把 $T$ 替换为某个 $\tau$-覆盖的各成员后，可假设存在
$r' \in R'(T)$，使 $a' = t(r')$ 且 $s'(r') = f(b)$。
由于引理中的第一个图是笛卡尔的，可找到 $r \in R(T)$，
使 $s(r) = b$ 且 $f(r) = r'$。显然，
$a = t(r) \in U(T)$ 是映到 $(a', \overline{b})$ 的截面。
\end{proof}








\section{用群胚表述下降}
\label{groupoids-section-groupoids-descent}

\noindent
笛卡尔态射定义如下。

\begin{definition}
\label{groupoids-definition-cartesian-morphism}
设 $S$ 为概形，$f : (U', R', s', t', c') \to (U, R, s, t, c)$
为 $S$ 上群胚概形的态射。若图
$$
\xymatrix{
R' \ar[r]_f \ar[d]_{s'} & R \ar[d]^s \\
U' \ar[r]^f & U
}
$$
是概形范畴中的纤维积方块，则称 $f$ 是{\it 笛卡尔的}，
或称{\it $(U', R', s', t', c')$ 在 $(U, R, s, t, c)$ 上是笛卡尔的}。
所谓{\it 在 $(U, R, s, t, c)$ 上笛卡尔的群胚概形态射}，
是指与指向 $(U, R, s, t, c)$ 的结构态射相容的群胚概形态射。
\end{definition}

\noindent
笛卡尔态射与下降数据相关。先证明一个一般引理，
描述固定群胚概形上的笛卡尔群胚概形范畴。

\begin{lemma}
\label{groupoids-lemma-characterize-cartesian-schemes}
设 $S$ 为概形，$(U, R, s, t, c)$ 为 $S$ 上的群胚概形。
在 $(U, R, s, t, c)$ 上笛卡尔的群胚概形范畴等价于
如下偶对 $(V, \varphi)$ 的范畴：$V$ 是 $U$ 上的概形，而
$$
\varphi :
V \times_{U, t} R
\longrightarrow
R \times_{s, U} V
$$
是 $R$ 上的同构，满足 $e^*\varphi = \text{id}_V$，并且作为
$R \times_{s, U, t} R$ 上的概形态射有
$$
c^*\varphi = \text{pr}_1^*\varphi \circ \text{pr}_0^*\varphi
$$
\end{lemma}

\begin{proof}
引理中的拉回记号表示基变换。所示公式有意义，因为作为
$R \times_{s, U, t} R$ 上的概形有
$$
(R \times_{s, U, t} R) \times_{\text{pr}_1, R, \text{pr}_1} (V \times_{U, t} R)
=
(R \times_{s, U, t} R) \times_{\text{pr}_0, R, \text{pr}_0} (R \times_{s, U} V)
$$

\medskip\noindent
给定 $(V, \varphi)$，令 $U' = V$、$R' = V \times_{U, t} R$。
令 $t' : R' \to U'$ 为投影 $V \times_{U, t} R \to V$；
令 $s'$ 为 $\varphi$ 后接投影 $R \times_{s, U} V \to V$；
令 $c'$ 为复合
\begin{align*}
R' \times_{s', U', t'} R'
& \xrightarrow{\varphi, 1}
(R \times_{s, U} V) \times_V (V \times_{U, t} R) \\
& \xrightarrow{}
R \times_{s, U} V \times_{U, t} R \\
& \xrightarrow{\varphi^{-1}, 1}
V \times_{U, t} (R \times_{s, U, t} R) \\
& \xrightarrow{1, c}
V \times_{U, t} R = R'
\end{align*}
略去的计算表明，这给出 $(U, R, s, t, c)$ 上的群胚概形；
显然，该群胚概形在 $(U, R, s, t, c)$ 上是笛卡尔的。

\medskip\noindent
反之，给定笛卡尔态射
$f : (U', R', s', t', c') \to (U, R, s, t, c)$，
则态射
$$
U' \times_{U, t} R \xleftarrow{t', f} R' \xrightarrow{f, s'} R \times_{s, U} U'
$$
都是同构。可以令 $V = U'$，并令 $\varphi$ 为复合
$(f, s') \circ (t', f)^{-1}$。我们略去验证 $\varphi$
满足引理中的条件，也略去验证这两个构造互逆。
\end{proof}

\noindent
设 $S$ 为概形，$f : X \to Y$ 为 $S$ 上的概形态射。
于是得到 $S$ 上的群胚概形
$(X, X \times_Y X, \text{pr}_1, \text{pr}_0, c)$。
事实上，$j : X \times_Y X \to X \times_S X$ 是等价关系，
可取其关联群胚；见引理 \ref{groupoids-lemma-equivalence-groupoid}。

\begin{lemma}
\label{groupoids-lemma-cartesian-equivalent-descent-datum}
设 $S$ 为概形，$f : X \to Y$ 为 $S$ 上的概形态射。
引理 \ref{groupoids-lemma-characterize-cartesian-schemes} 的构造给出等价
$$
\begin{matrix}
\text{群胚概形范畴} \\
\text{在下列对象上 Cartesian： } (X, X \times_Y X, \ldots)
\end{matrix}
\longrightarrow
\begin{matrix}
\text{ 下降数据范畴} \\
\text{ 相对于 } X/Y
\end{matrix}
$$
\end{lemma}

\begin{proof}
这直接源于引理 \ref{groupoids-lemma-characterize-cartesian-schemes}
以及《下降》定义 \ref{descent-definition-descent-datum}
中概形上下降数据的定义。
\end{proof}







\section{分离性条件}
\label{groupoids-section-separation}

\noindent
给定 $S$ 上的群胚 $(U, R, s, t, c)$ 时，这实际上是对态射
$j : R \to U \times_S U$ 的条件。与前一节一样，先写出相应的图。

\begin{lemma}
\label{groupoids-lemma-diagram-diagonal}
设 $S$ 为概形，$(U, R, s, t, c)$ 为 $S$ 上的群胚，
$G \to U$ 为稳定子群概形。交换图
$$
\xymatrix{
R \ar[d]^{\Delta_{R/U \times_S U}} \ar[rrr]_{f \mapsto (f, s(f))} & & &
R \times_{s, U} U \ar[d] \ar[r] & U \ar[d] \\
R \times_{(U \times_S U)} R \ar[rrr]^{(f, g) \mapsto (f, f^{-1} \circ g)} & & &
R \times_{s, U} G \ar[r] & G
}
$$
中的两个左侧水平箭头都是同构，右侧方块是纤维积方块。
\end{lemma}

\begin{proof}
证明略。这是关于定义以及代数几何函子观点的练习。
\end{proof}

\begin{lemma}
\label{groupoids-lemma-diagonal}
设 $S$ 为概形，$(U, R, s, t, c)$ 为 $S$ 上的群胚，
$G \to U$ 为稳定子群概形。
\begin{enumerate}
\item 下列条件等价：
\begin{enumerate}
\item $j : R \to U \times_S U$ 分离；
\item $G \to U$ 分离；
\item $e : U \to G$ 是闭浸入。
\end{enumerate}
\item 下列条件等价：
\begin{enumerate}
\item $j : R \to U \times_S U$ 拟分离；
\item $G \to U$ 拟分离；
\item $e : U \to G$ 拟紧。
\end{enumerate}
\end{enumerate}
\end{lemma}

\begin{proof}
群概形 $G \to U$ 是 $R \to U \times_S U$
沿对角态射 $U \to U \times_S U$ 的基变换，见
引理 \ref{groupoids-lemma-groupoid-stabilizer}。因此，若
$j$ 分离（resp.\ 拟分离），则
$G \to U$ 分离（resp.\ 拟分离）
（见《概形》引理 \ref{schemes-lemma-separated-permanence}）。
故在 (1) 与 (2) 中均有 (a) $\Rightarrow$ (b)。

\medskip\noindent
若 $G \to U$ 分离（resp.\ 拟分离），则态射
$U \to G$ 作为结构态射 $G \to U$ 的截面，是闭浸入
（resp.\ 拟紧）；见《概形》引理 \ref{schemes-lemma-section-immersion}。
故在 (1) 与 (2) 中均有 (b) $\Rightarrow$ (a)。

\medskip\noindent
由引理 \ref{groupoids-lemma-diagram-diagonal} 的结果
（以及《概形》引理 \ref{schemes-lemma-base-change-immersion}
和 \ref{schemes-lemma-quasi-compact-preserved-base-change}），
若 $e$ 是闭浸入（resp.\ 拟紧），则
$\Delta_{R/U \times_S U}$ 是闭浸入（resp.\ 拟紧）。
故在 (1) 与 (2) 中均有 (c) $\Rightarrow$ (a)。
\end{proof}









\section{有限平坦群胚：仿射情形}
\label{groupoids-section-finite-flat}

\noindent
设 $S$ 为概形，$(U, R, s, t, c)$ 为 $S$ 上的群胚概形。
假设 $U = \Spec(A)$ 与 $R = \Spec(B)$ 仿射。此时得到两个环同态
$s^\sharp, t^\sharp : A \longrightarrow B$。
令 $C$ 为 $s^\sharp$ 与 $t^\sharp$ 的等化子，即
\begin{equation}
\label{groupoids-equation-invariants}
C = \{a \in A \mid t^\sharp(a) = s^\sharp(a) \}.
\end{equation}
有时称它为{\it $U$ 上 $R$-不变函数环}。
$M = \Spec(C)$ 有什么性质？首先，图
$$
\xymatrix{
R \ar[r]_s \ar[d]_t & U \ar[d] \\
U \ar[r] & M
}
$$
交换；换言之，态射 $U \to M$ 等化 $s, t$。
此外，若 $T$ 为任意仿射概形，且态射 $U \to T$ 等化 $s, t$，
则 $U \to T$ 通过 $U \to M$ 分解。也就是说，
$U \to M$ 是仿射概形范畴中的余等化子。

\medskip\noindent
我们希望找出保证态射 $U \to M$ 在概形范畴中确为“商”的条件。
这一问题将在别处详论（此处插入未来的引用）；这里仅讨论若干特殊情形，
主要关注 $s, t$ 有限局部自由的情形。

\begin{example}
\label{groupoids-example-quotient-projective-line}
设 $k$ 为域，$U = \text{GL}_{2, k}$，而
$B \subset \text{GL}_2$ 是上三角矩阵的闭子群概形。
则商层 $\text{GL}_{2, k}/B$（在 Zariski、\'etale 或 fppf 拓扑下，
见定义 \ref{groupoids-definition-quotient-sheaf}）由射影直线表示：
$\mathbf{P}^1 = \text{GL}_{2, k}/B$。（细节略。）
另一方面，此情形下的不变函数环仅为 $k$。注意，此时态射
$s, t : R = \text{GL}_{2, k} \times_k B \to \text{GL}_{2, k} = U$
是相对维数为 $3$ 的光滑态射。
\end{example}

\noindent
回忆在《习题》习题 \ref{exercises-exercise-trace-det} 和
\ref{exercises-exercise-trace-det-rings} 中，我们已经定义了
有限局部自由模及有限局部自由环扩张的行列式和范数。
若 $\varphi : A \to B$ 是有限局部自由环同态，则以
$\text{Norm}_\varphi(b) \in A$ 表示 $b \in B$ 的范数。
对于有限局部自由概形态射，范数构造见
《除子》引理 \ref{divisors-lemma-finite-locally-free-has-norm}。

\begin{lemma}
\label{groupoids-lemma-determinant-trick}
设 $S$ 为概形，$(U, R, s, t, c)$ 为 $S$ 上的群胚概形。
假设 $U = \Spec(A)$ 与 $R = \Spec(B)$ 仿射，并且
$s, t : R \to U$ 有限局部自由。令 $C$ 如
(\ref{groupoids-equation-invariants})，并取 $f \in A$。则
$\text{Norm}_{s^\sharp}(t^\sharp(f)) \in C$。
\end{lemma}

\begin{proof}
考虑引理 \ref{groupoids-lemma-diagram} 的交换图
$$
\xymatrix{
& U & \\
R \ar[d]_s \ar[ru]^t &
R \times_{s, U, t} R
\ar[l]^-{\text{pr}_0} \ar[d]^{\text{pr}_1} \ar[r]_-c &
R \ar[d]^s \ar[lu]_t \\
U & R \ar[l]_t \ar[r]^s & U
}
$$
把 $f \in \Gamma(U, \mathcal{O}_U)$ 看作函数。
图上半部分的交换性说明，作为
$\Gamma(R \times_{S, U, t} R, \mathcal{O})$ 中的元素，有
$\text{pr}_0^\sharp(t^\sharp(f)) = c^\sharp(t^\sharp(f))$。
观察右下笛卡尔方块，范数构造与基变换的相容性给出
$s^\sharp(\text{Norm}_{s^\sharp}(t^\sharp(f))) =
\text{Norm}_{\text{pr}_1^\sharp}(c^\sharp(t^\sharp(f)))$。
类似地得到
$t^\sharp(\text{Norm}_{s^\sharp}(t^\sharp(f))) =
\text{Norm}_{\text{pr}_1^\sharp}(\text{pr}_0^\sharp(t^\sharp(f)))$。
因此由本证明中的第一个等式，
$s^\sharp(\text{Norm}_{s^\sharp}(t^\sharp(f))) =
t^\sharp(\text{Norm}_{s^\sharp}(t^\sharp(f)))$，如所需。
\end{proof}

\begin{lemma}
\label{groupoids-lemma-finite-locally-free-disjoint-free}
设 $S$ 为概形，$(U, R, s, t, c)$ 为 $S$ 上的群胚概形。
假设 $s, t : R \to U$ 有限局部自由。则
$$
U = \coprod\nolimits_{r \geq 1} U_r
$$
是 $R$-不变开集的不交并，并且 $R$ 在 $U_r$ 上的限制
$R_r$ 满足 $s, t : R_r \to U_r$ 有限局部自由、秩为 $r$。
\end{lemma}

\begin{proof}
由《概形态射》引理 \ref{morphisms-lemma-finite-locally-free}，
存在分解 $U = \coprod\nolimits_{r \geq 0} U_r$，使
$s : s^{-1}(U_r) \to U_r$ 有限局部自由、秩为 $r$。
由于 $s$ 满，$U_0 = \emptyset$。注意，
$u \in U_r \Leftrightarrow$ 当且仅当概形论纤维
$s^{-1}(u)$ 在 $\kappa(u)$ 上次数为 $r$。若
$z \in R$ 满足 $s(z) = u$、$t(z) = u'$，则采用
引理 \ref{groupoids-lemma-diagram} 的记号，
$$
\text{pr}_1^{-1}(z) \to \Spec(\kappa(z))
$$
由所引引理知，它同时是
$s^{-1}(u) \to \Spec(\kappa(u))$ 与
$s^{-1}(u') \to \Spec(\kappa(u'))$ 的基变换。因此
$u \in U_r \Leftrightarrow u' \in U_r$；换言之，诸开子集
$U_r$ 都是 $R$-不变的。特别地，$R$ 在 $U_r$ 上的限制就是
$s^{-1}(U_r)$，且 $s : R_r \to U_r$ 有限局部自由、秩为 $r$。
由 $R_r$ 的逆，$t : R_r \to U_r$ 同构于 $s$，故其秩也为 $r$。
\end{proof}

\begin{lemma}
\label{groupoids-lemma-integral-over-invariants}
设 $S$ 为概形，$(U, R, s, t, c)$ 为 $S$ 上的群胚概形。
假设 $U = \Spec(A)$ 与 $R = \Spec(B)$ 仿射，并且
$s, t : R \to U$ 有限局部自由。令 $C \subset A$ 如
(\ref{groupoids-equation-invariants})。则 $A$ 在 $C$ 上整。
\end{lemma}

\begin{proof}
首先，由引理 \ref{groupoids-lemma-finite-locally-free-disjoint-free}，
$(U, R, s, t, c)$ 是群胚概形 $(U_r, R_r, s, t, c)$ 的不交并，
其中每个 $s, t : R_r \to U_r$ 的秩恒为 $r$。
由于 $U$ 拟紧，除有限多个 $r$ 外均有 $U_r = \emptyset$。
只需对每个 $(U_r, R_r, s, t, c)$ 证明该引理，故可假设
$s, t$ 有限局部自由、秩为 $r$。

\medskip\noindent
假设 $s, t$ 有限局部自由、秩为 $r$。取 $f \in A$，
考虑元素 $x - f \in A[x]$，其中把 $x$ 看作
$\mathbf{A}^1$ 上的坐标。由于
$$
(U \times \mathbf{A}^1, R \times \mathbf{A}^1,
s \times \text{id}_{\mathbf{A}^1},
t \times \text{id}_{\mathbf{A}^1},
c \times \text{id}_{\mathbf{A}^1})
$$
也是源和靶均有限的群胚概形，可对它应用
引理 \ref{groupoids-lemma-determinant-trick}，得到
$P(x) = \text{Norm}_{s^\sharp}(t^\sharp(x - f))$ 属于
$C[x]$。由于 $s^\sharp : A \to B$ 有限局部自由、秩为 $r$，
$P$ 是次数为 $r$ 的首一多项式。此外，由 Cayley--Hamilton 定理
有 $P(f) = 0$（《交换代数》引理 \ref{algebra-lemma-charpoly}）。
更准确地，Cayley--Hamilton 定理说明多项式
$s^\sharp(P) \in B[x]$ 零化元素 $t^\sharp(f)$。
由于 $P$ 的系数在 $C$ 中，得到 $t^\sharp(P(f)) = 0$ 于 $B$ 中。
又因 $t$ 忠实平坦，得到 $P(f) = 0$。
\end{proof}

\begin{lemma}
\label{groupoids-lemma-invariants-base-change}
设 $S$ 为概形，$(U, R, s, t, c)$ 为 $S$ 上的群胚概形。
假设 $U = \Spec(A)$ 与 $R = \Spec(B)$ 仿射，
$s, t : R \to U$ 有限局部自由。令 $C \subset A$ 如
(\ref{groupoids-equation-invariants})。给定环同态 $C \to C'$，并令
$U' = \Spec(A \otimes_C C')$、
$R' = \Spec(B \otimes_C C')$。则：
\begin{enumerate}
\item 映射 $s, t, c$ 诱导映射 $s', t', c'$，使
$(U', R', s', t', c')$ 成为群胚概形。令 $C^1 \subset A'$
为 $U'$ 上的 $R'$-不变函数。
\item 典范映射 $\varphi : C' \to C^1$ 满足：
\begin{enumerate}
\item 对每个 $f \in C^1$，存在 $n > 0$ 及多项式
$P \in C'[x]$，使其在 $C^1[x]$ 中的像为 $(x - f)^n$；
\item 对每个 $f \in \Ker(\varphi)$，存在 $n > 0$ 使
$f^n = 0$。
\end{enumerate}
\item 若 $C \to C'$ 平坦，则 $\varphi$ 是同构。
\end{enumerate}
\end{lemma}

\begin{proof}
第 (1) 点的证明略。记 $A' = A \otimes_C C'$、
$B' = B \otimes_C C'$。则
$$
C^1
= \{a \in A' \mid (t')^\sharp(a) = (s')^\sharp(a) \}
= \{a \in A \otimes_C C' \mid t^\sharp \otimes 1(a) = s^\sharp \otimes 1(a) \}.
$$
换言之，$C^1$ 是差映射
$(t^\sharp - s^\sharp) \otimes 1$ 的核，而后者正是
$C$-线性映射 $t^\sharp - s^\sharp : A \to B$
沿 $C \to C'$ 的基变换。因此 (3) 成立。

\medskip\noindent
证明 (2)(b)。由于 $C \to A$ 整
（引理 \ref{groupoids-lemma-integral-over-invariants}）且单，
$\Spec(A) \to \Spec(C)$ 满，见
《交换代数》引理 \ref{algebra-lemma-integral-overring-surjective}。
因此其满态射基变换 $\Spec(A') \to \Spec(C')$ 也满
（《概形态射》引理 \ref{morphisms-lemma-base-change-surjective}）。
于是 $\Spec(C^1) \to \Spec(C')$ 也满。这例如由
《交换代数》引理 \ref{algebra-lemma-image-dense-generic-points}
蕴含 (2)(b)。

\medskip\noindent
证明 (2)(a)。由引理 \ref{groupoids-lemma-finite-locally-free-disjoint-free}，
群胚概形 $(U, R, s, t, c)$ 分解为有限多个群胚概形
$(U_r, R_r, s, t, c)$ 的不交并，其中
$s, t : R_r \to U_r$ 有限局部自由、秩为 $r$。
沿 $U' = \Spec(C') \to U$ 拉回，得到 $U'$ 及
$U^1 = \Spec(C^1)$ 的类似分解。下一段将证明，对相应环系
$A_r, B_r, C_r, C'_r, C^1_r$，取 $n = r$ 时 (2)(a) 成立。
给定 $f \in C^1$，令 $P_r \in C_r[x]$ 为如下多项式：
其在 $C^1_r[x]$ 中的像是 $(x - f)^r$ 的像。
选择充分可除的整数 $n$，可得多项式 $P \in C'[x]$，
其在 $C^1[x]$ 中的像为 $(x - f)^n$；具体地，取
$P$ 为 $C'[x]$ 中唯一满足其在 $C'_r[x]$ 中的像为
$P_r^{n/r}$ 的元素。

\medskip\noindent
本段在环同态 $s^\sharp, t^\sharp : A \to B$
有限局部自由且秩恒为 $r$ 的情形下证明 (2)(a)。
取 $f \in C^1 \subset A' = A \otimes_C C'$。
选择平坦 $C$-代数 $D$ 以及满射 $D \to C'$，并选择
$f$ 的提升 $g \in A \otimes_C D$。考虑多项式
$$
P = \text{Norm}_{s^\sharp \otimes 1}(t^\sharp \otimes 1(x - g))
$$
于 $(A \otimes_C D)[x]$ 中。由
引理 \ref{groupoids-lemma-determinant-trick} 及本引理的 (3)，
$P$ 的系数属于 $D$（比较引理 \ref{groupoids-lemma-integral-over-invariants} 的证明）。
另一方面，$P$ 在 $(A \otimes_C C')[x]$ 中的像是
$(x - f)^r$，因为
$t^\sharp \otimes 1(x - f) = s^\sharp(x - f)$，且
$s^\sharp$ 有限局部自由、秩为 $r$。
取陈述 (2)(a) 中的 $P$ 为上述 $P$ 在映射
$D[x] \to C'[x]$ 下的像，即得所需。
\end{proof}

\begin{lemma}
\label{groupoids-lemma-points}
设 $S$ 为概形，$(U, R, s, t, c)$ 为 $S$ 上的群胚概形。
假设 $U = \Spec(A)$ 与 $R = \Spec(B)$ 仿射，
$s, t : R \to U$ 有限局部自由，并令 $C \subset A$ 如
(\ref{groupoids-equation-invariants})。则 $U \to M = \Spec(C)$ 具有下列性质：
\begin{enumerate}
\item 点集映射 $|U| \to |M|$ 是满射，并且
$u_0, u_1 \in |U|$ 映到同一点当且仅当存在
$r \in |R|$，满足 $t(r) = u_0$ 且 $s(r) = u_1$；即
$$
|M| = |U|/|R|
$$
\item 对任意代数闭域 $k$，有
$$
M(k) = U(k)/R(k)
$$
\end{enumerate}
\end{lemma}

\begin{proof}
由于 $C \to A$ 整（引理 \ref{groupoids-lemma-integral-over-invariants}）且单，
$\Spec(A) \to \Spec(C)$ 满，见
《交换代数》引理 \ref{algebra-lemma-integral-overring-surjective}。
因此 $|U| \to |M|$ 满。

\medskip\noindent
设 $k$ 为代数闭域，$C \to k$ 为环同态。
满态射在基变换下保持
（《概形态射》引理 \ref{morphisms-lemma-base-change-surjective}），
故 $A \otimes_C k$ 非零。此时 $k \subset A \otimes_C k$
是非零整扩张。因此 $A \otimes_C k$ 的任意剩余域都是
$k$ 的代数扩张，从而等于 $k$。于是 $U(k) \to M(k)$ 满。

\medskip\noindent
设 $a_0, a_1 : A \to k$ 为两个环同态。若存在环同态
$b : B \to k$，使 $a_0 = b \circ t^\sharp$ 且
$a_1 = b \circ s^\sharp$，则按定义
$a_0|_C = a_1|_C$。因此映射 $U(k) \to M(k)$
等化两个映射 $R(k) \to U(k)$。反之，假设
$a_0|_C = a_1|_C$，并将此代数同态记作 $c : C \to k$。
考虑图
$$
\xymatrix{
& &
B \ar@{-->}[lld] \\
k & &
A
\ar@<0.5ex>[ll]^{a_0}
\ar@<-0.5ex>[ll]_{a_1}
\ar@<1ex>[u]
\ar@<-1ex>[u] \\
& &
C \ar[u] \ar[llu]^c
}
$$
若能构造使该图交换的虚线箭头，则引理第 (2) 点的证明完成。
由于 $s : A \to B$ 有限，存在有限多个环同态
$b_1, \ldots, b_n : B \to k$，使
$b_i \circ s^\sharp = a_1$。若虚线箭头不存在，则
$a'_i = b_i \circ t^\sharp$（$i = 1, \ldots, n$）
均不等于 $a_0$。于是 $A \otimes_C k$ 的极大理想
$$
\mathfrak m'_i = \Ker(a_i' \otimes 1 : A \otimes_C k \to k)
$$
均不同于
$\mathfrak m = \Ker(a_0 \otimes 1 : A \otimes_C k \to k)$。
由《交换代数》引理 \ref{algebra-lemma-silly}，可得元素
$f \in A \otimes_C k$，满足 $f \in \mathfrak m$，但对
$i = 1, \ldots, n$ 有 $f \not \in \mathfrak m_i'$。
考虑范数
$$
g = \text{Norm}_{s^\sharp \otimes 1}(t^\sharp \otimes 1(f))
\in
A \otimes_C k
$$
由引理 \ref{groupoids-lemma-determinant-trick}，它属于经映射
$c : C \to k$ 基变换所得群胚的不变函数
$C^1 \subset A \otimes_C k$。一方面，$a_1(g) \in k^*$，
因为 $t^\sharp(f)$ 在 $a_1$ 上方所有点
（对应于 $b_1, \ldots, b_n$）处的值都可逆
（此处插入关于行列式性质的未来引用）。
另一方面，由于 $f \in \mathfrak m$，$f$ 不是单位，
故 $t^\sharp(f)$ 也不是单位（因为
$t^\sharp \otimes 1$ 忠实平坦），所以它的范数不是单位
（此处插入关于行列式性质的未来引用）。
于是 $C^1$ 含有一个既非幂零又非单位的元素。
下面说明这导致矛盾。把引理 \ref{groupoids-lemma-invariants-base-change}
应用于映射 $c : C \to C' = k$，可知 $k$ 到不变函数
$C^1$ 的映射是单射；而且对任意元素 $x \in C^1$，
存在整数 $n > 0$ 使 $x^n \in k$。因此 $C^1$ 的每个元素
要么是单位，要么幂零。

\medskip\noindent
还需完成 (1) 的证明。已知 $|U| \to |M|$ 满，且显然
$|U| \to |M|$ 是 $|R|$-不变的。最后，假设
$u_0, u_1 \in U$ 映到同一点 $m \in M$。
诱导的域扩张 $\kappa(u_0)/\kappa(m)$ 和
$\kappa(u_1)/\kappa(m)$ 都是代数扩张（如上所用，因为
$A$ 在 $C$ 上整）。若 $k$ 是 $\kappa(m)$ 的代数闭包，
则可找到 $\kappa(m)$-嵌入
$\overline{u}_0 : \kappa(u_0) \to k$ 和
$\overline{u}_1 : \kappa(u_1) \to k$。它们确定映到
$M(k)$ 中同一点的 $k$-值点
$\overline{u}_0, \overline{u}_1 \in U(k)$。
由 (2)，存在点 $\overline{r} \in R(k)$，满足
$s(\overline{r}) = \overline{u}_0$ 且
$t(\overline{r}) = \overline{u}_1$。
$\overline{r}$ 的像 $r \in R$ 满足
$s(r) = u_0$ 且 $t(r) = u_1$，如所需。
\end{proof}

\begin{lemma}
\label{groupoids-lemma-etale}
设 $S$ 为概形，$f : (U', R', s', t') \to (U, R, s, t, c)$
为 $S$ 上群胚概形的态射。假设：
\begin{enumerate}
\item $U$, $R$, $U'$, $R'$ 都仿射；
\item $s, t, s', t'$ 有限局部自由；
\item 图
$$
\xymatrix{
R' \ar[d]_{s'} \ar[r]_f & R \ar[d]^s \\
U' \ar[r]^f & U
}
\quad
\quad
\xymatrix{
R' \ar[d]_{t'} \ar[r]_f & R \ar[d]^t \\
U' \ar[r]^f & U
}
\quad
\quad
\xymatrix{
G' \ar[d] \ar[r]_f & G \ar[d] \\
U' \ar[r]^f & U
}
$$
均为笛卡尔的，其中 $G$ 与 $G'$ 为稳定子群概形；
\item $f : U' \to U$ 是 \'etale 态射。
\end{enumerate}
则从 $U$ 上 $R$-不变函数到 $U'$ 上 $R'$-不变函数的映射
$C \to C'$ 是 \'etale 的，并且
$U' = \Spec(C') \times_{\Spec(C)} U$。
\end{lemma}

\begin{proof}
令 $M = \Spec(C)$、$M' = \Spec(C')$。写
$U = \Spec(A)$、$U' = \Spec(A')$、$R = \Spec(B)$ 及
$R' = \Spec(B')$。下文将直接使用引理 \ref{groupoids-lemma-integral-over-invariants}、
\ref{groupoids-lemma-invariants-base-change} 和
\ref{groupoids-lemma-points} 的结果。

\medskip\noindent
假设 $C$ 是严格 Hensel 局部环。设 $p \in M$ 为闭点，
$p' \in M'$ 映到 $p$。断言：此时存在
$(U, R, s, t, c)$ 上的不交并分解
$(U', R', s', t', c') = (U, R, s, t, c) \amalg (U'', R'', s'', t'', c'')$，
使得对 $M$ 上相应的不交并分解
$M' = M \amalg M''$，点 $p'$ 对应于 $p \in M$。

\medskip\noindent
该断言蕴含引理。假设 $M_1 \to M$ 是仿射概形的平坦态射。
可将一切基变换到 $M_1$，而不影响假设 (1) -- (4)。
由引理 \ref{groupoids-lemma-invariants-base-change}，
$M_1$、resp.\ $M_1'$ 分别是 $U_1$ 上 $R_1$-不变函数、
resp.\ $U'_1$ 上 $R'_1$-不变函数的谱。
假设 $p' \in M'$ 映到 $p \in M$。令 $M_1$ 为
$\mathcal{O}_{M, p}$ 的严格 Hensel 化之谱，其闭点为
$p_1 \in M_1$。选择映到 $p_1$ 与 $p'$ 的点
$p'_1 \in M'_1$。由该断言得到
$$
(U'_1, R'_1, s'_1, t'_1, c'_1) =
(U_1, R_1, s_1, t_1, c_1) \amalg
(U''_1, R''_1, s''_1, t''_1, c''_1)
$$
并且相应地，作为 $M_1$ 上的概形有
$M'_1 = M_1 \amalg M''_1$。写 $M_1 = \Spec(C_1)$，
并把 $C_1 = \colim C_i$ 写成 \'etale $C$-代数的滤过余极限。
令 $M_i = \Spec(C_i)$。则 $M_1 = \lim M_i$，其他概形亦然。
由《概形的极限》引理 \ref{limits-lemma-descend-opens} 和
\ref{limits-lemma-descend-isomorphism}，可找到 $i$，使
$$
(U'_i, R'_i, s'_i, t'_i, c'_i) =
(U_i, R_i, s_i, t_i, c_i) \amalg
(U''_i, R''_i, s''_i, t''_i, c''_i)
$$
故 $M'_i = M_i \amalg M''_i$。特别地，经过 \'etale 基变换后，
$M' \to M$ 在 $p'$ 上方某点成为 \'etale 的。
这蕴含 $M' \to M$ 在 $p'$ 处为 \'etale
（例如由《概形态射》引理 \ref{morphisms-lemma-set-points-where-fibres-etale}）。
证明断言之后，我们将证明 $U' \cong M' \times_M U$。

\medskip\noindent
证明该断言。注意 $U_p$ 与 $U'_{p'}$ 都只有有限多个点。
对 $u \in U_p$，扩张 $\kappa(u)/\kappa(p)$ 是代数的，
故 $\kappa(u)$ 可分闭。由于 $U' \to U$ 为 \'etale，
态射 $U'_{p'} \to U_p$ 在剩余域扩张上诱导同构。
取映到 $u \in U_p$ 的 $u' \in U'_{p'}$。
由假设 (3)，概形论纤维的态射
$(s')^{-1}(u') \to s^{-1}(u)$、
$(t')^{-1}(u') \to t^{-1}(u)$ 及
$G'_{u'} \to G_u$ 都是同构。注意
$U_p = t(s^{-1}(u))$（集合论意义下），可知
$U'_{p'}$ 的点满射到 $U_p$ 的点。
假设 $U'_{p'}$ 中两点 $u'_1$ 与 $u'_2$ 映到
$U_p$ 的同一点 $u$。则存在点
$r' \in R'_{p'}$，满足 $s'(r') = u'_1$ 且 $t'(r') = u'_2$。
考虑两个域塔
$$
\kappa(r')/\kappa(u'_1)/\kappa(u)/\kappa(p) \quad
\kappa(r')/\kappa(u'_2)/\kappa(u)/\kappa(p)
$$
它们的两个“端点”分别与下面两个域塔的两个“端点”相同：
$$
\kappa(r')/\kappa(u'_1)/\kappa(p')/\kappa(p) \quad
\kappa(r')/\kappa(u'_2)/\kappa(p')/\kappa(p)
$$
由于 $(U'_{p'}, R'_{p'}, s', t', c')$ 是 $p'$ 上的群胚，
后两者诱导相同映射 $\kappa(p') \to \kappa(r')$。
又因 $\kappa(u)/\kappa(p)$ 纯不可分，所诱导的两个映射
$\kappa(u) \to \kappa(r')$ 相同。因此 $r'$ 映到纤维
$G_u$ 的一点。由假设 (3)，$r' \in (G')_{u'_1}$。
具体地，可把 $G$ 看成 $R$ 的闭子概形，其中后者经
$s$ 视为 $U$ 上的概形；再利用其向 $U'$ 的基变换给出
$G' \subset R'$。特别地，$u'_1 = u'_2$。
故 $U'_{p'} \to U_p$ 在点上双射，并在剩余域上诱导同构。
于是 $U'_{p'}$ 是有限闭点集
（《交换代数》引理 \ref{algebra-lemma-finite-residue-extension-closed}），
因而 $U'_{p'}$ 在 $U'$ 中闭。令 $J' \subset A'$ 为在集合论意义下
截出 $U'_{p'}$ 的根理想。

\medskip\noindent
断言证明的第二部分。令 $\mathfrak m \subset C$ 为极大理想。
由《代数进阶》引理 \ref{more-algebra-lemma-integral-over-henselian-pair}，
$(A, \mathfrak m A)$ 是 Hensel 对。令
$J = \sqrt{\mathfrak m A}$。由《代数进阶》引理 \ref{more-algebra-lemma-change-ideal-henselian-pair}，
$(A, J)$ 是 Hensel 对；根据上面的讨论，\'etale 环同态
$A \to A'$ 诱导同构 $A/J \to A'/J'$。
由《代数进阶》引理 \ref{more-algebra-lemma-characterize-henselian-pair}，
$A' = A \times A''$。考虑相应的不交并分解
$U' = U \amalg U''$。开集 $(s')^{-1}(U)$ 是 $R'$
中可特化到 $R'_{p'}$ 一点的点集；$(t')^{-1}(U)$ 亦然。
类似地，$(s')^{-1}(U'') = (t')^{-1}(U'')$，
因为这是不能特化到 $R'_{p'}$ 的点集。因此得到不交并分解
$$
(U', R', s', t', c') =
(U, R, s, t, c) \amalg
(U'', R'', s'', t'', c'')
$$
这立即给出 $M' = M \amalg M''$，断言得证。

\medskip\noindent
最后还需证明典范映射 $U' \to M' \times_M U$ 是同构。
它是 \'etale 态射
（《概形态射》引理 \ref{morphisms-lemma-etale-permanence}）。
另一方面，基变换到严格 Hensel 局部环
（如证明第三段所做），并利用断言证明过程中得到的双射
$U'_{p'} \to U_p$，可知 $U' \to M' \times_M U$
普遍双射（略去若干细节）。然而，普遍双射的 \'etale 态射
是同构（《下降》引理 \ref{descent-lemma-universally-injective-etale-open-immersion}），
证明完成。
\end{proof}

\begin{lemma}
\label{groupoids-lemma-basis}
设 $S$ 为概形，$(U, R, s, t, c)$ 为 $S$ 上的群胚概形。假设：
\begin{enumerate}
\item $U = \Spec(A)$ 与 $R = \Spec(B)$ 仿射；
\item 存在元素 $x_i \in A$（$i \in I$），使
$B = \bigoplus_{i \in I} s^\sharp(A)t^\sharp(x_i)$。
\end{enumerate}
则 $A = \bigoplus_{i\in I} Cx_i$，且 $B \cong A \otimes_C A$，
其中 $C \subset A$ 是 (\ref{groupoids-equation-invariants}) 中
$U$ 上的 $R$-不变函数。
\end{lemma}

\begin{proof}
本证明中以 $s, t : A \to B$ 代替
$s^\sharp, t^\sharp$，类似地以
$c : B \to B \otimes_{s, A, t} B$ 表示相应映射。
记 $p_0 : B \to B \otimes_{s, A, t} B$、
$b \mapsto b \otimes 1$，并记
$p_1 : B \to B \otimes_{s, A, t} B$、
$b \mapsto 1 \otimes b$。由引理 \ref{groupoids-lemma-diagram-pull}
及 $C$ 的定义，有交换图
$$
\xymatrix{
B \otimes_{s, A, t} B &
B \ar@<-1ex>[l]_-c \ar@<1ex>[l]^-{p_0} &
A \ar[l]^t \\
B \ar[u]^{p_1} &
A \ar@<-1ex>[l]_s \ar@<1ex>[l]^t \ar[u]_s &
C \ar[u] \ar[l]
}
$$
此外，左侧两个方块在环范畴中是余笛卡尔的，而顶行同构于图
$$
\xymatrix{
B \otimes_{t, A, t} B &
B \ar@<-1ex>[l]_-{p_1} \ar@<1ex>[l]^-{p_0} &
A \ar[l]^t
}
$$
由《下降》引理 \ref{descent-lemma-ff-exact}，
它是等化子图，因为条件 (2) 特别蕴含 $s$
（以及与之同构的箭头 $t$）忠实平坦。
底行按 $C$ 的定义也是等化子图。
利用 $x_i$ 可得到交换图
$$
\xymatrix{
B \otimes_{s, A, t} B &
B \ar@<-1ex>[l]_-c \ar@<1ex>[l]^-{p_0} &
A \ar[l]^t \\
\bigoplus_{i \in I} B x_i \ar[u]^{p_1} &
\bigoplus_{i \in I} A x_i \ar@<-1ex>[l]_s \ar@<1ex>[l]^t \ar[u]_s &
\bigoplus_{i \in I} C x_i \ar[u] \ar[l]
}
$$
其中，右侧竖直箭头把 $x_i$ 映到 $x_i$，
中间竖直箭头把 $x_i$ 映到 $t(x_i)$，左侧竖直箭头把
$x_i$ 映到
$c(t(x_i)) = t(x_i) \otimes 1 = p_0(t(x_i))$
（该等式来自引理 \ref{groupoids-lemma-diagram} 中图上半部分的交换性）。
于是该图交换。由假设，中间竖直箭头是同构。
由于左侧两个方块余笛卡尔，左侧竖直箭头也是同构。
另一方面，水平行正合（即均为等化子），故右侧竖直箭头也是同构。
\end{proof}

\begin{proposition}
\label{groupoids-proposition-finite-flat-equivalence}
设 $S$ 为概形，$(U, R, s, t, c)$ 为 $S$ 上的群胚概形。假设：
\begin{enumerate}
\item $U = \Spec(A)$ 与 $R = \Spec(B)$ 仿射；
\item $s, t : R \to U$ 有限局部自由；
\item $j = (t, s)$ 是等价关系。
\end{enumerate}
在此情形下，令 $C \subset A$ 如
(\ref{groupoids-equation-invariants})。则 $U \to M = \Spec(C)$
有限局部自由，且 $R = U \times_M U$。此外，$M$
表示 fppf 拓扑下的商层 $U/R$
（见定义 \ref{groupoids-definition-quotient-sheaf}）。
\end{proposition}

\begin{proof}
本证明中以 $s, t : A \to B$ 代替
$s^\sharp, t^\sharp$。由
引理 \ref{groupoids-lemma-criterion-quotient-representable}，
只需证明 $C \to A$ 有限局部自由，并且映射
$$
t \otimes s : A \otimes_C A \longrightarrow B
$$
是同构。首先，$j$ 是单态射，而且有限
（因为 $s$ 和 $t$ 已经有限）。因此由
《概形态射》引理 \ref{morphisms-lemma-finite-monomorphism-closed}，
$j$ 是闭浸入。故 $A \otimes_C A \to B$ 满。

\medskip\noindent
我们将如引理 \ref{groupoids-lemma-invariants-base-change} 那样，
沿平坦环同态 $C \to C'$ 作基变换，并使用不变函数的形成与平坦基变换交换
这一事实，见所引引理的 (3)。下面将证明：对每个素理想
$\mathfrak p \subset C$，存在局部平坦环同态
$C_{\mathfrak p} \to C_{\mathfrak p}'$，使得基变换到
$C_{\mathfrak p}'$ 后结论成立。这立即蕴含
$A \otimes_C A \to B$ 是单射（使用
《交换代数》引理 \ref{algebra-lemma-characterize-zero-local}）。
结合《交换代数》引理 \ref{algebra-lemma-local-flat-ff}、
\ref{algebra-lemma-flat-localization} 及
\ref{algebra-lemma-flatness-descends}，还可推出
$C \to A$ 平坦。又因 $U \to \Spec(C)$ 满
（引理 \ref{groupoids-lemma-points}），故 $C \to A$ 忠实平坦。
同构 $B \cong A \otimes_C A$ 蕴含 $A$ 是有限表现
$C$-模，见《交换代数》引理 \ref{algebra-lemma-descend-properties-modules}。
于是 $A$ 在 $C$ 上有限局部自由，见
《交换代数》引理 \ref{algebra-lemma-finite-projective}。

\medskip\noindent
由引理 \ref{groupoids-lemma-finite-locally-free-disjoint-free}，
$A$ 是有限多个环 $A_r$ 的积，$B$ 是有限多个环 $B_r$ 的积，
且群胚概形相应分解（见引理 \ref{groupoids-lemma-integral-over-invariants} 的证明）。
于是 $C$ 也是诸环 $C_r$ 的积，$C'$ 也相应分解。
因此可以并且确实假设环同态
$s, t : A \to B$ 有限局部自由、恒秩为 $r$。

\medskip\noindent
将使用的局部环同态 $C_{\mathfrak p} \to C_{\mathfrak p}'$
可以取为任意使 $C_{\mathfrak p}'$ 的剩余域为无限域的
局部平坦环同态。由《交换代数》引理 \ref{algebra-lemma-flat-local-given-residue-field}，这样的同态存在。

\medskip\noindent
假设 $C$ 是极大理想为 $\mathfrak m$、剩余域无限的局部环，
并假设 $s, t : A \to B$ 有限局部自由、恒秩为 $r > 0$。
由于 $C \subset A$ 整
（引理 \ref{groupoids-lemma-integral-over-invariants}），
位于 $\mathfrak m$ 上方的所有素理想都是极大理想，
且 $A$ 的所有极大理想都位于 $\mathfrak m$ 上方。
$C \subset B$ 也同样如此。选取 $A$ 中位于 $\mathfrak m$
上方的极大理想 $\mathfrak m'$（其存在性由
引理 \ref{groupoids-lemma-points}）。因为 $t : A \to B$ 有限局部自由，
$B$ 中位于 $\mathfrak m'$ 上方的极大理想至多有限多个。
再次由引理 \ref{groupoids-lemma-points}，$A$ 只有有限多个极大理想，
即 $A$ 半局部；进而 $B$ 也半局部。现在，因为
$t \otimes s : A \otimes_C A \to B$ 满，可把
《交换代数》引理 \ref{algebra-lemma-semi-local-module-basis-in-submodule}
应用于环同态 $C \to A$、$A$-模 $M = B$
（经 $t$ 视作 $A$-模）以及 $C$-子模 $s(A) \subset B$。
该引理给出 $x_1, \ldots, x_r \in A$，使 $M$ 是以
$s(x_1), \ldots, s(x_r)$ 为基的自由 $A$-模。
因此应用引理 \ref{groupoids-lemma-basis} 可知
$C \to A$ 有限自由，且 $B \cong A \otimes_C A$。
\end{proof}



\section{有限平坦群胚}
\label{groupoids-section-finite-flat-general}

\noindent
本节证明一个引理，用来说明：若每个等价类都包含于某个仿射开集，
则概形对有限平坦等价关系取商仍是概形。见
《代数空间的性质》命题 \ref{spaces-properties-proposition-finite-flat-equivalence-global}。

\begin{lemma}
\label{groupoids-lemma-find-invariant-affine}
设 $S$ 为概形，$(U, R, s, t, c)$ 为 $S$ 上的群胚概形。
假设 $s$, $t$ 有限局部自由。设 $u \in U$ 为一点，并且
$t(s^{-1}(\{u\}))$ 包含于 $U$ 的某个仿射开集。
则 $u$ 在 $U$ 中有一个 $R$-不变仿射开邻域。
\end{lemma}

\begin{proof}
由于 $s$ 有限局部自由，其纤维有限。因此
$t(s^{-1}(\{u\})) = \{u_1, \ldots, u_n\}$ 是有限集。
注意 $u \in \{u_1, \ldots, u_n\}$。
取包含 $\{u_1, \ldots, u_n\}$ 的仿射开集
$W \subset U$；特别地，$u \in W$。考虑
$Z = R \setminus s^{-1}(W) \cap t^{-1}(W)$。这是 $R$ 的闭子集。
像 $t(Z)$ 是 $U$ 的闭子集；粗略地说，它由 $U$ 中与
$U \setminus W$ 某一点 $R$-等价的点组成。因此
$W' = U \setminus t(Z)$ 是 $U$ 中包含于 $W$ 的
$R$-不变开子概形，并且
$\{u_1, \ldots, u_n\} \subset W'$。图示如下：
$$
\{u_1, \ldots, u_n\} \subset W' \subset W \subset U.
$$
取元素 $f \in \Gamma(W, \mathcal{O}_W)$，满足
$\{u_1, \ldots, u_n\} \subset D(f) \subset W'$。
这样的 $f$ 由《交换代数》引理 \ref{algebra-lemma-silly} 存在。
根据 $W'$ 的选择，有 $s^{-1}(W') \subset t^{-1}(W)$，
因而得到图
$$
\xymatrix{
s^{-1}(W') \ar[d]_s \ar[r]_-t & W \\
W'
}
$$
由假设，竖直箭头有限局部自由。令
$$
g = \text{Norm}_s(t^\sharp f) \in \Gamma(W', \mathcal{O}_{W'})
$$
按构造，$g$ 是 $W'$ 上在 $u$ 处非零的函数，因为
$t^\sharp(f)$ 在 $R$ 中位于 $u$ 上方的每一点处都非零，
这是由于 $f$ 在 $u_1, \ldots, u_n$ 处非零。
类似地，$D(g) \subset W'$ 恰由如下点 $w$ 构成：
$f$ 在与 $w$ 等价的任一点处都不为零。这说明 $D(g)$
是 $W'$ 的 $R$-不变仿射开集。最终图示为
$$
\{u_1, \ldots, u_n\} \subset D(g) \subset D(f) \subset W' \subset W \subset U
$$
故结论成立。
\end{proof}







\section{拟射影概形的下降}
\label{groupoids-section-quasi-projective}

\noindent
可用引理 \ref{groupoids-lemma-find-invariant-affine}
说明某一类下降数据是有效的。

\begin{lemma}
\label{groupoids-lemma-descend-along-finite-with-bound}
设 $X \to Y$ 为满的有限局部自由态射，$d$ 为正整数。
假设对每个几何点 $\bar{y} : \mathrm{Spec}(k) \to Y$，
纤维 $X_{\bar{y}}$ 至多有 $d$ 个点。设 $V$ 为 $X$ 上的概形，
并且对所有 $(y, v_1, \ldots, v_d)$，其中 $y \in Y$ 且
$v_1, \ldots, v_d \in V_y$，都存在仿射开集
$U \subset V$，使 $v_1, \ldots, v_d \in U$。
则 $V/X/Y$ 上的任意下降数据都是有效的。
\end{lemma}

\begin{proof}
设 $\varphi$ 为《下降》定义 \ref{descent-definition-descent-datum} 意义下的下降数据。
回忆从 $Y$ 上概形到相对于 $\{X \to Y\}$ 的下降数据的函子全忠实，
见《下降》引理 \ref{descent-lemma-refine-coverings-fully-faithful}。
因此利用《概形的构造》引理 \ref{constructions-lemma-relative-glueing}，只需在 $Y$ 仿射时证明。
略去若干细节（若 $Y$ 分离或其对角仿射，则可避免这一论证，
因为此时从仿射概形到 $X$ 的每个态射都仿射）。

\medskip\noindent
假设 $Y$ 仿射。若 $V$ 也仿射，则有效性由
《下降》引理 \ref{descent-lemma-affine} 得出。因此由
《下降》引理 \ref{descent-lemma-effective-for-fpqc-is-local-upstairs}，
只需证明 $V$ 的每个点 $v$ 都有一个
$\varphi$-不变仿射开邻域。考虑群胚
$(X, X \times_Y X, \text{pr}_1, \text{pr}_0, \text{pr}_{02})$。
由引理 \ref{groupoids-lemma-cartesian-equivalent-descent-datum}，
下降数据 $\varphi$ 确定且被如下群胚概形的笛卡尔态射确定：
$$
(V, R, s, t, c)
\longrightarrow
(X, X \times_Y X, \text{pr}_1, \text{pr}_0, \text{pr}_{02})
$$
该态射在 $\Spec(\mathbf{Z})$ 上。由于 $X \to Y$ 有限局部自由，
$\text{pr}_i : X \times_Y X \to X$，从而 $s$ 与 $t$
都有限局部自由，且其几何纤维至多有 $d$ 个点。特别地，
点 $v \in V$ 的 $R$-轨道 $t(s^{-1}(\{v\}))$
至多有 $d$ 个点。再次使用引理 \ref{groupoids-lemma-cartesian-equivalent-descent-datum} 的范畴等价，
可知 $V$ 的 $\varphi$-不变开集恰是 $V$ 的 $R$-不变开集。
由假设，存在包含轨道 $t(s^{-1}(\{v\}))$ 的 $V$ 的仿射开集，
因为该轨道的所有点都映到 $Y$ 的同一点。因此引理 \ref{groupoids-lemma-find-invariant-affine} 给出包含 $v$ 的
$R$-不变仿射开集。
\end{proof}

\begin{lemma}
\label{groupoids-lemma-descend-along-finite}
设 $X \to Y$ 为满的有限局部自由态射。 
设 $V$ 为 $X$ 上的概形，并且对所有 
$(y, v_1, \ldots, v_d)$，其中 $y \in Y$ 且
$v_1, \ldots, v_d \in V_y$，都存在仿射开集
$U \subset V$，使 $v_1, \ldots, v_d \in U$。
则 $V/X/Y$ 上的任意下降数据都是有效的。
\end{lemma}

\begin{proof}
与引理 \ref{groupoids-lemma-descend-along-finite-with-bound} 的证明一样，
可假设 $Y$ 仿射，特别是拟紧。于是存在整数 $d$，使
$X \to Y$ 的所有几何纤维至多有 $d$ 个点，见
《概形态射》第 \ref{morphisms-section-universally-bounded} 节中的讨论。
关于 $V$ 的假设显然蕴含可应用
引理 \ref{groupoids-lemma-descend-along-finite-with-bound}，证明完成。
\end{proof}

\begin{lemma}
\label{groupoids-lemma-descend-along-finite-quasi-projective}
设 $X \to Y$ 为满的有限局部自由态射，$V$ 为 $X$ 上的概形，
并假设下列条件之一成立：
\begin{enumerate}
\item $V \to X$ 射影；
\item $V \to X$ 拟射影；
\item $V$ 上存在充足可逆层；
\item $V$ 上存在 $X$-充足可逆层；
\item $V$ 上存在 $X$-很充足可逆层。
\end{enumerate}
则 $V/X/Y$ 上的任意下降数据都是有效的。
\end{lemma}

\begin{proof}
验证引理 \ref{groupoids-lemma-descend-along-finite} 中的条件。
设 $y \in Y$，且 $v_1, \ldots, v_d \in V$ 是 $y$ 上方的点。
情形 (1) 是 (2) 的特殊情形，见
《概形态射》引理 \ref{morphisms-lemma-projective-quasi-projective}。
情形 (2) 是 (4) 的特殊情形，见
《概形态射》定义 \ref{morphisms-definition-quasi-projective}。
若 $V$ 上存在充足可逆层，则由《概形的性质》引理 \ref{properties-lemma-ample-finite-set-in-affine}，
存在包含 $v_1, \ldots, v_d$ 的仿射开集。因此 (3) 成立。
在情形 (4) 与 (5) 中，可无妨把 $Y$ 替换为 $y$ 的仿射开邻域；
此时 $X$ 也仿射。在情形 (4) 中，由《概形态射》定义 \ref{morphisms-definition-relatively-ample}，$V$ 有充足可逆层，
故结论由情形 (3) 得出。在情形 (5) 中，可把 $V$
替换为包含 $v_1, \ldots, v_d$ 的拟紧开集，再由
《概形态射》引理 \ref{morphisms-lemma-ample-very-ample}
归约到情形 (4)。
\end{proof}

\begin{lemma}
\label{groupoids-lemma-descend-along-finite-radicial}
设 $X \to Y$ 为满的有限局部自由根态射。 
设 $V$ 为 $X$ 上的概形。则 $V/X/Y$ 上的任意下降数据都是有效的。
\end{lemma}

\begin{proof}
取 $d = 1$，应用引理 \ref{groupoids-lemma-descend-along-finite-with-bound}。
\end{proof}
% Generated mechanically from the frozen more-groupoids.tex target.
% Do not edit this generated file; edit the bound locale source instead.

% OK, start here.
%

\setcounter{chapter}{39}
\chapter{群胚概形进阶}



\phantomsection
\label{more-groupoids-section-phantom}


\section{引言}
\label{more-groupoids-section-introduction}

\noindent
本章讨论群胚概形的若干进阶论题。尽管各项结果都以群胚概形来表述，
读者仍应始终将如下 $2$-笛卡尔图牢记在心：
\begin{equation}
\label{more-groupoids-equation-quotient-stack}
\vcenter{
\xymatrix{
R \ar[r] \ar[d] & U \ar[d] \\
U \ar[r] & [U/R]
}
}
\end{equation}
其中 $[U/R]$ 是商叠，见《代数空间中的群胚》注 \ref{spaces-groupoids-remark-fundamental-square}。
许多结果的动机都来自对这一图表的思考。例如，可参阅
Keel 和 Mori 的优美论文 \cite{K-M}。





\section{记号}
\label{more-groupoids-section-notation}

\noindent
我们继续沿用《群胚概形》第 \ref{groupoids-section-notation} 节中引入的约定与记号。






\section{若干常用图}
\label{more-groupoids-section-diagrams}

\noindent
为便于本章引用，我们简要重述《群胚概形》引理 \ref{groupoids-lemma-diagram} 与
\ref{groupoids-lemma-diagram-pull} 的结果。
设 $S$ 为概形，$(U, R, s, t, c)$ 为 $S$ 上的群胚概形。
在交换图
\begin{equation}
\label{more-groupoids-equation-diagram}
\vcenter{
\xymatrix{
& U & \\
R \ar[d]_s \ar[ru]^t &
R \times_{s, U, t} R
\ar[l]^-{\text{pr}_0} \ar[d]^{\text{pr}_1} \ar[r]_-c &
R \ar[d]^s \ar[lu]_t \\
U & R \ar[l]_t \ar[r]^s & U
}
}
\end{equation}
中，下方两个方块都是纤维积方块。此外，上方的三角形
（实际上是一个方块）也是笛卡尔的。

\medskip\noindent
图
\begin{equation}
\label{more-groupoids-equation-pull}
\vcenter{
\xymatrix{
R \times_{t, U, t} R
\ar@<1ex>[r]^-{\text{pr}_1} \ar@<-1ex>[r]_-{\text{pr}_0}
\ar[d]_{\text{pr}_0 \times c \circ (i, 1)} &
R \ar[r]^t \ar[d]^{\text{id}_R} &
U \ar[d]^{\text{id}_U} \\
R \times_{s, U, t} R
\ar@<1ex>[r]^-c \ar@<-1ex>[r]_-{\text{pr}_0} \ar[d]_{\text{pr}_1} &
R \ar[r]^t \ar[d]^s &
U \\
R \ar@<1ex>[r]^s \ar@<-1ex>[r]_t &
U
}
}
\end{equation}
交换。由给定的竖直映射，上方两行彼此同构；
左下方的两个方块均为笛卡尔方块。





\section{微分层}
\label{more-groupoids-section-differentials}

\noindent
下面的引理对应于《群胚概形》引理 \ref{groupoids-lemma-group-scheme-module-differentials}。

\begin{lemma}
\label{more-groupoids-lemma-sheaf-differentials}
设 $S$ 为概形，$(U, R, s, t, c)$ 为 $S$ 上的群胚概形。
经 $t$ 把 $R$ 看成 $U$ 上的概形，则其微分层是浸入
$e : U \to R$ 的余法层经 $t$ 拉回所得层的商。换言之，
存在标准满射 $t^*\mathcal{C}_{U/R} \to \Omega_{R/U}$。
若 $s$ 平坦，则该映射为同构。
\end{lemma}

\begin{proof}
注意，$e : U \to R$ 是态射 $s$ 的截面，因而是浸入；
见《概形》引理 \ref{schemes-lemma-section-immersion}。
考虑图
$$
\xymatrix{
R \ar[r]_-{(1, i)} \ar[d]_t &
R \times_{s, U, t} R \ar[d]^c \ar[rr]_{(\text{pr}_0, i \circ \text{pr}_1)} & &
R \times_{t, U, t} R \\
U \ar[r]^e &
R
}
$$
左侧方块是笛卡尔的，因为若 $a \circ b = e$，则
$b = i(a)$。两个水平映射的复合是 $t : R \to U$ 的对角态射，
而右上方的水平箭头是同构。因此，由于 $\Omega_{R/U}$
是该复合的余法层，它同构于 $(1, i)$ 的余法层。
由《概形态射》引理 \ref{morphisms-lemma-conormal-functorial-flat}，得到满射
$t^*\mathcal{C}_{U/R} \to \Omega_{R/U}$；若 $c$ 平坦，
该映射即为同构。根据图 (\ref{more-groupoids-equation-pull}) 的性质，
$c$ 是 $s$ 的基变换。因此若 $s$ 平坦，则 $c$ 平坦；
见《概形态射》引理 \ref{morphisms-lemma-base-change-flat}。
\end{proof}






\section{局部结构}
\label{more-groupoids-section-local}

\noindent
设 $S$ 为概形，$(U, R, s, t, c, e, i)$ 为 $S$ 上的群胚概形，
并设 $u \in U$ 为一点。本节说明局部环
$$
A = \mathcal{O}_{U, u}
\quad\text{且}\quad
B = \mathcal{O}_{R, e(u)}
$$
上得到的结构。约定仍以相应字母表示由态射
$s, t, c, e, i$ 诱导的局部环同态。特别地，有局部环的交换图
$$
\xymatrix{
A \ar[rd]_t \ar[rrd]^1 \\
& B \ar[r]^e & A \\
A \ar[ru]^s \ar[rru]_1
}
$$
因此，若 $I \subset B$ 表示 $e : B \to A$ 的核，则
$B = s(A) \oplus I = t(A) \oplus I$。记
$$
C = \mathcal{O}_{R \times_{s, U, t} R, (e(u), e(u))}
$$
则
$$
C = (B \otimes_{s, A, t} B)_{\mathfrak m_B \otimes B + B \otimes \mathfrak m_B}
$$
令 $J \subset C$ 为 $C$ 中由 $I \otimes B + B \otimes I$
生成的理想。于是 $J$ 也是局部环同态
$$
(e, e) : C \longrightarrow A
$$
的核。合成律 $c : R \times_{s, U, t} R \to R$ 对应于环映射
$$
c : B \longrightarrow C
$$
且该映射把 $I$ 映入 $J$。

\begin{lemma}
\label{more-groupoids-lemma-first-order-structure-c}
由 $c$ 诱导的映射 $I/I^2 \to J/J^2$ 是复合
$$
I/I^2 \xrightarrow{(1, 1)} I/I^2 \oplus I/I^2 \to J/J^2
$$
其中第二个箭头来自等式 $J = (I \otimes B + B \otimes I)C$。
映射 $i : B \to B$ 诱导映射 $-1 : I/I^2 \to I/I^2$。
\end{lemma}

\begin{proof}
要描述从 $C$ 到另一局部环的局部同态，只需说明它在
$b_1 \otimes b_2$ 形元素上的作用。依此观察，有两个标准映射
$$
e_2 : C \to B,\ b_1 \otimes b_2 \mapsto b_1s(e(b_2)),\quad
e_1 : C \to B,\ b_1 \otimes b_2 \mapsto t(e(b_1))b_2
$$
它们对应于嵌入
$R \to R \times_{s, U, t} R$，分别由
$r \mapsto (r, e(s(r)))$ 与 $r \mapsto (e(t(r)), r)$ 给出。
这些映射诱导 $J/J^2 \to I/I^2$，合在一起给出引理中映射
$I/I^2 \oplus I/I^2 \to J/J^2$ 的逆。因此，只需证明
$e_1 \circ c : B \to B$ 与 $e_2 \circ c : B \to B$
均为恒等映射。这是因为两个复合
$R \to R \times_{s, U, t} R \to R$ 都是恒等映射。

\medskip\noindent
关于 $i$ 的断言由关于 $c$ 的断言以及等式
$c \circ (1, i) = e \circ t$ 得出。其余细节从略。
\end{proof}






\section{群胚的性质}
\label{more-groupoids-section-technical-lemma}

\noindent
设 $(U, R, s, t, c)$ 为群胚概形。本节结果背后的思路是：
$s: R \to U$ 应当是态射 $U \to [U/R]$ 的基变换
（见图 (\ref{more-groupoids-equation-quotient-stack})）。因此，
$s : R \to U$ 的局部性质应当反映态射 $U \to [U/R]$ 的局部性质。
这一思路不能直接照搬，因为 $[U/R]$ 未必是代数叠，
从而无法谈论 $U \to [U/R]$ 的几何或代数性质。
不过，即使完全不提商叠，仍可实现其中的一部分。

\medskip\noindent
下面是此类结果的第一个例子。粗略地说，引理中得到的开集
$W \subset U'$ 是态射 $U' \to [U/R]$ 具有性质
$\mathcal{P}$ 的轨迹。

\begin{lemma}
\label{more-groupoids-lemma-local-source}
设 $S$ 为概形，$(U, R, s, t, c, e, i)$ 为 $S$ 上的群胚，
且 $g : U' \to U$ 为概形态射。记 $h$ 为复合
$$
\xymatrix{
h : U' \times_{g, U, t} R \ar[r]_-{\text{pr}_1} & R \ar[r]_s & U.
}
$$
设 $\mathcal{P}, \mathcal{Q}, \mathcal{R}$ 为概形态射的性质。
假设
\begin{enumerate}
\item $\mathcal{R} \Rightarrow \mathcal{Q}$；
\item $\mathcal{Q}$ 在基变换与复合下保持；
\item 对任意具有 $\mathcal{Q}$ 的态射 $f : X \to Y$，
存在最大的开集 $W(\mathcal{P}, f) \subset X$，使得
$f|_{W(\mathcal{P}, f)}$ 具有 $\mathcal{P}$；且
\item 对任意具有 $\mathcal{Q}$ 的态射 $f : X \to Y$
以及任意具有 $\mathcal{R}$ 的态射 $Y' \to Y$，有
$Y' \times_Y W(\mathcal{P}, f) = W(\mathcal{P}, f')$，
其中 $f' : X_{Y'} \to Y'$ 是 $f$ 的基变换。
\end{enumerate}
若 $s, t$ 具有 $\mathcal{R}$，且 $g$ 具有 $\mathcal{Q}$，
则存在开子概形 $W \subset U'$，使得
$W \times_{g, U, t} R = W(\mathcal{P}, h)$。
\end{lemma}

\begin{proof}
注意下图交换：
$$
\xymatrix{
U' \times_{g, U, t} R \times_{t, U, t} R
\ar[rr]_-{\text{pr}_{12}}
\ar@<1ex>[d]^-{\text{pr}_{02}} \ar@<-1ex>[d]_-{\text{pr}_{01}} & &
R \times_{t, U, t} R
\ar@<1ex>[d]^-{\text{pr}_1} \ar@<-1ex>[d]_-{\text{pr}_0}
\\
U' \times_{g, U, t} R \ar[rr]^{\text{pr}_1} & & R
}
$$
且两个方块均为笛卡尔方块（这里用到了两个映射
$t \circ \text{pr}_i : R \times_{t, U, t} R \to U$ 相等）。
将此事实与图 (\ref{more-groupoids-equation-pull}) 的性质结合，得到交换图
$$
\xymatrix{
U' \times_{g, U, t} R \times_{t, U, t} R
\ar[rr]_-{c \circ (i, 1)}
\ar@<1ex>[d]^-{\text{pr}_{02}} \ar@<-1ex>[d]_-{\text{pr}_{01}} & &
R
\ar@<1ex>[d]^-{s} \ar@<-1ex>[d]_-{t}
\\
U' \times_{g, U, t} R \ar[rr]^h & & U
}
$$
其中两个方块均为笛卡尔方块。

\medskip\noindent
假设 $s, t$ 具有 $\mathcal{R}$，且 $g$ 具有 $\mathcal{Q}$。
态射 $h$ 是 $s$（它具有 $\mathcal{R}$，因而具有
$\mathcal{Q}$）与 $g$ 的一个基变换（它具有
$\mathcal{Q}$）的复合，故该态射具有 $\mathcal{Q}$。
因此 $W(\mathcal{P}, h) \subset U' \times_{g, U, t} R$ 存在。
由假设，
$\text{pr}_{01}^{-1}(W(\mathcal{P}, h)) =
\text{pr}_{02}^{-1}(W(\mathcal{P}, h))$，
因为二者都是使 $c \circ (i, 1)$ 具有 $\mathcal{P}$ 的最大开集。
注意投影 $U' \times_{g, U, t} R \to U'$ 有截面
$\sigma : U' \to U' \times_{g, U, t} R$，
$u' \mapsto (u', e(g(u')))$。此外，在同构
$$
(U' \times_{g, U, t} R) \times_{U'} (U' \times_{g, U, t} R)
=
U' \times_{g, U, t} R \times_{t, U, t} R
$$
下，左侧的两个投影到 $U' \times_{g, U, t} R$
分别与右侧的态射 $\text{pr}_{01}$ 和 $\text{pr}_{02}$ 一致。
由于
$\text{pr}_{01}^{-1}(W(\mathcal{P}, h)) =
\text{pr}_{02}^{-1}(W(\mathcal{P}, h))$，
可知 $W(\mathcal{P}, h)$ 是 $U$ 的某个子集的逆像；
这个子集必为开集
$W = \sigma^{-1}(W(\mathcal{P}, h))$。
\end{proof}

\begin{remark}
\label{more-groupoids-remark-local-source-warning}
警告：使用引理 \ref{more-groupoids-lemma-local-source} 时必须谨慎。
例如，它适用于 $\mathcal{P}=$“平坦”、$\mathcal{Q}=$“空条件”，
以及 $\mathcal{R}=$“平坦且局部有限表示”。然而，
给定概形态射 $f : X \to Y$，使 $f|_W$ 平坦的最大开集
$W \subset X$ {\it 并不是} $f$ 平坦的点所成的集合！
\end{remark}

\begin{remark}
\label{more-groupoids-remark-local-source-apply}
尽管有注 \ref{more-groupoids-remark-local-source-warning} 中的警告，
在某些情形下仍可使用引理 \ref{more-groupoids-lemma-local-source}，
而不会造成太大歧义。下面列出这些情形。每种情形中，
我们略去假设 (1) 与 (2) 的验证，并给出蕴含 (3) 与 (4)
的参考文献：
\begin{enumerate}
\item $\mathcal{Q} = \mathcal{R} =$“局部有限型”，且
$\mathcal{P} =$“相对维数 $\leq d$”。
见《概形态射》定义 \ref{morphisms-definition-relative-dimension-d} 以及
《概形态射》引理 \ref{morphisms-lemma-openness-bounded-dimension-fibres} 和
\ref{morphisms-lemma-dimension-fibre-after-base-change}。
\item $\mathcal{Q} = \mathcal{R} =$“局部有限型”，且
$\mathcal{P} =$“局部拟有限”。
这是上一项中 $d = 0$ 的情形，见《概形态射》引理 \ref{morphisms-lemma-locally-quasi-finite-rel-dimension-0}。
\item $\mathcal{Q} = \mathcal{R} =$“局部有限型”，且
$\mathcal{P} =$“非分歧”。
见《概形态射》引理 \ref{morphisms-lemma-unramified-characterize} 和
\ref{morphisms-lemma-set-points-where-fibres-unramified}。
\end{enumerate}
上述情形的有趣之处在于：无需假设 $s, t$ 平坦，
便可对态射 $h$ 具有性质 $\mathcal{P}$ 的轨迹作出结论。
继续列举：
\begin{enumerate}
\item[(4)] $\mathcal{Q} =$“局部有限表示”，
$\mathcal{R} =$“平坦且局部有限表示”，且
$\mathcal{P} =$“平坦”。见《态射进阶》定理 \ref{more-morphisms-theorem-openness-flatness} 和
引理 \ref{more-morphisms-lemma-flat-locus-base-change}。
\item[(5)] $\mathcal{Q} =$“局部有限表示”，
$\mathcal{R} =$“平坦且局部有限表示”，且
$\mathcal{P}=$“Cohen--Macaulay”。见《态射进阶》定义 \ref{more-morphisms-definition-CM} 以及《态射进阶》引理 \ref{more-morphisms-lemma-base-change-CM} 和
\ref{more-morphisms-lemma-flat-finite-presentation-CM-open}。
\item[(6)] $\mathcal{Q} =$“局部有限表示”，
$\mathcal{R} =$“平坦且局部有限表示”，且
$\mathcal{P}=$“syntomic”。使用《概形态射》引理 \ref{morphisms-lemma-set-points-where-fibres-lci}
（该轨迹自动为开集）。
\item[(7)] $\mathcal{Q} =$“局部有限表示”，
$\mathcal{R} =$“平坦且局部有限表示”，且
$\mathcal{P}=$“光滑”。见《概形态射》引理 \ref{morphisms-lemma-set-points-where-fibres-smooth}
（该轨迹自动为开集）。
\item[(8)] $\mathcal{Q} =$“局部有限表示”，
$\mathcal{R} =$“平坦且局部有限表示”，且
$\mathcal{P}=$“\'etale”。见《概形态射》引理 \ref{morphisms-lemma-set-points-where-fibres-etale}
（该轨迹自动为开集）。
\end{enumerate}
\end{remark}

\noindent
下面是第二个结果。应把 $R$-不变开集 $W \subset U$
看作 $[U/R]$ 的最大开集的逆像；在该开集上，
态射 $U \to [U/R]$ 具有性质 $\mathcal{P}$。

\begin{lemma}
\label{more-groupoids-lemma-property-invariant}
设 $S$ 为概形，$(U, R, s, t, c)$ 为 $S$ 上的群胚。
令 $\tau \in \{Zariski, \linebreak[0] fppf,
\linebreak[0] \etale, \linebreak[0]
smooth, \linebreak[0] syntomic\}$\footnote{$fpqc$ 未列入其中并非排印错误。}。
设 $\mathcal{P}$ 为概形态射的一项性质，且它在目标上关于
$\tau$ 为局部的（《下降》定义 \ref{descent-definition-property-morphisms-local}）。
假设 $\{s : R \to U\}$ 与 $\{t : R \to U\}$
都是 $\tau$-拓扑的覆盖。令 $W \subset U$ 为最大的开子概形，
使 $s|_{s^{-1}(W)} : s^{-1}(W) \to W$ 具有性质
$\mathcal{P}$。则 $W$ 是 $R$-不变的，见《群胚概形》定义 \ref{groupoids-definition-invariant-open}。
\end{lemma}

\begin{proof}
开集 $W \subset U$ 的存在性与性质见《下降》引理 \ref{descent-lemma-largest-open-of-the-base}。
在图 (\ref{more-groupoids-equation-diagram}) 中，令 $W_1 \subset R$
为最大的开子概形，使态射
$\text{pr}_1 : R \times_{s, U, t} R \to R$
在其上具有性质 $\mathcal{P}$。由上述《下降》引理 \ref{descent-lemma-largest-open-of-the-base} 以及
$\{s : R \to U\}$ 和 $\{t : R \to U\}$
都是 $\tau$-拓扑覆盖这一假设，得到所需等式
$t^{-1}(W) = W_1 = s^{-1}(W)$。
\end{proof}

\begin{lemma}
\label{more-groupoids-lemma-property-G-invariant}
设 $S$ 为概形，$(U, R, s, t, c)$ 为 $S$ 上的群胚，
并设 $G \to U$ 为其稳定子群概形。
令 $\tau \in \{fppf, \linebreak[0] \etale, \linebreak[0]
smooth, \linebreak[0] syntomic\}$。
设 $\mathcal{P}$ 为在目标上关于 $\tau$ 为局部的态射性质。
假设 $\{s : R \to U\}$ 与 $\{t : R \to U\}$
都是 $\tau$-拓扑的覆盖。令 $W \subset U$ 为最大的开子概形，
使 $G_W \to W$ 具有性质 $\mathcal{P}$。则 $W$ 是
$R$-不变的（见《群胚概形》定义 \ref{groupoids-definition-invariant-open}）。
\end{lemma}

\begin{proof}
开集 $W \subset U$ 的存在性与性质见《下降》引理 \ref{descent-lemma-largest-open-of-the-base}。态射
$$
G \times_{U, t} R \longrightarrow R \times_{s, U} G, \quad
(g, r) \longmapsto (r, r^{-1} \circ g \circ r)
$$
是 $R$ 上的同构（其中 $\circ$ 表示群胚中的合成）。
因此，由上述《下降》引理 \ref{descent-lemma-largest-open-of-the-base} 所证明的
$W$ 的性质，得到 $s^{-1}(W) = t^{-1}(W)$。
\end{proof}



\section{纤维的比较}
\label{more-groupoids-section-fibres}

\noindent
设 $(U, R, s, t, c, e, i)$ 为 $S$ 上的群胚概形。
图 (\ref{more-groupoids-equation-diagram}) 提供了比较群胚中映射
$s : R \to U$ 的各纤维的方法。对点 $u \in U$，
以 $F_u = s^{-1}(u)$ 表示 $s : R \to U$ 在 $u$ 上方的概形论纤维。
例如，该图说明：若 $u, u' \in U$ 为满足
$s(r) = u$ 且 $t(r) = u'$ 的点，则
$(F_u)_{\kappa(r)} \cong (F_{u'})_{\kappa(r)}$。
这是下面更一般且更精确的引理 \ref{more-groupoids-lemma-two-fibres}
的特殊情形；取 $r' = i(r)$ 即可看出。

\medskip\noindent
由概形 $X$ 与点 $x \in X$ 组成的一对 $(X, x)$，
有时称为{\it $X$ 在 $x$ 处的芽}。
一个{\it 芽的态射} $f : (X, x) \to (S, s)$
是定义在 $x$ 的某个开邻域上的态射 $f : U \to S$，
并满足 $f(x) = s$。称两个这样的 $f$, $f'$
给出同一个芽态射，当且仅当 $f$ 与 $f'$
在 $x$ 的某个开邻域上一致。
令 $\tau \in \{Zariski, \etale, smooth, syntomic, fppf\}$。
我们暂时引入如下概念：称两个芽态射
$f : (X, x) \to (S, s)$ 与 $f' : (X', x') \to (S', s')$
{\it 在基底上关于 $\tau$-拓扑局部同构}，若存在带基点概形
$(S'', s'')$ 及芽态射
$g : (S'', s'') \to (S, s)$、$g' : (S'', s'') \to (S', s')$，
使得
\begin{enumerate}
\item $g$ 与 $g'$ 在 $s''$ 处为开浸入（相应地为\ \'etale、光滑、
syntomic，或平坦且局部有限表示）；
\item 对位于 $(s'', x)$ 与 $(s'', x')$ 上方的某些点
$\tilde x$ 与 $\tilde x'$，存在 $(S'', s'')$ 上的芽同构
$$
(S'' \times_{g, S, f} X, \tilde x)
\cong
(S'' \times_{g', S', f'} X', \tilde  x')
$$
。
\end{enumerate}
最后，若存在如上的 $S'', s'', g, g'$，但把 (1) 换成
$g$ 与 $g'$ 在 $s''$ 处平坦这一条件，则简称芽态射
$f : (X, x) \to (S, s)$ 与 $f' : (X', x') \to (S', s')$
{\it 在基底上平坦局部同构}。这比上面的任何
$\tau$ 条件都弱得多，因为平坦态射未必是开映射。

\begin{lemma}
\label{more-groupoids-lemma-two-fibres}
设 $S$ 为概形，$(U, R, s, t, c)$ 为 $S$ 上的群胚。
设 $r, r' \in R$ 且在 $U$ 中有 $t(r) = t(r')$。
令 $u = s(r)$、$u' = s(r')$，并以
$F_u = s^{-1}(u)$ 与 $F_{u'} = s^{-1}(u')$ 表示概形论纤维。
\begin{enumerate}
\item 存在公共域扩张 $\kappa(u) \subset k$、
$\kappa(u') \subset k$ 以及同构
$(F_u)_k \cong (F_{u'})_k$。
\item 可选取 (1) 中的同构，使得 $r$ 上方的某个点映到
$r'$ 上方的某个点。
\item 若态射 $s$, $t$ 平坦，则芽态射
$s : (R, r) \to (U, u)$ 与 $s : (R, r') \to (U, u')$
在基底上平坦局部同构。
\item 若态射 $s$, $t$ 为 \'etale
（相应地为\ 光滑、syntomic，或平坦且局部有限表示），则芽态射
$s : (R, r) \to (U, u)$ 与 $s : (R, r') \to (U, u')$
在基底上关于 \'etale（相应地为\ 光滑、syntomic 或 fppf）
拓扑局部同构。
\end{enumerate}
\end{lemma}

\begin{proof}
我们反复使用图 (\ref{more-groupoids-equation-diagram}) 的存在性与性质。
由该图的性质以及《概形》引理 \ref{schemes-lemma-points-fibre-product}，存在
$R \times_{s, U, t} R$ 的一点 $\xi$，使
$\text{pr}_0(\xi) = r$ 且 $c(\xi) = r'$。令
$\tilde r = \text{pr}_1(\xi) \in R$。

\medskip\noindent
证明 (1)。令 $k = \kappa(\tilde r)$。由于
$t(\tilde r) = u$ 且 $s(\tilde r) = u'$，
$k$ 是 $\kappa(u)$ 与 $\kappa(u')$ 的公共扩张；
事实上，$(F_u)_k$ 与 $(F_{u'})_k$ 都同构于
$\text{pr}_1 : R \times_{s, U, t} R \to R$
在 $\tilde r$ 上方的纤维。因此 (1) 得证。

\medskip\noindent
(2) 由点 $\xi$ 映到 $r$，\ 相应地映到 $r'$ 立即得出。

\medskip\noindent
(3) 由上述结论与定义显然成立（以点 $\xi$
选取 $\tilde u$ 与 $\tilde u'$）。

\medskip\noindent
若 $s$ 与 $t$ 平坦且有限表示，则它们是开态射
（《概形态射》引理 \ref{morphisms-lemma-fppf-open}）。
因此，$\tilde r$ 的某个仿射开邻域 $V''$ 的像覆盖
$u$ 的某个开邻域 $V$，相应地\ 覆盖 $u'$ 的某个开邻域 $V'$。
利用这些邻域即可验证“在基底上关于
$\tau$-拓扑局部同构”定义中的性质 (1) 与 (2)。
\end{proof}







\section{Cohen--Macaulay 呈示}
\label{more-groupoids-section-CM}

\noindent
给定任意满足 $s, t$ 平坦且局部有限表示的群胚
$(U, R, s, t, c)$，存在一个“等价”的群胚
$(U', R', s', t', c')$，使 $s'$ 与 $t'$ 都是
Cohen--Macaulay 态射（并且局部有限表示）。
关于 Cohen--Macaulay 态射的更多内容，见
《态射进阶》第 \ref{more-morphisms-section-CM} 节。
这里的“等价”可以理解为商叠 $[U/R]$ 与 $[U'/R']$
等价；见《代数空间中的群胚》第 \ref{spaces-groupoids-section-stacks} 节以及第 \ref{spaces-groupoids-section-quotient-stack-restrict} 节。

\begin{lemma}
\label{more-groupoids-lemma-make-CM}
设 $S$ 为概形，$(U, R, s, t, c)$ 为 $S$ 上的群胚。
假设 $s$ 与 $t$ 平坦且局部有限表示。则存在开集
$U' \subset U$，使得
\begin{enumerate}
\item $t^{-1}(U') \subset R$ 是 $R$ 中使态射 $s$
为 Cohen--Macaulay 的最大开子概形；
\item $s^{-1}(U') \subset R$ 是 $R$ 中使态射 $t$
为 Cohen--Macaulay 的最大开子概形；
\item 态射 $t|_{s^{-1}(U')} : s^{-1}(U') \to U$ 满；
\item 态射 $s|_{t^{-1}(U')} : t^{-1}(U') \to U$ 满；且
\item $R$ 在 $U'$ 上的限制
$R' = s^{-1}(U') \cap t^{-1}(U')$ 定义一个群胚
$(U', R', s', t', c')$，其态射 $s'$ 与 $t'$
均为 Cohen--Macaulay 且局部有限表示。
\end{enumerate}
\end{lemma}

\begin{proof}
在引理 \ref{more-groupoids-lemma-local-source} 中取
$g = \text{id}$、
$\mathcal{Q} =$“局部有限表示”、
$\mathcal{R} =$“平坦且局部有限表示”，以及
$\mathcal{P}=$“Cohen--Macaulay”；见注 \ref{more-groupoids-remark-local-source-apply}。由此得到开集 $U' \subset U$，
使 $t^{-1}(U') \subset R$ 是 $R$ 中使态射 $s$
为 Cohen--Macaulay 的最大开子概形。这证明了 (1)。
令 $i : R \to R$ 为群胚的取逆映射。由于 $i$ 是同构，
且 $s \circ i = t$、$t \circ i = s$，可知
$s^{-1}(U')$ 也是 $R$ 中使 $t$ 为 Cohen--Macaulay
的最大开集。这证明了 (2)。
由《态射进阶》引理 \ref{more-morphisms-lemma-flat-finite-presentation-CM-open}，
开子集 $t^{-1}(U')$ 在 $s : R \to U$ 的每个纤维中稠密。
这证明了 (3)；(4) 同理。(5) 是 (1)、(2) 以及《群胚概形》第 \ref{groupoids-section-restrict-groupoid} 节中关于限制的讨论
的形式推论。
\end{proof}








\section{群胚的限制}
\label{more-groupoids-section-restricting-groupoids}

\noindent
本节汇集若干引理，讨论群胚的哪些性质会由限制继承。
这些引理大多可由考察限制的定义图
\begin{equation}
\label{more-groupoids-equation-restriction}
\vcenter{
\xymatrix{
R' \ar[d] \ar[r] \ar@/_3pc/[dd]_{t'} \ar@/^1pc/[rr]^{s'}&
R \times_{s, U} U' \ar[r] \ar[d] &
U' \ar[d]^g \\
U' \times_{U, t} R \ar[d] \ar[r] &
R \ar[r]^s \ar[d]_t &
U \\
U' \ar[r]^g &
U
}
}
\end{equation}
证明。见《群胚概形》引理 \ref{groupoids-lemma-restrict-groupoid}。

\begin{lemma}
\label{more-groupoids-lemma-restrict-preserves-type}
设 $S$ 为概形，$(U, R, s, t, c)$ 为 $S$ 上的群胚概形，
$g : U' \to U$ 为概形态射，并设
$(U', R', s', t', c')$ 为 $(U, R, s, t, c)$ 经 $g$ 的限制。
\begin{enumerate}
\item 若 $s, t$ 局部有限型，且 $g$ 局部有限型，则
$s', t'$ 局部有限型。
\item 若 $s, t$ 局部有限表示，且 $g$ 局部有限表示，则
$s', t'$ 局部有限表示。
\item 若 $s, t$ 平坦，且 $g$ 平坦，则 $s', t'$ 平坦。
\item 此处再添加若干条目。
\end{enumerate}
\end{lemma}

\begin{proof}
局部有限型这一性质在复合与任意基变换下稳定，见
《概形态射》引理 \ref{morphisms-lemma-composition-finite-type}
和 \ref{morphisms-lemma-base-change-finite-type}。
因此 (1) 由图 (\ref{more-groupoids-equation-restriction}) 立即得出。
其余情形见《概形态射》引理 \ref{morphisms-lemma-composition-finite-presentation}、
\ref{morphisms-lemma-base-change-finite-presentation}、
\ref{morphisms-lemma-composition-flat} 和
\ref{morphisms-lemma-base-change-flat}。
\end{proof}

\noindent
下面的引理本可用来以更统一的方式证明上一引理的结果。

\begin{lemma}
\label{more-groupoids-lemma-restrict-property}
设 $S$ 为概形，$(U, R, s, t, c)$ 为 $S$ 上的群胚概形，
$g : U' \to U$ 为概形态射。设
$(U', R', s', t', c')$ 为 $(U, R, s, t, c)$ 经 $g$ 的限制，并令
$h = s \circ \text{pr}_1 : U' \times_{g, U, t} R \to U$。
若 $\mathcal{P}$ 是概形态射的一项性质，且
\begin{enumerate}
\item $h$ 具有性质 $\mathcal{P}$；并且
\item $\mathcal{P}$ 在基变换下保持，
\end{enumerate}
则 $s', t'$ 具有性质 $\mathcal{P}$。
\end{lemma}

\begin{proof}
由图 (\ref{more-groupoids-equation-restriction})，$s'$ 是 $h$ 的基变换；
而 $t'$ 作为概形态射同构于 $s'$，故结论显然。
\end{proof}

\begin{lemma}
\label{more-groupoids-lemma-double-restrict}
设 $S$ 为概形，$(U, R, s, t, c)$ 为 $S$ 上的群胚概形。
设 $g : U' \to U$ 与 $g' : U'' \to U'$ 为概形态射，
并令 $g'' = g \circ g'$。设 $(U', R', s', t', c')$
为 $R$ 在 $U'$ 上的限制。令
$h = s \circ \text{pr}_1 : U' \times_{g, U, t} R \to U$、
$h' = s' \circ \text{pr}_1 : U'' \times_{g', U', t} R \to U'$，以及
$h'' = s \circ \text{pr}_1 : U'' \times_{g'', U, t} R \to U$。
下图交换：
$$
\xymatrix{
U'' \times_{g', U', t} R' \ar[d]^{h'} &
(U' \times_{g, U, t} R) \times_U (U'' \times_{g'', U, t} R)
\ar[l] \ar[r] \ar[d] &
U'' \times_{g'', U, t} R \ar[d]_{h''} \\
U' &
U' \times_{g, U, t} R \ar[l]_{\text{pr}_0} \ar[r]^h &
U
}
$$
其中两个方块均为笛卡尔方块，而左上方的水平箭头由规则
$$
\begin{matrix}
(U' \times_{g, U, t} R) \times_U (U'' \times_{g'', U, t} R) &
\longrightarrow &
U'' \times_{g', U', t} R' \\
((u', r_0), (u'', r_1)) &
\longmapsto &
(u'', (c(r_1, i(r_0)), (g'(u''), u')))
\end{matrix}
$$
给出；记号在证明中解释。
\end{lemma}

\begin{proof}
我们利用函子观点作出具体计算，把引理归约为群胚范畴限制中
关于箭头的断言。在引理最后的公式中，记号
$((u', r_0), (u'', r_1))$ 表示
$(U' \times_{g, U, t} R) \times_U (U'' \times_{g'', U, t} R)$
的一个 $T$-值点。这意味着 $u', u'', r_0, r_1$ 分别是
$U', U'', R, R$ 的 $T$-值点，并且
$g(u') = t(r_0)$、$g(g'(u'')) = g''(u'') = t(r_1)$，以及
$s(r_0) = s(r_1)$。严格地说，这里应写
$g \circ u' = t \circ r_0$ 等式，但那会使记号更难阅读。
若把 $r_1$ 与 $r_0$ 看作群胚范畴中的箭头，则可用下图表示：
$$
\xymatrix{
t(r_0) = g(u') &
s(r_0) = s(r_1) \ar[l]_{r_0} \ar[r]^-{r_1} &
t(r_1) = g(g'(u''))
}
$$
该图尤其表明复合 $c(r_1, i(r_0))$ 有定义。回忆
$$
R' = R \times_{(t, s), U \times_S U, g \times g} U' \times_S U'
$$
因此，$R'$ 的一个 $T$-值点形如 $(r, (u'_0, u'_1))$，
且 $t(r) = g(u'_0)$、$s(r) = g(u'_1)$。特别地，
给定如上的 $((u', r_0), (u'', r_1))$，便得到 $R'$ 的
$T$-值点 $(c(r_1, i(r_0)), (g'(u''), u'))$，因为
$t(c(r_1, i(r_0))) = t(r_1) = g(g'(u''))$ 且
$s(c(r_1, i(r_0))) = s(i(r_0)) = t(r_0) = g(u')$。
由读者按此定义验证左侧方块交换。

\medskip\noindent
为证明左侧方块是笛卡尔的，设给定 $(v'', p')$ 与 $(v', p)$，
它们分别是 $U'' \times_{g', U', t} R'$ 与
$U' \times_{g, U, t} R$ 的 $T$-值点，并满足
$v' = s'(p')$。这还意味着 $g'(v'') = t'(p')$
且 $g(v') = t(p)$。由以上讨论，可写
$p' = (r, (u_0', u_1'))$，其中 $t(r) = g(u'_0)$ 且
$s(r) = g(u'_1)$。因而
$v' = s'(p') = u_1'$ 且 $g'(v'') = t'(p') = u_0'$。
图示如下：
$$
\xymatrix{
s(p) \ar[r]^-p &
g(v') = g(u'_1) \ar[r]^-r &
g(u'_0) = g(g'(v''))
}
$$
我们必须证明存在唯一的上述 $T$-值点
$((u', r_0), (u'', r_1))$，使
$v' = u'$、$p = r_0$、$v'' = u''$ 且
$p' = (c(r_1, i(r_0)), (g'(u''), u'))$。
比较上面的两个图可知，只能取
$$
((u', r_0), (u'', r_1)) = ((v', p), (v'', c(r, p))
$$
其余细节从略。
\end{proof}

\begin{lemma}
\label{more-groupoids-lemma-double-restrict-property}
设 $S$ 为概形，$(U, R, s, t, c)$ 为 $S$ 上的群胚概形。
设 $g : U' \to U$ 与 $g' : U'' \to U'$ 为概形态射，
令 $g'' = g \circ g'$，并设 $(U', R', s', t', c')$
为 $R$ 在 $U'$ 上的限制。令
$h = s \circ \text{pr}_1 : U' \times_{g, U, t} R \to U$、
$h' = s' \circ \text{pr}_1 : U'' \times_{g', U', t} R \to U'$，以及
$h'' = s \circ \text{pr}_1 : U'' \times_{g'', U, t} R \to U$。
令 $\tau \in \{Zariski, \linebreak[0] \etale, \linebreak[0]
smooth, \linebreak[0] syntomic, \linebreak[0] fppf, \linebreak[0] fpqc\}$。
设 $\mathcal{P}$ 为概形态射的一项性质，它在基变换下保持，
并且在目标上关于 $\tau$-拓扑为局部的。若
\begin{enumerate}
\item $h(U' \times_U R)$ 在 $U$ 中为开集；
\item $\{h : U' \times_U R \to h(U' \times_U R)\}$
是一个 $\tau$-覆盖；
\item $h'$ 具有性质 $\mathcal{P}$，
\end{enumerate}
则 $h''$ 具有性质 $\mathcal{P}$。反之，若
\begin{enumerate}
\item[(a)] $\{t : R \to U\}$ 是一个 $\tau$-覆盖；
\item[(d)] $h''$ 具有性质 $\mathcal{P}$，
\end{enumerate}
则 $h'$ 具有性质 $\mathcal{P}$。
\end{lemma}

\begin{proof}
这由引理 \ref{more-groupoids-lemma-double-restrict} 中图的性质形式地得出。
在第一种情形中，由于 $g'' = g \circ g'$，态射 $h''$ 的像
包含于 $h$ 的像。因此，可把图右下角的 $U$ 替换为
$h(U' \times_U R)$；这解释了引理中条件 (1) 与 (2) 的作用。
在第二种情形中，注意
$\{\text{pr}_0 : U' \times_{g, U, t} R \to U'\}$
作为 $\tau$ 的基变换，并由条件 (a)，是一个 $\tau$-覆盖。
\end{proof}






\section{域上群胚的性质}
\label{more-groupoids-section-properties-groupoids-on-fields}

\noindent
“域上的群胚”是指群胚概形 $(U, R, s, t, c)$，其中
$U$ 是某个域的谱。它{\bf 并不}表示
$(U, R, s, t, c)$ 定义在某个域上；更确切地说，
它{\bf 并不}表示态射 $s, t : R \to U$ 相等。
给定任意域 $k$、抽象群 $G$ 以及群同态
$\varphi : G \to \text{Aut}(k)$，令
\begin{align*}
U & = \Spec(k) \\
R & = \coprod\nolimits_{g \in G} \Spec(k) \\
s & = \coprod\nolimits_{g \in G} \Spec(\text{id}_k) \\
t & = \coprod\nolimits_{g \in G} \Spec(\varphi(g)) \\
c & = \text{下列范畴中的复合： }G
\end{align*}
便得到 $\mathbf{Z}$ 上的群胚概形 $(U, R, s, t, c)$。
这个例子仍是 $\Spec(k^G)$ 上的群胚概形。因此，若 $G$ 有限，
则 $U = \Spec(k)$ 在 $\Spec(k^G)$ 上有限。
在某种意义上，本节的目标是证明：对 $s, t$ 施加适当的有限性条件，
会迫使任意域上群胚定义在一个有限指数子域 $k' \subset k$ 上。

\medskip\noindent
若 $k$ 为域，$(G, m)$ 是 $k$ 上的群概形，结构态射为
$p : G \to \Spec(k)$，则 $(\Spec(k), G, p, p, m)$
是域上群胚的一个例子（在这种情形中，整个结构当然定义在域上）。
因此，本节可视为《群胚概形》第 \ref{groupoids-section-properties-group-schemes-field} 节的对应版本。

\begin{lemma}
\label{more-groupoids-lemma-groupoid-on-field-open-multiplication}
设 $S$ 为概形，$(U, R, s, t, c)$ 为 $S$ 上的群胚概形。
若 $U$ 是某个域的谱，则合成态射
$c : R \times_{s, U, t} R \to R$ 是开映射。
\end{lemma}

\begin{proof}
由图 (\ref{more-groupoids-equation-pull})，合成态射同构于投影
$\text{pr}_1 : R \times_{t, U, t} R \to R$。
该投影是开映射，见《概形态射》引理 \ref{morphisms-lemma-scheme-over-field-universally-open}。
\end{proof}

\begin{lemma}
\label{more-groupoids-lemma-groupoid-on-field-separated}
设 $S$ 为概形，$(U, R, s, t, c)$ 为 $S$ 上的群胚概形。
若 $U$ 是某个域的谱，则 $R$ 是分离概形。
\end{lemma}

\begin{proof}
由《群胚概形》引理 \ref{groupoids-lemma-group-scheme-over-field-separated}，
稳定子群概形 $G \to U$ 分离。由《群胚概形》引理 \ref{groupoids-lemma-diagonal}，态射
$j = (t, s) : R \to U \times_S U$ 分离。
由于 $U$ 是域的谱，概形 $U \times_S U$ 仿射
（见《概形》第 \ref{schemes-section-fibre-products} 节
中纤维积的构造）。因此 $R$ 分离，见《概形》引理 \ref{schemes-lemma-separated-permanence}。
\end{proof}

\begin{lemma}
\label{more-groupoids-lemma-groupoid-on-field-homogeneous}
设 $S$ 为概形，$(U, R, s, t, c)$ 为 $S$ 上的群胚概形。
假设 $U = \Spec(k)$，其中 $k$ 为域。对任意点
$r, r' \in R$，存在域扩张 $k'/k$、点
$r_1, r_2 \in R \times_{s, \Spec(k)} \Spec(k')$ 以及图
$$
\xymatrix{
R &
R \times_{s, \Spec(k)} \Spec(k')
\ar[l]_-{\text{pr}_0} \ar[r]^\varphi &
R \times_{s, \Spec(k)} \Spec(k')
\ar[r]^-{\text{pr}_0} &
R
}
$$
使得 $\varphi$ 是 $\Spec(k')$ 上的概形同构，并且
$\varphi(r_1) = r_2$、$\text{pr}_0(r_1) = r$、以及
$\text{pr}_0(r_2) = r'$。
\end{lemma}

\begin{proof}
这是引理 \ref{more-groupoids-lemma-two-fibres} 的 (1)、(2) 两项的特殊情形。
\end{proof}

\begin{lemma}
\label{more-groupoids-lemma-restrict-groupoid-on-field}
设 $S$ 为概形，$(U, R, s, t, c)$ 为 $S$ 上的群胚概形。
假设 $U = \Spec(k)$，其中 $k$ 为域。设 $k'/k$ 为域扩张，
$U' = \Spec(k')$，并设 $(U', R', s', t', c')$
为 $(U, R, s, t, c)$ 经 $U' \to U$ 的限制。在定义图
$$
\xymatrix{
R' \ar[d] \ar[r] \ar@/_3pc/[dd]_{t'} \ar@/^1pc/[rr]^{s'} \ar@{..>}[rd] &
R \times_{s, U} U' \ar[r] \ar[d] &
U' \ar[d] \\
U' \times_{U, t} R \ar[d] \ar[r] &
R \ar[r]^s \ar[d]_t &
U \\
U' \ar[r] &
U
}
$$
中，所有态射均满、平坦且普遍开；虚线箭头
$R' \to R$ 还仿射。
\end{lemma}

\begin{proof}
态射 $U' \to U$ 等于 $\Spec(k') \to \Spec(k)$，
因而仿射、满且平坦。由《概形态射》引理 \ref{morphisms-lemma-scheme-over-field-universally-open}，
态射 $s, t : R \to U$ 与 $U' \to U$ 均普遍开。
由于 $R$ 非空且 $U$ 是域的谱，态射 $s, t : R \to U$
满且平坦。最后应用《概形态射》引理 \ref{morphisms-lemma-base-change-surjective}、
\ref{morphisms-lemma-composition-surjective}、
\ref{morphisms-lemma-composition-open}、
\ref{morphisms-lemma-base-change-affine}、
\ref{morphisms-lemma-composition-affine}、
\ref{morphisms-lemma-base-change-flat} 和
\ref{morphisms-lemma-composition-flat} 即得结论。
\end{proof}

\begin{lemma}
\label{more-groupoids-lemma-groupoid-on-field-explain-points}
设 $S$ 为概形，$(U, R, s, t, c)$ 为 $S$ 上的群胚概形。
假设 $U = \Spec(k)$，其中 $k$ 为域。
对任意点 $r \in R$，存在
\begin{enumerate}
\item 域扩张 $k'/k$，其中 $k'$ 代数闭；
\item 点 $r' \in R'$，其中 $(U', R', s', t', c')$ 是
$(U, R, s, t, c)$ 经 $\Spec(k') \to \Spec(k)$ 的限制，
\end{enumerate}
使得
\begin{enumerate}
\item 点 $r'$ 在态射 $R' \to R$ 下映到 $r$；并且
\item 映射 $s', t' : R' \to \Spec(k')$ 诱导同构
$k' \to \kappa(r')$。
\end{enumerate}
\end{lemma}

\begin{proof}
将这个几何陈述转写为域上的陈述，就是说我们可以找到图
$$
\xymatrix{
k' & k' \ar[l]^1 & \\
k' \ar[u]^\tau & \kappa(r) \ar[lu]^\sigma & k \ar[l]^-s \ar[lu]_i \\
& k \ar[lu]^i \ar[u]_t
}
$$
其中 $i : k \to k'$ 是 $k$ 到 $k'$ 的嵌入，
映射 $s, t : k \to \kappa(r)$ 由 $s, t : R \to U$ 诱导，
而映射 $\tau : k' \to k'$ 是一个自同构。为构造这样的图，
可以按如下步骤进行：
\begin{enumerate}
\item 选取域映射 $i : k \to k'$，其中 $k'$ 代数闭，
且在 $k$ 上的超越次数非常大。
\item 选取嵌入 $\sigma : \kappa(r) \to k'$，使得
$\sigma \circ s = i$。这样的 $\sigma$ 确实存在：
只需选取 $\kappa(r)$ 在 $k$ 上的一组超越基
$\{x_\alpha\}_{\alpha \in A}$，并找到在 $i(k)$ 上代数无关的
$y_\alpha \in k'$（$\alpha \in A$），再按如下规则将
$s(k)(\{x_\alpha\})$ 映入 $k'$：
对 $\lambda \in k$ 令 $s(\lambda) \mapsto i(\lambda)$，
并对 $\alpha \in A$ 令 $x_\alpha \mapsto y_\alpha$。
然后利用 $k'$ 代数闭，将其延拓为
$\tau : \kappa(\alpha) \to k'$。
\item 选取自同构 $\tau : k' \to k'$，使得
$\tau \circ i = \sigma \circ t$。为此，选取 $k$ 在其素域上的
一组超越基 $\{x_\alpha\}_{\alpha \in A}$。
一方面，向 $\{i(x_\alpha)\}$ 添加
$\{y_\beta\}_{\beta \in B}$，将其扩充为 $k'$ 的一组超越基；
另一方面，向 $\{\sigma(t(x_\alpha))\}$ 添加
$\{z_\gamma\}_{\gamma \in C}$，将其扩充为 $k'$ 的一组超越基。
由于 $k'$ 代数闭，我们可以把同构
$\sigma \circ t \circ i^{-1} : i(k) \to \sigma(t(k))$
延拓为它们在 $k'$ 中的代数闭包之间的同构
$\tau' : \overline{i(k)} \to \overline{\sigma(t(k))}$。
因为 $k'$ 的超越次数很大，可见集合 $B$ 和 $C$ 具有相同基数。
因此可利用一个双射 $B \to C$，将 $\tau'$ 延拓为同构
$$
\overline{i(k)}(\{y_\beta\})
\longrightarrow
\overline{\sigma(t(k))}(\{z_\gamma\})
$$
又因为 $k'$ 是等号两边的代数闭包，可见该同构进一步延拓为
自同构 $\tau : k' \to k'$，正合所需。
\end{enumerate}
这就证明了引理。
\end{proof}

\begin{lemma}
\label{more-groupoids-lemma-groupoid-on-field-move-point}
设 $S$ 为概形，$(U, R, s, t, c)$ 为 $S$ 上的群胚概形。
假设 $U = \Spec(k)$，其中 $k$ 为域。
若 $r \in R$ 是一点，并且 $s, t$ 均诱导同构
$k \to \kappa(r)$，则映射
$$
R \longrightarrow R, \quad
x \longmapsto c(r, x)
$$
（精确记号见证明）是一个将 $e$ 映到 $r$ 的自同构 $R \to R$。
\end{lemma}

\begin{proof}
若从函子角度思考群胚，这一点完全显然。不过我们也把细节完整写出。
记 $a : U \to R$ 为像为 $r$ 且满足
$s \circ a = \text{id}_U$ 的态射；这样的态射由
$s : k \to \kappa(r)$ 是同构这一假设而存在。
类似地，记 $b : U \to R$ 为像为 $r$ 且满足
$t \circ b = \text{id}_U$ 的态射。注意
$b = a \circ (t \circ a)^{-1}$，特别地
$a \circ s \circ b = b$。

\medskip\noindent
考虑态射 $\Psi : R \to R$，它在 $T$ 值点上由
$$
(f : T \to R) \longmapsto (c(a \circ t \circ f, f) : T \to R)
$$
给出。要看它确有定义，需要验证
$s \circ a \circ t \circ f = t \circ f$；这由 $s \circ a = 1$
显然成立。注意 $\Phi(e) = a$，因此要证明该引理，只需证明
$\Phi$ 是 $R$ 的自同构。令 $\Phi : R \to R$ 为在 $T$ 值点上由
$$
(g : T \to R) \longmapsto (c(i \circ b \circ t \circ g, g) : T \to R).
$$
给出的态射。它确有定义，因为
$s \circ i \circ b \circ t \circ g = t \circ b \circ t \circ g =
t \circ g$。我们断言 $\Phi$ 与 $\Psi$ 互为逆。为此计算
\begin{align*}
& c(a \circ t \circ c(i \circ b \circ t \circ g, g),
c(i \circ b \circ t \circ g, g)) \\
& =
c(a \circ t \circ i \circ b \circ t \circ g,
c(i \circ b \circ t \circ g, g)) \\
& =
c(a \circ s \circ b \circ t \circ g,
c(i \circ b \circ t \circ g, g)) \\
& =
c(b \circ t \circ g, c(i \circ b \circ t \circ g, g)) \\
& =
c(c(b \circ t \circ g, i \circ b \circ t \circ g), g)) \\
& =
c(e, g) \\
& = g
\end{align*}
其中使用了上面得到的关系 $a \circ s \circ b = b$。
反过来，
\begin{align*}
& c(i \circ b \circ t \circ c(a \circ t \circ f, f), c(a \circ t \circ f, f)) \\
& =
c(i \circ b \circ t \circ a \circ t \circ f, c(a \circ t \circ f, f)) \\
& =
c(i \circ a \circ (t \circ a)^{-1} \circ t \circ a \circ t \circ f,
c(a \circ t \circ f, f)) \\
& =
c(i \circ a \circ t \circ f, c(a \circ t \circ f, f)) \\
& =
c(c(i \circ a \circ t \circ f, a \circ t \circ f), f) \\
& =
c(e, f) \\
& = f
\end{align*}
引理得证。
\end{proof}

\begin{lemma}
\label{more-groupoids-lemma-groupoid-on-field-translate-open}
设 $S$ 为概形，$(U, R, s, t, c)$ 为 $S$ 上的群胚概形。
若 $U$ 是一个域的谱，$W \subset R$ 是开集，
且 $Z \to R$ 是概形态射，则复合
$Z \times_{s, U, t} W \to R \times_{s, U, t} R \to R$
的像是开集。
\end{lemma}

\begin{proof}
写成 $U = \Spec(k)$。考虑域扩张 $k'/k$，并记
$U' = \Spec(k')$。设 $R'$ 为 $R$ 经 $U' \to U$ 的限制。
令 $Z' = Z \times_R R'$ 及 $W' = R' \times_R W$。
考虑 $Z \times_{s, U, t} W$ 的一点 $\xi = (z, w)$。
令 $r \in R$ 为 $z$ 在 $Z \to R$ 下的像。
按引理 \ref{more-groupoids-lemma-groupoid-on-field-explain-points} 选取
$k' \supset k$ 及 $r' \in R'$。可以选取映到 $z$ 和 $r'$ 的
$z' \in Z'$。于是可以找到
$\xi' \in Z' \times_{s', U', t'} W'$，它映到 $z'$ 和 $\xi$。
开集 $c(r', W')$（引理
\ref{more-groupoids-lemma-groupoid-on-field-move-point}）包含在
$Z' \times_{s', U', t'} W' \to R'$ 的像中。注意
$Z' \times_{s', U', t'} W' = (Z \times_{s, U, t} W)
\times_{R \times_{s, U, t} R} (R' \times_{s', U', t'} R')$。
因此 $Z' \times_{s', U', t'} W' \to R' \to R$ 的像包含在
$Z \times_{s, U, t} W \to R$ 的像中。由于 $R' \to R$
是开映射（引理 \ref{more-groupoids-lemma-restrict-groupoid-on-field}），
可知该像包含 $\xi$ 的像的一个开邻域，正合所需。
\end{proof}

\begin{lemma}
\label{more-groupoids-lemma-groupoid-on-field-geometrically-irreducible}
设 $S$ 为概形，$(U, R, s, t, c)$ 为 $S$ 上的群胚概形。
假设 $U = \Spec(k)$，其中 $k$ 为域。滥用记号，以 $e \in R$
表示恒等态射 $e : U \to R$ 的像。则
\begin{enumerate}
\item $R$ 的每个局部环 $\mathcal{O}_{R, r}$ 都有唯一的极小素理想；
\item 恰有一个经过 $e$ 的 $R$ 的不可约分支 $Z$；并且
\item 无论经由 $s$ 还是 $t$，$Z$ 在 $k$ 上都几何不可约。
\end{enumerate}
\end{lemma}

\begin{proof}
设 $r \in R$ 为一点。在本证明中，我们将直接使用经过点 $r$
的 $R$ 的不可约分支与局部环 $\mathcal{O}_{R, r}$ 的极小素理想
之间的对应，而不再另行说明。按引理
\ref{more-groupoids-lemma-groupoid-on-field-explain-points} 选取
$k \subset k'$ 和 $r' \in R'$。注意
$\mathcal{O}_{R, r} \to \mathcal{O}_{R', r'}$ 忠实平坦且为局部同态，
见引理 \ref{more-groupoids-lemma-restrict-groupoid-on-field}。
因此对 $r' \in R'$ 的结论蕴含对 $r \in R$ 的结论。
换言之，可以假设 $s, t : k \to \kappa(r)$ 是同构。
由引理 \ref{more-groupoids-lemma-groupoid-on-field-move-point}，存在一个把 $e$
移到 $r$ 的自同构。因此可以假设 $r = e$，也就是说，
(1) 由 (2) 推出。

\medskip\noindent
先在 $k$ 可分代数闭的情形证明 (2)。设 $X, Y \subset R$
是经过 $e$ 的不可约分支。由《代数簇》引理 \ref{varieties-lemma-bijection-irreducible-components} 和
\ref{varieties-lemma-separably-closed-irreducible}，
概形 $X \times_{s, U, t} Y$ 也不可约。
因此 $c(X \times_{s, U, t} Y) \subset R$ 是不可约子集。
我们断言它同时包含 $X$ 和 $Y$（作为 $R$ 的子集）。
确实，设 $T$ 为一个域的谱。若 $x : T \to X$ 是 $X$ 的一个
$T$ 值点，则 $c(x, e \circ s \circ x) = x$；又因为 $e \in Y$，
故 $e \circ s \circ x$ 经过 $Y$ 分解。对 $Y$ 的点也类似。
这显然蕴含 $X = Y$，即恰有一个经过 $e$ 的 $R$ 的不可约分支。

\medskip\noindent
现在一般地证明 (2) 和 (3)。设 $k \subset k'$ 为一个可分代数闭包，
并令 $(U', R', s', t', c')$ 为 $(U, R, s, t, c)$
经 $\Spec(k') \to \Spec(k)$ 的限制。由上一段，恰有一个经过 $e'$
的 $R'$ 的不可约分支 $Z'$。记 $e'' \in R \times_{s, U} U'$
为 $e$ 的基变换。由于
$R' \to R \times_{s, U} U'$ 忠实平坦，见引理
\ref{more-groupoids-lemma-restrict-groupoid-on-field}，且 $e' \mapsto e''$，
可见恰有一个经过 $e''$ 的
$R \times_{s, k} k'$ 的不可约分支 $Z''$。
又因 $R \times_k k' \to R$ 忠实平坦，这蕴含恰有一个经过 $e$
的 $R$ 的不可约分支 $Z$。这证明了 (2)。

\medskip\noindent
为证明 (3)，设 $Z''' \subset R \times_k k'$ 是
$Z \times_k k'$ 的任意不可约分支。由《代数簇》引理 \ref{varieties-lemma-orbit-irreducible-components}，
对某个 $\sigma \in \text{Gal}(k'/k)$ 有
$Z''' = \sigma(Z'')$。由于 $\sigma(e'') = e''$，
可见 $e'' \in Z'''$，从而 $Z''' = Z''$。
这说明经由态射 $s$，$Z$ 在 $\Spec(k)$ 上几何不可约。
同样的论证说明，经由态射 $t$，$Z$ 在 $\Spec(k)$ 上也几何不可约。
\end{proof}

\begin{lemma}
\label{more-groupoids-lemma-groupoid-on-field-locally-finite-type-dimension}
设 $S$ 为概形，$(U, R, s, t, c)$ 为 $S$ 上的群胚概形。
假设 $U = \Spec(k)$，其中 $k$ 为域，并假设 $s, t$ 局部有限型。则
\begin{enumerate}
\item $R$ 等维；
\item 对所有 $r \in R$，有 $\dim(R) = \dim_r(R)$；
\item 对任意 $r \in R$，有
$\text{trdeg}_{s(k)}(\kappa(r)) = \text{trdeg}_{t(k)}(\kappa(r))$；并且
\item 对任意闭点 $r \in R$，有
$\dim(R) = \dim(\mathcal{O}_{R, r})$。
\end{enumerate}
\end{lemma}

\begin{proof}
设 $r, r' \in R$。由引理
\ref{more-groupoids-lemma-groupoid-on-field-homogeneous} 及
《概形态射》引理 \ref{morphisms-lemma-dimension-fibre-after-base-change}，
有 $\dim_r(R) = \dim_{r'}(R)$。由《概形态射》引理 \ref{morphisms-lemma-dimension-fibre-at-a-point}，
$$
\dim_r(R) =
\dim(\mathcal{O}_{R, r}) + \text{trdeg}_{s(k)}(\kappa(r)) =
\dim(\mathcal{O}_{R, r}) + \text{trdeg}_{t(k)}(\kappa(r)).
$$
另一方面，$R$（或 $R$ 的任意开子集）的维数是 $R$ 的局部环维数的
上确界，见《概形的性质》引理 \ref{properties-lemma-codimension-local-ring}。
显然，该值在闭点 $r$ 处达到最大；在这种情形，
$\text{trdeg}_k(\kappa(r)) = 0$（由希尔伯特零点定理，见
《概形态射》第 \ref{morphisms-section-points-finite-type} 节）。
引理由此得证。
\end{proof}

\begin{lemma}
\label{more-groupoids-lemma-groupoid-on-field-dimension-equal-stabilizer}
设 $S$ 为概形，$(U, R, s, t, c)$ 为 $S$ 上的群胚概形。
假设 $U = \Spec(k)$，其中 $k$ 为域，并假设 $s, t$ 局部有限型。
则 $\dim(R) = \dim(G)$，其中 $G$ 是 $R$ 的稳定子群概形。
\end{lemma}

\begin{proof}
设 $Z \subset R$ 为经过 $e$ 的不可约分支（见引理
\ref{more-groupoids-lemma-groupoid-on-field-geometrically-irreducible}），
并把它视为 $R$ 的一个整闭子概形。令 $k'_s$，相应地，\ $k'_t$
为 $s(k)$，相应地，\ $t(k)$ 在 $\Gamma(Z, \mathcal{O}_Z)$ 中的整闭包。
回忆 $k'_s$ 和 $k'_t$ 都是域，见《代数簇》引理 \ref{varieties-lemma-integral-closure-ground-field}。
由《代数簇》命题 \ref{varieties-proposition-unique-base-field}，
作为 $\Gamma(Z, \mathcal{O}_Z)$ 的子环有 $k'_s = k'_t$。
由于 $e$ 经过 $Z$ 分解，得到交换图
$$
\xymatrix{
k \ar[rd]_t \ar[rrd]^1 \\
& \Gamma(Z, \mathcal{O}_Z) \ar[r]^e & k \\
k \ar[ru]^s \ar[rru]_1
}
$$
一方面，这说明 $k'_s = s(k)$、$k'_t = t(k)$，故
$s(k) = t(k)$；结合上图，这竟蕴含 $s = t$！
换言之，我们得出 $Z$ 是
$G = R \times_{(t, s), U \times_S U, \Delta} U$ 的闭子概形。
由于 $G$ 和 $R$ 都等维，引理随即成立；见引理
\ref{more-groupoids-lemma-groupoid-on-field-locally-finite-type-dimension} 及
《群胚概形》引理 \ref{groupoids-lemma-group-scheme-finite-type-field}。
\end{proof}

\begin{remark}
\label{more-groupoids-remark-warn-dimension-groupoid-on-field}
警告：若不要求 $s$ 和 $t$ 局部有限型，引理
\ref{more-groupoids-lemma-groupoid-on-field-dimension-equal-stabilizer}
便是错误的。一个简单例子是从作用
$$
\mathbf{G}_{m, \mathbf{Q}} \times_{\mathbf{Q}} \mathbf{A}^1_{\mathbf{Q}}
\to \mathbf{A}^1_{\mathbf{Q}}
$$
出发，把相应的群胚概形限制到
$\mathbf{A}^1_{\mathbf{Q}}$ 的一般点上。换言之，经由态射
$\Spec(\mathbf{Q}(x)) \to
\Spec(\mathbf{Q}[x]) = \mathbf{A}^1_{\mathbf{Q}}$
作限制。于是得到群胚概形 $(U, R, s, t, c)$，其中
$U = \Spec(\mathbf{Q}(x))$ 且
$$
R = \Spec\left(
\mathbf{Q}(x)[y]\left[
\frac{1}{P(xy)}, P \in \mathbf{Q}[T], P \not = 0
\right]
\right)
$$
。
在这种情形，$\dim(R) = 1$ 而 $\dim(G) = 0$。
\end{remark}

\begin{lemma}
\label{more-groupoids-lemma-groupoid-characteristic-zero-smooth}
设 $S$ 为概形，$(U, R, s, t, c)$ 为 $S$ 上的群胚概形。假设
\begin{enumerate}
\item $U = \Spec(k)$，其中 $k$ 为域；
\item $s, t$ 局部有限型；并且
\item $k$ 的特征为零。
\end{enumerate}
则 $s, t : R \to U$ 光滑。
\end{lemma}

\begin{proof}
由引理 \ref{more-groupoids-lemma-sheaf-differentials}，
$R \to U$ 的微分层自由。因此光滑性由《代数簇》引理 \ref{varieties-lemma-char-zero-differentials-free-smooth} 得出。
\end{proof}

\begin{lemma}
\label{more-groupoids-lemma-reduced-group-scheme-perfect-field-characteristic-p-smooth}
设 $S$ 为概形，$(U, R, s, t, c)$ 为 $S$ 上的群胚概形。假设
\begin{enumerate}
\item $U = \Spec(k)$，其中 $k$ 为域；
\item $s, t$ 局部有限型；
\item $R$ 既约；并且
\item $k$ 完美。
\end{enumerate}
则 $s, t : R \to U$ 光滑。
\end{lemma}

\begin{proof}
由引理 \ref{more-groupoids-lemma-sheaf-differentials}，
层 $\Omega_{R/U}$ 自由。因此引理由《代数簇》引理 \ref{varieties-lemma-char-p-differentials-free-smooth} 得出。
\end{proof}










\section{域上群胚的态射}
\label{more-groupoids-section-morphisms-groupoids-on-fields}

\noindent
本节研究域上群胚之间的态射。这比研究域上的群概形的态射稍为一般，
但二者非常相近。

\begin{situation}
\label{more-groupoids-situation-morphism-groupoids-on-field}
设 $S$ 为概形。设 $U = \Spec(k)$ 是 $S$ 上的概形，其中 $k$ 为域。
设 $(U, R_1, s_1, t_1, c_1)$、
$(U, R_2, s_2, t_2, c_2)$ 为 $S$ 上具有相同第一分量的群胚概形。
设 $a : R_1 \to R_2$ 为一个态射，使得 $(\text{id}_U, a)$
定义 $S$ 上群胚概形的态射，见《群胚概形》定义 \ref{groupoids-definition-groupoid}。特别地，下列图交换：
$$
\vcenter{
\xymatrix{
R_1 \ar[rrd]^{t_1} \ar[rdd]_{s_1} \ar[rd]_a \\
& R_2 \ar[d]^{t_2} \ar[r]_{s_2} & U \\
& U
}
}
\quad\quad
\vcenter{
\xymatrix{
R_1 \times_{s_1, U, t_1} R_1 \ar[r]_-{c_1} \ar[d]_{a \times a} &
R_1 \ar[d]^a \\
R_2 \times_{s_2, U, t_2} R_2 \ar[r]^-{c_2} &
R_2
}
}
$$
\end{situation}

\noindent
下一个引理是《群胚概形》引理 \ref{groupoids-lemma-open-subgroup-closed-over-field} 的推广。

\begin{lemma}
\label{more-groupoids-lemma-open-image-is-closed}
记号和假设同情形 \ref{more-groupoids-situation-morphism-groupoids-on-field}。
若 $a(R_1)$ 在 $R_2$ 中开，则 $a(R_1)$ 在 $R_2$ 中闭。
\end{lemma}

\begin{proof}
设 $r_2 \in R_2$ 为 $a(R_1)$ 的闭包中的一点。我们要证明
$r_2 \in a(R_1)$。按引理
\ref{more-groupoids-lemma-groupoid-on-field-explain-points}，选取适配于
$(U, R_2, s_2, t_2, c_2)$ 和 $r_2$ 的
$k \subset k'$ 与 $r_2' \in R'_2$。
设 $R_i'$ 为 $R_i$ 经态射
$U' = \Spec(k') \to U = \Spec(k)$ 的限制。
设 $a' : R'_1 \to R_2'$ 为 $a$ 的基变换。图
$$
\xymatrix{
R'_1 \ar[r]_{a'} \ar[d]_{p_1} & R'_2 \ar[d]^{p_2} \\
R_1 \ar[r]^a & R_2
}
$$
是纤维方形。因此 $a'$ 的像是 $a$ 的像在态射
$p_2 : R'_2 \to R_2$ 下的逆像。由引理
\ref{more-groupoids-lemma-restrict-groupoid-on-field}，映射 $p_2$ 满且开。
所以由《拓扑学》引理 \ref{topology-lemma-open-morphism-quotient-topology}，
可见 $r_2'$ 属于 $a'(R'_1)$ 的闭包。
这意味着可以假设 $r_2 \in R_2$ 具有如下性质：
$s_2$ 和 $t_2$ 诱导的映射
$k \to \kappa(r_2)$ 都是同构。

\medskip\noindent
在这种情形可以使用引理
\ref{more-groupoids-lemma-groupoid-on-field-move-point}。
该引理说明 $c(r_2, a(R_1))$ 是 $r_2$ 的一个开邻域。
由于假设 $r_2$ 属于 $a(R_1)$ 的闭包，故
$a(R_1) \cap c(r_2, a(R_1)) \not = \emptyset$。
利用 $R_2$ 和 $R_1$ 中的逆，可见这意味着
$c_2(a(R_1), a(R_1))$ 包含 $r_2$。
而 $c_2(a(R_1), a(R_1)) \subset a(c_1(R_1, R_1)) = a(R_1)$，
故如欲证地得到 $r_2 \in a(R_1)$。
\end{proof}

\begin{lemma}
\label{more-groupoids-lemma-map-groupoids-on-field-image}
记号和假设同情形 \ref{more-groupoids-situation-morphism-groupoids-on-field}。
设 $Z \subset R_2$ 为约化闭子概形（见《概形》定义 \ref{schemes-definition-reduced-induced-scheme}），其底拓扑空间是
$a : R_1 \to R_2$ 的像的闭包。则在集合论意义下，
$c_2(Z \times_{s_2, U, t_2} Z) \subset Z$。
\end{lemma}

\begin{proof}
考虑交换图
$$
\xymatrix{
R_1 \times_{s_1, U, t_1} R_1 \ar[r] \ar[d] & R_1 \ar[d] \\
R_2 \times_{s_2, U, t_2} R_2 \ar[r] & R_2
}
$$
由《代数簇》引理 \ref{varieties-lemma-closure-image-product-map}，
左侧竖直箭头之像的闭包在集合论意义下是
$Z \times_{s_2, U, t_2} Z$。结论随即成立。
\end{proof}

\begin{lemma}
\label{more-groupoids-lemma-map-groupoids-on-perfect-field-image}
记号和假设同情形 \ref{more-groupoids-situation-morphism-groupoids-on-field}。
假设 $k$ 完美。设 $Z \subset R_2$ 为约化闭子概形（见
《概形》定义 \ref{schemes-definition-reduced-induced-scheme}），
其底拓扑空间是 $a : R_1 \to R_2$ 的像的闭包。则
$$
(U, Z, s_2|_Z, t_2|_Z, c_2|_Z)
$$
是 $S$ 上的群胚概形。
\end{lemma}

\begin{proof}
先说明该陈述为何有意义。由于 $U$ 是完美域 $k$ 的谱，
无论经哪个投影，概形 $Z$ 在 $k$ 上都几何约化，见
《代数簇》引理 \ref{varieties-lemma-perfect-reduced}。
因此概形 $Z \times_{s_2, U, t_2} Z \subset Z$ 约化，见
《代数簇》引理 \ref{varieties-lemma-geometrically-reduced-any-base-change}。
于是由引理 \ref{more-groupoids-lemma-map-groupoids-on-field-image}，
$c$ 诱导态射 $Z \times_{s_2, U, t_2} Z \to Z$。
最后，显然 $e_2$ 经过 $Z$ 分解，而且映射
$i_2 : R_2 \to R_2$ 保持 $Z$。由于七元组
$(U, R_2, s_2, t_2, c_2, e_2, i_2)$ 的态射满足群胚公理，
限制到 $Z$ 后仍满足这些公理。
\end{proof}

\begin{lemma}
\label{more-groupoids-lemma-locally-closed-image-is-closed}
记号和假设同情形 \ref{more-groupoids-situation-morphism-groupoids-on-field}。
若像 $a(R_1)$ 是 $R_2$ 的局部闭子集，则它是闭子集。
\end{lemma}

\begin{proof}
设 $k \subset k'$ 为域 $k$ 的一个完美闭包。
设 $R_i'$ 为 $R_i$ 经态射
$U' = \Spec(k') \to \Spec(k)$ 的限制。注意，态射
$R_i' \to R_i$ 是普遍同胚，因为它们是普遍同胚
$U' \to U$ 的基变换的复合（见引理
\ref{more-groupoids-lemma-restrict-groupoid-on-field} 陈述中的图）。
因此只需证明 $a'(R_1')$ 在 $R_2'$ 中闭。换言之，可以假设 $k$ 完美。

\medskip\noindent
若 $k$ 完美，则由引理
\ref{more-groupoids-lemma-map-groupoids-on-perfect-field-image}，
像的闭包是群胚概形 $Z \subset R_2$。
把同一引理应用于
$\text{id}_{R_1} : R_1 \to R_1$，可见
$(R_2)_{red}$ 是群胚概形。因此可以把引理
\ref{more-groupoids-lemma-open-image-is-closed} 应用于态射
$a|_{(R_2)_{red}} : (R_2)_{red} \to Z$，
从而得到 $Z$ 等于 $a$ 的像。
\end{proof}

\begin{lemma}
\label{more-groupoids-lemma-quasi-compact-map-groupoids-on-field-image}
记号和假设同情形 \ref{more-groupoids-situation-morphism-groupoids-on-field}。
假设 $a : R_1 \to R_2$ 是拟紧态射。设 $Z \subset R_2$ 为
$a : R_1 \to R_2$ 的概形论像（见《概形态射》定义 \ref{morphisms-definition-scheme-theoretic-image}）。则
$$
(U, Z, s_2|_Z, t_2|_Z, c_2|_Z)
$$
是 $S$ 上的群胚概形。
\end{lemma}

\begin{proof}
主要困难是证明
$c_2|_{Z \times_{s_2, U, t_2} Z}$ 映入 $Z$。考虑交换图
$$
\xymatrix{
R_1 \times_{s_1, U, t_1} R_1 \ar[r] \ar[d]^{a \times a} & R_1 \ar[d] \\
R_2 \times_{s_2, U, t_2} R_2 \ar[r] & R_2
}
$$
由《代数簇》引理 \ref{varieties-lemma-scheme-theoretic-image-product-map}，
$a \times a$ 的概形论像是
$Z \times_{s_2, U, t_2} Z$。由图的交换性可知，
底部水平箭头把 $Z \times_{s_2, U, t_2} Z$ 映入 $Z$。
与引理 \ref{more-groupoids-lemma-map-groupoids-on-perfect-field-image}
的证明相同，还有 $i_2(Z) \subset Z$，且 $e_2$ 经过 $Z$ 分解。
因此可同该引理的证明一样得出结论。
\end{proof}

\begin{lemma}
\label{more-groupoids-lemma-groupoid-on-field-image}
设 $S$ 为概形，$(U, R, s, t, c)$ 为 $S$ 上的群胚概形。
假设 $U$ 是一个域的谱。设 $Z \subset U \times_S U$
为约化闭子概形（见《概形》定义 \ref{schemes-definition-reduced-induced-scheme}），其底拓扑空间是
$j = (t, s) : R \to U \times_S U$ 的像的闭包。则在集合论意义下，
$\text{pr}_{02}(Z \times_{\text{pr}_1, U, \text{pr}_0} Z) \subset Z$。
\end{lemma}

\begin{proof}
由于 $(U, U \times_S U, \text{pr}_1, \text{pr}_0, \text{pr}_{02})$
是 $S$ 上的群胚概形，这是引理
\ref{more-groupoids-lemma-map-groupoids-on-field-image} 的特殊情形。
不过也可如下直接证明。

\medskip\noindent
写成 $U = \Spec(k)$。记
$R_s$（相应地，\ $Z_s$；相应地，\ $U^2_s$）为概形
$R$（相应地，\ $Z$；相应地，\ $U \times_S U$），
把它经由 $s$（相应地，\ $\text{pr}_1|_Z$；相应地，\ 
$\text{pr}_1$）视为 $k$ 上的概形。类似地，记
${}_tR$（相应地，\ ${}_tZ$；相应地，\ ${}_tU^2$）为概形
$R$（相应地，\ $Z$；相应地，\ $U \times_S U$），
把它经由 $t$（相应地，\ $\text{pr}_0|_Z$；相应地，\ 
$\text{pr}_0$）视为 $k$ 上的概形。
态射 $j$ 诱导 $k$ 上的概形态射
$j_s : R_s \to U^2_s$ 和 ${}_tj : {}_tR \to {}_tU^2$。
考虑交换图
$$
\xymatrix{
R_s \times_k {}_tR \ar[r]^c \ar[d]_{j_s \times {}_tj} &  R \ar[d]^j \\
U^2_s \times_k {}_tU^2 \ar[r] & U \times_S U
}
$$
由《代数簇》引理 \ref{varieties-lemma-closure-image-product-map}，
$j_s \times {}_tj$ 之像的闭包是 $Z_s \times_k {}_tZ$。
由图的交换性可知，底部水平箭头把
$Z_s \times_k {}_tZ$ 映入 $Z$。
\end{proof}

\begin{lemma}
\label{more-groupoids-lemma-groupoid-on-perfect-field-image}
设 $S$ 为概形，$(U, R, s, t, c)$ 为 $S$ 上的群胚概形。
假设 $U$ 是一个完美域的谱。设 $Z \subset U \times_S U$
为约化闭子概形（见《概形》定义 \ref{schemes-definition-reduced-induced-scheme}），其底拓扑空间是
$j = (t, s) : R \to U \times_S U$ 的像的闭包。则
$$
(U, Z, \text{pr}_0|_Z, \text{pr}_1|_Z,
\text{pr}_{02}|_{Z \times_{\text{pr}_1, U, \text{pr}_0} Z})
$$
是 $S$ 上的群胚概形。
\end{lemma}

\begin{proof}
由于 $(U, U \times_S U, \text{pr}_1, \text{pr}_0, \text{pr}_{02})$
是 $S$ 上的群胚概形，这是引理
\ref{more-groupoids-lemma-map-groupoids-on-perfect-field-image} 的特殊情形。
不过也可如下直接证明。

\medskip\noindent
先说明该陈述为何有意义。由于 $U$ 是完美域 $k$ 的谱，
无论经哪个投影，概形 $Z$ 在 $k$ 上都几何约化，见
《代数簇》引理 \ref{varieties-lemma-perfect-reduced}。
因此概形
$Z \times_{\text{pr}_1, U, \text{pr}_0} Z \subset Z$
约化，见《代数簇》引理 \ref{varieties-lemma-geometrically-reduced-any-base-change}。
于是由引理 \ref{more-groupoids-lemma-groupoid-on-field-image}，
$\text{pr}_{02}$ 诱导态射
$Z \times_{\text{pr}_1, U, \text{pr}_0} Z \to Z$。
最后，显然 $\Delta_{U/S}$ 经过 $Z$ 分解，而且映射
$\sigma : U \times_S U \to U \times_S U$、
$(x, y) \mapsto (y, x)$ 保持 $Z$。由于
$(U, U \times_S U, \text{pr}_0, \text{pr}_1, \text{pr}_{02},
\Delta_{U/S}, \sigma)$
满足群胚公理，限制到 $Z$ 后仍满足这些公理。
\end{proof}

\begin{lemma}
\label{more-groupoids-lemma-quasi-compact-groupoid-on-field-image}
设 $S$ 为概形，$(U, R, s, t, c)$ 为 $S$ 上的群胚概形。
假设 $U$ 是一个域的谱，并假设 $R$ 拟紧
（等价地，$s, t$ 拟紧）。设 $Z \subset U \times_S U$
为 $j = (t, s) : R \to U \times_S U$ 的概形论像
（见《概形态射》定义 \ref{morphisms-definition-scheme-theoretic-image}）。则
$$
(U, Z, \text{pr}_0|_Z, \text{pr}_1|_Z,
\text{pr}_{02}|_{Z \times_{\text{pr}_1, U, \text{pr}_0} Z})
$$
是 $S$ 上的群胚概形。
\end{lemma}

\begin{proof}
由于 $(U, U \times_S U, \text{pr}_1, \text{pr}_0, \text{pr}_{02})$
是 $S$ 上的群胚概形，这是引理
\ref{more-groupoids-lemma-quasi-compact-map-groupoids-on-field-image} 的特殊情形。
不过也可如下直接证明。

\medskip\noindent
主要困难是证明
$\text{pr}_{02}|_{Z \times_{\text{pr}_1, U, \text{pr}_0} Z}$
映入 $Z$。写成 $U = \Spec(k)$。记
$R_s$（相应地，\ $Z_s$；相应地，\ $U^2_s$）为概形
$R$（相应地，\ $Z$；相应地，\ $U \times_S U$），
把它经由 $s$（相应地，\ $\text{pr}_1|_Z$；相应地，\ 
$\text{pr}_1$）视为 $k$ 上的概形。类似地，记
${}_tR$（相应地，\ ${}_tZ$；相应地，\ ${}_tU^2$）为概形
$R$（相应地，\ $Z$；相应地，\ $U \times_S U$），
把它经由 $t$（相应地，\ $\text{pr}_0|_Z$；相应地，\ 
$\text{pr}_0$）视为 $k$ 上的概形。
态射 $j$ 诱导 $k$ 上的概形态射
$j_s : R_s \to U^2_s$ 和 ${}_tj : {}_tR \to {}_tU^2$。
考虑交换图
$$
\xymatrix{
R_s \times_k {}_tR \ar[r]^c \ar[d]_{j_s \times {}_tj} &  R \ar[d]^j \\
U^2_s \times_k {}_tU^2 \ar[r] & U \times_S U
}
$$
由《代数簇》引理 \ref{varieties-lemma-scheme-theoretic-image-product-map}，
$j_s \times {}_tj$ 的概形论像是 $Z_s \times_k {}_tZ$。
由图的交换性可知，底部水平箭头把
$Z_s \times_k {}_tZ$ 映入 $Z$。
与引理 \ref{more-groupoids-lemma-groupoid-on-perfect-field-image} 的证明相同，
还有 $\sigma(Z) \subset Z$，且 $\Delta_{U/S}$ 经过 $Z$ 分解。
因此可同该引理的证明一样得出结论。
\end{proof}







\section{群胚的切片}
\label{more-groupoids-section-slicing}

\noindent
下一个引理表明，只要稳定子的维数足够小，就可以对一个
Cohen--Macaulay 群胚概形作切片，以降低纤维的维数。
这是改进商叠的给定表现过程中的关键一步。

\begin{situation}
\label{more-groupoids-situation-slice}
设 $S$ 为概形。设 $(U, R, s, t, c)$ 为 $S$ 上的群胚概形。
设 $g : U' \to U$ 为概形态射。设 $u \in U$ 为一点，
并设 $u' \in U'$ 为满足 $g(u') = u$ 的一点。
给定这些数据，记 $(U', R', s', t', c')$ 为
$(U, R, s, t, c)$ 经态射 $g$ 的限制。
记 $G \to U$ 为 $R$ 的稳定子群概形；它是 $R$ 的局部闭子概形。
记 $h$ 为复合
$$
h = s \circ \text{pr}_1 : U' \times_{g, U, t} R \longrightarrow U.
$$
记 $F_u = s^{-1}(u)$（概形论纤维），并记 $G_u$ 为 $G$ 在 $u$
上的概形论纤维。类似地，对 $R'$ 记
$F'_{u'} = (s')^{-1}(u')$。由于 $g(u') = u$，有
$$
F'_{u'} = h^{-1}(u) \times_{\Spec(\kappa(u))} \Spec(\kappa(u')).
$$
点 $e(u) \in R$ 也可视为 $G_u$ 和 $F_u$ 上的点；
而 $e'(u')$ 是 $R'$（相应地，\ $G'_{u'}$；相应地，\ 
$F'_{u'}$）的一点，它映到 $R$（相应地，\ $G_u$；相应地，\ 
$F_u$）中的 $e(u)$。
\end{situation}

\begin{lemma}
\label{more-groupoids-lemma-slice}
设 $S$ 为概形。设 $(U, R, s, t, c, e, i)$ 为 $S$ 上的群胚概形。
设 $G \to U$ 为稳定子群概形。假设 $s$ 和 $t$ 是
Cohen--Macaulay 且局部有限表现的。设 $u \in U$ 是概形 $U$
的有限型点，见《概形态射》定义 \ref{morphisms-definition-finite-type-point}。采用情形
\ref{more-groupoids-situation-slice} 的记号，令
$$
d_1 = \dim(G_u), \quad
d_2 = \dim_{e(u)}(F_u).
$$
若 $d_2 > d_1$，则存在仿射概形 $U'$ 和态射 $g : U' \to U$，
使得（采用情形 \ref{more-groupoids-situation-slice} 的记号）
\begin{enumerate}
\item $g$ 是浸入；
\item $u \in U'$；
\item $g$ 局部有限表现；
\item 态射 $h : U' \times_{g, U, t} R \longrightarrow U$
在 $(u, e(u))$ 处是 Cohen--Macaulay 的；并且
\item $\dim_{e'(u)}(F'_u) = d_2 - 1$。
\end{enumerate}
\end{lemma}

\begin{proof}
取 $u$ 的仿射邻域 $\Spec(A) \subset U$，使得 $u$ 对应于 $U$
的闭点，见《概形态射》引理 \ref{morphisms-lemma-identify-finite-type-points}。
取 $e(u)$ 的仿射邻域 $\Spec(B) \subset R$，使其在 $j$ 下映入开集
$\Spec(A) \times_S \Spec(A) \subset U \times_S U$。
设 $\mathfrak m \subset A$ 为对应于 $u$ 的极大理想，
并设 $\mathfrak q \subset B$ 为对应于 $e(u)$ 的素理想。图示为
$$
\vcenter{
\xymatrix{
B & A \ar[l]^s \\
A \ar[u]^t
}
}
\quad\text{且}\quad
\vcenter{
\xymatrix{
B_{\mathfrak q} & A_{\mathfrak m} \ar[l]^s \\
A_{\mathfrak m} \ar[u]^t
}
}
$$
注意，两个诱导映射
$s, t : \kappa(\mathfrak m) \to \kappa(\mathfrak q)$
相等且均为同构，因为
$s \circ e = t \circ e = \text{id}_U$。
特别地，$\mathfrak q$ 也是极大理想。
环同态 $s, t : A \to B$ 有限表现且平坦。由假设，环
$$
\mathcal{O}_{F_u, e(u)} = B_{\mathfrak q}/s(\mathfrak m)B_{\mathfrak q}
$$
是维数为 $d_2$ 的 Cohen--Macaulay 环。该维数等式由
《概形态射》引理 \ref{morphisms-lemma-dimension-fibre-at-a-point} 给出。


\medskip\noindent
设 $R''$ 为 $R$ 经态射
$\Spec(\kappa(u)) \to U$ 向
$u = \Spec(\kappa(u))$ 的限制。由于 $u \to U$ 局部有限型，
可见 $(\Spec(\kappa(u)), R'', s'', t'', c'')$
是一个 $s'', t''$ 局部有限型的群胚概形，见引理
\ref{more-groupoids-lemma-restrict-preserves-type}。
由引理 \ref{more-groupoids-lemma-groupoid-on-field-dimension-equal-stabilizer}，
这蕴含 $\dim(G'') = \dim(R'')$。又由引理
\ref{more-groupoids-lemma-groupoid-on-field-locally-finite-type-dimension}，
有 $\dim(R'') = \dim_{e''}(R'') = \dim(\mathcal{O}_{R'', e''})$。
由《群胚概形》引理 \ref{groupoids-lemma-restrict-stabilizer}，有 $G'' = G_u$。
因此得到 $\dim(\mathcal{O}_{R'', e''}) = d_1$。

\medskip\noindent
作为概形，$R''$ 为
$$
R'' =
R \times_{(U \times_S U)}
\Big(
\Spec(\kappa(\mathfrak m)) \times_S \Spec(\kappa(\mathfrak m))
\Big)
$$
因此 $e''$ 的一个仿射开邻域是下列环的谱：
$$
B \otimes_{(A \otimes A)} (\kappa(\mathfrak m) \otimes \kappa(\mathfrak m))
=
B/s(\mathfrak m)B + t(\mathfrak m)B
$$
由此得到
$$
\mathcal{O}_{R'', e''} =
B_{\mathfrak q}/s(\mathfrak m)B_{\mathfrak q} + t(\mathfrak m)B_{\mathfrak q}
$$
于是现在知道此环的维数为 $d_1$。

\medskip\noindent
我们断言，这蕴含可以找到元素 $f \in \mathfrak m$，使得
$$
\dim(B_{\mathfrak q}/(s(\mathfrak m)B_{\mathfrak q} + fB_{\mathfrak q}) < d_2
$$
确实，设 $\mathfrak n_j \supset s(\mathfrak m)B_{\mathfrak q}$，
$j = 1, \ldots, m$，对应于局部环
$B_{\mathfrak q}/s(\mathfrak m)B_{\mathfrak q}$ 的极小素理想。
这样的理想只有有限多个，因为该环是 Noether 环
（它本质上是域上的有限型环；也可因为 Cohen--Macaulay 环是
Noether 环）。由 Cohen--Macaulay 条件，
$\dim(B_{\mathfrak q}/\mathfrak n_j) = d_2$，例如可见
《交换代数》引理 \ref{algebra-lemma-CM-dim-formula}。注意
$\dim(B_{\mathfrak q}/(\mathfrak n_j + t(\mathfrak m)B_{\mathfrak q}))
\leq d_1$，
因为它是环
$\mathcal{O}_{R'', e''} =
B_{\mathfrak q}/s(\mathfrak m)B_{\mathfrak q} + t(\mathfrak m)B_{\mathfrak q}$
的商，而后者维数为 $d_1$。由于 $d_1 < d_2$，这蕴含
$\mathfrak m \not \subset t^{-1}(\mathfrak n_i)$。
由素理想回避，见《交换代数》引理 \ref{algebra-lemma-silly}，
可以找到 $f \in \mathfrak m$，使得对
$j = 1, \ldots, m$ 都有 $t(f) \not \in \mathfrak n_j$。
对这样选取的 $f$，上述显示不等式成立，见
《交换代数》引理 \ref{algebra-lemma-one-equation}。


\medskip\noindent
令 $A' = A/fA$、$U' = \Spec(A')$。显然 $U' \to U$ 是浸入，
局部有限表现，且 $u \in U'$。因此引理的 (1)、(2)、(3) 成立。
态射
$$
U' \times_{g, U, t} R \longrightarrow U
$$
经过 $\Spec(A)$ 分解，并对应于环同态
$$
\xymatrix{
B/t(f)B \ar@{=}[r] & A/(f) \otimes_{A, t} B & A \ar[l]_-s
}
$$
现在，由于
$B_{\mathfrak q}/s(\mathfrak m)B_{\mathfrak q}$
是正维 Cohen--Macaulay 环，且 $f$ 不属于任何极小素理想，
可见 $t(f)$ 在其上不是零因子；例如见《交换代数》引理 \ref{algebra-lemma-reformulate-CM}。因此由《交换代数》引理 \ref{algebra-lemma-grothendieck-general}，可知
$s : A_{\mathfrak m} \to B_{\mathfrak q}/t(f)B_{\mathfrak q}$
平坦，其纤维环为
$B_{\mathfrak q}/(s(\mathfrak m)B_{\mathfrak q} + t(f)B_{\mathfrak q})$；
由《交换代数》引理 \ref{algebra-lemma-reformulate-CM}，
该环仍为 Cohen--Macaulay。这蕴含引理的 (4)。
为看出 (5)，注意由图 (\ref{more-groupoids-equation-restriction})，
纤维 $F'_u$ 等于 $h$ 在 $u$ 上的纤维。因此由
《概形态射》引理 \ref{morphisms-lemma-dimension-fibre-at-a-point}，
$\dim_{e'(u)}(F'_u) =
\dim(B_{\mathfrak q}/(s(\mathfrak m)B_{\mathfrak q} + t(f)B_{\mathfrak q}))$
而由《交换代数》引理 \ref{algebra-lemma-reformulate-CM}，
此环的维数为 $d_2 - 1$。这证明了引理的最后一个断言，证明完成。
\end{proof}

\noindent
既然已经知道如何作切片，就可以把它同前面的材料结合起来，
得到下列“最优”结论。这里所谓最优，是因为 $G_u$ 是 $F_u$
的局部闭子概形，总有不等式
$\dim(G_u) = \dim_{e(u)}(G_u) \leq \dim_{e(u)}(F_u)$，
所以不可能比引理所给出的切片更进一步。

\begin{lemma}
\label{more-groupoids-lemma-max-slice}
设 $S$ 为概形。设 $(U, R, s, t, c, e, i)$ 为 $S$ 上的群胚概形。
设 $G \to U$ 为稳定子群概形。假设 $s$ 和 $t$ 是
Cohen--Macaulay 且局部有限表现的。设 $u \in U$ 是概形 $U$
的有限型点，见《概形态射》定义 \ref{morphisms-definition-finite-type-point}。
采用情形 \ref{more-groupoids-situation-slice} 的记号，存在仿射概形 $U'$
和态射 $g : U' \to U$，使得
\begin{enumerate}
\item $g$ 是浸入；
\item $u \in U'$；
\item $g$ 局部有限表现；
\item 态射 $h : U' \times_{g, U, t} R \longrightarrow U$
是 Cohen--Macaulay 且局部有限表现的；
\item 态射 $s', t' : R' \to U'$ 是 Cohen--Macaulay
且局部有限表现的；并且
\item $\dim_{e(u)}(F'_u) = \dim(G'_u)$。
\end{enumerate}
\end{lemma}

\begin{proof}
由于 $s$ 局部有限表现，概形 $F_u$ 在 $\kappa(u)$ 上局部有限型。
因此 $\dim_{e(u)}(F_u) < \infty$，可以对
$\dim_{e(u)}(F_u)$ 作归纳。

\medskip\noindent
若 $\dim_{e(u)}(F_u) = \dim(G_u)$，则无事可证。
假设 $\dim_{e(u)}(F_u) > \dim(G_u)$。这意味着可以应用引理
\ref{more-groupoids-lemma-slice}，从而得到态射 $g : U' \to U$，它满足
(1)、(2)、(3)；以
$\dim_{e(u)}(F'_u) < \dim_{e(u)}(F_u)$ 代替 (6)；
并以复合
$$
h = s \circ \text{pr}_1 : U' \times_{g, U, t} R \longrightarrow U
$$
在点 $(u, e(u))$ 处为 Cohen--Macaulay 来代替 (4) 和 (5)。
应用注 \ref{more-groupoids-remark-local-source-apply}，得到开子概形
$U'' \subset U'$，使得
$U'' \times_{g, U, t} R \subset U' \times_{g, U, t} R$
是 $h$ 为 Cohen--Macaulay 的最大开子概形。
由于 $(u, e(u)) \in U'' \times_{g, U, t} R$，可见
$u \in U''$。因此可以用 $U''$ 代替 $U'$，并假设 $h$
处处为 Cohen--Macaulay！由引理 \ref{more-groupoids-lemma-restrict-property}，
可知 $s', t'$ 局部有限表现且为 Cohen--Macaulay
（使用《概形态射》引理 \ref{morphisms-lemma-base-change-finite-presentation} 及
《态射进阶》引理 \ref{more-morphisms-lemma-base-change-CM}）。

\medskip\noindent
由构造，$\dim_{e'(u)}(F'_u) < \dim_{e(u)}(F_u)$，
故可将归纳假设应用于 $(U', R', s', t', c')$ 及点 $u \in U'$。
注意 $u$ 也是 $U'$ 的有限型点（例如可使用《概形态射》引理 \ref{morphisms-lemma-identify-finite-type-points}
对有限型点的刻画看出这一点）。设
$g' : U'' \to U'$ 和 $(U'', R'', s'', t'', c'')$
是从 $(U', R', s', t', c')$ 及点 $u \in U'$ 出发所得相应问题的解。
我们断言复合
$$
g'' = g \circ g' : U'' \longrightarrow U
$$
是原问题的一个解。性质 (1)、(2)、(3)、(5)、(6) 显然成立。
为看出 (4)，注意态射
$$
h'' = s \circ \text{pr}_1 : U'' \times_{g'', U, t} R \longrightarrow U
$$
局部有限表现且为 Cohen--Macaulay；这可应用引理
\ref{more-groupoids-lemma-double-restrict-property} 得出
（使用《态射进阶》引理 \ref{more-morphisms-lemma-CM-local-source-and-target}
来说明 Cohen--Macaulay 态射在目标上是 fppf 局部的）。
\end{proof}

\noindent
当稳定子群概形的纤维维数为 0 时，这给出下列切片引理。

\begin{lemma}
\label{more-groupoids-lemma-max-slice-quasi-finite}
设 $S$ 为概形。设 $(U, R, s, t, c, e, i)$ 为 $S$ 上的群胚概形。
设 $G \to U$ 为稳定子群概形。假设 $s$ 和 $t$ 是
Cohen--Macaulay 且局部有限表现的。设 $u \in U$ 是概形 $U$
的有限型点，见《概形态射》定义 \ref{morphisms-definition-finite-type-point}。
假设 $G \to U$ 局部拟有限。采用情形
\ref{more-groupoids-situation-slice} 的记号，存在仿射概形 $U'$
和态射 $g : U' \to U$，使得
\begin{enumerate}
\item $g$ 是浸入；
\item $u \in U'$；
\item $g$ 局部有限表现；
\item 态射 $h : U' \times_{g, U, t} R \longrightarrow U$
平坦、局部有限表现且局部拟有限；并且
\item 态射 $s', t' : R' \to U'$ 平坦、局部有限表现且局部拟有限。
\end{enumerate}
\end{lemma}

\begin{proof}
取引理 \ref{more-groupoids-lemma-max-slice} 中的 $g : U' \to U$。
由于 $h^{-1}(u) = F'_u$，可见 $h$ 在 $(u, e(u))$ 处相对维数
$\leq 0$。因此由注 \ref{more-groupoids-remark-local-source-apply}，
得到开子概形 $U'' \subset U'$，使得 $u \in U''$，且
$U'' \times_{g, U, t} R$ 是
$U' \times_{g, U, t} R$ 中 $h$ 相对维数 $\leq 0$ 的最大开子概形。
用 $U''$ 代替 $U'$ 后，可见 $h$ 的相对维数 $\leq 0$。
由《概形态射》引理 \ref{morphisms-lemma-locally-quasi-finite-rel-dimension-0}，
这蕴含 $h$ 局部拟有限。由于它仍局部有限表现且为
Cohen--Macaulay，可见它平坦、局部有限表现且局部拟有限，
即上面的 (4) 成立。这蕴含 $s'$ 作为 $h$ 的基变换也平坦、
局部有限表现且局部拟有限，见引理 \ref{more-groupoids-lemma-restrict-property}。
\end{proof}







\section{群胚的平展（\'Etale）局部化}
\label{more-groupoids-section-etale-localize}

\noindent
本节开始把《态射进阶》第 \ref{more-morphisms-section-etale-localization} 节
中的平展（\'etale）局部化技术应用于群胚概形。
这方面更深入的材料见《代数空间中的群胚进阶》第 \ref{spaces-more-groupoids-section-etale-localize} 节。
引理 \ref{more-groupoids-lemma-quasi-finite-over-base-j-proper}
将用于证明关于在一个概形上分离且拟有限的代数空间的结果，
即《代数空间的态射》命题 \ref{spaces-morphisms-proposition-locally-quasi-finite-separated-over-scheme}
及其推论《代数空间的态射》引理 \ref{spaces-morphisms-lemma-locally-quasi-finite-separated-representable}。

\begin{lemma}
\label{more-groupoids-lemma-quasi-finite-over-base}
设 $S$ 为概形。设 $(U, R, s, t, c)$ 为 $S$ 上的群胚概形。
设 $p \in S$ 为一点，并设 $u \in U$ 为位于 $p$ 上方的一点。假设
\begin{enumerate}
\item $U \to S$ 局部有限型；
\item $U \to S$ 在 $u$ 处拟有限；
\item $U \to S$ 分离；
\item $R \to S$ 分离；
\item $s$, $t$ 平坦且局部有限表现；并且
\item $s^{-1}(\{u\})$ 有限。
\end{enumerate}
则存在平展（\'etale）邻域 $(S', p') \to (S, p)$，满足
$\kappa(p) = \kappa(p')$，并有基变换图
$$
\xymatrix{
R' \amalg W'
\ar@{=}[r] &
S' \times_S R
\ar[r] \ar@<2ex>[d]^{s'} \ar@<-2ex>[d]_{t'} &
R \ar@<1ex>[d]^s \ar@<-1ex>[d]_t \\
U' \amalg W
\ar@{=}[r] &
S' \times_S U
\ar[r] \ar[d] &
U \ar[d] \\
 &
S' \ar[r] &
S
}
$$
其中等号表示分解为开闭子概形，并满足
\begin{enumerate}
\item[(a)] 存在 $U'$ 的一点 $u'$，它在 $U$ 中映到 $u$；
\item[(b)] 在集合论意义下，纤维
$(U')_{p'}$ 等于 $t'\big((s')^{-1}(\{u'\})\big)$；
\item[(c)] 在集合论意义下，纤维
$(R')_{p'}$ 等于 $(s')^{-1}\big((U')_{p'}\big)$；
\item[(d)] 概形 $U'$ 和 $R'$ 在 $S'$ 上有限；
\item[(e)] 有 $s'(R') \subset U'$ 和 $t'(R') \subset U'$；
\item[(f)] 有
$c'(R' \times_{s', U', t'} R') \subset R'$，
其中 $c'$ 是 $c$ 的基变换；并且
\item[(g)] 态射 $s', t', c'$ 通过系统
$(U', R', s'|_{R'}, t'|_{R'}, c'|_{R' \times_{s', U', t'} R'})$
确定一个群胚结构。
\end{enumerate}
\end{lemma}

\begin{proof}
记 $f : U \to S$ 为 $U$ 的结构态射。由假设 (6)，可写
$s^{-1}(\{u\}) = \{r_1, \ldots, r_n\}$。
由于该集合有限，可见 $s$ 在这有限多个逆像中的每一个处拟有限，
见《概形态射》引理 \ref{morphisms-lemma-finite-fibre}。
因此 $f \circ s : R \to S$ 在每个 $r_i$ 处拟有限
（《概形态射》引理 \ref{morphisms-lemma-composition-quasi-finite}）。
所以 $r_i$ 在纤维 $R_p$ 中是孤立点，见《概形态射》引理 \ref{morphisms-lemma-quasi-finite-at-point-characterize}。
写成 $t(\{r_1, \ldots, r_n\}) = \{u_1, \ldots, u_m\}$。
注意可能有 $m < n$，并且
$u \in \{u_1, \ldots, u_m\}$。
由于 $t$ 平坦且局部有限表现，纤维态射
$t_p : R_p \to U_p$ 平坦且局部有限表现
（《概形态射》引理 \ref{morphisms-lemma-base-change-flat} 和
\ref{morphisms-lemma-base-change-finite-presentation}），
因而为开映射（《概形态射》引理 \ref{morphisms-lemma-fppf-open}）。
每个 $r_i$ 在 $R_p$ 中孤立这一事实蕴含每个
$u_j = t(r_i)$ 在 $U_p$ 中孤立。再次使用《概形态射》引理 \ref{morphisms-lemma-quasi-finite-at-point-characterize}，
可见 $f$ 在 $u_1, \ldots, u_m$ 处拟有限。


\medskip\noindent
记 $F_u = s^{-1}(u)$ 和 $F_{u_j} = s^{-1}(u_j)$ 为概形论纤维。
注意 $F_u$ 在 $\kappa(u)$ 上有限，因为它在 $\kappa(u)$ 上局部有限型
且只有有限多个点（例如，这可由远为一般的《概形态射》引理 \ref{morphisms-lemma-locally-quasi-finite-qc-source-universally-bounded}
得出）。由引理 \ref{more-groupoids-lemma-two-fibres}，
$F_u$ 和 $F_{u_j}$ 在 $\kappa(u)$ 与 $\kappa(u_j)$ 的某个公共域扩张上
变为同构。因此 $F_{u_j}$ 在 $\kappa(u_j)$ 上有限。特别地，
对每个 $j = 1, \ldots, m$，集合 $s^{-1}(\{u_j\})$ 有限。
于是对每个 $u_j$，假设 (2) 和 (6) 也成立
（上面已看到 $U \to S$ 在 $u_j$ 处拟有限）。
因此第一段的论证适用于每个 $u_j$，从而 $R \to U$
在下列集合的每个点处拟有限：
$$
\{r_1, \ldots, r_N\} = s^{-1}(\{u_1, \ldots, u_m\})
$$
注意
$t(\{r_1, \ldots, r_N\}) = \{u_1, \ldots, u_m\}$ 且
$t^{-1}(\{u_1, \ldots, u_m\}) = \{r_1, \ldots, r_N\}$，
因为 $R$ 是群胚\footnote{用群胚语言解释：
原集合 $\{r_1, \ldots, r_n\}$ 是以 $u$ 为源的箭头集合；
集合 $\{u_1, \ldots, u_m\}$ 是与 $u$ 同构的对象集合；
而 $\{r_1, \ldots, r_N\}$ 是所有与 $u$ 等价的对象之间的全部箭头集合。}。
此外，
$\text{pr}_0(c^{-1}(\{r_1, \ldots, r_N\})) = \{r_1, \ldots, r_N\}$
且
$\text{pr}_1(c^{-1}(\{r_1, \ldots, r_N\})) = \{r_1, \ldots, r_N\}$。
类似地，还得到
$e(\{u_1, \ldots, u_m\}) \subset \{r_1, \ldots, r_N\}$
以及
$i(\{r_1, \ldots, r_N\}) = \{r_1, \ldots, r_N\}$。


\medskip\noindent
可以把《态射进阶》引理 \ref{more-morphisms-lemma-etale-splits-off-quasi-finite-part-technical}
应用于两对
$(U \to S, \{u_1, \ldots, u_m\})$ 和
$(R \to S, \{r_1, \ldots, r_N\})$，
从而得到一个平展（\'etale）邻域
$(S', p') \to (S, p)$，它诱导等同
$\kappa(p) = \kappa(p')$，并使 $S' \times_S U$ 和
$S' \times_S R$ 分解为
$$
S' \times_S U = U' \amalg W, \quad
S' \times_S R = R' \amalg W'
$$
其中 $U' \to S'$ 有限，且 $(U')_{p'}$ 双射地映到
$\{u_1, \ldots, u_m\}$；$R' \to S'$ 有限，且
$(R')_{p'}$ 双射地映到 $\{r_1, \ldots, r_N\}$。
此外，$W_{p'}$（相应地，\ $(W')_{p'}$）的任何点都不映到
任何 $u_j$（相应地，\ $r_i$）。此时，引理的 (a)、(b)、(c)、(d)
均已满足。此外，(e) 和 (f) 中的包含关系在 $p'$ 上的纤维上成立，
即 $s'((R')_{p'}) \subset (U')_{p'}$、
$t'((R')_{p'}) \subset (U')_{p'}$，以及
$c'((R' \times_{s', U', t'} R')_{p'}) \subset (R')_{p'}$。


\medskip\noindent
我们断言，可以把 $S'$ 换成 $p'$ 的一个 Zariski 开邻域，
使 (e) 和 (f) 中的包含关系成立。例如，考虑集合
$E = (s'|_{R'})^{-1}(W)$。这是 $R'$ 的开闭子集，且不包含
$R'$ 中位于 $p'$ 上方的任何点。由于 $R' \to S'$ 是闭映射，
用 $S' \setminus (R' \to S')(E)$ 代替 $S'$ 后，
就得到 $E$ 为空的情形。换言之，$s'$ 把 $R'$ 映入 $U'$。
注意，进一步缩小 $S'$ 时此性质仍然保持。
类似地，可以安排使 $t'$ 把 $R'$ 映入 $U'$。此时 (e) 成立。
以同样方式考虑集合
$E = (c'|_{R' \times_{s', U', t'} R'})^{-1}(W')$。
它是概形 $R' \times_{s', U', t'} R'$ 的开闭子集；
该概形在 $S'$ 上有限，而且它不包含位于 $p'$ 上方的任何点。
因此，用
$S' \setminus (R' \times_{s', U', t'} R' \to S')(E)$
代替 $S'$ 后，就得到 $E$ 为空的情形。换言之，得到 (f) 中的包含关系。
还可以对恒等态射
$e' : S' \times_S U \to S' \times_S R$
和逆态射 $i' : S' \times_S R \to S' \times_S R$
重复此论证，从而可以假设（再把 $S'$ 缩小一些后）
$(e'|_{U'})^{-1}(W') = \emptyset$ 且
$(i'|_{R'})^{-1}(W') = \emptyset$。

\medskip\noindent
此时可以考虑结构
$$
(U', R', s'|_{R'}, t'|_{R'}, c'|_{R' \times_{t', U', s'} R'},
e'|_{U'}, i'|_{R'}).
$$
$S'$ 上群胚概形的公理成立，因为这些公理对群胚概形
$(S' \times_S U, S' \times_S R, s', t', c', e', i')$
成立。
\end{proof}

\begin{lemma}
\label{more-groupoids-lemma-quasi-finite-over-base-j-proper}
设 $S$ 为概形。设 $(U, R, s, t, c)$ 为 $S$ 上的群胚概形。
设 $p \in S$ 为一点，并设 $u \in U$ 为位于 $p$ 上方的一点。
假设引理 \ref{more-groupoids-lemma-quasi-finite-over-base} 的假设 (1) -- (6)
成立，并且还有
\begin{enumerate}
\item[(7)] $j : R \to U \times_S U$ 普遍闭\footnote{
鉴于其余条件，这等价于要求 $j$ 为固有态射。}。
\end{enumerate}
则可以选取 $(S', p') \to (S, p)$、分解
$S' \times_S U = U' \amalg W$ 和
$S' \times_S R = R' \amalg W'$，以及 $u' \in U'$，
使引理 \ref{more-groupoids-lemma-quasi-finite-over-base} 的 (a) -- (g) 成立，
并且还有
\begin{enumerate}
\item[(h)] $R'$ 是 $S' \times_S R$ 向 $U'$ 的限制。
\end{enumerate}
\end{lemma}

\begin{proof}
把引理 \ref{more-groupoids-lemma-quasi-finite-over-base} 应用于概形 $S$ 上的群胚
$(U, R, s, t, c)$ 及点 $p$、$u$。由此得到平展（\'etale）邻域
$(S', p') \to (S, p)$、不交并分解
$$
S' \times_S U = U' \amalg W, \quad
S' \times_S R = R' \amalg W'
$$
以及满足结论 (a)、(b)、(c)、(d)、(e)、(f)、(g) 的 $u' \in U'$。
可以把 $S'$ 缩小为 $p'$ 的更小邻域，而不影响结论 (a) -- (g)。
下面证明适当缩小后结论 (h) 也成立。记 $j'$ 为 $j$ 到 $S'$ 的基变换。
由结论 (e)，显然
$$
j'^{-1}(U' \times_{S'} U') = R' \amalg Rest
$$
其中 $Rest$ 是某个开闭部分。由结论 (d)，$U' \to S'$ 有限，
故 $U' \times_{S'} U'$ 在 $S'$ 上有限。
由于 $j$ 普遍闭，$j'$ 也普遍闭，从而 $j'|_{Rest}$
也普遍闭。由结论 (b) 和 (c)，可见态射
$$
(U' \times_{S'} U' \to S') \circ j'|_{Rest} :
Rest
\longrightarrow
S'
$$
在 $p'$ 上的纤维为空。因此，由于 $Rest \to S'$ 是闭态射的复合，
故为闭态射；用
$S' \setminus \Im(Rest \to S')$ 代替 $S'$ 后，
可以假设 $Rest = \emptyset$。这恰好是说 $R'$ 是
$S' \times_S R$ 向开子概形 $U' \subset S' \times_S U$ 的限制，
见《群胚概形》引理 \ref{groupoids-lemma-restrict-groupoid-relation} 及其证明。
\end{proof}






\section{有限群胚}
\label{more-groupoids-section-finite-groupoids}

\noindent
若态射 $s$ 和 $t$ 有限，群胚概形 $(U, R, s, t, c)$
有时称为{\it 有限的}。这个说法可能造成混淆，因为它并不蕴含
$U$、$R$ 或商层 $U/R$ 在任何对象上有限。

\begin{lemma}
\label{more-groupoids-lemma-finite-stratify}
设 $(U, R, s, t, c)$ 为概形 $S$ 上的群胚概形，并假设 $s, t$ 有限。
存在一列 $R$-不变闭子概形
$$
U = Z_0 \supset Z_1 \supset Z_2 \supset \ldots
$$
使得 $\bigcap Z_r = \emptyset$，并且
$s^{-1}(Z_{r - 1}) \setminus s^{-1}(Z_r) \to Z_{r - 1} \setminus Z_r$
有限局部自由，秩为 $r$。
\end{lemma}

\begin{proof}
设 $\{Z_r\}$ 为有限型拟凝聚模
$s_*\mathcal{O}_R$ 的 Fitting 理想在 $U$ 上给出的分层，见
《除子》引理 \ref{divisors-lemma-locally-free-rank-r-pullback}。
由于恒等态射 $e : U \to R$ 是 $s$ 的截面，可见
$s_*\mathcal{O}_R$ 以 $\mathcal{O}_S$ 为直和项。
因此 $U = Z_{-1} = Z_0$（略去细节）。
由于 Fitting 理想的形成与基变换交换
（《代数进阶》引理 \ref{more-algebra-lemma-fitting-ideal-basics}），
并且图 (\ref{more-groupoids-equation-diagram}) 的右下方块为 Cartesian 方块，
可见 $s^{-1}(Z_r)$ 对应于
$\text{pr}_{1, *}\mathcal{O}_{R \times_{s, U, t} R}$
的第 $r$ 个 Fitting 理想。再利用左下方块也为 Cartesian 方块，
得到 $s^{-1}(Z_r) = t^{-1}(Z_r)$；换言之，$Z_r$ 为 $R$-不变。
态射
$s^{-1}(Z_{r - 1}) \setminus s^{-1}(Z_r) \to Z_{r - 1} \setminus Z_r$
有限局部自由且秩为 $r$，因为由《除子》引理 \ref{divisors-lemma-locally-free-rank-r-pullback}，
模 $s_*\mathcal{O}_R$ 在 $Z_{r - 1} \setminus Z_r$ 上拉回为
秩 $r$ 的有限局部自由模。
\end{proof}

\begin{lemma}
\label{more-groupoids-lemma-finite-flat-over-almost-dense-subscheme}
设 $(U, R, s, t, c)$ 为概形 $S$ 上的群胚概形，并假设 $s, t$ 有限。
存在开子概形 $W \subset U$ 和闭子概形 $W' \subset W$，使得
\begin{enumerate}
\item $W$ 和 $W'$ 均为 $R$-不变；
\item 在集合论意义下，$U = t(s^{-1}(\overline{W}))$；
\item $W$ 是 $W'$ 的加厚；并且
\item 限制 $(W', R', s', t', c')$ 的映射 $s'$, $t'$
有限局部自由。
\end{enumerate}
\end{lemma}

\begin{proof}
考虑引理 \ref{more-groupoids-lemma-finite-stratify} 的分层
$U = Z_0 \supset Z_1 \supset Z_2 \supset \ldots$。

\medskip\noindent
我们将构造不交并
$W = \coprod_{r \geq 1} W_r$ 和
$W' = \coprod_{r \geq 1} W'_r$，使每个
$W'_r \to W_r$ 都是 $U$ 的 $R$-不变子概形的加厚，
且限制 $(W_r', R_r', s_r', t_r', c_r')$ 的态射
$s_r', t_r'$ 有限局部自由，秩为 $r$。
首先令 $W_1 = W'_1 = U \setminus Z_1$。
这是 $U$ 的 $R$-不变开子概形；$W_0$ 确实是 $W'_0$ 的加厚；
而限制 $(W_1', R_1', s_1', t_1', c_1')$ 的映射
$s_1'$, $t_1'$ 为同构，也就是秩 $1$ 的有限局部自由态射。
此外，$U \setminus Z_1$ 的每个点都属于
$t(s^{-1}(\overline{W_1}))$。

\medskip\noindent
假设对 $r \leq n$ 已找到子概形 $W'_r \subset W_r \subset U$，使得
\begin{enumerate}
\item $W_1, \ldots, W_n$ 两两不交；
\item $W_r$ 和 $W_r'$ 为 $R$-不变；
\item 在集合论意义下，
$U \setminus Z_n \subset \bigcup_{r \leq n} t(s^{-1}(\overline{W_r}))$；
\item $W_r$ 是 $W'_r$ 的加厚；
\item 限制 $(W_r', R_r', s_r', t_r', c_r')$ 的映射
$s_r'$, $t_r'$ 有限局部自由，秩为 $r$。
\end{enumerate}
然后在集合论意义下令
$$
W_{n + 1} = Z_n \setminus
\left(
Z_{n + 1} \cup \bigcup\nolimits_{r \leq n} t(s^{-1}(\overline{W_r}))
\right)
$$
而在概形论意义下令
$$
W'_{n + 1} = Z_n \setminus
\left(
Z_{n + 1} \cup \bigcup\nolimits_{r \leq n} t(s^{-1}(\overline{W_r}))
\right)
$$
则 $W_{n + 1}$ 是 $U$ 的 $R$-不变开子概形。事实上，
$Z_{n + 1} \setminus \overline{U \setminus Z_{n + 1}}$ 在 $U$ 中开，
而由性质 (3) 及 $t$ 为闭态射这一事实，
$\overline{U \setminus Z_{n + 1}}$ 包含于我们所删去的闭子集
$\bigcup\nolimits_{r \leq n} t(s^{-1}(\overline{W_r}))$。
显然，$W'_{n + 1}$ 是 $W_{n + 1}$ 中具有相同底拓扑空间的闭子概形。
最后，性质 (1)、(2)、(3) 显然，而性质 (5) 由引理
\ref{more-groupoids-lemma-finite-stratify} 得出。

\medskip\noindent
由引理 \ref{more-groupoids-lemma-finite-stratify}，有
$\bigcap Z_r = \emptyset$。因此 $U$ 的每个点都属于
$U \setminus Z_n$，其中 $n$ 适当选取。所以在集合论意义下
$U = \bigcup_{r \geq 1} t(s^{-1}(\overline{W_r}))$，
从而 (2) 成立。于是 $W' \subset W$ 满足 (1)、(2)、(3)、(4)。
\end{proof}

\noindent
设 $(U, R, s, t, c)$ 为群胚概形。给定点 $u \in U$，
$u$ 的{\it $R$-轨道}是 $U$ 的子集 $t(s^{-1}(\{u\}))$。

\begin{lemma}
\label{more-groupoids-lemma-finite-flat-over-almost-dense-subscheme-addendum}
在引理 \ref{more-groupoids-lemma-finite-flat-over-almost-dense-subscheme} 中，
再假设 $s$ 和 $t$ 有限表现。则
\begin{enumerate}
\item 态射 $W' \to W$ 有限表现；并且
\item 若 $u \in U$ 是一点，其 $R$-轨道由 $U$ 的不可约分支的一般点
组成，则 $u \in W$。
\end{enumerate}
\end{lemma}

\begin{proof}
在这种情形，引理 \ref{more-groupoids-lemma-finite-stratify} 的分层
$U = Z_0 \supset Z_1 \supset Z_2 \supset \ldots$
由有限表现的闭浸入 $Z_k \to U$ 给出，见《除子》引理 \ref{divisors-lemma-locally-free-rank-r-pullback}。
由此立即得到 (1)，因为 $W' \to W$ 局部由开集 $W$ 与 $Z_r$
相交给出。为证明 (2)，设 $\{u_1, \ldots, u_n\}$ 为 $u$ 的轨道。
由于闭子概形 $Z_k$ 为 $R$-不变且
$\bigcap Z_k = \emptyset$，可以找到 $k$，使所有 $i$ 均满足
$u_i \in Z_k$ 且 $u_i \not \in Z_{k + 1}$。
$Z_k \to U$ 和 $Z_{k + 1} \to U$ 的像局部可构造
（《概形态射》定理 \ref{morphisms-theorem-chevalley}）。
由于 $u_i \in U$ 是 $U$ 的一个不可约分支的一般点，
存在 $u_i$ 的开邻域 $U_i$，它在集合论意义下包含于
$Z_k \setminus Z_{k + 1}$（《概形的性质》引理 \ref{properties-lemma-generic-point-in-constructible}）。
在引理 \ref{more-groupoids-lemma-finite-flat-over-almost-dense-subscheme} 的证明中，
我们把 $W$ 构造为不交并 $\coprod W_r$，其中
$W_r \subset Z_{r - 1} \setminus Z_r$，且
$U = \bigcup t(s^{-1}(\overline{W_r}))$。
由于 $\{u_1, \ldots, u_n\}$ 是一个 $R$-轨道，
若 $u \in t(s^{-1}(\overline{W_r}))$，则对某个 $i$ 有
$u_i \in \overline{W_r}$，从而
$U_i \cap W_r \not = \emptyset$，继而 $r = k$。
因此如欲证地得到 $u$ 属于
$$
W_{k + 1} = Z_k \setminus
\left(
Z_{k + 1} \cup \bigcup\nolimits_{r \leq k} t(s^{-1}(\overline{W_r}))
\right)
$$
。
\end{proof}

\begin{lemma}
\label{more-groupoids-lemma-invariant-affine-open-around-generic-point}
设 $(U, R, s, t, c)$ 为概形 $S$ 上的群胚概形。假设 $s, t$
有限且有限表现，并且 $U$ 拟分离。设
$u_1, \ldots, u_m \in U$ 为若干点，它们的轨道均由 $U$
的不可约分支的一般点组成。则存在 $R$-不变子概形
$V' \subset V \subset U$，使得
\begin{enumerate}
\item $u_1, \ldots, u_m \in V'$；
\item $V$ 在 $U$ 中开；
\item $V'$ 和 $V$ 均仿射；
\item $V' \subset V$ 是有限表现的加厚；
\item 限制 $(V', R', s', t', c')$ 的态射 $s', t'$ 有限局部自由。
\end{enumerate}
\end{lemma}

\begin{proof}
取引理 \ref{more-groupoids-lemma-finite-flat-over-almost-dense-subscheme} 中的
$W' \subset W \subset U$。由引理
\ref{more-groupoids-lemma-finite-flat-over-almost-dense-subscheme-addendum}，
有 $u_j \in W$，且 $W' \to W$ 是有限表现的加厚。
由《概形的极限》引理 \ref{limits-lemma-affines-glued-in-closed-affine}，
只需找到 $W'$ 的一个包含 $u_j$ 的 $R$-不变仿射开子概形 $V'$
（因为这时可令 $V \subset W$ 为对应的开子概形，而它将是仿射的）。
因此可以用向 $W'$ 的限制
$(W', R', s', t', c')$ 代替 $(U, R, s, t, c)$。
换言之，可以假设群胚概形 $(U, R, s, t, c)$ 的态射 $s$、$t$
有限局部自由。由《概形的性质》引理 \ref{properties-lemma-maximal-points-affine}，
可以找到一个包含 $u_1, \ldots, u_m$ 各轨道之并的仿射开集。
最后应用《群胚概形》引理 \ref{groupoids-lemma-find-invariant-affine} 即得结论。
\end{proof}

\noindent
下一个引理是引理
\ref{more-groupoids-lemma-invariant-affine-open-around-generic-point} 的特殊情形，
但在本情形论证稍为容易，因此我们重作证明
（这样可避免使用引理
\ref{more-groupoids-lemma-finite-flat-over-almost-dense-subscheme-addendum}）。

\begin{lemma}
\label{more-groupoids-lemma-invariant-affine-open-around-generic-point-Noetherian}
设 $(U, R, s, t, c)$ 为概形 $S$ 上的群胚概形。
假设 $s, t$ 有限，$U$ 局部 Noether，并且
$u_1, \ldots, u_m \in U$ 的轨道均由 $U$ 的不可约分支的一般点组成。
则存在 $R$-不变子概形 $V' \subset V \subset U$，使得
\begin{enumerate}
\item $u_1, \ldots, u_m \in V'$；
\item $V$ 在 $U$ 中开；
\item $V'$ 和 $V$ 均仿射；
\item $V' \subset V$ 是加厚；
\item 限制 $(V', R', s', t', c')$ 的态射 $s', t'$ 有限局部自由。
\end{enumerate}
\end{lemma}

\begin{proof}
设 $\{u_{j1}, \ldots, u_{jn_j}\}$ 为 $u_j$ 的轨道。
取引理 \ref{more-groupoids-lemma-finite-flat-over-almost-dense-subscheme} 中的
$W' \subset W \subset U$。由于
$U = t(s^{-1}(\overline{W}))$，可见至少有一个
$u_{ji} \in \overline{W}$。由于 $u_{ji}$ 是不可约分支的一般点，
且 $U$ 局部 Noether，这蕴含 $u_{ji} \in W$。
又因 $W$ 为 $R$-不变，得到 $u_j \in W$，事实上整个轨道都包含于 $W$。
由《概形上同调》引理 \ref{coherent-lemma-image-affine-finite-morphism-affine-Noetherian}，
只需找到 $W'$ 的一个包含 $u_1, \ldots, u_m$ 的
$R$-不变仿射开子概形 $V'$
（因为这时可令 $V \subset W$ 为对应的开子概形，而它将是仿射的）。
因此可以用向 $W'$ 的限制
$(W', R', s', t', c')$ 代替 $(U, R, s, t, c)$。
换言之，可以假设群胚概形 $(U, R, s, t, c)$ 的态射 $s$、$t$
有限局部自由。由《概形的性质》引理 \ref{properties-lemma-maximal-points-affine}，
可以找到包含 $\{u_{ij}\}$ 的仿射开集
（由《概形的性质》引理 \ref{properties-lemma-locally-Noetherian-quasi-separated}，
局部 Noether 概形拟分离）。
最后应用《群胚概形》引理 \ref{groupoids-lemma-find-invariant-affine} 即得结论。
\end{proof}

\begin{lemma}
\label{more-groupoids-lemma-find-affine-integral}
设 $(U, R, s, t, c)$ 为概形 $S$ 上的群胚概形，其中 $s, t$ 整。
设 $g : U' \to U$ 为整态射，并使 $U$ 中的每个 $R$-轨道都与
$g(U')$ 相交。设 $(U', R', s', t', c')$ 为 $R$ 向 $U'$ 的限制。
若 $u' \in U'$ 包含于一个 $R'$-不变仿射开集中，则其像
$u \in U$ 包含于 $U$ 的一个 $R$-不变仿射开集中。
\end{lemma}

\begin{proof}
设 $W' \subset U'$ 为 $R'$-不变仿射开集。令
$\tilde R = U' \times_{g, U, t} R$，并带有映射
$\text{pr}_0 : \tilde R \to U'$ 及
$h = s \circ \text{pr}_1 : \tilde R \to U$。
注意 $\text{pr}_0$ 和 $h$ 都是整态射。
因此 $\tilde W = \text{pr}_0^{-1}(W')$ 仿射。
由于 $W'$ 为 $R'$-不变，像 $W = h(\tilde W)$
在集合论意义下为 $R$-不变，并且在集合论意义下
$\tilde W = h^{-1}(W)$（略去细节）。
所以，若能证明 $W$ 为开集，则 $W$ 是概形，且态射
$\tilde W \to W$ 整且满；由《概形的极限》命题 \ref{limits-proposition-affine}，这蕴含 $W$ 仿射。
然而，每个轨道与 $U'$ 相交这一假设蕴含
$h : \tilde R \to U$ 满。整满态射为商映射
（《拓扑学》引理 \ref{topology-lemma-closed-morphism-quotient-topology} 及
《概形态射》引理 \ref{morphisms-lemma-integral-universally-closed}），
所以 $W$ 为开集。
\end{proof}

\noindent
下一个技术引理在拟仿射概形上的有限群胚情形中构造“几乎”不变函数。

\begin{lemma}
\label{more-groupoids-lemma-find-almost-invariant-function}
设 $(U, R, s, t, c)$ 为群胚概形，其中 $s, t$ 有限且有限表现。
设 $u_1, \ldots, u_m \in U$ 为若干点，其 $R$-轨道均由 $U$
的不可约分支的一般点组成。设
$j : U \to \Spec(A)$ 为浸入。设 $I \subset A$ 为理想，使得
$j(U) \cap V(I) = \emptyset$ 且
$V(I) \cup j(U)$ 在 $\Spec(A)$ 中闭。则存在 $h \in I$，使
$j^{-1}D(h)$ 是 $U$ 的一个包含 $u_1, \ldots, u_m$ 的
$R$-不变仿射开子概形。
\end{lemma}

\begin{proof}
取引理 \ref{more-groupoids-lemma-invariant-affine-open-around-generic-point} 中的
$u_1, \ldots, u_m \in V' \subset V \subset U$。
由于 $U \setminus V$ 在 $U$ 中闭，$j$ 为浸入，且
$V(I) \cup j(U)$ 在 $\Spec(A)$ 中闭，可以找到理想
$J \subset I$，使
$V(J) = V(I) \cup j(U \setminus V)$。
例如，可以取 $I$ 中在 $j(U \setminus V)$ 上消失的元素所成的理想。
因此可以用 $(V', R', s', t', c')$、
$j|_{V'} : V' \to \Spec(A)$ 以及 $J$ 分别代替
$(U, R, s, t, c)$、$j : U \to \Spec(A)$ 以及 $I$。
换言之，可以假设 $U$ 仿射且 $s$、$t$ 有限局部自由。
取任意 $f \in I$，使其在 $u_1, \ldots, u_m$ 的各个
$R$-轨道的所有点处均不消失
（《交换代数》引理 \ref{algebra-lemma-silly}）。考虑
$$
g = \text{Norm}_s(t^\sharp(j^\sharp(f))) \in \Gamma(U, \mathcal{O}_U)
$$
由于 $f \in I$，且 $V(I) \cup j(U)$ 闭，可见
$U \cap D(f) \to D(f)$ 为闭浸入。因此对某个 $n > 0$，
$f^ng$ 是某个元素 $h \in I$ 的像。我们断言 $h$ 合用。
确实，在《群胚概形》引理 \ref{groupoids-lemma-determinant-trick} 中已经看到，
$g$ 是 $R$-不变函数，故 $D(g) \subset U$ 为 $R$-不变。
由于 $f$ 在 $u_j$ 的轨道上不消失，函数 $g$ 在 $u_j$ 处不消失。
此外，$V(g) \supset V(j^\sharp(f))$，故
$j^{-1}D(h) = D(g)$。
\end{proof}

\begin{lemma}
\label{more-groupoids-lemma-no-specializations-map-to-same-point}
设 $(U, R, s, t, c)$ 为群胚概形。若 $s, t$ 有限，且
$u, u' \in R$ 是同一轨道中的不同点，则 $u'$ 不是 $u$ 的特殊化。
\end{lemma}

\begin{proof}
取 $r \in R$，使 $s(r) = u$ 且 $t(r) = u'$。
若 $u \leadsto u'$，则可以找到非平凡特殊化
$r \leadsto r'$，满足 $s(r') = u'$，见《概形》引理 \ref{schemes-lemma-quasi-compact-closed}。
令 $u'' = t(r')$。注意 $u'' \not = u'$，因为有限态射的纤维中
不存在特殊化。因此可以继续，找到非平凡特殊化
$r' \leadsto r''$，满足 $s(r'') = u''$，如此等等。
这表明 $u$ 的轨道包含特殊化的无穷序列
$u \leadsto u' \leadsto u'' \leadsto \ldots$；
但这是荒谬的，因为轨道 $t(s^{-1}(\{u\}))$ 有限。
\end{proof}

\begin{lemma}
\label{more-groupoids-lemma-get-affine}
设 $j : V \to \Spec(A)$ 为概形的拟紧浸入。
设 $f \in A$，使 $j^{-1}D(f)$ 仿射，且
$j(V) \cap V(f)$ 闭。则 $V$ 仿射。
\end{lemma}

\begin{proof}
这可由《概形态射》引理 \ref{morphisms-lemma-get-affine} 得出，但我们也给出直接证明。
令 $A' = \Gamma(V, \mathcal{O}_V)$。则
$j' : V \to \Spec(A')$ 是拟紧开浸入，见《概形的性质》引理 \ref{properties-lemma-quasi-affine}。设 $f' \in A'$ 为 $f$ 的像。
则 $(j')^{-1}D(f') = j^{-1}D(f)$ 仿射。另一方面，
$j'(V) \cap V(f')$ 是 $\Spec(A')$ 的子概形，它同
$\Spec(A)$ 的闭子概形 $j(V) \cap V(f)$ 同构。
因此它在 $\Spec(A')$ 中闭，例如可见《概形》引理 \ref{schemes-lemma-section-immersion}。
故可以用 $A'$ 代替 $A$，并假设 $j$ 为开浸入且
$A = \Gamma(V, \mathcal{O}_V)$。

\medskip\noindent
在这种情形，我们断言 $j(V) = \Spec(A)$，这将完成证明。
否则，可以找到主仿射开集 $D(g) \subset \Spec(A)$，
它同补集相交而避开闭子集 $j(V) \cap V(f)$。
注意 $j$ 把 $j^{-1}D(f)$ 同构地映到 $D(f)$，见
《概形的性质》引理 \ref{properties-lemma-invert-f-affine}。
因此 $D(g)$ 同 $V(f)$ 相交。另一方面，$j^{-1}D(g)$ 是仿射开集
$j^{-1}D(f)$ 的主开集，因而仿射。于是再次由《概形的性质》引理 \ref{properties-lemma-invert-f-affine}，可见 $D(g)$ 同构于
$j^{-1}D(g) \subset j^{-1}D(f)$，这蕴含
$D(g) \subset D(f)$。这一矛盾完成证明。
\end{proof}

\begin{lemma}
\label{more-groupoids-lemma-find-affine-codimension-1}
设 $(U, R, s, t, c)$ 为群胚概形，并设 $u \in U$。假设
\begin{enumerate}
\item $s, t$ 是有限态射；
\item $U$ 分离且局部 Noether；
\item 对 $u$ 的轨道中的每个点 $u'$，有
$\dim(\mathcal{O}_{U, u'}) \leq 1$。
\end{enumerate}
则 $u$ 包含于 $U$ 的一个 $R$-不变仿射开集中。
\end{lemma}

\begin{proof}
$u$ 的 $R$-轨道有限。由条件 (2) 和 (3)，它包含于 $U$
的一个仿射开集 $U'$，见《代数簇》命题 \ref{varieties-proposition-finite-set-of-points-of-codim-1-in-affine}。
于是 $t(s^{-1}(U \setminus U'))$ 是 $U$ 的不包含 $u$
的 $R$-不变闭子集。因此
$U \setminus t(s^{-1}(U \setminus U'))$ 是 $U'$ 中包含 $u$
的 $R$-不变开集。用这个开集代替 $U$ 后，可以假设 $U$ 拟仿射。

\medskip\noindent
由引理 \ref{more-groupoids-lemma-find-affine-integral}，可以用其约化代替 $U$，
并假设 $U$ 约化。这意味着引理
\ref{more-groupoids-lemma-finite-flat-over-almost-dense-subscheme} 的
$R$-不变子概形 $W' \subset W \subset U$ 相等，即 $W' = W$。
由于 $U = t(s^{-1}(\overline{W}))$，$u$ 的 $R$-轨道中的某点
$u'$ 包含于 $\overline{W}$；再由引理
\ref{more-groupoids-lemma-find-affine-integral}，可以用 $\overline{W}$ 代替 $U$，
并用 $u'$ 代替 $u$。因此可以假设存在稠密开 $R$-不变子概形
$W \subset U$，使限制 $(W, R_W, s_W, t_W, c_W)$ 的态射
$s_W, t_W$ 有限局部自由。

\medskip\noindent
若 $u \in W$，则由《群胚概形》引理 \ref{groupoids-lemma-find-invariant-affine} 已得结论
（因为 $W$ 拟仿射，所以 $W$ 的任意有限点集都包含于仿射开集中，
见《概形的性质》引理 \ref{properties-lemma-ample-finite-set-in-affine}）。
故假设 $u \not \in W$，于是 $u$ 的轨道中没有点属于 $W$。
设 $\xi \in U$ 是一个具有非平凡特殊化到 $u$ 的轨道中某点
$u'$ 的点。由于 $u$ 的轨道内的点之间不存在特殊化
（引理 \ref{more-groupoids-lemma-no-specializations-map-to-same-point}），
可见 $\xi$ 不在该轨道中。由假设 (3)，$\xi$ 是 $U$ 的一般点，
因而 $\xi \in W$。由于 $U$ 为 Noether 概形，这样的点只有有限多个，
记为 $\xi_1, \ldots, \xi_m \in W$。由于 $s_W, t_W$ 平坦，
每个 $\xi_j$ 的轨道都由 $W$（因而也由 $U$）的不可约分支的一般点组成。

\medskip\noindent
设 $j : U \to \Spec(A)$ 为 $U$ 到一个仿射概形的浸入
（这是可能的，因为 $U$ 拟仿射）。设 $J \subset A$ 为理想，使
$V(J) \cap j(W) = \emptyset$ 且
$V(J) \cup j(W)$ 闭。把引理
\ref{more-groupoids-lemma-find-almost-invariant-function}
应用于群胚概形 $(W, R_W, s_W, t_W, c_W)$、态射
$j|_W : W \to \Spec(A)$、诸点 $\xi_j$ 以及理想 $J$，
得到 $f \in J$，使 $(j|_W)^{-1}D(f)$ 是包含 $\xi_j$（对所有
$j$）的 $R_W$-不变仿射开集。由于 $f \in J$，可见
$j^{-1}D(f) \subset W$；也就是说，$j^{-1}D(f)$ 是
$U$ 的包含所有 $\xi_j$ 且包含于 $W$ 的 $R$-不变仿射开集。

\medskip\noindent
在
$$
U \setminus j^{-1}D(f) = j^{-1}V(f).
$$
上赋予约化诱导闭子概形结构，并记为 $Z$。
则 $Z$ 在集合论意义下为 $R$-不变
（但在概形论意义下未必为 $R$-不变）。
设 $(Z, R_Z, s_Z, t_Z, c_Z)$ 为 $R$ 向 $Z$ 的限制。
由于 $Z \to U$ 有限，可知 $s_Z$ 和 $t_Z$ 有限。
由于 $u \in Z$，$u$ 的轨道包含于 $Z$，并且它与把 $u$ 视为 $Z$
的一点所得的 $R_Z$-轨道一致。由于
$\dim(\mathcal{O}_{U, u'}) \leq 1$，且所有 $j$ 均有
$\xi_j \not \in Z$，可见对 $u$ 的轨道中的所有 $u'$，
有 $\dim(\mathcal{O}_{Z, u'}) \leq 0$。
换言之，$u$ 的 $R_Z$-轨道由 $Z$ 的不可约分支的一般点组成。

\medskip\noindent
设 $I \subset A$ 为理想，使 $V(I) \cap j(U) =\emptyset$
且 $V(I) \cup j(U)$ 闭。把引理
\ref{more-groupoids-lemma-find-almost-invariant-function} 应用于群胚概形
$(Z, R_Z, s_Z, t_Z, c_Z)$、限制 $j|_Z$、理想 $I$
以及点 $u \in Z$，得到 $h \in I$，使
$j^{-1}D(h) \cap Z$ 是包含 $u$ 的 $R_Z$-不变仿射开集。


\medskip\noindent
考虑 $R_W$-不变函数（《群胚概形》引理 \ref{groupoids-lemma-determinant-trick}）
$$
g = 
\text{Norm}_{s_W}(t_W^\sharp(j^\sharp(h)|_W)) \in \Gamma(W, \mathcal{O}_W)
$$
（下文只需要 $g$ 在 $j^{-1}D(f)$ 上的限制；在这种情形，
范数沿仿射概形之间的有限局部自由态射取得。）
我们断言
$$
V = (W_g \cap j^{-1}D(f)) \cup (j^{-1}D(h) \cap Z)
$$
是 $U$ 的 $R$-不变仿射开集，这将完成引理的证明。
由构造，它在集合论意义下为 $R$-不变。
由于 $V$ 是可构造集，要证明它开，只需证明它在 $U$ 中对一般化封闭
（《拓扑学》引理 \ref{topology-lemma-characterize-closed-Noetherian}，
或更一般的《拓扑学》引理 \ref{topology-lemma-constructible-stable-specialization-closed}）。
由于 $W_g \cap j^{-1}D(f)$ 在 $U$ 中开，只需考虑 $U$ 中的特殊化
$u_1 \leadsto u_2$，其中 $u_2 \in j^{-1}D(h) \cap Z$。
这意味着 $h$ 在 $j(u_2)$ 中非零，且 $u_2 \in Z$。
若 $u_1 \in Z$，则 $j(u_1) \leadsto j(u_2)$；由于 $h$
在 $j(u_2)$ 中非零，它在 $j(u_1)$ 中也非零，从而 $u_1 \in V$。
若 $u_1 \not \in Z$ 且也不属于
$W_g \cap j^{-1}D(f)$，则 $u_1 \in W$、$u_1 \not \in W_g$，
因为 $Z = j^{-1}V(f)$ 的补集包含于 $W \cap j^{-1}D(f)$。
因此存在点 $r_1 \in R$，使 $s(r_1) = u_1$，且 $h$
在 $t(r_1)$ 中为零。由于 $s$ 有限，可以找到特殊化
$r_1 \leadsto r_2$，满足 $s(r_2) = u_2$。
但这时可知 $h$ 在 $u'_2 = t(r_2)$ 中为零，
这同 $j^{-1}D(h) \cap Z$ 为 $R$-不变且 $u_2$ 属于其中相矛盾。
因此 $V$ 为开集。


\medskip\noindent
注意，对于函数 $h \in I$，有 $V \subset j^{-1}D(h)$。
因此得到浸入
$$
j' : V \longrightarrow \Spec(A_h)
$$
设 $f' \in A_h$ 为 $f$ 的像。则 $(j')^{-1}D(f')$
是 $U$ 的仿射开集 $j^{-1}D(f)$ 中由 $g$ 确定的主开集。
因此 $(j')^{-1}D(f)$ 仿射。最后，由 $h \in I$ 及理想 $I$ 的选取，
$j'(V) \cap V(f') = j'(j^{-1}D(h) \cap Z)$
在
$\Spec(A_h/(f')) = \Spec((A/f)_h) = D(h) \cap V(f)$ 中闭。
故可应用引理 \ref{more-groupoids-lemma-get-affine}，得出如上所断言的 $V$ 仿射。
\end{proof}







\section{ind-拟仿射态射的下降}
\label{more-groupoids-section-ind-quasi-affine}

\noindent
ind-拟仿射态射定义于《态射进阶》第 \ref{more-morphisms-section-ind-quasi-affine} 节。
本节是《下降》第 \ref{descent-section-quasi-affine} 节对 ind-拟仿射态射的对应版本。

\medskip\noindent
设 $X$ 为拟分离概形。设 $E \subset X$ 为 $X$ 的某个非空族拟紧开集
之交。写成 $E = \bigcap_{i \in I} U_i$，其中
$U_i \subset X$ 是拟紧开集且 $I$ 非空。
加入有限交后，可以假设对 $i, j \in I$，存在 $k \in I$，使
$U_k \subset U_i \cap U_j$。在这种情形，对定义于 $X$ 上的任意层
$\mathcal{F}$，有
\begin{equation}
\label{more-groupoids-equation-sections-of-intersection}
\Gamma(E, \mathcal{F}|_E) = \colim \Gamma(U_i, \mathcal{F}|_{U_i})
\end{equation}
确实，固定 $i_0 \in I$，并分别用 $U_{i_0}$ 和
$\{i \in I \mid U_i \subset U_{i_0}\}$ 代替 $X$ 和 $I$。
于是 $X$ 拟紧且拟分离，因而是谱空间，见《概形的性质》引理 \ref{properties-lemma-quasi-compact-quasi-separated-spectral}。
再由《拓扑学》引理 \ref{topology-lemma-make-spectral-space} 及《空间上的层》引理 \ref{sheaves-lemma-descend-opens} 可见该等式成立。
（事实上，若 $\mathcal{F}$ 为 Abel 层，该公式对高阶上同调群也成立；
见《层上同调》引理 \ref{cohomology-lemma-colimit}。）


\begin{lemma}
\label{more-groupoids-lemma-sits-in-functions}
设 $X$ 为 ind-拟仿射概形。设 $E \subset X$ 为 $X$
的某个非空族拟紧开集之交。令
$A = \Gamma(E, \mathcal{O}_X|_E)$、$Y = \Spec(A)$。
则《概形》引理 \ref{schemes-lemma-morphism-into-affine} 的典范态射
$$
j : (E, \mathcal{O}_X|_E) \longrightarrow (Y, \mathcal{O}_Y)
$$
确定同构
$(E, \mathcal{O}_X|_E) \to (E', \mathcal{O}_Y|_{E'})$，
其中 $E' \subset Y$ 是若干拟紧开集之交。
若 $W \subset E$ 在 $X$ 中开，则 $j(W)$ 在 $Y$ 中开。
\end{lemma}

\begin{proof}
注意 $(E, \mathcal{O}_X|_E)$ 是局部环化空间，所以《概形》引理 \ref{schemes-lemma-morphism-into-affine}
可应用于 $A \to \Gamma(E, \mathcal{O}_X|_E)$。
写成 $E = \bigcap_{i \in I} U_i$，其中 $I \not = \emptyset$，
且 $U_i \subset X$ 为拟紧开集。可以且确实假设，对
$i, j \in I$，存在 $k \in I$，使 $U_k \subset U_i \cap U_j$。
令 $A_i = \Gamma(U_i, \mathcal{O}_{U_i})$。得到交换图
$$
\xymatrix{
(E, \mathcal{O}_X|_E) \ar[r] \ar[d] &
(\Spec(A), \mathcal{O}_{\Spec(A)}) \ar[d] \\
(U_i, \mathcal{O}_{U_i}) \ar[r] &
(\Spec(A_i), \mathcal{O}_{\Spec(A_i)})
}
$$
由于 $U_i$ 拟仿射，可见 $U_i \to \Spec(A_i)$ 是拟紧开浸入。
另一方面，$A = \colim A_i$。因此作为拓扑空间，
$\Spec(A) = \lim \Spec(A_i)$（《概形的极限》引理 \ref{limits-lemma-topology-limit}）。由于
$E = \lim U_i$（由《拓扑学》引理 \ref{topology-lemma-make-spectral-space}），可见
$E \to \Spec(A)$ 是到其像 $E'$ 的同胚，而且 $E'$ 是
$\Spec(A)$ 中各开集 $U_i \subset \Spec(A_i)$ 的逆像之交。
对任意 $e \in E$，局部环 $\mathcal{O}_{X, e}$
是 $\mathcal{O}_{U_i, e}$ 的值，而后者与 $\Spec(A)$ 上的值相同。

\medskip\noindent
为证明引理的最后一个断言，论证如下。选取
$i, j \in I$，使 $U_i \subset U_j$。考虑交换图
$$
\xymatrix{
U_i \ar[r] \ar[d] & \Spec(A_i) \ar[d] \\
U_i \ar[r] & \Spec(A_j)
}
\quad\quad
\xymatrix{
W \ar[r] \ar[d] & \Spec(A_i) \ar[d] \\
W \ar[r] & \Spec(A_j)
}
\quad\quad
\xymatrix{
W \ar[r] \ar[d] & \Spec(A) \ar[d] \\
W \ar[r] & \Spec(A_j)
}
$$
由《概形的性质》引理 \ref{properties-lemma-cartesian-diagram-quasi-affine}，
第一个图为 Cartesian 方块。因此第二个图也为 Cartesian 方块。
取极限后，得到第三个图为 Cartesian 方块，所以该图的顶部水平箭头
是开浸入。
\end{proof}

\begin{lemma}
\label{more-groupoids-lemma-affine-base-change}
设给定概形的 Cartesian 方块
$$
\xymatrix{
X \ar[d]_f \ar[r] & \Spec(B) \ar[d] \\
Y \ar[r] & \Spec(A)
}
$$
设 $E \subset Y$ 为 $Y$ 的某个非空族拟紧开集之交。若 $Y$ 拟分离，
且 $A \to B$ 平坦，则
$$
\Gamma(f^{-1}(E), \mathcal{O}_X|_{f^{-1}(E)}) =
\Gamma(E, \mathcal{O}_Y|_E) \otimes_A B
$$
\end{lemma}

\begin{proof}
写成 $E = \bigcap_{i \in I} V_i$，其中 $V_i \subset Y$ 为拟紧开集。
可以且确实假设，对 $i, j \in I$，存在 $k \in I$，使
$V_k \subset V_i \cap V_j$。于是类似地，在 $X$ 中有
$f^{-1}(E) = \bigcap_{i \in I} f^{-1}(V_i)$。
所以结论由等式 (\ref{more-groupoids-equation-sections-of-intersection})
及其对 $V_i$、$f^{-1}(V_i)$ 的相应结论得出；后者即
《概形上同调》引理 \ref{coherent-lemma-flat-base-change-cohomology}。
\end{proof}

\begin{lemma}[Gabber]
\label{more-groupoids-lemma-ind-quasi-affine}
设 $S$ 为概形。设 $\{X_i \to S\}_{i\in I}$ 为 fpqc 覆盖。
设 $(V_i/X_i, \varphi_{ij})$ 为关于
$\{X_i \to S\}$ 的下降数据，见《下降》定义 \ref{descent-definition-descent-datum-for-family-of-morphisms}。
若每个态射 $V_i \to X_i$ 都 ind-拟仿射，则该下降数据有效。
\end{lemma}

\begin{proof}
ind-拟仿射是概形态射的一项性质，且在任意基变换下保持，见
《态射进阶》引理 \ref{more-morphisms-lemma-base-change-ind-quasi-affine}。
因此《下降》引理 \ref{descent-lemma-descending-types-morphisms} 适用；
只需在 fpqc 覆盖由单个仿射概形之间的平坦满态射
$\{X \to S\}$ 给出时证明该引理。写成
$X = \Spec(A)$、$S = \Spec(R)$，于是
$R \to A$ 是忠实平坦环同态。设 $(V, \varphi)$ 为 $X$ 在 $S$
上的下降数据，并假设 $V \to X$ ind-拟仿射；换言之，
$V$ ind-拟仿射。

\medskip\noindent
设 $(U, R, s, t, c)$ 为 $S$ 上的群胚概形，其中
$U = X$、$R = X \times_S X$，而 $s$、$t$、$c$ 取通常的定义。
由《群胚概形》引理 \ref{groupoids-lemma-cartesian-equivalent-descent-datum}，
二元组 $(V, \varphi)$ 对应于群胚概形的 Cartesian 态射
$(U', R', s', t', c') \to (U, R, s, t, c)$。
任取一点 $u' \in U'$。由《群胚概形》引理 \ref{groupoids-lemma-constructing-invariant-opens}、
\ref{groupoids-lemma-first-observation} 和
\ref{groupoids-lemma-second-observation}，
可以选取 $u' \in W \subset E \subset U'$，其中
$W$ 开且为 $R'$-不变，而 $E$ 在集合论意义下为 $R'$-不变，
并且是某个非空族拟紧开集之交。

\medskip\noindent
转回 $(V, \varphi)$ 的语言，对任意 $v \in V$，可以找到
$v \in W \subset E \subset V$，满足以下性质：
(a) $W$ 开，且 $\varphi(W \times_S X) = X \times_S W$；
(b) $E$ 是若干拟紧开集之交，且在集合论意义下
$\varphi(E \times_S X) = X \times_S E$。
这里用 $E \times_S X$ 表示 $E$ 在 $V \times_S X$
中关于投影态射的逆像，并对 $X \times_S E$ 使用类似记号。
由引理 \ref{more-groupoids-lemma-affine-base-change}，这蕴含
$\varphi$ 定义 $A \otimes_R A$-代数的同构
\begin{align*}
\Gamma(E, \mathcal{O}_V|_E) \otimes_R A
& =
\Gamma(E \times_S X, \mathcal{O}_{V \times_S X}|_{E \times_S X}) \\
& \to
\Gamma(X \times_S E, \mathcal{O}_{X \times_S V}|_{X \times_S E}) \\
& =
A \otimes_R \Gamma(E, \mathcal{O}_V|_E)
\end{align*}
称之为 $\psi$。$\varphi$ 的余圈条件转化为
$\psi$ 的余圈条件，如《下降》定义 \ref{descent-definition-descent-datum-modules} 所述
（略去细节）。由《下降》命题 \ref{descent-proposition-descent-module}，得到 $R$-代数 $R'$
及 $A$-代数同构
$\chi : R' \otimes_R A \to \Gamma(E, \mathcal{O}_V|_E)$，
它与 $\psi$ 以及 $R' \otimes_R A$ 上的典范下降数据相容。

\medskip\noindent
由引理 \ref{more-groupoids-lemma-sits-in-functions}，得到局部环化空间的典范“嵌入”
$$
j : (E, \mathcal{O}_V|_E) \longrightarrow
\Spec(\Gamma(E, \mathcal{O}_V|_E)) = \Spec(R' \otimes_R A)
$$
该映射的构造是典范的，并得到交换图
$$
\xymatrix{
& E \times_S X \ar[rr]_\varphi \ar[ld] \ar[rd]^{j'} & &
X \times_S E \ar[rd] \ar[ld]_{j''} \\
E \ar[rd]^j  & &
\Spec(R' \otimes_R A \otimes_R A) \ar[ld] \ar[rd] & &
E \ar[ld]_j \\
& \Spec(R' \otimes_R A) \ar[rd] && \Spec(R' \otimes_R A) \ar[ld] \\
& & \Spec(R')
}
$$
其中 $j'$ 和 $j''$ 来自把同一构造分别应用于
$E \times_S X \subset V \times_S X$ 和
$X \times_S E \subset X \times_S V$，并使用 $\chi$
以及构造 $\psi$ 时所用的等同。于是 $j(W)$ 是
$\Spec(R' \otimes_R A)$ 的开子概形，而且它在两个投影
$\Spec(R' \otimes_R A \otimes_R A) \to \Spec(R' \otimes_R A)$
下的逆像相等。由《下降》引理 \ref{descent-lemma-open-fpqc-covering}，得到开集
$W_0 \subset \Spec(R')$，其向 $\Spec(A)$ 的基变换为 $j(W)$。
考察上图，可见下降数据
$(W, \varphi|_{W \times_S X})$ 有效。
由《下降》引理 \ref{descent-lemma-effective-for-fpqc-is-local-upstairs}，
可见原下降数据有效。
\end{proof}
% Generated mechanically from the frozen etale.tex target.
% Do not edit this generated file; edit the bound locale source instead.

% OK, start here.
%

\setcounter{chapter}{40}
\chapter{概形的平展（\'Etale）态射}


\phantomsection
\label{etale-section-phantom}




\section{引言}
\label{etale-section-introduction}

\noindent
本章讨论概形的平展（\'etale）态射。我们通过研究诺特情形来阐明其中若干较为重要的概念。
主要目标是为读者汇集足够的交换代数结果，使其能够开始阅读平展（\'etale）上同调的专著。
另一个辅助目标，则是给出足够的证据，使那些竟把“概形的平展（\'etale）拓扑”称作一般性
胡言乱语的读者停止这种亵渎。

\medskip\noindent
对于代数几何中的标准结果（包括概形与交换代数方面的结果），我们将引用
Stacks project 的其他章节。对于本章陈述的新结果，我们会给出详细证明。




\section{约定}
\label{etale-section-conventions}

\noindent
本章常常假设概形为局部诺特概形，也常常假设环为诺特环。不过，在每一项陈述中，
我们都会在必要时重申这些假设，并确保列出全部条件。另一方面，下面是一些经常
使用、值得牢记的一般事实：
\begin{enumerate}
\item 若 $A$ 为诺特环，则有限型环同态 $A \to B$ 是有限呈示的。见《交换代数》引理 \ref{algebra-lemma-Noetherian-finite-type-is-finite-presentation}。
\item 局部诺特概形之间（局部）有限型的态射自动是（局部）有限呈示的。
见《概形态射》引理 \ref{morphisms-lemma-noetherian-finite-type-finite-presentation}。
\item 此处还应增列更多类似事实。
\end{enumerate}




\section{非分歧态射}
\label{etale-section-unramified-definition}

\noindent
我们先为诺特局部环定义“局部环的非分歧同态”。这里不能简单沿用“非分歧”
一词，因为《交换代数》第 \ref{algebra-section-unramified} 节中已经有非分歧
环同态的概念，而两者并不相同。稍作讨论之后，我们将把这一概念整体化，
以描述局部诺特概形之间的非分歧态射。

\begin{definition}
\label{etale-definition-unramified-rings}
设 $A$、$B$ 为诺特局部环。若局部同态 $A \to B$ 满足
\begin{enumerate}
\item $\mathfrak m_AB = \mathfrak m_B$；
\item $\kappa(\mathfrak m_B)$ 是 $\kappa(\mathfrak m_A)$ 的有限可分扩张；且
\item $B$ 本质上是 $A$ 上的有限型代数（即 $B$ 是某个有限型
$A$-代数在一个素理想处的局部化），
\end{enumerate}
则称它为{\it 局部环的非分歧同态}。
\end{definition}

\noindent
这是《交换代数》第 \ref{algebra-section-unramified} 节中定义的局部版本。
该节把环同态 $R \to S$ 定义为非分歧的，当且仅当它是有限型的且
$\Omega_{S/R} = 0$。若存在 $g \in S$、$g \not \in \mathfrak q$，
使得 $R \to S_g$ 是非分歧环同态，则称 $R \to S$ 在素理想
$\mathfrak q \subset S$ 处非分歧。《交换代数》引理 \ref{algebra-lemma-unramified-at-prime} 和
\ref{algebra-lemma-characterize-unramified} 证明了：给定有限型环同态
$R \to S$，以及 $S$ 的素理想 $\mathfrak q$，它位于
$\mathfrak p \subset R$ 之上，则有
$$
R \to S\text{ 在下列点非分歧： }\mathfrak q
\Leftrightarrow
\mathfrak pS_{\mathfrak q} = \mathfrak q S_{\mathfrak q}
\text{ 且 }
\kappa(\mathfrak p) \subset \kappa(\mathfrak q)\text{ 有限可分}
$$
因此，对于局部环之间的局部同态，上述定义中的性质与非分歧性密切相关。
事实上，我们已经证明了下述引理。

\begin{lemma}
\label{etale-lemma-characterize-unramified-Noetherian}
\begin{slogan}
非分歧性是茎上的局部条件。
\end{slogan}
设 $A$ 为诺特环，且 $A \to B$ 为有限型同态。设 $\mathfrak q$ 是
$B$ 的素理想，并位于 $\mathfrak p \subset A$ 之上。则
$A \to B$ 在 $\mathfrak q$ 处非分歧，当且仅当
$A_{\mathfrak p} \to B_{\mathfrak q}$ 是局部环的非分歧同态。
\end{lemma}

\begin{proof}
见上面的讨论。
\end{proof}

\noindent
下面用完备化刻画非分歧性。对于诺特局部环 $A$，以 $A^\wedge$
表示 $A$ 关于其极大理想的完备化。它仍是诺特局部环，见《交换代数》引理 \ref{algebra-lemma-completion-Noetherian-Noetherian}。

\begin{lemma}
\label{etale-lemma-unramified-completions}
设 $A$、$B$ 为诺特局部环，且 $A \to B$ 为局部同态。
\begin{enumerate}
\item 若 $A \to B$ 是局部环的非分歧同态，则 $B^\wedge$ 是有限
$A^\wedge$-模；
\item 若 $A \to B$ 是局部环的非分歧同态，且
$\kappa(\mathfrak m_A) = \kappa(\mathfrak m_B)$，则
$A^\wedge \to B^\wedge$ 满射；
\item 若 $A \to B$ 是局部环的非分歧同态，且
$\kappa(\mathfrak m_A)$ 可分闭，则 $A^\wedge \to B^\wedge$ 满射；
\item 若 $A$ 和 $B$ 是完备离散赋值环，则 $A \to B$ 是局部环的
非分歧同态，当且仅当 $A$ 的一致化参数映为 $B$ 的一致化参数，
剩余域扩张为有限可分扩张（并且 $B$ 本质上是 $A$ 上的有限型代数）。
\end{enumerate}
\end{lemma}

\begin{proof}
(1) 是《交换代数》引理 \ref{algebra-lemma-finite-after-completion}
的特例。对于 (2)，注意 $\kappa(\mathfrak m_A)$-向量空间
$B^\wedge/\mathfrak m_{A^\wedge}B^\wedge$ 由 $1$ 生成。
因此，由 Nakayama 引理（《交换代数》引理 \ref{algebra-lemma-NAK}），
映射 $A^\wedge \to B^\wedge$ 满射。(3) 是 (2) 的特例。
(4) 直接由定义得到。
\end{proof}

\begin{lemma}
\label{etale-lemma-characterize-unramified-completions}
设 $A$、$B$ 为诺特局部环，并设 $A \to B$ 为局部同态，其中
$B$ 本质上是 $A$ 上的有限型代数。下列条件等价：
\begin{enumerate}
\item $A \to B$ 是局部环的非分歧同态；
\item $A^\wedge \to B^\wedge$ 是局部环的非分歧同态；且
\item $A^\wedge \to B^\wedge$ 是非分歧环同态。
\end{enumerate}
\end{lemma}

\begin{proof}
(1) 与 (2) 的等价性来自如下事实：$\mathfrak m_AA^\wedge$ 是
$A^\wedge$ 的极大理想（$B$ 也类似），而 $B \to B^\wedge$ 忠实平坦。
例如，若 $A^\wedge \to B^\wedge$ 非分歧，则
$\mathfrak m_AB^\wedge = (\mathfrak m_AB)B^\wedge = \mathfrak m_BB^\wedge$，
因而 $\mathfrak m_AB = \mathfrak m_B$。

\medskip\noindent
假设等价条件 (1) 和 (2) 成立。由引理 \ref{etale-lemma-unramified-completions}，可知 $A^\wedge \to B^\wedge$
是有限的。因此 $A^\wedge \to B^\wedge$ 是有限呈示的；再由
《交换代数》引理 \ref{algebra-lemma-characterize-unramified}，可知
$A^\wedge \to B^\wedge$ 在 $\mathfrak m_{B^\wedge}$ 处非分歧。
由于 $B^\wedge$ 是局部环，可知 $A^\wedge \to B^\wedge$ 非分歧。

\medskip\noindent
假设 (3) 成立。由《交换代数》引理 \ref{algebra-lemma-unramified-at-prime}，$A^\wedge \to B^\wedge$
是局部环的非分歧同态，即 (2) 成立。
\end{proof}

\begin{definition}
\label{etale-definition-unramified-schemes}
（一般情形的定义见《概形态射》定义 \ref{morphisms-definition-unramified}。）
设 $Y$ 为局部诺特概形，$f : X \to Y$ 为局部有限型态射，并设
$x \in X$。
\begin{enumerate}
\item 若 $\mathcal{O}_{Y, f(x)} \to \mathcal{O}_{X, x}$ 是局部环的
非分歧同态，则称 $f$ {\it 在 $x$ 处非分歧}。
\item 若这一态射在 $X$ 的每一点都非分歧，则称态射
$f : X \to Y$ {\it 非分歧}。
\end{enumerate}
\end{definition}

\noindent
下面证明此定义与概形态射一章中的定义一致。特别地，这保证了一个态射的
非分歧点所成集合是开集。

\begin{lemma}
\label{etale-lemma-unramified-definition}
设 $Y$ 为局部诺特概形，$f : X \to Y$ 为局部有限型态射，并设
$x \in X$。则 $f$ 按定义 \ref{etale-definition-unramified-schemes} 在 $x$ 处非分歧，当且仅当它按
《概形态射》定义 \ref{morphisms-definition-unramified} 非分歧。
\end{lemma}

\begin{proof}
这由引理 \ref{etale-lemma-characterize-unramified-Noetherian} 和定义得到。
\end{proof}

\noindent
下面列出非分歧态射的一些结果。此处的表述仅适用于局部诺特概形之间的
局部有限型态射。每一项都引用本项目先前证明的一般结果，不过其中一些结果
在诺特情形下可以更容易地证明：
\begin{enumerate}
\item 非分歧性在源和目标上对于 Zariski 拓扑都是局部的。
\item 非分歧态射在基变换与复合下稳定。见《概形态射》引理 \ref{morphisms-lemma-base-change-unramified} 和
\ref{morphisms-lemma-composition-unramified}。
\item 概形的非分歧态射是局部拟有限的；拟紧的非分歧态射是拟有限的。
见《概形态射》引理 \ref{morphisms-lemma-unramified-quasi-finite}。
\item 非分歧态射的相对维数为 $0$。见《概形态射》定义 \ref{morphisms-definition-relative-dimension-d} 和《概形态射》引理 \ref{morphisms-lemma-locally-quasi-finite-rel-dimension-0}。
\item 一个态射非分歧，当且仅当它的所有纤维均非分歧。换言之，
可以在概形论纤维上检验非分歧性。见《概形态射》引理 \ref{morphisms-lemma-unramified-etale-fibres}。
\item 设 $X$ 和 $Y$ 在基概形 $S$ 上非分歧。则从 $X$ 到 $Y$ 的任意
$S$-态射都非分歧。见《概形态射》引理 \ref{morphisms-lemma-unramified-permanence}。
\end{enumerate}

\section{非分歧态射的另外三种刻画}
\label{etale-section-three-other}

\noindent
下述定理给出态射在一点处非分歧的三种等价刻画。一般概形情形的（部分）陈述见
《概形态射》引理 \ref{morphisms-lemma-unramified-at-point}。

\begin{theorem}
\label{etale-theorem-unramified-equivalence}
设 $Y$ 为局部诺特概形，$f : X \to Y$ 为局部有限型概形态射，并设
$x$ 为 $X$ 的一点。下列条件等价：
\begin{enumerate}
\item $f$ 在 $x$ 处非分歧；
\item 相对微分模在 $x$ 处的茎 $\Omega_{X/Y, x}$ 为零；
\item 存在 $x$ 的开邻域 $U$、$f(x)$ 的开邻域 $V$，以及交换图
$$
\xymatrix{
U \ar[rr]_i \ar[rd] & & \mathbf{A}^n_V \ar[ld] \\
& V
}
$$
其中 $i$ 是由拟凝聚理想层 $\mathcal{I}$ 定义的闭浸入，并且
$g \in \mathcal{I}_{i(x)}$ 的微分 $\text{d}g$ 生成
$\Omega_{\mathbf{A}^n_V/V, i(x)}$；且
\item 对角态射 $\Delta_{X/Y} : X \to X \times_Y X$ 在 $x$ 处为局部同构。
\end{enumerate}
\end{theorem}

\begin{proof}
(1) 与 (2) 的等价性已在《概形态射》引理 \ref{morphisms-lemma-unramified-at-point} 中证明。

\medskip\noindent
若 $f$ 在 $x$ 处非分歧，则 $f$ 在 $x$ 的某个开邻域上非分歧；
这一点不能直接由本章定义 \ref{etale-definition-unramified-schemes}
推出，但可以由《概形态射》定义 \ref{morphisms-definition-unramified} 推出，而我们已在引理 \ref{etale-lemma-unramified-definition} 中证明两种定义等价。选取仿射开集
$V \subset Y$、$U \subset X$，使得 $f(U) \subset V$、$x \in U$，
并使 $f$ 在 $U$ 上非分歧，即 $f|_U : U \to V$ 非分歧。
由《概形态射》引理 \ref{morphisms-lemma-diagonal-unramified-morphism}，
态射 $U \to U \times_V U$ 是开浸入。这证明了 (1) 蕴含 (4)。

\medskip\noindent
若 $\Delta_{X/Y}$ 在 $x$ 处为局部同构，则由《概形态射》引理 \ref{morphisms-lemma-differentials-diagonal}，
$\Omega_{X/Y, x} = 0$。因此 (4) 蕴含 (2)。至此，(1)、(2) 与 (4)
彼此等价。

\medskip\noindent
假设 (3) 成立。图表条件与《概形态射》引理 \ref{morphisms-lemma-differentials-relative-immersion} 一起表明
$\Omega_{U/V, x} = 0$。由于 $\Omega_{U/V, x} = \Omega_{X/Y, x}$，
可知 (2) 成立。

\medskip\noindent
最后假设 (2) 成立。为证明 (3)，可在 $X$ 和 $Y$ 上局部化，并假设
$X$ 和 $Y$ 均为仿射概形。记 $X = \Spec(B)$、$Y = \Spec(A)$。
点 $x \in X$ 对应于素理想 $\mathfrak q \subset B$。我们的假设为
$\Omega_{B/A, \mathfrak q} = 0$（概形上的微分与交换代数中的微分模
之间的关系见《概形态射》引理 \ref{morphisms-lemma-differentials-affine}）。由于 $Y$ 局部诺特且
$f$ 局部有限型，可知 $A$ 是诺特环，并且
$B \cong A[x_1, \ldots, x_n]/(f_1, \ldots, f_m)$；见《概形的性质》引理 \ref{properties-lemma-locally-Noetherian} 和《概形态射》引理 \ref{morphisms-lemma-locally-finite-type-characterize}。
特别地，$\Omega_{B/A}$ 是有限 $B$-模。因此可以找到一个
$g \in B$、$g \not \in \mathfrak q$，使得主局部化
$(\Omega_{B/A})_g$ 为零。于是用 $B_g$ 代替 $B$ 后，有
$\Omega_{B/A} = 0$（微分模的形成与局部化交换，见《交换代数》引理 \ref{algebra-lemma-differentials-localize}）。这意味着
$\text{d}(f_j)$ 生成典范映射
$\Omega_{A[x_1, \ldots, x_n]/A} \otimes_A B \to \Omega_{B/A}$
的核。因此，$A$-代数的满射 $A[x_1, \ldots, x_n] \to B$
给出 (3) 中的交换图，定理得证。
\end{proof}

\noindent
这个定理如何使用？下面给出几点说明：
\begin{enumerate}
\item 设 $f : X \to Y$ 和 $g : Y \to Z$ 是局部诺特概形之间的两个
局部有限型态射。存在典范短正合列
$$
f^*(\Omega_{Y/Z}) \to \Omega_{X/Z} \to \Omega_{X/Y} \to 0
$$
见《概形态射》引理 \ref{morphisms-lemma-triangle-differentials}。
因此，该定理说明：若 $g \circ f$ 非分歧，则 $f$ 也非分歧。
这正是《概形态射》引理 \ref{morphisms-lemma-unramified-permanence}。
\item 由于 $\Omega_{X/Y}$ 同构于对角态射的余法层
（《概形态射》引理 \ref{morphisms-lemma-differentials-diagonal}），
若 $X \to Y$ 是局部诺特概形之间的局部有限型单态射，则
$X \to Y$ 非分歧。特别地，局部诺特概形的开浸入和闭浸入均非分歧。
见《概形态射》引理 \ref{morphisms-lemma-open-immersion-unramified} 和
\ref{morphisms-lemma-closed-immersion-unramified}。
\item 该定理还说明，局部诺特概形之间的局部有限型态射
$f : X \to Y$ 的非非分歧点集合，正是凝聚层 $\Omega_{X/Y}$ 的支集。
这使我们能够给“分歧轨迹”给出概形论定义。
\end{enumerate}

\section{非分歧态射的函子性刻画}
\label{etale-section-functorial-unramified}

\noindent
在基础代数几何中，我们知道某些态射类可以通过函子性质来刻画，
而且这类描述十分有用。非分歧态射也有这样的刻画。

\begin{theorem}
\label{etale-theorem-formally-unramified}
设 $f : X \to S$ 为概形态射。假设 $S$ 是局部诺特概形，且 $f$
局部有限型。则下列条件等价：
\begin{enumerate}
\item $f$ 非分歧；
\item 态射 $f$ 形式非分歧：对于任意仿射 $S$-概形 $T$，以及由一个
平方为零的理想定义的 $T$ 的子概形 $T_0$，自然映射
$$
\Hom_S(T, X) \longrightarrow \Hom_S(T_0, X)
$$
是单射。
\end{enumerate}
\end{theorem}

\begin{proof}
更一般的陈述与证明见《态射进阶》引理 \ref{more-morphisms-lemma-unramified-formally-unramified}。
下面概述当前情形的证明。

\medskip\noindent
首先检验这两个性质在源和目标上都是局部的。因此可假设 $S$ 和 $X$
均为仿射概形。记 $X = \Spec(B)$、$S = \Spec(R)$，并记
$T = \Spec(C)$。设 $J$ 为 $C$ 的平方为零的理想，且
$T_0 = \Spec(C/J)$。假设给定图表
$$
\xymatrix{
& B \ar[d]^\phi \ar[rd]^{\bar{\phi}}
& \\
R \ar[r] \ar[ur] & C \ar[r]
& C/J
}
$$
其次，检验对应 $\phi' \mapsto \phi' - \phi$ 给出
$\bar{\phi}$ 的提升集合与模 $\text{Der}_R(B, J)$ 之间的双射。
由非分歧态射的微分模为零这一刻画，可得 (1) $\Rightarrow$ (2)；
见定理 \ref{etale-theorem-unramified-equivalence}。

\medskip\noindent
为得到反向蕴含，考虑由平方为零的理想 $J = I/I^2$ 定义的满射
$q : C = (B \otimes_R B)/I^2 \to B = C/J$，其中 $I$ 是乘法映射
$B \otimes_R B \to B$ 的核。我们已有一个提升 $B \to C$，例如由
$b \mapsto b \otimes 1$ 给出。于是，按照上面的同一论证，
$\text{id} : B \to C/J$ 的提升与 $\text{Der}_R(B, J)$ 之间存在双射。
因此假设蕴含后一模为零。但我们知道 $J \cong \Omega_{B/R}$。
所以 $B/R$ 非分歧。
\end{proof}



\section{非分歧态射的拓扑性质}
\label{etale-section-topological-unramified}

\noindent
第一个对我们有用的拓扑结果是：非分歧且分离的态射具有性质良好的截面。
本节材料不需要任何诺特性假设。

\begin{proposition}
\label{etale-proposition-properties-sections}
非分歧态射的截面。
\begin{enumerate}
\item 非分歧态射的任意截面都是开浸入。
\item 分离态射的任意截面都是闭浸入。
\item 非分歧分离态射的任意截面既开又闭。
\end{enumerate}
\end{proposition}

\begin{proof}
固定基概形 $S$。若 $f : X' \to X$ 是任意 $S$-态射，则图态射
$\Gamma_f : X' \to X' \times_S X$ 是对角态射
$\Delta_{X/S} : X \to X \times_S X$ 沿投影
$X' \times_S X \to X \times_S X$ 的基变换。若
$g : X \to S$ 分离（相应地，\ 非分歧），则由《概形》定义 \ref{schemes-definition-separated}（相应地，由《概形态射》引理 \ref{morphisms-lemma-diagonal-unramified-morphism}），对角态射是闭浸入
（相应地，开浸入）。因此，作为基变换，图态射也是如此（见《概形》引理 \ref{schemes-lemma-base-change-immersion}）。在特例
$X' = S$ 中，分别得到 (1)，相应地\ 得到 (2)。结合 (1) 和 (2) 即得 (3)。
\end{proof}

\noindent
现在可以明确描述非分歧态射的截面。

\begin{theorem}
\label{etale-theorem-sections-unramified-maps}
设 $Y$ 为连通概形，且 $f : X \to Y$ 非分歧并分离。
$f$ 的每个截面都给出到某个连通分支的同构。存在双射对应
$$
\text{下列对象的截面： }f
\leftrightarrow
\left\{
\begin{matrix}
\text{连通分支 }X'\text{ 属于 }X\text{ 使得}\\
\text{诱导映射 }X' \to Y\text{ 为同构}
\end{matrix}
\right\}
$$
特别地，给定 $x \in X$，至多存在一个经过 $x$ 的截面。
\end{theorem}

\begin{proof}
直接由命题 \ref{etale-proposition-properties-sections} 的 (3) 得到。
\end{proof}

\noindent
上一定理让我们初步看到非分歧态射的“刚性”。下述命题提供了进一步说明；
它本身很有意义，并且也用于代数基本群理论（见
\cite[Expos\'e V]{SGA1}）。另见更一般的《概形态射》引理 \ref{morphisms-lemma-value-at-one-point}。

\begin{proposition}
\label{etale-proposition-equality}
设 $S$ 为概形。设 $\pi : X \to S$ 非分歧并分离。设 $Y$ 为
$S$-概形，$y \in Y$ 为一点。设 $f, g : Y \to X$ 为两个
$S$-态射，并假设
\begin{enumerate}
\item $Y$ 连通；
\item $x = f(y) = g(y)$；且
\item 剩余域上的诱导映射
$f^\sharp, g^\sharp : \kappa(x) \to \kappa(y)$ 相等。
\end{enumerate}
则 $f = g$。
\end{proposition}

\begin{proof}
映射 $f, g : Y \to X$ 定义了映射
$f', g' : Y \to X_Y = Y \times_S X$，它们都是结构映射
$X_Y \to Y$ 的截面。注意 $f = g$ 当且仅当 $f' = g'$。
结构映射 $X_Y \to Y$ 是 $\pi$ 的基变换，因而仍非分歧且分离
（见《概形态射》引理 \ref{morphisms-lemma-base-change-unramified}
和《概形》引理 \ref{schemes-lemma-separated-permanence}）。
因此，由定理 \ref{etale-theorem-sections-unramified-maps}，只需证明
$f'$ 与 $g'$ 经过 $X_Y$ 的同一点。这恰由假设 (2) 和 (3) 保证，
即 $f'(y) = g'(y) \in X_Y$。
\end{proof}

\begin{lemma}
\label{etale-lemma-finitely-many-maps-to-unramified}
设 $S$ 为诺特概形。设 $X \to S$ 为拟紧非分歧态射。设
$Y \to S$ 为态射，其中 $Y$ 诺特。则 $\Mor_S(Y, X)$ 是有限集。
\end{lemma}

\begin{proof}
先假设 $X \to S$ 分离（实践中通常如此）。由于 $Y$ 诺特，
它只有有限多个连通分支，故可假设 $Y$ 连通。选取点 $y \in Y$，
其像为 $s \in S$。由于 $X \to S$ 非分歧且拟紧，纤维 $X_s$
是有限的；记 $X_s = \{x_1, \ldots, x_n\}$，并且
$\kappa(x_i)/\kappa(s)$ 是有限域扩张。见《概形态射》引理 \ref{morphisms-lemma-unramified-quasi-finite}、
\ref{morphisms-lemma-residue-field-quasi-finite} 和
\ref{morphisms-lemma-quasi-finite}。对于每个 $i$，
$\kappa(s)$-代数同态 $\kappa(x_i) \to \kappa(y)$ 至多只有有限多个
（由初等域论）。因此，由命题 \ref{etale-proposition-equality}，
$\Mor_S(Y, X)$ 有限。

\medskip\noindent
一般情形。存在非空开集 $U \subset S$，使得 $X_U \to U$ 有限
（特别地，分离）；见《概形态射》引理 \ref{morphisms-lemma-generically-finite}（该引理适用，因为我们已经
看到拟紧非分歧态射是拟有限的，而且由《概形态射》引理 \ref{morphisms-lemma-finite-type-Noetherian-quasi-separated}，
$X \to S$ 拟分离）。设 $Z \subset S$ 为支集等于 $U$ 的补集的
既约闭子概形。由诺特归纳，$\Mor_Z(Y_Z, X_Z)$ 有限（略去细节）。
由第一段的结果，集合 $\Mor_U(Y_U, X_U)$ 也有限。因此只需证明映射
$$
\Mor_S(Y, X) \longrightarrow \Mor_Z(Y_Z, X_Z) \times \Mor_U(Y_U, X_U)
$$
是单射。两个态射 $a, b : Y \to X$ 相等的点集在 $Y$ 中是开集，
因为 $\Delta : X \to X \times_S X$ 是开态射；见《概形态射》引理 \ref{morphisms-lemma-diagonal-unramified-morphism}。由此得到所需结论。
\end{proof}







\section{泛单射的非分歧态射}
\label{etale-section-universally-injective-unramified}

\noindent
回忆：若概形态射 $f : X \to Y$ 的任意基变换在底拓扑空间上都是单射，
则称 $f$ 泛单射；见《概形态射》定义 \ref{morphisms-definition-universally-injective}。泛单射的非分歧态射
可刻画如下。

\begin{lemma}
\label{etale-lemma-universally-injective-unramified}
设 $f : X \to S$ 为概形态射。下列条件等价：
\begin{enumerate}
\item $f$ 非分歧且为单态射；
\item $f$ 非分歧且泛单射；
\item $f$ 局部有限型且为单态射；
\item $f$ 泛单射、局部有限型且形式非分歧；
\item $f$ 局部有限型，并且对所有 $s \in S$，$X_s$ 或为空，
或 $X_s \to s$ 为同构。
\end{enumerate}
\end{lemma}

\begin{proof}
《态射进阶》引理 \ref{more-morphisms-lemma-unramified-formally-unramified} 已表明：
形式非分歧且局部有限型等价于非分歧。因此 (4) 等价于 (2)。
单态射显然泛单射且形式非分歧，故 (3) 蕴含 (4)。(1) 蕴含 (3)
也很清楚。最后，若 (2) 成立，则
$\Delta : X \to X \times_S X$ 既是开浸入
（《概形态射》引理 \ref{morphisms-lemma-diagonal-unramified-morphism}）
又是满射（《概形态射》引理 \ref{morphisms-lemma-universally-injective}），因而是同构；也就是说，
$f$ 是单态射。因此 (2) 蕴含 (1)。

\medskip\noindent
条件 (3) 蕴含 (5)，因为单态射在基变换下保持（《概形》引理 \ref{schemes-lemma-base-change-monomorphism}），并且《概形》引理 \ref{schemes-lemma-mono-towards-spec-field} 描述了到域的谱的单态射。
由《概形态射》引理 \ref{morphisms-lemma-universally-injective} 和
\ref{morphisms-lemma-unramified-etale-fibres}，条件 (5) 蕴含 (4)。
\end{proof}

\noindent
由此得到闭浸入的如下有用刻画。

\begin{lemma}
\label{etale-lemma-characterize-closed-immersion}
设 $f : X \to S$ 为概形态射。下列条件等价：
\begin{enumerate}
\item $f$ 是闭浸入；
\item $f$ 是固有单态射；
\item $f$ 固有、非分歧且泛单射；
\item $f$ 泛闭、非分歧且为单态射；
\item $f$ 泛闭、非分歧且泛单射；
\item $f$ 泛闭、局部有限型且为单态射；
\item $f$ 泛闭、泛单射、局部有限型且形式非分歧。
\end{enumerate}
\end{lemma}

\begin{proof}
(4)--(7) 的等价性直接由引理 \ref{etale-lemma-universally-injective-unramified} 得到。

\medskip\noindent
设 $f : X \to S$ 满足 (6)。则 $f$ 分离，见《概形》引理 \ref{schemes-lemma-monomorphism-separated}，并且其纤维有限。
因此《态射进阶》引理 \ref{more-morphisms-lemma-characterize-finite} 表明 $f$ 有限。
于是《概形态射》引理 \ref{morphisms-lemma-finite-monomorphism-closed}
蕴含 $f$ 是闭浸入，即 (1) 成立。

\medskip\noindent
注意 (1) $\Rightarrow$ (2)，因为闭浸入是固有单态射
（《概形态射》引理 \ref{morphisms-lemma-closed-immersion-proper}
及《概形》引理 \ref{schemes-lemma-immersions-monomorphisms}）。
由引理 \ref{etale-lemma-universally-injective-unramified} 可知 (2) 蕴含
(3)，而 (3) 蕴含 (5) 是显然的。
\end{proof}

\noindent
下面是另一个性质相近的结果。

\begin{lemma}
\label{etale-lemma-finite-unramified-one-point}
设 $\pi : X \to S$ 为概形态射，$s \in S$。假设
\begin{enumerate}
\item $\pi$ 有限；
\item $\pi$ 非分歧；
\item $\pi^{-1}(\{s\}) = \{x\}$；且
\item $\kappa(s) \subset \kappa(x)$ 为纯不可分扩张
\footnote{由条件 (2)，这等价于 $\kappa(s) = \kappa(x)$。}。
\end{enumerate}
则存在 $s$ 的开邻域 $U$，使得
$\pi|_{\pi^{-1}(U)} : \pi^{-1}(U) \to U$ 是闭浸入。
\end{lemma}

\begin{proof}
这一问题在 $S$ 上是局部的，故可假设 $S = \Spec(A)$。由有限态射的定义，
这蕴含 $X = \Spec(B)$。注意，定义 $\pi$ 的环同态
$\varphi : A \to B$ 是有限非分歧环同态。设
$\mathfrak p \subset A$ 为对应于 $s$ 的素理想，设
$\mathfrak q \subset B$ 为对应于 $x$ 的素理想。条件 (2)、(3) 和 (4)
蕴含 $B_{\mathfrak q}/\mathfrak pB_{\mathfrak q} = \kappa(\mathfrak p)$。
由《交换代数》引理 \ref{algebra-lemma-unique-prime-over-localize-below}，
$B_{\mathfrak q} = B_{\mathfrak p}$（注意有限环同态满足上升性质，见
《交换代数》第 \ref{algebra-section-going-up} 节）。因此
$B_{\mathfrak p}/\mathfrak pB_{\mathfrak p} = \kappa(\mathfrak p)$。
由于 $B$ 是有限 $A$-模，由 Nakayama 引理（见《交换代数》引理 \ref{algebra-lemma-NAK}）可知
$B_{\mathfrak p} = \varphi(A_{\mathfrak p})$。再次利用 $B$ 作为
$A$-模的有限性，存在 $f \in A$、$f \not \in \mathfrak p$，使得
$B_f = \varphi(A_f)$，如所需。
\end{proof}

\noindent
以上拓扑结果将用于给出平展（\'etale）态射的函子性刻画，
类似于定理 \ref{etale-theorem-formally-unramified}。




\section{非分歧态射的例子}
\label{etale-section-examples}

\noindent
下面给出几个例子。

\begin{example}
\label{etale-example-etale-field-extensions}
设 $k$ 为域。拟紧非分歧态射 $X \to \Spec(k)$ 是仿射的。
这是因为 $X$ 维数为 $0$ 且诺特，因而是有限离散集；每一点都给出一个
仿射开集，所以 $X$ 是有限多个仿射概形的不交并，故为仿射概形。
诺特正规化迫使 $X$ 成为有限 $k$-代数 $A$ 的谱。该代数是 $k$ 的
有限可分域扩张的乘积。因此，到 $\Spec(k)$ 的拟紧非分歧态射对应于
有限多个 $k$ 的有限可分域扩张。特别地，源连通而目标只有一点的
非分歧态射必为有限可分域扩张。稍后将看到，
$X \to \Spec(k)$ 平展（\'etale）当且仅当它非分歧。因此至少在此情形，
我们得到概形平展（\'etale）拓扑的极为简单的描述。当然，该拓扑的上同调
则是另一回事。
\end{example}

\begin{example}
\label{etale-example-standard-etale}
定理 \ref{etale-theorem-unramified-equivalence} 中的性质 (3)
为非分歧态射提供了一类典范例子。固定环 $R$ 与整数 $n$。
设 $I = (g_1, \ldots, g_m)$ 是 $R[x_1, \ldots, x_n]$ 中的理想。
设 $\mathfrak q \subset R[x_1, \ldots, x_n]$ 为素理想。假设
$I \subset \mathfrak q$，且矩阵
$$
\left(\frac{\partial g_i}{\partial x_j}\right) \bmod \mathfrak q
\quad\in\quad
\text{Mat}(n \times m, \kappa(\mathfrak q))
$$
的秩为 $n$。则态射
$f : Z = \Spec(R[x_1, \ldots, x_n]/I) \to \Spec(R)$
在对应于 $\mathfrak q$ 的点 $x \in Z \subset \mathbf{A}^n_R$
处非分歧。显然必须有 $m \geq n$。在极端情形 $m = n$ 下，
即由各 $g_i$ 定义的映射
$\mathbf{A}^n_R \to \mathbf{A}^n_R$ 的微分给出切空间的同构时，
$f$ 在 $x$ 处也平坦，因而是平展（\'etale）态射（见《交换代数》定义 \ref{algebra-definition-standard-smooth}、引理 \ref{algebra-lemma-standard-smooth} 和例 \ref{algebra-example-make-standard-smooth}）。
\end{example}

\begin{example}
\label{etale-example-number-theory-etale}
固定数域扩张 $L/K$，其整数环分别为 $\mathcal{O}_L$ 与
$\mathcal{O}_K$。嵌入 $K \to L$ 定义态射
$f : \Spec(\mathcal{O}_L) \to \Spec(\mathcal{O}_K)$。
如上所述，按本章意义 $f$ 非分歧的点，恰对应于按经典意义 $f$
非分歧的点。按经典意义，$\Spec(\mathcal{O}_L)$ 中的分歧轨迹可由
不同理想的零点集定义；它是 $\mathcal{O}_L$ 中的理想。事实上，
不同理想正是模 $\Omega_{\mathcal{O}_L/\mathcal{O}_K}$ 的零化子。
类似地，判别式是 $\mathcal{O}_K$ 中的理想，即不同理想的范数。
判别式的零点集恰为 $K$ 中在 $L$ 内分歧的点集。因此，记 $X$
为 $\Spec(\mathcal{O}_L)$ 中由不同理想定义的闭子集之补，
便得到非分歧态射 $X \to \Spec(\mathcal{O}_K)$。此外，该态射也平坦，
因为离散赋值环之间的任意局部同态都是平坦的，所以它实际上平展
（\'etale）。若 $L/K$ 为有限 Galois 扩张，并记 $Y$ 为
$\Spec(\mathcal{O}_K)$ 中由判别式定义的闭子集之补，则甚至得到有限
平展（\'etale）态射 $X \to Y$。因此，这是有限平展（\'etale）覆盖的例子。
\end{example}





\section{平坦态射}
\label{etale-section-flat-morphisms}

\noindent
本节只用于汇总后文所需的平坦性性质。因此，我们满足于精确陈述各定理，
并给出证明的参考文献。

\medskip\noindent
先简要回顾诺特环上平坦模的必要事实，随后陈述 Grothendieck 的一个定理；
它给出某些模的“超平面截面”平坦的充分条件。

\begin{definition}
\label{etale-definition-flat-rings}
模与环的平坦性。
\begin{enumerate}
\item 设 $N$ 是环 $A$ 上的模。若函子 $M \mapsto M \otimes_A N$
正合，则称该模{\it 平坦}。
\item 若该函子还忠实，则称 $N$ 在 $A$ 上{\it 忠实平坦}。
\item 若环同态 $f : A \to B$ 所定义的函子
$M \mapsto M \otimes_A B$ 正合（相应地，忠实且正合），则称
该同态{\it 平坦（相应地，忠实平坦）}。
\end{enumerate}
\end{definition}

\noindent
下面列出若干事实，并给出代数一章中的参考文献。
\begin{enumerate}
\item 自由模与投射模都是平坦模。对于自由模这是显然的；对于投射模，
这来自它们是自由模的直和因子，以及 $\otimes$ 与直和交换。
\item 平坦性是局部性质；也就是说，$M$ 在 $A$ 上平坦，当且仅当对所有
$\mathfrak p \in \Spec(A)$，$M_{\mathfrak p}$ 在
$A_{\mathfrak p}$ 上平坦。见《交换代数》引理 \ref{algebra-lemma-flat-localization}。
\item 若 $M$ 是平坦 $A$-模，且 $A \to B$ 是环同态，则
$M \otimes_A B$ 是平坦 $B$-模。见《交换代数》引理 \ref{algebra-lemma-flat-base-change}。
\item 局部环上的有限平坦模是自由模。见《交换代数》引理 \ref{algebra-lemma-finite-flat-local}。
\item 若 $f : A \to B$ 是任意环之间的同态，则 $f$ 平坦，当且仅当
对所有 $\mathfrak q \in \Spec(B)$，诱导映射
$A_{f^{-1}(\mathfrak q)} \to B_{\mathfrak q}$ 平坦。见《交换代数》引理 \ref{algebra-lemma-flat-localization}。
\item 若 $f : A \to B$ 是局部环之间的局部同态，则 $f$ 平坦，当且仅当
它忠实平坦。见《交换代数》引理 \ref{algebra-lemma-local-flat-ff}。
\item 环同态 $A \to B$ 忠实平坦，当且仅当它平坦且诱导的谱上态射满射。
见《交换代数》引理 \ref{algebra-lemma-ff-rings}。
\item 若 $A$ 是诺特局部环，则完备化 $A^\wedge$ 在 $A$ 上忠实平坦。
见《交换代数》引理 \ref{algebra-lemma-completion-faithfully-flat}。
\item 设 $A$ 为诺特局部环，$M$ 为 $A$-模。则 $M$ 在 $A$ 上平坦，
当且仅当 $M \otimes_A A^\wedge$ 在 $A^\wedge$ 上平坦。
（结合上一项与《交换代数》引理 \ref{algebra-lemma-flatness-descends}。）
\end{enumerate}
在转向几何范畴之前，先给出 Grothendieck 定理；它提供了一种构造平坦模的
便利方法。

\begin{theorem}
\label{etale-theorem-flatness-grothendieck}
设 $A$、$B$ 为诺特局部环，$f : A \to B$ 为局部同态。
设 $M$ 为有限 $B$-模，并且作为 $A$-模是平坦的。若
$t \in \mathfrak m_B$，且乘以 $t$ 在 $M/\mathfrak m_AM$ 上是单射，
则 $M/tM$ 也是平坦 $A$-模。
\end{theorem}

\begin{proof}
见《交换代数》引理 \ref{algebra-lemma-mod-injective}。另见
\cite[Section 20]{MatCA}。
\end{proof}

\begin{definition}
\label{etale-definition-flat-schemes}
（见《概形态射》定义 \ref{morphisms-definition-flat}。）
设 $f : X \to Y$ 为概形态射，$\mathcal{F}$ 为拟凝聚
$\mathcal{O}_X$-模。
\begin{enumerate}
\item 设 $x \in X$。若 $\mathcal{F}_x$ 是平坦
$\mathcal{O}_{Y, f(x)}$-模，则称 $\mathcal{F}$
{\it 在 $x \in X$ 处关于 $Y$ 平坦}。这里借助映射
$\mathcal{O}_{Y, f(x)} \to \mathcal{O}_{X, x}$，把
$\mathcal{F}_x$ 看作 $\mathcal{O}_{Y, f(x)}$-模。
\item 设 $x \in X$。若
$\mathcal{O}_{Y, f(x)} \to \mathcal{O}_{X, x}$ 平坦，则称 $f$
{\it 在 $x \in X$ 处平坦}。
\item 若该态射在 $X$ 的所有点处平坦，则称 $f$ {\it 平坦}。
\item 平坦且满射的态射 $f : X \to Y$ 有时称为{\it 忠实平坦}态射。
\end{enumerate}
\end{definition}

\noindent
再列出若干结果：
\begin{enumerate}
\item 按定义，态射的平坦性在源和目标上对于 Zariski 拓扑都是局部性质。
\item 开浸入平坦（这是显然的，因为它在局部环上诱导同构）。
\item 平坦态射在基变换与复合下稳定。见《概形态射》引理 \ref{morphisms-lemma-base-change-flat} 和
\ref{morphisms-lemma-composition-flat}。
\item 若 $f : X \to Y$ 平坦，则拉回函子
$\QCoh(\mathcal{O}_Y) \to \QCoh(\mathcal{O}_X)$ 正合。
考察各茎即可立即得到这一点。
\item 设 $f : X \to Y$ 为概形态射，并假设 $Y$ 拟紧且拟分离。
在此情形，若函子 $f^*$ 正合，则 $f$ 平坦。（略去证明。提示：
利用《概形的性质》引理 \ref{properties-lemma-extend-trivial}
说明 $Y$ 有“足够多”的理想层，再使用《交换代数》引理 \ref{algebra-lemma-flat} 中的平坦性刻画。）
\end{enumerate}



\section{平坦态射的拓扑性质}
\label{etale-section-topological-flat}

\noindent
下面“回顾”平坦态射所具有的一些开性性质。

\begin{theorem}
\label{etale-theorem-flat-open}
设 $Y$ 为局部诺特概形，$f : X \to Y$ 为局部有限型态射，
$\mathcal{F}$ 为凝聚 $\mathcal{O}_X$-模。则 $X$ 中
$\mathcal{F}$ 关于 $Y$ 平坦的点所成集合是开集。特别地，
$f$ 的平坦点集在 $X$ 中是开集。
\end{theorem}

\begin{proof}
见《态射进阶》定理 \ref{more-morphisms-theorem-openness-flatness}。
\end{proof}

\begin{theorem}
\label{etale-theorem-flat-map-open}
设 $Y$ 为局部诺特概形，$f : X \to Y$ 为平坦局部有限型态射。
则 $f$ 泛开。
\end{theorem}

\begin{proof}
见《概形态射》引理 \ref{morphisms-lemma-fppf-open}。
\end{proof}

\begin{theorem}
\label{etale-theorem-flat-is-quotient}
忠实平坦拟紧态射对于 Zariski 拓扑是商映射。
\end{theorem}

\begin{proof}
见《概形态射》引理 \ref{morphisms-lemma-fpqc-quotient-topology}。
\end{proof}

\noindent
研究平坦态射的一个重要原因是：它们为刻画由另一概形的点参数化的概形族
提供了适当框架。朴素地看，任意态射 $f : X \to S$ 都可视为由 $S$ 的点
参数化的族。然而，若不要求 $f$ 平坦，这个所谓的族可能出现非常反常的现象。
例如，族中的相对维数（纤维的维数）未必恒定；Hilbert 多项式等其他数值不变量
也可能随纤维而改变。平坦性排除了这类现象，从而赋予纤维某种“连续性”。


\section{平展（\'Etale）态射}
\label{etale-section-etale-morphisms}

\noindent
本节定义平展（\'etale）态射，并证明它们的一系列重要性质。其中无疑最重要的
是定理 \ref{etale-theorem-formally-etale} 中的函子性刻画。随后还将讨论若干
对平展（\'etale）扩张不敏感的环性质（即一个环具有这些性质，当且仅当其所有
平展（\'etale）扩张都具有这些性质），以说明平展（\'etale）上同调的基本信条：
平展（\'etale）态射是局部同构的代数对应物。

\medskip\noindent
如标题所示，我们将定义平展（\'etale）态射类。为了定义基概形 $S$ 上概形范畴
中的平展（\'etale）站点 $S_\etale$，我们把这类态射的满射族视为覆盖。
直观上，平展（\'etale）态射应当表达覆盖空间的观念，因此应接近局部同构。
对于代数闭域上的簇，只要把“局部同构”替换为“形式局部同构”（即完备化后的
同构），这一说法就可作为定义。随后可对任意基域规定：到代数闭包的基变换按
上述意义平展（\'etale）。不过，我们不采用这种审美上不尽如人意的构造，
而采用更简洁、虽略为抽象的代数方法。

\medskip\noindent
首先为诺特局部环定义“局部环的平展（\'etale）同态”。这里不能直接使用
“平展（\'etale）”一词，因为《交换代数》第 \ref{algebra-section-etale} 节中已有平展（\'etale）环同态的概念，
而两者并不相同。

\begin{definition}
\label{etale-definition-etale-ring}
设 $A$、$B$ 为诺特局部环。若局部同态 $f : A \to B$ 平坦，
并且是局部环的非分歧同态（见定义 \ref{etale-definition-unramified-rings}），则称它为
{\it 局部环的平展（\'etale）同态}。
\end{definition}

\noindent
这是《交换代数》第 \ref{algebra-section-etale} 节中平展（\'etale）
环同态定义的局部版本。该节的精确定义是相对维数为 $0$ 的光滑环同态。
《交换代数》引理 \ref{algebra-lemma-etale-standard-smooth} 证明：
平展（\'etale）$R$-代数 $S$ 总有呈示
$$
S = R[x_1, \ldots, x_n]/(f_1, \ldots, f_n)
$$
使得
$$
g =
\det
\left(
\begin{matrix}
\partial f_1/\partial x_1 &
\partial f_2/\partial x_1 &
\ldots &
\partial f_n/\partial x_1 \\
\partial f_1/\partial x_2 &
\partial f_2/\partial x_2 &
\ldots &
\partial f_n/\partial x_2 \\
\ldots & \ldots & \ldots & \ldots \\
\partial f_1/\partial x_n &
\partial f_2/\partial x_n &
\ldots &
\partial f_n/\partial x_n
\end{matrix}
\right)
$$
在 $S$ 中的像可逆。下面两个引理联系这两种概念。

\begin{lemma}
\label{etale-lemma-characterize-etale-Noetherian}
设 $A$ 为诺特环，$A \to B$ 为有限型同态。设 $\mathfrak q$ 为
$B$ 的素理想，并位于 $\mathfrak p \subset A$ 之上。则
$A \to B$ 在 $\mathfrak q$ 处平展（\'etale），当且仅当
$A_{\mathfrak p} \to B_{\mathfrak q}$ 是局部环的平展（\'etale）同态。
\end{lemma}

\begin{proof}
见《交换代数》引理 \ref{algebra-lemma-etale}（平展（\'etale）同态平坦）、
\ref{algebra-lemma-etale-at-prime}（平展（\'etale）同态非分歧）和
\ref{algebra-lemma-characterize-etale}（平坦非分歧同态平展（\'etale））。
\end{proof}

\begin{lemma}
\label{etale-lemma-characterize-etale-completions}
设 $A$、$B$ 为诺特局部环。设 $A \to B$ 为局部同态，并且 $B$
本质上是 $A$ 上的有限型代数。下列条件等价：
\begin{enumerate}
\item $A \to B$ 是局部环的平展（\'etale）同态；
\item $A^\wedge \to B^\wedge$ 是局部环的平展（\'etale）同态；且
\item $A^\wedge \to B^\wedge$ 是平展（\'etale）环同态。
\end{enumerate}
此外，在此情形，对某个 $n \geq 1$，作为 $A^\wedge$-模有
$B^\wedge \cong (A^\wedge)^{\oplus n}$。
\end{lemma}

\begin{proof}
关于非分歧环同态，我们已有对应结果（引理 \ref{etale-lemma-characterize-unramified-completions}）。因此，为证明
(1)、(2)、(3) 等价，只需证明 $A \to B$ 平坦，当且仅当
$A^\wedge \to B^\wedge$ 平坦。根据上面列出的平坦同态性质，这是显然的，
因为环同态 $A \to A^\wedge$ 和 $B \to B^\wedge$ 都忠实平坦。
对于最后一项陈述，由引理 \ref{etale-lemma-unramified-completions}，
$B^\wedge$ 是有限平坦 $A^\wedge$-模。因此，根据第 \ref{etale-section-flat-morphisms} 节中列出的平坦模性质，它是有限自由模。
\end{proof}

\noindent
上引理中的整数 $n$ 正是剩余域扩张的次数
$[\kappa(\mathfrak m_B) : \kappa(\mathfrak m_A)]$。特别地，若
$\kappa(\mathfrak m_A)$ 可分闭，则 $A^\wedge \to B^\wedge$
是同构，这也证实了前面的说法。

\begin{definition}
\label{etale-definition-etale-schemes-1}
（见《概形态射》定义 \ref{morphisms-definition-etale}。）
设 $Y$ 为局部诺特概形，$f : X \to Y$ 为局部有限型概形态射。
\begin{enumerate}
\item 设 $x \in X$。若
$\mathcal{O}_{Y, f(x)} \to \mathcal{O}_{X, x}$ 是局部环的
平展（\'etale）同态，则称 $f$ {\it 在 $x \in X$ 处平展（\'etale）}。
\item 若该态射在所有点处都平展（\'etale），则称它{\it 平展（\'etale）}。
\end{enumerate}
\end{definition}

\noindent
下面证明此定义与概形态射一章中的定义一致。特别地，这保证了一个态射的
平展（\'etale）点所成集合是开集。

\begin{lemma}
\label{etale-lemma-etale-definition}
设 $Y$ 为局部诺特概形，$f : X \to Y$ 局部有限型，并设
$x \in X$。则态射 $f$ 按定义 \ref{etale-definition-etale-schemes-1}
在 $x$ 处平展（\'etale），当且仅当它按《概形态射》定义 \ref{morphisms-definition-etale} 在 $x$ 处平展（\'etale）。
\end{lemma}

\begin{proof}
这由引理 \ref{etale-lemma-characterize-etale-Noetherian} 和定义得到。
\end{proof}

\noindent
下面列出平展（\'etale）态射的一些结果。此处表述仅适用于局部诺特概形之间的
局部有限型态射。每一项都给出本项目先前证明的一般结果作为参考，不过其中
一些结果在诺特情形下可更容易地证明：
\begin{enumerate}
\item 平展（\'etale）态射非分歧（由定义即得）。
\item 平展（\'Etale）性在源和目标上对于 Zariski 拓扑都是局部的。
\item 平展（\'Etale）态射在基变换与复合下稳定。见《概形态射》引理 \ref{morphisms-lemma-base-change-etale} 和
\ref{morphisms-lemma-composition-etale}。
\item 概形的平展（\'Etale）态射局部拟有限；拟紧平展（\'etale）态射
拟有限（因为前面已知非分歧态射具有这一性质）。
\item 平展（\'Etale）态射的相对维数为 $0$。见《概形态射》定义 \ref{morphisms-definition-relative-dimension-d} 和《概形态射》引理 \ref{morphisms-lemma-locally-quasi-finite-rel-dimension-0}。
\item 一个态射平展（\'etale），当且仅当它平坦且所有纤维都平展
（\'etale）。见《概形态射》引理 \ref{morphisms-lemma-etale-flat-etale-fibres}。
\item 平展（\'Etale）态射是开态射。这是因为平展（\'etale）态射平坦，
再应用定理 \ref{etale-theorem-flat-map-open}。
\item 设 $X$ 和 $Y$ 在基概形 $S$ 上平展（\'etale）。
则从 $X$ 到 $Y$ 的任意 $S$-态射都平展（\'etale）。见《概形态射》引理 \ref{morphisms-lemma-etale-permanence}。
\end{enumerate}






\section{结构定理}
\label{etale-section-structure-etale-map}

\noindent
下面给出描述平展（\'etale）态射与非分歧态射局部结构的定理。
除其本身显然重要外，这一定理还使我们能够转向平展（\'etale）态射的
另一种定义；与目前使用的定义相比，它更能体现几何直观。

\medskip\noindent
为陈述定理，需要{\it 标准平展（\'etale）环同态}的概念，见《交换代数》定义 \ref{algebra-definition-standard-etale}。具体地，设 $R$ 为环，
$f, g \in R[t]$ 为多项式，并满足
\begin{enumerate}
\item[(a)] $f$ 为首一多项式；且
\item[(b)] $f' = \text{d}f/\text{d}t$ 在局部化 $R[t]_g/(f)$ 中可逆。
\end{enumerate}
则同态
$$
R \longrightarrow R[t]_g/(f) = R[t, 1/g]/(f)
$$
给出标准平展（\'etale）代数；任意标准平展（\'etale）代数都同构于这种
代数之一。证明这类环同态平坦、非分歧，因而平展（\'etale），是一个愉快的
练习（这当然正如所料）。标准平展（\'etale）环同态的一个特例是任意同态
$$
R \longrightarrow R[t]_{f'}/(f) = R[t, 1/f']/(f)
$$
其中 $f$ 为首一多项式；任意标准平展（\'etale）代数都同构于其中某个代数的
主局部化。

\begin{theorem}
\label{etale-theorem-structure-etale}
设 $f : A \to B$ 为局部环的平展（\'etale）同态。则存在
$f, g \in A[t]$，使得
\begin{enumerate}
\item $B' = A[t]_g/(f)$ 是标准平展（\'etale）代数，见上面的 (a)、(b)；且
\item $B$ 同构于 $B'$ 在某个素理想处的局部化。
\end{enumerate}
\end{theorem}

\begin{proof}
取某个有限型 $A$-代数 $B'$，使其在素理想处的局部化为
$B = B'_{\mathfrak q}$（这是因为 $B$ 本质上是 $A$ 上的有限型代数）。
由引理 \ref{etale-lemma-characterize-etale-Noetherian}，$A \to B'$
在 $\mathfrak q$ 处平展（\'etale）。应用《交换代数》命题 \ref{algebra-proposition-etale-locally-standard}，可知 $B'$ 的某个
主局部化是标准平展（\'etale）代数。
\end{proof}

\noindent
下面是局部环的非分歧同态版本。

\begin{theorem}
\label{etale-theorem-structure-unramified}
设 $f : A \to B$ 为局部环的非分歧态射。则存在 $f, g \in A[t]$，使得
\begin{enumerate}
\item $B' = A[t]_g/(f)$ 是标准平展（\'etale）代数，见上面的 (a)、(b)；且
\item $B$ 同构于 $B'$ 在某个素理想处局部化所得环的商。
\end{enumerate}
\end{theorem}

\begin{proof}
取某个有限型 $A$-代数 $B'$，使其在素理想处的局部化为
$B = B'_{\mathfrak q}$（这是因为 $B$ 本质上是 $A$ 上的有限型代数）。
由引理 \ref{etale-lemma-characterize-unramified-Noetherian}，$A \to B'$
在 $\mathfrak q$ 处非分歧。应用《交换代数》命题 \ref{algebra-proposition-unramified-locally-standard}，可知 $B'$
的某个主局部化是标准平展（\'etale）$A$-代数的商。
\end{proof}

\noindent
再通过标准提升论证，得到下面这个对后文至关重要的几何陈述。

\begin{theorem}
\label{etale-theorem-geometric-structure}
设 $\varphi : X \to Y$ 为概形态射，$x \in X$。设 $V \subset Y$
为 $\varphi(x)$ 的仿射开邻域。若 $\varphi$ 在 $x$ 处平展
（\'etale），则存在仿射开集 $U \subset X$，满足 $x \in U$ 与
$\varphi(U) \subset V$，并有图表
$$
\xymatrix{
X \ar[d] & U \ar[l] \ar[d] \ar[r]_-j & \Spec(R[t]_{f'}/(f)) \ar[d] \\
Y & V \ar[l] \ar@{=}[r] & \Spec(R)
}
$$
其中 $j$ 是开浸入，$f \in R[t]$ 为首一多项式。
\end{theorem}

\begin{proof}
这等价于《概形态射》引理 \ref{morphisms-lemma-etale-locally-standard-etale}，尽管两个陈述略有差异。
另见《代数簇》引理 \ref{varieties-lemma-geometric-structure-unramified} 中非分歧态射的变体。
\end{proof}


\section{平展（\'Etale）态射与光滑态射}
\label{etale-section-etale-smooth}

\noindent
平展（\'etale）态射是相对维数为零的光滑态射。投影
$\mathbf{A}^n_S \to S$ 是相对维数为 $n$ 的光滑态射的标准例子。
事实上，任意光滑态射在平展（\'etale）局部都具有这种形式。
精确陈述如下。

\begin{theorem}
\label{etale-theorem-smooth-etale-over-n-space}
设 $\varphi : X \to Y$ 为概形态射，$x \in X$。若 $\varphi$
在 $x$ 处光滑，则存在整数 $n \geq 0$ 以及仿射开集
$V \subset Y$、$U \subset X$，满足 $x \in U$ 与
$\varphi(U) \subset V$，并存在交换图
$$
\xymatrix{
X \ar[d] & U \ar[l] \ar[d] \ar[r]_-\pi &
\mathbf{A}^n_R \ar[d] \ar@{=}[r] &  \Spec(R[x_1, \ldots, x_n]) \ar[dl] \\
Y & V \ar[l] \ar@{=}[r] & \Spec(R)
}
$$
其中 $\pi$ 平展（\'etale）。
\end{theorem}

\begin{proof}
见《概形态射》引理 \ref{morphisms-lemma-smooth-etale-over-affine-space}。
\end{proof}




\section{平展（\'Etale）态射的拓扑性质}
\label{etale-section-topological-etale}

\noindent
下面给出平展（\'etale）态射与非分歧态射的一些拓扑性质。首先给出
Grothendieck 所称的{\it 平展（\'etale）态射的基本性质}，见
\cite[Expos\'e I.5]{SGA1}。

\begin{theorem}
\label{etale-theorem-etale-radicial-open}
设 $f : X \to Y$ 为概形态射。下列条件等价：
\begin{enumerate}
\item $f$ 是开浸入；
\item $f$ 泛单射且平展（\'etale）；且
\item $f$ 是局部有限呈示的平坦单态射。
\end{enumerate}
\end{theorem}

\begin{proof}
开浸入泛单射，因为开浸入的任意基变换仍是开浸入。此外，由《概形态射》引理 \ref{morphisms-lemma-open-immersion-etale}，它平展
（\'etale）。故 (1) 蕴含 (2)。

\medskip\noindent
假设 $f$ 泛单射且平展（\'etale）。由于 $f$ 平展（\'etale），
它平坦且局部有限呈示；见《概形态射》引理 \ref{morphisms-lemma-etale-flat} 和
\ref{morphisms-lemma-etale-locally-finite-presentation}。
由引理 \ref{etale-lemma-universally-injective-unramified}，$f$ 是单态射。
故 (2) 蕴含 (3)。

\medskip\noindent
假设 $f$ 平坦、局部有限呈示且为单态射。由《概形态射》引理 \ref{morphisms-lemma-fppf-open}，$f$ 是开态射。因此可用 $f(X)$
代替 $Y$，并假设 $f$ 满射。此时 $f$ 既开又双射，因而是同胚；
于是 $f$ 拟紧。《下降》引理 \ref{descent-lemma-flat-surjective-quasi-compact-monomorphism-isomorphism}
表明 $f$ 是同构，结论得证。
\end{proof}

\noindent
下面是另一个性质相近的结果。

\begin{lemma}
\label{etale-lemma-finite-etale-one-point}
设 $\pi : X \to S$ 为概形态射，$s \in S$。假设
\begin{enumerate}
\item $\pi$ 有限；
\item $\pi$ 平展（\'etale）；
\item $\pi^{-1}(\{s\}) = \{x\}$；且
\item $\kappa(s) \subset \kappa(x)$ 为纯不可分扩张
\footnote{由条件 (2)，这等价于 $\kappa(s) = \kappa(x)$。}。
\end{enumerate}
则存在 $s$ 的开邻域 $U$，使得
$\pi|_{\pi^{-1}(U)} : \pi^{-1}(U) \to U$ 是同构。
\end{lemma}

\begin{proof}
由引理 \ref{etale-lemma-finite-unramified-one-point}，存在 $s$ 的开邻域
$U$，使得 $\pi|_{\pi^{-1}(U)} : \pi^{-1}(U) \to U$ 是闭浸入。
但平展（\'etale）的闭浸入是开浸入（例如由定理 \ref{etale-theorem-etale-radicial-open}）。因此缩小 $U$ 后得到同构。
\end{proof}

\begin{lemma}
\label{etale-lemma-relative-frobenius-etale}
设 $U \to X$ 为平展（\'etale）概形态射，其中 $X$ 是特征 $p$
的概形。则相对 Frobenius 态射
$F_{U/X} : U \to U \times_{X, F_X} X$ 是同构。
\end{lemma}

\begin{proof}
由《代数簇》引理 \ref{varieties-lemma-relative-frobenius}，
态射 $F_{U/X}$ 是泛同胚。作为 $X$ 上两个平展（\'etale）概形之间的态射，
$F_{U/X}$ 平展（\'etale）（《概形态射》引理 \ref{morphisms-lemma-etale-permanence}）。因此，由定理 \ref{etale-theorem-etale-radicial-open}，$F_{U/X}$ 是同构。
\end{proof}






\section{平展（\'etale）拓扑的拓扑不变性}
\label{etale-section-topological-invariance}

\noindent
下面给出一个极其关键的定理；粗略地说，它断言平展（\'etale）性是一种拓扑性质。

\begin{theorem}
\label{etale-theorem-etale-topological}
设 $X$ 与 $Y$ 是基概形 $S$ 上的两个概形。设 $S_0$ 是 $S$ 的闭子概形，
且二者具有相同的底拓扑空间（例如，若 $S_0$ 在 $S$ 中的理想层的平方为零）。
记 $X_0$（相应地为\ $Y_0$）为基变换 $S_0 \times_S X$
（相应地为\ $S_0 \times_S Y$）。
若 $X$ 在 $S$ 上平展（\'etale），则映射
$$
\Mor_S(Y, X) \longrightarrow \Mor_{S_0}(Y_0, X_0)
$$
是双射。
\end{theorem}

\begin{proof}
经由 $Y \to S$ 作基变换后，可假设 $Y = S$。在此情形，定理断言：
任意 $S$-态射 $\sigma_0 : S_0 \to X$ 都唯一地经由平展（\'etale）
结构态射 $f : X \to S$ 的一个截面 $S \to X$ 分解。

\medskip\noindent
唯一性。假设有两个使 $\sigma_0$ 经其分解的截面 $\sigma, \sigma'$。
由于 $X \to S$ 平展（\'etale），
$\Delta : X \to X \times_S X$ 是开浸入
（《概形态射》引理 \ref{morphisms-lemma-diagonal-unramified-morphism}）。
态射 $(\sigma, \sigma') : S \to X \times_S X$ 经由这个开子概形分解，
因为对任意 $s \in S$ 都有
$(\sigma, \sigma')(s) = (\sigma_0(s), \sigma_0(s))$。故
$\sigma = \sigma'$。

\medskip\noindent
为证明存在性，先将问题约化到仿射情形
（建议读者略过这一步）。
取仿射开覆盖 $X = \bigcup X_i$，使每个 $X_i$ 都映入 $S$ 的某个
仿射开集 $S_i$。对每个 $s \in S$，可选取 $i$ 使得
$\sigma_0(s) \in X_i$。
选取 $s$ 的仿射开邻域 $U \subset S_i$，使得
$\sigma_0(U_0) \subset X_{i, 0}$。注意
$X' = X_i \times_S U = X_i \times_{S_i} U$ 是仿射的。
若能将 $\sigma_0|_{U_0} : U_0 \to X'_0$ 提升为
$U \to X'$，则由唯一性，这些局部提升可黏合成整体态射
$S \to X$。因此可假设 $S$ 与 $X$ 均为仿射概形。

\medskip\noindent
$S$ 与 $X$ 均仿射时的存在性。写成 $S = \Spec(A)$ 与
$X = \Spec(B)$。于是 $A \to B$ 平展（\'etale），特别地，
它光滑（相对维数为 $0$）。由 $|S_0| = |S|$ 可知
$S_0 = \Spec(A/I)$，其中 $I \subset A$ 局部幂零。
故存在性由《交换代数》引理 \ref{algebra-lemma-smooth-strong-lift} 得出。
\end{proof}

\noindent
由前一定理的证明，还可得到所许诺的平展（\'etale）态射函子刻画的
一个方向。下述定理将在
《平展（\'Etale）上同调》定理 \ref{etale-cohomology-theorem-topological-invariance}
中得到加强。

\begin{theorem}[范畴的非凡等价（Une equivalence remarquable de cat\'egories）]
\label{etale-theorem-remarkable-equivalence}
\begin{reference}
\cite[IV, Theorem 18.1.2]{EGA}
\end{reference}
设 $S$ 为概形。设 $S_0 \subset S$ 为具有相同底拓扑空间的闭子概形
（例如，若 $S_0$ 在 $S$ 中的理想层的平方为零）。函子
$$
X \longmapsto X_0 = S_0 \times_S X
$$
给出范畴等价
$$
\{
X\text{ 是在 }S\text{ 上的平展概形}
\}
\leftrightarrow
\{
X_0\text{ 是在 }S_0\text{ 上的平展概形}
\}
$$
\end{theorem}

\begin{proof}
由定理 \ref{etale-theorem-etale-topological}，该函子是全忠实的。
尚需证明它本质满。设 $Y \to S_0$ 是平展（\'etale）概形态射。

\medskip\noindent
先假设当 $S$ 与 $Y$ 均仿射时结论成立。此时选取仿射开覆盖
$Y = \bigcup V_j$，使每个 $V_j$ 都映入 $S$ 的某个仿射开集。
由仿射情形的假设，可找到平展（\'etale）态射 $W_j \to S$，
使 $W_{j, 0} \cong V_j$（作为 $S_0$ 上的概形）。
令 $W_{j, j'} \subset W_j$ 为这样的开子概形：其底拓扑空间对应于
$V_j \cap V_{j'}$。由于存在 $S_0$ 上的概形同构
$$
W_{j, j', 0} \cong V_j \cap V_{j'} \cong W_{j', j, 0}
$$
由全忠实性可得 $S$ 上的概形同构
$\theta_{j, j'} : W_{j, j'} \to W_{j', j}$。
这里略去验证这些同构满足《概形》第 \ref{schemes-section-glueing-schemes} 节中的余圈条件。
应用《概形》引理 \ref{schemes-lemma-glue-schemes}，
沿着同一化 $\theta_{j, j'}$ 黏合诸概形 $W_j$，得到概形
$X \to S$。由构造显然可知 $X \to S$ 平展（\'etale），且
$X_0 \cong Y$。

\medskip\noindent
因此，只需证明 $S$ 与 $Y$ 均仿射的情形。记
$S = \Spec(R)$，$S_0 = \Spec(R/I)$，其中 $I$ 局部幂零。
由《交换代数》引理 \ref{algebra-lemma-etale-standard-smooth}，
$Y$ 是某个环 $\overline{A}$ 的谱，并且
$$
\overline{A} = (R/I)[x_1, \ldots, x_n]/(\overline{f}_1, \ldots, \overline{f}_n)
$$
且
$$
\overline{g} =
\det
\left(
\begin{matrix}
\partial \overline{f}_1/\partial x_1 &
\partial \overline{f}_2/\partial x_1 &
\ldots &
\partial \overline{f}_n/\partial x_1 \\
\partial \overline{f}_1/\partial x_2 &
\partial \overline{f}_2/\partial x_2 &
\ldots &
\partial \overline{f}_n/\partial x_2 \\
\ldots & \ldots & \ldots & \ldots \\
\partial \overline{f}_1/\partial x_n &
\partial \overline{f}_2/\partial x_n &
\ldots &
\partial \overline{f}_n/\partial x_n
\end{matrix}
\right)
$$
在 $\overline{A}$ 中的像可逆。任取提升
$f_i \in R[x_1, \ldots, x_n]$，并令
$$
A = R[x_1, \ldots, x_n]/(f_1, \ldots, f_n)
$$
由于 $I$ 局部幂零，理想 $IA$ 也局部幂零
（《交换代数》引理 \ref{algebra-lemma-locally-nilpotent}）。
注意 $\overline{A} = A/IA$。由《交换代数》引理 \ref{algebra-lemma-locally-nilpotent-unit}，$f_i$ 的偏导数矩阵之
行列式在代数 $A$ 中可逆。故 $R \to A$ 平展（\'etale），证明完成。
\end{proof}



\section{函子刻画}
\label{etale-section-functorial-etale}

\noindent
最后给出所许诺的函子刻画。因此，可以用四种方式理解概形的平展
（\'etale）态射：
\begin{enumerate}
\item 将其视为相对维数为 $0$ 的光滑态射；
\item 将其视为局部有限呈示、平坦且非分歧的态射；
\item 使用结构定理；以及
\item 使用函子刻画。
\end{enumerate}

\begin{theorem}
\label{etale-theorem-formally-etale}
设 $f : X \to S$ 是局部有限呈示的态射。下列条件等价：
\begin{enumerate}
\item $f$ 平展（\'etale）；
\item 对所有仿射 $S$-概形 $Y$ 以及由平方零理想定义的闭子概形
$Y_0 \subset Y$，自然映射
$$
\Mor_S(Y, X) \longrightarrow \Mor_S(Y_0, X)
$$
是双射。
\end{enumerate}
\end{theorem}

\begin{proof}
这正是《态射进阶》引理 \ref{more-morphisms-lemma-etale-formally-etale}。
\end{proof}

\noindent
这一刻画表明，定义 $X$ 的方程组的解可以唯一地穿过幂零增厚而提升。



\section{非分歧态射的平展（\'Etale）局部结构}
\label{etale-section-unramified-etale-local}

\noindent
在《态射进阶》第 \ref{more-morphisms-section-etale-localization} 节中，读者可以找到一些关于
拟有限态射之平展（\'etale）局部结构的结果。本节要把这些结果与本章已经
得到的非分歧态射的拓扑性质结合起来。应当记住的基本整体图景是
$$
\xymatrix{
V \ar[r] \ar[dr] & X_U \ar[d] \ar[r] & X \ar[d]^f \\
& U \ar[r] & S
}
$$
见 More on Morphisms, Equation
(\ref{more-morphisms-equation-basic-diagram})。先从一个非常一般的情形开始。

\begin{lemma}
\label{etale-lemma-unramified-etale-local}
设 $f : X \to S$ 为概形态射。设
$x_1, \ldots, x_n \in X$ 是在 $S$ 中具有同一像 $s$ 的点。
假设 $f$ 在每个 $x_i$ 处非分歧。则存在平展（\'etale）邻域
$(U, u) \to (S, s)$ 以及开子概形
$V_{i, j} \subset X_U$，$i = 1, \ldots, n$，$j = 1, \ldots, m_i$，
使得
\begin{enumerate}
\item $V_{i, j} \to U$ 是经过 $u$ 的闭浸入；
\item 除非 $i = i'$ 且 $j = j'$，否则 $u$ 不在
$V_{i, j} \cap V_{i', j'}$ 的像中；且
\item $(X_U)_u$ 中映到 $x_i$ 的任意点都属于某个 $V_{i, j}$。
\end{enumerate}
\end{lemma}

\begin{proof}
由《概形态射》定义 \ref{morphisms-definition-unramified}，
每个 $x_i$ 都有一个在 $S$ 上局部有限型的开邻域。以
$\{x_1, \ldots, x_n\}$ 的某个开邻域代替 $X$，可假设 $f$
局部有限型。应用《态射进阶》引理 \ref{more-morphisms-lemma-etale-makes-quasi-finite-finite-multiple-points-var}，
得到平展（\'etale）邻域 $(U, u)$ 以及在 $U$ 上有限的开子概形
$V_{i, j}$。由引理 \ref{etale-lemma-finite-unramified-one-point}，
必要时缩小 $U$ 后可使 $V_{i, j} \to U$ 成为闭浸入。
\end{proof}

\begin{lemma}
\label{etale-lemma-unramified-etale-local-technical}
设 $f : X \to S$ 为概形态射。设
$x_1, \ldots, x_n \in X$ 是在 $S$ 中具有同一像 $s$ 的点。
假设 $f$ 分离，并且 $f$ 在每个 $x_i$ 处非分歧。则存在平展
（\'etale）邻域 $(U, u) \to (S, s)$ 以及不交并分解
$$
X_U =
W \amalg \coprod\nolimits_{i, j} V_{i, j}
$$
使得
\begin{enumerate}
\item $V_{i, j} \to U$ 是经过 $u$ 的闭浸入；
\item 纤维 $W_u$ 不含映到任意 $x_i$ 的点。
\end{enumerate}
特别地，若 $f^{-1}(\{s\}) = \{x_1, \ldots, x_n\}$，则纤维
$W_u$ 为空。
\end{lemma}

\begin{proof}
应用引理 \ref{etale-lemma-unramified-etale-local}。可假设 $U$ 仿射，
因而 $X_U$ 分离。于是 $V_{i, j} \to X_U$ 是闭映射，见
《概形态射》引理 \ref{morphisms-lemma-image-proper-scheme-closed}。
假设 $(i, j) \not = (i', j')$。则
$V_{i, j} \cap V_{i', j'}$ 是 $V_{i, j}$ 中的闭集，而且其在
$U$ 中的像不含 $u$。故缩小 $U$ 后，可假设
$V_{i, j} \cap V_{i', j'} = \emptyset$。此外，
$\bigcup V_{i, j}$ 是 $X_U$ 的开闭子概形，因而具有开闭补子概形
$W$。证明完成。
\end{proof}

\noindent
下述引理在某种意义上比前一引理弱得多，但在这里把它明确陈述出来可能
很有用。它断言：有限非分歧态射在基上平展（\'etale）局部地是闭浸入。

\begin{lemma}
\label{etale-lemma-finite-unramified-etale-local}
设 $f : X \to S$ 为有限非分歧概形态射，$s \in S$。
存在平展（\'etale）邻域 $(U, u) \to (S, s)$ 以及有限不交并分解
$$
X_U = \coprod\nolimits_j V_j
$$
使得每个 $V_j \to U$ 都是闭浸入。
\end{lemma}

\begin{proof}
由于 $X \to S$ 有限，$s$ 上的纤维是由 $X$ 中的点组成的有限集
$\{x_1, \ldots, x_n\}$。对该集合应用引理 \ref{etale-lemma-unramified-etale-local-technical}（有限态射分离，见
《概形态射》第 \ref{morphisms-section-integral} 节）。
$W$ 在 $U$ 中的像是闭子集（因为 $X_U \to U$ 有限，因而固有），
且不含 $u$。从 $U$ 中删去该像后，便有 $W = \emptyset$，如所需。
\end{proof}




\section{平展（\'Etale）态射的平展（\'etale）局部结构}
\label{etale-section-etale-local-etale}

\noindent
这看起来有些多余，但也许有助于形成对平展（\'etale）态射的直观认识。
只需照搬第 \ref{etale-section-unramified-etale-local} 节的结果，
并把``闭浸入''改成``同构''。

\begin{lemma}
\label{etale-lemma-etale-etale-local}
设 $f : X \to S$ 为概形态射。设
$x_1, \ldots, x_n \in X$ 是在 $S$ 中具有同一像 $s$ 的点。
假设 $f$ 在每个 $x_i$ 处平展（\'etale）。则存在平展
（\'etale）邻域 $(U, u) \to (S, s)$ 以及开子概形
$V_{i, j} \subset X_U$，$i = 1, \ldots, n$，$j = 1, \ldots, m_i$，
使得
\begin{enumerate}
\item $V_{i, j} \to U$ 是同构；
\item 除非 $i = i'$ 且 $j = j'$，否则 $u$ 不在
$V_{i, j} \cap V_{i', j'}$ 的像中；且
\item $(X_U)_u$ 中映到 $x_i$ 的任意点都属于某个 $V_{i, j}$。
\end{enumerate}
\end{lemma}

\begin{proof}
平展（\'etale）态射是非分歧态射，故可应用引理 \ref{etale-lemma-unramified-etale-local}。此时 $V_{i, j} \to U$
既是闭浸入又平展（\'etale），因而是开浸入，例如可由定理 \ref{etale-theorem-etale-radicial-open} 得出。用诸
$V_{i, j} \to U$ 的像的交代替 $U$，即得结论。
\end{proof}

\begin{lemma}
\label{etale-lemma-etale-etale-local-technical}
设 $f : X \to S$ 为概形态射。设
$x_1, \ldots, x_n \in X$ 是在 $S$ 中具有同一像 $s$ 的点。
假设 $f$ 分离，且 $f$ 在每个 $x_i$ 处平展（\'etale）。则存在平展
（\'etale）邻域 $(U, u) \to (S, s)$ 以及概形的有限不交并分解
$$
X_U =
W \amalg \coprod\nolimits_{i, j} V_{i, j}
$$
使得
\begin{enumerate}
\item $V_{i, j} \to U$ 是同构；
\item 纤维 $W_u$ 不含映到任意 $x_i$ 的点。
\end{enumerate}
特别地，若 $f^{-1}(\{s\}) = \{x_1, \ldots, x_n\}$，则纤维
$W_u$ 为空。
\end{lemma}

\begin{proof}
平展（\'etale）态射是非分歧态射，故可应用引理 \ref{etale-lemma-unramified-etale-local-technical}。与引理 \ref{etale-lemma-etale-etale-local} 的证明一样，诸态射
$V_{i, j} \to U$ 均为开浸入；用它们的像的交代替 $U$ 后即得结论。
\end{proof}

\noindent
下述引理在某种意义上比前一引理弱得多，但在这里把它明确陈述出来可能
很有用。它断言：有限平展（\'etale）态射在基上平展（\'etale）局部地
是一个``拓扑覆盖空间''，即基的若干个副本的有限乘积。

\begin{lemma}
\label{etale-lemma-finite-etale-etale-local}
设 $f : X \to S$ 为有限平展（\'etale）概形态射，$s \in S$。
存在平展（\'etale）邻域 $(U, u) \to (S, s)$ 以及有限不交并分解
$$
X_U = \coprod\nolimits_j V_j
$$
使得每个 $V_j \to U$ 都是同构。
\end{lemma}

\begin{proof}
平展（\'etale）态射是非分歧态射，故可应用引理 \ref{etale-lemma-finite-unramified-etale-local}。与引理 \ref{etale-lemma-etale-etale-local} 的证明一样，可知
$V_{i, j} \to U$ 是开浸入；用它们的像的交代替 $U$ 后即得结论。
\end{proof}




\section{保持性质}
\label{etale-section-properties-permanence}

\noindent
下面给出 Noether 局部环的平展（\'etale）同态（定义见定义 \ref{etale-definition-etale-ring}）的若干``保持''性质。关于 Noether
局部环的完备化与 Hensel 化的相应材料，见《代数进阶》第 \ref{more-algebra-section-permanence-completion} 节与
\ref{more-algebra-section-permanence-henselization}。

\begin{lemma}
\label{etale-lemma-etale-dimension}
设 $A$, $B$ 为 Noether 局部环，且 $A \to B$ 为局部环的平展
（\'etale）同态。则 $\dim(A) = \dim(B)$。
\end{lemma}

\begin{proof}
例如见《交换代数》引理 \ref{algebra-lemma-dimension-base-fibre-equals-total}。
\end{proof}

\begin{proposition}
\label{etale-proposition-etale-depth}
设 $A$, $B$ 为 Noether 局部环，且 $f : A \to B$ 为局部环的平展
（\'etale）同态。则 $\text{depth}(A) = \text{depth}(B)$。
\end{proposition}

\begin{proof}
见《交换代数》引理 \ref{algebra-lemma-apply-grothendieck}。
\end{proof}

\begin{proposition}
\label{etale-proposition-etale-CM}
\begin{slogan}
Cohen-Macaulay 性质沿平展（\'etale）映射上升且下降。
\end{slogan}
设 $A$, $B$ 为 Noether 局部环，且 $f : A \to B$ 为局部环的平展
（\'etale）同态。则 $A$ 是 Cohen-Macaulay 环，当且仅当 $B$ 也是。
\end{proposition}

\begin{proof}
局部环 $A$ 是 Cohen-Macaulay 环，当且仅当
$\dim(A) = \text{depth}(A)$。由于这两个不变量在平展（\'etale）
扩张下均保持不变，结论随即成立。
\end{proof}

\begin{proposition}
\label{etale-proposition-etale-regular}
设 $A$, $B$ 为 Noether 局部环，且 $f : A \to B$ 为局部环的平展
（\'etale）同态。则 $A$ 正则，当且仅当 $B$ 正则。
\end{proposition}

\begin{proof}
若 $B$ 正则，则由《交换代数》引理 \ref{algebra-lemma-flat-under-regular}，$A$ 正则。反过来假设 $A$
正则。令 $\mathfrak m$ 为 $A$ 的极大理想。于是
$\dim_{\kappa(\mathfrak m)} \mathfrak m/\mathfrak m^2 =
\dim(A) = \dim(B)$（见引理 \ref{etale-lemma-etale-dimension}）。
另一方面，$\mathfrak mB$ 是 $B$ 的极大理想，因而
$\mathfrak m_B/\mathfrak m_B = \mathfrak mB/\mathfrak m^2B$
至多由 $\dim(B)$ 个元素生成。因此 $B$ 正则。
（也可使用稍一般的《交换代数》引理 \ref{algebra-lemma-flat-over-regular-with-regular-fibre}。）
\end{proof}


\begin{proposition}
\label{etale-proposition-etale-reduced}
设 $A$, $B$ 为 Noether 局部环，且 $f : A \to B$ 为局部环的平展
（\'etale）同态。则 $A$ 约化，当且仅当 $B$ 约化。
\end{proposition}

\begin{proof}
由 $A \to B$ 忠实平坦可立即看出：若 $B$ 约化，则 $A$ 也约化。
另见《交换代数》引理 \ref{algebra-lemma-descent-reduced}。
反过来假设 $A$ 约化。由假设，$B$ 是某个有限型 $A$-代数 $B'$
在素理想 $\mathfrak q$ 处的局部化。以一个局部化代替 $B'$ 后，
可假设 $B'$ 在 $A$ 上平展（\'etale），见引理 \ref{etale-lemma-characterize-etale-Noetherian}。于是《交换代数》引理 \ref{algebra-lemma-reduced-goes-up} 可用于 $A \to B'$，从而 $B'$
约化。因此 $B$ 约化。
\end{proof}

\begin{remark}
\label{etale-remark-technicality-needed}
如果采用更弱的平展（\'etale）局部环映射 $A \to B$ 定义，即去掉
$B$ 在 $A$ 上本质有限型这一假设，则上述关于``约化性''的结论不再成立。
确实，某个 Noether 局部整环 $A$ 可能具有非约化的完备化
$A^\wedge$，见《例子》第 \ref{examples-section-local-completion-nonreduced} 节。但环映射
$A \to A^\wedge$ 平坦，$\mathfrak m_AA^\wedge$ 是
$A^\wedge$ 的极大理想，而且 $A$ 与 $A^\wedge$ 当然具有相同的
剩余域。这就是为什么必须只对本质有限型的环扩张考虑这个概念
（若 $A$ 非 Noether，则应为本质有限呈示）。
\end{remark}

\begin{proposition}
\label{etale-proposition-etale-normal}
\begin{reference}
\cite[Expose I, Theorem 9.5 part (i)]{SGA1}
\end{reference}
设 $A$, $B$ 为 Noether 局部环，且 $f : A \to B$ 为局部环的平展
（\'etale）同态。则 $A$ 是整闭整环，当且仅当 $B$ 也是。
\end{proposition}

\begin{proof}
关于整闭性的下降，见《交换代数》引理 \ref{algebra-lemma-descent-normal}。反过来假设 $A$ 整闭。由假设，
$B$ 是某个有限型 $A$-代数 $B'$ 在素理想 $\mathfrak q$ 处的局部化。
以一个局部化代替 $B'$ 后，可假设 $B'$ 在 $A$ 上平展
（\'etale），见引理 \ref{etale-lemma-characterize-etale-Noetherian}。
于是《交换代数》引理 \ref{algebra-lemma-normal-goes-up} 可用于
$A \to B'$，从而 $B'$ 整闭。因此 $B$ 是整闭整环。
\end{proof}

\noindent
以上命题说明了为什么我们希望把平展（\'etale）映射理解为``局部同构''。
还有一项性质极好地说明了这个定义确实是``正确''的：对于有限型
$\mathbf{C}$-概形，一个态射平展（\'etale），当且仅当相应解析空间上的
态射（即赋予复拓扑的 $\mathbf{C}$-值点之间的态射）在解析意义下为
局部同构（即在源上局部地为开嵌入）。这可借助结构定理以及解析化与
完备局部环的形成相交换这一事实来证明——细节留给读者。








\section{平展（\'etale）态射的下降}
\label{etale-section-descending-etale}

\noindent
为理解本节所用语言，建议读者先参阅《下降》第 \ref{descent-section-descent-datum} 节。设 $f : X \to S$ 为概形态射。
考虑拉回函子
\begin{equation}
\label{etale-equation-descent-etale}
U\text{ 是在 }S\text{ 上的平展概形} \longrightarrow
\begin{matrix}
\text{下降数据 }(V, \varphi)\text{ 相对于 }X/S \\
\text{ 其中 }V\text{ 在 }X\text{ 上平展}
\end{matrix}
\end{equation}
它把 $U$ 送到典范下降数据 $(X \times_S U, can)$。

\begin{lemma}
\label{etale-lemma-faithful}
若 $f : X \to S$ 满射，则函子
(\ref{etale-equation-descent-etale}) 忠实。
\end{lemma}

\begin{proof}
设 $a, b : U_1 \to U_2$ 是两个 $S$ 上平展（\'etale）概形之间的
态射，并假设 $a$ 与 $b$ 到 $X$ 的基变换相同。须证 $a = b$。
由命题 \ref{etale-proposition-equality}，只需证明 $a$ 与 $b$
在点和剩余域上相同。这是显然的，因为对每个 $u \in U_1$，
都可找到映到 $u$ 的点 $v \in X \times_S U_1$。
\end{proof}

\begin{lemma}
\label{etale-lemma-fully-faithful}
假设 $f : X \to S$ 是商映射，且 $f$ 的任意平展（\'etale）基变换
也是商映射。则函子 (\ref{etale-equation-descent-etale}) 全忠实。
\end{lemma}

\begin{proof}
由引理 \ref{etale-lemma-faithful}，该函子忠实。设 $U_1 \to S$ 与
$U_2 \to S$ 为平展（\'etale）态射，并设
$a : X \times_S U_1 \to X \times_S U_2$ 是与典范下降数据相容的
态射。我们将证明 $a$ 是某个态射 $U_1 \to U_2$ 的基变换。

\medskip\noindent
设 $U'_2 \subset U_2$ 为开子概形，并考虑
$W = a^{-1}(X \times_S U'_2)$。这是 $X \times_S U_1$ 的开子概形，
且与 $V_1 = X \times_S U_1$ 上的典范下降数据相容。这意味着，
$W$ 在两个投影 $V_1 \times_{U_1} V_1 \to V_1$ 下的逆像相同。
由于 $V_1 \to U_1$ 满射（它是 $X \to S$ 的基变换），
$W$ 是某个子集 $U'_1 \subset U_1$ 的逆像。又因 $W$ 开，
对 $f$ 的假设蕴含 $U'_1 \subset U_1$ 也开。

\medskip\noindent
设 $U_2 = \bigcup U_{2, i}$ 是仿射开覆盖。由上一段得到开覆盖
$U_1 = \bigcup U_{1, i}$，使得
$X \times_S U_{1, i} = a^{-1}(X \times_S U_{2, i})$。
若能证明存在态射 $U_{1, i} \to U_{2, i}$，其基变换是
$a_i : X \times_S U_{1, i} \to X \times_S U_{2, i}$，
则可利用忠实性把这些态射黏合成 $U_1 \to U_2$。这样便约化到
$U_2$ 仿射的情形。特别地，$U_2 \to S$ 分离
（《概形》引理 \ref{schemes-lemma-compose-after-separated}）。

\medskip\noindent
假设 $U_2 \to S$ 分离。由《概形》引理 \ref{schemes-lemma-semi-diagonal}，$a$ 的图 $\Gamma_a$ 是
$$
V = (X \times_S U_1) \times_X (X \times_S U_2) = X \times_S U_1 \times_S U_2
$$
的闭子概形。另一方面，该图也是开的，例如因为它是某个平展
（\'etale）态射的截面（命题 \ref{etale-proposition-properties-sections}）。由于 $a$ 是下降数据的态射，
$\Gamma_a \subset V$ 在两个投影
$V \times_{U_1 \times_S U_2} V \to V$ 下的逆像相同。
因此，与证明第二段中的论证相同，可以找到开闭子概形
$\Gamma \subset U_1 \times_S U_2$，其到 $X$ 的基变换为
$\Gamma_a$。于是 $\Gamma \to U_1$ 是平展（\'etale）态射，
且到 $X$ 的基变换是同构。这意味着 $\Gamma \to U_1$ 泛双射，
故由定理 \ref{etale-theorem-etale-radicial-open} 可知它是同构。
因此 $\Gamma$ 是某个态射 $U_1 \to U_2$ 的图，而该态射的基变换
正是 $a$，如所需。
\end{proof}

\begin{lemma}
\label{etale-lemma-fully-faithful-cases}
设 $f : X \to S$ 为概形态射。在下列情形中，函子
(\ref{etale-equation-descent-etale}) 全忠实：
\begin{enumerate}
\item $f$ 满射且泛闭（例如有限、整或固有）；
\item $f$ 满射且泛开（例如局部有限呈示且平坦、光滑或平展）；
\item $f$ 满射、拟紧且平坦。
\end{enumerate}
\end{lemma}

\begin{proof}
这由引理 \ref{etale-lemma-fully-faithful} 得出。例如，拓扑空间的闭满射是
商映射（《拓扑学》引理 \ref{topology-lemma-closed-morphism-quotient-topology}）。
有限、整和固有态射都是泛闭的；见《概形态射》引理 \ref{morphisms-lemma-integral-universally-closed}、
\ref{morphisms-lemma-finite-proper} 以及定义 \ref{morphisms-definition-proper}。另一方面，拓扑空间的开满射是商映射
（《拓扑学》引理 \ref{topology-lemma-open-morphism-quotient-topology}）。
平坦且局部有限呈示的态射、光滑态射和平展（\'etale）态射都是泛开的；
见《概形态射》引理 \ref{morphisms-lemma-fppf-open}、
\ref{morphisms-lemma-smooth-open} 与
\ref{morphisms-lemma-etale-open}。满射、拟紧、平坦态射的情形由
《概形态射》引理 \ref{morphisms-lemma-fpqc-quotient-topology} 得出。
\end{proof}

\begin{lemma}
\label{etale-lemma-reduce-to-affine}
设 $f : X \to S$ 为概形态射，$(V, \varphi)$ 为相对于 $X/S$ 的
下降数据，其中 $V \to X$ 平展（\'etale）。设
$S = \bigcup S_i$ 为开覆盖。假设
\begin{enumerate}
\item 下降数据 $(V, \varphi)$ 到 $X \times_S S_i/S_i$ 的拉回有效；
\item 对 $X \times_S (S_i \cap S_j) \to (S_i \cap S_j)$，
函子 (\ref{etale-equation-descent-etale}) 全忠实；且
\item 对 $X \times_S (S_i \cap S_j \cap S_k) \to (S_i \cap S_j \cap S_k)$，
函子 (\ref{etale-equation-descent-etale}) 忠实。
\end{enumerate}
则 $(V, \varphi)$ 有效。
\end{lemma}

\begin{proof}
（回忆：下降数据的拉回定义于《下降》定义 \ref{descent-definition-pullback-functor}。）
令 $X_i = X \times_S S_i$，并记 $(V_i, \varphi_i)$ 为
$(V, \varphi)$ 到 $X_i/S_i$ 的拉回。由假设 (1)，可找到平展
（\'etale）态射 $U_i \to S_i$ 以及同构
$X_i \times_{S_i} U_i \to V_i$，它与 $can$ 和 $\varphi_i$ 相容。
由假设 (2)，得到同构
$\psi_{ij} : U_i \times_{S_i} (S_i \cap S_j) \to
U_j \times_{S_j} (S_i \cap S_j)$。由假设 (3)，这些同构满足余圈条件，
故 $(U_i, \psi_{ij})$ 是 Zariski 覆盖 $\{S_i \to S\}$ 的下降数据。
于是《下降》引理 \ref{descent-lemma-Zariski-refinement-coverings-equivalence}
（本质上只是《概形》第 \ref{schemes-section-glueing-schemes} 节的重述）表明：存在概形态射
$U \to S$ 以及与 $\psi_{ij}$ 相容的同构
$U \times_S S_i \to U_i$。这些同构 $U \times_S S_i \to U_i$ 确定相应的同构
$X_i \times_S U \to V_i$；后者黏合成与典范下降数据和
$\varphi$ 相容的态射 $X \times_S U \to V$。
\end{proof}

\begin{lemma}
\label{etale-lemma-split-henselian}
设 $(A, I)$ 为 Hensel 对。设 $U \to \Spec(A)$ 为拟紧、分离且平展
（\'etale）的态射，并且
$U \times_{\Spec(A)} \Spec(A/I) \to \Spec(A/I)$ 有限。则
$$
U = U_{fin} \amalg U_{away}
$$
其中 $U_{fin} \to \Spec(A)$ 有限，而 $U_{away}$ 不含位于 $Z$
上方的点。
\end{lemma}

\begin{proof}
由 Zariski 主定理，概形 $U$ 拟仿射。事实上，可以找到开浸入
$U \to T$，其中 $T$ 仿射且 $T \to \Spec(A)$ 有限；见
《态射进阶》引理 \ref{more-morphisms-lemma-quasi-finite-separated-pass-through-finite}。
写 $Z = \Spec(A/I)$，并记 $U_Z \to T_Z$ 为其基变换。
由于 $U_Z \to Z$ 有限，$U_Z \to T_Z$ 既闭又开。因此，由
《代数进阶》引理 \ref{more-algebra-lemma-characterize-henselian-pair}，得到唯一分解
$T = T' \amalg T''$，其中 $T'_Z = U_Z$。令
$U_{fin} = U \cap T'$ 与 $U_{away} = U \cap T''$。由于
$T'_Z \subset U_Z$，$T'$ 的所有闭点都在 $U$ 中，故
$T' \subset U$，从而 $U_{fin} = T'$，于是
$U_{fin} \to \Spec(A)$ 有限。这里略去该分解唯一性的证明。
\end{proof}

\begin{proposition}
\label{etale-proposition-effective}
设 $f : X \to S$ 为满的整态射。函子
(\ref{etale-equation-descent-etale}) 诱导等价
$$
\begin{matrix}
\text{拟紧概形，}\\
\text{分离且在 }S\text{ 上平展}
\end{matrix}
\longrightarrow
\begin{matrix}
\text{下降数据 }(V, \varphi)\text{ 相对于 }X/S\text{ 其中}\\
V\text{ 拟紧、分离且在 }X\text{ 上平展}
\end{matrix}
$$
\end{proposition}

\begin{proof}
由引理 \ref{etale-lemma-fully-faithful-cases}，函子
(\ref{etale-equation-descent-etale}) 全忠实，而且在任意基变换
$S \to S'$ 后仍然如此。设 $(V, \varphi)$ 是相对于 $X/S$ 的下降数据，
其中 $V \to X$ 拟紧、分离且平展（\'etale）。由引理 \ref{etale-lemma-reduce-to-affine} 可知，只需在 $S$ 上 Zariski 局部地证明
有效性。特别地，可以并且确实假设 $S$ 仿射。

\medskip\noindent
若 $S$ 仿射，则存在有向集 $\Lambda$ 以及有限型
$\Spec(\mathbf{Z})$-仿射概形之间有限态射所成的逆系统
$X_\lambda \to S_\lambda$，使得
$(X \to S) = \lim (X_\lambda \to S_\lambda)$；见《交换代数》引理 \ref{algebra-lemma-limit-integral}。由于极限与极限交换，可得
$X \times_S X = \lim X_\lambda \times_{S_\lambda} X_\lambda$
以及
$X \times_S X \times_S X = \lim
X_\lambda \times_{S_\lambda} X_\lambda \times_{S_\lambda} X_\lambda$。
注意 $V \to X$ 是有限呈示态射。使用《概形的极限》引理 \ref{limits-lemma-descend-finite-presentation}，可找到
$\lambda$ 以及相对于 $X_\lambda/S_\lambda$ 的下降数据
$(V_\lambda, \varphi_\lambda)$，其到 $X/S$ 的拉回为
$(V, \varphi)$。当然，只需证明 $(V_\lambda, \varphi_\lambda)$ 有效。
由构造，$V_\lambda$ 拟紧。必要时增大 $\lambda$ 后，可假设
$V_\lambda \to X_\lambda$ 分离且平展（\'etale）；见《概形的极限》引理 \ref{limits-lemma-descend-separated-finite-presentation} 与
\ref{limits-lemma-descend-etale}。因此，可假设 $f$ 有限且满，
而 $S$ 是有限型 $\mathbf{Z}$-仿射概形。

\medskip\noindent
考虑开子集 $S' \subset S$，使得 $(V, \varphi)$ 到
$X' = X \times_S S'$ 的拉回 $(V', \varphi')$ 有效。下面将证明：
若 $S' \not = S$，则存在一个严格更大的开集，使下降数据在其上有效。
由于 $S$ 是 Noether 概形（因而其底拓扑空间也是 Noether 的），
这将完成证明。设 $\xi \in S$ 是闭子集 $Z = S \setminus S'$
某个不可约分支的一般点。若 $\xi \in S'' \subset S$ 是一个使下降数据
在其上有效的开集，则由第一段的黏合论证，下降数据在
$S' \cup S''$ 上有效。因此，在证明余下部分中，可以用 $\xi$ 的
一个仿射开邻域代替 $S$。

\medskip\noindent
第一次这样代替后，可假设 $Z$ 不可约，且以 $Z$ 为一般点。赋予 $Z$
约化的诱导闭子概形结构。再次缩小后，可假设
$X_Z = X \times_S Z = f^{-1}(Z) \to Z$ 平坦，见《概形态射》命题 \ref{morphisms-proposition-generic-flatness}。
令 $(V_Z, \varphi_Z)$ 为下降数据到 $X_Z/Z$ 的拉回。由
《态射进阶》引理 \ref{more-morphisms-lemma-separated-locally-quasi-finite-morphisms-fppf-descend}，
该下降数据有效，并得到平展（\'etale）态射 $U_Z \to Z$，
其基变换以与下降数据相容的方式同构于 $V_Z$。当然，
$U_Z \to Z$ 拟紧且分离（《下降》引理 \ref{descent-lemma-descending-property-quasi-compact} 与
\ref{descent-lemma-descending-property-separated}）。因此，再次缩小后，
可假设 $U_Z \to Z$ 有限，见《概形态射》引理 \ref{morphisms-lemma-generically-finite}。

\medskip\noindent
设 $S = \Spec(A)$，并设 $I \subset A$ 是与 $Z \subset S$ 对应的
素理想。令 $(A^h, IA^h)$ 为对 $(A, I)$ 的 Hensel 化。记
$S^h = \Spec(A^h)$，$Z^h = V(IA^h) \cong Z$。我们断言，只需在
基变换到 $S^h$ 后证明有效性。事实上，
$\{S^h \to S, S' \to S\}$ 是 fpqc 覆盖（由《代数进阶》引理 \ref{more-algebra-lemma-henselization-flat}，
$A \to A^h$ 平坦），而由《态射进阶》引理 \ref{more-morphisms-lemma-separated-locally-quasi-finite-morphisms-fppf-descend}，
分离平展（\'etale）态射满足 fpqc 下降。具体地，若
$U^h \to S^h$ 与 $U' \to S'$ 分别是对应于拉回
$(V^h, \varphi^h)$ 与 $(V', \varphi')$ 的对象，则所需同构
$$
U^h \times_S S^h \to S^h \times_S V^h
\quad\text{且}\quad
U^h \times_S S' \to S^h \times_S U'
$$
由第一段指出的全忠实性得到。这样便约化到下一段所述情形。

\medskip\noindent
现在 $S = \Spec(A)$，$Z = V(I)$，$S' = S \setminus Z$，其中
$(A, I)$ 为 Hensel 对；存在与下降数据 $(V', \varphi')$ 对应的
$U' \to S'$，以及与下降数据 $(V_Z, \varphi_Z)$ 对应的有限平展
（\'etale）态射 $U_Z \to Z$。此时不再有 $A$ 在
$\mathbf{Z}$ 上有限型这一条件，但余下论证甚至不会用到 $A$
是 Noether 环。由《代数进阶》引理 \ref{more-algebra-lemma-finite-etale-equivalence}，可找到有限平展
（\'etale）态射 $U_{fin} \to S$，其在 $Z$ 上的限制同构于
$U_Z \to Z$。写 $X = \Spec(B)$ 和 $Y = V(IB)$。由于 $(B, IB)$
为 Hensel 对（《代数进阶》引理 \ref{more-algebra-lemma-integral-over-henselian-pair}），且
$V \to X$ 到 $Y$ 的限制有限（因为它是 $U_Z \to Z$ 的基变换），
可知存在典范不交并分解
$$
V = V_{fin} \amalg V_{away}
$$
其中 $V_{fin} \to X$ 有限，而 $V_{away}$ 不含位于 $Y$ 上方的点；
见引理 \ref{etale-lemma-split-henselian}。利用该分解在
$X \times_S X$ 上的唯一性，可知 $\varphi$ 保持此分解，从而在
下降数据范畴中有
$$
(V, \varphi) = (V_{fin}, \varphi_{fin}) \amalg (V_{away}, \varphi_{away})
$$
。由《代数进阶》引理 \ref{more-algebra-lemma-finite-etale-equivalence}，存在唯一同构
$$
X \times_S U_{fin} \longrightarrow V_{fin}
$$
与 $Y$ 上给定的同构 $Y \times_Z U_Z \to V \times_X Y$ 相容。
由唯一性可知，该同构与下降数据相容，即
$(X \times_S U_{fin}, can) \cong (V_{fin}, \varphi_{fin})$。
记 $U'_{fin} = U_{fin} \times_S S'$。由全忠实性，得到态射
$U'_{fin} \to U'$，它是一个开闭子概形的嵌入。然后令
$U = U_{fin} \amalg_{U'_{fin}} U'$（按《概形》第 \ref{schemes-section-glueing-schemes} 节黏合概形）。
态射 $X \times_S U_{fin} \to V$ 与 $X \times_S U' \to V$
黏合成态射 $X \times_S U \to V$，它正是所需同构。
\end{proof}





\section{正规交叉除子}
\label{etale-section-normal-crossings}

\noindent
先给出定义。

\begin{definition}
\label{etale-definition-strict-normal-crossings}
设 $X$ 为局部 Noether 概形。$X$ 上的{\it 严格正规交叉除子}是一个
有效 Cartier 除子 $D \subset X$，使得对每个 $p \in D$，
局部环 $\mathcal{O}_{X, p}$ 正则，并且存在正则参数系
$x_1, \ldots, x_d \in \mathfrak m_p$ 以及 $1 \leq r \leq d$，
使 $D$ 在 $\mathcal{O}_{X, p}$ 中由 $x_1 \ldots x_r$ 截出。
\end{definition}

\noindent
我们经常遇到局部 Noether 概形 $X$ 上的有效 Cartier 除子 $E$，
其满足：存在严格正规交叉除子 $D$，使得在集合意义下
$E \subset D$。此时有
$E = \sum a_i D_i$，其中 $a_i \geq 0$，而
$D = \bigcup_{i \in I} D_i$ 是 $D$ 的不可约分支分解。
注意 $D' = \bigcup_{a_i > 0} D_i$ 是严格正规交叉除子，且在集合
意义下 $E = D'$。出现上述情形时，称 $E$
{\it 支撑于一个严格正规交叉除子上}。

\begin{lemma}
\label{etale-lemma-strict-normal-crossings}
设 $X$ 为局部 Noether 概形，$D \subset X$ 为有效 Cartier 除子。
设 $D_i \subset D$，$i \in I$，为其不可约分支，并把它们视为
$X$ 的约化闭子概形。下列条件等价：
\begin{enumerate}
\item $D$ 是严格正规交叉除子；且
\item $D$ 约化，每个 $D_i$ 都是有效 Cartier 除子，并且对有限子集
$J \subset I$，概形论交
$D_J = \bigcap_{j \in J} D_j$ 是正则概形，其每个不可约分支在
$X$ 中的余维数均为 $|J|$。
\end{enumerate}
\end{lemma}

\begin{proof}
假设 $D$ 是严格正规交叉除子。取 $p \in D$，并按定义 \ref{etale-definition-strict-normal-crossings} 选取正则参数系
$x_1, \ldots, x_d \in \mathfrak m_p$ 以及
$1 \leq r \leq d$。由于 $\mathcal{O}_{X, p}/(x_i)$ 是正则局部环
（特别地，是整环），穿过 $p$ 的 $D$ 的不可约分支
$D_1, \ldots, D_r$ 与 $\mathcal{O}_{X, p}$ 的高为一的素理想
$(x_1), \ldots, (x_r)$ 按 $1$ 对 $1$ 对应。由《交换代数》引理 \ref{algebra-lemma-regular-ring-CM}，诸交
$D_{i_1} \cap \ldots \cap D_{i_s}$ 在 $p$ 的某个开邻域中余维数为
$s$，而且该交在 $p$ 处具有正则局部环。由于这对每个
$p \in D$ 都成立，故条件 (2) 成立。

\medskip\noindent
反过来假设 (2)。设 $p \in D$。由于 $\mathcal{O}_{X, p}$
有限维，点 $p$ 至多属于 $\dim(\mathcal{O}_{X, p})$ 个分支
$D_i$。设对某个 $r \geq 1$ 有 $p \in D_1, \ldots, D_r$。
令 $x_1, \ldots, x_r \in \mathfrak m_p$ 分别为
$D_1, \ldots, D_r$ 的局部方程。则 $x_1$ 是
$\mathcal{O}_{X, p}$ 中的非零因子，且
$\mathcal{O}_{X, p}/(x_1) = \mathcal{O}_{D_1, p}$ 正则。
因此 $\mathcal{O}_{X, p}$ 正则，见《交换代数》引理 \ref{algebra-lemma-regular-mod-x}。由于
$D_1 \cap \ldots \cap D_r$ 是正则（因而整闭）概形，它是其不可约
分支的不交并（《概形的性质》引理 \ref{properties-lemma-normal-Noetherian}）。令
$Z \subset D_1 \cap \ldots \cap D_r$ 为含 $p$ 的不可约分支。
于是 $\mathcal{O}_{Z, p} = \mathcal{O}_{X, p}/(x_1, \ldots, x_r)$
正则且余维数为 $r$（注意，已经知道 $\mathcal{O}_{X, p}$ 正则，
从而是 Cohen-Macaulay 环；又因为该环链状，余维数没有歧义，见
《交换代数》引理 \ref{algebra-lemma-regular-ring-CM} 与
\ref{algebra-lemma-CM-dim-formula}）。因此
$\dim(\mathcal{O}_{Z, p}) = \dim(\mathcal{O}_{X, p}) - r$。
再选取 $x_{r + 1}, \ldots, x_n \in \mathfrak m_p$，使其像构成
$\mathfrak m_{Z, p}$ 的极小生成系。由 Nakayama 引理，
$\mathfrak m_p = (x_1, \ldots, x_n)$，从而 $D$ 是正规交叉除子。
\end{proof}

\begin{lemma}
\label{etale-lemma-smooth-pullback-strict-normal-crossings}
\begin{slogan}
严格正规交叉除子沿光滑态射的拉回仍是严格正规交叉除子。
\end{slogan}
设 $X$ 为局部 Noether 概形，$D \subset X$ 为严格正规交叉除子。
若 $f : Y \to X$ 是光滑概形态射，则拉回 $f^*D$ 是 $Y$ 上的
严格正规交叉除子。
\end{lemma}

\begin{proof}
由于 $f$ 平坦，拉回由《除子》引理 \ref{divisors-lemma-pullback-effective-Cartier-defined} 定义，
故该陈述有意义。设 $q \in f^*D$ 映到 $p \in D$。按定义 \ref{etale-definition-strict-normal-crossings} 选取正则参数系
$x_1, \ldots, x_d \in \mathfrak m_p$ 以及
$1 \leq r \leq d$。由于 $f$ 光滑，局部环同态
$\mathcal{O}_{X, p} \to \mathcal{O}_{Y, q}$ 平坦，而且纤维环
$$
\mathcal{O}_{Y, q}/\mathfrak m_p \mathcal{O}_{Y, q} =
\mathcal{O}_{Y_p, q}
$$
是正则局部环（例如见《交换代数》引理 \ref{algebra-lemma-characterize-smooth-over-field}）。选取
$y_1, \ldots, y_n \in \mathfrak m_q$，使其像构成
$\mathcal{O}_{Y_p, q}$ 的正则参数系。于是
$x_1, \ldots, x_d, y_1, \ldots, y_n$ 生成极大理想
$\mathfrak m_q$。由《交换代数》引理 \ref{algebra-lemma-dimension-base-fibre-equals-total}，
$\mathcal{O}_{Y, q}$ 是维数为 $d + n$ 的正则局部环，而
$x_1, \ldots, x_d, y_1, \ldots, y_n$ 是其正则参数系。由于
$f^*D$ 在 $\mathcal{O}_{Y, q}$ 中由 $x_1 \ldots x_r$ 截出，
引理得证。
\end{proof}

\noindent
下面定义正规交叉除子。

\begin{definition}
\label{etale-definition-normal-crossings}
设 $X$ 为局部 Noether 概形。$X$ 上的{\it 正规交叉除子}是一个
有效 Cartier 除子 $D \subset X$，使得对每个 $p \in D$，
都存在像中含 $p$ 的平展（\'etale）态射 $U \to X$，并且
$D \times_X U$ 是 $U$ 上的严格正规交叉除子。
\end{definition}

\noindent
例如，$D = V(x^2 + y^2)$ 是 $\Spec(\mathbf{R}[x, y])$ 上的
正规交叉除子（但不是严格正规交叉除子），因为拉回到平展
（\'etale）覆盖 $\Spec(\mathbf{C}[x, y])$ 后得到
$(x - iy)(x + iy) = 0$。

\begin{lemma}
\label{etale-lemma-smooth-pullback-normal-crossings}
\begin{slogan}
正规交叉除子沿光滑态射的拉回仍是正规交叉除子。
\end{slogan}
设 $X$ 为局部 Noether 概形，$D \subset X$ 为正规交叉除子。
若 $f : Y \to X$ 是光滑概形态射，则拉回 $f^*D$ 是 $Y$ 上的
正规交叉除子。
\end{lemma}

\begin{proof}
由于 $f$ 平坦，拉回由《除子》引理 \ref{divisors-lemma-pullback-effective-Cartier-defined} 定义，
故该陈述有意义。设 $q \in f^*D$ 映到 $p \in D$。
按定义 \ref{etale-definition-normal-crossings}，选取像中含 $p$ 的
平展（\'etale）态射 $U \to X$，使得
$D \times_X U \subset U$ 为严格正规交叉除子。
令 $V = Y \times_X U$。于是 $V \to Y$ 是 $U \to X$ 基变换所得
的平展（\'etale）态射（《概形态射》引理 \ref{morphisms-lemma-base-change-etale}），而由引理 \ref{etale-lemma-smooth-pullback-strict-normal-crossings}，
拉回 $D \times_X V$ 是 $V$ 上的严格正规交叉除子。
因此已对 $q \in f^*D$ 验证定义 \ref{etale-definition-normal-crossings} 的条件，结论成立。
\end{proof}

\begin{lemma}
\label{etale-lemma-characterize-normal-crossings-normalization}
设 $X$ 为局部 Noether 概形，$D \subset X$ 为闭子概形。下列条件等价：
\begin{enumerate}
\item $D$ 是 $X$ 中的正规交叉除子；
\item $D$ 约化，正规化 $\nu : D^\nu \to D$ 非分歧，并且对任意
$n \geq 1$，概形
$$
Z_n = D^\nu \times_D \ldots \times_D D^\nu
\setminus \{(p_1, \ldots, p_n) \mid p_i = p_j\text{ 对某个 }i\not = j\}
$$
正则，而态射 $Z_n \to X$ 是局部完全交态射，其余法层局部自由且
秩为 $n$。
\end{enumerate}
\end{lemma}

\begin{proof}
先说明如何理解条件 (2)。非分歧态射的对角态射是开的
（《概形态射》引理 \ref{morphisms-lemma-diagonal-unramified-morphism}）。
另一方面，$D^\nu \to D$ 分离，故对角态射
$D^\nu \to D^\nu \times_D D^\nu$ 是闭的。因此 $Z_n$ 是
$D^\nu \times_D \ldots \times_D D^\nu$ 的开闭子概形。
另一方面，$Z_n \to X$ 非分歧，因为它是复合
$$
Z_n \to D^\nu \times_D \ldots \times_D D^\nu \to \ldots \to
D^\nu \times_D D^\nu \to D^\nu \to D \to X
$$
而其中每个箭头都非分歧。由于非分歧态射形式非分歧
（《态射进阶》引理 \ref{more-morphisms-lemma-unramified-formally-unramified}），
$Z_n \to X$ 具有余法层
$\mathcal{C}_n = \mathcal{C}_{Z_n/X}$，见《态射进阶》定义 \ref{more-morphisms-definition-universal-thickening}。

\medskip\noindent
由《态射进阶》引理 \ref{more-morphisms-lemma-normalization-and-smooth}，正规化的形成与
平展（\'etale）局部化交换。以下性质都可以平展（\'etale）局部地
检验：局部环正则；态射非分歧；态射为局部完全交；或者态射非分歧且
具有给定秩的局部自由余法层（见《代数进阶》引理 \ref{more-algebra-lemma-regular-etale-extension}，《下降》引理 \ref{descent-lemma-descending-property-unramified}，《态射进阶》引理 \ref{more-morphisms-lemma-descending-property-lci}，以及《下降》引理 \ref{descent-lemma-finite-locally-free-descends}）。

\medskip\noindent
由上一段的说明和正规交叉除子的定义，只需证明严格正规交叉除子
$D = \bigcup_{i \in I} D_i$ 满足 (2)。在此情形，
$D^\nu = \coprod D_i$，且 $D^\nu \to D$ 非分歧（非分歧性在源上
是局部的，而 $D_i \to D$ 是非分歧的闭浸入）。类似地，
$Z_1 = D^\nu \to X$ 是局部完全交态射，因为这可以在源上局部检验，
而每个态射 $D_i \to X$ 都是正则浸入：它是 Cartier 除子的包含
（见引理 \ref{etale-lemma-strict-normal-crossings} 与《态射进阶》引理 \ref{more-morphisms-lemma-regular-immersion-lci}）。
由于有效 Cartier 除子的余法层可逆，余法层的要求得以满足。
同样，对 $n \geq 2$，概形 $Z_n$ 是诸概形
$D_J = \bigcap_{j \in J} D_j$ 的不交并，其中 $J \subset I$
遍历所有基数为 $n$ 的子集。由于 $D_J \to X$ 是余维数为 $n$
的正则浸入（由严格正规交叉的定义，以及可通过《除子》引理 \ref{divisors-lemma-Noetherian-scheme-regular-ideal} 在茎上检验
这一事实），同理可知 $Z_n \to X$ 具有所需性质。略去若干细节。

\medskip\noindent
假设 (2)。设 $p \in D$。由于 $D^\nu \to D$ 非分歧，它是有限的
（由《概形态射》引理 \ref{morphisms-lemma-finite-integral}）。
故 $D^\nu \to X$ 有限且非分歧。由引理 \ref{etale-lemma-finite-unramified-etale-local} 和平展（\'etale）局部化
（根据第二段的讨论和正规交叉除子的定义，这是允许的），可约化到如下
情形：
$D^\nu = \coprod_{i \in I} D_i$，其中 $I$ 有限，且
$D_i \to U$ 为闭浸入。必要时缩小 $X$ 后，可假设对所有
$i \in I$ 都有 $p \in D_i$。条件
$Z_1 =  D^\nu \to X$ 是非分歧局部完全交态射，且其余法层局部自由、
秩为 $1$；这蕴含 $D_i \subset X$ 是有效 Cartier 除子，见
《态射进阶》引理 \ref{more-morphisms-lemma-lci} 与《除子》引理 \ref{divisors-lemma-regular-immersion-noetherian}。
为完成证明，可假设 $X = \Spec(A)$ 仿射，且 $D_i = V(f_i)$，
其中 $f_i \in A$ 是非零因子。若 $I = \{1, \ldots, r\}$，则
$p \in Z_r = V(f_1, \ldots, f_r)$。上面所引的同一结果表明，
$(f_1, \ldots, f_r)$ 是 $A$ 中的 Koszul 正则理想。由于余法层
秩为 $r$，$f_1, \ldots, f_r$ 是定义
$\mathcal{O}_{X, p}$ 中 $Z_r$ 之理想的极小生成系。这蕴含
$f_1, \ldots, f_r$ 是 $\mathcal{O}_{X, p}$ 中的正则序列，且
$\mathcal{O}_{X, p}/(f_1, \ldots, f_r)$ 正则。因此，由《交换代数》引理 \ref{algebra-lemma-regular-mod-x}，
$f_1, \ldots, f_r$ 可扩充成 $\mathcal{O}_{X, p}$ 的正则参数系，
证明完成。
\end{proof}

\begin{lemma}
\label{etale-lemma-characterize-normal-crossings}
设 $X$ 为局部 Noether 概形，$D \subset X$ 为闭子概形。若 $X$
是 J-2 或 Nagata 概形，则下列条件等价：
\begin{enumerate}
\item $D$ 是 $X$ 中的正规交叉除子；
\item 对每个 $p \in D$，$D$ 到严格 Hensel 化
$\mathcal{O}_{X, p}^{sh}$ 的谱的拉回是严格正规交叉除子。
\end{enumerate}
\end{lemma}

\begin{proof}
蕴含 (1) $\Rightarrow$ (2) 是直接的，而且不需要 $X$ 为 J-2 或
Nagata 概形这一假设。事实上，设 $p \in D$，选取平展（\'etale）
邻域 $(U, u) \to (X, p)$，使得 $D$ 的拉回是 $U$ 上的严格正规
交叉除子。于是 $\mathcal{O}_{X, p}^{sh} = \mathcal{O}_{U, u}^{sh}$；
由于在 $\mathcal{O}_{U, u}$ 中已经如此，可知 $D$ 在
$\Spec(\mathcal{O}_{U, u}^{sh})$ 上的迹由某个正则参数系的一部分截出。

\medskip\noindent
为证明反向蕴含，将使用引理 \ref{etale-lemma-characterize-normal-crossings-normalization} 的判据。
注意正规化 $D^\nu \to D$ 的形成与严格 Hensel 化交换，见
《态射进阶》引理 \ref{more-morphisms-lemma-normalization-and-henselization}。
若能证明 $D^\nu \to D$ 有限，则 $D^\nu \to D$ 和诸概形
$Z_n$ 都满足所需性质，因为这些性质都可以在局部环层面检验
（但态射 $D^\nu \to D$ 的有限性不能在局部环上检验）。
略去详细验证。

\medskip\noindent
若 $X$ 为 Nagata 概形，则由《概形态射》引理 \ref{morphisms-lemma-nagata-normalization}，
$D^\nu \to D$ 有限。

\medskip\noindent
假设 $X$ 是 J-2 概形。选取点 $p \in D$。我们将证明
$D^\nu \to D$ 在 $p$ 的某个邻域上有限。由假设，
$\mathcal{O}_{X, p}^{sh}$ 存在正则参数系
$f_1, \ldots, f_d$ 以及 $1 \leq r \leq d$，使得 $D$ 在
$\Spec(\mathcal{O}_{X, p}^{sh})$ 上的迹由 $f_1 \ldots f_r$ 截出。
于是
$$
D^\nu \times_X \Spec(\mathcal{O}_{X, p}^{sh}) = 
\coprod\nolimits_{i = 1, \ldots, r} V(f_i)
$$
选取仿射平展（\'etale）邻域 $(U, u) \to (X, p)$，使得 $f_i$
来自 $f_i \in \mathcal{O}_U(U)$。令 $D_i = V(f_i) \subset U$。
$\mathcal{O}_{D_i, u}$ 的严格 Hensel 化为
$\mathcal{O}_{X, p}^{sh}/(f_i)$，而后者正则。因此
$\mathcal{O}_{D_i, u}$ 正则（例如由《代数进阶》引理 \ref{more-algebra-lemma-henselization-regular}）。由于 $X$ 是 J-2
概形，正则轨迹在 $D_i$ 中开。因此，以某个 Zariski 开集代替
$U$ 后，可假设对每个 $i$，$D_i$ 都正则。于是
$$
\coprod\nolimits_{i = 1, \ldots, r} D_i = D^\nu \times_X U
\longrightarrow D \times_X U
$$
是正规化态射，并且显然有限。换言之，已经找到 $(X, p)$ 的平展
（\'etale）邻域 $(U, u)$，使得 $D^\nu \to D$ 到该邻域的基变换
有限。由下降（《下降》引理 \ref{descent-lemma-descending-property-finite}），这蕴含
$D^\nu \to D$ 有限，证明完成。
\end{proof}
% Generated mechanically from the frozen chow.tex target.
% Do not edit this generated file; edit the bound locale source instead.

% OK, start here.
%

\setcounter{chapter}{41}
\chapter{Chow 同调与陈类}


\phantomsection
\label{chow-section-phantom}




\section{引言}
\label{chow-section-introduction}

\noindent
本章讨论 Chow 同调群，以及把向量丛的陈类构造为运算 Chow 上同调群中元素的过程
（系数一律取自 $\mathbf{Z}$）。

\medskip\noindent
本章首先给出关键引理 \ref{chow-lemma-milnor-gersten-low-degree} 的一个尽可能简短的代数证明。
我们先定义 Herbrand 商（第 \ref{chow-section-periodic-complexes} 节），并在若干情形下计算它
（第 \ref{chow-section-calculation} 节）。随后证明一些简单的代数引理，说明经修改后存在适当的分解
（第 \ref{chow-section-preparation-tame-symbol} 节）。利用这些结果，我们在第
\ref{chow-section-tame-symbol} 节构造并定义驯顺符号。证明关键引理只需要驯顺符号最基本的性质；
该证明见第 \ref{chow-section-key-lemma} 节。

\medskip\noindent
接着，我们在第 \ref{chow-section-setup} 节引入本章其余部分所采用的基本设置。为了使材料更具挑战性，
我们决定处理一个比通常情形稍为一般的情况。具体而言，假设概形 $X$ 在一个固定的局部诺特基概形上
局部有限型，并且该基概形泛链状且配有一个维数函数。这些假设足以定义 Chow 同调群 $\CH_*(X)$，
以及以陈类同它们作帽积的作用。这表明，对这种基概形上局部有限型的代数栈，我们也应当能够作出同样的定义。

\medskip\noindent
随后我们依照 \cite{F} 的开头几章定义循环、平坦拉回、固有推前和有理等价；但对于所涉及循环的支撑，
我们的处理没有那么精细。

\medskip\noindent
我们与 \cite{F} 的论述方式有所不同：在第 \ref{chow-section-key} 节利用上述关键引理证明一个基本的交换关系。
由此可证，与可逆层取交这一运算可下降到有理等价，并且满足交换律；见第
\ref{chow-section-commutativity} 节。再次应用关键引理可证，有效 Cartier 除子的 Gysin 映射可下降到有理等价；
见第 \ref{chow-section-gysin} 节。证明这些结果以后，就可以直接定义向量丛的陈类，证明其可加性和分裂原理，
引入 Chern 特征与 Todd 类，并陈述 Grothendieck--Riemann--Roch 定理。

\medskip\noindent
本章有两个附录。附录 A（第 \ref{chow-section-appendix-A} 节）讨论驯顺符号的另一种较长构造，
以及与之相应的关键引理证明。最后，在附录 B（第 \ref{chow-section-appendix-chow} 节）中，
我们简要讨论它同凝聚层 $K$-理论的关系，并讨论若干爆破引理。建议读者参阅这两个附录的引言以了解详情。

\medskip\noindent
下一章将重新研究代数闭域上光滑射影簇的 Chow 群 $\CH_*(X)$。利用 \cite{Samuel}、
\cite{ChevalleyI} 和 \cite{ChevalleyII} 中的移动引理，以及 Serre 的 Tor 公式
（见 \cite{Serre_local_algebra} 或 \cite{Serre_algebre_locale}），我们将在 $\CH_*(X)$ 上定义环结构。
参见 Intersection Theory，第 \ref{intersection-section-introduction} 节及以下各节。








\section{周期复形与 Herbrand 商}
\label{chow-section-periodic-complexes}

\noindent
周期复形当然有一个非常一般的概念。可以要求映射具有周期性，也可以要求对象具有周期性。
我们会在需要时补充这些情形；目前只需要以下几种情况。

\begin{definition}
\label{chow-definition-periodic-complex}
设 $R$ 为环。
\begin{enumerate}
\item $R$ 上的一个 {\it $2$-周期复形}由四元组 $(M, N, \varphi, \psi)$ 给出，
其中 $M$、$N$ 是 $R$-模，$\varphi : M \to N$、$\psi : N \to M$ 是 $R$-模同态，并且
$$
\xymatrix{
\ldots \ar[r] &
M \ar[r]^\varphi &
N \ar[r]^\psi &
M \ar[r]^\varphi &
N \ar[r] & \ldots
}
$$
构成一个复形。在这种情形下，定义该复形的{\it 上同调模}为 $R$-模
$$
H^0(M, N, \varphi, \psi) = \Ker(\varphi)/\Im(\psi)
\quad\text{且}\quad
H^1(M, N, \varphi, \psi) = \Ker(\psi)/\Im(\varphi).
$$
如果上同调群均为零，就称这个 $2$-周期复形是{\it 正合的}。
\item $R$ 上的一个 {\it $(2, 1)$-周期复形}由三元组 $(M, \varphi, \psi)$ 给出，
其中 $M$ 是 $R$-模，$\varphi : M \to M$、$\psi : M \to M$ 是 $R$-模同态，并且
$$
\xymatrix{
\ldots \ar[r] &
M \ar[r]^\varphi &
M \ar[r]^\psi &
M \ar[r]^\varphi &
M \ar[r] & \ldots
}
$$
构成一个复形。由于这是 $2$-周期复形的一种特殊情形，我们有它的{\it 上同调模}
$H^0(M, \varphi, \psi)$、$H^1(M, \varphi, \psi)$，以及相应的正合性概念。
\end{enumerate}
\end{definition}

\noindent
下文凡是对 $2$-周期复形证明的结果，都将不再另行说明而直接用于 $(2, 1)$-周期复形。
显然，所有 $2$-周期复形组成一个范畴，其中态射
$(f, g) : (M, N, \varphi, \psi) \to (M', N', \varphi', \psi')$
是满足 $\varphi' \circ f = g \circ \varphi$ 和 $\psi' \circ g = f \circ \psi$ 的一对态射
$f : M \to M'$ 与 $g : N \to N'$。这样得到一个阿贝尔范畴；其中的核与余核如《同调代数》引理 \ref{homology-lemma-cat-chain-abelian} 所述。

\begin{definition}
\label{chow-definition-periodic-length}
设 $(M, N, \varphi, \psi)$ 是环 $R$ 上的一个 $2$-周期复形，并且它的上同调模均具有有限长度。
在这种情形下，定义 $(M, N, \varphi, \psi)$ 的{\it 重数}为整数
$$
e_R(M, N, \varphi, \psi) =
\text{length}_R(H^0(M, N, \varphi, \psi))
-
\text{length}_R(H^1(M, N, \varphi, \psi))
$$
对于 $(2, 1)$-周期复形 $(M, \varphi, \psi)$，将这个整数记为
$e_R(M, \varphi, \psi)$，有时也称它为{\it（加法）Herbrand 商}。
\end{definition}

\noindent
如果 $(M, \varphi, \psi)$ 的上同调群是有限阿贝尔群，则通常把
$$
q(M, \varphi, \psi) =
\frac{\# H^0(M, \varphi, \psi)}{\# H^1(M, \varphi, \psi)}
$$
称为{\it（乘法）Herbrand 商}。换言之，乘法 Herbrand 商就是 $H^0$ 的元素个数除以
$H^1$ 的元素个数。如果 $R$ 是局部环，而且 $R$ 的剩余域含有 $q$ 个元素，那么
$$
q(M, \varphi, \psi) = q^{e_R(M, \varphi, \psi)}
$$

\medskip\noindent
环 $R$ 上 $(2, 1)$-周期复形的一个例子是任意形如 $(M, 0, \psi)$ 的三元组，
其中 $M$ 是 $R$-模，$\psi$ 是 $R$-线性映射。如果 $\psi$ 的核和余核都具有有限长度，则有
\begin{equation}
\label{chow-equation-multiplicity-coker-ker}
e_R(M, 0, \psi) = \text{length}_R(\Coker(\psi)) - \text{length}_R(\Ker(\psi))
\end{equation}
下面陈述并证明关于这些记号的几个必要引理。

\begin{lemma}
\label{chow-lemma-additivity-periodic-length}
设 $R$ 为环，并设有一个 $2$-周期复形的短正合列
$$
0 \to (M_1, N_1, \varphi_1, \psi_1)
\to (M_2, N_2, \varphi_2, \psi_2)
\to (M_3, N_3, \varphi_3, \psi_3)
\to 0
$$
如果其中两个复形的上同调模具有有限长度，那么第三个复形也具有这个性质，并且
$$
e_R(M_2, N_2, \varphi_2, \psi_2) =
e_R(M_1, N_1, \varphi_1, \psi_1) +
e_R(M_3, N_3, \varphi_3, \psi_3).
$$
\end{lemma}

\begin{proof}
简记 $A = (M_1, N_1, \varphi_1, \psi_1)$、
$B = (M_2, N_2, \varphi_2, \psi_2)$ 和 $C = (M_3, N_3, \varphi_3, \psi_3)$。
我们有上同调长正合列
$$
\ldots
\to H^1(C)
\to H^0(A)
\to H^0(B)
\to H^0(C)
\to H^1(A)
\to H^1(B)
\to H^1(C)
\to \ldots
$$
由此得到有限正合列
$$
0 \to I
\to H^0(A)
\to H^0(B)
\to H^0(C)
\to H^1(A)
\to H^1(B)
\to K \to 0
$$
以及滤过 $0 \to K \to H^1(C) \to I \to 0$。由长度函数的可加性
（《交换代数》引理 \ref{algebra-lemma-length-additive}）即得结论。
\end{proof}

\begin{lemma}
\label{chow-lemma-finite-periodic-length}
设 $R$ 为环。如果 $(M, N, \varphi, \psi)$ 是一个 $2$-周期复形，且 $M$、$N$ 具有有限长度，那么
$e_R(M, N, \varphi, \psi) = \text{length}_R(M) - \text{length}_R(N)$。
特别地，如果 $(M, \varphi, \psi)$ 是一个 $(2, 1)$-周期复形，且 $M$ 具有有限长度，那么
$e_R(M, \varphi, \psi) = 0$。
\end{lemma}

\begin{proof}
在由 $M$、$N$ 均具有有限长度的 $2$-周期复形所成的范畴中，量
``$\text{length}_R(M) - \text{length}_R(N)$'' 对短正合列具有可加性
（精确陈述留给读者）。考虑短正合列
$$
0 \to (M, \Im(\varphi), \varphi, 0) \to
(M, N, \varphi, \psi) \to (0, N/\Im(\varphi), 0, 0) \to 0
$$
结合开头的观察与引理 \ref{chow-lemma-additivity-periodic-length} 的可加性，
可将问题约化为以下两种情形：(a) $M = 0$；(b) $\varphi$ 是满射。这两种情形留给读者。
\end{proof}

\begin{lemma}
\label{chow-lemma-compare-periodic-lengths}
设 $R$ 为环。设
$f : (M, \varphi, \psi) \to (M', \varphi', \psi')$
是两个 $(2, 1)$-周期复形之间的映射，并且这两个复形的上同调模均具有有限长度。
如果 $\Ker(f)$ 和 $\Coker(f)$ 具有有限长度，那么
$e_R(M, \varphi, \psi) = e_R(M', \varphi', \psi')$。
\end{lemma}

\begin{proof}
应用引理 \ref{chow-lemma-additivity-periodic-length} 的可加性，并注意由
引理 \ref{chow-lemma-finite-periodic-length}，
$(\Ker(f), \varphi, \psi)$ 和 $(\Coker(f), \varphi', \psi')$ 的重数均为零。
\end{proof}




\section{若干重数的计算}
\label{chow-section-calculation}

\noindent
为了在后文证明某些循环相等，我们需要计算若干重数。除上一节关于重数的初等引理外，
主要工具是《交换代数》引理 \ref{algebra-lemma-order-vanishing-determinant}。

\begin{lemma}
\label{chow-lemma-length-multiplication}
设 $R$ 为诺特局部环，$M$ 为有限 $R$-模，并设 $x \in R$。假设
\begin{enumerate}
\item $\dim(\text{Supp}(M)) \leq 1$，并且
\item $\dim(\text{Supp}(M/xM)) \leq 0$。
\end{enumerate}
写成
$\text{Supp}(M) = \{\mathfrak m, \mathfrak q_1, \ldots, \mathfrak q_t\}$。
那么
$$
e_R(M, 0, x) =
\sum\nolimits_{i = 1, \ldots, t}
\text{ord}_{R/\mathfrak q_i}(x)
\text{length}_{R_{\mathfrak q_i}}(M_{\mathfrak q_i}).
$$
\end{lemma}

\begin{proof}
先作一些准备性说明。如果 $M$ 具有有限长度，即 $t = 0$，则引理成立；
因为此时等式两边都为零，见引理 \ref{chow-lemma-finite-periodic-length}。
此外，若满足 (1)、(2) 的模之间有短正合列 $0 \to M \to M' \to M'' \to 0$，
则只要引理对其中两个模成立，由长度的可加性（《交换代数》引理 \ref{algebra-lemma-length-additive}）和重数的可加性（引理 \ref{chow-lemma-additivity-periodic-length}），它对第三个模也成立。

\medskip\noindent
以 $M_i$ 表示 $M$ 在 $M_{\mathfrak q_i}$ 中的像，于是
$\text{Supp}(M_i) = \{\mathfrak m, \mathfrak q_i\}$。
映射 $M \to \bigoplus M_i$ 的核与余核的支撑均为 $\{\mathfrak m\}$，因而具有有限长度。
由上述准备性说明，只需对每个 $M_i$ 证明引理。因此可以假设
$\text{Supp}(M) = \{\mathfrak m, \mathfrak q\}$。
此时，由《交换代数》引理 \ref{algebra-lemma-Noetherian-power-ideal-kills-module}，有有限滤过
$M \supset \mathfrak qM \supset \mathfrak q^2M \supset \ldots \supset
\mathfrak q^nM = 0$。再次利用可加性可知，只需证明 $M$ 被 $\mathfrak q$ 零化的情形。
这时可把 $M$ 看成 $R/\mathfrak q$-模；换言之，可以假设 $R$ 是维数为 $1$、分式域为 $K$ 的诺特局部整环。
再商去挠子模，即商去 $M \to M \otimes_R K = V$ 的核（挠子模具有有限长度，
故可由准备性说明处理），便可假设 $M \subset V$ 是一个格
（《交换代数》定义 \ref{algebra-definition-lattice}）。于是 $x : M \to M$ 是单射，并且
$\text{length}_R(M/xM) = d(M, xM)$
（《交换代数》定义 \ref{algebra-definition-distance}）。由于
$\text{length}_K(V) = \dim_K(V)$，可知
$\det(x : V \to V) = x^{\dim_K(V)}$，且
$\text{ord}_R(\det(x : V \to V)) = \dim_K(V) \text{ord}_R(x)$。
因而在此情形下，所需等式由《交换代数》引理 \ref{algebra-lemma-order-vanishing-determinant} 得出。
\end{proof}

\begin{lemma}
\label{chow-lemma-additivity-divisors-restricted}
设 $R$ 为诺特局部环，$x \in R$。如果 $M$ 是有限 Cohen--Macaulay $R$-模，且
$\dim(\text{Supp}(M)) = 1$、$\dim(\text{Supp}(M/xM)) = 0$，那么
$$
\text{length}_R(M/xM)
=
\sum\nolimits_i \text{length}_R(R/(x, \mathfrak q_i))
\text{length}_{R_{\mathfrak q_i}}(M_{\mathfrak q_i}).
$$
这里 $\mathfrak q_1, \ldots, \mathfrak q_t$ 是 $M$ 的支撑的极小素理想。如果 $I \subset R$ 是一个理想，
$x$ 在 $R/I$ 上是非零因子，并且 $\dim(R/I) = 1$，那么
$$
\text{length}_R(R/(x, I))
=
\sum\nolimits_i \text{length}_R(R/(x, \mathfrak q_i))
\text{length}_{R_{\mathfrak q_i}}((R/I)_{\mathfrak q_i})
$$
这里 $\mathfrak q_1, \ldots, \mathfrak q_n$ 是包含 $I$ 的极小素理想。
\end{lemma}

\begin{proof}
这些都是引理 \ref{chow-lemma-length-multiplication} 的特殊情形。
\end{proof}

\noindent
下面是在另一种可以确定重数取值的情形。

\begin{lemma}
\label{chow-lemma-powers-period-length-zero}
设 $R$ 为环，$M$ 为 $R$-模。设 $\varphi : M \to M$ 是自同态，$n > 0$，满足
$\varphi^n = 0$，并且 $\Ker(\varphi)/\Im(\varphi^{n - 1})$ 作为 $R$-模具有有限长度。
那么对 $i = 0, \ldots, n$ 均有
$$
e_R(M, \varphi^i, \varphi^{n - i}) = 0
$$
\end{lemma}

\begin{proof}
按约定 $\varphi^0 = \text{id}_M$，所以 $i = 0, n$ 的情形是显然的。
把 $M$ 看成 $R[t]$-模，其中乘以 $t$ 的作用由 $\varphi$ 给出。记
$K_i = \Ker(t^i : M \to M)$，并令
$$
a_i = \text{length}_R(K_i/t^{n - i}M),\quad
b_i = \text{length}_R(K_i/tK_{i + 1}),\quad
c_i = \text{length}_R(K_1/t^iK_{i + 1})
$$
边界值为 $a_0 = a_n = b_0 = c_0 = 0$。当 $i < n$ 时，$c_i$ 是整数，
因为 $K_1/t^iK_{i + 1}$ 是按假设具有有限长度的 $K_1/t^{n - 1}M$ 的商。
下面将反复使用 $K_i \cap t^jM = t^jK_{i + j}$。对 $0 < i < n - 1$，有正合列
$$
0 \to
K_1/t^{n - i - 1}K_{n - i} \to
K_{i + 1}/t^{n - i - 1}M \xrightarrow{t} K_i/t^{n - i}M
\to K_i/tK_{i + 1} \to 0
$$
对 $i$ 作归纳可知，当 $i < n$ 时 $a_i$、$b_i$ 都是整数，并且
$$
c_{n - i - 1} - a_{i + 1} + a_i - b_i = 0
$$
对 $0 < i < n - 1$，还有短正合列
$$
0 \to
K_i/tK_{i + 1} \to
K_{i + 1}/tK_{i + 2} \xrightarrow{t^i}
K_1/t^{i + 1}K_{i + 2} \to
K_1/t^iK_{i + 1} \to 0
$$
因而
$$
b_i - b_{i + 1} + c_{i + 1} - c_i = 0
$$
由于 $b_0 = c_0$，可得对 $i < n$ 均有 $b_i = c_i$。于是
$$
a_2 = a_1 + b_{n - 2} - b_1,\quad
a_3 = a_2 + b_{n - 3} - b_2,\quad \ldots
$$
直接可见这蕴含所需的 $a_i = a_{n - i}$。
\end{proof}

\begin{lemma}
\label{chow-lemma-multiply-period-length}
设 $(R, \mathfrak m)$ 为诺特局部环。设 $(M, \varphi, \psi)$ 是 $R$ 上的一个
$(2, 1)$-周期复形，其中 $M$ 有限，上同调群作为 $R$-模具有有限长度。
设 $x \in R$ 满足 $\dim(\text{Supp}(M/xM)) \leq 0$。那么
$$
e_R(M, x\varphi, \psi) = e_R(M, \varphi, \psi) - e_R(\Im(\varphi), 0, x)
$$
并且
$$
e_R(M, \varphi, x\psi) = e_R(M, \varphi, \psi) + e_R(\Im(\psi), 0, x)
$$
\end{lemma}

\begin{proof}
只证明第一个公式；第二个公式的证明完全相同。
设 $M' = M[x^\infty]$ 为 $M$ 的 $x$-幂挠子模，并考虑短正合列
$0 \to M' \to M \to M'' \to 0$。那么 $M''$ 无 $x$-幂挠
（《代数进阶》引理 \ref{more-algebra-lemma-divide-by-torsion}）。
由于 $\varphi$、$\psi$ 都把 $M'$ 映入 $M'$，可得 $(2, 1)$-周期复形的短正合列
$$
0 \to (M', \varphi', \psi') \to (M, \varphi, \psi) \to
(M'', \varphi'', \psi'') \to 0
$$
此外还有短正合列
$0 \to M' \cap \Im(\varphi) \to \Im(\varphi) \to \Im(\varphi'') \to 0$。
由引理 \ref{chow-lemma-compare-periodic-lengths}，有
$e_R(M', \varphi, \psi) = e_R(M', x\varphi, \psi) =
e_R(M' \cap \Im(\varphi), 0, x) = 0$。
再由可加性（引理 \ref{chow-lemma-additivity-periodic-length}）可知，
只需对 $(M'', \varphi'', \psi'')$ 证明引理。这就约化到下一段所讨论的情形。

\medskip\noindent
假设 $x : M \to M$ 是单射。此时 $\Ker(x\varphi) = \Ker(\varphi)$。
另一方面，有短正合列
$$
0 \to \Im(\varphi)/x\Im(\varphi) \to
\Ker(\psi)/\Im(x\varphi) \to \Ker(\psi)/\Im(\varphi) \to 0
$$
这与 (\ref{chow-equation-multiplicity-coker-ker}) 一起证明了该公式。
\end{proof}







\section{驯顺符号的准备工作}
\label{chow-section-preparation-tame-symbol}

\noindent
本节给出一些引理，以帮助我们在下一节定义驯顺符号。

\begin{lemma}
\label{chow-lemma-glue-at-max}
设 $A$ 为诺特环，$\mathfrak m_1, \ldots, \mathfrak m_r$ 为 $A$ 中两两不同的极大理想。
对 $i = 1, \ldots, r$，设 $\varphi_i : A_{\mathfrak m_i} \to B_i$ 是环同态，
其核与余核均被 $\mathfrak m_i$ 的某个幂零化。那么存在环同态 $\varphi : A \to B$，使得
\begin{enumerate}
\item $\varphi$ 在 $\mathfrak m_i$ 处的局部化同构于 $\varphi_i$；并且
\item $\Ker(\varphi)$ 与 $\Coker(\varphi)$ 均被
$\mathfrak m_1 \cap \ldots \cap \mathfrak m_r$ 的某个幂零化。
\end{enumerate}
此外，如果每个 $\varphi_i$ 都是有限的、单射的或满射的，那么 $\varphi$ 也分别具有相应性质。
\end{lemma}

\begin{proof}
令 $I = \mathfrak m_1 \cap \ldots \cap \mathfrak m_r$，并令
$A_i = A_{\mathfrak m_i}$、$A' = \prod A_i$。那么
$IA' = \prod \mathfrak m_i A_i$，且 $A \to A'$ 是满足 $A/I \cong A'/IA'$ 的平坦环同态。
因此，由《代数进阶》引理 \ref{more-algebra-lemma-application-formal-glueing}，存在 $A$-模同态
$\varphi : A \to B$，使 $\varphi_i$ 同构于 $\varphi$ 在 $\mathfrak m_i$ 处的局部化。
再利用《代数进阶》注 \ref{more-algebra-remark-formal-glueing-algebras}
中的讨论，可在 $B$ 上赋予一个 $A$-代数结构，使其与 $B_i$ 上给定的 $A$-代数结构相符。
由 $\Ker(\varphi)_{\mathfrak m_i} \cong \Ker(\varphi_i)$ 和
$\Coker(\varphi)_{\mathfrak m_i} \cong \Coker(\varphi_i)$，
很容易得到引理的最后一个断言。
\end{proof}

\noindent
下面的引理与《交换代数》引理 \ref{algebra-lemma-nonregular-dimension-one} 非常相似。

\begin{lemma}
\label{chow-lemma-Noetherian-domain-dim-1-two-elements}
设 $(R, \mathfrak m)$ 是维数为 $1$ 的诺特局部环，并设 $a, b \in R$ 为非零因子。
那么存在有限环扩张 $R \subset R'$，其中 $R'/R$ 被 $\mathfrak m$ 的某个幂零化；
并且存在非零因子 $t, a', b' \in R'$，使得
$a = ta'$、$b = tb'$ 且 $R' = a'R' + b'R'$。
\end{lemma}

\begin{proof}
如果 $a$ 或 $b$ 是单位，则取 $R = R'$ 即得引理。因此可以假设 $a, b \in \mathfrak m$。
令 $I = (a, b)$。基本思路是沿 $I$ 爆破 $R$；我们不采用代数论证，而采用几何论证。
令 $X = \text{Proj}(\bigoplus_{d \geq 0} I^d)$。由《除子》引理 \ref{divisors-lemma-blowing-up-gives-effective-Cartier-divisor}，态射
$X \to \Spec(R)$ 在去心谱 $U = \Spec(R) \setminus \{\mathfrak m\}$ 上是同构。
因此可以而且将把 $U$ 看作 $X$ 的开子概形。由《除子》引理 \ref{divisors-lemma-blowing-up-projective}，态射 $X \to \Spec(R)$ 是射影的。
此外，例如由《除子》引理 \ref{divisors-lemma-blow-up-and-irreducible-components}，$X$ 的每个一般点都位于 $U$ 中。
于是由《代数簇》引理 \ref{varieties-lemma-finite-in-codim-1}，
$X \to \Spec(R)$ 是有限态射。因此 $X = \Spec(R')$ 是仿射的，且 $R \to R'$ 有限。
由于 $U = D(a)$，有 $R_a \cong R'_a$。故 $a$ 的某个幂零化有限 $R$-模 $R'/R$。
又因 $\mathfrak m = \sqrt{(a)}$，可知 $R'/R$ 被 $\mathfrak m$ 的某个幂零化。
由《除子》引理 \ref{divisors-lemma-blowing-up-gives-effective-Cartier-divisor}，
$IR'$ 是局部主理想。由于 $R'$ 是半局部环，$IR'$ 是主理想；见《交换代数》引理 \ref{algebra-lemma-locally-free-semi-local-free}。设 $IR' = (t)$，则
$a = a't$、$b = b't$，其余结论均显然。
\end{proof}

\begin{lemma}
\label{chow-lemma-not-infinitely-divisible}
设 $(R, \mathfrak m)$ 是维数为 $1$ 的诺特局部环。设 $a, b \in R$ 是非零因子，且
$a \in \mathfrak m$。存在整数 $n = n(R, a, b)$，使得对于任意有限环扩张
$R \subset R'$，如果对某个 $c \in R'$ 有 $b = a^m c$，那么 $m \leq n$。
\end{lemma}

\begin{proof}
取极小素理想 $\mathfrak q \subset R$。注意 $\dim(R/\mathfrak q) = 1$，特别地，
$R/\mathfrak q$ 不是域。可以选取一个与 $R/\mathfrak q$ 有相同分式域并支配它的离散赋值环 $A$；
见《交换代数》引理 \ref{algebra-lemma-dominate-by-dimension-1}。
由于 $R$ 中的非零因子不属于 $\mathfrak q$，$a$、$b$ 在 $A$ 中的像均非零。
设 $v$ 为 $A$ 上的离散赋值。因 $a \in \mathfrak m$，有 $v(a) > 0$。
我们断言 $n = v(b)/v(a)$ 可用。

\medskip\noindent
给定 $R \subset R'$，令 $A' = A \otimes_R R'$。由于
$\Spec(R') \to \Spec(R)$ 是满射
（《交换代数》引理 \ref{algebra-lemma-integral-overring-surjective}），
$\Spec(A') \to \Spec(A)$ 也为满射
（《交换代数》引理 \ref{algebra-lemma-surjective-spec-radical-ideal}）。
取位于 $(0) \subset A$ 上方的素理想 $\mathfrak q' \subset A'$。
那么 $A \subset A'' = A'/\mathfrak q'$ 是有限环扩张（并且仍诱导谱上的满射）。
取位于 $A$ 的极大理想上方的极大理想 $\mathfrak m'' \subset A''$，再取一个离散赋值环
$A'''$，它支配 $A''_{\mathfrak m''}$（见上引引理）。于是 $A \to A'''$ 是离散赋值环的扩张，并且等式 $b = a^m c$ 在 $A'''$ 中成立。因此 $v'''(b) \geq mv'''(a)$。由于 $v''' = ev$，其中 $e$ 是 $A'''/A$ 的分歧指数，便得到所需的 $m \leq n$。
\end{proof}

\begin{lemma}
\label{chow-lemma-prepare-tame-symbol}
设 $(A, \mathfrak m)$ 是维数为 $1$ 的诺特局部环。设 $r \geq 2$，并设
$a_1, \ldots, a_r \in A$ 是不全为单位的非零因子。那么存在
\begin{enumerate}
\item 有限环扩张 $A \subset B$，其中 $B/A$ 被 $\mathfrak m$ 的某个幂零化；
\item 对每个极大理想 $\mathfrak m_j \subset B$，存在非零因子
$\pi_j \in B_j = B_{\mathfrak m_j}$；并且
\item 在 $B_j$ 中存在分解 $a_i = u_{i, j} \pi_j^{e_{i, j}}$，
其中 $u_{i, j} \in B_j$ 为单位，$e_{i, j} \geq 0$。
\end{enumerate}
\end{lemma}

\begin{proof}
由于至少有一个 $a_i$ 不是单位，$\mathfrak m$ 不是 $A$ 的伴随素理想。
此外，对陈述中任意的 $A \subset B$，$\mathfrak m$ 不是 $B$ 的伴随素理想，
$\mathfrak m_j$ 也不是 $B_j$ 的伴随素理想。记住这一点有助于核查下面的论证。

\medskip\noindent
首先断言，只需对 $r = 2$ 证明引理。我们对 $r$ 作归纳；建议读者跳过这段证明。
设已经给定 $A \subset B$、$B_j = B_{\mathfrak m_j}$ 中的 $\pi_j$，以及在 $B_j$ 中的分解
$a_i = u_{i, j} \pi_j^{e_{i, j}}$（$i = 1, \ldots, r - 1$），
其中 $u_{i, j} \in B_j$ 为单位且 $e_{i, j} \geq 0$。
把 $r = 2$ 的情形用于 $B_j$ 中的 $\pi_j$ 与 $a_r$，便可找到扩张 $B_j \subset C_j$，
并且对 $C_j$ 的每个极大理想 $\mathfrak m_{j, k}$，找到非零因子
$\pi_{j, k} \in C_{j, k} = (C_j)_{\mathfrak m_{j, k}}$ 和引理所述的分解
$$
\pi_j = v_{j, k} \pi_{j, k}^{f_{j, k}}
\quad\text{且}\quad
a_r = w_{j, k} \pi_{j, k}^{g_{j, k}}
$$
由引理 \ref{chow-lemma-glue-at-max}，存在唯一的有限扩张 $B \subset C$，其中
$C/B$ 被 $\mathfrak m$ 的某个幂零化，并且对所有 $j$ 都有
$C_j \cong C_{\mathfrak m_j}$。$C$ 的极大理想与各局部化中的极大理想
$\mathfrak m_{j, k}$ 构成 $1$-对-$1$ 对应；在这些局部化中，对 $i \leq r - 1$ 有
$$
a_i = u_{i, j} \pi_j^{e_{i, j}} =
u_{i, j} v_{j, k}^{e_{i, j}} \pi_{j, k}^{e_{i, j}f_{j, k}}
$$
而 $a_r$ 也有正确的分解，故归纳步骤得证。

\medskip\noindent
下面证明 $r = 2$ 的情形。对
$$
\ell = \min(\text{length}_A(A/a_1A),\ \text{length}_A(A/a_2A)).
$$
作归纳。如果 $\ell = 0$，则 $a_1$ 或 $a_2$ 是单位，取 $A = B$ 即得引理。
因此可以而且将假设 $\ell > 0$。

\medskip\noindent
设有有限环扩张 $A \subset A'$，其中 $A'/A$ 被 $\mathfrak m$ 的某个幂零化，
且 $\mathfrak m$ 不是 $A'$ 的伴随素理想。设
$\mathfrak m_1, \ldots, \mathfrak m_r \subset A'$ 为极大理想，并令
$A'_i = A'_{\mathfrak m_i}$。如果能在每个 $A'_i$ 中解决 $a_1, a_2$ 的问题，
就可应用引理 \ref{chow-lemma-glue-at-max} 得到 $A$ 中 $a_1, a_2$ 的解。
选取 $x \in \{a_1, a_2\}$，使 $\ell = \text{length}_A(A/xA)$。
由引理 \ref{chow-lemma-compare-periodic-lengths} 
与 (\ref{chow-equation-multiplicity-coker-ker})，有
$\text{length}_A(A/xA) = \text{length}_A(A'/xA')$。另一方面，由《交换代数》引理 \ref{algebra-lemma-pushdown-module}，有
$$
\text{length}_A(A'/xA') =
\sum [\kappa(\mathfrak m_i) : \kappa(\mathfrak m)]
\text{length}_{A'_i}(A'_i/xA'_i)
$$
由于 $x \in \mathfrak m$，右边每一项都为正。因此，如果 $r > 1$，或者 $r = 1$ 且
$[\kappa(\mathfrak m_1) : \kappa(\mathfrak m)] > 1$，归纳假设可用于每个 $A'_i$ 中的
$a_1, a_2$。由此可以假设，上述每个 $A'$ 都是与 $A$ 有相同剩余域的局部环。

\medskip\noindent
应用上一段的讨论，可以把 $A$ 替换为引理 \ref{chow-lemma-Noetherian-domain-dim-1-two-elements} 对 $a_1, a_2 \in A$ 构造的环。
由于 $A$ 是局部环，必要时交换 $a_1$、$a_2$ 后，可设 $a_2 \in (a_1)$。
写成 $a_2 = a_1^m c$，其中 $m > 0$ 取最大值。事实上，由引理 \ref{chow-lemma-not-infinitely-divisible}，即使把 $A$ 换成上一段所述的任意有限扩张
$A \subset A'$，也可假设 $m$ 仍为最大值。如果 $c$ 是单位，证明结束；否则，把 $A$ 替换为
引理 \ref{chow-lemma-Noetherian-domain-dim-1-two-elements} 对 $a_1, c \in A$ 构造的环。
此时或者 (1) $c = a_1 c'$，或者 (2) $a_1 = c a'_1$。第一种情形不可能发生，
因为它给出 $a_2 = a_1^{m + 1} c'$，与 $m$ 的最大性矛盾。第二种情形给出
$a_1 = c a'_1$ 和 $a_2 = c^{m + 1} (a'_1)^m$。
因此只需对 $A$ 以及 $c, a'_1$ 证明引理。如果 $a'_1$ 是单位，证明结束；否则，
由于 $cA$ 严格大于 $a_1A$，有 $\text{length}_A(A/cA) < \ell$，
于是由归纳假设得到结论。
\end{proof}






\section{驯顺符号}
\label{chow-section-tame-symbol}

\noindent
考虑维数为 $1$ 的诺特局部环 $(A, \mathfrak m)$。以 $Q(A)$ 表示 $A$ 的全分式环；见
《交换代数》例 \ref{algebra-example-localize-at-prime}。
{\it 驯顺符号}将是一个映射
$$
\partial_A(-, -) : Q(A)^* \times Q(A)^* \longrightarrow \kappa(\mathfrak m)^*
$$
并满足以下性质：
\begin{enumerate}
\item $\partial_A(f, gh) = \partial_A(f, g) \partial_A(f, h)$
\label{chow-item-bilinear}
对 $f, g, h \in Q(A)^*$ 成立；
\item $\partial_A(f, g) \partial_A(g, f) = 1$
\label{chow-item-skew}
对 $f, g \in Q(A)^*$ 成立；
\item $\partial_A(f, 1 - f) = 1$
\label{chow-item-1-x}
对满足 $1 - f \in Q(A)^*$ 的 $f \in Q(A)^*$ 成立；
\item $\partial_A(aa', b) = \partial_A(a, b)\partial_A(a', b)$
\label{chow-item-bilinear-better}
和 $\partial_A(a, bb') = \partial_A(a, b)\partial_A(a, b')$
对非零因子 $a, a', b, b' \in A$ 成立；
\item $\partial_A(b, b) = (-1)^m$
\label{chow-item-skew-better}
对非零因子 $b \in A$ 成立，其中 $m = \text{length}_A(A/bA)$；
\item $\partial_A(u, b) = u^m \bmod \mathfrak m$
\label{chow-item-normalization}
对单位 $u \in A$ 和非零因子 $b \in A$ 成立，其中
$m = \text{length}_A(A/bA)$；并且
\item
\label{chow-item-1-x-better}
$\partial_A(a, b - a)\partial_A(b, b) = \partial_A(b, b - a)\partial_A(a, b)$
对使 $a, b, b - a$ 都是非零因子的 $a, b \in A$ 成立。
\end{enumerate}
由于处理 $A$ 的元素比较方便，我们常把 $\partial_A$ 看作定义在 $A$ 的非零因子对上的映射，
满足 (\ref{chow-item-bilinear-better})、(\ref{chow-item-skew-better})、
(\ref{chow-item-normalization}) 和 (\ref{chow-item-1-x-better})。
读者不难验证，令
$$
\partial_A(\frac{a}{b}, \frac{c}{d}) =
\partial_A(a, c) \partial_A(a, d)^{-1} \partial_A(b, c)^{-1} \partial_A(b, d)
$$
便得到良定义映射 $Q(A)^* \times Q(A)^* \to \kappa(\mathfrak m)^*$，
它满足 (\ref{chow-item-bilinear})、(\ref{chow-item-skew})、(\ref{chow-item-1-x}) 以及上述其他性质。

\medskip\noindent
我们并不声称满足这些性质的映射是唯一的，而是给出构造这种映射的方法。
具体而言，给定非零因子 $a_1, a_2 \in A$，选择如引理 \ref{chow-lemma-prepare-tame-symbol} 中的环扩张 $A \subset B$ 与局部分解。定义
\begin{equation}
\label{chow-equation-tame-symbol}
\partial_A(a_1, a_2) = \prod\nolimits_j
\text{Norm}_{\kappa(\mathfrak m_j)/\kappa(\mathfrak m)}
((-1)^{e_{1, j}e_{2, j}}u_{1, j}^{e_{2, j}}u_{2, j}^{-e_{1, j}}
\bmod \mathfrak m_j)^{m_j}
\end{equation}
其中 $m_j = \text{length}_{B_j}(B_j/\pi_j B_j)$，乘积遍历 $B$ 的极大理想
$\mathfrak m_1, \ldots, \mathfrak m_r$。

\begin{lemma}
\label{chow-lemma-well-defined-tame-symbol}
公式 (\ref{chow-equation-tame-symbol}) 确定 $\kappa(\mathfrak m)^*$ 中一个良定义的元素。
换言之，右边不依赖于局部分解的选择，也不依赖于 $B$ 的选择。
\end{lemma}

\begin{proof}
先证与分解选择无关。设 $B$ 是一个诺特 $1$-维局部环，$a_1, a_2 \in B$，
并设 $\pi, \theta$ 为非零因子，使得
$$
a_1 = u_1 \pi^{e_1} = v_1 \theta^{f_1},\quad
a_2 = u_2 \pi^{e_2} = v_2 \theta^{f_2}
$$
其中 $e_i, f_i \geq 0$ 是整数，$u_i, v_i$ 是 $B$ 中的单位。由此可得
$$
a_1^{e_2} = u_1^{e_2}u_2^{-e_1}a_2^{e_1},\quad
a_1^{f_2} = v_1^{f_2}v_2^{-f_1}a_2^{f_1}
$$
另一方面，令 $m = \text{length}_B(B/\pi B)$、$k = \text{length}_B(B/\theta B)$，
便有 $e_2 m = \text{length}_B(B/a_2 B) = f_2 k$。
用上式展开 $a_1^{e_2m} = a_1^{f_2 k}$，得到
$$
(u_1^{e_2}u_2^{-e_1})^m =  (v_1^{f_2}v_2^{-f_1})^k
$$
这已证明除符号外的所需等式。为了验证符号一致，只须证明 $me_1e_2$ 为偶数当且仅当
$kf_1f_2$ 为偶数。由于 $me_2 = kf_2$ 且 $me_1 = kf_1$（论证同上），结论成立。

\medskip\noindent
再证与 $B$ 的选择无关。设有引理 \ref{chow-lemma-prepare-tame-symbol} 中的两个扩张
$A \subset B$ 和 $A \subset B'$。那么
$$
C = (B \otimes_A B')/(\mathfrak m\text{-幂挠子})
$$
给出第三个这样的扩张。因此可以假设 $A \subset B \subset C$，并已在 $B$ 的各局部环中给定分解；
需要证明，在 $C$ 的各局部环中使用同一组分解会得到 $\kappa(\mathfrak m)$ 中的同一元素。
由范数的传递性（《域》引理 \ref{fields-lemma-trace-and-norm-tower}），
问题归结为：若 $B$ 是维数为 $1$ 的诺特局部环，且 $\pi \in B$ 是非零因子，则
$$
\lambda^m = \prod \text{Norm}_{\kappa_k/\kappa}(\lambda)^{m_k}
$$
这里使用如下记号：
(1) $\kappa$ 是 $B$ 的剩余域；
(2) $\lambda$ 是 $\kappa$ 的一个元素；
(3) $\mathfrak m_k \subset C$ 是 $C$ 的极大理想；
(4) $\kappa_k = \kappa(\mathfrak m_k)$ 是 $C_k = C_{\mathfrak m_k}$ 的剩余域；
(5) $m = \text{length}_B(B/\pi B)$；以及
(6) $m_k = \text{length}_{C_k}(C_k/\pi C_k)$。
上述等式成立，因为对 $\lambda \in \kappa$ 有
$\text{Norm}_{\kappa_k/\kappa}(\lambda) = \lambda^{[\kappa_k : \kappa]}$，
且 $m = \sum m_k[\kappa_k:\kappa]$。首先，由引理 \ref{chow-lemma-compare-periodic-lengths} 与 (\ref{chow-equation-multiplicity-coker-ker})，有
$m = \text{length}_B(B/xB) = \text{length}_B(C/\pi C)$。
最后，由《交换代数》引理 \ref{algebra-lemma-pushdown-module}，有
$\text{length}_B(C/\pi C) = \sum m_k[\kappa_k:\kappa]$。
\end{proof}

\begin{lemma}
\label{chow-lemma-tame-symbol}
驯顺符号 (\ref{chow-equation-tame-symbol}) 满足
(\ref{chow-item-bilinear-better})、(\ref{chow-item-skew-better})、
(\ref{chow-item-normalization}) 和 (\ref{chow-item-1-x-better})，因而给出映射
$\partial_A : Q(A)^* \times Q(A)^* \to \kappa(\mathfrak m)^*$，
它满足 (\ref{chow-item-bilinear})、(\ref{chow-item-skew}) 和 (\ref{chow-item-1-x})。
\end{lemma}

\begin{proof}
先证明 (\ref{chow-item-bilinear-better})。设非零因子 $a_1, a_2, a_3 \in A$。
为 $a_1, a_2, a_3$ 选择引理 \ref{chow-lemma-prepare-tame-symbol} 中的
$A \subset B$。等式
$$
\partial_A(a_1a_2, a_3) = \partial_A(a_1, a_3) \partial_A(a_2, a_3)
$$
由 $B_j$ 中的等式
$$
(-1)^{(e_{1, j} + e_{2, j})e_{3, j}}
(u_{1, j}u_{2, j})^{e_{3, j}}u_{3, j}^{-e_{1, j} - e_{2, j}} =
(-1)^{e_{1, j}e_{3, j}}
u_{1, j}^{e_{3, j}}u_{3, j}^{-e_{1, j}}
(-1)^{e_{2, j}e_{3, j}}
u_{2, j}^{e_{3, j}}u_{3, j}^{-e_{2, j}}
$$
推出。性质 (\ref{chow-item-skew-better}) 和 (\ref{chow-item-normalization}) 同样是立即可得的。

\medskip\noindent
下面证明 (\ref{chow-item-1-x-better})。设 $a_1, a_2, a_1 - a_2 \in A$ 为非零因子，
并令 $a_3 = a_1 - a_2$。为 $a_1, a_2, a_3$ 选择引理 \ref{chow-lemma-prepare-tame-symbol} 中的 $A \subset B$。只需证明
$$
(-1)^{e_{1, j}e_{2, j} + e_{1, j}e_{3, j} + e_{2, j}e_{3, j} + e_{2, j}}
u_{1, j}^{e_{2, j} - e_{3, j}}
u_{2, j}^{e_{3, j} - e_{1, j}}
u_{3, j}^{e_{1, j} - e_{2, j}} \bmod \mathfrak m_j = 1
$$
如果 $e_{1, j} = e_{2, j} = e_{3, j}$，这很显然。设 $e_{1, j} > e_{2, j}$。
由 $a_3 = a_1 - a_2$ 可知 $e_{3, j} = e_{2, j}$，且 $u_{3, j}$ 与
$-u_{2, j}$ 有相同的剩余类。因此公式成立，符号也恰好相符；
这项核查正是 (\ref{chow-equation-tame-symbol}) 中如此选择符号的原因。
其他情形的处理完全相同。
\end{proof}

\begin{lemma}
\label{chow-lemma-norm-down-tame-symbol}
设 $(A, \mathfrak m)$ 为维数为 $1$ 的诺特局部环。设 $A \subset B$ 是有限环扩张，
$B/A$ 被 $\mathfrak m$ 的某个幂零化，并且 $\mathfrak m$ 不是 $B$ 的伴随素理想。
对非零因子 $a, b \in A$，有
$$
\partial_A(a, b) = \prod
\text{Norm}_{\kappa(\mathfrak m_j)/\kappa(\mathfrak m)}(\partial_{B_j}(a, b))
$$
其中乘积遍历 $B$ 的极大理想 $\mathfrak m_j$，且 $B_j = B_{\mathfrak m_j}$。
\end{lemma}

\begin{proof}
为 $a, b$ 选择引理 \ref{chow-lemma-prepare-tame-symbol} 中的 $B_j \subset C_j$。
由引理 \ref{chow-lemma-glue-at-max}，可选取有限环扩张 $B \subset C$，
使得对所有 $j$ 都有 $C_j \cong C_{\mathfrak m_j}$。
设 $\mathfrak m_{j, k} \subset C$ 是位于 $\mathfrak m_j$ 上方的 $C$ 的极大理想。设
$$
a = u_{j, k}\pi_{j, k}^{f_{j, k}},\quad
b = v_{j, k}\pi_{j, k}^{g_{j, k}}
$$
为选择 $C_j \cong C_{\mathfrak m_j}$ 后得到的局部分解。由定义，
$$
\partial_A(a, b) = 
\prod\nolimits_{j, k}
\text{Norm}_{\kappa(\mathfrak m_{j, k})/\kappa(\mathfrak m)}
((-1)^{f_{j, k}g_{j, k}}u_{j, k}^{g_{j, k}}v_{j, k}^{-f_{j, k}}
\bmod \mathfrak m_{j, k})^{m_{j, k}}
$$
并且
$$
\partial_{B_j}(a, b) = 
\prod\nolimits_k
\text{Norm}_{\kappa(\mathfrak m_{j, k})/\kappa(\mathfrak m_j)}
((-1)^{f_{j, k}g_{j, k}}u_{j, k}^{g_{j, k}}v_{j, k}^{-f_{j, k}}
\bmod \mathfrak m_{j, k})^{m_{j, k}}
$$
对域塔 $\kappa(\mathfrak m_{j, k})/\kappa(\mathfrak m_j)/\kappa(\mathfrak m)$
应用范数的传递性即得结论；见《域》引理 \ref{fields-lemma-trace-and-norm-tower}。
\end{proof}

\begin{lemma}
\label{chow-lemma-tame-symbol-formally-smooth}
设 $(A, \mathfrak m, \kappa) \to (A', \mathfrak m', \kappa')$
是诺特局部环的局部同态。假设 $A \to A'$ 平坦，且 $\dim(A) = \dim(A') = 1$。
令 $m = \text{length}_{A'}(A'/\mathfrak mA')$。
对非零因子 $a_1, a_2 \in A$，$\partial_A(a_1, a_2)^m$ 经
$\kappa \to \kappa'$ 映为 $\partial_{A'}(a_1, a_2)$。
\end{lemma}

\begin{proof}
如果 $a_1, a_2$ 都是单位，那么 $\partial_A(a_1, a_2) = 1$ 且
$\partial_{A'}(a_1, a_2) = 1$，结论成立。否则，可以选择引理 \ref{chow-lemma-prepare-tame-symbol} 中的环扩张 $A \subset B$ 与局部分解。
以 $\mathfrak m_1, \ldots, \mathfrak m_m$ 表示 $B$ 的极大理想。
设 $\mathfrak m_1, \ldots, \mathfrak m_m$ 是 $B$ 的极大理想，其剩余域为
$\kappa_1, \ldots, \kappa_m$。对每个 $j \in \{1, \ldots, m\}$，
以 $\pi_j \in B_j = B_{\mathfrak m_j}$ 表示一个非零因子，使得有引理中的分解
$a_i = u_{i, j}\pi_j^{e_{i, j}}$。由定义，
$$
\partial_A(a_1, a_2) = \prod\nolimits_j
\text{Norm}_{\kappa_j/\kappa}
((-1)^{e_{1, j}e_{2, j}}u_{1, j}^{e_{2, j}}u_{2, j}^{-e_{1, j}}
\bmod \mathfrak m_j)^{m_j}
$$
其中 $m_j = \text{length}_{B_j}(B_j/\pi_j B_j)$。

\medskip\noindent
令 $B' = A' \otimes_A B$。由于 $A'$ 在 $A$ 上平坦，$A' \subset B'$ 是环扩张，
且 $B'/A'$ 被 $\mathfrak m'$ 的某个幂零化。设
$$
\mathfrak m'_{j, l},\quad l = 1, \ldots, n_j
$$
为 $B'$ 中位于 $\mathfrak m_j$ 上方的极大理想。以 $\kappa'_{j, l}$ 表示
$\mathfrak m'_{j, l}$ 的剩余域，以 $B'_{j, l}$ 表示 $B'$ 在 $\mathfrak m'_{j, l}$ 处的局部化。
在 $B'_{j, l}$ 中分解 $a_1$、$a_2$ 时，采用 $B_j$ 中给定分解
$a_i = u_{i, j} \pi_j^{e_{i, j}}$ 的像。由定义，
$$
\partial_{A'}(a_1, a_2) = \prod\nolimits_{j, l}
\text{Norm}_{\kappa'_{j, l}/\kappa'}
((-1)^{e_{1, j}e_{2, j}}u_{1, j}^{e_{2, j}}u_{2, j}^{-e_{1, j}}
\bmod \mathfrak m'_{j, l})^{m'_{j, l}}
$$
其中 $m'_{j, l} = \text{length}_{B'_{j, l}}(B'_{j, l}/\pi_j B'_{j, l})$。

\medskip\noindent
比较上述公式可知，只需证明：对每个 $j$ 和任意单位 $u \in B_j$，都有
\begin{equation}
\label{chow-equation-to-prove}
\left(\text{Norm}_{\kappa_j/\kappa}(u \bmod \mathfrak m_j)^{m_j}\right)^m
=
\prod\nolimits_l
\text{Norm}_{\kappa'_{j, l}/\kappa'}(u \bmod \mathfrak m'_{j, l})^{m'_{j, l}}
\end{equation}
在 $\kappa'$ 中成立。我们用《代数进阶》第 \ref{more-algebra-section-determinants-finite-length} 节中有限长度模的自同态行列式构造证明这一点。
令 $M = B_j/\pi_j B_j$。由《代数进阶》引理 \ref{more-algebra-lemma-multiplication}，有
$$
\text{Norm}_{\kappa_j/\kappa}(u \bmod \mathfrak m_j)^{m_j} =
\det\nolimits_\kappa(u : M \to M)
$$
因此，由《代数进阶》引理 \ref{more-algebra-lemma-flat-base-change-det}，
(\ref{chow-equation-to-prove}) 的左边等于
$\det_{\kappa'}(u : M \otimes_A A' \to M \otimes_A A')$。
有 $A'$-模同构
$$
M \otimes_A A' = (B_j/\pi_j B_j) \otimes_A A' =
\bigoplus\nolimits_l B'_{j, l}/\pi_j B'_{j, l}
$$
令 $M'_l = B'_{j, l}/\pi_j B'_{j, l}$。再次由《代数进阶》引理 \ref{more-algebra-lemma-multiplication}，可得
$\text{Norm}_{\kappa'_{j, l}/\kappa'}(u \bmod \mathfrak m'_{j, l})^{m'_{j, l}}
= \det_{\kappa'}(u_j : M'_l \to M'_l)$。
所以 (\ref{chow-equation-to-prove}) 由行列式构造的乘法性成立；见《代数进阶》引理 \ref{more-algebra-lemma-ses}。
\end{proof}






\section{一个关键引理}
\label{chow-section-key-lemma}

\noindent
本节应用上述结果证明引理 \ref{chow-lemma-milnor-gersten-low-degree}。
该引理是如下陈述的低次数情形：Milnor K-理论存在一个类似于 Quillen K-理论中
Gersten--Quillen 复形的复形。见注 \ref{chow-remark-gersten-complex-milnor}。

\begin{lemma}
\label{chow-lemma-perpare-key}
设 $(A, \mathfrak m)$ 是一个 $2$-维诺特局部环。设 $t \in \mathfrak m$ 为非零因子，并写成
$V(t) = \{\mathfrak m, \mathfrak q_1, \ldots, \mathfrak q_r\}$。
设 $A_{\mathfrak q_i} \subset B_i$ 是有限环扩张，其中
$B_i/A_{\mathfrak q_i}$ 被 $t$ 的某个幂零化。那么存在保持剩余域相同的有限局部环扩张
$A \subset B$，使得 $B_i \cong B_{\mathfrak q_i}$，并且 $B/A$ 被 $t$ 的某个幂零化。
\end{lemma}

\begin{proof}
选取 $n > 0$，使 $B_i \subset t^{-n}A_{\mathfrak q_i}$。
令 $M \subset t^{-n}A$ 和 $M' \subset t^{-2n}A$ 分别为\ 由如下元素组成的 $A$-子模：
它们在 $t^{-n}A_{\mathfrak q_i}$ 和 $t^{-2n}A_{\mathfrak q_i}$ 中分别\ 映入 $B_i$。
由于 $A$ 诺特，$M \subset M'$ 是有限 $A$-模；又因局部化正合，
$M_{\mathfrak q_i} = M'_{\mathfrak q_i} = B_i$。因此对某个 $c > 0$，
$M'/M$ 被 $\mathfrak m^c$ 零化。注意，在乘法
$t^{-n}A \times t^{-n}A \to t^{-2n}A$ 下，有 $M \cdot M \subset M'$。
故 $B = A + \mathfrak m^{c + 1}M$ 是有限 $A$-代数，并具有所需的局部化。
我们略去如下验证：$B$ 是局部环，其极大理想为
$\mathfrak m + \mathfrak m^{c + 1}M$。
\end{proof}

\begin{lemma}
\label{chow-lemma-key-nonzerodivisors}
设 $(A, \mathfrak m)$ 是一个 $2$-维诺特局部环，$a, b \in A$ 是非零因子。那么
$$
\sum
\text{ord}_{A/\mathfrak q}(\partial_{A_{\mathfrak q}}(a, b))
=
0
$$
其中求和遍历 $A$ 中所有高度为 $1$ 的素理想 $\mathfrak q$。
\end{lemma}

\begin{proof}
若 $\mathfrak q$ 是 $A$ 中高度为 $1$ 的素理想，并且 $a, b$ 在 $A_\mathfrak q$ 中的像是单位，
则 $\partial_{A_\mathfrak q}(a, b) = 1$。因此该和为有限和。事实上，若
$V(ab) = \{\mathfrak m, \mathfrak q_1, \ldots, \mathfrak q_r\}$，
则求和遍历 $i = 1, \ldots, r$。对每个 $i$，为 $a, b$ 选取引理 \ref{chow-lemma-prepare-tame-symbol} 中的扩张 $A_{\mathfrak q_i} \subset B_i$。
把引理 \ref{chow-lemma-perpare-key} 用于 $t = ab$ 与给定的素理想列表，
可假设存在有限局部扩张 $A \subset B$，其中 $B/A$ 被 $ab$ 的某个幂零化，
并且对每个 $i$ 都有 $B_{\mathfrak q_i} \cong B_i$。
注意，若 $\mathfrak q_{i, j}$ 是 $B$ 中位于 $\mathfrak q_i$ 上方的素理想，则由
引理 \ref{chow-lemma-norm-down-tame-symbol} 和《交换代数》引理 \ref{algebra-lemma-finite-extension-dim-1}，有
$$
\text{ord}_{A/\mathfrak q_i}(\partial_{A_{\mathfrak q_i}}(a, b))
=
\sum\nolimits_j
\text{ord}_{B/\mathfrak q_{i, j}}(\partial_{B_{\mathfrak q_{i, j}}}(a, b))
$$
因此可以用 $B$ 代替 $A$，并约化到下一段讨论的情形。

\medskip\noindent
假设对每个 $i$，存在非零因子 $\pi_i \in A_{\mathfrak q_i}$ 和单位
$u_i, v_i \in A_{\mathfrak q_i}$，使得对某些整数 $e_i, f_i \geq 0$，在
$A_{\mathfrak q_i}$ 中有
$$
a = u_i \pi_i^{e_i},\quad b = v_i \pi_i^{f_i}
$$
令 $m_i = \text{length}_{A_{\mathfrak q_i}}(A_{\mathfrak q_i}/\pi_i)$。
根据定义，
$\partial_{A_{\mathfrak q_i}}(a, b) =
((-1)^{e_if_i}u_i^{f_i}v_i^{-e_i})^{m_i}$。
由于 $a, b$ 是非零因子，$(2, 1)$-周期复形 $(A/(ab), a, b)$ 的上同调为零。
以 $M_i$ 表示 $A/(ab)$ 在 $A_{\mathfrak q_i}/(ab)$ 中的像。于是有映射
$$
A/(ab) \longrightarrow \bigoplus M_i
$$
其核与余核均支撑在 $\{\mathfrak m\}$ 上，因而具有有限长度。因此，由引理 \ref{chow-lemma-compare-periodic-lengths}，有
$$
\sum e_A(M_i, a, b) = 0
$$
故只须证明
$e_A(M_i, a, b) =
- \text{ord}_{A/\mathfrak q_i}(\partial_{A_{\mathfrak q_i}}(a, b))$。

\medskip\noindent
先证明 $\pi_i, u_i, v_i$ 是某些元素 $\pi_i, u_i, v_i \in A$ 的像这一情形
（使用相同符号不会引起混淆）。此时
\begin{align*}
e_A(M_i, a, b) & =
e_A(M_i, u_i\pi_i^{e_i}, v_i\pi_i^{f_i}) \\
& =
e_A(M_i, \pi_i^{e_i}, \pi_i^{f_i}) -
e_A(\pi_i^{e_i}M_i, 0, u_i) +
e_A(\pi_i^{f_i}M_i, 0, v_i) \\
& =
0 -
f_im_i\text{ord}_{A/\mathfrak q_i}(u_i) +
e_im_i\text{ord}_{A/\mathfrak q_i}(v_i) \\
& =
-m_i\text{ord}_{A/\mathfrak q_i}(u_i^{f_i}v_i^{-e_i}) =
-\text{ord}_{A/\mathfrak q_i}(\partial_{A_{\mathfrak q_i}}(a, b))
\end{align*}
第二个等式来自引理 \ref{chow-lemma-multiply-period-length}。注意
$M_i \subset (M_i)_{\mathfrak q_i} = A_{\mathfrak q_i}/(\pi_i^{e_i + f_i})$，并且
$(\pi_i^{e_i}M_i)_{\mathfrak q_i} \cong A_{\mathfrak q_i}/\pi_i^{f_i}$、
$(\pi_i^{f_i}M_i)_{\mathfrak q_i} \cong A_{\mathfrak q_i}/\pi_i^{e_i}$。
第三个等式中的 $0$ 来自引理 \ref{chow-lemma-powers-period-length-zero}，
另外两项来自引理 \ref{chow-lemma-length-multiplication}。
最后两个等式由阶函数的乘法性和驯顺符号的定义得到。

\medskip\noindent
在一般情形下，先选取 $c \in A$、$c \not \in \mathfrak q_i$，使 $c\pi_i \in A$。
以 $c\pi_i$ 代替 $\pi_i$，以 $c^{-e_i}u_i$ 代替 $u_i$，并以
$c^{-f_i}v_i$ 代替 $v_i$ 后，可以而且将假设 $\pi_i$ 属于 $A$。
接着选取 $c \in A$、$c \not \in \mathfrak q_i$，使 $cu_i, cv_i \in A$。
由引理 \ref{chow-lemma-length-multiplication}，有
$$
e_A(M_i, ca, cb) = e_A(M_i, a, b) - e_A(aM_i, 0, c) + e_A(bM_i, 0, c)
$$
另一方面，由于 $c$ 是 $A_{\mathfrak q_i}$ 中的单位，在 $\kappa(\mathfrak q_i)^*$ 中有
$$
\partial_{A_{\mathfrak q_i}}(ca, cb) =
c^{m_i(f_i - e_i)}\partial_{A_{\mathfrak q_i}}(a, b)
$$
上一段的论证表明
$e_A(M_i, ca, cb) = -
\text{ord}_{A/\mathfrak q_i}(\partial_{A_{\mathfrak q_i}}(ca, cb))$。
因此只须证明
$$
e_A(aM_i, 0, c) = \text{ord}_{A/\mathfrak q_i}(c^{m_if_i})
\quad\text{且}\quad
e_A(bM_i, 0, c) = \text{ord}_{A/\mathfrak q_i}(c^{m_ie_i})
$$
这由引理 \ref{chow-lemma-length-multiplication} 以及上文对在 $\mathfrak q_i$ 处局部化所得情形的描述推出。
\end{proof}

\begin{lemma}[关键引理]
\label{chow-lemma-milnor-gersten-low-degree}
\begin{reference}
当 $A$ 是优秀环时，这是 \cite[Proposition 1]{Kato-Milnor-K}。
\end{reference}
设 $A$ 是一个 $2$-维诺特局部整环，其分式域为 $K$。设 $f, g \in K^*$。
设 $\mathfrak q_1, \ldots, \mathfrak q_t$ 是 $A$ 中满足如下条件的高度为 $1$ 的素理想
$\mathfrak q$：$f$ 或 $g$ 至少有一个不属于 $A^*_{\mathfrak q}$。那么
$$
\sum\nolimits_{i = 1, \ldots, t}
\text{ord}_{A/\mathfrak q_i}(\partial_{A_{\mathfrak q_i}}(f, g))
=
0
$$
也可以写成
$$
\sum\nolimits_{\text{高度}(\mathfrak q) = 1}
\text{ord}_{A/\mathfrak q}(\partial_{A_{\mathfrak q}}(f, g))
=
0
$$
因为对 $A$ 中任意高度为 $1$ 的素理想 $\mathfrak q$，若
$f, g \in A^*_{\mathfrak q}$，则 $\partial_{A_{\mathfrak q}}(f, g) = 1$。
\end{lemma}

\begin{proof}
驯顺符号 $\partial_{A_{\mathfrak q}}(f, g)$ 是双线性的，而阶函数
$\text{ord}_{A/\mathfrak q}$ 是可加的，所以只需在 $f$、$g$ 都属于 $A$ 时证明该公式。
这一情形已在引理 \ref{chow-lemma-key-nonzerodivisors} 中证明。
\end{proof}

\begin{remark}[Milnor K-理论]
\label{chow-remark-gersten-complex-milnor}
对域 $k$，以 $K^M_*(k)$ 表示 $k^*$ 上的张量代数对如下双边理想的商：
该理想由所有 $x \in k \setminus \{0, 1\}$ 对应的元素
$x \otimes 1 - x$ 生成。因此 $K^M_0(k) = \mathbf{Z}$、$K_1^M(k) = k^*$，并且
$$
K^M_2(k) = k^* \otimes_\mathbf{Z} k^* / \langle x \otimes 1 - x \rangle
$$
若 $A$ 是分式域为 $F = \text{Frac}(A)$、剩余域为 $\kappa$ 的离散赋值环，
则有如第 \ref{chow-section-tame-symbol} 节定义的驯顺符号
$$
\partial_A : K_{i + 1}^M(F) \to K_i^M(\kappa)
$$
见 \cite{Kato-Milnor-K}。更一般地，利用正规化以及 $K_i^M$ 上的范数映射，可以把这一映射推广到
$A$ 为维数为 $1$ 的优秀局部整环的情形；见 \cite{Kato-Milnor-K}。
可以推测，第 \ref{chow-section-tame-symbol} 节的方法能够把驯顺符号 $\partial_A$ 的构造推广到任意维数为
$1$ 的诺特局部整环 $A$。接着，设 $X$ 是带有维数函数 $\delta$ 的诺特概形。
利用这些驯顺符号，可以得到下列箭头：
$$
\bigoplus\nolimits_{\delta(x) = j + 1} K^M_{i + 1}(\kappa(x))
\longrightarrow
\bigoplus\nolimits_{\delta(x) = j} K^M_i(\kappa(x))
\longrightarrow
\bigoplus\nolimits_{\delta(x) = j - 1} K^M_{i - 1}(\kappa(x))
$$
但尚不清楚该复合是否为零，亦即是否由此得到一个阿贝尔群复形。
对优秀概形 $X$，这在 \cite{Kato-Milnor-K} 中得到证明。
当 $i = 1$ 而 $j$ 任意时，结论来自引理 \ref{chow-lemma-milnor-gersten-low-degree}。
\end{remark}





















\section{基本设置}
\label{chow-section-setup}

\noindent
在本章中，我们始终在一个配有维数函数 $\delta$ 的局部诺特、泛链状基概形 $S$ 上工作。
尽管本章的讨论（至少一部分）很可能可以推广，但这似乎是一个良好的第一步近似；
它恰好就是 \cite{Thorup} 所讨论的一般程度。通常不假设所研究的概形是分离的或拟紧的。
许多有趣的代数栈并不分离和/或并不拟紧，因此这里也是研究如何为这些对象发展合理理论的一个好例子。
为了便于引用这些假设，我们为其编号。

\begin{situation}
\label{chow-situation-setup}
这里 $S$ 是一个局部诺特且泛链状的概形。此外，假设 $S$ 配有维数函数
$\delta : S \longrightarrow \mathbf{Z}$。
\end{situation}

\noindent
泛链状概形的概念见《概形态射》定义 \ref{morphisms-definition-universally-catenary}；维数函数的概念见《拓扑学》定义 \ref{topology-definition-dimension-function}。回忆，任意局部诺特链状概形在局部都有维数函数；
见《概形的性质》引理 \ref{properties-lemma-catenary-dimension-function}。
此外，泛链状概形相当常见；见《概形态射》引理 \ref{morphisms-lemma-ubiquity-uc}。

\medskip\noindent
设 $(S, \delta)$ 如情形 \ref{chow-situation-setup}。任意在 $S$ 上局部有限型的概形 $X$
都是局部诺特且链状的。事实上，$X$ 有一个与 $(f : X \to S, \delta)$ 相伴的典范维数函数
$$
\delta = \delta_{X/S} : X \longrightarrow \mathbf{Z}
$$
其规则为
$\delta_{X/S}(x) = \delta(f(x)) + \text{trdeg}_{\kappa(f(x))}\kappa(x)$。
见《概形态射》引理 \ref{morphisms-lemma-dimension-function-propagates}。
此外，若 $h : X \to Y$ 是 $S$ 上局部有限型概形之间的态射，且
$x \in X$、$y = h(x)$，则显然
$\delta_{X/S}(x) = \delta_{Y/S}(y) + \text{trdeg}_{\kappa(y)}\kappa(x)$。
下文将自由使用这个函数及其性质。

\medskip\noindent
下面给出上述设置的基本例子。事实上，最值得关注的是基为一个域的谱，
或者基为 Dedekind 环的谱（例如\ $\mathbf{Z}$ 或离散赋值环）的情形。

\begin{example}
\label{chow-example-field}
这里 $S = \Spec(k)$，其中 $k$ 是域。令 $\delta(pt) = 0$，其中 $pt$ 表示 $S$ 的唯一点。
$(S, \delta)$ 是情形 \ref{chow-situation-setup} 中设置的一个例子；
这来自《概形态射》引理 \ref{morphisms-lemma-ubiquity-uc}。
\end{example}

\begin{example}
\label{chow-example-domain-dimension-1}
这里 $S = \Spec(A)$，其中 $A$ 是维数为 $1$ 的诺特整环。例如可以取 $A = \mathbf{Z}$。
若 $\mathfrak p$ 是极大理想，令 $\delta(\mathfrak p) = 0$；若
$\mathfrak p = (0)$ 对应一般点，令 $\delta(\mathfrak p) = 1$。
这是情形 \ref{chow-situation-setup} 的一个例子；
结论来自《概形态射》引理 \ref{morphisms-lemma-ubiquity-uc}。
\end{example}

\begin{example}
\label{chow-example-CM-irreducible}
这里 $S$ 是 Cohen--Macaulay 概形。由《概形态射》引理 \ref{morphisms-lemma-ubiquity-uc}，$S$ 泛链状。令
$\delta(s) = -\dim(\mathcal{O}_{S, s})$。若 $s' \leadsto s$ 是 $S$ 中点的非平凡特殊化，
则 $\mathcal{O}_{S, s'}$ 是 $\mathcal{O}_{S, s}$ 在非极大素理想
$\mathfrak p \subset \mathcal{O}_{S, s}$ 处的局部化；见《概形》引理 \ref{schemes-lemma-specialize-points}。因此由《交换代数》引理 \ref{algebra-lemma-CM-dim-formula}，有
$\dim(\mathcal{O}_{S, s}) =
\dim(\mathcal{O}_{S, s'}) + \dim(\mathcal{O}_{S, s}/\mathfrak p) >
\dim(\mathcal{O}_{S, s'})$。故 $\delta(s') > \delta(s)$。
若 $s' \leadsto s$ 是紧邻特殊化，则 $\mathfrak p$ 与 $\mathfrak m_s$ 之间没有严格介于二者的素理想，
从而 $\delta(s') = \delta(s) + 1$。所以 $\delta$ 是维数函数。
换言之，$(S, \delta)$ 是情形 \ref{chow-situation-setup} 的一个例子。
\end{example}

\noindent
如果 $S$ 是 Jacobson 概形，而 $\delta$ 把闭点映为零，那么 $\delta$ 就是把每个点映到其闭包维数的函数。

\begin{lemma}
\label{chow-lemma-delta-is-dimension}
设 $(S, \delta)$ 如情形 \ref{chow-situation-setup}。进一步假设 $S$ 是 Jacobson 概形，
并且对 $S$ 的每个闭点 $s$ 都有 $\delta(s) = 0$。设 $X$ 在 $S$ 上局部有限型。
设 $Z \subset X$ 是整闭子概形，$\xi \in Z$ 是它的一般点。下列整数相同：
\begin{enumerate}
\item $\delta_{X/S}(\xi)$；
\item $\dim(Z)$；以及
\item $\dim(\mathcal{O}_{Z, z})$，其中 $z$ 是 $Z$ 的闭点。
\end{enumerate}
\end{lemma}

\begin{proof}
设 $X \to S$、$\xi \in Z \subset X$ 如引理所述。由于 $X$ 在 $S$ 上局部有限型，
$X$ 是 Jacobson 概形；见《概形态射》引理 \ref{morphisms-lemma-Jacobson-universally-Jacobson}。因此，$X$ 的闭点在 $Z$ 的每个闭子集中稠密，
并映到 $S$ 的闭点。所以，给定 $Z$ 中不可约闭子集的任意链，可以用 $Z$ 的某个闭点将其终结。
从而，当 $z \in Z$ 遍历 $Z$ 的闭点时，
$\dim(Z) = \sup_z(\dim(\mathcal{O}_{Z, z})$；
见《概形的性质》引理 \ref{properties-lemma-codimension-local-ring}。
由维数函数的性质，$\dim(\mathcal{O}_{Z, z}) = \delta(\xi) - \delta(z)$。
对每个闭点 $z \in Z$，域扩张 $\kappa(z)/\kappa(f(z))$ 有限；见《概形态射》引理 \ref{morphisms-lemma-jacobson-finite-type-points}。
因此，对闭点 $z \in Z$，有 $\delta_{X/S}(z) = \delta(f(z)) = 0$。
由此可知三个整数全都相等。
\end{proof}

\noindent
在上述引理的情形下，$\delta$ 在一个闭不可约子集的一般点处的值，就是该不可约闭子集的维数。
但一般来说不能期待这一等式成立。例如，若 $S = \Spec(\mathbf{C}[[t]])$ 且
$X = \Spec(\mathbf{C}((t)))$，则对 $X$ 的唯一点有 $\delta(x) = 1$，但
$\dim(X) = 0$。尽管如此，我们仍希望把 $\delta_{X/S}$ 视作给出不可约闭子概形的维数。
因此引入以下术语。

\begin{definition}
\label{chow-definition-delta-dimension}
设 $(S, \delta)$ 如情形 \ref{chow-situation-setup}。对任意在 $S$ 上局部有限型的概形 $X$
和任意不可约闭子集 $Z \subset X$，定义
$$
\dim_\delta(Z) = \delta(\xi)
$$
其中 $\xi \in Z$ 是 $Z$ 的一般点。称其为{\it $Z$ 的 $\delta$-维数}。
如果 $Z$ 是 $X$ 的闭子概形，则定义 $\dim_\delta(Z)$ 为其不可约分支的
$\delta$-维数的上确界。
\end{definition}







\section{循环}
\label{chow-section-cycles}

\noindent
由于不假设所研究的概形拟紧，定义循环时必须稍加谨慎。例如，有理函数可能有无穷多个极点，
因此必须允许无穷和。不过，如果 $X$ 拟紧，那么循环仍像通常一样是有限和。

\begin{definition}
\label{chow-definition-cycles}
设 $(S, \delta)$ 如情形 \ref{chow-situation-setup}，设 $X$ 在 $S$ 上局部有限型，并设
$k \in \mathbf{Z}$。
\begin{enumerate}
\item $X$ 上的一个{\it 循环}是形式和
$$
\alpha = \sum n_Z [Z]
$$
其中求和遍历整闭子概形 $Z \subset X$，每个 $n_Z \in \mathbf{Z}$，并且集合
$\{Z; n_Z \not = 0\}$ 局部有限
（《拓扑学》定义 \ref{topology-definition-locally-finite}）。
\item $X$ 上的一个 {\it $k$-循环}是满足下述条件的循环
$$
\alpha = \sum n_Z [Z]
$$
即 $n_Z \not = 0 \Rightarrow \dim_\delta(Z) = k$。
\item $X$ 上所有 $k$-循环构成的阿贝尔群记为 $Z_k(X)$。
\end{enumerate}
\end{definition}

\noindent
换言之，$X$ 上的 $k$-循环是 $\delta$-维数为 $k$ 的整闭子概形的局部有限形式
$\mathbf{Z}$-线性组合。$k$-循环 $\alpha = \sum n_Z[Z]$ 与
$\beta = \sum m_Z[Z]$ 的加法定义为
$$
\alpha + \beta = \sum (n_Z + m_Z)[Z],
$$
也就是把系数相加。

\begin{remark}
\label{chow-remark-cycles-pointwise}
设 $(S, \delta)$ 如情形 \ref{chow-situation-setup}，设 $X$ 在 $S$ 上局部有限型，并设
$k \in \mathbf{Z}$。可以写成
$$
Z_k(X) = \bigoplus\nolimits_{\delta(x) = k}' K_0^M(\kappa(x))
\quad\subset\quad
\bigoplus\nolimits_{\delta(x) = k} K_0^M(\kappa(x))
$$
其中采用以下记号与约定：
\begin{enumerate}
\item $K_0^M(\kappa(x)) = \mathbf{Z}$ 是点 $x \in X$ 的剩余域 $\kappa(x)$ 的
Milnor K-理论的次数 $0$ 部分（见注 \ref{chow-remark-gersten-complex-milnor}）；
\item 右边的直和遍历满足 $\delta(x) = k$ 的所有点 $x \in X$；
\item 记号 $\bigoplus'_x$ 表示只取由局部有限元素构成的子群；也就是取满足如下条件的元素
$\sum_x n_x$：对每个拟紧开集 $U \subset X$，满足 $n_x \not = 0$ 的点 $x \in U$ 只有有限多个。
\end{enumerate}
\end{remark}

\begin{definition}
\label{chow-definition-support-cycle}
设 $(S, \delta)$ 如情形 \ref{chow-situation-setup}，并设 $X$ 在 $S$ 上局部有限型。
$X$ 上循环 $\alpha = \sum n_Z [Z]$ 的{\it 支撑}为
$$
\text{Supp}(\alpha) = \bigcup\nolimits_{n_Z \not = 0} Z \subset X
$$
\end{definition}

\noindent
由于集合 $\{Z; n_Z \not = 0\}$ 局部有限，$\text{Supp}(\alpha)$ 是 $X$ 的闭子集。
若 $\alpha$ 是 $k$-循环，则 $\text{Supp}(\alpha)$ 的每个不可约分支 $Z$ 的
$\delta$-维数都为 $k$。

\begin{definition}
\label{chow-definition-effective-cycle}
设 $(S, \delta)$ 如情形 \ref{chow-situation-setup}，并设 $X$ 在 $S$ 上局部有限型。
$X$ 上的循环 $\alpha$ 如果可以写成 $\alpha =\sum n_Z [Z]$，且对所有 $Z$ 都有 $n_Z \geq 0$，
就称该循环是{\it 有效的}。
\end{definition}

\noindent
所有有效循环构成一个幺半群，因为两个有效循环之和仍有效；但它不是群
（除非 $X = \emptyset$）。




\section{与闭子概形相伴的循环}
\label{chow-section-cycle-of-closed-subscheme}

\begin{lemma}
\label{chow-lemma-multiplicity-finite}
设 $(S, \delta)$ 如情形 \ref{chow-situation-setup}，设 $X$ 在 $S$ 上局部有限型，
并设 $Z \subset X$ 是闭子概形。
\begin{enumerate}
\item 设 $Z' \subset Z$ 是一个不可约分支，$\xi \in Z'$ 是它的一般点。那么
$$
\text{length}_{\mathcal{O}_{X, \xi}} \mathcal{O}_{Z, \xi} < \infty
$$
\item 如果 $\dim_\delta(Z) \leq k$，并且 $\xi \in Z$ 满足
$\delta(\xi) = k$，那么 $\xi$ 是 $Z$ 的某个不可约分支的一般点。
\end{enumerate}
\end{lemma}

\begin{proof}
设 $Z' \subset Z$、$\xi \in Z'$ 如 (1) 所述。那么
$\dim(\mathcal{O}_{Z, \xi}) = 0$；例如见《概形的性质》引理 \ref{properties-lemma-codimension-local-ring}。因此 $\mathcal{O}_{Z, \xi}$ 是维数为零的诺特局部环，
从而作为自身的模具有有限长度（见《交换代数》命题 \ref{algebra-proposition-dimension-zero-ring}）。所以它作为 $\mathcal{O}_{X, \xi}$-模也具有有限长度；
见《交换代数》引理 \ref{algebra-lemma-length-independent}。

\medskip\noindent
假设 $\xi \in Z$ 且 $\delta(\xi) = k$。考虑闭包
$Z' = \overline{\{\xi\}}$。根据定义，它是满足 $\dim_\delta(Z') = k$ 的不可约闭子概形。
由于 $\dim_\delta(Z) = k$，它必为 $Z$ 的一个不可约分支。因此 (2) 成立。
\end{proof}

\begin{definition}
\label{chow-definition-cycle-associated-to-closed-subscheme}
设 $(S, \delta)$ 如情形 \ref{chow-situation-setup}，设 $X$ 在 $S$ 上局部有限型，
并设 $Z \subset X$ 是闭子概形。
\begin{enumerate}
\item 对任意一般点为 $\xi$ 的不可约分支 $Z' \subset Z$，称整数
$m_{Z', Z} = \text{length}_{\mathcal{O}_{X, \xi}} \mathcal{O}_{Z, \xi}$
（引理 \ref{chow-lemma-multiplicity-finite}）为 {\it $Z'$ 在 $Z$ 中的重数}。
\item 假设 $\dim_\delta(Z) \leq k$。定义{\it 与 $Z$ 相伴的 $k$-循环}为
$$
[Z]_k
=
\sum m_{Z', Z}[Z']
$$
其中求和遍历 $Z$ 中 $\delta$-维数为 $k$ 的不可约分支。
（由《除子》引理 \ref{divisors-lemma-components-locally-finite}，这是一个 $k$-循环。）
\end{enumerate}
\end{definition}

\noindent
必须注意，只有在 $Z$ 的 $\delta$-维数不超过 $k$ 时才定义 $[Z]_k$。
换言之，按约定，只要写下 $[Z]_k$，就隐含着 $\dim_\delta(Z) \leq k$。




\section{与凝聚层相伴的循环}
\label{chow-section-cycle-of-coherent-sheaf}



\begin{lemma}
\label{chow-lemma-length-finite}
设 $(S, \delta)$ 如情形 \ref{chow-situation-setup}，设 $X$ 在 $S$ 上局部有限型，
并设 $\mathcal{F}$ 是凝聚 $\mathcal{O}_X$-模。
\begin{enumerate}
\item $\mathcal{F}$ 的支撑的不可约分支所成集合局部有限。
\item 设 $Z' \subset \text{Supp}(\mathcal{F})$ 是一个不可约分支，
$\xi \in Z'$ 是它的一般点。那么
$$
\text{length}_{\mathcal{O}_{X, \xi}} \mathcal{F}_\xi < \infty
$$
\item 如果 $\dim_\delta(\text{Supp}(\mathcal{F})) \leq k$，并且
$\xi \in \text{Supp}(\mathcal{F})$ 满足 $\delta(\xi) = k$，那么 $\xi$ 是
$\text{Supp}(\mathcal{F})$ 的某个不可约分支的一般点。
\end{enumerate}
\end{lemma}

\begin{proof}
由《概形上同调》引理 \ref{coherent-lemma-coherent-support-closed}，
$\mathcal{F}$ 的支撑 $Z$ 是 $X$ 的闭子集。可以把 $Z$ 看成 $X$ 的约化闭子概形
（《概形》引理 \ref{schemes-lemma-reduced-closed-subscheme}）。
因此，把《除子》引理 \ref{divisors-lemma-components-locally-finite} 用于 $Z$ 即得 (1)，
把引理 \ref{chow-lemma-multiplicity-finite} 用于 $Z$ 即得 (3)。

\medskip\noindent
设 $\xi \in Z'$ 如 (2) 所述。在此情形下，对 $X$ 中任意特殊化
$\xi' \leadsto \xi$，都有 $\mathcal{F}_{\xi'} = 0$。回忆，
$\mathcal{O}_{X, \xi}$ 的非极大素理想对应于 $X$ 中特殊化到 $\xi$ 的点
（《概形》引理 \ref{schemes-lemma-specialize-points}）。
因此 $\mathcal{F}_\xi$ 是有限 $\mathcal{O}_{X, \xi}$-模，其支撑为
$\{\mathfrak m_\xi\}$。所以由《交换代数》引理 \ref{algebra-lemma-support-point}，
它具有有限长度。
\end{proof}

\begin{definition}
\label{chow-definition-cycle-associated-to-coherent-sheaf}
设 $(S, \delta)$ 如情形 \ref{chow-situation-setup}，设 $X$ 在 $S$ 上局部有限型，
并设 $\mathcal{F}$ 是凝聚 $\mathcal{O}_X$-模。
\begin{enumerate}
\item 对任意一般点为 $\xi$ 的不可约分支
$Z' \subset \text{Supp}(\mathcal{F})$，称整数
$m_{Z', \mathcal{F}} = \text{length}_{\mathcal{O}_{X, \xi}} \mathcal{F}_\xi$
（引理 \ref{chow-lemma-length-finite}）为 {\it $Z'$ 在 $\mathcal{F}$ 中的重数}。
\item 假设 $\dim_\delta(\text{Supp}(\mathcal{F})) \leq k$。定义
{\it 与 $\mathcal{F}$ 相伴的 $k$-循环}为
$$
[\mathcal{F}]_k
=
\sum m_{Z', \mathcal{F}}[Z']
$$
其中求和遍历 $\text{Supp}(\mathcal{F})$ 中 $\delta$-维数为 $k$ 的不可约分支。
（由引理 \ref{chow-lemma-length-finite}，这是一个 $k$-循环。）
\end{enumerate}
\end{definition}

\noindent
必须注意，只有当 $\mathcal{F}$ 凝聚，且 $\text{Supp}(\mathcal{F})$ 的
$\delta$-维数不超过 $k$ 时，才定义 $[\mathcal{F}]_k$。换言之，按约定，只要写下
$[\mathcal{F}]_k$，就隐含着 $\mathcal{F}$ 在 $X$ 上凝聚，且
$\dim_\delta(\text{Supp}(\mathcal{F})) \leq k$。

\begin{lemma}
\label{chow-lemma-cycle-closed-coherent}
设 $(S, \delta)$ 如情形 \ref{chow-situation-setup}，设 $X$ 在 $S$ 上局部有限型，
并设 $Z \subset X$ 是闭子概形。如果 $\dim_\delta(Z) \leq k$，那么
$[Z]_k = [{\mathcal O}_Z]_k$。
\end{lemma}

\begin{proof}
这是因为在此情形下，重数 $m_{Z', Z}$ 与 $m_{Z', \mathcal{O}_Z}$ 依定义相等。
\end{proof}

\begin{lemma}
\label{chow-lemma-additivity-sheaf-cycle}
设 $(S, \delta)$ 如情形 \ref{chow-situation-setup}，设 $X$ 在 $S$ 上局部有限型，
并设 $0 \to \mathcal{F} \to \mathcal{G} \to \mathcal{H} \to 0$
是 $X$ 上凝聚层的短正合列。假设 $\mathcal{F}$、$\mathcal{G}$ 与 $\mathcal{H}$ 的支撑的
$\delta$-维数都满足 $\leq k$。那么
$[\mathcal{G}]_k = [\mathcal{F}]_k + [\mathcal{H}]_k$。
\end{lemma}

\begin{proof}
这立即来自长度的可加性；见《交换代数》引理 \ref{algebra-lemma-length-additive}。
\end{proof}









\section{固有推前的准备工作}
\label{chow-section-preparation-pushforward}

\begin{lemma}
\label{chow-lemma-equal-dimension}
设 $(S, \delta)$ 如情形 \ref{chow-situation-setup}，设 $X$、$Y$ 在 $S$ 上局部有限型，
并设 $f : X \to Y$ 是态射。假设 $X$、$Y$ 都是整概形，且
$\dim_\delta(X) = \dim_\delta(Y)$。那么，或者 $f(X)$ 包含在 $Y$ 的某个真闭子概形中，
或者 $f$ 是支配态射，并且函数域扩张 $R(X)/R(Y)$ 有限。
\end{lemma}

\begin{proof}
由于 $X$ 不可约，闭包 $\overline{f(X)} \subset Y$ 不可约
（《拓扑学》引理 \ref{topology-lemma-image-irreducible-space} 和
\ref{topology-lemma-irreducible}）。如果 $\overline{f(X)} \not = Y$，结论已得。
如果 $\overline{f(X)} = Y$，则 $f$ 支配；由《概形态射》引理 \ref{morphisms-lemma-dominant-finite-number-irreducible-components}，
$Y$ 的一般点 $\eta_Y$ 位于 $f$ 的像中。这当然蕴含 $f(\eta_X) = \eta_Y$，
其中 $\eta_X \in X$ 是 $X$ 的一般点。由于 $\delta(\eta_X) = \delta(\eta_Y)$，
$R(Y) = \kappa(\eta_Y) \subset \kappa(\eta_X) = R(X)$ 是超越次数为 $0$ 的扩张。
所以由《概形态射》引理 \ref{morphisms-lemma-finite-degree}，
$R(Y) \subset R(X)$ 是有限扩张；该引理适用是由《概形态射》引理 \ref{morphisms-lemma-permanence-finite-type} 保证的。
\end{proof}

\begin{lemma}
\label{chow-lemma-quasi-compact-locally-finite}
设 $(S, \delta)$ 如情形 \ref{chow-situation-setup}，设 $X$、$Y$ 在 $S$ 上局部有限型，
并设 $f : X \to Y$ 是态射。假设 $f$ 拟紧，且 $\{Z_i\}_{i \in I}$ 是 $X$ 中闭子集的局部有限族。
那么 $\{\overline{f(Z_i)}\}_{i \in I}$ 是 $Y$ 中闭子集的局部有限族。
\end{lemma}

\begin{proof}
设 $V \subset Y$ 是拟紧开子集。由于 $f$ 拟紧，开集 $f^{-1}(V)$ 拟紧。
因此集合 $\{i \in I \mid Z_i \cap f^{-1}(V) \not = \emptyset \}$ 有限；
这里略去一个简单的拓扑论证。这个集合与
$$
\{i \in I \mid f(Z_i) \cap V \not = \emptyset \} =
\{i \in I \mid \overline{f(Z_i)} \cap V \not = \emptyset \}
$$
相同，故引理得证。
\end{proof}









\section{固有推前}
\label{chow-section-proper-pushforward}

\begin{definition}
\label{chow-definition-proper-pushforward}
设 $(S, \delta)$ 如情形 \ref{chow-situation-setup}，设 $X$、$Y$ 在 $S$ 上局部有限型，
并设 $f : X \to Y$ 是态射。假设 $f$ 固有。
\begin{enumerate}
\item 设 $Z \subset X$ 是满足 $\dim_\delta(Z) = k$ 的整闭子概形。定义
$$
f_*[Z] =
\left\{
\begin{matrix}
0 & \text{若} & \dim_\delta(f(Z))< k, \\
\deg(Z/f(Z)) [f(Z)] & \text{若} & \dim_\delta(f(Z)) = k.
\end{matrix}
\right.
$$
这里把 $f(Z) \subset Y$ 看作整闭子概形。若
$\dim_\delta(f(Z)) = \dim_\delta(Z)$，则由引理 \ref{chow-lemma-equal-dimension}，
$Z$ 在 $f(Z)$ 上的次数有限。
\item 设 $\alpha = \sum n_Z [Z]$ 是 $X$ 上的 $k$-循环。定义 $\alpha$ 的{\it 推前}为和
$$
f_* \alpha = \sum n_Z f_*[Z]
$$
其中每个 $f_*[Z]$ 如上定义。由引理 \ref{chow-lemma-quasi-compact-locally-finite}，这个和局部有限。
\end{enumerate}
\end{definition}

\noindent
依定义，循环的固有推前
$$
f_* : Z_k(X) \longrightarrow Z_k(Y)
$$
是阿贝尔群同态。它把 $X \mapsto Z_k(X)$ 变成如下范畴上的协变函子：
对象是在 $S$ 上局部有限型的概形，态射为固有态射。

\begin{lemma}
\label{chow-lemma-compose-pushforward}
设 $(S, \delta)$ 如情形 \ref{chow-situation-setup}，设 $X$、$Y$、$Z$ 在 $S$ 上局部有限型，
并设 $f : X \to Y$、$g : Y \to Z$ 为固有态射。那么作为映射
$Z_k(X) \to Z_k(Z)$，有 $g_* \circ f_* = (g \circ f)_*$。
\end{lemma}

\begin{proof}
设 $W \subset X$ 是维数为 $k$ 的整闭子概形。考虑
$W' = f(W) \subset Y$ 与 $W'' = g(f(W)) \subset Z$。
由于 $f$、$g$ 固有，$W'$（以及\ $W''$）分别是 $Y$（以及\ $Z$）的整闭子概形。
需要证明 $g_*(f_*[W]) = (g \circ f)_*[W]$。若 $\dim_\delta(W'') < k$，
则两边均为零。若 $\dim_\delta(W'') = k$，则所诱导的两个态射
$$
W \longrightarrow
W' \longrightarrow
W''
$$
都满足引理 \ref{chow-lemma-equal-dimension} 的假设。因此
$$
g_*(f_*[W]) = \deg(W/W')\deg(W'/W'')[W''],
\quad
(g \circ f)_*[W] = \deg(W/W'')[W''].
$$
再应用《概形态射》引理 \ref{morphisms-lemma-degree-composition} 即得结论。
\end{proof}

\noindent
闭浸入是固有态射。如果 $i : Z \to X$ 是闭浸入，那么映射
$$
i_* : Z_k(Z) \longrightarrow Z_k(X)
$$
全都是{\it 单射}。

\begin{lemma}
\label{chow-lemma-exact-sequence-closed}
设 $(S, \delta)$ 如情形 \ref{chow-situation-setup}。设 $X$ 在 $S$ 上局部有限型，并设
$X_1, X_2 \subset X$ 是闭子概形，使得在集合意义下 $X = X_1 \cup X_2$。
对每个 $k \in \mathbf{Z}$，阿贝尔群序列
$$
\xymatrix{
Z_k(X_1 \cap X_2) \ar[r] &
Z_k(X_1) \oplus Z_k(X_2) \ar[r] &
Z_k(X) \ar[r] &
0
}
$$
正合。这里 $X_1 \cap X_2$ 是概形论交，映射都是推前映射，其中一个乘以 $-1$。
\end{lemma}

\begin{proof}
先假设 $X$ 拟紧。那么 $Z_k(X)$ 是自由 $\mathbf{Z}$-模；它的一组基为元素 $[Z]$，
其中 $Z \subset X$ 是 $\delta$-维数为 $k$ 的整闭子概形。群
$Z_k(X_1)$、$Z_k(X_2)$、$Z_k(X_1 \cap X_2)$ 分别由这些 $Z$ 中满足
$Z \subset X_1$、$Z \subset X_2$、$Z \subset X_1 \cap X_2$ 的子集自由生成。
这立即证明了此情形下的引理。一般情形相似，略去证明。
\end{proof}

\begin{lemma}
\label{chow-lemma-cycle-push-sheaf}
设 $(S, \delta)$ 如情形 \ref{chow-situation-setup}，并设 $f : X \to Y$ 是在 $S$ 上局部有限型的概形之间的固有态射。
\begin{enumerate}
\item 设 $Z \subset X$ 是满足 $\dim_\delta(Z) \leq k$ 的闭子概形。那么
$$
f_*[Z]_k = [f_*{\mathcal O}_Z]_k.
$$
\item 设 $\mathcal{F}$ 是 $X$ 上的凝聚层，且
$\dim_\delta(\text{Supp}(\mathcal{F})) \leq k$。那么
$$
f_*[\mathcal{F}]_k = [f_*{\mathcal F}]_k.
$$
\end{enumerate}
该陈述有意义，因为由《概形上同调》命题 \ref{coherent-proposition-proper-pushforward-coherent}，$f_*\mathcal{F}$ 和
$f_*\mathcal{O}_Z$ 都是凝聚 $\mathcal{O}_Y$-模。
\end{lemma}

\begin{proof}
(1) 由 (2) 和引理 \ref{chow-lemma-cycle-closed-coherent} 推出。
设 $\mathcal{F}$ 是 $X$ 上的凝聚层，并假设
$\dim_\delta(\text{Supp}(\mathcal{F})) \leq k$。
由《概形上同调》引理 \ref{coherent-lemma-coherent-support-closed}，
存在闭子概形 $i : Z \to X$ 和凝聚 $\mathcal{O}_Z$-模 $\mathcal{G}$，使得
$i_*\mathcal{G} \cong \mathcal{F}$，并且 $\mathcal{F}$ 的支撑为 $Z$。
设 $Z' \subset Y$ 是 $f|_Z : Z \to Y$ 的概形论像。考虑概形的交换图
$$
\xymatrix{
Z \ar[r]_i \ar[d]_{f|_Z} &
X \ar[d]^f \\
Z' \ar[r]^{i'} & Y
}
$$
沿该图的两条路径可得
$f_*\mathcal{F} = f_*i_*\mathcal{G} = i'_*(f|_Z)_*\mathcal{G}$。
假设已知结果对闭浸入以及 $f|_Z$ 成立，则
$$
f_*[\mathcal{F}]_k = f_*i_*[\mathcal{G}]_k
= (i')_*(f|_Z)_*[\mathcal{G}]_k =
(i')_*[(f|_Z)_*\mathcal{G}]_k =
[(i')_*(f|_Z)_*\mathcal{G}]_k = [f_*\mathcal{F}]_k
$$
如所需。闭浸入的情形直接可得，略去证明。注意
$f|_Z : Z \to Z'$ 是支配态射；见《概形态射》引理 \ref{morphisms-lemma-quasi-compact-scheme-theoretic-image}。
因此已经约化到 $\dim_\delta(X) \leq k$ 且 $f : X \to Y$ 固有并支配的情形。

\medskip\noindent
假设 $\dim_\delta(X) \leq k$，且 $f : X \to Y$ 固有并支配。
由于 $f$ 支配，对一般点为 $\eta$ 的每个不可约分支 $Z \subset Y$，
都存在点 $\xi \in X$ 使 $f(\xi) = \eta$。因此
$\delta(\eta) \leq \delta(\xi) \leq k$。所以在表达式
$$
f_*[\mathcal{F}]_k = \sum n_Z[Z],
\quad
\text{且}
\quad
[f_*\mathcal{F}]_k = \sum m_Z[Z].
$$
中，只要 $n_Z \not = 0$ 或 $m_Z \not = 0$，整闭子概形 $Z$ 实际上就是 $Y$ 中
$\delta$-维数为 $k$ 的一个不可约分支。选取这样的整闭子概形 $Z \subset Y$，
并以 $\eta$ 表示其一般点。注意，对任何满足 $f(\xi) = \eta$ 的 $\xi \in X$，
都有 $\delta(\xi) \geq k$，因而 $\xi$ 也是 $X$ 中某个 $\delta$-维数为 $k$ 的不可约分支的一般点
（见引理 \ref{chow-lemma-multiplicity-finite}）。由于 $f$ 拟紧且 $X$ 局部诺特，
这样的点只有有限多个，故 $f^{-1}(\{\eta\})$ 有限。
由《概形态射》引理 \ref{morphisms-lemma-generically-finite}，存在开邻域
$\eta \in V \subset Y$，使 $f^{-1}(V) \to V$ 有限。以 $V$ 代替 $Y$，
以 $f^{-1}(V)$ 代替 $X$，便约化到 $Y$ 仿射且 $f$ 有限的情形。

\medskip\noindent
写成 $Y = \Spec(R)$、$X = \Spec(A)$；这是可能的，因为有限态射是仿射的。
那么 $R$、$A$ 是诺特环，且 $A$ 在 $R$ 上有限。此外，对某个有限 $A$-模 $M$，有
$\mathcal{F} = \widetilde{M}$。注意，$f_*\mathcal{F}$ 对应于被视作 $R$-模的 $M$。
设 $\mathfrak p \subset R$ 是对应于 $\eta \in Y$ 的极小素理想。
$Z$ 在 $[f_*\mathcal{F}]_k$ 中的系数显然为
$\text{length}_{R_{\mathfrak p}}(M_{\mathfrak p})$。
设 $\mathfrak q_i$（$i = 1, \ldots, t$）是 $A$ 中位于 $\mathfrak p$ 上方的素理想。
由于 $A_{\mathfrak p}$ 在维数为零的诺特局部环 $R_{\mathfrak p}$ 上有限，
它是 Artin 环，因而 $A_{\mathfrak p} = \prod A_{\mathfrak q_i}$。
$Z$ 在 $f_*[\mathcal{F}]_k$ 中的系数显然为
$$
\sum\nolimits_{i = 1, \ldots, t}
[\kappa(\mathfrak q_i) : \kappa(\mathfrak p)]
\text{length}_{A_{\mathfrak q_i}}(M_{\mathfrak q_i})
$$
所需等式于是来自《交换代数》引理 \ref{algebra-lemma-pushdown-module}。
\end{proof}




















\section{平坦拉回的准备工作}
\label{chow-section-preparation-flat-pullback}


\noindent
回忆，若局部有限型态射 $f : X \to Y$ 的每个非空纤维都是维数为 $r$ 的等维概形，
就称它的相对维数为 $r$。见《概形态射》定义 \ref{morphisms-definition-relative-dimension-d}。

\begin{lemma}
\label{chow-lemma-flat-inverse-image-dimension}
设 $(S, \delta)$ 如情形 \ref{chow-situation-setup}，设 $X$、$Y$ 在 $S$ 上局部有限型，
并设 $f : X \to Y$ 是态射。假设 $f$ 是相对维数为 $r$ 的平坦态射。
对任何闭子集 $Z \subset Y$，只要 $f^{-1}(Z)$ 非空，就有
$$
\dim_\delta(f^{-1}(Z)) = \dim_\delta(Z) + r.
$$
若 $Z$ 不可约，而 $Z' \subset f^{-1}(Z)$ 是一个不可约分支，
则 $Z'$ 支配 $Z$，并且 $\dim_\delta(Z') = \dim_\delta(Z) + r$。
\end{lemma}

\begin{proof}
只需证明最后一个断言。可以用整闭子概形 $Z$ 代替 $Y$，
并用概形论逆像 $f^{-1}(Z) = Z \times_Y X$ 代替 $X$。
因此可以假设 $Z = Y$ 为整概形，且 $f$ 是相对维数为 $r$ 的平坦态射。
由于 $Y$ 局部诺特，局部有限型态射 $f$ 实际上局部有限呈示。
故《概形态射》引理 \ref{morphisms-lemma-fppf-open} 适用，并表明 $f$ 是开态射。
设 $\xi \in X$ 是 $X$ 的某个不可约分支的一般点。由 $f$ 的开性，
$f(\xi)$ 是 $Z = Y$ 的一般点 $\eta$。因为 $f$ 的相对维数为 $r$，有
$\dim_\xi(X_\eta) = r$。另一方面，由于 $\xi$ 是 $X$ 的一般点，
$\mathcal{O}_{X, \xi} = \mathcal{O}_{X_\eta, \xi}$ 只有一个素理想，因而维数为 $0$。
所以由《概形态射》引理 \ref{morphisms-lemma-dimension-fibre-at-a-point}，
$\kappa(\xi)$ 在 $\kappa(\eta)$ 上的超越次数为 $r$。
换言之，$\delta(\xi) = \delta(\eta) + r$，如所需。
\end{proof}

\noindent
下面的引理将用于证明：闭子概形的局部有限族经平坦拉回后仍局部有限。

\begin{lemma}
\label{chow-lemma-inverse-image-locally-finite}
设 $(S, \delta)$ 如情形 \ref{chow-situation-setup}，设 $X$、$Y$ 在 $S$ 上局部有限型，
并设 $f : X \to Y$ 是态射。假设 $\{Z_i\}_{i \in I}$ 是 $Y$ 中闭子集的局部有限族，
那么 $\{f^{-1}(Z_i)\}_{i \in I}$ 是 $X$ 中闭子集的局部有限族。
\end{lemma}

\begin{proof}
设 $U \subset X$ 是拟紧开子集。由于像 $f(U) \subset Y$ 是拟紧子集，
存在拟紧开集 $V \subset Y$ 使 $f(U) \subset V$。注意
$$
\{i \in I \mid f^{-1}(Z_i) \cap U \not = \emptyset \}
\subset
\{i \in I \mid Z_i \cap V \not = \emptyset \}.
$$
右边依假设有限，故结论成立。
\end{proof}



\section{平坦拉回}
\label{chow-section-flat-pullback}

\noindent
下文以 $f^{-1}(Z)$ 表示如下概形论逆像：对概形态射 $f : X \to Y$ 与闭子概形
$Z \subset Y$，取该闭子概形的{\it 概形论逆像}。回忆，概形论逆像是纤维积
$$
\xymatrix{
f^{-1}(Z) \ar[r] \ar[d] & X \ar[d] \\
Z \ar[r] & Y
}
$$
如果 $\mathcal{I} \subset \mathcal{O}_Y$ 是对应于 $Y$ 中 $Z$ 的拟凝聚理想层，
它也就是由拟凝聚理想层 $f^{-1}(\mathcal{I})\mathcal{O}_X$ 截出的 $X$ 的闭子概形。
对此的讨论见《概形》第 \ref{schemes-section-closed-immersion} 节、
引理 \ref{schemes-lemma-fibre-product-immersion} 和定义 \ref{schemes-definition-inverse-image-closed-subscheme}。

\begin{definition}
\label{chow-definition-flat-pullback}
设 $(S, \delta)$ 如情形 \ref{chow-situation-setup}，设 $X$、$Y$ 在 $S$ 上局部有限型，
并设 $f : X \to Y$ 是态射。假设 $f$ 是相对维数为 $r$ 的平坦态射。
\begin{enumerate}
\item 设 $Z \subset Y$ 是 $\delta$-维数为 $k$ 的整闭子概形。定义 $f^*[Z]$ 为
$X$ 上与概形论逆像相伴的 $(k+r)$-循环
$$
f^*[Z] = [f^{-1}(Z)]_{k+r}.
$$
这是有意义的，因为由引理 \ref{chow-lemma-flat-inverse-image-dimension}，
$\dim_\delta(f^{-1}(Z)) = k + r$。
\item 设 $\alpha = \sum n_i [Z_i]$ 是 $Y$ 上的 $k$-循环。定义
{\it $\alpha$ 沿 $f$ 的平坦拉回}为和
$$
f^* \alpha = \sum n_i f^*[Z_i]
$$
其中每个 $f^*[Z_i]$ 如上定义。由引理 \ref{chow-lemma-inverse-image-locally-finite}，
该和局部有限。
\item 将如此得到的阿贝尔群映射记为 $f^* : Z_k(Y) \to Z_{k + r}(X)$。
\end{enumerate}
\end{definition}

\noindent
开浸入是平坦的。这是平坦态射的一个重要但平凡的特殊情形。如果 $U \subset X$ 是开集，
有时把循环沿 $j : U \to X$ 的拉回称为该循环在 $U$ 上的{\it 限制}。注意，此时映射
$$
j^* : Z_k(X) \longrightarrow Z_k(U)
$$
全都是{\it 满射}。理由如下：给定任意整闭子概形 $Z' \subset U$，
可以在 $X$ 中取 $Z'$ 的闭包 $Z$，并把它看成 $X$ 的约化闭子概形
（见《概形》引理 \ref{schemes-lemma-reduced-closed-subscheme}）。
显然 $Z \cap U = Z'$，换言之 $j^*[Z] = [Z']$，故得满射性。事实上还有稍强的结论。

\begin{lemma}
\label{chow-lemma-exact-sequence-open}
设 $(S, \delta)$ 如情形 \ref{chow-situation-setup}，设 $X$ 在 $S$ 上局部有限型。
设 $U \subset X$ 是开子概形，并把 $i : Y = X \setminus U \to X$ 视作 $X$ 的约化闭子概形。
对每个 $k \in \mathbf{Z}$，阿贝尔群复形
$$
\xymatrix{
Z_k(Y) \ar[r]^{i_*} & Z_k(X) \ar[r]^{j^*} & Z_k(U) \ar[r] & 0
}
$$
正合。
\end{lemma}

\begin{proof}
先假设 $X$ 拟紧。那么 $Z_k(X)$ 是自由 $\mathbf{Z}$-模；它的一组基由元素 $[Z]$ 给出，
其中 $Z \subset X$ 是 $\delta$-维数为 $k$ 的整闭子概形。
这样的基元素映到基元素 $[Z \cap U]$，或者在 $Z \subset Y$ 时映为零。
故此情形下引理显然。一般情形相似，略去证明。
\end{proof}

\begin{lemma}
\label{chow-lemma-compose-flat-pullback}
设 $(S, \delta)$ 如情形 \ref{chow-situation-setup}，设 $X, Y, Z$ 在 $S$ 上局部有限型。
设 $f : X \to Y$、$g : Y \to Z$ 分别是相对维数为 $r$、$s$ 的平坦态射。
那么 $g \circ f$ 是相对维数为 $r + s$ 的平坦态射，并且
$$
f^* \circ g^* = (g \circ f)^*
$$
作为映射 $Z_k(Z) \to Z_{k + r + s}(X)$ 成立。
\end{lemma}

\begin{proof}
由《概形态射》引理 \ref{morphisms-lemma-composition-relative-dimension-d}，
复合是相对维数为 $r + s$ 的平坦态射。假设
\begin{enumerate}
\item $W \subset Z$ 是 $\delta$-维数为 $k$ 的闭整子概形；
\item $W' \subset Y$ 是 $\delta$-维数为 $k + s$ 的闭整子概形，且
$W' \subset g^{-1}(W)$；以及
\item $W'' \subset Y$ 是 $\delta$-维数为 $k + s + r$ 的闭整子概形，且
$W'' \subset f^{-1}(W')$。
\end{enumerate}
需要证明，$[W'']$ 在 $(g \circ f)^*[W]$ 中的系数 $n$，等于
$[W'']$ 在 $f^*(g^*[W])$ 中的系数 $m$。由引理 \ref{chow-lemma-flat-inverse-image-dimension}，只核查这些情形已经足够。
设 $\xi'' \in W''$、$\xi' \in W'$、$\xi \in W$ 分别为一般点。考虑局部环
$A = \mathcal{O}_{Z, \xi}$、$B = \mathcal{O}_{Y, \xi'}$ 和
$C = \mathcal{O}_{X, \xi''}$。于是有平坦局部环同态 $A \to B$、$B \to C$，并且
$$
n = \text{length}_C(C/\mathfrak m_AC),
\quad
\text{且}
\quad
m = \text{length}_C(C/\mathfrak m_BC) \text{length}_B(B/\mathfrak m_AB)
$$
所需等式由《交换代数》引理 \ref{algebra-lemma-pullback-transitive} 推出。
\end{proof}

\begin{lemma}
\label{chow-lemma-pullback-coherent}
设 $(S, \delta)$ 如情形 \ref{chow-situation-setup}，设 $X, Y$ 在 $S$ 上局部有限型，
并设 $f : X \to Y$ 是相对维数为 $r$ 的平坦态射。
\begin{enumerate}
\item 设 $Z \subset Y$ 是满足 $\dim_\delta(Z) \leq k$ 的闭子概形。那么
$\dim_\delta(f^{-1}(Z)) \leq k + r$，并且在 $Z_{k + r}(X)$ 中有
$[f^{-1}(Z)]_{k + r} = f^*[Z]_k$。
\item 设 $\mathcal{F}$ 是 $Y$ 上满足
$\dim_\delta(\text{Supp}(\mathcal{F})) \leq k$ 的凝聚层。那么
$\dim_\delta(\text{Supp}(f^*\mathcal{F})) \leq k + r$，并且
$$
f^*[{\mathcal F}]_k = [f^*{\mathcal F}]_{k+r}
$$
在 $Z_{k + r}(X)$ 中成立。
\end{enumerate}
\end{lemma}

\begin{proof}
关于维数的断言立即来自引理 \ref{chow-lemma-flat-inverse-image-dimension}。
(1) 由 (2)、引理 \ref{chow-lemma-cycle-closed-coherent} 以及等式
$f^*\mathcal{O}_Z = \mathcal{O}_{f^{-1}(Z)}$ 推出。

\medskip\noindent
下面证明 (2)。由于 $X$、$Y$ 局部诺特，可以应用《概形上同调》引理 \ref{coherent-lemma-coherent-Noetherian}，得知 $\mathcal{F}$ 有限型；所以
$f^*\mathcal{F}$ 有限型（《模层》引理 \ref{modules-lemma-pullback-finite-type}），
进而 $f^*\mathcal{F}$ 凝聚（再次应用《概形上同调》引理 \ref{coherent-lemma-coherent-Noetherian}）。因此该引理有意义。
设 $W \subset Y$ 是 $\delta$-维数为 $k$ 的整闭子概形，并设 $W' \subset X$ 是维数为
$k + r$、经 $f$ 映入 $W$ 的整闭子概形。需要证明，$[W']$ 在
$f^*[{\mathcal F}]_k$ 中的系数 $n$ 等于 $[W']$ 在 $[f^*{\mathcal F}]_{k+r}$ 中的系数 $m$。
设 $\xi \in W$、$\xi' \in W'$ 为一般点。令
$A = \mathcal{O}_{Y, \xi}$、$B = \mathcal{O}_{X, \xi'}$，并把
$M = \mathcal{F}_\xi$ 看作 $A$-模。注意，由维数假设，$M$ 具有有限长度，
但实际上无需验证这一点；见引理 \ref{chow-lemma-length-finite}。
有 $f^*\mathcal{F}_{\xi'} = B \otimes_A M$，从而
$$
n = \text{length}_B(B \otimes_A M)
\quad
\text{且}
\quad
m = \text{length}_A(M) \text{length}_B(B/\mathfrak m_AB)
$$
所需等式来自《交换代数》引理 \ref{algebra-lemma-pullback-module}。
\end{proof}



\section{推前与拉回}
\label{chow-section-push-pull}

\noindent
本节验证，在有意义时，固有推前与平坦拉回彼此相容。由上面的工作，这归结为上同调与基变换。

\begin{lemma}
\label{chow-lemma-flat-pullback-proper-pushforward}
设 $(S, \delta)$ 如情形 \ref{chow-situation-setup}。设
$$
\xymatrix{
X' \ar[r]_{g'} \ar[d]_{f'} & X \ar[d]^f \\
Y' \ar[r]^g & Y
}
$$
是在 $S$ 上局部有限型的概形的纤维积图。假设 $f : X \to Y$ 固有，
而 $g : Y' \to Y$ 是相对维数为 $r$ 的平坦态射。那么 $f'$ 也固有，
$g'$ 也是相对维数为 $r$ 的平坦态射。对 $X$ 上任意 $k$-循环 $\alpha$，有
$$
g^*f_*\alpha = f'_*(g')^*\alpha
$$
在 $Z_{k + r}(Y')$ 中成立。
\end{lemma}

\begin{proof}
$f'$ 固有的断言来自《概形态射》引理 \ref{morphisms-lemma-base-change-proper}。
$g'$ 平坦且相对维数为 $r$ 的断言来自《概形态射》引理 \ref{morphisms-lemma-base-change-relative-dimension-d} 和
\ref{morphisms-lemma-base-change-flat}。只须在 $\alpha = [W]$ 时证明循环等式，
其中 $W \subset X$ 是 $\delta$-维数为 $k$ 的整闭子概形。
注意，此时有 $\alpha = [\mathcal{O}_W]_k$；见引理 \ref{chow-lemma-cycle-closed-coherent}。所以由引理 \ref{chow-lemma-cycle-push-sheaf} 和 \ref{chow-lemma-pullback-coherent}，
只须证明 $f'_*(g')^*\mathcal{O}_W$ 同构于 $g^*f_*\mathcal{O}_W$。
这来自上同调与基变换；见《概形上同调》引理 \ref{coherent-lemma-flat-base-change-cohomology}。
\end{proof}

\begin{lemma}
\label{chow-lemma-finite-flat}
设 $(S, \delta)$ 如情形 \ref{chow-situation-setup}，设 $X$、$Y$ 在 $S$ 上局部有限型。
设 $f : X \to Y$ 是次数为 $d$ 的有限局部自由态射
（见《概形态射》定义 \ref{morphisms-definition-finite-locally-free}）。
那么 $f$ 既固有，又是相对维数为 $0$ 的平坦态射，并且对每个
$\alpha \in Z_k(Y)$，有
$$
f_*f^*\alpha = d\alpha
$$
\end{lemma}

\begin{proof}
由《概形态射》引理 \ref{morphisms-lemma-finite-flat}，有限局部自由态射平坦且有限；
由《概形态射》引理 \ref{morphisms-lemma-finite-proper}，有限态射固有。
我们略去有限态射的相对维数为 $0$ 的证明。因此该公式有意义。
为证明它，设 $Z \subset Y$ 是 $\delta$-维数为 $k$ 的整闭子概形。
只须对 $\alpha = [Z]$ 证明公式。有限局部自由态射的基变换仍有限局部自由
（《概形态射》引理 \ref{morphisms-lemma-base-change-finite-locally-free}），
故 $f_*f^*\mathcal{O}_Z$ 是 $Z$ 上秩为 $d$ 的有限局部自由层。因此
$$
f_*f^*[Z] = f_*f^*[\mathcal{O}_Z]_k =
[f_*f^*\mathcal{O}_Z]_k = d[Z]
$$
这里使用了引理 \ref{chow-lemma-pullback-coherent} 和
\ref{chow-lemma-cycle-push-sheaf}。
\end{proof}








\section{主除子的准备}
\label{chow-section-preparation-principal-divisors}

\noindent
本节的一部分材料与《除子》第 \ref{divisors-section-Weil-divisors} 节
中的讨论有所重叠。

\begin{lemma}
\label{chow-lemma-divisor-delta-dimension}
设 $(S, \delta)$ 如情形 \ref{chow-situation-setup}。
设 $X$ 在 $S$ 上局部有限型，并假设 $X$ 是整概形。
\begin{enumerate}
\item 若 $Z \subset X$ 是整闭子概形，则下列条件等价：
\begin{enumerate}
\item $Z$ 是素除子；
\item $Z$ 在 $X$ 中的余维数为 $1$；
\item $\dim_\delta(Z) = \dim_\delta(X) - 1$。
\end{enumerate}
\item 若 $Z$ 是 $X$ 上某个有效 Cartier 除子的不可约分支，则
$\dim_\delta(Z) = \dim_\delta(X) - 1$。
\end{enumerate}
\end{lemma}

\begin{proof}
(1) 由素除子的定义（《除子》定义 \ref{divisors-definition-Weil-divisor}）和维数函数的定义（《拓扑学》定义 \ref{topology-definition-dimension-function}）立即得到。
设 $\xi \in Z$ 是有效 Cartier 除子 $D \subset X$ 的某个不可约分支 $Z$ 的一般点。
那么 $\dim(\mathcal{O}_{D, \xi}) = 0$，并且对某个非零因子
$f \in \mathcal{O}_{X, \xi}$，有
$\mathcal{O}_{D, \xi} = \mathcal{O}_{X, \xi}/(f)$（《除子》引理 \ref{divisors-lemma-effective-Cartier-in-points}）。
由《交换代数》引理 \ref{algebra-lemma-one-equation}，
$\dim(\mathcal{O}_{X, \xi}) = 1$。于是《概形的性质》引理 \ref{properties-lemma-codimension-local-ring} 表明 $Z$ 满足 (1)，证明完毕。
\end{proof}

\begin{lemma}
\label{chow-lemma-finite-in-codimension-one}
设 $f : X \to Y$ 是概形态射，$\xi \in Y$ 是一点。假设
\begin{enumerate}
\item $X$、$Y$ 都是整概形；
\item $Y$ 局部诺特；
\item $f$ 固有且支配，并且 $R(Y) \subset R(X)$ 是有限扩张；
\item $\dim(\mathcal{O}_{Y, \xi}) = 1$。
\end{enumerate}
那么存在 $\xi$ 的开邻域 $V \subset Y$，使得
$f|_{f^{-1}(V)} : f^{-1}(V) \to V$ 是有限态射。
\end{lemma}

\begin{proof}
本引理是《代数簇》引理 \ref{varieties-lemma-finite-in-codim-1} 的特例。
这里给出本情形下的直接论证。由《概形上同调》引理 \ref{coherent-lemma-proper-finite-fibre-finite-in-neighbourhood}，
只需证明 $f^{-1}(\{\xi\})$ 是有限集。用一个仿射开集替换 $Y$，记为
$Y = \Spec(R)$。由于假设 $Y$ 局部诺特，$R$ 是诺特环。
由于 $f$ 固有，它是拟紧的；因而可以取有限仿射开覆盖
$X = U_1 \cup \ldots \cup U_n$，其中每个 $U_i = \Spec(A_i)$。
注意，$R \to A_i$ 是整环之间有限型的单同态，且其诱导的分式域扩张有限。
因此该引理由《交换代数》引理 \ref{algebra-lemma-finite-in-codim-1} 得出。
\end{proof}


\section{主除子}
\label{chow-section-principal-divisors}

\noindent
下面的定义是《除子》定义 \ref{divisors-definition-principal-divisor} 在当前设定下的对应形式。

\begin{definition}
\label{chow-definition-principal-divisor}
设 $(S, \delta)$ 如情形 \ref{chow-situation-setup}。
设 $X$ 在 $S$ 上局部有限型，并假设 $X$ 是整概形且
$\dim_\delta(X) = n$。设 $f \in R(X)^*$。定义{\it 与 $f$ 相伴的主除子}为
$(n - 1)$-循环
$$
\text{div}(f) = \text{div}_X(f) = \sum \text{ord}_Z(f) [Z]
$$
其中右端定义见《除子》定义 \ref{divisors-definition-principal-divisor}。由引理 \ref{chow-lemma-divisor-delta-dimension}，素除子的 $\delta$-维数为 $n - 1$，
故此定义有意义。
\end{definition}

\noindent
在该定义的情形下，对 $f, g \in R(X)^*$，在 $Z_{n - 1}(X)$ 中有
$$
\text{div}_X(fg) = \text{div}_X(f) + \text{div}_X(g)
$$
见《除子》引理 \ref{divisors-lemma-div-additive}。
下面的引理将被更一般的引理 \ref{chow-lemma-flat-pullback-rational-equivalence} 所涵盖。

\begin{lemma}
\label{chow-lemma-flat-pullback-principal-divisor}
设 $(S, \delta)$ 如情形 \ref{chow-situation-setup}。
设 $X$、$Y$ 在 $S$ 上局部有限型，并假设 $X$、$Y$ 都是整概形，且
$n = \dim_\delta(Y)$。设 $f : X \to Y$ 是相对维数为 $r$ 的平坦态射，
$g \in R(Y)^*$。那么在 $Z_{n + r - 1}(X)$ 中有
$$
f^*(\text{div}_Y(g)) = \text{div}_X(g)
$$
\end{lemma}

\begin{proof}
由于 $f$ 平坦，它是支配态射，故 $f$ 诱导嵌入
$R(Y) \subset R(X)$；因而可以把 $g$ 看成 $R(X)^*$ 的元素。
设 $Z \subset X$ 是 $\delta$-维数为 $n + r - 1$ 的整闭子概形，
$\xi \in Z$ 是其一般点。若 $\dim_\delta(f(Z)) > n - 1$，
则方程两端 $[Z]$ 的系数都为零。因此可以假设
$Z' = \overline{f(Z)}$ 是 $Y$ 中 $\delta$-维数为 $n - 1$ 的整闭子概形。
令 $\xi' = f(\xi)$；它是 $Z'$ 的一般点。置
$A = \mathcal{O}_{Y, \xi'}$、$B = \mathcal{O}_{X, \xi}$。
环同态 $A \to B$ 是维数为 $1$ 的诺特局部整环之间的平坦局部同态。
$g$ 属于 $A$ 的分式域。所需证明的是
$$
\text{ord}_A(g) \text{length}_B(B/\mathfrak m_AB)
=
\text{ord}_B(g).
$$
这由《交换代数》引理 \ref{algebra-lemma-pullback-module} 得出
（细节从略）。
\end{proof}











\section{主除子与推前}
\label{chow-section-two-fun}

\noindent
第一个引理表明，沿一般有限态射推前主除子，所得仍是主除子。

\begin{lemma}
\label{chow-lemma-proper-pushforward-alteration}
设 $(S, \delta)$ 如情形 \ref{chow-situation-setup}。
设 $X$、$Y$ 在 $S$ 上局部有限型，并假设 $X$、$Y$ 都是整概形，且
$n = \dim_\delta(X) = \dim_\delta(Y)$。设 $p : X \to Y$ 是支配的固有态射，
$f \in R(X)^*$。置
$$
g = \text{Nm}_{R(X)/R(Y)}(f).
$$
那么
$p_*\text{div}(f) = \text{div}(g)$。
\end{lemma}

\begin{proof}
设 $Z \subset Y$ 是 $\delta$-维数为 $n - 1$ 的整闭子概形。
需要证明 $[Z]$ 在 $p_*\text{div}(f)$ 与 $\text{div}(g)$ 中的系数相等。
可以把引理 \ref{chow-lemma-finite-in-codimension-one} 用于态射
$p : X \to Y$ 和一般点 $\xi \in Z$。因此可以用 $\xi$ 的仿射开邻域替换 $Y$，
并假设 $p : X \to Y$ 是有限态射。写成 $Y = \Spec(R)$、$X = \Spec(A)$，
其中 $p$ 由诺特整环间的有限同态 $R \to A$ 诱导；该同态诱导分式域的有限扩张
$L/K$。现在有 $f \in L$、$g = \text{Nm}(f) \in K$，以及满足
$\dim(R_{\mathfrak p}) = 1$ 的素理想 $\mathfrak p \subset R$。
$[Z]$ 在 $\text{div}_Y(g)$ 中的系数为 $\text{ord}_{R_\mathfrak p}(g)$，
而 $[Z]$ 在 $p_*\text{div}_X(f)$ 中的系数为
$$
\sum\nolimits_{\mathfrak q\text{ 位于 }\mathfrak p\text{ 之上}}
[\kappa(\mathfrak q) : \kappa(\mathfrak p)]
\text{ord}_{A_{\mathfrak q}}(f)
$$
所需等式于是由《交换代数》引理 \ref{algebra-lemma-finite-extension-dim-1} 得出。
\end{proof}

\noindent
在主除子的讨论中，$\mathbf{P}^1_S$ 上的“泛”主除子
$[0] - [\infty]$ 起着重要作用。为更精确地表述这一点，以
\begin{equation}
\label{chow-equation-zero-infty}
D_0, D_\infty \subset
\mathbf{P}^1_S = \underline{\text{Proj}}_S(\mathcal{O}_S[T_0, T_1])
\end{equation}
表示分别由 $\mathcal{O}(1)$ 的截面 $T_1$ 与\ $T_0$ 截出的闭子概形。
它们是有效 Cartier 除子；见《除子》定义 \ref{divisors-definition-effective-Cartier-divisor} 和引理 \ref{divisors-lemma-characterize-OD}。下面的引理说，粗略而言，
“$\text{div}(T_1/T_0) = [D_0] - [D_1]$”，并且这是泛主除子。

\begin{lemma}
\label{chow-lemma-rational-function}
设 $(S, \delta)$ 如情形 \ref{chow-situation-setup}。
设 $X$ 在 $S$ 上局部有限型，并假设 $X$ 是整概形，且
$n = \dim_\delta(X)$。设 $f \in R(X)^*$。取非空开集 $U \subset X$，
使 $f$ 对应于一个截面 $f \in \Gamma(U, \mathcal{O}_X^*)$。
令 $Y \subset X \times_S \mathbf{P}^1_S$ 为
$f : U \to \mathbf{P}^1_S$ 的图的闭包。那么
\begin{enumerate}
\item 投影态射 $p : Y \to X$ 是固有的；
\item $p|_{p^{-1}(U)} : p^{-1}(U) \to U$ 是同构；
\item 沿态射 $q : Y \to \mathbf{P}^1_S$ 的拉回
$Y_0 = q^{-1}D_0$ 和 $Y_\infty = q^{-1}D_\infty$ 有定义
（《除子》定义 \ref{divisors-definition-pullback-effective-Cartier-divisor}）；
\item 有
$$
\text{div}_Y(f) = [Y_0]_{n - 1} - [Y_\infty]_{n - 1}
$$
\item 有
$$
\text{div}_X(f) = p_*\text{div}_Y(f)
$$
\item 若通过态射 $p$ 把 $Y_0$ 和 $Y_\infty$ 看成 $X$ 的闭子概形，则有
$$
\text{div}_X(f) = [Y_0]_{n - 1} - [Y_\infty]_{n - 1}
$$
\end{enumerate}
\end{lemma}

\begin{proof}
由于 $X$ 是整概形，$U$ 也是整概形。因此 $Y$ 是整概形，且
$(1, f)(U) \subset Y$ 是稠密开子概形。还要注意，闭子概形
$Y \subset X \times_S \mathbf{P}^1_S$ 与开集 $U$ 的选择无关，
因为它终究是单点集 $\{\eta'\} = \{(1, f)(\eta)\}$ 的闭包，
其中 $\eta \in X$ 是一般点。下面证明引理的各个断言。

\medskip\noindent
对于 (1)，注意 $p$ 是闭浸入
$Y \to X \times_S \mathbf{P}^1_S = \mathbf{P}^1_X$ 与固有态射
$\mathbf{P}^1_X \to X$ 的复合。固有态射的复合仍固有
（《概形态射》引理 \ref{morphisms-lemma-composition-proper}），故结论成立。

\medskip\noindent
显然 $Y \cap U \times_S \mathbf{P}^1_S = (1, f)(U)$，所以 (2) 成立；
同时也得到 $\dim_\delta(Y) = n$。

\medskip\noindent
由于 $f \in R(X)^*$，$q(\eta') = f(\eta)$ 不在 $D_0$ 或 $D_\infty$ 中。
因此 (3) 由《除子》引理 \ref{divisors-lemma-pullback-effective-Cartier-defined} 得到。
再由引理 \ref{chow-lemma-divisor-delta-dimension} 得到
$\dim_\delta(Y_0) = n - 1$ 和 $\dim_\delta(Y_\infty) = n - 1$。

\medskip\noindent
考虑有效 Cartier 除子 $Y_0$。在每一点 $\xi \in Y_0$，都有
$f \in \mathcal{O}_{Y, \xi}$，且 $Y_0$ 的局部方程由 $f$ 给出。
特别地，若 $\delta(\xi) = n - 1$，从而 $\xi$ 是某个
$\delta$-维数为 $n - 1$ 的整闭子概形 $Z$ 的一般点，则 $[Z]$ 在
$\text{div}_Y(f)$ 中的系数为
$$
\text{ord}_Z(f) =
\text{length}_{\mathcal{O}_{Y, \xi}}
(\mathcal{O}_{Y, \xi}/f\mathcal{O}_{Y, \xi}) =
\text{length}_{\mathcal{O}_{Y, \xi}}
(\mathcal{O}_{Y_0, \xi})
$$
这正是 $[Z]$ 在 $[Y_0]_{n - 1}$ 中的系数。对有理函数 $1/f$ 作类似论证，
可知 $-[Y_\infty]$ 对应于 $\text{div}_Y(f)$ 表达式中系数为负的各项。
故 (4) 成立。

\medskip\noindent
注意 $D_0 \to S$ 是同构，因而 $X \times_S D_0 \to X$ 也是同构。
显然，在 $X \times_S \mathbf{P}^1_S$ 中有
$Y_0 = Y \cap X \times_S D_0$（概形论交），故 $Y_0 \to X$ 确实是闭浸入。
于是
$$
p_*\mathcal{O}_{Y_0} = \mathcal{O}_{Y'_0}
$$
其中 $Y'_0 \subset X$ 是 $Y_0 \to X$ 的像。由引理 \ref{chow-lemma-cycle-push-sheaf}，有
$p_*[Y_0]_{n - 1} = [Y'_0]_{n - 1}$。对 $D_\infty$ 和 $Y_\infty$ 同样成立，
故 (6) 是 (5) 的推论。最后，(5) 立即由引理 \ref{chow-lemma-proper-pushforward-alteration} 得到。
\end{proof}

\noindent
下面的引理表明，固有曲线上主除子的次数为零。

\begin{lemma}
\label{chow-lemma-curve-principal-divisor}
设 $K$ 是任意域。设 $X$ 是一个 $1$-维整概形，并带有固有态射
$c : X \to \Spec(K)$。设 $f \in K(X)^*$ 是可逆有理函数。那么
$$
\sum\nolimits_{x \in X \text{ 为闭点}}
[\kappa(x) : K] \text{ord}_{\mathcal{O}_{X, x}}(f)
=
0
$$
其中 $\text{ord}$ 如《交换代数》定义 \ref{algebra-definition-ord} 所定义。
换言之，$c_*\text{div}(f) = 0$。
\end{lemma}

\begin{proof}
考虑图表
$$
\xymatrix{
Y \ar[r]_p \ar[d]_q & X \ar[d]^c \\
\mathbf{P}^1_K \ar[r]^-{c'} & \Spec(K)
}
$$
这是在引理 \ref{chow-lemma-rational-function} 中，从 $S = \Spec(K)$ 上的
$X$ 和有理函数 $f$ 出发构造的。下文不再说明而直接使用该引理的全部结论。
需要证明 $c_*\text{div}_X(f) = c_*p_*\text{div}_Y(f) = 0$，
这等价于证明 $c'_*q_*\text{div}_Y(f) = 0$。
若 $q(Y)$ 是 $\mathbf{P}^1_K$ 的闭点，则 $\text{div}_X(f) = 0$，引理成立。
所以可以假设 $q$ 是支配态射。假设能证明
$q : Y \to \mathbf{P}^1_K$ 是次数为 $d$ 的有限局部自由态射
（见《概形态射》定义 \ref{morphisms-definition-finite-locally-free}）。
由于 $\text{div}_Y(f) = [q^{-1}D_0]_0 - [q^{-1}D_\infty]_0$，
由平坦拉回的定义可知
$\text{div}_Y(f) = q^*([D_0]_0 - [D_\infty]_0)$。
再由引理 \ref{chow-lemma-finite-flat} 得到
$q_*\text{div}_Y(f) = d([D_0]_0 - [D_\infty]_0)$。
而显然 $c'_*[D_0]_0 = c'_*[D_\infty]_0$，结论即得。

\medskip\noindent
还须证明 $q$ 是有限局部自由态射。（由于 $\mathbf{P}^1_K$ 连通，
它会自动具有某个固定次数。）因为 $\dim(\mathbf{P}^1_K) = 1$，
例如由引理 \ref{chow-lemma-finite-in-codimension-one} 可知 $q$ 是有限态射。
$\mathbf{P}^1_K$ 在各闭点处的局部环都是维数为 $1$ 的正则局部环
（即离散赋值环），因为它们是 $K[T]$ 的局部化；见《交换代数》引理 \ref{algebra-lemma-dim-affine-space}。因此，对闭点 $y\in Y$，只要
$\mathcal{O}_{Y, y}$ 无挠，它在 $\mathcal{O}_{\mathbf{P}^1_K, q(y)}$ 上就是平坦的
（《代数进阶》引理 \ref{more-algebra-lemma-dedekind-torsion-free-flat}）。
这显然成立，因为 $\mathcal{O}_{Y, y}$ 是整环且 $q$ 支配。
所以 $q$ 平坦。最后由《概形态射》引理 \ref{morphisms-lemma-finite-flat}，$q$ 是有限局部自由态射。
\end{proof}





\section{有理等价}
\label{chow-section-rational-equivalence}

\noindent
本节定义 $k$-循环上的{\it 有理等价}。我们将允许主除子的像
（沿闭浸入）的局部有限和。这会产生一些颇为奇特的现象；见例 \ref{chow-example-weird}。然而，若不允许这种和，我们便不知道如何证明线丛陈类的帽积
可下降到有理等价类。

\begin{definition}
\label{chow-definition-rational-equivalence}
设 $(S, \delta)$ 如情形 \ref{chow-situation-setup}。
设概形 $X$ 在 $S$ 上局部有限型，并设 $k \in \mathbf{Z}$。
\begin{enumerate}
\item 给定任意局部有限族 $\{W_j \subset X\}$，其中各成员是满足
$\dim_\delta(W_j) = k + 1$ 的整闭子概形，并给定任意
$f_j \in R(W_j)^*$，可以考虑
$$
\sum (i_j)_*\text{div}(f_j) \in Z_k(X)
$$
其中 $i_j : W_j \to X$ 是包含态射。该表达式有意义，因为态射
$\coprod i_j : \coprod W_j \to X$ 是固有的。
\item 若 $\alpha$ 是上述形式的循环，就称 $\alpha \in Z_k(X)$
{\it 有理等价于零}。
\item 若 $\alpha - \beta$ 有理等价于零，就称
$\alpha, \beta \in Z_k(X)$ {\it 有理等价}，并记作
$\alpha \sim_{rat} \beta$。
\item 定义
$$
\CH_k(X) = Z_k(X) / \sim_{rat}
$$
为{\it $X$ 上 $k$-循环的 Chow 群}。它有时也称为
{\it $X$ 上模有理等价的 $k$-循环 Chow 群}。
\end{enumerate}
\end{definition}

\noindent
还有许多其他有趣的（适当的）等价关系；有理等价是其中最粗的一个。

\begin{remark}
\label{chow-remark-chow-group-pointwise}
设 $(S, \delta)$ 如情形 \ref{chow-situation-setup}。
设 $X$ 在 $S$ 上局部有限型，并设 $k \in \mathbf{Z}$。下面说明存在表现
$$
\bigoplus\nolimits_{\delta(x) = k + 1}' K_1^M(\kappa(x))
\xrightarrow{\partial}
\bigoplus\nolimits_{\delta(x) = k}' K_0^M(\kappa(x)) \to
\CH_k(X) \to 0
$$
这里采用注 \ref{chow-remark-cycles-pointwise} 中引入的记号与约定，此外：
\begin{enumerate}
\item $K_1^M(\kappa(x)) = \kappa(x)^*$ 是点 $x \in X$ 的剩余域
$\kappa(x)$ 的 Milnor K-理论的次数 $1$ 部分；见注 \ref{chow-remark-gersten-complex-milnor}；
\item 微分 $\partial$ 定义如下：给定元素 $\xi = \sum_x f_x$，
以 $W_x = \overline{x}$ 表示 $X$ 中以 $x$ 为一般点的整闭子概形，并置
$$
\partial(\xi) = \sum (W_x \to X)_*\text{div}(f_x)
$$
其属于 $Z_k(X)$。这是有意义的，因为注 \ref{chow-remark-cycles-pointwise} 已说明该复形的第二项等于 $Z_k(X)$。
\end{enumerate}
与定义 \ref{chow-definition-rational-equivalence} 比较，立即可知这确实给出
$\CH_k(X)$ 的一个表现。
\end{remark}

\noindent
下面是一个非常简单但很重要的引理。

\begin{lemma}
\label{chow-lemma-restrict-to-open}
设 $(S, \delta)$ 如情形 \ref{chow-situation-setup}。
设概形 $X$ 在 $S$ 上局部有限型。设 $U \subset X$ 是开子概形，并把
$i : Y = X \setminus U \to X$ 记为 $X$ 的约化闭子概形。设
$k \in \mathbf{Z}$，并设 $\alpha, \beta \in Z_k(X)$。
若 $\alpha|_U \sim_{rat} \beta|_U$，则存在循环
$\gamma \in Z_k(Y)$，使得
$$
\alpha \sim_{rat} \beta + i_*\gamma.
$$
换言之，Abel 群序列
$$
\xymatrix{
\CH_k(Y) \ar[r]^{i_*} & \CH_k(X) \ar[r]^{j^*} & \CH_k(U) \ar[r] & 0
}
$$
正合。
\end{lemma}

\begin{proof}
设 $\{W_j\}_{j \in J}$ 是 $U$ 中 $\delta$-维数为 $k + 1$ 的整闭子概形的
局部有限族，并设 $f_j \in R(W_j)^*$ 满足定义中的等式
$(\alpha - \beta)|_U = \sum (i_j)_*\text{div}(f_j)$。
令 $W_j' \subset X$ 为 $W_j$ 的闭包。设 $V \subset X$ 是拟紧开集。
由于 $V$ 是诺特的，$V \cap U$ 也是 $U$ 中的拟紧开集。因此集合
$\{j \in J \mid W_j \cap V \not = \emptyset\}
= \{j \in J \mid W'_j \cap V \not = \emptyset\}$
是有限的，因为 $\{W_j\}$ 局部有限。换言之，$\{W'_j\}$ 也局部有限。
由于 $R(W_j) = R(W'_j)$，可知
$$
\alpha - \beta - \sum (i'_j)_*\text{div}(f_j)
$$
是支撑在 $Y$ 上的循环，引理即得；见引理 \ref{chow-lemma-exact-sequence-open}。
\end{proof}

\begin{lemma}
\label{chow-lemma-exact-sequence-closed-chow}
设 $(S, \delta)$ 如情形 \ref{chow-situation-setup}。
设 $X$ 在 $S$ 上局部有限型。设 $X_1, X_2 \subset X$ 为闭子概形，
并且集合论上 $X = X_1 \cup X_2$。对每个 $k \in \mathbf{Z}$，Abel 群序列
$$
\xymatrix{
\CH_k(X_1 \cap X_2) \ar[r] &
\CH_k(X_1) \oplus \CH_k(X_2) \ar[r] &
\CH_k(X) \ar[r] &
0
}
$$
正合。这里 $X_1 \cap X_2$ 是概形论交；各映射均为推前映射，其中一个乘以
$-1$。
\end{lemma}

\begin{proof}
由引理 \ref{chow-lemma-exact-sequence-closed}，箭头
$\CH_k(X_1) \oplus \CH_k(X_2) \to \CH_k(X)$ 是满射。
假设 $(\alpha_1, \alpha_2)$ 在此映射下为零。写成
$\alpha_1 = \sum n_{1, i}[W_{1, i}]$ 和
$\alpha_2 = \sum n_{2, i}[W_{2, i}]$。于是得到 $X$ 中
$\delta$-维数为 $k + 1$ 的整闭子概形的局部有限族
$\{W_j\}_{j \in J}$，以及 $f_j \in R(W_j)^*$，使得
$$
\sum n_{1, i}[W_{1, i}] + \sum n_{2, i}[W_{2, i}] = \sum (i_j)_*\text{div}(f_j)
$$
作为 $X$ 上的循环成立，其中 $i_j : W_j \to X$ 为包含态射。
选择不交并分解 $J = J_1 \amalg J_2$，使得当 $j \in J_1$ 时
$W_j \subset X_1$，当 $j \in J_2$ 时 $W_j \subset X_2$。
这是可能的，因为各 $W_j$ 都是整概形。于是上述等式可写成
$$
\sum n_{1, i}[W_{1, i}] - \sum\nolimits_{j \in J_1} (i_j)_*\text{div}(f_j) =
- \sum n_{2, i}[W_{2, i}] + \sum\nolimits_{j \in J_2} (i_j)_*\text{div}(f_j)
$$
因此这个表达式是 $X_1 \cap X_2$ 上的一个循环（！）。换言之，
元素 $(\alpha_1, \alpha_2)$ 属于第一个箭头的像，证明完毕。
\end{proof}

\begin{example}
\label{chow-example-weird}
下面是一个“奇特”的例子。设 $S$ 是域 $k$ 的谱，$\delta$ 如例 \ref{chow-example-field}。设 $X = C_1 \cup C_2 \cup \ldots$ 是曲线
$C_j \cong \mathbf{P}^1_k$ 的无穷并，并按如下方式粘合：对
$j = 1, 2, 3, \ldots$，把点 $\infty \in C_j$ 与点
$0 \in C_{j + 1}$ 横截地粘合。取点 $0 \in C_1$，它给出零维循环
$[0] \in Z_0(X)$。本情形的“奇特”之处在于，事实上
$[0] \sim_{rat} 0$！具体地，可以选择有理函数 $f_j \in R(C_j)$，
使其在 $0$ 处有一个一阶零点，在 $\infty$ 处有一个一阶极点，且没有其他零点或极点。
于是和 $\sum (i_j)_*\text{div}(f_j)$ 恰好就是 $0$-循环 $[0]$。
事实上，在此例中 $\CH_0(X) = 0$。如果你觉得这过于离奇，只要始终要求所用空间拟紧即可
（这样 $X$ 对你而言根本不会出现）。
\end{example}

\begin{remark}
\label{chow-remark-infinite-sums-rational-equivalences}
设 $(S, \delta)$ 如情形 \ref{chow-situation-setup}。
设概形 $X$ 在 $S$ 上局部有限型。假设有 $X$ 上 $k$-循环的无穷族
$\alpha_i, \beta_i \in Z_k(X)$，$i \in I$。假设 $\alpha_i$ 与
$\beta_i$ 的支撑构成 $X$ 的闭子集的局部有限族，从而
$\sum \alpha_i$ 和 $\sum \beta_i$ 都定义了循环；再假设对每个 $i$，
$\alpha_i \sim_{rat} \beta_i$。此时并不能直接断定
$\sum \alpha_i \sim_{rat} \sum \beta_i$。问题在于，这些有理等价可能由局部有限族
$\{W_{i, j}, f_{i, j} \in R(W_{i, j})^*\}_{j \in J_i}$ 给出，
但它们的并 $\{W_{i, j}\}_{i \in I, j\in J_i}$ 未必局部有限。

\medskip\noindent
在许多实际情形中，存在闭子集的局部有限族 $\{T_i\}_{i \in I}$，
使 $\alpha_i, \beta_i$ 都支撑在 $T_i$ 上，并且在 $\CH_k(T_i)$ 中
$\alpha_i = \beta_i$；换言之，各族
$\{W_{i, j}, f_{i, j} \in R(W_{i, j})^*\}_{j \in J_i}$
由子概形 $W_{i, j} \subset T_i$ 组成。在这种情况下，$X$ 上确实有
$\sum \alpha_i \sim_{rat} \sum \beta_i$，因为此时族
$\{W_{i, j}\}_{i \in I, j\in J_i}$ 自动局部有限。
\end{remark}






\section{有理等价与推前、拉回}
\label{chow-section-properties-rational-equivalence}

\noindent
本节证明平坦拉回与固有推前都保持有理等价。

\begin{lemma}
\label{chow-lemma-prepare-flat-pullback-rational-equivalence}
设 $(S, \delta)$ 如情形 \ref{chow-situation-setup}。
设概形 $X$、$Y$ 在 $S$ 上局部有限型。假设 $Y$ 是整概形且
$\dim_\delta(Y) = k$。设 $f : X \to Y$ 是相对维数为 $r$ 的平坦态射。
那么对 $g \in R(Y)^*$，有
$$
f^*\text{div}_Y(g) =
\sum n_j i_{j, *}\text{div}_{X_j}(g \circ f|_{X_j})
$$
作为 $X$ 上的 $(k + r - 1)$-循环成立；其中求和遍历 $X$ 的不可约分支
$X_j$，而 $n_j$ 是 $X_j$ 在 $X$ 中的重数。
\end{lemma}

\begin{proof}
设 $Z \subset X$ 是 $\delta$-维数为 $k + r - 1$ 的整闭子概形。
需要证明 $[Z]$ 在 $f^*\text{div}(g)$ 中的系数 $n$，等于 $[Z]$ 在
$\sum i_{j, *} \text{div}(g \circ f|_{X_j})$ 中的系数 $m$。
设 $Z'$ 是 $f(Z)$ 的闭包；它是 $Y$ 的整闭子概形。由引理 \ref{chow-lemma-flat-inverse-image-dimension}，有
$\dim_\delta(Z') \geq k - 1$。因此，要么 $Z' = Y$，要么 $Z'$ 是
$Y$ 上的素除子。若 $Z' = Y$，则系数 $n$ 和 $m$ 都为零：
对 $n$，这由 $f^*$ 的定义显然；对 $m$，这是因为
$g \circ f|_{X_j}$ 在 $X_j$ 中任何映到 $Y$ 一般点的点处都是单位。
因此可以假设 $Z' \subset Y$ 是素除子。

\medskip\noindent
下面把 $n$ 与 $m$ 相等这一断言转化为代数等式。具体地，令 $\xi' \in Z'$ 和
$\xi \in Z$ 为一般点，置 $A = \mathcal{O}_{Y, \xi'}$、
$B = \mathcal{O}_{X, \xi}$。注意 $A$、$B$ 都是诺特环，$A \to B$
是平坦局部同态，$A$ 是整环，而 $\mathfrak m_AB$ 是局部环 $B$ 的定义理想。
有理函数 $g$ 是 $A$ 的分式域 $Q(A)$ 的元素。按构造，与 $\xi$ 相交的闭子概形
$X_j$ 以 $1$ 对 $1$ 的方式对应于极小素理想
$$
\mathfrak q_1, \ldots, \mathfrak q_s \subset B
$$
相应的整数 $n_j$ 是长度
$$
n_i = \text{length}_{B_{\mathfrak q_i}}(B_{\mathfrak q_i})
$$
有理函数 $g \circ f|_{X_j}$ 对应于 $g \in Q(A)$ 的像
$g_i \in \kappa(\mathfrak q_i)^*$。综合起来，有
$$
n = \text{ord}_A(g) \text{length}_B(B/\mathfrak m_AB)
$$
并且
$$
m = \sum \text{ord}_{B/\mathfrak q_i}(g_i)
\text{length}_{B_{\mathfrak q_i}}(B_{\mathfrak q_i})
$$
将该有理函数写成 $g = x/y$，其中非零的 $x, y \in A$。于是只需证明
$$
\text{length}_A(A/(x)) \text{length}_B(B/\mathfrak m_AB) =
\text{length}_B(B/xB)
$$
（此等式使用《交换代数》引理 \ref{algebra-lemma-pullback-module}）等于
$$
\sum\nolimits_{i = 1, \ldots, s}
\text{length}_{B/\mathfrak q_i}(B/(x, \mathfrak q_i))
\text{length}_{B_{\mathfrak q_i}}(B_{\mathfrak q_i})
$$
并对 $y$ 证明类似等式。由于 $A \to B$ 平坦，$x$ 是 $B$ 中的非零因子。
故所需等式由引理 \ref{chow-lemma-additivity-divisors-restricted} 得出。
\end{proof}

\begin{lemma}
\label{chow-lemma-flat-pullback-rational-equivalence}
设 $(S, \delta)$ 如情形 \ref{chow-situation-setup}。
设概形 $X$、$Y$ 在 $S$ 上局部有限型。设 $f : X \to Y$ 是相对维数为
$r$ 的平坦态射。设 $\alpha \sim_{rat} \beta$ 是 $Y$ 上有理等价的
$k$-循环。那么，作为 $X$ 上的 $(k + r)$-循环，
$f^*\alpha \sim_{rat} f^*\beta$。
\end{lemma}

\begin{proof}
需要证明什么？假设给定一族闭浸入
$$
i_j : W_j \longrightarrow Y
$$
其中每个 $W_j$ 都是 $\delta$-维数为 $k + 1$ 的整概形，并给定有理函数
$g_j \in R(W_j)^*$。再假设族 $\{i_j(W_j)\}_{j \in J}$ 在 $Y$ 上局部有限。
那么需要证明
$$
f^*(\sum i_{j, *}\text{div}(g_j)) = \sum f^*i_{j, *}\text{div}(g_j)
$$
在 $X$ 上有理等价于零。由引理 \ref{chow-lemma-inverse-image-locally-finite}，$\{W_j\}$ 在 $X$ 中局部有限，
故右端的和有意义。

\medskip\noindent
考虑纤维积
$$
i'_j : W'_j = W_j \times_Y X \longrightarrow X.
$$
并以 $f_j : W'_j \to W_j$ 表示第一投影。由引理 \ref{chow-lemma-flat-pullback-proper-pushforward}，上述和可写为
$$
\sum i'_{j, *}(f_j^*\text{div}(g_j))
$$
由引理 \ref{chow-lemma-prepare-flat-pullback-rational-equivalence}，
每个 $f_j^*\text{div}(g_j)$ 在 $W'_j$ 上有理等价于零。
因此每个 $i'_{j, *}(f_j^*\text{div}(g_j))$ 都有理等价于零。
由注 \ref{chow-remark-infinite-sums-rational-equivalences} 中的讨论，
展示的和也有同样的性质。
\end{proof}

\begin{lemma}
\label{chow-lemma-proper-pushforward-rational-equivalence}
设 $(S, \delta)$ 如情形 \ref{chow-situation-setup}。
设概形 $X$、$Y$ 在 $S$ 上局部有限型，并设 $p : X \to Y$ 是固有态射。
假设 $\alpha, \beta \in Z_k(X)$ 有理等价。那么 $p_*\alpha$ 与
$p_*\beta$ 有理等价。
\end{lemma}

\begin{proof}
需要证明什么？假设给定一族闭浸入
$$
i_j : W_j \longrightarrow X
$$
其中每个 $W_j$ 都是 $\delta$-维数为 $k + 1$ 的整概形，并给定有理函数
$f_j \in R(W_j)^*$。再假设族 $\{i_j(W_j)\}_{j \in J}$ 在 $X$ 上局部有限。
那么需要证明
$$
p_*\left(\sum i_{j, *}\text{div}(f_j)\right)
$$
在 $X$ 上有理等价于零。

\medskip\noindent
注意该和等于
$$
\sum p_*i_{j, *}\text{div}(f_j).
$$
设 $W'_j \subset Y$ 是 $p \circ i_j$ 的像这一整闭子概形。
由引理 \ref{chow-lemma-quasi-compact-locally-finite}，族 $\{W'_j\}$ 在 $Y$
中局部有限。因此，只须对给定的 $j$ 证明：要么
$p_*i_{j, *}\text{div}(f_j) = 0$，要么它等于
$i'_{j, *}\text{div}(g_j)$，其中某个 $g_j \in R(W'_j)^*$。

\medskip\noindent
上述论证把问题约化到一个 $\delta$-维数为 $k + 1$ 的整闭子概形
$W \subset X$。设 $f \in R(W)^*$，并如上令 $W' = p(W)$。
得到态射的交换图
$$
\xymatrix{
W \ar[r]_i \ar[d]_{p'} & X \ar[d]^p \\
W' \ar[r]^{i'} & Y
}
$$
由引理 \ref{chow-lemma-compose-pushforward}，
$p_*i_*\text{div}(f) = i'_*(p')_*\text{div}(f)$。如上所述，
需要证明 $(p')_*\text{div}(f)$ 要么是 $W'$ 上某个有理函数的除子，要么为零。
分三种情形。

\medskip\noindent
情形 $\dim_\delta(W') < k$。此时自动有 $(p')_*\text{div}(f) = 0$，
无须再证。

\medskip\noindent
情形 $\dim_\delta(W') = k$。下面证明此时 $(p')_*\text{div}(f) = 0$。
设 $\eta \in W'$ 为一般点。注意 $c : W_\eta \to \Spec(K)$ 是域
$K = \kappa(\eta)$ 上的固有整曲线，其函数域 $K(W_\eta)$ 与 $R(W)$ 等同。
有图表
$$
\xymatrix{
W_\eta \ar[r] \ar[d]_c & W \ar[d]^{p'} \\
\Spec(K) \ar[r] & W'
}
$$
以 $f_\eta \in K(W_\eta)^*$ 表示对应于 $f \in R(W)^*$ 的有理函数。
此外，$W_\eta$ 的闭点 $\xi$ 按 $1 - 1$ 对应于其中满足
$p'(Z) = W'$ 的 $\delta$-维数为 $k$ 的闭整子概形
$Z = Z_\xi \subset W$。注意 $Z_\xi$ 在 $\text{div}(f)$ 中的重数等于
$\text{ord}_{\mathcal{O}_{W_\eta, \xi}}(f_\eta)$，这只是因为局部环
$\mathcal{O}_{W_\eta, \xi}$ 与 $\mathcal{O}_{W, \xi}$ 作为各自分式域的子环彼此等同。
因此 $[W']$ 在 $(p')_*\text{div}(f)$ 中的重数等于
$[\Spec(K)]$ 在 $c_*\text{div}(f_\eta)$ 中的重数。
由引理 \ref{chow-lemma-curve-principal-divisor}，该重数为零。

\medskip\noindent
情形 $\dim_\delta(W') = k + 1$。此时可应用引理 \ref{chow-lemma-proper-pushforward-alteration}，从而确有
$p'_*\text{div}(f) = \text{div}(g)$，其中某个 $g \in R(W')^*$，正合所需。
\end{proof}















\section{有理等价与射影直线}
\label{chow-section-different-rational-equivalence}

\noindent
设 $(S, \delta)$ 如情形 \ref{chow-situation-setup}，并设概形 $X$ 在 $S$
上局部有限型。给定任意闭子概形
$Z \subset X \times_S \mathbf{P}^1_S = X \times \mathbf{P}^1$，
分别令 $Z_0$、\ $Z_\infty$ 为概形论闭子概形
$Z_0 = \text{pr}_2^{-1}(D_0)$、\ $Z_\infty = \text{pr}_2^{-1}(D_\infty)$。
这里 $D_0$、$D_\infty$ 如 (\ref{chow-equation-zero-infty}) 中所定义。

\begin{lemma}
\label{chow-lemma-rational-equivalence-family}
设 $(S, \delta)$ 如情形 \ref{chow-situation-setup}。
设概形 $X$ 在 $S$ 上局部有限型。设
$W \subset X \times_S \mathbf{P}^1_S$ 是 $\delta$-维数为 $k + 1$ 的整闭子概形。
假设 $W \not = W_0$ 且 $W \not = W_\infty$。那么
\begin{enumerate}
\item $W_0$、$W_\infty$ 是 $W$ 的有效 Cartier 除子；
\item 可以把 $W_0$、$W_\infty$ 看成 $X$ 的闭子概形，并且
$$
[W_0]_k \sim_{rat} [W_\infty]_k,
$$
\item 对任意局部有限族的整闭子概形
$W_i \subset X \times_S \mathbf{P}^1_S$，若其 $\delta$-维数为 $k + 1$，
且 $W_i \not = (W_i)_0$、$W_i \not = (W_i)_\infty$，则在 $X$ 上有
$\sum ([(W_i)_0]_k - [(W_i)_\infty]_k) \sim_{rat} 0$；
\item 对任意满足 $\alpha \sim_{rat} 0$ 的 $\alpha \in Z_k(X)$，
存在如上的整闭子概形局部有限族
$W_i \subset X \times_S \mathbf{P}^1_S$，使得
$\alpha = \sum ([(W_i)_0]_k - [(W_i)_\infty]_k)$。
\end{enumerate}
\end{lemma}

\begin{proof}
由假设，$W$ 的一般点在投影
$X \times_S \mathbf{P}^1_S \to \mathbf{P}^1_S$ 下不映入
$D_0$ 或 $D_\infty$，所以 (1) 由《除子》引理 \ref{divisors-lemma-pullback-effective-Cartier-defined} 得出。

\medskip\noindent
由于 $X \times_S D_0 \to X$ 是闭浸入，$W_0$ 同构于 $X$ 的某个闭子概形；
$W_\infty$ 同理。态射 $p : W \to X$ 是闭浸入
$W \to X \times_S \mathbf{P}^1_S$ 与固有态射
$X \times_S \mathbf{P}^1_S \to X$ 的复合，故它固有。
由引理 \ref{chow-lemma-rational-function}，作为 $W$ 上的循环，有
$[W_0]_k \sim_{rat} [W_\infty]_k$。而显然
$p_*[W_0]_k = [W_0]_k$，对 $W_\infty$ 也一样，所以 (2) 由引理 \ref{chow-lemma-proper-pushforward-rational-equivalence} 得出。

\medskip\noindent
在已知 (1)、(2) 后，断言 (3) 唯一需要说明的是：族
$\{(W_i)_0, (W_i)_\infty\}$ 是 $X$ 的闭子概形的局部有限族。这是显然的。

\medskip\noindent
假设 $\alpha \sim_{rat} 0$。按定义，存在 $\delta$-维数为 $k + 1$ 的整闭子概形
$V_i \subset X$ 和有理函数 $f_i \in R(V_i)^*$，使得族
$\{V_i\}_{i \in I}$ 在 $X$ 中局部有限，并且
$\alpha = \sum (V_i \to X)_*\text{div}(f_i)$。令
$$
W_i \subset V_i \times_S \mathbf{P}^1_S \subset X \times_S \mathbf{P}^1_S
$$
为引理 \ref{chow-lemma-rational-function} 中有理映射 $f_i$ 的图的闭包。
由同一引理，$(V_i \to X)_*\text{div}(f_i)$ 等于
$[(W_i)_0]_k - [(W_i)_\infty]_k$，结论显然。
\end{proof}

\begin{lemma}
\label{chow-lemma-closed-subscheme-cross-p1}
设 $(S, \delta)$ 如情形 \ref{chow-situation-setup}。
设概形 $X$ 在 $S$ 上局部有限型，并设 $Z$ 是
$X \times \mathbf{P}^1$ 的闭子概形。假设
\begin{enumerate}
\item $\dim_\delta(Z) \leq k + 1$；
\item $\dim_\delta(Z_0) \leq k$、$\dim_\delta(Z_\infty) \leq k$；
\item 对 $Z$ 的任意嵌入点 $\xi$（《除子》定义 \ref{divisors-definition-embedded}），要么
$\xi \not \in Z_0 \cup Z_\infty$，要么 $\delta(\xi) < k$。
\end{enumerate}
那么，作为 $X$ 上的 $k$-循环，有
$[Z_0]_k \sim_{rat} [Z_\infty]_k$。
\end{lemma}

\begin{proof}
设 $\{W_i\}_{i \in I}$ 是 $Z$ 中 $\delta$-维数为 $k + 1$ 的不可约分支的集合。
按定义写成
$$
[Z]_{k + 1} = \sum n_i[W_i]
$$
其中 $n_i > 0$。由《除子》引理 \ref{divisors-lemma-components-locally-finite}，$\{W_i\}$ 是
$X \times_S \mathbf{P}^1_S$ 的闭子集的局部有限族。我们断言
$$
[Z_0]_k = \sum n_i[(W_i)_0]_k
$$
并且对 $[Z_\infty]_k$ 也有类似等式。一旦证明此点，引理即由引理 \ref{chow-lemma-rational-equivalence-family} 得出。

\medskip\noindent
设 $Z' \subset X$ 是 $\delta$-维数为 $k$ 的整闭子概形。
为证明上述等式，只须证明 $[Z']$ 在 $[Z_0]_k$ 中的系数 $n$，
等于 $[Z']$ 在 $\sum n_i[(W_i)_0]_k$ 中的系数 $m$。
设 $\xi' \in Z'$ 为一般点，并置
$\xi = (\xi', 0) \in  X \times_S \mathbf{P}^1_S$。
考虑局部环 $A = \mathcal{O}_{X \times_S \mathbf{P}^1_S, \xi}$。
设 $I \subset A$ 是截出 $Z$ 的理想，即 $A/I = \mathcal{O}_{Z, \xi}$。
设 $t \in A$ 是截出 $X \times_S D_0$ 的元素（即把 $\mathbf{P}^1$
在零点处的坐标拉回）。由 $\xi' \in Z'$ 的选择，
$\delta(\xi) = k$，故 $\dim(A/I) = 1$。由假设 (3)，$\xi$ 不是嵌入点，
所以 $A/I$ 是 Cohen--Macaulay 环。由于 $\dim_\delta(Z_0)
= k$，
有 $\dim(A/(t, I)) = 0$，这蕴含 $t$ 是 $A/I$ 上的非零因子。
最后，经过 $\xi$ 的不可约闭子概形 $W_i$ 对应于位于 $I$ 上方的极小素理想
$I \subset \mathfrak q_i$。重数 $n_i$ 对应于长度
$\text{length}_{A_{\mathfrak q_i}}(A/I)_{\mathfrak q_i}$。因此
$$
n = \text{length}_A(A/(t, I))
$$
并且
$$
m = \sum
\text{length}_A(A/(t, \mathfrak q_i))
\text{length}_{A_{\mathfrak q_i}}(A/I)_{\mathfrak q_i}
$$
结论于是由引理 \ref{chow-lemma-additivity-divisors-restricted} 得出。
\end{proof}

\begin{lemma}
\label{chow-lemma-coherent-sheaf-cross-p1}
设 $(S, \delta)$ 如情形 \ref{chow-situation-setup}。
设概形 $X$ 在 $S$ 上局部有限型。设 $\mathcal{F}$ 是
$X \times \mathbf{P}^1$ 上的凝聚层。设
$i_0, i_\infty : X \to X \times \mathbf{P}^1$ 是满足
$i_t(x) = (x, t)$ 的闭浸入。记
$\mathcal{F}_0 = i_0^*\mathcal{F}$、$\mathcal{F}_\infty = i_\infty^*\mathcal{F}$。
假设
\begin{enumerate}
\item $\dim_\delta(\text{Supp}(\mathcal{F})) \leq k + 1$；
\item $\dim_\delta(\text{Supp}(\mathcal{F}_0)) \leq k$、
$\dim_\delta(\text{Supp}(\mathcal{F}_\infty)) \leq k$；
\item 对 $\mathcal{F}$ 的任意嵌入伴随点 $\xi$，要么
$\xi \not \in (X \times \mathbf{P}^1)_0 \cup (X \times \mathbf{P}^1)_\infty$，
要么 $\delta(\xi) < k$。
\end{enumerate}
那么，作为 $X$ 上的 $k$-循环，有
$[\mathcal{F}_0]_k \sim_{rat} [\mathcal{F}_\infty]_k$。
\end{lemma}

\begin{proof}
设 $\{W_i\}_{i \in I}$ 是 $\text{Supp}(\mathcal{F})$ 中
$\delta$-维数为 $k + 1$ 的不可约分支的集合。按定义写成
$$
[\mathcal{F}]_{k + 1} = \sum n_i[W_i]
$$
其中 $n_i > 0$。由引理 \ref{chow-lemma-length-finite}，$\{W_i\}$ 是
$X \times_S \mathbf{P}^1_S$ 的闭子集的局部有限族。我们断言
$$
[\mathcal{F}_0]_k = \sum n_i[(W_i)_0]_k
$$
并且对 $[\mathcal{F}_\infty]_k$ 也有类似等式。一旦证明此点，引理即由引理 \ref{chow-lemma-rational-equivalence-family} 得出。

\medskip\noindent
设 $Z' \subset X$ 是 $\delta$-维数为 $k$ 的整闭子概形。
为证明上述等式，只须证明 $[Z']$ 在 $[\mathcal{F}_0]_k$ 中的系数 $n$，
等于 $[Z']$ 在 $\sum n_i[(W_i)_0]_k$ 中的系数 $m$。
设 $\xi' \in Z'$ 为一般点，并置
$\xi = (\xi', 0) \in  X \times_S \mathbf{P}^1_S$。
考虑局部环 $A = \mathcal{O}_{X \times_S \mathbf{P}^1_S, \xi}$，
并把 $M = \mathcal{F}_\xi$ 看成 $A$-模。设 $t \in A$ 是截出
$X \times_S D_0$ 的元素（即把 $\mathbf{P}^1$ 在零点处的坐标拉回）。
由 $\xi' \in Z'$ 的选择，$\delta(\xi) = k$，故
$\dim(\text{Supp}(M)) = 1$。由假设 (3)，$\xi$ 不是 $\mathcal{F}$ 的伴随点，
所以 $M$ 是 Cohen--Macaulay 模。由于
$\dim_\delta(\text{Supp}(\mathcal{F}_0)) = k$，有
$\dim(\text{Supp}(M/tM)) = 0$，这蕴含 $t$ 是 $M$ 上的非零因子。
最后，经过 $\xi$ 的不可约闭子概形 $W_i$ 对应于
$\text{Ass}(M)$ 的极小素理想 $\mathfrak q_i$。重数 $n_i$ 对应于长度
$\text{length}_{A_{\mathfrak q_i}}M_{\mathfrak q_i}$。因此
$$
n = \text{length}_A(M/tM)
$$
并且
$$
m = \sum
\text{length}_A(A/(t, \mathfrak q_i)A)
\text{length}_{A_{\mathfrak q_i}}M_{\mathfrak q_i}
$$
结论于是由引理 \ref{chow-lemma-additivity-divisors-restricted} 得出。
\end{proof}







\section{Chow 群与包络}
\label{chow-section-envelopes}

\noindent
定义如下。

\begin{definition}
\label{chow-definition-envelope}
\begin{reference}
\cite[Definition 18.3]{F}
\end{reference}
设 $X$ 为概形。所谓{\it 包络}，是指完全分解的固有态射 $f : Y \to X$
（《态射进阶》定义 \ref{more-morphisms-definition-cd-morphism}）。
\end{definition}

\noindent
引理 \ref{chow-lemma-envelope} 的正合序列是采用该定义的主要动机。

\begin{lemma}
\label{chow-lemma-composition-envelope}
设 $(S, \delta)$ 如情形 \ref{chow-situation-setup}。
设概形 $X$ 在 $S$ 上局部有限型。若 $f : Y \to X$ 和
$g : Z \to Y$ 都是包络，则 $f \circ g$ 也是包络。
\end{lemma}

\begin{proof}
这由《概形态射》引理 \ref{morphisms-lemma-composition-proper} 和
《态射进阶》引理 \ref{more-morphisms-lemma-composition-cd} 得出。
\end{proof}

\begin{lemma}
\label{chow-lemma-base-change-envelope}
设 $(S, \delta)$ 如情形 \ref{chow-situation-setup}。
设 $X' \to X$ 是在 $S$ 上局部有限型的概形之间的态射。
若 $f : Y \to X$ 是包络，则 $f$ 的基变换 $f' : Y' \to X'$ 也是包络。
\end{lemma}

\begin{proof}
这由《概形态射》引理 \ref{morphisms-lemma-base-change-proper} 和
《态射进阶》引理 \ref{more-morphisms-lemma-base-change-cd} 得出。
\end{proof}

\begin{lemma}
\label{chow-lemma-envelope}
设 $(S, \delta)$ 如情形 \ref{chow-situation-setup}。
设概形 $X$ 在 $S$ 上局部有限型，并设 $f : Y \to X$ 是包络。
那么对所有 $k \in \mathbf{Z}$，有正合序列
$$
\CH_k(Y \times_X Y) \xrightarrow{p_* - q_*}
\CH_k(Y) \xrightarrow{f_*}
\CH_k(X) \to 0
$$
其中 $p, q : Y \times_X Y \to Y$ 是投影。
\end{lemma}

\begin{proof}
由于 $f$ 是包络，$f$ 是固有态射，故循环及循环类的推前都有定义；见第 \ref{chow-section-proper-pushforward} 节和 \ref{chow-section-push-pull}。
类似地，态射 $p$ 和 $q$ 是 $f$ 的基变换，所以也固有。两个箭头的复合为零，因为
$f_* \circ p_* = (p \circ f)_* = (q \circ f)_* = f_* \circ q_*$；
见引理 \ref{chow-lemma-compose-pushforward}。

\medskip\noindent
先证明 $f_* : Z_k(Y) \to Z_k(X)$ 是满射。设
$\alpha = \sum n_i[Z_i] \in Z_k(X)$，其中 $Z_i \subset X$ 是整闭子概形的
局部有限族。令 $x_i \in Z_i$ 为一般点。由于 $f$ 是包络，因而完全分解，
存在点 $y_i \in Y$，满足 $f(y_i) = x_i$，且
$\kappa(y_i)/\kappa(x_i)$ 是平凡扩张。令 $W_i \subset Y$ 为以 $y_i$ 为一般点的
整闭子概形。由于 $f$ 是闭映射，$f(W_i) = Z_i$。由此，闭子概形族 $W_i$
在 $Y$ 上局部有限。由于 $\kappa(y_i)/\kappa(x_i)$ 平凡，有
$\dim_\delta(W_i) = \dim_\delta(Z_i) = k$。因此
$\beta = \sum n_i[W_i]$ 属于 $Z_k(Y)$。最后，因为
$\kappa(y_i)/\kappa(x_i)$ 平凡，支配态射
$f|_{W_i} : W_i \to Z_i$ 的次数为 $1$，故 $f_*\beta = \alpha$。

\medskip\noindent
既然 $f_* : Z_k(Y) \to Z_k(X)$ 是满射，映射
$f_* : \CH_k(Y) \to \CH_k(X)$ 自然更是满射。

\medskip\noindent
设 $\beta \in Z_k(Y)$，且 $f_*\beta$ 在 $\CH_k(X)$ 中为零。
这意味着可以找到整闭子概形的局部有限族 $Z_j \subset X$，满足
$\dim_\delta(Z_j) = k + 1$，以及 $f_j \in R(Z_j)^*$，使得作为循环有
$$
f_*\beta = \sum (Z_j \to X)_*\text{div}(f_j)
$$
其中 $i_j : Z_j \to X$ 是给定的闭浸入。完全照上述方式，可以找到整闭子概形的
局部有限族 $W_j \subset Y$，使得 $f(W_j) = Z_j$ 且
$W_j \to Z_j$ 是双有理态射，即诱导同构 $R(Z_j) = R(W_j)$。
以 $g_j \in R(W_j)^*$ 表示对应于 $f_j$ 的元素。注意 $W_j \to Z_j$ 固有，
且在 $Z_j$ 上作为循环有
$(W_j \to Z_j)_*\text{div}(g_j) = \text{div}(f_j)$。
由此，若把 $\beta$ 替换为与其有理等价的循环
$$
\beta' = \beta - \sum (W_j \to Y)_*\text{div}(g_j)
$$
则 $f_*\beta' = 0$。（这里使用引理 \ref{chow-lemma-compose-pushforward}。）
因此，要完成引理的证明，只须证明下一段中的断言。

\medskip\noindent
断言：若 $\beta \in Z_k(Y)$ 且 $f_*\beta = 0$，则对某个
$\gamma \in Z_k(Y \times_X Y)$，在 $Z_k(Y)$ 中有
$\beta = \delta + p_*\gamma - q_*\gamma$。具体地，写成
$\beta = \sum_{j \in J} n_j[W_j]$，其中 $\{W_j\}_{j \in J}$ 是 $Y$ 中
满足 $\dim_\delta(W_j) = k$ 的整闭子概形的局部有限族。
固定一个整闭子概形 $Z \subset X$，并考虑子集
$J_Z = \{j \in J : f(W_j) = Z\}$。这是有限集。分三种情形：
\begin{enumerate}
\item $J_Z = \emptyset$。此时置 $\gamma_Z = 0$。
\item $J_Z \not = \emptyset$ 且 $\dim_\delta(Z) = k$。
考察 $Z$ 的系数，条件 $f_*\beta = 0$ 蕴含
$\sum_{j \in J_Z} n_j\deg(W_j/Z) = 0$。此时选择双有理地映到 $Z$ 的整闭子概形
$W \subset Y$（见上文）。考察一般点可知，$W_j \times_Z W$ 有唯一一个
不可约分支 $W'_j \subset W_j \times_Z W \subset Y \times_X Y$，
它双有理地映到 $W_j$。于是 $W'_j \to W$ 支配，且
$\deg(W'_j/W) = \deg(W_j/W)$。因此，若置
$\gamma_Z = \sum_{j \in J_Z} n_j[W'_j]$，则
$p_*\gamma_Z = \sum_{j \in J_Z} n_j[W_j]$，并且
$q_*\gamma_Z = \sum_{j \in J_Z} n_j\deg(W'_j/W)[W] = 0$。
\item $J_Z \not = \emptyset$ 且 $\dim_\delta(Z) < k$。
此时选择双有理地映到 $Z$ 的整闭子概形 $W \subset Y$（见上文）。
考察一般点可知，$W_j \times_Z W$ 有唯一一个不可约分支
$W'_j \subset W_j \times_Z W \subset Y \times_X Y$，它双有理地映到 $W_j$。
于是 $W'_j \to W$ 支配，且
$k = \dim_\delta(W'_j) > \dim_\delta(W) = \dim_\delta(Z)$。
因此，若置 $\gamma_Z = \sum_{j \in J_Z} n_j[W'_j]$，则
$p_*\gamma_Z = \sum_{j \in J_Z} n_j[W_j]$，并且 $q_*\gamma_Z = 0$。
\end{enumerate}
由于整闭子概形族 $\{f(W_j)\}$ 在 $X$ 上局部有限
（引理 \ref{chow-lemma-quasi-compact-locally-finite}），$Y \times_X Y$ 上的
$k$-循环
$$
\gamma = \sum\nolimits_{Z \subset X\text{ 为闭整子概形}} \gamma_Z
$$
有良好定义。由上述计算，$p_*\gamma_Z = \beta$ 且 $q_*\gamma_Z = 0$，
这就蕴含所需结论。
\end{proof}













\section{Chow 群与 K-群}
\label{chow-section-chow-and-K}

\noindent
本节比较凝聚层范畴的 $K_0$ 与 Chow 群。

\medskip\noindent
设 $(S, \delta)$ 如情形 \ref{chow-situation-setup}。
设概形 $X$ 在 $S$ 上局部有限型。以
$\textit{Coh}(X) = \textit{Coh}(\mathcal{O}_X)$ 表示 $X$ 上凝聚层的范畴。
这是 Abel 范畴；见《概形上同调》引理 \ref{coherent-lemma-coherent-abelian-Noetherian}。对任意
$k \in \mathbf{Z}$，以 $\textit{Coh}_{\leq k}(X)$ 表示
$\textit{Coh}(X)$ 的全子范畴，其对象是满足
$\dim_\delta(\text{Supp}(\mathcal{F})) \leq k$ 的凝聚层 $\mathcal{F}$。

\begin{lemma}
\label{chow-lemma-Serre-subcategories}
设 $(S, \delta)$ 如情形 \ref{chow-situation-setup}。
设概形 $X$ 在 $S$ 上局部有限型。范畴
$\textit{Coh}_{\leq k}(X)$ 是 Abel 范畴 $\textit{Coh}(X)$ 的 Serre 子范畴。
\end{lemma}

\begin{proof}
Serre 子范畴的定义见《同调代数》定义 \ref{homology-definition-serre-subcategory}。该引理的证明直接，故略去。
\end{proof}

\begin{lemma}
\label{chow-lemma-cycles-k-group}
设 $(S, \delta)$ 如情形 \ref{chow-situation-setup}。
设概形 $X$ 在 $S$ 上局部有限型。映射
$$
Z_k(X)
\longrightarrow
K_0(\textit{Coh}_{\leq k}(X)/\textit{Coh}_{\leq k - 1}(X)),
\quad
\sum n_Z[Z] \mapsto
\left[\bigoplus\nolimits_{n_Z > 0} \mathcal{O}_Z^{\oplus n_Z}\right]
-
\left[\bigoplus\nolimits_{n_Z < 0} \mathcal{O}_Z^{\oplus -n_Z}\right]
$$
与映射
$$
K_0(\textit{Coh}_{\leq k}(X)/\textit{Coh}_{\leq k - 1}(X))
\longrightarrow
Z_k(X),\quad
\mathcal{F} \longmapsto [\mathcal{F}]_k
$$
是互逆同构。
\end{lemma}

\begin{proof}
注意，若 $\sum n_Z[Z]$ 属于 $Z_k(X)$，则直和
$\bigoplus\nolimits_{n_Z > 0} \mathcal{O}_Z^{\oplus n_Z}$ 与
$\bigoplus\nolimits_{n_Z < 0} \mathcal{O}_Z^{\oplus -n_Z}$
都是 $X$ 上的凝聚层，因为族 $\{Z \mid n_Z > 0\}$ 在 $X$ 上局部有限。
映射 $\mathcal{F} \to [\mathcal{F}]_k$ 在
$\textit{Coh}_{\leq k}(X)$ 上可加；见引理 \ref{chow-lemma-additivity-sheaf-cycle}。而若
$\mathcal{F} \in \textit{Coh}_{\leq k - 1}(X)$，则
$[\mathcal{F}]_k = 0$。由《同调代数》引理 \ref{homology-lemma-serre-subcategory-K-groups} 的 (1)，第二个映射也有良好定义。
显然，先应用第一个映射再应用第二个映射所得复合是恒等映射。

\medskip\noindent
反过来，设从 $X$ 上的凝聚层 $\mathcal{F}$ 出发。写成
$[\mathcal{F}]_k = \sum_{i \in I} n_i[Z_i]$，其中 $n_i > 0$，而
$Z_i \subset X$、$i \in I$ 是两两不同、$\delta$-维数为 $k$ 的整闭子概形。
需要证明
$$
[\mathcal{F}] = [\bigoplus\nolimits_{i \in I} \mathcal{O}_{Z_i}^{\oplus n_i}]
$$
在 $K_0(\textit{Coh}_{\leq k}(X)/\textit{Coh}_{\leq k - 1}(X))$ 中成立。
以 $\xi_i \in Z_i$ 表示一般点。若置
$$
\mathcal{F}' = \Ker(\mathcal{F} \to \bigoplus \xi_{i, *}\mathcal{F}_{\xi_i})
$$
则 $\mathcal{F}'$ 是 $\mathcal{F}$ 的、支撑维数 $\leq k - 1$ 的最大凝聚子模。
特别地，$\mathcal{F}$ 与 $\mathcal{F}/\mathcal{F}'$ 在
$K_0(\textit{Coh}_{\leq k}(X)/\textit{Coh}_{\leq k - 1}(X))$ 中具有同一类。
因此，以 $\mathcal{F}/\mathcal{F}'$ 替换 $\mathcal{F}$ 后，可以并且将假设
上面展示的核 $\mathcal{F}'$ 为零。

\medskip\noindent
对每个 $i \in I$，选择滤过
$$
\mathcal{F}_{\xi_i} = \mathcal{F}_i^0 \supset \mathcal{F}_i^1 \supset
\ldots \supset \mathcal{F}_i^{n_i} = 0
$$
使各相继商在 $\xi_i$ 处的剩余域上维数为 $1$。这是可能的，因为
$\mathcal{F}_{\xi_i}$ 作为 $\mathcal{O}_{X, \xi_i}$-模的长度为 $n_i$。
对 $p > n_i$，置 $\mathcal{F}_i^p = 0$。对 $p \geq 0$，记
$$
\mathcal{F}^p =
\Ker\left(\mathcal{F} \longrightarrow \bigoplus
\xi_{i, *}(\mathcal{F}_{\xi_i}/\mathcal{F}_i^p)\right)
$$
那么 $\mathcal{F}^p$ 凝聚，$\mathcal{F}^0 = \mathcal{F}$；并且在
$\xi_i$ 的某个开邻域内，$\mathcal{F}^p/\mathcal{F}^{p + 1}$ 同构于秩为
$1$ 的自由 $\mathcal{O}_{Z_i}$-模（若 $n_i > p$），或同构于 $0$
（若 $n_i \leq p$）。此外，$\mathcal{F}' = \bigcap \mathcal{F}^p = 0$。
由于每个拟紧开集 $U \subset X$ 只包含有限多个 $\xi_i$，可知当
$p \gg 0$ 时 $\mathcal{F}^p|_U$ 为零。因此
$\bigoplus_{p \geq 0} \mathcal{F}^p$ 是凝聚 $\mathcal{O}_X$-模。
考虑短正合列
$$
0 \to
\bigoplus\nolimits_{p > 0} \mathcal{F}^p \to
\bigoplus\nolimits_{p \geq 0} \mathcal{F}^p \to
\bigoplus\nolimits_{p > 0} \mathcal{F}^p/\mathcal{F}^{p + 1} \to 0
$$
和
$$
0 \to
\bigoplus\nolimits_{p > 0} \mathcal{F}^p \to
\bigoplus\nolimits_{p \geq 0} \mathcal{F}^p \to
\mathcal{F} \to 0
$$
它们都是凝聚 $\mathcal{O}_X$-模的短正合列。这已经表明
$$
[\mathcal{F}] = [\bigoplus \mathcal{F}^p/\mathcal{F}^{p + 1}]
$$
在 $K_0(\textit{Coh}_{\leq k}(X)/\textit{Coh}_{\leq k - 1}(X))$ 中成立。
接下来，对每个满足 $n_i > p$ 的 $p \geq 0$ 与 $i \in I$，
选择非零理想层 $\mathcal{I}_{i, p} \subset \mathcal{O}_{Z_i}$，并选择 $X$ 上的映射
$\mathcal{I}_{i, p} \to \mathcal{F}^p/\mathcal{F}^{p + 1}$，使其在上述
$\xi_i$ 的开邻域上为同构。《概形上同调》引理 \ref{coherent-lemma-extend-coherent} 保证可以这样选择。然后考虑短正合列
$$
0 \to
\bigoplus\nolimits_{p \geq 0, i \in I, n_i > p} \mathcal{I}_{i, p}
\to
\bigoplus \mathcal{F}^p/\mathcal{F}^{p + 1} \to
\mathcal{Q} \to 0
$$
以及短正合列
$$
0 \to
\bigoplus\nolimits_{p \geq 0, i \in I, n_i > p} \mathcal{I}_{i, p}
\to
\bigoplus\nolimits_{p \geq 0, i \in I, n_i > p} \mathcal{O}_{Z_i}
\to
\mathcal{Q}' \to 0
$$
注意 $\mathcal{Q}$ 和 $\mathcal{Q}'$ 在各点 $\xi_i$ 的某个邻域内都为零，
且都支撑在 $\bigcup Z_i$ 上。因此 $\mathcal{Q}$ 和 $\mathcal{Q}'$ 属于
$\textit{Coh}_{\leq k - 1}(X)$。由于
$$
\bigoplus\nolimits_{i \in I} \mathcal{O}_{Z_i}^{\oplus n_i} \cong
\bigoplus\nolimits_{p \geq 0, i \in I, n_i > p} \mathcal{O}_{Z_i}
$$
证明完成。
\end{proof}

\begin{lemma}
\label{chow-lemma-finite-cycles-k-group}
设 $\pi : X \to Y$ 是在情形 \ref{chow-situation-setup} 中
$(S, \delta)$ 上局部有限型的概形之间的有限态射。那么
$\pi_* : \textit{Coh}(X) \to \textit{Coh}(Y)$ 是正合函子；它把
$\textit{Coh}_{\leq k}(X)$ 送入 $\textit{Coh}_{\leq k}(Y)$，并在这些范畴及其商范畴的
$K_0$ 上诱导同态。引理 \ref{chow-lemma-cycles-k-group} 的映射构成交换图
$$
\xymatrix{
Z_k(X) \ar[d]^{\pi_*} \ar[r] &
K_0(\textit{Coh}_{\leq k}(X)/\textit{Coh}_{\leq k - 1}(X))
\ar[d]^{\pi_*} \ar[r] &
Z_k(X) \ar[d]^{\pi_*} \\
Z_k(Y) \ar[r] &
K_0(\textit{Coh}_{\leq k}(Y)/\textit{Coh}_{\leq k - 1}(Y)) \ar[r] &
Z_k(Y)
}
$$
\end{lemma}

\begin{proof}
有限态射是仿射态射，所以由《概形上同调》引理 \ref{coherent-lemma-relative-affine-vanishing}，沿 $\pi$ 推前拟凝聚模是正合函子。
有限态射是固有态射，所以 $\pi_*$ 把凝聚层送到凝聚层；见《概形上同调》命题 \ref{coherent-proposition-proper-pushforward-coherent}。
关于支撑维数的断言显然。右侧的交换性立即由引理 \ref{chow-lemma-cycle-push-sheaf} 得到。由于各水平箭头是双射，左侧也交换。
\end{proof}

\begin{lemma}
\label{chow-lemma-from-chow-to-K}
设 $X$ 是情形 \ref{chow-situation-setup} 中在 $(S, \delta)$ 上局部有限型的概形。
存在典范映射
$$
\CH_k(X)
\longrightarrow
K_0(\textit{Coh}_{\leq k + 1}(X)/\textit{Coh}_{\leq k - 1}(X))
$$
它由引理 \ref{chow-lemma-cycles-k-group} 中的映射
$Z_k(X) \to K_0(\textit{Coh}_{\leq k}(X)/\textit{Coh}_{\leq k - 1}(X))$
诱导。
\end{lemma}

\begin{proof}
需要证明：$Z_k(X)$ 中有理等价于零的元素 $\alpha$，在
$K_0(\textit{Coh}_{\leq k + 1}(X)/\textit{Coh}_{\leq k - 1}(X))$
中映到零。按照定义 \ref{chow-definition-rational-equivalence} 写成
$\alpha = \sum (i_j)_*\text{div}(f_j)$。注意
$$
\pi = \coprod i_j : W = \coprod W_j \longrightarrow X
$$
是有限态射，因为每个 $i_j : W_j \to X$ 都是闭浸入，且 $W_j$ 的族在
$X$ 中局部有限。因此可以使用引理 \ref{chow-lemma-finite-cycles-k-group}
把问题约化到 $W$ 的情形。由于 $W$ 是整概形的不交并，进一步约化到下一段讨论的情形。

\medskip\noindent
假设 $X$ 是 $\delta$-维数为 $k + 1$ 的整概形。设 $f$ 是 $X$ 上的非零有理函数，
并令 $\alpha = \text{div}(f)$。需要证明 $\alpha$ 在
$K_0(\textit{Coh}_{\leq k + 1}(X)/\textit{Coh}_{\leq k - 1}(X))$
中映到零。设 $\mathcal{I} \subset \mathcal{O}_X$ 是 $f$ 的分母理想；见《除子》定义 \ref{divisors-definition-regular-meromorphic-ideal-denominators}。
于是有短正合列
$$
0 \to \mathcal{I} \to \mathcal{O}_X \to \mathcal{O}_X/\mathcal{I} \to 0
$$
和
$$
0 \to \mathcal{I} \xrightarrow{f} \mathcal{O}_X \to
\mathcal{O}_X/f\mathcal{I} \to 0
$$
见《除子》引理 \ref{divisors-lemma-regular-meromorphic-ideal-denominators}。我们断言
$$
[\mathcal{O}_X/\mathcal{I}]_k - [\mathcal{O}_X/f\mathcal{I}]_k =
\text{div}(f)
$$
该断言表明，元素 $\alpha = \text{div}(f)$ 在
$K_0(\textit{Coh}_{\leq k}(X)/\textit{Coh}_{\leq k - 1}(X))$ 中由
$[\mathcal{O}_X/\mathcal{I}] - [\mathcal{O}_X/f\mathcal{I}]$ 表示。
上述短正合列于是表明，该元素在
$K_0(\textit{Coh}_{\leq k + 1}(X)/\textit{Coh}_{\leq k - 1}(X))$ 中映到零。

\medskip\noindent
为证明该断言，设 $Z \subset X$ 是 $\delta$-维数为 $k$ 的整闭子概形，
$\xi \in Z$ 是其一般点。那么
$I = \mathcal{I}_\xi \subset A = \mathcal{O}_{X, \xi}$ 是满足
$fI \subset A$ 的理想。$[Z]$ 在 $\text{div}(f)$ 中的系数是
$\text{ord}_A(f)$。（当然，与往常一样，把 $X$ 的函数域与 $A$ 的分式域等同。）
另一方面，$[Z]$ 在
$[\mathcal{O}_X/\mathcal{I}] - [\mathcal{O}_X/f\mathcal{I}]$ 中的系数为
$$
\text{length}_A(A/I) - \text{length}_A(A/fI)
$$
利用《交换代数》定义 \ref{algebra-definition-distance} 的距离函数，
可将其改写为
$$
d(A, I) - d(A, fI) = d(I, fI) = \text{ord}_A(f)
$$
这些等式由《交换代数》引理 \ref{algebra-lemma-properties-distance-function} 和
\ref{algebra-lemma-order-vanishing-determinant} 得出。
使用这些引理并非必要，但很方便。
\end{proof}

\begin{remark}
\label{chow-remark-good-cases-K-A}
设 $(S, \delta)$ 如情形 \ref{chow-situation-setup}。
设概形 $X$ 在 $S$ 上局部有限型。稍后将在引理 \ref{chow-lemma-cycles-rational-equivalence-K-group} 中看到，引理 \ref{chow-lemma-from-chow-to-K} 的映射
$$
\CH_k(X)
\longrightarrow
K_0(\textit{Coh}_{k + 1}(X)/\textit{Coh}_{\leq k - 1}(X))
$$
是单射。将它与典范映射
$$
K_0(\textit{Coh}_{k + 1}(X)/\textit{Coh}_{\leq k - 1}(X))
\longrightarrow
K_0(\textit{Coh}(X)/\textit{Coh}_{\leq k - 1}(X))
$$
复合，得到典范映射
$$
\CH_k(X)
\longrightarrow
K_0(\textit{Coh}(X)/\textit{Coh}_{\leq k - 1}(X)).
$$
关于这个映射是否应为单射，我们未能在文献中找到相关断言或猜想。
合理的预期似乎是其核为挠群。稍后将回到这个问题（插入未来引用）。
\end{remark}

\begin{lemma}
\label{chow-lemma-K-coherent-supported-on-closed}
设 $X$ 是局部诺特概形，$Z \subset X$ 是闭子概形。以
$\textit{Coh}_Z(X) \subset \textit{Coh}(X)$ 表示这样的凝聚
$\mathcal{O}_X$-模构成的 Serre 子范畴：其集合论支撑包含在 $Z$ 中。
那么正合包含函子 $\textit{Coh}(Z) \to \textit{Coh}_Z(X)$ 诱导同构
$$
K'_0(Z) = K_0(\textit{Coh}(Z)) \longrightarrow K_0(\textit{Coh}_Z(X))
$$
\end{lemma}

\begin{proof}
设 $\mathcal{F}$ 是 $\textit{Coh}_Z(X)$ 的对象。设
$\mathcal{I} \subset \mathcal{O}_X$ 是 $Z$ 的拟凝聚理想层。考虑降滤过
$$
\ldots \subset
\mathcal{F}^p = \mathcal{I}^p \mathcal{F} \subset
\mathcal{F}^{p - 1} \subset \ldots \subset \mathcal{F}^0 = \mathcal{F}
$$
完全如引理 \ref{chow-lemma-from-chow-to-K} 的证明，这个滤过局部有限，因而
$\bigoplus_{p \geq 0} \mathcal{F}^p$、
$\bigoplus_{p \geq 1} \mathcal{F}^p$ 与
$\bigoplus_{p \geq 0} \mathcal{F}^p/\mathcal{F}^{p + 1}$
都是支撑在 $Z$ 上的凝聚 $\mathcal{O}_X$-模。因此，与引理 \ref{chow-lemma-from-chow-to-K} 的证明完全一样，在 $K_0(\textit{Coh}_Z(X))$ 中有
$$
[\mathcal{F}] =
[\bigoplus\nolimits_{p \geq 0} \mathcal{F}^p/\mathcal{F}^{p + 1}]
$$
凝聚模 $\bigoplus_{p \geq 0} \mathcal{F}^p/\mathcal{F}^{p + 1}$ 被
$\mathcal{I}$ 零化，所以 $[\mathcal{F}]$ 属于该像。事实上，我们断言映射
$$
\mathcal{F} \longmapsto 
c(\mathcal{F}) =
[\bigoplus\nolimits_{p \geq 0} \mathcal{F}^p/\mathcal{F}^{p + 1}]
$$
可通过 $K_0(\textit{Coh}_Z(X))$ 分解，并且是引理所述映射的逆。
为此只须证明：若
$$
0 \to \mathcal{F} \to \mathcal{G} \to \mathcal{H} \to 0
$$
是 $\textit{Coh}_Z(X)$ 中的短正合列，则
$c(\mathcal{G}) = c(\mathcal{F}) + c(\mathcal{H})$。
注意，对所有 $q \geq 0$，有短正合列
$$
0 \to
(\mathcal{F} \cap \mathcal{I}^q\mathcal{G})/
(\mathcal{F} \cap \mathcal{I}^{q + 1}\mathcal{G}) \to
\mathcal{G}^q/\mathcal{G}^{q + 1} \to
\mathcal{H}^q/\mathcal{H}^{q + 1} \to 0
$$
对 $p, q \geq 0$，考虑凝聚子模
$$
\mathcal{F}^{p, q} = \mathcal{I}^p\mathcal{F} \cap \mathcal{I}^q\mathcal{G}
$$
完全照上述方式，并利用滤过 $\mathcal{F}^p = \mathcal{I}^p\mathcal{F}$ 与
$\mathcal{F} \cap \mathcal{I}^q\mathcal{G}$ 都局部有限，得到
$$
[\bigoplus\nolimits_{p \geq 0} \mathcal{F}^p/\mathcal{F}^{p + 1}] =
[\bigoplus\nolimits_{p, q \geq 0}
\mathcal{F}^{p, q}/(\mathcal{F}^{p + 1, q} + \mathcal{F}^{p, q + 1})] =
[\bigoplus\nolimits_{q \geq 0}
(\mathcal{F} \cap \mathcal{I}^q\mathcal{G})/
(\mathcal{F} \cap \mathcal{I}^{q + 1}\mathcal{G})]
$$
其位于 $K_0(\textit{Coh}(Z))$ 中。结合上述正合列即可得到所需结论。
若干细节从略。
\end{proof}













\section{与可逆层相伴的除子}
\label{chow-section-divisor-invertible-sheaf}

\noindent
下面的定义是《除子》定义 \ref{divisors-definition-divisor-invertible-sheaf} 在当前设定下的对应形式。

\begin{definition}
\label{chow-definition-divisor-invertible-sheaf}
设 $(S, \delta)$ 如情形 \ref{chow-situation-setup}。
设 $X$ 在 $S$ 上局部有限型，并假设 $X$ 是整概形且
$n = \dim_\delta(X)$。设 $\mathcal{L}$ 是可逆 $\mathcal{O}_X$-模。
\begin{enumerate}
\item 对 $\mathcal{L}$ 的任意非零亚纯截面 $s$，定义{\it 与 $s$ 相伴的 Weil 除子}为
$(n - 1)$-循环
$$
\text{div}_\mathcal{L}(s) =
\sum \text{ord}_{Z, \mathcal{L}}(s) [Z]
$$
其定义见《除子》定义 \ref{divisors-definition-divisor-invertible-sheaf}。由引理 \ref{chow-lemma-divisor-delta-dimension}，Weil 除子的 $\delta$-维数为 $n - 1$，
故此定义有意义。
\item 定义{\it 与 $\mathcal{L}$ 相伴的 Weil 除子}为
$$
c_1(\mathcal{L}) \cap [X] =
\text{div}_\mathcal{L}(s)\text{ 的类} \in \CH_{n - 1}(X)
$$
其中 $s$ 是 $X$ 上 $\mathcal{L}$ 的任意非零亚纯截面。由《除子》引理 \ref{divisors-lemma-divisor-meromorphic-well-defined}，此定义与选择无关。
\end{enumerate}
\end{definition}

\noindent
设 $X$、$S$ 如上面的定义 \ref{chow-definition-divisor-invertible-sheaf}，
并置 $n = \dim_\delta(X)$。由定义显然有
$Cl(X) = \CH_{n - 1}(X)$，其中 $Cl(X)$ 是《除子》定义 \ref{divisors-definition-class-group} 中定义的 $X$ 的 Weil 除子类群。映射
$$
\Pic(X) \longrightarrow \CH_{n - 1}(X), \quad
\mathcal{L} \longmapsto c_1(\mathcal{L}) \cap [X]
$$
与 Divisors, Equation (\ref{divisors-equation-c1}) 对任意局部诺特整概形构造的映射
$\Pic(X) \to Cl(X)$ 相同。特别地，该映射是 Abel 群同态；若 $X$ 正规，则它是单射；
若 $X$ 的所有局部环都是 UFD，则它是同构。见《除子》引理 \ref{divisors-lemma-normal-c1-injective} 和
\ref{divisors-lemma-local-rings-UFD-c1-bijective}。
在某些情形下，与可逆层相伴的 Weil 除子很容易计算。

\begin{lemma}
\label{chow-lemma-compute-c1}
设 $(S, \delta)$ 如情形 \ref{chow-situation-setup}。
设 $X$ 在 $S$ 上局部有限型，并假设 $X$ 是整概形且
$n = \dim_\delta(X)$。设 $\mathcal{L}$ 是可逆 $\mathcal{O}_X$-模，
$s \in \Gamma(X, \mathcal{L})$ 是非零整体截面。那么在 $Z_{n - 1}(X)$ 中
$$
\text{div}_\mathcal{L}(s) = [Z(s)]_{n - 1}
$$
并且在 $\CH_{n - 1}(X)$ 中
$$
c_1(\mathcal{L}) \cap [X] = [Z(s)]_{n - 1}
$$
\end{lemma}

\begin{proof}
设 $Z \subset X$ 是 $\delta$-维数为 $n - 1$ 的整闭子概形，
$\xi \in Z$ 是其一般点。选择生成元 $s_\xi \in \mathcal{L}_\xi$，并对某个
$f \in \mathcal{O}_{X, \xi}$ 写成 $s = fs_\xi$。由 $Z(s)$ 的定义
（见《除子》定义 \ref{divisors-definition-zero-scheme-s}），
$Z(s)$ 由拟凝聚理想层 $\mathcal{I} \subset \mathcal{O}_X$ 截出，且
$\mathcal{I}_\xi = (f)$。因此
$\text{length}_{\mathcal{O}_{X, \xi}}(\mathcal{O}_{Z(s), \xi})
=
\text{length}_{\mathcal{O}_{X, \xi}}(\mathcal{O}_{X, \xi}/(f))
=
\text{ord}_{\mathcal{O}_{X, \xi}}(f)$，正合所需。
\end{proof}

\noindent
下面的引理将被更一般的引理 \ref{chow-lemma-flat-pullback-cap-c1} 所涵盖。

\begin{lemma}
\label{chow-lemma-flat-pullback-divisor-invertible-sheaf}
设 $(S, \delta)$ 如情形 \ref{chow-situation-setup}。
设 $X$、$Y$ 在 $S$ 上局部有限型，并假设 $X$、$Y$ 都是整概形且
$n = \dim_\delta(Y)$。设 $\mathcal{L}$ 是可逆 $\mathcal{O}_Y$-模，
$f : X \to Y$ 是相对维数为 $r$ 的平坦态射。那么在
$\CH_{n + r - 1}(X)$ 中
$$
f^*(c_1(\mathcal{L}) \cap [Y]) = c_1(f^*\mathcal{L}) \cap [X]
$$
\end{lemma}

\begin{proof}
设 $s$ 是 $\mathcal{L}$ 的非零亚纯截面。我们将证明实际上有
$f^*\text{div}_\mathcal{L}(s) = \text{div}_{f^*\mathcal{L}}(f^*s)$，
从而引理成立。为此，设 $\xi \in Y$ 是一点，且
$s_\xi \in \mathcal{L}_\xi$ 是生成元。对某个 $g \in R(Y)^*$ 写成
$s = gs_\xi$。存在 $\xi$ 的开邻域 $V \subset Y$，使得
$s_\xi \in \mathcal{L}(V)$，且 $s_\xi$ 生成 $\mathcal{L}|_V$。因此
$$
\text{div}_\mathcal{L}(s)|_V = \text{div}_Y(g)|_V.
$$
完全类似地，由于 $f^*s_\xi$ 在 $f^{-1}(V)$ 上生成 $f^*\mathcal{L}$，
且 $f^*s = g f^*s_\xi$，还有
$$
\text{div}_\mathcal{L}(f^*s)|_{f^{-1}(V)}
=
\text{div}_X(g)|_{f^{-1}(V)}.
$$
于是，$f^{-1}(V)$ 上所需的循环等式由主除子拉回的相应结论得到；见引理 \ref{chow-lemma-flat-pullback-principal-divisor}。
\end{proof}



\section{与可逆层相交}
\label{chow-section-intersecting-with-divisors}

\noindent
本节研究如下构造。

\begin{definition}
\label{chow-definition-cap-c1}
设 $(S, \delta)$ 如情形 \ref{chow-situation-setup}。
设 $X$ 在 $S$ 上局部有限型，并设 $\mathcal{L}$ 是可逆
$\mathcal{O}_X$-模。对每个整数 $k$，定义运算
$$
c_1(\mathcal{L}) \cap - :
Z_{k + 1}(X) \to \CH_k(X)
$$
称为{\it 与 $\mathcal{L}$ 的第一陈类相交}。
\begin{enumerate}
\item 给定满足 $\dim_\delta(W) = k + 1$ 的整闭子概形
$i : W \to X$，定义
$$
c_1(\mathcal{L}) \cap [W] = i_*(c_1({i^*\mathcal{L}}) \cap [W])
$$
其中右端由定义 \ref{chow-definition-divisor-invertible-sheaf} 定义。
\item 对一般的 $(k + 1)$-循环 $\alpha = \sum n_i [W_i]$，置
$$
c_1(\mathcal{L}) \cap \alpha = \sum n_i c_1(\mathcal{L}) \cap [W_i]
$$
\end{enumerate}
\end{definition}

\noindent
将每个 $c_1(\mathcal{L}) \cap W_i = \sum_j n_{i, j} [Z_{i, j}]$ 写成上述形式，
其中 $\{Z_{i, j}\}_j$ 是 $W_i$ 的整闭子概形的局部有限和。
由于 $\{W_i\}$ 是 $X$ 上整闭子概形的局部有限族，容易看出
$\{Z_{i, j}\}_{i, j}$ 是 $X$ 的闭子概形的局部有限族。因此
$c_1(\mathcal{L}) \cap \alpha = \sum n_in_{i, j}[Z_{i, j}]$ 是循环。
另一种更方便的理解方式是注意到态射 $\coprod W_i \to X$ 固有。
因此可将 $c_1(\mathcal{L}) \cap \alpha$ 看作
$\CH_k(\coprod W_i) = \prod \CH_k(W_i)$ 中某个类的推前。
这也说明所得结果在 $X$ 上模有理等价有良好定义。

\medskip\noindent
接下来几节的主要目标，是证明与 $c_1(\mathcal{L})$ 相交可下降到有理等价类。
这并非平凡事实。

\begin{lemma}
\label{chow-lemma-c1-cap-additive}
设 $(S, \delta)$ 如情形 \ref{chow-situation-setup}。
设 $X$ 在 $S$ 上局部有限型。设 $\mathcal{L}$、$\mathcal{N}$ 是 $X$ 上的可逆层。
那么对每个 $\alpha \in Z_{k + 1}(X)$，在 $\CH_k(X)$ 中有
$$
c_1(\mathcal{L}) \cap \alpha  + c_1(\mathcal{N}) \cap \alpha =
c_1(\mathcal{L} \otimes_{\mathcal{O}_X} \mathcal{N}) \cap \alpha
$$
此外，对所有 $\alpha$，有 $c_1(\mathcal{O}_X) \cap \alpha = 0$。
\end{lemma}

\begin{proof}
可加性直接由《除子》引理 \ref{divisors-lemma-c1-additive} 与定义得到。
为说明 $c_1(\mathcal{O}_X) \cap \alpha = 0$，考虑截面
$1 \in \Gamma(X, \mathcal{O}_X)$。它在任意整闭子概形
$W \subset X$ 上限制为处处非零的截面。因此
$c_1(\mathcal{O}_X) \cap [W] = 0$，正合所需。
\end{proof}

\noindent
回忆，$Z(s) \subset X$ 表示概形 $X$ 上某个可逆层的整体截面 $s$ 的零点概形；
见《除子》定义 \ref{divisors-definition-zero-scheme-s}。

\begin{lemma}
\label{chow-lemma-prepare-geometric-cap}
设 $(S, \delta)$ 如情形 \ref{chow-situation-setup}。
设 $Y$ 在 $S$ 上局部有限型。设 $\mathcal{L}$ 是可逆
$\mathcal{O}_Y$-模，并设 $s \in \Gamma(Y, \mathcal{L})$。假设
\begin{enumerate}
\item $\dim_\delta(Y) \leq k + 1$；
\item $\dim_\delta(Z(s)) \leq k$；
\item 对 $Z(s)$ 中 $\delta$-维数为 $k$ 的任意不可约分支的一般点 $\xi$，
乘以 $s$ 诱导单射 $\mathcal{O}_{Y, \xi} \to \mathcal{L}_\xi$。
\end{enumerate}
写成 $[Y]_{k + 1} = \sum n_i[Y_i]$，其中 $Y_i \subset Y$ 是 $Y$ 中
$\delta$-维数为 $k + 1$ 的不可约分支。置
$s_i = s|_{Y_i} \in \Gamma(Y_i, \mathcal{L}|_{Y_i})$。那么，作为 $Y$ 上的
$k$-循环，有
\begin{equation}
\label{chow-equation-equal-as-cycles}
[Z(s)]_k =  \sum n_i[Z(s_i)]_k
\end{equation}
\end{lemma}

\begin{proof}
设 $Z \subset Y$ 是 $\delta$-维数为 $k$ 的整闭子概形，
$\xi \in Z$ 是其一般点。要比较 $[Z]$ 在表达式
$\sum n_i[Z(s_i)]_k$ 中的系数 $n$ 与 $[Z]$ 在表达式 $[Z(s)]_k$ 中的系数 $m$。
选择生成元 $s_\xi \in \mathcal{L}_\xi$。写成
$A = \mathcal{O}_{Y, \xi}$、$L = \mathcal{L}_\xi$。那么
$L = As_\xi$。对某个唯一的 $f \in A$ 写成 $s = f s_\xi$。
假设 (3) 意味着 $f : A \to A$ 是单射。由于
$\dim_\delta(Y) \leq k + 1$ 且 $\dim_\delta(Z) = k$，有
$\dim(A) = 0$ 或 $1$。并且
$$
m = \text{length}_A(A/(f))
$$
在两种情形下都是有限的。

\medskip\noindent
若 $\dim(A) = 0$，则 $f : A \to A$ 为单射蕴含 $f \in A^*$。
所以此时 $m$ 为零。此外，条件 $\dim(A) = 0$ 意味着 $\xi$ 不位于任何
$\delta$-维数为 $k + 1$ 的不可约分支上，即 $n = 0$。

\medskip\noindent
现在设 $\dim(A) = 1$。由于 $A$ 是诺特局部环，它只有有限多个极小素理想
$\mathfrak q_1, \ldots, \mathfrak q_t$。它们与经过 $\xi'$ 的 $Y_i$ 一一对应。
此外，$n_i = \text{length}_{A_{\mathfrak q_i}}(A_{\mathfrak q_i})$。
$[Z]$ 在 $[Z(s_i)]_k$ 中的重数为
$\text{length}_A(A/(f, \mathfrak q_i))$。因此，此时需要证明的等式是
$$
\text{length}_A(A/(f))
=
\sum \text{length}_{A_{\mathfrak q_i}}(A_{\mathfrak q_i})
\text{length}_A(A/(f, \mathfrak q_i))
$$
它由引理 \ref{chow-lemma-additivity-divisors-restricted} 得出。
\end{proof}

\noindent
下面的引理可用于计算可逆层的 $c_1$ 与闭子概形所对应循环的交积。
回忆，$Z(s) \subset X$ 表示概形 $X$ 上某个可逆层的整体截面 $s$ 的零点概形；
见《除子》定义 \ref{divisors-definition-zero-scheme-s}。

\begin{lemma}
\label{chow-lemma-geometric-cap}
设 $(S, \delta)$ 如情形 \ref{chow-situation-setup}。
设 $X$ 在 $S$ 上局部有限型。设 $\mathcal{L}$ 是可逆
$\mathcal{O}_X$-模，$Y \subset X$ 是闭子概形，且
$s \in \Gamma(Y, \mathcal{L}|_Y)$。假设
\begin{enumerate}
\item $\dim_\delta(Y) \leq k + 1$；
\item $\dim_\delta(Z(s)) \leq k$；
\item 对 $Z(s)$ 中 $\delta$-维数为 $k$ 的任意不可约分支的一般点 $\xi$，
乘以 $s$ 诱导单射
$\mathcal{O}_{Y, \xi} \to (\mathcal{L}|_Y)_\xi$\footnote{例如，当
$s$ 是 $\mathcal{L}|_Y$ 的正则截面时，这一条件成立。}。
\end{enumerate}
那么在 $\CH_k(X)$ 中有
$$
c_1(\mathcal{L}) \cap [Y]_{k + 1} = [Z(s)]_k
$$
\end{lemma}

\begin{proof}
写成
$$
[Y]_{k + 1} = \sum n_i[Y_i]
$$
其中 $Y_i \subset Y$ 是 $Y$ 中 $\delta$-维数为 $k + 1$ 的不可约分支，
且 $n_i > 0$。由假设，限制
$s_i = s|_{Y_i} \in \Gamma(Y_i, \mathcal{L}|_{Y_i})$ 非零，因而是正则截面。
由引理 \ref{chow-lemma-compute-c1}，$[Z(s_i)]_k$ 表示
$c_1(\mathcal{L}|_{Y_i})$。所以按定义
$$
c_1(\mathcal{L}) \cap [Y]_{k + 1} = \sum n_i[Z(s_i)]_k
$$
结论于是由引理 \ref{chow-lemma-prepare-geometric-cap} 得出。
\end{proof}




\section{与可逆层相交及推前、拉回}
\label{chow-section-intersecting-with-divisors-push-pull}

\noindent
本节证明运算 $c_1(\mathcal{L}) \cap -$ 与平坦拉回、固有推前交换。

\begin{lemma}
\label{chow-lemma-prepare-flat-pullback-cap-c1}
设 $(S, \delta)$ 如情形 \ref{chow-situation-setup}。
设 $X$、$Y$ 在 $S$ 上局部有限型。设 $f : X \to Y$ 是相对维数为
$r$ 的平坦态射，$\mathcal{L}$ 是 $Y$ 上的可逆层。假设 $Y$ 是整概形且
$n = \dim_\delta(Y)$。设 $s$ 是 $\mathcal{L}$ 的非零亚纯截面。那么
$$
f^*\text{div}_\mathcal{L}(s) = \sum n_i\text{div}_{f^*\mathcal{L}|_{X_i}}(s_i)
$$
在 $Z_{n + r - 1}(X)$ 中成立。这里求和遍历其中
$\delta$-维数为 $n + r$ 的不可约分支 $X_i \subset X$；截面
$s_i = f|_{X_i}^*(s)$ 是 $s$ 的拉回；而
$n_i = m_{X_i, X}$ 是 $X_i$ 在 $X$ 中的重数。
\end{lemma}

\begin{proof}
为证明这个循环等式，可以在 $Y$ 上局部工作。因此可以假设 $Y$ 仿射，且对某些非零截面
$p \in \Gamma(Y, \mathcal{L})$ 和 $q \in \Gamma(Y, \mathcal{O})$，有
$s = p/q$。若能同时证明
$$
f^*\text{div}_\mathcal{L}(p) =
\sum n_i\text{div}_{f^*\mathcal{L}|_{X_i}}(p_i)
\quad\text{且}\quad
f^*\text{div}_\mathcal{O}(q) =
\sum n_i\text{div}_{\mathcal{O}_{X_i}}(q_i)
$$
（记号含义显然），则由可加性即得结论；见《除子》引理 \ref{divisors-lemma-c1-additive}。所以可以假设
$s \in \Gamma(Y, \mathcal{L})$。此时应用等式
(\ref{chow-equation-equal-as-cycles})，得到
$$
[Z(f^*(s))]_{k + r - 1} =
\sum n_i\text{div}_{f^*\mathcal{L}|_{X_i}}(s_i)
$$
其中 $f^*(s) \in f^*\mathcal{L}$ 表示 $s$ 到 $X$ 的拉回。另一方面，有
$$
f^*\text{div}_\mathcal{L}(s) = f^*[Z(s)]_{k - 1}
= [f^{-1}(Z(s))]_{k + r - 1},
$$
这由引理 \ref{chow-lemma-compute-c1} 和 \ref{chow-lemma-pullback-coherent} 得出。
由于 $Z(f^*(s)) = f^{-1}(Z(s))$，结论成立。
\end{proof}

\begin{lemma}
\label{chow-lemma-flat-pullback-cap-c1}
设 $(S, \delta)$ 如情形 \ref{chow-situation-setup}。
设 $X$、$Y$ 在 $S$ 上局部有限型。设 $f : X \to Y$ 是相对维数为
$r$ 的平坦态射，$\mathcal{L}$ 是 $Y$ 上的可逆层，$\alpha$ 是 $Y$ 上的
$k$-循环。那么在 $\CH_{k + r - 1}(X)$ 中
$$
f^*(c_1(\mathcal{L}) \cap \alpha) = c_1(f^*\mathcal{L}) \cap f^*\alpha
$$
\end{lemma}

\begin{proof}
写成 $\alpha = \sum n_i[W_i]$。我们将通过在 $X$ 的闭子概形
$f^{-1}(W_i)$ 上构造有理等价，证明在 $\CH_{k + r - 1}(X)$ 中
$$
f^*(c_1(\mathcal{L}) \cap [W_i]) = c_1(f^*\mathcal{L}) \cap f^*[W_i]
$$
注 \ref{chow-remark-infinite-sums-rational-equivalences} 中的讨论于是表明，
引理的等式成立。

\medskip\noindent
设 $W \subset Y$ 是 $\delta$-维数为 $k$ 的整闭子概形。
考虑闭子概形 $W' = f^{-1}(W) = W \times_Y X$，从而有纤维积图
$$
\xymatrix{
W' \ar[r] \ar[d]_h & X \ar[d]^f \\
W \ar[r] & Y
}
$$
需要证明
$f^*(c_1(\mathcal{L}) \cap [W]) = c_1(f^*\mathcal{L}) \cap f^*[W]$。
选择 $\mathcal{L}|_W$ 的非零亚纯截面 $s$。设
$W'_i \subset W'$ 是 $\delta$-维数为 $k + r$ 的不可约分支。
按定义写成 $[W']_{k + r} = \sum n_i[W'_i]$，其中 $n_i$ 是
$W'_i$ 在 $W'$ 中的重数。因此在 $Z_{k + r}(X)$ 中有
$f^*[W] = \sum n_i[W'_i]$。由于每个 $W'_i \to W$ 都支配，
对每个 $i$，$s_i = s|_{W'_i}$ 都是非零亚纯截面。由引理 \ref{chow-lemma-prepare-flat-pullback-cap-c1}，有循环等式
$$
h^*\text{div}_{\mathcal{L}|_W}(s) =
\sum n_i\text{div}_{f^*\mathcal{L}|_{W'_i}}(s_i)
$$
其位于 $Z_{k + r - 1}(W')$ 中。证明由此完成：左端是 $W'$ 上的循环，
其在 $\CH_{k + r - 1}(X)$ 中推前为
$f^*(c_1(\mathcal{L}) \cap [W])$；右端也是 $W'$ 上的循环，其在
$\CH_{k + r - 1}(X)$ 中推前为
$c_1(f^*\mathcal{L}) \cap f^*[W]$。
\end{proof}

\begin{lemma}
\label{chow-lemma-equal-c1-as-cycles}
设 $(S, \delta)$ 如情形 \ref{chow-situation-setup}。
设 $X$、$Y$ 在 $S$ 上局部有限型。设 $f : X \to Y$ 是固有态射，
$\mathcal{L}$ 是 $Y$ 上的可逆层。设 $s$ 是 $Y$ 上 $\mathcal{L}$ 的非零亚纯截面，
并将该截面记为 $s$。假设 $X$、$Y$ 都是整概形，$f$ 支配，且
$\dim_\delta(X) = \dim_\delta(Y)$。那么，作为 $Y$ 上的循环，
$$
f_*\left(\text{div}_{f^*\mathcal{L}}(f^*s)\right) =
[R(X) : R(Y)]\text{div}_\mathcal{L}(s).
$$
特别地，
$$
f_*(c_1(f^*\mathcal{L}) \cap [X]) =
[R(X) : R(Y)] c_1(\mathcal{L}) \cap [Y] =
c_1(\mathcal{L}) \cap f_*[X]
$$
\end{lemma}

\begin{proof}
由定义，$f_*[X] = [R(X) : R(Y)][Y]$，所以最后一个等式由第一个等式得出。
可以再次使用引理 \ref{chow-lemma-proper-pushforward-alteration} 来证明第一个等式。
因为所证的是循环等式，可以在 $Y$ 上局部工作。因此可以假设
$\mathcal{L} = \mathcal{O}_Y$。此时 $s$ 对应于有理函数 $g \in R(Y)$，
而所需证明的只是
$$
f_*\left(\text{div}_X(g)\right) =
[R(X) : R(Y)]\text{div}_Y(g).
$$
与上述引理 \ref{chow-lemma-proper-pushforward-alteration} 的结果比较，
由 $g \in R(Y)^*$ 时
$\text{Nm}_{R(X)/R(Y)}(g) = g^{[R(X) : R(Y)]}$，即得该等式。
\end{proof}

\begin{lemma}
\label{chow-lemma-pushforward-cap-c1}
设 $(S, \delta)$ 如情形 \ref{chow-situation-setup}。
设 $X$、$Y$ 在 $S$ 上局部有限型。设 $p : X \to Y$ 是固有态射，
$\alpha \in Z_{k + 1}(X)$，而 $\mathcal{L}$ 是 $Y$ 上的可逆层。
那么在 $\CH_k(Y)$ 中有
$$
p_*(c_1(p^*\mathcal{L}) \cap \alpha) = c_1(\mathcal{L}) \cap p_*\alpha
$$
\end{lemma}

\begin{proof}
先假设 $p$ 具有如下性质：对每个整闭子概形 $W \subset X$，映射
$p|_W : W \to Y$ 都是闭浸入。那么由与 $c_1(\mathcal{L})$ 作帽积的定义，
引理成立。

\medskip\noindent
利用这一观察约化到特殊情形。写成 $\alpha = \sum n_i[W_i]$，其中
$n_i \not = 0$，且各 $W_i$ 两两不同。设 $W'_i \subset Y$ 是 $W_i$ 的像
（视为整闭子概形）。考虑图表
$$
\xymatrix{
X' = \coprod W_i \ar[r]_-q \ar[d]_{p'} & X \ar[d]^p \\
Y' = \coprod W'_i \ar[r]^-{q'} & Y.
}
$$
由于 $\{W_i\}$ 在 $X$ 上局部有限，且 $p$ 固有，可知
$\{W'_i\}$ 在 $Y$ 上局部有限，并且 $q, q', p'$ 也都是固有态射。
还可以把 $\sum n_i[W_i]$ 看成 $k$-循环
$\alpha' \in Z_k(X')$。显然 $q_*\alpha' = \alpha$。由证明开头的观察，有
$q_*(c_1(q^*p^*\mathcal{L}) \cap \alpha')
= c_1(p^*\mathcal{L}) \cap q_*\alpha'$
以及
$(q')_*(c_1((q')^*\mathcal{L}) \cap p'_*\alpha') =
c_1(\mathcal{L}) \cap q'_*p'_*\alpha'$。因此只须对态射 $p'$ 和循环
$\sum n_i[W_i]$ 证明引理。显然，这意味着可以假设 $X$、$Y$ 都是整概形，
$f : X \to Y$ 支配，且 $\alpha = [X]$。此时结论由引理 \ref{chow-lemma-equal-c1-as-cycles} 得出。
\end{proof}






\section{关键公式}
\label{chow-section-key}

\noindent
设 $(S, \delta)$ 如情形 \ref{chow-situation-setup}。
设 $X$ 在 $S$ 上局部有限型，并假设
$X$ 是整概形且 $\dim_\delta(X) = n$。
设 $\mathcal{L}$ 和 $\mathcal{N}$ 是 $X$ 上的可逆层。
设 $s$ 是 $\mathcal{L}$ 的非零亚纯截面，而
$t$ 是 $\mathcal{N}$ 的非零亚纯截面。
设 $Z_i \subset X$（$i \in I$）是一族局部有限的余维 $1$
不可约闭子集，并具有如下性质：
若 $Z \not \in \{Z_i\}$ 且 $\xi$ 是其一般点，则 $s$ 是
$\mathcal{L}_\xi$ 的生成元，而 $t$ 是 $\mathcal{N}_\xi$ 的生成元。
由《除子》引理 \ref{divisors-lemma-divisor-meromorphic-locally-finite}，这样的集合存在。
于是
$$
\text{div}_\mathcal{L}(s) = \sum \text{ord}_{Z_i, \mathcal{L}}(s) [Z_i]
$$
并且类似地，
$$
\text{div}_\mathcal{N}(t) = \sum \text{ord}_{Z_i, \mathcal{N}}(t) [Z_i]
$$
进一步展开定义：对每个 $i$ 选取生成元
$s_i \in \mathcal{L}_{\xi_i}$ 和 $t_i \in \mathcal{N}_{\xi_i}$，
其中 $\xi_i$ 是 $Z_i$ 的一般点。于是可以写成
$$
s = f_i s_i
\quad\text{且}\quad
t = g_i t_i
$$
置 $B_i = \mathcal{O}_{X, \xi_i}$。由定义，
$$
\text{ord}_{Z_i, \mathcal{L}}(s) = \text{ord}_{B_i}(f_i)
\quad\text{且}\quad
\text{ord}_{Z_i, \mathcal{N}}(t) = \text{ord}_{B_i}(g_i)
$$
因为 $t_i$ 是 $\mathcal{N}_{\xi_i}$ 的生成元，它在纤维
$\mathcal{N}_{\xi_i} \otimes \kappa(\xi_i)$ 中的像是
$\mathcal{N}|_{Z_i}$ 的非零亚纯截面。将此像记为
$t_i|_{Z_i}$。由定义可得
$$
c_1(\mathcal{N}) \cap \text{div}_\mathcal{L}(s) =
\sum \text{ord}_{B_i}(f_i)
(Z_i \to X)_*\text{div}_{\mathcal{N}|_{Z_i}}(t_i|_{Z_i})
$$
以及类似的等式
$$
c_1(\mathcal{L}) \cap \text{div}_\mathcal{N}(t) =
\sum \text{ord}_{B_i}(g_i)
(Z_i \to X)_*\text{div}_{\mathcal{L}|_{Z_i}}(s_i|_{Z_i})
$$
它们位于 $\CH_{n - 2}(X)$ 中。下面要找出这两个循环之间的有理等价。
为此，考虑驯顺符号
$$
\partial_{B_i}(f_i, g_i) \in \kappa(\xi_i)^*
$$
见第 \ref{chow-section-tame-symbol} 节。

\begin{lemma}[关键公式]
\label{chow-lemma-key-formula}
在上述情形中，循环
$$
\sum
(Z_i \to X)_*\left(
\text{ord}_{B_i}(f_i) \text{div}_{\mathcal{N}|_{Z_i}}(t_i|_{Z_i}) -
\text{ord}_{B_i}(g_i) \text{div}_{\mathcal{L}|_{Z_i}}(s_i|_{Z_i}) \right)
$$
等于循环
$$
\sum (Z_i \to X)_*\text{div}(\partial_{B_i}(f_i, g_i))
$$
\end{lemma}

\begin{proof}
先考察用 $us_i$ 替换 $s_i$ 时会发生什么，其中 $u$ 是 $B_i$ 中的单位。
此时 $f_i$ 被 $u^{-1} f_i$ 替换。因此，引理第一个表达式的第一部分不变，
而在第二部分中增加
$$
-\text{ord}_{B_i}(g_i)\text{div}(u|_{Z_i})
$$
（其中 $u|_{Z_i}$ 是 $u$ 在剩余域中的像）；这是
《除子》引理 \ref{divisors-lemma-divisor-meromorphic-well-defined}
的结论。第二个表达式中则增加
$$
\text{div}(\partial_{B_i}(u^{-1}, g_i))
$$
这是驯顺符号双线性的结果。由驯顺符号的性质
(\ref{chow-item-normalization})，这两项相等。

\medskip\noindent
设 $Z \subset X$ 是满足 $\dim_\delta(Z) = n - 2$ 的不可约闭子集。
为证明引理中两个循环的 $Z$ 系数相同，可以像上一段那样作替换
$s_i \mapsto us_i$。完全类似地，也可以用 $t_i \mapsto vt_i$
作替换，其中 $v$ 是 $B_i$ 的单位。

\medskip\noindent
由于所证的是循环相等，可以逐个系数论证。因此，选取满足
$\dim_\delta(Z) = n - 2$ 的不可约闭子集 $Z \subset X$，并比较其系数。
设 $\xi \in Z$ 为一般点，并置 $A = \mathcal{O}_{X, \xi}$。
这是一个维数为 $2$ 的 Noether 局部整环。选取
$\mathcal{L}_\xi$ 和 $\mathcal{N}_\xi$ 的生成元 $\sigma$ 和 $\tau$。
缩小 $X$ 后，可以并且确实假设 $\sigma$ 和 $\tau$ 在整个 $X$ 上
给出可逆层 $\mathcal{L}$ 和 $\mathcal{N}$ 的平凡化。
由于缩小 $X$ 后 $Z_i$ 局部有限，可以假设对所有 $i \in I$ 都有
$Z \subset Z_i$，并且 $I$ 有限。此时 $\xi_i$ 对应于高度为 $1$ 的素理想
$\mathfrak q_i \subset A$。可以写成 $s_i = a_i \sigma$ 和
$t_i = b_i \tau$，其中 $a_i$ 和 $b_i$ 是 $A_{\mathfrak q_i}$ 中的单位。
由上述讨论，只须在所有 $i$ 都满足 $a_i = b_i = 1$ 时证明引理。

\medskip\noindent
假设所有 $i$ 都有 $a_i = b_i = 1$。由于所选 $\sigma$ 和 $\tau$
是平凡化截面，引理的第一个表达式为零。写
$s = f\sigma$ 和 $t = g \tau$，其中 $f$、$g$ 属于 $A$ 的分式域。
由上一段，问题已约化到对所有 $i$ 都有 $f_i = f$ 和 $g_i = g$ 的情形。
此外，若 $A$ 的高度 $1$ 素理想 $\mathfrak q$ 不在
$\{\mathfrak q_i\}$ 中，则 $f$、$g$ 都是 $A_\mathfrak q$ 中的单位
（这是由引理之前讨论中对族 $\{Z_i\}$ 的选取所得）。因此，引理第二个
表达式中 $Z$ 的系数为
$$
\sum\nolimits_i \text{ord}_{A/\mathfrak q_i}(\partial_{B_i}(f, g))
$$
而关键引理引理 \ref{chow-lemma-milnor-gersten-low-degree} 说明它为零。
\end{proof}

\begin{remark}
\label{chow-remark-higher-chow-pointwise}
设 $(S, \delta)$ 如情形 \ref{chow-situation-setup}。
设 $X$ 在 $S$ 上局部有限型，并设 $k \in \mathbf{Z}$。
我们断言存在复形
$$
\bigoplus\nolimits_{\delta(x) = k + 2}' K_2^M(\kappa(x))
\xrightarrow{\partial}
\bigoplus\nolimits_{\delta(x) = k + 1}' K_1^M(\kappa(x))
\xrightarrow{\partial}
\bigoplus\nolimits_{\delta(x) = k}' K_0^M(\kappa(x))
$$
这里采用注 \ref{chow-remark-chow-group-pointwise} 中引入的记号和约定，
并另作如下约定：
\begin{enumerate}
\item $K_2^M(\kappa(x))$ 是点 $x \in X$ 的剩余域 $\kappa(x)$ 的
Milnor K-理论的次数 $2$ 部分（见注 \ref{chow-remark-gersten-complex-milnor}）；它是
$\kappa(x)^* \otimes_\mathbf{Z} \kappa(x)^*$ 关于由形如
$\lambda \otimes (1 - \lambda)$ 的元素生成的子群所取的商，
其中 $\lambda \in \kappa(x) \setminus \{0, 1\}$；以及
\item 第一个微分 $\partial$ 定义如下：给定第一项中的元素
$\xi = \sum_x \alpha_x$，置
$$
\partial(\xi) = \sum\nolimits_{x \leadsto x',\ \delta(x') = k + 1}
\partial_{\mathcal{O}_{W_x, x'}}(\alpha_x)
$$
其中
$\partial_{\mathcal{O}_{W_x, x'}} : K_2^M(\kappa(x)) \to K_1^M(\kappa(x))$
是第 \ref{chow-section-tame-symbol} 节中构造的驯顺符号。
\end{enumerate}
我们断言所得确为复形，即 $\partial \circ \partial = 0$。
为此，只须取上述元素 $\xi$ 以及满足 $\delta(x'') = k$ 的点
$x'' \in X$，并验证元素 $\partial(\partial(\xi))$ 中 $x''$ 的系数为零。
由于 $\xi = \sum \alpha_x$ 是局部有限和，利用可加性，实际上可以假设
$\xi = \alpha_x$，其中某个 $x \in X$ 满足 $\delta(x) = k + 2$，且
$\alpha_x \in K_2^M(\kappa(x))$。再由线性性，可以假设
$\alpha_x = f \otimes g$，其中 $f, g \in \kappa(x)^*$。
以 $W \subset X$ 表示一般点为 $x$ 的整闭子概形。若
$x'' \not \in W$，则立即可见，$\partial(\partial(\xi))$ 中 $x$ 的系数为零。
若 $x'' \in W$，则 $\partial(\partial(x))$ 中 $x''$ 的系数等于
$$
\sum\nolimits_{x \leadsto x' \leadsto x'',\ \delta(x') = k + 1}
\text{ord}_{\mathcal{O}_{\overline{\{x'\}}, x''}}(
\partial_{\mathcal{O}_{W, x'}}(f, g))
$$
关键代数引理引理 \ref{chow-lemma-milnor-gersten-low-degree}
恰好说明这个和为零。
\end{remark}

\begin{remark}
\label{chow-remark-higher-chow}
设 $(S, \delta)$ 如情形 \ref{chow-situation-setup}。
设 $X$ 在 $S$ 上局部有限型，并设 $k \in \mathbf{Z}$。
注 \ref{chow-remark-higher-chow-pointwise} 中的复形以及注 \ref{chow-remark-chow-group-pointwise} 中 $\CH_k(X)$ 的表示提示我们，
可以定义第一高阶 Chow 群
$$
\CH^M_k(X, 1) =
H_1(\text{the complex of Remark \ref{chow-remark-higher-chow-pointwise}})
$$
我们使用上标 ${}^M$，以区别于文献中定义的高阶 Chow 群，
例如 Spencer Bloch 的论文（\cite{Bloch} 和 \cite{Bloch-moving}）中的定义。
设 $U \subset X$ 是开子集，其补集为 $Y \subset X$（视为既约闭子概形）。
于是对 $i = 2, 1, 0$ 有分裂短正合列
$$
0 \to
\bigoplus\nolimits_{y \in Y, \delta(y) = k + i}' K_i^M(\kappa(y)) \to
\bigoplus\nolimits_{x \in X, \delta(x) = k + i}' K_i^M(\kappa(x)) \to
\bigoplus\nolimits_{u \in U, \delta(u) = k + i}' K_i^M(\kappa(u)) \to 0
$$
它与注 \ref{chow-remark-higher-chow-pointwise} 中各复形的边界映射相容。
应用蛇形引理（见《同调代数》引理 \ref{homology-lemma-long-exact-sequence-chain}），得到六项正合列
$$
\CH^M_k(Y, 1) \to \CH^M_k(X, 1) \to \CH^M_k(U, 1) \to
\CH_k(Y) \to \CH_k(X) \to \CH_k(U) \to 0
$$
它延拓了引理 \ref{chow-lemma-restrict-to-open} 的典范正合列。
稍作进一步工作，还可以为第一高阶 Chow 群
$\CH^M_k(X, 1)$ 定义平坦拉回与固有推前。我们将在后文回到这一点
（此处插入未来引用）。
\end{remark}






\section{与可逆层相交及有理等价}
\label{chow-section-commutativity}

\noindent
应用关键引理，可以得到与可逆层相交的基本性质。特别地，我们将看到
$c_1(\mathcal{L}) \cap -$ 下降到有理等价类，并且不同可逆层所给出的
这类运算彼此交换。

\begin{lemma}
\label{chow-lemma-commutativity-on-integral}
设 $(S, \delta)$ 如情形 \ref{chow-situation-setup}。
设 $X$ 在 $S$ 上局部有限型，并假设 $X$ 是整概形且
$\dim_\delta(X) = n$。设 $\mathcal{L}$、$\mathcal{N}$ 是 $X$ 上的可逆层。
选取 $\mathcal{L}$ 的非零亚纯截面 $s$ 和 $\mathcal{N}$ 的非零亚纯截面
$t$。置 $\alpha = \text{div}_\mathcal{L}(s)$ 及
$\beta = \text{div}_\mathcal{N}(t)$。则
$$
c_1(\mathcal{N}) \cap \alpha
=
c_1(\mathcal{L}) \cap \beta
$$
在 $\CH_{n - 2}(X)$ 中成立。
\end{lemma}

\begin{proof}
这立即由关键引理引理 \ref{chow-lemma-key-formula} 及其之前的讨论得出。
\end{proof}

\begin{lemma}
\label{chow-lemma-factors}
设 $(S, \delta)$ 如情形 \ref{chow-situation-setup}。
设 $X$ 在 $S$ 上局部有限型，$\mathcal{L}$ 是 $X$ 上的可逆层。
运算 $\alpha \mapsto c_1(\mathcal{L}) \cap \alpha$ 下降到有理等价类，
从而给出运算
$$
c_1(\mathcal{L}) \cap - : \CH_{k + 1}(X) \to \CH_k(X)
$$
\end{lemma}

\begin{proof}
设 $\alpha \in Z_{k + 1}(X)$ 且 $\alpha \sim_{rat} 0$。
需要证明定义 \ref{chow-definition-cap-c1} 所定义的
$c_1(\mathcal{L}) \cap \alpha$ 为零。
由定义 \ref{chow-definition-rational-equivalence}，存在局部有限族
$\{W_j\}$，其成员都是满足 $\dim_\delta(W_j) = k + 2$
的整闭子概形；还存在有理函数 $f_j \in R(W_j)^*$，使得
$$
\alpha = \sum (i_j)_*\text{div}_{W_j}(f_j)
$$
注意，$p : \coprod W_j \to X$ 是固有态射。因此
$\alpha = p_*\alpha'$，其中 $\alpha' \in Z_{k + 1}(\coprod W_j)$
是主除子 $\text{div}_{W_j}(f_j)$ 之和。
由引理 \ref{chow-lemma-pushforward-cap-c1}，
$c_1(\mathcal{L}) \cap \alpha = p_*(c_1(p^*\mathcal{L}) \cap \alpha')$。
故只须证明每个
$c_1(\mathcal{L}|_{W_j}) \cap \text{div}_{W_j}(f_j)$ 都为零。
换言之，可以假设 $X$ 是整概形，并且对某个 $f \in R(X)^*$ 有
$\alpha = \text{div}_X(f)$。

\medskip\noindent
假设 $X$ 是整概形，且对某个 $f \in R(X)^*$ 有
$\alpha = \text{div}_X(f)$。可以把 $f$ 看成可逆层
$\mathcal{N} = \mathcal{O}_X$ 的正则亚纯截面。选取 $\mathcal{L}$ 的
一个亚纯截面 $s$，并记 $\beta = \text{div}_\mathcal{L}(s)$。
由引理 \ref{chow-lemma-commutativity-on-integral} 可得
$$
c_1(\mathcal{L}) \cap \alpha = c_1(\mathcal{O}_X) \cap \beta.
$$
然而引理 \ref{chow-lemma-c1-cap-additive} 说明右边在 $\CH_k(X)$ 中为零，
正合所需。
\end{proof}

\noindent
设 $(S, \delta)$ 如情形 \ref{chow-situation-setup}。
设 $X$ 在 $S$ 上局部有限型，$\mathcal{L}$ 是 $X$ 上的可逆层。
将运算 $c_1(\mathcal{L}) \cap - $ 记为
$$
c_1(\mathcal{L}) \cap - : \CH_{k + 1}(X) \to \CH_k(X)
$$
由引理 \ref{chow-lemma-factors}，这个记法有意义。对所有 $s \geq 0$，
以 $c_1(\mathcal{L})^s \cap -$ 表示该运算的 $s$ 次迭代。

\begin{lemma}
\label{chow-lemma-cap-commutative}
设 $(S, \delta)$ 如情形 \ref{chow-situation-setup}。
设 $X$ 在 $S$ 上局部有限型，$\mathcal{L}$、$\mathcal{N}$ 是 $X$ 上的
可逆层。对任意 $\alpha \in \CH_{k + 2}(X)$，都有
$$
c_1(\mathcal{L}) \cap c_1(\mathcal{N}) \cap \alpha
=
c_1(\mathcal{N}) \cap c_1(\mathcal{L}) \cap \alpha
$$
作为 $\CH_k(X)$ 中的元素。
\end{lemma}

\begin{proof}
将其写成 $\alpha = \sum m_j[Z_j]$，其中
$Z_j \subset X$ 是满足 $\dim_\delta(Z_j) = k + 2$ 的一族局部有限的
整闭子概形。考虑固有态射 $p : \coprod Z_j \to X$。
把 $\alpha' = \sum m_j[Z_j]$ 视为 $\coprod Z_j$ 上的
$(k + 2)$-循环。反复应用引理 \ref{chow-lemma-pushforward-cap-c1}，得到
$c_1(\mathcal{L}) \cap c_1(\mathcal{N}) \cap \alpha
= p_*(c_1(p^*\mathcal{L}) \cap c_1(p^*\mathcal{N}) \cap \alpha')$
以及
$c_1(\mathcal{N}) \cap c_1(\mathcal{L}) \cap \alpha
= p_*(c_1(p^*\mathcal{N}) \cap c_1(p^*\mathcal{L}) \cap \alpha')$。
因此，只须在 $X$ 为整概形且 $\alpha = [X]$ 的情形证明该公式。
此时结论由引理 \ref{chow-lemma-commutativity-on-integral} 和定义得出。
\end{proof}






\section{Gysin 同态}
\label{chow-section-intersecting-effective-Cartier}

\noindent
本节为可逆层的一个截面的零点概形 $D$ 定义 Gysin 映射。
当 $D$ 是有效 Cartier 除子时会出现一个重要情形；但把定义推广到任意
$D$，可以更灵活地表述各种相容性，见注 \ref{chow-remark-pullback-pairs} 以及引理 \ref{chow-lemma-closed-in-X-gysin}、\ref{chow-lemma-gysin-flat-pullback} 和
\ref{chow-lemma-gysin-commutes-gysin}。
这些结果还可以推广到带有虚法丛的局部主闭子概形
（注 \ref{chow-remark-generalize-to-virtual}），或推广到伪除子
（注 \ref{chow-remark-generalize-to-pseudo-divisor}）。

\medskip\noindent
回忆，有效 Cartier 除子与配对 $(\mathcal{L}, s)$ 的同构类按 $1$ 对 $1$ 对应，
其中 $\mathcal{L}$ 是可逆层，而 $s$ 是正则整体截面；见
《除子》引理 \ref{divisors-lemma-characterize-OD}。
若 $D$ 对应于 $(\mathcal{L}, s)$，则
$\mathcal{L} = \mathcal{O}_X(D)$。阅读本节时请记住这一点。

\begin{definition}
\label{chow-definition-gysin-homomorphism}
设 $(S, \delta)$ 如情形 \ref{chow-situation-setup}。
设 $X$ 在 $S$ 上局部有限型。设 $(\mathcal{L}, s)$ 是由一个可逆层
和一个整体截面 $s \in \Gamma(X, \mathcal{L})$ 组成的配对。
设 $D = Z(s)$ 是 $s$ 的零点概形，并以
$i : D \to X$ 表示闭浸入。对每个整数 $k$，按如下规则定义一个
{\it Gysin 同态}
$$
i^* : Z_{k + 1}(X) \to \CH_k(D).
$$
\begin{enumerate}
\item 给定满足 $\dim_\delta(W) = k + 1$ 的整闭子概形
$W \subset X$，定义
\begin{enumerate}
\item 若 $W \not \subset D$，则令 $i^*[W] = [D \cap W]_k$，
把它视为 $D$ 上的 $k$-循环；以及
\item 若 $W \subset D$，则令
$i^*[W] = i'_*(c_1(\mathcal{L}|_W) \cap [W])$，
其中 $i' : W \to D$ 是诱导的闭浸入。
\end{enumerate}
\item 对一般的 $(k + 1)$-循环 $\alpha = \sum n_j[W_j]$，置
$$
i^*\alpha = \sum n_j i^*[W_j]
$$
\item 若 $D$ 是有效 Cartier 除子，则以
$D \cdot \alpha = i_*i^*\alpha$ 表示类 $i^*\alpha$ 推前后所得的
$X$ 上的类。
\end{enumerate}
\end{definition}

\noindent
事实上，正如下文将看到的，这个 Gysin 同态 $i^*$ 可以视为非平坦拉回
的一个例子。因此，有时也非正式地把类 $i^*\alpha$ 称为类 $\alpha$ 的
{\it 拉回}。

\begin{remark}
\label{chow-remark-generalize-to-virtual}
设 $X$ 是情形 \ref{chow-situation-setup} 中那样在 $S$ 上局部有限型
的概形。设 $(D, \mathcal{N}, \sigma)$ 是如下三元组：局部主闭子概形
（《除子》定义 \ref{divisors-definition-effective-Cartier-divisor}）
$i : D \to X$、可逆 $\mathcal{O}_D$-模 $\mathcal{N}$，以及
$\mathcal{O}_D$-模的满射
$\sigma : \mathcal{N}^{\otimes -1} \to i^*\mathcal{I}_D$
\footnote{这个条件保证：若 $D$ 是有效 Cartier 除子，则
$\mathcal{N} = \mathcal{O}_X(D)|_D$。}。
应把 $\mathcal{N}$ 看作{\it $D$ 在 $X$ 中的虚法丛}。
定义 \ref{chow-definition-gysin-homomorphism} 中的构造
$i^* : Z_{k + 1}(X) \to \CH_k(D)$ 可以推广到这类三元组，见
第 \ref{chow-section-gysin-higher-codimension} 节。
\end{remark}

\begin{remark}
\label{chow-remark-generalize-to-pseudo-divisor}
设 $X$ 是情形 \ref{chow-situation-setup} 中那样在 $S$ 上局部有限型
的概形。在 \cite{F} 中，$X$ 上的{\it 伪除子}定义为三元组
$D = (\mathcal{L}, Z, s)$，其中 $\mathcal{L}$ 是可逆
$\mathcal{O}_X$-模，$Z \subset X$ 是闭子集，而
$s \in \Gamma(X \setminus Z, \mathcal{L})$ 是处处非零截面。
与上述情形类似，对 $\CH_{k + 1}(X)$ 中每个 $\alpha$，
可以在 $\CH_k(Z \cap |\alpha|)$ 中定义乘积 $D \cdot \alpha$，
其中 $|\alpha|$ 是 $\alpha$ 的支撑。
\end{remark}

\begin{lemma}
\label{chow-lemma-support-cap-effective-Cartier}
设 $(S, \delta)$ 如情形 \ref{chow-situation-setup}，并设 $X$ 在 $S$ 上
局部有限型。设 $(\mathcal{L}, s, i : D \to X)$ 如定义 \ref{chow-definition-gysin-homomorphism}，而 $\alpha$ 是 $X$ 上的
$(k + 1)$-循环。那么在 $\CH_k(X)$ 中，
$i_*i^*\alpha = c_1(\mathcal{L}) \cap \alpha$。
特别地，若 $D$ 是有效 Cartier 除子，则
$D \cdot \alpha = c_1(\mathcal{O}_X(D)) \cap \alpha$。
\end{lemma}

\begin{proof}
写 $\alpha = \sum n_j[W_j]$，其中 $i_j : W_j \to X$ 是满足
$\dim_\delta(W_j) = k$ 的整闭子概形。
由于 $D$ 是 $s$ 的零点概形，$D \cap W_j$ 是限制 $s|_{W_j}$ 的零点概形。
因此，对每个满足 $W_j \not \subset D$ 的 $j$，由引理 \ref{chow-lemma-geometric-cap} 有
$c_1(\mathcal{L}) \cap [W_j] = [D \cap W_j]_k$。于是，由定义 \ref{chow-definition-cap-c1}，在 $\CH_k(X)$ 中有
$$
c_1(\mathcal{L}) \cap \alpha
=
\sum\nolimits_{W_j \not \subset D} n_j[D \cap W_j]_k
+
\sum\nolimits_{W_j \subset D}
n_j i_{j, *}(c_1(\mathcal{L})|_{W_j}) \cap [W_j])
$$
右边逐项等于定义 \ref{chow-definition-gysin-homomorphism} 中
$D$ 上的类 $i^*\alpha$ 的推前。因此结论成立。
\end{proof}

\begin{lemma}
\label{chow-lemma-easy-gysin}
设 $(S, \delta)$ 如情形 \ref{chow-situation-setup}。
设 $X$ 在 $S$ 上局部有限型。设
$(\mathcal{L}, s, i : D \to X)$ 如定义 \ref{chow-definition-gysin-homomorphism}。
\begin{enumerate}
\item 设 $Z \subset X$ 是满足 $\dim_\delta(Z) \leq k + 1$ 的闭子概形，
并且 $D \cap Z$ 是 $Z$ 上的有效 Cartier 除子。则
$i^*[Z]_{k + 1} = [D \cap Z]_k$。
\item 设 $\mathcal{F}$ 是 $X$ 上的凝聚层，满足
$\dim_\delta(\text{Supp}(\mathcal{F})) \leq k + 1$，并且
$s : \mathcal{F} \to \mathcal{F} \otimes_{\mathcal{O}_X} \mathcal{L}$
是单射。则
$$
i^*[\mathcal{F}]_{k + 1} = [i^*\mathcal{F}]_k
$$
在 $\CH_k(D)$ 中成立。
\end{enumerate}
\end{lemma}

\begin{proof}
设 $Z \subset X$ 如 (1)，并置 $\mathcal{F} = \mathcal{O}_Z$。
$D \cap Z$ 是有效 Cartier 除子的假设等价于
$s : \mathcal{F} \to \mathcal{F} \otimes_{\mathcal{O}_X} \mathcal{L}$
为单射的假设。此外，
$[Z]_{k + 1} = [\mathcal{F}]_{k + 1}]$，并且
$[D \cap Z]_k = [\mathcal{O}_{D \cap Z}]_k = [i^*\mathcal{F}]_k$。
见引理 \ref{chow-lemma-cycle-closed-coherent}。因此 (1) 由 (2) 得出。

\medskip\noindent
写 $[\mathcal{F}]_{k + 1} = \sum m_j[W_j]$，其中 $m_j > 0$，
而 $W_j \subset X$ 是两两不同、$\delta$-维数为 $k + 1$ 的整闭子概形。
$s : \mathcal{F} \to \mathcal{F} \otimes_{\mathcal{O}_X} \mathcal{L}$
为单射这一假设蕴含对所有 $j$ 都有 $W_j \not \subset D$。
由定义，
$$
i^*[\mathcal{F}]_{k + 1} = \sum m_j [D \cap W_j]_k.
$$
我们断言循环等式
$$
\sum [D \cap W_j]_k = [i^*\mathcal{F}]_k
$$
成立。
设 $Z \subset D$ 是 $\delta$-维数为 $k$ 的整闭子概形。
设 $\xi \in Z$ 为其一般点，$A = \mathcal{O}_{X, \xi}$，
$M = \mathcal{F}_\xi$。设 $f \in A$ 生成 $D$ 的理想，即满足
$\mathcal{O}_{D, \xi} = A/fA$。根据假设，
$\dim(\text{Supp}(M)) = 1$，映射 $f : M \to M$ 是单射，并且
$\text{length}_A(M/fM) < \infty$。此外，$\text{length}_A(M/fM)$
是 $[i^*\mathcal{F}]_k$ 中 $[Z]$ 的系数。另一方面，设
$\mathfrak q_1, \ldots, \mathfrak q_t$ 是 $M$ 的支撑中的极小素理想。
那么
$$
\sum
\text{length}_{A_{\mathfrak q_i}}(M_{\mathfrak q_i})
\text{ord}_{A/\mathfrak q_i}(f)
$$
是 $\sum [D \cap W_j]_k$ 中 $[Z]$ 的系数。
因此，该等式由引理 \ref{chow-lemma-additivity-divisors-restricted} 得出。
\end{proof}

\begin{remark}
\label{chow-remark-gysin-on-cycles}
设 $X \to S$、$\mathcal{L}$、$s$、$i : D \to X$ 如定义 \ref{chow-definition-gysin-homomorphism}，并假设
$\mathcal{L}|_D \cong \mathcal{O}_D$。在这种情形，可以在循环上定义
典范映射 $i^* : Z_{k + 1}(X) \to Z_k(D)$：每当
$W \subset D$ 是整闭子概形时，要求 $i^*[W] = 0$。
这种做法将在后文有用。
\end{remark}

\begin{remark}
\label{chow-remark-pullback-pairs}
设 $f : X' \to X$ 是情形 \ref{chow-situation-setup} 中那样的
$S$ 上局部有限型概形间的态射。设
$(\mathcal{L}, s, i : D \to X)$ 是定义 \ref{chow-definition-gysin-homomorphism} 中的三元组。可以置
$\mathcal{L}' = f^*\mathcal{L}$、$s' = f^*s$ 以及
$D' = X' \times_X D = Z(s')$。这给出交换图
$$
\xymatrix{
D' \ar[d]_g \ar[r]_{i'} & X' \ar[d]^f \\
D \ar[r]^i & X
}
$$
于是可以考察 $i^*$ 与 $(i')^*$ 之间的各种相容性。
\end{remark}

\begin{lemma}
\label{chow-lemma-closed-in-X-gysin}
设 $(S, \delta)$ 如情形 \ref{chow-situation-setup}。
设 $f : X' \to X$ 是 $S$ 上局部有限型概形间的固有态射。
设 $(\mathcal{L}, s, i : D \to X)$ 如定义 \ref{chow-definition-gysin-homomorphism}。依注 \ref{chow-remark-pullback-pairs} 构成图表
$$
\xymatrix{
D' \ar[d]_g \ar[r]_{i'} & X' \ar[d]^f \\
D \ar[r]^i & X
}
$$
对 $X'$ 上任意 $(k + 1)$-循环 $\alpha'$，在 $\CH_k(D)$ 中有
$i^*f_*\alpha' = g_*(i')^*\alpha'$
（这是有意义的，因为 $f_*$ 已在循环层次上定义）。
\end{lemma}

\begin{proof}
设对某个整闭子概形 $W' \subset X'$ 有 $\alpha = [W']$。
令 $W = f(W') \subset X$。若 $W' \not \subset D'$，则
$W \not \subset D$，并且
$$
[W' \cap D']_k = \text{div}_{\mathcal{L}'|_{W'}}({s'|_{W'}})
\quad\text{且}\quad
[W \cap D]_k = \text{div}_{\mathcal{L}|_W}(s|_W)
$$
因此，由引理 \ref{chow-lemma-equal-c1-as-cycles}，第一个循环沿
$f_*$ 的像等于第二个循环，故等式作为循环成立。
若 $W' \subset D'$，则 $W \subset D$；由引理 \ref{chow-lemma-equal-c1-as-cycles} 的第二个断言，
$f_*(c_1(\mathcal{L}|_{W'}) \cap [W'])$ 在 $\CH_k(W)$ 中等于
$c_1(\mathcal{L}|_W) \cap [W]$。
由注 \ref{chow-remark-infinite-sums-rational-equivalences}，
对一般的 $\alpha'$ 也得到结论。
\end{proof}

\begin{lemma}
\label{chow-lemma-gysin-flat-pullback}
设 $(S, \delta)$ 如情形 \ref{chow-situation-setup}。设
$f : X' \to X$ 是 $S$ 上局部有限型概形间相对维数为 $r$ 的平坦态射。
设 $(\mathcal{L}, s, i : D \to X)$ 如定义 \ref{chow-definition-gysin-homomorphism}。依注 \ref{chow-remark-pullback-pairs} 构成图表
$$
\xymatrix{
D' \ar[d]_g \ar[r]_{i'} & X' \ar[d]^f \\
D \ar[r]^i & X
}
$$
对 $X$ 上任意 $(k + 1)$-循环 $\alpha$，在 $\CH_{k + r}(D')$ 中有
$(i')^*f^*\alpha = g^*i^*\alpha$
（这是有意义的，因为 $f^*$ 已在循环层次上定义）。
\end{lemma}

\begin{proof}
设对某个整闭子概形 $W \subset X$ 有 $\alpha = [W]$。
令 $W' = f^{-1}(W) \subset X'$。若 $W \not \subset D$，则
$W' \not \subset D'$，并且
$$
W' \cap D' = g^{-1}(W \cap D)
$$
作为 $D'$ 的闭子概形。故等式作为循环成立；见
引理 \ref{chow-lemma-pullback-coherent}。
若 $W \subset D$，则 $W' \subset D'$ 且 $W' = g^{-1}(W)$，
并有 $[W']_{k + 1 + r} = g^*[W]$；此时等式由引理 \ref{chow-lemma-flat-pullback-cap-c1} 在 $\CH_{k + r}(D')$ 中成立。
由注 \ref{chow-remark-infinite-sums-rational-equivalences}，
对一般的 $\alpha'$ 也得到结论。
\end{proof}








\section{Gysin 同态与有理等价}
\label{chow-section-gysin}

\noindent
本节使用关键公式证明 Gysin 同态下降到有理等价类。还将证明一个重要的交换性。

\begin{lemma}
\label{chow-lemma-gysin-factors-general}
设 $(S, \delta)$ 如情形 \ref{chow-situation-setup}。
设 $X$ 在 $S$ 上局部有限型。设 $X$ 为整概形，并且
$n = \dim_\delta(X)$。设 $i : D \to X$ 是有效 Cartier 除子。
设 $\mathcal{N}$ 是可逆 $\mathcal{O}_X$-模，而 $t$ 是
$\mathcal{N}$ 的非零亚纯截面。则
$i^*\text{div}_\mathcal{N}(t) = c_1(\mathcal{N}|_D) \cap [D]_{n - 1}$
在 $\CH_{n - 2}(D)$ 中成立。
\end{lemma}

\begin{proof}
对某些 $\delta$-维数为 $n - 1$ 的整闭子概形 $Z_i \subset X$，写
$\text{div}_\mathcal{N}(t) = \sum \text{ord}_{Z_i, \mathcal{N}}(t)[Z_i]$。
可以假设族 $\{Z_i\}$ 局部有限；在
$U = X \setminus \bigcup Z_i$ 上，$t \in \Gamma(U, \mathcal{N}|_U)$
是生成元；而且 $D$ 的每个不可约分支都是某个 $Z_i$。见
《除子》引理 \ref{divisors-lemma-components-locally-finite}、
\ref{divisors-lemma-divisor-locally-finite} 和
\ref{divisors-lemma-divisor-meromorphic-locally-finite}。

\medskip\noindent
置 $\mathcal{L} = \mathcal{O}_X(D)$。以
$s \in \Gamma(X, \mathcal{O}_X(D)) = \Gamma(X, \mathcal{L})$
表示典范截面。将把第 \ref{chow-section-key} 节中的讨论应用于当前情形。
对每个 $i$，设 $\xi_i \in Z_i$ 是其一般点，并置
$B_i = \mathcal{O}_{X, \xi_i}$。对每个 $i$，选取 $B_i$ 上的生成元
$s_i \in \mathcal{L}_{\xi_i}$ 和 $t_i \in \mathcal{N}_{\xi_i}$；
但若 $Z_i \not \subset D$，则坚持选取 $s_i = s$。
写 $s = f_i s_i$ 及 $t = g_i t_i$，其中 $f_i, g_i \in B_i$。
于是 $\text{ord}_{Z_i, \mathcal{N}}(t) = \text{ord}_{B_i}(g_i)$。
另一方面，$f_i \in B_i$，并且由于对 $s_i$ 的选取，有
$$
[D]_{n - 1} = \sum \text{ord}_{B_i}(f_i)[Z_i]
$$
我们断言循环等式
$$
i^*\text{div}_\mathcal{N}(t) =
\sum \text{ord}_{B_i}(g_i) \text{div}_{\mathcal{L}|_{Z_i}}(s_i|_{Z_i})
$$
成立。更准确地说，右边是表示左边的循环。事实上，这由
$\text{div}_\mathcal{N}(t)$ 的公式及如下事实直接得到：当
$Z_i \not \subset D$ 时，
$\text{div}_{\mathcal{L}|_{Z_i}}(s_i|_{Z_i}) = [Z(s_i|_{Z_i})]_{n - 2} =
[Z_i \cap D]_{n - 2}$，因为在这种情形，
$s_i|_{Z_i} = s|_{Z_i}$ 是正则截面；见引理 \ref{chow-lemma-compute-c1}。类似地，
$$
c_1(\mathcal{N}) \cap [D]_{n - 1} =
\sum \text{ord}_{B_i}(f_i) \text{div}_{\mathcal{N}|_{Z_i}}(t_i|_{Z_i})
$$
关键公式（引理 \ref{chow-lemma-key-formula}）给出循环等式
$$
\sum \left(
\text{ord}_{B_i}(f_i) \text{div}_{\mathcal{N}|_{Z_i}}(t_i|_{Z_i}) -
\text{ord}_{B_i}(g_i) \text{div}_{\mathcal{L}|_{Z_i}}(s_i|_{Z_i}) \right) =
\sum \text{div}_{Z_i}(\partial_{B_i}(f_i, g_i))
$$
若 $Z_i \not \subset D$，则 $f_i = 1$，从而
$\text{div}_{Z_i}(\partial_{B_i}(f_i, g_i)) = 0$。因此，在 $D$ 上，
表示 $i^*\text{div}_\mathcal{N}(t)$ 和
$c_1(\mathcal{N}) \cap [D]_{n - 1}$ 的上述具体循环有理等价。
证明完毕。
\end{proof}

\begin{lemma}
\label{chow-lemma-gysin-factors}
设 $(S, \delta)$ 如情形 \ref{chow-situation-setup}。
设 $X$ 在 $S$ 上局部有限型。设
$(\mathcal{L}, s, i : D \to X)$ 如定义 \ref{chow-definition-gysin-homomorphism}。Gysin 同态下降到有理等价类，
从而给出映射 $i^* : \CH_{k + 1}(X) \to \CH_k(D)$。
\end{lemma}

\begin{proof}
设 $\alpha \in Z_{k + 1}(X)$，并假设 $\alpha \sim_{rat} 0$。
这意味着存在一族局部有限的整闭子概形 $W_j \subset X$，
其 $\delta$-维数为 $k + 2$；并存在 $f_j \in R(W_j)^*$，使得
$\alpha = \sum i_{j, *}\text{div}_{W_j}(f_j)$。
置 $X' = \coprod W_i$，并考虑注 \ref{chow-remark-pullback-pairs} 中的图表
$$
\xymatrix{
D' \ar[d]_q \ar[r]_{i'} & X' \ar[d]^p \\
D \ar[r]^i & X
}
$$
由于 $X' \to X$ 固有，由引理 \ref{chow-lemma-closed-in-X-gysin} 有
$i^*p_* = q_*(i')^*$。又因 $q_*$ 下降到有理等价类
（引理 \ref{chow-lemma-proper-pushforward-rational-equivalence}），
只须对 $X'$ 上的 $\alpha' = \sum \text{div}_{W_j}(f_j)$ 证明结论。
显然，这把问题约化到 $X$ 为整概形，且对某个 $f \in R(X)^*$ 有
$\alpha = \text{div}(f)$ 的情形。

\medskip\noindent
假设 $X$ 是整概形，且对某个 $f \in R(X)^*$ 有
$\alpha = \text{div}(f)$。若 $X = D$，则 $i^*\alpha$ 等于
$c_1(\mathcal{L}) \cap \alpha$，而由引理 \ref{chow-lemma-factors}，
它有理等价于零。若 $D \not = X$，则由引理 \ref{chow-lemma-gysin-factors-general}，$i^*\text{div}_X(f)$ 在
$\CH_{n - 2}(D)$ 中等于
$c_1(\mathcal{O}_D) \cap [D]_{n - 1}$。当然，与
$c_1(\mathcal{O}_D)$ 作帽积是零映射
（引理 \ref{chow-lemma-c1-cap-additive}）。
\end{proof}

\begin{lemma}
\label{chow-lemma-gysin-back}
设 $(S, \delta)$ 如情形 \ref{chow-situation-setup}，并设 $X$ 在 $S$ 上
局部有限型。设 $(\mathcal{L}, s, i : D \to X)$ 如定义 \ref{chow-definition-gysin-homomorphism}。则
$i^*i_* : \CH_k(D) \to \CH_{k - 1}(D)$ 把 $\alpha$ 映为
$c_1(\mathcal{L}|_D) \cap \alpha$。
\end{lemma}

\begin{proof}
这立即由循环上 $i_*$ 的定义以及定义 \ref{chow-definition-gysin-homomorphism} 中 $i^*$ 的定义得出。
\end{proof}

\begin{lemma}
\label{chow-lemma-gysin-commutes-cap-c1}
设 $(S, \delta)$ 如情形 \ref{chow-situation-setup}，并设 $X$ 在 $S$ 上
局部有限型。设 $(\mathcal{L}, s, i : D \to X)$ 是定义 \ref{chow-definition-gysin-homomorphism} 中的三元组。
设 $\mathcal{N}$ 是可逆 $\mathcal{O}_X$-模。则对所有
$\alpha \in \CH_k(X)$，在 $\CH_{k - 2}(D)$ 中有
$i^*(c_1(\mathcal{N}) \cap \alpha) = c_1(i^*\mathcal{N}) \cap i^*\alpha$。
\end{lemma}

\begin{proof}
采用与引理 \ref{chow-lemma-gysin-factors} 完全相同的证明，
结论由引理 \ref{chow-lemma-pushforward-cap-c1}、
\ref{chow-lemma-cap-commutative} 和
\ref{chow-lemma-gysin-factors-general} 得出。
\end{proof}

\begin{lemma}
\label{chow-lemma-gysin-commutes-gysin}
设 $(S, \delta)$ 如情形 \ref{chow-situation-setup}，并设 $X$ 在 $S$ 上
局部有限型。设 $(\mathcal{L}, s, i : D \to X)$ 和
$(\mathcal{L}', s', i' : D' \to X)$ 是定义 \ref{chow-definition-gysin-homomorphism} 中的两个三元组。则图表
$$
\xymatrix{
\CH_k(X) \ar[r]_{i^*} \ar[d]_{(i')^*} & \CH_{k - 1}(D) \ar[d]^{j^*} \\
\CH_{k - 1}(D') \ar[r]^{(j')^*} & \CH_{k - 2}(D \cap D')
}
$$
交换，其中每个映射都是 Gysin 映射。
\end{lemma}

\begin{proof}
以 $j : D \cap D' \to D$ 和 $j' : D \cap D' \to D'$ 表示分别对应于
$(\mathcal{L}|_{D'}, s|_{D'}$ 和 $(\mathcal{L}'_D, s|_D)$ 的闭浸入。
需要证明对所有 $\alpha \in \CH_k(X)$ 都有
$(j')^*i^*\alpha = j^* (i')^*\alpha$。
设 $W \subset X$ 是维数为 $k$ 的整闭子概形。先在
$\alpha = [W]$ 的情形证明等式，并由关键公式推出它。

\medskip\noindent
令 $\sigma$ 为 $\mathcal{L}|_W$ 的一个非零亚纯截面；若
$W \not \subset D$，则要求它等于 $s|_W$。
令 $\sigma'$ 为 $\mathcal{L}'|_W$ 的一个非零亚纯截面；若
$W \not \subset D'$，则要求它等于 $s'|_W$。写
$$
\text{div}_{\mathcal{L}|_W}(\sigma) =
\sum \text{ord}_{Z_i, \mathcal{L}|_W}(\sigma)[Z_i] = \sum n_i[Z_i]
$$
以及类似的等式
$$
\text{div}_{\mathcal{L}'|_W}(\sigma') =
\sum \text{ord}_{Z_i, \mathcal{L}'|_W}(\sigma')[Z_i] = \sum n'_i[Z_i]
$$
如第 \ref{chow-section-key} 节的讨论。于是，若 $n_i \not = 0$，
则 $Z_i \subset D$；若 $n'_i \not = 0$，则 $Z'_i \subset D'$。
对每个 $i$，设 $\xi_i \in Z_i$ 为一般点。与第 \ref{chow-section-key} 节中一样，对每个 $i$ 选取元素
$\sigma_i \in \mathcal{L}_{\xi_i}$，以及相应的\ 
$\sigma'_i \in \mathcal{L}'_{\xi_i}$；它们分别生成
$B_i = \mathcal{O}_{W, \xi_i}$ 上的相应模，并且若
$Z_i \not \subset D$，则前者等于 $s$ 的像；\ 相应地，若
$Z_i \not \subset D'$，则\ 后者等于 $s'$ 的像。
写 $\sigma = f_i \sigma_i$ 及 $\sigma' = f'_i\sigma'_i$，使得
$n_i = \text{ord}_{B_i}(f_i)$ 且
$n'_i = \text{ord}_{B_i}(f'_i)$。由定义，作为循环有
$$
(j')^*i^*[W] =
\sum \text{ord}_{B_i}(f_i) \text{div}_{\mathcal{L}'|_{Z_i}}(\sigma'_i|_{Z_i})
$$
以及
$$
j^*(i')^*[W] =
\sum \text{ord}_{B_i}(f'_i) \text{div}_{\mathcal{L}|_{Z_i}}(\sigma_i|_{Z_i})
$$
此时关键公式（引理 \ref{chow-lemma-key-formula}）给出循环等式
$$
\sum \left(
\text{ord}_{B_i}(f_i) \text{div}_{\mathcal{L}'|_{Z_i}}(\sigma'_i|_{Z_i}) -
\text{ord}_{B_i}(f'_i) \text{div}_{\mathcal{L}|_{Z_i}}(\sigma_i|_{Z_i})
\right) =
\sum \text{div}_{Z_i}(\partial_{B_i}(f_i, f'_i))
$$
注意，若 $Z_i \not \subset D \cap D'$，则 $f_i = 1$ 或
$f'_i = 1$，从而
$\text{div}_{Z_i}(\partial_{B_i}(f_i, f'_i)) = 0$。
因此，在 $D \cap D' \cap W$ 上，表示
$(j')^*i^*[W]$ 与 $j^*(i')^*[W]$ 的上述具体循环有理等价。
由注 \ref{chow-remark-infinite-sums-rational-equivalences}，
对一般的 $\alpha$ 也得到结论。
\end{proof}
















\section{相对有效 Cartier 除子}
\label{chow-section-relative-effective-cartier}

\noindent
相对有效 Cartier 除子在《除子》第 \ref{divisors-section-effective-Cartier-morphisms} 节中定义并研究。
为建立向量丛陈类的基本结果，只需考虑周围概形和有效 Cartier 除子
都在基底上平坦的情形。

\begin{lemma}
\label{chow-lemma-relative-effective-cartier}
设 $(S, \delta)$ 如情形 \ref{chow-situation-setup}。
设 $X$、$Y$ 在 $S$ 上局部有限型。设 $p : X \to Y$ 是相对维数为
$r$ 的平坦态射。设 $i : D \to X$ 是相对有效 Cartier 除子
（《除子》定义 \ref{divisors-definition-relative-effective-Cartier-divisor}）。
置 $\mathcal{L} = \mathcal{O}_X(D)$。对任意
$\alpha \in \CH_{k + 1}(Y)$，在 $\CH_{k + r}(D)$ 中有
$$
i^*p^*\alpha = (p|_D)^*\alpha
$$
并且在 $\CH_{k + r}(X)$ 中有
$$
c_1(\mathcal{L}) \cap p^*\alpha = i_* ((p|_D)^*\alpha)
$$
\end{lemma}

\begin{proof}
设 $W \subset Y$ 是 $\delta$-维数为 $k + 1$ 的整闭子概形。
由《除子》引理 \ref{divisors-lemma-relative-Cartier}，
$D \cap p^{-1}W$ 是 $p^{-1}W$ 上的有效 Cartier 除子。
由引理 \ref{chow-lemma-easy-gysin}，下式中的第一个等式成立：
$$
i^*[p^{-1}W]_{k + r + 1} =
[D \cap p^{-1}W]_{k + r} =
[(p|_D)^{-1}(W)]_{k + r}.
$$
第二个等式成立，是因为作为概形有
$D \cap p^{-1}(W) = (p|_D)^{-1}(W)$。
又因定义给出 $p^*[W] = [p^{-1}W]_{k + r + 1}$，可知作为循环有
$i^*p^*[W] = (p|_D)^*[W]$。若
$\alpha = \sum m_j[W_j]$ 是一般的 $k + 1$ 循环，则作为循环有
$i^*\alpha = \sum m_j i^*p^*[W_j] = \sum m_j(p|_D)^*[W_j]$。
这就证明了第一个等式。再应用引理 \ref{chow-lemma-support-cap-effective-Cartier}，即可由第一个等式推出第二个。
\end{proof}








\section{仿射丛}
\label{chow-section-affine-vector}

\noindent
对仿射丛而言，Chow 群上的拉回映射是满射。

\begin{lemma}
\label{chow-lemma-pullback-affine-fibres-surjective}
设 $(S, \delta)$ 如情形 \ref{chow-situation-setup}。
设 $X$、$Y$ 在 $S$ 上局部有限型。设 $f : X \to Y$ 是相对维数为
$r$ 的平坦态射。假设对每个 $y \in Y$，都存在开邻域
$U \subset Y$，使得 $f|_{f^{-1}(U)} : f^{-1}(U) \to U$
可与态射 $U \times \mathbf{A}^r \to U$ 等同。则对所有
$k \in \mathbf{Z}$，映射
$f^* : \CH_k(Y) \to \CH_{k + r}(X)$ 都是满射。
\end{lemma}

\begin{proof}
设 $\alpha \in \CH_{k + r}(X)$。写
$\alpha = \sum m_j[W_j]$，其中 $m_j \not = 0$，而 $W_j$
是两两不同、$\delta$-维数为 $k + r$ 的整闭子概形。
于是族 $\{W_j\}$ 在 $X$ 中局部有限。对任意拟紧开集
$V \subset Y$，只有有限多个 $j$ 使
$f^{-1}(V) \cap W_j$ 非空。因此，像的闭包所成的族
$Z_j = \overline{f(W_j)}$ 是 $Y$ 中局部有限的一族整闭子概形。

\medskip\noindent
考虑纤维积图表
$$
\xymatrix{
f^{-1}(Z_j) \ar[r] \ar[d]_{f_j} & X \ar[d]^f \\
Z_j \ar[r] & Y
}
$$
假设对某个 $k$-循环 $\beta_j \in \CH_k(Z_j)$，
$[W_j] \in Z_{k + r}(f^{-1}(Z_j))$ 与 $f_j^*\beta_j$ 有理等价。
那么 $\beta = \sum m_j \beta_j$ 是 $Y$ 上的 $k$-循环，而
$f^*\beta = \sum m_j f_j^*\beta_j$ 与 $\alpha$ 有理等价
（见注 \ref{chow-remark-infinite-sums-rational-equivalences}）。
这把问题约化到如下情形：$Y$ 是整概形，并且
$\alpha = [W]$，其中所指的是支配 $Y$ 的 $X$ 中整闭子概形。
特别地，可以假设 $d = \dim_\delta(Y) < \infty$。

\medskip\noindent
于是可以对 $d = \dim_\delta(Y)$ 作归纳。若 $d < k$，则
$\CH_{k + r}(X) = 0$，引理成立。由假设，存在稠密开集
$V \subset Y$，使得作为 $V$ 上的概形有
$f^{-1}(V) \cong V \times \mathbf{A}^r$。
假设能够证明对某个 $\beta \in Z_k(V)$ 有
$\alpha|_{f^{-1}(V)} = f^*\beta$。由引理 \ref{chow-lemma-exact-sequence-open}，对某个 $\beta' \in Z_k(Y)$ 有
$\beta = \beta'|_V$。由引理 \ref{chow-lemma-restrict-to-open} 的正合列
$\CH_k(f^{-1}(Y \setminus V)) \to \CH_k(X) \to \CH_k(f^{-1}(V))$，
可知 $\alpha - f^*\beta'$ 来自某个循环
$\alpha' \in \CH_{k + r}(f^{-1}(Y \setminus V))$。
由于 $\dim_\delta(Y \setminus V) < d$，对 $d$ 使用归纳即可。

\medskip\noindent
所以可以假设 $X = Y \times \mathbf{A}^r$。此时可以把 $f$ 分解为
$$
X = Y \times \mathbf{A}^r \to
Y \times \mathbf{A}^{r - 1} \to \ldots \to
Y \times \mathbf{A}^1 \to Y.
$$
故只须处理 $r = 1$ 的情形。由证明第二段中的论证，问题约化到
$\alpha = [W]$、$Y$ 为整概形且 $W \to Y$ 支配的情形。
仍可对 $d = \dim_\delta(Y)$ 作归纳。若
$W = Y \times \mathbf{A}^1$，则 $[W] = f^*[Y]$。
最后，若 $W \subset Y \times \mathbf{A}^1$ 是真包含，则
$W \to Y$ 诱导有限域扩张 $R(W)/R(Y)$。
设 $P(T) \in R(Y)[T]$ 是首一不可约多项式，使得 $W \to Y$
的一般纤维在 $\mathbf{A}^1_{R(Y)}$ 中由 $P$ 截出。
选取非空开集 $V \subset Y$，使 $P \in \Gamma(V, \mathcal{O}_Y)[T]$，
并且 $W \cap f^{-1}(V)$ 仍由 $P$ 截出。于是
$\alpha|_{f^{-1}(V)} \sim_{rat} 0$，故对
$(Y \setminus V) \times \mathbf{A}^1$ 上某个循环 $\alpha'$ 有
$\alpha \sim_{rat} \alpha'$。对维数应用归纳即可。
\end{proof}

\begin{lemma}
\label{chow-lemma-linebundle}
设 $(S, \delta)$ 如情形 \ref{chow-situation-setup}。
设 $X$ 在 $S$ 上局部有限型，$\mathcal{L}$ 是可逆
$\mathcal{O}_X$-模。设
$$
p :
L = \underline{\Spec}(\text{Sym}^*(\mathcal{L}))
\longrightarrow
X
$$
是 $X$ 上相应的向量丛。则对所有 $k$，
$p^* : \CH_k(X) \to \CH_{k + 1}(L)$ 都是同构。
\end{lemma}

\begin{proof}
满射性见引理 \ref{chow-lemma-pullback-affine-fibres-surjective}。
设 $o : X \to L$ 是 $L \to X$ 的零截面，即对应于满射
$\text{Sym}^*(\mathcal{L}) \to \mathcal{O}_X$ 的态射；该满射对所有
$n > 0$ 都把 $\mathcal{L}^{\otimes n}$ 映为零。
于是 $p \circ o = \text{id}_X$，且 $o(X)$ 是 $L$ 上的有效
Cartier 除子。故由引理 \ref{chow-lemma-relative-effective-cartier}，
$o^* \circ p^* = \text{id}$，从而 $p^*$ 也是单射。
\end{proof}

\begin{remark}
\label{chow-remark-when-isomorphism}
后文将看到（引理 \ref{chow-lemma-vectorbundle}），若 $X$ 是 $Y$ 上秩为
$r$ 的向量丛，则拉回映射 $\CH_k(Y) \to \CH_{k + r}(X)$ 是同构。
只要 $X \to Y$ 满足引理 \ref{chow-lemma-pullback-affine-fibres-surjective} 的假设，这一结论就成立；
见 \cite[Lemma 2.2]{Totaro-group}。在注 \ref{chow-remark-higher-chow-isomorphism} 中，我们将用高阶 Chow 群勾勒一个证明。
\end{remark}

\begin{lemma}
\label{chow-lemma-linebundle-formulae}
在引理 \ref{chow-lemma-linebundle} 的情形中，以 $o : X \to L$
表示零截面（见该引理的证明）。则有
\begin{enumerate}
\item $o(X)$ 是 $p^*\mathcal{L}^{\otimes -1}$ 的某个正则整体截面的
零点概形；
\item 因为 $o$ 是闭浸入，有 $o_* : \CH_k(X) \to \CH_k(L)$；
\item 因为 $o(X)$ 是有效 Cartier 除子，有
$o^* : \CH_{k + 1}(L) \to \CH_k(X)$；
\item $o^* p^* : \CH_k(X) \to \CH_k(X)$ 是恒等映射；
\item 对任意 $\alpha \in \CH_k(X)$，有
$o_*\alpha = - p^*(c_1(\mathcal{L}) \cap \alpha)$；以及
\item $o^* o_* : \CH_k(X) \to \CH_{k - 1}(X)$ 等于映射
$\alpha \mapsto - c_1(\mathcal{L}) \cap \alpha$。
\end{enumerate}
\end{lemma}

\begin{proof}
由于 $p_*\mathcal{O}_L = \text{Sym}^*(\mathcal{L})$，投影公式
（《层上同调》引理 \ref{cohomology-lemma-projection-formula}）给出
$p_*(p^*\mathcal{L}^{\otimes -1}) =
\text{Sym}^*(\mathcal{L}) \otimes_{\mathcal{O}_X} \mathcal{L}^{\otimes -1}$，
而 (1) 中提到的截面是典范平凡化
$\mathcal{O}_X \to
\mathcal{L} \otimes_{\mathcal{O}_X} \mathcal{L}^{\otimes -1}$。
这里略去证明该截面的零点轨迹恰为 $o(X)$。这证明了 (1)。

\medskip\noindent
(2)、(3) 和 (4) 已在引理 \ref{chow-lemma-linebundle} 的证明过程中看到。
当然，(4) 正是引理 \ref{chow-lemma-relative-effective-cartier} 中的第一个公式。

\medskip\noindent
(5) 由引理 \ref{chow-lemma-relative-effective-cartier} 中的第二个公式、
与 $c_1$ 作帽积的可加性（引理 \ref{chow-lemma-c1-cap-additive}），
以及与 $c_1$ 作帽积和平坦拉回交换这一事实
（引理 \ref{chow-lemma-flat-pullback-cap-c1}）得出。

\medskip\noindent
(6) 由引理 \ref{chow-lemma-gysin-back} 以及
$o^*p^*\mathcal{L} = \mathcal{L}$ 得出。
\end{proof}

\begin{lemma}
\label{chow-lemma-decompose-section}
设 $Y$ 是概形。设 $\mathcal{L}_i$（$i = 1, 2$）是可逆
$\mathcal{O}_Y$-模。设 $s$ 是
$\mathcal{L}_1 \otimes_{\mathcal{O}_Y} \mathcal{L}_2$ 的整体截面，
并以 $i : D \to Y$ 表示 $s$ 的零点概形。则存在交换图
$$
\xymatrix{
D_1 \ar[r]_{i_1} \ar[d]_{p_1} &
L \ar[d]^p &
D_2 \ar[l]^{i_2} \ar[d]^{p_2} \\
D \ar[r]^i &
Y &
D \ar[l]_i
}
$$
以及 $p^*\mathcal{L}_i$ 的截面 $s_i$，使得：
\begin{enumerate}
\item $p^*s = s_1 \otimes s_2$；
\item $p$ 有限型、平坦且相对维数为 $1$；
\item $D_i$ 是 $s_i$ 的零点概形；
\item 对 $i = 1, 2$，在 $D$ 上有
$D_i \cong
\underline{\Spec}(\text{Sym}^*(\mathcal{L}_{3 - i}^{\otimes -1})|_D))$；
\item $p^{-1}D = D_1 \cup D_2$（概形论并）；
\item $D_1 \cap D_2$（概形论交）同构地映到 $D$；以及
\item 对 $i = 1, 2$，$D_1 \cap D_2 \to D_i$ 是线丛
$D_i \to D$ 的零截面。
\end{enumerate}
此外，该图表和各截面 $s_i$ 的形成与任意基变换交换。
\end{lemma}

\begin{proof}
设 $p : X \to Y$ 是如下拟凝聚 $\mathcal{O}_Y$-代数层的相对谱：
$$
\mathcal{A} =
\left(\bigoplus\nolimits_{a_1, a_2 \geq 0}
\mathcal{L}_1^{\otimes -a_1} \otimes_{\mathcal{O}_Y}
\mathcal{L}_2^{\otimes -a_2}\right)/\mathcal{J}
$$
其中 $\mathcal{J}$ 是由形如 $st - t$ 的局部截面生成的理想；
这里 $t$ 是满足 $a_1, a_2 > 0$ 的任意直和项
$\mathcal{L}_1^{\otimes -a_1} \otimes \mathcal{L}_2^{\otimes -a_2}$
的局部截面。把截面 $s_i$ 看作映射
$p^*\mathcal{L}_i^{\otimes -1} \to \mathcal{O}_X$，则它们定义为映射
$\mathcal{L}_i^{\otimes -1} \to \mathcal{A} = p_*\mathcal{O}_X$
的伴随映射。对任意 $y \in Y$，可以选取包含 $y$ 的仿射开集
$V \subset Y$，设 $V = \Spec(A)$，并选取平凡化
$z_i : \mathcal{O}_V \to \mathcal{L}_i^{\otimes -1}|_V$。
注意，$f = s(z_1z_2) \in A$ 截出闭子概形 $D$。显然
$$
p^{-1}(V) = \Spec(B[z_1, z_2]/(z_1 z_2 - f))
$$
由于 $D_i$ 由 $z_i$ 截出，其余各点都很清楚。
\end{proof}

\begin{lemma}
\label{chow-lemma-decompose-section-formulae}
在引理 \ref{chow-lemma-decompose-section} 的情形中，假设 $Y$ 是
情形 \ref{chow-situation-setup} 中那样在 $(S, \delta)$ 上局部有限型。
则对所有 $\alpha \in \CH_k(Y)$，在 $\CH_k(D_1)$ 中有
$i_1^*p^*\alpha = p_1^*i^*\alpha$。
\end{lemma}

\begin{proof}
设 $W \subset Y$ 是 $\delta$-维数为 $k$ 的整闭子概形。
分两种情形讨论。

\medskip\noindent
假设 $W \subset D$。由 Gysin 同态的定义以及引理 \ref{chow-lemma-c1-cap-additive} 的可加性，在 $\CH_{k - 1}(D)$ 中有
$i^*[W] = c_1(\mathcal{L}_1) \cap [W] + c_1(\mathcal{L}_2) \cap [W]$。
因此
$p_1^*i^*[W] =
p_1^*(c_1(\mathcal{L}_1) \cap [W]) + p_1^*(c_1(\mathcal{L}_2) \cap [W])$。
另一方面，平坦拉回的构造给出
$p^*[W] = [p^{-1}(W)]_{k + 1}$。由该引理
$p^{-1}(W) = W_1 \cup W_2$（概形论地），其中
$W_i = p_i^{-1}(W)$ 是 $W$ 上的线丛
（因为图表的形成与基变换交换）。于是
$[p^{-1}(W)]_{k + 1} = [W_1] + [W_2]$，因为 $W_i$ 是 $L$ 中
$\delta$-维数为 $k + 1$ 的整闭子概形。因此
\begin{align*}
i_1^*p^*[W]
& =
i_1^*[p^{-1}(W)]_{k + 1} \\
& =
i_1^*([W_1] + [W_2]) \\
& =
c_1(p_1^*\mathcal{L}_1) \cap [W_1] + [W_1 \cap W_2]_k \\
& =
c_1(p_1^*\mathcal{L}_1) \cap p_1^*[W] + [W_1 \cap W_2]_k \\
& =
p_1^*(c_1(\mathcal{L}_1) \cap [W]) + [W_1 \cap W_2]_k
\end{align*}
这里使用了 Gysin 同态的构造、平坦拉回的定义（用于第二个等式），
以及 $c_1 \cap -$ 与平坦拉回的相容性
（引理 \ref{chow-lemma-flat-pullback-cap-c1}）。
由于 $W_1 \cap W_2$ 是线丛 $W_1 \to W$ 的零截面，引理 \ref{chow-lemma-linebundle-formulae} 给出
$[W_1 \cap W_2]_k = p_1^*(c_1(\mathcal{L}_2) \cap [W])$。
注意，这里使用了 $D_1$ 是 $\mathcal{L}_2$ 的逆的相对谱所给出的线丛
这一事实。因此所得结果与前式相同。

\medskip\noindent
假设 $W \not \subset D$。此时，$i_1^*p^*[W]$ 与
$p_1^*i^*[W]$ 都由 $W$ 在 $D_1$ 中的概形论逆像所关联的
$k - 1$ 循环表示。
\end{proof}

\begin{lemma}
\label{chow-lemma-normal-cone-effective-Cartier}
在情形 \ref{chow-situation-setup} 中，设 $X$ 是 $S$ 上局部有限型的概形。
设 $(\mathcal{L}, s, i : D \to X)$ 是定义 \ref{chow-definition-gysin-homomorphism} 中的三元组。则存在交换图
$$
\xymatrix{
D' \ar[r]_{i'} \ar[d]_p & X' \ar[d]^g \\
D \ar[r]^i & X
}
$$
满足
\begin{enumerate}
\item $p$ 和 $g$ 有限型、平坦且相对维数为 $1$；
\item 对所有 $k$，$p^* : \CH_k(D) \to \CH_{k + 1}(D')$ 都是单射；
\item $D' \subset X'$ 是某个整体截面
$s' \in \Gamma(X', \mathcal{O}_{X'})$ 的零点概形；
\item 作为映射 $\CH_k(X) \to \CH_k(D')$，有
$p^*i^* = (i')^*g^*$。
\end{enumerate}
此外，经由局部有限型态射 $Y \to X$ 作任意基变换后，这些性质仍成立。
\end{lemma}

\begin{proof}
注意，$(i')^*$ 已有定义，因为存在定义 \ref{chow-definition-gysin-homomorphism} 中那样的三元组
$(\mathcal{O}_{X'}, s', i' : D' \to X')$。因此该陈述有意义。

\medskip\noindent
置 $\mathcal{L}_1 = \mathcal{O}_X$、$\mathcal{L}_2 = \mathcal{L}$，
并把引理 \ref{chow-lemma-decompose-section} 应用于
$\mathcal{L} = \mathcal{L}_1 \otimes_{\mathcal{O}_X} \mathcal{L}_2$
的截面 $s$。取 $D' = D_1$。此时结论由该引理和引理 \ref{chow-lemma-decompose-section-formulae} 得出，而单射性由引理 \ref{chow-lemma-linebundle} 得出。
\end{proof}

\begin{remark}
\label{chow-remark-higher-chow-isomorphism}
设 $(S, \delta)$ 如情形 \ref{chow-situation-setup}。
设 $Y$ 在 $S$ 上局部有限型，$r \geq 0$，而
$f : X \to Y$ 是概形态射。假设每个 $y \in Y$ 都包含于某个开集
$V \subset Y$，并且作为 $V$ 上的概形有
$f^{-1}(V) \cong V \times \mathbf{A}^r$。
本注记勾勒事实
$f^* : \CH_k(Y) \to \CH_{k + r}(X)$ 为同构的一个证明。
首先，由引理 \ref{chow-lemma-pullback-affine-fibres-surjective}，
该映射是满射。设 $\alpha \in \CH_k(Y)$ 且 $f^*\alpha = 0$。
下面证明 $\alpha = 0$。

\medskip\noindent
第 1 步。可以假设 $\dim_\delta(Y) < \infty$。（实际遇到的所有情形中，
这一点都是立即成立的，所以建议读者略过本步。）事实上，在 $X$ 上见证
$f^*\alpha = 0$ 的任一有理等价会用到一族局部有限、维数为
$k + r + 1$ 的整闭子概形。取它们在 $Y$ 中各像的闭包之并，
得到满足 $\dim_\delta(Y') \leq k + r + 1$ 的闭子概形
$Y' \subset Y$，使得 $\alpha$ 是某个 $\alpha' \in \CH_k(Y')$ 的像，
并且 $(f')^*\alpha = 0$，其中 $f'$ 是 $f$ 到 $Y'$ 的基变换。

\medskip\noindent
第 2 步。假设 $d = \dim_\delta(Y) < \infty$，并对 $d$ 作归纳。
若 $d < k$，则 $\alpha = 0$，结论成立；这是归纳的初始情形。
一般地，对 $f$ 的假设表明，可以选取稠密开集
$V \subset Y$，使 $U = f^{-1}(V) = \mathbf{A}^r_V$。
以 $Y' \subset Y$ 表示 $V$ 的补集，视为既约闭子概形，并置
$X' = f^{-1}(Y')$。考虑
$$
\xymatrix{
\CH^M_{k + r}(U, 1) \ar[r] &
\CH_{k + r}(X') \ar[r] &
\CH_{k + r}(X) \ar[r] &
\CH_{k + r}(U) \ar[r] &
0 \\
\CH^M_k(V, 1) \ar[r] \ar[u] &
\CH_k(Y') \ar[r] \ar[u] &
\CH_k(Y) \ar[r] \ar[u] &
\CH_k(V) \ar[r] \ar[u] &
0
}
$$
这里使用了 $V$ 和 $U$ 的第一高阶 Chow 群、注 \ref{chow-remark-higher-chow} 中构造的六项正合列，以及这些高阶 Chow 群
上的平坦拉回和平坦拉回与六项正合列的相容性。
由于 $U = \mathbf{A}^r_V$，右侧竖直映射是同构。
根据对 $d$ 的归纳，映射 $\CH_k(Y') \to \CH_{k + r}(X')$ 是双射。
所以只须证明
$$
\CH^M_k(V, 1) \longrightarrow \CH^M_{k + r}(U, 1)
$$
为满射即可完成论证。仿照引理 \ref{chow-lemma-pullback-affine-fibres-surjective} 的证明进行论证，
这又约化到下面的第 3 步。

\medskip\noindent
第 3 步。设 $F$ 是域。则 $\CH^M_0(\mathbf{A}^1_F, 1) = 0$。
（在上述引理的证明中，我们类似地证明了
$\CH_0(\mathbf{A}^1_F) = 0$。）有
$$
\CH^M_0(\mathbf{A}^1_F, 1) = \Coker\left(
\partial : K^M_2(F(T)) \longrightarrow
\bigoplus\nolimits_{\mathfrak p \subset F[T]\text{ 为极大理想}}
\kappa(\mathfrak p)^*\right)
$$
证明该余核为零的经典论证，是对 $\kappa(\mathfrak p)/F$ 的次数作归纳，
证明对应于 $\mathfrak p$ 的直和项在像中。
若 $\mathfrak p$ 由首一不可约多项式 $P(T) \in F[T]$ 生成，而
$u \in \kappa(x)^*$ 是某个满足 $\deg(Q) < \deg(P)$ 的
$Q(T) \in F[T]$ 的剩余类，则可以证明 $\partial(Q, P)$ 在
$\mathfrak p$ 处给出元素 $u$，并且可能还在整除 $Q$ 的一些次数更低
的素理想处给出其他单位。这就完成了证明概要。
\end{remark}











\section{双变相交理论}
\label{chow-section-bivariant}

\noindent
为妥善讨论向量丛的高阶陈类，我们依照 \cite{F} 引入双变 Chow 类。
我们的定义与 \cite{F} 有两点不同：(1) 所处的框架不同；(2) 仅要求
双变类与可逆模截面的零点概形所给出的 Gysin 同态交换
（第 \ref{chow-section-intersecting-effective-Cartier} 节）。
后文将在引理 \ref{chow-lemma-gysin-commutes} 中看到，我们的双变类
与所有高余维 Gysin 同态交换，因而满足 \cite{F} 中要求的全部性质；
另见 \cite[Theorem 17.1]{F}。

\begin{definition}
\label{chow-definition-bivariant-class}
\begin{reference}
类似于 \cite[Definition 17.1]{F}
\end{reference}
设 $(S, \delta)$ 如情形 \ref{chow-situation-setup}。
设 $f : X \to Y$ 是 $S$ 上局部有限型概形间的态射，并设
$p \in \mathbf{Z}$。{\it $f$ 的次数为 $p$ 的双变类 $c$}由如下规则给出：
对每个局部有限型态射 $Y' \to Y$ 及每个 $k$，指定映射
$$
c \cap - : \CH_k(Y') \longrightarrow \CH_{k - p}(X')
$$
其中 $X' = Y' \times_Y X$，并满足以下条件：
\begin{enumerate}
\item 若 $Y'' \to Y'$ 是固有态射，则对 $Y''$ 上所有 $\alpha''$，
有 $c \cap (Y'' \to Y')_*\alpha'' = (X'' \to X')_*(c \cap \alpha'')$，
其中 $X'' = Y'' \times_Y X$；
\item 若 $Y'' \to Y'$ 是相对维数固定、局部有限型的平坦态射，
则对 $Y'$ 上所有 $\alpha'$，有
$c \cap (Y'' \to Y')^*\alpha' = (X'' \to X')^*(c \cap \alpha')$；以及
\item 若 $(\mathcal{L}', s', i' : D' \to Y')$ 如定义 \ref{chow-definition-gysin-homomorphism}，且其到 $X'$ 的拉回为
$(\mathcal{N}', t', j' : E' \to X')$，则对 $Y'$ 上所有 $\alpha'$，有
$c \cap (i')^*\alpha' = (j')^*(c \cap \alpha')$。
\end{enumerate}
$f$ 的所有次数为 $p$ 的双变类所成的集合记为 $A^p(X \to Y)$。
\end{definition}

\noindent
设 $(S, \delta)$ 如情形 \ref{chow-situation-setup}。设 $X \to Y$
和 $Y \to Z$ 是 $S$ 上局部有限型概形间的态射，并设
$p \in \mathbf{Z}$。显然，$A^p(X \to Y)$ 是 Abel 群。
此外，显然存在结合的双线性复合
$$
A^p(X \to Y) \times A^q(Y \to Z) \to A^{p + q}(X \to Z)
$$

\begin{lemma}
\label{chow-lemma-flat-pullback-bivariant}
设 $(S, \delta)$ 如情形 \ref{chow-situation-setup}。
设 $f : X \to Y$ 是 $S$ 上局部有限型概形间相对维数为 $r$ 的平坦态射。
对 $Y' \to Y$ 指定
$(f')^* : \CH_k(Y') \to \CH_{k + r}(X')$ 的规则是次数为 $-r$ 的双变类，
其中 $X' = X \times_Y Y'$。
\end{lemma}

\begin{proof}
这由引理 \ref{chow-lemma-flat-pullback-rational-equivalence}、
\ref{chow-lemma-compose-flat-pullback}、
\ref{chow-lemma-flat-pullback-proper-pushforward} 和
\ref{chow-lemma-gysin-flat-pullback} 得出。
\end{proof}

\begin{lemma}
\label{chow-lemma-gysin-bivariant}
设 $(S, \delta)$ 如情形 \ref{chow-situation-setup}。
设 $X$ 在 $S$ 上局部有限型。设
$(\mathcal{L}, s, i : D \to X)$ 是定义 \ref{chow-definition-gysin-homomorphism} 中的三元组。
对 $f : X' \to X$ 指定
$(i')^* : \CH_k(X') \to \CH_{k - 1}(D')$ 的规则是次数为 $1$ 的双变类，
其中 $D' = D \times_X X'$。
\end{lemma}

\begin{proof}
这由引理 \ref{chow-lemma-gysin-factors}、
\ref{chow-lemma-closed-in-X-gysin}、\ref{chow-lemma-gysin-flat-pullback} 和
\ref{chow-lemma-gysin-commutes-gysin} 得出。
\end{proof}

\begin{lemma}
\label{chow-lemma-push-proper-bivariant}
设 $(S, \delta)$ 如情形 \ref{chow-situation-setup}。
设 $f : X \to Y$ 和 $g : Y \to Z$ 是 $S$ 上局部有限型概形间的态射。
设 $c \in A^p(X \to Z)$，并假设 $f$ 固有。
对 $Z' \to Z$ 指定
$\alpha \longmapsto f'_*(c \cap \alpha)$ 的规则是双变类，
记为 $f_* \circ c \in A^p(Y \to Z)$。
\end{lemma}

\begin{proof}
这由引理 \ref{chow-lemma-compose-pushforward}、
\ref{chow-lemma-flat-pullback-proper-pushforward} 和
\ref{chow-lemma-closed-in-X-gysin} 得出。
\end{proof}

\begin{remark}
\label{chow-remark-restriction-bivariant}
设 $(S, \delta)$ 如情形 \ref{chow-situation-setup}。设 $X \to Y$
和 $Y' \to Y$ 是 $S$ 上局部有限型概形间的态射，并置
$X' = Y' \times_Y X$。则有明显的限制映射
$$
A^p(X \to Y) \longrightarrow A^p(X' \to Y'),\quad
c \longmapsto res(c)
$$
其构造如下：把 $Y'$ 上局部有限型的概形 $Y''$ 视为 $Y$ 上局部有限型
的概形，并对任意 $\alpha'' \in \CH_k(Y'')$ 置
$res(c) \cap \alpha'' = c \cap \alpha''$。这个限制运算显然与复合相容。
\end{remark}

\begin{remark}
\label{chow-remark-bivariant-commute}
设 $(S, \delta)$ 如情形 \ref{chow-situation-setup}，并设 $X$ 在 $S$ 上
局部有限型。对 $i = 1, 2$，设 $Z_i \to X$ 是局部有限型概形的态射。
设 $c_i \in A^{p_i}(Z_i \to X)$（$i = 1, 2$）是双变类。
对任意 $\alpha \in \CH_k(X)$，可以问在
$\CH_{k - p_1 - p_2}(Z_1 \times_X Z_2)$ 中是否有
$$
c_1 \cap c_2 \cap \alpha = c_2 \cap c_1 \cap \alpha
$$
若该等式成立，并且经任意局部有限型基变换 $X' \to X$ 后仍成立，
就称 $c_1$ 和 $c_2$ {\it 交换}。当然，这等价于在
$A^{p_1 + p_2}(Z_1 \times_X Z_2 \to X)$ 中有
$$
res(c_1) \circ c_2 = res(c_2) \circ c_1
$$
这里 $res(c_1) \in A^{p_1}(Z_1 \times_X Z_2 \to Z_2)$ 是注 \ref{chow-remark-restriction-bivariant} 中那样的 $c_1$ 的限制；
$res(c_2)$ 类似。
\end{remark}

\begin{example}
\label{chow-example-gysin-commute}
设 $(S, \delta)$ 如情形 \ref{chow-situation-setup}，并设 $X$ 在 $S$ 上
局部有限型。设 $(\mathcal{L}, s, i : D \to X)$ 是定义 \ref{chow-definition-gysin-homomorphism} 中的三元组。
设 $Z \to X$ 是局部有限型概形的态射，而 $c \in A^p(Z \to X)$
是双变类。则引理 \ref{chow-lemma-gysin-bivariant} 的双变 Gysin 类
$c' \in A^1(D \to X)$ 按注 \ref{chow-remark-bivariant-commute} 的意义与 $c$ 交换。
这正是定义 \ref{chow-definition-bivariant-class} 中条件 (3) 的重述。
\end{example}

\begin{remark}
\label{chow-remark-more-general-bivariant}
文献似乎没有考虑如下更一般类型的双变类。设给定图表
$$
X \longrightarrow Z \longleftarrow Y
$$
其中各概形如情形 \ref{chow-situation-setup} 中那样在
$(S, \delta)$ 上局部有限型。设 $p \in \mathbf{Z}$。
可以考虑如下规则 $c$：对每个局部有限型态射 $Z' \to Z$
及所有 $k \in \mathbf{Z}$，指定映射
$$
c \cap - : \CH_k(Y') \longrightarrow \CH_{k - p}(X')
$$
其中 $X' = X \times_Z Z'$ 且 $Y' = Z' \times_Z Y$，并与以下运算相容：
\begin{enumerate}
\item 给定固有态射 $Z'' \to Z'$ 时的固有推前；
\item 给定相对维数固定的平坦态射 $Z'' \to Z'$ 时的平坦拉回；以及
\item 给定定义 \ref{chow-definition-gysin-homomorphism} 中那样的
$D' \subset Z'$ 时的 Gysin 映射。
\end{enumerate}
这里略去详细表述。把所有这类运算所成的集合记为
$A^p(X \to Z \leftarrow Y)$。这个概念有用的一个简单例子，是存在固有态射
$f : X_2 \to X_1$ 的情形。此时，按定义 \ref{chow-definition-bivariant-class} 的意义，$f_*$ 不是双变运算；
但按上述推广意义它是，即
$f_* \in A^0(X_1 \to X_1 \leftarrow X_2)$。
\end{remark}





\section{Chow 上同调与第一陈类}
\label{chow-section-chow-cohomology}

\noindent
我们最关心的是 $A^p(X) = A^p(X \to X)$；它始终指
$\text{id}_X$ 的双变上同调类。陈类正是位于这里。

\begin{definition}
\label{chow-definition-chow-cohomology}
设 $(S, \delta)$ 如情形 \ref{chow-situation-setup}。
设 $X$ 在 $S$ 上局部有限型。$X$ 的{\it Chow 上同调}是分次
$\mathbf{Z}$-代数 $A^*(X)$，其次数 $p$ 分量为 $A^p(X \to X)$。
\end{definition}

\noindent
警告：尚不清楚 $A^*(X)$ 上的 $\mathbf{Z}$-代数结构是否交换；
但我们将看到陈类位于其中心。

\begin{remark}
\label{chow-remark-pullback-cohomology}
设 $(S, \delta)$ 如情形 \ref{chow-situation-setup}。
设 $f : Y' \to Y$ 是 $S$ 上局部有限型概形间的态射。
作为注 \ref{chow-remark-restriction-bivariant} 的一个特例，
存在典范 $\mathbf{Z}$-代数映射 $res : A^*(Y) \to A^*(Y')$。
文献中常把这个映射记为 $f^*$。
\end{remark}

\begin{lemma}
\label{chow-lemma-cap-c1-bivariant}
设 $(S, \delta)$ 如情形 \ref{chow-situation-setup}。
设 $X$ 在 $S$ 上局部有限型，$\mathcal{L}$ 是可逆
$\mathcal{O}_X$-模。对 $f : X' \to X$ 指定
$c_1(f^*\mathcal{L}) \cap - : \CH_k(X') \to \CH_{k - 1}(X')$
的规则是次数为 $1$ 的双变类。
\end{lemma}

\begin{proof}
这由引理 \ref{chow-lemma-factors}、\ref{chow-lemma-pushforward-cap-c1}、
\ref{chow-lemma-flat-pullback-cap-c1} 和
\ref{chow-lemma-gysin-commutes-cap-c1} 得出。
\end{proof}

\noindent
上述引理终于使我们能够作出如下定义。

\begin{definition}
\label{chow-definition-first-chern-class}
设 $(S, \delta)$ 如情形 \ref{chow-situation-setup}。设 $X$ 在 $S$ 上
局部有限型，$\mathcal{L}$ 是可逆 $\mathcal{O}_X$-模。
$\mathcal{L}$ 的{\it 第一陈类}
$c_1(\mathcal{L}) \in A^1(X)$ 是引理 \ref{chow-lemma-cap-c1-bivariant} 的双变类。
\end{definition}

\noindent
有限局部自由模的陈类将在第 \ref{chow-section-intersecting-chern-classes} 节中构造。下面证明
$c_1(\mathcal{L})$ 位于 $A^*(X)$ 的中心。

\begin{lemma}
\label{chow-lemma-c1-center}
设 $(S, \delta)$ 如情形 \ref{chow-situation-setup}。
设 $X$ 在 $S$ 上局部有限型，$\mathcal{L}$ 是可逆
$\mathcal{O}_X$-模。则
\begin{enumerate}
\item $c_1(\mathcal{L}) \in A^1(X)$ 位于 $A^*(X)$ 的中心；以及
\item 若 $f : X' \to X$ 局部有限型，且 $c \in A^*(X' \to X)$，
则 $c \circ c_1(\mathcal{L}) = c_1(f^*\mathcal{L}) \circ c$。
\end{enumerate}
\end{lemma}

\begin{proof}
当然，(2) 蕴含 (1)。设 $p : L \to X$ 如引理 \ref{chow-lemma-linebundle}，并设 $o : X \to L$ 为零截面。
以 $p' : L' \to X'$ 和 $o' : X' \to L'$ 表示它们的基变换。
由引理 \ref{chow-lemma-linebundle-formulae}，有
$$
p^*(c_1(\mathcal{L}) \cap \alpha) = - o_* \alpha
\quad\text{且}\quad
(p')^*(c_1(f^*\mathcal{L}) \cap \alpha') = - o'_* \alpha'
$$
由于 $c$ 是双变类，有
\begin{align*}
(p')^*(c \cap c_1(\mathcal{L}) \cap \alpha)
& =
c \cap p^*(c_1(\mathcal{L}) \cap \alpha) \\
& =
- c \cap o_* \alpha \\
& =
- o'_*(c \cap \alpha) \\
& =
(p')^*(c_1(f^*\mathcal{L}) \cap c \cap \alpha)
\end{align*}
由上述某个引理，$(p')^*$ 是单射，故
$c \cap c_1(\mathcal{L}) \cap \alpha =
c_1(f^*\mathcal{L}) \cap c \cap \alpha$。
经任意局部有限型基变换 $Y \to X$ 后，同一结论仍成立，
所以得到 (2) 所述双变类的等式。
\end{proof}

\begin{lemma}
\label{chow-lemma-vanish-above-dimension}
设 $(S, \delta)$ 如情形 \ref{chow-situation-setup}。
设 $X$ 是 $S$ 上有限型且具有充足可逆层的概形。假设
$d = \dim(X) < \infty$（这里确实指维数，而不是 $\delta$-维数）。
则对 $X$ 上任意可逆层
$\mathcal{L}_1, \ldots, \mathcal{L}_{d + 1}$，在 $A^{d + 1}(X)$ 中有
$c_1(\mathcal{L}_1) \circ \ldots \circ c_1(\mathcal{L}_{d + 1}) = 0$。
\end{lemma}

\begin{proof}
对 $d$ 作归纳。初始情形 $d = 0$ 成立，因为此时 $X$ 是有限个闭点所成的
集合，所以每个可逆模都平凡。假设 $d > 0$。由《除子》引理 \ref{divisors-lemma-quasi-projective-Noetherian-pic-effective-Cartier}，
对某些有效 Cartier 除子 $D, D' \subset X$，可以写成
$\mathcal{L}_{d + 1} \cong \mathcal{O}_X(D) \otimes
\mathcal{O}_X(D')^{\otimes -1}$。于是 $c_1(\mathcal{L}_{d + 1})$ 是
$c_1(\mathcal{O}_X(D))$ 与 $c_1(\mathcal{O}_X(D'))$ 之差，故可以假设
对某个有效 Cartier 除子有
$\mathcal{L}_{d + 1} = \mathcal{O}_X(D)$。

\medskip\noindent
以 $i : D \to X$ 表示包含态射，并以 $i^* \in A^1(D \to X)$
表示引理 \ref{chow-lemma-gysin-bivariant} 中由 Gysin 同态给出的双变类。
由引理 \ref{chow-lemma-support-cap-effective-Cartier}，在 $A^1(X)$ 中有
$i_* \circ i^* = c_1(\mathcal{L}_{d + 1})$
（左边的意义由引理 \ref{chow-lemma-push-proper-bivariant} 保证）。
由于 $c_1(\mathcal{L}_i)$ 同时与 $i_*$ 和 $i^*$ 交换
（根据双变类的定义），可得
$$
c_1(\mathcal{L}_1) \circ \ldots \circ c_1(\mathcal{L}_{d + 1}) =
i_* \circ c_1(\mathcal{L}_1) \circ \ldots \circ c_1(\mathcal{L}_d) \circ i^* =
i_* \circ c_1(\mathcal{L}_1|_D) \circ \ldots \circ c_1(\mathcal{L}_d|_D)
\circ i^*
$$
于是由对 $d$ 的归纳得到结论。事实上，因为 $X$ 的一般点都不在 $D$ 中，
有 $\dim(D) < d$。
\end{proof}

\begin{remark}
\label{chow-remark-ring-loc-classes}
设 $(S, \delta)$ 如情形 \ref{chow-situation-setup}。
设 $Z \to X$ 是 $S$ 上局部有限型概形间的闭浸入，且 $p \geq 0$。
在这种框架下定义
$$
A^{(p)}(Z \to X) =
\prod\nolimits_{i \leq p - 1} A^i(X) \times
\prod\nolimits_{i \geq p} A^i(Z \to X).
$$
于是 $A^{(p)}(Z \to X)$ 典范地带有分次代数结构。
事实上，更一般地存在乘法
$$
A^{(p)}(Z \to X) \times A^{(q)}(Z \to X)
\longrightarrow A^{(\max(p, q))}(Z \to X)
$$
为定义这些乘法，定义映射
\begin{align*}
A^i(Z \to X) \times A^j(X) & \to A^{i + j}(Z \to X) \\
A^i(X) \times A^j(Z \to X) & \to A^{i + j}(Z \to X) \\
A^i(Z \to X) \times A^j(Z \to X) & \to A^{i + j}(Z \to X)
\end{align*}
对第一个映射，使用双变类的复合。对第二个映射，使用限制
$A^i(X) \to A^i(Z)$（注 \ref{chow-remark-restriction-bivariant}）
以及复合 $A^i(Z) \times A^j(Z \to X) \to A^{i + j}(Z \to X)$。
对第三个映射，把 $(c, c')$ 映为 $res(c) \circ c'$，其中
$res : A^i(Z \to X) \to A^i(Z)$ 是限制映射
（见注 \ref{chow-remark-restriction-bivariant}）。
这里略去验证这些乘法在适当意义下具有结合性。
\end{remark}

\begin{remark}
\label{chow-remark-res-push}
设 $(S, \delta)$ 如情形 \ref{chow-situation-setup}。
设 $Z \to X$ 是 $S$ 上局部有限型概形间的闭浸入。
以 $res : A^p(Z \to X) \to A^p(Z)$ 表示注 \ref{chow-remark-restriction-bivariant} 的限制映射。
对 $c \in A^p(Z \to X)$ 和 $\alpha \in \CH_*(Z)$，有
$res(c) \cap \alpha = c \cap i_*\alpha$。事实上，
$res(c) \cap \alpha = c \cap \alpha$，而 $c$ 与固有推前的相容性给出
$(Z \to Z)_*(c \cap \alpha) = c \cap (Z \to X)_*\alpha$。
\end{remark}






\section{关于双变类的引理}
\label{chow-section-bivariant-II}

\noindent
本节证明若干关于双变类的初等结果。先给出判断一个运算能否下降到
有理等价类的判据。

\begin{lemma}
\label{chow-lemma-factors-through-rational-equivalence}
\begin{reference}
\cite[Theorem 17.1]{F} 的一个很弱的形式
\end{reference}
设 $(S, \delta)$ 如情形 \ref{chow-situation-setup}。
设 $f : X \to Y$ 是 $S$ 上局部有限型概形间的态射，并设
$p \in \mathbf{Z}$。假设给定如下规则：对每个局部有限型态射
$Y' \to Y$ 及每个 $k$，指定映射
$$
c \cap - : Z_k(Y') \longrightarrow \CH_{k - p}(X')
$$
其中 $Y' = X' \times_X Y$；并且每当
$\mathcal{L}'|_{D'} \cong \mathcal{O}_{D'}$ 时，该规则满足定义 \ref{chow-definition-bivariant-class} 的条件 (3)。则
$c \cap -$ 下降到有理等价类。
\end{lemma}

\begin{proof}
该陈述有意义：若给定定义 \ref{chow-definition-gysin-homomorphism} 中那样的三元组
$(\mathcal{L}, s, i : D \to X)$，且
$\mathcal{L}|_D \cong \mathcal{O}_D$，则运算 $i^*$
已在循环层次上定义；见注 \ref{chow-remark-gysin-on-cycles}。
设 $\alpha \in Z_k(X')$ 是有理等价于零的循环。需要证明
$c \cap \alpha = 0$。由引理 \ref{chow-lemma-rational-equivalence-family}，
存在循环 $\beta \in Z_{k + 1}(X' \times \mathbf{P}^1)$，使得
$\alpha = i_0^*\beta - i_\infty^*\beta$，其中
$i_0, i_\infty : X' \to X' \times \mathbf{P}^1$ 是位于
$0, \infty$ 上的 $X'$ 的闭浸入。它们是具有平凡法丛的有效
Cartier 除子的例子，所以
$c \cap i_0^*\beta = j_0^*(c \cap \beta)$ 且
$c \cap i_\infty^*\beta = j_\infty^*(c \cap \beta)$，其中
$j_0, j_\infty : Y' \to Y' \times \mathbf{P}^1$ 是与前述相同的闭浸入。
又因 $j_0^*(c \cap \beta) \sim_{rat} j_\infty^*(c \cap \beta)$
（由引理 \ref{chow-lemma-rational-equivalence-family}），结论成立。
\end{proof}

\begin{lemma}
\label{chow-lemma-bivariant-weaker}
\begin{reference}
\cite[Theorem 17.1]{F} 的弱形式
\end{reference}
设 $(S, \delta)$ 如情形 \ref{chow-situation-setup}。
设 $f : X \to Y$ 是 $S$ 上局部有限型概形间的态射，并设
$p \in \mathbf{Z}$。假设给定如下规则：对每个局部有限型态射
$Y' \to Y$ 及每个 $k$，指定映射
$$
c \cap - : \CH_k(Y') \longrightarrow \CH_{k - p}(X')
$$
其中 $Y' = X' \times_X Y$；该规则满足定义 \ref{chow-definition-bivariant-class} 的条件 (1)、(2)，并且每当
$\mathcal{L}'|_{D'} \cong \mathcal{O}_{D'}$ 时满足条件 (3)。
则 $c \cap -$ 是双变类。
\end{lemma}

\begin{proof}
设 $Y' \to Y$ 是局部有限型概形态射。设
$(\mathcal{L}', s', i' : D' \to Y')$ 如定义 \ref{chow-definition-gysin-homomorphism}，其到 $X'$ 的拉回为
$(\mathcal{N}', t', j' : E' \to X')$。
需要证明对所有 $\alpha' \in \CH_k(Y')$，都有
$c \cap (i')^*\alpha' = (j')^*(c \cap \alpha')$。

\medskip\noindent
以 $g : Y'' \to Y'$ 表示引理 \ref{chow-lemma-normal-cone-effective-Cartier} 中构造的相对维数为 $1$
的光滑态射，并同时采用其中的 $i'' : D'' \to Y''$ 和
$p : D'' \to D'$。（警告：$D''$ 不是 $D'$ 的完整逆像。）
以 $f : X'' \to X'$ 和 $E'' \subset X''$ 表示它们经
$X' \to Y'$ 的基变换。图示为
$$
\xymatrix{
& X'' \ar[rr] \ar'[d][dd]_h & & Y'' \ar[dd]^g \\
E'' \ar[rr] \ar[dd]_q \ar[ru]^{j''} & & D'' \ar[dd]^p \ar[ru]^{i''} & \\
& X' \ar'[r][rr] & & Y' \\
E' \ar[rr] \ar[ru]^{j'} & & D' \ar[ru]^{i'}
}
$$
由该引理所给性质，$\beta' = (i')^*\alpha'$ 是满足
$p^*\beta' = (i'')^*g^*\alpha'$ 的 $\CH_{k - 1}(D')$ 中唯一元素。
类似地，$\gamma' = (j')^*(c \cap \alpha')$ 是满足
$q^*\gamma' = (j'')^*h^*(c \cap \alpha')$ 的
$\CH_{k - 1 - p}(E')$ 中唯一元素。由对 $c$ 的假设，已知
$$
(j'')^*h^*(c \cap \alpha') =
(j'')^*(c \cap g^*\alpha') =
c \cap (i'')^*g^*\alpha'
$$
注意，引理陈述中所假设的条件 (3) 的修改版适用于 $i''$
及其基变换 $j''$。类似地，已知
$$
q^*(c \cap \beta') = c \cap p^*\beta'
$$
由上述唯一性，得到 $\gamma' = c \cap \beta'$。
\end{proof}

\noindent
下面给出双变类为零的判据。

\begin{lemma}
\label{chow-lemma-bivariant-zero}
设 $(S, \delta)$ 如情形 \ref{chow-situation-setup}。
设 $f : X \to Y$ 是 $S$ 上局部有限型概形间的态射。
设 $c \in A^p(X \to Y)$。对 $Y'' \to Y' \to Y$，置
$X'' = Y'' \times_Y X$ 及 $X' = Y' \times_Y X$。以下条件等价：
\begin{enumerate}
\item $c$ 为零；
\item 对每个在 $Y$ 上局部有限型的整概形 $Y'$，在
$\CH_*(X')$ 中有 $c \cap [Y'] = 0$；以及
\item 对每个在 $Y$ 上局部有限型的整概形 $Y'$，存在固有双有理态射
$Y'' \to Y'$，使得在 $\CH_*(X'')$ 中有 $c \cap [Y''] = 0$。
\end{enumerate}
\end{lemma}

\begin{proof}
蕴含关系 (1) $\Rightarrow$ (2) $\Rightarrow$ (3) 是显然的。
条件 (3) 蕴含 (2)，因为 $(Y'' \to Y')_*[Y''] = [Y']$，并且由于
$c$ 是双变类，有
$c \cap [Y'] = (X'' \to X')_*(c \cap [Y''])$。现在假设 (2)。
设 $Y' \to Y$ 局部有限型，并设 $\alpha \in \CH_k(Y')$。
写 $\alpha = \sum n_i [Y'_i]$，其中 $Y'_i \subset Y'$ 是一族局部有限、
$\delta$-维数为 $k$ 的整闭子概形。于是 $\alpha$ 是循环
$\alpha' = \sum n_i[Y'_i]$ 在 $Y'' = \coprod Y'_i$ 上沿固有态射
$Y'' \to Y'$ 的推前。由双变类的性质，只须证明
$c \cap \alpha' = 0$ 在 $\CH_{k - p}(X'')$ 中成立。
有 $\CH_{k - p}(X'') = \prod \CH_{k - p}(X'_i)$，其中
$X'_i = Y'_i \times_Y X$；这立即由定义得出。投影映射
$\CH_{k - p}(X'') \to \CH_{k - p}(X'_i)$ 由平坦拉回给出。
由于与 $c$ 作帽积和平坦拉回交换，只须证明
$c \cap [Y'_i]$ 在 $\CH_{k - p}(X'_i)$ 中为零，而这由假设成立。
\end{proof}

\begin{lemma}
\label{chow-lemma-disjoint-decomposition-bivariant}
设 $(S, \delta)$ 如情形 \ref{chow-situation-setup}。
设 $f : X \to Y$ 是 $S$ 上局部有限型概形间的态射。
假设由开闭子概形给出不交并分解
$X = \coprod_{i \in I} X_i$ 和 $Y = \coprod_{j \in J} Y_j$，
并且存在集合映射 $a : I \to J$，使得
$f(X_i) \subset Y_{a(i)}$。则
$$
A^p(X \to Y) = \prod\nolimits_{i \in I} A^p(X_i \to Y_{a(i)})
$$
\end{lemma}

\begin{proof}
设给定元素 $(c_i) \in \prod_i A^p(X_i \to Y_{a(i)})$。
对 $\beta \in \CH_k(Y)$，可把它映到 $\CH_{k - p}(X)$ 中的如下元素：
其到 $X_i$ 的限制为 $c_i \cap \beta|_{Y_{a(i)}}$。
这是可行的，因为 $\CH_{k - p}(X) = \prod_i \CH_{k - p}(X_i)$。
经任意局部有限型基变换 $Y' \to Y$ 后，同一构造仍可进行，
从而得到 $c \in A^p(X \to Y)$。于是得到从公式右边到左边的映射
$\Psi$。反之，给定 $c \in A^p(X \to Y)$ 和元素
$\beta_i \in \CH_k(Y_{a(i)})$，可以考虑 $\CH_{k - p}(X_i)$ 中的元素
$(c \cap (Y_{a(i)} \to Y)_*\beta_i)|_{X_i}$。
经任意局部有限型基变换 $Y' \to Y$ 后，同样可行，从而得到
$c_i \in A^p(X_i \to Y_{a(i)})$。于是得到从公式左边到右边的映射
$\Phi$。显然，$\Phi \circ \Psi = \text{id}$。
对反向复合，设 $c \in A^p(X \to Y)$ 且 $\beta \in \CH_k(Y)$。
设 $\Phi(c) = (c_i)$，并取 $j \in J$。因为 $c$ 与平坦拉回交换，有
$$
(c \cap \beta)|_{\coprod_{a(i) = j} X_i} =
c \cap \beta|_{Y_j}
$$
又因为 $c$ 与固有推前交换，有
$$
(\coprod\nolimits_{a(i) = j} X_i \to X)_*
((c \cap \beta)|_{\coprod_{a(i) = j} X_i})
=
c \cap (Y_j \to Y)_*\beta|_{Y_j}
$$
左边是 $X$ 上的循环：对满足 $a(i) = j$ 的 $i \in I$，
其到 $X_i$ 的限制为 $(c \cap \beta)|_{X_i}$，其余分量上为 $0$。
右边是 $X$ 上的循环：对满足 $a(i) = j$ 的 $i \in I$，
其到 $X_i$ 的限制为 $c_i \cap \beta|_{Y_j}$。
所以如所需有 $c \cap \beta = \Psi((c_i))$。
\end{proof}

\begin{remark}
\label{chow-remark-completion-bivariant}
设 $(S, \delta)$ 如情形 \ref{chow-situation-setup}。
设 $f : X \to Y$ 是 $S$ 上局部有限型概形间的态射。
设 $X = \coprod_{i \in I} X_i$ 和 $Y = \coprod_{j \in J} Y_j$
分别是 $X$、$Y$ 到各连通分支的分解
（因为 $X$、$Y$ 局部 Noether，连通分支是开的；见《拓扑学》引理 \ref{topology-lemma-locally-Noetherian-locally-connected} 和
《概形的性质》引理 \ref{properties-lemma-Noetherian-topology}）。
设 $a(i) \in J$ 是满足 $f(X_i) \subset Y_{a(i)}$ 的指标。
由引理 \ref{chow-lemma-disjoint-decomposition-bivariant}，
$A^p(X \to Y) = \prod A^p(X_i \to Y_{a(i)})$。
在此框架下，置
$$
A^*(X \to Y)^\wedge = \prod\nolimits_i A^*(X_i \to Y_{a(i)})
$$
较为方便。这个“完备化”的双变群是子集
$$
A^*(X \to Y)^\wedge \quad\subset\quad \prod\nolimits_{p \geq 0} A^p(X \to Y)
$$
它由如下元素组成：$c = (c_0, c_1, c_2, \ldots)$，且对每个连通分支
$X_i$，$c_p$ 在 $A^p(X_i \to Y_{a(i)})$ 中的像仅对有限多个 $p$
非零。若 $Y \to Z$ 是第二个态射，则复合
$A^*(X \to Y) \times A^*(Y \to Z) \to A^*(X \to Z)$
延拓为完备化之间的复合
$A^*(X \to Y)^\wedge \times A^*(Y \to Z)^\wedge \to A^*(X \to Z)^\wedge$。
有时把 $A^*(X)^\wedge = A^*(X \to X)^\wedge$ 称为 $X$ 的
{\it 完备双变上同调环}。
\end{remark}

\begin{lemma}
\label{chow-lemma-envelope-bivariant}
设 $(S, \delta)$ 如情形 \ref{chow-situation-setup}。
设 $f : X \to Y$ 是 $S$ 上局部有限型概形间的态射。
设 $g : Y' \to Y$ 是包络（定义 \ref{chow-definition-envelope}），
并记 $X' = Y' \times_Y X$。设 $p \in \mathbf{Z}$ 且
$c' \in A^p(X' \to Y')$。若两个限制
$$
res_1(c') = res_2(c') \in A^p(X' \times_X X' \to Y' \times_Y Y')
$$
相等（见证明），则存在唯一的 $c \in A^p(X \to Y)$，使其在
$A^p(X' \to Y')$ 中的限制满足 $res(c) = c'$。
\end{lemma}

\begin{proof}
有交换图
$$
\xymatrix{
X' \times_X X' \ar[d]^{f''} \ar@<1ex>[r]^-a \ar@<-1ex>[r]_-b &
X' \ar[d]^{f'} \ar[r]_h &
X \ar[d]^f \\
Y' \times_Y Y' \ar@<1ex>[r]^-p \ar@<-1ex>[r]_-q &
Y' \ar[r]^g &
Y
}
$$
元素 $res_1(c')$ 是 $c'$ 关于由态射 $a, f', p, f''$ 组成的 Cartesian
方块的限制（见注 \ref{chow-remark-restriction-bivariant}）；
元素 $res_2(c')$ 是 $c'$ 关于由态射 $b, f', q, f''$ 组成的 Cartesian
方块的限制。假设 $res_1(c') = res_2(c')$，并设
$\beta \in \CH_k(Y)$。由引理 \ref{chow-lemma-envelope}，可以找到
$\beta' \in \CH_k(Y')$，使 $g_*\beta' = \beta$。于是置
$$
c \cap \beta = h_*(c' \cap \beta')
$$
为证明这与 $\beta'$ 的选取无关，只须证明对
$\gamma \in \CH_k(Y' \times_Y Y')$，
$h_*(c' \cap (p_*\gamma - q_*\gamma))$ 为零。
因为 $c'$ 是双变类，有
$$
h_*(c' \cap (p_*\gamma - q_*\gamma)) =
h_*(a_*(c' \cap \gamma) - b_*(c' \cap \gamma)) = 0
$$
最后一个等式成立，是因为 $h \circ a = h \circ b$，从而
$h_* \circ a_* = h_* \circ b_*$。

\medskip\noindent
注意，要求 $res(c) = c'$ 以及双变类与固有推前相容，已经唯一地强制
了对 $c \cap \beta$ 的上述选取。

\medskip\noindent
当然，为定义双变类 $c$，还需要对任意局部有限型态射
$Y_1 \to Y$ 构造满足定义 \ref{chow-definition-bivariant-class}
所列条件的映射
$c \cap -: \CH_k(Y_1) \to \CH_{k + p}(Y_1 \times_Y X)$。
记 $Y'_1 = Y' \times_Y Y_1$、$X_1 = X \times_Y Y_1$。
由引理 \ref{chow-lemma-base-change-envelope}，态射
$Y'_1 \to Y_1$ 是包络。因此可以使用基变换后的图表
$$
\xymatrix{
X'_1 \times_{X_1} X'_1 \ar[d]^{f''_1} \ar@<1ex>[r]^-{a_1} \ar@<-1ex>[r]_-{b_1} &
X'_1 \ar[d]^{f'_1} \ar[r]_{h_1} &
X_1 \ar[d]^{f_1} \\
Y'_1 \times_{Y_1} Y'_1 \ar@<1ex>[r]^-{p_1} \ar@<-1ex>[r]_-{q_1} &
Y'_1 \ar[r]^{g_1} &
Y_1
}
$$
并用同样的论证得到如前良定义的映射
$c \cap - : \CH_k(Y_1) \to \CH_{k + p}(X_1)$。

\medskip\noindent
接下来，需要验证 $c$ 满足定义 \ref{chow-definition-bivariant-class} 的条件 (1)、(2)、(3)。
例如，假设 $t : Y_2 \to Y_1$ 是 $Y$ 上局部有限型概形间的固有态射。
像上面一样，以下标 ${}_1$、相应地\ ${}_2$ 表示第一幅图到
$Y_1$、相应地\ $Y_2$ 的基变换。以 $t' : Y'_2 \to Y'_1$、
$s : X_2 \to X_1$ 和 $s' : X'_2 \to X'_1$ 表示 $t$ 到
$Y'$、$X$ 和 $X'$ 的基变换。需要证明
$$
s_*(c \cap \beta_2) = c \cap t_*\beta_2
$$
对 $\beta_2 \in \CH_k(Y_2)$ 成立。选取
$\beta'_2 \in \CH_k(Y'_2)$，使 $g_{2, *}\beta'_2 = \beta_2$。
由于 $c'$ 是双变类，且图表
$$
\vcenter{
\xymatrix{
X'_2 \ar[d]_{s'} \ar[r]_{h_2} & X_2 \ar[d]^s \\
X'_1 \ar[r]^{h_1} & X_1
}
}
\quad\text{且}\quad
\vcenter{
\xymatrix{
X'_2 \ar[d]_{s'} \ar[r]_{f'_2} & Y'_2 \ar[d]^{t'} \\
X'_2 \ar[r]^{f'_1} & Y'_1
}
}
$$
都是 Cartesian 的，故有
$$
s_*(c \cap \beta_2) =
s_*(h_{2, *}(c' \cap \beta'_2)) =
h_{1, *}s'_*(c' \cap \beta'_2) =
h_{1, *}(c' \cap (t'_*\beta'_2))
$$
按构造，最后一个表达式计算 $c \cap t_*\beta_2$：类
$t'_*\beta'_2 \in \CH_k(Y'_1)$ 沿 $g_{1, *}$ 的像是
$t_*\beta_2$。这证明了条件 (1)。其余条件用同样方法证明，
这里略去详细论证。
\end{proof}















\section{射影空间丛公式}
\label{chow-section-projective-space-bundle-formula}

\noindent
设 $(S, \delta)$ 如情形 \ref{chow-situation-setup}。
设 $X$ 在 $S$ 上局部有限型。考虑秩为 $r$ 的有限局部自由
$\mathcal{O}_X$-模 $\mathcal{E}$。我们的约定是，{\it 与
$\mathcal{E}$ 关联的射影丛}是 $X$ 上的态射
$$
\xymatrix{
\mathbf{P}(\mathcal{E}) =
\underline{\text{Proj}}_X(\text{Sym}^*(\mathcal{E}))
\ar[r]^-\pi
& X
}
$$
其中 $\mathcal{O}_{\mathbf{P}(\mathcal{E})}(1)$ 规范化为满足
$\pi_*(\mathcal{O}_{\mathbf{P}(\mathcal{E})}(1)) = \mathcal{E}$。
特别地，存在满射
$\pi^*\mathcal{E} \to \mathcal{O}_{\mathbf{P}(\mathcal{E})}(1)$。
下文将非正式地说“设 $(\pi : P \to X, \mathcal{O}_P(1))$
是与 $\mathcal{E}$ 关联的射影丛”，以表示
$P = \mathbf{P}(\mathcal{E})$ 且
$\mathcal{O}_P(1) = \mathcal{O}_{\mathbf{P}(\mathcal{E})}(1)$ 的情形。

\begin{lemma}
\label{chow-lemma-cap-projective-bundle}
设 $(S, \delta)$ 如情形 \ref{chow-situation-setup}。
设 $X$ 在 $S$ 上局部有限型。设 $\mathcal{E}$ 是秩为 $r$ 的有限局部自由
$\mathcal{O}_X$-模 $\mathcal{E}$。设
$(\pi : P \to X, \mathcal{O}_P(1))$ 是与 $\mathcal{E}$ 关联的射影丛。
对任意 $\alpha \in \CH_k(X)$，元素
$$
\pi_*\left(
c_1(\mathcal{O}_P(1))^s \cap \pi^*\alpha
\right)
\in
\CH_{k + r - 1 - s}(X)
$$
在 $s < r - 1$ 时为 $0$，在 $s = r - 1$ 时等于 $\alpha$。
\end{lemma}

\begin{proof}
设 $Z \subset X$ 是 $\delta$-维数为 $k$ 的整闭子概形。
注意，$\pi^*[Z] = [\pi^{-1}(Z)]$，因为 $\pi^{-1}(Z)$ 是
$\delta$-维数为 $r - 1$ 的整概形。若 $s < r - 1$，则根据构造，
$c_1(\mathcal{O}_P(1))^s \cap \pi^*[Z]$ 由一个支撑在
$\pi^{-1}(Z)$ 上的 $(k + r - 1 - s)$-循环表示。
因此，由维数原因，该循环的推前为零。

\medskip\noindent
设 $s = r - 1$。由上述论证，对某个 $n \in \mathbf{Z}$ 有
$\pi_*(c_1(\mathcal{O}_P(1))^s \cap \pi^*\alpha) = n [Z]$。
需要证明 $n = 1$。由同样的维数原因，只须在用
$X \setminus T$ 替换 $X$ 后证明这一结果，其中
$T \subset Z$ 是真闭子集。设 $\xi$ 是 $Z$ 的一般点。
可以选取元素 $e_1, \ldots, e_{r - 1} \in \mathcal{E}_\xi$，
使它们构成 $\mathcal{E}_\xi$ 的某个基的一部分。
它们给出 $\mathcal{O}_P(1)|_{\pi^{-1}(Z)}$ 的有理截面
$s_1, \ldots, s_{r - 1}$；这些截面的公共零点集是态射
$\mathbf{P}(\mathcal{E}|_Z) \to Z$ 的一个有理截面之像的闭包，
并上一个支撑映到 $Z$ 的真闭子集 $T$ 的闭子集。
从 $X$ 中移除 $T$（相应地从 $P$ 中移除 $\pi^{-1}(T)$）后，
$s_1, \ldots, s_n$ 构成一列整体截面
$s_i \in \Gamma(\pi^{-1}(Z), \mathcal{O}_{\pi^{-1}(Z)}(1))$，
其公共零点集是某个截面 $Z \to \pi^{-1}(Z)$ 的像。
因此，反复应用引理 \ref{chow-lemma-geometric-cap}，依次得到
\begin{eqnarray*}
\pi^*[Z] & = & [\pi^{-1}(Z)] \\
c_1(\mathcal{O}_P(1)) \cap \pi^*[Z] & = & [Z(s_1)] \\
c_1(\mathcal{O}_P(1))^2 \cap \pi^*[Z] & = & [Z(s_1) \cap Z(s_2)] \\
\ldots & = & \ldots \\
c_1(\mathcal{O}_P(1))^{r - 1} \cap \pi^*[Z] & = &
[Z(s_1) \cap \ldots \cap Z(s_{r - 1})]
\end{eqnarray*}
由于 $\pi$ 在 $Z$ 上的一个截面的像沿 $\pi$ 的推前显然为 $[Z]$，
在 $\alpha = [Z]$ 时便得到结论。这里略去验证这些论证对一般循环
$\alpha = \sum n_j [Z_j]$ 也给出结论。
\end{proof}

\begin{lemma}[射影空间丛公式]
\label{chow-lemma-chow-ring-projective-bundle}
设 $(S, \delta)$ 如情形 \ref{chow-situation-setup}。
设 $X$ 在 $S$ 上局部有限型。
设 $\mathcal{E}$ 是秩为 $r$ 的有限局部自由 $\mathcal{O}_X$-模
$\mathcal{E}$。设 $(\pi : P \to X, \mathcal{O}_P(1))$
是与 $\mathcal{E}$ 关联的射影丛。映射
$$
\bigoplus\nolimits_{i = 0}^{r - 1}
\CH_{k + i}(X)
\longrightarrow
\CH_{k + r - 1}(P),
$$
$$
(\alpha_0, \ldots, \alpha_{r-1})
\longmapsto
\pi^*\alpha_0 +
c_1(\mathcal{O}_P(1)) \cap \pi^*\alpha_1
+ \ldots +
c_1(\mathcal{O}_P(1))^{r - 1} \cap \pi^*\alpha_{r-1}
$$
是同构。
\end{lemma}

\begin{proof}
固定 $k \in \mathbf{Z}$。先证该映射为单射。
设 $(\alpha_0, \ldots, \alpha_{r - 1})$ 是左端的一个元素，
且映到零。由引理 \ref{chow-lemma-cap-projective-bundle} 可知
$$
0 = \pi_*(\pi^*\alpha_0 +
c_1(\mathcal{O}_P(1)) \cap \pi^*\alpha_1
+ \ldots +
c_1(\mathcal{O}_P(1))^{r - 1} \cap \pi^*\alpha_{r-1})
= \alpha_{r - 1}
$$
接着可知
$$
0 = \pi_*(c_1(\mathcal{O}_P(1)) \cap (\pi^*\alpha_0 +
c_1(\mathcal{O}_P(1)) \cap \pi^*\alpha_1
+ \ldots +
c_1(\mathcal{O}_P(1))^{r - 2} \cap \pi^*\alpha_{r - 2}))
= \alpha_{r - 2}
$$
依此类推。因此该映射为单射。

\medskip\noindent
还须证明该映射为满射。
设 $X_i$（$i \in I$）为 $X$ 的不可约分支。
则 $P_i = \mathbf{P}(\mathcal{E}|_{X_i})$（$i \in I$）
是 $P$ 的不可约分支。考虑交换图
$$
\xymatrix{
\coprod P_i \ar[d]_{\coprod \pi_i} \ar[r]_p & P \ar[d]^\pi \\
\coprod X_i \ar[r]^q & X
}
$$
注意 $p_*$ 为满射。若 $\beta \in \CH_k(\coprod X_i)$，则
$\pi^* q_* \beta = p_*(\coprod \pi_i)^* \beta$，参见
引理 \ref{chow-lemma-flat-pullback-proper-pushforward}。由
引理 \ref{chow-lemma-pushforward-cap-c1}，与
$c_1(\mathcal{O}(1))$ 作帽积时也有类似结论。
因此，若引理中的映射对每个态射 $\pi_i : P_i \to X_i$ 都是满射，
则它对 $\pi : P \to X$ 也是满射。故可设 $X$ 不可约。
于是 $\dim_\delta(X) < \infty$，特别地，可以对
$\dim_\delta(X)$ 作归纳。

\medskip\noindent
若 $\dim_\delta(X) < k$，结论显然成立。
设 $\alpha \in \CH_{k + r - 1}(P)$。
对任意局部闭子概形 $T \subset X$，记
$\gamma_T : \bigoplus \CH_{k + i}(T) \to \CH_{k + r - 1}(\pi^{-1}(T))$
为映射
$$
\gamma_T(\alpha_0, \ldots, \alpha_{r - 1})
= \pi^*\alpha_0 + \ldots +
c_1(\mathcal{O}_{\pi^{-1}(T)}(1))^{r - 1} \cap \pi^*\alpha_{r - 1}.
$$
设对某个非空开集 $U \subset X$ 有
$\alpha|_{\pi^{-1}(U)} = \gamma_U(\alpha_0, \ldots, \alpha_{r - 1})$。
此时可以选取提升 $\alpha'_i \in \CH_{k + i}(X)$，并由
引理 \ref{chow-lemma-restrict-to-open} 看出，
$\alpha - \gamma_X(\alpha'_0, \ldots, \alpha'_{r - 1})$
有理等价于 $P_Y = \mathbf{P}(\mathcal{E}|_Y)$ 上的一个 $k$-循环，
其中 $Y = X \setminus U$ 赋予既约闭子概形结构。
注意 $\dim_\delta(Y) < \dim_\delta(X)$。
由归纳假设，结论对 $P_Y \to Y$ 成立，从而对 $\alpha$ 也成立。
因此，可以用 $X$ 的任意非空开集替换 $X$。

\medskip\noindent
特别地，可以设 $\mathcal{E} \cong \mathcal{O}_X^{\oplus r}$。
此时 $\mathbf{P}(\mathcal{E}) = X \times \mathbf{P}^{r - 1}$。
采用分层
$$
\mathbf{P}^{r - 1} = \mathbf{A}^{r - 1}
\amalg \mathbf{A}^{r - 2}
\amalg \ldots
\amalg \mathbf{A}^0
$$
每一层的闭包都是一个 $\mathbf{P}^{r - 1 - i}$，它表示
$c_1(\mathcal{O}(1))^i \cap [\mathbf{P}^{r - 1}]$。
因此 $P$ 有类似的分层
$$
P = U^{r - 1} \amalg U^{r - 2} \amalg \ldots \amalg U^0
$$
设 $P^i$ 为 $U^i$ 的闭包，设 $\pi^i : P^i \to X$
为 $\pi$ 在 $P^i$ 上的限制。
设 $\alpha \in \CH_{k + r - 1}(P)$。由
引理 \ref{chow-lemma-pullback-affine-fibres-surjective}，
可对某个 $\alpha_0 \in \CH_k(X)$ 写成
$\alpha|_{U^{r - 1}} = \pi^*\alpha_0|_{U^{r - 1}}$。
因此差 $\alpha - \pi^*\alpha_0$ 是某个
$\alpha' \in \CH_{k + r - 1}(P^{r - 2})$ 的像。
再次应用引理 \ref{chow-lemma-pullback-affine-fibres-surjective}，
可对某个 $\alpha_1 \in \CH_{k + 1}(X)$ 写成
$\alpha'|_{U^{r - 2}} = (\pi^{r - 2})^*\alpha_1|_{U^{r - 2}}$。
由引理 \ref{chow-lemma-relative-effective-cartier}，
$(\pi^{r - 2})^*\alpha_1$ 的像表示
$c_1(\mathcal{O}_P(1)) \cap \pi^*\alpha_1$。
还可看出，
$\alpha - \pi^*\alpha_0 - c_1(\mathcal{O}_P(1)) \cap \pi^*\alpha_1$
是某个 $\alpha'' \in \CH_{k + r - 1}(P^{r - 3})$ 的像。
依此类推。
\end{proof}

\begin{lemma}
\label{chow-lemma-vectorbundle}
设 $(S, \delta)$ 如情形 \ref{chow-situation-setup}。
设 $X$ 在 $S$ 上局部有限型。
设 $\mathcal{E}$ 是 $X$ 上秩为 $r$ 的有限局部自由层。令
$$
p :
E = \underline{\Spec}(\text{Sym}^*(\mathcal{E}))
\longrightarrow
X
$$
为 $X$ 上相应的向量丛。则对所有 $k$，
$p^* : \CH_k(X) \to \CH_{k + r}(E)$ 都是同构。
\end{lemma}

\begin{proof}
（秩一向量丛的情形见引理 \ref{chow-lemma-linebundle}。）
满射性见引理 \ref{chow-lemma-pullback-affine-fibres-surjective}。
设 $(\pi  : P \to X, \mathcal{O}_P(1))$ 是与有限局部自由层
$\mathcal{E} \oplus \mathcal{O}_X$ 关联的射影空间丛。
设 $s \in \Gamma(P, \mathcal{O}_P(1))$ 对应于整体截面
$(0, 1) \in \Gamma(X, \mathcal{E} \oplus \mathcal{O}_X)$。
令 $D = Z(s) \subset P$。注意，
$(\pi|_D : D \to X , \mathcal{O}_P(1)|_D)$
是与 $\mathcal{E}$ 关联的射影空间丛。记
$\pi_D = \pi|_D$ 且 $\mathcal{O}_D(1) = \mathcal{O}_P(1)|_D$。
此外，$D$ 是 $P$ 上的有效 Cartier 除子。因此
$\mathcal{O}_P(D) = \mathcal{O}_P(1)$
（见《除子》引理 \ref{divisors-lemma-characterize-OD}）。
还有一个同构 $E \cong P \setminus D$。记
$j : E \to P$ 为相应的开浸入。
为证明单射性，使用如下事实：
$$
j^* :
\CH_{k + r}(P)
\longrightarrow
\CH_{k + r}(E)
$$
的核由支撑在有效 Cartier 除子 $D$ 上的循环组成，见
引理 \ref{chow-lemma-restrict-to-open}。所以，若 $p^*\alpha = 0$，
则对某个 $\beta \in \CH_{k + r}(D)$ 有
$\pi^*\alpha = i_*\beta$。由
引理 \ref{chow-lemma-chow-ring-projective-bundle}，可写成
$$
\beta = \pi_D^*\beta_0 +
\ldots + c_1(\mathcal{O}_D(1))^{r - 1} \cap \pi_D^* \beta_{r - 1}.
$$
其中 $\beta_i \in \CH_{k + i}(X)$。由引理 \ref{chow-lemma-relative-effective-cartier} 和
\ref{chow-lemma-pushforward-cap-c1}，这推出
$$
\pi^*\alpha = i_*\beta =
c_1(\mathcal{O}_P(1)) \cap \pi^*\beta_0 +
\ldots +
c_1(\mathcal{O}_D(1))^r \cap \pi^*\beta_{r - 1}.
$$
由于 $\mathcal{E} \oplus \mathcal{O}_X$ 的秩为 $r + 1$，
除非 $\alpha$ 以及所有 $\beta_i$ 都为零，否则这与
引理 \ref{chow-lemma-pushforward-cap-c1} 矛盾。
\end{proof}








\section{向量丛的陈类}
\label{chow-section-chern-classes-vector-bundles}

\noindent
可以利用射影空间丛公式，通过把 $c_1(\mathcal{O}(1))^r$
展开成较低次幂来定义秩为 $r$ 的向量丛的陈类，见公式
(\ref{chow-equation-chern-classes})。其中符号的缘由稍后说明。

\begin{definition}
\label{chow-definition-chern-classes}
设 $(S, \delta)$ 如情形 \ref{chow-situation-setup}。
设 $X$ 在 $S$ 上局部有限型。假设 $X$ 整且
$n = \dim_\delta(X)$。设 $\mathcal{E}$ 是 $X$ 上秩为 $r$
的有限局部自由层。设 $(\pi : P \to X, \mathcal{O}_P(1))$
是与 $\mathcal{E}$ 关联的射影空间丛。
\begin{enumerate}
\item 由引理 \ref{chow-lemma-chow-ring-projective-bundle}，存在元素
$c_i \in \CH_{n - i}(X)$（$i = 0, \ldots, r$），
使得 $c_0 = [X]$，且
\begin{equation}
\label{chow-equation-chern-classes}
\sum\nolimits_{i = 0}^r
(-1)^i c_1(\mathcal{O}_P(1))^i \cap \pi^*c_{r - i}
= 0.
\end{equation}
\item 采用上述记号，令
$c_i(\mathcal{E}) \cap [X] = c_i$，
视为 $\CH_{n - i}(X)$ 的元素。称它们为
{\it $X$ 上 $\mathcal{E}$ 的陈类}。
\item {\it $X$ 上 $\mathcal{E}$ 的全陈类}是组合
$$
c({\mathcal E}) \cap [X] =
c_0({\mathcal E}) \cap [X]
+ c_1({\mathcal E}) \cap [X] + \ldots
+ c_r({\mathcal E}) \cap [X]
$$
它是 $\CH_*(X) = \bigoplus_{k \in \mathbf{Z}} \CH_k(X)$
的元素。
\end{enumerate}
\end{definition}

\noindent
下面验证，当向量丛秩为 $1$ 时，这并未给出新的概念。

\begin{lemma}
\label{chow-lemma-first-chern-class}
设 $(S, \delta)$ 如情形 \ref{chow-situation-setup}。
设 $X$ 在 $S$ 上局部有限型。假设 $X$ 整且
$n = \dim_\delta(X)$。设 $\mathcal{L}$ 是可逆
$\mathcal{O}_X$-模。定义 \ref{chow-definition-chern-classes}
所定义的 $X$ 上 $\mathcal{L}$ 的第一陈类，等于
定义 \ref{chow-definition-divisor-invertible-sheaf}
赋予 $\mathcal{L}$ 的 Weil 除子。
\end{lemma}

\begin{proof}
本证明用 $c_1(\mathcal{L}) \cap [X]$ 表示
定义 \ref{chow-definition-divisor-invertible-sheaf} 中的构造。
由于 $\mathcal{L}$ 的秩为 $1$，根据我们的规范化有
$\mathbf{P}(\mathcal{L}) = X$ 且
$\mathcal{O}_{\mathbf{P}(\mathcal{L})}(1) = \mathcal{L}$。
因此 (\ref{chow-equation-chern-classes}) 成为
$$
(-1)^1 c_1(\mathcal{L}) \cap c_0 + (-1)^0 c_1 = 0
$$
由于 $c_0 = [X]$，可得所需的
$c_1 = c_1(\mathcal{L}) \cap [X]$。
\end{proof}

\begin{remark}
\label{chow-remark-equation-signs}
也可以把方程 \ref{chow-equation-chern-classes} 改写为
\begin{equation}
\label{chow-equation-signs}
\sum\nolimits_{i = 0}^r
c_1(\mathcal{O}_P(-1))^i \cap \pi^*c_{r - i}
= 0.
\end{equation}
不过，相比其对偶，我们认为使用重言商层
$\mathcal{O}_P(1)$ 更为方便。
\end{remark}




\section{与陈类相交}
\label{chow-section-intersecting-chern-classes}

\noindent
本节把 $X$ 上向量丛的陈类定义为 $X$ 上的双变类，参见
引理 \ref{chow-lemma-cap-cp-bivariant} 及该引理之后的讨论。
我们的构造遵循熟知的模式：先在素循环上定义运算，再求和。
在引理 \ref{chow-lemma-determine-intersections} 中，将证明结果由
相应射影丛上的通常公式所确定。随后证明，与陈类作帽积可下降到
有理等价，并与正常推前、平坦拉回以及有效 Cartier 除子嵌入的
Gysin 映射交换。也可以直接使用
引理 \ref{chow-lemma-determine-intersections} 中的刻画和
引理 \ref{chow-lemma-push-proper-bivariant}，从而避开这些引理；
希望看到具体推导的读者可参阅《空间的 Chow 群》引理 \ref{spaces-chow-lemma-segre-classes}。

\begin{definition}
\label{chow-definition-cap-chern-classes}
设 $(S, \delta)$ 如情形 \ref{chow-situation-setup}。
设 $X$ 在 $S$ 上局部有限型。设 $\mathcal{E}$ 是 $X$ 上秩为
$r$ 的有限局部自由层。对每个整数 $k$ 及任意
$0 \leq j \leq r$，定义运算
$$
c_j(\mathcal{E}) \cap - : Z_k(X) \to \CH_{k - j}(X)
$$
称为{\it 与 $\mathcal{E}$ 的第 $j$ 陈类相交}。
\begin{enumerate}
\item 给定一个 $\delta$-维数为 $k$ 的整闭子概形
$i : W \to X$，定义
$$
c_j(\mathcal{E}) \cap [W] = i_*(c_j({i^*\mathcal{E}}) \cap [W])
\in
\CH_{k - j}(X)
$$
其中 $c_j({i^*\mathcal{E}}) \cap [W]$ 如
定义 \ref{chow-definition-chern-classes} 中所定义。
\item 对一般的 $k$-循环 $\alpha = \sum n_i [W_i]$，令
$$
c_j(\mathcal{E}) \cap \alpha = \sum n_i c_j(\mathcal{E}) \cap [W_i]
$$
\end{enumerate}
\end{definition}

\noindent
若 $\mathcal{E}$ 的秩为 $1$，则这与此前的定义
（定义 \ref{chow-definition-cap-c1}）一致，这是
引理 \ref{chow-lemma-first-chern-class} 的结论。

\begin{lemma}
\label{chow-lemma-determine-intersections}
设 $(S, \delta)$ 如情形 \ref{chow-situation-setup}。
设 $X$ 在 $S$ 上局部有限型。设 $\mathcal{E}$ 是 $X$ 上秩为
$r$ 的有限局部自由层。设
$(\pi : P \to X, \mathcal{O}_P(1))$ 是与 $\mathcal{E}$ 关联的射影丛。
对 $\alpha \in Z_k(X)$，元素
$c_j(\mathcal{E}) \cap \alpha$ 是 $\CH_{k - j}(X)$ 中满足
$\alpha_0 = \alpha$ 以及
$$
\sum\nolimits_{i = 0}^r
(-1)^i c_1(\mathcal{O}_P(1))^i \cap
\pi^*(\alpha_{r - i}) = 0
$$
在 $P$ 的 Chow 群中成立的唯一元素 $\alpha_j$。
\end{lemma}

\begin{proof}
满足 $\alpha_0 = \alpha$ 及所示方程的
$\alpha_0, \ldots, \alpha_r$ 的唯一性，来自射影空间丛公式
引理 \ref{chow-lemma-chow-ring-projective-bundle}。
当 $\alpha = [W]$ 且 $W$ 是 $X$ 的整闭子概形时，该恒等式依定义成立。
对 $X$ 上一般的 $k$-循环 $\alpha$，写成
$\alpha = \sum n_a[W_a]$，其中 $n_a \not = 0$，且
$i_a : W_a \to X$ 是两两不同的整闭子概形。于是族
$\{W_a\}$ 在 $X$ 上局部有限。令
$P_a = \pi^{-1}(W_a) = \mathbf{P}(\mathcal{E}|_{W_a})$，
并记 $i'_a : P_a \to P$ 为相应的闭浸入。考虑纤维积图
$$
\xymatrix{
P' \ar@{=}[r] \ar[d]_{\pi'} &
\coprod P_a \ar[d]_{\coprod \pi_a} \ar[r]_{\coprod i'_a} &
P \ar[d]^\pi \\
X' \ar@{=}[r] &
\coprod W_a \ar[r]^{\coprod i_a} &
X
}
$$
态射 $p : X' \to X$ 是正常的。此外，
$\pi' : P' \to X'$ 连同可逆层
$\mathcal{O}_{P'}(1) = \coprod \mathcal{O}_{P_a}(1)$
是与 $\mathcal{E}' = p^*\mathcal{E}$ 关联的射影丛；该可逆层也正是
$\mathcal{O}_P(1)$ 的拉回。依定义，
$$
c_j(\mathcal{E}) \cap [\alpha]
=
\sum i_{a, *}(c_j(\mathcal{E}|_{W_a}) \cap [W_a]).
$$
记 $\beta_{a, j} = c_j(\mathcal{E}|_{W_a}) \cap [W_a]$，
它是 $\CH_{k - j}(W_a)$ 的元素。根据定义，对每个 $a$ 有
$$
\sum\nolimits_{i = 0}^r
(-1)^i c_1(\mathcal{O}_{P_a}(1))^i \cap \pi_a^*(\beta_{a, r - i}) = 0
$$
因而显然有
$$
\sum\nolimits_{i = 0}^r
(-1)^i c_1(\mathcal{O}_{P'}(1))^i \cap (\pi')^*(\beta_{r - i}) = 0
$$
其中 $\beta_j = \sum n_a\beta_{a, j} \in \CH_{k - j}(X')$。
记 $p' : P' \to P$ 为态射 $\coprod i'_a$。由
引理 \ref{chow-lemma-flat-pullback-proper-pushforward}，有
$\pi^*p_*\beta_j = p'_*(\pi')^*\beta_j$。
再由引理 \ref{chow-lemma-pushforward-cap-c1} 的投影公式，得到
$$
\sum\nolimits_{i = 0}^r
(-1)^i c_1(\mathcal{O}_P(1))^i \cap \pi^*(p_*\beta_j) = 0
$$
由于 $p_*\beta_j$ 表示 $c_j(\mathcal{E}) \cap \alpha$，故得证。
\end{proof}

\noindent
下文将始终使用陈类的这一刻画来证明更多性质。

\begin{lemma}
\label{chow-lemma-cap-chern-class-factors-rational-equivalence}
设 $(S, \delta)$ 如情形 \ref{chow-situation-setup}。
设 $X$ 在 $S$ 上局部有限型。设 $\mathcal{E}$ 是 $X$ 上秩为
$r$ 的有限局部自由层。若 $\alpha \sim_{rat} \beta$ 是 $X$ 上
有理等价的 $k$-循环，则在 $\CH_{k - j}(X)$ 中有
$c_j(\mathcal{E}) \cap \alpha = c_j(\mathcal{E}) \cap \beta$。
\end{lemma}

\begin{proof}
由引理 \ref{chow-lemma-determine-intersections}，元素
$\alpha_j = c_j(\mathcal{E}) \cap \alpha$（$j \geq 1$）与
$\beta_j = c_j(\mathcal{E}) \cap \beta$（$j \geq 1$），由与
$\mathcal{E}$ 关联的射影丛的 Chow 群中的{\it 同一个}方程唯一确定。
（这当然用到了平坦拉回与有理等价相容这一事实，见
引理 \ref{chow-lemma-flat-pullback-rational-equivalence}。）因此它们相等。
\end{proof}

\noindent
换言之，与有限局部自由层的陈类作帽积可下降到有理等价，从而给出映射
$$
c_j(\mathcal{E}) \cap - : \CH_k(X) \to \CH_{k - j}(X).
$$
下一步是证明陈类是双变类，见
定义 \ref{chow-definition-bivariant-class}。

\begin{lemma}
\label{chow-lemma-pushforward-cap-cj}
设 $(S, \delta)$ 如情形 \ref{chow-situation-setup}。
设 $X$、$Y$ 在 $S$ 上局部有限型。设 $\mathcal{E}$ 是 $X$ 上秩为
$r$ 的有限局部自由层。设 $p : X \to Y$ 是正常态射，
设 $\alpha$ 是 $X$ 上的 $k$-循环，并设 $\mathcal{E}$ 是 $Y$ 上
的有限局部自由层。则
$$
p_*(c_j(p^*\mathcal{E}) \cap \alpha) = c_j(\mathcal{E}) \cap p_*\alpha
$$
\end{lemma}

\begin{proof}
设 $(\pi : P \to Y, \mathcal{O}_P(1))$ 是与 $\mathcal{E}$ 关联的射影丛。
则 $P_X = X \times_Y P$ 是与 $p^*\mathcal{E}$ 关联的射影丛，且
$\mathcal{O}_{P_X}(1)$ 是 $\mathcal{O}_P(1)$ 的拉回。记
$\alpha_j = c_j(p^*\mathcal{E}) \cap \alpha$，于是
$\alpha_0 = \alpha$。由引理 \ref{chow-lemma-determine-intersections}，有
$$
\sum\nolimits_{i = 0}^r
(-1)^i c_1(\mathcal{O}_P(1))^i \cap
\pi_X^*(\alpha_{r - i}) = 0
$$
在 $P_X$ 的 Chow 群中成立。考虑纤维积图
$$
\xymatrix{
P_X \ar[r]_-{p'} \ar[d]_{\pi_X} & P \ar[d]^\pi \\
X \ar[r]^p & Y
}
$$
对上述等式施以正常推前 $p'_*$
（引理 \ref{chow-lemma-proper-pushforward-rational-equivalence}）。
利用引理 \ref{chow-lemma-pushforward-cap-c1} 和
\ref{chow-lemma-flat-pullback-proper-pushforward}，得到
$$
\sum\nolimits_{i = 0}^r
(-1)^i c_1(\mathcal{O}_P(1))^i \cap
\pi^*(p_*\alpha_{r - i}) = 0
$$
在 $P$ 的 Chow 群中成立。由
引理 \ref{chow-lemma-determine-intersections} 的刻画即得结论。
\end{proof}

\begin{lemma}
\label{chow-lemma-flat-pullback-cap-cj}
设 $(S, \delta)$ 如情形 \ref{chow-situation-setup}。
设 $X$、$Y$ 在 $S$ 上局部有限型。设 $\mathcal{E}$ 是 $Y$ 上秩为
$r$ 的有限局部自由层。设 $f : X \to Y$ 是相对维数为 $r$
的平坦态射，设 $\alpha$ 是 $Y$ 上的 $k$-循环。则
$$
f^*(c_j(\mathcal{E}) \cap \alpha) = c_j(f^*\mathcal{E}) \cap f^*\alpha
$$
\end{lemma}

\begin{proof}
记 $\alpha_j = c_j(\mathcal{E}) \cap \alpha$，于是
$\alpha_0 = \alpha$。由引理 \ref{chow-lemma-determine-intersections}，有
$$
\sum\nolimits_{i = 0}^r
(-1)^i c_1(\mathcal{O}_P(1))^i \cap
\pi^*(\alpha_{r - i}) = 0
$$
在与 $\mathcal{E}$ 关联的射影丛
$(\pi : P \to Y, \mathcal{O}_P(1))$ 的 Chow 群中成立。
考虑纤维积图
$$
\xymatrix{
P_X = \mathbf{P}(f^*\mathcal{E}) \ar[r]_-{f'} \ar[d]_{\pi_X} &
P \ar[d]^\pi \\
X \ar[r]^f & Y
}
$$
注意 $\mathcal{O}_{P_X}(1)$ 是 $\mathcal{O}_P(1)$ 的拉回。
对上述方程施以平坦拉回 $(f')^*$
（引理 \ref{chow-lemma-flat-pullback-rational-equivalence}）。由
引理 \ref{chow-lemma-flat-pullback-cap-c1} 和
\ref{chow-lemma-compose-flat-pullback} 可知
$$
\sum\nolimits_{i = 0}^r
(-1)^i c_1(\mathcal{O}_{P_X}(1))^i \cap
\pi_X^*(f^*\alpha_{r - i}) = 0
$$
在 $P_X$ 的 Chow 群中成立。由
引理 \ref{chow-lemma-determine-intersections} 的刻画即得结论。
\end{proof}

\begin{lemma}
\label{chow-lemma-cap-chern-class-commutes-with-gysin}
设 $(S, \delta)$ 如情形 \ref{chow-situation-setup}。
设 $X$ 在 $S$ 上局部有限型。设 $\mathcal{E}$ 是 $X$ 上秩为
$r$ 的有限局部自由层。设 $(\mathcal{L}, s, i : D \to X)$
如定义 \ref{chow-definition-gysin-homomorphism}。
则对所有 $\alpha \in \CH_k(X)$ 有
$c_j(\mathcal{E}|_D) \cap i^*\alpha = i^*(c_j(\mathcal{E}) \cap \alpha)$。
\end{lemma}

\begin{proof}
记 $\alpha_j = c_j(\mathcal{E}) \cap \alpha$，于是
$\alpha_0 = \alpha$。由引理 \ref{chow-lemma-determine-intersections}，有
$$
\sum\nolimits_{i = 0}^r
(-1)^i c_1(\mathcal{O}_P(1))^i \cap
\pi^*(\alpha_{r - i}) = 0
$$
在与 $\mathcal{E}$ 关联的射影丛
$(\pi : P \to X, \mathcal{O}_P(1))$ 的 Chow 群中成立。
考虑纤维积图
$$
\xymatrix{
P_D = \mathbf{P}(\mathcal{E}|_D) \ar[r]_-{i'} \ar[d]_{\pi_D} &
P \ar[d]^\pi \\
D \ar[r]^i & X
}
$$
注意 $\mathcal{O}_{P_D}(1)$ 是 $\mathcal{O}_P(1)$ 的拉回。
对上述方程施以 Gysin 映射 $(i')^*$
（引理 \ref{chow-lemma-gysin-factors}）。应用引理 \ref{chow-lemma-gysin-commutes-cap-c1} 和
\ref{chow-lemma-gysin-flat-pullback}，得到
$$
\sum\nolimits_{i = 0}^r
(-1)^i c_1(\mathcal{O}_{P_D}(1))^i \cap
\pi_D^*(i^*\alpha_{r - i}) = 0
$$
在 $P_D$ 的 Chow 群中成立。由
引理 \ref{chow-lemma-determine-intersections} 的刻画即得结论。
\end{proof}

\noindent
至此已有足够材料证明，与陈类作帽积定义了一个双变类。

\begin{lemma}
\label{chow-lemma-cap-cp-bivariant}
设 $(S, \delta)$ 如情形 \ref{chow-situation-setup}。
设 $X$ 在 $S$ 上局部有限型。设 $\mathcal{E}$ 是秩为 $r$ 的局部自由
$\mathcal{O}_X$-模，并设 $0 \leq p \leq r$。则把
$f : X' \to X$ 映为
$c_p(f^*\mathcal{E}) \cap - : \CH_k(X') \to \CH_{k - p}(X')$
的规则，是一个次数为 $p$ 的双变类。
\end{lemma}

\begin{proof}
这立即来自引理 \ref{chow-lemma-cap-chern-class-factors-rational-equivalence}、
\ref{chow-lemma-pushforward-cap-cj}、
\ref{chow-lemma-flat-pullback-cap-cj}、
\ref{chow-lemma-cap-chern-class-commutes-with-gysin}
以及定义 \ref{chow-definition-bivariant-class}。
\end{proof}

\noindent
借助这个引理，可以如下定义有限局部自由模的陈类。

\begin{definition}
\label{chow-definition-chern-classes-final}
设 $(S, \delta)$ 如情形 \ref{chow-situation-setup}。
设 $X$ 在 $S$ 上局部有限型。设 $\mathcal{E}$ 是秩为 $r$ 的局部自由
$\mathcal{O}_X$-模。对 $i = 0, \ldots, r$，$\mathcal{E}$ 的
{\it 第 $i$ 陈类}是引理 \ref{chow-lemma-cap-cp-bivariant}
所构造的次数为 $i$ 的双变类
$c_i(\mathcal{E}) \in A^i(X)$。$\mathcal{E}$ 的
{\it 全陈类}是形式和
$$
c(\mathcal{E}) = 
c_0(\mathcal{E}) + c_1(\mathcal{E}) + \ldots + c_r(\mathcal{E})
$$
并将它视为 $X$ 上的非齐次双变类。
\end{definition}

\noindent
由定义 \ref{chow-definition-cap-chern-classes} 后的注记，
若 $\mathcal{E}$ 可逆，则这个定义与
定义 \ref{chow-definition-first-chern-class} 一致。
下面说明陈类位于双变 Chow 上同调环 $A^*(X)$ 的中心。

\begin{lemma}
\label{chow-lemma-cap-commutative-chern}
设 $(S, \delta)$ 如情形 \ref{chow-situation-setup}。
设 $X$ 在 $S$ 上局部有限型。设 $\mathcal{E}$ 是秩为 $r$ 的局部自由
$\mathcal{O}_X$-模。则
\begin{enumerate}
\item $c_j(\mathcal{E}) \in A^j(X)$ 位于 $A^*(X)$ 的中心；
\item 若 $f : X' \to X$ 局部有限型且 $c \in A^*(X' \to X)$，则
$c \circ c_j(\mathcal{E}) = c_j(f^*\mathcal{E}) \circ c$。
\end{enumerate}
特别地，若 $\mathcal{F}$ 是 $X$ 上另一个秩为 $s$ 的局部自由
$\mathcal{O}_X$-模，则对所有 $\alpha \in \CH_k(X)$，
$$
c_i(\mathcal{E}) \cap c_j(\mathcal{F}) \cap \alpha
=
c_j(\mathcal{F}) \cap c_i(\mathcal{E}) \cap \alpha
$$
作为 $\CH_{k - i - j}(X)$ 的元素成立。
\end{lemma}

\begin{proof}
(2) 推出 (1) 是立即的。设 $\alpha \in \CH_k(X)$，并记
$\alpha_j = c_j(\mathcal{E}) \cap \alpha$，于是
$\alpha_0 = \alpha$。由引理 \ref{chow-lemma-determine-intersections}，有
$$
\sum\nolimits_{i = 0}^r
(-1)^i c_1(\mathcal{O}_P(1))^i \cap
\pi^*(\alpha_{r - i}) = 0
$$
在与 $\mathcal{E}$ 关联的射影丛
$(\pi : P \to Y, \mathcal{O}_P(1))$ 的 Chow 群中成立。
记 $\pi' : P' \to X'$ 为 $\pi$ 沿 $f$ 的基变换。利用
引理 \ref{chow-lemma-c1-center} 及双变类的性质，得到
\begin{align*}
0 & = c \cap \left(\sum\nolimits_{i = 0}^r
(-1)^i c_1(\mathcal{O}_P(1))^i \cap
\pi^*(\alpha_{r - i})\right) \\
& =
\sum\nolimits_{i = 0}^r
(-1)^i c_1(\mathcal{O}_{P'}(1))^i \cap
(\pi')^*(c \cap \alpha_{r - i})
\end{align*}
在 $P'$ 的 Chow 群中成立（略去计算）。因此，由
引理 \ref{chow-lemma-determine-intersections} 的刻画，
$c \cap \alpha_j$ 等于
$c_j(f^*\mathcal{E}) \cap (c \cap \alpha)$。引理得证。
\end{proof}

\begin{remark}
\label{chow-remark-extend-to-finite-locally-free}
设 $(S, \delta)$ 如情形 \ref{chow-situation-setup}。
设 $X$ 在 $S$ 上局部有限型。设 $\mathcal{E}$ 是有限局部自由
$\mathcal{O}_X$-模。即使 $\mathcal{E}$ 的秩不是常数，仍可定义
$\mathcal{E}$ 的陈类。
事实上，此时可写成
$$
X = X_0 \amalg X_1 \amalg X_2 \amalg \ldots
$$
其中 $X_r \subset X$ 是 $\mathcal{E}$ 秩为 $r$ 的开闭子空间。
由引理 \ref{chow-lemma-disjoint-decomposition-bivariant}，有
$A^p(X) = \prod A^p(X_r)$。因此，可以把 $c_p(\mathcal{E})$
定义为各 $A^p(X_r)$ 中的类 $c_p(\mathcal{E}|_{X_r})$ 之积。
明确地说，若 $X' \to X$ 是局部有限型态射，则通过拉回得到
$X'$ 的相应分解，并由定义有
$$
\CH_*(X') = \prod\nolimits_{r \geq 0} \CH_*(X'_r)
$$
此时 $c_p(\mathcal{E}) \in A^p(X)$ 是保持这些直积分解，
并在各因子上由已经定义的运算
$c_i(\mathcal{E}|_{X_r}) \cap -$ 作用的双变类。
注意，在这种情形下，可能有无穷多个 $p$ 使
$c_p(\mathcal{E})$ 非零。因此，全陈类是完备化双变上同调环中的元素
$$
c(\mathcal{E}) =
c_0(\mathcal{E}) + c_1(\mathcal{E}) + c_2(\mathcal{E}) + \ldots
\in A^*(X)^\wedge
$$
参见注 \ref{chow-remark-completion-bivariant}。
在这里，把 $\mathcal{E}$ 的“秩”定义为元素
$r(\mathcal{E}) \in A^0(X)$，即把
$(\alpha_r) \in \prod \CH_*(X'_r)$ 映到
$(r\alpha_r) \in \prod \CH_*(X'_r)$ 的双变运算。
仍然有 $c_p(\mathcal{E})$ 和 $r(\mathcal{E})$ 位于
$A^*(X)$ 的中心。
\end{remark}

\begin{remark}
\label{chow-remark-top-chern-class}
设 $(S, \delta)$ 如情形 \ref{chow-situation-setup}，设 $X$
在 $S$ 上局部有限型，并设 $\mathcal{E}$ 是有限局部自由
$\mathcal{O}_X$-模。一般地，如
注 \ref{chow-remark-extend-to-finite-locally-free} 写成
$X = \coprod X_r$。若仅有有限多个 $X_r$ 非空，则可令
$$
c_{top}(\mathcal{E}) = \sum\nolimits_r c_r(\mathcal{E}|_{X_r})
\in A^*(X) = \bigoplus A^*(X_r)
$$
其中等式来自引理 \ref{chow-lemma-disjoint-decomposition-bivariant}。
若有无穷多个 $X_r$ 非空，则用同一记号表示
$$
c_{top}(\mathcal{E}) = \prod c_r(\mathcal{E}|_{X_r})
\in \prod A^r(X_r) \subset A^*(X)^\wedge
$$
记号参见注 \ref{chow-remark-completion-bivariant}。
\end{remark}











\section{陈类之间的多项式关系}
\label{chow-section-relations-chern-classes}

\noindent
设 $(S, \delta)$ 如情形 \ref{chow-situation-setup}，设 $X$ 在 $S$
上局部有限型，并设 $\mathcal{E}_i$ 是 $X$ 上有限局部自由层的有限族。
由引理 \ref{chow-lemma-cap-commutative-chern}，陈类
$$
c_j(\mathcal{E}_i) \in A^*(X)
$$
生成 Chow 上同调代数 $A^*(X)$ 的一个交换（甚至中心的）
$\mathbf{Z}$-子代数。因此，可以说明这些陈类的多项式为零，
或者两个多项式相等，具体是什么意思。例如，等式
$c_1(\mathcal{E}_1)^5 + c_2(\mathcal{E}_2)c_3(\mathcal{E}_3) = 0$
表示：对所有局部有限型态射 $f : Y \to X$，运算
$$
\CH_k(Y) \longrightarrow \CH_{k - 5}(Y), \quad
\alpha \longmapsto
c_1(\mathcal{E}_1)^5 \cap \alpha +
c_2(\mathcal{E}_2) \cap c_3(\mathcal{E}_3) \cap \alpha
$$
都为零。由引理 \ref{chow-lemma-bivariant-zero}，这等价于如下要求：
对任意态射 $f : Y \to X$，其中 $Y$ 是 $S$ 上局部有限型的整概形，循环
$$
c_1(\mathcal{E}_1)^5 \cap [Y] +
c_2(\mathcal{E}_2) \cap c_3(\mathcal{E}_3) \cap [Y]
$$
在 $\CH_{\dim(Y) - 5}(Y)$ 中为零。

\medskip\noindent
一个具体例子是引理 \ref{chow-lemma-c1-cap-additive} 中证明的关系
$$
c_1(\mathcal{L} \otimes_{\mathcal{O}_X} \mathcal{N})
=
c_1(\mathcal{L}) + c_1(\mathcal{N})
$$
更一般地，下面说明任意局部自由层与可逆层张量时的情形。

\begin{lemma}
\label{chow-lemma-chern-classes-E-tensor-L}
设 $(S, \delta)$ 如情形 \ref{chow-situation-setup}。
设 $X$ 在 $S$ 上局部有限型。设 $\mathcal{E}$ 是 $X$ 上秩为 $r$
的有限局部自由层，并设 $\mathcal{L}$ 是 $X$ 上的可逆层。则在
$A^*(X)$ 中有
\begin{equation}
\label{chow-equation-twist}
c_i({\mathcal E} \otimes {\mathcal L})
=
\sum\nolimits_{j = 0}^i
\binom{r - i + j}{j} c_{i - j}({\mathcal E}) c_1({\mathcal L})^j
\end{equation}
\end{lemma}

\begin{proof}
该式应对任意三元组 $(X, \mathcal{E}, \mathcal{L})$ 成立。
特别地，它应在 $X$ 整时成立；由引理 \ref{chow-lemma-bivariant-zero}，
只须对这样的 $X$ 证明与 $[X]$ 作帽积后的等式。因此设 $X$ 整。
令 $(\pi : P \to X, \mathcal{O}_P(1))$，resp.\ $(\pi' : P' \to X, \mathcal{O}_{P'}(1))$
分别为与
$\mathcal{E}$，resp.\ $\mathcal{E} \otimes \mathcal{L}$
关联的射影空间丛。考虑规范态射
$$
\xymatrix{
P \ar[rd]_\pi \ar[rr]_g & & P' \ar[ld]^{\pi'} \\
& X &
}
$$
参见《概形的构造》引理 \ref{constructions-lemma-twisting-and-proj}。
它满足
$g^*\mathcal{O}_{P'}(1)
= \mathcal{O}_P(1) \otimes \pi^* {\mathcal L}$。
这意味着
$$
\sum\nolimits_{i = 0}^r
(-1)^i
(\xi + x)^i \cap \pi^*(c_{r - i}(\mathcal{E} \otimes \mathcal{L}) \cap [X])
=
0
$$
在 $\CH_*(P)$ 中成立，其中 $\xi$ 表示
$c_1(\mathcal{O}_P(1))$，而 $x$ 表示
$c_1(\pi^*\mathcal{L})$。由简单的代数运算，这等价于
$$
\sum\nolimits_{i = 0}^r
(-1)^i \xi^i \left(
\sum\nolimits_{j = i}^r
(-1)^{j - i}
\binom{j}{i}
x^{j - i} \cap
\pi^*(c_{r - j}(\mathcal{E} \otimes \mathcal{L}) \cap [X])
\right)
=
0
$$
与 Equation (\ref{chow-equation-chern-classes}) 比较，可得
$$
c_{r - i}(\mathcal{E}) \cap [X] =
\sum\nolimits_{j = i}^r
\binom{j}{i}
(-c_1(\mathcal{L}))^{j - i} \cap
c_{r - j}(\mathcal{E} \otimes \mathcal{L}) \cap [X]
$$
重新整理（消去负号并重新编号）便得到所需关系。
\end{proof}

\noindent
(\ref{chow-equation-twist}) 的一些具体情形为
\begin{align*}
c_1(\mathcal{E} \otimes \mathcal{L})
& =
c_1(\mathcal{E}) +
r c_1(\mathcal{L}) \\
c_2(\mathcal{E} \otimes \mathcal{L})
& =
c_2(\mathcal{E}) +
(r - 1) c_1(\mathcal{E}) c_1(\mathcal{L}) +
\binom{r}{2} c_1(\mathcal{L})^2 \\
c_3(\mathcal{E} \otimes \mathcal{L})
& =
c_3(\mathcal{E}) +
(r - 2) c_2(\mathcal{E})c_1(\mathcal{L}) +
\binom{r - 1}{2} c_1(\mathcal{E})c_1(\mathcal{L})^2 +
\binom{r}{3} c_1(\mathcal{L})^3
\end{align*}








\section{陈类的可加性}
\label{chow-section-additivity-chern-classes}

\noindent
所有预备引理都从最终结果立即推出。

\begin{lemma}
\label{chow-lemma-get-rid-of-trivial-subbundle}
设 $(S, \delta)$ 如情形 \ref{chow-situation-setup}。
设 $X$ 在 $S$ 上局部有限型。设 $\mathcal{E}$、$\mathcal{F}$ 是
$X$ 上秩分别为 $r$、$r - 1$ 的有限局部自由层，且嵌入短正合列
$$
0 \to \mathcal{O}_X \to \mathcal{E} \to \mathcal{F} \to 0
$$
则在 $A^*(X)$ 中有
$$
c_r(\mathcal{E}) = 0, \quad
c_j(\mathcal{E}) = c_j(\mathcal{F}), \quad j = 0, \ldots, r - 1
$$
\end{lemma}

\begin{proof}
由引理 \ref{chow-lemma-bivariant-zero}，只须证明：若 $X$ 整，则
$c_j(\mathcal{E}) \cap [X] = c_j(\mathcal{F}) \cap [X]$。
令 $(\pi : P \to X, \mathcal{O}_P(1))$，resp.\ $(\pi' : P' \to X, \mathcal{O}_{P'}(1))$
分别表示与 $\mathcal{E}$，resp.\ $\mathcal{F}$ 关联的射影空间丛。
满射 $\mathcal{E} \to \mathcal{F}$ 给出 $X$ 上的闭浸入
$$
i : P' \longrightarrow P
$$
此外，元素
$1 \in \Gamma(X, \mathcal{O}_X) \subset \Gamma(X, \mathcal{E})$
给出整体截面 $s \in \Gamma(P, \mathcal{O}_P(1))$，其零点集恰为
$P'$。因此 $P'$ 是 $P$ 上的有效 Cartier 除子，并满足
$\mathcal{O}_P(P') \cong \mathcal{O}_P(1)$。故由
引理 \ref{chow-lemma-relative-effective-cartier}，对 $X$ 上任意循环类
$\alpha$ 有
$$
c_1(\mathcal{O}_P(1)) \cap \pi^*\alpha = i_*((\pi')^*\alpha)
$$
由引理 \ref{chow-lemma-determine-intersections}，
$\alpha_j = c_j(\mathcal{F}) \cap [X]$（$j = 0, \ldots, r - 1$）满足
$$
\sum\nolimits_{j = 0}^{r - 1} (-1)^jc_1(\mathcal{O}_{P'}(1))^j
\cap (\pi')^*\alpha_j = 0
$$
把它推前到 $P$，并使用上述注记及
引理 \ref{chow-lemma-pushforward-cap-c1}，得到
$$
\sum\nolimits_{j = 0}^{r - 1}
(-1)^j c_1(\mathcal{O}_P(1))^{j + 1}
\cap \pi^*\alpha_j = 0
$$
由引理 \ref{chow-lemma-determine-intersections} 中的唯一性，可得
$c_r(\mathcal{E}) \cap [X] = 0$，并且
$c_j(\mathcal{E}) \cap [X] = \alpha_j = c_j(\mathcal{F}) \cap [X]$
对 $j = 0, \ldots, r - 1$ 成立。引理得证。
\end{proof}

\begin{lemma}
\label{chow-lemma-additivity-invertible-subsheaf}
设 $(S, \delta)$ 如情形 \ref{chow-situation-setup}。
设 $X$ 在 $S$ 上局部有限型。设 $\mathcal{E}$、$\mathcal{F}$ 是
$X$ 上秩分别为 $r$、$r - 1$ 的有限局部自由层，且嵌入短正合列
$$
0 \to \mathcal{L} \to \mathcal{E} \to \mathcal{F} \to 0
$$
其中 $\mathcal{L}$ 是可逆层。则在 $A^*(X)$ 中有
$$
c(\mathcal{E}) = c(\mathcal{L}) c(\mathcal{F})
$$
\end{lemma}

\begin{proof}
这个关系实际上就是
$c_i(\mathcal{E}) = c_i(\mathcal{F}) + c_1(\mathcal{L})c_{i - 1}(\mathcal{F})$。
由引理 \ref{chow-lemma-get-rid-of-trivial-subbundle}，对
$j = 0, \ldots, r$ 有
$c_j(\mathcal{E} \otimes \mathcal{L}^{\otimes -1})
= c_j(\mathcal{F} \otimes \mathcal{L}^{\otimes -1})$，其中约定
$c_r(\mathcal{F} \otimes \mathcal{L}^{-1}) = 0$。
应用引理 \ref{chow-lemma-chern-classes-E-tensor-L}，得到
$$
\sum_{j = 0}^i
\binom{r - i + j}{j} (-1)^j c_{i - j}({\mathcal E}) c_1({\mathcal L})^j
=
\sum_{j = 0}^i
\binom{r - 1 - i + j}{j} (-1)^j c_{i - j}({\mathcal F}) c_1({\mathcal L})^j
$$
令 $c_i(\mathcal{E}) = c_i(\mathcal{F}) + c_1(\mathcal{L})c_{i - 1}(\mathcal{F})$，
便给出这个方程的一个“解”。若证明这是唯一可能的解，引理即得。
这里略去验证。
\end{proof}

\begin{lemma}
\label{chow-lemma-additivity-chern-classes}
设 $(S, \delta)$ 如情形 \ref{chow-situation-setup}。
设 $X$ 是 $S$ 上局部有限型概形。假设 ${\mathcal E}$ 位于秩为
$r_i$ 的有限局部自由层 $\mathcal{E}_i$ 的正合列
$$
0
\to
{\mathcal E}_1
\to
{\mathcal E}
\to
{\mathcal E}_2
\to
0
$$
中。则全陈类在 $A^*(X)$ 中满足
$$
c({\mathcal E}) = c({\mathcal E}_1) c({\mathcal E}_2)
$$
\end{lemma}

\begin{proof}
由引理 \ref{chow-lemma-bivariant-zero}，可设 $X$ 整，并只须证明与
$[X]$ 作帽积后的恒等式。对 $r_1$ 归纳。当 $r_1 = 1$ 时，结论是
引理 \ref{chow-lemma-additivity-invertible-subsheaf}。设 $r_1 > 1$。
令 $(\pi : P \to X, \mathcal{O}_P(1))$ 表示与
$\mathcal{E}_1$ 关联的射影空间丛。注意
\begin{enumerate}
\item $\pi^* : \CH_*(X) \to \CH_*(P)$ 是单射；
\item $\pi^*\mathcal{E}_1$ 位于短正合列
$0 \to \mathcal{F} \to \pi^*\mathcal{E}_1 \to \mathcal{L} \to 0$
中，其中 $\mathcal{L}$ 可逆。
\end{enumerate}
第一点来自射影空间丛公式，第二点来自射影空间丛的定义。
（事实上 $\mathcal{L} = \mathcal{O}_P(1)$。）令
$Q = \pi^*\mathcal{E}/\mathcal{F}$，它位于正合列
$0 \to \mathcal{L} \to Q \to \pi^*\mathcal{E}_2 \to 0$ 中。
由归纳假设，
\begin{eqnarray*}
c(\pi^*\mathcal{E}) \cap [P]
& = &
c(\mathcal{F}) \cap c(\pi^*\mathcal{E}/\mathcal{F}) \cap [P] \\
& = &
c(\mathcal{F}) \cap c(\mathcal{L}) \cap c(\pi^*\mathcal{E}_2) \cap [P] \\
& = &
c(\pi^*\mathcal{E}_1) \cap c(\pi^*\mathcal{E}_2) \cap [P]
\end{eqnarray*}
由于 $[P] = \pi^*[X]$，由
引理 \ref{chow-lemma-flat-pullback-cap-cj} 即得结论。
\end{proof}

\begin{lemma}
\label{chow-lemma-chern-filter-by-linebundles}
设 $(S, \delta)$ 如情形 \ref{chow-situation-setup}。
设 $X$ 在 $S$ 上局部有限型。设 ${\mathcal L}_i$（$i = 1, \ldots, r$）
是 $X$ 上的可逆 $\mathcal{O}_X$-模。设 $\mathcal{E}$ 是赋有滤过
$$
0 = \mathcal{E}_0 \subset \mathcal{E}_1 \subset \mathcal{E}_2
\subset \ldots \subset \mathcal{E}_r = \mathcal{E}
$$
的局部自由秩 $\mathcal{O}_X$-模，且
$\mathcal{E}_i/\mathcal{E}_{i - 1} \cong \mathcal{L}_i$。
令 $c_1({\mathcal L}_i) = x_i$。则在 $A^*(X)$ 中有
$$
c(\mathcal{E})
=
\prod\nolimits_{i = 1}^r (1 + x_i)
$$
\end{lemma}

\begin{proof}
应用引理 \ref{chow-lemma-additivity-invertible-subsheaf} 并作归纳。
\end{proof}




\section{零维循环的次数}
\label{chow-section-degree-zero-cycles}

\noindent
先定义域上的正常概形上零维循环的次数。一种做法是如
引理 \ref{chow-lemma-spell-out-degree-zero-cycle} 直接定义，
再由引理 \ref{chow-lemma-curve-principal-divisor} 证明它良定义。
这里改用如下定义。

\begin{definition}
\label{chow-definition-degree-zero-cycle}
设 $k$ 是域（例 \ref{chow-example-field}），并设
$p : X \to \Spec(k)$ 正常。$X$ 上{\it 零维循环的次数}由正常推前
$$
p_* : \CH_0(X) \to \CH_0(\Spec(k))
$$
（引理 \ref{chow-lemma-proper-pushforward-rational-equivalence}）
与把 $[\Spec(k)]$ 映到 $1$ 的自然同构
$\CH_0(\Spec(k)) = \mathbf{Z}$ 复合给出。记作 $\deg(\alpha)$。
\end{definition}

\noindent
进一步把它写明如下。

\begin{lemma}
\label{chow-lemma-spell-out-degree-zero-cycle}
设 $k$ 是域，$X$ 在 $k$ 上正常，并设
$\alpha = \sum n_i[Z_i]$ 属于 $Z_0(X)$。则
$$
\deg(\alpha) = \sum n_i\deg(Z_i)
$$
其中 $\deg(Z_i)$ 是 $Z_i \to \Spec(k)$ 的次数，即
$\deg(Z_i) = \dim_k \Gamma(Z_i, \mathcal{O}_{Z_i})$。
\end{lemma}

\begin{proof}
这正是正常推前的定义
（定义 \ref{chow-definition-proper-pushforward}）。
\end{proof}

\noindent
下面把它与《代数簇》第 \ref{varieties-section-divisors-curves} 节
所定义的域上 $1$-维正常概形上的向量丛次数联系起来。

\begin{lemma}
\label{chow-lemma-degree-vector-bundle}
设 $k$ 是域。设 $X$ 是 $k$ 上维数 $\leq 1$ 的正常概形。
设 $\mathcal{E}$ 是常秩有限局部自由 $\mathcal{O}_X$-模。则
$$
\deg(\mathcal{E}) = \deg(c_1(\mathcal{E}) \cap [X]_1)
$$
其中左端由《代数簇》定义 \ref{varieties-definition-degree-invertible-sheaf} 定义。
\end{lemma}

\begin{proof}
设 $C_i \subset X$（$i = 1, \ldots, t$）是赋予诱导既约概形结构的
$1$ 维不可约分支，并设 $m_i$ 是 $C_i$ 在 $X$ 中的重数。则
$[X]_1 = \sum m_i[C_i]$，且 $c_1(\mathcal{E}) \cap [X]_1$
是循环 $m_i c_1(\mathcal{E}|_{C_i}) \cap [C_i]$ 的推前之和。
由《代数簇》引理 \ref{varieties-lemma-degree-in-terms-of-components}，
$\mathcal{E}$ 的次数也有类似分解，故只须证明 $X$ 是 $k$ 上正常曲线的情形。

\medskip\noindent
假设 $X$ 是 $k$ 上正常曲线。由《除子》引理 \ref{divisors-lemma-filter-after-modification}，存在修改
$f : X' \to X$，使 $f^*\mathcal{E}$ 有一个滤过，其逐次商是可逆
$\mathcal{O}_{X'}$-模。由于 $f_*[X']_1 = [X]_1$，由
引理 \ref{chow-lemma-pushforward-cap-cj} 可得
$$
\deg(c_1(\mathcal{E}) \cap [X]_1) = \deg(c_1(f^*\mathcal{E}) \cap [X']_1)
$$
由《代数簇》引理 \ref{varieties-lemma-degree-birational-pullback}，
次数也满足类似关系，因此归约到 $\mathcal{E}$ 有逐次商为可逆
$\mathcal{O}_X$-模的滤过这一情形。此时利用次数的可加性
（《代数簇》引理 \ref{varieties-lemma-degree-additive}）以及第一陈类的可加性
（引理 \ref{chow-lemma-additivity-chern-classes}），可归约到下一段的情形。

\medskip\noindent
假设 $X$ 是 $k$ 上正常曲线，且 $\mathcal{E}$ 是可逆
$\mathcal{O}_X$-模。由《除子》引理 \ref{divisors-lemma-quasi-projective-Noetherian-pic-effective-Cartier}，
$\mathcal{E}$ 同构于
$\mathcal{O}_X(D) \otimes \mathcal{O}_X(D')^{\otimes -1}$，
其中 $D, D'$ 是 $X$ 上的有效 Cartier 除子（这里还用到 $X$ 射影；
例如参见《代数簇》引理 \ref{varieties-lemma-dim-1-proper-projective}）。
利用可逆层张量积下次数的可加性
（《代数簇》引理 \ref{varieties-lemma-degree-tensor-product}）
以及可逆层张量积下 $c_1$ 的可加性
（引理 \ref{chow-lemma-c1-cap-additive} 或 \ref{chow-lemma-chern-classes-E-tensor-L}），
归约到 $\mathcal{E} = \mathcal{O}_X(D)$ 的情形。此时左端给出
$\deg(D)$（《代数簇》引理 \ref{varieties-lemma-degree-effective-Cartier-divisor}），而由
引理 \ref{chow-lemma-geometric-cap}，右端给出 $\deg([D]_0)$。
根据定义，
$$
[D]_0 = \sum\nolimits_{x \in D}
\text{length}_{\mathcal{O}_{X, x}}(\mathcal{O}_{D, x}) [x] =
\sum\nolimits_{x \in D}
\text{length}_{\mathcal{O}_{D, x}}(\mathcal{O}_{D, x}) [x]
$$
因此
$$
\deg([D]_0) = \sum\nolimits_{x \in D}
\text{length}_{\mathcal{O}_{D, x}}(\mathcal{O}_{D, x}) [\kappa(x) : k] =
\dim_k \Gamma(D, \mathcal{O}_D) = \deg(D)
$$
倒数第二个等式由《交换代数》引理 \ref{algebra-lemma-pushdown-module} 以及 $D$ 为仿射这一事实得到。
\end{proof}

\noindent
最后，把以上内容与《代数簇》第 \ref{varieties-section-num} 节所定义的数值交数联系起来。

\begin{lemma}
\label{chow-lemma-degrees-and-numerical-intersections}
设 $k$ 是域，$X$ 是 $k$ 上正常概形，$Z \subset X$ 是维数为
$d$ 的闭子概形，并设 $\mathcal{L}_1, \ldots, \mathcal{L}_d$
是可逆 $\mathcal{O}_X$-模。则
$$
(\mathcal{L}_1 \cdots \mathcal{L}_d \cdot Z) =
\deg(
c_1(\mathcal{L}_1) \cap \ldots \cap c_1(\mathcal{L}_d) \cap [Z]_d)
$$
其中左端由《代数簇》定义 \ref{varieties-definition-intersection-number} 定义。特别地，
若 $\mathcal{L}$ 是充足可逆 $\mathcal{O}_X$-模，则
$$
\deg_\mathcal{L}(Z) = \deg(c_1(\mathcal{L})^d \cap [Z]_d)
$$
\end{lemma}

\begin{proof}
对 $d$ 归纳。若 $d = 0$，结论由《代数簇》引理 \ref{varieties-lemma-chi-tensor-finite} 成立。设 $d > 0$。

\medskip\noindent
设 $Z_i \subset Z$（$i = 1, \ldots, t$）是赋予诱导既约概形结构的
$d$ 维不可约分支，并设 $m_i$ 是 $Z_i$ 在 $Z$ 中的重数。则
$[Z]_d = \sum m_i[Z_i]$，且
$c_1(\mathcal{L}_1) \cap \ldots \cap c_1(\mathcal{L}_d) \cap [Z]_d$
是循环
$m_i c_1(\mathcal{L}_1) \cap \ldots \cap c_1(\mathcal{L}_d) \cap [Z_i]$
之和。由《代数簇》引理 \ref{varieties-lemma-numerical-polynomial-leading-term}，
$(\mathcal{L}_1 \cdots \mathcal{L}_d \cdot Z)$ 也有类似分解，
故只须证明 $Z = X$ 是 $k$ 上 $d$ 维正常簇的情形。

\medskip\noindent
由 Chow 引理，存在双有理正常态射 $f : Y \to X$，其中 $Y$
在 $k$ 上 H-射影。参见《概形上同调》引理 \ref{coherent-lemma-chow-Noetherian} 和注 \ref{coherent-remark-chow-Noetherian}。此时，由《代数簇》引理 \ref{varieties-lemma-intersection-number-and-pullback} 有
$$
(f^*\mathcal{L}_1 \cdots f^*\mathcal{L}_d \cdot Y) =
(\mathcal{L}_1 \cdots \mathcal{L}_d \cdot X)
$$
并由引理 \ref{chow-lemma-pushforward-cap-c1} 有
$$
f_*(c_1(f^*\mathcal{L}_1) \cap \ldots \cap c_1(f^*\mathcal{L}_d) \cap [Y]) =
c_1(\mathcal{L}_1) \cap \ldots \cap c_1(\mathcal{L}_d) \cap [X]
$$
因此可用 $Y$ 替换 $X$，并设 $X$ 在 $k$ 上射影。

\medskip\noindent
若 $X$ 是正常的 $d$ 维射影簇，则由《除子》引理 \ref{divisors-lemma-quasi-projective-Noetherian-pic-effective-Cartier}，
可对某些有效 Cartier 除子 $D, D' \subset X$ 写成
$\mathcal{L}_1 = \mathcal{O}_X(D) \otimes \mathcal{O}_X(D')^{\otimes -1}$。
利用等式两端的可加性（《代数簇》引理 \ref{varieties-lemma-intersection-number-additive} 和
引理 \ref{chow-lemma-c1-cap-additive}），可归约到
$\mathcal{L}_1 = \mathcal{O}_X(D)$，其中 $D$ 是有效 Cartier 除子。
由《代数簇》引理 \ref{varieties-lemma-numerical-intersection-effective-Cartier-divisor}，有
$$
(\mathcal{L}_1 \cdots \mathcal{L}_d \cdot X) =
(\mathcal{L}_2 \cdots \mathcal{L}_d \cdot D)
$$
而由引理 \ref{chow-lemma-geometric-cap}，有
$$
c_1(\mathcal{L}_1) \cap \ldots \cap c_1(\mathcal{L}_d) \cap [X] =
c_1(\mathcal{L}_2) \cap \ldots \cap c_1(\mathcal{L}_d) \cap [D]_{d - 1}
$$
归纳假设即给出结论。
\end{proof}








\section{给定余维数的循环}
\label{chow-section-cycles-codimension}

\noindent
在某些情形下，所有循环的阿贝尔群还有一个由余维数给出的分次。

\begin{lemma}
\label{chow-lemma-locally-equidimensional}
设 $(S, \delta)$ 如情形 \ref{chow-situation-setup}。
设 $X$ 在 $S$ 上局部有限型。如第 \ref{chow-section-setup} 节，记
$\delta = \delta_{X/S}$。下列条件等价：
\begin{enumerate}
\item 存在分解 $X = \coprod_{n \in \mathbf{Z}} X_n$ 为开闭子概形，
使得只要 $\xi \in X_n$ 是 $X_n$ 某个不可约分支的一般点，
就有 $\delta(\xi) = n$。
\item 对所有 $x \in X$，存在 $x$ 的开邻域 $U \subset X$ 和整数
$n$，使得只要 $\xi \in U$ 是 $U$ 某个不可约分支的一般点，
就有 $\delta(\xi) = n$。
\item 对所有 $x \in X$，存在整数 $n_x$，使得对任意包含 $x$ 的
$X$ 的不可约分支的一般点 $\xi$，都有 $\delta(\xi) = n_x$。
\end{enumerate}
若 $X$ 正规或 Cohen--Macaulay，则这些条件成立\footnote{事实上，
只须 $X$ 是 $(S_2)$。可与《局部上同调》引理 \ref{local-cohomology-lemma-catenary-S2-equidimensional} 比较。}。
\end{lemma}

\begin{proof}
(1) $\Rightarrow$ (2) $\Rightarrow$ (3) 显然。反之，若 (3) 成立，
令 $X_n = \{x \in X \mid n_x = n\}$，便得到 (1) 中的分解。
事实上，$X_n$ 是开的：给定 $x$，所有经过 $x$ 的 $X$ 的不可约分支之并，
减去所有不经过 $x$ 的 $X$ 的不可约分支之并，是 $x$ 的一个开邻域。
若 $X$ 正规，则 $X$ 是整概形的不交并
（《概形的性质》引理 \ref{properties-lemma-normal-locally-Noetherian}），
故这些性质成立。若 $X$ 是 Cohen--Macaulay，则
$\delta' : X \to \mathbf{Z}$，$x \mapsto -\dim(\mathcal{O}_{X, x})$
是 $X$ 上的维数函数（见例 \ref{chow-example-CM-irreducible}）。
由于 $\delta - \delta'$ 局部常值
（《拓扑学》引理 \ref{topology-lemma-dimension-function-unique}），
且对 $X$ 的每个一般点 $\xi$ 都有 $\delta'(\xi) = 0$，故 (2) 成立。
\end{proof}

\noindent
设 $(S, \delta)$ 如情形 \ref{chow-situation-setup}。设 $X$ 在 $S$
上局部有限型，并满足引理 \ref{chow-lemma-locally-equidimensional}
中的等价条件。对整闭子概形 $Z \subset X$，有 $Z$ 在 $X$ 中的
余维数 $\text{codim}(Z, X)$，见《拓扑学》定义 \ref{topology-definition-codimension}。定义一个{\it 余维数 $p$ 的循环}
为 $X$ 上满足
$n_Z \not = 0 \Rightarrow \text{codim}(Z, X) = p$ 的循环
$\alpha = \sum n_Z[Z]$。所有余维数 $p$ 的循环所成的阿贝尔群记作
$Z^p(X)$。令 $X = \coprod X_n$ 为
引理 \ref{chow-lemma-locally-equidimensional} (1) 中的分解。
回忆循环被定义为局部有限和，显然
$$
Z^p(X) = \prod\nolimits_n Z_{n - p}(X_n)
$$
此外，$\prod_p Z^p(X) = \prod_k Z_k(X)$。现在可以完全照此前的方式
定义 $X$ 上余维数 $p$ 循环的有理等价；事实上，对满足
引理 \ref{chow-lemma-locally-equidimensional} 的 $X$，可以从头重新发展
给定余维数循环的整套理论。不过，这里直接把{\it 余维数 $p$ 循环的
Chow 群}定义为
$$
\CH^p(X) = \prod\nolimits_n \CH_{n - p}(X_n)
$$
如前，有 $\prod_p \CH^p(X) = \prod_k \CH_k(X)$。
若 $X$ 拟紧，则公式中的积为有限积（因而也是直和），且有
$\bigoplus_p \CH^p(X) = \bigoplus_k \CH_k(X)$。若 $X$ 还为有限维，
则这些群中仅有限多个非零。

\medskip\noindent
上文对 Chow 群证明的许多构造与结果，都有关于 Chow 群
$\CH^*(X)$ 的自然对应版本。证明方法是如
引理 \ref{chow-lemma-locally-equidimensional} 把有关概形分解成“等维”部分，
再对上述积分解中的各因子应用已经证明的结果。列举如下：
\begin{enumerate}
\item 若 $f : X \to Y$ 是 $S$ 上局部有限型概形之间的平坦态射，
且 $X$、$Y$ 满足引理 \ref{chow-lemma-locally-equidimensional} 的等价条件，
则平坦拉回给出映射
$$
f^* : \CH^p(Y) \to \CH^p(X)
$$
\item 若 $f : X \to Y$ 是 $S$ 上局部有限型概形之间的态射，
且 $X$、$Y$ 满足引理 \ref{chow-lemma-locally-equidimensional} 的等价条件，
则称 $f$ 的{\it 余维数}为 $r \in \mathbf{Z}$，如果对所有满足
$f(Z) \subset W$ 的不可约分支对 $Z \subset X$、$W \subset Y$，都有
$\dim_\delta(W) - \dim_\delta(Z) = r$。
\item 若 $f : X \to Y$ 是 $S$ 上局部有限型概形之间的正常态射，
且 $X$、$Y$ 满足引理 \ref{chow-lemma-locally-equidimensional} 的等价条件，
并且 $f$ 的余维数为 $r$，则正常推前是映射
$$
f_* : \CH^p(X) \to \CH^{p + r}(Y)
$$
\item 若 $f : X \to Y$ 是 $S$ 上局部有限型概形之间的态射，
且 $X$、$Y$ 满足引理 \ref{chow-lemma-locally-equidimensional} 的等价条件，
$f$ 的余维数为 $r$，并且 $c \in A^q(X \to Y)$，则 $c$ 诱导映射
$$
c \cap -  : \CH^p(Y) \to \CH^{p + q - r}(X)
$$
\item 若 $X$ 是 $S$ 上局部有限型概形，满足
引理 \ref{chow-lemma-locally-equidimensional} 的等价条件，且
$\mathcal{L}$ 是可逆 $\mathcal{O}_X$-模，则
$$
c_1(\mathcal{L}) \cap - : \CH^p(X) \to \CH^{p + 1}(X)
$$
\item 若 $X$ 是 $S$ 上局部有限型概形，满足
引理 \ref{chow-lemma-locally-equidimensional} 的等价条件，且
$\mathcal{E}$ 是有限局部自由 $\mathcal{O}_X$-模，则
$$
c_i(\mathcal{E}) \cap - : \CH^p(X) \to \CH^{p + i}(X)
$$
\end{enumerate}
警告：一个态射具有余维数 $r$ 的性质在基变换下并不保持。

\begin{remark}
\label{chow-remark-fundamental-class}
设 $(S, \delta)$ 如情形 \ref{chow-situation-setup}。
设 $X$ 在 $S$ 上局部有限型，并满足
引理 \ref{chow-lemma-locally-equidimensional} 的等价条件。
设 $X = \coprod X_n$ 是开闭子概形的分解，使 $X_n$ 的每个不可约分支
都具有 $\delta$-维数 $n$。在此情形下，有时令
$$
[X] = \sum\nolimits_n [X_n]_n \in \CH^0(X)
$$
这个类是 Chow 理论中 $X$ 的一种“基本类”。
\end{remark}




\section{分裂原理}
\label{chow-section-splitting-principle}

\noindent
在我们的设定中，不太容易精确说明分裂原理的内容。下面给出一种表述。

\begin{lemma}
\label{chow-lemma-splitting-principle}
设 $(S, \delta)$ 如情形 \ref{chow-situation-setup}。设 $X$ 在 $S$
上局部有限型，并设 $\mathcal{E}_i$ 是秩为 $r_i$ 的局部自由
$\mathcal{O}_X$-模的有限族。存在相对维数为 $d$ 的射影平坦态射
$\pi : P \to X$，使得
\begin{enumerate}
\item 对任意态射 $f : Y \to X$，映射
$\pi_Y^* : \CH_*(Y) \to \CH_{* + d}(Y \times_X P)$ 是单射；
\item 每个 $\pi^*\mathcal{E}_i$ 都有一个滤过，其逐次商
$\mathcal{L}_{i, 1}, \ldots, \mathcal{L}_{i, r_i}$ 是可逆
${\mathcal O}_P$-模。
\end{enumerate}
此外，当 (1) 成立时，限制映射 $A^*(X) \to A^*(P)$
（注 \ref{chow-remark-pullback-cohomology}）是单射。
\end{lemma}

\begin{proof}
可以设对所有 $i$ 都有 $r_i \geq 1$。对 $\sum (r_i - 1)$ 归纳。
若这个整数为 $0$，则对所有 $i$，$\mathcal{E}_i$ 都可逆，取
$\pi = \text{id}_X$ 即得结论。否则，可选取满足 $r_i > 1$ 的
$i$，并考虑态射 $\pi_i : P_i = \mathbf{P}(\mathcal{E}_i) \to X$。
有有限局部自由 $\mathcal{O}_{P_i}$-模的短正合列
$$
0 \to \mathcal{F} \to \pi_i^*\mathcal{E}_i \to \mathcal{O}_{P_i}(1) \to 0
$$
其中三项的秩分别为 $r_i - 1$、$r_i$、$1$。注意，由射影丛公式
（引理 \ref{chow-lemma-chow-ring-projective-bundle}），在任意基变换后，
$\pi_i^*$ 在 Chow 群上仍为单射。对有限局部自由
$\mathcal{O}_{P_i}$-模 $\mathcal{F}$ 以及
$\pi_{i'}^*\mathcal{E}_{i'}$（$i' \not = i$）应用归纳假设，
得到具有引理所述性质的态射 $\pi : P \to P_i$。
于是复合 $\pi_i \circ \pi : P \to X$ 满足要求。略去一些细节。
\end{proof}

\begin{remark}
\label{chow-remark-the-proof-shows-more}
引理 \ref{chow-lemma-splitting-principle} 的证明还表明态射
$\pi : P \to X$ 具有下列性质：
\begin{enumerate}
\item $\pi$ 是与有限常秩局部自由模关联的射影空间丛的有限复合；
\item 对每个 $\alpha \in \CH_k(X)$，都有
$\alpha = \pi_*(\xi_1 \cap \ldots \cap \xi_d \cap \pi^*\alpha)$，
其中 $\xi_i$ 是某个可逆 $\mathcal{O}_P$-模的第一陈类。
\end{enumerate}
第二点由第一点及引理 \ref{chow-lemma-cap-projective-bundle} 得出。
以后可按需要在此添加更多说明。
\end{remark}

\noindent
设 $(S, \delta)$、$X$ 和 $\mathcal{E}_i$ 如
引理 \ref{chow-lemma-splitting-principle}。所谓{\it 分裂原理}，是指形式地写成
$$
c(\mathcal{E}_i) = \prod (1 + x_{i, j})
$$
符号 $x_{i, 1}, \ldots, x_{i, r_i}$ 称为 $\mathcal{E}_i$ 的
{\it 陈根}。换言之，$\mathcal{E}_i$ 的第 $p$ 陈类是陈根的第 $p$
初等对称函数。分裂原理的用处在于：要证明 $\mathcal{E}_i$ 的陈类之间
的多项式关系，只须证明陈根之间对应的关系。

\medskip\noindent
具体地，设 $\pi : P \to X$ 如引理 \ref{chow-lemma-splitting-principle}。
回忆有规范的 $\mathbf{Z}$-代数映射
$\pi^* : A^*(X) \to A^*(P)$，见
注 \ref{chow-remark-pullback-cohomology}。对 $X$ 上每个 $Y$，
$\pi_Y^*$ 在 Chow 群上的单射性推出映射
$\pi^* : A^*(X) \to A^*(P)$ 是单射（略去细节）。由
引理 \ref{chow-lemma-chern-filter-by-linebundles}，有
$$
\pi^*c(\mathcal{E}_i) = \prod (1 + c_1(\mathcal{L}_{i, j}))
$$
因此，可以把陈根 $x_{i, j}$ 视为元素
$c_1(\mathcal{L}_{i, j}) \in A^*(P)$，并把所示方程理解为：
对方程左端应用单射 $\pi^* : A^*(X) \to A^*(P)$ 后，
方程在 $A^*(P)$ 中成立。

\medskip\noindent
要看清这一方法，最好给出若干例子。

\begin{lemma}
\label{chow-lemma-chern-classes-dual}
在情形 \ref{chow-situation-setup} 中，设 $X$ 在 $S$ 上局部有限型。
设 $\mathcal{E}$ 是有限局部自由 $\mathcal{O}_X$-模，其对偶为
$\mathcal{E}^\vee$。则在 $A^i(X)$ 中有
$$
c_i(\mathcal{E}^\vee) = (-1)^i c_i(\mathcal{E})
$$
\end{lemma}

\begin{proof}
选取引理 \ref{chow-lemma-splitting-principle} 中的态射
$\pi : P \to X$。由于 $\pi^*$ 在任意基变换后都是单射，只须在拉回到
$P$ 后证明 $\mathcal{E}$ 与 $\mathcal{E}^\vee$ 的陈类之间的关系。
因此可以设存在可逆 $\mathcal{O}_X$-模
${\mathcal L}_i$（$i = 1, \ldots, r$）以及滤过
$$
0 = \mathcal{E}_0 \subset \mathcal{E}_1 \subset \mathcal{E}_2
\subset \ldots \subset \mathcal{E}_r = \mathcal{E}
$$
使得 $\mathcal{E}_i/\mathcal{E}_{i - 1} \cong \mathcal{L}_i$。
于是得到对偶滤过
$$
0 = \mathcal{E}_r^\perp \subset \mathcal{E}_1^\perp \subset \mathcal{E}_2^\perp
\subset \ldots \subset \mathcal{E}_0^\perp = \mathcal{E}^\vee
$$
并且 $\mathcal{E}_{i - 1}^\perp/\mathcal{E}_i^\perp \cong
\mathcal{L}_i^{\otimes -1}$。令 $x_i = c_1(\mathcal{L}_i)$。
由引理 \ref{chow-lemma-c1-cap-additive}，有
$c_1(\mathcal{L}_i^{\otimes -1}) = - x_i$。由
引理 \ref{chow-lemma-chern-filter-by-linebundles}，在 $A^*(X)$ 中有
$$
c(\mathcal{E}) = \prod\nolimits_{i = 1}^r (1 + x_i)
\quad\text{且}\quad
c(\mathcal{E}^\vee) = \prod\nolimits_{i = 1}^r (1 - x_i)
$$
结论由一个形式计算得出，这里略去。
\end{proof}

\begin{lemma}
\label{chow-lemma-chern-classes-tensor-product}
在情形 \ref{chow-situation-setup} 中，设 $X$ 在 $S$ 上局部有限型。
设 $\mathcal{E}$ 和 $\mathcal{F}$ 是秩分别为 $r$ 和 $s$ 的有限局部自由
$\mathcal{O}_X$-模。则在 $A^*(X)$ 中有
$$
c_1(\mathcal{E} \otimes \mathcal{F})
=
r c_1(\mathcal{F}) + s c_1(\mathcal{E})
$$
$$
c_2(\mathcal{E} \otimes \mathcal{F})
=
r c_2(\mathcal{F}) + s c_2(\mathcal{E}) +
{r \choose 2} c_1(\mathcal{F})^2 +
(rs - 1) c_1(\mathcal{F})c_1(\mathcal{E}) +
{s \choose 2} c_1(\mathcal{E})^2
$$
依此类推。
\end{lemma}

\begin{proof}
完全仿照引理 \ref{chow-lemma-chern-classes-dual} 的证明，可以设有可逆
$\mathcal{O}_X$-模 ${\mathcal L}_i$（$i = 1, \ldots, r$）、
${\mathcal N}_i$（$i = 1, \ldots, s$）以及滤过
$$
0 = \mathcal{E}_0 \subset \mathcal{E}_1 \subset \mathcal{E}_2
\subset \ldots \subset \mathcal{E}_r = \mathcal{E}
\quad\text{且}\quad
0 = \mathcal{F}_0 \subset \mathcal{F}_1 \subset \mathcal{F}_2
\subset \ldots \subset \mathcal{F}_s = \mathcal{F}
$$
使得 $\mathcal{E}_i/\mathcal{E}_{i - 1} \cong \mathcal{L}_i$，且
$\mathcal{F}_j/\mathcal{F}_{j - 1} \cong \mathcal{N}_j$。
按字典序排列二元组 $(i, j)$，得到滤过
$$
0 \subset \ldots \subset
\mathcal{E}_i \otimes \mathcal{F}_j
+
\mathcal{E}_{i - 1} \otimes \mathcal{F}
\subset \ldots \subset \mathcal{E} \otimes \mathcal{F}
$$
其逐次商为
$$
\mathcal{L}_1 \otimes \mathcal{N}_1,
\mathcal{L}_1 \otimes \mathcal{N}_2,
\ldots,
\mathcal{L}_1 \otimes \mathcal{N}_s,
\mathcal{L}_2 \otimes \mathcal{N}_1,
\ldots,
\mathcal{L}_r \otimes \mathcal{N}_s
$$
由引理 \ref{chow-lemma-chern-filter-by-linebundles}，在 $A^*(X)$ 中有
$$
c(\mathcal{E}) = \prod (1 + x_i),
\quad
c(\mathcal{F}) = \prod (1 + y_j),
\quad\text{且}\quad
c(\mathcal{E} \otimes \mathcal{F}) = \prod (1 + x_i + y_j),
$$
结论由一个形式计算得出，这里略去。
\end{proof}

\begin{remark}
\label{chow-remark-equalities-nonconstant-rank}
即使考虑秩可以不是常数的有限局部自由 $\mathcal{O}_X$-模，
上面证明的等式仍然成立。事实上，可以如下使用秩和陈类的多项式。
考虑分次多项式环 $\mathbf{Z}[r, c_1, c_2, c_3, \ldots]$，
其中 $r$ 的次数为 $0$，$c_i$ 的次数为 $i$。设
$$
P \in \mathbf{Z}[r, c_1, c_2, c_3, \ldots]
$$
是次数为 $p$ 的齐次多项式。则对 $X$ 上任意有限局部自由
$\mathcal{O}_X$-模 $\mathcal{E}$，可以考虑
$$
P(\mathcal{E}) =
P(r(\mathcal{E}), c_1(\mathcal{E}), c_2(\mathcal{E}), c_3(\mathcal{E}), \ldots)
\in A^p(X)
$$
记号与约定见注 \ref{chow-remark-extend-to-finite-locally-free}。
要证明这些多项式之间的关系（涉及多个有限局部自由模时亦然），
可以在 $X$ 上局部工作，并使用上述分裂原理。例如，断言
$$
c_2(\SheafHom_{\mathcal{O}_X}(\mathcal{E}, \mathcal{E})) =
P(\mathcal{E})
$$
其中 $P = 2rc_2 - (r - 1)c_1^2$。事实上，由于
$\SheafHom_{\mathcal{O}_X}(\mathcal{E}, \mathcal{E}) =
\mathcal{E} \otimes \mathcal{E}^\vee$，该式容易从上述引理 \ref{chow-lemma-chern-classes-dual} 和
\ref{chow-lemma-chern-classes-tensor-product} 得出：只须把 $X$ 如
注 \ref{chow-remark-extend-to-finite-locally-free} 分解成
$\mathcal{E}$ 的秩为常数的部分。
\end{remark}

\begin{example}
\label{chow-example-power-sum}
对每个 $p \geq 1$，存在唯一的次数为 $p$ 的齐次多项式
$P_p \in \mathbf{Z}[c_1, c_2, c_3, \ldots]$，使得对任意
$n \geq p$，在 $\mathbf{Z}[x_1, \ldots, x_n]$ 中有
$$
P_p(s_1, s_2, \ldots, s_p) = \sum x_i^p
$$
其中 $s_1, \ldots, s_p$ 是 $x_1, \ldots, x_n$ 的初等对称多项式，即
$$
s_i = \sum\nolimits_{1 \leq j_1 < \ldots < j_i \leq n}
x_{j_1}x_{j_2} \ldots x_{j_i}
$$
$P_p$ 的存在性来自熟知的事实：初等对称函数生成整数上的全部对称函数环。
刻画 $P_p \in \mathbf{Z}[c_1, c_2, c_3, \ldots]$ 的另一种方式是，
作为形式幂级数有
$$
\log(1 + c_1 + c_2 + c_3 + \ldots) =
\sum\nolimits_{p \geq 1} (-1)^{p - 1}\frac{P_p}{p}
$$
这可由写成 $1 + c_1 + c_2 + \ldots = \prod (1 + x_i)$ 并应用
对数函数的幂级数立即看出。展开左端，得到
\begin{align*}
& (c_1 + c_2 + \ldots) - (1/2)(c_1 + c_2 + \ldots)^2 +
(1/3)(c_1 + c_2 + \ldots)^3 - \ldots \\
& =
c_1 + (c_2 - (1/2)c_1^2) + (c_3 - c_1c_2 + (1/3)c_1^3) + \ldots
\end{align*}
由此得到
\begin{align*}
P_1 & = c_1, \\
P_2 & = c_1^2 - 2c_2, \\
P_3 & = c_1^3 - 3c_1c_2 + 3c_3, \\
P_4 & = c_1^4 - 4c_1^2c_2 + 4c_1c_3 + 2c_2^2 - 4c_4,
\end{align*}
依此类推。由于有限局部自由 $\mathcal{O}_X$-模 $\mathcal{E}$ 的陈类
是陈根 $x_i$ 的初等对称多项式，可见
$$
P_p(\mathcal{E}) = \sum x_i^p
$$
为方便起见，在 $\mathbf{Z}[r, c_1, c_2, c_3, \ldots]$ 中令
$P_0 = r$，使得作为双变类有 $P_0(\mathcal{E}) = r(\mathcal{E})$
（如注 \ref{chow-remark-extend-to-finite-locally-free} 和
\ref{chow-remark-equalities-nonconstant-rank}）。
\end{example}






\section{陈类与截面}
\label{chow-section-top-chern-class}

\noindent
本节很短，主要结果是：可以用一个正则截面的零点轨迹计算有限局部自由模的顶陈类。

\medskip\noindent
设 $(S, \delta)$ 如情形 \ref{chow-situation-setup}。设 $X$ 是 $S$
上局部有限型概形，$\mathcal{E}$ 是有限局部自由 $\mathcal{O}_X$-模，
且 $f : X' \to X$ 局部有限型。设
$$
s \in \Gamma(X', f^*\mathcal{E})
$$
是 $\mathcal{E}$ 到 $X'$ 的拉回的整体截面。设 $Z(s) \subset X'$
为 $s$ 的零点概形。更精确地，定义 $Z(s)$ 为如下闭子概形：
其拟凝聚理想层是映射
$s : f^*\mathcal{E}^\vee \to \mathcal{O}_{X'}$ 的像。

\begin{lemma}
\label{chow-lemma-top-chern-class}
在上述情形中，假设 $\dim_\delta(X') = n$，$f^*\mathcal{E}$ 的常秩为
$r$，$\dim_\delta(Z(s)) \leq n - r$，并且对每个满足
$\delta(\xi) = n - r$ 的一般点 $\xi \in Z(s)$，$Z(s)$ 在
$\mathcal{O}_{X', \xi}$ 中的理想由长度为 $r$ 的正则序列生成。则在
$\CH_*(X')$ 中有
$$
c_r(\mathcal{E}) \cap [X']_n = [Z(s)]_{n - r}
$$
\end{lemma}

\begin{proof}
由于 $c_r(\mathcal{E})$ 是双变类
（引理 \ref{chow-lemma-cap-cp-bivariant}），可以设 $X = X'$，
并须证明 $\CH_{n - r}(X)$ 中的等式
$c_r(\mathcal{E}) \cap [X]_n = [Z(s)]_{n - r}$。
对 $r \geq 0$ 归纳。（$r = 0$ 的情形显然。）当 $r = 1$ 时，
结论由引理 \ref{chow-lemma-geometric-cap} 给出。设 $r > 1$。

\medskip\noindent
设 $\pi : P \to X$ 是与 $\mathcal{E}$ 关联的射影空间丛，
并考虑短正合列
$$
0 \to \mathcal{E}' \to \pi^*\mathcal{E} \to \mathcal{O}_P(1) \to 0
$$
由射影空间丛公式（引理 \ref{chow-lemma-chow-ring-projective-bundle}），
只须在沿 $\pi$ 拉回后证明该等式。注意，
$\pi^{-1}Z(s) = Z(\pi^*s)$ 的 $\delta$-维数 $\leq n - 1$，
且在 $\delta$-维数为 $n - 1$ 的一般点处关于正则序列的假设，
由平坦拉回仍成立；见《交换代数》引理 \ref{algebra-lemma-flat-increases-depth}。设
$t \in \Gamma(P, \mathcal{O}_P(1))$ 是 $\pi^*s$ 的像。断言
$$
[Z(t)]_{n + r - 2} = c_1(\mathcal{O}_P(1)) \cap [P]_{n + r - 1}
$$
假设该断言成立，则如下完成证明。限制 $\pi^*s|_{Z(t)}$ 在
$\mathcal{O}_P(1)|_{Z(t)}$ 中映到零，故来自唯一元素
$s' \in \Gamma(Z(t), \mathcal{E}'|_{Z(t)})$。注意，作为 $P$
的闭子概形，有 $Z(s') = Z(\pi^*s)$。若
$\xi \in Z(s')$ 是满足 $\delta(\xi) = n - 1$ 的一般点，则
$Z(s')$ 在 $\mathcal{O}_{Z(t), \xi}$ 中的理想可由长度为
$r - 1$ 的正则序列生成：它由 $\mathcal{O}_{P, \xi}$ 中
$r - 1$ 个元素的像生成，而这 $r - 1$ 个元素连同
$Z(t)$ 在 $\mathcal{O}_{P, \xi}$ 中理想的一个生成元，构成
$\mathcal{O}_{P, \xi}$ 中长度为 $r$ 的正则序列。因此可对
$Z(t)$ 上的 $s'$ 应用归纳假设，得到
$c_{r - 1}(\mathcal{E}') \cap [Z(t)]_{n + r - 2} = [Z(s')]_{n - 1}$。
综合以上各式，得到
\begin{align*}
c_r(\pi^*\mathcal{E}) \cap [P]_{n + r - 1}
& =
c_{r - 1}(\mathcal{E}') \cap c_1(\mathcal{O}_P(1)) \cap [P]_{n + r - 1} \\
& =
c_{r - 1}(\mathcal{E}') \cap [Z(t)]_{n + r - 2} \\
& =
[Z(s')]_{n - 1} \\
& = [Z(\pi^*s)]_{n - 1}
\end{align*}
这正是所需结论。

\medskip\noindent
证明断言。它来自对已经使用过的
引理 \ref{chow-lemma-geometric-cap} 的一次应用。我们有
$\pi^{-1}(Z(s)) = Z(\pi^*s) \subset Z(t)$。另一方面，对
$x \in X$，若 $P_x \subset Z(t)$，则 $t|_{P_x} = 0$；这推出
$s$ 在纤维 $\mathcal{E} \otimes \kappa(x)$ 中为零，进而推出
$x \in Z(s)$。因此 $\dim_\delta(Z(t)) \leq n + (r - 1) - 1$。
最后，设 $\xi \in Z(t)$ 是满足
$\delta(\xi) = n + r - 2$ 的一般点。若 $\xi$ 不是 $P \to X$
某条纤维的一般点，则 $Z(t)$ 的一个局部方程在
$\mathcal{O}_{P, \xi}$ 中是非零因子，这是立即的（因为由
《交换代数》引理 \ref{algebra-lemma-grothendieck}，可在纤维上检验）。
若 $\xi$ 是某条纤维的一般点，则 $x = \pi(\xi) \in Z(s)$，且
$\delta(x) = n + r - 2 - (r - 1) = n - 1$。由于 $r > 1$，
这与 $\dim_\delta(Z(s)) \leq n - r$ 矛盾，所以这种情形不会发生。
\end{proof}

\begin{lemma}
\label{chow-lemma-easy-virtual-class}
设 $(S, \delta)$ 如情形 \ref{chow-situation-setup}。设 $X$ 是 $S$
上局部有限型概形，并设
$$
0 \to \mathcal{N}' \to \mathcal{N} \to \mathcal{E} \to 0
$$
是有限局部自由 $\mathcal{O}_X$-模的短正合列。考虑闭嵌入
$$
i :
N' = \underline{\Spec}_X(\text{Sym}((\mathcal{N}')^\vee))
\longrightarrow
N = \underline{\Spec}_X(\text{Sym}(\mathcal{N}^\vee))
$$
对 $\alpha \in \CH_k(X)$，有
$$
i_*(p')^*\alpha = p^*(c_{top}(\mathcal{E}) \cap \alpha)
$$
其中 $p' : N' \to X$ 和 $p : N \to X$ 是结构态射。
\end{lemma}

\begin{proof}
这里 $c_{top}(\mathcal{E})$ 是注 \ref{chow-remark-top-chern-class}
中定义的双变类。由其定义，为验证该公式，可以设 $\mathcal{E}$
为常秩。类似地，可设 $\mathcal{N}'$ 和 $\mathcal{N}$ 的常秩分别为
$r'$ 和 $r$，于是 $\mathcal{E}$ 的秩为 $r - r'$，且
$c_{top}(\mathcal{E}) = c_{r - r'}(\mathcal{E})$。注意，
$p^*\mathcal{E}$ 有规范截面
$$
s \in \Gamma(N, p^*\mathcal{E}) = \Gamma(X, p_*p^*\mathcal{E}) =
\Gamma(X, \mathcal{E} \otimes_{\mathcal{O}_X} \text{Sym}(\mathcal{N}^\vee)
\supset \Gamma(X, \SheafHom(\mathcal{N}, \mathcal{E}))
$$
它对应于引理陈述中的满射 $\mathcal{N} \to \mathcal{E}$。
该截面的零点概形恰为 $N' \subset N$。设 $Y \subset X$ 是
$\delta$-维数为 $n$ 的整闭子概形。于是
\begin{enumerate}
\item $p^*[Y] = [p^{-1}(Y)]$，因为 $p^{-1}(Y)$ 是
$\delta$-维数为 $n + r$ 的整概形；
\item $(p')^*[Y] = [(p')^{-1}(Y)]$，因为 $(p')^{-1}(Y)$ 是
$\delta$-维数为 $n + r'$ 的整概形；
\item $s$ 在 $p^{-1}Y$ 上的限制以 $(p')^{-1}Y$ 为零点概形，
且闭浸入 $(p')^{-1}Y \to p^{-1}Y$ 是正则浸入（局部由正则序列截出）。
\end{enumerate}
由引理 \ref{chow-lemma-top-chern-class}，得到
$$
(p')^*[Y] = c_{r - r'}(p^*\mathcal{E}) \cap p^*[Y]
\quad\text{于}\quad \CH_*(N)
$$
引理得证。
\end{proof}










\section{陈特征与张量积}
\label{chow-section-chern-classes-tensor}

\noindent
设 $(S, \delta)$ 如情形 \ref{chow-situation-setup}，设 $X$ 在 $S$
上局部有限型。若 $x_i$ 是 ${\mathcal E}$ 的陈根，则把有限局部自由
$\mathcal{O}_X$-模的{\it 陈特征}定义为形式表达式
$$
ch({\mathcal E}) = \sum\nolimits_{i=1}^r e^{x_i}
$$
把它写成陈类的多项式，得到
\begin{align*}
ch(\mathcal{E})
& =
r(\mathcal{E})
+
c_1(\mathcal{E}) +
\frac{1}{2}(c_1(\mathcal{E})^2 - 2c_2(\mathcal{E}))
+
\frac{1}{6}(c_1(\mathcal{E})^3 - 3c_1(\mathcal{E})c_2(\mathcal{E}) + 3c_3(\mathcal{E})) \\
& \quad\quad +
\frac{1}{24}(c_1(\mathcal{E})^4 - 4c_1(\mathcal{E})^2c_2(\mathcal{E}) + 4c_1(\mathcal{E})c_3(\mathcal{E}) + 2c_2(\mathcal{E})^2 - 4c_4(\mathcal{E}))
+
\ldots \\
& =
\sum\nolimits_{p = 0, 1, 2, \ldots} \frac{P_p(\mathcal{E})}{p!}
\end{align*}
其中 $P_p$ 是例 \ref{chow-example-power-sum} 中陈类的多项式。
上述表达式的次数 $p$ 分量是
$$
ch_p(\mathcal{E}) = \frac{P_p(\mathcal{E})}{p!} \in A^p(X) \otimes \mathbf{Q}
$$
系数是有理数是什么意思？这只是说，把
$ch_p(\mathcal{E})$ 视为 $A^p(X) \otimes \mathbf{Q}$ 的元素。

\begin{remark}
\label{chow-remark-extend-chern-character-to-finite-locally-free}
在上述讨论中，即使 $\mathcal{E}$ 的秩不是常数，也已经定义了
$\mathcal{E}$ 的陈特征分量
$ch_p(\mathcal{E}) \in A^p(X) \otimes \mathbf{Q}$。
见注 \ref{chow-remark-extend-to-finite-locally-free} 和
\ref{chow-remark-equalities-nonconstant-rank}。因此 $\mathcal{E}$ 的完整陈特征
是 $\prod_{p \geq 0} (A^p(X) \otimes \mathbf{Q})$ 的元素。若 $X$
拟紧且 $\dim(X) < \infty$（通常维数），则利用
引理 \ref{chow-lemma-vanish-above-dimension} 和分裂原理可证明
$ch(\mathcal{E}) \in A^*(X) \otimes \mathbf{Q}$。
\end{remark}

\begin{lemma}
\label{chow-lemma-chern-character-additive}
设 $(S, \delta)$ 如情形 \ref{chow-situation-setup}。设 $X$ 在 $S$
上局部有限型，并设
$
0 \to \mathcal{E}_1 \to \mathcal{E} \to \mathcal{E}_2 \to 0
$
是有限局部自由 $\mathcal{O}_X$-模的短正合列。则有等式
$$
ch(\mathcal{E}) = ch(\mathcal{E}_1) + ch(\mathcal{E}_2)
$$
更精确地，在 $A^p(X)$ 中有
$P_p(\mathcal{E}) = P_p(\mathcal{E}_1) + P_p(\mathcal{E}_2)$，
其中 $P_p$ 如例 \ref{chow-example-power-sum}。
\end{lemma}

\begin{proof}
只须证明更精确的陈述。由第 \ref{chow-section-splitting-principle} 节，
这是因为：若 $x_{1, i}$（$i = 1, \ldots, r_1$）与
$x_{2, i}$（$i = 1, \ldots, r_2$）分别是
$\mathcal{E}_1$ 与 $\mathcal{E}_2$ 的陈根，则
$x_{1, 1}, \ldots, x_{1, r_1}, x_{2, 1}, \ldots, x_{2, r_2}$
是 $\mathcal{E}$ 的陈根。因此，结论来自例 \ref{chow-example-power-sum} 中对 $P_p$ 的选取。
\end{proof}

\begin{lemma}
\label{chow-lemma-chern-character-multiplicative}
设 $(S, \delta)$ 如情形 \ref{chow-situation-setup}。设 $X$ 在 $S$
上局部有限型，并设 $\mathcal{E}_1$ 和 $\mathcal{E}_2$ 是有限局部自由
$\mathcal{O}_X$-模。则有等式
$$
ch(\mathcal{E}_1 \otimes_{\mathcal{O}_X} \mathcal{E}_2) =
ch(\mathcal{E}_1) ch(\mathcal{E}_2)
$$
更精确地，在 $A^p(X)$ 中有
$$
P_p(\mathcal{E}_1 \otimes_{\mathcal{O}_X} \mathcal{E}_2) =
\sum\nolimits_{p_1 + p_2 = p}
{p \choose p_1} P_{p_1}(\mathcal{E}_1) P_{p_2}(\mathcal{E}_2)
$$
其中 $P_p$ 如例 \ref{chow-example-power-sum}。
\end{lemma}

\begin{proof}
只须证明更精确的陈述。由第 \ref{chow-section-splitting-principle} 节，
这是因为：若 $x_{1, i}$（$i = 1, \ldots, r_1$）与
$x_{2, i}$（$i = 1, \ldots, r_2$）分别是
$\mathcal{E}_1$ 与 $\mathcal{E}_2$ 的陈根，则
$x_{1, i} + x_{2, j}$（$1 \leq i \leq r_1$，$1 \leq j \leq r_2$）
是 $\mathcal{E}_1 \otimes \mathcal{E}_2$ 的陈根。因此，结论来自
$(x_{1, i} + x_{2, j})^p$ 的二项式公式以及例 \ref{chow-example-power-sum} 中多项式 $P_p$ 的形式。
\end{proof}

\begin{lemma}
\label{chow-lemma-chern-character-dual}
在情形 \ref{chow-situation-setup} 中，设 $X$ 在 $S$ 上局部有限型。
设 $\mathcal{E}$ 是有限局部自由 $\mathcal{O}_X$-模，其对偶为
$\mathcal{E}^\vee$。则在 $A^i(X) \otimes \mathbf{Q}$ 中有
$ch_i(\mathcal{E}^\vee) = (-1)^i ch_i(\mathcal{E})$。
\end{lemma}

\begin{proof}
结论来自陈类的相应结果
（引理 \ref{chow-lemma-chern-classes-dual}）。
\end{proof}











\section{陈类与导出范畴}
\label{chow-section-pre-derived}

\noindent
本节在如下假设下定义概形 $X$ 的导出范畴中完美对象 $E$ 的全陈类：
$E$ 可在 $X$ 的一个包络上由有限局部自由模的有限复形表示。

\medskip\noindent
设 $(S, \delta)$ 如情形 \ref{chow-situation-setup}，设 $X$ 在 $S$
上局部有限型，并设
$$
\mathcal{E}^a \to \mathcal{E}^{a + 1} \to \ldots \to \mathcal{E}^b
$$
是常秩有限局部自由 $\mathcal{O}_X$-模的有界复形。
用公式
$$
c(\mathcal{E}^\bullet) = \prod\nolimits_{n = a, \ldots, b}
c(\mathcal{E}^n)^{(-1)^n} \in \prod\nolimits_{p \geq 0} A^p(X)
$$
定义{\it 复形的全陈类}。这里的逆是形式逆，即
$$
(1 + c_1 + c_2 + c_3 + \ldots)^{-1} =
1 - c_1 + c_1^2 - c_2 - c_1^3 + 2c_1 c_2 - c_3 + \ldots
$$
用 $c_p(\mathcal{E}^\bullet) \in A^p(X)$ 表示
$c(\mathcal{E}^\bullet)$ 的次数 $p$ 部分。类似地，用公式
$$
ch(\mathcal{E}^\bullet) = \sum\nolimits_{n = a, \ldots, b}
(-1)^n ch(\mathcal{E}^n) \in
\prod\nolimits_{p \geq 0} (A^p(X) \otimes \mathbf{Q})
$$
定义{\it 复形的陈特征}。用
$ch_p(\mathcal{E}^\bullet) \in A^p(X) \otimes \mathbf{Q}$ 表示
$ch(\mathcal{E}^\bullet)$ 的次数 $p$ 部分。最后，对例 \ref{chow-example-power-sum} 中的
$P_p \in \mathbf{Z}[r, c_1, c_2, c_3, \ldots]$，定义
$$
P_p(\mathcal{E}^\bullet) = \sum\nolimits_{n = a, \ldots, b}
(-1)^n P_p(\mathcal{E}^n)
$$
它属于 $A^p(X)$。如常，有
$ch_p(\mathcal{E}^\bullet) = (1/p!)P_p(\mathcal{E}^\bullet)$。
下一个引理表明，这些构造只依赖于复形在导出范畴中的像。

\begin{lemma}
\label{chow-lemma-pre-derived-chern-class}
设 $(S, \delta)$ 如情形 \ref{chow-situation-setup}。
设 $X$ 在 $S$ 上局部有限型。设 $E \in D(\mathcal{O}_X)$ 是一个对象，
并存在表示 $E$ 的有限局部自由 $\mathcal{O}_X$-模的局部有界复形
$\mathcal{E}^\bullet$。则上述构造的轻微推广
$$
c(\mathcal{E}^\bullet) \in \prod\nolimits_{p \geq 0} A^p(X),\quad
ch(\mathcal{E}^\bullet) \in
\prod\nolimits_{p \geq 0} A^p(X) \otimes \mathbf{Q},\quad
P_p(\mathcal{E}^\bullet) \in A^p(X)
$$
与复形 $\mathcal{E}^\bullet$ 的选择无关。
\end{lemma}

\begin{proof}
先对全陈类证明；另两种情形用同样的论证，只须以
引理 \ref{chow-lemma-chern-character-additive} 代替
引理 \ref{chow-lemma-additivity-chern-classes}。

\medskip\noindent
如注 \ref{chow-remark-extend-to-finite-locally-free}，为定义全陈类
$c(\mathcal{E}^\bullet)$，把 $X$ 分解为开闭子概形
$$
X = \coprod\nolimits_{i \in I} X_i
$$
使对所有 $n$、$i$，$\mathcal{E}^n$ 在 $X_i$ 上的秩为常数。
（这些秩是 $X$ 上的局部常值函数，所以可以这样做。）由于
$\mathcal{E}^\bullet$ 局部有界，对固定的 $i$，层
$\mathcal{E}^n|_{X_i}$ 中仅有限多个非零。因此可以如上定义
$$
c(\mathcal{E}^\bullet|_{X_i}) =
\prod\nolimits_n c(\mathcal{E}^n|_{X_i})^{(-1)^n}
\in \prod\nolimits_{p \geq 0} A^p(X_i)
$$
由引理 \ref{chow-lemma-disjoint-decomposition-bivariant}，有
$A^p(X) = \prod_i A^p(X_i)$。因此，对每个 $p \in \mathbf{Z}$，
存在唯一元素 $c_p(\mathcal{E}^\bullet) \in A^p(X)$，使得对所有
$i$，其在 $X_i$ 上的限制为 $c_p(\mathcal{E}^\bullet|_{X_i})$。

\medskip\noindent
设有限局部自由 $\mathcal{O}_X$-模的另一个局部有界复形
$\mathcal{F}^\bullet$ 也表示 $E$。设 $g : Y \to X$ 局部有限型且
$Y$ 整。由引理 \ref{chow-lemma-bivariant-zero}，只须证明
$c(g^*\mathcal{E}^\bullet) \cap [Y]$ 与
$c(g^*\mathcal{F}^\bullet) \cap [Y]$ 相同；甚至只须在把 $Y$
替换为 $Y$ 上整、正常且双有理的概形后证明。首先可得
$g^*\mathcal{E}^\bullet$ 和 $g^*\mathcal{F}^\bullet$ 是常秩有限局部自由
$\mathcal{O}_Y$-模的有界复形。接着，由《平坦性进阶》引理 \ref{flat-lemma-blowup-complex-integral}，可设对所有
$i \in \mathbf{Z}$，$H^i(Lg^*E)$ 都是 Tor 维数 $\leq 1$ 的完美对象。
这就归约到下一段的情形。

\medskip\noindent
假设 $X$ 整，$\mathcal{E}^\bullet$ 和 $\mathcal{F}^\bullet$ 是常秩有限局部自由模的
有界复形，且对所有 $i \in \mathbf{Z}$，$H^i(E)$ 都是 Tor 维数
$\leq 1$ 的完美 $\mathcal{O}_X$-模。须证明
$c(\mathcal{E}^\bullet) \cap [X]$ 与
$c(\mathcal{F}^\bullet) \cap [X]$ 相同。记
$d_\mathcal{E}^i : \mathcal{E}^i \to \mathcal{E}^{i + 1}$ 和
$d_\mathcal{F}^i : \mathcal{F}^i \to \mathcal{F}^{i + 1}$
为这些复形的微分。由《平坦性进阶》注 \ref{flat-remark-when-you-have-a-complex}，对所有 $i$，
$\Im(d_\mathcal{E}^i)$、$\Ker(d_\mathcal{E}^i)$、
$\Im(d_\mathcal{F}^i)$ 和 $\Ker(d_\mathcal{F}^i)$ 都是有限局部自由
$\mathcal{O}_X$-模。由可加性（引理 \ref{chow-lemma-additivity-chern-classes}），
$$
c(\mathcal{E}^\bullet) = \prod\nolimits_i
c(\Ker(d_\mathcal{E}^i))^{(-1)^i} c(\Im(d_\mathcal{E}^i))^{(-1)^i}
$$
$\mathcal{F}^\bullet$ 也类似。由于有短正合列
$$
0 \to \Im(d_\mathcal{E}^i) \to \Ker(d_\mathcal{E}^i) \to H^i(E) \to 0
\quad\text{且}\quad
0 \to \Im(d_\mathcal{F}^i) \to \Ker(d_\mathcal{F}^i) \to H^i(E) \to 0
$$
便归约到下一段陈述并解决的问题。

\medskip\noindent
假设 $X$ 整，且有两个短正合列
$$
0 \to \mathcal{E}' \to \mathcal{E} \to \mathcal{Q} \to 0
\quad\text{且}\quad
0 \to \mathcal{F}' \to \mathcal{F} \to \mathcal{Q} \to 0
$$
其中 $\mathcal{E}$、$\mathcal{E}'$、$\mathcal{F}$、$\mathcal{F}'$
有限局部自由。问题是证明
$c(\mathcal{E})c(\mathcal{E}')^{-1} \cap [X] =
c(\mathcal{F})c(\mathcal{F}')^{-1} \cap [X]$。为此，考虑定义
$\mathcal{G}$ 的短正合列
$$
0 \to \mathcal{G} \to \mathcal{E} \oplus \mathcal{F} \to \mathcal{Q} \to 0
$$
由于 $\mathcal{Q}$ 的 Tor 维数 $\leq 1$，$\mathcal{G}$ 有限局部自由。
追图表明，满射 $\mathcal{G} \to \mathcal{F}$ 的核同构地映到
$\mathcal{E}$ 中的 $\mathcal{E}'$，而满射
$\mathcal{G} \to \mathcal{E}$ 的核同构地映到 $\mathcal{F}$ 中的
$\mathcal{F}'$。（在仿射局部，这来自 Schanuel 引理或等价于该引理，
见《交换代数》引理 \ref{algebra-lemma-Schanuel}。）因此
$$
c(\mathcal{E})c(\mathcal{F}') = c(\mathcal{G}) =
c(\mathcal{F})c(\mathcal{E}')
$$
正是所需结论。
\end{proof}

\begin{lemma}
\label{chow-lemma-defined-by-envelope}
设 $(S, \delta)$ 如情形 \ref{chow-situation-setup}。
设 $X$ 在 $S$ 上局部有限型，并设 $E \in D(\mathcal{O}_X)$
是完美对象。假设存在一个包络
$f : Y \to X$（定义 \ref{chow-definition-envelope}），
使得 $Lf^*E$ 在 $D(\mathcal{O}_Y)$ 中同构于有限局部自由
$\mathcal{O}_Y$-模组成的局部有界复形 $\mathcal{E}^\bullet$。
则存在唯一的双变类
$c(E) \in \prod_{p \geq 0} A^p(X)$、
$ch(E) \in \prod_{p \geq 0} A^p(X) \otimes \mathbf{Q}$ 以及
$P_p(E) \in A^p(X)$；它们与 $f : Y \to X$
及 $\mathcal{E}^\bullet$ 的选择无关，并且这些类在
$Y$ 上的限制分别等于 $c(\mathcal{E}^\bullet)$、
$ch(\mathcal{E}^\bullet)$ 和 $P_p(\mathcal{E}^\bullet)$。
\end{lemma}

\begin{proof}
固定 $p \in \mathbf{Z}$。我们对陈类
$c_p(E) \in A^p(X)$ 证明此引理；其余情形的论证从略。

\medskip\noindent
设 $g : T \to X$ 是局部有限型态射，并且存在有限局部自由
$\mathcal{O}_T$-模组成的局部有界复形 $\mathcal{E}^\bullet$，
它在 $D(\mathcal{O}_T)$ 中表示 $Lg^*E$。由
引理 \ref{chow-lemma-pre-derived-chern-class}，双变类
$c_p(\mathcal{E}^\bullet) \in A^p(T)$ 与
$\mathcal{E}^\bullet$ 的选择无关。将此类记为
$c_p(Lg^*E) \in A^p(T)$。对任意进一步的局部有限型态射
$h : T' \to T$，令 $g' = g \circ h$，则
$L(g')^*E = L(g \circ h)^*E = Lh^*Lg^*E$ 在
$D(\mathcal{O}_{T'})$ 中由 $h^*\mathcal{E}^\bullet$ 表示。
因此 $c_p(L(g')^*E)$ 有定义，并且等于
$c_p(Lg^*E)$ 在 $T'$ 上的限制
（注 \ref{chow-remark-restriction-bivariant}；严格来说，这里需要应用
引理 \ref{chow-lemma-cap-cp-bivariant}）。

\medskip\noindent
取引理陈述中的 $f : Y \to X$ 和 $\mathcal{E}^\bullet$。
把引理 \ref{chow-lemma-envelope-bivariant} 应用于
$f : Y \to X$ 以及上一段构造的类
$c' = c_p(Lf^*E)$，便得到双变类 $c_p(E) \in A^p(X)$。
该引理的假设成立，因为上一段的讨论给出
$res_1(c') = c_p(Lg^*E) = res_2(c')$，其中
$g = f \circ p = f \circ q : Y \times_X Y \to X$。

\medskip\noindent
最后，假设 $f' : Y' \to X$ 是另一个包络，并且
$L(f')^*E$ 由有限局部自由 $\mathcal{O}_{Y'}$-模组成的有界复形表示。
于是 $c_p(Lf^*E)$ 与 $c_p(L(f')^*E)$ 在
$Y \times_X Y'$ 上的限制相等。由于
$Y \times_X Y' \to X$ 是包络（引理 \ref{chow-lemma-base-change-envelope} 和
\ref{chow-lemma-composition-envelope}），由
引理 \ref{chow-lemma-envelope-bivariant} 中的唯一性可知，
$c_p(E)$ 的这两个候选者相同。
\end{proof}

\begin{definition}
\label{chow-definition-defined-on-perfect}
设 $(S, \delta)$ 如情形 \ref{chow-situation-setup}。
设 $X$ 在 $S$ 上局部有限型，并设 $E \in D(\mathcal{O}_X)$
是完美对象。
\begin{enumerate}
\item 若存在包络 $f : Y \to X$，使得 $Lf^*E$ 在
$D(\mathcal{O}_Y)$ 中同构于有限局部自由
$\mathcal{O}_Y$-模组成的局部有界复形，则称
{\it $E$ 的陈类已定义}\footnote{若干判据见
引理 \ref{chow-lemma-chern-classes-defined}。}。
\item 若 $E$ 的陈类已定义，则应用
引理 \ref{chow-lemma-defined-by-envelope} 定义
$$
c(E) \in \prod\nolimits_{p \geq 0} A^p(X),\quad
ch(E) \in
\prod\nolimits_{p \geq 0} A^p(X) \otimes \mathbf{Q},\quad
P_p(E) \in A^p(X)
$$
\end{enumerate}
\end{definition}

\noindent
此定义适用于本章所设想的许多情形，但并非全部情形，见
引理 \ref{chow-lemma-chern-classes-defined}。对于定义中一般的
$E/X/(S,\delta)$，也许存在这些双变类的初等构造；我们尚不知道。

\begin{lemma}
\label{chow-lemma-chern-classes-defined}
设 $(S, \delta)$ 如情形 \ref{chow-situation-setup}。
设 $X$ 在 $S$ 上局部有限型，并设 $E \in D(\mathcal{O}_X)$
是完美对象。若下列条件之一成立，则 $E$ 的陈类已定义：
\begin{enumerate}
\item 存在包络 $f : Y \to X$，使得 $Lf^*E$ 在
$D(\mathcal{O}_Y)$ 中同构于有限局部自由
$\mathcal{O}_Y$-模组成的局部有界复形；
\item $E$ 可由有限局部自由 $\mathcal{O}_X$-模组成的有界复形表示；
\item $X$ 的不可约分支都是拟紧的；
\item $X$ 是拟紧的；
\item 存在 $S$ 上局部有限型概形的态射 $X \to X'$，
使得 $E$ 是 $X'$ 上某个陈类已定义的完美对象 $E'$ 的拉回；或者
\item 此处增补更多情形。
\end{enumerate}
\end{lemma}

\begin{proof}
条件 (1) 正是定义 \ref{chow-definition-defined-on-perfect} 的第 (1) 部分。
条件 (2) 蕴含 (1)。

\medskip\noindent
在 (3) 的情形，假设 $X$ 的不可约分支 $X_i$ 都是拟紧的。
把 $X_i$ 视为 $X$ 上的既约整闭子概形。态射
$\coprod X_i \to X$ 是包络。对每个 $i$，存在包络
$X'_i \to X_i$，使得 $X'_i$ 具有充足的可逆模族；见《态射进阶》命题 \ref{more-morphisms-proposition-envelope-with-resolution-property}。
注意 $f : Y = \coprod X'_i \to X$ 是包络；此处略去一个小细节。
由《概形的导出范畴》引理 \ref{perfect-lemma-resolution-property-ample-family}，
每个 $X'_i$ 都具有消解性质。因此 $D(\mathcal{O}_{X'_i})$
中的完美对象 $L(f|_{X'_i})^*E$ 可由有限局部自由
$\mathcal{O}_{X'_i}$-模组成的有界复形表示；见
《概形的导出范畴》引理 \ref{perfect-lemma-resolution-property-perfect-complex}。
这证明了 (3) 蕴含 (1)。

\medskip\noindent
第 (4) 部分蕴含 (3)。

\medskip\noindent
设 $g : X \to X'$ 和 $E'$ 如第 (5) 部分。则存在包络
$f' : Y' \to X'$，使得 $L(f')^*E'$ 由一个局部有界的
$\mathcal{O}_{Y'}$-模复形 $(\mathcal{E}')^\bullet$ 表示。
由引理 \ref{chow-lemma-base-change-envelope}，基变换
$f : Y \to X$ 是包络。此外，拉回复形
$\mathcal{E}^\bullet = g^*(\mathcal{E}')^\bullet$ 表示 $Lf^*E$，
从而 $E$ 的陈类已定义。
\end{proof}

\begin{lemma}
\label{chow-lemma-chern-classes-computed}
设 $(S, \delta)$ 如情形 \ref{chow-situation-setup}。
设 $X$ 在 $S$ 上局部有限型，并设 $E \in D(\mathcal{O}_X)$
是完美对象。假设 $E$ 的陈类已定义。对局部有限型态射
$g : W \to X$，若 $W$ 整，则存在交换图
$$
\xymatrix{
W' \ar[rd]_{g'} \ar[rr]_b & & W \ar[ld]^g \\
& X
}
$$
其中 $W'$ 整，$b : W' \to W$ 是真双有理态射，
$L(g')^*E$ 由常秩局部自由 $\mathcal{O}_{W'}$-模组成的有界复形
$\mathcal{E}^\bullet$ 表示，并且在 $A^p(W')$ 中有
$res(c_p(E)) = c_p(\mathcal{E}^\bullet)$。
\end{lemma}

\begin{proof}
选择包络 $f : Y \to X$，使得 $Lf^*E$ 在
$D(\mathcal{O}_Y)$ 中同构于有限局部自由
$\mathcal{O}_Y$-模组成的局部有界复形 $\mathcal{E}^\bullet$。
由引理 \ref{chow-lemma-base-change-envelope}，$f$ 的基变换
$Y \times_X W \to W$ 是包络。选择一点
$\xi \in Y \times_X W$，它映到 $W$ 的泛点且剩余域相同。
考虑以 $\xi$ 为泛点的整闭子概形
$W' \subset Y \times_X W$。投影
$Y \times_X W \to W$ 在 $W'$ 上的限制是真双有理态射
$b : W' \to W$。令 $g' = g \circ b$。最后，考虑拉回
$(W' \to Y)^*\mathcal{E}^\bullet$。这是 $W'$ 上有限局部自由模
组成的局部有界复形。由于 $W'$ 整，它是有界的，且各项具有常秩。
最后，按构造，$(W' \to Y)^*\mathcal{E}^\bullet$ 表示
$L(g')^*E$，而其第 $p$ 个陈类按构造给出 $c_p(E)$
经 $W' \to X$ 的限制。证明完毕。
\end{proof}

\begin{lemma}
\label{chow-lemma-commutative-chern-perfect}
设 $(S, \delta)$ 如情形 \ref{chow-situation-setup}。
设 $X$ 在 $S$ 上局部有限型。设 $E \in D(\mathcal{O}_X)$
是完美对象。若 $E$ 的陈类已定义，则
\begin{enumerate}
\item $c_p(E)$ 属于代数 $A^*(X)$ 的中心；并且
\item 若 $g : X' \to X$ 局部有限型且 $c \in A^*(X' \to X)$，则
$c \circ c_p(E) = c_p(Lg^*E) \circ c$。
\end{enumerate}
\end{lemma}

\begin{proof}
第 (1) 部分立即由第 (2) 部分推出。设 $g : X' \to X$ 和
$c \in A^*(X' \to X)$ 如 (2)。为证明
$c \circ c_p(E) - c_p(Lg^*E) \circ c = 0$，使用
引理 \ref{chow-lemma-bivariant-zero} 的判据。因此可以假设 $X$ 整；由
引理 \ref{chow-lemma-chern-classes-computed}，甚至可以假设
$E$ 由常秩有限局部自由 $\mathcal{O}_X$-模组成的有界复形
$\mathcal{E}^\bullet$ 表示。此时只须证明
$$
c \cap c_p(\mathcal{E}^\bullet) \cap [X] =
c_p(\mathcal{E}^\bullet) \cap c \cap [X] 
$$
在 $\CH_*(X')$ 中成立。这立即由
引理 \ref{chow-lemma-cap-commutative-chern} 以及
$c_p(\mathcal{E}^\bullet)$ 作为局部自由模
$\mathcal{E}^n$ 的陈类之多项式的构造得出。
\end{proof}

\begin{lemma}
\label{chow-lemma-additivity-on-perfect}
设 $(S, \delta)$ 如情形 \ref{chow-situation-setup}。
设 $X$ 在 $S$ 上局部有限型，并设
$$
E_1 \to E_2 \to E_3 \to E_1[1]
$$
是 $D(\mathcal{O}_X)$ 中完美对象的可区别三角。若下列条件之一成立：
\begin{enumerate}
\item 存在包络 $f : Y \to X$，使得
$Lf^*E_1 \to Lf^*E_2$ 可由有限局部自由
$\mathcal{O}_Y$-模组成的局部有界复形之间的映射表示；
\item $E_1 \to E_2$ 可由有限局部自由 $\mathcal{O}_X$-模组成的
局部有界复形之间的映射表示；
\item $X$ 的不可约分支都是拟紧的；
\item $X$ 是拟紧的；或者
\item 此处增补更多情形；
\end{enumerate}
则 $E_1$、$E_2$、$E_3$ 的陈类均已定义，并且
$c(E_2) = c(E_1) c(E_3)$、$ch(E_2) = ch(E_1) + ch(E_3)$，以及
$P_p(E_2) = P_p(E_1) + P_p(E_3)$。
\end{lemma}

\begin{proof}
设 $f : Y \to X$ 是包络，并设
$\alpha^\bullet : \mathcal{E}_1^\bullet \to \mathcal{E}_2^\bullet$
是有限局部自由 $\mathcal{O}_Y$-模组成的局部有界复形之间的映射，
它表示 $Lf^*E_1 \to Lf^*E_2$。于是锥
$C(\alpha)^\bullet$ 表示 $Lf^*E_3$。由于
$C(\alpha)^n = \mathcal{E}_2^n \oplus \mathcal{E}_1^{n + 1}$，
$C(\alpha)^\bullet$ 是有限局部自由 $\mathcal{O}_Y$-模组成的
局部有界复形。因此 $E_1$、$E_2$、$E_3$ 的陈类均已定义。
此外，回忆 $c_p(E_1)$ 定义为 $A^p(X)$ 中唯一的元素，
其在 $A^p(Y)$ 中的限制是 $c_p(\mathcal{E}_1^\bullet)$；
$E_2$ 和 $E_3$ 也同理。因此只须在
$\prod_{p \geq 0} A^p(Y)$ 中证明
$c(\mathcal{E}_2^\bullet) =
c(\mathcal{E}_1^\bullet) c(C(\alpha)^\bullet)$。
进一步，只须在限制到 $Y$ 的每个连通分支后证明它。
因而可以假设复形 $\mathcal{E}_1^\bullet$ 和
$\mathcal{E}_2^\bullet$ 都是常秩有限局部自由
$\mathcal{O}_Y$-模组成的有界复形。此时所需等式由
引理 \ref{chow-lemma-additivity-chern-classes} 的乘法性推出。
对 $ch$ 或 $P_p$，使用引理 \ref{chow-lemma-chern-character-additive}。

\medskip\noindent
上一段已经说明，在条件 (1) 成立时引理成立。
由于 (2) 蕴含 (1)，第二种情形也已处理。现假设 (3)。
完全仿照引理 \ref{chow-lemma-chern-classes-defined} 的证明，
可找到包络 $f : Y \to X$，使得 $Y$ 是不交并
$Y = \coprod Y_i$，其中每个分支都是拟紧（且拟分离）的
具有消解性质的概形。于是 $Lf^*E_1 \to Lf^*E_2$
在 $Y_i$ 上的限制可由有限局部自由模组成的有界复形之间的映射表示；见
《概形的导出范畴》命题 \ref{perfect-proposition-perfect-resolution-property}。
由此可见条件 (3) 蕴含条件 (1)。当然，条件 (4) 蕴含条件 (3)，
证明完毕。
\end{proof}

\begin{remark}
\label{chow-remark-splitting-principle-perfect}
完美复形的陈类一旦有定义，就满足一种分裂原理。具体地，假设
$(S, \delta), X, E$ 如定义 \ref{chow-definition-defined-on-perfect}，并且 $E$ 的陈类已定义。
假设要证明双变类 $c_p(E)$、$P_p(E)$ 与 $ch_p(E)$ 之间的某个关系。
为此，可以选择包络 $f : Y \to X$ 以及表示 $E$ 的、由有限局部自由
$\mathcal{O}_X$-模组成的局部有界复形 $\mathcal{E}^\bullet$。
由引理 \ref{chow-lemma-defined-by-envelope} 中的唯一性，
只须证明双变类 $c_p(\mathcal{E}^\bullet)$、
$P_p(\mathcal{E}^\bullet)$ 与 $ch_p(\mathcal{E}^\bullet)$
之间所需的关系。因此可将 $X$ 替换为 $Y$ 的一个连通分支，
并假设 $E$ 由常秩有限局部自由模组成的有界复形
$\mathcal{E}^\bullet$ 表示。利用分裂原理
（引理 \ref{chow-lemma-splitting-principle}），可以假设每个
$\mathcal{E}^i$ 都有一个滤过，其逐次商
$\mathcal{L}_{i, j}$ 是可逆模。令
$x_{i, j} = c_1(\mathcal{L}_{i, j})$，则
$$
c(E) =
\prod\nolimits_{i\text{ 为偶数}} (1 + x_{i, j})
\prod\nolimits_{i\text{ 为奇数}} (1 + x_{i, j})^{-1}
$$
并且
$$
P_p(E) =  \sum\nolimits_{i\text{ 为偶数}} (x_{i, j})^p -
\sum\nolimits_{i\text{ 为奇数}} (x_{i, j})^p
$$
对上面 $c(E)$ 的表达式形式地取对数，得到
$$
\log(c(E)) = \sum (-1)^{p - 1}\frac{P_p(E)}{p}
$$
考察例 \ref{chow-example-power-sum} 中多项式
$P_p$ 的构造可知，$P_p(E)$ 关于 $E$ 的陈类的表达式
与向量丛情形完全相同。换言之，
\begin{align*}
P_1(E) & = c_1(E), \\
P_2(E) & = c_1(E)^2 - 2c_2(E), \\
P_3(E) & = c_1(E)^3 - 3c_1(E)c_2(E) + 3c_3(E), \\
P_4(E) & = c_1(E)^4 - 4c_1(E)^2c_2(E) + 4c_1(E)c_3(E) + 2c_2(E)^2 - 4c_4(E),
\end{align*}
等等。另一方面，双变类 $P_0(E) = r(E) = ch_0(E)$
不能从 $E$ 的陈类 $c(E)$ 恢复；陈类并不记录复形的秩。
\end{remark}

\begin{lemma}
\label{chow-lemma-chern-classes-perfect-dual}
在情形 \ref{chow-situation-setup} 中，设 $X$ 在 $S$ 上局部有限型。
设 $E \in D(\mathcal{O}_X)$ 是陈类已定义的完美对象。则在
$A^i(X)$ 中，$c_i(E^\vee) = (-1)^i c_i(E)$、
$P_i(E^\vee) = (-1)^iP_i(E)$，且
$ch_i(E^\vee) = (-1)^ich_i(E)$。
\end{lemma}

\begin{proof}
第一种证明：仿照引理 \ref{chow-lemma-commutative-chern-perfect}
的证明，把问题约化到 $E$ 由常秩有限局部自由模组成的有界复形表示的情形，
再应用引理 \ref{chow-lemma-chern-classes-dual}。
第二种证明：使用注 \ref{chow-remark-splitting-principle-perfect}
中讨论的分裂原理，以及 $E^\vee$ 的陈根是 $E$ 的陈根的相反数这一事实。
\end{proof}

\begin{lemma}
\label{chow-lemma-chern-class-perfect-tensor-invertible}
在情形 \ref{chow-situation-setup} 中，设 $X$ 在 $S$ 上局部有限型。
设 $E$ 是 $D(\mathcal{O}_X)$ 中陈类已定义的完美对象。
设 $\mathcal{L}$ 是可逆 $\mathcal{O}_X$-模。则
$$
c_i(E \otimes \mathcal{L}) =
\sum\nolimits_{j = 0}^i
\binom{r - i + j}{j} c_{i - j}(E) c_1(\mathcal{L})^j
$$
只要 $E$ 具有常秩 $r \in \mathbf{Z}$。
\end{lemma}

\begin{proof}
当 $E$ 是秩为 $r$ 的局部自由对象时，这就是
引理 \ref{chow-lemma-chern-classes-E-tensor-L}。读者可由这个特殊情形
通过形式计算推出本引理。另一种方法是使用
注 \ref{chow-remark-splitting-principle-perfect} 的分裂原理。
在这种方法中，最终需要证明如下代数事实：若形式地写成
$$
\frac{\prod_{a = 1, \ldots, n} (1 + x_a)}{\prod_{n = 1, \ldots, m} (1 + y_b)}
= 1 + c_1 + c_2 + c_3 + \ldots
$$
其中 $c_i$ 是 $\mathbf{Z}[x_i, y_j]$ 中次数为 $i$ 的齐次多项式，
则有
$$
\frac{\prod_{a = 1, \ldots, n} (1 + x_a + t)}{\prod_{b = 1, \ldots, m} (1 + y_b + t)}
= \sum\nolimits_{i \geq 0} \sum\nolimits_{j = 0}^i
\binom{r - i + j}{j} c_{i - j} t^j
$$
其中 $r = n - m$。细节从略。
\end{proof}

\begin{lemma}
\label{chow-lemma-chern-classes-perfect-tensor-product}
在情形 \ref{chow-situation-setup} 中，设 $X$ 在 $S$ 上局部有限型。
设 $E$ 与 $F$ 是 $D(\mathcal{O}_X)$ 中陈类已定义的完美对象。则
$$
c_1(E \otimes_{\mathcal{O}_X}^\mathbf{L} F) =
r(E) c_1(\mathcal{F}) + r(F) c_1(\mathcal{E})
$$
而 $c_2(E \otimes_{\mathcal{O}_X}^\mathbf{L} F)$ 的表达式为
$$
r(E) c_2(F) + r(F) c_2(E) + {r(E) \choose 2} c_1(F)^2 +
(r(E)r(F) - 1) c_1(F)c_1(E) + {r(F) \choose 2} c_1(E)^2
$$
$A^*(X)$ 中的更高陈类也有类似表达式。相应地，在
$A^*(X) \otimes \mathbf{Q}$ 中有
$ch(E \otimes_{\mathcal{O}_X}^\mathbf{L} F) = ch(E) ch(F)$。
更精确地，在 $A^p(X)$ 中有
$$
P_p(E \otimes_{\mathcal{O}_X}^\mathbf{L} F) = \sum\nolimits_{p_1 + p_2 = p}
{p \choose p_1} P_{p_1}(E) P_{p_2}(F)
$$
。
\end{lemma}

\begin{proof}
选择包络 $f : Y \to X$，使得 $Lf^*E$ 与 $Lf^*F$
可由有限局部自由 $\mathcal{O}_X$-模组成的局部有界复形表示后，
由向量丛的相应结论作计算即可得到本结论；见引理 \ref{chow-lemma-chern-classes-tensor-product} 和
\ref{chow-lemma-chern-character-multiplicative}。
更好的证明也许是像注 \ref{chow-remark-splitting-principle-perfect}
那样使用分裂原理，把本引理约化为多项式环中的计算；
下一段将说明这些计算。

\medskip\noindent
设 $A$ 是交换环（对我们而言，它是 $X$ 的双变 Chow 环中
由陈类生成的子环）。设 $S$ 是有限集，并给定映射
$\epsilon : S \to \{\pm 1\}$ 和 $f : S \to A$。定义
$$
P_p(S, f , \epsilon) = \sum\nolimits_{s \in S} \epsilon(s) f(s)^p
$$
为 $A$ 中的元素。给定第二个三元组 $(S', \epsilon', f')$，
关于 $P_p$ 所须证明的等式是
$$
P_p(S \times S', f + f' , \epsilon \epsilon') = 
\sum\nolimits_{p_1 + p_2 = p}
{p \choose p_1} P_{p_1}(S, f, \epsilon) P_{p_2}(S', f', \epsilon')
$$
为说明它成立，可约化到以 $S \amalg S'$ 为变量的多项式环，
并验证每一项 $f(s)^if'(s')^j$ 在等式两边的系数相同。
为验证 $c_1(E \otimes_{\mathcal{O}_X}^\mathbf{L} F)$ 和
$c_2(E \otimes_{\mathcal{O}_X}^\mathbf{L} F)$ 的公式，
利用分裂原理把问题约化为在无挠环中检验这些公式。
然后使用注 \ref{chow-remark-splitting-principle-perfect}
中证明的 $P_j(E)$ 与 $c_i(E)$ 之间的关系。例如
$$
c_1(E \otimes F) = P_1(E \otimes F) = r(F)P_1(E) + r(E)P_1(F) =
r(F)c_1(E) + r(E)c_1(F)
$$
其中中间的等式使用了按定义 $r(E) = P_0(E)$。类似地，
\begin{align*}
& 2c_2(E \otimes F) \\
& = c_1(E \otimes F)^2 - P_2(E \otimes F) \\
& =
(r(F)c_1(E) + r(E)c_1(F))^2 -
r(F)P_2(E) - P_1(E)P_1(F) - r(E)P_2(F) \\
& =
(r(F)c_1(E) + r(E)c_1(F))^2 -
r(F)(c_1(E)^2 - 2c_2(E)) - c_1(E)c_1(F) - \\
& \quad r(E)(c_1(F)^2 - 2c_2(F))
\end{align*}
读者可验证，这与引理陈述中的公式相差因子 $2$ 后一致。
\end{proof}







\section{局部化陈类的一个简单情形}
\label{chow-section-preparation-localized-chern}

\noindent
本节讨论下一引理所构造的双变类的一些性质；其中大多数性质
立即由引理给出的刻画推出。我们建议读者跳过本节其余部分。

\begin{lemma}
\label{chow-lemma-silly}
设 $(S, \delta)$ 如情形 \ref{chow-situation-setup}。
设 $X$ 在 $S$ 上局部有限型。设 $i_j : X_j \to X$（$j = 1, 2$）
是闭浸入，并且在集合论意义下 $X = X_1 \cup X_2$。
设 $E_2 \in D(\mathcal{O}_{X_2})$ 是完美对象。假设
\begin{enumerate}
\item $E_2$ 的陈类已定义；
\item 限制 $E_2|_{X_1 \cap X_2}$ 为零，
分别地\ 同构于置于上同调次数 $0$ 的、秩 $< p$ 的有限局部自由
$\mathcal{O}_{X_1 \cap X_2}$-模。
\end{enumerate}
则存在典范双变类
$$
P'_p(E_2),\text{ 分别 }c'_p(E_2) \in A^p(X_2 \to X)
$$
它由下列性质刻画：
$$
P'_p(E_2) \cap i_{2, *} \alpha_2 = P_p(E_2) \cap \alpha_2
\quad\text{且}\quad
P'_p(E_2) \cap i_{1, *} \alpha_1 = 0,
$$
相应地，
$$
c'_p(E_2) \cap i_{2, *} \alpha_2 = c_p(E_2) \cap \alpha_2
\quad\text{且}\quad
c'_p(E_2) \cap i_{1, *} \alpha_1 = 0
$$
对 $\alpha_i \in \CH_k(X_i)$ 成立，并且在任意局部有限型基变换
$X' \to X$ 后也同样成立。
\end{lemma}

\begin{proof}
下文将直接使用第 \ref{chow-section-pre-derived} 节的内容，不再另行说明。

\medskip\noindent
假设 $E_2|_{X_1 \cap X_2}$ 为零。考虑局部有限型概形态射
$X' \to X$，并记 $i'_j : X'_j \to X'$ 为 $i_j$ 的基变换。
由引理 \ref{chow-lemma-exact-sequence-closed-chow}，任意元素
$\alpha' \in \CH_k(X')$ 都可写成
$i'_{1, *}\alpha'_1 + i'_{2, *}\alpha'_2$，其中
$\alpha'_2 \in \CH_k(X'_2)$ 的不确定性至多相差一个由
$X'_1 \cap X'_2 \to X'_2$ 推前的像中的元素。于是可以令
$P'_p(E_2) \cap \alpha' = P_p(E_2) \cap \alpha'_2 \in \CH_{k - p}(X'_2)$。
由于假设 $E_2$ 在 $X_1 \cap X_2$ 上的限制为零，
这个定义是良定的。

\medskip\noindent
若 $E_2|_{X_1 \cap X_2}$ 同构于置于上同调次数 $0$ 的、秩
$< p$ 的有限局部自由 $\mathcal{O}_{X_1 \cap X_2}$-模，
则由秩的考量有 $c_p(E_2|_{X_1 \cap X_2}) = 0$，
因而可以完全相同的方式论证。
\end{proof}

\begin{lemma}
\label{chow-lemma-silly-independent}
引理 \ref{chow-lemma-silly} 中 $A^p(X_2 \to X)$ 内的双变类
$P'_p(E_2)$，分别地\ $c'_p(E_2)$，与 $X_1$ 的选择无关。
\end{lemma}

\begin{proof}
设 $X_1' \subset X$ 是另一个闭子概形，使得在集合论意义下
$X = X'_1 \cup X_2$，并且限制 $E_2|_{X'_1 \cap X_2}$ 为零，
分别地\ 同构于置于上同调次数 $0$ 的、秩 $< p$ 的有限局部自由
$\mathcal{O}_{X'_1 \cap X_2}$-模。于是
$X = (X_1 \cap X'_1) \cup X_2$。因此任意元素
$\alpha \in \CH_k(X)$ 都可写成 $i_*\beta + i_{2, *}\alpha_2$，其中
$\alpha_2 \in \CH_k(X'_2)$ 且 $\beta \in \CH_k(X_1 \cap X'_1)$。
由此显然，
$P'_p(E_2) \cap \alpha = P_p(E_2) \cap \alpha_2 \in \CH_{k - p}(X_2)$，
分别地\  $c'_p(E_2) \cap \alpha = c_p(E_2) \cap \alpha_2 \in \CH_{k - p}(X_2)$，
与使用 $X_1$ 还是 $X'_1$ 无关。任意基变换后也同样成立。
\end{proof}

\begin{lemma}
\label{chow-lemma-base-change-silly}
在引理 \ref{chow-lemma-silly} 中，设 $X' \to X$ 是局部有限型态射。
记 $X' = X'_1 \cup X'_2$ 以及
$E'_2 \in D(\mathcal{O}_{X'_2})$ 为拉回到 $X'$ 的对象。
则使用 $X' = X'_1 \cup X'_2$ 和 $E_2'$，由
引理 \ref{chow-lemma-silly} 构造的 $A^p(X_2' \to X')$ 中的类
$P'_p(E_2')$，分别地\ $c'_p(E_2')$，是
$A^p(X_2 \to X)$ 中的类 $P'_p(E_2)$，分别地\ $c'_p(E_2)$，
的限制（注 \ref{chow-remark-restriction-bivariant}）。
\end{lemma}

\begin{proof}
立即由引理 \ref{chow-lemma-silly} 对这些类的刻画得出。
\end{proof}

\begin{lemma}
\label{chow-lemma-silly-silly}
在引理 \ref{chow-lemma-silly} 中，设 $E_2$ 是某个完美对象
$E \in D(\mathcal{O}_X)$ 的限制，并且 $E|_{X_1}$ 为零，
分别地\ 同构于置于上同调次数 $0$ 的、秩 $< p$ 的有限局部自由
$\mathcal{O}_{X_1}$-模。若 $E$ 的陈类已定义，则
$i_{2, *} \circ P'_p(E_2) = P_p(E)$，
分别地\ $i_{2, *} \circ c'_p(E_2) = c_p(E)$
（其中 $\circ$ 如引理 \ref{chow-lemma-push-proper-bivariant}）。
\end{lemma}

\begin{proof}
先假设 $E|_{X_1}$ 为零。采用引理 \ref{chow-lemma-silly}
证明中的记号，此情形下本引理由下式推出：
\begin{align*}
P_p(E) \cap \alpha'
& =
i'_{1, *}(P_p(E) \cap \alpha'_1) +
i'_{2, *}(P_p(E) \cap \alpha'_2) \\
& =
i'_{1, *}(P_p(E|_{X_1}) \cap \alpha'_1) +
i'_{2, *}(P'_p(E_2) \cap \alpha') \\
& =
i'_{2, *}(P'_p(E_2) \cap \alpha')
\end{align*}
当 $E|_{X_1}$ 同构于置于上同调次数 $0$ 的、秩 $< p$ 的
有限局部自由 $\mathcal{O}_{X_1}$-模时，论证类似。
\end{proof}

\begin{lemma}
\label{chow-lemma-silly-shrink}
在引理 \ref{chow-lemma-silly} 中，假设存在闭子概形
$X'_2 \subset X_2$ 和 $X_1 \subset X'_1 \subset X$，使得在集合论意义下
$X = X'_1 \cup X'_2$。假设 $E_2|_{X'_1 \cap X_2}$ 为零，
分别地\ 同构于置于次数 $0$ 的、秩 $< p$ 的有限局部自由模。则
$(X'_2 \to X_2)_* \circ P'_p(E_2|_{X'_2}) = P'_p(E_2)$，
分别地\  $(X'_2 \to X_2)_* \circ c'_p(E_2|_{X'_2}) = c_p(E_2)$
（其中 $\circ$ 如引理 \ref{chow-lemma-push-proper-bivariant}）。
\end{lemma}

\begin{proof}
立即由引理 \ref{chow-lemma-silly} 对这些类的刻画得出。
\end{proof}

\begin{lemma}
\label{chow-lemma-silly-commutes}
在引理 \ref{chow-lemma-silly} 中，设 $f : Y \to X$ 局部有限型，
并设 $c \in A^*(Y \to X)$。则在 $A^*(Y_2 \to Y)$ 中有
$$
c \circ P'_p(E_2) = P'_p(Lf_2^*E_2) \circ c
\quad\text{分别}\quad
c \circ c'_p(E_2) = c'_p(Lf_2^*E_2) \circ c
$$
其中 $f_2 : Y_2 \to X_2$ 是 $f$ 的基变换。
\end{lemma}

\begin{proof}
设 $\alpha \in \CH_k(X)$。可写成
$$
\alpha = \alpha_1 + \alpha_2
$$
其中 $\alpha_i \in \CH_k(X_i)$；这里省略了闭浸入
$X_i \to X$ 的推前。读者于是可检验
$c'_p(E_2) \cap \alpha = c_p(E_2) \cap \alpha_2$、
$c \cap c'_p(E_2) \cap \alpha = c \cap c_p(E_2) \cap \alpha_2$、
$c \cap \alpha = c \cap \alpha_1 + c \cap \alpha_2$，以及
$c'_p(Lf_2^*E_2) \cap c \cap \alpha = c_p(Lf_2^*E_2) \cap c \cap \alpha_2$。
最后应用引理 \ref{chow-lemma-commutative-chern-perfect}。
\end{proof}

\begin{lemma}
\label{chow-lemma-silly-compose}
在引理 \ref{chow-lemma-silly} 中，假设 $E_2|_{X_1 \cap X_2}$ 为零。则
\begin{align*}
P'_1(E_2) & = c'_1(E_2), \\
P'_2(E_2) & = c'_1(E_2)^2 - 2c'_2(E_2), \\
P'_3(E_2) & = c'_1(E_2)^3 - 3c'_1(E_2)c'_2(E_2) + 3c'_3(E_2), \\
P'_4(E_2) & = c'_1(E_2)^4 - 4c'_1(E_2)^2c'_2(E_2) +
4c'_1(E_2)c'_3(E_2) + 2c'_2(E_2)^2 - 4c'_4(E_2),
\end{align*}
等等，其中乘法如注 \ref{chow-remark-ring-loc-classes}。
\end{lemma}

\begin{proof}
零层的秩 $< 1$，所以类 $c'_p(E_2)$ 对所有 $p \geq 1$
都有定义，因而该陈述有意义。这些等式立即由
引理 \ref{chow-lemma-silly} 所构造类的刻画，以及
注 \ref{chow-remark-splitting-principle-perfect} 给出的
与 $E_2$ 的陈类取帽积的相应结论推出。
\end{proof}

\begin{lemma}
\label{chow-lemma-silly-sum-c}
设 $(S, \delta)$ 如情形 \ref{chow-situation-setup}。
设 $X$ 在 $S$ 上局部有限型。设 $i_j : X_j \to X$（$j = 1, 2$）
是闭浸入，使得在集合论意义下 $X = X_1 \cup X_2$。
设 $E, F \in D(\mathcal{O}_X)$ 是完美对象。假设
\begin{enumerate}
\item $E$ 与 $F$ 的陈类已定义；
\item 限制 $E|_{X_1 \cap X_2}$ 与 $F|_{X_1 \cap X_2}$
同构于置于上同调次数 $0$ 的有限局部自由
$\mathcal{O}_{X_1}$-模，其秩分别 $< p$ 和 $< q$。
\end{enumerate}
采用注 \ref{chow-remark-ring-loc-classes} 的记号，令
$$
c^{(p)}(E) = 1 + c_1(E) + \ldots + c_{p - 1}(E) +
c'_p(E|_{X_2}) + c'_{p + 1}(E|_{X_2}) + \ldots \in A^{(p)}(X_2 \to X)
$$
其中 $c'_p(E|_{X_2})$ 如引理 \ref{chow-lemma-silly}。
$c^{(q)}(F)$ 与 $c^{(p + q)}(E \oplus F)$ 也类似定义。
则在 $A^{(p + q)}(X_2 \to X)$ 中有
$c^{(p + q)}(E \oplus F) = c^{(p)}(E)c^{(q)}(F)$。
\end{lemma}

\begin{proof}
立即由引理 \ref{chow-lemma-silly} 对这些类的刻画以及
引理 \ref{chow-lemma-additivity-on-perfect} 中的可加性得出。
\end{proof}

\begin{lemma}
\label{chow-lemma-silly-sum-P}
设 $(S, \delta)$ 如情形 \ref{chow-situation-setup}。
设 $X$ 在 $S$ 上局部有限型。设 $i_j : X_j \to X$（$j = 1, 2$）
是闭浸入，使得在集合论意义下 $X = X_1 \cup X_2$。
设 $E, F \in D(\mathcal{O}_{X_2})$ 是完美对象。假设
\begin{enumerate}
\item $E$ 与 $F$ 的陈类已定义；
\item 限制 $E|_{X_1 \cap X_2}$ 与 $F|_{X_1 \cap X_2}$ 均为零。
\end{enumerate}
对 $p \geq 0$，记由引理 \ref{chow-lemma-silly} 构造的类为
$P'_p(E), P'_p(F), P'_p(E \oplus F) \in A^p(X_2 \to X)$。则
$P'_p(E \oplus F) = P'_p(E) + P'_p(F)$。
\end{lemma}

\begin{proof}
立即由引理 \ref{chow-lemma-silly} 对这些类的刻画以及
引理 \ref{chow-lemma-additivity-on-perfect} 中的可加性得出。
\end{proof}

\begin{lemma}
\label{chow-lemma-silly-tensor-invertible}
在引理 \ref{chow-lemma-silly} 中，假设 $E_2$ 具有常秩 $0$。
设 $\mathcal{L}$ 是可逆 $\mathcal{O}_X$-模。则
$$
c'_i(E_2 \otimes \mathcal{L}) =
\sum\nolimits_{j = 0}^i
\binom{- i + j}{j} c'_{i - j}(E_2) c_1(\mathcal{L})^j
$$
\end{lemma}

\begin{proof}
关于秩的假设蕴含 $E_2|_{X_1 \cap X_2}$ 为零。
因此 $c'_i(E_2)$ 对所有 $i \geq 1$ 都有定义，陈述有意义。
所述等式立即由引理 \ref{chow-lemma-chern-class-perfect-tensor-invertible}
以及引理 \ref{chow-lemma-silly} 对 $c'_i$ 的刻画得出。
\end{proof}

\begin{lemma}
\label{chow-lemma-silly-tensor-product}
在情形 \ref{chow-situation-setup} 中，设 $X$ 在 $S$ 上局部有限型。设
$$
X = X_1 \cup X_2 = X'_1 \cup X'_2
$$
是在集合论意义下把 $X$ 写成闭子概形之并的两种方式。
设 $E$、$E'$ 是 $D(\mathcal{O}_X)$ 中陈类已定义的完美对象。
假设 $E|_{X_1}$ 和 $E'|_{X'_1}$ 为零\footnote{大概存在本引理的一个变体，
其中只假设这些限制同构于秩 $< p$ 和 $< p'$ 的有限局部自由模。}，
对 $i = 1, 2$，记
\begin{enumerate}
\item $r = P'_0(E) \in A^0(X_2 \to X)$ 以及
$r' = P'_0(E') \in A^0(X'_2 \to X)$；
\item $\gamma_p = c'_p(E|_{X_2}) \in A^p(X_2 \to X)$ 以及
$\gamma'_p = c'_p(E'|_{X'_2}) \in A^p(X'_2 \to X)$；
\item $\chi_p = P'_p(E|_{X_2}) \in A^p(X_2 \to X)$ 以及
$\chi'_p = P'_p(E'|_{X'_2}) \in A^p(X'_2 \to X)$
\end{enumerate}
为引理 \ref{chow-lemma-silly} 所构造的类。则有
$$
c'_1((E \otimes_{\mathcal{O}_X}^\mathbf{L} E')|_{X_2 \cap X'_2}) =
r \gamma'_1 + r' \gamma_1
$$
在 $A^1(X_2 \cap X'_2 \to X)$ 中成立，并且
$$
c'_2((E \otimes_{\mathcal{O}_X}^\mathbf{L} E')|_{X_2 \cap X'_2}) =
r \gamma'_2 + r' \gamma_2 + {r \choose 2} (\gamma'_1)^2 +
(rr' - 1) \gamma'_1\gamma_1 + {r' \choose 2} \gamma_1^2
$$
在 $A^2(X_2 \cap X'_2 \to X)$ 中成立；更高陈类也类似。
同样有
$$
P'_p((E \otimes_{\mathcal{O}_X}^\mathbf{L} E')|_{X_2 \cap X'_2}) =
\sum\nolimits_{p_1 + p_2 = p}
{p \choose p_1} \chi_{p_1} \chi'_{p_2}
$$
在 $A^p(X_2 \cap X'_2 \to X)$ 中成立。
\end{lemma}

\begin{proof}
先说明该陈述有意义。事实上，
$X = (X_2 \cap X'_2) \cup Y$，其中
$Y = (X_1 \cap X'_1) \cup (X_1 \cap X'_2) \cup (X_2 \cap X'_1)$，
而对象 $E \otimes_{\mathcal{O}_X}^\mathbf{L} E'$ 在 $Y$ 上的限制为零。
所述等式由引理 \ref{chow-lemma-silly} 对这些类的刻画以及
引理 \ref{chow-lemma-chern-classes-perfect-tensor-product} 的等式推出。
细节从略。
\end{proof}







\section{无穷远处的 Gysin 类}
\label{chow-section-gysin-at-infty}

\noindent
本节讨论下一引理所构造的双变类。我们建议读者跳过本节其余部分。

\begin{lemma}
\label{chow-lemma-gysin-at-infty}
设 $(S, \delta)$ 如情形 \ref{chow-situation-setup}。
设 $X$ 在 $S$ 上局部有限型。设
$b : W \to \mathbf{P}^1_X$ 是真概形态射，并且在
$\mathbf{A}^1_X$ 上是同构。记
$i_\infty : W_\infty \to W$ 为除子
$D_\infty \subset \mathbf{P}^1_X$ 的逆像，其中补集为
$\mathbf{A}^1_X$。则存在典范双变类
$$
C \in A^0(W_\infty \to X)
$$
具有如下性质：对 $\alpha \in \CH_k(X)$，
$i_{\infty, *}(C \cap \alpha) = i_{0, *}\alpha$；
任意局部有限型基变换 $X' \to X$ 后也同样成立。
\end{lemma}

\begin{proof}
给定 $\alpha \in \CH_k(X)$，存在 $\beta \in \CH_{k + 1}(W)$，
其在 $b^{-1}(\mathbf{A}^1_X)$ 上的限制是 $\alpha$ 的平坦拉回；见
引理 \ref{chow-lemma-exact-sequence-open}。$\beta$ 的另一个选择与
$\beta$ 相差一个支撑在 $W_\infty$ 上的循环；见
引理 \ref{chow-lemma-restrict-to-open}。由于
(\ref{chow-equation-zero-infty}) 中的有效 Cartier 除子
$D_\infty \subset \mathbf{P}^1_X$ 的正规丛平凡，
Gysin 同态 $i_\infty^*$ 消去支撑在 $W_\infty$ 上的循环类；见
注 \ref{chow-remark-gysin-on-cycles}。因此令
$C \cap \alpha = i_\infty^*\beta$ 是良定的。

\medskip\noindent
由于 $W_\infty$ 与 $W_0 = X \times \{0\}$ 分别是
$\mathbf{P}^1_X$ 中有理等价的有效 Cartier 除子
$D_0, D_\infty$ 的拉回，$i_\infty^*\beta$ 与
$i_0^*\beta$ 映到 $W$ 上同一个循环类；确切地说，由
引理 \ref{chow-lemma-support-cap-effective-Cartier}，二者都表示类
$c_1(\mathcal{O}_{\mathbf{P}^1_X}(1)) \cap \beta$。
由 $\beta$ 的选择，在 $W_0 = X \times \{0\}$ 上作为循环有
$i_0^*\beta = \alpha$；例如见
引理 \ref{chow-lemma-relative-effective-cartier}。
因此得到引理所述的
$i_{\infty, *}(C \cap \alpha) = i_{0, *}\alpha$。

\medskip\noindent
注意，对任意局部有限型基变换 $X' \to X$，关于 $b$ 的假设仍然成立。
因此由上述同一构造得到运算
$C \cap - : \CH_k(X') \to \CH_k(W'_\infty)$。
为说明这族运算定义了双变类，考虑图
$$
\xymatrix{
& & & \CH_*(X) \ar[d]^{\text{平坦拉回}} \\
\CH_{* + 1}(W_\infty) \ar[r] \ar[rd]^0 &
\CH_{* + 1}(W) \ar[d]^{i_\infty^*} \ar[rr]^{\text{平坦拉回}} & &
\CH_{* + 1}(\mathbf{A}^1_X) \ar[r] \ar@{..>}[lld]^{C \cap -} &
0 \\
& \CH_*(W_\infty)
}
$$
这里取图示的 $X$，并考虑对任意 $X' \to X$ 的基变换图。已知平坦拉回和
$i_\infty^*$ 都是双变运算；见引理 \ref{chow-lemma-flat-pullback-bivariant} 和 \ref{chow-lemma-gysin-bivariant}。
于是一个形式论证（其中涉及概形及其 Chow 群的巨大交换图）
表明虚线箭头是双变运算。
\end{proof}

\begin{lemma}
\label{chow-lemma-base-change-gysin-at-infty}
在引理 \ref{chow-lemma-gysin-at-infty} 中，设 $X' \to X$ 是局部有限型态射。
记 $b' : W' \to \mathbf{P}^1_{X'}$ 和
$i'_\infty : W'_\infty \to W'$ 分别为 $b$ 与 $i_\infty$ 的基变换。
则使用 $b'$ 按引理 \ref{chow-lemma-gysin-at-infty} 构造的类
$C' \in A^0(W'_\infty \to X')$ 是 $C$ 的限制
（注 \ref{chow-remark-restriction-bivariant}）。
\end{lemma}

\begin{proof}
立即由构造以及平坦拉回和 $i_\infty^*$ 也满足类似陈述这一事实得出。
\end{proof}

\begin{lemma}
\label{chow-lemma-gysin-at-infty-independent}
在引理 \ref{chow-lemma-gysin-at-infty} 中，设 $g : W' \to W$ 是真态射，
并且在 $\mathbf{A}^1_X$ 上是同构。设
$C' \in A^0(W'_\infty \to X)$ 和 $C \in A^0(W_\infty \to X)$
是引理 \ref{chow-lemma-gysin-at-infty} 所构造的类。则在
$A^0(W_\infty \to X)$ 中有 $g_{\infty, *} \circ C' = C$。
\end{lemma}

\begin{proof}
令 $b' = b \circ g : W' \to \mathbf{P}^1_X$。
记 $i'_\infty : W'_\infty \to W'$ 为包含态射，并记
$g_\infty : W'_\infty \to W_\infty$ 为 $g$ 的限制。
给定 $\alpha \in \CH_k(X)$，选择 $\beta' \in \CH_{k + 1}(W')$，
使其在 $(b')^{-1}\mathbf{A}^1_X$ 上的限制为 $\alpha$ 的平坦拉回。
于是 $\beta = g_*\beta' \in \CH_{k + 1}(W)$ 在
$b^{-1}\mathbf{A}^1_X$ 上的限制为 $\alpha$ 的平坦拉回。由
引理 \ref{chow-lemma-closed-in-X-gysin}，
$i_\infty^*\beta = g_{\infty, *}(i'_\infty)^*\beta'$。
此等式及其经任意局部有限型态射 $X' \to X$ 基变换后的相应事实，
正给出引理的断言。
\end{proof}

\begin{lemma}
\label{chow-lemma-homomorphism-pre}
在引理 \ref{chow-lemma-gysin-at-infty} 中有
$C \circ (W_\infty \to X)_* \circ i_\infty^* = i_\infty^*$。
\end{lemma}

\begin{proof}
设 $\beta \in \CH_{k + 1}(W)$。记
$i_0 : X = X \times \{0\} \to W$ 为 $\mathbf{P}^1$ 中位于
$0$ 上方的纤维的闭浸入。则在 $\CH_k(X)$ 中有
$(W_\infty \to X)_* i_\infty^* \beta = i_0^*\beta$，因为
$i_{\infty, *}i_\infty^*\beta$ 与 $i_{0, *}i_0^*\beta$
表示 $W$ 上的同一个类（例如由
引理 \ref{chow-lemma-support-cap-effective-Cartier}），
所以推前到 $X$ 上也得到同一个类。
$\beta$ 在 $b^{-1}(\mathbf{A}^1_X)$ 上的限制进一步限制为
$i_0^*\beta = (W_\infty \to X)_* i_\infty^* \beta$ 的平坦拉回，
因为可在由 $i_0$ 拉回后检验这一点；见引理 \ref{chow-lemma-linebundle} 和 \ref{chow-lemma-linebundle-formulae}。
因此在计算 $(W_\infty \to X)_*i_\infty^*\beta$ 在 $C$
下的像时可以使用 $\beta$，从而得到所需结论。
\end{proof}

\begin{lemma}
\label{chow-lemma-gysin-at-infty-commutes}
在引理 \ref{chow-lemma-gysin-at-infty} 中，设 $f : Y \to X$
是局部有限型态射且 $c \in A^*(Y \to X)$。则在
$A^*(W_\infty \times_X Y \to X)$ 中有 $C \circ c = c \circ C$。
\end{lemma}

\begin{proof}
考虑交换图
$$
\xymatrix{
W_\infty \times_X Y \ar@{=}[r] &
W_{Y, \infty} \ar[r]_{i_{Y, \infty}} \ar[d] &
W_Y \ar[r]_{b_Y} \ar[d] &
\mathbf{P}^1_Y \ar[r]_{p_Y} \ar[d] &
Y \ar[d]^f \\
& W_\infty \ar[r]^{i_\infty} &
W \ar[r]^b &
\mathbf{P}^1_X \ar[r]^p &
X
}
$$
其中各方块均为笛卡尔方块。对元素 $\alpha \in \CH_k(X)$，
选择 $\beta \in \CH_{k + 1}(W)$，使其在
$b^{-1}(\mathbf{A}^1_X)$ 上的限制为 $\alpha$ 的平坦拉回。
于是 $c \cap \beta$ 是 $\CH_*(W_Y)$ 中的类，其在
$b_Y^{-1}(\mathbf{A}^1_Y)$ 上的限制为 $c \cap \alpha$ 的平坦拉回。
其次，由于 $c$ 是双变类，有
$$
i_{Y, \infty}^*(c \cap \beta) = c \cap i_\infty^*\beta
$$
这正说明 $C \cap c \cap \alpha = c \cap C \cap \alpha$。
对任意局部有限型基变换 $X' \to X$，同一论证仍然适用。
这就证明了引理。
\end{proof}





\section{局部化陈类的准备}
\label{chow-section-preparation-localized-chern-II}

\noindent
本节讨论下一引理所构造的双变类的一些性质。
我们建议读者跳过本节其余部分。

\begin{lemma}
\label{chow-lemma-localized-chern-pre}
设 $(S, \delta)$ 如情形 \ref{chow-situation-setup}。
设 $X$ 在 $S$ 上局部有限型，并设 $Z \subset X$ 是闭子概形。设
$$
b : W \longrightarrow \mathbf{P}^1_X
$$
是真概形态射。设 $Q \in D(\mathcal{O}_W)$ 是完美对象。
记 $W_\infty \subset W$ 为除子
$D_\infty \subset \mathbf{P}^1_X$ 的逆像，其中补集为
$\mathbf{A}^1_X$。假设
\begin{enumerate}
\item[(A0)] $Q$ 的陈类已定义
（第 \ref{chow-section-pre-derived} 节）；
\item[(A1)] $b$ 在 $\mathbf{A}^1_X$ 上是同构；
\item[(A2)] 存在闭子概形 $T \subset W_\infty$，
它包含 $W_\infty$ 中所有位于 $X \setminus Z$ 上方的点，
并且 $Q|_T$ 为零，分别地\ 同构于置于上同调次数 $0$ 的、
秩 $< p$ 的有限局部自由 $\mathcal{O}_T$-模。
\end{enumerate}
则存在典范双变类
$$
P'_p(Q),\text{ 分别 }c'_p(Q) \in A^p(Z \to X)
$$
使得
$(Z \to X)_* \circ P'_p(Q) = P_p(Q|_{X \times \{0\}})$，
分别地\ $(Z \to X)_* \circ c'_p(Q) = c_p(Q|_{X \times \{0\}})$。
\end{lemma}

\begin{proof}
记 $E \subset W_\infty$ 为 $Z$ 的逆像。于是
$W_\infty = T \cup E$，且 $b$ 诱导真态射 $E \to Z$。
记 $C \in A^0(W_\infty \to X)$ 为
引理 \ref{chow-lemma-gysin-at-infty} 所构造的双变类。记
$A^p(E \to W_\infty)$ 中的 $P'_p(Q|_E)$，分别地\ $c'_p(Q|_E)$，
为引理 \ref{chow-lemma-silly} 所构造的双变类。这是有意义的，
因为由假设 (A2)，$(Q|_E)|_{T \cap E}$ 为零，分别地\ 同构于
置于上同调次数 $0$ 的、秩 $< p$ 的有限局部自由
$\mathcal{O}_{E \cap T}$-模。定义
$$
P'_p(Q) = (E \to Z)_* \circ P'_p(Q|_E) \circ C,\text{ 分别 }
c'_p(Q) = (E \to Z)_* \circ c'_p(Q|_E) \circ C
$$
这是双变类；见引理 \ref{chow-lemma-push-proper-bivariant}。
由于 $E \to Z \to X$ 等于 $E \to W_\infty \to W \to X$，有
\begin{align*}
(Z \to X)_* \circ c'_p(Q)
& =
(W \to X)_* \circ i_{\infty, *} \circ (E \to W_\infty)_*
\circ c'_p(Q|_E) \circ C \\
& =
(W \to X)_* \circ i_{\infty, *} \circ c_p(Q|_{W_\infty}) \circ C \\
& =
(W \to X)_* \circ c_p(Q) \circ i_{\infty, *} \circ C \\
& =
(W \to X)_*\circ c_p(Q) \circ i_{0, *} \\
& =
(W \to X)_* \circ i_{0, *} \circ c_p(Q|_{X \times \{0\}}) \\
& =
c_p(Q|_{X \times \{0\}})
\end{align*}
第二个等式由引理 \ref{chow-lemma-silly-silly} 得出。
第三个等式使用了 $c_p(Q)$ 是双变类。
第四个等式由引理 \ref{chow-lemma-gysin-at-infty} 得出。
第五个等式再次使用了 $c_p(Q)$ 是双变类。
最后一个等式是因为按上述方式将 $W_0$ 与 $X$ 识别后，
$(W_0 \to W) \circ (W \to X)$ 是 $X$ 上的恒等态射。
完全相同的一串等式证明关于 $P'_p(Q)$ 的性质。
\end{proof}

\begin{lemma}
\label{chow-lemma-base-change-localized-chern-pre}
在引理 \ref{chow-lemma-localized-chern-pre} 中，设
$X' \to X$ 是局部有限型态射。记
$Z'$、$b' : W' \to \mathbf{P}^1_{X'}$ 与
$T' \subset W'_\infty$ 分别为 $Z$、
$b : W \to \mathbf{P}^1_X$ 与 $T \subset W_\infty$ 的基变换。
令 $Q' = (W' \to W)^*Q$。则使用 $b'$、$Q'$ 与 $T'$，
按引理 \ref{chow-lemma-localized-chern-pre} 构造的
$A^p(Z' \to X')$ 中的类 $P'_p(Q')$，分别地\ $c'_p(Q')$，
是 $A^p(Z \to X)$ 中的类 $P'_p(Q)$，分别地\ $c'_p(Q)$，
的限制（注 \ref{chow-remark-restriction-bivariant}）。
\end{lemma}

\begin{proof}
回忆该构造为
$$
P'_p(Q) = (E \to Z)_* \circ P'_p(Q|_E) \circ C,\text{ 分别 }
c'_p(Q) = (E \to Z)_* \circ c'_p(Q|_E) \circ C
$$
因此本引理由 $C$ 的相应基变换性质
（引理 \ref{chow-lemma-base-change-gysin-at-infty}），以及
引理 \ref{chow-lemma-silly} 所构造的类也具有同样基变换性质这一事实推出
（省略一个小细节）。
\end{proof}

\begin{lemma}
\label{chow-lemma-localized-chern-pre-independent}
引理 \ref{chow-lemma-localized-chern-pre} 中的双变类
$P'_p(Q)$，分别地\ $c'_p(Q)$，与闭子概形 $T$ 的选择无关。
此外，给定在 $\mathbf{A}^1_X$ 上为同构的真态射
$g : W' \to W$，令 $Q' = g^*Q$，则有
$P'_p(Q) = P'_p(Q')$，分别地\ $c'_p(Q) = c'_p(Q')$。
\end{lemma}

\begin{proof}
关于 $T$ 的独立性立即由引理 \ref{chow-lemma-silly-independent} 得出。

\medskip\noindent
设 $g : W' \to W$ 是在 $\mathbf{A}^1_X$ 上为同构的真态射。
注意，取 $T' = g^{-1}(T) \subset W'_\infty$ 得到满足 (A2) 的闭子概形，
所以与 $b' = b \circ g : W' \to \mathbf{P}^1_X$ 和 $Q'$ 对应的
$A^p(Z \to X)$ 中的运算子 $P'_p(Q')$，分别地\ $c'_p(Q')$，
已有定义。记 $E' \subset W'_\infty$ 为 $Z$ 在
$W'_\infty$ 中的逆像。回忆
$$
c'_p(Q') = (E' \to Z)_* \circ c'_p(Q'|_{E'}) \circ C'
$$
其中 $C' \in A^0(W'_\infty \to X)$ 且
$c'_p(Q'|_{E'}) \in A^p(E' \to W'_\infty)$。
由引理 \ref{chow-lemma-gysin-at-infty-independent}，
$g_{\infty, *} \circ C' = C$。还要注意，$E'$ 也是 $E$ 经
$g_\infty$ 在 $W'_\infty$ 中的逆像。又因 $Q' = g^*Q$，
$c'_p(Q'|_{E'})$ 正是 $c'_p(Q|_E)$ 在位于
$W'_\infty$ 上方的概形上的限制；见
注 \ref{chow-remark-restriction-bivariant}。因此
\begin{align*}
c'_p(Q')
& = 
(E' \to Z)_* \circ c'_p(Q'|_{E'}) \circ C' \\
& =
(E \to Z)_* \circ (E' \to E)_* \circ c'_p(Q|_E) \circ C' \\
& =
(E \to Z)_* \circ c'_p(Q|_E) \circ g_{\infty, *} \circ C' \\
& =
(E \to Z)_* \circ c'_p(Q|_E) \circ C \\
& =
c'_p(Q)
\end{align*}
第三个等式使用了 $c'_p(Q|_E)$ 是双变类，因而与真推前交换。
等式 $P'_p(Q) = P'_p(Q')$ 完全同理证明。
\end{proof}

\begin{lemma}
\label{chow-lemma-homomorphism}
在引理 \ref{chow-lemma-localized-chern-pre} 中，假设 $Q|_T$
同构于秩 $< p$ 的有限局部自由 $\mathcal{O}_T$-模。
记 $C \in A^0(W_\infty \to X)$ 为
引理 \ref{chow-lemma-gysin-at-infty} 的类。则
$$
C \circ c_p(Q|_{X \times \{0\}}) =
C \circ (Z \to X)_* \circ c'_p(Q) = c_p(Q|_{W_\infty}) \circ C
$$
\end{lemma}

\begin{proof}
第一个等式成立，因为由引理 \ref{chow-lemma-localized-chern-pre}，
$c_p(Q|_{X \times \{0\}}) =
(Z \to X)_* \circ c'_p(Q)$。第二个等式可逐个循环类证明
（见引理 \ref{chow-lemma-bivariant-zero}）。由于引理中的双变类构造
与基变换相容，可取某个 $\alpha \in \CH_k(X)$，并只须证明
$C \cap (Z \to X)_*(c'_p(Q) \cap \alpha) =
c_p(Q|_{W_\infty}) \cap C \cap \alpha$。注意
\begin{align*}
C \cap (Z \to X)_*(c'_p(Q) \cap \alpha)
& =
C \cap (Z \to X)_* (E \to Z)_*(c'_p(Q|_E) \cap C \cap \alpha) \\
& =
C \cap (W_\infty \to X)_*(E \to W_\infty)_*(c'_p(Q|_E) \cap C \cap \alpha) \\
& =
C \cap (W_\infty \to X)_*(E \to W_\infty)_*(c'_p(Q|_E) \cap i_\infty^*\beta) \\
& =
C \cap (W_\infty \to X)_*(c_p(Q|_{W_\infty}) \cap i_\infty^*\beta) \\
& =
C \cap (W_\infty \to X)_*i_\infty^*(c_p(Q) \cap \beta) \\
& =
i_\infty^*(c_p(Q) \cap \beta) \\
& =
c_p(Q|_{W_\infty}) \cap i_\infty^*\beta \\
& =
c_p(Q|_{W_\infty}) \cap C \cap \alpha
\end{align*}
这正是所需结论。第一个等式使用了
$c'_p(Q) = (E \to Z)_* \circ c'_p(Q|_E) \circ C$，其中
$E \subset W_\infty$ 是 $Z$ 的逆像，而 $c'_p(Q|_E)$ 是
引理 \ref{chow-lemma-silly} 所构造的类。第二个等式只是说
$E \to Z \to X$ 等于 $E \to W_\infty \to X$。
在第三个等式中，选择 $\beta \in \CH_{k + 1}(W)$，使其在
$b^{-1}(\mathbf{A}^1_X)$ 上的限制是 $\alpha$ 的平坦拉回，
从而按构造有 $C \cap \alpha = i_\infty^*\beta$。
第四个等式是引理 \ref{chow-lemma-silly-silly}。
第五个等式使用了 $c_p(Q)$ 是双变类，因而与 $i_\infty^*$ 交换。
第六个等式是引理 \ref{chow-lemma-homomorphism-pre}。
第七个等式再次使用了 $c_p(Q)$ 是双变类。
最后一个等式由 $C \cap \alpha = i_\infty^*\beta$ 得出。
\end{proof}

\begin{lemma}
\label{chow-lemma-homomorphism-commute}
在引理 \ref{chow-lemma-localized-chern-pre} 中，设 $Y \to X$
是局部有限型态射，并设 $c \in A^*(Y \to X)$ 是双变类。则在
$A^*(Y \times_X Z \to X)$ 中有
$$
P'_p(Q) \circ c = c \circ P'_p(Q)
\quad\text{分别}\quad
c'_p(Q) \circ c = c \circ c'_p(Q)
$$
\end{lemma}

\begin{proof}
设 $E \subset W_\infty$ 为 $Z$ 的逆像。回忆
$P'_p(Q) = (E \to Z)_* \circ P'_p(Q|_E) \circ C$，
分别地\ $c'_p(Q) = (E \to Z)_* \circ c'_p(Q|_E) \circ C$，
其中 $C$ 如引理 \ref{chow-lemma-gysin-at-infty}，而
$P'_p(Q|_E)$，分别地\ $c'_p(Q|_E)$，如引理 \ref{chow-lemma-silly}。
由引理 \ref{chow-lemma-gysin-at-infty-commutes}，$C$ 与 $c$ 交换；
由引理 \ref{chow-lemma-silly-commutes}，$P'_p(Q|_E)$，
分别地\ $c'_p(Q|_E)$，与 $c$ 交换。由于 $c$ 是双变类，
按定义它与沿 $E \to Z$ 的真推前交换。证明完毕。
\end{proof}

\begin{lemma}
\label{chow-lemma-localized-chern-pre-compose}
在引理 \ref{chow-lemma-localized-chern-pre} 中，假设 $Q|_T$ 为零。
在 $A^*(Z \to X)$ 中有
\begin{align*}
P'_1(Q) & = c'_1(Q), \\
P'_2(Q) & = c'_1(Q)^2 - 2c'_2(Q), \\
P'_3(Q) & = c'_1(Q)^3 - 3c'_1(Q)c'_2(Q) + 3c'_3(Q), \\
P'_4(Q) & = c'_1(Q)^4 - 4c'_1(Q)^2c'_2(Q) +
4c'_1(Q)c'_3(Q) + 2c'_2(Q)^2 - 4c'_4(Q),
\end{align*}
等等，其中乘法如注 \ref{chow-remark-ring-loc-classes}。
\end{lemma}

\begin{proof}
零层的秩 $< 1$，所以类 $c'_p(Q)$ 对所有 $p \geq 1$
都有定义，因而该陈述有意义。在
引理 \ref{chow-lemma-localized-chern-pre} 的证明中，我们用
引理 \ref{chow-lemma-gysin-at-infty} 的双变类
$C \in A^0(W_\infty \to X)$，以及引理 \ref{chow-lemma-silly}
给出的双变类 $P'_p(Q|_E)$ 和 $c'_p(Q|_E)$，构造了类
$P'_p(Q)$ 与 $c'_p(Q)$；这里 $Q|_E$ 是 $Q$ 在
$W_\infty$ 中 $Z$ 的逆像 $E$ 上的限制。由
引理 \ref{chow-lemma-silly-compose}，$P'_p(Q|_E)$ 与
$c'_p(Q|_E)$ 之间满足所需关系。回忆
$$
P'_p(Q) = (E \to Z)_* \circ P'_p(Q|_E) \circ C
\quad\text{且}\quad
c'_p(Q) = (E \to Z)_* \circ c'_p(Q|_E) \circ C
$$
为完成证明，只须说明注 \ref{chow-remark-ring-loc-classes}
在各类 $a_p = c'_p(Q)$ 上定义的乘法，与在各类
$b_p = c'_p(Q|_E)$ 上定义的乘法相容：
$$
a_{p_1}a_{p_2} \ldots a_{p_r} =
(E \to Z)_* \circ b_{p_1}b_{p_2} \ldots b_{p_r} \circ C
$$
省略一些细节。当 $r = 1$ 时这显然成立。若 $r > 1$，则由
注 \ref{chow-remark-res-push}，注 \ref{chow-remark-ring-loc-classes}
中的乘法是在乘积 $a_{p_1}a_{p_2} \ldots a_{p_r}$，分别地\ 
$b_{p_1}b_{p_2} \ldots b_{p_r}$ 的各因子之间插入
$(Z \to X)_*$，分别地\ $(E \to W_\infty)_*$，
然后取双变类的复合。由引理 \ref{chow-lemma-silly}，
$$
(E \to W_\infty)_* \circ b_{p_i} = c_{p_i}(Q|_{W_\infty})
$$
而由引理 \ref{chow-lemma-homomorphism}，
$$
C \circ (Z \to X)_* \circ a_{p_i} = c_{p_i}(Q|_{W_\infty}) \circ C
$$
对 $i = 2, \ldots, r$ 成立。计算表明，所需等式的两边均化简为
$$
(E \to Z)_* \circ c'_{p_1}(Q|_E) \circ
c_{p_2}(Q|_{W_\infty}) \circ \ldots \circ
c_{p_r}(Q|_{W_\infty}) \circ C
$$
证明完毕。
\end{proof}

\begin{lemma}
\label{chow-lemma-localized-chern-pre-sum-c}
在引理 \ref{chow-lemma-localized-chern-pre} 中，假设 $Q|_T$
同构于秩 $< p$ 的有限局部自由 $\mathcal{O}_T$-模。
再假设存在另一个陈类已定义的完美对象
$Q' \in D(\mathcal{O}_W)$，且 $Q'|_T$ 同构于置于上同调次数
$0$ 的、秩 $< p'$ 的有限局部自由 $\mathcal{O}_T$-模。
采用注 \ref{chow-remark-ring-loc-classes} 的记号，令
$$
c^{(p)}(Q) = 1 + c_1(Q|_{X \times \{0\}}) + \ldots +
c_{p - 1}(Q|_{X \times \{0\}}) +
c'_{p}(Q) + c'_{p + 1}(Q) + \ldots
$$
为 $A^{(p)}(Z \to X)$ 中的元素，其中对 $i \geq p$，
$c'_i(Q)$ 如引理 \ref{chow-lemma-localized-chern-pre}。
$c^{(p')}(Q')$ 与 $c^{(p + p')}(Q \oplus Q')$ 也类似定义。
则在 $A^{(p + p')}(Z \to X)$ 中有
$c^{(p + p')}(Q \oplus Q') = c^{(p)}(Q)c^{(p')}(Q')$。
\end{lemma}

\begin{proof}
回忆，对 $i \geq p$，$c'_i(Q)$ 在 $A^p(X)$ 中的像等于
$c_i(Q|_{X \times \{0\}})$；$Q'$ 和 $Q \oplus Q'$ 的情形也类似，
见引理 \ref{chow-lemma-localized-chern-pre}。因此次数
$< p + p'$ 的等式由引理 \ref{chow-lemma-additivity-on-perfect}
的可加性推出。

\medskip\noindent
取 $n \geq p + p'$。仿照引理 \ref{chow-lemma-localized-chern-pre}
的证明，记 $E \subset W_\infty$ 为 $Z$ 的逆像。注意，由
引理 \ref{chow-lemma-silly-sum-c}，在
$A^{(p + p')}(E \to W_\infty)$ 中有
$$
c^{(p + p')}(Q|_E \oplus Q'|_E) =
c^{(p)}(Q|_E)c^{(p')}(Q'|_E)
$$
又因按构造
$$
c'_p(Q \oplus Q') = (E \to Z)_* \circ c'_p(Q|_E \oplus Q'|_E) \circ C
$$
所以只须对所有 $i + j = n$ 证明
$$
(E \to Z)_* \circ c^{(p)}_i(Q|_E)c^{(p')}_j(Q'|_E) \circ C
=
c^{(p)}_i(Q)c^{(p')}_j(Q')
$$
在 $A^n(Z \to X)$ 中成立，其中两边的乘法均为
注 \ref{chow-remark-ring-loc-classes} 中的乘法。依照
$i \geq p$、$j \geq p'$ 是否成立，共有三种情形。

\medskip\noindent
假设 $i \geq p$ 且 $j \geq p'$。此时乘积的定义是在两个因子之间
插入 $(E \to W_\infty)_*$，分别地\ $(Z \to X)_*$，
再取双变类的复合；见注 \ref{chow-remark-res-push}。
换言之，只须证明
$$
(E \to Z)_* \circ c'_i(Q|_E) \circ
(E \to W_\infty)_* \circ c'_j(Q'|_E) \circ C =
c'_i(Q) \circ (Z \to X)_* \circ c'_j(Q')
$$
由引理 \ref{chow-lemma-silly}，左边等于
$$
(E \to Z)_* \circ c'_i(Q|_E) \circ c_j(Q'|_{W_\infty}) \circ C
$$
由于 $c'_i(Q) = (E \to Z)_* \circ c'_i(Q|_E) \circ C$，
右边等于
$$
(E \to Z)_* \circ c'_i(Q|_E) \circ C \circ (Z \to X)_* \circ c'_j(Q')
$$
由引理 \ref{chow-lemma-homomorphism}，这立即等于上述左边。

\medskip\noindent
假设 $i \geq p$ 且 $j < p$。展开此情形中的乘积，只须证明
$$
(E \to Z)_* \circ c'_i(Q|_E) \circ c_j(Q'|_{W_\infty}) \circ C =
c'_i(Q) \circ c_j(Q'|_{X \times \{0\}})
$$
再次使用 $c'_i(Q) = (E \to Z)_* \circ c'_i(Q|_E) \circ C$，
只须证明 $c_j(Q'|_{W_\infty}) \circ C =
C \circ c_j(Q'|_{X \times \{0\}})$；这是
引理 \ref{chow-lemma-homomorphism} 的一部分。

\medskip\noindent
假设 $i < p$ 且 $j \geq p'$。展开此情形中的乘积，只须证明
$$
(E \to Z)_* \circ c_i(Q|_E) \circ c'_j(Q'|_E) \circ C =
c_i(Q|_{Z \times \{0\}}) \circ c'_j(Q')
$$
然而，由于 $c'_j(Q|_E)$ 与 $c'_j(Q')$ 是双变类，
它们同与陈类取帽积交换
（引理 \ref{chow-lemma-cap-commutative-chern}）。因此只须证明
$$
(E \to Z)_* \circ c'_j(Q'|_E) \circ c_i(Q|_{W_\infty}) \circ C =
c'_j(Q') \circ c_i(Q|_{X \times \{0\}})
$$
这就约化到上一段讨论的情形。
\end{proof}

\begin{lemma}
\label{chow-lemma-localized-chern-pre-sum-P}
在引理 \ref{chow-lemma-localized-chern-pre} 中，假设 $Q|_T$ 为零。
再假设有另一个陈类已定义的完美对象
$Q' \in D(\mathcal{O}_W)$，且其限制 $Q'|_T$ 为零。
此时由引理 \ref{chow-lemma-localized-chern-pre} 构造的类
$P'_p(Q), P'_p(Q'), P'_p(Q \oplus Q') \in A^p(Z \to X)$
满足 $P'_p(Q \oplus Q') = P'_p(Q) + P'_p(Q')$。
\end{lemma}

\begin{proof}
立即由这些类的构造以及引理 \ref{chow-lemma-silly-sum-P} 得出。
\end{proof}









 
\section{局部化陈类}
\label{chow-section-localized-chern}

\noindent
先概述构造。设 $F$ 是域，$X$ 是 $F$ 上的簇，
$E$ 是 $D(\mathcal{O}_X)$ 中的完美对象，并设
$Z \subset X$ 是闭子概形，使得 $E|_{X \setminus Z} = 0$。
目标是构造元素
$$
c_p(Z \to X, E) \in A^p(Z \to X)
$$
为此将构造图
$$
\xymatrix{
W \ar[d]_f \ar[r]_q & X \\
\mathbf{P}^1_F
}
$$
以及 $D(\mathcal{O}_W)$ 中的完美对象 $Q$，使得
\begin{enumerate}
\item $f$ 平坦，且 $f$、$q$ 都是真态射；对
$t \in \mathbf{P}^1_F$，记 $W_t$ 为 $f$ 的纤维，
$q_t : W_t \to X$ 为 $q$ 的限制，并令 $Q_t = Q|_{W_t}$；
\item 对 $t \in \mathbf{A}^1_F$，$q_t : W_t \to X$
是同构且 $Q_t = q_t^*E$；
\item $q_\infty : W_\infty \to X$ 在 $X \setminus Z$ 上是同构；
\item 若 $T \subset W_\infty$ 是
$q_\infty^{-1}(X \setminus Z)$ 的闭包，则 $Q_\infty|_T$ 为零。
\end{enumerate}
其思想是把它看成由 $t \in \mathbf{P}^1$ 参数化的族
$\{(W_t, Q_t)\}$。当 $t \not = \infty$ 时，在
$Q_t = X$ 的 Chow 群上，$c_p(Q_t)$ 就是 $c_p(E)$。
而当 $t = \infty$ 时，由于 $Q_\infty|_T = 0$，
$c_p(Q_\infty)$ 把 $Q_\infty$ 上的类送到支撑在
$E = q_\infty^{-1}(Z)$ 上的类。把 $E$ 看作
$q_\infty : W_\infty \to X$ 的例外轨迹。由于任意
$\alpha \in \CH_*(X)$ 都产生一个循环“族”
$\alpha_t \in \CH_*(W_t)$，可将
$c_p(Z \to X, E) \cap \alpha$ 定义为推前
$(E \to Z)_*(c_p(Q_\infty) \cap \alpha_\infty)$。

\medskip\noindent
使该方案成立有两个主要要素：(1) $W$ 与 $Q$ 的构造是某种代数的
Macpherson 图构造；它在《平坦性进阶》第 \ref{flat-section-blowup-complexes-III} 节中完成。
(2) 给定 $W$ 与 $Q$ 后实际类的构造在第 \ref{chow-section-preparation-localized-chern-II} 节中完成，并依赖于第 \ref{chow-section-gysin-at-infty} 节和
\ref{chow-section-preparation-localized-chern}。

\begin{situation}
\label{chow-situation-loc-chern}
设 $(S, \delta)$ 如情形 \ref{chow-situation-setup}。
设 $X$ 是 $S$ 上局部有限型概形，$i : Z \to X$ 是闭浸入，
$E \in D(\mathcal{O}_X)$ 是对象，并设 $p \geq 0$。假设
\begin{enumerate}
\item $E$ 是 $D(\mathcal{O}_X)$ 中的完美对象；
\item 限制 $E|_{X \setminus Z}$ 为零，分别地\ 同构于置于上同调次数
$0$ 的、秩 $< p$ 的有限局部自由
$\mathcal{O}_{X \setminus Z}$-模；并且
\item 下列条件至少有一个\footnote{初次阅读时请忽略此技术条件；
相关讨论见注 \ref{chow-remark-loc-chern}。}成立：
(a) $X$ 拟紧；
(b) $X$ 的不可约分支拟紧；
(c) 存在表示 $E$ 的有限局部自由 $\mathcal{O}_X$-模组成的
局部有界复形；或者
(d) 存在 $S$ 上局部有限型概形的态射 $X \to X'$，使得
$E$ 是 $X'$ 上某个完美对象的拉回，且 $X'$ 的不可约分支拟紧。
\end{enumerate}
\end{situation}

\begin{lemma}
\label{chow-lemma-independent-loc-chern}
在情形 \ref{chow-situation-loc-chern} 中，存在典范双变类
$$
P_p(Z \to X, E) \in A^p(Z \to X),
\quad\text{分别}\quad
c_p(Z \to X, E) \in A^p(Z \to X)
$$
其性质为，在 $X$ 上作为双变类有
\begin{equation}
\label{chow-equation-defining-property-localized-classes}
i_* \circ P_p(Z \to X, E) = P_p(E),
\quad\text{分别}\quad
i_* \circ c_p(Z \to X, E) = c_p(E)
\end{equation}
（其中 $\circ$ 如引理 \ref{chow-lemma-push-proper-bivariant}）。
\end{lemma}

\begin{proof}
这些双变类的构造如下。令
$$
b : W \longrightarrow \mathbf{P}^1_X
\quad\text{且}\quad
T \longrightarrow W_\infty
\quad\text{且}\quad
Q
$$
分别为《平坦性进阶》第 \ref{flat-section-blowup-complexes-III} 节和引理 \ref{flat-lemma-graph-construction} 中构造的爆破、
$D(\mathcal{O}_W)$ 中的完美对象 $Q$ 以及闭浸入。
令 $T' \subset T$ 为开闭子概形，使得 $Q|_{T'}$ 为零，
分别地\ 同构于置于上同调次数 $0$ 的、秩 $< p$ 的有限局部自由
$\mathcal{O}_{T'}$-模。由情形 \ref{chow-situation-loc-chern}
的条件 (2)，态射
$$
T' \to T \to W_\infty \to X
$$
在 $X$ 的开子概形 $X \setminus Z$ 上都是概形同构。
下面检验 $Q$ 的陈类已定义。回忆按构造
$Q|_{X \times \{0\}} \cong E$，于是使用 $W, b, Q, T'$，
由引理 \ref{chow-lemma-localized-chern-pre} 构造的双变类给出
$$
P_p(Z \to X, E) = P'_p(Q) \in A^p(Z \to X)
$$
以及
$$
c_p(Z \to X, E) = c'_p(Q) \in A^p(Z \to X)
$$
并满足 (\ref{chow-equation-defining-property-localized-classes})。

\medskip\noindent
本段证明 $Q$ 的陈类已定义
（定义 \ref{chow-definition-defined-on-perfect}）；建议读者跳过。
若情形 \ref{chow-situation-loc-chern} 的假设 (3)(a) 或 (3)(b) 成立，
即 $X$ 具有拟紧不可约分支，则 $W$ 也如此（因为
$W \to X$ 是真态射）。所以由
引理 \ref{chow-lemma-chern-classes-defined}，$D(\mathcal{O}_W)$
中任意完美对象的陈类均已定义。若 (3)(c) 成立，即 $E$
可由有限局部自由模组成的局部有界复形表示，则由《平坦性进阶》引理 \ref{flat-lemma-graph-construction} 的第 (5) 部分，
对象 $Q$ 可由有限局部自由 $\mathcal{O}_W$-模组成的局部有界复形表示。
故 $Q$ 的陈类已定义。最后假设 (3)(d) 成立，即存在
$S$ 上局部有限型概形的态射 $X \to X'$，使得 $E$ 是
$X'$ 上完美对象 $E'$ 的拉回，且 $X'$ 的不可约分支拟紧。
令 $b' : W' \to \mathbf{P}^1_{X'}$ 和
$Q' \in D(\mathcal{O}_{W'})$ 为从三元组
$(\mathbf{P}^1_{X'}, (\mathbf{P}^1_{X'})_\infty, L(p')^*E')$
出发，按《平坦性进阶》第 \ref{flat-section-blowup-complexes-III} 节构造的态射与完美对象。
由上述讨论，$Q'$ 的陈类已定义。由于 $b$ 和 $b'$ 是应用
《平坦性进阶》引理 \ref{flat-lemma-complex-and-divisor-blowup-pre} 构造的，
由《平坦性进阶》引理 \ref{flat-lemma-complex-and-divisor-blowup-base-change}，
存在态射 $W \to W'$，使得 $Q = L(W \to W')^*Q'$。
于是由引理 \ref{chow-lemma-chern-classes-defined}，
$Q$ 的陈类已定义。
\end{proof}

\begin{definition}
\label{chow-definition-localized-chern}
设 $(S, \delta)$、$X$、$E \in D(\mathcal{O}_X)$ 以及
$i : Z \to X$ 如情形 \ref{chow-situation-loc-chern}。
\begin{enumerate}
\item 若限制 $E|_{X \setminus Z}$ 为零，则对所有 $p \geq 0$，
由引理 \ref{chow-lemma-independent-loc-chern} 的构造定义
$$
P_p(Z \to X, E) \in A^p(Z \to X)
$$
并用公式
$$
ch(Z \to X, E) =
\sum\nolimits_{p = 0, 1, 2, \ldots} \frac{P_p(Z \to X, E)}{p!}
\quad\text{于}\quad \prod\nolimits_{p \geq 0} A^p(Z \to X) \otimes \mathbf{Q}
$$
定义{\it 局部化陈特征}。
\item 若限制 $E|_{X \setminus Z}$ 同构于置于上同调次数
$0$ 的、秩 $< p$ 的有限局部自由
$\mathcal{O}_{X \setminus Z}$-模，则由
引理 \ref{chow-lemma-independent-loc-chern} 的构造定义
{\it 局部化第 $p$ 陈类} $c_p(Z \to X, E)$。
\end{enumerate}
\end{definition}

\noindent
在该定义的情形中，假设 $E|_{X \setminus Z}$ 为零。
那么确有等式
$$
i_* \circ ch(Z \to X, E) = ch(E)
$$
在 $A^*(X) \otimes \mathbf{Q}$ 中成立，因为上面已经证明了等式
(\ref{chow-equation-defining-property-localized-classes})。

\medskip\noindent
下面给出一个重要的一致性检验。

\begin{lemma}
\label{chow-lemma-base-change-loc-chern}
在情形 \ref{chow-situation-loc-chern} 中，设
$f : X' \to X$ 是局部有限型概形态射。记
$E' = f^*E$ 且 $Z' = f^{-1}(Z)$。则定义 \ref{chow-definition-localized-chern} 中的双变类
$$
P_p(Z' \to X', E') \in A^p(Z' \to X'),
\quad\text{分别}\quad
c_p(Z' \to X', E') \in A^p(Z' \to X')
$$
即使用 $X', Z', E'$ 按引理 \ref{chow-lemma-independent-loc-chern}
构造的类，是双变类 $P_p(Z \to X, E) \in A^p(Z \to X)$，分别地\ 
$c_p(Z \to X, E) \in A^p(Z \to X)$，的限制
（注 \ref{chow-remark-restriction-bivariant}）。
\end{lemma}

\begin{proof}
记 $p : \mathbf{P}^1_X \to X$ 与
$p' : \mathbf{P}^1_{X'} \to X'$ 为结构态射。回忆
$b : W \to \mathbf{P}^1_X$ 与
$b' : W' \to \mathbf{P}^1_{X'}$ 是在《平坦性进阶》引理 \ref{flat-lemma-complex-and-divisor-blowup-pre} 中分别从三元组
$(\mathbf{P}^1_X, (\mathbf{P}^1_X)\infty, p^*E)$ 与
$(\mathbf{P}^1_{X'}, (\mathbf{P}^1_{X'})\infty, (p')^*E')$
构造的态射。此外
$Q = L\eta_{\mathcal{I}_\infty}p^*E$ 且
$Q = L\eta_{\mathcal{I}'_\infty}(p')^*E'$，其中
$\mathcal{I}_\infty \subset \mathcal{O}_W$ 是 $W_\infty$ 的理想层，
$\mathcal{I}'_\infty \subset \mathcal{O}_{W'}$ 是
$W'_\infty$ 的理想层。其次，
$h : \mathbf{P}^1_{X'} \to \mathbf{P}^1_X$ 是概形态射，使得有效
Cartier 除子 $(\mathbf{P}^1_X)_\infty$ 的拉回是有效 Cartier 除子
$(\mathbf{P}^1_{X'})_\infty$，且 $h^*p^*E = (p')^*E'$。
由《平坦性进阶》引理 \ref{flat-lemma-complex-and-divisor-blowup-base-change}，得到交换图
$$
\xymatrix{
W' \ar[rd]_{b'} \ar[r]_-g &
\mathbf{P}^1_{X'} \times_{\mathbf{P}^1_X} W \ar[d]_r \ar[r]_-q &
W \ar[d]^b \\
&
\mathbf{P}^1_{X'} \ar[r] &
\mathbf{P}^1_X
}
$$
其中 $W'$ 是 $\mathbf{P}^1_{X'}$ 关于 $b$ 的“严格变换”，
且 $Q' = (q \circ g)^*Q$。回忆
$P_p(Z \to X, E) = P'_p(Q)$，分别地\ 
$c_p(Z \to X, E) = c'_p(Q)$，其中 $P'_p(Q)$，分别地\ $c'_p(Q)$，
是在引理 \ref{chow-lemma-localized-chern-pre} 中使用 $b, Q, T'$ 构造的；
这里 $T'$ 是闭子概形 $T' \subset W_\infty$，满足以下两个性质：
(a) $T'$ 包含 $W_\infty$ 中所有位于 $X \setminus Z$ 上方的点；
(b) $Q|_{T'}$ 为零，分别地\ 同构于置于次数 $0$ 的、秩
$< p$ 的有限局部自由模。在引理 \ref{chow-lemma-localized-chern-pre} 的构造中，我们选定了一个具有
(a)、(b) 性质的闭子概形 $T'$，但 $T'$ 的具体选择无关紧要；见
引理 \ref{chow-lemma-localized-chern-pre-independent}。

\medskip\noindent
其次，由引理 \ref{chow-lemma-base-change-localized-chern-pre}，
双变类 $P_p(Z \to X, E) = P'_p(Q)$，分别地\ 
$c_p(Z \to X, E) = c_p(Q')$，在 $X'$ 上的限制对应于类
$P'_p(q^*Q)$，分别地\ $c'_p(q^*Q)$；后者是在
引理 \ref{chow-lemma-localized-chern-pre} 中使用
$r : \mathbf{P}^1_{X'} \times_{\mathbf{P}^1_X} W \to \mathbf{P}^1_{X'}$、
复形 $q^*Q$ 及逆像 $q^{-1}(T')$ 构造的。

\medskip\noindent
现在由引理 \ref{chow-lemma-localized-chern-pre-independent}
的第二个陈述，$P'_p(Q') = P'_p(q^*Q)$，分别地\ 
$c'_p(q^*Q) = c'_p(Q')$。又因
$P_p(Z' \to X', E') = P'_p(Q')$，分别地\ 
$c_p(Z' \to X', E') = c'_p(Q')$，引理得证。
\end{proof}

\begin{remark}
\label{chow-remark-loc-chern}
在情形 \ref{chow-situation-loc-chern} 中，更自然的做法是把假设 (3)
替换为假设“$E$ 的陈类已定义”。事实上，把引理 \ref{chow-lemma-independent-loc-chern} 和
\ref{chow-lemma-base-change-loc-chern} 与
引理 \ref{chow-lemma-envelope-bivariant} 结合起来，很容易把定义推广到
这个略微更一般的情形。若以后需要，我们将在此完成该推广。
\end{remark}

\begin{lemma}
\label{chow-lemma-loc-chern-after-pushforward}
在情形 \ref{chow-situation-loc-chern} 中，对任意
$\alpha \in \CH_*(Z)$，在 $\CH_*(Z)$ 中有
$$
P_p(Z \to X, E) \cap i_*\alpha = P_p(E|_Z) \cap \alpha,
\quad\text{分别}\quad
c_p(Z \to X, E) \cap i_*\alpha = c_p(E|_Z) \cap \alpha
$$
\end{lemma}

\begin{proof}
只证明第二个等式，第一个的证明从略。由于
$c_p(Z \to X, E)$ 是双变类，且 $Z \to X$ 沿 $Z \to X$
的基变换是 $\text{id} : Z \to Z$，有
$c_p(Z \to X, E) \cap i_*\alpha = c_p(Z \to X, E) \cap \alpha$。
由引理 \ref{chow-lemma-base-change-loc-chern}，
$c_p(Z \to X, E)$ 在 $Z$ 上的限制（！）是
$\text{id} : Z \to Z$ 与 $E|_Z$ 的局部化陈类。
因此取 $X = Z$，结论由
(\ref{chow-equation-defining-property-localized-classes}) 得出。
\end{proof}

\begin{lemma}
\label{chow-lemma-loc-chern-disjoint}
在情形 \ref{chow-situation-loc-chern} 中，若
$\alpha \in \CH_k(X)$ 的支撑与 $Z$ 不交，则
$P_p(Z \to X, E) \cap \alpha = 0$，分别地\ 
$c_p(Z \to X, E) \cap \alpha = 0$。
\end{lemma}

\begin{proof}
这立即由局部化陈类的构造得出。也可先把
$c_p(Z \to X, E)$ 限制到 $\alpha$ 的支撑，再用
引理 \ref{chow-lemma-base-change-loc-chern} 看出该限制为零，
从而计算 $c_p(Z \to X, E) \cap \alpha$。
\end{proof}

\begin{lemma}
\label{chow-lemma-loc-chern-shrink-Z}
在情形 \ref{chow-situation-loc-chern} 中，假设
$Z \subset Z' \subset X$，其中 $Z'$ 是 $X$ 的闭子概形。则
$P_p(Z' \to X, E) = (Z \to Z')_* \circ P_p(Z \to X, E)$，
分别地\ $c_p(Z' \to X, E) = (Z \to Z')_* \circ c_p(Z \to X, E)$
（其中 $\circ$ 如引理 \ref{chow-lemma-push-proper-bivariant}）。
\end{lemma}

\begin{proof}
引理 \ref{chow-lemma-independent-loc-chern} 中
$P_p(Z' \to X, E)$，分别地\ $c_p(Z' \to X, E)$ 的构造，
使用完全相同的态射 $b : W \to \mathbf{P}^1_X$ 与
$D(\mathcal{O}_W)$ 中的完美对象 $Q$。于是应用
引理 \ref{chow-lemma-silly-shrink} 即得结论。省略一些细节。
\end{proof}

\begin{lemma}
\label{chow-lemma-loc-chern-agree}
在引理 \ref{chow-lemma-silly} 中，设 $E_2$ 是某个完美对象
$E \in D(\mathcal{O}_X)$ 的限制，且该对象在 $X_1$ 上的限制为零，
分别地\ 同构于置于上同调次数 $0$ 的、秩 $< p$ 的有限局部自由
$\mathcal{O}_{X_1}$-模。则引理 \ref{chow-lemma-silly} 的类
$P'_p(E_2)$，分别地\ $c'_p(E_2)$，与定义 \ref{chow-definition-localized-chern} 的
$P_p(X_2 \to X, E)$，分别地\ $c_p(X_2 \to X, E)$ 相同，
只要 $E$ 满足情形 \ref{chow-situation-loc-chern} 的假设 (3)。
\end{lemma}

\begin{proof}
关于 $E$ 的假设蕴含存在包含 $X_1$ 的开集 $U \subset X$，
使得 $E|_U$ 为零，分别地\ 同构于秩 $< p$ 的有限局部自由
$\mathcal{O}_U$-模。见《代数进阶》引理 \ref{more-algebra-lemma-lift-perfect-from-residue-field}。
令 $Z \subset X$ 为 $U$ 在 $X$ 中的补集，并赋予诱导的既约闭子概形结构。
则由引理 \ref{chow-lemma-loc-chern-shrink-Z}，
$P_p(X_2 \to X, E) = (Z \to X_2)_* \circ P_p(Z \to X, E)$，
分别地\ $c_p(X_2 \to X, E) = (Z \to X_2)_* \circ c_p(Z \to X, E)$。
现在可证明 $P_p(X_2 \to X, E)$，分别地\ $c_p(X_2 \to X, E)$，
满足引理 \ref{chow-lemma-silly} 给出的对 $P'_p(E_2)$，分别地\ 
$c'_p(E_2)$ 的刻画。事实上，由刚刚证明的关系
$P_p(X_2 \to X, E) = (Z \to X_2)_* \circ P_p(Z \to X, E)$，
分别地\ $c_p(X_2 \to X, E) = (Z \to X_2)_* \circ c_p(Z \to X, E)$，
以及 $X_1 \cap Z = \emptyset$，由
引理 \ref{chow-lemma-loc-chern-disjoint}，复合
$P_p(X_2 \to X, E) \circ i_{1, *}$，分别地\ 
$c_p(X_2 \to X, E) \circ i_{1, *}$ 为零。另一方面，由
引理 \ref{chow-lemma-loc-chern-after-pushforward}，
$P_p(X_2 \to X, E) \circ i_{2, *} = P_p(E_2)$，
分别地\ $c_p(X_2 \to X, E) \circ i_{2, *} = c_p(E_2)$。
\end{proof}





\section{两个技术引理}
\label{chow-section-tools-loc-chern}

\noindent
本节建立一些附加工具，以便更方便地处理局部化陈类。
下一个引理是我们在引理 \ref{chow-lemma-localized-chern-pre-compose}
和 \ref{chow-lemma-localized-chern-pre-sum-c} 的证明中已经遇到的
某个事实的更精确版本。

\begin{lemma}
\label{chow-lemma-homomorphism-final}
设 $(S, \delta)$ 如情形 \ref{chow-situation-setup}。
设 $X$ 在 $S$ 上局部有限型，并设
$b : W \longrightarrow \mathbf{P}^1_X$ 是真概形态射。
设 $n \geq 1$。对 $i = 1, \ldots, n$，设
$Z_i \subset X$ 是闭子概形，$Q_i \in D(\mathcal{O}_W)$
是完美对象，$p_i \geq 0$ 是整数，并设
$T_i \subset W_\infty$（$i = 1, \ldots, n$）是闭子概形。
记 $W_i = b^{-1}(\mathbf{P}^1_{Z_i})$。假设
\begin{enumerate}
\item 对 $i = 1, \ldots, n$，数据 $b, Z_i, Q_i, T_i, p_i$
满足引理 \ref{chow-lemma-localized-chern-pre} 的假设；
\item $Q_i|_{W \setminus W_i}$ 为零，分别地\ 同构于置于上同调次数
$0$ 的、秩 $< p_i$ 的有限局部自由模；
\item $W$ 上的 $Q_i$ 满足情形 \ref{chow-situation-loc-chern}
的假设 (3)。
\end{enumerate}
则 $P'_{p_n}(Q_n) \circ \ldots \circ P'_{p_1}(Q_1)$ 等于
$$
(W_{n, \infty} \cap \ldots \cap W_{1, \infty} \to
Z_n \cap \ldots \cap Z_1)_* \circ
P'_{p_n}(Q_n|_{W_{n, \infty}}) \circ \ldots \circ P'_{p_1}(Q_1|_{W_{1, \infty}})
\circ C
$$
这是 $A^{p_n + \ldots + p_1}(Z_n \cap \ldots \cap Z_1 \to X)$
中的等式；分别地\ $c'_{p_n}(Q_n) \circ \ldots \circ c'_{p_1}(Q_1)$ 等于
$$
(W_{n, \infty} \cap \ldots \cap W_{1, \infty} \to
Z_n \cap \ldots \cap Z_1)_* \circ
c'_{p_n}(Q_n|_{W_{n, \infty}}) \circ \ldots \circ c'_{p_1}(Q_1|_{W_{1, \infty}})
\circ C
$$
这是 $A^{p_n + \ldots + p_1}(Z_n \cap \ldots \cap Z_1 \to X)$
中的等式。
\end{lemma}

\begin{proof}
对 $n$ 归纳证明关于陈类的陈述；关于 $P_p(-)$ 的陈述完全同理证明。
当 $n = 1$ 时，陈述就是 $c'_{p_1}(Q_1)$ 的构造，因为
$W_{1, \infty}$ 是 $Z_1$ 在 $W_\infty$ 中的逆像。
当 $n > 1$ 时，由归纳假设，
$c'_{p_n}(Q_n) \circ \ldots \circ c'_{p_1}(Q_1)$ 等于
$$
c'_{p_n}(Q_n) \circ
(W_{n - 1, \infty} \cap \ldots \cap W_{1, \infty} \to
Z_{n - 1} \cap \ldots \cap Z_1)_* \circ
c'_{p_{n - 1}}(Q_{n - 1}|_{W_{n - 1}, \infty}) \circ \ldots \circ
c'_{p_1}(Q_1|_{W_{1, \infty}})
\circ C
$$
由引理 \ref{chow-lemma-base-change-localized-chern-pre}，
$c'_{p_n}(Q_n)$ 在 $Z_{n - 1} \cap \ldots \cap Z_1$ 上的限制，
由闭子集 $Z_n \cap \ldots \cap Z_1$、态射
$b' : W_{n - 1} \cap \ldots \cap W_1 \to
\mathbf{P}^1_{Z_{n - 1} \cap \ldots \cap Z_1}$，以及
$Q_n$ 在 $W_{n - 1} \cap \ldots \cap W_1$ 上的限制计算。
注意 $(b')^{-1}(Z_n) = W_n \cap \ldots \cap W_1$，并且
$(W_n \cap \ldots \cap W_1)_\infty =
W_{n, \infty} \cap \ldots \cap W_{1, \infty}$。
记 $C_{n - 1} \in A^0(W_{n - 1, \infty} \cap \ldots \cap W_{1, \infty} \to
Z_{n - 1} \cap \ldots \cap Z_1)$ 为引理 \ref{chow-lemma-gysin-at-infty}
的类。于是 $c'_{p_n}(Q_n)$ 在
$Z_{n - 1} \cap \ldots \cap Z_1$ 上的限制为
\begin{align*}
&
(W_{n, \infty} \cap \ldots \cap W_{1, \infty} \to Z_n \cap \ldots \cap Z_1)_*
\circ
c'_{p_n}(Q_n|_{(W_n \cap \ldots \cap W_1)_\infty})
\circ
C_{n - 1} \\
& =
(W_{n, \infty} \cap \ldots \cap W_{1, \infty} \to Z_n \cap \ldots \cap Z_1)_*
\circ
c'_{p_n}(Q_n|_{W_{n, \infty}})
\circ
C_{n - 1}
\end{align*}
其中等式由引理 \ref{chow-lemma-base-change-silly} 得出
（右边的限制记号从略）。因此上式化为
\begin{align*}
(W_{n, \infty} \cap \ldots \cap W_{1, \infty} \to
Z_n \cap \ldots \cap Z_1)_* \circ
c'_{p_n}(Q_n|_{W_n, \infty}) \circ \\
C_{n - 1} \circ
(W_{n - 1, \infty} \cap \ldots \cap W_{1, \infty} \to
Z_{n - 1} \cap \ldots \cap Z_1)_* \\
\circ
c'_{p_{n - 1}}(Q_{n - 1}|_{W_{n - 1}, \infty}) \circ \ldots \circ
c'_{p_1}(Q_1|_{W_{1, \infty}})
\circ C
\end{align*}
由引理 \ref{chow-lemma-homomorphism-pre}，复合
$C_{n - 1} \circ (W_{n - 1, \infty} \cap \ldots \cap W_{1, \infty} \to
Z_{n - 1} \cap \ldots \cap Z_1)_*$
在 Gysin 映射
$$
(W_{n - 1, \infty} \cap \ldots \cap W_{1, \infty} \to
W_{n - 1} \cap \ldots \cap W_1)^*
$$
的像中的元素上是恒等映射。因此只须证明
$c'_{p_{n - 1}}(Q_{n - 1}|_{W_{n - 1}, \infty}) \circ \ldots \circ
c'_{p_1}(Q_1|_{W_{1, \infty}}) \circ C$
的像中每个元素都在该 Gysin 映射的像中。由
引理 \ref{chow-lemma-loc-chern-agree} 以及引理陈述中关于 $Q_i$
的假设 (2)、(3)，可以写成
$$
c'_{p_i}(Q_i|_{W_{i, \infty}}) = c_{p_i}(W_i \to W, Q_i)\text{ 在 } W_{i, \infty}\text{ 上的限制}
$$
因此，若 $\beta \in \CH_{k + 1}(W)$ 在
$b^{-1}(\mathbf{A}^1_X)$ 上限制为 $\alpha$ 的平坦拉回，则
\begin{align*}
& c'_{p_{n - 1}}(Q_{n - 1}|_{W_{n - 1}, \infty}) \cap \ldots \cap
c'_{p_1}(Q_1|_{W_{1, \infty}})
\cap C \cap \alpha \\
& =
c'_{p_{n - 1}}(Q_{n - 1}|_{W_{n - 1}, \infty}) \cap \ldots \cap
c'_{p_1}(Q_1|_{W_{1, \infty}})
\cap i_\infty^* \beta \\
& =
c_{p_{n - 1}}(W_{n - 1} \to W, Q_{n - 1}) \cap \ldots \cap
c_{p_{n - 1}}(W_1 \to W, Q_1) \cap i_\infty^* \beta \\
& =
(W_{n - 1, \infty} \cap \ldots \cap W_{1, \infty} \to
W_{n - 1} \cap \ldots \cap W_1)^*
\left(c_{p_{n - 1}}(W_{n - 1} \to W, Q_{n - 1}) \cap \ldots \cap
c_{p_1}(W_1 \to W, Q_1) \cap \beta\right)
\end{align*}
如所需。最后一个等式使用了
$c_{p_i}(W_i \to W, Q_i)$ 是双变类，因而按定义与
$i_\infty^*$ 交换。
\end{proof}

\noindent
若要计算复形的局部化陈类，下面的引理提供了极大的灵活性。

\begin{lemma}
\label{chow-lemma-independent-loc-chern-bQ}
假设 $(S, \delta), X, Z, b : W \to \mathbf{P}^1_X, Q, T, p$
满足引理 \ref{chow-lemma-localized-chern-pre} 的假设。
设 $F \in D(\mathcal{O}_X)$ 是完美对象，使得
\begin{enumerate}
\item $Q$ 在 $b^{-1}(\mathbf{A}^1_X)$ 上的限制同构于 $F$ 的拉回；
\item $F|_{X \setminus Z}$ 为零，分别地\ 同构于置于上同调次数
$0$ 的、秩 $< p$ 的有限局部自由
$\mathcal{O}_{X \setminus Z}$-模；并且
\item $W$ 上的 $Q$ 与 $X$ 上的 $F$ 满足
情形 \ref{chow-situation-loc-chern} 的假设 (3)。
\end{enumerate}
则引理 \ref{chow-lemma-localized-chern-pre} 所构造的
$A^p(Z \to X)$ 中的类 $P'_p(Q)$，分别地\ $c'_p(Q)$，
等于定义 \ref{chow-definition-localized-chern} 中的
$P_p(Z \to X, F)$，分别地\ $c_p(Z \to X, F)$。
\end{lemma}

\begin{proof}
这些假设在沿局部有限型态射 $X' \to X$ 作基变换后仍然成立。
因此只须对任意 $\alpha \in \CH_k(X)$ 证明
$P_p(Z \to X, F) \cap \alpha = P'_p(Q) \cap \alpha$，
分别地\ $c_p(Z \to X, F) \cap \alpha = c'_p(Q) \cap \alpha$。
依照引理 \ref{chow-lemma-gysin-at-infty} 中 $C$ 的构造，
选择 $\beta \in \CH_{k + 1}(W)$，使其在
$b^{-1}(\mathbf{A}^1_X)$ 上的限制等于 $\alpha$ 的平坦拉回。
记 $W' = b^{-1}(Z)$，并记
$E = W'_\infty \subset W_\infty$ 为 $Z$ 经
$W_\infty \to X$ 的逆像。引理由以下等式链推出
（$P_p$ 的情形类似）：
\begin{align*}
c'_p(Q) \cap \alpha
& =
(E \to Z)_*(c'_p(Q|_E) \cap i_\infty^*\beta) \\
& =
(E \to Z)_*(c_p(E \to W_\infty, Q|_{W_\infty}) \cap i_\infty^*\beta) \\
& =
(W'_\infty \to Z)_*(c_p(W' \to W, Q) \cap i_\infty^*\beta) \\
& =
(W'_\infty \to Z)_*((i'_\infty)^*(c_p(W' \to W, Q) \cap \beta)) \\
& =
(W'_\infty \to Z)_*((i'_\infty)^*(c_p(Z' \to X, F) \cap \beta)) \\
& =
(W'_0 \to Z)_*((i'_0)^*(c_p(Z' \to X, F) \cap \beta)) \\
& =
(W'_0 \to Z)_*(c_p(Z' \to X, F) \cap i_0^*\beta)) \\
& =
c_p(Z \to X, F) \cap \alpha
\end{align*}
第一个等式是引理 \ref{chow-lemma-localized-chern-pre} 中
$c'_p(Q)$ 的构造。第二个等式是引理 \ref{chow-lemma-loc-chern-agree}。
$W' \to W$ 沿 $W_\infty \to W$ 的基变换是态射
$E = W'_\infty \to W_\infty$，所以第三个等式由
引理 \ref{chow-lemma-base-change-loc-chern} 得出。第四个等式中
$i'_\infty : W'_\infty \to W'$ 是包含态射；该等式由
$c_p(W' \to W, Q)$ 是双变类这一事实得出。对于第五个等式，
注意由引理的假设 (1)，$Q$ 与 $F$ 在
$D(\mathcal{O}_{b^{-1}(\mathbf{A}^1_X)})$ 中限制为同一个对象；
应用引理 \ref{chow-lemma-base-change-loc-chern}，
$c_p(W' \to W, Q)$ 与 $c_p(Z' \to X, F)$ 在
$A^p((b')^{-1} \to b^{-1}(\mathbf{A}^1_X))$ 中限制为同一个双变类。
由于 $(i'_\infty)^*$ 消去支撑在 $W'_\infty$ 上的循环
（见注 \ref{chow-remark-gysin-on-cycles}），第五个等式成立。
第六个等式成立，因为 $W'_\infty$ 与 $W'_0$ 是
$\mathbf{P}^1_Z$ 中有理等价的有效 Cartier 除子
$D_0, D_\infty$ 的拉回，所以 $i_\infty^*\beta$ 与
$i_0^*\beta$ 映到 $W'$ 上同一个循环类；确切地说，由
引理 \ref{chow-lemma-support-cap-effective-Cartier}，二者均表示类
$c_1(\mathcal{O}_{\mathbf{P}^1_Z}(1)) \cap c_p(Z \to X, F_) \cap \beta$。
第七个等式成立，因为 $c_p(Z \to X, F)$ 是双变类。
按构造 $W'_0 = Z$ 且 $i_0^*\beta = \alpha$，这说明最后一个等式成立。
\end{proof}







\section{局部化陈类的性质}
\label{chow-section-properties-loc-chern}

\noindent
本节的主要结果是局部化陈类的可加性和乘法性。

\begin{lemma}
\label{chow-lemma-loc-chern-character}
在情形 \ref{chow-situation-loc-chern} 中，假设
$E|_{X \setminus Z}$ 为零。则
\begin{align*}
P_1(Z \to X, E) & = c_1(Z \to X, E), \\
P_2(Z \to X, E) & = c_1(Z \to X, E)^2 - 2c_2(Z \to X, E), \\
P_3(Z \to X, E) & = c_1(Z \to X, E)^3 - 3c_1(Z \to X, E)c_2(Z \to X, E)
+ 3c_3(Z \to X, E),
\end{align*}
依此类推，其中乘积取自注 \ref{chow-remark-ring-loc-classes}
中的代数 $A^{(1)}(Z \to X)$。
\end{lemma}

\begin{proof}
该陈述是有意义的，因为零层的秩 $< 1$，所以对所有
$p \geq 1$，类 $c_p(Z \to X, E)$ 都有定义。
结论本身立即由更一般的引理 \ref{chow-lemma-localized-chern-pre-compose} 得出，因为局部化陈类
是使用引理 \ref{chow-lemma-localized-chern-pre} 的步骤、在
第 \ref{chow-section-localized-chern} 节中定义的。
\end{proof}

\begin{lemma}
\label{chow-lemma-loc-chern-classes-commute}
在情形 \ref{chow-situation-loc-chern} 中，设
$Y \to X$ 局部有限型，且 $c \in A^*(Y \to X)$。则
$$
P_p(Z \to X, E) \circ c = c \circ P_p(Z \to X, E),
$$
分别地
$$
c_p(Z \to X, E) \circ c = c \circ c_p(Z \to X, E)
$$
在 $A^*(Y \times_X Z \to X)$ 中成立。
\end{lemma}

\begin{proof}
这由引理 \ref{chow-lemma-homomorphism-commute} 得出。
更确切地，令
$$
b : W \to \mathbf{P}^1_X
\quad\text{且}\quad
Q
\quad\text{且}\quad
T' \subset T \subset W_\infty
$$
如引理 \ref{chow-lemma-independent-loc-chern} 的证明中所取。
按定义，作为双变运算，
$c_p(Z \to X, E) = c'_p(Q)$，其中右边是在
引理 \ref{chow-lemma-localized-chern-pre} 中使用
$W, b, Q, T'$ 构造的双变类。由
引理 \ref{chow-lemma-homomorphism-commute}，在
$A^*(Y \times_X Z \to X)$ 中有
$P'_p(Q) \circ c = c \circ P'_p(Q)$，分别地\ 
$c'_p(Q) \circ c = c \circ c'_p(Q)$，结论随之成立。
\end{proof}

\begin{remark}
\label{chow-remark-loc-chern-classes}
在情形 \ref{chow-situation-loc-chern} 中，方便起见定义
$$
c^{(p)}(Z \to X, E) = 1 + c_1(E) + \ldots + c_{p - 1}(E) +
c_p(Z \to X, E) + c_{p + 1}(Z \to X, E) + \ldots
$$
将其视为注 \ref{chow-remark-ring-loc-classes} 中所考虑的代数
$A^{(p)}(Z \to X)$ 的元素。
\end{remark}

\begin{lemma}
\label{chow-lemma-additivity-loc-chern-c}
设 $(S, \delta)$ 如情形 \ref{chow-situation-setup}。
设 $X$ 在 $S$ 上局部有限型，且 $Z \to X$ 是闭浸入。设
$$
E_1 \to E_2 \to E_3 \to E_1[1]
$$
是 $D(\mathcal{O}_X)$ 中完美对象组成的可区别三角。假设
\begin{enumerate}
\item 限制 $E_1|_{X \setminus Z}$ 和 $E_3|_{X \setminus Z}$
分别同构于置于次数 $0$ 的、秩 $< p_1$ 和 $< p_3$ 的有限局部自由
$\mathcal{O}_{X \setminus Z}$-模；并且
\item 下列条件至少有一个成立：
(a) $X$ 拟紧；
(b) $X$ 的不可约分支拟紧；
(c) $E_3 \to E_1[1]$ 可由有限局部自由
$\mathcal{O}_X$-模组成的局部有界复形之间的映射表示；或者
(d) 存在包络 $f : Y \to X$，使得
$Lf^*E_3 \to Lf^*E_1[1]$ 可由有限局部自由
$\mathcal{O}_Y$-模组成的局部有界复形之间的映射表示。
\end{enumerate}
采用注 \ref{chow-remark-loc-chern-classes} 中的记号，在
$A^{(p_1 + p_3)}(Z \to X)$ 中有
$$
c^{(p_1 + p_3)}(Z \to X, E_2) = c^{(p_1)}(Z \to X, E_1)c^{(p_3)}(Z \to X, E_3)
$$
\end{lemma}

\begin{proof}
注意，这些假设蕴含 $E_2|_{X \setminus Z}$ 为零，
分别地\ 同构于秩 $< p_1 + p_3$ 的有限局部自由
$\mathcal{O}_{X \setminus Z}$-模。因此该陈述是有意义的。

\medskip\noindent
设 $f : Y \to X$ 是包络。展开引理中公式的左右两边可见，
需要证明 $A^*(X)$ 和 $A^*(Z \to X)$ 中若干类的等式。
由引理 \ref{chow-lemma-envelope-bivariant} 中的唯一性，只须证明
$A^*(Y)$ 和 $A^*(Z \to Y)$ 中对应的关系。又因所涉及各类的构造
与基变换相容（引理 \ref{chow-lemma-base-change-loc-chern}），
可将 $X$ 替换为 $Y$，并将可区别三角替换为其拉回。

\medskip\noindent
在引理 \ref{chow-lemma-additivity-on-perfect} 的证明中已经看到，
条件 (2)(a)、(2)(b) 和 (2)(c) 都蕴含条件 (2)(d)。
结合上一段的讨论，问题化归为下一段所述的情形。

\medskip\noindent
设 $\varphi^\bullet : \mathcal{E}_3^\bullet[-1] \to \mathcal{E}_1^\bullet$
是有限局部自由 $\mathcal{O}_X$-模组成的局部有界复形之间的映射，
它在导出范畴中表示映射 $E_3[-1] \to E_1$。考虑概形
$X' = \mathbf{A}^1 \times X$ 及投影
$g : X' \to X$。令 $Z' = g^{-1}(Z) = \mathbf{A}^1 \times Z$。
以 $t$ 表示 $\mathbf{A}^1$ 上的坐标。考虑 $X'$ 上复形映射
$$
t g^*\varphi^\bullet :
g^*\mathcal{E}_3^\bullet[-1]
\longrightarrow
g^*\mathcal{E}_1^\bullet
$$
的锥 $\mathcal{C}^\bullet$。得到可区别三角
$$
g^*\mathcal{E}_1^\bullet \to \mathcal{C}^\bullet \to
g^*\mathcal{E}_3^\bullet \to g^*\mathcal{E}_1^\bullet[1]
$$
其中前三项逐项组成分裂短正合复形列。显然
$\mathcal{C}^\bullet$ 是有限局部自由
$\mathcal{O}_{X'}$-模组成的有界复形；它在
$X' \setminus Z'$ 上的限制同构于置于次数 $0$ 的、秩
$< p_1 + p_3$ 的有限局部自由
$\mathcal{O}_{X' \setminus Z'}$-模。因此，对
$p \geq p_1 + p_3$，局部化陈类
$$
c_p(Z' \to X', \mathcal{C}^\bullet) \in A^p(Z' \to X')
$$
已有定义。对任意 $\alpha \in \CH_k(X)$，考虑
$$
c_p(Z' \to X', \mathcal{C}^\bullet) \cap g^*\alpha
\in \CH_{k + 1 - p}(\mathbf{A}^1 \times X)
$$
若限制到 $t = 0$，则映射 $t g^*\varphi^\bullet$ 限制为零，
而 $\mathcal{C}^\bullet|_{t = 0}$ 是
$\mathcal{E}_1^\bullet$ 与 $\mathcal{E}_3^\bullet$ 的直和。
由局部化陈类与基变换的相容性
（引理 \ref{chow-lemma-base-change-loc-chern}），得到
$$
i_0^* \circ c^{(p_1 + p_3)}(Z' \to X', \mathcal{C}^\bullet) \circ g^* =
c^{(p_1 + p_2)}(Z \to X, E_1 \oplus E_3)
$$
作为 $A^{(p_1 + p_3)}(Z \to X)$ 中的等式。另一方面，若限制到
$t = 1$，则映射 $t g^*\varphi^\bullet$ 限制为 $\varphi$，
而 $\mathcal{C}^\bullet|_{t = 1}$ 是表示 $E_2$ 的有限局部自由模
组成的有界复形。因此
$$
i_1^* \circ c^{(p_1 + p_3)}(Z' \to X', \mathcal{C}^\bullet) \circ g^* =
c^{(p_1 + p_2)}(Z \to X, E_2)
$$
在 $A^{(p_1 + p_3)}(Z \to X)$ 中成立。按有理等价的定义，
$i_0^* = i_1^*$（更确切地，这由
引理 \ref{chow-lemma-linebundle-formulae} 中的公式得出），所以
$$
c^{(p_1 + p_2)}(Z \to X, E_2) = c^{(p_1 + p_2)}(Z \to X, E_1 \oplus E_3)
$$
这把问题化归为下一段所讨论的情形。

\medskip\noindent
假设 $E_2 = E_1 \oplus E_3$，并且三元组 $(X, Z, E_i)$
如情形 \ref{chow-situation-loc-chern}。对 $i = 1, 3$，令
$$
b_i : W_i \to \mathbf{P}^1_X
\quad\text{且}\quad
Q_i
\quad\text{且}\quad
T'_i \subset T_i \subset W_{i, \infty}
$$
如引理 \ref{chow-lemma-independent-loc-chern} 的证明中所取。
按定义
$$
c_p(Z \to X, E_i) = c'_p(Q_i)
$$
其中右边是在引理 \ref{chow-lemma-localized-chern-pre} 中使用
$W_i, b_i, Q_i, T'_i$ 构造的双变类。令
$W = W_1 \times_{b_1, \mathbf{P}^1_X, b_2} W_2$，并考虑笛卡儿图
$$
\xymatrix{
W \ar[d]_{g_1} \ar[rd]^b \ar[r]_{g_3} & W_3 \ar[d]^{b_3} \\
W_1 \ar[r]^{b_1} & \mathbf{P}^1_X
}
$$
当然，$b^{-1}(\mathbf{A}^1)$ 同构地映到 $\mathbf{A}^1_X$。
注意，$T' = g_1^{-1}(T'_1) \cap g_2^{-1}(T'_2)$ 仍包含
$W_\infty$ 中所有位于 $X \setminus Z$ 上方的点。
由引理 \ref{chow-lemma-localized-chern-pre-independent}，可以使用
$W$、$b$、$g_i^*\mathcal{Q}_i$ 与
$T'$ 构造 $i = 1, 3$ 时的 $c_p(Z \to X, E_i)$。
此外，由引理 \ref{chow-lemma-independent-loc-chern-bQ}
给出的更强独立性，可以使用
$W$、$b$、$g_1^*Q_1 \oplus g_3^*Q_3$ 与 $T'$
计算各类 $c_p(Z \to X, E_2)$。于是所需等式由
引理 \ref{chow-lemma-localized-chern-pre-sum-c} 得出。
\end{proof}

\begin{lemma}
\label{chow-lemma-additivity-loc-chern-P}
设 $(S, \delta)$ 如情形 \ref{chow-situation-setup}。
设 $X$ 在 $S$ 上局部有限型，且 $Z \to X$ 是闭浸入。设
$$
E_1 \to E_2 \to E_3 \to E_1[1]
$$
是 $D(\mathcal{O}_X)$ 中完美对象组成的可区别三角。假设
\begin{enumerate}
\item 限制 $E_1|_{X \setminus Z}$ 和 $E_3|_{X \setminus Z}$
为零；并且
\item 下列条件至少有一个成立：
(a) $X$ 拟紧；
(b) $X$ 的不可约分支拟紧；
(c) $E_3 \to E_1[1]$ 可由有限局部自由
$\mathcal{O}_X$-模组成的局部有界复形之间的映射表示；或者
(d) 存在包络 $f : Y \to X$，使得
$Lf^*E_3 \to Lf^*E_1[1]$ 可由有限局部自由
$\mathcal{O}_Y$-模组成的局部有界复形之间的映射表示。
\end{enumerate}
则对所有 $p \in \mathbf{Z}$ 都有
$$
P_p(Z \to X, E_2) = P_p(Z \to X, E_1) + P_p(Z \to X, E_3)
$$
从而
$ch(Z \to X, E_2) = ch(Z \to X, E_1) + ch(Z \to X, E_3)$。
\end{lemma}

\begin{proof}
证明与引理 \ref{chow-lemma-additivity-loc-chern-c} 的证明完全相同，
只是最后使用引理 \ref{chow-lemma-localized-chern-pre-sum-P}。
当 $p > 0$ 时，可令 $p_1 = p_3 = 1$，由
引理 \ref{chow-lemma-additivity-loc-chern-c} 以及
引理 \ref{chow-lemma-loc-chern-character} 中
$P_p(Z \to X, E)$ 与 $c_p(Z \to X, E)$ 的关系推出本引理。
$p = 0$ 的情形可直接证明（只有当 $X$ 有某个连通分支完全包含在
$Z$ 中时才有实际内容）。
\end{proof}

\begin{lemma}
\label{chow-lemma-loc-chern-tensor-product}
在情形 \ref{chow-situation-setup} 中，设 $X$ 在 $S$ 上局部有限型。
设 $Z_i \subset X$（$i = 1, 2$）是闭子概形，且
$F_i$（$i = 1, 2$）是 $D(\mathcal{O}_X)$ 中的完美对象。
假设对 $i = 1, 2$，限制 $F_i|_{X \setminus Z_i}$
为零\footnote{想必本引理还有一个变体，其中只假设
$F_i|_{X \setminus Z_i}$ 同构于秩 $< p_i$ 的有限局部自由
$\mathcal{O}_{X \setminus Z_i}$-模。}，并且 $X$ 上的 $F_i$
满足情形 \ref{chow-situation-loc-chern} 的假设 (3)。记
$r_i = P_0(Z_i \to X, F_i) \in A^0(Z_i \to X)$。则在
$A^1(Z_1 \cap Z_2 \to X)$ 中有
$$
c_1(Z_1 \cap Z_2 \to X, F_1 \otimes_{\mathcal{O}_X}^\mathbf{L} F_2) =
r_1 c_1(Z_2 \to X, F_2) + r_2 c_1(Z_1 \to X, F_1)
$$
并且在 $A^2(Z_1 \cap Z_2 \to X)$ 中有
\begin{align*}
c_2(Z_1 \cap Z_2 \to X, F_1 \otimes_{\mathcal{O}_X}^\mathbf{L} F_2)
& =
r_1 c_2(Z_2 \to X, F_2) +
r_2 c_2(Z_1 \to X, F_1) + \\
& {r_1 \choose 2} c_1(Z_2 \to X, F_2)^2 + \\
& (r_1r_2 - 1) c_1(Z_2 \to X, F_2)c_1(Z_1 \to X, F_1) + \\
& {r_2 \choose 2} c_1(Z_1 \to X, F_1)^2
\end{align*}
高阶陈类依此类推。类似地，在
$\prod_{p \geq 0} A^p(Z_1 \cap Z_2 \to X) \otimes \mathbf{Q}$
中有
$$
ch(Z_1 \cap Z_2 \to X, F_1 \otimes_{\mathcal{O}_X}^\mathbf{L} F_2) =
ch(Z_1 \to X, F_1) ch(Z_2 \to X, F_2)
$$
更确切地，在 $A^p(Z_1 \cap Z_2 \to X)$ 中有
$$
P_p(Z_1 \cap Z_2 \to X, F_1 \otimes_{\mathcal{O}_X}^\mathbf{L} F_2) =
\sum\nolimits_{p_1 + p_2 = p}
{p \choose p_1} P_{p_1}(Z_1 \to X, F_1) P_{p_2}(Z_2 \to X, F_2)
$$
\end{lemma}

\begin{proof}
选取真态射 $b_i : W_i \to \mathbf{P}^1_X$、对象
$Q_i \in D(\mathcal{O}_{W_i})$ 以及闭子概形
$T_i \subset W_{i, \infty}$，如 $F_i$ 的局部化陈类构造中所取，
或更一般地如引理 \ref{chow-lemma-independent-loc-chern-bQ} 所取。
选取交换图
$$
\xymatrix{
W \ar[d]^{g_1} \ar[rd]^b \ar[r]_{g_2} & W_2 \ar[d]^{b_2} \\
W_1 \ar[r]^{b_1} & \mathbf{P}^1_X
}
$$
其中所有态射都是真态射，并且在 $\mathbf{A}^1_X$ 上为同构。
例如，可取 $W$ 为 $b_1^{-1}(\mathbf{A}^1_X)$ 与
$b_2^{-1}(\mathbf{A}^1_X)$ 之间同构之图的闭包。
由引理 \ref{chow-lemma-independent-loc-chern-bQ}，可以使用
$W$、$b = b_i \circ g_i$、$Lg_i^*Q_i$ 和
$g_i^{-1}(T_i)$ 构造局部化陈类
$c_p(Z_i \to X, F_i)$。于是化归为下一段所述情形。

\medskip\noindent
假设已有
\begin{enumerate}
\item 真态射 $b : W \to \mathbf{P}^1_X$，它在
$\mathbf{A}^1_X$ 上为同构；
\item $E_i \subset W_\infty$ 是 $Z_i$ 的逆像；
\item 陈类已有定义的完美对象 $Q_i \in D(\mathcal{O}_W)$，且满足
\begin{enumerate}
\item $Q_i$ 在 $b^{-1}(\mathbf{A}^1_X)$ 上的限制是
$F_i$ 的拉回；并且
\item 存在闭子概形 $T_i \subset W_\infty$，它包含
$W_\infty$ 中所有位于 $X \setminus Z_i$ 上方的点，并使
$Q_i|_{T_i}$ 为零。
\end{enumerate}
\end{enumerate}
由引理 \ref{chow-lemma-independent-loc-chern-bQ}，对 $i = 1, 2$ 有
$$
c_p(Z_i \to X, F_i) = c'_p(Q_i) =
(E_i \to Z_i)_* \circ c'_p(Q_i|_{E_i}) \circ C
$$
以及
$$
P_p(Z_i \to X, F_i) = P'_p(Q_i) =
(E_i \to Z_i)_* \circ P'_p(Q_i|_{E_i}) \circ C
$$
其次，注意 $Q = Q_1 \otimes_{\mathcal{O}_W}^\mathbf{L} Q_2$
对 $F_1 \otimes_{\mathcal{O}_X}^\mathbf{L} F_2$ 和
$T_1 \cup T_2$ 满足 (3)(a) 与 (3)(b)。因此由同一个引理可知
$$
c_p(Z_1 \cap Z_2 \to X, F_1 \otimes_{\mathcal{O}_X}^\mathbf{L} F_2) =
(E_1 \cap E_2 \to Z_1 \cap Z_2)_* \circ
c'_p(Q|_{E_1 \cap E_2}) \circ C
$$
以及
$$
P_p(Z_1 \cap Z_2 \to X, F_1 \otimes_{\mathcal{O}_X}^\mathbf{L} F_2) =
(E_1 \cap E_2 \to Z_1 \cap Z_2)_* \circ
P'_p(Q|_{E_1 \cap E_2}) \circ C
$$
由引理 \ref{chow-lemma-silly-tensor-product}，各类
$c'_p(Q|_{E_1 \cap E_2})$ 与 $P'_p(Q|_{E_1 \cap E_2})$
可用各类 $c'_p(Q_i|_{E_i})$ 与 $P'_p(Q_i|_{E_i})$
按正确方式展开。最后，引理 \ref{chow-lemma-homomorphism-final}
说明 $c'_p(Q_i|_{E_i})$ 与 $P'_p(Q_i|_{E_i})$ 的多项式，
同所需的 $c'_p(Q_i)$ 与 $P'_p(Q_i)$ 的对应多项式一致。
\end{proof}






\section{无穷远处的爆破}
\label{chow-section-blowup-Z-first}

\noindent
设 $X$ 是概形，$Z \subset X$ 是由有限型拟凝聚理想层截出的闭子概形。
记 $X' \to X$ 为以 $Z$ 为中心的爆破。设
$b : W \to \mathbf{P}^1_X$ 为以 $\infty(Z)$ 为中心的爆破，
并记 $E \subset W$ 为例外除子。存在交换图
$$
\xymatrix{
X' \ar[r] \ar[d] & W \ar[d]^b \\
X \ar[r]^\infty & \mathbf{P}^1_X
}
$$
其中水平箭头均为闭浸入（《除子》引理 \ref{divisors-lemma-strict-transform}）。记 $E \subset W$
为例外除子，$W_\infty \subset W$ 为
$(\mathbf{P}^1_X)_\infty$ 的逆像。则有
\begin{enumerate}
\item $b$ 在 $\mathbf{A}^1_X \cup \mathbf{P}^1_{X \setminus Z}$
上为同构；
\item $X'$ 是 $W$ 上的有效 Cartier 除子；
\item $X' \cap E$ 是 $X' \to X$ 的例外除子；
\item 作为 $W$ 上的有效 Cartier 除子，
$W_\infty = X' + E$；
\item $E = \underline{\text{Proj}}_Z(\mathcal{C}_{Z/X, *}[S])$，
其中 $S$ 是置于次数 $1$ 的变量；
\item $X' \cap E = \underline{\text{Proj}}_Z(\mathcal{C}_{Z/X, *})$；
\item
\label{chow-item-cone-is-open}
$E \setminus X' = E \setminus (X' \cap E) =
\underline{\Spec}_Z(\mathcal{C}_{Z/X, *}) = C_ZX$；
\item
\label{chow-item-find-Z-in-blowup}
存在闭浸入 $\mathbf{P}^1_Z \to W$，它与 $b$ 的复合是包含态射
$\mathbf{P}^1_Z \to \mathbf{P}^1_X$，而沿 $\infty$
的基变换是复合 $Z \to C_ZX \to E \to W_\infty$；
其中第一个箭头是锥顶。
\end{enumerate}
回忆 $\mathcal{C}_{Z/X, *}$ 是 $Z$ 在 $X$ 中的余法代数，见
《除子》定义 \ref{divisors-definition-conormal-sheaf}；
而 $C_ZX$ 是 $Z$ 在 $X$ 中的法锥，见
《除子》定义 \ref{divisors-definition-normal-cone}。

\medskip\noindent
现在证明上述各项。我们强烈建议读者亲自推演若干例子，而不是阅读
这些证明。

\medskip\noindent
第 (1) 项由《除子》引理 \ref{divisors-lemma-blowing-up-gives-effective-Cartier-divisor}
中的相应陈述得出。由同一引理，$E \subset W$ 是有效 Cartier 除子。

\medskip\noindent
由《除子》引理 \ref{divisors-lemma-blow-up-pullback-effective-Cartier}，
$W_\infty$ 是有效 Cartier 除子。由于 $E \subset W_\infty$，
可对某个有效 Cartier 除子 $D$ 写成 $W_\infty = D + E$；见
《除子》引理 \ref{divisors-lemma-difference-effective-Cartier-divisors}。
下文将看到 $D = X'$，这便证明 (2) 和 (4)。

\medskip\noindent
由于 $X'$ 是闭浸入 $\infty : X \to \mathbf{P}^1_X$
的严格变换（见上文），$X' \to X$ 的例外除子等于交
$X' \cap E$（例如因为二者都由 $Z$ 的理想层在 $X'$ 上的拉回截出）。
这证明了 (3)。

\medskip\noindent
$\infty(Z)$ 与 $\mathbf{P}^1_Z$ 的交是有效 Cartier 除子
$(\mathbf{P}^1_Z)_\infty$，所以 $\mathbf{P}^1_Z$ 经爆破 $b$
的严格变换同构地映到 $\mathbf{P}^1_Z$；见《除子》引理 \ref{divisors-lemma-strict-transform} 和
\ref{divisors-lemma-blow-up-effective-Cartier-divisor}。
这给出 (8) 中提到的态射 $\mathbf{P}^1_Z \to W$。
由于 $b$ 分离，该态射是闭浸入；见《概形》引理 \ref{schemes-lemma-section-immersion}。

\medskip\noindent
假设 $\Spec(A) \subset X$ 是仿射开集，并且
$Z \cap \Spec(A)$ 对应于有限生成理想 $I \subset A$。
$\infty(Z \cap \Spec(A))$ 的一个仿射邻域是以
$s = T_0/T_1$ 为坐标的 $A$ 上仿射空间。记
$J = (I, s) \subset A[s]$ 为由 $I$ 和 $s$ 生成的理想。
令 $B = A[s] \oplus J \oplus J^2 \oplus \ldots$ 为
$(A[s], J)$ 的 Rees 代数。注意，对所有 $n \geq 0$，
作为 $A[s]$ 的 $A$-子模，有
$$
J^n =
I^n \oplus sI^{n - 1} \oplus s^2I^{n - 2} \ldots \oplus s^nA
\oplus s^{n + 1}A \oplus \ldots
$$
考虑开子概形
$$
\text{Proj}(B) = \text{Proj}(A[s] \oplus J \oplus J^2 \oplus \ldots)
\subset W
$$
最后，将元素 $s \in J$ 视为 $B$ 的次数 $1$ 元素，并记作 $S$。

\medskip\noindent
由于 $\text{Proj}$ 的形成与基变换交换（《概形的构造》引理 \ref{constructions-lemma-base-change-map-proj}），可见
$$
E = \text{Proj}(B \otimes_{A[s]} A/I) =
\text{Proj}((A/I \oplus I/I^2 \oplus I^2/I^3 \oplus \ldots)[S])
$$
由上面对 $J^n$ 各次幂的描述，立即验证
$B \otimes_{A[s]} A/I = \bigoplus J^n/J^{n + 1}$
确实如上所示。这证明了 (5)，因为由《除子》引理 \ref{divisors-lemma-affine-conormal-sheaf}，
$Z \cap \Spec(A)$ 在 $\Spec(A)$ 中的余法代数对应于分次
$A$-代数 $A/I \oplus I/I^2 \oplus I^2/I^3 \oplus \ldots$。

\medskip\noindent
回忆 $\text{Proj}(B)$ 由仿射开集 $D_+(S)$ 以及
$f \in I$ 时的 $D_+(f^{(1)})$ 覆盖；它们分别是仿射爆破代数
$A[s][\frac{J}{s}]$ 和 $A[s][\frac{J}{f}]$ 的谱，见
《除子》引理 \ref{divisors-lemma-blowing-up-affine} 和
《交换代数》定义 \ref{algebra-definition-blow-up}。
逐一描述这些仿射开集即可完成证明。

\medskip\noindent
先考虑开集 $D_+(S)$，即 $A[s][\frac{J}{s}]$ 的谱。
由上面对 $J$ 的各次幂的描述，
$$
A[s][\textstyle{\frac{J}{s}}] = \sum s^{-n}I^n[s] \subset A[s, s^{-1}]
$$
该环中的元素 $s$ 是非零因子，并且既定义例外除子 $E$，
也定义 $W_\infty$。因此 $D \cap D_+(S) = \emptyset$。
最后，$A[s][\frac{J}{s}]$ 对 $s$ 的商是余法代数
$$
A/I \oplus I/I^2 \oplus I^2/I^3 \oplus \ldots
$$
这证明了 (7)。

\medskip\noindent
再考虑开集 $D_+(f^{(1)})$，即 $A[s][\frac{J}{f}]$ 的谱。
由上面对 $J$ 的各次幂的描述，
$$
A[s][\textstyle{\frac{J}{f}}] =
A[\textstyle{\frac{I}{f}}][\textstyle{\frac{s}{f}}]
$$
其中 $\frac{s}{f}$ 是变量。元素 $f$ 在该环中是非零因子，
其零点概形定义例外除子 $E$。由于 $s$ 定义 $W_\infty$，且
$s = f \cdot \frac{s}{f}$，可知 $\frac{s}{f}$ 定义上面构造的除子
$D$。于是
$$
D \cap D_+(f^{(1)}) = \Spec(A[\textstyle{\frac{I}{f}}])
$$
这是爆破 $X'$ 在 $\Spec(A)$ 上的相应开集。具体而言，到
$(A, I)$ 的 Rees 代数的满分次 $A[s]$-代数映射
$B \to A \oplus I \oplus I^2 \oplus \ldots$
对应于 $\Spec(A[s])$ 上的闭浸入 $X' \to W$。
这便如需地证明了 $D = X'$。

\medskip\noindent
下面证明 (6)。注意，上一段中 $\frac{s}{f}$ 的零点概形，
是 $S$ 的零点概形在仿射开集 $D_+(f^{(1)})$ 上的限制。
因此 $S = 0$ 在 $E$ 上定义 $X' \cap E$。由 (5) 即得 (6)。

\medskip\noindent
最后还须证明 (8) 的最后一部分。这是清楚的，因为态射
$\mathbf{P}^1_Z \to W$ 在仿射局部由满射
$$
B \to B \otimes_{A[s]} A/I =
(A/I \oplus I/I^2 \oplus I^2/I^3 \oplus \ldots)[S] \to
A/I[S]
$$
以及同一化 $\text{Proj}(A/I[S]) = \Spec(A/I)$ 给出。
省略一些细节。






\section{高余维 Gysin 同态}
\label{chow-section-gysin-higher-codimension}

\noindent
设 $(S, \delta)$ 如情形 \ref{chow-situation-setup}，且
$X$ 是 $S$ 上局部有限型概形。本节考虑三元组
$$
(Z \to X, \mathcal{N}, \sigma : \mathcal{N}^\vee \to \mathcal{C}_{Z/X})
$$
其中 $Z \to X$ 是闭浸入，$\mathcal{N}$ 是局部自由
$\mathcal{O}_Z$-模，而
$\sigma : \mathcal{N}^\vee \to \mathcal{C}_{Z/X}$
是从 $\mathcal{N}$ 的对偶到 $Z$ 在 $X$ 中的余法层的满射；见
《概形态射》第 \ref{morphisms-section-conormal-sheaf} 节。
我们称 $\mathcal{N}$ 是{\it $Z$ 在 $X$ 中的虚拟法层}。

\begin{lemma}
\label{chow-lemma-pullback-virtual-normal-sheaf}
设 $(S, \delta)$ 如情形 \ref{chow-situation-setup}。设
$$
\xymatrix{
Z' \ar[r] \ar[d]_g & X' \ar[d]^f \\
Z \ar[r] & X
}
$$
是 $S$ 上局部有限型概形的笛卡儿图，其中水平箭头均为闭浸入。
若 $\mathcal{N}$ 是 $Z$ 在 $X$ 中的虚拟法层，则
$\mathcal{N}' = g^*\mathcal{N}$ 是 $Z'$ 在 $X'$ 中的虚拟法层。
\end{lemma}

\begin{proof}
这由《概形态射》引理 \ref{morphisms-lemma-conormal-functorial-flat} 所证明的映射
$g^*\mathcal{C}_{Z/X} \to \mathcal{C}_{Z'/X'}$ 的满性得出。
\end{proof}

\noindent
设 $(S, \delta)$ 如情形 \ref{chow-situation-setup}，且
$X$ 是 $S$ 上局部有限型概形。设 $\mathcal{N}$ 是闭浸入
$Z \to X$ 的虚拟法丛。在此情形下，令
$$
p : N = \underline{\Spec}_Z(\text{Sym}(\mathcal{N}^\vee)) \longrightarrow Z
$$
为 $Z$ 上的向量丛，其截面对应于 $\mathcal{N}$ 的截面。
此时存在典范闭浸入
$$
C_ZX \longrightarrow N_ZX \longrightarrow N
$$
第一个闭浸入是 Divisors, Equation
(\ref{divisors-equation-normal-cone-in-normal-bundle})，
第二个闭浸入对应于 $\sigma$ 诱导的满射
$\text{Sym}(\mathcal{N}^\vee) \to \text{Sym}(\mathcal{C}_{Z/X})$。
令
$$
b : W \longrightarrow \mathbf{P}^1_X
$$
为第 \ref{chow-section-blowup-Z-first} 节中构造的沿
$\infty(Z)$ 的爆破。由引理 \ref{chow-lemma-gysin-at-infty}，
存在典范双变类
$$
C \in A^0(W_\infty \to X)
$$
考虑 (\ref{chow-item-cone-is-open}) 中的开浸入
$j : C_ZX \to W_\infty$，以及上面构造的闭浸入
$i : C_ZX \to N$。由引理 \ref{chow-lemma-vectorbundle}，
对每个 $\alpha \in \CH_k(X)$，存在唯一的
$\beta \in \CH_*(Z)$，使得
$$
i_*j^*(C \cap \alpha) = p^*\beta
$$
定义 $c(Z \to X, \mathcal{N}) \cap \alpha = \beta$。

\begin{lemma}
\label{chow-lemma-construction-gysin}
上述构造定义了双变类\footnote{记号
$A^*(Z \to X)^\wedge$ 见注 \ref{chow-remark-completion-bivariant} 的讨论。若 $X$ 拟紧，则
$A^*(Z \to X)^\wedge = A^*(Z \to X)$。}
$$
c(Z \to X, \mathcal{N}) \in A^*(Z \to X)^\wedge
$$
此外，该构造与引理 \ref{chow-lemma-pullback-virtual-normal-sheaf} 中的基变换相容。
若 $\mathcal{N}$ 的秩恒为 $r$，则
$c(Z \to X, \mathcal{N}) \in A^r(Z \to X)$。
\end{lemma}

\begin{proof}
由于 $i_* \circ j^* \circ C$ 和 $p^*$ 都是双变类
（见引理 \ref{chow-lemma-flat-pullback-bivariant} 和
\ref{chow-lemma-push-proper-bivariant}），可以用等式
$$
i_* \circ j^* \circ C  = p^* \circ c(Z \to X, \mathcal{N})
$$
（作适当解释）将 $c(Z \to X, \mathcal{N})$ 定义为双变类。
这是可行的，因为由引理 \ref{chow-lemma-vectorbundle}，
$p^*$ 在 Chow 群上总是双射。

\medskip\noindent
设 $X' \to X$、$Z' \to X'$ 与 $\mathcal{N}'$ 如
引理 \ref{chow-lemma-pullback-virtual-normal-sheaf}。记
$c = c(Z \to X, \mathcal{N})$ 且
$c' = c(Z' \to X', \mathcal{N}')$。引理的第二个陈述是说，
$c'$ 是注 \ref{chow-remark-restriction-bivariant} 意义下
$c$ 的限制。由于该断言声称对所有局部有限型的 $X'/X$ 都成立，
一个形式论证表明，只须对
$\alpha' \in \CH_k(X')$ 检验
$c' \cap \alpha' = c \cap \alpha'$。为此，注意存在交换图
$$
\xymatrix{
C_{Z'}X' \ar[d] \ar[r] &
W'_\infty \ar[d] \ar[r] &
W' \ar[d] \ar[r] &
\mathbf{P}^1_{X'} \ar[d] \\
C_ZX \ar[r] &
W_\infty \ar[r] &
W \ar[r] &
\mathbf{P}^1_X
}
$$
它诱导闭浸入
$$
W' \to W \times_{\mathbf{P}^1_X} \mathbf{P}^1_{X'},\quad
W'_\infty \to W_\infty \times_X X',\quad
C_{Z'}X' \to C_ZX \times_Z Z'
$$
为得到 $c \cap \alpha'$，使用在引理 \ref{chow-lemma-gysin-at-infty} 中由态射
$W \times_{\mathbf{P}^1_X} \mathbf{P}^1_{X'} \to \mathbf{P}^1_{X'}$
定义的类 $C \cap \alpha'$。另一方面，为得到
$c' \cap \alpha'$，使用由态射
$W' \to \mathbf{P}^1_{X'}$ 定义的类 $C' \cap \alpha'$。
由引理 \ref{chow-lemma-gysin-at-infty-independent}，
$C' \cap \alpha'$ 沿闭浸入
$W'_\infty \to (W \times_{\mathbf{P}^1_X} \mathbf{P}^1_{X'})_\infty$
的推前等于 $C \cap \alpha'$。因此，由引理 \ref{chow-lemma-flat-pullback-proper-pushforward}，二者到开集
$$
C_{Z'}X' \subset W'_\infty,\quad
C_ZX \times_Z Z' \subset (W \times_{\mathbf{P}^1_X} \mathbf{P}^1_{X'})_\infty
$$
的拉回也相同。又有交换图
$$
\xymatrix{
C_{Z'} X' \ar[d] \ar[r] & N' \ar@{=}[d] \\
C_ZX \times_Z Z' \ar[r] & N \times_Z Z'
}
$$
这些类推前到 $N'$ 上得到同一个类，故在 $\CH_*(Z')$ 中得到
同一元素 $c \cap \alpha' = c' \cap \alpha'$。
\end{proof}

\begin{lemma}
\label{chow-lemma-gysin-decompose}
设 $(S, \delta)$ 如情形 \ref{chow-situation-setup}。
设 $X$ 是 $S$ 上局部有限型概形，且 $\mathcal{N}$
是 $X$ 的闭子概形 $Z$ 的虚拟法层。假设存在有限局部自由
$\mathcal{O}_Z$-模的短正合列
$0 \to \mathcal{N}' \to \mathcal{N} \to \mathcal{E} \to 0$，
使得给定满射
$\sigma : \mathcal{N}^\vee \to \mathcal{C}_{Z/X}$
分解经过映射
$\sigma' : (\mathcal{N}')^\vee \to \mathcal{C}_{Z/X}$。
则作为双变类，有
$$
c(Z \to X, \mathcal{N}) = c_{top}(\mathcal{E}) \circ c(Z \to X, \mathcal{N}')
$$
\end{lemma}

\begin{proof}
记 $N' \to N$ 为对应于满射
$\mathcal{N}^\vee \to (\mathcal{N}')^\vee$ 的向量丛闭浸入。
于是存在闭浸入
$$
C_ZX \to N' \to N
$$
所需双变类关系立即由引理 \ref{chow-lemma-easy-virtual-class} 得出。
\end{proof}

\begin{lemma}
\label{chow-lemma-gysin-excess}
设 $(S, \delta)$ 如情形 \ref{chow-situation-setup}。考虑
$S$ 上局部有限型概形的笛卡儿图
$$
\xymatrix{
Z' \ar[r] \ar[d]_g & X' \ar[d]^f \\
Z \ar[r] & X
}
$$
其中水平箭头均为闭浸入。设 $\mathcal{N}$，分别地\ 
$\mathcal{N}'$，是 $Z \subset X$，分别地\ $Z' \to X'$，
的虚拟法层。假设给定 $Z'$ 上有限局部自由模的短正合列
$0 \to \mathcal{N}' \to g^*\mathcal{N} \to \mathcal{E} \to 0$，
使得图
$$
\xymatrix{
g^*\mathcal{N}^\vee \ar[r] \ar[d] &
(\mathcal{N}')^\vee \ar[d] \\
g^*\mathcal{C}_{Z/X} \ar[r] &
\mathcal{C}_{Z'/X'}
}
$$
交换。则在 $A^*(Z' \to X')^\wedge$ 中有
$$
res(c(Z \to X, \mathcal{N})) =
c_{top}(\mathcal{E}) \circ c(Z' \to X', \mathcal{N}')
$$
\end{lemma}

\begin{proof}
由引理 \ref{chow-lemma-construction-gysin}，
$res(c(Z \to X, \mathcal{N})) = c(Z' \to X', g^*\mathcal{N})$；
于是等式由引理 \ref{chow-lemma-gysin-decompose} 得出。
\end{proof}

\noindent
设 $(S, \delta)$ 如情形 \ref{chow-situation-setup}，且
$X$ 是 $S$ 上局部有限型概形。设 $\mathcal{N}$ 是
$X$ 的闭子概形 $Z$ 的虚拟法层。设 $Y \to X$ 是局部有限型态射。
假设 $Z \times_X Y \to Y$ 是正则闭浸入；见《除子》第 \ref{divisors-section-regular-immersions} 节。此时余法层
$\mathcal{C}_{Z \times_X Y/Y}$ 是有限局部自由
$\mathcal{O}_{Z \times_X Y}$-模，并得到短正合列
$$
0 \to \mathcal{E}^\vee \to
\mathcal{N}^\vee|_{Z \times_X Y} \to \mathcal{C}_{Z \times_X Y/Y} \to 0
$$
商 $\mathcal{N}|_{Y \times_X Z} \to \mathcal{E}$
称为此情形的{\it 超额法层}。

\begin{lemma}
\label{chow-lemma-gysin-fundamental}
在上述情形中，假设 $\dim_\delta(Y) = n$，且
$\mathcal{C}_{Y \times_X Z/Z}$ 的秩恒为 $r$。则在
$\CH_*(Z \times_X Y)$ 中有
$$
c(Z \to X, \mathcal{N}) \cap [Y]_n =
c_{top}(\mathcal{E}) \cap [Z \times_X Y]_{n - r}
$$
\end{lemma}

\begin{proof}
双变类 $c_{top}(\mathcal{E}) \in A^*(Z \times_X Y)$
在注 \ref{chow-remark-top-chern-class} 中定义。
由引理 \ref{chow-lemma-construction-gysin}，可将 $X$ 替换为 $Y$。
因此可假设 $Z \to X$ 是余维 $r$ 的正则闭浸入，
$\dim_\delta(X) = n$，并须在 $\CH_*(Z)$ 中证明
$c(Z \to X, \mathcal{N}) \cap [X]_n =
c_{top}(\mathcal{E}) \cap [Z]_{n - r}$。
由引理 \ref{chow-lemma-gysin-decompose}，甚至可假设
$\mathcal{N}^\vee \to \mathcal{C}_{Z/X}$ 是同构。
换言之，须在 $\CH_*(Z)$ 中证明
$c(Z \to X, \mathcal{C}_{Z/X}^\vee) \cap [X]_n = [Z]_{n - r}$。

\medskip\noindent
下面逐步追踪
$c(Z \to X, \mathcal{C}_{Z/X}^\vee) \cap [X]_n$ 的定义。
设 $b : W \to \mathbf{P}^1_X$ 是沿 $\infty(Z)$ 的爆破。
首先计算 $C \cap [X]_n$，其中
$C \in A^0(W_\infty \to X)$ 是
引理 \ref{chow-lemma-gysin-at-infty} 中的类。注意
$[W]_{n + 1}$ 是 $W$ 上的周，它在 $\mathbf{A}^1_X$
上的限制等于 $[X]_n$ 的平坦拉回。因此
$C \cap [X]_n$ 等于 $i_\infty^*[W]_{n + 1}$。
由于 $W_\infty$ 是 $W$ 上的有效 Cartier 除子，
$i_\infty^*[W]_{n + 1} = [W_\infty]_n$；见
引理 \ref{chow-lemma-easy-gysin}。该类在开集
$C_ZX \subset W_\infty$ 上的限制当然正是 $[C_ZX]_n$。
因为 $Z \subset X$ 是正则嵌入，有
$$
\mathcal{C}_{Z/X, *} = \text{Sym}(\mathcal{C}_{Z/X})
$$
作为分次 $\mathcal{O}_Z$-代数；见《除子》引理 \ref{divisors-lemma-quasi-regular-immersion}。因此
$p : N = C_ZX \to Z$ 是与秩 $r$ 的有限局部自由模
$\mathcal{C}_{Z/X}$ 对应之向量丛的结构态射。显然
$p^*[Z]_{n - r} = [C_ZX]_n$，证明完毕。
\end{proof}

\begin{lemma}
\label{chow-lemma-gysin-easy}
设 $(S, \delta)$ 如情形 \ref{chow-situation-setup}。
设 $X$ 是 $S$ 上局部有限型概形，且 $\mathcal{N}$
是 $X$ 的闭子概形 $Z$ 的虚拟法层。设
$Y \to X$ 是局部有限型态射。给定整数 $r$、$n$，假设
\begin{enumerate}
\item $\mathcal{N}$ 局部自由且秩为 $r$；
\item $Y$ 的每个不可约分支的 $\delta$-维数均为 $n$；
\item $\dim_\delta(Z \times_X Y) \leq n - r$；并且
\item 对满足 $\delta(\xi) = n - r$ 的
$\xi \in Z \times_X Y$，局部环
$\mathcal{O}_{Y, \xi}$ 是 Cohen--Macaulay 环。
\end{enumerate}
则在 $\CH_{n - r}(Z \times_X Y)$ 中有
$c(Z \to X, \mathcal{N}) \cap [Y]_n = [Z \times_X Y]_{n - r}$。
\end{lemma}

\begin{proof}
由于 $Z \times_X Y$ 是 $Y$ 的闭子概形，该陈述是有意义的。
因为 $\mathcal{N}$ 的秩为 $r$，可知
$c(Z \to X, \mathcal{N}) \cap [Y]_n$ 属于
$\CH_{n - r}(Z \times_X Y)$。由于
$\dim_\delta(Z \cap Y) \leq n - r$，Chow 群
$\CH_{n - r}(Z \times_X Y)$ 由 $\delta$-维数为
$n - r$ 的各不可约分支 $W \subset Z \times_X Y$
的周类自由生成。设 $\xi \in W$ 是泛点。
由假设 (2)，$\dim(\mathcal{O}_{Y, \xi}) = r$。
另一方面，由于 $\mathcal{N}$ 的秩为 $r$，且
$\mathcal{N}^\vee \to \mathcal{C}_{Z/X}$ 是满射，
$Z$ 的理想层在局部由 $r$ 个方程截出。因此
$Z \times_X Y$ 在 $Y$ 中的拟凝聚理想层
$\mathcal{I} \subset \mathcal{O}_Y$ 在局部由 $r$ 个元素生成。
由于 $\mathcal{O}_{Y, \xi}$ 是维数 $r$ 的 Cohen--Macaulay 环，
且 $\mathcal{I}_\xi$ 是定义理想（因为 $\xi$ 是
$Z \times_X Y$ 的泛点），$\mathcal{I}_\xi$ 由正则序列生成
（《交换代数》引理 \ref{algebra-lemma-reformulate-CM}）。
由《除子》引理 \ref{divisors-lemma-Noetherian-scheme-regular-ideal}，
在 $\xi$ 的某个开邻域 $V \subset Y$ 上，$\mathcal{I}$
由正则序列生成。根据上面对
$\CH_{n - r}(Z \times_X Y)$ 的描述，只须在
$\CH_{n - r}(Z \times_X V)$ 中证明
$c(Z \to X, \mathcal{N}) \cap [V]_n = [Z \times_X V]_{n - r}$。
这由引理 \ref{chow-lemma-gysin-fundamental} 得出，因为超额法层在
$V$ 上为 $0$。
\end{proof}

\begin{lemma}
\label{chow-lemma-gysin-agrees}
设 $(S, \delta)$ 如情形 \ref{chow-situation-setup}，且
$X$ 是 $S$ 上局部有限型概形。设
$(\mathcal{L}, s, i : D \to X)$ 是定义 \ref{chow-definition-gysin-homomorphism} 中的三元组。
将 Gysin 同态 $i^*$ 视为 $A^1(D \to X)$ 的元素
（见引理 \ref{chow-lemma-gysin-bivariant}），则它与双变类
$c(D \to X, \mathcal{N}) \in A^1(D \to X)$ 相同；
后者把 $\mathcal{N} = i^*\mathcal{L}$ 视为
$D$ 在 $X$ 中的虚拟法层而构造。
\end{lemma}

\begin{proof}
使用引理 \ref{chow-lemma-bivariant-zero} 的判据。
因此可假设 $X$ 是整概形，并须证明
$i^*[X]$ 等于 $c \cap [X]$。令
$n = \dim_\delta(X)$。照例分两种情形。

\medskip\noindent
若 $X = D$，则两个类都由
$c_1(\mathcal{N}) \cap [X]_n$ 表示。见
引理 \ref{chow-lemma-gysin-fundamental} 和定义 \ref{chow-definition-gysin-homomorphism}。

\medskip\noindent
若 $D \not = X$，则 $D \to X$ 是有效 Cartier 除子，
特别地，是余维 $1$ 的正则闭浸入。再次由
引理 \ref{chow-lemma-gysin-fundamental} 得到
$c(D \to X, \mathcal{N}) \cap [X]_n = [D]_{n - 1}$。
根据定义，Gysin 同态也具有同一性质，故仍得到结论。
\end{proof}

\begin{lemma}
\label{chow-lemma-gysin-commutes}
设 $(S, \delta)$ 如情形 \ref{chow-situation-setup}，且
$X$ 是 $S$ 上局部有限型概形。设 $Z \subset X$ 是具有虚拟法层
$\mathcal{N}$ 的闭子概形。设 $Y \to X$ 局部有限型，且
$c \in A^*(Y \to X)$。则 $c$ 与 $c(Z \to X, \mathcal{N})$
交换（注 \ref{chow-remark-bivariant-commute}）。
\end{lemma}

\begin{proof}
可用引理 \ref{chow-lemma-bivariant-zero} 检验此事。
因此可假设 $X$ 是整概形，并须在 $\CH_*(Z \times_X Y)$ 中证明
$c \cap c(Z \to X, \mathcal{N}) \cap [X] =
c(Z \to X, \mathcal{N}) \cap c \cap [X]$。

\medskip\noindent
若 $Z = X$，则由引理 \ref{chow-lemma-gysin-fundamental}，
$c(Z \to X, \mathcal{N}) = c_{top}(\mathcal{N})$；
它与双变类 $c$ 交换，见引理 \ref{chow-lemma-cap-commutative-chern}。

\medskip\noindent
假设 $Z$ 不等于 $X$。由引理 \ref{chow-lemma-bivariant-zero}，
甚至只须在对 $X$ 沿某个非零理想爆破后证明结论。
沿 $Z$ 的理想层爆破 $X$。这化归为 $Z$ 是有效 Cartier 除子的情形；
见《除子》引理 \ref{divisors-lemma-blowing-up-gives-effective-Cartier-divisor}。

\medskip\noindent
若 $Z$ 是有效 Cartier 除子，则
$$
c(Z \to X, \mathcal{N}) =
c_{top}(\mathcal{E}) \circ i^*
$$
其中 $i^* \in A^1(Z \to X)$ 是与 $i : Z \to X$ 对应的
Gysin 同态（引理 \ref{chow-lemma-gysin-bivariant}），而
$\mathcal{E}$ 是 $\mathcal{N}^\vee \to \mathcal{C}_{Z/X}$
之核的对偶；见引理 \ref{chow-lemma-gysin-decompose} 和
\ref{chow-lemma-gysin-agrees}。陈类位于双变环的中心
（以引理 \ref{chow-lemma-cap-commutative-chern} 所表述的强意义），
而按双变类的定义，$c$ 与 Gysin 同态 $i^*$ 交换，故结论成立。
\end{proof}

\noindent
设 $(S, \delta)$ 如情形 \ref{chow-situation-setup}。
设 $X$ 是 $S$ 上局部有限型、$\delta$-维数为 $n$ 的整概形。
设 $Z \subset Y \subset X$ 是闭子概形，且二者均为
$X$ 中的有效 Cartier 除子。记 $o : Y \to C_Y X$
为 $Y$ 在 $X$ 中的法线锥之零截面。由于
$C_YX$ 是 $Y$ 上的线丛，得到双变类
$o^* \in A^1(Y \to C_YX)$；见引理 \ref{chow-lemma-gysin-bivariant}。

\begin{lemma}
\label{chow-lemma-relation-normal-cones}
采用上述记号，在 $\CH_{n - 1}(Y \times_{o, C_Y X} C_ZX)$ 中有
$$
o^*[C_ZX]_n = [C_Z Y]_{n - 1}
$$
\end{lemma}

\begin{proof}
记 $W \to \mathbf{P}^1_X$ 为第 \ref{chow-section-blowup-Z-first} 节中沿 $\infty(Z)$ 的爆破。
类似地，记 $W' \to \mathbf{P}^1_X$ 为沿 $\infty(Y)$ 的爆破。
由于 $\infty(Z) \subset \infty(Y)$，得到理想层的反向包含，
从而得到定义这些爆破的分次代数之间的映射。这产生从
$W$ 到 $W'$ 的有理态射；事实上，它有典范代表
$$
W \supset U \longrightarrow W'
$$
见《概形的构造》引理 \ref{constructions-lemma-morphism-relative-proj}。
一次局部计算（省略）表明，$U$ 至少包含 $W$ 中所有不位于
$\infty$ 上方的点，以及特殊纤维的开子概形 $C_Z X$。
缩小 $U$ 后，可假设 $U_\infty = C_Z X$ 且
$\mathbf{A}^1_X \subset U$。另一次局部计算（省略）表明，
态射 $U_\infty \to W'_\infty$ 诱导典范态射
$C_Z X \to C_Y X \subset W'_\infty$；这是由
$Z \subset Y$ 所给的理想层包含诱导的法锥态射。
记 $W'' \subset W$ 为 $\mathbf{P}^1_Y \subset \mathbf{P}^1_X$
在 $W$ 中的严格变换。由《除子》引理 \ref{divisors-lemma-strict-transform}，$W''$ 是
$\mathbf{P}^1_Y$ 沿 $\infty(Z)$ 的爆破，因而
$(W'' \cap U)_\infty = C_ZY$。

\medskip\noindent
考虑 (\ref{chow-item-find-Z-in-blowup}) 中的有效 Cartier 除子
$i : \mathbf{P}^1_Y \to W'$，以及由引理 \ref{chow-lemma-gysin-bivariant} 得到的相应双变类
$i^* \in A^1(\mathbf{P}^1_Y \to W')$。
类似地，记 $(i'_\infty)^* \in A^1(W'_\infty \to W')$
为无穷远处的 Gysin 映射。注意，$i'_\infty$ 在 $U$ 上的限制
（注 \ref{chow-remark-restriction-bivariant}）等于
$i_\infty^* \in A^1(W_\infty \to W)$ 在 $U$ 上的限制。
一方面，
$$
(i'_\infty)^* i^* [U]_{n + 1} =
i_\infty^* i^* [U]_{n + 1} =
i_\infty^* [(W'' \cap U)_\infty]_{n + 1} =
[C_ZY]_n
$$
这是因为 $i_\infty^*$ 消灭所有支撑在 $\infty$ 上方的类，
$i^*[U]$ 与 $[W'']$ 作为 $\mathbf{A}^1$ 上的周相同，
并且 $C_ZY$ 是 $W'' \cap U$ 在 $\infty$ 上的纤维。
另一方面，
$$
(i'_\infty)^* i^* [U]_{n + 1} =
i^* i_\infty^*[U]_{n + 1} =
i^* [U_\infty] =
o^*[C_YX]_n
$$
这是因为 $(i'_\infty)^*$ 与 $i^*$ 交换
（引理 \ref{chow-lemma-gysin-commutes-gysin}），并且
$i : \mathbf{P}^1_Y \to W'$ 在 $\infty$ 上的纤维分解为
$o : Y \to C_YX$ 与开浸入 $C_YX \to W'_\infty$。
引理得证。
\end{proof}

\begin{lemma}
\label{chow-lemma-gysin-composition}
设 $(S, \delta)$ 如情形 \ref{chow-situation-setup}。
设 $Z \subset Y \subset X$ 是某个 $S$ 上局部有限型概形的闭子概形。
设 $\mathcal{N}$ 是 $Z \subset X$ 的虚拟法层，
$\mathcal{N}'$ 是 $Z \subset Y$ 的虚拟法层，而
$\mathcal{N}''$ 是 $Y \subset X$ 的虚拟法层。假设存在交换图
$$
\xymatrix{
(\mathcal{N}'')^\vee|_Z \ar[r] \ar[d] &
\mathcal{N}^\vee \ar[r] \ar[d] &
(\mathcal{N}')^\vee \ar[d] \\
\mathcal{C}_{Y/X}|_Z \ar[r] &
\mathcal{C}_{Z/X} \ar[r] &
\mathcal{C}_{Z/Y}
}
$$
其中下方序列来自《态射进阶》引理 \ref{more-morphisms-lemma-transitivity-conormal}，而上方序列是短正合列。
则在 $A^*(Z \to X)^\wedge$ 中有
$$
c(Z \to X, \mathcal{N}) =
c(Z \to Y, \mathcal{N}') \circ c(Y \to X, \mathcal{N}'')
$$
\end{lemma}

\begin{proof}
注意，沿任意局部有限型态射 $X' \to X$ 作基变换后，
这些假设仍然成立（虚拟法层的短正合列在局部分裂，
故沿任意基变换仍然正合）。因此可用
引理 \ref{chow-lemma-bivariant-zero} 检验双变类的等式。
于是可假设 $X$ 是整概形，并须证明
$c(Z \to X, \mathcal{N}) \cap [X] =
c(Z \to Y, \mathcal{N}') \cap c(Y \to X, \mathcal{N}'') \cap [X]$。

\medskip\noindent
若 $Y = X$，则
\begin{align*}
c(Z \to Y, \mathcal{N}') \cap c(Y \to X, \mathcal{N}'') \cap [X]
& =
c(Z \to Y, \mathcal{N}') \cap c_{top}(\mathcal{N}'') \cap [Y] \\
& =
c_{top}(\mathcal{N}''|_Z) \cap c(Z \to Y, \mathcal{N}') \cap [Y] \\
& =
c(Z \to X, \mathcal{N}) \cap [X] 
\end{align*}
第一个等式由引理 \ref{chow-lemma-gysin-decompose} 得出。
第二个等式是因为陈类与双变类交换
（引理 \ref{chow-lemma-cap-commutative-chern}）。
第三个等式由引理 \ref{chow-lemma-gysin-decompose} 得出。

\medskip\noindent
假设 $Y \not = X$。由引理 \ref{chow-lemma-bivariant-zero}，
甚至只须在对 $X$ 沿某个非零理想爆破后证明结论。
沿 $Y$ 的理想层与 $Z$ 的理想层之积爆破 $X$。
这化归为 $Y$ 和 $Z$ 均为 $X$ 上有效 Cartier 除子的情形；见
《除子》引理 \ref{divisors-lemma-blowing-up-gives-effective-Cartier-divisor} 和
\ref{divisors-lemma-blowing-up-two-ideals}。

\medskip\noindent
记 $\mathcal{N}'' \to \mathcal{E}$ 为有限局部自由
$\mathcal{O}_Z$-模的满射，使得
$0 \to \mathcal{E}^\vee \to (\mathcal{N}'')^\vee \to \mathcal{C}_{Y/X} \to 0$
是短正合列。于是 $\mathcal{N} \to \mathcal{E}|_Z$
也是满射。记 $\mathcal{N}_1$ 为该映射的有限局部自由核，并注意
$\mathcal{N}^\vee \to \mathcal{C}_{Z/X}$ 分解经过
$\mathcal{N}_1$。由引理 \ref{chow-lemma-gysin-decompose}，
$$
c(Y \to X, \mathcal{N}'') = c_{top}(\mathcal{E}) \circ
c(Y \to X, \mathcal{C}_{Y/X}^\vee)
$$
且
$$
c(Z \to X, \mathcal{N}) = c_{top}(\mathcal{E}|_Z) \circ
c(Z \to X, \mathcal{N}_1)
$$
由于丛的陈类与双变类交换
（引理 \ref{chow-lemma-cap-commutative-chern}），只须在
$A^*(Z \to X)$ 中证明
$$
c(Z \to X, \mathcal{N}_1) =
c(Z \to Y, \mathcal{N}') \circ c(Y \to X, \mathcal{C}_{Y/X}^\vee)
$$
为此可假设 $\mathcal{N}'' = \mathcal{C}_{Y/X}$。
这化归为下一段所讨论的情形。

\medskip\noindent
在本段中，$Z$ 和 $Y$ 是维数为 $n$ 的整概形 $X$
上的有效 Cartier 除子，且
$\mathcal{N}'' = \mathcal{C}_{Y/X}$。此时由
引理 \ref{chow-lemma-gysin-fundamental}，
$c(Y \to X, \mathcal{C}_{Y/X}^\vee) \cap [X] = [Y]_{n - 1}$。
因此须证明
$c(Z \to X, \mathcal{N}) \cap [X] = c(Z \to Y, \mathcal{N}') \cap [Y]_{n - 1}$。
记 $N$ 和 $N'$ 为与 $\mathcal{N}$ 和 $\mathcal{N}'$
对应的 $Z$ 上向量丛。考虑 $Z$ 上锥和向量丛的交换图
$$
\xymatrix{
N' \ar[r]_i &
N \ar[r] &
(C_Y X) \times_Y Z \\
C_Z Y \ar[r] \ar[u] &
C_Z X \ar[u]
}
$$
注意，$N'$ 是 $N$ 中在 $Z$ 上相对的有效 Cartier 除子，且
$$
\xymatrix{
N' \ar[d] \ar[r]_i & N \ar[d] \\
Z \ar[r]^-o & (C_Y X) \times_Y Z
}
$$
是笛卡儿图，其中 $o$ 是 $Y$ 上线丛
$C_Y X$ 的零截面。由引理 \ref{chow-lemma-relation-normal-cones}，有
$o^*[C_ZX]_n = [C_Z Y]_{n - 1}$，作为
$$
\CH_{n - 1}(Y \times_{o, C_Y X} C_ZX) =
\CH_{n - 1}(Z \times_{o, (C_Y X) \times_Y Z} C_ZX)
$$
中的等式。由上方正方形的笛卡儿性质，这蕴含
$$
i^*[C_ZX]_n = [C_Z Y]_{n - 1}
$$
在 $\CH_{n - 1}(N')$ 中成立。现在注意
$\gamma = c(Z \to X, \mathcal{N}) \cap [X]$ 与
$\gamma' = c(Z \to Y, \mathcal{N}') \cap [Y]_{n - 1}$
分别由 $\CH_n(N)$ 中的 $p^*\gamma = [C_Z X]_n$，
以及 $\CH_{n - 1}(N')$ 中的
$(p')^*\gamma' = [C_Z Y]_{n - 1}$ 刻画。
由引理 \ref{chow-lemma-relative-effective-cartier}，
$i^* \circ p^* = (p')^*$，证明完毕。
\end{proof}

\begin{remark}[浸入情形的变体]
\label{chow-remark-gysin-for-immersion}
设 $(S, \delta)$ 如情形 \ref{chow-situation-setup}。
设 $X$ 是 $S$ 上局部有限型概形，且
$i : Z \to X$ 是概形的浸入。在此情形下，
\begin{enumerate}
\item $Z$ 在 $X$ 中的余法层 $\mathcal{C}_{Z/X}$ 已有定义
（《概形态射》定义 \ref{morphisms-definition-conormal-sheaf}）；
\item 由有限局部自由 $\mathcal{O}_Z$-模 $\mathcal{N}$
和满射 $\sigma : \mathcal{N}^\vee \to \mathcal{C}_{Z/X}$
组成的对，称为浸入 $Z \to X$ 的虚拟法丛；
\item 选取开子概形 $U \subset X$，使得 $Z \to X$
分解经过闭浸入 $Z \to U$，并定义
$c(Z \to X, \mathcal{N}) = c(Z \to U, \mathcal{N}) \circ (U \to X)^*$。
\end{enumerate}
双变类 $c(Z \to X, \mathcal{N})$ 与开子概形 $U$ 的选择无关。
所有引理对这一稍为一般的构造都有直接对应的版本。省略细节。
\end{remark}







\section{计算若干类}
\label{chow-section-calculate}

\noindent
为进一步推进，需要计算上面构造的若干类的取值。

\begin{lemma}
\label{chow-lemma-compute-koszul}
设 $(S, \delta)$ 如情形 \ref{chow-situation-setup}。
设 $X$ 是 $S$ 上局部有限型概形，且 $\mathcal{E}$
是秩 $r$ 的局部自由 $\mathcal{O}_X$-模。则
$$
\prod\nolimits_{n = 0, \ldots, r} c(\wedge^n \mathcal{E})^{(-1)^n} =
1 - (r - 1)! c_r(\mathcal{E}) + \ldots
$$
\end{lemma}

\begin{proof}
由分裂原理，可将其化为在以 $\mathcal{E}$ 的陈根
$x_1, \ldots, x_r$ 为变量的多项式环中计算。见
第 \ref{chow-section-splitting-principle} 节。注意
$$
c(\wedge^n \mathcal{E}) =
\prod\nolimits_{1 \leq i_1 < \ldots < i_n \leq r}
(1 + x_{i_1} + \ldots + x_{i_n})
$$
因此，引理等式左边的对数为
$$
-
\sum\nolimits_{p \geq 1}
\sum\nolimits_{n = 0}^r
\sum\nolimits_{1 \leq i_1 < \ldots < i_n \leq r}
\frac{(-1)^{p + n}}{p}(x_{i_1} + \ldots + x_{i_n})^p
$$
请注意前面的负号。不过，
$$
\sum\nolimits_{p \geq 0}
\sum\nolimits_{n = 0}^r
\sum\nolimits_{1 \leq i_1 < \ldots < i_n \leq r}
\frac{(-1)^{p + n}}{p!}(x_{i_1} + \ldots + x_{i_n})^p
=
\prod (1 - e^{-x_i})
$$
由此可见，所得陈类的第一个非零项在次数 $r$，且等于所预言的值。
\end{proof}

\begin{lemma}
\label{chow-lemma-compute-section}
设 $(S, \delta)$ 如情形 \ref{chow-situation-setup}。
设 $X$ 是 $S$ 上局部有限型概形，且 $\mathcal{C}$
是秩 $r$ 的局部自由 $\mathcal{O}_X$-模。考虑态射
$$
X = \underline{\text{Proj}}_X(\mathcal{O}_X[T])
\xrightarrow{i}
E = \underline{\text{Proj}}_X(\text{Sym}^*(\mathcal{C})[T])
\xrightarrow{\pi}
X
$$
对 $t = 1, \ldots, r - 1$ 有
$c_t(i_*\mathcal{O}_X) = 0$，并且在 $A^0(C \to E)$ 中有
$$
p^* \circ \pi_* \circ c_r(i_*\mathcal{O}_X) = (-1)^{r - 1}(r - 1)! j^*
$$
其中 $j : C \to E$ 和 $p : C \to X$
分别是向量丛
$C = \underline{\Spec}(\text{Sym}^*(\mathcal{C}))$
的包含态射和结构态射。
\end{lemma}

\begin{proof}
典范映射 $\pi^*\mathcal{C} \to \mathcal{O}_E(1)$
恰好沿 $i(X)$ 消失。因此，映射
$$
\pi^*\mathcal{C} \otimes \mathcal{O}_E(-1) \to \mathcal{O}_E
$$
上的 Koszul 复形是 $i_*\mathcal{O}_X$ 的消解。
特别地，$i_*\mathcal{O}_X$ 是 $D(\mathcal{O}_E)$ 中陈类已有定义的
完美对象。对 $t = 1, \ldots, t - 1$，
$c_t(i_*\mathcal{O}_X)$ 的消失由
引理 \ref{chow-lemma-compute-koszul} 得出。该引理还给出
$$
c_r(i_*\mathcal{O}_X) = - (r - 1)!
c_r(\pi^*\mathcal{C} \otimes \mathcal{O}_E(-1))
$$
另一方面，由引理 \ref{chow-lemma-chern-classes-dual}，
$$
c_r(\pi^*\mathcal{C} \otimes \mathcal{O}_E(-1)) =
(-1)^r c_r(\pi^*\mathcal{C}^\vee \otimes \mathcal{O}_E(1))
$$
而 $\pi^*\mathcal{C}^\vee \otimes \mathcal{O}_E(1)$
有一个恰好沿 $i(X)$ 消失的截面 $s$。

\medskip\noindent
将 $X$ 替换为 $X$ 上局部有限型概形后，只须证明等式两边
对元素 $\alpha \in \CH_*(E)$ 的作用相同。由于
$C \to X$ 是向量丛，$C$ 上每个周类都形如
$p^*\beta$，其中 $\beta \in \CH_*(X)$
（引理 \ref{chow-lemma-vectorbundle}）。因此由
引理 \ref{chow-lemma-restrict-to-open}，可写
$\alpha = \pi^*\beta + \gamma$，其中 $\gamma$
支撑在 $E \setminus C$ 上。使用上述等式，只须证明
$$
p^*(\pi_*(c_r(\pi^*\mathcal{C}^\vee \otimes \mathcal{O}_E(1)) \cap [W])) =
j^*[W]
$$
其中 $W \subset E$ 是整闭子概形，并且或者 (a) 与 $C$ 不交，
或者 (b) 形如 $W = \pi^{-1}Y$，其中
$Y \subset X$ 是整闭子概形。使用截面 $s$ 和
引理 \ref{chow-lemma-top-chern-class}，在情形 (a) 有
$c_r(\pi^*\mathcal{C}^\vee \otimes \mathcal{O}_E(1)) \cap [W] = 0$，
在情形 (b) 有
$c_r(\pi^*\mathcal{C}^\vee \otimes \mathcal{O}_E(1)) \cap [W] = [i(Y)]$。
结论容易由此得出；省略细节。
\end{proof}

\begin{lemma}
\label{chow-lemma-agreement-with-loc-chern}
设 $(S, \delta)$ 如情形 \ref{chow-situation-setup}。
设 $i : Z \to X$ 是 $S$ 上局部有限型概形之间余维 $r$
的正则闭浸入。令 $\mathcal{N} = \mathcal{C}_{Z/X}^\vee$
为法层。若 $X$ 拟紧（或其不可约分支拟紧），则对
$t = 1, \ldots, r - 1$ 有
$c_t(Z \to X, i_*\mathcal{O}_Z) = 0$，并且
$$
c_r(Z \to X, i_*\mathcal{O}_Z) = (-1)^{r - 1} (r - 1)! c(Z \to X, \mathcal{N})
\quad\text{于}\quad
A^r(Z \to X)
$$
其中 $c_t(Z \to X, i_*\mathcal{O}_Z)$ 是定义 \ref{chow-definition-localized-chern} 中的局部化陈类。
\end{lemma}

\begin{proof}
对任意 $x \in Z$，可选取仿射开邻域
$\Spec(A) \subset X$，使得
$Z \cap \Spec(A) = V(f_1, \ldots, f_r)$，其中
$f_1, \ldots, f_r \in A$ 是正则序列。见《除子》定义 \ref{divisors-definition-regular-immersion} 和引理 \ref{divisors-lemma-Noetherian-scheme-regular-ideal}。
于是由例如《代数进阶》引理 \ref{more-algebra-lemma-regular-koszul-regular}，
$f_1, \ldots, f_r$ 上的 Koszul 复形是
$A/(f_1, \ldots, f_r)$ 的消解。因此
$A/(f_1, \ldots, f_r)$ 作为 $A$-模是完美的。
由此 $F = i_*\mathcal{O}_Z$ 是 $D(\mathcal{O}_X)$ 中的完美对象，
其在 $X \setminus Z$ 上的限制为零。$X$ 拟紧
（或其不可约分支拟紧）的假设意味着局部化陈类
$c_t(Z \to X, i_*\mathcal{O}_Z)$ 已有定义；见
情形 \ref{chow-situation-loc-chern} 和定义 \ref{chow-definition-localized-chern}。综上，该陈述是有意义的。

\medskip\noindent
记 $b : W \to \mathbf{P}^1_X$ 为第 \ref{chow-section-blowup-Z-first} 节中沿 $\infty(Z)$ 的爆破。
由 (\ref{chow-item-find-Z-in-blowup})，存在闭浸入
$$
i' : \mathbf{P}^1_Z \longrightarrow W
$$
我们断言 $Q = i'_*\mathcal{O}_{\mathbf{P}^1_Z}$
是 $D(\mathcal{O}_W)$ 中的完美对象，并且 $F$、$Q$
满足引理 \ref{chow-lemma-independent-loc-chern-bQ} 的假设。

\medskip\noindent
假设该断言成立。引理 \ref{chow-lemma-independent-loc-chern-bQ} 给出：对所有
$p \geq 1$，
$$
c_p(Z \to X, F) = c'_p(Q) = (E \to Z)_* \circ c'_p(Q|_E) \circ C
$$
注意，$Q|_E$ 等于 $Z$ 的结构层沿态射 $Z \to E$ 的推前，
而该态射是 $i'$ 沿 $\infty$ 的基变换。
因此，将引理 \ref{chow-lemma-compute-section} 应用于 $E \to Z$，
便得到 $1 \leq t \leq r - 1$ 时
$c_t(Z \to X, F)$ 的消失。由于
$\mathcal{C}_{Z/X} = \mathcal{N}^\vee$ 局部自由，双变类
$c(Z \to X, \mathcal{N})$ 由关系
$$
j^* \circ C = p^* \circ c(Z \to X, \mathcal{N})
$$
刻画，其中 $j : C_ZX \to W_\infty$ 和 $p : C_ZX \to Z$
为给定映射。（回忆 $C \in A^0(W_\infty \to X)$ 是
引理 \ref{chow-lemma-gysin-at-infty} 中的类。）
因此，引理陈述中的显示等式由
引理 \ref{chow-lemma-compute-section} 中的相应等式得出。

\medskip\noindent
证明该断言。设 $A$ 和 $f_1, \ldots, f_r$ 如上。
考虑第 \ref{chow-section-blowup-Z-first} 节中的仿射开集
$\Spec(A[s]) \subset \mathbf{P}^1_X$。回忆在该开集上，
$s = 0$ 定义 $(\mathbf{P}^1_X)_\infty$。因此在
$\Spec(A[s])$ 上，我们沿正则序列
$s, f_1, \ldots, f_r$ 生成的理想爆破。由《代数进阶》引理 \ref{more-algebra-lemma-blowup-regular-sequence}，这 $r + 1$
个仿射图均为 $A[s]$ 上的整体完全交。与仿射爆破代数
$$
A[s][f_1/s, \ldots, f_r/s] = A[s, y_1, \ldots, y_r]/(sy_i - f_i)
$$
对应的图包含 $i'(Z \cap \Spec(A))$，后者是由
$y_1, \ldots, y_r$ 截出的闭子概形。由于
$y_1, \ldots, y_r, sy_1 - f_1, \ldots, sy_r - f_r$
是多项式环 $A[s, y_1, \ldots, y_r]$ 中的正则序列，
可知 $i'$ 是正则浸入。省略一些细节。
如上可得 $Q = i'_*\mathcal{O}_{\mathbf{P}^1_Z}$
是 $D(\mathcal{O}_W)$ 中的完美对象。$F$ 与 $Q$ 的其他假设在
引理 \ref{chow-lemma-independent-loc-chern-bQ}
（以及引理 \ref{chow-lemma-localized-chern-pre}）中均立即得到验证。
\end{proof}

\begin{lemma}
\label{chow-lemma-actual-computation}
在引理 \ref{chow-lemma-agreement-with-loc-chern} 的情形中，设
$\dim_\delta(X) = n$。则
\begin{enumerate}
\item 当 $t = 1, \ldots, r - 1$ 时，
$c_t(Z \to X, i_*\mathcal{O}_Z) \cap [X]_n = 0$；
\item $c_r(Z \to X, i_*\mathcal{O}_Z) \cap [X]_n =
(-1)^{r - 1}(r - 1)![Z]_{n - r}$；
\item 当 $t = 0, \ldots, r - 1$ 时，
$ch_t(Z \to X, i_*\mathcal{O}_Z) \cap [X]_n = 0$；并且
\item $ch_r(Z \to X, i_*\mathcal{O}_Z) \cap [X]_n = [Z]_{n - r}$。
\end{enumerate}
\end{lemma}

\begin{proof}
第 (1)、(2) 项立即由引理 \ref{chow-lemma-agreement-with-loc-chern} 与引理 \ref{chow-lemma-gysin-fundamental} 合并得出。
然后利用引理 \ref{chow-lemma-loc-chern-character} 中
$ch_p = (1/p!)P_p$ 与 $c_p$ 之间的关系推出第 (3)、(4) 项。
（具体而言，若 $c_1 = c_2 = \ldots = c_{r - 1} = 0$，则
$(-1)^{r - 1}(r - 1)!ch_r = c_r$。）
\end{proof}







\section{一个 Adams 算子}
\label{chow-section-adams}

\noindent
我们只做定义第二个 Adams 算子所需的最少工作。
设 $X$ 是概形。回忆 $\textit{Vect}(X)$ 表示有限局部自由
$\mathcal{O}_X$-模的范畴。此外，回忆在
《概形的导出范畴》第 \ref{perfect-section-K-groups} 节中，已经构造了与该范畴对应的第零
$K$-群 $K_0(\textit{Vect}(X))$。最后，
$K_0(\textit{Vect}(X))$ 是环；见
《概形的导出范畴》注 \ref{perfect-remark-K-ring}。

\begin{lemma}
\label{chow-lemma-second-adams-operator}
设 $X$ 是概形。存在环映射
$$
\psi^2 :
K_0(\textit{Vect}(X))
\longrightarrow
K_0(\textit{Vect}(X))
$$
当 $\mathcal{L}$ 可逆时，它将 $[\mathcal{L}]$ 映到
$[\mathcal{L}^{\otimes 2}]$，并且与拉回相容。
\end{lemma}

\begin{proof}
设 $X$ 是概形，并设 $\mathcal{E}$ 是有限局部自由
$\mathcal{O}_X$-模。考虑 $K_0(\textit{Vect}(X))$ 的元素
$$
\psi^2(\mathcal{E}) = [\text{Sym}^2(\mathcal{E})] - [\wedge^2(\mathcal{E})]
$$

\medskip\noindent
设 $X$ 是概形，并考虑有限局部自由 $\mathcal{O}_X$-模的短正合列
$$
0 \to \mathcal{E} \to \mathcal{F} \to \mathcal{G} \to 0
$$
将其视为 $\mathcal{F}$ 上具有 $2$ 步的滤过。
$\text{Sym}^2(\mathcal{F})$ 上诱导的滤过具有 $3$ 步，其分次片为
$\text{Sym}^2(\mathcal{E})$、$\mathcal{E} \otimes \mathcal{F}$
以及 $\text{Sym}^2(\mathcal{G})$。因此
$$
[\text{Sym}^2(\mathcal{F})] =
[\text{Sym}^2(\mathcal{E})] +
[\mathcal{E} \otimes \mathcal{F}] +
[\text{Sym}^2(\mathcal{G})]
$$
完全类似地可证
$$
[\wedge^2(\mathcal{F})] =
[\wedge^2(\mathcal{E})] +
[\mathcal{E} \otimes \mathcal{F}] +
[\wedge^2(\mathcal{G})]
$$
由此
$\psi^2(\mathcal{F}) = \psi^2(\mathcal{E}) + \psi^2(\mathcal{G})$。
因而得到良定义的加法映射
$\psi^2 : K_0(\textit{Vect}(X)) \to K_0(\textit{Vect}(X))$。

\medskip\noindent
显然，该映射与拉回交换。

\medskip\noindent
还须证明 $\psi^2$ 是环映射。设 $X$ 是概形，且
$\mathcal{E}$、$\mathcal{F}$ 是有限局部自由
$\mathcal{O}_X$-模。注意存在短正合列
$$
0 \to \wedge^2(\mathcal{E}) \otimes \wedge^2(\mathcal{F}) \to
\text{Sym}^2(\mathcal{E} \otimes \mathcal{F}) \to
\text{Sym}^2(\mathcal{E}) \otimes \text{Sym}^2(\mathcal{F}) \to 0
$$
其中第一个映射将
$(e \wedge e') \otimes (f \wedge f')$ 映到
$(e \otimes f)(e' \otimes f') - (e' \otimes f)(e \otimes f')$，
第二个映射将 $(e \otimes f) (e' \otimes f')$ 映到
$ee' \otimes ff'$。类似地，存在短正合列
$$
0 \to \text{Sym}^2(\mathcal{E}) \otimes \wedge^2(\mathcal{F}) \to
\wedge^2(\mathcal{E} \otimes \mathcal{F}) \to
\wedge^2(\mathcal{E}) \otimes \text{Sym}^2(\mathcal{F}) \to 0
$$
其中第一个映射将 $e e' \otimes f \wedge f'$ 映到
$(e \otimes f) \wedge (e' \otimes f') + (e' \otimes f) \wedge (e \otimes f')$，
第二个映射将
$(e \otimes f) \wedge (e' \otimes f')$ 映到
$(e \wedge e') \otimes (f f')$。
如上可知映射 $\psi^2$ 具有乘法性。又显然
$\psi^2(1) = 1$，证明完毕。
\end{proof}

\begin{remark}
\label{chow-remark-adams-derived}
设 $X$ 是概形，并且 $2$ 在 $X$ 上可逆。
则可以直接在 $K$-群
$K_0(X) = K_0(D_{perf}(\mathcal{O}_X))$
上定义 Adams 算子 $\psi^2$（《概形的导出范畴》定义 \ref{perfect-definition-K-group}）。
具体而言，给定 $X$ 上的完美复形 $L$，交换因子便给出群
$\{\pm 1\}$ 在 $L \otimes^\mathbf{L} L$ 上的作用。于是可以定义
$$
\psi^2(L) = [(L \otimes^\mathbf{L} L)^+] -
[(L \otimes^\mathbf{L} L)^-]
$$
其中 $(-)^+$ 表示取不变量，$(-)^-$ 表示取反不变量
（作适当定义）。利用取不变量和反不变量的正合性，可以仿照
引理 \ref{chow-lemma-second-adams-operator} 的证明说明它良定义。
当 $2$ 在 $X$ 上不可逆时，情形复杂得多，必须采用另一种方法。
\end{remark}

\begin{lemma}
\label{chow-lemma-minus-adams-operator}
设 $X$ 是概形。存在环映射
$\psi^{-1} : K_0(\textit{Vect}(X)) \to K_0(\textit{Vect}(X))$；
当 $\mathcal{E}$ 有限局部自由时，它将 $[\mathcal{E}]$ 映到
$[\mathcal{E}^\vee]$，并与拉回相容。
\end{lemma}

\begin{proof}
只须检验取对偶与短正合列及拉回相容。这是清楚的。
\end{proof}

\begin{remark}
\label{chow-remark-chern-classes-K}
设 $(S, \delta)$ 如情形 \ref{chow-situation-setup}，且
$X$ 在 $S$ 上局部有限型。陈类映射定义典范映射
$$
c : K_0(\textit{Vect}(X)) \longrightarrow \prod\nolimits_{i \geq 0} A^i(X)
$$
方法是将左边的生成元 $[\mathcal{E}]$ 映到
$c(\mathcal{E}) = 1 + c_1(\mathcal{E}) + c_2(\mathcal{E}) + \ldots$，
再作乘法延拓。因此 $-[\mathcal{E}]$ 被映到形式逆元
$c(\mathcal{E})^{-1}$；这也解释了右边为何取无限乘积。
由引理 \ref{chow-lemma-additivity-chern-classes}，该映射良定义。
\end{remark}

\begin{remark}
\label{chow-remark-chern-character-K}
设 $(S, \delta)$ 如情形 \ref{chow-situation-setup}，且
$X$ 在 $S$ 上局部有限型。陈特征映射定义典范环映射
$$
ch : K_0(\textit{Vect}(X)) \longrightarrow
\prod\nolimits_{i \geq 0} A^i(X) \otimes \mathbf{Q}
$$
方法是将左边的生成元 $[\mathcal{E}]$ 映到
$ch(\mathcal{E})$，再作加法延拓。由引理 \ref{chow-lemma-chern-character-additive}，该映射良定义；由
引理 \ref{chow-lemma-chern-character-multiplicative}，它是环同态。
\end{remark}

\begin{lemma}
\label{chow-lemma-adams-and-chern}
设 $(S, \delta)$ 如情形 \ref{chow-situation-setup}，且
$X$ 在 $S$ 上局部有限型。设 $\psi^2$ 如引理 \ref{chow-lemma-second-adams-operator}，而 $c$ 与 $ch$ 如注 \ref{chow-remark-chern-classes-K} 和 \ref{chow-remark-chern-character-K}。
则对所有 $\alpha \in K_0(\textit{Vect}(X))$，有
$c_i(\psi^2(\alpha)) = 2^i c_i(\alpha)$ 以及
$ch_i(\psi^2(\alpha)) = 2^i ch_i(\alpha)$。
\end{lemma}

\begin{proof}
注意，在 $A^i(X)$ 上乘以 $2^i$ 的映射
$\prod_{i \geq 0} A^i(X) \to \prod_{i \geq 0} A^i(X)$
是环映射。由于 $\psi^2$ 也是环映射，只须对
$K_0(\textit{Vect}(X))$ 的加法生成元证明这些公式。
因此可假设 $\alpha = [\mathcal{E}]$，其中 $\mathcal{E}$
是某个有限局部自由 $\mathcal{O}_X$-模。
根据 $\mathcal{E}$ 的陈类构造，立即化归为
$\mathcal{E}$ 的秩恒为 $r$ 的情形；见注 \ref{chow-remark-extend-to-finite-locally-free}。
此时可选取投影光滑态射 $p : P \to X$，使得限制
$A^*(X) \to A^*(P)$ 是单射，并且 $p^*\mathcal{E}$
有有限滤过，其各分次片是可逆 $\mathcal{O}_P$-模
$\mathcal{L}_j$；见引理 \ref{chow-lemma-splitting-principle}。于是
$[p^*\mathcal{E}] = \sum [\mathcal{L}_j]$，从而按
$\psi^2$ 的定义，
$\psi^2([p^\mathcal{E}]) = \sum [\mathcal{L}_j^{\otimes 2}]$。
令 $x_j  = c_1(\mathcal{L}_j)$，则在 $\prod A^i(P)$ 中有
$$
c(\alpha) = \prod (1 + x_j)
\quad\text{且}\quad
c(\psi^2(\alpha)) = \prod (1 + 2 x_j)
$$
并且在 $\prod A^i(P)$ 中有
$$
ch(\alpha) = \sum \exp(x_j)
\quad\text{且}\quad
ch(\psi^2(\alpha)) = \sum \exp(2 x_j)
$$
所需结论由这些公式得出。
\end{proof}

\begin{remark}
\label{chow-remark-perf-Z-cohomology-K}
设 $X$ 是局部 Noether 概形，且 $Z \subset X$ 是闭子概形。
考虑严格满、饱和的三角子范畴
$$
D_{Z, perf}(\mathcal{O}_X) \subset D(\mathcal{O}_X)
$$
其对象是由 $\mathcal{O}_X$-模组成的完美复形，并且其上同调层
在集合论意义下支撑在 $Z$ 上。记
$\textit{Coh}_Z(X) \subset \textit{Coh}(X)$ 为凝聚
$\mathcal{O}_X$-模的 Serre 子范畴，其中各模的集合论支撑包含在
$Z$ 中。注意，给定
$E \in D_{Z, perf}(\mathcal{O}_X)$，在 $X$ 上 Zariski 局部，
只有有限多个上同调层 $H^i(E)$ 非零
（并且它们在集合论意义下都支撑在 $Z$ 上）。因此可按规则
$$
E \longmapsto
[\bigoplus\nolimits_{i \in \mathbf{Z}} H^{2i}(E)] -
[\bigoplus\nolimits_{i \in \mathbf{Z}} H^{2i + 1}(E)]
$$
定义映射
$$
K_0(D_{Z, perf}(\mathcal{O}_X))
\longrightarrow
K_0(\textit{Coh}_Z(X)) = K'_0(Z)
$$
（等号由引理 \ref{chow-lemma-K-coherent-supported-on-closed} 得出）。
这是可行的，因为 $D_{Z, perf}(\mathcal{O}_X)$ 中每个可区别三角
都给出上同调层的长正合列。
\end{remark}

\begin{remark}
\label{chow-remark-perf-Z-regular}
设 $X$、$Z$、$D_{Z, perf}(\mathcal{O}_X)$ 如注 \ref{chow-remark-perf-Z-cohomology-K}，并假设 $X$ 正则。
则存在典范映射
$$
K_0(\textit{Coh}(Z)) \longrightarrow K_0(D_{Z, perf}(\mathcal{O}_X))
$$
定义如下。对任意凝聚 $\mathcal{O}_Z$-模 $\mathcal{F}$，
以 $\mathcal{F}[0]$ 表示 $D(\mathcal{O}_X)$ 中次数 $0$
处为 $\mathcal{F}$、其他次数均为零的对象。
由《概形的导出范畴》引理 \ref{perfect-lemma-perfect-on-regular}，$\mathcal{F}[0]$
是 $X$ 上的完美复形。因此 $\mathcal{F}[0]$ 是
$D_{Z, perf}(\mathcal{O}_X)$ 的对象。另一方面，给定凝聚
$\mathcal{O}_Z$-模的短正合列
$0 \to \mathcal{F} \to \mathcal{F}' \to \mathcal{F}'' \to 0$，
得到可区别三角
$\mathcal{F}[0] \to \mathcal{F}'[0] \to \mathcal{F}''[0] \to \mathcal{F}[1]$；
见《导出范畴》第 \ref{derived-section-canonical-delta-functor} 节。这表明，将
$[\mathcal{F}]$ 映到 $[\mathcal{F}[0]]$ 定义了映射
$K_0(\textit{Coh}(Z)) \to K_0(D_{Z, perf}(\mathcal{O}_X))$。
为这个可怕的记号致歉。
\end{remark}

\begin{lemma}
\label{chow-lemma-perf-Z-regular}
设 $X$ 是 Noether 正则概形，且 $Z \subset X$ 是闭子概形。
注 \ref{chow-remark-perf-Z-cohomology-K} 和
\ref{chow-remark-perf-Z-regular} 中构造的映射互为逆映射，从而
$K'_0(Z) = K_0(D_{Z, perf}(\mathcal{O}_X))$。
\end{lemma}

\begin{proof}
显然，复合
$$
K_0(\textit{Coh}(Z)) \longrightarrow
K_0(D_{Z, perf}(\mathcal{O}_X)) \longrightarrow
K_0(\textit{Coh}(Z))
$$
是恒等映射。因此只须证明第一个箭头为满射。
设 $E$ 是 $D_{Z, perf}(\mathcal{O}_X)$ 的对象。
回忆由《概形的导出范畴》引理 \ref{perfect-lemma-perfect-on-regular}，
$D_{perf}(\mathcal{O}_X) = D^b_{\textit{Coh}}(\mathcal{O}_X)$。
因此各上同调 $H^i(E)$ 都凝聚，可视为
$D_{Z, perf}(\mathcal{O}_X)$ 的对象，并且仅有有限多个非零。
利用典范截断的可区别三角可见
$$
[E] = \sum (-1)^i[H^i(E)[0]]
$$
在 $K_0(D_{Z, perf}(\mathcal{O}_X))$ 中成立。
于是只须对任意在集合论意义下支撑在 $Z$ 上的凝聚
$\mathcal{O}_X$-模证明 $[\mathcal{F}[0]]$ 位于该映射的像中。
由于可在 $\mathcal{F}$ 上找到有限滤过，其各次商都是
$\mathcal{O}_Z$-模，证明完毕。
\end{proof}

\begin{remark}
\label{chow-remark-localized-chern-classes-K}
设 $(S, \delta)$ 如情形 \ref{chow-situation-setup}。
设 $X$ 在 $S$ 上局部有限型，$Z \subset X$ 是闭子概形，并设
$D_{Z, perf}(\mathcal{O}_X)$ 如注 \ref{chow-remark-perf-Z-cohomology-K}。若 $X$ 拟紧
（或更一般地，$X$ 的不可约分支拟紧），则局部化陈类定义典范映射
$$
c(Z \to X, -) : K_0(D_{Z, perf}(\mathcal{O}_X)) \longrightarrow
A^0(X) \times \prod\nolimits_{i \geq 1} A^i(Z \to X)
$$
方法是将左边的生成元 $[E]$ 映到
$$
c(Z \to X, E) = 1 + c_1(Z \to X, E) + c_2(Z \to X, E) + \ldots
$$
再作乘法延拓（右边的乘积如注 \ref{chow-remark-ring-loc-classes}）。$X$ 的拟紧条件保证
局部化陈类已有定义（情形 \ref{chow-situation-loc-chern} 和
定义 \ref{chow-definition-localized-chern}），并且这些局部化陈类
将可区别三角转化为双变 Chow 环中的相应乘积
（引理 \ref{chow-lemma-additivity-loc-chern-c}）。
\end{remark}

\begin{remark}
\label{chow-remark-localized-chern-character-K}
设 $(S, \delta)$ 如情形 \ref{chow-situation-setup}。
设 $X$ 在 $S$ 上局部有限型，$Z \subset X$ 是闭子概形，并设
$D_{Z, perf}(\mathcal{O}_X)$ 如注 \ref{chow-remark-perf-Z-cohomology-K}。若 $X$ 的不可约分支拟紧，
则局部化陈特征定义典范的加法且乘法映射
$$
ch(Z \to X, -) : K_0(D_{Z, perf}(\mathcal{O}_X)) \longrightarrow
\prod\nolimits_{i \geq 0} A^i(Z \to X) \otimes \mathbf{Q}
$$
方法是将左边的生成元 $[E]$ 映到
$ch(Z \to X, E)$，再作加法延拓。具体而言，$X$
的不可约分支满足的条件保证局部化陈特征已有定义
（情形 \ref{chow-situation-loc-chern} 和定义 \ref{chow-definition-localized-chern}），并且这些局部化陈特征
将可区别三角转化为双变 Chow 环中的相应和
（引理 \ref{chow-lemma-additivity-loc-chern-P}）。
$K_0(D_{Z, perf}(X))$ 上的乘法用导出张量积定义
（《概形的导出范畴》注 \ref{perfect-remark-perf-Z}），所以由引理 \ref{chow-lemma-loc-chern-tensor-product}，
$ch(Z \to X, \alpha \beta) = ch(Z \to X, \alpha) ch(Z \to X, \beta)$。
若 $X$ 拟紧，则映射 $ch(Z \to X, -)$ 的像包含在
$A^*(Z \to X) \otimes \mathbf{Q}$ 中；省略细节。
\end{remark}

\begin{remark}
\label{chow-remark-chern-classes-agree}
设 $(S, \delta)$ 如情形 \ref{chow-situation-setup}。
设 $X$ 在 $S$ 上局部有限型，并假设 $X$ 拟紧
（或更一般地，$X$ 的不可约分支拟紧）。令 $Z = X$，并采用
注 \ref{chow-remark-localized-chern-classes-K} 和
\ref{chow-remark-localized-chern-character-K} 中的记号，则
$D_{Z, perf}(\mathcal{O}_X) = D_{perf}(\mathcal{O}_X)$，且
$$
K_0(D_{Z, perf}(\mathcal{O}_X)) = K_0(D_{perf}(\mathcal{O}_X)) = K_0(X)
$$
见《概形的导出范畴》定义 \ref{perfect-definition-K-group}。因此，作为注 \ref{chow-remark-localized-chern-classes-K} 和
\ref{chow-remark-localized-chern-character-K} 的特殊情形，得到
$$
c : K_0(X) \to \prod A^i(X)
\quad\text{且}\quad
ch : K_0(X) \to \prod A^i(X) \otimes \mathbf{Q}
$$
当然，也可以直接使用定义 \ref{chow-definition-defined-on-perfect}、引理 \ref{chow-lemma-additivity-on-perfect} 和
\ref{chow-lemma-chern-classes-perfect-tensor-product} 构造这些映射
（立即可见这会产生相同的类）。
回忆存在典范映射 $K_0(\textit{Vect}(X)) \to K_0(X)$，
它将有限局部自由模映到作为完美复形的自身
（置于次数 $0$）；见《概形的导出范畴》第 \ref{perfect-section-K-groups} 节。于是图
$$
\xymatrix{
K_0((\textit{Vect}(X)) \ar[rd]_c \ar[rr] & &
K_0(D_{perf}(\mathcal{O}_X)) = K_0(X) \ar[ld]^c \\
& \prod A^i(X)
}
$$
交换，其中东南方向的箭头是在注 \ref{chow-remark-chern-classes-K} 中构造的。类似地，图
$$
\xymatrix{
K_0((\textit{Vect}(X)) \ar[rd]_{ch} \ar[rr] & &
K_0(D_{perf}(\mathcal{O}_X)) = K_0(X) \ar[ld]^{ch} \\
& \prod A^i(X) \otimes \mathbf{Q}
}
$$
交换，其中东南方向的箭头是在注 \ref{chow-remark-chern-character-K} 中构造的。
\end{remark}











\section{重访 Chow 群与 K 群}
\label{chow-section-chow-and-K-II}

\noindent
本节是第 \ref{chow-section-chow-and-K} 节的继续。
设 $(S, \delta)$ 如情形 \ref{chow-situation-setup}，且
$X$ 在 $S$ 上局部有限型。$X$ 上凝聚层的 K 群
$K'_0(X) = K_0(\textit{Coh}(X))$ 具有典范递增滤过
$$
F_kK'_0(X) =
\Im\Big(K_0(\textit{Coh}_{\leq k}(X)) \to K_0(\textit{Coh}(X)\Big)
$$
这称为按支撑维数的滤过。注意
$$
\text{gr}_k K'_0(X) \subset K'_0(X)/F_{k - 1}K'_0(X) =
K_0(\textit{Coh}(X)/\textit{Coh}_{\leq k - 1}(X))
$$
其中等号由《同调代数》引理 \ref{homology-lemma-serre-subcategory-K-groups} 得出。
注 \ref{chow-remark-good-cases-K-A} 中的讨论表明，存在典范映射
$$
\CH_k(X) \longrightarrow \text{gr}_k K'_0(X)
$$
它将 $\delta$-维数为 $k$ 的整闭子概形
$Z \subset X$ 的类映到右边的类 $[\mathcal{O}_Z]$。

\begin{proposition}
\label{chow-proposition-K-tensor-Q}
设 $(S, \delta)$ 如情形 \ref{chow-situation-setup}。
假设给定 $S$ 上局部有限型概形之间的闭浸入
$X \to Y$，其中 $Y$ 正则且拟紧。考虑以下三个映射的复合：
注 \ref{chow-remark-perf-Z-regular} 中的映射
$\mathcal{F} \mapsto \mathcal{F}[0]$，注 \ref{chow-remark-localized-chern-character-K} 中的映射
$ch(X \to Y, -)$，以及映射 $c \mapsto c \cap [Y]$；即
$$
K'_0(X) \to
K_0(D_{X, perf}(\mathcal{O}_Y)) \to
A^*(X \to Y) \otimes \mathbf{Q} \to
\CH_*(X) \otimes \mathbf{Q}
$$
该复合诱导同构
$$
K'_0(X) \otimes \mathbf{Q}
\longrightarrow
\CH_*(X) \otimes \mathbf{Q}
$$
此同构依赖于 $Y$ 的选择。此外，上述典范映射
$$
\CH_k(X) \otimes \mathbf{Q}
\longrightarrow
\text{gr}_k K'_0(X) \otimes \mathbf{Q}
$$
对所有 $k \in \mathbf{Z}$ 都是 $\mathbf{Q}$-向量空间同构。
\end{proposition}

\begin{proof}
由于 $Y$ 正则，可以应用注 \ref{chow-remark-perf-Z-regular} 中的构造。由于 $Y$ 拟紧，
可以应用注 \ref{chow-remark-localized-chern-character-K} 中的构造。
$Y$ 局部等维（引理 \ref{chow-lemma-locally-equidimensional}），
因而“基本周” $[Y]$ 已定义为 $\CH_*(Y)$ 的元素；见注 \ref{chow-remark-fundamental-class}。结合上面构造的映射
$\CH_k(X) \to \text{gr}_kK'_0(X)$，只须证明：
\begin{enumerate}
\item 若 $\mathcal{F}$ 是凝聚 $\mathcal{O}_X$-模，
其支撑的 $\delta$-维数 $\leq k$，则上述复合将
$[\mathcal{F}]$ 映入
$\bigoplus_{k' \leq k} \CH_{k'}(X) \otimes \mathbf{Q}$；
\item 若 $Z \subset X$ 是 $\delta$-维数为 $k$ 的整闭子概形，
则上述复合将 $[\mathcal{O}_Z]$ 映到某个元素，其次数 $k$
部分是 $\CH_k(X) \otimes \mathbf{Q}$ 中的类 $[Z]$。
\end{enumerate}
事实上，若这些断言成立，则我们的映射诱导映射
$\text{gr}_kK'_0(X) \otimes \mathbf{Q} \to CH_k(X) \otimes \mathbf{Q}$，
它们与命题上方给出的典范映射
$\CH_k(X) \otimes \mathbf{Q} \to \text{gr}_k K'_0(X) \otimes \mathbf{Q}$
互为逆映射。

\medskip\noindent
给定凝聚 $\mathcal{O}_X$-模 $\mathcal{F}$，上述复合将
$[\mathcal{F}]$ 映到
$$
ch(X \to Y, \mathcal{F}[0]) \cap [Y] \in \CH_*(X) \otimes \mathbf{Q}
$$
若 $\mathcal{F}$ 在集合论意义下支撑于闭子概形
$Z \subset X$，则由引理 \ref{chow-lemma-loc-chern-shrink-Z}，
$$
ch(X \to Y, \mathcal{F}[0]) = (Z \to X)_* \circ ch(Z \to Y, \mathcal{F}[0])
$$
因此在此情形下，所得元素位于
$\CH_*(Z) \to \CH_*(X)$ 的像中。这证明了条件 (1)。

\medskip\noindent
设 $Z \subset X$ 是 $\delta$-维数为 $k$ 的整闭子概形。
由与上面相同的论证，上述复合将 $[\mathcal{O}_Z]$ 映到元素
$$
ch(X \to Y, \mathcal{O}_Z[0]) \cap [Y] =
(Z \to X)_* ch(Z \to Y, \mathcal{O}_Z[0]) \cap [Y]
$$
因此只须证明
$ch(Z \to Y, \mathcal{O}_Z[0]) \cap [Y] \in
\CH_*(Z) \otimes \mathbf{Q}$ 的次数 $k$ 部分为 $[Z]$。
由于 $\CH_k(Z) = \mathbf{Z}$，为证明此事，可将 $Y$
替换为 $Z$ 的泛点 $\xi$ 的一个开邻域。由于正则局部环
$\mathcal{O}_{X, \xi}$ 的极大理想由正则序列生成
（《交换代数》引理 \ref{algebra-lemma-regular-ring-CM}），
可假设 $Z$ 的理想由正则序列生成；见《除子》引理 \ref{divisors-lemma-Noetherian-scheme-regular-ideal}。
于是结论由引理 \ref{chow-lemma-actual-computation} 得出。
\end{proof}







\section{正则概形上的有理交乘积}
\label{chow-section-intersection-regular}

\noindent
我们将证明：若 $X$ 是具有仿射对角态射的 Noether 正则有限维概形，
则 $\CH_*(X) \otimes \mathbf{Q}$ 上存在交乘积。
该构造以如下结果为基础（它是上一节命题的推论）。

\begin{lemma}
\label{chow-lemma-K-tensor-Q}
设 $(S, \delta)$ 如情形 \ref{chow-situation-setup}。
设 $X$ 是 $S$ 上有限型的拟紧正则概形，具有仿射对角态射，
且 $\delta_{X/S} : X \to \mathbf{Z}$ 有界。则
注 \ref{chow-remark-chern-character-K} 中的映射 $ch$
与映射 $c \mapsto c \cap [X]$ 的复合
$$
K_0(\textit{Vect}(X)) \otimes \mathbf{Q}
\longrightarrow
A^*(X) \otimes \mathbf{Q}
\longrightarrow
\CH_*(X) \otimes \mathbf{Q}
$$
是同构。
\end{lemma}

\begin{proof}
由《概形的导出范畴》引理 \ref{perfect-lemma-Kprime-K}、
\ref{perfect-lemma-regular-resolution-property} 和
\ref{perfect-lemma-K-is-old-K}，有
$K'_0(X) = K_0(X) = K_0(\textit{Vect}(X))$。
由注 \ref{chow-remark-chern-classes-agree}，当 $X = Y$ 时，
给定复合与命题 \ref{chow-proposition-K-tensor-Q} 中的映射一致。
故结论由该命题得出。
\end{proof}

\noindent
设 $X, S, \delta$ 如引理 \ref{chow-lemma-K-tensor-Q}。
为简单起见，考虑给定余维的周；见第 \ref{chow-section-cycles-codimension} 节。令 $[X]$ 为 $X$ 的基本周；
见注 \ref{chow-remark-fundamental-class}。
选取 $\alpha \in CH^i(X)$ 和 $\beta \in CH^j(X)$。
由该引理，存在唯一的
$\alpha' \in K_0(\textit{Vect}(X)) \otimes \mathbf{Q}$，使得
$ch(\alpha') \cap [X]  = \alpha$。这当然意味着：若
$i' \not = i$，则 $ch_{i'}(\alpha') \cap [X] = 0$；并且
$ch_i(\alpha') \cap [X] = \alpha$。
由引理 \ref{chow-lemma-adams-and-chern}，
$\alpha'' = 2^{-i}\psi^2(\alpha')$ 是另一个解。
由唯一性，$\alpha'' = \alpha'$，从而当 $i' \not = i$ 时，
$ch_{i'}(\alpha) = 0$ 在
$A^{i'}(X) \otimes \mathbf{Q}$ 中成立。于是可以定义
$$
\alpha \cdot \beta = ch(\alpha') \cap \beta =
ch_i(\alpha') \cap \beta
$$
作为 $\CH^{i + j}(X) \otimes \mathbf{Q}$ 的元素；
这是由上面观察到的 $\alpha'$ 的性质保证的。
该配对是对称的：事实上，若选取提升 $\beta$ 的
$\beta' \in K_0(\textit{Vect}(X)) \otimes \mathbf{Q}$，则
$$
\alpha \cdot \beta = ch(\alpha') \cap \beta =
ch(\alpha') \cap ch(\beta') \cap [X]
$$
而我们知道陈类彼此交换。出于同样理由，交乘积满足结合律：
$$
(\alpha \cdot \beta) \cdot \gamma =
ch(\alpha') \cap ch(\beta') \cap ch(\gamma') \cap [X]
$$
因为双变类的复合满足结合律。也许更好的表述如下：
$\CH^*(X) \otimes \mathbf{Q}$ 上存在唯一的、与分次相容的交换结合交乘积，
使得同构
$K_0(\textit{Vect}(X)) \otimes \mathbf{Q} \to \CH^*(X) \otimes \mathbf{Q}$
是环同构。






\section{局部完全交态射的 Gysin 映射}
\label{chow-section-koszul}

\noindent
阅读本节前，建议读者先了解正则浸入
（《除子》第 \ref{divisors-section-regular-immersions} 节）
以及局部完全交态射
（《态射进阶》第 \ref{more-morphisms-section-lci} 节）。

\medskip\noindent
设 $(S, \delta)$ 如情形 \ref{chow-situation-setup}。
设 $i : X \to Y$ 是 $S$ 上局部有限型概形之间的正则浸入\footnote{见
《除子》定义 \ref{divisors-definition-regular-immersion}。
请注意，对于局部 Noether 概形，正则浸入与
Koszul-正则浸入或拟正则浸入是同一概念；见
《除子》引理 \ref{divisors-lemma-regular-immersion-noetherian}。
本节中将不再另行说明地使用这一事实。}。
特别地，余法层 $\mathcal{C}_{X/Y}$ 是有限局部自由的
（见《除子》引理 \ref{divisors-lemma-quasi-regular-immersion}）。
因此法层
$$
\mathcal{N}_{X/Y} = \SheafHom_{\mathcal{O}_X}(\mathcal{C}_{X/Y}, \mathcal{O}_X)
$$
也是有限局部自由的，并且存在满射
$\mathcal{N}_{X/Y}^\vee \to \mathcal{C}_{X/Y}$
（因为同构当然也是满射）。
第 \ref{chow-section-gysin-higher-codimension} 节中的构造给出规范双变类
$$
i^! = c(X \to Y, \mathcal{N}_{X/Y}) \in A^*(X \to Y)^\wedge
$$
我们需要关于这一概念的几个引理。

\begin{lemma}
\label{chow-lemma-composition-regular-immersion}
设 $(S, \delta)$ 如情形 \ref{chow-situation-setup}。
设 $i : X \to Y$ 和 $j : Y \to Z$ 是 $S$ 上局部有限型概形之间的正则浸入。
则 $j \circ i$ 是正则浸入，并且
$(j \circ i)^! = i^! \circ j^!$。
\end{lemma}

\begin{proof}
第一项断言即《除子》引理 \ref{divisors-lemma-composition-regular-immersion}。
由《除子》引理 \ref{divisors-lemma-transitivity-conormal-quasi-regular}，
存在短正合列
$$
0 \to
i^*(\mathcal{C}_{Y/Z}) \to
\mathcal{C}_{X/Z} \to
\mathcal{C}_{X/Y} \to 0
$$
因此结论由更一般的引理 \ref{chow-lemma-gysin-composition} 得出。
\end{proof}

\begin{lemma}
\label{chow-lemma-section-smooth}
设 $(S, \delta)$ 如情形 \ref{chow-situation-setup}。
设 $p : P \to X$ 是 $S$ 上局部有限型概形之间的光滑态射，
且 $s : X \to P$ 是一个截面。则 $s$ 是正则浸入，并且在
$A^*(X)^\wedge$ 中有 $1 = s^! \circ p^*$，其中
$p^* \in A^*(P \to X)^\wedge$ 是引理 \ref{chow-lemma-flat-pullback-bivariant} 的双变类。
\end{lemma}

\begin{proof}
第一项断言即《除子》引理 \ref{divisors-lemma-section-smooth-regular-immersion}。
由于这些假设在沿任意局部有限型态射 $X' \to X$ 作基变换时保持不变，
只需对任意整数闭子概形 $Z \subset X$ 证明
$\CH_*(X)$ 中 $s^! \cap p^*[Z] = [Z]$。
将 $P$ 替换为 $s(Z)$ 的一个开邻域后，可以假设
$P \to X$ 是固定相对维数 $r$ 的光滑态射。
记 $\dim_\delta(Z) = n$。则 $p^{-1}(Z)$ 的每个不可约分支的维数均为
$r + n$，且 $p^*[Z]$ 由 $[p^{-1}(Z)]_{n + r}$ 给出。
请注意，按概形论意义有 $s(X) \cap p^{-1}(Z) = s(Z)$。
因此，由上面所用的同一文献结果，$s(X) \cap p^{-1}(Z)$ 是
$\overline{p}^{-1}(Z)$ 中余维数为 $r$ 的正则嵌入闭子概形。
结论由引理 \ref{chow-lemma-gysin-fundamental} 得出。
\end{proof}

\noindent
设 $(S, \delta)$ 如情形 \ref{chow-situation-setup}。考虑 $S$ 上局部有限型概形的交换图
$$
\xymatrix{
X \ar[rd]_f \ar[rr]_i & & P \ar[ld]^g \\
& Y
}
$$
其中 $g$ 光滑而 $i$ 是正则浸入。将上面讨论的双变类 $i^!$ 与
引理 \ref{chow-lemma-flat-pullback-bivariant} 的双变类
$g^* \in A^*(P \to Y)^\wedge$ 复合，得到
$$
f^! = i^! \circ g^* \in A^*(X \to Y)
$$
请注意，态射 $f$ 是局部完全交态射；见《态射进阶》定义 \ref{more-morphisms-definition-lci}。
反之，若 $f : X \to Y$ 是局部 Noether 概形之间的局部完全交态射，
且 $f = g \circ i$、$g$ 光滑，则 $i$ 是正则浸入。
我们断言，对 $f^!$ 的这一构造仅依赖于态射 $f$，
而不依赖于分解 $f = g \circ i$ 的选择。

\begin{lemma}
\label{chow-lemma-lci-gysin-well-defined}
设 $(S, \delta)$ 如情形 \ref{chow-situation-setup}。
设 $f : X \to Y$ 是 $S$ 上局部有限型概形之间的局部完全交态射。
双变类 $f^!$ 与满足 $g$ 光滑的分解 $f = g \circ i$ 的选择无关
（只要这种分解存在）。
\end{lemma}

\begin{proof}
给定另一个这样的分解 $f = g' \circ i'$，可以考虑光滑态射
$g'' : P \times_Y P' \to Y$、浸入
$i'' : X \to P \times_Y P'$ 以及分解 $f = g'' \circ i''$。
因此可以假设有图
$$
\xymatrix{
& P' \ar[d]^p \ar[rd]^{g'} \\
X \ar[r]^i \ar[ru]^{i'} & P \ar[r]^g & Y
}
$$
其中 $p$ 是光滑态射。于是 $(g')^* = p^* \circ g^*$
（引理 \ref{chow-lemma-compose-flat-pullback}），故只需在
$A^*(X \to P)$ 中证明 $i^! = (i')^! \circ p^*$。
考虑交换图
$$
\xymatrix{
& X \times_P P' \ar[d]^{\overline{p}} \ar[r]_j & P' \ar[d]^p \\
X \ar[ru]^s \ar[r]^1 & X \ar[r]^i & P
}
$$
其中 $s =(1, i')$。则 $s$ 和 $j$ 都是正则浸入
（分别由《除子》引理 \ref{divisors-lemma-section-smooth-regular-immersion}
和《除子》引理 \ref{divisors-lemma-flat-base-change-regular-immersion}），
且 $i' = j \circ s$。由引理 \ref{chow-lemma-composition-regular-immersion}，
有 $(i')^! = s^! \circ j^!$。
由于该方块是笛卡尔的，双变类 $j^!$ 是 $i^!$ 在 $P'$ 上的限制
（注 \ref{chow-remark-restriction-bivariant}）；见
引理 \ref{chow-lemma-construction-gysin}。
由于双变类与平坦拉回可交换，得到
$j^! \circ p^* = \overline{p}^* \circ i^!$。
因此只需证明 $s^! \circ \overline{p}^* = \text{id}$，
这已由引理 \ref{chow-lemma-section-smooth} 完成。
\end{proof}

\begin{definition}
\label{chow-definition-lci-gysin}
设 $(S, \delta)$ 如情形 \ref{chow-situation-setup}。
设 $f : X \to Y$ 是 $S$ 上局部有限型概形之间的局部完全交态射。
若可以写成 $f = g \circ i$，其中 $g$ 光滑而 $i$ 是浸入，
则称{\it $f$ 的 Gysin 映射存在}。
在这种情形下，如上定义{\it Gysin 映射}
$f^! = i^! \circ g^* \in A^*(X \to Y)$。
\end{definition}

\noindent
由定义可知：对于正则浸入，该定义与前面的构造一致；
对于光滑态射，它与平坦拉回一致。事实上，
对所有 syntomic 态射，这种一致性都成立。

\begin{lemma}
\label{chow-lemma-lci-gysin-flat}
设 $(S, \delta)$ 如情形 \ref{chow-situation-setup}。
设 $f : X \to Y$ 是 $S$ 上局部有限型概形之间的局部完全交态射。
若 $f$ 的 Gysin 映射存在且 $f$ 平坦，则 $f^!$ 等于
引理 \ref{chow-lemma-flat-pullback-bivariant} 的双变类。
\end{lemma}

\begin{proof}
选取分解 $f = g \circ i$，其中 $i : X \to P$ 是浸入，
$g : P \to Y$ 光滑。请注意，对于任意局部有限型态射 $Y' \to Y$，
基变换 $f'$、$g'$、$i'$ 满足相同假设
（见《概形态射》引理 \ref{morphisms-lemma-base-change-smooth}、
\ref{morphisms-lemma-base-change-syntomic} 以及《态射进阶》引理 \ref{more-morphisms-lemma-flat-lci}）。
因此，由引理 \ref{chow-lemma-bivariant-zero}，归结为当 $Y$ 整时证明
$f^*[Y] = i^!(g^*[Y])$。令 $n = \dim_\delta(Y)$。
将 $X$ 和 $P$ 分解为连通分支后，可以假设 $f$ 是相对维数为 $r$ 的平坦态射，
而 $g$ 是相对维数为 $t$ 的光滑态射。
此时 $f^*[Y] = [X]_{n + s}$ 且 $g^*[Y] = [P]_{n + t}$。
另一方面，$i$ 是余维数为 $t - s$ 的正则浸入。
因此 $i^![P]_{n + t} = [X]_{n + s}$
（引理 \ref{chow-lemma-gysin-fundamental}），证明完成。
\end{proof}

\begin{lemma}
\label{chow-lemma-lci-gysin-composition}
设 $(S, \delta)$ 如情形 \ref{chow-situation-setup}。
设 $f : X \to Y$ 和 $g : Y \to Z$ 是 $S$ 上局部有限型概形之间的局部完全交态射。
假设 $g \circ f$ 和 $g$ 的 Gysin 映射存在。
则 $f$ 的 Gysin 映射存在，并且
$(g \circ f)^! = f^! \circ g^!$。
\end{lemma}

\begin{proof}
由《态射进阶》引理 \ref{more-morphisms-lemma-composition-lci}，$g \circ f$ 是局部完全交态射，
故引理的陈述是有意义的。
若 $X \to P$ 是 $X$ 到某个 $Z$ 上光滑概形 $P$ 中的浸入，
则 $X \to P \times_Z Y$ 是 $X$ 到某个 $Y$ 上光滑概形中的浸入。
这证明了引理的第一项断言。
设 $Y \to P'$ 是 $Y$ 到某个 $Z$ 上光滑概形 $P'$ 中的浸入。
考虑交换图
$$
\xymatrix{
X \ar[r] \ar[d] &
P \times_Z Y \ar[r]_a \ar[ld]^p &
P \times_Z P' \ar[ld]^q \\
Y \ar[r]_b \ar[d] &
P' \ar[ld] \\
Z
}
$$
其中水平箭头是正则浸入，西南方向的箭头光滑，且该方块是笛卡尔的。
由于双变类与平坦拉回可交换，因而
$a^! \circ q^* = p^* \circ b^!$。
将此事实与引理 \ref{chow-lemma-composition-regular-immersion} 和
\ref{chow-lemma-compose-flat-pullback} 结合，读者即可验证引理的陈述成立。
这里省略一个小细节。
\end{proof}

\begin{lemma}
\label{chow-lemma-lci-gysin-commutes}
设 $(S, \delta)$ 如情形 \ref{chow-situation-setup}。
考虑 $S$ 上局部有限型概形的交换图
$$
\xymatrix{
X'' \ar[d] \ar[r] &
X' \ar[d] \ar[r] &
X \ar[d]^f \\
Y'' \ar[r] &
Y' \ar[r] &
Y
}
$$
其中两个方块都是笛卡尔的。假设 $f : X \to Y$ 是局部完全交态射，
且 $f$ 的 Gysin 映射存在。令 $c \in A^*(Y'' \to Y')$。
以 $res(f^!) \in A^*(X' \to Y')$ 表示 $f^!$ 在 $Y'$ 上的限制
（注 \ref{chow-remark-restriction-bivariant}）。
则 $c$ 与 $res(f^!)$ 可交换（注 \ref{chow-remark-bivariant-commute}）。
\end{lemma}

\begin{proof}
选取分解 $f = g \circ i$，其中 $g$ 光滑而 $i$ 是浸入。
由于 $f^! = i^! \circ g^!$，只需分别对 $g^!$（它由平坦拉回给出）
和 $i^!$ 证明该引理。关于平坦拉回的结论是双变类定义的一部分；
关于 $i^!$ 的情形立即由引理 \ref{chow-lemma-gysin-commutes} 得出。
\end{proof}

\begin{lemma}
\label{chow-lemma-lci-gysin-easy}
设 $(S, \delta)$ 如情形 \ref{chow-situation-setup}。
考虑 $S$ 上局部有限型概形的笛卡尔图
$$
\xymatrix{
X' \ar[d]_{f'} \ar[r] &
X \ar[d]^f \\
Y' \ar[r] &
Y
}
$$
假设
\begin{enumerate}
\item $f$ 是局部完全交态射，且 $f$ 的 Gysin 映射存在；
\item $X$、$X'$、$Y$、$Y'$ 满足引理 \ref{chow-lemma-locally-equidimensional} 中的等价条件；
\item 对任意 $x' \in X'$，若其在 $X$、$Y'$、$Y$ 中的像分别为
$x$、$y'$、$y$，则 $n_{x'} - n_{y'} = n_x - n_y$，
其中 $n_{x'}$、$n_x$、$n_{y'}$、$n_y$ 如该引理所定义；以及
\item 对每个一般点 $\xi \in X'$，局部环
$\mathcal{O}_{Y', f'(\xi)}$ 都是 Cohen--Macaulay 的。
\end{enumerate}
则 $f^![Y'] = [X']$，其中 $[Y']$ 和 $[X']$ 如
注 \ref{chow-remark-fundamental-class} 所定义。
\end{lemma}

\begin{proof}
回忆 $n_{x'}$ 是 $\delta(\xi)$ 的共同值，这里 $\xi$ 遍历经过 $x'$ 的
不可约分支的一般点。此外，函数
$x' \mapsto n_{x'}$、$x \mapsto n_x$、$y' \mapsto n_{y'}$ 和
$y \mapsto n_y$ 都是局部常值的。分别以 $X'_n$、$X_n$、$Y'_n$、
$Y_n$ 表示 $X'$、$X$、$Y'$、$Y$ 中相应函数取值为 $n$ 的开闭子概形。
回忆 $[X'] = \sum [X'_n]_n$ 且 $[Y'] = \sum [Y'_n]_n$。
由此显然，为证明该引理，可以将 $X'$ 替换为它的一个连通分支，
并将 $X$、$Y'$、$Y$ 替换为该分支所映入的连通分支。
于是 $X'$、$X$、$Y'$、$Y$ 都是 $\delta$-等维的，
即各自的每个不可约分支都具有相同的 $\delta$-维数。
分别记 $X'$、$X$、$Y'$、$Y$ 的这一共同值为
$n'$、$n$、$m'$、$m$。最后一项假设意味着
$n' - m' = n - m$。

\medskip\noindent
选取分解 $f = g \circ i$，其中 $i : X \to P$ 是浸入，
$g : P \to Y$ 光滑。由于 $X$ 连通，$P \to Y$ 在 $i(X)$ 各点处的相对维数恒定。
因此，将 $P$ 替换为 $i(X)$ 的一个开邻域后，可以假设
$P \to Y$ 具有恒定相对维数，且 $i : X \to P$ 是闭浸入。
以 $g' : Y' \times_Y P \to Y'$ 表示 $g$ 的基变换，
以 $i' : X' \to Y' \times_Y P$ 表示 $i$ 的基变换。
显然 $g^*[Y] = [P]$ 且 $(g')^*[Y'] = [Y' \times_Y P]$。
最后，若 $\xi' \in X'$ 是一般点，则
$\mathcal{O}_{Y' \times_Y P, i'(\xi)}$ 是 Cohen--Macaulay 的。
事实上，局部环映射
$\mathcal{O}_{Y', f'(\xi)} \to \mathcal{O}_{Y' \times_Y P, i'(\xi)}$
平坦且具有正则纤维
（见《交换代数》第 \ref{algebra-section-smooth-overview} 节）；
正则局部环是 Cohen--Macaulay 的
（《交换代数》引理 \ref{algebra-lemma-regular-ring-CM}）；
由假设 (4)，$\mathcal{O}_{Y', f'(\xi)}$ 是 Cohen--Macaulay 的；
于是《交换代数》引理 \ref{algebra-lemma-CM-goes-up} 给出所需结论。
因此归结为下一段所讨论的情形。

\medskip\noindent
假设 $f$ 是正则闭浸入，并且 $X'$、$X$、$Y'$、$Y$ 分别是
$\delta$-维数为 $n'$、$n$、$m'$、$m$ 的 $\delta$-等维概形，且
$m' - n' = m - n$。在这种情形下，结论立即由
引理 \ref{chow-lemma-gysin-easy} 得出。
\end{proof}

\begin{remark}
\label{chow-remark-gysin-chern-classes}
设 $(S, \delta)$ 如情形 \ref{chow-situation-setup}。
设 $f : X \to Y$ 是 $S$ 上局部有限型概形之间的局部完全交态射，
并假设 $f$ 的 Gysin 映射存在。则
$f^! \circ c_i(\mathcal{E}) = c_i(f^*\mathcal{E}) \circ f^!$；
对于 Chern 特征标也有类似结论；见
引理 \ref{chow-lemma-lci-gysin-commutes}。
例如，若 $X$ 和 $Y$ 满足引理 \ref{chow-lemma-locally-equidimensional} 的等价条件，且 $Y$ 是
Cohen--Macaulay 的，则由引理 \ref{chow-lemma-lci-gysin-easy} 有
$f^![Y] = [X]$。在这种情形下还有
$f^!(c_i(\mathcal{E}) \cap [Y]) = c_i(f^*\mathcal{E}) \cap [X]$；
对于 Chern 特征标同样如此。
\end{remark}

\begin{lemma}
\label{chow-lemma-compare-gysin-base-change}
设 $(S, \delta)$ 如情形 \ref{chow-situation-setup}。
考虑 $S$ 上局部有限型概形的笛卡尔方块
$$
\xymatrix{
X' \ar[d]_{f'} \ar[r]_{g'} &
X \ar[d]^f \\
Y' \ar[r]^g &
Y
}
$$
假设
\begin{enumerate}
\item $f$ 和 $f'$ 都是局部完全交态射；以及
\item $f$ 的 Gysin 映射存在。
\end{enumerate}
则 $\mathcal{C} = \Ker(H^{-1}((g')^*\NL_{X/Y}) \to H^{-1}(\NL_{X'/Y'}))$
是有限局部自由 $\mathcal{O}_{X'}$-模，$f'$ 的 Gysin 映射存在，且
$$
res(f^!) = c_{top}(\mathcal{C}^\vee) \circ (f')^!
$$
在 $A^*(X' \to Y')$ 中成立。
\end{lemma}

\begin{proof}
$\mathcal{C}$ 有限局部自由这一事实立即由《代数进阶》引理 \ref{more-algebra-lemma-base-change-lci-bis} 得出。
选取分解 $f = h \circ i$，其中 $h : P \to Y$ 光滑而 $i$ 是浸入。
则可分解 $f' = h' \circ i'$，其中 $h' : P' \to Y'$ 和
$i' : X' \to P'$ 是相应的基变换。图示为
$$
\xymatrix{
X' \ar[r] \ar[d] &
P' \ar[r] \ar[d] &
Y' \ar[d] \\
X \ar[r]^i &
P \ar[r]^h &
Y
}
$$
特别地，$f'$ 的 Gysin 映射存在。由《态射进阶》引理 \ref{more-morphisms-lemma-get-NL}，有
$$
\NL_{X/Y} = \left( \mathcal{C}_{X/P} \to i^*\Omega_{P/Y} \right)
$$
其中 $\mathcal{C}_{X/P}$ 是嵌入 $i$ 的余法层；带撇版本同理。
由《概形态射》引理 \ref{morphisms-lemma-base-change-differentials}，
$\Omega_{P/Y}$ 拉回为 $\Omega_{P'/Y'}$，故
$(g')^*i^*\Omega_{P/Y} = (i')^*\Omega_{P'/Y'}$。
还要回忆，$(g')^*\mathcal{C}_{X/P} \to \mathcal{C}_{X'/P'}$ 是满射；见
《概形态射》引理 \ref{morphisms-lemma-conormal-functorial-flat}。
由此推出，层 $\mathcal{C}$ 规范同构于有限局部自由模之间映射
$(g')^*\mathcal{C}_{X/P} \to \mathcal{C}_{X'/P'}$ 的核。
回忆 $i^!$ 是利用 $\mathcal{N} = \mathcal{C}_{Z/X}^\vee$ 定义的，
$(i')^!$ 亦然。因此应用引理 \ref{chow-lemma-gysin-excess} 得到
$$
res(i^!) = c_{top}(\mathcal{C}^\vee) \circ (i')^!
$$
在 $A^*(X' \to P')$ 中成立。最后，由
$f^! = i^! \circ h^*$、$(f')^! = (i')^! \circ (h')^*$ 以及
$(h')^* = res(h^*)$，即得结论。
\end{proof}

\begin{lemma}[爆破公式]
\label{chow-lemma-blow-up-formula}
设 $(S, \delta)$ 如情形 \ref{chow-situation-setup}。
设 $i : Z \to X$ 是 $S$ 上局部有限型概形之间的正则闭浸入，
$b : X' \to X$ 是以 $Z$ 为中心的爆破。图示为
$$
\xymatrix{
E \ar[r]_j \ar[d]_\pi & X' \ar[d]^b \\
Z \ar[r]^i & X
}
$$
假设 $b$ 的 Gysin 映射存在。则
$$
res(b^!) = c_{top}(\mathcal{F}^\vee) \circ \pi^*
$$
在 $A^*(E \to Z)$ 中成立，其中 $\mathcal{F}$ 是规范映射
$\pi^*\mathcal{C}_{Z/X} \to \mathcal{C}_{E/X'}$ 的核。
\end{lemma}

\begin{proof}
由《代数进阶》引理 \ref{more-algebra-lemma-blowup-regular-sequence}，
态射 $b$ 是局部完全交态射，故该陈述有意义。
由于 $Z \to X$ 是正则浸入（因而更是拟正则浸入），
$\mathcal{C}_{Z/X}$ 有限局部自由，且映射
$\text{Sym}^*(\mathcal{C}_{Z/X}) \to \mathcal{C}_{Z/X, *}$
是同构；见《除子》引理 \ref{divisors-lemma-quasi-regular-immersion}。
由 $E = \text{Proj}(\mathcal{C}_{Z/X, *})$ 可知
$E = \mathbf{P}(\mathcal{C}_{Z/X})$ 是 $Z$ 上的射影空间丛。
因此 $E \to Z$ 光滑，当然也是局部完全交态射。
于是引理 \ref{chow-lemma-compare-gysin-base-change} 适用，并给出
$$
res(b^!) = c_{top}(\mathcal{C}^\vee) \circ \pi^!
$$
其中 $\mathcal{C}$ 如该引理的陈述所定义。
当然，由引理 \ref{chow-lemma-lci-gysin-flat}，有 $\pi^* = \pi^!$。
只需再证明 $\mathcal{F}$ 等于映射
$H^{-1}(j^*\NL_{X'/X}) \to H^{-1}(\NL_{E/Z})$ 的核 $\mathcal{C}$。

\medskip\noindent
由于 $E \to Z$ 光滑，有 $H^{-1}(\NL_{E/Z}) = 0$；见
《态射进阶》引理 \ref{more-morphisms-lemma-NL-smooth}。
因此只需证明 $\mathcal{F}$ 可与 $H^{-1}(j^*\NL_{X'/X})$ 等同。
由《态射进阶》引理 \ref{more-morphisms-lemma-get-triangle-NL} 和
\ref{more-morphisms-lemma-NL-immersion}，有正合列
$$
0 \to H^{-1}(j^*\NL_{X'/X}) \to H^{-1}(\NL_{E/X}) \to
\mathcal{C}_{E/X'} \to \ldots
$$
将同样的引理应用于 $E \to Z \to X$，得到同构
$\pi^*\mathcal{C}_{Z/X} = H^{-1}(\pi^*\NL_{Z/X}) \to H^{-1}(\NL_{E/X})$。
于是结论成立。
\end{proof}

\begin{lemma}
\label{chow-lemma-lci-gysin-product-regular}
设 $(S, \delta)$ 如情形 \ref{chow-situation-setup}。
设 $f : X \to Y$ 是 $S$ 上局部有限型概形之间的态射，
并且 $X$ 和 $Y$ 都拟紧、正则、具有仿射对角态射且维数有限。
则 $f$ 是局部完全交态射。再假设 $f$ 的 Gysin 映射存在，则
$$
f^!(\alpha \cdot \beta) = f^!\alpha \cdot f^!\beta
$$
在 $\CH^*(X) \otimes \mathbf{Q}$ 中成立，其中交乘积如
第 \ref{chow-section-intersection-regular} 节所定义。
\end{lemma}

\begin{proof}
第一项断言由《态射进阶》引理 \ref{more-morphisms-lemma-morphism-regular-schemes-is-lci} 得出。
请注意，由引理 \ref{chow-lemma-lci-gysin-easy}，有 $f^![Y] = [X]$。
如第 \ref{chow-section-intersection-regular} 节，写成
$\alpha = ch(\alpha') \cap [Y]$ 和 $\beta = ch(\beta') \cap [Y]$，
其中 $\alpha', \beta' \in K_0(\textit{Vect}(X)) \otimes \mathbf{Q}$。
令 $c = ch(\alpha')$、$c' = ch(\beta')$；由构造可知
$\alpha \cdot \beta = c \cap c' \cap [Y]$。
由引理 \ref{chow-lemma-lci-gysin-commutes}，$f^!$ 与 $c$、$c'$ 均可交换。因此
\begin{align*}
f^!(\alpha \cdot \beta)
& =
f^!(c \cap c' \cap [Y]) \\
& =
c \cap c' \cap f^![Y] \\
& =
c \cap c' \cap [X] \\
& =
(c \cap [X]) \cdot (c' \cap [X]) \\
& =
(c \cap f^![Y]) \cdot (c' \cap f^![Y]) \\
& =
f^!(\alpha) \cdot f^!(\beta)
\end{align*}
这正是所需结论。
\end{proof}

\begin{lemma}
\label{chow-lemma-projection-formula-regular}
设 $(S, \delta)$ 如情形 \ref{chow-situation-setup}。
设 $f : X \to Y$ 是 $S$ 上局部有限型概形之间的态射，
并且 $X$ 和 $Y$ 都拟紧、正则、具有仿射对角态射且维数有限。
则 $f$ 是局部完全交态射。再假设 $f$ 的 Gysin 映射存在且 $f$ 固有，则
$$
f_*(\alpha \cdot f^!\beta) = f_*\alpha \cdot \beta
$$
在 $\CH^*(Y) \otimes \mathbf{Q}$ 中成立，其中交乘积如
第 \ref{chow-section-intersection-regular} 节所定义。
\end{lemma}

\begin{proof}
第一项断言由《态射进阶》引理 \ref{more-morphisms-lemma-morphism-regular-schemes-is-lci} 得出。
请注意，由引理 \ref{chow-lemma-lci-gysin-easy}，有 $f^![Y] = [X]$。
如第 \ref{chow-section-intersection-regular} 节，写成
$\alpha = ch(\alpha') \cap [X]$ 和 $\beta = ch(\beta') \cap [Y]$，
其中 $\alpha' \in K_0(\textit{Vect}(X)) \otimes \mathbf{Q}$ 且
$\beta' \in K_0(\textit{Vect}(Y)) \otimes \mathbf{Q}$。
令 $c = ch(\alpha')$、$c' = ch(\beta')$。于是
\begin{align*}
f_*(\alpha \cdot f^!\beta)
& =
f_*(c \cap f^!(c' \cap [Y]_e)) \\
& =
f_*(c \cap c' \cap f^![Y]_e) \\
& =
f_*(c \cap c' \cap [X]_d) \\
& =
f_*(c' \cap c \cap [X]_d) \\
& =
c' \cap f_*(c \cap [X]_d) \\
& =
\beta \cdot f_*(\alpha)
\end{align*}
第一项等式来自交乘积的构造。由引理 \ref{chow-lemma-lci-gysin-commutes}，$f^!$ 与 $c'$ 可交换。
Chern 类位于双变环的中心，这使得可以交换以 $c$ 和 $c'$ 对 $[X]$ 作帽积的次序。
由于 $c'$ 是双变类，允许将 $c'$ 与 $f_*$ 交换。最后一项等式再次来自交乘积的构造。
\end{proof}













\section{对角态射的 Gysin 映射}
\label{chow-section-gysin-for-diagonal}

\noindent
设 $(S, \delta)$ 如情形 \ref{chow-situation-setup}。
设 $f : X \to Y$ 是 $S$ 上局部有限型概形之间的光滑态射。
则对角态射 $\Delta : X \longrightarrow X \times_Y X$ 是正则浸入；见
《态射进阶》引理 \ref{more-morphisms-lemma-smooth-diagonal-perfect}。
因此有第 \ref{chow-section-koszul} 节中构造的 Gysin 映射
$$
\Delta^! \in A^*(X \to X \times_Y X)^\wedge
$$
若 $X \to Y$ 具有恒定相对维数 $d$，则
$\Delta^! \in A^d(X \to X \times_Y X)$。

\begin{lemma}
\label{chow-lemma-diagonal-identity}
在上述情形中，$A^0(X)$ 内有 $\Delta^! \circ \text{pr}_i^! = 1$。
\end{lemma}

\begin{proof}
请注意，投影 $\text{pr}_i : X \times_Y X \to X$ 都是光滑的，
因而这些投影也有 Gysin 映射。因此该引理的陈述有意义，
并且它是引理 \ref{chow-lemma-lci-gysin-composition} 的特殊情形。
\end{proof}

\begin{proposition}
\label{chow-proposition-compute-bivariant}
\begin{reference}
\cite[Proposition 17.4.2]{F}
\end{reference}
设 $(S, \delta)$ 如情形 \ref{chow-situation-setup}。
设 $f : X \to Y$ 和 $g : Y \to Z$ 是 $S$ 上局部有限型概形之间的态射。
若 $g$ 是相对维数为 $d$ 的光滑态射，则
$A^p(X \to Y) = A^{p - d}(X \to Z)$。
\end{proposition}

\begin{proof}
我们将使用如下事实：光滑态射是 Gysin 映射存在的局部完全交态射
（见第 \ref{chow-section-koszul} 节）。特别地，有
$g^! \in A^{-d}(Y \to Z)$。于是可以将
$c \in A^p(X \to Y)$ 映到 $c \circ g^! \in A^{p - d}(X \to Z)$。

\medskip\noindent
反之，令 $c' \in A^{p - d}(X \to Z)$。
以 $res(c')$ 表示 $c'$ 沿态射 $Y \to Z$ 的限制
（注 \ref{chow-remark-restriction-bivariant}）。由于图
$$
\xymatrix{
X \times_Z Y \ar[r]_{\text{pr}_2} \ar[d]_{\text{pr}_1} & Y \ar[d]^g \\
X \ar[r]^f & Z
}
$$
是笛卡尔的，得到 $res(c') \in A^{p - d}(X \times_Z Y \to Y)$。
令 $\Delta : Y \to Y \times_Z Y$ 为对角态射，
并以 $res(\Delta^!)$ 表示 $\Delta^!$ 沿态射
$X \times_Z Y \to Y \times_Z Y$ 在 $X \times_Z Y$ 上的限制。
由于图
$$
\xymatrix{
X \ar[r] \ar[d] & X \times_Z Y \ar[d] \\
Y \ar[r]^-\Delta & Y \times_Z Y
}
$$
是笛卡尔的，可见 $res(\Delta^!) \in A^d(X \to X \times_Z Y)$。
将这两个限制复合，得到
$$
res(\Delta^!) \circ res(c') \in A^p(X \to Y)
$$
于是构造了映射 $A^p(X \to Y) \to A^{p - d}(X \to Z)$ 和
$A^{p - d}(X \to Z) \to A^p(X \to Y)$。
为完成证明，将说明这两个映射互为逆映射。

\medskip\noindent
先从 $c \in A^p(X \to Y)$ 开始。考虑图
$$
\xymatrix{
X \ar[d] \ar[r] & Y \ar[d] \\
X \times_Z Y \ar[r] \ar[d]^{\text{pr}_1} &
Y \times_Z Y \ar[r]_{p_2} \ar[d]^{p_1} &
Y \ar[d]^g \\
X \ar[r]^f &
Y \ar[r]^g &
Z
}
$$
其中各方块都是笛卡尔的。该图下方的两个方块表明
$res(c \circ g^!) = res(c) \cap p_2^!$；在此公式中，
$res(c)$ 表示 $c$ 沿 $p_1$ 的限制。
考察该图上方的方块并使用引理 \ref{chow-lemma-lci-gysin-commutes}，得到
$c \circ \Delta^! = res(\Delta^!) \circ res(c)$。计算得
\begin{align*}
res(\Delta^!) \circ res(c \circ g^!)
& =
res(\Delta^!) \circ res(c) \circ p_2^! \\
& =
c \circ \Delta^! \circ p_2^! \\
& =
c
\end{align*}
最后一项等式由引理 \ref{chow-lemma-diagonal-identity} 得出。

\medskip\noindent
反之，从 $c' \in A^{p - d}(X \to Z)$ 开始。
考察上图下方的矩形，得到
$res(c') \circ g^! = \text{pr}_1^! \circ c'$。计算得
\begin{align*}
res(\Delta^!) \circ res(c') \circ g^!
& =
res(\Delta^!) \circ \text{pr}_1^! \circ c' \\
& =
c'
\end{align*}
最后一项等式成立，是因为该图左侧两个方块表明
$\text{id} = res(\Delta^! \circ p_1^!) = res(\Delta^!) \circ \text{pr}_1^!$。
证明完成。
\end{proof}









\section{外积}
\label{chow-section-exterior-product}

\noindent
设 $k$ 是域。本节在 $S = \Spec(k)$ 上工作，其中
$\delta : S \to \mathbf{Z}$ 将唯一的点映到 $0$；见
例 \ref{chow-example-field}。

\medskip\noindent
考虑 $k$ 上局部有限型概形的笛卡尔方块
$$
\xymatrix{
X \times_k Y \ar[r] \ar[d] & Y \ar[d] \\
X \ar[r] & \Spec(k) = S
}
$$
则存在规范映射
$$
\times :
\CH_n(X) \otimes_{\mathbf{Z}} \CH_m(Y)
\longrightarrow
\CH_{n + m}(X \times_k Y)
$$
它由如下规则唯一确定：给定维数分别为 $n$ 和 $m$ 的整闭子概形
$X' \subset X$ 与 $Y' \subset Y$，有
$$
[X'] \times [Y'] = [X' \times_k Y']_{n + m}
$$
在 $\CH_{n + m}(X \times_k Y)$ 中成立。

\begin{lemma}
\label{chow-lemma-exterior-product-well-defined}
映射
$\times : \CH_n(X) \otimes_{\mathbf{Z}} \CH_m(Y) \to \CH_{n + m}(X \times_k Y)$
是良定义的。
\end{lemma}

\begin{proof}
先注意，若 $\alpha = \sum n_i[X_i]$ 且 $\beta = \sum m_j[Y_j]$，
其中 $X_i \subset X$ 和 $Y_j \subset Y$ 分别是维数为 $n$ 和 $m$ 的
整闭子概形的局部有限族，则 $X_i \times_k Y_j$ 构成
$X \times_k Y$ 中维数为 $n + m$ 的闭子概形的局部有限族，因而确实可以将
$$
\alpha \times \beta = \sum n_i m_j [X_i \times_k Y_j]_{n + m}
$$
视为 $X \times_k Y$ 上的一个 $(n + m)$-周。
由此得到加性映射
$\times : Z_n(X) \otimes_{\mathbf{Z}} Z_m(Y) \to Z_{n + m}(X \times_k Y)$。
问题在于证明这一过程与有理等价相容。

\medskip\noindent
设 $i : X' \to X$ 是一个维数为 $n$ 的整闭子概形的包含态射。
由引理 \ref{chow-lemma-flat-pullback-bivariant}，沿态射
$p' : X' \to \Spec(k)$ 的平坦拉回是元素
$(p')^* \in A^{-n}(X' \to \Spec(k))$；因而由
引理 \ref{chow-lemma-push-proper-bivariant}，有
$c' = i_* \circ (p')^* \in A^{-n}(X \to \Spec(k))$。
这给出映射
$$
c' \cap - : \CH_m(Y) \longrightarrow \CH_{m + n}(X \times_k Y)
$$
读者容易验证，对于任意维数为 $m$ 的整闭子概形 $Y' \subset Y$，
该映射把 $[Y']$ 映到 $[X' \times_k Y']_{n + m}$。
因此，构造 $([X'], [Y']) \mapsto [X' \times_k Y']_{n + m}$
在第二个变量中可通过有理等价下降，也就是说，它给出良定义映射
$Z_n(X) \otimes_{\mathbf{Z}} \CH_m(Y) \to \CH_{n + m}(X \times_k Y)$。
由对称性，另一个变量也同样如此，结论得证。
\end{proof}

\begin{lemma}
\label{chow-lemma-chow-cohomology-towards-point}
设 $k$ 是域，$X$ 是 $k$ 上局部有限型概形。
则对所有 $p \in \mathbf{Z}$，都有规范等同
$$
A^p(X \to \Spec(k)) = \CH_{-p}(X)
$$
\end{lemma}

\begin{proof}
考虑元素 $[\Spec(k)] \in \CH_0(\Spec(k))$。
把 $c$ 映到 $c \cap [\Spec(k)]$，即得到映射
$A^p(X \to \Spec(k)) \to \CH_{-p}(X)$。

\medskip\noindent
反之，设有 $\alpha \in \CH_{-p}(X)$。
可按如下方式定义 $c_\alpha \in A^p(X \to \Spec(k))$：
给定 $X' \to \Spec(k)$ 和 $\alpha' \in \CH_n(X')$，令
$$
c_\alpha \cap \alpha' = \alpha \times \alpha'
$$
作为 $\CH_{n - p}(X \times_k X')$ 中的元素。
为证明这是双变类，如定义 \ref{chow-definition-cycles} 写成
$\alpha = \sum n_i[X_i]$。考虑复合
$$
\coprod X_i \xrightarrow{g} X \to \Spec(k)
$$
并以 $f : \coprod X_i \to \Spec(k)$ 表示该复合。
则 $g$ 固有，而 $f$ 是相对维数为 $-p$ 的平坦态射。
由引理 \ref{chow-lemma-flat-pullback-bivariant}，沿 $f$ 的拉回是双变类
$f^* \in A^p(\coprod X_i \to \Spec(k))$。
以 $\nu \in A^0(\coprod X_i)$ 表示如下双变类：
它在第 $i$ 个分支上将周乘以 $n_i$。于是
$\nu \circ f^* \in A^p(\coprod X_i \to X)$。
最后，由引理 \ref{chow-lemma-push-proper-bivariant}，有双变类
$$
g_* \circ \nu \circ f^*
$$
读者容易验证 $c_\alpha$ 等于该双变类，因而它本身也是双变类。

\medskip\noindent
为完成证明，还须说明这两个构造互为逆映射。
由于 $c_\alpha \cap [\Spec(k)] = \alpha$，其中一个方向是显然的。
对另一个方向，令 $c \in A^p(X \to \Spec(k))$，并置
$\alpha = c \cap [\Spec(k)]$。由引理 \ref{chow-lemma-bivariant-zero}，
只需在 $X'$ 是 $\Spec(k)$ 上局部有限型整概形时证明
$$
c \cap [X'] = c_\alpha \cap [X']
$$
但此时 $p' : X' \to \Spec(k)$ 是相对维数为 $\dim(X')$ 的平坦态射，
故 $[X'] = (p')^*[\Spec(k)]$。
因此，双变类 $c$ 和 $c_\alpha$ 在 $[\Spec(k)]$ 上相同，
就蕴含它们与 $[X']$ 作帽积时也相同，证明完成。
\end{proof}

\begin{lemma}
\label{chow-lemma-chow-cohomology-towards-point-commutes}
设 $k$ 是域，$X$ 是 $k$ 上局部有限型概形。
设 $c \in A^p(X \to \Spec(k))$，$Y \to Z$ 是 $k$ 上局部有限型概形之间的态射，
且 $c' \in A^q(Y \to Z)$。则在 $A^{p + q}(X \times_k Y \to Z)$ 中有
$c \circ c' = c' \circ c$。
\end{lemma}

\begin{proof}
在引理 \ref{chow-lemma-chow-cohomology-towards-point} 的证明中已经看到，
$c$ 由固有前推、在各连通分支上乘以整数以及平坦拉回组合而成。
依双变类的定义，$c'$ 与这些运算中的每一个都可交换，故结论成立。
这里省略若干细节。
\end{proof}

\begin{remark}
\label{chow-remark-commuting-exterior}
引理 \ref{chow-lemma-chow-cohomology-towards-point} 和
\ref{chow-lemma-chow-cohomology-towards-point-commutes} 的要点如下。
设 $k$ 是域，$X$ 是 $k$ 上局部有限型概形。
设 $\alpha \in \CH_*(X)$，$Y \to Z$ 是 $k$ 上局部有限型概形之间的态射，
且 $c' \in A^q(Y \to Z)$。则对任意 $\beta \in \CH_*(Z)$，
$$
\alpha \times (c' \cap \beta) = c' \cap (\alpha \times \beta)
$$
在 $\CH_*(X \times_k Y)$ 中成立。事实上，只需取与 $\alpha$ 对应的双变类
$c = c_\alpha \in A^*(X \to \Spec(k))$；见引理 \ref{chow-lemma-chow-cohomology-towards-point} 的证明。
\end{remark}

\begin{lemma}
\label{chow-lemma-exterior-product-associative}
外积满足结合律。更确切地，设 $k$ 是域，
$X, Y, Z$ 是 $k$ 上局部有限型概形，并设
$\alpha \in \CH_*(X)$、$\beta \in \CH_*(Y)$、$\gamma \in \CH_*(Z)$。
则在 $\CH_*(X \times_k Y \times_k Z)$ 中有
$(\alpha \times \beta) \times \gamma =
\alpha \times (\beta \times \gamma)$。
\end{lemma}

\begin{proof}
从略。提示：概形纤维积的结合律。
\end{proof}







\section{交乘积}
\label{chow-section-intersection-product}

\noindent
设 $k$ 是域。本节在 $S = \Spec(k)$ 上工作，其中
$\delta : S \to \mathbf{Z}$ 将唯一的点映到 $0$；见
例 \ref{chow-example-field}。

\medskip\noindent
设 $X$ 是 $k$ 上光滑概形。第 \ref{chow-section-gysin-for-diagonal} 节的双变类
$\Delta^!$ 使我们能够在 $X$ 上局部有限型概形的 Chow 群上定义一种交乘积。
具体地，设 $Y \to X$ 和 $Z \to X$ 是局部有限型概形之间的态射。注意到
$$
Y \times_X Z =  (Y \times_k Z) \times_{X \times_k X, \Delta} X
$$
因此可以考虑映射序列
$$
\CH_n(Y) \otimes_\mathbf{Z} \CH_m(Z)
\xrightarrow{\times}
\CH_{n + m}(Y \times_k Z)
\xrightarrow{\Delta^!}
\CH_{n + m - *}(Y \times_X Z)
$$
其中第一支箭头是第 \ref{chow-section-exterior-product} 节中构造的外积，
第二支箭头是第 \ref{chow-section-gysin-for-diagonal} 节中研究的对角态射 Gysin 映射。
若 $X$ 是维数为 $d$ 的等维概形，则所得元素落在
$\CH_{n + m - d}(Y \times_X Z)$ 中；一般而言，可依照 $X$ 的开闭子概形
（其上 $X$ 具有给定维数），将其分解为位于这些部分之上的分量。
给定 $\alpha \in \CH_*(Y)$ 和 $\beta \in \CH_*(Z)$，记
$$
\alpha \cdot \beta = \Delta^!(\alpha \times \beta)
\in \CH_*(Y \times_X Z)
$$
在特殊情形 $X = Y = Z$ 下，得到乘法
$$
\CH_*(X) \times \CH_*(X) \to \CH_*(X),\quad
(\alpha, \beta) \mapsto \alpha \cdot \beta
$$
称为{\it 交乘积}。显然，这一乘积是对称的。
结合律由下一个引理给出。

\begin{lemma}
\label{chow-lemma-associative}
上面定义的乘积满足结合律。更确切地，设 $k$ 是域，
$X$ 在 $k$ 上光滑，$Y, Z, W$ 是 $X$ 上局部有限型概形，并设
$\alpha \in \CH_*(Y)$、$\beta \in \CH_*(Z)$、$\gamma \in \CH_*(W)$。
则在 $\CH_*(Y \times_X Z \times_X W)$ 中有
$(\alpha \cdot \beta) \cdot \gamma =
\alpha \cdot (\beta \cdot \gamma)$。
\end{lemma}

\begin{proof}
由引理 \ref{chow-lemma-exterior-product-associative}，在
$\CH_*(Y \times_k Z \times_k W)$ 中有
$(\alpha \times \beta) \times \gamma =
\alpha \times (\beta \times \gamma)$。考虑闭浸入
$$
\Delta_{12} : X \times_k X \longrightarrow X \times_k X \times_k X,
\quad (x, x') \mapsto (x, x, x')
$$
以及
$$
\Delta_{23} : X \times_k X \longrightarrow X \times_k X \times_k X,
\quad (x, x') \mapsto (x, x', x')
$$
以 $\Delta_{12}^!$ 和 $\Delta_{23}^!$ 表示相应的双变类。
注意，$\Delta_{12}^!$ 是 $\Delta^!$ 沿映射 $\text{pr}_{12}$ 在
$X \times_k X \times_k X$ 上的限制
（注 \ref{chow-remark-restriction-bivariant}），而
$\Delta_{23}^!$ 是 $\Delta^!$ 沿映射 $\text{pr}_{23}$ 在
$X \times_k X \times_k X$ 上的限制。
因此，$\Delta_{12}^!$ 沿 $\Delta_{23}$ 的限制显然是 $\Delta^!$，
而 $\Delta_{23}^!$ 沿 $\Delta_{12}$ 的限制同样是 $\Delta^!$。
于是由引理 \ref{chow-lemma-gysin-commutes}，有
$$
\Delta^! \circ \Delta_{12}^! =
\Delta^! \circ \Delta_{23}^! 
$$
现在用如下等式链证明该引理：
\begin{align*}
(\alpha \cdot \beta) \cdot \gamma
& =
\Delta^!(\Delta^!(\alpha \times \beta) \times \gamma) \\
& =
\Delta^!(\Delta_{12}^!((\alpha \times \beta) \times \gamma)) \\
& =
\Delta^!(\Delta_{23}^!((\alpha \times \beta) \times \gamma)) \\
& =
\Delta^!(\Delta_{23}^!(\alpha \times (\beta \times \gamma)) \\
& =
\Delta^!(\alpha \times \Delta^!(\beta \times \gamma)) \\
& =
\alpha \cdot (\beta \cdot \gamma)
\end{align*}
除第二项和倒数第二项等式或许需要说明外，其余等式均由上文显然成立。
等式
$\Delta_{23}^!(\alpha \times (\beta \times \gamma)) =
\alpha \times \Delta^!(\beta \times \gamma)$
由注 \ref{chow-remark-commuting-exterior} 成立；第二项等式同理。
\end{proof}

\begin{lemma}
\label{chow-lemma-identify-chow-for-smooth}
设 $k$ 是域，$X$ 是 $k$ 上维数为 $d$ 的等维光滑概形。映射
$$
A^p(X) \longrightarrow \CH_{d - p}(X),\quad
c \longmapsto c \cap [X]_d
$$
是同构。在这一同构下，双变类的复合变为上面定义的交乘积。
\end{lemma}

\begin{proof}
以 $g : X \to \Spec(k)$ 表示结构态射。该映射是同构的复合
$$
A^p(X) \to A^{p - d}(X \to \Spec(k)) \to \CH_{d - p}(X)
$$
第一支是命题 \ref{chow-proposition-compute-bivariant} 的同构
$c \mapsto c \circ g^*$，第二支是引理 \ref{chow-lemma-chow-cohomology-towards-point} 的同构
$c \mapsto c \cap [\Spec(k)]$。
由引理 \ref{chow-lemma-chow-cohomology-towards-point} 的证明，
第二支箭头的逆把 $\alpha \in \CH_{d - p}(X)$ 映到双变类
$c_\alpha$；对于 $k$ 上局部有限型的 $Y$ 和
$\beta \in \CH_*(Y)$，该双变类把它映到
$\CH_*(X \times_k Y)$ 中的 $\alpha \times \beta$。
再由命题 \ref{chow-proposition-compute-bivariant} 的证明，
第一支箭头的逆把 $c_\alpha$ 映到如下双变类：
对于局部有限型的 $Y \to X$ 和 $\beta \in \CH_*(Y)$，它把后者映到
$\Delta^!(\alpha \times \beta) = \alpha \cdot \beta$。
由此得到引理的最后一项结论。
\end{proof}

\begin{lemma}
\label{chow-lemma-lci-gysin-product}
设 $k$ 是域，$f : X \to Y$ 是 $k$ 上光滑概形之间的态射。
则 $f$ 的 Gysin 映射存在，并且
$f^!(\alpha \cdot \beta) = f^!\alpha \cdot f^!\beta$。
\end{lemma}

\begin{proof}
请注意，$X \to X \times_k Y$ 是 $X$ 到某个 $Y$ 上光滑概形中的浸入。
故 $f$ 的 Gysin 映射存在（定义 \ref{chow-definition-lci-gysin}）。
为证明该公式，可以把 $X$ 和 $Y$ 分解为连通分支；
因此可以假设 $X$ 是 $k$ 上维数为 $d$ 的等维光滑概形，
而 $Y$ 是 $k$ 上维数为 $e$ 的等维光滑概形。
请注意 $f^![Y]_e = [X]_d$（例如见引理 \ref{chow-lemma-lci-gysin-easy}）。
写成 $\alpha = c \cap [Y]_e$ 和 $\beta = c' \cap [Y]_e$，
于是 $\alpha \cdot \beta = c \cap c' \cap [Y]_e$；见
引理 \ref{chow-lemma-identify-chow-for-smooth}。
由引理 \ref{chow-lemma-lci-gysin-commutes}，$f^!$ 与 $c$、$c'$ 均可交换。因此
\begin{align*}
f^!(\alpha \cdot \beta)
& =
f^!(c \cap c' \cap [Y]_e) \\
& =
c \cap c' \cap f^![Y]_e \\
& =
c \cap c' \cap [X]_d \\
& =
(c \cap [X]_d) \cdot (c' \cap [X]_d) \\
& =
(c \cap f^![Y]_e) \cdot (c' \cap f^![Y]_e) \\
& =
f^!(\alpha) \cdot f^!(\beta)
\end{align*}
这正是所需结论；这里也对 $X$ 使用了引理 \ref{chow-lemma-identify-chow-for-smooth}。

\medskip\noindent
另一种证明方法是先证明
$(f \times f)^!(\alpha \times \beta) = f^!\alpha \times f^!\beta$，
再使用引理 \ref{chow-lemma-lci-gysin-composition}。
\end{proof}

\begin{lemma}
\label{chow-lemma-projection-formula}
设 $k$ 是域，$f : X \to Y$ 是 $k$ 上光滑概形之间的固有态射。
则 $f$ 的 Gysin 映射存在，并且
$f_*(\alpha \cdot f^!\beta) = f_*\alpha \cdot \beta$。
\end{lemma}

\begin{proof}
请注意，$X \to X \times_k Y$ 是 $X$ 到某个 $Y$ 上光滑概形中的浸入。
故 $f$ 的 Gysin 映射存在（定义 \ref{chow-definition-lci-gysin}）。
为证明该公式，可以把 $X$ 和 $Y$ 分解为连通分支；
因此可以假设 $X$ 是 $k$ 上维数为 $d$ 的等维光滑概形，
而 $Y$ 是 $k$ 上维数为 $e$ 的等维光滑概形。
请注意 $f^![Y]_e = [X]_d$（例如见引理 \ref{chow-lemma-lci-gysin-easy}）。
写成 $\alpha = c \cap [X]_d$ 和 $\beta = c' \cap [Y]_e$；见
引理 \ref{chow-lemma-identify-chow-for-smooth}。于是
\begin{align*}
f_*(\alpha \cdot f^!\beta)
& =
f_*(c \cap f^!(c' \cap [Y]_e)) \\
& =
f_*(c \cap c' \cap f^![Y]_e) \\
& =
f_*(c \cap c' \cap [X]_d) \\
& =
f_*(c' \cap c \cap [X]_d) \\
& =
c' \cap f_*(c \cap [X]_d) \\
& =
\beta \cdot f_*(\alpha)
\end{align*}
第一项等式来自对 $X$ 应用引理 \ref{chow-lemma-identify-chow-for-smooth} 的结果。
由引理 \ref{chow-lemma-lci-gysin-commutes}，$f^!$ 与 $c'$ 可交换。
交乘积的交换律使得可以交换以 $c$ 和 $c'$ 对 $[X]_d$ 作帽积的次序
（通过该引理）。由于 $c'$ 是双变类，允许将 $c'$ 与 $f_*$ 交换。
最后一项等式仍由该引理给出。
\end{proof}

\begin{lemma}
\label{chow-lemma-intersect-properly}
设 $k$ 是域，$X$ 是 $k$ 上光滑整概形。
设 $Y, Z \subset X$ 是整闭子概形，并置
$d = \dim(Y) + \dim(Z) - \dim(X)$。假设
\begin{enumerate}
\item $\dim(Y \cap Z) \leq d$；以及
\item 对每个满足 $\xi \in Y \cap Z$ 和 $\delta(\xi) = d$ 的点，
$\mathcal{O}_{Y, \xi}$ 与 $\mathcal{O}_{Z, \xi}$ 都是 Cohen--Macaulay 的。
\end{enumerate}
则在 $\CH_d(X)$ 中有 $[Y] \cdot [Z] = [Y \cap Z]_d$。
\end{lemma}

\begin{proof}
回忆 $[Y] \cdot [Z] = \Delta^!([Y \times Z])$，其中
$\Delta^! = c(\Delta : X \to X \times X, \mathcal{T}_{X/k})$
是高余维 Gysin 映射（第 \ref{chow-section-gysin-higher-codimension} 节），且
$\mathcal{T}_{X/k} = \SheafHom(\Omega_{X/k}, \mathcal{O}_X)$
是秩为 $\dim(X)$ 的局部自由层。概形之间有等式
$$
Y \cap Z = X \times_{\Delta, (X \times X)} (Y \times Z)
$$
并且 $\dim(Y \times Z) = \dim(Y) + \dim(Z)$，
故引理 \ref{chow-lemma-gysin-easy} 的条件 (1)、(2)、(3) 成立。
最后，若 $\xi \in Y \cap Z$，则有平坦局部同态
$$
\mathcal{O}_{Y, \xi} \longrightarrow
\mathcal{O}_{Y \times Z, \xi}
$$
其“纤维”为 $\mathcal{O}_{Z, \xi}$。因此，若
$\mathcal{O}_{Y, \xi}$ 和 $\mathcal{O}_{Z, \xi}$ 都是 Cohen--Macaulay 的，
则 $\mathcal{O}_{Y \times Z, \xi}$ 也是；见
《交换代数》引理 \ref{algebra-lemma-CM-goes-up}。
由此可见，引理 \ref{chow-lemma-gysin-easy} 的所有假设都满足，结论随即得出。
\end{proof}

\begin{lemma}
\label{chow-lemma-intersect-regularly-embedded}
设 $k$ 是域，$X$ 是 $k$ 上光滑概形，$i : Y \to X$ 是正则闭浸入，
且 $\alpha \in \CH_*(X)$。若 $Y$ 是维数为 $e$ 的等维概形，则
$\alpha \cdot [Y]_e = i_*(i^!(\alpha))$ 在 $\CH_*(X)$ 中成立。
\end{lemma}

\begin{proof}
将 $X$ 分解为连通分支后，可以并且确实假设 $X$ 是维数为 $d$ 的等维概形。
写成 $\alpha = c \cap [X]_n$，其中 $x \in A^*(X)$；见
引理 \ref{chow-lemma-identify-chow-for-smooth}。于是
$$
i_*(i^!(\alpha)) = i_*(i^!(c \cap [X]_n)) =
i_*(c \cap i^![X]_n) = i_*(c \cap [Y]_e) =
c \cap i_*[Y]_e = \alpha \cdot [Y]_e
$$
第一项等式来自 $c$ 的选择；第二项等式由
引理 \ref{chow-lemma-lci-gysin-commutes} 给出；第三项等式是因为
$\CH_*(Y)$ 中 $i^![X]_d = [Y]_e$
（引理 \ref{chow-lemma-lci-gysin-easy}）；第四项等式是因为双变类与固有前推可交换；
最后一项等式由引理 \ref{chow-lemma-identify-chow-for-smooth} 给出。
\end{proof}

\begin{lemma}
\label{chow-lemma-intersection-regular-smooth}
设 $k$ 是域，$X$ 是 $k$ 上拟紧且具有仿射对角态射的光滑概形。
则本节在 $\CH^*(X)$ 上构造的交乘积，在与 $\mathbf{Q}$ 张量后，
与第 \ref{chow-section-intersection-regular} 节中构造的交乘积一致。
\end{lemma}

\begin{proof}
设 $\alpha \in \CH^i(X)$ 且 $\beta \in \CH^j(X)$。
如第 \ref{chow-section-intersection-regular} 节，写成
$\alpha = ch(\alpha') \cap [X]$ 和 $\beta = ch(\beta') \cap [X]$，
其中 $\alpha', \beta' \in K_0(\textit{Vect}(X)) \otimes \mathbf{Q}$。
令 $c = ch(\alpha')$、$c' = ch(\beta')$。
于是第 \ref{chow-section-intersection-regular} 节中的交乘积给出
$c \cap c' \cap [X]$。由引理 \ref{chow-lemma-identify-chow-for-smooth}，
这与 $\alpha \cdot \beta$ 相同（更确切地说，是使用其推广：
对任意域 $k$ 上光滑概形 $X$，映射
$A^i(X) \to \CH^i(X)$、$c \mapsto c \cap [X]$ 是同构）。
\end{proof}








\section{Dedekind 整环上的外积}
\label{chow-section-exterior-product-dim-1}

\noindent
设 $S$ 是局部 Noether 概形，并且存在由 Dedekind 整环的谱组成的开覆盖。
对闭点 $s \in S$ 置 $\delta(s) = 0$，否则置 $\delta(s) = 1$。
则 $(S, \delta)$ 是一般情形 \ref{chow-situation-setup} 的特殊情形；
见例 \ref{chow-example-domain-dimension-1}。
请注意，$S$ 正规（《交换代数》引理 \ref{algebra-lemma-characterize-Dedekind}），
因而是正规整概形的不交并
（《概形的性质》引理 \ref{properties-lemma-normal-locally-Noetherian}）。
所以以下所有论证都归结为 $S$ 不可约的情形。
另一方面，我们允许 $S$ 非分离（例如，$S$ 可以是原点 $0$ 加倍的仿射直线）。

\medskip\noindent
考虑 $S$ 上局部有限型概形的笛卡尔方块
$$
\xymatrix{
X \times_S Y \ar[r] \ar[d] & Y \ar[d] \\
X \ar[r] & S
}
$$
我们断言存在规范映射
$$
\times :
\CH_n(X) \otimes_{\mathbf{Z}} \CH_m(Y)
\longrightarrow
\CH_{n + m - 1}(X \times_S Y)
$$
它由如下规则唯一确定：给定 $\delta$-维数分别为 $n$ 和 $m$ 的整闭子概形
$X' \subset X$ 与 $Y' \subset Y$，置
\begin{enumerate}
\item 若 $X'$ 或 $Y'$ 支配 $S$ 的一个不可约分支，则
$[X'] \times [Y'] = [X' \times_S Y']_{n + m - 1}$；
\item 若 $X'$ 和 $Y'$ 都不支配 $S$ 的任何不可约分支，则
$[X'] \times [Y'] = 0$。
\end{enumerate}

\begin{lemma}
\label{chow-lemma-exterior-product-well-defined-dim-1}
映射 $\times : \CH_n(X) \otimes_{\mathbf{Z}} \CH_m(Y) \to
\CH_{n + m - 1}(X \times_S Y)$ 是良定义的。
\end{lemma}

\begin{proof}
考虑 $n$-周和 $m$-周 $\alpha = \sum_{i \in I} n_i[X_i]$ 与
$\beta = \sum_{j \in J} m_j[Y_j]$，其中 $X_i \subset X$ 和 $Y_j \subset Y$
分别是 $\delta$-维数为 $n$ 和 $m$ 的整闭子概形的局部有限族。
令 $K \subset I \times J$ 为所有满足如下条件的序对
$(i, j) \in I \times J$ 的集合：$X_i$ 或 $Y_j$ 支配 $S$ 的一个不可约分支。
则 $\{X_i \times_S Y_j\}_{(i, j) \in K}$ 是 $X \times_S Y$ 中
$\delta$-维数为 $n + m - 1$ 的闭子概形的局部有限族。
这意味着确实可以将
$$
\alpha \times \beta =
\sum\nolimits_{(i, j) \in K} n_i m_j [X_i \times_S Y_j]_{n + m - 1}
$$
视为 $X \times_S Y$ 上的一个 $(n + m - 1)$-周。
由此得到加性映射
$\times : Z_n(X) \otimes_{\mathbf{Z}} Z_m(Y) \to Z_{n + m}(X \times_S Y)$。
问题在于证明这一过程与有理等价相容。

\medskip\noindent
设 $i : X' \to X$ 是一个 $\delta$-维数为 $n$、
并支配 $S$ 某个不可约分支的整闭子概形的包含态射。
则 $p' : X' \to S$ 是相对维数为 $n - 1$ 的平坦态射；见
《代数进阶》引理 \ref{more-algebra-lemma-dedekind-torsion-free-flat}。
因此，由引理 \ref{chow-lemma-flat-pullback-bivariant}，沿 $p'$ 的平坦拉回是元素
$(p')^* \in A^{-n + 1}(X' \to S)$；进而由
引理 \ref{chow-lemma-push-proper-bivariant}，有
$c' = i_* \circ (p')^* \in A^{-n + 1}(X \to S)$。
这给出映射
$$
c' \cap - : \CH_m(Y) \longrightarrow \CH_{m + n - 1}(X \times_S Y)
$$
对于任意 $\delta$-维数为 $m$ 的整闭子概形 $Y' \subset Y$，
该映射把 $[Y']$ 映到 $[X' \times_S Y']_{n + m - 1}$。

\medskip\noindent
设 $i : X' \to X$ 是一个 $\delta$-维数为 $n$ 的整闭子概形的包含态射，
并且复合 $X' \to X \to S$ 通过某个闭点 $s \in S$ 分解。
由于 $s$ 是 Dedekind 整环之谱的闭点，$s$ 是 $S$ 上法丛平凡的有效
Cartier 除子。以 $c \in A^1(s \to S)$ 表示 Gysin 同态；见
引理 \ref{chow-lemma-gysin-bivariant}。态射 $p' : X' \to s$
是相对维数为 $n$ 的平坦态射。因此，由
引理 \ref{chow-lemma-flat-pullback-bivariant}，沿 $p'$ 的平坦拉回是元素
$(p')^* \in A^{-n}(X' \to S)$。于是
$$
c' = i_* \circ (p')^* \circ c \in A^{-n}(X \to S)
$$
由引理 \ref{chow-lemma-push-proper-bivariant}，这给出映射
$$
c' \cap - : \CH_m(Y) \longrightarrow \CH_{m + n - 1}(X \times_S Y)
$$
对于任意 $\delta$-维数为 $m$ 的整闭子概形 $Y' \subset Y$，
若 $Y'$ 支配 $S$ 的一个不可约分支，该映射把 $[Y']$ 映到
$[X' \times_S Y']_{n + m - 1}$；否则映到 $0$。

\medskip\noindent
由前两段可知，构造
$([X'], [Y']) \mapsto [X' \times_S Y']_{n + m - 1}$
在第二个变量中可通过有理等价下降，也就是说，它给出良定义映射
$Z_n(X) \otimes_{\mathbf{Z}} \CH_m(Y) \to \CH_{n + m - 1}(X \times_S Y)$。
由对称性，另一个变量也同样如此，结论得证。
\end{proof}

\begin{lemma}
\label{chow-lemma-chow-cohomology-towards-base-dim-1}
设 $(S, \delta)$ 如上，$X$ 是 $S$ 上局部有限型概形。
则对所有 $p \in \mathbf{Z}$，都有规范等同
$$
A^p(X \to S) = \CH_{1 - p}(X)
$$
\end{lemma}

\begin{proof}
考虑元素 $[S]_1 \in \CH_1(S)$。
把 $c$ 映到 $c \cap [S]_1$，即得到映射
$A^p(X \to S) \to \CH_{1 - p}(X)$。

\medskip\noindent
反之，设有 $\alpha \in \CH_{1 - p}(X)$。
可按如下方式定义 $c_\alpha \in A^p(X \to S)$：
给定 $X' \to S$ 和 $\alpha' \in \CH_n(X')$，令
$$
c_\alpha \cap \alpha' = \alpha \times \alpha'
$$
作为 $\CH_{n - p}(X \times_S X')$ 中的元素。
为证明这是双变类，如定义 \ref{chow-definition-cycles} 写成
$\alpha = \sum_{i \in I} n_i[X_i]$。特别地，态射
$$
g : \coprod\nolimits_{i \in I} X_i \longrightarrow X
$$
是固有的。选取 $i \in I$。若 $X_i$ 支配 $S$ 的一个不可约分支，
则结构态射 $p_i : X_i \to S$ 平坦，并且有
$\xi_i = p_i^* \in A^p(X_i \to S)$。
另一方面，若 $p_i$ 分解为 $p'_i : X_i \to s_i$，随后是闭点的包含
$s_i \to S$，则有
$\xi_i = (p'_i)^* \circ c_i \in A^p(X_i \to S)$，
其中 $c_i \in A^1(s_i \to S)$ 是 Gysin 同态，而
$(p'_i)^*$ 是平坦拉回。注意
$$
A^p(\coprod\nolimits_{i \in I} X_i \to S) =
\prod\nolimits_{i \in I} A^p(X_i \to S)
$$
因此有
$$
\xi = \sum n_i \xi_i \in A^p(\coprod\nolimits_{i \in I} X_i \to S)
$$
最后，由于 $g$ 固有，引理 \ref{chow-lemma-push-proper-bivariant} 给出双变类
$$
g_* \circ \xi \in A^p(X \to S)
$$
读者容易验证 $c_\alpha$ 等于该双变类
（请与引理 \ref{chow-lemma-exterior-product-well-defined-dim-1} 的证明比较），
因而它本身也是双变类。

\medskip\noindent
为完成证明，还须说明这两个构造互为逆映射。
由于 $c_\alpha \cap [S]_1 = \alpha$，其中一个方向是显然的。
对另一个方向，令 $c \in A^p(X \to S)$，并置 $\alpha = c \cap [S]_1$。
由引理 \ref{chow-lemma-bivariant-zero}，只需在 $X'$ 是 $S$ 上局部有限型整概形时证明
$$
c \cap [X'] = c_\alpha \cap [X']
$$
但或者 $p' : X' \to S$ 是相对维数为 $\dim_\delta(X') - 1$ 的平坦态射，
因而 $[X'] = (p')^*[S]_1$；或者 $X' \to S$ 分解为
$X' \to s \to S$，因而 $[X'] = (p')^*(s \to S)^*[S]_1$。
所以，双变类 $c$ 和 $c_\alpha$ 在 $[S]_1$ 上相同，
就蕴含它们与 $[X']$ 作帽积时也相同
（因为双变类与平坦拉回和 Gysin 映射可交换），证明完成。
\end{proof}

\begin{lemma}
\label{chow-lemma-chow-cohomology-towards-base-dim-1-commutes}
设 $(S, \delta)$ 如上，$X$ 是 $S$ 上局部有限型概形。
设 $c \in A^p(X \to S)$，$Y \to Z$ 是 $S$ 上局部有限型概形之间的态射，
且 $c' \in A^q(Y \to Z)$。则在
$A^{p + q}(X \times_S Y \to X \times_S Z)$ 中有
$c \circ c' = c' \circ c$。
\end{lemma}

\begin{proof}
在引理 \ref{chow-lemma-chow-cohomology-towards-base-dim-1} 的证明中已经看到，
$c$ 由固有前推、在各连通分支上乘以整数、平坦拉回以及 Gysin 映射组合而成。
依双变类的定义，$c'$ 与这些运算中的每一个都可交换，故结论成立。
这里省略若干细节。
\end{proof}

\begin{remark}
\label{chow-remark-commuting-exterior-dim-1}
引理 \ref{chow-lemma-chow-cohomology-towards-base-dim-1} 和
\ref{chow-lemma-chow-cohomology-towards-base-dim-1-commutes} 的要点如下。
设 $(S, \delta)$ 如上，$X$ 是 $S$ 上局部有限型概形。
设 $\alpha \in \CH_*(X)$，$Y \to Z$ 是 $S$ 上局部有限型概形之间的态射，
且 $c' \in A^q(Y \to Z)$。则对任意 $\beta \in \CH_*(Z)$，
$$
\alpha \times (c' \cap \beta) = c' \cap (\alpha \times \beta)
$$
在 $\CH_*(X \times_S Y)$ 中成立。事实上，只需取与 $\alpha$ 对应的双变类
$c = c_\alpha \in A^*(X \to S)$；见引理 \ref{chow-lemma-chow-cohomology-towards-base-dim-1} 的证明。
\end{remark}

\begin{lemma}
\label{chow-lemma-exterior-product-associative-dim-1}
外积满足结合律。更确切地，设 $(S, \delta)$ 如上，
$X, Y, Z$ 是 $S$ 上局部有限型概形，并设
$\alpha \in \CH_*(X)$、$\beta \in \CH_*(Y)$、$\gamma \in \CH_*(Z)$。
则在 $\CH_*(X \times_S Y \times_S Z)$ 中有
$(\alpha \times \beta) \times \gamma =
\alpha \times (\beta \times \gamma)$。
\end{lemma}

\begin{proof}
从略。提示：概形纤维积的结合律。
\end{proof}















\section{Dedekind 整环上的交乘积}
\label{chow-section-intersection-product-dim-1}

\noindent
设 $S$ 是局部 Noether 概形，并且存在由 Dedekind 整环的谱组成的开覆盖。
对闭点 $s \in S$ 置 $\delta(s) = 0$，否则置 $\delta(s) = 1$。
则 $(S, \delta)$ 是一般情形 \ref{chow-situation-setup} 的特殊情形；
见例 \ref{chow-example-domain-dimension-1} 以及
第 \ref{chow-section-exterior-product-dim-1} 节中的讨论。

\medskip\noindent
设 $X$ 是 $S$ 上光滑概形。第 \ref{chow-section-gysin-for-diagonal} 节的双变类
$\Delta^!$ 使我们能够在 $X$ 上局部有限型概形的 Chow 群上定义一种交乘积。
具体地，设 $Y \to X$ 和 $Z \to X$ 是局部有限型概形之间的态射。注意到
$$
Y \times_X Z =  (Y \times_S Z) \times_{X \times_S X, \Delta} X
$$
因此可以考虑映射序列
$$
\CH_n(Y) \otimes_\mathbf{Z} \CH_m(Y)
\xrightarrow{\times}
\CH_{n + m - 1}(Y \times_S Z)
\xrightarrow{\Delta^!}
\CH_{n + m - *}(Y \times_X Z)
$$
其中第一支箭头是第 \ref{chow-section-exterior-product-dim-1} 节中构造的外积，
第二支箭头是第 \ref{chow-section-gysin-for-diagonal} 节中研究的对角态射 Gysin 映射。
若 $X$ 是维数为 $d$ 的等维概形，则 $X \to S$ 是相对维数为 $d - 1$ 的光滑态射，
故所得元素落在 $\CH_{n + m - d}(Y \times_X Z)$ 中。
一般而言，可依照 $X$ 的开闭子概形（其上 $X$ 具有给定维数），
将其分解为位于这些部分之上的分量。
给定 $\alpha \in \CH_*(Y)$ 和 $\beta \in \CH_*(Z)$，记
$$
\alpha \cdot \beta = \Delta^!(\alpha \times \beta)
\in \CH_*(Y \times_X Z)
$$
在特殊情形 $X = Y = Z$ 下，得到乘法
$$
\CH_*(X) \times \CH_*(X) \to \CH_*(X),\quad
(\alpha, \beta) \mapsto \alpha \cdot \beta
$$
称为{\it 交乘积}。显然，这一乘积是对称的。
结合律由下一个引理给出。

\begin{lemma}
\label{chow-lemma-associative-dim-1}
上面定义的乘积满足结合律。更确切地，设 $(S, \delta)$ 如上，
$X$ 在 $S$ 上光滑，$Y, Z, W$ 是 $X$ 上局部有限型概形，并设
$\alpha \in \CH_*(Y)$、$\beta \in \CH_*(Z)$、$\gamma \in \CH_*(W)$。
则在 $\CH_*(Y \times_X Z \times_X W)$ 中有
$(\alpha \cdot \beta) \cdot \gamma =
\alpha \cdot (\beta \cdot \gamma)$。
\end{lemma}

\begin{proof}
由引理 \ref{chow-lemma-exterior-product-associative-dim-1}，在
$\CH_*(Y \times_S Z \times_S W)$ 中有
$(\alpha \times \beta) \times \gamma =
\alpha \times (\beta \times \gamma)$。考虑闭浸入
$$
\Delta_{12} : X \times_S X \longrightarrow X \times_S X \times_S X,
\quad (x, x') \mapsto (x, x, x')
$$
以及
$$
\Delta_{23} : X \times_S X \longrightarrow X \times_S X \times_S X,
\quad (x, x') \mapsto (x, x', x')
$$
以 $\Delta_{12}^!$ 和 $\Delta_{23}^!$ 表示相应的双变类。
注意，$\Delta_{12}^!$ 是 $\Delta^!$ 沿映射 $\text{pr}_{12}$ 在
$X \times_S X \times_S X$ 上的限制
（注 \ref{chow-remark-restriction-bivariant}），而
$\Delta_{23}^!$ 是 $\Delta^!$ 沿映射 $\text{pr}_{23}$ 在
$X \times_S X \times_S X$ 上的限制。
因此，$\Delta_{12}^!$ 沿 $\Delta_{23}$ 的限制显然是 $\Delta^!$，
而 $\Delta_{23}^!$ 沿 $\Delta_{12}$ 的限制同样是 $\Delta^!$。
于是由引理 \ref{chow-lemma-gysin-commutes}，有
$$
\Delta^! \circ \Delta_{12}^! =
\Delta^! \circ \Delta_{23}^! 
$$
现在用如下等式链证明该引理：
\begin{align*}
(\alpha \cdot \beta) \cdot \gamma
& =
\Delta^!(\Delta^!(\alpha \times \beta) \times \gamma) \\
& =
\Delta^!(\Delta_{12}^!((\alpha \times \beta) \times \gamma)) \\
& =
\Delta^!(\Delta_{23}^!((\alpha \times \beta) \times \gamma)) \\
& =
\Delta^!(\Delta_{23}^!(\alpha \times (\beta \times \gamma)) \\
& =
\Delta^!(\alpha \times \Delta^!(\beta \times \gamma)) \\
& =
\alpha \cdot (\beta \cdot \gamma)
\end{align*}
除第二项和倒数第二项等式或许需要说明外，其余等式均由上文显然成立。
等式
$\Delta_{23}^!(\alpha \times (\beta \times \gamma)) =
\alpha \times \Delta^!(\beta \times \gamma)$
由注 \ref{chow-remark-commuting-exterior} 成立；第二项等式同理。
\end{proof}

\begin{lemma}
\label{chow-lemma-identify-chow-for-smooth-dim-1}
设 $(S, \delta)$ 如上，$X$ 是 $S$ 上维数为 $d$ 的等维光滑概形。映射
$$
A^p(X) \longrightarrow \CH_{d - p}(X),\quad
c \longmapsto c \cap [X]_d
$$
是同构。在这一同构下，双变类的复合变为上面定义的交乘积。
\end{lemma}

\begin{proof}
以 $g : X \to S$ 表示结构态射。该映射是同构的复合
$$
A^p(X) \to A^{p - d + 1}(X \to S) \to \CH_{d - p}(X)
$$
第一支是命题 \ref{chow-proposition-compute-bivariant} 的同构
$c \mapsto c \circ g^*$，第二支是引理 \ref{chow-lemma-chow-cohomology-towards-base-dim-1} 的同构
$c \mapsto c \cap [S]_1$。
由引理 \ref{chow-lemma-chow-cohomology-towards-base-dim-1} 的证明，
第二支箭头的逆把 $\alpha \in \CH_{d - p}(X)$ 映到双变类
$c_\alpha$；对于 $k$ 上局部有限型的 $Y$ 和
$\beta \in \CH_*(Y)$，该双变类把它映到
$\CH_*(X \times_k Y)$ 中的 $\alpha \times \beta$。
再由命题 \ref{chow-proposition-compute-bivariant} 的证明，
第一支箭头的逆把 $c_\alpha$ 映到如下双变类：
对于局部有限型的 $Y \to X$ 和 $\beta \in \CH_*(Y)$，它把后者映到
$\Delta^!(\alpha \times \beta) = \alpha \cdot \beta$。
由此得到引理的最后一项结论。
\end{proof}










\section{Todd 类}
\label{chow-section-todd-classes}

\noindent
与秩为 $r$ 的向量丛 $\mathcal{E}$ 相关联的最后一种类是其
{\it Todd 类} $Todd(\mathcal{E})$。用 Chern 根
$x_1, \ldots, x_r$ 表示时，它定义为
$$
Todd(\mathcal{E})
=
\prod\nolimits_{i = 1}^r
\frac{x_i}{1 - e^{-x_i}}
$$
用 Chern 类 $c_i = c_i(\mathcal{E})$ 表示时，有
$$
Todd(\mathcal{E})
=
1
+
\frac{1}{2}c_1
+
\frac{1}{12}(c_1^2 + c_2)
+
\frac{1}{24}c_1c_2
+
\frac{1}{720}(-c_1^4 + 4c_1^2c_2 + 3c_2^2 + c_1c_3 - c_4)
+
\ldots
$$
上一节已经对分母作了适当说明。对于任意正合列
$$
0
\to
{\mathcal E}_1
\to
{\mathcal E}
\to
{\mathcal E}_2
\to
0
$$
都有
$$
Todd({\mathcal E}) = Todd({\mathcal E}_1) Todd({\mathcal E}_2).
$$







\section{Grothendieck--Riemann--Roch}
\label{chow-section-grr}

\noindent
设 $(S, \delta)$ 如情形 \ref{chow-situation-setup}，
$X, Y$ 在 $S$ 上局部有限型。设 $\mathcal{E}$ 是有限局部自由层；
具体地，${\mathcal E}$ 是 $X$ 上秩为 $r$ 的层。
设 $f : X \to Y$ 是固有光滑态射，并假设
$R^if_*\mathcal{E}$ 都是 $Y$ 上有限秩局部自由层。
在这种情形下，Grothendieck--Riemann--Roch 定理断言
$$
f_*(Todd(T_{X/Y}) ch(\mathcal{E}))
=
\sum (-1)^i ch(R^if_*\mathcal{E})
$$
其中
$$
T_{X/Y} = \SheafHom_{\mathcal{O}_X}(\Omega_{X/Y}, \mathcal{O}_X)
$$
是 $X$ 在 $Y$ 上的相对切丛。若 $Y = \Spec(k)$，其中 $k$ 是域，
则可将其改述为
$$
\chi(X, \mathcal{E}) = \deg(Todd(T_{X/k}) ch(\mathcal{E}))
$$
该定理还有更一般的形式，并且在正确的一般性下表述时更容易证明。
我们将在别处回到这个问题（在此插入未来引用）。







\section{基概形的变换}
\label{chow-section-change-base}

\noindent
本节说明如何比较不同基概形上的理论。

\begin{situation}
\label{chow-situation-setup-base-change}
这里 $(S, \delta)$ 和 $(S', \delta')$ 如情形 \ref{chow-situation-setup}。
此外，$g : S' \to S$ 是概形的平坦态射，$c \in \mathbf{Z}$ 是满足如下条件的整数：
对每个 $s \in S$，以及作为 $g^{-1}(\{s\})$ 某个不可约分支一般点的
$s' \in S'$，都有 $\delta'(s') = \delta(s) + c$。
\end{situation}

\noindent
我们将看到，对于 $S$ 上局部有限型概形 $X$，存在 Chow 群之间的良定义映射
$\CH_k(X) \to \CH_{k + c}(X \times_S S')$；它大体上与本章定义的运算可交换。

\begin{lemma}
\label{chow-lemma-dimension-base-change}
在情形 \ref{chow-situation-setup-base-change} 中，设 $X \to S$ 局部有限型，
并以 $X' \to S'$ 表示沿 $S' \to S$ 的基变换。
若 $X$ 是满足 $\dim_\delta(X) = k$ 的整概形，则
$X'$ 的每个不可约分支 $Z'$ 都满足 $\dim_{\delta'}(Z') = k + c$。
\end{lemma}

\begin{proof}
投影 $X' \to X$ 是平坦态射 $S' \to S$ 的基变换，因而平坦
（《概形态射》引理 \ref{morphisms-lemma-base-change-flat}）。
所以 $X'$ 任一不可约分支的一般点 $x'$ 都映到一般点 $x \in X$
（因为由《概形态射》引理 \ref{morphisms-lemma-generalizations-lift-flat}，
一般化可沿 $X' \to X$ 提升）。令 $s \in S$ 为 $x$ 的像。
回忆概形 $S'_s = S' \times_S s$ 与 $g^{-1}(\{s\})$ 具有相同的底拓扑空间
（《概形》引理 \ref{schemes-lemma-fibre-topological}）。
可将 $x'$ 视为概形 $S'_s \times_s x$ 的一个点；该概形带有单态射
$S'_s \times_s x \to S' \times_S X$。当然，$x'$ 也是
$S'_s \times_s x$ 某个不可约分支的一般点。
利用 $\Spec(\kappa(x)) \to \Spec(\kappa(s)) = s$ 的平坦性并重复上述论证，
可见 $x'$ 映到 $g^{-1}(\{s\})$ 某个不可约分支的一般点 $s'$。
因此由假设，$\delta'(s') = \delta(s) + c$。
由《概形态射》引理 \ref{morphisms-lemma-dimension-fibre-after-base-change}，有
$\dim_x(X_s) = \dim_{x'}(X_{s'})$。
由于 $x$ 是不可约分支 $X_s$ 的一般点
（这是不可约概形，不过我们不需要这一点），而 $x'$ 是 $X'_{s'}$ 某个不可约分支的一般点，
故由《概形态射》引理 \ref{morphisms-lemma-dimension-fibre-at-a-point}，得到
$\text{trdeg}_{\kappa(s)}(\kappa(x)) = 
\text{trdeg}_{\kappa(s')}(\kappa(x'))$。于是
$$
\delta_{X'/S'}(x') = \delta(s') + \text{trdeg}_{\kappa(s')}(\kappa(x')) =
\delta(s) + c + \text{trdeg}_{\kappa(s)}(\kappa(x)) = \delta_{X/S}(x) + c
$$
由定义 \ref{chow-definition-delta-dimension}，这就证明了所需结论。
\end{proof}

\noindent
在情形 \ref{chow-situation-setup-base-change} 中，设 $X \to S$ 局部有限型，
并以 $X' \to S'$ 表示沿 $g : S' \to S$ 的基变换。
存在唯一同态
$$
g^* : Z_k(X) \longrightarrow Z_{k + c}(X')
$$
它把任意 $\delta$-维数为 $k$ 的整闭子概形 $Z \subset X$ 的 $[Z]$
映到 $[Z \times_S S']_{k + c}$。由引理 \ref{chow-lemma-dimension-base-change}，这是有意义的。

\begin{lemma}
\label{chow-lemma-pullback-coherent-base-change}
在情形 \ref{chow-situation-setup-base-change} 中，设 $X \to S$ 局部有限型，
且 $X' \to S$ 是沿 $S' \to S$ 的基变换。
\begin{enumerate}
\item 设 $Z \subset X$ 是满足 $\dim_\delta(Z) \leq k$ 的闭子概形，
其基变换为 $Z' \subset X'$。则 $\dim_{\delta'}(Z')) \leq k + c$，
且在 $Z_{k + c}(X')$ 中有 $[Z']_{k + c} = g^*[Z]_k$。
\item 设 $\mathcal{F}$ 是 $X$ 上满足
$\dim_\delta(\text{Supp}(\mathcal{F})) \leq k$ 的凝聚层，
其在 $X'$ 上的基变换为 $\mathcal{F}'$。
则 $\dim_\delta(\text{Supp}(\mathcal{F}')) \leq k + c$，
且在 $Z_{k + c}(X')$ 中有
$g^*[\mathcal{F}]_k = [\mathcal{F}']_{k + c}$。
\end{enumerate}
\end{lemma}

\begin{proof}
该证明与引理 \ref{chow-lemma-pullback-coherent} 的证明完全相同，
建议读者略过。

\medskip\noindent
关于维数的断言由引理 \ref{chow-lemma-dimension-base-change} 得出。
由引理 \ref{chow-lemma-cycle-closed-coherent}，以及凝聚模
$\mathcal{O}_Z$ 的基变换是 $\mathcal{O}_{Z'}$ 这一事实，
(1) 由 (2) 得出。

\medskip\noindent
证明 (2)。由于 $X$、$X'$ 局部 Noether，可应用
《概形上同调》引理 \ref{coherent-lemma-coherent-Noetherian} 得知
$\mathcal{F}$ 有限型；所以 $\mathcal{F}'$ 有限型
（《模层》引理 \ref{modules-lemma-pullback-finite-type}），
进而 $\mathcal{F}'$ 凝聚
（再次使用《概形上同调》引理 \ref{coherent-lemma-coherent-Noetherian}）。故引理的陈述有意义。
设 $W \subset X$ 是 $\delta$-维数为 $k$ 的整闭子概形，
$W' \subset X'$ 是 $\delta'$-维数为 $k + c$ 的整闭子概形，
并且在 $X' \to X$ 下映入 $W$。须证明 $g^*[\mathcal{F}]_k$ 中
$[W']$ 的系数 $n$ 等于 $[\mathcal{F}']_{k + c}$ 中 $[W']$ 的系数 $m$。
设 $\xi \in W$ 和 $\xi' \in W'$ 为一般点。令
$A = \mathcal{O}_{X, \xi}$、$B = \mathcal{O}_{X', \xi'}$，
并将 $M = \mathcal{F}_\xi$ 视为 $A$-模。
（由维数假设，$M$ 的长度有限，但实际上不必验证这一点；
见引理 \ref{chow-lemma-length-finite}。）有
$\mathcal{F}'_{\xi'} = B \otimes_A M$。因此
$$
n = \text{length}_B(B \otimes_A M)
\quad
\text{且}
\quad
m = \text{length}_A(M) \text{length}_B(B/\mathfrak m_AB)
$$
所需等式遂由《交换代数》引理 \ref{algebra-lemma-pullback-module} 得出。
\end{proof}

\begin{lemma}
\label{chow-lemma-pullback-base-change}
在情形 \ref{chow-situation-setup-base-change} 中，设 $X \to S$ 局部有限型，
且 $X' \to S'$ 是沿 $S' \to S$ 的基变换。
上述映射 $g^* : Z_k(X) \to Z_{k + c}(X')$ 可通过有理等价下降，
从而给出 Chow 群之间的映射
$$
g^* : \CH_k(X) \longrightarrow \CH_{k + c}(X')
$$
\end{lemma}

\begin{proof}
设 $\alpha \in Z_k(X)$ 是有理等价于零的 $k$-周。
由引理 \ref{chow-lemma-rational-equivalence-family}，存在一个整闭子概形的局部有限族
$W_i \subset X \times \mathbf{P}^1$，其 $\delta$-维数为 $k$，
且不包含在 $X \times \mathbf{P}^1$ 的除子
$(X \times \mathbf{P}^1)_0$ 或 $(X \times \mathbf{P}^1)_\infty$ 中，并满足
$\alpha = \sum ([(W_i)_0]_k - [(W_i)_\infty]_k)$。
因此，只需对 $\delta$-维数为 $k$ 的整闭子概形
$W \subset X \times \mathbf{P}^1$ 证明如下两点，
其中该子概形不包含在 $X \times \mathbf{P}^1$ 的除子
$(X \times \mathbf{P}^1)_0$ 或 $(X \times \mathbf{P}^1)_\infty$ 中：
\begin{enumerate}
\item 基变换 $W' \subset X' \times \mathbf{P}^1$ 满足
引理 \ref{chow-lemma-closed-subscheme-cross-p1} 的假设，其中以 $k + c$ 代替 $k$；以及
\item $g^*[W_0]_k = [(W')_0]_{k + c}$ 且
$g^*[W_\infty]_k = [(W')_\infty]_{k + c}$。
\end{enumerate}
(2) 立即由引理 \ref{chow-lemma-pullback-coherent-base-change} 以及
$(W')_0$ 是 $W_0$ 的基变换这一事实得出（利用纤维积的结合律）。
对于 (1)，维数断言先由引理 \ref{chow-lemma-dimension-base-change} 得出。
再令 $w' \in (W')_0$，其像为 $w \in W_0$，并令
$z \in \mathbf{P}^1_S$。以 $t \in \mathcal{O}_{\mathbf{P}^1_S, z}$
表示 $0 : S \to \mathbf{P}^1_S$ 的通常方程。
由于 $\mathcal{O}_{W, w} \to \mathcal{O}_{W', w'}$ 平坦，
且 $t$ 是 $\mathcal{O}_{W, w}$ 上的非零因子
（因为 $W$ 整且 $W \not = W_0$），故 $t$ 也是
$\mathcal{O}_{W', w'}$ 中的非零因子。因此 $W'$ 没有落在 $W'_0$ 上的伴随点。
\end{proof}

\begin{lemma}
\label{chow-lemma-pullback-base-change-pullback}
在情形 \ref{chow-situation-setup-base-change} 中，设
$Y \to X \to S$ 局部有限型，并以 $Y' \to X' \to S'$ 表示沿
$S' \to S$ 的基变换。假设 $f : Y \to X$ 是相对维数为 $r$ 的平坦态射。
则 $f' : Y' \to X'$ 是相对维数为 $r$ 的平坦态射，并且周群与 Chow 群的图
$$
\vcenter{
\xymatrix{
Z_{k + r}(Y) \ar[r]_{g^*} & Z_{k + c + r}(Y') \\
Z_k(X) \ar[r]^{g^*} \ar[u]^{f^*} & Z_{k + c}(X') \ar[u]_{(f')^*}
}
}
\quad\text{且}\quad
\vcenter{
\xymatrix{
\CH_{k + r}(Y) \ar[r]_{g^*} & \CH_{k + c + r}(Y') \\
\CH_k(X) \ar[r]^{g^*} \ar[u]^{f^*} & \CH_{k + c}(X') \ar[u]_{(f')^*}
}
}
$$
均交换。
\end{lemma}

\begin{proof}
只需证明第一个图交换。为此，设 $Z \subset X$ 是
$\delta$-维数为 $k$ 的整闭子概形，并以 $Z' \subset X'$ 表示其基变换。
由构造，有 $g^*[Z] = [Z']_{k + c}$。由引理 \ref{chow-lemma-pullback-coherent}，有
$(f')^*g^*[Z] = [Z' \times_{X'} Y']_{k + c + r}$。
反过来，由定义 \ref{chow-definition-flat-pullback}，有
$f^*[Z] = [Z \times_X Y]_{k + r}$；由
引理 \ref{chow-lemma-pullback-coherent-base-change}，有
$g^*f^*[Z] = [(Z \times_X Y)']_{k + r + c}$。
由于纤维积的结合律给出
$(Z \times_X Y)' = Z' \times_{X'} Y'$，结论成立。
\end{proof}

\begin{lemma}
\label{chow-lemma-pullback-base-change-pushforward}
在情形 \ref{chow-situation-setup-base-change} 中，设
$Y \to X \to S$ 局部有限型，并以 $Y' \to X' \to S'$ 表示沿
$S' \to S$ 的基变换。假设 $f : Y \to X$ 固有。
则 $f' : Y' \to X'$ 固有，并且周群与 Chow 群的图
$$
\vcenter{
\xymatrix{
Z_k(Y) \ar[r]_{g^*} \ar[d]_{f_*} & Z_{k + c}(Y') \ar[d]^{f'_*} \\
Z_k(X) \ar[r]^{g^*} & Z_{k + c}(X')
}
}
\quad\text{且}\quad
\vcenter{
\xymatrix{
\CH_k(Y) \ar[r]_{g^*} \ar[d]_{f_*} & \CH_{k + c}(Y') \ar[d]^{f'_*} \\
\CH_k(X) \ar[r]^{g^*} & \CH_{k + c}(X')
}
}
$$
均交换。
\end{lemma}

\begin{proof}
只需证明第一个图交换。为此，设 $Z \subset Y$ 是
$\delta$-维数为 $k$ 的整闭子概形，并以 $Z' \subset X'$ 表示其基变换。
由构造，有 $g^*[Z] = [Z']_{k + c}$。由引理 \ref{chow-lemma-cycle-push-sheaf}，有
$(f')_*g^*[Z] = [f'_*\mathcal{O}_{Z'}]_{k + c}$。
由同一引理，有 $f_*[Z] = [f_*\mathcal{O}_Z]_k$；由
引理 \ref{chow-lemma-pullback-coherent-base-change}，有
$g^*f_*[Z] = [(X' \to X)^*f_*\mathcal{O}_Z]_{k + r}$。
因此只需证明 $X'$ 上凝聚模之间的同构
$$
(X' \to X)^*f_*\mathcal{O}_Z \cong f'_*\mathcal{O}_{Z'}
$$
由于 $X' \to X$ 平坦且
$\mathcal{O}_{Z'} = (Y' \to Y)^*\mathcal{O}_Z$，
该同构由平坦基变换得出；见《概形上同调》引理 \ref{coherent-lemma-flat-base-change-cohomology}。
\end{proof}

\begin{lemma}
\label{chow-lemma-pullback-base-change-c1}
在情形 \ref{chow-situation-setup-base-change} 中，设 $X \to S$ 局部有限型，
并以 $X' \to S'$ 表示沿 $S' \to S$ 的基变换。
设 $\mathcal{L}$ 是可逆 $\mathcal{O}_X$-模，其在 $X'$ 上的基变换为
$\mathcal{L}'$。则 Chow 群的图
$$
\xymatrix{
\CH_k(X) \ar[r]_{g^*} \ar[d]_{c_1(\mathcal{L}) \cap -} &
\CH_{k + c}(X') \ar[d]^{c_1(\mathcal{L}') \cap -} \\
\CH_{k - 1}(X) \ar[r]^{g^*} & \CH_{k + c - 1}(X')
}
$$
交换。
\end{lemma}

\begin{proof}
设 $p : L \to X$ 是与 $\mathcal{L}$ 相关联的线丛，
其零截面为 $o : X \to L$。对 $\alpha \in CH_k(X)$，
我们知道 $\beta = c_1(\mathcal{L}) \cap \alpha$ 是满足
$o_*\alpha = - p^*\beta$ 的唯一元素 $\CH_{k - 1}(X)$；见
引理 \ref{chow-lemma-linebundle} 和 \ref{chow-lemma-linebundle-formulae}。
拉回后仍有相同的刻画。因此引理由引理 \ref{chow-lemma-pullback-base-change-pullback} 和
\ref{chow-lemma-pullback-base-change-pushforward} 得出。
\end{proof}

\begin{lemma}
\label{chow-lemma-pullback-base-change-chern-classes}
在情形 \ref{chow-situation-setup-base-change} 中，设 $X \to S$ 局部有限型，
并以 $X' \to S'$ 表示沿 $S' \to S$ 的基变换。
设 $\mathcal{E}$ 是秩为 $r$ 的有限局部自由 $\mathcal{O}_X$-模，
其在 $X'$ 上的基变换为 $\mathcal{E}'$。则对所有 $i$，Chow 群的图
$$
\xymatrix{
\CH_k(X) \ar[r]_{g^*} \ar[d]_{c_i(\mathcal{E}) \cap -} &
\CH_{k + c}(X') \ar[d]^{c_i(\mathcal{E}') \cap -} \\
\CH_{k - i}(X) \ar[r]^{g^*} & \CH_{k + c - i}(X')
}
$$
交换。
\end{lemma}

\begin{proof}
置 $P = \mathbf{P}(\mathcal{E})$。$P$ 的基变换 $P'$ 等于
$\mathbf{P}(\mathcal{E}')$。我们已经知道，平坦拉回以及与可逆模的
$c_1$ 作杯积都与基变换可交换（引理 \ref{chow-lemma-pullback-base-change-pullback} 和
\ref{chow-lemma-pullback-base-change-c1}）。因此，该引理由
引理 \ref{chow-lemma-determine-intersections} 中对与
$c_i(\mathcal{E})$ 作帽积的刻画得出。
\end{proof}

\begin{lemma}
\label{chow-lemma-compose-base-change}
设 $(S, \delta)$、$(S', \delta')$、$(S'', \delta'')$ 如
情形 \ref{chow-situation-setup}。设 $g : S' \to S$ 和 $g' : S'' \to S'$
是概形的平坦态射，并设 $c, c' \in \mathbf{Z}$ 是整数，使得
$S, \delta, S', \delta', g, c$ 以及 $S', \delta', S'', g', c'$ 均如
情形 \ref{chow-situation-setup-base-change}。
设 $X \to S$ 局部有限型，并以 $X' \to S'$ 和 $X'' \to S''$
表示分别沿 $S' \to S$ 和 $S'' \to S$ 的基变换。则
\begin{enumerate}
\item $S, \delta, S'', \delta'', g \circ g', c + c'$ 如
情形 \ref{chow-situation-setup-base-change}；
\item 映射 $g^* : Z_k(X) \to Z_{k + c}(X')$ 与
$(g')^* : Z_{k + c}(X') \to Z_{k + c + c'}(X'')$ 的复合给出映射
$(g \circ g')^* : Z_k(X) \to Z_{k + c + c'}(X'')$；以及
\item 引理 \ref{chow-lemma-pullback-base-change} 的映射
$g^* : \CH_k(X) \to \CH_{k + c}(X')$ 与
$(g')^* : \CH_{k + c}(X') \to \CH_{k + c + c'}(X'')$ 的复合，
给出引理 \ref{chow-lemma-pullback-base-change} 的映射
$(g \circ g')^* : \CH_k(X) \to \CH_{k + c + c'}(X'')$。
\end{enumerate}
\end{lemma}

\begin{proof}
设 $s \in S$，并设 $s'' \in S''$ 是
$(g \circ g')^{-1}(\{s\})$ 某个不可约分支的一般点。置 $s' = g'(s'')$。
显然，$s''$ 是 $(g')^{-1}(\{s'\})$ 某个不可约分支的一般点。
此外，由于 $g'$ 平坦，一般化可沿 $g'$ 提升
（《概形态射》引理 \ref{morphisms-lemma-base-change-flat}），
所以 $s'$ 也是 $g^{-1}(\{s\})$ 某个不可约分支的一般点。
因此由假设，$\delta'(s') = \delta(s) + c$ 且
$\delta''(s'') = \delta'(s') + c'$。于是
$\delta''(s'') = \delta(s) + c + c'$，第一项断言成立。

\medskip\noindent
对于第二项，设 $Z \subset X$ 是 $\delta$-维数为 $k$ 的整闭子概形，
并以 $Z' \subset X'$ 和 $Z'' \subset X''$ 表示其基变换。
由定义，$g^*[Z] = [Z']_{k + c}$。由引理 \ref{chow-lemma-pullback-coherent-base-change}，有
$(g')^*[Z']_{k + c} = [Z'']_{k + c + c'}$。这证明了最后一项断言。
\end{proof}

\begin{lemma}
\label{chow-lemma-chow-limit}
在情形 \ref{chow-situation-setup-base-change} 中，假设 $c = 0$，
并假设 $S' = \lim_{i \in I} S_i$ 是 $S$ 上仿射概形 $S_i$ 的滤过极限，且
\begin{enumerate}
\item 当 $\delta_i$ 等于 $S_i \to S \xrightarrow{\delta} \mathbf{Z}$ 时，
序对 $(S_i, \delta_i)$ 如情形 \ref{chow-situation-setup}；
\item $S_i, \delta_i, S, \delta, S \to S_i, c = 0$ 如
情形 \ref{chow-situation-setup-base-change}；
\item 对 $i \geq i'$，$S_i, \delta_i, S_{i'}, \delta_{i'}, S_i \to S_{i'}, c = 0$
如情形 \ref{chow-situation-setup-base-change}。
\end{enumerate}
若 $X$ 是 $S$ 上拟紧有限型概形，并以 $X'$ 和 $X_i$ 表示分别沿
$S' \to S$ 和 $S_i \to S$ 的基变换，则
$Z_k(X') = \colim Z_k(X_i)$ 且 $\CH_k(X') = \colim \CH_k(X_i)$。
\end{lemma}

\begin{proof}
由引理 \ref{chow-lemma-compose-base-change}，得到周群系 $Z_k(X_i)$、
Chow 群系 $\CH_k(X_i)$，以及映射
$\colim Z_k(X_i) \to Z_k(X')$ 和 $\colim \CH_i(X_i) \to \CH_k(X')$。
可以把 $S$ 替换为 $X \to S$ 所通过的某个拟紧开子概形；
因此可以并且确实假设本证明中出现的所有概形都是 Noether 的
（从而拟紧且拟分离）。

\medskip\noindent
先证明映射 $\colim Z_k(X_i) \to Z_k(X')$ 满射。
设 $Z' \subset X'$ 是 $\delta'$-维数为 $k$ 的整闭子概形。
由《概形的极限》引理 \ref{limits-lemma-descend-finite-presentation}，
可找到某个 $i$ 以及有限表现态射 $Z_i \to X_i$，其基变换为 $Z'$。
增大 $i$ 后，可以假设 $Z_i$ 是 $X_i$ 的闭子概形；见
《概形的极限》引理 \ref{limits-lemma-descend-closed-immersion-finite-presentation}。
于是 $Z' \to X_i$ 通过 $Z_i$ 分解，并且可以用 $Z' \to X_i$ 的概形论像替换
$Z_i$。由此可见，可以假设 $Z_i$ 是 $X_i$ 的整闭子概形。
由引理 \ref{chow-lemma-dimension-base-change}，得到
$\dim_{\delta_i}(Z_i) = \dim_{\delta'}(Z') = k$。
因此，$Z_k(X_i) \to Z_k(X')$ 把 $[Z_i]$ 映到 $[Z']$，
从而满射性成立。

\medskip\noindent
下面证明映射 $\colim Z_k(X_i) \to Z_k(X')$ 单射。
设 $\alpha_i = \sum n_j[Z_j] \in Z_k(X_i)$ 是一个在
$Z_k(X')$ 中的像为零的周。可以并且确实假设：若 $j \not = j'$，
则 $Z_j \not = Z_{j'}$；并且对所有 $j$ 都有 $n_j \not = 0$。
以 $Z'_j \subset X'$ 表示 $Z_j$ 的基变换。
由引理 \ref{chow-lemma-dimension-base-change}，$Z'_j$ 的每个不可约分支都具有
$\delta'$-维数 $k$。此外，由于 $Z_j$ 不可约且 $Z'_j \to Z_j$ 平坦
（它是 $S' \to S$ 的基变换），$Z'_j \to Z_j$ 支配。
所以，若 $Z'_j$ 非空，则 $Z'_j$ 的某个不可约分支（记为 $Z'$）支配 $Z_j$。
于是当 $j' \not = j$ 时，$Z'$ 不可能是 $Z'_{j'}$ 的不可约分支。
因此，若 $Z'_j$ 非空，则
$(S' \to S_i)^*\alpha_i
= \sum [Z'_j]_r$ 非零（因为 $Z'$ 的系数将非零）。
故对所有 $j$ 都有 $Z'_j = \emptyset$。然而，由《概形的极限》引理 \ref{limits-lemma-limit-nonempty}，这意味着存在某个过渡态射
$S_{i'} \to S_i$，使 $Z_j$ 的基变换为空。
因此 $\alpha_i$ 在余极限中消失，正合所需。

\medskip\noindent
映射 $\colim Z_k(X_i) \to Z_k(X')$ 的满射性蕴含
$\colim \CH_k(X_i) \to \CH_k(X')$ 满射。
为完成证明，须说明后一个映射单射。
设 $\alpha_i \in \CH_k(X_i)$ 是一个周，其像
$\alpha' \in \CH_k(X')$ 为零。则存在 $\delta"$-维数为 $k + 1$ 的整闭子概形
$W_l' \subset X'$、$l = 1, \ldots, r$，以及 $W'_l$ 上的非零有理函数
$f'_l$，使得
$\alpha' = \sum_{l = 1, \ldots, r} \text{div}_{W'_l}(f'_l)$。
重复上述论证，可找到某个 $i$ 以及 $\delta_i$-维数为 $k + 1$ 的整闭子概形
$W_{i, l} \subset X_i$，其基变换为 $W'_l$。
增大 $i$ 后，可以假设在 $W_{i, l}$ 上有有理函数 $f_{i, l}$。
具体地，可将 $f'_l$ 视为非空开子概形 $U'_l \subset W'_l$ 上的结构层截面；
由《概形的极限》引理 \ref{limits-lemma-descend-opens}，这些开子概形可以下降；
再增大 $i$ 后，由《概形的极限》引理 \ref{limits-lemma-descend-section}，
$f'_l$ 也可以下降。我们断言，可能再次增大 $i$ 后，有
$$
\alpha_i = \sum\nolimits_{l = 1, \ldots, r} \text{div}_{W_{i, l}}(f_{i, l})
$$

\medskip\noindent
为证明该断言，取 $\delta'$-维数为 $k$ 的整闭子概形的有限族
$Z'_{l, j} \subset W'_l$，使 $f'_l$ 在
$\bigcup_j Y'_{l, j}$ 之外是可逆正则函数。
增大 $i$ 后（利用上述论证），可以假设存在 $\delta_i$-维数为 $k$ 的整闭子概形
$Z_{i, l, j} \subset W_i$，使 $f_{i, l}$ 在
$\bigcup_j Z_{i, l, j}$ 之外是可逆正则函数。于是可写成
$$
\text{div}_{W'_l}(f'_l) = \sum n_{l, j} [Z'_{l, j}]
$$
以及
$$
\text{div}_{W_{i, l}}(f_{i, l}) = \sum n_{i, l, j} [Z_{i, l, j}]
$$
为证明该断言，只需说明 $n_{l, i} = n_{i, l, j}$。
事实上，这将蕴含
$\beta_i =
\alpha_i - \sum\nolimits_{l = 1, \ldots, r} \text{div}_{W_{i, l}}(f_{i, l})$
是 $X_i$ 上的周，并且其拉回在 $X'$ 上作为周为零！
由一个从略的简单论证可知，存在某个 $i' \geq i$，使得
$\beta_i$ 在 $X_{i'}$ 上的拉回作为周为零。

\medskip\noindent
为证明等式 $n_{l, i} = n_{i, l, j}$，选取一般点
$\xi' \in Z'_{l, j}$，并以 $\xi \in Z_{i, l, j}$ 表示其像；后者也是一般点。
此时局部环映射
$$
\mathcal{O}_{W_{i, l}, \xi}
\longrightarrow
\mathcal{O}_{W'_l, \xi'}
$$
平坦，因为 $W'_l \to W_{i, l}$ 是平坦态射 $S' \to S_i$ 的基变换。
还因为 $Z_{i, l, j}$ 拉回为 $Z'_{l, j}$，有
$\mathfrak m_\xi \mathcal{O}_{W'_l, \xi'} = \mathfrak m_{\xi'}$！
因此，等式
$$
n_{l, j} = \text{ord}_{Z'_{l, j}}(f'_l) =
\text{ord}_{\mathcal{O}_{W'_l, \xi'}}(f'_l)
\quad\text{且}\quad
n_{i, l, j} = \text{ord}_{Z_{i, l, j}}(f_{i, l}) =
\text{ord}_{\mathcal{O}_{W_{i, l}, \xi}}(f_{i, l})
$$
由《交换代数》引理 \ref{algebra-lemma-pullback-module} 以及
《交换代数》第 \ref{algebra-section-orders-of-vanishing} 节中
$\text{ord}$ 的构造得出。
\end{proof}












\section{附录 A：关键引理的另一种处理方法}
\label{chow-section-appendix-A}

\noindent
在本附录中，我们首先定义局部环 $(R, \mathfrak m, \kappa)$ 上有限长度模
$M$ 的行列式 $\det_\kappa(M)$，见小节
\ref{chow-subsection-determinants-finite-length}。
行列式 $\det_\kappa(M)$ 是 $1$ 维 $\kappa$-向量空间。
在小节 \ref{chow-subsection-periodic-complexes-determinants} 中，我们用它定义
有限长度模 $M$ 上正合 $(2, 1)$-周期复形 $(M, \varphi, \psi)$ 的行列式
$\det_\kappa(M, \varphi, \psi) \in \kappa^*$。
在小节 \ref{chow-subsection-symbols} 中，当 $R$ 是 $1$ 维诺特环时，我们利用这些
行列式为一对非零因子 $a, b \in R$ 构造驯符号
$d_R(a, b) = \det_\kappa(R/ab, a, b)$。
尽管无疑有
$$
d_R(a, b) = \partial_R(a, b)
$$
其中 $\partial_R$ 如第 \ref{chow-section-tame-symbol} 节所定义，但我们（尚）未加入
验证。本附录所构造的驯符号的优点是，它可以推广到有限 $R$-模 $M$ 的一对
单射自同态 $\varphi, \psi$（例如），其中该模的维数为 $1$，且
$\varphi(\psi(M)) = \psi(\varphi(M))$。在小节
\ref{chow-subsection-length-determinant} 中，我们把 Herbrand 商与行列式联系起来。
主结果（命题 \ref{chow-proposition-length-determinant-periodic-complex}）的一种
易于陈述的形式是公式
$$
-e_R(M, \varphi, \psi) =
\text{ord}_R(\det\nolimits_K(M_K, \varphi, \psi))
$$
这里 $(M, \varphi, \psi)$ 是定义在 $1$ 维诺特局部整环 $R$ 上的
$(2, 1)$-周期复形，其分式域为 $K$，且该复形的 Herbrand 商 $e_R$
（定义 \ref{chow-definition-periodic-length}）有定义。
我们用这一命题，对本附录所构造的驯符号给出关键引理
（引理 \ref{chow-lemma-milnor-gersten-low-degree}）的另一种证明；见引理
\ref{chow-lemma-secondary-ramification}。


\subsection{有限长度模的行列式}
\label{chow-subsection-determinants-finite-length}

\noindent
本节内容与论文 \cite{determinant} 及学位论文 \cite{Joe} 中的内容有关。

\medskip\noindent
设 $(R, \mathfrak m, \kappa)$ 为局部环，设
$\varphi : M \to M$ 是有限长度 $R$-模 $M$ 的 $R$-线性自同态。
在《More on Algebra》第
\ref{more-algebra-section-determinants-finite-length} 节中，我们已经定义了
$\varphi$ 关于 $\kappa$ 的行列式 $\det_\kappa(\varphi)$
（以及迹和特征多项式）。本节将构造一个典范的 $1$ 维 $\kappa$-向量空间
$\det_\kappa(M)$，使得
$\det_\kappa(\varphi : M \to M) : \det_\kappa(M) \to \det_\kappa(M)$
等于乘以 $\det_\kappa(\varphi)$ 的映射。
若 $M$ 被 $\mathfrak m$ 零化，则可把 $M$ 视为有限维 $\kappa$-向量空间，
此时有 $\det_\kappa(M) = \wedge^n_\kappa(M)$，其中
$n = \dim_\kappa(M)$。我们的构造将把这一点推广到 $R$ 上所有有限长度模；
若 $R$ 包含其剩余域，则在适当意义下，行列式 $\det_\kappa(M)$ 由通常的
行列式给出，见注 \ref{chow-remark-explain-determinant}。

\begin{definition}
\label{chow-definition-determinant}
设 $R$ 为局部环，其极大理想为 $\mathfrak m$，剩余域为 $\kappa$。
设 $M$ 为有限长度 $R$-模，并记 $l = \text{length}_R(M)$。
\begin{enumerate}
\item 给定元素 $x_1, \ldots, x_r \in M$，我们用
$\langle x_1, \ldots, x_r \rangle = Rx_1 + \ldots + Rx_r$ 表示由
$x_1, \ldots, x_r$ 生成的 $M$ 的 $R$-子模。
\item 若 $M$ 中元素的一个 $l$-元组 $(e_1, \ldots, e_l)$ 对每个
$i = 1, \ldots, l$ 都满足
$\mathfrak m e_i \subset \langle e_1, \ldots, e_{i - 1} \rangle$，
则称它是{\it 容许的}。
\item 一个{\it 符号} $[e_1, \ldots, e_l]$ 意指
$(e_1, \ldots, e_l)$ 是容许 $l$-元组。
\item 符号之间的{\it 容许关系}是下列关系之一：
\begin{enumerate}
\item 若 $(e_1, \ldots, e_l)$ 是容许序列，且对某个
$1 \leq a \leq l$ 有
$e_a \in \langle e_1, \ldots, e_{a - 1}\rangle$，则
$[e_1, \ldots, e_l] = 0$；
\item 若 $(e_1, \ldots, e_l)$ 是容许序列，且对某个
$1 \leq a \leq l$ 有 $e_a = \lambda e'_a + x$，其中
$\lambda \in R^*$，并且
$x \in \langle e_1, \ldots, e_{a - 1}\rangle$，则
$$
[e_1, \ldots, e_l] =
\overline{\lambda} [e_1, \ldots, e_{a - 1}, e'_a, e_{a + 1}, \ldots, e_l]
$$
其中 $\overline{\lambda} \in \kappa^*$ 是 $\lambda$ 在剩余域中的像；
\item 若 $(e_1, \ldots, e_l)$ 是容许序列，且
$\mathfrak m e_a \subset \langle e_1, \ldots, e_{a - 2}\rangle$，则
$$
[e_1, \ldots, e_l] =
- [e_1, \ldots, e_{a - 2}, e_a, e_{a - 1}, e_{a + 1}, \ldots, e_l].
$$
\end{enumerate}
\item
我们把有限长度 $R$-模 $M$ 的{\it 行列式}定义为
$$
\det\nolimits_\kappa(M) =
\left\{
\frac{\text{由符号生成的 }\kappa\text{-向量空间}}
{\text{可容许关系的 }\kappa\text{-线性组合}}
\right\}
$$
\end{enumerate}
\end{definition}

\noindent
我们强调，始终有 $l = \text{length}_R(M)$。还要强调的是，这并不推出
符号 $[e_1, \ldots, e_l]$ 对各个分量具有可加性（通常并非如此）。
在证明行列式 $\det_\kappa(M)$ 的维数确实为 $1$ 之前，我们必须先证明
它的维数至多为 $1$。

\begin{lemma}
\label{chow-lemma-dimension-at-most-one}
沿用上述记号，有 $\dim_\kappa(\det_\kappa(M)) \leq 1$。
\end{lemma}

\begin{proof}
固定 $M$ 的一个容许序列 $(f_1, \ldots, f_l)$，使得
$$
\text{length}_R(\langle f_1, \ldots, f_i\rangle) = i
$$
对 $i = 1, \ldots, l$ 成立。这样的容许序列存在，恰是因为 $M$ 的长度为
$l$。我们将证明，$\det_\kappa(M)$ 的任意元素都是符号
$[f_1, \ldots, f_l]$ 的一个 $\kappa$-倍数。这就证明了引理。

\medskip\noindent
设 $(e_1, \ldots, e_l)$ 是 $M$ 的容许序列。只需证明
$[e_1, \ldots, e_l]$ 是 $[f_1, \ldots, f_l]$ 的倍数。
先假设 $\langle e_1, \ldots, e_l\rangle \not = M$。于是存在
$i \in [1, \ldots, l]$，使得
$e_i \in \langle e_1, \ldots, e_{i - 1}\rangle$。由第一条容许关系立即得到
$[e_1, \ldots, e_n] = 0$ 在 $\det_\kappa(M)$ 中成立。
因此可假设 $\langle e_1, \ldots, e_l\rangle = M$。特别地，存在最小指标
$i \in \{1, \ldots, l\}$，使得
$f_1 \in \langle e_1, \ldots, e_i\rangle$。这意味着
$e_i = \lambda f_1 + x$，其中
$x \in \langle e_1, \ldots, e_{i - 1}\rangle$ 且 $\lambda \in R^*$。
由第二条容许关系，这意味着
$[e_1, \ldots, e_l] =
\overline{\lambda}[e_1, \ldots, e_{i - 1}, f_1, e_{i + 1}, \ldots, e_l]$。
注意 $\mathfrak m f_1 = 0$。因此，应用第三条容许关系 $i - 1$ 次可得
$$
[e_1, \ldots, e_l] =
(-1)^{i - 1}\overline{\lambda}
[f_1, e_1, \ldots, e_{i - 1}, e_{i + 1}, \ldots, e_l].
$$
还要注意，
$ \langle f_1, e_1, \ldots, e_{i - 1}, e_{i + 1}, \ldots, e_l\rangle = M$
也成立。归纳假设已经证明，原符号等于某个标量乘以
$$
[f_1, \ldots, f_j, e_{j + 1}, \ldots, e_l]
$$
其中 $(f_1, \ldots, f_j, e_{j + 1}, \ldots, e_l)$ 是元素生成 $M$ 的
某个容许序列，即，\ 满足
$\langle f_1, \ldots, f_j, e_{j + 1}, \ldots, e_l\rangle = M$。
然后取最小的 $i$，使得
$f_{j + 1} \in \langle f_1, \ldots, f_j, e_{j + 1}, \ldots, e_i\rangle$，
并重复上述过程，得到
$$
[f_1, \ldots, f_j, e_{j + 1}, \ldots, e_l]
=
(\text{标量}) [f_1, \ldots, f_j, f_{j + 1}, e_{j + 1},
\ldots, \hat{e_i}, \ldots, e_l]
$$
如此继续，便得到所需结论。
\end{proof}

\noindent
在证明 $\det_\kappa(M)$ 的维数总是 $1$ 之前，先证明当该模是
$\kappa$ 上的向量空间时，它与通常的最高外幂一致。

\begin{lemma}
\label{chow-lemma-compare-det}
设 $R$ 为局部环，其极大理想为 $\mathfrak m$，剩余域为 $\kappa$。
设 $M$ 是被 $\mathfrak m$ 零化的有限长度 $R$-模，并设
$l = \dim_\kappa(M)$。则映射
$$
\det\nolimits_\kappa(M) \longrightarrow \wedge^l_\kappa(M),
\quad
[e_1, \ldots, e_l] \longmapsto e_1 \wedge \ldots \wedge e_l
$$
是同构。
\end{lemma}

\begin{proof}
引理中所述规则显然给出一个 $\kappa$-线性映射，因为通常的符号
$e_1 \wedge \ldots \wedge e_l$ 满足全部容许关系。该映射也显然满射。
另一方面，由引理 \ref{chow-lemma-dimension-at-most-one}，左端的维数至多为一，
故此映射是同构。
\end{proof}

\begin{lemma}
\label{chow-lemma-determinant-dimension-one}
设 $R$ 为局部环，其极大理想为 $\mathfrak m$，剩余域为 $\kappa$。
设 $M$ 为有限长度 $R$-模。上面定义的行列式 $\det_\kappa(M)$ 是维数为
$1$ 的 $\kappa$-向量空间。对任意满足
$\langle f_1, \ldots f_l \rangle = M$ 的容许序列，符号
$[f_1, \ldots, f_l]$ 都生成该空间。
\end{lemma}

\begin{proof}
由引理 \ref{chow-lemma-dimension-at-most-one} 及其证明可知，
$\det_\kappa(M)$ 的维数至多为 $1$，而且它实际上由
$[f_1, \ldots, f_l]$ 生成。我们将对 $l = \text{length}(M)$ 归纳，
证明它非零。当 $l = 1$ 时，结论由引理 \ref{chow-lemma-compare-det} 得出。
选取非零元素 $f \in M$，使得 $\mathfrak m f = 0$。令
$\overline{M} = M /\langle f \rangle$，并把商映射记为
$x \mapsto \overline{x}$。我们将定义一个满射
$$
\psi : \det\nolimits_k(M) \to \det\nolimits_\kappa(\overline{M})
$$
由归纳假设，$\overline{M}$ 的行列式非零，因而这将证明引理。

\medskip\noindent
我们如下在符号上定义 $\psi$。设 $(e_1, \ldots, e_l)$ 是容许序列。
若 $f \not \in \langle e_1, \ldots, e_l \rangle$，就令
$\psi([e_1, \ldots, e_l]) = 0$。若
$f \in \langle e_1, \ldots, e_l \rangle$，则选取使
$f \in \langle e_1, \ldots, e_i \rangle$ 成立的最小 $i$。
可写成 $e_i = \lambda f + x$，其中 $\lambda \in R$ 是单位，且
$x \in \langle e_1, \ldots, e_{i - 1} \rangle$。在此情形令
$$
\psi([e_1, \ldots, e_l]) =
(-1)^i
\overline{\lambda}[\overline{e}_1, \ldots,
\overline{e}_{i - 1},
\overline{e}_{i + 1}, \ldots, \overline{e}_l].
$$
确实，$(\overline{e}_1, \ldots,
\overline{e}_{i - 1},
\overline{e}_{i + 1}, \ldots, \overline{e}_l)$ 是
$\overline{M}$ 中的容许序列，所以该定义有意义。下面证明，将此规则按
$\kappa$-线性方式延拓到符号的线性组合，确实会在行列式上给出映射。
为此必须证明容许关系都被映为零。

\medskip\noindent
类型 (a) 的关系。设 $(e_1, \ldots, e_l)$ 是容许序列，且对某个
$1 \leq a \leq l$ 有
$e_a \in \langle e_1, \ldots, e_{a - 1}\rangle$。
设 $i$ 是使 $f \in \langle e_1, \ldots, e_i\rangle$ 成立的最小指标。
则 $i \not = a$，而且当 $i < a$ 时，
$\overline{e}_a \in \langle \overline{e}_1, \ldots,
\hat{\overline{e}_i}, \ldots, \overline{e}_{a - 1}\rangle$；
当 $i > a$ 时，
$\overline{e}_a \in \langle \overline{e}_1, \ldots,
\overline{e}_{a - 1}\rangle$。因此，$\det_\kappa(\overline{M})$ 中的
相应容许关系迫使符号 $[\overline{e}_1, \ldots,
\overline{e}_{i - 1},
\overline{e}_{i + 1}, \ldots, \overline{e}_l]$ 为零，正如所需。

\medskip\noindent
类型 (b) 的关系。设 $(e_1, \ldots, e_l)$ 是容许序列，且对某个
$1 \leq a \leq l$ 有 $e_a = \lambda e'_a + x$，其中
$\lambda \in R^*$，并且
$x \in \langle e_1, \ldots, e_{a - 1}\rangle$。
设 $i$ 是使 $f \in \langle e_1, \ldots, e_i\rangle$ 成立的最小指标，
并写成 $e_i = \mu f + y$，其中
$y \in \langle e_1, \ldots, e_{i - 1}\rangle$。若 $i < a$，则所需等式为
$$
(-1)^i
\overline{\lambda}
[\overline{e}_1,
\ldots,
\overline{e}_{i - 1},
\overline{e}_{i + 1},
\ldots,
\overline{e}_l]
=
(-1)^i
\overline{\lambda}
[\overline{e}_1,
\ldots,
\overline{e}_{i - 1},
\overline{e}_{i + 1},
\ldots,
\overline{e}_{a - 1},
\overline{e}'_a,
\overline{e}_{a + 1},
\ldots,
\overline{e}_l]
$$
它由 $\overline{e}_a = \lambda \overline{e}'_a + \overline{x}$ 以及
$\det_\kappa(\overline{M})$ 中相应的容许关系得出。若 $i > a$，
则所需等式为
$$
(-1)^i
\overline{\lambda}
[\overline{e}_1,
\ldots,
\overline{e}_{i - 1},
\overline{e}_{i + 1},
\ldots,
\overline{e}_l]
=
(-1)^i
\overline{\lambda}
[\overline{e}_1,
\ldots,
\overline{e}_{a - 1},
\overline{e}'_a,
\overline{e}_{a + 1},
\ldots,
\overline{e}_{i - 1},
\overline{e}_{i + 1},
\ldots,
\overline{e}_l]
$$
它同样由 $\overline{e}_a = \lambda \overline{e}'_a + \overline{x}$ 以及
$\det_\kappa(\overline{M})$ 中相应的容许关系得出。需要特别处理的是
$i = a$ 的情形。此时
$e_a = \lambda e'_a + x = \mu f + y$，从而还有
$e'_a = \lambda^{-1}(\mu f + y - x)$。于是
$$
\psi([e_1, \ldots, e_l])
= (-1)^i \overline{\mu}
[\overline{e}_1, \ldots,
\overline{e}_{i - 1},
\overline{e}_{i + 1}, \ldots, \overline{e}_l]
=
\psi(
\overline{\lambda}
[e_1, \ldots, e_{a - 1}, e'_a, e_{a + 1}, \ldots, e_l]
)
$$
正如所需。

\medskip\noindent
类型 (c) 的关系。设 $(e_1, \ldots, e_l)$ 是容许序列，且
$\mathfrak m e_a \subset \langle e_1, \ldots, e_{a - 2}\rangle$。
设 $i$ 是使 $f \in \langle e_1, \ldots, e_i\rangle$ 成立的最小指标。
写成 $e_i = \lambda f + x$，其中
$x \in \langle e_1, \ldots, e_{i - 1}\rangle$。分成 $4$ 种情形：

\medskip\noindent
情形 1：$i < a - 1$。所需等式为
\begin{align*}
& (-1)^i
\overline{\lambda}
[\overline{e}_1,
\ldots,
\overline{e}_{i - 1},
\overline{e}_{i + 1},
\ldots,
\overline{e}_l] \\
& =
(-1)^{i + 1}
\overline{\lambda}
[\overline{e}_1,
\ldots,
\overline{e}_{i - 1},
\overline{e}_{i + 1},
\ldots,
\overline{e}_{a - 2},
\overline{e}_a,
\overline{e}_{a - 1},
\overline{e}_{a + 1},
\ldots,
\overline{e}_l]
\end{align*}
它由 $\det_\kappa(\overline{M})$ 中类型 (c) 的容许关系得出。

\medskip\noindent
情形 2：$i > a$。所需等式为
\begin{align*}
& (-1)^i
\overline{\lambda}
[\overline{e}_1,
\ldots,
\overline{e}_{i - 1},
\overline{e}_{i + 1},
\ldots,
\overline{e}_l] \\
& =
(-1)^{i + 1}
\overline{\lambda}
[\overline{e}_1,
\ldots,
\overline{e}_{a - 2},
\overline{e}_a,
\overline{e}_{a - 1},
\overline{e}_{a + 1},
\ldots,
\overline{e}_{i - 1},
\overline{e}_{i + 1},
\ldots,
\overline{e}_l]
\end{align*}
它由 $\det_\kappa(\overline{M})$ 中类型 (c) 的容许关系得出。

\medskip\noindent
情形 3：$i = a$。写成 $e_a = \lambda f + \mu e_{a - 1} + y$，其中
$y \in \langle e_1, \ldots, e_{a - 2}\rangle$。于是根据定义，
$$
\psi([e_1, \ldots, e_l]) =
(-1)^a
\overline{\lambda}
[\overline{e}_1,
\ldots,
\overline{e}_{a - 1},
\overline{e}_{a + 1},
\ldots,
\overline{e}_l]
$$
若 $\overline{\mu}$ 非零，则
$e_{a - 1} = - \mu^{-1} \lambda f + \mu^{-1}e_a - \mu^{-1} y$，
因而根据定义得到
$$
\psi(-[e_1, \ldots, e_{a - 2}, e_a, e_{a - 1}, e_{a + 1}, \ldots, e_l]) =
(-1)^a
\overline{\mu^{-1}\lambda}
[\overline{e}_1,
\ldots,
\overline{e}_{a - 2},
\overline{e}_a,
\overline{e}_{a + 1},
\ldots,
\overline{e}_l]
$$
由于在 $\overline{M}$ 中有
$\overline{e}_a = \mu \overline{e}_{a - 1} + \overline{y}$，由
$\det_\kappa(\overline{M})$ 的关系 (a) 可知两个结果相等。
另一方面，若 $\overline{\mu}$ 为零，则可写成
$e_a = \lambda f + y$，其中
$y \in \langle e_1, \ldots, e_{a - 2}\rangle$，而且
$$
\psi(-[e_1, \ldots, e_{a - 2}, e_a, e_{a - 1}, e_{a + 1}, \ldots, e_l]) =
(-1)^a
\overline{\lambda}
[\overline{e}_1,
\ldots,
\overline{e}_{a - 1},
\overline{e}_{a + 1},
\ldots,
\overline{e}_l]
$$
它等于 $\psi([e_1, \ldots, e_l])$。

\medskip\noindent
情形 4：$i = a - 1$。此时根据定义，
$$
\psi([e_1, \ldots, e_l]) =
(-1)^{a - 1}
\overline{\lambda}
[\overline{e}_1,
\ldots,
\overline{e}_{a - 2},
\overline{e}_a,
\ldots,
\overline{e}_l]
$$
若 $f \not \in \langle e_1, \ldots, e_{a - 2}, e_a \rangle$，则
$$
\psi(-[e_1, \ldots, e_{a - 2}, e_a, e_{a - 1}, e_{a + 1}, \ldots, e_l]) =
(-1)^{a + 1}\overline{\lambda}
[\overline{e}_1,
\ldots,
\overline{e}_{a - 2},
\overline{e}_a,
\ldots,
\overline{e}_l]
$$
由于 $(-1)^{a - 1} = (-1)^{a + 1}$，两个表达式相同。
最后假设 $f \in \langle e_1, \ldots, e_{a - 2}, e_a \rangle$。
此时可见 $e_{a - 1} = \lambda f + x$，其中
$x \in \langle e_1, \ldots, e_{a - 2}\rangle$，并且
$e_a = \mu f + y$，其中
$y \in \langle e_1, \ldots, e_{a - 2}\rangle$，这里
$\lambda, \mu \in R$ 均为单位。于是同时有
$e_a \in \langle e_1, \ldots, e_{a - 1} \rangle$ 和
$e_{a - 1} \in \langle e_1, \ldots, e_{a - 2}, e_a\rangle$。
在这种情形，类型 (a) 的关系同时适用于
$[e_1, \ldots, e_l]$ 与
$[e_1, \ldots, e_{a - 2}, e_a, e_{a - 1}, e_{a + 1}, \ldots, e_l]$。
结合上面已经证明的 $\psi$ 与这些关系的相容性，可见
$$
\psi([e_1, \ldots, e_l])
\quad\text{且}\quad
\psi([e_1, \ldots, e_{a - 2}, e_a, e_{a - 1}, e_{a + 1}, \ldots, e_l])
$$
二者都为零，正如所需。

\medskip\noindent
至此已经证明 $\psi$ 良定义，只剩证明它是满射。为此，设
$(\overline{f}_2, \ldots, \overline{f}_l)$ 是 $\overline{M}$ 中的
容许序列。可以选取提升 $f_2, \ldots, f_l \in M$，于是
$(f, f_2, \ldots, f_l)$ 是 $M$ 中的容许序列。由于
$\psi([f, f_2, \ldots, f_l]) = [f_2, \ldots, f_l]$，结论成立。
\end{proof}

\noindent
设 $R$ 为局部环，其极大理想为 $\mathfrak m$，剩余域为 $\kappa$。
注意，若 $\varphi : M \to N$ 是有限长度 $R$-模的同构，则由规则
$$
\det\nolimits_\kappa(\varphi) :
\det\nolimits_\kappa(M)
\to
\det\nolimits_\kappa(N)
$$
以及
$$
\det\nolimits_\kappa(\varphi)([e_1, \ldots, e_l])
=
[\varphi(e_1), \ldots, \varphi(e_l)]
$$
对 $M$ 的任意符号 $[e_1, \ldots, e_l]$，得到一个同构。
因此可见，$\det\nolimits_\kappa$ 是函子
\begin{equation}
\label{chow-equation-functor}
\left\{
\begin{matrix}
\text{有限长度 }R\text{-模}\\
\text{（以同构为态射）}
\end{matrix}
\right\}
\longrightarrow
\left\{
\begin{matrix}
1\text{-维 }\kappa\text{-向量空间}\\
\text{（以同构为态射）}
\end{matrix}
\right\}
\end{equation}
这是“行列式函子”的典型情形（见 \cite{Knudsen}），下面的可加性也是如此。

\begin{lemma}
\label{chow-lemma-det-exact-sequences}
设 $(R, \mathfrak m, \kappa)$ 为局部环。对有限长度 $R$-模的每个短正合列
$$
0 \to K \to L \to M \to 0
$$
都存在一个典范同构
$$
\gamma_{K \to L \to M} :
\det\nolimits_\kappa(K) \otimes_\kappa \det\nolimits_\kappa(M)
\longrightarrow
\det\nolimits_\kappa(L)
$$
它在非零符号上由规则
$$
[e_1, \ldots, e_k]
\otimes
[\overline{f}_1, \ldots, \overline{f}_m]
\longrightarrow
[e_1, \ldots, e_k, f_1, \ldots, f_m]
$$
定义，并具有下列性质：
\begin{enumerate}
\item 对短正合列的每个同构，即对每个交换图
$$
\xymatrix{
0 \ar[r] &
K \ar[r] \ar[d]^u &
L \ar[r] \ar[d]^v &
M \ar[r] \ar[d]^w &
0 \\
0 \ar[r] &
K' \ar[r] &
L' \ar[r] &
M' \ar[r] &
0
}
$$
若两行短正合且 $u, v, w$ 均为同构，则
$$
\gamma_{K' \to L' \to M'} \circ
(\det\nolimits_\kappa(u) \otimes \det\nolimits_\kappa(w))
=
\det\nolimits_\kappa(v) \circ
\gamma_{K \to L \to M},
$$
\item 对有限长度 $R$-模的每个行列均正合的交换方图
$$
\xymatrix{
& 0 \ar[d] & 0 \ar[d] & 0 \ar[d] & \\
0 \ar[r] & A \ar[r] \ar[d] & B \ar[r] \ar[d] & C \ar[r] \ar[d] & 0 \\
0 \ar[r] & D \ar[r] \ar[d] & E \ar[r] \ar[d] & F \ar[r] \ar[d] & 0 \\
0 \ar[r] & G \ar[r] \ar[d] & H \ar[r] \ar[d] & I \ar[r] \ar[d] & 0 \\
& 0  & 0  & 0  &
}
$$
下图交换：
$$
\xymatrix{
\det\nolimits_\kappa(A) \otimes
\det\nolimits_\kappa(C) \otimes
\det\nolimits_\kappa(G) \otimes
\det\nolimits_\kappa(I)
\ar[dd]_{\epsilon}
\ar[rrr]_-{\gamma_{A \to B \to C} \otimes \gamma_{G \to H \to I}}
& & &
\det\nolimits_\kappa(B) \otimes
\det\nolimits_\kappa(H)
\ar[d]^{\gamma_{B \to E \to H}}
\\
& & & \det\nolimits_\kappa(E)
\\
\det\nolimits_\kappa(A) \otimes
\det\nolimits_\kappa(G) \otimes
\det\nolimits_\kappa(C) \otimes
\det\nolimits_\kappa(I)
\ar[rrr]^-{\gamma_{A \to D \to G} \otimes \gamma_{C \to F \to I}}
& & &
\det\nolimits_\kappa(D) \otimes
\det\nolimits_\kappa(F)
\ar[u]_{\gamma_{D \to E \to F}}
}
$$
其中 $\epsilon$ 是交换张量积因子的映射再乘以 $(-1)^{cg}$，这里
$c = \text{length}_R(C)$ 且 $g = \text{length}_R(G)$；
\item 若 $0 \to K \to L \to M \to 0$ 实际上是 $\kappa$-向量空间的
短正合列，则映射 $\gamma_{K \to L \to M}$ 与通常的同构一致。
\end{enumerate}
\end{lemma}

\begin{proof}
在映射 $\gamma_{K \to L \to M}$ 的显式描述中采用非零符号，原因如下。
若 $(e_1, \ldots, e_l)$ 是 $K$ 中的容许序列，而
$(\overline{f}_1, \ldots, \overline{f}_m)$ 是 $M$ 中的容许序列，
则 $(e_1, \ldots, e_l, f_1, \ldots, f_m)$ 不一定是 $L$ 中的容许序列
（这里 $f_i \in L$ 当然表示 $\overline{f}_i$ 的一个提升）。然而，若符号
$[e_1, \ldots, e_l]$ 在 $\det_\kappa(K)$ 中非零，则必有
$K = \langle e_1, \ldots, e_k\rangle$（见引理
\ref{chow-lemma-dimension-at-most-one} 的证明），在这种情形，
$(e_1, \ldots, e_k, f_1, \ldots, f_m)$ 的确是容许序列。
此外，由 $\det_\kappa(L)$ 中类型 (b) 的容许关系可见，
$[e_1, \ldots, e_k, f_1, \ldots, f_m]$ 在 $\det_\kappa(L)$ 中的值
也与提升 $f_i$ 的选取无关。根据这一点，很明显，$K$ 中
$e_1, \ldots, e_k$ 之间的容许关系会转化为 $L$ 中
$e_1, \ldots, e_k, f_1, \ldots, f_m$ 之间的容许关系；
$\overline{f}_1, \ldots, \overline{f}_m$ 之间的容许关系亦然。
因此，$\gamma$ 按引理所述定义了向量空间的线性映射。

\medskip\noindent
由引理 \ref{chow-lemma-determinant-dimension-one}，$\det_\kappa(L)$ 由任意一个
符号 $[x_1, \ldots, x_{k + m}]$ 生成，只要
$(x_1, \ldots, x_{k + m})$ 是满足
$L = \langle x_1, \ldots, x_{k + m}\rangle$ 的容许序列。
所以映射 $\gamma_{K \to L \to M}$ 显然满射，因而是同构。

\medskip\noindent
性质 (1) 成立，因为
\begin{eqnarray*}
& & \det\nolimits_\kappa(v)([e_1, \ldots, e_k, f_1, \ldots, f_m]) \\
& = &
[v(e_1), \ldots, v(e_k), v(f_1), \ldots, v(f_m)] \\
& = &
\gamma_{K' \to L' \to M'}([u(e_1), \ldots, u(e_k)]
\otimes [w(f_1), \ldots, w(f_m)]).
\end{eqnarray*}
性质 (2) 意味着：给定生成 $\det_\kappa(A)$ 的符号
$[\alpha_1, \ldots, \alpha_a]$、生成 $\det_\kappa(C)$ 的符号
$[\gamma_1, \ldots, \gamma_c]$、生成 $\det_\kappa(G)$ 的符号
$[\zeta_1, \ldots, \zeta_g]$，以及生成 $\det_\kappa(I)$ 的符号
$[\iota_1, \ldots, \iota_i]$，有
\begin{eqnarray*}
& & [\alpha_1, \ldots, \alpha_a, \tilde\gamma_1, \ldots, \tilde\gamma_c,
\tilde\zeta_1, \ldots, \tilde\zeta_g, \tilde\iota_1, \ldots, \tilde\iota_i] \\
& = &
(-1)^{cg} [\alpha_1, \ldots, \alpha_a, \tilde\zeta_1, \ldots, \tilde\zeta_g,
\tilde\gamma_1, \ldots, \tilde\gamma_c, \tilde\iota_1, \ldots, \tilde\iota_i]
\end{eqnarray*}
在 $\det_\kappa(E)$ 中成立（其中 $\tilde{x}$ 是 $E$ 中适当的提升）。
这是因为可以按下列次序使用类型 (c) 的容许关系 $cg$ 次：先将
$\tilde\zeta_1$ 越过各元素
$\tilde\gamma_c, \ldots, \tilde\gamma_1$
（这是允许的，因为 $\mathfrak m\tilde\zeta_1 \subset A$），
再将 $\tilde\zeta_2$ 越过各元素
$\tilde\gamma_c, \ldots, \tilde\gamma_1$
（这是允许的，因为 $\mathfrak m\tilde\zeta_2 \subset A + R\tilde\zeta_1$），
依此类推。

\medskip\noindent
引理的第 (3) 部分显然成立。证明完毕。
\end{proof}

\noindent
可以利用引理中的映射 $\gamma$ 如下定义更一般的映射 $\gamma$。
设 $(R, \mathfrak m, \kappa)$ 为局部环，$M$ 为有限长度 $R$-模，并设给定
一个有限滤过（见《Homology》定义 \ref{homology-definition-filtered}）
$$
0 = F^m \subset F^{m - 1} \subset \ldots \subset F^{n + 1} \subset F^n = M
$$
则有一个良定义的典范同构
$$
\gamma_{(M, F)} :
\det\nolimits_\kappa(F^{m - 1}/F^m) \otimes_\kappa \ldots \otimes_k 
\det\nolimits_\kappa(F^n/F^{n + 1})
\longrightarrow
\det\nolimits_\kappa(M)
$$
为构造它，我们使用由短正合列
$0 \to F^{i - 1}/F^i \to M/F^i \to M/F^{i - 1} \to 0$
给出的引理 \ref{chow-lemma-det-exact-sequences} 中的同构。
引理 \ref{chow-lemma-det-exact-sequences} 的第 (2) 部分取 $G = 0$ 表明，若改用
短正合列 $0 \to F^i \to F^{i - 1} \to F^{i - 1}/F^i \to 0$，
得到的是同一个同构。

\medskip\noindent
下面是行列式函子的另一个典型结果。它不难证明；通常困难之处是证明
行列式函子的存在性。

\begin{lemma}
\label{chow-lemma-uniqueness-det}
设 $(R, \mathfrak m, \kappa)$ 为任意局部环。赋有映射
$\gamma_{K \to L \to M}$ 的函子
$$
\det\nolimits_\kappa :
\left\{
\begin{matrix}
\text{有限长度 }R\text{-模} \\
\text{（以同构为态射）}
\end{matrix}
\right\}
\longrightarrow
\left\{
\begin{matrix}
1\text{-维 }\kappa\text{-向量空间} \\
\text{（以同构为态射）}
\end{matrix}
\right\}
$$
由下列性质刻画：
\begin{enumerate}
\item 它限制到被 $\mathfrak m$ 零化的模所成子范畴后，与通常的行列式
函子同构（见引理 \ref{chow-lemma-compare-det}）；
\item 引理 \ref{chow-lemma-det-exact-sequences} 的 (1)、(2)、(3) 均成立。
\end{enumerate}
\end{lemma}

\begin{proof}
从略。
\end{proof}

\begin{lemma}
\label{chow-lemma-determinant-quotient-ring}
设 $(R', \mathfrak m') \to (R, \mathfrak m)$ 为局部环同态，并且它在
剩余域 $\kappa$ 上诱导同构。则对每个有限长度 $R$-模，限制
$M_{R'}$ 是有限长度 $R'$-模，而且存在典范同构
$$
\det\nolimits_{R, \kappa}(M)
\longrightarrow
\det\nolimits_{R', \kappa}(M_{R'})
$$
该同构关于 $M$ 函子性成立，并与分别为 $\det_{R, \kappa}$ 和
$\det_{R', \kappa}$ 定义的引理 \ref{chow-lemma-det-exact-sequences} 中的同构
$\gamma_{K \to L \to M}$ 相容。
\end{lemma}

\begin{proof}
若 $M$ 作为 $R$-模的长度为 $l$，则 $M$ 作为 $R'$-模（即 $M_{R'}$）
的长度也为 $l$，见《Algebra》引理 \ref{algebra-lemma-pushdown-module}。
注意，$M$ 在 $R$ 上的容许序列 $x_1, \ldots, x_l$ 也是 $M$ 在
$R'$ 上的容许序列，因为 $\mathfrak m'$ 映入 $\mathfrak m$。
所需同构把符号
$[x_1, \ldots, x_l] \in \det\nolimits_{R, \kappa}(M)$
映到相应的符号
$[x_1, \ldots, x_l] \in \det\nolimits_{R', \kappa}(M)$。
可以立即验证，它关于同构函子性成立，并与引理
\ref{chow-lemma-det-exact-sequences} 中的同构 $\gamma$ 相容。
\end{proof}

\begin{remark}
\label{chow-remark-explain-determinant}
设 $(R, \mathfrak m, \kappa)$ 为局部环，并假设 $\kappa$ 的特征为零，
或者其特征为 $p$ 且 $p R = 0$。设 $M_1, \ldots, M_n$ 是有限长度
$R$-模。下面将证明，存在一个理想 $I \subset \mathfrak m$，对
$i = 1, \ldots, n$ 零化 $M_i$，并且存在典范满射
$R/I \to \kappa$ 的一个截面 $\sigma : \kappa \to R/I$。
通过 $\sigma$ 限制得到的 $M_i$ 的模 $M_{i, \kappa}$ 是维数为
$l_i = \text{length}_R(M_i)$ 的 $\kappa$-向量空间；由引理
\ref{chow-lemma-determinant-quotient-ring} 可见
$$
\det\nolimits_\kappa(M_i) = \wedge_\kappa^{l_i}(M_{i, \kappa})
$$
对于被 $I$ 零化的有限长度 $R$-模所成的短正合列，这些同构与引理
\ref{chow-lemma-det-exact-sequences} 中的同构
$\gamma_{K \to M \to L}$ 相容。由此可知，验证 $\det_\kappa$ 的一个性质，
往往可以归结为在有限维向量空间范畴上验证通常行列式的相应性质。

\medskip\noindent
作为 $I$，可以取模 $M = \bigoplus M_i$ 的零化子
（《Algebra》定义 \ref{algebra-definition-annihilator}）。此时
$R/I \subset \text{End}_R(M)$，因而它具有有限长度。因此 $R/I$ 是
剩余域为 $\kappa$ 的阿廷局部环。由于阿廷局部环完备，由 Cohen 结构定理
（《Algebra》定理 \ref{algebra-theorem-cohen-structure-theorem}）可知
$R/I$ 有系数环；根据我们对 $R$ 的假设，该系数环是域。
\end{remark}

\noindent
下面给出一种能够计算线性映射行列式的情形。事实上，任何情形都没有
神秘之处；一个任取的例子见例 \ref{chow-example-determinant-map}。

\begin{lemma}
\label{chow-lemma-times-u-determinant}
设 $R$ 为剩余域为 $\kappa$ 的局部环，$u \in R^*$ 为单位，
$M$ 为 $R$ 上有限长度模。把乘以 $u$ 的映射记为 $u_M : M \to M$。
则
$$
\det\nolimits_\kappa(u_M) :
\det\nolimits_\kappa(M)
\longrightarrow
\det\nolimits_\kappa(M)
$$
等于乘以 $\overline{u}^l$ 的映射，其中 $l = \text{length}_R(M)$，而
$\overline{u} \in \kappa^*$ 是 $u$ 的像。
\end{lemma}

\begin{proof}
把满足
$\det\nolimits_\kappa(u_M) = f_M \text{id}_{\det\nolimits_\kappa(M)}$
的元素记为 $f_M \in \kappa^*$。设
$0 \to K \to L \to M \to 0$ 是有限 $R$-模的短正合列。
于是 $u_k$、$u_L$、$u_M$ 给出短正合列的一个同构。因此由引理
\ref{chow-lemma-det-exact-sequences} (1) 得到 $f_K f_M = f_L$。
这表明，对长度归纳后，只需证明长度为 $1$ 的情形，而该情形显然成立。
\end{proof}

\begin{example}
\label{chow-example-determinant-map}
考虑局部环 $R = \mathbf{Z}_p$。令
$M = \mathbf{Z}_p/(p^2) \oplus \mathbf{Z}_p/(p^3)$，并设
$u : M \to M$ 是由矩阵
$$
u =
\left(
\begin{matrix}
a & b \\
pc & d
\end{matrix}
\right)
$$
给出的映射，其中 $a, b, c, d \in \mathbf{Z}_p$，且
$a, d \in \mathbf{Z}_p^*$。在这种情形，$\det_\kappa(u)$ 等于乘以
$a^2d^3 \bmod p \in \mathbf{F}_p^*$。考察 $u$ 在符号
$[p^2e, pe, pf, e, f]$ 上的作用即可容易看出这一点，其中
$e = (0 , 1) \in M$ 且 $f = (1, 0) \in M$。
\end{example}





\subsection{周期复形与行列式}
\label{chow-subsection-periodic-complexes-determinants}

\noindent
设 $R$ 为剩余域为 $\kappa$ 的局部环，$(M, \varphi, \psi)$ 为 $R$ 上的
$(2, 1)$-周期复形。假设 $M$ 具有有限长度，且
$(M, \varphi, \psi)$ 正合。我们将利用行列式构造定义这种情形的一个不变量；
见小节 \ref{chow-subsection-determinants-finite-length}。简记
$K_\varphi = \Ker(\varphi)$、
$I_\varphi = \Im(\varphi)$、
$K_\psi = \Ker(\psi)$ 以及
$I_\psi = \Im(\psi)$。短正合列
$$
0 \to K_\varphi \to M \to I_\varphi \to 0, \quad
0 \to K_\psi \to M \to I_\psi \to 0
$$
给出同构
$$
\gamma_\varphi :
\det\nolimits_\kappa(K_\varphi)
\otimes
\det\nolimits_\kappa(I_\varphi)
\longrightarrow
\det\nolimits_\kappa(M), \quad
\gamma_\psi :
\det\nolimits_\kappa(K_\psi)
\otimes
\det\nolimits_\kappa(I_\psi)
\longrightarrow
\det\nolimits_\kappa(M),
$$
见引理 \ref{chow-lemma-det-exact-sequences}。另一方面，复形的正合性给出等式
$K_\varphi = I_\psi$ 和 $K_\psi = I_\varphi$，从而通过交换因子得到同构
$$
\sigma :
\det\nolimits_\kappa(K_\varphi)
\otimes
\det\nolimits_\kappa(I_\varphi)
\longrightarrow
\det\nolimits_\kappa(K_\psi)
\otimes
\det\nolimits_\kappa(I_\psi)
$$
利用这些记号，可以定义所需不变量。

\begin{definition}
\label{chow-definition-periodic-determinant}
设 $R$ 为剩余域为 $\kappa$ 的局部环，$(M, \varphi, \psi)$ 为 $R$ 上的
$(2, 1)$-周期复形。假设 $M$ 具有有限长度，且
$(M, \varphi, \psi)$ 正合。$(M, \varphi, \psi)$ 的{\it 行列式}是元素
$$
\det\nolimits_\kappa(M, \varphi, \psi) \in \kappa^*
$$
它满足：复合
$$
\det\nolimits_\kappa(M)
\xrightarrow{\gamma_\psi \circ \sigma \circ \gamma_\varphi^{-1}}
\det\nolimits_\kappa(M)
$$
等于乘以
$(-1)^{\text{length}_R(I_\varphi)\text{length}_R(I_\psi)}
\det\nolimits_\kappa(M, \varphi, \psi)$ 的映射。
\end{definition}

\begin{remark}
\label{chow-remark-more-elementary}
下面给出上述行列式更具体的描述。设 $R$ 为剩余域为 $\kappa$ 的局部环，
$(M, \varphi, \psi)$ 为 $R$ 上的 $(2, 1)$-周期复形。假设 $M$ 具有
有限长度，且 $(M, \varphi, \psi)$ 正合。仍简记
$I_\varphi = \Im(\varphi)$、$I_\psi = \Im(\psi)$。
假设 $\text{length}_R(I_\varphi) = a$ 且
$\text{length}_R(I_\psi) = b$；由正合性，
$a + b = \text{length}_R(M)$。选取容许序列
$x_1, \ldots, x_a \in I_\varphi$ 和 $y_1, \ldots, y_b \in I_\psi$，
使符号 $[x_1, \ldots, x_a]$ 生成 $\det_\kappa(I_\varphi)$，而符号
$[x_1, \ldots, x_b]$ 生成 $\det_\kappa(I_\psi)$。选取
$\tilde x_i \in M$ 使得 $\varphi(\tilde x_i) = x_i$，并选取
$\tilde y_j \in M$ 使得 $\psi(\tilde y_j) = y_j$。于是
$\det_\kappa(M, \varphi, \psi)$ 由等式
$$
[x_1, \ldots, x_a, \tilde y_1, \ldots, \tilde y_b]
=
(-1)^{ab} \det\nolimits_\kappa(M, \varphi, \psi)
[y_1, \ldots, y_b, \tilde x_1, \ldots, \tilde x_a]
$$
在 $\det_\kappa(M)$ 中成立这一条件刻画。这也解释了符号因子。
\end{remark}

\begin{lemma}
\label{chow-lemma-periodic-determinant-shift}
设 $R$ 为剩余域为 $\kappa$ 的局部环，$(M, \varphi, \psi)$ 为 $R$ 上的
$(2, 1)$-周期复形。假设 $M$ 具有有限长度，且
$(M, \varphi, \psi)$ 正合。则
$$
\det\nolimits_\kappa(M, \varphi, \psi)
\det\nolimits_\kappa(M, \psi, \varphi)
= 1.
$$
\end{lemma}

\begin{proof}
从略。
\end{proof}

\begin{lemma}
\label{chow-lemma-periodic-determinant-sign}
设 $R$ 为剩余域为 $\kappa$ 的局部环，$(M, \varphi, \varphi)$ 为 $R$ 上的
$(2, 1)$-周期复形。假设 $M$ 具有有限长度，且
$(M, \varphi, \varphi)$ 正合。则
$\text{length}_R(M) = 2 \text{length}_R(\Im(\varphi))$，并且
$$
\det\nolimits_\kappa(M, \varphi, \varphi)
=
(-1)^{\text{length}_R(\Im(\varphi))}
=
(-1)^{\frac{1}{2}\text{length}_R(M)}
$$
\end{lemma}

\begin{proof}
这直接由定义中的符号规则得出。
\end{proof}

\begin{lemma}
\label{chow-lemma-periodic-determinant-easy-case}
设 $R$ 为剩余域为 $\kappa$ 的局部环，$M$ 为有限长度 $R$-模。
\begin{enumerate}
\item 若 $\varphi : M \to M$ 是同构，则
$\det_\kappa(M, \varphi, 0) = \det_\kappa(\varphi)$；
\item 若 $\psi : M \to M$ 是同构，则
$\det_\kappa(M, 0, \psi) = \det_\kappa(\psi)^{-1}$。
\end{enumerate}
\end{lemma}

\begin{proof}
证明 (1)。令 $\psi = 0$。使用定义 \ref{chow-definition-periodic-determinant}
上方的记号，可作认同
$K_\varphi = I_\psi = 0$、$I_\varphi = K_\psi = M$。
在这些认同下，映射
$$
\gamma_\varphi :
\kappa \otimes \det\nolimits_\kappa(M)
=
\det\nolimits_\kappa(K_\varphi)
\otimes
\det\nolimits_\kappa(I_\varphi)
\longrightarrow
\det\nolimits_\kappa(M)
$$
被认同为 $\det_\kappa(\varphi^{-1})$。另一方面，映射
$\gamma_\psi$ 被认同为恒等映射。因此在这种情形，
$\gamma_\psi \circ \sigma \circ \gamma_\varphi^{-1}$ 等于
$\det_\kappa(\varphi)$，结论随即得到。省略 (2) 的证明。
\end{proof}

\begin{lemma}
\label{chow-lemma-periodic-determinant}
设 $R$ 为剩余域为 $\kappa$ 的局部环。假设有 $(2, 1)$-周期复形的短正合列
$$
0 \to (M_1, \varphi_1, \psi_1)
\to (M_2, \varphi_2, \psi_2)
\to (M_3, \varphi_3, \psi_3)
\to 0
$$
其中所有 $M_i$ 都具有有限长度，且每个 $(M_1, \varphi_1, \psi_1)$
都正合。则
$$
\det\nolimits_\kappa(M_2, \varphi_2, \psi_2) =
\det\nolimits_\kappa(M_1, \varphi_1, \psi_1)
\det\nolimits_\kappa(M_3, \varphi_3, \psi_3).
$$
在 $\kappa^*$ 中成立。
\end{lemma}

\begin{proof}
简记
$I_{\varphi, i} = \Im(\varphi_i)$、
$K_{\varphi, i} = \Ker(\varphi_i)$、
$I_{\psi, i} = \Im(\psi_i)$ 以及
$K_{\psi, i} = \Ker(\psi_i)$。注意，有一个有限长度 $R$-模的行列均正合的
交换方图
$$
\xymatrix{
& 0 \ar[d] & 0 \ar[d] & 0 \ar[d] & \\
0 \ar[r] &
K_{\varphi, 1} \ar[r] \ar[d] &
K_{\varphi, 2} \ar[r] \ar[d] &
K_{\varphi, 3} \ar[r] \ar[d] &
0 \\
0 \ar[r] &
M_1 \ar[r] \ar[d] &
M_2 \ar[r] \ar[d] &
M_3 \ar[r] \ar[d] &
0 \\
0 \ar[r] &
I_{\varphi, 1} \ar[r] \ar[d] &
I_{\varphi, 2} \ar[r] \ar[d] &
I_{\varphi, 3} \ar[r] \ar[d] &
0 \\
& 0  & 0  & 0  &
}
$$
顶行正合，因为它可以与像所成的序列
$I_{\psi, 1} \to I_{\psi, 2} \to I_{\psi, 3} \to 0$ 认同；
底行同理。还存在一个涉及模 $I_{\psi, i}$ 与 $K_{\psi, i}$ 的类似图。
按照定义，$\det_\kappa(M_2, \varphi_2, \psi_2)$ 在不计符号时对应于
下图左侧竖直映射的复合：
$$
\xymatrix{
\det_\kappa(M_1) \otimes
\det_\kappa(M_3) \ar[r]^\gamma
\ar[d]^{\gamma^{-1} \otimes \gamma^{-1}} &
\det_\kappa(M_2)
\ar[d]^{\gamma^{-1}} \\
\det\nolimits_\kappa(K_{\varphi, 1})
\otimes
\det\nolimits_\kappa(I_{\varphi, 1})
\otimes
\det\nolimits_\kappa(K_{\varphi, 3})
\otimes
\det\nolimits_\kappa(I_{\varphi, 3})
\ar[d]^{\sigma \otimes \sigma}
\ar[r]^-{\gamma \otimes \gamma} &
\det\nolimits_\kappa(K_{\varphi, 2})
\otimes
\det\nolimits_\kappa(I_{\varphi, 2})
\ar[d]^\sigma
\\
\det\nolimits_\kappa(K_{\psi, 1})
\otimes
\det\nolimits_\kappa(I_{\psi, 1})
\otimes
\det\nolimits_\kappa(K_{\psi, 3})
\otimes
\det\nolimits_\kappa(I_{\psi, 3})
\ar[d]^{\gamma \otimes \gamma}
\ar[r]^-{\gamma \otimes \gamma}
&
\det\nolimits_\kappa(K_{\psi, 2})
\otimes
\det\nolimits_\kappa(I_{\psi, 2})
\ar[d]^\gamma \\
\det_\kappa(M_1)
\otimes
\det_\kappa(M_3) \ar[r]^\gamma
&
\det_\kappa(M_2)
}
$$
应用引理 \ref{chow-lemma-det-exact-sequences} (2)，可知顶部和底部方块在一个
符号因子意义下交换。中间方块显然交换（只是交换因子）。因此可见
$\det\nolimits_\kappa(M_2, \varphi_2, \psi_2) =
\epsilon \det\nolimits_\kappa(M_1, \varphi_1, \psi_1)
\det\nolimits_\kappa(M_3, \varphi_3, \psi_3)
$
对某个符号 $\epsilon$ 成立。该符号可以算出；具体而言，上图外侧矩形交换时
相差
\begin{eqnarray*}
\epsilon & = &
(-1)^{\text{length}(I_{\varphi, 1})\text{length}(K_{\varphi, 3})
+ \text{length}(I_{\psi, 1})\text{length}(K_{\psi, 3})} \\
& = &
(-1)^{\text{length}(I_{\varphi, 1})\text{length}(I_{\psi, 3})
+ \text{length}(I_{\psi, 1})\text{length}(I_{\varphi, 3})}
\end{eqnarray*}
（证明从略）。由此容易看出符号也完全吻合。
\end{proof}

\begin{example}
\label{chow-example-dual-numbers}
设 $k$ 为域。考虑 $k$ 上的对偶数环 $R = k[T]/(T^2)$，并把 $T$ 在
$R$ 中的类记为 $t$。令 $M = R$，并取 $\varphi = ut$、$\psi = vt$，
其中 $u, v \in k^*$。在这种情形，$\det_k(M)$ 有生成元
$e = [t, 1]$。我们作认同
$I_\varphi = K_\varphi = I_\psi = K_\psi = (t)$。于是
$\gamma_\varphi(t \otimes t) = u^{-1}[t, 1]$
（因为 $u^{-1} \in M$ 是 $t \in I_\varphi$ 的一个提升），并且
$\gamma_\psi(t \otimes t) = v^{-1}[t, 1]$（理由相同）。因此
$\det_k(M, \varphi, \psi) = -u/v \in k^*$。
\end{example}

\begin{example}
\label{chow-example-Zp}
设 $R = \mathbf{Z}_p$ 且 $M = \mathbf{Z}_p/(p^l)$。令
$\varphi = p^b u$ 且 $\varphi = p^a v$，其中 $a, b \geq 0$、
$a + b = l$，而 $u, v \in \mathbf{Z}_p^*$。与例
\ref{chow-example-dual-numbers} 中相同的计算表明
\begin{eqnarray*}
\det\nolimits_{\mathbf{F}_p}(\mathbf{Z}_p/(p^l), p^bu, p^av) & = &
(-1)^{ab}u^a/v^b \bmod p \\
& = &
(-1)^{\text{ord}_p(\alpha)\text{ord}_p(\beta)}
\frac{\alpha^{\text{ord}_p(\beta)}}{\beta^{\text{ord}_p(\alpha)}} \bmod p
\end{eqnarray*}
其中 $\alpha = p^bu, \beta = p^av \in \mathbf{Z}_p$。
更一般的情形（及其证明）见引理 \ref{chow-lemma-symbol-is-usual-tame-symbol}。
\end{example}

\begin{example}
\label{chow-example-generic-vector-space}
设 $R = k$ 为域，设 $M = k^{\oplus a} \oplus k^{\oplus b}$ 的维数为
$l = a + b$。令 $\varphi$ 和 $\psi$ 为下列对角矩阵
$$
\varphi = \text{diag}(u_1, \ldots, u_a, 0, \ldots, 0),
\quad
\psi = \text{diag}(0, \ldots, 0, v_1, \ldots, v_b)
$$
其中 $u_i, v_j \in k^*$。此时
$$
\det\nolimits_k(M, \varphi, \psi)
=
\frac{u_1 \ldots u_a}{v_1 \ldots v_b}.
$$
这可以直接计算得到，也可以先计算 $l = 1$ 的情形，再使用引理
\ref{chow-lemma-periodic-determinant} 的可加性。
\end{example}

\begin{example}
\label{chow-example-special-vector-space}
设 $R = k$ 为域，设 $M = k^{\oplus a} \oplus k^{\oplus a}$ 的维数为
$l = 2a$。令 $\varphi$ 和 $\psi$ 为下列分块矩阵
$$
\varphi =
\left(
\begin{matrix}
0 & U \\
0 & 0
\end{matrix}
\right),
\quad
\psi =
\left(
\begin{matrix}
0 & V \\
0 & 0
\end{matrix}
\right),
$$
其中 $U, V \in \text{Mat}(a \times a, k)$ 可逆。此时
$$
\det\nolimits_k(M, \varphi, \psi)
=
(-1)^a\frac{\det(U)}{\det(V)}.
$$
这可以直接计算得到。当 $a = 1$ 时，该计算与例
\ref{chow-example-dual-numbers} 中的计算类似。
\end{example}

\begin{example}
\label{chow-example-a-la-oort}
设 $R = k$ 为域，$M = k^{\oplus 4}$。令
$$
\varphi =
\left(
\begin{matrix}
  0 &   0 &   0 &   0 \\
u_1 &   0 &   0 &   0 \\
  0 &   0 &   0 &   0 \\
  0 &   0 & u_2 &   0
\end{matrix}
\right)
\quad
\varphi =
\left(
\begin{matrix}
  0 &   0 &   0 &   0 \\
  0 &   0 & v_2 &   0 \\
  0 &   0 &   0 &   0 \\
v_1 &   0 &   0 &   0
\end{matrix}
\right)
\quad
$$
其中 $u_1, u_2, v_1, v_2 \in k^*$。则
$$
\det\nolimits_k(M, \varphi, \psi) = -\frac{u_1u_2}{v_1v_2}.
$$
\end{example}

\noindent
下面讨论线性自同态复合的行列式等于各行列式之积这一事实的对应结论。
为避免公式过长，对任意 $R$-模映射 $\varphi : M \to M$，记
$I_\varphi = \Im(\varphi)$，并记
$K_\varphi = \Ker(\varphi)$。对一对态射
$\varphi, \psi : M \to M$，还记
$\varphi\psi = \varphi \circ \psi$。

\begin{lemma}
\label{chow-lemma-multiplicativity-determinant}
设 $R$ 为剩余域为 $\kappa$ 的局部环，$M$ 为有限长度 $R$-模，
$\alpha, \beta, \gamma$ 为 $M$ 的自同态。假设
\begin{enumerate}
\item $I_\alpha = K_{\beta\gamma}$，而且对 $\alpha, \beta, \gamma$
的任意置换也有类似等式；
\item $K_\alpha = I_{\beta\gamma}$，而且对 $\alpha, \beta, \gamma$
的任意置换也有类似等式。
\end{enumerate}
则
\begin{enumerate}
\item 三元组 $(M, \alpha, \beta\gamma)$ 是正合 $(2, 1)$-周期复形；
\item 三元组 $(I_\gamma, \alpha, \beta)$ 是正合 $(2, 1)$-周期复形；
\item 三元组 $(M/K_\beta, \alpha, \gamma)$ 是正合 $(2, 1)$-周期复形；
\item 有
$$
\det\nolimits_\kappa(M, \alpha, \beta\gamma)
=
\det\nolimits_\kappa(I_\gamma, \alpha, \beta)
\det\nolimits_\kappa(M/K_\beta, \alpha, \gamma).
$$
\end{enumerate}
\end{lemma}

\begin{proof}
假设显然蕴含引理的第 (1) 部分。

\medskip\noindent
为看出第 (1) 部分，注意假设蕴含
$I_{\gamma\alpha} = I_{\alpha\gamma}$；对核及任意其他一对态射也类似。
此外，$I_{\gamma\beta} =I_{\beta\gamma} = K_\alpha \subset I_\gamma$；
对任意其他一对也类似。特别地，得到短正合列
$$
0 \to I_{\beta\gamma} \to I_\gamma \xrightarrow{\alpha} I_{\alpha\gamma} \to 0
$$
以及类似的短正合列
$$
0 \to I_{\alpha\gamma} \to I_\gamma \xrightarrow{\beta} I_{\beta\gamma} \to 0.
$$
这证明 $(I_\gamma, \alpha, \beta)$ 是正合 $(2, 1)$-周期复形，故引理的
第 (2) 部分成立。

\medskip\noindent
为看出 $\alpha$、$\gamma$ 在 $M/K_\beta$ 上给出良定义的自同态，必须验证
$\alpha(K_\beta) \subset K_\beta$ 和 $\gamma(K_\beta) \subset K_\beta$。
这确实成立，因为
$\alpha(K_\beta) = \alpha(I_{\gamma\alpha}) = I_{\alpha\gamma\alpha}
\subset I_{\alpha\gamma} = K_\beta$；另一情形类似。映射
$\alpha : M/K_\beta \to M/K_\beta$ 的核为
$K_{\beta\alpha}/K_\beta = I_\gamma/K_\beta$。类似地，映射
$\gamma : M/K_\beta \to M/K_\beta$ 的核等于 $I_\alpha/K_\beta$。
因此 (3) 成立。

\medskip\noindent
记 $r = \text{length}_R(K_\alpha)$、
$s = \text{length}_R(K_\beta)$、$t = \text{length}_R(K_\gamma)$。
由上述正合列及假设，有
$\text{length}_R(I_\alpha) = s + t$、$\text{length}_R(I_\beta) = r + t$、
$\text{length}_R(I_\gamma) = r + s$，以及
$\text{length}(M) = r + s + t$。选取
\begin{enumerate}
\item 生成 $K_\alpha$ 的容许序列 $x_1, \ldots, x_r \in K_\alpha$；
\item 生成 $K_\beta$ 的容许序列 $y_1, \ldots, y_s \in K_\beta$；
\item 生成 $K_\gamma$ 的容许序列 $z_1, \ldots, z_t \in K_\gamma$；
\item 元素 $\tilde x_i \in M$，使得 $\beta\gamma\tilde x_i = x_i$；
\item 元素 $\tilde y_i \in M$，使得 $\alpha\gamma\tilde y_i = y_i$；
\item 元素 $\tilde z_i \in M$，使得 $\beta\alpha\tilde z_i = z_i$。
\end{enumerate}
在这些选取下，序列
$y_1, \ldots, y_s, \alpha\tilde z_1, \ldots, \alpha\tilde z_t$
是 $I_\alpha$ 中生成它的容许序列。因此由注
\ref{chow-remark-more-elementary}，行列式
$D = \det_\kappa(M, \alpha, \beta\gamma)$ 是满足下式的唯一元素
$\kappa^*$：
\begin{align*}
[y_1, \ldots, y_s,
\alpha\tilde z_1, \ldots, \alpha\tilde z_s,
\tilde x_1, \ldots, \tilde x_r] \\
= (-1)^{r(s + t)} D
[x_1, \ldots, x_r,
\gamma\tilde y_1, \ldots, \gamma\tilde y_s,
\tilde z_1, \ldots, \tilde z_t]
\end{align*}
由同一注，$D_1 = \det_\kappa(M/K_\beta, \alpha, \gamma)$ 由下式刻画：
$$
[y_1, \ldots, y_s,
\alpha\tilde z_1, \ldots, \alpha\tilde z_t,
\tilde x_1, \ldots, \tilde x_r]
=
(-1)^{rt} D_1
[y_1, \ldots, y_s,
\gamma\tilde x_1, \ldots, \gamma\tilde x_r,
\tilde z_1, \ldots, \tilde z_t]
$$
同样，$D_2 = \det_\kappa(I_\gamma, \alpha, \beta)$ 由下式刻画：
$$
[y_1, \ldots, y_s,
\gamma\tilde x_1, \ldots, \gamma\tilde x_r,
\tilde z_1, \ldots, \tilde z_t]
=
(-1)^{rs} D_2
[x_1, \ldots, x_r,
\gamma\tilde y_1, \ldots, \gamma\tilde y_s,
\tilde z_1, \ldots, \tilde z_t]
$$
合并上述公式即得 $D = D_1 D_2$，正如所需。
\end{proof}


\begin{lemma}
\label{chow-lemma-tricky}
设 $R$ 为剩余域为 $\kappa$ 的局部环。设
$\alpha : (M, \varphi, \psi) \to (M', \varphi', \psi')$
为 $R$ 上 $(2, 1)$-周期复形的态射。假设
\begin{enumerate}
\item $M$、$M'$ 具有有限长度；
\item $(M, \varphi, \psi)$、$(M', \varphi', \psi')$ 正合；
\item 映射 $\varphi$、$\psi$ 在 $K = \Ker(\alpha)$ 上诱导零映射；
\item 映射 $\varphi$、$\psi$ 在 $Q = \Coker(\alpha)$ 上诱导零映射。
\end{enumerate}
记 $N = \alpha(M) \subset M'$。得到 $(2, 1)$-周期复形的两个短正合列
$$
\begin{matrix}
0 \to (N, \varphi', \psi') \to (M', \varphi', \psi') \to (Q, 0, 0) \to 0 \\
0 \to (K, 0, 0) \to (M, \varphi, \psi) \to (N, \varphi', \psi') \to 0
\end{matrix}
$$
它们诱导两个同构 $\alpha_i : Q \to K$，$i = 0, 1$。则
$$
\det\nolimits_\kappa(M, \varphi, \psi)
=
\det\nolimits_\kappa(\alpha_0^{-1} \circ \alpha_1)
\det\nolimits_\kappa(M', \varphi', \psi')
$$
特别地，若 $\alpha_0 = \alpha_1$，则
$\det\nolimits_\kappa(M, \varphi, \psi) =
\det\nolimits_\kappa(M', \varphi', \psi')$。
\end{lemma}

\begin{proof}
证明此引理（至少）有两种方法。一种是利用行列式的性质构造一个庞大的交换图；
另一种是使用以元素的容许序列表述的行列式刻画。我们采用第二种方法。

\medskip\noindent
先精确说明映射 $\alpha_i$。具体而言，$\alpha_0$ 是复合
$$
\alpha_0 : Q = H^0(Q, 0, 0) \to H^1(N, \varphi', \psi') \to H^2(K, 0, 0) = K
$$
而 $\alpha_1$ 是复合
$$
\alpha_1 : Q = H^1(Q, 0, 0) \to H^2(N, \varphi', \psi') \to H^3(K, 0, 0) = K
$$
它们来自引理中所示复形短正合列的边界映射。复形
$(M, \varphi, \psi)$、$(M', \varphi', \psi')$ 的正合性蕴含这些映射
都是同构。

\medskip\noindent
我们采用记号 $I_\varphi = \Im(\varphi)$、
$K_\varphi = \Ker(\varphi)$，其他映射的记号类似。$M$ 与 $M'$ 的正合性
意味着 $K_\varphi = I_\psi$ 及另外三个类似等式。记
$k = \text{length}_R(K)$、$a = \text{length}_R(I_\varphi)$、
$b = \text{length}_R(I_\psi)$。于是
$\text{length}_R(M) = a + b$、
$\text{length}_R(N) = a + b - k$、$\text{length}_R(Q) = k$，并且
$\text{length}_R(M') = a + b$。下面的正合列还将表明
$\text{length}_R(I_{\varphi'}) = a$ 与
$\text{length}_R(I_{\psi'}) = b$。

\medskip\noindent
假设 $K \subset K_\varphi = I_\psi$ 意味着 $\varphi$ 经 $N$ 分解，给出
正合列
$$
0 \to \alpha(I_\psi) \to N \xrightarrow{\varphi\alpha^{-1}} I_\psi \to 0.
$$
这里 $\varphi\alpha^{-1}(x') = y$ 意指 $x' = \alpha(x)$ 且
$y = \varphi(x)$。类似地，有
$$
0 \to \alpha(I_\varphi) \to N \xrightarrow{\psi\alpha^{-1}} I_\varphi \to 0.
$$
$\psi'$ 在 $Q$ 上诱导零映射这一假设意味着
$I_{\psi'} = K_{\varphi'} \subset N$。这意味着商
$\varphi'(N) \subset I_{\varphi'}$ 与 $Q$ 认同。注意
$\varphi'(N) = \alpha(I_\varphi)$。因此得到同构
$$
\varphi' : Q \to I_{\varphi'}/\alpha(I_\varphi)
$$
它可简单描述为
$\varphi'(x' \bmod N) = \varphi'(x') \bmod \alpha(I_\varphi)$。
完全类似地得到
$$
\psi' : Q \to I_{\psi'}/\alpha(I_\psi)
$$
最后，注意 $\alpha_0$ 是复合
$$
\xymatrix{
Q \ar[r]^-{\varphi'} &
I_{\varphi'}/\alpha(I_\varphi)
\ar[rrr]^-{\psi\alpha^{-1}|_{I_{\varphi'}/\alpha(I_\varphi)}} & & &
K
}
$$
而类似地，
$\alpha_1 = \varphi\alpha^{-1}|_{I_{\psi'}/\alpha(I_\psi)} \circ \psi'$。

\medskip\noindent
为缩短下面的公式，以下将写 $\alpha x$ 而不写 $\alpha(x)$。这不会引起混淆，
因为所有映射都用希腊字母表示，而元素用罗马字母表示。选取
\begin{enumerate}
\item 生成 $K$ 的容许序列 $z_1, \ldots, z_k \in K$；
\item 元素 $z'_i \in M$，使得 $\varphi z'_i = z_i$；
\item 元素 $z''_i \in M$，使得 $\psi z''_i = z_i$；
\item 元素 $x_{k + 1}, \ldots, x_a \in I_\varphi$，使得
$z_1, \ldots, z_k, x_{k + 1}, \ldots, x_a$ 是生成 $I_\varphi$ 的
容许序列；
\item 元素 $\tilde x_i \in M$，使得 $\varphi \tilde x_i = x_i$；
\item 元素 $y_{k + 1}, \ldots, y_b \in I_\psi$，使得
$z_1, \ldots, z_k, y_{k + 1}, \ldots, y_b$ 是生成 $I_\psi$ 的
容许序列；
\item 元素 $\tilde y_i \in M$，使得 $\psi \tilde y_i = y_i$；
\item 元素 $w_1, \ldots, w_k \in M'$，使得
$w_1 \bmod N, \ldots, w_k \bmod N$ 是 $Q$ 中生成 $Q$ 的容许序列。
\end{enumerate}
由注 \ref{chow-remark-more-elementary}，元素
$D = \det_\kappa(M, \varphi, \psi) \in \kappa^*$ 由下式刻画：
\begin{eqnarray*}
& &
[z_1, \ldots, z_k,
x_{k + 1}, \ldots, x_a,
z''_1, \ldots, z''_k,
\tilde y_{k + 1}, \ldots, \tilde y_b] \\
& = &
(-1)^{ab} D
[z_1, \ldots, z_k,
y_{k + 1}, \ldots, y_b,
z'_1, \ldots, z'_k,
\tilde x_{k + 1}, \ldots, \tilde x_a]
\end{eqnarray*}
由上面的讨论，
$\alpha x_{k + 1}, \ldots, \alpha x_a, \varphi w_1, \ldots, \varphi w_k$
是生成 $I_{\varphi'}$ 的容许序列，而
$\alpha y_{k + 1}, \ldots, \alpha y_b, \psi w_1, \ldots, \psi w_k$
是生成 $I_{\psi'}$ 的容许序列。因此由注
\ref{chow-remark-more-elementary}，元素
$D' = \det_\kappa(M', \varphi', \psi') \in \kappa^*$ 由下式刻画：
\begin{eqnarray*}
& &
[\alpha x_{k + 1}, \ldots, \alpha x_a,
\varphi' w_1, \ldots, \varphi' w_k,
\alpha \tilde y_{k + 1}, \ldots, \alpha \tilde y_b,
w_1, \ldots, w_k]
\\
& = &
(-1)^{ab} D'
[\alpha y_{k + 1}, \ldots, \alpha y_b,
\psi' w_1, \ldots, \psi' w_k,
\alpha \tilde x_{k + 1}, \ldots, \alpha \tilde x_a,
w_1, \ldots, w_k]
\end{eqnarray*}
注意，在第一个，resp.\ 第二个显示公式中，两边符号的最前，resp.\ 最后
$k$ 个分量相同。因此这些公式实际上等价于下列等式：
\begin{eqnarray*}
& &
[\alpha x_{k + 1}, \ldots, \alpha x_a,
\alpha z''_1, \ldots, \alpha z''_k,
\alpha \tilde y_{k + 1}, \ldots, \alpha \tilde y_b] \\
& = &
(-1)^{ab} D
[\alpha y_{k + 1}, \ldots, \alpha y_b,
\alpha z'_1, \ldots, \alpha z'_k,
\alpha \tilde x_{k + 1}, \ldots, \alpha \tilde x_a]
\end{eqnarray*}
以及
\begin{eqnarray*}
& &
[\alpha x_{k + 1}, \ldots, \alpha x_a,
\varphi' w_1, \ldots, \varphi' w_k,
\alpha \tilde y_{k + 1}, \ldots, \alpha \tilde y_b]
\\
& = &
(-1)^{ab} D'
[\alpha y_{k + 1}, \ldots, \alpha y_b,
\psi' w_1, \ldots, \psi' w_k,
\alpha \tilde x_{k + 1}, \ldots, \alpha \tilde x_a]
\end{eqnarray*}
这些等式在 $\det_\kappa(N)$ 中成立。注意，
$\varphi' w_1, \ldots, \varphi' w_k$ 与
$\alpha z''_1, \ldots, z''_k$ 是生成模
$I_{\varphi'}/\alpha(I_\varphi)$ 的容许序列。写成
$$
[\varphi' w_1, \ldots, \varphi' w_k]
= \lambda_0 [\alpha z''_1, \ldots, \alpha z''_k]
$$
在 $\det_\kappa(I_{\varphi'}/\alpha(I_\varphi))$ 中成立，其中
$\lambda_0 \in \kappa^*$。类似地，写成
$$
[\psi' w_1, \ldots, \psi' w_k]
= \lambda_1 [\alpha z'_1, \ldots, \alpha z'_k]
$$
在 $\det_\kappa(I_{\psi'}/\alpha(I_\psi))$ 中成立，其中
$\lambda_1 \in \kappa^*$。一方面，由上面对 $\alpha_i$ 的描述，显然有
$$
\alpha_i([w_1, \ldots, w_k]) = \lambda_i[z_1, \ldots, z_k]
$$
对 $i = 0, 1$ 成立。这意味着
$$
\det\nolimits_\kappa(\alpha_0^{-1} \circ \alpha_1)
=
\lambda_1/\lambda_0
$$
另一方面，显然有
\begin{eqnarray*}
& &
\lambda_0 [\alpha x_{k + 1}, \ldots, \alpha x_a,
\alpha z''_1, \ldots, \alpha z''_k,
\alpha \tilde y_{k + 1}, \ldots, \alpha \tilde y_b] \\
& = &
[\alpha x_{k + 1}, \ldots, \alpha x_a,
\varphi' w_1, \ldots, \varphi' w_k,
\alpha \tilde y_{k + 1}, \ldots, \alpha \tilde y_b]
\end{eqnarray*}
以及
\begin{eqnarray*}
& &
\lambda_1[\alpha y_{k + 1}, \ldots, \alpha y_b,
\alpha z'_1, \ldots, \alpha z'_k,
\alpha \tilde x_{k + 1}, \ldots, \alpha \tilde x_a] \\
& = &
[\alpha y_{k + 1}, \ldots, \alpha y_b,
\psi' w_1, \ldots, \psi' w_k,
\alpha \tilde x_{k + 1}, \ldots, \alpha \tilde x_a]
\end{eqnarray*}
这些等式蕴含 $\lambda_0 D = \lambda_1 D'$，引理得证。
\end{proof}








\subsection{符号}
\label{chow-subsection-symbols}

\noindent
这一构造适当的一般性也许由下列引理中的情形给出。

\begin{lemma}
\label{chow-lemma-pre-symbol}
设 $A$ 为诺特局部环，$M$ 为维数为 $1$ 的有限 $A$-模。假设
$\varphi, \psi : M \to M$ 是两个单射 $A$-模映射，并且
$\varphi(\psi(M)) = \psi(\varphi(M))$；例如，$\varphi$ 与 $\psi$
交换时即如此。则 $\text{length}_A(M/\varphi\psi M) < \infty$，而且
$(M/\varphi\psi M, \varphi, \psi)$ 是正合 $(2, 1)$-周期复形。
\end{lemma}

\begin{proof}
设 $\mathfrak q$ 为 $M$ 的支集的一个极小素理想。则
$M_{\mathfrak q}$ 是有限长度 $A_{\mathfrak q}$-模，见《Algebra》引理
\ref{algebra-lemma-support-point}。因此 $\varphi$ 与 $\psi$ 都诱导同构
$M_{\mathfrak q} \to M_{\mathfrak q}$。所以 $M/\varphi\psi M$ 的支集为
$\{\mathfrak m_A\}$，从而它具有有限长度（见上述引理）。最后，
$\varphi$ 在 $M/\varphi\psi M$ 上的核显然为
$\psi M/\varphi\psi M$，所以 $\varphi$ 的核就是 $\psi$ 在
$M/\varphi\psi M$ 上的像。反向情形类似，因为根据假设，
$M/\varphi\psi M = M/\psi\varphi M$。
\end{proof}

\begin{lemma}
\label{chow-lemma-symbol-defined}
设 $A$ 为诺特局部环，$a, b \in A$。
\begin{enumerate}
\item 若 $M$ 是维数为 $1$ 的有限 $A$-模，且 $a, b$ 在 $M$ 上均为
非零因子，则 $\text{length}_A(M/abM) < \infty$，而且
$(M/abM, a, b)$ 是正合 $(2, 1)$-周期复形；
\item 若 $a, b$ 是非零因子且 $\dim(A) = 1$，则
$\text{length}_A(A/(ab)) < \infty$，而且
$(A/(ab), a, b)$ 是正合 $(2, 1)$-周期复形。
\end{enumerate}
特别地，在这些情形，
$\det_\kappa(M/abM, a, b) \in \kappa^*$，
resp.\ $\det_\kappa(A/(ab), a, b) \in \kappa^*$ 均有定义。
\end{lemma}

\begin{proof}
由引理 \ref{chow-lemma-pre-symbol} 得出。
\end{proof}

\begin{definition}
\label{chow-definition-symbol-M}
设 $A$ 为剩余域为 $\kappa$ 的诺特局部环，$a, b \in A$。
设 $M$ 为维数为 $1$ 的有限 $A$-模，且 $a, b$ 在 $M$ 上均为非零因子。
把{\it 与 $M, a, b$ 关联的符号}定义为元素
$$
d_M(a, b) =
\det\nolimits_\kappa(M/abM, a, b) \in \kappa^*
$$
\end{definition}

\begin{lemma}
\label{chow-lemma-multiplicativity-symbol}
设 $A$ 为诺特局部环，$a, b, c \in A$。设 $M$ 为满足
$\dim(\text{Supp}(M)) = 1$ 的有限 $A$-模，并假设 $a, b, c$ 在 $M$
上均为非零因子。则
$$
d_M(a, bc) = d_M(a, b) d_M(a, c)
$$
且 $d_M(a, b)d_M(b, a) = 1$。
\end{lemma}

\begin{proof}
把引理 \ref{chow-lemma-multiplicativity-determinant} 应用于 $M/abcM$ 以及
分别由乘以 $a, b, c$ 给出的自同态 $\alpha, \beta, \gamma$，即得第一个
断言。第二个断言来自引理 \ref{chow-lemma-periodic-determinant-shift}。
\end{proof}

\begin{definition}
\label{chow-definition-tame-symbol}
设 $A$ 为剩余域为 $\kappa$ 的 $1$ 维诺特局部整环，$K$ 为 $A$ 的分式域。
把 $A$ 的{\it 驯符号}定义为映射
$$
K^* \times K^* \longrightarrow \kappa^*,
\quad
(x, y) \longmapsto d_A(x, y)
$$
其中利用引理 \ref{chow-lemma-multiplicativity-symbol} 的乘法性，把
$d_A(x, y)$ 延拓到 $K^* \times K^*$。
\end{definition}

\noindent
显然，可以更一般地把 $d_M(-, -)$ 延拓到 $A$ 的某些分式环
（即使 $A$ 不是整环）。

\begin{lemma}
\label{chow-lemma-symbol-when-equal}
设 $A$ 为诺特局部环，$M$ 为维数为 $1$ 的有限 $A$-模。设
$a \in A$ 在 $M$ 上为非零因子。则
$d_M(a, a) = (-1)^{\text{length}_A(M/aM)}$。
\end{lemma}

\begin{proof}
立即由引理 \ref{chow-lemma-periodic-determinant-sign} 得出。
\end{proof}

\begin{lemma}
\label{chow-lemma-symbol-when-one-is-a-unit}
设 $A$ 为诺特局部环，$M$ 为维数为 $1$ 的有限 $A$-模。设
$b \in A$ 在 $M$ 上为非零因子，并设 $u \in A^*$。则
$$
d_M(u, b) = u^{\text{length}_A(M/bM)} \bmod \mathfrak m_A.
$$
特别地，若 $M = A$，则
$d_A(u, b) = u^{\text{ord}_A(b)} \bmod \mathfrak m_A$。
\end{lemma}

\begin{proof}
注意，在此情形 $M/ubM = M/bM$，而乘以 $b$ 在其上为零。因此由引理
\ref{chow-lemma-periodic-determinant-easy-case}，
$d_M(u, b) = \det_\kappa(u|_{M/bM})$。再由引理
\ref{chow-lemma-times-u-determinant} 得出结论。
\end{proof}

\begin{lemma}
\label{chow-lemma-symbol-short-exact-sequence}
设 $A$ 为诺特局部环，$a, b \in A$。设
$$
0 \to M \to M' \to M'' \to 0
$$
为维数为 $1$ 的 $A$-模的短正合列，并且 $a, b$ 在这三个 $A$-模上
均为非零因子。则
$$
d_{M'}(a, b) = d_M(a, b) d_{M''}(a, b)
$$
在 $\kappa^*$ 中成立。
\end{lemma}

\begin{proof}
容易看出，这给出正合 $(2, 1)$-周期复形的短正合列
$$
0 \to
(M/abM, a, b) \to
(M'/abM', a, b) \to
(M''/abM'', a, b) \to 0
$$
因此结论由引理 \ref{chow-lemma-periodic-determinant} 得出。
\end{proof}

\begin{lemma}
\label{chow-lemma-symbol-compare-modules}
设 $A$ 为诺特局部环，$\alpha : M \to M'$ 为维数为 $1$ 的有限
$A$-模之间的同态。设 $a, b \in A$，并假设
\begin{enumerate}
\item $a$、$b$ 在 $M$ 与 $M'$ 上均为非零因子；
\item $\dim(\Ker(\alpha)), \dim(\Coker(\alpha)) \leq 0$。
\end{enumerate}
则 $d_M(a, b) = d_{M'}(a, b)$。
\end{lemma}

\begin{proof}
若 $a \in A^*$，则所需等式由等式
$\text{length}(M/bM) = \text{length}(M'/bM')$ 与引理
\ref{chow-lemma-symbol-when-one-is-a-unit} 得出。类似地，若 $b$ 是单位，
引理也成立（由引理 \ref{chow-lemma-multiplicativity-symbol} 的对称性）。
因此可假设 $a, b \in \mathfrak m_A$。这特别蕴含 $\mathfrak m$ 不是
$M$ 的伴随素理想，因而 $\alpha : M \to M'$ 是单射。于是可把 $M$
视为 $M'$ 的子模。根据假设，$M'/M$ 是支集为 $\{\mathfrak m_A\}$ 的
有限 $A$-模，因此具有有限长度。注意，对任意满足
$M \subset M'' \subset M'$ 的第三个模 $M''$，映射 $M \to M''$ 与
$M'' \to M'$ 也满足引理的假设。因此对 $M'/M$ 的长度归纳，可归结到
$\text{length}_A(M'/M) = 1$ 的情形。最后，在这种情形考虑映射
$$
\overline{\alpha} : M/abM \longrightarrow M'/abM'.
$$
按构造，$\overline{\alpha}$ 的余核 $Q$ 的长度为 $1$。由于
$a, b \in \mathfrak m_A$，它们在 $Q$ 上作用为零。还可推出
$\overline{\alpha}$ 的核 $K$ 长度为 $1$，因而 $a$、$b$ 在 $K$ 上也
作用为零。因此可以应用引理 \ref{chow-lemma-tricky}。所以只需看出两个映射
$\alpha_i : Q \to K$ 相同。事实上，它们都等于映射
$q = x' \bmod \Im(\overline{\alpha}) \mapsto abx' \in K$。
验证从略。
\end{proof}

\begin{lemma}
\label{chow-lemma-compute-symbol-M}
设 $A$ 为诺特局部环，$M$ 为满足 $\dim(\text{Supp}(M)) = 1$ 的有限
$A$-模。设 $a, b \in A$ 在 $M$ 上为非零因子，并设
$\mathfrak q_1, \ldots, \mathfrak q_t$ 为 $M$ 的支集中的极小素理想。
则在 $\kappa^*$ 中有
$$
d_M(a, b)
=
\prod\nolimits_{i = 1, \ldots, t}
d_{A/\mathfrak q_i}(a, b)^{
\text{length}_{A_{\mathfrak q_i}}(M_{\mathfrak q_i})}
$$
\end{lemma}

\begin{proof}
选取由 $A$-子模组成的滤过
$$
0 = M_0 \subset M_1 \subset \ldots \subset M_n = M
$$
使得每个商 $M_j/M_{j - 1}$ 都与 $A/\mathfrak p_j$ 同构，其中
$\mathfrak p_j$ 是 $A$ 的某个素理想。见《Algebra》引理
\ref{algebra-lemma-filter-Noetherian-module}。对每个 $j$，或者对某个 $i$
有 $\mathfrak p_j = \mathfrak q_i$，或者
$\mathfrak p_j = \mathfrak m_A$。此外，对固定的 $i$，满足
$\mathfrak p_j = \mathfrak q_i$ 的 $j$ 的数目等于
$\text{length}_{A_{\mathfrak q_i}}(M_{\mathfrak q_i})$，见《Algebra》引理
\ref{algebra-lemma-filter-minimal-primes-in-support}。因此对每个 $j$，
$d_{M_j}(a, b)$ 均有定义，而且
$$
d_{M_j}(a, b)
=
\left\{
\begin{matrix}
d_{M_{j - 1}}(a, b) d_{A/\mathfrak q_i}(a, b) &
\text{若} & \mathfrak p_j = \mathfrak q_i \\
d_{M_{j - 1}}(a, b) & \text{若} & \mathfrak p_j = \mathfrak m_A
\end{matrix}
\right.
$$
第一种情形使用引理 \ref{chow-lemma-symbol-short-exact-sequence}，第二种情形
使用引理 \ref{chow-lemma-symbol-compare-modules}。引理得证。
\end{proof}

\begin{lemma}
\label{chow-lemma-symbol-is-usual-tame-symbol}
设 $A$ 为分式域为 $K$ 的离散赋值环。对非零 $x, y \in K$，有
$$
d_A(x, y)
=
(-1)^{\text{ord}_A(x)\text{ord}_A(y)}
\frac{x^{\text{ord}_A(y)}}{y^{\text{ord}_A(x)}} \bmod \mathfrak m_A,
$$
换言之，该符号等于通常的驯符号。
\end{lemma}

\begin{proof}
由乘法性，只需证明 $x, y \in A$ 的情形。设 $t \in A$ 为一致化元。
把 $x = t^bu$ 与 $y = t^bv$ 写成这种形式，其中 $a, b \geq 0$ 且
$u, v \in A^*$。令 $l = a + b$。于是
$t^{l - 1}, \ldots, t^b$ 是 $(x)/(xy)$ 中的容许序列，而
$t^{l - 1}, \ldots, t^a$ 是 $(y)/(xy)$ 中的容许序列。因此由注
\ref{chow-remark-more-elementary}，$d_A(x, y)$ 由等式
$$
[t^{l - 1}, \ldots, t^b, v^{-1}t^{b - 1}, \ldots, v^{-1}]
=
(-1)^{ab} d_A(x, y)
[t^{l - 1}, \ldots, t^a, u^{-1}t^{a - 1}, \ldots, u^{-1}].
$$
刻画。于是由符号 $[x_1, \ldots, x_l]$ 的容许关系可见
$$
d_A(x, y) = (-1)^{ab} u^a/v^b \bmod \mathfrak m_A
$$
正如所需。
\end{proof}

\begin{lemma}
\label{chow-lemma-symbol-is-steinberg-prepare}
设 $A$ 为诺特局部环，$a, b \in A$。设 $M$ 为维数为 $1$ 的有限
$A$-模，并且 $a$、$b$、$b - a$ 在该模上均为非零因子。则
$$
d_M(a, b - a)d_M(b, b) = d_M(b, b - a)d_M(a, b)
$$
在 $\kappa^*$ 中成立。
\end{lemma}

\begin{proof}
由引理 \ref{chow-lemma-compute-symbol-M}，只需在
$M = A/\mathfrak q$ 的情形证明该关系，其中
$\mathfrak q \subset A$ 是满足 $\dim(A/\mathfrak q) = 1$ 的某个素理想。

\medskip\noindent
当 $M = A/\mathfrak q$ 时，可以把 $A$ 换成 $A/\mathfrak q$，并把
$a, b$ 换成它们在 $A/\mathfrak q$ 中的像。因此可假设 $A = M$，而
$A$ 是维数为 $1$ 的诺特局部整环。原因是 $A$ 与 $A/\mathfrak q$ 的
剩余域 $\kappa$ 相同，而且对任意 $A/\mathfrak q$-模 $M$，在 $A$ 上
或 $A/\mathfrak q$ 上取得的行列式典范地认同。见引理
\ref{chow-lemma-determinant-quotient-ring}。

\medskip\noindent
只需在 $a, b$ 都属于极大理想时证明该关系。事实上，其中一个或两个
都是单位的情形由引理 \ref{chow-lemma-symbol-when-one-is-a-unit} 与
\ref{chow-lemma-symbol-when-equal} 得出。

\medskip\noindent
按照引理 \ref{chow-lemma-Noetherian-domain-dim-1-two-elements}，选取扩张
$A \subset A'$ 与分解 $a = ta'$、$b = tb'$。注意还有
$b - a = t(b' - a')$，并且
$A' = (a', b') = (a', b' - a') = (b' - a', b')$。
这里以及下文把 $A'$ 视为 $A$-模，并把 $a, b, a', b', t$ 视为 $A'$
的 $A$-模自同态。我们用记号 $d^A_{A'}(a', b')$ 等表示
$$
d^A_{A'}(a', b')
=
\det\nolimits_\kappa(A'/a'b'A', a', b')
$$
它由引理 \ref{chow-lemma-pre-symbol} 定义。上标 ${}^A$ 用于把它与已经定义的
符号 $d_{A'}(a', b')$ 区分开；后者不同，例如其取值在 $A'$ 的剩余域中，
而该剩余域可能与 $\kappa$ 不同。由引理
\ref{chow-lemma-symbol-compare-modules} 可见
$d_A(a, b) = d^A_{A'}(a, b)$，其他组合也类似。利用这一点与乘法性，
只需证明
$$
d^A_{A'}(a', b' - a') d^A_{A'}(b', b')
=
d^A_{A'}(b', b' - a') d^A_{A'}(a', b')
$$
现在，由于 $(a', b') = A'$ 等，有
$$
\begin{matrix}
A'/(a'(b' - a')) & \cong & A'/(a') \oplus A'/(b' - a') \\
A'/(b'(b' - a')) & \cong & A'/(b') \oplus A'/(b' - a') \\
A'/(a'b') & \cong & A'/(a') \oplus A'/(b')
\end{matrix}
$$
此外，注意乘以 $b' - a'$ 在 $A/(a')$ 上等于乘以 $b'$，而乘以
$b' - a'$ 在 $A/(b')$ 上等于乘以 $-a'$。使用引理
\ref{chow-lemma-periodic-determinant-easy-case} 与
\ref{chow-lemma-periodic-determinant}，得到
$$
\begin{matrix}
d^A_{A'}(a', b' - a') & = &
\det\nolimits_\kappa(b'|_{A'/(a')})^{-1}
\det\nolimits_\kappa(a'|_{A'/(b' - a')}) \\
d^A_{A'}(b', b' - a') & = &
\det\nolimits_\kappa(-a'|_{A'/(b')})^{-1}
\det\nolimits_\kappa(b'|_{A'/(b' - a')}) \\
d^A_{A'}(a', b') & = &
\det\nolimits_\kappa(b'|_{A'/(a')})^{-1}
\det\nolimits_\kappa(a'|_{A'/(b')})
\end{matrix}
$$
因此
$$
(-1)^{\text{length}_A(A'/(b'))}
d^A_{A'}(a', b' - a')
=
d^A_{A'}(b', b' - a') d^A_{A'}(a', b')
$$
其中符号来自上述第二个等式中的 $-a'$。另一方面，由引理
\ref{chow-lemma-periodic-determinant-sign}，有
$d^A_{A'}(b', b') = (-1)^{\text{length}_A(A'/(b'))}$，引理得证。
\end{proof}

\noindent
驯符号是 Steinberg 符号。

\begin{lemma}
\label{chow-lemma-symbol-is-steinberg}
设 $A$ 为分式域为 $K$ 的维数为 $1$ 的诺特局部整环。对
$x \in K \setminus \{0, 1\}$，有
$$
d_A(x, 1 -x) = 1
$$
\end{lemma}

\begin{proof}
把 $x = a/b$ 写成这种形式，其中 $a, b \in A$。由于
$1 - x = (b - a)/b$，假设还蕴含 $b - a \not = 0$。因此计算得
$$
d_A(x, 1 - x)
=
d_A(a, b - a)d_A(a, b)^{-1}d_A(b, b - a)^{-1}d_A(b, b)
$$
所以必须证明
$d_A(a, b - a) d_A(b, b) = d_A(b, b - a) d_A(a, b)$。
这正是引理 \ref{chow-lemma-symbol-is-steinberg-prepare}。
\end{proof}










\subsection{长度与行列式}
\label{chow-subsection-length-determinant}

\noindent
本节用行列式比较格。关键引理如下。

\begin{lemma}
\label{chow-lemma-key-lemma}
设 $R$ 为诺特局部环。设 $\mathfrak q \subset R$ 为满足
$\dim(R/\mathfrak q) = 1$ 的素理想。设
$\varphi : M \to N$ 为有限 $R$-模的同态。假设存在
$x_1, \ldots, x_l \in M$ 与 $y_1, \ldots, y_l \in M$，使得
\begin{enumerate}
\item $M = \langle x_1, \ldots, x_l\rangle$,
\item 对 $i = 1, \ldots, l$，有
$\langle x_1, \ldots, x_i\rangle / \langle x_1, \ldots, x_{i - 1}\rangle
\cong R/\mathfrak q$,
\item $N = \langle y_1, \ldots, y_l\rangle$,
\item 对 $i = 1, \ldots, l$，有
$\langle y_1, \ldots, y_i\rangle / \langle y_1, \ldots, y_{i - 1}\rangle
\cong R/\mathfrak q$.
\end{enumerate}
则 $\varphi$ 为单射当且仅当 $\varphi_{\mathfrak q}$ 为同构，而且此时
$$
\text{length}_R(\Coker(\varphi)) = \text{ord}_{R/\mathfrak q}(f)
$$
其中 $f \in \kappa(\mathfrak q)$ 是满足
$$
[\varphi(x_1), \ldots, \varphi(x_l)] = f [y_1, \ldots, y_l]
$$
的元，上式是 $\det_{\kappa(\mathfrak q)}(N_{\mathfrak q})$ 中的等式。
\end{lemma}

\begin{proof}
首先注意，$l = 1$ 时引理成立。事实上，此时 $x_1$ 是 $M$ 在
$R/\mathfrak q$ 上的一组基，$y_1$ 是 $N$ 在 $R/\mathfrak q$ 上的一组基，
并且对某个 $f \in R$ 有 $\varphi(x_1) = fy_1$。因此 $\varphi$ 为单射
当且仅当 $f \not \in \mathfrak q$。再者，
$\Coker(\varphi) = R/(f, \mathfrak q)$，故由 $\text{ord}_{R/q}(f)$ 的定义
（见代数，定义 \ref{algebra-definition-ord}），引理成立。

\medskip\noindent
实际上，更一般地，假设对某些 $f_i \in R$，$f_i \not \in \mathfrak q$，有
$\varphi(x_i) = f_iy_i$。则诱导映射
$$
\langle x_1, \ldots, x_i\rangle / \langle x_1, \ldots, x_{i - 1}\rangle
\longrightarrow
\langle y_1, \ldots, y_i\rangle / \langle y_1, \ldots, y_{i - 1}\rangle
$$
都是单射，且其余核同构于 $R/(f_i, \mathfrak q)$。从而
$$
\text{length}_R(\Coker(\varphi)) = \sum \text{ord}_{R/\mathfrak q}(f_i).
$$
另一方面，由符号的容许关系 (b)，此时显然有
$$
[\varphi(x_1), \ldots, \varphi(x_l)] = f_1 \ldots f_l [y_1, \ldots, y_l]
$$
因此结论在这种情形也成立。

\medskip\noindent
一般情形对 $l$ 用归纳法证明。假设 $l > 1$。取最小的
$i \in \{1, \ldots, l\}$，使得
$\varphi(x_1) \in \langle y_1, \ldots, y_i\rangle$。我们对 $i$ 作归纳。
若 $i = 1$，则得到交换图
$$
\xymatrix{
0 \ar[r] &
\langle x_1 \rangle \ar[r] \ar[d] &
\langle x_1, \ldots, x_l \rangle \ar[r] \ar[d] &
\langle x_1, \ldots, x_l \rangle / \langle x_1 \rangle \ar[r] \ar[d] &
0 \\
0 \ar[r] &
\langle y_1 \rangle \ar[r] &
\langle y_1, \ldots, y_l \rangle \ar[r] &
\langle y_1, \ldots, y_l \rangle / \langle y_1 \rangle \ar[r] &
0
}
$$
此时引理由蛇形引理和对 $l$ 的归纳得出。现假设 $i > 1$。写成
$\varphi(x_1) = a_1 y_1 + \ldots + a_{i - 1} y_{i - 1} + a y_i$，
其中 $a_j, a \in R$ 且 $a \not \in \mathfrak q$（否则 $i$ 不是最小的）。令
$$
x'_j =
\left\{
\begin{matrix}
x_j & \text{若} & j = 1 \\
ax_j & \text{若} & j \geq 2
\end{matrix}
\right.
\quad\text{且}\quad
y'_j =
\left\{
\begin{matrix}
y_j & \text{若} & j < i \\
ay_j & \text{若} & j \geq i
\end{matrix}
\right.
$$
并令 $M' = \langle x'_1, \ldots, x'_l \rangle$、
$N' = \langle y'_1, \ldots, y'_l \rangle$。由构造有
$\varphi(x'_1) = a_1 y'_1 + \ldots + a_{i - 1} y'_{i - 1} + y'_i$，
而对 $j > 1$ 有
$\varphi(x'_j) = a\varphi(x_i) \in \langle y'_1, \ldots, y'_l\rangle$，
故得到如下 $R$-模与映射的交换图：
$$
\xymatrix{
M' \ar[d] \ar[r]_{\varphi'} & N' \ar[d] \\
M \ar[r]^\varphi & N
}
$$
由证明第二段的结果，我们知道
$\text{length}_R(M/M') = (l - 1)\text{ord}_{R/\mathfrak q}(a)$，并且类似地
$\text{length}_R(M/M') = (l - i + 1)\text{ord}_{R/\mathfrak q}(a)$。
追图给出
$$
\text{length}_R(\Coker(\varphi')) =
\text{length}_R(\Coker(\varphi)) + i\ \text{ord}_{R/\mathfrak q}(a).
$$
另一方面，写作
$$
[\varphi(x_1), \ldots, \varphi(x_l)] = f [y_1, \ldots, y_l],
\quad
[\varphi'(x'_1), \ldots, \varphi(x'_l)] = f' [y'_1, \ldots, y'_l]
$$
时，显然 $f' = a^if$。因此只需对
$\varphi(x_1) = a_1y_1 + \ldots a_{i - 1}y_{i - 1} + y_i$ 的情形，
即 $a = 1$ 的情形证明引理。其次，回忆由符号的容许关系有
$$
[y_1, \ldots, y_l] = [y_1, \ldots, y_{i - 1},
a_1y_1 + \ldots a_{i - 1}y_{i - 1} + y_i, y_{i + 1}, \ldots, y_l]
$$
序列
$y_1, \ldots, y_{i - 1},
a_1y_1 + \ldots + a_{i - 1}y_{i - 1} + y_i, y_{i + 1}, \ldots, y_l$
也满足引理的条件 (3)、(4)。因此实际上可以假设
$\varphi(x_1) = y_i$。在这种情形，注意有 $\mathfrak q x_1 = 0$，
从而也有 $\mathfrak q y_i = 0$。又由定义
$\det_{\kappa(\mathfrak q)}(N_{\mathfrak q})$ 的第三个容许关系，有
$$
[y_1, \ldots, y_l] =
- [y_1, \ldots, y_{i - 2}, y_i, y_{i - 1}, y_{i + 1}, \ldots, y_l]
$$
因此可以把 $y_1, \ldots, y_l$ 换成序列
$y'_1, \ldots, y'_l =
y_1, \ldots, y_{i - 2}, y_i, y_{i - 1}, y_{i + 1}, \ldots, y_l$
（它也满足引理的条件 (3) 和 (4)）。这显然使不变量 $i$ 减少 $1$，
从而对 $i$ 的归纳完成证明。
\end{proof}

\noindent
为了应用前一引理，我们证明具有所需性质的元序列往往存在。

\begin{lemma}
\label{chow-lemma-good-sequence-exists}
设 $R$ 为诺特局部环，$\mathfrak q \subset R$ 为素理想。设 $M$ 为有限
$R$-模，且 $\mathfrak q$ 是 $M$ 的支集的一个极小素理想。则存在
$x_1, \ldots, x_l \in M$，使得
\begin{enumerate}
\item $M / \langle x_1, \ldots, x_l\rangle$ 的支集不包含 $\mathfrak q$；
\item 对 $i = 1, \ldots, l$，有
$\langle x_1, \ldots, x_i\rangle / \langle x_1, \ldots, x_{i - 1}\rangle
\cong R/\mathfrak q$。
\end{enumerate}
而且此时 $l = \text{length}_{R_\mathfrak q}(M_\mathfrak q)$。
\end{lemma}

\begin{proof}
$\mathfrak q$ 是 $M$ 的支集中的极小素理想这一条件意味着
$l = \text{length}_{R_\mathfrak q}(M_\mathfrak q)$ 有限（见代数，引理
\ref{algebra-lemma-support-point}）。因此可以找到
$y_1, \ldots, y_l \in M_{\mathfrak q}$，使得对 $i = 1, \ldots, l$ 有
$\langle y_1, \ldots, y_i\rangle / \langle y_1, \ldots, y_{i - 1}\rangle
\cong \kappa(\mathfrak q)$。可以找到 $f_i \in R$、$f_i \not \in \mathfrak q$，
使 $f_i y_i$ 是某个元 $z_i \in M$ 的像。另外，由于 $R$ 是诺特的，
可对某些 $g_j \in R$ 写成 $\mathfrak q = (g_1, \ldots, g_t)$。由假设，
在模 $M_{\mathfrak q}$ 中有
$g_j y_i \in \langle y_1, \ldots, y_{i - 1} \rangle$。由 $z_i$ 的选取，
可以再找到 $f_{ji} \in R$、$f_{ij} \not \in \mathfrak q$，使得
$f_{ij} g_j z_i \in \langle z_1, \ldots, z_{i - 1} \rangle$
（这是模 $M$ 中的等式）。取
$$
x_1 = f_{11}f_{12}\ldots f_{1t}z_1,
\quad
x_2 = f_{11}f_{12}\ldots f_{1t}f_{21}f_{22}\ldots f_{2t}z_2,
$$
并依此类推，即得引理。事实上，所有元 $f_i, f_{ij}$ 在
$R_{\mathfrak q}$ 中都可逆，因而仍有
$R_{\mathfrak q}x_1 + \ldots + R_{\mathfrak q}x_i /
R_{\mathfrak q}x_1 + \ldots + R_{\mathfrak q}x_{i - 1}
\cong \kappa(\mathfrak q)$，其中 $i = 1, \ldots, l$。由构造，
$\mathfrak q x_i \in \langle x_1, \ldots, x_{i - 1}\rangle$。因此
$\langle x_1, \ldots, x_i\rangle / \langle x_1, \ldots, x_{i - 1}\rangle$
是一个由一个元生成且被 $\mathfrak q$ 消去的 $R$-模，将其在
$\mathfrak q$ 处局部化会得到一个在 $\kappa(\mathfrak q)$ 上的 $q$-维向量空间。
故它同构于 $R/\mathfrak q$。
\end{proof}

\noindent
以下是本节的主要结果。下面将看到这一命题的各种不同推论。
建议读者先自行证明较容易的引理
\ref{chow-lemma-application-herbrand-quotient}。

\begin{proposition}
\label{chow-proposition-length-determinant-periodic-complex}
设 $R$ 为剩余域为 $\kappa$ 的诺特局部环。假设
$(M, \varphi, \psi)$ 是 $R$ 上的一个 $(2, 1)$-周期复形。假设
\begin{enumerate}
\item $M$ 是有限 $R$-模；
\item $(M, \varphi, \psi)$ 的上同调模都具有有限长度；
\item $\dim(\text{Supp}(M)) = 1$。
\end{enumerate}
设 $\mathfrak q_i$、$i = 1, \ldots, t$ 是 $M$ 的支集的极小素理想。则\footnote{
显然，可以通过把 $\det_\kappa(M, \varphi, \psi)$ 重新定义为它当前值的逆
而去掉负号；见定义 \ref{chow-definition-periodic-determinant}。}
$$
- e_R(M, \varphi, \psi) =
\sum\nolimits_{i = 1, \ldots, t}
\text{ord}_{R/\mathfrak q_i}\left(
\det\nolimits_{\kappa(\mathfrak q_i)}
(M_{\mathfrak q_i}, \varphi_{\mathfrak q_i}, \psi_{\mathfrak q_i})
\right)
$$
\end{proposition}

\begin{proof}
首先如下化归到 $t = 1$ 的情形。注意
$\text{Supp}(M) = \{\mathfrak m, \mathfrak q_1, \ldots, \mathfrak q_t\}$，
其中 $\mathfrak m \subset R$ 是极大理想。令 $M_i$ 表示
$M \to M_{\mathfrak q_i}$ 的像，从而
$\text{Supp}(M_i) = \{\mathfrak m, \mathfrak q_i\}$。映射 $\varphi$（分别地，\ $\psi$）
诱导出 $R$-模映射 $\varphi_i : M_i \to M_i$（分别地，\ 
$\psi_i : M_i \to M_i$）。因此得到 $(2, 1)$-周期复形的态射
$$
(M, \varphi, \psi) \longrightarrow
\bigoplus\nolimits_{i = 1, \ldots, t} (M_i, \varphi_i, \psi_i).
$$
该映射的核与余核的支集都包含在 $\{\mathfrak m\}$ 中。故由引理
\ref{chow-lemma-compare-periodic-lengths}，
$$
e_R(M, \varphi, \psi) =
\sum\nolimits_{i = 1, \ldots, t}
e_R(M_i, \varphi_i, \psi_i)
$$
另一方面，显然有 $M_{\mathfrak q_i} = M_{i, \mathfrak q_i}$，因此引理公式
右边的各项等于
$$
\text{ord}_{R/\mathfrak q_i}\left(
\det\nolimits_{\kappa(\mathfrak q_i)}
(M_{i, \mathfrak q_i}, \varphi_{i, \mathfrak q_i}, \psi_{i, \mathfrak q_i})
\right)
$$
换言之，若能对每个模 $M_i$ 证明引理，则引理成立。这便化归到
$t = 1$ 的情形。

\medskip\noindent
假设在诺特局部环上有一个 $(2, 1)$-周期复形
$(M, \varphi, \psi)$，其中 $M$ 为有限 $R$-模，
$\text{Supp}(M) = \{\mathfrak m, \mathfrak q\}$，且上同调模具有有限长度。
在这种情形，证明由引理 \ref{chow-lemma-key-lemma} 和仔细的计数得出。记
$K_\varphi = \Ker(\varphi)$、
$I_\varphi = \Im(\varphi)$、
$K_\psi = \Ker(\psi)$ 以及
$I_\psi = \Im(\psi)$。由于 $R$ 是诺特的，它们都是有限 $R$-模。令
$$
a = \text{length}_{R_{\mathfrak q}}(I_{\varphi, \mathfrak q})
= \text{length}_{R_{\mathfrak q}}(K_{\psi, \mathfrak q}),
\quad
b = \text{length}_{R_{\mathfrak q}}(I_{\psi, \mathfrak q})
= \text{length}_{R_{\mathfrak q}}(K_{\varphi, \mathfrak q}).
$$
等式成立是因为复形在 $\mathfrak q$ 处局部化后变得正合。注意
$l = \text{length}_{R_{\mathfrak q}}(M_{\mathfrak q})$ 等于 $l = a + b$。

\medskip\noindent
我们将应用引理 \ref{chow-lemma-good-sequence-exists}，在支集包含于
$\{\mathfrak m, \mathfrak q\}$ 的有限 $R$-模 $N$ 中选取元序列。在这种情形，
$N_{\mathfrak q}$ 有有限长度，记为 $n \in \mathbf{N}$。我们把满足引理
\ref{chow-lemma-good-sequence-exists} 的性质 (1) 和 (2) 的序列
$w_1, \ldots, w_n \in N$ 称为“良序列”。注意，$N$ 对由良序列生成的子模所取的商
$N/\langle w_1, \ldots, w_n \rangle$ 的支集包含于 $\{\mathfrak m\}$，
因而它具有有限长度（代数，引理 \ref{algebra-lemma-support-point}）。再者，符号
$[w_1, \ldots, w_n] \in \det_{\kappa(\mathfrak q)}(N_{\mathfrak q})$
是一个生成元；见引理 \ref{chow-lemma-determinant-dimension-one}。

\medskip\noindent
在此约定下，选取良序列
$$
\begin{matrix}
x_1, \ldots, x_b & \text{于} & K_\varphi, &
t_1, \ldots, t_a & \text{于} & K_\psi, \\
y_1, \ldots, y_a & \text{于} & I_\varphi \cap \langle t_1, \ldots t_a\rangle, &
s_1, \ldots, s_b & \text{于} & I_\psi \cap \langle x_1, \ldots, x_b\rangle.
\end{matrix}
$$
我们将如下对这些选择作少许调整。选取 $y_i \in I_\varphi$ 在 M 中的提升
$\tilde y_i \in M$，以及 $s_i \in I_\psi$ 在 M 中的提升
$\tilde s_i \in M$。未必有
$\mathfrak q \tilde y_1 \subset \langle x_1, \ldots, x_b\rangle$，也未必有
$\mathfrak q \tilde s_1 \subset \langle t_1, \ldots, t_a\rangle$。然而，利用
$\mathfrak q$ 有限生成（如引理 \ref{chow-lemma-good-sequence-exists} 的证明中那样），
可以找到 $d \in R$、$d \not \in \mathfrak q$，使得
$\mathfrak q d\tilde y_1 \subset \langle x_1, \ldots, x_b\rangle$
且
$\mathfrak q d\tilde s_1 \subset \langle t_1, \ldots, t_a\rangle$。
因此，用 $dy_i$ 替换 $y_i$，用 $d\tilde y_i$ 替换 $\tilde y_i$，
用 $ds_i$ 替换 $s_i$，并用 $d\tilde s_i$ 替换 $\tilde s_i$ 后，还可以假设
$x_1, \ldots, x_b, \tilde y_1, \ldots, \tilde y_b$
与 $t_1, \ldots, t_a, \tilde s_1, \ldots, \tilde s_b$
都是 $M$ 中的良序列。

\medskip\noindent
最后，在有限 $R$-模
$$
\langle
x_1, \ldots, x_b, \tilde y_1, \ldots, \tilde y_a
\rangle
\cap
\langle
t_1, \ldots, t_a, \tilde s_1, \ldots, \tilde s_b
\rangle.
$$
中选取一个良序列 $z_1, \ldots, z_l$。注意，这也是 $M$ 中的良序列。

\medskip\noindent
由于 $I_{\varphi, \mathfrak q} = K_{\psi, \mathfrak q}$，存在唯一元
$h \in \kappa(\mathfrak q)$，使得在
$\det_{\kappa(\mathfrak q)}(K_{\psi, \mathfrak q})$ 中有
$[y_1, \ldots, y_a] = h [t_1, \ldots, t_a]$。类似地，由于
$I_{\psi, \mathfrak q} = K_{\varphi, \mathfrak q}$，存在唯一元
$h \in \kappa(\mathfrak q)$，使得在
$\det_{\kappa(\mathfrak q)}(K_{\varphi, \mathfrak q})$ 中有
$[s_1, \ldots, s_b] = g [x_1, \ldots, x_b]$。对 $M$ 中的三个良序列也可以这样做。
总之，得到恒等式
\begin{align*}
[y_1, \ldots, y_a]
& =
h [t_1, \ldots, t_a] \\
[s_1, \ldots, s_b]
& =
g [x_1, \ldots, x_b] \\
[z_1, \ldots, z_l]
& =
f_\varphi [x_1, \ldots, x_b, \tilde y_1, \ldots, \tilde y_a] \\
[z_1, \ldots, z_l]
& =
f_\psi [t_1, \ldots, t_a, \tilde s_1, \ldots, \tilde s_b]
\end{align*}
其中 $g, h, f_\varphi, f_\psi \in \kappa(\mathfrak q)$。

\medskip\noindent
记号已全部建立，现在计算
$\det_{\kappa(\mathfrak q)}(M, \varphi, \psi)$。考虑元
$[z_1, \ldots, z_l]$。在定义 \ref{chow-definition-periodic-determinant} 的映射
$\gamma_\psi \circ \sigma \circ \gamma_\varphi^{-1}$ 下，有
\begin{eqnarray*}
[z_1, \ldots, z_l] & = &
f_\varphi [x_1, \ldots, x_b, \tilde y_1, \ldots, \tilde y_a] \\
& \mapsto & f_\varphi [x_1, \ldots, x_b] \otimes [y_1, \ldots, y_a] \\
& \mapsto &
f_\varphi h/g [t_1, \ldots, t_a] \otimes [s_1, \ldots, s_b] \\
& \mapsto &
f_\varphi h/g [t_1, \ldots, t_a, \tilde s_1, \ldots, \tilde s_b] \\
& = &
f_\varphi h/f_\psi g [z_1, \ldots, z_l]
\end{eqnarray*}
这意味着，除去一个符号后，
$\det_{\kappa(\mathfrak q)}
(M_{\mathfrak q}, \varphi_{\mathfrak q}, \psi_{\mathfrak q})$
等于 $f_\varphi h/f_\psi g$。

\medskip\noindent
对以下量作简记：
\begin{eqnarray*}
k_\varphi & = & \text{length}_R(K_\varphi/\langle x_1, \ldots, x_b\rangle) \\
k_\psi    & = & \text{length}_R(K_\psi/\langle t_1, \ldots, t_a\rangle) \\
i_\varphi & = & \text{length}_R(I_\varphi/\langle y_1, \ldots, y_a\rangle) \\
i_\psi    & = & \text{length}_R(I_\psi/\langle s_1, \ldots, s_a\rangle) \\
m_\varphi & = & \text{length}_R(M/
\langle x_1, \ldots, x_b, \tilde y_1, \ldots, \tilde y_a\rangle) \\
m_\psi    & = & \text{length}_R(M/
\langle t_1, \ldots, t_a, \tilde s_1, \ldots, \tilde s_b\rangle) \\
\delta_\varphi & = & \text{length}_R(
\langle x_1, \ldots, x_b, \tilde y_1, \ldots, \tilde y_a\rangle
\langle z_1, \ldots, z_l\rangle) \\
\delta_\psi & = & \text{length}_R(
\langle t_1, \ldots, t_a, \tilde s_1, \ldots, \tilde s_b\rangle
\langle z_1, \ldots, z_l\rangle)
\end{eqnarray*}
使用正合列 $0 \to K_\varphi \to M \to I_\varphi \to 0$，得
$m_\varphi = k_\varphi + i_\varphi$。类似地，有
$m_\psi = k_\psi + i_\psi$。又有
$\delta_\varphi + m_\varphi = \delta_\psi + m_\psi$，因为两边都等于
$\langle z_1, \ldots, z_l \rangle$ 在 $M$ 中的余长度。最后，首次应用关键引理
\ref{chow-lemma-key-lemma} 得
$$
\delta_\varphi = \text{ord}_{R/\mathfrak q}(f_\varphi),
\quad
\delta_\psi = \text{ord}_{R/\mathfrak q}(f_\psi)
$$

\medskip\noindent
接下来计算周期复形的重数：
\begin{eqnarray*}
e_R(M, \varphi, \psi) & = &
\text{length}_R(K_\varphi/I_\psi) - \text{length}_R(K_\psi/I_\varphi) \\
& = &
\text{length}_R(
\langle x_1, \ldots, x_b\rangle/
\langle s_1, \ldots, s_b\rangle)
+ k_\varphi - i_\psi \\
& & -
\text{length}_R(
\langle t_1, \ldots, t_a\rangle/
\langle y_1, \ldots, y_a\rangle)
- k_\psi + i_\varphi \\
& = &
\text{ord}_{R/\mathfrak q}(g/h) + k_\varphi - i_\psi - k_\psi + i_\varphi \\
& = &
\text{ord}_{R/\mathfrak q}(g/h) + m_\varphi - m_\psi \\
& = &
\text{ord}_{R/\mathfrak q}(g/h) + \delta_\psi - \delta_\varphi \\
& = &
\text{ord}_{R/\mathfrak q}(f_\psi g/f_\varphi h)
\end{eqnarray*}
其中在第三个等式中两次用到了关键引理
\ref{chow-lemma-key-lemma}。根据前面对
$\det_{\kappa(\mathfrak q)}
(M_{\mathfrak q}, \varphi_{\mathfrak q}, \psi_{\mathfrak q})$ 的计算，命题得证。
\end{proof}

\noindent
在多数应用中，以下引理已经足够。

\begin{lemma}
\label{chow-lemma-application-herbrand-quotient}
设 $R$ 为极大理想为 $\mathfrak m$ 的诺特局部环。设 $M$ 为有限
$R$-模，$\psi : M \to M$ 为 $R$-模映射。假设
\begin{enumerate}
\item $\Ker(\psi)$ 与 $\Coker(\psi)$ 具有有限长度；
\item $\dim(\text{Supp}(M)) \leq 1$。
\end{enumerate}
写成
$\text{Supp}(M) = \{\mathfrak m, \mathfrak q_1, \ldots, \mathfrak q_t\}$，
并用 $f_i \in \kappa(\mathfrak q_i)^*$ 表示使得映射
$\det_{\kappa(\mathfrak q_i)}(\psi_{\mathfrak q_i}) :
\det_{\kappa(\mathfrak q_i)}(M_{\mathfrak q_i})
\to \det_{\kappa(\mathfrak q_i)}(M_{\mathfrak q_i})$
等于乘以 $f_i$ 的那个元。则
$$
\text{length}_R(\Coker(\psi))
-
\text{length}_R(\Ker(\psi))
=
\sum\nolimits_{i = 1, \ldots, t}
\text{ord}_{R/\mathfrak q_i}(f_i).
$$
\end{lemma}

\begin{proof}
回忆 $H^0(M, 0, \psi) = \Coker(\psi)$ 且
$H^1(M, 0, \psi) = \Ker(\psi)$；见定义
\ref{chow-definition-periodic-length} 之前的注记。联合命题
\ref{chow-proposition-length-determinant-periodic-complex} 与引理
\ref{chow-lemma-periodic-determinant-easy-case}，即得本引理。

\medskip\noindent
另一证明。如命题 \ref{chow-proposition-length-determinant-periodic-complex} 的证明那样，
化归到 $\text{Supp}(M) = \{\mathfrak m, \mathfrak q\}$ 的情形。然后直接联合引理
\ref{chow-lemma-key-lemma} 与 \ref{chow-lemma-good-sequence-exists}，证明命题
\ref{chow-proposition-length-determinant-periodic-complex} 的这一特殊情形。此时需要的计数少得多，
建议读者自行完成。细节从略。
\end{proof}




\subsection{关键引理的应用}
\label{chow-subsection-application-tame-symbol}

\noindent
本节应用上述结果，利用上面构造的驯符号 $d_A$ 证明关键引理
（引理 \ref{chow-lemma-milnor-gersten-low-degree}）的类似结果。其与 Milnor $K$-理论的关系见注
\ref{chow-remark-gersten-complex-milnor}。

\begin{lemma}[关键引理]
\label{chow-lemma-secondary-ramification}
\begin{reference}
当 $A$ 为优良环时，这是 \cite[Proposition 1]{Kato-Milnor-K}。
\end{reference}
设 $A$ 为分式域为 $K$ 的 $2$-维诺特局部整环。设 $f, g \in K^*$。设
$\mathfrak q_1, \ldots, \mathfrak q_t$ 是 $A$ 的高为 $1$ 的素理想 $\mathfrak q$，
且 $f$ 或 $g$ 中至少有一个不是 $A^*_{\mathfrak q}$ 的元。则
$$
\sum\nolimits_{i = 1, \ldots, t}
\text{ord}_{A/\mathfrak q_i}(d_{A_{\mathfrak q_i}}(f, g))
=
0
$$
也可以写成
$$
\sum\nolimits_{\text{高度}(\mathfrak q) = 1}
\text{ord}_{A/\mathfrak q}(d_{A_{\mathfrak q}}(f, g))
=
0
$$
因为对 $A$ 的任意高为一的素理想 $\mathfrak q$，若
$f, g \in A^*_{\mathfrak q}$，则由引理
\ref{chow-lemma-symbol-when-one-is-a-unit} 有 $d_{A_{\mathfrak q}}(f, g) = 1$。
\end{lemma}

\begin{proof}
驯符号 $d_{A_{\mathfrak q}}(f, g)$ 具有可加性（引理
\ref{chow-lemma-multiplicativity-symbol}），而阶函数
$\text{ord}_{A/\mathfrak q}$ 也具有可加性（代数，引理
\ref{algebra-lemma-ord-additive}），故只需在 $f = a \in A$ 且
$g = b \in A$ 时证明公式。在这种情形，必须证明
$$
\sum\nolimits_{\text{高度}(\mathfrak q) = 1}
\text{ord}_{A/\mathfrak q}(\det\nolimits_\kappa(A_{\mathfrak q}/(ab), a, b))
= 0
$$
由命题 \ref{chow-proposition-length-determinant-periodic-complex}，这等价于证明
$$
e_A(A/(ab), a, b) = 0.
$$
由于复形
$A/(ab) \xrightarrow{a} A/(ab) \xrightarrow{b} A/(ab) \xrightarrow{a} A/(ab)$
是正合的，结论得证。
\end{proof}





\section{附录 B：其他方法}
\label{chow-section-appendix-chow}

\noindent
在本附录中，我们首先简要尝试将正文材料与相干层的 $K$-理论联系起来。
特别地，我们说明与可逆模的 $c_1$ 作杯积如何与张量乘这个可逆模相关；
见引理 \ref{chow-lemma-coherent-sheaf-cap-c1}。这些材料显然非常有趣，值得作更详细、
更充分的阐述。


\subsection{有理等价与 K-群}
\label{chow-subsection-rational-equivalence-K-groups}

\noindent
本节是第 \ref{chow-section-chow-and-K} 节的继续。以下引理的动机来自同调，引理
\ref{homology-lemma-serre-subcategory-K-groups}。

\begin{lemma}
\label{chow-lemma-maps-between-coherent-sheaves}
设 $(S, \delta)$ 如情形 \ref{chow-situation-setup} 中所述。设 $X$ 为 $S$ 上局部有限型的概形，
$\mathcal{F}$ 为 $X$ 上的相干层。设
$$
\xymatrix{
\ldots \ar[r] &
\mathcal{F} \ar[r]^\varphi &
\mathcal{F} \ar[r]^\psi &
\mathcal{F} \ar[r]^\varphi &
\mathcal{F} \ar[r] & \ldots
}
$$
是同调，等式 (\ref{homology-equation-cyclic-complex}) 中的复形。假设
\begin{enumerate}
\item $\dim_\delta(\text{Supp}(\mathcal{F})) \leq k + 1$。
\item 对 $i = 0, 1$，有
$\dim_\delta(\text{Supp}(H^i(\mathcal{F}, \varphi, \psi))) \leq k$。
\end{enumerate}
则作为 $X$ 上的 $k$-维循环，有
$$
[H^0(\mathcal{F}, \varphi, \psi)]_k
\sim_{rat}
[H^1(\mathcal{F}, \varphi, \psi)]_k
$$
\end{lemma}

\begin{proof}
令 $\{W_j\}_{j \in J}$ 为 $\text{Supp}(\mathcal{F})$ 中具有 $\delta$-维数
$k + 1$ 的不可约分支所成的集合。由引理 \ref{chow-lemma-length-finite}，
$\{W_j\}$ 是 $X$ 的闭子集的局部有限族。对每个 $j$，令
$\xi_j \in W_j$ 为一般点，并令
$$
f_j = \det\nolimits_{\kappa(\xi_j)}
(\mathcal{F}_{\xi_j}, \varphi_{\xi_j}, \psi_{\xi_j})
\in
R(W_j)^*.
$$
记号见定义 \ref{chow-definition-periodic-determinant}。我们声称
$$
- [H^0(\mathcal{F}, \varphi, \psi)]_k + [H^1(\mathcal{F}, \varphi, \psi)]_k
=
\sum (W_j \to X)_*\text{div}(f_j)
$$
若证明这一等式，引理即随之得出。

\medskip\noindent
设 $Z \subset X$ 为 $\delta$-维数为 $k$ 的整闭子概形。要证明上述等式，
只需证明 $[Z]$ 在
$
[H^0(\mathcal{F}, \varphi, \psi)]_k - [H^1(\mathcal{F}, \varphi, \psi)]_k
$
中的系数 $n$ 等于 $[Z]$ 在
$
\sum (W_j \to X)_*\text{div}(f_j)
$
中的系数 $m$。令 $\xi \in Z$ 为一般点。考虑局部环
$A = \mathcal{O}_{X, \xi}$。把 $M = \mathcal{F}_\xi$ 视为 $A$-模，并用
$\varphi, \psi : M \to M$ 表示 $\varphi, \psi$ 在茎上的作用。由 $\xi \in Z$ 的选取，
有 $\delta(\xi) = k$，从而 $\dim(\text{Supp}(M)) = 1$。最后，经过 $\xi$ 的整闭子概形
$W_j$ 与 $\text{Supp}(M)$ 的极小素理想 $\mathfrak q_i$ 对应。在每种情形，
元 $f_j \in R(W_j)^*$ 对应于 $\kappa(\mathfrak q_i)^*$ 中的元
$\det_{\kappa(\mathfrak q_i)}(M_{\mathfrak q_i}, \varphi, \psi)$。因此
$$
n = - e_A(M, \varphi, \psi)
$$
且
$$
m =
\sum
\text{ord}_{A/\mathfrak q_i}
(\det\nolimits_{\kappa(\mathfrak q_i)}(M_{\mathfrak q_i}, \varphi, \psi))
$$
于是结论由命题 \ref{chow-proposition-length-determinant-periodic-complex} 得出。
\end{proof}

\begin{lemma}
\label{chow-lemma-cycles-rational-equivalence-K-group}
设 $(S, \delta)$ 如情形 \ref{chow-situation-setup} 中所述，$X$ 为 $S$ 上局部有限型的概形。
引理 \ref{chow-lemma-from-chow-to-K} 中的映射
$$
\CH_k(X) \longrightarrow
K_0(\textit{Coh}_{\leq k + 1}(X)/\textit{Coh}_{\leq k - 1}(X))
$$
诱导出从 $\CH_k(X)$ 到下列映射的像 $B_k(X)$ 上的双射：
$$
K_0(\textit{Coh}_{\leq k}(X)/\textit{Coh}_{\leq k - 1}(X))
\longrightarrow
K_0(\textit{Coh}_{\leq k + 1}(X)/\textit{Coh}_{\leq k - 1}(X)).
$$
\end{lemma}

\begin{proof}
由引理 \ref{chow-lemma-cycles-k-group}，我们有
$Z_k(X) = K_0(\textit{Coh}_{\leq k}(X)/\textit{Coh}_{\leq k - 1}(X))$，
且这一等同与引理 \ref{chow-lemma-from-chow-to-K} 的映射相容。因此，假设
$K_0(\textit{Coh}_{\leq k}(X)/\textit{Coh}_{\leq k - 1}(X))$ 中有一个元
$[A] - [B]$，它在 $B_k(X)$ 中映到零，即在
$K_0(\textit{Coh}_{\leq k + 1}(X)/\textit{Coh}_{\leq k - 1}(X))$ 中映到零。
我们必须证明 $[A] - [B]$ 对应于 $X$ 上与零有理等价的循环。假设
$[A] = [\mathcal{A}]$ 且 $[B] = [\mathcal{B}]$，其中 $\mathcal{A}, \mathcal{B}$ 是 $X$ 上的相干层，
支撑在 $\delta$-维数 $\leq k$ 的部分上。元 $[A] - [B]$ 在群
$K_0(\textit{Coh}_{\leq k + 1}(X)/\textit{Coh}_{\leq k - 1}(X))$ 中映到零这一假设意味着：
存在 $X$ 上支撑在 $\delta$-维数 $\leq k - 1$ 的部分上的相干层
$\mathcal{A}', \mathcal{B}'$，使得
$[\mathcal{A} \oplus \mathcal{A}'] - [\mathcal{B} \oplus \mathcal{B}']$ 在
$K_0(\textit{Coh}_{k + 1}(X))$ 中为零（使用同调，引理
\ref{homology-lemma-serre-subcategory-K-groups} 的 (1)）。由同调，引理
\ref{homology-lemma-serre-subcategory-K-groups} 的 (2)，这意味着在范畴
$\textit{Coh}_{\leq k + 1}(X)$ 中存在一个 $(2, 1)$-周期复形
$(\mathcal{F}, \varphi, \psi)$，使得
$\mathcal{A} \oplus \mathcal{A}' = H^0(\mathcal{F}, \varphi, \psi)$
且 $\mathcal{B} \oplus \mathcal{B}' = H^1(\mathcal{F}, \varphi, \psi)$。
由引理 \ref{chow-lemma-maps-between-coherent-sheaves}，这意味着
$$
[\mathcal{A} \oplus \mathcal{A}']_k
\sim_{rat}
[\mathcal{B} \oplus \mathcal{B}']_k
$$
这证明了，在引理 \ref{chow-lemma-cycles-k-group} 的映射
$$
K_0(\textit{Coh}_{\leq k}(X)/\textit{Coh}_{\leq k - 1}(X))
\longrightarrow Z_k(X)
$$
下，$[A] - [B]$ 映到一个与零有理等价的循环。这正是所需证明的，证明完成。
\end{proof}














\subsection{Cartier 除子与 K-群}
\label{chow-subsection-cartier-coherent}

\noindent
本节说明：与可逆层 $\mathcal{L}$ 的第一 Chern 类相交，如何对应于在
$K$-群中张量乘 $\mathcal{L} - \mathcal{O}$。

\begin{lemma}
\label{chow-lemma-no-embedded-points-modules}
设 $A$ 为诺特局部环，$M$ 为有限 $A$-模，$a, b \in A$。假设
\begin{enumerate}
\item $\dim(A) = 1$,
\item $a$ 和 $b$ 都是 $A$ 中的非零因子；
\item $A$ 没有嵌入素理想；
\item $M$ 没有嵌入伴随素理想；
\item $\text{Supp}(M) = \Spec(A)$。
\end{enumerate}
令 $I = \{x \in A \mid x(a/b) \in A\}$。设
$\mathfrak q_1, \ldots, \mathfrak q_t$ 为 $A$ 的极小素理想。则
$(a/b)IM \subset M$，且
$$
\text{length}_A(M/(a/b)IM)
-
\text{length}_A(M/IM)
=
\sum\nolimits_i
\text{length}_{A_{\mathfrak q_i}}(M_{\mathfrak q_i})
\text{ord}_{A/\mathfrak q_i}(a/b)
$$
\end{lemma}

\begin{proof}
因为 $M$ 没有嵌入伴随素理想，且 $M$ 的支集是 $\Spec(A)$，故
$\text{Ass}(M) = \{\mathfrak q_1, \ldots, \mathfrak q_t\}$。因此 $a$、$b$ 在 $M$ 上都是非零因子。注意
\begin{align*}
& \text{length}_A(M/(a/b)IM) \\
& = \text{length}_A(bM/aIM) \\
& = \text{length}_A(M/aIM)
-
\text{length}_A(M/bM) \\
& = \text{length}_A(M/aM) + \text{length}_A(aM/aIM) - \text{length}_A(M/bM) \\
& = \text{length}_A(M/aM) + \text{length}_A(M/IM) - \text{length}_A(M/bM)
\end{align*}
这是因为单射 $b : M \to bM$ 把 $(a/b)IM$ 映到 $aIM$，而单射
$a : M \to aM$ 把 $IM$ 映到 $aIM$。因此引理等式的左边等于
$$
\text{length}_A(M/aM) - \text{length}_A(M/bM).
$$
分别以 $x = a, b$ 应用引理 \ref{chow-lemma-additivity-divisors-restricted} 的第二个公式，
并使用代数，定义 \ref{algebra-definition-ord} 中 $\text{ord}$-函数的定义，即得结论。
\end{proof}

\begin{lemma}
\label{chow-lemma-no-embedded-points}
设 $(S, \delta)$ 如情形 \ref{chow-situation-setup} 中所述。设 $X$ 为 $S$ 上局部有限型的概形，
$\mathcal{L}$ 为可逆 $\mathcal{O}_X$-模，$\mathcal{F}$ 为相干 $\mathcal{O}_X$-模，而
$s \in \Gamma(X, \mathcal{K}_X(\mathcal{L}))$ 为 $\mathcal{L}$ 的一个亚纯截面。假设
\begin{enumerate}
\item $\dim_\delta(X) \leq k + 1$,
\item $X$ 没有嵌入点；
\item $\mathcal{F}$ 没有嵌入伴随点；
\item $\mathcal{F}$ 的支集是 $X$；
\item 截面 $s$ 是正则亚纯的。
\end{enumerate}
在这种情形，令 $\mathcal{I} \subset \mathcal{O}_X$ 为 $s$ 的分母理想；见除子，定义
\ref{divisors-definition-regular-meromorphic-ideal-denominators}。则：
\begin{enumerate}
\item 存在短正合列
$$
\begin{matrix}
0 &
\to &
\mathcal{I}\mathcal{F} &
\xrightarrow{1} &
\mathcal{F} &
\to &
\mathcal{Q}_1 &
\to &
0 \\
0 &
\to &
\mathcal{I}\mathcal{F} &
\xrightarrow{s} &
\mathcal{F} \otimes_{\mathcal{O}_X} \mathcal{L} &
\to &
\mathcal{Q}_2 &
\to &
0
\end{matrix}
$$
\item 相干层 $\mathcal{Q}_1$、$\mathcal{Q}_2$ 支撑在 $\delta$-维数 $\leq k$ 的部分上；
\item 在 $X$ 的每个 $\delta$-维数为 $k + 1$ 的不可约分支 $X_i$ 上，截面 $s$ 限制为
一个正则亚纯截面 $s_i$；
\item 写成 $[\mathcal{F}]_{k + 1} = \sum m_i[X_i]$，则在 $Z_k(X)$ 中
$$
[\mathcal{Q}_2]_k - [\mathcal{Q}_1]_k
=
\sum m_i(X_i \to X)_*\text{div}_{\mathcal{L}|_{X_i}}(s_i)
$$
特别地，在 $\CH_k(X)$ 中
$$
[\mathcal{Q}_2]_k - [\mathcal{Q}_1]_k
=
c_1(\mathcal{L}) \cap [\mathcal{F}]_{k + 1}
$$
\end{enumerate}
\end{lemma}

\begin{proof}
回忆除子，引理 \ref{divisors-lemma-make-maps-regular-section} 给出了单射
$1 : \mathcal{I}\mathcal{F} \to \mathcal{F}$ 与
$s : \mathcal{I}\mathcal{F} \to \mathcal{F} \otimes_{\mathcal{O}_X}\mathcal{L}$，它们的余核支撑在闭的稀疏子集 $T$ 上。
按引理的记号，用 $\mathcal{Q}_i$ 表示这些余核。从而
$\dim_\delta(\text{Supp}(\mathcal{Q}_i)) \leq k$。由除子，引理
\ref{divisors-lemma-pullback-meromorphic-sections-defined} 与
\ref{divisors-lemma-meromorphic-sections-pullback}，拉回 $s_i$ 有定义，且是
$\mathcal{L}|_{X_i}$ 的正则亚纯截面。(4) 中循环的等式意味着 (4) 中循环类的等式。
因此只剩下证明
$$
[\mathcal{Q}_2]_k - [\mathcal{Q}_1]_k
=
\sum m_i(X_i \to X)_*\text{div}_{\mathcal{L}|_{X_i}}(s_i)
$$
在 $Z_k(X)$ 中成立。为此，设 $Z \subset X$ 为 $\delta$-维数为 $k$ 的整闭子概形，
$\xi \in Z$ 为一般点。令 $A = \mathcal{O}_{X, \xi}$ 且 $M = \mathcal{F}_\xi$。再选取生成元
$s_\xi \in \mathcal{L}_\xi$。则可写成 $s = (a/b) s_\xi$，其中 $a, b \in A$ 是非零因子。
此时 $I = \mathcal{I}_\xi = \{x \in A \mid x(a/b) \in A\}$。左边中 $[Z]$ 的系数为
$$
\text{length}_A(M/(a/b)IM) - \text{length}_A(M/IM)
$$
而右边中 $[Z]$ 的系数为
$$
\sum
\text{length}_{A_{\mathfrak q_i}}(M_{\mathfrak q_i})
\text{ord}_{A/\mathfrak q_i}(a/b)
$$
其中 $\mathfrak q_1, \ldots, \mathfrak q_t$ 是 $1$-维局部环 $A$ 的极小素理想。因此结论由引理
\ref{chow-lemma-no-embedded-points-modules} 得出。
\end{proof}

\begin{lemma}
\label{chow-lemma-coherent-sheaf-cap-c1}
设 $(S, \delta)$ 如情形 \ref{chow-situation-setup} 中所述。设 $X$ 为 $S$ 上局部有限型的概形，
$\mathcal{L}$ 为可逆 $\mathcal{O}_X$-模，$\mathcal{F}$ 为相干 $\mathcal{O}_X$-模。假设
$\dim_\delta(\text{Supp}(\mathcal{F})) \leq k + 1$。则元
$$
[\mathcal{F} \otimes_{\mathcal{O}_X} \mathcal{L}]
-
[\mathcal{F}]
\in
K_0(\textit{Coh}_{\leq k + 1}(X)/\textit{Coh}_{\leq k - 1}(X))
$$
属于引理 \ref{chow-lemma-cycles-rational-equivalence-K-group} 中的子群 $B_k(X)$，并且在映射
$B_k(X) \to \CH_k(X)$ 下映到元 $c_1(\mathcal{L}) \cap [\mathcal{F}]_{k + 1}$。
\end{lemma}

\begin{proof}
设
$$
0 \to \mathcal{K} \to \mathcal{F} \to \mathcal{F}' \to 0
$$
是除子，引理 \ref{divisors-lemma-remove-embedded-points} 中构造的短正合列。特别地，
$\mathcal{F}'$ 没有嵌入伴随点。由于 $\mathcal{K}$ 的支集在 $\mathcal{F}$ 的支集中稀疏，有
$\dim_\delta(\text{Supp}(\mathcal{K})) \leq k$。从 $\mathcal{K}$ 出发，再次应用除子，引理
\ref{divisors-lemma-remove-embedded-points}，得到短正合列
$$
0 \to \mathcal{K}'' \to \mathcal{K} \to \mathcal{K}' \to 0
$$
其中 $\dim_\delta(\text{Supp}(\mathcal{K}'')) < k$，而 $\mathcal{K}'$ 没有嵌入伴随点。
假设能对相干层 $\mathcal{F}'$ 和 $\mathcal{K}'$ 证明引理。使用等式
$$
[\mathcal{F}]_{k + 1}
=
[\mathcal{F}']_{k + 1}
+ [\mathcal{K}']_{k + 1}
+ [\mathcal{K}'']_{k + 1}
$$
（使用引理 \ref{chow-lemma-additivity-sheaf-cycle}），以及
$$
[\mathcal{F} \otimes_{\mathcal{O}_X} \mathcal{L}]
-
[\mathcal{F}]
=
[\mathcal{F}' \otimes_{\mathcal{O}_X} \mathcal{L}]
-
[\mathcal{F}']
+
[\mathcal{K}' \otimes_{\mathcal{O}_X} \mathcal{L}]
-
[\mathcal{K}']
+
[\mathcal{K}'' \otimes_{\mathcal{O}_X} \mathcal{L}]
-
[\mathcal{K}'']
$$
（使用 $\otimes \mathcal{L}$ 的正合性），再利用 $[\mathcal{K}'']_{k + 1}$ 以及
$[\mathcal{K}'' \otimes_{\mathcal{O}_X} \mathcal{L}]
- [\mathcal{K}'']$ 在
$K_0(\textit{Coh}_{\leq k + 1}(X)/\textit{Coh}_{\leq k - 1}(X))$ 中显然为零，可见结论对
$\mathcal{F}$ 也成立。这说明可以假设层 $\mathcal{F}$ 没有嵌入伴随点。

\medskip\noindent
假设 $X$、$\mathcal{F}$ 如引理中所述，并进一步假设 $\mathcal{F}$ 没有嵌入伴随点。考虑除子，
引理 \ref{divisors-lemma-no-embedded-points-endos} 中构造的理想层
$\mathcal{I} \subset \mathcal{O}_X$、相应的闭子概形 $i : Z \to X$ 以及相干 $\mathcal{O}_Z$-模
$\mathcal{G}$。回忆 $Z$ 是一个没有嵌入点的局部诺特概形，$\mathcal{G}$ 是没有嵌入伴随点的相干层，
且 $\text{Supp}(\mathcal{G}) = Z$、$i_*\mathcal{G} = \mathcal{F}$。再令
$\mathcal{N} = \mathcal{L}|_Z$。

\medskip\noindent
由除子，引理 \ref{divisors-lemma-regular-meromorphic-section-exists}，可逆层 $\mathcal{N}$ 在 $Z$ 上有一个
正则亚纯截面 $s$。用 $\mathcal{J} \subset \mathcal{O}_Z$ 表示 $s$ 的分母层。由引理
\ref{chow-lemma-no-embedded-points}，存在短正合列
$$
\begin{matrix}
0 &
\to &
\mathcal{J}\mathcal{G} &
\xrightarrow{1} &
\mathcal{G} &
\to &
\mathcal{Q}_1 &
\to &
0 \\
0 &
\to &
\mathcal{J}\mathcal{G} &
\xrightarrow{s} &
\mathcal{G} \otimes_{\mathcal{O}_Z} \mathcal{N} &
\to &
\mathcal{Q}_2 &
\to &
0
\end{matrix}
$$
使得 $\dim_\delta(\text{Supp}(\mathcal{Q}_i)) \leq k$，且循环
$
[\mathcal{Q}_2]_k - [\mathcal{Q}_1]_k
$
是 $c_1(\mathcal{N}) \cap [\mathcal{G}]_{k + 1}$ 的一个代表。由投影公式（见上同调，引理
\ref{cohomology-lemma-projection-formula}）有
$i_*(\mathcal{G} \otimes \mathcal{N}) = \mathcal{F} \otimes \mathcal{L}$，从而
$$
[\mathcal{F} \otimes_{\mathcal{O}_X} \mathcal{L}]
-
[\mathcal{F}]
=
[i_*\mathcal{Q}_2] - [i_*\mathcal{Q}_1]
$$
在 $K_0(\textit{Coh}_{\leq k + 1}(X)/\textit{Coh}_{\leq k - 1}(X))$ 中成立。这已经说明
$[\mathcal{F} \otimes_{\mathcal{O}_X} \mathcal{L}] - [\mathcal{F}]$ 是 $B_k(X)$ 的元。再者，
\begin{eqnarray*}
[i_*\mathcal{Q}_2]_k - [i_*\mathcal{Q}_1]_k
& = &
i_*\left( [\mathcal{Q}_2]_k - [\mathcal{Q}_1]_k \right) \\
& = &
i_*\left(c_1(\mathcal{N}) \cap [\mathcal{G}]_{k + 1} \right) \\
& = &
c_1(\mathcal{L}) \cap i_*[\mathcal{G}]_{k + 1} \\
& = &
c_1(\mathcal{L}) \cap [\mathcal{F}]_{k + 1}
\end{eqnarray*}
这由上述结果及引理 \ref{chow-lemma-pushforward-cap-c1} 和
\ref{chow-lemma-cycle-push-sheaf} 得出。按定义，它与该元在
$B_k(X) \to \CH_k(X)$ 下的像一致。因此引理得证。
\end{proof}
% Generated mechanically from the frozen intersection.tex target.
% Do not edit this generated file; edit the bound locale source instead.

% OK, start here.
%

\setcounter{chapter}{42}
\chapter{相交理论}



\phantomsection
\label{intersection-section-phantom}



\section{引言}
\label{intersection-section-introduction}

\noindent
本章在代数闭域上的非奇异射影簇上，构造模有理等价的 Chow 群上的交积。
我们采用的工具是 Serre 的 Tor 公式
（见 \cite[Chapter V]{Serre_algebre_locale}）、化归到对角线以及移动引理。

\medskip\noindent
我们先回顾循环，并说明如何构造循环的固有推前与平坦拉回。随后引入循环的有理等价，
由此得到 Chow 群 $\CH_*(X)$。固有推前与平坦拉回均可通过有理等价下降，从而给出
Chow 群上的运算。这些内容见第
\ref{intersection-section-cycles}、
\ref{intersection-section-cycle-of-closed}、
\ref{intersection-section-cycle-of-coherent-sheaf}、
\ref{intersection-section-pushforward}、
\ref{intersection-section-flat-pullback}、
\ref{intersection-section-rational-equivalence}、
\ref{intersection-section-alternative}、
\ref{intersection-section-pushforward-and-rational-equivalence} 及
\ref{intersection-section-flat-pullback-and-rational-equivalence} 节。
证明主要引用 Chow 同调一章；其中已在泛链状诺特基上局部有限型概形的框架下证明这些结果，
见 Chow 同调，第 \ref{chow-section-setup} 节及以下各节。

\medskip\noindent
由于我们研究的是非奇异射影簇 $X$，两个不可约闭子簇之交 $V \cap W$ 的任一不可约分支
的维数至少为 $\dim(V) + \dim(W) - \dim(X)$。若对每个不可约分支 $Z$ 都取等号，
则称 $V$ 与 $W$ 正常相交。在此情形下，用公式
$$
e_Z = \sum\nolimits_i
(-1)^i
\text{length}_{\mathcal{O}_{X, Z}}
\text{Tor}_i^{\mathcal{O}_{X, Z}}(\mathcal{O}_{W, Z}, \mathcal{O}_{V, Z})
$$
定义交重数 $e_Z = e(X, V \cdot W, Z)$。
我们需要运用少量交换代数来证明，在简单情形下这些交重数符合直观；具体而言，有时
$$
e_Z = \text{length}_{\mathcal{O}_{X, Z}} \mathcal{O}_{V \cap W, Z},
$$
也就是说，只有 $\text{Tor}_0$ 有贡献。当 $V$ 与 $W$ 在 $Z$ 的泛点处均为
Cohen-Macaulay，或当 $W$ 在 $\mathcal{O}_{X, Z}$ 中由一个正则序列截出，
且该序列在 $\mathcal{O}_{V, Z}$ 上也为正则序列时，便是如此。不过，
例 \ref{intersection-example-naive-multiplicity-wrong} 表明一般情形确实需要高阶 Tor。
此外，这与 Samuel 重数之间还有联系。这些内容见第
\ref{intersection-section-intersect-properly}、
\ref{intersection-section-tor-formula}、
\ref{intersection-section-multiplicities}、
\ref{intersection-section-computing-intersection-multiplicities} 及
\ref{intersection-section-intersection-product} 节。

\medskip\noindent
化归到对角线是指：可以通过让 $V \times W$ 与 $X \times X$ 中的对角线相交，
来求 $V$ 与 $W$ 的交。在概形论意义的交这一层面上，这个看似平淡的说法是显然的；
它把一般一对子概形的相交问题化归为其中一个子概形局部由正则序列截出的情形。
我们依照 Serre 的方法，由此得到交重数的正性。此外，化归到对角线还导出交重数的
可加性、结合律以及投影公式。这些内容见第
\ref{intersection-section-exterior-product}、
\ref{intersection-section-reduction-diagonal}、
\ref{intersection-section-associative}、
\ref{intersection-section-flat-pullback-and-intersection-products} 及
\ref{intersection-section-projection-formula-flat} 节。

\medskip\noindent
最后讨论移动引理及其应用。移动引理由两部分组成。第一部分说，给定闭子簇
$$
Z \subset X \subset \mathbf{P}^N
$$
且 $X$ 非奇异，可以找到与 $X$ 正常相交的子簇 $C \subset \mathbf{P}^N$，使得
$$
C \cdot X = [Z] + \sum m_j [Z_j]
$$
并且其余分支 $Z_j$ 比 $Z$ “更一般”。第二部分说，可以沿一条有理曲线移动
$C \subset \mathbf{P}^N$，使其移到相对于任意给定子簇列表处于一般位置的子簇。
两部分合在一起表明，只须对 $X$ 上正常相交的循环定义交积；上文已经完成这一步。
当然，只有证明这些乘积能通过有理等价下降，才能得到 $\CH_*(X)$ 上的交积；
正文将证明这一点。这些结果及若干应用见第
\ref{intersection-section-projection}、
\ref{intersection-section-moving-lemma}、
\ref{intersection-section-intersections-and-rational-equivalence}、
\ref{intersection-section-chow-rings}、
\ref{intersection-section-general-pullback} 及
\ref{intersection-section-pullback-cycles} 节。



\section{约定}
\label{intersection-section-conventions}

\noindent
固定任意特征的代数闭基域 $\mathbf{C}$。所有概形与簇均定义在 $\mathbf{C}$ 上，
所有态射也都是 $\mathbf{C}$ 上的态射。若簇 $X$ 是正则概形，则称 $X$
{\it 非奇异}（见《性质》，定义 \ref{properties-definition-regular}）。
在当前情形下，这等价于态射 $X \to \Spec(\mathbf{C})$ 光滑
（见《簇》，引理 \ref{varieties-lemma-geometrically-regular-smooth}）。


\section{循环}
\label{intersection-section-cycles}

\noindent
设 $X$ 为簇。$X$ 的一个{\it 闭子簇}是整闭子概形 $Z \subset X$。
$X$ 上的一个{\it $k$-循环}是有限形式和 $\sum n_i [Z_i]$，其中每个 $Z_i$
都是维数为 $k$ 的闭子簇。每当以 $\alpha = \sum n_i[Z_i]$ 表示一个
$k$-循环时，总假设各 $Z_i$ 两两不同，且对所有 $i$ 都有 $n_i \not = 0$。
此时，$\alpha$ 的{\it 支撑}是维数为 $k$ 的闭子集
$$
\text{Supp}(\alpha) = \bigcup Z_i \subset X
$$
$k$-循环群记为 $Z_k(X)$。见《Chow 同调》，第 \ref{chow-section-cycles} 节。


\section{闭子概形所对应的循环}
\label{intersection-section-cycle-of-closed}

\noindent
设 $X$ 为簇，$Z \subset X$ 为满足 $\dim(Z) \leq k$ 的闭子概形。
设 $Z_i$ 为 $Z$ 中维数为 $k$ 的不可约分支，并定义 $n_i$ 为
{\it $Z_i$ 在 $Z$ 中的重数}：
$$
n_i = \text{length}_{\mathcal{O}_{X, Z_i}} \mathcal{O}_{Z, Z_i}
$$
这里 $\mathcal{O}_{X, Z_i}$ 与 $\mathcal{O}_{Z, Z_i}$ 分别表示
\ $X$ 与\ $Z$ 在 $Z_i$ 的泛点处的局部环。把如下 $k$-循环定义为
$Z$ 所对应的 $k$-循环：
$$
[Z]_k = \sum n_i [Z_i].
$$
见《Chow 同调》，第 \ref{chow-section-cycle-of-closed-subscheme} 节。


\section{凝聚层所对应的循环}
\label{intersection-section-cycle-of-coherent-sheaf}

\noindent
设 $X$ 为簇，$\mathcal{F}$ 为满足
$\dim(\text{Supp}(\mathcal{F})) \leq k$ 的凝聚 $\mathcal{O}_X$-模。
设 $Z_i$ 为 $\text{Supp}(\mathcal{F})$ 中维数为 $k$ 的不可约分支，并定义
$n_i$ 为{\it $Z_i$ 在 $\mathcal{F}$ 中的重数}：
$$
n_i = \text{length}_{\mathcal{O}_{X, Z_i}} \mathcal{F}_{\xi_i}
$$
这里 $\mathcal{O}_{X, Z_i}$ 是 $X$ 在 $Z_i$ 的泛点 $\xi_i$ 处的局部环，
$\mathcal{F}_{\xi_i}$ 是 $\mathcal{F}$ 在该点的茎。把如下 $k$-循环定义为
$\mathcal{F}$ 所对应的 $k$-循环：
$$
[\mathcal{F}]_k = \sum n_i [Z_i].
$$
见《Chow 同调》，第 \ref{chow-section-cycle-of-coherent-sheaf} 节。
注意，若 $Z \subset X$ 是满足 $\dim(Z) \leq k$ 的闭子概形，则按定义有
$[Z]_k = [\mathcal{O}_Z]_k$。


\section{固有推前}
\label{intersection-section-pushforward}

\noindent
设 $f : X \to Y$ 为簇之间的固有态射，$Z \subset X$ 为 $k$ 维闭子簇。
若 $\dim(f(Z)) < k$，定义 $f_*[Z]$ 为 $0$；若 $\dim(f(Z)) = k$，则定义
其为 $d \cdot [f(Z)]$，其中
$$
d = [\mathbf{C}(Z) : \mathbf{C}(f(Z))] = \deg(Z/f(Z))
$$
是支配态射 $Z \to f(Z)$ 的次数；见《态射》，定义
\ref{morphisms-definition-degree}。设 $\alpha = \sum n_i [Z_i]$
为 $X$ 上的 $k$-循环。定义 $\alpha$ 的{\it 推前}为
$f_* \alpha = \sum n_i f_*[Z_i]$，其中每个 $f_*[Z_i]$ 如上定义。
这给出群同态
$$
f_* : Z_k(X) \longrightarrow Z_k(Y)
$$
见《Chow 同调》，第 \ref{chow-section-proper-pushforward} 节。

\begin{lemma}
\label{intersection-lemma-push-coherent}
\begin{reference}
见 \cite[Chapter V]{Serre_algebre_locale}。
\end{reference}
设 $f : X \to Y$ 为簇之间的固有态射。若 $\mathcal{F}$ 为满足
$\dim(\text{Supp}(\mathcal{F})) \leq k$ 的凝聚层，则
$f_*[\mathcal{F}]_k = [f_*\mathcal{F}]_k$。特别地，若
$Z \subset X$ 为维数 $\leq k$ 的闭子概形，则
$f_*[Z]_k = [f_*\mathcal{O}_Z]_k$。
\end{lemma}

\begin{proof}
见《Chow 同调》，引理 \ref{chow-lemma-cycle-push-sheaf}。
\end{proof}

\begin{lemma}
\label{intersection-lemma-compose-pushforward}
设 $f : X \to Y$ 与 $g : Y \to Z$ 为簇之间的固有态射。
则作为映射 $Z_k(X) \to Z_k(Z)$，有 $g_* \circ f_* = (g \circ f)_*$。
\end{lemma}

\begin{proof}
这是《Chow 同调》引理 \ref{chow-lemma-compose-pushforward} 的特例。
\end{proof}



\section{平坦拉回}
\label{intersection-section-flat-pullback}

\noindent
设 $f : X \to Y$ 为簇之间的平坦态射。由《态射》引理
\ref{morphisms-lemma-dimension-fibre-at-a-point-additive}，
$f$ 的每个纤维都具有维数 $r = \dim(X) - \dim(Y)$\footnote{反过来，
若 $f : X \to Y$ 是簇之间的支配态射，$X$ 为 Cohen-Macaulay，$Y$ 非奇异，
且所有纤维都具有相同维数 $r$，则 $f$ 平坦。这由《代数》引理
\ref{algebra-lemma-CM-over-regular-flat} 与《簇》引理
\ref{varieties-lemma-dimension-fibres-locally-algebraic} 得出；
后者给出 $\dim(X) = \dim(Y) + r$。}。
设 $Z \subset Y$ 为 $k$ 维闭子簇。定义 $f^*[Z]$ 为概形论逆像所对应的
$(k + r)$-循环，即 $f^*[Z] = [f^{-1}(Z)]_{k + r}$。设
$\alpha = \sum n_i [Z_i]$ 为 $Y$ 上的 $k$-循环。定义 $\alpha$ 的
{\it 拉回}为和 $f^* \alpha = \sum n_i f^*[Z_i]$，其中每个 $f^*[Z_i]$
如上定义。这给出群同态
$$
f^* : Z_k(Y) \longrightarrow Z_{k + r}(X)
$$
见《Chow 同调》，第 \ref{chow-section-flat-pullback} 节。

\begin{lemma}
\label{intersection-lemma-pullback}
设 $f : X \to Y$ 为簇之间的平坦态射，并置 $r = \dim(X) - \dim(Y)$。
则 $f^*[\mathcal{F}]_k = [f^*\mathcal{F}]_{k + r}$，其中
$\mathcal{F}$ 是 $Y$ 上的凝聚层，且 $\mathcal{F}$ 的支撑维数至多为 $k$。
\end{lemma}

\begin{proof}
见《Chow 同调》，引理 \ref{chow-lemma-pullback-coherent}。
\end{proof}

\begin{lemma}
\label{intersection-lemma-compose-flat-pullback}
设 $f : X \to Y$ 与 $g : Y \to Z$ 为簇之间的平坦态射。
则 $g \circ f$ 平坦，并且作为映射
$Z_k(Z) \to Z_{k + \dim(X) - \dim(Z)}(X)$，有
$f^* \circ g^* = (g \circ f)^*$。
\end{lemma}

\begin{proof}
这是《Chow 同调》引理 \ref{chow-lemma-compose-flat-pullback} 的特例。
\end{proof}


\section{有理等价}
\label{intersection-section-rational-equivalence}

\noindent
下面给出的有理等价定义，乍看之下可能不同于读者熟悉的定义，也不同于
\cite[Chapter I]{F} 或《Chow 同调》第
\ref{chow-section-rational-equivalence} 节中的表述。不过，我们将在第
\ref{intersection-section-alternative} 节证明这两个概念一致。

\medskip\noindent
设 $X$ 为簇，$W \subset X \times \mathbf{P}^1$ 为维数 $k + 1$
的闭子簇，$a, b$ 为 $\mathbf{P}^1$ 上两个不同的闭点。假设
$X \times a$、$X \times b$ 与 $W$ 正常相交，即
$$
\dim (W \cap X \times a) \leq k,\quad
\dim (W \cap X \times b) \leq k.
$$
只要 $W \to \mathbf{P}^1$ 是支配的，或者 $W$ 包含在投影位于某个不同于
$a$、$b$ 的闭点之上的纤维中，这一条件便成立（后一种是将被舍弃的无趣情形）。
此时，态射 $W \to \mathbf{P}^1$ 的概形论纤维 $W_a$ 等于
$X \times \mathbf{P}^1$ 中的概形论交 $W \cap X \times a$。
把 $X \times a$ 与 $X \times b$ 认同于 $X$，便可把纤维 $W_a$ 与 $W_b$
看成 $X$ 中维数 $\leq k$ 的闭子概形\footnote{有时把 $W_a$ 看成
$X \times \mathbf{P}^1$ 的闭子概形，有时把它看成 $X$ 的闭子概形；
所取观点应总能由上下文辨明。}。有理等价的一个基本例子是
$$
[W_a]_k \sim_{rat} [W_b]_k
$$
在给定 $W$ 时，循环 $[W_a]_k$ 与 $[W_b]_k$ 在实际中很容易计算，
因为它们由同一个 Cartier 除子的正常相交得到（见第
\ref{intersection-section-intersection-product} 节）。由于 $\mathbf{P}^1$ 的自同构群
是 $2$-重传递的，可以把闭点对 $a, b$ 移到任意希望的位置。传统选择是
$a = 0$ 与 $b = \infty$。

\medskip\noindent
更一般地，设 $\alpha = \sum n_i [W_i]$ 为
$X \times \mathbf{P}^1$ 上的 $(k + 1)$-循环，并在
$\mathbf{P}^1$ 上取不同闭点对 $a_i, b_i$。假设
$X \times a_i$、$X \times b_i$ 与 $W_i$ 正常相交；换言之，每组
$W_i, a_i, b_i$ 都满足上面的条件。所谓{\it 有理等价于零的循环}，
是指任意形如
$$
\sum n_i([W_{i, a_i}]_k - [W_{i, b_i}]_k).
$$
的循环。这确实是一个 $k$-循环。有理等价于零的 $k$-循环全体构成
$k$-循环群的一个加法子群。若 $\alpha - \alpha'$ 有理等价于零，
则称两个 $k$-循环{\it 有理等价}，记作
$\alpha \sim_{rat} \alpha'$。

\medskip\noindent
定义
$$
\CH_k(X) = Z_k(X)/ \sim_{rat}
$$
称其为{\it $X$ 上 $k$-循环的 Chow 群}。引理
\ref{intersection-lemma-rational-equivalence} 将表明，这与《Chow 同调》定义
\ref{chow-definition-rational-equivalence} 中的 Chow 群一致。


\section{有理等价与有理函数}
\label{intersection-section-alternative}

\noindent
设 $X$ 为簇，$W \subset X$ 为维数 $k + 1$ 的子簇，
$f \in \mathbf{C}(W)^*$ 为 $W$ 上的非零有理函数。对任意维数为 $k$ 的
子簇 $Z \subset W$，可定义 $f$ 在 $Z$ 上的消失阶
$\text{ord}_{W, Z}(f)$。若 $f$ 是局部环 $\mathcal{O}_{W, Z}$ 的元，则
$$
\text{ord}_{W, Z}(f) =
\text{length}_{\mathcal{O}_{X, Z}} \mathcal{O}_{W, Z}/f\mathcal{O}_{W, Z}
$$
其中\ $\mathcal{O}_{X, Z}$ 与\ $\mathcal{O}_{W, Z}$ 分别是 $X$ 与 $W$ 在 $Z$
的泛点处的局部环。一般情形下，由乘法性将此定义延拓。与 $f$
相关联的{\it 主除子}为
$$
\text{div}_W(f) = \sum \text{ord}_{W, Z}(f)[Z]
$$
它属于 $Z_k(W)$。因为 $W \subset X$ 是闭子簇，可以把
$\text{div}_W(f)$ 看作 $X$ 上的循环。
见《Chow 同调》第 \ref{chow-section-principal-divisors} 节。

\begin{lemma}
\label{intersection-lemma-rational-equivalence}
设 $X$ 为簇，$W \subset X$ 为维数 $k + 1$ 的子簇，
$f \in \mathbf{C}(W)^*$ 为 $W$ 上的非零有理函数。则
$\text{div}_W(f)$ 在 $X$ 上有理等价于零。反之，这些主除子生成
$X$ 上有理等价于零的循环所构成的阿贝尔群。
\end{lemma}

\begin{proof}
第一个断言由《Chow 同调》引理
\ref{chow-lemma-rational-function} 得出。更确切地，设
$W' \subset X \times \mathbf{P}^1$ 为 $f$ 的图像的闭包。则
$\text{div}_W(f) = [W'_0]_k - [W'_\infty]$ 在
$Z_k(W) \subset Z_k(X)$ 中成立；见《Chow 同调》引理
\ref{chow-lemma-rational-function} 的第 (6) 部分。

\medskip\noindent
对于第二个断言，设 $W' \subset X \times \mathbf{P}^1$ 是一个维数为
$k + 1$ 且支配 $\mathbf{P}^1$ 的闭子簇。我们将证明
$[W'_0]_k - [W'_\infty]_k$ 是主除子，从而完成证明。设
$W \subset X$ 为 $W'$ 在到 $X$ 的投影下的像。则
$W \subset X$ 是闭子簇，$W' \to W$ 是固有支配的，其纤维数为
$0$ 或 $1$。若 $\dim(W) = k$，则 $W' = W \times \mathbf{P}^1$，且
$[W'_0]_k - [W'_\infty]_k = [W] - [W] = 0$。
若 $\dim(W) = k + 1$，则
$W' \to W$ 在一般点上是有限的\footnote{若 $W' \to W$ 是双有理的，
则结论由《Chow 同调》引理
\ref{chow-lemma-rational-function} 得出。我们需要证明，即使 $W' \to W$
的次数 $> 1$，基本有理等价
$[W'_0]_k \sim_{rat} [W'_\infty]_k$ 仍来自 $X$ 的某个子簇上的主除子。}。
以 $f$ 表示投影 $W' \to \mathbf{P}^1$，并把它视为
$\mathbf{C}(W')^*$ 的元。令 $g = \text{Nm}(f) \in \mathbf{C}(W)^*$ 为范数。由
《Chow 同调》引理 \ref{chow-lemma-proper-pushforward-alteration}，有
$$
\text{div}_W(g) = \text{pr}_{X, *}\text{div}_{W'}(f)
$$
再由《Chow 同调》引理 \ref{chow-lemma-rational-function}，
$\text{div}_{W'}(f) = [W'_0]_k - [W'_\infty]_k$，证明完毕。
\end{proof}


\section{固有推前与有理等价}
\label{intersection-section-pushforward-and-rational-equivalence}

\noindent
设 $f : X \to Y$ 为簇之间的固有态射，$\alpha \sim_{rat} 0$ 为
$X$ 上有理等价于 $0$ 的 $k$-循环。则 $\alpha$ 的{推前}也有理等价于零：
$f_* \alpha \sim_{rat} 0$。见 \cite{F} 第 I 章或《Chow 同调》引理
\ref{chow-lemma-proper-pushforward-rational-equivalence}。

\medskip\noindent
因此得到交换图
$$
\xymatrix{
Z_k(X) \ar[r] \ar[d]_{f_*} & \CH_k(X) \ar[d]^{f_*} \\
Z_k(Y) \ar[r] & \CH_k(Y)
}
$$
其中各项均为 $k$-循环群。


\section{平坦拉回与有理等价}
\label{intersection-section-flat-pullback-and-rational-equivalence}

\noindent
设 $f : X \to Y$ 为簇之间的平坦态射，并置 $r = \dim(X) - \dim(Y)$。
设 $\alpha \sim_{rat} 0$ 为 $Y$ 上有理等价于 $0$ 的 $k$-循环。则
$\alpha$ 的拉回也有理等价于零：$f^* \alpha \sim_{rat} 0$。
见 \cite{F} 第 I 章或《Chow 同调》引理
\ref{chow-lemma-flat-pullback-rational-equivalence}。

\medskip\noindent
因此得到交换图
$$
\xymatrix{
Z_{k + r}(X) \ar[r] & \CH_{k + r}(X) \\
Z_k(Y) \ar[r] \ar[u]^{f^*} & \CH_k(Y) \ar[u]_{f^*}
}
$$
其中各项均为 $k$-循环群。


\section{开子簇的短正合列}
\label{intersection-section-ses}

\noindent
设 $X$ 为簇，$U \subset X$ 为开子簇。将
$X \setminus U = \bigcup Z_i$ 分解为不可约分支\footnote{本章只考虑簇的
Chow 群，因而不能直接取 $Z_k(X \setminus U)$ 与
$\CH_k(X \setminus U)$，故采用各簇 $Z_i$ 来表述。}。则对每个
$k \geq 0$ 都存在交换图
$$
\xymatrix{
\bigoplus Z_k(Z_i) \ar[r] \ar[d] &
Z_k(X) \ar[r] \ar[d] &
Z_k(U) \ar[d] \ar[r] &
0 \\
\bigoplus \CH_k(Z_i) \ar[r] &
\CH_k(X) \ar[r] &
\CH_k(U) \ar[r] &
0
}
$$
它的两行均正合。竖直箭头是典范商映射。左侧水平箭头由沿闭浸入
$Z_i \to X$ 的固有推前给出；右侧水平箭头由沿开浸入
$j : U \to X$ 的平坦拉回给出。既然已知这些映射可以经过有理等价分解，
方块的交换性随即得出。上行正合，是因为 $X$ 的每个子簇要么包含在某个
$Z_i$ 中，要么它与 $U$ 的交不可约。下行正合，是因为 $U$ 上的每个主除子
$\text{div}_W(f)$ 都是 $X$ 上某个主除子的限制。更确切地，若
$W \subset U$ 是维数为 $(k + 1)$ 的闭子簇，且
$f \in \mathbf{C}(W)^*$，以 $\overline{W}$ 表示 $W$ 在 $X$ 中的闭包。则
$W \subset \overline{W}$ 是开浸入，故
$\mathbf{C}(W) = \mathbf{C}(\overline{W})$，可把 $f$ 看作 $\overline{W}$ 上的非常值有理函数。
于是显然有
$$
j^*\text{div}_{\overline{W}}(f) = \text{div}_W(f)
$$
该等式在 $Z_k(X)$ 中成立。由此容易得到下行的正合性。详情见《Chow 同调》引理
\ref{chow-lemma-restrict-to-open}。


\section{正常相交}
\label{intersection-section-intersect-properly}

\noindent
先给出几个用于估计维数的引理。

\begin{lemma}
\label{intersection-lemma-dimension-product-varieties}
设 $X$ 与 $Y$ 为簇。则 $X \times Y$ 也是簇，且
$\dim(X \times Y) = \dim(X) + \dim(Y)$。
\end{lemma}

\begin{proof}
由《簇》引理 \ref{varieties-lemma-product-varieties}，概形
$X \times Y = X \times_{\Spec(\mathbf{C})} Y$ 是簇。维数断言是《簇》引理
\ref{varieties-lemma-dimension-product-locally-algebraic}。
\end{proof}

\noindent
回忆，概形的正则浸入 $i : X \to Y$ 是一个闭浸入，其对应理想层在局部
由正则序列生成；见《除子》第 \ref{divisors-section-regular-immersions} 节。
此外，余法层 $\mathcal{C}_{X/Y}$ 有限局部自由，其秩等于该正则序列的长度。
若 $\mathcal{C}_{X/Y}$ 是秩为 $c$ 的局部自由层，就称 $i$ 是
{\it 余维数为 $c$ 的正则浸入}。

\medskip\noindent
更一般地，回忆《态射进阶》第 \ref{more-morphisms-section-lci} 节中的定义：
若可以用开集 $U$ 覆盖 $X$，使 $f|_U$ 都能分解为
$$
\xymatrix{
U \ar[rr]_i \ar[rd] & & \mathbf{A}^n_Y \ar[ld] \\
& Y
}
$$
其中 $i$ 是 Koszul 正则浸入，则称 $f : X \to Y$ 为局部完全交态射。
当 $Y$ 局部 Noether 时，这等价于要求 $i$ 为正则浸入；见《除子》引理
\ref{divisors-lemma-regular-immersion-noetherian}。若对任意上述分解，闭浸入 $i$
的余法层秩都为 $n - r$，就称 $f$ 是{\it 相对维数为 $r$ 的局部完全交态射}。
换言之，$i$ 是余维数为 $n - r$ 的 Koszul-正则浸入；在 Noether 情形下，
这就是余维数为 $n - r$ 的正则浸入。

\begin{lemma}
\label{intersection-lemma-pullback-by-regular-immersion}
设 $f : X \to Y$ 为簇之间的态射。
\begin{enumerate}
\item 若 $Z \subset Y$ 是维数为 $d$ 的子簇，且 $f$ 是余维数为 $c$ 的正则浸入，则
$f^{-1}(Z)$ 的每个不可约分支的维数都 $\geq d - c$。
\item 若 $Z \subset Y$ 是维数为 $d$ 的子簇，且
$f$ 是相对维数为 $r$ 的局部完全交态射，则 $f^{-1}(Z)$ 的每个不可约分支
的维数都 $\geq d + r$。
\end{enumerate}
\end{lemma}

\begin{proof}
证明 (1)。可以局部地工作，故可设
$Y = \Spec(A)$ 且 $X = V(f_1, \ldots, f_c)$，其中 $f_1, \ldots, f_c$
是 $A$ 中的正则序列。若 $Z = \Spec(A/\mathfrak p)$，则
$f^{-1}(Z) = \Spec(A/\mathfrak p + (f_1, \ldots, f_c))$。
若 $V$ 是 $f^{-1}(Z)$ 的一个不可约分支，则可选取一个不属于
$f^{-1}(Z)$ 其他任何不可约分支的闭点 $v \in V$。于是
$$
\dim(Z) = \dim \mathcal{O}_{Z, v}
\quad\text{且}\quad
\dim(V) = \dim \mathcal{O}_{V, v} = \dim \mathcal{O}_{Z, v}/(f_1, \ldots, f_c)
$$
第一个等式例如由《代数》引理
\ref{algebra-lemma-dimension-prime-polynomial-ring} 得出，第二个等式则由闭点的选取得出。
最后，由极大理想中的一个元作商至多使维数减少 $1$，即得结论；见《代数》引理
\ref{algebra-lemma-one-equation}。

\medskip\noindent
证明 (2)。按局部完全交的定义选取一个分解，然后应用 (1)。略去若干细节。
\end{proof}

\begin{lemma}
\label{intersection-lemma-diagonal-regular-immersion}
设 $X$ 为非奇异簇。则对角态射
$\Delta : X \to X \times X$ 是余维数为 $\dim(X)$ 的正则浸入。
\end{lemma}

\begin{proof}
事实上，非奇异射影簇之间的任意闭浸入都是正则浸入；见《除子》引理
\ref{divisors-lemma-immersion-smooth-into-smooth-regular-immersion}。
\end{proof}

\noindent
下一引理说明如何进行化归到对角线。

\begin{lemma}
\label{intersection-lemma-intersect-in-smooth}
设 $X$ 为非奇异簇，$W,V \subset X$ 为闭子簇，且
$\dim(W) = s$、$\dim(V) = r$。则 $V \cap W$ 的每个不可约分支 $Z$
的维数都 $\geq r + s - \dim(X)$。
\end{lemma}

\begin{proof}
由于在概形论意义下 $V \cap W = \Delta^{-1}(V \times W)$，结论由引理
\ref{intersection-lemma-diagonal-regular-immersion} 与
\ref{intersection-lemma-pullback-by-regular-immersion} 得出。
\end{proof}

\noindent
这一引理引出下列定义。

\begin{definition}
\label{intersection-definition-proper-intersection}
设 $X$ 为非奇异簇。
\begin{enumerate}
\item 设 $W,V \subset X$ 为闭子簇，且
$\dim(W) = s$、$\dim(V) = r$。若
$\dim(V \cap W) \leq r + s - \dim(X)$，就称 $W$ 与 $V$ {\it 正常相交}。
\item 设 $\alpha = \sum n_i [W_i]$ 为 $s$-循环，
$\beta = \sum_j m_j [V_j]$ 为 $X$ 上的 $r$-循环。若对所有 $i$ 与 $j$，
$W_i$ 与 $V_j$ 都正常相交，就称 $\alpha$ 与 $\beta$ {\it 正常相交}。
\end{enumerate}
\end{definition}


\section{用 Tor 公式定义相交重数}
\label{intersection-section-tor-formula}

\noindent
我们将频繁使用如下基本事实。给定带环空间 $(X, \mathcal{O}_X)$ 上的模层
$\mathcal{F}$、$\mathcal{G}$ 以及一点 $x \in X$，有
$$
\text{Tor}_p^{\mathcal{O}_X}(\mathcal{F}, \mathcal{G})_x =
\text{Tor}_p^{\mathcal{O}_{X, x}}(\mathcal{F}_x, \mathcal{G}_x)
$$
它们是 $\mathcal{O}_{X, x}$-模。这一事实可以从《上同调》第
\ref{cohomology-section-flat} 节中对导出张量积的构造以多种方式看出；例如它由
《上同调》引理 \ref{cohomology-lemma-check-K-flat-stalks} 得出。
此外，若 $X$ 是概形，而 $\mathcal{F}$ 与 $\mathcal{G}$ 拟凝聚，则模
$\text{Tor}_p^{\mathcal{O}_X}(\mathcal{F}, \mathcal{G})$ 也拟凝聚；见《概形的导出范畴》引理
\ref{perfect-lemma-quasi-coherence-tensor-product}。对我们的用途更重要的是下列结果。

\begin{lemma}
\label{intersection-lemma-tensor-coherent}
设 $X$ 为局部 Noether 概形。
\begin{enumerate}
\item 若 $\mathcal{F}$ 与 $\mathcal{G}$ 是凝聚 $\mathcal{O}_X$-模，则
$\text{Tor}_p^{\mathcal{O}_X}(\mathcal{F}, \mathcal{G})$ 也是凝聚的。
\item 若 $L$ 与 $K$ 属于 $D^-_{\textit{Coh}}(\mathcal{O}_X)$，则
$L \otimes_{\mathcal{O}_X}^\mathbf{L} K$ 也属于该范畴。
\end{enumerate}
\end{lemma}

\begin{proof}
我们用较为初等的方法证明 (1)，并用先前建立的一般理论证明 (2)。

\medskip\noindent
证明 (1)。由于 $\text{Tor}$ 的形成与局部化可交换，可以设 $X$ 为仿射概形。
于是对某个 Noether 环 $A$，有 $X = \Spec(A)$，且 $\mathcal{F}$、$\mathcal{G}$
分别对应于有限 $A$-模 $M$ 与 $N$；见《概形的上同调》引理
\ref{coherent-lemma-coherent-Noetherian}。由《概形的导出范畴》引理
\ref{perfect-lemma-quasi-coherence-tensor-product}，可以先在 $A$ 上计算 $M$ 与 $N$
的 $\text{Tor}$，再取其所关联的 $\mathcal{O}_X$-模，以此计算 $\text{Tor}$。由《代数》引理
\ref{algebra-lemma-tor-noetherian}，模 $\text{Tor}_p^A(M, N)$ 都是有限的，因而结论成立。

\medskip\noindent
由《概形的导出范畴》引理
\ref{perfect-lemma-identify-pseudo-coherent-noetherian}，所设条件等价于要求 $L$ 与 $K$
都（局部地）拟凝聚。再由《上同调》引理
\ref{cohomology-lemma-tensor-pseudo-coherent}，
$L \otimes_{\mathcal{O}_X}^\mathbf{L} K$ 也拟凝聚。
\end{proof}

\begin{lemma}
\label{intersection-lemma-compute-tor-nonsingular}
设 $X$ 为非奇异簇，$\mathcal{F}$、$\mathcal{G}$ 为凝聚 $\mathcal{O}_X$-模。则
$\mathcal{O}_X$-模 $\text{Tor}_p^{\mathcal{O}_X}(\mathcal{F}, \mathcal{G})$
是凝聚的，它在 $x$ 处的茎等于
$\text{Tor}_p^{\mathcal{O}_{X, x}}(\mathcal{F}_x, \mathcal{G}_x)$，其支撑包含于
$\text{Supp}(\mathcal{F}) \cap \text{Supp}(\mathcal{G})$，并且它只可能在
$p \in \{0, \ldots, \dim(X)\}$ 时非零。
\end{lemma}

\begin{proof}
上文已讨论茎的断言，它也蕴含了支撑条件。由引理
\ref{intersection-lemma-tensor-coherent}，各 $\text{Tor}$ 都是凝聚的。负次 $\text{Tor}$
的消失直接来自构造。对 $p > \dim(X)$，$\text{Tor}_p$ 的消失可见如下：
局部环 $\mathcal{O}_{X, x}$ 因 $X$ 非奇异而正则，且其维数
$\leq \dim(X)$；见《代数》引理
\ref{algebra-lemma-dimension-prime-polynomial-ring}。因此 $\mathcal{O}_{X, x}$ 的有限整体维数
$\leq \dim(X)$；见《代数》引理
\ref{algebra-lemma-finite-gl-dim-finite-dim-regular}。这蕴含模的 $\text{Tor}$-群在超过此维数后消失；
见《代数进阶》引理 \ref{more-algebra-lemma-finite-gl-dim-tor-dimension}。
\end{proof}

\noindent
设 $X$ 为非奇异簇，$W, V \subset X$ 为闭子簇，且
$\dim(W) = s$、$\dim(V) = r$。假设 $V$ 与 $W$ 正常相交。
此时引理 \ref{intersection-lemma-intersect-in-smooth} 告诉我们，$V \cap W$ 的所有不可约分支
的维数都等于 $r + s - \dim(X)$。层
$\text{Tor}_j^{\mathcal{O}_X}(\mathcal{O}_W, \mathcal{O}_V)$ 是凝聚的，支撑于 $V \cap W$，
并且当 $j < 0$ 或 $j > \dim(X)$ 时为零；见引理
\ref{intersection-lemma-compute-tor-nonsingular}。定义{\it 交积}为
$$
W \cdot V = \sum\nolimits_i (-1)^i
[\text{Tor}_i^{\mathcal{O}_X}(\mathcal{O}_W, \mathcal{O}_V)]_{r + s - \dim(X)}.
$$
强调一下，这一定义只因我们假设 $V$ 与 $W$ 正常相交才有意义。
为了在一般情形下定义交积，这一事实将迫使我们使用移动引理。

\medskip\noindent
在此记号下，循环 $V \cdot W$ 是交 $V \cap W$ 的不可约分支 $Z$
的形式线性组合 $\sum e_Z Z$。整数 $e_Z$ 称为{\it 相交重数}：
$$
e_Z = e(X, V \cdot W, Z) =
\sum\nolimits_i
(-1)^i
\text{length}_{\mathcal{O}_{X, Z}}
\text{Tor}_i^{\mathcal{O}_{X, Z}}(\mathcal{O}_{W, Z}, \mathcal{O}_{V, Z})
$$
其中\ $\mathcal{O}_{X, Z}$、\ $\mathcal{O}_{W, Z}$ 与\ $\mathcal{O}_{V, Z}$ 分别表示 $X$、$W$ 与 $V$
在\ $Z$ 的泛点处的局部环。正如后文将看到的，这些 $\text{Tor}$ 的长度交错和
具有许多良好性质。

\medskip\noindent
在横截相交的情形下，相交数为 $1$。

\begin{lemma}
\label{intersection-lemma-transversal}
设 $X$ 为非奇异簇，$V, W \subset X$ 为正常相交的闭子簇。设 $Z$ 为
$V \cap W$ 的一个不可约分支，并假设闭子概形 $V \cap W$ 中 $Z$ 的重数
（按第 \ref{intersection-section-cycle-of-closed} 节的意义）为 $1$。则
$e(X, V \cdot W, Z) = 1$，且 $V$ 与 $W$ 在 $Z$ 的一般点处光滑。
\end{lemma}

\begin{proof}
设 $(A, \mathfrak m, \kappa) =
(\mathcal{O}_{X, \xi}, \mathfrak m_\xi, \kappa(\xi))$，其中 $\xi \in Z$ 是泛点。则
$\dim(A) = \dim(X) - \dim(Z)$；见《簇》引理
\ref{varieties-lemma-dimension-locally-algebraic}。设 $I, J \subset A$ 分别截出 $V$ 与 $W$
在 $\Spec(A)$ 中的迹。置 $\overline{I} = I + \mathfrak m^2/\mathfrak m^2$。则
$\dim_\kappa \overline{I} \leq \dim(X) - \dim(V)$，且等号成立当且仅当
$A/I$ 正则。这由上述引理与正则环的定义得出；见《代数》定义
\ref{algebra-definition-regular-local} 及其前的讨论。对 $\overline{J}$ 亦然。
若重数为 $1$，则 $\text{length}_A(A/I + J) = 1$，因而
$I + J = \mathfrak m$，从而 $\overline{I} + \overline{J} = \mathfrak m/\mathfrak m^2$。
于是，由于交是正常的，上述各不等式处处取等号。因此可找到
$f_1, \ldots, f_a \in I$ 与 $g_1, \ldots g_b \in J$，使
$\overline{f}_1, \ldots, \overline{g}_b$ 是 $\mathfrak m/\mathfrak m^2$ 的一组基。于是
$f_1, \ldots, g_b$ 是一组正则参数系，也是正则序列；见《代数》引理
\ref{algebra-lemma-regular-ring-CM}。同一引理说明
$A/(f_1, \ldots, f_a)$ 是维数为 $\dim(X) - \dim(V)$ 的正则局部环，故
$A/(f_1, \ldots, f_a) \to A/I$ 是同构。事实上，若核非零，则 $A/I$
的维数严格更小；见《代数》引理
\ref{algebra-lemma-regular-domain} 与 \ref{algebra-lemma-one-equation}。
利用对称性，得到 $I = (f_1, \ldots, f_a)$ 与 $J = (g_1, \ldots, g_b)$。
因此，关于 $f_1, \ldots, f_a$ 的 Koszul 复形
$K_\bullet(A, f_1, \ldots, f_a)$ 是 $A/I$ 的一个解析；见《代数进阶》引理
\ref{more-algebra-lemma-regular-koszul-regular}。从而
\begin{align*}
\text{Tor}_p^A(A/I, A/J)
& =
H_p(K_\bullet(A, f_1, \ldots, f_a) \otimes_A A/J) \\
& =
H_p(K_\bullet(A/J, f_1 \bmod J, \ldots, f_a \bmod J))
\end{align*}
上文已证明 $f_1 \bmod J, \ldots, f_a \bmod J$ 是正则局部环 $A/J$
中的一组正则参数系，故只有一个非零同调群，即
$H_0 = A/(I + J) = \kappa$。证明完毕。
\end{proof}

\begin{example}
\label{intersection-example-naive-multiplicity-wrong}
本例表明，在上述相交重数公式中有必要使用高次 tors。
设 $X$ 为维数 $4$ 的非奇异簇，$p \in X$ 为闭点，
$V, W \subset X$ 为 $X$ 中的闭子簇。假设存在同构
$$
\mathcal{O}_{X, p}^\wedge \cong \mathbf{C}[[x, y, z, w]]
$$
使 $V$ 的理想为 $(xz, xw, yz, yw)$，而 $W$ 的理想为
$(x - z, y - w)$。计算表明
$$
\text{length}\ \mathbf{C}[[x, y, z, w]]/
(xz, xw, yz, yw, x - z, y - w) = 3
$$
另一方面，重数 $e(X, V \cdot W, p) = 2$。这可从下述事实看出：在 $p$
的形式局部范围内，$V$ 是两个光滑平面 $x = y = 0$ 与 $z = w = 0$
的并，它们各自与平面 $x - z = y - w = 0$ 的相交重数都为 $1$；见引理
\ref{intersection-lemma-transversal}。为了得到一个实际例子，取由 $5$ 个次数 $> 1$
的齐次多项式给出的一般态射 $f : \mathbf{P}^2 \to \mathbf{P}^4$。其像
$V \subset \mathbf{P}^4 = X$ 将具有上述类型的奇点，因为将有
$p_1, p_2 \in \mathbf{P}^2$ 使 $f(p_1) = f(p_2)$。要得到 $W$，只需取一个经过这样一点的一般平面。
\end{example}






\section{代数重数}
\label{intersection-section-multiplicities}

\noindent
设 $(A, \mathfrak m, \kappa)$ 为 Noether 局部环，$M$ 为有限 $A$-模，
$I \subset A$ 为定义理想；见《代数》定义
\ref{algebra-definition-ideal-definition}。回忆，函数
$$
\chi_{I, M}(n) = \text{length}_A(M/I^nM) =
\sum\nolimits_{p = 0, \ldots, n - 1} \text{length}_A(I^pM/I^{p + 1}M)
$$
是数值多项式；见《代数》命题
\ref{algebra-proposition-hilbert-function-polynomial}。由《代数》引理
\ref{algebra-lemma-support-dimension-d}，这个多项式的次数等于
$\dim(\text{Supp}(M))$。

\begin{definition}
\label{intersection-definition-multiplicity}
在上述情形中，假设 $d \geq \dim(\text{Supp}(M))$。若
$d > \dim(\text{Supp}(M))$，则置 $e_I(M, d) = 0$；若
$d = \dim(\text{Supp}(M))$，则将 $e_I(M, d)$ 定义为数值多项式
$\chi_{I, M}$ 的首项系数的 $d!$ 倍。因此在两种情形下都有
$$
\chi_{I, M}(n) \sim e_I(M, d) \frac{n^d}{d!} + \text{低阶项}
$$
$M$ 对于定义理想 $I$ 的{\it 重数}是
$e_I(M) = e_I(M, \dim(\text{Supp}(M)))$。
\end{definition}

\noindent
这些重数具有下列性质。

\begin{lemma}
\label{intersection-lemma-multiplicity-ses}
设 $A$ 为 Noether 局部环，$I \subset A$ 为定义理想，
$0 \to M' \to M \to M'' \to 0$ 为有限 $A$-模的短正合列。若
$d \geq \dim(\text{Supp}(M))$，则
$$
e_I(M, d) = e_I(M', d) + e_I(M'', d)
$$
\end{lemma}

\begin{proof}
立即由定义与《代数》引理 \ref{algebra-lemma-hilbert-ses-chi} 得出。
\end{proof}

\begin{lemma}
\label{intersection-lemma-multiplicity-as-a-sum}
设 $A$ 为 Noether 局部环，$I \subset A$ 为定义理想，$M$ 为有限 $A$-模。
若 $d \geq \dim(\text{Supp}(M))$，则
$$
e_I(M, d) =
\sum \text{length}_{A_\mathfrak p}(M_\mathfrak p) e_I(A/\mathfrak p, d)
$$
其中对所有满足 $\dim(A/\mathfrak p) = d$ 的素理想
$\mathfrak p \subset A$ 求和。
\end{lemma}

\begin{proof}
对于维数 $\leq d$ 的模的短正合列，等式左边与右边都具有可加性；见引理
\ref{intersection-lemma-multiplicity-ses} 与《代数》引理
\ref{algebra-lemma-length-additive}。因此，由《代数》引理
\ref{algebra-lemma-filter-Noetherian-module}，只需证明 $M = A/\mathfrak q$ 的情形，其中
$\mathfrak q$ 是 $A$ 的素理想，且 $\dim(A/\mathfrak q) \leq d$。此时结论显然。
\end{proof}

\begin{lemma}
\label{intersection-lemma-leading-coefficient}
设 $P$ 为次数为 $r$、首项系数为 $a$ 的多项式。则对任意 $t$，有
$$
r! a = \sum\nolimits_{i = 0, \ldots, r} (-1)^i{r \choose i} P(t - i)
$$
\end{lemma}

\begin{proof}
以 $\Delta$ 表示把多项式 $P$ 映为多项式
$\Delta(P) = P(t) - P(t - 1)$ 的算子。断言
$$
\Delta^r(P) = \sum\nolimits_{i = 0, \ldots, r} (-1)^i {r \choose i} P(t - i)
$$
当 $r = 0, 1$ 时直接验证即可。假设它对 $r$ 成立，则计算得
\begin{align*}
\Delta^{r + 1}(P)
& =
\sum\nolimits_{i = 0, \ldots, r} (-1)^i {r \choose i} \Delta(P)(t - i) \\
& =
\sum\nolimits_{n = -r, \ldots, 0} (-1)^i {r \choose i}
(P(t - i) - P(t - i - 1))
\end{align*}
所以断言由等式
$$
{r + 1 \choose i} = {r \choose i} + {r \choose i - 1}
$$
得出。当 $P$ 的次数为 $r$ 时，$\Delta(P)$ 的次数为 $r - 1$，
首项系数为 $ra$，故引理成立。
\end{proof}

\noindent
一个重要事实是，重数可以用 Koszul 复形计算。回忆，若 $R$ 是环，且
$f_1, \ldots, f_r \in R$，则 $K_\bullet(f_1, \ldots, f_r)$ 表示 Koszul 复形；见
《代数进阶》第 \ref{more-algebra-section-koszul} 节。

\begin{theorem}
\label{intersection-theorem-multiplicity-with-koszul}
\begin{reference}
\cite[Theorem 1 in part B of Chapter IV]{Serre_algebre_locale}
\end{reference}
设 $A$ 为 Noether 局部环，$I = (f_1, \ldots, f_r) \subset A$ 为定义理想，
$M$ 为有限 $A$-模。则
$$
e_I(M, r) = \sum
(-1)^i\text{length}_A H_i(K_\bullet(f_1, \ldots, f_r) \otimes_A M)
$$
\end{theorem}

\begin{proof}
通过置 $K^n = K_{-n}(f_1, \ldots, f_r)$，把 Koszul 复形
$K_\bullet(f_1, \ldots, f_r)$ 改写为上链复形 $K^\bullet$。则 $K^\bullet$
位于次数 $-r, \ldots, 0$，且
$H^i(K^\bullet \otimes_A M) = H_{-i}(K_\bullet(f_1, \ldots, f_r) \otimes_A M)$。
定理的断言是有意义的，因为模 $H^i(K^\bullet \otimes M)$ 均被
$f_1, \ldots, f_r$ 零化；见《代数进阶》引理
\ref{more-algebra-lemma-homotopy-koszul}，因此它们具有有限长度。
在复形 $K^\bullet$ 上定义滤过：
$$
F^p(K^n \otimes_A M) =
I^{\max(0, p + n)}(K^n \otimes_A M),\quad p \in \mathbf{Z}
$$
由于 $f_i I^p \subset I^{p + 1}$，这是由子复形给出的滤过。因此得到一个滤过复形及其谱序列；
见《同调》第 \ref{homology-section-filtered-complex} 节。有
$$
E_0 = \bigoplus\nolimits_{p, q} E_0^{p, q} =
\bigoplus\nolimits_{p, q} \text{gr}^p(K^{p + q} \otimes_A M) =
\text{Gr}_I(K^\bullet \otimes_A M)
$$
因为 $K^n$ 有限自由，有
$$
\text{Gr}_I(K^\bullet \otimes_A M) =
\text{Gr}_I(K^\bullet) \otimes_{\text{Gr}_I(A)} \text{Gr}_I(M)
$$
注意，$\text{Gr}_I(K^\bullet)$ 是环 $\text{Gr}_I(A)$ 上关于各元
$\overline{f}_1, \ldots, \overline{f}_r \in I/I^2$ 的 Koszul 复形。一个简单计算（略去）表明，
$E_0$ 上的微分 $d_0$ 与来自 Koszul 复形的微分一致。由于
$\text{Gr}_I(M)$ 是有限 $\text{Gr}_I(A)$-模，而 $\text{Gr}_I(A)$ 是 Noether 环：它是
$A/I[x_1, \ldots, x_r]$ 在 $x_i \mapsto \overline{f}_i$ 下的商，故同调模
$E_1 = \bigoplus E_1^{p, q}$ 是有限 $\text{Gr}_I(A)$-模。然而与上文一样，
$E_1$ 被 $\overline{f}_1, \ldots, \overline{f}_r$ 零化，因而 $E_1$ 具有有限长度。
特别地，当 $p \gg 0$ 时，$\text{Gr}^p_F(K^\bullet \otimes M)$ 无同调。

\medskip\noindent
下面使用《同调》引理 \ref{homology-lemma-filtered-complex-ss-converges}
验证上述谱序列收敛。所需等式容易由 Artin--Rees 引理得出，其形式见《代数》引理
\ref{algebra-lemma-map-AR}。于是
\begin{align*}
\sum (-1)^i\text{length}_A(H^i(K^\bullet \otimes_A M))
& =
\sum (-1)^{p + q} \text{length}_A(E_\infty^{p, q}) \\
& =
\sum (-1)^{p + q} \text{length}_A(E_1^{p, q})
\end{align*}
因为上文已证明 $E_1$ 的长度有限（当然，这里使用了长度的可加性）。选取足够大的 $t$，
使得当 $p \geq t$ 时 $\text{Gr}^p_F(K^\bullet \otimes M)$ 无同调（见上文）。再次使用可加性得
$$
\sum (-1)^{p + q} \text{length}_A(E_1^{p, q}) =
\sum\nolimits_n \sum\nolimits_{p \leq t}
(-1)^n \text{length}_A(\text{gr}^p(K^n \otimes_A M))
$$
根据上述滤过的选取，以及《代数》第
\ref{algebra-section-Noetherian-local} 节中 $\chi_{I, M}$ 的定义，这等于
$$
\sum\nolimits_{n = -r, \ldots, 0} (-1)^n{r \choose |n|} \chi_{I, M}(t + n)
$$
定理由引理 \ref{intersection-lemma-leading-coefficient} 与 $e_I(M, r)$ 的定义得出。
\end{proof}

\begin{remark}[平凡推广]
\label{intersection-remark-trivial-generalization}
设 $(A, \mathfrak m, \kappa)$ 为 Noether 局部环，$M$ 为有限 $A$-模，$I \subset A$ 为理想。
下列条件等价：
\begin{enumerate}
\item $I' = I + \text{Ann}(M)$ 是定义理想；见《代数》定义
\ref{algebra-definition-ideal-definition}；
\item $I$ 在 $\overline{A} = A/\text{Ann}(M)$ 中的像 $\overline{I}$ 是定义理想；
\item $\text{Supp}(M/IM) \subset \{\mathfrak m\}$；
\item $\dim(\text{Supp}(M/IM)) \leq 0$；以及
\item $\text{length}_A(M/IM) < \infty$。
\end{enumerate}
这由《代数》引理 \ref{algebra-lemma-support-point} 得出（略去细节）。若这些条件成立，
则对所有 $n$ 都有 $M/I^nM = M/(I')^nM$；而当把 $M$ 看作 $\overline{A}$-模时，
对所有 $n$ 都有 $M/I^nM = M/\overline{I}^nM$。因此可定义
$$
\chi_{I, M}(n) = \text{length}_A(M/I^nM) =
\sum\nolimits_{p = 0, \ldots, n - 1} \text{length}_A(I^pM/I^{p + 1}M)
$$
并且由上述等式，对所有 $n$ 都有
$$
\chi_{I, M}(n) = \chi_{I', M}(n) = \chi_{\overline{I}, M}(n)
$$
《代数》第 \ref{algebra-section-Noetherian-local} 节的所有结果以及本节的所有结果，
在此设定中都有对应形式。特别地，对 $d \geq \dim(\text{Supp}(M))$，可以定义重数
$e_I(M, d)$，并且如 $I$ 为定义理想时一样，有
$$
\chi_{I, M}(n) \sim e_I(M, d) \frac{n^d}{d!} + \text{低阶项}
$$
\end{remark}



\section{相交重数的计算}
\label{intersection-section-computing-intersection-multiplicities}

\noindent
本节讨论若干可以用其他方法计算相交重数的情形。先看第一个例子。

\begin{lemma}
\label{intersection-lemma-intersection-multiplicity-CM}
设 $X$ 为非奇异簇，$W, V \subset X$ 为正常相交的闭子簇。设 $Z$ 为
$V \cap W$ 的一个不可约分支，其泛点为 $\xi$。假设 $\mathcal{O}_{W, \xi}$
与 $\mathcal{O}_{V, \xi}$ 均为 Cohen-Macaulay。则
$$
e(X, V \cdot W, Z) =
\text{length}_{\mathcal{O}_{X, \xi}}(\mathcal{O}_{V \cap W, \xi})
$$
其中 $V \cap W$ 是概形论之交。特别地，若 $V$ 与 $W$ 都是 Cohen-Macaulay，则
$V \cdot W = [V \cap W]_{\dim(V) + \dim(W) - \dim(X)}$。
\end{lemma}

\begin{proof}
置 $A = \mathcal{O}_{X, \xi}$、$B = \mathcal{O}_{V, \xi}$ 与
$C = \mathcal{O}_{W, \xi}$。由 Auslander--Buchsbaum 定理（《代数》命题
\ref{algebra-proposition-Auslander-Buchsbaum}），可找到它的一个有限自由解析
$F_\bullet \to B$，其长度为
$$
\text{depth}(A) - \text{depth}(B) =
\dim(A) - \dim(B) = \dim(C)
$$
第一个等式是因为 $A$ 与 $B$ 为 Cohen-Macaulay，第二个等式是因为
$V$ 与 $W$ 正常相交。于是 $F_\bullet \otimes_A C$ 是代表
$B \otimes_A^\mathbf{L} C$ 的有限自由模复形，故其同调模的支撑包含于
$\{\mathfrak m_A\}$。由于 $C$ 是 Cohen-Macaulay，可以应用无同调引理（《代数》引理
\ref{algebra-lemma-acyclic}），得到 $F_\bullet \otimes_A C$ 只在次数 $0$ 上有非零同调。证明完毕。
\end{proof}

\begin{lemma}
\label{intersection-lemma-one-ideal-ci}
设 $A$ 为 Noether 局部环，$I = (f_1, \ldots, f_r)$ 为由正则序列生成的理想，
$M$ 为有限 $A$-模。假设 $\dim(\text{Supp}(M/IM)) = 0$。则
$$
e_I(M, r) = \sum (-1)^i\text{length}_A(\text{Tor}_i^A(A/I, M))
$$
这里 $e_I(M, r)$ 如注 \ref{intersection-remark-trivial-generalization} 中所定义。
\end{lemma}

\begin{proof}
由于 $f_1, \ldots, f_r$ 是正则序列，Koszul 复形
$K_\bullet(f_1, \ldots, f_r)$ 是 $A/I$ 在 $A$ 上的一个解析；见《代数进阶》引理
\ref{more-algebra-lemma-noetherian-finite-all-equivalent}。因此等式右边等于
$$
\sum (-1)^i\text{length}_A H_i(K_\bullet(f_1, \ldots, f_r) \otimes_A M)
$$
若 $I$ 是定义理想，结论立即由定理 \ref{intersection-theorem-multiplicity-with-koszul} 得出。
在一般情形下，以 $\overline{A} = A/\text{Ann}(M)$ 取代 $A$，并以
$\overline{f}_1, \ldots, \overline{f}_r$ 取代 $f_1, \ldots, f_r$。这是允许的，因为
$$
K_\bullet(f_1, \ldots, f_r) \otimes_A M =
K_\bullet(\overline{f}_1, \ldots, \overline{f}_r) \otimes_{\overline{A}} M
$$
又因为 $e_I(M, r) = e_{\overline{I}}(M, r)$，其中
$\overline{I} = (\overline{f}_1, \ldots, \overline{f}_r) \subset \overline{A}$ 是定义理想，
在这种情形下再次应用定理 \ref{intersection-theorem-multiplicity-with-koszul} 即得结论。
\end{proof}

\begin{lemma}
\label{intersection-lemma-multiplicity-with-lci}
设 $X$ 为非奇异簇，$W,V \subset X$ 为正常相交的闭子簇。设 $Z$ 为
$V \cap W$ 的一个不可约分支，其泛点为 $\xi$。假设 $V$ 在
$\mathcal{O}_{X, \xi}$ 中的理想由正则序列
$f_1, \ldots, f_c \in \mathcal{O}_{X, \xi}$ 截出。则 $e(X, V\cdot W, Z)$ 等于下列 Hilbert 多项式
的首项系数的 $c!$ 倍：
$$
t \mapsto \text{length}_{\mathcal{O}_{X, \xi}}
\mathcal{O}_{W, \xi}/(f_1, \ldots, f_c)^t,\quad t \gg 0.
$$
特别地，该系数 $> 0$。
\end{lemma}

\begin{proof}
等式
$$
e(X, V\cdot W, Z) = e_{(f_1, \ldots, f_c)}(\mathcal{O}_{W, \xi}, c)
$$
由更一般的引理 \ref{intersection-lemma-one-ideal-ci} 得出。为了证明
$e_{(f_1, \ldots, f_c)}(\mathcal{O}_{W, \xi}, c)$ 为 $> 0$，或者等价地，为了证明
$e_{(f_1, \ldots, f_c)}(\mathcal{O}_{W, \xi}, c)$ 是 Hilbert 多项式的首项系数，
只需证明 $\mathcal{O}_{W, \xi}$ 的维数是 $c$，因为由《代数》命题
\ref{algebra-proposition-dimension}，Hilbert 多项式的次数等于该维数。设
$\dim(V) = r$、$\dim(W) = s$ 且 $\dim(X) = n$。由交的正常性，
$\dim(Z) = r + s - n$。因此 $\kappa(\xi)$ 对 $\mathbf{C}$ 的超越次数为
$r + s - n$；见《代数》引理 \ref{algebra-lemma-dimension-prime-polynomial-ring}。
由于 $V$ 在 $\xi$ 的一个邻域中由正则序列截出，有 $r + c = n$；见《除子》引理
\ref{divisors-lemma-Noetherian-scheme-regular-ideal}，此后例如可应用引理
\ref{intersection-lemma-pullback-by-regular-immersion}。因此
$$
\dim(\mathcal{O}_{W, \xi}) = s - (r + s - n) = s - ((n - c) + s - n) = c
$$
其中第一个等式由《代数》引理
\ref{algebra-lemma-dimension-at-a-point-finite-type-field} 得出。
\end{proof}

\begin{lemma}
\label{intersection-lemma-multiplicity-with-effective-Cartier-divisor}
在引理 \ref{intersection-lemma-multiplicity-with-lci} 中假设 $c = 1$，即 $V$ 是有效 Cartier 除子。则
$$
e(X, V \cdot W, Z) =
\text{length}_{\mathcal{O}_{X, \xi}}
(\mathcal{O}_{W, \xi}/f_1\mathcal{O}_{W, \xi}).
$$
\end{lemma}

\begin{proof}
此时，由交的正常性，$f_1$ 在 $\mathcal{O}_{W, \xi}$ 中的像非零，因而是非零因子。
此外，$\mathcal{O}_{W, \xi}$ 是维数为 $1$ 的 Noether 局部整环。所以对所有
$t \geq 1$，有
$$
\text{length}_{\mathcal{O}_{X, \xi}}
(\mathcal{O}_{W, \xi}/f_1^t\mathcal{O}_{W, \xi}) =
t \text{length}_{\mathcal{O}_{X, \xi}}
(\mathcal{O}_{W, \xi}/f_1\mathcal{O}_{W, \xi})
$$
见《代数》引理 \ref{algebra-lemma-ord-additive}。这就证明了引理。
\end{proof}

\begin{lemma}
\label{intersection-lemma-multiplicity-lci-CM}
在引理 \ref{intersection-lemma-multiplicity-with-lci} 中，假设局部环
$\mathcal{O}_{W, \xi}$ 是 Cohen-Macaulay。则
$$
e(X, V \cdot W, Z) =
\text{length}_{\mathcal{O}_{X, \xi}} (\mathcal{O}_{W, \xi}/
f_1\mathcal{O}_{W, \xi} + \ldots + f_c\mathcal{O}_{W, \xi}).
$$
\end{lemma}

\begin{proof}
这立即由引理 \ref{intersection-lemma-intersection-multiplicity-CM} 得出。另一种方法是从引理
\ref{intersection-lemma-multiplicity-with-lci} 推导。事实上，由《代数》引理
\ref{algebra-lemma-reformulate-CM}，$f_1, \ldots, f_c$ 是 $\mathcal{O}_{W, \xi}$ 中的正则序列。
再由《代数》引理 \ref{algebra-lemma-regular-quasi-regular}，
$f_1, \ldots, f_c$ 是拟正则序列。这容易蕴含
$\mathcal{O}_{W, \xi}/(f_1, \ldots, f_c)^t$ 的长度等于
$$
{c + t \choose c}
\text{length}_{\mathcal{O}_{X, \xi}} (\mathcal{O}_{W, \xi}/
f_1\mathcal{O}_{W, \xi} + \ldots + f_c\mathcal{O}_{W, \xi}).
$$
比较首项系数即得结论。
\end{proof}


\section{用 Tor 公式定义交积}
\label{intersection-section-intersection-product}

\noindent
设 $X$ 为非奇异簇，$\alpha = \sum n_i [W_i]$ 为 $r$-循环，
$\beta = \sum_j m_j [V_j]$ 为 $X$ 上的 $s$-循环。假设 $\alpha$ 与 $\beta$ 正常相交；
见定义 \ref{intersection-definition-proper-intersection}。此时定义
$$
\alpha \cdot \beta = \sum\nolimits_{i,j} n_i m_j W_i \cdot V_j.
$$
其中 $W_i \cdot V_j$ 如第 \ref{intersection-section-tor-formula} 节中所定义。若
$\beta = [V]$，其中 $V$ 是维数为 $s$ 的闭子簇，则有时将
$\alpha \cdot \beta = \alpha \cdot V$ 记作这一形式。

\begin{lemma}
\label{intersection-lemma-rational-equivalence-and-intersection}
设 $X$ 为非奇异簇，$a, b \in \mathbf{P}^1$ 为不同的闭点，且 $k \geq 0$。
\begin{enumerate}
\item 若 $W \subset X \times \mathbf{P}^1$ 是维数为 $k + 1$ 且与 $X \times a$ 正常相交的闭子簇，则
\begin{enumerate}
\item 作为 $X \times \mathbf{P}^1$ 上的循环，$[W_a]_k = W \cdot X \times a$；以及
\item 作为 $X$ 上的循环，$[W_a]_k = \text{pr}_{X, *}(W \cdot X \times a)$。
\end{enumerate}
\item 设 $\alpha$ 为 $X \times \mathbf{P}^1$ 上的 $(k + 1)$-循环，它与
$X \times a$ 及 $X \times b$ 都正常相交。则
$pr_{X,*}( \alpha \cdot X \times a - \alpha \cdot X \times b)$ 有理等价于零。
\item 反之，任意有理等价于 $0$ 的 $k$-循环都具有这种形式。
\end{enumerate}
\end{lemma}

\begin{proof}
首先注意，$X \times a$ 是 $X \times \mathbf{P}^1$ 中的有效 Cartier 除子，而 $W_a$
是 $W$ 与 $X \times a$ 的概形论之交。因此，(1)(a) 中的等式立即由定义以及引理
\ref{intersection-lemma-multiplicity-with-effective-Cartier-divisor} 中对 Cartier 除子情形的相交重数计算得出。
(1)(b) 成立，是因为 $W_a \to X \times \mathbf{P}^1 \to X$ 同构地映到其像上；
在第 \ref{intersection-section-rational-equivalence} 节中，我们正是以这种方式把 $W_a$ 看作 $X$ 的闭子概形。
(2) 与 (3) 是 (1) 与各定义的形式推论。
\end{proof}

\noindent
对于闭子概形的横截相交，相交重数为 $1$。

\begin{lemma}
\label{intersection-lemma-transversal-subschemes}
设 $X$ 为非奇异簇，$r, s \geq 0$，并设 $Y, Z \subset X$ 为闭子概形，其满足
$\dim(Y) \leq r$ 与 $\dim(Z) \leq s$。假设
$[Y]_r = \sum n_i[Y_i]$ 与 $[Z]_s = \sum m_j[Z_j]$ 正常相交。设 $T$ 是对某些
$i_0$ 与 $j_0$ 而言 $Y_{i_0} \cap Z_{j_0}$ 的一个不可约分支，并假设闭子概形
$Y \cap Z$ 中 $T$ 的重数（按第 \ref{intersection-section-cycle-of-closed} 节的意义）为 $1$。则
\begin{enumerate}
\item $[Y]_r \cdot [Z]_s$ 中 $T$ 的系数为 $1$；
\item $Y$ 与 $Z$ 在 $T$ 的泛点处非奇异；
\item $n_{i_0} = 1$、$m_{j_0} = 1$；以及
\item 当 $i \not = i_0$ 且 $j \not = j_0$ 时，$T$ 不包含于 $Y_i$ 或 $Z_j$ 中。
\end{enumerate}
\end{lemma}

\begin{proof}
置 $n = \dim(X)$、$a = n - r$、$b = n - s$。由交的正常性，注意到
$\dim(T) = r + s - n = n - a - b$。设 $(A, \mathfrak m, \kappa) =
(\mathcal{O}_{X, \xi}, \mathfrak m_\xi, \kappa(\xi))$，其中 $\xi \in T$ 是泛点。则
$\dim(A) = a + b$；见《簇》引理
\ref{varieties-lemma-dimension-locally-algebraic}。设 $I_0, I, J_0, J \subset A$ 分别截出
$Y_{i_0}$、$Y$、$Z_{j_0}$、$Z$ 在 $\Spec(A)$ 中的迹。由同一引用，
$\dim(A/I) = \dim(A/I_0) = b$ 且 $\dim(A/J) = \dim(A/J_0) = a$。
置 $\overline{I} = I + \mathfrak m^2/\mathfrak m^2$。则
$I \subset I_0 \subset \mathfrak m$、$J \subset J_0 \subset \mathfrak m$，且 $I + J = \mathfrak m$。
由引理 \ref{intersection-lemma-transversal} 及其证明，
$I_0 = (f_1, \ldots, f_a)$ 且 $J_0 = (g_1, \ldots, g_b)$，其中
$f_1, \ldots, g_b$ 是正则局部环 $A$ 的一组正则参数系。由于
$I + J = \mathfrak m$，映射
$$
I \oplus J \to
\mathfrak m/\mathfrak m^2 =
\kappa f_1
\oplus \ldots \oplus
\kappa f_a
\oplus
\kappa g_1
\oplus \ldots \oplus
\kappa g_b
$$
是满的。因此可找到
$f_1', \ldots, f_a' \in I$ 与 $g'_1, \ldots, g_b' \in J$，使它们在
$\mathfrak m/\mathfrak m^2$ 中的剩余类分别等于 $f_1, \ldots, f_a$ 与
$g_1, \ldots, g_b$ 的剩余类。于是 $f'_1, \ldots, g'_b$ 是 $A$ 的一组正则参数系。
由《代数》引理 \ref{algebra-lemma-regular-ring-CM}，
$A/(f'_1, \ldots, f'_a)$ 是维数为 $b$ 的正则局部环。因此
$A/(f'_1, \ldots, f'_a)$ 的任何非平凡商环
都具有严格更小的维数；见《代数》引理
\ref{algebra-lemma-regular-domain} 与 \ref{algebra-lemma-one-equation}。故
$I = (f'_1, \ldots, f'_a) = I_0$。由对称性，$J = J_0$。这证明了 (2)、(3) 与 (4)。
最后，$[Y]_r \cdot [Z]_s$ 中 $T$ 的系数等于
$Y_{i_0} \cdot Z_{j_0}$ 中 $T$ 的系数，由引理 \ref{intersection-lemma-transversal} 它等于 $1$。
\end{proof}



\section{外积}
\label{intersection-section-exterior-product}

\noindent
设 $X$ 与 $Y$ 为簇，而 $V$、\ $W$ 分别为 $X$、\ $Y$ 的闭子簇。
乘积 $V\times W$ 是 $X\times Y$ 的闭子簇；见引理
\ref{intersection-lemma-dimension-product-varieties}。对 $k$-循环
$\alpha = \sum n_i [V_i]$ 和 $Y$ 上的 $l$-循环 $\beta = \sum m_j [V_j]$，
定义 $\alpha$ 与 $\beta$ 的{\it 外积}为循环
$\alpha \times \beta = \sum n_i m_j [V_i \times W_j]$。外积定义了一个 $\mathbf{Z}$-线性映射
$$
Z_r(X) \otimes_\mathbf{Z} Z_s(Y) \longrightarrow Z_{r + s}(X \times Y)
$$
下面证明外积可以经过有理等价分解。

\begin{lemma}
\label{intersection-lemma-exterior-product-rational-equivalence}
设 $X$ 与 $Y$ 为簇，$\alpha \in Z_r(X)$、$\beta \in Z_s(Y)$。若
$\alpha \sim_{rat} 0$ 或 $\beta \sim_{rat} 0$，则 $\alpha \times \beta \sim_{rat} 0$。
\end{lemma}

\begin{proof}
由线性以及 $X$ 与 $Y$ 之间的对称性，只需证明如下情形：
$\alpha = [V]$，其中 $V \subset X$ 是维数为 $r$ 的子簇；并且
$\beta = [W_a]_s - [W_b]_s$，其中 $W \subset Y \times \mathbf{P}^1$ 是维数为
$s + 1$ 的闭子簇，且与 $Y \times a$、$Y \times b$ 均正常相交。此时，只要证明
$$
[(V \times W)_a]_{r + s} = [V] \times [W_a]_s
$$
以及将 $a$ 换为 $b$ 后的同类等式，引理就成立。事实上，此时
$\alpha \times \beta = [(V \times W)_a]_{r + s} - [(V \times W)_b]_{r + s}$，正如所需。
为证明上述显示等式，注意到概形的等式
$$
V \times W_a = (V \times W)_a
$$
投影 $V \times W_a \to W_a$ 诱导不可约分支之间的双射；例如见《簇》引理
\ref{varieties-lemma-bijection-irreducible-components}。设 $W' \subset W_a$ 是泛点为 $\zeta$
的不可约分支。则 $V \times W'$ 是 $V \times W_a$ 的对应不可约分支；见引理
\ref{intersection-lemma-dimension-product-varieties}。设 $\xi$ 为 $V \times W'$ 的泛点。我们需要证明
$$
\text{length}_{\mathcal{O}_{Y, \zeta}}(\mathcal{O}_{W_a, \zeta}) =
\text{length}_{\mathcal{O}_{X \times Y, \xi}}(
\mathcal{O}_{V \times W_a, \xi})
$$
在该公式中，可用 $\mathcal{O}_{W_a, \zeta}$ 取代 $\mathcal{O}_{Y, \zeta}$，并用
$\mathcal{O}_{V \times W_a, \zeta}$ 取代 $\mathcal{O}_{X \times Y, \zeta}$；见《代数》引理
\ref{algebra-lemma-length-independent}。由于
$\mathcal{O}_{W_a, \zeta} \to \mathcal{O}_{V \times W_a, \xi}$ 平坦，根据《代数》引理
\ref{algebra-lemma-pullback-module}，只需证明
$$
\text{length}_{\mathcal{O}_{V \times W_a, \xi}}(
\mathcal{O}_{V \times W_a, \xi}/
\mathfrak m_\zeta\mathcal{O}_{V \times W_a, \xi}) = 1
$$
这是正确的，因为右边的商是簇在一个泛点处的局部环
$\mathcal{O}_{V \times W', \xi}$，因而等于 $\kappa(\xi)$。
\end{proof}

\noindent
由此得出，对任意一对簇 $X$ 与 $Y$，外积定义交换图
$$
\xymatrix{
Z_r(X) \otimes_\mathbf{Z} Z_s(Y) \ar[r] \ar[d] &
Z_{r + s}(X \times Y) \ar[d] \\
\CH_r(X) \otimes_\mathbf{Z} \CH_s(Y) \ar[r] &
\CH_{r + s}(X \times Y)
}
$$
对于非奇异簇，可将外积视为若干拉回的交积。

\begin{lemma}
\label{intersection-lemma-exterior-product}
设 $X$ 与 $Y$ 为非奇异簇，$\alpha \in Z_r(X)$、$\beta \in Z_s(Y)$。则
\begin{enumerate}
\item $\text{pr}_Y^*(\beta) = [X] \times \beta$ 且
$\text{pr}_X^*(\alpha) = \alpha \times [Y]$；
\item $\alpha \times [Y]$ 与 $[X]\times \beta$ 在 $X\times Y$ 上正常相交；以及
\item 在 $Z_{r + s}(X \times Y)$ 中有
$\alpha \times \beta =
(\alpha \times [Y])\cdot ([X]\times\beta) =
pr_Y^*(\alpha) \cdot pr_X^*(\beta)$。
\end{enumerate}
\end{lemma}

\begin{proof}
由线性，可以假设 $\alpha = [V]$ 且 $\beta = [W]$。此时 (1) 说的是
$\text{pr}_Y^{-1}(W) = X \times W$ 与 $\text{pr}_X^{-1}(V) = V \times Y$，这是显然的。
(2) 成立，因为 $X \times W \cap V \times Y = V \times W$，而由引理
\ref{intersection-lemma-dimension-product-varieties}，$\dim(V \times W) = r + s$。

\medskip\noindent
证明 (3)。设 $\xi$ 为 $V \times W$ 的泛点。投影 $X \times W \to W$ 是
$X \to \Spec(\mathbf{C})$ 的基变换，因而光滑。所以 $X \times W$ 在位于 $W$
的泛点之上的每个点处都非奇异，特别是在 $\xi$ 处。对 $V \times Y$ 亦然。因此
$\mathcal{O}_{X \times W, \xi}$ 与 $\mathcal{O}_{V \times Y, \xi}$ 是 Cohen-Macaulay 局部环，可应用引理
\ref{intersection-lemma-intersection-multiplicity-CM}。由于在概形论意义下
$V \times Y \cap X \times W = V \times W$，证明完毕。
\end{proof}



\section{化归到对角线}
\label{intersection-section-reduction-diagonal}

\noindent
设 $X$ 为非奇异簇。我们用 $\Delta$ 既表示对角态射
$\Delta : X \to X \times X$，也表示其像 $\Delta \subset X \times X$。
所谓化归到对角线，是指 $X$ 上的交积可化为 $X \times X$ 上外积与对角线的交积。

\begin{lemma}
\label{intersection-lemma-tor-and-diagonal}
设 $X$ 为非奇异簇。
\begin{enumerate}
\item 若 $\mathcal{F}$ 与 $\mathcal{G}$ 是凝聚 $\mathcal{O}_X$-模，则有典范同构
$$
\text{Tor}_i^{\mathcal{O}_{X \times X}}(\mathcal{O}_\Delta,
\text{pr}_1^*\mathcal{F} \otimes_{\mathcal{O}_{X \times X}}
\text{pr}_2^*\mathcal{G})
=
\Delta_*\text{Tor}_i^{\mathcal{O}_X}(\mathcal{F}, \mathcal{G})
$$
\item 若 $K$ 与 $M$ 属于 $D_\QCoh(\mathcal{O}_X)$，则在 $D_\QCoh(\mathcal{O}_X)$ 中有典范同构
$$
L\Delta^* \left(
L\text{pr}_1^*K \otimes_{\mathcal{O}_{X \times X}}^\mathbf{L} L\text{pr}_2^*M
\right)
= K \otimes_{\mathcal{O}_X}^\mathbf{L} M
$$
并且在 $D_\QCoh(X \times X)$ 中有典范同构
$$
\mathcal{O}_\Delta \otimes_{\mathcal{O}_{X \times X}}^\mathbf{L}
L\text{pr}_1^*K \otimes_{\mathcal{O}_{X \times X}}^\mathbf{L} L\text{pr}_2^*M
= \Delta_*(K \otimes_{\mathcal{O}_X}^\mathbf{L} M)
$$
\end{enumerate}
\end{lemma}

\begin{proof}
我们用较为初等的方法证明 (1)，并用更一般的理论证明 (2)。由于 (2) 蕴含 (1)，
读者可以跳过 (1) 的证明。

\medskip\noindent
证明 (1)。选取仿射开集 $\Spec(A) \subset X$。则 $A$ 是 Noether $\mathbf{C}$-代数，
$\mathcal{F}$、$\mathcal{G}$ 分别对应于有限 $A$-模 $M$ 与 $N$；见《概形的上同调》引理
\ref{coherent-lemma-coherent-Noetherian}。由《概形的导出范畴》引理
\ref{perfect-lemma-quasi-coherence-tensor-product}，可以先在 $A$ 上计算 $M$ 与 $N$
的 $\text{Tor}$，再取其关联 $\mathcal{O}_X$-模，从而计算 $\mathcal{O}_X$ 上的 $\text{Tor}_i$。
对于 $\mathcal{O}_{X \times X}$ 上的 $\text{Tor}_i$，我们在 $A \otimes_\mathbf{C} A$ 上计算
$A$ 与 $M \otimes_\mathbf{C} N$ 的 tor，再取其关联 $\mathcal{O}_{X \times X}$-模。
因此，在这个仿射片上需证明
$$
\text{Tor}_i^{A \otimes_\mathbf{C} A}(A, M \otimes_\mathbf{C} N) =
\text{Tor}_i^A(M, N)
$$
为此，选取由有限自由 $A$-模组成的解析
$F_\bullet \to M$ 与 $G_\bullet \to M$；见《代数》引理
\ref{algebra-lemma-resolution-by-finite-free}。注意，
$\text{Tot}(F_\bullet \otimes_\mathbf{C} G_\bullet)$ 是 $M \otimes_\mathbf{C} N$ 的解析，因为它计算
$\mathbf{C}$ 上的 Tor 群。当然，$F_\bullet \otimes_\mathbf{C} G_\bullet$ 的各项都是有限自由
$A \otimes_\mathbf{C} A$-模。因此，显示等式的左边是模
$$
H_i(A \otimes_{A \otimes_\mathbf{C} A}
\text{Tot}(F_\bullet \otimes_\mathbf{C} G_\bullet))
$$
而右边是模
$$
H_i(\text{Tot}(F_\bullet \otimes_A G_\bullet))
$$
由于 $A \otimes_{A \otimes_\mathbf{C} A} (F_p \otimes_\mathbf{C} G_q)
= F_p \otimes_A G_q$，这些模相等。这在仿射开集
$\Spec(A) \times \Spec(A)$ 上定义了同构，对应用于相交数相等性而言这已经足够。
略去这些同构可粘合的证明。

\medskip\noindent
证明 (2)。由《概形的导出范畴》引理
\ref{perfect-lemma-cohomology-base-change} 中的投影公式，第二个断言由第一个断言得出。
为证明第一个断言，用 K-平坦复形 $\mathcal{K}^\bullet$ 与 $\mathcal{M}^\bullet$
分别代表 $K$ 与 $M$。由于拉回和张量积保持 K-平坦复形；见《上同调》引理
\ref{cohomology-lemma-tensor-product-K-flat} 与
\ref{cohomology-lemma-pullback-K-flat}，只需证明
$$
\Delta^*\text{Tot}(
\text{pr}_1^*\mathcal{K}^\bullet
\otimes_{\mathcal{O}_{X \times X}} \text{pr}_2^*\mathcal{M}^\bullet)
=
\text{Tot}(
\mathcal{K}^\bullet \otimes_{\mathcal{O}_X} \mathcal{M}^\bullet)
$$
因此，只需说明当 $\mathcal{K}$ 与 $\mathcal{M}$ 是 $\mathcal{O}_X$-模时，
无论它们是否拟凝聚或平坦，都有典范同构
$$
\Delta^*(\text{pr}_1^*\mathcal{K}
\otimes_{\mathcal{O}_{X \times X}} \text{pr}_2^*\mathcal{M})
\longrightarrow
\mathcal{K} \otimes_{\mathcal{O}_X} \mathcal{M}
$$
略去细节。
\end{proof}

\begin{lemma}
\label{intersection-lemma-reduction-diagonal}
设 $X$ 为非奇异簇。设 $\alpha$、\ $\beta$ 分别为 $X$ 上的 $r$-循环与\ $s$-循环。
假设 $\alpha$ 与 $\beta$ 正常相交。则
\begin{enumerate}
\item $\alpha \times \beta$ 与 $[\Delta]$ 正常相交；
\item 作为 $X \times X$ 上的循环，有
$\Delta_*(\alpha \cdot \beta) = [\Delta] \cdot \alpha\times\beta$；
\item 若 $X$ 是固有的，则
$\text{pr}_{1, *}([\Delta] \cdot \alpha\times\beta) = \alpha\cdot\beta$，
其中 $pr_1 : X\times X \to X$ 是投影。
\end{enumerate}
\end{lemma}

\begin{proof}
由线性，只需证明 $\alpha = [V]$ 且 $\beta = [W]$ 的情形，其中
$V \subset X$ 与 $W \subset Y$ 是正常相交的闭子簇。回忆，$V \times W$ 是维数为
$r + s$ 的闭子簇。注意，在概形论意义下既有
$V \cap W = \Delta^{-1}(V \times W)$，也有
$\Delta(V \cap W) = \Delta \cap V \times W$。这证明了 (1)。

\medskip\noindent
证明 (2)。设 $Z \subset V \cap W$ 是泛点为 $\xi$ 的不可约分支。
我们需要证明，$\alpha \cdot \beta$ 中 $Z$ 的系数等于
$[\Delta] \cdot \alpha \times \beta$ 中 $\Delta(Z)$ 的系数。前者是整数
$$
\sum (-1)^i
\text{length}_{\mathcal{O}_{X, \xi}}
\text{Tor}_i^{\mathcal{O}_X}(\mathcal{O}_V, \mathcal{O}_W)_\xi
$$
后者是整数
$$
\sum (-1)^i
\text{length}_{\mathcal{O}_{X \times Y, \Delta(\xi)}}
\text{Tor}_i^{\mathcal{O}_{X \times Y}}(
\mathcal{O}_\Delta, \mathcal{O}_{V \times W})_{\Delta(\xi)}
$$
然而，由引理 \ref{intersection-lemma-tor-and-diagonal}，有
$$
\text{Tor}_i^{\mathcal{O}_X}(\mathcal{O}_V, \mathcal{O}_W)_\xi \cong
\text{Tor}_i^{\mathcal{O}_{X \times Y}}(
\mathcal{O}_\Delta, \mathcal{O}_{V \times W})_{\Delta(\xi)}
$$
它们是 $\mathcal{O}_{X \times X, \Delta(\xi)}$-模。因此两者的长度相等；更确切地说，这里使用《代数》引理
\ref{algebra-lemma-length-independent}。

\medskip\noindent
(2) 蕴含 (3)，因为由引理 \ref{intersection-lemma-compose-pushforward}，有
$\text{pr}_{1, *} \circ \Delta_* = \text{id}$。
\end{proof}

\begin{proposition}
\label{intersection-proposition-positivity}
\begin{reference}
这是 \cite{Serre_algebre_locale} 的主要结果之一。
\end{reference}
设 $X$ 为非奇异簇，$V \subset X$ 与 $W \subset Y$ 为正常相交的闭子簇。
设 $Z \subset V \cap W$ 为不可约分支。则 $e(X, V \cdot W, Z) > 0$。
\end{proposition}

\begin{proof}
由引理 \ref{intersection-lemma-reduction-diagonal}，有
$$
e(X, V \cdot W, Z) = e(X \times X, \Delta \cdot V \times W, \Delta(Z))
$$
由于 $\Delta : X \to X \times X$ 是正则浸入（见引理
\ref{intersection-lemma-diagonal-regular-immersion}），根据引理
\ref{intersection-lemma-multiplicity-with-lci}，
$e(X \times X, \Delta \cdot V \times W, \Delta(Z))$ 是正整数。
\end{proof}

\noindent
下一引理是本章建立该理论时的一个关键引理。它本质上告诉我们，相交数具有适当的可加性。

\begin{lemma}
\label{intersection-lemma-tor-sheaf}
\begin{reference}
\cite[Chapter V]{Serre_algebre_locale}
\end{reference}
设 $X$ 为非奇异簇，$\mathcal{F}$ 与 $\mathcal{G}$ 为 $X$ 上的凝聚层，并满足
$\dim(\text{Supp}(\mathcal{F})) \leq r$、
$\dim(\text{Supp}(\mathcal{G})) \leq s$ 以及
$\dim(\text{Supp}(\mathcal{F}) \cap \text{Supp}(\mathcal{G}) )
\leq r + s - \dim X$。此时 $[\mathcal{F}]_r$ 与 $[\mathcal{G}]_s$ 正常相交，且
$$
[\mathcal{F}]_r \cdot [\mathcal{G}]_s =
\sum (-1)^p
[\text{Tor}_p^{\mathcal{O}_X}(\mathcal{F}, \mathcal{G})]_{r + s - \dim(X)}.
$$
\end{lemma}

\begin{proof}
$[\mathcal{F}]_r$ 与 $[\mathcal{G}]_s$ 正常相交的断言是立即的。我们证明的是循环等式，
故可在 $X$ 上局部地工作。注意，循环交积、$\text{Tor}$-层以及凝聚层所关联的循环的形成，
都与限制到开子概形可交换。因此可以并且就假设 $X$ 仿射。

\medskip\noindent
记
$$
RHS(\mathcal{F}, \mathcal{G}) = [\mathcal{F}]_r \cdot [\mathcal{G}]_s
\quad\text{且}\quad
LHS(\mathcal{F}, \mathcal{G}) = \sum (-1)^p
[\text{Tor}_p^{\mathcal{O}_X}(\mathcal{F}, \mathcal{G})]_{r + s - \dim(X)}
$$
考虑 $X$ 上凝聚层的短正合列
$$
0 \to \mathcal{F}_1 \to \mathcal{F}_2 \to \mathcal{F}_3 \to 0
$$
并假设 $\text{Supp}(\mathcal{F}_i) \subset \text{Supp}(\mathcal{F})$。则对
$i = 1, 2, 3$，$LHS(\mathcal{F}_i, \mathcal{G})$ 与
$RHS(\mathcal{F}_i, \mathcal{G})$ 都有定义，且
$$
RHS(\mathcal{F}_2, \mathcal{G}) =
RHS(\mathcal{F}_1, \mathcal{G}) + RHS(\mathcal{F}_3, \mathcal{G})
$$
LHS 亦然。具体地，支撑条件保证一切均有定义；短正合列与长度的可加性给出
$$
[\mathcal{F}_2]_r = [\mathcal{F}_1]_r  + [\mathcal{F}_3]_r
$$
见《Chow 同调》引理 \ref{chow-lemma-additivity-sheaf-cycle}，这蕴含 RHS 的可加性。
$\text{Tor}$ 的长正合列
$$
\ldots \to \text{Tor}_1(\mathcal{F}_3, \mathcal{G}) \to
\text{Tor}_0(\mathcal{F}_1, \mathcal{G}) \to
\text{Tor}_0(\mathcal{F}_2, \mathcal{G}) \to
\text{Tor}_0(\mathcal{F}_3, \mathcal{G}) \to 0
$$
以及长度的可加性，如上蕴含 LHS 的可加性。

\medskip\noindent
由《代数》引理 \ref{algebra-lemma-filter-Noetherian-module} 及 $X$ 的仿射性，
可找到 $\mathcal{F}$ 的一个滤过，其分次片是 $\text{Supp}(\mathcal{F})$ 的闭子簇的结构层。
由上段所证的可加性，只需在以 $\mathcal{O}_V$ 取代 $\mathcal{F}$ 时证明
$LHS = RHS$，其中 $V \subset \text{Supp}(\mathcal{F})$。由对称性，可对 $\mathcal{G}$ 作同样处理。
因而问题化为证明
$$
LHS(\mathcal{O}_V, \mathcal{O}_W) = RHS(\mathcal{O}_V, \mathcal{O}_W)
$$
其中 $W \subset \text{Supp}(\mathcal{G})$ 是闭子簇。若 $\dim(V) = r$ 且
$\dim(W) = s$，则该等式就是 $V \cdot W$ 的{\bf 定义}。另一方面，若
$\dim(V) < r$ 或 $\dim(W) < s$，即 $[V]_r = 0$ 或 $[W]_s = 0$，
则需证明 $RHS(\mathcal{O}_V, \mathcal{O}_W) = 0$
\footnote{取 $r = s = 1$，令 $X$ 为非奇异曲面，而 $V = W$ 为 $X$ 的一个闭点 $x$，
读者即可看出这并非平凡事实。此时在 $x$ 处有 $3$ 个非零 $\text{Tor}$，其长度分别为
$1, 2, 1$。}。

\medskip\noindent
设 $Z \subset V \cap W$ 为维数 $r + s - \dim(X)$ 的不可约分支。这是一个分支可能的最大维数，
因而只需证明 $RHS$ 中 $Z$ 的系数为零。设 $\xi \in Z$ 为泛点。记
$A = \mathcal{O}_{X, \xi}$、$B = \mathcal{O}_{X \times X, \Delta(\xi)}$ 与
$C = \mathcal{O}_{V \times W, \Delta(\xi)}$。由引理 \ref{intersection-lemma-tor-and-diagonal}，有
$$
\text{下列对象的系数： }Z\text{ 于 }
RHS(\mathcal{O}_V, \mathcal{O}_W) = 
\sum (-1)^i
\text{length}_B \text{Tor}_i^B(A, C)
$$
由于 $\dim(V) < r$ 或 $\dim(W) < s$，有 $\dim(V \times W) < r + s$，这蕴含
$\dim(C) < \dim(X)$（略去一个小细节）。此外，$B \to A$ 的核 $I$ 由长度为
$\dim(X)$ 的正则序列生成；见引理 \ref{intersection-lemma-diagonal-regular-immersion}。
由引理 \ref{intersection-lemma-one-ideal-ci} 可得消失，因为根据《代数》命题
\ref{algebra-proposition-dimension}，$C$ 对 $I$ 的 Hilbert 函数的次数为
$\dim(C) < n$。
\end{proof}

\begin{remark}
\label{intersection-remark-Serre-conjectures}
设 $(A, \mathfrak m, \kappa)$ 为正则局部环，$M$ 与 $N$ 为非零有限 $A$-模，且
$M \otimes_A N$ 支撑于 $\{\mathfrak m\}$。则
$$
\chi(M, N) = \sum (-1)^i \text{length}_A \text{Tor}_i^A(M, N)
$$
是有限的。置 $r = \dim(\text{Supp}(M))$ 与 $s = \dim(\text{Supp}(N))$。
\cite{Serre_algebre_locale} 证明了 $r + s \leq \dim(A)$，并提出下列猜想：
\begin{enumerate}
\item 若 $r + s < \dim(A)$，则 $\chi(M, N) = 0$；以及
\item 若 $r + s = \dim(A)$，则 $\chi(M, N) > 0$。
\end{enumerate}
证明引理 \ref{intersection-lemma-tor-sheaf} 与命题 \ref{intersection-proposition-positivity} 的论证可用于证明：
若 $A$ 包含一个域，则 (1) 与 (2) 成立，正如 Serre 的文本中所做的那样。
目前，猜想 (1) 在一般情形下已知，且已知一般地有 $\chi(M, N) \geq 0$（Gabber）。
据我们所知，正性仍是开放问题。
\end{remark}



\section{交的结合性}
\label{intersection-section-associative}

\noindent
显然，上面定义的正常相交具有交换性。利用关键引理
\ref{intersection-lemma-tor-sheaf}，可以证明（正常）交积具有结合性。

\begin{lemma}
\label{intersection-lemma-associative}
设 $X$ 为非奇异簇，$U, V, W$ 为闭子簇。假设 $U, V, W$ 两两正常相交，且
$\dim(U \cap V \cap W) \leq \dim(U) + \dim(V) + \dim(W) - 2\dim(X)$。则作为 $X$ 上的循环，有
$$
U \cdot (V \cdot W) = (U \cdot V) \cdot W
$$
\end{lemma}

\begin{proof}
下面不再特别说明地使用引理 \ref{intersection-lemma-tor-sheaf}。它给出
\begin{align*}
V \cdot W
& =
\sum (-1)^i [\text{Tor}_i(\mathcal{O}_V, \mathcal{O}_W)]_{b + c - n} \\
U \cdot (V \cdot W)
& =
\sum (-1)^{i + j}
[
\text{Tor}_j(\mathcal{O}_U, \text{Tor}_i(\mathcal{O}_V, \mathcal{O}_W))
]_{a + b + c - 2n} \\
U \cdot V
& =
\sum (-1)^i [\text{Tor}_i(\mathcal{O}_U, \mathcal{O}_V)]_{a + b - n} \\
(U \cdot V) \cdot W
& =
\sum (-1)^{i + j}
[
\text{Tor}_j(\text{Tor}_i(\mathcal{O}_U, \mathcal{O}_V), \mathcal{O}_W))
]_{a + b + c - 2n}
\end{align*}
其中 $\dim(U) = a$、$\dim(V) = b$、$\dim(W) = c$、$\dim(X) = n$。
引理中的假设保证上述公式中的凝聚层满足所需的支撑维数界，因而这些公式有意义。
现在考虑导出范畴 $D_{\textit{Coh}}(\mathcal{O}_X)$ 中的对象
$$
K =
\mathcal{O}_U \otimes^\mathbf{L}_{\mathcal{O}_X} \mathcal{O}_V
\otimes^\mathbf{L}_{\mathcal{O}_X} \mathcal{O}_W
$$
我们断言，上面得到的 $U \cdot (V \cdot W)$ 与 $(U \cdot V) \cdot W$ 的表达式都等于
$$
\sum (-1)^k [H^k(K)]_{a + b + c - 2n}
$$
这将证明引理。由对称性，只需证明其中一个等式。为此，用 K-平坦复形
$M^\bullet$ 与 $L^\bullet$ 分别代表 $\mathcal{O}_U$ 与
$\mathcal{O}_V \otimes_{\mathcal{O}_X}^\mathbf{L} \mathcal{O}_W$，并使用《同调》第
\ref{homology-section-double-complex} 节中与双复形
$M^\bullet \otimes_{\mathcal{O}_X} L^\bullet$ 相关联的谱序列。该谱序列的 $E_2$ 页为
$$
E_2^{p, q} =
\text{Tor}_{-p}(\mathcal{O}_U, \text{Tor}_{-q}(\mathcal{O}_V, \mathcal{O}_W))
$$
并收敛到 $H^{p + q}(K)$（略去细节；可与《代数进阶》例
\ref{more-algebra-example-tor} 比较）。由于长度对短正合列具有可加性，结论成立。
\end{proof}




\section{平坦拉回与交积}
\label{intersection-section-flat-pullback-and-intersection-products}

\noindent
简要讨论交与平坦拉回之间的相互作用。

\begin{lemma}
\label{intersection-lemma-flat-pull-back-and-intersections-sheaves}
设 $f : X \to Y$ 为非奇异簇之间的平坦态射。置 $e = \dim(X) - \dim(Y)$。
设 $\mathcal{F}$ 与 $\mathcal{G}$ 为 $Y$ 上的凝聚层，并满足
$\dim(\text{Supp}(\mathcal{F})) \leq r$、
$\dim(\text{Supp}(\mathcal{G})) \leq s$ 以及
$\dim(\text{Supp}(\mathcal{F}) \cap \text{Supp}(\mathcal{G}) )
\leq r + s - \dim(Y)$。此时循环
$[f^*\mathcal{F}]_{r + e}$ 与 $[f^*\mathcal{G}]_{s + e}$ 正常相交，且
$$
f^*([\mathcal{F}]_r \cdot [\mathcal{G}]_s) =
[f^*\mathcal{F}]_{r + e} \cdot [f^*\mathcal{G}]_{s + e}
$$
\end{lemma}

\begin{proof}
由 $f$ 的相对维数为 $e$ 这一假设，$[f^*\mathcal{F}]_{r + e}$ 与
$[f^*\mathcal{G}]_{s + e}$ 正常相交的断言是立即的。由引理
\ref{intersection-lemma-tor-sheaf} 与 \ref{intersection-lemma-pullback}，只需证明
$$
f^*\text{Tor}_i^{\mathcal{O}_Y}(\mathcal{F}, \mathcal{G}) =
\text{Tor}_i^{\mathcal{O}_X}(f^*\mathcal{F}, f^*\mathcal{G})
$$
它们是 $\mathcal{O}_X$-模。这由《上同调》引理
\ref{cohomology-lemma-pullback-tensor-product} 以及 $f^*$ 的正合性得出；因此
$Lf^*\mathcal{F} = f^*\mathcal{F}$，对 $\mathcal{G}$ 亦然。
\end{proof}

\begin{lemma}
\label{intersection-lemma-flat-pullback-and-intersections}
设 $f : X \to Y$ 为非奇异簇之间的平坦态射。设 $\alpha$ 为 $Y$ 上的 $r$-循环，
$\beta$ 为 $Y$ 上的 $s$-循环。假设 $\alpha$ 与 $\beta$ 正常相交。则 $f^*\alpha$
与 $f^*\beta$ 正常相交，且 $f^*( \alpha \cdot \beta ) = f^*\alpha \cdot f^*\beta$。
\end{lemma}

\begin{proof}
由线性，可以假设 $\alpha = [V]$ 且 $\beta = [W]$，其中 $V, W \subset Y$ 是维数为 $r, s$
的闭子簇。设 $f$ 的相对维数为 $e$。则引理是引理
\ref{intersection-lemma-flat-pull-back-and-intersections-sheaves} 的特例，因为
$[V] = [\mathcal{O}_V]_r$、$[W] = [\mathcal{O}_W]_r$、
$f^*[V] = [f^{-1}(V)]_{r + e} = [f^*\mathcal{O}_V]_{r + e}$，以及
$f^*[W] = [f^{-1}(W)]_{s + e} = [f^*\mathcal{O}_W]_{s + e}$。
\end{proof}


\section{平坦固有态射的投影公式}
\label{intersection-section-projection-formula-flat}

\noindent
简要讨论平坦固有态射的投影公式。

\begin{lemma}
\label{intersection-lemma-projection-formula-flat}
\begin{reference}
更一般的公式见 \cite[Chapter V, C), Section 7, formula (10)]{Serre_algebre_locale}。
\end{reference}
设 $f : X \to Y$ 为非奇异簇之间的平坦固有态射，并置
$e = \dim(X) - \dim(Y)$。设 $\alpha$ 为 $X$ 上的 $r$-循环，$\beta$ 为 $Y$ 上的 $s$-循环。
假设 $\alpha$ 与 $f^*(\beta)$ 正常相交。则 $f_*(\alpha)$ 与 $\beta$ 正常相交，且
$$
f_*(\alpha) \cdot \beta = f_*( \alpha \cdot f^*\beta)
$$
\end{lemma}

\begin{proof}
由线性，化归到 $\alpha = [V]$ 且 $\beta = [W]$ 的情形，其中
$V \subset X$ 与 $W \subset Y$ 是分别为维数 $r$ 与 $s$ 的闭子簇。则
$f^{-1}(W)$ 具有纯维数 $s + e$。假设循环 $[V]$ 与 $f^*[W]$ 正常相交。
下文不再特别说明地使用事实：$V \cap f^{-1}(W) \to f(V) \cap W$ 是满的。

\medskip\noindent
设 $a$ 为 $V \to f(V)$ 的一般纤维的维数。若 $a > 0$，则 $f_*[V] = 0$。
特别地，$f_*\alpha$ 与 $\beta$ 正常相交。为完成这种情形，需证明
$f_*([V] \cdot f^*[W]) = 0$。然而，$V \to f(V)$ 的每个纤维维数都 $\geq a$；见《态射》引理
\ref{morphisms-lemma-openness-bounded-dimension-fibres}。所以 $V \cap f^{-1}(W)$ 的每个不可约分支 $Z$
在 $f(Z)$ 上的纤维维数都 $\geq a$。这当然蕴含所需结论。

\medskip\noindent
现假设 $V \to f(V)$ 在一般点上有限。设 $Z \subset f(V) \cap W$ 为不可约分支，而
$Z_i \subset V \cap f^{-1}(W)$、$i = 1, \ldots, t$ 为 $V \cap f^{-1}(W)$ 中支配 $Z$ 的不可约分支。
根据假设，每个 $Z_i$ 的维数为
$r + s + e - \dim(X) = r + s - \dim(Y)$。因此
$\dim(Z) \leq r + s - \dim(Y)$。所以 $f(V)$ 与 $W$ 正常相交，
$\dim(Z) = r + s - \dim(Y)$，且每个 $Z_i \to Z$ 在一般点上有限。
特别地，$V \to f(V)$ 在 $Z$ 的泛点 $\xi$ 之上有有限纤维。因此 $V \to Y$
在 $\xi$ 的一个开邻域中有限；见《概形的上同调》引理
\ref{coherent-lemma-proper-finite-fibre-finite-in-neighbourhood}。
使用导出张量积的一个非常一般的投影公式，得到
$$
Rf_*(\mathcal{O}_V \otimes_{\mathcal{O}_X}^\mathbf{L} Lf^*\mathcal{O}_W) =
Rf_*\mathcal{O}_V \otimes_{\mathcal{O}_Y}^\mathbf{L} \mathcal{O}_W
$$
见《概形的导出范畴》引理 \ref{perfect-lemma-cohomology-base-change}。
由于 $f$ 平坦，有 $Lf^*\mathcal{O}_W = f^*\mathcal{O}_W$。由于 $f|_V$ 在 $\xi$ 的一个开邻域中有限，
对任意支撑包含于 $V$ 的 $X$ 上凝聚层，有
$$
(Rf_*\mathcal{F})_\xi = (f_*\mathcal{F})_\xi
$$
见《概形的上同调》引理
\ref{coherent-lemma-higher-direct-images-zero-finite-fibre}。所以对所有 $i$，有
\begin{equation}
\label{intersection-equation-stalks}
\left(
f_*\text{Tor}_i^{\mathcal{O}_X}(\mathcal{O}_V, f^*\mathcal{O}_W)
\right)_\xi =
\left(\text{Tor}_i^{\mathcal{O}_Y}(f_*\mathcal{O}_V, \mathcal{O}_W)\right)_\xi
\end{equation}
由引理 \ref{intersection-lemma-pullback}，$f^*[W] = [f^*\mathcal{O}_W]_{s + e}$，因而由引理
\ref{intersection-lemma-tor-sheaf}，有
$$
[V] \cdot f^*[W] =
\sum (-1)^i
[\text{Tor}_i^{\mathcal{O}_X}(\mathcal{O}_V,
f^*\mathcal{O}_W)]_{r + s - \dim(Y)}
$$
应用引理 \ref{intersection-lemma-push-coherent} 得
$$
f_*([V] \cdot f^*[W]) =
\sum (-1)^i
[f_*\text{Tor}_i^{\mathcal{O}_X}(\mathcal{O}_V,
f^*\mathcal{O}_W)]_{r + s - \dim(Y)}
$$
由引理 \ref{intersection-lemma-push-coherent}，$f_*[V] = [f_*\mathcal{O}_V]_r$，因而再次由引理
\ref{intersection-lemma-tor-sheaf}，有
$$
[f_*V] \cdot [W] =
\sum (-1)^i
[\text{Tor}_i^{\mathcal{O}_X}(f_*\mathcal{O}_V,
\mathcal{O}_W)]_{r + s - \dim(Y)}
$$
比较 $f_*([V] \cdot f^*[W])$ 与 $f_*[V] \cdot [W]$ 的公式，并通过取在 $\xi$ 处的茎的长度
来考察 $Z$ 的系数，可见等式 (\ref{intersection-equation-stalks}) 完成了证明。
\end{proof}

\begin{lemma}
\label{intersection-lemma-transfer}
设 $X \to P$ 为非奇异簇之间的闭浸入，$C' \subset P \times \mathbf{P}^1$ 为维数
$r + 1$ 的闭子簇。假设
\begin{enumerate}
\item 纤维 $C = C'_0$ 的维数为 $r$，即 $C' \to \mathbf{P}^1$ 是支配的；
\item $C'$ 与 $X \times \mathbf{P}^1$ 正常相交；
\item $[C]_r$ 与 $X$ 正常相交。
\end{enumerate}
把 $\alpha = [C]_r \cdot X$ 视为 $X$ 上的循环，并把
$\beta = C' \cdot X \times \mathbf{P}^1$ 视为 $X \times \mathbf{P}^1$ 上的循环。则作为 $X$ 上的循环，有
$$
\alpha = \text{pr}_{X, *}(\beta \cdot X \times 0)
$$
其中 $\text{pr}_X : X \times \mathbf{P}^1 \to X$ 是投影。
\end{lemma}

\begin{proof}
设 $\text{pr} : P \times \mathbf{P}^1 \to P$ 为投影。由于我们证明的是循环的等式，
只需把 $\alpha$、\ $\beta$ 分别看作 $P$、\ $P \times \mathbf{P}^1$ 上的循环，并对沿
$\text{pr}$ 的推前证明结论。因为
$\text{pr}^*X = X \times \mathbf{P}^1$，且 $\text{pr}$ 定义了 $C'_0$ 到 $C$ 的同构，
投影公式（引理 \ref{intersection-lemma-projection-formula-flat}）给出
$$
\text{pr}_*([C'_0]_r \cdot X \times \mathbf{P}^1) = [C]_r \cdot X = \alpha
$$
另一方面，由引理 \ref{intersection-lemma-rational-equivalence-and-intersection}，
作为 $P \times \mathbf{P}^1$ 上的循环，有 $[C'_0]_r = C' \cdot P \times 0$。因此
$$
[C'_0]_r \cdot X \times \mathbf{P}^1 =
(C' \cdot P \times 0) \cdot X \times \mathbf{P}^1 =
(C' \cdot X \times \mathbf{P}^1) \cdot P \times 0
$$
这里使用了结合性（引理 \ref{intersection-lemma-associative}）以及交积的交换性。
仍需证明：在 $P \times \mathbf{P}^1$ 上，$C' \cdot X \times \mathbf{P}^1$ 与 $P \times 0$
的交积，作为循环，等于在 $X \times \mathbf{P}^1$ 上 $\beta$ 与 $X \times 0$ 的交积。写成
$C' \cdot X \times \mathbf{P}^1 = \sum n_k[E_k]$，从而
$\beta = \sum n_k[E_k]$，其中 $E_k \subset X \times \mathbf{P}^1 \subset P \times \mathbf{P}^1$ 为子簇。
两次应用引理 \ref{intersection-lemma-rational-equivalence-and-intersection}，可知两个交都等于
$\sum m_k[E_{k, 0}]$。证明完毕。
\end{proof}




\section{投影}
\label{intersection-section-projection}

\noindent
回顾一下，我们始终在一个固定的代数闭基域 $\mathbf{C}$ 上工作。若 $V$ 是
$\mathbf{C}$ 上的有限维向量空间，则置
$$
\mathbf{P}(V) = \text{Proj}(\text{Sym}(V))
$$
其中 $\text{Sym}(V)$ 是 $V$ 在 $\mathbf{C}$ 上的对称代数。见《构造》例
\ref{constructions-example-projective-space}。这里采用的规范化使得
$V = \Gamma(\mathbf{P}(V), \mathcal{O}_{\mathbf{P}(V)}(1))$。
当然，若 $\dim(V) = n + 1$，则 $\mathbf{P}(V) \cong \mathbf{P}^n_{\mathbf{C}}$。
注意，$\mathbf{P}(V)$ 是非奇异射影簇。

\medskip\noindent
设 $p \in \mathbf{P}(V)$ 为闭点。点 $p$ 对应于向量空间的一个满射
$V \to L_p$，其中 $\dim(L_p) = 1$；见《构造》引理
\ref{constructions-lemma-apply}。记 $W_p = \Ker(V \to L_p)$。
{\it 以 $p$ 为中心的投影}是《构造》引理
\ref{constructions-lemma-morphism-proj} 中的态射
$$
r_p : \mathbf{P}(V) \setminus \{p\} \longrightarrow \mathbf{P}(W_p)
$$
先看下面这个预备引理。

\begin{lemma}
\label{intersection-lemma-projection-generically-finite}
设 $V$ 是维数为 $n + 1$ 的向量空间，$X \subset \mathbf{P}(V)$ 是闭子概形。
若 $X \not = \mathbf{P}(V)$，则存在非空 Zariski 开集
$U \subset \mathbf{P}(V)$，使得对每个闭点 $p \in U$，投影 $r_p$ 的限制定义有限态射
$r_p|_X : X \to \mathbf{P}(W_p)$。
\end{lemma}

\begin{proof}
我们断言取 $U = \mathbf{P}(V) \setminus X$ 即可。对 $U$ 的闭点 $p$，确实得到态射
$r_p|_X : X \to \mathbf{P}(W_p)$。由于 $X$ 是固有概形，这个态射是固有的
（《态射》引理 \ref{morphisms-lemma-locally-projective-proper} 和
\ref{morphisms-lemma-image-proper-scheme-closed}）。另一方面，直接计算可见 $r_p$ 的纤维是仿射直线。
因此 $r_p|X$ 的纤维既固有又仿射，从而是有限的
（《态射》引理 \ref{morphisms-lemma-finite-proper}）。最后，具有有限纤维的固有态射是有限态射
（《概形的上同调》引理 \ref{coherent-lemma-characterize-finite}）。
\end{proof}

\begin{lemma}
\label{intersection-lemma-projection-generically-immersion}
设 $V$ 是维数为 $n + 1$ 的向量空间，$X \subset \mathbf{P}(V)$ 是闭子簇，
$x \in X$ 是非奇异点。
\begin{enumerate}
\item 若 $\dim(X) < n - 1$，则存在非空 Zariski 开集
$U \subset \mathbf{P}(V)$，使得对每个闭点 $p \in U$，态射
$r_p|_X : X \to r_p(X)$ 在 $r_p(x)$ 的某个开邻域上是同构。
\item 若 $\dim(X) = n - 1$，则存在非空 Zariski 开集
$U \subset \mathbf{P}(V)$，使得对每个闭点 $p \in U$，态射
$r_p|_X : X \to \mathbf{P}(W_p)$ 在 $x$ 处为 \'etale 态射。
\end{enumerate}
\end{lemma}

\begin{proof}
先证明 (1)。注意，若 $x, y \in X$ 在 $r_p$ 下有相同的像，则 $p$ 位于直线
$\overline{xy}$ 上。考虑有限型概形
$$
T = \{(y, p) \mid
y \in X \setminus \{x\},\ p \in \mathbf{P}(V),\ p \in \overline{xy}\}
$$
以及由 $(y, p) \mapsto y$ 和 $(y, p) \mapsto p$ 给出的态射
$T \to X$ 和 $T \to \mathbf{P}(V)$。由于 $T \to X$ 的每个纤维都是一条直线，
可见 $T$ 的维数为 $\dim(X) + 1 < \dim(\mathbf{P}(V))$。因此
$T \to \mathbf{P}(V)$ 不是满射。另一方面，考虑有限型概形
$$
T' = \{p \mid
p \in \mathbf{P}(V) \setminus \{x\},
\ \overline{xp}\text{ 相切于 }X\text{ 在 }x\}
$$
此时 $T'$ 的维数为 $\dim(X) < \dim(\mathbf{P}(V))$，所以态射
$T' \to \mathbf{P}(V)$ 也不是满射。取非空开集
$U \subset \mathbf{P}(V) \setminus X$，使它与上述两个像均不相交。这样的 $U$ 存在，
因为由《态射》引理 \ref{morphisms-lemma-chevalley}，$T$ 和 $T'$ 在
$\mathbf{P}(V)$ 中的像都是可构造集。于是对闭点 $p \in U$，投影
$r_p|_X : X \to \mathbf{P}(W_p)$ 在 $x$ 的切空间上是单射，并且
$r_p^{-1}(\{r_p(x)\}) = \{x\}$。这说明 $r_p$ 在 $x$ 处非分歧
（《簇》引理 \ref{varieties-lemma-injective-tangent-spaces-unramified}）；
由引理 \ref{intersection-lemma-projection-generically-finite}，它是有限态射；又有
$r_p^{-1}(\{r_p(x)\}) = \{x\}$。因此可应用《\'Etale 态射》引理
\ref{etale-lemma-finite-unramified-one-point}：在 $\mathbf{P}(W_p)$ 中存在
$r_p(x)$ 的开邻域 $R$，使得 $(r_p|_X)^{-1}(R) \to R$ 是闭浸入。这证明了 (1)。

\medskip\noindent
再证明 (2)。在这种情形下，我们仍可断定态射
$T' \to \mathbf{P}(V)$ 不是满射。与上面相同地论证可知，若
$U \subset \mathbf{P}(V)$ 避开 $X$ 和 $T'$ 的像，则投影
$r_p|_X : X \to \mathbf{P}(W_p)$ 在 $x$ 处为 \'etale 态射且为有限态射。
\end{proof}

\begin{lemma}
\label{intersection-lemma-projection-generically-immersion-refined}
设 $V$ 是向量空间，$Z \subset X \subset \mathbf{P}(V)$ 是闭子簇，且
$Z \not = X \not = \mathbf{P}(V)$。设 $x \in Z$ 是 $X$ 上的非奇异点。
则存在非空 Zariski 开集 $U \subset \mathbf{P}(V)$，使得对每个闭点 $p \in U$，态射
$r_p|_Z : Z \to r_p(Z)$ 在 $r_p(x)$ 的某个开邻域上是同构。
\end{lemma}

\begin{proof}
仿照引理 \ref{intersection-lemma-projection-generically-immersion} 的证明，可找到非空 Zariski 开集
$U$，使得对 $p \in U$，映射 $r_p|_X : X \to r_p(X)$ 在 $x$ 处为 \'etale 态射，且
$r_p^{-1}(r_p(\{x\})) \cap Z = \{x\}$。细节从略。令 $Z' = r_p(Z)$，把它看作
$\mathbf{P}(W_p)$ 的闭子簇。则态射 $\pi : Z \to Z'$ 是有限的
（引理 \ref{intersection-lemma-projection-generically-finite}），在 $x$ 处非分歧（略去一个小细节），并且
$\pi^{-1}(\pi(\{x\})) = \{x\}$。结论由《\'Etale 态射》引理
\ref{etale-lemma-finite-unramified-one-point} 得出。
\end{proof}

\begin{lemma}
\label{intersection-lemma-projection-injective}
设 $V$ 是维数为 $n + 1$ 的向量空间，$Y, Z \subset \mathbf{P}(V)$ 是闭子簇。
则存在非空 Zariski 开集 $U \subset \mathbf{P}(V)$，使得对每个闭点 $p \in U$，有
$$
Y \cap r_p^{-1}(r_p(Z)) = (Y \cap Z) \cup E
$$
其中 $E \subset Y$ 是闭集，并且
$\dim(E) \leq \dim(Y) + \dim(Z) + 1 - n$。
\end{lemma}

\begin{proof}
令 $Y' = Y \setminus Y \cap Z$。设 $y \in Y'$、$z \in Z$ 是满足
$r_p(y) = r_p(z)$ 的闭点。于是 $p$ 位于通过 $y$ 和 $z$ 的直线 $\overline{yz}$ 上。
考虑有限型概形
$$
T = \{(y, z, p) \mid y \in Y', z \in Z, p \in \overline{yz}\}
$$
以及由 $(y, z, p) \mapsto p$ 给出的态射 $T \to \mathbf{P}(V)$。
注意，$T$ 不可约，且 $\dim(T) = \dim(Y) + \dim(Z) + 1$。
因此 $T \to \mathbf{P}(V)$ 的一般纤维维数至多为
$\dim(Y) + \dim(Z) + 1 - n$。更精确地说，存在非空开集
$U \subset \mathbf{P}(V) \setminus (Y \cup Z)$，使得其上的纤维维数至多为
$\dim(Y) + \dim(Z) + 1 - n$
（《簇》引理 \ref{varieties-lemma-dimension-fibres-locally-algebraic}）。
设 $p \in U$ 为闭点，$F \subset T$ 是 $T \to \mathbf{P}(V)$ 在 $p$ 上的纤维。则
$$
(Y \cap r_p^{-1}(r_p(Z))) \setminus (Y \cap Z)
$$
是 $F \to Y$ 在映射 $(y, z, p) \mapsto y$ 下的像。再次由《簇》引理
\ref{varieties-lemma-dimension-fibres-locally-algebraic}，$F \to Y$ 的像的闭包维数至多为
$\dim(Y) + \dim(Z) + 1 - n$。
\end{proof}

\begin{lemma}
\label{intersection-lemma-find-lines}
设 $V$ 是向量空间，$B \subset \mathbf{P}(V)$ 是余维 $\geq 2$ 的闭子簇。
设 $p \in \mathbf{P}(V)$ 是闭点，且 $p \not \in B$。则存在直线
$\ell \subset \mathbf{P}(V)$，使得 $p \in \ell$ 且 $\ell \cap B = \emptyset$。
此外，这些直线扫过 $\mathbf{P}(V)$ 的一个开子集。
\end{lemma}

\begin{proof}
考虑 $B$ 在投影 $r_p : \mathbf{P}(V) \to \mathbf{P}(W_p)$ 下的像。
由于 $\dim(W_p) = \dim(V) - 1$，可见 $r_p(B)$ 在 $\mathbf{P}(W_p)$ 中的余维
$\geq 1$。对任意满足 $r_p(q) \not \in r_p(B)$ 的 $q \in \mathbf{P}(V)$，
连接 $p$ 与 $q$ 的直线 $\ell = \overline{pq}$ 都满足要求。
\end{proof}

\begin{lemma}
\label{intersection-lemma-doubly-transitive}
设 $V$ 是向量空间，$G = \text{PGL}(V)$。则
$G \times \mathbf{P}(V) \to \mathbf{P}(V)$ 是双重传递的。
\end{lemma}

\begin{proof}
从略。提示：这可由 $\text{GL}(V)$ 在所有线性无关向量对上的作用是双重传递的这一事实推出。
\end{proof}

\begin{lemma}
\label{intersection-lemma-determinant}
设 $k$ 是域，$n \geq 1$ 是整数，$x_{ij}, 1 \leq i, j \leq n$ 是变元。则
$$
\det
\left(
\begin{matrix}
x_{11} & x_{12} & \ldots & x_{1n} \\
x_{21} & \ldots & \ldots & \ldots \\
\ldots & \ldots & \ldots & \ldots \\
x_{n1} & \ldots & \ldots & x_{nn}
\end{matrix}
\right)
$$
是多项式环 $k[x_{ij}]$ 中的不可约元。
\end{lemma}

\begin{proof}
设 $V$ 是 $n$ 维向量空间。将引理翻译成几何语言，它断言不可逆线性映射
$V \to V$ 的轨迹 $C$ 不可约（注意，由于行列式关于每个变元的次数都是 $1$，它是无平方因子的）。
设 $W$ 是维数为 $n - 1$ 的向量空间。由初等线性代数，态射
$$
\Hom(W, V) \times \Hom(V, W) \longrightarrow \Hom(V, V),\quad
(\psi, \varphi) \longmapsto \psi \circ \varphi
$$
的像为 $C$。由于源不可约，像也不可约。细节从略。
\end{proof}

\noindent
设 $V$ 是维数为 $n + 1$ 的向量空间。令 $E = \text{End}(V)$，并令
$E^\vee = \Hom(E, \mathbf{C})$ 为其对偶向量空间。记
$\mathbf{P} = \mathbf{P}(E^\vee)$。存在典范线性映射
$$
V \longrightarrow V \otimes_\mathbf{C} E^\vee = \Hom(E, V)
$$
它把 $v \in V$ 送到 $\Hom(E, V)$ 中的映射 $g \mapsto g(v)$。
回顾，存在典范映射
$E^\vee \to \Gamma(\mathbf{P}, \mathcal{O}_\mathbf{P}(1))$，且它是同构。
因此得到 $\mathbf{P}$ 上模层的典范映射
$$
\psi : V \otimes \mathcal{O}_\mathbf{P} \to V \otimes \mathcal{O}_\mathbf{P}(1)
$$
它在整体截面上恢复上述映射。回顾，射影丛 $\mathbf{P}(\mathcal{E})$ 定义为
$\mathcal{E}$ 的对称代数的相对 Proj；见《构造》定义
\ref{constructions-definition-projective-bundle}。我们要研究与 $\psi$ 相伴的、
$\mathbf{P}(V \otimes \mathcal{O}_\mathbf{P}(1))$ 与
$\mathbf{P}(V \otimes \mathcal{O}_\mathbf{P})$ 之间的有理映射。
由《构造》引理 \ref{constructions-lemma-relative-proj-base-change}，存在典范同构
$$
\mathbf{P}(V \otimes \mathcal{O}_\mathbf{P}) = \mathbf{P} \times \mathbf{P}(V)
$$
由《构造》引理 \ref{constructions-lemma-twisting-and-proj}，可见
$$
\mathbf{P}(V \otimes \mathcal{O}_\mathbf{P}(1)) =
\mathbf{P}(V \otimes \mathcal{O}_\mathbf{P}) = \mathbf{P} \times \mathbf{P}(V)
$$
再结合《构造》引理 \ref{constructions-lemma-morphism-relative-proj}，得到
\begin{equation}
\label{intersection-equation-r-psi}
\mathbf{P} \times \mathbf{P}(V) \supset
U(\psi) \xrightarrow{r_\psi} \mathbf{P} \times \mathbf{P}(V)
\end{equation}
为了更好地理解它，我们考察 $\mathbf{P}$ 上各纤维中的情形。设 $g \in E$ 非零。
它定义非零映射 $E^\vee \to \mathbf{C}$，从而定义一点 $[g] \in \mathbf{P}$。
另一方面，$g$ 定义 $\mathbf{C}$-线性映射 $g : V \to V$。因此由《构造》引理
\ref{constructions-lemma-morphism-proj} 得到映射
$$
\mathbf{P}(V) \supset U(g) \xrightarrow{r_g} \mathbf{P}(V)
$$
下面将用到：$U(g)$ 是纤维 $U(\psi)_{[g]}$，而 $r_g$ 是 $r_\psi$ 在点 $[g]$ 上的纤维。
我们还将用到另一个观察：$U(g)$ 在 $\mathbf{P}(V)$ 中的补集是闭浸入
$$
\mathbf{P}(\Coker(g)) \longrightarrow \mathbf{P}(V)
$$
的像，而 $r_g$ 的像是闭浸入
$$
\mathbf{P}(\Im(g)) \longrightarrow \mathbf{P}(V)
$$
的像。

\begin{lemma}
\label{intersection-lemma-make-family}
沿用上述记号。设 $X, Y$ 是 $\mathbf{P}(V)$ 中正常相交的闭子簇，并且
$X \not = \mathbf{P}(V)$ 且 $X \cap Y \not = \emptyset$。对满足
$[\text{id}_V] \in \ell$ 的一般直线 $\ell \subset \mathbf{P}$，有
\begin{enumerate}
\item 对所有 $[g] \in \ell$，都有 $X \subset U_g$；
\item 对所有 $[g] \in \ell$，$g(X)$ 与 $Y$ 正常相交。
\end{enumerate}
\end{lemma}

\begin{proof}
令 $B \subset \mathbf{P}$ 为“坏”点的集合，即违反 (1) 或 (2) 中任一项的点 $[g]$。
注意，由假设 $[\text{id}_V] \not \in B$。此外，$B$ 是闭集。因此只需证明
$\dim(B) \leq \dim(\mathbf{P}) - 2$（引理 \ref{intersection-lemma-find-lines}）。

\medskip\noindent
首先，考虑由使 $g : V \to V$ 可逆的点 $[g]$ 组成的开集
$G = \text{PGL}(V) \subset \mathbf{P}$。由于 $G$ 在 $\mathbf{P}(V)$ 上双重传递
（引理 \ref{intersection-lemma-doubly-transitive}），可见
$$
T = \{(x, y, [g]) \mid x \in X, y \in Y, [g] \in G, r_g(x) = y\}
$$
是 $X \times Y$ 上的局部平凡纤维化，其纤维等于 $G$ 中一点的稳定子。因此 $T$ 是簇。
注意，$T \to G$ 在 $[g]$ 上的纤维是 $r_g(X) \cap Y$。态射 $T \to G$ 是满射，
因为 $X$ 的任意平移都与 $Y$ 相交（事实上，由 $X$ 与 $Y$ 正常相交且
$X \cap Y \not = \emptyset$，可知
$\dim(X) + \dim(Y) \geq \dim(\mathbf{P}(V))$；于是《簇》引理
\ref{varieties-lemma-intersection-in-projective-space} 蕴含 $X$ 的所有平移都与 $Y$ 相交）。
由于簇之间支配态射的纤维维数不会在余维 $1$ 处跳跃
（《簇》引理 \ref{varieties-lemma-dimension-fibres-locally-algebraic}），
我们得到 $B \cap G$ 的余维 $\geq 2$。

\medskip\noindent
接着考察补集 $Z = \mathbf{P} \setminus G$。由于行列式是不可约多项式
（引理 \ref{intersection-lemma-determinant}），它是不可约簇。因此只需证明 $B$ 不包含 $Z$ 的泛点。
对一般点 $[g] \in Z$，余核 $V \to \Coker(g)$ 的维数为 $1$，所以 $U(g)$ 是一点的补集。
由于 $X \not = \mathbf{P}(V)$，可见对一般的 $[g] \in Z$ 有 $X \subset U(g)$。
此外，态射 $r_g|_X : X \to r_g(X)$ 是有限的，因而
$\dim(r_g(X)) = \dim(X)$。另一方面，对这样的 $g$，$r_g$ 的像是余维为 $1$ 的闭子空间
$H = \mathbf{P}(\Im(g)) \subset \mathbf{P}(V)$。对 $Z$ 的一般点，
$H \cap Y$ 的维数比 $Y$ 小 $1$（比较《簇》引理
\ref{varieties-lemma-exact-sequence-induction}）。所以我们需要证明 $r_g(X)$ 与
$H \cap Y$ 在 $H$ 中正常相交。固定 $H$ 后，通过在 $g$ 后复合一个自同构，
可以用 $H = \mathbf{P}(\Im(g))$ 的任意自同构移动 $r_g(X)$。因此可仿照上面的论证得出：
对满足给定 $H = \mathbf{P}(\Im(g))$ 的一般 $g$，$H \cap Y$ 与 $r_g(X)$ 正常相交。
略去一些细节。
\end{proof}







\section{移动引理}
\label{intersection-section-moving-lemma}

\noindent
移动引理断言：给定一个 $r$-循环 $\alpha$ 和一个 $s$-循环 $\beta$，存在
$\alpha'$，满足 $\alpha' \sim_{rat} \alpha$，且 $\alpha'$ 与 $\beta$ 正常相交
（引理 \ref{intersection-lemma-moving-move}）。见 \cite{Samuel}、\cite{ChevalleyI}、
\cite{ChevalleyII}。其中的关键是引理 \ref{intersection-lemma-moving}；读者可在
\cite[Example 11.4.1]{F} 中找到该引理的这一表述，并在 \cite{Roberts} 中找到证明。

\begin{lemma}
\label{intersection-lemma-moving}
\begin{reference}
见 \cite{Roberts}。
\end{reference}
设 $X \subset \mathbf{P}^N$ 是非奇异闭子簇。令 $n = \dim(X)$，且
$0 \leq d, d' < n$。设 $Z \subset X$ 是维数为 $d$ 的闭子簇，并设
$T_i \subset X$（$i \in I$）是有限多个维数为 $d'$ 的闭子簇。则存在子簇
$C \subset \mathbf{P}^N$，$C$ 与 $X$ 正常相交，并且
$$
C \cdot X = Z + \sum\nolimits_{j \in J} m_j Z_j
$$
其中 $Z_j \subset X$ 是维数为 $d$、彼此不同于 $Z$ 的不可约子簇，并且
$$
\dim(Z_j \cap T_i) \leq \dim(Z \cap T_i)
$$
若 $Z$ 与 $T_i$ 在 $X$ 中不正常相交，则此不等式严格成立。
\end{lemma}

\begin{proof}
写成 $\mathbf{P}^N = \mathbf{P}(V_N)$，于是 $\dim(V_N) = N + 1$。
我们要依次选取以下以点为中心的投影：
\begin{align*}
& r_N : 
\mathbf{P}(V_N) \setminus \{p_N\} \to \mathbf{P}(V_{N - 1}), \\
& r_{N - 1} :
\mathbf{P}(V_{N - 1}) \setminus \{p_{N - 1}\} \to \mathbf{P}(V_{N - 2}), \\
& \ldots,\\
& r_{n + 1} :
\mathbf{P}(V_{n + 1}) \setminus \{p_{n + 1}\} \to \mathbf{P}(V_n)
\end{align*}
如第 \ref{intersection-section-projection} 节所述。在每一步，我们都从一个适当的 Zariski 开集中选取
$p_N, p_{N - 1}, \ldots, p_{n + 1}$。选取闭点 $x \in Z \subset X$。
对每个 $i$，选取闭点 $x_{it} \in T_i \cap Z$，并在 $T_i \cap Z$ 的每个不可约分支中至少选取一个。
取这些投影的复合，得到态射
$$
\pi = (r_{n + 1} \circ \ldots \circ r_N)|_X :
X \longrightarrow \mathbf{P}(V_n)
$$
它具有以下性质：
\begin{enumerate}
\item $\pi$ 是有限态射；
\item $\pi$ 在 $x$ 以及所有 $x_{it}$ 处为 \'etale 态射；
\item $\pi|_Z : Z \to \pi(Z)$ 在 $\pi(x_{it})$ 的某个开邻域上是同构；
\item $T_i \cap \pi^{-1}(\pi(Z)) = (T_i \cap Z) \cup E_i$，其中
$E_i \subset T_i$ 是闭集，并且
$\dim(E_i) \leq d + d' + 1 - (n + 1) = d + d' - n$。
\end{enumerate}
由引理 \ref{intersection-lemma-projection-generically-finite}、
\ref{intersection-lemma-projection-generically-immersion}、
\ref{intersection-lemma-projection-generically-immersion-refined}、
\ref{intersection-lemma-projection-injective} 以及归纳法，直接可知可以如此选取。
注意，最后一次投影以 $\mathbf{P}(V_{n + 1})$ 中的点为中心，且
$\dim(V_{n + 1}) = n + 2$，这解释了 (4) 中的不等式。稍作展开，令
$X_j, Z_j, T_{i, j} \subset \mathbf{P}(V_j)$ 为 $X$、$Z$、$T_i$ 的像，其中
$j = N, \ldots, n$。则
$T_{i, j + 1} \cap r_{j + 1}^{-1}(Z_j) = (T_{i, j + 1} \cap Z_{j + 1})
\cup E_{i, j + 1}$，而引理 \ref{intersection-lemma-projection-injective} 给出了
$E_{i, j + 1}$ 的维数界。再注意，$T_i \cap \pi^{-1}(\pi(Z))$ 是
$T_i \cap Z$ 与各 $E_{i, j}$ 的逆像之并。由于各 $X_j$ 之间的过渡态射都是有限态射，
便得到所需的维数界。

\medskip\noindent
令 $C \subset \mathbf{P}(V_N)$ 为
$(r_{n + 1} \circ \ldots \circ r_N)^{-1}(\pi(Z))$ 的概形论闭包。
由于 $\pi$ 在 $Z$ 的点 $x$ 处为 \'etale 态射，可见闭子概形
$C \cap X = \pi^{-1}(\pi(Z))$ 以重数 $1$ 包含 $Z$（略去局部计算）。
因此由引理 \ref{intersection-lemma-transversal-subschemes}，得到
$$
C \cdot X = [Z] + \sum m_j[Z_j]
$$
其中 $Z_j \subset X$ 是一些维数为 $d$ 的子簇。注意，在集合论意义下
$$
C \cap X = \pi^{-1}(\pi(Z))
$$
因此
$T_i \cap Z_j \subset T_i \cap \pi^{-1}(\pi(Z)) \subset (T_i \cap Z) \cup E_i$。
对 $T_i \cap Z_j$ 中包含于 $E_i$ 的任一不可约分支，已经有了所需的维数界。
最后，设 $V$ 是 $T_i \cap Z_j$ 中包含于 $T_i \cap Z$ 的一个不可约分支。
为了完成证明，只需说明 $V$ 不包含任何点 $x_{it}$，因为这时
$\dim(V) < \dim(Z \cap T_i)$。为此，只需对所有 $i, t, j$ 证明
$x_{it} \not \in Z_j$。

\medskip\noindent
令 $Z' = \pi(Z)$，并在概形论意义下令 $Z'' = \pi^{-1}(Z')$。
由条件 (3)，可找到包含 $\pi(x_{it})$ 的开集 $U \subset \mathbf{P}(V_n)$，使得
$\pi^{-1}(U) \cap Z \to U \cap Z'$ 是同构。特别地，$Z \to Z'$ 在 $x_{it}$ 处是局部同构。
另一方面，由条件 (2)，$Z'' \to Z'$ 在 $x_{it}$ 处为 \'etale 态射。因此闭浸入
$Z \to Z''$ 在 $x_{it}$ 处为 \'etale 态射
（《态射》引理 \ref{morphisms-lemma-etale-permanence}）。
所以在 $x_{it}$ 的一个 Zariski 邻域内有 $Z = Z''$，这就证明了断言。
\end{proof}

\noindent
实际的移动由下面的引理完成。

\begin{lemma}
\label{intersection-lemma-move}
设 $C \subset \mathbf{P}^N$ 是闭子簇，$X \subset \mathbf{P}^N$ 是子簇，且
$T_i \subset X$ 是有限多个闭子簇。假设 $C$ 与 $X$ 正常相交。则存在闭子簇
$C' \subset \mathbf{P}^N \times \mathbf{P}^1$，使得
\begin{enumerate}
\item $C' \to \mathbf{P}^1$ 是支配的；
\item 在概形论意义下 $C'_0 = C$；
\item $C'$ 与 $X \times \mathbf{P}^1$ 正常相交；
\item $C'_\infty$ 与给定的每个 $T_i$ 正常相交。
\end{enumerate}
\end{lemma}

\begin{proof}
若 $C \cap X = \emptyset$，则取常值族 $C' = C \times \mathbf{P}^1$。
所以不妨并且确实假设 $C \cap X \not = \emptyset$。

\medskip\noindent
写成 $\mathbf{P}^N = \mathbf{P}(V)$，于是 $\dim(V) = N + 1$。令
$E = \text{End}(V)$，$E^\vee = \Hom(E, \mathbf{C})$，并如引理
\ref{intersection-lemma-make-family} 中那样置 $\mathbf{P} = \mathbf{P}(E^\vee)$。
选取通过 $\text{id}_V$ 的一般直线 $\ell \subset \mathbf{P}$。
令 $C' \subset \ell \times \mathbf{P}(V)$ 为在 $[g] \in \ell$ 上具有纤维
$r_g(C)$ 的闭子概形。更精确地说，$C'$ 是
$$
\ell \times C \subset \mathbf{P} \times \mathbf{P}(V)
$$
在态射 (\ref{intersection-equation-r-psi}) 下的像。由引理 \ref{intersection-lemma-make-family}，
这是有意义的，即 $\ell \times C \subset U(\psi)$。态射
$\ell \times C \to C'$ 是有限的，并且对所有 $[g] \in \ell$，在集合论意义下有
$C'_{[g]} = r_g(C)$。取 $0 = [\text{id}_V] \in \ell$ 时，(1) 与 (2) 显然成立。
(3) 来自如下事实：对所有 $[g] \in \ell$，$r_g(C)$ 与 $X$ 正常相交。
(4) 来自如下事实：一般点 $\infty = [g] \in \ell$ 是 $\mathbf{P}$ 的一般点，
而对这样的点以及 $\mathbf{P}(V)$ 的任意闭子簇 $T$，$r_g(C) \cap T$ 都是正常相交。
细节从略。
\end{proof}

\begin{lemma}
\label{intersection-lemma-moving-move}
设 $X$ 是非奇异射影簇，$\alpha$ 是一个 $r$-循环，$\beta$ 是 $X$ 上的
$s$-循环。则存在 $r$-循环 $\alpha'$，使得 $\alpha' \sim_{rat} \alpha$，且
$\alpha'$ 与 $\beta$ 正常相交。
\end{lemma}

\begin{proof}
对某些维数为 $s$ 的子簇 $T_i \subset X$，写成 $\beta = \sum n_i[T_i]$。
由线性性，可假设 $\alpha = [Z]$，其中 $Z \subset X$ 是维数为 $r$ 的不可约闭子簇。
我们对下列各整数的最大值 $e$ 作归纳：
$$
\dim(Z \cap T_i)
$$
基本情形是 $e = r + s - \dim(X)$。此时 $Z$ 与 $\beta$ 正常相交，引理显然成立。

\medskip\noindent
归纳步骤。假设 $e > r + s - \dim(X)$。选取嵌入
$X \subset \mathbf{P}^N$，应用引理 \ref{intersection-lemma-moving} 找到闭子簇
$C \subset \mathbf{P}^N$，使得 $C \cdot X = [Z] + \sum m_j[Z_j]$，且归纳假设适用于每个
$Z_j$。接着，对 $C$、$X$、$T_i$ 应用引理 \ref{intersection-lemma-move}，得到
$C' \subset \mathbf{P}^N \times \mathbf{P}^1$。把
$\gamma = C' \cdot X \times \mathbf{P}^1$ 看作 $X \times \mathbf{P}^1$ 上的循环。
由引理 \ref{intersection-lemma-transfer}，有
$$
[Z] + \sum m_j[Z_j] = \text{pr}_{X, *}(\gamma \cdot X \times 0)
$$
另一方面，循环
$\gamma_\infty = \text{pr}_{X, *}(\gamma \cdot X \times \infty)$
支撑于 $C'_\infty \cap X$，因而与 $\beta$ 正常相交。由引理
\ref{intersection-lemma-rational-equivalence-and-intersection}，可见
$[Z] \sim_{rat} - \sum m_j[Z_j] + \gamma_\infty$。
归纳假设说明每个 $[Z_j]$ 都有理等价于一个与 $\beta$ 正常相交的循环，故证明完毕。
\end{proof}



\section{交积与有理等价}
\label{intersection-section-intersections-and-rational-equivalence}

\noindent
沿用上述定义，我们证明交积模有理等价是良定义的。先处理一个特殊情形。

\begin{lemma}
\label{intersection-lemma-well-defined-special-case}
设 $X$ 是非奇异簇，$W \subset X \times \mathbf{P}^1$ 是一个支配
$\mathbf{P}^1$ 的 $(s + 1)$ 维子簇。设 $W_a$、\ $W_b$ 分别是
$W \to \mathbf{P}^1$ 在 $a$、\ $b$ 上的纤维。设 $V$ 是 $X$ 的 $r$ 维子簇，
且 $V$ 与 $W_a$、$W_b$ 都正常相交。则
$[V] \cdot [W_a]_r \sim_{rat} [V] \cdot [W_b]_r$。
\end{lemma}

\begin{proof}
有 $[W_a]_r = \text{pr}_{X,*}(W \cdot X \times a)$，对 $[W_b]_r$ 也类似；
见引理 \ref{intersection-lemma-rational-equivalence-and-intersection}。因此问题化为证明
$$
V \cdot \text{pr}_{X,*}( W \cdot X \times a) \sim_{rat} V \cdot
\text{pr}_{X,*}( W \cdot X\times b).
$$
应用投影公式引理 \ref{intersection-lemma-projection-formula-flat}，得到
$$
V \cdot \text{pr}_{X,*}( W \cdot X\times a) =
\text{pr}_{X,*}(V \times \mathbf{P}^1 \cdot (W \cdot X\times a))
$$
对 $b$ 也类似。因此问题化为证明
$$
\text{pr}_{X,*}(V \times \mathbf{P}^1 \cdot (W \cdot X\times a))
\sim_{rat}
\text{pr}_{X,*}(V \times \mathbf{P}^1 \cdot (W \cdot X\times b))
$$
若 $V \times \mathbf{P}^1$ 与 $W$ 正常相交，则交重数的结合性
（引理 \ref{intersection-lemma-associative}）给出
$V \times \mathbf{P}^1 \cdot (W \cdot X\times a) =
(V \times \mathbf{P}^1 \cdot W) \cdot X \times a$，对 $b$ 亦然。
于是问题化为证明
$$
\text{pr}_{X,*}((V \times \mathbf{P}^1 \cdot W) \cdot X\times a)
\sim_{rat}
\text{pr}_{X,*}((V \times \mathbf{P}^1 \cdot W) \cdot X\times b)
$$
这由引理 \ref{intersection-lemma-rational-equivalence-and-intersection} 成立。

\medskip\noindent
上面的论证并不完全成立。障碍在于，我们不知道 $V \times \mathbf{P}^1$ 与 $W$ 正常相交；
只知道 $V$ 与 $W_a$、以及 $V$ 与 $W_b$ 正常相交。令 $Z_i$（$i \in I$）为
$V \times \mathbf{P}^1 \cap W$ 的不可约分支。由引理
\ref{intersection-lemma-intersect-in-smooth}，若 $n = \dim(X)$，则
$\dim(Z_i) \geq r + 1 + s + 1 - n - 1 = r + s + 1 - n$。
由于假设了 $V$ 与 $W_a$ 正常相交，可知
$\dim(Z_{i, a}) = r + s - n$ 或 $Z_{i, a} = \emptyset$。
另一方面，若 $Z_{i, a} \not = \emptyset$，则
$\dim(Z_{i, a}) \geq \dim(Z_i) - 1 = r + s - n$。
所以，若 $Z_i$ 与 $X \times a$ 相交，则
$\dim(Z_i) = r + s + 1 - n$，并且此时 $Z_i \to \mathbf{P}^1$ 是满射。
于是可写 $I = I' \amalg I''$，其中 $I'$ 是使 $Z_i \to \mathbf{P}^1$ 为满射的
$i \in I$ 的集合，而 $I''$ 是使 $Z_i$ 位于某个闭点
$t_i \in \mathbf{P}^1$ 上方的 $i \in I$ 的集合；这里 $t_i \not = a$ 且 $t_i \not = b$。
考虑循环
$$
\gamma = \sum\nolimits_{i \in I'} e_i [Z_i]
$$
其中取
$$
e_i = \sum\nolimits_p (-1)^p
\text{length}_{\mathcal{O}_{X \times \mathbf{P}^1, Z_i}}
\text{Tor}_p^{\mathcal{O}_{X \times \mathbf{P}^1, Z_i}}(
\mathcal{O}_{V \times \mathbf{P}^1, Z_i}, \mathcal{O}_{W, Z_i})
$$
我们将证明，$\gamma$ 可以代替 $V \times \mathbf{P}^1$ 与 $W$ 的交积。

\medskip\noindent
与上面完全相同，我们利用交积的结合性来证明这一点。令
$U = \mathbf{P}^1 \setminus \{t_i, i \in I''\}$。注意，$X \times a$ 和
$X \times b$ 都包含于 $X \times U$。由 $U$ 的选取，子簇
$$
V \times U,\quad W_U,\quad X \times a\quad\text{属于}\quad X \times U
$$
两两横截相交；此外，$\dim(V \times U \cap W_U \cap X \times a) = \dim(V \cap W_a)$
具有预期维数。因此由引理 \ref{intersection-lemma-associative}，作为 $X \times U$ 上的循环，有
$$
V \times U \cdot (W_U \cdot X \times a) =
(V \times U \cdot W_U) \cdot X \times a
$$
由构造，$\gamma$ 在 $X \times U$ 上的限制是循环 $V \times U \cdot W_U$。
显然，$V \times \mathbf{P}^1 \cdot (W \times X \times a)$ 在 $X \times U$ 上的限制是
$V \times U \cdot (W_U \cdot X \times a)$。故
$$
V \times \mathbf{P}^1 \cdot (W \cdot X \times a) =
\gamma \cdot X \times a
$$
作为 $X \times \mathbf{P}^1$ 上的循环成立（因为等式两边都包含于 $X \times U$，
且根据前述结论，它们限制到 $X \times U$ 后相等）。对 $b$ 也有相同结论，于是
\begin{align*}
V \cdot [W_a]
& =
\text{pr}_{X,*}(V \times \mathbf{P}^1 \cdot (W \cdot X\times a)) \\
& =
\text{pr}_{X, *}(\gamma \cdot X \times a) \\
& \sim_{rat} 
\text{pr}_{X, *}(\gamma \cdot X \times b) \\
& =
\text{pr}_{X,*}(V \times \mathbf{P}^1 \cdot (W \cdot X\times b)) \\
& =
V \cdot [W_b]
\end{align*}
第一个和最后一个等号来自证明的第一段，第二个和倒数第二个等号已在本段中证明，
中间的等价关系就是引理 \ref{intersection-lemma-rational-equivalence-and-intersection}。
\end{proof}

\begin{theorem}
\label{intersection-theorem-well-defined}
设 $X$ 是非奇异射影簇，$\alpha$、\ $\beta$ 分别是 $X$ 上的 $r$、\ $s$-循环。
假设 $\alpha$ 与 $\beta$ 正常相交，因而 $\alpha \cdot \beta$ 有定义。
再假设 $\alpha \sim_{rat} 0$。则 $\alpha \cdot \beta \sim_{rat} 0$。
\end{theorem}

\begin{proof}
选取闭浸入 $X \subset \mathbf{P}^N$。由线性性，只需证明如下情形：
$\beta = [Z]$，其中 $Z \subset X$ 是一个与 $\alpha$ 正常相交的 $s$ 维闭子簇。
条件 $\alpha \sim_{rat} 0$ 意味着存在有限多个 $(r + 1)$ 维闭子簇
$W_i \subset X \times \mathbf{P}^1$，使得
$$
\alpha = \sum [W_{i, a_i}]_r - [W_{i, b_i}]_r
$$
其中 $a_i, b_i$ 是 $\mathbf{P}^1$ 中的若干点对。令
$W_{i, a_i}^t$ 与 $W_{i, b_i}^t$ 分别为 $W_{i, a_i}$ 与 $W_{i, b_i}$ 的不可约分支。
我们将对下列各整数的最大值 $d$ 作归纳：
$$
\dim(Z \cap W_{i, a_i}^t),\quad \dim(Z \cap W_{i, b_i}^t)
$$
证明余下部分的主要困难在于：尽管已知 $Z$ 与 $\alpha$ 正常相交，
$Z$ 却未必与“中间”簇 $W_{i, a_i}^t$ 和 $W_{i, b_i}^t$ 正常相交；
也就是说，可能有 $d >  r + s - \dim(X)$。

\medskip\noindent
基本情形：$d = r + s - \dim(X)$。此时 $Z$ 与所有 $W_{i, a_i}^t$、
$W_{i, b_i}^t$ 的交都是正常相交。所需结论由引理
\ref{intersection-lemma-well-defined-special-case} 得出，因为对每个 $i$，该引理给出
$[Z] \cdot [W_{i, a_i}]_r \sim_{rat} [Z] \cdot [W_{i, b_i}]_r$。

\medskip\noindent
归纳步骤：$d >  r + s - \dim(X)$。对 $Z \subset X$ 以及子簇族
$\{W_{i, a_i}^t, W_{i, b_i}^t\}$ 应用引理 \ref{intersection-lemma-moving}。
由此找到与 $X$ 正常相交的闭子簇 $C \subset \mathbf{P}^N$，使得
$$
C \cdot X = [Z] + \sum m_j [Z_j]
$$
并且
$$
\dim(Z_j \cap W_{i, a_i}^t) \leq \dim(Z \cap W_{i, a_i}^t),\quad
\dim(Z_j \cap W_{i, b_i}^t) \leq \dim(Z \cap W_{i, b_i}^t)
$$
当右端 $> r + s - \dim(X)$ 时，不等式严格成立。这蕴含两点：
(a) 归纳假设适用于每个 $Z_j$；(b) $C \cdot X$ 与 $\alpha$ 正常相交
（因为 $\alpha$ 是那些与 $Z$ 正常相交的 $[W_{i, a_i}^t]$ 和
$[W_{i, a_i}^t]$ 的线性组合）。
接着，按照引理 \ref{intersection-lemma-move}，对 $C$、$X$、$W_{i, a_i}^t$、
$W_{i, b_i}^t$ 选取 $C' \subset \mathbf{P}^N \times \mathbf{P}^1$。
对某些维数为 $s + 1$ 的子簇 $E_k \subset X \times \mathbf{P}^1$，写成
$C' \cdot X \times \mathbf{P}^1 = \sum n_k [E_k]$。
由命题 \ref{intersection-proposition-positivity}，对所有 $k$ 都有 $n_k > 0$。
由引理 \ref{intersection-lemma-transfer}，有
$$
[Z] + \sum m_j [Z_j] = \sum n_k[E_{k, 0}]_s
$$
由于 $E_{k, 0} \subset C \cap X$，可见 $[E_{k, 0}]_s$ 与 $\alpha$ 正常相交。
另一方面，循环
$$
\gamma = \sum n_k[E_{k, \infty}]_s
$$
支撑于 $C'_\infty \cap X$，因而与每个 $W_{i, a_i}^t$、$W_{i, b_i}^t$ 正常相交。
所以由基本情形和线性性，有
$$
\gamma \cdot \alpha \sim_{rat} 0
$$
我们已经看到 $E_{k, 0}$ 与 $E_{k, \infty}$ 都同 $\alpha$ 正常相交，
故把引理 \ref{intersection-lemma-well-defined-special-case} 应用于
$E_k \subset X \times \mathbf{P}^1$ 和 $\alpha$，得到
$$
[E_{k, 0}] \cdot \alpha \sim_{rat} [E_{k, \infty}] \cdot \alpha
$$
综上，有
\begin{align*}
[Z] \cdot \alpha
& =
(\sum n_k[E_{k, 0}]_r - \sum m_j[Z_j]) \cdot \alpha \\
& \sim_{rat}
\sum n_k [E_{k, 0}] \cdot \alpha \quad (\text{由归纳假设})\\
& \sim_{rat}
\sum n_k [E_{k, \infty}] \cdot \alpha \quad (\text{由该引理})\\
& =
\gamma \cdot \alpha \\
& \sim_{rat}
0 \quad (\text{由基础情形})
\end{align*}
证明完毕。
\end{proof}

\begin{remark}
\label{intersection-remark-quasi-projective}
引理 \ref{intersection-lemma-moving-move} 和定理 \ref{intersection-theorem-well-defined} 对非奇异拟射影簇也成立，
证明相同。唯一的变化是需要证明移动引理 \ref{intersection-lemma-moving} 的下述版本：
设 $X \subset \mathbf{P}^N$ 是闭子簇。令 $n = \dim(X)$，且
$0 \leq d, d' < n$。设 $X^{reg} \subset X$ 为非奇异点组成的开子集。
设 $Z \subset X^{reg}$ 是维数为 $d$ 的闭子簇，并设
$T_i \subset X^{reg}$（$i \in I$）是有限多个维数为 $d'$ 的闭子簇。
则存在子簇 $C \subset \mathbf{P}^N$，使得 $C$ 与 $X$ 正常相交，并且
$$
(C \cdot X)|_{X^{reg}} = Z + \sum\nolimits_{j \in J} m_j Z_j
$$
其中 $Z_j \subset X^{reg}$ 是维数为 $d$、彼此不同于 $Z$ 的不可约子簇，并且
$$
\dim(Z_j \cap T_i) \leq \dim(Z \cap T_i)
$$
若 $Z$ 与 $T_i$ 在 $X^{reg}$ 中不正常相交，则不等式严格成立。
\end{remark}



\section{Chow 环}
\label{intersection-section-chow-rings}

\noindent
设 $X$ 是非奇异射影簇。如下定义交积
$$
\CH_r(X) \times \CH_s(X) \longrightarrow \CH_{r + s - \dim(X)}(X),\quad
(\alpha, \beta) \longmapsto \alpha \cdot \beta
$$
设 $\alpha \in Z_r(X)$ 且 $\beta \in Z_s(X)$。若 $\alpha$ 与 $\beta$ 正常相交，
则采用第 \ref{intersection-section-intersection-product} 节给出的定义。否则，按照引理
\ref{intersection-lemma-moving-move} 选取 $\alpha \sim_{rat} \alpha'$，并令
$$
\alpha \cdot \beta =
\text{下列对象的类： }\alpha' \cdot \beta \in \CH_{r + s - \dim(X)}(X)
$$
由定理 \ref{intersection-theorem-well-defined}，这是良定义的，并且可下降到有理等价类。
$\CH_*(X)$ 上的交积是交换的（这一点显然）、结合的
（引理 \ref{intersection-lemma-associative}），且有单位元 $[X] \in \CH_{\dim(X)}(X)$。

\medskip\noindent
我们常用 $\CH^c(X) = \CH_{\dim X - c}(X)$ 表示余维为 $c$ 的循环的 Chow 群；
见《Chow 同调》第 \ref{chow-section-cycles-codimension} 节。交积定义了乘法
$$
\CH^k(X) \times \CH^l(X) \longrightarrow \CH^{k + l}(X)
$$
它是交换的、结合的，并且有单位元 $1 = [X] \in \CH^0(X)$。


\section{一般态射的拉回}
\label{intersection-section-general-pullback}

\noindent
设 $f : X \to Y$ 是非奇异射影簇之间的态射。定义
$$
f^* : \CH_k(Y) \to \CH_{k+\dim X - \dim Y}(X)
$$
其规则为
$$
f^*(\alpha) = pr_{X, *}(\Gamma_f \cdot pr_Y^*(\alpha))
$$
其中 $\Gamma_f \subset X\times Y$ 是 $f$ 的图。注意，在这种一般情形下，
它只定义在循环类上，而不定义在循环本身上。采用第 \ref{intersection-section-chow-rings} 节引入的
$\CH^*$ 记号，可以把拉回看作映射
$$
f^* : \CH^*(Y) \to \CH^*(X)
$$
换言之，这是分次 Abel 群之间的映射。

\begin{lemma}
\label{intersection-lemma-pullback-and-intersection-product}
设 $f : X \to Y$ 是非奇异射影簇之间的态射。Chow 群上的拉回映射满足：
\begin{enumerate}
\item $f^* : \CH^*(Y) \to \CH^*(X)$ 是环同态；
\item 对可复合的一对态射 $f, g$，有 $(g \circ f)^* = f^* \circ g^*$；
\item 投影公式成立：$f_*(\alpha) \cdot \beta =
f_*( \alpha \cdot f^*\beta)$；
\item 若 $f$ 平坦，则它与此前的定义一致。
\end{enumerate}
\end{lemma}

\begin{proof}
这些结论都容易由上面的结果推出。

\medskip\noindent
对 (1)，只需对 $X \times Y$ 上的循环 $\alpha$、$\beta$ 证明
$\text{pr}_{X,*}( \Gamma_f \cdot \alpha \cdot \beta) =
\text{pr}_{X,*}(\Gamma_f \cdot \alpha) \cdot
\text{pr}_{X,*}(\Gamma_f \cdot \beta)$。若 $\alpha$ 是 $X \times Y$ 上与
$\Gamma_f$ 正常相交的循环，则由于 $\Gamma_f$ 是图，容易看出作为循环有
$$
\Gamma_f \cdot \alpha =
\Gamma_f \cdot \text{pr}_X^*(\text{pr}_{X,*}(\Gamma_f \cdot \alpha))
$$
因此得到下列计算中的第一个等号：
\begin{align*}
\text{pr}_{X,*}(\Gamma_f \cdot \alpha \cdot \beta)
& =
\text{pr}_{X,*}(
\Gamma_f \cdot
\text{pr}_X^*(\text{pr}_{X,*}(\Gamma_f \cdot \alpha)) \cdot \beta) \\
& =
\text{pr}_{X,*}(\text{pr}_X^*(\text{pr}_{X,*}(\Gamma_f \cdot \alpha))
\cdot (\Gamma_f \cdot \beta)) \\
& =
\text{pr}_{X,*}(\Gamma_f \cdot \alpha) \cdot
\text{pr}_{X,*}(\Gamma_f \cdot \beta)
\end{align*}
最后一步使用了平坦情形下的投影公式
（引理 \ref{intersection-lemma-projection-formula-flat}）。

\medskip\noindent
若 $g : Y \to Z$，则性质 (2) 形式地由下述观察推出：
$$
\Gamma =
\text{pr}_{X \times Y}^*\Gamma_f \cdot
\text{pr}_{Y \times Z}^*\Gamma_g
$$
在 $Z_*(X \times Y \times Z)$ 中成立，其中
$\Gamma = \{(x, f(x), g(f(x))\}$，并且它同构地映到 $X \times Z$ 中的
$\Gamma_{g \circ f}$。该等式由概形论等式和引理 \ref{intersection-lemma-transversal} 得出。

\medskip\noindent
对 (3)，两次使用平坦映射的投影公式：
\begin{align*}
f_*(\alpha \cdot pr_{X, *}(\Gamma_f \cdot pr_Y^*(\beta)))
& =
f_*(pr_{X, *}(pr_X^*\alpha \cdot \Gamma_f \cdot pr_Y^*(\beta))) \\
& =
pr_{Y, *}(pr_X^*\alpha \cdot \Gamma_f \cdot pr_Y^*(\beta))) \\
& =
pt_{Y, *}(pr_X^*\alpha \cdot \Gamma_f) \cdot \beta \\
& =
f_*(\alpha) \cdot \beta
\end{align*}
在最后一个等式中使用了上面对图所作的观察。这证明了 (3)。

\medskip\noindent
性质 (4) 依赖于在 $f$ 平坦时识别交积
$\Gamma_f \cdot pr_Y^*\alpha$。具体地，在这种情形下，若 $V \subset Y$ 是闭子簇，
则概形 $f^{-1}(V) \cong \Gamma_f \cap pr_Y^{-1}(V)$ 的每个泛点 $\xi$
都位于 $V$ 的泛点之上。因此，$pr_Y^{-1}(V) = X \times V$ 在 $\xi$ 处的局部环是
Cohen--Macaulay 的。由于 $\Gamma_f \subset X \times Y$ 是正则浸入
（作为光滑射影簇之间的态射），可得
$$
\Gamma_f \cdot pr_Y^*[V] = [\Gamma_f \cap pr_Y^{-1}(V)]_d
$$
其中 $d$ 是 $\Gamma_f \cap pr_Y^{-1}(V)$ 的维数；见引理
\ref{intersection-lemma-multiplicity-lci-CM}。由于 $\Gamma_f \cap pr_Y^{-1}(V)$
同构地映到 $f^{-1}(V)$，结论得证。
\end{proof}


\section{循环的拉回}
\label{intersection-section-pullback-cycles}

\noindent
设 $X$ 和 $Y$ 是非奇异射影簇，$f : X \to Y$ 是态射。
设 $Z \subset Y$ 是闭子簇，并令 $f^{-1}(Z)$ 为概形论逆像：
$$
\xymatrix{
f^{-1}(Z) \ar[r] \ar[d] & Z \ar[d] \\
X \ar[r] & Y
}
$$
这是概形的纤维积图。特别地，$f^{-1}(Z) \subset X$ 是 $X$ 的闭子概形。
在这种情形下，总有
$$
\dim f^{-1}(Z) \geq \dim Z + \dim X - \dim Y.
$$
若上式取等号，并且概形 $Z$ 在 $f^{-1}(Z)$ 各泛点的像处都是 Cohen--Macaulay 的，
则
$f^*[Z] = [f^{-1}(Z)]_{\dim Z + \dim X - \dim Y}$。
这是因为可以把 $f^{-1}(Z)$ 识别为 $\Gamma_f$ 与 $X \times Z$ 的概形论交，
再应用引理 \ref{intersection-lemma-multiplicity-lci-CM}。细节与上面引理
\ref{intersection-lemma-pullback-and-intersection-product} 第 (4) 部分的证明相似。
% Generated mechanically from the frozen pic.tex target.
% Do not edit this generated file; edit the bound locale source instead.

% OK, start here.
%

\setcounter{chapter}{43}
\chapter{曲线的 Picard 概形}



\phantomsection
\label{pic-section-phantom}



\section{引言}
\label{pic-section-introduction}

\noindent
本章完成构造代数闭域上非奇异射影曲线的 Picard 概形所必需的工作。
更详尽的讨论和历史背景见 \cite{Kleiman-Picard}。

\medskip\noindent
在 Stacks 项目的后续部分，我们将在远为一般的情形下讨论 Hilbert 函子与 Quot 函子。


\section{点的 Hilbert 概形}
\label{pic-section-hilbert-scheme-points}

\noindent
设 $X \to S$ 是概形态射，$d \geq 0$ 是整数。对 $S$ 上的概形 $T$，令
$$
\Hilbfunctor^d_{X/S}(T) =
\left\{
\begin{matrix}
Z \subset X_T\text{ closed subscheme such that }\\
Z \to T\text{ is finite locally free of degree }d
\end{matrix}
\right\}
$$
若 $T' \to T$ 是 $S$ 上概形之间的态射，且
$Z \in \Hilbfunctor^d_{X/S}(T)$，则基变换
$Z_{T'} \subset X_{T'}$ 是 $\Hilbfunctor^d_{X/S}(T')$ 的一个元素。
由此得到函子
$$
\Hilbfunctor^d_{X/S} :
(\Sch/S)^{opp} \longrightarrow \textit{Sets},\quad
T \longrightarrow \Hilbfunctor^d_{X/S}(T)
$$
一般而言，$\Hilbfunctor^d_{X/S}$ 是代数空间（此处将来插入引用）。
本节将证明：若 $X \to S$ 的任一纤维中的任意有限多个点都包含于某个仿射开集中，
则 $\Hilbfunctor^d_{X/S}$ 可由概形表示。若 $\Hilbfunctor^d_{X/S}$ 可由概形表示，
我们常把这个概形记作 $\underline{\Hilbfunctor}^d_{X/S}$。

\begin{lemma}
\label{pic-lemma-hilb-d-sheaf}
设 $X \to S$ 是概形态射。函子 $\Hilbfunctor^d_{X/S}$ 满足 fpqc 拓扑的层性质
（《拓扑》定义 \ref{topologies-definition-sheaf-property-fpqc}）。
\end{lemma}

\begin{proof}
设 $\{T_i \to T\}_{i \in I}$ 是 $S$ 上概形的一个 fpqc 覆盖。令
$X_i = X_{T_i} = X \times_S T_i$。注意，$\{X_i \to X_T\}_{i \in I}$ 是
$X_T$ 的 fpqc 覆盖（《拓扑》引理 \ref{topologies-lemma-fpqc}），并且
$X_{T_i \times_T T_{i'}} = X_i \times_{X_T} X_{i'}$。
假设给定一族元素 $Z_i \in \Hilbfunctor^d_{X/S}(T_i)$，使得 $Z_i$ 与 $Z_{i'}$
映到 $\Hilbfunctor^d_{X/S}(T_i \times_T T_{i'})$ 中的同一元素。
由闭浸入的有效下降（《下降》引理 \ref{descent-lemma-closed-immersion}），
存在闭浸入 $Z \to X_T$，其沿 $X_i \to X_T$ 的基变换等于 $Z_i \to X_i$。
于是态射 $Z \to T$ 沿 $T_i$ 的基变换为 $Z_i \to T_i$。因此由《下降》引理
\ref{descent-lemma-descending-property-finite-locally-free}，$Z \to T$ 是次数为 $d$ 的局部自由有限态射。
\end{proof}

\begin{lemma}
\label{pic-lemma-hilb-d-limit-preserving}
设 $X \to S$ 是概形态射。若 $X \to S$ 是有限表现的，则函子
$\Hilbfunctor^d_{X/S}$ 保持极限（《极限》注 \ref{limits-remark-limit-preserving}）。
\end{lemma}

\begin{proof}
设 $T = \lim T_i$ 是 $S$ 上仿射概形的一个极限。需要证明
$\Hilbfunctor^d_{X/S}(T) = \colim \Hilbfunctor^d_{X/S}(T_i)$。
注意，若 $Z \to X_T$ 是 $\Hilbfunctor^d_{X/S}(T)$ 的一个元素，则
$Z \to T$ 是有限表现的。因此由《极限》引理
\ref{limits-lemma-descend-finite-presentation}，存在某个 $i$、一个在 $T_i$ 上有限表现的概形
$Z_i$，以及 $T_i$ 上的态射 $Z_i \to X_{T_i}$，使其基变换到 $T$ 后给出
$Z \to X_T$。应用《极限》引理
\ref{limits-lemma-descend-closed-immersion-finite-presentation} 可知，增大 $i$ 后可假设
$Z_i \to X_{T_i}$ 是闭浸入。再应用《极限》引理
\ref{limits-lemma-descend-finite-locally-free} 可知，必要时再次增大 $i$ 后，
$Z_i \to T_i$ 是次数为 $d$ 的局部自由有限态射。因此如愿有
$Z_i \in \Hilbfunctor^d_{X/S}(T_i)$。
\end{proof}

\noindent
设 $S$ 是概形，$i : X \to Y$ 是 $S$ 上概形之间的闭浸入。则存在函子变换
$$
\Hilbfunctor^d_{X/S} \longrightarrow \Hilbfunctor^d_{Y/S}
$$
它把元素 $Z \in \Hilbfunctor^d_{X/S}(T)$ 映到
$\Hilbfunctor^d_{Y/S}$ 中的 $i_T(Z) \subset Y_T$。这里
$i_T : X_T \to Y_T$ 是 $i$ 的基变换。

\begin{lemma}
\label{pic-lemma-hilb-d-of-closed}
设 $S$ 是概形，$i : X \to Y$ 是概形的闭浸入。若
$\Hilbfunctor^d_{Y/S}$ 可由概形表示，则 $\Hilbfunctor^d_{X/S}$ 也可由概形表示，
并且相应的概形态射
$\underline{\Hilbfunctor}^d_{X/S} \to \underline{\Hilbfunctor}^d_{Y/S}$
是闭浸入。
\end{lemma}

\begin{proof}
设 $T$ 是 $S$ 上的概形，且 $Z \in \Hilbfunctor^d_{Y/S}(T)$。
断言：存在闭子概形 $T_X \subset T$，使得概形态射 $T' \to T$ 经由 $T_X$ 分解，
当且仅当 $Z_{T'} \to Y_{T'}$ 经由 $X_{T'}$ 分解。
把这一断言应用于表示 $\Hilbfunctor^d_{Y/S}$ 的概形 $T_{univ}$ 及其泛对象\footnote{见
《范畴》第 \ref{categories-section-opposite} 节}
$Z_{univ} \in \Hilbfunctor^d_{Y/S}(T_{univ})$，得到闭子概形
$T_{univ, X} \subset T_{univ}$，使得
$Z_{univ, X} = Z_{univ} \times_{T_{univ}} T_{univ, X}$
是 $X \times_S T_{univ, X}$ 的闭子概形，从而定义
$\Hilbfunctor^d_{X/S}(T_{univ, X})$ 的一个元素。形式论证随即表明，
$T_{univ, X}$ 是表示 $\Hilbfunctor^d_{X/S}$ 的概形，其泛对象为 $Z_{univ, X}$。

\medskip\noindent
证明该断言。考虑 $Z' = X_T \times_{Y_T} Z$。给定 $T' \to T$，可见
$Z_{T'} \to Y_{T'}$ 经由 $X_{T'}$ 分解，当且仅当
$Z'_{T'} \to Z_{T'}$ 是同构。因此该断言由非常一般的《平坦性的进一步内容》引理
\ref{flat-lemma-Weil-restriction-closed-subschemes} 得出。
不过，在本特殊情形下也可以直接证明如下：先化归到
$T = \Spec(A)$ 且 $Z = \Spec(B)$ 的情形。进一步缩小 $T$ 后，
可假设作为 $A$-模存在同构 $\varphi : B \to A^{\oplus d}$。
于是对某个理想 $J \subset B$，有 $Z' = \Spec(B/J)$。取一组生成元
$g_\beta \in J$，并写成 $\varphi(g_\beta) = (g_\beta^1, \ldots, g_\beta^d)$。
显然，$T_X$ 由 $\Spec(A/(g_\beta^j))$ 给出。
\end{proof}

\begin{lemma}
\label{pic-lemma-hilb-d-separated}
设 $X \to S$ 是概形态射。若 $X \to S$ 是分离的，且
$\Hilbfunctor^d_{X/S}$ 可表示，则
$\underline{\Hilbfunctor}^d_{X/S} \to S$ 是分离的。
\end{lemma}

\begin{proof}
本证明中所有未标明基底的积都取在 $S$ 上。令
$H = \underline{\Hilbfunctor}^d_{X/S}$，并令
$Z \in \Hilbfunctor^d_{X/S}(H)$ 为泛对象。沿两个投影
$H \times H \to H$ 拉回 $Z$，得到两个对象
$Z_1, Z_2 \in \Hilbfunctor^d_{X/S}(H \times H)$。于是
$Z_1 = Z \times H \subset X_{H \times H}$，且
$Z_2 = H \times Z
\subset X_{H \times H}$。由于 $H$ 表示函子
$\Hilbfunctor^d_{X/S}$，对角态射 $\Delta : H \to H \times H$ 具有如下泛性质：
概形态射 $T \to H \times H$ 经由 $\Delta$ 分解，当且仅当作为
$\Hilbfunctor^d_{X/S}(T)$ 的元素有 $Z_{1, T} = Z_{2, T}$。
令 $Z = Z_1 \times_{X_{H \times H}} Z_2$。于是可见，
$T \to H \times H$ 经由 $\Delta$ 分解，当且仅当态射
$Z_T \to Z_{1, T}$ 和 $Z_T \to Z_{2, T}$ 都是同构。
由非常一般的《平坦性的进一步内容》引理
\ref{flat-lemma-Weil-restriction-closed-subschemes}，$\Delta$ 是闭浸入。
在引理 \ref{pic-lemma-hilb-d-of-closed} 的证明中，读者可以找到适用于本特殊情形的、
所需结论的另一种更简单证明。
\end{proof}

\begin{lemma}
\label{pic-lemma-hilb-d-An}
设 $X \to S$ 是仿射概形之间的态射，且 $d \geq 0$。则
$\Hilbfunctor^d_{X/S}$ 可表示。
\end{lemma}

\begin{proof}
设 $S = \Spec(R)$。对某个基数充分大的集合 $I$，可以选取 $X$ 到
$R[x_i; i \in I]$ 的谱中的闭浸入。因此由引理 \ref{pic-lemma-hilb-d-of-closed}，
可假设 $X = \Spec(A)$，其中 $A = R[x_i; i \in I]$。
在这种情形下，我们用《概形》引理 \ref{schemes-lemma-glue-functors} 证明本引理。

\medskip\noindent
该引理的条件 (1) 由引理 \ref{pic-lemma-hilb-d-sheaf} 得出。

\medskip\noindent
对每个基数为 $d$ 的子集 $W \subset A$，构造 $\Hilbfunctor^d_{X/S}$ 的子函子
$F_W$。（只考虑 $W$ 由 $x_i$ 中的一族单项式组成的情形已经足够，但这里并不需要这一点。）
具体地说，规定 $Z \in \Hilbfunctor^d_{X/S}(T)$ 属于 $F_W(T)$，当且仅当
$\mathcal{O}_T$-线性映射
$$
\bigoplus\nolimits_{f \in W} \mathcal{O}_T
\longrightarrow
(Z \to T)_*\mathcal{O}_Z,\quad
(g_f) \longmapsto \sum g_f f|_Z
$$
是满射（等价地，是同构）。这里，对 $f \in A$ 和
$Z \in \Hilbfunctor^d_{X/S}(T)$，用 $f|_Z$ 表示 $f$ 沿态射
$Z \to X_T \to X$ 的拉回。

\medskip\noindent
开性，即该引理的条件 (2)(b)。它由《代数》引理
\ref{algebra-lemma-cokernel-flat} 得出。

\medskip\noindent
覆盖性，即该引理的条件 (2)(c)。由于
$$
A \otimes_R \mathcal{O}_T =
(X_T \to T)_*\mathcal{O}_{X_T} \to (Z \to T)_*\mathcal{O}_Z
$$
是满射，且 $(Z \to T)_*\mathcal{O}_Z$ 是秩为 $d$ 的局部自由有限模，
所以对每个点 $t \in T$，可以找到一个基数为 $d$ 的有限子集 $W \subset A$，
其像构成 $d$ 维 $\kappa(t)$-向量空间
$((Z \to T)_*\mathcal{O}_Z)_t \otimes_{\mathcal{O}_{T, t}} \kappa(t)$ 的一组基。
由 Nakayama 引理，存在 $t$ 的开邻域 $V \subset T$，使得 $Z_V \in F_W(V)$。

\medskip\noindent
可表示性，即该引理的条件 (2)(a)。设 $W \subset A$ 的基数为 $d$。
我们断言，$F_W$ 可由 $R$ 上的一个仿射概形表示。这里将构造这个仿射概形，
不过也建议读者自行完成这一推演。把 $W$ 的元素编号为
$f_1, \ldots, f_d$。如下构造 $T_{univ} = \Spec(R_{univ})$ 上
$F_W$ 的泛元素 $Z_{univ} = \Spec(B_{univ})$。首先令
$$
R_{univ} = R[c_{kl}^m, b^l, b_i^l]/\mathfrak a_{univ}
$$
其中理想 $\mathfrak a_{univ}$ 将在下面描述。这里指标 $k, l, m, i$ 都遍历
$\{1, \ldots, d\}$。然后令
$$
B_{univ} = R_{univ}e_1 \oplus R_{univ}e_2 \oplus \ldots \oplus R_{univ}e_d
$$
作为 $R_{univ}$-模，并按下列规则赋予代数结构：
$$
e_ke_l = \sum c_{kl}^m e_m
$$
闭浸入 $Z_{univ} \to X_{T_{univ}}$ 由 $R_{univ}$-代数映射
$$
\Psi : A \otimes_R R_{univ} \longrightarrow B_{univ}
$$
给出；它把 $1 \otimes 1$ 映到 $\sum b^le_l$，把 $x_i$ 映到
$\sum b_i^le_l$（存在性见下文）。理想 $\mathfrak a_{univ}$ 按以下条件规定：
\begin{enumerate}
\item $B_{univ}$ 上的乘法是交换的，即应有
$c_{lk}^m - c_{kl}^m \in \mathfrak a_{univ}$；
\item $B_{univ}$ 上的乘法是结合的，即应有
$c_{lk}^m c_{m n}^p - c_{lq}^p c_{kn}^q \in \mathfrak a_{univ}$；
\item $\sum b^le_l$ 是 $B_{univ}$ 中的乘法单位元 $1$，换言之，对所有 $k$，
应有 $(\sum b^le_l)e_k = e_k$；这意味着
$\sum b^lc_{lk}^m - \delta_{km} \in \mathfrak a_{univ}$
（Kronecker delta）。
\end{enumerate}
既然已保证 $B_{univ}$ 是交换 $R_{univ}$-代数，且由 $\sum b^le_l$ 给出单位元 $1$，
同态 $\Psi$ 便存在。最后一个条件是：
\begin{enumerate}
\item[(5)] $f_l$ 映到 $B_{univ}$ 中的 $e_l$。对某些
$h_l^m \in R[c_{kl}^m, b^l, b_i^l]$，写成
$\Psi(f_l) - e_l \equiv \sum h_l^me_m$，则应有
$h_l^m \in \mathfrak a_{univ}$。
\end{enumerate}
令 $\mathfrak a_{univ} \subset R[c_{kl}^m, b^l, b_i^l]$ 由 (1) -- (5)
中列出的元素生成，则显然 $F_W$ 由 $\Spec(R_{univ})$ 表示。
\end{proof}

\begin{proposition}
\label{pic-proposition-hilb-d-representable}
设 $X \to S$ 是概形态射，且 $d \geq 0$。假设对所有
$(s, x_1, \ldots, x_d)$，其中 $s \in S$ 且 $x_1, \ldots, x_d \in X_s$，
都存在仿射开集 $U \subset X$，使得 $x_1, \ldots, x_d \in U$。
则 $\Hilbfunctor^d_{X/S}$ 可由概形表示。
\end{proposition}

\begin{proof}
使用相对粘合（《构造》第 \ref{constructions-section-relative-glueing} 节）或函子观点
（《概形》引理 \ref{schemes-lemma-glue-functors}），可化归到 $S$ 仿射的情形。
细节从略。

\medskip\noindent
假设 $S$ 仿射。对仿射开集 $U \subset X$，用
$F_U \subset \Hilbfunctor^d_{X/S}$ 表示如下子函子：对概形 $T/S$，元素
$Z \in \Hilbfunctor^d_{X/S}(T)$ 属于 $F_U(T)$，当且仅当 $Z \subset U_T$。
我们将应用《概形》引理 \ref{schemes-lemma-glue-functors} 和这些子函子 $F_U$ 得出结论。

\medskip\noindent
条件 (1) 即引理 \ref{pic-lemma-hilb-d-sheaf}。

\medskip\noindent
条件 (2)(a) 来自 $F_U = \Hilbfunctor^d_{U/S}$，以及后者由引理
\ref{pic-lemma-hilb-d-An} 可表示这一事实。具体地，若 $Z \in F_U(T)$，
则可把 $Z$ 看作 $U_T$ 的闭子概形；它在 $T$ 上是次数为 $d$ 的局部自由有限概形，
因而 $Z \in \Hilbfunctor^d_{U/S}(T)$。反过来，若
$Z \in \Hilbfunctor^d_{U/S}(T)$，则 $Z \to U_T \to X_T$ 是闭浸入\footnote{若
$X \to S$ 分离，这一点显然：在这种情形下，《态射》引理
\ref{morphisms-lemma-image-proper-scheme-closed} 告诉我们浸入
$\varphi : Z \to X_T$ 具有闭像，因而由《概形》引理
\ref{schemes-lemma-immersion-when-closed}，它是闭浸入。建议读者略过本脚注的其余部分，
因为我们不知道存在任何满足 $X \to S$ 的上述假设但 $X \to S$ 不分离的实例。
在一般情形下，设 $x \in X_T$ 是 $\varphi(Z)$ 的闭包中的一点。需要证明
$x \in \varphi(Z)$。设 $t \in T$ 为 $x$ 的像。由关于 $X \to S$ 的假设，
可以选取包含 $x$ 和 $\varphi(Z_t)$ 的仿射开集 $W \subset X_T$。
于是 $\varphi^{-1}(W)$ 是包含整个纤维 $Z_t$ 的开集；由于 $Z \to T$ 是闭态射，
把 $T$ 替换为 $t$ 的一个开邻域后，可以假设 $Z = \varphi^{-1}(W)$。
这时由分离情形（因为 $W \to T$ 分离），$\varphi(Z) \subset W$ 是闭集，
所以 $x \in \varphi(Z)$。}，从而可以把 $Z$ 看作 $F_U(T)$ 的元素。

\medskip\noindent
设对 $S$ 上的某个概形 $T$，有 $Z \in \Hilbfunctor^d_{X/S}(T)$。令
$$
B = (Z \to T)\left((Z \to X_T \to X)^{-1}(X \setminus U)\right)
$$
这是 $T$ 的闭子集，并且显然在开集 $T_{Z, U} = T \setminus B$ 上，
限制 $Z_{t'}$ 映入 $U_{T'}$。另一方面，对任意 $b \in B$，纤维 $Z_b$ 不映入 $U$。
所以对给定态射 $T' \to T$，有
$Z_{T'} \in F_U(T')$ $\Leftrightarrow$ $T' \to T$ 经由开集 $T_{Z, U}$ 分解。
这证明了条件 (2)(b)。

\medskip\noindent
条件 (2)(c) 来自关于 $X/S$ 的假设。只需证明如下事实：若 $T$ 是一个域的谱，
$Z \subset X_T$ 是在 $T$ 上次数为 $d$ 的有限平坦闭子概形，则
$Z \to X_T \to X$ 经由 $X$ 的某个仿射开集 $U$ 分解。
这一点显然，因为 $Z$ 至多有 $d$ 个点，而这些点都映入 $T \to S$ 的像点上方的 $X$ 的纤维中。
\end{proof}

\begin{remark}
\label{pic-remark-when-proposition-applies}
设 $f : X \to S$ 是概形态射。在下列每种情形下，命题
\ref{pic-proposition-hilb-d-representable} 的假设、因而其结论都成立：
\begin{enumerate}
\item $X$ 是拟仿射的；
\item $f$ 是拟仿射态射；
\item $f$ 是拟射影态射；
\item $f$ 是局部射影态射；
\item $X$ 上存在充足可逆层；
\item $X$ 上存在 $f$-充足可逆层；
\item $X$ 上存在 $f$-很充足可逆层。
\end{enumerate}
事实上，在每种情形下，纤维 $X_s$ 中的任意有限点集都包含于 $X$ 的某个拟紧开集 $U$；
该开集带有充足可逆层，或同构于某个仿射概形的开子集，或同构于某个分次环的
$\text{Proj}$ 的开子集（每种情形都可由展开定义得到）。因此由《性质》引理
\ref{properties-lemma-ample-finite-set-in-affine}，存在所需的仿射开集。
\end{remark}




\section{光滑曲线上的除子模空间}
\label{pic-section-divisors}

\noindent
对相对维数为 $1$ 的光滑态射 $X \to S$，函子 $\Hilbfunctor^d_{X/S}$ 参数化
《除子》第 \ref{divisors-section-effective-Cartier-morphisms} 节中定义的相对有效 Cartier 除子。

\begin{lemma}
\label{pic-lemma-divisors-on-curves}
设 $X \to S$ 是相对维数为 $1$ 的光滑概形态射，$D \subset X$ 是闭子概形。
考虑以下条件：
\begin{enumerate}
\item $D \to S$ 是局部自由有限态射；
\item $D$ 是 $X/S$ 上的相对有效 Cartier 除子；
\item $D \to S$ 是局部拟有限、平坦且局部有限表现的态射；
\item $D \to S$ 是局部拟有限且平坦的态射。
\end{enumerate}
总有蕴含关系
$$
(1) \Rightarrow (2) \Leftrightarrow (3) \Rightarrow (4)
$$
若 $S$ 局部 Noether，则最后一个箭头也是等价。若 $X \to S$ 固有
（而 $S$ 任意），则第一个箭头也是等价。
\end{lemma}

\begin{proof}
先证 (2) 与 (3) 等价。由《除子》引理 \ref{divisors-lemma-fibre-Cartier}，
只需在 $S$ 是域 $k$ 的谱时证明 (2) 与 (3) 等价。设 $x \in X$ 是闭点。
由于 $X$ 在 $k$ 上光滑且相对维数为 $1$，可见 $\mathcal{O}_{X, x}$ 是维数为 $1$ 的正则局部环
（见《簇》引理 \ref{varieties-lemma-smooth-regular}）。因此
$\mathcal{O}_{X, x}$ 是离散赋值环（《代数》引理
\ref{algebra-lemma-characterize-dvr}），从而是主理想整环。于是，任何在 $X$ 的所有泛点处均不为零的理想层
$\mathcal{I} \subset \mathcal{O}_X$ 都是可逆的
（《除子》引理 \ref{divisors-lemma-effective-Cartier-in-points}）。
换言之，$X$ 的任何不包含泛点的闭子概形都是有效 Cartier 除子。
由此 (2) 与 (3) 等价。

\medskip\noindent
若 $S$ 是 Noether 概形，则任意局部拟有限态射 $D \to S$ 都是局部有限表现的
（《态射》引理 \ref{morphisms-lemma-noetherian-finite-type-finite-presentation}），
故 (3) 与 (4) 等价。

\medskip\noindent
若 $X \to S$ 固有（且 $S$ 任意），则 $D \to S$ 也固有。
由于固有的局部拟有限态射是有限态射（《态射的进一步内容》引理
\ref{more-morphisms-lemma-characterize-finite}），而有限、平坦且有限表现的态射是局部自由有限态射
（《态射》引理 \ref{morphisms-lemma-finite-flat}），可见 (1) 与 (2) 等价。
\end{proof}

\begin{lemma}
\label{pic-lemma-sum-divisors-on-curves}
设 $X \to S$ 是相对维数为 $1$ 的光滑概形态射，$D_1, D_2 \subset X$ 是
$S$ 上次数分别为 $d_1$、$d_2$ 的局部自由有限闭子概形。则 $D_1 + D_2$ 是
$S$ 上次数为 $d_1 + d_2$ 的局部自由有限概形。
\end{lemma}

\begin{proof}
由引理 \ref{pic-lemma-divisors-on-curves}，$D_1$ 与 $D_2$ 是 $X/S$ 上的相对有效 Cartier 除子。
因此由《除子》引理 \ref{divisors-lemma-sum-relative-effective-Cartier-divisor}，
$D = D_1 + D_2$ 是 $X/S$ 上的相对有效 Cartier 除子。
再由引理 \ref{pic-lemma-divisors-on-curves}，$D \to S$ 是局部拟有限、平坦且局部有限表现的态射。
把《态射》引理 \ref{morphisms-lemma-image-universally-closed-separated}
应用于满的整态射 $D_1 \amalg D_2 \to D$，可知 $D \to S$ 分离。
于是《态射》引理 \ref{morphisms-lemma-image-proper-is-proper} 蕴含 $D \to S$ 固有。
这又蕴含 $D \to S$ 有限（《态射的进一步内容》引理
\ref{more-morphisms-lemma-characterize-finite}），进而可知 $D \to S$ 局部自由有限
（《态射》引理 \ref{morphisms-lemma-finite-flat}）。
所以只需证明 $D \to S$ 的次数为 $d_1 + d_2$。为此，可以基变换到 $X \to S$ 的一个纤维，
因而可假设对某个域 $k$ 有 $S = \Spec(k)$。在这种情形下，存在有限多个闭点
$x_1, \ldots, x_n \in X$，使 $D_1$ 与 $D_2$ 都支撑于
$\{x_1, \ldots, x_n\}$。事实上，存在非零因子
$f_{i, j} \in \mathcal{O}_{X, x_i}$，使得
$$
D_1 = \coprod \Spec(\mathcal{O}_{X, x_i}/(f_{i, 1}))
\quad\text{且}\quad
D_2 = \coprod \Spec(\mathcal{O}_{X, x_i}/(f_{i, 2}))
$$
于是
$$
D = \coprod \Spec(\mathcal{O}_{X, x_i}/(f_{i, 1}f_{i, 2}))
$$
由此容易看出，$D$ 在 $k$ 上的次数为 $d_1 + d_2$
（必要时使用《代数》引理 \ref{algebra-lemma-ord-additive}）。
\end{proof}

\begin{lemma}
\label{pic-lemma-difference-divisors-on-curves}
设 $X \to S$ 是相对维数为 $1$ 的光滑概形态射，$D_1, D_2 \subset X$ 是
$S$ 上次数分别为 $d_1$、$d_2$ 的局部自由有限闭子概形。若
$D_1 \subset D_2$（作为闭子概形），则存在闭子概形 $D \subset X$，
它在 $S$ 上局部自由有限且次数为 $d_2 - d_1$，并满足 $D_2 = D_1 + D$。
\end{lemma}

\begin{proof}
本证明几乎与引理 \ref{pic-lemma-sum-divisors-on-curves} 的证明完全相同。
由引理 \ref{pic-lemma-divisors-on-curves}，$D_1$ 与 $D_2$ 是 $X/S$ 上的相对有效 Cartier 除子。
由《除子》引理 \ref{divisors-lemma-difference-relative-effective-Cartier-divisor}，
存在相对有效 Cartier 除子 $D \subset X$，使得 $D_2 = D_1 + D$。
所以由引理 \ref{pic-lemma-divisors-on-curves}，$D \to S$ 局部拟有限、平坦且局部有限表现。
由于 $D$ 是 $D_2$ 的闭子概形，$D \to S$ 是有限态射。因此 $D \to S$ 局部自由有限
（《态射》引理 \ref{morphisms-lemma-finite-flat}）。
所以只需证明 $D \to S$ 的次数为 $d_2 - d_1$，而这由引理
\ref{pic-lemma-sum-divisors-on-curves} 得出。
\end{proof}

\noindent
设 $X \to S$ 是相对维数为 $1$ 的光滑概形态射。由引理
\ref{pic-lemma-divisors-on-curves}，对 $S$ 上的概形 $T$ 以及
$D \in \Hilbfunctor^d_{X/S}(T)$，可以把 $D$ 看作 $X_T/T$ 上的相对有效 Cartier 除子，
使得 $D \to T$ 是次数为 $d$ 的局部自由有限态射。因此由引理
\ref{pic-lemma-sum-divisors-on-curves}，得到函子变换
$$
\Hilbfunctor^{d_1}_{X/S} \times \Hilbfunctor^{d_2}_{X/S}
\longrightarrow
\Hilbfunctor^{d_1 + d_2}_{X/S},\quad
(D_1, D_2) \longmapsto D_1 + D_2
$$
若对所有次数 $d$，$\Hilbfunctor^d_{X/S}$ 都可表示，则这个函子变换对应于
$S$ 上的概形态射
$$
\underline{\Hilbfunctor}^{d_1}_{X/S}
\times_S
\underline{\Hilbfunctor}^{d_2}_{X/S}
\longrightarrow
\underline{\Hilbfunctor}^{d_1 + d_2}_{X/S}
$$
注意，$\underline{\Hilbfunctor}^0_{X/S} = S$ 且
$\underline{\Hilbfunctor}^1_{X/S} = X$。上述态射的一个特殊情形是
$$
\underline{\Hilbfunctor}^d_{X/S} \times_S X
\longrightarrow
\underline{\Hilbfunctor}^{d + 1}_{X/S},\quad
(D, x) \longmapsto D + x
$$

\begin{lemma}
\label{pic-lemma-universal-object}
设 $X \to S$ 是相对维数为 $1$ 的光滑概形态射，并且函子
$\Hilbfunctor^d_{X/S}$ 都可表示。则态射
$\underline{\Hilbfunctor}^d_{X/S} \times_S X \to
\underline{\Hilbfunctor}^{d + 1}_{X/S}$
是次数为 $d + 1$ 的局部自由有限态射。
\end{lemma}

\begin{proof}
令 $D_{univ} \subset X \times_S \underline{\Hilbfunctor}^{d + 1}_{X/S}$
为泛对象。存在交换图
$$
\xymatrix{
\underline{\Hilbfunctor}^d_{X/S} \times_S X \ar[rr] \ar[rd] & &
D_{univ} \ar[ld] \ar@{^{(}->}[r] &
\underline{\Hilbfunctor}^{d + 1}_{X/S} \times_S X \\
& \underline{\Hilbfunctor}^{d + 1}_{X/S}
}
$$
其中上方的水平箭头把 $(D', x)$ 映到 $(D' + x, x)$。
我们断言该态射是同构，这当然证明了引理。具体地，给定 $S$ 上概形 $T$，
$D_{univ}$ 的一个 $T$-值点 $\xi$ 由一对 $\xi = (D, x)$ 给出，其中
$D \subset X_T$ 是在 $T$ 上次数为 $d + 1$ 的局部自由有限闭子概形，而
$x : T \to X$ 是其图 $x : T \to X_T$ 经由 $D$ 分解的态射。
于是由引理 \ref{pic-lemma-difference-divisors-on-curves}，可对某个在 $T$ 上次数为 $d$ 的
局部自由有限子概形 $D' \subset X_T$ 写成 $D = D' + x$。
把 $\xi = (D, x)$ 送到 $(D', x)$，便得到所需的逆映射。
\end{proof}

\begin{lemma}
\label{pic-lemma-hilb-d-smooth}
设 $X \to S$ 是相对维数为 $1$ 的光滑概形态射，并且函子
$\Hilbfunctor^d_{X/S}$ 都可表示。则概形
$\underline{\Hilbfunctor}^d_{X/S}$ 在 $S$ 上光滑，且相对维数为 $d$。
\end{lemma}

\begin{proof}
有 $\underline{\Hilbfunctor}^0_{X/S} = S$ 且
$\underline{\Hilbfunctor}^1_{X/S} = X$，所以结论对 $d = 0, 1$ 成立。
假设结论对 $d$ 成立，则 $\underline{\Hilbfunctor}^d_{X/S} \times_S X$ 在 $S$ 上光滑
（《态射》引理 \ref{morphisms-lemma-base-change-smooth} 和
\ref{morphisms-lemma-composition-smooth}）。由引理 \ref{pic-lemma-universal-object}，
$\underline{\Hilbfunctor}^d_{X/S} \times_S X \to
\underline{\Hilbfunctor}^{d + 1}_{X/S}$
是次数为 $d + 1$ 的局部自由有限态射，所以结论由《下降》引理
\ref{descent-lemma-smooth-permanence} 得出。我们略去相对维数确如所述的验证
（可以考察纤维，也可以在上面的论证中逐步追踪维数）。
\end{proof}

\noindent
现在汇总域上固有光滑曲线情形中迄今所得的全部信息。

\begin{proposition}
\label{pic-proposition-hilb-d}
设 $X$ 是域 $k$ 上几何不可约的光滑固有曲线。
\begin{enumerate}
\item 函子 $\Hilbfunctor^d_{X/k}$ 由 $k$ 上维数为 $d$ 的光滑固有簇
$\underline{\Hilbfunctor}^d_{X/k}$ 表示。
\item 对域扩张 $k'/k$，$\underline{\Hilbfunctor}^d_{X/k}$ 的 $k'$-有理点
与 $X_{k'}$ 上次数为 $d$ 的有效 Cartier 除子成 $1$ 对 $1$ 的双射。
\item 对 $d_1, d_2 \geq 0$，存在态射
$$
\underline{\Hilbfunctor}^{d_1}_{X/k}
\times_k
\underline{\Hilbfunctor}^{d_2}_{X/k}
\longrightarrow
\underline{\Hilbfunctor}^{d_1 + d_2}_{X/k}
$$
它是次数为 ${d_1 + d_2 \choose d_1}$ 的局部自由有限态射。
\end{enumerate}
\end{proposition}

\begin{proof}
由命题 \ref{pic-proposition-hilb-d-representable}（另见注
\ref{pic-remark-when-proposition-applies}）以及 $X$ 射影这一事实
（《簇》引理 \ref{varieties-lemma-dim-1-proper-projective}），
函子 $\Hilbfunctor^d_{X/k}$ 可表示。由引理 \ref{pic-lemma-hilb-d-separated}，
概形 $\underline{\Hilbfunctor}^d_{X/k}$ 在 $k$ 上分离。由引理
\ref{pic-lemma-hilb-d-smooth}，概形 $\underline{\Hilbfunctor}^d_{X/k}$ 在 $k$ 上光滑。从
$X = \underline{\Hilbfunctor}^1_{X/k}$ 出发，利用引理
\ref{pic-lemma-universal-object} 中的态射和归纳法，得到态射
$$
X^d = X \times_k X \times_k \ldots \times_k X \longrightarrow
\underline{\Hilbfunctor}^d_{X/k},\quad
(x_1, \ldots, x_d) \longrightarrow x_1 + \ldots + x_d
$$
它是次数为 $d!$ 的局部自由有限态射。由于 $X$ 在 $k$ 上固有，$X^d$ 也固有，
所以由《态射》引理 \ref{morphisms-lemma-image-proper-is-proper}，
$\underline{\Hilbfunctor}^d_{X/k}$ 在 $k$ 上固有。由于 $X$ 在 $k$ 上几何不可约，
积 $X^d$ 不可约（《簇》引理
\ref{varieties-lemma-bijection-irreducible-components}），因而其像不可约
（事实上几何不可约）。这证明了 (1)。(2) 由定义得出。(3) 由交换图
$$
\xymatrix{
X^{d_1} \times_k X^{d_2} \ar[d] \ar@{=}[r] & X^{d_1 + d_2} \ar[d] \\
\underline{\Hilbfunctor}^{d_1}_{X/k}
\times_k
\underline{\Hilbfunctor}^{d_2}_{X/k}
\ar[r] &
\underline{\Hilbfunctor}^{d_1 + d_2}_{X/k}
}
$$
以及局部自由有限态射次数的乘法性得出。
\end{proof}

\begin{remark}
\label{pic-remark-universal-object-hilb-d}
设 $X$ 如命题 \ref{pic-proposition-hilb-d}，是域 $k$ 上几何不可约的光滑固有曲线，
并令 $d \geq 0$。由引理 \ref{pic-lemma-divisors-on-curves}，泛闭对象是在
$\underline{\Hilbfunctor}^{d + 1}_{X/k}$ 上的相对有效除子
$$
D_{univ} \subset \underline{\Hilbfunctor}^{d + 1}_{X/k} \times_k X
$$
事实上，作为概形，$D_{univ}$ 同构于
$\underline{\Hilbfunctor}^d_{X/k} \times_k X$；见引理
\ref{pic-lemma-universal-object} 的证明。特别地，$D_{univ}$ 是有效 Cartier 除子，
从而得到可逆模 $\mathcal{O}(D_{univ})$。若
$[D] \in \underline{\Hilbfunctor}^{d + 1}_{X/k}$ 表示与次数为 $d + 1$ 的有效
Cartier 除子 $D \subset X$ 对应的 $k$-有理点，则 $\mathcal{O}(D_{univ})$
在纤维 $[D] \times X$ 上的限制为 $\mathcal{O}_X(D)$。
\end{remark}


\section{Picard 函子}
\label{pic-section-picard-functor}

\noindent
对任意概形 $X$，以 $\Pic(X)$ 表示可逆 $\mathcal{O}_X$-模的同构类所成的集合。
参见《模》定义 \ref{modules-definition-pic}。给定概形态射 $f : X \to Y$，
拉回给出群同态 $\Pic(Y) \to \Pic(X)$。对应
$X \leadsto \Pic(X)$ 是从概形范畴到阿贝尔群范畴的反变函子。
这个函子不可表示，但这一构造的相对版本有时可表示。

\medskip\noindent
现在为概形态射 $f : X \to S$ 定义 Picard 函子。构造背后的想法是取层
$R^1f_*\mathbf{G}_m$，其中用 fppf 拓扑计算高阶直像。展开这些定义，便得到
如下更直接的定义。

\begin{definition}
\label{pic-definition-picard-functor}
设 $\Sch_{fppf}$ 是《拓扑》定义
\ref{topologies-definition-big-small-fppf} 中的大站点。设 $f : X \to S$
是该站点中的态射。{\it Picard 函子} $\Picardfunctor_{X/S}$ 定义为函子
$$
(\Sch/S)_{fppf} \longrightarrow \textit{Sets},\quad
T \longmapsto \Pic(X_T)
$$
的 fppf 层化。若此函子可表示，则以 $\underline{\Picardfunctor}_{X/S}$
表示一个表示它的概形。
\end{definition}

\noindent
一个常用的观察是：若 $T \in \Ob((\Sch/S)_{fppf})$，则
$\Picardfunctor_{X_T/T}$ 是 $\Picardfunctor_{X/S}$ 在
$(\Sch/T)_{fppf}$ 上的限制。要确定 $\Picardfunctor_{X/S}$ 在 $S$ 上概形
$T$ 处的取值并非平凡；下面的引理对此有所帮助。

\begin{lemma}
\label{pic-lemma-flat-geometrically-connected-fibres}
设 $f : X \to S$ 如定义 \ref{pic-definition-picard-functor}。若对所有
$T \in \Ob((\Sch/S)_{fppf})$，$\mathcal{O}_T \to f_{T, *}\mathcal{O}_{X_T}$
均为同构，则对所有 $T$，序列
$$
0 \to \Pic(T) \to \Pic(X_T) \to \Picardfunctor_{X/S}(T)
$$
正合。
\end{lemma}

\begin{proof}
为简化记号，可以用 $T$ 代替 $S$、用 $X_T$ 代替 $X$，并假设 $S = T$。
设 $\mathcal{N}$ 是可逆 $\mathcal{O}_S$-模。若
$f^*\mathcal{N} \cong \mathcal{O}_X$，则由假设有
$f_*f^*\mathcal{N} \cong f_*\mathcal{O}_X \cong \mathcal{O}_S$。
由于 $\mathcal{N}$ 局部平凡，典范映射
$\mathcal{N} \to f_*f^*\mathcal{N}$ 局部为同构（因为由假设，
$\mathcal{O}_S \to f_*f^*\mathcal{O}_S$ 是同构）。因此
$\mathcal{N} \to f_*f^*\mathcal{N} \to \mathcal{O}_S$ 是同构，故
$\mathcal{N}$ 平凡。这证明了第一个箭头是单射。

\medskip\noindent
设 $\mathcal{L}$ 是位于 $\Pic(X) \to \Picardfunctor_{X/S}(S)$ 核中的可逆
$\mathcal{O}_X$-模。于是存在 fppf 覆盖 $\{S_i \to S\}$，使得
$\mathcal{L}$ 在 $X_{S_i}$ 上的拉回是平凡可逆层。选取一个平凡化截面
$s_i$。此时 $\text{pr}_0^*s_i$ 和 $\text{pr}_1^*s_j$ 都是
$X_{S_i \times_S S_j}$ 上 $\mathcal{L}$ 的平凡化截面，因而二者相差一个乘法可逆元
$$
f_{ij} \in
\Gamma(X_{S_i \times_S S_j}, \mathcal{O}_{X_{S_i \times_S S_j}}^*) =
\Gamma(S_i \times_S S_j, \mathcal{O}_{S_i \times_S S_j}^*)
$$
（等式来自关于结构层推前的假设）。这些元素当然在
$S_i \times_S S_j \times_S S_k$ 上满足余圈条件，从而给出 fppf 覆盖
$\{S_i \to S\}$ 下可逆层的下降数据。由《下降》命题
\ref{descent-proposition-fpqc-descent-quasi-coherent}，存在可逆
$\mathcal{O}_S$-模 $\mathcal{N}$ 及其在 $S_i$ 上的平凡化，其相应下降数据为
$\{f_{ij}\}$。由于从下降数据到模的函子是全忠实的（见上引命题），故
$f^*\mathcal{N} \cong \mathcal{L}$。
\end{proof}

\begin{lemma}
\label{pic-lemma-flat-geometrically-connected-fibres-with-section}
设 $f : X \to S$ 如定义 \ref{pic-definition-picard-functor}。假设 $f$ 有截面
$\sigma$，且对所有 $T \in \Ob((\Sch/S)_{fppf})$，
$\mathcal{O}_T \to f_{T, *}\mathcal{O}_{X_T}$ 均为同构。则
$$
0 \to \Pic(T) \to \Pic(X_T) \to \Picardfunctor_{X/S}(T) \to 0
$$
是分裂正合序列，其分裂由 $\sigma_T^* : \Pic(X_T) \to \Pic(T)$ 给出。
\end{lemma}

\begin{proof}
记 $K(T) = \Ker(\sigma_T^* : \Pic(X_T) \to \Pic(T))$。由于 $\sigma$ 是
$f$ 的截面，$\Pic(X_T)$ 是 $\Pic(T)$ 与 $K(T)$ 的直和。因此由引理
\ref{pic-lemma-flat-geometrically-connected-fibres}，对所有 $T$ 均有
$K(T) \subset \Picardfunctor_{X/S}(T)$。此外，由构造显然可见，
$\Picardfunctor_{X/S}$ 是预层 $K$ 的层化。为完成证明，只需说明 $K$
对 fppf 覆盖满足层条件；这在下一段完成。

\medskip\noindent
设 $\{T_i \to T\}$ 是 fppf 覆盖。设 $\mathcal{L}_i$ 是 $K(T_i)$ 中的元素，
并且对所有 $i$ 和 $j$，它们在 $K(T_i \times_T T_j)$ 中的像相同。
对每个 $i$ 选取同构
$\alpha_i : \mathcal{O}_{T_i} \to \sigma_{T_i}^*\mathcal{L}_i$，
并选取同构
$$
\varphi_{ij} :
\mathcal{L}_i|_{X_{T_i \times_T T_j}}
\longrightarrow
\mathcal{L}_j|_{X_{T_i \times_T T_j}}
$$
若映射
$$
\alpha_j|_{T_i \times_T T_j} \circ
\sigma_{T_i \times_T T_j}^*\varphi_{ij} \circ
\alpha_i|_{T_i \times_T T_j} :
\mathcal{O}_{T_i \times_T T_j} \to \mathcal{O}_{T_i \times_T T_j}
$$
不是乘以 $1$，而是乘以某个 $u_{ij}$，则可将 $\varphi_{ij}$ 乘以
$u_{ij}^{-1}$ 来修正。完成此修正后，考虑自映射
$$
\varphi_{ki}|_{X_{T_i \times_T T_j \times_T T_k}} \circ
\varphi_{jk}|_{X_{T_i \times_T T_j \times_T T_k}} \circ
\varphi_{ij}|_{X_{T_i \times_T T_j \times_T T_k}}
\quad\text{定义于}\quad
\mathcal{L}_i|_{X_{T_i \times_T T_j \times_T T_k}}
$$
它由概形 $X_{T_i \times_T T_j \times_T T_k}$ 上某个正则函数 $f_{ijk}$
的乘法给出。由对 $\varphi_{ij}$ 的选择，此映射沿 $\sigma$ 的拉回等于
乘以 $1$。再由关于 $X$ 上函数的假设，有 $f_{ijk} = 1$。于是得到 fppf 覆盖
$\{X_{T_i} \to X\}$ 的下降数据。由《下降》命题
\ref{descent-proposition-fpqc-descent-quasi-coherent}，存在可逆
$\mathcal{O}_{X_T}$-模 $\mathcal{L}$ 以及同构
$\alpha : \mathcal{O}_T \to \sigma_T^*\mathcal{L}$，其向 $X_{T_i}$ 的拉回
恢复 $(\mathcal{L}_i, \alpha_i)$（略去一个小细节）。因此
$\mathcal{L}$ 如所需地定义了 $K(T)$ 的一个对象。
\end{proof}



\section{一个可表示性判据}
\label{pic-section-representability}

\noindent
为证明 Picard 函子可表示，将使用以下判据。

\begin{lemma}
\label{pic-lemma-criterion}
设 $k$ 是域，$G : (\Sch/k)^{opp} \to \textit{Groups}$ 是函子。采用
《概形》定义 \ref{schemes-definition-representable-by-open-immersions}
中的术语，假设：
\begin{enumerate}
\item $G$ 对 Zariski 拓扑满足层性质；
\item 存在子函子 $F \subset G$，使得
\begin{enumerate}
\item $F$ 可表示；
\item $F \subset G$ 由开浸入表示；
\item 对 $k$ 的每个域扩张 $K$ 及 $g \in G(K)$，存在 $g' \in G(k)$，
使得 $g'g \in F(K)$。
\end{enumerate}
\end{enumerate}
则 $G$ 由 $k$ 上的群概形表示。
\end{lemma}

\begin{proof}
这由《概形》引理 \ref{schemes-lemma-glue-functors} 得出。具体地，取
$I = G(k)$；对 $i = g' \in I$，令 $F_i \subset G$ 为如下子函子：它把
$k$ 上的 $T$ 对应到满足 $g'g \in F(T)$ 的元素 $g \in G(T)$ 所成的集合。
乘以 $g'$ 给出 $F_i \cong F$。在乘以 $g'$ 给出的同构下，映射
$F_i \to G$ 与映射 $F \to G$ 同构，故由开浸入表示。最后，由假设 (2)(c)，
族 $(F_i)_{i \in I}$ 覆盖 $G$。因此上述引理适用，证明完成。
\end{proof}



\section{曲线的 Picard 概形}
\label{pic-section-picard-curve}

\noindent
本节将在 $k$ 为代数闭域、$X$ 为 $k$ 上光滑射影曲线时应用引理
\ref{pic-lemma-criterion}，证明 $\Picardfunctor_{X/k}$ 可表示。为此要用到
《概形的导出范畴》一章中发展的少量上同调与基变换理论。

\begin{lemma}
\label{pic-lemma-check-conditions}
设 $k$ 是域，$X$ 是 $k$ 上具有 $k$-有理点的光滑射影曲线。则引理
\ref{pic-lemma-flat-geometrically-connected-fibres-with-section} 的假设成立。
\end{lemma}

\begin{proof}
“具有 $k$-有理点”恰好意味着结构态射 $f : X \to \Spec(k)$ 有截面，
这验证了第一个条件。由《簇》引理
\ref{varieties-lemma-regular-functions-proper-variety}，
$k' = H^0(X, \mathcal{O}_X)$ 是 $k$ 的域扩张。由于 $X$ 有 $k$-有理点，
存在 $k$-代数同态 $k' \to k$，从而 $k' = k$。因为 $k$ 是域，任意态射
$T \to \Spec(k)$ 都平坦。因此由上同调与基变换（《概形上同调》引理
\ref{coherent-lemma-flat-base-change-cohomology}），
$\mathcal{O}_T \to f_{T, *}\mathcal{O}_{X_T}$ 是同构。证明完成。
\end{proof}

\noindent
设 $X$ 是域 $k$ 上的光滑射影曲线，并有一个 $k$-有理点 $\sigma$。则由引理
\ref{pic-lemma-check-conditions} 和
\ref{pic-lemma-flat-geometrically-connected-fibres-with-section}，函子
$$
\Picardfunctor_{X/k, \sigma} : (\Sch/k)^{opp} \longrightarrow \textit{Ab},\quad
T \longmapsto \Ker(\Pic(X_T) \xrightarrow{\sigma_T^*} \Pic(T))
$$
在 $(\Sch/k)_{fppf}$ 上同构于 $\Picardfunctor_{X/k}$。因此只需证明
$\Picardfunctor_{X/k, \sigma}$ 可表示。以后用记号
“$\mathcal{L} \in \Picardfunctor_{X/k, \sigma}(T)$”表示：$T$ 是 $k$ 上概形，
$\mathcal{L}$ 是可逆 $\mathcal{O}_{X_T}$-模，并且它通过 $\sigma_T$ 在 $T$
上的限制同构于 $\mathcal{O}_T$。

\begin{lemma}
\label{pic-lemma-define-open}
设 $k$ 是域，$X$ 是 $k$ 上带 $k$-有理点 $\sigma$ 的光滑射影曲线。
对 $k$ 上概形 $T$，考虑由如下 $\mathcal{L}$ 组成的子集
$F(T) \subset \Picardfunctor_{X/k, \sigma}(T)$：
$Rf_{T, *}\mathcal{L}$ 同构于置于次数 $0$ 的可逆 $\mathcal{O}_T$-模。
则 $F \subset \Picardfunctor_{X/k, \sigma}$ 是子函子，且该包含由开浸入表示。
\end{lemma}

\begin{proof}
将《概形的导出范畴》引理
\ref{perfect-lemma-open-where-cohomology-in-degree-i-rank-r-geometric}
用于 $i = 0$ 和 $r = 1$，再结合《概形》定义
\ref{schemes-definition-representable-by-open-immersions}，即得结论。
\end{proof}

\noindent
为便于继续，作如下定义。

\begin{definition}
\label{pic-definition-genus}
设 $k$ 是域，$X$ 是 $k$ 上光滑射影几何不可约曲线。$X$ 的{\it 亏格}定义为
$g = \dim_k H^1(X, \mathcal{O}_X)$。
\end{definition}

\begin{lemma}
\label{pic-lemma-open-representable}
设 $k$ 是域，$X$ 是 $k$ 上亏格为 $g$、带 $k$-有理点 $\sigma$ 的光滑射影曲线。
引理 \ref{pic-lemma-define-open} 中定义的开子函子 $F$ 由
$\underline{\Hilbfunctor}^g_{X/k}$ 的一个开子概形表示。
\end{lemma}

\begin{proof}
本证明中未注明底的积均取在 $\Spec(k)$ 上。由命题
\ref{pic-proposition-hilb-d}，概形 $H = \underline{\Hilbfunctor}^g_{X/k}$ 存在。
考虑泛除子 $D_{univ} \subset H \times X$ 及相应的可逆层
$\mathcal{O}(D_{univ})$；见注 \ref{pic-remark-universal-object-hilb-d}。
再张量上经 $\sigma_H : H \to H \times X$ 的拉回来作归一化，得到
$$
\mathcal{L}_H =
\mathcal{O}(D_{univ})
\otimes_{\mathcal{O}_{H \times X}}
\text{pr}_H^*\sigma_H^*\mathcal{O}(D_{univ})^{\otimes -1}
\in
\Picardfunctor_{X/k, \sigma}(H)
$$
由 Yoneda 引理（《范畴》引理 \ref{categories-lemma-yoneda}），可逆层
$\mathcal{L}_H$ 定义自然变换
$$
h_H \longrightarrow \Picardfunctor_{X/k, \sigma}
$$
由于 $F$ 是开子函子，存在极大开集 $W \subset H$，使得
$\mathcal{L}_H|_{W \times X}$ 属于 $F(W)$。这个开集恰好是把《概形的导出范畴》
引理 \ref{perfect-lemma-open-where-cohomology-in-degree-i-rank-r-geometric}
以 $i = 0$、$r = 1$ 应用于态射 $H \times X \to H$ 和层
$\mathcal{F} = \mathcal{O}(D_{univ})$ 所构造的开子概形。再次应用 Yoneda 引理，
得到交换图
$$
\xymatrix{
h_W \ar[d] \ar[r] & F \ar[d] \\
h_H \ar[r] & \Picardfunctor_{X/k, \sigma}
}
$$
为完成证明，将说明上方水平箭头是同构。

\medskip\noindent
设 $\mathcal{L} \in F(T) \subset \Picardfunctor_{X/k, \sigma}(T)$。
设 $\mathcal{N}$ 是满足 $Rf_{T, *}\mathcal{L} \cong \mathcal{N}[0]$ 的可逆
$\mathcal{O}_T$-模。伴随映射
$$
f_T^*\mathcal{N} \longrightarrow \mathcal{L}
\quad\text{corresponds to a section }s\text{ 属于}\quad
\mathcal{L} \otimes f_T^*\mathcal{N}^{\otimes -1}
$$
在 $X_T$ 上给出一个截面。断言：$s$ 的零点概形是 $(T \times X)/T$ 上的
相对有效 Cartier 除子 $D$，并且在 $T$ 上是次数为 $g$ 的局部自由有限态射。

\medskip\noindent
先假定断言成立并完成引理的证明。$D$ 定义态射 $m : T \to H$，使 $D$ 是
$D_{univ}$ 的拉回。于是
$$
(m \times \text{id}_X)^*\mathcal{O}(D_{univ}) \cong
\mathcal{O}_{T \times X}(D)
$$
由于 $\mathcal{O}_{T \times X}(D) \cong
\mathcal{L} \otimes f_T^*\mathcal{N}^{\otimes -1}$，可知
$(m \times \text{id}_X)^*\mathcal{L}_H$ 与 $\mathcal{L}$ 相差一个 $T$ 上可逆模的拉回。
这说明 $m : T \to H$ 通过上述开集 $W \subset H$ 分解。此外，按上法以经
$\sigma_T$ 的拉回来调整后，这些可逆模定义 $\Picardfunctor_{X/k, \sigma}(T)$
中的同一元素。利用 Yoneda 引理追图可见，$m \in h_W(T)$ 映到
$\mathcal{L} \in F(T)$。略去如下验证：规则 $F(T) \to h_W(T)$、
$\mathcal{L} \mapsto m$ 给出上述函子变换的逆。

\medskip\noindent
现在证明断言。由于 $D$ 是 $T \times X$ 的局部主闭子概形，只需证明 $D$
在 $T$ 上的各纤维是有效 Cartier 除子；参见引理
\ref{pic-lemma-divisors-on-curves} 和《除子》引理
\ref{divisors-lemma-fibre-Cartier}。由于取 $\mathcal{L}$ 的上同调与基变换交换
（《概形的导出范畴》引理
\ref{perfect-lemma-flat-proper-perfect-direct-image-general}），可归约到
$T = \Spec(K)$，其中 $K/k$ 是域扩张。此时 $\mathcal{L}$ 是 $X_K$ 上的可逆层，
并满足 $H^0(X_K, \mathcal{L}) = K$ 和 $H^1(X_K, \mathcal{L}) = 0$。因此
$$
\deg(\mathcal{L}) = \chi(X_K, \mathcal{L}) - \chi(X_K, \mathcal{O}_{X_K})
= 1 - (1 - g) = g
$$
参见《簇》定义 \ref{varieties-definition-degree-invertible-sheaf}。最后只须说明
$\mathcal{L}$ 的非零截面在 $X_K$ 上定义有效 Cartier 除子；这是显然的。
\end{proof}

\begin{lemma}
\label{pic-lemma-twist-with-general-divisor}
设 $k$ 是可分闭域，$X$ 是 $k$ 上亏格为 $g$ 的光滑射影曲线。设 $K/k$ 是域扩张，
$\mathcal{L}$ 是 $X_K$ 上的可逆层。则存在 $X$ 上的可逆层 $\mathcal{L}_0$，使得
$\dim_K H^0(X_K,
\mathcal{L} \otimes_{\mathcal{O}_{X_K}} \mathcal{L}_0|_{X_K}) = 1$ 且
$\dim_K H^1(X_K,
\mathcal{L} \otimes_{\mathcal{O}_{X_K}} \mathcal{L}_0|_{X_K}) = 0$。
\end{lemma}

\begin{proof}
本证明是《簇》引理 \ref{varieties-lemma-general-degree-g-line-bundle}
证明的一个变式。建议读者先阅读该证明。

\medskip\noindent
先选取充足可逆层 $\mathcal{L}_0$，并对某个 $n \gg 0$，以
$\mathcal{L} \otimes_{\mathcal{O}_{X_K}} \mathcal{L}_0^{\otimes n}|_{X_K}$
代替 $\mathcal{L}$。这样便可假设
$H^0(X_K, \mathcal{L}) \not = 0$ 且 $H^1(X_K, \mathcal{L}) = 0$。
事实上，消灭性来自《概形上同调》引理
\ref{coherent-lemma-vanshing-gives-ample}，而非零性来自张量积的次数为
$\gg 0$。最后对 $t = \dim_K H^0(X_K, \mathcal{L})$ 作降阶归纳。
基本情形 $t = 1$ 显然。假设 $t > 1$。

\medskip\noindent
注意，对 $X$ 的 $k$-有理点 $x$，逆像 $x_K$ 是 $X_K$ 的 $K$-有理点。
此外，由《簇》引理 \ref{varieties-lemma-smooth-separable-closed-points-dense}，
$k$-有理点有无穷多个。因此各点 $x_K$ 在 $X_K$ 中构成 Zariski 稠密的点集。

\medskip\noindent
取非零 $s \in H^0(X_K, \mathcal{L})$。由上一段，存在 $k$-有理点 $x$，
使 $s$ 在 $x_K$ 处不消失。设 $\mathcal{I}$ 是
$i : x_K \to X_K$ 的理想层，如《簇》引理
\ref{varieties-lemma-regular-point-on-curve}。考察短正合序列
$$
0 \to \mathcal{I} \otimes_{\mathcal{O}_{X_K}} \mathcal{L} \to
\mathcal{L} \to i_*i^*\mathcal{L} \to 0
$$
注意 $H^0(X_K, i_*i^*\mathcal{L}) = H^0(x_K, i^*\mathcal{L})$ 在 $K$ 上维数为
$1$。由于 $s$ 在 $x$ 处不消失，映射
$$
H^0(X_K, \mathcal{L}) \longrightarrow H^0(X, i_*i^*\mathcal{L})
$$
是满射。因此
$\dim_K H^0(X_K, \mathcal{I} \otimes_{\mathcal{O}_{X_K}} \mathcal{L}) = t - 1$。
最后，上同调长正合序列还给出
$H^1(X_K, \mathcal{I} \otimes_{\mathcal{O}_{X_K}} \mathcal{L}) = 0$，
从而完成归纳步骤。
\end{proof}

\begin{proposition}
\label{pic-proposition-pic-curve}
设 $k$ 是可分闭域，$X$ 是 $k$ 上光滑射影曲线。Picard 函子
$\Picardfunctor_{X/k}$ 可表示。
\end{proposition}

\begin{proof}
由于 $k$ 可分闭，$X$ 有 $k$-有理点 $\sigma$；见《簇》引理
\ref{varieties-lemma-smooth-separable-closed-points-dense}。如上所述，只需证明
对沿 $\sigma$ 平凡的可逆模进行分类的函子 $\Picardfunctor_{X/k, \sigma}$
可表示。为此验证引理 \ref{pic-lemma-criterion} 的条件 (1)、(2)(a)、(2)(b)
和 (2)(c)。

\medskip\noindent
函子 $\Picardfunctor_{X/k, \sigma}$ 同构于 $\Picardfunctor_{X/k}$，故对 fppf
拓扑满足层条件。更准确地说，我们已在引理
\ref{pic-lemma-flat-geometrically-connected-fibres-with-section} 的证明中证明了
$\Picardfunctor_{X/k, \sigma}$ 的层条件；由引理
\ref{pic-lemma-check-conditions}，该证明适用。这就证明了条件 (1)。

\medskip\noindent
取引理 \ref{pic-lemma-define-open} 中定义的 $F$ 为子函子。条件 (2)(b) 随即成立；
条件 (2)(a) 是引理 \ref{pic-lemma-open-representable}；条件 (2)(c) 是引理
\ref{pic-lemma-twist-with-general-divisor}。
\end{proof}

\noindent
事实上，上述证明还给出更多信息，现汇总如下。

\begin{lemma}
\label{pic-lemma-picard-pieces}
设 $k$ 是可分闭域，$X$ 是 $k$ 上亏格为 $g$ 的光滑射影曲线。
\begin{enumerate}
\item $\underline{\Picardfunctor}_{X/k}$ 是一族 $g$ 维光滑固有簇
$\underline{\Picardfunctor}^d_{X/k}$ 的不交并；
\item $\underline{\Picardfunctor}^d_{X/k}$ 的 $k$-点对应于次数为 $d$ 的可逆
$\mathcal{O}_X$-模；
\item $\underline{\Picardfunctor}^0_{X/k}$ 是开闭子群概形；
\item 对 $d \geq 0$，存在典范态射
$\gamma_d :
\underline{\Hilbfunctor}^d_{X/k} \to \underline{\Picardfunctor}^d_{X/k}$；
\item 态射 $\gamma_d$ 在 $d \geq g$ 时满，在 $d \geq 2g - 1$ 时光滑；
\item 态射
$\underline{\Hilbfunctor}^g_{X/k} \to \underline{\Picardfunctor}^g_{X/k}$
是双有理的。
\end{enumerate}
\end{lemma}

\begin{proof}
选取 $X$ 的 $k$-有理点 $\sigma$。回忆 $\Picardfunctor_{X/k}$ 同构于函子
$\Picardfunctor_{X/k, \sigma}$。由《概形的导出范畴》引理
\ref{perfect-lemma-chi-locally-constant-geometric}，对每个
$d \in \mathbf{Z}$，存在开子函子
$$
\Picardfunctor^d_{X/k, \sigma} \subset \Picardfunctor_{X/k, \sigma}
$$
它在 $k$ 上概形 $T$ 处的取值由满足
$\chi(X_t, \mathcal{L}_t) = d + 1 - g$ 的
$\mathcal{L} \in \Picardfunctor_{X/k, \sigma}(T)$ 组成；而且作为 fppf 层有
$$
\Picardfunctor_{X/k, \sigma} =
\coprod\nolimits_{d \in \mathbf{Z}} \Picardfunctor^d_{X/k, \sigma}
$$
因此，概形 $\underline{\Picardfunctor}_{X/k}$（它由命题
\ref{pic-proposition-pic-curve} 存在）有相应分解
$$
\underline{\Picardfunctor}_{X/k, \sigma} =
\coprod\nolimits_{d \in \mathbf{Z}} \underline{\Picardfunctor}^d_{X/k, \sigma}
$$
其中 $\underline{\Picardfunctor}^d_{X/k, \sigma}$ 的点对应于 $X$ 上次数为 $d$
的可逆模的同构类。

\medskip\noindent
固定 $d \geq 0$。由 $\underline{\Hilbfunctor}^d_{X/k} \times_k X$ 上的可逆层
$\mathcal{O}(D_{univ})$（注 \ref{pic-remark-universal-object-hilb-d}），再由 Yoneda
引理（《范畴》引理 \ref{categories-lemma-yoneda}），得到态射
$$
\gamma_d :
\underline{\Hilbfunctor}^d_{X/k}
\longrightarrow
\underline{\Picardfunctor}^d_{X/k}
$$
命题 \ref{pic-proposition-pic-curve} 和引理 \ref{pic-lemma-open-representable}
中对 $X/k$ 的 Picard 函子可表示性的证明表明，$\gamma_g$ 在
$\underline{\Hilbfunctor}^g_{X/k}$ 的一个非空开集上诱导开浸入。此外，该证明还表明，
此开集在群概形 $\underline{\Picardfunctor}_{X/k}$ 的 $k$-有理点作用下的平移构成
开覆盖。由于 $\underline{\Hilbfunctor}^g_{X/K}$ 在 $k$ 上光滑且维数为 $g$
（命题 \ref{pic-proposition-hilb-d}），可知群概形
$\underline{\Picardfunctor}_{X/k}$ 在 $k$ 上光滑且维数为 $g$。

\medskip\noindent
由《群胚》引理 \ref{groupoids-lemma-group-scheme-over-field-separated}，
$\underline{\Picardfunctor}_{X/k}$ 是分离的。因此对每个 $d \geq 0$，
$\gamma_d$ 的像是 $k$ 上固有簇（《态射》引理
\ref{morphisms-lemma-scheme-theoretic-image-is-proper}）。

\medskip\noindent
设 $d \geq g$。对任意域扩张 $K/k$ 以及任意次数为 $d$ 的可逆
$\mathcal{O}_{X_K}$-模 $\mathcal{L}$，有
$\chi(X_K, \mathcal{L}) = d + 1 - g > 0$。故 $\mathcal{L}$ 有非零截面，
从而对某个次数为 $d$ 的除子 $D \subset X_K$，有
$\mathcal{L} = \mathcal{O}_{X_K}(D)$。因此 $\gamma_d$ 是满射。

\medskip\noindent
综合上述事实可见，当 $d \geq g$ 时，
$\underline{\Picardfunctor}^d_{X/k}$ 固有。这完成了 (2) 的证明，因为现在已知
$d \geq g$ 时 $\underline{\Picardfunctor}^d_{X/k}$ 固有，而通过平移可知所有
$\underline{\Picardfunctor}^d_{X/k}$ 都固有。

\medskip\noindent
还须证明 $d \geq 2g - 1$ 时 $\gamma_d$ 光滑。取次数为 $d$ 的可逆
$\mathcal{O}_X$-模 $\mathcal{L}$。与 $\mathcal{L}$ 对应的点的纤维是
$$
Z = \{D \subset X \mid \mathcal{O}_X(D) \cong \mathcal{L}\} \subset
\underline{\Hilbfunctor}^d_{X/k}
$$
并带有其自然概形结构。任意同构 $\mathcal{O}_X(D) \to \mathcal{L}$ 只在乘以
非零标量的意义下唯一，故典范截面 $1 \in \mathcal{O}_X(D)$ 映到截面
$s \in \Gamma(X, \mathcal{L})$，后者也只在乘以非零标量的意义下唯一。由此得到态射
$$
Z \longrightarrow
\text{Proj}(\text{Sym}(\Gamma(X, \mathcal{L})^*))
$$
（由于我们的约定，这里出现对偶）。此态射是同构，因为给定
$\mathcal{L}$ 的截面，可以取相应的有效 Cartier 除子；换言之，可以构造所示态射的逆。
略去精确表述与证明。对每个次数为 $d \geq 2g - 1$ 的 $\mathcal{L}$，由《簇》引理
\ref{varieties-lemma-vanishing-degree-2g-and-1-line-bundle} 有
$\dim H^0(X, \mathcal{L}) = d + 1 - g$，故
$\text{Proj}(\text{Sym}(\Gamma(X, \mathcal{L})^*))
\cong \mathbf{P}^{d - g}_k$。于是
$\dim(Z) = \dim(\mathbf{P}^{d - g}_k) = d - g$。因此态射 $\gamma_d$ 的所有纤维
的维数都等于 $\underline{\Hilbfunctor}^d_{X/k}$ 与
$\underline{\Picardfunctor}^d_{X/k}$ 的维数之差。由《代数》引理
\ref{algebra-lemma-CM-over-regular-flat}，$\gamma_d$ 平坦。又因其纤维光滑，
由《态射》引理 \ref{morphisms-lemma-smooth-flat-smooth-fibres}，
$\gamma_d$ 光滑。
\end{proof}





\section{关于 Picard 群的一些注记}
\label{pic-section-remarks-picard}

\noindent
本节继续《簇》第 \ref{varieties-section-picard-groups} 节的讨论；这一讨论还将在
《代数曲线》第 \ref{curves-section-torsion-in-pic} 节继续。

\begin{lemma}
\label{pic-lemma-pic-descends}
设 $k$ 是域，$X$ 是 $k$ 上拟紧且拟分离的概形，并满足
$H^0(X, \mathcal{O}_X) = k$。若 $X$ 有 $k$-有理点，则对任意 Galois 扩张
$k'/k$ 均有
$$
\Pic(X) = \Pic(X_{k'})^{\text{Gal}(k'/k)}
$$
此外，$\text{Gal}(k'/k)$ 在 $\Pic(X_{k'})$ 上的作用连续。
\end{lemma}

\begin{proof}
由于 $\text{Gal}(k'/k) = \text{Aut}(k'/k)$，它右作用于 $\Spec(k')$，从而右作用于
$X_{k'} = X \times_{\Spec(k)} \Spec(k')$；又因 $\Pic(-)$ 是反变函子，它左作用于
$\Pic(X_{k'})$。若 $k'/k$ 是无限 Galois 扩张，则把它写成有限 Galois 扩张的滤过余极限
$k' = \colim k'_\lambda$；见《域》引理
\ref{fields-lemma-infinite-galois-limit}。此时
$X_{k'} = \lim X_{k_\lambda}$（如《极限》第 \ref{limits-section-limits} 节），
并由《极限》引理 \ref{limits-lemma-descend-invertible-modules} 得到
$$
\Pic(X_{k'}) = \colim \Pic(X_{k_\lambda})
$$
此外，由《簇》引理 \ref{varieties-lemma-change-fields-pic}，这个阿贝尔群系统中的
过渡映射都是单射。因此 $\Pic(X_{k'})$ 的每个元素都被某个开子群
$\text{Gal}(k'/k_\lambda)$ 固定，这恰好说明作用连续。过渡映射的单射性还表明，
为证明关于不动点的断言，只需处理 $k'/k$ 为有限 Galois 扩张的情形。

\medskip\noindent
假设 $k'/k$ 是有限 Galois 扩张，Galois 群为
$G = \text{Gal}(k'/k)$。设 $\mathcal{L}$ 是 $\Pic(X_{k'})$ 中被 $G$ 固定的元素。
将用 Galois 下降（《下降》引理 \ref{descent-lemma-galois-descent}）证明
$\mathcal{L}$ 是 $X$ 上可逆层的拉回。回忆
$f_\sigma = \text{id}_X \times \Spec(\sigma) : X_{k'} \to X_{k'}$，并且
$\sigma$ 通过沿 $f_\sigma$ 拉回来作用于 $\Pic(X_{k'})$。由于 $\mathcal{L}$
被 $G$ 的作用固定，对每个 $\sigma \in G$ 可选取同构
$\varphi_\sigma : \mathcal{L} \to f_\sigma^*\mathcal{L}$。问题在于，尚不清楚能否选择
$\varphi_\sigma$ 使余圈条件
$\varphi_{\sigma\tau} = f_\sigma^*\varphi_\tau \circ \varphi_\sigma$
成立。为说明这是可能的，利用 $X$ 有 $k$-有理点 $x \in X(k)$。当然，$x$
相应地确定一个 $k'$-有理点 $x' \in X_{k'}$，且对所有 $\sigma$，$f_\sigma$
都固定它。在 $x'$ 处 $\mathcal{L}$ 的纤维中取非零元素 $s$；该纤维是
$1$ 维 $k' = \kappa(x')$-向量空间
$$
\mathcal{L}_{x'} \otimes_{\mathcal{O}_{X_{k'}, x'}} \kappa(x').
$$
于是 $f_\sigma^*s$ 是 $x'$ 处 $f_\sigma^*\mathcal{L}$ 的纤维中的非零元素。
由于可以把 $\varphi_\sigma$ 乘以 $(k')^*$ 中的元素，可假设
$\varphi_\sigma$ 把 $s$ 映到 $f_\sigma^*s$。此时
$\varphi_{\sigma\tau}$ 与 $f_\sigma^*\varphi_\tau \circ \varphi_\sigma$
都把 $s$ 映到 $f_{\sigma\tau}^*s = f_\tau^*f_\sigma^*s$。由于
$H^0(X_{k'}, \mathcal{O}_{X_{k'}}) = k'$，这两个同构必相同（一个是另一个乘以
全局可逆元，而它们在 $x'$ 处相同）。证明完成。
\end{proof}

\begin{lemma}
\label{pic-lemma-torsion-descends}
设 $k$ 是特征 $p > 0$ 的域，$X$ 是 $k$ 上拟紧且拟分离的概形，并满足
$H^0(X, \mathcal{O}_X) = k$。设整数 $n$ 与 $p$ 互素。则对任意纯不可分扩张
$k'/k$，映射
$$
\Pic(X)[n] \longrightarrow \Pic(X_{k'})[n]
$$
为双射。
\end{lemma}

\begin{proof}
首先，由《簇》引理 \ref{varieties-lemma-change-fields-pic}，映射
$\Pic(X) \to \Pic(X_{k'})$ 是单射。因此只须证明引理中的映射满。设
$\mathcal{L}$ 是可逆 $\mathcal{O}_{X_{k'}}$-模，并且它在
$\Pic(X_{k'})$ 中的阶整除 $n$。选取可逆模同构
$\alpha : \mathcal{L}^{\otimes n} \to \mathcal{O}_{X_{k'}}$。
将证明偶对 $(\mathcal{L}, \alpha)$ 可下降到 $X$。

\medskip\noindent
置 $A = k' \otimes_k k'$。由于 $k'/k$ 纯不可分，乘法映射 $A \to k'$
的核是 $A$ 的局部幂零理想 $I$。注意
$$
X_A = X \times_{\Spec(k)} \Spec(A) = X_{k'} \times_X X_{k'}
$$
带有两个投影 $\text{pr}_i : X_A \to X_{k'}$，$i = 0, 1$；它们在 $A/I$
上相同。因此，可逆模 $\mathcal{L}_i = \text{pr}_i^*\mathcal{L}$
在闭子概形 $X_{A/I} = X_{k'}$ 上相同。由于 $X_{A/I} \to X_A$ 是加厚，
且 $\mathcal{L}_i$ 是 $n$-挠的，由《态射进阶》引理
\ref{more-morphisms-lemma-torsion-pic-thickening}，存在同构
$\varphi : \mathcal{L}_0 \to \mathcal{L}_1$。可以选择 $\varphi$ 使其模 $I$
约化为恒等映射。事实上，由《概形上同调》引理
\ref{coherent-lemma-flat-base-change-cohomology}，
$H^0(X, \mathcal{O}_X) = k$ 蕴含
$H^0(X_{k'}, \mathcal{O}_{X_{k'}}) = k'$；又因 $A \to k'$ 满，可以把
$\varphi$ 乘以 $A$ 中适当元素来调整。考虑映射
$$
\lambda : \mathcal{O}_{X_A}
\xrightarrow{\text{pr}_0^*\alpha^{-1}}
\mathcal{L}_0^{\otimes n}
\xrightarrow{\varphi^{\otimes n}}
\mathcal{L}_1^{\otimes n}
\xrightarrow{\text{pr}_0^*\alpha}
\mathcal{O}_{X_A}
$$
由于 $H^0(X_A, \mathcal{O}_{X_A}) = A$（同一引理），可把 $\lambda$
视为 $A$ 的元素。因为 $\varphi$ 模 $I$ 约化为恒等映射，所以
$\lambda = 1 \bmod I$。于是 $\lambda$ 在 $1 + I$ 中有唯一的 $n$ 次方根
（《代数》引理 \ref{algebra-lemma-lift-nth-roots}）；把 $\varphi$ 乘以该根的逆，
便得到 $\lambda = 1$。断言 $(\mathcal{L}, \varphi)$ 是 fpqc 覆盖
$\{X_{k'} \to X\}$ 的下降数据（《下降》定义
\ref{descent-definition-descent-datum-quasi-coherent}）。若断言成立，则由《下降》命题
\ref{descent-proposition-fpqc-descent-quasi-coherent}，$\mathcal{L}$ 是某个可逆
$\mathcal{O}_X$-模 $\mathcal{N}$ 的拉回。Picard 群间映射的单射性表明，
$\mathcal{N}$ 是 $\Pic(X)$ 中与 $\mathcal{L}$ 同阶的挠元素。

\medskip\noindent
证明断言。须验证
$$
\text{pr}_{12}^*\varphi \circ
\text{pr}_{01}^*\varphi =
\text{pr}_{02}^*\varphi
\quad\text{定义于}\quad
X_{k'} \times_X X_{k'} \times_X X_{k'} = X_{k' \otimes_k k' \otimes_k k'}
$$
与前面一样，对角态射
$\Delta : X_{k'} \to X_{k' \otimes_k k' \otimes_k k'}$ 是加厚。等式的左右两边
是映射 $a, b : p_0^*\mathcal{L} \to p_2^*\mathcal{L}$，并分别与
$p_0^*\alpha$ 和 $p_2^*\alpha$ 相容；这里
$p_i : X_{k' \otimes_k k' \otimes_k k'} \to X_{k'}$ 是投影态射。
最后，$a, b$ 沿 $\Delta$ 的拉回是同一映射。仿射局部地（在局部平凡化中），
这意味着 $a, b$ 由乘以可逆函数给出；这些函数模某个局部幂零理想约化为同一函数，
并且有相同的 $n$ 次幂。于是由《代数》引理
\ref{algebra-lemma-lift-nth-roots}，这些函数相同。
\end{proof}
% Generated mechanically from the frozen weil.tex target.
% Do not edit this generated file; edit the bound locale source instead.

% OK, start here.
%

\setcounter{chapter}{44}
\chapter{Weil 上同调理论}



\phantomsection
\label{weil-section-phantom}




\section{引言}
\label{weil-section-introduction}

\noindent
本章讨论基域上光滑射影概形的 Weil 上同调理论。简言之，这里的上同调理论
$H^*$ 具有 K\"unneth 公式、Poincar\'e 对偶和周期类，并满足适当的相容性。
需要提醒读者，文献中对于何者构成“Weil 上同调理论”并无普遍一致的约定。

\medskip\noindent
阅读本章之前，读者应先参阅《范畴》第 \ref{categories-section-monoidal} 节和
《同调代数》第 \ref{homology-section-monoidal} 节；其中定义了（对称）幺半范畴，
并发展了本章所需的基本语言。借助这些语言，我们在第
\ref{weil-section-correspondences} 节构造一个对称幺半分次范畴，其对象是光滑射影概形，
态射是对应。在第 \ref{weil-section-chow-motives} 节加入投影算子的像并将 Lefschetz
动机可逆化，从而得到 Chow 动机的对称幺半 Karoubian 范畴 $M_k$。该范畴配有反变函子
$$
h : \{\text{smooth projective schemes over }k\} \longrightarrow M_k
$$
如下文所见，Weil 上同调理论的一个关键性质是它通过 $h$ 分解。

\medskip\noindent
首先，在基域代数闭的情形下，定义所谓“经典 Weil 上同调理论”；见第
\ref{weil-section-axioms-classical} 节。这个概念与
\cite[Section 1.2]{Kleiman-cycles} 中引入的概念相同，也与
\cite[page 65]{Kleiman-motives} 中引入的概念一致。不过，它并不先验地与
\cite[page 10]{Kleiman-standard} 中引入的概念一致，因为该处作者添加了两个
Lefschetz 型公理，而尚不知道本章所定义的任一经典 Weil 上同调理论是否满足这些公理。
在第 \ref{weil-section-axioms-classical} 节末，将证明经典 Weil 上同调理论具有形式
$H^* = G \circ h$，其中 $G$ 是从 $M_k$ 到 $H^*$ 的系数域上分次向量空间范畴的
对称幺半函子。

\medskip\noindent
在第 \ref{weil-section-cycles-nonclosed} 节证明若干关于非闭域上周期群的引理；
讨论任意域上光滑射影概形的 Weil 上同调理论时将会用到它们。

\medskip\noindent
在一般基域 $k$ 上，对 Weil 上同调理论 $H^*$ 采用这些公理的动机如下：
\begin{enumerate}
\item 对某个从 $M_k$ 到 $H^*$ 的系数域 $F$ 上分次向量空间范畴的对称幺半函子
$G$，有 $H^* = G \circ h$；
\item $G$ 应把 Tate 动机（Lefschetz 动机的逆）送到位于次数 $-2$ 的 $1$ 维向量空间
$F(1)$；
\item 当 $k$ 代数闭时，在选定 $F(1)$ 的一个基元素后，应恢复第
\ref{weil-section-axioms-classical} 节所讨论的概念。
\end{enumerate}
首先在第 \ref{weil-section-axioms} 节分析前两个条件。在第
\ref{weil-section-further} 节发展若干进一步结果后，于第 \ref{weil-section-old} 节加入
得到性质 (3) 所需的公理。

\medskip\noindent
最后在第 \ref{weil-section-c1} 节详述 Weil 上同调理论的另一种处理方式：使用第一
Chern 类映射代替周期类。后续章节中要证明某些上同调理论是 Weil 上同调理论时，
这种方式最为适用；见《de Rham 上同调》第 \ref{derham-section-weil} 节。





\section{约定与记号}
\label{weil-section-conventions}

\noindent
设 $F$ 是域。本章把 $F$-分次向量空间范畴视为 $F$-线性对称幺半范畴，
其结合约束取通常形式，交换约束则含有符号。参见《同调代数》例
\ref{homology-example-graded-vector-spaces}。

\medskip\noindent
设 $R$ 是环。本章所谓{\it 分次交换 $R$-代数} $A$，是微分为零的交换微分分次
$R$-代数（《微分分次代数》定义 \ref{dga-definition-dga} 和
\ref{dga-definition-cdga}）。因此，$A$ 是一个 $R$-模，并由 $R$-子模赋予分次
$A = \bigoplus_{n \in \mathbf{Z}} A^n$。本章把 $R$-双线性乘法
$$
A^n \times A^m \longrightarrow A^{n + m},\quad
\alpha \times \beta \longmapsto \alpha \cup \beta
$$
称为{\it 杯积}。当 $\alpha \in A^n$ 且 $\beta \in A^m$ 时，交换约束为
$\alpha \cup \beta = (-1)^{nm} \beta \cup \alpha$。最后，存在乘法单位元
$1 \in A^0$；等价地，存在与 $A$ 上 $R$-模结构相容的加性及乘性映射
$R \to A^0$。

\medskip\noindent
设 $k$ 是域，$X$ 是 $k$ 上有限型概形。《Chow 同调》定义
\ref{chow-definition-rational-equivalence} 已定义 $X$ 上 $k$ 维周期模有理等价所得的
Chow 群 $\CH_k(X)$。若 $X$ 正规或 Cohen--Macaulay，还可考虑余维数为 $p$ 的
周期所成的 Chow 群 $\CH^p(X)$（《Chow 同调》第
\ref{chow-section-cycles-codimension} 节）；此时 $[X] \in \CH^0(X)$ 表示
$X$ 的“基本类”，见《Chow 同调》注 \ref{chow-remark-fundamental-class}。
若 $X$ 光滑，且 $\alpha$、$\beta$ 是 $X$ 上周期，则
$\alpha \cdot \beta$ 表示 $\alpha$ 与 $\beta$ 的交积；见《Chow 同调》第
\ref{chow-section-intersection-product} 节。













\section{对应}
\label{weil-section-correspondences}

\noindent
设 $k$ 是域。对 $k$ 上概形 $X$、$Y$，以 $X \times Y$ 表示 $X$ 与 $Y$ 在 $k$ 上概形范畴
中的积。本节构造一个 $\mathbf{Q}$ 上分次范畴，其对象是 $k$ 上光滑射影概形，
态射是对应。

\medskip\noindent
设 $X$、$Y$ 是 $k$ 上光滑射影概形。设 $X = \coprod X_d$ 是 $X$ 分解为等维开闭
子概形的分解，并满足 $\dim(X_d) = d$。定义{\it 从 $X$ 到 $Y$ 的 $r$ 次对应}
所成的 $\mathbf{Q}$-向量空间为
$$
\text{Corr}^r(X, Y) =
\bigoplus\nolimits_d \CH^{d + r}(X_d \times Y) \otimes \mathbf{Q}
\subset
\CH^*(X \times Y) \otimes \mathbf{Q}
$$
给定 $c \in \text{Corr}^r(X, Y)$ 和 $\beta \in \CH_j(Y) \otimes \mathbf{Q}$，
可用公式
$$
c^*(\beta) = \text{pr}_{1, *}(c \cdot \text{pr}_2^*\beta)
\quad\text{于}\quad
\CH_{j - r}(X) \otimes \mathbf{Q}
$$
定义 $\beta$ 沿 $c$ 的{\it 拉回}。这是有意义的：$\text{pr}_2$ 在
$X_d \times Y$ 上平坦且相对维数为 $d$，所以 $\text{pr}_2^*\beta$ 是
$X_d \times Y$ 上维数为 $d + j$ 的周期；于是
$c \cdot \text{pr}_2^*\beta$ 是 $X_d \times Y$ 上维数为 $j - r$ 的周期，而沿固有态射
$\text{pr}_1$ 的推前仍是同维数周期。类似地，改用余维数分次后，给定
$\alpha \in \CH^i(X) \otimes \mathbf{Q}$，可用公式
$$
c_*(\alpha) = \text{pr}_{2, *}(c \cdot \text{pr}_1^*\alpha)
\quad\text{于}\quad
\CH^{i + r}(Y) \otimes \mathbf{Q}
$$
定义 $\alpha$ 沿 $c$ 的{\it 推前}。这是有意义的：$\text{pr}_1^*\alpha$ 是
$X \times Y$ 上余维数为 $i$ 的周期，故 $c \cdot \text{pr}_1^*\alpha$
是 $X_d \times Y$ 上余维数为 $i + d + r$ 的周期，其推前是 $Y$ 上余维数为
$i + r$ 的周期。

\medskip\noindent
给定 $k$ 上三个光滑射影概形 $X, Y, Z$，定义{\it 对应的复合}
$$
\text{Corr}^s(Y, Z)
\times
\text{Corr}^r(X, Y)
\longrightarrow
\text{Corr}^{r + s}(X, Z)
$$
为
$$
(c', c)
\longmapsto
c' \circ c =
\text{pr}_{13, *}(\text{pr}_{12}^*c \cdot \text{pr}_{23}^*c')
$$
其中 $\text{pr}_{12} : X \times Y \times Z \to X \times Y$ 是投影，
$\text{pr}_{13}$、$\text{pr}_{23}$ 类似。

\begin{lemma}
\label{weil-lemma-composition-correspondences}
对应具有以下性质：
\begin{enumerate}
\item 对应的复合是 $\mathbf{Q}$-双线性的，并满足结合律；
\item 存在典范同构
$$
\CH_{-r}(X) \otimes \mathbf{Q} = \text{Corr}^r(X, \Spec(k))
$$
使沿对应的拉回对应于复合；
\item 存在典范同构
$$
\CH^r(X) \otimes \mathbf{Q} = \text{Corr}^r(\Spec(k), X)
$$
使沿对应的推前对应于复合；
\item 对应的复合与周期的推前和拉回相容。
\end{enumerate}
\end{lemma}

\begin{proof}
双线性立即来自推前与拉回的线性以及交积的双线性。为证明结合律，设有
$X, Y, Z, W$ 以及 $c \in \text{Corr}(X, Y)$、$c' \in \text{Corr}(Y, Z)$、
$c'' \in \text{Corr}(Z, W)$。则
\begin{align*}
c'' \circ (c' \circ c)
& =
\text{pr}^{134}_{14, *}(
\text{pr}^{134, *}_{13}
\text{pr}^{123}_{13, *}(\text{pr}^{123, *}_{12}c \cdot
\text{pr}^{123, *}_{23}c')
\cdot \text{pr}^{134, *}_{34}c'') \\
& =
\text{pr}^{134}_{14, *}(
\text{pr}^{1234}_{134, *}
\text{pr}^{1234, *}_{123}(\text{pr}^{123, *}_{12}c \cdot
\text{pr}^{123, *}_{23}c')
\cdot \text{pr}^{134, *}_{34}c'') \\
& =
\text{pr}^{134}_{14, *}(
\text{pr}^{1234}_{134, *}
(\text{pr}^{1234, *}_{12}c \cdot
\text{pr}^{1234, *}_{23}c')
\cdot \text{pr}^{134, *}_{34}c'') \\
& =
\text{pr}^{134}_{14, *}
\text{pr}^{1234}_{134, *}
((\text{pr}^{1234, *}_{12}c \cdot
\text{pr}^{1234, *}_{23}c')
\cdot \text{pr}^{1234, *}_{34}c'') \\
& =
\text{pr}^{1234}_{14, *}(
(\text{pr}^{1234, *}_{12}c \cdot
\text{pr}^{1234, *}_{23}c') \cdot
\text{pr}^{1234, *}_{34}c'')
\end{align*}
这里用记号
$$
p^{1234}_{134} : X \times Y \times Z \times W \to X \times Z \times W
\quad\text{且}\quad
p^{134}_{14} : X \times Z \times W \to X \times W
$$
表示投影，其他指标同理。第一个等式是复合的定义。第二个等式来自
$\text{pr}^{134, *}_{13} \text{pr}^{123}_{13, *} =
\text{pr}^{1234}_{134, *} \text{pr}^{1234, *}_{123}$，见《Chow 同调》引理
\ref{chow-lemma-flat-pullback-proper-pushforward}。第三个等式成立，是因为交积与
$p^{1234}_{123}$ 的 Gysin 映射（即平坦拉回）交换；见《Chow 同调》引理
\ref{chow-lemma-lci-gysin-product}。第四个等式来自 $p^{1234}_{134}$ 的投影公式；
见《Chow 同调》引理 \ref{chow-lemma-projection-formula}。第四个等式表示固有推前
与复合相容；见《Chow 同调》引理 \ref{chow-lemma-compose-pushforward}。
又由《Chow 同调》引理 \ref{chow-lemma-associative}，交积满足结合律，
故对应复合的结合律得证。

\medskip\noindent
略去 (2)、(3) 的证明；其实质只是仔细记录各周期所在之处及其（余）维数。

\medskip\noindent
周期推前、拉回的断言意指
$(c' \circ c)^*(\alpha) = c^*((c')^*(\alpha))$ 和
$(c' \circ c)_*(\alpha) = (c')_*(c_*(\alpha))$。合并 (1)、(2)、(3) 即得。
\end{proof}

\begin{example}
\label{weil-example-graph-correspondence}
设 $f : Y \to X$ 是 $k$ 上光滑射影概形之间的态射。记
$\Gamma_f \subset X \times Y$ 为 $f$ 的图。更确切地，$\Gamma_f$ 是闭浸入
$$
(f, \text{id}_Y) : Y \longrightarrow X \times Y
$$
的像。设 $X = \coprod X_d$ 是 $X$ 分解为等维且维数为 $d$ 的开闭部分
$X_d$ 的分解。于是 $\Gamma_f \cap (X_d \times Y)$ 的纯余维数为 $d$。因此
$[\Gamma_f] \in \CH^*(X \times Y) \otimes \mathbf{Q}$ 属于
$\text{Corr}^0(X \times Y)$；换言之，$[\Gamma_f]$ 是从 $X$ 到 $Y$ 的
$0$ 次对应。
\end{example}

\begin{lemma}
\label{weil-lemma-category-correspondences}
$k$ 上光滑射影概形，以对应为态射、以上述对应复合为复合，构成一个
$\mathbf{Q}$ 上分次范畴
（《微分分次代数》定义 \ref{dga-definition-graded-category}）。
\end{lemma}

\begin{proof}
除恒等态射的存在性外，一切均由构造及引理
\ref{weil-lemma-composition-correspondences} 显然可得。给定光滑射影概形 $X$，
考虑对角线 $\Delta \subset X \times X$ 的类 $[\Delta]$，它属于
$\text{Corr}^0(X, X)$。注意，$\Delta$ 等于恒等态射
$\text{id}_X : X \to X$ 的图；下文将用到这一事实。

\medskip\noindent
为证明 $[\Delta]$ 可充当恒等态射，须对任意对应
$c \in \text{Corr}^r(Y, X)$ 和 $c' \in \text{Corr}^s(X, Y)$ 证明
$[\Delta] \circ c = c$ 以及 $c' \circ [\Delta] = c'$。对于后一等式，须证明
$$
c' = \text{pr}_{13, *}(\text{pr}_{12}^*[\Delta] \cdot \text{pr}_{23}^*c')
$$
其中 $\text{pr}_{12} : X \times X \times Y \to X \times X$ 是投影，
$\text{pr}_{13}$、$\text{pr}_{23}$ 类似。可将
$c' = \sum a_i [Z_i]$，其中 $Z_i \subset X \times Y$ 是整闭子概形，
$a_i$ 是有理数。因此，只需对 $X \times Y$ 的任意整闭子概形 $Z$，在
$X \times Y$ 的 Chow 群中证明
$$
[Z] = \text{pr}_{13, *}(\text{pr}_{12}^*[\Delta] \cdot \text{pr}_{23}^*[Z])
$$
将 $X$、$Y$ 分别替换为包含 $Z$ 在两个投影下之像的不可约分支后，可再假设
$X$、$Y$ 都是整概形。此时须证明
$$
[Z] = \text{pr}_{13, *}([\Delta \times Y] \cdot [X \times Z])
$$
记 $Z' \subset X \times X \times Y$ 为 $Z$ 在态射
$(\Delta, 1) : X \times Y \to X \times X \times Y$ 下的像。则 $Z'$ 是
$X \times X \times Y$ 的闭子概形，与 $Z$ 同构，并且按概形论意义有
$Z' = \Delta \times Y \cap X \times Z$。由《Chow 同调》引理
\ref{chow-lemma-intersect-properly}\footnote{读者可验证
$\dim(Z') = \dim(\Delta \times Y) + \dim(X \times Z) -
\dim(X \times X \times Y)$，并且 $Z'$ 有唯一泛点，它映到 $Z$ 的泛点
（该处局部环为 CM）以及 $X$ 的某一点（该处局部环为 CM）。因此，该引理的
全部假设确已满足。}
可得
$$
[Z'] = [\Delta \times Y] \cdot [X \times Z]
$$
又因 $Z'$ 沿 $\text{pr}_{13}$ 同构地映到 $Z$，结论随即成立。
$[\Delta] \circ c = c$ 的验证相同，故略去。
\end{proof}

\begin{lemma}
\label{weil-lemma-contravariant-functor}
从 $k$ 上光滑射影概形范畴到对应范畴存在一个反变函子；它在对象上恒等，
并把 $f : Y \to X$ 送到元素 $[\Gamma_f] \in \text{Corr}^0(X, Y)$。
\end{lemma}

\begin{proof}
在引理 \ref{weil-lemma-category-correspondences} 的证明中已经看到，该构造把恒等态射
送到恒等态射。为完成证明，还须说明：若 $g : Z \to Y$ 是 $k$ 上光滑射影概形
之间的另一态射，则在 $\text{Corr}^0(X, Z)$ 中有
$[\Gamma_g] \circ [\Gamma_f] = [\Gamma_{f \circ g}]$。仿照引理
\ref{weil-lemma-category-correspondences} 的证明可知，当 $X$、$Y$、$Z$ 都是整概形时，
只需在 $\CH^*(X \times Z)$ 中证明
$$
[\Gamma_{f \circ g}] =
\text{pr}_{13, *}([\Gamma_f \times Z] \cdot [X \times \Gamma_g])
$$
记 $Z' \subset X \times Y \times Z$ 为闭浸入
$(f \circ g, g, 1) : Z \to X \times Y \times Z$ 的像。按概形论意义有
$Z' = \Gamma_f \times Z \cap X \times \Gamma_g$，故由《Chow 同调》引理
\ref{chow-lemma-intersect-properly} 得
$$
[Z'] = [\Gamma_f \times Z] \cdot [X \times \Gamma_g]
$$
而 $\text{pr}_{13, *}([Z']) = [\Gamma_{f \circ g}]$ 显然成立，证明完成。
\end{proof}

\begin{remark}
\label{weil-remark-transpose}
设 $X$、$Y$ 是 $k$ 上光滑射影概形，并假设 $X$ 是维数为 $d$ 的等维概形，
$Y$ 是维数为 $e$ 的等维概形。则交换因子的同构
$X \times Y \to Y \times X$ 决定同构
$$
\text{Corr}^r(X, Y) \longrightarrow \text{Corr}^{d - e + r}(Y, X),\quad
c \longmapsto c^t
$$
称为{\it 转置}。它既作用于周期，也作用于周期类。有时有用的一个例子是态射
$f : Y \to X$ 之图的转置 $[\Gamma_f]^t = [\Gamma_f^t]$。
\end{remark}

\begin{lemma}
\label{weil-lemma-functor-and-cycles}
设 $f : Y \to X$ 是 $k$ 上光滑射影概形之间的态射，并令
$[\Gamma_f] \in \text{Corr}^0(X, Y)$ 如例
\ref{weil-example-graph-correspondence}。则
\begin{enumerate}
\item 沿对应 $[\Gamma_f]$ 的周期推前与 Gysin 映射
$f^! : \CH^*(X) \to \CH^*(Y)$ 一致；
\item 沿对应 $[\Gamma_f]$ 的周期拉回与推前映射
$f_* : \CH_*(Y) \to \CH_*(X)$ 一致；
\item 若 $X$、$Y$ 分别是维数为 $d$、$e$ 的等维概形，则
\begin{enumerate}
\item 沿注 \ref{weil-remark-transpose} 中对应 $[\Gamma_f^t]$ 的周期推前，
对应于沿 $f$ 的周期推前；
\item 沿注 \ref{weil-remark-transpose} 中对应 $[\Gamma_f^t]$ 的周期拉回，
对应于 Gysin 映射 $f^!$。
\end{enumerate}
\end{enumerate}
\end{lemma}

\begin{proof}
先证 (1)。回忆
$[\Gamma_f]_*(\alpha) =
\text{pr}_{2, *}([\Gamma_f] \cdot \text{pr}_1^*\alpha)$。
有
$$
[\Gamma_f] \cdot \text{pr}_1^*\alpha =
(f, 1)_*((f, 1)^! \text{pr}_1^*\alpha) =
(f, 1)_*((f, 1)^! \text{pr}_1^!\alpha) =
(f, 1)_*(f^!\alpha)
$$
第一个等式由《Chow 同调》引理
\ref{chow-lemma-intersect-regularly-embedded} 得到。第二个等式由《Chow 同调》
引理 \ref{chow-lemma-lci-gysin-flat} 得到。第三个等式来自
$\text{pr}_1 \circ (f, 1) = f$ 以及《Chow 同调》引理
\ref{chow-lemma-lci-gysin-composition}。最后，由
$\text{pr}_{2, *} \circ (f, 1)_* = 1_*$ 及《Chow 同调》引理
\ref{chow-lemma-compose-pushforward} 即得结论。

\medskip\noindent
再证 (2)。回忆
$[\Gamma_f]_*(\beta) =
\text{pr}_{1, *}([\Gamma_f] \cdot \text{pr}_2^*\beta)$。
完全仿照上面的论证可得
$$
[\Gamma_f] \cdot \text{pr}_2^*\beta = (f, 1)_*\beta
$$
故结论如前成立。

\medskip\noindent
(3) 也以完全相同的方式证明。
\end{proof}

\begin{example}
\label{weil-example-decompose-P1}
设 $X = \mathbf{P}^1_k$。则
$$
\text{Corr}^0(X, X) = \CH^1(X \times X) \otimes \mathbf{Q} =
\CH_1(X \times X) \otimes \mathbf{Q}
$$
选取一个 $k$-有理点 $x \in X$，并考虑周期
$c_0 = [x \times X]$ 和 $c_2 = [X \times x]$。计算表明，在
$\text{Corr}^0(X, X)$ 中有 $1 = [\Delta] = c_0 + c_2$，而复合法则为
$c_0 \circ c_0 = c_0$、
$c_0 \circ c_2 = 0$、
$c_2 \circ c_0 = 0$ 以及
$c_2 \circ c_2 = c_2$。
换言之，$c_0$、$c_2$ 是代数 $\text{Corr}^0(X, X)$ 中相互正交的幂等元；
事实上，作为 $\mathbf{Q}$-代数有
$$
\text{Corr}^0(X, X) = \mathbf{Q} \times \mathbf{Q}
$$
\end{example}

\noindent
对应范畴是一个对称幺半范畴。给定 $k$ 上光滑射影概形 $X$、$Y$，定义
$X \otimes Y = X \times Y$。给定 $k$ 上四个光滑射影概形
$X, X', Y, Y'$，定义张量积
$$
\otimes :
\text{Corr}^r(X, Y) \times \text{Corr}^{r'}(X', Y')
\longrightarrow
\text{Corr}^{r + r'}(X \times X', Y \times Y')
$$
如下：
$$
(c, c') \longmapsto
c \otimes c' = \text{pr}_{13}^*c \cdot \text{pr}_{24}^*c'
$$
其中 $\text{pr}_{13} : X \times X' \times Y \times Y' \to X \times Y$
和 $\text{pr}_{24} : X \times X' \times Y \times Y' \to X' \times Y'$
是投影。对于结合约束
$$
X \otimes (Y \otimes Z) = (X \otimes Y) \otimes Z
$$
采用概形之积的通常结合约束；交换约束则由交换因子的同构
$X \times Y \to Y \times X$ 给出。

\begin{lemma}
\label{weil-lemma-tensor-product}
上述对应张量积使对应范畴成为以 $\Spec(k)$ 为单位对象的对称幺半范畴。
\end{lemma}

\begin{proof}
略。
\end{proof}

\begin{lemma}
\label{weil-lemma-prep-dual}
设 $f : Y \to X$ 是 $k$ 上光滑射影概形之间的态射。假设 $X$、$Y$ 分别是
维数为 $d$、$e$ 的等维概形。记
$a = [\Gamma_f] \in \text{Corr}^0(X, Y)$ 及
$a^t = [\Gamma_f^t] \in \text{Corr}^{d - e}(Y, X)$。置
$\eta_X = [\Gamma_{X \to X \times X}] \in \text{Corr}^0(X \times X, X)$、
$\eta_Y = [\Gamma_{Y \to Y \times Y}] \in \text{Corr}^0(Y \times Y, Y)$、
$[X] \in \text{Corr}^{-d}(X, \Spec(k))$ 以及
$[Y] \in \text{Corr}^{-e}(Y, \Spec(k))$。下图在对应范畴中交换：
$$
\xymatrix{
X \otimes Y \ar[r]_{a \otimes \text{id}} \ar[d]_{\text{id} \otimes a^t} &
Y \otimes Y \ar[r]_{\eta_Y} &
Y \ar[d]^{[Y]} \\
X \otimes X \ar[r]^{\eta_X} &
X \ar[r]^{[X]} &
\Spec(k)
}
$$
\end{lemma}

\begin{proof}
回忆，对任意 $k$ 上光滑射影概形 $W$，有
$\text{Corr}^r(W, \Spec(k)) = \CH_{-r}(W)$；给定
$c \in \text{Corr}^s(W', W)$，与 $c$ 复合等同于沿 $c$ 的拉回映射
$\CH_{-r}(W) \to \CH_{-r - s}(W')$
（引理 \ref{weil-lemma-composition-correspondences}）。最后，引理
\ref{weil-lemma-functor-and-cycles} 说明如何将此转化为通常的周期推前与拉回。
有
$$
(a \otimes \text{id})^* \eta_Y^* [Y] =
(a \otimes \text{id})^* [\Delta_Y] =
(f \times \text{id})_*\Delta_Y = [\Gamma_f]
$$
沿图的另一条路径则得到
$$
(\text{id} \otimes a^t)^* \eta_X^* [X] =
(\text{id} \otimes a^t)^* [\Delta_X] =
(\text{id} \times f)^![\Delta_X] = [\Gamma_f]
$$
最后一个等式由《Chow 同调》引理
\ref{chow-lemma-lci-gysin-easy} 得到。换言之，沿图的任一路径都得到
$\text{Corr}^d(X \times Y, \Spec(k))$ 中与周期
$\Gamma_f \subset X \times Y$ 对应的元素。
\end{proof}







\section{Chow 动机}
\label{weil-section-chow-motives}

\noindent
固定一个基域 $k$。本节构造一个带对称幺半结构的加性 Karoubian
$\mathbf{Q}$-线性范畴 $M_k$，以及一个反变函子
$$
h : \{\text{smooth projective schemes over }k\} \longrightarrow M_k
$$
它把积映到张量积，把不交并映到直和。这个构造由下述性质刻画：$h$ 通过这样一个
对称幺半范畴分解，其对象是光滑射影簇，态射是 $0$ 次对应，并且例
\ref{weil-example-decompose-P1} 中 $h(\mathbf{P}^1_k)$ 上投影算子 $c_2$ 的像在
$M_k$ 中可逆；见引理 \ref{weil-lemma-characterize-motives}。本节末将证明每个动机，
即 $M_k$ 的每个对象，都有一个（左）对偶；见引理 \ref{weil-lemma-dual-general}。

\medskip\noindent
$k$ 上的一个{\it 动机}或{\it Chow 动机}是三元组 $(X, p, m)$，其中
\begin{enumerate}
\item $X$ 是 $k$ 上光滑射影概形；
\item $p \in \text{Corr}^0(X, X)$ 满足 $p \circ p = p$；
\item $m \in \mathbf{Z}$。
\end{enumerate}
给定另一个动机 $(Y, q, n)$，定义一个{\it 动机态射}或{\it Chow 动机态射}
为下式中的一个元素：
$$
\Hom((X, p, m), (Y, q, n)) =
q \circ \text{Corr}^{n - m}(X, Y) \circ p \subset \text{Corr}^{n - m}(X, Y)
$$
动机态射的复合用上面定义的对应复合来定义。

\begin{lemma}
\label{weil-lemma-motives}
以 $k$ 上动机为对象、以 $k$ 上动机态射为态射的范畴 $M_k$ 是
$\mathbf{Q}$-线性范畴。存在反变函子
$$
h : \{\text{smooth projective schemes over }k\} \longrightarrow M_k
$$
其定义为 $h(X) = (X, 1, 0)$ 和 $h(f) = [\Gamma_f]$。
\end{lemma}

\begin{proof}
由引理 \ref{weil-lemma-contravariant-functor} 立即得到。
\end{proof}

\begin{lemma}
\label{weil-lemma-Karoubian}
范畴 $M_k$ 是 Karoubian 范畴。
\end{lemma}

\begin{proof}
设 $M = (X, p, m)$ 是动机，$a \in \Mor(M, M)$ 是投影算子。则无论在
$\Mor(M, M)$ 还是在 $\text{Corr}^0(X, X)$ 中，均有
$a = a \circ a$。置 $N = (X, a, m)$。由于在
$\text{Corr}^0(X, X)$ 中有 $a = p \circ a \circ a$，可知
$a : N \to M$ 是 $M_k$ 中的态射。再设
$b : (Y, q, n) \to M$ 是满足 $(1 - a) \circ b = 0$ 的态射。则既有
$b = a \circ b$，又有 $b = b \circ q$，故 $b$ 是态射
$b : (Y, q, n) \to N$。因此，投影算子 $1 - a$ 有核 $N$，从而
$M_k$ 是 Karoubian 范畴；见《同调代数》定义
\ref{homology-definition-karoubian}。
\end{proof}

\noindent
定义函子
$$
\otimes : M_k \times M_k \longrightarrow M_k
$$
在对象上采用公式
$$
(X, p, m) \otimes (Y, q, n) = (X \times Y, p \otimes q, m + n)
$$
在态射上，采用映射
$$
\xymatrix{
\Mor((X, p, m), (Y, q, n)) \times
\Mor((X', p', m'), (Y', q', n')) \ar[d] \\
\Mor(
(X \times X', p \otimes p', m + m'),
(Y \times Y', q \otimes q', n + n'))
}
$$
其法则为 $(a, a') \longmapsto a \otimes a'$，其中对应上的
$\otimes$ 如第 \ref{weil-section-correspondences} 节所定义。这一定义有意义：
按动机态射的定义，可写
$a = q \circ c \circ p$、$a' = q' \circ c' \circ p'$，其中
$c \in \text{Corr}^{n - m}(X, Y)$ 且
$c' \in \text{Corr}^{n' - m'}(X', Y')$；于是
$$
a \otimes a' =
(q \circ c \circ p) \otimes (q' \circ c' \circ p') =
(q \otimes q') \circ (c \otimes c') \circ (p \otimes p')
$$
它确为从 $(X \times X', p \otimes p', m + m')$ 到
$(Y \times Y', q \otimes q', n + n')$ 的动机态射。

\begin{lemma}
\label{weil-lemma-motives-monoidal}
赋予上述张量积后，范畴 $M_k$ 连同显然的结合约束、交换约束及单位对象
$\mathbf{1} = (\Spec(k), 1, 0)$ 成为对称幺半范畴。
\end{lemma}

\begin{proof}
由引理 \ref{weil-lemma-tensor-product} 立即得到。细节略。
\end{proof}

\noindent
动机 $\mathbf{1}(n) = (\Spec(k), 1, n)$ 很有用。注意
$$
\mathbf{1} = \mathbf{1}(0)
\quad\text{且}\quad
\mathbf{1}(n + m) = \mathbf{1}(n) \otimes \mathbf{1}(m)
$$
因此，与 $\mathbf{1}(1)$ 作张量积给出动机范畴的一个自等价。给定动机 $M$，
有时记 $M(n) = M \otimes \mathbf{1}(n)$。注意，若 $M = (X, p, m)$，则
$M(n) = (X, p, m + n)$。

\begin{lemma}
\label{weil-lemma-inverse-h2}
采用例 \ref{weil-example-decompose-P1} 的记号，有
\begin{enumerate}
\item 动机 $(X, c_0, 0)$ 与动机
$\mathbf{1} = (\Spec(k), 1, 0)$ 同构；
\item 动机 $(X, c_2, 0)$ 与动机
$\mathbf{1}(-1) = (\Spec(k), 1, -1)$ 同构。
\end{enumerate}
\end{lemma}

\begin{proof}
以下不再特别说明地使用引理 \ref{weil-lemma-contravariant-functor}。结构态射
$X \to \Spec(k)$ 给出对应 $a \in \text{Corr}^0(\Spec(k), X)$。另一方面，
有理点 $x$ 是态射 $\Spec(k) \to X$，它给出对应
$b \in \text{Corr}^0(X, \Spec(k))$。作为 $\Spec(k)$ 上对应，有
$b \circ a = 1$。复合 $a \circ b$ 对应于复合 $X \to x \to X$ 的图，即
$c_0 = [x \times X]$。因此
$a = a \circ b \circ a = c_0 \circ a$ 且
$b = b \circ a \circ b = b \circ c_0$。展开定义可见，$a$、$b$ 是互逆态射
$a : (\Spec(k), 1, 0) \to (X, c_0, 0)$ 和
$b : (X, c_0, 0) \to (\Spec(k), 1, 0)$。

\medskip\noindent
完全仿照上述论证证明第二个断言。记
$$
a' \in \text{Corr}^1(\Spec(k), X) = \CH^1(X)
$$
为点 $x$ 的类，并记
$$
b' \in \text{Corr}^{-1}(X, \Spec(k)) = \CH_1(X)
$$
为 $[X]$ 的类。由于在
$X = \Spec(k) \times X \times \Spec(k)$ 上有
$[x] \cdot [X] = [x]$，故作为 $\Spec(k)$ 上对应有
$b' \circ a' = 1$。在 $X \times \Spec(k) \times X$ 上计算交积
$\text{pr}_{12}^*b' \cdot \text{pr}_{23}^*a'$，所得周期为
$X \times \Spec(k) \times x$。因此，作为 $X$ 上对应，复合
$a' \circ b'$ 等于 $c_2$。于是
$a' = a' \circ b \circ a' = c_2 \circ a'$ 且
$b' = b' \circ a' \circ b' = b' \circ c_2$。回忆
$$
\Mor((\Spec(k), 1, -1), (X, c_2, 0)) =
c_2 \circ \text{Corr}^1(\Spec(k), X)
\subset
\text{Corr}^1(\Spec(k), X)
$$
以及
$$
\Mor((X, c_2, 0), (\Spec(k), 1, -1)) =
\text{Corr}^{-1}(X, \Spec(k)) \circ c_2
\subset
\text{Corr}^{-1}(X, \Spec(k))
$$
故 $a'$、$b'$ 是互逆态射
$a' : (\Spec(k), 1, -1) \to (X, c_2, 0)$ 和
$b' : (X, c_2, 0) \to (\Spec(k), 1, -1)$。
\end{proof}

\begin{remark}[Lefschetz 动机与 Tate 动机]
\label{weil-remark-lefschetz-tate}
设 $X = \mathbf{P}^1_k$，并令 $c_2$ 如例
\ref{weil-example-decompose-P1}。文献中有时把动机 $(X, c_2, 0)$ 称为
{\it Lefschetz 动机}；依参考文献不同，也可能用
$L$、$\mathbf{L}$、$\mathbf{Q}(-1)$ 或 $h^2(\mathbf{P}^1_k)$ 表示它。
由引理 \ref{weil-lemma-inverse-h2}，Lefschetz 动机同构于
$\mathbf{1}(-1)$。因此 Lefschetz 动机可逆（《范畴》定义
\ref{categories-definition-invertible}），其逆为
$\mathbf{1}(1)$。动机 $\mathbf{1}(1)$ 有时称为{\it Tate 动机}；依参考文献不同，
也可能用 $L^{-1}$、$\mathbf{L}^{-1}$、$\mathbf{T}$ 或
$\mathbf{Q}(1)$ 表示它。
\end{remark}

\begin{lemma}
\label{weil-lemma-additive}
范畴 $M_k$ 是加性范畴。
\end{lemma}

\begin{proof}
设 $(Y, p, m)$、$(Z, q, n)$ 是动机。若 $n = m$，则在显然的记号下，
直和由 $(Y \amalg Z, p + q, m)$ 给出。细节略。

\medskip\noindent
设 $n < m$。令 $X$、$c_2$ 如例 \ref{weil-example-decompose-P1}。考虑
\begin{align*}
(Z, q, n)
& =
(Z, q, m) \otimes (\Spec(k), 1, -1) \otimes \ldots \otimes
(\Spec(k), 1, -1) \\
& \cong
(Z, q, m) \otimes (X, c_2, 0) \otimes \ldots \otimes (X, c_2, 0) \\
& \cong
(Z \times X^{m - n}, q \otimes c_2 \otimes \ldots \otimes c_2, m)
\end{align*}
这里使用了引理 \ref{weil-lemma-inverse-h2}，从而归约到第一段所讨论的情形。
\end{proof}

\begin{lemma}
\label{weil-lemma-decompose-P1}
在 $M_k$ 中有 $h(\mathbf{P}^1_k) \cong \mathbf{1} \oplus \mathbf{1}(-1)$。
\end{lemma}

\begin{proof}
由例 \ref{weil-example-decompose-P1} 和引理 \ref{weil-lemma-inverse-h2} 得到。
\end{proof}

\begin{lemma}
\label{weil-lemma-characterize-motives}
令 $X$、$c_2$ 如例 \ref{weil-example-decompose-P1}，并设
$\mathcal{C}$ 是一个 $\mathbf{Q}$-线性 Karoubian 对称幺半范畴。若对称幺半范畴
之间的任意 $\mathbf{Q}$-线性函子
$$
F :
\left\{
\begin{matrix}
\text{smooth projective schemes over }k\\
\text{morphisms are correspondences of degree }0
\end{matrix}
\right\}
\longrightarrow
\mathcal{C}
$$
满足 $F(c_2)$ 在 $F(X)$ 上的像是可逆对象，则它唯一地通过对称幺半范畴之间的函子
$F : M_k \to \mathcal{C}$ 分解。
\end{lemma}

\begin{proof}
以 $U$ 表示引理陈述中假定存在的 $\mathcal{C}$ 中可逆对象。规定
$$
F(X, p, m) = \left(\text{the image of
the projector }F(p)\text{ 于 }F(X)\right) \otimes U^{\otimes -m}
$$
从而把 $F$ 延拓到动机。由于 $U$ 可逆且 $\mathcal{C}$ 是 Karoubian 范畴，
这一定义有意义。这个选择的一个重要性质是 $F(X, c_2, 0) = U$。注意
\begin{align*}
F((X, p, m) \otimes (Y, q, n))
& =
F(X \times Y, p \otimes q, m + n) \\
& =
\left(\text{the image of }F(p \otimes q)\text{ 于 }F(X \times Y)\right)
\otimes U^{\otimes -m - n} \\
& =
F(X, p, m) \otimes F(Y, q, n)
\end{align*}
因此，该法则在对象层面与张量积相容（细节略）。

\medskip\noindent
下面把 $F$ 延拓到动机态射。设
$$
a \in
\Hom((Y, p, m), (Z, q, n)) =
q \circ \text{Corr}^{n - m}(Y, Z) \circ p \subset \text{Corr}^{n - m}(Y, Z)
$$
是一个态射。若 $n = m$，则 $a$ 是 $0$ 次对应，可以使用
$F(a) : F(Y) \to F(Z)$ 得到所需映射
$F(Y, p, m) \to F(Z, q, n)$。若 $n < m$，则得到典范等同
\begin{align*}
s : F((Z, q, n))
& \to
F(Z, q, m) \otimes U^{m - n} \\
& \to
F(Z, q, m) \otimes F(X, c_2, 0) \otimes \ldots \otimes F(X, c_2, 0) \\
& \to
F((Z, q, m) \otimes (X, c_2, 0) \otimes \ldots \otimes (X, c_2, 0)) \\
& \to
F((Z \times X^{m - n}, q \otimes c_2 \otimes \ldots \otimes c_2, m))
\end{align*}
具体而言，第一个同构使用上面对动机上 $F$ 的定义；第二个使用 $U$ 的选择；
第三个使用 $F$ 与动机张量积的相容性；第四个是动机张量积的定义。另一方面，
类似地有同构
$$
\sigma : (Z, q, n) \to
(Z \times X^{m - n}, q \otimes c_2 \otimes \ldots \otimes c_2, m)
$$
（见引理 \ref{weil-lemma-additive} 的证明）。将 $a$ 与此同构复合得到
$$
\sigma \circ a \in
\Hom((Y, p, m),
(Z \times X^{m - n}, q \otimes c_2 \otimes \ldots \otimes c_2, m))
$$
综合以上构造，得到
$$
s^{-1} \circ F(\sigma \circ a) :
F(Y, p, m) \to
F(Z, q, n)
$$
若 $n > m$，类似地定义同构
$$
t : F((Y, p, m)) \to
F((Y \times X^{n - m}, p \otimes c_2 \otimes \ldots \otimes c_2, n))
$$
以及
$$
\tau : (Y, p, m)) \to
(Y \times X^{n - m}, p \otimes c_2 \otimes \ldots \otimes c_2, n)
$$
并置 $F(a) = F(a \circ \tau^{-1}) \circ t$。略去验证该构造定义了对称幺半范畴
之间函子的过程。
\end{proof}

\begin{lemma}
\label{weil-lemma-dual}
设 $X$ 是 $k$ 上维数为 $d$ 的等维光滑射影概形。则在 $M_k$ 中，
$h(X)(d)$ 是 $h(X)$ 的左对偶。
\end{lemma}

\begin{proof}
以下不再特别说明地使用引理 \ref{weil-lemma-composition-correspondences}。计算得
$$
\Hom(\mathbf{1}, h(X) \otimes h(X)(d)) =
\text{Corr}^d(\Spec(k), X \times X) = \CH^d(X \times X)
$$
此处取 $\eta = [\Delta]$。另一方面，
$$
\Hom(h(X)(d) \otimes h(X), \mathbf{1}) =
\text{Corr}^{-d}(X \times X, \Spec(k)) = \CH_d(X \times X)
$$
此处同样取对角线的类 $\epsilon = [\Delta]$。对应
$[\Delta] \otimes 1$ 与 $1 \otimes [\Delta]$ 按任一顺序复合，均为
$\text{Corr}^0(X, X)$ 中的对应 $[\Delta] = 1$；这就证明《范畴》定义
\ref{categories-definition-dual} 所要求的图交换。具体地，注意
$$
[\Delta] \otimes 1 \in \text{Corr}^d(X, X \times X \times X) =
\CH^{2d}(X \times X \times X \times X)
$$
在显然的记号下由周期
$\text{pr}^{1234, -1}_{23}(\Delta) \cap \text{pr}^{1234, -1}_{14}(\Delta)$
的类给出。类似地，类
$$
1 \otimes [\Delta] \in \text{Corr}^{-d}(X \times X \times X, X) =
\CH^{2d}(X \times X \times X \times X)
$$
由周期
$\text{pr}^{1234, -1}_{23}(\Delta) \cap \text{pr}^{1234, -1}_{14}(\Delta)$
的类给出。按定义，复合
$(1 \otimes [\Delta]) \circ ([\Delta] \otimes 1)$ 是交积
$$
[\text{pr}^{12345, -1}_{23}(\Delta) \cap \text{pr}^{12345, -1}_{14}(\Delta)]
\cdot
[\text{pr}^{12345, -1}_{34}(\Delta) \cap \text{pr}^{12345, -1}_{15}(\Delta)]
=
[\text{small diagonal in } X^5]
$$
沿 $\text{pr}^{12345}_{15, *}$ 的推前，结果正是所需的 $\Delta$。
另一复合顺序的公式证明略去。
\end{proof}

\begin{lemma}
\label{weil-lemma-dual-general}
$M_k$ 的每个对象都有左对偶。
\end{lemma}

\begin{proof}
设 $M = (X, p, m)$ 是 $M_k$ 的对象，则 $M$ 是
$(X, 1, m) = h(X)(m)$ 的直和项。由《同调代数》引理
\ref{homology-lemma-Karoubian-dual}，只需证明
$h(X)(m) = h(X) \otimes \mathbf{1}(m)$ 有对偶。按构造，
$\mathbf{1}(-m)$ 是 $\mathbf{1}(m)$ 的左对偶。因此，只需证明
$h(X)$ 有左对偶；见《范畴》引理 \ref{categories-lemma-tensor-dual}。
设 $X = \coprod X_i$ 是 $X$ 的不可约分支分解。则
$h(X) = \bigoplus h(X_i)$，只需证明每个 $h(X_i)$ 有左对偶；见《同调代数》
引理 \ref{homology-lemma-additive-dual}。这由引理 \ref{weil-lemma-dual} 得到。
\end{proof}






\section{动机的 Chow 群}
\label{weil-section-chow-groups-motives}

\noindent
如下定义动机的 Chow 群。

\begin{definition}
\label{weil-definition-chow-group-motives}
设 $k$ 是基域，$M = (X, p, m)$ 是 $k$ 上 Chow 动机。对
$i \in \mathbf{Z}$，用公式
$$
\CH^i(M) = p\left(\CH^{i + m}(X) \otimes \mathbf{Q}\right)
$$
定义{\it $M$ 的第 $i$ 个 Chow 群}。
\end{definition}

\noindent
若 $X$ 是 $k$ 上光滑射影概形，则
$\CH^i(h(X)) = \CH^i(X) \otimes \mathbf{Q}$。

\medskip\noindent
注意，$\CH^i(-)$ 是从 $M_k$ 到 $\mathbf{Q}$-向量空间的函子。事实上，若
$c : M \to N$ 是动机 $M = (X, p, m)$ 与 $N = (Y, q, n)$ 之间的态射，
则 $c$ 是从 $X$ 到 $Y$ 的 $n - m$ 次对应；故沿 $c$ 的推前
（第 \ref{weil-section-correspondences} 节）是一族映射
$$
c_* :
\CH^{i + m}(X) \otimes \mathbf{Q}
\longrightarrow
\CH^{i + n}(Y) \otimes \mathbf{Q}
$$
按动机态射的定义，有 $c = q \circ c \circ p$，故对所有
$i \in \mathbf{Z}$ 确实得到
$$
c_* : \CH^i(M) \to \CH^i(N)
$$
由引理 \ref{weil-lemma-composition-correspondences}，这与动机态射的复合相容。
Chow 群的这一函子性也可由下述引理推出。

\begin{lemma}
\label{weil-lemma-chow-groups-representable}
设 $k$ 是基域。动机范畴 $M_k$ 上的函子 $\CH^i(-)$ 由
$\mathbf{1}(-i)$ 表示；换言之，在 $M_k$ 中关于 $M$ 函子性地有
$$
\CH^i(M) = \Hom_{M_k}(\mathbf{1}(-i), M)
$$
\end{lemma}

\begin{proof}
由定义和引理 \ref{weil-lemma-composition-correspondences} 立即得到。
\end{proof}

\noindent
由引理 \ref{weil-lemma-chow-groups-representable}、Yoneda 引理和引理
\ref{weil-lemma-dual} 中的对偶性，读者不难设想可得下述结果。

\begin{lemma}[Manin]
\label{weil-lemma-manin}
设 $k$ 是基域，$c : M \to N$ 是动机态射。若对每个 $k$ 上光滑射影概形 $X$，
映射 $c \otimes 1 : M \otimes h(X) \to N \otimes h(X)$ 在 Chow 群上诱导同构，
则 $c$ 是同构。
\end{lemma}

\begin{proof}
$M_k$ 的任一对象 $L$ 都是某个 $h(X)(m)$ 的直和项，其中 $X$ 是 $k$ 上光滑射影
概形，$m \in \mathbf{Z}$。注意，$M \otimes h(X)(m)$ 的 Chow 群与
$M \otimes h(X)$ 的 Chow 群仅相差次数平移。因此，假设蕴含：对 $M_k$ 的每个对象
$L$，映射 $c \otimes 1 : M \otimes L \to N \otimes L$ 在 Chow 群上诱导同构。
由引理 \ref{weil-lemma-chow-groups-representable} 可知，对每个 $L$，
$$
\Hom_{M_k}(\mathbf{1}, M \otimes L) \to
\Hom_{M_k}(\mathbf{1}, N \otimes L)
$$
是同构。由于 $M_k$ 的每个对象都有左对偶（引理
\ref{weil-lemma-dual-general}），故对 $M_k$ 的每个对象 $K$，
$$
\Hom_{M_k}(K, M) \to \Hom_{M_k}(K, N)
$$
是同构；见《范畴》引理 \ref{categories-lemma-left-dual}。最后由 Yoneda 引理
（《范畴》引理 \ref{categories-lemma-yoneda}）得出结论。
\end{proof}




\section{射影空间丛公式}
\label{weil-section-projective-space-bundle}

\noindent
设 $k$ 是基域，$X$ 是 $k$ 上光滑射影概形，$\mathcal{E}$ 是秩为 $r$ 的局部自由
$\mathcal{O}_X$-模。约定与 $\mathcal{E}$ {\it 相伴的射影丛}是 $X$ 上态射
$$
\xymatrix{
P = \mathbf{P}(\mathcal{E}) =
\underline{\text{Proj}}_X(\text{Sym}^*(\mathcal{E}))
\ar[r]^-p
& X
}
$$
并将 $\mathcal{O}_P(1)$ 规范化为满足
$p_*(\mathcal{O}_P(1)) = \mathcal{E}$。回忆
$$
[\Gamma_p] \in \text{Corr}^0(X, P) \subset \CH^*(X \times P) \otimes \mathbf{Q}
$$
参见例 \ref{weil-example-graph-correspondence}。对
$i = 0, \ldots, r - 1$，考虑对应
$$
c_i = c_1(\text{pr}_2^*\mathcal{O}_P(1))^i \cap [\Gamma_p]
\in \text{Corr}^i(X, P)
$$
可将且确将 $c_i$ 视为态射 $h(X)(-i) \to h(P)$。

\begin{lemma}[射影空间丛公式]
\label{weil-lemma-projective-space-bundle-formula}
在上述情形中，映射
$$
\sum\nolimits_{i = 0, \ldots, r - 1} c_i :
\bigoplus\nolimits_{i = 0, \ldots, r - 1} h(X)(-i)
\longrightarrow
h(P)
$$
是动机范畴中的同构。
\end{lemma}

\begin{proof}
由引理 \ref{weil-lemma-manin}，只需证明：与任意光滑射影概形 $Z$ 取积后，该映射在
动机的 Chow 群上给出同构。注意，$P \times Z \to X \times Z$ 是与
$\mathcal{E}$ 到 $X \times Z$ 的拉回相伴的射影丛。因此，Chow 群上的断言由
《Chow 同调》引理 \ref{chow-lemma-chow-ring-projective-bundle} 中的射影空间丛
公式成立。具体而言，由引理 \ref{weil-lemma-functor-and-cycles} 和《Chow 同调》引理
\ref{chow-lemma-lci-gysin-flat}，沿 $[\Gamma_p]$ 的周期推前由沿 $p$ 的周期拉回
给出。因此，沿 $c_i$ 的推前把 $\alpha$ 送到
$c_1(\mathcal{O}_P(1))^i \cap p^*\alpha$。略去若干细节。
\end{proof}

\noindent
在上述情形中，对 $j = 0, \ldots, r - 1$ 考虑对应
$$
c'_j = c_1(\text{pr}_1^*\mathcal{O}_P(1))^{r - 1 - j} \cap [\Gamma_p^t] \in
\text{Corr}^{-j}(P, X)
$$
对 $i, j \in \{0, \ldots, r - 1\}$，有
$$
c'_j \circ c_i =
\text{pr}_{13, *}\left(
c_1(\text{pr}_2^*\mathcal{O}_P(1))^{i + r - 1 - j} \cap
(\text{pr}_{12}^*[\Gamma_p] \cdot \text{pr}_{23}^*[\Gamma_p^t])
\right)
$$
周期 $\text{pr}_{12}^{-1}\Gamma_p$ 与
$\text{pr}_{23}^{-1}\Gamma_p^t$ 横截相交，其交为
$(p, 1, p) : P \to X \times P \times X$ 的像。注意，
$(p, p) = \text{pr}_{13} \circ (p, 1, p) : P \to X \times X$ 的纤维维数为
$r - 1$。由此立即可得，当 $i + r - 1 - j < r - 1$，亦即 $i < j$ 时，
$c'_j \circ c_i = 0$。另一方面，由射影空间丛公式
（《Chow 同调》引理 \ref{chow-lemma-chow-ring-projective-bundle}），周期
$c_1(\mathcal{O}_P(1))^{r - 1} \cap [P]$ 在 $X$ 中映到 $[X]$。因此，当
$i = j$ 时，上述推前给出对角线的类，从而对所有
$i \in \{0, \ldots, r - 1\}$ 都有
$$
c'_i \circ c_i = 1 \in \text{Corr}^0(X, X)
$$
于是复合
$$
\bigoplus h(X)(-i)
\xrightarrow{\bigoplus c_i}
h(P)
\xrightarrow{\bigoplus c'_j}
\bigoplus h(X)(-j)
$$
的矩阵可逆（它是对角元均为 $1$ 的上三角矩阵）。由射影空间丛公式
（引理 \ref{weil-lemma-projective-space-bundle-formula}）可知，反向复合也可逆，
但直接证明这一点似乎稍难。

\begin{lemma}
\label{weil-lemma-diagonal-projective-bundle}
令 $p : P \to X$ 如引理 \ref{weil-lemma-projective-space-bundle-formula}。
$P$ 的对角线在 $\CH^*(P \times P)$ 中的类 $[\Delta_P]$ 可写为
$$
[\Delta_P] =
\left(\sum\nolimits_{i = 0, \ldots, r - 1}
{r - 1 \choose i} c_{r - 1 - i}(\text{pr}_1^*\mathcal{S}^\vee) \cap
c_1(\text{pr}_2^*\mathcal{O}_P(1))^i\right)
\cap
(p \times p)^*[\Delta_X]
$$
其中 $\mathcal{S}$ 是典范满射
$p^*\mathcal{E} \to \mathcal{O}_P(1)$ 的核。
\end{lemma}

\begin{proof}
注意 $(p \times p)^*[\Delta_X] = [P \times_X P]$。由于
$\Delta_P \subset P \times_X P \subset P \times P$，且与 Chern 类作帽积同固有推前
交换（《Chow 同调》引理 \ref{chow-lemma-pushforward-cap-cj}），只需证明
$\Delta_P \subset P \times_X P$ 在 $\CH^*(P \times_X P)$ 中的类等于
$$
\left(\sum\nolimits_{i = 0, \ldots, r - 1}
{r - 1 \choose i} c_{r - 1 - i}(q_1^*\mathcal{S}^\vee) \cap
c_1(q_2^*\mathcal{O}_P(1))^i\right)
\cap
[P \times_X P]
$$
其中 $q_i : P \times_X P \to P$（$i = 1, 2$）是投影。置
$q = p \circ q_1 = p \circ q_2 : P \times_X P \to X$。考虑映射
$$
q_1^*\mathcal{S} \otimes q_2^*\mathcal{O}_P(-1) \to
q^*\mathcal{E} \otimes q^*\mathcal{E}^\vee \to
\mathcal{O}_{P \times_X P}
$$
其中最后一个箭头是评价映射
$\mathcal{E} \otimes_{\mathcal{O}_X} \mathcal{E}^\vee \to \mathcal{O}_X$
沿 $q$ 的拉回。该复合的源是秩为 $r - 1$ 的局部自由模；局部计算表明，此映射
恰在 $\Delta_P$ 上消失。由《Chow 同调》引理
\ref{chow-lemma-top-chern-class}，类 $[\Delta_P]$ 是对偶
$$
q_1^*\mathcal{S}^\vee \otimes q_2^*\mathcal{O}_P(1)
$$
的最高 Chern 类。所求结论由《Chow 同调》引理
\ref{chow-lemma-chern-classes-E-tensor-L} 得到。
\end{proof}








\section{经典 Weil 上同调理论}
\label{weil-section-axioms-classical}

\noindent
本节定义所谓经典 Weil 上同调理论。这恰是
\cite[Section 1.2]{Kleiman-cycles} 中所谓的 Weil 上同调理论。

\medskip\noindent
固定一个代数闭域 $k$（基域）。本节所谓{\it 簇}均指 $k$ 上簇；见《簇》第
\ref{varieties-section-varieties} 节。再固定一个特征为 $0$ 的域 $F$（系数域）。
Weil 上同调理论由数据 (D1)、(D2)、(D3) 给出，并须满足公理 (A)、(B)、(C)。

\medskip\noindent
数据如下：
\begin{enumerate}
\item[(D1)] 从光滑射影簇范畴到分次交换 $F$-代数范畴的反变函子 $H^*$。
\item[(D2)] 对每个光滑射影簇 $X$，一个群同态
$\gamma : \CH^i(X) \to H^{2i}(X)$。
\item[(D3)] 对每个维数为 $d$ 的光滑射影簇 $X$，一个映射
$\int_X : H^{2d}(X) \to F$。
\end{enumerate}
下面作若干说明，以解释这些数据的含义并引入有关术语。

\medskip\noindent
关于 (D1)。给定光滑射影簇 $X$，称 $H^*(X)$ 为 $X$ 的{\it 上同调}。给定光滑
射影簇之间的态射 $f : X \to Y$，以
$f^* : H^*(Y) \to H^*(X)$ 表示映射 $H^*(f)$，并称其为{\it 拉回映射}。

\medskip\noindent
关于 (D2)。映射 $\gamma$ 称为{\it 周期类映射}，并称
$\gamma(\alpha)$ 为 $\alpha$ 的{\it 上同调类}。若
$Z \subset Y \subset X$ 是闭子概形，其中 $Y$、$X$ 是光滑射影簇而 $Z$ 是整概形，
则 $[Z]$ 既可能表示 $\CH^*(Y)$ 中周期 $[Z]$ 的类，也可能表示其在
$\CH^*(X)$ 中的类。在这种情形下，记号 $\gamma([Z])$ 有歧义，须由上下文判断
其含义。

\medskip\noindent
关于 (D3)。映射 $\int_X$ 有时称为{\it 迹映射}，也有时记作
$\text{Tr}_X$。

\medskip\noindent
第一个公理通常称为{\it Poincar\'e 对偶}：
\begin{enumerate}
\item[(A)] 设 $X$ 是维数为 $d$ 的光滑射影簇。则
\begin{enumerate}
\item 对所有 $i$，有 $\dim_F H^i(X) < \infty$；
\item 对所有 $i$，
$H^i(X) \times H^{2d - i}(X) \rightarrow H^{2d}(X) \rightarrow F$
是完美配对，其中最后的映射是迹映射 $\int_X$；
\item 除非 $i \in [0, 2d]$，否则 $H^i(X) = 0$；
\item $\int_X : H^{2d}(X) \to F$ 是同构。
\end{enumerate}
\end{enumerate}
设 $f : X \to Y$ 是光滑射影簇之间的态射，且
$\dim(X) = d$、$\dim(Y) = e$。利用 Poincar\'e 对偶，可把{\it 推前}
$$
f_* : H^{2d - i}(X) \longrightarrow H^{2e - i}(Y)
$$
定义为线性映射 $f^* : H^i(Y) \to H^i(X)$ 的对偶映射。用公式表述：对
$a \in H^{2d - i}(X)$，元素 $f_*a \in H^{2e - i}(Y)$ 由下式刻画：
$$
\int_X f^*b \cup a = \int_Y b \cup f_*a
$$
其中 $b \in H^i(Y)$ 任意。

\begin{lemma}
\label{weil-lemma-pushforward-classical}
假设给定满足 (A) 的 (D1)、(D3)。对光滑射影簇之间的态射
$f : X \to Y$，有 $f_*(f^*b \cup a) = b \cup f_*a$。若
$g : Y \to Z$ 是光滑射影簇之间的另一态射，则
$g_* \circ f_* = (g \circ f)_*$。
\end{lemma}

\begin{proof}
第一个等式成立，因为
$$
\int_Y c \cup b \cup f_*a =
\int_X f^*c \cup f^*b \cup a =
\int_Y c \cup f_*(f^*b \cup a).
$$
第二个等式成立，因为
$$
\int_Z c \cup (g \circ f)_*a = \int_X (g \circ f)^*c \cup a =
\int_X f^* g^* c \cup a = \int_Y g^*c \cup f_*a = \int_Z c \cup g_*f_*a
$$
证明完毕。
\end{proof}

\noindent
第二个公理通过{\it K\"unneth 公式}断言 $H^*$ 保持由积给出的幺半结构：
\begin{enumerate}
\item[(B)] 设 $X$、$Y$ 是光滑射影簇。映射
$$
H^*(X) \otimes_F H^*(Y) \to H^*(X \times Y),\quad
a \otimes b \mapsto \text{pr}_1^*a \cup \text{pr}_2^*b
$$
是同构。
\end{enumerate}
第三个公理涉及周期类映射：
\begin{enumerate}
\item[(C)] 周期类映射满足下列规则：
\begin{enumerate}
\item 对光滑射影簇之间的态射 $f : X \to Y$ 及
$\beta \in \CH^*(Y)$，有 $\gamma(f^!\beta) = f^*\gamma(\beta)$；
\item 对光滑射影簇之间的态射 $f : X \to Y$ 及
$\alpha \in \CH^*(X)$，有 $\gamma(f_*\alpha) = f_*\gamma(\alpha)$；
\item 对任意光滑射影簇 $X$ 及 $\alpha, \beta \in \CH^*(X)$，有
$\gamma(\alpha \cdot \beta) = \gamma(\alpha) \cup \gamma(\beta)$；
\item $\int_{\Spec(k)} \gamma([\Spec(k)]) = 1$。
\end{enumerate}
\end{enumerate}

\begin{remark}
\label{weil-remark-replace-cup-product-classical}
设 $X$ 是光滑射影簇。得到映射
$$
H^*(X) \otimes_F H^*(X) \longrightarrow H^*(X \times X)
\xrightarrow{\Delta^*} H^*(X)
$$
其中第一个箭头如公理 (B)，而 $\Delta^*$ 是沿对角态射
$\Delta : X \to X \times X$ 的拉回。由于拉回是代数同态且
$\text{pr}_i \circ \Delta = \text{id}$，该复合即杯积。另一方面，给定
$X$ 上周期 $\alpha, \beta$，交积由公式
$$
\alpha \cdot \beta =
\Delta^!(\alpha \times \beta)
$$
定义。换言之，$\alpha \cdot \beta$ 是 $X \times X$ 上外积
$\alpha \times \beta$ 沿对角线的拉回。还要注意，在
$\CH^*(X \times X)$ 中有
$\alpha \times \beta = \text{pr}_1^*\alpha \cdot \text{pr}_2^*\beta$
（证明略）。因此，在给定公理 (C)(a) 的前提下，公理 (C)(c) 等价于
$\gamma$ 与外积相容，即 $\gamma(\alpha \times \beta)$ 等于
$\text{pr}_1^*\gamma(\alpha) \cup \text{pr}_2^*\gamma(\beta)$。
\cite{Kleiman-cycles} 正是这样表述公理 (C)(c) 的。
\end{remark}

\begin{definition}
\label{weil-definition-weil-cohomology-theory-classical}
设 $k$ 是代数闭域，$F$ 是特征为 $0$ 的域。系数在 $F$ 中的 $k$ 上一个
{\it 经典 Weil 上同调理论}，是满足 Poincar\'e 对偶、K\"unneth 公式及与周期类
相容性的数据 (D1)、(D2)、(D3)；更确切地说，它们满足 (A)、(B)、(C)。
\end{definition}

\noindent
下面作少量准备工作。

\begin{lemma}
\label{weil-lemma-degrees-cycles-classical}
设 $H^*$ 是经典 Weil 上同调理论
（定义 \ref{weil-definition-weil-cohomology-theory-classical}），$X$ 是维数为
$d$ 的光滑射影簇。则图
$$
\xymatrix{
\CH^d(X) \ar[r]_-\gamma \ar@{=}[d] &
H^{2d}(X) \ar[d]^{\int_X} \\
\CH_0(X) \ar[r]^\deg & F
}
$$
交换，其中 $\deg : \CH_0(X) \to \mathbf{Z}$ 是《Chow 同调》第
\ref{chow-section-degree-zero-cycles} 节讨论的零维周期的次数。
\end{lemma}

\begin{proof}
由公理 (C)(d)，结论对 $\Spec(k)$ 成立。设
$x : \Spec(k) \to X$ 是 $X$ 的闭点。由公理 (C)(b)，在
$H^{2d}(X)$ 中有 $\gamma([x]) = x_*\gamma([\Spec(k)])$。于是由 $x_*$ 的定义，
$\int_X \gamma([x]) = 1$。
\end{proof}

\begin{lemma}
\label{weil-lemma-trace-product-classical}
设 $H^*$ 是经典 Weil 上同调理论
（定义 \ref{weil-definition-weil-cohomology-theory-classical}），$X$、$Y$ 是光滑
射影簇。则 $\int_{X \times Y} = \int_X \otimes \int_Y$。
\end{lemma}

\begin{proof}
设 $\dim(X) = d$、$\dim(Y) = e$。由公理 (B)，
$H^{2d + 2e}(X \times Y) = H^{2d}(X) \otimes H^{2e}(Y)$；由公理 (A)(d)，
此空间为 $1$ 维。由引理 \ref{weil-lemma-degrees-cycles-classical}，这个 $1$ 维
向量空间由闭点 $(x, y)$ 的类 $\gamma([x \times y])$ 生成，并且
$\int_{X \times Y} \gamma([x \times y]) = 1$。由公理 (C)(a)、(C)(c)，
$\gamma([x \times y]) = \gamma([x]) \otimes \gamma([y])$；又有
$\int_X \gamma([x]) = 1$ 和 $\int_Y \gamma([y]) = 1$，故结论成立。
\end{proof}

\begin{lemma}
\label{weil-lemma-pr2star-classical}
设 $H^*$ 是经典 Weil 上同调理论
（定义 \ref{weil-definition-weil-cohomology-theory-classical}），$X$、$Y$ 是光滑
射影簇。则 $\text{pr}_{2, *} : H^*(X \times Y) \to H^*(Y)$ 把
$a \otimes b$ 送到 $(\int_X a) b$。
\end{lemma}

\begin{proof}
这等价于引理 \ref{weil-lemma-trace-product-classical} 的结论。
\end{proof}

\begin{lemma}
\label{weil-lemma-class-diagonal-classical}
设 $H^*$ 是经典 Weil 上同调理论
（定义 \ref{weil-definition-weil-cohomology-theory-classical}），$X$ 是维数为
$d$ 的光滑射影簇。选取 $H^i(X)$ 在 $F$ 上的基
$e_{i, j}, j = 1, \ldots, \beta_i$。利用 K\"unneth 公式写成
$$
\gamma([\Delta]) =
\sum\nolimits_{i = 0, \ldots, 2d}
\sum\nolimits_j e_{i, j} \otimes e'_{2d - i , j}
\quad\text{于}\quad
\bigoplus\nolimits_i H^i(X) \otimes_F H^{2d - i}(X)
$$
其中 $e'_{2d - i, j} \in H^{2d - i}(X)$。则
$\int_X e_{i, j} \cup e'_{2d - i, j'} = (-1)^i\delta_{jj'}$。
\end{lemma}

\begin{proof}
回忆，$\Delta^* : H^*(X \times X) \to H^*(X)$ 等于杯积映射
$H^*(X) \otimes_F H^*(X) \to H^*(X)$；见注
\ref{weil-remark-replace-cup-product-classical}。另一方面，由公理 (C)(b) 及
$\gamma([X]) = 1$，有
$\gamma([\Delta]) = \Delta_*\gamma([X]) = \Delta_*1$。事实上，
$[X] \cdot [X] = [X]$，故由公理 (C)(c)，上同调类 $\gamma([X])$ 在
$1$ 维 $F$-代数 $H^0(X)$ 中等于 $0$ 或 $1$；这里还使用了公理 (A)(d)、
(A)(b)。但 $\gamma([X])$ 不可能为零：对 $X$ 的闭点 $x$，有
$[X] \cdot [x] = [x]$，而引理 \ref{weil-lemma-degrees-cycles-classical} 给出
$\gamma([x])$ 非零。因此，由 $\Delta_*$ 的定义，
$$
\int_{X \times X} \gamma([\Delta]) \cup a \otimes b =
\int_{X \times X} \Delta_*1 \cup a \otimes b =
\int_X a \cup b
$$
另一方面，由引理 \ref{weil-lemma-trace-product-classical}，
$$
\int_{X \times X} (\sum e_{i, j} \otimes e'_{2d -i , j}) \cup a \otimes b =
\sum (\int_X a \cup e_{i, j})(\int_X e'_{2d - i, j} \cup b)
$$
注意，这里进行了两次顺序交换，故符号为 $1$。于是，若选取 $a$ 使得
$\int_X a \cup e_{i, j} = 1$ 而其他所有配对均为零，则对所有 $b$ 有
$\int_X e'_{2d - i, j} \cup b = \int_X a \cup b$；亦即
$e'_{2d - i, j} = a$。引理得证。
\end{proof}

\begin{lemma}
\label{weil-lemma-square-diagonal-classical}
设 $H^*$ 是经典 Weil 上同调理论
（定义 \ref{weil-definition-weil-cohomology-theory-classical}），$X$ 是光滑射影簇。
则
$$
\sum\nolimits_{i = 0, \ldots, 2\dim(X)} (-1)^i\dim_F H^i(X) =
\deg([\Delta] \cdot [\Delta]) = \deg(c_d(\mathcal{T}_X) \cap [X])
$$
\end{lemma}

\begin{proof}
先证右侧等式。由《Chow 同调》引理
\ref{chow-lemma-intersect-regularly-embedded}，
$[\Delta] \cdot [\Delta] = \Delta_*(\Delta^![\Delta])$。由于
$\Delta_*$ 保持 $0$-维周期的次数，只需计算 $\Delta^![\Delta]$ 的次数。类
$\Delta^![\Delta]$ 由 $[\Delta]$ 与
$\Delta \subset X \times X$ 的法层之最高 Chern 类作帽积给出（《Chow 同调》
引理 \ref{chow-lemma-gysin-fundamental}）。由于 $\Delta$ 的余法层是
$\Omega_{X/k}$（《态射》引理
\ref{morphisms-lemma-differentials-diagonal}），可知法层正是切层
$\mathcal{T}_X = \SheafHom_{\mathcal{O}_X}(\Omega_{X/k}, \mathcal{O}_X)$，
如所需。

\medskip\noindent
再证左侧等式。由引理 \ref{weil-lemma-degrees-cycles-classical}，
\begin{align*}
\deg([\Delta] \cdot [\Delta])
& =
\int_{X \times X} \gamma([\Delta]) \cup \gamma([\Delta]) \\
& =
\int_{X \times X} \Delta_*1 \cup \gamma([\Delta]) \\
& =
\int_{X \times X} \Delta_*(\Delta^*\gamma([\Delta])) \\
& =
\int_X \Delta^*\gamma([\Delta])
\end{align*}
按引理 \ref{weil-lemma-class-diagonal-classical} 写
$\gamma([\Delta]) = \sum  e_{i, j} \otimes e'_{2d - i , j}$。
回忆 $\Delta^*$ 由杯积给出，得到
$$
\int_X \sum\nolimits_{i, j} e_{i, j} \cup e'_{2d - i, j} =
\sum\nolimits_{i, j} \int_X e_{i, j} \cup e'_{2d - i, j} =
\sum\nolimits_{i, j} (-1)^i = \sum (-1)^i\beta_i
$$
正是所求。
\end{proof}




\noindent
现在如下把经典 Weil 上同调理论与动机联系起来。

\begin{lemma}
\label{weil-lemma-from-functor-to-weil-classical}
设 $k$ 是代数闭域，$F$ 是特征为 $0$ 的域。考虑对称幺半范畴之间的
$\mathbf{Q}$-线性函子
$$
G : M_k \longrightarrow \text{graded }F\text{-向量空间}
$$
并假设 $G(\mathbf{1}(1))$ 仅在次数 $-2$ 非零。则得到数据
(D1)、(D2)、(D3)，它们满足 (A)、(B)、(C) 的全部条件，唯有 (A)(c)、
(A)(d) 可能除外。
\end{lemma}

\begin{proof}
置 $H^*(X) = G(h(X))$，从而得到从光滑射影簇范畴到分次 $F$-向量空间范畴的
反变函子。由假设，存在与拉回相容的典范同构
$$
H^*(X \times Y) = G(h(X \times Y)) = G(h(X) \otimes h(Y)) =
G(h(X)) \otimes G(h(Y)) = H^*(X) \otimes H^*(Y)
$$
利用沿对角线 $\Delta : X \to X \times X$ 的拉回，得到与拉回相容的分次向量
空间典范映射
$$
H^*(X) \otimes H^*(X) = H^*(X \times X) \to H^*(X)
$$
这在 $H^*(X)$ 上定义了函子性的分次 $F$-代数结构。由于 $\Delta$ 与交换约束
$h(X) \otimes h(X) \to h(X) \otimes h(X)$（交换两个因子）交换，而 $G$ 是
对称幺半范畴之间的函子（故与交换约束相容），再由《同调代数》例
\ref{homology-example-graded-vector-spaces} 中的约定，可知 $H^*(X)$ 是分次交换
代数。因此得到数据 (D1)。

\medskip\noindent
由于 $\mathbf{1}(1)$ 在动机范畴中可逆，$G(\mathbf{1}(1))$ 在分次
$F$-向量空间范畴中也可逆。因此
$\sum_i \dim_F G^i(\mathbf{1}(1)) = 1$。按假设，它仅在次数 $-2$ 非零，故可选取
分次 $F$-向量空间的同构 $F[2] \to G(\mathbf{1}(1))$。这里及下文，
$F[n]$ 表示在次数 $-n$ 为 $F$、其他次数为零的分次 $F$-向量空间。利用与
张量积的相容性，对每个 $n \in \mathbf{Z}$ 得到与张量积相容的同构
$F[2n] \to G(\mathbf{1}(n))$。

\medskip\noindent
设 $X$ 是光滑射影簇。由引理 \ref{weil-lemma-composition-correspondences}，
$$
\CH^r(X) \otimes \mathbf{Q} = \text{Corr}^r(\Spec(k), X) =
\Hom(\mathbf{1}(-r), h(X))
$$
施用函子 $G$，得到
$$
\gamma :
\CH^r(X) \otimes \mathbf{Q} \longrightarrow
\Hom(G(\mathbf{1}(-r)), H^*(X)) = H^{2r}(X)
$$
这就是数据 (D2)。

\medskip\noindent
设 $X$ 是维数为 $d$ 的光滑射影簇。由引理
\ref{weil-lemma-composition-correspondences}，
$$
\Mor(h(X)(d), \mathbf{1}) = \Mor((X, 1, d), (\Spec(k), 1, 0)) =
\text{Corr}^{-d}(X, \Spec(k)) = \CH_d(X)
$$
因此，$\CH_d(X)$ 中周期 $[X]$ 的类定义态射
$h(X)(d) \to \mathbf{1}$。施用 $G$，得到
$$
H^*(X) \otimes F[-2d] = G(h(X)(d)) \longrightarrow G(\mathbf{1}) = F
$$
此映射除次数 $0$ 外均为零；在该次数上得到
$\int_X : H^{2d}(X) \to F$。这就是数据 (D3)。

\medskip\noindent
设 $X$ 是维数为 $d$ 的光滑射影簇。由引理 \ref{weil-lemma-dual}，
$h(X)(d)$ 是 $h(X)$ 的左对偶。因此，在分次 $F$-向量空间范畴中，
$G(h(X)(d)) = H^*(X) \otimes F[-2d]$ 是 $H^*(X)$ 的左对偶。由《同调代数》
引理 \ref{homology-lemma-left-dual-graded-vector-spaces}，
$\sum_i \dim_F H^i(X) < \infty$，且
$\epsilon : h(X)(d) \otimes h(X) \to \mathbf{1}$ 给出非退化配对
$H^{2d - i}(X) \otimes_F H^i(X) \to F$。在引理 \ref{weil-lemma-dual} 的证明中，
已经看到在等同
$$
\Hom(h(X)(d) \otimes h(X), \mathbf{1}) =
\text{Corr}^{-d}(X \times X, \Spec(k)) =
\CH_d(X \times X)
$$
下，$\epsilon$ 由 $[\Delta]$ 给出。因此，$\epsilon$ 是
$h(\Delta)(d) : h(X)(d) \otimes h(X) \to h(X)(d)$ 与
$[X] : h(X)(d) \to \mathbf{1}$ 的复合。故上述配对由杯积后接
$\int_X$ 给出。这证明了公理 (A) 的 (a)、(b)。

\medskip\noindent
公理 (B) 来自 $G$ 与张量结构相容的假设以及上面对杯积的构造。

\medskip\noindent
现在验证公理 (C)。构造 $\gamma$ 时，把 $X$ 上周期 $\alpha$ 解释为从
$\Spec(k)$ 到 $X$ 的某次对应 $a$，再施用 $G$。若
$f : Y \to X$ 是光滑射影簇之间的态射，则 $f^!\alpha$ 是 $\alpha$ 沿从
$X$ 到 $Y$ 的对应 $[\Gamma_f]$ 的推前（注意方向）；见引理
\ref{weil-lemma-functor-and-cycles}。因此，把 $f^!\alpha$ 视为从 $\Spec(k)$ 到
$Y$ 的对应时，它等于 $a \circ [\Gamma_f]$；见引理
\ref{weil-lemma-composition-correspondences}。由于 $G$ 是函子，可知 $\gamma$ 与
拉回相容，即公理 (C)(a) 成立。

\medskip\noindent
设 $f : Y \to X$ 是光滑射影簇之间的态射，
$\beta \in \CH^r(Y)$ 是 $Y$ 上周期。须对所有 $c \in H^*(X)$ 证明
$$
\int_Y \gamma(\beta) \cup f^*c = \int_X \gamma(f_*\beta) \cup c
$$
令 $a, a^t, \eta_X, \eta_Y, [X], [Y]$ 如引理
\ref{weil-lemma-prep-dual}。把 $\beta$ 视为从 $\Spec(k)$ 到 $Y$ 的 $r$ 次对应，
记为 $b$。由引理 \ref{weil-lemma-functor-and-cycles} 和
\ref{weil-lemma-composition-correspondences}，把 $f_*\beta$ 视为从
$\Spec(k)$ 到 $X$ 的对应时，它等于 $a^t \circ b$。若能证明复合
$$
h(X) = \mathbf{1} \otimes h(X)
\xrightarrow{b \otimes 1}
h(Y)(r) \otimes h(X)
\xrightarrow{1 \otimes a}
h(Y)(r) \otimes h(Y)
\xrightarrow{\eta_Y}
h(Y)(r)
\xrightarrow{[Y]}
\mathbf{1}(r - e)
$$
等于复合
$$
h(X) = \mathbf{1} \otimes h(X)
\xrightarrow{a^t \circ b \otimes 1}
h(X)(r + d - e) \otimes h(X)
\xrightarrow{\eta_X}
h(X)(r + d - e)
\xrightarrow{[X]}
\mathbf{1}(r - e)
$$
则上述等式成立。这立即来自引理 \ref{weil-lemma-prep-dual}。因此公理 (C)(b) 成立。

\medskip\noindent
为证明公理 (C)(c)，使用注
\ref{weil-remark-replace-cup-product-classical} 中的讨论。因此，只需证明
$\gamma$ 与外积相容。设 $X$、$Y$ 是光滑射影簇，$\alpha$、$\beta$ 分别是
其上的周期。以 $a$、$b$ 表示从 $\Spec(k)$ 到 $X$、$Y$ 的相应对应。则
$\alpha \times \beta$ 对应于从 $\Spec(k)$ 到
$X \otimes Y = X \times Y$ 的对应 $a \otimes b$。所求相容性因此来自 $G$ 与
两边张量结构的相容性。

\medskip\noindent
公理 (C)(d) 成立，因为周期 $[\Spec(k)]$ 对应于 $h(\Spec(k))$ 上恒等态射。
引理得证。
\end{proof}

\begin{lemma}
\label{weil-lemma-from-weil-to-functor-classical}
设 $k$ 是代数闭域，$F$ 是特征为 $0$ 的域，$H^*$ 是经典 Weil 上同调理论。
则可构造对称幺半范畴之间的 $\mathbf{Q}$-线性函子
$$
G : M_k \longrightarrow \text{graded }F\text{-向量空间}
$$
使得 $H^*(X) = G(h(X))$。
\end{lemma}

\begin{proof}
由引理 \ref{weil-lemma-characterize-motives}，只需在 $k$ 上光滑射影概形所成的范畴上
构造函子 $G$，其中态射为 $0$ 次对应，并使 $G(c_2)$ 在
$G(\mathbf{P}^1)$ 上的像是可逆分次 $F$-向量空间。由于每个光滑射影概形都典范地
分解为光滑射影簇的不交并，只需在以光滑射影簇为对象、以 $0$ 次对应为态射的范畴
上构造 $G$（略去若干细节）。

\medskip\noindent
给定光滑射影簇 $X$，置 $G(X) = H^*(X)$。

\medskip\noindent
给定光滑射影簇之间的对应 $c \in \text{Corr}^0(X, Y)$，考虑按下述法则定义的映射
$G(c) : G(X) = H^*(X) \to G(Y) = H^*(Y)$：
$$
a \longmapsto
G(c)(a) = \text{pr}_{2, *}(\gamma(c) \cup \text{pr}_1^*a)
$$
显然，$G(c)$ 关于 $c$ 加性，故为 $\mathbf{Q}$-线性。公理 (C)(a)、(C)(b)、
(C)(c) 所给出的 $\gamma$ 与拉回、推前、交积的相容性表明：若
$c' \in \text{Corr}^0(Y, Z)$，则
$G(c' \circ c) = G(c') \circ G(c)$。具体地，对 $a \in H^*(X)$ 有
\begin{align*}
(G(c') \circ G(c))(a)
& =
\text{pr}^{23}_{3, *}(\gamma(c') \cup
\text{pr}^{23, *}_2(\text{pr}^{12}_{2, *}(\gamma(c) \cup
\text{pr}^{12, *}_1a))) \\
& =
\text{pr}^{23}_{3, *}(\gamma(c') \cup
\text{pr}^{123}_{23, *}(\text{pr}^{123, *}_{12}(\gamma(c) \cup
\text{pr}^{12, *}_1 a))) \\
& =
\text{pr}^{23}_{3, *}
\text{pr}^{123}_{23, *}(
\text{pr}^{123, *}_{23}\gamma(c') \cup
\text{pr}^{123, *}_{12}\gamma(c) \cup
\text{pr}^{123, *}_1 a) \\
& =
\text{pr}^{23}_{3, *}
\text{pr}^{123}_{23, *}(
\gamma(\text{pr}^{123, *}_{23}c') \cup
\gamma(\text{pr}^{123, *}_{12}c) \cup
\text{pr}^{123, *}_1 a) \\
& =
\text{pr}^{13}_{3, *}
\text{pr}^{123}_{13, *}(
\gamma(\text{pr}^{123, *}_{23}c' \cdot \text{pr}^{123, *}_{12}c) \cup
\text{pr}^{123, *}_1 a) \\
& =
\text{pr}^{13}_{3, *}(
\gamma(\text{pr}^{123}_{13, *}(
\text{pr}^{123, *}_{23}c' \cdot \text{pr}^{123, *}_{12}c)) \cup
\text{pr}^{13, *}_1 a) \\
& =
G(c' \circ c)(a)
\end{align*}
其中使用显然的记号。第一个等式来自定义。第二个等式成立，因为
$\text{pr}^{23, *}_2 \circ \text{pr}^{12}_{2, *} =
\text{pr}^{123}_{23, *} \circ \text{pr}^{123, *}_{12}$；这立即来自引理
\ref{weil-lemma-pr2star-classical} 对沿投影推前的描述。第三个等式来自引理
\ref{weil-lemma-pushforward-classical} 以及 $H^*$ 的函子性。第四个等式来自公理
(C)(a)，以及平坦态射的 Gysin 映射与平坦拉回一致这一事实（《Chow 同调》引理
\ref{chow-lemma-lci-gysin-flat}）。第五个等式使用公理 (C)(c) 及引理
\ref{weil-lemma-pushforward-classical}，从而得到
$\text{pr}^{23}_{3, *} \circ \text{pr}^{123}_{23, *} =
\text{pr}^{13}_{3, *} \circ \text{pr}^{123}_{13, *}$。
第六个等式使用引理 \ref{weil-lemma-pushforward-classical} 中的投影公式及公理
(C)(b)，从而有 $
\text{pr}^{123}_{13, *}
\gamma(\text{pr}^{123, *}_{23}c' \cdot \text{pr}^{123, *}_{12}c) =
\gamma(\text{pr}^{123}_{13, *}(
\text{pr}^{123, *}_{23}c' \cdot \text{pr}^{123, *}_{12}c))$。
最后一个等式即定义。

\medskip\noindent
为完成 $G$ 是函子的证明，还须说明它保持恒等态射。换言之，若
$1 = [\Delta] \in \text{Corr}^0(X, X)$ 是对应范畴中的恒等态射
（见引理 \ref{weil-lemma-category-correspondences} 及其证明），则须证明
$G([\Delta]) = \text{id}$。这由引理
\ref{weil-lemma-class-diagonal-classical} 对 $\gamma([\Delta])$ 的确定以及引理
\ref{weil-lemma-pr2star-classical} 得到。至此完成在光滑射影簇及 $0$ 次对应上函子
$G$ 的构造。

\medskip\noindent
由公理 (A)(c)、(A)(d)，$G(\Spec(k)) = H^*(\Spec(k))$ 作为
$F$-代数典范同构于 $F$。K\"unneth 公理 (B) 表明该函子与张量积相容。因此，
它是对称幺半范畴之间的函子。

\medskip\noindent
还须验证 $G(c_2)$ 在 $G(\mathbf{P}^1)$ 上的像是可逆分次 $F$-向量空间
（特别地，此时尚不知道 $G$ 可延拓到 $M_k$）。由公理 (A)(d)，映射
$\int_{\mathbf{P}^1} : H^2(\mathbf{P}^1) \to F$ 是同构。由公理 (A)(b)，
$\dim_F H^0(\mathbf{P}^1) = 1$。由引理
\ref{weil-lemma-square-diagonal-classical} 和公理 (A)(c)，
$2 - \dim_F H^1(\mathbf{P}^1) = c_1(T_{\mathbf{P}^1}) = 2$，故
$H^1(\mathbf{P}^1) = 0$。因此
$$
G(\mathbf{P}^1) = H^0(\mathbf{P}^1) \oplus H^2(\mathbf{P}^1)
$$
回忆，$1 = c_0 + c_2$ 是 $\text{Corr}^0(\mathbf{P}^1, \mathbf{P}^1)$
中恒等元分解为两个正交幂等元之和；见例
\ref{weil-example-decompose-P1}。由引理 \ref{weil-lemma-inverse-h2} 的证明，有
$c_0 = a \circ b$，其中
$a \in \text{Corr}^0(\Spec(k), \mathbf{P}^1)$、
$b \in \text{Corr}^0(\mathbf{P}^1, \Spec(k))$，且在
$\text{Corr}^0(\Spec(k), \Spec(k))$ 中有 $b \circ a = 1$。由于
$F = G(\Spec(k))$，由函子性可知 $G(c_0)$ 是到直和项
$H^0(\mathbf{P}^1) \subset G(\mathbf{P}^1)$ 的投影算子。因此，
$G(c_2)$ 必为到 $H^2(\mathbf{P}^1)$ 的投影，证明完成。
\end{proof}

\begin{proposition}
\label{weil-proposition-weil-cohomology-theory-classical}
设 $k$ 是代数闭域，$F$ 是特征为 $0$ 的域。一个经典 Weil 上同调理论等价于
下述数据：对称幺半范畴之间的 $\mathbf{Q}$-线性函子
$$
G : M_k \longrightarrow \text{graded }F\text{-向量空间}
$$
连同分次 $F$-向量空间的同构
$F[2] \to G(\mathbf{1}(1))$，并且对任意光滑射影簇 $X$ 还满足
\begin{enumerate}
\item $G(h(X))$ 集中在非负次数；
\item $\dim_F G^0(h(X)) = 1$。
\end{enumerate}
\end{proposition}

\begin{proof}
给定 $G$ 和 $F[2] \to G(\mathbf{1}(1))$，置
$H^*(X) = G(h(X))$，得到数据 (D1)、(D2)、(D3)，它们满足 (A)、(B)、(C)
的全部条件，唯有 (A)(c)、(A)(d) 可能除外；见引理
\ref{weil-lemma-from-functor-to-weil-classical} 及其证明。注意，在已知公理
(A)(a)、(A)(b) 的前提下，假设 (1)、(2) 蕴含公理 (A)(c)、(A)(d)。

\medskip\noindent
反之，给定 $H^*$，由引理
\ref{weil-lemma-from-weil-to-functor-classical} 的构造得到函子 $G$。令
$X = \mathbf{P}^1, c_0, c_2$ 如例 \ref{weil-example-decompose-P1}。引理
\ref{weil-lemma-inverse-h2} 已构造动机同构
$1(-1) \to (X, c_2, 0)$。在引理
\ref{weil-lemma-from-weil-to-functor-classical} 的证明中已经看到
$G(1(-1)) = G(X, c_2, 0) = H^2(\mathbf{P}^1)[-2]$。因此，公理 (A)(d)
中的同构 $\int_{\mathbf{P}^1} : H^2(\mathbf{P}^1) \to F$ 给出同构
$G(1(-1)) \to F[-2]$，从而确定同构
$F[2] \to G(\mathbf{1}(1))$。最后，由
$G(h(X)) = H^*(X)$，假设 (1)、(2) 来自公理 (A)。
\end{proof}











\section{非闭域上的周期}
\label{weil-section-cycles-nonclosed}

\noindent
下面给出若干引理，以用于研究非代数闭基域上的动机。

\begin{lemma}
\label{weil-lemma-generated-by-separable}
设 $k$ 是域，$X$ 是 $k$ 上光滑射影概形。则 $\CH_0(X)$ 由剩余域在 $k$ 上
可分的闭点之类生成。
\end{lemma}

\begin{proof}
若 $k$ 的特征为 $0$ 或为完美域，结论立即成立。因此可设 $k$ 是特征
$p > 0$ 的无限域。

\medskip\noindent
可设 $X$ 是维数为 $d$ 的不可约概形。于是
$k' = H^0(X, \mathcal{O}_X)$ 是 $k$ 的有限可分域扩张，且 $X$ 在 $k'$ 上
几何整。见《簇》引理 \ref{varieties-lemma-smooth-geometrically-normal}、
\ref{varieties-lemma-proper-geometrically-reduced-global-sections} 和
\ref{varieties-lemma-baby-stein}。以 $k'$ 替换 $k$，从而可且确可假设
$X$ 几何整。

\medskip\noindent
设 $x \in X$ 是闭点。为证明引理，将说明 $[x] \in \CH_0(X)$ 有理等价于
若干剩余域在 $k$ 上可分的闭点之类的整数线性组合。选取一个丰沛可逆
$\mathcal{O}_X$-模 $\mathcal{L}$，置
$$
V = \{s \in H^0(X, \mathcal{L}) \mid s(x) = 0 \}
$$
用 $\mathcal{L}$ 的某个幂替换它后，可假设：(a) $\mathcal{L}$ 非常丰沛；
(b) $V$ 在 $X \setminus x$ 上生成 $\mathcal{L}$；(c) 态射
$X \setminus x \to \mathbf{P}(V)$ 是浸入；(d) 映射
$V \to \mathfrak m_x\mathcal{L}_x/\mathfrak m_x^2\mathcal{L}_x$
满射。见《态射》引理
\ref{morphisms-lemma-finite-type-ample-very-ample}、《簇》引理
\ref{varieties-lemma-generate-over-complement} 以及《性质》命题
\ref{properties-proposition-characterize-ample}。考虑集合
$$
V^d \supset U =
\{
(s_1, \ldots, s_d) \in V^d \mid s_1, \ldots, s_d
\text{ generate }
\mathfrak m_x\mathcal{L}_x/\mathfrak m_x^2\mathcal{L}_x
\text{ 相对于 }\kappa(x)
\}
$$
由于 $\mathcal{O}_{X, x}$ 是维数为 $d$ 的正则局部环，有
$\dim_{\kappa(x)}(\mathfrak m_x/\mathfrak m_x^2) = d$；故 $U$ 是 $V^d$
的非空（Zariski）开集。对 $(s_1, \ldots, s_d) \in U$，置
$H_i = Z(s_i)$。由于 $s_1, \ldots, s_d$ 生成
$\mathfrak m_x\mathcal{L}_x$，存在闭子概形 $Z \subset X$，使得按概形论意义有
$$
H_1 \cap \ldots \cap H_d = x \amalg Z
$$
由 Bertini 定理（取《簇》引理 \ref{varieties-lemma-bertini} 的形式），对一般的
$s_1 \in V$，概形 $H_1 \cap (X \setminus x)$ 在 $k$ 上光滑且维数为
$d - 1$。选定 $s_1$ 后，对一般的 $s_2 \in V$，概形
$H_1 \cap H_2 \cap (X \setminus x)$ 在 $k$ 上光滑且维数为
$d - 2$；依此继续。因此，对充分一般的
$(s_1, \ldots, s_d) \in U$，概形 $Z$ 在 $\Spec(k)$ 上 \'etale。特别地，
$H_1 \cap \ldots \cap H_d$ 的维数为 $0$，从而反复应用《Chow 同调》引理
\ref{chow-lemma-intersect-properly}（细节略）可得，在 $\CH_0(X)$ 中
$$
[H_1] \cdot \ldots \cdot [H_d] = [x] + [Z]
$$
这说明 $[x] \sim_{rat} - [Z] + [Z']$，故证明完成；其中
$Z' = H'_1 \cap \ldots \cap H'_d$ 是 $\mathcal{L}$ 的充分一般截面之消失轨迹的
一般完全交，按同样论证它在 $k$ 上 \'etale。
\end{proof}

\begin{lemma}
\label{weil-lemma-chow-limit}
设 $K/k$ 是代数域扩张，$X$ 是 $k$ 上有限型概形。则
$\CH_i(X_K) = \colim \CH_i(X_{k'})$，其中余极限取遍满足 $k'/k$ 有限的子扩张
$K/k'/k$。
\end{lemma}

\begin{proof}
这是《Chow 同调》引理 \ref{chow-lemma-chow-limit} 的特殊情形。
\end{proof}

\begin{lemma}
\label{weil-lemma-divide-difference-points}
设 $k$ 是域，$X$ 是 $k$ 上几何不可约光滑射影概形，
$x, x' \in X$ 是 $k$-有理点，$n$ 是在 $k$ 中可逆的整数。则存在有限可分扩张
$k'/k$，使 $[x] - [x']$ 到 $X_{k'}$ 的拉回在
$\CH_0(X_{k'})$ 中可被 $n$ 整除。
\end{lemma}

\begin{proof}
设 $k'$ 是 $k$ 的可分代数闭包。若能证明 $[x] - [x']$ 到 $X_{k'}$ 的拉回在
$\CH_0(X_{k'})$ 中可被 $n$ 整除，则由引理 \ref{weil-lemma-chow-limit} 即得结论。
因此，可且确可假设 $k$ 可分闭。

\medskip\noindent
假设 $\dim(X) > 1$。设 $\mathcal{L}$ 是 $X$ 上丰沛可逆层，置
$$
V = \{s \in H^0(X, \mathcal{L}) \mid s(x) = 0\text{ 且 }s(x') = 0 \}
$$
以 $\mathcal{L}$ 的某个幂替换它后，由《簇》引理
\ref{varieties-lemma-generate-over-complement} 和
\ref{varieties-lemma-bertini}，对一般的 $v \in V$，相应除子
$H_v \subset X$ 在 $x$、$x'$ 之外光滑。为找到 $v$，使用 $k$ 为无限域这一事实
（因其可分闭）。若选取一般的 $s$，则 $s$ 在
$\mathfrak m_x\mathcal{L}_x/\mathfrak m_x^2\mathcal{L}_x$ 中的像非零，
这蕴含 $H_v$ 在 $x$ 处光滑（细节略）；对 $x'$ 同理。因此 $H_v$ 光滑。
将一切基变换到 $k$ 的代数闭包后应用《簇》引理
\ref{varieties-lemma-connectedness-ample-divisor}，可知 $H_v$ 几何连通。
只需把 $[x] - [x']$ 视为 $\CH_0(H_v)$ 的元素来证明结论。这样就归约到曲线情形。

\medskip\noindent
假设 $X$ 是曲线。则 $\mathcal{O}_X(x - x')$ 定义
$J = \underline{\Pic}^0_{X/k}$ 的一个 $k$-有理点 $g$；见《曲线的 Picard 概形》
引理 \ref{pic-lemma-picard-pieces}。回忆，$J$ 是 $k$ 上固有光滑簇，同时也是
$k$ 上群概形（同一引文），故 $J$ 几何整（见《簇》引理
\ref{varieties-lemma-geometrically-connected-criterion} 和
\ref{varieties-lemma-smooth-geometrically-normal}）。换言之，$J$ 是阿贝尔簇；
见《群胚》定义 \ref{groupoids-definition-abelian-variety}。由《群胚》命题
\ref{groupoids-proposition-review-abelian-varieties}，
$[n] : J \to J$ 有限 \'etale（这里使用了 $n$ 在 $k$ 中可逆）。由于 $k$ 可分闭，
存在 $g' \in J(k)$ 使 $g = [n](g')$。若 $\mathcal{L}$ 是与 $g'$ 对应的
$X$ 上次数为 $0$ 的可逆模，则
$\mathcal{O}_X(x - x') \cong \mathcal{L}^{\otimes n}$，正是所求。
\end{proof}

\begin{lemma}
\label{weil-lemma-kernel-to-closure}
设 $K/k$ 是代数域扩张，$X$ 是 $k$ 上有限型概形。引理
\ref{weil-lemma-chow-limit} 所构造映射
$\CH_i(X) \to \CH_i(X_K)$ 的核是挠群。
\end{lemma}

\begin{proof}
显然，只需对任意有限扩张 $k'/k$ 证明：沿
$\pi : X_{k'} \to X$ 的平坦拉回
$\CH_i(X) \to \CH_i(X_{k'})$ 之核为挠群。由《Chow 同调》引理
\ref{chow-lemma-finite-flat}，
$\pi_* \pi^* \alpha = [k' : k] \alpha$，故结论显然。
\end{proof}

\begin{lemma}[Voevodsky]
\label{weil-lemma-smash-nilpotence}
\begin{reference}
\cite{nilpotence}
\end{reference}
设 $k$ 是域，$X$ 是 $k$ 上几何不可约光滑射影概形，
$x, x' \in X$ 是 $k$-有理点。当 $n$ 充分大时，零维周期
$$
([x] - [x']) \times \ldots \times ([x] - [x']) \in
\CH_0(X^n)
$$
的类是挠元。
\end{lemma}

\begin{proof}
若能在基变换到 $k$ 的代数闭包后证明此结论，则由引理
\ref{weil-lemma-kernel-to-closure}，拉回之核为挠群，故原域 $k$ 上的结论随之成立。
因此可且确可假设 $k$ 代数闭。

\medskip\noindent
利用 Bertini 定理，可选取一条经过 $x$、$x'$ 的光滑曲线
$C \subset X$；见引理 \ref{weil-lemma-divide-difference-points} 的证明。因此可假设
$X$ 是曲线。

\medskip\noindent
设 $X$ 是曲线且 $k$ 代数闭。采用《曲线的 Picard 概形》第
\ref{pic-section-hilbert-scheme-points}、\ref{pic-section-divisors} 节的记号，
写 $S^n(X) = \underline{\Hilbfunctor}^n_{X/k}$。存在典范态射
$$
\pi : X^n \longrightarrow S^n(X)
$$
它把 $k$-有理点 $(x_1, \ldots, x_n)$ 送到与 $X$ 上除子
$[x_1] + \ldots + [x_n]$ 对应的 $k$-有理点。对称群 $S_n$ 忠实地作用在
$X^n$ 上。态射 $\pi$ 是 $S_n$-不变的，且 $\pi$ 的纤维在集合论意义下是
$S_n$-轨道。最后，$\pi$ 是次数为 $n!$ 的有限平坦态射；见《曲线的 Picard 概形》
引理 \ref{pic-lemma-universal-object}。

\medskip\noindent
设 $\alpha_n$ 是引理陈述中公式给出的 $X^n$ 上零维周期，并令
$\mathcal{L} = \mathcal{O}_X(x - x')$。则
$c_1(\mathcal{L}) \cap [X] = [x] - [x']$，故
$$
\alpha_n = c_1(\mathcal{L}_1) \cap \ldots \cap c_1(\mathcal{L}_n) \cap [X^n]
$$
其中 $\mathcal{L}_i = \text{pr}_i^*\mathcal{L}$，而
$\text{pr}_i : X^n \to X$ 是第 $i$ 个投影。由《除子》引理
\ref{divisors-lemma-finite-locally-free-has-norm} 或
\ref{divisors-lemma-norm-in-normal-case}，$\pi$ 有范数。置
$\mathcal{N} = \text{Norm}_\pi(\mathcal{L}_1)$；见《除子》引理
\ref{divisors-lemma-norm-invertible}。计算表明，在 $\Pic(X^n)$ 中有
$$
\pi^*\mathcal{N} =
(\mathcal{L}_1 \otimes \ldots \otimes \mathcal{L}_n)^{\otimes (n - 1)!}
$$
细节略；提示：这是因为
$\text{Norm}_\pi : \pi_*\mathcal{O}_{X^n} \to \mathcal{O}_{S^n(X)}$
与自然映射 $\pi_*\mathcal{O}_{S^n(X)} \to \mathcal{O}_{X^n}$ 的复合，等于
遍历所有 $\sigma \in S_n$ 后，$\sigma$ 在
$\pi_*\mathcal{O}_{X^n}$ 上作用的乘积。考虑 $\CH_0(S^n(X))$ 中的
$$
\beta_n = c_1(\mathcal{N})^n \cap [S^n(X)]
$$
注意，$c_1(\mathcal{L}_i) \cap c_1(\mathcal{L}_i) = 0$，因为
$\mathcal{L}_i$ 从曲线拉回；见《Chow 同调》引理
\ref{chow-lemma-vanish-above-dimension}。于是
\begin{align*}
\pi^*\beta_n
& =
((n - 1)!)^n
(\sum\nolimits_{i = 1, \ldots, n} c_1(\mathcal{L}_i))^n \cap [X^n] \\
& =
((n - 1)!)^n n^n 
c_1(\mathcal{L}_1) \cap \ldots \cap c_1(\mathcal{L}_n) \cap [X^n] \\
& =
(n!)^n \alpha_n
\end{align*}
因此只需证明 $\beta_n$ 是挠元。

\medskip\noindent
存在典范态射
$$
f : S^n(X) \longrightarrow \underline{\Picardfunctor}^n_{X/k}
$$
见《曲线的 Picard 概形》引理 \ref{pic-lemma-picard-pieces}。当
$n \geq 2g - 1$ 时，该态射是射影空间丛（细节略；比较《曲线的 Picard 概形》
引理 \ref{pic-lemma-picard-pieces} 的证明）。如下文所证，可逆层
$\mathcal{N}$ 在 $f$ 的纤维上平凡。因此，由射影空间丛公式（《Chow 同调》引理
\ref{chow-lemma-chow-ring-projective-bundle}），存在
$\underline{\Picardfunctor}^n_{X/k}$ 上可逆模 $\mathcal{M}$，使
$\mathcal{N} = f^*\mathcal{M}$。于是显然
$$
c_1(\mathcal{N})^n = f^*(c_1(\mathcal{M})^n)
$$
为零，因为 $n > g = \dim(\underline{\Picardfunctor}^n_{X/k})$，并可如前使用
《Chow 同调》引理 \ref{chow-lemma-vanish-above-dimension}。

\medskip\noindent
还须证明 $\mathcal{N}$ 在 $f$ 的一个纤维 $F$ 上平凡。由于 $f$ 的纤维是
射影空间，且 $\Pic(\mathbf{P}^m_k) = \mathbf{Z}$（《除子》引理
\ref{divisors-lemma-Pic-projective-space-UFD}），可通过计算
$\mathcal{N}$ 在纤维内一条直线上的次数来证明。这里改为证明
$\mathcal{N}$ 代数等价于零。首先断言：存在 $k$ 上连通有限型概形 $T$、
$T \times X$ 上可逆模 $\mathcal{L}'$ 以及 $k$-有理点
$p, q \in T$，使得
$\mathcal{M}_p \cong \mathcal{O}_X$ 且 $\mathcal{M}_q = \mathcal{L}$。
事实上，由于 $\mathcal{L} = \mathcal{O}_X(x - x')$，可取
$T = X$、$p = x'$、$q = x$ 以及
$\mathcal{L}' = \mathcal{O}_{X \times X}(\Delta)
\otimes \text{pr}_2^*\mathcal{O}_X(-x')$。
对 $i = 1, \ldots, n$，令 $T \times X^n$ 上的 $\mathcal{L}'_i$ 为
$\mathcal{L}'$ 沿
$\text{id}_T \times \text{pr}_i : T \times X^n \to T \times X$
的拉回。最后，令 $T \times S^n(X)$ 上
$\mathcal{N}' = \text{Norm}_{\text{id}_T \times \pi}(\mathcal{L}'_1)$。
按构造，有 $\mathcal{N}'_p = \mathcal{O}_{S^n(X)}$ 及
$\mathcal{N}'_q = \mathcal{N}$。由此
$$
\mathcal{N}'|_{T \times F}
$$
是 $T \times F \cong T \times \mathbf{P}^m_k$ 上可逆模，其在 $p$ 上的纤维
是平凡可逆模，在 $q$ 上的纤维为 $\mathcal{N}|_F$。平凡丛的欧拉示性数为
$1$，且该欧拉示性数在族中局部常值（《概形的导出范畴》引理
\ref{perfect-lemma-chi-locally-constant-geometric}），故对所有
$s \in \mathbf{Z}$，有 $\chi(F, \mathcal{N}^{\otimes s}|_F) = 1$。
这只有在 $\mathcal{N}|_F \cong \mathcal{O}_F$ 时才可能（见《概形上同调》引理
\ref{coherent-lemma-cohomology-projective-space-over-ring}），证明完成。
略去若干细节。
\end{proof}















\section{Weil 上同调理论（一）}
\label{weil-section-axioms}

\noindent
本节给出第 \ref{weil-section-axioms-classical} 节在任意域上的对应形式。换言之，
我们要确定怎样的数据与公理对应于从动机范畴到分次向量空间范畴的对称幺半函子
$G$，并要求 $G(\mathbf{1}(1))$ 集中在次数 $-2$。在第 \ref{weil-section-old} 节中，
我们将再加入一个附加条件，从而定义 Weil 上同调理论。

\medskip\noindent
固定一个域 $k$（基域）。
固定一个特征为 $0$ 的域 $F$（系数域）。
所需数据如下：
\begin{enumerate}
\item[(D0)] 一个 $1$ 维 $F$-向量空间 $F(1)$。
\item[(D1)] 从 $k$ 上光滑射影概形范畴到分次交换 $F$-代数范畴的反变函子 $H^*$。
\item[(D2)] 对每个 $k$ 上光滑射影概形 $X$，一个群同态
$\gamma : \CH^i(X) \to H^{2i}(X)(i)$。
\item[(D3)] 对每个 $k$ 上非空、等维且维数为 $d$ 的光滑射影概形 $X$，一个映射
$\int_X : H^{2d}(X)(d) \to F$。
\end{enumerate}
下面作若干说明，以阐明这些数据的含义并引入相关术语。

\medskip\noindent
关于 (D0) 的说明。
向量空间 $F(1)$ 在 $F$-向量空间范畴上给出 {\it Tate 扭转}。具体地，
对 $n \in \mathbf{Z}$，若 $n \geq 0$，置 $F(n) = F(1)^{\otimes n}$；
置 $F(-1) = \Hom_F(F(1), F)$；若 $n < 0$，则置
$F(n) = F(-1)^{\otimes - n}$。请与《进一步的代数》第
\ref{more-algebra-section-picard} 节比较。对 $F$-向量空间 $V$，定义其
{\it 第 $n$ 个 Tate 扭转}为
$$
V(n) = V \otimes_F F(n)
$$
下文采用自明的记号。例如，给定 $F$-向量空间 $U$、$V$、$W$ 以及线性映射
$U \otimes_F V \to W$，对 $n, m \in \mathbf{Z}$ 可得线性映射
$U(n) \otimes_F V(m) \to W(n + m)$。

\medskip\noindent
关于 (D1) 的说明。
给定 $k$ 上光滑射影概形 $X$，称 $H^*(X)$ 为 $X$ 的{\it 上同调}。
给定 $k$ 上光滑射影概形之间的态射 $f : X \to Y$，将映射 $H^*(f)$ 记为
$f^* : H^*(Y) \to H^*(X)$，并称之为{\it 拉回映射}。

\medskip\noindent
关于 (D2) 的说明。映射 $\gamma$ 称为{\it 周期类映射}，
$\gamma(\alpha)$ 称为 $\alpha$ 的{\it 上同调类}。
若 $Z \subset Y \subset X$ 是闭子概形，$Y$ 与 $X$ 均为 $k$ 上光滑射影概形，
且 $Z$ 整，则 $[Z]$ 既可能表示周期 $[Z]$ 在 $\CH^*(Y)$ 中的类，
也可能表示其在 $\CH^*(X)$ 中的类。此时记号 $\gamma([Z])$ 有歧义，
其含义须由上下文判定。

\medskip\noindent
关于 (D3) 的说明。映射 $\int_X$ 有时称为{\it 迹映射}，
也记作 $\text{Tr}_X$。

\medskip\noindent
第一条公理通常称为 {\it Poincar\'e 对偶}：
\begin{enumerate}
\item[(A)] 设 $X$ 是 $k$ 上非空、等维且维数为 $d$ 的光滑射影概形。则
\begin{enumerate}
\item 对所有 $i$，均有 $\dim_F H^i(X) < \infty$；
\item 对所有 $i$，配对
$H^i(X) \times H^{2d - i}(X)(d) \rightarrow
H^{2d}(X)(d) \rightarrow F$
是完美配对，其中最后一个映射是迹映射 $\int_X$。
\end{enumerate}
\end{enumerate}
设 $f : X \to Y$ 是非空光滑射影概形之间的态射，其中 $X$ 等维且维数为 $d$，
$Y$ 等维且维数为 $e$。利用 Poincar\'e 对偶，可把{\it 推前}
$$
f_* : H^{2d - i}(X)(d) \longrightarrow H^{2e - i}(Y)(e)
$$
定义为线性映射 $f^* : H^i(Y) \to H^i(X)$ 的对偶映射。用公式说，
对 $a \in H^{2d - i}(X)(d)$，元素 $f_*a \in H^{2e - i}(Y)(e)$ 由下式刻画：
$$
\int_X f^*b \cup a = \int_Y b \cup f_*a
$$
其中 $b \in H^i(Y)$ 任意。

\begin{lemma}
\label{weil-lemma-pushforward}
设给定满足 (A) 的数据 (D0)、(D1) 与 (D3)。若 $f : X \to Y$ 是
$k$ 上非空等维光滑射影概形之间的态射，则
$f_*(f^*b \cup a) = b \cup f_*a$。若 $g : Y \to Z$ 是第二个态射，且
$Z$ 非空、光滑、射影并等维，则 $g_* \circ f_* = (g \circ f)_*$。
\end{lemma}

\begin{proof}
第一个等式成立，因为
$$
\int_Y c \cup b \cup f_*a =
\int_X f^*c \cup f^*b \cup a =
\int_Y c \cup f_*(f^*b \cup a).
$$
第二个等式成立，因为
$$
\int_Z c \cup (g \circ f)_*a = \int_X (g \circ f)^*c \cup a =
\int_X f^* g^* c \cup a = \int_Y g^*c \cup f_*a = \int_Z c \cup g_*f_*a
$$
证明完毕。
\end{proof}

\noindent
第二条公理通过 {\it K\"unneth 公式}说明，$H^*$ 保持由乘积给出的幺半结构：
\begin{enumerate}
\item[(B)] 设 $X$ 与 $Y$ 是 $k$ 上光滑射影概形。
\begin{enumerate}
\item 映射
$H^*(X) \otimes_F H^*(Y) \to H^*(X \times Y)$，
$\alpha \otimes \beta \mapsto \text{pr}_1^*\alpha \cup \text{pr}_2^*\beta$
是同构；
\item 若 $X$ 与 $Y$ 非空且等维，则经由 (a) 有
$\int_{X \times Y} = \int_X \otimes \int_Y$。
\end{enumerate}
\end{enumerate}
利用公理 (B)(b)，可以计算沿投影的推前。

\begin{lemma}
\label{weil-lemma-pr2star}
设给定满足 (A) 与 (B) 的数据 (D0)、(D1) 与 (D3)。
设 $X$ 与 $Y$ 是 $k$ 上非空等维光滑射影概形，维数分别为 $d$ 与 $e$。则
$\text{pr}_{2, *} : H^*(X \times Y)(d + e) \to H^*(Y)(e)$ 把
$a \otimes b$ 送到 $(\int_X a) b$。
\end{lemma}

\begin{proof}
这由公理 (B)(a) 与 (B)(b) 立即得到。
\end{proof}

\noindent
第三条公理涉及周期类映射：
\begin{enumerate}
\item[(C)] 周期类映射满足下列规则：
\begin{enumerate}
\item 对 $k$ 上光滑射影概形之间的态射 $f : X \to Y$ 及
$\beta \in \CH^*(Y)$，有 $\gamma(f^!\beta) = f^*\gamma(\beta)$；
\item 对 $k$ 上非空等维光滑射影概形之间的态射 $f : X \to Y$ 及
$\alpha \in \CH^*(X)$，有 $\gamma(f_*\alpha) = f_*\gamma(\alpha)$；
\item 对任意 $k$ 上光滑射影概形 $X$ 及
$\alpha, \beta \in \CH^*(X)$，有
$\gamma(\alpha \cdot \beta) = \gamma(\alpha) \cup \gamma(\beta)$；以及
\item $\int_{\Spec(k)} \gamma([\Spec(k)]) = 1$。
\end{enumerate}
\end{enumerate}
下面说明公理 (C)(b)。设 $f : X \to Y$ 如 (C)(b)，且
$\dim(X) = d$、$\dim(Y) = e$。Chow 群上的推前给出
$$
f_* : \CH^{d - i}(X) = \CH_i(X) \to \CH_i(Y) = \CH^{e - i}(Y)
$$
设 $\alpha \in \CH^{d - i}(X)$。一方面，
$f_*\alpha \in \CH^{e - i}(Y)$，因而
$\gamma(f_*\alpha) \in H^{2e - 2i}(Y)(e - i)$。
另一方面，
$\gamma(\alpha) \in H^{2d - 2i}(X)(d - i)$，因而同样有
$f_*\gamma(\alpha) \in H^{2e - 2i}(Y)(e - i)$。
所以条件 $\gamma(f_*\alpha) = f_*\gamma(\alpha)$ 确有意义。

\begin{remark}
\label{weil-remark-replace-cup-product}
设给定满足 (A)、(B) 与 (C)(a) 的数据 (D0)、(D1)、(D2) 与 (D3)。
设 $X$ 是 $k$ 上光滑射影概形。得到映射
$$
H^*(X) \otimes_F H^*(X) \longrightarrow H^*(X \times X)
\xrightarrow{\Delta^*} H^*(X)
$$
其中第一支箭头如公理 (B)，而 $\Delta^*$ 是沿对角态射
$\Delta : X \to X \times X$ 的拉回。由于拉回是代数同态，且
$\text{pr}_i \circ \Delta = \text{id}$，这一复合就是杯积。
另一方面，给定 $X$ 上的周期 $\alpha, \beta$，相交积由公式
$$
\alpha \cdot \beta =
\Delta^!(\alpha \times \beta)
$$
定义。换言之，$\alpha \cdot \beta$ 是把 $X \times X$ 上的外积
$\alpha \times \beta$ 沿对角线拉回所得。还要注意，
$\CH^*(X \times X)$ 中有
$\alpha \times \beta = \text{pr}_1^*\alpha \cdot \text{pr}_2^*\beta$
（证明略）。因此，在公理 (C)(a) 成立时，公理 (C)(c) 等价于
$\gamma$ 与外积相容，即 $\gamma(\alpha \times \beta)$ 等于
$\text{pr}_1^*\gamma(\alpha) \cup \text{pr}_2^*\gamma(\beta)$。
\end{remark}

\begin{lemma}
\label{weil-lemma-base}
设给定满足 (A)、(B) 与 (C) 的数据 (D0)、(D1)、(D2) 与 (D3)。
则当 $i \not = 0$ 时 $H^i(\Spec(k)) = 0$，且存在唯一的 $F$-代数同构
$F = H^0(\Spec(k))$。此外，$\gamma([\Spec(k)]) = 1$ 且
$\int_{\Spec(k)} 1 = 1$。
\end{lemma}

\begin{proof}
由公理 (C)(d)，$H^0(\Spec(k))$ 非零，甚至 $\gamma([\Spec(k)])$ 也非零。
由于 $\Spec(k) \times \Spec(k) = \Spec(k)$，公理 (B)(a) 给出
$$
H^*(\Spec(k)) \otimes_F H^*(\Spec(k)) = H^*(\Spec(k))
$$
比较维数可知，只有 $H^0$ 非零，且其维数为 $1$。于是存在唯一的
$F$-代数同构 $F = H^0(\Spec(k))$，它把 $1 \in F$ 送到
$1 \in H^0(\Spec(k))$。在 $\Spec(k)$ 的 Chow 环中，
$[\Spec(k)] \cdot [\Spec(k)] = [\Spec(k)]$，故由公理 (C)(c) 得
$\gamma([\Spec(k)) \cup \gamma([\Spec(k)]) = \gamma([\Spec(k)])$。
既已知道 $\gamma([\Spec(k)])$ 非零，它就必等于 $1$。最后，公理 (C)(d) 给出
$\int_{\Spec(k)} \gamma([\Spec(k)]) = 1$，故
$\int_{\Spec(k)} 1 = 1$。
\end{proof}

\begin{lemma}
\label{weil-lemma-unit}
设给定满足 (A)、(B) 与 (C) 的数据 (D0)、(D1)、(D2) 与 (D3)。
设 $X$ 是 $k$ 上光滑射影概形。若 $X = \emptyset$，则 $H^*(X) = 0$。
若 $X$ 非空，则 $\gamma([X]) = 1$，且在 $H^0(X)$ 中 $1 \not = 0$。
\end{lemma}

\begin{proof}
先设 $X$ 非空。注意，$[X]$ 是 $[\Spec(k)]$ 沿结构态射
$p : X \to \Spec(k)$ 的拉回。因此，由公理 (C)(a) 与引理
\ref{weil-lemma-base}，有 $\gamma([X]) = 1$。设 $X' \subset X$ 是一个不可约分支。
由函子性，只需证明 $H^0(X')$ 中 $1 \not = 0$。所以可且确可假设
$X$ 不可约；特别地，它非空且等维，设其维数为 $d$。
要证明 $1 \not = 0$，只需证明 $H^*(X)$ 非零。

\medskip\noindent
取闭点 $x \in X$，使其剩余域 $k'$ 在 $k$ 上可分；见《簇》引理
\ref{varieties-lemma-smooth-separable-closed-points-dense}。
设 $i : \Spec(k') \to X$ 为包含态射，并记结构态射为
$p : X \to \Spec(k)$。注意，在 $\CH_0(\Spec(k))$ 中有
$p_*i_*[\Spec(k')] = [k' : k][\Spec(k)]$。两次使用公理 (C)(b)，再用引理
\ref{weil-lemma-base}，得到
$$
p_*i_*\gamma([\Spec(k')]) = \gamma([k' : k][\Spec(k)]) = [k' : k]
\in F = H^0(\Spec(k))
$$
非零。因此 $i_*\gamma([\Spec(k)]) \in H^{2d}(X)(d)$ 非零
（因为它经 $p_*$ 映到非零元素）。这就证明了 $X$ 非空的情形。

\medskip\noindent
最后考虑空概形。公理 (B)(a) 给出
$H^*(\emptyset) \otimes H^*(\emptyset) = H^*(\emptyset)$，从而
$H^*(\emptyset)$ 要么为零，要么在次数 $0$ 上为 $1$ 维。再次使用公理 (B)(a)，
可知对所有 $k$ 上光滑射影概形 $X$，均有
$H^*(\emptyset) \otimes H^*(X) = H^*(\emptyset)$。由公理 (A)(b) 及上面已经证明的
$H^0(X)$ 非零可知，当 $\dim(X) > 0$ 时，$H^*(X)$ 至少在两个次数上非零。
这迫使 $H^*(\emptyset)$ 为零。
\end{proof}

\begin{lemma}
\label{weil-lemma-push-unit}
设给定满足 (A)、(B) 与 (C) 的数据 (D0)、(D1)、(D2) 与 (D3)。
设 $i : X \to Y$ 是 $k$ 上非空等维光滑射影概形之间的闭浸入。则在
$H^{2c}(Y)(c)$ 中有 $\gamma([X]) = i_*1$，其中
$c = \dim(Y) - \dim(X)$。
\end{lemma}

\begin{proof}
由引理 \ref{weil-lemma-unit}，在 $H^0(X)$ 中有 $1 = \gamma([X])$；
再应用公理 (C)(b) 即得。
\end{proof}

\begin{lemma}
\label{weil-lemma-class-diagonal}
设给定满足 (A)、(B) 与 (C) 的数据 (D0)、(D1)、(D2) 与 (D3)。
设 $X$ 是 $k$ 上非空等维光滑射影概形，维数为 $d$。选取
$H^i(X)$ 在 $F$ 上的一组基 $e_{i, j}, j = 1, \ldots, \beta_i$。
利用 K\"unneth 公式写成
$$
\gamma([\Delta]) =
\sum\nolimits_i
\sum\nolimits_j e_{i, j} \otimes e'_{2d - i , j}
\quad\text{于}\quad
\bigoplus\nolimits_i H^i(X) \otimes_F H^{2d - i}(X)(d)
$$
其中 $e'_{2d - i, j} \in H^{2d - i}(X)(d)$。则
$\int_X e_{i, j} \cup e'_{2d - i, j'} = (-1)^i\delta_{jj'}$。
\end{lemma}

\begin{proof}
回忆，$\Delta^* : H^*(X \times X) \to H^*(X)$ 等于杯积映射
$H^*(X) \otimes_F H^*(X) \to H^*(X)$；见注记
\ref{weil-remark-replace-cup-product}。另一方面，由引理 \ref{weil-lemma-push-unit}，
$\gamma([\Delta]) = \Delta_*1$，所以引理 \ref{weil-lemma-pushforward} 给出
$$
\int_{X \times X} \gamma([\Delta]) \cup a \otimes b =
\int_{X \times X} \Delta_*1 \cup a \otimes b =
\int_X a \cup b
$$
另一方面，由公理 (B)(b)，
$$
\int_{X \times X} (\sum e_{i, j} \otimes e'_{2d -i , j}) \cup a \otimes b =
\sum (\int_X a \cup e_{i, j})(\int_X e'_{2d - i, j} \cup b)
$$
这里交换了两次次序，故每一项的符号均为 $1$。于是，若选取 $a$，使
$\int_X a \cup e_{i, j} = 1$，而所有其他配对均为零，则对所有 $b$ 有
$\int_X e'_{2d - i, j} \cup b = \int_X a \cup b$；也就是说，
$e'_{2d - i, j} = a$。引理得证。
\end{proof}

\begin{lemma}
\label{weil-lemma-cohomology-P1}
设给定满足 (A)、(B) 与 (C) 的数据 (D0)、(D1)、(D2) 与 (D3)。
则 $H^*(\mathbf{P}^1_k)$ 在次数 $0$ 与 $2$ 上均为 $1$ 维，在其他次数上为零。
\end{lemma}

\begin{proof}
设 $x \in \mathbf{P}^1_k$ 是一个 $k$-有理点。注意，在
$\mathbf{P}^1_k \times \mathbf{P}^1_k$ 上作为除子有
$\Delta = \text{pr}_1^*x + \text{pr}_2^*x$。利用公理 (C)(a) 与
$\gamma$ 的可加性，得到
$$
\gamma([\Delta]) =
\text{pr}_1^*\gamma([x]) +
\text{pr}_2^*\gamma([x]) =
\gamma([x]) \otimes 1 + 1 \otimes \gamma([x])
$$
其位于
$H^*(\mathbf{P}^1_k \times \mathbf{P}^1_k) =
H^*(\mathbf{P}^1_k) \otimes_F H^*(\mathbf{P}^1_k)$ 中。
然而，由引理 \ref{weil-lemma-class-diagonal}，$\gamma([\Delta])$ 不可能写成少于
$\sum \beta_i$ 个纯张量之和，其中
$\beta_i = \dim_F H^i(\mathbf{P}^1_k)$。因此
$\sum \beta_i \leq 2$。由引理 \ref{weil-lemma-unit}，有
$H^0(\mathbf{P}^1_k) \not = 0$。由 Poincar\'e 对偶，更准确地说由公理 (A)(b)，
$\beta_0 = \beta_2$。故结论成立。
\end{proof}

\begin{lemma}
\label{weil-lemma-weil-additive}
设给定满足 (A)、(B) 与 (C) 的数据 (D0)、(D1)、(D2) 与 (D3)。
若 $X$ 与 $Y$ 是 $k$ 上光滑射影概形，则
$H^*(X \amalg Y) \to H^*(X) \times H^*(Y)$，
$a \mapsto (i^*a, j^*a)$ 是同构，其中 $i$、$j$ 是余积注入。
\end{lemma}

\begin{proof}
若 $X$ 或 $Y$ 为空，则由引理 \ref{weil-lemma-unit} 的
$H^*(\emptyset) = 0$，结论成立。因此可假设 $X$ 与 $Y$ 均非空。

\medskip\noindent
先证该映射为单射。注意，可以找到光滑射影概形的态射
$X' \to X$ 与 $Y' \to Y$，使 $X'$ 与 $Y'$ 等维且维数相同，并且
$X' \to X$ 与 $Y' \to Y$ 各有一个截面。事实上，把
$X = \coprod X_d$ 与 $Y = \coprod Y_e$ 分解为等维数 $d$ 与 $e$ 的开闭子概形，
再对某个充分大的 $n$ 取
$X' = \coprod X_d \times \mathbf{P}^{n - d}$ 和
$Y' = \coprod Y_e \times \mathbf{P}^{n - e}$。于是，沿
$X' \amalg Y' \to X \amalg Y$ 的拉回为单射（因为该态射有截面），
故只需对 $X', Y'$ 证明单射性；下一段正是如此处理。

\medskip\noindent
现在设 $X$ 与 $Y$ 等维且维数同为 $d$，证明该映射为单射。
在 $\CH^0(X \amalg Y)$ 中有 $[X \amalg Y] = [X] + [Y]$，且
$[X]$ 与 $[Y]$ 是 $\CH^0(X \amalg Y)$ 中的正交幂等元。因此
$$
1 = \gamma([X \amalg Y] = \gamma([X]) + \gamma([Y]) = i_*1 + j_*1
$$
是分解为正交幂等元。这里使用了引理 \ref{weil-lemma-unit}、
\ref{weil-lemma-push-unit} 以及公理 (C)(c)。由投影公式（引理
\ref{weil-lemma-pushforward}），
$$
a = a \cup 1 = a \cup i_*1 + a \cup j_*1 =
i_*(i^*a) + j_*(j^*a)
$$
所以该映射为单射。

\medskip\noindent
下面证明满射性。把 $e = \gamma([X])$ 与 $f = \gamma([Y])$ 视为
$H^0(X \amalg Y)$ 中的元素。由公理 (C)(a)，有
$i^*e = 1$、$i^*f = 0$、$j^*e = 0$、$j^*f = 1$。
所以，若 $i^* : H^*(X \amalg Y) \to  H^*(X)$ 与
$j^* : H^*(X \amalg Y) \to H^*(Y)$ 均为满射，则
$(i^*, j^*)$ 也是满射。事实上，对 $a, a' \in H^*(X \amalg Y)$ 有
$$
(i^*a, j^*a') = (i^*(a \cup e + a' \cup f), j^*(a \cup e + a' \cup f))
$$
由对称性，只需证明 $i^* : H^*(X \amalg Y) \to  H^*(X)$ 为满射。
若存在态射 $Y \to X$，则存在态射 $g : X \amalg Y \to X$，满足
$g \circ i = \text{id}_X$，从而结论成立。最后注意，要证明
$i^*$ 为满射，只需在张量上一个非零分次 $F$-向量空间后证明。于是，由公理 (B)(b)
以及上同调的非零性（引理 \ref{weil-lemma-unit}），只需对某个有限可分扩张
$k'/k$，把 $X$、$Y$ 分别换为 $X \times \Spec(k')$、
$Y \times \Spec(k')$ 后证明 $i^*$ 为满射。选取 $k'$，使存在闭点
$x \in X$ 满足 $\kappa(x) = k'$；这由《簇》引理
\ref{varieties-lemma-smooth-separable-closed-points-dense} 可行。此时存在态射
$Y \times \Spec(k') \to X \times \Spec(k')$，证明完成。
\end{proof}

\begin{lemma}
\label{weil-lemma-from-functor-to-weil}
设 $k$ 是域，$F$ 是特征为 $0$ 的域。设给定对称幺半范畴之间的
$\mathbf{Q}$-线性函子
$$
G : M_k \longrightarrow \text{graded }F\text{-向量空间}
$$
并且 $G(\mathbf{1}(1))$ 仅在次数 $-2$ 上非零。则可得到满足上述
(A)、(B)、(C) 全部公理的数据 (D0)、(D1)、(D2)、(D3)。
\end{lemma}

\begin{proof}
本证明与引理 \ref{weil-lemma-from-functor-to-weil-classical} 的证明相同；
我们建议读者改读该引理的证明。

\medskip\noindent
置 $H^*(X) = G(h(X))$，便得到从 $k$ 上光滑射影概形范畴到分次
$F$-向量空间范畴的反变函子。由假设，存在与拉回相容的典范同构
$$
H^*(X \times Y) = G(h(X \times Y)) = G(h(X) \otimes h(Y)) =
G(h(X)) \otimes G(h(Y)) = H^*(X) \otimes H^*(Y)
$$
沿对角态射 $\Delta : X \to X \times X$ 拉回，得到与拉回相容的分次向量空间典范映射
$$
H^*(X) \otimes H^*(X) = H^*(X \times X) \to H^*(X)
$$
这在 $H^*(X)$ 上定义了函子性的分次 $F$-代数结构。由于 $\Delta$ 与交换约束
$h(X) \otimes h(X) \to h(X) \otimes h(X)$（交换两个因子）相容，
$G$ 又是对称幺半范畴之间的函子（因而与交换约束相容），结合《同调》例
\ref{homology-example-graded-vector-spaces} 中的约定，可知 $H^*(X)$ 是分次交换代数。
由此得到数据 (D1)。

\medskip\noindent
由于 $\mathbf{1}(1)$ 在动机范畴中可逆，$G(\mathbf{1}(1))$ 在分次
$F$-向量空间范畴中也可逆。因此
$\sum_i \dim_F G^i(\mathbf{1}(1)) = 1$。根据假设，它仅在次数 $-2$ 上非零。
数据 (D0) 取为向量空间 $F(1) = G^{-2}(\mathbf{1}(1))$。
由于 $G$ 是对称幺半函子，对所有 $n \in \mathbf{Z}$ 有
$F(n) = G^{-2n}(\mathbf{1}(n))$。于是
$$
H^{2r}(X)(r) = G^{2r}(h(X)) \otimes G^{-2r}(\mathbf{1}(r)) =
G^0(h(X)(r))
$$
这个公式下文将反复使用。

\medskip\noindent
设 $X$ 是 $k$ 上光滑射影概形。由引理
\ref{weil-lemma-composition-correspondences}，
$$
\CH^r(X) \otimes \mathbf{Q} = \text{Corr}^r(\Spec(k), X) =
\Hom(\mathbf{1}(-r), h(X)) = \Hom(\mathbf{1}, h(X)(r))
$$
应用函子 $G$，得到到 $\Hom(G(\mathbf{1}), G(h(X)(r)))$ 的映射。
取 $G^0(\mathbf{1}) = F$ 中 $1$ 在
$G^0(h(X)(r)) = H^{2r}(X)(r)$ 中的像，便得到
$$
\gamma :
\CH^r(X) \otimes \mathbf{Q} \longrightarrow H^{2r}(X)(r)
$$
这就是数据 (D2)。

\medskip\noindent
设 $X$ 是 $k$ 上非空等维光滑射影概形，维数为 $d$。由引理
\ref{weil-lemma-composition-correspondences}，
$$
\Mor(h(X)(d), \mathbf{1}) = \Mor((X, 1, d), (\Spec(k), 1, 0)) =
\text{Corr}^{-d}(X, \Spec(k)) = \CH_d(X)
$$
因此，$\CH_d(X)$ 中周期 $[X]$ 的类定义态射
$h(X)(d) \to \mathbf{1}$。应用 $G$ 并取次数 $0$ 部分，得到
$$
H^{2d}(X)(d) = G^0(h(X)(d)) \longrightarrow G^0(\mathbf{1}) = F
$$
此映射 $\int_X : H^{2d}(X)(d) \to F$ 就是数据 (D3)。

\medskip\noindent
设 $X$ 是 $k$ 上非空等维光滑射影概形，维数为 $d$。由引理
\ref{weil-lemma-dual}，$h(X)(d)$ 是 $h(X)$ 的左对偶。因此
$G(h(X)(d)) = H^*(X) \otimes_F F(d)[2d]$ 是分次 $F$-向量空间范畴中
$H^*(X)$ 的左对偶。这里 $[n]$ 表示分次向量空间上的平移函子。
由《同调》引理 \ref{homology-lemma-left-dual-graded-vector-spaces}，
$\sum_i \dim_F H^i(X) < \infty$，并且
$\epsilon : h(X)(d) \otimes h(X) \to \mathbf{1}$ 给出非退化配对
$H^{2d - i}(X)(d) \otimes_F H^i(X) \to F$。
在引理 \ref{weil-lemma-dual} 的证明中已经看到，经由同一
$$
\Hom(h(X)(d) \otimes h(X), \mathbf{1}) =
\text{Corr}^{-d}(X \times X, \Spec(k)) =
\CH_d(X \times X)
$$
$\epsilon$ 由 $[\Delta]$ 给出。因此，$\epsilon$ 是
$[X] : h(X)(d) \to \mathbf{1}$ 与
$h(\Delta)(d) : h(X)(d) \otimes h(X) \to h(X)(d)$ 的复合。
所以，上述配对就是先作杯积，再应用 $\int_X$。这证明了公理 (A)。

\medskip\noindent
公理 (B) 由 $G$ 与张量结构相容的假设以及上面对杯积的构造立即得到。

\medskip\noindent
公理 (C)。在我们对 $\gamma$ 的构造中，把 $X$ 上的周期 $\alpha$ 视为从
$\Spec(k)$ 到 $X$ 的某一次数的对应 $a$，然后应用 $G$。若
$f : Y \to X$ 是 $k$ 上非空等维光滑射影概形之间的态射，则
$f^!\alpha$ 是由从 $X$ 到 $Y$ 的对应 $[\Gamma_f]$ 对 $\alpha$ 所作的推前（！）；
见引理 \ref{weil-lemma-functor-and-cycles}。因此，把 $f^!\alpha$ 视为从
$\Spec(k)$ 到 $Y$ 的对应时，它等于 $a \circ [\Gamma_f]$；见引理
\ref{weil-lemma-composition-correspondences}。由于 $G$ 是函子，可知 $\gamma$ 与拉回相容，
即公理 (C)(a) 成立。

\medskip\noindent
设 $f : Y \to X$ 是 $k$ 上非空等维光滑射影概形之间的态射，并设
$\beta \in \CH^r(Y)$ 是 $Y$ 上的周期。须证明对所有 $c \in H^*(X)$ 均有
$$
\int_Y \gamma(\beta) \cup f^*c = \int_X \gamma(f_*\beta) \cup c
$$
令 $a, a^t, \eta_X, \eta_Y, [X], [Y]$ 如引理 \ref{weil-lemma-prep-dual}。
把 $b$ 取为视作从 $\Spec(k)$ 到 $Y$ 的次数为 $r$ 的对应的 $\beta$。
由引理 \ref{weil-lemma-functor-and-cycles} 与
\ref{weil-lemma-composition-correspondences}，把 $f_*\beta$ 视为从
$\Spec(k)$ 到 $X$ 的对应时，它等于 $a^t \circ b$。若能证明
$$
h(X) = \mathbf{1} \otimes h(X)
\xrightarrow{b \otimes 1}
h(Y)(r) \otimes h(X)
\xrightarrow{1 \otimes a}
h(Y)(r) \otimes h(Y)
\xrightarrow{\eta_Y}
h(Y)(r)
\xrightarrow{[Y]}
\mathbf{1}(r - e)
$$
等于
$$
h(X) = \mathbf{1} \otimes h(X)
\xrightarrow{a^t \circ b \otimes 1}
h(X)(r + d - e) \otimes h(X)
\xrightarrow{\eta_X}
h(X)(r + d - e)
\xrightarrow{[X]}
\mathbf{1}(r - e)
$$
则上面的等式成立。这立即由引理 \ref{weil-lemma-prep-dual} 得到。
因此公理 (C)(b) 成立。

\medskip\noindent
为证明公理 (C)(c)，使用注记 \ref{weil-remark-replace-cup-product-classical} 中的讨论。
所以只需证明 $\gamma$ 与外积相容。设 $X$、$Y$ 是 $k$ 上非空光滑射影概形，
$\alpha$、$\beta$ 分别是其上的周期。以 $a$、$b$ 表示相应的从
$\Spec(k)$ 到 $X$、$Y$ 的对应。则 $\alpha \times \beta$ 对应于从
$\Spec(k)$ 到 $X \otimes Y = X \times Y$ 的对应 $a \otimes b$。
所需结论遂由 $G$ 与两侧的张量结构相容得到。

\medskip\noindent
公理 (C)(d) 成立，是因为周期 $[\Spec(k)]$ 对应于 $h(\Spec(k))$ 上的恒等态射。
引理证毕。
\end{proof}

\begin{lemma}
\label{weil-lemma-from-weil-to-functor}
设 $k$ 是域，$F$ 是特征为 $0$ 的域。给定满足 (A)、(B)、(C) 的数据
(D0)、(D1)、(D2)、(D3)，可以构造对称幺半范畴之间的 $\mathbf{Q}$-线性函子
$$
G : M_k \longrightarrow \text{graded }F\text{-向量空间}
$$
使得 $H^*(X) = G(h(X))$。
\end{lemma}

\begin{proof}
本引理的证明与引理 \ref{weil-lemma-from-weil-to-functor-classical} 的证明相同；
我们建议读者改读该引理的证明。

\medskip\noindent
由引理 \ref{weil-lemma-characterize-motives}，只需在以 $k$ 上光滑射影概形为对象、
以次数为 $0$ 的对应为态射的范畴上构造函子 $G$，并使 $G(c_2)$ 在
$G(\mathbf{P}^1_k)$ 上的像为可逆分次 $F$-向量空间。

\medskip\noindent
设 $X$ 是 $k$ 上光滑射影概形。存在典范分解
$$
X = \coprod\nolimits_{0 \leq d \le \dim(X)} X_d
$$
为开闭子概形，其中 $X_d$ 等维且维数为 $d$。相应地，由引理
\ref{weil-lemma-weil-additive} 有
$$
H^*(X) \longrightarrow \prod\nolimits_{0 \leq d \le \dim(X)} H^*(X_d)
$$
若 $Y$ 是另一个 $k$ 上光滑射影概形，并同样分解
$Y = \coprod Y_e$，则
$$
\text{Corr}^0(X, Y) = \bigoplus \text{Corr}^0(X_d, Y_e)
$$
在对应范畴中还有 $X \otimes Y = \coprod X_d \otimes Y_e$。
由这些观察可知，只需在以 $k$ 上等维光滑射影概形为对象、
以次数为 $0$ 的对应为态射的范畴上构造 $G$。（略去若干细节。）

\medskip\noindent
给定 $k$ 上等维光滑射影概形 $X$，置 $G(X) = H^*(X)$。
若 $X = \emptyset$，则 $G(X) = 0$（引理 \ref{weil-lemma-unit}）。
所以从或到 $G(\emptyset)$ 的映射均为零，下文可且确可假设所讨论的概形非空。

\medskip\noindent
给定 $k$ 上非空等维光滑射影概形之间的对应
$c \in \text{Corr}^0(X, Y)$，考虑映射
$G(c) : G(X) = H^*(X) \to G(Y) = H^*(Y)$，其规则为
$$
a \longmapsto
G(c)(a) = \text{pr}_{2, *}(\gamma(c) \cup \text{pr}_1^*a)
$$
显然，$G(c)$ 关于 $c$ 可加，因而是 $\mathbf{Q}$-线性的。
公理 (C)(a)、(C)(b)、(C)(c) 分别给出的 $\gamma$ 与拉回、推前、相交积的相容性表明，
若 $c' \in \text{Corr}^0(Y, Z)$，则
$G(c' \circ c) = G(c') \circ G(c)$。事实上，对 $a \in H^*(X)$ 有
\begin{align*}
(G(c') \circ G(c))(a)
& =
\text{pr}^{23}_{3, *}(\gamma(c') \cup
\text{pr}^{23, *}_2(\text{pr}^{12}_{2, *}(\gamma(c) \cup
\text{pr}^{12, *}_1a))) \\
& =
\text{pr}^{23}_{3, *}(\gamma(c') \cup
\text{pr}^{123}_{23, *}(\text{pr}^{123, *}_{12}(\gamma(c) \cup
\text{pr}^{12, *}_1 a))) \\
& =
\text{pr}^{23}_{3, *}
\text{pr}^{123}_{23, *}(
\text{pr}^{123, *}_{23}\gamma(c') \cup
\text{pr}^{123, *}_{12}\gamma(c) \cup
\text{pr}^{123, *}_1 a) \\
& =
\text{pr}^{23}_{3, *}
\text{pr}^{123}_{23, *}(
\gamma(\text{pr}^{123, *}_{23}c') \cup
\gamma(\text{pr}^{123, *}_{12}c) \cup
\text{pr}^{123, *}_1 a) \\
& =
\text{pr}^{13}_{3, *}
\text{pr}^{123}_{13, *}(
\gamma(\text{pr}^{123, *}_{23}c' \cdot \text{pr}^{123, *}_{12}c) \cup
\text{pr}^{123, *}_1 a) \\
& =
\text{pr}^{13}_{3, *}(
\gamma(\text{pr}^{123}_{13, *}(\text{pr}^{123, *}_{23}c' \cdot
\text{pr}^{123, *}_{12}c)) \cup
\text{pr}^{13, *}_1 a) \\
& =
G(c' \circ c)(a)
\end{align*}
这里采用自明的记号。第一个等式来自定义。第二个等式成立，是因为
$\text{pr}^{23, *}_2 \circ \text{pr}^{12}_{2, *} =
\text{pr}^{123}_{23, *} \circ \text{pr}^{123, *}_{12}$；
这立即由引理 \ref{weil-lemma-pr2star} 中沿投影推前的描述得到。
第三个等式来自引理 \ref{weil-lemma-pushforward} 以及 $H^*$ 是函子这一事实。
第四个等式来自公理 (C)(a)，以及对平坦态射 Gysin 映射与平坦拉回一致这一事实
（《Chow 同调》引理 \ref{chow-lemma-lci-gysin-flat}）。
第五个等式使用公理 (C)(c)，并由引理 \ref{weil-lemma-pushforward} 得到
$\text{pr}^{23}_{3, *} \circ \text{pr}^{123}_{23, *} =
\text{pr}^{13}_{3, *} \circ \text{pr}^{123}_{13, *}$。
第六个等式使用引理 \ref{weil-lemma-pushforward} 的投影公式以及公理 (C)(b)，从而有 $
\text{pr}^{123}_{13, *}
\gamma(\text{pr}^{123, *}_{23}c' \cdot \text{pr}^{123, *}_{12}c) =
\gamma(\text{pr}^{123}_{13, *}(
\text{pr}^{123, *}_{23}c' \cdot \text{pr}^{123, *}_{12}c))$。
最后一个等式就是定义。

\medskip\noindent
为完成 $G$ 为函子的证明，还须验证恒等态射得到保持。换言之，若
$1 = [\Delta] \in \text{Corr}^0(X, X)$ 是对应范畴中的恒等态射
（引理 \ref{weil-lemma-category-correspondences}），则须证明
$G([\Delta]) = \text{id}$。这由引理 \ref{weil-lemma-class-diagonal} 对
$\gamma([\Delta])$ 的确定以及引理 \ref{weil-lemma-pr2star} 得到。
至此，$G$ 已构造为以 $k$ 上光滑射影概形为对象、以次数为 $0$ 的对应为态射之范畴上的函子。

\medskip\noindent
由引理 \ref{weil-lemma-base}，$G(\Spec(k)) = H^*(\Spec(k))$ 作为
$F$-代数典范同构于 $F$。K\"unneth 公理 (B)(a) 表明该函子与张量积相容。
因此，它是对称幺半范畴之间的函子。

\medskip\noindent
还须验证 $G(c_2)$ 在 $G(\mathbf{P}^1_k) = H^*(\mathbf{P}^1_k)$ 上的像是
可逆分次 $F$-向量空间（特别地，此时尚不知道 $G$ 能延拓到 $M_k$）。
由引理 \ref{weil-lemma-cohomology-P1}，上同调仅在次数 $0$ 与 $2$ 上非零，且两者维数均为
$1$。由例 \ref{weil-example-decompose-P1}，$1 = c_0 + c_2$ 是
$\text{Corr}^0(\mathbf{P}^1_k, \mathbf{P}^1_k)$ 中恒等态射分解为正交幂等元之和。
此外，$c_0 = a \circ b$，其中
$a \in \text{Corr}^0(\Spec(k), \mathbf{P}^1_k)$、
$b \in \text{Corr}^0(\mathbf{P}^1_k, \Spec(k))$，且
$b \circ a = 1$ 在 $\text{Corr}^0(\Spec(k), \Spec(k))$ 中成立；
见引理 \ref{weil-lemma-inverse-h2} 的证明。因此，$G(c_0)$ 是到次数 $0$ 部分的投影算子，
从而 $G(c_2)$ 必为到次数 $2$ 部分的投影算子。证明完成。
\end{proof}

\begin{proposition}
\label{weil-proposition-weil-cohomology-theory}
设 $k$ 是域，$F$ 是特征为 $0$ 的域。下列两者之间存在 $1$ 对 $1$ 的对应：
\begin{enumerate}
\item 满足 (A)、(B)、(C) 的数据 (D0)、(D1)、(D2)、(D3)；以及
\item $\mathbf{Q}$-线性对称幺半函子
$$
G : M_k \longrightarrow \text{graded }F\text{-向量空间}
$$
且 $G(\mathbf{1}(1))$ 仅在次数 $-2$ 上非零。
\end{enumerate}
\end{proposition}

\begin{proof}
给定 (2) 中的 $G$，置 $H^*(X) = G(h(X))$，便得到满足 (A)、(B)、(C) 的数据
(D0)、(D1)、(D2)、(D3)；见引理 \ref{weil-lemma-from-functor-to-weil} 及其证明。

\medskip\noindent
反之，给定满足 (A)、(B)、(C) 的数据 (D0)、(D1)、(D2)、(D3)，
由引理 \ref{weil-lemma-from-weil-to-functor} 证明中的构造，得到 (2) 中的函子 $G$。

\medskip\noindent
略去验证这两个构造互逆的详细证明。
\end{proof}











\section{进一步的性质}
\label{weil-section-further}

\noindent
本节证明若给定第 \ref{weil-section-axioms} 节中满足 (A)、(B)、(C) 的数据
(D0)、(D1)、(D2)、(D3)，还可得到的若干结果。

\begin{lemma}
\label{weil-lemma-trace-disjoint-union}
设给定满足 (A)、(B)、(C) 的数据 (D0)、(D1)、(D2)、(D3)。设 $X, Y$ 是
$k$ 上非空光滑射影概形，二者均等维且维数为 $d$。则
$\int_{X \amalg Y} = \int_X + \int_Y$。
\end{lemma}

\begin{proof}
以 $i : X \to X \amalg Y$ 和 $j : Y \to X \amalg Y$ 表示余积注入。
由引理 \ref{weil-lemma-weil-additive}，映射
$(i^*, j^*) : H^*(X \amalg Y) \to H^*(X) \times H^*(Y)$ 是同构。
引理的断言是说，在同构
$(i^*, j^*) : H^{2d}(X \amalg Y)(d) \to H^{2d}(X)(d) \oplus H^{2d}(Y)(d)$
之下，映射 $\int_X + \int_Y$ 变为 $\int_{X \amalg Y}$。这是因为
$$
\int_{X \amalg Y} a =
\int_{X \amalg Y} i_*(i^*a) + j_*(j^*a) =
\int_X i^*a + \int_Y j^*a
$$
其中等式 $a = i_*(i^*a) + j_*(j^*a)$ 已在引理
\ref{weil-lemma-weil-additive} 的证明中得到。
\end{proof}

\begin{lemma}
\label{weil-lemma-dim-0}
设给定满足 (A)、(B)、(C) 的数据 (D0)、(D1)、(D2)、(D3)。
设 $X$ 是 $k$ 上零维光滑射影概形。则
\begin{enumerate}
\item 当 $i \not = 0$ 时，$H^i(X) = 0$；
\item $H^0(X)$ 是 $F$ 上有限可分代数；
\item $\dim_F H^0(X) = \deg(X \to \Spec(F))$；
\item $\int_X : H^0(X) \to F$ 是迹映射；
\item $\gamma([X]) = 1$；以及
\item $\int_X \gamma([X]) = \deg(X \to \Spec(k))$。
\end{enumerate}
\end{lemma}

\begin{proof}
可写 $X = \Spec(k')$，其中 $k'$ 是 $k$ 上有限可分代数。注意，
$\deg(X \to \Spec(k)) = [k' : k]$。选取有限 Galois 扩张 $k''/k$，
使其包含 $k'$ 的每个因子。（回忆，有限可分 $k$-代数是 $k$ 的有限可分域扩张之积。）
置 $\Sigma = \Hom_k(k', k'')$。于是
$$
k' \otimes_k k'' = \prod\nolimits_{\sigma \in \Sigma} k''
$$
置 $Y = \Spec(k'')$，公理 (B)(a) 与引理 \ref{weil-lemma-weil-additive} 给出分次交换
$F$-代数的等式
$$
H^*(X) \otimes_F H^*(Y) =
\prod\nolimits_{\sigma \in \Sigma} H^*(Y)
$$
由引理 \ref{weil-lemma-unit}，$F$-代数 $H^*(Y)$ 非零。比较上述等式两侧的维数，
可知 $H^*(X)$ 仅集中在次数 $0$，且 $\dim_F H^0(X) = [k' : k]$。
把此结论用于 $Y$，得到 $H^*(Y) = H^0(Y)$。由于作为 $F$-代数有
$$
H^0(X) \otimes_F H^0(Y) = H^0(Y) \times \ldots \times H^0(Y)
$$
而可在忠实平坦基变换 $F \to H^0(Y)$ 后检验可分性，故 $H^0(X)$ 是可分
$F$-代数。

\medskip\noindent
上述同构由映射
$$
H^0(X) \otimes_F H^0(Y) \longrightarrow
\prod\nolimits_{\sigma \in \Sigma} H^0(Y),\quad
a \otimes b \longmapsto \prod\nolimits_\sigma \Spec(\sigma)^*a \cup b
$$
给出。经此同构，由引理 \ref{weil-lemma-trace-disjoint-union} 有
$\int_{X \times Y} = \sum_\sigma \int_Y$。因此，在 $H^0(Y)$ 中有
$$
\int_X a = \text{pr}_{1, *}(a \otimes 1) = \sum \Spec(\sigma)^*a
$$
第一个等式来自引理 \ref{weil-lemma-pr2star}，第二个等式来自刚才的观察。
选取代数闭包 $\overline{F}$ 以及 $F$-代数映射
$\tau : H^0(Y) \to \overline{F}$。上述同构经基变换成为同构
$$
H^0(X) \otimes_F \overline{F} \longrightarrow
\prod\nolimits_{\sigma \in \Sigma} \overline{F},\quad
a \otimes b \longmapsto \prod\nolimits_\sigma \tau(\Spec(\sigma)^*a) b
$$
由此，$a \mapsto \tau(\Spec(\sigma)^*a)$ 构成 $H^0(X)$ 到
$\overline{F}$ 的全部嵌入。把 $\tau$ 应用于上面得到的 $\int_X a$ 公式，
可知 $\int_X$ 是迹映射。由引理 \ref{weil-lemma-unit}，有
$\gamma([X]) = 1$。最后，
$\int_X \gamma([X]) = \deg(X \to \Spec(k))$，因为
$\gamma([X]) = 1$，而 $1$ 的迹等于 $[k' : k]$。
\end{proof}

\begin{lemma}
\label{weil-lemma-degrees-cycles}
设给定满足 (A)、(B)、(C) 的数据 (D0)、(D1)、(D2)、(D3)。设 $X$ 是
$k$ 上非空等维光滑射影概形，维数为 $d$。则图表
$$
\xymatrix{
\CH^d(X) \ar[r]_-\gamma \ar@{=}[d] &
H^{2d}(X)(d) \ar[d]^{\int_X} \\
\CH_0(X) \ar[r]^\deg & F
}
$$
交换，其中 $\deg : \CH_0(X) \to \mathbf{Z}$ 是《Chow 同调》第
\ref{chow-section-degree-zero-cycles} 节所讨论的零维周期次数。
\end{lemma}

\begin{proof}
设 $x$ 是 $X$ 的闭点，其剩余域在 $k$ 上可分。把 $x$ 视为概形，并以
$i : x \to X$ 表示包含态射。为避免混淆，以
$\gamma' : \CH_0(x) \to H^0(x)$ 表示 $x$ 的周期类映射。则
$$
\int_X \gamma([x]) = \int_X \gamma(i_*[x]) =
\int_X i_*\gamma'([x]) = \int_x \gamma'([x]) = \deg(x \to \Spec(k))
$$
第二个等式是公理 (C)(b)，第三个等式是上同调上 $i_*$ 的定义，最后一个等式来自引理
\ref{weil-lemma-dim-0}。由引理 \ref{weil-lemma-generated-by-separable}，
$\CH_0(X)$ 由上述点 $x$ 的类生成，故引理得证。
\end{proof}

\begin{lemma}
\label{weil-lemma-square-diagonal}
设给定满足 (A)、(B)、(C) 的数据 (D0)、(D1)、(D2)、(D3)。设 $X$ 是
$k$ 上非空等维光滑射影概形，维数为 $d$。则
$$
\sum\nolimits_i (-1)^i\dim_F H^i(X) =
\deg(\Delta \cdot \Delta) = \deg(c_d(\mathcal{T}_{X/k}))
$$
\end{lemma}

\begin{proof}
先证右边的等式。由《Chow 同调》引理
\ref{chow-lemma-intersect-regularly-embedded}，有
$[\Delta] \cdot [\Delta] = \Delta_*(\Delta^![\Delta])$。
由于 $\Delta_*$ 保持 $0$-维周期的次数，只需计算 $\Delta^![\Delta]$ 的次数。
类 $\Delta^![\Delta]$ 由 $[\Delta]$ 与
$\Delta \subset X \times X$ 的法层之最高 Chern 类作帽积给出
（《Chow 同调》引理 \ref{chow-lemma-gysin-fundamental}）。
由于 $\Delta$ 的余法层是 $\Omega_{X/k}$（《态射》引理
\ref{morphisms-lemma-differentials-diagonal}），法层正是切层
$\mathcal{T}_{X/k} = \SheafHom_{\mathcal{O}_X}(\Omega_{X/k}, \mathcal{O}_X)$，
这就得到所需结论。

\medskip\noindent
再证左边的等式。由引理 \ref{weil-lemma-degrees-cycles}，
\begin{align*}
\deg([\Delta] \cdot [\Delta])
& =
\int_{X \times X} \gamma([\Delta]) \cup \gamma([\Delta]) \\
& =
\int_{X \times X} \Delta_*1 \cup \gamma([\Delta]) \\
& =
\int_{X \times X} \Delta_*(\Delta^*\gamma([\Delta])) \\
& =
\int_X \Delta^*\gamma([\Delta])
\end{align*}
这里使用了引理 \ref{weil-lemma-push-unit} 与 \ref{weil-lemma-pushforward}。
如引理 \ref{weil-lemma-class-diagonal}，写
$\gamma([\Delta]) = \sum  e_{i, j} \otimes e'_{2d - i , j}$。
回忆，$\Delta^*$ 由杯积给出（注记 \ref{weil-remark-replace-cup-product}），于是
$$
\int_X \sum\nolimits_{i, j} e_{i, j} \cup e'_{2d - i, j} =
\sum\nolimits_{i, j} \int_X e_{i, j} \cup e'_{2d - i, j} =
\sum\nolimits_{i, j} (-1)^i = \sum (-1)^i\beta_i
$$
如所需。
\end{proof}

\begin{lemma}
\label{weil-lemma-algebra-relations}
设 $F$ 是特征为 $0$ 的域。设 $F'$ 以及 $F_i$、$i = 1, \ldots, r$ 是有限可分
$F$-代数，$A$ 是有限 $F$-代数。设
$\sigma, \sigma' : A \to F'$ 与 $\sigma_i : A \to F_i$ 是 $F$-代数映射，
且 $\sigma$、$\sigma'$ 均为满射。若存在关系
$$
\text{Tr}_{F'/F} \circ \sigma - \text{Tr}_{F'/F} \circ \sigma' =
n(\sum m_i \text{Tr}_{F_i/F} \circ \sigma_i)
$$
其中 $n > 1$，而 $m_i$ 为整数，则 $\sigma = \sigma'$。
\end{lemma}

\begin{proof}
可以把 $A = \prod A_j$ 写成局部 Artin $F$-代数
$(A_j, \mathfrak m_j, \kappa_j)$ 的有限积；见《代数》引理
\ref{algebra-lemma-finite-dimensional-algebra} 与命题
\ref{algebra-proposition-dimension-zero-ring}。以 $A' = \prod \kappa_j$ 表示对所有满足
$\kappa_j/k$ 可分的 $j$ 所取之积。于是映射 $\sigma, \sigma', \sigma_i$
全都经由 $A \to A'$ 分解。用 $A'$ 代替 $A$ 后，可假设 $A$ 是有限可分
$F$-代数。

\medskip\noindent
选取代数闭包 $\overline{F}$。置
$\overline{A} = A \otimes_F \overline{F}$、
$\overline{F}' = F' \otimes_F \overline{F}$ 以及
$\overline{F}_i = F_i \otimes_F \overline{F}$。
把 $\sigma$、$\sigma'$、$\sigma_i$ 作基变换，得到 $\overline{F}$-代数映射
$\overline{A} \to \overline{F}'$ 与 $\overline{A} \to \overline{F}_i$。
此外，$\text{Tr}_{\overline{F}'/\overline{F}}$ 是
$\text{Tr}_{F'/F}$ 的基变换，$\text{Tr}_{F_i/F}$ 亦然。因此可用
$\overline{F}$ 代替 $F$，归结到下一段讨论的情形。

\medskip\noindent
假设 $F$ 代数闭，且 $A$ 是有限可分 $F$-代数。则 $A$、$F'$、$F_i$
各自都是若干个 $F$ 的直积。称 $F$ 的若干副本之积
$F \times \ldots \times F$ 中的元素 $e$ 为最小幂等元，如果它生成其中一个因子，
即 $e = (0, \ldots, 0, 1, 0, \ldots, 0)$。设 $e \in A$ 是最小幂等元。
由于 $\sigma$ 与 $\sigma'$ 为满射，$\sigma(e)$ 与 $\sigma'(e)$ 各自要么是最小幂等元，
要么为零。若 $\sigma \not = \sigma'$，则可选取最小幂等元 $e \in A$，使
$\sigma(e) = 0$ 且 $\sigma'(e) \not = 0$，或反之。于是
$\text{Tr}_{F'/F}(\sigma(e)) = 0$ 且
$\text{Tr}_{F'/F}(\sigma'(e)) = 1$，或反之。
另一方面，$\sigma_i(e)$ 是幂等元，故
$\text{Tr}_{F_i/F}(\sigma_i(e)) = r_i$ 是整数。于是
$$
-1 = \sum n m_i r_i = n (\sum m_i r_i)
\quad\text{或}\quad
1 = \sum n m_i r_i = n (\sum m_i r_i)
$$
这不可能。
\end{proof}

\begin{lemma}
\label{weil-lemma-relations-classes-points}
设给定满足 (A)、(B)、(C) 的数据 (D0)、(D1)、(D2)、(D3)。设
$k'/k$ 是有限可分扩张，$X$ 是 $k'$ 上光滑射影概形，
$x, x' \in X$ 是 $k'$-有理点。若 $\gamma(x) \not = \gamma(x')$，则在
$\CH_0(X)$ 中，$[x] - [x']$ 不能被任何整数 $n > 1$ 整除。
\end{lemma}

\begin{proof}
若 $x$ 与 $x'$ 位于 $X$ 的不同不可约分支上，结论显然成立。因此可假设
$X$ 不可约，维数为 $d$。反设 $[x] - [x']$ 在 $\CH_0(X)$ 中能被
$n > 1$ 整除。由引理 \ref{weil-lemma-generated-by-separable}，可在
$\CH_0(X)$ 中写成 $[x] - [x'] = n(\sum m_i [x_i])$，其中
$x_i \in X$ 是剩余域在 $k$ 上可分的闭点。于是，在 $H^{2d}(X)(d)$ 中有
$$
\gamma([x]) - \gamma([x']) = n (\sum m_i \gamma([x_i]))
$$
以 $i^*, (i')^*, i_i^*$ 分别表示拉回映射
$H^0(X) \to H^0(x)$、$H^0(X) \to H^0(x')$、$H^0(X) \to H^0(x_i)$。
回忆，$H^0(x)$ 是有限可分 $F$-代数，而
$\int_x : H^0(x) \to F$ 是迹映射（引理 \ref{weil-lemma-dim-0}），记为
$\text{Tr}_x$。对 $x'$ 与 $x_i$ 亦然。于是，由公理 (A)(b) 形式的
Poincar\'e 对偶，上述等式对偶于
$$
\text{Tr}_x \circ i^* - \text{Tr}_{x'} \circ (i')^* =
n(\sum m_i \text{Tr}_{x_i} \circ i_i^*)
$$
此等式位于 $\Hom_F(H^0(X), F)$ 中。最后，由于 $x$ 与 $x'$ 是
$k'$-有理点，$i^*$ 与 $(i')^*$ 均为满射；事实上，复合
$H^0(\Spec(k')) \to H^0(X) \to H^0(x)$ 及
$H^0(\Spec(k')) \to H^0(X) \to H^0(x')$ 都是同构。
由引理 \ref{weil-lemma-algebra-relations}，得到 $i^* = (i')^*$，
这与 $\gamma([x]) \not = \gamma([x'])$ 的假设矛盾。
\end{proof}

\begin{lemma}
\label{weil-lemma-classes-points}
设给定满足 (A)、(B)、(C) 的数据 (D0)、(D1)、(D2)、(D3)。设
$k'/k$ 是有限可分扩张，$X$ 是 $k'$ 上维数为 $d$ 的几何不可约光滑射影概形。
则 $\gamma : \CH_0(X) \to H^{2d}(X)(d)$ 经由
$\deg : \CH_0(X) \to \mathbf{Z}$ 分解。
\end{lemma}

\begin{proof}
由引理 \ref{weil-lemma-generated-by-separable}，只需证明：若闭点
$x, x' \in X$ 的剩余域在 $k$ 上可分，则
$\deg(x') \gamma([x]) = \deg(x) \gamma([x'])$。

\medskip\noindent
先归结到 $k'$-有理点的情形。取 Galois 扩张 $k''/k'$，使
$\kappa(x)$ 与 $\kappa(x')$ 均可在 $k$ 上嵌入 $k''$。置
$Y = X \times_{\Spec(k')} \Spec(k'')$，并以 $p : Y \to X$ 表示投影。
由 $k''/k'$ 的选取，$Y$ 上存在分别映到 $x$、\ $x'$ 的 $k''$-有理点
$y$、\ $y'$。在 $\CH_0(X)$ 中有
$p_*[y] = [k'' : \kappa(x)][x]$ 以及
$p_*[y'] = [k'' : \kappa(x')][x']$。由公理 (C)(b) 给出的推前相容性，
只需在 $\CH^{2d}(Y)(d)$ 中证明 $\gamma([y]) = \gamma([y'])$。
这就归结到下一段的讨论。

\medskip\noindent
假设 $x$ 与 $x'$ 是 $k'$-有理点。由引理
\ref{weil-lemma-divide-difference-points}，存在有限可分域扩张 $k''/k'$，
使差 $[x] - [x']$ 的拉回 $[y] - [y']$ 在
$Y = X \times_{\Spec(k')} \Spec(k'')$ 上能被某个整数 $n > 1$ 整除。
（注意，$y, y' \in Y$ 是 $k''$-有理点。）由引理
\ref{weil-lemma-relations-classes-points}，在 $H^{2d}(Y)(d)$ 中有
$\gamma([y]) = \gamma([y'])$。再由公理 (C)(b) 中与推前的相容性，
对 $x$ 与 $x'$ 得到同一结论。
\end{proof}

\begin{lemma}
\label{weil-lemma-injective}
设给定满足 (A)、(B)、(C) 的数据 (D0)、(D1)、(D2)、(D3)。设
$f : X \to Y$ 是 $k$ 上不可约光滑射影概形之间的支配态射。则
$H^*(Y) \to H^*(X)$ 是单射。
\end{lemma}

\begin{proof}
存在整闭子概形 $Z \subset X$，其维数与 $Y$ 相同，且映满 $Y$。
所以对某个 $m > 0$ 有 $f_*[Z] = m[Y]$。由引理 \ref{weil-lemma-unit}，
在 $H^*(Y)$ 中有 $f_* \gamma([Z]) = m \gamma([Y]) = m$。
因此，由投影公式（引理 \ref{weil-lemma-pushforward}），
$f_*(f^*a \cup \gamma([Z])) = m a$，结论随即成立。
\end{proof}

\begin{lemma}
\label{weil-lemma-otimes}
设给定满足 (A)、(B)、(C) 的数据 (D0)、(D1)、(D2)、(D3)。设
$k''/k'/k$ 是有限可分代数，$X$ 是 $k'$ 上光滑射影概形。则
$$
H^*(X) \otimes_{H^0(\Spec(k'))} H^0(\Spec(k'')) =
H^*(X \times_{\Spec(k')} \Spec(k''))
$$
\end{lemma}

\begin{proof}
下文直接使用引理 \ref{weil-lemma-dim-0} 的结果而不再说明。对某个有限可分
$k'$-代数 $l$，写成
$$
k' \otimes_k k'' = k'' \times l
$$
并记 $F' = H^0(\Spec(k'))$、$F'' = H^0(\Spec(k''))$、
$G = H^0(\Spec(l))$。由于
$\Spec(k') \times \Spec(k'') = \Spec(k'') \amalg \Spec(l)$，
由公理 (B)(a) 与引理 \ref{weil-lemma-weil-additive} 得到
$$
F' \otimes_F F'' = F'' \times G
$$
从左到右的映射把 $F''$ 认同为 $F' \otimes_{F'} F''$。同理，作为
$F' \otimes_F F'' = F'' \times G$ 上的模，有
$$
H^*(X) \otimes_F F'' = H^*(X \times_{\Spec(k')} \Spec(k''))
\times H^*(X \times_{\Spec(k')} \Spec(l))
$$
引理得证。
\end{proof}






















\section{Weil 上同调理论（二）}
\label{weil-section-old}

\noindent
对我们而言，当基域 $k$ 不代数闭时，Weil 上同调理论应是经典 Weil 上同调理论
（第 \ref{weil-section-axioms-classical} 节）的对应形式。在第 \ref{weil-section-axioms} 节中，
我们列出了保证上同调理论来自 $k$ 上动机范畴的一个对称幺半函子的公理。
迄今为止，公理中尚缺少条件：当 $i < 0$ 时 $H^i(X) = 0$；还缺少关于
$H^{2d}(X)(d)$ 的一个条件，其中 $X$ 等维且维数为 $d$，该条件对应于经典公理
(A)(c) 与 (A)(d)。先来说明为何必须施加这些条件。

\begin{example}
\label{weil-example-weird-weil}
设 $k = \mathbf{C}$ 与 $F = \mathbf{C}$ 均为复数域。对 $k$ 上光滑射影概形
$X$，记 $H^{p, q}(X) = H^q(X, \Omega^p_{X/k})$。设 $(H')^*$ 是把
$X$ 送到带通常杯积的 $(H')^*(X) = \bigoplus H^{p, q}(X)$ 的函子。
这是一个经典 Weil 上同调理论（此处待插入将来的引用）。由命题
\ref{weil-proposition-weil-cohomology-theory-classical}，得到从 $M_k$ 到分次
$F$-向量空间范畴的 $\mathbf{Q}$-线性对称幺半函子 $G'$。当然，在此情形下，
对 $M_k$ 中每个 $M$，$G'(M)$ 自然带双分次，即
$$
(G')(M) = \bigoplus (G')^{p, q}(M),\quad
(G')^n = \bigoplus\nolimits_{n = p + q} (G')^{p, q}(M)
$$
其中如式中所示，$(G')^{p, q}$ 位于总次数 $p + q$。现在置
$$
G^n(M) = \bigoplus\nolimits_{n = 3p - q} (G')^{p, q}(M)
$$
以构造到分次 $F$-向量空间范畴的 $\mathbf{Q}$-线性对称幺半函子 $G$。
略去验证这确实定义对称幺半函子；其中一个技术要点是，由于上面选取了奇数
$3$ 与 $-1$，函子 $G$ 与交换约束相容。注意，
$G(\mathbf{1}(1))$ 仍集中在次数 $-2$。因此，由引理
\ref{weil-lemma-from-functor-to-weil-classical}，得到函子 $H^*$、周期类 $\gamma$ 与迹映射，
它们满足经典公理 (A)、(B)、(C)，可能例外的只有经典公理 (A)(a) 与 (A)(d)。
然而，若 $E$ 是 $k$ 上椭圆曲线，则 $\dim H^{-1}(E) = 1$；
也就是说，公理 (A)(a) 的确不成立。
\end{example}

\begin{lemma}
\label{weil-lemma-H-0-separable}
设给定满足 (A)、(B)、(C) 的数据 (D0)、(D1)、(D2)、(D3)。设 $X$ 是
$k$ 上光滑射影概形，并置 $k' = \Gamma(X, \mathcal{O}_X)$。下列条件等价：
\begin{enumerate}
\item 存在有限多个闭点 $x_1, \ldots, x_r \in X$，其剩余域在 $k$ 上可分，
使 $H^0(X) \to H^0(x_1) \oplus \ldots \oplus H^0(x_r)$ 为单射；
\item 映射 $H^0(\Spec(k')) \to H^0(X)$ 是同构。
\end{enumerate}
若这些条件成立，则 $H^0(X)$ 是 $F$ 上有限可分代数。若 $X$ 等维且维数为
$d$，则 (1)、(2) 还等价于：
\begin{enumerate}
\item[(3)] 闭点的类作为 $H^0(X)$-模生成 $H^{2d}(X)(d)$。
\end{enumerate}
\end{lemma}

\begin{proof}
此断言确有意义，因为 $k'$ 是 $k$ 上有限可分代数（《簇》引理
\ref{varieties-lemma-proper-geometrically-reduced-global-sections}），
因而 $\Spec(k')$ 在 $k$ 上光滑且射影。由 $H^*$ 与直和的相容性
（引理 \ref{weil-lemma-weil-additive} 与 \ref{weil-lemma-trace-disjoint-union}），
只需在 $X$ 连通时证明引理。因此可假设 $X$ 不可约，并须证明 (1)、(2)、(3) 等价。
置 $d = \dim(X)$。这意味着 $k'$ 是 $k$ 上有限可分域，而 $X$ 在 $k'$ 上几何不可约；
见《簇》引理 \ref{varieties-lemma-proper-geometrically-reduced-global-sections} 与
\ref{varieties-lemma-baby-stein}。

\medskip\noindent
由引理 \ref{weil-lemma-generated-by-separable}，可假设 (3) 中闭点的剩余域在 $k$ 上可分。
由公理 (A)(a) 与 (A)(b)，条件 (1) 与 (3) 等价。

\medskip\noindent
若 (2) 成立，任取闭点 $x \in X$，使其剩余域在 $k'$ 上有限可分。则例如由引理
\ref{weil-lemma-injective}，$H^0(\Spec(k')) = H^0(X) \to H^0(x)$ 为单射。

\medskip\noindent
假设等价条件 (1)、(3) 成立。按 (1) 选取 $x_1, \ldots, x_r \in X$，
再选取有限可分扩张 $k''/k'$。由引理 \ref{weil-lemma-otimes}，
$$
H^0(X) \otimes_{H^0(\Spec(k'))} H^0(\Spec(k'')) =
H^0(X \times_{\Spec(k')} \Spec(k''))
$$
因此，为证明 $H^0(\Spec(k')) \to H^0(X)$ 是同构，可以用 $k''$ 代替 $k'$。
于是可假设 $x_1, \ldots, x_r$ 都是 $k'$-有理点（此步骤把每个 $x_i$ 替换为多个点，
所以 $r$ 会增大）。由引理 \ref{weil-lemma-classes-points}，
$\gamma(x_1) = \gamma(x_2) = \ldots = \gamma(x_r)$。
由公理 (A)(b)，所有映射 $H^0(X) \to H^0(x_i)$ 均相同。这就证明了 (2)。

\medskip\noindent
最后，由引理 \ref{weil-lemma-dim-0}，若 (1) 成立，则 $H^0(X)$ 是可分 $F$-代数。
\end{proof}

\begin{lemma}
\label{weil-lemma-negative-cohomology}
设给定满足 (A)、(B)、(C) 的数据 (D0)、(D1)、(D2)、(D3)。若存在
$k$ 上光滑射影概形 $Y$，使某个 $i < 0$ 的 $H^i(Y)$ 非零，则存在
$k$ 上等维光滑射影概形 $X$，使引理 \ref{weil-lemma-H-0-separable} 的等价条件对
$X$ 不成立。
\end{lemma}

\begin{proof}
由引理 \ref{weil-lemma-weil-additive}，可假设 $Y$ 不可约，因而更是等维的。
若 $i$ 为奇数，则用 $Y \times Y$ 代替 $Y$ 后，得到一个 $Y$ 等维且
$i = -2l$ 的例子，其中 $l > 0$。置 $X = Y \times (\mathbf{P}^1_k)^l$。
利用公理 (B)(a)，得到
$$
H^0(X) \supset H^0(Y) \oplus
H^i(Y) \otimes_F H^2(\mathbf{P}^1_k)^{\otimes_F l}
$$
两个直和项都非零。因此，$H^0(X)$ 显然不可能同构于
$\Gamma(X, \mathcal{O}_X) = \Gamma(Y, \mathcal{O}_Y)$ 的谱的 $H^0$，
因为后者落在第一个直和项中。
\end{proof}

\noindent
因此，现在作出如下定义是合理的。

\begin{definition}
\label{weil-definition-weil-cohomology-theory}
设 $k$ 是域，$F$ 是特征为 $0$ 的域。以 $F$ 为系数的 $k$ 上
{\it Weil 上同调理论}由满足 Poincar\'e 对偶、K\"unneth 公式以及与周期类相容性的
数据 (D0)、(D1)、(D2)、(D3) 给出；更准确地说，它满足第
\ref{weil-section-axioms} 节的公理 (A)、(B)、(C)，并且对每个 $k$ 上光滑射影概形
$X$，引理 \ref{weil-lemma-H-0-separable} 的等价条件 (1)、(2) 都成立。
\end{definition}

\noindent
由引理 \ref{weil-lemma-negative-cohomology}，这还意味着不存在非零的负次数上同调群。
特别地，若 $k$ 代数闭，则上述 Weil 上同调理论连同一个同构
$F \to F(1)$，等同于一个经典 Weil 上同调理论。

\begin{remark}
\label{weil-remark-betti-numbers-in-some-sense}
设 $H^*$ 是 Weil 上同调理论（定义 \ref{weil-definition-weil-cohomology-theory}）。
设 $k'/k$ 是有限可分域扩张，$X$ 是 $k'$ 上维数为 $d$ 的几何不可约光滑射影概形。
假设对某些域 $F_i$ 有
$$
H^0(\Spec(k')) = F_1 \times \ldots \times F_r
$$
相应地，可以写成
$$
H^*(X) = \prod\nolimits_{i = 1, \ldots, r}
H^*(X) \otimes_{H^0(\Spec(k'))} F_i
$$
定义 \ref{weil-definition-weil-cohomology-theory} 中最后一个假设表明，
$H^0(X)$ 是 $\prod F_i$ 上秩为 $1$ 的自由模。换言之，每个因子
$H^0(X) \otimes_{H^0(\Spec(k'))} F_i$ 在 $F_i$ 上维数为 $1$。
Poincar\'e 对偶于是表明，次数 $2d$ 的上同调也具有相同性质。
然而，尚不清楚其他次数上是否也成立。具体而言，给定
$0 < n < \dim(X)$，我们不知道整数
$$
\dim_{F_i} H^n(X) \otimes_{H^0(\Spec(k'))} F_i
$$
是否与 $i$ 无关。这个问题与下述公开问题密切相关：给定代数闭基域
$\overline{k}$、特征为零的域 $F$、$\overline{k}$ 上系数域为 $F$ 的经典 Weil
上同调理论 $H^*$，以及 $\overline{k}$ 上光滑射影簇 $X$，$X$ 的 Betti 数
$$
\beta_i = \dim_F H^i(X)
$$
是否与 $F$ 及 Weil 上同调理论 $H^*$ 无关？
\end{remark}

\begin{proposition}
\label{weil-proposition-weil-cohomology-theory-again}
设 $k$ 是域，$F$ 是特征为 $0$ 的域。Weil 上同调理论等同于一个
$\mathbf{Q}$-线性对称幺半函子
$$
G : M_k \longrightarrow \text{graded }F\text{-向量空间}
$$
并满足：
\begin{enumerate}
\item $G(\mathbf{1}(1))$ 仅在次数 $-2$ 上非零；以及
\item 对每个 $k$ 上光滑射影概形 $X$，若
$k' = \Gamma(X, \mathcal{O}_X)$，则分次 $F$-向量空间的同态
$G(h(\Spec(k'))) \to G(h(X))$ 在次数 $0$ 上是同构。
\end{enumerate}
\end{proposition}

\begin{proof}
这立即由命题 \ref{weil-proposition-weil-cohomology-theory} 与定义
\ref{weil-definition-weil-cohomology-theory} 得到。当然，也可把 (2) 换成如下条件：
存在剩余域在 $k$ 上可分的闭点 $x_1, \ldots, x_r \in X$，使
$G(h(X)) \to \bigoplus G(h(x_i))$ 在次数 $0$ 上为单射。
\end{proof}







\section{Chern 类}
\label{weil-section-chern}

\noindent
本节讨论如何从第一 Chern 类与射影空间丛公式得到全部 Chern 类。
本节可参见 \cite{Grothendieck-chern}，不过我们的公理略有不同。

\medskip\noindent
设 $\mathcal{C}$ 是满足下列性质的概形范畴：
\begin{enumerate}
\item 每个 $X \in \Ob(\mathcal{C})$ 都拟紧且拟分离。
\item 若 $X \in \Ob(\mathcal{C})$，而 $U \subset X$ 开且闭，则
$U \to X$ 是 $\mathcal{C}$ 的态射。若 $\mathcal{C}$ 的态射
$X' \to X$ 经由 $U$ 分解，则 $X' \to U$ 是 $\mathcal{C}$ 的态射。
\item 若 $X \in \Ob(\mathcal{C})$，而 $\mathcal{E}$ 是有限局部自由
$\mathcal{O}_X$-模，则：
\begin{enumerate}
\item $p : \mathbf{P}(\mathcal{E}) \to X$ 是 $\mathcal{C}$ 的态射；
\item 对 $\mathcal{C}$ 中的态射 $f : X' \to X$，诱导态射
$\mathbf{P}(f^*\mathcal{E}) \to \mathbf{P}(\mathcal{E})$ 是 $\mathcal{C}$ 的态射；
\item 若 $\mathcal{E} \to \mathcal{F}$ 是到另一个有限局部自由
$\mathcal{O}_X$-模的满射，则闭浸入
$\mathbf{P}(\mathcal{F}) \to \mathbf{P}(\mathcal{E})$ 是 $\mathcal{C}$ 的态射。
\end{enumerate}
\end{enumerate}
再设给定从范畴 $\mathcal{C}$ 到分次代数范畴的反变函子 $A$。
这里，分次代数 $A$ 是带幺、结合、不必交换的 $\mathbf{Z}$-代数 $A$，并带有分次
$A = \bigoplus_{i \geq 0} A^i$。给定 $\mathcal{C}$ 的态射
$f : X' \to X$，以 $f^* : A(X) \to A(X')$ 表示诱导的代数映射。
以 $a \cup b$ 表示 $a, b \in A(X)$ 的乘积。最后，假设对 $\mathcal{C}$ 的每个对象
$X$，给定可加映射
$$
c_1^A : \Pic(X) \longrightarrow A^1(X)
$$
并假设满足下列公理：
\begin{enumerate}
\item 给定 $X \in \Ob(\mathcal{C})$ 与 $\mathcal{L} \in \Pic(X)$，元素
$c_1^A(\mathcal{L})$ 位于代数 $A(X)$ 的中心。
\item 若 $X \in \Ob(\mathcal{C})$ 且 $X = U \amalg V$，其中 $U$、$V$ 开且闭，
则经诱导映射 $A(X) \to A(U)$ 与 $A(X) \to A(V)$ 有
$A(X) = A(U) \times A(V)$。
\item 若 $f : X' \to X$ 是 $\mathcal{C}$ 的态射，而 $\mathcal{L}$ 是可逆
$\mathcal{O}_X$-模，则 $f^*c_1^A(\mathcal{L}) =
c_1^A(f^*\mathcal{L})$。
\item 给定 $X \in \Ob(\mathcal{C})$ 及常秩为 $r$ 的局部自由
$\mathcal{O}_X$-模 $\mathcal{E}$，考虑 $\mathcal{C}$ 的态射
$p : P = \mathbf{P}(\mathcal{E}) \to X$。则映射
$$
\bigoplus\nolimits_{i = 0, \ldots, r - 1} A(X)
\longrightarrow A(P),\quad
(a_0, \ldots, a_{r - 1}) \longmapsto
\sum c_1^A(\mathcal{O}_P(1))^i \cup p^*(a_i)
$$
是双射。
\item 设 $X \in \Ob(\mathcal{C})$，而 $\mathcal{E} \to \mathcal{F}$ 是秩分别为
$r + 1$、$r$ 的有限局部自由 $\mathcal{O}_X$-模之间的满射。以
$i : P' = \mathbf{P}(\mathcal{F}) \to \mathbf{P}(\mathcal{E}) = P$
表示相应的包含态射。它是 $\mathcal{C}$ 的态射，并把 $P'$ 实现为 $P$ 上有效
Cartier 除子。若 $a \in A(P)$ 且 $i^*a = 0$，则
$a \cup c_1^A(\mathcal{O}_P(P')) = 0$。
\end{enumerate}
为叙述结果，回忆 $\textit{Vect}(X)$ 表示有限局部自由 $\mathcal{O}_X$-模的
（正合）范畴。在《概形的导出范畴》第 \ref{perfect-section-K-groups} 节中，
我们定义了该范畴的第零 $K$-群 $K_0(\textit{Vect}(X))$。此外，已经看到
$K_0(\textit{Vect}(X))$ 是环；见《概形的导出范畴》注记
\ref{perfect-remark-K-ring}。

\begin{proposition}
\label{weil-proposition-chern-class}
在上述情形中，存在唯一规则，对每个 $X \in \Ob(\mathcal{C})$ 指定一个“全 Chern 类”
$$
c^A : K_0(\textit{Vect}(X)) \longrightarrow  \prod\nolimits_{i \geq 0} A^i(X)
$$
并满足下列性质：
\begin{enumerate}
\item 对 $X \in \Ob(\mathcal{C})$，有
$c^A(\alpha + \beta) = c^A(\alpha) c^A(\beta)$ 及 $c^A(0) = 1$。
\item 若 $f : X' \to X$ 是 $\mathcal{C}$ 的态射，则
$f^* \circ c^A =  c^A \circ f^*$。
\item 给定 $X \in \Ob(\mathcal{C})$ 与 $\mathcal{L} \in \Pic(X)$，有
$c^A([\mathcal{L}]) = 1 + c_1^A(\mathcal{L})$。
\end{enumerate}
\end{proposition}

\begin{proof}
设 $X \in \Ob(\mathcal{C})$，而 $\mathcal{E}$ 是有限局部自由
$\mathcal{O}_X$-模。先说明如何定义元素 $c^A(\mathcal{E}) \in A(X)$。

\medskip\noindent
第一步，设 $X = \bigcup X_r$ 是开闭子概形的分解，使
$\mathcal{E}|_{X_r}$ 常秩为 $r$。由于 $X$ 拟紧，此分解有限，故
$A(X) = \prod A(X_r)$。所以，只需在 $\mathcal{E}$ 常秩为 $r$ 时定义
$c^A(\mathcal{E})$。此时，设 $p : P \to X$ 是 $\mathcal{E}$ 的射影丛。
可以唯一地定义元素 $c_i^A(\mathcal{E}) \in A^i(X)$（$i \geq 0$），
使 $c_0^A(\mathcal{E}) = 1$，且等式
\begin{equation}
\label{weil-equation-chern-classes}
\sum\nolimits_{i = 0}^r
(-1)^i c_1(\mathcal{O}_P(1))^i \cup p^*c^A_{r - i}(\mathcal{E})
= 0
\end{equation}
成立。照例，在 $A(X)$ 中置
$c^A(\mathcal{E}) = c_0^A(\mathcal{E}) + c_1^A(\mathcal{E}) + \ldots
+ c_r^A(\mathcal{E})$。

\medskip\noindent
若 $\mathcal{E}$ 可逆，则 $c^A(\mathcal{E}) = 1 + c_1^A(\mathcal{L})$。
这立即由上述构造得到。

\medskip\noindent
元素 $c_i^A(\mathcal{E})$ 位于 $A(X)$ 的中心。为证明这一点，可假设
$\mathcal{E}$ 常秩为 $r$。设 $p : P \to X$ 是相应的射影丛。若
$a \in A(X)$，则
$p^*a \cup (-1)^r c_1(\mathcal{O}_P(1))^r =
(-1)^r c_1(\mathcal{O}_P(1))^r \cup p^*a$；因此，定义
$c_i^A(\mathcal{E})$ 的表达式中其他各项也必具有相同性质，结论得证。

\medskip\noindent
若 $f : X' \to X$ 是 $\mathcal{C}$ 的态射，则
$f^*c_i^A(\mathcal{E}) = c_i^A(f^*\mathcal{E})$。为证明这一点，可假设
$\mathcal{E}$ 常秩为 $r$。设 $p : P \to X$ 与 $p' : P' \to X'$ 分别为
$\mathcal{E}$ 与 $f^*\mathcal{E}$ 对应的射影丛。诱导态射
$g : P' \to P$ 是 $\mathcal{C}$ 的态射。把定义 $c_i^A(\mathcal{E})$ 的等式沿
$g$ 拉回，便得到 $f^*\mathcal{E}$ 的相应等式，结论随即成立。

\medskip\noindent
设 $X \in \Ob(\mathcal{C})$。考虑有限局部自由 $\mathcal{O}_X$-模的短正合列
$$
0 \to \mathcal{L} \to \mathcal{E} \to \mathcal{F} \to 0
$$
其中 $\mathcal{L}$ 可逆。则
$$
c^A(\mathcal{E}) = c^A(\mathcal{L}) c^A(\mathcal{F})
$$
事实上，由 $c^A_i$ 的构造，可假设 $\mathcal{E}$ 常秩为 $r + 1$，
$\mathcal{F}$ 常秩为 $r$。包含
$$
i : P' = \mathbf{P}(\mathcal{F}) \longrightarrow \mathbf{P}(\mathcal{E}) = P
$$
是 $\mathcal{C}$ 的态射，并且是可逆模
$\mathcal{L}^{\otimes -1} \otimes \mathcal{O}_P(1)$ 的一个正则截面之零概形。
按定义，元素
$$
\sum\nolimits_{i = 0}^r (-1)^i c_1^A(\mathcal{O}_P(1))^i \cup
p^*c^A_i(\mathcal{F})
$$
在 $P'$ 上的拉回为零。因此，由关于上同调 $A$ 的假设 (5)，在 $A^*(P)$ 中有
$$
\left(c_1^A(\mathcal{O}_P(1)) - c_1^A(\mathcal{L})\right) \cup
\left(\sum\nolimits_{i = 0}^r (-1)^i c_1^A(\mathcal{O}_P(1))^i \cup
p^*c^A_i(\mathcal{F})\right) = 0
$$
由 $c_1^A(\mathcal{E})$ 的定义，这就给出所需等式。

\medskip\noindent
设 $X \in \Ob(\mathcal{C})$。考虑有限局部自由 $\mathcal{O}_X$-模的短正合列
$$
0 \to \mathcal{E} \to \mathcal{F} \to \mathcal{G} \to 0
$$
则
$$
c^A(\mathcal{F}) = c^A(\mathcal{E}) c^A(\mathcal{G})
$$
事实上，由 $c^A_i$ 的构造，可假设 $\mathcal{E}$、$\mathcal{F}$、
$\mathcal{G}$ 分别常秩为 $r$、$s$、$t$。对 $r$ 作归纳。
$r = 1$ 的情形已在上面处理。若 $r > 1$，只需在沿态射
$\mathbf{P}(\mathcal{E}^\vee) \to X$ 拉回后验证。因此可假设存在可逆子模
$\mathcal{L} \subset \mathcal{E}$，使
$\mathcal{E}' = \mathcal{E}/\mathcal{L}$ 与
$\mathcal{F}' = \mathcal{E}/\mathcal{L}$ 都有限局部自由
（秩分别为 $s - 1$ 与 $t - 1$）。于是
$$
c^A(\mathcal{E}) = c^A(\mathcal{L}) c^A(\mathcal{E}')
\quad\text{且}\quad
c^A(\mathcal{F}) = c^A(\mathcal{L}) c^A(\mathcal{F}')
$$
又有短正合列
$$
0 \to \mathcal{E}' \to \mathcal{F}' \to \mathcal{G} \to 0
$$
由归纳假设，
$$
c^A(\mathcal{F}') = c^A(\mathcal{E}') c^A(\mathcal{G})
$$
故结论由形式计算得到。

\medskip\noindent
至此，对 $X \in \Ob(\mathcal{C})$，可以定义
$c^A : K_0(\textit{Vect}(X)) \to A(X)$。具体地，把生成元
$[\mathcal{E}]$ 送到 $c^A(\mathcal{E})$，并乘法延拓。例如，
$c^A(-[\mathcal{E}]) = c^A(\mathcal{E})^{-1}$ 是
$a^A([\mathcal{E}])$ 的形式逆元。上面证明的短正合列上的乘法性保证此定义有效。

\medskip\noindent
唯一性。设 $X \in \Ob(\mathcal{C})$，而 $\mathcal{E}$ 是有限局部自由
$\mathcal{O}_X$-模。要证明引理的条件 (1)、(2)、(3) 唯一确定
$c^A([\mathcal{E}])$。可假设 $\mathcal{E}$ 常秩为 $r$；此处已经使用 (2)。
再对 $r$ 作归纳。若 $r = 1$，唯一性由 (3) 得到。若 $r > 1$，使用 (2)
拉回到射影丛 $p : P \to X$，从而可假设存在短正合列
$0 \to \mathcal{E}' \to \mathcal{E} \to \mathcal{E}'' \to 0$，
其中 $\mathcal{E}'$ 与 $\mathcal{E}''$ 的秩都较低。由归纳假设，
$c^A(\mathcal{E}')$ 与 $c^A(\mathcal{E}'')$ 都被唯一确定。
于是 $\mathcal{E}$ 的唯一性由公理 (1) 得到。
\end{proof}

\begin{lemma}
\label{weil-lemma-splitting-principle}
在上述情形中，设 $X \in \Ob(\mathcal{C})$，而 $\mathcal{E}_i$ 是有限多个秩为
$r_i$ 的局部自由 $\mathcal{O}_X$-模。则存在 $\mathcal{C}$ 中的态射
$p : P \to X$，使得：
\begin{enumerate}
\item $p^* : A(X) \to A(P)$ 为单射；
\item 每个 $p^*\mathcal{E}_i$ 都有一个滤过，其逐次商
$\mathcal{L}_{i, 1}, \ldots, \mathcal{L}_{i, r_i}$ 是可逆
$\mathcal{O}_P$-模。
\end{enumerate}
\end{lemma}

\begin{proof}
可假设对所有 $i$ 均有 $r_i \geq 1$。对 $\sum (r_i - 1)$ 作归纳。
若此整数为 $0$，则对所有 $i$，$\mathcal{E}_i$ 都可逆，取
$\pi = \text{id}_X$ 即得结论。否则，可选取满足 $r_i > 1$ 的一个 $i$，
并考虑与 $\mathcal{E}_i$ 对应的射影丛 $p : P \to X$。存在有限局部自由
$\mathcal{O}_P$-模的短正合列
$$
0 \to \mathcal{F} \to p^*\mathcal{E}_i \to \mathcal{O}_P(1) \to 0
$$
三者的秩分别为 $r_i - 1$、$r_i$、$1$。由假设，
$p^* : A(X) \to A(P)$ 为单射。把归纳假设用于有限局部自由
$\mathcal{O}_P$-模 $\mathcal{F}$ 以及所有满足 $i' \not = i$ 的
$p^*\mathcal{E}_{i'}$，得到满足引理所述性质的态射 $p' : P' \to P$。
于是复合 $p \circ p' : P' \to X$ 满足要求。
\end{proof}

\begin{lemma}
\label{weil-lemma-chern-classes-E-tensor-L}
设 $X \in \Ob(\mathcal{C})$。设 $\mathcal{E}$ 是有限局部自由
$\mathcal{O}_X$-模，$\mathcal{L}$ 是可逆 $\mathcal{O}_X$-模。则
$$
c^A_i({\mathcal E} \otimes {\mathcal L})
=
\sum\nolimits_{j = 0}^i
\binom{r - i + j}{j} c^A_{i - j}({\mathcal E}) \cup c^A_1({\mathcal L})^j
$$
\end{lemma}

\begin{proof}
由 $c^A_i$ 的构造，可假设 $\mathcal{E}$ 常秩为 $r$。设
$p : P \to X$ 与 $p' : P' \to X$ 分别为 $\mathcal{E}$ 和
$\mathcal{E} \otimes \mathcal{L}$ 对应的射影丛。存在同构
$g : P \to P'$，满足
$g^*\mathcal{O}_{P'}(1) = \mathcal{O}_P(1) \otimes p^*\mathcal{L}$；
见《构造》引理 \ref{constructions-lemma-twisting-and-proj}。因此
$$
g^*c_1^A(\mathcal{O}_{P'}(1)) =
c_1^A(\mathcal{O}_P(1)) + p^*c_1^A(\mathcal{L})
$$
所需等式由此以及用等式 (\ref{weil-equation-chern-classes}) 给出的 Chern 类定义形式地推出。
\end{proof}

\begin{proposition}
\label{weil-proposition-chern-character}
在上述情形中，假设对所有 $X \in \Ob(\mathcal{C})$，$A(X)$ 都是
$\mathbf{Q}$-代数。则存在唯一规则，对每个 $X \in \Ob(\mathcal{C})$ 指定一个
“Chern 特征”
$$
ch^A : K_0(\textit{Vect}(X)) \longrightarrow
\prod\nolimits_{i \geq 0} A^i(X)
$$
并满足：
\begin{enumerate}
\item 对所有 $X \in \Ob(\mathcal{C})$，$ch^A$ 是环映射。
\item 若 $f : X' \to X$ 是 $\mathcal{C}$ 的态射，则
$f^* \circ ch^A =  ch^A \circ f^*$。
\item 给定 $X \in \Ob(\mathcal{C})$ 与 $\mathcal{L} \in \Pic(X)$，有
$ch^A([\mathcal{L}]) = \exp(c_1^A(\mathcal{L}))$。
\end{enumerate}
\end{proposition}

\begin{proof}
设 $X \in \Ob(\mathcal{C})$，而 $\mathcal{E}$ 是有限局部自由
$\mathcal{O}_X$-模。先说明如何定义秩 $r^A(\mathcal{E}) \in A^0(X)$。
设 $X = \bigcup X_r$ 是开闭子概形的分解，使
$\mathcal{E}|_{X_r}$ 常秩为 $r$。由于 $X$ 拟紧，此分解有限，设为
$X = X_0 \amalg X_1 \amalg \ldots \amalg X_n$。于是
$A(X) = A(X_0) \times A(X_1) \times \ldots \times A(X_n)$，故可定义
$r^A(\mathcal{E}) = (0, 1, \ldots, n) \in A^0(X)$。

\medskip\noindent
设 $P_p(c_1, \ldots, c_p)$ 是《Chow 同调》例
\ref{chow-example-power-sum} 中构造的多项式。定义
$$
ch^A(\mathcal{E}) = r^A(\mathcal{E}) +
\sum\nolimits_{i \geq 1} (1/i!)
P_i(c^A_1(\mathcal{E}), \ldots, c^A_i(\mathcal{E}))
\in \prod\nolimits_{i \geq 0} A^i(X)
$$
其中 $ci^A$ 是命题 \ref{weil-proposition-chern-class} 的 Chern 类。
由此立即得到引理的性质 (2)、(3)。

\medskip\noindent
还须证明以下三个断言：
\begin{enumerate}
\item 若 $0 \to \mathcal{E}_1 \to \mathcal{E} \to \mathcal{E}_2 \to 0$
是 $X \in \Ob(\mathcal{C})$ 上有限局部自由 $\mathcal{O}_X$-模的短正合列，则
$ch^A(\mathcal{E}) = ch^A(\mathcal{E}_1) + ch^A(\mathcal{E}_2)$。
\item 若 $\mathcal{E}_1$ 与 $\mathcal{E}_2 \to 0$ 是
$X \in \Ob(\mathcal{C})$ 上有限局部自由 $\mathcal{O}_X$-模，则
$ch^A(\mathcal{E}_1 \otimes \mathcal{E}_2) =
ch^A(\mathcal{E}_1) ch^A(\mathcal{E}_2)$。
\end{enumerate}
第一个断言将证明 $ch^A$ 经由 $K_0(\textit{Vect}(X))$ 分解；
第一、第二个断言合在一起则表明 $ch^A$ 是环映射。

\medskip\noindent
为证明这些断言，可归结到 $\mathcal{E}_1$、$\mathcal{E}_2$ 分别常秩为
$r_1$、$r_2$ 的情形。此时，$A^0(X)$ 中的等式立即成立。为证明更高次数中的等式，
由引理 \ref{weil-lemma-splitting-principle}，可假设 $\mathcal{E}_1$ 与
$\mathcal{E}_2$ 有滤过，其分次部分是可逆模
$\mathcal{L}_{1, j}$、$j = 1, \ldots, r_1$ 以及
$\mathcal{L}_{2, j}$、$j = 1, \ldots, r_2$。
利用 Chern 类的乘法性，得到
$$
c_i^A(\mathcal{E}_1) =
s_i(c_1^A(\mathcal{L}_{1, 1}), \ldots, c_1^A(\mathcal{L}_{1, r_1}))
$$
其中 $s_i$ 是《Chow 同调》例 \ref{chow-example-power-sum} 中的第 $i$ 个初等对称函数。
对 $c_i^A(\mathcal{E}_2)$ 亦然。在情形 (1) 中，得到
$$
c_i^A(\mathcal{E}) =
s_i(c_1^A(\mathcal{L}_{1, 1}), \ldots, c_1^A(\mathcal{L}_{1, r_1}),
c_1^A(\mathcal{L}_{2, 1}), \ldots, c_1^A(\mathcal{L}_{2, r_2}))
$$
而在情形 (2) 中，得到
$$
c_i^A(\mathcal{E}_1 \otimes \mathcal{E}_2) =
s_i(c_1^A(\mathcal{L}_{1, 1}) + c_1^A(\mathcal{L}_{2, 1}),
\ldots, c_1^A(\mathcal{L}_{1, r_1}) + c_1^A(\mathcal{L}_{2, r_2}))
$$
由多项式 $P_i$ 的定义，这意味着
$$
P_i(c^A_1(\mathcal{E}_1), \ldots, c^A_i(\mathcal{E}_1)) =
\sum\nolimits_{j = 1, \ldots, r_1} c_1^A(\mathcal{L}_{1, j})^i
$$
对 $\mathcal{E}_2$ 亦然。在情形 (1) 中还有
$$
P_i(c^A_1(\mathcal{E}), \ldots, c^A_i(\mathcal{E})) =
\sum\nolimits_{j = 1, \ldots, r_1} c_1^A(\mathcal{L}_{1, j})^i +
\sum\nolimits_{j = 1, \ldots, r_2} c_1^A(\mathcal{L}_{2, j})^i
$$
在情形 (2) 中，相应地得到
$$
P_i(c^A_1(\mathcal{E}_1 \otimes \mathcal{E}_2), \ldots,
c^A_i(\mathcal{E}_1 \otimes \mathcal{E}_2)) =
\sum\nolimits_{j = 1, \ldots, r_1}
\sum\nolimits_{j' = 1, \ldots, r_2}
(c_1^A(\mathcal{L}_{1, j}) + c_1^A(\mathcal{L}_{2, j'}))^i
$$
所以，所需等式现在由对称多项式之间的初等恒等式推出。

\medskip\noindent
略去唯一性的证明。
\end{proof}

\begin{lemma}
\label{weil-lemma-adams-and-chern}
在上述情形中，设 $X \in \Ob(\mathcal{C})$。若 $\psi^2$ 如《Chow 同调》引理
\ref{chow-lemma-second-adams-operator}，而 $c^A$、$ch^A$ 如命题
\ref{weil-proposition-chern-class}、\ref{weil-proposition-chern-character}，则对所有
$\alpha \in K_0(\textit{Vect}(X))$ 有
$c^A_i(\psi^2(\alpha)) = 2^i c^A_i(\alpha)$ 以及
$ch^A_i(\psi^2(\alpha)) = 2^i ch^A_i(\alpha)$。
\end{lemma}

\begin{proof}
注意，映射 $\prod_{i \geq 0} A^i(X) \to \prod_{i \geq 0} A^i(X)$，
其在 $A^i(X)$ 上乘以 $2^i$，是环映射。由于 $\psi^2$ 也是环映射，
只需对 $K_0(\textit{Vect}(X))$ 的加法生成元证明这些公式。因此，可假设
$\alpha = [\mathcal{E}]$，其中 $\mathcal{E}$ 是某个有限局部自由
$\mathcal{O}_X$-模。由 $\mathcal{E}$ 的 Chern 类的构造，立即归结到
$\mathcal{E}$ 常秩为 $r$ 的情形。此时，可以选取射影光滑态射
$p : P \to X$，使限制映射 $A^*(X) \to A^*(P)$ 为单射，并使
$p^*\mathcal{E}$ 有有限滤过，其分次部分是可逆 $\mathcal{O}_P$-模
$\mathcal{L}_j$；见引理 \ref{weil-lemma-splitting-principle}。于是
$[p^*\mathcal{E}] = \sum [\mathcal{L}_j]$，从而由 $\psi^2$ 的定义，
$\psi^2([p^*\mathcal{E}]) = \sum [\mathcal{L}_j^{\otimes 2}]$。
置 $x_j  = c^A_1(\mathcal{L}_j)$，则在 $\prod A^i(P)$ 中有
$$
c^A(\alpha) = \prod (1 + x_j)
\quad\text{且}\quad
c^A(\psi^2(\alpha)) = \prod (1 + 2 x_j)
$$
并且在 $\prod A^i(P)$ 中有
$$
ch^A(\alpha) = \sum \exp(x_j)
\quad\text{且}\quad
ch^A(\psi^2(\alpha)) = \sum \exp(2 x_j)
$$
所需结论由这些公式推出。
\end{proof}







\section{外幂与 K-群}
\label{weil-section-lambda-operations}

\noindent
下面只作定义 lambda 算子所需的最低限度工作。设 $X$ 是概形。回忆，
$\textit{Vect}(X)$ 表示有限局部自由 $\mathcal{O}_X$-模的范畴。此外，在
《概形的导出范畴》第 \ref{perfect-section-K-groups} 节中，已经构造了与此范畴相关的
第零 $K$-群 $K_0(\textit{Vect}(X))$。最后，$K_0(\textit{Vect}(X))$ 是环；
见《概形的导出范畴》注记 \ref{perfect-remark-K-ring}。

\begin{lemma}
\label{weil-lemma-lambda-operations}
设 $X$ 是概形。存在映射
$$
\lambda^r : K_0(\textit{Vect}(X)) \longrightarrow K_0(\textit{Vect}(X))
$$
当 $\mathcal{E}$ 是有限局部自由 $\mathcal{O}_X$-模时，它把
$[\mathcal{E}]$ 送到 $[\wedge^r(\mathcal{E})]$，并且这些映射与拉回相容。
\end{lemma}

\begin{proof}
考虑环 $R = K_0(\textit{Vect}(X))[[t]]$，其中 $t$ 是变量。
对有限局部自由 $\mathcal{O}_X$-模 $\mathcal{E}$，在 $R$ 中置
$$
c(\mathcal{E}) = \sum\nolimits_{i = 0}^\infty [\wedge^i(\mathcal{E})] t^i
$$
断言：给定有限局部自由 $\mathcal{O}_X$-模的短正合列
$$
0 \to \mathcal{E}' \to \mathcal{E} \to \mathcal{E}'' \to 0
$$
有 $c(\mathcal{E}) = c(\mathcal{E}') c(\mathcal{E}'')$。
该断言表明，$c$ 延拓为映射
$$
c :  K_0(\textit{Vect}(X)) \longrightarrow R
$$
并把 $K_0(\textit{Vect}(X))$ 中的加法变为 $R$ 中的乘法。
写 $c(\alpha) = \sum \lambda^i(\alpha) t^i$，便得到所需算子
$\lambda^i$。

\medskip\noindent
为证明断言，把短正合列视为 $\mathcal{E}$ 上有 $2$ 步的滤过。
它在 $\wedge^r(\mathcal{E})$ 上诱导一个有 $r + 1$ 步的滤过，其子商为
$$
\wedge^r(\mathcal{E}'),
\wedge^{r - 1}(\mathcal{E}') \otimes \mathcal{E}'',
\wedge^{r - 2}(\mathcal{E}') \otimes \wedge^2(\mathcal{E}''), \ldots
\wedge^r(\mathcal{E}'')
$$
因此，$[\wedge^r(\mathcal{E})]$ 等于
$$
\sum\nolimits_{i = 0}^r
[\wedge^{r - i}(\mathcal{E}')] [\wedge^i(\mathcal{E}'')]
$$
结论容易由此及初等代数得到。
\end{proof}










\section{Weil 上同调理论，III}
\label{weil-section-c1}

\noindent
设 $k$ 是域，$F$ 是特征为零的域。假设给定下列数据：
\begin{enumerate}
\item[(D0)] 一个 $1$ 维 $F$-向量空间 $F(1)$。
\item[(D1)] 从 $k$ 上光滑射影概形范畴到分次交换 $F$-代数范畴的反变函子
$H^*(-)$。
\item[(D2')] 对每个 $k$ 上光滑射影概形 $X$，一个阿贝尔群同态
$c_1^H : \Pic(X) \to H^2(X)(1)$。
\end{enumerate}
关于 (D0) 与 (D1)，采用第 \ref{weil-section-axioms} 节所述的术语、记号与约定。
给定 $k$ 上光滑射影概形 $X$ 及可逆 $\mathcal{O}_X$-模 $\mathcal{L}$，
(D2') 中的上同调类 $c_1^H(\mathcal{L}) \in H^2(X)(1)$ 有时称为
{\it $\mathcal{L}$ 在上同调中的第一 Chern 类}。

\medskip\noindent
公理如下。
\begin{enumerate}
\item[(A1)] $H^*$ 与有限余积相容。
\item[(A2)] $c_1^H$ 与拉回相容。
\item[(A3)] 设 $X$ 是 $k$ 上光滑射影概形，$\mathcal{E}$ 是秩为
$r \geq 1$ 的局部自由 $\mathcal{O}_X$-模。考虑态射
$p : P = \mathbf{P}(\mathcal{E}) \to X$。则映射
$$
\bigoplus\nolimits_{i = 0, \ldots, r - 1} H^*(X)(-i)
\longrightarrow H^*(P),\quad
(a_0, \ldots, a_{r - 1}) \longmapsto
\sum c_1^H(\mathcal{O}_P(1))^i \cup p^*(a_i)
$$
是 $F$-向量空间同构。
\item[(A4)] 设 $i : Y \to X$ 是 $k$ 上有效 Cartier 除子的包含，且
$X$ 与 $Y$ 都是 $k$ 上光滑射影概形。若 $a \in H^*(X)$ 且
$i^*a = 0$，则 $a \cup c_1^H(\mathcal{O}_X(Y)) = 0$。
\item[(A5)] $H^*$ 与有限积相容。
\item[(A6)] 设 $X$ 是 $k$ 上非空、光滑、射影且维数为 $d$ 的等维概形。
则存在 $F$-线性映射 $\lambda : H^{2d}(X)(d) \to F$，使得
$(\text{id} \otimes \lambda) \gamma([\Delta]) =  1$ 在 $H^*(X)$ 中成立。
\item[(A7)] 若 $b : X' \to X$ 是 $k$ 上光滑射影概形 $X$ 沿光滑中心的
爆破\footnote{此时 $X'$ 也在 $k$ 上光滑且射影；见《态射进阶》引理
\ref{more-morphisms-lemma-blowup}。}，则
$b^* : H^*(X) \to H^*(X')$ 是单射。
\item[(A8)] 若 $X$ 是 $k$ 上光滑射影概形，且
$k' = \Gamma(X, \mathcal{O}_X)$，则映射 $H^0(\Spec(k')) \to H^0(X)$
是同构。
\item[(A9)] 设 $X$ 是 $k$ 上非空光滑射影、维数为 $d$ 的等维概形。
设 $i : Y \to X$ 是在 $k$ 上光滑的非空有效 Cartier 除子。对
$a \in H^{2d - 2}(X)(d - 1)$，有
$\lambda_Y(i^*(a)) = \lambda_X(a \cup c_1^H(\mathcal{O}_X(Y))$，其中
$\lambda_Y$ 与 $\lambda_X$ 分别是公理 (A6) 对 $X$ 与 $Y$ 给出的映射。
\end{enumerate}
下面更精确地说明各条公理的含义。公理 (A3)、(A4) 与 (A7) 按照其陈述已经清楚。

\medskip\noindent
关于 (A1)。这意味着 $H^*(\emptyset) = 0$，并且
$(i^*, j^*) : H^*(X \amalg Y) \to H^*(X) \times H^*(Y)$
是同构，其中 $i$ 与 $j$ 是余投射。

\medskip\noindent
关于 (A2)。这意味着：给定 $k$ 上光滑射影概形之间的态射
$f : X \to Y$ 及可逆 $\mathcal{O}_Y$-模 $\mathcal{N}$，有
$f^*c_1^H(\mathcal{L}) = c_1^H(f^*\mathcal{L})$。

\medskip\noindent
关于 (A5)。这意味着 $H^*(\Spec(k)) = F$，并且对 $k$ 上光滑射影的
$X$ 与 $Y$，映射 $H^*(X) \otimes_F H^*(Y) \to H^*(X \times Y)$，
$a \otimes b \mapsto p^*(a) \cup q^*(b)$ 是同构，其中 $p$ 与 $q$ 是投射。

\medskip\noindent
关于 (A6)。设 $X$ 是 $k$ 上非空光滑射影、维数为 $d$ 的等维概形。
若已有公理 (A1)--(A4)，则由引理 \ref{weil-lemma-cycle-classes} 可考虑对角线的类
$$
\gamma([\Delta]) \in
H^{2d}(X \times X)(d) = \bigoplus\nolimits_i H^i(X) \otimes_F H^{2d - i}(X)(d)
$$
其中张量分解来自公理 (A5)。给定 $F$-线性映射
$\lambda : H^{2d}(X)(d) \to F$，也可先投射到 $H^{2d}(X)(d)$，从而把
$\lambda$ 看作 $F$-线性映射 $\lambda : H^*(X)(d) \to F$。
如此，公理 (A6) 便有意义。

\medskip\noindent
关于 (A8)。设 $X$ 是 $k$ 上光滑射影概形。则
$k' = \Gamma(X, \mathcal{O}_X)$ 是有限可分 $k$-代数（《簇》引理
\ref{varieties-lemma-proper-geometrically-reduced-global-sections}），
因而 $\Spec(k')$ 在 $k$ 上光滑且射影。因此可以把 $H^*$ 应用于
$\Spec(k')$，公理 (A8) 也就有意义。

\medskip\noindent
关于 (A9)。由注记 \ref{weil-remark-trace} 将会看到，若公理 (A1)--(A7)
成立，则公理 (A6) 中的映射 $\lambda$ 是唯一的。

\begin{lemma}
\label{weil-lemma-chern-classes}
假设给定满足公理 (A1)、(A2)、(A3) 与 (A4) 的数据 (D0)、(D1) 与 (D2')。
则存在唯一规则，对每个 $k$ 上光滑射影概形 $X$ 指定一个环同态
$$
ch^H :
K_0(\textit{Vect}(X))
\longrightarrow
\prod\nolimits_{i \geq 0} H^{2i}(X)(i)
$$
它与拉回相容，并且对任意可逆 $\mathcal{O}_X$-模 $\mathcal{L}$ 有
$ch^H(\mathcal{L}) = \exp(c_1^H(\mathcal{L}))$。
\end{lemma}

\begin{proof}
把命题 \ref{weil-proposition-chern-character} 应用于 $k$ 上光滑射影概形范畴、
函子 $A : X \mapsto \bigoplus_{i \geq 0} H^{2i}(X)(i)$ 以及映射
$c_1^H$，立即得到结论。
\end{proof}

\begin{lemma}
\label{weil-lemma-cycle-classes}
假设给定满足公理 (A1)、(A2)、(A3) 与 (A4) 的数据 (D0)、(D1) 与 (D2')。
则存在唯一规则，对每个 $k$ 上光滑射影概形 $X$ 指定一个分次环同态
$$
\gamma : \CH^*(X) \longrightarrow \bigoplus\nolimits_{i \geq 0} H^{2i}(X)(i)
$$
它与拉回相容，并且对 $K_0(\textit{Vect}(X))$ 中的 $\alpha$ 有
$ch^H(\alpha) = \gamma(ch(\alpha))$。
\end{lemma}

\begin{proof}
回忆有同构
$$
K_0(\textit{Vect}(X)) \otimes \mathbf{Q}
\longrightarrow \CH^*(X) \otimes \mathbf{Q},\quad
\alpha \longmapsto ch(\alpha) \cap [X]
$$
见《Chow 同调》引理 \ref{chow-lemma-K-tensor-Q}。由《Chow 同调》注记
\ref{chow-remark-chern-character-K}，这是环同构。如下定义 $\gamma$：
$\gamma(\alpha) = ch^H(\alpha')$，其中 $ch^H$ 如引理
\ref{weil-lemma-chern-classes}，而 $\alpha' \in K_0(\textit{Vect}(X))$ 满足
$ch(\alpha') \cap [X] = \alpha$ 在 $\CH^*(X) \otimes \mathbf{Q}$ 中成立。

\medskip\noindent
构造 $\alpha \mapsto \gamma(\alpha)$ 与拉回相容，因为 $ch^H$ 与取 Chern 类
都与拉回相容；见引理 \ref{weil-lemma-chern-classes} 及《Chow 同调》注记
\ref{chow-remark-gysin-chern-classes}。

\medskip\noindent
还须证明 $\gamma$ 是分次的。令
$\psi^2 : K_0(\textit{Vect}(X)) \to K_0(\textit{Vect}(X))$ 为第二 Adams
算子；见《Chow 同调》引理 \ref{chow-lemma-second-adams-operator}。
若 $\alpha \in \CH^i(X)$，且
$\alpha' \in K_0(\textit{Vect}(X)) \otimes \mathbf{Q}$ 是满足
$ch(\alpha') \cap [X] = \alpha$ 的唯一元素，则在《Chow 同调》第
\ref{chow-section-intersection-regular} 节已经看到
$\psi^2(\alpha') = 2^i \alpha'$。因此由引理 \ref{weil-lemma-adams-and-chern}
可得 $ch^H(\alpha') \in H^{2i}(X)(i)$，正合所需。
\end{proof}

\begin{lemma}
\label{weil-lemma-divide-pullback-good-blowing-up}
设 $b : X' \to X$ 是 $k$ 上光滑射影概形沿光滑闭子概形
$Z \subset X$ 的爆破。考虑图表
$$
\xymatrix{
E \ar[r]_j \ar[d]_\pi & X' \ar[d]^b \\
Z \ar[r]^i & X
}
$$
假设存在 $K_0(X)$ 中的元素，其在 $Z$ 上的限制等于
$K_0(Z)$ 中 $\mathcal{C}_{Z/X}$ 的类。再假设 $Z$ 的每个不可约分支在
$X$ 中的余维都是 $r$。则存在周期 $\theta \in \CH^{r - 1}(X')$，使得
$b^![Z] = [E] \cdot \theta$ 在 $\CH^r(X')$ 中成立，并且
$\pi_*j^!(\theta) = [Z]$ 在 $\CH^r(Z)$ 中成立。
\end{lemma}

\begin{proof}
概形 $X$ 在 $k$ 上光滑且射影，故
$K_0(X) = K_0(\textit{Vect}(X))$。见《概形的导出范畴》引理
\ref{perfect-lemma-resolution-property-ample} 与
\ref{perfect-lemma-K-is-old-K}。取 $\alpha \in K_0(\text{Vect}(X))$，
使其在 $Z$ 上的限制为 $\mathcal{C}_{Z/X}$。由《Chow 同调》引理
\ref{chow-lemma-minus-adams-operator}，存在元素 $\alpha^\vee$，其限制为
$\mathcal{C}_{Z/X}^\vee$。由爆破公式（《Chow 同调》引理
\ref{chow-lemma-blow-up-formula}），有
$$
b^![Z] = b^!i_*[Z] = j_* res(b^!)([Z]) =
j_*(c_{r - 1}(\mathcal{F}^\vee) \cap \pi^*[Z]) =
j_*(c_{r - 1}(\mathcal{F}^\vee) \cap [E])
$$
其中 $\mathcal{F}$ 是满射
$\pi^*\mathcal{C}_{Z/X} \to \mathcal{C}_{E/X'}$ 的核。注意
$b^*\alpha^\vee - [\mathcal{O}_{X'}(E)]$ 是
$K_0(\text{Vect}(X'))$ 的元素，其在 $E$ 上的限制为
$[\pi^*\mathcal{C}_{Z/X}^\vee] - [\mathcal{C}_{E/X'}^\vee] =
[\mathcal{F}^\vee]$。由于以 Chern 类作帽积与 $j_*$ 可交换，上式等于
$$
c_{r - 1}(b^*\alpha^\vee - [\mathcal{O}_{X'}(E)]) \cap [E]
$$
作为 $X'$ 的 Chow 群中的元素。因此令
$$
\theta = c_{r - 1}(b^*\alpha^\vee - [\mathcal{O}_{X'}(E)]) \cap [X']
$$
便得到第一个关系 $\theta \cdot [E] = b^![Z]$；例如可用《Chow 同调》引理
\ref{chow-lemma-identify-chow-for-smooth}。对第二个关系，注意
$$
j^!\theta = j^!(c_{r - 1}(b^*\alpha^\vee - [\mathcal{O}_{X'}(E)]) \cap [X'])
= c_{r - 1}(\mathcal{F}^\vee) \cap j^![X'] =
c_{r - 1}(\mathcal{F}^\vee) \cap [E]
$$
在 $E$ 的 Chow 群中成立。要证明其 $\pi_*$ 等于 $[Z]$，只需证明
余维 $r - 1$ 的周期
$(-1)^{r - 1}c_{r - 1}(\mathcal{F}) \cap [E]$ 在 $\pi$ 的纤维上的次数为
$1$。略去这一计算。
\end{proof}

\begin{lemma}
\label{weil-lemma-A5-A6-imply}
假设给定满足公理 (A1)--(A4) 与 (A7) 的数据 (D0)、(D1) 与 (D2')。
设 $X$ 是 $k$ 上光滑射影概形，$Z \subset X$ 是光滑闭子概形，且
$Z$ 的每个不可约分支在 $X$ 中的余维都是 $r$。假设
$K_0(Z)$ 中 $\mathcal{C}_{Z/X}$ 的类是 $K_0(X)$ 中某元素的限制。
若 $a \in H^*(X)$ 且 $a|_Z = 0$ 在 $H^*(Z)$ 中成立，则
$\gamma([Z]) \cup a = 0$。
\end{lemma}

\begin{proof}
令 $b : X' \to X$ 为该爆破。由 (A7)，只需证明
$$
b^*(\gamma([Z]) \cup a) = b^*\gamma([Z]) \cup b^*a = 0
$$
由引理 \ref{weil-lemma-divide-pullback-good-blowing-up}，有
$$
b^*\gamma([Z]) = \gamma(b^*[Z]) =
\gamma([E] \cdot \theta) =
\gamma([E]) \cup \gamma(\theta)
$$
由于 $b^*a$ 在 $E$ 上限制为零，且
$\gamma([E]) = c^H_1(\mathcal{O}_{X'}(E))$，所需结论由 (A4) 得到。
\end{proof}

\begin{lemma}
\label{weil-lemma-poincare-duality}
假设给定满足公理 (A1)--(A7) 的数据 (D0)、(D1) 与 (D2')。
则第 \ref{weil-section-axioms} 节的公理 (A) 成立，其中
$\int_X = \lambda$ 如公理 (A6)。
\end{lemma}

\begin{proof}
设 $X$ 是 $k$ 上非空光滑射影、维数为 $d$ 的等维概形。将证明分次
$F$-向量空间 $H^*(X)(d)[2d]$ 是 $H^*(X)$ 的左对偶。由《同调》引理
\ref{homology-lemma-left-dual-graded-vector-spaces}，这将证明所需结论。
我们将使用公理 (A5)，它特别给出
$$
H^*(X \times X)(d) =
\bigoplus H^i(X) \otimes H^j(X)(d) =
\bigoplus H^i(X)(d) \otimes H^j(X)
$$
定义映射
$$
\eta : F \longrightarrow H^*(X \times X)(d)
$$
为乘以 $\gamma([\Delta]) \in H^{2d}(X \times X)(d)$。另一方面，定义映射
$$
\epsilon :
H^*(X \times X)(d) \longrightarrow H^*(X)(d) \xrightarrow{\lambda} F
$$
如下：先沿对角态射 $\Delta : X \to X \times X$ 作拉回 $\Delta^*$，
再应用公理 (A6) 中的 $F$-线性映射
$\lambda : H^{2d}(X)(d) \to F$，这里预先与投射
$H^*(X)(d) \to H^{2d}(X)(d)$ 复合。

为证明 $H^*(X)(d)$ 是 $H^*(X)$ 的左对偶，须证明映射
$$
\eta \otimes 1 :
H^*(X) \longrightarrow H^*(X \times X \times X)(d)
$$
与
$$
1 \otimes \epsilon : H^*(X \times X \times X)(d) \longrightarrow H^*(X)
$$
的复合是恒等映射。若 $a \in H^*(X)$，则该复合把 $a$ 映到
$$
(1 \otimes \lambda)(\Delta_{23}^*(q_{12}^*\gamma([\Delta]) \cup q_3^*a)) =
(1 \otimes \lambda)(\gamma([\Delta]) \cup p_2^*a)
$$
其中 $q_i : X \times X \times X \to X$ 与
$q_{ij} : X \times X \times X \to X \times X$ 是投射，
$\Delta_{23} : X \times X \to X \times X \times X$ 是对角态射，而
$p_i : X \times X \to X$ 是投射。等式成立，是因为
$\Delta_{23}^*(q_{12}^*\gamma([\Delta]) =
\Delta_{23}^*\gamma([\Delta \times X]) = \gamma([\Delta])$，并且
$\Delta_{23}^* q_3^*a = p_2^*a$。由于
$\gamma([\Delta]) \cup p_1^*a = \gamma([\Delta]) \cup p_2^*a$（见下文），
上式化为
$$
(1 \otimes \lambda)(\gamma([\Delta]) \cup p_1^*a) = a
$$
这正是对 $\lambda$ 的选择所保证的。《范畴》定义
\ref{categories-definition-dual} 中的第二个条件
$(\epsilon \otimes 1) \circ (1 \otimes \eta) = \text{id}$ 以完全相同的方式证明。

\medskip\noindent
注意 $p_1^*a$ 与 $\text{pr}_2^*a$ 在
$\Delta \subset X \times X$ 上限制为同一个上同调类。此外，
$\mathcal{C}_{\Delta/X \times X} = \Omega^1_\Delta$ 是
$p_1^*\Omega^1_X$ 的限制。因此引理 \ref{weil-lemma-A5-A6-imply} 给出
$\gamma([\Delta]) \cup p_1^*a = \gamma([\Delta]) \cup p_2^*a$，证明完成。
\end{proof}

\begin{remark}[迹映射的唯一性]
\label{weil-remark-trace}
假设给定满足公理 (A1)--(A7) 的数据 (D0)、(D1) 与 (D2')。
设 $X$ 是 $k$ 上非空光滑射影、维数为 $d$ 的等维概形。结合引理
\ref{weil-lemma-poincare-duality} 与《同调》引理
\ref{homology-lemma-left-dual-graded-vector-spaces} 的证明可见，
$$
\gamma([\Delta]) \in \bigoplus\nolimits_i H^i(X) \otimes H^{2d - i}(X)(d)
$$
对每个 $i$ 都在 $H^i(X)$ 与 $H^{2d - i}(X)(d)$ 之间定义完美对偶。
特别地，公理 (A6) 中的线性映射
$\int_X = \lambda : H^{2d}(X)(d) \to F$ 是唯一的！从现在起，把
$\int_X$ 称为 $X$ 的迹映射。
\end{remark}

\begin{lemma}
\label{weil-lemma-trace-product}
假设给定满足公理 (A1)--(A7) 的数据 (D0)、(D1) 与 (D2')。
则第 \ref{weil-section-axioms} 节的公理 (B) 成立。
\end{lemma}

\begin{proof}
公理 (B)(a) 由公理 (A5) 立即得到。设 $X$ 与 $Y$ 是 $k$ 上非空光滑射影概形，
其维数分别为 $d$ 与 $e$。为证明公理 (B)(b)，注意 $X \times Y$ 的对角线
$\Delta_{X \times Y}$ 是以下两个对角线之拉回的交积：$X$ 的
$\Delta_X$ 与 $Y$ 的 $\Delta_Y$，相应投射为
$p : X \times Y \times X \times Y \to X \times X$ 及
$q : X \times Y \times X \times Y \to Y \times Y$。
$\gamma$ 与交积的相容性于是给出
$$
\gamma([\Delta_{X \times Y}]) \in
H^{2d + 2e}(X \times Y \times X \times Y)(d + e)
$$
是 $\gamma([\Delta_X])$ 与 $\gamma([\Delta_Y])$ 分别沿 $p$ 与 $q$ 拉回后的杯积。
写
$$
\gamma([\Delta_{X \times Y}]) = \sum \eta_{X \times Y, i}
\text{ 其中 }
\eta_{X \times Y, i} \in
H^i(X \times Y) \otimes H^{2d + 2e - i}(X \times Y)(d + e)
$$
并类似地写 $\gamma([\Delta_X]) = \sum \eta_{X, i}$ 及
$\gamma([\Delta_Y]) = \sum \eta_{Y, i}$。上述观察表明
$$
\eta_{X \times Y, 0} =
\sum\nolimits_{i \in \mathbf{Z}} p^*\eta_{X, i} \cup q^*\eta_{Y, -i}
$$
（若该上同调理论在负次数消失——几乎所有情形都如此——则只有 $i = 0$ 项有贡献，
而 $\eta_{X \times Y, 0}$ 如预期位于
$H^0(X) \otimes H^0(Y) \otimes H^{2d}(X)(d) \otimes H^{2e}(Y)(e)$；
但这里不需要这一点。）由于 $\lambda_X : H^{2d}(X)(d) \to F$ 与
$\lambda_Y : H^{2e}(Y)(e) \to F$ 分别把 $\eta_{X, 0}$ 与 $\eta_{Y, 0}$
送到 $H^*(X)$ 与 $H^*(Y)$ 中的 $1$，可见 $\lambda_X \otimes \lambda_Y$
把 $\eta_{X \times Y, 0}$ 送到
$H^*(X) \otimes H^*(Y) = H^*(X \times Y)$ 中的 $1$，证明完成。
\end{proof}

\begin{lemma}
\label{weil-lemma-trace-base}
假设给定满足公理 (A1)--(A7) 的数据 (D0)、(D1) 与 (D2')。
则第 \ref{weil-section-axioms} 节的公理 (C)(d) 成立。
\end{lemma}

\begin{proof}
由构造，$\gamma([\Spec(k)]) = 1 \in H^*(\Spec(k))$。由于
$$
H^0(\Spec(k)) = F,\quad
H^0(\Spec(k) \times \Spec(k)) = H^0(\Spec(k)) \otimes_F H^0(\Spec(k))
$$
公理 (A6) 中的映射 $\int_{\Spec(k)} = \lambda$ 必把 $1$ 送到 $1$，
因为引理 \ref{weil-lemma-trace-product} 已经给出
$\int_{\Spec(k) \times \Spec(k)} = \int_{\Spec(k)} \int_{\Spec(k)}$。
\end{proof}

\noindent
假设给定满足公理 (A1)--(A7) 的数据 (D0)、(D1) 与 (D2')。
于是得到第 \ref{weil-section-axioms} 节中的数据 (D0)、(D1)、(D2) 与 (D3)，
并满足第 \ref{weil-section-axioms} 节的公理 (A)、(B)、(C)(a)、(C)(c) 与 (C)(d)；见引理
\ref{weil-lemma-poincare-duality}、\ref{weil-lemma-trace-product} 与
\ref{weil-lemma-trace-base}。此外，还有第 \ref{weil-section-axioms} 节中构造的推前
$f_* : H^*(X) \to H^*(Y)$。第 \ref{weil-section-axioms} 节尚未明确的唯一公理是 (C)(b)。

\begin{lemma}
\label{weil-lemma-ok-for-projective-bundle}
假设给定满足公理 (A1)--(A7) 的数据 (D0)、(D1) 与 (D2')。
设 $p : P \to X$ 如公理 (A3)，其中 $X$ 非空且等维。则 $\gamma$ 与沿
$p$ 的推前可交换。
\end{lemma}

\begin{proof}
只需在 $\CH_*(P)$ 的生成元上证明。因而只需对形如
$\xi^i \cdot p^*\alpha$ 的周期类证明，其中 $0 \leq i \leq r - 1$ 且
$\alpha \in \CH_*(X)$。注意，若 $i < r - 1$，则
$p_*(\xi^i \cdot p^*\alpha) = 0$；而
$p_*(\xi^{r - 1} \cdot p^*\alpha) = \alpha$。另一方面，
$\gamma(\xi^i \cdot p^*\alpha) = c^i \cup p^*\gamma(\alpha)$；由投影公式
（引理 \ref{weil-lemma-pushforward}），有
$$
p_*\gamma(\xi^i \cdot p^*\alpha) = p_*(c^i) \cup \gamma(\alpha)
$$
所以只需证明：当 $i < r - 1$ 时 $p_*c^i = 0$，且
$p_*c^{r - 1} = 1$。等价地，只需证明由下列规则定义的
$\lambda_P : H^{2d + 2r - 2}(P)(d + r - 1) \to F$
$$
\lambda_P(c^i \cup p^*(a)) =
\left\{
\begin{matrix}
0 & \text{若} & i < r - 1 \\
\int_X(a) & \text{若} & i = r - 1
\end{matrix}
\right.
$$
满足公理 (A5) 的条件。这由引理 \ref{weil-lemma-diagonal-projective-bundle}
中对 $P$ 的对角线类的计算得出。
\end{proof}

\begin{lemma}
\label{weil-lemma-integrate-1}
假设给定满足公理 (A1)--(A7) 的数据 (D0)、(D1) 与 (D2')。
若 $k'/k$ 是 Galois 扩张，则
$\int_{\Spec(k')} 1 = [k' : k]$。
\end{lemma}

\begin{proof}
注意，作为 $\Spec(k') \times \Spec(k')$ 上的周期，有
$$
\Spec(k') \times \Spec(k') =
\coprod\nolimits_{\sigma \in \text{Gal}(k'/k)}
(\Spec(\sigma) \times \text{id})^{-1} \Delta
$$
对 $\gamma$ 的构造总是把 $[X]$ 送到 $H^0(X)$ 中的 $1$。因此
$1 \otimes 1 = 1 = \sum (\Spec(\sigma) \times \text{id})^*\gamma([\Delta])$。
以 $\lambda : H^0(\Spec(k')) \to F$ 表示公理 (A6) 给出的映射；换言之，
$(\text{id} \otimes \lambda)(\gamma(\Delta)) = 1$ 在
$H^0(\Spec(k'))$ 中成立。于是
\begin{align*}
\lambda(1) 1
& =
(\text{id} \otimes \lambda)(1 \otimes 1) \\
& =
(\text{id} \otimes \lambda)(
\sum (\Spec(\sigma) \times \text{id})^*\gamma([\Delta])) \\
& =
\sum (\Spec(\sigma) \times \text{id})^*(
(\text{id} \otimes \lambda)(\gamma([\Delta])) \\
& =
\sum (\Spec(\sigma) \times \text{id})^*(1) \\
& =
[k' : k]
\end{align*}
由于 $\lambda$ 是 $\int_{\Spec(k')}$ 的另一个名称（注记
\ref{weil-remark-trace}），证明完成。
\end{proof}

\begin{lemma}
\label{weil-lemma-enough}
假设给定满足公理 (A1)--(A7) 的数据 (D0)、(D1) 与 (D2')。
为证明 $\gamma$ 与推前可交换，只需证明：若
$i : Z \to X$ 是 $k$ 上非空光滑射影等维概形之间的闭浸入，则
$i_*(1) = \gamma([Z])$。
\end{lemma}

\begin{proof}
下面不再另行说明地使用这些事实：引理 \ref{weil-lemma-cycle-classes} 中构造的
周期类映射 $\gamma$ 与交积及拉回相容；并且第 \ref{weil-section-axioms} 节的公理
(A)、(B)、(C)(a)、(C)(c) 与 (C)(d) 已经成立，见引理
\ref{weil-lemma-poincare-duality}、注记 \ref{weil-remark-trace} 以及引理
\ref{weil-lemma-trace-product} 与 \ref{weil-lemma-trace-base}。特别地，例如可用引理
\ref{weil-lemma-pushforward} 看出 $H^*$ 上的推前与复合相容并满足投影公式。

\medskip\noindent
设 $f : X \to Y$ 是 $k$ 上非空等维光滑射影概形之间的态射。
目标是对 $X$ 上任意周期类 $\alpha$ 证明
$f_*\gamma(\alpha) = \gamma(f_*\alpha)$。可把 $\alpha$ 写成局部自由
$\mathcal{O}_X$-模的 Chern 类之乘积的 $\mathbf{Q}$-线性组合
（《Chow 同调》引理 \ref{chow-lemma-K-tensor-Q}）。因此可假设
$\alpha$ 是有限局部自由 $\mathcal{O}_X$-模
$\mathcal{E}_1, \ldots, \mathcal{E}_r$ 的 Chern 类之乘积。
按分裂原理（《Chow 同调》引理 \ref{chow-lemma-splitting-principle}）选择
$p : P \to X$。由《Chow 同调》注记
\ref{chow-remark-the-proof-shows-more}，$p$ 是若干射影空间丛的复合，并且
$\alpha = p_*(\xi_1 \cap \ldots \cap \xi_d \cap \cdot p^*\alpha)$，其中
$\xi_i$ 是可逆模的第一 Chern 类。由引理
\ref{weil-lemma-ok-for-projective-bundle}，$p_*$ 与周期类可交换。因此只需对复合
$f \circ p$ 证明此性质。由于
$p^*\mathcal{E}_1, \ldots, p^*\mathcal{E}_r$ 都有逐次商为可逆模的滤过，
问题化归到 $\alpha$ 形如
$\xi_1 \cap \ldots \cap \xi_t \cap [X]$ 的情形，其中 $\xi_i$ 是可逆模
$\mathcal{L}_i$ 的第一 Chern 类。

\medskip\noindent
假设
$\alpha = c_1(\mathcal{L}_1) \cap \ldots \cap c_1(\mathcal{L}_t) \cap [X]$，
其中 $\mathcal{L}_i$ 是 $X$ 上可逆模。令 $\mathcal{L}$ 为一个丰沛可逆
$\mathcal{O}_X$-模。当 $n \gg 0$ 时，可逆 $\mathcal{O}_X$-模
$\mathcal{L}^{\otimes n}$ 与 $\mathcal{L}_1 \otimes \mathcal{L}^{\otimes n}$
在 $k$ 上的 $X$ 上都非常丰沛；见《态射》引理
\ref{morphisms-lemma-invertible-add-enough-ample-very-ample}。由于
$c_1(\mathcal{L}_1) = c_1(\mathcal{L}_1 \otimes \mathcal{L}^{\otimes n})
- c_1(\mathcal{L}^{\otimes n})$，问题化归到 $\mathcal{L}_1$ 非常丰沛的情形。
对 $i = 2, \ldots, t$ 的 $\mathcal{L}_i$ 重复此过程，便可假设所有
$i = 1, \ldots, t$ 的 $\mathcal{L}_i$ 都在 $k$ 上的 $X$ 上非常丰沛。

\medskip\noindent
假设 $k$ 无限，且
$\alpha = 
c_1(\mathcal{L}_1) \cap \ldots \cap c_1(\mathcal{L}_t) \cap [X]$，
其中 $\mathcal{L}_i$ 是 $X$ 上相对于 $k$ 的非常丰沛可逆模。利用《簇》引理
\ref{varieties-lemma-bertini} 形式的 Bertini 定理，可以逐次找到
$\mathcal{L}_i$ 的正则截面 $s_i$，使概形
$Z(s_1) \cap \ldots \cap Z(s_i)$ 在 $k$ 上光滑且在 $X$ 中余维为 $i$。
由以可逆模的第一类作帽积的构造（追溯至《Chow 同调》定义
\ref{chow-definition-divisor-invertible-sheaf}），问题化归到
$\alpha = [Z]$ 的情形，其中 $Z \subset X$ 是非空、光滑、等维的闭子概形。

\medskip\noindent
假设 $\alpha = [Z]$，其中 $Z \subset X$ 是光滑闭子概形。选择闭嵌入
$X \to \mathbf{P}^n$。可将 $f$ 分解为
$$
X \to Y \times \mathbf{P}^n \to Y
$$
由引理 \ref{weil-lemma-ok-for-projective-bundle}，第二个态射的结论已知；所以只需处理
$\alpha = [Z]$、$i : Z \to X$ 为闭浸入且 $f$ 也为闭浸入的情形。
此时 $j = f \circ i$ 也是闭嵌入。对 $i$ 与 $j$ 使用假设即得结论。

\medskip\noindent
还须在 $k$ 为有限域时证明本引理。建议读者跳过余下证明。以上论证仍然成立，
唯一例外是：由于 $k$ 有限，不知道能否在 $k$ 上选择截面 $s_i$ 截出
$Z$。然而，由同一引理可知，在 $k$ 的代数闭包 $\overline{k}$ 上可以找到
$s_i$。因此存在有限扩张 $k'/k$，使这些截面 $s_i$ 定义在 $k'$ 上。
以 $\pi : X_{k'} \to X$ 表示投射。上述论证表明，由引理中的假设，
对周期 $\pi^*\alpha$ 及态射 $f \circ \pi : X_{k'} \to Y$ 可得所需结论。
又有 $\pi_*\pi^*\alpha = [k' : k] \alpha$；见《Chow 同调》引理
\ref{chow-lemma-finite-flat}。另一方面，
$$
\pi_*\gamma(\pi^*\alpha) = \pi_*\pi^*\gamma(\alpha) =
\gamma(\alpha) \pi_*1
$$
由该上同调理论的投影公式成立。注意 $\pi$ 是一个投射（！），所以由引理
\ref{weil-lemma-pr2star} 有 $\pi_*(1) = \int_{\Spec(k')}(1) 1$。
因此，要完成有限域情形的证明，只需证明
$\int_{\Spec(k')}(1) = [k' : k]$；这在引理
\ref{weil-lemma-integrate-1} 中完成。
\end{proof}

\noindent
以下引理使用《构造》第 \ref{constructions-section-grassmannian} 节中定义并构造的
Grassmann 簇。

\begin{lemma}
\label{weil-lemma-grassmanian}
假设给定满足公理 (A1)--(A7) 的数据 (D0)、(D1) 与 (D2')。
给定整数 $0 < l < n$ 以及 $k$ 上非空等维光滑射影概形 $X$，考虑投射态射
$p : X \times \mathbf{G}(l, n) \to X$。则 $\gamma$ 与沿 $p$ 的推前可交换。
\end{lemma}

\begin{proof}
若 $l = 1$ 或 $l = n - 1$，则 $p$ 是射影丛，结论由引理
\ref{weil-lemma-ok-for-projective-bundle} 得到。一般而言，存在态射
$$
h : Y \to X \times \mathbf{G}(l, n)
$$
使 $h$ 与 $p \circ h$ 都是若干射影空间丛的复合。具体而言，以
$\mathbf{G}(1, 2, \ldots, l; n)$ 表示部分旗簇。则态射
$$
\mathbf{G}(1, 2, \ldots, l; n) \to \mathbf{G}(l, n)
$$
是若干射影空间丛的复合，结构态射
$\mathbf{G}(1, 2, \ldots, l; n) \to \Spec(k)$ 亦然。因此可令
$Y = X \times \mathbf{G}(1, 2, \ldots, l; n)$。由于
$X \times \mathbf{G}(l, n)$ 上每个周期都是 $Y$ 上某周期的推前，
$Y \to X$ 的结论与 $Y \to X \times \mathbf{G}(l, n)$ 的结论共同推出
$p$ 的结论。
\end{proof}

\begin{lemma}
\label{weil-lemma-enough-better}
假设给定满足公理 (A1)--(A7) 的数据 (D0)、(D1) 与 (D2')。
为证明 $\gamma$ 与推前可交换，只需证明：若
$i : Z \to X$ 是 $k$ 上非空光滑射影等维概形之间的闭浸入，且
$K_0(Z)$ 中 $\mathcal{C}_{Z/X}$ 的类是 $K_0(X)$ 中某个类的拉回，则
$i_*(1) = \gamma([Z])$。
\end{lemma}

\begin{proof}
由引理 \ref{weil-lemma-enough}，只需对 $k$ 上非空光滑射影等维概形之间的闭浸入
$i : Z \to X$ 证明 $i_*(1) = \gamma([Z])$。设 $Z$ 在 $X$ 中余维为
$r$。令 $\mathcal{L}$ 是 $X$ 上充分丰沛的可逆模。选择 $n > 0$ 及满射
$$
\mathcal{O}_Z^{\oplus n} \to \mathcal{C}_{Z/X} \otimes \mathcal{L}|_Z
$$
这给出到 $k$ 上 Grassmann 簇的态射
$g : Z \to \mathbf{G}(n - r, n)$；见《构造》第
\ref{constructions-section-grassmannian} 节。考虑复合
$$
Z \to X \times \mathbf{G}(n - r, n) \to X
$$
由引理 \ref{weil-lemma-grassmanian}，沿第二个态射的推前与周期类相容。
闭浸入 $Z \to X \times \mathbf{G}(n - r, n)$ 的余法层 $\mathcal{C}$
位于短正合列
$$
0 \to \mathcal{C}_{Z/X} \to \mathcal{C} \to
g^*\Omega_{\mathbf{G}(n - r, n)} \to 0
$$
中。见《态射进阶》引理
\ref{more-morphisms-lemma-two-unramified-morphisms-formally-smooth}。
由于 $\mathcal{C}_{Z/X} \otimes \mathcal{L}|_Z$ 是
$\mathbf{G}(n - r, n)$ 上某个有限局部自由层的拉回，可知
$K_0(Z)$ 中 $\mathcal{C}$ 的类是
$K_0(X \times \mathbf{G}(n - r, n))$ 中某个类的拉回。因此，对于
$Z \to X \times \mathbf{G}(n - r, n)$ 已有该性质，结论随即得到。
\end{proof}

\begin{lemma}
\label{weil-lemma-injective-H0}
假设给定满足公理 (A1)--(A7) 的数据 (D0)、(D1) 与 (D2')。
若 $k''/k'/k$ 是有限可分域扩张，则
$H^0(\Spec(k')) \to H^0(\Spec(k''))$ 是单射。
\end{lemma}

\begin{proof}
可用 $k''$ 在 $k$ 上的正规闭包作替换；该闭包在 $k$ 上是 Galois 的，
见《域》引理 \ref{fields-lemma-normal-closure-galois}。于是 $k''$ 在 $k'$ 上也
是 Galois 的；见《域》引理 \ref{fields-lemma-galois-goes-up}。由此得到同构
$$
k' \otimes_k k'' \longrightarrow
\prod\nolimits_{\sigma \in \text{Gal}(k''/k')} k'',\quad
\eta \otimes \zeta \longmapsto (\eta \sigma(\zeta))_\sigma
$$
它给出同构
$\coprod_\sigma \Spec(k'') \to \Spec(k') \times \Spec(k'')$，
而在上同调上给出同构
$$
H^*(\Spec(k')) \otimes_F H^*(\Spec(k''))
\to
\prod\nolimits_\sigma H^*(\Spec(k'')),\quad
a' \otimes a'' \longmapsto (\pi^*a' \cup \Spec(\sigma)^*a'')_\sigma
$$
其中 $\pi : \Spec(k'') \to \Spec(k')$ 是由 $k'$ 到 $k''$ 的包含对应的态射。
取 $a'' = 1$ 即得本引理。
\end{proof}

\begin{lemma}
\label{weil-lemma-pushforward-blowup}
假设给定满足公理 (A1)--(A8) 的数据 (D0)、(D1) 与 (D2')。
设 $b : X' \to X$ 是 $k$ 上非空等维、维数为 $d$ 的光滑射影概形
$X$ 沿无处稠密光滑中心 $Z$ 的爆破。则 $b_*(1) = 1$。
\end{lemma}

\begin{proof}
可用 $X$ 的一个连通分支替换 $X$（略去一些细节）。因此可假设 $X$ 连通，
从而不可约。置
$k' = \Gamma(X, \mathcal{O}_X) = \Gamma(X', \mathcal{O}_{X'})$；
略去该等式的证明。选择不在例外除子中且剩余域 $k''$ 在 $k$ 上可分的闭点
$x' \in X'$；这可由《簇》引理
\ref{varieties-lemma-smooth-separable-closed-points-dense} 实现。以
$x \in X$ 表示其像（其剩余域当然也等于 $k''$）。考虑图表
$$
\xymatrix{
x' \times X' \ar[r] \ar[d] & X' \times X' \ar[d] \\
x \times X \ar[r] & X \times X
}
$$
对角线 $\Delta = \Delta_X$ 的类拉回为“对角点”
$\delta_x : x \to x \times X$ 的类；对角线 $\Delta'$ 的类亦然。
另一方面，对角点 $\delta_x$ 沿左侧竖直箭头拉回为对角点 $\delta_{x'}$。
写 $\gamma([\Delta]) = \sum \eta_i$，其中
$\eta_i \in H^i(X) \otimes H^{2d - i}(X)(d)$；并写
$\gamma([\Delta']) = \sum \eta'_i$，其中
$\eta'_i \in H^i(X') \otimes H^{2d - i}(X')(d)$。以上论证表明，
$\eta_0$ 与 $\eta'_0$ 映到
$$
H^0(x') \otimes_F H^{2d}(X')(d)
$$
中的同一个类。由公理 (A8)，
$H^0(\Spec(k')) = H^0(X) = H^0(X')$。由引理
\ref{weil-lemma-injective-H0}，这个公共值单射地映入 $H^0(x')$。因而
$\eta_0$ 通过映射
$$
H^0(X) \otimes_F H^{2d}(X)(d)
\longrightarrow
H^0(X') \otimes_F H^{2d}(X')(d)
$$
映到 $\eta'_0$。这意味着 $\int_X$ 等于 $\int_{X'}$ 与拉回映射的复合，
故本引理得证。
\end{proof}

\begin{lemma}
\label{weil-lemma-done}
假设给定满足公理 (A1)--(A8) 的数据 (D0)、(D1) 与 (D2')。
则周期类映射 $\gamma$ 与推前可交换。
\end{lemma}

\begin{proof}
令 $i : Z \to X$ 如引理 \ref{weil-lemma-enough-better}。考虑图表
$$
\xymatrix{
E \ar[r]_j \ar[d]_\pi & X' \ar[d]^b \\
Z \ar[r]^i & X
}
$$
令 $\theta \in \CH^{r - 1}(X')$ 如引理
\ref{weil-lemma-divide-pullback-good-blowing-up}。因为 $\pi$ 是射影空间丛，
由引理 \ref{weil-lemma-ok-for-projective-bundle}，关系
$\pi_*j^!\theta = [Z]$ 在 $\CH_*(Z)$ 中成立蕴含
$\pi_*\gamma(j^!\theta) = 1$。因此
$$
i_*(1) = i_*(\pi_*(\gamma(j^!\theta))) =
b_*j_*(j^*\gamma(\theta)) =
b_*(j_*(1) \cup \gamma(\theta))
$$
由 (A9)，$j_*(1) = \gamma([E])$。所以这等于
$$
b_*(\gamma([E]) \cup \gamma(\theta)) =
b_*(\gamma([E] \cdot \theta)) =
b_*(\gamma(b^*[Z])) =
b_*b^*\gamma([Z]) = b_*(1) \cup \gamma([Z])
$$
由引理 \ref{weil-lemma-pushforward-blowup}，$b_*(1) = 1$，证明完成。
\end{proof}

\begin{proposition}
\label{weil-proposition-get-weil}
假设给定满足公理 (A1)--(A8) 的数据 (D0)、(D1) 与 (D2')。
则得到一个 Weil 上同调理论。
\end{proposition}

\begin{proof}
第 \ref{weil-section-axioms} 节的公理 (A)、(B)、(C)(a)、(C)(c) 与 (C)(d)
由引理 \ref{weil-lemma-poincare-duality}、\ref{weil-lemma-trace-product} 与
\ref{weil-lemma-trace-base} 得到。由引理 \ref{weil-lemma-done}，
公理 (C)(b) 成立。最后，定义 \ref{weil-definition-weil-cohomology-theory}
中的附加条件成立，因为它与公理 (A8) 相同。
\end{proof}

\noindent
下面的引理有时可用于把问题化归到闭域，从而证明在非闭域上得到 Weil 上同调理论。

\begin{lemma}
\label{weil-lemma-check-over-extension}
设 $k'/k$ 是域扩张，$F'/F$ 是特征为 $0$ 的域扩张。假设给定：
\begin{enumerate}
\item 对 $k$ 与 $F$ 的数据 (D0)、(D1)、(D2')，记为
$F(1), H^*, c_1^H$；
\item 对 $k'$ 与 $F'$ 的数据 (D0)、(D1)、(D2')，记为
$F'(1), (H')^*, c_1^{H'}$；以及
\item 一个同构 $F(1) \otimes_F F' \to F'(1)$，以及在 $k$ 上光滑射影概形
$X$ 的范畴上的函子性同构
$H^*(X) \otimes_F F' \to (H')^*(X_{k'})$，使图表
$$
\xymatrix{
\Pic(X) \ar[r]_{c_1^H} \ar[d] & H^2(X)(1) \ar[d] \\
\Pic(X_{k'}) \ar[r]^{c_1^{H'}} & (H')^2(X_{k'})(1)
}
$$
交换。
\end{enumerate}
在此情形下，若 $F'(1), (H')^*, c_1^{H'}$ 满足公理 (A1)--(A9)，
则 $F(1), H^*, c_1^H$ 也满足这些公理。
\end{lemma}

\begin{proof}
逐条验证各公理。

\medskip\noindent
公理 (A1)。须证明 $H^*(\emptyset) = 0$，并且
$(i^*, j^*) : H^*(X \amalg Y) \to H^*(X) \times H^*(Y)$
是同构，其中 $i$ 与 $j$ 是余投射。由同构
$H^*(X) \otimes_F F' \to (H')^*(X_{k'})$ 的函子性，这归结为如下线性代数事实：
若 $\varphi : V \to W$ 是 $F$-向量空间之间的 $F$-线性映射，则
$\varphi$ 是同构，当且仅当
$\varphi_{F'} : V \otimes_F F' \to W \otimes_F F'$ 是同构。

\medskip\noindent
公理 (A2)。这意味着：给定 $k$ 上光滑射影概形之间的态射
$f : X \to Y$ 及可逆 $\mathcal{O}_Y$-模 $\mathcal{N}$，有
$f^*c_1^H(\mathcal{L}) = c_1^H(f^*\mathcal{L})$。这立即由
$c_1^{H'}$ 的相应性质、本引理中的交换图表以及如下事实得到：对任意
$F$-向量空间 $V$，典范映射 $V \to V \otimes_F F'$ 是单射。

\medskip\noindent
公理 (A3)。这由公理 (A1) 的证明中所述原理以及 $c_1^H$ 与
$c_1^{H'}$ 的相容性得到。

\medskip\noindent
公理 (A4)。设 $i : Y \to X$ 是 $k$ 上有效 Cartier 除子的包含，且
$X$ 与 $Y$ 都在 $k$ 上光滑且射影。对满足 $i^*a = 0$ 的
$a \in H^*(X)$，须证明 $a \cup c_1^H(\mathcal{O}_X(Y)) = 0$。
以 $a' \in (H')^*(X_{k'})$ 表示 $a$ 的像。该假设蕴含
$(i')^*a' = 0$，其中 $i' : Y_{k'} \to X_{k'}$ 是 $i$ 的基变换。
因此由 $(H')^*$ 的公理可得
$a' \cup c_1^{H'}(\mathcal{O}_{X_{k'}}(Y_{k'})) = 0$。由于
$a' \cup c_1^{H'}(\mathcal{O}_{X_{k'}}(Y_{k'}))$ 是
$a \cup c_1^H(\mathcal{O}_X(Y))$ 的像，结论由公理 (A2) 的证明中所述原理得到。

\medskip\noindent
公理 (A5)。这意味着 $H^*(\Spec(k)) = F$，并且对 $k$ 上光滑射影的
$X$ 与 $Y$，映射 $H^*(X) \otimes_F H^*(Y) \to H^*(X \times Y)$，
$a \otimes b \mapsto p^*(a) \cup q^*(b)$ 是同构，其中 $p$ 与 $q$ 是投射。
这由公理 (A1) 的证明中所述原理得到。

\medskip\noindent
暂时中断逐条论证，先证明对每个 $k$ 上光滑射影概形 $X$，图表
$$
\xymatrix{
\CH^*(X) \ar[r]_-\gamma \ar[d]_{g^*} & \bigoplus H^{2i}(X)(i) \ar[d] \\
\CH^*(X_{k'}) \ar[r]^-{\gamma'} & \bigoplus (H')^{2i}(X_{k'})(i)
}
$$
交换。注意公理 (A1)--(A4) 已知，故已有 $\gamma$；见引理
\ref{weil-lemma-cycle-classes}。左侧竖直箭头是《Chow 同调》第
\ref{chow-section-change-base} 节对基概形态射
$g : \Spec(k') \to \Spec(k)$ 讨论的映射。更确切地说，它是《Chow 同调》引理
\ref{chow-lemma-pullback-base-change} 给出的映射。取
$\alpha \in \CH^*(X)$。在 $\CH^*(X) \otimes \mathbf{Q}$ 中写
$\alpha = ch(\beta) \cap [X]$，其中
$\beta \in K_0(\textit{Vect}(X)) \otimes \mathbf{Q}$，从而
$\gamma(\alpha) = ch^{H}(\beta)$；这正是对 $\gamma$ 的构造。
《Chow 同调》引理 \ref{chow-lemma-pullback-base-change} 中的映射由引理
\ref{chow-lemma-pullback-base-change-chern-classes} 可知与以 Chern 类作帽积相容，
所以 $g^*\alpha = ch((X_{k'} \to X)^*\beta) \cap [X_{k'}]$。于是
$\gamma'(g^*\alpha) = ch^{H'}((X_{k'} \to X)^*\beta)$。

因此，只要对任意秩为 $r$ 的局部自由 $\mathcal{O}_X$-模
$\mathcal{E}$ 及 $0 \leq i \leq r$，$H^{2i}(X)(i)$ 中的元素
$c_i^H(\mathcal{E})$ 映到 $(H')^{2i}(X_{k'})(i)$ 中的元素
$c_i^{H'}(\mathcal{E}_{k'})$，图表便交换。由于对 $X$ 与 $X'$ 都有射影空间丛公式，
可用 $X$ 上的射影空间丛有限次替换 $X$ 来证明此事。因此可假设
$\mathcal{E}$ 有一个滤过，其分次片是可逆 $\mathcal{O}_X$-模
$\mathcal{L}_1, \ldots, \mathcal{L}_r$。见《Chow 同调》引理
\ref{chow-lemma-splitting-principle} 与注记
\ref{chow-remark-the-proof-shows-more}。此时 $c^H_i(\mathcal{E}$ 是
$c^H_1(\mathcal{L}_1), \ldots, c^H_1(\mathcal{L}_r)$ 的第 $i$ 个初等对称多项式；
由第一 Chern 类的一致性假设即得结论。

\medskip\noindent
公理 (A6)。设给定 $F$-向量空间 $V$、$W$，元素 $v \in V$ 以及张量
$\xi \in V \otimes_F W$。记 $V' = V \otimes_F F'$、
$W' = W \otimes_F F'$，并以 $v'$、$\xi'$ 表示 $v$、$\xi$ 在
$V'$、$V' \otimes_{F'} W'$ 中的像。证明公理 (A6) 所用的线性代数原理如下：
存在满足 $(1 \otimes \lambda)\xi = v$ 的 $F$-线性映射
$\lambda : W \to F$，当且仅当存在满足
$(1 \otimes \lambda')\xi' = v'$ 的 $F'$-线性映射
$\lambda' : W \otimes_F F' \to F'$。

\medskip\noindent
设 $X$ 是 $k$ 上非空等维光滑射影、维数为 $d$ 的概形。以
$\gamma = \gamma([\Delta])$ 表示 $H^{2d}(X \times X)(d)$ 中的元素
（未标注的纤维积均取在 $k$ 上）。回忆，公理 (A1)--(A4) 已知，故这有意义；
见引理 \ref{weil-lemma-cycle-classes}。在 K\"unneth 分解中可写
$$
\gamma = \sum \gamma_i, \quad
\gamma_i \in H^i(X) \otimes_F H^{2d - i}(X)(d)
$$
类似地，在 $(H')^{2d}(X_{k'} \times_{k'} X_{k'})(d)$ 中写
$\gamma' = \gamma([\Delta']) = \sum \gamma'_i$。由上述线性代数原理，只需证明
$\gamma_0$ 在 $(H')^0(X) \otimes_{F'} (H')^{2d}(X')(d)$ 中映到
$\gamma'_0$。由 K\"unneth 分解的相容性，只需证明 $\gamma$ 在
$$
(H')^{2d}(X_{k'} \times_{k'} X_{k'})(d) = (H')^{2d}((X \times X)_{k'})(d)
$$
中映到 $\gamma'$。由于 $\Delta_{k'} = \Delta'$，这由上面的讨论得到。

\medskip\noindent
公理 (A7)。这由如下线性代数事实得到：$F$-向量空间之间的线性映射
$V \to W$ 是单射，当且仅当
$V \otimes_F F' \to W \otimes_F F'$ 是单射。

\medskip\noindent
公理 (A8)。这由公理 (A1) 的证明中使用的线性代数事实得到。

\medskip\noindent
公理 (A9)。设 $X$ 是 $k$ 上非空光滑射影、维数为 $d$ 的等维概形。
设 $i : Y \to X$ 是在 $k$ 上光滑的非空有效 Cartier 除子。
令 $\lambda_Y$ 与 $\lambda_X$ 分别如公理 (A6) 对 $X$ 与 $Y$ 所给。
须证明：对 $a \in H^{2d - 2}(X)(d - 1)$，有
$\lambda_Y(i^*(a)) = \lambda_X(a \cup c_1^H(\mathcal{O}_X(Y))$。
由注记 \ref{weil-remark-trace}，映射
$\lambda_X : H^{2d}(X)(d) \to F$ 与
$\lambda_Y : H^{2d - 2}(Y)(d - 1)$ 由公理 (A6) 的要求唯一确定。
由上面对 $H^*$ 证明公理 (A6) 的过程可知，经给定的识别
$H^{2d}(X)(d) \otimes_F F' = H^{2d}(X_{k'})(d)$，
$\lambda_X \otimes \text{id}_{F'}$ 对应于 $\lambda_{X_{k'}}$。
因此，已知 $F'(1), (H')^*, c_1^{H'}$ 满足公理 (A9)，作图追踪便推出
$F(1), H^*, c_1^H$ 满足该公理。本引理证毕。
\end{proof}
% Generated mechanically from the frozen adequate.tex target.
% Do not edit this generated file; edit the bound locale source instead.

% OK, start here.
%

\setcounter{chapter}{45}
\chapter{适足模}



\phantomsection
\label{adequate-section-phantom}


\section{引言}
\label{adequate-section-introduction}

\noindent
对任意概形 $X$，准凝聚模范畴 $\QCoh(\mathcal{O}_X)$ 是阿贝尔范畴，并且是所有
$\mathcal{O}_X$-模所成阿贝尔范畴的一个弱 Serre 子范畴。把代数空间 $X$ 上的准凝聚模范畴
视为 $X$ 的小 \'etale 位点上所有 $\mathcal{O}_X$-模所成范畴的子范畴时，同样的结论成立。
此外，对于拟紧且拟分离的态射 $f : X \to Y$，推前 $f_*$ 及高阶直像都保持准凝聚性。

\medskip\noindent
接下来，设 $X$ 为概形，$\mathcal{O}$ 为 $X$ 的某个大位点（例如大 fppf 位点）上的结构层。
准凝聚 $\mathcal{O}$-模范畴是阿贝尔范畴（事实上，它等价于上面提到的概形 $X$ 上通常的
准凝聚 $\mathcal{O}_X$-模范畴），但它嵌入 $\textit{Mod}(\mathcal{O})$ 的函子并不正合。
例如，在 $\mathbf{A}^1_k = \Spec(k[x])$ 上考虑由乘以 $x$ 给出的准凝聚模映射
$$
\mathcal{O}_{\mathbf{A}^1_k}
\longrightarrow
\mathcal{O}_{\mathbf{A}^1_k}
$$
在准凝聚层的阿贝尔范畴中，此映射是单射；然而在 $\mathbf{A}^1_k$ 的大位点上所有
$\mathcal{O}$-模所成的阿贝尔范畴中，它有非平凡核。这一点可通过在
$\Spec(k[x]/(x)) = \Spec(k)$ 上取截面看出。此外，对于拟紧且拟分离的态射
$f : X \to Y$，函子 $f_{big, *}$ 并不保持准凝聚性。

\medskip\noindent
本章引入一类称为适足模的对象。它们与准凝聚模密切相关，并“修复”上述两个问题。
另一种方案将在讨论代数栈上的准凝聚模时采用：考虑局部准凝聚且满足平坦基变换性质的
$\mathcal{O}$-模。参见《栈上同调》第
\ref{stacks-cohomology-section-loc-qcoh-flat-base-change} 节、《栈上同调》评注
\ref{stacks-cohomology-remark-bousfield-colocalization}，以及《栈的导出范畴》第
\ref{stacks-perfect-section-derived} 节。






\section{约定}
\label{adequate-section-conventions}

\noindent
本章固定
$\tau \in \{Zar, \etale, smooth, syntomic, fppf\}$，
并固定《拓扑》第 \ref{topologies-section-procedure} 节所述的一个大
$\tau$-位点 $\Sch_\tau$。所有概形均视为 $\Sch_\tau$ 的对象。特别地，给定概形 $S$，
可得位点
$(\textit{Aff}/S)_\tau \subset (\Sch/S)_\tau$。
这些位点上的结构层 $\mathcal{O}$ 由规则
$\mathcal{O}(T) = \Gamma(T, \mathcal{O}_T)$ 定义。

\medskip\noindent
以下所有环 $A$ 都满足 $\Spec(A)$（同构于）$\Sch_\tau$ 的一个对象。给定环 $A$，
以 $\textit{Alg}_A$ 表示如下 $A$-代数范畴：其对象是形如
$B = \Gamma(U, \mathcal{O}_U)$ 的 $A$-代数 $B$，其中 $U$ 是 $\Sch_\tau$ 的仿射对象。
因此，给定仿射概形 $S = \Spec(A)$，函子
$$
(\textit{Aff}/S)_\tau \longrightarrow \textit{Alg}_A,
\quad
U \longmapsto \mathcal{O}(U)
$$
是一个等价。





\section{适足函子}
\label{adequate-section-quasi-coherent}

\noindent
本节讨论一个与准凝聚层的直像密切相关的主题。其中大部分材料取自论文
\cite{Jaffe}。

\begin{definition}
\label{adequate-definition-module-valued-functor}
设 $A$ 为环。所谓{\it 模值函子}，是满足下列条件的函子
$F : \textit{Alg}_A \to \textit{Ab}$：
\begin{enumerate}
\item 对 $\textit{Alg}_A$ 的每个对象 $B$，群 $F(B)$ 都赋有 $B$-模结构；
\item 对 $\textit{Alg}_A$ 的任意态射 $B \to B'$，映射
$F(B) \to F(B')$ 是 $B$-线性的。
\end{enumerate}
所谓{\it 模值函子之间的态射}，是函子变换 $\varphi : F \to G$，使得对所有
$B \in \Ob(\textit{Alg}_A)$，映射 $F(B) \to G(B)$ 都是 $B$-线性的。
\end{definition}

\noindent
设 $S = \Spec(A)$ 为仿射概形。$\textit{Alg}_A$ 上的模值函子范畴等价于

$\mathcal{O}$-模预层范畴
$\textit{PMod}((\textit{Aff}/S)_\tau, \mathcal{O})$。
该等价把模值函子 $F$ 送到由规则
$\mathcal{F}(U) = F(\mathcal{O}(U))$
定义的预层 $\mathcal{F}$。这直接来自第 \ref{adequate-section-conventions} 节给出的等价
$(\textit{Aff}/S)_\tau \to \textit{Alg}_A$, $U \mapsto \mathcal{O}(U)$。
其拟逆由 $F(B) = \mathcal{F}(\Spec(B))$ 给出。

\medskip\noindent
模值函子的一个重要特例构造如下。设 $M$ 为 $A$-模。用 $\underline{M}$ 表示
模值函子 $B \mapsto M \otimes_A B$（带有显然的 $B$-模结构）。注意，若
$M \to N$ 是 $A$-模映射，则它诱导模值函子之间的态射
$\underline{M} \to \underline{N}$。反过来，任意模值函子态射
$\underline{M} \to \underline{N}$ 都来自某个 $A$-模映射 $M \to N$；
在 $B = A$ 处取值即可看出这一点。换言之，$\text{Mod}_A$ 是
$\textit{Alg}_A$ 上模值函子范畴的满子范畴。

\medskip\noindent
给定 $A$-模映射 $\varphi : M \to N$，由于 $\otimes$ 右正合，有
$\Coker(\underline{M} \to \underline{N}) =
\underline{Q}$，其中 $Q = \Coker(M \to N)$。但核一般没有相同的性质：
例如，$A$-模的单射在基变换后未必仍是单射。因此下述定义是自然的。

\begin{definition}
\label{adequate-definition-adequate-functor}
设 $A$ 为环。称 $\textit{Alg}_A$ 上的模值函子 $F$
\begin{enumerate}
\item 为{\it 适足的}，若存在 $A$-模映射 $M \to N$，使得 $F$ 同构于
$\Ker(\underline{M} \to \underline{N})$；
\item 为{\it 线性适足的}，若 $F$ 同构于某个映射
$\underline{A^{\oplus n}} \to \underline{A^{\oplus m}}$ 的核。
\end{enumerate}
\end{definition}

\noindent
注意，$F$ 适足当且仅当存在正合列
$0 \to F \to \underline{M} \to \underline{N}$；而 $F$ 线性适足当且仅当存在正合列
$0 \to F \to \underline{A^{\oplus n}} \to \underline{A^{\oplus m}}$。

\medskip\noindent
设 $A$ 为环。本节将证明，$\textit{Alg}_A$ 上的适足函子构成阿贝尔范畴
（引理 \ref{adequate-lemma-cokernel-adequate} 与 \ref{adequate-lemma-kernel-adequate}），
并且有一族生成元（引理 \ref{adequate-lemma-adequate-surjection-from-linear}）。
我们还将看到，它是 $\textit{Alg}_A$ 上所有模值函子所成范畴的弱 Serre 子范畴
（引理 \ref{adequate-lemma-extension-adequate}），并且具有任意余极限
（引理 \ref{adequate-lemma-colimit-adequate}）。

\begin{lemma}
\label{adequate-lemma-adequate-finite-presentation}
设 $A$ 为环，$F$ 为 $\textit{Alg}_A$ 上的适足函子。若
$B = \colim B_i$ 是 $A$-代数的滤过余极限，则
$F(B) = \colim F(B_i)$。
\end{lemma}

\begin{proof}
这是因为，对任意 $A$-模 $M$，都有
$M \otimes_A B = \colim M \otimes_A B_i$（参见《代数》引理
\ref{algebra-lemma-tensor-products-commute-with-limits}），
并且滤过余极限与正合列可交换；参见《代数》引理
\ref{algebra-lemma-directed-colimit-exact}。
\end{proof}

 
\begin{remark}
\label{adequate-remark-settheoretic}
考虑范畴 $\textit{Alg}_{fp, A}$：其对象是形如
$B = A[x_1, \ldots, x_n]/(f_1, \ldots, f_m)$ 的 $A$-代数 $B$，
态射为 $A$-代数映射。每个 $A$-代数 $B$ 都是有限表示
$A$-代数的滤过余极限，也就是 $\textit{Alg}_{fp, A}$ 中对象的滤过余极限。
由引理 \ref{adequate-lemma-adequate-finite-presentation} 可知，每个适足函子 $F$
都由其在 $\textit{Alg}_{fp, A}$ 上的限制决定。因此，对于某些问题，
可以只考虑 $\textit{Alg}_{fp, A}$ 上的函子。例如，适足函子范畴与第
\ref{adequate-section-conventions} 节所选大 $\tau$-位点无关。
\end{remark}

\begin{lemma}
\label{adequate-lemma-adequate-flat}
设 $A$ 为环，$F$ 为 $\textit{Alg}_A$ 上的适足函子。若
$B \to B'$ 平坦，则
$F(B) \otimes_B B' \to F(B')$
是同构。
\end{lemma}

\begin{proof}
选取正合列 $0 \to F \to \underline{M} \to \underline{N}$。
由此得到图表
$$
\xymatrix{
0 \ar[r] & F(B) \otimes_B B' \ar[r] \ar[d] &
(M \otimes_A B)\otimes_B B' \ar[r] \ar[d] &
(N \otimes_A B)\otimes_B B' \ar[d] \\
0 \ar[r] & F(B') \ar[r] &
M \otimes_A B' \ar[r] &
N \otimes_A B'
}
$$
其中两行均正合（上行正合是因为 $B \to B'$ 平坦）。右侧两个竖直箭头均为同构，
故左侧箭头也是同构。
\end{proof}

\begin{lemma}
\label{adequate-lemma-adequate-surjection-from-linear}
设 $A$ 为环，$F$ 为 $\textit{Alg}_A$ 上的适足函子。则存在满射
$L \to F$，其中 $L$ 是若干线性适足函子的直和。
\end{lemma}

\begin{proof}
选取正合列 $0 \to F \to \underline{M} \to \underline{N}$，其中
$\underline{M} \to \underline{N}$ 由 $\varphi : M \to N$ 给出。
由引理 \ref{adequate-lemma-adequate-finite-presentation}，只需构造 $L \to F$，
使得对每个有限表示 $A$-代数 $B$，映射 $L(B) \to F(B)$ 都是满射。
因此，只需对给定的有限表示 $A$-代数 $B$ 和元素 $\xi \in F(B)$，
构造线性适足函子 $L$ 及映射 $L \to F$，使 $\xi$ 属于
$L(B) \to F(B)$ 的像。（这是因为此类对 $(B, \xi)$ 的同构类构成一个集合。）

\medskip\noindent
为此，把 $\xi$ 在 $\underline{M}(B) = M \otimes_A B$ 中的像写成
$\sum_{i = 1, \ldots, n} m_i \otimes b_i$。已知
$\sum \varphi(m_i) \otimes b_i = 0$ 于 $N \otimes_A B$ 中。
由于 $N$ 是有限表示 $A$-模的滤过余极限，可以找到有限表示 $A$-模 $N'$
以及 $A$-模的交换图
$$
\xymatrix{
A^{\oplus n} \ar[r] \ar[d]_{m_1, \ldots, m_n} & N' \ar[d] \\
M \ar[r] & N
}
$$
使得 $(b_1, \ldots, b_n)$ 在 $N' \otimes_A B$ 中映为零。
选取一个表示 $A^{\oplus l} \to A^{\oplus k} \to N' \to 0$，
再选取图中映射 $A^{\oplus n} \to N'$ 的一个提升
$A^{\oplus n} \to A^{\oplus k}$。于是存在
$(c_1, \ldots, c_l) \in B^{\oplus l}$，使
$(b_1, \ldots, b_n, c_1, \ldots, c_l)$ 在 $B^{\oplus k}$ 中映为零，
其中所用映射为 $B^{\oplus n} \oplus B^{\oplus l} \to B^{\oplus k}$。
考虑交换图
$$
\xymatrix{
A^{\oplus n} \oplus A^{\oplus l} \ar[r] \ar[d] & A^{\oplus k} \ar[d] \\
M \ar[r] & N
}
$$
其中左侧竖直箭头在直和项 $A^{\oplus l}$ 上为零。令 $L$ 为映射
$\underline{A^{\oplus n + l}}
\to \underline{A^{\oplus k}}$ 的核，
则它满足要求，因为元素
$(b_1, \ldots, b_n, c_1, \ldots, c_l) \in L(B)$ 映到 $\xi$。
\end{proof}

\noindent
考虑分次 $A$-代数 $B = \bigoplus_{d \geq 0} B_d$。有两个 $A$-代数映射
$p, a : B \to B[t, t^{-1}]$，即 $p : b \mapsto b$，以及当 $b$ 齐次时
$a : b \mapsto t^{\deg(b)} b$。若 $F$ 是 $\textit{Alg}_A$ 上的模值函子，定义
\begin{equation}
\label{adequate-equation-weight-k}
F(B)^{(k)} = \{\xi \in F(B) \mid t^k F(p)(\xi) = F(a)(\xi)\}.
\end{equation}
对于在平坦环扩张下表现良好的函子，这给出一个直和分解。这正是
$\mathbf{G}_m$ 的表示完全可约这一事实的体现。

\begin{lemma}
\label{adequate-lemma-flat-functor-split}
设 $A$ 为环，$F$ 为 $\textit{Alg}_A$ 上的模值函子。假设对每个平坦映射
$B \to B'$，映射
$F(B) \otimes_B B' \to F(B')$ 都是同构。设 $B$ 为分次 $A$-代数。则
\begin{enumerate}
\item $F(B) = \bigoplus_{k \in \mathbf{Z}} F(B)^{(k)}$；
\item 映射 $B \to B_0 \to B$ 诱导的映射 $F(B) \to F(B)$
的像包含于 $F(B)^{(0)}$。
\end{enumerate}
\end{lemma}

\begin{proof}
设 $x \in F(B)$。映射 $p : B \to B[t, t^{-1}]$ 是自由的，故
$$
F(B[t, t^{-1}]) =
\bigoplus\nolimits_{k \in \mathbf{Z}} F(p)(F(B)) \cdot t^k =
\bigoplus\nolimits_{k \in \mathbf{Z}} F(B) \cdot t^k
$$
如下所示，在证明余下部分中省略记号 $F(p)$。对某些 $x_k \in F(B)$，写
$F(a)(x) = \sum t^k x_k$。记
$\epsilon : B[t, t^{-1}] \to B$
为把 $t \mapsto 1$ 的 $B$-代数映射。注意，复合
$\epsilon \circ p, \epsilon \circ a : B \to B[t, t^{-1}] \to B$
均为恒等映射。因此
$$
x = F(\epsilon)(F(a)(x)) = F(\epsilon)(\sum t^k x_k) = \sum x_k.
$$
另一方面，断言 $x_k \in F(B)^{(k)}$。为此考虑交换图
$$
\xymatrix{
B \ar[r]_a \ar[d]_{a'} &
B[t, t^{-1}] \ar[d]^f \\
B[s, s^{-1}] \ar[r]^-g &
B[t, s, t^{-1}, s^{-1}]
}
$$
其中 $a'(b) = s^{\deg(b)}b$，$f(b) = b$，$f(t) = st$，
$g(b) = t^{\deg(b)}b$ 且 $g(s) = s$。于是
$$
F(g)(F(a'))(x) = F(g)(\sum s^k x_k) =
\sum s^k F(a)(x_k)
$$
沿另一条路径则有
$$
F(f)(F(a))(x) = F(f)(\sum t^k x_k) = \sum (st)^k x_k.
$$
由于 $B \to B[s, t, s^{-1}, t^{-1}]$ 是自由的，故
$F(B[t, s, t^{-1}, s^{-1}]) =
\bigoplus_{k, l \in \mathbf{Z}} F(B) \cdot t^ks^l$。
比较上述表达式的系数，得到所需的 $F(a)(x_k) = t^k x_k$。

\medskip\noindent
最后，$F(B_0) \to F(B)$ 的像包含于 $F(B)^{(0)}$，因为
$B_0 \to B \xrightarrow{a} B[t, t^{-1}]$ 等于
$B_0 \to B \xrightarrow{p} B[t, t^{-1}]$。
\end{proof}

\noindent
作为引理 \ref{adequate-lemma-flat-functor-split} 的一个特例，注意
$$
\underline{M}(B)^{(k)} = M \otimes_A B_k
$$
其中 $B_k$ 是分次 $A$-代数 $B$ 的 $k$ 次部分。

\begin{lemma}
\label{adequate-lemma-lift-map}
设 $A$ 为环。给定 $\textit{Alg}_A$ 上模值函子的实线图表
$$
\xymatrix{
0 \ar[r] &
L \ar[d]_\varphi \ar[r] &
\underline{A^{\oplus n}} \ar[r] \ar@{..>}[ld] &
\underline{A^{\oplus m}} \\
& \underline{M}
}
$$
且其行正合，则存在使图表交换的虚线箭头。
\end{lemma}

\begin{proof}
假设映射 $A^{\oplus n} \to A^{\oplus m}$ 由 $m \times n$ 矩阵
$(a_{ij})$ 给出。考虑环
$B = A[x_1, \ldots, x_n]/(\sum a_{ij}x_j)$。元素
$(x_1, \ldots, x_n) \in \underline{A^{\oplus n}}(B)$ 在
$\underline{A^{\oplus m}}(B)$ 中映为零，因而是唯一元素
$\xi \in L(B)$ 的像。注意，$\xi$ 具有如下泛性质：对任意 $A$-代数 $C$
和任意 $\xi' \in L(C)$，存在 $A$-代数映射 $B \to C$，使得
$\xi$ 经映射 $L(B) \to L(C)$ 映到 $\xi'$。

\medskip\noindent
注意，$B$ 是分次 $A$-代数，因而可用引理
\ref{adequate-lemma-flat-functor-split} 与 \ref{adequate-lemma-adequate-flat}
把这些函子在 $B$ 上的取值分解为分次部分。由于
$(x_1, \ldots, x_n)$ 是 $\underline{A^{\oplus n}}(B)$ 中的一次元素，
故 $\xi \in L(B)^{(1)}$。于是
$\varphi(\xi) \in \underline{M}(B)^{(1)} = M \otimes_A B_1$。
由于 $B_1$ 作为 $A$-模由 $x_1, \ldots, x_n$ 生成，可以写成
$\varphi(\xi) = \sum m_i \otimes x_i$。考虑映射
$A^{\oplus n} \to M$，它把第 $i$ 个基向量送到 $m_i$。
按构造，所诱导的映射
$\underline{A^{\oplus n}} \to \underline{M}$
把元素 $\xi$ 送到 $\varphi(\xi)$。由上述泛性质可知该图表交换。
\end{proof}

\begin{lemma}
\label{adequate-lemma-cokernel-into-module}
设 $A$ 为环，$\varphi : F \to \underline{M}$ 为 $\textit{Alg}_A$
上模值函子之间的映射，并设 $F$ 适足。则 $\Coker(\varphi)$ 适足。
\end{lemma}

\begin{proof}
由引理 \ref{adequate-lemma-adequate-surjection-from-linear}，可设
$F = \bigoplus L_i$ 是若干线性适足函子的直和。选取正合列
$0 \to L_i \to \underline{A^{\oplus n_i}} \to \underline{A^{\oplus m_i}}$。
对每个 $i$，依引理 \ref{adequate-lemma-lift-map} 选取映射
$A^{\oplus n_i} \to M$。考虑图表
$$
\xymatrix{
0 \ar[r] &
\bigoplus L_i  \ar[r] \ar[d] &
\bigoplus \underline{A^{\oplus n_i}} \ar[r] \ar[ld] &
\bigoplus \underline{A^{\oplus m_i}} \\
& \underline{M}
}
$$
考虑 $A$-模
$$
Q =
\Coker(\bigoplus A^{\oplus n_i} \to M \oplus \bigoplus A^{\oplus m_i})
\quad\text{且}\quad
P = \Coker(\bigoplus A^{\oplus n_i} \to \bigoplus A^{\oplus m_i}).
$$
于是 $\Coker(\varphi)$ 同构于 $\underline{Q} \to \underline{P}$ 的核。
\end{proof}

\begin{lemma}
\label{adequate-lemma-cokernel-adequate}
\begin{slogan}
环上的代数范畴中，适足函子之间映射的余核仍然适足。
\end{slogan}
设 $A$ 为环，$\varphi : F \to G$ 为 $\textit{Alg}_A$ 上适足函子之间的映射。
则 $\Coker(\varphi)$ 适足。
\end{lemma}

\begin{proof}
选取单射 $G \to \underline{M}$。于是有单射
$G/F \to \underline{M}/F$。由引理 \ref{adequate-lemma-cokernel-into-module}，
$\underline{M}/F$ 适足，因而可找到单射
$\underline{M}/F \to \underline{N}$。
复合后得到单射 $G/F \to \underline{N}$。再由引理
\ref{adequate-lemma-cokernel-into-module}，所诱导映射
$G \to \underline{N}$ 的余核适足，故可找到单射
$\underline{N}/G \to \underline{K}$。于是
$0 \to G/F \to \underline{N} \to \underline{K}$ 正合，结论成立。
\end{proof}

\begin{lemma}
\label{adequate-lemma-kernel-adequate}
设 $A$ 为环，$\varphi : F \to G$ 为 $\textit{Alg}_A$ 上适足函子之间的映射。
则 $\Ker(\varphi)$ 适足。
\end{lemma}

\begin{proof}
选取单射 $F \to \underline{M}$ 和单射
$G \to \underline{N}$。把 $F \to \underline{M \oplus N}$
记作对角映射，使图表
$$
\xymatrix{
F \ar[d] \ar[r] & G \ar[d] \\
\underline{M \oplus N} \ar[r] & \underline{N}
}
$$
交换。由引理 \ref{adequate-lemma-cokernel-adequate}，可找到模映射
$M \oplus N \to K$，使得 $F$ 是
$\underline{M \oplus N} \to \underline{K}$ 的核。
于是 $\Ker(\varphi)$ 是
$\underline{M \oplus N} \to \underline{K \oplus N}$ 的核。
\end{proof}

\begin{lemma}
\label{adequate-lemma-colimit-adequate}
设 $A$ 为环。$\textit{Alg}_A$ 上适足函子的任意直和仍然适足；
适足函子的余极限也适足。
\end{lemma}

\begin{proof}
关于直和的断言是直接的。一般余极限可写成两个直和之间某个映射的核；参见
《范畴》引理 \ref{categories-lemma-colimits-coproducts-coequalizers}。
因此结论由引理 \ref{adequate-lemma-kernel-adequate} 得出。
\end{proof}

\begin{lemma}
\label{adequate-lemma-flat-linear-functor}
设 $A$ 为环，$F, G$ 为 $\textit{Alg}_A$ 上的模值函子，
$\varphi : F \to G$ 为函子变换。假设
\begin{enumerate}
\item $\varphi$ 可加；
\item 对每个 $A$-代数 $B$、$\xi \in F(B)$ 以及单位
$u \in B^*$，在 $G(B)$ 中有 $\varphi(u\xi) = u\varphi(\xi)$；
\item 对任意平坦环映射 $B \to B'$，有
$G(B) \otimes_B B' = G(B')$。
\end{enumerate}
则 $\varphi$ 是模值函子之间的态射。
\end{lemma}

\begin{proof}
设 $B$ 为 $A$-代数，$\xi \in F(B)$ 且 $b \in B$。需证明
$\varphi(b \xi) = b \varphi(\xi)$。考虑环映射
$$
B \to B' = B[x, y, x^{-1}, y^{-1}]/(x + y - b).
$$
此环映射忠实平坦，故 $G(B) \subset G(B')$。另一方面，
$$
\varphi(b\xi) = \varphi((x + y)\xi) =
\varphi(x\xi) + \varphi(y\xi) = x\varphi(\xi) + y\varphi(\xi)
= (x + y)\varphi(\xi) = b\varphi(\xi)
$$
这是因为 $x, y$ 在 $B'$ 中都是单位。结论成立。
\end{proof}

\begin{lemma}
\label{adequate-lemma-extension-adequate-key}
设 $A$ 为环，并设
$0 \to \underline{M} \to G \to L \to 0$
是 $\textit{Alg}_A$ 上模值函子的短正合列，其中 $L$ 线性适足。则 $G$ 适足。
\end{lemma}

\begin{proof}
先指出：对任意平坦 $A$-代数映射 $B \to B'$，映射
$G(B) \otimes_B B' \to G(B')$ 是同构。事实上，这对
$\underline{M}$ 与 $L$ 成立，见引理 \ref{adequate-lemma-adequate-flat}；
于是由五引理对 $G$ 也成立。特别地，由引理
\ref{adequate-lemma-flat-functor-split}，对任意分次 $A$-代数 $B$，有
$G(B) = \bigoplus_{k \in \mathbf{Z}} G(B)^{(k)}$。

\medskip\noindent
选取正合列
$0 \to L \to \underline{A^{\oplus n}} \to \underline{A^{\oplus m}}$。
设映射 $A^{\oplus n} \to A^{\oplus m}$ 由 $m \times n$ 矩阵
$(a_{ij})$ 给出。考虑分次 $A$-代数
$B = A[x_1, \ldots, x_n]/(\sum a_{ij}x_j)$。元素
$(x_1, \ldots, x_n) \in \underline{A^{\oplus n}}(B)$ 在
$\underline{A^{\oplus m}}(B)$ 中映为零，故它是唯一元素
$\xi \in L(B)$ 的像。注意 $\xi \in L(B)^{(1)}$。映射
$$
\Hom_A(B, C) \longrightarrow L(C), \quad
f \longmapsto L(f)(\xi)
$$
定义一个函子同构。原因是 $f$ 由像
$c_i = f(x_i) \in C$ 决定，而这些像必须满足关系
$\sum a_{ij}c_j = 0$；另一方面，$L(C)$ 正是满足关系
$\sum a_{ij} c_j = 0$ 的 $n$-元组 $(c_1, \ldots, c_n)$ 所成集合。

\medskip\noindent
根据上述引理，函子 $\underline{M}$、$G$、$L$ 在 $B$ 上的取值均为各权空间的直和。
因此，$0 \to \underline{M} \to G \to L \to 0$ 的正合性蕴含序列
$0 \to \underline{M}(B)^{(1)} \to G(B)^{(1)} \to L(B)^{(1)} \to 0$
正合。于是可以选取映到 $\xi$ 的元素 $\theta \in G(B)^{(1)}$。

\medskip\noindent
考虑分次 $A$-代数
$$
C = A[x_1, \ldots, x_n, y_1, \ldots, y_n]/
(\sum a_{ij}x_j, \sum a_{ij}y_j)
$$
有三个分次 $A$-代数同态 $p_1, p_2, m : B \to C$，由下列规则定义：
$$
p_1(x_i) = x_i, \quad
p_1(x_i) = y_i, \quad
m(x_i) = x_i + y_i.
$$
将证明元素
$$
\tau = G(m)(\theta) - G(p_1)(\theta) - G(p_2)(\theta) \in G(C)
$$
为零。首先，直接计算可见 $\tau$ 在 $L(C)$ 中映为零，故
$\tau$ 是 $\underline{M}(C)$ 的元素。此外，由于
$m$, $p_1$, $p_2$ 都是分次代数映射，有
$\tau \in G(C)^{(1)}$；又因 $\underline{M} \subset G$，可得
$$
\tau \in \underline{M}(C)^{(1)} = M \otimes_A C_1.
$$
由于 $C_1 = p_1(B_1) \oplus p_2(B_1)$，可唯一地写成
$\tau = \underline{M}(p_1)(\tau_1) + \underline{M}(p_2)(\tau_2)$，
其中 $\tau_i \in M \otimes_A B_1 = \underline{M}(B)^{(1)}$。
考虑由 $x_i \mapsto x_i$、$y_i \mapsto 0$ 定义的环映射
$q_1 : C \to B$。于是
$\underline{M}(q_1)(\tau) =
\underline{M}(q_1)(\underline{M}(p_1)(\tau_1) + \underline{M}(p_2)(\tau_2)) =
\tau_1$。
另一方面，由于 $q_1 \circ m = q_1 \circ p_1$，有
$G(q_1)(\tau) = - G(q_1 \circ p_2)(\tau)$。又因
$q_1 \circ p_2$ 分解为 $B \to A \to B$，故
$G(q_1 \circ p_2)(\tau)$ 属于 $G(B)^{(0)}$；见引理
\ref{adequate-lemma-flat-functor-split}。因此 $\tau_1 = 0$，因为它属于
$G(B)^{(0)} \cap \underline{M}(B)^{(1)} \subset
G(B)^{(0)} \cap G(B)^{(1)} = 0$。
同理 $\tau_2 = 0$，从而 $\tau = 0$。

\medskip\noindent
由 $\theta \in G(B)$ 得到函子变换
$$
\psi : L(-) = \Hom_A(B, - ) \longrightarrow G(-)
$$
它把 $f : B \to C$ 送到 $G(f)(\theta)$。由于 $\theta$ 是 $\xi$ 的提升，
映射 $\psi$ 是 $G \to L$ 的右逆。用 $\psi$ 表述，上面证明的结论意味着：
$\tau = 0$ 表明 $\psi$ 可加；而 $\theta \in G(B)^{(1)}$ 蕴含，
对任意 $A$-代数 $D$、$l \in L(D)$ 及单位 $u \in D^*$，在 $G(D)$ 中有
$\psi(ul) = u\psi(l)$。由引理 \ref{adequate-lemma-flat-linear-functor}，
$\psi$ 是模值函子之间的态射。显然，这蕴含
$G \cong \underline{M} \oplus L$，结论成立。
\end{proof}

\begin{remark}
\label{adequate-remark-linearly-adequate}
设 $A$ 为环。引理 \ref{adequate-lemma-extension-adequate-key} 的证明表明：
若 $\textit{Alg}_A$ 上模值函子的扩张
$0 \to \underline{M} \to E \to L \to 0$
中 $L$ 线性适足，则该扩张分裂。证明只用到模值函子
$F = \underline{M}$ 的下列性质：
\begin{enumerate}
\item 对平坦环映射 $B \to B'$，映射
$F(B) \otimes_B B' \to F(B')$ 是同构；
\item
$F(C)^{(1)} = F(p_1)(F(B)^{(1)}) \oplus F(p_2)(F(B)^{(1)})$，
其中 $B = A[x_1, \ldots, x_n]/(\sum a_{ij}x_j)$ 且
$C = A[x_1, \ldots, x_n, y_1, \ldots, y_n]/
(\sum a_{ij}x_j, \sum a_{ij}y_j)$。
\end{enumerate}
任意适足函子 $F$ 都满足这两个性质；略去细节。因此，$L$ 是适足函子阿贝尔范畴中的
投射对象。
\end{remark}

\begin{lemma}
\label{adequate-lemma-extension-adequate}
设 $A$ 为环，并设
$0 \to F \to G \to H \to 0$
是 $\textit{Alg}_A$ 上模值函子的短正合列。若 $F$ 与 $H$ 适足，则 $G$ 也适足。
\end{lemma}

\begin{proof}
选取正合列 $0 \to F \to \underline{M} \to \underline{N}$。
若能证明 $(\underline{M} \oplus G)/F$ 适足，则 $G$ 是适足函子映射
$(\underline{M} \oplus G)/F \to \underline{N}$ 的核，
故由引理 \ref{adequate-lemma-kernel-adequate} 可知它适足。
因此可设 $F = \underline{M}$。

\medskip\noindent
由引理 \ref{adequate-lemma-adequate-surjection-from-linear}，可选取满射
$L \to H$，其中 $L$ 是若干线性适足函子的直和。
若能证明拉回 $G \times_H L$ 适足，则 $G$ 是映射
$\Ker(L \to H) \to G \times_H L$ 的余核，
故由引理 \ref{adequate-lemma-cokernel-adequate} 可知它适足。
因此可设 $H = \bigoplus L_i$ 是若干线性适足函子的直和。
由引理 \ref{adequate-lemma-extension-adequate-key}，每个拉回
$G \times_H L_i$ 都适足。再由引理
\ref{adequate-lemma-colimit-adequate}，$\bigoplus G \times_H L_i$ 适足。
而 $G$ 是下列映射的余核：
$$
\bigoplus\nolimits_{i \not = i'} F \longrightarrow
\bigoplus G \times_H L_i
$$
其中，直和项 $(i, i')$ 中的 $\xi$ 映到
$(0, \ldots, 0, \xi, 0, \ldots, 0, -\xi, 0, \ldots, 0)$，
其非零项位于直和项 $i$ 与 $i'$ 中。因此由引理
\ref{adequate-lemma-cokernel-adequate}，$G$ 适足。
\end{proof}

\begin{lemma}
\label{adequate-lemma-base-change-adequate}
设 $A \to A'$ 为环映射。若 $F$ 是 $\textit{Alg}_A$ 上的适足函子，
则其在 $\textit{Alg}_{A'}$ 上的限制 $F'$ 也适足。
\end{lemma}

\begin{proof}
选取正合列 $0 \to F \to \underline{M} \to \underline{N}$。则
$F'(B') = F(B') = \Ker(M \otimes_A B' \to N \otimes_A B')$。
由于 $M \otimes_A B' = M \otimes_A A' \otimes_{A'} B'$，且 $N$ 也同样，
可见 $F'$ 是
$\underline{M \otimes_A A'} \to \underline{N \otimes_A A'}$ 的核。
\end{proof}

\begin{lemma}
\label{adequate-lemma-pushforward-adequate}
设 $A \to A'$ 为环映射。若 $F'$ 是 $\textit{Alg}_{A'}$ 上的适足函子，
则 $\textit{Alg}_A$ 上的模值函子
$F : B \mapsto F'(A' \otimes_A B)$ 也适足。
\end{lemma}

\begin{proof}
选取正合列 $0 \to F' \to \underline{M'} \to \underline{N'}$。则
\begin{align*}
F(B) & = F'(A' \otimes_A B) \\
& = \Ker(M' \otimes_{A'} (
A' \otimes_A B) \to N' \otimes_{A'} (A' \otimes_A B)) \\
& = \Ker(M' \otimes_A B \to N' \otimes_A B)
\end{align*}
因此，$F$ 是 $\underline{M} \to \underline{N}$ 的核，其中
$M = M'$、$N = N'$ 均视为 $A$-模。
\end{proof}

\begin{lemma}
\label{adequate-lemma-adequate-product}
设 $A = A_1 \times \ldots \times A_n$ 为若干环的乘积。$A$ 上的一个适足函子，
等价于一列适足函子 $F_1, \ldots, F_n$，其中 $F_i$ 定义在 $A_i$ 上。
\end{lemma}

\begin{proof}
这是因为 $A$-代数 $B$ 典范地分解为乘积
$B_1 \times \ldots \times B_n$，而 $A$-模也有同样的分解。
令 $F(B) = \coprod F_i(B_i)$ 即得该对应。略去细节。
\end{proof}

\begin{lemma}
\label{adequate-lemma-adequate-descent}
设 $A \to A'$ 为环映射，$F$ 为 $\textit{Alg}_A$ 上的模值函子，并满足
\begin{enumerate}
\item $F$ 在 $A'$-代数范畴上的限制 $F'$ 适足；
\item 对任意 $A$-代数 $B$，序列
$$
0 \to F(B) \to F(B \otimes_A A') \to F(B \otimes_A A' \otimes_A A')
$$
正合。
\end{enumerate}
则 $F$ 适足。
\end{lemma}

\begin{proof}
函子 $B \to F(B \otimes_A A')$ 与
$B \mapsto F(B \otimes_A A' \otimes_A A')$ 都适足；见引理
\ref{adequate-lemma-pushforward-adequate} 与
\ref{adequate-lemma-base-change-adequate}。因此，$F$ 作为适足函子之间某个映射的核而适足；
见引理 \ref{adequate-lemma-kernel-adequate}。
\end{proof}






\section{适足函子的高阶 Ext 群}
\label{adequate-section-higher-ext}

\noindent
设 $A$ 为环。引理 \ref{adequate-lemma-extension-adequate} 已经表明：
在 $\textit{Alg}_A$ 上模值函子的范畴中，适足函子的任意扩张仍然适足。
本节将证明，高阶 Ext 群也具有同样的性质。

\begin{lemma}
\label{adequate-lemma-adjoint}
设 $A$ 为环。对 $\textit{Alg}_A$ 上的每个模值函子 $F$，
都存在 $\textit{Alg}_A$ 上模值函子的态射 $Q(F) \to F$，使得
(1) $Q(F)$ 适足，并且
(2) 对每个适足函子 $G$，映射
$\Hom(G, Q(F)) \to \Hom(G, F)$ 是双射。
\end{lemma}

\begin{proof}
选取一族线性适足函子 $\{L_i\}_{i \in I}$，使每个线性适足函子
都同构于某个 $L_i$。这是可以做到的。假设可以找到
$Q(F) \to F$，它满足 (1)，并满足下列性质 (2)'：
对每个 $i \in I$，映射
$\Hom(L_i, Q(F)) \to \Hom(L_i, F)$ 是双射。则 (2) 成立。
事实上，结合引理 \ref{adequate-lemma-adequate-surjection-from-linear} 与
\ref{adequate-lemma-kernel-adequate} 可知，每个适足函子 $G$ 都位于正合列
$$
K \to L \to G \to 0
$$
中，其中 $K$ 与 $L$ 都是若干线性适足函子的直和。因此，(2)'
蕴含映射
$\Hom(L, Q(F)) \to \Hom(L, F)$
与
$\Hom(K, Q(F)) \to \Hom(K, F)$
均为双射，从而对 $G$ 也有相同结论。

\medskip\noindent
考虑范畴 $\mathcal{I}$，其对象为二元组
$(i, \varphi)$，其中 $i \in I$ 且 $\varphi : L_i \to F$ 为态射。
态射 $(i, \varphi) \to (i', \varphi')$ 是满足
$\varphi' \circ \psi = \varphi$ 的映射
$\psi : L_i \to L_{i'}$。令
$$
Q(F) = \colim_{(i, \varphi) \in \Ob(\mathcal{I})} L_i
$$
存在自然映射 $Q(F) \to F$；由引理
\ref{adequate-lemma-colimit-adequate}，它适足；而且由构造可知它满足性质 (2)'。
\end{proof}

\begin{lemma}
\label{adequate-lemma-enough-injectives}
设 $A$ 为环。记 $\mathcal{P}$ 为 $\textit{Alg}_A$ 上模值函子的范畴，
记 $\mathcal{A}$ 为 $\textit{Alg}_A$ 上适足函子的范畴。
记 $i : \mathcal{A} \to \mathcal{P}$ 为包含函子，
并记 $Q : \mathcal{P} \to \mathcal{A}$ 为引理
\ref{adequate-lemma-adjoint} 中的构造。则
\begin{enumerate}
\item $i$ 全忠实且正合，其像为弱 Serre 子范畴；
\item $\mathcal{P}$ 有足够多的内射对象；
\item 函子 $Q$ 是 $i$ 的右伴随，因而左正合；
\item $Q$ 将内射对象映为内射对象；
\item $\mathcal{A}$ 有足够多的内射对象。
\end{enumerate}
\end{lemma}

\begin{proof}
本引理只是汇集前面已经得到的一些事实。由定义、弱 Serre 子范畴的刻画
（见《同调代数》引理 \ref{homology-lemma-characterize-weak-serre-subcategory}），以及引理
\ref{adequate-lemma-cokernel-adequate}、\ref{adequate-lemma-kernel-adequate} 和
\ref{adequate-lemma-extension-adequate}，(1) 是清楚的。
回忆 $\mathcal{P}$ 等价于范畴
$\textit{PMod}((\textit{Aff}/\Spec(A))_\tau, \mathcal{O})$。
因此由《内射对象》命题 \ref{injectives-proposition-presheaves-modules} 得到 (2)。
由引理 \ref{adequate-lemma-adjoint} 和《范畴》引理 \ref{categories-lemma-adjoint-exact} 得到 (3)。
由《同调代数》引理 \ref{homology-lemma-adjoint-preserve-injectives} 和
\ref{homology-lemma-adjoint-enough-injectives} 得到 (4) 与 (5)。
\end{proof}

\noindent
设 $A$ 为环。与《形式形变理论》第 \ref{formal-defos-section-tangent-spaces-functors} 节中一样，
给定 $A$-代数 $B$ 与 $B$-模 $N$，令 $B[N]$ 为底层 $B$-模为
$B \oplus N$ 的 $R$-代数，其乘法由
$(b, m)(b', m ') = (bb', bm' + b'm)$ 给出。
注意，该构造关于二元组 $(B, N)$ 具有函子性：态射
$(B, N) \to (B', N')$ 由一个 $A$-代数映射
$B \to B'$ 和一个 $B$-模映射 $N \to N'$ 给出。
在某种意义下，下一引理所定义的二元组函子 $TF$ 是 $F$ 的切空间。
以下只考虑使 $B[N]$ 成为 $\textit{Alg}_A$ 的对象的二元组 $(B, N)$。

\begin{lemma}
\label{adequate-lemma-tangent-functor}
设 $A$ 为环，$F$ 为模值函子。对每个
$B \in \Ob(\textit{Alg}_A)$ 及 $B$-模 $N$，都有典范分解
$$
F(B[N]) = F(B) \oplus TF(B, N)
$$
它由下列性质刻画：
\begin{enumerate}
\item $TF(B, N) = \Ker(F(B[N]) \to F(B))$；
\item $TF(B, N)$ 上有一个 $B$-模结构，它与 $F(B[N])$ 上的
$B[N]$-模结构相容；
\item $TF$ 是二元组 $(B, N)$ 的范畴上的函子；
\item
\label{adequate-item-mult-map}
存在典范映射 $N \otimes_B F(B) \to TF(B, N)$，
它们给出定义在二元组 $(B, N)$ 的范畴上的函子之间的变换；
\item $TF(B, 0) = 0$；当 $N \to N'$ 为零映射时，
映射 $TF(B, N) \to TF(B, N')$ 为零。
\end{enumerate}
\end{lemma}

\begin{proof}
由于 $B \to B[N] \to B$ 是恒等映射，可见 $F(B) \to F(B[N])$
是直和项，而它的补空间是 (1) 中定义的 $TF(N, B)$。
该构造关于二元组 $(B, N)$ 具有函子性，这是因为给定二元组的态射
$(B, N) \to (B', N')$，便在 $\textit{Alg}_A$ 中得到交换图
$$
\xymatrix{
B' \ar[r] & B'[N'] \ar[r] & B' \\
B \ar[r] \ar[u] & B[N] \ar[r] \ar[u] & B \ar[u]
}
$$
$B$-模结构来自 $B[N]$-模结构与环映射 $B \to B[N]$。
(4) 中的映射是复合
$$
N \otimes_B F(B) \longrightarrow
B[N] \otimes_{B[N]} F(B[N]) \longrightarrow F(B[N])
$$
其像包含于 $TF(B, N)$。（第一支箭头使用包含映射
$N \to B[N]$ 与 $F(B) \to F(B[N])$，第二支箭头是乘法映射。）
若 $N = 0$，则 $B = B[N]$，故 $TF(B, 0) = 0$。
若 $N \to N'$ 为零映射，则它分解为
$N \to 0 \to N'$；由于 $TF(B, 0) = 0$，诱导映射为零。
\end{proof}

\noindent
设 $A$ 为环，设 $M$ 为 $A$-模。则取值于模的函子
$\underline{M}$ 的切空间 $T\underline{M}$ 由规则
$T\underline{M}(B, N) = N \otimes_A M$ 给出。特别地，对给定的 $B$，
函子 $N \mapsto T\underline{M}(B, N)$ 是可加且右正合的。
事实证明，对内射的取值于模的函子也成立。

\begin{lemma}
\label{adequate-lemma-tangent-injective}
设 $A$ 为环，设 $I$ 为取值于模的函子范畴中的内射对象。
则对任意 $B \in \Ob(\textit{Alg}_A)$ 以及短正合序列
$0 \to N_1 \to N \to N_2 \to 0$
（其中为 $B$-模），序列
$$
TI(B, N_1) \to TI(B, N) \to TI(B, N_2) \to 0
$$
是正合的。
\end{lemma}

\begin{proof}
以下使用引理 \ref{adequate-lemma-tangent-functor} 的结果，不再另行说明。
记 $h : \textit{Alg}_A \to \textit{Sets}$ 为函子
$h(C) = \Mor_A(B[N], C)$。对 $h_1$ 与 $h_2$ 作同样记号。
与满射 $N \to N_2$ 对应的映射 $B[N] \to B[N_2]$
是满射。它对应映射 $h_2 \to h$，使得对所有 $A$-代数 $C$，
$h_2(C) \to h(C)$ 是单射。另一方面，有两个映射
$p, q : h \to h_1$，分别对应零映射 $N_1 \to N$ 与嵌入 $N_1 \to N$。
注意
$$
\xymatrix{
h_2 \ar[r] & h \ar@<1ex>[r] \ar@<-1ex>[r] & h_1
}
$$
是等化子图表。记 $\mathcal{O}_h$ 为取值于模的函子
$C \mapsto \bigoplus_{h(C)} C$。对 $\mathcal{O}_{h_1}$ 与
$\mathcal{O}_{h_2}$ 作同样记号。注意
$$
\Hom_\mathcal{P}(\mathcal{O}_h, F) = F(B[N])
$$
其中 $\mathcal{P}$ 是 $\textit{Alg}_A$ 上取值于模的函子范畴。
我们断言在 $\mathcal{P}$ 中有等化子图表
$$
\xymatrix{
\mathcal{O}_{h_2} \ar[r] &
\mathcal{O}_h \ar@<1ex>[r] \ar@<-1ex>[r] &
\mathcal{O}_{h_1}
}
$$
事实上，设 $C \in \Ob(\textit{Alg}_A)$，且
$\xi = \sum_{i = 1, \ldots, n} c_i \cdot f_i$，其中 $c_i \in C$，
$f_i : B[N] \to C$ 是 $\mathcal{O}_h(C)$ 的元素。
若 $p(\xi) = q(\xi)$，则
$$
\sum c_i \cdot f_i \circ z = \sum c_i \cdot f_i \circ y
$$
其中 $z, y : B[N_1] \to B[N]$ 为映射
$z : (b, m_1) \mapsto (b, 0)$ 及
$y : (b, m_1) \mapsto (b, m_1)$。这意味着对每个 $i$，
存在某个 $j$ 使 $f_j \circ z = f_i \circ y$。
显然这意味着 $f_i(N_1) = 0$，即 $f_i$ 经由唯一映射
$\overline{f}_i : B[N_2] \to C$ 分解。因此 $\xi$ 是
$\overline{\xi} = \sum c_i \cdot \overline{f}_i$ 的像。
由于 $I$ 是内射对象，它将此等化子图表变为余等化子图表
$$
\xymatrix{
I(B[N_1]) \ar@<1ex>[r] \ar@<-1ex>[r] &
I(B[N]) \ar[r] &
I(B[N_2])
}
$$
该图表与直和分解
$I(B[N]) = I(B) \oplus TI(B, N)$ 及
$I(B[N_i]) = I(B) \oplus TI(B, N_i)$ 相容。
零映射 $N \to N_1$ 诱导零映射
$TI(B, N) \to TI(B, N_1)$。
因此上述余等化子性质意味着序列
$TI(B, N_1) \to TI(B, N) \to TI(B, N_2) \to 0$
正合，如所需。
\end{proof}

\begin{lemma}
\label{adequate-lemma-exactness-implies}
设 $A$ 为环。设 $F$ 为取值于模的函子，满足：对任意
$B \in \Ob(\textit{Alg}_A)$，$B$-模上的函子 $TF(B, -)$
将 $B$-模的短正合序列变为右正合序列。则
\begin{enumerate}
\item $TF(B, N_1 \oplus N_2) = TF(B, N_1) \oplus TF(B, N_2)$；
\item $TF(B, N)$ 上存在第二个函子性的 $B$-模结构，
通过对 $x \in TF(B, N)$ 与 $b \in B$ 置
$x \cdot b = TF(B, b\cdot 1_N)(x)$ 定义；
\item
\label{adequate-item-mult-map-linear}
引理 \ref{adequate-lemma-tangent-functor} 中的典范映射
$N \otimes_B F(B) \to TF(B, N)$，
相对于第二个 $B$-模结构也是 $B$-线性的；
\item
\label{adequate-item-tangent-right-exact}
给定有限表示的 $B$-模 $N$，存在典范同构
$TF(B, B) \otimes_B N \to TF(B, N)$，
其中张量积使用 $TF(B, B)$ 上的第二个 $B$-模结构。
\end{enumerate}
\end{lemma}

\begin{proof}
以下使用引理 \ref{adequate-lemma-tangent-functor} 的结果，不再另行说明。
映射 $N_1 \to N_1 \oplus N_2$ 与 $N_2 \to N_1 \oplus N_2$
给出映射
$TF(B, N_1) \oplus TF(B, N_2) \to TF(B, N_1 \oplus N_2)$，
它是单射，因为映射 $N_1 \oplus N_2 \to N_1$ 与
$N_1 \oplus N_2 \to N_2$ 诱导逆映射。
由于 $TF$ 右正合，我们看到
$TF(B, N_1) \to TF(B, N_1 \oplus N_2) \to TF(B, N_2) \to 0$
正合。因此
$TF(B, N_1) \oplus TF(B, N_2) \to TF(B, N_1 \oplus N_2)$
是同构。这证明了 (1)。

\medskip\noindent
为证明 (2)，只需证明
$x \cdot (b_1 + b_2) = x \cdot b_1 + x \cdot b_2$。
（结合律与可加性显然。）为此考虑
$$
N \xrightarrow{(b_1, b_2)} N \oplus N \xrightarrow{+} N
$$
并应用 $TF(B, -)$。

\medskip\noindent
(3) 立即由映射
$N \otimes_B F(B) \to TF(B, N)$ 关于 $(B, N)$ 的函子性得到。

\medskip\noindent
设 $N$ 为有限表示的 $B$-模。取表示
$B^{\oplus m} \to B^{\oplus n} \to N \to 0$。这给出正合序列
$$
TF(B, B^{\oplus m}) \to TF(B, B^{\oplus n}) \to TF(B, N) \to 0
$$
由 $TF(B, -)$ 的右正合性。由 (1) 可写成
$TF(B, B^{\oplus m}) = TF(B, B)^{\oplus m}$ 且
$TF(B, B^{\oplus n}) = TF(B, B)^{\oplus n}$。接下来设
$B^{\oplus m} \to B^{\oplus n}$ 由矩阵 $T = (b_{ij})$ 给出。
则诱导映射 $TF(B, B)^{\oplus m} \to TF(B, B)^{\oplus n}$
由元素为 $TF(B, b_{ij} \cdot 1_B)$ 的矩阵给出。
结合 $\otimes$ 的右正合性即可证明 (4)。
\end{proof}

\begin{example}
\label{adequate-example-module-structure-different}
设 $F$ 为如引理 \ref{adequate-lemma-exactness-implies} 中的取值于模的函子。
$TF(B, N)$ 上的两个模结构不总是一致。下面给出一个例子。
设 $A = \mathbf{F}_p$，其中 $p$ 为素数。令 $F(B) = B$，
但其 $B$-模结构由 $b \cdot x = b^px$ 给出。则
$TF(B, N) = N$，其 $B$-模结构由
$b \cdot x = b^px$ 给出，其中 $x \in N$。然而，第二个 $B$-模
结构由 $x \cdot b = bx$ 给出。注意，在此情形下典范映射
$N \otimes_B F(B) \to TF(B, N)$ 为零，因为将元素
$n \in B[N]$ 提升到 $p$ 次方所得为零。
\end{example}

\noindent
在下一个引理中，我们将经常使用如下观察：若
$0 \to F \to G \to H \to 0$ 是 $\textit{Alg}_A$ 上模值函子的正合列，
则对任意二元组 $(B, N)$，序列
$0 \to TF(B, N) \to TG(B, N) \to TH(B, N) \to 0$ 也是正合的。
这是因为
$0 \to F(B[N]) \to G(B[N]) \to H(B[N]) \to 0$ 是正合的。

\begin{lemma}
\label{adequate-lemma-exactness-permanence}
设 $A$ 为环。对于模值函子 $F$（定义在 $\textit{Alg}_A$ 上），称性质
$(*)$ 成立，是指对所有 $B \in \Ob(\textit{Alg}_A)$，函子
$TF(B, -)$ 在 $B$-模上将 $B$-模的短正合列变为右正合列。设
$0 \to F \to G \to H \to 0$ 是 $\textit{Alg}_A$ 上模值函子的短正合列。
\begin{enumerate}
\item 若 $(*)$ 对 $F, G$ 成立，则 $(*)$ 对 $H$ 成立；
\item 若 $(*)$ 对 $F, H$ 成立，则 $(*)$ 对 $G$ 成立；
\item 若 $H' \to H$ 是 $\textit{Alg}_A$ 上模值函子的态射，且 $(*)$
对 $F$、$G$、$H$ 与 $H'$ 都成立，则 $(*)$ 对 $G \times_H H'$ 成立。
\end{enumerate}
\end{lemma}

\begin{proof}
固定 $B$。设
$0 \to N_1 \to N_2 \to N_3 \to 0$ 是 $B$-模的短正合列。
(1) 由下图的图追逐得到：
$$
\xymatrix{
0 \ar[r] &
TF(B, N_1) \ar[r] \ar[d] &
TG(B, N_1) \ar[r] \ar[d] &
TH(B, N_1) \ar[r] \ar[d] & 0 \\
0 \ar[r] &
TF(B, N_2) \ar[r] \ar[d] &
TG(B, N_2) \ar[r] \ar[d] &
TH(B, N_2) \ar[r] \ar[d] & 0 \\
0 \ar[r] &
TF(B, N_3) \ar[r] \ar[d] &
TG(B, N_3) \ar[r] \ar[d] &
TH(B, N_3) \ar[r] & 0 \\
& 0 & 0
}
$$
其中横行正合，涉及 $TF$ 与 $TG$ 的列也正合。
为证明 (2)，在下图中作图追逐：
$$
\xymatrix{
0 \ar[r] &
TF(B, N_1) \ar[r] \ar[d] &
TG(B, N_1) \ar[r] \ar[d] &
TH(B, N_1) \ar[r] \ar[d] & 0 \\
0 \ar[r] &
TF(B, N_2) \ar[r] \ar[d] &
TG(B, N_2) \ar[r] \ar[d] &
TH(B, N_2) \ar[r] \ar[d] & 0 \\
0 \ar[r] &
TF(B, N_3) \ar[r] \ar[d] &
TG(B, N_3) \ar[r] &
TH(B, N_3) \ar[r] \ar[d] & 0 \\
& 0 & & 0
}
$$
其中横行正合，涉及 $TF$ 与 $TH$ 的列也正合。
(3) 由 (2) 得到，因为 $G \times_H H'$ 位于正合列
$0 \to F \to G \times_H H' \to H' \to 0$ 中。
\end{proof}

\noindent
本节的大部分工作都是为了证明下面这个关键的消失结果。

\begin{lemma}
\label{adequate-lemma-ext-group-zero-key}
设 $A$ 为环，$M$、$P$ 为 $A$-模，且 $P$ 有限表示。则
$\Ext^i_\mathcal{P}(\underline{P}, \underline{M}) = 0$
对 $i > 0$ 成立，其中 $\mathcal{P}$ 是模值函子范畴，
这些函子定义在 $\textit{Alg}_A$ 上。
\end{lemma}

\begin{proof}
选取 $\underline{M} \to I^\bullet$ 在 $\mathcal{P}$ 中的内射分解，见引理
\ref{adequate-lemma-enough-injectives}。由《导出范畴》引理
\ref{derived-lemma-compute-ext-resolutions}，$\Ext^i_\mathcal{P}(\underline{P}, \underline{M})$
中的任意元素都来自态射 $\varphi : \underline{P} \to I^i$，且
$d^i \circ \varphi = 0$。我们将证明下列 Yoneda 扩张
$$
E : 0 \to \underline{M} \to I^0 \to \ldots \to
I^{i - 1} \times_{\Ker(d^i)} \underline{P} \to \underline{P} \to 0
$$
是 $\underline{P}$ 被 $\underline{M}$ 扩张、由 $\varphi$ 关联的平凡 Yoneda 扩张；
由《导出范畴》引理 \ref{derived-lemma-yoneda-extension}，这就证明了引理。

\medskip\noindent
对于模值函子 $F$（定义在 $\textit{Alg}_A$ 上），称性质 $(*)$ 成立，是指对所有
$B \in \Ob(\textit{Alg}_A)$，函子 $TF(B, -)$ 在 $B$-模上将 $B$-模的短正合列
变为右正合列。回忆模值函子 $\underline{M}, I^n, \underline{P}$
都满足性质 $(*)$，见引理 \ref{adequate-lemma-tangent-injective} 及其前面的注记。
将 $0 \to \underline{M} \to I^\bullet$ 分裂为短正合列后，由引理
\ref{adequate-lemma-exactness-permanence} 可知每个函子
$\Im(d^{n - 1}) = \Ker(d^n) \subset I^n$ 都满足 $(*)$，
而 $I^{i - 1} \times_{\Ker(d^i)} \underline{P}$ 也满足
$(*)$。

\medskip\noindent
因此可设 Yoneda 扩张写成
$$
E : 0 \to \underline{M} \to F_{i - 1} \to \ldots \to
F_0 \to \underline{P} \to 0
$$
其中每个模值函子 $F_j$ 都满足 $(*)$。令
$G_j(B) = TF_j(B, B)$，并通过引理
\ref{adequate-lemma-exactness-implies} 中定义的{\it 第二个} $B$-模结构
将其视为 $B$-模。由于 $TF_j$ 是二元组上的函子，$G_j$ 是
$\textit{Alg}_A$ 上的模值函子。此外，由于 $E$ 是正合列，
序列 $G_{j + 1} \to G_j \to G_{j - 1}$ 正合（见引理
\ref{adequate-lemma-exactness-permanence} 前的注记）。
注意 $T\underline{M}(B, B) = M \otimes_A B = \underline{M}(B)$，
且这里两个 $B$-模结构一致。因此得到 Yoneda 扩张
$$
E' :  0 \to \underline{M} \to G_{i - 1} \to \ldots \to
G_0 \to \underline{P} \to 0
$$
此外，引理 \ref{adequate-lemma-tangent-functor} (\ref{adequate-item-mult-map}) 的典范映射
$$
F_j(B) = B \otimes_B F_j(B) \longrightarrow TF_j(B, B) = G_j(B)
$$
由引理 \ref{adequate-lemma-exactness-implies} (\ref{adequate-item-mult-map-linear}) 可知是
$B$-线性的，并且关于 $B$ 具有函子性。因此得到 Yoneda 扩张的映射
$$
\xymatrix{
0 \ar[r] &
\underline{M} \ar[r] \ar[d]^1 &
F_{i - 1} \ar[r] \ar[d] &
\ldots \ar[r] &
F_0 \ar[r] \ar[d] &
\underline{P} \ar[r] \ar[d]^1 & 0 \\
0 \ar[r] &
\underline{M} \ar[r] &
G_{i - 1} \ar[r] &
\ldots \ar[r] &
G_0 \ar[r] &
\underline{P} \ar[r] & 0
}
$$
特别地，由上面提到的 Yoneda Ext 引理，$E$ 与 $E'$ 在
$\Ext^i_\mathcal{P}(\underline{P}, \underline{M})$ 中具有相同的类。
最后，令 $N$ 为有限表示的 $A$-模。则
$$
0 \to T\underline{M}(A, N) \to TF_{i - 1}(A, N) \to \ldots \to
TF_0(A, N) \to T\underline{P}(A, N) \to 0
$$
正合。由引理 \ref{adequate-lemma-exactness-implies} (\ref{adequate-item-tangent-right-exact})，
取 $B = A$，这转化为下列 $A$-模序列的正合性：
$$
0 \to M \otimes_A N \to G_{i - 1}(A) \otimes_A N \to \ldots \to
G_0(A) \otimes_A N \to P \otimes_A N \to 0
$$
因此 $A$-模序列
$0 \to M \to G_{i - 1}(A) \to \ldots \to G_0(A) \to P \to 0$
是泛正合的，即与任意有限表示的 $A$-模 $N$ 张量后仍正合。令
$K = \Ker(G_0(A) \to P)$，于是有正合列
$$
0 \to K \to G_0(A) \to P \to 0
\quad\text{且}\quad
G_2(A) \to G_1(A) \to K \to 0
$$
将第二个序列与 $N$ 张量，得到
$K \otimes_A N = \Coker(G_2(A) \otimes_A N \to G_1(A) \otimes_A N)$。
而 $G_2(A) \otimes_A N \to G_1(A) \otimes_A N \to G_0(A) \otimes_A N$
的正合性蕴含 $K \otimes_A N \to G_0(A) \otimes_A N$ 是单射。
由《代数》定理 \ref{algebra-theorem-universally-exact-criteria}，
这意味着 $A$-模扩张 $0 \to K \to G_0(A) \to P \to 0$ 是正合的；
又因假设 $P$ 有限表示，故该序列分裂，见《代数》引理
\ref{algebra-lemma-universally-exact-split}。
任一分裂 $P \to G_0(A)$ 都定义了映射 $\underline{P} \to G_0$，
它分裂满射 $G_0 \to \underline{P}$。因此 Yoneda 扩张 $E'$ 等价于
平凡的 Yoneda 扩张，证毕。
\end{proof}

\begin{lemma}
\label{adequate-lemma-ext-group-zero}
设 $A$ 为环，$M$ 为 $A$-模，$L$ 为 $\textit{Alg}_A$ 上的线性适足函子。则
$\Ext^i_\mathcal{P}(L, \underline{M}) = 0$
对 $i > 0$ 成立，其中 $\mathcal{P}$ 是模值函子范畴，
这些函子定义在 $\textit{Alg}_A$ 上。
\end{lemma}

\begin{proof}
由于 $L$ 线性适足，存在正合列
$$
0 \to L \to \underline{A^{\oplus m}} \to \underline{A^{\oplus n}} \to
\underline{P} \to 0
$$
其中 $P = \Coker(A^{\oplus m} \to A^{\oplus n})$ 是由线性适足函子定义
给出的有限自由 $A$-模映射的余核。由引理
\ref{adequate-lemma-ext-group-zero-key}，我们有
$\Ext^i_\mathcal{P}(\underline{P}, \underline{M})$
和
$\Ext^i_\mathcal{P}(\underline{A}, \underline{M})$
在 $i > 0$ 时均消失。令
$K = \Ker(\underline{A^{\oplus n}} \to \underline{P})$。
由与正合列
$0 \to K \to \underline{A^{\oplus n}} \to \underline{P} \to 0$
关联的 Ext 群长正合列，可得
$\Ext^i_\mathcal{P}(K, \underline{M}) = 0$ 对 $i > 0$ 成立。
再对正合列
$0 \to L \to \underline{A^{\oplus m}} \to K \to 0$
重复此论证即可。
\end{proof}

\begin{lemma}
\label{adequate-lemma-RQ-zero}
沿用引理 \ref{adequate-lemma-enough-injectives} 的记号，对任意适足函子 $F$，都有
$R^pQ(F) = 0$ 对所有 $p > 0$ 成立。
\end{lemma}

\begin{proof}
选取正合列 $0 \to F \to \underline{M^0} \to \underline{M^1}$。
令 $M^2 = \Coker(M^0 \to M^1)$，于是
$0 \to F \to \underline{M^0} \to \underline{M^1}
\to \underline{M^2} \to 0$ 是一个分解。由《导出范畴》引理
\ref{derived-lemma-two-ss-complex-functor}，得到谱序列
$$
R^pQ(\underline{M^q}) \Rightarrow R^{p + q}Q(F)
$$
由于 $Q(\underline{M^q}) = \underline{M^q}$，只需证明对任意
$A$-模 $M$，$R^pQ(\underline{M}) = 0$ 且 $p > 0$。

\medskip\noindent
在范畴 $\mathcal{P}$ 中选取内射分解 $\underline{M} \to I^\bullet$。
若 $R^iQ(\underline{M})$ 非零，则
$\Ker(Q(I^i) \to Q(I^{i + 1}))$ 严格大于
$Q(I^{i - 1}) \to Q(I^i)$ 的像。因此由引理
\ref{adequate-lemma-adequate-surjection-from-linear}，存在一个线性适足函子 $L$ 及映射
$\varphi : L \to Q(I^i)$，其像落在
$Q(I^i) \to Q(I^{i + 1})$ 的核中，但不经过
$Q(I^{i - 1}) \to Q(I^i)$ 的像分解。
由于 $Q$ 是包含函子的左伴随，映射
$\varphi$ 对应于映射 $\varphi' : L \to I^i$，且具有相同性质。
于是 $\varphi'$ 给出
$\Ext^i_\mathcal{P}(L, \underline{M})$ 中的非零元素，这与引理
\ref{adequate-lemma-ext-group-zero} 矛盾。
\end{proof}














\section{适足模}
\label{adequate-section-adequate}

\noindent
在《下降》一书的 \ref{descent-section-quasi-coherent-sheaves} 节中，
我们看到，概形 $S$ 上的拟凝聚模，正是与 $S$ 关联的任意大位点
$(\Sch/S)_\tau$ 上的拟凝聚模。我们也看到，这一识别有两个问题：
\begin{enumerate}
\item $\QCoh(\mathcal{O}_S) \to
\textit{Mod}((\Sch/S)_\tau, \mathcal{O})$、
$\mathcal{F} \mapsto \mathcal{F}^a$ 一般不是正合的
（《下降》引理 \ref{descent-lemma-equivalence-quasi-coherent-limits}），并且
\item 给定拟紧且拟分离的态射 $f : X \to S$，函子 $f_*$ 一般不保持
大位点上的拟凝聚层（《下降》命题
\ref{descent-proposition-equivalence-quasi-coherent-functorial}）。
\end{enumerate}
(1) 意味着我们不能在 $D(\mathcal{O})$ 中定义由上同调层拟凝聚的复形组成的三角子范畴。
(2) 意味着即使 $\mathcal{F}$ 拟凝聚且 $f$ 拟紧拟分离，
$Rf_*\mathcal{F}$ 也不是上同调层拟凝聚的复形。
此外，《下降》引理
\ref{descent-lemma-equivalence-quasi-coherent-limits}
和命题
\ref{descent-proposition-equivalence-quasi-coherent-functorial}
的证明中给出的例子并非病态的。

\medskip\noindent
本节讨论 $(\Sch/S)_\tau$ 上一个稍大的 $\mathcal{O}$-模范畴，
它包含拟凝聚模，是阿贝尔的，并且当 $f$ 拟紧拟分离时在 $f_*$ 下保持。
为此，设 $S$ 为概形，令 $\mathcal{F}$ 为 $(\Sch/S)_\tau$ 上的
$\mathcal{O}$-模预层。对 $(\Sch/S)_\tau$ 的任意仿射对象
$U = \Spec(A)$，可将 $\mathcal{F}$ 限制到 $(\textit{Aff}/U)_\tau$，
得到该位点上的 $\mathcal{O}$-模预层。相应的模值函子（见
\ref{adequate-section-quasi-coherent} 节）记为
$$
F = F_{\mathcal{F}, A} :
\textit{Alg}_A \longrightarrow \textit{Ab},
\quad
B \longmapsto \mathcal{F}(\Spec(B))
$$
赋值 $\mathcal{F} \mapsto F_{\mathcal{F}, A}$ 是阿贝尔范畴之间的正合函子。

\begin{definition}
\label{adequate-definition-adequate}
若 $(\Sch/S)_\tau$ 上的 $\mathcal{O}$-模层 $\mathcal{F}$ 存在一个
$\tau$-覆盖
$\{\Spec(A_i) \to S\}_{i \in I}$，使得对所有 $i \in I$，
$F_{\mathcal{F}, A_i}$ 适足，则称该层 {\it 适足}。
\end{definition}

\noindent
下面将看到，适足 $\mathcal{O}$-模的范畴与所选拓扑 $\tau$ 无关。

\begin{lemma}
\label{adequate-lemma-adequate-local}
设 $S$ 为概形，$\mathcal{F}$ 为 $(\Sch/S)_\tau$ 上的适足
$\mathcal{O}$-模。对 $S$ 上任意仿射概形 $\Spec(A)$，
函子 $F_{\mathcal{F}, A}$ 都适足。
\end{lemma}

\begin{proof}
设 $\{\Spec(A_i) \to S\}_{i \in I}$ 为 $\tau$-覆盖，且对所有 $i \in I$，
$F_{\mathcal{F}, A_i}$ 适足。可以找到标准仿射 $\tau$-覆盖
$\{\Spec(A'_j) \to \Spec(A)\}_{j = 1, \ldots, m}$，
使得 $\Spec(A'_j) \to \Spec(A) \to S$ 经过某个
$\Spec(A_{i(j)})$，其中 $i(j) \in I$。于是 $F_{\mathcal{F}, A'_j}$
是 $F_{\mathcal{F}, A_{i(j)}}$ 在 $A'_j$-代数范畴上的限制。
由此，依据引理 \ref{adequate-lemma-base-change-adequate}，
$F_{\mathcal{F}, A'_j}$ 适足。由引理
\ref{adequate-lemma-adequate-product}，序列
$F_{\mathcal{F}, A'_j}$ 对应于 $A' = A'_1 \times \ldots \times A'_m$
上的适足“乘积”函子 $F'$。由于 $\mathcal{F}$ 是层（对 Zariski 拓扑），
该乘积函子 $F'$ 等于 $F_{\mathcal{F}, A'}$，即 $F$ 在 $A'$-代数上的限制。
最后，$\{\Spec(A') \to \Spec(A)\}$ 是 $\tau$-覆盖。
由引理 \ref{adequate-lemma-adequate-descent} 可知 $F_{\mathcal{F}, A}$ 适足。
\end{proof}

\begin{lemma}
\label{adequate-lemma-adequate-affine}
设 $S = \Spec(A)$ 为仿射概形。$(\Sch/S)_\tau$ 上的适足
$\mathcal{O}$-模范畴与 $\textit{Alg}_A$ 上取值于模的适足函子范畴等价。
\end{lemma}

\begin{proof}
给定适足模 $\mathcal{F}$，函子 $F_{\mathcal{F}, A}$
由引理 \ref{adequate-lemma-adequate-local} 可知是适足的。
给定适足函子 $F$，取正合列
$0 \to F \to \underline{M} \to \underline{N}$，并考虑
$\mathcal{O}$-模 $\mathcal{F} = \Ker(M^a \to N^a)$，其中
$M^a$ 表示 $(\Sch/S)_\tau$ 上与 $S$ 上拟凝聚层
$\widetilde{M}$ 相关联的拟凝聚 $\mathcal{O}$-模。注意
$F = F_{\mathcal{F}, A}$；特别地，按定义 $\mathcal{F}$ 是适足的。
我们省略证明这些构造给出互逆的范畴等价。
\end{proof}

\begin{lemma}
\label{adequate-lemma-pullback-adequate}
设 $f : T \to S$ 为概形态射。
适足 $\mathcal{O}$-模 $\mathcal{F}$ 在 $(\Sch/S)_\tau$ 上的拉回
$f^*\mathcal{F}$，是 $(\Sch/T)_\tau$ 上的适足
$\mathcal{O}$-模。
\end{lemma}

\begin{proof}
拉回映射
$f^* : \textit{Mod}((\Sch/S)_\tau, \mathcal{O}) \to
\textit{Mod}((\Sch/T)_\tau, \mathcal{O})$
由限制给出，即对 $T$ 上任意概形 $V$，有
$f^*\mathcal{F}(V) = \mathcal{F}(V)$。
因此本引理立即由引理 \ref{adequate-lemma-adequate-local}
及定义得到。
\end{proof}

\noindent
下面给出适足 $\mathcal{O}$-模范畴的刻画。
为理解其意义，考虑概形 $S$ 上拟凝聚 $\mathcal{O}_S$-模之间的映射
$\mathcal{G} \to \mathcal{H}$。
所关联的 $\mathcal{O}$-模映射 $\mathcal{G}^a \to \mathcal{H}^a$
的余核是拟凝聚的，因为它等于 $(\mathcal{H}/\mathcal{G})^a$。
但一般而言，$\mathcal{G}^a \to \mathcal{H}^a$ 的核并非拟凝聚。
不过，它是适足的。

\begin{lemma}
\label{adequate-lemma-adequate-characterize}
设 $S$ 为概形，且 $\mathcal{F}$ 是 $(\Sch/S)_\tau$ 上的
$\mathcal{O}$-模。以下命题等价：
\begin{enumerate}
\item $\mathcal{F}$ 适足；
\item 存在仿射开覆盖 $S = \bigcup S_i$，以及拟凝聚
$\mathcal{O}_{S_i}$-模映射 $\mathcal{G}_i \to \mathcal{H}_i$，使得
$\mathcal{F}|_{(\Sch/S_i)_\tau}$ 是
$\mathcal{G}_i^a \to \mathcal{H}_i^a$ 的核；
\item 存在 $\tau$-覆盖 $\{S_i \to S\}_{i \in I}$，以及
$\mathcal{O}_{S_i}$-拟凝聚模映射 $\mathcal{G}_i \to \mathcal{H}_i$，使得
$\mathcal{F}|_{(\Sch/S_i)_\tau}$ 是
$\mathcal{G}_i^a \to \mathcal{H}_i^a$ 的核；
\item 存在 $\tau$-覆盖 $\{f_i : S_i \to S\}_{i \in I}$，使得每个
$f_i^*\mathcal{F}$ 都适足；
\item 对 $S$ 上任意仿射概形 $U$，限制
$\mathcal{F}|_{(\Sch/U)_\tau}$ 是拟凝聚
$\mathcal{O}_U$-模映射 $\mathcal{G}^a \to \mathcal{H}^a$ 的核。
\end{enumerate}
\end{lemma}

\begin{proof}
设 $U = \Spec(A)$ 为 $S$ 上的仿射概形。
令 $F = F_{\mathcal{F}, A}$。按定义，函子 $F$ 适足，当且仅当存在
$A$-模映射 $M \to N$，使得 $F = \Ker(\underline{M} \to \underline{N})$。
结合引理 \ref{adequate-lemma-adequate-local} 与
\ref{adequate-lemma-adequate-affine}，可见 (1) 与 (5) 等价。

\medskip\noindent
显然 (5) 蕴含 (2)，而 (2) 蕴含 (3)。若 (3) 成立，则可细化覆盖
$\{S_i \to S\}$，使每个 $S_i = \Spec(A_i)$ 都是仿射的。
于是根据证明的预备说明，$F_{\mathcal{F}, A_i}$ 适足。因此 $\mathcal{F}$
按定义适足，故 (3) 蕴含 (1)。

\medskip\noindent
最后，(4) 与 (1) 等价：一个方向使用引理
\ref{adequate-lemma-pullback-adequate}；另一个方向使用
$\tau$-覆盖的复合仍为 $\tau$-覆盖。
\end{proof}

\noindent
与拟凝聚层的情形一样，适足模范畴与所选拓扑无关。

\begin{lemma}
\label{adequate-lemma-adequate-fpqc}
设 $\mathcal{F}$ 是 $(\Sch/S)_\tau$ 上的适足 $\mathcal{O}$-模。
对 $S$ 上仿射概形之间的任意满射平坦态射
$\Spec(B) \to \Spec(A)$，扩张的 {\v C}ech 复形
$$
0 \to \mathcal{F}(\Spec(A)) \to
\mathcal{F}(\Spec(B)) \to
\mathcal{F}(\Spec(B \otimes_A B)) \to \ldots
$$
都是正合的。特别地，$\mathcal{F}$ 满足 fpqc 覆盖的层条件，
并且是 $(\Sch/S)_{fppf}$ 上的 $\mathcal{O}$-模层。
\end{lemma}

\begin{proof}
令 $A \to B$ 如引理所述，并设 $F = F_{\mathcal{F}, A}$。此函子由
引理 \ref{adequate-lemma-adequate-local} 可知是适足的。
由引理 \ref{adequate-lemma-adequate-flat}，由于
$A \to B$、$A \to B \otimes_A B$ 等都是平坦的，有
$F(B) = F(A) \otimes_A B$，
$F(B \otimes_A B) = F(A) \otimes_A B \otimes_A B$，等等。
正合性由《下降》引理 \ref{descent-lemma-ff-exact} 得出。

\medskip\noindent
因此，$\mathcal{F}$ 满足 $\tau$-覆盖（特别是 Zariski 覆盖）以及
仿射概形到仿射概形的任意忠实平坦覆盖的层条件。仿照《下降》引理
\ref{descent-lemma-standard-fpqc-covering} 及命题
\ref{descent-proposition-fpqc-descent-quasi-coherent} 的证明，
可知 $\mathcal{F}$ 满足所有 fpqc 覆盖的层条件（这些覆盖由
$(\Sch/S)_\tau$ 的对象组成）。细节略。
\end{proof}

\noindent
引理 \ref{adequate-lemma-adequate-fpqc} 特别表明，对任意两个拓扑 $\tau, \tau'$，
关于 $\tau$-拓扑和 $\tau'$-拓扑的适足模集合相同
（作为底层范畴 $\Sch/S$ 上的模预层）。

\begin{definition}
\label{adequate-definition-category-adequate-modules}
设 $S$ 为概形。$(\Sch/S)_\tau$ 上的适足 $\mathcal{O}$-模范畴记为
{\it $\textit{Adeq}(\mathcal{O})$} 或
{\it $\textit{Adeq}((\Sch/S)_\tau, \mathcal{O})$}。若不选定拓扑，
只考虑适足模的阿贝尔范畴，则简记为 {\it $\textit{Adeq}(S)$}。
\end{definition}

\begin{lemma}
\label{adequate-lemma-same-cohomology-adequate}
设 $S$ 为概形，且 $\mathcal{F}$ 是 $(\Sch/S)_\tau$ 上的适足
$\mathcal{O}$-模。
\begin{enumerate}
\item 限制 $\mathcal{F}|_{S_{Zar}}$ 是概形 $S$ 上的拟凝聚
$\mathcal{O}_S$-模。
\item 限制 $\mathcal{F}|_{S_\etale}$ 是与
$\mathcal{F}|_{S_{Zar}}$ 关联的拟凝聚模。
\item 对 $S$ 上任意仿射概形 $U$，都有 $H^q(U, \mathcal{F}) = 0$
（对所有 $q > 0$）。
\item 存在典范同构
$$
H^q(S, \mathcal{F}|_{S_{Zar}}) =
H^q((\Sch/S)_\tau, \mathcal{F}).
$$
\end{enumerate}
\end{lemma}

\begin{proof}
由引理 \ref{adequate-lemma-adequate-flat} 和引理
\ref{adequate-lemma-adequate-local} 可知，对 $S$ 上仿射概形之间的任意平坦态射
$U \to V$，都有
$\mathcal{F}(U) = \mathcal{F}(V) \otimes_{\mathcal{O}(V)} \mathcal{O}(U)$。
特别地，若 $U \subset V \subset S$ 是 $S$ 的仿射开子集，也成立，
因此 $\mathcal{F}|_{S_{Zar}}$ 拟凝聚。故 (1) 成立。

\medskip\noindent
设 $S' \to S$ 为概形的 \'etale 态射。对 $U \subset S'$ 中映到仿射开集
$V \subset S$ 的仿射开集，我们有
$\mathcal{F}(U) = \mathcal{F}(V) \otimes_{\mathcal{O}(V)} \mathcal{O}(U)$，
因为 $U \to V$ 为 \'etale，因而平坦。所以
$\mathcal{F}|_{S'_{Zar}}$ 是 $\mathcal{F}|_{S_{Zar}}$ 的拉回。
这就证明了 (2)。

\medskip\noindent
我们将把《位点上的上同调》引理
\ref{sites-cohomology-lemma-cech-vanish-collection}
应用于位点 $(\Sch/S)_\tau$，其中 $\mathcal{B}$ 是 $S$ 上仿射概形的集合，
而 $\text{Cov}$ 是标准仿射 $\tau$-覆盖的集合。
该引理的假设 (3) 由《下降》引理
\ref{descent-lemma-standard-covering-Cech} 以及
在单个仿射覆盖情形下的引理 \ref{adequate-lemma-adequate-fpqc} 满足。
因此，对 $S$ 上每个仿射概形 $U$，都有 $H^p(U, \mathcal{F}) = 0$。
这证明了 (3)。完全仿照《下降》命题
\ref{descent-proposition-same-cohomology-quasi-coherent} 的证明，
即得到上同调的相等性 (4)。
\end{proof}

\begin{remark}
\label{adequate-remark-compare}
设 $S$ 为概形。我们有函子
$u : \QCoh(\mathcal{O}_S) \to \textit{Adeq}(\mathcal{O})$
以及
$v : \textit{Adeq}(\mathcal{O}) \to \QCoh(\mathcal{O}_S)$。
具体而言，函子 $u : \mathcal{F} \mapsto \mathcal{F}^a$
来自取相应的 $\mathcal{O}$-模；根据引理
\ref{adequate-lemma-adequate-characterize}，该模是适足的。
反之，函子 $v$ 来自限制
$v : \mathcal{G} \mapsto \mathcal{G}|_{S_{Zar}}$，见引理
\ref{adequate-lemma-same-cohomology-adequate}。
由于 $\mathcal{F}^a$ 可描述为在环化拓扑斯之间的态射
$((\Sch/S)_\tau, \mathcal{O}) \to (S_{Zar}, \mathcal{O}_S)$
下对 $\mathcal{F}$ 的拉回，见《下降》注记
\ref{descent-remark-change-topologies-ringed-sites}；又因为限制是推前，
故 $u$ 与 $v$ 如下互为伴随：
$$
\SheafHom_{\mathcal{O}_S}(\mathcal{F}, v\mathcal{G})
=
\SheafHom_\mathcal{O}(u\mathcal{F}, \mathcal{G})
$$
其中 $\mathcal{O}$ 表示大位点上的结构层。
从上述描述立即可见，对所有拟凝聚层，伴随映射
$\mathcal{F} \to vu\mathcal{F}$ 都是同构。
\end{remark}

\begin{lemma}
\label{adequate-lemma-sheafification-adequate}
设 $S$ 为概形，且 $\mathcal{F}$ 是 $(\Sch/S)_\tau$ 上的
$\mathcal{O}$-模预层。若对 $S$ 上每个仿射概形
$\Spec(A)$，函子 $F_{\mathcal{F}, A}$ 都适足，则 $\mathcal{F}$ 的层化是适足的
$\mathcal{O}$-模。
\end{lemma}

\begin{proof}
设 $U = \Spec(A)$ 为 $S$ 上的仿射概形。
令 $F = F_{\mathcal{F}, A}$。
层化为 $\mathcal{F}^\# = (\mathcal{F}^+)^+$，见《位点》
第 \ref{sites-section-sheafification} 节。
按构造
$$
(\mathcal{F})^+(U) =
\colim_\mathcal{U} \check{H}^0(\mathcal{U}, \mathcal{F})
$$
其中余极限取遍位点 $(\Sch/S)_\tau$ 中的覆盖。
由于 $U$ 仿射，只需对 $U$ 的标准仿射 $\tau$-覆盖
$\mathcal{U} = \{U_i \to U\}_{i \in I} =
\{\Spec(A_i) \to \Spec(A)\}_{i \in I}$ 取极限。
由于每个 $A \to A_i$ 及 $A \to A_i \otimes_A A_j$ 都平坦，有
$$
\check{H}^0(\mathcal{U}, \mathcal{F}) =
\Ker(\prod F(A) \otimes_A A_i \to \prod F(A) \otimes_A A_i \otimes_A A_j)
$$
由引理 \ref{adequate-lemma-adequate-flat} 得出。
由于 $A \to \prod A_i$ 忠实平坦，根据《下降》引理
\ref{descent-lemma-ff-exact}，上述对象总是典范同构于 $F(A)$。
因此，在 $S$ 上所有仿射概形处，预层 $(\mathcal{F})^+$ 与
$\mathcal{F}$ 取相同的值。再次重复该论证，可知
$\mathcal{F}^\# = ((\mathcal{F})^+)^+$ 也有同样性质。
于是 $F_{\mathcal{F}, A} = F_{\mathcal{F}^\#, A}$，从而
$\mathcal{F}^\#$ 是适足的。
\end{proof}

\begin{lemma}
\label{adequate-lemma-abelian-adequate}
设 $S$ 为概形。
\begin{enumerate}
\item 范畴 $\textit{Adeq}(\mathcal{O})$ 是阿贝尔范畴。
\item 函子
$\textit{Adeq}(\mathcal{O}) \to
\textit{Mod}((\Sch/S)_\tau, \mathcal{O})$
是正合的。
\item 若 $0 \to \mathcal{F}_1 \to \mathcal{F}_2 \to \mathcal{F}_3 \to 0$
是 $\mathcal{O}$-模的短正合列，且
$\mathcal{F}_1$ 与 $\mathcal{F}_3$ 适足，则
$\mathcal{F}_2$ 适足。
\item 范畴 $\textit{Adeq}(\mathcal{O})$ 有余极限，并且函子
$\textit{Adeq}(\mathcal{O}) \to
\textit{Mod}((\Sch/S)_\tau, \mathcal{O})$
与余极限交换。
\end{enumerate}
\end{lemma}

\begin{proof}
设 $\varphi : \mathcal{F} \to \mathcal{G}$ 是适足 $\mathcal{O}$-模之间的映射。
为证明 (1) 和 (2)，只需证明在
$\textit{Mod}((\Sch/S)_\tau, \mathcal{O})$ 中计算的
$\mathcal{K} = \Ker(\varphi)$ 与 $\mathcal{Q} = \Coker(\varphi)$ 都适足。
设 $U = \Spec(A)$ 为 $S$ 上的仿射概形。
令 $F = F_{\mathcal{F}, A}$，$G = F_{\mathcal{G}, A}$。
由引理 \ref{adequate-lemma-kernel-adequate} 与
\ref{adequate-lemma-cokernel-adequate}，诱导映射
$F \to G$ 的核 $K$ 与余核 $Q$ 是适足函子。
由于核是在预层层面计算的，$K = F_{\mathcal{K}, A}$，
故 $\mathcal{K}$ 适足。为证明余核的结论，记 $\mathcal{Q}'$ 为
$\varphi$ 的预层余核。于是 $Q = F_{\mathcal{Q}', A}$ 且
$\mathcal{Q} = (\mathcal{Q}')^\#$。因此由引理
\ref{adequate-lemma-sheafification-adequate}，$\mathcal{Q}$ 适足。

\medskip\noindent
设 $0 \to \mathcal{F}_1 \to \mathcal{F}_2 \to \mathcal{F}_3 \to 0$
是 $\mathcal{O}$-模的短正合列，且
$\mathcal{F}_1$ 与 $\mathcal{F}_3$ 适足。
设 $U = \Spec(A)$ 为 $S$ 上的仿射概形。
令 $F_i = F_{\mathcal{F}_i, A}$。函子列
$$
0 \to F_1 \to F_2 \to F_3 \to 0
$$
是正合的，因为对 $U$ 上的仿射概形 $V = \Spec(B)$，由引理
\ref{adequate-lemma-same-cohomology-adequate} 有 $H^1(V, \mathcal{F}_1) = 0$。
由于由引理 \ref{adequate-lemma-adequate-local} 可知 $F_1$ 与 $F_3$ 是适足函子，
故由引理 \ref{adequate-lemma-extension-adequate}，$F_2$ 适足。
因此 $\mathcal{F}_2$ 适足。

\medskip\noindent
设 $\mathcal{I} \to \textit{Adeq}(\mathcal{O})$，$i \mapsto \mathcal{F}_i$
为一个图表。记 $\mathcal{F} = \colim_i \mathcal{F}_i$，其余极限在
$\textit{Mod}((\Sch/S)_\tau, \mathcal{O})$ 中计算。
为证明 (4)，只需证明 $\mathcal{F}$ 适足。
令 $\mathcal{F}' = \colim_i \mathcal{F}_i$ 为在 $\mathcal{O}$-模预层中计算的余极限。
于是 $\mathcal{F} = (\mathcal{F}')^\#$。
设 $U = \Spec(A)$ 为 $S$ 上的仿射概形。
令 $F_i = F_{\mathcal{F}_i, A}$。由引理
\ref{adequate-lemma-colimit-adequate}，函子
$\colim_i F_i = F_{\mathcal{F}', A}$ 适足。
引理 \ref{adequate-lemma-sheafification-adequate} 表明 $\mathcal{F}$ 适足。
\end{proof}

\noindent
下述引理说明，拟紧且拟分离态射下适足模的总直像
$Rf_*\mathcal{F}$ 是一个上同调层均适足的复形。

\begin{lemma}
\label{adequate-lemma-direct-image-adequate}
设 $f : T \to S$ 为拟紧且拟分离的概形态射。对
$(\Sch/T)_\tau$ 上任意适足 $\mathcal{O}_T$-模，推前
$f_*\mathcal{F}$ 以及高阶直像 $R^if_*\mathcal{F}$ 都是
$(\Sch/S)_\tau$ 上的适足 $\mathcal{O}_S$-模。
\end{lemma}

\begin{proof}
首先说明高阶直像的计算方式。取内射分解
$\mathcal{F} \to \mathcal{I}^\bullet$。
于是 $R^if_*\mathcal{F}$ 是复形 $f_*\mathcal{I}^\bullet$ 的第 $i$ 个上同调层。
因此，$R^if_*\mathcal{F}$ 是如下预层所关联的层：它把
$(\Sch/S)_\tau$ 的对象 $U/S$ 送到模
\begin{align*}
\frac{\Ker(f_*\mathcal{I}^i(U) \to f_*\mathcal{I}^{i + 1}(U))}
{\Im(f_*\mathcal{I}^{i - 1}(U) \to f_*\mathcal{I}^i(U))}
& =
\frac{\Ker(\mathcal{I}^i(U \times_S T) \to
\mathcal{I}^{i + 1}(U \times_S T))}
{\Im(\mathcal{I}^{i - 1}(U \times_S T) \to \mathcal{I}^i(U \times_S T))}
\\
& =
H^i(U \times_S T, \mathcal{F}) \\
& = H^i((\Sch/U \times_S T)_\tau,
\mathcal{F}|_{(\Sch/U \times_S T)_\tau}) \\
& = H^i(U \times_S T, \mathcal{F}|_{(U \times_S T)_{Zar}})
\end{align*}
第一个等号由《拓扑》引理
\ref{topologies-lemma-morphism-big-fppf}（以及其他拓扑的相应结论）给出；
第二个等号来自在 $(\Sch/T)_\tau$ 的对象上计算 $\mathcal{F}$ 的上同调的定义；
第三个等号由《位点上的上同调》引理
\ref{sites-cohomology-lemma-cohomology-of-open} 给出；最后一个等号由
引理 \ref{adequate-lemma-same-cohomology-adequate} 给出。
因此根据引理 \ref{adequate-lemma-sheafification-adequate}，只需证明下一段中的断言。

\medskip\noindent
设 $A$ 为环，$T$ 为在 $A$ 上拟紧且拟分离的概形，且 $\mathcal{F}$ 是
$(\Sch/T)_\tau$ 上的适足 $\mathcal{O}_T$-模。对 $A$-代数 $B$，令
$T_B = T \times_{\Spec(A)} \Spec(B)$，并记
$\mathcal{F}_B = \mathcal{F}|_{(T_B)_{Zar}}$，即 $\mathcal{F}$ 到
$T_B$ 的小 Zariski 位点的限制。
（回顾一下，这是概形 $T_B$ 上的“通常”拟凝聚层，见引理
\ref{adequate-lemma-same-cohomology-adequate}。）
断言：函子
$$
B \longmapsto H^q(T_B, \mathcal{F}_B)
$$
是适足的。我们将按通常的切分 $T$ 的方法证明本引理。

\medskip\noindent
情形 I：$T$ 为仿射。在此情形所有概形 $T_B$ 都仿射，且对所有 $q \geq 1$，
$H^q(T_B, \mathcal{F}_B) = 0$。
函子 $B \mapsto H^0(T_B, \mathcal{F}_B)$ 由引理
\ref{adequate-lemma-pushforward-adequate} 可知适足。

\medskip\noindent
情形 II：$T$ 分离。令 $n$ 为覆盖 $T$ 所需仿射概形的最小数目。
对 $n$ 归纳。基例为情形 I。取仿射开覆盖 $T = V_1 \cup \ldots \cup V_n$。
令 $V = V_1 \cup \ldots \cup V_{n - 1}$，$U = V_n$。注意
$$
U \cap V = (V_1 \cap V_n) \cup \ldots \cup (V_{n - 1} \cap V_n)
$$
也是 $n - 1$ 个仿射开集的并，因为 $T$ 分离；见《概形》引理
\ref{schemes-lemma-characterize-separated}。
对每个 $B$，基变换 $U_B$、$V_B$ 及
$(U \cap V)_B = U_B \cap V_B$ 的行为相同。因此对每个 $B$ 有长正合列
$$
0 \to
H^0(T_B, \mathcal{F}_B) \to
H^0(U_B, \mathcal{F}_B) \oplus H^0(V_B, \mathcal{F}_B) \to
H^0((U \cap V)_B, \mathcal{F}_B) \to
H^1(T_B, \mathcal{F}_B) \to \ldots
$$
其关于 $B$ 函子性见《上同调》引理
\ref{cohomology-lemma-mayer-vietoris}。
由归纳假设，函子
$B \mapsto H^q(U_B, \mathcal{F}_B)$、
$B \mapsto H^q(V_B, \mathcal{F}_B)$ 以及
$B \mapsto H^q((U \cap V)_B, \mathcal{F}_B)$
都是适足的。利用引理 \ref{adequate-lemma-kernel-adequate} 与
\ref{adequate-lemma-cokernel-adequate}，可见函子
$B \mapsto H^q(T_B, \mathcal{F}_B)$
位于外项均适足的短正合列中。因此断言由引理
\ref{adequate-lemma-extension-adequate} 得出。

\medskip\noindent
情形 III：一般的拟紧拟分离情形。仍按覆盖 $T$ 所需仿射概形数目 $n$ 归纳。
基例 $n = 1$ 是情形 I。取仿射开覆盖 $T = V_1 \cup \ldots \cup V_n$。
令 $V = V_1 \cup \ldots \cup V_{n - 1}$，$U = V_n$。注意，由于 $T$ 拟分离，
$U \cap V$ 是某个仿射概形的拟紧开子集，故情形 II 可应用于它。
其余论证与上段完全相同，略去。
\end{proof}







\section{寄生适足模}
\label{adequate-section-parasitic}

\noindent
本节开始比较概形 $S$ 上的适足模与拟凝聚模。回顾，有函子
$u : \QCoh(\mathcal{O}_S) \to \textit{Adeq}(\mathcal{O})$
以及
$v : \textit{Adeq}(\mathcal{O}) \to \QCoh(\mathcal{O}_S)$，
它们满足伴随关系
$$
\SheafHom_{\QCoh(\mathcal{O}_S)}(\mathcal{F}, v\mathcal{G})
=
\SheafHom_{\textit{Adeq}(\mathcal{O})}(u\mathcal{F}, \mathcal{G})
$$
并且对每个拟凝聚层 $\mathcal{F}$，$\mathcal{F} \to vu\mathcal{F}$ 是同构，见
注记 \ref{adequate-remark-compare}。
因此 $u$ 是全忠实嵌入，我们可将 $\QCoh(\mathcal{O}_S)$
识别为 $\textit{Adeq}(\mathcal{O})$ 的满子范畴。
函子 $v$ 正合，但一般而言 $u$ 不左正合。
$v$ 的核是寄生适足模构成的子范畴。

\medskip\noindent
在《下降》定义 \ref{descent-definition-parasitic} 中给出了寄生模的定义。
对于适足模，该概念与所选拓扑无关。

\begin{lemma}
\label{adequate-lemma-parasitic-adequate}
设 $S$ 为概形，且 $\mathcal{F}$ 是 $(\Sch/S)_\tau$ 上的适足
$\mathcal{O}$-模。以下条件等价：
\begin{enumerate}
\item $v\mathcal{F} = 0$；
\item $\mathcal{F}$ 是寄生的；
\item $\mathcal{F}$ 关于 $\tau$-拓扑是寄生的；
\item 对所有开子集 $U \subset S$，有 $\mathcal{F}(U) = 0$；
\item 存在仿射开覆盖 $S = \bigcup U_i$，使得对所有 $i$ 有
$\mathcal{F}(U_i) = 0$。
\end{enumerate}
\end{lemma}

\begin{proof}
由定义，(2) $\Rightarrow$ (3) $\Rightarrow$ (4) $\Rightarrow$ (5) 立即成立。
假设 (5)。设 $S = \bigcup U_i$ 是仿射开覆盖，且对所有 $i$ 有
$\mathcal{F}(U_i) = 0$。令 $V \to S$ 为平坦态射。存在仿射开覆盖
$V = \bigcup V_j$，使每个 $V_j$ 映入某个 $U_i$。由于态射
$V_j \to S$ 平坦，$V_j \to U_i$ 也平坦。因此相应的环映射
$A_i = \mathcal{O}(U_i) \to \mathcal{O}(V_j) = B_j$ 平坦。于是由引理
\ref{adequate-lemma-adequate-local} 及引理 \ref{adequate-lemma-adequate-flat}，有
$\mathcal{F}(U_i) \otimes_{A_i} B_j \to \mathcal{F}(V_j)$
为同构。因此 $\mathcal{F}(V_j) = 0$。由于 $\mathcal{F}$ 是 Zariski 拓扑下的层，
可得 $\mathcal{F}(V) = 0$。这样就证明了 (5) 蕴含 (2)。

\medskip\noindent
这证明了 (2)、(3)、(4)、(5) 等价。由于 (1) 与 (3) 等价（见注记
\ref{adequate-remark-compare}），故五个条件彼此等价。
\end{proof}

\noindent
设 $S$ 为概形。寄生适足模构成 $\textit{Adeq}(\mathcal{O})$ 的 Serre 子范畴。
其商范畴是拟凝聚模范畴。

\begin{lemma}
\label{adequate-lemma-adequate-by-parasitic}
设 $S$ 为概形。寄生适足模的子范畴
$\mathcal{C} \subset \textit{Adeq}(\mathcal{O})$ 是 Serre 子范畴。此外，函子 $v$
诱导范畴等价
$$
\textit{Adeq}(\mathcal{O}) / \mathcal{C} = \QCoh(\mathcal{O}_S).
$$
\end{lemma}

\begin{proof}
范畴 $\mathcal{C}$ 是正合函子
$v : \textit{Adeq}(\mathcal{O}) \to \QCoh(\mathcal{O}_S)$ 的核，见引理
\ref{adequate-lemma-parasitic-adequate}。
因此由《同调》引理 \ref{homology-lemma-kernel-exact-functor}，它是 Serre 子范畴。
由《同调》引理 \ref{homology-lemma-serre-subcategory-is-kernel}，得到诱导的正合函子
$\overline{v} :
\textit{Adeq}(\mathcal{O}) / \mathcal{C}
\to
\QCoh(\mathcal{O}_S)$。
由于 $u$ 是 $v$ 的右逆，立即可见 $\overline{v}$ 本质满射。
由《同调》引理 \ref{homology-lemma-quotient-by-kernel-exact-functor}，
$\overline{v}$ 是忠实的。又由于 $u$ 是 $v$ 的右逆，最终得到
$\overline{v}$ 全忠实。
\end{proof}

\begin{lemma}
\label{adequate-lemma-direct-image-parasitic-adequate}
设 $f : T \to S$ 为拟紧且拟分离的概形态射。对
$(\Sch/T)_\tau$ 上任意寄生适足 $\mathcal{O}_T$-模，推前
$f_*\mathcal{F}$ 及高阶直像 $R^if_*\mathcal{F}$ 都是
$(\Sch/S)_\tau$ 上的寄生适足 $\mathcal{O}_S$-模。
\end{lemma}

\begin{proof}
我们已经在引理 \ref{adequate-lemma-direct-image-adequate} 中看到，这些高阶直像是适足的。
因此只需证明，对任意开集 $\tau$-覆盖 $\{U_i \to S\}$，有
$(R^if_*\mathcal{F})(U_i) = 0$。而由《下降》引理
\ref{descent-lemma-direct-image-parasitic}，$R^if_*\mathcal{F}$ 是寄生的。
\end{proof}














\section{适足模的导出范畴 I}
\label{adequate-section-comparison}


\noindent
设 $S$ 为概形。继续讨论第 \ref{adequate-section-parasitic} 节的内容。
正合函子 $v$ 诱导函子
$$
D(\textit{Adeq}(\mathcal{O}))
\longrightarrow
D(\QCoh(\mathcal{O}_S))
$$
有界版本也同样如此。

\begin{lemma}
\label{adequate-lemma-quotient-easy}
设 $S$ 为概形。令
$\mathcal{C} \subset \textit{Adeq}(\mathcal{O})$ 为由寄生适足模组成的
满子范畴。则
$$
D(\textit{Adeq}(\mathcal{O}))/D_\mathcal{C}(\textit{Adeq}(\mathcal{O}))
= D(\QCoh(\mathcal{O}_S))
$$
有界版本也同样成立。
\end{lemma}

\begin{proof}
立即由《导出范畴》引理
\ref{derived-lemma-quotient-by-serre-easy} 得出。
\end{proof}

\noindent
接下来反向寻找描述，考察函子链
$$
K^+(\QCoh(\mathcal{O}_S))
\longrightarrow
K^+(\textit{Adeq}(\mathcal{O}))
\longrightarrow
D^+(\textit{Adeq}(\mathcal{O}))
\longrightarrow
D^+(\QCoh(\mathcal{O}_S)).
$$
在某些情形下，适足模的导出范畴是拟凝聚模复形的同伦范畴
在通用拟同构处的局部化。设 $S$ 为概形。若拟凝聚
$\mathcal{O}_S$-模复形之间的映射
$\varphi : \mathcal{F}^\bullet \to \mathcal{G}^\bullet$
满足：对每个概形态射 $f : T \to S$，拉回 $f^*\varphi$ 都是拟同构，
则称其为{\it 通用拟同构}。

\begin{lemma}
\label{adequate-lemma-describe-Dplus-adequate}
设 $U = \Spec(A)$ 为仿射概形。
有界下导出范畴 $D^+(\textit{Adeq}(\mathcal{O}))$ 是
$K^+(\QCoh(\mathcal{O}_U))$ 在通用拟同构乘法子集处的局部化。
\end{lemma}

\begin{proof}
若 $\varphi : \mathcal{F}^\bullet \to \mathcal{G}^\bullet$
是拟凝聚 $\mathcal{O}_U$-模复形之间的态射，则
$u\varphi : u\mathcal{F}^\bullet \to
u\mathcal{G}^\bullet$ 是拟同构，当且仅当 $\varphi$ 是通用拟同构。
因此，由《导出范畴》引理
\ref{derived-lemma-triangle-functor-localize}，通用拟同构的集合 $S$
是与三角结构相容的饱和乘法系统。
于是 $S^{-1}K^+(\QCoh(\mathcal{O}_U))$ 存在且为三角范畴，见
《导出范畴》命题
\ref{derived-proposition-construct-localization}。
由《导出范畴》引理
\ref{derived-lemma-universal-property-localization}，得到典范函子
$can : S^{-1}K^+(\QCoh(\mathcal{O}_U)) \to
D^{+}(\textit{Adeq}(\mathcal{O}))$。

\medskip\noindent
注意，几乎按定义，$U$ 上每个适足模都能嵌入某个拟凝聚层，见
引理 \ref{adequate-lemma-adequate-characterize}。因此由《导出范畴》引理
\ref{derived-lemma-subcategory-right-resolution}，对任意
$\mathcal{F}^\bullet \in \Ob(K^+(\textit{Adeq}(\mathcal{O})))$，
存在拟同构
$\mathcal{F}^\bullet \to u\mathcal{G}^\bullet$，
其中 $\mathcal{G}^\bullet \in \Ob(K^+(\QCoh(\mathcal{O}_U)))$。
这证明了 $can$ 本质满射。

\medskip\noindent
类似地，设 $\mathcal{F}^\bullet$ 与 $\mathcal{G}^\bullet$
是拟凝聚 $\mathcal{O}_U$-模的有界下复形。
$D^+(\textit{Adeq}(\mathcal{O}))$ 中这两个对象之间的态射，
由一对
$f : u\mathcal{F}^\bullet \to \mathcal{H}^\bullet$ 与
$s : u\mathcal{G}^\bullet \to \mathcal{H}^\bullet$ 组成，其中 $s$
是拟同构。取拟同构
$s' : \mathcal{H}^\bullet \to u\mathcal{E}^\bullet$。
于是
$s' \circ f : \mathcal{F} \to \mathcal{E}^\bullet$ 与通用拟同构
$s' \circ s : \mathcal{G}^\bullet \to \mathcal{E}^\bullet$
给出 $S^{-1}K^{+}(\QCoh(\mathcal{O}_U))$ 中的态射，
其映到给定的态射。这证明了全忠实中的“满”部分。
“忠实”部分同理可证。
\end{proof}

\begin{lemma}
\label{adequate-lemma-right-adjoint-adequate}
设 $U = \Spec(A)$ 为仿射概形。包含函子
$$
\textit{Adeq}(\mathcal{O}) \to
\textit{Mod}((\Sch/U)_\tau, \mathcal{O})
$$
具有右伴随 $A$\footnote{这就是“适足化子”。}。
此外，对每个适足模 $\mathcal{F}$，伴随映射
$A(\mathcal{F}) \to \mathcal{F}$ 都是同构。
\end{lemma}

\begin{proof}
由《拓扑》引理
\ref{topologies-lemma-affine-big-site-fppf}
（其他拓扑同理），我们可以在 $(\textit{Aff}/U)_\tau$ 上的
$\mathcal{O}$-模中工作。
记 $\mathcal{P}$ 为 $\textit{Alg}_A$ 上取值于模的函子范畴，
记 $\mathcal{A}$ 为 $\textit{Alg}_A$ 上的适足函子范畴。记
$i : \mathcal{A} \to \mathcal{P}$ 为包含函子。记
$Q : \mathcal{P} \to \mathcal{A}$ 为引理
\ref{adequate-lemma-adjoint} 的构造。
有交换图
\begin{equation}
\label{adequate-equation-categories}
\vcenter{
\xymatrix{
\textit{Adeq}(\mathcal{O}) \ar[r]_-k \ar@{=}[d] &
\textit{Mod}((\textit{Aff}/U)_\tau, \mathcal{O}) \ar[r]_-j &
\textit{PMod}((\textit{Aff}/U)_\tau, \mathcal{O}) \ar@{=}[d] \\
\mathcal{A} \ar[rr]^-i & & \mathcal{P}
}
}
\end{equation}
左侧纵向等号是引理 \ref{adequate-lemma-adequate-affine}，
右侧纵向等号已在第 \ref{adequate-section-quasi-coherent} 节说明。
定义 $A(\mathcal{F}) = Q(j(\mathcal{F}))$。
由于 $j$ 是全忠实函子，立即得到 $A$ 是包含函子 $k$ 的右伴随。
同样，由于 $k$ 也是全忠实函子，最后的断言形式上成立。
\end{proof}

\noindent
函子 $A$ 是右伴随，因而是左正合的。由于包含函子是正合的，见
引理 \ref{adequate-lemma-abelian-adequate}，我们得到 $A$ 将内射对象变为内射对象，
且范畴 $\textit{Adeq}(\mathcal{O})$ 有足够多的内射对象，见
《同调》引理 \ref{homology-lemma-adjoint-enough-injectives}
以及
《内射对象》定理
\ref{injectives-theorem-sheaves-modules-injectives}。
这也可由
(\ref{adequate-equation-categories}) 中的等价性、
以及引理 \ref{adequate-lemma-enough-injectives} 得到。

\begin{lemma}
\label{adequate-lemma-RA-zero}
设 $U = \Spec(A)$ 为仿射概形。对
$\textit{Adeq}(\mathcal{O})$ 的任意对象 $\mathcal{F}$，都有
$R^pA(\mathcal{F}) = 0$（对所有 $p > 0$），其中 $A$ 如
引理 \ref{adequate-lemma-right-adjoint-adequate} 中所定义。
\end{lemma}

\begin{proof}
使用引理 \ref{adequate-lemma-right-adjoint-adequate} 证明中的记号，
在 $(\textit{Aff}/U)_\tau$ 上的 $\mathcal{O}$-模范畴中取内射分解
$k(\mathcal{F}) \to \mathcal{I}^\bullet$。
由
《层上同调》引理 \ref{sites-cohomology-lemma-include-modules}
以及
引理 \ref{adequate-lemma-same-cohomology-adequate}，
复形 $j(\mathcal{I}^\bullet)$ 是正合的。
另一方面，由《层上同调》引理
\ref{sites-cohomology-lemma-injective-module-injective-presheaf}，
每个 $j(\mathcal{I}^n)$ 都是模预层范畴中的内射对象。
因此 $R^pA(\mathcal{F}) = R^pQ(j(k(\mathcal{F})))$。
结论现由引理 \ref{adequate-lemma-RQ-zero} 得出。
\end{proof}

\noindent
设 $S$ 为概形。根据第 \ref{adequate-section-adequate} 节的讨论，嵌入
$\textit{Adeq}(\mathcal{O}) \subset
\textit{Mod}((\Sch/S)_\tau, \mathcal{O})$
将 $\textit{Adeq}(\mathcal{O})$ 显示为所有 $\mathcal{O}$-模范畴中的弱 Serre 子范畴。记
$$
D_{\textit{Adeq}}(\mathcal{O}) \subset
D(\mathcal{O}) = D(\textit{Mod}((\Sch/S)_\tau, \mathcal{O}))
$$
为上同调层适足的复形所组成的三角子范畴，见
《导出范畴》第 \ref{derived-section-triangulated-sub} 节。
于是得到典范函子
$$
D(\textit{Adeq}(\mathcal{O}))
\longrightarrow
D_{\textit{Adeq}}(\mathcal{O})
$$
见《导出范畴》公式 (\ref{derived-equation-compare})。

\begin{lemma}
\label{adequate-lemma-bounded-below}
若 $U = \Spec(A)$ 为仿射概形，则上述函子的有界下版本
\begin{equation}
\label{adequate-equation-compare-bounded-adequate}
D^+(\textit{Adeq}(\mathcal{O}))
\longrightarrow
D^+_{\textit{Adeq}}(\mathcal{O})
\end{equation}
是一个等价。
\end{lemma}

\begin{proof}
令 $A : \textit{Mod}(\mathcal{O}) \to \textit{Adeq}(\mathcal{O})$
为在引理 \ref{adequate-lemma-right-adjoint-adequate} 中构造的包含函子的右伴随。
由于 $A$ 左正合且 $\textit{Mod}(\mathcal{O})$ 有足够多内射对象，
$A$ 具有右导出函子
$RA : D^+_{\textit{Adeq}}(\mathcal{O}) \to D^+(\textit{Adeq}(\mathcal{O}))$。
我们断言 $RA$ 是
(\ref{adequate-equation-compare-bounded-adequate}) 的拟逆。
关键事实是：若 $\mathcal{F}$ 为适足模，则伴随映射
$\mathcal{F} \to RA(\mathcal{F})$ 由引理 \ref{adequate-lemma-RA-zero} 是拟同构。
 
\medskip\noindent
也就是说，要完整证明本引理，只需证明：
\begin{enumerate}
\item 对 $\mathcal{F}^\bullet \in K^+(\textit{Adeq}(\mathcal{O}))$，
典范映射 $\mathcal{F}^\bullet \to RA(\mathcal{F}^\bullet)$
是拟同构；以及
\item 对 $\mathcal{G}^\bullet \in K^+(\textit{Mod}(\mathcal{O}))$，
典范映射 $RA(\mathcal{G}^\bullet) \to \mathcal{G}^\bullet$
是拟同构。
\end{enumerate}
(1) 与 (2) 都由关键事实以及谱序列论证得到，所用谱序列是
《导出范畴》引理 \ref{derived-lemma-two-ss-complex-functor} 中的某一个谱序列。
细节从略。
\end{proof}

\begin{lemma}
\label{adequate-lemma-ext-adequate}
设 $U = \Spec(A)$ 为仿射概形。
设 $\mathcal{F}$ 与 $\mathcal{G}$ 为适足 $\mathcal{O}$-模。对任意 $i \geq 0$，自然映射
$$
\Ext^i_{\textit{Adeq}(\mathcal{O})}(\mathcal{F}, \mathcal{G})
\longrightarrow
\Ext^i_{\textit{Mod}(\mathcal{O})}(\mathcal{F}, \mathcal{G})
$$
是同构。
\end{lemma}

\begin{proof}
按定义，这些 Ext 群由导出范畴中的 Hom 集计算。
因此立即由引理 \ref{adequate-lemma-bounded-below} 得出。
\end{proof}









\section{纯扩张}
\label{adequate-section-pure}

\noindent
我们希望用代数刻画仿射概形的大站点上拟凝聚层的扩张。
为此引入如下概念。

\begin{definition}
\label{adequate-definition-pure}
设 $A$ 为环。
\begin{enumerate}
\item 若对每个 $A$-模的普遍正合序列
$0 \to K \to M \to N \to 0$，序列
$0 \to \Hom_A(P, K) \to \Hom_A(P, M) \to \Hom_A(P, N) \to 0$
都是正合的，则称 $A$-模 $P$ 为{\it 纯射影}。
\item 若对每个 $A$-模的普遍正合序列
$0 \to K \to M \to N \to 0$，序列
$0 \to \Hom_A(N, I) \to \Hom_A(M, I) \to \Hom_A(K, I) \to 0$
都是正合的，则称 $A$-模 $I$ 为{\it 纯内射}。
\end{enumerate}
\end{definition}

\noindent
我们来刻画纯射影模。

\begin{lemma}
\label{adequate-lemma-pure-projective}
设 $A$ 为环。
\begin{enumerate}
\item 模是纯射影的，当且仅当它是有限表示 $A$-模直和的直和项。
\item 对任意模 $M$，存在一个普遍正合序列
$0 \to N \to P \to M \to 0$。
其中 $P$ 为纯射影模。
\end{enumerate}
\end{lemma}

\begin{proof}
首先注意，由《代数》定理
\ref{algebra-theorem-universally-exact-criteria}，
有限表示 $A$-模是纯射影的。
因此有限表示 $A$-模直和的直和项确实是纯射影的。设 $M$ 为任意
$A$-模。将其写成有限表示 $A$-模的滤过余极限
$M = \colim_{i \in I} P_i$。考虑序列
$$
0 \to N \to \bigoplus P_i \to M \to 0.
$$
对任意有限表示 $A$-模 $P$，映射
$\Hom_A(P, \bigoplus P_i) \to \Hom_A(P, M)$
是满射，因为任意 $P \to M$ 的映射都经过某个 $P_i$ 分解。
因此由
《代数》定理 \ref{algebra-theorem-universally-exact-criteria}，
该序列是普遍正合的。这证明了 (2)。
若 $M$ 现在是纯射影的，则该序列分裂，从而可见
$M$ 是 $\bigoplus P_i$ 的直和项。
\end{proof}

\noindent
我们来刻画纯内射模。

\begin{lemma}
\label{adequate-lemma-pure-injective}
设 $A$ 为环。对任意 $A$-模 $M$，置
$M^\vee = \Hom_\mathbf{Z}(M, \mathbf{Q}/\mathbf{Z})$。
\begin{enumerate}
\item 对任意 $A$-模 $M$，$A$-模 $M^\vee$ 是纯内射的。
\item $A$-模 $I$ 是纯内射的，当且仅当映射
$I \to (I^\vee)^\vee$ 分裂。
\item 对任意模 $M$，存在普遍正合序列
$0 \to M \to I \to N \to 0$，其中 $I$ 是纯内射的。
\end{enumerate}
\end{lemma}

\begin{proof}
以下使用函子 $M \mapsto M^\vee$ 的下列性质，见
《代数》补充内容第 \ref{more-algebra-section-injectives-modules} 节，
不再另行说明。第 (1) 项成立，因为
$\Hom_A(N, M^\vee) =
\Hom_\mathbf{Z}(N \otimes_A M, \mathbf{Q}/\mathbf{Z})$
且 $\mathbf{Q}/\mathbf{Z}$ 在阿贝尔群范畴中是内射对象。
因此若 $I \to (I^\vee)^\vee$ 分裂，则 $I$ 是纯内射的。
我们断言，对任意 $A$-模 $M$，评价映射
$ev : M \to (M^\vee)^\vee$ 是普遍单射。
为此注意，$ev^\vee : ((M^\vee)^\vee)^\vee \to M^\vee$
有右逆，即 $ev' : M^\vee \to ((M^\vee)^\vee)^\vee$。
于是对任意 $A$-模 $N$，将正合忠实函子
${}^\vee$ 应用于映射
$N \otimes_A M \to N \otimes_A (M^\vee)^\vee$，得到
$$
\Hom_A(N, ((M^\vee)^\vee)^\vee) =
\Big(N \otimes_A (M^\vee)^\vee\Big)^\vee
\to
\Big(N \otimes_A M\Big)^\vee =
\Hom_A(N, M^\vee)
$$
该映射由右逆的存在性而为满射。断言得证。
该断言蕴含 (3) 以及 (2) 中条件的必要性。
\end{proof}

\noindent
在继续之前，先作如下观察；本节其余部分将频繁使用它。

\begin{lemma}
\label{adequate-lemma-split-universally-exact-sequence}
设 $A$ 为环。
\begin{enumerate}
\item 设 $L \to M \to N$ 为普遍正合序列（其中对象为 $A$-模），并令
$K = \Im(M \to N)$。
则 $K \to N$ 是普遍单射。
\item 任意普遍正合复形都可以拆分为普遍正合短正合序列。
\end{enumerate}
\end{lemma}

\begin{proof}
(1) 的证明。对任意 $A$-模 $T$，由张量积 $\otimes$ 的右正合性，序列
$L \otimes_A T \to M \otimes_A T \to K \otimes_A T \to 0$ 是正合的。
由假设，序列
$L \otimes_A T \to M \otimes_A T \to N \otimes_A T$ 是正合的。
合并这两点可知 $K \otimes_A T \to N \otimes_A T$ 是单射。

\medskip\noindent
(2) 的含义如下：设 $M^\bullet$ 为 $A$-模的普遍正合复形。置
$K^i = \Ker(d^i) \subset M^i$。
则短正合序列
$0 \to K^i \to M^i \to K^{i + 1} \to 0$
都是普遍正合的。这立即由 (1) 得出。
\end{proof}

\begin{definition}
\label{adequate-definition-pure-resolution}
设 $A$ 为环，设 $M$ 为 $A$-模。
\begin{enumerate}
\item {\it 纯射影分解} $P_\bullet \to M$ 是普遍正合序列
$$
\ldots \to P_1 \to P_0 \to M \to 0
$$
且每个 $P_i$ 都是纯射影模。
\item {\it 纯内射分解} $M \to I^\bullet$ 是普遍正合序列
$$
0 \to M \to I^0 \to I^1 \to \ldots
$$
且每个 $I^i$ 都是纯内射模。
\end{enumerate}
\end{definition}

\noindent
这些分解在所有普遍正合的左分解或右分解类中具有通常的唯一性性质。

\begin{lemma}
\label{adequate-lemma-pure-projective-resolutions}
设 $A$ 为环。
\begin{enumerate}
\item 任意 $A$-模都有纯射影分解。
\end{enumerate}
设 $M \to N$ 为 $A$-模的态射。
设 $P_\bullet \to M$ 为纯射影分解，
设 $N_\bullet \to N$ 为普遍正合分解。
\begin{enumerate}
\item[(2)] 存在复形态射 $P_\bullet \to N_\bullet$，
诱导给定的映射
$$
M = \Coker(P_1 \to P_0) \to \Coker(N_1 \to N_0) = N
$$
\item[(3)] 两个复形态射 $\alpha, \beta : P_\bullet \to N_\bullet$，
若诱导相同的映射 $M \to N$，则它们同伦。
\end{enumerate}
\end{lemma}

\begin{proof}
(1) 立即由引理 \ref{adequate-lemma-pure-projective} 得出。
在证明 (2) 与 (3) 之前注意，由引理
\ref{adequate-lemma-split-universally-exact-sequence}，
可以将普遍正合复形 $N_\bullet \to N \to 0$
拆分为普遍正合短正合序列
$0 \to K_0 \to N_0 \to N \to 0$
以及 $0 \to K_i \to N_i \to K_{i - 1} \to 0$。

\medskip\noindent
(2) 的证明。由于 $P_0$ 是纯射影的，
可以找到一个映射 $P_0 \to N_0$，提升映射 $P_0 \to M \to N$。
得到诱导映射 $P_1 \to F_0 \to N_0$，其最终落入 $K_0$。
由于 $P_1$ 是纯射影的，可以将其提升为映射 $P_1 \to N_1$。
这又诱导映射 $P_2 \to P_1 \to N_1$，它映到 $N_0$（即映到 $K_1$）为零。
因此可以提升得到映射 $P_2 \to N_2$。如此重复。

\medskip\noindent
(3) 的证明。要证明 $\alpha, \beta$ 同伦，只需证明差
$\gamma = \alpha - \beta$ 同伦于零。注意 $\gamma_0 : P_0 \to N_0$
的像包含于 $K_0$。因此可以将 $\gamma_0$ 提升为映射
$h_0 : P_0 \to N_1$。考虑映射
$\gamma_1' = \gamma_1 - h_0 \circ d_{P, 1} : P_1 \to N_1$。
由 $h_0$ 的选取可知 $\gamma_1'$ 的像包含于 $K_1$。
由于 $P_1$ 是纯射影的，可以将 $\gamma_1'$ 提升为映射
$h_1 : P_1 \to N_2$。此时有
$\gamma_1 = h_0 \circ d_{F, 1} + d_{N, 2} \circ h_1$。重复上述过程。
\end{proof}

\begin{lemma}
\label{adequate-lemma-pure-injective-resolutions}
设 $A$ 为环。
\begin{enumerate}
\item 任意 $A$-模都有纯内射分解。
\end{enumerate}
设 $M \to N$ 为 $A$-模的态射。
设 $M \to M^\bullet$ 为普遍正合分解，且
设 $N \to I^\bullet$ 为纯内射分解。
\begin{enumerate}
\item[(2)] 存在复形态射 $M^\bullet \to I^\bullet$，
诱导给定的映射
$$
M = \Ker(M^0 \to M^1) \to \Ker(I^0 \to I^1) = N
$$
\item[(3)] 两个复形态射 $\alpha, \beta : M^\bullet \to I^\bullet$，
若诱导相同的映射 $M \to N$，则它们同伦。
\end{enumerate}
\end{lemma}

\begin{proof}
本引理是引理 \ref{adequate-lemma-pure-projective-resolutions} 的对偶命题。
证明相同，只需反转所有箭头。
\end{proof}

\noindent
利用以上材料可如下定义纯扩张群。设 $A$ 为环，设 $M$、$N$ 为 $A$-模。
取纯内射分解 $N \to I^\bullet$。由
引理 \ref{adequate-lemma-pure-injective-resolutions}，复形
$$
\Hom_A(M, I^\bullet)
$$
在同伦意义下良定义。因此其第 $i$ 个上同调模
是 $M$ 与 $N$ 的良定义不变量。

\begin{definition}
\label{adequate-definition-pure-ext}
设 $A$ 为环，设 $M$、$N$ 为 $A$-模。
第 $i$ 个{\it 纯扩张模} $\text{Pext}^i_A(M, N)$
定义为复形 $\Hom_A(M, I^\bullet)$ 的第 $i$ 个上同调模，
其中 $I^\bullet$ 是 $N$ 的纯内射分解。
\end{definition}

\noindent
警告：$A$-模的正合序列不一定导出纯扩张群的长正合序列。
（为此需要普遍正合序列。）
我们收集一些由以上材料显然得到的事实。

\begin{lemma}
\label{adequate-lemma-facts-pext}
设 $A$ 为环。
\begin{enumerate}
\item $\text{Pext}^i_A(M, N) = 0$ 对 $i > 0$ 成立，只要 $N$ 为纯内射模；
\item $\text{Pext}^i_A(M, N) = 0$ 对 $i > 0$ 成立，只要 $M$ 为纯射影模；
特别地，当 $M$ 为有限表示的 $A$-模时也成立；
\item $\text{Pext}^i_A(M, N)$ 也是复形
$\Hom_A(P_\bullet, N)$ 的第 $i$ 个上同调模，其中 $P_\bullet$
是 $M$ 的纯射影分解。
\end{enumerate}
\end{lemma}

\begin{proof}
为证明 (3)，考虑双复形
$$
A^{\bullet, \bullet} = \Hom_A(P_\bullet, I^\bullet)
$$
其每一行在除第 $0$ 次之外的次数上正合，而该行的上同调为
$\Hom_A(M, I^q)$。其每一列除第 $0$ 次外均为正合，
而该列的上同调为 $\Hom_A(P_p, N)$。
因此，与该复形相关的两个谱序列（见
《同调》第 \ref{homology-section-double-complex} 节）
退化，从而得到所述等式。
\end{proof}





\section{大站点上拟凝聚层的高阶 Ext}
\label{adequate-section-big}

\noindent
事实证明，与纯内射模 $I$ 关联的取值于模的函子
$\underline{I}$，在 $\textit{Alg}_A$ 上的适足函子范畴中给出一个内射对象。
警告：纯射影模并不一定在适足函子范畴中给出射影对象。
不过我们确实有大量射影对象，即线性适足函子。

\begin{lemma}
\label{adequate-lemma-pure-injective-injective-adequate}
设 $A$ 为环。
设 $\mathcal{A}$ 为 $\textit{Alg}_A$ 上适足函子的范畴。
$\mathcal{A}$ 的内射对象恰好是函子
$\underline{I}$，其中 $I$ 为纯内射 $A$-模。
\end{lemma}

\begin{proof}
设 $I$ 为 $\mathcal{A}$ 的内射对象。
取某个 $A$-模 $M$ 的嵌入 $I \to \underline{M}$。
由于 $I$ 是内射对象，可知 $\underline{M} = I \oplus F$，其中 $F$ 为取值于模的函子。
于是 $M = I(A) \oplus F(A)$，从而
$I = \underline{I(A)}$。因此任意内射对象
都形如某个 $A$-模 $I$ 的 $\underline{I}$。
显然 $I$ 必须是纯内射的，因为每个普遍正合序列
$0 \to M \to N \to L \to 0$
都给出 $\mathcal{A}$ 中的正合序列
$0 \to \underline{M} \to \underline{N} \to \underline{L} \to 0$。

\medskip\noindent
最后，设 $I$ 为纯内射 $A$-模。取嵌入
$\underline{I} \to J$，其中 J 为 $\mathcal{A}$ 的内射对象（见
引理 \ref{adequate-lemma-enough-injectives}）。
由上文可知 $J = \underline{I'}$，其中某个 $A$-模 $I'$ 为纯内射。
由于 $\underline{I} \to \underline{I'}$ 是内射的，
映射 $I \to I'$ 是普遍单射。由 $I$ 的假设，该映射分裂。
因此 $\underline{I}$ 是 $J = \underline{I'}$ 的直和项，
从而是范畴 $\mathcal{A}$ 的内射对象。
\end{proof}

\noindent
设 $U = \Spec(A)$ 为仿射概形，设 $M$ 为 $A$-模。
我们用记号 $M^a$ 表示 $(\Sch/U)_\tau$ 上与 $U$ 上拟凝聚层
$\widetilde{M}$ 关联的 $\mathcal{O}$-模拟凝聚层。
现在记号已齐备，可以表述如下引理。

\begin{lemma}
\label{adequate-lemma-big-ext}
设 $U = \Spec(A)$ 为仿射概形。设 $M$、$N$ 为 $A$-模。
对所有 $i$，都有典范同构
$$
\Ext^i_{\textit{Mod}(\mathcal{O})}(M^a, N^a) = \text{Pext}^i_A(M, N)
$$
且该同构关于 $M$ 与 $N$ 具有函子性。
\end{lemma}

\begin{proof}
我们构造一个从右向左的典范箭头。若
$N \to I^\bullet$ 是纯内射分解，则
$M^a \to (I^\bullet)^a$ 是（适足的）$\mathcal{O}$-模的正合复形。
因此 $\text{Pext}^i_A(M, N)$ 的任意元素
都给出 $D(\mathcal{O})$ 中的映射 $N^a \to M^a[i]$，即
左侧群的一个元素。

\medskip\noindent
为证明该映射是同构，注意可将
$\Ext^i_{\textit{Mod}(\mathcal{O})}(M^a, N^a)$ 替换为
$\Ext^i_{\textit{Adeq}(\mathcal{O})}(M^a, N^a)$，见
引理 \ref{adequate-lemma-ext-adequate}。
令 $\mathcal{A}$ 为 $\textit{Alg}_A$ 上适足函子的范畴。
我们已经看到 $\mathcal{A}$ 等价于
$\textit{Adeq}(\mathcal{O})$，见
引理 \ref{adequate-lemma-adequate-affine}；另见
引理 \ref{adequate-lemma-right-adjoint-adequate} 的证明。
因此现在只需证明
$$
\Ext^i_\mathcal{A}(\underline{M}, \underline{N}) =
\text{Pext}^i_A(M, N)
$$
然而，这由引理
\ref{adequate-lemma-pure-injective-injective-adequate} 立即可见，
因为纯内射分解 $N \to I^\bullet$ 恰好对应于
$\mathcal{A}$ 中 $\underline{N}$ 的内射分解。
\end{proof}







\section{适足模的导出范畴 II}
\label{adequate-section-derived-categories}

\noindent
设 $S$ 为概形。记 $\mathcal{O}_S$ 为 $S$ 的结构层，
记 $\mathcal{O}$ 为大站点 $(\Sch/S)_\tau$ 的结构层。
在《下降》备注 \ref{descent-remark-change-topologies-ringed} 中，
我们构造了带环空间的站点态射
\begin{equation}
\label{adequate-equation-compare-big-small}
f :
((\Sch/S)_\tau, \mathcal{O})
\longrightarrow
(S_{Zar}, \mathcal{O}_S).
\end{equation}
在前几节中我们看到，函子
$f_* : \textit{Mod}(\mathcal{O}) \to \textit{Mod}(\mathcal{O}_S)$
将适足层变为拟凝聚层，并诱导正合函子
$v : \textit{Adeq}(\mathcal{O}) \to \QCoh(\mathcal{O}_S)$；事实上，
$f_* = v$ 诱导等价
$\textit{Adeq}(\mathcal{O})/\mathcal{C} \to \QCoh(\mathcal{O}_S)$，
其中 $\mathcal{C}$ 为寄生适足模构成的子范畴。
此外，函子 $f^*$ 将拟凝聚模变为适足模，并诱导函子
$u : \QCoh(\mathcal{O}_S) \to \textit{Adeq}(\mathcal{O})$，
它是 $v$ 的左伴随。
\medskip\noindent
$D_{\textit{Adeq}}(\mathcal{O})$ 与 $D_\QCoh(S)$ 之间也存在非常相似的关系。
首先说明范畴 $D_{\textit{Adeq}}(\mathcal{O})$
与所选拓扑无关。

\begin{remark}
\label{adequate-remark-D-adeq-independence-topology}
设 $S$ 为概形。
设 $\tau, \tau' \in \{Zar, \etale, smooth, syntomic, fppf\}$。
记 $\mathcal{O}_\tau$，resp.\ $\mathcal{O}_{\tau'}$ 为结构层 $\mathcal{O}$
在 $(\Sch/S)_\tau$，resp.\ $(\Sch/S)_{\tau'}$ 上视为层所得的对象。
则 $D_{\textit{Adeq}}(\mathcal{O}_\tau)$ 与
$D_{\textit{Adeq}}(\mathcal{O}_{\tau'})$ 典范同构。
这由《层上同调》引理
\ref{sites-cohomology-lemma-compare-topologies-derived-adequate-modules} 得出。
具体地，假设 $\tau$ 比拓扑 $\tau'$ 更强，令
$\mathcal{C} = (\Sch/S)_{fppf}$，令 $\mathcal{B}$ 为
$S$ 上仿射概形的集合。上文已经看到了假设 (1) 与 (2)。
假设 (3) 显然成立，假设 (4) 由引理
\ref{adequate-lemma-same-cohomology-adequate} 得出。
\end{remark}

\begin{remark}
\label{adequate-remark-D-adeq-and-D-QCoh}
设 $S$ 为概形。态射 $f$ 见
(\ref{adequate-equation-compare-big-small})，它诱导伴随函子
$Rf_* : D_{\textit{Adeq}}(\mathcal{O}) \to D_\QCoh(S)$
和
$Lf^* : D_\QCoh(S) \to D_{\textit{Adeq}}(\mathcal{O})$。
此外 $Rf_* Lf^* \cong \text{id}_{D_\QCoh(S)}$。

\medskip\noindent
下面略述证明。由
备注 \ref{adequate-remark-D-adeq-independence-topology}，
可假设拓扑 $\tau$ 为 Zariski 拓扑。
我们将使用 $\mathcal{O}$-模上无界全导出函子
$Lf^*$ 与 $Rf_*$ 的存在性及其伴随性，见
《层上同调》引理 \ref{sites-cohomology-lemma-adjoint}。
此时 $f_*$ 只是 $(\Sch/S)_{Zar}$ 到其子范畴
$S_{Zar}$ 的限制。因此显然
$Rf_* = f_*$ 诱导
$Rf_* : D_{\textit{Adeq}}(\mathcal{O}) \to D_\QCoh(S)$。
设 $\mathcal{G}^\bullet$ 为 $D_\QCoh(S)$ 中的对象。
可取一个系统
$\mathcal{K}_1^\bullet \to \mathcal{K}_2^\bullet \to \ldots$
以及图表，其中各复形为有界上、平坦的 $\mathcal{O}_S$-模复形，
转移映射逐项分裂单射：
$$
\xymatrix{
\mathcal{K}_1^\bullet \ar[d] \ar[r] &
\mathcal{K}_2^\bullet \ar[d] \ar[r] & \ldots \\
\tau_{\leq 1}\mathcal{G}^\bullet \ar[r] &
\tau_{\leq 2}\mathcal{G}^\bullet \ar[r] & \ldots
}
$$
其具有《导出范畴》引理
\ref{derived-lemma-special-direct-system} 中列出的性质 (1)、(2)、(3)，
其中 $\mathcal{P}$ 是平坦 $\mathcal{O}_S$-模的集合。
于是 $Lf^*\mathcal{G}^\bullet$ 由
$\colim f^*\mathcal{K}_n^\bullet$ 计算，见
《层上同调》引理 \ref{sites-cohomology-lemma-pullback-K-flat} 及
\ref{sites-cohomology-lemma-derived-base-change}
（注意，由于《\'Etale 上同调》引理
\ref{etale-cohomology-lemma-points-fppf}，我们的站点有足够多点）。
我们需要证明 $H^i(Lf^*\mathcal{G}^\bullet) =
\colim H^i(f^*\mathcal{K}_n^\bullet)$ 对每个 $i$ 都是适足的。由
引理 \ref{adequate-lemma-abelian-adequate}，
只需证明每个 $H^i(f^*\mathcal{K}_n^\bullet)$ 都是适足的。

\medskip\noindent
$H^i(f^*\mathcal{K}_n^\bullet)$ 的适足性在 $S$ 上是局部的，
因此可假设 $S = \Spec(A)$ 为仿射概形。由于 $S$ 仿射，
$D_\QCoh(S) = D(\QCoh(\mathcal{O}_S))$，见
《概形的导出范畴》第
\ref{perfect-section-derived-quasi-coherent} 节的讨论。
因此存在拟同构
$\mathcal{F}^\bullet \to \mathcal{K}_n^\bullet$，
其中 $\mathcal{F}^\bullet$ 是平坦拟凝聚模的有界上复形。
于是 $f^*\mathcal{F}^\bullet \to f^*\mathcal{K}_n^\bullet$ 是
拟同构，且 $f^*\mathcal{F}^\bullet$ 的上同调层是适足的。

\medskip\noindent
最后的断言
$Rf_* Lf^* \cong \text{id}_{D_\QCoh(S)}$
由上述函子的显式描述得出。
（直白地说：若 $\mathcal{F}$ 是拟凝聚的且 $p > 0$，则
$L_pf^*\mathcal{F}$ 是寄生适足模。）
\end{remark}


\begin{remark}
\label{adequate-remark-conclusion}
上面的备注 \ref{adequate-remark-D-adeq-and-D-QCoh}
蕴含导出范畴之间的等价
$$
D_{\textit{Adeq}}(\mathcal{O})/D_\mathcal{C}(\mathcal{O})
\longrightarrow
D_\QCoh(S)
$$
其中 $\mathcal{C}$ 为寄生适足模构成的范畴。
具体地，显然 $D_\mathcal{C}(\mathcal{O})$ 是 $Rf_*$ 的核，
因而得到所示函子。对 $D_{\textit{Adeq}}(\mathcal{O})$ 的任意对象 $X$，
映射 $Lf^*Rf_*X \to X$ 在 $D_\QCoh(S)$ 中映为拟同构，因此
$Lf^*Rf_*X \to X$ 是
$D_{\textit{Adeq}}(\mathcal{O})/D_\mathcal{C}(\mathcal{O})$ 中的同构。
最后，对 $D_{\textit{Adeq}}(\mathcal{O})$ 中的对象 $X, Y$，映射
$$
Rf_* :
\Hom_{D_{\textit{Adeq}}(\mathcal{O})/D_\mathcal{C}(\mathcal{O})}(X, Y)
\to
\Hom_{D_\QCoh(S)}(Rf_*X, Rf_*Y)
$$
是双射，因为 $Lf^*$ 给出其逆映射（由上述备注）。
\end{remark}
% Generated mechanically from the frozen dualizing.tex target.
% Do not edit this generated file; edit the bound locale source instead.

% OK, start here.
%

\setcounter{chapter}{46}
\chapter{对偶化复形}



\phantomsection
\label{dualizing-section-phantom}


\section{引言}
\label{dualizing-section-introduction}

\noindent
本章讨论交换代数中的对偶化复形。参考文献为 \cite{RD}。

\medskip\noindent
我们首先讨论本质满射与本质单射、射影覆盖、内射包、
Artinian 环上的对偶，以及剩余域的内射包，并很快导出 Matlis 对偶定理的证明。
参见第 \ref{dualizing-section-essential} 节、
\ref{dualizing-section-injective-modules} 节、
\ref{dualizing-section-projective-cover} 节、
\ref{dualizing-section-injective-hull} 节、
\ref{dualizing-section-artinian} 节、
\ref{dualizing-section-hull-residue-field} 节，以及命题 \ref{dualizing-proposition-matlis}。

\medskip\noindent
接下来三节以很大的普遍性讨论局部上同调，参见
\ref{dualizing-section-bad-local-cohomology}、\ref{dualizing-section-local-cohomology} 和
\ref{dualizing-section-local-cohomology-noetherian} 节。我们在
\ref{dualizing-section-depth} 节将其中部分内容用于讨论深度。另一个应用是：
给定环 $A$ 的有限生成理想 $I$，$A$ 的导出范畴中的
``$I$-完备''对象与 ``$I$-挠''对象等价，参见
\ref{dualizing-section-torsion-and-complete} 节。
若要进一步了解局部上同调（例如依赖局部对偶的有限性定理，见下文），
请参阅《局部上同调》第 \ref{local-cohomology-section-introduction} 节。

\medskip\noindent
本章主要致力于环映射的对偶与对偶化复形。参见
\ref{dualizing-section-trivial}、\ref{dualizing-section-base-change-trivial-duality}、
\ref{dualizing-section-dualizing}、\ref{dualizing-section-dualizing-local}、
\ref{dualizing-section-dimension-function}、\ref{dualizing-section-local-duality}、
\ref{dualizing-section-dualizing-module}、\ref{dualizing-section-CM}、
\ref{dualizing-section-gorenstein}、\ref{dualizing-section-ubiquity-dualizing} 和
\ref{dualizing-section-formal-fibres} 节。
关键定义是：对 Noetherian 环 $A$，对偶化复形
$\omega_A^\bullet$ 是一个对象，满足
$\omega_A^\bullet \in D^{+}(A)$，其上同调模
$H^i(\omega_A^\bullet)$ 都是有限 $A$-模，具有有限内射维数，并且映射
$$
A \longrightarrow R\Hom_A(\omega_A^\bullet, \omega_A^\bullet)
$$
是拟同构。在建立对偶化复形的一些初等性质后，我们证明对偶化复形导出一个维数函数。
接着证明 Grothendieck 局部对偶定理。简要讨论对偶化模与 Cohen--Macaulay 环后，
我们引入 Gorenstein 环，并说明许多熟悉的 Noetherian 环具有对偶化复形。
最后一节应用上述材料，说明 Noetherian 局部环的形式纤维为 Gorenstein 环或局部完全交时，
存在良好的理论。

\medskip\noindent
最后几节描述代数几何中使用的``上惊叹号函子''的代数构造，例如书籍
\cite{RD} 中所用者。本主题将在关于概形对偶的章节中继续讨论。
参见《概形的对偶》第 \ref{duality-section-introduction} 节。







\section{本质满射与本质单射}
\label{dualizing-section-essential}

\noindent
我们主要将在模范畴中工作，但不妨在一般情形下给出定义。

\begin{definition}
\label{dualizing-definition-essential}
设 $\mathcal{A}$ 为阿贝尔范畴。
\begin{enumerate}
\item $\mathcal{A}$ 中的单射 $A \subset B$ 称为{\it 本质的}，
或称 $B$ 是 $A$ 的{\it 本质扩张}，
若每个非零子对象 $B' \subset B$ 与 $A$ 的交均非零。
\item $\mathcal{A}$ 中的满射 $f : A \to B$ 称为{\it 本质的}，
若对每个真子对象 $A' \subset A$ 都有 $f(A') \not = B$。
\end{enumerate}
\end{definition}

\noindent
关于这一概念有如下引理。

\begin{lemma}
\label{dualizing-lemma-essential}
设 $\mathcal{A}$ 为阿贝尔范畴。
\begin{enumerate}
\item 若 $A \subset B$ 与 $B \subset C$ 是本质扩张，则
$A \subset C$ 是本质扩张。
\item 若 $A \subset B$ 是本质扩张且 $C \subset B$ 是子对象，
则 $A \cap C \subset C$ 是本质扩张。
\item 若 $A \to B$ 与 $B \to C$ 是本质满射，则
$A \to C$ 是本质满射。
\item 给定本质满射 $f : A \to B$ 及核为 $K$ 的满射
$A \to C$，则态射 $C \to B/f(K)$ 是本质满射。
\end{enumerate}
\end{lemma}

\begin{proof}
证明略。
\end{proof}

\begin{lemma}
\label{dualizing-lemma-union-essential-extensions}
设 $R$ 为环，$M$ 为 $R$-模，并设 $E = \colim E_i$ 是
$R$-模的滤过余极限。给定相容的 $R$-模本质单射族
$M \to E_i$。则 $M \to E$ 是本质单射。
\end{lemma}

\begin{proof}
这由定义以及滤过余极限的正合性立即推出（代数，引理
\ref{algebra-lemma-directed-colimit-exact}）。
\end{proof}

\begin{lemma}
\label{dualizing-lemma-essential-extension}
设 $R$ 为环，$M \subset N$ 为 $R$-模。下列条件等价：
\begin{enumerate}
\item $M \subset N$ 是本质扩张；
\item 对每个非零 $x \in N$，存在 $f \in R$，使得 $fx \in M$
且 $fx \not = 0$。
\end{enumerate}
\end{lemma}

\begin{proof}
设 (1) 成立，并取非零元素 $x \in N$。由 (1) 有
$Rx \cap M \not = 0$，故 (2) 成立。

\medskip\noindent
设 (2) 成立。令 $N' \subset N$ 为非零子模，并取非零
$x \in N'$。由 (2)，可找到 $f \in R$，使得 $fx \in M$
且 $fx \not = 0$。因此 $N' \cap M \not = 0$。
\end{proof}




\section{内射模}
\label{dualizing-section-injective-modules}

\noindent
下面给出一些关于环上内射模的结果。

\begin{lemma}
\label{dualizing-lemma-product-injectives}
设 $R$ 为环。任意内射 $R$-模的直积都是内射的。
\end{lemma}

\begin{proof}
这是《同调》引理 \ref{homology-lemma-product-injectives} 的特例。
\end{proof}

\begin{lemma}
\label{dualizing-lemma-injective-flat}
设 $R \to S$ 为平坦环映射。若 $E$ 是内射 $S$-模，
则 $E$ 作为 $R$-模也是内射的。
\end{lemma}

\begin{proof}
这是因为由《代数》引理 \ref{algebra-lemma-adjoint-tensor-restrict}，
以及张量 $S$ 保持正合这一事实，有
$\Hom_R(M, E) = \Hom_S(M \otimes_R S, E)$。
\end{proof}

\begin{lemma}
\label{dualizing-lemma-injective-epimorphism}
设 $R \to S$ 为环的满同态。设 $E$ 为 $S$-模。
若 $E$ 作为 $R$-模是内射的，则 $E$ 作为 $S$-模也是内射的。
\end{lemma}

\begin{proof}
这是因为对任意 $S$-模 $N$ 都有
$\Hom_R(N, E) = \Hom_S(N, E)$，见《代数》引理
\ref{algebra-lemma-epimorphism-modules}。
\end{proof}

\begin{lemma}
\label{dualizing-lemma-hom-injective}
设 $R \to S$ 为环映射。若 $E$ 是内射 $R$-模，
则 $\Hom_R(S, E)$ 是内射 $S$-模。
\end{lemma}

\begin{proof}
这是因为由《代数》引理
\ref{algebra-lemma-adjoint-hom-restrict}，有
$\Hom_S(N, \Hom_R(S, E)) = \Hom_R(N, E)$。
\end{proof}

\begin{lemma}
\label{dualizing-lemma-essential-extensions-in-injective}
设 $R$ 为环。设 $I$ 为内射 $R$-模。设 $E \subset I$ 为子模。
以下条件等价：
\begin{enumerate}
\item $E$ 是内射的；
\item 对所有 $E \subset E' \subset I$，若 $E \subset E'$ 是本质的，
则 $E = E'$。
\end{enumerate}
特别地，$R$-模是内射的，当且仅当每个本质扩张都是平凡的。
\end{lemma}

\begin{proof}
最后的断言由第一项及 $R$-模范畴有足够多内射对象这一事实得出
（《更多代数》第 \ref{more-algebra-section-injectives-modules} 节）。

\medskip\noindent
假设 (1)。令 $E \subset E' \subset I$ 如 (2) 中所述。
则恒等映射 $\text{id}_E : E \to E$ 可延拓为映射
$\alpha : E' \to E$。$\alpha$ 的核必为零，因为它与 $E$ 的交为零，
而 $E'$ 是本质扩张。因此 $E = E'$。

\medskip\noindent
假设 (2)。设 $M \subset N$ 为 $R$-模，并设 $\varphi : M \to E$
为 $R$-模映射。为证明 (1)，需证明 $\varphi$ 可延拓为态射
$N \to E$。考虑所有对
$(M', \varphi')$ 所成的集合 $\mathcal{S}$，其中
$M \subset M' \subset N$，且 $\varphi' : M' \to E$ 为在 $M$ 上
与 $\varphi$ 一致的 $R$-模映射。按规则
$(M', \varphi') \leq (M'', \varphi'')$ 当且仅当
$M' \subset M''$ 且 $\varphi''|_{M'} = \varphi'$，
在 $\mathcal{S}$ 上定义序。显然，$\mathcal{S}$ 的全序子集可取其上界。
由 Zorn 引理，不妨设 $(M, \varphi)$ 是极大元。

\medskip\noindent
选取 $\varphi$ 与包含映射 $E \to I$ 的复合
$\psi : N \to I$ 的延拓。这是可能的，因为 $I$ 是内射的。
若 $\psi(N) \subset E$，则 $\psi$ 即为所需延拓。
若 $\psi(N)$ 不包含于 $E$，则由 (2)，包含
$E \subset E + \psi(N)$ 不是本质的。因此可以找到非零子模
$K \subset E + \psi(N)$，使其与 $E$ 的交为 $0$。
这意味着 $M' = \psi^{-1}(E + K)$ 严格包含 $M$。
于是可利用下式将 $\varphi$ 延拓到 $M'$：
$$
M' \xrightarrow{\psi|_{M'}} E + K \to (E + K)/K = E
$$
这与 $(M, \varphi)$ 的极大性矛盾。
\end{proof}

\begin{example}
\label{dualizing-example-reduced-ring-injective}
设 $R$ 为约化环。设 $\mathfrak p \subset R$ 为极小素理想，
使得 $K = R_\mathfrak p$ 是域（《代数》引理
\ref{algebra-lemma-minimal-prime-reduced-ring}）。
则 $K$ 是内射 $R$-模。事实上，对任意 $R$-模 $M$，有
$\Hom_R(M, K) = \Hom_K(M_\mathfrak p, K)$。
由于局部化是正合函子，且对 $K$-向量空间取对偶是正合函子，
可知 $\Hom_R(-, K)$ 是正合函子，即 $K$ 是内射 $R$-模。
\end{example}

\begin{lemma}
\label{dualizing-lemma-sum-injective-modules}
设 $R$ 为 Noetherian 环。内射模的直和是内射的。
\end{lemma}

\begin{proof}
设 $E_i$ 是由集合 $I$ 参数化的一族内射模。
令 $E = \bigoplus E_i$。为证明 $E$ 是内射的，使用《内射模》引理
\ref{injectives-lemma-criterion-baer}。设
$\varphi : J \to E$ 为从 $R$ 的理想到 $E$ 的模映射。
由于 $J$ 是有限生成 $R$-模（因为 $R$ 是 Noetherian 环），
可找到有限多个元素 $i_1, \ldots, i_r \in I$，使得
$\varphi$ 的像包含于 $\bigoplus_{j = 1, \ldots, r} E_{i_j}$。
于是利用各模 $E_{i_j}$ 的内射性，可将 $\varphi$ 延拓到
$\bigoplus_{j = 1, \ldots, r} E_{i_j}$。
\end{proof}

\begin{lemma}
\label{dualizing-lemma-localization-injective-modules}
设 $R$ 为 Noetherian 环。设 $S \subset R$ 为乘法闭子集。
若 $E$ 是内射 $R$-模，则 $S^{-1}E$ 是内射 $S^{-1}R$-模。
\end{lemma}

\begin{proof}
由于 $R \to S^{-1}R$ 是环的满同态，只需证明 $S^{-1}E$
作为 $R$-模是内射的，见引理 \ref{dualizing-lemma-injective-epimorphism}。
为此使用《内射模》引理 \ref{injectives-lemma-criterion-baer}。
设 $I \subset R$ 为理想，且 $\varphi : I \to S^{-1} E$ 为
$R$-模映射。由于 $I$ 是有限表示 $R$-模（因为 $R$ 是 Noetherian 环），
可找到 $f \in S$ 及 $R$-模映射 $I \to E$，使得
$f\varphi$ 是复合 $I \to E \to S^{-1}E$
（《代数》引理 \ref{algebra-lemma-hom-from-finitely-presented}）。
然后可将 $I \to E$ 延拓为同态 $R \to E$。于是复合
$$
R \to E \to S^{-1}E \xrightarrow{f^{-1}} S^{-1}E
$$
就是 $\varphi$ 到 $R$ 的所需延拓。
\end{proof}

\begin{lemma}
\label{dualizing-lemma-injective-module-divide}
设 $R$ 为 Noetherian 环。设 $I$ 为内射 $R$-模。
\begin{enumerate}
\item 设 $f \in R$。则 $E = \bigcup I[f^n] = I[f^\infty]$
是 $I$ 的内射子模。
\item 设 $J \subset R$ 为理想。则 $J$-幂挠子模
$I[J^\infty]$ 是 $I$ 的内射子模。
\end{enumerate}
\end{lemma}

\begin{proof}
我们使用引理 \ref{dualizing-lemma-essential-extensions-in-injective} 证明 (1)。
设 $E \subset E' \subset I$，且 $E'$ 是 $E$ 的本质扩张。
我们将证明 $E' = E$。否则可取 $x \in E'$ 且 $x \not \in E$。
令 $J = \{ a \in R \mid ax \in E\}$。由于 $R$ 是 Noetherian 环，
可写 $J = (g_1, \ldots, g_t)$，其中 $g_i \in R$。
按定义，$E$ 是 $I$ 中被 $f$ 的某个幂湮灭的元素所成的集合，
故可取整数 $n_i$ 使 $f^{n_i}g_ix = 0$。
令 $n = \mathrm{max}\{ n_i \}$。则 $x' = f^n x$ 属于 $E'$，
但不属于 $E$，且被 $J$ 湮灭。令 $J' = \{ a \in R \mid ax' \in E \}$，
则 $J \subset J'$。反之，$a \in J'$ 当且仅当 $ax' \in E$，
当且仅当对某个 $m \geq 0$ 有 $f^m a x' = 0$。但
$f^m a x' = f^{m + n} a x$，故 $ax \in E$，即 $a \in J$。
于是 $J = J'$。因此 $J = J' = \text{Ann}(x')$，从而 $Rx' \cap E  = 0$。
所以 $E'$ 不是 $E$ 的本质扩张，矛盾。

\medskip\noindent
为证明 (2)，写成 $J = (f_1, \ldots, f_t)$。则
$I[J^\infty]$ 等于
$$
(\ldots((I[f_1^\infty])[f_2^\infty])\ldots)[f_t^\infty]
$$
结论由 (1) 及归纳法得出。
\end{proof}

\begin{lemma}
\label{dualizing-lemma-injective-dimension-over-polynomial-ring}
设 $A$ 为 Noetherian 环。设 $E$ 为内射 $A$-模。
则 $E \otimes_A A[x]$ 作为 $D(A[x])$ 中的对象具有内射振幅 $[0, 1]$。
特别地，$E \otimes_A A[x]$ 作为 $A[x]$-模具有有限内射维数。
\end{lemma}

\begin{proof}
记 $E[x] = E \otimes_A A[x]$。考虑 $A[x]$-模的短正合序列
$$
0 \to E[x] \to \Hom_A(A[x], E[x]) \to \Hom_A(A[x], E[x]) \to 0
$$
其中第一个映射将 $p \in E[x]$ 送到 $f \mapsto fp$，第二个映射将
$\varphi$ 送到 $f \mapsto \varphi(xf) - x\varphi(f)$。
第二个映射是满射，因为作为阿贝尔群
$\Hom_A(A[x], E[x]) = \prod_{n \geq 0} E[x]$，且该映射将
$(e_n)$ 送到 $(e_{n + 1} - xe_n)$，这是满射。
作为 $A$-模，有 $E[x] \cong \bigoplus_{n \geq 0} E$，
由引理 \ref{dualizing-lemma-sum-injective-modules} 可知它是内射的。
因此 $A[x]$-模 $\Hom_A(A[x], E[x])$ 由引理
\ref{dualizing-lemma-hom-injective} 是内射的，证明完毕。
\end{proof}



\section{射影覆盖}
\label{dualizing-section-projective-cover}

\noindent
本节简要讨论射影覆盖。

\begin{definition}
\label{dualizing-definition-projective-cover}
设 $R$ 为环。$R$-模的满射 $P \to M$ 称为{\it 射影覆盖}，
有时也称为{\it 射影包络}，若 $P$ 是射影 $R$-模且 $P \to M$ 是本质满射。
\end{definition}

\noindent
射影覆盖并不总是存在。例如，若 $k$ 是域且 $R = k[x]$ 是 $k$ 上的多项式环，
则模 $M = R/(x)$ 没有射影覆盖。事实上，对任意满射 $f : P \to M$，
其中 $P$ 在 $R$ 上是射影的，真子模 $(x - 1)P$ 也满射到 $M$。
因此 $f$ 不是本质满射。

\begin{lemma}
\label{dualizing-lemma-projective-cover-unique}
设 $R$ 为环且 $M$ 为 $R$-模。若 $M$ 存在射影覆盖，则其在同构意义下唯一。
\end{lemma}

\begin{proof}
设 $P \to M$ 与 $P' \to M$ 是射影覆盖。由于 $P$ 是射影 $R$-模，
且 $P' \to M$ 为满射，可取与到 $M$ 的映射相容的 $R$-模映射
$\alpha : P \to P'$。由于 $P' \to M$ 是本质满射，故 $\alpha$ 是满射。
由于 $P'$ 是射影 $R$-模，可取直和分解
$P = \Ker(\alpha) \oplus P'$。由于 $P' \to M$ 为满射，
且 $P \to M$ 是本质满射，得 $\Ker(\alpha)$ 为零，如所需。
\end{proof}

\noindent
下面给出射影覆盖存在的一个例子。

\begin{lemma}
\label{dualizing-lemma-projective-covers-local}
设 $(R, \mathfrak m, \kappa)$ 为局部环。任意有限 $R$-模都有射影覆盖。
\end{lemma}

\begin{proof}
设 $M$ 为有限 $R$-模。令 $r = \dim_\kappa(M/\mathfrak m M)$。
取 $x_1, \ldots, x_r \in M$，其在 $M/\mathfrak m M$ 中映为一组基。
考虑映射 $f : R^{\oplus r} \to M$。由 Nakayama 引理，该映射为满射
（《代数》引理 \ref{algebra-lemma-NAK}）。若
$N \subset R^{\oplus r}$ 是真子模，则
$N/\mathfrak m N \to \kappa^{\oplus r}$ 不满射（再次由 Nakayama 引理），
从而 $N/\mathfrak m N \to M/\mathfrak m M$ 不满射。因此 $f$ 是本质满射。
\end{proof}







\section{内射包}
\label{dualizing-section-injective-hull}

\noindent
本节简要讨论内射包。

\begin{definition}
\label{dualizing-definition-injective-hull}
设 $R$ 为环。$R$-模的单射 $M \to I$ 称为{\it 内射包}，若 $I$ 是内射
$R$-模且 $M \to I$ 是本质单射。
\end{definition}

\noindent
内射包总是存在。

\begin{lemma}
\label{dualizing-lemma-injective-hull}
设 $R$ 为环。任意 $R$-模都有内射包。
\end{lemma}

\begin{proof}
设 $M$ 为 $R$-模。由《更多代数》第
\ref{more-algebra-section-injectives-modules} 节，$R$-模范畴有足够多内射对象。
取单射 $M \to I$，其中 $I$ 为内射 $R$-模。
考虑由满足 $M \subset E \subset I$ 的子模组成的集合 $\mathcal{S}$，
其中 $E$ 是 $M$ 的本质扩张。按包含关系给 $\mathcal{S}$ 排序。
若 $\{E_\alpha\}$ 是 $\mathcal{S}$ 的全序子集，则
$\bigcup E_\alpha$ 也是 $M$ 的本质扩张
（引理 \ref{dualizing-lemma-union-essential-extensions}）。
因此可应用 Zorn 引理，取极大元 $E \in \mathcal{S}$。
我们断言 $M \subset E$ 是内射包，即 $E$ 是内射 $R$-模。
这由引理 \ref{dualizing-lemma-essential-extensions-in-injective} 得出。
\end{proof}

\begin{lemma}
\label{dualizing-lemma-injective-hull-unique}
设 $R$ 为环。设 $M$、$N$ 为 $R$-模，且 $M \to E$ 与 $N \to E'$ 是内射包。
则
\begin{enumerate}
\item 对任意 $R$-模映射 $\varphi : M \to N$，存在 $R$-模映射
$\psi : E \to E'$，使下图交换：
$$
\xymatrix{
M \ar[r] \ar[d]_\varphi & E \ar[d]^\psi \\
N \ar[r] & E'
}
$$
\item 若 $\varphi$ 是单射，则 $\psi$ 是单射；
\item 若 $\varphi$ 是本质单射，则 $\psi$ 是同构；
\item 若 $\varphi$ 是同构，则 $\psi$ 是同构；
\item 若 $M \to I$ 是将 $M$ 嵌入内射 $R$-模的嵌入，则存在同构
$I \cong E \oplus I'$，且它与 $M$ 的嵌入相容。
\end{enumerate}
特别地，$M$ 的内射包 $E$ 在同构意义下唯一。
\end{lemma}

\begin{proof}
(1) 由 $E'$ 为内射 $R$-模这一事实得出。(2) 因为
$\Ker(\psi) \cap M = 0$，且 $E$ 是 $M$ 的本质扩张。
设 $\varphi$ 是本质单射。由 (2)，$E \cong \psi(E) \subset E'$，
从而因 $E$ 内射，有 $E' = \psi(E) \oplus E''$。
由于 $E'$ 是 $M$ 的本质扩张（引理 \ref{dualizing-lemma-essential}），得 $E'' = 0$。
(4) 是 (3) 的特例。设 $M \to I$ 如 (5) 所述。
取延拓 $M \to I$ 的映射 $\alpha : E \to I$。如前所述，$\alpha$ 是单射。
因此同样地，$\alpha(E)$ 是 $I$ 的直和项。这就证明了 (5)。
\end{proof}

\begin{example}
\label{dualizing-example-injective-hull-domain}
设 $R$ 为整环，分式域为 $K$。则 $R \subset K$ 是 $R$ 的内射包。
事实上，由例 \ref{dualizing-example-reduced-ring-injective}，$K$ 是内射 $R$-模；
由引理 \ref{dualizing-lemma-essential-extension}，$R \subset K$ 是本质扩张。
\end{example}

\begin{definition}
\label{dualizing-definition-indecomposable}
加性范畴中的对象 $X$ 称为{\it 不可分解的}，若它非零，且当
$X = Y \oplus Z$ 时，$Y = 0$ 或 $Z = 0$。
\end{definition}

\begin{lemma}
\label{dualizing-lemma-indecomposable-injective}
设 $R$ 为环。设 $E$ 为不可分解的内射 $R$-模。则
\begin{enumerate}
\item $E$ 是 $E$ 的任意非零子模的内射包；
\item $E$ 的任意两个非零子模的交非零；
\item $\text{End}_R(E)$ 是非交换局部环，其极大理想由核非零的
$\varphi : E \to E$ 组成；
\item $E$ 上的零因子集合是 $R$ 的素理想 $\mathfrak p$，且 $E$ 是内射
$R_\mathfrak p$-模。
\end{enumerate}
\end{lemma}

\begin{proof}
(1) 由引理 \ref{dualizing-lemma-injective-hull-unique} 得出。(2) 由 (1) 及内射包的定义得出。

\medskip\noindent
(3) 的证明。令 $A = \text{End}_R(E)$，并令
$I = \{\varphi \in A \mid \Ker(\varphi) \not = 0\}$。
这意味着 $I$ 是双边理想，且任意 $\varphi \in A$、$\varphi \not \in I$ 都可逆。
设 $\varphi$ 与 $\psi$ 都不是单射。由 (2)，
$\Ker(\varphi) \cap \Ker(\psi)$ 非零。因此 $\varphi + \psi \in I$，
从而 $I$ 是双边理想。若 $\varphi \in A$ 且 $\varphi \not \in I$，则
$E \cong \varphi(E) \subset E$ 是内射子模，故因 $E$ 不可分解而有
$E = \varphi(E)$。

\medskip\noindent
(4) 的证明。考虑环映射 $R \to A$，令 $\mathfrak p \subset R$ 为极大理想
$I$ 的逆像。显然 $\mathfrak p$ 是素理想，且 $R \to A$ 延拓为
$R_\mathfrak p \to A$。因此 $E$ 是 $R_\mathfrak p$-模。
由引理 \ref{dualizing-lemma-injective-epimorphism}，$E$ 作为 $R_\mathfrak p$-模是内射的。
\end{proof}

\begin{lemma}
\label{dualizing-lemma-injective-hull-indecomposable}
设 $\mathfrak p \subset R$ 是环 $R$ 的素理想。
设 $E$ 是 $R/\mathfrak p$ 的内射包。则
\begin{enumerate}
\item $E$ 是不可分解的；
\item $E$ 是 $\kappa(\mathfrak p)$ 的内射包；
\item $E$ 是环 $R_\mathfrak p$ 上 $\kappa(\mathfrak p)$ 的内射包。
\end{enumerate}
\end{lemma}

\begin{proof}
由引理 \ref{dualizing-lemma-essential-extension}，包含
$R/\mathfrak p \subset \kappa(\mathfrak p)$ 是本质扩张。
于是引理 \ref{dualizing-lemma-injective-hull-unique} 表明 (2) 成立。
对 $f \in R$ 且 $f \not \in \mathfrak p$，映射
$f : \kappa(\mathfrak p) \to \kappa(\mathfrak p)$ 是同构，
故映射 $f : E \to E$ 是同构，见引理 \ref{dualizing-lemma-injective-hull-unique}。
因此 $E$ 是 $R_\mathfrak p$-模。由引理
\ref{dualizing-lemma-injective-epimorphism}，它作为 $R_\mathfrak p$-模是内射的。
最后，设 $E' \subset E$ 是非零内射 $R$-子模。
则 $J = (R/\mathfrak p) \cap E'$ 非零。缩小 $E'$ 后，可设 $E'$ 是 $J$ 的内射包
（例如见引理 \ref{dualizing-lemma-injective-hull-unique}）。
注意，由引理 \ref{dualizing-lemma-essential-extension}，$R/\mathfrak p$ 是 $J$ 的本质扩张。
因此由引理 \ref{dualizing-lemma-injective-hull-unique} 的 (3)，$E' \to E$ 是同构。
所以 $E$ 不可分解。
\end{proof}

\begin{lemma}
\label{dualizing-lemma-indecomposable-injective-noetherian}
设 $R$ 为 Noetherian 环。设 $E$ 为不可分解的内射 $R$-模。
则存在 $R$ 的某个素理想 $\mathfrak p$，使 $E$ 是
$\kappa(\mathfrak p)$ 的内射包。
\end{lemma}

\begin{proof}
令 $\mathfrak p$ 为引理 \ref{dualizing-lemma-indecomposable-injective} 中得到的素理想。
设 $\mathfrak p = (f_1, \ldots, f_r)$。取非零元素
$x \in \bigcap \Ker(f_i : E \to E)$，见引理 \ref{dualizing-lemma-indecomposable-injective}。
则 $(R_\mathfrak p)x$ 是 $E$ 中与 $\kappa(\mathfrak p)$ 同构的模。
由引理 \ref{dualizing-lemma-indecomposable-injective} 得出结论。
\end{proof}

\begin{proposition}[Noetherian 环上内射模的结构]
\label{dualizing-proposition-structure-injectives-noetherian}
设 $R$ 为 Noetherian 环。
每个内射模都是不可分解内射模的直和。
每个不可分解内射模都是某个素理想处剩余域的内射包。
\end{proposition}

\begin{proof}
第二个断言就是引理 \ref{dualizing-lemma-indecomposable-injective-noetherian}。
对第一个断言，设 $I$ 为内射 $R$-模。我们使用超限归纳构造
$I_\alpha \subset I$：对序数 $\alpha$，它们是 $\beta < \alpha$ 时的不可分解内射
$R$-模 $E_{\beta + 1}$ 的直和。
当 $\alpha = 0$ 时令 $I_0 = 0$。设序数 $\alpha$ 已给定且 $I_\alpha$ 已构造。
由引理 \ref{dualizing-lemma-sum-injective-modules}，$I_\alpha$ 是内射 $R$-模。
因此 $I \cong I_\alpha \oplus I'$。若 $I' = 0$，则完成。
否则由《代数》引理 \ref{algebra-lemma-ass-zero}，$I'$ 有一个伴随素理想。
因此对某个素理想 $\mathfrak p$，$I'$ 含有 $R/\mathfrak p$ 的一个副本。
于是由引理 \ref{dualizing-lemma-injective-hull-unique} 与
\ref{dualizing-lemma-injective-hull-indecomposable}，$I'$ 含有不可分解子模 $E$。
令 $I_{\alpha + 1} = I_\alpha \oplus E_\alpha$。
若 $\alpha$ 是极限序数，且对 $\beta < \alpha$ 已构造 $I_\beta$，则令
$I_\alpha = \bigcup_{\beta < \alpha} I_\beta$。
注意 $I_\alpha = \bigoplus_{\beta < \alpha} E_{\beta + 1}$。
证明完毕。
\end{proof}



\section{Artinian 局部环上的对偶}
\label{dualizing-section-artinian}

\noindent
设 $(R, \mathfrak m, \kappa)$ 为 Artinian 局部环。
回忆这意味着 $R$ 是 Noetherian 环，且 $R$ 作为 $R$-模具有有限长度。
此外，$R$-模是有限的当且仅当它具有有限长度。本节中将不再重复说明这些事实。
更多细节见《代数》第 \ref{algebra-section-length} 节与
\ref{algebra-section-artinian} 节，以及《代数》命题
\ref{algebra-proposition-dimension-zero-ring}。

\begin{lemma}
\label{dualizing-lemma-finite}
设 $(R, \mathfrak m, \kappa)$ 为 Artinian 局部环。
设 $E$ 为 $\kappa$ 的内射包。对每个有限 $R$-模 $M$，有
$$
\text{length}_R(M) = \text{length}_R(\Hom_R(M, E))
$$
特别地，$\kappa$ 的内射包 $E$ 是有限 $R$-模。
\end{lemma}

\begin{proof}
由于 $E$ 是 $\kappa$ 的本质扩张，有
$\kappa = E[\mathfrak m]$，其中 $E[\mathfrak m]$ 是 $E$ 中的
$\mathfrak m$-挠部分（记号见《更多代数》第
\ref{more-algebra-section-torsion} 节）。
于是 $\Hom_R(\kappa, E) \cong \kappa$，长度等式对 $M = \kappa$ 成立。
我们对 $M$ 的长度归纳证明引理中的显示等式。若 $M$ 非零，则存在核为 $M'$ 的满射
$M \to \kappa$。由于函子 $M \mapsto \Hom_R(M, E)$ 正合，可得短正合序列
$$
0 \to \Hom_R(\kappa, E) \to \Hom_R(M, E) \to \Hom_R(M', E) \to 0.
$$
该序列以及序列 $0 \to M' \to M \to \kappa \to 0$ 的长度可加性，
结合归纳假设对 $M'$ 与 $\kappa$ 的等式，推出 $M$ 的等式。
引理的最后一个断言由 $E = \Hom_R(R, E)$ 得出。
\end{proof}

\begin{lemma}
\label{dualizing-lemma-evaluate}
设 $(R, \mathfrak m, \kappa)$ 为 Artinian 局部环。
设 $E$ 为 $\kappa$ 的内射包。对任意有限 $R$-模 $M$，求值映射
$$
M \longrightarrow \Hom_R(\Hom_R(M, E), E)
$$
是同构。特别地，$R = \Hom_R(E, E)$。
\end{lemma}

\begin{proof}
注意显示的箭头是单射。事实上，若 $x \in M$ 为非零元素，
则存在非零映射 $Rx \to \kappa$，可将其延拓为在 $x$ 上不为零的映射
$\varphi : M \to E$。由引理 \ref{dualizing-lemma-finite}，箭头的源与靶具有相同长度，
故它是同构。最后一个断言由取 $M = R$ 得出。
\end{proof}

\noindent
记 $\text{Mod}^{fg}_R$ 为环 $R$ 上有限 $R$-模所成的范畴，则下一个引理可表述为：

\begin{lemma}
\label{dualizing-lemma-duality}
设 $(R, \mathfrak m, \kappa)$ 为 Artinian 局部环。
设 $E$ 为 $\kappa$ 的内射包。函子 $D(-) = \Hom_R(-, E)$
诱导正合反等价
$\text{Mod}^{fg}_R \to \text{Mod}^{fg}_R$，并且
$D \circ D \cong \text{id}$。
\end{lemma}

\begin{proof}
由引理 \ref{dualizing-lemma-evaluate}，我们已知 $D \circ D = \text{id}$ 在
$\text{Mod}^{fg}_R$ 上成立。因此立即得到 $D$ 是反等价。
\end{proof}

\begin{lemma}
\label{dualizing-lemma-duality-torsion-cotorsion}
假设与记号同引理 \ref{dualizing-lemma-duality}。设 $I \subset R$ 为理想，$M$ 为有限 $R$-模。
则
$$
D(M[I]) = D(M)/ID(M) \quad\text{且}\quad D(M/IM) = D(M)[I]
$$
\end{lemma}

\begin{proof}
设 $I = (f_1, \ldots, f_t)$。考虑映射
$$
M^{\oplus t} \xrightarrow{f_1, \ldots, f_t} M
$$
其余核为 $M/IM$。应用正合函子 $D$，得到 $D(M/IM)$ 是 $D(M)[I]$。
另一种情形同理可证。
\end{proof}



\section{剩余域的内射包}
\label{dualizing-section-hull-residue-field}

\noindent
本节中的大多数结果针对 Noetherian 局部环。

\begin{lemma}
\label{dualizing-lemma-quotient}
设 $R \to S$ 是局部环的满射，核为 $I$。
设 $E$ 为 $R$ 上 $R$ 的剩余域的内射包。
则 $E[I]$ 是 $S$ 上 $S$ 的剩余域的内射包。
\end{lemma}

\begin{proof}
注意，由 $S = R/I$，有 $E[I] = \Hom_R(S, E)$。
由引理 \ref{dualizing-lemma-hom-injective}，$E[I]$ 是内射 $S$-模。
由于 $E$ 是 $\kappa = R/\mathfrak m_R$ 的本质扩张，$E[I]$ 同样是
$\kappa$ 的本质扩张。结论得出。
\end{proof}

\begin{lemma}
\label{dualizing-lemma-torsion-submodule-sum-injective-hulls}
设 $(R, \mathfrak m, \kappa)$ 为局部环。
设 $E$ 为 $\kappa$ 的内射包。设 $M$ 为 $\mathfrak m$-幂挠 $R$-模，
且 $n = \dim_\kappa(M[\mathfrak m]) < \infty$。
则 $M$ 同构于 $E^{\oplus n}$ 的一个子模。
\end{lemma}

\begin{proof}
注意 $E^{\oplus n}$ 是 $\kappa^{\oplus n} = M[\mathfrak m]$ 的内射包。
因此存在 $R$-模映射 $M \to E^{\oplus n}$，其在 $M[\mathfrak m]$ 上为单射。
由于 $M$ 是 $\mathfrak m$-幂挠的，包含 $M[\mathfrak m] \subset M$ 是本质扩张
（例如由引理 \ref{dualizing-lemma-essential-extension}），故
$M \to E^{\oplus n}$ 的核为零。
\end{proof}

\begin{lemma}
\label{dualizing-lemma-union-artinian}
设 $(R, \mathfrak m, \kappa)$ 为 Noetherian 局部环。
设 $E$ 为 $R$ 上 $\kappa$ 的内射包。
设 $E_n$ 为 $R/\mathfrak m^n$ 上 $\kappa$ 的内射包。
则 $E = \bigcup E_n$ 且 $E_n = E[\mathfrak m^n]$。
\end{lemma}

\begin{proof}
由引理 \ref{dualizing-lemma-quotient}，$E_n = E[\mathfrak m^n]$。
又因 $\bigcup E_n = E[\mathfrak m^\infty]$ 是包含 $\kappa$ 的内射
$R$-子模，故 $E = \bigcup E_n$，见引理 \ref{dualizing-lemma-injective-module-divide}。
\end{proof}

\noindent
下一个引理说明，Noetherian 局部环的剩余域的内射包只取决于完备化。

\begin{lemma}
\label{dualizing-lemma-compare}
设 $R \to S$ 是局部 Noetherian 环的平坦局部同态，满足
$R/\mathfrak m_R \cong S/\mathfrak m_R S$。
则 $R$ 的剩余域的内射包也是 $S$ 的剩余域的内射包。
\end{lemma}

\begin{proof}
注意 $\mathfrak m_RS = \mathfrak m_S$，因为前者的商是域。
令 $\kappa = R/\mathfrak m_R = S/\mathfrak m_S$。
令 $E_R$ 为 $R$ 上 $\kappa$ 的内射包，令 $E_S$ 为 $S$ 上 $\kappa$ 的内射包。
由引理 \ref{dualizing-lemma-injective-flat}，$E_S$ 是内射 $R$-模。
取剩余域恒等的延拓 $E_R \to E_S$。由引理
\ref{dualizing-lemma-union-artinian}，该映射是同构，因为对每个 $n$，$R \to S$ 诱导同构
$R/\mathfrak m_R^n \to S/\mathfrak m_S^n$。
\end{proof}

\begin{lemma}
\label{dualizing-lemma-endos}
设 $(R, \mathfrak m, \kappa)$ 为 Noetherian 局部环。
设 $E$ 为 $R$ 上 $\kappa$ 的内射包。则 $\Hom_R(E, E)$
与 $R$ 的完备化有典范同构。
\end{lemma}

\begin{proof}
按引理 \ref{dualizing-lemma-union-artinian}，写 $E = \bigcup E_n$，其中
$E_n = E[\mathfrak m^n]$。$E$ 的任意自同态都保持这一滤过。因此
$$
\Hom_R(E, E) = \lim \Hom_R(E_n, E_n)
$$
由引理 \ref{dualizing-lemma-evaluate}，
$\Hom_R(E_n, E_n) = \Hom_{R/\mathfrak m^n}(E_n, E_n) = R/\mathfrak m^n$，
故得证。
\end{proof}

\begin{lemma}
\label{dualizing-lemma-injective-hull-has-dcc}
设 $(R, \mathfrak m, \kappa)$ 为 Noetherian 局部环。
设 $E$ 为 $R$ 上 $\kappa$ 的内射包。则 $E$ 满足降链条件。
\end{lemma}

\begin{proof}
若 $E \supset M_1 \supset M_2 \supset \ldots$ 是子模序列，则
$$
\Hom_R(E, E) \to \Hom_R(M_1, E) \to \Hom_R(M_2, E) \to \ldots
$$
是满射序列。由引理 \ref{dualizing-lemma-endos}，其中每一项都是完备化
$R^\wedge = \Hom_R(E, E)$ 上的模。
由于 $R^\wedge$ 是 Noetherian 环（《代数》引理
\ref{algebra-lemma-completion-Noetherian-Noetherian}），序列稳定，即
$\Hom_R(M_n, E) = \Hom_R(M_{n + 1}, E) = \ldots$。
由于 $E$ 内射，这只能在 $\Hom_R(M_n/M_{n + 1}, E)$ 为零时发生。
然而若 $M_n/M_{n + 1}$ 非零，则其中含有被 $\mathfrak m$ 湮灭的非零元素，
因为由引理 \ref{dualizing-lemma-union-artinian}，$E$ 是 $\mathfrak m$-幂挠的。
此时 $M_n/M_{n + 1}$ 有非零映射到 $E$，与假定的消失矛盾。证明完毕。
\end{proof}

\begin{lemma}
\label{dualizing-lemma-describe-categories}
设 $(R, \mathfrak m, \kappa)$ 为 Noetherian 局部环。
设 $E$ 为 $\kappa$ 的内射包。
\begin{enumerate}
\item 对 $R$-模 $M$，以下条件等价：
\begin{enumerate}
\item $M$ 满足升链条件；
\item $M$ 是有限 $R$-模；
\item 存在 $n, m$ 及正合序列
$R^{\oplus m} \to R^{\oplus n} \to M \to 0$。
\end{enumerate}
\item 对 $R$-模 $M$，以下条件等价：
\begin{enumerate}
\item $M$ 满足降链条件；
\item $M$ 是 $\mathfrak m$-幂挠的，且
$\dim_\kappa(M[\mathfrak m]) < \infty$；
\item 存在 $n, m$ 及正合序列
$0 \to M \to E^{\oplus n} \to E^{\oplus m}$。
\end{enumerate}
\end{enumerate}
\end{lemma}

\begin{proof}
(1) 的证明略。

\medskip\noindent
设 $M$ 为满足降链条件的 $R$-模，取 $x \in M$。
$\mathfrak m^n x$ 是子模降链，因而稳定。所以对某个 $n$，
$\mathfrak m^nx = \mathfrak m^{n + 1}x$。由 Nakayama 引理
（《代数》引理 \ref{algebra-lemma-NAK}），这意味着 $\mathfrak m^n x = 0$，
即 $x$ 是 $\mathfrak m$-幂挠的。由于 $M[\mathfrak m]$ 是 $\kappa$ 上的向量空间，
要满足降链条件，它必须是有限维的。

\medskip\noindent
设 $M$ 是 $\mathfrak m$-幂挠的，且其 $\mathfrak m$-挠子模 $M[\mathfrak m]$ 有限维。
由引理 \ref{dualizing-lemma-torsion-submodule-sum-injective-hulls}，对某个 $n$，$M$ 是
$E^{\oplus n}$ 的子模。考虑商 $N = E^{\oplus n}/M$。由引理
\ref{dualizing-lemma-injective-hull-has-dcc}，$E$ 满足降链条件，故 $E^{\oplus n}$ 与 $N$ 也满足。
因此 $N$ 满足 (2)(a)，由证明的第二段可知 $N$ 满足 (2)(b)。
再由引理 \ref{dualizing-lemma-torsion-submodule-sum-injective-hulls}，对某个 $m$，$N$ 是
$E^{\oplus m}$ 的子模。因此有短正合序列
$0 \to M \to E^{\oplus n} \to E^{\oplus m}$。

\medskip\noindent
设有短正合序列 $0 \to M \to E^{\oplus n} \to E^{\oplus m}$。
由于由引理 \ref{dualizing-lemma-injective-hull-has-dcc}，$E$ 满足降链条件，
故 $M$ 也满足降链条件。
\end{proof}

\begin{proposition}[Matlis 对偶]
\label{dualizing-proposition-matlis}
设 $(R, \mathfrak m, \kappa)$ 为完备局部 Noetherian 环。
设 $E$ 为 $R$ 上 $\kappa$ 的内射包。函子
$D(-) = \Hom_R(-, E)$ 诱导反等价
$$
\left\{
\begin{matrix}
R\text{-modules with the} \\
\text{descending chain condition}
\end{matrix}
\right\}
\longleftrightarrow
\left\{
\begin{matrix}
R\text{-modules with the} \\
\text{ascending chain condition}
\end{matrix}
\right\}
$$
并且等价两侧都有 $D \circ D = \text{id}$。
\end{proposition}

\begin{proof}
由引理 \ref{dualizing-lemma-endos}，有 $R = \Hom_R(E, E) = D(E)$。
显然 $E = \Hom_R(R, E) = D(R)$。由于 $E$ 是内射的，函子 $D$ 正合。
结论立即由引理 \ref{dualizing-lemma-describe-categories} 中对这些范畴的描述得出。
\end{proof}

\begin{remark}
\label{dualizing-remark-matlis}
设 $(R, \mathfrak m, \kappa)$ 为 Noetherian 局部环。
设 $E$ 为 $R$ 上 $\kappa$ 的内射包。这里给出 Matlis 对偶的补充：若 $N$ 是
$\mathfrak m$-幂挠模，且 $M = \Hom_R(N, E)$ 是 $R$ 的完备化上的有限模，
则 $N$ 满足降链条件。事实上，对任意满足 $N'' \subset N' \subset N$ 且
$N'' \not = N'$ 的子模，可取嵌入 $\kappa \subset N''/N'$，从而得到湮灭
$N''$ 的非零映射 $N' \to E$，并将其延拓为湮灭 $N''$ 的映射 $N \to E$。
因此 $N \supset N' \mapsto M' = \Hom_R(N/N', E) \subset M$
是从 $N$ 的子模到 $M$ 的子模的保持包含映射，结论随即得到。
\end{remark}




















\section{导出挠部分}
\label{dualizing-section-bad-local-cohomology}

\noindent
设 $A$ 为环，$I \subset A$ 为有限生成理想
（若 $I$ 不是有限生成的，或许应采用不同定义）。令 $Z = V(I) \subset \Spec(A)$。
回忆，$I$-幂挠模所成的范畴 $I^\infty\text{-torsion}$
只取决于闭子集 $Z$，而不取决于满足 $Z = V(I)$ 的有限生成理想 $I$ 的选择，见
《更多代数》引理 \ref{more-algebra-lemma-local-cohomology-closed}。
本节考虑函子
$$
H^0_{I} : \text{Mod}_A \longrightarrow I^\infty\text{-torsion},\quad
M \longmapsto M[I^\infty] = \bigcup M[I^n]
$$
它将 $M$ 送到其 $I$-幂挠子模。

\medskip\noindent
设 $A$ 为环，$I$ 为有限生成理想。注意 $I^\infty\text{-torsion}$ 是 Grothendieck
阿贝尔范畴（直和存在，滤余极限正合，且由
《更多代数》引理 \ref{more-algebra-lemma-I-power-torsion-presentation}，
$\bigoplus A/I^n$ 是生成元）。因此导出范畴
$D(I^\infty\text{-torsion})$ 存在，见《内射模》备注
\ref{injectives-remark-existence-D}。
函子 $H^0_I$ 左正合，其导出延拓记为
$$
R\Gamma_I : D(A) \longrightarrow D(I^\infty\text{-torsion}).
$$
{\bf Warning:} 除非环 $A$ 是 Noetherian 环，否则此函子不配称为局部上同调。
函子 $H^0_I$、$R\Gamma_I$ 及卫星函子 $H^p_I$ 只取决于闭子集
$Z \subset \Spec(A)$，而不取决于满足 $V(I) = Z$ 的有限生成理想 $I$ 的选择。
不过，我们坚持在上述函子中使用下标 $I$，因为记号 $R\Gamma_Z$ 将用于另一函子，见
（\ref{dualizing-equation-local-cohomology}）；据我们所知，该函子与 $R\Gamma_I$ 一致
仅在 $A$ 为 Noetherian 环时成立（见引理 \ref{dualizing-lemma-local-cohomology-noetherian}）。

\begin{lemma}
\label{dualizing-lemma-adjoint}
设 $A$ 为环，$I \subset A$ 为有限生成理想。
函子 $R\Gamma_I$ 是函子 $D(I^\infty\text{-torsion}) \to D(A)$ 的右伴随。
\end{lemma}

\begin{proof}
这是因为取 $I$-幂挠子模是包含函子
$I^\infty\text{-torsion} \to \text{Mod}_A$ 的右伴随。见《导出范畴》引理
\ref{derived-lemma-derived-adjoint-functors}。
\end{proof}

\begin{lemma}
\label{dualizing-lemma-local-cohomology-ext}
设 $A$ 为环，$I \subset A$ 为有限生成理想。对 $D(A)$ 中任意对象 $K$，有
$$
R\Gamma_I(K) = \text{hocolim}\ R\Hom_A(A/I^n, K)
$$
在 $D(A)$ 中，且
$$
R^q\Gamma_I(K) = \colim_n \Ext_A^q(A/I^n, K)
$$
作为模对所有 $q \in \mathbf{Z}$ 成立。
\end{lemma}

\begin{proof}
设 $J^\bullet$ 为表示 $K$ 的 K-内射复形。则
$$
R\Gamma_I(K) = J^\bullet[I^\infty] = \colim J^\bullet[I^n] =
\colim \Hom_A(A/I^n, J^\bullet)
$$
其中第一等式是 $R\Gamma_I(K)$ 的定义。由《导出范畴》引理
\ref{derived-lemma-colim-hocolim}，得到引理陈述中的第一个显示等式。
第二个显示等式由
$H^q(\Hom_A(A/I^n, J^\bullet)) = \Ext^q_A(A/I^n, K)$
以及阿贝尔群范畴中的滤余极限正合而得。
\end{proof}

\begin{lemma}
\label{dualizing-lemma-bad-local-cohomology-vanishes}
设 $A$ 为环，$I \subset A$ 为有限生成理想。
设 $K^\bullet$ 为 $A$-模复形，且对某个 $f \in I$，
$f : K^\bullet \to K^\bullet$ 是同构，即 $K^\bullet$ 是 $A_f$-模复形。
则 $R\Gamma_I(K^\bullet) = 0$。
\end{lemma}

\begin{proof}
事实上，此时 $R\Gamma_I(K^\bullet)$ 的上同调模既是 $f$-幂挠的，
又由 $f$ 的自同构作用。因此上同调模为零，故对象为零。
\end{proof}

\noindent
设 $A$ 为环，$I \subset A$ 为有限生成理想。
由《更多代数》引理 \ref{more-algebra-lemma-I-power-torsion}，
$I$-幂挠模所成的范畴是全体 $A$-模范畴的 Serre 子范畴。因此存在函子
\begin{equation}
\label{dualizing-equation-compare-torsion}
D(I^\infty\text{-torsion}) \longrightarrow D_{I^\infty\text{-torsion}}(A)
\end{equation}
其中右侧是 $D(A)$ 的全子范畴，由上同调为 $I$-幂挠模的对象组成。
该函子与该范畴都在《导出范畴》第
\ref{derived-section-triangulated-sub} 节中讨论。

\begin{lemma}
\label{dualizing-lemma-not-equal}
设 $A$ 为环，$I$ 为有限生成理想。设 $M$ 与 $N$ 为 $I$-幂挠模。
\begin{enumerate}
\item $\Hom_{D(A)}(M, N) = \Hom_{D({I^\infty\text{-torsion}})}(M, N)$；
\item $\Ext^1_{D(A)}(M, N) =
\Ext^1_{D({I^\infty\text{-torsion}})}(M, N)$；
\item $\Ext^2_{D({I^\infty\text{-torsion}})}(M, N) \to
\Ext^2_{D(A)}(M, N)$ 一般不是满射；
\item （\ref{dualizing-equation-compare-torsion}）一般不是等价。
\end{enumerate}
\end{lemma}

\begin{proof}
(1) 与 (2) 立即由 $I$-幂挠模构成 $\text{Mod}_A$ 的 Serre 子范畴这一事实得出。
(4) 由 (3) 得出。

\medskip\noindent
为证明 (3)，设环 $A$ 中有元素 $f \in A$，使得 $A[f]$ 含有一个被 $f$
湮灭的非零元素 $x$，且 $A$ 含有满足 $f^nx_n = x$ 的元素 $x_n$。
这样的环 $A$ 存在，因为可取
$$
A = \mathbf{Z}[f, x, x_n]/(fx, f^nx_n - x)
$$
给定 $A$，令 $I = (f)$。则正合序列
$$
0 \to A[f] \to A \xrightarrow{f} A \to A/fA \to 0
$$
定义 $\Ext^2_A(A/fA, A[f])$ 中的一个元素。我们断言该元素不来自
$\Ext^2_{D(f^\infty\text{-torsion})}(A/fA, A[f])$ 中的元素。
否则将存在正合序列
$$
0 \to A[f] \to M \to N \to A/fA \to 0
$$
其中 $M$ 与 $N$ 是定义同一个 $2$-扩张类的 $f$-幂挠模。
由于 $A \to A$ 是自由模复形，且两个 $2$-扩张类相同，可找到映射
$$
\xymatrix{
0 \ar[r] &
A[f] \ar[r] \ar[d] &
A \ar[r] \ar[d]_\varphi &
A \ar[r] \ar[d]_\psi &
A/fA \ar[r] \ar[d] & 0 \\
0 \ar[r] &
A[f] \ar[r] &
M \ar[r] &
N \ar[r] &
A/fA \ar[r] & 0
}
$$
（略去若干细节）。然后可用 $\varphi$ 的像替换 $M$，用 $\psi$ 的像替换 $N$。
此时 $M$ 是循环模，故对某个 $n$ 有 $f^n M = 0$。
考察 $\varphi(x_{n + 1})$，便与 $f^{n + 1}x_n = x$ 在 $A[f]$ 中非零这一事实矛盾。
\end{proof}









\section{局部上同调}
\label{dualizing-section-local-cohomology}

\noindent
设 $A$ 为环，$I \subset A$ 为有限生成理想。令
$Z = V(I) \subset \Spec(A)$。我们将构造函子
\begin{equation}
\label{dualizing-equation-local-cohomology}
R\Gamma_Z : D(A) \longrightarrow D_{I^\infty\text{-torsion}}(A).
\end{equation}
它是包含函子的右伴随。记号见第 \ref{dualizing-section-bad-local-cohomology} 节。
$R\Gamma_Z(K)$ 的上同调模称为 $K$ 关于 $Z$ 的{\it 局部上同调群}。
由引理 \ref{dualizing-lemma-not-equal}，即使将二者都视为取值于 $D(A)$ 的函子，
此函子一般也{\bf 不}等于 $R\Gamma_I( - )$。
第 \ref{dualizing-section-local-cohomology-noetherian} 节将证明，当 $A$ 为 Noetherian 环时，
二者一致。

\medskip\noindent
我们将在局部上同调一章继续讨论局部上同调，见《局部上同调》第
\ref{local-cohomology-section-introduction} 节。例如，那里将证明
$R\Gamma_Z$ 计算 $\Spec(A)$ 上相伴拟凝聚层复形的支撑在 $Z$ 中的上同调。
见《局部上同调》引理
\ref{local-cohomology-lemma-local-cohomology-is-local-cohomology}。


\begin{lemma}
\label{dualizing-lemma-local-cohomology-adjoint}
设 $A$ 为环，$I \subset A$ 为有限生成理想。
包含函子 $D_{I^\infty\text{-torsion}}(A) \to D(A)$ 存在右伴随
$R\Gamma_Z$（\ref{dualizing-equation-local-cohomology}）。
事实上，若 $I$ 由 $f_1, \ldots, f_r \in A$ 生成，则有
$$
R\Gamma_Z(K) =
(A \to \prod\nolimits_{i_0} A_{f_{i_0}} \to
\prod\nolimits_{i_0 < i_1} A_{f_{i_0}f_{i_1}}
\to \ldots \to A_{f_1\ldots f_r}) \otimes_A^\mathbf{L} K
$$
并且对 $K \in D(A)$ 函子性成立。
\end{lemma}

\begin{proof}
设 $I = (f_1, \ldots, f_r)$ 为理想。设 $K^\bullet$ 为 $A$-模复形。
存在复形的典范映射
$$
(A \to \prod\nolimits_{i_0} A_{f_{i_0}} \to
\prod\nolimits_{i_0 < i_1} A_{f_{i_0}f_{i_1}} \to
\ldots \to A_{f_1\ldots f_r}) \longrightarrow A.
$$
即从扩充的 {\v C}ech 复形到 $A$ 的映射。
与 $K^\bullet$ 张量并取相伴全复形，得到映射
$$
\text{Tot}\left(
K^\bullet \otimes_A
(A \to \prod\nolimits_{i_0} A_{f_{i_0}} \to
\prod\nolimits_{i_0 < i_1} A_{f_{i_0}f_{i_1}} \to
\ldots \to A_{f_1\ldots f_r})\right)
\longrightarrow
K^\bullet
$$
于 $D(A)$ 中。我们断言左侧复形的上同调模是 $I$-幂挠的，
即 LHS 是 $D_{I^\infty\text{-torsion}}(A)$ 中的对象。事实上，有
$$
(A \to \prod\nolimits_{i_0} A_{f_{i_0}} \to
\prod\nolimits_{i_0 < i_1} A_{f_{i_0}f_{i_1}} \to
\ldots \to A_{f_1\ldots f_r}) = \colim K(A, f_1^n, \ldots, f_r^n)
$$
见《更多代数》引理
\ref{more-algebra-lemma-extended-alternating-Cech-is-colimit-koszul}。
此外，由《更多代数》引理 \ref{more-algebra-lemma-homotopy-koszul}，
复形 $K(A, f_1^n, \ldots, f_r^n)$ 上乘以 $f_i^n$ 零伦。
由于
$$
H^q\left( LHS \right) =
\colim H^q(\text{Tot}(K^\bullet \otimes_A K(A, f_1^n, \ldots, f_r^n)))
$$
断言得证。另一方面，若 $K^\bullet$ 是
$D_{I^\infty\text{-torsion}}(A)$ 中的对象，则复形
$K^\bullet \otimes_A A_{f_{i_0} \ldots f_{i_p}}$ 的上同调消失。
因此此时映射 $LHS \to K^\bullet$ 在 $D(A)$ 中是同构。构造
$$
R\Gamma_Z(K^\bullet) =
\text{Tot}\left(
K^\bullet \otimes_A
(A \to \prod\nolimits_{i_0} A_{f_{i_0}} \to
\prod\nolimits_{i_0 < i_1} A_{f_{i_0}f_{i_1}} \to
\ldots \to A_{f_1\ldots f_r})\right)
$$
对 $K^\bullet$ 函子性成立，并定义三角范畴之间的正合函子
$D(A) \to D_{I^\infty\text{-torsion}}(A)$。
由上述自然变换 $R\Gamma_Z \to \text{id}$ 的存在，以及它在子范畴中的
$K^\bullet$ 上取同构这一事实，形式地得出 $R\Gamma_Z$ 是所需右伴随。
\end{proof}

\begin{lemma}
\label{dualizing-lemma-local-cohomology-and-restriction}
设 $A \to B$ 为环同态，$I \subset A$ 为有限生成理想。令 $J = IB$，
再令 $Z = V(I)$、$Y = V(J)$。则
$$
R\Gamma_Z(M_A) = R\Gamma_Y(M)_A
$$
对 $M \in D(B)$ 函子性成立。这里 $(-)_A$ 表示标量限制函子
$D(B) \to D(A)$ 以及
$D_{J^\infty\text{-torsion}}(B) \to D_{I^\infty\text{-torsion}}(A)$。
\end{lemma}

\begin{proof}
这由伴随函子的唯一性得出，因为 $R\Gamma_Z((-)_A)$ 与
$R\Gamma_Y(-)_A$ 都是函子
$D_{I^\infty\text{-torsion}}(A) \to D(B)$，
$K \mapsto K \otimes_A^\mathbf{L} B$ 的右伴随。
也可使用交错 {\v C}ech 复形给出的 $R\Gamma_Z$ 与 $R\Gamma_Y$ 的描述
（引理 \ref{dualizing-lemma-local-cohomology-adjoint}）。
具体地，若 $I = (f_1, \ldots, f_r)$，则 $J$ 由 $f_1, \ldots, f_r$
的像 $g_1, \ldots, g_r \in B$ 生成。
此时该引理的陈述由典范同构
\begin{align*}
& M_A \otimes_A (A \to \prod\nolimits_{i_0} A_{f_{i_0}} \to
\prod\nolimits_{i_0 < i_1} A_{f_{i_0}f_{i_1}}
\to \ldots \to A_{f_1\ldots f_r}) \\
& = 
M \otimes_B (B \to \prod\nolimits_{i_0} B_{g_{i_0}} \to
\prod\nolimits_{i_0 < i_1} B_{g_{i_0}g_{i_1}}
\to \ldots \to B_{g_1\ldots g_r})
\end{align*}
的存在得出，其中 $M$ 为任意 $B$-模。
\end{proof}

\begin{lemma}
\label{dualizing-lemma-torsion-change-rings}
设 $A \to B$ 为环同态，$I \subset A$ 为有限生成理想。令 $J = IB$，
$Z = V(I)$ 且 $Y = V(J)$。则
$$
R\Gamma_Z(K) \otimes_A^\mathbf{L} B = R\Gamma_Y(K \otimes_A^\mathbf{L} B)
$$
对 $K \in D(A)$ 函子性成立。
\end{lemma}

\begin{proof}
写 $I = (f_1, \ldots, f_r)$。则 $J$ 由 $f_1, \ldots, f_r$
的像 $g_1, \ldots, g_r \in B$ 生成，并且作为 $B$-模复形有
$$
(A \to \prod A_{f_{i_0}} \to \ldots \to A_{f_1\ldots f_r}) \otimes_A B =
(B \to \prod B_{g_{i_0}} \to \ldots \to B_{g_1\ldots g_r})
$$
以 $A$-模的 K-平坦复形 $K^\bullet$ 表示 $K$。由于与
$$
K^\bullet \otimes_A
(A \to \prod A_{f_{i_0}} \to \ldots \to A_{f_1\ldots f_r}) \otimes_A B
$$
以及
$$
K^\bullet \otimes_A B \otimes_B
(B \to \prod B_{g_{i_0}} \to \ldots \to B_{g_1\ldots g_r})
$$
相伴的全复形分别表示该引理所示公式的左侧与右侧
（见引理 \ref{dualizing-lemma-local-cohomology-adjoint}），故结论成立。
\end{proof}

\begin{lemma}
\label{dualizing-lemma-local-cohomology-vanishes}
设 $A$ 为环，$I \subset A$ 为有限生成理想。设 $K^\bullet$ 为
$A$-模复形，并且对某个 $f \in I$，映射
$f : K^\bullet \to K^\bullet$ 是同构；换言之，$K^\bullet$ 是
$A_f$-模复形。则 $R\Gamma_Z(K^\bullet) = 0$。
\end{lemma}

\begin{proof}
事实上，此时 $R\Gamma_Z(K^\bullet)$ 的上同调模既是 $f$-幂挠的，
而 $f$ 又通过自同构作用。因此这些上同调模为零，从而该对象为零。
\end{proof}

\begin{lemma}
\label{dualizing-lemma-torsion-tensor-product}
设 $A$ 为环，$I \subset A$ 为有限生成理想。对 $K, L \in D(A)$ 有
$$
R\Gamma_Z(K \otimes_A^\mathbf{L} L) =
K \otimes_A^\mathbf{L} R\Gamma_Z(L) =
R\Gamma_Z(K) \otimes_A^\mathbf{L} L =
R\Gamma_Z(K) \otimes_A^\mathbf{L} R\Gamma_Z(L)
$$
若 $K$ 或 $L$ 属于 $D_{I^\infty\text{-torsion}}(A)$，则
$K \otimes_A^\mathbf{L} L$ 也属于其中。
\end{lemma}

\begin{proof}
由引理 \ref{dualizing-lemma-local-cohomology-adjoint}，对某个 $C \in D(A)$，
$R\Gamma_Z$ 由 $C \otimes^\mathbf{L} -$ 给出。因此对一般的
$K, L \in D(A)$ 有
$$
R\Gamma_Z(K \otimes_A^\mathbf{L} L) =
K \otimes^\mathbf{L} L \otimes_A^\mathbf{L} C =
K \otimes_A^\mathbf{L} R\Gamma_Z(L)
$$
其余等式由此形式地得出；该引理的最后一个陈述也随之成立。
\end{proof}

\begin{lemma}
\label{dualizing-lemma-local-cohomology-ss}
设 $A$ 为环，$I, J \subset A$ 为有限生成理想。令 $Z = V(I)$、
$Y = V(J)$。则 $Z \cap Y = V(I + J)$，并且作为函子
$D(A) \to D_{(I + J)^\infty\text{-torsion}}(A)$ 有
$R\Gamma_Y \circ R\Gamma_Z = R\Gamma_{Y \cap Z}$。对 $K \in D^+(A)$，
存在谱序列
$$
E_2^{p, q} = H^p_Y(H^q_Z(K)) \Rightarrow H^{p + q}_{Y \cap Z}(K)
$$
如《导出范畴》引理
\ref{derived-lemma-grothendieck-spectral-sequence} 所述。
\end{lemma}

\begin{proof}
严格说来，$R\Gamma_Y$ 与 $R\Gamma_Z$ 不能直接复合，所以引理中的记号
略有滥用。该陈述的含义只是：先将 $R\Gamma_Z$ 与包含函子
$D_{I^\infty\text{-torsion}}(A) \to D(A)$ 复合，再与
$R\Gamma_Y$ 复合。于是等式
$R\Gamma_Y \circ R\Gamma_Z = R\Gamma_{Y \cap Z}$
由下述事实得出：函子
$$
D_{I^\infty\text{-torsion}}(A) \to D(A) \xrightarrow{R\Gamma_Y}
D_{(I + J)^\infty\text{-torsion}}(A)
$$
是包含函子
$D_{(I + J)^\infty\text{-torsion}}(A) \to D_{I^\infty\text{-torsion}}(A)$
的右伴随。也可使用引理 \ref{dualizing-lemma-local-cohomology-adjoint}
以及如下事实证明该公式：关于 $f_1, \ldots, f_r$ 与
$g_1, \ldots, g_m$ 的扩充 {\v C}ech 复形的张量积，正是关于
$f_1, \ldots, f_n. g_1, \ldots, g_m$ 的扩充 {\v C}ech 复形。
最后一个断言由此及所引引理得出。
\end{proof}

\noindent
以下引理是《更多代数》引理
\ref{more-algebra-lemma-restriction-derived-complete-equivalence}
对上同调为挠模的复形的类似结果。

\begin{lemma}
\label{dualizing-lemma-torsion-flat-change-rings}
设 $A \to B$ 为平坦环同态，$I \subset A$ 为有限生成理想，且
$A/I = B/IB$。则换基与标量限制诱导互为拟逆的等价
$D_{I^\infty\text{-torsion}}(A) = D_{(IB)^\infty\text{-torsion}}(B)$。
\end{lemma}

\begin{proof}
更准确地说，这两个函子分别为：对
$D_{I^\infty\text{-torsion}}(A)$ 中的 $K$，取
$K \mapsto K \otimes_A^\mathbf{L} B$；对
$D_{(IB)^\infty\text{-torsion}}(B)$ 中的 $M$，取 $M \mapsto M_A$。
其成立的原因是
$H^i(K \otimes_A^\mathbf{L} B) = H^i(K) \otimes_A B = H^i(K)$。
第一个等式由 $A \to B$ 平坦得出，第二个等式由《更多代数》引理
\ref{more-algebra-lemma-neighbourhood-isomorphism} 得出。
\end{proof}

\noindent
《更多代数》引理 \ref{more-algebra-lemma-neighbourhood-equivalence} 与
\ref{more-algebra-lemma-neighbourhood-extensions}
已经对模的 $\Hom$ 和 $\Ext^1$ 证明了以下引理。

\begin{lemma}
\label{dualizing-lemma-neighbourhood-extensions}
设 $A \to B$ 为平坦环同态，$I \subset A$ 为有限生成理想，且
$A/I \to B/IB$ 是同构。对 $K \in D_{I^\infty\text{-torsion}}(A)$
及 $L \in D(A)$，映射
$$
R\Hom_A(K, L) \longrightarrow R\Hom_B(K \otimes_A B, L \otimes_A B)
$$
是拟同构。特别地，若 $M$、$N$ 为 $A$-模，且 $M$ 是 $I$-幂挠的，
则对所有 $i$，典范映射
$$
\Ext^i_A(M, N)
\longrightarrow
\Ext^i_B(M \otimes_A B, N \otimes_A B)
$$
都是同构。
\end{lemma}

\begin{proof}
令 $Z = V(I) \subset \Spec(A)$、$Y = V(IB) \subset \Spec(B)$。
由于 $K$ 的上同调模是 $I$-幂挠的，典范映射
$R\Gamma_Z(L) \to L$ 在 $D(A)$ 中诱导同构
$$
R\Hom_A(K, R\Gamma_Z(L)) \to R\Hom_A(K, L)
$$
类似地，$K \otimes_A B$ 的上同调模是 $IB$-幂挠的，并且在
$D(B)$ 中有同构
$$
R\Hom_B(K \otimes_A B, R\Gamma_Y(L \otimes_A B)) \to 
R\Hom_B(K \otimes_A B, L \otimes_A B)
$$
由引理 \ref{dualizing-lemma-torsion-change-rings}，
$R\Gamma_Z(L) \otimes_A B = R\Gamma_Y(L \otimes_A B)$。
因此只需证明映射
$$
R\Hom_A(K, R\Gamma_Z(L)) \to R\Hom_B(K \otimes_A B, R\Gamma_Z(L) \otimes_A B)
$$
是拟同构。这由引理 \ref{dualizing-lemma-torsion-flat-change-rings} 得出。
\end{proof}




\section{Noetherian 环上的局部上同调}
\label{dualizing-section-local-cohomology-noetherian}

\noindent
设 $A$ 为环，$I \subset A$ 为有限生成理想。令
$Z = V(I) \subset \Spec(A)$。回忆（\ref{dualizing-equation-compare-torsion}）是函子
$$
D(I^\infty\text{-torsion}) \to D_{I^\infty\text{-torsion}}(A)
$$
事实上，存在函子的自然变换
\begin{equation}
\label{dualizing-equation-compare-torsion-functors}
(\ref{dualizing-equation-compare-torsion}) \circ R\Gamma_I(-)
\longrightarrow
R\Gamma_Z(-)
\end{equation}
具体地，给定 $A$-模复形 $K^\bullet$，$D(A)$ 中的典范映射
$R\Gamma_I(K^\bullet) \to K^\bullet$ 唯一地经由
$R\Gamma_Z(K^\bullet)$ 分解，因为 $R\Gamma_I(K^\bullet)$ 的上同调模
是 $I$-幂挠的（见引理 \ref{dualizing-lemma-adjoint}）。
一般而言，此映射不是同构（引理 \ref{dualizing-lemma-not-equal} 已给出这一点）。

\begin{lemma}
\label{dualizing-lemma-local-cohomology-noetherian}
设 $A$ 为 Noetherian 环，$I \subset A$ 为理想。
\begin{enumerate}
\item 对 $K \in D_{I^\infty\text{-torsion}}(A)$，伴随映射
$R\Gamma_I(K) \to K$ 是同构；
\item 函子（\ref{dualizing-equation-compare-torsion}）
$D(I^\infty\text{-torsion}) \to D_{I^\infty\text{-torsion}}(A)$
是等价；
\item 函子变换（\ref{dualizing-equation-compare-torsion-functors}）是同构；
换言之，对 $K \in D(A)$ 有 $R\Gamma_I(K) = R\Gamma_Z(K)$。
\end{enumerate}
\end{lemma}

\begin{proof}
一个省略的形式论证表明，只需证明 (1)。

\medskip\noindent
设 $M$ 为 $I$-幂挠 $A$-模。选取到内射 $A$-模中的嵌入
$M \to J$。由引理 \ref{dualizing-lemma-injective-module-divide}，
$J[I^\infty]$ 是内射 $A$-模，因而得到嵌入
$M \to J[I^\infty]$。所以每个 $I$-幂挠模都有内射分解
$M \to J^\bullet$，其中 $J^n$ 也都是 $I$-幂挠的。因此
$R\Gamma_I(M) = M$（这并非显然；若 $A$ 不是 Noetherian 环，
一般也不成立）。接着设
$K \in D_{I^\infty\text{-torsion}}^+(A)$。此时谱序列
$$
R^q\Gamma_I(H^p(K)) \Rightarrow R^{p + q}\Gamma_I(K)
$$
（《导出范畴》引理 \ref{derived-lemma-two-ss-complex-functor}）
收敛，并且由上文可知，只有 $q = 0$ 的项可能非零。
故 $R\Gamma_I(K) \to K$ 是同构。

\medskip\noindent
再设 $K$ 是 $D_{I^\infty\text{-torsion}}(A)$ 的任意对象。
由引理 \ref{dualizing-lemma-local-cohomology-ext} 有
$$
R^q\Gamma_I(K) = \colim \Ext^q_A(A/I^n, K)
$$
选取生成 $I$ 的 $f_1, \ldots, f_r \in A$。令
$K_n^\bullet = K(A, f_1^n, \ldots, f_r^n)$ 为次数位于
$-r, \ldots, 0$ 的 Koszul 复形。由《更多代数》引理
\ref{more-algebra-lemma-sequence-Koszul-complexes}，
$D(A)$ 中的 pro-对象 $\{A/I^n\}$ 与 $\{K_n^\bullet\}$ 相同，故
$$
R^q\Gamma_I(K) = \colim \Ext^q_A(K_n^\bullet, K)
$$
任取表示 $K$ 的 $A$-模复形 $K^\bullet$。由于 $K_n^\bullet$
是由有限自由模组成的有限复形，有
$$
\Ext^q_A(K_n, K) =
H^q(\text{Tot}((K_n^\bullet)^\vee \otimes_A K^\bullet))
$$
其中 $(K_n^\bullet)^\vee$ 是复形 $K_n^\bullet$ 的对偶。
见《更多代数》引理 \ref{more-algebra-lemma-RHom-out-of-projective}。
由于 $(K_n^\bullet)^\vee$ 是次数位于 $0, \ldots, r$、由有限自由
$A$-模组成的复形，复形
$\text{Tot}((K_n^\bullet)^\vee \otimes_A K^\bullet)$ 与复形
$\text{Tot}((K_n^\bullet)^\vee \otimes_A \tau_{\geq q - r - 2} K^\bullet)$
在次数 $q - 1$ 及更高次数上的项相同。因此
$$
\Ext^q_A(K_n, K) = \text{Ext}^q_A(K_n, \tau_{\geq q - r - 2}K)
$$
对所有 $n$ 成立。于是
$$
R^q\Gamma_I(K) = R^q\Gamma_I(\tau_{\geq q - r - 2}K) =
H^q(\tau_{\geq q - r - 2}K) = H^q(K)
$$
故映射 $R\Gamma_I(K) \to K$ 是同构。
\end{proof}

\begin{lemma}
\label{dualizing-lemma-compute-local-cohomology-noetherian}
设 $A$ 为 Noetherian 环，$I = (f_1, \ldots, f_r)$ 为 $A$ 的理想。
令 $Z = V(I) \subset \Spec(A)$。在 $D(A)$ 中存在典范同构
$$
R\Gamma_I(A) \to
(A \to \prod\nolimits_{i_0} A_{f_{i_0}} \to
\prod\nolimits_{i_0 < i_1} A_{f_{i_0}f_{i_1}} \to
\ldots \to A_{f_1\ldots f_r}) \to R\Gamma_Z(A)
$$
若 $M$ 为 $A$-模，则类似地，在 $D(A)$ 中有
$$
R\Gamma_I(M) \cong
(M \to \prod\nolimits_{i_0} M_{f_{i_0}} \to
\prod\nolimits_{i_0 < i_1} M_{f_{i_0}f_{i_1}} \to
\ldots \to M_{f_1\ldots f_r}) \cong R\Gamma_Z(M)
$$
\end{lemma}

\begin{proof}
这由引理 \ref{dualizing-lemma-local-cohomology-noetherian} 以及引理
\ref{dualizing-lemma-local-cohomology-adjoint} 中对函子 $R\Gamma_Z$ 的计算得出。
\end{proof}

\begin{lemma}
\label{dualizing-lemma-local-cohomology-change-rings}
若 $A \to B$ 是 Noetherian 环的同态，且 $I \subset A$ 为理想，
则在 $D(B)$ 中有
$$
R\Gamma_I(A) \otimes_A^\mathbf{L} B =
R\Gamma_Z(A) \otimes_A^\mathbf{L} B =
R\Gamma_Y(B) = R\Gamma_{IB}(B)
$$
其中 $Y = V(IB) \subset \Spec(B)$。
\end{lemma}

\begin{proof}
合并引理 \ref{dualizing-lemma-compute-local-cohomology-noetherian} 与
\ref{dualizing-lemma-torsion-change-rings} 即得。
\end{proof}






\section{深度}
\label{dualizing-section-depth}

\noindent
本节重新考察《代数》第 \ref{algebra-section-depth} 节引入的深度概念。

\begin{lemma}
\label{dualizing-lemma-depth}
设 $A$ 为 Noetherian 环，$I \subset A$ 为理想，$M$ 为有限
$A$-模，且 $IM \not = M$。则以下整数相等：
\begin{enumerate}
\item $\text{depth}_I(M)$；
\item 使 $\Ext_A^i(A/I, M)$ 非零的最小整数 $i$；
\item 使 $H^i_I(M)$ 非零的最小整数 $i$。
\end{enumerate}
此外，对任意被 $I$ 的某次幂零化的有限 $A$-模 $N$，当
$i < \text{depth}_I(M)$ 时有 $\Ext^i_A(N, M) = 0$。
\end{lemma}

\begin{proof}
由《代数》引理 \ref{algebra-lemma-depth-finite-noetherian}，可以对
$\text{depth}_I(M)$ 归纳来证明 (1) 与 (2) 相等。

\medskip\noindent
先看基础情形。若 $\text{depth}_I(M) = 0$，则 $I$ 包含于 $M$
的相伴素理想之并中（《代数》引理
\ref{algebra-lemma-ass-zero-divisors}）。
由素理想避开引理（《代数》引理 \ref{algebra-lemma-silly}），
对某个相伴素理想 $\mathfrak p$ 有 $I \subset \mathfrak p$。
因此 $\Hom_A(A/I, M)$ 非零，故此时等式成立。

\medskip\noindent
设 $\text{depth}_I(M) > 0$。取 $f \in I$，使它是 $M$ 上的非零因子，
并且 $\text{depth}_I(M/fM) = \text{depth}_I(M) - 1$。考察短正合序列
$$
0 \to M \to M \to M/fM \to 0
$$
以及相伴的 $\Ext^*_A(A/I, -)$ 长正合序列。注意
$\Ext^i_A(A/I, M)$ 是有限 $A/I$-模（《代数》引理
\ref{algebra-lemma-ext-noetherian} 与
\ref{algebra-lemma-annihilate-ext}）。因此得到
$$
\Hom_A(A/I, M/fM) = \Ext^1_A(A/I, M)
$$
以及短正合序列
$$
0 \to \Ext^i_A(A/I, M) \to \text{Ext}^i_A(A/I, M/fM) \to
\Ext^{i + 1}_A(A/I, M) \to 0
$$
于是由归纳得到 (1) 与 (2) 相等。

\medskip\noindent
注意，对所有 $n \geq 1$ 都有
$\text{depth}_I(M) = \text{depth}_{I^n}(M)$；例如这由《代数》引理
\ref{algebra-lemma-regular-sequence-powers} 得出。因此由 (1) 与 (2)
相等可知，对所有 $n$ 以及 $i < \text{depth}_I(M)$，有
$\Ext^i_A(A/I^n, M) = 0$。设 $N$ 为被 $I$ 的某次幂零化的有限
$A$-模。则对某些 $n, m \geq 0$，可选取短正合序列
$$
0 \to N' \to (A/I^n)^{\oplus m} \to N \to 0
$$
此时
$\Hom_A(N, M) \subset \Hom_A((A/I^n)^{\oplus m}, M)$，
并且当 $i < \text{depth}_I(M)$ 时，
$\Ext^i_A(N, M) \subset \text{Ext}^{i - 1}_A(N', M)$。
于是一个简单的归纳论证表明，该引理的最后一个陈述成立。

\medskip\noindent
最后证明 (3) 等于 (1) 与 (2)。由引理
\ref{dualizing-lemma-local-cohomology-ext} 有
$H^p_I(M) = \colim \Ext^p_A(A/I^n, M)$。
所以当 $i < \text{depth}_I(M)$ 时 $H^i_I(M) = 0$。
当 $i = \text{depth}_I(M)$ 时，利用
$\Ext_A^{i - 1}(I/I^n, M)$ 的消失可知，映射
$\Ext_A^i(A/I, M) \to H_I^i(M)$ 是单射，从而证明在正确次数上非零。
\end{proof}

\begin{lemma}
\label{dualizing-lemma-depth-in-ses}
设 $A$ 为 Noetherian 环，$0 \to N' \to N \to N'' \to 0$
为有限 $A$-模的短正合序列，$I \subset A$ 为理想。则
\begin{enumerate}
\item
$\text{depth}_I(N) \geq \min\{\text{depth}_I(N'), \text{depth}_I(N'')\}$；
\item
$\text{depth}_I(N'') \geq \min\{\text{depth}_I(N), \text{depth}_I(N') - 1\}$；
\item
$\text{depth}_I(N') \geq \min\{\text{depth}_I(N), \text{depth}_I(N'') + 1\}$。
\end{enumerate}
\end{lemma}

\begin{proof}
设 $IN \not = N$、$IN' \not = N'$ 且 $IN'' \not = N''$。
此时可以使用引理 \ref{dualizing-lemma-depth} 中由 Ext 群
$\Ext^i(A/I, N)$ 给出的深度刻画，并使用《代数》引理
\ref{algebra-lemma-long-exact-seq-ext} 的上同调长正合序列
$$
\begin{matrix}
0
\to \Hom_A(A/I, N')
\to \Hom_A(A/I, N)
\to \Hom_A(A/I, N'')
\\
\phantom{0\ }
\to \Ext^1_A(A/I, N')
\to \Ext^1_A(A/I, N)
\to \Ext^1_A(A/I, N'')
\to \ldots
\end{matrix}
$$
若 $IN = N$，该论证也成立，因为此时对所有 $i$ 都有
$\Ext^i_A(A/I, N) = 0$。在 $IN' \not = N'$ 和/或
$IN'' \not = N''$ 的情形同理。
\end{proof}

\begin{lemma}
\label{dualizing-lemma-depth-drops-by-one}
设 $A$ 为 Noetherian 环，$I \subset A$ 为理想，$M$ 为有限
$A$-模，且 $IM \not = M$。
\begin{enumerate}
\item 若 $x \in I$ 是 $M$ 上的非零因子，则
$\text{depth}_I(M/xM) = \text{depth}_I(M) - 1$。
\item $I$ 中任意 $M$-正则序列 $x_1, \ldots, x_r$ 都可扩充为
$I$ 中长度为 $\text{depth}_I(M)$ 的 $M$-正则序列。
\end{enumerate}
\end{lemma}

\begin{proof}
(2) 是 (1) 的形式推论。设 $x \in I$ 如 (1) 所述。由短正合序列
$0 \to M \to M \to M/xM \to 0$ 及引理 \ref{dualizing-lemma-depth-in-ses}，
有 $\text{depth}_I(M/xM) \geq \text{depth}_I(M) - 1$。
另一方面，若 $x_1, \ldots, x_r \in I$ 是 $M/xM$ 的正则序列，
则 $x, x_1, \ldots, x_r$ 是 $M$ 的正则序列。因此 (1) 成立。
\end{proof}

\begin{lemma}
\label{dualizing-lemma-depth-CM}
设 $R$ 为 Noetherian 局部环。若 $M$ 是有限 Cohen--Macaulay
$R$-模，且 $I \subset R$ 为非平凡理想，则
$$
\text{depth}_I(M) = \dim(\text{Supp}(M)) - \dim(\text{Supp}(M/IM)).
$$
\end{lemma}

\begin{proof}
对 $\text{depth}_I(M)$ 归纳。

\medskip\noindent
若 $\text{depth}_I(M) = 0$，则 $I$ 包含于 $M$ 的某个相伴素理想
$\mathfrak p$ 中（《代数》引理
\ref{algebra-lemma-ideal-nonzerodivisor}）。于是
$\mathfrak p \in \text{Supp}(M/IM)$，从而
$\dim(\text{Supp}(M/IM)) \geq \dim(R/\mathfrak p) = \dim(\text{Supp}(M))$，
其中等式由《代数》引理
\ref{algebra-lemma-CM-ass-minimal-support} 成立。故此时引理成立。

\medskip\noindent
若 $\text{depth}_I(M) > 0$，取 $x \in I$，使它是 $M$ 上的非零因子。
注意 $(M/xM)/I(M/xM) = M/IM$。另一方面，由引理
\ref{dualizing-lemma-depth-drops-by-one} 有
$\text{depth}_I(M/xM) = \text{depth}_I(M) - 1$；由《代数》引理
\ref{algebra-lemma-one-equation-module} 有
$\dim(\text{Supp}(M/xM)) = \dim(\text{Supp}(M)) -  1$。
因此由归纳假设得到结论。
\end{proof}

\begin{lemma}
\label{dualizing-lemma-depth-flat-CM}
设 $R \to S$ 为 Noetherian 局部环之间的平坦局部环同态。
以 $\mathfrak m \subset R$ 表示极大理想，设 $I \subset S$ 为理想。
若 $S/\mathfrak mS$ 是 Cohen--Macaulay 的，则
$$
\text{depth}_I(S) \geq \dim(S/\mathfrak mS) - \dim(S/\mathfrak mS + I)
$$
\end{lemma}

\begin{proof}
由《代数》引理 \ref{algebra-lemma-grothendieck-regular-sequence}，
$S$ 中任何在 $S/\mathfrak mS$ 中映为正则序列的序列，都是 $S$
中的正则序列。因此只需在 $R$ 为域的情形证明该引理；
这是引理 \ref{dualizing-lemma-depth-CM} 的特殊情形。
\end{proof}

\begin{lemma}
\label{dualizing-lemma-divide-by-torsion}
设 $A$ 为环，$I \subset A$ 为有限生成理想，$M$ 为 $A$-模。
令 $Z = V(I)$。则 $H^0_I(M) = H^0_Z(M)$。以 $N$ 表示这个共同的值，
并令 $M' = M/N$。则
\begin{enumerate}
\item 对所有 $p > 0$，有 $H^0_I(M') = 0$、
$H^p_I(M) = H^p_I(M')$ 及 $H^p_I(N) = 0$；
\item 对所有 $p > 0$，有 $H^0_Z(M') = 0$、
$H^p_Z(M) = H^p_Z(M')$ 及 $H^p_Z(N) = 0$。
\end{enumerate}
\end{lemma}

\begin{proof}
按定义，$H^0_I(M) = M[I^\infty]$ 是 $I$-幂挠的。
由引理 \ref{dualizing-lemma-local-cohomology-adjoint}，若
$I = (f_1, \ldots, f_r)$，则
$$
H^0_Z(M) = \Ker(M \longrightarrow M_{f_1} \times \ldots \times M_{f_r})
$$
因此 $H^0_I(M) \subset H^0_Z(M)$。反之，若 $x \in H^0_Z(M)$，
则对某个 $e_i \geq 1$，它被某个 $f_i^{e_i}$ 零化，因而被 $I$
的某次幂零化。这证明了第一个等式，并且 $N$ 是 $I$-幂挠的。
由引理 \ref{dualizing-lemma-adjoint} 有 $R\Gamma_I(N) = N$；
由引理 \ref{dualizing-lemma-local-cohomology-adjoint} 有 $R\Gamma_Z(N) = N$。
这证明了 (1) 与 (2) 中 $H^p_I(N)$ 和 $H^p_Z(N)$ 在高次的消失。
由上述论证以及《更多代数》引理
\ref{more-algebra-lemma-divide-by-torsion} 所给出的 $M'$ 为
$I$-幂无挠这一事实，得到 $H^0_I(M')$ 与 $H^0_Z(M')$ 的消失。
$M$ 与 $M'$ 的高次上同调相等，则立即由上同调长正合序列得出。
\end{proof}









\section{挠模与完备模}
\label{dualizing-section-torsion-and-complete}

\noindent
设 $A$ 为环，$I$ 为有限生成理想。此时可以考察上同调模为
$I$-幂挠模的复形所构成的导出范畴
$D_{I^\infty\text{-torsion}}(A)$（第 \ref{dualizing-section-local-cohomology} 节），
以及导出完备复形的导出范畴 $D_{comp}(A, I)$（《更多代数》第
\ref{more-algebra-section-derived-completion} 节）。本节证明这两个范畴等价。
更一般的陈述见 \cite{Dwyer-Greenlees}。

\begin{lemma}
\label{dualizing-lemma-complete-and-local}
\begin{slogan}
这类结果有时称为 Greenlees--May 对偶。
\end{slogan}
设 $A$ 为环，$I$ 为有限生成理想。令 $R\Gamma_Z$ 如引理
\ref{dualizing-lemma-local-cohomology-adjoint} 所述；令 ${\ }^\wedge$
表示《更多代数》引理 \ref{more-algebra-lemma-derived-completion}
中的导出完备化。对 $D(A)$ 中的对象 $K$，在 $D(A)$ 中有
$$
R\Gamma_Z(K^\wedge) = R\Gamma_Z(K)
\quad\text{且}\quad
(R\Gamma_Z(K))^\wedge = K^\wedge
$$
\end{lemma}

\begin{proof}
选取生成 $I$ 的 $f_1, \ldots, f_r \in A$。回忆由《更多代数》引理
\ref{more-algebra-lemma-derived-completion} 有
$$
K^\wedge = R\Hom_A\left((A \to \prod A_{f_{i_0}}
\to \prod A_{f_{i_0i_1}} \to \ldots \to A_{f_1 \ldots f_r}), K\right)
$$
因此锥 $C = \text{Cone}(K \to K^\wedge)$ 由
$$
R\Hom_A\left((\prod A_{f_{i_0}}
\to \prod A_{f_{i_0i_1}} \to \ldots \to A_{f_1 \ldots f_r}), K\right)
$$
给出。它可由一个带有限滤过的复形表示，其相继商同构于
$$
R\Hom_A(A_{f_{i_0} \ldots f_{i_p}}, K), \quad p > 0
$$
对这些复形施用 $R\Gamma_Z$ 后均消失，见引理
\ref{dualizing-lemma-local-cohomology-vanishes}。对可区别三角
$K \to K^\wedge \to C \to K[1]$ 施用 $R\Gamma_Z$，
即知该引理的第一个公式成立。

\medskip\noindent
回忆由引理 \ref{dualizing-lemma-local-cohomology-adjoint} 有
$$
R\Gamma_Z(K) =
K \otimes^\mathbf{L} (A \to \prod A_{f_{i_0}}
\to \prod A_{f_{i_0i_1}} \to \ldots \to A_{f_1 \ldots f_r})
$$
因此锥 $C = \text{Cone}(R\Gamma_Z(K) \to K)$ 可由一个带有限滤过的
复形表示，其相继商同构于
$$
K \otimes_A A_{f_{i_0} \ldots f_{i_p}}, \quad p > 0
$$
对这些复形施用 ${\ }^\wedge$ 后均消失，见《更多代数》引理
\ref{more-algebra-lemma-derived-completion-vanishes}。
对可区别三角 $R\Gamma_Z(K) \to K \to C \to R\Gamma_Z(K)[1]$
施用导出完备化，即知该引理的第二个公式成立。
\end{proof}

\noindent
以下结果是关于三角范畴的容许子范畴这一很一般现象的一个特殊情形。

\begin{proposition}
\label{dualizing-proposition-torsion-complete}
\begin{reference}
这是 \cite[Theorem 1.1]{Porta-Liran-Yekutieli} 的一个特殊情形。
\end{reference}
设 $A$ 为环，$I \subset A$ 为有限生成理想。函子
$R\Gamma_Z$ 与 ${\ }^\wedge$ 定义互为拟逆的范畴等价
$$
D_{I^\infty\text{-torsion}}(A) \leftrightarrow D_{comp}(A, I)
$$
\end{proposition}

\begin{proof}
立即由引理 \ref{dualizing-lemma-complete-and-local} 得出。
\end{proof}

\noindent
上述命题的以下补充更精确地描述了态射之间的对应。

\begin{lemma}
\label{dualizing-lemma-compare-RHom}
采用引理 \ref{dualizing-lemma-complete-and-local} 的记号。对 $D(A)$ 中的对象
$K, L$，在 $D(A)$ 中存在典范同构
$$
R\Hom_A(K^\wedge, L^\wedge) \longrightarrow R\Hom_A(R\Gamma_Z(K), R\Gamma_Z(L))
$$
\end{lemma}

\begin{proof}
设 $I = (f_1, \ldots, f_r)$。以
$C = (A \to \prod A_{f_i} \to \ldots \to A_{f_1 \ldots f_r})$
表示交错 {\v C}ech 复形。则导出完备化由
$R\Hom_A(C, -)$ 给出（《更多代数》引理
\ref{more-algebra-lemma-derived-completion}），局部上同调由
$C \otimes^\mathbf{L} -$ 给出（引理
\ref{dualizing-lemma-local-cohomology-adjoint}）。结合同构
$$
R\Hom_A(K \otimes^\mathbf{L} C, L \otimes^\mathbf{L} C) =
R\Hom_A(K, R\Hom_A(C,  L \otimes^\mathbf{L} C))
$$
（《更多代数》引理 \ref{more-algebra-lemma-internal-hom}）与映射
$$
L \to R\Hom_A(C,  L \otimes^\mathbf{L} C)
$$
（《更多代数》引理 \ref{more-algebra-lemma-internal-hom-diagonal}），
得到映射
$$
\gamma :
R\Hom_A(K, L)
\longrightarrow
R\Hom_A(K \otimes^\mathbf{L} C, L \otimes^\mathbf{L} C)
$$
另一方面，右侧等于
$$
R\Hom_A(C, R\Hom_A(K, L \otimes^\mathbf{L} C)).
$$
因而是导出完备的。由导出完备化的泛性质，$\gamma$ 经由
$R\Hom_A(K, L)$ 的导出完备化分解。然而，由《更多代数》引理
\ref{more-algebra-lemma-completion-RHom}，导出完备化可移入
$R\Hom_A$ 内，于是得到所需映射。

\medskip\noindent
为证明该引理的映射是同构，可以假设 $K$ 与 $L$ 都是导出完备的，
即 $K = K^\wedge$ 且 $L = L^\wedge$。此时考察映射
$$
\gamma : R\Hom_A(K, L) \longrightarrow R\Hom_A(R\Gamma_Z(K), R\Gamma_Z(L))
$$
由命题 \ref{dualizing-proposition-torsion-complete}，左侧与右侧的上同调群相同。
换言之，只需验证映射 $\gamma$ 将 $D(A)$ 中的态射
$\alpha : K \to L$ 送到态射
$R\Gamma_Z(\alpha) : R\Gamma_Z(K) \to R\Gamma_Z(L)$。
我们省略验证（提示：注意 $R\Gamma_Z(\alpha)$ 正是映射
$\alpha \otimes \text{id}_C :
K \otimes^\mathbf{L} C
\to
L \otimes^\mathbf{L} C$；它与《更多代数》引理
\ref{more-algebra-lemma-internal-hom-diagonal} 中映射的构造几乎相同）。
\end{proof}

\begin{lemma}
\label{dualizing-lemma-completion-local}
设 $I$ 与 $J$ 为 Noetherian 环 $A$ 的理想，$M$ 为有限 $A$-模。
令 $Z =V(J)$。考察局部上同调的导出 $I$-进完备化
$R\Gamma_Z(M)^\wedge$。则
\begin{enumerate}
\item 有 $R\Gamma_Z(M)^\wedge = R\lim R\Gamma_Z(M/I^nM)$；
\item 存在短正合序列
$$
0 \to R^1\lim H^{i - 1}_Z(M/I^nM) \to H^i(R\Gamma_Z(M)^\wedge) \to
\lim H^i_Z(M/I^nM) \to 0
$$
\end{enumerate}
特别地，$R\Gamma_Z(M)^\wedge$ 在负次数上的上同调消失。
\end{lemma}

\begin{proof}
设 $J = (g_1, \ldots, g_m)$。由引理
\ref{dualizing-lemma-local-cohomology-adjoint}，$R\Gamma_Z(M)$ 由复形
$$
M \to \prod M_{g_{j_0}} \to \prod M_{g_{j_0}g_{j_1}} \to
\ldots \to M_{g_1g_2\ldots g_m}
$$
计算。由《更多代数》引理
\ref{more-algebra-lemma-when-derived-completion-is-completion}，
此复形的导出 $I$-进完备化由通常完备化构成的复形
$$
\lim M/I^nM \to \prod \lim (M/I^nM)_{g_{j_0}} \to
\ldots \to \lim (M/I^nM)_{g_1g_2\ldots g_m}
$$
给出。又因为 $R\Gamma_Z(M/I^nM)$ 由复形
$ M/I^nM \to \prod (M/I^nM)_{g_{j_0}} \to
\ldots \to (M/I^nM)_{g_1g_2\ldots g_m}$ 计算，且这些复形之间的
转移映射都是满射，所以由《更多代数》引理
\ref{more-algebra-lemma-compute-Rlim-modules} 得到 (1)。
(2) 随后由《更多代数》引理
\ref{more-algebra-lemma-break-long-exact-sequence-modules} 得出。
\end{proof}

\begin{lemma}
\label{dualizing-lemma-completion-local-H0}
采用引理 \ref{dualizing-lemma-completion-local} 的记号与假设，并设 $A$ 是
$I$-进完备的。则
$$
H^0(R\Gamma_Z(M)^\wedge) = \colim H^0_{V(J')}(M)
$$
其中滤过余极限取遍满足 $V(J') \cap V(I) = V(J) \cap V(I)$ 的
$J' \subset J$。
\end{lemma}

\begin{proof}
由于 $M$ 是有限 $A$-模，$M$ 是 $I$-进完备的。
引理 \ref{dualizing-lemma-completion-local} 的证明表明
$$
H^0(R\Gamma_Z(M)^\wedge) =
\Ker(M^\wedge \to \prod M_{g_j}^\wedge) =
\Ker(M \to \prod M_{g_j}^\wedge)
$$
其中右侧采用通常的 $I$-进完备化。映射
$M_{g_j} \to M_{g_j}^\wedge$ 的核 $K_j$ 是
$\bigcap I^n M_{g_j}$。由《代数》引理
\ref{algebra-lemma-intersection-powers-ideal-module}，对每个
$\mathfrak p \in V(IA_{g_j})$，都可找到
$f \in A_{g_j}$、$f \not \in \mathfrak p$，使得 $(K_j)_f = 0$。

\medskip\noindent
设 $s \in H^0(R\Gamma_Z(M)^\wedge)$。由上文，可将 $s$ 视为 $M$
的元素。由上述论证，$s$ 的支撑 $Z'$ 与 $D(g_j)$ 的交同
$D(g_j) \cap V(I)$ 不交。因此 $Z'$ 是 $\Spec(A)$ 的闭子集，且
$Z' \cap V(I) \subset V(J)$。于是对某个满足 $J' \subset J$ 且
$V(J') \cap V(I) \subset V(J)$ 的理想，有
$Z' \cup V(J) = V(J')$，并且 $s \in H^0_{V(J')}(M)$。
反之，任意满足 $J' \subset J$ 与 $V(J') \cap V(I) \subset V(J)$ 的
$s \in H^0_{V(J')}(M)$，对所有 $j$，在 $M_{g_j}^\wedge$ 中的像都为零。
这就证明了该引理。
\end{proof}









\section{环同态的平凡对偶}
\label{dualizing-section-trivial}

\noindent
设 $A \to B$ 为环同态。考察函子
$$
\Hom_A(B, -) : \text{Mod}_A \longrightarrow \text{Mod}_B,\quad
M \longmapsto \Hom_A(B, M)
$$
此函子左正合，并有导出扩张
$R\Hom(B, -) : D(A) \to D(B)$。

\begin{lemma}
\label{dualizing-lemma-right-adjoint}
设 $A \to B$ 为环同态。上面构造的函子 $R\Hom(B, -)$
是标量限制函子 $D(B) \to D(A)$ 的右伴随。
\end{lemma}

\begin{proof}
由《代数》引理 \ref{algebra-lemma-adjoint-hom-restrict}，标量限制与
$\Hom_A(B, -)$ 是伴随函子，故结论成立。另见《导出范畴》引理
\ref{derived-lemma-derived-adjoint-functors}。
\end{proof}

\begin{lemma}
\label{dualizing-lemma-composition-right-adjoints}
设 $A \to B \to C$ 为环同态。则函子
$R\Hom(C, -) \circ R\Hom(B, -) : D(A) \to D(C)$
就是 $R\Hom(C, -) : D(A) \to D(C)$。
\end{lemma}

\begin{proof}
由右伴随的唯一性及引理 \ref{dualizing-lemma-right-adjoint} 得出。
\end{proof}

\begin{lemma}
\label{dualizing-lemma-RHom-ext}
设 $\varphi : A \to B$ 为环同态。对 $D(A)$ 中的 $K$ 有
$$
\varphi_*R\Hom(B, K) = R\Hom_A(B, K)
$$
其中 $\varphi_* : D(B) \to D(A)$ 是标量限制。特别地，
$R^q\Hom(B, K) = \Ext_A^q(B, K)$。
\end{lemma}

\begin{proof}
选取表示 $K$ 的 K-内射复形 $I^\bullet$。则 $R\Hom(B, K)$
由 $B$-模复形 $\Hom_A(B, I^\bullet)$ 表示。将它视为 $A$-模复形时，
它表示 $R\Hom_A(B, K)$，故引理成立。
\end{proof}

\noindent
设 $A$ 为 Noetherian 环。以
$$
D_{\textit{Coh}}(A) \subset D(A)
$$
表示如下全子范畴：其对象是 $D(A)$ 中上同调模均为有限 $A$-模的对象
$K$。这一定义成立，见《导出范畴》第
\ref{derived-section-triangulated-sub} 节；因为 $A$ 是 Noetherian 环，
有限 $A$-模构成 $\text{Mod}_A$ 的 Serre 子范畴。

\begin{lemma}
\label{dualizing-lemma-exact-support-coherent}
采用上述记号，并设 $A \to B$ 是 Noetherian 环之间的有限环同态。
则 $R\Hom(B, -)$ 将 $D^+_{\textit{Coh}}(A)$ 映入
$D^+_{\textit{Coh}}(B)$。
\end{lemma}

\begin{proof}
需要证明：若 $K \in D^+(A)$ 的上同调模有限，则复形
$R\Hom(B, K)$ 的上同调模也有限。由引理 \ref{dualizing-lemma-RHom-ext}，
只要证明 Ext 模 $\Ext^i_A(B, K)$ 是有限 $A$-模即可。
由于 $K$ 下有界，存在收敛谱序列
$$
\Ext^p_A(B, H^q(K)) \Rightarrow \text{Ext}^{p + q}_A(B, K)
$$
而由《代数》引理 \ref{algebra-lemma-ext-noetherian}，
模 $\Ext^p_A(B, H^q(K))$ 有限，故证明完成。
\end{proof}

\begin{remark}
\label{dualizing-remark-exact-support}
设 $A$ 为环，$I \subset A$ 为理想，并令 $B = A/I$。
此时函子 $\Hom_A(B, -)$ 等于函子
$$
\text{Mod}_A \longrightarrow \text{Mod}_B,\quad M \longmapsto M[I]
$$
它将 $M$ 送到其 $I$-挠子模。
\end{remark}

\begin{situation}
\label{dualizing-situation-resolution}
设 $R \to A$ 为环同态。我们给出 $R\Hom(A, -)$ 的另一种构造，
它将在本章后文中很有用。具体地，设有 $R$ 上的微分分次代数
$(E, d)$ 以及拟同构 $E \to A$，其中将 $A$ 视为微分为零的
$R$ 上微分分次代数。则有交换图
$$
\vcenter{
\xymatrix{
D(E, \text{d}) \ar[rd] & &  D(A) \ar[ll] \ar[ld] \\
& D(R)
}
}
\quad\text{且}\quad
\vcenter{
\xymatrix{
D(E, \text{d}) \ar[rr]_{- \otimes_E^\mathbf{L} A} & &  D(A) \\
& D(R) \ar[lu]^{- \otimes_R^\mathbf{L} E} \ar[ru]_{- \otimes_R^\mathbf{L} A}
}
}
$$
其中水平箭头是范畴等价（《微分分次代数》引理
\ref{dga-lemma-qis-equivalence}）。第一个图显然交换。
第二个图也交换：一方面，第一个图交换，而这里的函子是其中函子的左伴随
（《微分分次代数》例 \ref{dga-example-map-hom-tensor}）；另一方面，
也可使用 $E \otimes^\mathbf{L}_E A = E \otimes_E A$
及《微分分次代数》引理
\ref{dga-lemma-compose-tensor-functors-general}。
\end{situation}

\begin{lemma}
\label{dualizing-lemma-RHom-dga}
在情形 \ref{dualizing-situation-resolution} 中，函子 $R\Hom(A, -)$
等于函子 $R\Hom(E, -) : D(R) \to D(E, \text{d})$
与等价 $- \otimes^\mathbf{L}_E A : D(E, \text{d}) \to D(A)$ 的复合。
\end{lemma}

\begin{proof}
这是因为 $R\Hom(E, -)$ 是 $- \otimes^\mathbf{L}_R E$ 的右伴随，
见《微分分次代数》引理 \ref{dga-lemma-tensor-hom-adjoint}。
因此，此函子对环同态 $R \to A$ 所起的作用与函子
$R\Hom(A, -)$ 相同（引理 \ref{dualizing-lemma-right-adjoint}），故二者必经由等价
$- \otimes^\mathbf{L}_E A : D(E, \text{d}) \to D(A)$ 对应。
\end{proof}

\begin{lemma}
\label{dualizing-lemma-RHom-is-tensor}
在情形 \ref{dualizing-situation-resolution} 中，假设
\begin{enumerate}
\item 将 $E$ 视为 $D(R)$ 的对象时，它是紧对象；
\item $N = \Hom^\bullet_R(E^\bullet, R)$ 计算 $R\Hom(E, R)$。
\end{enumerate}
则 $R\Hom(E, -) : D(R) \to D(E)$ 同构于
$K \mapsto K \otimes_R^\mathbf{L} N$。
\end{lemma}

\begin{proof}
这是《微分分次代数》引理
\ref{dga-lemma-RHom-is-tensor} 的特殊情形。
\end{proof}

\begin{lemma}
\label{dualizing-lemma-RHom-is-tensor-special}
在情形 \ref{dualizing-situation-resolution} 中，设 $A$ 是完美 $R$-模。则
$$
R\Hom(A, -) : D(R) \to D(A)
$$
由 $K \mapsto K \otimes_R^\mathbf{L} M$ 给出，其中
$M = R\Hom(A, R) \in D(A)$。
\end{lemma}

\begin{proof}
应用《分次幂代数》引理
\ref{dpa-lemma-tate-resoluton-pseudo-coherent-ring-map}，
选取 $A$ 在 $R$ 上的 Tate 分解 $(E, \text{d})$。
注意，当 $i > 0$ 时 $E^i = 0$；$E^0 = R[x_1, \ldots, x_n]$
是多项式代数；当 $i < 0$ 时，$E^i$ 是有限自由 $E^0$-模。
因此，将 $E$ 视为 $R$-模复形时，它是自由 $R$-模的上有界复形。
验证引理 \ref{dualizing-lemma-RHom-is-tensor} 的假设。第一条成立是因为
$A$ 完美（故由《更多代数》命题
\ref{more-algebra-proposition-perfect-is-compact}，它也是紧对象）；
第二条由《更多代数》引理
\ref{more-algebra-lemma-RHom-out-of-projective} 成立。
由该引理，对某个微分分次 $E$-模 $N$，函子
$K \mapsto R\Hom(E, K)$ 同构于
$K \mapsto K \otimes_R^\mathbf{L} N$。注意在 $D(A)$ 中
$$
(R \otimes_R E) \otimes_E^\mathbf{L} A = R \otimes_E E \otimes_E A
$$
因此由《微分分次代数》引理
\ref{dga-lemma-compose-tensor-functors-general-algebra}，
$- \otimes_R^\mathbf{L} N$ 与 $- \otimes_R^\mathbf{L} A$ 的复合
具有 $- \otimes_R M$ 的形式，其中 $M \in D(A)$。
最后应用引理 \ref{dualizing-lemma-RHom-dga} 即完成证明。
\end{proof}

\begin{lemma}
\label{dualizing-lemma-compute-for-effective-Cartier-algebraic}
设 $R \to A$ 为满环同态，其核 $I$ 是可逆 $R$-模。函子
$R\Hom(A, -) : D(R) \to D(A)$ 同构于
$K \mapsto K \otimes_R^\mathbf{L} N[-1]$，其中 $N$ 是可逆
$A$-模 $I \otimes_R A$ 的逆。
\end{lemma}

\begin{proof}
由于 $A$ 有有限射影分解
$$
0 \to I \to R \to A \to 0
$$
可知 $A$ 是完美 $R$-模。由引理
\ref{dualizing-lemma-RHom-is-tensor-special}，只需证明 $D(A)$ 中
$R\Hom(A, R)$ 由 $N[-1]$ 表示。这意味着 $R\Hom(A, R)$ 只有一个
非零上同调模，即次数 $1$ 上的 $N$。由于
$\text{Mod}_A \to \text{Mod}_R$ 全忠实，只需在施用标量限制函子
$i_* : D(A) \to D(R)$ 后证明。由引理 \ref{dualizing-lemma-RHom-ext} 有
$$
i_*R\Hom(A, R) = R\Hom_R(A, R)
$$
利用上述有限射影分解可知，后者由复形
$R \to I^{\otimes -1}$ 表示，其中 $R$ 位于次数 $0$。
映射 $R \to I^{\otimes -1}$ 是单射，其余核为 $N$。
\end{proof}





\section{平凡对偶的换基}
\label{dualizing-section-base-change-trivial-duality}

\noindent
本节考察环的余笛卡尔方块
$$
\xymatrix{
A \ar[r]_\alpha & A' \\
R \ar[u]^\varphi \ar[r]^\rho & R' \ar[u]_{\varphi'}
}
$$
换言之，$A' = A \otimes_R R'$。若 $A$ 与 $R'$ 在 $R$ 上
{\bf Tor 独立}，则存在典范换基映射
\begin{equation}
\label{dualizing-equation-base-change}
R\Hom(A, K) \otimes_A^\mathbf{L} A'
\longrightarrow
R\Hom(A', K \otimes_R^\mathbf{L} R')
\end{equation}
它位于 $D(A')$ 中，并对 $D(R)$ 中的 $K$ 函子性成立。
事实上，由引理 \ref{dualizing-lemma-right-adjoint} 的伴随性，这样的箭头等价于
$D(R')$ 中的映射
$$
\varphi'_*\left(R\Hom(A, K) \otimes_A^\mathbf{L} A'\right)
\longrightarrow
K \otimes_R^\mathbf{L} R'
$$
其中 $\varphi'_* : D(A') \to D(R')$ 为标量限制函子。
对左侧应用《更多代数》引理
\ref{more-algebra-lemma-base-change-comparison}，可知这又等价于
$D(R')$ 中的映射
$$
\varphi_*(R\Hom(A, K)) \otimes_R^\mathbf{L} R'
\longrightarrow
K \otimes_R^\mathbf{L} R'
$$
其中 $\varphi_* : D(A) \to D(R)$ 为标量限制函子。
可取 $can \otimes^\mathbf{L} \text{id}_{R'}$，其中
$can : \varphi_*(R\Hom(A, K)) \to K$ 是 $R\Hom(A, -)$ 与
$\varphi_*$ 之间伴随的余单位。

\begin{lemma}
\label{dualizing-lemma-check-base-change-is-iso}
在上述情形中，映射（\ref{dualizing-equation-base-change}）是同构，当且仅当
《更多代数》引理
\ref{more-algebra-lemma-internal-hom-diagonal-better} 中的映射
$$
R\Hom_R(A, K) \otimes_R^\mathbf{L} R'
\longrightarrow
R\Hom_R(A, K \otimes_R^\mathbf{L} R')
$$
是同构。
\end{lemma}

\begin{proof}
为证明该映射是同构，只需证明施用 $\varphi'_*$ 后它是同构。
对（\ref{dualizing-equation-base-change}）施用函子 $\varphi'_*$，并利用
$A' = A \otimes_R^\mathbf{L} R'$，得到《更多代数》公式
（\ref{more-algebra-equation-base-change-RHom}）中的导出 Hom 换基映射
$R\Hom_R(A, K) \otimes_R^\mathbf{L} R' \to
R\Hom_{R'}(A \otimes_R^\mathbf{L} R', K \otimes_R^\mathbf{L} R')$。
完全依照《更多代数》引理
\ref{more-algebra-lemma-base-change-RHom} 的证明展开左右两侧，
特别是使用《更多代数》引理
\ref{more-algebra-lemma-upgrade-adjoint-tensor-RHom}，即得所需结果。
\end{proof}

\begin{lemma}
\label{dualizing-lemma-flat-bc-surjection}
设 $R \to A$ 与 $R \to R'$ 为环同态，并令
$A' = A \otimes_R R'$。假设
\begin{enumerate}
\item 作为 $R$-模，$A$ 是伪凝聚的；
\item 作为 $R$-模，$R'$ 有有限 Tor 维数（例如 $R \to R'$ 平坦）；
\item $A$ 与 $R'$ 在 $R$ 上 Tor 独立。
\end{enumerate}
则对 $K \in D^+(R)$，映射（\ref{dualizing-equation-base-change}）是同构。
\end{lemma}

\begin{proof}
由引理 \ref{dualizing-lemma-check-base-change-is-iso} 及《更多代数》引理
\ref{more-algebra-lemma-internal-hom-evaluate-tensor-isomorphism}
的第 (4) 部分得出。
\end{proof}

\begin{lemma}
\label{dualizing-lemma-bc-surjection}
设 $R \to A$ 与 $R \to R'$ 为环同态，并令
$A' = A \otimes_R R'$。假设
\begin{enumerate}
\item 作为 $R$-模，$A$ 是完美的；
\item $A$ 与 $R'$ 在 $R$ 上 Tor 独立。
\end{enumerate}
则对所有 $K \in D(R)$，映射（\ref{dualizing-equation-base-change}）都是同构。
\end{lemma}

\begin{proof}
由引理 \ref{dualizing-lemma-check-base-change-is-iso} 及《更多代数》引理
\ref{more-algebra-lemma-internal-hom-evaluate-tensor-isomorphism}
的第 (1) 部分得出。
\end{proof}







\section{对偶复形}
\label{dualizing-section-dualizing}

\noindent
本节定义 Noetherian 环的对偶复形。

\begin{definition}
\label{dualizing-definition-dualizing}
设 $A$ 为 Noetherian 环。一个{\it 对偶复形}是 $A$-模复形
$\omega_A^\bullet$，满足
\begin{enumerate}
\item $\omega_A^\bullet$ 有有限内射维数；
\item 对所有 $i$，$H^i(\omega_A^\bullet)$ 是有限 $A$-模；
\item $A \to R\Hom_A(\omega_A^\bullet, \omega_A^\bullet)$ 是拟同构。
\end{enumerate}
\end{definition}

\noindent
这一概念需要一些时间才能熟悉。读者不妨先不看证明，自行证明下面若干引理。

\begin{lemma}
\label{dualizing-lemma-finite-ext-into-bounded-injective}
设 $A$ 为 Noetherian 环，$K, L \in D_{\textit{Coh}}(A)$，
并设 $L$ 有有限内射维数。则 $R\Hom_A(K, L)$ 属于
$D_{\textit{Coh}}(A)$。
\end{lemma}

\begin{proof}
取整数 $n$，考察可区别三角
$$
\tau_{\leq n}K \to K \to \tau_{\geq n + 1}K \to \tau_{\leq n}K[1]
$$
见《导出范畴》评注
\ref{derived-remark-truncation-distinguished-triangle}。
由于 $L$ 有有限内射维数，对某个常数 $c$，
$R\Hom_A(\tau_{\geq n + 1}K, L)$ 在次数 $\geq c - n$ 上的
上同调消失。因此，给定 $i$，对某个 $n \gg - i$，映射
$\Ext^i_A(K, L) \to \Ext^i_A(\tau_{\leq n}K, L)$ 是同构。
将《概形的导出范畴》引理
\ref{perfect-lemma-coherent-internal-hom} 应用于
$\tau_{\leq n}K$ 与 $L$，可知对所有 $i$，
$\Ext^i_A(K, L)$ 都是有限 $A$-模。因此 $R\Hom_A(K, L)$
确为 $D_{\textit{Coh}}(A)$ 的对象。
\end{proof}

\begin{lemma}
\label{dualizing-lemma-dualizing}
设 $A$ 为 Noetherian 环。若 $\omega_A^\bullet$ 为对偶复形，则函子
$$
D : K \longmapsto R\Hom_A(K, \omega_A^\bullet)
$$
是反等价 $D_{\textit{Coh}}(A) \to D_{\textit{Coh}}(A)$；
它交换 $D^+_{\textit{Coh}}(A)$ 与 $D^-_{\textit{Coh}}(A)$，
并诱导反等价
$D^b_{\textit{Coh}}(A) \to D^b_{\textit{Coh}}(A)$。
此外，$D \circ D$ 同构于恒等函子。
\end{lemma}

\begin{proof}
设 $K$ 为 $D_{\textit{Coh}}(A)$ 的对象。由引理
\ref{dualizing-lemma-finite-ext-into-bounded-injective}，
$R\Hom_A(K, \omega_A^\bullet)$ 是 $D_{\textit{Coh}}(A)$ 的对象。
由《更多代数》引理
\ref{more-algebra-lemma-internal-hom-evaluate-isomorphism-technical}
及对偶复形的假设，得到典范同构
$$
K = R\Hom_A(\omega_A^\bullet, \omega_A^\bullet) \otimes_A^\mathbf{L} K
\longrightarrow
R\Hom_A(R\Hom_A(K, \omega_A^\bullet), \omega_A^\bullet)
$$
因此该函子有拟逆，证明完成。
\end{proof}

\noindent
设 $R$ 为环。回忆：若 $D(R)$ 的对象 $L$ 关于导出张量积赋予
$D(R)$ 的对称幺半结构是可逆对象，则称其 {\it 可逆}。
由《更多代数》引理 \ref{more-algebra-lemma-invertible-derived}，
这等价于：$L$ 完美，$L = \bigoplus H^n(L)[-n]$ 是有限直和，
每个 $H^n(L)$ 都有限射影，并且存在开覆盖
$\Spec(R) = \bigcup D(f_i)$，使得对某些整数 $n_i$ 有
$L \otimes_R R_{f_i} \cong R_{f_i}[-n_i]$。

\begin{lemma}
\label{dualizing-lemma-equivalence-comes-from-invertible}
设 $A$ 为 Noetherian 环，且
$F : D^b_{\textit{Coh}}(A) \to D^b_{\textit{Coh}}(A)$
为 $A$-线性范畴等价。则 $F(A)$ 是 $D(A)$ 的可逆对象。
\end{lemma}

\begin{proof}
设 $\mathfrak m \subset A$ 为极大理想，余域为 $\kappa$。
考察对象 $F(\kappa)$。由于
$\kappa = \Hom_{D(A)}(\kappa, \kappa)$，可知 $F(\kappa)$
的所有上同调群都被 $\mathfrak m$ 零化。又有
$$
\Ext^i_A(\kappa, \kappa) = \text{Ext}^i_A(F(\kappa), F(\kappa))
= \Hom_{D(A)}(F(\kappa), F(\kappa)[i])
$$
当 $i < 0$ 时为零。设 $H^a(F(\kappa)) \not = 0$、
$H^b(F(\kappa)) \not = 0$，其中 $a$ 最小而 $b$ 最大
（特别地 $a \leq b$）。则 $D(A)$ 中存在非零映射
$$
F(\kappa) \to H^b(F(\kappa))[-b] \to H^a(F(\kappa))[-b]
\to F(\kappa)[a - b]
$$
（它非零是因为在上同调上诱导非零映射）。故 $b = a$，于是
$F(\kappa) = \kappa[-a]$。

\medskip\noindent
设 $G$ 是 $F$ 的拟逆。按上述论证，存在整数 $b$，使
$G(\kappa) = \kappa[-b]$。复合后得到 $a + b = 0$。
设 $E$ 为被 $\mathfrak m$ 的某次幂零化的有限 $A$-模。
对 $E$ 的长度归纳，可知对某个被 $\mathfrak m$ 的某次幂零化的
有限 $A$-模 $E'$，有 $G(E) = E'[-b]$。于是
$E[-a] = F(E')$。接着考察群
$$
\Ext^i_A(A, E') = \text{Ext}^i_A(F(A), F(E')) =
\Hom_{D(A)}(F(A), E[-a + i])
$$
左侧非零当且仅当 $i = 0$，此时得到 $E'$。令
$E = E' = \kappa$ 并使用 Nakayama 引理，可知当 $j > a$ 时
$H^j(F(A))_\mathfrak m$ 为零，而当 $j = a$ 时它由 $1$ 个元素生成。
另一方面，若对某个 $j < a$，$H^j(F(A))_\mathfrak m$ 非零，
则对某个 $i < 0$ 及某个 $E$，存在映射
$F(A) \to E[-a + i]$（《更多代数》引理
\ref{more-algebra-lemma-detect-cohomology}），这给出矛盾。
因此对某个由 $1$ 个元素生成的 $A_\mathfrak m$-模 $M$，有
$F(A)_\mathfrak m = M[-a]$。然而
$$
A_\mathfrak m = \Hom_{D(A)}(A, A)_\mathfrak m =
\Hom_{D(A)}(F(A), F(A))_\mathfrak m = \Hom_{A_\mathfrak m}(M, M)
$$
所以 $M \cong A_\mathfrak m$。由此存在
$f \in A$、$f \not \in \mathfrak m$，使得
$F(A)_f$ 同构于 $A_f[-a]$。证明完成。
\end{proof}

\begin{lemma}
\label{dualizing-lemma-dualizing-unique}
设 $A$ 为 Noetherian 环。若 $\omega_A^\bullet$ 与
$(\omega'_A)^\bullet$ 都是对偶复形，则对 $D(A)$ 的某个可逆对象
$L$，$(\omega'_A)^\bullet$ 拟同构于
$\omega_A^\bullet \otimes_A^\mathbf{L} L$。
\end{lemma}

\begin{proof}
由引理 \ref{dualizing-lemma-dualizing} 与
\ref{dualizing-lemma-equivalence-comes-from-invertible}，函子
$K \mapsto R\Hom_A(R\Hom_A(K, \omega_A^\bullet), (\omega_A')^\bullet)$
将 $A$ 映到某个可逆对象 $L$。换言之，存在同构
$$
L \longrightarrow R\Hom_A(\omega_A^\bullet, (\omega_A')^\bullet)
$$
由于 $L$ 有有限 Tor 维数，可应用《更多代数》引理
\ref{more-algebra-lemma-internal-hom-evaluate-isomorphism-technical}
看出，对 $D^b_{\textit{Coh}}(A)$ 中的 $K$，映射
$$
R\Hom_A(\omega_A^\bullet, (\omega'_A)^\bullet) \otimes_A^\mathbf{L} K
\longrightarrow
R\Hom_A(R\Hom_A(K, \omega_A^\bullet), (\omega_A')^\bullet)
$$
是同构。特别地，令 $K = \omega_A^\bullet$ 即完成证明。
\end{proof}

\begin{lemma}
\label{dualizing-lemma-dualizing-localize}
设 $A$ 为 Noetherian 环，$B = S^{-1}A$ 为局部化。
若 $\omega_A^\bullet$ 是对偶复形，则
$\omega_A^\bullet \otimes_A B$ 是 $B$ 的对偶复形。
\end{lemma}

\begin{proof}
设 $\omega_A^\bullet \to I^\bullet$ 为拟同构，其中 $I^\bullet$
是内射模的有界复形。则 $S^{-1}I^\bullet$ 是内射
$B = S^{-1}A$-模的有界复形（引理
\ref{dualizing-lemma-localization-injective-modules}），并表示
$\omega_A^\bullet \otimes_A B$。因此
$\omega_A^\bullet \otimes_A B$ 有有限内射维数。
由 $A \to B$ 的平坦性，
$H^i(\omega_A^\bullet \otimes_A B) = H^i(\omega_A^\bullet) \otimes_A B$，
故 $\omega_A^\bullet \otimes_A B$ 的上同调模有限。最后，映射
$$
B \longrightarrow
R\Hom_A(\omega_A^\bullet \otimes_A B, \omega_A^\bullet \otimes_A B)
$$
是拟同构，因为在此情形中内 Hom 的形成与平坦换基交换，见
《更多代数》引理 \ref{more-algebra-lemma-base-change-RHom}。
\end{proof}

\begin{lemma}
\label{dualizing-lemma-dualizing-glue}
设 $A$ 为 Noetherian 环，$f_1, \ldots, f_n \in A$ 生成单位理想。
若 $\omega_A^\bullet$ 是 $A$-模复形，且对所有 $i$，
$(\omega_A^\bullet)_{f_i}$ 都是 $A_{f_i}$ 的对偶复形，则
$\omega_A^\bullet$ 是 $A$ 的对偶复形。
\end{lemma}

\begin{proof}
考察双复形
$$
\prod\nolimits_{i_0} (\omega_A^\bullet)_{f_{i_0}}
\to
\prod\nolimits_{i_0 < i_1} (\omega_A^\bullet)_{f_{i_0}f_{i_1}}
\to \ldots
$$
相伴全复形拟同构于 $\omega_A^\bullet$；例如可用《下降》评注
\ref{descent-remark-standard-covering}，或《概形的导出范畴》引理
\ref{perfect-lemma-alternating-cech-complex-complex-computes-cohomology}。
由假设，将复形 $(\omega_A^\bullet)_{f_i}$ 视为
$A_{f_i}$-模复形时，它们有有限内射维数。因此，对每个
$(\omega_A^\bullet)_{f_{i_0} \ldots f_{i_p}}$、$p > 0$，
它在 $A_{f_{i_0} \ldots f_{i_p}}$ 上有有限内射维数，见引理
\ref{dualizing-lemma-localization-injective-modules}。进而，每个复形
$(\omega_A^\bullet)_{f_{i_0} \ldots f_{i_p}}$、$p > 0$
在 $A$ 上也有有限内射维数，见引理 \ref{dualizing-lemma-injective-flat}。
于是 $\omega_A^\bullet$ 作为 $A$-模复形有有限内射维数
（因为它可由带有限滤过的复形表示，而各分次部分都有有限内射维数）。
对每个 $i$，$H^n(\omega_A^\bullet)_{f_i}$ 是有限
$A_{f_i}$ 模，故由《代数》引理 \ref{algebra-lemma-cover}，
$H^i(\omega_A^\bullet)$ 是有限 $A$-模。最后，映射
$A \to R\Hom_A(\omega_A^\bullet, \omega_A^\bullet)$ 到 $A_{f_i}$
的导出换基是映射
$A_{f_i} \to R\Hom_A((\omega_A^\bullet)_{f_i}, (\omega_A^\bullet)_{f_i})$，
见《更多代数》引理 \ref{more-algebra-lemma-base-change-RHom}。
所以 $A \to R\Hom_A(\omega_A^\bullet, \omega_A^\bullet)$
是同构，证明完成。
\end{proof}

\begin{lemma}
\label{dualizing-lemma-dualizing-finite}
设 $A \to B$ 为 Noetherian 环之间的有限环同态，且
$\omega_A^\bullet$ 为对偶复形。则
$R\Hom(B, \omega_A^\bullet)$ 是 $B$ 的对偶复形。
\end{lemma}

\begin{proof}
设 $\omega_A^\bullet \to I^\bullet$ 为拟同构，其中 $I^\bullet$
是内射模的有界复形。则 $\Hom_A(B, I^\bullet)$ 是内射
$B$-模的有界复形（引理 \ref{dualizing-lemma-hom-injective}），并表示
$R\Hom(B, \omega_A^\bullet)$。所以
$R\Hom(B, \omega_A^\bullet)$ 有有限内射维数。
由引理 \ref{dualizing-lemma-exact-support-coherent}，它是
$D_{\textit{Coh}}(B)$ 的对象。最后计算
$$
\Hom_{D(B)}(R\Hom(B, \omega_A^\bullet), R\Hom(B, \omega_A^\bullet)) =
\Hom_{D(A)}(R\Hom(B, \omega_A^\bullet), \omega_A^\bullet) = B
$$
并且当 $n \not = 0$ 时，
$$
\Hom_{D(B)}(R\Hom(B, \omega_A^\bullet), R\Hom(B, \omega_A^\bullet)[n]) =
\Hom_{D(A)}(R\Hom(B, \omega_A^\bullet), \omega_A^\bullet[n]) = 0
$$
这证明了对偶复形的最后一条性质。在所示等式中，第一个等式由引理
\ref{dualizing-lemma-right-adjoint} 成立，第二个等式由引理
\ref{dualizing-lemma-dualizing} 成立。
\end{proof}

\begin{lemma}
\label{dualizing-lemma-dualizing-quotient}
设 $A \to B$ 为 Noetherian 环之间的满同态，且
$\omega_A^\bullet$ 为对偶复形。则
$R\Hom(B, \omega_A^\bullet)$ 是 $B$ 的对偶复形。
\end{lemma}

\begin{proof}
这是引理 \ref{dualizing-lemma-dualizing-finite} 的特殊情形。
\end{proof}

\begin{lemma}
\label{dualizing-lemma-dualizing-polynomial-ring}
设 $A$ 为 Noetherian 环。若 $\omega_A^\bullet$ 为对偶复形，
则 $\omega_A^\bullet \otimes_A A[x]$ 是 $A[x]$ 的对偶复形。
\end{lemma}

\begin{proof}
令 $B = A[x]$ 且 $\omega_B^\bullet = \omega_A^\bullet \otimes_A B$。
由引理 \ref{dualizing-lemma-injective-dimension-over-polynomial-ring} 与
《更多代数》引理 \ref{more-algebra-lemma-finite-injective-dimension}，
$\omega_B^\bullet$ 有有限内射维数。由 $A \to B$ 的平坦性，
$H^i(\omega_B^\bullet) = H^i(\omega_A^\bullet) \otimes_A B$，
故 $\omega_A^\bullet \otimes_A B$ 的上同调模有限。最后，映射
$$
B \longrightarrow R\Hom_B(\omega_B^\bullet, \omega_B^\bullet)
$$
是拟同构，因为在此情形中内 Hom 的形成与平坦换基交换，见
《更多代数》引理 \ref{more-algebra-lemma-base-change-RHom}。
\end{proof}

\begin{proposition}
\label{dualizing-proposition-dualizing-essentially-finite-type}
设 $A$ 为有对偶复形的 Noetherian 环。则任意在 $A$ 上本质有限型的
$A$-代数都有对偶复形。
\end{proposition}

\begin{proof}
这由引理 \ref{dualizing-lemma-dualizing-localize}、
\ref{dualizing-lemma-dualizing-quotient} 与
\ref{dualizing-lemma-dualizing-polynomial-ring} 合并得出。
\end{proof}

\begin{lemma}
\label{dualizing-lemma-find-function}
设 $A$ 为 Noetherian 环，$\omega_A^\bullet$ 为对偶复形。
设 $\mathfrak m \subset A$ 为极大理想，并令
$\kappa = A/\mathfrak m$。则对某个 $n \in \mathbf{Z}$，
$R\Hom_A(\kappa, \omega_A^\bullet) \cong \kappa[n]$。
\end{lemma}

\begin{proof}
这是因为 $R\Hom_A(\kappa, \omega_A^\bullet)$ 是 $\kappa$
上的对偶复形（引理 \ref{dualizing-lemma-dualizing-quotient}），而
$\kappa$ 上的对偶复形在平移意义下唯一（引理
\ref{dualizing-lemma-dualizing-unique}），并且 $\kappa$ 本身是
$\kappa$ 上的对偶复形。
\end{proof}




\section{局部环上的对偶复形}
\label{dualizing-section-dualizing-local}

\noindent
本节中，$(A, \mathfrak m, \kappa)$ 表示配备对偶复形
$\omega_A^\bullet$ 的 Noetherian 局部环，并且引理
\ref{dualizing-lemma-find-function} 中的整数 $n$ 为零。更准确地说，假设
$R\Hom_A(\kappa, \omega_A^\bullet) = \kappa[0]$。
此时称该对偶复形{\it 规范化}。注意，规范化对偶复形在同构意义下唯一，
并且 $A$ 的任意其他对偶复形都同构于某个规范化对偶复形的平移
（引理 \ref{dualizing-lemma-dualizing-unique}）。

\begin{lemma}
\label{dualizing-lemma-normalized-finite}
设 $(A, \mathfrak m, \kappa) \to (B, \mathfrak m', \kappa')$
为 Noetherian 局部环之间的有限局部同态，$\omega_A^\bullet$
为规范化对偶复形。则
$\omega_B^\bullet = R\Hom(B, \omega_A^\bullet)$ 是 $B$
的规范化对偶复形。
\end{lemma}

\begin{proof}
由引理 \ref{dualizing-lemma-dualizing-finite}，复形
$\omega_B^\bullet$ 是 $B$ 的对偶复形。由引理
\ref{dualizing-lemma-right-adjoint} 有
$$
R\Hom_B(\kappa', \omega_B^\bullet) =
R\Hom_B(\kappa', R\Hom(B, \omega_A^\bullet)) =
R\Hom_A(\kappa', \omega_A^\bullet)
$$
作为 $A$-模，$\kappa'$ 同构于若干份 $\kappa$ 的有限直和；
又因为 $\omega_A^\bullet$ 规范化，所以此复形的上同调只位于次数
$0$。因此 $\omega_B^\bullet$ 也是规范化对偶复形。
\end{proof}

\begin{lemma}
\label{dualizing-lemma-normalized-quotient}
设 $(A, \mathfrak m, \kappa)$ 为带规范化对偶复形
$\omega_A^\bullet$ 的 Noetherian 局部环。设 $A \to B$ 为满射。
则 $\omega_B^\bullet = R\Hom_A(B, \omega_A^\bullet)$
是 $B$ 的规范化对偶复形。
\end{lemma}

\begin{proof}
这是引理 \ref{dualizing-lemma-normalized-finite} 的特殊情形。
\end{proof}

\begin{lemma}
\label{dualizing-lemma-equivalence-finite-length}
设 $(A, \mathfrak m, \kappa)$ 为 Noetherian 局部环。
设 $F$ 是有限长度 $A$-模范畴的 $A$-线性自等价。
则 $F$ 同构于恒等函子。
\end{lemma}

\begin{proof}
由于 $\kappa$ 是该范畴唯一的单对象，有
$F(\kappa) \cong \kappa$。该范畴是 Abel 范畴，故 $F$ 正合。
因此对所有有限长度模 $E$，$F(E)$ 与 $E$ 的长度相同。
由
$\Hom(E, \kappa) = \Hom(F(E), F(\kappa)) \cong \Hom(F(E), \kappa)$
及 Nakayama 引理，$E$ 与 $F(E)$ 有相同数目的生成元。
所以 $F(A/\mathfrak m^n)$ 是循环 $A$-模。选取生成元
$e \in F(A/\mathfrak m^n)$。由 $F$ 的 $A$-线性，
$\mathfrak m^n e = 0$。由于长度相等，映射
$A/\mathfrak m^n \to F(A/\mathfrak m^n)$ 必为同构。选取元素
$$
e \in \lim F(A/\mathfrak m^n)
$$
使它对所有 $n$ 都映到生成元（略去一个简短论证）。
于是得到同构系统 $A/\mathfrak m^n \to F(A/\mathfrak m^n)$，
它与所有 $A$-模映射
$A/\mathfrak m^n \to A/\mathfrak m^{n'}$ 相容
（再次使用 $F$ 的 $A$-线性）。任意有限长度模都是循环模直和之间
某个映射的余核，因此得到该引理的同构。
\end{proof}

\begin{lemma}
\label{dualizing-lemma-dualizing-finite-length}
设 $(A, \mathfrak m, \kappa)$ 为带规范化对偶复形
$\omega_A^\bullet$ 的 Noetherian 局部环。设 $E$ 为 $\kappa$
的内射包。则当 $N$ 遍历有限长度 $A$-模时，存在函子性同构
$$
R\Hom_A(N, \omega_A^\bullet) = \Hom_A(N, E)[0]
$$
\end{lemma}

\begin{proof}
对 $N$ 的长度归纳可知，$R\Hom_A(N, \omega_A^\bullet)$
是位于次数 $0$ 的有限长度模。因此
$R\Hom_A(-, \omega_A^\bullet)$ 在有限长度模范畴上诱导反等价。
由命题 \ref{dualizing-proposition-matlis}，$\Hom_A(-, E)$ 也如此，故
$$
N \longmapsto \Hom_A(R\Hom_A(N, \omega_A^\bullet), E)
$$
是引理 \ref{dualizing-lemma-equivalence-finite-length} 中那样的等价，
因而同构于恒等函子。又由命题 \ref{dualizing-proposition-matlis}，
$\Hom_A(-, E)$ 施用两次就是恒等函子，故得到该引理的陈述。
\end{proof}

\begin{lemma}
\label{dualizing-lemma-sitting-in-degrees}
设 $(A, \mathfrak m, \kappa)$ 为带规范化对偶复形
$\omega_A^\bullet$ 的 Noetherian 局部环。设 $M$ 为有限
$A$-模，并令 $d = \dim(\text{Supp}(M))$。则
\begin{enumerate}
\item 若 $\Ext^i_A(M, \omega_A^\bullet)$ 非零，则
$i \in \{-d, \ldots, 0\}$；
\item $\Ext^i_A(M, \omega_A^\bullet)$ 的支集维数至多为 $-i$；
\item $\text{depth}(M)$ 是满足
$\Ext^{-\delta}_A(M, \omega_A^\bullet) \not = 0$ 的最小整数
$\delta \geq 0$。
\end{enumerate}
\end{lemma}

\begin{proof}
对 $d$ 作归纳。若 $d = 0$，则结论由引理
\ref{dualizing-lemma-dualizing-finite-length} 与 Matlis 对偶
（命题 \ref{dualizing-proposition-matlis}）得出；后者保证当 $M$ 非零时，
$\Hom_A(M, E)$ 非零。

\medskip\noindent
假设结论对支集维数 $< d$ 的模成立，并且 $M$ 的深度 $> 0$。
选取一个在 $M$ 上为非零因子的 $f \in \mathfrak m$，并考虑短正合列
$$
0 \to M \to M \to M/fM \to 0
$$
由于 $\dim(\text{Supp}(M/fM)) = d - 1$
（代数，引理 \ref{algebra-lemma-one-equation-module}），
可以应用归纳假设。记
$E^i = \Ext^i_A(M, \omega_A^\bullet)$ 以及
$F^i = \Ext^i_A(M/fM, \omega_A^\bullet)$，得到长正合列
$$
\ldots \to F^i \to E^i \xrightarrow{f} E^i \to F^{i + 1} \to \ldots
$$
由归纳假设，当
$i + 1 \not \in \{-\dim(\text{Supp}(M/fM)), \ldots, -\text{depth}(M/fM)\}$
时，$E^i/fE^i = 0$。由 Nakayama 引理（代数，引理
\ref{algebra-lemma-NAK}）和代数，引理
\ref{algebra-lemma-depth-drops-by-one}，可知当
$i \not \in \{-\dim(\text{Supp}(M)), \ldots, -\text{depth}(M)\}$
时，$E^i = 0$。此外，在边界情形
$i = - \text{depth}(M)$，由归纳假设 $F^{i + 1}$ 非零，
从而可知 $E^i$ 非零。由于 $E^i/fE^i \subset F^{i + 1}$，有
$$
\dim(\text{Supp}(F^{i + 1})) \geq \dim(\text{Supp}(E^i/fE^i))
\geq \dim(\text{Supp}(E^i)) - 1
$$
（参见上面所用的引理），由此还得到维数估计 (2)。

\medskip\noindent
若 $M$ 的深度为 $0$ 且 $d > 0$，令
$N = M[\mathfrak m^\infty]$，并置
$M' = M/N$（比较引理 \ref{dualizing-lemma-divide-by-torsion}）。
则 $M'$ 的深度 $> 0$，且 $\dim(\text{Supp}(M')) = d$。
因此结论对 $M'$ 成立；又因为
$R\Hom_A(N, \omega_A^\bullet) = \Hom_A(N, E)$
（引理 \ref{dualizing-lemma-dualizing-finite-length}），
$\Ext$ 的长正合上同调列蕴含结论对 $M$ 也成立。
\end{proof}

\begin{remark}
\label{dualizing-remark-vanishing-for-arbitrary-modules}
设 $(A, \mathfrak m)$ 和 $\omega_A^\bullet$ 如引理
\ref{dualizing-lemma-sitting-in-degrees} 中所述。由更多代数，引理
\ref{more-algebra-lemma-injective-amplitude}，该引理的第 (3) 部分适用，
故 $\omega_A^\bullet$ 的内射振幅位于 $[-d, 0]$。
特别地，对任意 $A$-模 $M$（不必有限），当
$i \not \in \{-d, \ldots, 0\}$ 时，有
$\Ext^i_A(M, \omega_A^\bullet) = 0$。
\end{remark}

\begin{lemma}
\label{dualizing-lemma-local-CM}
设 $(A, \mathfrak m, \kappa)$ 为带规范化对偶复形
$\omega_A^\bullet$ 的 Noetherian 局部环。设 $M$ 为有限
$A$-模。下列条件等价：
\begin{enumerate}
\item $M$ 是 Cohen--Macaulay 模；
\item $\Ext^i_A(M, \omega_A^\bullet)$ 至多对一个 $i$ 非零；
\item 当 $i \not = \dim(\text{Supp}(M))$ 时，
$\Ext^{-i}_A(M, \omega_A^\bullet)$ 为零。
\end{enumerate}
以 $CM_d$ 表示深度为 $d$ 的有限 Cohen--Macaulay $A$-模范畴。
则 $M \mapsto \Ext^{-d}_A(M, \omega_A^\bullet)$
定义了 $CM_d$ 的一个反自等价。
\end{lemma}

\begin{proof}
以下不再另行说明地使用引理 \ref{dualizing-lemma-sitting-in-degrees} 的结论。
固定一个有限模 $M$。若 $M$ 是 Cohen--Macaulay 模，则只有
$\Ext^{-d}_A(M, \omega_A^\bullet)$ 可能非零，故
(1) $\Rightarrow$ (3)。(3) $\Rightarrow$ (2) 是立即的。
假设 (2)，并令 $N = \Ext^{-\delta}_A(M, \omega_A^\bullet)$
为那个非零的 $\Ext$，其中 $\delta = \text{depth}(M)$。由于
$$
M[0] = R\Hom_A(R\Hom_A(M, \omega_A^\bullet), \omega_A^\bullet) =
R\Hom_A(N[\delta], \omega_A^\bullet)
$$
（引理 \ref{dualizing-lemma-dualizing}），可知
$M = \Ext_A^{-\delta}(N, \omega_A^\bullet)$。
因此 $\delta \geq \dim(\text{Supp}(M))$。另一方面，
还知道 $\delta \leq \dim(\text{Supp}(M))$
（代数，引理 \ref{algebra-lemma-bound-depth}），故 $M$ 是
Cohen--Macaulay 模。

\medskip\noindent
为证明最后一个陈述，只需对 $CM_d$ 中的 $M$ 证明
$N = \Ext^{-d}_A(M, \omega_A^\bullet)$ 属于 $CM_d$。
上面已经看到 $M[0] = R\Hom_A(N[d], \omega_A^\bullet)$，
由 (1) 与 (3) 的等价性即得所需结论。
\end{proof}

\begin{lemma}
\label{dualizing-lemma-dualizing-artinian}
设 $(A, \mathfrak m, \kappa)$ 为带规范化对偶复形
$\omega_A^\bullet$ 的 Noetherian 局部环。若 $\dim(A) = 0$，
则 $\omega_A^\bullet \cong E[0]$，其中 $E$ 是剩余域的内射包。
\end{lemma}

\begin{proof}
由引理 \ref{dualizing-lemma-dualizing-finite-length} 立即得到。
\end{proof}

\begin{lemma}
\label{dualizing-lemma-divide-by-finite-length-ideal}
设 $(A, \mathfrak m, \kappa)$ 为带规范化对偶复形的
Noetherian 局部环。设 $I \subset \mathfrak m$ 为有限长度理想，
并置 $B = A/I$。则在 $D(A)$ 中有特异三角
$$
\omega_B^\bullet \to \omega_A^\bullet \to \Hom_A(I, E)[0] \to
\omega_B^\bullet[1]
$$
其中 $E$ 是 $\kappa$ 的内射包，$\omega_B^\bullet$ 是 $B$
的规范化对偶复形。
\end{lemma}

\begin{proof}
使用短正合列 $0 \to I \to A \to B \to 0$ 以及引理
\ref{dualizing-lemma-dualizing-finite-length} 和
\ref{dualizing-lemma-normalized-quotient}。
\end{proof}

\begin{lemma}
\label{dualizing-lemma-divide-by-nonzerodivisor}
设 $(A, \mathfrak m, \kappa)$ 为带规范化对偶复形
$\omega_A^\bullet$ 的 Noetherian 局部环。设
$f \in \mathfrak m$ 为非零因子，并置 $B = A/(f)$。
则在 $D(A)$ 中有特异三角
$$
\omega_B^\bullet \to \omega_A^\bullet \to \omega_A^\bullet \to
\omega_B^\bullet[1]
$$
其中 $\omega_B^\bullet$ 是 $B$ 的规范化对偶复形。
\end{lemma}

\begin{proof}
使用短正合列 $0 \to A \to A \to B \to 0$ 和引理
\ref{dualizing-lemma-normalized-quotient}。
\end{proof}

\begin{lemma}
\label{dualizing-lemma-nonvanishing-generically-local}
设 $(A, \mathfrak m, \kappa)$ 为带规范化对偶复形
$\omega_A^\bullet$ 的 Noetherian 局部环。设 $\mathfrak p$ 是 $A$
的极小素理想，且 $\dim(A/\mathfrak p) = e$。则
$H^i(\omega_A^\bullet)_\mathfrak p$ 非零当且仅当 $i = -e$。
\end{lemma}

\begin{proof}
由于 $A_\mathfrak p$ 的维数为零，存在整数 $n > 0$，使得
$\mathfrak p^nA_\mathfrak p$ 为零。置 $B = A/\mathfrak p^n$ 以及
$\omega_B^\bullet = R\Hom_A(B, \omega_A^\bullet)$。
由 $B_\mathfrak p = A_\mathfrak p$ 可知
$$
(\omega_B^\bullet)_\mathfrak p =
R\Hom_A(B, \omega_A^\bullet) \otimes_A^\mathbf{L} A_\mathfrak p =
R\Hom_{A_\mathfrak p}(B_\mathfrak p, (\omega_A^\bullet)_\mathfrak p) =
(\omega_A^\bullet)_\mathfrak p
$$
第二个等号由更多代数，引理
\ref{more-algebra-lemma-base-change-RHom} 得到。由引理
\ref{dualizing-lemma-normalized-quotient}，可以用 $B$ 替换 $A$。
替换后有 $\dim(A) = e$。于是由引理
\ref{dualizing-lemma-sitting-in-degrees} 的第 (1)、(2) 部分，
$H^i(\omega_A^\bullet)_\mathfrak p$ 仅当 $i = -e$ 时才可能非零。
另一方面，$(\omega_A^\bullet)_\mathfrak p$ 是非零环
$A_\mathfrak p$ 的对偶复形（引理 \ref{dualizing-lemma-dualizing-localize}），
故剩下的这个模必为非零。
\end{proof}





\section{对偶复形与维数函数}
\label{dualizing-section-dimension-function}

\noindent
局部情形中的结果有如下推论：具有对偶复形的 Noetherian 环
是有限维的普遍等长链环。

\begin{lemma}
\label{dualizing-lemma-nonvanishing-generically}
设 $A$ 为 Noetherian 环，$\mathfrak p$ 为 $A$ 的极小素理想。
则 $H^i(\omega_A^\bullet)_\mathfrak p$ 恰好对一个 $i$ 非零。
\end{lemma}

\begin{proof}
复形 $\omega_A^\bullet \otimes_A A_\mathfrak p$ 是
$A_\mathfrak p$ 的对偶复形（引理 \ref{dualizing-lemma-dualizing-localize}）。
由于 $\mathfrak p$ 极小，$A_\mathfrak p$ 的维数为零。
故结论由引理 \ref{dualizing-lemma-dualizing-artinian} 得出。
\end{proof}

\noindent
设 $A$ 为 Noetherian 环，$\omega_A^\bullet$ 为对偶复形。
引理 \ref{dualizing-lemma-find-function} 使我们能够定义函数
$$
\delta = \delta_{\omega_A^\bullet} : \Spec(A) \longrightarrow \mathbf{Z}
$$
它把 $\mathfrak p$ 映到引理 \ref{dualizing-lemma-find-function} 针对
$A_\mathfrak p$ 上的对偶复形
$(\omega_A^\bullet)_\mathfrak p$（引理
\ref{dualizing-lemma-dualizing-localize}）及剩余域 $\kappa(\mathfrak p)$
所给出的整数。更准确地说，把 $\delta(\mathfrak p)$ 定义为使
$$
(\omega_A^\bullet)_\mathfrak p[-\delta(\mathfrak p)]
$$
成为 Noetherian 局部环 $A_\mathfrak p$ 上规范化对偶复形的唯一整数。

\begin{lemma}
\label{dualizing-lemma-quotient-function}
设 $A$ 为 Noetherian 环，$\omega_A^\bullet$ 为对偶复形。
设 $A \to B$ 为满环同态，并令
$\omega_B^\bullet = R\Hom(B, \omega_A^\bullet)$ 为引理
\ref{dualizing-lemma-dualizing-quotient} 所给出的 $B$ 的对偶复形。则
$$
\delta_{\omega_B^\bullet} = \delta_{\omega_A^\bullet}|_{\Spec(B)}
$$
\end{lemma}

\begin{proof}
这由这些函数的定义和引理 \ref{dualizing-lemma-normalized-quotient} 得出。
\end{proof}

\begin{lemma}
\label{dualizing-lemma-dimension-function}
设 $A$ 为 Noetherian 环，$\omega_A^\bullet$ 为对偶复形。
上面定义的函数 $\delta = \delta_{\omega_A^\bullet}$ 是维数函数
（拓扑，定义 \ref{topology-definition-dimension-function}）。
\end{lemma}

\begin{proof}
设 $\mathfrak p \subset \mathfrak q$ 为一个直接特化。
需要证明 $\delta(\mathfrak p) = \delta(\mathfrak q) + 1$。
可以把 $A$ 替换为 $A/\mathfrak p$，把复形 $\omega_A^\bullet$
替换为
$\omega_{A/\mathfrak p}^\bullet = R\Hom(A/\mathfrak p, \omega_A^\bullet)$，
把素理想 $\mathfrak p$ 替换为 $(0)$，并把素理想 $\mathfrak q$
替换为 $\mathfrak q/\mathfrak p$；参见引理
\ref{dualizing-lemma-quotient-function}。因此可以假设 $A$ 是整环，
$\mathfrak p = (0)$，而 $\mathfrak q$ 是高度为 $1$ 的素理想。

\medskip\noindent
由引理 \ref{dualizing-lemma-nonvanishing-generically}，
$H^i(\omega_A^\bullet)_{(0)}$ 恰好对一个 $i$ 非零，将其记为
$i_0$。事实上 $i_0 = -\delta((0))$，因为
$(\omega_A^\bullet)_{(0)}[-\delta((0))]$ 是域 $A_{(0)}$
上的规范化对偶复形。

\medskip\noindent
另一方面，$(\omega_A^\bullet)_\mathfrak q[-\delta(\mathfrak q)]$
是 $A_\mathfrak q$ 的规范化对偶复形。由引理
\ref{dualizing-lemma-nonvanishing-generically-local} 可知
$$
H^e((\omega_A^\bullet)_\mathfrak q[-\delta(\mathfrak q)])_{(0)} =
H^{e - \delta(\mathfrak q)}(\omega_A^\bullet)_{(0)}
$$
仅当 $e = -\dim(A_\mathfrak q) = -1$ 时非零。因此
$$
-\delta((0)) = -1 - \delta(\mathfrak q)
$$
这正是所需等式。
\end{proof}

\begin{lemma}
\label{dualizing-lemma-universally-catenary}
设 $A$ 为具有对偶复形的 Noetherian 环。则 $A$ 是有限维的
普遍等长链环。
\end{lemma}

\begin{proof}
由引理 \ref{dualizing-lemma-dimension-function}，$\Spec(A)$ 具有维数函数，
因而是等长链的；参见拓扑，引理
\ref{topology-lemma-dimension-function-catenary}。
所以 $A$ 是等长链环；参见代数，引理
\ref{algebra-lemma-catenary}。由命题
\ref{dualizing-proposition-dualizing-essentially-finite-type}，
$A$ 是普遍等长链环。

\medskip\noindent
任意对偶复形 $\omega_A^\bullet$ 均属于
$D^b_{\textit{Coh}}(A)$，故由引理
\ref{dualizing-lemma-nonvanishing-generically}，函数
$\delta_{\omega_A^\bullet}$ 在极小素理想处的取值有界。
另一方面，对剩余域为 $\kappa$ 的极大理想 $\mathfrak m$，整数
$i = -\delta(\mathfrak m)$ 是使
$\Ext_A^i(\kappa, \omega_A^\bullet)$ 非零的唯一整数
（引理 \ref{dualizing-lemma-find-function}）。由于 $\omega_A^\bullet$
具有有限内射维数，这些值也有界。$A$ 的维数等于
$\delta(\mathfrak p) - \delta(\mathfrak m)$ 的最大值，
其中 $\mathfrak p \subset \mathfrak m$ 是由一个极小素理想与一个
极大素理想组成的一对。因此 $\Spec(A)$ 的维数有界。
\end{proof}

\begin{lemma}
\label{dualizing-lemma-depth-dualizing-module}
设 $(A, \mathfrak m, \kappa)$ 为带规范化对偶复形
$\omega_A^\bullet$ 的 Noetherian 局部环。令 $d = \dim(A)$ 以及
$\omega_A = H^{-d}(\omega_A^\bullet)$。则
\begin{enumerate}
\item $\omega_A$ 的支集是 $\Spec(A)$ 中所有维数为 $d$
的不可约分支之并；
\item $\omega_A$ 满足 $(S_2)$；参见代数，定义
\ref{algebra-definition-conditions}。
\end{enumerate}
\end{lemma}

\begin{proof}
以下不再另行说明地使用引理 \ref{dualizing-lemma-sitting-in-degrees}。
由引理 \ref{dualizing-lemma-nonvanishing-generically-local}，$\omega_A$
的支集包含所有维数为 $d$ 的不可约分支。设
$\mathfrak p \subset A$ 为素理想。由引理
\ref{dualizing-lemma-dimension-function}，复形
$(\omega_A^\bullet)_{\mathfrak p}[-\dim(A/\mathfrak p)]$
是 $A_\mathfrak p$ 的规范化对偶复形。因此，若
$\dim(A/\mathfrak p) + \dim(A_\mathfrak p) < d$，则
$(\omega_A)_\mathfrak p = 0$。这证明 $\omega_A$ 的支集是所有维数为
$d$ 的不可约分支之并：由于 $A$ 是等长链环
（引理 \ref{dualizing-lemma-universally-catenary}），这个并的补集恰由满足
$\dim(A/\mathfrak p) + \dim(A_\mathfrak p) < d$ 的 $A$ 的素理想
$\mathfrak p$ 构成。另一方面，若
$\dim(A/\mathfrak p) + \dim(A_\mathfrak p) = d$，则
$$
(\omega_A)_\mathfrak p =
H^{-\dim(A_\mathfrak p)}\left(
(\omega_A^\bullet)_{\mathfrak p}[-\dim(A/\mathfrak p)] \right)
$$
因此，为证明 $\omega_A$ 满足 $(S_2)$，只需证明 $\omega_A$
的深度至少为 $\min(\dim(A), 2)$。对 $\dim(A)$ 作归纳。
$\dim(A) = 0$ 的情形显然。

\medskip\noindent
假设 $\text{depth}(A) > 0$。选取非零因子
$f \in \mathfrak m$，并置 $B = A/fA$。则
$\dim(B) = \dim(A) - 1$，可以对 $B$ 应用归纳假设。
由引理 \ref{dualizing-lemma-divide-by-nonzerodivisor}，乘以 $f$ 在
$\omega_A$ 上是单射，且有
$\omega_A/f\omega_A \subset \omega_B$。这证明 $\omega_A$ 的深度
至少为 $1$。若 $\dim(A) > 1$，则 $\dim(B) > 0$，而 $\omega_B$
的深度 $ > 0$。所以 $\omega_A$ 的深度 $> 1$，此情形得证。

\medskip\noindent
假设 $\dim(A) > 0$ 且 $\text{depth}(A) = 0$。令
$I = A[\mathfrak m^\infty]$，并置 $B = A/I$。则 $B$ 的深度
$\geq 1$，且由引理 \ref{dualizing-lemma-divide-by-finite-length-ideal}，
$\omega_A = \omega_B$。上面已经对 $\omega_B$ 证明了结论，
故证明完成。
\end{proof}





\section{局部对偶定理}
\label{dualizing-section-local-duality}

\noindent
本节的主要结果归功于 Grothendieck。

\begin{lemma}
\label{dualizing-lemma-local-cohomology-of-dualizing}
设 $(A, \mathfrak m, \kappa)$ 为 Noetherian 局部环，
$\omega_A^\bullet$ 为规范化对偶复形，并令
$Z = V(\mathfrak m) \subset \Spec(A)$。则
$E = R^0\Gamma_Z(\omega_A^\bullet)$ 是 $\kappa$ 的内射包，且
$R\Gamma_Z(\omega_A^\bullet) = E[0]$。
\end{lemma}

\begin{proof}
由引理 \ref{dualizing-lemma-local-cohomology-noetherian}，有
$R\Gamma_{\mathfrak m} = R\Gamma_Z$。因此，由引理
\ref{dualizing-lemma-local-cohomology-ext}，
$$
R\Gamma_Z(\omega_A^\bullet) =
R\Gamma_{\mathfrak m}(\omega_A^\bullet) =
\text{hocolim}\ R\Hom_A(A/\mathfrak m^n, \omega_A^\bullet)
$$
设 $E'$ 为剩余域的一个内射包。由引理
\ref{dualizing-lemma-dualizing-finite-length}，可以找到与转移映射相容的同构
$$
R\Hom_A(A/\mathfrak m^n, \omega_A^\bullet) \cong \Hom_A(A/\mathfrak m^n, E')[0]
$$
由引理 \ref{dualizing-lemma-union-artinian}，
$E' = \bigcup E'[\mathfrak m^n] = \colim \Hom_A(A/\mathfrak m^n, E')$，
故 $E \cong E'$，而复形 $R\Gamma_Z(\omega_A^\bullet)$
的所有其他上同调群均为零。
\end{proof}

\begin{remark}
\label{dualizing-remark-specific-injective-hull}
设 $(A, \mathfrak m, \kappa)$ 为带规范化对偶复形
$\omega_A^\bullet$ 的 Noetherian 局部环。由上面的引理
\ref{dualizing-lemma-local-cohomology-of-dualizing}，
$R\Gamma_Z(\omega_A^\bullet)$ 是置于次数 $0$ 的剩余域内射包。
事实上，这给出了内射包的一种“构造”或“实现”，它比任意选取一个
内射包略微更典范。具体而言，规范化对偶复形在同构意义下唯一，
其自同构群为 $A$ 的单位群；而 $\kappa$ 的内射包也在同构意义下
唯一，其自同构群为 $A$ 关于 $\mathfrak m$ 的完备化
$A^\wedge$ 的单位群。
\end{remark}

\noindent
以下是本节的主要结果。

\begin{theorem}
\label{dualizing-theorem-local-duality}
设 $(A, \mathfrak m, \kappa)$ 为 Noetherian 局部环，
$\omega_A^\bullet$ 为规范化对偶复形，$E$ 为剩余域的一个内射包，
并令 $Z = V(\mathfrak m) \subset \Spec(A)$。以 ${}^\wedge$
表示关于 $\mathfrak m$ 的导出完备化。则对 $D(A)$ 中的 $K$，有
$$
R\Hom_A(K, \omega_A^\bullet)^\wedge \cong R\Hom_A(R\Gamma_Z(K), E[0])
$$
\end{theorem}

\begin{proof}
由引理 \ref{dualizing-lemma-local-cohomology-of-dualizing}，注意到
$E[0] \cong R\Gamma_Z(\omega_A^\bullet)$。由更多代数，引理
\ref{more-algebra-lemma-completion-RHom}，左端的完备化可以移入内部。
因此只需证明
$$
R\Hom_A(K^\wedge, (\omega_A^\bullet)^\wedge)
=
R\Hom_A(R\Gamma_Z(K), R\Gamma_Z(\omega_A^\bullet))
$$
这由命题 \ref{dualizing-proposition-torsion-complete} 给出的
$D_{comp}(A, \mathfrak m)$ 与
$D_{\mathfrak m^\infty\text{-torsion}}(A)$ 之间的等价得出。
更准确地说，这是引理 \ref{dualizing-lemma-compare-RHom} 的特殊情形。
\end{proof}

\noindent
下面给出上述定理的一个特殊情形。

\begin{lemma}
\label{dualizing-lemma-special-case-local-duality}
设 $(A, \mathfrak m, \kappa)$ 为 Noetherian 局部环，
$\omega_A^\bullet$ 为规范化对偶复形，$E$ 为剩余域的一个内射包。
设 $K \in D_{\textit{Coh}}(A)$。则
$$
\Ext^{-i}_A(K, \omega_A^\bullet)^\wedge =
\Hom_A(H^i_{\mathfrak m}(K), E)
$$
其中 ${}^\wedge$ 表示 $\mathfrak m$-进完备化。
\end{lemma}

\begin{proof}
由引理 \ref{dualizing-lemma-dualizing}，$R\Hom_A(K, \omega_A^\bullet)$
是 $D_{\textit{Coh}}(A)$ 的对象。由更多代数，引理
\ref{more-algebra-lemma-derived-completion-pseudo-coherent}，
$R\Hom_A(K, \omega_A^\bullet)$ 的导出完备化之上同调模等于通常的完备化
$\Ext^i_A(K, \omega_A^\bullet)^\wedge$。另一方面，由引理
\ref{dualizing-lemma-local-cohomology-noetherian}，对
$Z = V(\mathfrak m)$ 有 $R\Gamma_{\mathfrak m} = R\Gamma_Z$。
此外，函子 $\Hom_A(-, E)$ 正合，故可通过上同调分解。
因此该引理是定理 \ref{dualizing-theorem-local-duality} 的推论。
\end{proof}





\section{对偶模}
\label{dualizing-section-dualizing-module}

\noindent
设 $(A, \mathfrak m, \kappa)$ 为 Noetherian 局部环，且
$\omega_A^\bullet$ 为规范化对偶复形。称引理
\ref{dualizing-lemma-depth-dualizing-module} 中所述的模
$\omega_A = H^{-\dim(A)}(\omega_A^\bullet)$ 为 $A$ 的一个
{\it 对偶模}。这个模是 $A$ 的一个典范模。通常约定把 Noetherian
局部环 $(A, \mathfrak m, \kappa)$ 的{\it 典范模}定义为满足
$$
\Hom_A(K, E) \cong H^{\dim(A)}_\mathfrak m(A)
$$
的有限 $A$-模 $K$，其中 $E$ 是剩余域的一个内射包。对偶模之所以
是典范模，是因为由引理 \ref{dualizing-lemma-special-case-local-duality}，
$$
\Hom_A(H^{\dim(A)}_\mathfrak m(A), E) = (\omega_A)^\wedge
$$
进而施用 $\Hom_A(-, E)$ 得
\begin{align*}
\Hom_A(\omega_A, E)
& =
\Hom_A((\omega_A)^\wedge, E) \\
& =
\Hom_A(\Hom_A(H^{\dim(A)}_\mathfrak m(A), E), E) \\
& = H^{\dim(A)}_\mathfrak m(A)
\end{align*}
第一个等号成立是因为 $E$ 是 $\mathfrak m$-幂挠的，第二个等号由
上式得到，第三个等号由 Matlis 对偶（命题
\ref{dualizing-proposition-matlis}）得到。上述典范模定义的用处在于：
即使 $A$ 没有对偶复形，该定义仍然有意义。




\section{Cohen--Macaulay 环}
\label{dualizing-section-CM}

\noindent
代数，第 \ref{algebra-section-CM} 节和第
\ref{algebra-section-CM-ring} 节研究了 Cohen--Macaulay 模与环。

\begin{lemma}
\label{dualizing-lemma-depth-in-terms-dualizing-complex}
设 $(A, \mathfrak m, \kappa)$ 为带规范化对偶复形
$\omega_A^\bullet$ 的 Noetherian 局部环。则 $\text{depth}(A)$
等于满足 $H^{-\delta}(\omega_A^\bullet) \not = 0$ 的最小整数
$\delta \geq 0$。
\end{lemma}

\begin{proof}
这由引理 \ref{dualizing-lemma-sitting-in-degrees} 立即得出。
下面再给出两种证明方式。

\medskip\noindent
第一种。由 Nakayama 引理，$\delta$ 是满足
$\Hom_A(H^{-\delta}(\omega_A^\bullet), \kappa) \not = 0$
的最小整数。换言之，它是使
$\Ext_A^{-\delta}(\omega_A^\bullet, \kappa)$ 非零的最小整数。
利用引理 \ref{dualizing-lemma-dualizing} 以及 $\omega_A^\bullet$ 的规范化，
这等于使 $\Ext_A^\delta(\kappa, A)$ 非零的最小整数。
由代数，引理 \ref{algebra-lemma-depth-ext}，这正是 $A$ 的深度。

\medskip\noindent
第二种。由局部对偶定理（采用引理
\ref{dualizing-lemma-special-case-local-duality} 的形式），$\delta$
是使 $H^\delta_\mathfrak m(A)$ 非零的最小整数。由引理
\ref{dualizing-lemma-depth}，这正是 $A$ 的深度。
\end{proof}

\begin{lemma}
\label{dualizing-lemma-apply-CM}
设 $(A, \mathfrak m, \kappa)$ 为带规范化对偶复形
$\omega_A^\bullet$ 的 Noetherian 局部环，其对偶模为
$\omega_A = H^{-\dim(A)}(\omega_A^\bullet)$。下列条件等价：
\begin{enumerate}
\item $A$ 是 Cohen--Macaulay 环；
\item $\omega_A^\bullet$ 集中在单一次数上；
\item $\omega_A^\bullet = \omega_A[\dim(A)]$。
\end{enumerate}
在这种情况下，$\omega_A$ 是极大 Cohen--Macaulay 模。
\end{lemma}

\begin{proof}
由引理 \ref{dualizing-lemma-local-CM} 立即得到。
\end{proof}

\begin{lemma}
\label{dualizing-lemma-has-dualizing-module-CM}
设 $A$ 为 Noetherian 环。若存在有限 $A$-模 $\omega_A$，使得
$\omega_A[0]$ 是对偶复形，则 $A$ 是 Cohen--Macaulay 环。
\end{lemma}

\begin{proof}
可以把 $A$ 替换为它在某个素理想处的局部化
（引理 \ref{dualizing-lemma-dualizing-localize} 及代数，定义
\ref{algebra-definition-ring-CM}）。此时结论由引理
\ref{dualizing-lemma-apply-CM} 立即得出。
\end{proof}

\begin{lemma}
\label{dualizing-lemma-CM-open}
设 $A$ 为带对偶复形 $\omega_A^\bullet$ 的 Noetherian 环，
$M$ 为有限 $A$-模。则
$$
U = \{\mathfrak p \in \Spec(A) \mid M_\mathfrak p\text{ 为 Cohen--Macaulay}\}
$$
是 $\Spec(A)$ 的开子集，并且它与 $\text{Supp}(M)$ 的交是稠密的。
\end{lemma}

\begin{proof}
若 $\mathfrak p$ 是 $\text{Supp}(M)$ 的一般点，则
$\text{depth}(M_\mathfrak p) = \dim(M_\mathfrak p) = 0$，
故 $\mathfrak p \in U$。这证明了稠密性。若 $\mathfrak p \in U$，
则由引理 \ref{dualizing-lemma-local-CM}，
$$
R\Hom_A(M, \omega_A^\bullet)_\mathfrak p =
R\Hom_{A_\mathfrak p}(M_\mathfrak p, (\omega_A^\bullet)_\mathfrak p)
$$
恰有一个非零上同调模，设它位于次数 $i_0$。由于
$R\Hom_A(M, \omega_A^\bullet)$ 只有有限多个非零上同调模 $H^i$，
而且每个都是有限 $A$-模，所以可以找到
$f \in A$、$f \not \in \mathfrak p$，使得当 $i \not = i_0$ 时
$(H^i)_f = 0$。于是 $R\Hom_A(M, \omega_A^\bullet)_f$
恰有一个非零上同调模；反向应用刚才的论证，得到
$D(f) \subset U$。
\end{proof}

\begin{lemma}
\label{dualizing-lemma-CM}
设 $A$ 为 Noetherian 环。若 $A$ 具有对偶复形
$\omega_A^\bullet$，则
$\{\mathfrak p \in \Spec(A) \mid A_\mathfrak p\text{ 为 Cohen--Macaulay}\}$
是 $\Spec(A)$ 的稠密开子集。
\end{lemma}

\begin{proof}
这是引理 \ref{dualizing-lemma-CM-open} 与定义的直接推论。
\end{proof}






\section{Gorenstein 环}
\label{dualizing-section-gorenstein}

\noindent
到目前为止，我们见到的唯一显式对偶复形，是域 $\kappa$ 上由
$\kappa$ 给出的复形，其中 $\kappa$ 为域；参见引理
\ref{dualizing-lemma-find-function} 的证明。由命题
\ref{dualizing-proposition-dualizing-essentially-finite-type}，任意域上的有限型
代数都有对偶复形。然而，确实存在没有对偶复形的 Noetherian
（局部）环。我们已经看到，具有对偶复形的环是普遍等长链环
（引理 \ref{dualizing-lemma-universally-catenary}），但也存在不是等长链环的
Noetherian 局部环；参见例子，第
\ref{examples-section-non-catenary-Noetherian-local} 节。

\medskip\noindent
尽管如此，代数几何中的许多环都具有对偶复形，原因仅仅是它们是
Gorenstein 环的商。事实上，这个条件既必要又充分：Noetherian 环
具有对偶复形，当且仅当它是某个有限维 Gorenstein 环的商。
这就是 Sharp 猜想（\cite{Sharp}）；在文献中可见
\cite[Corollary 1.4]{Kawasaki}。回到当前主题，下面给出
Gorenstein 环的定义。

\begin{definition}
\label{dualizing-definition-gorenstein}
Gorenstein 环。
\begin{enumerate}
\item 设 $A$ 为 Noetherian 局部环。若 $A[0]$ 是 $A$
的对偶复形，则称 $A$ 是{\it Gorenstein} 环。
\item 设 $A$ 为 Noetherian 环。若对 $A$ 的每个素理想
$\mathfrak p$，$A_\mathfrak p$ 都是 Gorenstein 环，则称 $A$
是{\it Gorenstein} 环。
\end{enumerate}
\end{definition}

\noindent
这个定义是合理的，因为若 $A[0]$ 是 $A$ 的对偶复形，则由引理
\ref{dualizing-lemma-dualizing-localize}，$S^{-1}A[0]$ 是 $S^{-1}A$
的对偶复形。若有限维 Noetherian 环作为自身上的模具有有限内射维数，
则它是 Gorenstein 环；参见引理 \ref{dualizing-lemma-characterize-gorenstein}。

\begin{lemma}
\label{dualizing-lemma-gorenstein-CM}
Gorenstein 环是 Cohen--Macaulay 环。
\end{lemma}

\begin{proof}
由引理 \ref{dualizing-lemma-apply-CM} 得出。
\end{proof}

\noindent
正则环是 Gorenstein 环的一个例子。

\begin{lemma}
\label{dualizing-lemma-regular-gorenstein}
正则局部环是 Gorenstein 环。正则环是 Gorenstein 环。
\end{lemma}

\begin{proof}
设 $A$ 为维数为 $d$ 的正则环。则 $A$ 的整体维数有限且等于
$d$；参见代数，引理
\ref{algebra-lemma-finite-gl-dim-finite-dim-regular}。
因此对所有 $A$-模 $M$，有 $\Ext^{d + 1}_A(M, A) = 0$；
参见代数，引理 \ref{algebra-lemma-projective-dimension-ext}。
由更多代数，引理 \ref{more-algebra-lemma-injective-amplitude}，
$A$ 作为 $A$-模具有有限内射维数。故 $A[0]$ 是对偶复形，
于是由定义后的评注，$A$ 是 Gorenstein 环。
\end{proof}

\begin{lemma}
\label{dualizing-lemma-gorenstein}
设 $A$ 为 Noetherian 环。
\begin{enumerate}
\item 若 $A$ 具有对偶复形 $\omega_A^\bullet$，则
\begin{enumerate}
\item $A$ 是 Gorenstein 环 $\Leftrightarrow$ $\omega_A^\bullet$
是 $D(A)$ 的可逆对象；
\item $A_\mathfrak p$ 是 Gorenstein 环 $\Leftrightarrow$
$(\omega_A^\bullet)_\mathfrak p$ 是 $D(A_\mathfrak p)$ 的可逆对象；
\item $\{\mathfrak p \in \Spec(A) \mid A_\mathfrak p\text{ is Gorenstein}\}$
是开子集。
\end{enumerate}
\item 若 $A$ 是 Gorenstein 环，则 $A$ 具有对偶复形当且仅当
$A[0]$ 是对偶复形。
\end{enumerate}
\end{lemma}

\begin{proof}
关于 $D(A)$ 的可逆对象，参见更多代数，引理
\ref{more-algebra-lemma-invertible-derived} 以及第
\ref{dualizing-section-dualizing} 节中的讨论。

\medskip\noindent
由引理 \ref{dualizing-lemma-dualizing-localize}，对每个 $\mathfrak p$，
复形 $(\omega_A^\bullet)_\mathfrak p$ 都是 $A_\mathfrak p$
上的对偶复形。由定义及对偶复形的唯一性（引理
\ref{dualizing-lemma-dualizing-unique}），得到 (1)(b)。

\medskip\noindent
为证明 (1)(c)，假设 $A_\mathfrak p$ 是 Gorenstein 环。
令 $n_x$ 为使 $H^{n_{x}}((\omega_A^\bullet)_\mathfrak p)$
非零且同构于 $A_\mathfrak p$ 的唯一整数。由于
$\omega_A^\bullet$ 属于 $D^b_{\textit{Coh}}(A)$，其非零上同调
$H^i(\omega_A^\bullet)$ 只有有限多个，且都是有限 $A$-模。
因此存在 $f \in A$、$f \not \in \mathfrak p$，使得只有
$H^{n_x}((\omega_A^\bullet)_f)$ 非零，并且它在 $A_f$ 上由
$1$ 个元素生成。由于对偶复形是忠实的（依定义），可知
$A_f \cong H^{n_x}((\omega_A^\bullet)_f)$。于是对每个
$\mathfrak q \in D(f)$，$A_\mathfrak q$ 都是 Gorenstein 环。
这证明了 (1)(c) 中的集合是开的。

\medskip\noindent
证明 (1)(a)。蕴含 $\Leftarrow$ 由 (1)(b) 得出。
蕴含 $\Rightarrow$ 由上一段的讨论得出：那里证明了，若
$A_\mathfrak p$ 是 Gorenstein 环，则存在
$f \in A$、$f \not \in \mathfrak p$，使得复形
$(\omega_A^\bullet)_f$ 只有一个非零上同调模，而且该模可逆。

\medskip\noindent
若 $A[0]$ 是对偶复形，则由第 (1) 部分，$A$ 是 Gorenstein 环。
反之，第 (1) 部分表明 $\omega_A^\bullet$ 局部同构于 $A$
的某个平移。成为对偶复形是局部性质（引理
\ref{dualizing-lemma-dualizing-glue}），故结论显然。
\end{proof}

\begin{lemma}
\label{dualizing-lemma-gorenstein-ext}
设 $(A, \mathfrak m, \kappa)$ 为 Noetherian 局部环。
则 $A$ 是 Gorenstein 环，当且仅当对 $i \gg 0$，
$\Ext^i_A(\kappa, A)$ 为零。
\end{lemma}

\begin{proof}
注意到，$A[0]$ 是 $A$ 的对偶复形，当且仅当 $A$ 作为
$A$-模具有有限内射维数（这由定义 \ref{dualizing-definition-dualizing}
立即得出）。因此该引理由更多代数，引理
\ref{more-algebra-lemma-finite-injective-dimension-Noetherian-local} 得出。
\end{proof}

\begin{lemma}
\label{dualizing-lemma-characterize-gorenstein}
设 $A$ 为 Noetherian 环。下列条件等价：
\begin{enumerate}
\item $A$ 作为自身上的模具有有限内射维数；
\item $A$ 是有限维 Gorenstein 环。
\end{enumerate}
在这种情况下，$A[0]$ 是对偶复形。
\end{lemma}

\begin{proof}
若 $A$ 作为 $A$-模具有有限内射维数，则 $A[0]$ 是 $A$
上的对偶复形（参见定义 \ref{dualizing-definition-dualizing}）。由引理
\ref{dualizing-lemma-universally-catenary}，$A$ 有限维；由引理
\ref{dualizing-lemma-gorenstein}，该环是 Gorenstein 环。
反之，假设 $A$ 是 Gorenstein 环且
$\dim(A) \leq d < \infty$。设 $I \subset A$ 为任意理想。
断言当 $i > d$ 时，$\Ext^i_A(A/I, A) = 0$；由更多代数，引理
\ref{more-algebra-lemma-injective-amplitude}，这将证明 $A$
作为自身上的模具有有限内射维数。构造 $\Ext$ 群与局部化可交换
（例如参见更多代数，引理
\ref{more-algebra-lemma-base-change-RHom}）。所以为证明消失，
可以假设 $A$ 是维数 $e < d$ 的局部环。由于 $A$ 是 Gorenstein
环，$\omega_A^\bullet = A[e]$ 是规范化对偶复形。
于是所断言的消失例如由引理 \ref{dualizing-lemma-sitting-in-degrees} 得出。
\end{proof}

\begin{lemma}
\label{dualizing-lemma-gorenstein-divide-by-nonzerodivisor}
设 $(A, \mathfrak m, \kappa)$ 为 Noetherian 局部环，
$f \in \mathfrak m$ 为非零因子，并置 $B = A/(f)$。
则 $A$ 是 Gorenstein 环，当且仅当 $B$ 是 Gorenstein 环。
\end{lemma}

\begin{proof}
若 $A$ 是 Gorenstein 环，则由引理
\ref{dualizing-lemma-divide-by-nonzerodivisor}，$B$ 是 Gorenstein 环。
反之，假设 $B$ 是 Gorenstein 环。则当 $i \gg 0$ 时，
$\Ext^i_B(\kappa, B)$ 为零（引理 \ref{dualizing-lemma-gorenstein-ext}）。
回忆 $R\Hom(B, -) : D(A) \to D(B)$ 是限制函子的右伴随
（引理 \ref{dualizing-lemma-right-adjoint}）。因此
$$
R\Hom_A(\kappa, A) = R\Hom_B(\kappa, R\Hom(B, A)) =
R\Hom_B(\kappa, B[1])
$$
最后一个等号可由直接计算或引理
\ref{dualizing-lemma-compute-for-effective-Cartier-algebraic} 得到。
所以当 $i \gg 0$ 时，$\Ext^i_A(\kappa, A)$ 为零，因而 $A$
是 Gorenstein 环（引理 \ref{dualizing-lemma-gorenstein-ext}）。
\end{proof}

\begin{lemma}
\label{dualizing-lemma-gorenstein-lci}
若 $A \to B$ 是环的局部完全交同态，且 $A$ 是 Noetherian
Gorenstein 环，则 $B$ 是 Gorenstein 环。
\end{lemma}

\begin{proof}
由更多代数，定义
\ref{more-algebra-definition-local-complete-intersection}，可写成
$B = A[x_1, \ldots, x_n]/I$，其中 $I$ 是 Koszul 正则理想。
注意 Gorenstein 环 $A$ 上的多项式环是 Gorenstein 环：先归约到
$A$ 为局部环，再使用引理 \ref{dualizing-lemma-dualizing-polynomial-ring}
和 \ref{dualizing-lemma-gorenstein}。依定义，Koszul 正则理想局部由
Koszul 正则序列生成；参见更多代数，第
\ref{more-algebra-section-ideals} 节。考察
$A[x_1, \ldots, x_n]$ 的局部环，可知只需证明：若 $R$
是 Noetherian 局部 Gorenstein 环，并且
$f_1, \ldots, f_c \in \mathfrak m_R$ 是 Koszul 正则序列，
则 $R/(f_1, \ldots, f_c)$ 是 Gorenstein 环。这由引理
\ref{dualizing-lemma-gorenstein-divide-by-nonzerodivisor} 以及如下事实得出：
$R$ 中的 Koszul 正则序列就是正则序列（更多代数，引理
\ref{more-algebra-lemma-noetherian-finite-all-equivalent}）。
\end{proof}

\begin{lemma}
\label{dualizing-lemma-flat-under-gorenstein}
设 $A \to B$ 为 Noetherian 局部环之间的平坦局部同态。
下列条件等价：
\begin{enumerate}
\item $B$ 是 Gorenstein 环；
\item $A$ 和 $B/\mathfrak m_A B$ 都是 Gorenstein 环。
\end{enumerate}
\end{lemma}

\begin{proof}
以下将不再另行说明地使用局部 Gorenstein 环具有有限内射维数
这一事实以及引理 \ref{dualizing-lemma-gorenstein-ext}。由更多代数，引理
\ref{more-algebra-lemma-pseudo-coherence-and-base-change-ext}，
对所有 $i$ 有
$$
\Ext^i_A(\kappa_A, A) \otimes_A B =
\Ext^i_B(B/\mathfrak m_A B, B)
$$

\medskip\noindent
假设 (2)。利用 $R\Hom(B/\mathfrak m_A B, -) : D(B) \to D(B/\mathfrak m_A B)$
是限制函子的右伴随（引理 \ref{dualizing-lemma-right-adjoint}），得到
$$
R\Hom_B(\kappa_B, B) =
R\Hom_{B/\mathfrak m_A B}(\kappa_B, R\Hom(B/\mathfrak m_A B, B))
$$
$R\Hom(B/\mathfrak m_A B, B)$ 的上同调模为
$\Ext^i_B(B/\mathfrak m_A B, B) =
\Ext^i_A(\kappa_A, A) \otimes_A B$。由于 $A$ 是 Gorenstein 环，
这些模中只有有限多个非零，并且每个都同构于若干份
$B/\mathfrak m_A B$ 的直和。又因为 $B/\mathfrak m_A B$
是 Gorenstein 环，可知 $R\Hom_B(B/\mathfrak m_B, B)$
只有有限多个非零上同调模。因此 $B$ 是 Gorenstein 环。

\medskip\noindent
假设 (1)。由于 $B$ 具有有限内射维数，当 $i \gg 0$ 时，
$\Ext^i_B(B/\mathfrak m_A B, B)$ 为 $0$。由于 $A \to B$
忠实平坦，当 $i \gg 0$ 时，$\Ext^i_A(\kappa_A, A)$ 为 $0$。
故 $A$ 是 Gorenstein 环。这意味着 $\Ext^i_A(\kappa_A, A)$
恰好对一个 $i$ 非零，即 $i = \dim(A)$，并且
$\Ext^{\dim(A)}_A(\kappa_A, A) \cong \kappa_A$
（参见引理 \ref{dualizing-lemma-normalized-finite}、\ref{dualizing-lemma-apply-CM}
及 \ref{dualizing-lemma-gorenstein-CM}）。所以
$\Ext^i_B(B/\mathfrak m_A B, B)$ 除了一个 $i$ 以外均为零，
即 $i = \dim(A)$，且
$\Ext^{\dim(A)}_B(B/\mathfrak m_A B, B) \cong B/\mathfrak m_A B$。
由引理 \ref{dualizing-lemma-normalized-finite}，$B/\mathfrak m_A B$
是 Gorenstein 环。
\end{proof}

\begin{lemma}
\label{dualizing-lemma-tor-injective-hull}
设 $(A, \mathfrak m, \kappa)$ 为维数 $d$ 的 Noetherian 局部
Gorenstein 环，$E$ 为 $\kappa$ 的内射包。则当 $i \not = d$ 时，
$\text{Tor}_i^A(E, \kappa)$ 为零，并且
$\text{Tor}_d^A(E, \kappa) = \kappa$。
\end{lemma}

\begin{proof}
由于 $A$ 是 Gorenstein 环，$\omega_A^\bullet = A[d]$ 是 $A$
的规范化对偶复形。此外，$E$ 是
$R\Gamma_\mathfrak m(\omega_A^\bullet)$ 唯一的非零上同调模，
并且位于次数 $0$；参见引理
\ref{dualizing-lemma-local-cohomology-of-dualizing}。由引理
\ref{dualizing-lemma-torsion-tensor-product}，有
$$
E \otimes_A^\mathbf{L} \kappa =
R\Gamma_\mathfrak m(\omega_A^\bullet) \otimes_A^\mathbf{L} \kappa =
R\Gamma_\mathfrak m(\omega_A^\bullet \otimes_A^\mathbf{L} \kappa) =
R\Gamma_\mathfrak m(\kappa[d]) = \kappa[d]
$$
引理随即得证。
\end{proof}




\section{对偶复形的普遍存在性}
\label{dualizing-section-ubiquity-dualizing}

\noindent
许多 Noetherian 环都有对偶复形。

\begin{lemma}
\label{dualizing-lemma-flat-unramified}
设 $A \to B$ 为 Noetherian 局部环之间的局部同态，
$\omega_A^\bullet$ 为规范化对偶复形。若 $A \to B$ 平坦且
$\mathfrak m_A B = \mathfrak m_B$，则
$\omega_A^\bullet \otimes_A B$ 是 $B$ 的规范化对偶复形。
\end{lemma}

\begin{proof}
显然 $\omega_A^\bullet \otimes_A B$ 属于
$D^b_{\textit{Coh}}(B)$。设 $\kappa_A$ 和 $\kappa_B$ 分别为
$A$ 和 $B$ 的剩余域。由更多代数，引理
\ref{more-algebra-lemma-base-change-RHom}，有
$$
R\Hom_B(\kappa_B, \omega_A^\bullet \otimes_A B) =
R\Hom_A(\kappa_A, \omega_A^\bullet) \otimes_A B =
\kappa_A[0] \otimes_A B = \kappa_B[0]
$$
因此由更多代数，引理
\ref{more-algebra-lemma-finite-injective-dimension-Noetherian-local}，
$\omega_A^\bullet \otimes_A B$ 具有有限内射维数。最后，用相同的论证可得
$$
R\Hom_B(\omega_A^\bullet \otimes_A B, \omega_A^\bullet \otimes_A B) =
R\Hom_A(\omega_A^\bullet, \omega_A^\bullet) \otimes_A B = A \otimes_A B = B
$$
这正是所需结论。
\end{proof}

\begin{lemma}
\label{dualizing-lemma-flat-iso-mod-I}
设 $A \to B$ 为 Noetherian 环之间的平坦映射。设
$I \subset A$ 为理想，满足 $A/I = B/IB$，并且 $IB$ 包含于
$B$ 的 Jacobson 根中。设 $\omega_A^\bullet$ 为对偶复形。
则 $\omega_A^\bullet \otimes_A B$ 是 $B$ 的对偶复形。
\end{lemma}

\begin{proof}
显然 $\omega_A^\bullet \otimes_A B$ 属于
$D^b_{\textit{Coh}}(B)$。由更多代数，引理
\ref{more-algebra-lemma-base-change-RHom}，对任意
$K \in D^b_{\textit{Coh}}(A)$ 有
$$
R\Hom_B(K \otimes_A B, \omega_A^\bullet \otimes_A B) =
R\Hom_A(K, \omega_A^\bullet) \otimes_A B
$$
对任意理想 $IB \subset J \subset B$，存在唯一理想
$I \subset J' \subset A$，使得 $A/J' \otimes_A B = B/J$。
因此由更多代数，引理
\ref{more-algebra-lemma-finite-injective-dimension-Noetherian-radical}，
$\omega_A^\bullet \otimes_A B$ 具有有限内射维数。最后还有
$$
R\Hom_B(\omega_A^\bullet \otimes_A B, \omega_A^\bullet \otimes_A B) =
R\Hom_A(\omega_A^\bullet, \omega_A^\bullet) \otimes_A B = A \otimes_A B = B
$$
这正是所需结论。
\end{proof}

\begin{lemma}
\label{dualizing-lemma-completion-henselization-dualizing}
设 $A$ 为 Noetherian 环，$I \subset A$ 为理想，
$\omega_A^\bullet$ 为对偶复形。
\begin{enumerate}
\item $\omega_A^\bullet \otimes_A A^h$ 是对偶复形，它定义在环对
$(A, I)$ 的 Hensel 化 $(A^h, I^h)$ 上；
\item $\omega_A^\bullet \otimes_A A^\wedge$ 是 $I$-进完备化
$A^\wedge$ 上的对偶复形；
\item 若 $A$ 局部，则 $\omega_A^\bullet \otimes_A A^h$，相应地，\ 
$\omega_A^\bullet \otimes_A A^{sh}$，分别是 $A$ 的 Hensel 化，
相应地，\ 严格 Hensel 化上的对偶复形。
\end{enumerate}
\end{lemma}

\begin{proof}
由引理 \ref{dualizing-lemma-flat-unramified} 和
\ref{dualizing-lemma-flat-iso-mod-I} 立即得出。关于完备化与 Hensel 化，
参见更多代数，第 \ref{more-algebra-section-henselian-pairs}、
\ref{more-algebra-section-permanence-completion}、
\ref{more-algebra-section-permanence-henselization} 节，以及代数，第
\ref{algebra-section-completion} 和
\ref{algebra-section-completion-noetherian} 节。
\end{proof}

\begin{lemma}
\label{dualizing-lemma-ubiquity-dualizing}
下列类型的环具有对偶复形：
\begin{enumerate}
\item 域；
\item Noetherian 完备局部环；
\item $\mathbf{Z}$；
\item Dedekind 整环；
\item 由上述某个环通过取本质有限型代数、取理想进完备化、
取 Hensel 化或取严格 Hensel 化而得到的任意环。
\end{enumerate}
\end{lemma}

\begin{proof}
第 (5) 部分由命题
\ref{dualizing-proposition-dualizing-essentially-finite-type} 和引理
\ref{dualizing-lemma-completion-henselization-dualizing} 得出。
由引理 \ref{dualizing-lemma-regular-gorenstein}，正则局部环具有对偶复形。
由 Cohen 结构定理（代数，定理
\ref{algebra-theorem-cohen-structure-theorem}），Noetherian 完备局部环
是某个正则局部环的商。设 $A$ 为 Dedekind 整环。则每个理想
$I$ 都是有限投射 $A$-模（这由代数，引理
\ref{algebra-lemma-finite-projective} 以及如下事实得出：$A$
的局部环都是离散赋值环，因而都是主理想整环）。于是由更多代数，引理
\ref{more-algebra-lemma-injective-amplitude}，每个 $A$-模的内射维数
都有限且至多为 $1$。容易推出 $A[0]$ 是对偶复形。
\end{proof}





\section{形式纤维}
\label{dualizing-section-formal-fibres}

\noindent
本节是更多代数，第
\ref{more-algebra-section-properties-formal-fibres} 节的继续。
那里对局部环具有 Cohen--Macaulay 形式纤维的 Noetherian 环 $A$
建立了一套相当完善的理论。具体而言，我们证明了：(1) 只需检验在
极大理想处的局部化之形式纤维是 Cohen--Macaulay 的；(2) $A$
上的有限型环继承此性质；(3) 对 $A$ 的任意完备化 $A^\wedge$，
$A \to A^\wedge$ 的纤维都是 Cohen--Macaulay 的；(4) $A$
的 Hensel 化继承此性质。参见更多代数，引理
\ref{more-algebra-lemma-check-P-ring-maximal-ideals}、命题
\ref{more-algebra-proposition-finite-type-over-P-ring}、引理
\ref{more-algebra-lemma-map-P-ring-to-completion-P} 和引理
\ref{more-algebra-lemma-henselization-pair-P-ring}。对于局部环具有
几何约化、几何正规、$(S_n)$ 或几何 $(R_n)$ 形式纤维的
Noetherian 环，也有类似结论。本节将看到，对于局部环具有
Gorenstein 或局部完全交形式纤维的 Noetherian 环，同样如此。
这与本章相关，因为具有对偶复形的 Noetherian 环就是一个例子。

\begin{lemma}
\label{dualizing-lemma-formal-fibres-gorenstein}
更多代数，第 \ref{more-algebra-section-properties-formal-fibres} 节的
性质 (A)、(B)、(C)、(D)、(E) 对
$P(k \to R) =$“$R$ 是 Gorenstein 环”成立。
\end{lemma}

\begin{proof}
我们已经知道把 Gorenstein 换成 Cohen--Macaulay 时结论成立，
所以每一步中都可以假设所处理的环是 Cohen--Macaulay 环。
这对证明并无多少实质帮助，但有助于直觉。

\medskip\noindent
第 (A) 部分。设 $K/k$ 为有限生成域扩张，$R$ 为 Gorenstein
$k$-代数。可以找到 $k$ 上的整体完全交
$A = k[x_1, \ldots, x_n]/(f_1, \ldots, f_c)$，使 $K$
同构于 $A$ 的分式域；参见代数，引理
\ref{algebra-lemma-colimit-syntomic}。于是
$R \to R \otimes_k A$ 是相对整体完全交。由引理
\ref{dualizing-lemma-gorenstein-lci}，$R \otimes_k A$ 是 Gorenstein 环。
所以作为其局部化，$R \otimes_k K$ 也是 Gorenstein 环。

\medskip\noindent
证明 (B)。这是显然的，因为一个环是 Gorenstein 环，当且仅当
它的所有局部环都是 Gorenstein 环。

\medskip\noindent
第 (C) 部分。设 $A \to B \to C$ 为 Noetherian 环之间的平坦映射。
假设 $A \to B$ 的纤维是 Gorenstein 的，且 $B \to C$ 正则。
需要证明 $A \to C$ 的纤维是 Gorenstein 的。显然可以假设
$A = k$ 为域，进而假设 $B \to C$ 是 Noetherian 局部环之间的
正则局部同态。此时 $B$ 是 Gorenstein 环，而
$C/\mathfrak m_B C$ 正则，特别地，它是 Gorenstein 环
（引理 \ref{dualizing-lemma-regular-gorenstein}）。由引理
\ref{dualizing-lemma-flat-under-gorenstein}，$C$ 是 Gorenstein 环。

\medskip\noindent
第 (D) 部分由引理 \ref{dualizing-lemma-flat-under-gorenstein} 得出。
第 (E) 部分是立即的，因为该条件不涉及底域。
\end{proof}

\begin{lemma}
\label{dualizing-lemma-dualizing-gorenstein-formal-fibres}
设 $A$ 为 Noetherian 局部环。若 $A$ 具有对偶复形，则 $A$
的形式纤维都是 Gorenstein 的。
\end{lemma}

\begin{proof}
设 $\mathfrak p$ 为 $A$ 的素理想。$A$ 在 $\mathfrak p$ 处的形式纤维
同构于 $A/\mathfrak p$ 在 $(0)$ 处的形式纤维。商环
$A/\mathfrak p$ 具有对偶复形（引理
\ref{dualizing-lemma-dualizing-quotient}）。因此只需检验 $A$ 为局部整环且
$\mathfrak p = (0)$ 时的陈述。设 $\omega_A^\bullet$ 为 $A$
的对偶复形。则 $\omega_A^\bullet \otimes_A A^\wedge$ 是完备化
$A^\wedge$ 的对偶复形（引理 \ref{dualizing-lemma-flat-unramified}）。
又 $\omega_A^\bullet \otimes_A K$ 是 $A$ 的分式域 $K$
的对偶复形（引理 \ref{dualizing-lemma-dualizing-localize}）。所以对某个
$n \in \mathbf{Z}$，$\omega_A^\bullet \otimes_A K$
同构于 $K[n]$。类似地，可知形式纤维 $A^\wedge \otimes_A K$
的一个对偶复形为
$$
\omega_A^\bullet \otimes_A A^\wedge \otimes_{A^\wedge} (A^\wedge \otimes_A K) =
(\omega_A^\bullet \otimes_A K) \otimes_K (A^\wedge \otimes_A K) \cong
(A^\wedge \otimes_A K)[n]
$$
这正是所需结论。
\end{proof}

\noindent
下面验证除幂代数，评注 \ref{dpa-remark-no-good-ci-map} 中所承诺的结论。

\begin{lemma}
\label{dualizing-lemma-formal-fibres-lci}
更多代数，第 \ref{more-algebra-section-properties-formal-fibres} 节的
性质 (A)、(B)、(C)、(D)、(E) 对
$P(k \to R) =$“$R$ 是局部完全交”成立。参见除幂代数，定义
\ref{dpa-definition-lci}。
\end{lemma}

\begin{proof}
第 (A) 部分。设 $K/k$ 为有限生成域扩张，$R$ 为一个是局部完全交的
$k$-代数。可以找到 $k$ 上的整体完全交
$A = k[x_1, \ldots, x_n]/(f_1, \ldots, f_c)$，使 $K$
同构于 $A$ 的分式域；参见代数，引理
\ref{algebra-lemma-colimit-syntomic}。于是
$R \to R \otimes_k A$ 是相对整体完全交。由除幂代数，引理
\ref{dpa-lemma-avramov}，$R \otimes_k A$ 是局部完全交。

\medskip\noindent
证明 (B)。这是显然的，因为一个环是局部完全交，当且仅当它的所有
局部环都是完全交。

\medskip\noindent
第 (C) 部分。设 $A \to B \to C$ 为 Noetherian 环之间的平坦映射。
假设 $A \to B$ 的纤维是局部完全交，并且 $B \to C$ 正则。
需要证明 $A \to C$ 的纤维是局部完全交。显然可以假设
$A = k$ 为域，进而假设 $B \to C$ 是 Noetherian 局部环之间的
正则局部同态。此时 $B$ 是完全交，而 $C/\mathfrak m_B C$ 正则，
特别地，依定义它是完全交。由除幂代数，引理
\ref{dpa-lemma-avramov}，$C$ 是完全交。

\medskip\noindent
第 (D) 部分由 (C) 中相同的论证以及除幂代数，引理
\ref{dpa-lemma-avramov} 的另一个方向得出。第 (E) 部分是立即的，
因为该条件不涉及底域。
\end{proof}






\section{代数意义下的上叹号函子}
\label{dualizing-section-relative-dualizing-complex-algebraic}

\noindent
对 Noetherian 环的有限型同态 $R \to A$，我们将构造一个在非唯一同构
意义下良定义的函子 $\varphi^! : D(R) \to D(A)$。正如将在概形的
对偶理论，评注 \ref{duality-remark-local-calculation-shriek} 中看到的，
它在同构意义下与概形对偶理论中出现的上叹号函子一致。
为说明这一构造的动机，先提及另外两个性质：
\begin{enumerate}
\item $\varphi^!$ 把 $R$ 的对偶复形（若存在）送到 $A$ 的对偶复形；
\item $\omega_{A/R}^\bullet = \varphi^!(R)$ 是一种相对对偶复形：
它属于 $D^b_{\textit{Coh}}(A)$；若 $R \to A$ 平坦，则它限制到
纤维上后是对偶复形。
\end{enumerate}
这些陈述就是引理 \ref{dualizing-lemma-shriek-dualizing-algebraic} 和
\ref{dualizing-lemma-relative-dualizing-algebraic}。

\medskip\noindent
设 $\varphi : R \to A$ 为 Noetherian 环的有限型同态。
如下定义函子 $\varphi^! : D(R) \to D(A)$：
\begin{enumerate}
\item 若 $\varphi : R \to A$ 满，则置
$\varphi^!(K) = R\Hom(A, K)$。这里使用第 \ref{dualizing-section-trivial} 节的函子
$R\Hom(A, -) : D(R) \to D(A)$；
\item 一般地，选择满射 $\psi : P \to A$，其中
$P = R[x_1, \ldots, x_n]$，并置
$\varphi^!(K) = \psi^!(K \otimes_R^\mathbf{L} P)[n]$。
这里使用更多代数，第 \ref{more-algebra-section-derived-base-change} 节的函子
$- \otimes_R^\mathbf{L} P : D(R) \to D(P)$。
\end{enumerate}
注意按多项式环中的变量数作平移 $[n]$。这一构造{\bf 不是}典范的，
而函子 $\varphi^!$ 仅在函子的（非唯一）同构意义下良定义
\footnote{可以使这一构造成为典范的：在构造中以
$\Omega^n_{P/R}[n]$ 代替 $P[n]$，并在引理
\ref{dualizing-lemma-well-defined} 中使用它。若要这样做，本节的内容会复杂得多。}。

\begin{lemma}
\label{dualizing-lemma-well-defined}
设 $\varphi : R \to A$ 为 Noetherian 环的有限型同态。
则函子 $\varphi^!$ 在同构意义下良定义。
\end{lemma}

\begin{proof}
假设 $\psi_1 : P_1 = R[x_1, \ldots, x_n] \to A$ 与
$\psi_2 : P_2 = R[y_1, \ldots, y_m] \to A$ 是两个从多项式环到
$A$ 的满射。于是得到交换图
$$
\xymatrix{
R[x_1, \ldots, x_n, y_1, \ldots, y_m]
\ar[d]^{x_i \mapsto g_i} \ar[rr]_-{y_j \mapsto f_j} & &
R[x_1, \ldots, x_n] \ar[d] \\
R[y_1, \ldots, y_m] \ar[rr] & & A
}
$$
其中选择 $f_j$ 与 $g_i$，使得 $\psi_1(f_j) = \psi_2(y_j)$ 以及
$\psi_2(g_i) = \psi_1(x_i)$。由对称性，只需证明用
$P \to A$ 与 $P[y_1, \ldots, y_m] \to A$ 定义的函子同构。
由归纳可以假设 $m = 1$，从而归约到下一段讨论的情形。

\medskip\noindent
给定 $\psi : P \to A$，并设 $\chi : P[y] \to A$ 在 $P$
上诱导 $\psi$。记 $Q = P[y]$。选取 $g \in P$，使
$\psi(g) = \chi(y)$。以 $\pi : Q \to P$ 表示满足
$\pi(y) = g$ 的 $P$-代数映射。于是 $\chi = \psi \circ \pi$；
又因为两者都是限制函子 $D(A) \to D(Q)$ 的右伴随
（见第 \ref{dualizing-section-trivial} 节），故
$\chi^! = \psi^! \circ \pi^!$。因此
$$
\chi^!\left(K \otimes_R^\mathbf{L} Q\right)[n + 1] =
\psi^!\left(\pi^!\left(K \otimes_R^\mathbf{L} Q\right)[1]\right)[n]
$$
所以只需证明
$\pi^!(K \otimes_R^\mathbf{L} Q[1]) = K \otimes_R^\mathbf{L} P$。
因此只需证明函子 $\pi^!(-) : D(Q) \to D(P)$ 同构于
$K \mapsto K \otimes_Q^\mathbf{L} P[-1]$。这由引理
\ref{dualizing-lemma-compute-for-effective-Cartier-algebraic} 得出。
\end{proof}

\begin{lemma}
\label{dualizing-lemma-shriek-boundedness}
设 $\varphi : R \to A$ 为 Noetherian 环的有限型同态。
\begin{enumerate}
\item $\varphi^!$ 把 $D^+(R)$ 映入 $D^+(A)$，并把
$D^+_{\textit{Coh}}(R)$ 映入 $D^+_{\textit{Coh}}(A)$；
\item 若 $\varphi$ 完美，则 $\varphi^!$ 把 $D^-(R)$ 映入
$D^-(A)$，把 $D^-_{\textit{Coh}}(R)$ 映入
$D^-_{\textit{Coh}}(A)$，并把 $D^b_{\textit{Coh}}(R)$ 映入
$D^b_{\textit{Coh}}(A)$。
\end{enumerate}
\end{lemma}

\begin{proof}
按 $\varphi^!$ 的定义选取分解 $R \to P \to A$。函子
$- \otimes_R^\mathbf{L} : D(R) \to D(P)$ 保持子范畴
$D^+, D^+_{\textit{Coh}}, D^-, D^-_{\textit{Coh}}, D^b_{\textit{Coh}}$。
由引理 \ref{dualizing-lemma-exact-support-coherent}，函子
$R\Hom(A, -) : D(P) \to D(A)$ 保持 $D^+$ 与
$D^+_{\textit{Coh}}$。若 $R \to A$ 完美，则 $A$ 作为
$P$-模是完美的；参见更多代数，引理
\ref{more-algebra-lemma-perfect-ring-map}。回忆 $R\Hom(A, K)$
限制到 $D(P)$ 后就是 $R\Hom_P(A, K)$。由更多代数，引理
\ref{more-algebra-lemma-dual-perfect-complex}，对某个完美对象
$E \in D(P)$ 有 $R\Hom_P(A, K) = E \otimes_P^\mathbf{L} K$。
由于可以用有限投射 $P$-模的有限复形表示 $E$，所以一旦 $K$
属于 $D^-(P), D^-_{\textit{Coh}}(P), D^b_{\textit{Coh}}(P)$，
$R\Hom_P(A, K)$ 也属于相应范畴。限制函子
$D(A) \to D(P)$ 反映这些子范畴，故证明完成。
\end{proof}

\begin{lemma}
\label{dualizing-lemma-shriek-dualizing-algebraic}
设 $\varphi$ 为 Noetherian 环的有限型同态。若
$\omega_R^\bullet$ 是 $R$ 的对偶复形，则
$\varphi^!(\omega_R^\bullet)$ 是 $A$ 的对偶复形。
\end{lemma}

\begin{proof}
由引理 \ref{dualizing-lemma-dualizing-polynomial-ring} 和
\ref{dualizing-lemma-dualizing-quotient} 得出。
\end{proof}

\begin{lemma}
\label{dualizing-lemma-flat-bc}
设 $R \to R'$ 为 Noetherian 环的平坦同态，
$\varphi : R \to A$ 为有限型环映射，并设
$\varphi' : R' \to A' = A \otimes_R R'$ 为 $\varphi$ 诱导的映射。
则对 $D(R)$ 中的 $K$，有函子性映射
$$
\varphi^!(K) \otimes_A^\mathbf{L} A' \longrightarrow
(\varphi')^!(K \otimes_R^\mathbf{L} R')
$$
并且当 $K \in D^+(R)$ 时，这些映射是同构。
\end{lemma}

\begin{proof}
选取分解 $R \to P \to A$，其中 $P$ 是 $R$ 上的多项式环。
基变换给出相应分解 $R' \to P' \to A'$。由更多代数，引理
\ref{more-algebra-lemma-double-base-change}，有
$(K \otimes_R^\mathbf{L} P) \otimes_P^\mathbf{L} P' =
(K \otimes_R^\mathbf{L} R') \otimes_{R'}^\mathbf{L} P'$，
故只需构造关于 $K$ 函子性的映射
$$
R\Hom(A, K \otimes_R^\mathbf{L} P[n]) \otimes_A^\mathbf{L} A'
\longrightarrow
R\Hom(A', (K \otimes_R^\mathbf{L} P[n]) \otimes_P^\mathbf{L} P')
$$
为此，使用映射 (\ref{dualizing-equation-base-change})；它是在第
\ref{dualizing-section-base-change-trivial-duality} 节中针对 $P, A, P', A'$
构造的。
由引理 \ref{dualizing-lemma-flat-bc-surjection}，当 $K \in D^+(R)$ 时，
该映射为同构。
\end{proof}

\begin{lemma}
\label{dualizing-lemma-bc}
设 $R \to R'$ 为 Noetherian 环的同态。设
$\varphi : R \to A$ 为完美环映射（更多代数，定义
\ref{more-algebra-definition-pseudo-coherent-perfect}），并且 $R'$ 与
$A$ 在 $R$ 上 Tor 独立。设
$\varphi' : R' \to A' = A \otimes_R R'$ 为 $\varphi$ 诱导的映射。
则对 $D(R)$ 中的 $K$，有函子性同构
$$
\varphi^!(K) \otimes_A^\mathbf{L} A' =
(\varphi')^!(K \otimes_R^\mathbf{L} R')
$$
\end{lemma}

\begin{proof}
可以选择分解 $R \to P \to A$，其中 $P$ 是 $R$ 上的多项式环，
而 $A$ 是完美 $P$-模；参见更多代数，引理
\ref{more-algebra-lemma-perfect-ring-map}。基变换给出相应分解
$R' \to P' \to A'$。由更多代数，引理
\ref{more-algebra-lemma-double-base-change}，有
$(K \otimes_R^\mathbf{L} P) \otimes_P^\mathbf{L} P' =
(K \otimes_R^\mathbf{L} R') \otimes_{R'}^\mathbf{L} P'$，
所以只需构造关于 $K$ 函子性的映射
$$
R\Hom(A, K \otimes_R^\mathbf{L} P[n]) \otimes_A^\mathbf{L} A'
\longrightarrow
R\Hom(A', (K \otimes_R^\mathbf{L} P[n]) \otimes_P^\mathbf{L} P')
$$
我们有
$$
A \otimes_P^\mathbf{L} P' = A \otimes_R^\mathbf{L} R' = A'
$$
第一个等号由更多代数，引理
\ref{more-algebra-lemma-base-change-comparison} 应用于
$R, R', P, P'$ 得出。第二个等号成立是因为 $A$ 与 $R'$ 在
$R$ 上 Tor 独立。因此 $A$ 与 $P'$ 在 $P$ 上 Tor 独立，
可以使用映射 (\ref{dualizing-equation-base-change})；它是在第
\ref{dualizing-section-base-change-trivial-duality} 节中针对
$P, A, P', A'$ 构造的，由此得到所需箭头。由引理
\ref{dualizing-lemma-bc-surjection}，为完成证明，只需证明 $A$ 是完美
$P$-模，而这已在上面说明。
\end{proof}

\begin{lemma}
\label{dualizing-lemma-bc-flat}
设 $R \to R'$ 为 Noetherian 环的同态，
$\varphi : R \to A$ 为平坦有限型映射，并设
$\varphi' : R' \to A' = A \otimes_R R'$ 为 $\varphi$ 诱导的映射。
则对 $D(R)$ 中的 $K$，有函子性同构
$$
\varphi^!(K) \otimes_A^\mathbf{L} A' =
(\varphi')^!(K \otimes_R^\mathbf{L} R')
$$
\end{lemma}

\begin{proof}
这是引理 \ref{dualizing-lemma-bc} 的特殊情形；依据是更多代数，引理
\ref{more-algebra-lemma-flat-finite-presentation-perfect}。
\end{proof}

\begin{lemma}
\label{dualizing-lemma-composition-shriek-algebraic}
设 $A \xrightarrow{a} B \xrightarrow{b} C$ 为 Noetherian 环的有限型
同态。则存在函子变换 $b^! \circ a^! \to (b \circ a)^!$，
它在 $D^+(A)$ 上为同构。
\end{lemma}

\begin{proof}
选取 $A$ 上的多项式环 $P = A[x_1, \ldots, x_n]$ 以及满射
$P \to B$。选取在 $B$ 上生成 $C$ 的元素
$c_1, \ldots, c_m \in C$。置 $Q = P[y_1, \ldots, y_m]$，并记
$Q' = Q \otimes_P B = B[y_1, \ldots, y_m]$。设
$\chi : Q' \to C$ 为把 $y_j$ 送到 $c_j$ 的满射。图如下：
$$
\xymatrix{
& Q \ar[r]_{\psi'} & Q' \ar[r]_\chi & C \\
A \ar[r] & P \ar[r]^\psi \ar[u] & B \ar[u]
}
$$
由引理 \ref{dualizing-lemma-flat-bc-surjection}，对 $M \in D(P)$ 有箭头
$\psi^!(M) \otimes_B^\mathbf{L} Q' \to (\psi')^!(M \otimes_P^\mathbf{L} Q)$，
当 $M$ 下有界时它是同构。此外，由第 \ref{dualizing-section-trivial} 节，
两个函子都是限制函子 $D(C) \to D(Q)$ 的右伴随，故
$\chi^! \circ (\psi')^! = (\chi \circ \psi')^!$。于是
\begin{align*}
b^!(a^!(K))
& =
\chi^!(\psi^!(K \otimes_A^\mathbf{L} P)[n] \otimes_B^\mathbf{L} Q)[m] \\
& \to
\chi^!((\psi')^!(K \otimes_A^\mathbf{L} P \otimes_P^\mathbf{L} Q))[n + m] \\
& =
(\chi \circ \psi')^!(K\otimes_A^\mathbf{L} Q)[n + m] \\
& =
(b \circ a)^!(K)
\end{align*}
除上述结论外，这里还使用了更多代数，引理
\ref{more-algebra-lemma-double-base-change}。
\end{proof}

\begin{lemma}
\label{dualizing-lemma-upper-shriek-finite}
设 $\varphi : R \to A$ 为 Noetherian 环之间的有限映射。
则 $\varphi^!$ 同构于第 \ref{dualizing-section-trivial} 节的函子
$R\Hom(A, -) : D(R) \to D(A)$。
\end{lemma}

\begin{proof}
假设 $A$ 在 $R$ 上由 $n > 1$ 个元素生成。则可以把
$R \to A$ 分解为两个有限环映射的复合，并且每一步的生成元数都
$< n$。利用引理 \ref{dualizing-lemma-composition-shriek-algebraic} 和引理
\ref{dualizing-lemma-composition-right-adjoints}，只需证明 $A$ 在 $R$
上由一个元素生成时的引理。由于 $A$ 在 $R$ 上有限，$A$
是 $B = R[x]/(f)$ 的商，其中 $f$ 是关于 $x$ 的首一多项式
（代数，引理 \ref{algebra-lemma-finite-is-integral}）。再次使用关于
复合的引理，以及依定义对满射已有一致性这一事实，只需对
$R \to B = R[x]/(f)$ 证明引理。在这种情况下，函子
$\varphi^!$ 同构于 $K \mapsto K \otimes_R^\mathbf{L} B$；
这可把引理 \ref{dualizing-lemma-compute-for-effective-Cartier-algebraic}
用于映射 $R[x] \to B$ 来证明（注意 $\varphi^!$ 的定义与该引理中的
平移相加为零）。对于函子 $R\Hom(B, -) : D(R) \to D(B)$，
由引理 \ref{dualizing-lemma-RHom-is-tensor-special}，只需证明作为 $B$-模有
$\Hom_R(B, R) \cong B$。假设 $f$ 的次数为 $d$。则
$1, x, \ldots, x^{d - 1}$ 给出 $B$ 的一组 $R$-基。
令 $\delta_i : B \to R$、$i = 0, \ldots, d - 1$ 为关于这组基
取出 $x^i$ 系数的 $R$-线性映射。于是
$\delta_0, \ldots, \delta_{d - 1}$ 是 $\Hom_R(B, R)$ 的一组基。
最后，对 $0 \leq i \leq d - 1$，计算表明
$$
x^i \delta_{d - 1} =
\delta_{d - 1 - i} + b_1 \delta_{d - i} + \ldots + b_i \delta_{d - 1}
$$
其中某些 $c_1, \ldots, c_d \in R$ 满足上述关系
\footnote{若 $f = x^d + a_1 x^{d - 1} + \ldots + a_d$，则
$c_1 = -a_1$、$c_2 = a_1^2 - a_2$、
$c_3 = -a_1^3 + 2a_1a_2 -a_3$，等等。}。
所以 $\Hom_R(B, R)$ 是由 $\delta_{d - 1}$ 生成的主 $B$-模。
考察秩可知，它是秩为 $1$ 的自由 $B$-模。
\end{proof}

\begin{lemma}
\label{dualizing-lemma-upper-shriek-localize}
设 $R$ 为 Noetherian 环，$f \in R$。若 $\varphi$ 表示映射
$R \to R_f$，则 $\varphi^!$ 同构于
$- \otimes_R^\mathbf{L} R_f$。更一般地，若
$\varphi : R \to R'$ 满足 $\Spec(R') \to \Spec(R)$ 为开浸入，
则 $\varphi^!$ 同构于 $- \otimes_R^\mathbf{L} R'$。
\end{lemma}

\begin{proof}
选取表示 $R \to R[x] \to R[x]/(fx - 1) = R_f$，并注意到
$fx - 1$ 是 $R[x]$ 中的非零因子。因此可以用引理
\ref{dualizing-lemma-compute-for-effective-Cartier-algebraic} 计算函子
$\varphi^!$。略去细节；注意 $\varphi^!$ 的定义与该引理中的平移
相加为零。

\medskip\noindent
一般情形下，注意到 $R' \otimes_R R' = R'$。所以结论由上面的
基变换结果得出。使用引理 \ref{dualizing-lemma-flat-bc} 或引理
\ref{dualizing-lemma-bc} 均可。
\end{proof}

\begin{lemma}
\label{dualizing-lemma-upper-shriek-is-tensor-functor}
设 $\varphi : R \to A$ 为 Noetherian 环之间的完美同态
（例如 $\varphi$ 是平坦有限型的）。则对 $K \in D(R)$，有
$\varphi^!(K) = K \otimes_R^\mathbf{L} \varphi^!(R)$。
\end{lemma}

\begin{proof}
（括号内的陈述由更多代数，引理
\ref{more-algebra-lemma-flat-finite-presentation-perfect} 得出。）
可以选择分解 $R \to P \to A$，其中 $P$ 是 $R$ 上 $n$ 个变量的
多项式环，而 $A$ 是完美 $P$-模；参见更多代数，引理
\ref{more-algebra-lemma-perfect-ring-map}。回忆
$\varphi^!(K) = R\Hom(A, K \otimes_R^\mathbf{L} P[n])$。
因此结论由引理 \ref{dualizing-lemma-RHom-is-tensor-special} 与更多代数，引理
\ref{more-algebra-lemma-double-base-change} 得出。
\end{proof}

\begin{lemma}
\label{dualizing-lemma-relative-dualizing-if-have-omega}
设 $\varphi : A \to B$ 为 Noetherian 环的有限型同态，
$\omega_A^\bullet$ 为 $A$ 的对偶复形，并置
$\omega_B^\bullet = \varphi^!(\omega_A^\bullet)$。对
$K \in D_{\textit{Coh}}(A)$ 记
$D_A(K) = R\Hom_A(K, \omega_A^\bullet)$，对
$L \in D_{\textit{Coh}}(B)$ 记
$D_B(L) = R\Hom_B(L, \omega_B^\bullet)$。则对
$K \in D_{\textit{Coh}}(A)$，有函子性同构
$$
\varphi^!(K) = D_B(D_A(K) \otimes_A^\mathbf{L} B)
$$
\end{lemma}

\begin{proof}
由引理 \ref{dualizing-lemma-shriek-dualizing-algebraic}，注意到
$\omega_B^\bullet$ 是 $B$ 的对偶复形。设 $A \to B \to C$
为 Noetherian 环的有限型同态。若引理对 $A \to B$ 和
$B \to C$ 成立，则它对 $A \to C$ 也成立。这由引理
\ref{dualizing-lemma-composition-shriek-algebraic} 以及如下事实得出：由引理
\ref{dualizing-lemma-dualizing}，$D_B \circ D_B \cong \text{id}$。
因此只需在 $A \to B$ 为满射以及 $B$ 是 $A$ 上的多项式环时
证明引理。

\medskip\noindent
假设 $B = A[x_1, \ldots, x_n]$。由于
$D_A \circ D_A \cong \text{id}$，只需对
$D_{\textit{Coh}}(A)$ 中的 $K$ 证明
$D_B(K \otimes_A B) \cong D_A(K) \otimes_A B[n]$。
选取表示 $\omega_A^\bullet$ 的内射对象有界复形 $I^\bullet$。
选取拟同构 $I^\bullet \otimes_A B \to J^\bullet$，其中
$J^\bullet$ 是 $B$-模的有界复形。给定 $A$-模复形
$K^\bullet$，考虑显然的复形映射
$$
\Hom^\bullet(K^\bullet, I^\bullet) \otimes_A B[n]
\longrightarrow
\Hom^\bullet(K^\bullet \otimes_A B, J^\bullet[n])
$$
左端表示 $D_A(K) \otimes_A B[n]$，右端表示
$D_B(K \otimes_A B)$。因此，只需证明当 $K^\bullet$
的上同调模是有限 $A$-模时，该映射是拟同构。注意在任一次数
$r$，两边复形的上同调都只依赖于有限多个 $K^i$。所以可以把
$K^\bullet$ 替换为截断，也就是说，可以假设 $K^\bullet$
表示 $D^-_{\textit{Coh}}(A)$ 的对象。此时 $K^\bullet$
拟同构于有限自由 $A$-模的上有界复形。因此可以假设
$K^\bullet$ 本身就是有限自由 $A$-模的上有界复形。
在这种情况下，容易看出所显示的映射是复形的同构，
从而完成此情形的证明。

\medskip\noindent
假设 $A \to B$ 满。以 $i_* : D(B) \to D(A)$ 表示限制函子，
并回忆 $\varphi^!(-) = R\Hom(A, -)$ 是 $i_*$ 的右伴随
（引理 \ref{dualizing-lemma-right-adjoint}）。对 $F \in D(B)$，有
\begin{align*}
\Hom_B(F, D_B(D_A(K) \otimes_A^\mathbf{L} B))
& =
\Hom_B((D_A(K) \otimes_A^\mathbf{L} B) \otimes_B^\mathbf{L} F,
\omega_B^\bullet) \\
& =
\Hom_A(D_A(K) \otimes_A^\mathbf{L} i_*F, \omega_A^\bullet) \\
& =
\Hom_A(i_*F, D_A(D_A(K))) \\
& =
\Hom_A(i_*F, K) \\
& =
\Hom_B(F, \varphi^!(K))
\end{align*}
第一个等号由更多代数，引理 \ref{more-algebra-lemma-internal-hom}
及 $D_B$ 的定义得到。第二个等号由上述伴随性及等式
$i_*((D_A(K) \otimes_A^\mathbf{L} B) \otimes_B^\mathbf{L} F) =
D_A(K) \otimes_A^\mathbf{L} i_*F$ 得到（更多代数，引理
\ref{more-algebra-lemma-derived-base-change}）。第三个等号由更多代数，
引理 \ref{more-algebra-lemma-internal-hom} 得出。第四个等号成立是因为
$D_A \circ D_A = \text{id}$。最后一个等号再次由伴随性得到。
因此由 Yoneda 引理，结论成立。
\end{proof}






\section{Noetherian 情形中的相对对偶复形}
\label{dualizing-section-relative-dualizing-complexes-Noetherian}

\noindent
设 $\varphi : R \to A$ 为 Noetherian 环的有限型同态。
把 $D(A)$ 的对象
$$
\omega_{A/R}^\bullet = \varphi^!(R)
$$
定义为{\it $A$ 在 $R$ 上的相对对偶复形}。这里的 $\varphi^!$
如第 \ref{dualizing-section-relative-dualizing-complex-algebraic} 节所述。
由该节内容可知，$\omega_{A/R}^\bullet$ 在（非唯一）同构意义下良定义。

\begin{lemma}
\label{dualizing-lemma-base-change-relative-algebraic}
设 $R \to R'$ 为 Noetherian 环的同态，$R \to A$ 为有限型映射，
并置 $A' = A \otimes_R R'$。若
\begin{enumerate}
\item $R \to R'$ 平坦；或
\item $R \to A$ 平坦；或
\item $R \to A$ 完美，并且 $R'$ 与 $A$ 在 $R$ 上 Tor 独立，
\end{enumerate}
则在 $D(A')$ 中有同构
$\omega_{A/R}^\bullet \otimes_A^\mathbf{L} A' \to \omega^\bullet_{A'/R'}$。
\end{lemma}

\begin{proof}
由引理 \ref{dualizing-lemma-flat-bc}、\ref{dualizing-lemma-bc-flat}、
\ref{dualizing-lemma-bc} 及定义得出。
\end{proof}

\begin{lemma}
\label{dualizing-lemma-relative-dualizing-algebraic}
设 $\varphi : R \to A$ 为 Noetherian 环的有限型映射。
\begin{enumerate}
\item 若 $\varphi$ 完美，则
\begin{enumerate}
\item $\omega_{A/R}^\bullet$ 属于 $D^b_{\textit{Coh}}(A)$；
\item $\omega_{A/R}^\bullet$ 在 $R$ 上具有有限 Tor 维数；
\item $A \to R\Hom_A(\omega_{A/R}^\bullet, \omega_{A/R}^\bullet)$
是同构。
\end{enumerate}
\item 若 $\varphi$ 平坦，则 (1)(a)、(1)(b)、(1)(c) 成立，并且
\begin{enumerate}
\item $\omega_{A/R}^\bullet$ 是 $R$-完美的
（更多代数，定义 \ref{more-algebra-definition-relatively-perfect}）；
\item 对每个到域的映射 $R \to k$，基变换
$\omega_{A/R}^\bullet \otimes_A^\mathbf{L} (A \otimes_R k)$
是 $A \otimes_R k$ 的对偶复形。
\end{enumerate}
\end{enumerate}
\end{lemma}

\begin{proof}
假设 $R \to A$ 完美。按照 $\varphi^!$ 的定义选取
$R \to P \to A$。则 $A$ 是完美 $P$-模（更多代数，引理
\ref{more-algebra-lemma-perfect-ring-map}）。由引理
\ref{dualizing-lemma-shriek-boundedness}，这表明 $\omega_{A/R}^\bullet$
属于 $D^b_{\textit{Coh}}(A)$，从而证明 (1)(a)。为证明
$\omega_{A/R}^\bullet$ 作为 $R$-模复形具有有限 Tor 维数，注意到
$\omega_{A/R}^\bullet = \varphi^!(R) = R\Hom(A, P)[n]$
在 $D(P)$ 中映到 $R\Hom_P(A, P)[n]$；后者在 $D(P)$ 中完美
（更多代数，引理 \ref{more-algebra-lemma-dual-perfect-complex}）。
由于 $R \to P$ 平坦，它在 $D(R)$ 中具有有限 Tor 维数。
这证明 (1)(b)。$D(A)$ 的对象
$R\Hom_A(\omega_{A/R}^\bullet, \omega_{A/R}^\bullet)$
在 $D(P)$ 中映到
\begin{align*}
R\Hom_P(\omega_{A/R}^\bullet, R\Hom(A, P)[n])
& =
R\Hom_P(R\Hom_P(A, P)[n], P)[n] \\
& =
R\Hom_P(R\Hom_P(A, P), P)
\end{align*}
由已经使用过的更多代数，引理
\ref{more-algebra-lemma-dual-perfect-complex}，这等于 $A$。
这证明 (1)(c)。

\medskip\noindent
假设 $\varphi$ 平坦。则 $R \to A$ 是完美环映射
（更多代数，引理
\ref{more-algebra-lemma-flat-finite-presentation-perfect}），
故 (1)(a)、(1)(b)、(1)(c) 成立。于是由 (1)(a)、(1)(b) 及定义，
$\omega_{A/R}^\bullet$ 当然是 $R$-完美的。设 $R \to k$ 如 (2)(b)
中所述。由引理 \ref{dualizing-lemma-base-change-relative-algebraic}，有同构
$$
\omega_{A/R}^\bullet \otimes_A^\mathbf{L} (A \otimes_R k) \cong
\omega^\bullet_{A \otimes_R k/k}
$$
而由引理 \ref{dualizing-lemma-shriek-dualizing-algebraic}，右端是对偶复形。
证明完成。
\end{proof}

\begin{lemma}
\label{dualizing-lemma-base-change-dualizing-over-field}
设 $K/k$ 为域扩张，$A$ 为有限型 $k$-代数，并令
$A_K = A \otimes_k K$。若 $\omega_A^\bullet$ 是 $A$ 的对偶复形，
则 $\omega_A^\bullet \otimes_A A_K$ 是 $A_K$ 的对偶复形。
\end{lemma}

\begin{proof}
由对偶复形的唯一性，为 $A$ 选择哪个对偶复形并不重要；
略去详细证明。以 $\varphi : k \to A$ 表示代数结构。由引理
\ref{dualizing-lemma-shriek-dualizing-algebraic}，可以取
$\omega_A^\bullet = \varphi^!(k[0])$。结论由引理
\ref{dualizing-lemma-relative-dualizing-algebraic} 得出。
\end{proof}

\begin{lemma}
\label{dualizing-lemma-lci-shriek}
设 $\varphi : R \to A$ 为 Noetherian 环的局部完全交同态。
则 $\omega_{A/R}^\bullet$ 是 $D(A)$ 的可逆对象，并且对所有
$K \in D(R)$，有
$\varphi^!(K) = K \otimes_R^\mathbf{L} \omega_{A/R}^\bullet$。
\end{lemma}

\begin{proof}
回忆由更多代数，引理 \ref{more-algebra-lemma-lci-perfect}，
局部完全交同态是完美环映射。因此最后一个陈述由引理
\ref{dualizing-lemma-upper-shriek-is-tensor-functor} 得出。由更多代数，定义
\ref{more-algebra-definition-local-complete-intersection}，可写成
$A = R[x_1, \ldots, x_n]/I$，其中 $I$ 是 Koszul 正则理想。
第 \ref{dualizing-section-relative-dualizing-complex-algebraic} 节中
$\varphi^!$ 的构造表明，只需在 $A = R/I$ 且
$I \subset R$ 为 Koszul 正则理想时证明引理。由更多代数，引理
\ref{more-algebra-lemma-invertible-derived}，检验
$\omega_{A/R}^\bullet$ 在 $D(A)$ 中可逆是 $\Spec(A)$ 上的局部问题。
此外，由引理 \ref{dualizing-lemma-flat-bc}，$\omega_{A/R}^\bullet$
的形成与 $R$ 上的局部化可交换。结合更多代数，定义
\ref{more-algebra-definition-regular-ideal}、引理
\ref{more-algebra-lemma-noetherian-finite-all-equivalent} 和代数，引理
\ref{algebra-lemma-regular-sequence-in-neighbourhood}，可以找到
$g_1, \ldots, g_r \in R$，它们在 $A$ 中生成单位理想，并且
$I_{g_j} \subset R_{g_j}$ 由正则序列生成。因此可以假设
$A = R/(f_1, \ldots, f_c)$，其中 $f_1, \ldots, f_c$ 是 $R$
中的正则序列。于是考虑环映射
$$
R \to R/(f_1) \to R/(f_1, f_2) \to \ldots \to R/(f_1, \ldots, f_c) = A
$$
使用引理 \ref{dualizing-lemma-composition-shriek-algebraic}
（以及已经证明的最后一个陈述），可知只需对每一步证明引理。
最后，若对某个非零因子 $f$ 有 $A = R/(f)$，则引理成立，
因为由引理 \ref{dualizing-lemma-compute-for-effective-Cartier-algebraic}，
$\varphi^!(R) = R\Hom(A, R)$ 可逆。
\end{proof}

\begin{lemma}
\label{dualizing-lemma-gorenstein-shriek}
设 $\varphi : R \to A$ 为 Noetherian 环的平坦有限型同态。
下列条件等价：
\begin{enumerate}
\item 对所有素理想 $\mathfrak p \subset R$，纤维
$A \otimes_R \kappa(\mathfrak p)$ 都是 Gorenstein 环；
\item $\omega_{A/R}^\bullet$ 是 $D(A)$ 的可逆对象；参见更多代数，
引理 \ref{more-algebra-lemma-invertible-derived}。
\end{enumerate}
\end{lemma}

\begin{proof}
若 (2) 成立，则纤维环 $A \otimes_R \kappa(\mathfrak p)$
具有可逆的对偶复形，因而是 Gorenstein 环。参见引理
\ref{dualizing-lemma-relative-dualizing-algebraic} 和 \ref{dualizing-lemma-gorenstein}。

\medskip\noindent
反之，假设 (1)。由引理 \ref{dualizing-lemma-shriek-boundedness}，
注意到 $\omega_{A/R}^\bullet$ 属于 $D^b_{\textit{Coh}}(A)$
（因为 Noetherian 环的平坦有限型同态完美；参见更多代数，引理
\ref{more-algebra-lemma-flat-finite-presentation-perfect}）。
选取位于 $\mathfrak p \subset R$ 之上的素理想
$\mathfrak q \subset A$。则
$$
\omega_{A/R}^\bullet \otimes_A^\mathbf{L} \kappa(\mathfrak q) =
\omega_{A/R}^\bullet \otimes_A^\mathbf{L}
(A \otimes_R \kappa(\mathfrak p))
\otimes_{(A \otimes_R \kappa(\mathfrak p))}^\mathbf{L}
\kappa(\mathfrak q)
$$
应用引理 \ref{dualizing-lemma-relative-dualizing-algebraic}、
\ref{dualizing-lemma-gorenstein} 和假设 (1)，可知此复形有 $1$ 个非零上同调群，
它是 $1$ 维 $\kappa(\mathfrak q)$-向量空间。由更多代数，引理
\ref{more-algebra-lemma-lift-bounded-pseudo-coherent-to-perfect}，
存在 $f \in A$、$f \not \in \mathfrak q$，使
$(\omega_{A/R}^\bullet)_f$ 是 $D(A_f)$ 的可逆对象。
这证明 (2) 成立。
\end{proof}

\noindent
下面的引理有助于看清：从 Noetherian 环传到其有限型代数时，
维数函数如何变化。

\begin{lemma}
\label{dualizing-lemma-shriek-normalized}
设 $\varphi : R \to A$ 为 Noetherian 环的有限型同态。
假设 $R$ 局部，并设 $\mathfrak m \subset A$ 为位于 $R$
的极大理想之上的极大理想。若 $\omega_R^\bullet$ 是 $R$
的规范化对偶复形，则
$\varphi^!(\omega_R^\bullet)_\mathfrak m$ 是 $A_\mathfrak m$
的规范化对偶复形。
\end{lemma}

\begin{proof}
已经知道 $\varphi^!(\omega_R^\bullet)$ 是 $A$ 的对偶复形；
参见引理 \ref{dualizing-lemma-shriek-dualizing-algebraic}。按照
$\varphi^!$ 的构造选择分解 $R \to P \to A$，其中
$P = R[x_1, \ldots, x_n]$。若能对 $R \to P$ 以及 $P$
中对应于 $\mathfrak m$ 的极大理想 $\mathfrak m'$ 证明引理，
则把引理 \ref{dualizing-lemma-normalized-quotient} 用于
$P_{\mathfrak m'} \to A_\mathfrak m$，或把引理
\ref{dualizing-lemma-quotient-function} 用于 $P \to A$，即可得到
$R \to A$ 的结论。在 $A = R[x_1, \ldots, x_n]$ 的情形，
例如由代数，引理 \ref{algebra-lemma-dimension-base-fibre-equals-total}
（并结合代数，引理 \ref{algebra-lemma-dim-affine-space} 计算纤维维数），
有 $\dim(A_\mathfrak m) = \dim(R) + n$。$\omega_R^\bullet$
规范化意味着 $i = -\dim(R)$ 是使 $H^i(\omega_R^\bullet)$ 非零的
最小指标（由引理 \ref{dualizing-lemma-sitting-in-degrees} 和
\ref{dualizing-lemma-nonvanishing-generically-local} 得出）。于是
$\varphi^!(\omega_R^\bullet)_\mathfrak m =
\omega_R^\bullet \otimes_R A_\mathfrak m[n]$ 的第一个非零上同调模
位于次数 $-\dim(R) - n$，因此它是 $A_\mathfrak m$
的规范化对偶复形。
\end{proof}

\begin{lemma}
\label{dualizing-lemma-relative-dualizing-trivial-vanishing}
设 $R \to A$ 为 Noetherian 环的有限型同态，
$\mathfrak q \subset A$ 为位于 $\mathfrak p \subset R$ 之上的素理想。
则
$$
H^i(\omega_{A/R}^\bullet)_\mathfrak q \not = 0
\Rightarrow - d \leq i
$$
其中 $d$ 是 $\Spec(A) \to \Spec(R)$ 在 $\mathfrak p$ 上方的纤维
于点 $\mathfrak q$ 处的维数。
\end{lemma}

\begin{proof}
如第 \ref{dualizing-section-relative-dualizing-complex-algebraic} 节所述，
选择分解 $R \to P \to A$，其中
$P = R[x_1, \ldots, x_n]$，使得
$\omega_{A/R}^\bullet = R\Hom(A, P)[n]$。需要证明
$R\Hom(A, P)_\mathfrak q$ 在次数 $< n - d$ 上的上同调消失。
由引理 \ref{dualizing-lemma-RHom-ext}，这意味着需要证明：当 $i < n - d$
时，$\Ext_P^i(P/I, P)_{\mathfrak r} = 0$，其中
$\mathfrak r \subset P$ 是对应于 $\mathfrak q$ 的素理想，而 $I$
是 $P \to A$ 的核。由更多代数，引理
\ref{more-algebra-lemma-pseudo-coherence-and-base-change-ext}，
可以把它改写为
$\Ext_{P_\mathfrak r}^i(P_\mathfrak r/IP_\mathfrak r, P_\mathfrak r)$。
因此，由引理 \ref{dualizing-lemma-depth}，需要证明
$$
\text{depth}_{IP_\mathfrak r}(P_\mathfrak r) \geq n - d
$$
由引理 \ref{dualizing-lemma-depth-flat-CM}，有
$$
\text{depth}_{IP_\mathfrak r}(P_\mathfrak r) \geq
\dim((P \otimes_R \kappa(\mathfrak p))_\mathfrak r) -
\dim((P/I \otimes_R \kappa(\mathfrak p))_\mathfrak r)
$$
由代数，引理 \ref{algebra-lemma-codimension}，右端的两个量相符。
\end{proof}

\begin{lemma}
\label{dualizing-lemma-relative-dualizing-flat-vanishing}
设 $R \to A$ 为 Noetherian 环的平坦有限型同态，
$\mathfrak q \subset A$ 为位于 $\mathfrak p \subset R$ 之上的素理想。
则
$$
H^i(\omega_{A/R}^\bullet)_\mathfrak q \not = 0
\Rightarrow - d \leq i \leq 0
$$
其中 $d$ 是 $\Spec(A) \to \Spec(R)$ 在 $\mathfrak p$ 上方的纤维
于点 $\mathfrak q$ 处的维数。若 $\Spec(A) \to \Spec(R)$
的所有纤维维数都 $\leq d$，则 $\omega_{A/R}^\bullet$ 作为
$R$-模复形的 Tor 振幅位于 $[-d, 0]$。
\end{lemma}

\begin{proof}
下界已在引理 \ref{dualizing-lemma-relative-dualizing-trivial-vanishing} 中证明。
如第 \ref{dualizing-section-relative-dualizing-complex-algebraic} 节所述，
选择分解 $R \to P \to A$，其中
$P = R[x_1, \ldots, x_n]$，使得
$\omega_{A/R}^\bullet = R\Hom(A, P)[n]$。上界意味着当 $i > n$
时，$\Ext^i_P(A, P)$ 为零。这由更多代数，引理
\ref{more-algebra-lemma-perfect-over-polynomial-ring} 得出；该引理表明
$A$ 是完美 $P$-模，其 Tor 振幅位于 $[-n, 0]$。

\medskip\noindent
证明最后一个陈述。设 $R \to R'$ 为 Noetherian 环的同态，并置
$A' = A \otimes_R R'$。则
$$
\omega_{A'/R'}^\bullet =
\omega_{A/R}^\bullet \otimes_A^\mathbf{L} A' =
\omega_{A/R}^\bullet \otimes_R^\mathbf{L} R'
$$
第一个同构由引理 \ref{dualizing-lemma-base-change-relative-algebraic} 得出；
第二个同构发生在 $D(R')$ 中，由更多代数，引理
\ref{more-algebra-lemma-base-change-comparison} 得出。由证明的第一部分
（注意 $\Spec(A') \to \Spec(R')$ 的纤维维数都 $\leq d$），
可知 $\omega_{A/R}^\bullet \otimes_R^\mathbf{L} R'$
只有次数 $[-d, 0]$ 上的上同调。取 $R' = R \oplus M$ 为用有限
$R$-模 $M$ 对 $R$ 作的平方零加厚，可见对任意有限 $R$-模 $M$，
$R\Hom(A, P) \otimes_R^\mathbf{L} M$ 只有区间 $[-d, 0]$
内的上同调。任意 $R$-模都是有限 $R$-模的滤过余极限，且张量积
与余极限可交换，结论随即得出。
\end{proof}

\begin{lemma}
\label{dualizing-lemma-relative-dualizing-CM-vanishing}
设 $R \to A$ 为 Noetherian 环的有限型同态，
$\mathfrak p \subset R$ 为素理想。假设
\begin{enumerate}
\item $R_\mathfrak p$ 是 Cohen--Macaulay 环；
\item 对任意极小素理想 $\mathfrak q \subset A$，有
$\text{trdeg}_{\kappa(R \cap \mathfrak q)} \kappa(\mathfrak q) \leq r$。
\end{enumerate}
则
$$
H^i(\omega_{A/R}^\bullet)_\mathfrak p \not = 0 \Rightarrow - r \leq i
$$
而且 $H^{-r}(\omega_{A/R}^\bullet)_\mathfrak p$ 作为
$A_\mathfrak p$-模满足 $(S_2)$。
\end{lemma}

\begin{proof}
由引理 \ref{dualizing-lemma-base-change-relative-algebraic}，可以把 $R$
替换为 $R_\mathfrak p$。因此可以假设 $R$ 是 Cohen--Macaulay
局部环；需要对 $A$-模 $H^i(\omega_{A/R}^\bullet)$ 证明引理的陈述。

\medskip\noindent
设 $R^\wedge$ 为 $R$ 的完备化。映射 $R \to R^\wedge$ 平坦，
且 $R^\wedge$ 是 Cohen--Macaulay 环（更多代数，引理
\ref{more-algebra-lemma-completion-CM}）。由于
$A \to A \otimes_R R^\wedge$ 平坦，$A \otimes_R R^\wedge$
的极小素理想位于 $A$ 的极小素理想之上（平坦映射满足下降性；
参见代数，引理 \ref{algebra-lemma-flat-going-down}）。因此由态射，引理
\ref{morphisms-lemma-dimension-fibre-after-base-change}，条件 (2)
对有限型环映射 $R^\wedge \to A \otimes_R R^\wedge$ 成立。
再次应用引理 \ref{dualizing-lemma-base-change-relative-algebraic}，只需对
$R^\wedge \to A \otimes_R R^\wedge$ 证明引理。于是利用引理
\ref{dualizing-lemma-ubiquity-dualizing}，可以假设 $R$ 是具有对偶复形
$\omega_R^\bullet$ 的 Noetherian 局部 Cohen--Macaulay 环。

\medskip\noindent
设 $\mathfrak m \subset A$ 为极大理想。只需对
$H^i(\omega_{A/R}^\bullet)_\mathfrak m$ 证明引理的陈述。
若 $\mathfrak m$ 不位于 $R$ 的极大理想之上，则把 $R$
替换为某个局部化即可归约到这一情形（略去一个小细节）。

\medskip\noindent
可以假设 $\omega_R^\bullet$ 规范化。置 $d = \dim(R)$，
则对某个 $R$-模 $\omega_R$，有
$\omega_R^\bullet = \omega_R[d]$；参见引理
\ref{dualizing-lemma-apply-CM}。置
$\omega_A^\bullet = \varphi^!(\omega_R^\bullet)$。
由引理 \ref{dualizing-lemma-relative-dualizing-if-have-omega}，有
$$
\omega_{A/R}^\bullet =
R\Hom_A(\omega_R[d] \otimes_R^\mathbf{L} A, \omega_A^\bullet)
$$
由维数公式，$\dim(A_\mathfrak m) \leq d + r$；参见态射，引理
\ref{morphisms-lemma-dimension-formula-general}，并利用 Hilbert 零点定理
所给出的事实：$\kappa(\mathfrak m)$ 在 $R$ 的剩余域上有限。
由引理 \ref{dualizing-lemma-shriek-normalized}，
$(\omega_A^\bullet)_\mathfrak m$ 是 $A_\mathfrak m$
的规范化对偶复形。因此只有当 $-d - r \leq i \leq 0$ 时，
$H^i((\omega_A^\bullet)_\mathfrak m)$ 才可能非零；参见引理
\ref{dualizing-lemma-sitting-in-degrees}。由于
$\omega_R[d] \otimes_R^\mathbf{L} A$ 位于次数 $\leq -d$，
所述消失结论成立。最后还有
$$
H^{-r}(\omega_{A/R}^\bullet)_\mathfrak m =
\Hom_A(\omega_R \otimes_R A, H^{-d - r}(\omega_A^\bullet))_\mathfrak m
$$
由引理 \ref{dualizing-lemma-depth-dualizing-module}，
$H^{-d - r}(\omega_A^\bullet)_\mathfrak m$ 满足 $(S_2)$。
所以由更多代数，引理 \ref{more-algebra-lemma-hom-into-S2}，
最后一个陈述成立。
\end{proof}



\section{关于对偶复形的更多结果}
\label{dualizing-section-more-dualizing}

\noindent
这里汇集若干不太适合放在其他地方的引理。

\begin{lemma}
\label{dualizing-lemma-descent}
设 $A \to B$ 为 Noetherian 环之间的忠实平坦映射。
若 $K \in D(A)$，且 $K \otimes_A^\mathbf{L} B$ 是 $B$
的对偶复形，则 $K$ 是 $A$ 的对偶复形。
\end{lemma}

\begin{proof}
由于 $A \to B$ 平坦，有
$H^i(K) \otimes_A B = H^i(K \otimes_A^\mathbf{L} B)$。
由于 $K \otimes_A^\mathbf{L} B$ 属于 $D^b_{\textit{Coh}}(B)$，
先得到 $K$ 属于 $D^b(A)$，再由代数，引理
\ref{algebra-lemma-descend-properties-modules} 可知 $H^i(K)$
是有限 $A$-模。设 $M$ 为有限 $A$-模。由更多代数，引理
\ref{more-algebra-lemma-base-change-RHom}，有
$$
R\Hom_A(M, K) \otimes_A B = R\Hom_B(M \otimes_A B, K \otimes_A^\mathbf{L} B)
$$
由于 $K \otimes_A^\mathbf{L} B$ 具有有限内射维数，设其内射振幅
位于 $[a, b]$，右端在次数 $> b$ 上的上同调消失。
又因为 $A \to B$ 忠实平坦，$R\Hom_A(M, K)$ 在次数 $> b$
上的上同调也消失。因此由更多代数，引理
\ref{more-algebra-lemma-injective-amplitude}，$K$ 具有有限内射维数。
为完成证明，需要说明映射 $A \to R\Hom_A(K, K)$ 是同构。
为此再次使用更多代数，引理
\ref{more-algebra-lemma-base-change-RHom}，以及
$B \to R\Hom_B(K \otimes_A^\mathbf{L} B, K \otimes_A^\mathbf{L} B)$
是同构这一事实。
\end{proof}

\begin{lemma}
\label{dualizing-lemma-descent-ascent}
设 $\varphi : A \to B$ 为 Noetherian 环的同态。假设
\begin{enumerate}
\item $A \to B$ 是 syntomic 的，并在谱上诱导满射；或
\item $A \to B$ 是忠实平坦的局部完全交；或
\item $A \to B$ 是具有 Gorenstein 纤维的忠实平坦有限型映射。
\end{enumerate}
则 $K \in D(A)$ 是 $A$ 的对偶复形，当且仅当
$K \otimes_A^\mathbf{L} B$ 是 $B$ 的对偶复形。
\end{lemma}

\begin{proof}
由更多代数，引理 \ref{more-algebra-lemma-syntomic-lci}，
$A \to B$ 满足 (1) 当且仅当 $A \to B$ 满足 (2)。注意在 (2)、(3) 中，
相对对偶复形 $\varphi^!(A) = \omega_{B/A}^\bullet$ 都是
$D(B)$ 的可逆对象；参见引理 \ref{dualizing-lemma-lci-shriek} 和
\ref{dualizing-lemma-gorenstein-shriek}。此外，两种情形下都有
$\varphi^!(K) = K \otimes_A^\mathbf{L} \omega_{B/A}^\bullet$；
对于情形 (3)，参见引理
\ref{dualizing-lemma-upper-shriek-is-tensor-functor}。因此，除张量一个
$D(B)$ 的可逆对象以外，$\varphi^!(K)$ 与
$K \otimes_A^\mathbf{L} B$ 相同。所以 $\varphi^!(K)$ 是 $B$
的对偶复形，当且仅当 $K \otimes_A^\mathbf{L} B$ 是对偶复形
（成为对偶复形是局部性质，并且在平移下不变）。由此，若 $K$
是 $A$ 的对偶复形，则由引理
\ref{dualizing-lemma-shriek-dualizing-algebraic}，$K \otimes_A^\mathbf{L} B$
是 $B$ 的对偶复形。关于此性质的下降，参见引理
\ref{dualizing-lemma-descent}。
\end{proof}

\begin{lemma}
\label{dualizing-lemma-injective-hull-goes-up}
设 $(A, \mathfrak m, \kappa) \to (B, \mathfrak n, l)$ 为
Noetherian 环的平坦局部同态，且 $\mathfrak n = \mathfrak m B$。
若 $E$ 是 $\kappa$ 的内射包，则 $E \otimes_A B$ 是 $l$ 的内射包。
\end{lemma}

\begin{proof}
如引理 \ref{dualizing-lemma-union-artinian} 中那样，写成
$E = \bigcup E_n$。只需证明
$E_n \otimes_{A/\mathfrak m^n} B/\mathfrak n^n$ 是 $l$
在 $B/\mathfrak n$ 上的内射包。这把问题归约到 $A$ 与 $B$
都是 Artinian 局部环的情形。由代数，引理
\ref{algebra-lemma-pullback-module}，注意到
$\text{length}_A(A) = \text{length}_B(B)$ 以及
$\text{length}_A(E) = \text{length}_B(E \otimes_A B)$。
由引理 \ref{dualizing-lemma-finite}，有
$\text{length}_A(E) = \text{length}_A(A)$ 以及
$\text{length}_B(E') = \text{length}_B(B)$，其中 $E'$ 是 $l$
在 $B$ 上的内射包。因此
$\text{length}_B(E') = \text{length}_B(E \otimes_A B)$。又有
$$
\dim_l((E \otimes_A B)[\mathfrak n]) =
\dim_l(E[\mathfrak m] \otimes_A B) =
\dim_\kappa(E[\mathfrak m]) = 1
$$
这里使用了 $A \to B$ 的平坦性及 $\mathfrak n = \mathfrak mB$。
所以由引理 \ref{dualizing-lemma-torsion-submodule-sum-injective-hulls}，
存在单射的 $B$-模映射 $E \otimes_A B \to E'$。
由上面证明的长度相等性，它是同构。
\end{proof}

\begin{lemma}
\label{dualizing-lemma-injective-goes-up}
设 $\varphi : A \to B$ 为 Noetherian 环的平坦同态，并且对所有
素理想 $\mathfrak q \subset B$，都有
$\mathfrak p B_\mathfrak q = \mathfrak qB_\mathfrak q$，其中
$\mathfrak p = \varphi^{-1}(\mathfrak q)$；例如，$\varphi$ 可以是
\'etale 的。若 $I$ 是内射 $A$-模，则 $I \otimes_A B$ 是内射
$B$-模。
\end{lemma}

\begin{proof}
由代数，引理 \ref{algebra-lemma-etale-at-prime}，\'Etale 映射满足该假设。
由引理 \ref{dualizing-lemma-sum-injective-modules} 和命题
\ref{dualizing-proposition-structure-injectives-noetherian}，可以假设 $I$
是某个素理想 $\mathfrak p \subset A$ 的剩余域
$\kappa(\mathfrak p)$ 的内射包。于是 $I$ 是 $A_\mathfrak p$-模。
只需证明 $I \otimes_A B = I \otimes_{A_\mathfrak p} B_\mathfrak p$
作为 $B_\mathfrak p$-模内射；参见引理
\ref{dualizing-lemma-injective-flat}。因此可以假设
$(A, \mathfrak m, \kappa)$ 是 Noetherian 局部环，且 $I = E$
是剩余域 $\kappa$ 的内射包。我们的假设蕴含 Noetherian 环
$B/\mathfrak m B$ 是若干域的乘积（略去细节）。因此 $B$
中位于 $\mathfrak m$ 之上的素理想只有有限多个，记为
$\mathfrak m_1, \ldots, \mathfrak m_n$，而且它们都是极大理想。
如引理 \ref{dualizing-lemma-union-artinian} 中那样，写成
$E = \bigcup E_n$。则 $E \otimes_A B = \bigcup E_n \otimes_A B$，
而 $E_n \otimes_A B$ 是支集为
$\{\mathfrak m_1, \ldots, \mathfrak m_n\}$ 的有限 $B$-模，
因而分解为其在 $\mathfrak m_i$ 处局部化的乘积。所以
$E \otimes_A B = \prod (E \otimes_A B)_{\mathfrak m_i}$。
又由引理 \ref{dualizing-lemma-injective-hull-goes-up}，
$(E \otimes_A B)_{\mathfrak m_i} = E \otimes_A B_{\mathfrak m_i}$
是 $\mathfrak m_i$ 的剩余域之内射包，结论成立。
\end{proof}






\section{相对对偶复形}
\label{dualizing-section-relative-dualizing-complexes}

\noindent
对 Noetherian 环的有限型环映射 $\varphi : R \to A$，
第 \ref{dualizing-section-relative-dualizing-complexes-Noetherian} 节考虑了相对对偶复形
$\omega_{A/R}^\bullet = \varphi^!(R)$。若 $R$ 不是 Noetherian 环，
类似构造的复形一般不具有良好性质。本节对非 Noetherian 环的平坦有限
表现环映射 $R \to A$ 定义相对对偶复形。该定义被选择为能够推广到
概形的平坦有限表现态射；参见概形的对偶理论，第
\ref{duality-section-relative-dualizing-complexes} 节。我们将证明：
在该定义适用时，相对对偶复形存在；它在（非典范）同构意义下唯一；
并且在 Noetherian 情形中，恢复第
\ref{dualizing-section-relative-dualizing-complexes-Noetherian} 节的复形。

\medskip\noindent
只关心 Noetherian 情形的读者可以放心跳过本节！

\begin{definition}
\label{dualizing-definition-relative-dualizing-complex}
设 $R \to A$ 为平坦有限表现环映射。若对象 $K \in D(A)$ 满足
\begin{enumerate}
\item $K$ 是 $R$-完美的（更多代数，定义
\ref{more-algebra-definition-relatively-perfect}）；
\item $R\Hom_{A \otimes_R A}(A, K \otimes_A^\mathbf{L} (A \otimes_R A))$
同构于 $A$，
\end{enumerate}
则称这个对象为一个{\it 相对对偶复形}。
\end{definition}

\noindent
为理解这个定义，读者可能需要阅读并理解下面的一些引理。引理
\ref{dualizing-lemma-relative-dualizing-noetherian} 和
\ref{dualizing-lemma-uniqueness-relative-dualizing} 表明，这个定义与第
\ref{dualizing-section-relative-dualizing-complexes-Noetherian} 节中的定义并不冲突。

\begin{lemma}
\label{dualizing-lemma-uniqueness-relative-dualizing}
设 $R \to A$ 为平坦有限表现环映射。$R \to A$
的任意两个相对对偶复形均同构。
\end{lemma}

\begin{proof}
设 $K$ 与 $L$ 是 $R \to A$ 的两个相对对偶复形。记
$K_1 = K \otimes_A^\mathbf{L} (A \otimes_R A)$ 以及
$L_2 = (A \otimes_R A) \otimes_A^\mathbf{L} L$，它们分别是经由
第一、第二余投影 $A \to A \otimes_R A$ 所作的导出基变换。
由对称性，对 $L_2$ 的假设蕴含
$R\Hom_{A \otimes_R A}(A, L_2)$ 同构于 $A$。把更多代数，引理
\ref{more-algebra-lemma-internal-hom-evaluate-tensor-isomorphism}
的第 (3) 部分应用两次，得到
$$
A \otimes_{A \otimes_R A}^\mathbf{L} L_2 \cong
R\Hom_{A \otimes_R A}(A, K_1 \otimes_{A \otimes_R A}^\mathbf{L} L_2) \cong
A \otimes_{A \otimes_R A}^\mathbf{L} K_1
$$
对任一余投影应用限制函子 $D(A \otimes_R A) \to D(A)$，
即得所需结论。
\end{proof}

\begin{lemma}
\label{dualizing-lemma-relative-dualizing-noetherian}
设 $\varphi : R \to A$ 为 Noetherian 环的平坦有限型环映射。
则第 \ref{dualizing-section-relative-dualizing-complexes-Noetherian} 节的相对对偶复形
$\omega_{A/R}^\bullet = \varphi^!(R)$ 是定义
\ref{dualizing-definition-relative-dualizing-complex} 意义下的相对对偶复形。
\end{lemma}

\begin{proof}
由引理 \ref{dualizing-lemma-relative-dualizing-algebraic}，
$\varphi^!(R)$ 是 $R$-完美的。以
$\delta : A \otimes_R A \to A$ 表示乘法映射，以
$p_1, p_2 : A \to A \otimes_R A$ 表示两个余投影。则由引理
\ref{dualizing-lemma-flat-bc}，
$$
\varphi^!(R) \otimes_A^\mathbf{L} (A \otimes_R A) =
\varphi^!(R) \otimes_{A, p_1}^\mathbf{L} (A \otimes_R A) =
p_2^!(A)
$$
回忆
$
R\Hom_{A \otimes_R A}(A, \varphi^!(R) \otimes_A^\mathbf{L} (A \otimes_R A))
$
是 $\delta^!(\varphi^!(R) \otimes_A^\mathbf{L} (A \otimes_R A))$
在限制映射 $\delta_* : D(A) \to D(A \otimes_R A)$ 下的像。
这里使用第 \ref{dualizing-section-relative-dualizing-complex-algebraic} 节中
$\delta^!$ 的定义和引理 \ref{dualizing-lemma-RHom-ext}。由引理
\ref{dualizing-lemma-composition-shriek-algebraic}，
$\delta^!(p_2^!(A)) \cong A$，结论成立。
\end{proof}

\begin{lemma}
\label{dualizing-lemma-base-change-relative-dualizing}
设 $R \to A$ 为平坦有限表现环映射。则
\begin{enumerate}
\item 在 $D(A)$ 中存在相对对偶复形 $K$；
\item 对任意环映射 $R \to R'$，令 $A' = A \otimes_R R'$ 且
$K' = K \otimes_A^\mathbf{L} A'$，则 $K'$ 是 $R' \to A'$
的相对对偶复形。
\end{enumerate}
此外，若
$$
\xi : A \longrightarrow K \otimes_A^\mathbf{L} (A \otimes_R A)
$$
是循环模
$\Hom_{D(A \otimes_R A)}(A, K \otimes_A^\mathbf{L} (A \otimes_R A))$
的一个生成元，则在 (2) 中，$\xi$ 关于
$A \otimes_R A \to A' \otimes_{R'} A'$ 的导出基变换是循环模
$\Hom_{D(A' \otimes_{R'} A')}(A',
K' \otimes_{A'}^\mathbf{L} (A' \otimes_{R'} A'))$
的一个生成元。
\end{lemma}

\begin{proof}
先归约到 Noetherian 情形。由代数，引理
\ref{algebra-lemma-flat-finite-presentation-limit-flat}，存在有限型
$\mathbf{Z}$-子代数 $R_0 \subset R$ 和平坦有限型环映射
$R_0 \to A_0$，使得 $A = A_0 \otimes_{R_0} R$。由引理
\ref{dualizing-lemma-relative-dualizing-noetherian}，存在相对对偶复形
$K_0 \in D(A_0)$。因此，只要对 $K_0$ 证明 (2)，便可知
$K_0 \otimes_{A_0}^\mathbf{L} A$ 是 $R \to A$ 的对偶复形，
并由导出基变换的传递性知它也满足 (2)。相对对偶复形的唯一性
（引理 \ref{dualizing-lemma-uniqueness-relative-dualizing}）于是表明，
这对任意相对对偶复形都成立。

\medskip\noindent
设 $R$ 为 Noetherian 环，且设 $K$ 是 $R \to A$ 的相对对偶复形。
给定环映射 $R \to R'$，令 $A' = A \otimes_R R'$ 且
$K' = K \otimes_A^\mathbf{L} A'$。为完成证明，须证明 $K'$
是 $R' \to A'$ 的相对对偶复形。由更多代数，引理
\ref{more-algebra-lemma-base-change-relatively-perfect}，
$K'$ 在所有情形中均为 $R'$-完美的。由引理
\ref{dualizing-lemma-base-change-relative-algebraic} 和
\ref{dualizing-lemma-relative-dualizing-noetherian}，若 $R'$ 为 Noetherian 环，
则 $K'$ 是 $R' \to A'$ 的相对对偶复形（按任一种意义）。
导出张量积的传递性给出
$K \otimes_A^\mathbf{L} (A \otimes_R A)
\otimes_{A \otimes_R A}^\mathbf{L} (A' \otimes_{R'} A') =
K' \otimes_{A'}^\mathbf{L} (A' \otimes_{R'} A')$。
$R \to A$ 的平坦性保证
$A \otimes_{A \otimes_R A}^\mathbf{L} (A' \otimes_{R'} A') = A'$；
事实上，$A \otimes_R A$ 与 $R'$ 在 $R$ 上 Tor 独立，故可应用
更多代数，引理 \ref{more-algebra-lemma-base-change-comparison}。
最后，由更多代数，引理
\ref{more-algebra-lemma-more-relative-pseudo-coherent-is-moot}，
$A$ 作为 $A \otimes_R A$-模是伪凝聚的。因此，更多代数，引理
\ref{more-algebra-lemma-compute-RHom-relatively-perfect}
的全部假设均已验证。于是存在一个由 $R$-平坦有限表现
$A \otimes_R A$-模组成的下有界复形 $E^\bullet$，使得
$E^\bullet \otimes_R R'$ 表示
$R\Hom_{A' \otimes_{R'} A'}(A',
K' \otimes_{A'}^\mathbf{L} (A' \otimes_{R'} A'))$，
并且这些识别与导出基变换相容。
取 $n \in \mathbf{Z}$，$n \not = 0$。用正合列
$$
E^{n - 1} \to E^n \to Q^n \to 0
$$
定义 $Q^n$。由于 $\kappa(\mathfrak p)$ 是 Noetherian 环，
根据上面的论述，
$H^n(E^\bullet \otimes_R \kappa(\mathfrak p)) = 0$。
追踪交换图可知
$$
Q^n \otimes_R \kappa(\mathfrak p) \to E^{n + 1} \otimes_R \kappa(\mathfrak p)
$$
是单射。因此，对 $A \otimes_R A$ 中位于 $\mathfrak p$ 上方的素理想
$\mathfrak q$，$Q^n_\mathfrak q$ 是 $R_\mathfrak p$-平坦的，且
$Q^n_\mathfrak p \to E^{n + 1}_\mathfrak q$ 是
$R_\mathfrak p$-普遍单射；参见代数，引理
\ref{algebra-lemma-mod-injective}。因为这对所有素理想都成立，
故 $Q^n$ 是 $R$-平坦的，且 $Q^n \to E^{n + 1}$ 是
$R$-普遍单射。特别地，对任意环映射 $R \to R'$，
$H^n(E^\bullet \otimes_R R') = 0$。
令 $Z^0 = \Ker(E^0 \to E^1)$。由正合列
$0 \to Z^0 \to E^0 \to E^1 \to Q^1 \to 0$ 可知，
$Z^0$ 是 $R$-平坦的，并且对所有 $R \to R'$ 均有
$Z^0 \otimes_R R' = \Ker(E^0 \otimes_R R' \to E^1 \otimes_R R')$。
于是短正合列
$0 \to Q^{-1} \to Z^0 \to H^0(E^\bullet) \to 0$
表明
$$
H^0(E^\bullet \otimes_R R') = H^0(E^\bullet) \otimes_R R'
= A \otimes_R R' = A'
$$
正如所需。这个等式还给出了引理的最后一个断言。
\end{proof}

\begin{lemma}
\label{dualizing-lemma-relative-dualizing-RHom}
设 $R \to A$ 为平坦有限表现环映射，且设 $K$ 为相对对偶复形。
则 $A \to R\Hom_A(K, K)$ 是同构。
\end{lemma}

\begin{proof}
由代数，引理 \ref{algebra-lemma-flat-finite-presentation-limit-flat}，
存在有限型 $\mathbf{Z}$-子代数 $R_0 \subset R$ 和平坦有限型环映射
$R_0 \to A_0$，使得 $A = A_0 \otimes_{R_0} R$。由引理
\ref{dualizing-lemma-uniqueness-relative-dualizing}、
\ref{dualizing-lemma-relative-dualizing-noetherian} 和
\ref{dualizing-lemma-base-change-relative-dualizing}，存在相对对偶复形
$K_0 \in D(A_0)$，其导出基变换就是 $K$。
这把问题归约到下一段所讨论的情形。

\medskip\noindent
设 $R$ 为 Noetherian 环，且设 $K$ 是 $R \to A$ 的相对对偶复形。
给定环映射 $R \to R'$，令 $A' = A \otimes_R R'$ 且
$K' = K \otimes_A^\mathbf{L} A'$。为完成证明，只需证明
$R\Hom_{A'}(K', K') = A'$。由引理
\ref{dualizing-lemma-relative-dualizing-algebraic}，当 $R'$ 为 Noetherian 环时，
这已经成立。一般的 $R'$ 是 Noetherian $R$-代数的滤过余极限，
故由更多代数，引理
\ref{more-algebra-lemma-colimit-relatively-perfect}，结论成立。
\end{proof}

\begin{lemma}
\label{dualizing-lemma-relative-dualizing-composition}
设 $R \to A \to B$ 为平坦有限表现环映射。设 $K_{A/R}$ 和
$K_{B/A}$ 分别是 $R \to A$ 和 $A \to B$ 的相对对偶复形。
则 $K = K_{A/R} \otimes_A^\mathbf{L} K_{B/A}$ 是
$R \to B$ 的相对对偶复形。
\end{lemma}

\begin{proof}
我们将归约到 Noetherian 情形。具体而言，由代数，引理
\ref{algebra-lemma-flat-finite-presentation-limit-flat}，存在有限型
$\mathbf{Z}$-子代数 $R_0 \subset R$ 和平坦有限型环映射
$R_0 \to A_0$，使得 $A = A_0 \otimes_{R_0} R$。增大 $R_0$
并相应地替换 $A_0$ 后，可设存在平坦有限型环映射
$A_0 \to B_0$，使得 $B = B_0 \otimes_{R_0} R$（应用同一引理）。
若对 $R_0 \to A_0 \to B_0$ 证明本引理，则由引理
\ref{dualizing-lemma-uniqueness-relative-dualizing}、
\ref{dualizing-lemma-relative-dualizing-noetherian} 和
\ref{dualizing-lemma-base-change-relative-dualizing} 即得一般情形。
这把问题归约到下一段所讨论的情形。

\medskip\noindent
设 $R$ 为 Noetherian 环，并以 $\varphi : R \to A$ 和
$\psi : A \to B$ 表示给定的环映射。由上面所引的结果，
$K_{A/R} \cong \varphi^!(R)$ 且 $K_{B/A} \cong \psi^!(A)$。于是
$$
K = K_{A/R} \otimes_A^\mathbf{L} K_{B/A} \cong
\varphi^!(R) \otimes_A^\mathbf{L} \psi^!(A) \cong
\psi^!(\varphi^!(R)) \cong (\psi \circ \varphi)^!(R)
$$
由引理 \ref{dualizing-lemma-upper-shriek-is-tensor-functor} 和
\ref{dualizing-lemma-composition-shriek-algebraic}，$K$ 因而是
$R \to B$ 的相对对偶复形。
\end{proof}
% Generated mechanically from the frozen duality.tex target.
% Do not edit this generated file; edit the bound locale source instead.

% OK, start here.
%

\setcounter{chapter}{47}
\chapter{概形的对偶理论}



\phantomsection
\label{duality-section-phantom}


\section{引言}
\label{duality-section-introduction}

\noindent
本章研究概形态射的相对对偶理论以及概形上的对偶复形。
参考文献为 \cite{RD}。

\medskip\noindent
Noetherian 环的对偶复形已在《对偶化复形》第
\ref{dualizing-section-dualizing} 节及其后各节中定义并研究。
本章继续研究概形上的对偶复形；见第
\ref{duality-section-dualizing-schemes} 节。

\medskip\noindent
本章的主要篇幅用于研究如下情形中的推前函子之右伴随：
模层的导出范畴，且其上同调层为拟凝聚层。见第
\ref{duality-section-twisted-inverse-image}、
\ref{duality-section-restriction-to-opens}、
\ref{duality-section-base-change-map}、
\ref{duality-section-base-change-II}、
\ref{duality-section-trace}、
\ref{duality-section-compare-with-pullback}、
\ref{duality-section-sections-with-exact-support}、
\ref{duality-section-duality-finite}、
\ref{duality-section-perfect-proper}、
\ref{duality-section-dualizing-Cartier} 和
\ref{duality-section-examples} 节。
这里遵循论文
\cite{Neeman-Grothendieck}、\cite{LN}、
\cite{Lipman-notes} 和 \cite{Neeman-improvement}。

\medskip\noindent
第 \ref{duality-section-upper-shriek}、
\ref{duality-section-upper-shriek-properties} 和
\ref{duality-section-base-change-shriek} 节讨论 Noetherian 概形之间
有限型分离态射的重要而实用的上叹号函子 $f^!$，并最终在总览性的
第 \ref{duality-section-duality} 节汇总。

\medskip\noindent
第 \ref{duality-section-glue} 节说明另一种对偶理论与对偶复形理论：
在一个固定的、带有对偶复形的局部 Noetherian 基上工作
（该节对应 Hartshorne 书中的一条注记）。

\medskip\noindent
其余各节给出若干应用。

\medskip\noindent
下一章把本章内容延伸到代数空间上的对偶理论；见
《空间的对偶理论》第 \ref{spaces-duality-section-introduction} 节。







\section{概形上的对偶复形}
\label{duality-section-dualizing-schemes}

\noindent
局部 Noetherian 概形上的对偶复形被定义为这样一个复形：
仿射局部地，它来自相应环上的对偶复形。这个定义并非完全标准，
但在有限维 Noetherian 概形上与文献中的所有定义一致。

\begin{lemma}
\label{duality-lemma-equivalent-definitions}
设 $X$ 为局部 Noetherian 概形，$K$ 为 $D(\mathcal{O}_X)$ 的对象。
下列条件等价：
\begin{enumerate}
\item 对每个仿射开集 $U = \Spec(A) \subset X$，存在 $A$ 的对偶复形
$\omega_A^\bullet$，使得 $K|_U$ 同构于 $\omega_A^\bullet$ 在函子
$\widetilde{} : D(A) \to D(\mathcal{O}_U)$ 下的像。
\item 存在仿射开覆盖 $X = \bigcup U_i$，$U_i = \Spec(A_i)$，
使得对每个 $i$，均存在 $A_i$ 的对偶复形 $\omega_i^\bullet$，
且 $K|_{U_i}$ 同构于 $\omega_i^\bullet$ 在函子
$\widetilde{} : D(A_i) \to D(\mathcal{O}_{U_i})$ 下的像。
\end{enumerate}
\end{lemma}

\begin{proof}
设 (2) 成立，并令 $U = \Spec(A)$ 为 $X$ 的仿射开集。
条件 (2) 蕴含 $K$ 属于 $D_\QCoh(\mathcal{O}_X)$，因此在 $D(A)$
中存在对象 $\omega_A^\bullet$，其相伴拟凝聚层复形同构于 $K|_U$；
见《概形的导出范畴》，引理
\ref{perfect-lemma-affine-compare-bounded}。只需证明
$\omega_A^\bullet$ 是 $A$ 的对偶复形。

\medskip\noindent
由于 $X = \bigcup U_i$ 是开覆盖，可找到标准开覆盖
$U = D(f_1) \cup \ldots \cup D(f_m)$，使每个 $D(f_j)$
都是某个仿射开集 $U_i$ 中的标准开集；见《概形》，引理
\ref{schemes-lemma-standard-open-two-affines}。设
$D(f_j) = D(g_j)$，其中 $g_j \in A_{i_j}$。则
$A_{f_j} \cong (A_{i_j})_{g_j}$，并且由《概形的导出范畴》，引理
\ref{perfect-lemma-affine-compare-bounded}，在导出范畴中有
$$
(\omega_A^\bullet)_{f_j} \cong (\omega_i^\bullet)_{g_j}
$$
由《对偶化复形》，引理 \ref{dualizing-lemma-dualizing-localize}，
对 $j = 1, \ldots, m$，复形 $(\omega_A^\bullet)_{f_j}$
是 $A_{f_j}$ 上的对偶复形。于是由《对偶化复形》，引理
\ref{dualizing-lemma-dualizing-glue}，$\omega_A^\bullet$ 是对偶复形。
\end{proof}

\begin{definition}
\label{duality-definition-dualizing-scheme}
设 $X$ 为局部 Noetherian 概形。若 $D(\mathcal{O}_X)$ 的对象 $K$
满足引理 \ref{duality-lemma-equivalent-definitions} 中的等价条件，
则称 $K$ 为一个{\it 对偶复形}。
\end{definition}

\noindent
请参见本节开头的说明。

\begin{lemma}
\label{duality-lemma-affine-duality}
设 $A$ 为 Noetherian 环，$X = \Spec(A)$。设 $K, L$ 为 $D(A)$
的对象。若 $K \in D_{\textit{Coh}}(A)$ 且 $L$ 具有有限内射维数，
则在 $D(\mathcal{O}_X)$ 中有
$$
R\SheafHom_{\mathcal{O}_X}(\widetilde{K}, \widetilde{L})
=
\widetilde{R\Hom_A(K, L)}
$$
\end{lemma}

\begin{proof}
可设 $L$ 由内射 $A$-模组成的有限复形 $I^\bullet$ 给出。
对 $I^\bullet$ 的长度作归纳，并利用各构造与可区别三角的相容性，
归约到 $L = I[0]$ 的情形，其中 $I$ 为内射 $A$-模。在此情形下，
《概形的导出范畴》，引理
\ref{perfect-lemma-quasi-coherence-internal-hom} 表明，
$R\SheafHom_{\mathcal{O}_X}(\widetilde{K}, \widetilde{L})$
的第 $n$ 个上同调层是与预层
$$
D(f) \longmapsto \Ext^n_{A_f}(K \otimes_A A_f, I \otimes_A A_f)
$$
相伴的层。由于 $A$ 为 Noetherian 环，$A_f$-模
$I \otimes_A A_f$ 是内射的（《对偶化复形》，引理
\ref{dualizing-lemma-localization-injective-modules}）。因此
\begin{align*}
\Ext^n_{A_f}(K \otimes_A A_f, I \otimes_A A_f)
& =
\Hom_{A_f}(H^{-n}(K \otimes_A A_f), I \otimes_A A_f) \\
& =
\Hom_{A_f}(H^{-n}(K) \otimes_A A_f, I \otimes_A A_f) \\
& =
\Hom_A(H^{-n}(K), I) \otimes_A A_f
\end{align*}
最后一个等式成立，是因为 $H^{-n}(K)$ 是有限 $A$-模；见代数，引理
\ref{algebra-lemma-hom-from-finitely-presented}。这证明了典范映射
$$
\widetilde{R\Hom_A(K, L)}
\longrightarrow
R\SheafHom_{\mathcal{O}_X}(\widetilde{K}, \widetilde{L})
$$
在此情形下为拟同构，证明完成。
\end{proof}

\begin{lemma}
\label{duality-lemma-internal-hom-evaluate-isom}
设 $X$ 为 Noetherian 概形，$K, L, M \in D_\QCoh(\mathcal{O}_X)$。
则《上同调》，引理 \ref{cohomology-lemma-internal-hom-evaluate} 中的映射
$$
R\SheafHom(L, M) \otimes_{\mathcal{O}_X}^\mathbf{L} K
\longrightarrow
R\SheafHom(R\SheafHom(K, L), M)
$$
在下列两种情形中是同构：
\begin{enumerate}
\item $K \in D^-_{\textit{Coh}}(\mathcal{O}_X)$，
$L \in D^+_{\textit{Coh}}(\mathcal{O}_X)$，并且 $M$ 仿射局部地具有
有限内射维数（见证明）；或者
\item $K$ 和 $L$ 属于 $D_{\textit{Coh}}(\mathcal{O}_X)$，
对象 $R\SheafHom(L, M)$ 具有有限 Tor 维数，并且
$L$ 和 $M$ 仿射局部地具有有限内射维数
（特别地，$L$ 和 $M$ 是有界的）。
\end{enumerate}
\end{lemma}

\begin{proof}
先证 (1)。所谓 $M$ 仿射局部地具有有限内射维数，是指 $X$ 有一个
由仿射开集 $U = \Spec(A)$ 构成的开覆盖，使得与 $M|_U$ 对应的
$D(A)$ 中的对象（见《概形的导出范畴》，引理
\ref{perfect-lemma-affine-compare-bounded}）具有有限内射维数
\footnote{此条件不依赖于 Noetherian 概形 $X$ 的仿射开覆盖的选取。
细节从略。}。为证明本引理，可以把 $X$ 替换为 $U$，即可以假设
对某个 Noetherian 环 $A$ 有 $X = \Spec(A)$。注意，由《概形的导出范畴》，
引理 \ref{perfect-lemma-coherent-internal-hom}，
$R\SheafHom(K, L)$ 属于 $D^+_{\textit{Coh}}(\mathcal{O}_X)$。
此外，由引理 \ref{duality-lemma-affine-duality} 及《概形的导出范畴》，引理
\ref{perfect-lemma-quasi-coherence-internal-hom}，箭头左右两端的构造均与函子
$D(A) \to D_\QCoh(\mathcal{O}_X)$ 交换
（严格说来，这里用到了对 $K$、$L$、$M$ 的假设，以及刚刚对
$R\SheafHom(K, L)$ 证明的结论）。最后，由《代数进阶》，引理
\ref{more-algebra-lemma-internal-hom-evaluate-isomorphism} 的第 (2) 部分，
该箭头是同构。

\medskip\noindent
再证 (2)。论证同上。一个小差别是，此时
$R\SheafHom(K, L)$ 属于 $D_{\textit{Coh}}(\mathcal{O}_X)$，因为可以
仿射局部地考察（由引理 \ref{duality-lemma-affine-duality}，这是允许的），
并应用《对偶化复形》，引理
\ref{dualizing-lemma-finite-ext-into-bounded-injective}。
最后由《代数进阶》，引理
\ref{more-algebra-lemma-internal-hom-evaluate-isomorphism-technical} 得出结论。
\end{proof}

\begin{lemma}
\label{duality-lemma-dualizing-schemes}
设 $K$ 是局部 Noetherian 概形 $X$ 上的对偶复形。则 $K$ 是
$D_{\textit{Coh}}(\mathcal{O}_X)$ 的对象，并且
$D = R\SheafHom_{\mathcal{O}_X}(-, K)$ 诱导反等价
$$
D :
D_{\textit{Coh}}(\mathcal{O}_X)
\longrightarrow
D_{\textit{Coh}}(\mathcal{O}_X)
$$
它带有一个典范同构 $\text{id} \to D \circ D$。若 $X$ 拟紧，
则 $D$ 交换 $D^+_{\textit{Coh}}(\mathcal{O}_X)$ 与
$D^-_{\textit{Coh}}(\mathcal{O}_X)$，并诱导反等价
$D^b_{\textit{Coh}}(\mathcal{O}_X) \to D^b_{\textit{Coh}}(\mathcal{O}_X)$。
\end{lemma}

\begin{proof}
设 $U \subset X$ 为仿射开集，记 $U = \Spec(A)$，并设
$\omega_A^\bullet$ 为与 $K|_U$ 对应的 $A$ 的对偶复形，
如引理 \ref{duality-lemma-equivalent-definitions} 中所述。由引理
\ref{duality-lemma-affine-duality}，图表
$$
\xymatrix{
D_{\textit{Coh}}(A) \ar[r] \ar[d]_{R\Hom_A(-, \omega_A^\bullet)} &
D_{\textit{Coh}}(\mathcal{O}_U) \ar[d]^{R\SheafHom_{\mathcal{O}_X}(-, K|_U)} \\
D_{\textit{Coh}}(A) \ar[r] &
D(\mathcal{O}_U)
}
$$
交换。因此 $D$ 将 $D_{\textit{Coh}}(\mathcal{O}_X)$ 映入
$D_{\textit{Coh}}(\mathcal{O}_X)$。此外，对所有 $L$，典范映射
$$
L
\longrightarrow
R\SheafHom_{\mathcal{O}_X}(K, K) \otimes_{\mathcal{O}_X}^\mathbf{L} L
\longrightarrow
R\SheafHom_{\mathcal{O}_X}(R\SheafHom_{\mathcal{O}_X}(L, K), K)
$$
都是同构，其中第二个箭头使用了《上同调》，引理
\ref{cohomology-lemma-internal-hom-evaluate}；这是因为由《对偶化复形》，
引理 \ref{dualizing-lemma-dualizing}，该结论在仿射概形上成立
\footnote{另一种方法是先仿射局部地证明
$R\SheafHom_{\mathcal{O}_X}(K, K) = \mathcal{O}_X$，再用引理
\ref{duality-lemma-internal-hom-evaluate-isom} 的第 (2) 部分说明该映射为同构。}，
而且我们已经看到，在仿射概形上这恢复了代数中的情形。
拟紧情形下函子 $D$ 的有界性陈述，同样由《对偶化复形》，引理
\ref{dualizing-lemma-dualizing} 的相应陈述得到。
\end{proof}

\noindent
设 $X$ 为局部赋环空间。回忆：若 $D(\mathcal{O}_X)$ 的对象 $L$
关于由导出张量积给出的 $D(\mathcal{O}_X)$ 上的对称幺半结构是可逆对象，
则称该对象是{\it 可逆的}。我们在《上同调》，引理
\ref{cohomology-lemma-invertible-derived} 中已经看到，这意味着 $L$ 是完美的，
且存在开覆盖 $X = \bigcup U_i$，使得对某些整数 $n_i$ 有
$L|_{U_i} \cong \mathcal{O}_{U_i}[-n_i]$。在此情形下，函数
$$
x \mapsto n_x,\quad
\text{其中 }n_x\text{ is the unique integer such that }
H^{n_x}(L_x) \not = 0
$$
在 $X$ 上局部常值。特别地，有 $L = \bigoplus H^n(L)[-n]$；
这给出了一个表示 $L$ 的良定义的 $\mathcal{O}_X$-模复形
（其微分全为零）。

\begin{lemma}
\label{duality-lemma-dualizing-unique-schemes}
设 $X$ 为局部 Noetherian 概形。若 $K$ 与 $K'$ 是 $X$ 上的对偶复形，
则对 $D(\mathcal{O}_X)$ 的某个可逆对象 $L$，$K'$ 同构于
$K \otimes_{\mathcal{O}_X}^\mathbf{L} L$。
\end{lemma}

\begin{proof}
令
$$
L = R\SheafHom_{\mathcal{O}_X}(K, K')
$$
这是 $D(\mathcal{O}_X)$ 的可逆对象，因为此结论仿射局部地成立；
见《对偶化复形》，引理 \ref{dualizing-lemma-dualizing-unique} 及其证明。
同理，求值映射 $L \otimes_{\mathcal{O}_X}^\mathbf{L} K \to K'$
是同构。
\end{proof}

\begin{lemma}
\label{duality-lemma-dimension-function-scheme}
设 $X$ 为局部 Noetherian 概形，$\omega_X^\bullet$ 为 $X$ 上的对偶复形。
则 $X$ 是普遍等长链的，并且由
$X \to \mathbf{Z}$ 定义的函数
$$
x \longmapsto \delta(x)\text{ 使得 }
\omega_{X, x}^\bullet[-\delta(x)]
\text{ is a normalized dualizing complex over }
\mathcal{O}_{X, x}
$$
是维数函数。
\end{lemma}

\begin{proof}
由仿射情形的《对偶化复形》，引理
\ref{dualizing-lemma-dimension-function} 及各定义立即得到。
\end{proof}

\begin{lemma}
\label{duality-lemma-sitting-in-degrees}
设 $X$ 为局部 Noetherian 概形，$\omega_X^\bullet$ 为 $X$ 上的对偶复形，
其相伴维数函数为 $\delta$。设 $\mathcal{F}$ 为凝聚 $\mathcal{O}_X$-模。令
$\mathcal{E}^i = \SheafExt^{-i}_{\mathcal{O}_X}(\mathcal{F}, \omega_X^\bullet)$。
则 $\mathcal{E}^i$ 是凝聚 $\mathcal{O}_X$-模，并且对 $x \in X$ 有：
\begin{enumerate}
\item $\mathcal{E}^i_x$ 仅当
$\delta(x) \leq i \leq \delta(x) + \dim(\text{Supp}(\mathcal{F}_x))$ 时可能非零；
\item $\dim(\text{Supp}(\mathcal{E}^{i + \delta(x)}_x)) \leq i$；
\item $\text{depth}(\mathcal{F}_x)$ 是满足
$\mathcal{E}_x^{i + \delta(x)} \not = 0$ 的最小整数 $i \geq 0$；
\item 有
$x \in \text{Supp}(\bigoplus_{j \leq i} \mathcal{E}^j)
\Leftrightarrow
\text{depth}_{\mathcal{O}_{X, x}}(\mathcal{F}_x) + \delta(x) \leq i$。
\end{enumerate}
\end{lemma}

\begin{proof}
由引理 \ref{duality-lemma-dualizing-schemes}，$\mathcal{E}^i$ 是凝聚的。
选取 $x$ 的一个仿射邻域，并应用《概形的导出范畴》，引理
\ref{perfect-lemma-quasi-coherence-internal-hom} 以及《代数进阶》，引理
\ref{more-algebra-lemma-base-change-RHom} 的第 (3) 部分，可得
$$
\mathcal{E}^i_x =
\SheafExt^{-i}_{\mathcal{O}_X}(\mathcal{F}, \omega_X^\bullet)_x =
\Ext^{-i}_{\mathcal{O}_{X, x}}(\mathcal{F}_x,
\omega_{X, x}^\bullet) =
\Ext^{\delta(x) - i}_{\mathcal{O}_{X, x}}(\mathcal{F}_x,
\omega_{X, x}^\bullet[-\delta(x)])
$$
根据引理 \ref{duality-lemma-dimension-function-scheme} 中 $\delta$ 的构造，
第 (1)、(2)、(3) 项便归结为《对偶化复形》，引理
\ref{dualizing-lemma-sitting-in-degrees}。第 (4) 项是第 (3) 项和第 (1) 项的形式推论。
\end{proof}




\section{推前的右伴随}
\label{duality-section-twisted-inverse-image}

\noindent
本节及下一节的参考文献为
\cite{Neeman-Grothendieck}、\cite{LN}、
\cite{Lipman-notes} 和 \cite{Neeman-improvement}。

\medskip\noindent
设 $f : X \to Y$ 为概形态射。本节考察函子
$Rf_* : D_\QCoh(\mathcal{O}_X) \to D_\QCoh(\mathcal{O}_Y)$ 的右伴随。
文献中，当这个函子存在时，有时将其记作 $f^{\times}$。
这种记号并未得到普遍接受，所以这里不采用。我们也不把这样的函子记作
$f^!$，因为对一般态射 $f$，这会与 \cite{RD} 中的记号冲突。

\begin{lemma}
\label{duality-lemma-twisted-inverse-image}
\begin{reference}
这与 \cite[例 4.2]{Neeman-Grothendieck} 几乎相同。
\end{reference}
设 $f : X \to Y$ 为拟分离拟紧概形之间的态射。函子
$Rf_* : D_\QCoh(\mathcal{O}_X) \to D_\QCoh(\mathcal{O}_Y)$ 有右伴随。
\end{lemma}

\begin{proof}
验证《导出范畴》，命题 \ref{derived-proposition-brown} 的假设，即可证明右伴随存在。
首先，范畴 $D_\QCoh(\mathcal{O}_X)$ 有直和；见《概形的导出范畴》，引理
\ref{perfect-lemma-quasi-coherence-direct-sums}。
由《概形的导出范畴》，定理
\ref{perfect-theorem-bondal-van-den-Bergh}，范畴
$D_\QCoh(\mathcal{O}_X)$ 是紧生成的。由于 $X$ 和 $Y$ 均拟紧且拟分离，
$f$ 也如此；见《概形》，引理 \ref{schemes-lemma-compose-after-separated} 和
\ref{schemes-lemma-quasi-compact-permanence}。因此函子 $Rf_*$ 与直和交换；
见《概形的导出范畴》，引理
\ref{perfect-lemma-quasi-coherence-pushforward-direct-sums}。证毕。
\end{proof}

\begin{example}
\label{duality-example-affine-twisted-inverse-image}
设 $A \to B$ 为环同态，$Y = \Spec(A)$、$X = \Spec(B)$，
并设 $f : X \to Y$ 为与 $A \to B$ 对应的态射。通过等价
$D(B) \to D_\QCoh(\mathcal{O}_X)$ 和
$D(A) \to D_\QCoh(\mathcal{O}_Y)$，函子
$Rf_* : D_\QCoh(\mathcal{O}_X) \to D_\QCoh(\mathcal{O}_Y)$
对应于限制函子 $D(B) \to D(A)$。因此，它的右伴随对应于
《对偶化复形》，第 \ref{dualizing-section-trivial} 节中的函子
$K \longmapsto R\Hom(B, K)$。
\end{example}

\begin{example}
\label{duality-example-does-not-preserve-coherent}
若 $f : X \to Y$ 是 Noetherian 概形之间的分离有限型态射，则
$Rf_* : D_\QCoh(\mathcal{O}_X) \to D_\QCoh(\mathcal{O}_Y)$ 的右伴随
并不把 $D_{\textit{Coh}}(\mathcal{O}_Y)$ 映入
$D_{\textit{Coh}}(\mathcal{O}_X)$。事实上，设 $k$ 为域，并考虑态射
$f : \mathbf{A}^1_k \to \Spec(k)$。由例
\ref{duality-example-affine-twisted-inverse-image}，这归结为下述问题：当
$A = k$ 且 $B = k[x]$ 时，$R\Hom(B, -)$ 是否把
$D_{\textit{Coh}}(A)$ 映入 $D_{\textit{Coh}}(B)$。答案是否定的，因为
$$
R\Hom(k[x], k) = \left(\prod\nolimits_{n \geq 0} k\right)[0]
$$
并不是有限 $k[x]$-模。因此 $a(\mathcal{O}_Y)$ 没有凝聚上同调层。
\end{example}

\begin{example}
\label{duality-example-does-not-preserve-bounded-above}
若 $f : X \to Y$ 是 Noetherian 概形之间的固有态射，甚至是有限态射，则
$Rf_* : D_\QCoh(\mathcal{O}_X) \to D_\QCoh(\mathcal{O}_Y)$ 的右伴随
并不把 $D_\QCoh^-(\mathcal{O}_Y)$ 映入
$D_\QCoh^-(\mathcal{O}_X)$。事实上，设 $k$ 为域，$k[\epsilon]$
为 $k$ 上的对偶数，令 $X = \Spec(k)$、$Y = \Spec(k[\epsilon])$。
则对所有 $i \geq 0$，$\Ext^i_{k[\epsilon]}(k, k)$ 均非零。
因此，由例 \ref{duality-example-affine-twisted-inverse-image}，
$a(\mathcal{O}_Y)$ 不是上有界的。
\end{example}

\begin{lemma}
\label{duality-lemma-twisted-inverse-image-bounded-below}
设 $f : X \to Y$ 为拟紧拟分离概形之间的态射。设
$a : D_\QCoh(\mathcal{O}_Y) \to D_\QCoh(\mathcal{O}_X)$ 为引理
\ref{duality-lemma-twisted-inverse-image} 中 $Rf_*$ 的右伴随。则 $a$ 把
$D^+_\QCoh(\mathcal{O}_Y)$ 映入 $D^+_\QCoh(\mathcal{O}_X)$。
事实上，存在整数 $N$，使得 $i \leq c$ 时 $H^i(K) = 0$ 蕴含
$i \leq c - N$ 时 $H^i(a(K)) = 0$。
\end{lemma}

\begin{proof}
由《概形的导出范畴》，引理
\ref{perfect-lemma-quasi-coherence-direct-image}，函子 $Rf_*$
具有有限上同调维数。换言之，存在整数 $N$，使得若 $i \geq c$ 时
$H^i(L) = 0$，则 $i \geq N + c$ 时 $H^i(Rf_*L) = 0$。
设 $K \in D^+_\QCoh(\mathcal{O}_Y)$ 且 $i \leq c$ 时 $H^i(K) = 0$。则
$$
\Hom_{D(\mathcal{O}_X)}(\tau_{\leq c - N}a(K), a(K)) =
\Hom_{D(\mathcal{O}_Y)}(Rf_*\tau_{\leq c - N}a(K), K) = 0
$$
由上述结论，显然这蕴含 $i \leq c - N$ 时 $H^i(a(K)) = 0$。
\end{proof}

\noindent
设 $f : X \to Y$ 为拟分离拟紧概形之间的态射，并以 $a$ 表示
$Rf_* : D_\QCoh(\mathcal{O}_X) \to D_\QCoh(\mathcal{O}_Y)$ 的右伴随。
对每个 $K \in D_\QCoh(\mathcal{O}_Y)$ 和
$L \in D_\QCoh(\mathcal{O}_X)$，得到典范映射
\begin{equation}
\label{duality-equation-sheafy-trace}
Rf_*R\SheafHom_{\mathcal{O}_X}(L, a(K))
\longrightarrow
R\SheafHom_{\mathcal{O}_Y}(Rf_*L, K)
\end{equation}
具体地说，此映射构造为复合
$$
Rf_*R\SheafHom_{\mathcal{O}_X}(L, a(K)) \to
R\SheafHom_{\mathcal{O}_Y}(Rf_*L, Rf_*a(K)) \to
R\SheafHom_{\mathcal{O}_Y}(Rf_*L, K)
$$
其中第一个箭头见《上同调》，注
\ref{cohomology-remark-projection-formula-for-internal-hom}，
第二个箭头则是伴随的余单位 $Rf_*a(K) \to K$。

\begin{lemma}
\label{duality-lemma-iso-on-RSheafHom}
设 $f : X \to Y$ 为拟紧拟分离概形之间的态射，$a$ 为
$Rf_* : D_\QCoh(\mathcal{O}_X) \to D_\QCoh(\mathcal{O}_Y)$ 的右伴随。
设 $L \in D_\QCoh(\mathcal{O}_X)$ 且 $K \in D_\QCoh(\mathcal{O}_Y)$。
对映射 (\ref{duality-equation-sheafy-trace})
$$
Rf_*R\SheafHom_{\mathcal{O}_X}(L, a(K))
\longrightarrow
R\SheafHom_{\mathcal{O}_Y}(Rf_*L, K)
$$
应用《概形的导出范畴》，第 \ref{perfect-section-better-coherator} 节中讨论的函子
$DQ_Y : D(\mathcal{O}_Y) \to D_\QCoh(\mathcal{O}_Y)$ 后，所得映射为同构。
\end{lemma}

\begin{proof}
由《概形的导出范畴》，引理 \ref{perfect-lemma-better-coherator}，
$DQ_Y$ 存在，故此陈述有意义。由于 $DQ_Y$ 是嵌入函子
$D_\QCoh(\mathcal{O}_Y) \to D(\mathcal{O}_Y)$ 的右伴随，
为证明本引理，只需说明对任意 $M \in D_\QCoh(\mathcal{O}_Y)$，
映射 (\ref{duality-equation-sheafy-trace}) 诱导双射
$$
\Hom_Y(M, Rf_*R\SheafHom_{\mathcal{O}_X}(L, a(K)))
\longrightarrow
\Hom_Y(M, R\SheafHom_{\mathcal{O}_Y}(Rf_*L, K))
$$
为此，使用如下等式链：
\begin{align*}
\Hom_Y(M, Rf_*R\SheafHom_{\mathcal{O}_X}(L, a(K)))
& =
\Hom_X(Lf^*M, R\SheafHom_{\mathcal{O}_X}(L, a(K))) \\
& =
\Hom_X(Lf^*M \otimes_{\mathcal{O}_X}^\mathbf{L} L, a(K)) \\
& =
\Hom_Y(Rf_*(Lf^*M \otimes_{\mathcal{O}_X}^\mathbf{L} L), K) \\
& =
\Hom_Y(M \otimes_{\mathcal{O}_Y}^\mathbf{L} Rf_*L, K) \\
& =
\Hom_Y(M, R\SheafHom_{\mathcal{O}_Y}(Rf_*L, K))
\end{align*}
第一个等式由《上同调》，引理 \ref{cohomology-lemma-adjoint}；
第二个等式由《上同调》，引理 \ref{cohomology-lemma-internal-hom}；
第三个等式由 $a$ 的构造；第四个等式由《概形的导出范畴》，引理
\ref{perfect-lemma-cohomology-base-change}（这是关键一步）；
第五个等式由《上同调》，引理 \ref{cohomology-lemma-internal-hom}。
\end{proof}

\begin{example}
\label{duality-example-iso-on-RSheafHom}
如果不应用“拟凝聚化函子” $DQ_Y$，引理 \ref{duality-lemma-iso-on-RSheafHom}
的陈述并不成立。事实上，设 $Y = \Spec(R)$ 且 $X = \mathbf{A}^1_R$，
取 $L = \mathcal{O}_X$ 和 $K = \mathcal{O}_Y$。箭头左端属于
$D_\QCoh(\mathcal{O}_Y)$，但箭头右端同构于
$\prod_{n \geq 0} \mathcal{O}_Y$，后者不是拟凝聚的。
\end{example}

\begin{remark}
\label{duality-remark-iso-on-RSheafHom}
在引理 \ref{duality-lemma-iso-on-RSheafHom} 的情形下，由《概形的导出范畴》，引理
\ref{perfect-lemma-pushforward-better-coherator}，有
$$
DQ_Y(Rf_*R\SheafHom_{\mathcal{O}_X}(L, a(K))) =
Rf_* DQ_X(R\SheafHom_{\mathcal{O}_X}(L, a(K)))
$$
因此，如果 $R\SheafHom_{\mathcal{O}_X}(L, a(K)) \in D_\QCoh(\mathcal{O}_X)$，
便可从箭头左端“删去” $DQ_Y$。另一方面，如果已知
$R\SheafHom_{\mathcal{O}_Y}(Rf_*L, K) \in D_\QCoh(\mathcal{O}_Y)$，
便可从箭头右端“删去” $DQ_Y$。如果两者都成立，就得到
(\ref{duality-equation-sheafy-trace}) 是同构。结合《概形的导出范畴》，引理
\ref{perfect-lemma-quasi-coherence-internal-hom} 可知，在下列情形中
$Rf_*R\SheafHom_{\mathcal{O}_X}(L, a(K)) \to
R\SheafHom_{\mathcal{O}_Y}(Rf_*L, K)$ 是同构：
\begin{enumerate}
\item $L$ 和 $Rf_*L$ 都是完美的；或者
\item $K$ 下有界，且 $L$ 和 $Rf_*L$ 都是伪凝聚的。
\end{enumerate}
对 (2)，用到：若 $K$ 下有界，则 $a(K)$ 也下有界；见引理
\ref{duality-lemma-twisted-inverse-image-bounded-below}。
\end{remark}

\begin{example}
\label{duality-example-iso-on-RSheafHom-noetherian}
设 $f : X \to Y$ 为 Noetherian 概形之间的固有态射，
$L \in D^-_{\textit{Coh}}(X)$ 且 $K \in D^+_{\QCoh}(\mathcal{O}_Y)$。
则映射 $Rf_*R\SheafHom_{\mathcal{O}_X}(L, a(K)) \to
R\SheafHom_{\mathcal{O}_Y}(Rf_*L, K)$ 是同构。事实上，由
《概形的导出范畴》，引理
\ref{perfect-lemma-identify-pseudo-coherent-noetherian} 和
\ref{perfect-lemma-direct-image-coherent}，复形 $L$ 与 $Rf_*L$ 均为伪凝聚的，
故注 \ref{duality-remark-iso-on-RSheafHom} 中的讨论适用。
\end{example}

\begin{lemma}
\label{duality-lemma-iso-global-hom}
设 $f : X \to Y$ 为拟分离拟紧概形之间的态射。对所有
$L \in D_\QCoh(\mathcal{O}_X)$ 和 $K \in D_\QCoh(\mathcal{O}_Y)$，
(\ref{duality-equation-sheafy-trace}) 诱导全局导出 Hom 的同构
$R\Hom_X(L, a(K)) \to R\Hom_Y(Rf_*L, K)$。
\end{lemma}

\begin{proof}
由《上同调》，第 \ref{cohomology-section-global-RHom} 节中的构造，有
$$
R\Hom_X(L, a(K)) =
R\Gamma(X, R\SheafHom_{\mathcal{O}_X}(L, a(K))) =
R\Gamma(Y, Rf_*R\SheafHom_{\mathcal{O}_X}(L, a(K)))
$$
以及
$$
R\Hom_Y(Rf_*L, K) = R\Gamma(Y, R\SheafHom_{\mathcal{O}_Y}(Rf_*L, K))
$$
因此本引理是引理 \ref{duality-lemma-iso-on-RSheafHom} 的推论。具体而言，
$D(\mathcal{O}_Y)$ 中的映射 $E \to E'$ 若诱导同构
$DQ_Y(E) \to DQ_Y(E')$，便诱导拟同构
$R\Gamma(Y, E) \to R\Gamma(Y, E')$。事实上，
$H^i(Y, E) = \Ext^i_Y(\mathcal{O}_Y, E) = \Hom(\mathcal{O}_Y[-i], E) =
\Hom(\mathcal{O}_Y[-i], DQ_Y(E))$，因为 $\mathcal{O}_Y[-i]$ 属于
$D_\QCoh(\mathcal{O}_Y)$，而 $DQ_Y$ 是嵌入函子
$D_\QCoh(\mathcal{O}_Y) \to D(\mathcal{O}_Y)$ 的右伴随。
\end{proof}







\section{推前的右伴随与开集限制}
\label{duality-section-restriction-to-opens}

\noindent
本节研究推前的右伴随在多大程度上与到开子概形的限制交换。
这是一个基变换问题，所以先在较一般的情形下讨论。

\medskip\noindent
我们经常需要判断推前的右伴随是否与基变换交换。为此，考虑拟紧拟分离概形的笛卡尔方块
\begin{equation}
\label{duality-equation-base-change}
\vcenter{
\xymatrix{
X' \ar[r]_{g'} \ar[d]_{f'} & X \ar[d]^f \\
Y' \ar[r]^g & Y
}
}
\end{equation}
以
$$
a  : D_\QCoh(\mathcal{O}_Y) \to D_\QCoh(\mathcal{O}_X)
\quad\text{且}\quad
a'  : D_\QCoh(\mathcal{O}_{Y'}) \to D_\QCoh(\mathcal{O}_{X'})
$$
表示 $Rf_*$ 与 $Rf'_*$ 的右伴随（引理
\ref{duality-lemma-twisted-inverse-image}）。考虑《上同调》，注
\ref{cohomology-remark-base-change} 中的基变换映射。它在具有拟凝聚上同调层的层的导出范畴上诱导函子变换
$$
Lg^* \circ Rf_* \longrightarrow Rf'_* \circ L(g')^*
$$
从而在相反方向诱导右伴随之间的变换
$$
a \circ Rg_* \longleftarrow Rg'_* \circ a'
$$

\begin{lemma}
\label{duality-lemma-flat-precompose-pus}
在图 (\ref{duality-equation-base-change}) 中，假设 $g$ 平坦；更一般地，假设
$f$ 与 $g$ Tor 独立。则 $a \circ Rg_* \leftarrow Rg'_* \circ a'$ 是同构。
\end{lemma}

\begin{proof}
在此情形下，由《概形的导出范畴》，引理
\ref{perfect-lemma-compare-base-change}，对 $D_\QCoh(\mathcal{O}_X)$ 中每个
$K$，基变换映射
$Lg^* \circ Rf_* K \longrightarrow Rf'_* \circ L(g')^*K$ 都是同构。
因此，相应的伴随函子之间的变换也是同构。
\end{proof}

\noindent
设 $f : X \to Y$ 为拟紧拟分离概形之间的态射，$V \subset Y$
为拟紧开子概形，并令 $U = f^{-1}(V)$。这给出如
(\ref{duality-equation-base-change}) 中的笛卡尔方块
$$
\xymatrix{
U \ar[r]_{j'} \ar[d]_{f|_U} & X \ar[d]^f \\
V \ar[r]^j & Y
}
$$
由引理 \ref{duality-lemma-flat-precompose-pus}，映射
$\xi : a \circ Rj_* \leftarrow Rj'_* \circ a'$ 是同构，
其中 $a$ 和 $a'$ 分别是 $Rf_*$ 与 $R(f|_U)_*$ 的右伴随。
于是得到函子 $D_\QCoh(\mathcal{O}_Y) \to D_\QCoh(\mathcal{O}_U)$ 的变换
\begin{equation}
\label{duality-equation-sheafy}
(j')^* \circ a \to
(j')^* \circ a \circ Rj_* \circ j^* \xrightarrow{\xi^{-1}}
(j')^* \circ Rj'_* \circ a' \circ j^* \to a' \circ j^*
\end{equation}
其中第一个箭头来自 $\text{id} \to Rj_* \circ j^*$，最后一个箭头来自同构
$(j')^* \circ Rj'_* \to \text{id}$。特别地，(\ref{duality-equation-sheafy})
在 $K$ 上取值时为同构，当且仅当
$a(K)|_U \to a(Rj_*(K|_V))|_U$ 为同构。

\begin{example}
\label{duality-example-not-supported-on-inverse-image}
存在 Noetherian 概形的有限态射 $f : X \to Y$，使得
(\ref{duality-equation-sheafy}) 在某个
$K \in D_{\textit{Coh}}(\mathcal{O}_Y)$ 上取值时不是同构。
具体地，设 $X = \Spec(B) \to Y = \Spec(A)$，其中
$A = k[x, \epsilon]$，$k$ 是域且 $\epsilon^2 = 0$，并且
$B = k[x] = A/(\epsilon)$。对 $n \in \mathbf{N}$，令
$M_n = A/(\epsilon, x^n)$。注意
$$
\Ext^i_A(B, M_n) = M_n,\quad i \geq 0
$$
因为 $B$ 有自由周期分解 $\ldots \to A \to A \to A$，
其中各映射均由乘以 $\epsilon$ 给出。考虑
$D_{\textit{Coh}}(A)$ 的对象 $K = \bigoplus M_n[n] = \prod M_n[n]$；
这里在 $D(A)$ 中相等，由《导出范畴》，引理
\ref{derived-lemma-direct-sums} 和 \ref{derived-lemma-products}。
由例 \ref{duality-example-affine-twisted-inverse-image}，$a(K)$ 对应于
$R\Hom(B, K)$，而由上述结论
$$
H^0(R\Hom(B, K)) = \Ext^0_A(B, K) =
\prod\nolimits_{n \geq 1} \Ext^n_A(B. M_n) = 
\prod\nolimits_{n \geq 1} M_n
$$
然而，这个模含有不被 $x$ 的任何幂零化的元素，而复形 $K$
的上同调中的每个元素都被 $x$ 的某个幂零化。换言之，取
$V = D(x)$、$U = D(x)$ 和复形 $K$ 时，映射
(\ref{duality-equation-sheafy}) 不可能是同构，因为 $(j')^*(a(K))$ 非零，
而 $a'(j^*K)$ 为零。
\end{example}

\begin{lemma}
\label{duality-lemma-when-sheafy}
设 $f : X \to Y$ 为拟紧拟分离概形之间的态射，$a$ 为
$Rf_* : D_\QCoh(\mathcal{O}_X) \to D_\QCoh(\mathcal{O}_Y)$ 的右伴随。
设 $V \subset Y$ 为拟紧开集，其逆像为 $U \subset X$。
\begin{enumerate}
\item 对每个支撑在 $Y \setminus V$ 上的
$Q \in D_\QCoh^+(\mathcal{O}_Y)$，像 $a(Q)$ 支撑在
$X \setminus U$ 上，当且仅当 (\ref{duality-equation-sheafy}) 在
$D_\QCoh^+(\mathcal{O}_Y)$ 的所有 $K$ 上都是同构。
\item 对每个支撑在 $Y \setminus V$ 上的
$Q \in D_\QCoh(\mathcal{O}_Y)$，像 $a(Q)$ 支撑在
$X \setminus U$ 上，当且仅当 (\ref{duality-equation-sheafy}) 在
$D_\QCoh(\mathcal{O}_Y)$ 的所有 $K$ 上都是同构。
\item 若 $a$ 与直和交换，则 (1) 的等价条件蕴含 (2) 的等价条件。
\end{enumerate}
\end{lemma}

\begin{proof}
先证 (1)。设 $K \in D_\QCoh^+(\mathcal{O}_Y)$，选取可区别三角
$$
K \to Rj_*K|_V \to Q \to K[1]
$$
注意 $Q$ 属于 $D_\QCoh^+(\mathcal{O}_Y)$（《概形的导出范畴》，引理
\ref{perfect-lemma-quasi-coherence-direct-image}），且支撑在
$Y \setminus V$ 上（《概形的导出范畴》，定义
\ref{perfect-definition-supported-on}）。应用 $a$，在 $X$ 上得到可区别三角
$$
a(K) \to a(Rj_*K|_V) \to a(Q) \to a(K)[1]
$$
若 $a(Q)$ 支撑在 $X \setminus U$ 上，则限制到 $U$ 后，映射
$a(K)|_U \to a(Rj_*K|_V)|_U$ 是同构，即 (\ref{duality-equation-sheafy})
在 $K$ 上是同构。逆命题立即成立。

\medskip\noindent
(2) 的证明与 (1) 的证明完全相同。

\medskip\noindent
再证 (3)。假设 (1) 的等价条件成立，令 $T = Y \setminus V$。
用记号 $D_{\QCoh, T}(\mathcal{O}_Y)$ 和
$D_{\QCoh, f^{-1}(T)}(\mathcal{O}_X)$ 分别表示上同调层支撑在
$T$ 和 $f^{-1}(T)$ 上的复形。由于 $a$ 与直和交换，以
$$
\{Q \in D_{\QCoh, T}(\mathcal{O}_Y) \mid
a(Q) \in D_{\QCoh, f^{-1}(T)}(\mathcal{O}_X)\}
$$
为对象的严格满、饱和三角子范畴 $\mathcal{D}$ 在直和下封闭，
因而在导出余极限下也封闭。另一方面，范畴
$D_{\QCoh, T}(\mathcal{O}_Y)$ 由一个完美对象 $E$ 生成；
见《概形的导出范畴》，引理
\ref{perfect-lemma-generator-with-support}。由假设，$E \in \mathcal{D}$。
由《导出范畴》，引理 \ref{derived-lemma-write-as-colimit}，
$D_{\QCoh, T}(\mathcal{O}_Y)$ 的每个对象 $Q$ 都是某个系统
$Q_1 \to Q_2 \to Q_3 \to \ldots$ 的导出余极限，
且各转移映射的锥都是 $E$ 的平移的直和。归纳可知，对所有 $n$，
$Q_n \in \mathcal{D}$，最终 $Q$ 也属于 $\mathcal{D}$。
因此 (2) 的等价条件成立。
\end{proof}

\begin{lemma}
\label{duality-lemma-proper-noetherian}
设 $Y$ 为拟紧拟分离概形，$f : X \to Y$ 为固有态射。若
\footnote{此证明适用于满足下述条件的拟紧拟分离概形态射：
对 $X$ 上每个完美的 $P$，$Rf_*P$ 都是伪凝聚的。由 Kiehl 的一个定理
\cite{Kiehl} 容易推出，当 $f$ 固有且伪凝聚时，此条件成立。
这才是本引理及本章其他若干结果的正确一般性。}
\begin{enumerate}
\item $f$ 平坦且有限表现；或者
\item $Y$ 为 Noetherian 概形，
\end{enumerate}
则对 $Y$ 的所有拟紧开集 $V$，引理 \ref{duality-lemma-when-sheafy}
第 (1) 部分的等价条件均成立。
\end{lemma}

\begin{proof}
设 $Q \in D^+_\QCoh(\mathcal{O}_Y)$ 支撑在 $Y \setminus V$ 上。
反设 $a(Q)$ 不支撑在 $X \setminus U$ 上。则可以找到 $U$ 上的完美复形
$P_U$ 以及非零映射 $P_U \to a(Q)|_U$；这由《概形的导出范畴》，定理
\ref{perfect-theorem-bondal-van-den-Bergh} 得到。再应用《概形的导出范畴》，引理
\ref{perfect-lemma-lift-map-from-perfect-complex-with-support}，
可以假设存在 $X$ 上的完美复形 $P$ 以及映射 $P \to a(Q)$，
其到 $U$ 的限制非零。由 $a$ 的定义，此映射伴随于映射 $Rf_*P \to Q$。

\medskip\noindent
复形 $Rf_*P$ 是伪凝聚的。在情形 (1) 中，这由《概形的导出范畴》，引理
\ref{perfect-lemma-flat-proper-pseudo-coherent-direct-image-general} 得到；
在情形 (2) 中，这由《概形的导出范畴》，引理
\ref{perfect-lemma-direct-image-coherent} 和
\ref{perfect-lemma-identify-pseudo-coherent-noetherian} 得到。
因此可应用《概形的导出范畴》，引理
\ref{perfect-lemma-map-from-pseudo-coherent-to-complex-with-support}，
得到完美复形的映射 $I \to \mathcal{O}_Y$，其到 $V$ 的限制为同构，
并且复合 $I \otimes^\mathbf{L}_{\mathcal{O}_Y} Rf_*P \to Rf_*P \to Q$ 为零。
由《概形的导出范畴》，引理 \ref{perfect-lemma-cohomology-base-change}，有
$I \otimes^\mathbf{L}_{\mathcal{O}_Y} Rf_*P =
Rf_*(Lf^*I \otimes^\mathbf{L}_{\mathcal{O}_X} P)$。
从而复合
$$
Lf^*I \otimes^\mathbf{L}_{\mathcal{O}_X} P \to P \to a(Q)
$$
为零。然而，它到 $U$ 的限制正是映射 $P|_U \to a(Q)|_U$，
而我们假设该映射非零。此矛盾完成证明。
\end{proof}












\section{推前的右伴随与基变换，I}
\label{duality-section-base-change-map}

\noindent
映射 (\ref{duality-equation-sheafy}) 是基变换映射的一种特殊情形。
具体地，假设有拟紧拟分离概形的笛卡尔图
$$
\xymatrix{
X' \ar[r]_{g'} \ar[d]_{f'} & X \ar[d]^f \\
Y' \ar[r]^g & Y
}
$$
即 (\ref{duality-equation-base-change}) 中那样的图。假设 $f$ 与 $g$
{\bf Tor 独立}。于是可考虑由复合
\begin{equation}
\label{duality-equation-base-change-map}
L(g')^* \circ a \to
L(g')^* \circ a \circ Rg_* \circ Lg^* \leftarrow
L(g')^* \circ Rg'_* \circ a' \circ Lg^* \to a' \circ Lg^*
\end{equation}
给出的函子 $D_\QCoh(\mathcal{O}_Y) \to D_\QCoh(\mathcal{O}_{X'})$ 的态射。
第一个箭头来自伴随映射 $\text{id} \to Rg_* Lg^*$，
最后一个箭头来自伴随映射 $L(g')^*Rg'_* \to \text{id}$。
为使中间的箭头可逆，需要 Tor 独立假设；见引理
\ref{duality-lemma-flat-precompose-pus}。也可以利用 $L(g')^*$ 与 $R(g')_*$
的伴随性，把 (\ref{duality-equation-base-change-map}) 看作自然变换
$$
a \to a \circ Rg_* \circ Lg^* \leftarrow Rg'_* \circ a' \circ Lg^*
$$
其中第二个箭头同样可逆。若 $M \in D_\QCoh(\mathcal{O}_X)$ 且
$K \in D_\QCoh(\mathcal{O}_Y)$，则此映射在 Yoneda 函子上由下列映射给出：
\begin{align*}
\Hom_X(M, a(K))
& =
\Hom_Y(Rf_*M, K) \\
& \to
\Hom_Y(Rf_*M, Rg_* Lg^*K) \\
& =
\Hom_{Y'}(Lg^*Rf_*M, Lg^*K) \\
& \leftarrow
\Hom_{Y'}(Rf'_* L(g')^*M, Lg^*K) \\
& =
\Hom_{X'}(L(g')^*M, a'(Lg^*K)) \\
& =
\Hom_X(M, Rg'_*a'(Lg^*K))
\end{align*}
其中向左的箭头由《概形的导出范畴》，引理
\ref{perfect-lemma-compare-base-change} 中的基变换定理可逆；
这使构造更加明确。

\medskip\noindent
本节先证明，该基变换映射在叠置方块时满足若干自然相容性，
正如通常的基变换映射在《上同调》，注
\ref{cohomology-remark-compose-base-change} 和
\ref{cohomology-remark-compose-base-change-horizontal} 中那样。
建议初次阅读时跳过本节余下内容。

\begin{lemma}
\label{duality-lemma-compose-base-change-maps}
考虑拟紧拟分离概形的交换图
$$
\xymatrix{
X' \ar[r]_k \ar[d]_{f'} & X \ar[d]^f \\
Y' \ar[r]^l \ar[d]_{g'} & Y \ar[d]^g \\
Z' \ar[r]^m & Z
}
$$
其中两个方块都是笛卡尔的，并且 $f$ 与 $l$ Tor 独立，
$g$ 与 $m$ 也 Tor 独立。则两个方块的映射
(\ref{duality-equation-base-change-map}) 复合后给出外侧矩形的基变换映射
（精确陈述见证明）。
\end{lemma}

\begin{proof}
由假设可知 $g \circ f$ 与 $m$ Tor 独立（细节从略），故此陈述有意义。
在本证明中，以 $k^*$ 代替 $Lk^*$，以 $f_*$ 代替 $Rf_*$。
令 $a$、$b$、$c$ 分别为引理 \ref{duality-lemma-twisted-inverse-image}
对 $f$、$g$、$g \circ f$ 给出的右伴随，并对带撇版本作类似约定。
对应于上方方块的箭头是复合
$$
\gamma_{top} :
k^* \circ a \to k^* \circ a \circ l_* \circ l^*
\xleftarrow{\xi_{top}} k^* \circ k_* \circ a' \circ l^* \to a' \circ l^*
$$
其中 $\xi_{top} : k_* \circ a' \to a \circ l_*$ 是同构，因而可逆；
它是基变换映射 $l^* \circ f_* \to f'_* \circ k^*$ 的“对偶”箭头。
外侧两个箭头来自典范映射 $1 \to l_* \circ l^*$ 和
$k^* \circ k_* \to 1$。类似地，对第二个方块有
$$
\gamma_{bot} :
l^* \circ b \to l^* \circ b \circ m_* \circ m^*
\xleftarrow{\xi_{bot}} l^* \circ l_* \circ b' \circ m^* \to b' \circ m^*
$$
对外侧矩形则得到
$$
\gamma_{rect} :
k^* \circ c \to k^* \circ c \circ m_* \circ m^*
\xleftarrow{\xi_{rect}} k^* \circ k_* \circ c' \circ m^* \to c' \circ m^*
$$
有 $(g \circ f)_* = g_* \circ f_*$，因而 $c = a \circ b$，
类似地 $c' = a' \circ b'$。本引理断言 $\gamma_{rect}$ 等于复合
$$
k^* \circ c = k^* \circ a \circ b \xrightarrow{\gamma_{top}}
a' \circ l^* \circ b \xrightarrow{\gamma_{bot}}
a' \circ b' \circ m^* = c' \circ m^*
$$
为此，考察下图：
$$
\xymatrix{
& & k^* \circ a \circ b \ar[d] \ar[lldd] \\
& & k^* \circ a \circ l_* \circ l^* \circ b \ar[ld] \\
k^* \circ a \circ b \circ m_* \circ m^* \ar[r] &
k^* \circ a \circ l_* \circ l^* \circ b \circ m_* \circ m^* &
k^* \circ k_* \circ a' \circ l^* \circ b \ar[u]_{\xi_{top}} \ar[d] \ar[ld] \\
& k^*\circ k_* \circ a' \circ l^* \circ b \circ m_* \circ m^*
\ar[u]_{\xi_{top}} \ar[rd] &
a' \circ l^* \circ b \ar[d] \\
k^* \circ k_* \circ a' \circ b' \circ m^* \ar[uu]_{\xi_{rect}} \ar[ddrr] &
k^*\circ k_* \circ a' \circ l^* \circ l_* \circ b' \circ m^*
\ar[u]_{\xi_{bot}} \ar[l] \ar[dr] &
a' \circ l^* \circ b \circ m_* \circ m^* \\
& & a' \circ l^* \circ l_* \circ b' \circ m^* \ar[u]_{\xi_{bot}} \ar[d] \\
& & a' \circ b' \circ m^*
}
$$
沿右侧向下得到上述复合，沿左侧向下得到 $\gamma_{rect}$。
图右侧的所有四边形都交换，这由《范畴》，引理
\ref{categories-lemma-properties-2-cat-cats} 得到；也可更直接地使用
《范畴》，定义 \ref{categories-definition-horizontal-composition} 之前的讨论。
因此，只需说明在将箭头 $\xi_{top}$、$\xi_{bot}$、$\xi_{rect}$
取逆之后，图
$$
\xymatrix{
a \circ l_* \circ l^* \circ b \circ m_* &
a \circ b \circ m_* \ar[l] \\
k_* \circ a' \circ l^* \circ b \circ m_* \ar[u]_{\xi_{top}} & \\
k_* \circ a' \circ l^* \circ l_* \circ b' \ar[u]_{\xi_{bot}} \ar[r] &
k_* \circ a' \circ b' \ar[uu]_{\xi_{rect}}
}
$$
交换；注意，这不同于要求原图交换。然而，由《范畴》，引理
\ref{categories-lemma-properties-2-cat-cats}，图
$$
\xymatrix{
& a \circ l_* \circ l^* \circ b \circ m_* \\
a \circ l_* \circ l^* \circ l_* \circ b'
\ar[ru]^{\xi_{bot}} & &
k_* \circ a' \circ l^* \circ b \circ m_* \ar[ul]_{\xi_{top}} \\
& k_* \circ a' \circ l^* \circ l_* \circ b'
\ar[ul]^{\xi_{top}} \ar[ur]_{\xi_{bot}}
}
$$
交换。又因为图
$$
\vcenter{
\xymatrix{
a \circ l_* \circ l^* \circ b \circ m_* & a \circ b \circ m \ar[l] \\
a \circ l_* \circ l^* \circ l_* \circ b' \ar[u] &
a \circ l_* \circ b' \ar[l] \ar[u]
}
}
\quad\text{且}\quad
\vcenter{
\xymatrix{
a \circ l_* \circ l^* \circ l_* \circ b' \ar[r] & a \circ l_* \circ b' \\
k_* \circ a' \circ l^* \circ l_* \circ b' \ar[u] \ar[r] &
k_* \circ a' \circ b' \ar[u]
}
}
$$
交换（见所引参考），且复合 $l_* \to l_* \circ l^* \circ l_* \to l_*$
是恒等，所以只需证明
$$
k \circ a' \circ b' \xrightarrow{\xi_{bot}} a \circ l_* \circ b
\xrightarrow{\xi_{top}} a \circ b \circ m_*
$$
等于 $\xi_{rect}$（通过识别 $a \circ b = c$ 和
$a' \circ b' = c'$）。这正是《上同调》，注
\ref{cohomology-remark-compose-base-change} 中陈述的对偶，证明完成。
\end{proof}

\begin{lemma}
\label{duality-lemma-compose-base-change-maps-horizontal}
考虑拟紧拟分离概形的交换图
$$
\xymatrix{
X'' \ar[r]_{g'} \ar[d]_{f''} & X' \ar[r]_g \ar[d]_{f'} & X \ar[d]^f \\
Y'' \ar[r]^{h'} & Y' \ar[r]^h & Y
}
$$
其中两个方块都是笛卡尔的，并且 $f$ 与 $h$ Tor 独立，
$f'$ 与 $h'$ 也 Tor 独立。则两个方块的映射
(\ref{duality-equation-base-change-map}) 复合后给出外侧矩形的基变换映射
（精确陈述见证明）。
\end{lemma}

\begin{proof}
由假设可知 $f$ 与 $h \circ h'$ Tor 独立（细节从略），故此陈述有意义。
在本证明中，以 $g^*$ 代替 $Lg^*$，以 $f_*$ 代替 $Rf_*$。
令 $a$、$a'$、$a''$ 分别为引理 \ref{duality-lemma-twisted-inverse-image}
对 $f$、$f'$、$f''$ 给出的右伴随。对应于右侧方块的箭头是复合
$$
\gamma_{right} :
g^* \circ a \to g^* \circ a \circ h_* \circ h^*
\xleftarrow{\xi_{right}} g^* \circ g_* \circ a' \circ h^* \to a' \circ h^*
$$
其中 $\xi_{right} : g_* \circ a' \to a \circ h_*$ 是同构，因而可逆；
它是基变换映射 $h^* \circ f_* \to f'_* \circ g^*$ 的“对偶”箭头。
外侧两个箭头来自典范映射 $1 \to h_* \circ h^*$ 和
$g^* \circ g_* \to 1$。类似地，对左侧方块有
$$
\gamma_{left} :
(g')^* \circ a' \to (g')^* \circ a' \circ (h')_* \circ (h')^*
\xleftarrow{\xi_{left}}
(g')^* \circ (g')_* \circ a'' \circ (h')^* \to a'' \circ (h')^*
$$
对外侧矩形则得到
$$
\gamma_{rect} :
k^* \circ a \to
k^* \circ a \circ m_* \circ m^* \xleftarrow{\xi_{rect}}
k^* \circ k_* \circ a'' \circ m^* \to
a'' \circ m^*
$$
其中 $k = g \circ g'$ 且 $m = h \circ h'$。有
$k^* = (g')^* \circ g^*$ 和 $m^* = (h')^* \circ h^*$。
本引理断言 $\gamma_{rect}$ 等于复合
$$
k^* \circ a =
(g')^* \circ g^* \circ a \xrightarrow{\gamma_{right}}
(g')^* \circ a' \circ h^* \xrightarrow{\gamma_{left}}
a'' \circ (h')^* \circ h^* = a'' \circ m^*
$$
为此，考察下图：
$$
\xymatrix{
& (g')^* \circ g^* \circ a \ar[d] \ar[ddl] \\
& (g')^* \circ g^* \circ a \circ h_* \circ h^* \ar[ld] \\
(g')^* \circ g^* \circ a \circ h_* \circ (h')_* \circ (h')^* \circ h^* &
(g')^* \circ g^* \circ g_* \circ a' \circ h^*
\ar[u]_{\xi_{right}} \ar[d] \ar[ld] \\
(g')^* \circ g^* \circ g_* \circ a' \circ (h')_* \circ (h')^* \circ h^*
\ar[u]_{\xi_{right}} \ar[dr] &
(g')^* \circ a' \circ h^* \ar[d] \\
(g')^* \circ g^* \circ g_* \circ (g')_* \circ a'' \circ (h')^* \circ h^*
\ar[u]_{\xi_{left}} \ar[ddr] \ar[dr] &
(g')^* \circ a' \circ (h')_* \circ (h')^* \circ h^* \\
& (g')^*\circ (g')_* \circ a'' \circ (h')^* \circ h^*
\ar[u]_{\xi_{left}} \ar[d] \\
& a'' \circ (h')^* \circ h^*
}
$$
沿右侧向下得到上述复合，沿左侧向下得到 $\gamma_{rect}$。
图右侧的所有四边形都交换，这由《范畴》，引理
\ref{categories-lemma-properties-2-cat-cats} 得到；也可更直接地使用
《范畴》，定义 \ref{categories-definition-horizontal-composition} 之前的讨论。
因此只需说明
$$
g_* \circ (g')_* \circ a'' \xrightarrow{\xi_{left}}
g_* \circ a' \circ (h')_* \xrightarrow{\xi_{right}}
a \circ h_* \circ (h')_*
$$
等于 $\xi_{rect}$。这正是《上同调》，注
\ref{cohomology-remark-compose-base-change-horizontal} 中陈述的对偶，
证明完成。
\end{proof}

\begin{remark}
\label{duality-remark-going-around}
考虑拟紧拟分离概形的交换图
$$
\xymatrix{
X'' \ar[r]_{k'} \ar[d]_{f''} & X' \ar[r]_k \ar[d]_{f'} & X \ar[d]^f \\
Y'' \ar[r]^{l'} \ar[d]_{g''} & Y' \ar[r]^l \ar[d]_{g'} & Y \ar[d]^g \\
Z'' \ar[r]^{m'} & Z' \ar[r]^m & Z
}
$$
其中所有方块均为笛卡尔的，并且 $(f, l)$、$(g, m)$、
$(f', l')$、$(g', m')$ 都是 Tor 独立的态射对。
令 $a$、$a'$、$a''$、$b$、$b'$、$b''$ 分别为引理
\ref{duality-lemma-twisted-inverse-image} 对 $f$、$f'$、$f''$、$g$、$g'$、$g''$
给出的右伴随。如下把图中的方块标记为 $A$、$B$、$C$、$D$：
$$
\begin{matrix}
A & B \\
C & D
\end{matrix}
$$
则各方块的映射 (\ref{duality-equation-base-change-map}) 为
（这里使用 $k^* = Lk^*$ 等简记）
$$
\begin{matrix}
\gamma_A : (k')^* \circ a' \to a'' \circ (l')^* &
\gamma_B : k^* \circ a \to a' \circ l^* \\
\gamma_C : (l')^* \circ b' \to b'' \circ (m')^* &
\gamma_D : l^* \circ b \to b' \circ m^*
\end{matrix}
$$
对 $2 \times 1$ 和 $1 \times 2$ 的矩形，还有四个基变换映射
$$
\begin{matrix}
\gamma_{A + B} : (k \circ k')^* \circ a \to a'' \circ (l \circ l')^* \\
\gamma_{C + D} : (l \circ l')^* \circ b \to b'' \circ (m \circ m')^* \\
\gamma_{A + C} : (k')^* \circ (a' \circ b') \to (a'' \circ b'') \circ (m')^* \\
\gamma_{B + D} : k^* \circ (a \circ b) \to (a' \circ b') \circ m^*
\end{matrix}
$$
由引理 \ref{duality-lemma-compose-base-change-maps-horizontal}，有
$$
\gamma_{A + B} = \gamma_A \circ \gamma_B, \quad
\gamma_{C + D} = \gamma_C \circ \gamma_D
$$
由引理 \ref{duality-lemma-compose-base-change-maps}，有
$$
\gamma_{A + C} = \gamma_C \circ \gamma_A, \quad
\gamma_{B + D} = \gamma_D \circ \gamma_B
$$
严格地说，按《范畴》，第 \ref{categories-section-formal-cat-cat} 节的记号，
应写成 $\gamma_{A + B} = (\gamma_A \star \text{id}_{l^*}) \circ
(\text{id}_{(k')^*} \star \gamma_B)$，其他情形类似。不过，这里继续沿用
引理 \ref{duality-lemma-compose-base-change-maps} 和
\ref{duality-lemma-compose-base-change-maps-horizontal} 证明中的记号滥用，
省略与恒等变换的 $\star$ 积；只要知道变换的源和目标，便可确定应补入哪些项。
由以上讨论，先验地得到两个变换
$$
(k')^* \circ k^* \circ a \circ b
\longrightarrow
a'' \circ b'' \circ (m')^* \circ m^*
$$
即
$$
\gamma_C \circ \gamma_A \circ \gamma_D \circ \gamma_B =
\gamma_{A + C} \circ \gamma_{B + D}
$$
和
$$
\gamma_C \circ \gamma_D \circ \gamma_A \circ \gamma_B =
\gamma_{C + D} \circ \gamma_{A + B}
$$
本注的要点是指出这两个变换相等。为此，只需证明图
$$
\xymatrix{
(k')^* \circ a' \circ l^* \circ b \ar[r]_{\gamma_D} \ar[d]_{\gamma_A} &
(k')^* \circ a' \circ b' \circ m^* \ar[d]^{\gamma_A} \\
a'' \circ (l')^* \circ l^* \circ b \ar[r]^{\gamma_D} &
a'' \circ (l')^* \circ b' \circ m^*
}
$$
交换。这由《范畴》，引理 \ref{categories-lemma-properties-2-cat-cats}
得到；也可更直接地使用《范畴》，定义
\ref{categories-definition-horizontal-composition} 之前的讨论。
\end{remark}







\section{推前的右伴随与基变换，II}
\label{duality-section-base-change-II}

\noindent
本节证明，第 \ref{duality-section-base-change-map} 节的基变换映射在若干情形下是同构。
首先注意，只要推前的右伴随的形成与开集限制交换，就只需在仿射开集上检验。

\begin{remark}
\label{duality-remark-check-over-affines}
考虑拟紧拟分离概形的笛卡尔图
$$
\xymatrix{
X' \ar[r]_{g'} \ar[d]_{f'} & X \ar[d]^f \\
Y' \ar[r]^g & Y
}
$$
其中 $(g, f)$ Tor 独立。设 $V \subset Y$ 和 $V' \subset Y'$
为满足 $g(V') \subset V$ 的仿射开集。作笛卡尔图
$$
\vcenter{
\xymatrix{
U \ar[r] \ar[d] & X \ar[d] \\
V \ar[r] & Y
}
}
\quad\text{且}\quad
\vcenter{
\xymatrix{
U' \ar[r] \ar[d] & X' \ar[d] \\
V' \ar[r] & Y'
}
}
$$
假设关于 $K$ 和第一个图的 (\ref{duality-equation-sheafy})，以及关于
$Lg^*K$ 和第二个图的 (\ref{duality-equation-sheafy})，都是同构。
则基变换映射 (\ref{duality-equation-base-change-map})
$$
L(g')^*a(K) \longrightarrow a'(Lg^*K)
$$
到 $U'$ 的限制，同构于关于 $K|_V$ 及笛卡尔图
$$
\xymatrix{
U' \ar[r] \ar[d] & U \ar[d] \\
V' \ar[r] & V
}
$$
的基变换映射 (\ref{duality-equation-base-change-map})。这是因为
(\ref{duality-equation-sheafy}) 是基变换映射 (\ref{duality-equation-base-change-map})
的特殊情形，而且水平叠置方块时，基变换映射正确复合；见引理
\ref{duality-lemma-compose-base-change-maps-horizontal}。因此，为检验基变换映射
限制到 $U'$ 后是否为同构，只需处理最后一个图。
\end{remark}

\begin{lemma}
\label{duality-lemma-more-base-change}
在图 (\ref{duality-equation-base-change}) 中，假设
\begin{enumerate}
\item $g : Y' \to Y$ 是仿射概形之间的态射；
\item $f : X \to Y$ 是固有的；并且
\item $f$ 与 $g$ Tor 独立。
\end{enumerate}
则在下列情形中，基变换映射 (\ref{duality-equation-base-change-map}) 诱导同构
$$
L(g')^*a(K) \longrightarrow a'(Lg^*K)
$$
\begin{enumerate}
\item 若 $f$ 平坦且有限表现，则对所有 $K \in D_\QCoh(\mathcal{O}_Y)$；
\item 若 $f$ 完美且 $Y$ 为 Noetherian 概形，则对所有
$K \in D_\QCoh(\mathcal{O}_Y)$；
\item 若 $g$ 具有有限 Tor 维数且 $Y$ 为 Noetherian 概形，则对
$K \in D_\QCoh^+(\mathcal{O}_Y)$。
\end{enumerate}
\end{lemma}

\begin{proof}
记 $Y = \Spec(A)$、$Y' = \Spec(A')$。作为仿射态射的基变换，态射
$g'$ 是仿射的。设 $M$ 为 $D_\QCoh(\mathcal{O}_X)$ 的一个完美生成元；
见《概形的导出范畴》，定理
\ref{perfect-theorem-bondal-van-den-Bergh}。于是 $L(g')^*M$ 是
$D_\QCoh(\mathcal{O}_{X'})$ 的生成元；见《概形的导出范畴》，注
\ref{perfect-remark-pullback-generator}。因此，只需证明
(\ref{duality-equation-base-change-map}) 诱导全局 Hom 复形的同构
\begin{equation}
\label{duality-equation-iso}
R\Hom_{X'}(L(g')^*M, L(g')^*a(K))
\longrightarrow
R\Hom_{X'}(L(g')^*M, a'(Lg^*K))
\end{equation}
见《上同调》，第 \ref{cohomology-section-global-RHom} 节；
因为这将蕴含 $L(g')^*a(K) \to a'(Lg^*K)$ 的锥为零。
证明结构如下：先说明这两个 Hom 复形同构，再在证明最后说明该同构确由
(\ref{duality-equation-iso}) 诱导。

\medskip\noindent
左端。由于 $M$ 是完美的，由《概形的导出范畴》，引理
\ref{perfect-lemma-affine-morphism-and-hom-out-of-perfect}，典范映射
$$
R\Hom_X(M, a(K)) \otimes^\mathbf{L}_A A'
\longrightarrow
R\Hom_{X'}(L(g')^*M, L(g')^*a(K))
$$
是同构。将其与引理 \ref{duality-lemma-iso-global-hom} 的同构
$R\Hom_Y(Rf_*M, K) = R\Hom_X(M, a(K))$ 结合，可知左端等于
$R\Hom_Y(Rf_*M, K) \otimes^\mathbf{L}_A A'$。

\medskip\noindent
右端。先使用引理 \ref{duality-lemma-iso-global-hom} 的同构
$$
R\Hom_{X'}(L(g')^*M, a'(Lg^*K)) = R\Hom_{Y'}(Rf'_*L(g')^*M, Lg^*K)
$$
再由《概形的导出范畴》，引理 \ref{perfect-lemma-compare-base-change}，
基变换映射 $Lg^*Rf_*M \to Rf'_*L(g')^*M$ 是同构，故可将右端改写为
$R\Hom_{Y'}(Lg^*Rf_*M, Lg^*K)$。由于 $Y$、$Y'$ 仿射，且
$K$、$Rf_*M$ 属于 $D_\QCoh(\mathcal{O}_Y)$（《概形的导出范畴》，引理
\ref{perfect-lemma-quasi-coherence-direct-image}），在 $D(A')$ 中有典范映射
$$
\beta :
R\Hom_Y(Rf_*M, K) \otimes^\mathbf{L}_A A'
\longrightarrow
R\Hom_{Y'}(Lg^*Rf_*M, Lg^*K)
$$
这就是《代数进阶》，等式 (\ref{more-algebra-equation-base-change-RHom})
中的箭头；这里使用《概形的导出范畴》，引理
\ref{perfect-lemma-affine-compare-bounded} 和
\ref{perfect-lemma-quasi-coherence-internal-hom} 在代数表述之间来回转换。
\begin{enumerate}
\item 若 $f$ 平坦且有限表现，则由《概形的导出范畴》，引理
\ref{perfect-lemma-flat-proper-perfect-direct-image-general}，复形 $Rf_*M$
在 $Y$ 上完美；再由《代数进阶》，引理
\ref{more-algebra-lemma-base-change-RHom} 的第 (1) 部分，$\beta$ 是同构。
\item 若 $f$ 完美且 $Y$ 为 Noetherian 概形，则由《态射进阶》，引理
\ref{more-morphisms-lemma-perfect-proper-perfect-direct-image}，复形
$Rf_*M$ 在 $Y$ 上完美；于是 $\beta$ 同前为同构。
\item 若 $g$ 具有有限 Tor 维数且 $Y$ 为 Noetherian 概形，则复形
$Rf_*M$ 在 $Y$ 上伪凝聚（《概形的导出范畴》，引理
\ref{perfect-lemma-direct-image-coherent} 和
\ref{perfect-lemma-identify-pseudo-coherent-noetherian}）；由《代数进阶》，
引理 \ref{more-algebra-lemma-base-change-RHom} 的第 (4) 部分，$\beta$ 是同构。
\end{enumerate}
因此得到与上一段相同的答案。

\medskip\noindent
在证明余下部分中，我们说明第二、三段给出的
(\ref{duality-equation-iso}) 左右两端的识别，实际上正由
(\ref{duality-equation-iso}) 给出。为使公式便于处理，记
$(-, -)_X = R\Hom_X(-, -)$，以 $- \otimes A'$ 代替
$- \otimes_A^\mathbf{L} A'$，并简记 $g^* = Lg^*$、$f_* = Rf_*$。
考虑交换图
$$
\xymatrix{
((g')^*M, (g')^*a(K))_{X'} \ar[d] &
(M, a(K))_X \otimes A' \ar[l]^-\alpha \ar[d] &
(f_*M, K)_Y \otimes A' \ar@{=}[l] \ar[d] \\
((g')^*M, (g')^*a(g_*g^*K))_{X'} &
(M, a(g_*g^*K))_X \otimes A' \ar[l]^-\alpha &
(f_*M, g_*g^*K)_Y \otimes A' \ar@{=}[l] \ar@/_4pc/[dd]_{\mu'} \\
((g')^*M, (g')^*g'_*a'(g^*K))_{X'} \ar[u] \ar[d] &
(M, g'_*a'(g^*K))_X \otimes A' \ar[u] \ar[l]^-\alpha \ar[ld]^\mu &
(f_*M, K) \otimes A' \ar[d]^\beta \\
((g')^*M, a'(g^*K))_{X'} &
(f'_*(g')^*M, g^*K)_{Y'} \ar@{=}[l] \ar[r] &
(g^*f_*M, g^*K)_{Y'}
}
$$
标记为 $\alpha$ 的箭头，是《概形的导出范畴》，引理
\ref{perfect-lemma-affine-morphism-and-hom-out-of-perfect}
对四角为 $X', X, Y', Y$ 的图给出的映射。由于水平箭头对各变量函子性，
图的上部交换。中间的竖直箭头来自引理
\ref{duality-lemma-flat-precompose-pus} 的可逆变换
$g'_* \circ a' \to a \circ g_*$，故中间方块交换。沿左侧向下正是
(\ref{duality-equation-iso})。上方的水平箭头给出证明第二段所用的识别；
下方的水平箭头（包括 $\beta$）给出第三段所用的识别。
给定 $E \in D(A)$、$E' \in D(A')$ 及 $D(A)$ 中的
$c : E \to E'$，用 $\mu_c : E \otimes A' \to E'$ 表示由 $c$
以及限制与基变换的伴随性诱导的映射；若 $c$ 明确，就写
$\mu = \mu_c$，即从记号中省略 $c$。图中的 $\mu$ 正是这种映射，
其中 $c$ 由识别
$(M, g'_*a(g^*K))_X = ((g')^*M, a'(g^*K))_{X'}$ 给出；
含 $\mu$ 的三角形由《概形的导出范畴》，注
\ref{perfect-remark-multiplication-map} 交换。

\medskip\noindent
注意图
$$
\xymatrix{
(M, a(g_*g^*K))_X &
(f_*M, g_* g^*K)_Y \ar@{=}[l] &
(g^*f_*M, g^*K)_{Y'} \ar@{=}[l] \\
(M, g'_* a'(g^*K))_X \ar[u] &
((g')^*M, a'(g^*K))_{X'} \ar@{=}[l] &
(f'_*(g')^*M, g^*K)_{Y'} \ar@{=}[l] \ar[u]
}
$$
依变换 $g'_* \circ a' \to a \circ g_*$ 的定义交换。令 $\mu'$
如上对应于识别 $(f_*M, g_*g^*K)_X = (g^*f_*M, g^*K)_{Y'}$，
则六边形也交换。因此，只需证明 $\beta$ 等于
$(f_*M, K)_Y \otimes A' \to (f_*M, g_*g^*K)_X \otimes A'$
与 $\mu'$ 的复合。为此，只需证明所诱导的两个映射
$(f_*M, K)_Y \to (g^*f_*M, g^*K)_{Y'}$ 相同。换言之，
只需证明对所有 $E, K \in D(A)$，图
$$
\xymatrix{
R\Hom_A(E, K) \ar[rr]_{\text{由下列对象诱导： }\beta} \ar[rd] & &
R\Hom_{A'}(E \otimes_A^\mathbf{L} A', K \otimes_A^\mathbf{L} A') \\
& R\Hom_A(E, K \otimes_A^\mathbf{L} A') \ar[ru]
}
$$
交换。而《代数进阶》，第 \ref{more-algebra-section-base-change-RHom} 节
正是如此构造 $\beta$ 的，证明完成。
\end{proof}








\section{推前的右伴随与迹映射}
\label{duality-section-trace}

\noindent
设 $f : X \to Y$ 为拟紧拟分离概形之间的态射，并设
$a : D_\QCoh(\mathcal{O}_Y) \to D_\QCoh(\mathcal{O}_X)$ 为引理
\ref{duality-lemma-twisted-inverse-image} 中的右伴随。由《范畴》，第
\ref{categories-section-adjoint} 节，得到函子变换
$$
\text{Tr}_f : Rf_* \circ a \longrightarrow \text{id}
$$
对 $K \in D_\QCoh(\mathcal{O}_Y)$，相应的映射
$\text{Tr}_{f, K} : Rf_*a(K) \longrightarrow K$ 有时称为{\it 迹映射}。
它具有如下性质：刻画右伴随的双射
$$
\Hom_X(L, a(K)) \longrightarrow \Hom_Y(Rf_*L, K)
$$
对 $L \in D_\QCoh(\mathcal{O}_X)$ 由
$$
\varphi \longmapsto \text{Tr}_{f, K} \circ Rf_*\varphi
$$
给出。映射 (\ref{duality-equation-sheafy-trace})
$$
Rf_*R\SheafHom_{\mathcal{O}_X}(L, a(K))
\longrightarrow
R\SheafHom_{\mathcal{O}_Y}(Rf_*L, K)
$$
由与 $\text{Tr}_{f, K}$ 复合得到。本节将考察的每个迹映射都是这个迹映射的特殊情形。
在讨论若干特殊情形之前，先证明迹映射的形成与基变换交换。

\begin{lemma}[迹映射与基变换]
\label{duality-lemma-trace-map-and-base-change}
假设有图 (\ref{duality-equation-base-change})，其中 $f$ 与 $g$ Tor 独立。
则映射 $1 \star \text{Tr}_f : Lg^* \circ Rf_* \circ a \to Lg^*$ 和
$\text{Tr}_{f'} \star 1 : Rf'_* \circ a' \circ Lg^* \to Lg^*$
通过基变换映射
$\beta : Lg^* \circ Rf_* \to Rf'_* \circ L(g')^*$
（《上同调》，注 \ref{cohomology-remark-base-change}）以及
$\alpha : L(g')^* \circ a \to a' \circ Lg^*$
（\ref{duality-equation-base-change-map}）相一致。更精确地，函子变换的图
$$
\xymatrix{
Lg^* \circ Rf_* \circ a
\ar[d]_{\beta \star 1} \ar[r]_-{1 \star \text{Tr}_f} &
Lg^* \\
Rf'_* \circ L(g')^* \circ a \ar[r]^{1 \star \alpha} &
Rf'_* \circ a' \circ Lg^* \ar[u]_{\text{Tr}_{f'} \star 1}
}
$$
交换。
\end{lemma}

\begin{proof}
本证明中，以 $f_*$ 代替 $Rf_*$，以 $g^*$ 代替 $Lg^*$，
并省略与恒等变换的 $\star$ 积；只要知道变换的源和目标，
便可确定应补入哪些项。回忆，
$\beta : g^* \circ f_* \to f'_* \circ (g')^*$ 是同构，而 $\alpha$
利用同构 $\beta^\vee : g'_* \circ a' \to a \circ g_*$ 定义；
后者是 $\beta$ 的伴随，见引理 \ref{duality-lemma-flat-precompose-pus} 及其证明。
先注意，引理所述图的上方水平箭头等于复合
$$
g^* \circ f_* \circ a \to
g^* \circ f_* \circ a \circ g_* \circ g^* \to
g^* \circ g_* \circ g^* \to g^*
$$
其中第一个箭头是 $(g^*, g_*)$ 的单位，第二个箭头是
$\text{Tr}_f$，第三个箭头是 $(g^*, g_*)$ 的余单位。
这是单位与余单位的复合
$g^* \to g^* \circ g_* \circ g^* \to g^*$ 为恒等这一事实的直接推论。
考虑图
$$
\xymatrix{
& g^* \circ f_* \circ a \ar[ld]_\beta \ar[d] \ar[r]_{\text{Tr}_f} & g^* \\
f'_* \circ (g')^* \circ a \ar[dr] &
g^* \circ f_* \circ a \circ g_* \circ g^* \ar[d]_\beta \ar[ru] &
g^* \circ f_* \circ g'_* \circ a' \circ g^* \ar[l]_{\beta^\vee} \ar[d]_\beta &
f'_* \circ a' \circ g^* \ar[lu]_{\text{Tr}_{f'}} \\
& f'_* \circ (g')^* \circ a \circ g_* \circ g^* &
f'_* \circ (g')^* \circ g'_* \circ a' \circ g^* \ar[ru] \ar[l]_{\beta^\vee}
}
$$
其中两个方块交换，这由《范畴》，引理
\ref{categories-lemma-properties-2-cat-cats} 得到；也可更直接地使用
《范畴》，定义 \ref{categories-definition-horizontal-composition} 之前的讨论。
三角形由上述讨论交换。由《范畴》，引理
\ref{categories-lemma-transformation-between-functors-and-adjoints}，方块
$$
\xymatrix{
g^* \circ f_* \circ g'_* \circ a' \ar[d]_{\beta^\vee} \ar[r]_-\beta &
f'_* \circ (g')^* \circ g'_* \circ a' \ar[d] \\
g^* \circ f_* \circ a \circ g_* \ar[r] &
\text{id}
}
$$
交换，从而大图中的五边形交换。由于 $\beta$ 和 $\beta^\vee$ 均为同构，
并且依定义沿大图外侧的复合等于
$\text{Tr}_f \circ \alpha \circ \beta$，故本引理得证。
\end{proof}

\noindent
设 $f : X \to Y$ 为拟紧拟分离概形之间的态射，并设
$a : D_\QCoh(\mathcal{O}_Y) \to D_\QCoh(\mathcal{O}_X)$ 为引理
\ref{duality-lemma-twisted-inverse-image} 中 $Rf_*$ 的右伴随。由《范畴》，
第 \ref{categories-section-adjoint} 节，得到函子变换
$$
\eta_f : \text{id} \to  a \circ Rf_*
$$
称为该伴随的单位。

\begin{lemma}
\label{duality-lemma-unit-and-base-change}
假设有图 (\ref{duality-equation-base-change})，其中 $f$ 与 $g$ Tor 独立。
则映射 $1 \star \eta_f : L(g')^* \to L(g')^* \circ a \circ Rf_*$ 和
$\eta_{f'} \star 1 : L(g')^* \to a' \circ Rf'_* \circ L(g')^*$
通过基变换映射
$\beta : Lg^* \circ Rf_* \to Rf'_* \circ L(g')^*$
（《上同调》，注 \ref{cohomology-remark-base-change}）以及
$\alpha : L(g')^* \circ a \to a' \circ Lg^*$
（\ref{duality-equation-base-change-map}）相一致。更精确地，函子变换的图
$$
\xymatrix{
L(g')^* \ar[r]_-{1 \star \eta_f} \ar[d]_{\eta_{f'} \star 1} &
L(g')^* \circ a \circ Rf_* \ar[d]^\alpha \\
a' \circ Rf'_* \circ L(g')^* &
a' \circ Lg^* \circ Rf_* \ar[l]_-\beta
}
$$
交换。
\end{lemma}

\begin{proof}
此证明对偶于引理 \ref{duality-lemma-trace-map-and-base-change} 的证明。
本证明中，以 $f_*$ 代替 $Rf_*$，以 $g^*$ 代替 $Lg^*$，
并省略与恒等变换的 $\star$ 积；只要知道变换的源和目标，
便可确定应补入哪些项。回忆，
$\beta : g^* \circ f_* \to f'_* \circ (g')^*$ 是同构，而 $\alpha$
利用同构 $\beta^\vee : g'_* \circ a' \to a \circ g_*$ 定义；
后者是 $\beta$ 的伴随，见引理 \ref{duality-lemma-flat-precompose-pus} 及其证明。
先注意，引理所述图的左侧竖直箭头等于复合
$$
(g')^* \to (g')^* \circ g'_* \circ (g')^* \to
(g')^* \circ g'_* \circ a' \circ f'_* \circ (g')^* \to
a' \circ f'_* \circ (g')^*
$$
其中第一个箭头是 $((g')^*, g'_*)$ 的单位，第二个箭头是
$\eta_{f'}$，第三个箭头是 $((g')^*, g'_*)$ 的余单位。
这是单位与余单位的复合
$(g')^* \to (g')^* \circ (g')_* \circ (g')^* \to (g')^*$
为恒等这一事实的直接推论。考虑图
$$
\xymatrix{
& (g')^* \circ a \circ f_* \ar[r] &
(g')^* \circ a \circ g_* \circ g^* \circ f_*
\ar[ld]_\beta \\
(g')^* \ar[ru]^{\eta_f} \ar[dd]_{\eta_{f'}} \ar[rd] &
(g')^* \circ a \circ g_* \circ f'_* \circ (g')^* &
(g')^* \circ g'_* \circ a' \circ g^* \circ f_*
\ar[u]_{\beta^\vee} \ar[ld]_\beta \ar[d] \\
& (g')^* \circ g'_* \circ a' \circ f'_* \circ (g')^*
\ar[ld] \ar[u]_{\beta^\vee} &
a' \circ g^* \circ f_* \ar[lld]^\beta \\
a' \circ f'_* \circ (g')^*
}
$$
其中两个方块交换，这由《范畴》，引理
\ref{categories-lemma-properties-2-cat-cats} 得到；也可更直接地使用
《范畴》，定义 \ref{categories-definition-horizontal-composition} 之前的讨论。
三角形由上述讨论交换。由《范畴》，引理
\ref{categories-lemma-transformation-between-functors-and-adjoints} 的对偶，方块
$$
\xymatrix{
\text{id} \ar[r] \ar[d] &
g'_* \circ a' \circ g^* \circ f_* \ar[d]^\beta \\
g'_* \circ a' \circ g^* \circ f_* \ar[r]^{\beta^\vee} &
a \circ g_* \circ f'_* \circ (g')^*
}
$$
交换，从而大图中的五边形交换。由于 $\beta$ 和 $\beta^\vee$ 均为同构，
并且依定义沿大图外侧的复合等于
$\beta \circ \alpha \circ \eta_f$，故本引理得证。
\end{proof}

\begin{example}
\label{duality-example-trace-affine}
设 $A \to B$ 为环同态，$Y = \Spec(A)$、$X = \Spec(B)$，
并设 $f : X \to Y$ 为与 $A \to B$ 对应的态射。如例
\ref{duality-example-affine-twisted-inverse-image} 中所见，
$Rf_* : D_\QCoh(\mathcal{O}_X) \to D_\QCoh(\mathcal{O}_Y)$ 的右伴随
把 $D(A) = D_\QCoh(\mathcal{O}_Y)$ 的对象 $K$ 送到
$D(B) = D_\QCoh(\mathcal{O}_X)$ 中的 $R\Hom(B, K)$。迹映射为
$$
\text{Tr}_{f, K} : R\Hom(B, K) \longrightarrow R\Hom(A, K) = K
$$
它由 $A$-模映射 $A \to B$ 诱导。
\end{example}




\section{推前的右伴随与拉回}
\label{duality-section-compare-with-pullback}

\noindent
设 $f : X \to Y$ 为拟紧拟分离概形之间的态射，$a$ 为引理
\ref{duality-lemma-twisted-inverse-image} 中推前的右伴随。对
$K, L \in D_\QCoh(\mathcal{O}_Y)$，有典范映射
$$
Lf^*K \otimes^\mathbf{L}_{\mathcal{O}_X} a(L)
\longrightarrow
a(K \otimes_{\mathcal{O}_Y}^\mathbf{L} L)
$$
具体地，此映射伴随于映射
$$
Rf_*(Lf^*K \otimes^\mathbf{L}_{\mathcal{O}_X} a(L)) =
K \otimes^\mathbf{L}_{\mathcal{O}_Y} Rf_*(a(L))
\longrightarrow
K \otimes^\mathbf{L}_{\mathcal{O}_Y} L
$$
其中等式由《概形的导出范畴》，引理
\ref{perfect-lemma-cohomology-base-change}，而箭头使用迹映射
$Rf_*a(L) \to L$。取 $L = \mathcal{O}_Y$，得到映射
\begin{equation}
\label{duality-equation-compare-with-pullback}
Lf^*K \otimes^\mathbf{L}_{\mathcal{O}_X} a(\mathcal{O}_Y) \longrightarrow a(K)
\end{equation}
它关于 $K$ 函子性，并与可区别三角相容。

\begin{lemma}
\label{duality-lemma-compare-with-pullback-perfect}
设 $f : X \to Y$ 为拟紧拟分离概形之间的态射。对
$K, L \in D_\QCoh(\mathcal{O}_Y)$，若 $K$ 完美，则上面定义的映射
$Lf^*K \otimes^\mathbf{L}_{\mathcal{O}_X} a(L) \to
a(K \otimes_{\mathcal{O}_Y}^\mathbf{L} L)$ 是同构。特别地，若
$K$ 完美，则 (\ref{duality-equation-compare-with-pullback}) 是同构。
\end{lemma}

\begin{proof}
令 $K^\vee$ 为 $K$ 的“对偶”；见《上同调》，引理
\ref{cohomology-lemma-dual-perfect-complex}。对
$M \in D_\QCoh(\mathcal{O}_X)$，有
\begin{align*}
\Hom_{D(\mathcal{O}_Y)}(Rf_*M, K \otimes^\mathbf{L}_{\mathcal{O}_Y} L)
& =
\Hom_{D(\mathcal{O}_Y)}(
Rf_*M \otimes^\mathbf{L}_{\mathcal{O}_Y} K^\vee, L) \\
& =
\Hom_{D(\mathcal{O}_X)}(
M \otimes^\mathbf{L}_{\mathcal{O}_X} Lf^*K^\vee, a(L)) \\
& =
\Hom_{D(\mathcal{O}_X)}(M,
Lf^*K \otimes^\mathbf{L}_{\mathcal{O}_X} a(L))
\end{align*}
第二个等式由 $a$ 的定义和投影公式（《上同调》，引理
\ref{cohomology-lemma-projection-formula-perfect}）得到；也可使用更一般的
《概形的导出范畴》，引理 \ref{perfect-lemma-cohomology-base-change}。
于是由 Yoneda 引理得到结论。
\end{proof}

\begin{lemma}
\label{duality-lemma-restriction-compare-with-pullback}
假设有图 (\ref{duality-equation-base-change})，其中 $f$ 与 $g$ Tor 独立。
设 $K \in D_\QCoh(\mathcal{O}_Y)$。则图
$$
\xymatrix{
L(g')^*(Lf^*K \otimes^\mathbf{L}_{\mathcal{O}_X} a(\mathcal{O}_Y))
\ar[r] \ar[d] & L(g')^*a(K) \ar[d] \\
L(f')^*Lg^*K \otimes_{\mathcal{O}_{X'}}^\mathbf{L} a'(\mathcal{O}_{Y'})
\ar[r] & a'(Lg^*K)
}
$$
交换，其中水平箭头是关于 $K$ 和 $Lg^*K$ 的映射
(\ref{duality-equation-compare-with-pullback})，竖直映射则使用《上同调》，注
\ref{cohomology-remark-base-change} 和
(\ref{duality-equation-base-change-map}) 构造。
\end{lemma}

\begin{proof}
本证明中，以 $f_*$ 代替 $Rf_*$，以 $f^*$ 代替 $Lf^*$ 等，
并以 $\otimes$ 代替 $\otimes^\mathbf{L}_{\mathcal{O}_X}$ 等。
把 (\ref{duality-equation-compare-with-pullback}) 写成复合
\begin{align*}
f^*K \otimes a(\mathcal{O}_Y)
& \to
a(f_*(f^*K \otimes a(\mathcal{O}_Y))) \\
& \leftarrow
a(K \otimes f_*a(\mathcal{O}_K)) \\
& \to
a(K \otimes \mathcal{O}_Y) \\
& \to
a(K)
\end{align*}
其中第一个箭头是单位 $\eta_f$；第二个箭头是将 $a$ 应用于
《上同调》，等式 (\ref{cohomology-equation-projection-formula-map})，
它由《概形的导出范畴》，引理
\ref{perfect-lemma-cohomology-base-change} 为同构；第三个箭头是将
$a$ 应用于 $\text{id}_K \otimes \text{Tr}_f$；第四个箭头是将
$a$ 应用于同构 $K \otimes \mathcal{O}_Y = K$。只需证明这些映射
分别给出引理陈述中的交换方块。对 $\eta_f$ 和 $\text{Tr}_f$，
这分别是引理 \ref{duality-lemma-unit-and-base-change} 和
\ref{duality-lemma-trace-map-and-base-change}。对使用《上同调》，等式
(\ref{cohomology-equation-projection-formula-map}) 的箭头，
这是《上同调》，注 \ref{cohomology-remark-compatible-with-diagram}。
对乘法映射则显然。证明完成。
\end{proof}

\begin{lemma}
\label{duality-lemma-compare-on-open}
设 $f : X \to Y$ 为 Noetherian 概形之间的固有态射。设 $V \subset Y$
为开集，使得 $f^{-1}(V) \to V$ 是同构。则对
$K \in D_\QCoh^+(\mathcal{O}_Y)$，映射
(\ref{duality-equation-compare-with-pullback}) 限制到 $f^{-1}(V)$ 后为同构。
\end{lemma}

\begin{proof}
由引理 \ref{duality-lemma-proper-noetherian}，映射 (\ref{duality-equation-sheafy})
对 $D_\QCoh^+(\mathcal{O}_Y)$ 的对象为同构。因此，由引理
\ref{duality-lemma-restriction-compare-with-pullback}，关于 $K$ 的
(\ref{duality-equation-compare-with-pullback}) 到 $f^{-1}(V)$ 的限制，
就是关于 $K|_V$ 和 $f^{-1}(V) \to V$ 的映射
(\ref{duality-equation-compare-with-pullback})。故只需在 $f$ 为恒等态射时证明该映射为同构；这很清楚。
\end{proof}

\begin{lemma}
\label{duality-lemma-transitivity-compare-with-pullback}
设 $f : X \to Y$ 和 $g : Y \to Z$ 为拟紧拟分离概形的可复合态射，
并令 $h = g \circ f$。设 $a, b, c$ 分别为引理
\ref{duality-lemma-twisted-inverse-image} 对 $f, g, h$ 给出的伴随。
则对任意 $K \in D_\QCoh(\mathcal{O}_Z)$，图
$$
\xymatrix{
Lf^*(Lg^*K \otimes_{\mathcal{O}_Y}^\mathbf{L}
b(\mathcal{O}_Z)) \otimes_{\mathcal{O}_X}^\mathbf{L} a(\mathcal{O}_Y)
\ar@{=}[d] \ar[r] &
a(Lg^*K \otimes_{\mathcal{O}_Y}^\mathbf{L} b(\mathcal{O}_Z)) \ar[r] &
a(b(K)) \ar@{=}[d] \\
Lh^*K \otimes_{\mathcal{O}_X}^\mathbf{L} Lf^*b(\mathcal{O}_Z)
\otimes_{\mathcal{O}_X}^\mathbf{L} a(\mathcal{O}_Y) \ar[r] &
Lh^*K \otimes_{\mathcal{O}_X}^\mathbf{L} c(\mathcal{O}_Z) \ar[r] &
c(K)
}
$$
交换，其中各箭头为 (\ref{duality-equation-compare-with-pullback})，
并使用了 $Lh^* = Lf^* \circ Lg^*$ 和 $c = a \circ b$。
\end{lemma}

\begin{proof}
本证明中，以 $f_*$ 代替 $Rf_*$，以 $f^*$ 代替 $Lf^*$ 等，
并以 $\otimes$ 代替 $\otimes^\mathbf{L}_{\mathcal{O}_X}$ 等。
上方箭头的复合伴随于映射
$$
g_*f_*(f^*(g^*K \otimes b(\mathcal{O}_Z)) \otimes a(\mathcal{O}_Y)) \to K
$$
由《概形的导出范畴》，引理 \ref{perfect-lemma-cohomology-base-change}，
其左端等于 $K \otimes g_*f_*(f^*b(\mathcal{O}_Z) \otimes a(\mathcal{O}_Y))$。
检视各定义可知，该映射来自映射
$$
g_*f_*(f^*b(\mathcal{O}_Z) \otimes a(\mathcal{O}_Y))
\xleftarrow{g_*\epsilon}
g_*(b(\mathcal{O}_Z) \otimes f_*a(\mathcal{O}_Y)) \xrightarrow{g_*\alpha}
g_*(b(\mathcal{O}_Z)) \xrightarrow{\beta} \mathcal{O}_Z
$$
与 $\text{id}_K$ 的张量积。这里 $\epsilon$ 是《概形的导出范畴》，引理
\ref{perfect-lemma-cohomology-base-change} 中的同构，而 $\beta$
来自余单位映射 $g_*b \to \text{id}$。类似地，下方水平箭头的复合
伴随于 $\text{id}_K$ 与下列复合的张量积：
$$
g_*f_*(f^*b(\mathcal{O}_Z) \otimes a(\mathcal{O}_Y)) \xrightarrow{g_*f_*\delta}
g_*f_*(ab(\mathcal{O}_Z)) \xrightarrow{g_*\gamma}
g_*(b(\mathcal{O}_Z)) \xrightarrow{\beta}
\mathcal{O}_Z
$$
其中 $\gamma$ 来自余单位映射 $f_*a \to \text{id}$，而 $\delta$
是以复合
$$
f_*(f^*b(\mathcal{O}_Z) \otimes a(\mathcal{O}_Y))
\xleftarrow{\epsilon}
b(\mathcal{O}_Z) \otimes f_*a(\mathcal{O}_Y) \xrightarrow{\alpha}
b(\mathcal{O}_Z)
$$
为伴随的映射。由伴随函子、伴随映射和余单位的一般性质
（见《范畴》，第 \ref{categories-section-adjoint} 节），
如所需有 $\gamma \circ f_*\delta = \alpha \circ \epsilon^{-1}$。
\end{proof}





\section{闭浸入的推前右伴随}
\label{duality-section-sections-with-exact-support}

\noindent
设 $i : (Z, \mathcal{O}_Z) \to (X, \mathcal{O}_X)$ 为赋环空间的态射，
使得 $i$ 同胚地映到一个闭子集上，并且
$i^\sharp : \mathcal{O}_X \to i_*\mathcal{O}_Z$ 满射
（例如概形的闭浸入）。令 $\mathcal{I} = \Ker(i^\sharp)$。
对 $\mathcal{O}_X$-模层 $\mathcal{F}$，层
$$
\SheafHom_{\mathcal{O}_X}(i_*\mathcal{O}_Z, \mathcal{F})
$$
是被 $\mathcal{I}$ 零化的 $\mathcal{O}_X$-模层。因此，由《模》，引理
\ref{modules-lemma-i-star-equivalence}，存在一个 $\mathcal{O}_Z$-模层，
记作 $\SheafHom(\mathcal{O}_Z, \mathcal{F})$，使得作为
$\mathcal{O}_X$-模有
$$
i_*\SheafHom(\mathcal{O}_Z, \mathcal{F}) =
\SheafHom_{\mathcal{O}_X}(i_*\mathcal{O}_Z, \mathcal{F})
$$
下面详述其含义。

\begin{lemma}
\label{duality-lemma-compute-sheaf-with-exact-support}
采用上述记号。函子 $\SheafHom(\mathcal{O}_Z, -)$ 是函子
$i_* : \textit{Mod}(\mathcal{O}_Z) \to \textit{Mod}(\mathcal{O}_X)$
的右伴随。对开集 $V \subset Z$，有
$$
\Gamma(V, \SheafHom(\mathcal{O}_Z, \mathcal{F})) =
\{s \in \Gamma(U, \mathcal{F}) \mid \mathcal{I}s = 0\}
$$
其中 $U \subset X$ 为与 $Z$ 的交等于 $V$ 的开集。
\end{lemma}

\begin{proof}
设 $\mathcal{G}$ 为 $\mathcal{O}_Z$-模层。则
$$
\Hom_{\mathcal{O}_X}(i_*\mathcal{G}, \mathcal{F}) =
\Hom_{i_*\mathcal{O}_Z}(i_*\mathcal{G},
\SheafHom_{\mathcal{O}_X}(i_*\mathcal{O}_Z, \mathcal{F})) =
\Hom_{\mathcal{O}_Z}(\mathcal{G}, \SheafHom(\mathcal{O}_Z, \mathcal{F}))
$$
第一个等式由《模》，引理 \ref{modules-lemma-adjoint-tensor-restrict}，
第二个等式由 $i_*$ 的全忠实性；见《模》，引理
\ref{modules-lemma-i-star-equivalence}。截面的描述留给读者验证。
\end{proof}

\noindent
函子
$$
\textit{Mod}(\mathcal{O}_X)
\longrightarrow
\textit{Mod}(\mathcal{O}_Z),
\quad
\mathcal{F} \longmapsto \SheafHom(\mathcal{O}_Z, \mathcal{F})
$$
左正合，并有导出延拓
$$
R\SheafHom(\mathcal{O}_Z, -) : D(\mathcal{O}_X) \to D(\mathcal{O}_Z).
$$

\begin{lemma}
\label{duality-lemma-sheaf-with-exact-support-adjoint}
采用上述记号。函子 $R\SheafHom(\mathcal{O}_Z, -)$ 是函子
$Ri_* : D(\mathcal{O}_Z) \to D(\mathcal{O}_X)$ 的右伴随。
\end{lemma}

\begin{proof}
由引理 \ref{duality-lemma-compute-sheaf-with-exact-support}，$i_*$ 与
$\SheafHom(\mathcal{O}_Z, -)$ 是伴随函子，故结论成立；见《导出范畴》，
引理 \ref{derived-lemma-derived-adjoint-functors}。
\end{proof}

\begin{lemma}
\label{duality-lemma-sheaf-with-exact-support-ext}
采用上述记号。对 $D(\mathcal{O}_X)$ 中所有 $K$，在
$D(\mathcal{O}_X)$ 中有
$$
Ri_*R\SheafHom(\mathcal{O}_Z, K) =
R\SheafHom_{\mathcal{O}_X}(i_*\mathcal{O}_Z, K)
$$
\end{lemma}

\begin{proof}
由函子 $R\SheafHom(\mathcal{O}_Z, -)$ 的构造立即得到。
\end{proof}

\begin{lemma}
\label{duality-lemma-sheaf-with-exact-support-internal-home}
采用上述记号。对 $M \in D(\mathcal{O}_Z)$ 及
$D(\mathcal{O}_X)$ 中所有 $K$，在 $D(\mathcal{O}_Z)$ 中有
$$
R\SheafHom_{\mathcal{O}_X}(Ri_*M, K) =
Ri_*R\SheafHom_{\mathcal{O}_Z}(M, R\SheafHom(\mathcal{O}_Z, K))
$$
\end{lemma}

\begin{proof}
这由函子 $R\SheafHom(\mathcal{O}_Z, -)$ 的构造以及下述事实立即得到：
若 $\mathcal{K}^\bullet$ 是 $\mathcal{O}_X$-模的 K-内射复形，
则 $\SheafHom(\mathcal{O}_Z, \mathcal{K}^\bullet)$ 是
$\mathcal{O}_Z$-模的 K-内射复形；见《导出范畴》，引理
\ref{derived-lemma-adjoint-preserve-K-injectives}。
\end{proof}

\begin{lemma}
\label{duality-lemma-sheaf-with-exact-support-quasi-coherent}
设 $i : Z \to X$ 为概形的伪凝聚闭浸入
（若 $X$ 局部 Noetherian，则可取任意闭浸入）。则
\begin{enumerate}
\item $R\SheafHom(\mathcal{O}_Z, -)$ 把 $D^+_\QCoh(\mathcal{O}_X)$
映入 $D^+_\QCoh(\mathcal{O}_Z)$；并且
\item 若 $X = \Spec(A)$ 且 $Z = \Spec(B)$，则图
$$
\xymatrix{
D^+(B) \ar[r] & D_\QCoh^+(\mathcal{O}_Z) \\
D^+(A) \ar[r] \ar[u]^{R\Hom(B, -)} &
D_\QCoh^+(\mathcal{O}_X) \ar[u]_{R\SheafHom(\mathcal{O}_Z, -)}
}
$$
交换。
\end{enumerate}
\end{lemma}

\begin{proof}
先解释括号中的说明。若 $X$ 局部 Noetherian，则由《态射进阶》，引理
\ref{more-morphisms-lemma-Noetherian-pseudo-coherent}，$i$ 是伪凝聚的。

\medskip\noindent
设 $K$ 为 $D^+_\QCoh(\mathcal{O}_X)$ 的对象。为证明 (1)，由《态射》，
引理 \ref{morphisms-lemma-i-star-equivalence}，只需证明对
$H^n(R\SheafHom(\mathcal{O}_Z, K))$ 应用 $i_*$ 后，得到 $X$
上的拟凝聚模。由引理 \ref{duality-lemma-sheaf-with-exact-support-ext}，
这等价于证明 $R\SheafHom_{\mathcal{O}_X}(i_*\mathcal{O}_Z, K)$
属于 $D_\QCoh(\mathcal{O}_X)$。由于 $i$ 伪凝聚，层
$\mathcal{O}_Z$ 是伪凝聚 $\mathcal{O}_X$-模。因此结论由
《概形的导出范畴》，引理
\ref{perfect-lemma-quasi-coherence-internal-hom} 得到。

\medskip\noindent
如 (2) 中假设 $X = \Spec(A)$、$Z = \Spec(B)$。设 $I^\bullet$
为表示 $D^+(A)$ 中对象 $K$ 的下有界内射 $A$-模复形。
则作为 $B$-模复形，有 $R\Hom(B, K) = \Hom_A(B, I^\bullet)$。
选取拟同构
$$
\widetilde{I^\bullet} \longrightarrow \mathcal{I}^\bullet
$$
其中 $\mathcal{I}^\bullet$ 是下有界内射 $\mathcal{O}_X$-模复形。
由引理 \ref{duality-lemma-compute-sheaf-with-exact-support} 中函子
$\SheafHom(\mathcal{O}_Z, -)$ 的描述，存在映射
$$
\Hom_A(B, I^\bullet)
\longrightarrow
\Gamma(Z, \SheafHom(\mathcal{O}_Z, \mathcal{I}^\bullet))
$$
注意 $\SheafHom(\mathcal{O}_Z, \mathcal{I}^\bullet)$ 表示
$R\SheafHom(\mathcal{O}_Z, \widetilde{K})$。应用
$\widetilde{\ }$ 函子的泛性质，在 $D(\mathcal{O}_Z)$ 中得到映射
$$
\widetilde{\Hom_A(B, I^\bullet)}
\longrightarrow
R\SheafHom(\mathcal{O}_Z, \widetilde{K})
$$
可以在应用 $i_*$ 后检验此映射是 $D(\mathcal{O}_Z)$ 中的同构。
而一旦应用 $i_*$，便得到《概形的导出范畴》，引理
\ref{perfect-lemma-quasi-coherence-internal-hom} 中的同构；这里使用了
引理 \ref{duality-lemma-sheaf-with-exact-support-ext} 的识别。
\end{proof}

\begin{lemma}
\label{duality-lemma-sheaf-with-exact-support-coherent}
设 $i : Z \to X$ 为概形的闭浸入，并假设 $X$ 局部 Noetherian。
则 $R\SheafHom(\mathcal{O}_Z, -)$ 把
$D^+_{\textit{Coh}}(\mathcal{O}_X)$ 映入
$D^+_{\textit{Coh}}(\mathcal{O}_Z)$。
\end{lemma}

\begin{proof}
问题关于 $X$ 是局部的，故可假设 $X$ 仿射。记
$X = \Spec(A)$、$Z = \Spec(B)$，其中 $A$ Noetherian，且
$A \to B$ 满射。此时可应用引理
\ref{duality-lemma-sheaf-with-exact-support-quasi-coherent} 把问题转化为代数问题。
相应的代数结论是《对偶化复形》，引理
\ref{dualizing-lemma-exact-support-coherent} 的推论。
\end{proof}

\begin{lemma}
\label{duality-lemma-twisted-inverse-image-closed}
设 $X$ 为拟紧拟分离概形，$i : Z \to X$ 为伪凝聚闭浸入
（若 $X$ 为 Noetherian 概形，则任意闭浸入均为伪凝聚的）。
设 $a : D_\QCoh(\mathcal{O}_X) \to D_\QCoh(\mathcal{O}_Z)$ 为
$Ri_*$ 的右伴随。则对 $K \in D_\QCoh^+(\mathcal{O}_X)$，有函子性同构
$$
a(K) = R\SheafHom(\mathcal{O}_Z, K)
$$
\end{lemma}

\begin{proof}
括号中的陈述由《态射进阶》，引理
\ref{more-morphisms-lemma-Noetherian-pseudo-coherent} 得到。
由引理 \ref{duality-lemma-sheaf-with-exact-support-adjoint}，函子
$R\SheafHom(\mathcal{O}_Z, -)$ 是
$Ri_* : D(\mathcal{O}_Z) \to D(\mathcal{O}_X)$ 的右伴随。
此外，由引理 \ref{duality-lemma-sheaf-with-exact-support-quasi-coherent}
和引理 \ref{duality-lemma-twisted-inverse-image-bounded-below}，
$R\SheafHom(\mathcal{O}_Z, -)$ 与 $a$ 都把
$D_\QCoh^+(\mathcal{O}_X)$ 映入 $D_\QCoh^+(\mathcal{O}_Z)$。
因此，由伴随函子的唯一性得到该同构。
\end{proof}

\begin{example}
\label{duality-example-trace-closed-immersion}
若 $i : Z \to X$ 是 Noetherian 概形的闭浸入，则对
$K \in D_\QCoh^+(\mathcal{O}_X)$，图
$$
\xymatrix{
i_*a(K) \ar[rr]_-{\text{Tr}_{i, K}} \ar@{=}[d] & &
K \ar@{=}[d] \\
i_*R\SheafHom(\mathcal{O}_Z, K) \ar@{=}[r] &
R\SheafHom_{\mathcal{O}_X}(i_*\mathcal{O}_Z, K)
\ar[r] & K
}
$$
交换。这里水平等号由引理 \ref{duality-lemma-sheaf-with-exact-support-ext}，
下方水平箭头由映射 $\mathcal{O}_X \to i_*\mathcal{O}_Z$ 诱导。
图的交换性是引理 \ref{duality-lemma-twisted-inverse-image-closed} 的推论。
\end{example}




\section{闭浸入的推前右伴随与基变换}
\label{duality-section-sections-with-exact-support-base-change}

\noindent
考虑概形的笛卡尔图
$$
\xymatrix{
Z' \ar[r]_{i'} \ar[d]_g & X' \ar[d]^f \\
Z \ar[r]^i & X
}
$$
其中 $i$ 为闭浸入。若 $Z$ 与 $X'$ 在 $X$ 上 Tor 独立，则存在典范基变换映射
\begin{equation}
\label{duality-equation-base-change-exact-support}
Lg^*R\SheafHom(\mathcal{O}_Z, K)
\longrightarrow
R\SheafHom(\mathcal{O}_{Z'}, Lf^*K)
\end{equation}
它在 $D(\mathcal{O}_{Z'})$ 中关于 $D(\mathcal{O}_X)$ 的 $K$ 函子性。
具体地，由引理 \ref{duality-lemma-sheaf-with-exact-support-adjoint} 的伴随性，
这样的箭头等价于 $D(\mathcal{O}_{X'})$ 中的映射
$$
Ri'_*Lg^*R\SheafHom(\mathcal{O}_Z, K)
\longrightarrow
Lf^*K
$$
由 Tor 独立性，有 $Ri'_* \circ Lg^* = Lf^* \circ Ri_*$；
见《概形的导出范畴》，引理
\ref{perfect-lemma-compare-base-change-closed-immersion}。
因此，这又等价于映射
$$
Lf^*Ri_*R\SheafHom(\mathcal{O}_Z, K)
\longrightarrow
Lf^*K
$$
可取 $Lf^*(can)$，其中
$can : Ri_* R\SheafHom(\mathcal{O}_Z, K) \to K$ 是该伴随的余单位。

\begin{lemma}
\label{duality-lemma-check-base-change-is-iso}
在上述情形下，映射 (\ref{duality-equation-base-change-exact-support}) 是同构，
当且仅当《上同调》，注
\ref{cohomology-remark-prepare-fancy-base-change} 中的基变换映射
$$
Lf^*R\SheafHom_{\mathcal{O}_X}(\mathcal{O}_Z, K)
\longrightarrow
R\SheafHom_{\mathcal{O}_{X'}}(\mathcal{O}_{Z'}, Lf^*K)
$$
是同构。
\end{lemma}

\begin{proof}
由所假设的 Tor 独立性，$\mathcal{O}_{Z'} = Lf^*\mathcal{O}_Z$，
故此陈述有意义。由于 $i'_*$ 正合且忠实，只需证明在应用 $Ri'_*$ 后，
映射 (\ref{duality-equation-base-change-exact-support}) 是同构。由所假设的 Tor 独立性及
《概形的导出范畴》，引理
\ref{perfect-lemma-compare-base-change-closed-immersion}，有
$Ri'_* \circ Lg^* = Lf^* \circ Ri_*$，从而得到映射
$$
Lf^*Ri_*R\SheafHom(\mathcal{O}_Z, K)
\longrightarrow
Ri'_*R\SheafHom(\mathcal{O}_{Z'}, Lf^*K)
$$
由引理 \ref{duality-lemma-sheaf-with-exact-support-ext}，其源和目标正如本引理陈述中所示。
这里略去验证它与《上同调》，注
\ref{cohomology-remark-prepare-fancy-base-change} 中构造的是同一个映射。
\end{proof}

\begin{lemma}
\label{duality-lemma-flat-bc-sheaf-with-exact-support}
在上述情形中，假设 $f$ 平坦且 $i$ 伪凝聚。则对
$D^+_\QCoh(\mathcal{O}_X)$ 中的 $K$，
(\ref{duality-equation-base-change-exact-support}) 是同构。
\end{lemma}

\begin{proof}
第一种证明。为证明此映射是同构，可以局部地工作。因此可假设
$X$、$X'$、$Z$、$Z'$ 均为仿射概形，分别对应于环
$A$、$A'$、$B$、$B'$。于是 $B$ 与 $A'$ 在 $A$ 上 Tor 独立。
由引理 \ref{duality-lemma-check-base-change-is-iso}，只需验证
$$
R\Hom_A(B, K) \otimes_A^\mathbf{L} A' =
R\Hom_{A'}(B', K \otimes_A^\mathbf{L} A')
$$
对所有 $K \in D^+(A)$ 在 $D(A')$ 中成立。这里使用《概形的导出范畴》，
引理 \ref{perfect-lemma-quasi-coherence-internal-hom}，以及
$B$（相应地\ $B'$）作为 $A$-模（相应地\ $A'$-模）是伪凝聚的事实，
以便在环和概形两个层面比较导出 Hom。上式由《代数进阶》，引理
\ref{more-algebra-lemma-internal-hom-evaluate-tensor-isomorphism}
第 (3) 部分得到。另见《对偶化复形》，节
\ref{dualizing-section-base-change-trivial-duality} 中的讨论。

\medskip\noindent
第二种证明\footnote{此证明表明，只需假设 $K$ 属于
$D^+(\mathcal{O}_X)$。}。取 $z' \in Z'$，其像为 $z \in Z$。
先证明 (\ref{duality-equation-base-change-exact-support}) 在 $z'$ 处的茎上诱导映射
$$
R\Hom(\mathcal{O}_{Z, z}, K_z)
\otimes_{\mathcal{O}_{Z, x}}^\mathbf{L} \mathcal{O}_{Z', z'}
\longrightarrow
R\Hom(\mathcal{O}_{Z', z'},
K_z \otimes_{\mathcal{O}_{X, z}}^\mathbf{L} \mathcal{O}_{X', z'})
$$
即《对偶化复形》，方程 (\ref{dualizing-equation-base-change}) 中的映射。
事实上，这些映射的构造完全相同。随后应用《对偶化复形》，引理
\ref{dualizing-lemma-flat-bc-surjection}。
\end{proof}

\begin{lemma}
\label{duality-lemma-sheaf-with-exact-support-tensor}
设 $i : Z \to X$ 为概形的伪凝聚闭浸入。设
$M \in D_\QCoh(\mathcal{O}_X)$ 局部具有位于 $[a, \infty)$ 中的
Tor 振幅，并设 $K \in D_\QCoh^+(\mathcal{O}_X)$。则在
$D(\mathcal{O}_Z)$ 中存在典范同构
$$
R\SheafHom(\mathcal{O}_Z, K) \otimes_{\mathcal{O}_Z}^\mathbf{L} Li^*M =
R\SheafHom(\mathcal{O}_Z, K \otimes_{\mathcal{O}_X}^\mathbf{L} M)
$$
\end{lemma}

\begin{proof}
由引理 \ref{duality-lemma-sheaf-with-exact-support-adjoint} 和
\ref{duality-lemma-sheaf-with-exact-support-ext}，从左边到右边的映射等价于映射
$$
Ri_*R\SheafHom(\mathcal{O}_Z, K) \otimes_{\mathcal{O}_X}^\mathbf{L} M
\longrightarrow
K \otimes_{\mathcal{O}_X}^\mathbf{L} M
$$
对此映射，取余单位 $Ri_*R\SheafHom(\mathcal{O}_Z, K) \to K$
与 $\text{id}_M$ 的张量积。为证明在所给假设下此映射是同构，
可用引理 \ref{duality-lemma-sheaf-with-exact-support-quasi-coherent} 将问题转化为代数，
再例如应用《代数进阶》，引理
\ref{more-algebra-lemma-internal-hom-evaluate-tensor-isomorphism} 第 (3) 部分。
也可以不使用引理 \ref{duality-lemma-sheaf-with-exact-support-quasi-coherent}，
而像引理 \ref{duality-lemma-flat-bc-sheaf-with-exact-support} 的第二种证明那样考察茎。
\end{proof}



\section{有限态射的推前右伴随}
\label{duality-section-duality-finite}

\noindent
若 $i : Z \to X$ 为概形的闭浸入，则函子
$i_* : \textit{Mod}(\mathcal{O}_Z) \to \textit{Mod}(\mathcal{O}_X)$
有右伴随 $\SheafHom(\mathcal{O}_Z, -)$，其导出延拓
$R\SheafHom(\mathcal{O}_Z, -)$ 是
$Ri_* : D(\mathcal{O}_Z) \to D(\mathcal{O}_X)$ 的右伴随；见节
\ref{duality-section-sections-with-exact-support}。对于有限态射
$f : Y \to X$，这种策略不可行，因为函子
$f_* : \textit{Mod}(\mathcal{O}_Y) \to \textit{Mod}(\mathcal{O}_X)$
一般并不正合，因而没有右伴随。替代办法是考虑正合函子
$\textit{Mod}(f_*\mathcal{O}_Y) \to \textit{Mod}(\mathcal{O}_X)$，
再考虑相应的右伴随及其导出延拓。

\medskip\noindent
设 $f : Y \to X$ 为概形的仿射态射。对一个
$\mathcal{O}_X$-模层 $\mathcal{F}$，层
$$
\SheafHom_{\mathcal{O}_X}(f_*\mathcal{O}_Y, \mathcal{F})
$$
是一个 $f_*\mathcal{O}_Y$-模层。由此得到函子
$\textit{Mod}(\mathcal{O}_X) \to \textit{Mod}(f_*\mathcal{O}_Y)$，
记作 $\SheafHom(f_*\mathcal{O}_Y, -)$。

\begin{lemma}
\label{duality-lemma-compute-sheafhom-affine}
采用上述记号。函子 $\SheafHom(f_*\mathcal{O}_Y, -)$ 是限制函子
$\textit{Mod}(f_*\mathcal{O}_Y) \to \textit{Mod}(\mathcal{O}_X)$
的右伴随。对仿射开集 $U \subset X$，有
$$
\Gamma(U, \SheafHom(f_*\mathcal{O}_Y, \mathcal{F})) =
\Hom_A(B, \mathcal{F}(U))
$$
其中 $A = \mathcal{O}_X(U)$，且 $B = \mathcal{O}_Y(f^{-1}(U))$。
\end{lemma}

\begin{proof}
伴随性由《模》，引理 \ref{modules-lemma-adjoint-tensor-restrict} 得到。
由于 $f$ 仿射，$f_*\mathcal{O}_Y$ 是由视为 $A$-模的 $B$
所对应的拟凝聚层。因此，$U$ 上截面的描述由《概形》，引理
\ref{schemes-lemma-compare-constructions} 得到。
\end{proof}

\noindent
函子 $\SheafHom(f_*\mathcal{O}_Y, -)$ 左正合。令
$$
R\SheafHom(f_*\mathcal{O}_Y, -) :
D(\mathcal{O}_X)
\longrightarrow
D(f_*\mathcal{O}_Y)
$$
为其导出延拓。

\begin{lemma}
\label{duality-lemma-sheafhom-affine-adjoint}
采用上述记号。函子 $R\SheafHom(f_*\mathcal{O}_Y, -)$ 是函子
$D(f_*\mathcal{O}_Y) \to D(\mathcal{O}_X)$ 的右伴随。
\end{lemma}

\begin{proof}
由引理 \ref{duality-lemma-compute-sheafhom-affine} 和《导出范畴》，引理
\ref{derived-lemma-derived-adjoint-functors} 得到。
\end{proof}

\begin{lemma}
\label{duality-lemma-sheafhom-affine-ext}
采用上述记号。复合函子
$$
D(\mathcal{O}_X) \xrightarrow{R\SheafHom(f_*\mathcal{O}_Y, -)}
D(f_*\mathcal{O}_Y) \to D(\mathcal{O}_X)
$$
即函子 $K \mapsto R\SheafHom_{\mathcal{O}_X}(f_*\mathcal{O}_Y, K)$。
\end{lemma}

\begin{proof}
这由构造立即得到。
\end{proof}

\begin{lemma}
\label{duality-lemma-finite-twisted}
设 $f : Y \to X$ 为概形的有限伪凝聚态射
（Noetherian 概形之间的有限态射是伪凝聚的）。函子
$R\SheafHom(f_*\mathcal{O}_Y, -)$ 把 $D_\QCoh^+(\mathcal{O}_X)$
映入 $D_\QCoh^+(f_*\mathcal{O}_Y)$。若 $X$ 拟紧且拟分离，则图
$$
\xymatrix{
D_\QCoh^+(\mathcal{O}_X) \ar[rr]_a \ar[rd]_{R\SheafHom(f_*\mathcal{O}_Y, -)}
& & D_\QCoh^+(\mathcal{O}_Y) \ar[ld]^\Phi \\
& D_\QCoh^+(f_*\mathcal{O}_Y)
}
$$
交换，其中 $a$ 是引理 \ref{duality-lemma-twisted-inverse-image} 中对应于
$f$ 的右伴随，而 $\Phi$ 是《概形的导出范畴》，引理
\ref{perfect-lemma-affine-morphism-equivalence} 中的等价。
\end{lemma}

\begin{proof}
括号中的说明由《态射进阶》，引理
\ref{more-morphisms-lemma-Noetherian-pseudo-coherent} 得到。
由于 $f$ 伪凝聚，$\mathcal{O}_X$-模 $f_*\mathcal{O}_Y$ 伪凝聚；
见《态射进阶》，引理
\ref{more-morphisms-lemma-finite-pseudo-coherent}。因此
$R\SheafHom(f_*\mathcal{O}_Y, -)$ 把 $D_\QCoh^+(\mathcal{O}_X)$
映入 $D_\QCoh^+(f_*\mathcal{O}_Y)$；见《概形的导出范畴》，引理
\ref{perfect-lemma-quasi-coherence-internal-hom}。
于是 $\Phi \circ a$ 与 $R\SheafHom(f_*\mathcal{O}_Y, -)$ 在
$D_\QCoh^+(\mathcal{O}_X)$ 上相同，因为这两个函子都是限制函子
$D_\QCoh^+(f_*\mathcal{O}_Y) \to D_\QCoh^+(\mathcal{O}_X)$ 的右伴随。
为此使用引理 \ref{duality-lemma-twisted-inverse-image-bounded-below} 和
\ref{duality-lemma-sheafhom-affine-adjoint}。
\end{proof}

\begin{remark}
\label{duality-remark-trace-map-finite}
若 $f : Y \to X$ 是 Noetherian 概形之间的有限态射，则对
$K \in D_\QCoh^+(\mathcal{O}_X)$，图
$$
\xymatrix{
Rf_*a(K) \ar[r]_-{\text{Tr}_{f, K}} \ar@{=}[d] & K \ar@{=}[d] \\
R\SheafHom_{\mathcal{O}_X}(f_*\mathcal{O}_Y, K) \ar[r] & K
}
$$
交换。这由引理 \ref{duality-lemma-finite-twisted} 得到。下方水平箭头由映射
$\mathcal{O}_X \to f_*\mathcal{O}_Y$ 诱导，上方水平箭头是节
\ref{duality-section-trace} 中讨论的迹映射。
\end{remark}






\section{真平坦态射的推前右伴随}
\label{duality-section-proper-flat}

\noindent
对于拟紧拟分离概形之间的真、平坦且有限表示的态射，
推前的右伴随具有若干显著性质。

\begin{lemma}
\label{duality-lemma-proper-flat}
设 $Y$ 为拟紧拟分离概形，$f : X \to Y$ 为真、平坦且有限表示的概形态射。
设 $a$ 为引理 \ref{duality-lemma-twisted-inverse-image} 中
$Rf_* : D_\QCoh(\mathcal{O}_X) \to D_\QCoh(\mathcal{O}_Y)$ 的右伴随。
则 $a$ 与直和交换。
\end{lemma}

\begin{proof}
设 $P$ 为 $D(\mathcal{O}_X)$ 的完美对象。由《概形的导出范畴》，引理
\ref{perfect-lemma-flat-proper-perfect-direct-image-general}，
复形 $Rf_*P$ 在 $Y$ 上完美。设 $K_i$ 为
$D_\QCoh(\mathcal{O}_Y)$ 中的一族对象。于是
\begin{align*}
\Hom_{D(\mathcal{O}_X)}(P, a(\bigoplus K_i))
& =
\Hom_{D(\mathcal{O}_Y)}(Rf_*P, \bigoplus K_i) \\
& =
\bigoplus \Hom_{D(\mathcal{O}_Y)}(Rf_*P, K_i) \\
& =
\bigoplus \Hom_{D(\mathcal{O}_X)}(P, a(K_i))
\end{align*}
因为完美对象是紧对象（《概形的导出范畴》，命题
\ref{perfect-proposition-compact-is-perfect}）。由于
$D_\QCoh(\mathcal{O}_X)$ 有完美生成元（《概形的导出范畴》，定理
\ref{perfect-theorem-bondal-van-den-Bergh}），可知映射
$\bigoplus a(K_i) \to a(\bigoplus K_i)$ 是同构，即 $a$ 与直和交换。
\end{proof}

\begin{lemma}
\label{duality-lemma-proper-flat-relative}
设 $Y$ 为拟紧拟分离概形，$f : X \to Y$ 为真、平坦且有限表示的概形态射。
设 $a$ 为引理 \ref{duality-lemma-twisted-inverse-image} 中
$Rf_* : D_\QCoh(\mathcal{O}_X) \to D_\QCoh(\mathcal{O}_Y)$ 的右伴随。则
\begin{enumerate}
\item 对每个闭集 $T \subset Y$，若 $Q \in D_\QCoh(Y)$ 支撑于 $T$，
则 $a(Q)$ 支撑于 $f^{-1}(T)$；
\item 对每个拟紧开集 $V \subset Y$ 及任意
$K \in D_\QCoh(\mathcal{O}_Y)$，映射 (\ref{duality-equation-sheafy}) 是同构。
\end{enumerate}
\end{lemma}

\begin{proof}
这由引理 \ref{duality-lemma-when-sheafy}、\ref{duality-lemma-proper-noetherian}
和 \ref{duality-lemma-proper-flat} 得到。严格说来，(1) 首先只对补集为拟紧开集的
$T$ 成立；但一般情形也随之成立，因为任意闭子集都是此类闭集的交。
\end{proof}

\begin{lemma}
\label{duality-lemma-compare-with-pullback-flat-proper}
设 $Y$ 为拟紧拟分离概形，$f : X \to Y$ 为真、平坦且有限表示的概形态射。
则对 $D_\QCoh(\mathcal{O}_Y)$ 的每个对象 $K$，映射
(\ref{duality-equation-compare-with-pullback}) 都是同构。
\end{lemma}

\begin{proof}
由引理 \ref{duality-lemma-proper-flat}，$a$ 与直和交换。因此，
$D_\QCoh(\mathcal{O}_Y)$ 中使 (\ref{duality-equation-compare-with-pullback})
成为同构的对象构成 $D_\QCoh(\mathcal{O}_Y)$ 的严格满、饱和三角子范畴，
且该子范畴在取直和下封闭。由于 $D_\QCoh(\mathcal{O}_Y)$ 是由单个完美对象
生成的模范畴（分别见《概形的导出范畴》，定理
\ref{perfect-theorem-DQCoh-is-Ddga} 和定理
\ref{perfect-theorem-bondal-van-den-Bergh}），可如《代数进阶》，注
\ref{more-algebra-remark-P-resolution} 那样论证，说明只需对一个完美对象证明
(\ref{duality-equation-compare-with-pullback}) 是同构。但该结论对完美对象成立；见引理
\ref{duality-lemma-compare-with-pullback-perfect}。
\end{proof}

\noindent
下述引理表明，对有限表示的真平坦态射，基变换映射
(\ref{duality-equation-base-change-map}) 是同构。例
\ref{duality-example-base-change-wrong} 将说明，对完美真态射，这一结论不再总是成立；
此时必须加入 Tor 独立性条件。

\begin{lemma}
\label{duality-lemma-proper-flat-base-change}
设 $g : Y' \to Y$ 为拟紧拟分离概形之间的态射，
$f : X \to Y$ 为有限表示的真平坦态射。则对所有
$K \in D_\QCoh(\mathcal{O}_Y)$，基变换映射
(\ref{duality-equation-base-change-map}) 是同构。
\end{lemma}

\begin{proof}
由引理 \ref{duality-lemma-proper-flat-relative}，函子 $a$ 和 $a'$ 的构造
与到 $Y$ 和 $Y'$ 的开集的限制交换。因此可假设
$Y' \to Y$ 是仿射概形之间的态射；见注
\ref{duality-remark-check-over-affines}。此时结论由引理
\ref{duality-lemma-more-base-change} 得到。
\end{proof}

\begin{remark}
\label{duality-remark-relative-dualizing-complex}
设 $Y$ 为拟紧拟分离概形，$f : X \to Y$ 为有限表示的真平坦态射。
设 $a$ 为引理 \ref{duality-lemma-twisted-inverse-image} 中对应于 $f$ 的伴随。
在此情形下，$\omega_{X/Y}^\bullet = a(\mathcal{O}_Y)$ 有时称为
{\it 相对对偶化复形}。由引理
\ref{duality-lemma-compare-with-pullback-flat-proper}，对
$K \in D_\QCoh(\mathcal{O}_Y)$ 有函子性同构
$a(K) = Lf^*K \otimes_{\mathcal{O}_X}^\mathbf{L} \omega_{X/Y}^\bullet$。
此外，节 \ref{duality-section-trace} 中的迹映射
$$
\text{Tr}_{f, \mathcal{O}_Y} : Rf_*\omega_{X/Y}^\bullet \to \mathcal{O}_Y
$$
诱导 $D_\QCoh(\mathcal{O}_Y)$ 中所有 $K$ 的迹映射。更确切地，图
$$
\xymatrix{
Rf_*a(K) \ar[rrr]_{\text{Tr}_{f, K}} \ar@{=}[d] & & &
K \ar@{=}[d] \\
Rf_*(Lf^*K \otimes_{\mathcal{O}_X}^\mathbf{L} \omega_{X/Y}^\bullet)
\ar@{=}[r] &
K \otimes_{\mathcal{O}_Y}^\mathbf{L} Rf_*\omega_{X/Y}^\bullet
\ar[rr]^-{\text{id}_K \otimes \text{Tr}_{f, \mathcal{O}_Y}} & & K
}
$$
交换，其中下方右侧的等式由《概形的导出范畴》，引理
\ref{perfect-lemma-cohomology-base-change} 给出。若
$g : Y' \to Y$ 是拟紧拟分离概形之间的态射，且
$X' = Y' \times_Y X$，则由引理 \ref{duality-lemma-proper-flat-base-change}，有
$\omega_{X'/Y'}^\bullet = L(g')^*\omega_{X/Y}^\bullet$，其中
$g' : X' \to X$ 为投影；并且由引理
\ref{duality-lemma-trace-map-and-base-change}，对 $f' : X' \to Y'$ 的迹映射
$$
\text{Tr}_{f', \mathcal{O}_{Y'}} :
Rf'_*\omega_{X'/Y'}^\bullet \to \mathcal{O}_{Y'}
$$
是在基变换同构下由 $\text{Tr}_{f, \mathcal{O}_Y}$ 基变换得到的。
\end{remark}

\begin{remark}
\label{duality-remark-relative-dualizing-complex-relative-cup-product}
设 $f : X \to Y$、$\omega^\bullet_{X/Y}$ 和
$\text{Tr}_{f, \mathcal{O}_Y}$ 如注
\ref{duality-remark-relative-dualizing-complex} 中所述。设 $K$ 和 $M$ 属于
$D_\QCoh(\mathcal{O}_X)$，且 $M$ 伪凝聚（例如完美）。假设给定映射
$K \otimes_{\mathcal{O}_X}^\mathbf{L} M \to \omega^\bullet_{X/Y}$，
它通过《上同调》，方程 (\ref{cohomology-equation-internal-hom}) 对应于同构
$K \to R\SheafHom_{\mathcal{O}_X}(M, \omega^\bullet_{X/Y})$。
则相对杯积（《上同调》，注 \ref{cohomology-remark-cup-product}）
$$
Rf_*K \otimes_{\mathcal{O}_Y}^\mathbf{L} Rf_*M
\to
Rf_*(K \otimes_{\mathcal{O}_X}^\mathbf{L} M)
\to
Rf_*\omega^\bullet_{X/Y}
\xrightarrow{\text{Tr}_{f, \mathcal{O}_Y}}
\mathcal{O}_Y
$$
确定一个同构
$Rf_*K \to R\SheafHom_{\mathcal{O}_Y}(Rf_*M, \mathcal{O}_Y)$。
事实上，由于 $\omega^\bullet_{X/Y} = a(\mathcal{O}_Y)$，典范映射
(\ref{duality-equation-sheafy-trace})
$$
Rf_*R\SheafHom_{\mathcal{O}_X}(M, \omega^\bullet_{X/Y}) \to
R\SheafHom_{\mathcal{O}_Y}(Rf_*M, \mathcal{O}_Y)
$$
由引理 \ref{duality-lemma-iso-on-RSheafHom}、注
\ref{duality-remark-iso-on-RSheafHom} 以及 $M$ 与 $Rf_*M$ 伪凝聚这一事实而为同构；
见《概形的导出范畴》，引理
\ref{perfect-lemma-flat-proper-pseudo-coherent-direct-image-general}。
要看出相对杯积诱导的正是此同构，使用《上同调》，注
\ref{cohomology-remark-relative-cup-and-composition} 中图的交换性。
\end{remark}

\begin{lemma}
\label{duality-lemma-properties-relative-dualizing}
设 $Y$ 为拟紧拟分离概形，$f : X \to Y$ 为真、平坦且有限表示的概形态射，
其相对对偶化复形为 $\omega_{X/Y}^\bullet$
（注 \ref{duality-remark-relative-dualizing-complex}）。则
\begin{enumerate}
\item $\omega_{X/Y}^\bullet$ 是 $D(\mathcal{O}_X)$ 中的 $Y$-完美对象；
\item $Rf_*\omega_{X/Y}^\bullet$ 的正次数上同调层均为零；
\item $\mathcal{O}_X \to
R\SheafHom_{\mathcal{O}_X}(\omega_{X/Y}^\bullet, \omega_{X/Y}^\bullet)$
是同构。
\end{enumerate}
\end{lemma}

\begin{proof}
由于 $\omega_{X/Y}^\bullet$ 的构造与基变换交换（见注
\ref{duality-remark-relative-dualizing-complex}），可以并且确实假设 $Y$ 仿射。
对 $D(\mathcal{O}_X)$ 中的完美对象 $E$，有
\begin{align*}
Rf_*(E \otimes_{\mathcal{O}_X}^\mathbf{L} \omega_{X/Y}^\bullet)
& =
Rf_*R\SheafHom_{\mathcal{O}_X}(E^\vee, \omega_{X/Y}^\bullet) \\
& =
R\SheafHom_{\mathcal{O}_Y}(Rf_*E^\vee, \mathcal{O}_Y) \\
& =
(Rf_*E^\vee)^\vee
\end{align*}
第一个等式见《上同调》，引理
\ref{cohomology-lemma-dual-perfect-complex}。第二个等式见引理
\ref{duality-lemma-iso-on-RSheafHom}、注 \ref{duality-remark-iso-on-RSheafHom} 以及
《概形的导出范畴》，引理
\ref{perfect-lemma-flat-proper-perfect-direct-image-general}。
第三个等式是对偶的定义。尤其，这些引文也表明所得对象是
$D(\mathcal{O}_Y)$ 中的完美对象。由《态射进阶》，引理
\ref{more-morphisms-lemma-characterize-relatively-perfect}，可知
$\omega_{X/Y}^\bullet$ 是 $Y$-完美的。这证明了 (1)。

\medskip\noindent
设 $M$ 为 $D_\QCoh(\mathcal{O}_Y)$ 的对象。则
\begin{align*}
\Hom_Y(M, Rf_*\omega_{X/Y}^\bullet) & =
\Hom_X(Lf^*M, \omega_{X/Y}^\bullet) \\
& =
\Hom_Y(Rf_*Lf^*M, \mathcal{O}_Y) \\
& =
\Hom_Y(M \otimes_{\mathcal{O}_Y}^\mathbf{L} Rf_*\mathcal{O}_X, \mathcal{O}_Y)
\end{align*}
第一个等式由《上同调》，引理 \ref{cohomology-lemma-adjoint}。
第二个等式由 $a$ 的构造。第三个等式由《概形的导出范畴》，引理
\ref{perfect-lemma-cohomology-base-change}。回忆
$Rf_*\mathcal{O}_X$ 是 Tor 振幅位于 $[0, N]$ 的完美对象，其中
$N$ 为某个整数；见《概形的导出范畴》，引理
\ref{perfect-lemma-flat-proper-perfect-direct-image-general}。
因此可用位于次数 $[0, N]$ 的有限投射模复形表示
$Rf_*\mathcal{O}_X$（使用《代数进阶》，引理
\ref{more-algebra-lemma-perfect} 以及 $Y$ 仿射这一事实）。
故若对某个 $i > 0$ 有 $M = \mathcal{O}_Y[-i]$，则最后一个群为零。
由于 $Y$ 仿射，可知对 $i > 0$，
$H^i(Rf_*\omega_{X/Y}^\bullet) = 0$。这证明了 (2)。

\medskip\noindent
设 $E$ 为 $D_\QCoh(\mathcal{O}_X)$ 的完美对象。则有
\begin{align*}
\Hom_X(E, R\SheafHom_{\mathcal{O}_X}(\omega_{X/Y}^\bullet, \omega_{X/Y}^\bullet)
& =
\Hom_X(E \otimes_{\mathcal{O}_X}^\mathbf{L} \omega_{X/Y}^\bullet,
\omega_{X/Y}^\bullet) \\
& =
\Hom_Y(Rf_*(E \otimes_{\mathcal{O}_X}^\mathbf{L} \omega_{X/Y}^\bullet),
\mathcal{O}_Y) \\
& =
\Hom_Y(Rf_*(R\SheafHom_{\mathcal{O}_X}(E^\vee, \omega_{X/Y}^\bullet)),
\mathcal{O}_Y) \\
& =
\Hom_Y(R\SheafHom_{\mathcal{O}_Y}(Rf_*E^\vee, \mathcal{O}_Y),
\mathcal{O}_Y) \\
& =
R\Gamma(Y, Rf_*E^\vee) \\
& =
\Hom_X(E, \mathcal{O}_X)
\end{align*}
第一个等式由《上同调》，引理
\ref{cohomology-lemma-internal-hom}。第二个等式是
$\omega_{X/Y}^\bullet$ 的定义。第三个等式来自对偶完美复形
$E^\vee$ 的构造；见《上同调》，引理
\ref{cohomology-lemma-dual-perfect-complex}。第四个等式来自证明第一段中所示的等式
$Rf_*R\SheafHom_{\mathcal{O}_X}(E^\vee, \omega_{X/Y}^\bullet) =
R\SheafHom_{\mathcal{O}_Y}(Rf_*E^\vee, \mathcal{O}_Y)$。
第五个等式由完美复形的双重对偶性（《上同调》，引理
\ref{cohomology-lemma-dual-perfect-complex}）以及 $Rf_*E$ 完美这一事实成立；
后者见《概形的导出范畴》，引理
\ref{perfect-lemma-flat-proper-perfect-direct-image-general}。
最后一个等式是 $f$ 的 Leray 等式。这一串等式本质上由 Yoneda 引理表明 (3)
成立。具体地，由《概形的导出范畴》，引理
\ref{perfect-lemma-quasi-coherence-internal-hom}，对象
$R\SheafHom(\omega_{X/Y}^\bullet, \omega_{X/Y}^\bullet)$ 属于
$D_\QCoh(\mathcal{O}_X)$。在上式中取 $E = \mathcal{O}_X$，得到映射
$\alpha : \mathcal{O}_X \to
R\SheafHom_{\mathcal{O}_X}(\omega_{X/Y}^\bullet, \omega_{X/Y}^\bullet)$，
它对应于 $\text{id}_{\mathcal{O}_X} \in \Hom_X(\mathcal{O}_X, \mathcal{O}_X)$。
由于上述所有同构都关于 $E$ 函子性，可知 $\alpha$ 的锥是
$D_\QCoh(\mathcal{O}_X)$ 的一个对象 $C$，且对所有完美对象 $E$ 都有
$\Hom(E, C) = 0$。由于完美对象生成该范畴（《概形的导出范畴》，定理
\ref{perfect-theorem-bondal-van-den-Bergh}），可知 $\alpha$ 是同构。
\end{proof}

\begin{lemma}[刚性]
\label{duality-lemma-van-den-bergh}
设 $Y$ 为拟紧拟分离概形，$f : X \to Y$ 为有限表示的真平坦态射，
其相对对偶化复形为 $\omega_{X/Y}^\bullet$
（注 \ref{duality-remark-relative-dualizing-complex}）。存在典范同构
\begin{equation}
\label{duality-equation-pre-rigid}
\mathcal{O}_X =
c(L\text{pr}_1^*\omega_{X/Y}^\bullet) =
c(L\text{pr}_2^*\omega_{X/Y}^\bullet)
\end{equation}
以及典范同构
\begin{equation}
\label{duality-equation-rigid}
\omega_{X/Y}^\bullet =
c\left(L\text{pr}_1^*\omega_{X/Y}^\bullet
\otimes_{\mathcal{O}_{X \times_Y X}}^\mathbf{L}
L\text{pr}_2^*\omega_{X/Y}^\bullet\right)
\end{equation}
其中 $c$ 是引理 \ref{duality-lemma-twisted-inverse-image} 中对角态射
$\Delta : X \to X \times_Y X$ 的右伴随。
\end{lemma}

\begin{proof}
设 $a$ 为引理 \ref{duality-lemma-twisted-inverse-image} 中 $Rf_*$ 的右伴随。
考虑笛卡尔方块
$$
\xymatrix{
X \times_Y X \ar[r]_q \ar[d]_p & X \ar[d]_f \\
X \ar[r]^f & Y
}
$$
设 $b$ 为引理 \ref{duality-lemma-twisted-inverse-image} 中对应于 $p$ 的右伴随。则
\begin{align*}
\omega_{X/Y}^\bullet
& =
c(b(\omega_{X/Y}^\bullet)) \\
& =
c(Lp^*\omega_{X/Y}^\bullet
\otimes_{\mathcal{O}_{X \times_Y X}}^\mathbf{L} b(\mathcal{O}_X)) \\
& =
c(Lp^*\omega_{X/Y}^\bullet
\otimes_{\mathcal{O}_{X \times_Y X}}^\mathbf{L}
Lq^*a(\mathcal{O}_Y)) \\
& =
c(Lp^*\omega_{X/Y}^\bullet
\otimes_{\mathcal{O}_{X \times_Y X}}^\mathbf{L}
Lq^*\omega_{X/Y}^\bullet)
\end{align*}
如 (\ref{duality-equation-rigid}) 所示。解释如下：
\begin{enumerate}
\item 第一个等式成立，因为由 $\text{id}_X = p \circ \Delta$ 有
$\text{id} = c \circ b$。
\item 第二个等式由引理
\ref{duality-lemma-compare-with-pullback-flat-proper} 成立。
\item 第三个等式由引理 \ref{duality-lemma-proper-flat-base-change} 以及
$\mathcal{O}_X = Lf^*\mathcal{O}_Y$ 成立。
\item 第四个等式成立，因为 $\omega_{X/Y}^\bullet = a(\mathcal{O}_Y)$。
\end{enumerate}
方程 (\ref{duality-equation-pre-rigid}) 的证明完全相同。
\end{proof}

\begin{remark}
\label{duality-remark-van-den-bergh}
引理 \ref{duality-lemma-van-den-bergh} 表明，我们的相对对偶化复形在与
\cite{vdB-rigid} 所引入概念类似的意义下是{\it 刚性的}。
具体地，(\ref{duality-equation-rigid}) 右侧的函子关于
$\omega_{X/Y}^\bullet$ 是“二次的”，而 (\ref{duality-equation-rigid})
左侧的函子是“线性的”，
因而这在一定程度上“确定”了复形 $\omega_{X/Y}^\bullet$。
对偶理论还有一种使用“刚性”相对对偶化复形的处理方式；例如见
\cite{Neeman-rigid}、\cite{Yekutieli-rigid} 和 \cite{Yekutieli-Zhang}。
我们将在节 \ref{duality-section-relative-dualizing-complexes} 中回到这一点。
\end{remark}





\section{完美真态射的推前右伴随}
\label{duality-section-perfect-proper}

\noindent
本节恰当的一般性应当是考虑拟紧拟分离概形之间的完美真态射；见
\cite{LN}。

\begin{lemma}
\label{duality-lemma-proper-flat-noetherian}
设 $f : X \to Y$ 为 Noetherian 概形之间的完美真态射。
设 $a$ 为引理 \ref{duality-lemma-twisted-inverse-image} 中
$Rf_* : D_\QCoh(\mathcal{O}_X) \to D_\QCoh(\mathcal{O}_Y)$ 的右伴随。
则 $a$ 与直和交换。
\end{lemma}

\begin{proof}
设 $P$ 为 $D(\mathcal{O}_X)$ 的完美对象。由《态射进阶》，引理
\ref{more-morphisms-lemma-perfect-proper-perfect-direct-image}，
复形 $Rf_*P$ 在 $Y$ 上完美。设 $K_i$ 为
$D_\QCoh(\mathcal{O}_Y)$ 中的一族对象。则
\begin{align*}
\Hom_{D(\mathcal{O}_X)}(P, a(\bigoplus K_i))
& =
\Hom_{D(\mathcal{O}_Y)}(Rf_*P, \bigoplus K_i) \\
& =
\bigoplus \Hom_{D(\mathcal{O}_Y)}(Rf_*P, K_i) \\
& =
\bigoplus \Hom_{D(\mathcal{O}_X)}(P, a(K_i))
\end{align*}
因为完美对象是紧对象（《概形的导出范畴》，命题
\ref{perfect-proposition-compact-is-perfect}）。由于
$D_\QCoh(\mathcal{O}_X)$ 有完美生成元（《概形的导出范畴》，定理
\ref{perfect-theorem-bondal-van-den-Bergh}），可知映射
$\bigoplus a(K_i) \to a(\bigoplus K_i)$ 是同构，即 $a$ 与直和交换。
\end{proof}

\begin{lemma}
\label{duality-lemma-proper-flat-noetherian-relative}
设 $f : X \to Y$ 为 Noetherian 概形之间的完美真态射。
设 $a$ 为引理 \ref{duality-lemma-twisted-inverse-image} 中
$Rf_* : D_\QCoh(\mathcal{O}_X) \to D_\QCoh(\mathcal{O}_Y)$ 的右伴随。则
\begin{enumerate}
\item 对每个闭集 $T \subset Y$，若 $Q \in D_\QCoh(Y)$ 支撑于 $T$，
则 $a(Q)$ 支撑于 $f^{-1}(T)$；
\item 对每个开集 $V \subset Y$ 及任意 $K \in D_\QCoh(\mathcal{O}_Y)$，
映射 (\ref{duality-equation-sheafy}) 是同构；并且
\end{enumerate}
\end{lemma}

\begin{proof}
这由引理 \ref{duality-lemma-when-sheafy}、\ref{duality-lemma-proper-noetherian}
和 \ref{duality-lemma-proper-flat-noetherian} 得到。
\end{proof}

\begin{lemma}
\label{duality-lemma-compare-with-pullback-flat-proper-noetherian}
设 $f : X \to Y$ 为 Noetherian 概形之间的完美真态射。则对
$D_\QCoh(\mathcal{O}_Y)$ 的每个对象 $K$，映射
(\ref{duality-equation-compare-with-pullback}) 都是同构。
\end{lemma}

\begin{proof}
由引理 \ref{duality-lemma-proper-flat-noetherian}，$a$ 与直和交换。因此，
$D_\QCoh(\mathcal{O}_Y)$ 中使 (\ref{duality-equation-compare-with-pullback})
成为同构的对象构成 $D_\QCoh(\mathcal{O}_Y)$ 的严格满、饱和三角子范畴，
并且该子范畴在取直和下封闭。由于 $D_\QCoh(\mathcal{O}_Y)$ 是由单个
完美对象生成的模范畴（分别见《概形的导出范畴》，定理
\ref{perfect-theorem-DQCoh-is-Ddga} 和定理
\ref{perfect-theorem-bondal-van-den-Bergh}），可如《代数进阶》，注
\ref{more-algebra-remark-P-resolution} 那样论证，说明只需对一个完美对象证明
(\ref{duality-equation-compare-with-pullback}) 是同构。但该结论对完美对象成立；
见引理 \ref{duality-lemma-compare-with-pullback-perfect}。
\end{proof}

\begin{lemma}
\label{duality-lemma-proper-perfect-base-change}
设 $f : X \to Y$ 为 Noetherian 概形之间的完美真态射。
设 $g : Y' \to Y$ 为态射，且 $Y'$ Noetherian。若 $X$ 与
$Y'$ 在 $Y$ 上 Tor 独立，则对所有 $K \in D_\QCoh(\mathcal{O}_Y)$，
基变换映射 (\ref{duality-equation-base-change-map}) 是同构。
\end{lemma}

\begin{proof}
由引理 \ref{duality-lemma-proper-flat-noetherian-relative}，函子 $a$ 和 $a'$
的构造与到 $Y$ 和 $Y'$ 的开集的限制交换。因此可假设
$Y' \to Y$ 是仿射概形之间的态射；见注
\ref{duality-remark-check-over-affines}。此时结论由引理
\ref{duality-lemma-more-base-change} 得到。
\end{proof}



\section{有效 Cartier 除子的推前右伴随}
\label{duality-section-dualizing-Cartier}

\noindent
设 $X$ 为概形，$i : D \to X$ 为一个有效 Cartier 除子的嵌入。
记 $\mathcal{N} = i^*\mathcal{O}_X(D)$ 为 $i$ 的法层；见《态射》，节
\ref{morphisms-section-conormal-sheaf} 及《除子》，节
\ref{divisors-section-effective-Cartier-divisors}。回忆
$R\SheafHom(\mathcal{O}_D, -)$ 表示
$i_* : D(\mathcal{O}_D) \to D(\mathcal{O}_X)$ 的右伴随，并满足
$i_*R\SheafHom(\mathcal{O}_D, -) =
R\SheafHom_{\mathcal{O}_X}(i_*\mathcal{O}_D, -)$；见节
\ref{duality-section-sections-with-exact-support}。

\begin{lemma}
\label{duality-lemma-compute-for-effective-Cartier}
如上，设 $X$ 为概形，$D \subset X$ 为有效 Cartier 除子。
则在 $D(\mathcal{O}_D)$ 中存在典范同构
$R\SheafHom(\mathcal{O}_D, \mathcal{O}_X) = \mathcal{N}[-1]$。
\end{lemma}

\begin{proof}
等价地，这是说 $R\SheafHom(\mathcal{O}_D, \mathcal{O}_X)$
唯一非零的上同调层位于次数 $1$，且该层同构于 $\mathcal{N}$。
由于 $i_*$ 正合且全忠实，只需证明
$i_*R\SheafHom(\mathcal{O}_D, \mathcal{O}_X)$ 同构于
$i_*\mathcal{N}[-1]$。由引理
\ref{duality-lemma-sheaf-with-exact-support-ext}，有
$i_*R\SheafHom(\mathcal{O}_D, \mathcal{O}_X) =
R\SheafHom_{\mathcal{O}_X}(i_*\mathcal{O}_D, \mathcal{O}_X)$。
可使用分解
$$
0 \to \mathcal{I} \to \mathcal{O}_X \to i_*\mathcal{O}_D \to 0
$$
作计算，其中 $\mathcal{I}$ 是 $D$ 的理想层。由于局部计算给出
$R\SheafHom_{\mathcal{O}_X}(\mathcal{O}_X, \mathcal{O}_X) = \mathcal{O}_X$ 和
$R\SheafHom_{\mathcal{O}_X}(\mathcal{I}, \mathcal{O}_X) = \mathcal{O}_X(D)$，
可知
$$
R\SheafHom_{\mathcal{O}_X}(i_*\mathcal{O}_D, \mathcal{O}_X) =
(\mathcal{O}_X \to \mathcal{O}_X(D))
$$
其中右边的 $\mathcal{O}_X$ 位于次数 $0$，$\mathcal{O}_X(D)$ 位于次数 $1$。
结论由短正合列
$$
0 \to \mathcal{O}_X \to \mathcal{O}_X(D) \to i_*\mathcal{N} \to 0
$$
得到；该列来自 $D$ 是 $\mathcal{O}_X(D)$ 的典范截面的零点概形，
以及 $\mathcal{N} = i^*\mathcal{O}_X(D)$ 这一事实。
\end{proof}

\noindent
对 $D(\mathcal{O}_X)$ 的每个对象 $K$，存在典范映射
\begin{equation}
\label{duality-equation-map-effective-Cartier}
Li^*K
\otimes_{\mathcal{O}_D}^\mathbf{L}
R\SheafHom(\mathcal{O}_D, \mathcal{O}_X)
\longrightarrow
R\SheafHom(\mathcal{O}_D, K)
\end{equation}
它在 $D(\mathcal{O}_D)$ 中关于 $K$ 函子性，并与可区别三角相容。
具体地，此映射伴随于映射
$$
i_*(Li^*K \otimes^\mathbf{L}_{\mathcal{O}_D}
R\SheafHom(\mathcal{O}_D, \mathcal{O}_X)) =
K \otimes^\mathbf{L}_{\mathcal{O}_X}
R\SheafHom_{\mathcal{O}_X}(i_*\mathcal{O}_D, \mathcal{O}_X)
\longrightarrow K
$$
其中等式由《上同调》，引理
\ref{cohomology-lemma-projection-formula-closed-immersion} 给出，
而箭头来自映射 $\mathcal{O}_X \to i_*\mathcal{O}_D$ 所诱导的典范映射
$R\SheafHom_{\mathcal{O}_X}(i_*\mathcal{O}_D, \mathcal{O}_X) \to \mathcal{O}_X$。

\medskip\noindent
若 $K \in D_\QCoh(\mathcal{O}_X)$，则映射
(\ref{duality-equation-map-effective-Cartier}) 等于
(\ref{duality-equation-compare-with-pullback})；这里使用引理
\ref{duality-lemma-twisted-inverse-image-closed} 中的识别
$a(K) = R\SheafHom(\mathcal{O}_D, K)$。若
$K \in D_\QCoh(\mathcal{O}_X)$ 且 $X$ Noetherian，则下述引理是引理
\ref{duality-lemma-compare-with-pullback-flat-proper-noetherian} 的特殊情形。

\begin{lemma}
\label{duality-lemma-sheaf-with-exact-support-effective-Cartier}
如上，设 $X$ 为概形，$D \subset X$ 为有效 Cartier 除子。则
(\ref{duality-equation-map-effective-Cartier}) 与引理
\ref{duality-lemma-compute-for-effective-Cartier} 合在一起，定义了一个关于
$D(\mathcal{O}_X)$ 中的 $K$ 函子性的同构
$$
Li^*K \otimes_{\mathcal{O}_D}^\mathbf{L} \mathcal{N}[-1]
\longrightarrow
R\SheafHom(\mathcal{O}_D, K)
$$
\end{lemma}

\begin{proof}
由于 $i_*$ 在模上正合且全忠实，为证明该映射是同构，只需证明应用
$i_*$ 后它是同构。以下将不再说明地使用引理
\ref{duality-lemma-compute-for-effective-Cartier} 的证明中所用的短正合列
$0 \to \mathcal{I} \to \mathcal{O}_X \to i_*\mathcal{O}_D \to 0$
和
$0 \to \mathcal{O}_X \to \mathcal{O}_X(D) \to i_*\mathcal{N} \to 0$。
由《上同调》，引理
\ref{cohomology-lemma-projection-formula-closed-immersion}
（它用于定义映射 (\ref{duality-equation-map-effective-Cartier})），左边变为
$$
K \otimes_{\mathcal{O}_X}^\mathbf{L} i_*\mathcal{N}[-1] =
K \otimes_{\mathcal{O}_X}^\mathbf{L} (\mathcal{O}_X \to \mathcal{O}_X(D))
$$
右边变为
\begin{align*}
R\SheafHom_{\mathcal{O}_X}(i_*\mathcal{O}_D, K) & =
R\SheafHom_{\mathcal{O}_X}((\mathcal{I} \to \mathcal{O}_X), K) \\
& =
R\SheafHom_{\mathcal{O}_X}((\mathcal{I} \to \mathcal{O}_X), \mathcal{O}_X)
\otimes_{\mathcal{O}_X}^\mathbf{L} K
\end{align*}
最后一个等式由《上同调》，引理
\ref{cohomology-lemma-dual-perfect-complex} 得到。由于该映射来自同构
$$
R\SheafHom_{\mathcal{O}_X}((\mathcal{I} \to \mathcal{O}_X), \mathcal{O}_X)
= (\mathcal{O}_X \to \mathcal{O}_X(D))
$$
结论显然。
\end{proof}





\section{例中的推前右伴随}
\label{duality-section-examples}

\noindent
本节在若干例子中计算推前的右伴随。这些同构是典范的，但仅在最弱的
可能意义下如此；也就是说，我们既不证明也不声称这些同构与基变换、
态射复合等各种运算相容。关于这类问题有大量文献。读者可从
\cite{RD}、\cite{Conrad-GD} 中的材料开始（这些文献采用不同的对偶理论
出发点，但讨论了为相对对偶化复形构造典范代表的问题），然后继续查阅
Joseph Lipman 及其合作者的工作。

\begin{lemma}
\label{duality-lemma-upper-shriek-P1}
设 $Y$ 为拟紧拟分离概形。设 $\mathcal{E}$ 为秩为 $n + 1$ 的有限局部自由
$\mathcal{O}_Y$-模，其行列式为
$\mathcal{L} = \wedge^{n + 1}(\mathcal{E})$。设投影为
$f : X = \mathbf{P}(\mathcal{E}) \to Y$。设 $a$ 为引理
\ref{duality-lemma-twisted-inverse-image} 中
$Rf_* : D_\QCoh(\mathcal{O}_X) \to D_\QCoh(\mathcal{O}_Y)$ 的右伴随。
则存在同构
$$
c : f^*\mathcal{L}(-n - 1)[n] \longrightarrow a(\mathcal{O}_Y)
$$
特别地，若 $\mathcal{E} = \mathcal{O}_Y^{\oplus n + 1}$，则
$X = \mathbf{P}^n_Y$，并得到
$a(\mathcal{O}_Y) = \mathcal{O}_X(-n - 1)[n]$。
\end{lemma}

\begin{proof}
在《概形的上同调》，引理
\ref{coherent-lemma-cohomology-projective-bundle} 的证明中，我们构造了典范同构
$$
R^nf_*(f^*\mathcal{L}(-n - 1)) \longrightarrow \mathcal{O}_Y
$$
此外，$Rf_*(f^*\mathcal{L}(-n - 1))[n] = R^nf_*(f^*\mathcal{L}(-n - 1))$，
即其他高阶正像均为零。因此得到同构
$$
Rf_*(f^*\mathcal{L}(-n - 1)[n]) \longrightarrow \mathcal{O}_Y
$$
由于 $a$ 是 $Rf_*$ 的右伴随，此同构确定了引理陈述中的 $c$。
由引理 \ref{duality-lemma-proper-noetherian}，$a$ 的构造关于基底是局部的。
特别地，为检验 $c$ 是同构，可以在 $Y$ 上局部工作。换言之，可假设
$Y$ 仿射且 $\mathcal{E} = \mathcal{O}_Y^{\oplus n + 1}$。
此时层 $\mathcal{O}_X, \mathcal{O}_X(-1), \ldots,
\mathcal{O}_X(-n)$ 生成 $D_\QCoh(X)$；见《概形的导出范畴》，引理
\ref{perfect-lemma-generator-P1}。因此，只需证明映射
$c : \mathcal{O}_X(-n - 1)[n] \to a(\mathcal{O}_Y)$ 在函子
$$
F_{i, p}(-) = \Hom_{D(\mathcal{O}_X)}(\mathcal{O}_X(i), (-)[p])
$$
下变为同构，其中 $i \in \{-n, \ldots, 0\}$ 且
$p \in \mathbf{Z}$。对 $F_{0, p}$，这由箭头 $c$ 的构造成立！
对 $i \in \{-n, \ldots, -1\}$，有
$$
\Hom_{D(\mathcal{O}_X)}(\mathcal{O}_X(i), \mathcal{O}_X(-n - 1)[n + p]) =
H^p(X, \mathcal{O}_X(-n - 1 - i)) = 0
$$
这是射影空间上同调计算的结果（《概形的上同调》，引理
\ref{coherent-lemma-cohomology-projective-space-over-ring}）；并且有
$$
\Hom_{D(\mathcal{O}_X)}(\mathcal{O}_X(i), a(\mathcal{O}_Y)[p]) =
\Hom_{D(\mathcal{O}_Y)}(Rf_*\mathcal{O}_X(i), \mathcal{O}_Y[p]) = 0
$$
因为同一引理给出 $Rf_*\mathcal{O}_X(i) = 0$。因此
$F_{i, p}(c)$ 的源和目标均为零，故 $F_{i, p}(c)$ 必为同构。
证明完成。
\end{proof}

\begin{example}
\label{duality-example-base-change-wrong}
若 $f$ 为完美真态射而 $g$ 完美，则基变换映射
(\ref{duality-equation-base-change-map}) 不一定是同构。设 $k$ 为域。
令 $Y = \mathbf{A}^2_k$，并令 $f : X \to Y$ 为 $Y$ 在原点处的吹起。
记 $E \subset X$ 为例外除子。则 $f$ 可分解为
$$
X \xrightarrow{i} \mathbf{P}^1_Y \xrightarrow{p} Y
$$
由此得到分解 $a = c \circ b$，其中 $a$、$b$、$c$ 分别是引理
\ref{duality-lemma-twisted-inverse-image} 中 $Rf_*$、$Rp_*$、$Ri_*$ 的右伴随。
记 $\mathcal{O}(n)$ 为 $\mathbf{P}^1_Y$ 上结构层的 Serre 扭曲，
记 $\mathcal{O}_X(n)$ 为其到 $X$ 的限制。注意
$X \subset \mathbf{P}^1_Y$ 由一次方程截出，故
$\mathcal{O}(X) = \mathcal{O}(1)$。由引理
\ref{duality-lemma-upper-shriek-P1}，有
$b(\mathcal{O}_Y) = \mathcal{O}(-2)[1]$。由引理
\ref{duality-lemma-twisted-inverse-image-closed}，有
$$
a(\mathcal{O}_Y) = c(b(\mathcal{O}_Y)) =
c(\mathcal{O}(-2)[1]) =
R\SheafHom(\mathcal{O}_X, \mathcal{O}(-2)[1]) =
\mathcal{O}_X(-1)
$$
最后一个等式由引理
\ref{duality-lemma-sheaf-with-exact-support-effective-Cartier}。令
$Y' = \Spec(k)$ 为 $Y$ 中的原点。$a(\mathcal{O}_Y)$ 到
$X' = E = \mathbf{P}^1_k$ 的限制是次数为 $-1$、置于上同调次数
$0$ 的可逆层。另一方面，
$a'(\mathcal{O}_{\Spec(k)}) = \mathcal{O}_E(-2)[1]$
是次数为 $-2$、置于上同调次数 $-1$ 的可逆层，因而二者不同。
在此例中，引理 \ref{duality-lemma-more-base-change} 的 Tor 独立性假设不成立。
\end{example}

\begin{lemma}
\label{duality-lemma-ext}
设 $Y$ 为赋环空间，$\mathcal{I} \subset \mathcal{O}_Y$ 为理想层。
令 $\mathcal{O}_X = \mathcal{O}_Y/\mathcal{I}$ 及
$\mathcal{N} =
\SheafHom_{\mathcal{O}_Y}(\mathcal{I}/\mathcal{I}^2, \mathcal{O}_X)$。
则存在典范同构
$c : \mathcal{N} \to
\SheafExt^1_{\mathcal{O}_Y}(\mathcal{O}_X, \mathcal{O}_X)
$。
\end{lemma}

\begin{proof}
考虑典范短正合列
\begin{equation}
\label{duality-equation-second-order-thickening}
0 \to \mathcal{I}/\mathcal{I}^2 \to \mathcal{O}_Y/\mathcal{I}^2 \to
\mathcal{O}_X \to 0
\end{equation}
设 $U \subset X$ 为开集，$s \in \mathcal{N}(U)$。沿 $s$ 对
(\ref{duality-equation-second-order-thickening}) 作推出，可得到
$\mathcal{O}_X|_U$ 被 $\mathcal{O}_X|_U$ 的扩张 $E_s$。
这又定义了 $\SheafExt^1_{\mathcal{O}_Y}(\mathcal{O}_X, \mathcal{O}_X)$
在 $U$ 上的截面 $c(s)$。见《上同调》，引理
\ref{cohomology-lemma-section-RHom-over-U} 和《导出范畴》，引理
\ref{derived-lemma-ext-1}。反过来，给定 $\mathcal{O}_U$-模扩张
$$
0 \to \mathcal{O}_X|_U \to \mathcal{E} \to \mathcal{O}_X|_U \to 0
$$
可找到开覆盖 $U = \bigcup U_i$ 及截面
$e_i \in \mathcal{E}(U_i)$，它们映到 $1 \in \mathcal{O}_X(U_i)$。
于是 $e_i$ 定义映射
$\mathcal{O}_Y|_{U_i} \to \mathcal{E}|_{U_i}$，其核包含
$\mathcal{I}^2$。由此可见，$\mathcal{E}|_{U_i}$ 来自上述推出。
这说明 $c$ 满射。单射性的证明从略。
\end{proof}

\begin{lemma}
\label{duality-lemma-regular-ideal-ext}
设 $Y$ 为赋环空间，$\mathcal{I} \subset \mathcal{O}_Y$ 为理想层，
令 $\mathcal{O}_X = \mathcal{O}_Y/\mathcal{I}$。若 $\mathcal{I}$
是 Koszul-正则的（《除子》，定义
\ref{divisors-definition-regular-ideal-sheaf}），则
$R\SheafHom_{\mathcal{O}_Y}(\mathcal{O}_X, \mathcal{O}_X)$ 上的复合定义同构
$$
\wedge^i(\SheafExt^1_{\mathcal{O}_Y}(\mathcal{O}_X, \mathcal{O}_X))
\longrightarrow
\SheafExt^i_{\mathcal{O}_Y}(\mathcal{O}_X, \mathcal{O}_X)
$$
对所有 $i$ 均成立。
\end{lemma}

\begin{proof}
这里的复合是指《上同调》，引理
\ref{cohomology-lemma-internal-hom-composition} 中的映射
$$
R\SheafHom_{\mathcal{O}_Y}(\mathcal{O}_X, \mathcal{O}_X)
\otimes_{\mathcal{O}_Y}^\mathbf{L}
R\SheafHom_{\mathcal{O}_Y}(\mathcal{O}_X, \mathcal{O}_X)
\longrightarrow
R\SheafHom_{\mathcal{O}_Y}(\mathcal{O}_X, \mathcal{O}_X)
$$
它诱导乘法映射
$$
\SheafExt^a_{\mathcal{O}_Y}(\mathcal{O}_X, \mathcal{O}_X)
\otimes_{\mathcal{O}_Y}
\SheafExt^b_{\mathcal{O}_Y}(\mathcal{O}_X, \mathcal{O}_X)
\longrightarrow
\SheafExt^{a + b}_{\mathcal{O}_Y}(\mathcal{O}_X, \mathcal{O}_X)
$$
可与《代数进阶》，方程
(\ref{more-algebra-equation-simple-tor-product}) 比较。引理陈述的含义是诱导映射
$$
\SheafExt^1_{\mathcal{O}_Y}(\mathcal{O}_X, \mathcal{O}_X)
\otimes \ldots \otimes
\SheafExt^1_{\mathcal{O}_Y}(\mathcal{O}_X, \mathcal{O}_X)
\longrightarrow
\SheafExt^i_{\mathcal{O}_Y}(\mathcal{O}_X, \mathcal{O}_X)
$$
经由外积因子分解，并由此诱导同构。为证明这一点，可以在 $Y$ 上局部工作。
因此可假设存在 $\mathcal{O}_Y$ 的整体截面
$f_1, \ldots, f_r$，它们生成 $\mathcal{I}$ 并构成 Koszul 正则序列。记
$$
\mathcal{A} = \mathcal{O}_Y\langle \xi_1, \ldots, \xi_r\rangle
$$
为严格交换微分分次 $\mathcal{O}_Y$-代数层；它是在
$\mathcal{O}_Y$ 上由次数 $-1$ 的 $\xi_1, \ldots, \xi_r$
生成的（除幂）多项式代数，微分 $\text{d}$ 由规则
$\text{d}\xi_i = f_i$ 给出。记 $\mathcal{A}^\bullet$ 为其底层
$\mathcal{O}_Y$-模复形，即上述 Koszul 复形。于是典范映射
$\mathcal{A}^\bullet \to \mathcal{O}_X$ 是拟同构。由《上同调》，引理
\ref{cohomology-lemma-Rhom-strictly-perfect}，得到拟同构
$$
R\SheafHom_{\mathcal{O}_Y}(\mathcal{O}_X, \mathcal{O}_X) \to
\SheafHom^\bullet(\mathcal{A}^\bullet, \mathcal{A}^\bullet) \to
\SheafHom^\bullet(\mathcal{A}^\bullet, \mathcal{O}_X)
$$
后一个复形的微分为零，因而
$$
\SheafExt^i_{\mathcal{O}_Y}(\mathcal{O}_X, \mathcal{O}_X)
\cong \SheafHom_{\mathcal{O}_Y}(\mathcal{A}^{-i}, \mathcal{O}_X)
$$
对 $j \in \{1, \ldots, r\}$，令
$\delta_j : \mathcal{A} \to \mathcal{A}$ 为次数 $1$ 的导子，满足
$\delta_j(\xi_i) = \delta_{ij}$（Kronecker 符号）。计算表明
$\delta_j \circ \text{d} = - \text{d} \circ \delta_j$，故得到复形态射
$$
\delta_j : \mathcal{A}^\bullet \to \mathcal{A}^\bullet[1].
$$
因此 $\delta_j$ 定义了相应 $\SheafExt$-层的一个截面。另一次计算表明，
$\delta_1, \ldots, \delta_r$ 映到
$\SheafHom_{\mathcal{O}_Y}(\mathcal{A}^{-1}, \mathcal{O}_X)$
在 $\mathcal{O}_X$ 上的一组基。显然，作为 $\mathcal{A}$ 的自同态，
$\delta_j \circ \delta_j = 0$，且
$\delta_j \circ \delta_{j'} = - \delta_{j'} \circ \delta_j$；
故在 $\SheafExt$-层中亦如此。这说明上述映射经由外幂因子分解。
为看出所得映射是所需同构，读者只需验证元素
$$
\delta_{j_1} \circ \ldots \circ \delta_{j_i}
$$
对 $j_1 < \ldots < j_i$ 映到层
$\SheafHom_{\mathcal{O}_Y}(\mathcal{A}^{-i}, \mathcal{O}_X)$
在 $\mathcal{O}_X$ 上的一组基。
\end{proof}

\begin{lemma}
\label{duality-lemma-regular-immersion-ext}
设 $Y$ 为赋环空间，$\mathcal{I} \subset \mathcal{O}_Y$ 为理想层。
令 $\mathcal{O}_X = \mathcal{O}_Y/\mathcal{I}$ 及
$\mathcal{N} =
\SheafHom_{\mathcal{O}_Y}(\mathcal{I}/\mathcal{I}^2, \mathcal{O}_X)$。
若 $\mathcal{I}$ 是 Koszul-正则的（《除子》，定义
\ref{divisors-definition-regular-ideal-sheaf}），则
$$
R\SheafHom_{\mathcal{O}_Y}(\mathcal{O}_X, \mathcal{O}_Y) =
\wedge^r \mathcal{N}[-r]
$$
其中 $r : Y \to \{1, 2, 3, \ldots \}$ 把 $y$ 映到 $y$ 的某个邻域中
生成 $\mathcal{I}$ 所需的最少生成元数。
\end{lemma}

\begin{proof}
由引理 \ref{duality-lemma-ext} 和 \ref{duality-lemma-regular-ideal-ext}，可知对
$i \geq 0$ 有同构
$\wedge^i\mathcal{N} \to
\SheafExt^i_{\mathcal{O}_Y}(\mathcal{O}_X, \mathcal{O}_X)$。
因此，只需证明映射 $\mathcal{O}_Y \to \mathcal{O}_X$ 诱导同构
$$
\SheafExt^r_{\mathcal{O}_Y}(\mathcal{O}_X, \mathcal{O}_Y)
\longrightarrow
\SheafExt^r_{\mathcal{O}_Y}(\mathcal{O}_X, \mathcal{O}_X)
$$
并证明当 $i \not = r$ 时，
$\SheafExt^i_{\mathcal{O}_Y}(\mathcal{O}_X, \mathcal{O}_Y)$ 为零。
这些陈述关于 $Y$ 是局部的。因此可假设有 $\mathcal{O}_Y$ 的整体截面
$f_1, \ldots, f_r$，它们生成 $\mathcal{I}$ 并构成 Koszul 正则序列。
令 $\mathcal{A}^\bullet$ 为引理
\ref{duality-lemma-regular-ideal-ext} 的证明中引入的关于
$f_1, \ldots, f_r$ 的 Koszul 复形。则由《上同调》，引理
\ref{cohomology-lemma-Rhom-strictly-perfect}，有
$$
R\SheafHom_{\mathcal{O}_Y}(\mathcal{O}_X, \mathcal{O}_Y) =
\SheafHom^\bullet(\mathcal{A}^\bullet, \mathcal{O}_Y)
$$
记 $1 : \mathcal{A}^\bullet \to \mathcal{O}_Y$ 为微分分次
$\mathcal{O}_Y$-代数的映射，它在次数 $0$ 上由恒等映射
$\mathcal{A}^0 = \mathcal{O}_Y \to \mathcal{O}_Y$ 给出。
取引理 \ref{duality-lemma-regular-ideal-ext} 证明中的 $\delta_j$，通过把
$\delta_{j_1} \ldots \delta_{j_i}$ 映到次数 $i$ 中的
$1 \circ \delta_{j_1} \circ \ldots \circ \delta_{j_i}$，得到分次
$\mathcal{O}_Y$-模同构
$$
\mathcal{O}_Y\langle \delta_1, \ldots, \delta_r\rangle
\longrightarrow
\SheafHom^\bullet(\mathcal{A}^\bullet, \mathcal{O}_Y)
$$
在此同构下，右边的微分在左边诱导微分 $\text{d}$。由符号规则，
$\text{d}(1) = - \sum f_j \delta_j$。由于
$\delta_j : \mathcal{A}^\bullet \to \mathcal{A}^\bullet[1]$
是复形态射，故
$$
\text{d}(\delta_{j_1} \ldots \delta_{j_i}) =
(- \sum f_j \delta_j )\delta_{j_1} \ldots \delta_{j_i}
$$
注意在微分分次代数 $\mathcal{A}$ 上有
$\text{d} = \sum f_j \delta_j$。因此，由规则
$$
1 \circ \delta_{j_1} \ldots \delta_{j_i} \longmapsto
(\delta_{j_1} \circ \ldots \circ \delta_{j_i})(\xi_1 \ldots \xi_r)
$$
定义的映射在 $r$ 为奇数时给出复形同构
$$
\SheafHom^\bullet(\mathcal{A}^\bullet, \mathcal{O}_Y)
\longrightarrow \mathcal{A}^\bullet[-r]
$$
而在 $r$ 为偶数时与微分相差一个符号而交换。无论如何，这些复形有同构的
上同调，从而得到所需消没。在映射
$$
\SheafHom^\bullet(\mathcal{A}^\bullet, \mathcal{O}_Y)
\longrightarrow
\SheafHom^\bullet(\mathcal{A}^\bullet, \mathcal{O}_X)
$$
下次数 $r$ 的上同调同构也由此直接得到。
（注意，我们对 Koszul 复形的约定确实参与了同构
$R\SheafHom_{\mathcal{O}_X}(\mathcal{O}_X, \mathcal{O}_Y) =
\wedge^r \mathcal{N}[-r]$ 的定义。）
\end{proof}

\begin{lemma}
\label{duality-lemma-regular-immersion}
设 $Y$ 为拟紧拟分离概形，$i : X \to Y$ 为 Koszul-正则闭浸入。
设 $a$ 为引理 \ref{duality-lemma-twisted-inverse-image} 中
$Ri_* : D_\QCoh(\mathcal{O}_X) \to D_\QCoh(\mathcal{O}_Y)$ 的右伴随。
则存在同构
$$
\wedge^r\mathcal{N}[-r] \longrightarrow a(\mathcal{O}_Y)
$$
其中
$\mathcal{N} = \SheafHom_{\mathcal{O}_X}(\mathcal{C}_{X/Y}, \mathcal{O}_X)$
是 $i$ 的法层（《态射》，节
\ref{morphisms-section-conormal-sheaf}），而 $r$ 是将其秩视为
$X$ 上局部常值函数所得的函数。
\end{lemma}

\begin{proof}
由引理 \ref{duality-lemma-twisted-inverse-image-closed} 和
\ref{duality-lemma-sheaf-with-exact-support-ext}，回忆 $a(\mathcal{O}_Y)$ 是
$D_\QCoh(\mathcal{O}_X)$ 的对象，其到 $Y$ 的推前为
$R\SheafHom_{\mathcal{O}_Y}(i_*\mathcal{O}_X, \mathcal{O}_Y)$。
因此结论由引理 \ref{duality-lemma-regular-immersion-ext} 得到。
\end{proof}

\begin{lemma}
\label{duality-lemma-smooth-proper}
设 $S$ 为 Noetherian 概形，$f : X \to S$ 为相对维数 $d$ 的光滑真态射。
设 $a$ 为引理 \ref{duality-lemma-twisted-inverse-image} 中
$Rf_* : D_\QCoh(\mathcal{O}_X) \to D_\QCoh(\mathcal{O}_S)$ 的右伴随。
则在 $D(\mathcal{O}_X)$ 中存在同构
$$
\wedge^d \Omega_{X/S}[d] \longrightarrow a(\mathcal{O}_S)
$$
\end{lemma}

\begin{proof}
如注 \ref{duality-remark-relative-dualizing-complex}，令
$\omega_{X/S}^\bullet = a(\mathcal{O}_S)$。令 $c$ 为引理
\ref{duality-lemma-twisted-inverse-image} 中对应于
$\Delta : X \to X \times_S X$ 的右伴随。由于 $\Delta$ 是光滑态射的对角，
它是 Koszul-正则浸入；见《除子》，引理
\ref{divisors-lemma-immersion-smooth-into-smooth-regular-immersion}。
特别地，$\Delta$ 是完美真态射（《态射进阶》，引理
\ref{more-morphisms-lemma-regular-immersion-perfect}），从而得到
\begin{align*}
\mathcal{O}_X
& =
c(L\text{pr}_1^*\omega_{X/S}^\bullet) \\
& =
L\Delta^*(L\text{pr}_1^*\omega_{X/S}^\bullet)
\otimes_{\mathcal{O}_X}^\mathbf{L}
c(\mathcal{O}_{X \times_S X}) \\
& =
\omega_{X/S}^\bullet \otimes_{\mathcal{O}_X}^\mathbf{L}
c(\mathcal{O}_{X \times_S X}) \\
& =
\omega_{X/S}^\bullet
\otimes_{\mathcal{O}_X}^\mathbf{L}
\wedge^d(\mathcal{N}_\Delta)[-d]
\end{align*}
第一个等式是 (\ref{duality-equation-pre-rigid})，因为
$\omega_{X/S}^\bullet = a(\mathcal{O}_S)$。第二个等式由引理
\ref{duality-lemma-compare-with-pullback-flat-proper-noetherian}。第三个等式成立，因为
$\text{pr}_1 \circ \Delta = \text{id}_X$。第四个等式由引理
\ref{duality-lemma-regular-immersion}。注意 $\wedge^d(\mathcal{N}_\Delta)$
是可逆 $\mathcal{O}_X$-模。因此
$\wedge^d(\mathcal{N}_\Delta)[-d]$ 是 $D(\mathcal{O}_X)$ 中的可逆对象，
从而 $a(\mathcal{O}_S) = \omega_{X/S}^\bullet = \wedge^d(\mathcal{C}_\Delta)[d]$。
由于 $\Delta$ 的余法层 $\mathcal{C}_\Delta$ 由《态射》，引理
\ref{morphisms-lemma-differentials-diagonal} 等于 $\Omega_{X/S}$，证明完成。
\end{proof}







\section{上叹号函子}
\label{duality-section-upper-shriek}

\noindent
本节利用紧化，为在一个固定 Noetherian 基底上有限型且分离的概形之间的
态射构造函子 $f^!$。如凝聚对偶理论中通常的情形，需要证明若干图交换。
建议读者在读完构造后跳过这些引理的验证，直接进入下一节对上叹号函子性质的讨论。

\begin{situation}
\label{duality-situation-shriek}
这里 $S$ 是 Noetherian 概形，而 $\textit{FTS}_S$ 是如下范畴：
\begin{enumerate}
\item 对象是 $S$ 上的概形 $X$，其结构态射 $X \to S$ 既分离又有限型；
\item 对象之间的态射 $f : X \to Y$ 是 $S$ 上的概形态射。
\end{enumerate}
\end{situation}

\noindent
在情形 \ref{duality-situation-shriek} 中，给定 $\textit{FTS}_S$ 中的态射
$f : X \to Y$，我们将定义三角范畴间的正合函子
$$
f^! : D_\QCoh^+(\mathcal{O}_Y) \to D_\QCoh^+(\mathcal{O}_X)
$$
具体地，选取在 $Y$ 上的紧化 $X \to \overline{X}$；其存在性由
《平坦性进阶》，定理 \ref{flat-theorem-nagata} 和引理
\ref{flat-lemma-compactifyable}。记结构态射为
$\overline{f} : \overline{X} \to Y$。令
$\overline{a} : D_\QCoh(\mathcal{O}_Y) \to D_\QCoh(\mathcal{O}_{\overline{X}})$
为引理 \ref{duality-lemma-twisted-inverse-image} 中所构造的
$R\overline{f}_*$ 的右伴随。对 $K \in D_\QCoh^+(\mathcal{O}_Y)$，令
$$
f^!K  = \overline{a}(K)|_X
$$
由引理 \ref{duality-lemma-twisted-inverse-image-bounded-below}，所得对象属于
$D_\QCoh^+(\mathcal{O}_X)$。

\begin{lemma}
\label{duality-lemma-shriek-well-defined}
在情形 \ref{duality-situation-shriek} 中，设 $f : X \to Y$ 为
$\textit{FTS}_S$ 的态射。函子 $f^!$ 在典范同构意义下与紧化的选择无关。
\end{lemma}

\begin{proof}
$X$ 在 $Y$ 上的紧化范畴定义于《平坦性进阶》，节
\ref{flat-section-compactify}。由《平坦性进阶》，定理
\ref{flat-theorem-nagata} 和引理 \ref{flat-lemma-compactifyable}，
该范畴非空。对紧化的每个选择
$$
j : X \to \overline{X},\quad \overline{f} : \overline{X} \to Y
$$
上述构造配以函子
$j^* \circ \overline{a} :
D_\QCoh^+(\mathcal{O}_Y) \to D_\QCoh^+(\mathcal{O}_X)$，
其中 $\overline{a}$ 是引理 \ref{duality-lemma-twisted-inverse-image} 中构造的
$R\overline{f}_*$ 的右伴随。

\medskip\noindent
假设给定 $Y$ 上紧化 $j_i : X \to \overline{X}_i$ 之间的态射
$g : \overline{X}_1 \to \overline{X}_2$，满足
$g^{-1}(j_2(X)) = j_1(X)$\footnote{按我们对紧化的定义，这一点可能不成立。
见《平坦性进阶》，节 \ref{flat-section-compactify}。}。
令 $\overline{c}$ 为引理 \ref{duality-lemma-twisted-inverse-image} 中对应于 $g$ 的右伴随。
则 $\overline{c} \circ \overline{a}_2 = \overline{a}_1$，因为这些函子伴随于
$R\overline{f}_{2, *} \circ Rg_* = R(\overline{f}_2 \circ g)_*$。
由 (\ref{duality-equation-sheafy})，存在函子的典范变换
$$
j_1^* \circ \overline{c} \longrightarrow j_2^*
$$
其定义域和值域为
$D^+_\QCoh(\mathcal{O}_{\overline{X}_2}) \to D^+_\QCoh(\mathcal{O}_X)$；
由引理 \ref{duality-lemma-proper-noetherian}，它是同构。复合
$$
j_1^* \circ \overline{a}_1 \longrightarrow
j_1^* \circ \overline{c} \circ \overline{a}_2 \longrightarrow
j_2^* \circ \overline{a}_2
$$
是函子同构，记作 $\alpha_g$。

\medskip\noindent
考虑 $X$ 在 $Y$ 上的两个紧化
$j_i : X \to \overline{X}_i$，$i = 1, 2$。由《平坦性进阶》，引理
\ref{flat-lemma-compactifications-cofiltered} 第 (b) 部分，可找到具有稠密像的紧化
$j : X \to \overline{X}$ 以及紧化态射
$g_i : \overline{X} \to \overline{X}_i$。由《平坦性进阶》，引理
\ref{flat-lemma-compactifications-cofiltered} 第 (c) 部分，有
$g_i^{-1}(j_i(X)) = j(X)$。因此由上一段得到同构
$$
\alpha_{g_i} :
j^* \circ \overline{a}
\longrightarrow
j_i^* \circ \overline{a}_i
$$
从而得到同构
$$
\alpha_{g_2} \circ \alpha_{g_1}^{-1} :
j_1^* \circ \overline{a}_1 \to j_2^* \circ \overline{a}_2
$$
为完成证明，必须说明这些同构定义良好。我们断言，只需证明上一段中所构造的
同构的复合仍是这种同构（精确陈述见下一段）。建议读者在纸片上检验这一点；
不过我们也在本段余下部分完整写出细节。具体地，再选取一个具有稠密像的紧化
$j' : X \to \overline{X}'$ 以及紧化态射
$g'_i : \overline{X}' \to \overline{X}_i$。由《平坦性进阶》，引理
\ref{flat-lemma-compactifications-cofiltered}，可找到具有稠密像的紧化
$j'' : X \to \overline{X}''$ 以及紧化态射
$h : \overline{X}'' \to \overline{X}$ 和
$h' : \overline{X}'' \to \overline{X}'$。甚至可以假设
$g_1 \circ h = g'_1 \circ h'$ 且 $g_2 \circ h = g'_2 \circ h'$。
下一段的结果给出
$$
\alpha_{g_i} \circ \alpha_h = \alpha_{g_i \circ h} =
\alpha_{g'_i \circ h'} = \alpha_{g'_i} \circ \alpha_{h'}
$$
其中 $i = 1, 2$。由于这些全都是函子同构，可知所需等式
$\alpha_{g_2} \circ \alpha_{g_1}^{-1} =
\alpha_{g'_2} \circ \alpha_{g'_1}^{-1}$ 成立。

\medskip\noindent
假设给定紧化 $j_i : X \to \overline{X}_i$，其中
$i = 1, 2, 3$，以及紧化态射
$g : \overline{X}_1 \to \overline{X}_2$ 和
$h : \overline{X}_2 \to \overline{X}_3$，满足
$g^{-1}(j_2(X)) = j_1(X)$ 且 $h^{-1}(j_2(X)) = j_3(X)$。
令 $\overline{a}_i$ 如上。上述断言是指
$$
\alpha_g \circ \alpha_h = \alpha_{g \circ h} :
j_1^* \circ \overline{a}_1 \to j_3^* \circ \overline{a}_3
$$
令 $\overline{c}$（相应地\ $\overline{d}$）为引理
\ref{duality-lemma-twisted-inverse-image} 中对应于 $g$（相应地\ $h$）的右伴随。
则 $\overline{c} \circ \overline{a}_2 = \overline{a}_1$ 且
$\overline{d} \circ \overline{a}_3 = \overline{a}_2$，并存在典范变换
$$
j_1^* \circ \overline{c} \longrightarrow j_2^*
\quad\text{且}\quad
j_2^* \circ \overline{d} \longrightarrow j_3^*
$$
其函子分别为
$D^+_\QCoh(\mathcal{O}_{\overline{X}_2}) \to D^+_\QCoh(\mathcal{O}_X)$
和
$D^+_\QCoh(\mathcal{O}_{\overline{X}_3}) \to D^+_\QCoh(\mathcal{O}_X)$；
理由同上。记 $\overline{e}$ 为引理
\ref{duality-lemma-twisted-inverse-image} 中对应于 $h \circ g$ 的右伴随。
由 (\ref{duality-equation-sheafy})，存在函子的典范变换
$$
j_1^* \circ \overline{e} \longrightarrow j_3^*
$$
其定义域和值域为
$D^+_\QCoh(\mathcal{O}_{\overline{X}_3}) \to D^+_\QCoh(\mathcal{O}_X)$。
展开来说，必须证明复合
$$
\alpha_h \circ \alpha_g :
j_1^* \circ \overline{a}_1 \to
j_1^* \circ \overline{c} \circ \overline{a}_2 \to
j_2^* \circ \overline{a}_2 \to
j_2^* \circ \overline{d} \circ \overline{a}_3 \to
j_3^* \circ \overline{a}_3
$$
与复合
$$
\alpha_{h \circ g} :
j_1^* \circ \overline{a}_1 \to
j_1^* \circ \overline{e} \circ \overline{a}_3 \to
j_3^* \circ \overline{a}_3
$$
相同。将其分成两部分。第一部分是证明图
$$
\xymatrix{
\overline{a}_1 \ar[r] \ar[d] & \overline{c} \circ \overline{a}_2 \ar[d] \\
\overline{e} \circ \overline{a}_3 \ar[r] &
\overline{c} \circ \overline{d} \circ \overline{a}_3
}
$$
交换，其中下方水平箭头来自识别
$\overline{e} = \overline{c} \circ \overline{d}$。这是因为相应的总导出正像函子图
$$
\xymatrix{
R\overline{f}_{1, *} \ar[r] \ar[d] & Rg_* \circ R\overline{f}_{2, *} \ar[d] \\
R(h \circ g)_* \circ R\overline{f}_{3, *} \ar[r] &
Rg_* \circ Rh_* \circ R\overline{f}_{3, *}
}
$$
交换（此处插入将来引用）。第二部分是证明复合
$$
j_1^* \circ \overline{c} \circ \overline{d} \to
j_2^* \circ \overline{d} \to j_3^*
$$
通过识别 $\overline{e} = \overline{c} \circ \overline{d}$ 等于映射
$$
j_1^* \circ \overline{e} \to j_3^*
$$
这已在引理 \ref{duality-lemma-compose-base-change-maps} 中证明
（注意，本情形中该引理的态射 $f', g'$ 等于 $\text{id}_X$）。
\end{proof}

\begin{lemma}
\label{duality-lemma-upper-shriek-composition}
在情形 \ref{duality-situation-shriek} 中，设 $f : X \to Y$ 和
$g : Y \to Z$ 为 $\textit{FTS}_S$ 中可复合的态射。则存在典范同构
$(g \circ f)^! \to f^! \circ g^!$。
\end{lemma}

\begin{proof}
选取 $Y$ 在 $Z$ 上的紧化 $i : Y \to \overline{Y}$，再选取
$X$ 在 $\overline{Y}$ 上的紧化 $X \to \overline{X}$。这里两次使用了
《平坦性进阶》，定理 \ref{flat-theorem-nagata} 和引理
\ref{flat-lemma-compactifyable}。令 $\overline{a}$ 为引理
\ref{duality-lemma-twisted-inverse-image} 中对应于
$\overline{X} \to \overline{Y}$ 的右伴随，令 $\overline{b}$ 为引理
\ref{duality-lemma-twisted-inverse-image} 中对应于 $\overline{Y} \to Z$ 的右伴随。
则 $\overline{a} \circ \overline{b}$ 是引理
\ref{duality-lemma-twisted-inverse-image} 中对应于复合
$\overline{X} \to Z$ 的右伴随。因此
$g^! = i^* \circ \overline{b}$，且
$(g \circ f)^! = (X \to \overline{X})^* \circ \overline{a} \circ \overline{b}$。
令 $U$ 为 $Y$ 在 $\overline{X}$ 中的逆像，从而得到交换图
$$
\xymatrix{
X \ar[r]_j \ar[d] & U \ar[dl] \ar[r]_{j'} & \overline{X} \ar[dl] \\
Y \ar[r]_i \ar[d] & \overline{Y} \ar[dl] \\
Z
}
$$
令 $\overline{a}'$ 为引理 \ref{duality-lemma-twisted-inverse-image} 中对应于
$U \to Y$ 的右伴随。则 $f^! = j^* \circ \overline{a}'$。
由 (\ref{duality-equation-sheafy}) 得到
$$
\gamma : (j')^* \circ \overline{a} \to \overline{a}' \circ i^*
$$
并用它定义
$$
(g \circ f)^! =
(j' \circ j)^* \circ \overline{a} \circ \overline{b} =
j^* \circ (j')^* \circ \overline{a} \circ \overline{b}
\to
j^* \circ \overline{a}' \circ i^* \circ \overline{b} =
f^! \circ g^!
$$
由引理 \ref{duality-lemma-proper-noetherian}，它在
$D_\QCoh^+(\mathcal{O}_Z)$ 的对象上是同构。为完成证明，需说明此同构与所作选择无关。

\medskip\noindent
假设有两个图
$$
\vcenter{
\xymatrix{
X \ar[r]_{j_1} \ar[d] & U_1 \ar[dl] \ar[r]_{j'_1} & \overline{X}_1 \ar[dl] \\
Y \ar[r]_{i_1} \ar[d] & \overline{Y}_1 \ar[dl] \\
Z
}
}
\quad\text{且}\quad
\vcenter{
\xymatrix{
X \ar[r]_{j_2} \ar[d] & U_2 \ar[dl] \ar[r]_{j'_2} & \overline{X}_2 \ar[dl] \\
Y \ar[r]_{i_2} \ar[d] & \overline{Y}_2 \ar[dl] \\
Z
}
}
$$
首先可选取 $Y$ 在 $Z$ 上、具有稠密像且同时支配
$\overline{Y}_1$ 和 $\overline{Y}_2$ 的紧化
$i : Y \to \overline{Y}$；见《平坦性进阶》，引理
\ref{flat-lemma-compactifications-cofiltered}。由《平坦性进阶》，引理
\ref{flat-lemma-right-multiplicative-system} 及《范畴》，引理
\ref{categories-lemma-morphisms-right-localization} 和
\ref{categories-lemma-equality-morphisms-right-localization}，可选取
$X$ 在 $\overline{Y}$ 上、具有稠密像的紧化
$X \to \overline{X}$，以及态射
$\overline{X} \to \overline{X}_1$ 和
$\overline{X} \to \overline{X}_2$，使得复合
$\overline{X} \to \overline{Y} \to \overline{Y}_1$ 等于复合
$\overline{X} \to \overline{X}_1 \to \overline{Y}_1$，并使复合
$\overline{X} \to \overline{Y} \to \overline{Y}_2$ 等于复合
$\overline{X} \to \overline{X}_2 \to \overline{Y}_2$。
因此，只需在有如下交换图时比较这些图所确定的映射：
$$
\xymatrix{
X \ar[rr]_{j_1} \ar@{=}[d] & &
U_1 \ar[d] \ar[ddll] \ar[rr]_{j'_1} & &
\overline{X}_1 \ar[d] \ar[ddll] \\
X \ar'[r][rr]^-{j_2} \ar[d] & &
U_2 \ar'[dl][ddll] \ar'[r][rr]^-{j'_2} & &
\overline{X}_2 \ar[ddll] \\
Y \ar[rr]^{i_1} \ar@{=}[d] & & \overline{Y}_1 \ar[d] \\
Y \ar[rr]^{i_2} \ar[d] & & \overline{Y}_2 \ar[dll] \\
Z
}
$$
并且紧化 $X \to \overline{X}_1$ 和
$Y \to \overline{Y}_2$ 具有稠密像。分别用
$\overline{a}_i$、$\overline{a}'_i$、$\overline{c}$ 和
$\overline{c}'$ 表示引理 \ref{duality-lemma-twisted-inverse-image} 中对应于
$\overline{X}_i \to \overline{Y}_i$、$U_i \to Y$、
$\overline{X}_1 \to \overline{X}_2$ 和 $U_1 \to U_2$ 的右伴随。
以下每个方块
$$
\xymatrix{
X \ar[r] \ar[d] \ar@{}[dr]|A & U_1 \ar[d] \\
X \ar[r] & U_2
}
\quad
\xymatrix{
U_2 \ar[r] \ar[d] \ar@{}[dr]|B & \overline{X}_2 \ar[d] \\
Y \ar[r] & \overline{Y}_2
}
\quad
\xymatrix{
U_1 \ar[r] \ar[d] \ar@{}[dr]|C & \overline{X}_1 \ar[d] \\
Y \ar[r] & \overline{Y}_1
}
\quad
\xymatrix{
Y \ar[r] \ar[d] \ar@{}[dr]|D & \overline{Y}_1 \ar[d] \\
Y \ar[r] & \overline{Y}_2
}
\quad
\xymatrix{
X \ar[r] \ar[d] \ar@{}[dr]|E & \overline{X}_1 \ar[d] \\
X \ar[r] & \overline{X}_2
}
$$
都是笛卡尔的（对 A、D、E，见《平坦性进阶》，引理
\ref{flat-lemma-compactifications-cofiltered} 第 (c) 部分；对 B、C，回忆
$U_i$ 是 $Y$ 沿 $\overline{X}_i \to \overline{Y}_i$ 的逆像），
因而分别产生如下基变换映射 (\ref{duality-equation-sheafy})：
$$
\begin{matrix}
\gamma_A : j_1^* \circ \overline{c}' \to j_2^* &
\gamma_B : (j_2')^* \circ \overline{a}_2 \to \overline{a}'_2 \circ i_2^* &
\gamma_C : (j_1')^* \circ \overline{a}_1 \to \overline{a}'_1 \circ i_1^* \\
\gamma_D : i_1^* \circ \overline{d} \to i_2^* &
\gamma_E : (j'_1 \circ j_1)^* \circ \overline{c} \to (j'_2 \circ j_2)^*
\end{matrix}
$$
记 $f_1^! = j_1^* \circ \overline{a}'_1$、
$f_2^! = j_2^* \circ \overline{a}'_2$、
$g_1^! = i_1^* \circ \overline{b}_1$、
$g_2^! = i_2^* \circ \overline{b}_2$、
$(g \circ f)_1^! =
(j_1' \circ j_1)^* \circ \overline{a}_1 \circ \overline{b}_1$，以及
$(g \circ f)^!_2 =
(j_2' \circ j_2)^* \circ \overline{a}_2 \circ \overline{b}_2$。
证明第一段及引理 \ref{duality-lemma-shriek-well-defined} 中给出的构造使用
\begin{enumerate}
\item $\gamma_C$ 作为映射 $(g \circ f)^!_1 \to f_1^! \circ g_1^!$；
\item $\gamma_B$ 作为映射 $(g \circ f)^!_2 \to f_2^! \circ g_2^!$；
\item $\gamma_A$ 作为映射 $f_1^! \to f_2^!$；
\item $\gamma_D$ 作为映射 $g_1^! \to g_2^!$；并且
\item $\gamma_E$ 作为映射 $(g \circ f)^!_1 \to (g \circ f)^!_2$。
\end{enumerate}
必须证明图
$$
\xymatrix{
(g \circ f)^!_1 \ar[r]_{\gamma_E} \ar[d]_{\gamma_C} &
(g \circ f)^!_2 \ar[d]_{\gamma_B} \\
f_1^! \circ g_1^! \ar[r]^{\gamma_A \circ \gamma_D} & f_2^! \circ g_2^!
}
$$
交换。 将使用引理 \ref{duality-lemma-compose-base-change-maps}、
\ref{duality-lemma-compose-base-change-maps-horizontal}，以及注
\ref{duality-remark-going-around} 中的记号滥用（特别地，从记号中省略与恒等变换的
$\star$ 积）。可写成 $\gamma_E = \gamma_A \circ \gamma_F$，其中
$$
\xymatrix{
U_1 \ar[r] \ar[d] \ar@{}[rd]|F & \overline{X}_1 \ar[d] \\
U_2 \ar[r] & \overline{X}_2
}
$$
于是
$$
\gamma_B \circ \gamma_E = \gamma_B \circ \gamma_A  \circ \gamma_F
= \gamma_A \circ \gamma_B \circ \gamma_F
$$
最后一个等式成立，是因为方块 $A$ 与 $B$ 只在一个点相交
（类似于注 \ref{duality-remark-going-around} 中最后一个论证）。因此只需证明
$\gamma_D \circ \gamma_C = \gamma_B \circ \gamma_F$。由于二者都等于方块
$$
\xymatrix{
U_1 \ar[r] \ar[d] & \overline{X}_1 \ar[d] \\
Y \ar[r] & \overline{Y}_2
}
$$
的映射 (\ref{duality-equation-sheafy})，结论成立。
\end{proof}

\begin{lemma}
\label{duality-lemma-pseudo-functor}
在情形 \ref{duality-situation-shriek} 中，引理
\ref{duality-lemma-shriek-well-defined} 和
\ref{duality-lemma-upper-shriek-composition} 的构造定义了从范畴
$\textit{FTS}_S$ 到范畴的 $2$-范畴的伪函子（见《范畴》，定义
\ref{categories-definition-functor-into-2-category}）。
\end{lemma}

\begin{proof}
为证明这一点，必须证明对态射
$f : X \to Y$、$g : Y \to Z$、$h : Z \to T$，图
$$
\xymatrix{
(h \circ g \circ f)^! \ar[r]_{\gamma_{A + B}} \ar[d]_{\gamma_{B + C}} &
f^! \circ (h \circ g)^! \ar[d]^{\gamma_C} \\
(g \circ f)^! \circ h^! \ar[r]^{\gamma_A} & f^! \circ g^! \circ h^!
}
$$
交换（$\gamma$ 的含义见下文）。为此，先选取 $Z$ 在 $T$ 上的紧化
$\overline{Z}$，再选取 $Y$ 在 $\overline{Z}$ 上的紧化
$\overline{Y}$，最后选取 $X$ 在 $\overline{Y}$ 上的紧化
$\overline{X}$。这里使用《平坦性进阶》，定理
\ref{flat-theorem-nagata} 和引理 \ref{flat-lemma-compactifyable}。
令 $W \subset \overline{Y}$ 为 $Z$ 沿
$\overline{Y} \to \overline{Z}$ 的逆像，并令
$U \subset V \subset \overline{X}$ 为 $Y \subset W$ 沿
$\overline{X} \to \overline{Y}$ 的逆像。得到图
$$
\xymatrix{
X \ar[d]_f \ar[r] & U \ar[r] \ar[d] \ar@{}[dr]|A &
V \ar[d] \ar[r] \ar@{}[rd]|B & \overline{X} \ar[d] \\
Y \ar[d]_g \ar[r] & Y \ar[r] \ar[d] & W \ar[r] \ar[d] \ar@{}[rd]|C &
\overline{Y} \ar[d] \\
Z \ar[d]_h \ar[r] & Z \ar[d] \ar[r] & Z \ar[d] \ar[r] & \overline{Z} \ar[d] \\
T \ar[r] & T \ar[r] & T \ar[r] & T
}
$$
不引入大量记号，而完全像引理
\ref{duality-lemma-upper-shriek-composition} 的证明那样论证，可知第一个图中的映射
按所标示方式使用矩形 $A + B$、$B + C$、$A$、$C$ 的映射
(\ref{duality-equation-sheafy})。由引理
\ref{duality-lemma-compose-base-change-maps} 和
\ref{duality-lemma-compose-base-change-maps-horizontal}，有
$\gamma_{A + B} = \gamma_A \circ \gamma_B$ 且
$\gamma_{B + C} = \gamma_C \circ \gamma_B$，故只要
$\gamma_A \circ \gamma_C = \gamma_C \circ \gamma_A$，所需等式即成立。
这确实成立，因为方块 $A$ 与 $C$ 只在一个点相交
（类似于注 \ref{duality-remark-going-around} 中最后一个论证）。
\end{proof}

\begin{lemma}
\label{duality-lemma-map-pullback-to-shriek-well-defined}
在情形 \ref{duality-situation-shriek} 中，设
$f : X \to Y$ 是 $\textit{FTS}_S$ 的态射。存在关于
$D^+_\QCoh(\mathcal{O}_Y)$ 中的 $K$ 具有函子性的典范映射
$$
\mu_{f, K} :
Lf^*K \otimes_{\mathcal{O}_X}^\mathbf{L} f^!\mathcal{O}_Y
\longrightarrow
f^!K
$$
。若 $g : Y \to Z$ 是 $\textit{FTS}_S$ 的另一个态射，则对所有
$K \in D^+_\QCoh(\mathcal{O}_Z)$，图
$$
\xymatrix{
Lf^*(Lg^*K \otimes_{\mathcal{O}_Y}^\mathbf{L} g^!\mathcal{O}_Z)
\otimes_{\mathcal{O}_X}^\mathbf{L} f^!\mathcal{O}_Y
\ar@{=}[d] \ar[r]_-{\mu_f} &
f^!(Lg^*K \otimes_{\mathcal{O}_Y}^\mathbf{L} g^!\mathcal{O}_Z)
\ar[r]_-{f^!\mu_g} &
f^!g^!K \ar@{=}[d] \\
Lf^*Lg^*K \otimes_{\mathcal{O}_X}^\mathbf{L} Lf^* g^!\mathcal{O}_Z
\otimes_{\mathcal{O}_X}^\mathbf{L} f^!\mathcal{O}_Y \ar[r]^-{\mu_f} &
Lf^*Lg^*K \otimes_{\mathcal{O}_X}^\mathbf{L} f^!g^!\mathcal{O}_Z
\ar[r]^-{\mu_{g \circ f}} & f^!g^!K
}
$$
交换。
\end{lemma}

\begin{proof}
若 $f$ 是真态射，则 $f^! = a$，可使用
(\ref{duality-equation-compare-with-pullback})；若 $g$ 也是真态射，则引理
\ref{duality-lemma-transitivity-compare-with-pullback}（在更一般的情形下）证明该图交换。

\medskip\noindent
先定义映射 $\mu_{f, K}$。选取 $X$ 在 $Y$ 上的一个紧化
$j : X \to \overline{X}$。由于 $f^!$ 定义为
$j^* \circ \overline{a}$，将映射 (\ref{duality-equation-compare-with-pullback})
$$
L\overline{f}^*K \otimes_{\mathcal{O}_{\overline{X}}}^\mathbf{L}
\overline{a}(\mathcal{O}_Y)
\longrightarrow
\overline{a}(K)
$$
限制到 $X$，便得到 $\mu_{f, K}$。为说明它与紧化的选择无关，像引理
\ref{duality-lemma-shriek-well-defined} 的证明那样论证。建议读者先阅读该引理的证明。

\medskip\noindent
设给定 $Y$ 上两个紧化 $j_i : X \to \overline{X}_i$ 之间的态射
$g : \overline{X}_1 \to \overline{X}_2$，且
$g^{-1}(j_2(X)) = j_1(X)$。记 $\overline{c}$ 为引理
\ref{duality-lemma-twisted-inverse-image} 对态射 $g$ 给出的推前的右伴随。映射
$$
L\overline{f}_1^*K \otimes_{\mathcal{O}_{\overline{X}}}^\mathbf{L}
\overline{a}_1(\mathcal{O}_Y)
\longrightarrow
\overline{a}_1(K)
\quad\text{且}\quad
L\overline{f}_2^*K \otimes_{\mathcal{O}_{\overline{X}}}^\mathbf{L}
\overline{a}_2(\mathcal{O}_Y)
\longrightarrow
\overline{a}_2(K)
$$
嵌入交换图
$$
\xymatrix{
Lg^*(L\overline{f}_2^*K \otimes^\mathbf{L}
\overline{a}_2(\mathcal{O}_Y))
\otimes^\mathbf{L} \overline{c}(\mathcal{O}_{\overline{X}_2})
\ar@{=}[d] \ar[r]_-\sigma &
\overline{c}(L\overline{f}_2^*K \otimes^\mathbf{L}
\overline{a}_2(\mathcal{O}_Y)) \ar[r] &
\overline{c}(\overline{a}_2(K)) \ar@{=}[d] \\
L\overline{f}_1^*K \otimes^\mathbf{L} Lg^*\overline{a}_2(\mathcal{O}_Y)
\otimes^\mathbf{L} \overline{c}(\mathcal{O}_{\overline{X}_2})
\ar[r]^-{1 \otimes \tau} &
L\overline{f}_1^*K \otimes^\mathbf{L} \overline{a}_1(\mathcal{O}_Y) \ar[r] &
\overline{a}_1(K)
}
$$
，这由引理 \ref{duality-lemma-transitivity-compare-with-pullback} 得到。由引理
\ref{duality-lemma-compare-on-open}，映射 $\sigma$ 和 $\tau$ 限制到 $X$ 后是同构。
事实上还可说得更多。回忆在引理 \ref{duality-lemma-shriek-well-defined} 的证明中，
用映射 (\ref{duality-equation-sheafy})
$\gamma : j_1^* \circ \overline{c} \to j_2^*$ 构造了同构
$\alpha_g : j_1^* \circ \overline{a}_1 \to j_2^* \circ \overline{a}_2$。
用 $j_1$ 拉回映射 $\sigma$；若通过
$j_1^* \circ \overline{c} \to j_2^*$ 将
$j_1^*\overline{c}(\mathcal{O}_{\overline{X}_2})$ 与 $\mathcal{O}_X$ 等同，
则所得映射是
$j_2^*\left(L\overline{f}_2^*K \otimes^\mathbf{L}
\overline{a}_2(\mathcal{O}_Y)\right)$ 上的恒等映射；见引理
\ref{duality-lemma-restriction-compare-with-pullback}。类似地，映射
$\tau : Lg^*\overline{a}_2(\mathcal{O}_Y)
\otimes^\mathbf{L} \overline{c}(\mathcal{O}_{\overline{X}_2}) \to
\overline{a}_1(\mathcal{O}_Y) = \overline{c}(\overline{a}_2(\mathcal{O}_Y))$
拉回后，是 $j_2^*\overline{a}_2(\mathcal{O}_Y)$ 上的恒等映射。
所以用 $j_1$ 拉回，并在可以之处应用 $\gamma$，得到交换图
$$
\xymatrix{
j_2^*\left(L\overline{f}_2^*K \otimes^\mathbf{L}
\overline{a}_2(\mathcal{O}_Y)\right) \ar[r] \ar[d] &
j_2^*\overline{a}_2(K) \\
j_1^*L\overline{f}_1^*K \otimes^\mathbf{L} j_2^*\overline{a}_2(\mathcal{O}_Y) &
j_1^*(L\overline{f}_1^*K \otimes^\mathbf{L} \overline{a}_1(\mathcal{O}_Y))
\ar[r] \ar[l]_{1 \otimes \alpha_g} &
j_1^* \overline{a}_1(K) \ar[lu]_{\alpha_g}
}
$$
该图的交换性恰好说明：用紧化 $\overline{X}_1$ 构造的
$\mu_{f, K}$，经由引理 \ref{duality-lemma-shriek-well-defined} 的证明所用的等同
$\alpha_g$，与用紧化 $\overline{X}_2$ 构造的 $\mu_{f, K}$ 相同。
再完全像引理 \ref{duality-lemma-shriek-well-defined} 的证明那样使用若干范畴论论证，
便知 $\mu_{f, K}$ 定义良好（略去一个小细节）。

\medskip\noindent
由此，引理陈述中图的交换性来自同构
$(g \circ f)^! \to f^! \circ g^!$ 的构造（即引理
\ref{duality-lemma-upper-shriek-composition} 证明的第一部分，其中使用
$\overline{X} \to \overline{Y} \to Z$），以及引理
\ref{duality-lemma-transitivity-compare-with-pullback} 对
$\overline{X} \to \overline{Y} \to Z$ 的结论。
\end{proof}



\section{上叹号函子的性质}
\label{duality-section-upper-shriek-properties}

\noindent
下面给出上叹号函子的一些性质。

\begin{lemma}
\label{duality-lemma-shriek-open-immersion}
在情形 \ref{duality-situation-shriek} 中，设 $Y$ 是
$\textit{FTS}_S$ 的对象，且 $j : X \to Y$ 是开浸入。
则存在函子的典范同构 $j^! = j^*$。
\end{lemma}

\noindent
对 $\textit{FTS}_S$ 的 \'etale 态射 $f : X \to Y$，
也有 $f^* \cong f^!$；见引理 \ref{duality-lemma-shriek-etale}。

\begin{proof}
在此情形下，可选取 $\overline{X} = Y$ 作为紧化。
引理 \ref{duality-lemma-twisted-inverse-image} 对
$\text{id} : Y \to Y$ 给出的右伴随是恒等函子，故由定义
$j^! = j^*$。
\end{proof}

\begin{lemma}
\label{duality-lemma-restrict-before-or-after}
在情形 \ref{duality-situation-shriek} 中，设
$$
\xymatrix{
U \ar[r]_j \ar[d]_g & X \ar[d]^f \\
V \ar[r]^{j'} & Y
}
$$
是 $\textit{FTS}_S$ 中的交换图，其中 $j$ 和 $j'$ 是开浸入。
则作为函子
$D^+_\QCoh(\mathcal{O}_Y) \to D^+(\mathcal{O}_U)$，有
$j^* \circ f^! = g^! \circ (j')^*$。
\end{lemma}

\begin{proof}
令 $h = f \circ j = j' \circ g$。由引理
\ref{duality-lemma-upper-shriek-composition}，有
$h^! = j^! \circ f^! = g^! \circ (j')^!$。由引理
\ref{duality-lemma-shriek-open-immersion}，有
$j^! = j^*$ 且 $(j')^! = (j')^*$。
\end{proof}

\begin{lemma}
\label{duality-lemma-shriek-affine-line}
在情形 \ref{duality-situation-shriek} 中，设 $Y$ 是 $\textit{FTS}_S$ 的对象，
且 $f : X = \mathbf{A}^1_Y \to Y$ 是投影。则存在函子的一个
（非典范）同构 $f^!(-) \cong Lf^*(-) [1]$。
\end{lemma}

\begin{proof}
由于 $X = \mathbf{A}^1_Y \subset \mathbf{P}^1_Y$，且
$\mathcal{O}_{\mathbf{P}^1_Y}(-2)|_X \cong \mathcal{O}_X$，
结论由引理 \ref{duality-lemma-upper-shriek-P1} 和
\ref{duality-lemma-compare-with-pullback-flat-proper-noetherian} 得到。
\end{proof}

\begin{lemma}
\label{duality-lemma-shriek-closed-immersion}
在情形 \ref{duality-situation-shriek} 中，设 $Y$ 是
$\textit{FTS}_S$ 的对象，且 $i : X \to Y$ 是闭浸入。
则存在函子的典范同构
$i^!(-) = R\SheafHom(\mathcal{O}_X, -)$。
\end{lemma}

\begin{proof}
这只是引理 \ref{duality-lemma-twisted-inverse-image-closed} 的改述。
\end{proof}

\begin{remark}[上叹号的局部描述]
\label{duality-remark-local-calculation-shriek}
在情形 \ref{duality-situation-shriek} 中，设 $f : X \to Y$ 是
$\textit{FTS}_S$ 的态射。利用上述引理，可以如下局部计算 $f^!$。
设给定仿射开子集
$$
\xymatrix{
U \ar[r]_j \ar[d]_g & X \ar[d]^f \\
V \ar[r]^i & Y
}
$$
由于 $j^! \circ f^! = g^! \circ i^!$
（引理 \ref{duality-lemma-upper-shriek-composition}），且 $j^!$ 和 $i^!$
由限制给出（引理 \ref{duality-lemma-shriek-open-immersion}），可知对任意
$E \in D^+_\QCoh(\mathcal{O}_X)$，有
$$
(f^!E)|_U = g^!(E|_V)
$$
。写成 $U = \Spec(A)$、$V = \Spec(R)$，并令
$\varphi : R \to A$ 为对应于 $g$ 的有限型环同态。选取一个表示
$A = P/I$，其中 $P = R[x_1, \ldots, x_n]$ 是 $R$ 上 $n$ 个变量的
多项式代数。选取与 $E|_V$ 对应的对象 $K \in D^+(R)$
（《概形的导出范畴》，引理
\ref{perfect-lemma-affine-compare-bounded}）。那么断言 $f^!E|_U$ 对应于
$$
\varphi^!(K) = R\Hom(A, K \otimes_R^\mathbf{L} P)[n]
$$
，其中 $R\Hom(A, -) : D(P) \to D(A)$ 是《对偶复形》，第
\ref{dualizing-section-trivial} 节的函子，而
$\varphi^! : D(R) \to D(A)$ 是《对偶复形》，第
\ref{dualizing-section-relative-dualizing-complex-algebraic} 节的函子。
事实上，所选表示给出分解
$$
U \rightarrow \mathbf{A}^n_V \to \mathbf{A}^{n - 1}_V \to \ldots \to
\mathbf{A}^1_V \to V
$$
。恰好应用引理 \ref{duality-lemma-shriek-affine-line} $n$ 次，可知
$(\mathbf{A}^n_V \to V)^!(E|_V)$ 对应于
$K \otimes_R^\mathbf{L} P[n]$。由引理
\ref{duality-lemma-sheaf-with-exact-support-quasi-coherent} 和
\ref{duality-lemma-shriek-closed-immersion}，最后一步对应于应用
$R\Hom(A, -)$。
\end{remark}

\begin{lemma}
\label{duality-lemma-shriek-coherent}
在情形 \ref{duality-situation-shriek} 中，设 $f : X \to Y$ 是
$\textit{FTS}_S$ 的态射。则 $f^!$ 将
$D_{\textit{Coh}}^+(\mathcal{O}_Y)$ 映入
$D_{\textit{Coh}}^+(\mathcal{O}_X)$。
\end{lemma}

\begin{proof}
这个问题在 $X$ 上是局部的，故可设 $X$ 和 $Y$ 是仿射概形。
此时可将 $f : X \to Y$ 分解为
$$
X \xrightarrow{i} \mathbf{A}^n_Y \to \mathbf{A}^{n - 1}_Y \to \ldots \to
\mathbf{A}^1_Y \to Y
$$
，其中 $i$ 是闭浸入。该引理由以下各项推出：由引理
\ref{duality-lemma-shriek-affine-line}、
\ref{duality-lemma-sheaf-with-exact-support-coherent}、《对偶复形》，引理
\ref{dualizing-lemma-dualizing-polynomial-ring}，以及归纳法。
\end{proof}

\begin{lemma}
\label{duality-lemma-shriek-dualizing}
在情形 \ref{duality-situation-shriek} 中，设 $f : X \to Y$ 是
$\textit{FTS}_S$ 的态射。若 $K$ 是 $Y$ 的对偶复形，
则 $f^!K$ 是 $X$ 的对偶复形。
\end{lemma}

\begin{proof}
这个问题在 $X$ 上是局部的，故可设 $X$ 和 $Y$ 是仿射概形。
此时可将 $f : X \to Y$ 分解为
$$
X \xrightarrow{i} \mathbf{A}^n_Y \to \mathbf{A}^{n - 1}_Y \to \ldots \to
\mathbf{A}^1_Y \to Y
$$
，其中 $i$ 是闭浸入。由引理 \ref{duality-lemma-shriek-affine-line}、
《对偶复形》，引理 \ref{dualizing-lemma-dualizing-polynomial-ring}
以及归纳法，可知 $p^!K$ 是 $\mathbf{A}^n_Y$ 上的对偶复形，
其中 $p : \mathbf{A}^n_Y \to Y$ 是投影。类似地，由《对偶复形》，引理
\ref{dualizing-lemma-dualizing-quotient} 以及引理
\ref{duality-lemma-sheaf-with-exact-support-quasi-coherent} 和
\ref{duality-lemma-shriek-closed-immersion}，可知 $i^!$ 将对偶复形变为对偶复形。
\end{proof}

\begin{lemma}
\label{duality-lemma-shriek-via-duality}
在情形 \ref{duality-situation-shriek} 中，设 $f : X \to Y$ 是
$\textit{FTS}_S$ 的态射。设 $K$ 是 $Y$ 上的对偶复形。
对 $M \in D_{\textit{Coh}}(\mathcal{O}_Y)$，令
$D_Y(M) = R\SheafHom_{\mathcal{O}_Y}(M, K)$；对
$E \in D_{\textit{Coh}}(\mathcal{O}_X)$，令
$D_X(E) = R\SheafHom_{\mathcal{O}_X}(E, f^!K)$。
则对 $M \in D_{\textit{Coh}}^+(\mathcal{O}_Y)$，存在典范同构
$$
f^!M \longrightarrow D_X(Lf^*D_Y(M))
$$
。
\end{lemma}

\begin{proof}
选取 $X$ 在 $Y$ 上的紧化 $j : X \subset \overline{X}$
（《平坦性进阶》，定理 \ref{flat-theorem-nagata} 和引理
\ref{flat-lemma-compactifyable}）。令 $a$ 为引理
\ref{duality-lemma-twisted-inverse-image} 对 $\overline{X} \to Y$ 给出的右伴随。
对 $E \in D_{\textit{Coh}}(\mathcal{O}_{\overline{X}})$，令
$D_{\overline{X}}(E) = R\SheafHom_{\mathcal{O}_{\overline{X}}}(E, a(K))$。
由于 $R\SheafHom$ 的形成与限制到开集可交换，且
$f^! = j^* \circ a$，故只需证明：对
$M \in D_{\textit{Coh}}(\mathcal{O}_Y)$，存在典范同构
$$
a(M) \longrightarrow D_{\overline{X}}(L\overline{f}^*D_Y(M))
$$
。对 $F \in D_\QCoh(\mathcal{O}_X)$，有
\begin{align*}
\Hom_{\overline{X}}(
F, D_{\overline{X}}(L\overline{f}^*D_Y(M)))
& =
\Hom_{\overline{X}}(
F \otimes_{\mathcal{O}_X}^\mathbf{L} L\overline{f}^*D_Y(M), a(K)) \\
& =
\Hom_Y(
R\overline{f}_*(F \otimes_{\mathcal{O}_X}^\mathbf{L} L\overline{f}^*D_Y(M)),
K) \\
& =
\Hom_Y(
R\overline{f}_*(F) \otimes_{\mathcal{O}_Y}^\mathbf{L} D_Y(M),
K) \\
& =
\Hom_Y(
R\overline{f}_*(F), D_Y(D_Y(M))) \\
& =
\Hom_Y(R\overline{f}_*(F), M) \\
& = \Hom_{\overline{X}}(F, a(M))
\end{align*}
第一个等式由《上同调》，引理 \ref{cohomology-lemma-internal-hom} 得到。
第二个等式来自 $a$ 的定义。第三个等式由《概形的导出范畴》，引理
\ref{perfect-lemma-cohomology-base-change} 得到。第四个等式由《上同调》，
引理 \ref{cohomology-lemma-internal-hom} 和 $D_Y$ 的定义得到。
第五个等式由引理 \ref{duality-lemma-dualizing-schemes} 得到。
最后一个等式来自 $a$ 的定义。因此由 Yoneda 引理可知
$a(M) = D_{\overline{X}}(L\overline{f}^*D_Y(M))$。
\end{proof}

\begin{lemma}
\label{duality-lemma-perfect-comparison-shriek}
在情形 \ref{duality-situation-shriek} 中，设 $f : X \to Y$ 是
$\textit{FTS}_S$ 的态射。假设 $f$ 是完美态射（例如平坦态射）。则
\begin{enumerate}
\item[(a)] $f^!\mathcal{O}_Y$ 属于 $D_{\textit{Coh}}^b(\mathcal{O}_X)$；
\item[(b)] $f^!\mathcal{O}_Y$ 在 $D(f^{-1}\mathcal{O}_Y)$ 中具有有限 Tor 维数；
\item[(c)] $\mathcal{O}_X \to
R\SheafHom_{\mathcal{O}_X}(f^!\mathcal{O}_Y, f^!\mathcal{O}_Y)$
是同构；
\item[(d)] $f^!$ 将 $D_{\textit{Coh}}^b(\mathcal{O}_Y)$ 映入
$D_{\textit{Coh}}^b(\mathcal{O}_X)$；
\item[(e)] 引理 \ref{duality-lemma-map-pullback-to-shriek-well-defined} 的映射
$\mu_{f,  K} :
Lf^*K \otimes_{\mathcal{O}_X}^\mathbf{L} f^!\mathcal{O}_Y
\to
f^!K$
对所有 $K \in D_\QCoh^+(\mathcal{O}_Y)$ 都是同构。
\end{enumerate}
\end{lemma}

\begin{proof}
（有限表示的平坦态射是完美态射；见《态射进阶》，引理
\ref{more-morphisms-lemma-flat-finite-presentation-perfect}。）
断言 (a)、(b)、(c) 在 $X$ 上是局部的，故可设 $X$ 和 $Y$ 是仿射的。
此时注 \ref{duality-remark-local-calculation-shriek} 将 (a)、(b)、(c) 分别化为
《对偶复形》，引理 \ref{dualizing-lemma-relative-dualizing-algebraic}
的 (1)(a)、(1)(b)、(1)(c)。（使用《概形的导出范畴》，引理
\ref{perfect-lemma-tor-dimension-rel-affine}、
\ref{perfect-lemma-pseudo-coherent}、
\ref{perfect-lemma-identify-pseudo-coherent-noetherian}、
\ref{perfect-lemma-quasi-coherence-internal-hom} 以使各断言的表述相符。）

\medskip\noindent
为证明 (d) 和 (e)，先作一系列预备说明。
\begin{enumerate}
\item 已知 $f^!$ 将 $D_{\textit{Coh}}^+(\mathcal{O}_Y)$ 映入
$D_{\textit{Coh}}^+(\mathcal{O}_X)$；见引理
\ref{duality-lemma-shriek-coherent}。
\item 若 $f$ 是开浸入，则 $f^! = f^*$，且 (d)、(e) 成立，
因为在 $f^!$ 和 $\mu_f$ 的构造中可取 $\overline{X} = Y$；见引理
\ref{duality-lemma-shriek-open-immersion}。
\item 若 $f$ 是完美真态射，则由引理
\ref{duality-lemma-compare-with-pullback-flat-proper-noetherian}，(e) 成立。
\item 若存在开覆盖 $X = \bigcup U_i$，且 (d) 对 $U_i \to Y$ 成立，
则 (d) 对 $X \to Y$ 成立；(e) 也一样。这是因为 $f^!$ 和 $\mu_f$
的构造与转到开子概形可交换。
\item 若 $g : Y \to Z$ 是 $\textit{FTS}_S$ 中的第二个完美态射，
且 (e) 对 $f$ 和 $g$ 成立，则
$f^!g^!\mathcal{O}_Z =
Lf^*g^!\mathcal{O}_Z \otimes_{\mathcal{O}_X}^\mathbf{L} f^!\mathcal{O}_Y$，
并且由引理 \ref{duality-lemma-map-pullback-to-shriek-well-defined} 的交换图，
(e) 对 $g \circ f$ 成立。
\item 若 (d)、(e) 对 $f$ 和 $g$ 都成立，则它们对 $g \circ f$ 也成立。
事实上，由上一点，$f^!g^!\mathcal{O}_Z$ 上有界，而
$L(g \circ f)^*$ 具有有限上同调维数；于是 (d) 由上面已经看到的 (e) 推出。
\end{enumerate}
由此可知，只需在 $X$ 仿射时证明 (d)、(e)。选取浸入
$X \to \mathbf{A}^n_Y$（《态射》，引理
\ref{morphisms-lemma-quasi-affine-finite-type-over-S}），并将其分解为
$X \to U \to \mathbf{A}^n_Y \to Y$，其中 $X \to U$ 是闭浸入，
$U \subset \mathbf{A}^n_Y$ 是开集。注意由《态射进阶》，引理
\ref{more-morphisms-lemma-perfect-permanence}，$X \to U$ 是完美闭浸入。
所以只需对完美闭浸入以及投影 $\mathbf{A}^n_Y \to Y$ 证明该引理。

\medskip\noindent
设 $f : X \to Y$ 是完美闭浸入。已知 (d) 成立。设
$K \in D^b_{\textit{Coh}}(\mathcal{O}_Y)$。则
$f^!K = R\SheafHom(\mathcal{O}_X, K)$
（引理 \ref{duality-lemma-shriek-closed-immersion}），且
$f_*f^!K = R\SheafHom_{\mathcal{O}_Y}(f_*\mathcal{O}_X, K)$。
因为 $f$ 是完美的，复形 $f_*\mathcal{O}_X$ 是完美的，故
$R\SheafHom_{\mathcal{O}_Y}(f_*\mathcal{O}_X, K)$ 上有界。
这证明 (d) 成立。略去若干细节。

\medskip\noindent
设 $f : \mathbf{A}^n_Y \to Y$ 是投影。反复应用引理
\ref{duality-lemma-shriek-affine-line}，可知 (d) 成立。最后，(e) 成立，
因为它对 $\mathbf{P}^n_Y \to Y$ 成立（该态射平坦且真），而
$\mathbf{A}^n_Y \subset \mathbf{P}^n_Y$ 是开集。
\end{proof}

\begin{lemma}
\label{duality-lemma-flat-shriek-relatively-perfect}
在情形 \ref{duality-situation-shriek} 中，设 $f : X \to Y$ 是
$\textit{FTS}_S$ 的态射。若 $f$ 平坦，则
$f^!\mathcal{O}_Y$ 是 $D(\mathcal{O}_X)$ 中的 $Y$-完美对象，且
$\mathcal{O}_X \to
R\SheafHom_{\mathcal{O}_X}(f^!\mathcal{O}_Y, f^!\mathcal{O}_Y)$
是同构。
\end{lemma}

\begin{proof}
两个断言在 $X$ 上都是局部的，故可设 $X$ 和 $Y$ 是仿射的。
此时注 \ref{duality-remark-local-calculation-shriek} 将该引理化为一个代数引理，
即《对偶复形》，引理 \ref{dualizing-lemma-relative-dualizing-algebraic}。
（使用《概形的导出范畴》，引理
\ref{perfect-lemma-affine-locally-rel-perfect} 以使表述相符。）
\end{proof}

\begin{lemma}
\label{duality-lemma-lci-shriek}
在情形 \ref{duality-situation-shriek} 中，设 $f : X \to Y$ 是
$\textit{FTS}_S$ 的态射。假设 $f : X \to Y$ 是局部完全交态射。则
\begin{enumerate}
\item $f^!\mathcal{O}_Y$ 是 $D(\mathcal{O}_X)$ 中的可逆对象；并且
\item $f^!$ 将完美复形映为完美复形。
\end{enumerate}
\end{lemma}

\begin{proof}
回忆局部完全交态射是完美态射；见《态射进阶》，引理
\ref{more-morphisms-lemma-lci-properties}。由引理
\ref{duality-lemma-perfect-comparison-shriek}，只需证明 $f^!\mathcal{O}_Y$
是 $D(\mathcal{O}_X)$ 中的可逆对象。这个问题在 $X$ 和 $Y$ 上是局部的，
故可设 $X \to Y$ 分解为 $X \to \mathbf{A}^n_Y \to Y$，
其中第一个箭头是 Koszul 正则浸入；见《态射进阶》，第
\ref{more-morphisms-section-lci} 节。由引理
\ref{duality-lemma-shriek-affine-line}，结论对 $\mathbf{A}^n_Y \to Y$ 成立。
因此只需在 $f$ 是 Koszul 正则浸入时证明该引理。
再次局部地工作，归约到 $X = \Spec(A)$、$Y = \Spec(B)$ 的情形，
其中 $A = B/(f_1, \ldots, f_r)$，且
$f_1, \ldots, f_r \in B$ 是正则序列（这里使用：对 Noetherian 局部环，
Koszul 正则与正则这两个概念相同；见《代数进阶》，引理
\ref{more-algebra-lemma-noetherian-finite-all-equivalent}）。
所以 $X \to Y$ 是复合
$$
X = X_r \to X_{r - 1} \to \ldots \to X_1 \to X_0 = Y
$$
，其中每个箭头都是一个有效 Cartier 除子的包含。这样便归约到有效
Cartier 除子 $i : D \to X$ 的包含情形。此时由引理
\ref{duality-lemma-compute-for-effective-Cartier}，有
$i^!\mathcal{O}_X = \mathcal{N}[1]$，证明完成。
\end{proof}







\section{上叹号的基变换}
\label{duality-section-base-change-shriek}

\noindent
在情形 \ref{duality-situation-shriek} 中，设
$$
\xymatrix{
X' \ar[r]_{g'} \ar[d]_{f'} & X \ar[d]^f \\
Y' \ar[r]^g & Y
}
$$
是 $\textit{FTS}_S$ 中的笛卡尔图，且 $X$ 和 $Y'$ 在 $Y$ 上
Tor 独立。目前的设置不足以在这种一般性下构造基变换映射
$L(g')^* \circ f^! \to (f')^! \circ Lg^*$。原因是一般无法选取
$Y$ 上的紧化 $j : X \to \overline{X}$，使得
$\overline{X}$ 和 $Y'$ 在 $Y$ 上 Tor 独立；因此第
\ref{duality-section-base-change-map} 节对基变换映射的构造不适用\footnote{
熟悉导出代数几何的读者会认识到，这并不是一个“真正的”问题。
事实上，取 $\overline{X}'$ 为 $\overline{X}$ 和 $Y'$ 在 $Y$ 上的
导出纤维积，就可以完全像引理 \ref{duality-lemma-base-change-shriek-flat}
的证明那样论证并定义该映射。毕竟，$X$ 与 $Y'$ 的 Tor 独立性保证
$X'$ 是导出概形 $\overline{X}'$ 的开子概形。}。

\medskip\noindent
第 \ref{duality-section-relative-dualizing-complexes} 节将给出一种部分补救。
也就是说，若态射 $f$ 平坦，就有合适的相对对偶复形概念，并可利用引理
\ref{duality-lemma-compactifyable-relative-dualizing}、
\ref{duality-lemma-base-change-relative-dualizing} 和
\ref{duality-lemma-perfect-comparison-shriek} 构造典范基变换同构。
如果确实需要使用这一点，将在本章后文补充精确陈述与证明。

\begin{lemma}
\label{duality-lemma-base-change-shriek-flat}
在情形 \ref{duality-situation-shriek} 中，设
$$
\xymatrix{
X' \ar[r]_{g'} \ar[d]_{f'} & X \ar[d]^f \\
Y' \ar[r]^g & Y
}
$$
是 $\textit{FTS}_S$ 中的笛卡尔图，且 $g$ 平坦。则在
$D_\QCoh^+(\mathcal{O}_Y)$ 上存在同构
$L(g')^* \circ f^! \to (f')^! \circ Lg^*$。
\end{lemma}

\begin{proof}
事实上，因为 $g$ 平坦，所以对 $X$ 在 $Y$ 上的紧化
$j : X \to \overline{X}$ 的每个选择，概形 $\overline{X}$ 都与
$Y'$ Tor 独立。记 $j' : X' \to \overline{X}'$ 为 $j$ 的基变换，
并记 $\overline{g}' : \overline{X}' \to \overline{X}$ 为投影。
将基变换映射定义为复合
$$
L(g')^* \circ f^! = L(g')^* \circ j^* \circ a =
(j')^* \circ L(\overline{g}')^* \circ a \longrightarrow
(j')^* \circ a' \circ Lg^* = (f')^! \circ Lg^*
$$
，其中中间箭头是基变换映射 (\ref{duality-equation-base-change-map})，
而 $a$ 和 $a'$ 是引理 \ref{duality-lemma-twisted-inverse-image}
分别对 $\overline{X} \to Y$ 和 $\overline{X}' \to Y'$ 给出的推前的右伴随。
此构造与紧化的选择无关（如果确实需要这一结果，将陈述一个精确引理并予以证明）。

\medskip\noindent
为完成证明，只需说明基变换映射
$L(g')^* \circ a \to a' \circ Lg^*$ 在
$D_\QCoh^+(\mathcal{O}_Y)$ 上是同构。由引理
\ref{duality-lemma-proper-noetherian}，$a$ 和 $a'$ 的形成与限制到
$Y$ 和 $Y'$ 的仿射开集可交换。因此由注
\ref{duality-remark-check-over-affines}，可设 $Y$ 和 $Y'$ 是仿射的。
因此结论由引理 \ref{duality-lemma-more-base-change} 得到。
\end{proof}

\begin{lemma}
\label{duality-lemma-shriek-etale}
在情形 \ref{duality-situation-shriek} 中，设 $f : X \to Y$ 是
$\textit{FTS}_S$ 的 \'etale 态射。则在
$D^+_\QCoh(\mathcal{O}_Y)$ 上，函子满足 $f^! \cong f^*$。
\end{lemma}

\begin{proof}
将使用如下事实：\'etale 态射是平坦的、syntomic 的，并且是局部完全交态射
（《态射》，引理 \ref{morphisms-lemma-etale-syntomic} 和
\ref{morphisms-lemma-etale-flat}，以及《态射进阶》，引理
\ref{more-morphisms-lemma-flat-lci}）。由引理
\ref{duality-lemma-perfect-comparison-shriek}，只需证明
$f^!\mathcal{O}_Y = \mathcal{O}_X$。由引理
\ref{duality-lemma-lci-shriek}，已知 $f^!\mathcal{O}_Y$ 是可逆模。考虑交换图
$$
\xymatrix{
X \times_Y X \ar[r]_{p_2} \ar[d]_{p_1} & X \ar[d]^f \\
X \ar[r]^f & Y
}
$$
和对角态射 $\Delta : X \to X \times_Y X$。由于 $\Delta$ 是开浸入
（由《态射》，引理
\ref{morphisms-lemma-diagonal-unramified-morphism} 和
\ref{morphisms-lemma-etale-smooth-unramified}），故由引理
\ref{duality-lemma-shriek-open-immersion}，有 $\Delta^! = \Delta^*$。
由引理 \ref{duality-lemma-upper-shriek-composition}，有
$\Delta^! \circ p_1^! \circ f^! = f^!$。将引理
\ref{duality-lemma-base-change-shriek-flat} 应用于该图，得到
$p_1^!\mathcal{O}_X = p_2^*f^!\mathcal{O}_Y$。因此
$$
f^!\mathcal{O}_Y = \Delta^!p_1^!f^!\mathcal{O}_Y =
\Delta^*(p_1^*f^!\mathcal{O}_Y \otimes p_1^!\mathcal{O}_X) =
\Delta^*(p_2^*f^!\mathcal{O}_Y \otimes p_1^*f^!\mathcal{O}_Y) =
(f^!\mathcal{O}_Y)^{\otimes 2}
$$
，其中第二步再次使用了引理
\ref{duality-lemma-perfect-comparison-shriek}。故如所需，
$f^!\mathcal{O}_Y = \mathcal{O}_X$。
\end{proof}


\noindent
在本节余下部分，将陈述若干容易证明的结果；它们应当是良好的基变换映射理论的推论。

\begin{lemma}[权宜基变换]
\label{duality-lemma-base-change-locally}
在情形 \ref{duality-situation-shriek} 中，设
$$
\xymatrix{
X' \ar[r]_{g'} \ar[d]_{f'} & X \ar[d]^f \\
Y' \ar[r]^g & Y
}
$$
是 $\textit{FTS}_S$ 中的笛卡尔图。设
$E \in D^+_\QCoh(\mathcal{O}_Y)$ 是对象，且
$Lg^*E$ 属于 $D^+(\mathcal{O}_Y)$。若 $f$ 平坦，则对映入
$Y$、$Y'$ 和 $X$ 的仿射开集的仿射开集 $U' \subset X'$，
$L(g')^*f^!E$ 与 $(f')^!Lg^*E$ 限制后是
$D(\mathcal{O}_{U'})$ 中的同构对象。
\end{lemma}

\begin{proof}
由假设，立即归约到 $X$、$Y$、$Y'$、$X'$ 都仿射的情形。
设 $Y = \Spec(R)$、$Y' = \Spec(R')$、$X = \Spec(A)$、
$X' = \Spec(A')$。则 $A' = A \otimes_R R'$。设 $E$ 对应于
$K \in D^+(R)$。记给定映射为 $\varphi : R \to A$ 和
$\varphi' : R' \to A'$。由注
\ref{duality-remark-local-calculation-shriek}，可知 $L(g')^*f^!E$ 和
$(f')^!Lg^*E$ 分别对应于
$\varphi^!(K) \otimes_A^\mathbf{L} A'$ 和
$(\varphi')^!(K \otimes_R^\mathbf{L} R')$，其中
$\varphi^!$ 和 $(\varphi')^!$ 是《对偶复形》，第
\ref{dualizing-section-relative-dualizing-complex-algebraic} 节的函子。
结论由《对偶复形》，引理 \ref{dualizing-lemma-bc-flat} 得到。
\end{proof}

\begin{lemma}
\label{duality-lemma-relative-dualizing-fibres}
在情形 \ref{duality-situation-shriek} 中，设 $f : X \to Y$ 是
$\textit{FTS}_S$ 的态射。假设 $f$ 平坦。在
$D^b_{\textit{Coh}}(X)$ 中令
$\omega_{X/Y}^\bullet = f^!\mathcal{O}_Y$。设 $y \in Y$，且
$h : X_y \to X$ 是投影。则 $Lh^*\omega_{X/Y}^\bullet$
是 $X_y$ 上的对偶复形。
\end{lemma}

\begin{proof}
由引理 \ref{duality-lemma-perfect-comparison-shriek}，复形
$\omega_{X/Y}^\bullet$ 属于 $D^b_{\textit{Coh}}$。
成为对偶复形是局部性质。因此由引理
\ref{duality-lemma-base-change-locally}，只需证明
$(X_y \to y)^!\mathcal{O}_y$ 是 $X_y$ 上的对偶复形。
这由引理 \ref{duality-lemma-shriek-dualizing} 得到。
\end{proof}








\section{一种对偶性理论}
\label{duality-section-duality}

\noindent
本节具体说明：上述非常一般的结果对固定 Noetherian 基概形上的有限型分离概形
给出何种对偶性理论。

\medskip\noindent
回忆 Noetherian 概形 $X$ 上的对偶复形，是 $D(\mathcal{O}_X)$
中的对象，并在每个仿射局部给出相应环的对偶复形；见定义
\ref{duality-definition-dualizing-scheme}。

\medskip\noindent
给定 Noetherian 概形 $S$，用 $\textit{FTS}_S$ 表示 $S$ 上有限型且分离的
概形所成的范畴。则：
\begin{enumerate}
\item 函子 $f^!$ 使 $D_\QCoh^+$ 成为 $\textit{FTS}_S$ 上的伪函子；
\item 若 $f : X \to Y$ 是 $\textit{FTS}_S$ 中的真态射，则 $f^!$ 是
$Rf_* : D_\QCoh(\mathcal{O}_X) \to D_\QCoh(\mathcal{O}_Y)$
的右伴随在 $D_\QCoh^+(\mathcal{O}_Y)$ 上的限制，并且对所有
$K \in D_{\textit{Coh}}^-(\mathcal{O}_X)$ 和
$M \in D_\QCoh^+(\mathcal{O}_Y)$，存在典范同构
$$
Rf_*R\SheafHom_{\mathcal{O}_X}(K, f^!M)
\to
R\SheafHom_{\mathcal{O}_Y}(Rf_*K, M)
$$
；
\item 若 $\textit{FTS}_S$ 的对象 $X$ 有对偶复形
$\omega_X^\bullet$，则函子
$D_X = R\SheafHom_{\mathcal{O}_X}(-, \omega_X^\bullet)$
定义 $D_{\textit{Coh}}(\mathcal{O}_X)$ 的一个对合，它交换
$D_{\textit{Coh}}^+(\mathcal{O}_X)$ 与
$D_{\textit{Coh}}^-(\mathcal{O}_X)$，并固定
$D_{\textit{Coh}}^b(\mathcal{O}_X)$；
\item 若 $f : X \to Y$ 是 $\textit{FTS}_S$ 的态射，且
$\omega_Y^\bullet$ 是 $Y$ 上的对偶复形，则
\begin{enumerate}
\item $\omega_X^\bullet = f^!\omega_Y^\bullet$ 是 $X$ 的对偶复形；
\item 对 $M \in D_{\textit{Coh}}^+(\mathcal{O}_Y)$，典范地有
$f^!M = D_X(Lf^*D_Y(M))$；并且
\item 若另外 $f$ 是真态射，则对 $K$，若其属于
$D^-_{\textit{Coh}}(\mathcal{O}_X)$，有
$$
Rf_*R\SheafHom_{\mathcal{O}_X}(K, \omega_X^\bullet) =
R\SheafHom_{\mathcal{O}_Y}(Rf_*K, \omega_Y^\bullet)
$$
；
\end{enumerate}
\item 若 $f : X \to Y$ 是 $\textit{FTS}_S$ 中的闭浸入，则
$f^!(-) = R\SheafHom(\mathcal{O}_X, -)$；
\item 若 $f : Y \to X$ 是 $\textit{FTS}_S$ 中的有限态射，则
$f_*f^!(-) = R\SheafHom_{\mathcal{O}_X}(f_*\mathcal{O}_Y, -)$；
\item 若 $f : X \to Y$ 是一个有效 Cartier 除子到
$\textit{FTS}_S$ 某对象中的包含，则
$f^!(-) = Lf^*(-) \otimes_{\mathcal{O}_X} f^*\mathcal{O}_Y(X)[-1]$；
\item 若 $f : X \to Y$ 是余维数为 $c$ 的 Koszul 正则浸入，
其目标是 $\textit{FTS}_S$ 的对象，则
$f^!(-) \cong Lf^*(-) \otimes_{\mathcal{O}_X} \wedge^c\mathcal{N}[-c]$；并且
\item 若 $f : X \to Y$ 是 $\textit{FTS}_S$ 中相对维数为 $d$ 的光滑真态射，则
$f^!(-) \cong Lf^*(-)  \otimes_{\mathcal{O}_X} \Omega^d_{X/Y}[d]$。
\end{enumerate}
这些结论由引理
\ref{duality-lemma-dualizing-schemes}、
\ref{duality-lemma-iso-on-RSheafHom}、
\ref{duality-lemma-twisted-inverse-image-closed}、
\ref{duality-lemma-finite-twisted}、
\ref{duality-lemma-sheaf-with-exact-support-effective-Cartier}、
\ref{duality-lemma-regular-immersion}、
\ref{duality-lemma-smooth-proper}、
\ref{duality-lemma-upper-shriek-composition}、
\ref{duality-lemma-pseudo-functor}、
\ref{duality-lemma-shriek-closed-immersion}、
\ref{duality-lemma-shriek-dualizing}、
\ref{duality-lemma-shriek-via-duality}、
\ref{duality-lemma-perfect-comparison-shriek}，以及例
\ref{duality-example-iso-on-RSheafHom-noetherian} 得到。
这些函子是通过一个非常抽象的程序得到的，而该程序最终依赖于调用一个存在性定理
（《导出范畴》，命题 \ref{derived-proposition-brown}）。这意味着一般而言，
没有函子 $f^!$ 的显式描述。这有时会造成问题。但事实上，知道对偶复形的存在性和
对偶同构，往往已经足以唯一确定 $f^!$。






\section{对偶复形的黏合}
\label{duality-section-glue}

\noindent
现在利用对偶复形的黏合，为 $S$ 上所有有限型概形得到一种理论，
其中给定情形 \ref{duality-situation-dualizing} 中那样的偶对
$(S, \omega_S^\bullet)$。这类似于 \cite[Remark on page 310]{RD}。

\begin{situation}
\label{duality-situation-dualizing}
这里 $S$ 是 Noetherian 概形，$\omega_S^\bullet$ 是对偶复形。
\end{situation}

\noindent
在情形 \ref{duality-situation-dualizing} 中，设 $X$ 是 $S$ 上的有限型概形。
设 $\mathcal{U} : X = \bigcup_{i = 1, \ldots, n} U_i$ 是由
$\textit{FTS}_S$ 的对象组成的 $X$ 的有限开覆盖；见情形
\ref{duality-situation-shriek}。这仅仅意味着态射 $U_i \to S$ 是分离的
（它们已经是有限型的）。由《概形》，引理
\ref{schemes-lemma-compose-after-separated}，$S$ 上每个有限型仿射概形
都是 $\textit{FTS}_S$ 的对象，故这样的开覆盖当然存在。
于是对每个 $i, j, k \in \{1, \ldots, n\}$，态射
$p_i : U_i \to S$、$p_{ij} : U_i \cap U_j \to S$、
$p_{ijk} : U_i \cap U_j \cap U_k \to S$ 都是分离的，
而且这些概形各自都是 $\textit{FTS}_S$ 的对象。由这样的开覆盖得到
\begin{enumerate}
\item $\omega_i^\bullet = p_i^!\omega_S^\bullet$，它是 $U_i$ 上的对偶复形；
见第 \ref{duality-section-duality} 节；
\item 对每个 $i, j$，有典范同构
$\varphi_{ij} :
\omega_i^\bullet|_{U_i \cap U_j} \to \omega_j^\bullet|_{U_i \cap U_j}$；并且
\item
\label{duality-item-cocycle-glueing}
对每个 $i, j, k$，有
$$
\varphi_{ik}|_{U_i \cap U_j \cap U_k} =
\varphi_{jk}|_{U_i \cap U_j \cap U_k} \circ
\varphi_{ij}|_{U_i \cap U_j \cap U_k}
$$
，该等式位于 $D(\mathcal{O}_{U_i \cap U_j \cap U_k})$ 中。
\end{enumerate}
这里，在 (2) 中使用了 $(U_i \cap U_j \to U_i)^!$ 由限制给出
（引理 \ref{duality-lemma-shriek-open-immersion}），以及由引理
\ref{duality-lemma-upper-shriek-composition} 得到的典范同构
$$
(U_i \cap U_j \to U_i)^! \circ p_i^! = p_{ij}^! =
(U_i \cap U_j \to U_j)^! \circ p_j^!
$$
；为得到 (3)，使用了引理 \ref{duality-lemma-pseudo-functor} 所说明的上叹号函子构成伪函子。

\medskip\noindent
在刚刚描述的情形中，一个{\it 相对于 $\omega_S^\bullet$ 和
$\mathcal{U}$ 规范化的对偶复形}是偶对 $(K, \alpha_i)$，其中
$K \in D(\mathcal{O}_X)$，而
$\alpha_i : K|_{U_i} \to \omega_i^\bullet$ 是同构，并且
$\varphi_{ij}$ 由
$\alpha_j|_{U_i \cap U_j} \circ \alpha_i^{-1}|_{U_i \cap U_j}$ 给出。
由于成为概形上的对偶复形是局部性质，可知相对于
$\omega_S^\bullet$ 和 $\mathcal{U}$ 规范化的对偶复形确实是对偶复形。

\begin{lemma}
\label{duality-lemma-good-dualizing-unique}
在情形 \ref{duality-situation-dualizing} 中，设 $X$ 是 $S$ 上的有限型概形，
而 $\mathcal{U}$ 是由 $S$ 上分离概形组成的 $X$ 的有限开覆盖。
若存在相对于 $\omega_S^\bullet$ 和 $\mathcal{U}$ 规范化的对偶复形，
则它在唯一同构意义下唯一。
\end{lemma}

\begin{proof}
若 $(K, \alpha_i)$ 和 $(K', \alpha_i')$ 是两个这样的复形，则考虑
$L = R\SheafHom_{\mathcal{O}_X}(K, K')$。由引理
\ref{duality-lemma-dualizing-unique-schemes} 及其证明，这是
$D(\mathcal{O}_X)$ 中的可逆对象。利用 $\alpha_i$ 和 $\alpha'_i$，得到同构
$$
\alpha_i^t \otimes \alpha'_i :
L|_{U_i} \longrightarrow
R\SheafHom_{\mathcal{O}_X}(\omega_i^\bullet, \omega_i^\bullet) =
\mathcal{O}_{U_i}[0]
$$
。这已经蕴含在 $D(\mathcal{O}_X)$ 中有 $L = H^0(L)[0]$。
此外，$H^0(L)$ 是可逆层，并在 $X$ 的开集 $U_i$ 上带有给定的平凡化。
最后，
$\alpha_j|_{U_i \cap U_j} \circ \alpha_i^{-1}|_{U_i \cap U_j}$
和
$\alpha'_j|_{U_i \cap U_j} \circ (\alpha'_i)^{-1}|_{U_i \cap U_j}$
都给出 $\varphi_{ij}$ 这一条件，蕴含转移映射均为 $1$，因而得到同构
$H^0(L) = \mathcal{O}_X$。
\end{proof}

\begin{lemma}
\label{duality-lemma-good-dualizing-independence-covering}
在情形 \ref{duality-situation-dualizing} 中，设 $X$ 是 $S$ 上的有限型概形，
而 $\mathcal{U}$、$\mathcal{V}$ 是由 $S$ 上分离概形组成的
$X$ 的两个有限开覆盖。若存在相对于 $\omega_S^\bullet$ 和
$\mathcal{U}$ 规范化的对偶复形，则存在相对于
$\omega_S^\bullet$ 和 $\mathcal{V}$ 规范化的对偶复形，
并且这些复形典范同构。
\end{lemma}

\begin{proof}
只需在 $\mathcal{U}$ 由开集 $U_1, \ldots, U_n$ 给出、而
$\mathcal{V}$ 由开集 $U_1, \ldots, U_{n + m}$ 给出时证明。
事实上，可以并且确实假设 $m = 1$。从相对于
$\omega_S^\bullet$ 和 $\mathcal{V}$ 规范化的对偶复形
$(K, \alpha_i)$ 得到相对于 $\omega_S^\bullet$ 和 $\mathcal{U}$
规范化的对偶复形，只需忘掉 $i = n + 1$ 时的 $\alpha_i$。
反过来，设 $(K, \alpha_i)$ 是相对于 $\omega_S^\bullet$ 和
$\mathcal{U}$ 规范化的对偶复形。为完成证明，需要构造满足所需条件的映射
$\alpha_{n + 1} : K|_{U_{n + 1}} \to \omega_{n + 1}^\bullet$。
为此注意到 $U_{n + 1} = \bigcup U_i \cap U_{n + 1}$ 是开覆盖。
显然，$(K|_{U_{n + 1}}, \alpha_i|_{U_i \cap U_{n + 1}})$ 是相对于
$\omega_S^\bullet$ 和覆盖
$U_{n + 1} = \bigcup U_i \cap U_{n + 1}$ 规范化的对偶复形。
另一方面，由条件 (\ref{duality-item-cocycle-glueing})，偶对
$(\omega_{n + 1}^\bullet|_{U_{n + 1}}, \varphi_{n + 1i})$
是另一个相对于 $\omega_S^\bullet$ 和覆盖
$U_{n + 1} = \bigcup U_i \cap U_{n + 1}$ 规范化的对偶复形。
由引理 \ref{duality-lemma-good-dualizing-unique}，得到与给定局部同构相容的唯一同构
$$
\alpha_{n + 1} : K|_{U_{n + 1}} \longrightarrow \omega_{n + 1}^\bullet
$$
。证明这意味着它满足所需性质，是一个令人愉快的练习。
\end{proof}

\begin{lemma}
\label{duality-lemma-existence-good-dualizing}
在情形 \ref{duality-situation-dualizing} 中，设 $X$ 是 $S$ 上的有限型概形，
而 $\mathcal{U}$ 是由 $S$ 上分离概形组成的 $X$ 的有限开覆盖。
则存在相对于 $\omega_S^\bullet$ 和 $\mathcal{U}$ 规范化的对偶复形。
\end{lemma}

\begin{proof}
设 $\mathcal{U} : X = \bigcup_{i = 1, \ldots, n} U_i$。
对 $n$ 作归纳来证明该引理。基础情形 $n = 1$ 是立即的。
假设 $n > 1$。令 $X' = U_1 \cup \ldots \cup U_{n - 1}$，
并令 $(K', \{\alpha'_i\}_{i = 1, \ldots, n - 1})$ 是相对于
$\omega_S^\bullet$ 和
$\mathcal{U}' : X' = \bigcup_{i = 1, \ldots, n - 1} U_i$
规范化的对偶复形。显然，
$(K'|_{X' \cap U_n}, \alpha'_i|_{U_i \cap U_n})$ 是相对于
$\omega_S^\bullet$ 和覆盖
$X' \cap U_n = \bigcup_{i = 1, \ldots, n - 1} U_i \cap U_n$
规范化的对偶复形。另一方面，由条件
(\ref{duality-item-cocycle-glueing})，偶对
$(\omega_n^\bullet|_{X' \cap U_n}, \varphi_{ni})$ 是另一个相对于
$\omega_S^\bullet$ 和覆盖
$X' \cap U_n = \bigcup_{i = 1, \ldots, n - 1} U_i \cap U_n$
规范化的对偶复形。由引理 \ref{duality-lemma-good-dualizing-unique}，
得到与给定局部同构相容的唯一同构
$$
\epsilon : K'|_{X' \cap U_n} \longrightarrow \omega_i^\bullet|_{X' \cap U_n}
$$
。由《上同调》，引理 \ref{cohomology-lemma-glue}，得到
$K \in D(\mathcal{O}_X)$，以及同构
$\beta : K|_{X'} \to K'$ 和
$\gamma : K|_{U_n} \to \omega_n^\bullet$，使得
$\epsilon = \gamma|_{X'\cap U_n} \circ \beta|_{X' \cap U_n}^{-1}$。
然后定义
$$
\alpha_i = \alpha'_i \circ \beta|_{U_i}, i = 1, \ldots, n - 1,
\text{ 且 }
\alpha_n = \gamma
$$
。还需验证 $\varphi_{ij}$ 由
$\alpha_j|_{U_i \cap U_j} \circ \alpha_i^{-1}|_{U_i \cap U_j}$ 给出。
当 $i, j \leq n - 1$ 时，这由 $\alpha_i'$ 的相应条件得到；
当 $i = j = n$ 时也很清楚。若 $i < j = n$，则
$$
\alpha_n|_{U_i \cap U_n} \circ \alpha_i^{-1}|_{U_i \cap U_n} =
\gamma|_{U_i \cap U_n} \circ \beta^{-1}|_{U_i \cap U_n}
\circ (\alpha'_i)^{-1}|_{U_i \cap U_n} =
\epsilon|_{U_i \cap U_n} \circ (\alpha'_i)^{-1}|_{U_i \cap U_n}
$$
。这恰好等于 $\alpha_{in}$，因为 $\epsilon$ 是与映射
$\alpha_i'$ 和 $\alpha_{ni}$ 相容的唯一映射。
\end{proof}

\noindent
设 $(S, \omega_S^\bullet)$ 如情形 \ref{duality-situation-dualizing} 所述。
上述引理的要点是：给定任意 $S$ 上的有限型概形 $X$，存在在唯一同构意义下
给定的偶对 $(K, \alpha_U)$；它由对象 $K \in D(\mathcal{O}_X)$，
以及对每个在 $S$ 上分离的开子概形 $U \subset X$ 的同构
$\alpha_U : K|_U \to \omega_U^\bullet$ 组成。这里
$\omega_U^\bullet = (U \to S)^!\omega_S^\bullet$ 是 $U$ 上的对偶复形；
见第 \ref{duality-section-duality} 节。此外，若
$\mathcal{U} : X = \bigcup U_i$ 是由在 $S$ 上分离的开集组成的有限开覆盖，
则 $(K, \alpha_{U_i})$ 是相对于 $\omega_S^\bullet$ 和
$\mathcal{U}$ 规范化的对偶复形。事实上，在唯一同构意义下的唯一性由引理
\ref{duality-lemma-good-dualizing-unique} 给出；对一个开覆盖的存在性由引理
\ref{duality-lemma-existence-good-dualizing} 给出；而 $K$ 随后适用于所有开覆盖这一事实
由引理 \ref{duality-lemma-good-dualizing-independence-covering} 给出。

\begin{definition}
\label{duality-definition-good-dualizing}
设 $S$ 是 Noetherian 概形，$\omega_S^\bullet$ 是 $S$ 上的对偶复形。
设 $X$ 是 $S$ 上的有限型概形。上面构造的复形 $K$ 称为
{\it 相对于 $\omega_S^\bullet$ 规范化的对偶复形}，并记作
$\omega_X^\bullet$。
\end{definition}

\noindent
正如术语所示，相对于 $\omega_S^\bullet$ 规范化的对偶复形不仅是
$X$ 的导出范畴中的对象，还配备了上面描述的局部同构。若 $X$ 在 $S$ 上分离，
这与令 $\omega_X^\bullet = p^!\omega_S^\bullet$ 并不冲突，
其中 $p : X \to S$ 是结构态射。更一般地，有下面的合理性检验。

\begin{lemma}
\label{duality-lemma-good-over-both}
设 $(S, \omega_S^\bullet)$ 如情形 \ref{duality-situation-dualizing} 所述。
设 $f : X \to Y$ 是 $S$ 上有限型概形之间的态射。
设 $\omega_X^\bullet$ 和 $\omega_Y^\bullet$ 是相对于
$\omega_S^\bullet$ 规范化的对偶复形。则 $\omega_X^\bullet$
是相对于 $\omega_Y^\bullet$ 规范化的对偶复形。
\end{lemma}

\begin{proof}
这只是一个记账问题。选取有限仿射开覆盖
$\mathcal{V} : Y = \bigcup V_j$。对每个 $j$，选取有限仿射开覆盖
$f^{-1}(V_j) = U_{ji}$。令 $\mathcal{U} : X = \bigcup U_{ji}$。
概形 $V_j$ 和 $U_{ji}$ 在 $S$ 上分离，故对
$q_j : V_j \to S$、$p_{ji} : U_{ji} \to S$、
$f_{ji} : U_{ji} \to V_j$、$f_{ji}' : U_{ji} \to Y$
都有上叹号函子。令 $(L, \beta_j)$ 是相对于
$\omega_S^\bullet$ 和 $\mathcal{V}$ 规范化的对偶复形。
令 $(K, \gamma_{ji})$ 是相对于 $\omega_S^\bullet$ 和
$\mathcal{U}$ 规范化的对偶复形。（换言之，
$L = \omega_Y^\bullet$ 且 $K = \omega_X^\bullet$。）可定义
$$
\alpha_{ji} :
K|_{U_{ji}} \xrightarrow{\gamma_{ji}}
p_{ji}^!\omega_S^\bullet = f_{ji}^!q_j^!\omega_S^\bullet
\xrightarrow{f_{ji}^!\beta_j^{-1}} f_{ji}^!(L|_{V_j}) =
(f_{ji}')^!(L)
$$
。为完成证明，必须说明
$\alpha_{ji}|_{U_{ji} \cap U_{j'i'}}
\circ \alpha_{j'i'}^{-1}|_{U_{ji} \cap U_{j'i'}}$
是典范同构
$(f_{ji}')^!(L)|_{U_{ji} \cap U_{j'i'}} \to
(f_{j'i'}')^!(L)|_{U_{ji} \cap U_{j'i'}}$。这是形式的，略去细节。
\end{proof}

\begin{lemma}
\label{duality-lemma-open-immersion-good-dualizing-complex}
设 $(S, \omega_S^\bullet)$ 如情形 \ref{duality-situation-dualizing} 所述。
设 $j : X \to Y$ 是 $S$ 上有限型概形之间的开浸入。
设 $\omega_X^\bullet$ 和 $\omega_Y^\bullet$ 是相对于
$\omega_S^\bullet$ 规范化的对偶复形。则存在典范同构
$\omega_X^\bullet = \omega_Y^\bullet|_X$。
\end{lemma}

\begin{proof}
这由定义 \ref{duality-definition-good-dualizing} 紧前面给出的规范化对偶复形构造立即得到。
\end{proof}

\begin{lemma}
\label{duality-lemma-proper-map-good-dualizing-complex}
设 $(S, \omega_S^\bullet)$ 如情形 \ref{duality-situation-dualizing} 所述。
设 $f : X \to Y$ 是 $S$ 上有限型概形之间的真态射。
设 $\omega_X^\bullet$ 和 $\omega_Y^\bullet$ 是相对于
$\omega_S^\bullet$ 规范化的对偶复形。令 $a$ 为引理
\ref{duality-lemma-twisted-inverse-image} 对 $f$ 给出的右伴随。
则存在典范同构 $a(\omega_Y^\bullet) = \omega_X^\bullet$。
\end{lemma}

\begin{proof}
设 $p : X \to S$ 和 $q : Y \to S$ 是结构态射。若 $X$ 和 $Y$
在 $S$ 上分离，则结论来自
$\omega_X^\bullet = p^!\omega_S^\bullet$、
$\omega_Y^\bullet = q^!\omega_S^\bullet$、$f^! = a$，以及
$f^! \circ q^! = p^!$（引理 \ref{duality-lemma-upper-shriek-composition}）。
在一般情形下，先使用引理 \ref{duality-lemma-good-over-both} 归约到
$Y = S$。此时 $X$ 和 $Y$ 在 $S$ 上分离，而刚才已经看到了结论。
\end{proof}

\noindent
设 $(S, \omega_S^\bullet)$ 如情形 \ref{duality-situation-dualizing} 所述。
对 $S$ 上的有限型概形 $X$，用 $\omega_X^\bullet$ 表示 $X$ 的
相对于 $\omega_S^\bullet$ 规范化的对偶复形。像引理
\ref{duality-lemma-dualizing-schemes} 中那样定义
$D_X(-) = R\SheafHom_{\mathcal{O}_X}(-, \omega_X^\bullet)$。
设 $f : X \to Y$ 是 $S$ 上有限型概形之间的态射。定义
$$
f_{new}^! = D_X \circ Lf^* \circ D_Y :
D_{\textit{Coh}}^+(\mathcal{O}_Y)
\to
D_{\textit{Coh}}^+(\mathcal{O}_X)
$$
。若 $f : X \to Y$ 和 $g : Y \to Z$ 是 $S$ 上有限型概形之间
可复合的态射，则定义
\begin{align*}
(g \circ f)^!_{new} & = D_X \circ L(g \circ f)^* \circ D_Z \\
& = D_X \circ Lf^* \circ Lg^* \circ D_Z \\
& \to D_X \circ Lf^* \circ D_Y \circ D_Y \circ Lg^* \circ D_Z \\
& = f^!_{new} \circ g^!_{new}
\end{align*}
，其中箭头在引理 \ref{duality-lemma-dualizing-schemes} 中定义。
下面的引理汇集这些结果。

\begin{lemma}
\label{duality-lemma-duality-bootstrap}
设 $(S, \omega_S^\bullet)$ 如情形 \ref{duality-situation-dualizing} 所述。
按上面的方式，对 $S$ 上所有有限型概形（的态射）定义
$f^!_{new}$ 和 $\omega_X^\bullet$。则：
\begin{enumerate}
\item 函子 $f^!_{new}$ 和箭头
$(g \circ f)^!_{new} \to f^!_{new} \circ g^!_{new}$
使 $D_{\textit{Coh}}^+$ 成为从 $S$ 上有限型概形的范畴到范畴的
$2$-范畴的伪函子；
\item $\omega_X^\bullet = (X \to S)^!_{new} \omega_S^\bullet$；
\item 函子 $D_X$ 定义 $D_{\textit{Coh}}(\mathcal{O}_X)$ 的一个对合，
它交换 $D_{\textit{Coh}}^+(\mathcal{O}_X)$ 与
$D_{\textit{Coh}}^-(\mathcal{O}_X)$，并固定
$D_{\textit{Coh}}^b(\mathcal{O}_X)$；
\item 对 $S$ 上有限型概形之间的态射 $f : X \to Y$，有
$\omega_X^\bullet = f^!_{new}\omega_Y^\bullet$；
\item 对 $M \in D_{\textit{Coh}}^+(\mathcal{O}_Y)$，有
$f^!_{new}M = D_X(Lf^*D_Y(M))$；并且
\item 若另外 $f$ 是真态射，则 $f^!_{new}$ 同构于
$Rf_* : D_\QCoh(\mathcal{O}_X) \to D_\QCoh(\mathcal{O}_Y)$
的右伴随在 $D_{\textit{Coh}}^+(\mathcal{O}_Y)$ 上的限制，
而且对 $K \in D^-_{\textit{Coh}}(\mathcal{O}_X)$ 和
$M \in D_{\textit{Coh}}^+(\mathcal{O}_Y)$，存在典范同构
$$
Rf_*R\SheafHom_{\mathcal{O}_X}(K, f^!_{new}M)
\to
R\SheafHom_{\mathcal{O}_Y}(Rf_*K, M)
$$
，并且对 $K \in D^-_{\textit{Coh}}(\mathcal{O}_X)$，有
$$
Rf_*R\SheafHom_{\mathcal{O}_X}(K, \omega_X^\bullet) =
R\SheafHom_{\mathcal{O}_Y}(Rf_*K, \omega_Y^\bullet)
$$
，并且
\end{enumerate}
若 $X$ 在 $S$ 上分离，则 $\omega_X^\bullet$ 典范同构于
$(X \to S)^!\omega_S^\bullet$；若 $f$ 是 $S$ 上分离概形之间的态射，
则存在典范同构\footnote{尚未检查这些同构与
$(g \circ f)^! \to f^! \circ g^!$ 和
$(g \circ f)^!_{new} \to f^!_{new} \circ g^!_{new}$ 相容。
若后文需要，将在这里完成检查。}
$f_{new}^!K = f^!K$，其中 $K$ 属于 $D_{\textit{Coh}}^+$。
\end{lemma}

\begin{proof}
设 $f : X \to Y$、$g : Y \to Z$、$h : Z \to T$ 是
$S$ 上有限型概形之间的态射。必须证明图
$$
\xymatrix{
(h \circ g \circ f)^!_{new} \ar[r] \ar[d] &
f^!_{new} \circ (h \circ g)^!_{new} \ar[d] \\
(g \circ f)^!_{new} \circ h^!_{new} \ar[r] &
f^!_{new} \circ g^!_{new} \circ h^!_{new}
}
$$
交换。令 $\eta_Y : \text{id} \to D_Y^2$ 和
$\eta_Z : \text{id} \to D_Z^2$ 为引理
\ref{duality-lemma-dualizing-schemes} 的典范同构。然后利用《范畴》，引理
\ref{categories-lemma-properties-2-cat-cats}，一个略去的计算表明，
两个箭头
$(h \circ g \circ f)^!_{new} \to f^!_{new} \circ g^!_{new} \circ h^!_{new}$
都由
$$
1 \star \eta_Y \star 1 \star \eta_Z \star 1 :
D_X \circ Lf^* \circ Lg^* \circ Lh^* \circ D_T
\longrightarrow
D_X \circ Lf^* \circ D_Y^2 \circ Lg^* \circ D_Z^2 \circ Lh^* \circ D_T
$$
给出。这证明了 (1)。(2) 由 $(X \to S)^!_{new}$ 的定义以及
$D_S(\omega_S^\bullet) = \mathcal{O}_S$ 这一事实立即得到。
(3) 是引理 \ref{duality-lemma-dualizing-schemes}。(4) 由与 (2) 相同的论证得到。
(5) 是 $f^!_{new}$ 的定义。

\medskip\noindent
证明 (6)。令 $a$ 为引理 \ref{duality-lemma-twisted-inverse-image}
对 $S$ 上有限型概形之间的真态射 $f : X \to Y$ 给出的右伴随。
问题在于不知道 $X$ 或 $Y$ 在 $S$ 上是否分离（一般确实不分离），
所以不能立即把引理 \ref{duality-lemma-shriek-via-duality} 应用于基概形
$S$ 上的 $f$。为绕过这一点，使用引理
\ref{duality-lemma-proper-map-good-dualizing-complex} 的典范等同
$\omega_X^\bullet = a(\omega_Y^\bullet)$。因此，将引理
\ref{duality-lemma-shriek-via-duality} 应用于以 $Y$ 为基概形的
$f : X \to Y$，可知 $f^!_{new}$ 是 $a$ 在
$D_{\textit{Coh}}^+(\mathcal{O}_Y)$ 上的限制！所显示的等式由例
\ref{duality-example-iso-on-RSheafHom-noetherian} 得到。

\medskip\noindent
最后的断言由规范化对偶复形的构造以及已经使用过的引理
\ref{duality-lemma-shriek-via-duality} 得到。
\end{proof}

\begin{remark}
\label{duality-remark-independent-omega-S}
设 $S$ 是具有对偶复形的 Noetherian 概形。设 $f : X \to Y$ 是
$S$ 上有限型概形之间的态射。则函子
$$
f_{new}^! : D^+_{Coh}(\mathcal{O}_Y) \to D^+_{Coh}(\mathcal{O}_X)
$$
在典范同构意义下与对偶复形 $\omega_S^\bullet$ 的选择无关。
下面略述证明。任意第二个对偶复形都形如
$\omega_S^\bullet \otimes_{\mathcal{O}_S}^\mathbf{L} \mathcal{L}$，
其中 $\mathcal{L}$ 是 $D(\mathcal{O}_S)$ 中的可逆对象；见引理
\ref{duality-lemma-dualizing-unique-schemes}。对任意有限型分离态射
$p : U \to S$，由引理 \ref{duality-lemma-compare-with-pullback-perfect}，有
$p^!(\omega_S^\bullet \otimes^\mathbf{L}_{\mathcal{O}_S} \mathcal{L}) =
p^!(\omega_S^\bullet) \otimes^\mathbf{L}_{\mathcal{O}_U} Lp^*\mathcal{L}$。
因此，若 $\omega_X^\bullet$ 和 $\omega_Y^\bullet$ 是相对于
$\omega_S^\bullet$ 规范化的对偶复形，则
$\omega_X^\bullet \otimes_{\mathcal{O}_X}^\mathbf{L} La^*\mathcal{L}$ 和
$\omega_Y^\bullet \otimes_{\mathcal{O}_Y}^\mathbf{L} Lb^*\mathcal{L}$
是相对于
$\omega_S^\bullet \otimes_{\mathcal{O}_S}^\mathbf{L} \mathcal{L}$
规范化的对偶复形（其中 $a : X \to S$ 和 $b : Y \to S$ 是结构态射）。
于是结论由下列等式得到：
\begin{align*}
& R\SheafHom_{\mathcal{O}_X}(Lf^*R\SheafHom_{\mathcal{O}_Y}(K,
\omega_Y^\bullet \otimes_{\mathcal{O}_Y}^\mathbf{L} Lb^*\mathcal{L}),
\omega_X^\bullet \otimes_{\mathcal{O}_X}^\mathbf{L} La^*\mathcal{L}) \\
& = R\SheafHom_{\mathcal{O}_X}(Lf^*R(\SheafHom_{\mathcal{O}_Y}(K,
\omega_Y^\bullet) \otimes_{\mathcal{O}_Y}^\mathbf{L} Lb^*\mathcal{L}),
\omega_X^\bullet \otimes_{\mathcal{O}_X}^\mathbf{L} La^*\mathcal{L}) \\
& = R\SheafHom_{\mathcal{O}_X}(Lf^*R\SheafHom_{\mathcal{O}_Y}(K,
\omega_Y^\bullet) \otimes_{\mathcal{O}_X}^\mathbf{L} La^*\mathcal{L},
\omega_X^\bullet \otimes_{\mathcal{O}_X}^\mathbf{L} La^*\mathcal{L}) \\
& = R\SheafHom_{\mathcal{O}_X}(Lf^*R\SheafHom_{\mathcal{O}_Y}(K,
\omega_Y^\bullet), \omega_X^\bullet)
\end{align*}
，其中 $K \in D^+_{Coh}(\mathcal{O}_Y)$。
最后一个等式成立，是因为 $La^*\mathcal{L}$ 在
$D(\mathcal{O}_X)$ 中可逆。
\end{remark}


\begin{example}
\label{duality-example-trace-proper}
设 $S$ 是 Noetherian 概形，$\omega_S^\bullet$ 是对偶复形。
设 $f : X \to Y$ 是 $S$ 上有限型概形之间的真态射。
设 $\omega_X^\bullet$ 和 $\omega_Y^\bullet$ 是相对于
$\omega_S^\bullet$ 规范化的对偶复形。在此情形下，有
$a(\omega_Y^\bullet) = \omega_X^\bullet$
（引理 \ref{duality-lemma-proper-map-good-dualizing-complex}），故迹映射
（第 \ref{duality-section-trace} 节）是典范箭头
$$
\text{Tr}_f : Rf_*\omega_X^\bullet \longrightarrow \omega_Y^\bullet
$$
，它产生同构（引理 \ref{duality-lemma-duality-bootstrap}）
$$
\Hom_X(L, \omega_X^\bullet) = \Hom_Y(Rf_*L, \omega_Y^\bullet)
$$
和
$$
Rf_*R\SheafHom_{\mathcal{O}_X}(L, \omega_X^\bullet) =
R\SheafHom_{\mathcal{O}_Y}(Rf_*L, \omega_Y^\bullet)
$$
，其中 $L$ 属于 $D_\QCoh(\mathcal{O}_X)$。
\end{example}

\begin{remark}
\label{duality-remark-dualizing-finite}
设 $S$ 是 Noetherian 概形，$\omega_S^\bullet$ 是对偶复形。
设 $f : X \to Y$ 是 $S$ 上有限型概形之间的有限态射。
设 $\omega_X^\bullet$ 和 $\omega_Y^\bullet$ 是相对于
$\omega_S^\bullet$ 规范化的对偶复形。由引理
\ref{duality-lemma-finite-twisted} 和
\ref{duality-lemma-proper-map-good-dualizing-complex}，在
$D_\QCoh^+(f_*\mathcal{O}_X)$ 中有
$$
f_*\omega_X^\bullet = R\SheafHom(f_*\mathcal{O}_X, \omega_Y^\bullet)
$$
，而例 \ref{duality-example-trace-proper} 的迹映射是映射
$$
\text{Tr}_f : Rf_*\omega_X^\bullet = f_*\omega_X^\bullet =
R\SheafHom(f_*\mathcal{O}_X, \omega_Y^\bullet) \longrightarrow
\omega_Y^\bullet
$$
，它常称为“在 $1$ 处求值”。
\end{remark}

\begin{remark}
\label{duality-remark-relative-dualizing-complex-shriek}
设 $f : X \to Y$ 是情形 \ref{duality-situation-dualizing} 中那样的偶对
$(S, \omega_S^\bullet)$ 上有限型概形之间的平坦真态射。
相对对偶复形（注 \ref{duality-remark-relative-dualizing-complex}）是
$\omega_{X/Y}^\bullet = a(\mathcal{O}_Y)$。由引理
\ref{duality-lemma-proper-map-good-dualizing-complex}，在
$$
\omega_X^\bullet = a(\omega_Y^\bullet) =
Lf^*\omega_Y^\bullet \otimes_{\mathcal{O}_X}^\mathbf{L} \omega_{X/Y}^\bullet
$$
中有第一个典范同构。该等式位于 $D(\mathcal{O}_X)$ 中。
第二个典范同构由注 \ref{duality-remark-relative-dualizing-complex} 中的讨论得到。
\end{remark}




\section{维数函数}
\label{duality-section-dimension-functions}

\noindent
需要进一步了解转到另一个概形上的有限型概形时，维数函数如何变化。

\begin{lemma}
\label{duality-lemma-good-dualizing-normalized}
设 $S$ 是 Noetherian 概形，$\omega_S^\bullet$ 是对偶复形。
设 $X$ 是 $S$ 上的有限型概形，$\omega_X^\bullet$ 是相对于
$\omega_S^\bullet$ 规范化的对偶复形。若 $x \in X$ 是位于
$S$ 的闭点 $s$ 之上的闭点，且 $\omega_{S, s}^\bullet$ 是
$\mathcal{O}_{S, s}$ 上的规范化对偶复形，则
$\omega_{X, x}^\bullet$ 是 $\mathcal{O}_{X, x}$ 上的规范化对偶复形。
\end{lemma}

\begin{proof}
可用 $x$ 的一个仿射邻域替换 $X$，故可以并且确实假设
$f : X \to S$ 分离。于是 $\omega_X^\bullet = f^!\omega_S^\bullet$。
必须说明 $R\Hom_{\mathcal{O}_{X, x}}(\kappa(x), \omega_{X, x}^\bullet)$
集中于次数 $0$。令 $i_x : x \to X$ 表示包含态射；由于 $x$ 是闭点，
它是闭浸入。因此由引理 \ref{duality-lemma-shriek-closed-immersion}，
$R\Hom_{\mathcal{O}_{X, x}}(\kappa(x), \omega_{X, x}^\bullet)$
表示 $i_x^!\omega_X^\bullet$。考虑交换图
$$
\xymatrix{
x \ar[r]_{i_x} \ar[d]_\pi & X \ar[d]^f \\
s \ar[r]^{i_s} & S
}
$$
。由《态射》，引理
\ref{morphisms-lemma-closed-point-fibre-locally-finite-type}，
扩张 $\kappa(x)/\kappa(s)$ 是有限的，因而 $\pi$ 是有限态射。所以
$$
i_x^!\omega_X^\bullet = i_x^! f^! \omega_S^\bullet =
\pi^! i_s^! \omega_S^\bullet
$$
。于是若 $\omega_{S, s}^\bullet$ 是 $\mathcal{O}_{S, s}$ 上的
规范化对偶复形，则由上面的同一论证，
$i_s^!\omega_S^\bullet = \kappa(s)[0]$。有
$$
R\pi_*(\pi^!(\kappa(s)[0])) =
R\SheafHom_{\mathcal{O}_s}(R\pi_*(\kappa(x)[0]), \kappa(s)[0]) =
\widetilde{\Hom_{\kappa(s)}(\kappa(x), \kappa(s))}
$$
。第一个等式来自例 \ref{duality-example-iso-on-RSheafHom-noetherian}，
其中取 $L = \kappa(x)[0]$。第二个等式成立，是因为 $\pi_*$ 正合。
因此 $\pi^!(\kappa(s)[0])$ 支撑在次数 $0$，结论成立。
\end{proof}

\begin{lemma}
\label{duality-lemma-good-dualizing-dimension-function}
设 $S$ 是 Noetherian 概形，$\omega_S^\bullet$ 是对偶复形。
设 $f : X \to S$ 是有限型态射，$\omega_X^\bullet$ 是相对于
$\omega_S^\bullet$ 规范化的对偶复形。对所有 $x \in X$，有
$$
\delta_X(x) - \delta_S(f(x)) = \text{trdeg}_{\kappa(f(x))}(\kappa(x))
$$
，其中 $\delta_S$\ 和 $\delta_X$\ 分别是
$\omega_S^\bullet$ 和 $\omega_X^\bullet$ 的维数函数；见引理
\ref{duality-lemma-dimension-function-scheme}。
\end{lemma}

\begin{proof}
可用 $x$ 的一个仿射邻域替换 $X$。故可以并且确实假设存在
$S$ 上的紧化 $X \subset \overline{X}$。再用
$\overline{X}$ 替换 $X$，可设 $X$ 在 $S$ 上为真。还可用含 $x$ 的
$X$ 的连通分支替换 $X$，从而假设 $X$ 连通。其次，回忆
$\delta_X$ 和函数
$x \mapsto \delta_S(f(x)) + \text{trdeg}_{\kappa(f(x))}(\kappa(x))$
都是 $X$ 上的维数函数；见《态射》，引理
\ref{morphisms-lemma-dimension-function-propagates}
（并使用由引理 \ref{duality-lemma-dimension-function-scheme} 得到的
$S$ 是普遍链状的事实）。由《拓扑》，引理
\ref{topology-lemma-dimension-function-unique}，二者之差局部常值；
由于 $X$ 连通，它是常值的。因此只需在 $X$ 的任一点证明等式。
由《性质》，引理 \ref{properties-lemma-locally-Noetherian-closed-point}，
概形 $X$ 有闭点 $x$。由于 $X \to S$ 是真态射，$x$ 的像 $s$
在 $S$ 中闭。于是应用引理 \ref{duality-lemma-good-dualizing-normalized} 即得结论。
\end{proof}

\begin{lemma}
\label{duality-lemma-shriek}
在情形 \ref{duality-situation-shriek} 中，设 $f : X \to Y$ 是
$\textit{FTS}_S$ 的态射。设 $x \in X$ 的像为 $y \in Y$。则
$$
H^i(f^!\mathcal{O}_Y)_x \not = 0
\Rightarrow - \dim_x(X_y) \leq i.
$$
\end{lemma}

\begin{proof}
因为该陈述在 $X$ 上是局部的，可设 $X$ 和 $Y$ 是仿射概形。
写成 $X = \Spec(A)$ 和 $Y = \Spec(R)$。
$f^!\mathcal{O}_Y$ 对应于《对偶复形》，第
\ref{dualizing-section-relative-dualizing-complexes-Noetherian} 节的
相对对偶复形 $\omega_{A/R}^\bullet$；这由注
\ref{duality-remark-local-calculation-shriek} 得到。因此该引理由《对偶复形》，引理
\ref{dualizing-lemma-relative-dualizing-trivial-vanishing} 得到。
\end{proof}

\begin{lemma}
\label{duality-lemma-flat-shriek}
在情形 \ref{duality-situation-shriek} 中，设 $f : X \to Y$ 是
$\textit{FTS}_S$ 的态射。设 $x \in X$ 的像为 $y \in Y$。
若 $f$ 平坦，则
$$
H^i(f^!\mathcal{O}_Y)_x \not = 0
\Rightarrow - \dim_x(X_y) \leq i \leq 0.
$$
事实上，若 $f$ 的所有纤维维数都 $\leq d$，则作为
$D(X, f^{-1}\mathcal{O}_Y)$ 的对象，$f^!\mathcal{O}_Y$
的 Tor 振幅位于 $[-d, 0]$。
\end{lemma}

\begin{proof}
完全像引理 \ref{duality-lemma-shriek} 的证明中那样论证，结论由《对偶复形》，引理
\ref{dualizing-lemma-relative-dualizing-flat-vanishing} 得到。
\end{proof}

\begin{lemma}
\label{duality-lemma-shriek-over-CM}
在情形 \ref{duality-situation-shriek} 中，设 $f : X \to Y$ 是
$\textit{FTS}_S$ 的态射。设 $x \in X$ 的像为 $y \in Y$。假设
\begin{enumerate}
\item $\mathcal{O}_{Y, y}$ 是 Cohen-Macaulay 的；并且
\item 对含 $x$ 的 $X$ 的任意不可约分支的泛点 $\xi$，有
$\text{trdeg}_{\kappa(f(\xi))}(\kappa(\xi)) \leq r$。
\end{enumerate}
则
$$
H^i(f^!\mathcal{O}_Y)_x \not = 0
\Rightarrow - r \leq i
$$
，而且作为 $\mathcal{O}_{X, x}$-模，茎
$H^{-r}(f^!\mathcal{O}_Y)_x$ 满足 $(S_2)$。
\end{lemma}

\begin{proof}
用 $x$ 的一个开邻域替换 $X$ 后，可设 $X$ 的每个不可约分支都经过 $x$。
于是完全像引理 \ref{duality-lemma-shriek} 的证明中那样论证，结论由《对偶复形》，引理
\ref{dualizing-lemma-relative-dualizing-CM-vanishing} 得到。
\end{proof}

\begin{lemma}
\label{duality-lemma-flat-quasi-finite-shriek}
在情形 \ref{duality-situation-shriek} 中，设 $f : X \to Y$ 是
$\textit{FTS}_S$ 的态射。若 $f$ 平坦且拟有限，则
$$
f^!\mathcal{O}_Y = \omega_{X/Y}[0]
$$
，其中 $\omega_{X/Y}$ 是某个在 $Y$ 上平坦的凝聚
$\mathcal{O}_X$-模。
\end{lemma}

\begin{proof}
这是引理 \ref{duality-lemma-flat-shriek} 以及如下事实的推论：由引理
\ref{duality-lemma-shriek-coherent}，$f^!\mathcal{O}_Y$ 的上同调层是凝聚的。
\end{proof}

\begin{lemma}
\label{duality-lemma-CM-shriek}
在情形 \ref{duality-situation-shriek} 中，设 $f : X \to Y$ 是
$\textit{FTS}_S$ 的态射。若 $f$ 是 Cohen-Macaulay 态射
（《态射进阶》，定义 \ref{more-morphisms-definition-CM}），则
$$
f^!\mathcal{O}_Y = \omega_{X/Y}[d]
$$
，其中 $\omega_{X/Y}$ 是某个在 $Y$ 上平坦的凝聚
$\mathcal{O}_X$-模，而 $d$ 是 $X$ 上给出 $X$ 在 $Y$ 上相对维数的
局部常值函数。
\end{lemma}

\begin{proof}
由《态射》，引理
\ref{morphisms-lemma-flat-finite-presentation-CM-fibres-relative-dimension}，
相对维数 $d$ 定义良好且局部常值。由引理
\ref{duality-lemma-shriek-coherent}，$f^!\mathcal{O}_Y$ 的上同调层是凝聚的。
若能说明 $f^!\mathcal{O}_Y$ 的其他上同调层均为零，则由引理
\ref{duality-lemma-flat-shriek} 得到 $\omega_{X/Y}$ 的平坦性。

\medskip\noindent
这个问题在 $X$ 上是局部的，故可设 $X$ 和 $Y$ 是仿射的，且态射的相对维数为
$d$。若 $d = 0$，结论直接由引理
\ref{duality-lemma-flat-quasi-finite-shriek} 得到。若 $d > 0$，则可设存在分解
$$
X \xrightarrow{g} \mathbf{A}^d_Y \xrightarrow{p} Y
$$
，其中 $g$ 拟有限且平坦；见《态射进阶》，引理
\ref{more-morphisms-lemma-flat-finite-presentation-characterize-CM}。
于是 $f^! = g^! \circ p^!$。由引理
\ref{duality-lemma-shriek-affine-line}，有
$p^!\mathcal{O}_Y \cong \mathcal{O}_{\mathbf{A}^d_Y}[-d]$。
最后用 $d = 0$ 的情形即得结论。
\end{proof}

\begin{remark}
\label{duality-remark-the-same-is-true}
设 $S$ 是配备对偶复形 $\omega_S^\bullet$ 的 Noetherian 概形。
在此情形下，引理 \ref{duality-lemma-shriek}、
\ref{duality-lemma-flat-shriek}、
\ref{duality-lemma-flat-quasi-finite-shriek} 和 \ref{duality-lemma-CM-shriek}
对 $S$ 上有限型概形之间的任意态射 $f : X \to Y$ 都成立，
但要用 $f_{new}^!$ 替换 $f^!$。这是清楚的，因为每种情形的证明都立即归约到
仿射情形，而此时由引理 \ref{duality-lemma-duality-bootstrap}，有
$f^! = f_{new}^!$。
\end{remark}




\section{对偶模}
\label{duality-section-dualizing-module}


\noindent
本节是《对偶复形》，第 \ref{dualizing-section-dualizing-module} 节的延续。

\medskip\noindent
设 $X$ 是 Noetherian 概形，$\omega_X^\bullet$ 是对偶复形。
设 $n \in \mathbf{Z}$ 是使 $H^n(\omega_X^\bullet)$ 非零的最小整数。
换言之，$-n$ 是与 $\omega_X^\bullet$ 相伴的维数函数的最大值
（引理 \ref{duality-lemma-dimension-function-scheme}）。有时
$H^n(\omega_X^\bullet)$ 称为 $X$ 的{\it 对偶模}或{\it 对偶层}，
并常记作 $\omega_X$。下文说“设 $\omega_X$ 为一个对偶模”时，
意指上述情形。

\medskip\noindent
在 Noetherian 概形 $X$ 上使用对偶模 $\omega_X$ 时必须谨慎：
\begin{enumerate}
\item 从 $X$ 转到 $X$ 的开集 $U$ 时，整数 $n$ 可能改变；
此时 $\omega_X|_U = \omega_U$ 不成立；
\item 对偶复形不唯一；对偶模仅在张量一个可逆模的意义下唯一。
\end{enumerate}
第二个问题往往无关紧要，因为将研究基变换 $S$ 上的有限型概形
$X$，其中该基变换配备固定的对偶复形 $\omega_S^\bullet$，而
$\omega_X^\bullet$ 将取为相对于 $\omega_S^\bullet$ 规范化的对偶复形。
若 $X$ 等维，则第一个问题不会发生；更精确地说，只要与
$\omega_X^\bullet$ 相伴的维数函数（引理
\ref{duality-lemma-dimension-function-scheme}）将 $X$ 的每个泛点映到同一个整数即可。

\begin{example}
\label{duality-example-proper-over-local}
设 $S = \Spec(A)$，其中 $(A, \mathfrak m, \kappa)$ 是 Noetherian 局部环，
而 $\omega_S^\bullet$ 对应于规范化对偶复形
$\omega_A^\bullet$。若 $f : X \to S$ 在 $S$ 上为真，且
$\omega_X^\bullet = f^!\omega_S^\bullet$，则凝聚层
$$
\omega_X = H^{-\dim(X)}(\omega_X^\bullet)
$$
是对偶模，并常称为 $X$ 的对偶模（这里默认已指定 $S$ 和
$\omega_S^\bullet$）。将会看到它具有良好性质。
\end{example}

\begin{example}
\label{duality-example-equidimensional-over-field}
设 $X$ 是域 $k$ 上的等维有限型概形。通常取
$\omega_X^\bullet$ 为相对于 $k[0]$ 规范化的对偶复形，并把
$$
\omega_X = H^{-\dim(X)}(\omega_X^\bullet)
$$
称为 $X$ 的对偶模。若 $X$ 在 $k$ 上分离，则
$\omega_X^\bullet = f^!\mathcal{O}_{\Spec(k)}$，其中
$f : X \to \Spec(k)$ 是结构态射；这由引理
\ref{duality-lemma-duality-bootstrap} 得到。若 $X$ 在 $k$ 上为真，
则这是例 \ref{duality-example-proper-over-local} 的特殊情形。
\end{example}

\begin{lemma}
\label{duality-lemma-dualizing-module}
设 $X$ 是连通 Noetherian 概形，$\omega_X$ 是 $X$ 上的对偶模。
$\omega_X$ 的支集是相对于任意维数函数具有最大维数的不可约分支之并，
而且 $\omega_X$ 是满足性质 $(S_2)$ 的凝聚 $\mathcal{O}_X$-模。
\end{lemma}

\begin{proof}
按上面讨论的约定，存在对偶复形 $\omega_X^\bullet$，使得
$\omega_X$ 是其最左边的非零上同调层。由于 $X$ 连通，任意两个维数函数
相差一个常数（《拓扑》，引理
\ref{topology-lemma-dimension-function-unique}）。因此可使用与
$\omega_X^\bullet$ 相伴的维数函数（引理
\ref{duality-lemma-dimension-function-scheme}）。有了这些说明，该引理由
《对偶复形》，引理 \ref{dualizing-lemma-depth-dualizing-module}
和各定义得到（特别是《概形的上同调》，定义
\ref{coherent-definition-depth}）。
\end{proof}

\begin{lemma}
\label{duality-lemma-vanishing-good-dualizing}
设 $X/A$、$\omega_X^\bullet$、$\omega_X$ 如例
\ref{duality-example-proper-over-local} 所述。则
\begin{enumerate}
\item $H^i(\omega_X^\bullet) \not = 0 \Rightarrow
i \in \{-\dim(X), \ldots, 0\}$；
\item $H^i(\omega_X^\bullet)$ 的支集维数至多为 $-i$；
\item $\text{Supp}(\omega_X)$ 是维数为 $\dim(X)$ 的分支之并；并且
\item $\omega_X$ 满足性质 $(S_2)$。
\end{enumerate}
\end{lemma}

\begin{proof}
令 $\delta_X$ 和 $\delta_S$ 为与 $\omega_X^\bullet$ 和
$\omega_S^\bullet$ 相伴的维数函数，如引理
\ref{duality-lemma-good-dualizing-dimension-function} 所述。由于 $X$ 在 $A$ 上为真，
$X$ 的每个闭子概形都含有一个映到闭点 $s \in S$ 的闭点 $x$，并且
$\delta_X(x) = \delta_S(s) = 0$。因此对任意点 $\xi \in X$，有
$\delta_X(\xi) = \dim(\overline{\{\xi\}})$。所以可通过考察
$\Spec(\mathcal{O}_{X, x})$ 上的情形来检验引理的每个陈述；此时结论由
《对偶复形》，引理 \ref{dualizing-lemma-sitting-in-degrees} 和
\ref{dualizing-lemma-depth-dualizing-module} 得到。略去若干细节。
最后两个陈述也可由引理 \ref{duality-lemma-dualizing-module} 推出。
\end{proof}

\begin{lemma}
\label{duality-lemma-dualizing-module-proper-over-A}
设带对偶模 $\omega_X$ 的 $X/A$ 如例
\ref{duality-example-proper-over-local} 所述。令 $d = \dim(X_s)$ 为闭纤维的维数。
若 $\dim(X) = d + \dim(A)$，则对偶模 $\omega_X$ 表示凝聚
$\mathcal{O}_X$-模范畴上的函子
$$
\mathcal{F} \longmapsto \Hom_A(H^d(X, \mathcal{F}), \omega_A)
$$
。
\end{lemma}

\begin{proof}
有
\begin{align*}
\Hom_X(\mathcal{F}, \omega_X)
& =
\Ext^{-\dim(X)}_X(\mathcal{F}, \omega_X^\bullet) \\
& =
\Hom_X(\mathcal{F}[\dim(X)], \omega_X^\bullet) \\
& =
\Hom_X(\mathcal{F}[\dim(X)], f^!(\omega_A^\bullet)) \\
& =
\Hom_S(Rf_*\mathcal{F}[\dim(X)], \omega_A^\bullet) \\
& =
\Hom_A(H^d(X, \mathcal{F}), \omega_A)
\end{align*}
第一个等式成立，是因为当 $i < -\dim(X)$ 时
$H^i(\omega_X^\bullet) = 0$；见引理
\ref{duality-lemma-vanishing-good-dualizing} 和《导出范畴》，引理
\ref{derived-lemma-negative-exts}。第二个等式由 Ext 群的定义得到。
第三个等式来自对 $\omega_X^\bullet$ 的选择。第四个等式成立，
是因为 $f^!$ 是引理 \ref{duality-lemma-twisted-inverse-image} 对 $f$ 给出的右伴随；
见第 \ref{duality-section-duality} 节。最后一个等式成立，是因为当 $i > d$ 时
$R^if_*\mathcal{F}$ 为零（《概形的上同调》，引理
\ref{coherent-lemma-higher-direct-images-zero-above-dimension-fibre}），
并且当 $j < -\dim(A)$ 时 $H^j(\omega_A^\bullet)$ 为零。
\end{proof}








\section{Cohen-Macaulay 概形}
\label{duality-section-CM}

\noindent
本节是《对偶复形》，第 \ref{dualizing-section-CM} 节的延续。
对 Cohen-Macaulay 概形，对偶性具有特别简单的形式。

\begin{lemma}
\label{duality-lemma-dualizing-module-CM-scheme}
设 $X$ 是带对偶复形 $\omega_X^\bullet$ 的局部 Noetherian 概形。
\begin{enumerate}
\item $X$ 是 Cohen-Macaulay 的 $\Leftrightarrow$
$\omega_X^\bullet$ 局部仅有一个非零上同调层；
\item $\mathcal{O}_{X, x}$ 是 Cohen-Macaulay 的 $\Leftrightarrow$
$\omega_{X, x}^\bullet$ 仅有一个非零上同调；
\item $U = \{x \in X \mid \mathcal{O}_{X, x}\text{ 为 Cohen--Macaulay}\}$
是开集且为 Cohen-Macaulay 概形。
\end{enumerate}
若 $X$ 连通且为 Cohen-Macaulay 概形，则存在整数 $n$ 和凝聚
Cohen-Macaulay $\mathcal{O}_X$-模 $\omega_X$，使得
$\omega_X^\bullet = \omega_X[-n]$。
\end{lemma}

\begin{proof}
由定义和《对偶复形》，引理
\ref{dualizing-lemma-dualizing-localize}，对每个 $x \in X$，
复形 $\omega_{X, x}^\bullet$ 是 $\mathcal{O}_{X, x}$ 上的对偶复形。
由《对偶复形》，引理 \ref{dualizing-lemma-apply-CM}，可知 (2) 成立。

\medskip\noindent
为证明 (3)，假设 $\mathcal{O}_{X, x}$ 是 Cohen-Macaulay 的。
令 $n_x$ 为使 $H^{n_{x}}(\omega_{X, x}^\bullet)$ 非零的唯一整数。
对 $x$ 的仿射邻域 $V \subset X$，有
$\omega_X^\bullet|_V$ 属于 $D^b_{\textit{Coh}}(\mathcal{O}_V)$，
故只有有限多个凝聚模 $H^i(\omega_X^\bullet)|_V$ 非零。
因此缩小 $V$ 后，可设只有 $H^{n_x}$ 非零；见《模》，引理
\ref{modules-lemma-finite-type-stalk-zero}。于是可知对每个
$v \in V$，$\mathcal{O}_{X, v}$ 都是 Cohen-Macaulay 的。
这同时证明 $U$ 是开集和 Cohen-Macaulay 概形。

\medskip\noindent
证明 (1)。蕴含 $\Leftarrow$ 由 (2) 得到。蕴含 $\Rightarrow$
由上一段的讨论得到：那里已经说明，若 $\mathcal{O}_{X, x}$
是 Cohen-Macaulay 的，则在 $x$ 的一个邻域中，复形
$\omega_X^\bullet$ 只有一个非零上同调层。

\medskip\noindent
假设 $X$ 连通且为 Cohen-Macaulay 概形。上面说明映射
$x \mapsto n_x$ 局部常值。由于 $X$ 连通，它是常值的，设其值为 $n$。
令 $\omega_X = H^n(\omega_X^\bullet)$。由《对偶复形》，引理
\ref{dualizing-lemma-apply-CM}（以及《概形的上同调》，定义
\ref{coherent-definition-Cohen-Macaulay}），$\omega_X$ 是
Cohen-Macaulay 的，故引理成立。
\end{proof}

\begin{lemma}
\label{duality-lemma-has-dualizing-module-CM-scheme}
设 $X$ 是局部 Noetherian 概形。若存在凝聚层 $\omega_X$，使得
$\omega_X[0]$ 是 $X$ 上的对偶复形，则 $X$ 是 Cohen-Macaulay 概形。
\end{lemma}

\begin{proof}
这由《对偶复形》，引理
\ref{dualizing-lemma-has-dualizing-module-CM} 和我们的定义立即得到。
\end{proof}

\begin{lemma}
\label{duality-lemma-affine-flat-Noetherian-CM}
在情形 \ref{duality-situation-shriek} 中，设 $f : X \to Y$ 是
$\textit{FTS}_S$ 的态射，$x \in X$。若 $f$ 平坦，则下列条件等价：
\begin{enumerate}
\item $f$ 在 $x$ 处是 Cohen-Macaulay 的；
\item $f^!\mathcal{O}_Y$ 在 $x$ 的一个邻域中只有一个非零上同调层。
\end{enumerate}
\end{lemma}

\begin{proof}
引理的一个方向由引理 \ref{duality-lemma-CM-shriek} 得到。为证明逆向，
可设 $f^!\mathcal{O}_Y$ 只有一个非零上同调层。令 $y = f(x)$。
设 $\xi_1, \ldots, \xi_n \in X_y$ 是纤维 $X_y$ 中特化到 $x$ 的泛点。
令 $d_1, \ldots, d_n$ 为 $X_y$ 相应不可约分支的维数。
由《态射进阶》，引理
\ref{more-morphisms-lemma-flat-finite-presentation-CM-open}，
态射 $f : X \to Y$ 在 $\eta_i$ 处是 Cohen-Macaulay 的。
所以由引理 \ref{duality-lemma-CM-shriek}，有 $d_1 = \ldots = d_n$。
若 $d$ 表示其共同值，则 $d = \dim_x(X_y)$。缩小 $X$ 后，可设所有纤维
维数至多为 $d$（《态射》，引理
\ref{morphisms-lemma-openness-bounded-dimension-fibres}）。于是由引理
\ref{duality-lemma-flat-shriek}，唯一的非零上同调层
$\omega = H^{-d}(f^!\mathcal{O}_Y)$ 在 $Y$ 上平坦。因此若
$h : X_y \to X$ 表示典范态射，则由《概形的导出范畴》，引理
\ref{perfect-lemma-tor-independence-and-tor-amplitude}，有
$Lh^*(f^!\mathcal{O}_Y) = Lh^*(\omega[d]) = (h^*\omega)[d]$。
所以由引理 \ref{duality-lemma-relative-dualizing-fibres}，
$h^*\omega[d]$ 是 $X_y$ 的对偶复形。进而由引理
\ref{duality-lemma-dualizing-module-CM-scheme}，$X_y$ 是 Cohen-Macaulay 的。
这如所需地证明 $f$ 在 $x$ 处是 Cohen-Macaulay 的。
\end{proof}

\begin{remark}
\label{duality-remark-CM-morphism-compare-dualizing}
在情形 \ref{duality-situation-shriek} 中，设 $f : X \to Y$ 是
$\textit{FTS}_S$ 的态射。假设 $f$ 是相对维数为 $d$ 的
Cohen-Macaulay 态射。令 $\omega_{X/Y} = H^{-d}(f^!\mathcal{O}_Y)$
为 $f^!\mathcal{O}_Y$ 的唯一非零上同调层；见引理
\ref{duality-lemma-CM-shriek}。则存在典范同构
$$
f^!K = Lf^*K \otimes_{\mathcal{O}_X}^\mathbf{L} \omega_{X/Y}[d]
$$
，其中 $K \in D^+_\QCoh(\mathcal{O}_Y)$；见引理
\ref{duality-lemma-perfect-comparison-shriek}。特别地，若 $S$
有对偶复形 $\omega_S^\bullet$，
$\omega_Y^\bullet = (Y \to S)^!\omega_S^\bullet$，且
$\omega_X^\bullet = (X \to S)^!\omega_S^\bullet$，则
$$
\omega_X^\bullet =
Lf^*\omega_Y^\bullet \otimes_{\mathcal{O}_X}^\mathbf{L} \omega_{X/Y}[d]
$$
。进一步，若 $X$ 和 $Y$ 连通且为 Cohen-Macaulay 概形，
而 $\omega_Y$ 和 $\omega_X$ 分别表示 $\omega_Y^\bullet$ 和
$\omega_X^\bullet$ 的唯一非零上同调层，则
$$
\omega_X = f^*\omega_Y \otimes_{\mathcal{O}_X} \omega_{X/Y}.
$$
类似结论对 $S$ 上任意有限型概形 $X$ 和 $Y$ 也成立
（即不必在 $S$ 上分离），其中对偶复形像第
\ref{duality-section-glue} 节那样相对于 $\omega_S^\bullet$ 规范化。
\end{remark}





\section{Gorenstein 概形}
\label{duality-section-gorenstein}

\noindent
本节是《对偶复形》，第 \ref{dualizing-section-gorenstein} 节的延续。

\begin{definition}
\label{duality-definition-gorenstein}
设 $X$ 是概形。若 $X$ 局部 Noetherian，且对所有 $x \in X$，
$\mathcal{O}_{X, x}$ 都是 Gorenstein 环，则称 $X$ 为{\it Gorenstein} 概形。
\end{definition}

\noindent
这个定义是合理的，因为一个 Noetherian 环称为 Gorenstein 环，当且仅当
其所有局部环都是 Gorenstein 环；见《对偶复形》，定义
\ref{dualizing-definition-gorenstein}。

\begin{lemma}
\label{duality-lemma-gorenstein-CM}
Gorenstein 概形是 Cohen-Macaulay 的。
\end{lemma}

\begin{proof}
在仿射局部考察，结论由代数中的相应结果得到，即《对偶复形》，引理
\ref{dualizing-lemma-gorenstein-CM}。
\end{proof}

\begin{lemma}
\label{duality-lemma-regular-gorenstein}
正则概形是 Gorenstein 的。
\end{lemma}

\begin{proof}
在仿射局部考察，结论由代数中的相应结果得到，即《对偶复形》，引理
\ref{dualizing-lemma-regular-gorenstein}。
\end{proof}

\begin{lemma}
\label{duality-lemma-gorenstein}
设 $X$ 是局部 Noetherian 概形。
\begin{enumerate}
\item 若 $X$ 有对偶复形 $\omega_X^\bullet$，则
\begin{enumerate}
\item $X$ 是 Gorenstein 的 $\Leftrightarrow$
$\omega_X^\bullet$ 是 $D(\mathcal{O}_X)$ 中的可逆对象；
\item $\mathcal{O}_{X, x}$ 是 Gorenstein 的 $\Leftrightarrow$
$\omega_{X, x}^\bullet$ 是 $D(\mathcal{O}_{X, x})$ 中的可逆对象；
\item $U = \{x \in X \mid \mathcal{O}_{X, x}\text{ is Gorenstein}\}$
是 Gorenstein 开子概形。
\end{enumerate}
\item 若 $X$ 是 Gorenstein 的，则 $X$ 有对偶复形，当且仅当
$\mathcal{O}_X[0]$ 是对偶复形。
\end{enumerate}
\end{lemma}

\begin{proof}
在仿射局部考察，结论由代数中的相应结果得到，即《对偶复形》，引理
\ref{dualizing-lemma-gorenstein}。
\end{proof}

\begin{lemma}
\label{duality-lemma-gorenstein-lci}
若 $f : Y \to X$ 是局部完全交态射，且 $X$ 是 Gorenstein 概形，
则 $Y$ 是 Gorenstein 概形。
\end{lemma}

\begin{proof}
由《态射进阶》，引理 \ref{more-morphisms-lemma-affine-lci}，
只需证明关于环同态的相应陈述。这就是《对偶复形》，引理
\ref{dualizing-lemma-gorenstein-lci}。
\end{proof}

\begin{lemma}
\label{duality-lemma-gorenstein-local-syntomic}
性质 $\mathcal{P}(S) =$``$S$ is Gorenstein'' 在 syntomic 拓扑中是局部的。
\end{lemma}

\begin{proof}
设 $\{S_i \to S\}$ 是 syntomic 覆盖。概形 $S$ 局部 Noetherian，
当且仅当每个 $S_i$ 都是 Noetherian 的；见《下降》，引理
\ref{descent-lemma-Noetherian-local-fppf}。因此现在可设 $S$ 和 $S_i$
都局部 Noetherian。若 $S$ 是 Gorenstein 的，则由引理
\ref{duality-lemma-gorenstein-lci}，每个 $S_i$ 都是 Gorenstein 的。
反过来，若每个 $S_i$ 都是 Gorenstein 的，则对每个点
$s \in S$，可选取 $i$ 和映到 $s$ 的 $t \in S_i$。
于是 $\mathcal{O}_{S, s} \to \mathcal{O}_{S_i, t}$ 是平坦局部环同态，
且 $\mathcal{O}_{S_i, t}$ 是 Gorenstein 的。因此由《对偶复形》，引理
\ref{dualizing-lemma-flat-under-gorenstein}，$\mathcal{O}_{S, s}$
是 Gorenstein 的。
\end{proof}






\section{Gorenstein 态射}
\label{duality-section-gorenstein-morphisms}

\noindent
本节是一个系列的一部分。关于正规态射、正则态射和 Cohen-Macaulay 态射的
相应各节见《态射进阶》，第
\ref{more-morphisms-section-normal}、
\ref{more-morphisms-section-regular} 和
\ref{more-morphisms-section-CM} 节。

\medskip\noindent
下面的引理说明，定义几何 Gorenstein 概形没有意义，因为它们与
Gorenstein 概形相同。

\begin{lemma}
\label{duality-lemma-gorenstein-base-change}
设 $X$ 是域 $k$ 上的局部 Noetherian 概形，$k'/k$ 是有限生成域扩张。
设 $x \in X$ 是点，$x' \in X_{k'}$ 是位于 $x$ 上方的点。则
$$
\mathcal{O}_{X, x}\text{ is Gorenstein}
\Leftrightarrow
\mathcal{O}_{X_{k'}, x'}\text{ is Gorenstein}
$$
。若 $X$ 在 $k$ 上局部有限型，则对任意域扩张 $k'/k$ 也有同样结论。
\end{lemma}

\begin{proof}
在两种情形下，环同态
$\mathcal{O}_{X, x} \to \mathcal{O}_{X_{k'}, x'}$ 都是 Noetherian
局部环之间的忠实平坦局部同态。因此若
$\mathcal{O}_{X_{k'}, x'}$ 是 Gorenstein 的，则由《对偶复形》，引理
\ref{dualizing-lemma-flat-under-gorenstein}，$\mathcal{O}_{X, x}$
也是 Gorenstein 的。反向上升也使用《对偶复形》，引理
\ref{dualizing-lemma-flat-under-gorenstein}。所以必须证明
$$
\mathcal{O}_{X_{k'}, x'}/\mathfrak m_x \mathcal{O}_{X_{k'}, x'} =
\kappa(x) \otimes_k k'
$$
是 Gorenstein 的。注意在第一种情形下 $k \to k'$ 有限生成，
在第二种情形下 $k \to \kappa(x)$ 有限生成。因此结论来自
Gorenstein 满足性质 (A)；见《对偶复形》，引理
\ref{dualizing-lemma-formal-fibres-gorenstein}。
\end{proof}

\noindent
上述引理保证下面确实是 Gorenstein 态射的正确定义。

\begin{definition}
\label{duality-definition-gorenstein-morphism}
设 $f : X \to Y$ 是概形态射。假设所有纤维 $X_y$ 都是局部
Noetherian 概形。
\begin{enumerate}
\item 设 $x \in X$ 且 $y = f(x)$。若 $f$ 在 $x$ 处平坦，
而概形 $X_y$ 在 $x$ 处的局部环是 Gorenstein 的，则称 $f$
{\it 在 $x$ 处为 Gorenstein}；
\item 若 $f$ 在 $X$ 的每一点都是 Gorenstein 的，则称 $f$
为{\it Gorenstein 态射}。
\end{enumerate}
\end{definition}

\noindent
下面给出一个等价表述。

\begin{lemma}
\label{duality-lemma-gorenstein-morphism}
设 $f : X \to Y$ 是概形态射。假设 $f$ 的所有纤维都局部 Noetherian。
则下列条件等价：
\begin{enumerate}
\item $f$ 是 Gorenstein 的；并且
\item $f$ 平坦，且其纤维都是 Gorenstein 概形。
\end{enumerate}
\end{lemma}

\begin{proof}
这直接由定义得到。
\end{proof}

\begin{lemma}
\label{duality-lemma-gorenstein-CM-morphism}
Gorenstein 态射是 Cohen-Macaulay 的。
\end{lemma}

\begin{proof}
结论由引理 \ref{duality-lemma-gorenstein-CM} 和定义得到。
\end{proof}

\begin{lemma}
\label{duality-lemma-lci-gorenstein}
syntomic 态射是 Gorenstein 的。等价地，平坦局部完全交态射是
Gorenstein 的。
\end{lemma}

\begin{proof}
回忆 syntomic 态射平坦，且其纤维是域上的局部完全交；见《态射》，引理
\ref{morphisms-lemma-syntomic-flat-fibres}。由引理
\ref{duality-lemma-gorenstein-lci}，域上的局部完全交是 Gorenstein 概形，
故结论成立。性质“syntomic”和“平坦且为局部完全交态射”由《态射进阶》，引理
\ref{more-morphisms-lemma-flat-lci} 可知等价。
\end{proof}

\begin{lemma}
\label{duality-lemma-composition-gorenstein}
设 $f : X \to Y$ 和 $g : Y \to Z$ 是态射。假设 $f$、$g$ 和
$g \circ f$ 的纤维 $X_y$、$Y_z$、$X_z$ 都局部 Noetherian。
\begin{enumerate}
\item 若 $f$ 在 $x$ 处为 Gorenstein，且 $g$ 在 $f(x)$ 处为
Gorenstein，则 $g \circ f$ 在 $x$ 处为 Gorenstein；
\item 若 $f$ 和 $g$ 都是 Gorenstein 的，则 $g \circ f$ 是 Gorenstein 的；
\item 若 $g \circ f$ 在 $x$ 处为 Gorenstein，且 $f$ 在 $x$ 处平坦，
则 $f$ 在 $x$ 处为 Gorenstein，而 $g$ 在 $f(x)$ 处为 Gorenstein；
\item 若 $g \circ f$ 是 Gorenstein 的且 $f$ 平坦，则 $f$ 是
Gorenstein 的，而 $g$ 在 $f$ 的像中的每一点都是 Gorenstein 的。
\end{enumerate}
\end{lemma}

\begin{proof}
翻译为代数后，结论由《对偶复形》，引理
\ref{dualizing-lemma-flat-under-gorenstein} 得到。
\end{proof}

\begin{lemma}
\label{duality-lemma-flat-morphism-from-gorenstein-scheme}
\begin{slogan}
平坦纤维化的总空间具有 Gorenstein 性，蕴含基与纤维也具有 Gorenstein 性
\end{slogan}
设 $f : X \to Y$ 是局部 Noetherian 概形之间的平坦态射。
若 $X$ 是 Gorenstein 的，则 $f$ 是 Gorenstein 的，并且对所有
$x \in X$，$\mathcal{O}_{Y, f(x)}$ 都是 Gorenstein 的。
\end{lemma}

\begin{proof}
翻译为代数后，结论由《对偶复形》，引理
\ref{dualizing-lemma-flat-under-gorenstein} 得到。
\end{proof}

\begin{lemma}
\label{duality-lemma-base-change-gorenstein}
设 $f : X \to Y$ 是概形态射。假设所有纤维 $X_y$ 都是局部
Noetherian 概形。设 $Y' \to Y$ 局部有限型，并设
$f' : X' \to Y'$ 是 $f$ 的基变换。设 $x' \in X'$ 的像为 $x \in X$。
\begin{enumerate}
\item 若 $f$ 在 $x$ 处为 Gorenstein，则
$f' : X' \to Y'$ 在 $x'$ 处为 Gorenstein；
\item 若 $f$ 在 $x$ 处平坦，且 $f'$ 在 $x'$ 处为 Gorenstein，
则 $f$ 在 $x$ 处为 Gorenstein；
\item 若 $Y' \to Y$ 在 $f'(x')$ 处平坦，且 $f'$ 在 $x'$ 处为
Gorenstein，则 $f$ 在 $x$ 处为 Gorenstein。
\end{enumerate}
\end{lemma}

\begin{proof}
注意，对 $Y' \to Y$ 的假设蕴含：若 $y' \in Y'$ 映到 $y \in Y$，
则域扩张 $\kappa(y')/\kappa(y)$ 有限生成。因此所有纤维
$X'_{y'} = (X_y)_{\kappa(y')}$ 也都局部 Noetherian；见《簇》，引理
\ref{varieties-lemma-locally-Noetherian-base-change}。所以该引理有意义。
令 $y' = f'(x')$ 和 $y = f(x)$。于是得到局部环的交换图
$$
\xymatrix{
\mathcal{O}_{X', x'} & \mathcal{O}_{X, x} \ar[l] \\
\mathcal{O}_{Y', y'} \ar[u] & \mathcal{O}_{Y, y} \ar[l] \ar[u]
}
$$
，其中左上角是右上角与左下角在右下角上的张量积的一个局部化。

\medskip\noindent
假设 $f$ 在 $x$ 处为 Gorenstein。
$\mathcal{O}_{Y, y} \to \mathcal{O}_{X, x}$ 的平坦性蕴含
$\mathcal{O}_{Y', y'} \to \mathcal{O}_{X', x'}$ 的平坦性；见《代数》，引理
\ref{algebra-lemma-base-change-flat-up-down}。
$\mathcal{O}_{X, x}/\mathfrak m_y\mathcal{O}_{X, x}$ 是 Gorenstein 的
这一事实蕴含
$\mathcal{O}_{X', x'}/\mathfrak m_{y'}\mathcal{O}_{X', x'}$
是 Gorenstein 的；见引理 \ref{duality-lemma-gorenstein-base-change}。
所以 $f'$ 在 $x'$ 处为 Gorenstein。

\medskip\noindent
假设 $f$ 在 $x$ 处平坦，且 $f'$ 在 $x'$ 处为 Gorenstein。
$\mathcal{O}_{X', x'}/\mathfrak m_{y'}\mathcal{O}_{X', x'}$
是 Gorenstein 的这一事实蕴含
$\mathcal{O}_{X, x}/\mathfrak m_y\mathcal{O}_{X, x}$ 是 Gorenstein 的；
见引理 \ref{duality-lemma-gorenstein-base-change}。所以 $f$ 在 $x$ 处为 Gorenstein。

\medskip\noindent
假设 $Y' \to Y$ 在 $y'$ 处平坦，且 $f'$ 在 $x'$ 处为 Gorenstein。
$\mathcal{O}_{Y', y'} \to \mathcal{O}_{X', x'}$ 和
$\mathcal{O}_{Y, y} \to \mathcal{O}_{Y', y'}$ 的平坦性蕴含
$\mathcal{O}_{Y, y} \to \mathcal{O}_{X, x}$ 的平坦性；见《代数》，引理
\ref{algebra-lemma-base-change-flat-up-down}。
$\mathcal{O}_{X', x'}/\mathfrak m_{y'}\mathcal{O}_{X', x'}$
是 Gorenstein 的这一事实蕴含
$\mathcal{O}_{X, x}/\mathfrak m_y\mathcal{O}_{X, x}$ 是 Gorenstein 的；
见引理 \ref{duality-lemma-gorenstein-base-change}。所以 $f$ 在 $x$ 处为 Gorenstein。
\end{proof}

\begin{lemma}
\label{duality-lemma-flat-lft-base-change-gorenstein}
设 $f : X \to Y$ 是平坦且局部有限型的概形态射。则集合
$\{x \in X \mid f\text{ is Gorenstein at }x\}$
的形成与任意基变换交换。
\end{lemma}

\begin{proof}
该假设蕴含 $f$ 的任意纤维都在某个域上局部有限型，因而局部
Noetherian；任意基变换也同样如此。因此该陈述有意义。考察纤维后，
问题归结为：设 $X$ 是域 $k$ 上局部有限型的概形，设 $K/k$ 是域扩张，
并设点 $x_K \in X_K$ 的像为 $x \in X$。要证明：
$\mathcal{O}_{X_K, x_K}$ 是 Gorenstein 的，当且仅当
$\mathcal{O}_{X, x}$ 是 Gorenstein 的。

\medskip\noindent
可用少量代数如下解决该问题。选取包含 $x$ 的仿射开集
$\Spec(A) \subset X$。设 $x$ 对应于 $\mathfrak p \subset A$。
令 $A_K = A \otimes_k K$，则包含 $x_K$ 的仿射开集为
$\Spec(A_K) \subset X_K$。设 $x_K$ 对应于 $\mathfrak p_K \subset A_K$。
令 $\omega_A^\bullet$ 为 $A$ 的一个对偶复形。由《对偶复形》，引理
\ref{dualizing-lemma-base-change-dualizing-over-field}，
$\omega_A^\bullet \otimes_A A_K$ 是 $A_K$ 的一个对偶复形。
现在结论成立，因为
$A_\mathfrak p \to (A_K)_{\mathfrak p_K}$ 是 Noetherian 环的平坦局部同态，
所以 $(\omega_A^\bullet)_\mathfrak p$ 是 $D(A_\mathfrak p)$ 中的可逆对象，
当且仅当
$(\omega_A^\bullet)_\mathfrak p \otimes_{A_\mathfrak p} (A_K)_{\mathfrak p_K}$
是 $D((A_K)_{\mathfrak p_K})$ 中的可逆对象。
略去了一些细节；提示：考察上同调模。
\end{proof}

\begin{lemma}
\label{duality-lemma-affine-flat-Noetherian-gorenstein}
在情形 \ref{duality-situation-shriek} 中，设 $f : X \to Y$ 是
$\textit{FTS}_S$ 的态射，并设 $x \in X$。若 $f$ 平坦，则下列条件等价：
\begin{enumerate}
\item $f$ 在 $x$ 处为 Gorenstein；
\item $f^!\mathcal{O}_Y$ 在 $x$ 的某个邻域内与一个可逆对象同构。
\end{enumerate}
特别地，$f$ 为 Gorenstein 的点组成 $X$ 中的开集。
\end{lemma}

\begin{proof}
令 $\omega^\bullet = f^!\mathcal{O}_Y$。由引理
\ref{duality-lemma-relative-dualizing-fibres}，这是一个上同调层相干的有界复形，
且它在纤维 $X_y$ 上的导出限制 $Lh^*\omega^\bullet$ 是 $X_y$ 上的对偶复形。
以 $i : x \to X_y$ 表示该点的嵌入。则下列条件等价：
\begin{enumerate}
\item $f$ 在 $x$ 处为 Gorenstein；
\item $\mathcal{O}_{X_y, x}$ 是 Gorenstein 的；
\item $Lh^*\omega^\bullet$ 在 $x$ 的某个邻域内可逆；
\item $Li^* Lh^* \omega^\bullet$ 恰有一个非零上同调，且它在
$\kappa(x)$ 上的维数为 $1$；
\item $L(h \circ i)^* \omega^\bullet$ 恰有一个非零上同调，且它在
$\kappa(x)$ 上的维数为 $1$；
\item $\omega^\bullet$ 在 $x$ 的某个邻域内可逆。
\end{enumerate}
(1) 与 (2) 的等价性由定义得到（因为 $f$ 平坦）。
(2) 与 (3) 的等价性来自引理 \ref{duality-lemma-gorenstein}。
(3) 与 (4) 的等价性来自《代数进阶》，引理
\ref{more-algebra-lemma-lift-bounded-pseudo-coherent-to-perfect}。
(4) 与 (5) 等价，是因为 $Li^* Lh^* = L(h \circ i)^*$。
(5) 与 (6) 的等价性来自《代数进阶》，引理
\ref{more-algebra-lemma-lift-bounded-pseudo-coherent-to-perfect}。
由此引理得证。
\end{proof}

\begin{lemma}
\label{duality-lemma-flat-finite-presentation-characterize-gorenstein}
设 $f : X \to S$ 是平坦且局部有限表示的概形态射。设 $x \in X$ 的像为
$s \in S$，并令 $d = \dim_x(X_s)$。则下列条件等价：
\begin{enumerate}
\item $f$ 在 $x$ 处为 Gorenstein；
\item 存在 $x$ 的开邻域 $U \subset X$ 以及 $S$ 上的局部拟有限态射
$U \to \mathbf{A}^d_S$，它在 $x$ 处为 Gorenstein；
\item 存在 $x$ 的开邻域 $U \subset X$ 以及 $S$ 上的局部拟有限
Gorenstein 态射 $U \to \mathbf{A}^d_S$；
\item 对 $x$ 的任意开邻域 $U \subset X$ 及任意 $S$-态射
$g : U \to \mathbf{A}^d_S$，都有：
$g$ 在 $x$ 处拟有限 $\Rightarrow$ $g$ 在 $x$ 处为 Gorenstein。
\end{enumerate}
特别地，$f$ 为 Gorenstein 的点组成 $X$ 中的开集。
\end{lemma}

\begin{proof}
选取满足 $x \in U$ 的仿射开集 $U = \Spec(A) \subset X$，以及满足
$f(U) \subset V$ 的仿射开集 $V = \Spec(R) \subset S$。则 $R \to A$
是有限表示的平坦环映射。设 $\mathfrak p \subset A$ 是对应于 $x$ 的素理想。
用一个主局部化替换 $A$ 后，可以假设存在拟有限映射
$R[x_1, \ldots, x_d]  \to A$；见《代数》，引理
\ref{algebra-lemma-quasi-finite-over-polynomial-algebra}。
因此至少存在一对 $(U, g)$，其中 $U \subset X$ 是 $x$ 的开邻域，且
$g : U \to \mathbf{A}^d_S$ 是局部\footnote{若 $S$ 拟分离，则 $g$
是拟有限的。}拟有限态射。

\medskip\noindent
由此，该引理转化为下面的代数问题（略去转化过程）。给定有限表示的平坦映射
$R \to A$、素理想 $\mathfrak p \subset A$，以及在 $\mathfrak p$ 处拟有限的
$\varphi : R[x_1, \ldots, x_d] \to A$，则下列条件等价：
\begin{enumerate}
\item[(a)] $\Spec(\varphi)$ 在 $\mathfrak p$ 处为 Gorenstein；并且
\item[(b)] $\Spec(A) \to \Spec(R)$ 在 $\mathfrak p$ 处为 Gorenstein；
\item[(c)] $\Spec(A) \to \Spec(R)$ 在 $\mathfrak p$ 的某个开邻域内为
Gorenstein。
\end{enumerate}
在每种情形下，$R[x_1, \ldots, x_n] \to A$ 在 $\mathfrak p$ 处平坦，
故由平坦性的开性（《代数》，定理
\ref{algebra-theorem-openness-flatness}），可以假设
$R[x_1, \ldots, x_n] \to A$ 平坦（用适当的主局部化替换 $A$）。
由《代数》，引理 \ref{algebra-lemma-flat-finite-presentation-limit-flat}，
存在 $R_0 \subset R$ 和 $R_0[x_1, \ldots, x_n] \to A_0$，使得
$R_0$ 在 $\mathbf{Z}$ 上有限型，$R_0 \to A_0$ 有限型，并且
$R_0[x_1, \ldots, x_n] \to A_0$ 平坦。注意，由引理
\ref{duality-lemma-base-change-gorenstein}，有限型平坦态射为 Gorenstein 的点集
与基变换交换。如此便归结到 $R$ 为 Noetherian 的情形。

\medskip\noindent
于是可以假设 $X$ 和 $S$ 都是仿射的，且 $f$ 有如下分解：
$$
X \xrightarrow{g} \mathbf{A}^n_S \xrightarrow{p} S
$$
，其中 $g$ 平坦且拟有限，而 $S$ 是 Noetherian 的。于是 $X$ 和
$\mathbf{A}^n_S$ 在 $S$ 上分离，并且
$$
f^!\mathcal{O}_S = g^!p^!\mathcal{O}_S = g^!\mathcal{O}_{\mathbf{A}^n_S}[n]
$$
；这是上叹号函子的已知性质（引理
\ref{duality-lemma-upper-shriek-composition} 和
\ref{duality-lemma-shriek-affine-line}）。因此由引理
\ref{duality-lemma-affine-flat-Noetherian-gorenstein}，(a)、(b)、(c) 等价。
\end{proof}

\begin{lemma}
\label{duality-lemma-gorenstein-local-source-and-target}
性质 $\mathcal{P}(f)=$“$f$ 的纤维局部 Noetherian，且 $f$ 为
Gorenstein”在目标的 fppf 拓扑下是局部的，在源的 syntomic 拓扑下也是局部的。
\end{lemma}

\begin{proof}
有
$\mathcal{P}(f) =
\mathcal{P}_1(f) \wedge \mathcal{P}_2(f)$，
其中 $\mathcal{P}_1(f)=$“$f$ 平坦”，而
$\mathcal{P}_2(f)=$“$f$ 的纤维局部 Noetherian 且为 Gorenstein”。
已知 $\mathcal{P}_1$ 在源和目标的 fppf 拓扑下都是局部的；见《下降》，引理
\ref{descent-lemma-descending-property-flat} 和
\ref{descent-lemma-flat-fpqc-local-source}。所以只须处理 $\mathcal{P}_2$。

\medskip\noindent
设 $f : X \to Y$ 是概形态射，且
$\{\varphi_i : Y_i \to Y\}_{i \in I}$ 是 $Y$ 的一个 fppf 覆盖。
以 $f_i : X_i \to Y_i$ 表示 $f$ 沿 $\varphi_i$ 的基变换。
取 $i \in I$ 和一点 $y_i \in Y_i$，令 $y = \varphi_i(y_i)$。注意
$$
X_{i, y_i} = \Spec(\kappa(y_i)) \times_{\Spec(\kappa(y))} X_y.
$$
并且 $\kappa(y_i)/\kappa(y)$ 是有限生成域扩张。因此若 $X_y$ 局部
Noetherian，则 $X_{i, y_i}$ 也局部 Noetherian；见《簇》，引理
\ref{varieties-lemma-locally-Noetherian-base-change}。若进一步 $X_y$ 为
Gorenstein，则 $X_{i, y_i}$ 也为 Gorenstein；见引理
\ref{duality-lemma-gorenstein-base-change}。因此 $\mathcal{P}_2$ 在目标上是 fppf 局部的。

\medskip\noindent
设 $\{X_i \to X\}$ 是 $X$ 的一个 syntomic 覆盖，并设 $y \in Y$。
此时 $\{X_{i, y} \to X_y\}$ 是该纤维的一个 syntomic 覆盖。因此，由引理
\ref{duality-lemma-gorenstein-local-syntomic}，$\mathcal{P}_2$ 在源的 syntomic
拓扑下是局部的。
\end{proof}




\section{对偶复形进阶}
\label{duality-section-more-dualizing}

\noindent
下面是一些不太适合放在别处的引理。

\begin{lemma}
\label{duality-lemma-descent-ascent}
设 $f : X \to Y$ 是局部 Noetherian 概形的态射。假设下列条件之一成立：
\begin{enumerate}
\item $f$ 是满的 syntomic 态射；或
\item $f$ 是满的平坦局部完全交态射；或
\item $f$ 是有限型的满 Gorenstein 态射。
\end{enumerate}
则 $K \in D_\QCoh(\mathcal{O}_Y)$ 是 $Y$ 上的对偶复形，当且仅当
$Lf^*K$ 是 $X$ 上的对偶复形。
\end{lemma}

\begin{proof}
取仿射开集并使用《概形的导出范畴》，引理
\ref{perfect-lemma-affine-compare-bounded}，该结论转化为《对偶复形》，引理
\ref{dualizing-lemma-descent-ascent}。
\end{proof}







\section{域上固有概形的对偶性}
\label{duality-section-duality-proper-over-field}

\noindent
本节阐明上面关于对偶复形的一般材料对域上固有概形的对偶性所产生的结论。

\begin{lemma}
\label{duality-lemma-duality-proper-over-field}
设 $X$ 是域 $k$ 上的固有概形。存在一个具有下列性质的对偶复形
$\omega_X^\bullet$：
\begin{enumerate}
\item 仅当 $i \in [-\dim(X), 0]$ 时，$H^i(\omega_X^\bullet)$ 才可能非零；
\item $\omega_X = H^{-\dim(X)}(\omega_X^\bullet)$ 是相干
$(S_2)$-模，其支撑由维数为 $\dim(X)$ 的不可约分支组成；
\item $H^i(\omega_X^\bullet)$ 的支撑维数至多为 $-i$；
\item 对闭点 $x \in X$，模
$H^i(\omega_{X, x}^\bullet) \oplus \ldots \oplus H^0(\omega_{X, x}^\bullet)$
非零，当且仅当 $\text{depth}(\mathcal{O}_{X, x}) \leq -i$；
\item 对 $K \in D_\QCoh(\mathcal{O}_X)$，存在与平移和可区别三角相容的函子性同构
\footnote{由 Yoneda 引理，该性质在 $D_\QCoh(\mathcal{O}_X)$ 中将
$\omega_X^\bullet$ 刻画到唯一同构。由于 $\omega_X^\bullet$ 属于
$D^b_{\textit{Coh}}(\mathcal{O}_X)$，事实上只须考虑
$K \in D^b_{\textit{Coh}}(\mathcal{O}_X)$。}
$$
\Ext^i_X(K, \omega_X^\bullet) = \Hom_k(H^{-i}(X, K), k)
$$
；
\item 对 $X$ 上的拟相干层 $\mathcal{F}$，存在函子性同构
$\Hom(\mathcal{F}, \omega_X) = \Hom_k(H^{\dim(X)}(X, \mathcal{F}), k)$；并且
\item 若 $X \to \Spec(k)$ 光滑且相对维数为 $d$，则
$\omega_X^\bullet \cong \wedge^d\Omega_{X/k}[d]$ 且
$\omega_X \cong \wedge^d\Omega_{X/k}$。
\end{enumerate}
\end{lemma}

\begin{proof}
以 $f : X \to \Spec(k)$ 表示结构态射。设 $a$ 为该态射的推前的右伴随；
见引理 \ref{duality-lemma-twisted-inverse-image}。考虑相对对偶复形
$$
\omega_X^\bullet = a(\mathcal{O}_{\Spec(k)})
$$
。参见注 \ref{duality-remark-relative-dualizing-complex}。由于 $f$ 固有，按定义有
$f^!(\mathcal{O}_{\Spec(k)}) = a(\mathcal{O}_{\Spec(k)})$；见第
\ref{duality-section-upper-shriek} 节。应用引理 \ref{duality-lemma-shriek-dualizing}，
可知 $\omega_X^\bullet$ 是对偶复形。此外，$\omega_X^\bullet$ 和
$\omega_X$ 分别如例 \ref{duality-example-proper-over-local} 和例
\ref{duality-example-equidimensional-over-field} 中所述。

\medskip\noindent
(1)、(2)、(3) 来自引理 \ref{duality-lemma-vanishing-good-dualizing}。

\medskip\noindent
对闭点 $x \in X$，$\omega_{X, x}^\bullet$ 是
$\mathcal{O}_{X, x}$ 上的规范化对偶复形；见引理
\ref{duality-lemma-good-dualizing-normalized}。于是 (4) 来自《对偶复形》，引理
\ref{dualizing-lemma-depth-in-terms-dualizing-complex}。

\medskip\noindent
(5) 由构造成立，因为 $a$ 是
$Rf_* : D_\QCoh(\mathcal{O}_X) \to D(\mathcal{O}_{\Spec(k)}) = D(k)$
的右伴随，而后者可与 $K \mapsto R\Gamma(X, K)$ 等同。这里还使用了：
$k$-模的导出范畴 $D(k)$ 等同于分次 $k$-向量空间范畴。

\medskip\noindent
当 $\mathcal{F}$ 相干时，(6) 来自引理
\ref{duality-lemma-dualizing-module-proper-over-A}；一般情形则把 (5) 用于
$K = \mathcal{F}[0]$ 和 $i = -\dim(X)$ 并展开定义即可。

\medskip\noindent
(7) 来自引理 \ref{duality-lemma-smooth-proper}。
\end{proof}

\begin{remark}
\label{duality-remark-duality-proper-over-field}
设 $k$、$X$、$\omega_X^\bullet$ 如引理
\ref{duality-lemma-duality-proper-over-field} 中所述。经由 (5) 中的函子性同构，
复形 $\omega_X^\bullet$ 上的恒等映射对应于一个映射
$$
t : H^0(X, \omega_X^\bullet) \longrightarrow k
$$
。对 $D_\QCoh(\mathcal{O}_X)$ 中的任意 $K$，(5) 中
$\Hom(K, \omega_X^\bullet)$ 与 $H^0(X, K)^\vee$ 的等同对应于配对
$$
\Hom_X(K, \omega_X^\bullet) \times H^0(X, K) \longrightarrow k,\quad
(\alpha, \beta) \longmapsto t(\alpha(\beta))
$$
。这来自 (5) 中同构的函子性。类似地，对任意 $i \in \mathbf{Z}$ 得到配对
$$
\Ext^i_X(K, \omega_X^\bullet) \times H^{-i}(X, K) \longrightarrow k,\quad
(\alpha, \beta) \longmapsto t(\alpha(\beta))
$$
。为定义 $\alpha(\beta)$，这里把 $\alpha$ 看作态射
$K[-i] \to \omega_X^\bullet$，把 $\beta$ 看作 $H^0(X, K[-i])$ 的元素。
注意，当 $K$ 任意时，只知道该配对在一侧非退化：该配对分别给出
$\Hom_X(K, \omega_X^\bullet)$ 和\ $\Ext^i_X(K, \omega_X^\bullet)$，
与 $H^0(X, K)$ 的 $k$-线性对偶及\ $H^{-i}(X, K)$ 的线性对偶之间的同构；
反向一般不成立。若 $K$
属于 $D^b_{\textit{Coh}}(\mathcal{O}_X)$，则
$\Hom_X(K, \omega_X^\bullet)$、$\Ext_X(K, \omega_X^\bullet)$、
$H^0(X, K)$ 和 $H^i(X, K)$ 都是有限维 $k$-向量空间（由《概形的导出范畴》，引理
\ref{perfect-lemma-coherent-internal-hom} 和
\ref{perfect-lemma-direct-image-coherent-bdd-below}），
并且这些配对在通常意义下是完美的。
\end{remark}

\begin{remark}
\label{duality-remark-coherent-duality-proper-over-field}
继续注 \ref{duality-remark-duality-proper-over-field} 中的讨论，并沿用记号
$k$、$X$、$\omega_X^\bullet$、$t$。若 $\mathcal{F}$ 是相干
$\mathcal{O}_X$-模，则得到有限维 $k$-向量空间之间的完美配对
$$
\langle -, - \rangle :
\Ext^i_X(\mathcal{F}, \omega_X^\bullet) \times H^{-i}(X,\mathcal{F})
\longrightarrow k,\quad
(\alpha, \beta) \longmapsto t(\alpha(\beta))
$$
。这些配对满足下列（显然的）函子性：若
$\varphi : \mathcal{F} \to \mathcal{G}$ 是相干 $\mathcal{O}_X$-模的同态，则
$$
\langle \alpha \circ \varphi, \beta \rangle =
\langle \alpha, \varphi(\beta) \rangle
$$
对 $\alpha \in \Ext^i_X(\mathcal{G}, \omega_X^\bullet)$ 和
$\beta \in H^{-i}(X, \mathcal{F})$ 成立。换言之，经由这些配对，由
$\varphi$ 诱导的 $k$-线性映射
$\Ext^i_X(\mathcal{G}, \omega_X^\bullet) \to
\Ext^i_X(\mathcal{F}, \omega_X^\bullet)$，是由 $\varphi$ 诱导的
$k$-线性映射 $H^{-i}(X, \mathcal{F}) \to H^{-i}(X, \mathcal{G})$
的 $k$-线性对偶。按这种方式表述后，当 $\varphi$ 是拟相干
$\mathcal{O}_X$-模的同态时，结论仍然成立。
\end{remark}

\begin{lemma}
\label{duality-lemma-duality-proper-over-field-perfect}
设 $k$、$X$、$\omega_X^\bullet$ 如引理
\ref{duality-lemma-duality-proper-over-field} 中所述，并设
$t : H^0(X, \omega_X^\bullet) \to k$ 如注
\ref{duality-remark-duality-proper-over-field} 中所述。设
$E \in D(\mathcal{O}_X)$ 是完美的。则对所有 $i$，配对
$$
H^i(X, \omega_X^\bullet \otimes_{\mathcal{O}_X}^\mathbf{L} E^\vee)
\times
H^{-i}(X, E) \longrightarrow k, \quad
(\xi, \eta) \longmapsto
t((1_{\omega_X^\bullet} \otimes \epsilon)(\xi \cup \eta))
$$
都是完美的。这里 $\cup$ 表示《上同调》，第
\ref{cohomology-section-cup-product} 节的杯积，而
$\epsilon : E^\vee \otimes_{\mathcal{O}_X}^\mathbf{L} E \to \mathcal{O}_X$
如《上同调》，例 \ref{cohomology-example-dual-derived} 中所述。
\end{lemma}

\begin{proof}
用 $E[-i]$ 替换 $E$，便归结到 $i = 0$ 的情形。由《上同调》，引理
\ref{cohomology-lemma-ext-composition-is-cup}，该配对就是注
\ref{duality-remark-duality-proper-over-field} 中讨论的配对，故结论来自该注中的讨论。
\end{proof}

\begin{lemma}
\label{duality-lemma-duality-proper-over-field-CM}
设 $X$ 是域 $k$ 上的固有概形，它是 Cohen-Macaulay 的、等维的，且维数为
$d$。则引理 \ref{duality-lemma-duality-proper-over-field} 中的模 $\omega_X$
具有下列性质：
\begin{enumerate}
\item $\omega_X$ 是 $X$ 上的对偶模（第 \ref{duality-section-dualizing-module} 节）；
\item $\omega_X$ 是支撑为 $X$ 的相干 Cohen-Macaulay 模；
\item 对 $K \in D_\QCoh(X)$，存在与平移和可区别三角相容的函子性同构
$\Ext^i_X(K, \omega_X[d]) = \Hom_k(H^{-i}(X, K), k)$；
\item 对 $X$ 上的拟相干层 $\mathcal{F}$，存在函子性同构
$\Ext^{d - i}(\mathcal{F}, \omega_X) = \Hom_k(H^i(X, \mathcal{F}), k)$。
\end{enumerate}
\end{lemma}

\begin{proof}
由引理 \ref{duality-lemma-duality-proper-over-field}，$\omega_X$ 显然是对偶模
（因为它是对偶复形最左侧的非零上同调层）。有
$\omega_X^\bullet = \omega_X[d]$；并且因为 $X$ 是 Cohen-Macaulay 的，
$\omega_X$ 也是 Cohen-Macaulay 的；见引理
\ref{duality-lemma-dualizing-module-CM-scheme}。其余陈述由此与引理
\ref{duality-lemma-duality-proper-over-field} 中的相应陈述合并得到。
\end{proof}

\begin{remark}
\label{duality-remark-rework-duality-locally-free-CM}
设 $X$ 是域 $k$ 上的固有 Cohen-Macaulay 概形，它等维且维数为 $d$。
设 $\omega_X^\bullet$ 和 $\omega_X$ 如引理
\ref{duality-lemma-duality-proper-over-field} 中所述。由引理
\ref{duality-lemma-duality-proper-over-field-CM}，有
$\omega_X^\bullet = \omega_X[d]$。设
$t : H^d(X, \omega_X) \to k$ 是注
\ref{duality-remark-duality-proper-over-field} 中的映射。设 $\mathcal{E}$ 是有限局部自由
$\mathcal{O}_X$-模，其对偶为 $\mathcal{E}^\vee$。则有完美配对
$$
H^i(X, \omega_X \otimes_{\mathcal{O}_X} \mathcal{E}^\vee)
\times
H^{d - i}(X, \mathcal{E})
\longrightarrow
k,\quad
(\xi, \eta) \longmapsto t(1 \otimes \epsilon)(\xi \cup \eta))
$$
，其中 $\cup$ 是杯积，而
$\epsilon : \mathcal{E}^\vee \otimes_{\mathcal{O}_X} \mathcal{E}
\to \mathcal{O}_X$ 是求值映射。这是引理
\ref{duality-lemma-duality-proper-over-field-perfect} 的特殊情形。
\end{remark}

\noindent
下面对该对偶复形作一个相容性检验。

\begin{lemma}
\label{duality-lemma-sanity-check-duality}
设 $X$ 是域 $k$ 上的固有概形，且 $\omega_X^\bullet$、$\omega_X$
如引理 \ref{duality-lemma-duality-proper-over-field} 中所述。
\begin{enumerate}
\item 若对某个域 $k'$，$X \to \Spec(k)$ 分解为
$X \to \Spec(k') \to \Spec(k)$，则对态射 $X \to \Spec(k')$，
$\omega_X^\bullet$ 和 $\omega_X$ 仍如引理
\ref{duality-lemma-duality-proper-over-field} 中所述；
\item 若 $K/k$ 是域扩张，则 $\omega_X^\bullet$ 和 $\omega_X$
到基变换 $X_K$ 上的拉回，对态射 $X_K \to \Spec(K)$ 而言仍如引理
\ref{duality-lemma-duality-proper-over-field} 中所述。
\end{enumerate}
\end{lemma}

\begin{proof}
以 $f : X \to \Spec(k)$ 表示结构态射，并以
$f' : X \to \Spec(k')$ 表示给定的分解。在引理
\ref{duality-lemma-duality-proper-over-field} 的证明中，取了
$\omega_X^\bullet = a(\mathcal{O}_{\Spec(k)})$，其中 $a$ 是引理
\ref{duality-lemma-twisted-inverse-image} 对 $f$ 给出的右伴随。因此须证明
$a(\mathcal{O}_{\Spec(k)}) \cong a'(\mathcal{O}_{\Spec(k)})$，
其中 $a'$ 是引理 \ref{duality-lemma-twisted-inverse-image} 对 $f'$ 给出的右伴随。
由于 $k' \subset H^0(X, \mathcal{O}_X)$，可知 $k'/k$ 是有限扩张
（《概形的上同调》，引理
\ref{coherent-lemma-proper-over-affine-cohomology-finite}）。
由伴随的唯一性，有 $a = a' \circ b$，其中 $b$ 是引理
\ref{duality-lemma-twisted-inverse-image} 对 $g : \Spec(k') \to \Spec(k)$ 给出的右伴随。
换言之，有 $f^! = (f')^! \circ g^!$。所以只须证明作为 $k'$-模有
$\Hom_k(k', k) \cong k'$；见例
\ref{duality-example-affine-twisted-inverse-image}。这是因为两者都是相同维数的
$k'$-向量空间（即维数为 $1$）。

\medskip\noindent
证明 (2)。结论成立，因为由引理 \ref{duality-lemma-more-base-change}，$a$
满足基变换。参见注 \ref{duality-remark-relative-dualizing-complex} 中的讨论。
\end{proof}








\section{相对对偶复形}
\label{duality-section-relative-dualizing-complexes}

\noindent
对有限表示的固有平坦态射，已有刚性相对对偶复形；见注
\ref{duality-remark-relative-dualizing-complex} 和引理 \ref{duality-lemma-van-den-bergh}。
对 Noetherian 概形之间的分离有限型态射 $f : X \to Y$，可以考虑
$f^!\mathcal{O}_Y$。本节为概形（不必局部 Noetherian）之间的平坦且局部
有限表示态射（不必拟分离或拟紧）定义相对对偶复形。我们将证明这类复形存在、
在唯一同构意义下唯一，并与上述情形一致。阅读本节前，请先阅读《对偶复形》，
第 \ref{dualizing-section-relative-dualizing-complexes} 节。

\begin{definition}
\label{duality-definition-relative-dualizing-complex}
设 $X \to S$ 是平坦且局部有限表示的概形态射。任取开集
$W \subset X \times_S X$，使得对角态射
$\Delta_{X/S} : X \to X \times_S X$ 经过闭浸入
$\Delta : X \to W$ 分解。所谓一个{\it 相对对偶复形}，是一个偶对
$(K, \xi)$，其中 $K \in D(\mathcal{O}_X)$ 是对象，且有
$D(\mathcal{O}_W)$ 中的映射
$$
\xi : \Delta_*\mathcal{O}_X \longrightarrow L\text{pr}_1^*K|_W
$$
，满足：
\begin{enumerate}
\item $K$ 是 $S$-完美的（《概形的导出范畴》，定义
\ref{perfect-definition-relatively-perfect}）；并且
\item $\xi$ 给出 $\Delta_*\mathcal{O}_X$ 与
$R\SheafHom_{\mathcal{O}_W}(
\Delta_*\mathcal{O}_X, L\text{pr}_1^*K|_W)$ 之间的同构。
\end{enumerate}
\end{definition}

\noindent
由引理 \ref{duality-lemma-sheaf-with-exact-support-ext}，条件 (2) 等价于存在
$D(\mathcal{O}_X)$ 中的同构
$$
\mathcal{O}_X \longrightarrow
R\SheafHom(\mathcal{O}_X, L\text{pr}_1^*K|_W)
$$
，其沿 $\Delta$ 的推前等于 $\xi$。由于
$R\SheafHom(\mathcal{O}_X, L\text{pr}_1^*K|_W)$ 与开集 $W$ 的选择无关，
偶对 $(K, \xi)$ 所成的范畴也与此选择无关。若 $X \to S$ 分离，则可取
$W = X \times_S X$。下面的引理将把许多论证归结到环的情形。

\begin{lemma}
\label{duality-lemma-relative-dualizing-complex-algebra}
设 $X \to S$ 是平坦且局部有限表示的概形态射，并设 $(K, \xi)$
是一个相对对偶复形。则对任意交换图
$$
\xymatrix{
\Spec(A) \ar[d] \ar[r] & X \ar[d] \\
\Spec(R) \ar[r] & S
}
$$
，若水平箭头都是开浸入，则 $K$ 在 $\Spec(A)$ 上的限制，经由《概形的导出范畴》，
引理 \ref{perfect-lemma-affine-compare-bounded}，对应于《对偶复形》，定义
\ref{dualizing-definition-relative-dualizing-complex} 意义下 $R \to A$
的一个相对对偶复形。
\end{lemma}

\begin{proof}
由于 $R\SheafHom$ 的形成与限制到开集交换，可以直接假设
$X = \Spec(A)$ 且 $S = \Spec(R)$。注意，$D(\mathcal{O}_X)$ 中的相对完美对象
是伪凝聚的，因而属于 $D_\QCoh(\mathcal{O}_X)$（《概形的导出范畴》，引理
\ref{perfect-lemma-pseudo-coherent}）。所以该陈述有意义。
由《概形的导出范畴》，引理
\ref{perfect-lemma-quasi-coherence-pushforward}、
\ref{perfect-lemma-quasi-coherence-pullback}、
\ref{perfect-lemma-quasi-coherence-internal-hom}，取 $\Delta_*$、
$L\text{pr}_1^*$、$R\SheafHom$ 与代数侧的相应运算相容。对于最后一项，注意
$L\text{pr}_1^*K$ 是 $S$-完美的（因而下有界），而
$\Delta_*\mathcal{O}_X$ 是 $D(\mathcal{O}_W)$ 中的伪凝聚对象；翻译为代数，
这表示 $A$ 作为 $A \otimes_R A$-模是伪凝聚的。这由《代数进阶》，引理
\ref{more-algebra-lemma-more-relative-pseudo-coherent-is-moot}
应用于 $R \to A \otimes_R A \to A$ 得到。因此恰好恢复《对偶复形》，定义
\ref{dualizing-definition-relative-dualizing-complex} 中的条件。
\end{proof}

\begin{lemma}
\label{duality-lemma-relative-dualizing-RHom}
设 $X \to S$ 是平坦且局部有限表示的概形态射，并设 $(K, \xi)$
是一个相对对偶复形。则
$\mathcal{O}_X \to R\SheafHom_{\mathcal{O}_X}(K, K)$
是同构。
\end{lemma}

\begin{proof}
在仿射局部考察，并使用引理 \ref{duality-lemma-relative-dualizing-complex-algebra}，
便归结到《对偶复形》，引理
\ref{dualizing-lemma-relative-dualizing-RHom} 的代数情形。
\end{proof}

\begin{lemma}
\label{duality-lemma-uniqueness-relative-dualizing}
设 $X \to S$ 是平坦且局部有限表示的概形态射。若 $(K, \xi)$ 和
$(L, \eta)$ 是 $X/S$ 上的两个相对对偶复形，则存在唯一的同构
$K \to L$，它把 $\xi$ 送到 $\eta$。
\end{lemma}

\begin{proof}
设 $U \subset X$ 是映入 $S$ 中某个仿射开集的仿射开集。由引理
\ref{duality-lemma-relative-dualizing-complex-algebra} 及《对偶复形》，引理
\ref{dualizing-lemma-uniqueness-relative-dualizing}，存在同构
$K|_U \to L|_U$。读者可以在概形情形中复用该引理的论证来证明本情形。
这里改用粘合论证。

\medskip\noindent
假设已有同构 $\alpha : K \to L$。则对某个可逆截面
$u \in H^0(W, \Delta_*\mathcal{O}_X) = H^0(X, \mathcal{O}_X)$，有
$\alpha(\xi) = u \eta$（因为按假设，$\eta$ 和 $\alpha(\xi)$
都是某个可逆 $\Delta_*\mathcal{O}_X$-模的生成元）。因此以
$u^{-1}\alpha$ 替换 $\alpha$ 后，有 $\alpha(\xi) = \eta$。
由引理 \ref{duality-lemma-relative-dualizing-RHom}，$K$ 的自同构群是
$H^0(X, \mathcal{O}_X^*)$，所以这样的 $\alpha$ 至多有一个。

\medskip\noindent
令 $\mathcal{B}$ 为 $X$ 中映入 $S$ 的某个仿射开集的全体仿射开集。
由前两段的讨论，对每个 $U \in \mathcal{B}$，都有唯一同构
$\alpha_U : K|_U \to L|_U$ 把 $\xi$ 送到 $\eta$。由引理
\ref{duality-lemma-relative-dualizing-RHom}，对 $X$ 的任意开集 $U$ 及
$i < 0$，有 $\text{Ext}^i(K|_U, K|_U) = 0$。把《上同调》，引理
\ref{cohomology-lemma-vanishing-and-glueing} 应用于
$\text{id} : X \to X$，得到唯一态射 $\alpha : K \to L$，
它对所有 $U \in \mathcal{B}$ 都与 $\alpha_U$ 相符。由于局部如此，
$\alpha$ 把 $\xi$ 送到 $\eta$。
\end{proof}

\begin{lemma}
\label{duality-lemma-existence-relative-dualizing}
设 $X \to S$ 是平坦且局部有限表示的概形态射。则存在相对对偶复形
$(K, \xi)$。
\end{lemma}

\begin{proof}
令 $\mathcal{B}$ 为 $X$ 中映入 $S$ 的某个仿射开集的全体仿射开集。
对每个 $U$，都有 $U$ 在 $S$ 上的相对对偶复形 $(K_U, \xi_U)$。
具体地，选取仿射开集 $V \subset S$，使得 $U \to X \to S$ 经过 $V$ 分解。
写成 $U = \Spec(A)$ 和 $V = \Spec(R)$。由《对偶复形》，引理
\ref{dualizing-lemma-base-change-relative-dualizing}，存在 $R \to A$
的相对对偶复形 $K_A \in D(A)$。反向应用引理
\ref{duality-lemma-relative-dualizing-complex-algebra} 的证明，这决定了一个
$V$-完美对象 $K_U \in D(\mathcal{O}_U)$，以及
$D(\mathcal{O}_{U \times_V U})$ 中的映射
$$
\xi : \Delta_*\mathcal{O}_U \to L\text{pr}_1^*K_U
$$
。由于 $V$-完美等同于 $S$-完美，且
$U \times_V U = U \times_S U$，可知 $(K_U, \xi_U)$ 满足要求。

\medskip\noindent
若 $U' \subset U \subset X$ 且 $U', U \in \mathcal{B}$，则由引理
\ref{duality-lemma-uniqueness-relative-dualizing}，有唯一同构
$\rho_{U'}^U : K_U|_{U'} \to K_{U'}$ 属于 $D(\mathcal{O}_{U'})$，并把
$\xi_U|_{U' \times_S U'}$ 送到 $\xi_{U'}$（显然，相对对偶复形到开集的限制
仍是相对对偶复形）。唯一性保证，对 $\mathcal{B}$ 中的
$U'' \subset U' \subset U$，有
$\rho^U_{U''} = \rho^V_{U''} \circ \rho ^U_{U'}|_{U''}$。
把引理 \ref{duality-lemma-relative-dualizing-RHom} 应用于 $U/S$ 和 $K_U$，可知对
$U \in \mathcal{B}$ 及 $i < 0$，有 $\text{Ext}^i(K_U, K_U) = 0$。
因此 BBD 粘合引理（《上同调》，定理
\ref{cohomology-theorem-glueing-bbd-general}）说明存在唯一的解，即对象
$K \in D(\mathcal{O}_X)$ 和同构 $\rho_U : K|_U \to K_U$，使得对所有
$U' \subset U$、$U, U' \in \mathcal{B}$，都有
$\rho^U_{U'} \circ \rho_U|_{U'} = \rho_{U'}$。

\medskip\noindent
为完成证明，还须构造 $D(\mathcal{O}_W)$ 中的映射
$$
\xi : \Delta_*\mathcal{O}_X \longrightarrow L\text{pr}_1^*K|_W
$$
，它诱导从 $\Delta_*\mathcal{O}_X$ 到
$R\SheafHom_{\mathcal{O}_W}(\Delta_*\mathcal{O}_X, L\text{pr}_1^*K|_W)$
的同构。由于可以改变 $W$，取
$W = \bigcup_{U \in \mathcal{B}} U \times_S U$。利用 $\rho_U$ 得到同构
$$
R\SheafHom_{\mathcal{O}_W}(
\Delta_*\mathcal{O}_X, L\text{pr}_1^*K|_W)|_{U \times_S U}
\xrightarrow{\rho_U}
R\SheafHom_{\mathcal{O}_{U \times_S U}}(
\Delta_*\mathcal{O}_U, L\text{pr}_1^*K_U)
$$
。由于这些开集 $U \times_S U$ 覆盖 $W$，可知
$R\SheafHom_{\mathcal{O}_W}(\Delta_*\mathcal{O}_X, L\text{pr}_1^*K|_W)$
的上同调层除次数 $0$ 外均为零。此外，得到同构
$$
H^0\left(U \times_S U, R\SheafHom_{\mathcal{O}_W}(\Delta_*\mathcal{O}_X,
L\text{pr}_1^*K|_W)\right)
\xrightarrow{\rho_U}
H^0\left((R\SheafHom_{\mathcal{O}_{U \times_S U}}(
\Delta_*\mathcal{O}_U, L\text{pr}_1^*K_U)\right)
$$
。令左端的元素 $\tau_U$ 在此映射下映到 $\xi_U$。当
$U' \subset U \subset X$ 且开集 $U', U \in \mathcal{B}$ 时，
$\rho^U_{U'}$、$\xi_U$、$\xi_{U'}$、$\rho_U$、$\rho_{U'}$ 之间的相容性
蕴含 $\tau_U|_{U' \times_S U'} = \tau_{U'}$。因此得到第 $0$ 上同调层
$H^0(R\SheafHom_{\mathcal{O}_W}(\Delta_*\mathcal{O}_X, L\text{pr}_1^*K|_W))$
的整体截面 $\tau$。由于
$R\SheafHom_{\mathcal{O}_W}(\Delta_*\mathcal{O}_X, L\text{pr}_1^*K|_W)$
的其余上同调层均为零，该整体截面 $\tau$ 决定所需的态射 $\xi$。
由于 $\xi$ 在 $U \times_S U$ 上的限制给出 $\xi_U$，可知它满足定义
\ref{duality-definition-relative-dualizing-complex} 的最后一个条件。
\end{proof}

\begin{lemma}
\label{duality-lemma-base-change-relative-dualizing}
考虑概形的笛卡尔方块
$$
\xymatrix{
X' \ar[d]_{f'} \ar[r]_{g'} & X \ar[d]^f \\
S' \ar[r]^g & S
}
$$
。假设 $X \to S$ 平坦且局部有限表示。设 $(K, \xi)$ 是 $f$
的一个相对对偶复形。令 $K' = L(g')^*K$，并令 $\xi'$ 为 $\xi$
的导出基变换（见证明）。则 $(K', \xi')$ 是 $f'$ 的相对对偶复形。
\end{lemma}

\begin{proof}
考虑笛卡尔方块
$$
\xymatrix{
X' \ar[d]_{\Delta_{X'/S'}} \ar[r] & X \ar[d]^{\Delta_{X/S}} \\
X' \times_{S'} X' \ar[r]^{g' \times g'} & X \times_S X
}
$$
。选取开集 $W \subset X \times_S X$，使 $\Delta_{X/S}$ 经过闭浸入
$\Delta : X \to W$ 分解。选取开集 $W' \subset X' \times_{S'} X'$，
使 $\Delta_{X'/S'}$ 经过闭浸入 $\Delta' : X \to W'$ 分解，且
$(g' \times g')(W') \subset W$。仍以
$g' \times g' : W' \to W$ 表示所诱导的态射。于是
$$
L(g' \times g')^*\Delta_*\mathcal{O}_X =
\Delta'_*\mathcal{O}_{X'}
\quad\text{且}\quad
L(g' \times g')^*L\text{pr}_1^*K|_W =
L\text{pr}_1^*K'|_{W'}
$$
。第一个等式成立，是因为 $X$ 与 $X' \times_{S'} X'$ 在
$X \times_S X$ 上 Tor 独立（例如见《态射进阶》，引理
\ref{more-morphisms-lemma-case-of-tor-independence}）。第二个等式来自导出拉回的传递性
（《上同调》，引理 \ref{cohomology-lemma-derived-pullback-composition}）。
因此 $\xi' = L(g' \times g')^*\xi$ 可视为映射
$$
\xi' : \Delta'_*\mathcal{O}_{X'} \longrightarrow L\text{pr}_1^*K'|_{W'}
$$
。至此，该引理的证明是直接的。首先，由《概形的导出范畴》，引理
\ref{perfect-lemma-base-change-relatively-perfect}，$K'$ 是 $S'$-完美的。
为检验 $\xi'$ 诱导从 $\Delta'_*\mathcal{O}_{X'}$ 到
$R\SheafHom_{\mathcal{O}_{W'}}(
\Delta'_*\mathcal{O}_{X'}, L\text{pr}_1^*K'|_{W'})$ 的同构，
可以在仿射局部工作。由引理
\ref{duality-lemma-relative-dualizing-complex-algebra}，归结到《对偶复形》，引理
\ref{dualizing-lemma-base-change-relative-dualizing} 中已证的相应代数陈述。
\end{proof}

\begin{lemma}
\label{duality-lemma-flat-proper-relative-dualizing}
设 $S$ 是拟紧拟分离概形，$f : X \to S$ 是有限表示的固有平坦态射。
注 \ref{duality-remark-relative-dualizing-complex} 中的相对对偶复形
$\omega_{X/S}^\bullet$，连同 (\ref{duality-equation-pre-rigid})，是定义
\ref{duality-definition-relative-dualizing-complex} 意义下的相对对偶复形。
\end{lemma}

\begin{proof}
引理 \ref{duality-lemma-properties-relative-dualizing} 已证明
$\omega_{X/S}^\bullet$ 是 $S$-完美的。令 $c$ 为引理
\ref{duality-lemma-twisted-inverse-image} 对对角态射
$\Delta : X \to X \times_S X$ 给出的右伴随。把 $\Delta_*$ 应用于
(\ref{duality-equation-pre-rigid})，得到同构
$$
\Delta_*\mathcal{O}_X \to
\Delta_*(c(L\text{pr}_1^*\omega_{X/S}^\bullet)) =
R\SheafHom_{\mathcal{O}_{X \times_S X}}(
\Delta_*\mathcal{O}_X, L\text{pr}_1^*\omega_{X/S}^\bullet)
$$
。该等式由引理 \ref{duality-lemma-twisted-inverse-image-closed} 和
\ref{duality-lemma-sheaf-with-exact-support-ext} 得到。证明完成。
\end{proof}

\begin{remark}
\label{duality-remark-relative-dualizing-complex-bis}
设 $X \to S$ 是有限表示的平坦固有概形态射。由引理
\ref{duality-lemma-existence-relative-dualizing}，存在定义
\ref{duality-definition-relative-dualizing-complex} 意义下的相对对偶复形
$(\omega_{X/S}^\bullet, \xi)$。考虑任意态射 $g : S' \to S$，其中
$S'$ 拟紧拟分离（例如 $S$ 的一个仿射开集）。由引理
\ref{duality-lemma-base-change-relative-dualizing}，
$(L(g')^*\omega_{X/S}^\bullet, L(g')^*\xi)$ 是基变换
$f' : X' \to S'$ 的相对对偶复形，取定义
\ref{duality-definition-relative-dualizing-complex} 的意义。设
$\omega_{X'/S'}^\bullet$ 是注 \ref{duality-remark-relative-dualizing-complex}
意义下 $X' \to S'$ 的相对对偶复形。结合引理
\ref{duality-lemma-flat-proper-relative-dualizing} 和
\ref{duality-lemma-uniqueness-relative-dualizing}，可知存在唯一同构
$$
\omega_{X'/S'}^\bullet \longrightarrow L(g')^*\omega_{X/S}^\bullet
$$
，它与 (\ref{duality-equation-pre-rigid}) 和 $L(g')^*\xi$ 相容。这些同构与
$S$ 上拟紧拟分离概形之间的态射，以及引理
\ref{duality-lemma-proper-flat-base-change} 的基变换同构相容（如果以后需要这种相容性，
我们将在这里仔细陈述并证明它）。
\end{remark}

\begin{lemma}
\label{duality-lemma-compactifyable-relative-dualizing}
在情形 \ref{duality-situation-shriek} 中，设 $f : X \to Y$ 是
$\textit{FTS}_S$ 的态射。若 $f$ 平坦，则 $f^!\mathcal{O}_Y$
是定义 \ref{duality-definition-relative-dualizing-complex} 意义下 $X$ 在 $Y$
上的某个相对对偶复形的第一分量。
\end{lemma}

\begin{proof}
由引理 \ref{duality-lemma-flat-shriek-relatively-perfect}，
$f^!\mathcal{O}_Y$ 是 $Y$-完美的。由于 $f$ 分离，对角态射
$\Delta : X \to X \times_Y X$ 是闭浸入，并且
$\Delta_*\Delta^!(-) =
R\SheafHom_{\mathcal{O}_{X \times_Y X}}(\mathcal{O}_X, -)$；见引理
\ref{duality-lemma-twisted-inverse-image-closed} 和
\ref{duality-lemma-sheaf-with-exact-support-ext}。因此，为完成证明，只须证明
$\Delta^!(L\text{pr}_1^*f^!(\mathcal{O}_Y)) \cong \mathcal{O}_X$，
其中 $\text{pr}_1 : X \times_Y X \to X$ 是第一投影。有
$$
\mathcal{O}_X = \Delta^! \text{pr}_1^!\mathcal{O}_X =
\Delta^! \text{pr}_1^! L\text{pr}_2^*\mathcal{O}_Y =
\Delta^!(L\text{pr}_1^* f^!\mathcal{O}_Y)
$$
，其中 $\text{pr}_2 : X \times_Y X \to X$ 是第二投影，并使用了引理
\ref{duality-lemma-base-change-shriek-flat} 的基变换同构
$\text{pr}_1^! \circ L\text{pr}_2^* = L\text{pr}_1^* \circ f^!$。
\end{proof}

\begin{lemma}
\label{duality-lemma-relative-dualizing-composition}
设 $f : Y \to X$ 和 $X \to S$ 是有限表示的平坦概形态射。设
$(K, \xi)$ 和 $(M, \eta)$ 分别是 $X \to S$ 和 $Y \to X$
的相对对偶复形。令 $E = M \otimes_{\mathcal{O}_Y}^\mathbf{L} Lf^*K$。
则对适当的 $\zeta$，$(E, \zeta)$ 是 $Y \to S$ 的相对对偶复形。
\end{lemma}

\begin{proof}
使用引理 \ref{duality-lemma-relative-dualizing-complex-algebra} 及该引理的代数版本
（《对偶复形》，引理
\ref{dualizing-lemma-relative-dualizing-composition}），可知在仿射局部，
$E$ 是某个相对对偶复形的第一分量。特别地，$E$ 是 $S$-完美的，
因为这可在仿射局部检验；见《概形的导出范畴》，引理
\ref{perfect-lemma-affine-locally-rel-perfect}。

\medskip\noindent
先在态射 $X \to S$ 和 $Y \to X$ 都分离的情形中证明 $\zeta$ 存在；
此时 $\Delta_{X/S}$、$\Delta_{Y/X}$、$\Delta_{Y/S}$ 都是闭浸入。
考虑下图
$$
\xymatrix{
& & Y \ar@{=}[r] & Y \ar[d]^f \\
Y \ar[r]_{\Delta_{Y/X}} &
Y \times_X Y \ar[d]_m \ar[r]_\delta \ar[ru]_q &
Y \times_S Y \ar[d]^{f \times f} \ar[ru]_p & X\\
& X \ar[r]^{\Delta_{X/S}} & X \times_S X \ar[ru]_r
}
$$
，其中 $p$、$q$、$r$ 是第一投影。由引理
\ref{duality-lemma-sheaf-with-exact-support-internal-home}，有
$$
R\SheafHom_{\mathcal{O}_{Y \times_S Y}}(
\Delta_{Y/S, *}\mathcal{O}_Y, Lp^*E) =
R\delta_*\left(R\SheafHom_{\mathcal{O}_{Y \times_X Y}}(
\Delta_{Y/X, *}\mathcal{O}_Y,
R\SheafHom(\mathcal{O}_{Y \times_X Y}, Lp^*E))\right)
$$
。由引理 \ref{duality-lemma-sheaf-with-exact-support-tensor}，有
$$
R\SheafHom(\mathcal{O}_{Y \times_X Y}, Lp^*E) =
R\SheafHom(\mathcal{O}_{Y \times_X Y}, L(f \times f)^*Lr^*K)
\otimes_{\mathcal{O}_{Y \times_S Y}}^\mathbf{L} Lq^*M
$$
。由引理 \ref{duality-lemma-flat-bc-sheaf-with-exact-support}，有
$$
R\SheafHom(\mathcal{O}_{Y \times_X Y}, L(f \times f)^*Lr^*K) =
Lm^*R\SheafHom(\mathcal{O}_X, Lr^*K)
$$
。最后一个表达式经 $\xi$ 同构于
$Lm^*\mathcal{O}_X = \mathcal{O}_{Y \times_X Y}$，故前一个表达式同构于
$Lq^*M$。于是
$$
R\SheafHom_{\mathcal{O}_{Y \times_S Y}}(
\Delta_{Y/S, *}\mathcal{O}_Y, Lp^*E) =
R\delta_*\left(R\SheafHom_{\mathcal{O}_{Y \times_X Y}}(
\Delta_{Y/X, *}\mathcal{O}_Y, Lq^*M)\right)
$$
。括号内的对象经 $\eta$ 同构于
$\Delta_{Y/X, *}*\mathcal{O}_X$。这就完成了分离情形的证明。

\medskip\noindent
在一般情形中，选取开集 $W \subset X \times_S X$，使
$\Delta_{X/S}$ 经过闭浸入 $\Delta : X \to W$ 分解；并选取开集
$V \subset Y \times_X Y$，使 $\Delta_{Y/X}$ 经过闭浸入
$\Delta' : Y \to V$ 分解。最后，选取开集 $W' \subset Y \times_S Y$，
使它与 $Y \times_X Y$ 的交为 $V$，并映入 $W$。于是考虑图
$$
\xymatrix{
& & Y \ar@{=}[r] & Y \ar[d]^f \\
Y \ar[r]_{\Delta'} &
V \ar[d]_m \ar[r]_\delta \ar[ru]_q &
W' \ar[d]^{f \times f} \ar[ru]_p & X\\
& X \ar[r]^\Delta & W \ar[ru]_r
}
$$
，并使用与前面完全相同的论证。
\end{proof}






\section{lci 态射的基本类}
\label{duality-section-fundamental-class}

\noindent
本节将使用第 \ref{duality-section-examples} 节中的计算。因此，本节的结果也具有
与该节相同类型的非唯一性。

\begin{lemma}
\label{duality-lemma-determinant}
设 $X$ 是局部赋环空间，并设
$$
\mathcal{E}_1 \xrightarrow{\alpha} \mathcal{E}_0 \to \mathcal{F} \to 0
$$
是 $\mathcal{O}_X$-模的短正合序列。假设 $\mathcal{E}_1$ 和
$\mathcal{E}_0$ 分别是秩 $r_1, r_0$ 的局部自由模。则存在典范映射
$$
\wedge^{r_0 - r_1}\mathcal{F}
\longrightarrow
\wedge^{r_1}(\mathcal{E}_1^\vee) \otimes \wedge^{r_0}\mathcal{E}_0
$$
；它在 $x \in X$ 处的茎上是同构，当且仅当 $\mathcal{F}$ 在 $x$
的某个开邻域内局部自由且秩为 $r_0 - r_1$。
\end{lemma}

\begin{proof}
若 $r_1 > r_0$，则按约定 $\wedge^{r_0 - r_1}\mathcal{F} = 0$，
唯一的映射不可能是同构。因此可以假设 $r = r_0 - r_1 \geq 0$。
用下式定义该映射：
$$
s_1 \wedge \ldots \wedge s_r \mapsto
t_1^\vee \wedge \ldots \wedge t_{r_1}^\vee \otimes
\alpha(t_1) \wedge \ldots \wedge \alpha(t_{r_1}) \wedge
\tilde s_1 \wedge \ldots \wedge \tilde s_r
$$
，其中 $t_1, \ldots, t_{r_1}$ 是 $\mathcal{E}_1$ 的一个局部基，
相应地 $t_1^\vee, \ldots, t_{r_1}^\vee$ 是 $\mathcal{E}_1^\vee$ 的对偶基，
而 $s'_i$ 是 $s_i$ 到 $\mathcal{E}_0$ 的一个局部截面的提升。
略去该定义为良定义的证明。

\medskip\noindent
若 $\mathcal{F}$ 局部自由且秩为 $r$，直接可验该映射是同构。反之，
假设该映射在 $x$ 处的茎上是同构。则 $\wedge^r\mathcal{F}_x$ 可逆，
因而 $\mathcal{F}_x$ 可由至多 $r$ 个元素生成。只有当 $\alpha$
模 $\mathfrak m_x$ 的秩为 $r$ 时才可能如此；换言之，$\alpha$
模 $\mathfrak m_x$ 具有最大秩。这蕴含 $\alpha$ 在 $x$ 的某个邻域内
具有最大秩，故 $\mathcal{F}$ 在该邻域内局部自由且秩为 $r$，正合所求。
\end{proof}

\begin{lemma}
\label{duality-lemma-fundamental-class-lci}
设 $Y$ 是 Noetherian 概形，$f : X \to Y$ 是局部完全交态射，且它分解为
浸入 $X \to P$ 后接固有光滑态射 $P \to Y$。设 $r$ 是 $X$
上的局部常值函数，使得
$\omega_{X/Y} = H^{-r}(f^!\mathcal{O}_Y)$ 是 $f^!\mathcal{O}_Y$
唯一的非零上同调层；见引理 \ref{duality-lemma-lci-shriek}。则存在映射
$$
\wedge^r\Omega_{X/Y} \longrightarrow \omega_{X/Y}
$$
；它在点 $x$ 处的茎上是同构，当且仅当 $f$ 在 $x$ 处光滑。
\end{lemma}

\begin{proof}
该假设蕴含 $X$ 在 $Y$ 上可紧化，故 $f^!$ 有定义；见第
\ref{duality-section-upper-shriek} 节。设 $j : W \to P$ 是开子概形，使
$X \to P$ 经过闭浸入 $i : X \to W$ 分解。此外，若
$g : P \to Y$ 为给定态射，则 $f^! = i^! \circ j^! \circ g^!$。
由引理 \ref{duality-lemma-smooth-proper}，有
$g^!\mathcal{O}_Y = \wedge^d\Omega_{P/Y}[d]$，其中 $d$ 是给出 $P$
在 $Y$ 上相对维数的局部常值函数。还有 $j^! = j^*$，并且
$i^!\mathcal{O}_W = \wedge^c\mathcal{N}[-c]$，其中 $c$ 是 $X$ 在 $W$
中的余维（$X$ 上的局部常值函数），而 $\mathcal{N}$ 是 Koszul 正则浸入
$i$ 的法层；见引理 \ref{duality-lemma-regular-immersion}。合并以上各式，得到
$$
f^!\mathcal{O}_Y =
\left(\wedge^c\mathcal{N} \otimes_{\mathcal{O}_X}
\wedge^d\Omega_{P/Y}|_X\right)[d - c]
$$
，这里还使用了引理 \ref{duality-lemma-perfect-comparison-shriek}。
因此作为 $X$ 上的局部常值函数，有 $r = d|_X - c$。态射
$X \to P$ 的余法层是模 $\mathcal{I}/\mathcal{I}^2$，其中
$\mathcal{I} \subset \mathcal{O}_W$ 是 $i$ 的理想层；见《态射》，第
\ref{morphisms-section-conormal-sheaf} 节。考虑《态射》，引理
\ref{morphisms-lemma-differentials-relative-immersion} 的典范正合序列
$$
\mathcal{I}/\mathcal{I}^2 \to
\Omega_{P/Y}|_X \to \Omega_{X/Y} \to 0
$$
。应用引理 \ref{duality-lemma-determinant} 即得到所需映射。

\medskip\noindent
若 $f$ 在 $x$ 处光滑，则由引理 \ref{duality-lemma-determinant} 以及
$\Omega_{X/Y}$ 在 $x$ 处局部自由且秩为 $r$，该映射是同构。反之，
假设所构造的映射在 $x$ 处的茎上是同构。该引理说明，用 $x$
的一个开邻域替换 $X$ 后，$\Omega_{X/Y}$ 自由且秩为 $r$。另一方面，
还可假设 $X = \Spec(A)$、$Y = \Spec(R)$，其中
$A = R[x_1, \ldots, x_n]/(f_1, \ldots, f_m)$，而
$f_1, \ldots, f_m$ 是 Koszul 正则序列（这来自局部完全交态射的定义）。
显然这蕴含 $r = n - m$。于是偏导数矩阵
$\partial f_j/\partial x_i$ 在 $x$ 处剩余域上的秩为 $m$。重排变量后，
可假设 $(\partial f_j/\partial x_i)_{1 \leq i, j \leq m}$ 的行列式
在 $x$ 的某个开邻域内可逆。因此 $R \to A$ 在该点光滑；例如见《代数》，
例 \ref{algebra-example-make-standard-smooth}。
\end{proof}

\begin{lemma}
\label{duality-lemma-fundamental-class-almost-lci}
设 $f : X \to Y$ 是概形态射，并设 $r \geq 0$。假设：
\begin{enumerate}
\item $Y$ 是 Cohen-Macaulay 的（《性质》，定义
\ref{properties-definition-Cohen-Macaulay}）；
\item $f$ 分解为 $X \to P \to Y$，其中第一个态射是浸入，第二个态射光滑且固有；
\item 若 $x \in X$ 且 $\dim(\mathcal{O}_{X, x}) \leq 1$，则 $f$
在 $x$ 处为 Koszul（《态射进阶》，定义
\ref{more-morphisms-definition-lci}）；并且
\item 若 $\xi$ 是 $X$ 的某个不可约分支的泛点，则
$\text{trdeg}_{\kappa(f(\xi))} \kappa(\xi) = r$。
\end{enumerate}
那么，令 $\omega_{X/Y} = H^{-r}(f^!\mathcal{O}_Y)$，存在映射
$$
\wedge^r\Omega_{X/Y} \longrightarrow \omega_{X/Y}
$$
，它在 $f$ 光滑的轨迹上是同构。
\end{lemma}

\begin{proof}
设 $U \subset X$ 是 $f$ 为局部完全交态射的开子概形。由假设 (4)，
$f$ 在所有泛点处的相对维数都是 $r$，故引理
\ref{duality-lemma-fundamental-class-lci} 中的局部常值函数恒等于 $r$，并得到映射
$$
\wedge^r\Omega_{X/Y}|_U = \wedge^r \Omega_{U/Y}
\longrightarrow
\omega_{U/Y} = \omega_{X/Y}|_U
$$
；它在 $f$ 的光滑点处是同构（该轨迹包含于 $U$，因为光滑态射是局部完全交态射）。
由引理 \ref{duality-lemma-shriek-over-CM} 及 $Y$ 为 Cohen-Macaulay 的假设，
模 $\omega_{X/Y}$ 满足 $(S_2)$。由条件 (3)，$U$ 包含所有余维为 $1$
的点；再用《除子》，引理 \ref{divisors-lemma-depth-2-hartog}，得到
$j_*\omega_{U/Y} = \omega_{X/Y}$。因此 $U$ 上的映射延拓到 $X$，证明完成。
\end{proof}






\section{相干模的零延拓}
\label{duality-section-extension-by-zero}

\noindent
本节及随后几节的材料可见 \cite{RD} 中 Deligne 所写的附录。

\medskip\noindent
本节中，$j : U \to X$ 是 Noetherian 概形的开浸入。考虑如下构造的
$D^b_{\textit{Coh}}(\mathcal{O}_X)$ 中的逆系统 $(K_n)$。设
$\mathcal{F}^\bullet$ 是相干 $\mathcal{O}_X$-模的有界复形，并设
$\mathcal{I} \subset \mathcal{O}_X$ 是拟相干理想层，满足
$V(\mathcal{I}) = X \setminus U$。可令
$$
K_n = \mathcal{I}^n\mathcal{F}^\bullet
$$
。更精确地，$K_n$ 是 $D^b_{\textit{Coh}}(\mathcal{O}_X)$ 中由如下复形
表示的对象：它的次数 $q$ 项是 $\mathcal{F}^q$ 的相干子模
$\mathcal{I}^n\mathcal{F}^q$。注意，映射
$\ldots \to K_3 \to K_2 \to K_1$ 限制到 $U$ 后诱导同构。称这样的系统为
{\it Deligne 系统}。

\begin{lemma}
\label{duality-lemma-lift-map}
设 $j : U \to X$ 是 Noetherian 概形的开浸入。设 $(K_n)$ 是 Deligne 系统，
并以 $K \in D^b_{\textit{Coh}}(\mathcal{O}_U)$ 表示常值系统
$(K_n|_U)$ 的值。设 $L$ 是 $D^b_{\textit{Coh}}(\mathcal{O}_X)$ 中的对象。
则 $\colim \Hom_X(K_n, L) = \Hom_U(K, L|_U)$。
\end{lemma}

\begin{proof}
设 $L \to M \to N \to L[1]$ 是
$D^b_{\textit{Coh}}(\mathcal{O}_X)$ 中的可区别三角。得到交换图
$$
\xymatrix{
\ldots \ar[r] &
\colim \Hom_X(K_n, L) \ar[r] \ar[d] &
\colim \Hom_X(K_n, M) \ar[r] \ar[d] &
\colim \Hom_X(K_n, N) \ar[r] \ar[d] &
\ldots \\
\ldots \ar[r] &
\Hom_U(K, L|_U) \ar[r] &
\Hom_U(K, M|_U) \ar[r] &
\Hom_U(K, N|_U) \ar[r] &
\ldots
}
$$
。由《导出范畴》，引理 \ref{derived-lemma-representable-homological} 和
《代数》，引理 \ref{algebra-lemma-directed-colimit-exact}，两行都正合。
因此，若引理的陈述对 $N[-1]$、$L$、$N$、$L[1]$ 成立，则由五引理，
它对 $M$ 也成立。由此，利用 $L$ 的典范截断的可区别三角（见《导出范畴》，
注 \ref{derived-remark-truncation-distinguished-triangle}），可归结到 $L$
仅有一个非零上同调层的情形。

\medskip\noindent
选取相干 $\mathcal{O}_X$-模的有界复形 $\mathcal{F}^\bullet$，以及定义
$X \setminus U$ 的拟相干理想 $\mathcal{I} \subset \mathcal{O}_X$，使
$K_n$ 由 $\mathcal{I}^n\mathcal{F}^\bullet$ 表示。利用“朴素”截断，得到逐项
分裂且彼此相容的复形短正合序列
$$
0 \to \sigma_{\geq a + 1} \mathcal{I}^n\mathcal{F}^\bullet \to
\mathcal{I}^n\mathcal{F}^\bullet \to
\sigma_{\leq a} \mathcal{I}^n\mathcal{F}^\bullet \to 0
$$
，它们进而对应于 $D^b_{\textit{Coh}}(\mathcal{O}_X)$ 中彼此相容的
可区别三角系统。重复上述论证，归结到 $\mathcal{F}^\bullet$
仅有一个非零项的情形，也就是下一段所讨论的情形。

\medskip\noindent
给定相干 $\mathcal{O}_X$-模 $\mathcal{F}$ 和相干 $\mathcal{O}_X$-模
$\mathcal{G}$，须证明对所有 $i \geq 0$，典范映射
$$
\colim \Ext^i_X(\mathcal{I}^n\mathcal{F}, \mathcal{G})
\longrightarrow
\Ext^i_U(\mathcal{F}|_U, \mathcal{G}|_U)
$$
都是同构。当 $i = 0$ 时，这就是《概形的上同调》，引理
\ref{coherent-lemma-homs-over-open}。以下假设 $i > 0$。

\medskip\noindent
单射性。设 $\xi \in \Ext^i_X(\mathcal{I}^n\mathcal{F}, \mathcal{G})$
的限制在 $U$ 上为零。须证明存在 $m \geq n$，使 $\xi$ 到
$\mathcal{I}^m\mathcal{F} = \mathcal{I}^{m - n}\mathcal{I}^n\mathcal{F}$
上的限制为零。以 $\mathcal{I}^n\mathcal{F}$ 替换 $\mathcal{F}$ 后，
可以假设 $n = 0$；即有
$\xi \in \Ext^i_X(\mathcal{F}, \mathcal{G})$，其限制在 $U$ 上为零。
由《概形的导出范畴》，命题 \ref{perfect-proposition-DCoh}，有
$D^b_{\textit{Coh}}(\mathcal{O}_X) = D^b(\textit{Coh}(\mathcal{O}_X))$。
因此可以在相干 $\mathcal{O}_X$-模的 Abel 范畴中计算 $\Ext$ 群。
这说明存在满射 $\alpha : \mathcal{F}'' \to \mathcal{F}$，使得
$\xi \circ \alpha = 0$（这里使用了 $i > 0$）。令
$\mathcal{F}' = \Ker(\alpha)$，于是有短正合序列
$$
0 \to \mathcal{F}' \to \mathcal{F}'' \to \mathcal{F} \to 0
$$
。所以 $\xi$ 是某个元素
$\xi' \in \Ext^{i - 1}_X(\mathcal{F}', \mathcal{G})$ 的像，且该元素
到 $U$ 上的限制属于映射
$\Ext^{i - 1}_U(\mathcal{F}''|_U, \mathcal{G}|_U) \to
\Ext^{i - 1}_U(\mathcal{F}'|_U, \mathcal{G}|_U)$ 的像。
由 Artin--Rees，逆系统 $(\mathcal{I}^n\mathcal{F}')$ 和
$(\mathcal{I}^n \mathcal{F}'' \cap \mathcal{F}')$ pro-同构；见《概形的上同调》，
引理 \ref{coherent-lemma-Artin-Rees}。由于有相容的短正合序列系统
$$
0 \to
\mathcal{F}' \cap \mathcal{I}^n\mathcal{F}'' \to
\mathcal{I}^n\mathcal{F}'' \to
\mathcal{I}^n\mathcal{F} \to 0
$$
，得到交换图
$$
\xymatrix{
\colim \Ext^{i - 1}_X(\mathcal{I}^n\mathcal{F}'', \mathcal{G})
\ar[r] \ar[d] &
\colim \Ext^{i - 1}_X(\mathcal{F}' \cap \mathcal{I}^n\mathcal{F}'', \mathcal{G})
\ar[r] \ar[d] &
\colim \Ext^i_X(\mathcal{I}^n\mathcal{F}, \mathcal{G})
\ar[d] \\
\Ext^{i - 1}_U(\mathcal{F}''|_U, \mathcal{G}|_U) \ar[r] &
\Ext^{i - 1}_U(\mathcal{F}'|_U, \mathcal{G}|_U) \ar[r] &
\Ext^{i - 1}_U(\mathcal{F}|_U, \mathcal{G}|_U)
}
$$
，两行都正合。对 $i$ 作归纳并使用上面对逆系统的说明，可知左边两个竖直箭头
都是同构。现在 $\xi$ 给出右上群中的元素，它是中上群中 $\xi'$ 的像；
后者又映到中下群中的一个元素，而该元素来自左下群中的某个元素。因此，对某个
$n$，$\xi$ 在 $\Ext^i_X(\mathcal{I}^n\mathcal{F}, \mathcal{G})$
中的像为零，正合所求。

\medskip\noindent
满射性。设 $\xi \in \Ext^i_U(\mathcal{F}|_U, \mathcal{G}|_U)$。
像上面一样论证并使用 $i > 0$，可找到相干 $\mathcal{O}_U$-模的满射
$\mathcal{H} \to \mathcal{F}|_U$，使 $\xi$ 在
$\Ext^i_U(\mathcal{H}, \mathcal{G}|_U)$ 中的像为零。然后可找到相干
$\mathcal{O}_X$-模的映射 $\varphi : \mathcal{F}'' \to \mathcal{F}$，
它限制到 $U$ 后就是 $\mathcal{H} \to \mathcal{F}|_U$；见《性质》，引理
\ref{properties-lemma-extend-finite-presentation}。注意，该引理并不保证
$\varphi$ 是满射，但这不影响论证（多做一点工作也可以选到满的 $\varphi$）。
记 $\mathcal{F}' = \Ker(\varphi)$。短正合序列
$$
0 \to \mathcal{F}'|_U \to \mathcal{F}''|_U \to \mathcal{F}|_U \to 0
$$
说明 $\xi$ 是
$\Ext^{i - 1}_U(\mathcal{F}'|_U, \mathcal{G}|_U)$ 中 $\xi'$ 的像。
对 $i$ 作归纳，可找到 $n$，使 $\xi'$ 是
$\Ext^{i - 1}_X(\mathcal{I}^n\mathcal{F}', \mathcal{G})$ 中某个
$\xi'_n$ 的像。由 Artin--Rees，可找到 $m \geq n$，使
$\mathcal{F}' \cap \mathcal{I}^m\mathcal{F}'' \subset
\mathcal{I}^n\mathcal{F}'$。利用短正合序列
$$
0 \to \mathcal{F}' \cap \mathcal{I}^m\mathcal{F}'' \to
\mathcal{I}^m\mathcal{F}'' \to \mathcal{I}^m\Im(\varphi) \to 0
$$
，$\xi'_n$ 在
$\Ext^{i - 1}_X(\mathcal{F}' \cap \mathcal{I}^m\mathcal{F}'', \mathcal{G})$
中的像，经边界映射映到
$\Ext^i_X(\mathcal{I}^m\Im(\varphi), \mathcal{G})$ 中的元素 $\xi_m$，
后者映到 $\xi$。由于 $\Im(\varphi)$ 与 $\mathcal{F}$ 在 $U$ 上相同，
$\mathcal{F}/\mathcal{I}^m\Im(\varphi)$ 支撑于 $X \setminus U$。
因此存在 $l \geq m$，使
$\mathcal{I}^l\mathcal{F} \subset \mathcal{I}^m\Im(\varphi)$；见《概形的上同调》，
引理 \ref{coherent-lemma-power-ideal-kills-sheaf}。取 $\xi_m$ 在
$\Ext^i_X(\mathcal{I}^l\mathcal{F}, \mathcal{G})$ 中的像即得结论。
\end{proof}

\begin{lemma}
\label{duality-lemma-lift-map-plus}
引理 \ref{duality-lemma-lift-map} 的结论对
$L \in D^+_{\textit{Coh}}(\mathcal{O}_X)$ 仍然成立。
\end{lemma}

\begin{proof}
事实上，若 $(K_n)$ 是 Deligne 系统，则存在 $b \in \mathbf{Z}$，使得当
$i > b$ 时 $H^i(K_n) = 0$。于是
$\Hom(K_n, L) = \Hom(K_n, \tau_{\leq b}L)$ 且
$\Hom(K, L) = \Hom(K, \tau_{\leq b}L)$。把该引理的结论用于
$\tau_{\leq b}L$ 即得。
\end{proof}

\begin{lemma}
\label{duality-lemma-extension-by-zero}
设 $j : U \to X$ 是 Noetherian 概形的开浸入。
\begin{enumerate}
\item 设 $(K_n)$ 和 $(L_n)$ 是 Deligne 系统，并设 $K$ 和 $L$ 分别是
常值系统 $(K_n|_U)$ 和 $(L_n|_U)$ 的值。给定 $D(\mathcal{O}_U)$ 中的态射
$\alpha : K \to L$，存在唯一的 $D^b_{\textit{Coh}}(\mathcal{O}_X)$
的 pro-系统态射 $(K_n) \to (L_n)$，其限制到 $U$ 后为 $\alpha$；
\item 给定 $K \in D^b_{\textit{Coh}}(\mathcal{O}_U)$，存在 Deligne 系统
$(K_n)$，使 $(K_n|_U)$ 是值为 $K$ 的常值系统；
\item (2) 中 $D^b_{\textit{Coh}}(\mathcal{O}_X)$ 的 pro-对象 $(K_n)$
在唯一同构意义下唯一（作为 pro-对象）。
\end{enumerate}
\end{lemma}

\begin{proof}
(1) 直接来自引理 \ref{duality-lemma-lift-map}，以及 pro-系统之间的态射等同于它们
所共表示的函子之间的态射这一事实；见《范畴》，注
\ref{categories-remark-pro-category-copresheaves}。

\medskip\noindent
设 $K$ 如 (2) 所述。可以选取 $K' \in D^b_{\textit{Coh}}(\mathcal{O}_X)$，
使其限制到 $U$ 后同构于 $K$；见《概形的导出范畴》，引理
\ref{perfect-lemma-lift-coherent}。由《概形的导出范畴》，命题
\ref{perfect-proposition-DCoh}，可用相干 $\mathcal{O}_X$-模的有界复形
$\mathcal{F}^\bullet$ 表示 $K'$。选取拟相干理想层
$\mathcal{I} \subset \mathcal{O}_X$，其零化轨迹为 $X \setminus U$
（例如令 $\mathcal{I}$ 对应于 $X \setminus U$ 上的约化诱导子概形结构）。
于是可像本节开头那样，令 $K_n$ 为复形
$\mathcal{I}^n\mathcal{F}^\bullet$ 所表示的对象。

\medskip\noindent
(3) 直接由 (1) 和 (2) 得到。
\end{proof}

\begin{lemma}
\label{duality-lemma-extension-by-zero-triangle}
设 $j : U \to X$ 是 Noetherian 概形的开浸入，并设
$$
K \to L \to M \to K[1]
$$
是 $D^b_{\textit{Coh}}(\mathcal{O}_U)$ 中的可区别三角。则存在
$D^b_{\textit{Coh}}(\mathcal{O}_X)$ 中的可区别三角的逆系统
$$
K_n \to L_n \to M_n \to K_n[1]
$$
，使 $(K_n)$、$(L_n)$、$(M_n)$ 都是 Deligne 系统，并且这些可区别三角
限制到 $U$ 后同构于原来的可区别三角。
\end{lemma}

\begin{proof}
取引理 \ref{duality-lemma-extension-by-zero} 的 (2) 中那样的 $(K_n)$。
选取 $D^b_{\textit{Coh}}(\mathcal{O}_X)$ 的对象 $L'$，使其限制到 $U$
后为 $L$（该引理说明这是可以做到的）。由引理 \ref{duality-lemma-lift-map}，
可找到 $n$ 以及 $X$ 上的态射 $K_n \to L'$，其限制到 $U$ 后为给定箭头
$K \to L$。因此存在 $D^b_{\textit{Coh}}(\mathcal{O}_X)$ 中的态射
$K' \to L'$，其限制到 $U$ 后为给定箭头 $K \to L$。

\medskip\noindent
由《概形的导出范畴》，命题 \ref{perfect-proposition-DCoh}，可找到相干
$\mathcal{O}_X$-模的有界复形之间的态射
$\alpha^\bullet : \mathcal{F}^\bullet \to
\mathcal{G}^\bullet$，它表示 $K' \to L'$。选取拟相干理想层
$\mathcal{I} \subset \mathcal{O}_X$，其零化轨迹为 $X \setminus U$。
令 $K_n = \mathcal{I}^n\mathcal{F}^\bullet$，且
$L_n = \mathcal{I}^n\mathcal{G}^\bullet$。注意，$\alpha^\bullet$ 诱导复形态射
$\alpha_n^\bullet : \mathcal{I}^n\mathcal{F}^\bullet \to
\mathcal{I}^n\mathcal{G}^\bullet$。由《导出范畴》，第
\ref{derived-section-cones} 节中的锥构造，显然有
$$
C(\alpha_n)^\bullet = \mathcal{I}^nC(\alpha^\bullet)
$$
，故可令 $M_n = C(\alpha_n)^\bullet$。也就是说，得到相容的可区别三角系统
（见《导出范畴》，第 \ref{derived-section-canonical-delta-functor} 节中的讨论）
$$
K_n \to L_n \to M_n \to K_n[1]
$$
；由公理 TR3 及《导出范畴》，引理
\ref{derived-lemma-third-isomorphism-triangle}，它限制到 $U$ 后同构于原来的
可区别三角。
\end{proof}

\begin{remark}
\label{duality-remark-extension-by-zero}
设 $j : U \to X$ 是 Noetherian 概形的开浸入。把
$K \in D^b_{\textit{Coh}}(\mathcal{O}_U)$ 送到其限制在 $U$ 上为 $K$
的 Deligne 系统，这决定了函子
$$
Rj_! :
D^b_{\textit{Coh}}(\mathcal{O}_U)
\longrightarrow
\text{Pro-}D^b_{\textit{Coh}}(\mathcal{O}_X)
$$
。由引理 \ref{duality-lemma-extension-by-zero-triangle}，该函子是“正合的”；
由引理 \ref{duality-lemma-lift-map}，它是函子
$j^* : D^b_{\textit{Coh}}(\mathcal{O}_X) \to
D^b_{\textit{Coh}}(\mathcal{O}_U)$ 的“左伴随”。
\end{remark}

\begin{remark}
\label{duality-remark-extension-by-zero-linear-pro-system}
设 $(A_n)$ 和 $(B_n)$ 是范畴 $\mathcal{C}$ 的逆系统。所谓从 $(A_n)$ 到
$(B_n)$ 的线性 pro-态射，是指对某些固定整数 $c, d \geq 1$，给定对所有
$n \geq 1$ 都彼此相容的态射族 $\varphi_n : A_{cn + d} \to B_n$。
若存在 $c'' \geq \max(c, c')$ 和 $d'' \geq \max(d, d')$，使得对所有 $n$，
所诱导的两个态射 $A_{c'' n + d''} \to B_n$ 都相同，则称
$(\varphi_n : A_{cn + d} \to B_n)$ 与
$(\psi_n : A_{c'n + d'} \to B_n)$ 决定同一个态射。给定 $U$ 上的值后，
Deligne 系统 $(K_n)$ 很可能在线性 pro-同构意义下是良定义的。如果以后需要，
我们将在这里仔细陈述并证明这一点。
\end{remark}

\begin{lemma}
\label{duality-lemma-deligne-system-2-out-of-3}
设 $j : U \to X$ 是 Noetherian 概形的开浸入。设
$$
K_n \to L_n \to M_n \to K_n[1]
$$
是 $D^b_{\textit{Coh}}(\mathcal{O}_X)$ 中可区别三角的逆系统。若
$(K_n)$ 和 $(M_n)$ 都 pro-同构于 Deligne 系统，则 $(L_n)$ 也是如此。
\end{lemma}

\begin{proof}
由于 $(K_n|_U)$ 和 $(M_n|_U)$ 都 pro-同构于常值系统，它们本质常值。
以 $K$ 和 $M$ 表示其值。由《导出范畴》，引理
\ref{derived-lemma-essentially-constant-2-out-of-3}，逆系统 $L_n|_U$
也本质常值；以 $L$ 表示其值。取
$N \in D^b_{\textit{Coh}}(\mathcal{O}_X)$，考虑交换图
$$
\xymatrix{
\ldots \ar[r] &
\colim \Hom_X(M_n, N) \ar[r] \ar[d] &
\colim \Hom_X(L_n, N) \ar[r] \ar[d] &
\colim \Hom_X(K_n, N) \ar[r] \ar[d] &
\ldots \\
\ldots \ar[r] &
\Hom_U(M, N|_U) \ar[r] &
\Hom_U(L, N|_U) \ar[r] &
\Hom_U(K, N|_U) \ar[r] &
\ldots
}
$$
。由引理 \ref{duality-lemma-lift-map} 以及同构的 ind-系统有相同余极限这一事实，
中间箭头左右各两个竖直箭头都是同构。由五引理，中间的竖直箭头也是同构。
现在，设 $(L'_n)$ 是限制到 $U$ 后常值为 $L$ 的 Deligne 系统
（由引理 \ref{duality-lemma-extension-by-zero}，它存在），则同样有
$\colim \Hom_X(L'_n, N) = \Hom_U(L, N|_U)$。因此由《范畴》，注
\ref{categories-remark-pro-category-copresheaves}，pro-系统 $(L_n)$ 和
$(L'_n)$ pro-同构。
\end{proof}

\begin{lemma}
\label{duality-lemma-consequence-Artin-Rees-bis}
设 $X$ 是 Noetherian 概形，$\mathcal{I} \subset \mathcal{O}_X$
是拟相干理想层，$\mathcal{F}^\bullet$ 是相干 $\mathcal{O}_X$-模的复形，
且 $p \in \mathbf{Z}$。令 $\mathcal{H} = H^p(\mathcal{F}^\bullet)$ 和
$\mathcal{H}_n = H^p(\mathcal{I}^n\mathcal{F}^\bullet)$。则有典范的
$\mathcal{O}_X$-模映射
$$
\ldots \to \mathcal{H}_3 \to \mathcal{H}_2 \to \mathcal{H}_1 \to \mathcal{H}
$$
。存在 $c > 0$，使得对 $n \geq c$，$\mathcal{H}_n \to \mathcal{H}$
的像包含于 $\mathcal{I}^{n - c}\mathcal{H}$，并有典范的
$\mathcal{O}_X$-模映射
$\mathcal{I}^n\mathcal{H} \to \mathcal{H}_{n - c}$，使复合
$$
\mathcal{I}^n \mathcal{H} \to \mathcal{H}_{n - c} \to
\mathcal{I}^{n - 2c}\mathcal{H}
\quad\text{且}\quad
\mathcal{H}_n \to \mathcal{I}^{n - c}\mathcal{H} \to \mathcal{H}_{n - 2c}
$$
都是典范映射。特别地，逆系统 $(\mathcal{H}_n)$ 和
$(\mathcal{I}^n\mathcal{H})$ 作为 $\textit{Mod}(\mathcal{O}_X)$
的 pro-对象同构。
\end{lemma}

\begin{proof}
若 $X$ 仿射，翻译为代数后，这就是《代数进阶》，引理
\ref{more-algebra-lemma-consequence-Artin-Rees-bis}。一般情形中，完全照搬
该引理的证明，只把代数中的 Artin--Rees 引用替换为《概形的上同调》，引理
\ref{coherent-lemma-Artin-Rees}。略去细节。
\end{proof}

\begin{lemma}
\label{duality-lemma-characterize-extension-by-zero-algebra}
设 $j : U \to X$ 是 Noetherian 概形的开浸入，$a \leq b$ 是整数，
且 $(K_n)$ 是 $D^b_{\textit{Coh}}(\mathcal{O}_X)$ 的逆系统，满足当
$i \not \in [a, b]$ 时 $H^i(K_n) = 0$。则下列条件等价：
\begin{enumerate}
\item $(K_n)$ pro-同构于一个 Deligne 系统；
\item 对每个 $p \in \mathbf{Z}$，存在相干 $\mathcal{O}_X$-模
$\mathcal{F}$，使 pro-系统 $(H^p(K_n))$ 和
$(\mathcal{I}^n\mathcal{F})$ pro-同构。
\end{enumerate}
\end{lemma}

\begin{proof}
假设 (1)。为证明 (2)，可以假设 $(K_n)$ 本身是 Deligne 系统。
按定义，可以选取相干 $\mathcal{O}_X$-模的有界复形
$\mathcal{F}^\bullet$，以及定义 $X \setminus U$ 的拟相干理想层，使
$K_n$ 由 $\mathcal{I}^n\mathcal{F}^\bullet$ 表示。结论于是来自引理
\ref{duality-lemma-consequence-Artin-Rees-bis}。

\medskip\noindent
假设 (2)。对 $b - a$ 作归纳，证明 $(K_n)$ 满足 (1)。若 $a = b$，
则 (1) 本质上就是假设。若 $a < b$，考虑相容的可区别三角系统
$$
\tau_{\leq a}K_n \to K_n \to \tau_{\geq a + 1}K_n \to (\tau_{\leq a}K_n)[1]
$$
；见《导出范畴》，注
\ref{derived-remark-truncation-distinguished-triangle}。对 $b - a$ 作归纳，
可知 $\tau_{\leq a}K_n$ 和 $\tau_{\geq a + 1}K_n$ 都 pro-同构于
Deligne 系统。再由引理 \ref{duality-lemma-deligne-system-2-out-of-3} 即得结论。
\end{proof}

\begin{lemma}
\label{duality-lemma-characterize-extension-by-zero}
设 $j : U \to X$ 是 Noetherian 概形的开浸入，$(K_n)$ 是
$D^b_{\textit{Coh}}(\mathcal{O}_X)$ 中的逆系统，并设
$X = W_1 \cup \ldots \cup W_r$ 是开覆盖。则下列条件等价：
\begin{enumerate}
\item $(K_n)$ pro-同构于一个 Deligne 系统；
\item 对每个 $i$，限制 $(K_n|_{W_i})$ 都 pro-同构于开浸入
$U \cap W_i \to W_i$ 的一个 Deligne 系统。
\end{enumerate}
\end{lemma}

\begin{proof}
对 $r$ 作归纳。若 $r = 1$，结论显然。假设 $r > 1$，令
$V = W_1 \cup \ldots \cup W_{r - 1}$。由归纳假设，$(K_n|_V)$
是 Deligne 系统。于是归结到下一段的讨论。

\medskip\noindent
假设 $X = V \cup W$ 是开覆盖，且 $(K_n|_W)$、$(K_n|_V)$ 都
pro-同构于 Deligne 系统。须证明 $(K_n)$ pro-同构于一个 Deligne 系统。
注意，$(K_n|_{V \cap W})$ pro-同构于一个 Deligne 系统
（由 Deligne 系统的构造，限制到开集显然保持这项性质）。特别地，pro-系统
$(K_n|_{U \cap V})$、$(K_n|_{U \cap W})$、
$(K_n|_{U \cap V \cap W})$ 都本质常值。由《上同调》，引理
\ref{cohomology-lemma-exact-sequence-j-star} 中的可区别三角和《导出范畴》，
引理 \ref{derived-lemma-essentially-constant-2-out-of-3}，可知
$(K_n|_U)$ 本质常值。以
$K \in D^b_{\textit{Coh}}(\mathcal{O}_U)$ 表示该系统的值。设 $L$
是 $D^b_{\textit{Coh}}(\mathcal{O}_X)$ 中的对象，考虑图
$$
\xymatrix{
\colim \Ext^{-1}(K_n|_V, L|_V) \oplus
\colim \Ext^{-1}(K_n|_W, L|_W) \ar[r] \ar[d] &
\Ext^{-1}(K|_{U \cap V}, L|_{U \cap V}) \oplus
\Ext^{-1}(K|_{U \cap W}, L|_{U \cap W}) \ar[d] \\
\colim \Ext^{-1}(K_n|_{V \cap W}, L|_{V \cap W}) \ar[r] \ar[d] &
\Ext^{-1}(K|_{U \cap V \cap W}, L|_{U \cap V \cap W}) \ar[d] \\
\colim \Hom(K_n, L) \ar[d] \ar[r] &
\Hom(K|_U, L|_U) \ar[d] \\
\colim \Hom(K_n|_V, L|_V) \oplus \colim \Hom(K_n|_W, L|_W) \ar[r] \ar[d] &
\Hom(K|_{U \cap V}, L|_{U \cap V}) \oplus
\Hom(K|_{U \cap W}, L|_{U \cap W}) \ar[d] \\
\colim \Hom(K_n|_{V \cap W}, L|_{V \cap W}) \ar[r] &
\Hom(K|_{U \cap V \cap W}, L|_{U \cap V \cap W})
}
$$
。由《上同调》，引理 \ref{cohomology-lemma-mayer-vietoris-hom}，
以及滤过余极限正合这一事实，两个竖直序列都正合。由引理
\ref{duality-lemma-lift-map} 以及 pro-同构的系统有相同余极限这一事实，除中间箭头外
的所有水平箭头都是同构。由五引理，中间箭头也是同构。因此 $(K_n)$
pro-同构于 $K$ 的一个 Deligne 系统。即若 $(K'_n)$ 是限制到 $U$
后常值为 $K$ 的 Deligne 系统（由引理 \ref{duality-lemma-extension-by-zero}，
它存在），则同样有 $\colim \Hom_X(K'_n, L) = \Hom_U(K, L|_U)$。
故由《范畴》，注 \ref{categories-remark-pro-category-copresheaves}，
pro-系统 $(K_n)$ 和 $(K'_n)$ pro-同构。
\end{proof}

\begin{lemma}
\label{duality-lemma-tensoring-Deligne-system}
设 $j : U \to X$ 是 Noetherian 概形的开浸入，并设拟相干理想层
$\mathcal{I} \subset \mathcal{O}_X$ 满足 $V(\mathcal{I}) = X \setminus U$。
设 $K$ 属于 $D^b_{\textit{Coh}}(\mathcal{O}_X)$。则
$$
K \otimes_{\mathcal{O}_X}^\mathbf{L} \mathcal{I}^n
$$
pro-同构于一个 Deligne 系统，它在 $U$ 上的常值为 $K|_U$。
\end{lemma}

\begin{proof}
由引理 \ref{duality-lemma-characterize-extension-by-zero}，该问题在 $X$ 上是局部的。
因此可以假设 $X$ 是某个 Noetherian 环的谱。此时结论来自代数版本，
即《代数进阶》，引理 \ref{more-algebra-lemma-tensoring-Deligne-system}。
\end{proof}







\section{紧支撑上同调的预备知识}
\label{duality-section-preliminaries-compactly-supported}

\noindent
在情形 \ref{duality-situation-shriek} 中，设 $f : X \to Y$ 是范畴
$\textit{FTS}_S$ 中的态射。利用上一节的构造，将构造函子
$$
Rf_! :
D^b_{\textit{Coh}}(\mathcal{O}_X)
\longrightarrow
\text{Pro-}D^b_{\textit{Coh}}(\mathcal{O}_Y)
$$
。若 $f$ 是开浸入，它退化为注 \ref{duality-remark-extension-by-zero} 中的函子；
一般地，它使用 $f$ 的紧化来构造。在此之前，需要下面的引理来证明该构造良定义。

\begin{lemma}
\label{duality-lemma-well-defined-pre}
设 $f : X \to Y$ 是 Noetherian 概形的固有态射，$V \subset Y$
是开子概形，并令 $U = f^{-1}(V)$。图示为
$$
\xymatrix{
U \ar[r]_j \ar[d]_g & X \ar[d]^f \\
V \ar[r]^{j'} & Y
}
$$
。则有函子的典范同构 $Rj'_! \circ Rg_* \to Rf_* \circ Rj_!$，其定义域和值域为
$D^b_{\textit{Coh}}(\mathcal{O}_U) \to
\text{Pro-}D^b_{\textit{Coh}}(\mathcal{O}_Y)$；其中 $Rj_!$ 和
$Rj'_!$ 如注 \ref{duality-remark-extension-by-zero} 中所述。
\end{lemma}

\begin{proof}[第一个证明]
设 $K$ 是 $D^b_{\textit{Coh}}(\mathcal{O}_U)$ 中的对象。设 $(K_n)$
是 $U \to X$ 的 Deligne 系统，其限制到 $U$ 后常值为 $K$。
这当然意味着 $(K_n)$ 在 $\text{Pro-}D^b_{\textit{Coh}}(\mathcal{O}_X)$
中表示 $Rj_!K$。注意，$Rj'_!Rg_*K$ 和 $Rf_*Rj_!K$ 限制到 $V$
后都是值为 $Rg_*K$ 的常值 pro-对象。对前者这是直接的；对后者，它来自
$(Rf_*K_n)|_V = Rg_*(K_n|_U) = Rg_*K$。由引理
\ref{duality-lemma-extension-by-zero} 中 Deligne 系统的唯一性，只须证明
$(Rf_*K_n)$ pro-同构于一个 Deligne 系统。所引引理还会说明所得同构具有函子性。

\medskip\noindent
证明 $(Rf_*K_n)$ pro-同构于一个 Deligne 系统。首先，由上述唯一性，
该问题与对应于 $K$ 的 Deligne 系统 $(K_n)$ 的选择无关。由引理
\ref{duality-lemma-extension-by-zero-triangle} 和
\ref{duality-lemma-deligne-system-2-out-of-3}，若
$D^b_{\textit{Coh}}(\mathcal{O}_U)$ 中有可区别三角
$$
K \to L \to M \to K[1]
$$
，且结论对 $K$ 和 $M$ 成立，则它对 $L$ 也成立。利用典范截断的可区别三角
（《导出范畴》，注 \ref{derived-remark-truncation-distinguished-triangle}），
归结到下一段研究的问题。

\medskip\noindent
设 $\mathcal{F}$ 是相干 $\mathcal{O}_X$-模，并设
$\mathcal{J} \subset \mathcal{O}_Y$ 是定义 $Y \setminus V$ 的拟相干理想层。
以 $\mathcal{J}^n\mathcal{F}$ 表示
$f^*\mathcal{J}^n \otimes \mathcal{F} \to \mathcal{F}$ 的像。
须证明 $(Rf_*(\mathcal{J}^n\mathcal{F}))$ 是 Deligne 系统。由引理
\ref{duality-lemma-characterize-extension-by-zero}，问题在 $Y$ 上是局部的。
因此可以假设 $Y = \Spec(A)$ 仿射，且 $\mathcal{J}$ 对应于理想
$I \subset A$。由引理 \ref{duality-lemma-characterize-extension-by-zero-algebra}，
只须证明对某个有限 $A$-模 $M$，上同调模的逆系统
$(H^p(X, I^n\mathcal{F}))$ pro-同构于逆系统 $(I^n M)$。
这已在《概形的上同调》，引理
\ref{coherent-lemma-cohomology-powers-ideal-application} 中证明。
\end{proof}

\begin{proof}[第二个证明]
设 $K$ 是 $D^b_{\textit{Coh}}(\mathcal{O}_U)$ 中的对象，$L$ 是
$D^b_{\textit{Coh}}(\mathcal{O}_Y)$ 中的对象。将构造双射
$$
\Hom_{\text{Pro-}D^b_{\textit{Coh}}(\mathcal{O}_Y)}(Rj'_!Rg_*K, L)
\longrightarrow
\Hom_{\text{Pro-}D^b_{\textit{Coh}}(\mathcal{O}_Y)}(Rf_*Rj_!K, L)
$$
，它对 $K$ 和 $L$ 都具有函子性。固定 $K$，由《范畴》，注
\ref{categories-remark-pro-category-copresheaves}，这将决定 pro-对象的同构
$Rf_*Rj_!K \to Rj'_!Rg_*K$；让 $K$ 变化，便得到函子的同构。
为实际构造该同构，使用下列函子性等式链：
\begin{align*}
\Hom_{\text{Pro-}D^b_{\textit{Coh}}(\mathcal{O}_Y)}(Rj'_!Rg_*K, L)
& =
\Hom_V(Rg_*K, L|_V) \\
& =
\Hom_U(K, g^!(L|_V)) \\
& =
\Hom_U(K, f^!L|_U)) \\
& =
\Hom_{\text{Pro-}D^b_{\textit{Coh}}(\mathcal{O}_X)}(Rj_!K, f^!L) \\
& =
\Hom_{\text{Pro-}D^b_{\textit{Coh}}(\mathcal{O}_Y)}(Rf_*Rj_!K, L)
\end{align*}
。第一个等式来自引理 \ref{duality-lemma-lift-map}。第二个等式成立，是因为 $g$
作为 $f$ 到 $V$ 的基变换是固有的，故按构造 $g^!$ 是推前的右伴随；见第
\ref{duality-section-upper-shriek} 节。第三个等式成立，是因为由引理
\ref{duality-lemma-restrict-before-or-after}，有 $g^!(L|_V) = f^!L|_U$。
由引理 \ref{duality-lemma-shriek-coherent}，$f^!L$ 属于
$D^+_{\textit{Coh}}(\mathcal{O}_X)$，故第四个等式来自引理
\ref{duality-lemma-lift-map-plus}。第五个等式再次来自：$f$ 固有时，$f^!$
是 $Rf_*$ 的右伴随。
\end{proof}

\begin{lemma}
\label{duality-lemma-well-defined}
设 $j : U \to X$ 是 Noetherian 概形的开浸入，$j' : U \to X'$
是 $U$ 在 $X$ 上的一个紧化（见证明），并以 $f : X' \to X$
表示结构态射。则有函子的典范同构 $Rj_! \to Rf_* \circ R(j')_!$，
其定义域和值域为 $D^b_{\textit{Coh}}(\mathcal{O}_U) \to
\text{Pro-}D^b_{\textit{Coh}}(\mathcal{O}_X)$；其中 $Rj_!$ 和
$Rj'_!$ 如注 \ref{duality-remark-extension-by-zero} 中所述。
\end{lemma}

\begin{proof}
$X'$ 是 $U$ 在 $X$ 上的紧化，恰好意味着 $f : X' \to X$ 固有、
$j'$ 是开浸入且 $j = f \circ j'$；见《平坦性进阶》，第
\ref{flat-section-compactify} 节。若 $j'(U) = f^{-1}(j(U))$，
则结论直接来自引理 \ref{duality-lemma-well-defined-pre}。若
$j'(U) \not = f^{-1}(j(U))$，以 $X'' \subset X'$ 表示
$j' : U \to X'$ 的概形论闭包，并以 $j'' : U \to X''$ 表示相应的开浸入。
图示为
$$
\xymatrix{
& & X'' \ar[d]^{f'} \\
& & X' \ar[d]^f \\
U \ar[rr]^j \ar[rru]^{j'} \ar[rruu]^{j''} & & X
}
$$
。由《平坦性进阶》，引理 \ref{flat-lemma-compactifications-cofiltered}
的 (c) 及上述讨论，有同构
$Rf'_* \circ Rj''_! = Rj'_!$ 和 $R(f \circ f')_* \circ Rj''_! = Rj_!$。
由于 $R(f \circ f')_* = Rf_* \circ Rf'_*$，结论成立。
\end{proof}

\begin{remark}
\label{duality-remark-covariance-open-j-lower-shriek}
设 $X \supset U \supset U'$ 是 Noetherian 概形 $X$ 的开子概形，
并以 $j : U \to X$ 和 $j' : U' \to X$ 表示包含态射。断言存在典范映射
$$
Rj'_!(K|_{U'}) \longrightarrow Rj_!K
$$
，它对 $D^b_{\textit{Coh}}(\mathcal{O}_U)$ 中的 $K$ 具有函子性。
事实上，由引理 \ref{duality-lemma-lift-map}，对
$D^b_{\textit{Coh}}(\mathcal{O}_X)$ 中任意 $L$，有映射
\begin{align*}
\Hom_{\text{Pro-}D^b_{\textit{Coh}}(\mathcal{O}_X)}(Rj_!K, L)
& =
\Hom_U(K, L|_U) \\
& \to
\Hom_{U'}(K|_{U'}, L|_{U'}) \\
& =
\Hom_{\text{Pro-}D^b_{\textit{Coh}}(\mathcal{O}_X)}(Rj'_!(K|_{U'}), L)
\end{align*}
，它对 $L$ 和 $K'$ 具有函子性。由对 $L$ 的函子性和《范畴》，注
\ref{categories-remark-pro-category-copresheaves}，得到典范映射
$Rj'_!(K|_{U'}) \to Rj_!K$；上述箭头对 $K$ 的函子性又说明该映射对 $K$
具有函子性。

\medskip\noindent
下面显式构造该箭头。设 $\mathcal{F}^\bullet$ 是相干
$\mathcal{O}_X$-模的有界复形，其限制到 $U$ 后在导出范畴中表示 $K$。
引理 \ref{duality-lemma-extension-by-zero} 的证明已说明这样的复形总是存在。
令 $\mathcal{I}$ 和\ $\mathcal{I}'$ 是 $X$ 上的拟相干理想层，分别满足
$V(\mathcal{I}) = X \setminus U$ 和\ $V(\mathcal{I}') = X \setminus U'$。
以 $\mathcal{I} + \mathcal{I}'$ 替换 $\mathcal{I}$ 后，可以假设
$\mathcal{I}' \subset \mathcal{I}$。按构造，$Rj_!K$ 和\ $Rj'_!(K|_{U'})$
分别由 $D^b_{\textit{Coh}}(\mathcal{O}_X)$ 的逆系统
$(K_n)$ 和\ $(K'_n)$ 表示，其中
$$
K_n = \mathcal{I}^n\mathcal{F}^\bullet
\quad\text{分别}\quad
K'_n = (\mathcal{I}')^n\mathcal{F}^\bullet
$$
。显然，上面构造的映射由包含
$(\mathcal{I}')^n \subset \mathcal{I}^n$ 所给出的映射 $K'_n \to K_n$ 给出。
\end{remark}






\section{相干模的紧支撑上同调}
\label{duality-section-compactly-supported}

\noindent
在情形 \ref{duality-situation-shriek} 中，给定 $\textit{FTS}_S$ 中的态射
$f : X \to Y$，定义函子
$$
Rf_! : D^b_{\textit{Coh}}(\mathcal{O}_X)
\longrightarrow
\text{Pro-}D^b_{\textit{Coh}}(\mathcal{O}_Y)
$$
。具体地，选取在 $Y$ 上的紧化 $j : X \to \overline{X}$；
由《平坦性进阶》，定理 \ref{flat-theorem-nagata} 和引理
\ref{flat-lemma-compactifyable}，这是可以做到的。以
$\overline{f} : \overline{X} \to Y$ 表示结构态射。对
$K \in D^b_{\textit{Coh}}(\mathcal{O}_X)$，令
$$
Rf_!K  = R\overline{f}_* Rj_! K
$$
，其中 $Rj_!$ 如注 \ref{duality-remark-extension-by-zero} 中所述。

\begin{lemma}
\label{duality-lemma-lower-shriek-well-defined}
函子 $Rf_!$ 在同构意义下与紧化的选择无关。
\end{lemma}

\noindent
事实上，函子 $Rf_!$ 将被刻画为 $f^!$ 的“左伴随”，从而在唯一同构意义下确定。

\begin{proof}
考虑 $X$ 在 $Y$ 上的紧化所成的范畴。由《平坦性进阶》，定理
\ref{flat-theorem-nagata} 及引理
\ref{flat-lemma-compactifications-cofiltered}、
\ref{flat-lemma-compactifyable}，该范畴余滤过。对紧化的每个选择
$$
j : X \to \overline{X},\quad \overline{f} : \overline{X} \to Y
$$
，上述构造都给出函子 $R\overline{f}_* \circ Rj_!$。设
$g : \overline{X}_1 \to \overline{X}_2$ 是 $Y$ 上两个紧化
$j_i : X \to \overline{X}_i$ 之间的态射。则得到同构
$$
R\overline{f}_{2, *} \circ Rj_{2, !} =
R\overline{f}_{2, *} \circ Rg_* \circ j_{1, !} =
R\overline{f}_{1, *} \circ Rj_{1, !}
$$
，其中第一个等式使用了引理 \ref{duality-lemma-well-defined}。由此可知，
该函子在同构意义下与紧化的选择无关。
\end{proof}

\begin{proposition}
\label{duality-proposition-duality-compactly-supported}
在情形 \ref{duality-situation-shriek} 中，设 $f : X \to Y$ 是
$\textit{FTS}_S$ 的态射。则函子 $Rf_!$ 与 $f^!$ 在下述意义下伴随：
对所有 $K \in D^b_{\textit{Coh}}(\mathcal{O}_X)$ 和
$L \in D^+_{\textit{Coh}}(\mathcal{O}_Y)$，有
$$
\Hom_X(K, f^!L) =
\Hom_{\text{Pro-}D^+_{\textit{Coh}}(\mathcal{O}_Y)}(Rf_!K, L)
$$
，并且它对 $K$ 和 $L$ 双函子性地成立。
\end{proposition}

\begin{proof}
选取在 $Y$ 上的紧化 $j : X \to \overline{X}$，并以
$\overline{f} : \overline{X} \to Y$ 表示结构态射。于是
\begin{align*}
\Hom_X(K, f^!L)
& =
\Hom_X(K, j^*\overline{f}{}^!L) \\
& =
\Hom_{\text{Pro-}D^+_{\textit{Coh}}(\mathcal{O}_{\overline{X}})}
(Rj_!K, \overline{f}{}^!L) \\
& =
\Hom_{\text{Pro-}D^+_{\textit{Coh}}(\mathcal{O}_Y)}(Rf_*Rj_!K, L) \\
& =
\Hom_{\text{Pro-}D^+_{\textit{Coh}}(\mathcal{O}_Y)}(Rf_!K, L)
\end{align*}
。第一个等式直接来自第 \ref{duality-section-upper-shriek} 节中 $f^!$ 的构造。
由引理 \ref{duality-lemma-shriek-coherent}，$\overline{f}{}^!L$ 属于
$D^+_{\textit{Coh}}(\mathcal{O}_{\overline{X}})$，故第二个等式来自引理
\ref{duality-lemma-lift-map-plus}。由于 $\overline{f}$ 固有，按构造，函子
$\overline{f}{}^!$ 是推前的右伴随；这给出第三个等式。第四个等式成立，
是因为 $Rf_! = Rf_* Rj_!$。
\end{proof}

\begin{lemma}
\label{duality-lemma-compactly-supported-triangle}
在情形 \ref{duality-situation-shriek} 中，设 $f : X \to Y$ 是
$\textit{FTS}_S$ 的态射，并设
$$
K \to L \to M \to K[1]
$$
是 $D^b_{\textit{Coh}}(\mathcal{O}_X)$ 中的可区别三角。则存在
$D^b_{\textit{Coh}}(\mathcal{O}_Y)$ 中的可区别三角的逆系统
$$
K_n \to L_n \to M_n \to K_n[1]
$$
，使 pro-系统 $(K_n)$、$(L_n)$、$(M_n)$ 分别给出
$Rf_!K$、$Rf_!L$、$Rf_!M$。
\end{lemma}

\begin{proof}
选取在 $Y$ 上的紧化 $j : X \to \overline{X}$，并以
$\overline{f} : \overline{X} \to Y$ 表示结构态射。按照引理
\ref{duality-lemma-extension-by-zero-triangle}，选取对应于开浸入 $j$
和给定可区别三角的逆系统
$$
\overline{K}_n \to \overline{L}_n \to \overline{M}_n \to \overline{K}_n[1]
$$
，它位于 $D^b_{\textit{Coh}}(\mathcal{O}_{\overline{X}})$ 中。
取 $K_n = R\overline{f}_*\overline{K}_n$，并类似地定义 $L_n$ 和 $M_n$。
由 $Rf_!$ 的定义，这满足要求。
\end{proof}

\begin{remark}
\label{duality-remark-compose-inverse-systems}
设 $\mathcal{C}$ 是一个范畴。假设给定 $\mathcal{C}$ 的 pro-对象范畴中的
逆系统的逆系统
$$
\ldots \xrightarrow{\alpha_4} (M_{3, n}) \xrightarrow{\alpha_3} (M_{2, n})
\xrightarrow{\alpha_2} (M_{1, n})
$$
。换言之，箭头 $\alpha_i$ 是 pro-对象的态射。由《范畴》，例
\ref{categories-example-pro-morphism-inverse-systems}，每个 $\alpha_i$
都可由偶对 $(m_i, a_i)$ 表示，其中
$m_i : \mathbf{N} \to \mathbf{N}$ 是递增函数，而
$a_{i, n} : M_{i, m_i(n)} \to M_{i - 1, n}$ 是 $\mathcal{C}$ 中的态射，
使图
$$
\xymatrix{
\ldots \ar[r] &
M_{i, m_i(3)} \ar[d]^{a_{i, 3}} \ar[r] &
M_{i, m_i(2)} \ar[d]^{a_{i, 2}} \ar[r] &
M_{i, m_i(1)} \ar[d]^{a_{i, 1}} \\
\ldots \ar[r] &
M_{i - 1, 3} \ar[r] &
M_{i - 1, 2} \ar[r] &
M_{i - 1, 1}
}
$$
交换。用 $\max(n, m_i(n))$ 替换 $m_i(n)$，并相应调整态射 $a_i(n)$
（如所引例中那样），可以假设 $m_i(n) \geq n$。此时考虑逆系统
$$
\ldots \to
M_{4, m_4(m_3(m_2(4)))} \to
M_{3, m_3(m_2(3))} \to
M_{2, m_2(2)} \to
M_{1, 1}
$$
，其一般项为
$$
M_k = M_{k, m_k(m_{k - 1}(\ldots (m_2(k))\ldots))}
$$
。对 $\mathcal{C}$ 的任意对象 $N$，有
$$
\colim_i \colim_n \Mor_\mathcal{C}(M_{i, n}, N) =
\colim_k \Mor_\mathcal{C}(M_k, N) 
$$
。略去细节。换言之，逆系统 $(M_k)$ 具有性质
$$
\colim_i \Mor_{\text{Pro-}\mathcal{C}}((M_{i, n}), N) =
\Mor_{\text{Pro-}\mathcal{C}}((M_k), N)
$$
。由《范畴》，注 \ref{categories-remark-pro-category-copresheaves} 中的讨论，
该性质在 pro-同构意义下确定逆系统 $(M_k)$。这样，便可把
$\text{Pro-}\mathcal{C}$ 中的某些逆系统化为具有可数指标范畴的 pro-对象。
\end{remark}

\begin{remark}
\label{duality-remark-composition-lower-shriek}
在情形 \ref{duality-situation-shriek} 中，设 $f : X \to Y$ 和 $g : Y \to Z$
是 $\textit{FTS}_S$ 的可复合态射。定义复合函子
$$
Rg_! \circ Rf_! :
D^b_{\textit{Coh}}(\mathcal{O}_X)
\longrightarrow
\text{Pro-}D^b_{\textit{Coh}}(\mathcal{O}_Z)
$$
。具体地，由 $Rf_!$ 的构造，对
$D^b_{\textit{Coh}}(\mathcal{O}_X)$ 中的 $K$，输出 $Rf_!K$ 是
$D^b_{\textit{Coh}}(\mathcal{O}_Y)$ 中某个逆系统 $(M_n)$ 的
pro-同构类。由于 $Rg_!$ 以相同方式构造，
$$
\ldots \to Rg_!M_3 \to Rg_!M_2 \to Rg_!M_1
$$
是 $\text{Pro-}D^b_{\textit{Coh}}(\mathcal{O}_Y)$ 的逆系统。
由注 \ref{duality-remark-compose-inverse-systems} 中的讨论，存在唯一的
$D^b_{\textit{Coh}}(\mathcal{O}_Z)$ 中逆系统的 pro-同构类，记为
$Rg_! Rf_! K$，使得
$$
\Hom_{\text{Pro-}D^b_{\textit{Coh}}(\mathcal{O}_Z)}(Rg_!Rf_!K, L) =
\colim_n \Hom_{\text{Pro-}D^b_{\textit{Coh}}(\mathcal{O}_Z)}(Rg_!M_n, L)
$$
。略去验证该构造对 $K$ 具有函子性所需的讨论，因为这会直接由下一个引理得到。
\end{remark}

\begin{lemma}
\label{duality-lemma-composition-lower-shriek}
在情形 \ref{duality-situation-shriek} 中，设 $f : X \to Y$ 和 $g : Y \to Z$
是 $\textit{FTS}_S$ 的可复合态射。沿用注
\ref{duality-remark-composition-lower-shriek} 的记号，有
$Rg_! \circ Rf_! = R(g \circ f)_!$。
\end{lemma}

\begin{proof}
由《范畴》，注 \ref{categories-remark-pro-category-copresheaves} 中的讨论，
只须证明对 $D^b_{\textit{Coh}}(\mathcal{O}_Z)$ 中的 $L$ 计算 $\Hom$
时得到相同答案。沿用注 \ref{duality-remark-composition-lower-shriek} 中的记号，计算如下：
\begin{align*}
\Hom_Z(Rg_!Rf_!K, L)
& =
\colim_n \Hom_Z(Rg_!M_n, L) \\
& =
\colim_n \Hom_Y(M_n, g^!L) \\
& =
\Hom_Y(Rf_!K, g^!L) \\
& =
\Hom_X(K, f^!g^!L) \\
& =
\Hom_X(K, (g \circ f)^!L) \\
& =
\Hom_Z(R(g \circ f)_!K, L)
\end{align*}
。第一个等式是 $Rg_!Rf_!K$ 的定义。第二个等式是命题
\ref{duality-proposition-duality-compactly-supported} 对 $g$ 的应用。
第三个等式使用 $Rf_!K$ 由 $(M_n)$ 给出这一事实。第四个等式是命题
\ref{duality-proposition-duality-compactly-supported} 对 $f$ 的应用。
第五个等式是引理 \ref{duality-lemma-upper-shriek-composition}。第六个等式是命题
\ref{duality-proposition-duality-compactly-supported} 对 $g \circ f$ 的应用。
\end{proof}

\begin{remark}
\label{duality-remark-covariance-open-lower-shriek}
在情形 \ref{duality-situation-shriek} 中，设 $f : X \to Y$ 是
$\textit{FTS}_S$ 的态射，且 $U \subset X$ 是开集。令
$g = f|_U : U \to Y$。则存在典范态射
$$
Rg_!(K|_U) \longrightarrow Rf_!K
$$
，它对 $D^b_{\textit{Coh}}(\mathcal{O}_X)$ 中的 $K$ 具有函子性，
并且至少可以用三种方式定义。
\begin{enumerate}
\item 以 $i : U \to X$ 表示包含态射。由引理
\ref{duality-lemma-composition-lower-shriek}，有 $Rg_! = Rf_! \circ Ri_!$。
把 $Rf_!$ 应用于映射 $Ri_!(K|_U) \to K$ 即可；这是注
\ref{duality-remark-covariance-open-j-lower-shriek} 的特殊情形。
\item 选取 $X$ 在 $Y$ 上的紧化 $j : X \to \overline{X}$，其结构态射为
$\overline{f} : \overline{X} \to Y$。令
$j' = j \circ i : U \to \overline{X}$。利用
$Rf_! = R\overline{f}_* \circ Rj_!$ 和
$Rg_! = R\overline{f}_* \circ Rj'_!$，并把 $R\overline{f}_*$ 应用于注
\ref{duality-remark-covariance-open-j-lower-shriek} 中的映射
$Rj'_!(K|_U) \to Rj_!K$。
\item 使用
\begin{align*}
\Hom_{\text{Pro-}D^b_{\textit{Coh}}(\mathcal{O}_Y)}(Rf_!K, L)
& =
\Hom_X(K, f^!L) \\
& \to
\Hom_U(K|_U, f^!L|_U) \\
& =
\Hom_U(K|_U, g^!L) \\
& =
\Hom_{\text{Pro-}D^b_{\textit{Coh}}(\mathcal{O}_Y)}(Rg_!(K|_U), L)
\end{align*}
。它对 $L$ 和 $K$ 具有函子性。这里两次使用了命题
\ref{duality-proposition-duality-compactly-supported}，还使用了上叹号函子的构造所给出的
$g^! = i^* \circ f^!$。由对 $L$ 的函子性及《范畴》，注
\ref{categories-remark-pro-category-copresheaves}，得到
$\text{Pro-}D^b_{\textit{Coh}}(\mathcal{O}_Y)$ 中的典范映射
$Rg_!(K|_U) \to Rf_!K$；上述箭头对 $K$ 的函子性说明该映射也对 $K$
具有函子性。
\end{enumerate}
这三种构造给出同一个箭头；略去细节。
\end{remark}

\begin{remark}
\label{duality-remark-covariance-etale-lower-shriek}
把注 \ref{duality-remark-covariance-open-lower-shriek} 中紧支撑上同调的协变性推广到
\'etale 态射。具体地，在情形 \ref{duality-situation-shriek} 中，假设给定
$\textit{FTS}_S$ 中的交换图
$$
\xymatrix{
U \ar[rr]_h \ar[rd]_g & & X \ar[ld]^f \\
& Y
}
$$
，其中 $h$ 为 \'etale。则存在典范态射
$$
Rg_!(h^*K) \longrightarrow Rf_!K
$$
，它对 $D^b_{\textit{Coh}}(\mathcal{O}_X)$ 中的 $K$ 具有函子性。
用下列映射链定义该变换：
\begin{align*}
\Hom_{\text{Pro-}D^b_{\textit{Coh}}(\mathcal{O}_Y)}(Rf_!K, L)
& =
\Hom_X(K, f^!L) \\
& \to
\Hom_U(h^*K, h^*(f^!L)) \\
& =
\Hom_U(h^*K, h^!f^!L) \\
& =
\Hom_U(h^*K, g^!L) \\
& =
\Hom_{\text{Pro-}D^b_{\textit{Coh}}(\mathcal{O}_Y)}(Rg_!(h^*K), L)
\end{align*}
。它对 $L$ 和 $K$ 具有函子性。这里两次使用了命题
\ref{duality-proposition-duality-compactly-supported}，使用了引理
\ref{duality-lemma-shriek-etale} 的等式 $h^* = h^!$，还使用了引理
\ref{duality-lemma-upper-shriek-composition} 的等式 $h^! \circ f^! = g^!$。
由对 $L$ 的函子性及《范畴》，注
\ref{categories-remark-pro-category-copresheaves}，得到
$\text{Pro-}D^b_{\textit{Coh}}(\mathcal{O}_Y)$ 中的典范映射
$Rg_!(h^*K) \to Rf_!K$；上述箭头对 $K$ 的函子性说明该映射也对 $K$
具有函子性。
\end{remark}

\begin{remark}
\label{duality-remark-covariance-lower-shriek}
在注 \ref{duality-remark-covariance-open-lower-shriek} 和
\ref{duality-remark-covariance-etale-lower-shriek} 中已经看到，紧支撑上同调的构造
对开浸入和 \'etale 态射具有协变性。事实上，正确的一般性如下：给定
$\textit{FTS}_S$ 中的交换图
$$
\xymatrix{
U \ar[rr]_h \ar[rd]_g & & X \ar[ld]^f \\
& Y
}
$$
，其中 $h$ 平坦且拟有限，则存在典范变换
$$
Rg_! \circ h^* \longrightarrow Rf_!
$$
。像注 \ref{duality-remark-covariance-etale-lower-shriek} 中一样，可以使用
$D^+_{\textit{Coh}}(\mathcal{O}_X)$ 上的函子变换 $h^* \to h^!$
来构造该映射。回忆
$h^!K = h^*K \otimes \omega_{U/X}$，其中
$\omega_{U/X} = h^!\mathcal{O}_X$ 是平坦拟有限态射 $h$ 的相对对偶层
（见引理 \ref{duality-lemma-perfect-comparison-shriek} 和
\ref{duality-lemma-flat-quasi-finite-shriek}）。还记得，$\omega_{U/X}$
等于《判别式》，注
\ref{discriminant-remark-relative-dualizing-for-quasi-finite}
中由《判别式》，引理 \ref{discriminant-lemma-compare-dualizing}
构造的相对对偶模。因此，可以使用将在《判别式》，注
\ref{discriminant-remark-relative-dualizing-for-flat-quasi-finite}
中构造的迹元素 $\tau_{U/X} : \mathcal{O}_U \to \omega_{U/X}$
来定义该变换。如果以后需要，我们将在这里精确陈述并证明该结果。
\end{remark}







\section{紧支撑上同调的对偶性}
\label{duality-section-duality-compactly-supported}

\noindent
设 $k$ 是域，$U$ 是 $k$ 上有限型的分离概形，$K$ 是
$D^b_{\textit{Coh}}(\mathcal{O}_U)$ 中的对象。如下定义 $K$
的紧支撑上同调 $H^i_c(U, K)$。选取到 $k$ 上某个固有概形中的开浸入
$j : U \to X$，以及 $j : U \to X$ 的一个 Deligne 系统 $(K_n)$，
使其限制到 $U$ 后常值为 $K$。令
$$
H^i_c(U, K) = \lim H^i(X, K_n)
$$
。赋予它逆极限拓扑，并把它视为拓扑 $k$-向量空间（见《代数进阶》，第
\ref{more-algebra-section-topological-ring} 节）。这里有几点需要说明。

\medskip\noindent
第一，该定义与 $X$ 和 $(K_n)$ 的选择无关。事实上，若
$p : U \to \Spec(k)$ 表示结构态射，则已经知道
$Rp_!K = (R\Gamma(X, K_n))$ 在 $D(k)$ 中于 pro-同构意义下良定义，
所以定义 $H^i_c(U, K)$ 的逆极限也良定义。

\medskip\noindent
第二，使用表达式
$$
H^i(R\lim R\Gamma(X, K_n)) = R\Gamma(X, R\lim K_n)
$$
看起来可能更自然，但所得答案相同：由于 $k$-向量空间
$H^j(X, K_n)$ 都有限维，这些逆系统满足 Mittag--Leffler 条件，因而《上同调》，
引理 \ref{cohomology-lemma-RGamma-commutes-with-Rlim} 中的
$R^1\lim$ 项消失。

\medskip\noindent
若 $U' \subset U$ 是开子概形，则存在典范映射
$$
H^i_c(U', K|_{U'}) \longrightarrow H^i_c(U, K)
$$
，它对 $D^b_{\textit{Coh}}(\mathcal{O}_U)$ 中的 $K$ 具有函子性。
例如见注 \ref{duality-remark-covariance-open-lower-shriek}。事实上，由注
\ref{duality-remark-covariance-etale-lower-shriek}，更一般地，对 $k$
上有限型分离概形之间的 \'etale 态射 $U' \to U$，也存在这样的映射。

\medskip\noindent
若 $V$ 是 $k$-向量空间，如下在 $\Hom_k(V, k)$ 上赋予拓扑：把
$V = \bigcup V_i$ 写成其有限维 $k$-子向量空间的滤过并，并在
$\Hom_k(V, k) = \lim \Hom_k(V_i, k)$ 上使用逆极限拓扑。若
$\dim_k V < \infty$，则 $\Hom_k(V, k)$ 上的拓扑是离散的。更一般地，
若 $V = \colim_n V_n$ 写成有限维向量空间的有向余极限，则作为拓扑向量空间，
$\Hom_k(V, k) = \lim \Hom_k(V_n, k)$。

\begin{lemma}
\label{duality-lemma-duality-compact-support}
设 $p : U \to \Spec(k)$ 是有限型分离态射，其中 $k$ 是域。令
$\omega_{U/k}^\bullet = p^!\mathcal{O}_{\Spec(k)}$。则有拓扑
$k$-向量空间的典范同构
$$
\Hom_k(H^i(U, K), k) =
H^{-i}_c(U, R\SheafHom_{\mathcal{O}_U}(K, \omega_{U/k}^\bullet))
$$
，它对 $D^b_{\textit{Coh}}(\mathcal{O}_U)$ 中的 $K$ 具有函子性。
\end{lemma}

\begin{proof}
选取 $k$ 上的紧化 $j : U \to X$。设拟相干理想层
$\mathcal{I} \subset \mathcal{O}_X$ 满足 $V(\mathcal{I}) = X \setminus U$。
由《概形的导出范畴》，命题 \ref{perfect-proposition-DCoh}，可选取
$M \in D^b_{\textit{Coh}}(\mathcal{O}_X)$，使 $K = M|_U$。由引理
\ref{duality-lemma-lift-map}，有
$$
H^i(U, K) =
\Ext^i_U(\mathcal{O}_U, M|_U) =
\colim \Ext^i_X(\mathcal{I}^n, M) =
\colim H^i(X, R\SheafHom_{\mathcal{O}_X}(\mathcal{I}^n, M))
$$
。由于 $\mathcal{I}^n$ 是相干 $\mathcal{O}_X$-模，故 $\mathcal{I}^n$
属于 $D^-_{\textit{Coh}}(\mathcal{O}_X)$；因而由《概形的导出范畴》，引理
\ref{perfect-lemma-coherent-internal-hom}，
$R\SheafHom_{\mathcal{O}_X}(\mathcal{I}^n, M)$ 属于
$D^+_{\textit{Coh}}(\mathcal{O}_X)$。

\medskip\noindent
令 $\omega_{X/k}^\bullet = q^!\mathcal{O}_{\Spec(k)}$，其中
$q : X \to \Spec(k)$ 是结构态射；见第
\ref{duality-section-duality-proper-over-field} 节。于是
\begin{align*}
\Hom_k(
&
H^i(X, R\SheafHom_{\mathcal{O}_X}(\mathcal{I}^n, M)), k) \\
& =
\Ext^{-i}_X(R\SheafHom_{\mathcal{O}_X}(\mathcal{I}^n, M),
\omega_{X/k}^\bullet) \\
& =
H^{-i}(X, R\SheafHom_{\mathcal{O}_X}(R\SheafHom_{\mathcal{O}_X}(
\mathcal{I}^n, M), \omega_{X/k}^\bullet))
\end{align*}
；见引理 \ref{duality-lemma-duality-proper-over-field}。由引理
\ref{duality-lemma-internal-hom-evaluate-isom} 的 (1)，典范映射
$$
R\SheafHom_{\mathcal{O}_X}(M, \omega_{X/k}^\bullet)
\otimes_{\mathcal{O}_X}^\mathbf{L} \mathcal{I}^n
\longrightarrow
R\SheafHom_{\mathcal{O}_X}(R\SheafHom_{\mathcal{O}_X}(
\mathcal{I}^n, M), \omega_{X/k}^\bullet)
$$
是同构。注意，$\omega^\bullet_{U/k} = \omega^\bullet_{X/k}|_U$，
因为 $p^!$ 构造为 $q^!$ 后接限制到 $U$。所以
$R\SheafHom_{\mathcal{O}_X}(M, \omega_{X/k}^\bullet)$ 是
$D^b_{\textit{Coh}}(\mathcal{O}_X)$ 中的对象，且限制到 $U$ 后为
$R\SheafHom_{\mathcal{O}_U}(K, \omega_{U/k}^\bullet)$。由引理
\ref{duality-lemma-tensoring-Deligne-system}，可知
$$
\lim H^{-i}(X, R\SheafHom_{\mathcal{O}_X}(M, \omega_{X/k}^\bullet)
\otimes_{\mathcal{O}_X}^\mathbf{L} \mathcal{I}^n)
$$
是该引理等式右端的一种表示。合并以上所得全部同构，即得到该引理的同构。
\end{proof}

\begin{lemma}
\label{duality-lemma-duality-compact-support-restrict-open}
沿用引理 \ref{duality-lemma-duality-compact-support} 的记号，假设
$U' \subset U$ 是开子概形。则图
$$
\xymatrix{
\Hom_k(H^i(U, K), k) \ar[rr] & &
H^{-i}_c(U, R\SheafHom_{\mathcal{O}_U}(K, \omega_{U/k}^\bullet)) \\
\Hom_k(H^i(U', K|_{U'}), k) \ar[rr] \ar[u] & &
H^{-i}_c(U', R\SheafHom_{\mathcal{O}_{U'}}(K, \omega_{U'/k}^\bullet)) \ar[u]
}
$$
交换。这里，水平箭头是引理 \ref{duality-lemma-duality-compact-support} 的同构；
左侧竖直箭头是限制映射 $H^i(U, K) \to H^i(U', K|_{U'})$
的对偶映射；右侧竖直箭头来自注 \ref{duality-remark-covariance-open-lower-shriek}
（见该引理前面的讨论）。
\end{lemma}

\begin{proof}
强烈建议读者跳过这个证明。像引理 \ref{duality-lemma-duality-compact-support}
的证明那样选取 $X$ 和 $M$。以下从 $R\SheafHom$ 和
$\otimes^\mathbf{L}$ 中省略下标 $\mathcal{O}_X$。写成
$$
H^i(U, K) = \colim H^i(X, R\SheafHom(\mathcal{I}^n, M))
$$
以及
$$
H^i(U', K|_{U'}) = \colim H^i(X, R\SheafHom((\mathcal{I}')^n, M))
$$
，如引理 \ref{duality-lemma-duality-compact-support} 的证明中那样；这里按照注
\ref{duality-remark-covariance-open-j-lower-shriek} 中的讨论选取
$\mathcal{I}' \subset \mathcal{I}$，使映射
$H^i(U, K) \to H^i(U', K|_{U'})$ 由映射
$(\mathcal{I}')^n \to \mathcal{I}^n$ 诱导。类似地写成
$$
H^i_c(U, R\SheafHom(K, \omega_{U/k}^\bullet)) =
\lim H^i(X,  R\SheafHom(M, \omega_{X/k}^\bullet)
\otimes^\mathbf{L} \mathcal{I}^n)
$$
以及
$$
H^i_c(U', R\SheafHom(K|_{U'}, \omega_{U'/k}^\bullet)) =
\lim H^i(X,  R\SheafHom(M, \omega_{X/k}^\bullet)
\otimes^\mathbf{L} (\mathcal{I}')^n)
$$
，从而箭头 $H^i_c(U', R\SheafHom(K|_{U'}, \omega_{U'/k}^\bullet))
\to H^i_c(U, R\SheafHom(K, \omega_{U/k}^\bullet))$ 也由映射
$(\mathcal{I}')^n \to \mathcal{I}^n$ 导出。图
$$
\xymatrix{
R\SheafHom(M, \omega_{X/k}^\bullet)
\otimes^\mathbf{L} \mathcal{I}^n
\ar[rr] & &
R\SheafHom(R\SheafHom(\mathcal{I}^n, M), \omega_{X/k}^\bullet) \\
R\SheafHom(M, \omega_{X/k}^\bullet) \otimes^\mathbf{L} (\mathcal{I}')^n
\ar[rr] \ar[u] & &
R\SheafHom(R\SheafHom((\mathcal{I}')^n, M), \omega_{X/k}^\bullet) \ar[u]
}
$$
交换，因为《上同调》，引理
\ref{cohomology-lemma-internal-hom-evaluate} 中水平箭头的构造对三个变元
都具有函子性。因此最终归结为断言图
$$
\xymatrix{
\Hom_k(H^i(X, R\SheafHom(\mathcal{I}^n, M)), k) \ar[r] &
H^{-i}(X, R\SheafHom(R\SheafHom(
\mathcal{I}^n, M), \omega_{X/k}^\bullet)) \\
\Hom_k(H^i(X, R\SheafHom((\mathcal{I}')^n, M)), k) \ar[r] \ar[u] &
H^{-i}(X, R\SheafHom(R\SheafHom(
(\mathcal{I}')^n, M), \omega_{X/k}^\bullet)) \ar[u]
}
$$
交换。这成立是因为对偶同构
$$
\Hom_k(H^i(X, L), k) = \Ext^{-i}_X(L, \omega_{X/k}^\bullet) =
H^{-i}(X, R\SheafHom(L, \omega_{X/k}^\bullet))
$$
对 $D_\QCoh(\mathcal{O}_X)$ 中的 $L$ 具有函子性。
\end{proof}

\begin{lemma}
\label{duality-lemma-h0-compactly-supported}
设 $X$ 是域 $k$ 上的完备概形。设
$K \in D^b_{\textit{Coh}}(\mathcal{O}_X)$ 满足对 $i < 0$ 有
$H^i(K) = 0$。令 $\mathcal{F} = H^0(K)$。设 $Z \subset X$
是闭子集，其补集为 $U = X \setminus U$。则
$$
H^0_c(U, K|_U) \subset H^0(X, \mathcal{F})
$$
恰由 $\mathcal{F}$ 中在 $Z$ 的某个开邻域上消失的整体截面组成。
\end{lemma}

\begin{proof}
考虑注 \ref{duality-remark-covariance-open-lower-shriek} 中的映射
$H^0_c(U, K|_U) \to H^0_X(X, K) = H^0(X, K) = H^0(X, \mathcal{F})$。
为研究此映射，用一个有界复形 $\mathcal{F}^\bullet$ 表示 $K$，其中
对 $i < 0$ 有 $\mathcal{F}^i = 0$。根据定义，
$$
H^0_c(U, K|_U) = \lim H^0(X, \mathcal{I}^n\mathcal{F}^\bullet)
= \lim \Ker(
H^0(X, \mathcal{I}^n\mathcal{F}^0) \to
H^0(X, \mathcal{I}^n\mathcal{F}^1))
$$
由 Artin--Rees 引理（《概形上同调》，引理
\ref{coherent-lemma-Artin-Rees}），这等于
$\lim H^0(X, \mathcal{I}^n\mathcal{F})$。因此箭头
$H^0_c(U, K|_U) \to H^0(X, \mathcal{F})$ 是单射，其像由
$\mathcal{F}$ 中对每个 $n$ 都属于子层 $\mathcal{I}^n\mathcal{F}$
的整体截面组成。考察 $\mathcal{F}$ 的茎并应用 Krull 交定理
（《代数》，引理 \ref{algebra-lemma-intersect-powers-ideal-module-zero}），
即可把这些截面刻画为在 $Z$ 的某个邻域上消失的截面。函数情形见
《代数》，注 \ref{algebra-remark-intersection-powers-ideal} 中的讨论。
\end{proof}




\section{Lichtenbaum 定理}
\label{duality-section-lichtenbaum}

\noindent
以下定理由 Lichtenbaum 猜想，并由 Grothendieck 证明
（见 \cite{Hartshorne-local-cohomology}）。Kleiman 在
\cite{Kleiman-Lichtenbaum} 中给出了一个非常漂亮的证明。
该定理对带支集上同调情形的推广见
\cite{Lyubeznik-Lichtenbaum}。论证中最有趣的部分包含在下述引理的证明中。

\begin{lemma}
\label{duality-lemma-lichtenbaum}
设 $U$ 是代数簇，$\mathcal{F}$ 是凝聚 $\mathcal{O}_U$-模。
若 $H^d(U, \mathcal{F})$ 非零，则 $\dim(U) \geq d$；若等号成立，
则 $U$ 是完备的。
\end{lemma}

\begin{proof}
由《上同调》，命题 \ref{cohomology-proposition-vanishing-Noetherian}
中的 Grothendieck 消失结果，得到 $\dim(U) \geq d$。
假设 $\dim(U) = d$。选取一个使 $U$ 在 $X$ 中稠密的紧化
$U \to X$。（由《平坦性的进一步结果》，定理
\ref{flat-theorem-nagata} 和引理 \ref{flat-lemma-compactifyable}，
这样的紧化存在。）以其约化替换 $X$ 后，可知 $X$ 是维数为
$d$ 的完备代数簇，并且 $U$ 完备当且仅当 $U = X$。令
$Z = X \setminus U$。下面证明：若 $Z$ 非空，则
$H^d(U, \mathcal{F})$ 为零。

\medskip\noindent
选取一个凝聚 $\mathcal{O}_X$-模 $\mathcal{G}$，使其在 $U$ 上的限制为
$\mathcal{F}$；见《性质》，引理
\ref{properties-lemma-lift-finite-presentation}。令
$\omega_X^\bullet$ 表示第 \ref{duality-section-duality-proper-over-field} 节所述
$X$ 的对偶化复形，并令 $\omega_U^\bullet = \omega_X^\bullet|_U$。
由引理 \ref{duality-lemma-duality-compact-support}，$H^d(U, \mathcal{F})$
与
$$
H^{-d}_c(U, R\SheafHom_{\mathcal{O}_U}(\mathcal{F}, \omega_U^\bullet))
$$
互为对偶。由引理 \ref{duality-lemma-duality-proper-over-field}，
$\omega_X^\bullet$ 的上同调层在次数 $< -d$ 时消失，并且
$H^{-d}(\omega_X^\bullet) = \omega_X$ 是一个凝聚
$\mathcal{O}_X$-模，它满足 $(S_2)$，且支集为 $X$。特别地，
$\omega_X$ 是无挠的；见《除子》，引理
\ref{divisors-lemma-torsion-free-finite-noetherian-domain}。因此上同调层
$$
H^{-d}(R\SheafHom_{\mathcal{O}_X}(\mathcal{G}, \omega_X^\bullet)) =
\SheafHom(\mathcal{G}, \omega_X)
$$
是无挠的；见《除子》，引理
\ref{divisors-lemma-hom-into-torsion-free}。所以该层不存在在 $X$
的任何非空开集上消失的非零截面（这样的截面将是挠截面）。于是由引理
\ref{duality-lemma-h0-compactly-supported}，
$H^{-d}_c(U, R\SheafHom_{\mathcal{O}_U}(\mathcal{F}, \omega_U^\bullet))$
为零，因而如所欲，$H^d(U, \mathcal{F})$ 为零。
\end{proof}

\begin{theorem}
\label{duality-theorem-lichtenbaum}
设 $X$ 是域 $k$ 上的非空有限型分离概形。令 $d = \dim(X)$。
则下列条件等价：
\begin{enumerate}
\item 对 $X$ 上每个凝聚 $\mathcal{O}_X$-模 $\mathcal{F}$，都有
$H^d(X, \mathcal{F}) = 0$；
\item 对 $X$ 上每个拟凝聚 $\mathcal{O}_X$-模 $\mathcal{F}$，都有
$H^d(X, \mathcal{F}) = 0$；
\item 不存在维数为 $d$ 且在 $k$ 上完备的不可约分支
$X' \subset X$。
\end{enumerate}
\end{theorem}

\begin{proof}
假设存在不可约分支 $X' \subset X$，它是完备的且维数为 $d$；
这里把它视为整闭子概形。令 $\omega_{X'}$ 为 $X'$ 在 $k$ 上的一个
对偶化模；见引理 \ref{duality-lemma-duality-proper-over-field}。由该引理，
$H^d(X', \omega_{X'})$ 与 $H^0(X', \mathcal{O}_{X'})$ 互为对偶，
故前者非零。因此 $H^d(X, \omega_{X'}) = H^d(X', \omega_{X'})$
非零，从而条件 (1) 不成立。由此可见 (1) 蕴含 (3)。

\medskip\noindent
下面证明 (3) 蕴含 (1)。设 $\mathcal{F}$ 是凝聚
$\mathcal{O}_X$-模，且 $H^d(X, \mathcal{F})$ 非零。选取《概形上同调》，
引理 \ref{coherent-lemma-coherent-filter} 中的滤过
$$
0 = \mathcal{F}_0 \subset \mathcal{F}_1 \subset
\ldots \subset \mathcal{F}_m = \mathcal{F}
$$
得到正合序列
$$
H^d(X, \mathcal{F}_i) \to H^d(X, \mathcal{F}_{i + 1}) \to
H^d(X, \mathcal{F}_{i + 1}/\mathcal{F}_i)
$$
于是对某个 $i \in \{1, \ldots, m\}$，
$H^d(X, \mathcal{F}_{i + 1}/\mathcal{F}_i)$ 非零。由滤过的选取方式，
这意味着存在一个整闭子概形 $Z \subset X$ 以及一个非零凝聚理想层
$\mathcal{I} \subset \mathcal{O}_Z$，使 $H^d(Z, \mathcal{I})$ 非零。
由引理 \ref{duality-lemma-lichtenbaum}，得到 $\dim(Z) = d$ 且 $Z$ 在 $k$
上完备，这与 (3) 矛盾。因此 (3) 蕴含 (1)。

\medskip\noindent
最后证明：对任意诺特概形 $X$，(1) 与 (2) 等价。显然 (2) 蕴含
(1)。反过来，假设 (1) 成立，并设 $\mathcal{F}$ 是拟凝聚
$\mathcal{O}_X$-模。可将 $\mathcal{F} = \colim \mathcal{F}_i$
写成其凝聚子模的滤过余极限；见《性质》，引理
\ref{properties-lemma-quasi-coherent-colimit-finite-type}。于是由《上同调》，
引理 \ref{cohomology-lemma-quasi-separated-cohomology-colimit}，有
$H^d(X, \mathcal{F}) = \colim H^d(X, \mathcal{F}_i) = 0$。
所以 (2) 成立。
\end{proof}
% Generated mechanically from the frozen discriminant.tex target.
% Do not edit this generated file; edit the bound locale source instead.

% OK, start here.
%

\setcounter{chapter}{48}
\chapter{判别式与不同}



\phantomsection
\label{discriminant-section-phantom}


\section{引言}
\label{discriminant-section-introduction}

\noindent
本章研究概形的局部拟有限态射的不同理想与判别式。
其中部分内容可参考 \cite{Kunz}。

\medskip\noindent
给定 Noetherian 概形的拟有限态射 $f : Y \to X$，
存在相对对偶化模 $\omega_{Y/X}$。
第 \ref{discriminant-section-quasi-finite-dualizing} 节利用
Zariski 主定理和 \'etale 局部化方法从头构造这个模。
关键性质如下：给定交换图
$$
\xymatrix{
Y' \ar[d]_{f'} \ar[r]_{g'} & Y \ar[d]^f \\
X' \ar[r]^g & X
}
$$
其中 $g : X' \to X$ 平坦，$Y' \subset X' \times_X Y$ 为开子集，且
$f' : Y' \to X'$ 有限，则存在典范同构
$$
f'_*(g')^*\omega_{Y/X} =
\SheafHom_{\mathcal{O}_{X'}}(f'_*\mathcal{O}_{Y'}, \mathcal{O}_{X'})
$$
作为 $f'_*\mathcal{O}_{Y'}$-模层。第
\ref{discriminant-section-quasi-finite-traces} 节证明：若 $f$ 平坦，则存在典范整体截面
$\tau_{Y/X} \in H^0(Y, \omega_{Y/X})$，使得对每个如上的交换图，
$(g')^*\tau_{Y/X}$ 映到有限局部自由态射 $f'$ 在第
\ref{discriminant-section-discriminant} 节中的迹映射。
第 \ref{discriminant-section-different} 节把 Noetherian 概形的平坦拟有限态射之
不同理想定义为
$\tau_{Y/X} : \mathcal{O}_X \to \omega_{Y/X}$ 的余核的零化子。

\medskip\noindent
本章的主要目标是证明：对拟有限 syntomic
态射\footnote{亦即平坦且为 lci。}$f$，其不同理想与 K\"ahler 不同理想一致。
K\"ahler 不同理想是 $\Omega_{Y/X}$ 的第零 Fitting 理想；见第
\ref{discriminant-section-kahler-different} 节。
这一相等并非显然；第 \ref{discriminant-section-formula-different} 节采用 Tate
提出的一个精巧论证。沿途还将讨论 Noether 不同理想与 Dedekind 不同理想。

\medskip\noindent
直到本章末尾的第 \ref{discriminant-section-comparison} 和
\ref{discriminant-section-gorenstein-lci} 节，我们才把这些内容与概形对偶理论中
较为深入的材料联系起来。








\section{拟有限环同态的对偶化模}
\label{discriminant-section-quasi-finite-dualizing}

\noindent
设 $A \to B$ 是 Noetherian 环的拟有限同态。由 Zariski 主定理
（《代数》，引理 \ref{algebra-lemma-quasi-finite-open-integral-closure}），
存在分解 $A \to B' \to B$，其中 $A \to B'$ 有限，而 $B' \to B$
在谱上诱导开浸入。此时令
\begin{equation}
\label{discriminant-equation-dualizing}
\omega_{B/A} = \Hom_A(B', A) \otimes_{B'} B
\end{equation}
可将它看成一种相对对偶化模；见引理 \ref{discriminant-lemma-compare-dualizing} 和
\ref{discriminant-lemma-compare-dualizing-algebraic}。
本节将用初等交换代数方法证明：$\omega_{B/A}$ 不依赖于分解的选取，
并且 $\omega_{B/A}$ 的形成与平坦基变换可交换。
为证明分解无关性，先比较任意给定的两个分解。

\begin{lemma}
\label{discriminant-lemma-dominate-factorizations}
设 $A \to B$ 是拟有限环同态。给定两个分解
$A \to B' \to B$ 与 $A \to B'' \to B$，其中
$A \to B'$ 和 $A \to B''$ 有限，而 $\Spec(B) \to \Spec(B')$
与 $\Spec(B) \to \Spec(B'')$ 均为开浸入。则存在
一个 $A$-子代数 $B''' \subset B$，它在 $A$ 上有限，并且
$\Spec(B) \to \Spec(B''')$ 是开浸入，且 $B' \to B$ 与
$B'' \to B$ 均经过 $B'''$.
\end{lemma}

\begin{proof}
令 $B''' \subset B$ 为 $B' \to B$ 与 $B'' \to B$ 的像所生成的
$A$-子代数。由于 $B'$ 和 $B''$ 都由有限多个在 $A$ 上整的元素生成，
可知 $B'''$ 也由有限多个在 $A$ 上整的元素生成，因而
$B'''$ 在 $A$ 上有限
（《代数》，引理 \ref{algebra-lemma-characterize-finite-in-terms-of-integral}）。
考虑映射
$$
B = B' \otimes_{B'} B \to B''' \otimes_{B'} B \to B \otimes_{B'} B = B
$$
末一个等号成立是因为 $\Spec(B) \to \Spec(B')$ 是开浸入，
因而是单态射。第二个箭头为单射，因为 $B' \to B$ 平坦。
因此两个箭头均为同构。这意味着
$$
\xymatrix{
\Spec(B''') \ar[d] & \Spec(B) \ar[d] \ar[l] \\
\Spec(B') & \Spec(B) \ar[l]
}
$$
是笛卡尔图。开浸入的基变换仍为开浸入，故得结论。
\end{proof}

\begin{lemma}
\label{discriminant-lemma-dualizing-well-defined}
模 (\ref{discriminant-equation-dualizing}) 是良定义的，即与分解的选取无关。
\end{lemma}

\begin{proof}
令 $B', B'', B'''$ 如引理 \ref{discriminant-lemma-dominate-factorizations} 所述。
得到典范映射
$$
\omega''' = \Hom_A(B''', A) \otimes_{B'''} B \longrightarrow
\Hom_A(B', A) \otimes_{B'} B = \omega'
$$
以及一个涉及 $B''$ 的类似映射。只要证明这些映射都是同构即可。
取 $g \in B'$，使 $B'_g \to B_g$ 是同构；于是 $B'_g \to (B''')_g \to B_g$ 都是同构。只需证明 $(\omega''')_g \to \omega'_g$ 是同构。环同态 $B' \to B'''$ 的核和余核都是有限 $A$-模，且为 $g$-幂挠模，
因而被 $g$ 的某个幂零化。由此容易推出结论。
\end{proof}

\begin{lemma}
\label{discriminant-lemma-localize-dualizing}
设 $A \to B$ 是 Noetherian 环的拟有限同态。
\begin{enumerate}
\item 若对某个 $f \in A$，$A \to B$ 分解为 $A \to A_f \to B$，
则 $\omega_{B/A} = \omega_{B/A_f}$。
\item 若 $g \in B$，则 $(\omega_{B/A})_g = \omega_{B_g/A}$。
\item 若 $f \in A$，则 $\omega_{B_f/A_f} = (\omega_{B/A})_f$。
\end{enumerate}
\end{lemma}

\begin{proof}
设 $A \to B' \to B$ 是一个分解，其中 $A \to B'$ 有限，且
$\Spec(B) \to \Spec(B')$ 是开浸入。在情形 (1) 中，可以用分解
$A_f \to B'_f \to B$ 计算 $\omega_{B/A_f}$，再应用《代数》，引理
\ref{algebra-lemma-hom-from-finitely-presented}。
在情形 (2) 中，用分解 $A \to B' \to B_g$ 即得结论。
(3) 由 (1) 与 (2) 结合而得。
\end{proof}

\noindent
设 $A \to B$ 是 Noetherian 环的拟有限同态，$A \to A_1$ 是
Noetherian 环的任意同态，并令 $B_1 = B \otimes_A A_1$。得到余笛卡尔图
$$
\xymatrix{
B \ar[r] & B_1 \\
A \ar[u] \ar[r] & A_1 \ar[u]
}
$$
注意 $A_1 \to B_1$ 也是拟有限的（《代数》，引理
\ref{algebra-lemma-quasi-finite-base-change}）。
在此情形下定义一个典范的 $B$-线性基变换映射
\begin{equation}
\label{discriminant-equation-bc-dualizing}
\omega_{B/A} \longrightarrow \omega_{B_1/A_1}
\end{equation}
具体地，按照 $\omega_{B/A}$ 的构造选取分解 $A \to B' \to B$。
此时 $B'_1 = B' \otimes_A A_1$ 在 $A_1$ 上有限，而在构造
$\omega_{B_1/A_1}$ 时可采用分解 $A_1 \to B'_1 \to B_1$。
因此需要构造映射
$$
\Hom_A(B', A) \otimes_{B'} B
\longrightarrow
\Hom_{A_1}(B' \otimes_A A_1, A_1) \otimes_{B'_1} B_1
$$
只需构造一个 $B'$-线性映射
$\Hom_A(B', A) \to \Hom_{A_1}(B' \otimes_A A_1, A_1)$，
记为 $\varphi \mapsto \varphi_1$。给定 $A$-线性映射
$\varphi : B' \to A$，令 $\varphi_1$ 满足
$\varphi_1(b' \otimes a_1) = \varphi(b')a_1$。
这显然是 $A_1$-线性的，构造完成。

\begin{lemma}
\label{discriminant-lemma-bc-map-dualizing}
基变换映射 (\ref{discriminant-equation-bc-dualizing}) 与分解
$A \to B' \to B$ 的选取无关。给定环同态 $A \to A_1 \to A_2$，
$A \to A_1$ 与 $A_1 \to A_2$ 的基变换映射之复合，等于
$A \to A_2$ 的基变换映射。
\end{lemma}

\begin{proof}
略。提示：完全仿照引理 \ref{discriminant-lemma-dualizing-well-defined} 的论证，
并使用引理 \ref{discriminant-lemma-dominate-factorizations}。
\end{proof}

\begin{lemma}
\label{discriminant-lemma-dualizing-flat-base-change}
若 $A \to A_1$ 平坦，则基变换映射 (\ref{discriminant-equation-bc-dualizing}) 诱导同构
$\omega_{B/A} \otimes_B B_1 \to \omega_{B_1/A_1}$。
\end{lemma}

\begin{proof}
假设 $A \to A_1$ 平坦。由 $\omega_{B/A}$ 的构造，可设 $A \to B$ 有限。
此时 $\omega_{B/A} = \Hom_A(B, A)$，且
$\omega_{B_1/A_1} = \Hom_{A_1}(B_1, A_1)$。由于
$B_1 = B \otimes_A A_1$，结论由《代数进阶》，引理
\ref{more-algebra-lemma-pseudo-coherence-and-base-change-ext} 得出。
\end{proof}

\begin{lemma}
\label{discriminant-lemma-dualizing-composition}
设 $A \to B \to C$ 是 Noetherian 环的拟有限同态。则存在典范映射
$\omega_{B/A} \otimes_B \omega_{C/B} \to \omega_{C/A}$。
\end{lemma}

\begin{proof}
选取 $A \to B' \to B$，使 $A \to B'$ 有限，且
$\Spec(B) \to \Spec(B')$ 是开浸入。此时 $B' \to C$ 也拟有限。
选取 $B' \to C' \to C$，使 $B' \to C'$ 有限，且
$\Spec(C) \to \Spec(C')$ 是开浸入。于是该箭头的源为
$$
\Hom_A(B', A) \otimes_{B'} B \otimes_B
\Hom_B(B \otimes_{B'} C', B) \otimes_{B \otimes_{B'} C'} C
$$
它等于
$$
\Hom_A(B', A) \otimes_{B'}
\Hom_{B'}(C', B) \otimes_{C'} C
$$
而映射的复合
$\Hom_A(B', A) \times \Hom_{B'}(C', B) \to \Hom_A(C', A)$
确实给出它到
$\Hom_A(C', A) \otimes_{C'} C = \omega_{C/A}$ 的典范映射。
\end{proof}

\begin{lemma}
\label{discriminant-lemma-dualizing-product}
设 $A \to B$ 和 $A \to C$ 是 Noetherian 环的拟有限同态。
则作为 $B \times C$ 上的模，
$\omega_{B \times C/A} = \omega_{B/A} \times \omega_{C/A}$。
\end{lemma}

\begin{proof}
选取分解 $A \to B' \to B$ 和 $A \to C' \to C$，使得
$A \to B'$ 与 $A \to C'$ 有限，而 $\Spec(B) \to \Spec(B')$
与 $\Spec(C) \to \Spec(C')$ 是开浸入。则
$A \to B' \times C' \to B \times C$ 是同类分解。
用它计算 $\omega_{B \times C/A}$ 即得结论。
\end{proof}

\begin{lemma}
\label{discriminant-lemma-dualizing-associated-primes}
设 $A \to B$ 是 Noetherian 环的拟有限同态。则
$\text{Ass}_B(\omega_{B/A})$ 是 $B$ 中位于 $A$ 的伴随素理想之上的
所有素理想所成的集合。
\end{lemma}

\begin{proof}
选取分解 $A \to B' \to B$，其中 $A \to B'$ 有限，而 $B' \to B$
在谱上诱导开浸入。由于
$\omega_{B/A} = \omega_{B'/A} \otimes_{B'} B$，只需对
$\omega_{B'/A}$ 证明该断言。因此可设 $A \to B$ 有限。

\medskip\noindent
设 $\mathfrak p \in \text{Ass}(A)$，且 $\mathfrak q$ 是 $B$ 中位于
$\mathfrak p$ 之上的素理想。取 $x \in A$，使其零化子为
$\mathfrak p$。选取非零的 $\kappa(\mathfrak p)$-线性映射
$\lambda : \kappa(\mathfrak q) \to \kappa(\mathfrak p)$。
由于 $A/\mathfrak p \subset B/\mathfrak q$ 是有限环扩张，存在
$f \in A$、$f \not \in \mathfrak p$，使 $f\lambda$ 把
$B/\mathfrak q$ 映入 $A/\mathfrak p$。于是得到非零的 $A$-线性映射
$$
B \to B/\mathfrak q \to A/\mathfrak p \to A,\quad
b \mapsto f\lambda(b)x
$$
直接计算可知，$\omega_{B/A}$ 的这个元素的零化子是
$\mathfrak q$，故 $\mathfrak q \in \text{Ass}(\omega_{B/A})$。

\medskip\noindent
反之，设素理想 $\mathfrak q \subset B$ 位于素理想
$\mathfrak p \subset A$ 之上，且后者不是 $A$ 的伴随素理想。
需证明 $\mathfrak q \not \in \text{Ass}_B(\omega_{B/A})$。
把 $A$ 换成 $A_\mathfrak p$，并把 $B$ 换成 $B_\mathfrak p$，
可设 $\mathfrak p$ 是 $A$ 的极大理想。引理
\ref{discriminant-lemma-dualizing-flat-base-change} 与《代数》，引理
\ref{algebra-lemma-localize-ass} 说明这一步是允许的。
于是存在 $f \in \mathfrak m$，它是 $A$ 上的非零因子。
那么 $f$ 也是 $\omega_{B/A}$ 上的非零因子，故 $\mathfrak q$
不是这个模的伴随素理想。
\end{proof}

\begin{lemma}
\label{discriminant-lemma-dualizing-base-flat-flat}
设 $A \to B$ 是 Noetherian 环的平坦拟有限同态。则
$\omega_{B/A}$ 是平坦 $A$-模。
\end{lemma}

\begin{proof}
设素理想 $\mathfrak q \subset B$ 位于 $\mathfrak p \subset A$ 之上。
我们证明局部化 $\omega_{B/A, \mathfrak q}$ 在 $A_\mathfrak p$ 上平坦。
由《代数》，引理 \ref{algebra-lemma-flat-localization}，这就足够了。
由《代数》，引理
\ref{algebra-lemma-etale-makes-quasi-finite-finite-one-prime}，
可找到 \'etale 环同态 $A \to A'$ 和位于 $\mathfrak p$ 之上的素理想
$\mathfrak p' \subset A'$，使得 $\kappa(\mathfrak p') = \kappa(\mathfrak p)$，
并且
$$
B' = B \otimes_A A' = C \times D
$$
其中 $A' \to C$ 有限，且 $B \otimes_A A'$ 中同时位于
$\mathfrak q$ 与 $\mathfrak p'$ 之上的唯一素理想 $\mathfrak q'$
对应于 $C$ 的一个素理想。由引理
\ref{discriminant-lemma-dualizing-flat-base-change} 和《代数》，引理
\ref{algebra-lemma-base-change-flat-up-down}，只需证明
$\omega_{B'/A', \mathfrak q'}$ 在 $A'_{\mathfrak p'}$ 上平坦。
又由引理 \ref{discriminant-lemma-dualizing-product}，
$\omega_{B'/A'} = \omega_{C/A'} \times \omega_{D/A'}$，
故归结为 $B$ 在 $A$ 上有限平坦的情形。此时 $B$ 是有限局部自由
$A$-模，而 $\omega_{B/A} = \Hom_A(B, A)$ 是其对偶有限局部自由
$A$-模。
\end{proof}

\begin{lemma}
\label{discriminant-lemma-dualizing-base-change-of-flat}
若 $A \to B$ 平坦，则基变换映射 (\ref{discriminant-equation-bc-dualizing}) 诱导同构
$\omega_{B/A} \otimes_B B_1 \to \omega_{B_1/A_1}$。
\end{lemma}

\begin{proof}
若 $A \to B$ 有限平坦，则 $B$ 是有限局部自由 $A$-模。
此时 $\omega_{B/A} = \Hom_A(B, A)$ 是其对偶有限局部自由 $A$-模，
而这个模的形成与任意基变换可交换，故引理在此情形下成立。
下一段把一般的拟有限平坦情形归结为刚讨论的有限平坦情形。

\medskip\noindent
设 $\mathfrak q_1 \subset B_1$ 是素理想。我们证明该映射在
$\mathfrak q_1$ 处的局部化是同构；由《代数》，引理
\ref{algebra-lemma-characterize-zero-local}，这就足够了。
令 $\mathfrak q \subset B$ 与 $\mathfrak p \subset A$ 为
$\mathfrak q_1$ 下方的素理想。由《代数》，引理
\ref{algebra-lemma-etale-makes-quasi-finite-finite-one-prime}，
可找到 \'etale 环同态 $A \to A'$ 和位于 $\mathfrak p$ 之上的素理想
$\mathfrak p' \subset A'$，使得 $\kappa(\mathfrak p') = \kappa(\mathfrak p)$，
并且
$$
B' = B \otimes_A A' = C \times D
$$
其中 $A' \to C$ 有限，且 $B \otimes_A A'$ 中同时位于
$\mathfrak q$ 与 $\mathfrak p'$ 之上的唯一素理想 $\mathfrak q'$
对应于 $C$ 的一个素理想。令 $A'_1 = A' \otimes_A A_1$，并考察环同态
$A \to A' \to A'_1$ 与 $A \to A_1 \to A'_1$ 的基变换映射
(\ref{discriminant-equation-bc-dualizing})，如下图所示：
$$
\xymatrix{
\omega_{B'/A'} \otimes_{B'} B'_1 \ar[r] & \omega_{B'_1/A'_1} \\
\omega_{B/A} \otimes_B B'_1 \ar[r] \ar[u] &
\omega_{B_1/A_1} \otimes_{B_1} B'_1 \ar[u]
}
$$
其中 $B' = B \otimes_A A'$、$B_1 = B \otimes_A A_1$，且
$B_1' = B \otimes_A (A' \otimes_A A_1)$。由引理
\ref{discriminant-lemma-bc-map-dualizing}，该图交换。由引理
\ref{discriminant-lemma-dualizing-flat-base-change}，竖直箭头都是同构。
由于 $B_1 \to B'_1$ 是 \'etale 的，因而平坦，只需证明上方水平箭头
在 $B'_1$ 中位于 $\mathfrak q$ 之上的某个素理想 $\mathfrak q'_1$
处局部化后是同构（这样的素理想存在，并应用《代数》，引理
\ref{algebra-lemma-local-flat-ff}）。
因此可设 $B = C \times D$，其中 $A \to C$ 有限，且 $\mathfrak q$
对应于 $C$ 的一个素理想。此时对偶化模 $\omega_{B/A}$ 也有相应分解
（引理 \ref{discriminant-lemma-dualizing-product}），于是问题归结为上面已处理的
有限平坦情形 $A \to C$。
\end{proof}

\begin{remark}
\label{discriminant-remark-relative-dualizing-for-quasi-finite}
设 $f : Y \to X$ 是局部 Noetherian 概形之间的局部拟有限态射。
由引理 \ref{discriminant-lemma-localize-dualizing} 显然可知，$Y$ 上存在唯一的凝聚
$\mathcal{O}_Y$-模 $\omega_{Y/X}$，使得对每一对满足
$f(V) \subset U$ 的仿射开集 $\Spec(B) = V \subset Y$、
$\Spec(A) = U \subset X$，都存在典范同构
$$
H^0(V, \omega_{Y/X}) = \omega_{B/A}
$$
且这些同构与限制映射相容。
\end{remark}

\begin{lemma}
\label{discriminant-lemma-compare-dualizing-algebraic}
设 $A \to B$ 是 Noetherian 环的拟有限同态。令
$\omega_{B/A}^\bullet \in D(B)$ 为《对偶化复形》第
\ref{dualizing-section-relative-dualizing-complexes-Noetherian} 节讨论的
代数相对对偶化复形。则存在一个（非唯一的）同构
$\omega_{B/A} = H^0(\omega_{B/A}^\bullet)$。
\end{lemma}

\begin{proof}
选取分解 $A \to B' \to B$，其中 $A \to B'$ 有限，且
$\Spec(B') \to \Spec(B)$ 是开浸入。由《对偶化复形》，引理
\ref{dualizing-lemma-composition-shriek-algebraic}、
\ref{dualizing-lemma-upper-shriek-localize} 以及
$\omega_{B/A}^\bullet$ 的定义，有
$\omega_{B/A}^\bullet = \omega_{B'/A}^\bullet \otimes_B^\mathbf{L} B'$。
因此只需在 $A \to B$ 有限时证明存在该同构。在这种情形下，
《对偶化复形》，引理 \ref{dualizing-lemma-upper-shriek-finite} 给出
$\omega_{B/A}^\bullet = R\Hom(B, A)$，从而如所需有
$H^0(\omega^\bullet_{B/A}) = \Hom_A(B, A)$。
\end{proof}






\section{有限局部自由态射的判别式}
\label{discriminant-section-discriminant}

\noindent
设 $X$ 是概形，$\mathcal{F}$ 是有限局部自由 $\mathcal{O}_X$-模。
则存在典范的{\it 迹}映射
$$
\text{Trace} :
\SheafHom_{\mathcal{O}_X}(\mathcal{F}, \mathcal{F})
\longrightarrow
\mathcal{O}_X
$$
见《练习》，练习 \ref{exercises-exercise-trace-det}。这个映射满足：
$\text{Trace}(\text{id})$ 是 $\mathcal{O}_X$ 上对应于
$\mathcal{F}$ 的秩的局部常值函数。

\medskip\noindent
设 $\pi : X \to Y$ 是有限局部自由概形态射。则存在典范的
{\it $\pi$ 的迹}，即一个 $\mathcal{O}_Y$-线性映射
$$
\text{Trace}_\pi : \pi_*\mathcal{O}_X \longrightarrow \mathcal{O}_Y
$$
它把 $\pi_*\mathcal{O}_X$ 的局部截面 $f$ 映到
$\pi_*\mathcal{O}_X$ 上乘以 $f$ 的映射之迹。在仿射开集上，
这恢复《练习》，练习 \ref{exercises-exercise-trace-det-rings} 中的构造。
复合
$$
\mathcal{O}_Y \xrightarrow{\pi^\sharp} \pi_*\mathcal{O}_X
\xrightarrow{\text{Trace}_\pi} \mathcal{O}_Y
$$
等于乘以 $\pi$ 的次数（这是 $Y$ 上的局部常值函数）。仿照
《域论》，第 \ref{fields-section-trace-pairing} 节，可定义迹配对
$$
Q_\pi :
\pi_*\mathcal{O}_X \times \pi_*\mathcal{O}_X
\longrightarrow
\mathcal{O}_Y
$$
其规则为 $(f, g) \mapsto \text{Trace}_\pi(fg)$。可把 $Q_\pi$ 看成
相同秩的局部自由模之间的线性映射
$\pi_*\mathcal{O}_X \to
\SheafHom_{\mathcal{O}_Y}(\pi_*\mathcal{O}_X, \mathcal{O}_Y)$，
因而得到行列式
$$
\det(Q_\pi) :
\wedge^{top}(\pi_*\mathcal{O}_X)
\longrightarrow
\wedge^{top}(\pi_*\mathcal{O}_X)^{\otimes -1}
$$
换言之，得到整体截面
$$
\det(Q_\pi) \in \Gamma(Y, \wedge^{top}(\pi_*\mathcal{O}_X)^{\otimes -2})
$$
按定义，{\it $\pi$ 的判别式}是由这个整体截面定义的闭子概形
$D_\pi \subset Y$。显然，$D_\pi$ 是 $Y$ 的局部主闭子概形。

\begin{lemma}
\label{discriminant-lemma-discriminant}
设 $\pi : X \to Y$ 是有限局部自由概形态射。则 $\pi$ 是 \'etale 的，
当且仅当其判别式为空。
\end{lemma}

\begin{proof}
由《态射》，引理 \ref{morphisms-lemma-etale-flat-etale-fibres}，
只需检验 $\pi$ 的纤维为 \'etale。迹配对的构造与基变换可交换，
故问题归结为：设 $k$ 是域，$A$ 是有限维 $k$-代数，证明
$A$ 在 $k$ 上为 \'etale，当且仅当迹配对
$Q_{A/k} : A \times A \to k$、$(a, b) \mapsto \text{Trace}_{A/k}(ab)$
非退化。

\medskip\noindent
假设 $Q_{A/k}$ 非退化。若 $a \in A$ 是幂零元，则对每个 $b \in A$，
$ab$ 都是幂零的，从而 $Q_{A/k}(a, -)$ 恒为零。因此 $A$ 是既约的。
于是可写成乘积 $A = K_1 \times \ldots \times K_n$，其中每个 $K_i$
都是域（见《代数》，引理
\ref{algebra-lemma-finite-dimensional-algebra}、
\ref{algebra-lemma-artinian-finite-length} 和
\ref{algebra-lemma-minimal-prime-reduced-ring}）。
此时二次空间 $(A, Q_{A/k})$ 是各空间 $(K_i, Q_{K_i/k})$
的正交直和。由《域论》，引理
\ref{fields-lemma-separable-trace-pairing}，每个 $K_i$ 在 $k$ 上可分。
由《代数》，引理 \ref{algebra-lemma-etale-over-field}，这说明
$A$ 在 $k$ 上为 \'etale。反向结论可将上述论证倒读得到。
\end{proof}





\section{平坦拟有限环同态的迹}
\label{discriminant-section-quasi-finite-traces}

\noindent
本节标题中的迹，与《概形的对偶理论》第
\ref{duality-section-trace} 节讨论的迹性质完全不同。这里所指的是
《域论》第 \ref{fields-section-trace-pairing} 节讨论并在《练习》，练习
\ref{exercises-exercise-trace-det} 和
\ref{exercises-exercise-trace-det-rings} 中推广的迹。

\medskip\noindent
设 $A \to B$ 是 Noetherian 环的有限平坦同态。则 $B$ 作为 $A$-模
有限平坦，因而有限局部自由（《代数》，引理
\ref{algebra-lemma-finite-projective}）。给定 $b \in B$，可考察
$B$ 上乘以 $b$ 所给出的 $A$-线性映射 $B \to B$ 的
{\it 迹} $\text{Trace}_{B/A}(b)$。由上述参考资料，这定义了
$A$-线性映射 $\text{Trace}_{B/A} : B \to A$。由于 $A \to B$ 有限，
$\omega_{B/A} = \Hom_A(B, A)$，故
$\text{Trace}_{B/A} \in \omega_{B/A}$。

\medskip\noindent
对一般的平坦拟有限环同态，按如下方式定义迹的概念。

\begin{definition}
\label{discriminant-definition-trace-element}
设 $A \to B$ 是 Noetherian 环的平坦拟有限同态。
{\it 迹元}是唯一的\footnote{唯一性与存在性将在引理
\ref{discriminant-lemma-trace-unique} 和 \ref{discriminant-lemma-dualizing-tau} 中说明。}元素
$\tau_{B/A} \in \omega_{B/A}$，它满足如下性质：对任意 Noetherian
$A$-代数 $A_1$，若 $B_1 = B \otimes_A A_1$ 有乘积分解
$B_1 = C \times D$，且 $A_1 \to C$ 有限，则 $\tau_{B/A}$ 在
$\omega_{C/A_1}$ 中的像为 $\text{Trace}_{C/A_1}$。
这里利用基变换映射 (\ref{discriminant-equation-bc-dualizing}) 和引理
\ref{discriminant-lemma-dualizing-product} 得到
$\omega_{B/A} \to \omega_{B_1/A_1} \to \omega_{C/A_1}$。
\end{definition}

\noindent
先证明迹元的唯一性，再证明其存在性。

\begin{lemma}
\label{discriminant-lemma-trace-unique}
设 $A \to B$ 是 Noetherian 环的平坦拟有限同态。则
$\omega_{B/A}$ 中至多有一个迹元。
\end{lemma}

\begin{proof}
设素理想 $\mathfrak q \subset B$ 位于素理想
$\mathfrak p \subset A$ 之上。由《代数》，引理
\ref{algebra-lemma-etale-makes-quasi-finite-finite-one-prime}，
可找到 \'etale 环同态 $A \to A_1$ 和位于 $\mathfrak p$ 之上的
素理想 $\mathfrak p_1 \subset A_1$，使得
$\kappa(\mathfrak p_1) = \kappa(\mathfrak p)$，并且
$$
B_1 = B \otimes_A A_1 = C \times D
$$
其中 $A_1 \to C$ 有限，且 $B \otimes_A A_1$ 中同时位于
$\mathfrak q$ 与 $\mathfrak p_1$ 之上的唯一素理想 $\mathfrak q_1$
对应于 $C$ 的一个素理想。注意
$\omega_{C/A_1} = \omega_{B/A} \otimes_B C$
（结合引理 \ref{discriminant-lemma-dualizing-flat-base-change} 和
\ref{discriminant-lemma-dualizing-product}）。按此方式得到的环同态族 $B \to C$
是由平坦映射组成的联合单射族，而 $\tau_{B/A}$ 在
$\omega_{C/A_1}$ 中的像已被规定，故唯一性成立。
\end{proof}

\noindent
下面作一个合理性检验。

\begin{lemma}
\label{discriminant-lemma-finite-flat-trace}
设 $A \to B$ 是 Noetherian 环的有限平坦同态。则
$\text{Trace}_{B/A} \in \omega_{B/A}$ 是迹元。
\end{lemma}

\begin{proof}
设有环同态 $A \to A_1$，其中 $A_1$ 为 Noetherian，并有乘积分解
$B \otimes_A A_1 = C \times D$，其中 $A_1 \to C$ 有限。
当然，此时 $A_1 \to D$ 也有限。令 $B_1 = B \otimes_A A_1$。
迹的构造与基变换可交换，故 $\text{Trace}_{B/A}$ 映到
$\text{Trace}_{B_1/A_1}$。再注意，在引理
\ref{discriminant-lemma-dualizing-product} 的同构
$\omega_{B_1/A_1} = \omega_{C/A_1} \times \omega_{D/A_1}$ 下，有
$\text{Trace}_{B_1/A_1} = (\text{Trace}_{C/A_1}, \text{Trace}_{D/A_1})$，
证明即告完成。
\end{proof}

\begin{lemma}
\label{discriminant-lemma-trace-base-change}
设 $A \to B$ 是 Noetherian 环的平坦拟有限同态，
$\tau \in \omega_{B/A}$ 是迹元。
\begin{enumerate}
\item 若 $A \to A_1$ 是环同态，且 $A_1$ 为 Noetherian，则令
$B_1 = A_1 \otimes_A B$ 后，$\tau$ 在 $\omega_{B_1/A_1}$ 中的像是迹元。
\item 若对某个环 $R$ 与 $f \in R$ 有 $A = R_f$，则
$\tau$ 是 $\omega_{B/R}$ 中的迹元。
\item 若 $g \in B$，则 $\tau$ 在 $\omega_{B_g/A}$ 中的像是迹元。
\item 若 $B = B_1 \times B_2$，则 $\tau$ 在
$\omega_{B_1/A}$ 和 $\omega_{B_2/A}$ 中的像都是迹元。
\end{enumerate}
\end{lemma}

\begin{proof}
(1) 是定义的形式推论。

\medskip\noindent
由引理 \ref{discriminant-lemma-localize-dualizing}，
$\omega_{B/R} = \omega_{B/A}$，故 (2) 的陈述有意义。把 $\tau$
视为 $\omega_{B/R}$ 的元素，并将其记作 $\tau'$。为证明 (2)，设有
$R \to R_1$，其中 $R_1$ 为 Noetherian，并有乘积分解
$B \otimes_R R_1 = C \times D$，其中 $R_1 \to C$ 有限。
令 $A_1 = (R_1)_f$，则 $B \otimes_A A_1 = C \times D$。
由于 $R_1 \to C$ 有限，当然 $A_1 \to C$ 也有限。因此，
由 $\tau$ 的定义性质可得 $\tau'$ 的相应性质。

\medskip\noindent
由引理 \ref{discriminant-lemma-localize-dualizing}，
$\omega_{B_g/A} = (\omega_{B/A})_g$，故 (3) 的陈述有意义。
证明与 (2) 类似。设有 $A \to A_1$，其中 $A_1$ 为 Noetherian，
并有乘积分解 $B_g \otimes_A A_1 = C \times D$，其中
$A_1 \to C$ 有限。令 $B_1 = B \otimes_A A_1$。由于
$B_g \otimes_A A_1 = (B_1)_g$，$\Spec(C) \to \Spec(B_1)$ 是开浸入；
又因 $A_1 \to C$ 有限，所以 $B_1 \to C$ 有限，其像是闭的。
因此 $B_1 = C \times D_1$ 且 $D = (D_1)_g$。
应用 $\tau$ 的定义性质，即得 $\tau$ 在 $\omega_{B_g/A}$ 中之像的相应性质。

\medskip\noindent
由引理 \ref{discriminant-lemma-dualizing-product}，
$\omega_{B/A} = \omega_{B_1/A} \times \omega_{B_2/A}$，故 (4) 有意义。
设有 $A \to A'$，其中 $A'$ 为 Noetherian，并有乘积分解
$B \otimes_A A' = C \times D$，其中 $A' \to C$ 有限。
显然可把该乘积分解细化为
$B \otimes_A A' = C_1 \times C_2 \times D_1 \times D_2$，
其中 $A' \to C_i$ 有限，且 $B_i \otimes_A A' = C_i \times D_i$。
由 $\tau$ 的定义性质可得 $\tau$ 在 $\omega_{B_i/A}$ 中之像的相应性质。
这里用到显然的事实：在分解
$\omega_{C/A'} = \omega_{C_1/A'} \times \omega_{C_2/A'}$ 下，
$\text{Trace}_{C/A'} = (\text{Trace}_{C_1/A'}, \text{Trace}_{C_2/A'})$。
\end{proof}

\begin{lemma}
\label{discriminant-lemma-glue-trace}
设 $A \to B$ 是 Noetherian 环的平坦拟有限同态。
设 $g_1, \ldots, g_m \in B$ 生成单位理想，并设
$\tau \in \omega_{B/A}$ 的像在每个 $\omega_{B_{g_i}/A}$ 中
都是 $A \to B_{g_i}$ 的迹元。则 $\tau$ 是迹元。
\end{lemma}

\begin{proof}
设有 $A \to A_1$，其中 $A_1$ 为 Noetherian，并有乘积分解
$B \otimes_A A_1 = C \times D$，其中 $A_1 \to C$ 有限。
需证明 $\tau$ 在 $\omega_{C/A_1}$ 中的像为
$\text{Trace}_{C/A_1}$。注意 $g_1, \ldots, g_m$ 在
$B_1 = B \otimes_A A_1$ 中生成单位理想；由引理
\ref{discriminant-lemma-trace-base-change}，$\tau$ 映到
$\omega_{(B_1)_{g_i}/A_1}$ 中的迹元。因此可以用 $A_1$ 替换 $A$，
并用 $B_1$ 替换 $B$，归结为下一段所述情形。

\medskip\noindent
这里假设 $B = C \times D$，且 $A \to C$ 有限。令 $\tau_C$
为 $\tau$ 在 $\omega_{C/A}$ 中的像。需证明在 $\omega_{C/A}$ 中有
$\tau_C = \text{Trace}_{C/A}$。由迹元与乘积的相容性
（引理 \ref{discriminant-lemma-trace-base-change}），$\tau_C$ 映到
$\omega_{C_{g_i}/A}$ 中的迹元。因此，用 $C$ 替换 $B$ 后，
可设 $A \to B$ 有限平坦。

\medskip\noindent
假设 $A \to B$ 有限平坦。由引理 \ref{discriminant-lemma-finite-flat-trace}，
$\text{Trace}_{B/A}$ 是迹元。因此，由引理
\ref{discriminant-lemma-trace-base-change}，$\text{Trace}_{B/A}$ 映到
$\omega_{B_{g_i}/A}$ 中的迹元。迹元是唯一的（引理
\ref{discriminant-lemma-trace-unique}），故 $\text{Trace}_{B/A}$ 与 $\tau$
在 $\omega_{B_{g_i}/A} = (\omega_{B/A})_{g_i}$ 中的像相同。
由于 $g_1, \ldots, g_m$ 生成 $B$ 的单位理想，映射
$\omega_{B/A} \to \prod \omega_{B_{g_i}/A}$ 为单射，
从而如所需有 $\tau_C = \text{Trace}_{B/A}$。
\end{proof}

\begin{lemma}
\label{discriminant-lemma-dualizing-tau}
设 $A \to B$ 是 Noetherian 环的平坦拟有限同态。则存在迹元
$\tau \in \omega_{B/A}$。
\end{lemma}

\begin{proof}
选取分解 $A \to B' \to B$，其中 $A \to B'$ 有限，且
$\Spec(B) \to \Spec(B')$ 是开浸入。取 $g_1, \ldots, g_n \in B'$，
使得作为 $\Spec(B')$ 的开集有
$\Spec(B) = \bigcup D(g_i)$。假设能证明拟有限平坦环同态
$A \to B_{g_i}$ 的迹元 $\tau_i$ 存在。由引理
\ref{discriminant-lemma-trace-base-change}，对所有 $i, j$，元素 $\tau_i$ 与
$\tau_j$ 都映到 $\omega_{B_{g_ig_j}/A}$ 中的迹元。由迹元的唯一性
（引理 \ref{discriminant-lemma-trace-unique}），它们映到同一元素。因此，
$\omega_{B/A}$ 所对应拟凝聚模的层条件（见《代数》，引理
\ref{algebra-lemma-cover-module}）给出元素
$\tau \in \omega_{B/A}$。由引理 \ref{discriminant-lemma-glue-trace}，
$\tau$ 是迹元。这样便归结为下一段处理的情形。

\medskip\noindent
假设 $A \to B'$ 有限，$g \in B'$，且 $B = B'_g$ 在 $A$ 上平坦。
任务是在
$\omega_{B/A} = \Hom_A(B', A) \otimes_{B'} B$ 中构造迹元。
选取由有限自由 $A$-模 $F_0$ 与 $F_1$ 构成的 $B'$ 的解消
$F_1 \to F_0 \to B' \to 0$。于是有正合列
$$
0 \to \Hom_A(B', A) \to F_0^\vee \to F_1^\vee
$$
其中 $F_i^\vee = \Hom_A(F_i, A)$ 是对偶有限自由模。类似地有正合列
$$
0 \to \Hom_A(B', B') \to F_0^\vee \otimes_A B' \to F_1^\vee \otimes_A B'
$$
构造 $\tau$ 的想法是利用图
$$
B' \xrightarrow{\mu} \Hom_A(B', B')
\leftarrow \Hom_A(B', A) \otimes_A B'
\xrightarrow{ev} A
$$
其中第一个箭头把 $b' \in B'$ 映到乘以 $b'$ 所给出的 $A$-线性算子，
末一个箭头是取值映射。问题在于中间的箭头——它把
$\lambda' \otimes b'$ 映到 $b'' \mapsto \lambda'(b'')b'$——并非同构。
若 $B'$ 在 $A$ 上平坦，上述正合列说明它是同构，并且从左到右的复合
就是通常的迹 $\text{Trace}_{B'/A}$。一般情形下，考察图
$$
\xymatrix{
& \Hom_A(B', A) \otimes_A B' \ar[r] \ar[d] &
\Hom_A(B', A) \otimes_A B'_g \ar[d] \\
B' \ar[r]_-\mu \ar@{..>}[rru] \ar@{..>}[ru]^\psi &
\Hom_A(B', B') \ar[r] &
\Ker(F_0^\vee \otimes_A B'_g \to F_1^\vee \otimes_A B'_g)
}
$$
由 $A \to B'_g$ 的平坦性，右侧竖直箭头是同构。因此得到未标记的虚箭头。
由于 $B'_g = \colim \frac{1}{g^n}B'$，余极限与张量积可交换，
且 $B'$ 是有限表现 $A$-模，所以可找到 $n \geq 0$ 以及一个
$B'$-线性映射（对右 $B'$-模结构而言）
$\psi : B' \to \Hom_A(B', A) \otimes_A B'$，
使它与左侧竖直箭头的复合为 $g^n\mu$。再与 $ev$ 复合，得到元素
$ev \circ \psi \in \Hom_A(B', A)$。令
$$
\tau = (ev \circ \psi) \otimes g^{-n} \in
\Hom_A(B', A) \otimes_{B'} B'_g = \omega_{B'_g/A} = \omega_{B/A}
$$
略去这个元素与上述 $n$ 和 $\psi$ 的选取无关的直接验证。

\medskip\noindent
先在一个特殊情形中证明上一段构造的 $\tau$ 具有所需性质。
具体地，设 $B' = C' \times D'$，$g = (f, h)$，其中
$A \to C'$ 平坦，$D'_h$ 平坦，且 $f$ 是 $C'$ 中的单位。
需证明：$\tau$ 在 $\omega_{C'/A}$ 中映到 $\text{Trace}_{C'/A}$。
先如上为一对 $(D', h)$ 选取 $n_D$ 以及
$\psi_D : D' \to \Hom_A(D', A) \otimes_A D'$，并令
$\psi_C : C' \to \Hom_A(C', A) \otimes_A C' = \Hom_A(C', C')$
把 $c' \in C'$ 映到乘以 $c'$。然后取
$n = n_D$ 和 $\psi = (f^{n_D} \psi_C, \psi_D)$。
由上述 $\text{Trace}_{C'/A} = ev \circ \psi_C$，所需相容性显然成立。

\medskip\noindent
为证明一般情形下的所需性质，设给定 $A \to A_1$，其中 $A_1$
为 Noetherian，并有乘积分解 $B'_g \otimes_A A_1 = C \times D$，
其中 $A_1 \to C$ 有限。令 $B'_1 = B' \otimes_A A_1$。
由于 $B'_g \otimes_A A_1 = (B'_1)_g$，
$\Spec(C) \to \Spec(B'_1)$ 是开浸入；又因 $A_1 \to C$ 有限，
$B'_1 \to C$ 也有限，故其像为闭集。因此
$B'_1 = C \times D'$ 且 $D'_g = D$。可见 $A_1$ 上的
$B'_1 = C \times D'$ 与 $g$ 正是上一段的情形。
上述显示图的形成与基变换可交换，因此 $\tau$ 的形成与基变换
$A \to A_1$ 可交换（略去细节；可用解消
$F_1 \otimes_A A_1 \to F_0 \otimes_A A_1 \to B'_1 \to 0$
看出这一点）。所需相容性遂由上一段的结论得出。
\end{proof}

\begin{remark}
\label{discriminant-remark-relative-dualizing-for-flat-quasi-finite}
设 $f : Y \to X$ 是局部 Noetherian 概形之间的平坦局部拟有限态射，
并令 $\omega_{Y/X}$ 如注记
\ref{discriminant-remark-relative-dualizing-for-quasi-finite} 所述。由迹元的唯一性、
存在性以及与局部化的相容性（引理 \ref{discriminant-lemma-trace-unique}、
\ref{discriminant-lemma-dualizing-tau} 和 \ref{discriminant-lemma-trace-base-change}），
显然存在整体截面
$$
\tau_{Y/X} \in \Gamma(Y, \omega_{Y/X})
$$
使得对每一对满足 $f(V) \subset U$ 的仿射开集
$\Spec(B) = V \subset Y$、$\Spec(A) = U \subset X$，
元素 $\tau_{Y/X}$ 在典范同构
$H^0(V, \omega_{Y/X}) = \omega_{B/A}$ 下映到 $\tau_{B/A}$。
\end{remark}

\begin{lemma}
\label{discriminant-lemma-tau-nonzero}
设 $k$ 是域，$A$ 是有限 $k$-代数。假设 $A$ 为局部环，其剩余域为
$k'$。则下列条件等价：
\begin{enumerate}
\item $\text{Trace}_{A/k}$ 非零；
\item $\tau_{A/k} \in \omega_{A/k}$ 非零；
\item $k'/k$ 可分，且 $\text{length}_A(A)$ 与 $k$ 的特征互素。
\end{enumerate}
\end{lemma}

\begin{proof}
由引理 \ref{discriminant-lemma-finite-flat-trace}，条件 (1) 与 (2) 等价。
设 $\mathfrak m \subset A$。由于 $\dim_k(A) < \infty$，显然 $A$
作为 $A$-模具有有限长度。选取理想滤过
$$
A = I_0 \supset \mathfrak m = I_1 \supset I_2 \supset \ldots I_n = 0
$$
使得作为 $A$-模有 $I_i/I_{i + 1} \cong k'$。参见《代数》，引理
\ref{algebra-lemma-simple-pieces}；该引理还说明
$n = \text{length}_A(A)$。若 $a \in \mathfrak m$，则
$aI_i \subset I_{i + 1}$，从而显然有 $\text{Trace}_{A/k}(a) = 0$。
若 $a \not \in \mathfrak m$，其像为 $\lambda \in k'$，则得到
$$
\text{Trace}_{A/k}(a) =
\sum\nolimits_{i = 0, \ldots, n - 1}
\text{Trace}_k(a : I_i/I_{i - 1} \to I_i/I_{i - 1}) =
n \text{Trace}_{k'/k}(\lambda)
$$
最后应用《域论》，引理 \ref{fields-lemma-separable-trace-pairing}，
即得引理。
\end{proof}






\section{有限态射}
\label{discriminant-section-finite-morphisms}

\noindent
本节汇集前面各节的构造在有限态射情形下的若干说明。
设 $f : Y \to X$ 是局部 Noetherian 概形的有限态射，并令
$\omega_{Y/X}$ 如注记 \ref{discriminant-remark-relative-dualizing-for-quasi-finite} 所述。

\medskip\noindent
第一点说明是，作为 $f_*\mathcal{O}_Y$-模层有
$$
f_*\omega_{Y/X} = \SheafHom_{\mathcal{O}_X}(f_*\mathcal{O}_Y, \mathcal{O}_X)
$$
由于 $f$ 是仿射的，这个公式唯一刻画 $\omega_{Y/X}$；见《态射》，引理
\ref{morphisms-lemma-affine-equivalence-modules}。该公式成立是因为：
若 $\Spec(A) = U \subset X$ 是仿射开集，则逆像
$V = f^{-1}(U)$ 是某个有限 $A$-代数 $B$ 的谱，因而按构造有
$$
H^0(U, f_*\omega_{Y/X}) =
H^0(V, \omega_{Y/X}) =
\omega_{B/A} =
\Hom_A(B, A) =
H^0(U, \SheafHom_{\mathcal{O}_X}(f_*\mathcal{O}_Y, \mathcal{O}_X))
$$
特别地，得到典范取值映射
$$
f_*\omega_{Y/X} \longrightarrow \mathcal{O}_X
$$
若把 $f_*\omega_{Y/X}$ 看成层
$\SheafHom_{\mathcal{O}_X}(f_*\mathcal{O}_Y, \mathcal{O}_X)$，
这个映射由在 $1$ 处取值得到。

\medskip\noindent
第二点说明是，利用取值映射，可以典范地得到等同
$$
\Hom_Y(\mathcal{F}, f^*\mathcal{G} \otimes_{\mathcal{O}_Y} \omega_{Y/X})
=
\Hom_X(f_*\mathcal{F}, \mathcal{G})
$$
它关于 $Y$ 上的拟凝聚模 $\mathcal{F}$ 和 $X$ 上的有限局部自由模
$\mathcal{G}$ 具有函子性。若 $\mathcal{G} = \mathcal{O}_X$，
这立即由上述内容和《代数》，引理
\ref{algebra-lemma-adjoint-hom-restrict} 得出。对一般的 $\mathcal{G}$，
可应用同一引理以及 $f_*\mathcal{O}_Y$-模的同构
$$
f_*(f^*\mathcal{G} \otimes_{\mathcal{O}_Y} \omega_{Y/X}) =
\mathcal{G} \otimes_{\mathcal{O}_X}
\SheafHom_{\mathcal{O}_X}(f_*\mathcal{O}_Y, \mathcal{O}_X) =
\SheafHom_{\mathcal{O}_X}(f_*\mathcal{O}_Y, \mathcal{G})
$$
其中第一个等号是投影公式
（《上同调》，引理 \ref{cohomology-lemma-projection-formula}）。
也可以在仿射局部作直接计算来证明该公式。

\medskip\noindent
第三点说明是，若 $f$ 还平坦，则复合
$$
f_*\mathcal{O}_Y \xrightarrow{f_*\tau_{Y/X}} f_*\omega_{Y/X}
\longrightarrow \mathcal{O}_X
$$
等于第 \ref{discriminant-section-discriminant} 节讨论的迹映射
$\text{Trace}_f$。在仿射开集上考察即可立即看出这一点。

\medskip\noindent
第四点说明是，若 $f$ 平坦且 $X$ 为 Noetherian，则对任意
$D_\QCoh(\mathcal{O}_Y)$ 中的 $K$ 和 $D_\QCoh(\mathcal{O}_X)$ 中的 $M$，有
$$
\Hom_Y(K, Lf^*M \otimes_{\mathcal{O}_Y} \omega_{Y/X})
=
\Hom_X(Rf_*K, M)
$$
这可由《概形的对偶理论》第 \ref{duality-section-proper-flat} 节的内容得出，
但在这里也可直接证明如下。首先，若 $X$ 仿射，则由《对偶化复形》，引理
\ref{dualizing-lemma-right-adjoint} 和
\ref{dualizing-lemma-RHom-is-tensor-special}\footnote{在本情形下，
这个引理有更简单的证明。}以及《概形的导出范畴》，引理
\ref{perfect-lemma-affine-compare-bounded}，结论成立。然后使用归纳原理
（《概形的上同调》，引理 \ref{coherent-lemma-induction-principle}）与
Mayer--Vietoris（采用《上同调》，引理
\ref{cohomology-lemma-mayer-vietoris-hom} 的形式）完成证明。







\section{Noether 不同理想}
\label{discriminant-section-noether-different}

\noindent
文献中有许多不同的“不同理想”。本节及下节列出其中一些；
更多信息建议读者参阅 \cite{Kunz}。

\medskip\noindent
设 $A \to B$ 是环同态。将乘法映射记为
$$
\mu : B \otimes_A B \longrightarrow B,\quad
b \otimes b' \longmapsto bb'
$$
并令 $I = \Ker(\mu)$。显然，$I$ 由所有 $b \in B$ 对应的元素
$b \otimes 1 - 1 \otimes b$ 生成。因此，$I$ 的零化子
$J \subset B \otimes_A B$ 典范地成为 $B$-模。
$B$ 在 $A$ 上的 {\it Noether 不同理想}是 $J$ 在映射
$\mu : B \otimes_A B \to B$ 下的像。等价地，Noether 不同理想是映射
$$
J = \Hom_{B \otimes_A B}(B, B \otimes_A B) \longrightarrow B,\quad
\varphi \longmapsto \mu(\varphi(1))
$$
的像。先证明若干必要的引理。

\begin{lemma}
\label{discriminant-lemma-noether-different-product}
设 $A \to B_i$（$i = 1, 2$）为环同态，并令 $B = B_1 \times B_2$。
\begin{enumerate}
\item $\Ker(B \otimes_A B \to B)$ 的零化子 $J$ 等于
$J_1 \times J_2$，其中 $J_i$ 是
$\Ker(B_i \otimes_A B_i \to B_i)$ 的零化子。
\item $B$ 在 $A$ 上的 Noether 不同理想 $\mathfrak{D}$ 等于
$\mathfrak{D}_1 \times \mathfrak{D}_2$，其中 $\mathfrak{D}_i$
是 $B_i$ 在 $A$ 上的 Noether 不同理想。
\end{enumerate}
\end{lemma}

\begin{proof}
略。
\end{proof}

\begin{lemma}
\label{discriminant-lemma-noether-different-base-change}
设 $A \to B$ 是有限型环同态，$A \to A'$ 是平坦环同态，并令
$B' = B \otimes_A A'$。
\begin{enumerate}
\item $\Ker(B' \otimes_{A'} B' \to B')$ 的零化子 $J'$ 等于
$J \otimes_A A'$，其中 $J$ 是 $\Ker(B \otimes_A B \to B)$ 的零化子。
\item $B'$ 在 $A'$ 上的 Noether 不同理想 $\mathfrak{D}'$ 等于
$\mathfrak{D}B'$，其中 $\mathfrak{D}$ 是 $B$ 在 $A$ 上的 Noether 不同理想。
\end{enumerate}
\end{lemma}

\begin{proof}
选取 $B$ 作为 $A$-代数的一组生成元 $b_1, \ldots, b_n$。则
$$
J = \Ker(B \otimes_A B \xrightarrow{b_i \otimes 1 - 1 \otimes b_i}
(B \otimes_A B)^{\oplus n})
$$
因此 $J$ 的形成与平坦基变换可交换。关于 Noether 不同理想的结论
立即由此得出。
\end{proof}

\begin{lemma}
\label{discriminant-lemma-noether-different-localization}
设 $A \to B' \to B$ 是环同态，其中 $A \to B'$ 为有限型，
而 $B' \to B$ 在谱上诱导开浸入。
\begin{enumerate}
\item $\Ker(B \otimes_A B \to B)$ 的零化子 $J$ 等于
$J' \otimes_{B'} B$，其中 $J'$ 是
$\Ker(B' \otimes_A B' \to B')$ 的零化子。
\item $B$ 在 $A$ 上的 Noether 不同理想 $\mathfrak{D}$ 等于
$\mathfrak{D}'B$，其中 $\mathfrak{D}'$ 是 $B'$ 在 $A$ 上的
Noether 不同理想。
\end{enumerate}
\end{lemma}

\begin{proof}
记 $I = \Ker(B \otimes_A B \to B)$，并记
$I' = \Ker(B' \otimes_A B' \to B')$。由于
$\Spec(B) \to \Spec(B')$ 是开浸入，得到
$B = (B \otimes_A B) \otimes_{B' \otimes_A B'} B'$。
因此 $I = I'(B \otimes_A B)$。由于 $I'$ 有限生成，且
$B' \otimes_A B' \to B \otimes_A B$ 平坦，有
$J = J'(B \otimes_A B)$；见《代数》，引理
\ref{algebra-lemma-annihilator-flat-base-change}。
由于 $J'$ 的 $B' \otimes_A B'$-模结构经过
$B' \otimes_A B' \to B'$ 分解，故 (1) 成立。(2) 是 (1) 的推论。
\end{proof}

\begin{remark}
\label{discriminant-remark-construction-pairing}
设 $A \to B$ 是 Noetherian 环的拟有限同态，并令 $J$ 为
$\Ker(B \otimes_A B \to B)$ 的零化子。存在典范的 $B$-双线性配对
\begin{equation}
\label{discriminant-equation-pairing-noether}
\omega_{B/A} \times J \longrightarrow B
\end{equation}
其定义如下。选取分解 $A \to B' \to B$，其中 $A \to B'$ 有限，
而 $B' \to B$ 在谱上诱导开浸入。令 $J'$ 为
$\Ker(B' \otimes_A B' \to B')$ 的零化子。先定义
$$
\Hom_A(B', A) \times J' \longrightarrow B',\quad
(\lambda, \sum b_i \otimes c_i) \longmapsto \sum \lambda(b_i)c_i
$$
这恰为 $B'$-双线性的，因为对 $\xi \in J'$ 与 $b \in B'$，有
$(b \otimes 1)\xi = (1 \otimes b)\xi$。由引理
\ref{discriminant-lemma-noether-different-localization} 以及
$\omega_{B/A} = \Hom_A(B', A) \otimes_{B'} B$，
可将它扩张为上面所示的 $B$-双线性配对。
\end{remark}

\begin{lemma}
\label{discriminant-lemma-noether-pairing-compatibilities}
设 $A \to B$ 是 Noetherian 环的拟有限同态。
\begin{enumerate}
\item 若 $A \to A'$ 是 Noetherian 环的平坦同态，则
$$
\xymatrix{
\omega_{B/A} \times J \ar[r] \ar[d] & B \ar[d]  \\
\omega_{B'/A'} \times J' \ar[r] & B'
}
$$
交换，其中记号同引理 \ref{discriminant-lemma-noether-different-base-change}，
水平箭头由 (\ref{discriminant-equation-pairing-noether}) 给出。
\item 若 $B = B_1 \times B_2$，则
$$
\xymatrix{
\omega_{B/A} \times J \ar[r] \ar[d] & B \ar[d]  \\
\omega_{B_i/A} \times J_i \ar[r] & B_i
}
$$
对 $i = 1, 2$ 交换，其中记号同引理
\ref{discriminant-lemma-noether-different-product}，水平箭头由
(\ref{discriminant-equation-pairing-noether}) 给出。
\end{enumerate}
\end{lemma}

\begin{proof}
由注记 \ref{discriminant-remark-construction-pairing} 中配对的构造，
(1) 与 (2) 都归结为 $A \to B$ 有限的情形。此时 (1) 由收缩映射
$\Hom_A(M, A) \otimes_A M \otimes_A M \to M$、
$\lambda \otimes m \otimes m' \mapsto \lambda(m)m'$
与基变换可交换这一事实得出。为证明 (2)，使用
$J = J_1 \times J_2$ 包含于 $B \otimes_A B$ 的直和分量
$B_1 \otimes_A B_1$ 与 $B_2 \otimes_A B_2$ 中这一事实。
\end{proof}

\begin{lemma}
\label{discriminant-lemma-noether-pairing-flat-quasi-finite}
设 $A \to B$ 是 Noetherian 环的平坦拟有限同态。注记
\ref{discriminant-remark-construction-pairing} 的配对诱导同构
$J \to \Hom_B(\omega_{B/A}, B)$。
\end{lemma}

\begin{proof}
先在 $A \to B$ 有限平坦时证明。此时可在 $A$ 上局部化，并假设
$B$ 作为 $A$-模是有限自由的。设 $b_1, \ldots, b_n$ 是 $B$
作为 $A$-模的一组基，并将 $\omega_{B/A}$ 的对偶基记为
$b_1^\vee, \ldots, b_n^\vee$。注意
$\sum b_i \otimes c_i \in J$ 映到 $\Hom_B(\omega_{B/A}, B)$ 中
把 $b_i^\vee$ 送到 $c_i$ 的元素。设
$\varphi : \omega_{B/A} \to B$ 为 $B$-线性映射。断言
$\xi = \sum b_i \otimes \varphi(b_i^\vee)$ 属于 $J$。
事实上，$\varphi$ 的 $B$-线性恰好蕴含：对每个 $b \in B$，
$(b \otimes 1)\xi = (1 \otimes b)\xi$。因此该映射有逆，故为同构。

\medskip\noindent
设素理想 $\mathfrak q \subset B$ 位于 $\mathfrak p \subset A$ 之上。
我们证明局部化
$$
J_\mathfrak q
\longrightarrow
\Hom_B(\omega_B/A, B)_\mathfrak q
$$
是同构。由《代数》，引理 \ref{algebra-lemma-characterize-zero-local}，
这就足够了。由《代数》，引理
\ref{algebra-lemma-etale-makes-quasi-finite-finite-one-prime}，
可找到 \'etale 环同态 $A \to A'$ 以及位于 $\mathfrak p$ 之上的素理想
$\mathfrak p' \subset A'$，使得 $\kappa(\mathfrak p') = \kappa(\mathfrak p)$，
并且
$$
B' = B \otimes_A A' = C \times D
$$
其中 $A' \to C$ 有限，且 $B \otimes_A A'$ 中同时位于
$\mathfrak q$ 与 $\mathfrak p'$ 之上的唯一素理想 $\mathfrak q'$
对应于 $C$ 的一个素理想。令 $J'$ 为
$\Ker(B' \otimes_{A'} B' \to B')$ 的零化子。由引理
\ref{discriminant-lemma-dualizing-flat-base-change}、
\ref{discriminant-lemma-noether-different-base-change} 和
\ref{discriminant-lemma-noether-pairing-compatibilities}，映射
$J' \to \Hom_{B'}(\omega_{B'/A'}, B')$ 是对映射
$J \to \Hom_B(\omega_{B/A}, B)$ 应用函子 $- \otimes_B B'$ 得到的。
由于 $B_\mathfrak q \to B'_{\mathfrak q'}$ 忠实平坦，只需对
$(A' \to B', \mathfrak q')$ 证明结论。再由引理
\ref{discriminant-lemma-dualizing-product}、
\ref{discriminant-lemma-noether-different-product} 和
\ref{discriminant-lemma-noether-pairing-compatibilities}，归结为证明第一段已经处理的情形。
\end{proof}

\begin{lemma}
\label{discriminant-lemma-noether-different-flat-quasi-finite}
设 $A \to B$ 是 Noetherian 环的平坦拟有限同态。图
$$
\xymatrix{
J \ar[rr] \ar[rd]_\mu & &
\Hom_B(\omega_{B/A}, B) \ar[ld]^{\varphi \mapsto \varphi(\tau_{B/A})} \\
& B
}
$$
交换，其中水平箭头是引理
\ref{discriminant-lemma-noether-pairing-flat-quasi-finite} 的同构。
因此，$B$ 在 $A$ 上的 Noether 不同理想是映射
$\Hom_B(\omega_{B/A}, B) \to B$ 的像。
\end{lemma}

\begin{proof}
完全仿照引理 \ref{discriminant-lemma-noether-pairing-flat-quasi-finite} 的证明，
归结为有限自由同态 $A \to B$ 的情形。此时
$\tau_{B/A} = \text{Trace}_{B/A}$。选取 $B$ 作为 $A$-模的一组基
$b_1, \ldots, b_n$。令 $\xi = \sum b_i \otimes c_i \in J$，则
$\mu(\xi) = \sum b_i c_i$。另一方面，$\xi$ 在
$\Hom_B(\omega_{B/A}, B)$ 中的像把 $\text{Trace}_{B/A}$ 映到
$\sum \text{Trace}_{B/A}(b_i)c_i$。因此需证明
$$
\sum b_ic_i = \sum \text{Trace}_{B/A}(b_i)c_i
$$
其中 $\xi = \sum b_i \otimes c_i \in J$。写成
$b_i b_j = \sum_k a_{ij}^k b_k$，其中 $a_{ij}^k \in A$。
则右端为 $\sum_{i, j} a_{ij}^j c_i$。另一方面，$\xi \in J$ 蕴含
$$
(b_j \otimes 1)(\sum\nolimits_i b_i \otimes c_i) =
(1 \otimes b_j)(\sum\nolimits_i b_i \otimes c_i)
$$
从而 $b_j c_i = \sum_k a_{jk}^i c_k$。因此左端为
$\sum_{i, j} a_{ij}^i c_j$。由于 $a_{ij}^k = a_{ji}^k$，等式成立。
\end{proof}

\begin{lemma}
\label{discriminant-lemma-noether-different}
设 $A \to B$ 是有限型环同态，$\mathfrak{D} \subset B$ 是 Noether
不同理想。则 $V(\mathfrak{D})$ 是所有使 $A \to B$ 在
$\mathfrak q$ 处分歧的素理想 $\mathfrak q \subset B$ 所成的集合。
\end{lemma}

\begin{proof}
假设 $A \to B$ 在 $\mathfrak q$ 处非分歧。对某个
$g \in B$、$g \not \in \mathfrak q$，用 $B_g$ 替换 $B$ 后，
可设 $A \to B$ 非分歧（《代数》，定义
\ref{algebra-definition-unramified} 以及引理
\ref{discriminant-lemma-noether-different-localization}）。此时
$\Omega_{B/A} = 0$。因此，若 $I = \Ker(B \otimes_A B \to B)$，
则由《代数》，引理 \ref{algebra-lemma-differentials-diagonal}，
$I/I^2 = 0$。由于 $A \to B$ 为有限型，$I$ 有限生成。
于是由 Nakayama 引理（《代数》，引理 \ref{algebra-lemma-NAK}），
存在形如 $1 + i$ 的元素零化 $I$。从而 $\mathfrak{D} = B$。

\medskip\noindent
反之，假设 $\mathfrak{D} \not \subset \mathfrak q$。
按上述方式用一个主局部化替换 $B$ 后，可设 $\mathfrak{D} = B$。
这意味着 $I$ 的零化子中存在形如 $1 + i$ 的元素。
反过来，这又蕴含 $I/I^2 = \Omega_{B/A}$ 为零，故得结论。
\end{proof}







\section{K\"ahler 不同理想}
\label{discriminant-section-kahler-different}

\noindent
设 $A \to B$ 是有限型环同态。{\it K\"ahler 不同理想}是
$\Omega_{B/A}$ 作为 $B$-模的第零 Fitting 理想。将该定义整体化如下。

\begin{definition}
\label{discriminant-definition-kahler-different}
设 $f : Y \to X$ 是局部有限型概形态射。{\it K\"ahler 不同理想}
是 $\Omega_{Y/X}$ 的第 $0$ 个 Fitting 理想。
\end{definition}

\noindent
K\"ahler 不同理想是 $Y$ 上的拟凝聚理想层。

\begin{lemma}
\label{discriminant-lemma-base-change-kahler-different}
考虑概形的笛卡尔图
$$
\xymatrix{
Y' \ar[d]_{f'} \ar[r] & Y \ar[d]^f \\
X' \ar[r]^g & X
}
$$
其中 $f$ 局部为有限型。令 $R \subset Y$、相应地\ 令
$R' \subset Y'$ 为分别由 $f$、相应地\ 由 $f'$ 的 K\"ahler 不同理想
定义的闭子概形。则 $Y' \to Y$ 诱导同构 $R' \to R \times_Y Y'$。
\end{lemma}

\begin{proof}
这是因为 $\Omega_{Y'/X'}$ 是 $\Omega_{Y/X}$ 的拉回
（《态射》，引理 \ref{morphisms-lemma-base-change-differentials}），
然后可应用《代数进阶》，引理
\ref{more-algebra-lemma-fitting-ideal-basics}。
\end{proof}

\begin{lemma}
\label{discriminant-lemma-kahler-different}
设 $f : Y \to X$ 是局部有限型概形态射，并令 $R \subset Y$
为 K\"ahler 不同理想所定义的闭子概形。则 $R \subset Y$
恰为 $f$ 分歧的点集。
\end{lemma}

\begin{proof}
这是《除子》，引理
\ref{divisors-lemma-zero-fitting-ideal-omega-unramified} 的同一证明。
\end{proof}

\begin{lemma}
\label{discriminant-lemma-kahler-different-complete-intersection}
设 $A$ 是环，$n \geq 1$，且
$f_1, \ldots, f_n \in A[x_1, \ldots, x_n]$。令
$B = A[x_1, \ldots, x_n]/(f_1, \ldots, f_n)$。
则 $B$ 在 $A$ 上的 K\"ahler 不同理想是由
$\det(\partial f_i/\partial x_j)$ 生成的 $B$ 的理想。
\end{lemma}

\begin{proof}
这是因为由《代数》，引理 \ref{algebra-lemma-differential-seq}，
$\Omega_{B/A}$ 有表现
$$
\bigoplus\nolimits_{i = 1, \ldots, n} B f_i
\xrightarrow{\text{d}}
\bigoplus\nolimits_{j = 1, \ldots, n} B \text{d}x_j
\rightarrow \Omega_{B/A} \rightarrow 0
$$
\end{proof}



\section{Dedekind 不同理想}
\label{discriminant-section-dedekind-different}

\noindent
设 $A \to B$ 是环同态。若 $A$ 为 Noetherian、$A \to B$ 有限、
$A$ 上的每个非零因子在 $B$ 上仍为非零因子，并且当
$K = Q(A)$、$L = B \otimes_A K$ 时 $K \to L$ 为 \'etale，
则称{\it Dedekind 不同理想有定义}。此时 $K \subset L$ 有限 \'etale，且
$$
\mathcal{L}_{B/A} = \{x \in L \mid \text{Trace}_{L/K}(bx) \in A
\text{ 对所有 }b \in B\}
$$
是 Dedekind 补模。在这种情形下，{\it Dedekind 不同理想}是
$$
\mathfrak{D}_{B/A} = \{x \in L \mid x\mathcal{L}_{B/A} \subset B\}
$$
把它看成 $L$ 的 $B$-子模。由引理
\ref{discriminant-lemma-dedekind-different-ideal}，若 $A$ 正规，或 $B$ 在 $A$ 上平坦，
则 Dedekind 不同理想是 $B$ 的理想。

\begin{lemma}
\label{discriminant-lemma-dedekind-different-ideal}
假设 $A \to B$ 的 Dedekind 不同理想有定义。考虑下列陈述：
\begin{enumerate}
\item $A \to B$ 平坦；
\item $A$ 是正规环；
\item $\text{Trace}_{L/K}(B) \subset A$；
\item $1 \in \mathcal{L}_{B/A}$；
\item Dedekind 不同理想 $\mathfrak{D}_{B/A}$ 是 $B$ 的理想。
\end{enumerate}
则有 (1) $\Rightarrow$ (3)、(2) $\Rightarrow$ (3)、
(3) $\Leftrightarrow$ (4) 以及 (4) $\Rightarrow$ (5)。
\end{lemma}

\begin{proof}
(3) 与 (4) 的等价以及蕴含 (4) $\Rightarrow$ (5) 都是立即的。

\medskip\noindent
若 $A \to B$ 平坦，则 $\text{Trace}_{B/A} : B \to A$ 有定义，
而 $\text{Trace}_{L/K}$ 是其基变换。因此 (3) 成立。

\medskip\noindent
若 $A$ 正规，则 $A$ 是有限多个正规整环的乘积，故归结为正规整环情形。
此时 $K$ 是 $A$ 的分式域，而 $L = \prod L_i$ 是有限多个
$K$ 的有限可分域扩张之乘积。于是
$\text{Trace}_{L/K}(b) = \sum \text{Trace}_{L_i/K}(b_i)$，
其中 $b_i \in L_i$ 是 $b$ 的像。由于 $B$ 在 $A$ 上有限，
$b$ 在 $A$ 上整，故这些迹属于 $A$。事实上，$b_i$ 在 $K$ 上的
极小多项式的系数属于 $A$（《代数》，引理
\ref{algebra-lemma-minimal-polynomial-normal-domain}），并且
$\text{Trace}_{L_i/K}(b_i)$ 是其中某个系数的整数倍
（《域论》，引理
\ref{fields-lemma-trace-and-norm-from-minimal-polynomial}）。
\end{proof}

\begin{lemma}
\label{discriminant-lemma-dedekind-complementary-module}
若 $A \to B$ 的 Dedekind 不同理想有定义，则存在典范同构
$\mathcal{L}_{B/A} \to \omega_{B/A}$。
\end{lemma}

\begin{proof}
回忆 $A \to B$ 有限，故 $\omega_{B/A} = \Hom_A(B, A)$。
把 $x \in \mathcal{L}_{B/A}$ 映到映射
$b \mapsto \text{Trace}_{L/K}(bx)$。反之，给定 $A$-线性映射
$\varphi : B \to A$，得到 $K$-线性映射 $\varphi_K : L \to K$。
由于 $K \to L$ 有限 \'etale，迹配对非退化（引理
\ref{discriminant-lemma-discriminant}），故存在 $x \in L$，使得对每个 $y \in L$，
$\varphi_K(y) = \text{Trace}_{L/K}(xy)$。于是
$x \in \mathcal{L}_{B/A}$ 在 $\omega_{B/A}$ 中映到 $\varphi$。
\end{proof}

\begin{lemma}
\label{discriminant-lemma-flat-dedekind-complementary-module-trace}
若 $A \to B$ 的 Dedekind 不同理想有定义，且 $A \to B$ 平坦，则
\begin{enumerate}
\item 典范同构 $\mathcal{L}_{B/A} \to \omega_{B/A}$ 把
$1 \in \mathcal{L}_{B/A}$ 映到迹元
$\tau_{B/A} \in \omega_{B/A}$；
\item Dedekind 不同理想为
$\mathfrak{D}_{B/A} = \{b \in B \mid b\omega_{B/A} \subset B\tau_{B/A}\}$。
\end{enumerate}
\end{lemma}

\begin{proof}
第一个断言由引理 \ref{discriminant-lemma-dedekind-different-ideal} 的证明和引理
\ref{discriminant-lemma-finite-flat-trace} 得出。第二个断言立即由第一个断言和定义得出。
\end{proof}



\section{不同理想}
\label{discriminant-section-different}

\noindent
下面定义的动机在于：正如下文将看到的，它在有限平坦情形下恢复
Dedekind 不同理想。

\begin{definition}
\label{discriminant-definition-different}
设 $f : Y \to X$ 是局部 Noetherian 概形之间的平坦局部拟有限态射。
令 $\omega_{Y/X}$ 为相对对偶化模，并令
$\tau_{Y/X} \in \Gamma(Y, \omega_{Y/X})$ 为迹元
（注记 \ref{discriminant-remark-relative-dualizing-for-quasi-finite} 和
\ref{discriminant-remark-relative-dualizing-for-flat-quasi-finite}）。
$$
\Coker(\mathcal{O}_Y \xrightarrow{\tau_{Y/X}} \omega_{Y/X})
$$
的零化子称为 $Y/X$ 的{\it 不同理想}。它是凝聚理想
$\mathfrak{D}_f \subset \mathcal{O}_Y$。
\end{definition}

\noindent
下文注记 \ref{discriminant-remark-different-generalization} 将推广这个定义。
注意，若 $\omega_{Y/X}$ 是可逆 $\mathcal{O}_Y$-模，则
$\mathfrak{D}_f$ 局部由一个元素生成。先陈述它与 Dedekind 不同理想的一致性。

\begin{lemma}
\label{discriminant-lemma-flat-agree-dedekind}
设 $f : Y \to X$ 是 Noetherian 概形的平坦拟有限态射。
设 $V = \Spec(B) \subset Y$、$U = \Spec(A) \subset X$ 是满足
$f(V) \subset U$ 的仿射开子概形。若 $A \to B$ 的 Dedekind
不同理想有定义，则作为 $V$ 上的凝聚理想层有
$$
\mathfrak{D}_f|_V = \widetilde{\mathfrak{D}_{B/A}}
$$
\end{lemma}

\begin{proof}
这由引理 \ref{discriminant-lemma-dedekind-different-ideal} 和
\ref{discriminant-lemma-flat-dedekind-complementary-module-trace} 显然成立。
\end{proof}

\begin{lemma}
\label{discriminant-lemma-flat-gorenstein-agree-noether}
设 $f : Y \to X$ 是 Noetherian 概形的平坦拟有限态射。
设 $V = \Spec(B) \subset Y$、$U = \Spec(A) \subset X$ 是满足
$f(V) \subset U$ 的仿射开子概形。若 $\omega_{Y/X}|_V$ 可逆，
即若 $\omega_{B/A}$ 是可逆 $B$-模，则作为 $V$ 上的凝聚理想层有
$$
\mathfrak{D}_f|_V = \widetilde{\mathfrak{D}}
$$
其中 $\mathfrak{D} \subset B$ 是 $B$ 在 $A$ 上的 Noether 不同理想。
\end{lemma}

\begin{proof}
考虑映射
$$
\SheafHom_{\mathcal{O}_Y}(\omega_{Y/X}, \mathcal{O}_Y)
\longrightarrow
\mathcal{O}_Y,\quad
\varphi \longmapsto \varphi(\tau_{Y/X})
$$
在仿射开集上，该映射的像对应于 Noether 不同理想；见引理
\ref{discriminant-lemma-noether-different-flat-quasi-finite}。因此结论归结为如下初等事实：
给定可逆模 $\omega$ 与整体截面 $\tau$，映射
$\tau : \SheafHom(\omega, \mathcal{O}) = \omega^{\otimes -1} \to \mathcal{O}$
的像等于 $\Coker(\tau : \mathcal{O} \to \omega)$ 的零化子。
\end{proof}

\begin{lemma}
\label{discriminant-lemma-base-change-different}
考虑 Noetherian 概形的笛卡尔图
$$
\xymatrix{
Y' \ar[d]_{f'} \ar[r] & Y \ar[d]^f \\
X' \ar[r]^g & X
}
$$
其中 $f$ 平坦且拟有限。令 $R \subset Y$、相应地\ 令
$R' \subset Y'$ 为由不同理想 $\mathfrak{D}_f$、相应地\ 由
$\mathfrak{D}_{f'}$ 定义的闭子概形。则 $Y' \to Y$ 诱导双射闭浸入
$R' \to R \times_Y Y'$。若 $g$ 平坦，或者 $\omega_{Y/X}$ 可逆，
则 $R' = R \times_Y Y'$。
\end{lemma}

\begin{proof}
立即归结为 $X$、$X'$、$Y$、$Y'$ 都仿射的情形。换言之，
有 Noetherian 环的余笛卡尔图
$$
\xymatrix{
B' & B \ar[l] \\
A' \ar[u] & A \ar[l] \ar[u]
}
$$
其中 $A \to B$ 平坦且拟有限。基变换映射
$\omega_{B/A} \otimes_B B' \to \omega_{B'/A'}$ 是同构
（引理 \ref{discriminant-lemma-dualizing-base-change-of-flat}），并把迹元
$\tau_{B/A}$ 映到迹元 $\tau_{B'/A'}$
（引理 \ref{discriminant-lemma-trace-base-change}）。因此有限 $B$-模
$Q = \Coker(\tau_{B/A} : B \to \omega_{B/A})$ 满足
$Q \otimes_B B' = \Coker(\tau_{B'/A'} : B' \to \omega_{B'/A'})$。
所以 $\mathfrak{D}_{B/A}B' \subset \mathfrak{D}_{B'/A'}$，
即得到闭浸入 $R' \to R \times_Y Y'$。由于
$R = \text{Supp}(Q)$ 且 $R' = \text{Supp}(Q \otimes_B B')$
（《代数》，引理 \ref{algebra-lemma-support-closed}），由《代数》，引理
\ref{algebra-lemma-support-base-change}，
$R' \to R \times_Y Y'$ 为双射。若 $B \to B'$ 平坦，例如
$A \to A'$ 平坦，则等式
$\mathfrak{D}_{B/A}B' = \mathfrak{D}_{B'/A'}$ 成立；见《代数》，引理
\ref{algebra-lemma-annihilator-flat-base-change}。
最后，若 $\omega_{B/A}$ 可逆，则可局部化并设
$\omega_{B/A} = B \lambda$。写成 $\tau_{B/A} = b\lambda$，可见
$Q = B/bB$ 且 $\mathfrak{D}_{B/A} = bB$。在 $B'$ 上作同样论证，
得到 $\mathfrak{D}_{B'/A'} = bB'$，引理得证。
\end{proof}

\begin{lemma}
\label{discriminant-lemma-norm-different-in-discriminant}
设 $f : Y \to X$ 是 Noetherian 概形的有限平坦态射。则
$\text{Norm}_f : f_*\mathcal{O}_Y \to \mathcal{O}_X$ 把
$f_*\mathfrak{D}_f$ 映入判别式 $D_f$ 的理想层。
\end{lemma}

\begin{proof}
范数映射构造于《除子》，引理
\ref{divisors-lemma-finite-locally-free-has-norm}，而 $f$ 的判别式
构造于第 \ref{discriminant-section-discriminant} 节。问题是仿射局部的，故可设
$X = \Spec(A)$、$Y = \Spec(B)$，且 $f$ 由有限局部自由环同态
$A \to B$ 给出。进一步局部化后，可设 $B$ 作为 $A$-模是有限自由的。
选取 $B$ 作为 $A$-模的一组基 $b_1, \ldots, b_n \in B$，并将
$\omega_{B/A} = \Hom_A(B, A)$ 作为 $A$-模的对偶基记为
$b_1^\vee, \ldots, b_n^\vee$。由于 $b$ 的范数是 $A$-线性映射
$b : B \to B$ 的行列式，故
$\text{Norm}_{B/A}(b) = \det(b_i^\vee(bb_j))$。
判别式是 $\Spec(A)$ 中由 $\det(\text{Trace}_{B/A}(b_ib_j))$
定义的主闭子概形。若 $b \in \mathfrak{D}_{B/A}$，则存在
$c_i \in B$，使得
$b \cdot b_i^\vee = c_i \cdot \text{Trace}_{B/A}$；这里用圆点表示
$\omega_{B/A}$ 上的 $B$-模结构。写成 $c_i = \sum a_{il} b_l$。则
\begin{align*}
\text{Norm}_{B/A}(b)
& =
\det(b_i^\vee(bb_j)) \\
& =
\det( (b \cdot b_i^\vee)(b_j)) \\
& =
\det((c_i \cdot \text{Trace}_{B/A})(b_j)) \\
& =
\det(\text{Trace}_{B/A}(c_ib_j)) \\
& =
\det(a_{il}) \det(\text{Trace}_{B/A}(b_l b_j))
\end{align*}
这就证明了引理。
\end{proof}

\begin{lemma}
\label{discriminant-lemma-different-ramification}
设 $f : Y \to X$ 是 Noetherian 概形的平坦拟有限态射。
由不同理想 $\mathfrak{D}_f$ 定义的闭子概形 $R \subset Y$
恰为 $f$ 非 \'etale（等价地，分歧）的点集。
\end{lemma}

\begin{proof}
由于 $f$ 为有限表现且平坦，它在一点为 \'etale，当且仅当在该点非分歧。
此外，分歧点轨迹的形成与基变换可交换。见《态射》，第
\ref{morphisms-section-etale} 节，特别是《态射》，引理
\ref{morphisms-lemma-set-points-where-fibres-etale}。由引理
\ref{discriminant-lemma-base-change-different}，$R$ 的形成在集合意义下与基变换可交换。
故只需在 $X$ 是某个域的谱时证明。另一方面，
$(\omega_{Y/X}, \tau_{Y/X})$ 的构造在 $Y$ 上是局部的。
由于 $Y$ 在域上拟有限，所以是有限离散空间，故可设 $Y$ 只有一个点。

\medskip\noindent
设 $X = \Spec(k)$、$Y = \Spec(B)$，其中 $k$ 是域，而 $B$ 是
有限局部 $k$-代数。若 $Y \to X$ 为 \'etale，则 $B$ 是 $k$
的有限可分扩张。由《域论》，引理
\ref{fields-lemma-separable-trace-pairing}，迹元
$\text{Trace}_{B/k}$ 是 $\omega_{B/k}$ 的一个基元素。因此此时
$\mathfrak{D}_{B/k} = B$。反之，若 $\mathfrak{D}_{B/k} = B$，
则由引理 \ref{discriminant-lemma-norm-different-in-discriminant} 以及 $1$ 的范数为
$1$ 这一事实可知，判别式为空。因此由引理
\ref{discriminant-lemma-discriminant}，$Y \to X$ 为 \'etale。
\end{proof}

\begin{lemma}
\label{discriminant-lemma-norm-different-is-discriminant}
设 $f : Y \to X$ 是 Noetherian 概形的平坦拟有限态射，并令
$R \subset Y$ 为 $\mathfrak{D}_f$ 所定义的闭子概形。
\begin{enumerate}
\item 若 $\omega_{Y/X}$ 可逆，则 $R$ 是 $Y$ 的局部主闭子概形。
\item 若 $\omega_{Y/X}$ 可逆且 $f$ 有限，则 $R$ 的范数是 $f$
的判别式 $D_f$。
\item 若 $\omega_{Y/X}$ 可逆，且 $f$ 在 $Y$ 的伴随点处为 \'etale，
则 $R$ 是有效 Cartier 除子，并存在同构
$\mathcal{O}_Y(R) = \omega_{Y/X}$。
\end{enumerate}
\end{lemma}

\begin{proof}
(1) 的证明。可以在 $Y$ 上局部工作，故可设 $\omega_{Y/X}$ 是秩为
$1$ 的自由模。设 $\omega_{Y/X} = \mathcal{O}_Y\lambda$。
写成 $\tau_{Y/X} = h \lambda$，则 $R$ 由 $h$ 定义，即 $R$ 局部为主的。

\medskip\noindent
(2) 的证明。可设 $Y \to X$ 由有限自由环同态 $A \to B$ 给出，
且 $\omega_{B/A}$ 作为 $B$-模是秩为 $1$ 的自由模。选取
$\omega_{B/A}$ 的一个 $B$-基元素 $\lambda$，并对某个 $b \in B$ 写成
$\text{Trace}_{B/A} = b \cdot \lambda$。则
$\mathfrak{D}_{B/A} = (b)$，而 $D_f$ 由
$\det(\text{Trace}_{B/A}(b_ib_j))$ 定义，其中
$b_1, \ldots, b_n$ 是 $B$ 作为 $A$-模的一组基。令
$b_1^\vee, \ldots, b_n^\vee$ 为对偶基。写成
$b_i^\vee = c_i \cdot \lambda$，可见 $c_1, \ldots, c_n$
也是 $B$ 的一组基。因此，当 $c_i = \sum a_{il}b_l$ 时，
$\det(a_{il})$ 是 $A$ 中的单位。显然
$b \cdot b_i^\vee = c_i \cdot \text{Trace}_{B/A}$，
故由引理 \ref{discriminant-lemma-norm-different-in-discriminant} 证明中的计算可知，
$\text{Norm}_{B/A}(b)$ 等于 $\det(\text{Trace}_{B/A}(b_ib_j))$
乘以一个单位。

\medskip\noindent
(3) 的证明。在上述记号下，由引理 \ref{discriminant-lemma-different-ramification}
和假设可知，$h$ 在 $Y$ 的伴随点处不消失，这蕴含 $h$ 是非零因子。
典范同构把 $1$ 映到 $\tau_{Y/X}$；见《除子》，引理
\ref{divisors-lemma-characterize-OD}。
\end{proof}







\section{拟有限 syntomic 态射}
\label{discriminant-section-quasi-finite-syntomic}

\noindent
本节讨论拟有限 syntomic 态射具有可逆相对对偶化模这一事实。

\begin{lemma}
\label{discriminant-lemma-syntomic-quasi-finite}
设 $f : Y \to X$ 是概形态射。下列条件等价：
\begin{enumerate}
\item $f$ 局部拟有限且 syntomic；
\item $f$ 局部拟有限、平坦且为局部完全交态射；
\item $f$ 局部拟有限、平坦、局部有限表现，且 $f$ 的纤维为局部完全交；
\item $f$ 局部拟有限，并且对每个 $y \in Y$，存在仿射开集
$y \in V = \Spec(B) \subset Y$、$U = \Spec(A) \subset X$，
满足 $f(V) \subset U$，还存在整数 $n$ 以及
$h, f_1, \ldots, f_n \in A[x_1, \ldots, x_n]$，使得
$B = A[x_1, \ldots, x_n, 1/h]/(f_1, \ldots, f_n)$；
\item 对每个 $y \in Y$，存在仿射开集
$y \in V = \Spec(B) \subset Y$、$U = \Spec(A) \subset X$，
满足 $f(V) \subset U$，且 $A \to B$ 是形如
$B = A[x_1, \ldots, x_n]/(f_1, \ldots, f_n)$ 的相对整体完全交；
\item $f$ 局部拟有限、平坦、局部有限表现，且 $\NL_{Y/X}$
的 Tor 振幅在 $[-1, 0]$ 中；
\item $f$ 平坦、局部有限表现，$\NL_{Y/X}$ 是秩为 $0$、
Tor 振幅在 $[-1, 0]$ 中的完美复形。
\end{enumerate}
\end{lemma}

\begin{proof}
(1) 与 (2) 的等价性是《态射进阶》，引理
\ref{more-morphisms-lemma-flat-lci}。(1) 与 (3) 的等价性是《态射》，引理
\ref{morphisms-lemma-syntomic-flat-fibres}。

\medskip\noindent
若 $A \to B$ 如 (4) 所述，则
$B = A[x, x_1, \ldots, x_n]/(xh - 1, f_1, \ldots, f_n)$
是相对整体完全交；见《代数》，定义
\ref{algebra-definition-relative-global-complete-intersection}。
因此 (4) 蕴含 (5)，而 (5) 蕴含 (4) 是显然的。

\medskip\noindent
条件 (5) 蕴含 (1)：由《代数》，引理
\ref{algebra-lemma-relative-global-complete-intersection}，
相对整体完全交是 syntomic 的；而相对整体完全交的定义保证，
含 $n$ 个变量和 $n$ 个方程的相对整体完全交是拟有限的；见《代数》，定义
\ref{algebra-definition-relative-global-complete-intersection} 和引理
\ref{algebra-lemma-isolated-point-fibre}。

\medskip\noindent
《代数》，引理 \ref{algebra-lemma-syntomic} 或《态射》，引理
\ref{morphisms-lemma-syntomic-locally-standard-syntomic}
均说明 (1) 蕴含 (5)。

\medskip\noindent
《态射进阶》，引理 \ref{more-morphisms-lemma-flat-fp-NL-lci}
说明 (6) 与 (1) 等价。若等价条件 (1)--(6) 成立，则在仿射局部，
$Y \to X$ 由相对整体完全交
$B = A[x_1, \ldots, x_n]/(f_1, \ldots, f_n)$ 给出，
其变量数与方程数相同。利用这个表现可见
$$
\NL_{B/A} =\left(
(f_1, \ldots, f_n)/(f_1, \ldots, f_n)^2
\longrightarrow
\bigoplus\nolimits_{i = 1, \ldots, n} B \text{d} x_i\right)
$$
由《代数》，引理
\ref{algebra-lemma-relative-global-complete-intersection-conormal}，
模 $(f_1, \ldots, f_n)/(f_1, \ldots, f_n)^2$ 是自由的，
其生成元为元素 $f_1, \ldots, f_n$ 的同余类。因此
$\NL_{B/A}$ 的秩为 $0$，$\NL_{Y/X}$ 亦然。于是 (1)--(6) 蕴含 (7)。

\medskip\noindent
最后假设 (7)。由《态射进阶》，引理
\ref{more-morphisms-lemma-flat-fp-NL-lci}，$f$ 是 syntomic 的。
因此在适当的仿射开集上，$f$ 由相对整体完全交
$A \to B = A[x_1, \ldots, x_n]/(f_1, \ldots, f_m)$ 给出；见《态射》，引理
\ref{morphisms-lemma-syntomic-locally-standard-syntomic}。
完全仿照上述论证，$\NL_{B/A}$ 是秩为 $n - m$ 的完美复形。
因此 $n = m$，故 (5) 成立。证明完成。
\end{proof}

\begin{lemma}
\label{discriminant-lemma-characterize-invertible}
相对对偶化模的可逆性。
\begin{enumerate}
\item 若 $A \to B$ 是 Noetherian 环的拟有限平坦同态，则
$\omega_{B/A}$ 是可逆 $B$-模，当且仅当对所有素理想
$\mathfrak p \subset A$，$\omega_{B \otimes_A \kappa(\mathfrak p)/\kappa(\mathfrak p)}$
是可逆 $B \otimes_A \kappa(\mathfrak p)$-模。
\item 若 $Y \to X$ 是 Noetherian 概形的拟有限平坦态射，则
$\omega_{Y/X}$ 可逆，当且仅当对所有 $x \in X$，
$\omega_{Y_x/x}$ 可逆。
\end{enumerate}
\end{lemma}

\begin{proof}
(1) 的证明。由于 $A \to B$ 平坦，模 $\omega_{B/A}$ 为 $A$-平坦；
见引理 \ref{discriminant-lemma-dualizing-base-flat-flat}。因此，
$\omega_{B/A}$ 是可逆 $B$-模，当且仅当对每个素理想
$\mathfrak p \subset A$，$\omega_{B/A} \otimes_A \kappa(\mathfrak p)$
是可逆 $B \otimes_A \kappa(\mathfrak p)$-模；见《态射进阶》，引理
\ref{more-morphisms-lemma-flat-and-free-at-point-fibre}。
仍利用 $A \to B$ 平坦这一点，$\omega_{B/A}$ 的形成与基变换可交换；
见引理 \ref{discriminant-lemma-dualizing-base-change-of-flat}。因此在平坦性成立时，
相对对偶化模可逆，等价于对映射
$\kappa(\mathfrak p) \to B \otimes_A \kappa(\mathfrak p)$，
其相对对偶化模可逆。

\medskip\noindent
(2) 由 (1) 以及如下事实得出：在仿射局部，对偶化模由其代数对应物给出；
见注记 \ref{discriminant-remark-relative-dualizing-for-quasi-finite}。
\end{proof}

\begin{lemma}
\label{discriminant-lemma-dim-zero-global-complete-intersection-over-field}
设 $k$ 是域，并设 $B = k[x_1, \ldots, x_n]/(f_1, \ldots, f_n)$
是 $k$ 上维数为 $0$ 的整体完全交。则 $\omega_{B/k}$ 可逆。
\end{lemma}

\begin{proof}
由 Noether 正规化，见《代数》，引理
\ref{algebra-lemma-Noether-normalization}，存在有限单射 $k \to B$，
即 $\dim_k(B) < \infty$。因此作为 $B$-模有
$\omega_{B/k} = \Hom_k(B, k)$。由《对偶化复形》，引理
\ref{dualizing-lemma-dualizing-finite}，$R\Hom(B, k)$ 是 $B$ 的对偶复形；
又由《对偶化复形》，引理 \ref{dualizing-lemma-RHom-ext}，
$R\Hom(B, k)$ 等于放在次数 $0$ 的 $\omega_{B/k}$。
故只需证明 $B$ 为 Gorenstein（《对偶化复形》，引理
\ref{dualizing-lemma-gorenstein}）。这由《对偶化复形》，引理
\ref{dualizing-lemma-gorenstein-lci} 成立。
\end{proof}

\begin{lemma}
\label{discriminant-lemma-dualizing-syntomic-quasi-finite}
设 $f : Y \to X$ 是局部 Noetherian 概形的态射。若 $f$ 满足引理
\ref{discriminant-lemma-syntomic-quasi-finite} 的等价条件，则 $\omega_{Y/X}$
是可逆 $\mathcal{O}_Y$-模。
\end{lemma}

\begin{proof}
可设 $A \to B$ 是形如
$B = A[x_1, \ldots, x_n]/(f_1, \ldots, f_n)$ 的相对整体完全交，
而需证明 $\omega_{B/A}$ 可逆。这由结合引理
\ref{discriminant-lemma-characterize-invertible} 与
\ref{discriminant-lemma-dim-zero-global-complete-intersection-over-field} 得出。
\end{proof}

\begin{example}
\label{discriminant-example-universal-quasi-finite-syntomic}
设 $n \geq 1$ 与 $d \geq 1$ 为整数。令 $T$ 为满足
$e_i \geq 0$ 且 $\sum e_i \leq d$ 的多重指标
$E = (e_1, \ldots, e_n)$ 所成的集合。考虑环
$$
A = \mathbf{Z}[a_{i, E} ; 1 \leq i \leq n, E \in T]
$$
在 $A[x_1, \ldots, x_n]$ 中考虑元素
$f_i = \sum_{E \in T} a_{i, E} x^E$，其中按惯例
$x^E = x_1^{e_1} \ldots x_n^{e_n}$。考虑 $A$-代数
$$
B = A[x_1, \ldots, x_n]/(f_1, \ldots, f_n)
$$
记 $X_{n, d} = \Spec(A)$，并令 $Y_{n, d} \subset \Spec(B)$
为使态射 $\Spec(B) \to \Spec(A) = X_{n, d}$ 的限制拟有限的最大开子概形；
见《代数》，引理 \ref{algebra-lemma-quasi-finite-open}。
\end{example}

\begin{lemma}
\label{discriminant-lemma-universal-quasi-finite-syntomic-etale}
沿用例 \ref{discriminant-example-universal-quasi-finite-syntomic} 的记号，
概形 $X_{n, d}$ 与 $Y_{n, d}$ 都正则且不可约，态射
$Y_{n, d} \to X_{n, d}$ 局部拟有限且 syntomic，并存在稠密开子概形
$V \subset Y_{n, d}$，使 $Y_{n, d} \to X_{n, d}$ 在其上限制为
\'etale 态射 $V \to X_{n, d}$。
\end{lemma}

\begin{proof}
概形 $X_{n, d}$ 是多项式环 $A$ 的谱。因此 $X_{n, d}$ 正则且不可约。由于可写成
$$
f_i = a_{i, (0, \ldots, 0)} +
\sum\nolimits_{E \in T, E \not = (0, \ldots, 0)} a_{i, E} x^E
$$
环 $B$ 同构于以 $x_1, \ldots, x_n$ 以及所有满足
$E \not = (0, \ldots, 0)$ 的 $a_{i, E}$ 为变量的多项式环。
因此 $\Spec(B)$ 是不可约正则概形，其开集 $Y_{n, d}$ 亦然。
由引理 \ref{discriminant-lemma-syntomic-quasi-finite}，态射
$Y_{n, d} \to X_{n, d}$ 局部拟有限且 syntomic。为找到 $V$，
只需找到一点，使 $Y_{n, d} \to X_{n, d}$ 在该点为 \'etale
（按定义，态射为 \'etale 的点轨迹为开集）。因此只需找到
$X_{n, d}$ 的一点，使 $Y_{n, d} \to X_{n, d}$ 的纤维非空且为
\'etale；见《态射》，引理 \ref{morphisms-lemma-etale-at-point}。
选择对应于环同态 $\chi : A \to \mathbf{Q}$ 的点，该同态把
$f_i$ 映到 $x_i^d - 1$。于是
$$
B \otimes_{A, \chi} \mathbf{Q} =
\mathbf{Q}[x_1, \ldots, x_n]/(x_1^d - 1, \ldots, x_n^d - 1)
$$
这是 $\mathbf{Q}$ 上的非零 \'etale 代数。
\end{proof}

\begin{lemma}
\label{discriminant-lemma-locally-comes-from-universal}
设 $f : Y \to X$ 是概形态射。若 $f$ 满足引理
\ref{discriminant-lemma-syntomic-quasi-finite} 的等价条件，则对每个 $y \in Y$，
存在 $n, d$ 以及交换图
$$
\xymatrix{
Y \ar[d] &
V \ar[d] \ar[l] \ar[r] &
Y_{n, d} \ar[d] \\
X & U \ar[l] \ar[r] &
X_{n, d}
}
$$
其中 $U \subset X$ 与 $V \subset Y$ 为满足 $y \in V$ 的开集，
$Y_{n, d} \to X_{n, d}$ 如例
\ref{discriminant-example-universal-quasi-finite-syntomic} 所述，而右侧方块为笛卡尔图。
\end{lemma}

\begin{proof}
由引理 \ref{discriminant-lemma-syntomic-quasi-finite}，可选取仿射的 $U$ 与 $V$，
使 $U = \Spec(R)$、$V = \Spec(S)$，其中
$S = R[y_1, \ldots, y_n]/(g_1, \ldots, g_n)$。
沿用例 \ref{discriminant-example-universal-quasi-finite-syntomic} 的记号，
若选取充分大的 $d$，则每个 $g_i$ 可写成
$g_i = \sum_{E \in T} g_{i, E}y^E$，其中 $g_{i, E} \in R$。
将 $a_{i, E}$ 映到 $g_{i, E}$ 的映射 $A \to R$，以及将
$x_i \to y_i$ 的映射 $B \to S$，给出环的余笛卡尔图
$$
\xymatrix{
S & B \ar[l] \\
R \ar[u] & A \ar[l] \ar[u]
}
$$
引理得证。
\end{proof}











\section{有限 syntomic 态射}
\label{discriminant-section-finite-syntomic}

\noindent
本节给出第 \ref{discriminant-section-quasi-finite-syntomic} 节关于有限 syntomic 态射的对应版本。

\begin{lemma}
\label{discriminant-lemma-syntomic-finite}
设 $f : Y \to X$ 是概形态射。下列条件等价：
\begin{enumerate}
\item $f$ 有限且 syntomic；
\item $f$ 有限、平坦，且为局部完全交态射；
\item $f$ 有限、平坦、局部有限表现，且 $f$ 的纤维为局部完全交；
\item $f$ 有限，并且对每个 $x \in X$，存在仿射开集
$x \in U = \Spec(A) \subset X$、整数 $n$ 以及
$f_1, \ldots, f_n \in A[x_1, \ldots, x_n]$，使得
$f^{-1}(U)$ 同构于
$A[x_1, \ldots, x_n]/(f_1, \ldots, f_n)$ 的谱；
\item $f$ 有限、平坦、局部有限表现，且 $\NL_{X/Y}$ 的 Tor 振幅在 $[-1, 0]$ 中；
\item $f$ 有限、平坦、局部有限表现，且 $\NL_{X/Y}$ 是秩为 $0$、
Tor 振幅在 $[-1, 0]$ 中的完美复形。
\end{enumerate}
\end{lemma}

\begin{proof}
(1)、(2)、(3)、(5)、(6) 的等价性以及 (4) $\Rightarrow$ (1)，
均立即由引理 \ref{discriminant-lemma-syntomic-quasi-finite} 得出。
假设等价条件 (1)、(2)、(3)、(5)、(6) 成立。
取一点 $x \in X$ 以及 $X$ 中包含 $x$ 的仿射开集 $U = \Spec(A)$，
并设 $x$ 对应于素理想 $\mathfrak p \subset A$。写成
$f^{-1}(U) = \Spec(B)$ 及 $B = A[x_1, \ldots, x_n]/I$。
由 (6)，$\NL_{B/A}$ 是 Tor 振幅在 $[-1, 0]$ 中的完美复形，
故 $I/I^2$ 是秩为 $n$ 的有限局部自由 $B$-模。
由于 $B_\mathfrak p$ 是半局部环，$(I/I^2)_\mathfrak p$
是秩为 $n$ 的自由模；见《代数》，引理
\ref{algebra-lemma-locally-free-semi-local-free}。
因此用不属于 $\mathfrak p$ 的某个元素作主局部化并以所得环替换 $A$ 后，
可设 $I/I^2$ 是秩为 $n$ 的自由 $B$-模。
于是由《代数》，引理 \ref{algebra-lemma-huber}，
可找到 $B$ 在 $A$ 上变量数与方程数相同的一个表现。换言之，可设
$B = A[x_1, \ldots, x_n]/(f_1, \ldots, f_n)$。这证明了 (4)。
\end{proof}

\begin{example}
\label{discriminant-example-universal-finite-syntomic}
设 $d \geq 1$ 为整数。考虑变量 $a_{ij}^l$，其中
$1 \leq i, j, l \leq d$，并记
$$
A_d = \mathbf{Z}[a_{ij}^k]/J
$$
其中 $J$ 是由下列元素生成的理想：
$$
\left\{
\begin{matrix}
\sum_l a_{ij}^la_{lk}^m - \sum_l a_{il}^ma_{jk}^l & \forall i, j, k, m \\
a_{ij}^k - a_{ji}^k & \forall i, j, k \\
a_{i1}^j - \delta_{ij} & \forall i, j
\end{matrix}
\right.
$$
其中 $\delta_{ij}$ 表示 Kronecker delta 函数。如下定义一个
$A_d$-代数 $B_d$：作为 $A_d$-模，令
$$
B_d = A_d e_1 \oplus \ldots \oplus A_d e_d
$$
代数结构由把 $1$ 映到 $e_1$ 的映射 $A_d \to B_d$ 给出。
$B_d$ 上的乘法是 $A_d$-双线性映射
$$
m : B_d \times B_d \longrightarrow B_d, \quad
m(e_i, e_j) = \sum a_{ij}^k e_k
$$
直接验证可知，上述关系恰好保证这给出一个 $A_d$-代数结构。态射
$$
\pi_d : Y_d = \Spec(B_d) \longrightarrow \Spec(A_d) = X_d
$$
是秩为 $d$ 的“泛”有限自由态射。
\end{example}

\begin{lemma}
\label{discriminant-lemma-universal-finite-syntomic}
沿用例 \ref{discriminant-example-universal-finite-syntomic} 的记号，存在开子概形
$U_d \subset X_d$，具有如下性质：概形态射 $X \to X_d$ 通过
$U_d$ 分解，当且仅当 $Y_d \times_{X_d} X \to X$ 为 syntomic。
\end{lemma}

\begin{proof}
回忆：syntomic 等价于平坦且为局部完全交态射；见《态射进阶》，引理
\ref{more-morphisms-lemma-flat-lci}。令 $W_d \subset Y_d$ 为
$\pi_d$ 为 Koszul 的点集。它在 $Y_d$ 中是开集，而且其形成与任意基变换可交换；
见《态射进阶》，引理
\ref{more-morphisms-lemma-base-change-lci-fibres}。
由于 $\pi_d$ 有限，从而是闭态射，故 $Z = \pi_d(Y_d \setminus W_d)$
为闭集。显然 $U_d = X_d \setminus Z$，而且其形成与基变换可交换，
故引理成立。
\end{proof}

\begin{lemma}
\label{discriminant-lemma-universal-finite-syntomic-smooth}
沿用例 \ref{discriminant-example-universal-finite-syntomic} 的记号，并令 $U_d$
如引理 \ref{discriminant-lemma-universal-finite-syntomic} 所述，则 $U_d$
在 $\Spec(\mathbf{Z})$ 上光滑。
\end{lemma}

\begin{proof}
利用《态射进阶》，引理
\ref{more-morphisms-lemma-lifting-along-artinian-at-point} 来证明
$U_d \to \Spec(\mathbf{Z})$ 光滑。具体地，设有态射
$\Spec(A) \to U_d$，且 $A' \to A$ 是小扩张。则
$B = A \otimes_{A_d} B_d$ 是在 $A$ 上 syntomic 的有限自由
$A$-代数（由 $U_d$ 的构造）。由《环映射的光滑化》，命题
\ref{smoothing-proposition-lift-smooth}，存在 syntomic 环同态
$A' \to B'$，使得 $B \cong B' \otimes_{A'} A$。
令 $e'_1 = 1 \in B'$。对 $1 < i \leq d$，选取元素
$1 \otimes e_i \in A \otimes_{A_d} B_d = B$ 的提升 $e'_i \in B'$。
于是 $e'_1, \ldots, e'_d$ 是 $B'$ 在 $A'$ 上的一组基
（例如见《代数》，引理
\ref{algebra-lemma-local-artinian-basis-when-flat}）。
因此可唯一地写成 $e'_i e'_j = \sum \alpha_{ij}^l e'_l$，其中
$\alpha_{ij}^l \in A'$ 满足 $A'$ 中的关系
$\sum_l \alpha_{ij}^l \alpha_{lk}^m = \sum_l \alpha_{il}^m \alpha _{jk}^l$、
$\alpha_{ij}^k = \alpha_{ji}^k$ 以及 $\alpha_{i1}^j = \delta_{ij}$。
把 $a_{ij}^l \in A_d$ 映到 $\alpha_{ij}^l \in A'$，便确定了态射
$\Spec(A') \to X_d$。该态射与给定态射 $\Spec(A) \to U_d$ 相容。
由于 $\Spec(A')$ 与 $\Spec(A)$ 具有相同的底层拓扑空间，
由此得到所需提升 $\Spec(A') \to U_d$，并得出 $U_d$ 在
$\mathbf{Z}$ 上光滑。
\end{proof}

\begin{lemma}
\label{discriminant-lemma-universal-finite-syntomic-etale}
沿用例 \ref{discriminant-example-universal-finite-syntomic} 的记号，考虑使
$\pi_d$ 为 \'etale 的开子概形 $U'_d \subset X_d$。则 $U'_d$
是引理 \ref{discriminant-lemma-universal-finite-syntomic} 中开集 $U_d$ 的稠密子集。
\end{lemma}

\begin{proof}
完全仿照引理 \ref{discriminant-lemma-universal-finite-syntomic} 的证明，并应用《态射》，引理
\ref{morphisms-lemma-set-points-where-fibres-etale}，可得一个最大的开集
$U'_d \subset X_d$，使得 $\pi_d$ 在其上为 \'etale。此外，由于
\'etale 态射是 syntomic 的，有 $U'_d \subset U_d$。
为完成证明，需证 $U'_d \subset U_d$ 稠密。设
$u : \Spec(k) \to U_d$ 为态射，其中 $k$ 是域。令
$B = k \otimes_{A_d} B_d$，如引理
\ref{discriminant-lemma-universal-finite-syntomic-smooth} 的证明所述。
我们将证明：存在剩余域为 $k$ 的局部整环 $A'$，以及有限 syntomic
$A'$-代数 $B'$，满足 $B = k \otimes_{A'} B'$，且其泛纤维为 \'etale。
完全仿照上一段，这将确定态射 $\Spec(A') \to U_d$，它把泛点映入
$U'_d$，把闭点映到 $u$，从而完成证明。

\medskip\noindent
由引理 \ref{discriminant-lemma-syntomic-finite} 的 (4)，可选取表现
$B = k[x_1, \ldots, x_n]/(f_1, \ldots, f_n)$。
令 $d'$ 为多项式 $f_1, \ldots, f_n$ 的最大总次数，并令
$Y_{n, d'} \to X_{n, d'}$ 如例
\ref{discriminant-example-universal-quasi-finite-syntomic} 所述。由构造，存在态射
$u' : \Spec(k) \to X_{n, d'}$，使得
$$
\Spec(B) \cong Y_{n, d'} \times_{X_{n, d'}, u'} \Spec(k)
$$
记 $A = \mathcal{O}_{X_{n, d'}, u'}^h$ 为 $X_{n, d'}$ 在 $u'$ 的像处的
局部环之 Hensel 化。则可写成
$$
Y_{n, d'} \times_{X_{n, d'}} \Spec(A) = Z \amalg W
$$
其中 $Z \to \Spec(A)$ 有限，而 $W \to \Spec(A)$ 的闭纤维为空；
见《代数》，引理 \ref{algebra-lemma-characterize-henselian} 的 (13)，
或《态射进阶》，第 \ref{more-morphisms-section-etale-localization} 节中的讨论。
由引理 \ref{discriminant-lemma-universal-quasi-finite-syntomic-etale}，局部环 $A$ 正则
（这里还使用《代数进阶》，引理
\ref{more-algebra-lemma-henselization-regular}），且态射
$Z \to \Spec(A)$ 在 $\Spec(A)$ 的泛点上为 \'etale
（因为它映到 $X_{d, n'}$ 的泛点）。由构造，
$Z \times_{\Spec(A)} \Spec(k) \cong \Spec(B)$。
这已证明所需结论，唯一的例外是 $A$ 的剩余域到 $k$ 的映射未必为同构。
由《代数》，引理 \ref{algebra-lemma-flat-local-given-residue-field}，
存在平坦局部环同态 $A \to A'$，使 $A'$ 的剩余域为 $k$。
若 $A'$ 不是整环，则取极小素理想 $\mathfrak p \subset A'$；
由平坦性，它位于 $A$ 的唯一极小素理想之上。然后以 $A'/\mathfrak p$
替换 $A'$。令 $B'$ 为唯一的 $A'$-代数，使得
$Z \times_{\Spec(A)} \Spec(A') = \Spec(B')$。证明完成。
\end{proof}

\begin{remark}
\label{discriminant-remark-universal-finite-syntomic-smooth-top}
设 $\pi_d : Y_d \to X_d$ 如例
\ref{discriminant-example-universal-finite-syntomic} 所述。令 $U_d \subset X_d$
为使 $V_d = \pi_d^{-1}(U_d)$ 如引理
\ref{discriminant-lemma-universal-finite-syntomic} 中那样有限 syntomic 的最大开集。
则 $V_d$ 在 $\mathbf{Z}$ 上也光滑。
（当然，当 $d \geq 2$ 时，态射 $V_d \to U_d$ 并不光滑。）
仿照引理 \ref{discriminant-lemma-universal-finite-syntomic-smooth} 的证明，
这对应于如下形变问题：给定小扩张 $C' \to C$，以及带有截面
$B \to C$ 的有限 syntomic $C$-代数 $B$，求一个有限 syntomic
$C'$-代数 $B'$ 及截面 $B' \to C'$，使其与 $C$ 的张量积恢复
$B \to C$。由引理 \ref{discriminant-lemma-syntomic-finite}，可把
$B = C[x_1, \ldots, x_n]/(f_1, \ldots, f_n)$ 写成相对整体完全交。
作坐标变换后，可设 $x_1, \ldots, x_n$ 都在 $B \to C$ 的核中，
于是多项式 $f_i$ 的常数项为零。任取 $f_i$ 的提升
$f'_i \in C'[x_1, \ldots, x_n]$，并使其常数项为零。则
$B' = C'[x_1, \ldots, x_n]/(f'_1, \ldots, f'_n)$ 配上把 $x_i$
映到零的截面 $B' \to C'$ 即满足要求。
\end{remark}

\begin{lemma}
\label{discriminant-lemma-locally-comes-from-universal-finite}
设 $f : Y \to X$ 是概形态射。若 $f$ 满足引理
\ref{discriminant-lemma-syntomic-finite} 的等价条件，则对每个 $x \in X$，
存在 $d$ 以及交换图
$$
\xymatrix{
Y \ar[d] &
V \ar[d] \ar[l] \ar[r] &
V_d \ar[d] \ar[r] &
Y_d \ar[d]^{\pi_d}\\
X &
U \ar[l] \ar[r] &
U_d \ar[r] &
X_d
}
$$
具有如下性质：
\begin{enumerate}
\item $U \subset X$ 为开集，$x \in U$，且 $V = f^{-1}(U)$；
\item $\pi_d : Y_d \to X_d$ 如例
\ref{discriminant-example-universal-finite-syntomic} 所述；
\item $U_d \subset X_d$ 如引理 \ref{discriminant-lemma-universal-finite-syntomic} 所述，
且 $V_d = \pi_d^{-1}(U_d) \subset Y_d$；
\item 中间方块为笛卡尔图。
\end{enumerate}
\end{lemma}

\begin{proof}
选取 $x$ 的仿射开邻域 $U = \Spec(A) \subset X$，并写成
$V = f^{-1}(U) = \Spec(B)$。则 $B$ 是有限局部自由 $A$-模，
而包含关系 $A \subset B$ 在局部上是直和项。缩小 $U$ 后，
可选取 $B$ 作为 $A$-模的一组基 $1 = e_1, e_2, \ldots, e_d$。
唯一地写成 $e_i e_j = \sum \alpha_{ij}^l e_l$，其中
$\alpha_{ij}^l \in A$ 满足 $A$ 中的关系
$\sum_l \alpha_{ij}^l \alpha_{lk}^m = \sum_l \alpha_{il}^m \alpha _{jk}^l$、
$\alpha_{ij}^k = \alpha_{ji}^k$ 以及 $\alpha_{i1}^j = \delta_{ij}$。
把 $a_{ij}^l \in A_d$ 映到 $\alpha_{ij}^l \in A$，便确定态射
$\Spec(A) \to X_d$。由构造，$V \cong \Spec(A) \times_{X_d} Y_d$。
由 $U_d$ 的定义，$\Spec(A) \to X_d$ 通过 $U_d$ 分解。证明完成。
\end{proof}










\section{不同理想的一个公式}
\label{discriminant-section-formula-different}

\noindent
本节讨论 \cite[Appendix A]{Mazur-Roberts} 附录 A 中 Tate 的材料。
用本章的语言来说，这将证明：对平坦、拟有限的局部完全交态射，
不同理想等于 K\"ahler 不同理想。首先在一个特殊情形中计算 Noether 不同理想。

\begin{lemma}
\label{discriminant-lemma-tate}
\begin{reference}
\cite[Appendix]{Mazur-Roberts}
\end{reference}
设 $A \to P$ 是环同态，且 $f_1, \ldots, f_n \in P$ 是 Koszul 正则序列。
假设 $B = P/(f_1, \ldots, f_n)$ 在 $A$ 上平坦。设
$g_1, \ldots, g_n \in P \otimes_A B$ 是生成乘法映射
$P \otimes_A B \to B$ 之核的 Koszul 正则序列。写成
$f_i \otimes 1 = \sum g_{ij} g_j$。则
$\Ker(B \otimes_A B \to B)$ 的零化子是由 $\det(g_{ij})$ 的像生成的主理想。
\end{lemma}

\begin{proof}
Koszul 复形 $K_\bullet = K(P, f_1, \ldots, f_n)$ 是由有限自由
$P$-模组成的 $B$ 的消解。Koszul 复形
$M_\bullet = K(P \otimes_A B, g_1, \ldots, g_n)$ 是由有限自由
$P \otimes_A B$-模组成的 $B$ 的消解。存在复形映射
$$
K_\bullet \longrightarrow M_\bullet
$$
它在次数 $1$ 由矩阵 $(g_{ij})$ 给出，在次数 $n$ 由
$\det(g_{ij})$ 给出；见《代数进阶》，引理
\ref{more-algebra-lemma-functorial}。由于 $B$ 是平坦 $A$-模，
可把 $M_\bullet$ 看作平坦 $P$-模的复形
（通过 $P \to P \otimes_A B$，$p \mapsto p \otimes 1$）。
因此可用这两个复形计算 $\text{Tor}_*^P(B, B)$，从而显示中的映射
与 $B$ 张量后给出拟同构。显然 $H_n(K_\bullet \otimes_P B) = B$。
另一方面，$H_n(M_\bullet \otimes_P B)$ 是下列映射的核：
$$
B \otimes_A B \xrightarrow{g_1, \ldots, g_n} (B \otimes_A B)^{\oplus n}
$$
由于 $g_1, \ldots, g_n$ 生成 $B \otimes_A B \to B$ 的核，引理得证。
\end{proof}

\begin{lemma}
\label{discriminant-lemma-quasi-finite-complete-intersection}
设 $A$ 是环，$n \geq 1$，且
$h, f_1, \ldots, f_n \in A[x_1, \ldots, x_n]$。令
$B = A[x_1, \ldots, x_n, 1/h]/(f_1, \ldots, f_n)$，
并假设 $B$ 在 $A$ 上拟有限。则
\begin{enumerate}
\item $B$ 在 $A$ 上平坦，且 $A \to B$ 是相对局部完全交；
\item $I = \Ker(B \otimes_A B \to B)$ 的零化子 $J$
作为 $B$-模自由且秩为 $1$；
\item $B$ 在 $A$ 上的 Noether 不同理想由 $B$ 中的
$\det(\partial f_i/\partial x_j)$ 生成。
\end{enumerate}
\end{lemma}

\begin{proof}
注意
$B = A[x, x_1, \ldots, x_n]/(xh - 1, f_1, \ldots, f_n)$
是 $A$ 上的相对整体完全交；见《代数》，定义
\ref{algebra-definition-relative-global-complete-intersection}。
由《代数》，引理 \ref{algebra-lemma-relative-global-complete-intersection}，
$B$ 在 $A$ 上平坦。

\medskip\noindent
写成 $P' = A[x, x_1, \ldots, x_n]$ 以及
$P = P'/(xh - 1) = A[x_1, \ldots, x_n, 1/h]$。于是有
$P' \to P \to B$。由《代数进阶》，引理
\ref{more-algebra-lemma-relative-global-complete-intersection-koszul}，
$xh - 1, f_1, \ldots, f_n$ 是 $P'$ 中的 Koszul 正则序列。
又因 $xh - 1$ 是 $P'$ 中长度为一的 Koszul 正则序列
（例如同样由该引理），故由《代数进阶》，引理
\ref{more-algebra-lemma-truncate-koszul-regular}，
$f_1, \ldots, f_n$ 是 $P$ 中的 Koszul 正则序列。

\medskip\noindent
令 $g_i \in P \otimes_A B$ 为 $x_i \otimes 1 - 1 \otimes x_i$ 的像。
在 $A[x_1, \ldots, x_n] \otimes_A A[x_1, \ldots, x_n]$ 中简记
$y_i = x_i \otimes 1$、$z_i = 1 \otimes x_i$，于是 $g_i$ 是
$y_i - z_i$ 的像。对多项式 $f \in A[x_1, \ldots, x_n]$，
在上述张量积中记 $f(y) = f \otimes 1$ 及 $f(z) = 1 \otimes f$。
于是
$$
P \otimes_A B/(g_1, \ldots, g_n) =
\frac{A[y_1, \ldots, y_n, z_1, \ldots, z_n, \frac{1}{h(y)h(z)}]}
{(f_1(z), \ldots, f_n(z), y_1 - z_1, \ldots, y_n - z_n)}
$$
显然同构于 $B$。因此仿照上述论证，
$f_1(z), \ldots, f_n(z), y_1 - z_1, \ldots, y_n - z_n$
是
$A[y_1, \ldots, y_n, z_1, \ldots, z_n, \frac{1}{h(y)h(z)}]$
中的 Koszul 正则序列。序列 $f_1(z), \ldots, f_n(z)$ 也是
$A[y_1, \ldots, y_n, z_1, \ldots, z_n, \frac{1}{h(y)h(z)}]$
中的 Koszul 正则序列，因为映射
$$
P \longrightarrow A[y_1, \ldots, y_n, z_1, \ldots, z_n,
\textstyle{\frac{1}{h(y)h(z)}}],\quad x_i \longmapsto z_i
$$
平坦，并可应用《代数进阶》，引理
\ref{more-algebra-lemma-koszul-regular-flat-base-change}。
再由《代数进阶》，引理
\ref{more-algebra-lemma-truncate-koszul-regular}，得出
$g_1, \ldots, g_n$ 是 $P \otimes_A B$ 中的正则序列。

\medskip\noindent
至此，以上述 $P$、$f_1, \ldots, f_n$ 及
$g_i \in P \otimes_A B$，已验证引理 \ref{discriminant-lemma-tate} 的全部假设。
特别地，$I$ 的零化子 $J$ 在 $B$ 上由一个元素 $\delta$ 自由生成。
令 $f_{ij} = \partial f_i/\partial x_j \in A[x_1, \ldots, x_n]$。
初等计算表明，可写成
$$
f_i(y) =
f_i(z_1 + g_1, \ldots, z_n + g_n) =
f_i(z) + \sum\nolimits_j f_{ij}(z) g_j +
\sum\nolimits_{j, j'} F_{ijj'}g_jg_{j'}
$$
其中 $F_{ijj'} \in A[y_1, \ldots, y_n, z_1, \ldots, z_n]$。
取其在 $P \otimes_A B$ 中的像，项 $f_i(z)$ 映到零，因而得到
$$
f_i \otimes 1 = \sum\nolimits_j
\left(1 \otimes f_{ij} + \sum\nolimits_{j'} F_{ijj'}g_{j'}\right)g_j
$$
于是由引理 \ref{discriminant-lemma-tate}，有 $\delta = \det(g_{ij})$，其中
$g_{ij} = 1 \otimes f_{ij} + \sum_{j'} F_{ijj'}g_{j'}$。
由于 $g_{j'}$ 在 $B$ 中的像为零，故
$\det(\partial f_i/\partial x_j)$ 在 $B$ 中的像生成 $B$ 在 $A$ 上的
Noether 不同理想。
\end{proof}

\begin{lemma}
\label{discriminant-lemma-different-syntomic-quasi-finite}
设 $f : Y \to X$ 是 Noetherian 概形的态射。若 $f$ 满足引理
\ref{discriminant-lemma-syntomic-quasi-finite} 的等价条件，则 $f$ 的不同理想
$\mathfrak{D}_f$ 等于 $f$ 的 K\"ahler 不同理想。
\end{lemma}

\begin{proof}
由引理 \ref{discriminant-lemma-flat-gorenstein-agree-noether} 及
\ref{discriminant-lemma-dualizing-syntomic-quasi-finite}，$f$ 的不同理想在仿射局部
等于 Noether 不同理想。再应用引理
\ref{discriminant-lemma-kahler-different-complete-intersection} 与
\ref{discriminant-lemma-quasi-finite-complete-intersection} 对标准仿射块上的
Noether 不同理想和 K\"ahler 不同理想所作的计算，即得结论。
\end{proof}

\begin{lemma}
\label{discriminant-lemma-different-quasi-finite-complete-intersection}
设 $A$ 是环，$n \geq 1$，且
$h, f_1, \ldots, f_n \in A[x_1, \ldots, x_n]$。令
$B = A[x_1, \ldots, x_n, 1/h]/(f_1, \ldots, f_n)$，
并假设 $B$ 在 $A$ 上拟有限。则存在同构 $B \to \omega_{B/A}$，
把 $\det(\partial f_i/\partial x_j)$ 映到 $\tau_{B/A}$。
\end{lemma}

\begin{proof}
令 $J$ 为 $\Ker(B \otimes_A B \to B)$ 的零化子。由引理
\ref{discriminant-lemma-quasi-finite-complete-intersection}，映射 $A \to B$ 平坦，
且 $J$ 是自由 $B$-模，其生成元 $\xi$ 在 $B$ 中映到
$\det(\partial f_i/\partial x_j)$。于是结论由引理
\ref{discriminant-lemma-noether-different-flat-quasi-finite} 以及如下事实得出：
$\omega_{B/A}$ 是可逆 $B$-模（引理
\ref{discriminant-lemma-dualizing-syntomic-quasi-finite}）。
（警告：必须证明 $\omega_{B/A}$ 可逆，因为满足
$\Hom_B(M, B) \cong B$ 的有限 $B$-模 $M$ 未必自由。）
\end{proof}

\begin{example}
\label{discriminant-example-different-for-monogenic}
设 $A$ 是 Noetherian 环，且 $f, h \in A[x]$ 满足
$$
B = (A[x]/(f))_h = A[x, 1/h]/(f)
$$
在 $A$ 上拟有限。令 $f' \in A[x]$ 为 $f$ 关于 $x$ 的导数。
理想 $\mathfrak{D} = (f') \subset B$ 既是 $B$ 在 $A$ 上的
Noether 不同理想，也是 $B$ 在 $A$ 上的 K\"ahler 不同理想；
此外，它所对应的拟凝聚理想层正是 $\Spec(B)$ 在 $\Spec(A)$ 上的不同理想。
\end{example}

\begin{lemma}
\label{discriminant-lemma-discriminant-quasi-finite-morphism-smooth}
设 $S$ 是 Noetherian 概形，$X$、$Y$ 是 $S$ 上相对维数为 $n$ 的光滑概形，
而 $f : Y \to X$ 是 $S$ 上局部拟有限的态射。则 $f$ 平坦，且由 $f$ 的
不同理想切出的闭子概形 $R \subset Y$，就是由下列截面切出的局部主闭子概形：
$$
\wedge^n(\text{d}f) \in
\Gamma(Y,
(f^*\Omega^n_{X/S})^{\otimes -1} \otimes_{\mathcal{O}_Y} \Omega^n_{Y/S})
$$
若 $f$ 在 $Y$ 的伴随点处为 \'etale，则 $R$ 是有效 Cartier 除子，并且
$$
f^*\Omega^n_{X/S} \otimes_{\mathcal{O}_Y} \mathcal{O}(R) =
\Omega^n_{Y/S}
$$
是 $Y$ 上可逆层的等式。
\end{lemma}

\begin{proof}
为证明 $f$ 平坦，只需对所有 $s \in S$ 证明 $Y_s \to X_s$ 平坦
（《态射进阶》，引理
\ref{more-morphisms-lemma-morphism-between-flat-Noetherian}）。
$Y_s \to X_s$ 的平坦性由《代数》，引理
\ref{algebra-lemma-CM-over-regular-flat} 得出。由《态射进阶》，引理
\ref{more-morphisms-lemma-lci-permanence}，态射 $f$ 是局部完全交态射。
因此关于不同理想的断言，由引理
\ref{discriminant-lemma-different-syntomic-quasi-finite} 及 K\"ahler 不同理想的相应断言得出。
最后，由《态射》，引理 \ref{morphisms-lemma-triangle-differentials}，有正合序列
$$
f^*\Omega_{X/S} \xrightarrow{\text{d}f} \Omega_{Y/S} \to \Omega_{Y/X} \to 0
$$
而 $\Omega_{X/S}$ 与 $\Omega_{Y/S}$ 都是秩为 $n$ 的有限局部自由模
（《态射》，引理
\ref{morphisms-lemma-smooth-omega-finite-locally-free}），
故由第零 Fitting 理想的定义，关于 K\"ahler 不同理想的断言是显然的。
若 $f$ 在 $Y$ 的伴随点处为 \'etale，则 $\wedge^n\text{d}f$ 在
$Y$ 的伴随点处不消失，因而 $R$ 的局部方程是非零因子。
所以 $R$ 是有效 Cartier 除子。典范同构把 $1$ 映到
$\wedge^n\text{d}f$；见《除子》，引理
\ref{divisors-lemma-characterize-OD}。
\end{proof}






\section{Tate 映射}
\label{discriminant-section-tate-map}

\noindent
本节对局部 Noetherian 概形的局部拟有限 syntomic 态射，构造相对余切复形的
行列式与相对对偶化模之间的同构。依照 \cite[1.4.4]{Garel}，
我们把这个同构称为 Tate 映射。我们的做法是尽量避免局部计算。

\medskip\noindent
设 $Y \to X$ 是概形的局部拟有限 syntomic 态射。本节将不再说明地使用
引理 \ref{discriminant-lemma-syntomic-quasi-finite} 中关于这一概念的全部等价条件。
特别地，$\NL_{Y/X}$ 是 $D(\mathcal{O}_Y)$ 中 Tor 振幅在 $[-1, 0]$
内的完美对象。因此在 $Y$ 上有典范可逆模 $\det(\NL_{Y/X})$ 以及整体截面
$$
\delta(\NL_{Y/X}) \in \Gamma(Y, \det(\NL_{Y/X}))
$$
见《概形的导出范畴》，引理
\ref{perfect-lemma-determinant-two-term-complexes}。设给定概形交换图
$$
\xymatrix{
Y' \ar[r]_b \ar[d] & Y \ar[d] \\
X' \ar[r] & X
}
$$
其中竖直箭头均局部拟有限且 syntomic，并且该图把 $Y'$ 同构到
$X' \times_X Y$ 的一个开子概形。则典范映射
$$
Lb^*\NL_{Y/X} \longrightarrow \NL_{Y'/X'}
$$
由《态射进阶》，引理 \ref{more-morphisms-lemma-base-change-NL-flat}
可知是拟同构。因此得到典范同构
$b^*\det(\NL_{Y/X}) \to \det(\NL_{Y'/X'})$，它把典范截面
$\delta(\NL_{Y/X})$ 映到 $\delta(\NL_{Y'/ X'})$；见《概形的导出范畴》，注记
\ref{perfect-remark-functorial-det}。

\begin{remark}
\label{discriminant-remark-local-description-delta}
设 $Y \to X$ 是概形的局部拟有限 syntomic 态射。序对
$(\det(\NL_{Y/X}), \delta(\NL_{Y/X}))$ 在局部上是什么样？选取仿射开集
$V = \Spec(B) \subset Y$、$U = \Spec(A) \subset X$，满足
$f(V) \subset U$；再取整数 $n$ 以及
$f_1, \ldots, f_n \in A[x_1, \ldots, x_n]$，使得
$B = A[x_1, \ldots, x_n]/(f_1, \ldots, f_n)$。则
$$
\NL_{B/A} = \left(
(f_1, \ldots, f_n)/(f_1, \ldots, f_n)^2
\longrightarrow
\bigoplus\nolimits_{i = 1, \ldots, n} B \text{d} x_i\right)
$$
且 $(f_1, \ldots, f_n)/(f_1, \ldots, f_n)^2$ 是自由模，
其生成元为同余类 $\overline{f}_i$；见引理
\ref{discriminant-lemma-syntomic-quasi-finite} 的证明。因此 $\det(L_{B/A})$
由下列生成元自由生成：
$$
\text{d}x_1 \wedge \ldots \wedge \text{d}x_n
\otimes
(\overline{f}_1 \wedge \ldots \wedge \overline{f}_n)^{\otimes -1}
$$
而截面 $\delta(\NL_{B/A})$ 按定义为元素
$$
\delta(\NL_{B/A}) =
\det(\partial f_j/ \partial x_i) \cdot
\text{d}x_1 \wedge \ldots \wedge \text{d}x_n
\otimes
(\overline{f}_1 \wedge \ldots \wedge \overline{f}_n)^{\otimes -1}
$$
\end{remark}

\noindent
设 $Y \to X$ 是局部 Noetherian 概形的局部拟有限 syntomic 态射。
由注记 \ref{discriminant-remark-relative-dualizing-for-quasi-finite} 及
\ref{discriminant-remark-relative-dualizing-for-flat-quasi-finite}，有凝聚
$\mathcal{O}_Y$-模 $\omega_{Y/X}$ 以及典范整体截面
$$
\tau_{Y/X} \in \Gamma(Y, \omega_{Y/X})
$$
它在仿射局部恢复序对 $\omega_{B/A}, \tau_{B/A}$。
由引理 \ref{discriminant-lemma-dualizing-syntomic-quasi-finite}，模
$\omega_{Y/X}$ 可逆。设给定局部 Noetherian 概形交换图
$$
\xymatrix{
Y' \ar[r]_b \ar[d] & Y \ar[d] \\
X' \ar[r] & X
}
$$
其中竖直箭头均局部拟有限且 syntomic，并且该图把 $Y'$ 同构到
$X' \times_X Y$ 的一个开子概形。则有典范基变换映射
$$
b^*\omega_{Y/X} \longrightarrow \omega_{Y'/X'}
$$
这是把 $\tau_{Y/X}$ 映到 $\tau_{Y'/X'}$ 的同构。事实上，仿射情形中的
基变换映射就是 (\ref{discriminant-equation-bc-dualizing})；由引理
\ref{discriminant-lemma-dualizing-base-change-of-flat}，它是同构；而由引理
\ref{discriminant-lemma-trace-base-change} 的 (1)，它把 $\tau_{Y/X}$ 映到
$\tau_{Y'/X'}$。

\begin{proposition}
\label{discriminant-proposition-tate-map}
存在唯一的规则，它对局部 Noetherian 概形的每个局部拟有限 syntomic 态射
$Y \to X$ 指派同构
$$
c_{Y/X} : \det(\NL_{Y/X}) \longrightarrow \omega_{Y/X}
$$
并满足下列两个性质：
\begin{enumerate}
\item 截面 $\delta(\NL_{Y/X})$ 被映到 $\tau_{Y/X}$；
\item 该规则与限制到开集以及基变换相容。
\end{enumerate}
\end{proposition}

\begin{proof}
先重述命题。考虑范畴 $\mathcal{C}$：其对象记为 $Y/X$，
是局部 Noetherian 概形的局部拟有限 syntomic 态射 $Y \to X$；
其态射 $b/a : Y'/X' \to Y/X$ 是交换图
$$
\xymatrix{
Y' \ar[d] \ar[r]_b & Y \ar[d] \\
X' \ar[r]^a & X
}
$$
并且该图把 $Y'$ 同构到 $X' \times_X Y$ 的一个开子概形。
命题意味着：对 $\mathcal{C}$ 的每个对象 $Y/X$，有同构
$c_{Y/X} : \det(\NL_{Y/X}) \to \omega_{Y/X}$，满足
$c_{Y/X}(\delta(\NL_{Y/X})) = \tau_{Y/X}$；而对 $\mathcal{C}$ 的每个态射
$b/a : Y'/X' \to Y/X$，在上述识别
$b^*\det(\NL_{Y/X}) = \det(\NL_{Y'/X'})$ 及
$b^*\omega_{Y/X} = \omega_{Y'/X'}$ 下，有
$b^*c_{Y/X} = c_{Y'/X'}$。

\medskip\noindent
给定 $\mathcal{C}$ 中的 $Y/X$ 及 $y \in Y$，可找到仿射开集
$V \subset Y$、$U \subset X$，满足 $y \in V$ 及 $f(V) \subset U$，
使得存在同构
$$
\det(\NL_{Y/X})|_V \longrightarrow \omega_{Y/X}|_V
$$
把 $\delta(\NL_{Y/X})|_V$ 映到 $\tau_{Y/X}|_V$。
这是因为可按引理 \ref{discriminant-lemma-syntomic-quasi-finite} 的 (5) 选取仿射开集，
再应用注记 \ref{discriminant-remark-local-description-delta} 中
$\delta(\NL_{Y/X})$ 的仿射局部描述以及引理
\ref{discriminant-lemma-different-quasi-finite-complete-intersection}。
若截面 $\tau_{Y/X}$ 的零化子为零，则这些局部映射唯一，且自动可以粘合。
因此，只要 $\tau_{Y/X}$ 的零化子为零，就存在唯一同构
$c_{Y/X} : \det(\NL_{Y/X}) \to \omega_{Y/X}$，满足
$c_{Y/X}(\delta(\NL_{Y/X})) = \tau_{Y/X}$。
若 $b/a : Y'/X' \to Y/X$ 是 $\mathcal{C}$ 中的态射，且
$\tau_{Y'/X'}$ 的零化子也为零，则 $b^*c_{Y/X}$ 是唯一同构
$c_{Y'/X'} : \det(\NL_{Y'/X'}) \to \omega_{Y'/X'}$，满足
$c_{Y'/X'}(\delta(\NL_{Y'/X'})) = \tau_{Y'/X'}$。
这由等式 $b^*\delta(\NL_{Y/X}) = \delta(\NL_{Y'/X'})$ 及
$b^*\tau_{Y/X} = \tau_{Y'/X'}$ 形式地得出。

\medskip\noindent
上一段可概括如下。令 $\mathcal{C}_{nice} \subset \mathcal{C}$
表示由满足 $\tau_{Y/X}$ 的零化子为零的 $Y/X$ 构成的全子范畴。
这样就在 $\mathcal{C}_{nice}$ 上解决了问题。对
$\mathcal{C}_{nice}$ 中的 $Y/X$，仍以 $c_{Y/X}$ 表示刚得到的解。

\medskip\noindent
考虑 $\mathcal{C}$ 中的态射
$$
Y_1/X_1 \xleftarrow{b_1/a_1} Y/X \xrightarrow{b_2/a_2} Y_2/X_2
$$
并设 $Y_1/X_1$ 与 $Y_2/X_2$ 是 $\mathcal{C}_{nice}$ 的对象。
{\bf 断言。}$b_1^*c_{Y_1/X_1} = b_2^*c_{Y_2/X_2}$。
先证明该断言蕴含命题，然后再证明断言。

\medskip\noindent
设 $d, n \geq 1$，并考虑例
\ref{discriminant-example-universal-quasi-finite-syntomic} 中构造的局部拟有限 syntomic 态射
$Y_{n, d} \to X_{n, d}$。则 $Y_{n, d}$ 是不可约正则概形，且态射
$Y_{n, d} \to X_{n, d}$ 局部拟有限、syntomic，并在某个稠密开集上为
\'etale；见引理 \ref{discriminant-lemma-universal-quasi-finite-syntomic-etale}。
例如由引理 \ref{discriminant-lemma-different-ramification}，
$\tau_{Y_{n, d}/X_{n, d}}$ 非零。而不可约正则概形上的可逆模之非零截面
其零化子为零。因此 $Y_{n, d}/X_{n, d}$ 是 $\mathcal{C}_{nice}$ 的对象。

\medskip\noindent
设 $Y/X$ 是 $\mathcal{C}$ 的任意对象，并取 $y \in Y$。
由引理 \ref{discriminant-lemma-locally-comes-from-universal}，可找到
$n, d \geq 1$ 以及 $\mathcal{C}$ 中的态射
$$
Y/X \leftarrow V/U \xrightarrow{b/a} Y_{n, d}/X_{n, d}
$$
其中 $V \subset Y$ 与 $U \subset X$ 为开集。因此可用 $b$ 把上面构造的
典范态射 $c_{Y_{n, d}/X_{n, d}}$ 拉回到 $V$。该断言保证这些局部同构能够粘合，
从而得到良定义的整体同构 $c_{Y/X} : \det(\NL_{Y/X}) \to \omega_{Y/X}$，
满足 $c_{Y/X}(\delta(\NL_{Y/X})) = \tau_{Y/X}$。
若 $b/a : Y'/X' \to Y/X$ 是 $\mathcal{C}$ 中的态射，则该断言还表明，
以同样方式构造的映射 $c_{Y'/X'}$ 是局部构造的映射 $c_{Y/X}$ 经 $b$ 的拉回。
故只需证明断言。

\medskip\noindent
在证明的其余部分证明该断言。可取一点 $y \in Y$，并证明这些映射在
$y$ 的某个开邻域中一致。因此，可用 $y$ 在 $Y_1$、$Y_2$ 中的像之开邻域
分别替换 $Y_1$、$Y_2$。于是可设 $Y, X, Y_1, X_1, Y_2, X_2$ 都仿射，
即分别为环 $B, A, B_1, A_1, B_2, A_2$ 的谱。图示为
$$
\xymatrix{
B_1 \ar[r] & B & B_2 \ar[l] \\
A_1 \ar[u] \ar[r] & A \ar[u] & A_2 \ar[l] \ar[u]
}
$$
由假设，$B$ 的谱同时是 $A \otimes_{A_1} B_1$ 的谱以及
$A \otimes_{A_2} B_2$ 的谱中的仿射开集。进一步缩小后，可设存在元素
$g_i \in A \otimes_{A_i} B_i$，使给定映射诱导同构
$(A \otimes_{A_i} B_i)_{g_i} = B$；见《性质》，引理
\ref{properties-lemma-standard-open-two-affines}。取足够多的变量
$x_\alpha$、$y_\beta$，使得可以选取 $A_1$-代数和 $A_2$-代数的满射
$$
A'_1 = A_1[x_\alpha] \to A
\quad\text{且}\quad
A'_2 = A_2[y_\beta] \to A
$$
再选取 $g_i \in A \otimes_{A_i} B_i$ 的提升
$h_i \in A'_i \otimes_{A_i} B_i$，并考虑图
$$
\xymatrix{
(A'_1 \otimes_{A_1} B_1)_{h_1} \ar[r] & B &
(A'_2 \otimes_{A_1} B_2)_{h_2} \ar[l] \\
A'_1 \ar[u] \ar[r] & A \ar[u] & A'_2 \ar[l] \ar[u]
}
$$
由构造，两个方块均为余笛卡尔图。接着考虑环同态
$$
A' = A'_1 \times_A A'_2 \longrightarrow
B' =
(A'_1 \otimes_{A_1} B_1)_{h_1} \times_B
(A'_2 \otimes_{A_1} B_2)_{h_2}
$$
由《代数进阶》，引理
\ref{more-algebra-lemma-module-over-fibre-product}，有
$A'_1 \otimes_{A'} B' = (A'_1 \otimes_{A_1} B_1)_{h_1}$ 以及
$A'_2 \otimes_{A'} B' = (A'_2 \otimes_{A_2} B_2)_{h_2}$。
特别地，态射 $Y' = \Spec(B') \to \Spec(A') = X'$ 的纤维是态射
$Y_i \to X_i$ 的纤维经基变换所得概形的开子概形，因而是有限局部完全交
（见引理 \ref{discriminant-lemma-syntomic-quasi-finite}）。
由《代数进阶》，引理
\ref{more-algebra-lemma-properties-algebras-over-fibre-product}，
环同态 $A' \to B'$ 平坦且有限表现。因此由引理
\ref{discriminant-lemma-syntomic-quasi-finite} 的 (3)，$Y' \to X'$ 局部拟有限且 syntomic。
从而得到交换图
$$
\xymatrix{
& & Y/X \ar[ld] \ar@/_2pc/[lldd]_{b_1/a_1} \ar[rd] \ar@/^2pc/[rrdd]^{b_2/a_2} \\
& Y'_1/X'_1 \ar[rd] \ar[ld] & & Y'_2/X'_2 \ar[ld] \ar[rd] \\
Y_1/X_1 & & Y'/X' & & Y_2/X_2
}
$$
其中 $Y'_i/X'_i$ 对应于 $A'_i \to (A'_i \otimes_{A_i} B_i)_{h_i}$。

\medskip\noindent
本段用极限论证处理 $A'$ 与 $A'_i$ 可能不是 Noetherian 环的问题；
建议读者略过本段。可找到有限生成的 $\mathbf{Z}$-子代数
$A'' \subset A'$ 以及拟有限 syntomic 环同态 $A'' \to B''$，使得
$B' = A' \otimes_{A''} B''$。这由《极限》，引理
\ref{limits-lemma-descend-finite-presentation}、
\ref{limits-lemma-descend-syntomic} 及
\ref{limits-lemma-descend-quasi-finite} 得出。
于是环同态 $A'' \to A'_1 = A_1[x_\alpha]$ 将通过某个有限多项式子代数
$A_i[x_1, \ldots, x_n$ 分解，而环同态
$A'' \to A'_2 = A_2[y_\beta]$ 将通过 $A_2[y_1, \ldots, y_m]$ 分解。
扩大变量集合后，还可设
$h_1 \in A_1[x_1, \ldots, x_n] \otimes_{A_1} B_1$ 及
$h_2 \in A_2[y_1, \ldots, y_m] \otimes_{A_2} B_2$。
再扩大一次变量集合，可设映射
$B'' \to (A'_i \otimes_{A_i} B_i)_{h_i}$ 分别通过
$(A_1[x_1, \ldots, x_n] \otimes_{A_1} B_1)_{h_1}$ 与
$(A_2[y_1, \ldots, y_m] \otimes_{A_2} B_2)_{h_2}$ 分解。
分别用下列环的谱替换 $Y'$、$X'$、$Y'_1$、$X'_1$、$Y'_2$、$X'_2$：
$B''$、$A''$、$(A_1[x_1, \ldots, x_n] \otimes_{A_1} B_1)_{h_1}$、
$A_1[x_1, \ldots, x_n]$、$(A_2[y_1, \ldots, y_m] \otimes_{A_2} B_2)_{h_2}$、
$A_2[y_1, \ldots, y_m]$。这样得到如上图，其中所有概形均为 Noetherian。

\medskip\noindent
注意，$Y'_i/X'_i$ 是 $\mathcal{C}_{nice}$ 的对象，因为它通过在某个仿射空间与
$Y_i/X_i$ 的乘积中取开子概形而得到。特别地，$c_{Y_i/X_i}$ 经
$(b_i/a_i)$ 的拉回，等于 $c_{Y'_i/X'_i}$ 经
$Y/X \to Y'_i/X'_i$ 的拉回。若 $Y'/X'$ 是
$\mathcal{C}_{nice}$ 的对象，则证明已经完成（但这很可能不成立）：
沿方块两边拉回 $c_{Y'/X'}$ 即得到所需结论。
为绕过 $Y'/X'$ 不属于 $\mathcal{C}_{nice}$ 的问题，注意上述论证表明，
在必要时缩小全部概形 $X, Y, X'_1, Y'_1, X'_2, Y'_2, X', Y'$ 后，
可找到 $n, d \geq 1$，并把图扩充为
$$
\xymatrix{
& Y/X \ar[ld] \ar[rd] \\
Y'_1/X'_1 \ar[rd] & & Y'_2/X'_2 \ar[ld] \\
& Y'/X' \ar[d] \\
& Y_{n, d}/X_{n, d}
}
$$
于是可从 $c_{Y_{n, d}/X_{n, d}}$ 拉回，并应用已经给出的论证。证明完成。
\end{proof}













\section{不同理想的一种推广}
\label{discriminant-section-different-generalization}

\noindent
本节推广定义 \ref{discriminant-definition-different}，以涵盖 Dedekind 不同理想有定义且
$1 \in \mathcal{L}_{B/A}$ 的所有环同态 $A \to B$。先解释条件：
“$A \to B$ 把非零因子映到非零因子，并诱导平坦映射
$Q(A) \to Q(A) \otimes_A B$”。

\begin{lemma}
\label{discriminant-lemma-explain-condition}
设 $A \to B$ 是 Noetherian 环同态。考虑条件：
\begin{enumerate}
\item $A$ 的非零因子映到 $B$ 的非零因子；
\item (1) 成立，且 $Q(A) \to Q(A) \otimes_A B$ 平坦；
\item 对每个 $\mathfrak q \in \text{Ass}(B)$，$A \to B_\mathfrak q$ 平坦；
\item (3) 成立，且对位于 $\text{Ass}(A)$ 中某元素之上的每个
$\mathfrak q$，$A \to B_\mathfrak q$ 平坦。
\end{enumerate}
则有蕴含关系
$$
\xymatrix{
(1) & (2) \ar@{=>}[l] \ar@{=>}[d] \\
(3) \ar@{=>}[u] & (4) \ar@{=>}[l]
}
$$
若 $A \to B$ 满足上升性质，则 (2) 与 (4) 等价。
\end{lemma}

\begin{proof}
图中的水平蕴含是平凡的。令 $S \subset A$ 为非零因子的集合，于是
$Q(A) = S^{-1}A$ 且 $Q(A) \otimes_A B = S^{-1}B$。回忆
$S = A \setminus \bigcup_{\mathfrak p \in \text{Ass}(A)} \mathfrak p$；
见《代数》，引理 \ref{algebra-lemma-ass-zero-divisors}。
设素理想 $\mathfrak q \subset B$ 位于 $\mathfrak p \subset A$ 之上。

\medskip\noindent
假设 (2)。若 $\mathfrak q \in \text{Ass}(B)$，则 $\mathfrak q$
由零因子组成，所以 (1) 表明 $\mathfrak p$ 也如此。因此
$\mathfrak p$ 对应于 $S^{-1}A$ 的一个素理想。由假设 (2)，
$A \to B_\mathfrak q$ 平坦。若 $\mathfrak q$ 位于 $A$ 的伴随素理想
$\mathfrak p$ 之上，则当然 $\mathfrak p \in \Spec(S^{-1}A)$，同一论证适用。

\medskip\noindent
假设 (3)。令 $f \in A$ 为非零因子。若 $f$ 是 $B$ 上的零因子，
则 $f$ 包含在 $B$ 的某个伴随素理想 $\mathfrak q$ 中。
由假设，$A \to B_\mathfrak q$ 平坦；再由《代数》，引理
\ref{algebra-lemma-bourbaki}，得出 $\mathfrak p$ 是 $A$ 的伴随素理想。
这将推出 $f$ 是 $A$ 上的零因子，矛盾。

\medskip\noindent
假设 (4)，并且 $A \to B$ 满足上升性质。已经知道 (1) 成立。
若 $\mathfrak q$ 对应于 $S^{-1}B$ 的一个素理想，则 $\mathfrak p$
包含在 $A$ 的某个伴随素理想 $\mathfrak p'$ 中。由上升性质，存在包含
$\mathfrak q$ 且位于 $\mathfrak p$ 之上的素理想 $\mathfrak q'$。
由 (4)，$A \to B_{\mathfrak q'}$ 平坦。因此作为局部化，
$A \to B_{\mathfrak q}$ 平坦。故 $A \to S^{-1}B$ 平坦，
$S^{-1}A \to S^{-1}B$ 也平坦；见《代数》，引理
\ref{algebra-lemma-flat-localization}。
\end{proof}

\begin{remark}
\label{discriminant-remark-different-generalization}
可以推广定义 \ref{discriminant-definition-different}。设 $f : Y \to X$
是 Noetherian 概形的拟有限态射，具有如下性质：
\begin{enumerate}
\item $f$ 平坦的开集 $V \subset Y$ 包含
$\text{Ass}(\mathcal{O}_Y)$ 与 $f^{-1}(\text{Ass}(\mathcal{O}_X))$；
\item 迹元 $\tau_{V/X}$ 来自某个截面
$\tau \in \Gamma(Y, \omega_{Y/X})$。
\end{enumerate}
由引理 \ref{discriminant-lemma-dualizing-associated-primes}，条件 (1) 蕴含
$V$ 包含 $\omega_{Y/X}$ 的伴随点。特别地，若 $\tau$ 存在，则它唯一
（《除子》，引理
\ref{divisors-lemma-restriction-injective-open-contains-ass}）。
给定 $\tau$，可把不同理想 $\mathfrak{D}_f$ 定义为
$\Coker(\tau : \mathcal{O}_Y \to \omega_{Y/X})$ 的零化子。
在许多情形中，这与 Dedekind 不同理想一致（引理
\ref{discriminant-lemma-agree-dedekind}）。不过，对非正规环之间的非平坦映射，
这一推广不再度量态射的分歧；见例 \ref{discriminant-example-no-different}。
\end{remark}

\begin{lemma}
\label{discriminant-lemma-agree-dedekind}
假设 $A \to B$ 的 Dedekind 不同理想有定义。令
$X = \Spec(A)$、$Y = \Spec(B)$。注记
\ref{discriminant-remark-different-generalization} 的推广适用于态射
$f : Y \to X$，当且仅当 $1 \in \mathcal{L}_{B/A}$
（例如，当 $A$ 正规时如此；见引理
\ref{discriminant-lemma-dedekind-different-ideal}）。在此情形中，
$\mathfrak{D}_{B/A}$ 是 $B$ 的理想，而且作为 $Y$ 上的凝聚理想层有
$$
\mathfrak{D}_f = \widetilde{\mathfrak{D}_{B/A}}
$$
\end{lemma}

\begin{proof}
由于 $A \to B$ 的 Dedekind 不同理想有定义，可应用引理
\ref{discriminant-lemma-explain-condition}，得知 $Y \to X$ 满足注记
\ref{discriminant-remark-different-generalization} 的条件 (1)。回忆存在典范同构
$c : \mathcal{L}_{B/A} \to \omega_{B/A}$；见引理
\ref{discriminant-lemma-dedekind-complementary-module}。如上令
$K = Q(A)$ 及 $L = K \otimes_A B$。由构造，映射 $c$ 嵌入交换图
$$
\xymatrix{
\mathcal{L}_{B/A} \ar[r] \ar[d]_c & L \ar[d] \\
\omega_{B/A} \ar[r] & \Hom_K(L, K)
}
$$
其中右侧竖直箭头把 $x \in L$ 映到映射
$y \mapsto \text{Trace}_{L/K}(xy)$，下方水平箭头是
$\omega_{B/A}$ 的基变换映射 (\ref{discriminant-equation-bc-dualizing})。
可把下方水平映射分解为
$$
\omega_{B/A} = \Gamma(Y, \omega_{Y/X})
\to \Gamma(V, \omega_{V/X}) \to \Hom_K(L, K)
$$
由于 $\omega_{V/X}$ 的所有伴随点都映到 $A$ 的伴随素理想
（引理 \ref{discriminant-lemma-dualizing-associated-primes}），第二个映射为单射。
由迹元的定义（定义 \ref{discriminant-definition-trace-element}），元素
$\tau_{V/X}$ 在 $\Hom_K(L, K)$ 中映到 $\text{Trace}_{L/K}$。
因此，注记 \ref{discriminant-remark-different-generalization} 条件 (2) 中的
$\tau$ 存在，当且仅当 $1 \in \mathcal{L}_{B/A}$；此时
$\tau = c(1)$。在这种情形下，由引理
\ref{discriminant-lemma-dedekind-different-ideal}，有
$\mathfrak{D}_{B/A} \subset B$。最后，从定义以及上面得到的
$\tau = c(1)$ 立即可知 $\mathfrak{D}_f$ 与 $\mathfrak{D}_{B/A}$ 一致。
\end{proof}

\begin{example}
\label{discriminant-example-no-different}
设 $k$ 是域，$A = k[x, y]/(xy)$、$B = k[u, v]/(uv)$，并令
$A \to B$ 由 $x \mapsto u^n$ 及 $y \mapsto v^m$ 给出，其中
$n, m \in \mathbf{N}$ 与 $k$ 的特征互素。则
$A_{x + y} \to B_{x + y}$ 是（有限）\'etale，故处在 Dedekind 不同理想
有定义的情形。计算表明
$$
\text{Trace}_{L/K}(1) = (nx + my)/(x + y),\quad
\text{Trace}_{L/K}(u^i) = 0,\quad \text{Trace}_{L/K}(v^j) = 0
$$
其中 $1 \leq i < n$ 且 $1 \leq j < m$。由此
$1 \in \mathcal{L}_{B/A}$ 当且仅当 $n = m$。此外，计算表明，若
$n = m$，则 $\mathcal{L}_{B/A} = B$，且 Dedekind 不同理想也等于
$B$。换言之，注记 \ref{discriminant-remark-different-generalization} 中的不同理想
对 $\Spec(B) \to \Spec(A)$ 有定义，当且仅当 $n = m$；在此情形中，
不同理想为单位理想。因此在非平坦情形中，不同理想非零并不能保证态射为
\'etale 或非分歧。
\end{example}







\section{与对偶理论的比较}
\label{discriminant-section-comparison}

\noindent
本节把上述初等代数构造与概形对偶理论一章中的构造作比较。

\begin{lemma}
\label{discriminant-lemma-compare-dualizing}
设 $f : Y \to X$ 是 Noetherian 概形之间的拟有限分离态射。
对每一对仿射开集 $\Spec(B) = V \subset Y$、
$\Spec(A) = U \subset X$，若 $f(V) \subset U$，则存在同构
$$
H^0(V, f^!\mathcal{O}_X) = \omega_{B/A}
$$
其中 $f^!$ 如《概形对偶》第
\ref{duality-section-upper-shriek} 节所述。这些同构与限制映射相容，
并给出典范同构 $H^0(f^!\mathcal{O}_X) = \omega_{Y/X}$，
其中 $\omega_{Y/X}$ 如注 \ref{discriminant-remark-relative-dualizing-for-quasi-finite}
所述。类似地，若 $f : Y \to X$ 是具有对偶复形
$\omega_S^\bullet$ 的 Noetherian 基概形 $S$ 上有限型概形之间的拟有限态射，
则 $H^0(f_{new}^!\mathcal{O}_X) = \omega_{Y/X}$。
\end{lemma}

\begin{proof}
由 Zariski 主定理，可选取分解 $f = f' \circ j$，其中
$j : Y \to Y'$ 是开浸入，而 $f' : Y' \to X$ 是有限态射；见
《态射进阶》，引理
\ref{more-morphisms-lemma-quasi-finite-separated-pass-through-finite}。
由《概形对偶》，引理 \ref{duality-lemma-shriek-well-defined}
中的构造，有 $f^! = j^* \circ a'$，其中
$a' : D_\QCoh(\mathcal{O}_X) \to D_\QCoh(\mathcal{O}_{Y'})$
是《概形对偶》，引理 \ref{duality-lemma-twisted-inverse-image}
中 $Rf'_*$ 的右伴随。由《概形对偶》，引理
\ref{duality-lemma-finite-twisted}，可知在
$D_\QCoh^+(f'_*\mathcal{O}_{Y'})$ 中
$\Phi(a'(\mathcal{O}_X)) = R\SheafHom(f'_*\mathcal{O}_{Y'}, \mathcal{O}_X)$。
特别地，$a'(\mathcal{O}_X)$ 在次数 $< 0$ 的上同调层均为零。
其第零上同调层由下列同构确定：
$$
f'_*H^0(a'(\mathcal{O}_X)) =
\SheafHom_{\mathcal{O}_X}(f'_*\mathcal{O}_{Y'}, \mathcal{O}_X)
$$
这里借助《态射》，引理
\ref{morphisms-lemma-affine-equivalence-modules} 的等价，
把两边视为 $f'_*\mathcal{O}_{Y'}$-模。记
$(f')^{-1}U = V' = \Spec(B')$，便得到
$$
H^0(V', a'(\mathcal{O}_X)) = \Hom_A(B', A).
$$
由于 $a'(\mathcal{O}_X)$ 的第零上同调层是拟凝聚模，
其在 $V$ 上的限制如所需由
$\omega_{B/A} = \Hom_A(B', A) \otimes_{B'} B$ 给出。

\medskip\noindent
关于限制映射的断言是指：经由上述同构，拟凝聚
$\mathcal{O}_{Y'}$-模 $H^0(a'(\mathcal{O}_X))$ 对 $Y'$ 中开集的限制映射，
与引理 \ref{discriminant-lemma-localize-dualizing} 为模 $\omega_{B/A}$
定义的映射一致。这是显然的。

\medskip\noindent
设 $f : Y \to X$ 是具有对偶复形 $\omega_S^\bullet$ 的 Noetherian
基概形 $S$ 上有限型概形之间的拟有限态射。取开集
$V \subset Y$ 与 $U \subset X$，使得 $f(V) \subset U$，且
$V$、$U$ 均在 $S$ 上分离。以 $f|_V : V \to U$ 表示 $f$ 的限制。
由上面的讨论以及《概形对偶》，引理
\ref{duality-lemma-duality-bootstrap}，存在典范同构
$$
H^0(f_{new}^!\mathcal{O}_X)|_V = H^0((f|_V)^!\mathcal{O}_U) = \omega_{V/U} =
\omega_{Y/X}|_V
$$
我们略去验证这些同构可粘合成整体同构
$H^0(f_{new}^!\mathcal{O}_X) \to \omega_{Y/X}$。
\end{proof}

\begin{lemma}
\label{discriminant-lemma-compare-trace}
设 $f : Y \to X$ 是 Noetherian 概形之间的有限平坦态射。第
\ref{discriminant-section-discriminant} 节的映射
$$
\text{Trace}_f : f_*\mathcal{O}_Y \longrightarrow \mathcal{O}_X
$$
对应于映射 $\mathcal{O}_Y \to f^!\mathcal{O}_X$（见证明）。
以 $\tau_{Y/X} \in H^0(Y, f^!\mathcal{O}_X)$ 表示 $1$ 的像。
经引理 \ref{discriminant-lemma-compare-dualizing} 的同构
$H^0(f^!\mathcal{O}_X) = \omega_{X/Y}$，这与注
\ref{discriminant-remark-relative-dualizing-for-flat-quasi-finite} 中的构造一致。
\end{lemma}

\begin{proof}
函子 $f^!$ 定义于《概形对偶》第
\ref{duality-section-upper-shriek} 节。由于 $f$ 有限（因而为恰当态射），
$f^!$ 由 $f$ 的推前函子的右伴随给出。在《概形对偶》第
\ref{duality-section-duality-finite} 节中，我们已具体写出这一伴随。
特别地，对象 $f^!\mathcal{O}_X$ 仅有一个置于次数 $0$ 的上同调层，
并且对该层有
$$
f_*f^!\mathcal{O}_X =
\SheafHom_{\mathcal{O}_X}(f_*\mathcal{O}_Y, \mathcal{O}_X)
$$
这里还用到：由于 $f$ 是 Noetherian 概形之间的有限平坦态射，
$f_*\mathcal{O}_Y$ 有限局部自由，故所有高次 Ext 层均为零。
略去若干细节。最终得到
$$
\text{Trace}_f \in
\Hom_{\mathcal{O}_X}(f_*\mathcal{O}_Y, \mathcal{O}_X) =
\Gamma(X, f_*f^!\mathcal{O}_X) =
\Gamma(Y, f^!\mathcal{O}_X)
$$
另一方面，由引理 \ref{discriminant-lemma-compare-dualizing} 的识别，有
$f^!\mathcal{O}_X = \omega_{Y/X}$。因此在
$$
\Gamma(Y, f^!\mathcal{O}_X) = \Gamma(Y, \omega_{Y/X})
$$
中现在有两个元素：$\text{Trace}_f$，以及注
\ref{discriminant-remark-relative-dualizing-for-flat-quasi-finite} 中的
$\tau_{Y/X}$；本引理断言二者相同。

\medskip\noindent
设 $U = \Spec(A) \subset X$ 是仿射开集，其逆像为
$V = \Spec(B) \subset Y$。由于 $f$ 有限，$A \to B$ 有限，
故由构造有 $\omega_{Y/X}(V) = \Hom_A(B,A)$；这一同构与上面所述
把 $f_*f^!\mathcal{O}_Y$ 识别为
$\SheafHom_{\mathcal{O}_X}(f_*\mathcal{O}_Y, \mathcal{O}_X)$ 的识别一致。
于是 $\text{Trace}_f$ 与 $\tau_{Y/X}$ 的一致性，由引理
\ref{discriminant-lemma-finite-flat-trace} 中
$\tau_{B/A} = \text{Trace}_{B/A}$ 这一事实得出。
\end{proof}








\section{拟有限 Gorenstein 态射}
\label{discriminant-section-gorenstein-lci}

\noindent
本节讨论拟有限 Gorenstein 态射。

\begin{lemma}
\label{discriminant-lemma-gorenstein-quasi-finite}
设 $f : Y \to X$ 是 Noetherian 概形之间的拟有限态射。下列条件等价：
\begin{enumerate}
\item $f$ 是 Gorenstein 态射；
\item $f$ 平坦，且 $f$ 的纤维均为 Gorenstein；
\item $f$ 平坦，且 $\omega_{Y/X}$ 可逆
（注 \ref{discriminant-remark-relative-dualizing-for-quasi-finite}）；
\item 对每个 $y \in Y$，存在仿射开集
$y \in V = \Spec(B) \subset Y$、$U = \Spec(A) \subset X$，
使得 $f(V) \subset U$、$A \to B$ 平坦，且
$\omega_{B/A}$ 是可逆 $B$-模。
\end{enumerate}
\end{lemma}

\begin{proof}
由定义，(1) 与 (2) 等价。由注
\ref{discriminant-remark-relative-dualizing-for-quasi-finite} 中
$\omega_{Y/X}$ 的构造，(3) 与 (4) 等价。因此只需证明
(1)--(2) 与 (3)--(4) 等价。

\medskip\noindent
第一种证明。仿射局部地工作，可假设 $f$ 是分离态射。
应用引理 \ref{discriminant-lemma-compare-dualizing} 可知，
$\omega_{Y/X}$ 是 $f^!\mathcal{O}_X$ 的第零上同调层。
在两组条件下，$f$ 都平坦且拟有限，故
$f^!\mathcal{O}_X$ 同构于 $\omega_{Y/X}[0]$；见《概形对偶》，引理
\ref{duality-lemma-flat-quasi-finite-shriek}。
于是所需等价性由《概形对偶》，引理
\ref{duality-lemma-affine-flat-Noetherian-gorenstein} 得出。

\medskip\noindent
第二种证明。由引理 \ref{discriminant-lemma-characterize-invertible}，
只需在 $X$ 是域 $k$ 的谱时证明 (2) 与 (3) 等价。
此时 $Y = \Spec(B)$，其中 $B$ 是有限 $k$-代数。
在这种情形下，置于次数 $0$ 的
$\omega_{B/A} = \omega_{B/k} = \Hom_k(B, k)$ 是 $B$ 的对偶复形；
见《对偶复形》，引理 \ref{dualizing-lemma-dualizing-finite}。
因此所需等价性由《对偶复形》，引理
\ref{dualizing-lemma-gorenstein} 得出。
\end{proof}

\begin{remark}
\label{discriminant-remark-collect-results-qf-gorenstein}
设 $f : Y \to X$ 是 Noetherian 概形之间的拟有限 Gorenstein 态射。
令 $\mathfrak D_f \subset \mathcal{O}_Y$ 为不同理想，并令
$R \subset Y$ 为由 $\mathfrak D_f$ 切出的闭子概形。则有：
\begin{enumerate}
\item $\mathfrak D_f$ 是局部主理想；
\item $R$ 是局部主闭子概形；
\item 仿射局部地，$\mathfrak D_f$ 与 Noether 不同理想相同；
\item $R$ 的形成与基变换相容；
\item 若 $f$ 有限，则 $R$ 的范数是 $f$ 的判别式；
\item 若 $f$ 在 $Y$ 的伴随点处为 \'etale，则
$R$ 是有效 Cartier 除子，且 $\omega_{Y/X} = \mathcal{O}_Y(R)$。
\end{enumerate}
这些结论由引理 \ref{discriminant-lemma-flat-gorenstein-agree-noether}、
\ref{discriminant-lemma-base-change-different} 和
\ref{discriminant-lemma-norm-different-is-discriminant} 得出。
\end{remark}

\begin{remark}
\label{discriminant-remark-collect-results-qf-gorenstein-two}
设 $S$ 是配备对偶复形 $\omega_S^\bullet$ 的 Noetherian 概形。
设 $f : Y \to X$ 是 $S$ 上可紧化概形之间的拟有限 Gorenstein 态射。
再假设 $Y$、$X$ 均为 Cohen--Macaulay，且 $f$ 在 $Y$ 的一般点处为
\'etale。结合《概形对偶》，注
\ref{duality-remark-CM-morphism-compare-dualizing} 与注
\ref{discriminant-remark-collect-results-qf-gorenstein}，可知有典范同构
$$
\omega_Y = f^*\omega_X \otimes_{\mathcal{O}_Y} \omega_{Y/X} =
f^*\omega_X \otimes_{\mathcal{O}_Y} \mathcal{O}_Y(R)
$$
这是 $\mathcal{O}_Y$-模的同构。若进一步 $f$ 有限，则同构
$\mathcal{O}_Y(R) = \omega_{Y/X}$ 来自整体截面
$\tau_{Y/X} \in H^0(Y, \omega_{Y/X})$；经对偶，该截面对应于映射
$\text{Trace}_f : f_*\mathcal{O}_Y \to \mathcal{O}_X$；见引理
\ref{discriminant-lemma-compare-trace}。
\end{remark}
% Generated mechanically from the frozen derham.tex target.
% Do not edit this generated file; edit the bound locale source instead.

% OK, start here
%

\setcounter{chapter}{49}
\chapter{德拉姆上同调}



\phantomsection
\label{derham-section-phantom}


\section{引言}
\label{derham-section-introduction}

\noindent
本章先讨论概形态射的德拉姆复形，最后证明：当基域的特征为零时，
德拉姆上同调定义了一种韦伊上同调理论。




\section{德拉姆复形}
\label{derham-section-de-rham-complex}

\noindent
设 $p : X \to S$ 是概形态射。存在复形
$$
\Omega^\bullet_{X/S} =
\mathcal{O}_{X/S} \to \Omega^1_{X/S} \to \Omega^2_{X/S} \to \ldots
$$
它是 $p^{-1}\mathcal{O}_S$-模复形，其中
$\Omega^i_{X/S} = \wedge^i(\Omega_{X/S})$
置于次数 $i$，其微分由局部截面上的法则
$\text{d}(g_0 \text{d}g_1 \wedge \ldots \wedge \text{d}g_p) =
\text{d}g_0 \wedge \text{d}g_1 \wedge \ldots \wedge \text{d}g_p$
决定。参见《模》，第 \ref{modules-section-de-rham-complex} 节。

\medskip\noindent
给定概形的交换图
$$
\xymatrix{
X' \ar[r]_f \ar[d] & X \ar[d] \\
S' \ar[r] & S
}
$$
则存在复形的典范映射
$f^{-1}\Omega_{X/S}^\bullet \to \Omega^\bullet_{X'/S'}$ 以及
$\Omega_{X/S}^\bullet \to f_*\Omega^\bullet_{X'/S'}$。
参见《模》，第 \ref{modules-section-de-rham-complex} 节。
将其线性化后，对每个 $p$ 都得到线性映射
$f^*\Omega^p_{X/S} \to \Omega^p_{X'/S'}$。

\medskip\noindent
特别地，若 $f : Y \to X$ 是基概形 $S$ 上的概形态射，则有复形映射
$$
\Omega^\bullet_{X/S} \longrightarrow f_*\Omega^\bullet_{Y/S}
$$
将其线性化可知，对每个 $p \geq 0$ 都得到典范映射
$$
\Omega^p_{X/S} \otimes_{\mathcal{O}_X} f_*\mathcal{O}_Y
\longrightarrow
f_*\Omega^p_{Y/S}
$$

\begin{lemma}
\label{derham-lemma-base-change-de-rham}
设
$$
\xymatrix{
X' \ar[r]_f \ar[d] & X \ar[d] \\
S' \ar[r] & S
}
$$
为概形的笛卡儿图。则上述映射诱导同构
$f^*\Omega^p_{X/S} \to \Omega^p_{X'/S'}$。
\end{lemma}

\begin{proof}
将《态射》，引理 \ref{morphisms-lemma-base-change-differentials}
与外幂构造和基变换相容这一事实结合即可。
\end{proof}

\begin{lemma}
\label{derham-lemma-etale}
考虑概形的交换图
$$
\xymatrix{
X' \ar[r]_f \ar[d] & X \ar[d] \\
S' \ar[r] & S
}
$$
若 $X' \to X$ 与 $S' \to S$ 均为平展态射，则上述映射诱导同构
$f^*\Omega^p_{X/S} \to \Omega^p_{X'/S'}$。
\end{lemma}

\begin{proof}
我们有 $\Omega_{S'/S} = 0$ 和 $\Omega_{X'/X} = 0$；例如参见
《态射》，引理 \ref{morphisms-lemma-etale-at-point}。于是，由《态射》中的引理
\ref{morphisms-lemma-triangle-differentials} 与
\ref{morphisms-lemma-triangle-differentials-smooth}
所给出的短正合列可知
$\Omega_{X'/S'} = \Omega_{X'/S} = f^*\Omega_{X/S}$。
取外幂即得结论。
\end{proof}






\section{德拉姆上同调}
\label{derham-section-de-rham-cohomology}

\noindent
设 $p : X \to S$ 是概形态射。我们把
{\it $X$ 在 $S$ 上的德拉姆上同调}定义为上同调群
$$
H^i_{dR}(X/S) = H^i(R\Gamma(X, \Omega^\bullet_{X/S}))
$$
由于 $\Omega^\bullet_{X/S}$ 是 $p^{-1}\mathcal{O}_S$-模复形，
这些上同调群自然是 $H^0(S, \mathcal{O}_S)$-模。

\medskip\noindent
给定概形的交换图
$$
\xymatrix{
X' \ar[r]_f \ar[d] & X \ar[d] \\
S' \ar[r] & S
}
$$
利用第 \ref{derham-section-de-rham-complex} 节中的典范映射，我们得到拉回映射
$$
f^* :
R\Gamma(X, \Omega^\bullet_{X/S})
\longrightarrow
R\Gamma(X', \Omega^\bullet_{X'/S'})
$$
以及
$$
f^* : H^i_{dR}(X/S) \longrightarrow H^i_{dR}(X'/S')
$$
这些拉回满足显然的复合法则。
特别地，若固定基概形 $S$，则德拉姆上同调是 $S$ 上概形范畴的反变函子。

\begin{lemma}
\label{derham-lemma-de-rham-affine}
设 $X \to S$ 是由环同态 $R \to A$ 给出的仿射概形态射。则
$R\Gamma(X, \Omega^\bullet_{X/S}) = \Omega^\bullet_{A/R}$ 在 $D(R)$ 中成立，且
$H^i_{dR}(X/S) = H^i(\Omega^\bullet_{A/R})$。
\end{lemma}

\begin{proof}
这由《概形上同调》，引理
\ref{coherent-lemma-quasi-coherent-affine-cohomology-zero}
和勒雷无环引理
（《导出范畴》，引理 \ref{derived-lemma-leray-acyclicity}）推出。
\end{proof}

\begin{lemma}
\label{derham-lemma-quasi-coherence-relative}
设 $p : X \to S$ 是概形态射。若 $p$ 拟紧且拟分离，则
$Rp_*\Omega^\bullet_{X/S}$ 是 $D_\QCoh(\mathcal{O}_S)$ 中的对象。
\end{lemma}

\begin{proof}
存在一个第一叶为 $E_1^{a, b} = R^bp_*\Omega^a_{X/S}$、
收敛到 $Rp_*\Omega^\bullet_{X/S}$ 的上同调的谱序列
（参见《导出范畴》，引理 \ref{derived-lemma-two-ss-complex-functor}）。
因此，由《同调》，引理 \ref{homology-lemma-first-quadrant-ss}，
只需证明 $R^bp_*\Omega^a_{X/S}$ 拟凝聚。
这由《概形上同调》，引理
\ref{coherent-lemma-quasi-coherence-higher-direct-images} 推出。
\end{proof}

\begin{lemma}
\label{derham-lemma-coherence-relative}
设 $p : X \to S$ 是真概形态射，且 $S$ 局部诺特。则
$Rp_*\Omega^\bullet_{X/S}$ 是 $D_{\textit{Coh}}(\mathcal{O}_S)$ 中的对象。
\end{lemma}

\begin{proof}
在此情形下，由《态射》，引理 \ref{morphisms-lemma-finite-type-differentials}，
模 $\Omega^i_{X/S}$ 是凝聚的。因此，可以逐字采用引理
\ref{derham-lemma-quasi-coherence-relative} 证明中的论证，并使用《概形上同调》，命题
\ref{coherent-proposition-proper-pushforward-coherent}。
\end{proof}

\begin{lemma}
\label{derham-lemma-finite-de-Rham}
设 $A$ 是诺特环，并设 $X$ 是 $S = \Spec(A)$ 上的真概形。
则 $H^i_{dR}(X/S)$ 是有限 $A$-模（对所有 $i$）。
\end{lemma}

\begin{proof}
这是引理 \ref{derham-lemma-coherence-relative} 的特例。
\end{proof}

\begin{lemma}
\label{derham-lemma-proper-smooth-de-Rham}
设 $f : X \to S$ 是真光滑概形态射。则
$Rf_*\Omega^p_{X/S}$（$p \geq 0$）和 $Rf_*\Omega^\bullet_{X/S}$
都是 $D(\mathcal{O}_S)$ 中的完美对象，且其构造与任意基变换相容。
\end{lemma}

\begin{proof}
由于 $f$ 光滑，模 $\Omega^p_{X/S}$ 是有限局部自由
$\mathcal{O}_X$-模；参见《态射》，引理
\ref{morphisms-lemma-smooth-omega-finite-locally-free}。
由引理 \ref{derham-lemma-base-change-de-rham}，其构造与任意基变换相容。因此，
$Rf_*\Omega^p_{X/S}$ 是 $D(\mathcal{O}_S)$ 中的完美对象，
且其构造与任意基变换相容；参见《概形的导出范畴》，引理
\ref{perfect-lemma-flat-proper-perfect-direct-image-general}。
这证明了本引理的第一个断言。

\medskip\noindent
为了证明 $Rf_*\Omega^\bullet_{X/S}$ 在 $S$ 上完美，可以在 $S$ 上局部地论证。
因此可设 $S$ 拟紧。这意味着可以假设 $\Omega^n_{X/S}$ 在 $n$ 充分大时为零。
对每个 $p \geq 0$，我们断言
$Rf_*\sigma_{\geq p}\Omega^\bullet_{X/S}$ 是
$D(\mathcal{O}_S)$ 中的完美对象，且其构造与任意基变换相容。
由上述论证可知，当 $p \gg 0$ 时断言成立。假设该断言对 $p$ 成立，
并考虑可区别三角
$$
\sigma_{\geq p}\Omega^\bullet_{X/S} \to
\sigma_{\geq p - 1}\Omega^\bullet_{X/S} \to
\Omega^{p - 1}_{X/S}[-(p - 1)] \to
(\sigma_{\geq p}\Omega^\bullet_{X/S})[1]
$$
它位于 $D(f^{-1}\mathcal{O}_S)$ 中。
应用正合函子 $Rf_*$，在 $D(\mathcal{O}_S)$ 中得到一个可区别三角。
由于完美性具有三取二性质
（《上同调》，引理 \ref{cohomology-lemma-two-out-of-three-perfect}），
可得 $Rf_*\sigma_{\geq p - 1}\Omega^\bullet_{X/S}$ 是
$D(\mathcal{O}_S)$ 中的完美对象。其构造与任意基变换相容的断言同理。
\end{proof}





\section{杯积}
\label{derham-section-cup-product}

\noindent
考虑映射
$\Omega^p_{X/S} \times \Omega^q_{X/S} \to \Omega^{p + q}_{X/S}$
其定义为 $(\omega , \eta) \longmapsto \omega \wedge \eta$。
利用第 \ref{derham-section-de-rham-complex} 节给出的 $\text{d}$ 的公式，
以及 $\text{d} : \mathcal{O}_X \to \Omega_{X/S}$ 的莱布尼茨法则，可知
$\text{d}(\omega \wedge \eta) = \text{d}(\omega) \wedge \eta +
(-1)^{\deg(\omega)} \omega \wedge \text{d}(\eta)$。这说明
$\wedge$ 定义了态射
\begin{equation}
\label{derham-equation-wedge}
\wedge :
\text{Tot}(
\Omega^\bullet_{X/S} \otimes_{p^{-1}\mathcal{O}_S} \Omega^\bullet_{X/S})
\longrightarrow
\Omega^\bullet_{X/S}
\end{equation}
它是 $p^{-1}\mathcal{O}_S$-模复形的态射。

\medskip\noindent
把《上同调》，第 \ref{cohomology-section-cup-product} 节中的杯积
与（\ref{derham-equation-wedge}）结合，得到一个
$H^0(S, \mathcal{O}_S)$-双线性杯积映射
$$
\cup : H^i_{dR}(X/S) \times H^j_{dR}(X/S) \longrightarrow H^{i + j}_{dR}(X/S)
$$
例如，若 $\omega \in \Gamma(X, \Omega^i_{X/S})$ 和
$\eta \in \Gamma(X, \Omega^j_{X/S})$ 都是闭的，则
$\omega$ 与 $\eta$ 的德拉姆上同调类之杯积，就是 $\omega \wedge \eta$ 的
德拉姆上同调类；参见《上同调》，第 \ref{cohomology-section-cup-product} 节中的讨论。

\medskip\noindent
给定概形的交换图
$$
\xymatrix{
X' \ar[r]_f \ar[d] & X \ar[d] \\
S' \ar[r] & S
}
$$
则拉回映射
$f^* : R\Gamma(X, \Omega^\bullet_{X/S}) \to R\Gamma(X', \Omega^\bullet_{X'/S'})$
以及
$f^* : H^i_{dR}(X/S) \longrightarrow H^i_{dR}(X'/S')$
与上述杯积相容。

\begin{lemma}
\label{derham-lemma-cup-product-graded-commutative}
设 $p : X \to S$ 是概形态射。
$H^*_{dR}(X/S)$ 上的杯积满足结合律和分次交换律。
\end{lemma}

\begin{proof}
这由《上同调》中的引理 \ref{cohomology-lemma-cup-product-associative} 与
\ref{cohomology-lemma-cup-product-commutative}
以及 $\wedge$ 满足结合律和分次交换律这一事实推出。
\end{proof}

\begin{remark}
\label{derham-remark-relative-cup-product}
设 $p : X \to S$ 是概形态射。可以把
$\Omega^\bullet_{X/S}$ 看作微分分次 $p^{-1}\mathcal{O}_S$-代数层；参见
《微分分次层》，定义 \ref{sdga-definition-dga}。
特别地，《微分分次层》，第 \ref{sdga-section-misc} 节中的讨论适用。
例如，这意味着对任意交换图
$$
\xymatrix{
X \ar[d]_p \ar[r]_f & Y \ar[d]^q \\
S \ar[r]^h & T
}
$$
都存在典范的相对杯积
$$
\mu :
Rf_*\Omega^\bullet_{X/S}
\otimes_{q^{-1}\mathcal{O}_T}^\mathbf{L}
Rf_*\Omega^\bullet_{X/S}
\longrightarrow
Rf_*\Omega^\bullet_{X/S}
$$
它位于 $D(Y, q^{-1}\mathcal{O}_T)$ 中，满足结合律，并且在上同调上给出上述杯积。
\end{remark}

\begin{remark}
\label{derham-remark-cup-product-as-a-map}
设 $f : X \to S$ 是概形态射，并设 $\xi \in H_{dR}^n(X/S)$。
根据《微分分次层》，第 \ref{sdga-section-misc} 节中的讨论，存在典范态射
$$
\xi' : \Omega^\bullet_{X/S} \to \Omega^\bullet_{X/S}[n]
$$
它位于 $D(f^{-1}\mathcal{O}_S)$ 中，并由下列性质中的 (1) 和 (2) 唯一刻画：
\begin{enumerate}
\item $\xi'$ 可提升为右微分分次 $\Omega^\bullet_{X/S}$-模的导出范畴中的态射，并且
\item $\xi'(1) = \xi$ 在
$H^0(X, \Omega^\bullet_{X/S}[n]) = H^n_{dR}(X/S)$ 中成立；
\item 态射 $\xi'$ 把 $\eta \in H^m_{dR}(X/S)$ 送到
$\xi \cup \eta$；这是 $H^{n + m}_{dR}(X/S)$ 中的元素；
\item $\xi'$ 的构造与到开集的限制相容：对开集 $U \subset X$，
限制 $\xi'|_U$ 是与像 $\xi|_U \in H^n_{dR}(U/S)$ 对应的态射；
\item 对备注 \ref{derham-remark-relative-cup-product} 中的任意图，得到交换图
$$
\xymatrix{
Rf_*\Omega^\bullet_{X/S}
\otimes_{q^{-1}\mathcal{O}_T}^\mathbf{L}
Rf_*\Omega^\bullet_{X/S} \ar[d]_{\xi' \otimes \text{id}}
\ar[r]_-\mu &
Rf_*\Omega^\bullet_{X/S} \ar[d]^{\xi'} \\
Rf_*\Omega^\bullet_{X/S}[n]
\otimes_{q^{-1}\mathcal{O}_T}^\mathbf{L}
Rf_*\Omega^\bullet_{X/S}
\ar[r]^-\mu &
Rf_*\Omega^\bullet_{X/S}[n]
}
$$
它位于 $D(Y, q^{-1}\mathcal{O}_T)$ 中。
\end{enumerate}
\end{remark}




\section{霍奇上同调}
\label{derham-section-hodge-cohomology}

\noindent
设 $p : X \to S$ 是概形态射。我们把
{\it $X$ 在 $S$ 上的霍奇上同调}定义为上同调群
$$
H^n_{Hodge}(X/S) = \bigoplus\nolimits_{n = p + q} H^q(X, \Omega^p_{X/S})
$$
并将其视为分次 $H^0(X, \mathcal{O}_X)$-模。把微分形式的楔积
与《上同调》，第 \ref{cohomology-section-cup-product} 节中的杯积结合，
便定义了一个 $H^0(X, \mathcal{O}_X)$-双线性杯积
$$
\cup :
H^i_{Hodge}(X/S) \times H^j_{Hodge}(X/S)
\longrightarrow
H^{i + j}_{Hodge}(X/S)
$$
当然，若 $\xi \in H^q(X, \Omega^p_{X/S})$ 且
$\xi' \in H^{q'}(X, \Omega^{p'}_{X/S})$，则 $\xi \cup \xi' \in
H^{q + q'}(X, \Omega^{p + p'}_{X/S})$。

\begin{lemma}
\label{derham-lemma-cup-product-hodge-graded-commutative}
设 $p : X \to S$ 是概形态射。
$H^*_{Hodge}(X/S)$ 上的杯积满足结合律和分次交换律。
\end{lemma}

\begin{proof}
其证明与引理 \ref{derham-lemma-cup-product-graded-commutative} 的证明完全相同。
\end{proof}

\noindent
给定概形的交换图
$$
\xymatrix{
X' \ar[r]_f \ar[d] & X \ar[d] \\
S' \ar[r] & S
}
$$
则有拉回映射
$f^* : H^i_{Hodge}(X/S) \longrightarrow H^i_{Hodge}(X'/S')$
它与分次以及上述杯积相容。







\section{两个谱序列}
\label{derham-section-hodge-to-de-rham}

\noindent
设 $p : X \to S$ 是概形态射。由于 $p^{-1}\mathcal{O}_S$-模在 $X$ 上的范畴
有足够多的内射对象，
$\Omega^\bullet_{X/S}$ 存在嘉当--艾伦伯格解消。
参见《导出范畴》，引理 \ref{derived-lemma-cartan-eilenberg}。
因此可以应用《导出范畴》，引理
\ref{derived-lemma-two-ss-complex-functor}，得到两个都收敛到 $X$ 在 $S$
上之德拉姆上同调的谱序列。

\medskip\noindent
第一个通常称为{\it 霍奇--德拉姆谱序列}。该谱序列的第一叶为
$$
E_1^{p, q} = H^q(X, \Omega^p_{X/S})
$$
这些就是 $X/S$ 的霍奇上同调群（名称由此而来）。这一叶上的微分
$d_1$ 由映射
$d_1^{p, q} : H^q(X, \Omega^p_{X/S}) \to H^q(X, \Omega^{p + 1}_{X/S})$
给出；这些映射由微分
$\text{d} : \Omega^p_{X/S} \to \Omega^{p + 1}_{X/S}$ 诱导。图示如下：
$$
\xymatrix{
H^2(X, \mathcal{O}_X) \ar[r] \ar@{-->}[rrd] \ar@{..>}[rrrdd] &
H^2(X, \Omega^1_{X/S}) \ar[r] \ar@{-->}[rrd] &
H^2(X, \Omega^2_{X/S}) \ar[r] &
H^2(X, \Omega^3_{X/S}) \\
H^1(X, \mathcal{O}_X) \ar[r] \ar@{-->}[rrd] &
H^1(X, \Omega^1_{X/S}) \ar[r] \ar@{-->}[rrd] &
H^1(X, \Omega^2_{X/S}) \ar[r] &
H^1(X, \Omega^3_{X/S}) \\
H^0(X, \mathcal{O}_X) \ar[r] &
H^0(X, \Omega^1_{X/S}) \ar[r] &
H^0(X, \Omega^2_{X/S}) \ar[r] &
H^0(X, \Omega^3_{X/S})
}
$$
其中用虚线箭头标出 $E_2$ 叶上微分的源与靶，并用点线箭头表示
$E_3$ 叶上的一个微分。考察次数 $0$，得到
$$
H^0_{dR}(X/S) =
\Ker(\text{d} : H^0(X, \mathcal{O}_X) \to H^0(X, \Omega^1_{X/S}))
$$
当然，从德拉姆复形的次数 $0$ 项起始为
$\mathcal{O}_X \to \Omega^1_{X/S}$ 这一事实也能立刻看出这一点。

\medskip\noindent
第二个谱序列通常称为{\it 共轭谱序列}。该谱序列的第二叶为
$$
E_2^{p, q} = H^p(X, H^q(\Omega^\bullet_{X/S})) = H^p(X, \mathcal{H}^q)
$$
其中 $\mathcal{H}^q = H^q(\Omega^\bullet_{X/S})$ 是德拉姆复形的第 $q$
个上同调层；这里的德拉姆复形是 $X/S$ 的德拉姆复形。
这一叶上的微分由 $E_2^{p, q} \to E_2^{p + 2, q - 1}$ 给出。
图示如下：
$$
\xymatrix{
H^0(X, \mathcal{H}^2) \ar[rrd] \ar@{..>}[rrrdd] &
H^1(X, \mathcal{H}^2) \ar[rrd] &
H^2(X, \mathcal{H}^2) &
H^3(X, \mathcal{H}^2) \\
H^0(X, \mathcal{H}^1) \ar[rrd] &
H^1(X, \mathcal{H}^1) \ar[rrd] &
H^2(X, \mathcal{H}^1) &
H^3(X, \mathcal{H}^1) \\
H^0(X, \mathcal{H}^0) &
H^1(X, \mathcal{H}^0) &
H^2(X, \mathcal{H}^0) &
H^3(X, \mathcal{H}^0)
}
$$
考察次数 $0$，得到
$$
H^0_{dR}(X/S) = H^0(X, \mathcal{H}^0)
$$
细想之下，这是显然的。在次数 $1$，得到正合列
$$
0 \to H^1(X, \mathcal{H}^0) \to H^1_{dR}(X/S) \to
H^0(X, \mathcal{H}^1) \to H^2(X, \mathcal{H}^0) \to H^2_{dR}(X/S)
$$
事实表明，若 $X \to S$ 光滑且 $S$ 处于特征 $p$，则层
$\mathcal{H}^q$ 可以借助某个微分形式层来计算，而共轭谱序列是一个很有价值的
工具（在此插入将来的参考文献）。



\section{霍奇滤过}
\label{derham-section-hodge-filtration}

\noindent
设 $X \to S$ 是概形态射。$H^n_{dR}(X/S)$ 上的霍奇滤过，是由
霍奇--德拉姆谱序列诱导的滤过（《同调》，定义
\ref{homology-definition-filtration-cohomology-filtered-complex}）。
为避免误解，下面将它明确地定义出来。

\begin{definition}
\label{derham-definition-hodge-filtration}
设 $X \to S$ 是概形态射。$H^n_{dR}(X/S)$ 上的{\it 霍奇滤过}，
是由下列各项组成的滤过：
$$
F^pH^n_{dR}(X/S) = \Im\left(H^n(X, \sigma_{\geq p}\Omega^\bullet_{X/S})
\longrightarrow H^n_{dR}(X/S)\right)
$$
其中 $\sigma_{\geq p}\Omega^\bullet_{X/S}$ 如《同调》，
第 \ref{homology-section-truncations} 节所述。
\end{definition}

\noindent
当然，$\sigma_{\geq p}\Omega^\bullet_{X/S}$ 是相对德拉姆复形的子复形，
于是得到相对德拉姆复形的滤过
$$
\Omega^\bullet_{X/S} = \sigma_{\geq 0}\Omega^\bullet_{X/S} \supset
\sigma_{\geq 1}\Omega^\bullet_{X/S} \supset
\sigma_{\geq 2}\Omega^\bullet_{X/S} \supset
\sigma_{\geq 3}\Omega^\bullet_{X/S} \supset \ldots
$$
其满足
$\text{gr}^p(\Omega^\bullet_{X/S}) = \Omega^p_{X/S}[-p]$。
把 $\Omega^\bullet_{X/S}$ 视为滤过的层复形，则《上同调》，引理
\ref{cohomology-lemma-spectral-sequence-filtered-object} 所构造的谱序列，
与第 \ref{derham-section-hodge-to-de-rham} 节中借助《上同调》，例
\ref{cohomology-example-spectral-sequence-bis} 构造的霍奇--德拉姆谱序列相同。
此外，楔积（\ref{derham-equation-wedge}）把
$\text{Tot}(\sigma_{\geq i}\Omega^\bullet_{X/S} \otimes_{p^{-1}\mathcal{O}_S}
\sigma_{\geq j}\Omega^\bullet_{X/S})$ 送入
$\sigma_{\geq i + j}\Omega^\bullet_{X/S}$。因此得到交换图
$$
\xymatrix{
H^n(X, \sigma_{\geq i}\Omega^\bullet_{X/S}))
\times 
H^m(X, \sigma_{\geq j}\Omega^\bullet_{X/S}))
\ar[r] \ar[d] &
H^{n + m}(X, \sigma_{\geq i + j}\Omega^\bullet_{X/S})) \ar[d] \\
H^n_{dR}(X/S) \times
H^m_{dR}(X/S)
\ar[r]^\cup &
H^{n + m}_{dR}(X/S)
}
$$
特别地，可以得到
$$
F^iH^n_{dR}(X/S) \cup F^jH^m_{dR}(X/S) \subset F^{i + j}H^{n + m}_{dR}(X/S)
$$






\section{K\"unneth 公式}
\label{derham-section-kunneth}

\noindent
德拉姆上同调的一个重要特征是它具有 K\"unneth 公式。

\medskip\noindent
设 $a : X \to S$ 和 $b : Y \to S$ 是具有相同靶的概形态射。
设 $p : X \times_S Y \to X$ 和 $q : X \times_S Y \to Y$ 为投影态射，并令
$f = a \circ p = b \circ q$。图示如下：
$$
\xymatrix{
& X \times_S Y \ar[ld]^p \ar[rd]_q \ar[dd]^f \\
X \ar[rd]_a & & Y \ar[ld]^b \\
& S
}
$$
在本节中，给定 $\mathcal{O}_X$-模 $\mathcal{F}$ 和
$\mathcal{O}_Y$-模 $\mathcal{G}$，记
$$
\mathcal{F} \boxtimes \mathcal{G} =
p^*\mathcal{F} \otimes_{\mathcal{O}_{X \times_S Y}} q^*\mathcal{G}
$$
拟凝聚模上的双函子
$(\mathcal{F}, \mathcal{G}) \mapsto \mathcal{F} \boxtimes \mathcal{G}$
可扩张为以拟凝聚模为对象、以 $S$ 上有限阶微分算子为态射的双函子；
参见《态射》，备注 \ref{morphisms-remark-base-change-differential-operators}。
由《模》，引理
\ref{modules-lemma-differentials-relative-de-rham-complex-order-1}，
德拉姆复形 $\Omega^\bullet_{X/S}$ 和 $\Omega^\bullet_{Y/S}$ 的微分是
$1$ 阶的、定义在 $S$ 上的微分算子。因此可以考虑复形
$$
\text{Tot}(\Omega^\bullet_{X/S} \boxtimes \Omega^\bullet_{Y/S})
$$
请参见《概形的导出范畴》，第
\ref{perfect-section-kunneth-complexes} 节中的讨论。

\begin{lemma}
\label{derham-lemma-de-rham-complex-product}
在上述情形下，存在典范同构
$$
\text{Tot}(\Omega^\bullet_{X/S} \boxtimes \Omega^\bullet_{Y/S})
\longrightarrow
\Omega^\bullet_{X \times_S Y/S}
$$
它是 $f^{-1}\mathcal{O}_S$-模复形的同构。
\end{lemma}

\begin{proof}
由《态射》，引理 \ref{morphisms-lemma-differential-product}，可知
$
\Omega_{X \times_S Y/S} = p^*\Omega_{X/S} \oplus q^*\Omega_{Y/S}
$
。取外幂得到
$$
\Omega^n_{X \times_S Y/S} =
\bigoplus\nolimits_{i + j = n}
p^*\Omega^i_{X/S} \otimes_{\mathcal{O}_{X \times_S Y}} q^*\Omega^j_{Y/S} =
\bigoplus\nolimits_{i + j = n}
\Omega^i_{X/S} \boxtimes \Omega^j_{Y/S}
$$
这是外幂的初等性质。这些等同给出了本引理箭头两侧复形各项之间的同构。
我们略去验证这些映射与微分相容的步骤。
\end{proof}

\noindent
置 $A = \Gamma(S, \mathcal{O}_S)$。把引理
\ref{derham-lemma-de-rham-complex-product} 的结论与《概形的导出范畴》中的映射
（方程 \ref{perfect-equation-de-rham-kunneth}）结合，得到杯积
$$
R\Gamma(X, \Omega^\bullet_{X/S})
\otimes_A^\mathbf{L}
R\Gamma(Y, \Omega^\bullet_{Y/S})
\longrightarrow
R\Gamma(X \times_S Y, \Omega^\bullet_{X \times_S Y/S})
$$
在上同调层面，利用《更多代数》，第
\ref{more-algebra-section-products-tor} 节中的讨论，得到典范映射
$$
H^i_{dR}(X/S) \otimes_A H^j_{dR}(Y/S)
\longrightarrow
H^{i + j}_{dR}(X \times_S Y/S),\quad
(\xi, \zeta) \longmapsto p^*\xi \cup q^*\zeta
$$
注意，上述构造确实是先作拉回，再取杯积。

\begin{lemma}
\label{derham-lemma-kunneth-de-rham}
假设 $X$ 和 $Y$ 在 $S = \Spec(A)$ 上光滑、拟紧，且对角态射仿射。则映射
$$
R\Gamma(X, \Omega^\bullet_{X/S})
\otimes_A^\mathbf{L}
R\Gamma(Y, \Omega^\bullet_{Y/S})
\longrightarrow
R\Gamma(X \times_S Y, \Omega^\bullet_{X \times_S Y/S})
$$
是 $D(A)$ 中的同构。
\end{lemma}

\begin{proof}
由《态射》，引理 \ref{morphisms-lemma-smooth-omega-finite-locally-free}，
层 $\Omega^n_{X/S}$ 和 $\Omega^m_{Y/S}$ 分别是有限局部自由
$\mathcal{O}_X$-模和 $\mathcal{O}_Y$-模。另一方面，$X$ 和 $Y$ 在 $S$ 上平坦
（《态射》，引理 \ref{morphisms-lemma-smooth-flat}），
因而 $\Omega^n_{X/S}$ 和 $\Omega^m_{Y/S}$ 在 $S$ 上平坦。
还要注意，$\Omega^\bullet_{X/S}$ 局部有界。因此，结论由引理
\ref{derham-lemma-de-rham-complex-product} 和《概形的导出范畴》，引理
\ref{perfect-lemma-kunneth-special} 推出。
\end{proof}

\noindent
杯积还有相对版本，即映射
$$
Ra_*\Omega^\bullet_{X/S}
\otimes_{\mathcal{O}_S}^\mathbf{L}
Rb_*\Omega^\bullet_{Y/S}
\longrightarrow
Rf_*\Omega^\bullet_{X \times_S Y/S}
$$
它位于 $D(\mathcal{O}_S)$ 中。该构造把引理
\ref{derham-lemma-de-rham-complex-product} 与《概形的导出范畴》中的映射
（方程 \ref{perfect-equation-relative-de-rham-kunneth}）结合起来。
由此构造可见，该映射由下图给出：
$$
\xymatrix{
Ra_*\Omega^\bullet_{X/S}
\otimes_{\mathcal{O}_S}^\mathbf{L}
Rb_*\Omega^\bullet_{Y/S}
\ar[d]^{\text{伴随单位}}  \\
Rf_*(p^{-1}\Omega^\bullet_{X/S})
\otimes_{\mathcal{O}_S}^\mathbf{L}
Rf_*(q^{-1}\Omega^\bullet_{Y/S}) \ar[r] \ar[d]^{\text{相对杯积}} &
Rf_*(\Omega^\bullet_{X \times_S Y/S})
\otimes_{\mathcal{O}_S}^\mathbf{L}
Rf_*(\Omega^\bullet_{X \times_S Y/S}) \ar[d]^{\text{相对杯积}} \\
Rf_*(p^{-1}\Omega^\bullet_{X/S}
\otimes_{f^{-1}\mathcal{O}_S}^\mathbf{L}
q^{-1}\Omega^\bullet_{Y/S})
\ar[d]^{\text{从导出函子到通常函子}} \ar[r] &
Rf_*(\Omega^\bullet_{X \times_S Y/S}
\otimes_{f^{-1}\mathcal{O}_S}^\mathbf{L}
\Omega^\bullet_{X \times_S Y/S})
\ar[d]^{\text{从导出函子到通常函子}} \\
Rf_*\text{Tot}(p^{-1}\Omega^\bullet_{X/S}
\otimes_{f^{-1}\mathcal{O}_S}
q^{-1}\Omega^\bullet_{Y/S}) \ar[r] \ar[d]^{\text{典范映射}} &
Rf_*\text{Tot}(\Omega^\bullet_{X \times_S Y/S}
\otimes_{f^{-1}\mathcal{O}_S}
\Omega^\bullet_{X \times_S Y/S})
\ar[d]^{\eta \otimes \omega \mapsto \eta \wedge \omega} \\
Rf_*\text{Tot}(\Omega^\bullet_{X/S} \boxtimes \Omega^\bullet_{Y/S})
\ar@{=}[r]
&
Rf_*\Omega^\bullet_{X \times_S Y/S}
}
$$
这里，第一个箭头使用伴随的单位 $\text{id} \to Rp_* p^{-1}$
和 $\text{id} \to Rq_* q^{-1}$，以及等同
$Rf_* p^{-1} = Ra_* Rp_* p^{-1}$ 和
$Rf_* q^{-1} = Rb_* Rq_* q^{-1}$。
第二个箭头是《上同调》，备注 \ref{cohomology-remark-cup-product} 中的相对杯积。
第三个箭头把复形的导出张量积送到复形张量积的全复形。
最后一个等号来自引理 \ref{derham-lemma-de-rham-complex-product}。
在整体截面上，该构造恢复了前面给出的构造。

\begin{lemma}
\label{derham-lemma-kunneth-de-rham-relative}
假设 $X \to S$ 和 $Y \to S$ 光滑且拟紧，并且态射
$X \to X \times_S X$ 和 $Y \to Y \times_S Y$ 仿射。则相对杯积
$$
Ra_*\Omega^\bullet_{X/S}
\otimes_{\mathcal{O}_S}^\mathbf{L}
Rb_*\Omega^\bullet_{Y/S}
\longrightarrow
Rf_*\Omega^\bullet_{X \times_S Y/S}
$$
是 $D(\mathcal{O}_S)$ 中的同构。
\end{lemma}

\begin{proof}
这是引理 \ref{derham-lemma-kunneth-de-rham} 的直接推论。
\end{proof}







\section{德拉姆上同调中的第一陈类}
\label{derham-section-first-chern-class}

\noindent
设 $X \to S$ 是概形态射。存在复形映射
$$
\text{d}\log : \mathcal{O}_X^*[-1] \longrightarrow \Omega^\bullet_{X/S}
$$
它把截面 $g \in \mathcal{O}_X^*(U)$ 送到截面
$\text{d}\log(g) = g^{-1}\text{d}g$，后者属于 $\Omega^1_{X/S}(U)$。
因此可以考虑映射
$$
\Pic(X) = H^1(X, \mathcal{O}_X^*) =
H^2(X, \mathcal{O}_X^*[-1]) \longrightarrow H^2_{dR}(X/S)
$$
其中第一个等号来自《上同调》，引理
\ref{cohomology-lemma-h1-invertible}。
可逆模 $\mathcal{L}$ 的同构类的像记为
$c^{dR}_1(\mathcal{L}) \in H^2_{dR}(X/S)$。

\medskip\noindent
还可以利用映射 $\text{d}\log : \mathcal{O}_X^* \to \Omega^1_{X/S}$
在霍奇上同调中定义陈类
$$
c_1^{Hodge} : \Pic(X) \longrightarrow H^1(X, \Omega^1_{X/S})
\subset H^2_{Hodge}(X/S)
$$
这些构造与拉回相容。

\begin{lemma}
\label{derham-lemma-pullback-c1}
给定概形的交换图
$$
\xymatrix{
X' \ar[r]_f \ar[d] & X \ar[d] \\
S' \ar[r] & S
}
$$
则下列图
$$
\xymatrix{
\Pic(X') \ar[d]_{c_1^{dR}} &
\Pic(X) \ar[d]^{c_1^{dR}} \ar[l]^{f^*} \\
H^2_{dR}(X'/S') &
H^2_{dR}(X/S) \ar[l]_{f^*}
}
\quad
\xymatrix{
\Pic(X') \ar[d]_{c_1^{Hodge}} &
\Pic(X) \ar[d]^{c_1^{Hodge}} \ar[l]^{f^*} \\
H^1(X', \Omega^1_{X'/S'}) &
H^1(X, \Omega^1_{X/S}) \ar[l]_{f^*}
}
$$
均交换。
\end{lemma}

\begin{proof}
略。
\end{proof}

\noindent
让我们在 {\v C}ech 上同调中“计算”元素 $c^{dR}_1(\mathcal{L})$
（{\v C}ech 微分的符号规定同《上同调》，第
\ref{cohomology-section-cech-cohomology-of-complexes} 节）。
具体地，选取开覆盖
$\mathcal{U} : X = \bigcup_{i \in I} U_i$，使得有平凡化截面
$s_i$，它是 $\mathcal{L}|_{U_i}$ 的截面；这对所有 $i$ 都成立。
在交集 $U_{i_0i_1} = U_{i_0} \cap U_{i_1}$ 上，存在可逆函数
$f_{i_0i_1}$，使得
$f_{i_0i_1} = s_{i_1}|_{U_{i_0i_1}} s_{i_0}|_{U_{i_0i_1}}^{-1}$\footnote{
$0$-上闭链 $\{a_{i_0}\}$ 的 {\v C}ech 微分为
$a_{i_1} - a_{i_0}$，它定义在 $U_{i_0i_1}$ 上。}。
当然有
$$
f_{i_1i_2}|_{U_{i_0i_1i_2}}
f_{i_0i_2}^{-1}|_{U_{i_0i_1i_2}}
f_{i_0i_1}|_{U_{i_0i_1i_2}} = 1
$$
$\mathcal{L}$ 在 $H^1(X, \mathcal{O}_X^*)$ 中的上同调类，
是上闭链 $\{f_{i_0i_1}\}$ 在
$\check{\mathcal{C}}^\bullet(\mathcal{U}, \mathcal{O}_X^*)$
中的 {\v C}ech 上同调类之像。因此，$c_1^{dR}(\mathcal{L})$ 是
与 {\v C}ech 上闭链 $\{\alpha_{i_0 \ldots i_p}\}$ 相伴的上同调类之像；
该上闭链位于
$\text{Tot}(\check{\mathcal{C}}^\bullet(\mathcal{U}, \Omega_{X/S}^\bullet))$
中、次数为 $2$，并由下列各项给出：
\begin{enumerate}
\item $\alpha_{i_0} = 0$ 属于 $\Omega^2_{X/S}(U_{i_0})$；
\item $\alpha_{i_0i_1} = f_{i_0i_1}^{-1}\text{d}f_{i_0i_1}$ 属于
$\Omega^1_{X/S}(U_{i_0i_1})$；并且
\item $\alpha_{i_0i_1i_2} = 0$ 属于 $\mathcal{O}_{X/S}(U_{i_0i_1i_2})$。
\end{enumerate}
假设有可逆模 $\mathcal{L}_k$（$k = 1, \ldots, a$），
每个模在所有 $U_i$（$i \in I$）上都已平凡化，并由此得到上述上闭链
$f_{k, i_0i_1}$ 和 $\alpha_k = \{\alpha_{k, i_0 \ldots i_p}\}$。
利用《上同调》，第
\ref{cohomology-section-cech-cohomology-of-complexes} 节中的法则，可以算出
$$
\beta = \alpha_1 \cup \alpha_2 \cup \ldots \cup \alpha_a
$$
它由上闭链 $\beta = \{\beta_{i_0 \ldots i_p}\}$ 给出；后者描述如下：
\begin{enumerate}
\item $\beta_{i_0 \ldots i_p} = 0$ 属于
$\Omega^{2a - p}_{X/S}(U_{i_0 \ldots i_p})$，除非 $p = a$；并且
\item $\beta_{i_0 \ldots i_a} = (-1)^{a(a - 1)/2}
\alpha_{1, i_0i_1} \wedge \alpha_{2, i_1 i_2} \wedge \ldots \wedge
\alpha_{a, i_{a - 1}i_a}$ 属于
$\Omega^a_{X/S}(U_{i_0 \ldots i_a})$。
\end{enumerate}
因此，这是表示
$c_1^{dR}(\mathcal{L}_1) \cup \ldots \cup c_1^{dR}(\mathcal{L}_a)$ 的上闭链。
当然，同一计算还表明：上闭链
$\{\beta_{i_0 \ldots i_a}\}$ 位于
$\check{\mathcal{C}}^a(\mathcal{U}, \Omega_{X/S}^a))$ 中，并表示上同调类
$c_1^{Hodge}(\mathcal{L}_1) \cup \ldots \cup c_1^{Hodge}(\mathcal{L}_a)$。

\begin{remark}
\label{derham-remark-truncations}
下面用更抽象的语言重新表述上述计算。
设 $p : X \to S$ 是概形态射，并设 $\mathcal{L}$ 是可逆
$\mathcal{O}_X$-模。若把 $\text{d}\log$ 看作映射
$$
\mathcal{O}_X^*[-1] \to \sigma_{\geq 1}\Omega^\bullet_{X/S}
$$
则如上利用 $\Pic(X) = H^1(X, \mathcal{O}_X^*)$，得到上同调类
$$
\gamma_1(\mathcal{L}) \in H^2(X, \sigma_{\geq 1}\Omega^\bullet_{X/S})
$$
$\gamma_1(\mathcal{L})$ 在映射
$\sigma_{\geq 1}\Omega^\bullet_{X/S} \to \Omega^\bullet_{X/S}$
下的像恢复 $c_1^{dR}(\mathcal{L})$。特别地，
$c_1^{dR}(\mathcal{L}) \in F^1H^2_{dR}(X/S)$, see
参见第 \ref{derham-section-hodge-filtration} 节。$\gamma_1(\mathcal{L})$ 在映射
$\sigma_{\geq 1}\Omega^\bullet_{X/S} \to \Omega^1_{X/S}[-1]$ 下的像
恢复 $c_1^{Hodge}(\mathcal{L})$。
取杯积（参见第 \ref{derham-section-hodge-filtration} 节），得到
$$
\xi = \gamma_1(\mathcal{L}_1) \cup \ldots \cup \gamma_1(\mathcal{L}_a) \in
H^{2a}(X, \sigma_{\geq a}\Omega^\bullet_{X/S})
$$
第 \ref{derham-section-hodge-filtration} 节中的交换图表明，$\xi$ 被映到
$c_1^{dR}(\mathcal{L}_1) \cup \ldots \cup c_1^{dR}(\mathcal{L}_a)$
，后者位于 $H^{2a}_{dR}(X/S)$ 中；这里所用的映射为
$\sigma_{\geq a}\Omega^\bullet_{X/S} \to \Omega^\bullet_{X/S}$。
还可知
$c_1^{dR}(\mathcal{L}_1) \cup \ldots \cup c_1^{dR}(\mathcal{L}_a)$
属于 $F^a H^{2a}_{dR}(X/S)$。类似地，映射
$\sigma_{\geq a}\Omega^\bullet_{X/S} \to \Omega^a_{X/S}[-a]$
把 $\xi$ 送到
$c_1^{Hodge}(\mathcal{L}_1) \cup \ldots \cup c_1^{Hodge}(\mathcal{L}_a)$
；后者位于 $H^a(X, \Omega^a_{X/S})$ 中。
\end{remark}

\begin{remark}
\label{derham-remark-log-forms}
设 $p : X \to S$ 是概形态射。对 $i > 0$，以
$\Omega^i_{X/S, log} \subset \Omega^i_{X/S}$ 表示由下述形式的局部截面
生成的阿贝尔子层：
$$
\text{d}\log(u_1) \wedge \ldots \wedge \text{d}\log(u_i)
$$
其中 $u_1, \ldots, u_n$ 是 $\mathcal{O}_X$ 的可逆局部截面。
当 $i = 0$ 时，子层 $\Omega^0_{X/S, log} \subset \mathcal{O}_X$
是 $\mathbf{Z} \to \mathcal{O}_X$ 的像。对每个 $i \geq 0$，都有复形映射
$$
\Omega^i_{X/S, log}[-i] \longrightarrow \Omega^\bullet_{X/S}
$$
因为对数形式的微分为零。此外，对数形式的楔积仍是对数形式，因而得到双线性映射
$$
\wedge :  \Omega^i_{X/S, log} \times
\Omega^j_{X/S, log} \longrightarrow \Omega^{i + j}_{X/S, log}
$$
它们与（\ref{derham-equation-wedge}）及上述映射相容。
设 $\mathcal{L}$ 是可逆 $\mathcal{O}_X$-模。利用阿贝尔层的映射
$\text{d}\log : \mathcal{O}_X^* \to \Omega^1_{X/S, log}$
和等同 $\Pic(X) = H^1(X, \mathcal{O}_X^*)$，得到典范上同调类
$$
\tilde \gamma_1(\mathcal{L}) \in H^1(X, \Omega^1_{X/S, log})
$$
这些类具有下列性质：
\begin{enumerate}
\item $\tilde \gamma_1(\mathcal{L})$ 在典范映射
$\Omega^1_{X/S, log}[-1] \to \sigma_{\geq 1}\Omega^\bullet_{X/S}$
下的像，是该映射把 $\tilde \gamma_1(\mathcal{L})$ 送到的类
$\gamma_1(\mathcal{L}) \in 
H^2(X, \sigma_{\geq 1}\Omega^\bullet_{X/S})$；
这就是备注 \ref{derham-remark-truncations} 中的类；
\item $\tilde \gamma_1(\mathcal{L})$ 在典范映射
$\Omega^1_{X/S, log}[-1] \to \Omega^\bullet_{X/S}$
下的像，是该映射把 $\tilde \gamma_1(\mathcal{L})$ 送到的
$c_1^{dR}(\mathcal{L})$，它位于 $H^2_{dR}(X/S)$ 中；
\item $\tilde \gamma_1(\mathcal{L})$ 在典范映射
$\Omega^1_{X/S, log} \to \Omega^1_{X/S}$
下的像，是该映射把 $\tilde \gamma_1(\mathcal{L})$ 送到的
$c_1^{Hodge}(\mathcal{L})$，它位于 $H^1(X, \Omega^1_{X/S})$ 中；
\item 这些类的构造与拉回相容；
\item 在此添加更多内容。
\end{enumerate}
\end{remark}





\section{线丛的德拉姆上同调}
\label{derham-section-line-bundle}

\noindent
线丛是向量丛的特例，而向量丛又是带有某些附加结构的锥。
为了恰当地讨论这些对象的德拉姆复形，先讨论分次环的德拉姆复形是合适的。

\begin{remark}[分次环的德拉姆复形]
\label{derham-remark-de-rham-complex-graded}
设 $G$ 是以加法记号书写、幺元为 $0$ 的阿贝尔幺半群。
设 $R \to A$ 是环同态，并假设 $A$ 带有分次
$A = \bigoplus_{g \in G} A_g$；它由 $R$-模给出，使得 $R$ 映入 $A_0$，且
$A_g \cdot A_{g'} \subset A_{g + g'}$。于是微分模带有分次
$$
\Omega_{A/R} = \bigoplus\nolimits_{g \in G} \Omega_{A/R, g}
$$
其中 $\Omega_{A/R, g}$ 是 $R$-子模；它包含于 $\Omega_{A/R}$，并
由 $a_0 \text{d}a_1$ 生成；这里 $a_i \in A_{g_i}$，且
$g = g_0 + g_1$。类似地，得到
$$
\Omega^p_{A/R} = \bigoplus\nolimits_{g \in G} \Omega^p_{A/R, g}
$$
其中 $\Omega^p_{A/R, g}$ 是 $R$-子模；它包含于 $\Omega^p_{A/R}$，并
由 $a_0 \text{d}a_1 \wedge \ldots \wedge \text{d}a_p$ 生成；
这里 $a_i \in A_{g_i}$，且 $g = g_0 + g_1 + \ldots + g_p$。
显然，微分保持分次，而楔积以显然的方式与分次相容。
\end{remark}

\noindent
设 $f : X \to S$ 是概形态射，并设 $\pi : C \to X$ 是一个锥；参见
《构造》，定义 \ref{constructions-definition-abstract-cone}。
回忆这意味着 $\pi$ 仿射，且有分次
$\pi_*\mathcal{O}_C = \bigoplus_{n \geq 0} \mathcal{A}_n$，其中
$\mathcal{A}_0 = \mathcal{O}_X$。
在仿射开集上应用备注 \ref{derham-remark-de-rham-complex-graded} 中的讨论，可得
\footnote{请原谅这里的记号！}
$$
\pi_*(\Omega^\bullet_{C/S}) =
\bigoplus\nolimits_{n \geq 0} \Omega^\bullet_{C/S, n}
$$
典范地分解为子复形的直和。此外，还有分解
$$
\Omega^\bullet_{X/S} \to \Omega^\bullet_{C/S, 0} \to
\pi_*(\Omega^\bullet_{C/S})
$$
并且知道 $\omega \wedge \eta \in \Omega^{p + q}_{C/S, n + m}$，
只要 $\omega \in \Omega^p_{C/S, n}$ 且
$\eta \in \Omega^q_{C/S, m}$。

\medskip\noindent
设 $f : X \to S$ 是概形态射，并设 $\pi : L \to X$ 是与可逆
$\mathcal{O}_X$-模 $\mathcal{L}$ 相伴的线丛。这意味着 $\pi$
是满足下式的唯一仿射态射：
$$
\pi_*\mathcal{O}_L = 
\bigoplus\nolimits_{n \geq 0} \mathcal{L}^{\otimes n}
$$
这里等式指 $\mathcal{O}_X$-代数的等式。因此 $L$ 是 $X$ 上的锥。
由上述讨论得到典范直和分解
$$
\pi_*(\Omega^\bullet_{L/S}) =
\bigoplus\nolimits_{n \geq 0} \Omega^\bullet_{L/S, n}
$$
它与楔积以及上述 $\pi_*\mathcal{O}_L$ 的分解相容，并使
$\Omega_{X/S}$ 映入部分 $\Omega_{L/S, 0}$；该部分的次数为 $0$。

\medskip\noindent
还有一种情形对我们有用。考虑补集\footnote{概形 $L^\star$ 是
$\mathbf{G}_m$-扭子；它位于 $X$ 上并与 $L$ 相伴。因此下面所得分次是
$\mathbf{Z}$-分次；比较《群胚》，例
\ref{groupoids-example-Gm-on-affine}，以及引理
\ref{groupoids-lemma-complete-reducibility-Gm} 和
\ref{groupoids-lemma-Gm-equivariant-module}。}
$L^\star \subset L$，它是线丛中零截面 $o : X \to L$ 的补，而该线丛为
$L$。局部计算表明，存在典范同构
$$
(L^\star \to X)_*\mathcal{O}_{L^\star} =
\bigoplus\nolimits_{n \in \mathbf{Z}} \mathcal{L}^{\otimes n}
$$
这是 $\mathcal{O}_X$-代数的同构。右端是 $\mathbf{Z}$-分次拟凝聚
$\mathcal{O}_X$-代数。在仿射开集上使用备注
\ref{derham-remark-de-rham-complex-graded} 中的讨论，得到
$$
(L^\star \to X)_*(\Omega^\bullet_{L^\star/S}) =
\bigoplus\nolimits_{n \in \mathbf{Z}} \Omega^\bullet_{L^\star/S, n}
$$
它与楔积以及上述 $(L^\star \to X)_*\mathcal{O}_{L^\star}$ 的分解相容，
并使 $\Omega_{X/S}$ 映入部分 $\Omega_{L^\star/S, 0}$；该部分的次数为 $0$。
复形 $\Omega^\bullet_{L^\star/S, 0}$ 对我们尤为重要。

\begin{lemma}
\label{derham-lemma-the-complex-for-L-star}
沿用上述记号，存在复形的短正合列
$$
0 \to \Omega^\bullet_{X/S} \to
\Omega^\bullet_{L^\star/S, 0} \to
\Omega^\bullet_{X/S}[-1] \to 0
$$
\end{lemma}

\begin{proof}
上面已经构造了映射
$\Omega^\bullet_{X/S} \to \Omega^\bullet_{L^\star/S, 0}$。

\medskip\noindent
构造
$\text{Res} : \Omega^\bullet_{L^\star/S, 0} \to \Omega^\bullet_{X/S}[-1]$。
设 $U \subset X$ 是开集，并设 $s \in \mathcal{L}(U)$ 和
$s' \in \mathcal{L}^{\otimes -1}(U)$ 是满足 $s' s = 1$ 的截面。
于是 $s$ 给出代数层 $(L^\star \to X)_*\mathcal{O}_{L^\star}$
在 $U$ 上的可逆截面，其逆为 $s' = s^{-1}$。因而可以考虑 $1$-形式
$\text{d}\log(s) = s' \text{d}(s)$；它属于
$\Omega^1_{L^\star/S, 0}(U)$，这是由我们在
$\Omega^1_{L^\star/S}$ 上构造的分次得到的。下面的仿射计算将表明，$1$ 和
$\text{d}\log(s)$ 自由生成 $\Omega^\bullet_{L^\star/S, 0}|_U$；
这里采用其作为 $\Omega^\bullet_{X/S}|_U$ 上右模的结构。
因此，可以按下列法则定义 $\text{Res}$；它定义在 $U$ 上：
$$
\text{Res}(\omega' + \text{d}\log(s) \wedge \omega) = \omega
$$
这里 $\omega', \omega \in \Omega^\bullet_{X/S}(U)$ 任意。
该映射与局部生成元 $s$ 的选择无关，因而可以黏合成整体映射。
事实上，$s$ 的另一个选择必为 $gs$ 的形式，其中
$g \in \mathcal{O}_X(U)$ 可逆；此时
$\text{d}\log(gs) = g^{-1}\text{d}(g) + \text{d}\log(s)$，
由此容易看出无关性。最后注意，$\text{Res}$ 的定义法则与微分相容，因为
$\text{d}(\omega' + \text{d}\log(s) \wedge \omega) =
\text{d}(\omega') - \text{d}\log(s) \wedge \text{d}(\omega)$
，并且按照《同调》，定义
\ref{homology-definition-shift-cochain} 中的符号约定，
$\Omega^\bullet_{X/S}[-1]$ 上的微分把 $\omega'$ 送到
$-\text{d}(\omega')$。

\medskip\noindent
局部计算。可以用仿射开集覆盖 $X$，这些开集记为 $U \subset X$，
并满足 $\mathcal{L}|_U \cong \mathcal{O}_U$，而且映入某个仿射开集
$V \subset S$。写成 $U = \Spec(A)$、$V = \Spec(R)$，并选取
生成元 $s$；它属于 $\mathcal{L}$。于是有
$$
L^\star \times_X U = \Spec(A[s, s^{-1}])
$$
计算微分可得
$$
\Omega^1_{A[s, s^{-1}]/R} =
A[s, s^{-1}] \otimes_A \Omega^1_{A/R} \oplus A[s, s^{-1}] \text{d}\log(s)
$$
因而取外幂得到
$$
\Omega^p_{A[s, s^{-1}]/R} =
A[s, s^{-1}] \otimes_A \Omega^p_{A/R}
\oplus
A[s, s^{-1}] \text{d}\log(s) \otimes_A \Omega^{p - 1}_{A/R}
$$
取次数为 $0$ 的部分，得到
$$
\Omega^p_{A[s, s^{-1}]/R, 0} =
\Omega^p_{A/R} \oplus \text{d}\log(s) \otimes_A \Omega^{p - 1}_{A/R}
$$
引理得证。
\end{proof}

\begin{lemma}
\label{derham-lemma-the-complex-for-L-star-gives-chern-class}
“边缘”映射
$\delta : \Omega^\bullet_{X/S} \to \Omega^\bullet_{X/S}[2]$
位于 $D(X, f^{-1}\mathcal{O}_S)$ 中，它来自引理
\ref{derham-lemma-the-complex-for-L-star} 的短正合列；它就是备注
\ref{derham-remark-cup-product-as-a-map} 中对应于
$\xi = c_1^{dR}(\mathcal{L})$ 的映射。
\end{lemma}

\begin{proof}
准确地说，考虑短正合列的平移
$$
0 \to \Omega^\bullet_{X/S}[1] \to
\Omega^\bullet_{L^\star/S, 0}[1] \to
\Omega^\bullet_{X/S} \to 0
$$
该短正合列来自引理 \ref{derham-lemma-the-complex-for-L-star}。
作为下列分次的次数零部分：
$(L^\star \to X)_*\Omega^\bullet_{L^\star/S}$
，$\Omega^\bullet_{L^\star/S, 0}$ 是微分分次
$\mathcal{O}_X$-代数，并且映射
$\Omega^\bullet_{X/S} \to \Omega^\bullet_{L^\star/S, 0}$
是微分分次 $\mathcal{O}_X$-代数的同态。
因此可以把 $\Omega^\bullet_{X/S}[1] \to
\Omega^\bullet_{L^\star/S, 0}[1]$ 视为右微分分次
$\Omega^\bullet_{X/S}$-模的映射；它定义在 $X$ 上。映射
$\text{Res} : \Omega^\bullet_{L^\star/S, 0}[1] \to \Omega^\bullet_{X/S}$
是右微分分次 $\Omega^\bullet_{X/S}$-模的映射，因为它局部地由法则
$\text{Res}(\omega' + \text{d}\log(s) \wedge \omega) = \omega$ 定义；
定义；参见引理 \ref{derham-lemma-the-complex-for-L-star} 的证明。
因此，由《微分分次层》，第 \ref{sdga-section-misc} 节中的讨论可知，
$\delta$ 来自映射
$\delta' : \Omega^\bullet_{X/S} \to \Omega^\bullet_{X/S}[2]$
；它位于德拉姆复形上的右微分分次模的导出范畴
$D(\Omega^\bullet_{X/S}, \text{d})$ 中。备注
\ref{derham-remark-cup-product-as-a-map} 所述的唯一性表明，只需证明
$\delta(1) = c_1^{dR}(\mathcal{L})$。

\medskip\noindent
我们断言存在交换图
$$
\xymatrix{
0 \ar[r] &
\mathcal{O}_X^* \ar[r] \ar[d]_{\text{d}\log} &
E \ar[r] \ar[d] &
\underline{\mathbf{Z}} \ar[d] \ar[r] &
0 \\
0 \ar[r] &
\Omega^\bullet_{X/S}[1] \ar[r] &
\Omega^\bullet_{L^\star/S, 0}[1] \ar[r] &
\Omega^\bullet_{X/S} \ar[r] &
0
}
$$
其中上行是阿贝尔层的短正合列，其边缘映射把 $1$ 送到
$\mathcal{L}$ 在 $H^1(X, \mathcal{O}_X^*)$ 中的类。
由边缘映射同短正合列之间映射的相容性，只需证明该断言。
把 $E$ 定义为下述法则的层化：
$$
U \longmapsto \{(s, n) \mid
n \in \mathbf{Z},\ s \in \mathcal{L}^{\otimes n}(U)\text{ 生成元}\}
$$
其群结构由 $(s, n) \cdot (t, m) = (s \otimes t, n + m)$ 给出。
中间的竖直映射把 $(s, n)$ 送到 $\text{d}\log(s)$。
这给出短正合列之间的映射，因为证明引理
\ref{derham-lemma-the-complex-for-L-star} 时构造的映射
$Res : \Omega^1_{L^\star/S, 0} \to \mathcal{O}_X$ 把
$\text{d}\log(s)$ 送到 $1$，只要 $s$ 是 $\mathcal{L}$ 的局部生成元。
为了计算上行中 $1$ 的边缘，按照第
\ref{derham-section-first-chern-class} 节，选取局部平凡化 $s_i$；
它们来自 $\mathcal{L}$，并定义在开集 $U_i$ 上。在交集
$U_{i_0i_1} = U_{i_0} \cap U_{i_1}$
上，存在可逆函数 $f_{i_0i_1}$，使得
$f_{i_0i_1} = s_{i_1}|_{U_{i_0i_1}} s_{i_0}|_{U_{i_0i_1}}^{-1}$
，而 $\mathcal{L}$ 的上同调类由 {\v C}ech 上闭链
$\{f_{i_0i_1}\}$ 给出。于是当然有
$$
(f_{i_0i_1}, 0) = (s_{i_1}, 1)|_{U_{i_0i_1}} \cdot
(s_{i_0}, 1)|_{U_{i_0i_1}}^{-1}
$$
它们是 $E$ 的截面。这就完成了证明。
\end{proof}

\begin{lemma}
\label{derham-lemma-push-omega-a}
沿用上述记号，有
\begin{enumerate}
\item $\Omega^p_{L^\star/S, n} =
\Omega^p_{L^\star/S, 0} \otimes_{\mathcal{O}_X} \mathcal{L}^{\otimes n}$
该等式对所有 $n \in \mathbf{Z}$ 成立；这里把两边视为拟凝聚
$\mathcal{O}_X$-模；
\item $\Omega^\bullet_{X/S} = \Omega^\bullet_{L/X, 0}$
这是复形的等式；并且
\item 对 $n > 0$ 和 $p \geq 0$，有
$\Omega^p_{L/X, n} = \Omega^p_{L^\star/S, n}$.
\end{enumerate}
\end{lemma}

\begin{proof}
在每一种情形下，都存在整体定义的典范映射；略去的局部计算表明它是同构。
\end{proof}

\begin{lemma}
\label{derham-lemma-line-bundle-characteristic-zero}
在上述情形下，假设存在态射 $S \to \Spec(\mathbf{Q})$。
则 $\Omega^\bullet_{X/S} \to \pi_*\Omega^\bullet_{L/S}$ 是拟同构，并且
$H_{dR}^*(X/S) = H_{dR}^*(L/S)$。
\end{lemma}

\begin{proof}
设 $R$ 是 $\mathbf{Q}$-代数，并设 $A$ 是 $R$-代数。
仿射局部的陈述是：映射
$$
\Omega^\bullet_{A/R} \longrightarrow \Omega^\bullet_{A[t]/R}
$$
是 $R$-模复形的拟同构。事实上，它是同伦等价；其同伦逆由把
$g \omega + g' \text{d}t \wedge \omega'$ 送到 $g(0)\omega$ 的映射给出；
这里 $g, g' \in A[t]$，且 $\omega, \omega' \in \Omega^\bullet_{A/R}$。
该同伦把 $g \omega + g' \text{d}t \wedge \omega'$ 送到
$(\int g') \omega'$；其中 $\int g' \in A[t]$ 是常数项为零、
关于 $t$ 的导数为 $g'$ 的多项式。当然，这里用到 $R$ 包含
$\mathbf{Q}$，因为 $\int t^n = (1/n)t^{n + 1}$。
\end{proof}

\begin{example}
\label{derham-example-affine-line}
引理 \ref{derham-lemma-line-bundle-characteristic-zero} 在正特征下不成立。
$\mathbf{A}^1_k = \Spec(k[x])$ 在域 $k$ 上的德拉姆复形具有如下直和形式：
$$
k \oplus
\bigoplus\nolimits_{n \geq 1}
(k \cdot t^n \xrightarrow{n}
k \cdot t^{n - 1} \text{d}t)
$$
因此，若 $k$ 的特征为 $p > 0$，则
$H^0_{dR}(\mathbf{A}^1_k/k)$ 和
$H^1_{dR}(\mathbf{A}^1_k/k)$
作为 $k$-向量空间都是无限维的。
\end{example}








\section{射影空间的德拉姆上同调}
\label{derham-section-projective-space}

\noindent
设 $A$ 是环，且 $n \geq 1$。结构态射
$\mathbf{P}^n_A \to \Spec(A)$ 是相对维数为 $n$ 的真光滑态射。
它光滑、相对维数为 $n$ 且为有限型，因为 $\mathbf{P}^n_A$ 有有限仿射开覆盖，
其中每个概形都同构于 $\mathbf{A}^n_A$；参见《构造》，引理
\ref{constructions-lemma-standard-covering-projective-space}。
由《构造》，引理 \ref{constructions-lemma-projective-space-separated}，
它还分离且普遍闭，因而为真。
用 $\mathcal{O}$ 和 $\mathcal{O}(d)$ 分别表示结构层
$\mathcal{O}_{\mathbf{P}^n_A}$ 与 Serre 扭曲
$\mathcal{O}_{\mathbf{P}^n_A}(d)$。
用 $\Omega = \Omega_{\mathbf{P}^n_A/A}$ 表示相对微分层，并用
$\Omega^p$ 表示其外幂。

\begin{lemma}
\label{derham-lemma-euler-sequence}
存在短正合列
$$
0 \to \Omega \to \mathcal{O}(-1)^{\oplus n + 1} \to \mathcal{O} \to 0
$$
\end{lemma}

\begin{proof}
为说明这一点，回忆
$\mathbf{P}^n_A = \text{Proj}(A[T_0, \ldots, T_n])$,
并采用形式记号
$$
\mathcal{O}(-1)^{\oplus n + 1} =
\bigoplus\nolimits_{j = 0, \ldots, n} \mathcal{O}(-1) \text{d}T_j
$$
上述短正合列中的第一个箭头
$$
\Omega \to
\bigoplus\nolimits_{j = 0, \ldots, n} \mathcal{O}(-1) \text{d}T_j
$$
在每个标准开集
$D_+(T_i) = \Spec(A[T_0/T_i, \ldots, T_n/T_i])$
上由下列法则给出：
$$
\sum\nolimits_{j \not = i} g_j \text{d}(T_j/T_i)
\longmapsto
\sum\nolimits_{j \not = i} g_j/T_i \text{d}T_j
- (\sum\nolimits_{j \not = i} g_jT_j/T_i^2) \text{d}T_i
$$
这是有意义的，因为 $1/T_i$ 是 $\mathcal{O}(-1)$ 在
$D_+(T_i)$ 上的截面。映射
$$
\bigoplus\nolimits_{j = 0, \ldots, n} \mathcal{O}(-1) \text{d}T_j
\to
\mathcal{O}
$$
把 $\text{d}T_j$ 送到 $T_j$。更准确地说，在 $D_+(T_i)$ 上，
把截面 $\sum g_j \text{d}T_j$ 送到 $\sum T_jg_j$。
略去验证这确实给出短正合列。
\end{proof}

\noindent
给定整数 $k \in \mathbf{Z}$ 和拟凝聚
$\mathcal{O}_{\mathbf{P}^n_A}$-模 $\mathcal{F}$，照例用
$\mathcal{F}(k)$ 表示第 $k$ 个 Serre 扭曲；被扭曲的模是 $\mathcal{F}$。
参见《构造》，定义 \ref{constructions-definition-twist}。

\begin{lemma}
\label{derham-lemma-twisted-hodge-cohomology-projective-space}
在上述情形下，上同调群满足：
\begin{enumerate}
\item $H^q(\mathbf{P}^n_A, \Omega^p) = 0$，除非
$0 \leq p = q \leq n$；
\item 对 $0 \leq p \leq n$，$A$-模
$H^p(\mathbf{P}^n_A, \Omega^p)$ 是秩为 $1$ 的自由模；
\item 对 $q > 0$、$k > 0$ 和任意 $p$，有
$H^q(\mathbf{P}^n_A, \Omega^p(k)) = 0$；并且
\item 在此添加更多内容。
\end{enumerate}
\end{lemma}

\begin{proof}
以下将不再逐次说明地使用《概形上同调》，引理
\ref{coherent-lemma-cohomology-projective-space-over-ring} 的结论。
特别地，相应陈述对 $H^q(\mathbf{P}^n_A, \mathcal{O}(k))$ 成立。

\medskip\noindent
先证明 $p = 1$ 的情形。考虑短正合列
$$
0 \to \Omega \to \mathcal{O}(-1)^{\oplus n + 1} \to \mathcal{O} \to 0
$$
它来自引理 \ref{derham-lemma-euler-sequence}。由于 $\mathcal{O}(-1)$
在所有次数上的上同调均为零，所以
$H^q(\mathbf{P}^n_A, \Omega)$ 仅在次数 $1$ 可能非零；
在该次数，它由 $1$ 的边缘自由生成；该元素位于
$H^0(\mathbf{P}^n_A, \mathcal{O})$ 中。

\medskip\noindent
假设 $p > 1$。把上述短正合列看作定义了一个 $2$ 步滤过；
该滤过位于 $\mathcal{O}(-1)^{\oplus n + 1}$ 上。
它在 $\wedge^p\mathcal{O}(-1)^{\oplus n + 1}$ 上诱导的滤过如下：
$$
0 \to \Omega^p \to \wedge^p\left(\mathcal{O}(-1)^{\oplus n + 1}\right)
\to \Omega^{p - 1} \to 0
$$
注意，$\wedge^p\mathcal{O}(-1)^{\oplus n + 1}$ 同构于
$n + 1$ 取 $p$ 个 $\mathcal{O}(-p)$ 副本的直和，
因而在所有次数上的上同调均为零。由归纳假设可知，
$H^q(\mathbf{P}^n_A, \Omega^p)$ 为零，除非 $q = p$；而
$H^p(\mathbf{P}^n_A, \Omega^p)$ 是秩为 $1$ 的自由模，其生成元是
$H^{p - 1}(\mathbf{P}^n_A, \Omega^{p - 1})$ 中生成元的边缘。

\medskip\noindent
设 $k > 0$。注意 $\Omega^n = \mathcal{O}(-n - 1)$；例如，
这可由上述短正合列取 $p = n + 1$ 得到。因此 $\Omega^n(k)$
在正次数上的上同调均为零。利用短正合列
$$
0 \to \Omega^p(k) \to \wedge^p\left(\mathcal{O}(-1)^{\oplus n + 1}\right)(k)
\to \Omega^{p - 1}(k) \to 0
$$
并对 $p$ 作{\it 降序}归纳，得到 $\Omega^p(k)$
在正次数上的上同调均为零；这对所有 $p$ 成立。
\end{proof}

\begin{lemma}
\label{derham-lemma-hodge-cohomology-projective-space}
有 $H^q(\mathbf{P}^n_A, \Omega^p) = 0$，除非
$0 \leq p = q \leq n$。对 $0 \leq p \leq n$，$A$-模
$H^p(\mathbf{P}^n_A, \Omega^p)$ 是秩为 $1$ 的自由模，其基元素为
$c_1^{Hodge}(\mathcal{O}(1))^p$。
\end{lemma}

\begin{proof}
消失性与自由性由引理
\ref{derham-lemma-twisted-hodge-cohomology-projective-space} 给出。
当 $p = 0$ 时，$1 \in H^0(\mathbf{P}^n_A, \mathcal{O})$
显然是生成元。

\medskip\noindent
先证明 $p = 1$ 的情形。考虑短正合列
$$
0 \to \Omega \to \mathcal{O}(-1)^{\oplus n + 1} \to \mathcal{O} \to 0
$$
它来自引理 \ref{derham-lemma-euler-sequence}。在引理
\ref{derham-lemma-twisted-hodge-cohomology-projective-space} 的证明中已经看到，
$H^1(\mathbf{P}^n_A, \Omega)$ 的生成元是边缘 $\xi$；
它来自 $1 \in H^0(\mathbf{P}^n_A, \mathcal{O})$。
如引理 \ref{derham-lemma-euler-sequence} 的证明那样，把
$\mathcal{O}(-1)^{\oplus n + 1}$ 与
$\bigoplus_{j = 0, \ldots, n} \mathcal{O}(-1)\text{d}T_j$ 等同。
考虑开覆盖
$$
\mathcal{U} : 
\mathbf{P}^n_A =
\bigcup\nolimits_{i = 0, \ldots, n} D_{+}(T_i)
$$
整体截面 $1$（它属于 $\mathcal{O}$）在
$U_i = D_+(T_i)$ 上的限制，可以提升为截面
$T_i^{-1} \text{d}T_i$；它是
$\bigoplus \mathcal{O}(-1)\text{d}T_j$ 在 $U_i$ 上的截面。
因此，表示 $\xi$ 的上闭链由下式给出：
$$
T_{i_1}^{-1} \text{d}T_{i_1} - T_{i_0}^{-1} \text{d}T_{i_0} =
\text{d}\log(T_{i_1}/T_{i_0}) \in \Omega(U_{i_0i_1})
$$
另一方面，对每个 $i$，截面 $T_i$ 是 $\mathcal{O}(1)$
在 $U_i$ 上的平凡化截面。因此可见
$f_{i_0i_1} = T_{i_1}/T_{i_0} \in \mathcal{O}^*(U_{i_0i_1})$
是表示 $\mathcal{O}(1)$ 的上闭链；它位于 $\Pic(\mathbf{P}^n_A)$ 中。
参见第 \ref{derham-section-first-chern-class} 节。因此
$c_1^{Hodge}(\mathcal{O}(1))$ 由上闭链
$\text{d}\log(T_{i_1}/T_{i_0})$ 给出，这与上面对 $\xi$ 的结果一致。

\medskip\noindent
对一般的 $p$ 作归纳证明。基本情形 $p = 0, 1$ 已在上面处理。
假设 $p > 1$。在引理
\ref{derham-lemma-twisted-hodge-cohomology-projective-space} 的证明中已经看到，
$H^p(\mathbf{P}^n_A, \Omega^p)$ 的生成元，是
$c_1^{Hodge}(\mathcal{O}(1))^{p - 1}$ 在与下列短正合列相伴的
上同调长正合列中的边缘：
$$
0 \to \Omega^p \to \wedge^p\left(\mathcal{O}(-1)^{\oplus n + 1}\right)
\to \Omega^{p - 1} \to 0
$$
由第 \ref{derham-section-first-chern-class} 节中的计算，上同调类
$c_1^{Hodge}(\mathcal{O}(1))^{p - 1}$ 在相差一个符号的意义下，
由各项如下的上闭链表示：
$$
\beta_{i_0 \ldots i_{p - 1}} =
\text{d}\log(T_{i_1}/T_{i_0}) \wedge
\text{d}\log(T_{i_2}/T_{i_1}) \wedge \ldots \wedge
\text{d}\log(T_{i_{p - 1}}/T_{i_{p - 2}})
$$
它位于 $\Omega^{p - 1}(U_{i_0 \ldots i_{p - 1}})$ 中。各个
$\beta_{i_0 \ldots i_{p - 1}}$ 可以提升为截面
$\tilde \beta_{i_0 \ldots i_{p -1}} =
T_{i_0}^{-1}\text{d}T_{i_0} \wedge \beta_{i_0 \ldots i_{p - 1}}$
；这是 $\wedge^p(\bigoplus \mathcal{O}(-1) \text{d}T_j)$
在 $U_{i_0 \ldots i_{p - 1}}$ 上的截面。由此可知，
$H^p(\mathbf{P}^n_A, \Omega^p)$ 的生成元由下列上闭链给出，
其分量为
\begin{align*}
\sum\nolimits_{a = 0}^p (-1)^a
\tilde \beta_{i_0 \ldots \hat{i_a} \ldots i_p}
& =
T_{i_1}^{-1}\text{d}T_{i_1} \wedge \beta_{i_1 \ldots i_p}
+ \sum\nolimits_{a = 1}^p (-1)^a
T_{i_0}^{-1}\text{d}T_{i_0} \wedge
\beta_{i_0 \ldots \hat{i_a} \ldots i_p} \\
& =
(T_{i_1}^{-1}\text{d}T_{i_1} - T_{i_0}^{-1}\text{d}T_{i_0}) \wedge
\beta_{i_1 \ldots i_p} +
T_{i_0}^{-1}\text{d}T_{i_0} \wedge \text{d}(\beta)_{i_0 \ldots i_p} \\
& =
\text{d}\log(T_{i_1}/T_{i_0}) \wedge \beta_{i_1 \ldots i_p}
\end{align*}
把它视为 $\Omega^p$ 在 $U_{i_0 \ldots i_p}$ 上的截面。
除去符号，它与表示 $c_1^{Hodge}(\mathcal{O}(1))^p$ 的上闭链相同。
证明完成。
\end{proof}

\begin{lemma}
\label{derham-lemma-de-rham-cohomology-projective-space}
对 $0 \leq i \leq n$，德拉姆上同调
$H^{2i}_{dR}(\mathbf{P}^n_A/A)$ 是自由 $A$-模，秩为 $1$，
其基元素为 $c_1^{dR}(\mathcal{O}(1))^i$。
$\mathbf{P}^n_A$ 在 $A$ 上的德拉姆上同调在所有其他次数均为零。
\end{lemma}

\begin{proof}
考虑第 \ref{derham-section-hodge-to-de-rham} 节中的霍奇--德拉姆谱序列。
由引理 \ref{derham-lemma-hodge-cohomology-projective-space} 中对
$\mathbf{P}^n_A$ 在 $A$ 上的霍奇上同调所作的计算可知，
该谱序列在 $E_1$ 叶退化。因此
$H^{2i}_{dR}(\mathbf{P}^n_A/A)$ 是自由 $A$-模，秩为 $1$；
这对 $0 \leq i \leq n$ 成立，而其他次数为零。
注意 $c_1^{dR}(\mathcal{O}(1))^i \in H^{2i}_{dR}(\mathbf{P}^n_A/A)$
（$i = 0, \ldots, n$），且当 $i = n$ 时，该元素是
$c_1^{Hodge}(\mathcal{L})^n$ 在下列复形映射下的像：
$$
\Omega^n_{\mathbf{P}^n_A/A}[-n]
\longrightarrow
\Omega^\bullet_{\mathbf{P}^n_A/A}
$$
例如，这可由备注 \ref{derham-remark-truncations} 中的讨论得到，也可由第
\ref{derham-section-first-chern-class} 节中表示这些类的上闭链之显式描述得到。
谱序列表明，诱导映射
$$
H^n(\mathbf{P}^n_A, \Omega^n_{\mathbf{P}^n_A/A}) \longrightarrow
H^{2n}_{dR}(\mathbf{P}^n_A/A)
$$
是同构。由于 $c_1^{Hodge}(\mathcal{L})^n$ 是源的生成元
（引理 \ref{derham-lemma-hodge-cohomology-projective-space}），
可知 $c_1^{dR}(\mathcal{L})^n$ 是靶的生成元。
由杯积的 $A$-双线性性，进一步得到
$c_1^{dR}(\mathcal{L})^i$ 是
$H^{2i}_{dR}(\mathbf{P}^n_A/A)$ 的生成元；这对
$0 \leq i \leq n$ 成立。
\end{proof}









\section{光滑态射的谱序列}
\label{derham-section-relative-spectral-sequence}

\noindent
考虑概形的交换图
$$
\xymatrix{
X \ar[rr]_f \ar[rd]_p & & Y \ar[ld]^q \\
& S
}
$$
其中 $f$ 是光滑态射。于是得到局部分裂的短正合列
$$
0 \to f^*\Omega_{Y/S} \to \Omega_{X/S} \to \Omega_{X/Y} \to 0
$$
，这是《态射》，引理 \ref{morphisms-lemma-triangle-differentials-smooth} 的结论。
把它看作递降滤过 $F$；该滤过定义在 $\Omega_{X/S}$ 上，其中
$F^0\Omega_{X/S} = \Omega_{X/S}$、$F^1\Omega_{X/S} = f^*\Omega_{Y/S}$，且
$F^2\Omega_{X/S} = 0$。应用函子 $\wedge^p$，对每个 $p$ 得到诱导滤过
$$
\Omega^p_{X/S} = F^0\Omega^p_{X/S} \supset
F^1\Omega^p_{X/S} \supset
F^2\Omega^p_{X/S} \supset \ldots \supset F^{p + 1}\Omega^p_{X/S} = 0
$$
其相继商为
$$
\text{gr}^k\Omega^p_{X/S} =
F^k\Omega^p_{X/S}/F^{k + 1}\Omega^p_{X/S} =
f^*\Omega^k_{Y/S} \otimes_{\mathcal{O}_X} \Omega^{p - k}_{X/Y} =
f^{-1}\Omega^k_{Y/S} \otimes_{f^{-1}\mathcal{O}_Y} \Omega^{p - k}_{X/Y}
$$
这里 $k = 0, \ldots, p$。事实上，读者可用莱布尼茨法则验证：
$F^k\Omega^\bullet_{X/S}$ 是 $\Omega^\bullet_{X/S}$ 的子复形。
由此，$\Omega^\bullet_{X/S}$ 具有滤过复形的结构。也可以通过观察下式看出：
$$
F^k\Omega^\bullet_{X/S} =
\Im\left(\wedge :
\text{Tot}(
f^{-1}\sigma_{\geq k}\Omega^\bullet_{Y/S} \otimes_{p^{-1}\mathcal{O}_S}
\Omega^\bullet_{X/S})
\longrightarrow
\Omega^\bullet_{X/S}\right)
$$
它是 $X$ 上复形映射的像。滤过复形
$$
\Omega^\bullet_{X/S} = F^0\Omega^\bullet_{X/S} \supset
F^1\Omega^\bullet_{X/S} \supset F^2\Omega^\bullet_{X/S} \supset \ldots
$$
具有如下相伴分次部分：
$$
\text{gr}^k\Omega^\bullet_{X/S} = 
f^{-1}\Omega^k_{Y/S}[-k] \otimes_{f^{-1}\mathcal{O}_Y} \Omega^\bullet_{X/Y}
$$
这由上述讨论得到。

\begin{lemma}
\label{derham-lemma-spectral-sequence-smooth}
设 $f : X \to Y$ 是基概形 $S$ 上拟紧、拟分离且光滑的概形态射。
存在有界谱序列，其第一叶为
$$
E_1^{p, q} =
H^q(\Omega^p_{Y/S} \otimes_{\mathcal{O}_Y}^\mathbf{L} Rf_*\Omega^\bullet_{X/Y})
$$
它收敛到 $R^{p + q}f_*\Omega^\bullet_{X/S}$。
\end{lemma}

\begin{proof}
把 $\Omega^\bullet_{X/S}$ 配以上述滤过，视为滤过复形。
该谱序列就是《上同调》，引理
\ref{cohomology-lemma-relative-spectral-sequence-filtered-object} 的谱序列。
由《概形的导出范畴》，引理
\ref{perfect-lemma-cohomology-de-rham-base-change}，有
$$
Rf_*\text{gr}^k\Omega^\bullet_{X/S} =
\Omega^k_{Y/S}[-k] \otimes_{\mathcal{O}_Y}^\mathbf{L} Rf_*\Omega^\bullet_{X/Y}
$$
从而得到结论。
\end{proof}

\begin{remark}
\label{derham-remark-gauss-manin}
在引理 \ref{derham-lemma-spectral-sequence-smooth} 中，考虑上同调层
$$
\mathcal{H}^q_{dR}(X/Y) = H^q(Rf_*\Omega^\bullet_{X/Y})
$$
若 $f$ 除光滑外还为真，且 $S$ 是 $\mathbf{Q}$ 上的概形，则
$\mathcal{H}^q_{dR}(X/Y)$ 有限局部自由（在此插入将来的参考文献）。
若只假设 $\mathcal{H}^q_{dR}(X/Y)$ 是平坦
$\mathcal{O}_Y$-模，则得到（略去一个很短的论证）
$$
E_1^{p, q} =
\Omega^p_{Y/S} \otimes_{\mathcal{O}_Y} \mathcal{H}^q_{dR}(X/Y)
$$
而谱序列中的微分为映射
$$
d_1^{p, q} :
\Omega^p_{Y/S} \otimes_{\mathcal{O}_Y} \mathcal{H}^q_{dR}(X/Y)
\longrightarrow
\Omega^{p + 1}_{Y/S} \otimes_{\mathcal{O}_Y} \mathcal{H}^q_{dR}(X/Y)
$$
特别地，当 $p = 0$ 时，得到映射
$d_1^{0, q} : \mathcal{H}^q_{dR}(X/Y) \to
\Omega^1_{Y/S} \otimes_{\mathcal{O}_Y} \mathcal{H}^q_{dR}(X/Y)$
；它实际上是可积联络 $\nabla$（在此插入将来的参考文献）。
复形
$$
\mathcal{H}^q_{dR}(X/Y) \to
\Omega^1_{Y/S} \otimes_{\mathcal{O}_Y} \mathcal{H}^q_{dR}(X/Y) \to
\Omega^2_{Y/S} \otimes_{\mathcal{O}_Y} \mathcal{H}^q_{dR}(X/Y) \to \ldots
$$
的微分由 $d_1^{\bullet, q}$ 给出，它是 $\nabla$ 的德拉姆复形。
联络 $\nabla$ 称为{\it 高斯--曼宁联络}。
\end{remark}






\section{Leray--Hirsch 型定理}
\label{derham-section-leray-hirsch}

\noindent
本节证明：对光滑真态射，有时可以用下方的德拉姆上同调表示上方的德拉姆上同调。

\begin{lemma}
\label{derham-lemma-relative-global-generation-on-fibres}
设 $f : X \to Y$ 是光滑真概形态射。设 $N$ 和
$n_1, \ldots, n_N \geq 0$ 为整数，并设
$\xi_i \in H^{n_i}_{dR}(X/Y)$（$1 \leq i \leq N$）。
假设对所有点 $y \in Y$，$\xi_1, \ldots, \xi_N$ 在
$H^*_{dR}(X_y/y)$ 中的像构成 $\kappa(y)$ 上的基。则与
$$
\bigoplus\nolimits_{i = 1}^N \mathcal{O}_Y[-n_i]
\longrightarrow
Rf_*\Omega^\bullet_{X/Y}
$$
$\xi_1, \ldots, \xi_N$ 相伴的上述映射是同构。
\end{lemma}

\begin{proof}
由引理 \ref{derham-lemma-proper-smooth-de-Rham}，
$Rf_*\Omega^\bullet_{X/Y}$ 是 $D(\mathcal{O}_Y)$ 中的完美对象，
且其构造与任意基变换相容。因此，本引理的映射是
完美对象之间的映射 $a : K \to L$；这些对象位于
$D(\mathcal{O}_Y)$ 中。根据对纤维的假设，它到任一点的导出限制都是同构。
于是锥 $C$（即 $a$ 的锥）是 $D(\mathcal{O}_Y)$ 中的完美对象
（《上同调》，引理 \ref{cohomology-lemma-two-out-of-three-perfect}），
且其到任一点的导出限制为零。因此，由《更多代数》，引理
\ref{more-algebra-lemma-lift-perfect-from-residue-field}，$C$ 为零，
而 $a$ 是同构。（这里还使用《概形的导出范畴》中的引理
\ref{perfect-lemma-affine-compare-bounded} 和
\ref{perfect-lemma-perfect-affine}，以转化到代数情形。）
\end{proof}

\noindent
我们先在如下特殊情形下证明本节的主要结果。

\begin{lemma}
\label{derham-lemma-global-generation-on-fibres}
设 $f : X \to Y$ 是基概形 $S$ 上的光滑真概形态射。假设
\begin{enumerate}
\item $Y$ 和 $S$ 都是仿射的；
\item 存在整数 $N$、$n_1, \ldots, n_N \geq 0$，以及
$\xi_i \in H^{n_i}_{dR}(X/S)$（$1 \leq i \leq N$），使得对所有点
$y \in Y$，$\xi_1, \ldots, \xi_N$ 在 $H^*_{dR}(X_y/y)$ 中的像构成
$\kappa(y)$ 上的基。
\end{enumerate}
则映射
$$
\bigoplus\nolimits_{i = 1}^N H^*_{dR}(Y/S) \longrightarrow
H^*_{dR}(X/S), \quad
(a_1, \ldots, a_N) \longmapsto  \sum \xi_i \cup f^*a_i
$$
是同构。
\end{lemma}

\begin{proof}
记 $Y = \Spec(A)$ 且 $S = \Spec(R)$。在此情形下，由引理
\ref{derham-lemma-de-rham-affine}，$\Omega^\bullet_{A/R}$ 计算
$R\Gamma(Y, \Omega^\bullet_{Y/S})$。取有限仿射开覆盖
$\mathcal{U} : X = \bigcup_{i \in I} U_i$。考虑复形
$$
K^\bullet =
\text{Tot}(\check{\mathcal{C}}^\bullet(\mathcal{U}, \Omega_{X/S}^\bullet))
$$
，如《上同调》第
\ref{cohomology-section-cech-cohomology-of-complexes} 节所述。
我们汇集关于这个复形的一些事实；其中大部分可在刚才给出的参考文献中找到：
\begin{enumerate}
\item $K^\bullet$ 是一个 $R$-模复形，其各项都是 $A$-模；
\item $K^\bullet$ 表示 $R\Gamma(X, \Omega^\bullet_{X/S})$ 这个 $D(R)$ 中的对象
（《概形上同调》，引理
\ref{coherent-lemma-quasi-coherent-affine-cohomology-zero}，以及
《上同调》，引理 \ref{cohomology-lemma-cech-complex-complex-computes}）；
\item 存在自然映射 $\Omega^\bullet_{A/R} \to K^\bullet$，
它是 $R$-模复形的映射，并且逐项是 $A$-线性的，
并在上同调上诱导拉回映射
$H^*_{dR}(Y/S) \to H^*_{dR}(X/S)$；
\item $K^\bullet$ 上有一个记作 $\wedge$ 的乘法，使其成为微分分次
$R$-代数；
\item $K^\bullet$ 上的乘法诱导 $H^*_{dR}(X/S)$ 上的杯积
（《上同调》，第 \ref{cohomology-section-cup-product} 节）；
\item 滤过 $F$ 在 $\Omega^*_{X/S}$ 上，并诱导如下滤过
$$
K^\bullet =
F^0K^\bullet \supset F^1K^\bullet \supset F^2K^\bullet \supset \ldots
$$
；它由 $K^\bullet$ 的子复形构成，并满足
\begin{enumerate}
\item $F^kK^n \subset K^n$ 是 $A$-子模；
\item $F^kK^\bullet \wedge F^lK^\bullet \subset F^{k + l}K^\bullet$；
\item $\text{gr}^kK^\bullet$ 是 $A$-模复形；
\item $\text{gr}^0K^\bullet =
\text{Tot}(\check{\mathcal{C}}^\bullet(\mathcal{U}, \Omega_{X/Y}^\bullet))$
，且表示 $R\Gamma(X, \Omega^\bullet_{X/Y})$ 这个 $D(A)$ 中的对象；
\item 乘法诱导同构
$\Omega^k_{A/R}[-k] \otimes_A \text{gr}^0K^\bullet \to \text{gr}^kK^\bullet$
\end{enumerate}
\end{enumerate}
我们略去这些陈述的详细证明；请参阅引理
\ref{derham-lemma-spectral-sequence-smooth} 中构造谱序列之前的讨论。

\medskip\noindent
对每个 $i = 1, \ldots, N$，选取上闭链 $x_i \in K^{n_i}$
来表示 $\xi_i$。接着考虑复形的映射
$$
\tilde x :
M^\bullet = \bigoplus\nolimits_{i = 1, \ldots, N}
\Omega^\bullet_{A/R}[-n_i]
\longrightarrow
K^\bullet
$$
，它把 $\omega$（位于第 $i$ 个直和项）映到 $x_i \wedge \omega$。
只需证明这个映射是拟同构。我们按如下规则赋予 $M^\bullet$ 滤过复形结构：
$$
F^kM^\bullet =
\bigoplus\nolimits_{i = 1, \ldots, N}
(\sigma_{\geq k}\Omega^\bullet_{A/R})[-n_i]
$$
。在此选择下，映射 $\tilde x$ 是滤过复形的态射。注意
$\text{gr}^0M^\bullet = \bigoplus A[-n_i]$，且乘法诱导同构
$\Omega^k_{A/R}[-k] \otimes_A \text{gr}^0M^\bullet \to \text{gr}^kM^\bullet$.
由构造及引理 \ref{derham-lemma-relative-global-generation-on-fibres} 可知
$$
\text{gr}^0\tilde x :
\text{gr}^0M^\bullet \longrightarrow
\text{gr}^0K^\bullet
$$
在 $D(A)$ 中是同构。因而对所有 $k \geq 0$，得到同构
$$
\text{gr}^k \tilde x :
\text{gr}^kM^\bullet = \Omega^k_{A/R}[-k] \otimes_A \text{gr}^0M^\bullet
\longrightarrow
\Omega^k_{A/R}[-k] \otimes_A \text{gr}^0K^\bullet =
\text{gr}^kK^\bullet
$$
于 $D(A)$ 中。事实上，复形
$\text{gr}^0K^\bullet =
\text{Tot}(\check{\mathcal{C}}^\bullet(\mathcal{U}, \Omega_{X/Y}^\bullet))$
由《概形的导出范畴》，引理 \ref{perfect-lemma-K-flat}，作为
$A$-模复形是 K-平坦的。因此右边的张量积就是导出张量积；
左边直接考察即可得到同一结论。最后，取导出张量积
$\Omega^k_{A/R}[-k] \otimes_A^\mathbf{L} -$ 是 $D(A)$ 上的函子，
故将同构映到同构。对 $k$ 作归纳，推出
$$
\tilde x : M^\bullet/F^kM^\bullet \to K^\bullet/F^kK^\bullet
$$
在 $D(R)$ 中是同构，因为有短正合列
$$
0 \to F^kM^\bullet/F^{k + 1}M^\bullet \to
M^\bullet/F^{k + 1}M^\bullet \to
\text{gr}^kM^\bullet \to 0
$$
，对 $K^\bullet$ 也有类似的短正合列。由于这些滤过在任一给定次数上
都是有限的，这就证明了 $\tilde x$ 是拟同构。
\end{proof}

\begin{proposition}
\label{derham-proposition-global-generation-on-fibres}
设 $f : X \to Y$ 是基概形 $S$ 上的光滑真概形态射。设 $N$ 和
$n_1, \ldots, n_N \geq 0$ 为整数，并设
$\xi_i \in H^{n_i}_{dR}(X/S)$（$1 \leq i \leq N$）。假设对所有点
$y \in Y$，$\xi_1, \ldots, \xi_N$ 在 $H^*_{dR}(X_y/y)$ 中的像构成
$\kappa(y)$ 上的基。则映射
$$
\tilde \xi = \bigoplus \tilde \xi_i[-n_i] :
\bigoplus \Omega^\bullet_{Y/S}[-n_i]
\longrightarrow
Rf_*\Omega^\bullet_{X/S}
$$
（见证明）在 $D(Y, (Y \to S)^{-1}\mathcal{O}_S)$ 中是同构；相应地，映射
$$
\bigoplus\nolimits_{i = 1}^N H^*_{dR}(Y/S) \longrightarrow
H^*_{dR}(X/S), \quad
(a_1, \ldots, a_N) \longmapsto  \sum \xi_i \cup f^*a_i
$$
是同构。
\end{proposition}

\begin{proof}
记结构态射为 $p : X \to S$ 和 $q : Y \to S$。令
$\xi'_i : \Omega^\bullet_{X/S} \to \Omega^\bullet_{X/S}[n_i]$
为评注 \ref{derham-remark-cup-product-as-a-map} 中与 $\xi_i$ 对应的映射。记
$$
\tilde \xi_i :
\Omega^\bullet_{Y/S} \to Rf_*\Omega^\bullet_{X/S}[n_i]
$$
为 $\xi'_i$ 与典范映射
$\Omega^\bullet_{Y/S} \to Rf_*\Omega^\bullet_{X/S}$ 的复合。利用
$$
R\Gamma(Y, Rf_*\Omega^\bullet_{X/S}) = R\Gamma(X, \Omega^\bullet_{X/S})
$$
，在上同调上，$\tilde \xi_i$ 是映射 $\eta \mapsto \xi_i \cup f^*\eta$，
它从 $H^m_{dR}(Y/S)$ 映到 $H^{m + n}_{dR}(X/S)$。
此外，由于 $\xi'_i$ 的构造与到开集的限制可交换，$\tilde \xi_i$
的构造也与到开集的限制可交换。

\medskip\noindent
因此可以考虑映射
$$
\tilde \xi = \bigoplus \tilde \xi_i[-n_i] :
\bigoplus \Omega^\bullet_{Y/S}[-n_i]
\longrightarrow
Rf_*\Omega^\bullet_{X/S}
$$
。要证明该引理，只需证明它在 $D(Y, q^{-1}\mathcal{O}_S)$ 中是同构。
若能证明 $\tilde \xi$ 来自滤过复形的映射（配备适当滤过），
便可诉诸引理 \ref{derham-lemma-spectral-sequence-smooth} 的谱序列完成证明。
这比所需工作更多，因此我们改为约化到仿射情形（其证明在某种意义上
确实使用了谱序列）。

\medskip\noindent
事实上，若 $Y' \subset Y$ 是任意开集，其逆像为 $X' \subset X$，
则 $\tilde \xi|_{X'}$ 诱导映射
$$
\bigoplus\nolimits_{i = 1}^N H^*_{dR}(Y'/S) \longrightarrow
H^*_{dR}(X'/S), \quad
(a_1, \ldots, a_N) \longmapsto  \sum \xi_i|_{X'} \cup f^*a_i
$$
于 $Y'$ 上的上同调中，见上文讨论。因此，只需为 $Y$ 上的拓扑找一个基，使该命题对基中的
各成员成立（特别地，此时可以不再考虑映射 $\tilde \xi$）。这就约化到
$Y$ 和 $S$ 都是仿射的情形；该情形由引理
\ref{derham-lemma-global-generation-on-fibres} 处理，证明完成。
\end{proof}





\section{射影空间丛公式}
\label{derham-section-projective-space-bundle-formula}

\noindent
标题已说明一切。

\begin{proposition}
\label{derham-proposition-projective-space-bundle-formula}
设 $X \to S$ 是概形态射。设 $\mathcal{E}$ 是局部自由
$\mathcal{O}_X$-模，且其秩恒为 $r$。考虑态射
$p : P = \mathbf{P}(\mathcal{E}) \to X$。
则映射
$$
\bigoplus\nolimits_{i = 0, \ldots, r - 1} H^*_{dR}(X/S)
\longrightarrow
H^*_{dR}(P/S)
$$
由规则
$$
(a_0, \ldots, a_{r - 1}) \longmapsto
\sum\nolimits_{i = 0, \ldots, r - 1} c_1^{dR}(\mathcal{O}_P(1))^i \cup p^*(a_i)
$$
给出，并且是同构。
\end{proposition}

\begin{proof}
取仿射开集 $\Spec(A) \subset X$，使 $\mathcal{E}$ 限制为平凡局部自由模
$\mathcal{O}_{\Spec(A)}^{\oplus r}$。于是
$P \times_X \Spec(A) = \mathbf{P}^{r - 1}_A$。由此可见 $p$ 真且光滑，
见第 \ref{derham-section-projective-space} 节。此外，诸类
$c_1^{dR}(\mathcal{O}_P(1))^i$（$i = 0, 1, \ldots, r - 1$）限制到纤维
$X_y = \mathbf{P}^{r - 1}_y$ 后，自由生成德拉姆上同调
$H^*_{dR}(X_y/y)$ 这一 $\kappa(y)$-空间，见引理
\ref{derham-lemma-de-rham-cohomology-projective-space}。因此，我们已经验证了命题
\ref{derham-proposition-global-generation-on-fibres} 的条件，结论得证。
\end{proof}

\begin{remark}
\label{derham-remark-projective-space-bundle-formula}
在命题 \ref{derham-proposition-projective-space-bundle-formula} 的情形下，进一步得到映射
$$
\tilde \xi :
\bigoplus\nolimits_{t = 0, \ldots, r - 1}
\Omega^\bullet_{X/S}[-2t]
\longrightarrow
Rp_*\Omega^\bullet_{P/S}
$$
在 $D(X, (X \to S)^{-1}\mathcal{O}_X)$ 中是同构；这直接由应用命题
\ref{derham-proposition-global-generation-on-fibres} 得到。注意，对 $t = 0$，
相应箭头恰为典范映射
$c_{P/X} : \Omega^\bullet_{X/S} \to Rp_*\Omega^\bullet_{P/S}$
，见第 \ref{derham-section-de-rham-complex} 节。事实上，在这个特殊情形下，
还可进一步明确该映射。考虑典范映射
$$
\xi' : \Omega^\bullet_{P/S} \to \Omega^\bullet_{P/S}[2]
$$
，它是评注 \ref{derham-remark-cup-product-as-a-map} 中对应于
$c_1^{dR}(\mathcal{O}_P(1))$ 的映射。于是
$$
\xi'[2(t - 1)] \circ \ldots \circ \xi'[2] \circ \xi' : 
\Omega^\bullet_{P/S} \to \Omega^\bullet_{P/S}[2t]
$$
是评注 \ref{derham-remark-cup-product-as-a-map} 中对应于
$c_1^{dR}(\mathcal{O}_P(1))^t$ 的映射。追踪命题
\ref{derham-proposition-global-generation-on-fibres} 的证明中所作的选择，得到
$$
\tilde \xi|_{\Omega^\bullet_{X/S}[-2t]} =
Rp_*\xi'[-2] \circ \ldots \circ Rp_*\xi'[-2(t - 1)] \circ
Rp_*\xi'[-2t] \circ c_{P/X}[-2t]
$$
，这就是该同构到直和项 $\Omega^\bullet_{X/S}[-2t]$ 的限制。
由此得到下面将用到的简单推论：令
$$
M = \bigoplus\nolimits_{t = 1, \ldots, r - 1} \Omega^\bullet_{X/S}[-2t]
\quad\text{且}\quad
K = \bigoplus\nolimits_{t = 0, \ldots, r - 2} \Omega^\bullet_{X/S}[-2t]
$$
；将它们视为箭头 $\tilde \xi$ 的源中的子复形。由上文讨论形式地推出
$$
c_{P/X} \oplus
\tilde \xi|_M :
\Omega^\bullet_{X/S} \oplus M \longrightarrow
Rp_*\Omega^\bullet_{P/S}
$$
是同构，且图表
$$
\xymatrix{
K \ar[d]_{\tilde \xi|_K} \ar[r]_{\text{id}} &
M[2] \ar[d]^{(\tilde \xi|_M)[2]} \\
Rp_*\Omega^\bullet_{P/S} \ar[r]^{Rp_*\xi'} &
Rp_*\Omega^\bullet_{P/S}[2]
}
$$
交换；其中 $\text{id} : K \to M[2]$ 把对应于 $t$ 的直和项
（在 $K$ 的分解中）与对应于 $t + 1$ 的直和项（在 $M$ 的分解中）等同。
\end{remark}









\section{沿除子的对数极点}
\label{derham-section-divisor}

\noindent
设 $X \to S$ 是概形态射。设 $Y \subset X$ 是有效 Cartier 除子。
若 $X$ 沿 $Y$ 在平展局部看起来像 $Y \times \mathbf{A}^1$，
则存在复形的典范短正合列
$$
0 \to \Omega^\bullet_{X/S} \to
\Omega^\bullet_{X/S}(\log Y) \to
\Omega^\bullet_{Y/S}[-1] \to 0
$$
，它具有许多良好性质，本节将加以讨论。这个构造还有一个从正规交叉除子
出发的变体（《平展态射》，定义
\ref{etale-definition-strict-normal-crossings}），我们将在别处讨论
（此处插入未来的参考文献）。

\begin{definition}
\label{derham-definition-local-product}
设 $X \to S$ 是概形态射。设 $Y \subset X$ 是有效 Cartier 除子。
称{\it 对 $Y \subset X$（在 $S$ 上）定义了对数极点德拉姆复形}，
如果对所有 $y \in Y$ 以及局部方程 $f \in \mathcal{O}_{X, y}$
（它定义 $Y$），都有
\begin{enumerate}
\item $\mathcal{O}_{X, y} \to \Omega_{X/S, y}$，$g \mapsto g \text{d}f$
是分裂单射；
\item $\Omega^p_{X/S, y}$ 无 $f$-挠，对所有 $p$ 均如此。
\end{enumerate}
\end{definition}

\noindent
一个简单的局部计算表明，对每个 $y \in Y$，只需找到一个满足条件
(1) 和 (2) 的局部方程 $f$。

\begin{lemma}
\label{derham-lemma-log-complex}
设 $X \to S$ 是概形态射。设 $Y \subset X$ 是有效 Cartier 除子。
假设对 $Y \subset X$（在 $S$ 上）定义了对数极点德拉姆复形。
则存在复形的典范短正合列
$$
0 \to \Omega^\bullet_{X/S} \to
\Omega^\bullet_{X/S}(\log Y) \to
\Omega^\bullet_{Y/S}[-1] \to 0
$$
\end{lemma}

\begin{proof}
我们的假设是：对每个 $y \in Y$ 以及局部方程
$f \in \mathcal{O}_{X, y}$（它定义 $Y$），都有
$$
\Omega_{X/S, y} = \mathcal{O}_{X, y}\text{d}f \oplus M
\quad\text{且}\quad
\Omega^p_{X/S, y} = \wedge^{p - 1}(M)\text{d}f \oplus \wedge^p(M)
$$
，其中某个模为 $M$，并且对 $f$ 而言，$\wedge^p(M)$ 无挠。因而
$$
\Omega^p_{Y/S, y} = \wedge^p(M/fM) = \wedge^p(M)/f\wedge^p(M)
$$
。下文将默认使用这些事实。特别地，层 $\Omega^p_{X/S}$ 没有支撑在
$Y$ 上的非零局部截面，并且有典范包含
$$
\Omega^p_{X/S} \subset \Omega^p_{X/S}(Y)
$$
，见《平坦性的进一步结果》，第 \ref{flat-section-eta} 节。令
$U = \Spec(A)$ 为仿射开子概形，使得 $Y \cap U = V(f)$，其中
$f \in A$ 是非零因子。考虑 $\mathcal{O}_U$-子模，它位于
$\Omega^p_{X/S}(Y)|_U$ 中，
它由 $\Omega^p_{X/S}|_U$ 和
$\text{d}\log(f) \wedge \Omega^{p - 1}_{X/S}$ 生成，其中
$\text{d}\log(f) = f^{-1}\text{d}(f)$。这与 $f$ 的选择无关：
$Y$ 在 $U$ 上的理想的另一个生成元等于 $uf$，其中 $u \in A$ 是单位，且
$$
\text{d}\log(uf) - \text{d}\log(f) = \text{d}\log(u) = u^{-1}\text{d}u
$$
是 $\Omega_{X/S}$ 在 $U$ 上的截面。这些局部层可粘合成拟凝聚子模
$$
\Omega^p_{X/S} \subset \Omega^p_{X/S}(\log Y) \subset \Omega^p_{X/S}(Y)
$$
。约定把 $\Omega^p_{Y/S}$ 看作拟凝聚 $\mathcal{O}_X$-模。
存在唯一的满射 $\mathcal{O}_X$-线性映射
$$
\text{Res} : \Omega^p_{X/S}(\log Y) \to \Omega^{p - 1}_{Y/S}
$$
，其定义规则为
$$
\text{Res}(\eta' + \text{d}\log(f) \wedge \eta) = \eta|_{Y \cap U}
$$
；这里 $U$ 遍历如上的开集，且
$\eta' \in \Omega^p_{X/S}(U)$ 和 $\eta \in \Omega^{p - 1}_{X/S}(U)$ 任取。
若微分形式 $\eta$ 在 $U$ 上，并且限制到 $Y \cap U$ 后为零，则有
$\eta = \text{d}f \wedge \eta' + f\eta''$，其中 $\eta'$ 和 $\eta''$
是 $U$ 上的某些微分形式。因此得到短正合列
$$
0 \to \Omega^p_{X/S} \to \Omega^p_{X/S}(\log Y) \to \Omega^{p - 1}_{Y/S} \to 0
$$
，对所有 $p$ 均成立。还需定义微分
$\Omega^p_{X/S}(\log Y) \to \Omega^{p + 1}_{X/S}(\log Y)$。
在子层 $\Omega^p_{X/S}$ 上，使用 $X$ 在 $S$ 上的德拉姆复形的微分。
最后定义
$\text{d}(\text{d}\log(f) \wedge \eta) = -\text{d}\log(f) \wedge \text{d}\eta$.
这个符号由 Leibniz 法则（在 $\Omega^\bullet_{X/S}$ 上）所迫使，
并与 $\Omega^\bullet_{Y/S}[-1]$ 上的微分相容；按照《同调》，定义
\ref{homology-definition-shift-cochain} 中的符号约定，后者正是
$-\text{d}_{Y/S}$。这样便得到引理所述的复形短正合列。
\end{proof}

\begin{definition}
\label{derham-definition-log-complex}
设 $X \to S$ 是概形态射。设 $Y \subset X$ 是有效 Cartier 除子。
假设对 $Y \subset X$（在 $S$ 上）定义了对数极点德拉姆复形。则复形
$$
\Omega^\bullet_{X/S}(\log Y)
$$
（构造见引理 \ref{derham-lemma-log-complex}）称为
{\it $Y \subset X$（在 $S$ 上）的对数极点德拉姆复形}。
\end{definition}

\noindent
这个复形具有许多良好性质。

\begin{lemma}
\label{derham-lemma-multiplication-log}
设 $p : X \to S$ 是概形态射。设 $Y \subset X$ 是有效 Cartier 除子。
假设对 $Y \subset X$（在 $S$ 上）定义了对数极点德拉姆复形。
\begin{enumerate}
\item 映射
$\wedge : \Omega^p_{X/S} \times \Omega^q_{X/S} \to \Omega^{p + q}_{X/S}$
唯一延拓为 $\mathcal{O}_X$-双线性映射
$$
\wedge : \Omega^p_{X/S}(\log Y) \times \Omega^q_{X/S}(\log Y)
\to \Omega^{p + q}_{X/S}(\log Y)
$$
，并满足 Leibniz 法则
$
\text{d}(\omega \wedge \eta) = \text{d}(\omega) \wedge \eta +
(-1)^{\deg(\omega)} \omega \wedge \text{d}(\eta)$,
\item 对于 (1) 中的乘法，映射
$\Omega^\bullet_{X/S} \to \Omega^\bullet_{X/S}(\log(Y)$
是微分分次 $\mathcal{O}_S$-代数的同态；
\item 经由 (1) 中的映射，有 $\Omega^p_{X/S}(\log Y) =
\wedge^p(\Omega^1_{X/S}(\log Y))$；
\item 映射
$\text{Res} : \Omega^\bullet_{X/S}(\log Y) \to \Omega^\bullet_{Y/S}[-1]$
满足
$$
\text{Res}(\omega \wedge \eta) = \text{Res}(\omega) \wedge \eta|_Y
$$
，其中 $\omega$ 是 $\Omega^p_{X/S}(\log Y)$ 的局部截面，而 $\eta$
是 $\Omega^q_{X/S}$ 的局部截面。
\end{enumerate}
\end{lemma}

\begin{proof}
这由引理 \ref{derham-lemma-log-complex} 的证明中该复形的局部构造直接计算得到。
细节从略。
\end{proof}

\noindent
考虑交换图表
$$
\xymatrix{
X' \ar[r]_f \ar[d] & X \ar[d] \\
S' \ar[r] & S
}
$$
，其中各对象都是概形。设 $Y \subset X$ 是有效 Cartier 除子，且其拉回
$Y' = f^*Y$ 有定义（《除子》，定义
\ref{divisors-definition-pullback-effective-Cartier-divisor}）。
假设对 $Y \subset X$（在 $S$ 上）定义了对数极点德拉姆复形，并且
对 $Y' \subset X'$（在 $S'$ 上）定义了对数极点德拉姆复形。
在这种情况下，得到复形短正合列之间的映射
$$
\xymatrix{
0 \ar[r] &
f^{-1}\Omega^\bullet_{X/S} \ar[r] \ar[d] &
f^{-1}\Omega^\bullet_{X/S}(\log Y) \ar[r] \ar[d] &
f^{-1}\Omega^\bullet_{Y/S}[-1] \ar[r] \ar[d] &
0 \\
0 \ar[r] &
\Omega^\bullet_{X'/S'} \ar[r] &
\Omega^\bullet_{X'/S'}(\log Y') \ar[r] &
\Omega^\bullet_{Y'/S'}[-1] \ar[r] &
0
}
$$
。线性化后，对每个 $p$ 得到线性映射
$f^*\Omega^p_{X/S}(\log Y) \to \Omega^p_{X'/S'}(\log Y')$.

\begin{lemma}
\label{derham-lemma-gysin-via-log-complex}
设 $f : X \to S$ 是概形态射。设 $Y \subset X$ 是有效 Cartier 除子。
假设对 $Y \subset X$（在 $S$ 上）定义了对数极点德拉姆复形。记
$$
\delta : \Omega^\bullet_{Y/S} \to \Omega^\bullet_{X/S}[2]
$$
为 $D(X, f^{-1}\mathcal{O}_S)$ 中来自引理
\ref{derham-lemma-log-complex} 的短正合列的“边缘”映射。记
$$
\xi' : \Omega^\bullet_{X/S} \to \Omega^\bullet_{X/S}[2]
$$
为 $D(X, f^{-1}\mathcal{O}_S)$ 中评注
\ref{derham-remark-cup-product-as-a-map} 的映射，它对应于
$\xi = c_1^{dR}(\mathcal{O}_X(-Y))$。记
$$
\zeta' : \Omega^\bullet_{Y/S} \to \Omega^\bullet_{Y/S}[2]
$$
为 $D(Y, f|_Y^{-1}\mathcal{O}_S)$ 中评注
\ref{derham-remark-cup-product-as-a-map} 的映射，它对应于
$\zeta = c_1^{dR}(\mathcal{O}_X(-Y)|_Y)$。则图表
$$
\xymatrix{
\Omega^\bullet_{X/S} \ar[d]_{\xi'} \ar[r] &
\Omega^\bullet_{Y/S} \ar[d]^{\zeta'} \ar[ld]_\delta \\
\Omega^\bullet_{X/S}[2] \ar[r] &
\Omega^\bullet_{Y/S}[2]
}
$$
在 $D(X, f^{-1}\mathcal{O}_S)$ 中交换。
\end{lemma}

\begin{proof}
更准确地说，我们把 $\delta$ 定义为与如下移位短正合列相伴的边缘映射：
$$
0 \to \Omega^\bullet_{X/S}[1] \to
\Omega^\bullet_{X/S}(\log Y)[1] \to
\Omega^\bullet_{Y/S} \to 0
$$
。只需证明每个三角形都交换。置
$\mathcal{L} = \mathcal{O}_X(-Y)$。记 $\pi : L \to X$ 为满足
$\pi_*\mathcal{O}_L = \bigoplus_{n \geq 0} \mathcal{L}^{\otimes n}$ 的线丛。

\medskip\noindent
左上三角形的交换性。由引理
\ref{derham-lemma-the-complex-for-L-star-gives-chern-class}，映射 $\xi'$
是引理 \ref{derham-lemma-the-complex-for-L-star} 所给三角形的边缘映射。
由函子性，只需证明存在短正合列的态射
$$
\xymatrix{
0 \ar[r] &
\Omega^\bullet_{X/S}[1] \ar[r] \ar[d] &
\Omega^\bullet_{L^\star/S, 0}[1] \ar[r] \ar[d] &
\Omega^\bullet_{X/S} \ar[r] \ar[d] &
0 \\
0 \ar[r] &
\Omega^\bullet_{X/S}[1] \ar[r] &
\Omega^\bullet_{X/S}(\log Y)[1] \ar[r] &
\Omega^\bullet_{Y/S} \ar[r] &
0
}
$$
，其中左、右竖直箭头是显然的。可用如下规则定义中间的竖直箭头：
$$
\omega' + \text{d}\log(s) \wedge \omega \longmapsto
\omega' + \text{d}\log(f) \wedge \omega
$$
。其中 $\omega', \omega$ 是 $\Omega^\bullet_{X/S}$ 的局部截面，
$s$ 是 $\mathcal{L}$ 的局部生成元，而 $f \in \mathcal{O}_X(-Y)$
是相应的截面；它属于 $Y$ 在 $X$ 中的理想层。由于引理
\ref{derham-lemma-the-complex-for-L-star} 与 \ref{derham-lemma-log-complex}
中映射的构造完全吻合，此定义可行。

\medskip\noindent
右下三角形的交换性。记 $\overline{L}$ 为 $L$ 到 $Y$ 的限制。
由引理 \ref{derham-lemma-the-complex-for-L-star-gives-chern-class}，映射
$\zeta'$ 是引理 \ref{derham-lemma-the-complex-for-L-star} 中使用线丛
$\overline{L}$（定义在 $Y$ 上）所给三角形的边缘映射。由函子性，
只需证明存在短正合列的态射
$$
\xymatrix{
0 \ar[r] &
\Omega^\bullet_{X/S}[1] \ar[r] \ar[d] &
\Omega^\bullet_{X/S}(\log Y)[1] \ar[r] \ar[d] &
\Omega^\bullet_{Y/S} \ar[r] \ar[d] &
0 \\
0 \ar[r] &
\Omega^\bullet_{Y/S}[1] \ar[r] &
\Omega^\bullet_{\overline{L}^\star/S, 0}[1] \ar[r] &
\Omega^\bullet_{Y/S} \ar[r] &
0 \\
}
$$
，其中左、右竖直箭头是显然的。可用如下规则定义中间的竖直箭头：
$$
\omega' + \text{d}\log(f) \wedge \omega \longmapsto
\omega'|_Y + \text{d}\log(\overline{s}) \wedge \omega|_Y
$$
。其中 $\omega', \omega$ 是 $\Omega^\bullet_{X/S}$ 的局部截面，
$f$ 是 $\mathcal{O}_X(-Y)$ 的局部生成元，并被视为 $X$ 上的函数；
$\overline{s}$ 是 $f|_Y$，并被视为
$\mathcal{L}|_Y = \mathcal{O}_X(-Y)|_Y$ 的截面。由于引理
\ref{derham-lemma-the-complex-for-L-star} 与 \ref{derham-lemma-log-complex}
中映射的构造完全吻合，此定义可行。
\end{proof}

\begin{lemma}
\label{derham-lemma-log-complex-consequence}
设 $X \to S$ 是概形态射。设 $Y \subset X$ 是有效 Cartier 除子。
假设对 $Y \subset X$（在 $S$ 上）定义了对数极点德拉姆复形。
设 $b \in H^m_{dR}(X/S)$ 是德拉姆上同调类，且其到 $Y$ 的限制为零。
则 $c_1^{dR}(\mathcal{O}_X(Y)) \cup b = 0$ 于 $H^{m + 2}_{dR}(X/S)$ 中。
\end{lemma}

\begin{proof}
这直接由引理 \ref{derham-lemma-gysin-via-log-complex} 得到。事实上，
$$
c_1^{dR}(\mathcal{O}_X(Y)) \cup b =
-c_1^{dR}(\mathcal{O}_X(-Y)) \cup b = -\xi'(b) = -\delta(b|_Y) = 0
$$
，正如所需。第二个等号见评注 \ref{derham-remark-cup-product-as-a-map}。
\end{proof}

\begin{lemma}
\label{derham-lemma-check-log-smooth}
设 $X \to T \to S$ 是概形态射。设 $Y \subset X$ 是有效 Cartier 除子。
若 $X \to T$ 与 $Y \to T$ 都光滑，则对 $Y \subset X$（在 $S$ 上）
定义了对数极点德拉姆复形。
\end{lemma}

\begin{proof}
设 $y \in Y$ 是一点。由《态射的进一步结果》，引理
\ref{more-morphisms-lemma-etale-local-structure}，存在整数 $0 \geq m$
及交换图表
$$
\xymatrix{
Y \ar[d] &
V \ar[l] \ar[d] \ar[r] &
\mathbf{A}^m_T
\ar[d]^{(a_1, \ldots, a_m) \mapsto (a_1, \ldots, a_m, 0)} \\
X &
U \ar[l] \ar[r]^-\pi &
\mathbf{A}^{m + 1}_T
}
$$
，其中 $U \subset X$ 是开集，$V = Y \cap U$，$\pi$ 是平展的，
$V = \pi^{-1}(\mathbf{A}^m_T)$，且 $y \in V$。记
$z \in \mathbf{A}^m_T$ 为 $y$ 的像。于是
$$
\Omega^p_{X/S, y} = \Omega^p_{\mathbf{A}^{m + 1}_T/S, z}
\otimes_{\mathcal{O}_{\mathbf{A}^{m + 1}_T, z}} \mathcal{O}_{X, x}
$$
，由引理 \ref{derham-lemma-etale}。记 $x_1, \ldots, x_{m + 1}$ 为
$\mathbf{A}^{m + 1}_T$ 上的坐标函数。由于定义
\ref{derham-definition-local-product} 中的条件 (1) 和 (2) 不依赖于局部坐标的选择，
只需在 $f$ 为 $x_{m + 1}$ 经平坦局部环同态
$\mathcal{O}_{\mathbf{A}^{m + 1}_T, z} \to \mathcal{O}_{X, x}$
所得的像时检验条件 (1) 和 (2)。这样可见，只需对
$\mathbf{A}^m_T \subset \mathbf{A}^{m + 1}_T$ 及点 $z$ 检验条件 (1) 和 (2)。
为证明这一情形，可以假设 $S = \Spec(A)$ 和 $T = \Spec(B)$ 都是仿射的。
设 $A \to B$ 是态射 $T \to S$ 所对应的环同态，并置
$P = B[x_1, \ldots, x_{m + 1}]$，从而
$\mathbf{A}^{m + 1}_T = \Spec(P)$。有
$$
\Omega_{P/A} = \Omega_{B/A} \otimes_B P \oplus
\bigoplus\nolimits_{j = 1, \ldots, m} P \text{d}x_j \oplus
P \text{d}x_{m + 1}
$$
。因此映射 $P \to \Omega_{P/A}$，$g \mapsto g \text{d}x_{m + 1}$
是分裂单射，且 $x_{m + 1}$ 在 $\Omega^p_{P/A}$ 上是非零因子，
对所有 $p \geq 0$ 均如此。关于与 $z$ 对应的素理想局部化，证明即完成。
\end{proof}

\begin{remark}
\label{derham-remark-check-log-completion-1}
设 $S$ 是局部 Noether 概形。设 $X$ 在 $S$ 上局部有限型。
设 $Y \subset X$ 是有效 Cartier 除子。若映射
$$
\mathcal{O}_{X, y}^\wedge \longrightarrow \mathcal{O}_{Y, y}^\wedge
$$
对所有 $y \in Y$ 都有截面，则对 $Y \subset X$（在 $S$ 上）
定义了对数极点德拉姆复形。若将来需要此结果，我们会给出精确陈述，
并在此补上证明。
\end{remark}

\begin{remark}
\label{derham-remark-check-log-completion-2}
设 $S$ 是局部 Noether 概形。设 $X$ 在 $S$ 上局部有限型。
设 $Y \subset X$ 是有效 Cartier 除子。若对每个 $y \in Y$，
都能找到 $S$ 上概形的图表
$$
X \xleftarrow{\varphi} U \xrightarrow{\psi} V
$$
，其中 $\varphi$ 平展，且
$\psi|_{\varphi^{-1}(Y)} : \varphi^{-1}(Y) \to V$ 平展，
则对 $Y \subset X$（在 $S$ 上）定义了对数极点德拉姆复形。
一个特殊情形是序对 $(X, Y)$ 在平展局部看起来像
$(V \times \mathbf{A}^1, V \times \{0\})$。若将来需要此结果，
我们会给出精确陈述，并在此补上证明。
\end{remark}

















\section{计算}
\label{derham-section-calculations}

\noindent
本节计算一个标准吹起的若干霍奇上同调群和德拉姆上同调群。

\medskip\noindent
固定环 $R$，并置 $S = \Spec(R)$。固定整数 $0 \leq m$ 和 $1 \leq n$。
考虑闭浸入
$$
Z = \mathbf{A}^m_S \longrightarrow \mathbf{A}^{m + n}_S = X,\quad
(a_1, \ldots, a_m) \mapsto (a_1, \ldots, a_m, 0, \ldots 0).
$$
我们要考虑吹起 $L$：它是 $X$ 沿闭子概形 $Z$ 的吹起。写作
$$
X =
\mathbf{A}^{m + n}_S =
\Spec(A)
\quad\text{其中}\quad
A = R[x_1, \ldots, x_m, y_1, \ldots, y_n]
$$
。令 $A = R[x_1, \ldots, x_m, y_1, \ldots, y_n]$ 成为分次 $R$-代数，
其中规定 $\deg(x_i) = 0$ 且 $\deg(y_j) = 1$。对于这个分次，有
$$
P =
\text{Proj}(A) =
\mathbf{A}^m_S \times_S \mathbf{P}^{n - 1}_S =
Z \times_S \mathbf{P}^{n - 1}_S =
\mathbf{P}^{n - 1}_Z
$$
。注意，截出 $Z$（在 $X$ 中）的理想是 $A_+$。因此 $L$ 是 Rees 代数
$$
A \oplus A_+ \oplus (A_+)^2 \oplus \ldots =
\bigoplus\nolimits_{d \geq 0} A_{\geq d}
$$
的 Proj。因此 $L$ 是《态射的进一步结果》第
\ref{more-morphisms-section-proj-spec} 节中更一般地研究的现象的一个例子；
我们将不再说明而直接使用那里得到的观察。特别地，有交换图表
$$
\xymatrix{
P \ar[r]_0 \ar[d]_p &
L \ar[r]_-\pi \ar[d]^b &
P \ar[d]^p \\
Z \ar[r]^i &
X \ar[r] &
Z
}
$$
，使得 $\pi : L \to P$ 是 $P = Z \times_S \mathbf{P}^{n - 1}_S$
上的线丛，其零截面为 $0$；它的像 $E = 0(P) \subset L$
是吹起 $b$ 的例外除子。

\begin{lemma}
\label{derham-lemma-comparison}
对 $a \geq 0$，有
\begin{enumerate}
\item 映射 $\Omega^a_{X/S} \to b_*\Omega^a_{L/S}$ 是同构；
\item 映射 $\Omega^a_{Z/S} \to p_*\Omega^a_{P/S}$ 是同构；
\item 映射 $Rb_*\Omega^a_{L/S} \to i_*Rp_*\Omega^a_{P/S}$
在次数 $\geq 1$ 的上同调层上是同构。
\end{enumerate}
\end{lemma}

\begin{proof}
先证明 (2)。由于
$P = Z \times_S \mathbf{P}^{n - 1}_S$
，有
$$
\Omega^a_{P/S} = \bigoplus\nolimits_{a = r + s}
\text{pr}_1^*\Omega^r_{Z/S} \otimes
\text{pr}_2^*\Omega^s_{\mathbf{P}^{n - 1}_S/S}
$$
。回忆 $p = \text{pr}_1$，由投影公式（《上同调》，引理
\ref{cohomology-lemma-projection-formula}），得到
$$
p_*\Omega^a_{P/S} = \bigoplus\nolimits_{a = r + s}
\Omega^r_{Z/S} \otimes
\text{pr}_{1, *}\text{pr}_2^*\Omega^s_{\mathbf{P}^{n - 1}_S/S}
$$
。由第 \ref{derham-section-projective-space} 节中的计算，尤其是引理
\ref{derham-lemma-hodge-cohomology-projective-space} 的证明，有
$\text{pr}_{1, *}\text{pr}_2^*\Omega^s_{\mathbf{P}^{n - 1}_S/S} = 0$；
唯一的例外是 $s = 0$，此时得到
$\text{pr}_{1, *}\mathcal{O}_P = \mathcal{O}_Z$。这证明了 (2)。

\medskip\noindent
由第 \ref{derham-section-line-bundle} 节中的材料，尤其是引理
\ref{derham-lemma-push-omega-a}，有
$\pi_*\Omega^a_{L/S} = \Omega^a_{P/S} \oplus
\bigoplus_{k \geq 1} \Omega^a_{L/S, k}$.
由于上图中的复合 $\pi \circ 0$ 是 $P$ 上的恒等态射，为证明 (3)，
只需证明 $\Omega^a_{L/S, k}$ 对 $k > 0$ 有消失的高阶上同调。
由引理 \ref{derham-lemma-the-complex-for-L-star} 和 \ref{derham-lemma-push-omega-a}，
存在短正合列
$$
0 \to \Omega^a_{P/S} \otimes \mathcal{O}_P(k)
\to \Omega^a_{L/S, k} \to
\Omega^{a - 1}_{P/S} \otimes \mathcal{O}_P(k) \to 0
$$
，其中 $\Omega^{a - 1}_{P/S} = 0$，若 $a = 0$。由于
$P = Z \times_S \mathbf{P}^{n - 1}_S$，有
$$
\Omega^a_{P/S} = \bigoplus\nolimits_{i + j = a}
\Omega^i_{Z/S} \boxtimes \Omega^j_{\mathbf{P}^{n - 1}_S/S}
$$
，由引理 \ref{derham-lemma-de-rham-complex-product}。由于 $\Omega^i_{Z/S}$
是有限秩自由模，只需证明
$\mathcal{O}_Z \boxtimes \Omega^j_{\mathbf{P}^{n - 1}_S/S}(k)$
对 $k > 0$ 的高阶上同调为零。这由引理
\ref{derham-lemma-twisted-hodge-cohomology-projective-space} 应用于
$P = Z \times_S \mathbf{P}^{n - 1}_S = \mathbf{P}^{n - 1}_Z$ 得到，
从而 (3) 的证明完成。

\medskip\noindent
还需证明 (1)。若 $n = 1$，则我们吹起的是一个有效 Cartier 除子，
$b$ 是同构，故 (1) 成立。若 $n > 1$，则复合
$$
\Gamma(X, \Omega^a_{X/S})
\to
\Gamma(L, \Omega^a_{L/S})
\to
\Gamma(L \setminus E, \Omega^a_{L/S})
=
\Gamma(X \setminus Z, \Omega^a_{X/S})
$$
是同构，因为 $\Omega^a_{X/S}$ 是有限自由的（略去一个小细节）。
因此 (1) 失败的唯一可能是：$\Gamma(L, \Omega^a_{L/S})$ 中存在非零元素，
它在 $E = 0(P)$ 外消失。由于 $L$ 是 $P$ 上的线丛，零截面为
$0 : P \to L$，只需证明：在线丛上，微分形式层不存在这样的非零截面，
即它在零截面外恒为零。读者可以通过局部计算（这是较好的方法），
或利用 $\Omega_{L/S, k} \subset \Omega_{L^\star/S, k}$（见上述参考）
看出这一点。
\end{proof}

\noindent
此时我们建议读者跳到下一节。

\begin{lemma}
\label{derham-lemma-comparison-bis}
对 $a \geq 0$，存在典范映射
$$
b^*\Omega^a_{X/S} \longrightarrow
\Omega^a_{L/S} \longrightarrow
b^*\Omega^a_{X/S} \otimes_{\mathcal{O}_L} \mathcal{O}_L((n - 1)E)
$$
，其复合由包含
$\mathcal{O}_L \subset \mathcal{O}_L((n - 1)E)$.
诱导。
\end{lemma}

\begin{proof}
显示公式中的第一个箭头见第 \ref{derham-section-de-rham-complex} 节。
为得到第二个箭头，必须证明：若把 $\Omega^a_{L/S}$ 的局部截面视为
$b^*\Omega^a_{X/S}$ 的“亚纯截面”，则它的极点阶至多为
$n - 1$（沿 $E$）。为此，在 $L$ 上的仿射局部坐标图中工作。回忆 $L$ 被仿射吹起代数
$A[\frac{I}{y_i}]$ 的谱覆盖，其中 $I = A_{+}$ 是由
$y_1, \ldots, y_n$ 生成的理想。见《代数》第
\ref{algebra-section-blow-up} 节及《除子》，引理
\ref{divisors-lemma-blowing-up-affine}。由对称性，只需在对应于 $i = 1$
的坐标图上工作。此时
$$
A[\frac{I}{y_1}] = R[x_1, \ldots, x_m, y_1, t_2, \ldots, t_n]
$$
，其中 $t_i = y_i/y_1$，见《代数的进一步结果》，引理
\ref{more-algebra-lemma-blowup-regular-sequence}。因此模 $\Omega^1_{L/S}$
在相应仿射开集上由
$\text{d}x_1, \ldots, \text{d}x_m$、$\text{d}y_1$ 和
$\text{d}t_1, \ldots, \text{d}t_n$ 自由生成。当然，这些生成元中的前
$m + 1$ 个来自 $b^*\Omega^1_{X/S}$；对余下的 $n - 1$ 个，有
$$
\text{d}t_j =
\text{d}\frac{y_j}{y_1} =
\frac{1}{y_1}\text{d}y_j - \frac{y_j}{y_1^2}\text{d}y_1 =
\frac{\text{d}y_j - t_j \text{d}y_1}{y_1}
$$
，它有 $1$ 阶极点（沿 $E$），因为在此坐标图上 $E$ 由 $y_1$ 截出。
由于这些元素中任取 $a$ 个的外积给出 $\Omega^a_{L/S}$ 在此坐标图上的一组基，
且其中至多涉及 $n - 1$ 个 $\text{d}t_j$，证明完成。
\end{proof}

\begin{lemma}
\label{derham-lemma-blowup-twist-same-cohomology}
设 $E = 0(P)$ 是吹起 $b$ 的例外除子。对任意局部自由
$\mathcal{O}_X$-模 $\mathcal{E}$ 以及 $0 \leq i \leq n - 1$，映射
$$
\mathcal{E}
\longrightarrow
Rb_*(b^*\mathcal{E} \otimes_{\mathcal{O}_L} \mathcal{O}_L(iE))
$$
在 $D(\mathcal{O}_X)$ 中是同构。
\end{lemma}

\begin{proof}
由投影公式，只需对 $\mathcal{E} = \mathcal{O}_X$ 证明此结论，
见《上同调》，引理 \ref{cohomology-lemma-projection-formula}。
由于 $X$ 是仿射的，只需证明映射
$$
H^0(X, \mathcal{O}_X) \to
H^0(L, \mathcal{O}_L) \to
H^0(L, \mathcal{O}_L(iE))
$$
都是同构，并且 $H^j(X, \mathcal{O}_L(iE)) = 0$，其中 $j > 0$ 且
$0 \leq i \leq n - 1$；见《概形上同调》，引理
\ref{coherent-lemma-quasi-coherence-higher-direct-images-application}。
由于 $\pi$ 是仿射的，取 $\pi_*$ 后即可计算整体截面和上同调，
见《概形上同调》，引理 \ref{coherent-lemma-relative-affine-cohomology}。
若 $n = 1$，则 $L \to X$ 是同构且 $i = 0$，所以第一个陈述成立。
若 $n > 1$，则考虑复合
$$
H^0(X, \mathcal{O}_X) \to H^0(L, \mathcal{O}_L) \to
H^0(L, \mathcal{O}_L(iE)) \to H^0(L \setminus E, \mathcal{O}_L) =
H^0(X \setminus Z, \mathcal{O}_X)
$$
。由于
$H^0(X \setminus Z, \mathcal{O}_X) = H^0(X, \mathcal{O}_X)$
在此情形成立，因为 $Z$ 的余维为 $n \geq 2$（在 $X$ 中；细节从略），
从而第一个陈述成立。对于第二个陈述，回忆
$\mathcal{O}_L(E) = \mathcal{O}_L(-1)$，见《除子》，引理
\ref{divisors-lemma-blowing-up-gives-effective-Cartier-divisor}。因此
$$
\pi_*\mathcal{O}_L(iE) =
\pi_*\mathcal{O}_L(-i) =
\bigoplus\nolimits_{k \geq -i} \mathcal{O}_P(k)
$$
，如《态射的进一步结果》第 \ref{more-morphisms-section-proj-spec} 节所述。
最后，由 $P = \mathbf{P}^{n - 1}_Z$ 上结构层扭子的上同调消失即可得出结论；
该消失见《概形上同调》，引理
\ref{coherent-lemma-cohomology-projective-space-over-ring}。
\end{proof}























\section{吹起与德拉姆上同调}
\label{derham-section-blowing-up}

\noindent
固定基概形 $S$、光滑态射 $X \to S$，以及闭子概形
$Z \subset X$；它在 $S$ 上也光滑。记 $b : X' \to X$ 为 $X$ 沿 $Z$ 的吹起。
记 $E \subset X'$ 为例外除子。图示为
\begin{equation}
\label{derham-equation-blowup}
\vcenter{
\xymatrix{
E \ar[r]_j \ar[d]_p & X' \ar[d]^b \\
Z \ar[r]^i & X
}
}
\end{equation}
本节的目标是证明映射
$b^* : H_{dR}^*(X/S) \longrightarrow H_{dR}^*(X'/S)$
是单射（其实还可得到更多结论）。

\begin{lemma}
\label{derham-lemma-comparison-general}
采用《态射的进一步结果》，引理 \ref{more-morphisms-lemma-blowup} 中的记号。
对 $a \geq 0$，有
\begin{enumerate}
\item 映射 $\Omega^a_{X/S} \to b_*\Omega^a_{X'/S}$ 是同构；
\item 映射 $\Omega^a_{Z/S} \to p_*\Omega^a_{E/S}$ 是同构；
\item 映射 $Rb_*\Omega^a_{X'/S} \to i_*Rp_*\Omega^a_{E/S}$
在次数 $\geq 1$ 的上同调层上是同构。
\end{enumerate}
\end{lemma}

\begin{proof}
设 $\epsilon : X_1 \to X$ 是满平展态射。记
$i_1 : Z_1 \to X_1$、$b_1 : X'_1 \to X_1$、$E_1 \subset X'_1$ 以及
$p_1 : E_1 \to Z_1$ 为《态射的进一步结果》，引理
\ref{more-morphisms-lemma-blowup} 中所考虑对象的基变换。注意，$i_1$
是 $S$ 上光滑概形之间的闭浸入，并且由《除子》，引理
\ref{divisors-lemma-flat-base-change-blowing-up}，$b_1$ 是以 $Z_1$
为中心的吹起。假设能对态射 $b_1$、$p_1$ 和 $i_1$ 证明 (1)、(2)、(3)。
那么由引理 \ref{derham-lemma-etale}，(1)、(2)、(3) 中的映射经 $\epsilon$
拉回后都是同构。由于 $\epsilon$ 是满平坦态射，结论随即成立。
因此，由《态射的进一步结果》，引理
\ref{more-morphisms-lemma-etale-local-structure}，可以在平展局部工作，
并假设处于第 \ref{derham-section-calculations} 节讨论的情形。此时本引理就是
引理 \ref{derham-lemma-comparison}。
\end{proof}

\begin{lemma}
\label{derham-lemma-distinguished-triangle-blowup}
采用《态射的进一步结果》，引理 \ref{more-morphisms-lemma-blowup} 中的记号，
并记结构态射为 $f : X \to S$。存在典范可区别三角
$$
\Omega^\bullet_{X/S} \to
Rb_*(\Omega^\bullet_{X'/S}) \oplus i_*\Omega^\bullet_{Z/S} \to
i_*Rp_*(\Omega^\bullet_{E/S}) \to
\Omega^\bullet_{X/S}[1]
$$
于 $D(X, f^{-1}\mathcal{O}_S)$ 中，其中四个映射
$$
\begin{matrix}
\Omega^\bullet_{X/S} & \to & Rb_*(\Omega^\bullet_{X'/S}), \\
\Omega^\bullet_{X/S} & \to & i_*\Omega^\bullet_{Z/S}, \\
Rb_*(\Omega^\bullet_{X'/S}) & \to & i_*Rp_*(\Omega^\bullet_{E/S}), \\
i_*\Omega^\bullet_{Z/S} & \to & i_*Rp_*(\Omega^\bullet_{E/S})
\end{matrix}
$$
都是典范映射（第 \ref{derham-section-de-rham-complex} 节），
但其中一个的符号反转。
\end{lemma}

\begin{proof}
选取可区别三角
$$
C \to Rb_*\Omega^\bullet_{X'/S} \oplus i_*\Omega^\bullet_{Z/S}
\to i_*Rp_*\Omega^\bullet_{E/S} \to C[1]
$$
于 $D(X, f^{-1}\mathcal{O}_S)$ 中。只需证明
$\Omega^\bullet_{X/S}$ 与 $C$ 同构，并且该同构与典范映射相容。
由三角范畴公理，存在可区别三角之间的映射
$$
\xymatrix{
C' \ar[r] \ar[d] &
b_*\Omega^\bullet_{X'/S} \oplus i_*\Omega^\bullet_{Z/S} \ar[r] \ar[d] &
i_*p_*\Omega^\bullet_{E/S} \ar[r] \ar[d] &
C'[1] \ar[d] \\
C \ar[r] &
Rb_*\Omega^\bullet_{X'/S} \oplus i_*\Omega^\bullet_{Z/S} \ar[r] &
i_*Rp_*\Omega^\bullet_{E/S} \ar[r] &
C[1]
}
$$
由引理 \ref{derham-lemma-comparison-general} 的 (3) 以及《导出范畴》，命题
\ref{derived-proposition-9}，推出 $C' \to C$ 是同构。由引理
\ref{derham-lemma-comparison-general} 的 (2)，映射
$i_*\Omega^\bullet_{Z/S} \to i_*p_*\Omega^\bullet_{E/S}$ 是同构。
因此在导出范畴中 $C' = b_*\Omega^\bullet_{X'/S}$。最后，引理
\ref{derham-lemma-comparison-general} 的 (1) 告诉我们，它等于
$\Omega^\bullet_{X/S}$。我们略去其与典范映射相容的验证。
\end{proof}

\begin{proposition}
\label{derham-proposition-blowup-split}
采用《态射的进一步结果》，引理 \ref{more-morphisms-lemma-blowup} 中的记号。
映射 $\Omega^\bullet_{X/S} \to Rb_*\Omega^\bullet_{X'/S}$
在 $D(X, (X \to S)^{-1}\mathcal{O}_S)$ 中有分裂。
\end{proposition}

\begin{proof}
考虑引理 \ref{derham-lemma-distinguished-triangle-blowup} 中构造的三角形。
我们声称映射
$$
Rb_*(\Omega^\bullet_{X'/S}) \oplus i_*\Omega^\bullet_{Z/S} \to
i_*Rp_*(\Omega^\bullet_{E/S})
$$
有一个分裂，其像包含直和项 $i_*\Omega^\bullet_{Z/S}$。由《导出范畴》，
引理 \ref{derived-lemma-split}，这表明三角形的第一个箭头有一个分裂，
且该分裂在直和项 $i_*\Omega^\bullet_{Z/S}$ 上为零，从而证明该引理。
我们将把 $Rp_*\Omega^\bullet_{E/S}$ 分解为直和来证明此声称：
第一部分对应于 $\Omega^\bullet_{Z/S}$，第二部分可以通过
$Rb_*\Omega^\bullet_{X'/S}$ 提升。

\medskip\noindent
该声称的证明。可以把 $X$ 分解为相对 $S$ 具有固定维数的开闭子概形，
见《态射》，引理 \ref{morphisms-lemma-smooth-omega-finite-locally-free}。
由于导出范畴 $D(X, f^{-1}\mathcal{O})_S)$ 相应地分解为范畴的积，
可以假设 $X$ 具有固定相对维数 $N$（在 $S$ 上）。可以把
$Z = \coprod Z_m$ 分解为相对维数为 $m \geq 0$（在 $S$ 上）的开闭子概形。
态射 $i_m : Z_m \to X$ 是 $i$ 到 $Z_m$ 的限制，并且是余维为
$N - m$ 的正则浸入，
见《除子》，引理
\ref{divisors-lemma-immersion-smooth-into-smooth-regular-immersion}。
令 $E = \coprod E_m$ 为相应分解，即置 $E_m = p^{-1}(Z_m)$。
若 $p_m : E_m \to Z_m$ 表示 $p$ 到 $E_m$ 的限制，则有典范同构
$$
\tilde \xi_m :
\bigoplus\nolimits_{t = 0, \ldots, N - m - 1}
\Omega^\bullet_{Z_m/S}[-2t]
\longrightarrow
Rp_{m, *}\Omega^\bullet_{E_m/S}
$$
于 $D(Z_m, (Z_m \to S)^{-1}\mathcal{O}_S)$ 中；在次数 $0$ 上，它是典范映射
$\Omega^\bullet_{Z_m/S} \to Rp_{m, *}\Omega^\bullet_{E_m/S}$.
见评注 \ref{derham-remark-projective-space-bundle-formula}。因此有同构
$$
\tilde \xi :
\bigoplus\nolimits_m
\bigoplus\nolimits_{t = 0, \ldots, N - m - 1}
\Omega^\bullet_{Z_m/S}[-2t]
\longrightarrow
Rp_*(\Omega^\bullet_{E/S})
$$
于 $D(Z, (Z \to S)^{-1}\mathcal{O}_S)$ 中；它到源的直和项
$\Omega^\bullet_{Z/S} = \bigoplus \Omega^\bullet_{Z_m/S}$ 的限制
是典范映射 $\Omega^\bullet_{Z/S} \to Rp_*(\Omega^\bullet_{E/S})$。
考虑复形的子复形 $M_m$ 和 $K_m$，这里该复形为
$\bigoplus\nolimits_{t = 0, \ldots, N - m - 1} \Omega^\bullet_{Z_m/S}[-2t]$
，它们在评注 \ref{derham-remark-projective-space-bundle-formula} 中引入。置
$$
M = \bigoplus M_m
\quad\text{且}\quad
K = \bigoplus K_m
$$
。有 $M = K[-2]$，并且由构造，映射
$$
c_{E/Z} \oplus \tilde \xi|_M :
\Omega^\bullet_{Z/S} \oplus M
\longrightarrow
Rp_*(\Omega^\bullet_{E/S})
$$
是同构（见上述评注）。

\medskip\noindent
考虑映射
$$
\delta : \Omega^\bullet_{E/S}[-2] \longrightarrow \Omega^\bullet_{X'/S}
$$
，它是 $D(X', (X' \to S)^{-1}\mathcal{O}_S)$ 中引理
\ref{derham-lemma-gysin-via-log-complex} 的映射，并具有如下性质：复合
$$
\Omega^\bullet_{E/S}[-2] \longrightarrow \Omega^\bullet_{X'/S}
\longrightarrow
\Omega^\bullet_{E/S}
$$
是映射 $\theta'$；它见评注 \ref{derham-remark-cup-product-as-a-map}，并对应于
$c_1^{dR}(\mathcal{O}_{X'}(-E))|_E) = c_1^{dR}(\mathcal{O}_E(1))$。
评注 \ref{derham-remark-projective-space-bundle-formula}
的最后一个断言表明图表
$$
\xymatrix{
K[-2] \ar[d]_{(\tilde \xi|_K)[-2]} \ar[r]_{\text{id}} &
M \ar[d]^{\tilde x|_M} \\
Rp_*\Omega^\bullet_{E/S}[-2] \ar[r]^-{Rp_*\theta'} &
Rp_*\Omega^\bullet_{E/S}
}
$$
交换。因此，可用如下映射得到该声称所需的分裂：
\begin{align*}
Rp_*(\Omega^\bullet_{E/S})
& \xrightarrow{(c_{E/Z} \oplus \tilde \xi|_M)^{-1}}
\Omega^\bullet_{Z/S} \oplus M \\
& \xrightarrow{\text{id} \oplus \text{id}^{-1}}
\Omega^\bullet_{Z/S} \oplus K[-2] \\
& \xrightarrow{\text{id} \oplus (\tilde \xi|_K)[-2]}
\Omega^\bullet_{Z/S} \oplus Rp_*\Omega^\bullet_{E/S}[-2] \\
& \xrightarrow{\text{id} \oplus Rb_*\delta}
\Omega^\bullet_{Z/S} \oplus Rb_*\Omega^\bullet_{X'/S}
\end{align*}
上面所述 $\theta'$ 与 $\delta$ 的关系，再加上上面包含
$\theta'$、$\tilde \xi|_K$ 和 $\tilde \xi|_M$ 的交换图表，恰好说明
这是典范映射
$\Omega^\bullet_{Z/S} \oplus Rb_*(\Omega^\bullet_{X'/S}) \to
Rp_*(\Omega^\bullet_{E/S})$ 的一个截面。该声称的证明完成。
\end{proof}

\noindent
引理 \ref{derham-lemma-splitting-on-omega-a} 表明，在霍奇上同调上构造分裂，
比命题 \ref{derham-proposition-blowup-split} 的结果容易得多。
我们力劝读者跳到下一节。

\begin{lemma}
\label{derham-lemma-ext-zero}
设 $i : Z \to X$ 是余维为 $c$ 的正则概形闭浸入。则
$\Ext^q_{\mathcal{O}_X}(i_*\mathcal{F}, \mathcal{E}) = 0$，只要 $q < c$；
这里 $\mathcal{E}$ 在 $X$ 上局部自由，而 $\mathcal{F}$ 是任意
$\mathcal{O}_Z$-模。
\end{lemma}

\begin{proof}
由 $\Ext$ 的局部到整体谱序列，只需在 $X$ 上仿射局部地证明，
见《上同调》第 \ref{cohomology-section-ext} 节。因此可以假设
$X = \Spec(A)$，且存在正则序列 $f_1, \ldots, f_c$，它位于 $A$ 中，使得
$Z = \Spec(A/(f_1, \ldots, f_c))$。可以假设 $c \geq 1$。于是
$f_1 : \mathcal{E} \to \mathcal{E}$ 是单射。由于 $i_*\mathcal{F}$
被 $f_1$ 零化，这表明引理对 $i = 0$ 成立，并且有满射
$$
\Ext^{q - 1}_{\mathcal{O}_X}(i_*\mathcal{F}, \mathcal{E}/f_1\mathcal{E})
\longrightarrow
\Ext^q_{\mathcal{O}_X}(i_*\mathcal{F}, \mathcal{E})
$$
。所以只需证明该箭头的源为零。接着重复这个论证：若 $c \geq 2$，则映射
$f_2 : \mathcal{E}/f_1\mathcal{E} \to \mathcal{E}/f_1\mathcal{E}$
是单射。由于 $i_*\mathcal{F}$ 被 $f_2$ 零化，这表明引理对 $q = 1$
成立，并且有满射
$$
\Ext^{q - 2}_{\mathcal{O}_X}(i_*\mathcal{F},
\mathcal{E}/f_1\mathcal{E} + f_2\mathcal{E})
\longrightarrow
\Ext^{q - 1}_{\mathcal{O}_X}(i_*\mathcal{F}, \mathcal{E}/f_1\mathcal{E})
$$
如此继续即可证明引理。
\end{proof}

\begin{lemma}
\label{derham-lemma-splitting-on-omega-a}
采用《态射的进一步结果》，引理 \ref{more-morphisms-lemma-blowup} 中的记号。
对 $a \geq 0$，存在唯一箭头
$Rb_*\Omega^a_{X'/S} \to \Omega^a_{X/S}$ 于 $D(\mathcal{O}_X)$ 中；它与
$\Omega^a_{X/S} \to Rb_*\Omega^a_{X'/S}$ 的复合是
$\Omega^a_{X/S}$ 上的恒等态射。
\end{lemma}

\begin{proof}
可以把 $X$ 分解为相对 $S$ 具有固定维数的开闭子概形，见《态射》，引理
\ref{morphisms-lemma-smooth-omega-finite-locally-free}。由于导出范畴
$D(X, f^{-1}\mathcal{O})_S)$ 相应地分解为范畴的积，可以假设
$X$ 具有固定相对维数 $N$（在 $S$ 上）。可以把
$Z = \coprod Z_m$ 分解为相对维数为 $m \geq 0$（在 $S$ 上）的开闭子概形。
态射 $i_m : Z_m \to X$ 是 $i$ 到 $Z_m$ 的限制，并且是余维为
$N - m$ 的正则浸入，见《除子》，引理
\ref{divisors-lemma-immersion-smooth-into-smooth-regular-immersion}。
令 $E = \coprod E_m$ 为相应分解，即置 $E_m = p^{-1}(Z_m)$。
我们声称存在自然映射
$$
b^*\Omega^a_{X/S} \to \Omega^a_{X'/S} \to
b^*\Omega^a_{X/S} \otimes_{\mathcal{O}_{X'}}
\mathcal{O}_{X'}(\sum (N - m - 1)E_m)
$$
，其复合由包含
$\mathcal{O}_{X'} \to \mathcal{O}_{X'}(\sum (N - m - 1)E_m)$.
诱导。为证明这一点，只需说明第一个箭头的余核在 $X'$ 上局部地被有效
Cartier 除子 $\sum (N - m - 1)E_m$ 的一个局部方程零化。为此，
可以像引理 \ref{derham-lemma-comparison-general} 的证明那样在 $X$ 上平展局部地
工作，并应用引理 \ref{derham-lemma-comparison-bis}。利用引理
\ref{derham-lemma-blowup-twist-same-cohomology} 在平展局部计算，可见所诱导的复合
$$
\Omega^a_{X/S} \to Rb_*\Omega^a_{X'/S} \to
Rb_*\left(b^*\Omega^a_{X/S} \otimes_{\mathcal{O}_{X'}}
\mathcal{O}_{X'}(\sum (N - m - 1)E_m)\right)
$$
在 $D(\mathcal{O}_X)$ 中是同构；由此得到引理中映射的存在性。

\medskip\noindent
对于唯一性，只需证明不存在从 $\tau_{\geq 1}Rb_*\Omega_{X'/S}$
到 $\Omega^a_{X/S}$ 的非零映射（在 $D(\mathcal{O}_X)$ 中）。
为此，只需再证明不存在从 $R^qb_*\Omega_{X'/s}[-q]$
到 $\Omega^a_{X/S}$ 的非零映射（在 $D(\mathcal{O}_X)$ 中），其中
$q \geq 1$（细节从略）。由引理 \ref{derham-lemma-comparison-general} 可见，
$R^qb_*\Omega_{X'/s} \cong i_*R^qp_*\Omega^a_{E/S}$ 是定义在
$Z = \coprod Z_m$ 上的一个模的推前。此外，注意
$R^qp_*\Omega^a_{E/S}$ 到 $Z_m$ 的限制只有在 $q < N - m$ 时才可能非零。
事实上，$E_m \to Z_m$ 的纤维维数为 $N - m - 1$，可以应用《极限》，引理
\ref{limits-lemma-higher-direct-images-zero-above-dimension-fibre}。
因此所需消失由引理 \ref{derham-lemma-ext-zero} 得到。
\end{proof}














\section{比较微分形式层}
\label{derham-section-quasi-finite-syntomic}

\noindent
本节的目标是：对每个局部拟有限的 syntomic 概形态射 $f : Y \to X$，
构造映射
$$
c^p_{Y/X} :
\Omega^p_{Y/\mathbf{Z}}
\longrightarrow
f^*\Omega^p_{X/\mathbf{Z}} \otimes_{\mathcal{O}_Y} \det(\NL_{Y/X})
$$
，对所有 $p \geq 0$ 均满足下列性质：
\begin{enumerate}
\item
\label{derham-item-degree-zero}
$c_{Y/X}^0(1) = \delta(\NL_{Y/X})$，见《判别式》第
\ref{discriminant-section-tate-map} 节；
\item
\label{derham-item-multiplicative}
$c_{Y/X}^{q + p}(\omega \wedge \eta) = \omega \wedge c_{Y/X}^p(\eta)$
，其中 $\omega$ 是 $f^*\Omega^q_{X/\mathbf{Z}}$ 的局部截面，
而 $\eta$ 是 $\Omega^p_{Y/\mathbf{Z}}$ 的局部截面；
\item
\label{derham-item-base-change}
$c^p_{Y/X}$ 的构造与到开集的限制以及基变换相容
（表述见引理 \ref{derham-lemma-base-change-Garel-upstairs}）。
\end{enumerate}
注意，这些条件蕴含映射 $c^p_{Y/X}$ 是 $\mathcal{O}_Y$-线性的，
并且它与
$f^*\Omega^p_{X/\mathbf{Z}} \to \Omega^p_{Y/\mathbf{Z}}$
的复合是乘以 $\delta(\NL_{Y/X})$。事实上，我们将证明存在唯一的一族算子
$c^p_{Y/X}$ 具有性质 (1)、(2)、(3)。若 $f$ 是有限的，第
\ref{derham-section-trace} 节会把这个结果用于构造德拉姆复形上的迹映射。

\begin{lemma}
\label{derham-lemma-funny-map}
设 $R$ 是环，并考虑交换图表
$$
\xymatrix{
0 \ar[r] &
K^0 \ar[r] &
L^0 \ar[r] &
M^0 \ar[r] & 0 \\
& & L^{-1} \ar[u]_\partial \ar@{=}[r] &
M^{-1} \ar[u]
}
$$
，其中各项为 $R$-模，顶行正合，而 $M^0$ 和 $M^{-1}$
是秩相同的有限自由模。对 $p \geq 0$，存在典范映射
$$
c^p : 
\wedge^p(H^0(L^\bullet))
\longrightarrow
\wedge^p(K^0) \otimes_R \det(M^\bullet)
$$
，具有下列性质：
\begin{enumerate}
\item $c^0(1) = \delta(M^\bullet)$；
\item $c^{q + p}(\omega \wedge \eta) = \omega \wedge c^p(\eta)$
，其中 $\omega \in \wedge^q(K^0)$ 且
$\eta \in \wedge^p(H^0(L^\bullet))$。
\end{enumerate}
特别地，$c^p$ 与映射
$\wedge^p(K^0) \to \wedge^p(H^0(L^\bullet))$ 的复合是乘以
$\delta(M^\bullet)$。
\end{lemma}

\begin{proof}
设 $M^0$ 和 $M^{-1}$ 都是秩为 $n$ 的自由模。对每个 $p \geq 0$，
存在唯一的满射
$$
\pi_p :
\wedge^{p + n}(L^0)
\longrightarrow
\wedge^p(K^0) \otimes \wedge^n(M^0)
$$
，具有下列两个性质；(i) 有
$
\pi_p(k_1 \wedge \ldots \wedge k_p \wedge l_1 \wedge \ldots \wedge l_n) =
k_1 \wedge \ldots \wedge k_p \otimes
m_1 \wedge \ldots \wedge m_n
$
，如果 $k_1, \ldots, k_p \in K^0$，并且
$l_1, \ldots, l_n \in L^0$ 映到 $m_1, \ldots, m_n \in M^0$；
(ii) $\pi_p$ 的核是由外积 $l_1 \wedge \ldots \wedge l_{p + n}$
生成的子模，其中 $> p$ 个 $l_j$ 属于 $K^0$。取一组基
$e_1, \ldots, e_n$，它是 $L^{-1} = M^{-1}$ 的基。然后按如下规则定义映射：
$$
\eta \mapsto
\pi_p(\tilde \eta \wedge \partial(e_1) \wedge \ldots \wedge \partial(e_n))
\otimes (e_1 \wedge \ldots \wedge e_n)^{\otimes -1}
$$
，其中 $\eta \in \wedge^p(H^0(L^\bullet))$，而
$\tilde \eta \in \wedge^p(L^0)$ 是 $\eta$ 的提升。右边的表达式有意义，因为
$\det(M^\bullet) = \wedge^n(M^0) \otimes \wedge^n(M^{-1})^{\otimes -1}$.
右边的表达式与 $\tilde \eta$ 的选择无关，因为
$\partial(\xi) \wedge \partial(e_1) \wedge \ldots \wedge \partial(e_n) = 0$
，对任意 $\xi \in L^{-1}$。直接计算表明，该映射与
$L^{-1} = M^{-1}$ 的基的选择无关。由《代数的进一步结果》第
\ref{more-algebra-section-determinants-complexes} 节中
$\delta(M^\bullet)$ 的定义，立即得到
$c^0(1) = \delta(M^\bullet)$。最后，陈述 (2) 中的规则可由定义直接验证。
\end{proof}

\begin{lemma}
\label{derham-lemma-funny-map-functorial}
设 $R_1 \to R_2$ 是环同态。对 $i = 1, 2$，考虑交换图表
$$
\xymatrix{
0 \ar[r] &
K_i^0 \ar[r] &
L_i^0 \ar[r] &
M_i^0 \ar[r] & 0 \\
& & L_i^{-1} \ar[u]_{\partial_i} \ar@{=}[r] &
M_i^{-1} \ar[u]
}
$$
，其中各项为 $R_i$-模，如引理 \ref{derham-lemma-funny-map} 所述。假设有映射
$K_1^0 \to K_2^0$、$L_1^j \to L_2^j$、$M_1^j \to M_2^j$，
它们与给定环同态 $R_1 \to R_2$ 以及显示图表中的映射相容。若映射
$M_1^j \otimes_{R_1} R_2 \to M_2^j$ 都是同构，则图表
$$
\xymatrix{ 
\wedge^p(H^0(L_1^\bullet)) \ar[d]
\ar[r]_-{c^p} &
\wedge^p(K_1^0) \otimes_{R_1} \det(M_1^\bullet) \ar[d] \\
\wedge^p(H^0(L_2^\bullet))
\ar[r]^-{c^p} &
\wedge^p(K_2^0) \otimes_{R_2} \det(M_2^\bullet)
}
$$
交换。
\end{lemma}

\begin{proof}
这由引理 \ref{derham-lemma-funny-map} 的证明中给出的映射的显式描述得到。
注意，需要假设 $M_1^j \otimes_{R_1} R_2 \to M_2^j$ 都是同构，
才能看出我们用于 $M_1^{-1}$ 的基映到 $M_2^{-1}$ 的一组基
（因而在检验两个构造的相容性时可以使用“同一”组基）。
\end{proof}

\begin{remark}
\label{derham-remark-local-description}
设 $A$ 是一个环，并令 $P = A[x_1, \ldots, x_n]$。取
$f_1, \ldots, f_n \in P$，并置 $B = P/(f_1, \ldots, f_n)$。
假设 $A \to B$ 是拟有限的。于是
$B$ 是 $A$ 上的相对全局完全交（代数，定义
\ref{algebra-definition-relative-global-complete-intersection}），且由代数，引理
\ref{algebra-lemma-relative-global-complete-intersection-conormal}，
$(f_1, \ldots, f_n)/(f_1, \ldots, f_n)^2$ 是以各类
$\overline{f}_i$ 为生成元的自由模。考虑下图
$$
\xymatrix{
\Omega_{A/\mathbf{Z}} \otimes_A B \ar[r] &
\Omega_{P/\mathbf{Z}} \otimes_P B \ar[r] &
\Omega_{P/A} \otimes_P B \\
&
(f_1, \ldots, f_n)/(f_1, \ldots, f_n)^2 \ar[u] \ar@{=}[r] &
(f_1, \ldots, f_n)/(f_1, \ldots, f_n)^2 \ar[u]
}
$$
右列表示 $\NL_{B/A}$ 在 $D(B)$ 中的一个代表，因而有上同调
$\Omega_{B/A}$ 位于次数 $0$。顶行是分裂短正合列
$0 \to \Omega_{A/\mathbf{Z}} \otimes_A B \to
\Omega_{P/\mathbf{Z}} \otimes_P B \to \Omega_{P/A} \otimes_P B \to 0$.
由代数，引理 \ref{algebra-lemma-differential-seq}，中列有上同调
$\Omega_{B/\mathbf{Z}}$ 位于次数 $0$。
因此由引理 \ref{derham-lemma-funny-map}，我们得到典范的 $B$-模映射
$$
\Omega^p_{B/\mathbf{Z}} \longrightarrow
\Omega^p_{A/\mathbf{Z}} \otimes_A \det(\NL_{B/A})
$$
其中 $p \geq 0$，这些映射满足性质 (1) 和 (2)；特别地，它们与映射
$\Omega^p_{A/\mathbf{Z}} \to \Omega^p_{B/\mathbf{Z}}$
的复合是乘以 $\delta(\NL_{B/A})$。
\end{remark}

\begin{lemma}
\label{derham-lemma-base-change-Garel-upstairs}
考虑概形的交换图
$$
\xymatrix{
Y' \ar[d]_{f'} \ar[r]_b & Y \ar[d]^f \\
X' \ar[r]^a & X
}
$$
它将 $Y'$ 同构到 $X' \times_X Y$ 的一个开子概形。假设 $f$ 是
局部拟有限的 syntomic 态射，并且给定映射
$c^p_{Y/X} : \Omega^p_{Y/\mathbf{Z}} \to
f^*\Omega^p_{X/\mathbf{Z}} \otimes_{\mathcal{O}_Y} \det(\NL_{Y/X})$
（$p \geq 0$），满足
(\ref{derham-item-degree-zero}) 和 (\ref{derham-item-multiplicative})。
那么至多存在一族映射
$c^p_{Y'/X'} : \Omega^p_{Y'/\mathbf{Z}} \to
(f')^*\Omega^p_{X'/\mathbf{Z}} \otimes_{\mathcal{O}_{Y'}} \det(\NL_{Y'/X'})$
（$p \geq 0$），满足 (\ref{derham-item-degree-zero}) 和
(\ref{derham-item-multiplicative})，并使诸图
$$
\xymatrix{
b^*\Omega^p_{Y/\mathbf{Z}} \ar[rr]_-{b^*c^p_{Y/X}} \ar[d] & &
b^*(f^*\Omega^p_{X/\mathbf{Z}} \otimes_{\mathcal{O}_Y} \det(\NL_{Y/X}))
\ar@{=}[r] &
(f')^*a^*\Omega^p_{X/\mathbf{Z}} \otimes_{\mathcal{O}_Y} b^*\det(\NL_{Y/X}))
\ar[d] \\
\Omega^p_{Y'/\mathbf{Z}} \ar[rrr]^-{c^p_{Y'/X'}} & & &
(f')^*\Omega^p_{X'/\mathbf{Z}} \otimes_{\mathcal{O}_{Y'}} \det(\NL_{Y'/X'})
}
$$
对所有 $p \geq 0$ 交换。这里的竖直箭头使用第
\ref{derham-section-de-rham-complex} 节中的映射
$b^*\Omega^p_{Y/\mathbf{Z}} \to \Omega^p_{Y'/\mathbf{Z}}$ 和
$a^*\Omega^p_{X/\mathbf{Z}} \to \Omega^p_{X'/\mathbf{Z}}$
以及判别式，第 \ref{discriminant-section-tate-map} 节中的等同
$b^*\det(\NL_{Y/X}) = \det(\NL_{Y'/X'})$。
\end{lemma}

\begin{proof}
映射
$$
(f')^*\Omega_{X'/\mathbf{Z}} \oplus b^*\Omega_{Y/\mathbf{Z}}
\longrightarrow
\Omega_{Y'/\mathbf{Z}}
$$
是满射，因为 $Y'$ 是该纤维积的一个开子概形；相应的代数陈述是：
$\Omega_{B \otimes_A C/\mathbf{Z}}$ 作为模由
$\Omega_{B/\mathbf{Z}}$ 和 $\Omega_{C/\mathbf{Z}}$ 的像生成。
因此对所有 $p$，映射
$$
\bigoplus\nolimits_{i = 0, \ldots, p}
(f')^*\Omega^i_{X'/\mathbf{Z}} \otimes
b^*\Omega^{p - i}_{Y/\mathbf{Z}}
\longrightarrow
\Omega^p_{Y'/\mathbf{Z}}
$$
是满射。条件 (\ref{derham-item-degree-zero})、(\ref{derham-item-multiplicative})
连同引理中图的交换性表明，映射 $c^p_{Y'/X'}$
已被唯一确定（但其存在性并非立即可得）。
\end{proof}

\begin{lemma}
\label{derham-lemma-existence-base-change-Garel-upstairs}
考虑交换图
$$
\xymatrix{
Y' \ar[d]_{f'} \ar[r]_b & Y \ar[d]^f \\
X' \ar[r]^a & X
}
$$
它将 $Y'$ 同构到 $X' \times_X Y$ 的一个开子概形。假设 $f$ 是局部拟有限的
syntomic 态射。那么对每个 $y' \in Y'$，都可以找到开集
$V' \subset Y'$、$V \subset Y$、$U' \subset X$、$U \subset X$，
使得 $y' \in V'$，并且
$f'(V') \subset U'$, $b(V') \subset V$, $a(U') \subset U$,
$f(V) \subset U$；而且对 $p \geq 0$ 存在映射
$c^p_{V/U}$ 和 $c^p_{V'/U'}$，它们满足 (\ref{derham-item-degree-zero})
和 (\ref{derham-item-multiplicative})，并与图
$$
\xymatrix{
V' \ar[d] \ar[r]_b & V \ar[d] \\
U' \ar[r] & U
}
$$
相容，其含义如引理 \ref{derham-lemma-base-change-Garel-upstairs} 中所述。
\end{lemma}

\begin{proof}
显然，为了证明这一点，可以把 $Y'$ 替换为 $X' \times_X Y$。
按照判别式，引理 \ref{discriminant-lemma-syntomic-quasi-finite} 的 (5)，
取仿射开集 $V \subset Y$ 和 $U \subset X$，其中 $V$ 包含 $y'$ 的像。
取仿射开集 $U' \subset X'$，它包含 $y'$ 的像并映入 $U$。
把 $X, Y, X', Y'$ 分别替换为 $U, V, U', U' \times_U V$ 后，
我们可以假设处于下一段所述的情形。

\medskip\noindent
假设 $X' = \Spec(A')$、$X = \Spec(A)$、$Y = \Spec(B)$，且
$Y' = \Spec(A' \otimes_A B)$，其中
$B = A[x_1, \ldots, x_n]/(f_1, \ldots, f_n)$
是 $A$ 上的相对全局完全交。那么
$B' = A'[x_1, \ldots, x_n]/(f'_1, \ldots, f'_n)$
也是 $A'$ 上的相对全局完全交，其中 $f'_i$ 是 $f_i$ 在
$A'[x_1, \ldots, x_n]$ 中的像；见代数，引理
\ref{algebra-lemma-base-change-relative-global-complete-intersection}。
备注 \ref{derham-remark-local-description} 中的构造给出了映射
$c^p_{V/U}$ 和 $c^p_{V'/U'}$。
由引理 \ref{derham-lemma-funny-map-functorial} 以及 $B$、$B'$ 分别在
$A$、$A'$ 上的表示之间的相容性，这些映射彼此相容。略去若干细节。
\end{proof}

\begin{lemma}
\label{derham-lemma-Garel-upstairs}
存在唯一的规则，它为每个概形的局部拟有限 syntomic 态射
$f : Y \to X$ 指定 $\mathcal{O}_Y$-模映射
$$
c^p_{Y/X} :
\Omega^p_{Y/\mathbf{Z}}
\longrightarrow
f^*\Omega^p_{X/\mathbf{Z}} \otimes_{\mathcal{O}_Y} \det(\NL_{Y/X})
$$
其中 $p \geq 0$，并满足 (\ref{derham-item-degree-zero})、
(\ref{derham-item-multiplicative}) 和 (\ref{derham-item-base-change})。特别地，
$c^p_{Y/X}$ 与 $f^*\Omega^p_{X/\mathbf{Z}} \to \Omega^p_{Y/\mathbf{Z}}$
的复合是乘以 $\delta(\NL_{Y/X})$。
\end{lemma}

\begin{proof}
此证明与判别式，命题 \ref{discriminant-proposition-tate-map}
的证明非常相似，建议读者先阅读该证明。

\medskip\noindent
我们先重新表述该陈述。考虑范畴 $\mathcal{C}$：其对象记为
$Y/X$，是概形的局部拟有限 syntomic 态射 $f : Y \to X$；
其态射 $b/a : Y'/X' \to Y/X$ 是交换图
$$
\xymatrix{
Y' \ar[d]_{f'} \ar[r]_b & Y \ar[d]^f \\
X' \ar[r]^a & X
}
$$
它将 $Y'$ 同构到 $X' \times_X Y$ 的一个开子概形。引理的含义是：
对每个对象 $Y/X$（属于 $\mathcal{C}$），都有映射
$c^p_{Y/X}$（$p \geq 0$）满足性质 (\ref{derham-item-degree-zero})
和 (\ref{derham-item-multiplicative})；并且对每个态射
$b/a : Y'/X' \to Y/X$（属于 $\mathcal{C}$），映射
$c^p_{Y/X}$ 和 $c^p_{Y'/X'}$
按引理 \ref{derham-lemma-base-change-Garel-upstairs} 的意义相容。

\medskip\noindent
给定 $Y/X$（属于 $\mathcal{C}$）以及 $y \in Y$，可以找到
仿射开集 $V \subset Y$ 和 $U \subset X$，满足 $f(V) \subset U$，
使得存在一些映射
$$
\Omega^p_{Y/\mathbf{Z}}|_V
\longrightarrow
\left(
f^*\Omega^p_{X/\mathbf{Z}} \otimes_{\mathcal{O}_Y} \det(\NL_{Y/X})
\right)|_V
$$
满足性质 (\ref{derham-item-degree-zero}) 和 (\ref{derham-item-multiplicative})。
这是因为可以按照判别式，引理
\ref{discriminant-lemma-syntomic-quasi-finite} 的 (5) 选取仿射开集，
再使用备注 \ref{derham-remark-local-description} 中的构造。

\medskip\noindent
注意，$f$ 的平展轨迹恰好是 $\delta(\NL_{Y/X})$ 不消失的点集；
见判别式，第 \ref{discriminant-section-tate-map} 节中的讨论。特别地，
$f^*\Omega^p_{X/\mathbf{Z}} \to \Omega^p_{Y/\mathbf{Z}}$
在逆转 $\delta(\NL_{Y/X})$ 后成为同构。由此可知，如果
$\Omega^p_{X/\mathbf{Z}}$ 是有限局部自由的，且
$\delta(\NL_{Y/X})$ 在 $\mathcal{O}_Y$ 中的零化子为零，
那么上一段构造的局部映射是唯一的，并会自动粘合！

\medskip\noindent
令 $\mathcal{C}_{nice} \subset \mathcal{C}$ 表示由满足下列条件的
$Y/X$ 构成的全子范畴：
\begin{enumerate}
\item $\Omega_{X/\mathbf{Z}}$ 是局部自由的；
\item $\delta(\NL_{Y/X})$ 在 $\mathcal{O}_Y$ 中的零化子为零。
\end{enumerate}
由上一段的说明，对每个对象 $Y/X$（属于 $\mathcal{C}_{nice}$），
都有唯一的映射 $c^p_{Y/X}$（$p \geq 0$）满足条件
(\ref{derham-item-degree-zero}) 和 (\ref{derham-item-multiplicative})。
如果 $b/a : Y'/X' \to Y/X$ 是 $\mathcal{C}_{nice}$ 的一个态射，
那么映射 $c^p_{Y/X}$ 和 $c^p_{Y'/X'}$ 按照引理
\ref{derham-lemma-base-change-Garel-upstairs} 的意义相容：事实上，由引理
\ref{derham-lemma-existence-base-change-Garel-upstairs}，局部上确实存在相容映射；
而刚才的唯一性说明这些映射分别等于
$c^p_{Y/X}$ 和 $c^p_{Y'/X'}$。
换言之，我们已经在全子范畴 $\mathcal{C}_{nice}$ 上解决了该问题。
对 $Y/X$（属于 $\mathcal{C}_{nice}$），仍用 $c^p_{Y/X}$
表示刚刚得到的解。

\medskip\noindent
事实上，更一般地，假设有态射 $b/a : Y'/X' \to Y/X$
（属于 $\mathcal{C}$），并且 $Y/X$ 是 $\mathcal{C}_{nice}$ 的对象。
同一论证（这次使用引理 \ref{derham-lemma-base-change-Garel-upstairs}
中的唯一性）告诉我们，存在映射 $c^p_{Y'/X'}$（$p \geq 0$），
它们满足条件 (\ref{derham-item-degree-zero}) 和 (\ref{derham-item-multiplicative})，
并与已经构造的映射 $c^p_{Y/X}$ 相容。我们把它称为基变换
$(b/a)^*c^p_{Y/X}$。

\medskip\noindent
考虑态射
$$
Y_1/X_1 \xleftarrow{b_1/a_1} Y/X \xrightarrow{b_2/a_2} Y_2/X_2
$$
（属于 $\mathcal{C}$），并假设 $Y_1/X_1$ 和 $Y_2/X_2$
是 $\mathcal{C}_{nice}$ 的对象。
{\bf 断言。} 两个基变换 $(b_i/a_i)^*c^p_{Y_i/X_i}$（$i = 1, 2$）
相等。我们先说明该断言蕴含引理，然后再证明该断言。

\medskip\noindent
令 $d, n \geq 1$，并考虑判别式，例
\ref{discriminant-example-universal-quasi-finite-syntomic}
中构造的局部拟有限 syntomic 态射 $Y_{n, d} \to X_{n, d}$。
那么 $Y_{n, d}$ 和 $Y_{n, d}$ 是 $\mathbf{Z}$ 上光滑的
有限型不可约概形。事实上，$X_{n, d}$ 是 $\mathbf{Z}$ 上某个多项式环的谱，
而 $Y_{n, d}$ 是这种谱的一个开子概形。态射
$Y_{n, d} \to X_{n, d}$ 是局部拟有限的 syntomic 态射，并在某个稠密开集上平展；
见判别式，引理
\ref{discriminant-lemma-universal-quasi-finite-syntomic-etale}。
因此 $\delta(\NL_{Y_{n, d}/X_{n, d}})$ 非零：例如，判别式，备注
\ref{discriminant-remark-local-description-delta} 给出了
$\delta(\NL_{Y/X})$ 的局部描述，而态射，引理
\ref{morphisms-lemma-etale-at-point} 的 (8) 给出了平展态射的局部描述。
不可约正则概形上的可逆模的非零截面具有零零化子。因此
$Y_{n, d}/X_{n, d}$ 是 $\mathcal{C}_{nice}$ 的对象。

\medskip\noindent
令 $Y/X$ 是 $\mathcal{C}$ 的任意对象，并令 $y \in Y$。
由判别式，引理 \ref{discriminant-lemma-locally-comes-from-universal}，
可以找到 $n, d \geq 1$ 以及态射
$$
Y/X \leftarrow V/U \xrightarrow{b/a} Y_{n, d}/X_{n, d}
$$
（属于 $\mathcal{C}$），使得 $V \subset Y$ 和 $U \subset X$ 都是开集。
于是有基变换 $c^p_{V/U} = (b/a)^*c^p_{Y_{n, d}/X_{n, d}}$。
该断言保证这些局部构造的映射 $c^p_{V/U}$ 能够粘合！
因此得到满足性质 (\ref{derham-item-degree-zero}) 和
(\ref{derham-item-multiplicative}) 的良定义全局映射 $c^p_{Y/X}$。
最后，令 $b/a : Y'/X' \to Y/X$ 是 $\mathcal{C}$ 的任意态射。
本段已经构造了映射 $c^p_{Y'/X'}$（$p \geq 0$）和
$c^p_{Y/X}$（$p \geq 0$）。为了检验它们按引理
\ref{derham-lemma-base-change-Garel-upstairs} 的意义相容，可以在
$Y'/X'$ 和 $Y/X$ 上局部地工作。因此可以假设存在态射
$Y/X \to Y_{n, d}/X_{n, d}$。依构造，族
$c^p_{Y'/X'}$ 和族 $c^p_{X/Y}$ 都与指向
$Y_{n, d}/X_{n, d}$ 的映射相容。由此（以及引理
\ref{derham-lemma-base-change-Garel-upstairs} 中的唯一性）容易推出
$c^p_{Y'/X'}$ 与 $c^p_{Y/X}$ 相容。因此只需证明该断言。

\medskip\noindent
在证明的余下部分中，我们证明该断言。可以取一点 $y \in Y$，
并证明这些映射在 $y$ 的某个开邻域中相等。因此，可以用开邻域替换
$Y_1$、$Y_2$；这些开邻域分别包含 $y$ 在 $Y_1$ 和 $Y_2$ 中的像。
从而可以假设 $Y, X, Y_1, X_1, Y_2, X_2$ 都是仿射的，例如分别是环
$B, A, B_1, A_1, B_2, A_2$ 的谱。图示为
$$
\xymatrix{
B_1 \ar[r] & B & B_2 \ar[l] \\
A_1 \ar[u] \ar[r] & A \ar[u] & A_2 \ar[l] \ar[u]
}
$$
依假设，$B$ 的谱同时是 $A \otimes_{A_1} B_1$ 的谱和
$A \otimes_{A_2} B_2$ 的谱的仿射开集。进一步缩小后，可以假设存在元素
$g_i \in A \otimes_{A_i} B_i$，使我们的映射给出同构
$(A \otimes_{A_i} B_i)_{g_i} = B$；见
性质，引理 \ref{properties-lemma-standard-open-two-affines}。
令 $x_\alpha$、$y_\beta$ 为足够大的一组变量，使得可以选取满射
$$
A'_1 = A_1[x_\alpha] \to A
\quad\text{且}\quad
A'_2 = A_2[y_\beta] \to A
$$
，它们分别是 $A_1$-代数和 $A_2$-代数的满射。随后可以选取提升
$h_i \in A'_i \otimes_{A_i} B_i$，它提升
$g_i \in A \otimes_{A_i} B_i$；并考虑图
$$
\xymatrix{
(A'_1 \otimes_{A_1} B_1)_{h_1} \ar[r] & B &
(A'_2 \otimes_{A_1} B_2)_{h_2} \ar[l] \\
A'_1 \ar[u] \ar[r] & A \ar[u] & A'_2 \ar[l] \ar[u]
}
$$
依构造，两个方块都是余笛卡尔的。接着考虑环映射
$$
A' = A'_1 \times_A A'_2 \longrightarrow
B' =
(A'_1 \otimes_{A_1} B_1)_{h_1} \times_B
(A'_2 \otimes_{A_1} B_2)_{h_2}
$$
由更多代数，引理
\ref{more-algebra-lemma-module-over-fibre-product}，有
$A'_1 \otimes_{A'} B' = (A'_1 \otimes_{A_1} B_1)_{h_1}$ 和
$A'_2 \otimes_{A'} B' = (A'_2 \otimes_{A_2} B_2)_{h_2}$。特别地，态射
$Y' = \Spec(B') \to \Spec(A') = X'$ 的纤维是态射
$Y_i \to X_i$ 的纤维经基变换所得概形的开子概形，因而是有限完全局部交；
见判别式，引理 \ref{discriminant-lemma-syntomic-quasi-finite}。
由更多代数，引理
\ref{more-algebra-lemma-properties-algebras-over-fibre-product}，
环映射 $A' \to B'$ 是平坦且有限表示的。因此由判别式，引理
\ref{discriminant-lemma-syntomic-quasi-finite} 的 (3)，
可知 $Y'/X'$ 是 $\mathcal{C}$ 的对象。现在考虑交换图
$$
\xymatrix{
& & Y/X \ar[ld] \ar@/_2pc/[lldd]_{b_1/a_1} \ar[rd] \ar@/^2pc/[rrdd]^{b_2/a_2} \\
& Y'_1/X'_1 \ar[rd] \ar[ld] & & Y'_2/X'_2 \ar[ld] \ar[rd] \\
Y_1/X_1 & & Y'/X' & & Y_2/X_2
}
$$
其中 $Y'_i/X'_i$ 对应于
$A'_i \to (A'_i \otimes_{A_i} B_i)_{h_i}$。注意，
$Y'_i/X'_i$ 是 $\mathcal{C}_{nice}$ 的对象，因为它由一个无穷维仿射空间与
$Y_i/X_i$ 的积取开子概形而得；略去一个小细节。特别地，
$c^p_{Y_i/X_i}$ 经 $(b_i/a_i)$ 拉回后，与
$c^p_{Y'_i/X'_i}$ 经 $Y/X \to Y'_i/X'_i$ 拉回后所得的映射相同。
如果 $Y'/X'$ 是 $\mathcal{C}_{nice}$ 的对象，我们现在就已经完成了，
但实际几乎从不如此。因为在这种情况下，沿方块的两边拉回
$c^p_{Y'/X'}$ 就会给出所需结论。为了绕开
$Y'/X'$ 不属于 $\mathcal{C}_{nice}$ 的问题，注意上述论证表明：
在可能缩小所有概形 $X, Y, X'_1, Y'_1, X'_2, Y'_2, X', Y'$ 后，
可以找到某些 $n, d \geq 1$，并把图扩张如下：
$$
\xymatrix{
& Y/X \ar[ld] \ar[rd] \\
Y'_1/X'_1 \ar[rd] & & Y'_2/X'_2 \ar[ld] \\
& Y'/X' \ar[d] \\
& Y_{n, d}/X_{n, d}
}
$$
随后可以通过从 $c^p_{Y_{n, d}/X_{n, d}}$ 拉回来使用已经给出的论证。
证明完毕。
\end{proof}








\section{德拉姆复形上的迹映射}
\label{derham-section-trace}

\noindent
本节部分内容可参见 \cite{Garel}。设 $S$ 是一个概形，并设
$f : Y \to X$ 是 $S$ 上概形之间的有限局部自由态射。于是存在迹映射
$\text{Trace}_f : f_*\mathcal{O}_Y \to \mathcal{O}_X$；见判别式，第
\ref{discriminant-section-discriminant} 节。在这种情形下，
德拉姆复形上的迹映射是复形映射
$$
\Theta_{Y/X} : f_*\Omega^\bullet_{Y/S} \longrightarrow \Omega^\bullet_{X/S}
$$
使得 $\Theta_{Y/X}$ 等于 $\text{Trace}_f$（在次数 $0$ 上），并且满足
$$
\Theta_{Y/X}(\omega \wedge \eta) = \omega \wedge \Theta_{Y/X}(\eta)
$$
其中 $\omega$ 是 $\Omega^\bullet_{X/S}$ 的局部截面，$\eta$ 是
$f_*\Omega^\bullet_{Y/S}$ 的局部截面。我们尚不清楚这样的迹映射
$\Theta_{Y/X}$ 是否对每个有限局部自由态射 $Y \to X$ 都存在；如果你有反例或证明，
请发邮件至
\href{mailto:stacks.project@gmail.com}{stacks.project@gmail.com}
。

\begin{example}
\label{derham-example-no-trace}
下面给出一个德拉姆复形上不存在迹映射的例子。考虑
$\mathbf{C}$-代数 $B = \mathbf{C}[x, y]$，群 $G = \{\pm 1\}$
通过 $x \mapsto -x$ 和 $y \mapsto -y$ 作用于它。不变量
$A = B^G$ 构成 $\mathbf{C}$ 上的有限型正规整环，并由
$x^2, xy, y^2$ 生成。我们断言，对于包含 $A \subset B$，不存在合理的
迹映射
$\Omega_{B/\mathbf{C}} \to \Omega_{A/\mathbf{C}}$
作用在 $1$-形式上。事实上，考虑元素
$\omega = x \text{d} y \in \Omega_{B/\mathbf{C}}$.
由于 $\omega$ 在 $G$ 的作用下不变，如果存在“合理的”迹映射，
那么 $2\omega$ 应当属于
$\Omega_{A/\mathbf{C}} \to \Omega_{B/\mathbf{C}}$ 的像。但事实并非如此：
无法把 $2\omega$ 写成
$\text{d}(x^2)$、$\text{d}(xy)$ 和 $\text{d}(y^2)$ 的线性组合，
即使允许系数取自 $B$。
这个例子与 \cite{Zannier} 中的主定理相矛盾。
\end{example}

\begin{lemma}
\label{derham-lemma-Garel}
存在唯一的规则，它为概形的每个有限 syntomic 态射
$f : Y \to X$ 指定 $\mathcal{O}_X$-模映射
$$
\Theta^p_{Y/X} :
f_*\Omega^p_{Y/\mathbf{Z}}
\longrightarrow
\Omega^p_{X/\mathbf{Z}}
$$
并满足下列性质：
\begin{enumerate}
\item 与映射
$\Omega^p_{X/\mathbf{Z}} \otimes_{\mathcal{O}_X} f_*\mathcal{O}_Y
\to f_*\Omega^p_{Y/\mathbf{Z}}$ 的复合等于
$\text{id} \otimes \text{Trace}_f$
，其中 $\text{Trace}_f : f_*\mathcal{O}_Y \to \mathcal{O}_X$
是判别式，第 \ref{discriminant-section-discriminant} 节中的映射；
\item 该规则与基变换相容。
\end{enumerate}
\end{lemma}

\begin{proof}
首先，假设 $X$ 是局部诺特的。由引理
\ref{derham-lemma-Garel-upstairs}，有典范映射
$$
c^p_{Y/X} : \Omega_{Y/\mathbf{Z}}^p
\longrightarrow
f^*\Omega_{X/\mathbf{Z}}^p \otimes_{\mathcal{O}_Y} \det(\NL_{Y/X})
$$
由判别式，命题 \ref{discriminant-proposition-tate-map}，有典范同构
$$
c_{Y/X} : \det(\NL_{Y/X}) \to \omega_{Y/X}
$$
它把 $\delta(\NL_{Y/X})$ 映到 $\tau_{Y/X}$。合并这些映射得到
$$
c^p_{Y/X} \otimes c_{Y/X} :
\Omega_{Y/\mathbf{Z}}^p
\longrightarrow
f^*\Omega_{X/\mathbf{Z}}^p \otimes_{\mathcal{O}_Y} \omega_{Y/X}
$$
由判别式，第 \ref{discriminant-section-finite-morphisms} 节，
这等价于给出映射
$$
\Theta_{Y/X}^p :
f_*\Omega_{Y/\mathbf{Z}}^p
\longrightarrow
\Omega_{X/\mathbf{Z}}^p
$$
回忆，$c^p_{Y/X} \otimes c_{Y/X}$ 与 $\Theta_{Y/X}^p$
之间的关系使用求值映射
$f_*\omega_{Y/X} \to \mathcal{O}_X$
，它把 $\tau_{Y/X}$ 映到 $\text{Trace}_f(1)$；见判别式，第
\ref{discriminant-section-finite-morphisms} 节。因此性质 (1) 成立。
性质 (2) 对沿 $X' \to X$ 的基变换成立（其中 $X'$ 局部诺特），
因为 $c^p_{Y/X}$ 和 $c_{Y/X}$ 都与这种基变换相容。对有限 syntomic 态射
$f : Y \to X$，若 $X$ 局部诺特，我们仍用 $\Theta^p_{Y/X}$
表示刚刚得到的解。

\medskip\noindent
唯一性。假设有有限 syntomic 态射 $f: Y \to X$，使得
$X$ 在 $\Spec(\mathbf{Z})$ 上光滑，并且 $f$ 在 $X$ 的某个稠密开集上平展。
我们断言，在这种情形下，性质 (1) 唯一确定 $\Theta^p_{Y/X}$。
事实上，考虑映射
$$
\Omega^p_{X/\mathbf{Z}} \otimes_{\mathcal{O}_X} f_*\mathcal{O}_Y \to
f_*\Omega^p_{Y/\mathbf{Z}} \to
\Omega^p_{X/\mathbf{Z}}
$$
层 $\Omega^p_{X/\mathbf{Z}}$ 是无挠的（由所假设的光滑性），
所以只需检验 $\Theta^p_{Y/X}$ 在 $f$ 为平展的稠密开集上的限制被唯一确定；
也就是说，可以假设 $f$ 是平展的。然而，如果 $f$ 是平展的，那么
$f^*\Omega_{X/\mathbf{Z}} = \Omega_{Y/\mathbf{Z}}$
，因而所示等式中的第一个箭头是同构。由于复合已经确定，这就保证了唯一性。

\medskip\noindent
令 $f : Y \to X$ 是局部诺特概形之间的有限 syntomic 态射，并令
$x \in X$。由判别式，引理
\ref{discriminant-lemma-locally-comes-from-universal-finite}
，可以找到 $d \geq 1$ 以及交换图
$$
\xymatrix{
Y \ar[d] &
V \ar[d] \ar[l] \ar[r] &
V_d \ar[d] \\
X &
U \ar[l] \ar[r] &
U_d
}
$$
使得 $x \in U \subset X$ 是开集，$V = f^{-1}(U)$，并且
$V = U \times_{U_d} V_d$。因此 $\Theta^p_{Y/X}|_V$
是映射 $\Theta^p_{V_d/U_d}$ 的拉回。不过，由上面对唯一性的讨论以及
判别式，引理
\ref{discriminant-lemma-universal-finite-syntomic-smooth} 和
\ref{discriminant-lemma-universal-finite-syntomic-etale}
，映射 $\Theta^p_{V_d/U_d}$ 被要求 (1) 唯一确定。因此唯一性成立。

\medskip\noindent
至此，我们已经知道：对于所有有限 syntomic 态射 $Y \to X$，
只要 $X$ 局部诺特，解就存在且唯一。现在可以给出与引理
\ref{derham-lemma-Garel-upstairs} 的证明类似的论证，将结果推广到一般的 $X$。
不过，也可以直接使用绝对诺特逼近来完成证明。事实上，为构造
$\Theta^p_{Y/X}$，只需在 $X$ 上扎里斯基局部地构造它
（同时还需证明唯一性）。因此可以假设 $X$ 是仿射的（略去一个小细节）。
然后可将 $X = \lim_{i \in I} X_i$ 写成由有向集 $I$ 标号的诺特仿射概形的极限。
由代数，引理 \ref{algebra-lemma-colimit-category-fp-algebras}，
可以找到 $0 \in I$ 以及仿射概形之间的有限表示态射
$f_0 : Y_0 \to X_0$，它在 $X$ 上的基变换是 $Y \to X$。
增大 $0$ 后，可以假设 $Y_0 \to X_0$ 是有限 syntomic 态射；见代数，引理
\ref{algebra-lemma-colimit-lci} 和
\ref{algebra-lemma-colimit-finite}。对于 $i \geq 0$，基变换
$f_i : Y_i = Y_0 \times_{X_0} X_i \to X_i$ 也为有限 syntomic 态射。于是
$$
\Gamma(X, f_*\Omega^p_{Y/\mathbf{Z}}) =
\Gamma(Y, \Omega^p_{Y/\mathbf{Z}}) =
\colim_{i \geq 0} \Gamma(Y_i, \Omega^p_{Y_i/\mathbf{Z}}) =
\colim_{i \geq 0} \Gamma(X_i, f_{i, *}\Omega^p_{Y_i/\mathbf{Z}})
$$
因此我们可以（而且必须）把 $\Theta^p_{Y/X}$ 定义为映射
$\Theta^p_{Y_i/X_i}$ 的余极限。此映射与任意笛卡尔图
$$
\xymatrix{
Y' \ar[r] \ar[d] & Y \ar[d] \\
X' \ar[r] & X
}
$$
相容，其中 $X'$ 为仿射概形；因为由上述论证，我们已知这对诺特仿射概形成立
（略去一个小细节；提示：如果还写成
$X' = \lim_{j \in J} X'_j$，那么对每个 $i \in I$，存在 $j \in J$
以及态射 $X'_j \to X_i$，它与态射 $X' \to X$ 相容）。证明完毕。
\end{proof}

\begin{proposition}
\label{derham-proposition-Garel}
\begin{reference}
\cite{Garel}
\end{reference}
设 $f : Y \to X$ 是概形的有限 syntomic 态射。引理
\ref{derham-lemma-Garel} 中的映射 $\Theta^p_{Y/X}$ 定义复形映射
$$
\Theta_{Y/X} :
f_*\Omega^\bullet_{Y/\mathbf{Z}}
\longrightarrow
\Omega^\bullet_{X/\mathbf{Z}}
$$
并具有下列性质：
\begin{enumerate}
\item 在次数 $0$ 上得到
$\text{Trace}_f : f_*\mathcal{O}_Y \to \mathcal{O}_X$；见判别式，第
\ref{discriminant-section-discriminant} 节；
\item 有
$\Theta_{Y/X}(\omega \wedge \eta) = \omega \wedge \Theta_{Y/X}(\eta)$
，其中 $\omega$ 属于 $\Omega^\bullet_{X/\mathbf{Z}}$，$\eta$
属于 $f_*\Omega^\bullet_{Y/\mathbf{Z}}$；
\item 如果 $f$ 是基概形 $S$ 上的态射，那么
$\Theta_{Y/X}$ 诱导复形映射
$f_*\Omega^\bullet_{Y/S} \to \Omega^\bullet_{X/S}$.
\end{enumerate}
\end{proposition}

\begin{proof}
由判别式，引理
\ref{discriminant-lemma-locally-comes-from-universal-finite}
，对每个 $x \in X$，可以找到 $d \geq 1$ 以及交换图
$$
\xymatrix{
Y \ar[d] &
V \ar[d] \ar[l] \ar[r] &
V_d \ar[d] \ar[r] &
Y_d = \Spec(B_d) \ar[d] \\
X &
U \ar[l] \ar[r] &
U_d \ar[r] &
X_d = \Spec(A_d)
}
$$
使得 $x \in U \subset X$ 是仿射开集，$V = f^{-1}(U)$，并且
$V = U \times_{U_d} V_d$。写成 $U = \Spec(A)$ 和 $V = \Spec(B)$；
注意 $B = A \otimes_{A_d} B_d$，并回忆
$B_d = A_d e_1 \oplus \ldots \oplus A_d e_d$。假设给定
$a_1, \ldots, a_r \in A$ 和 $b_1, \ldots, b_s \in B$。
可以写成 $b_j = \sum a_{j, l} e_l$，其中 $a_{j, l} \in A$。
置 $N = r + sd$，并考虑分解
$$
\xymatrix{
V \ar[r] \ar[d] &
V' = \mathbf{A}^N \times V_d \ar[r] \ar[d] &
V_d \ar[d] \\
U \ar[r]&
U' = \mathbf{A}^N \times U_d \ar[r] &
U_d
}
$$
这里右下方的水平箭头由态射 $U \to U_d$（来自前面的图）和态射
$U \to \mathbf{A}^N$ 给出，后者由
$a_1, \ldots, a_r, a_{1, 1}, \ldots, a_{s, d}$ 给定。
于是函数 $a_1, \ldots, a_r$ 属于映射
$\Gamma(U', \mathcal{O}_{U'}) \to \Gamma(U, \mathcal{O}_U)$
的像，而函数 $b_1, \ldots, b_s$ 属于映射
$\Gamma(V', \mathcal{O}_{V'}) \to \Gamma(V, \mathcal{O}_V)$.
的像。这样可见，对这些群中任意有限多个元素\footnote{毕竟，
所有这些元素都是如下形式的元素的有限和：
$a_0 \text{d}a_1 \wedge \ldots \wedge \text{d}a_i$，其中
$a_0, \ldots, a_i \in A$，或者如下形式的元素的有限和：
$b_0 \text{d}b_1 \wedge \ldots \wedge \text{d}b_j$，其中
$b_0, \ldots, b_j \in B$。}：
$$
\Gamma(V, \Omega^i_{Y/\mathbf{Z}}),\quad i = 0, 1, 2, \ldots
\quad\text{且}\quad
\Gamma(U, \Omega^j_{X/\mathbf{Z}}),\quad j = 0, 1, 2, \ldots
$$
可以找到如上的分解 $V \to V' \to V_d$ 和
$U \to U' \to U_d$，其中 $V' = \mathbf{A}^N \times V_d$ 且
$U' = \mathbf{A}^N \times U_d$，使得这些截面都是下列群中截面的拉回：
$$
\Gamma(V', \Omega^i_{V'/\mathbf{Z}}),\quad i = 0, 1, 2, \ldots
\quad\text{且}\quad
\Gamma(U', \Omega^j_{U'/\mathbf{Z}}),\quad j = 0, 1, 2, \ldots
$$
其结果是，为检验
$\text{d} \circ \Theta_{Y/X} = \Theta_{Y/X} \circ \text{d}$
，只需检验这对 $\Theta_{V'/U'}$ 成立。引理的性质 (2) 也同样如此。

\medskip\noindent
由判别式，引理
\ref{discriminant-lemma-universal-finite-syntomic-smooth} 和
\ref{discriminant-lemma-universal-finite-syntomic-etale}
，概形 $U_d$ 是光滑的，且态射 $V_d \to U_d$ 在 $U_d$
的某个稠密开集上平展。因此态射 $V' \to U'$ 也具有同样的性质。
由于 $\Omega_{U'/\mathbf{Z}}$ 是局部自由的，故
$\Omega^p_{U'/\mathbf{Z}}$ 是无挠的；所以只需把所需关系限制到
$V'$ 为有限平展的开集后加以检验。随后可以在满平展基变换之后检验这些关系。
因而可以分裂该有限平展覆盖，并假设所考察的态射形如
$$
\coprod\nolimits_{i = 1, \ldots, d} W \longrightarrow W
$$
，其中 $W$ 在 $\mathbf{Z}$ 上光滑。在这种情形下，我们构造的任何局部性质
（只要确实成立）都可直接检验。这就完成了 (1) 和 (2) 的证明。

\medskip\noindent
最后注意到，(3) 由 (2) 推出，因为 $\Omega_{Y/S}$ 是
$\Omega_{Y/\mathbf{Z}}$ 关于如下子模的商：该子模由
$\Omega_{S/\mathbf{Z}}$ 的局部截面的拉回生成。
\end{proof}

\begin{example}
\label{derham-example-Garel}
设 $A$ 是一个环，并令
$f = x^d + \sum_{0 \leq i < d} a_{d - i} x^i \in A[x]$。
令 $B = A[x]/(f)$。由命题 \ref{derham-proposition-Garel}，有复形态射
$$
\Theta_{B/A} : \Omega^\bullet_B \longrightarrow \Omega^\bullet_A
$$
特别地，如果 $t \in B$ 表示 $x \in A[x]$ 的像，就可以考虑元素
$$
\Theta_{B/A}(t^i\text{d}t) \in \Omega^1_A,\quad i = 0, \ldots, d - 1
$$
这些元素是什么？按照命题 \ref{derham-proposition-Garel} 的证明中使用的同一原理，
只需在泛情形下计算，即 $A = \mathbf{Z}[a_1, \ldots, a_d]$ 时；
甚至当 $A$ 被分式域 $\mathbf{Q}(a_1, \ldots, a_d)$ 代替时。
形式地写成
$$
f = \prod\nolimits_{i = 1, \ldots, d} (x - \alpha_i)
$$
可见在 $\mathbf{Q}(\alpha_1, \ldots, \alpha_d)$ 上，代数 $B$ 分裂为：
$$
\mathbf{Q}(a_0, \ldots, a_{d - 1})[x]/(f) 
\longrightarrow
\prod\nolimits_{i = 1, \ldots, d} \mathbf{Q}(\alpha_1, \ldots, \alpha_d),
\quad
t \longmapsto (\alpha_1, \ldots, \alpha_d)
$$
因此，例如
$$
\Theta(\text{d}t) = \sum \text{d} \alpha_i = - \text{d}a_1
$$
接下来有
$$
\Theta(t\text{d}t) = \sum \alpha_i \text{d}\alpha_i =
a_1 \text{d} a_1 - \text{d}a_2
$$
接下来有
$$
\Theta(t^2\text{d}t) = \sum \alpha_i^2 \text{d}\alpha_i =
- a_1^2 \text{d} a_1 + a_1 \text{d}a_2 + a_2 \text{d}a_1 - \text{d}a_3
$$
（不计可能的计算错误），依此类推。这提示：如果 $f(x) = x^d - a$，那么
$$
\Theta_{B/A}(t^i\text{d}t) =
\left\{
\begin{matrix}
0 & \text{若} & i = 0, \ldots, d - 2 \\
\text{d}a & \text{若} & i = d - 1
\end{matrix}
\right.
$$
在 $\Omega_A$ 中成立。这确实如此，因为在这个特例中，可以对扩张
$\mathbf{Q}(a)[x]/(x^d - a)$ 进行计算来直接验证。
\end{example}

\begin{lemma}
\label{derham-lemma-Garel-map-frobenius-smooth-char-p}
设 $p$ 是素数。设 $X \to S$ 是相对维数为 $d$ 的光滑态射，
其中这些概形的特征为 $p$。相对 Frobenius 态射
$F_{X/S} : X \to X^{(p)}$（属于 $X/S$；簇，定义
\ref{varieties-definition-relative-frobenius}）是有限 syntomic 态射，
且相应的映射
$$
\Theta_{X/X^{(p)}} :
F_{X/S, *}\Omega^\bullet_{X/S} \to \Omega^\bullet_{X^{(p)}/S}
$$
除次数 $d$ 外在所有次数上都为零，而在该次数上定义一个满射。
\end{lemma}

\begin{proof}
由簇，引理 \ref{varieties-lemma-relative-frobenius-finite}，
$F_{X/S}$ 是有限态射。为了证明 $F_{X/S}$ 平坦，只需证明所有纤维间的态射
$F_{X/S, s} : X_s \to X^{(p)}_s$ 对每个 $s \in S$ 都平坦；
见更多态射，定理 \ref{more-morphisms-theorem-criterion-flatness-fibre}。
$X_s \to X^{(p)}_s$ 的平坦性由代数，引理
\ref{algebra-lemma-CM-over-regular-flat} 得出（并使用已经证明的有限性）。
由更多态射，引理 \ref{more-morphisms-lemma-lci-permanence}，
态射 $F_{X/S}$ 是局部完全交态射。因此 $F_{X/S}$ 是有限 syntomic 态射
（见更多态射，引理 \ref{more-morphisms-lemma-flat-lci}）。

\medskip\noindent
对每一点 $x \in X$，可以选取交换图
$$
\xymatrix{
X \ar[d] & U \ar[l] \ar[d]_\pi \\
S & \mathbf{A}^d_S \ar[l]
}
$$
其中 $\pi$ 是平展的，且 $x \in U$ 是 $X$ 中的开集；见态射，引理
\ref{morphisms-lemma-smooth-etale-over-affine-space}。注意，
$\mathbf{A}^d_S \to \mathbf{A}^d_S$, $(x_1, \ldots, x_d) \mapsto
(x_1^p, \ldots, x_d^p)$ 是 $\mathcal{A}^d_S$ 在 $S$ 上的相对 Frobenius。
交换图
$$
\xymatrix{
U \ar[d]_\pi \ar[r]_{F_{X/S}} & U^{(p)} \ar[d]^{\pi^{(p)}} \\
\mathbf{A}^d_S \ar[r]^{x_i \mapsto x_i^p} & \mathbf{A}^d_S
}
$$
来自簇，引理 \ref{varieties-lemma-relative-frobenius-endomorphism-identity}；
对 $\pi : U \to \mathbf{A}^d_S$，由平展态射，引理
\ref{etale-lemma-relative-frobenius-etale}，该图是笛卡尔图。
由于 $\Theta$ 的构造与基变换相容，且
$\Omega_{U/S} = \pi^*\Omega_{\mathbf{A}^d_S/S}$
（引理 \ref{derham-lemma-etale}），所以只需对 $\mathbf{A}^d_S$ 证明本引理。

\medskip\noindent
设 $A$ 是特征 $p$ 的环。考虑唯一的 $A$-代数同态
$A[y_1, \ldots, y_d] \to A[x_1, \ldots, x_d]$，它把 $y_i$ 映到
$x_i^p$。上述论证把问题化为计算映射
$$
\Theta^i : \Omega^i_{A[x_1, \ldots, x_d]/A} \to
\Omega^i_{A[y_1, \ldots, y_d]/A}
$$
我们建议读者亲自完成这一情形下的计算。如同例 \ref{derham-example-Garel}，
可以把它化为在泛情形下计算 $\Theta^i$ 的公式
$$
\mathbf{Z}[y_1, \ldots, y_d] \to \mathbf{Z}[x_1, \ldots, x_d],\quad
y_i \mapsto x_i^p
$$
进而，可以在复数情形下计算 $\Theta^i$ 的公式，即对 $\mathbf{C}$-代数映射
$$
\mathbf{C}[y_1, \ldots, y_d] \to \mathbf{C}[x_1, \ldots, x_d],\quad
y_i \mapsto x_i^p
$$
甚至可以逆转 $x_1, \ldots, x_d$ 和 $y_1, \ldots, y_d$。
在这种情形下，有 $\text{d}x_i = p^{-1} x_i^{- p + 1}\text{d}y_i$。
因此可见
\begin{align*}
\Theta^i(
x_1^{e_1} \ldots x_d^{e_d} \text{d}x_1 \wedge \ldots \wedge \text{d}x_i)
& =
p^{-i} \Theta^i(
x_1^{e_1 - p + 1} \ldots x_i^{e_i - p + 1} x_{i + 1}^{e_{i + 1}} \ldots
x_d^{e_d} \text{d}y_1 \wedge \ldots \wedge \text{d}y_i ) \\
& =
p^{-i} \text{Trace}(x_1^{e_1 - p + 1} \ldots x_i^{e_i - p + 1}
x_{i + 1}^{e_{i + 1}} \ldots x_d^{e_d})
\text{d}y_1 \wedge \ldots \wedge \text{d}y_i
\end{align*}
，这里使用了 $\Theta^i$ 的性质。初等计算表明，上式中的迹为零，除非
$e_1, \ldots, e_i$ 都同余于 $-1$（模 $p$），且
$e_{i + 1}, \ldots, e_d$ 都能被 $p$ 整除。此外，在这种情形下得到
$$
p^{d - i} y_1^{(e_1 - p + 1)/p} \ldots y_i^{(e_i - p + 1)/p}
y_{i + 1}^{e_{i + 1}/p} \ldots y_d^{e_d/p}
\text{d}y_1 \wedge \ldots \wedge \text{d}y_i
$$
由此可知，在特征 $p$ 中，除非 $d = i$，所得结果为零；
而在这种情况下，可以得到每个可能的 $d$-形式。
\end{proof}








\section{庞加莱对偶}
\label{derham-section-poincare-duality}

\noindent
本节证明域上的光滑固有概形的德拉姆上同调所满足的庞加莱对偶。
我们先说明霍奇上同调的情形。

\begin{lemma}
\label{derham-lemma-duality-hodge}
设 $k$ 是一个域。设 $X$ 是 $k$ 上非空、光滑、固有且等维的概形，
其维数为 $d$。存在 $k$-线性映射
$$
t : H^d(X, \Omega^d_{X/k}) \longrightarrow k
$$
；它在先与 $H^0(X, \mathcal{O}_X)$ 中某个单位的乘法复合的意义下唯一，
并具有如下性质：对所有 $p, q$，配对
$$
H^q(X, \Omega^p_{X/k}) \times H^{d - q}(X, \Omega^{d - p}_{X/k})
\longrightarrow
k, \quad
(\xi, \xi') \longmapsto t(\xi \cup \xi')
$$
是完美配对。
\end{lemma}

\begin{proof}
由概形对偶，引理 \ref{duality-lemma-duality-proper-over-field}，有
$\omega_X^\bullet = \Omega^d_{X/k}[d]$。由于 $\Omega_{X/k}$
是秩为 $d$ 的局部自由模（态射，引理
\ref{morphisms-lemma-smooth-omega-finite-locally-free}），有
$$
\Omega^d_{X/k} \otimes_{\mathcal{O}_X} (\Omega^p_{X/k})^\vee
\cong
\Omega^{d - p}_{X/k}
$$
因此，由概形对偶，引理
\ref{duality-lemma-duality-proper-over-field-perfect}，得到使该陈述成立的
$k$-线性映射 $t : H^d(X, \Omega^d_{X/k}) \to k$。特别地，配对
$H^0(X, \mathcal{O}_X) \times H^d(X, \Omega^d_{X/k}) \to k$
是完美的。这意味着任意 $k$-线性映射
$t' : H^d(X, \Omega^d_{X/k}) \to k$ 都形如
$\xi \mapsto t(g\xi)$，其中 $g \in H^0(X, \mathcal{O}_X)$。
当然，为使 $t'$ 仍在 $H^0(X, \mathcal{O}_X)$ 与
$H^d(X, \Omega^d_{X/k})$ 之间给出对偶，$g$ 必须是单位。
把配对 $\langle -, - \rangle_{p, q}$ 定义为用 $t$ 构造的配对，
把配对 $\langle -, - \rangle'_{p, q}$ 定义为用 $t'$ 构造的配对。显然有
$$
\langle \xi, \xi' \rangle'_{p, q} =
\langle g\xi, \xi' \rangle_{p, q}
$$
，其中 $\xi \in H^q(X, \Omega^p_{X/k})$ 且
$\xi' \in H^{d - q}(X, \Omega^{d - p}_{X/k})$。由于 $g$ 是单位，
即为可逆元，所以用 $t'$ 代替 $t$，对所有 $p, q$
仍得到完美配对。
\end{proof}

\begin{lemma}
\label{derham-lemma-bottom-part-degenerates}
设 $k$ 是一个域，设 $X$ 是 $k$ 上光滑固有概形。映射
$$
\text{d} : H^0(X, \mathcal{O}_X) \to H^0(X, \Omega^1_{X/k})
$$
为零。
\end{lemma}

\begin{proof}
由于 $X$ 在 $k$ 上光滑，它在 $k$ 上几何既约；见簇，引理
\ref{varieties-lemma-smooth-geometrically-normal}。因此
$H^0(X, \mathcal{O}_X) = \prod k_i$ 是有限多个有限可分域扩张
$k_i/k$ 的乘积；见簇，引理
\ref{varieties-lemma-proper-geometrically-reduced-global-sections}。由此
$\Omega_{H^0(X, \mathcal{O}_X)/k} = \prod \Omega_{k_i/k} = 0$
（例如见代数，引理
\ref{algebra-lemma-characterize-separable-algebraic-field-extensions}）。
由于引理中的映射分解为
$$
H^0(X, \mathcal{O}_X) \to
\Omega_{H^0(X, \mathcal{O}_X)/k} \to
H^0(X, \Omega_{X/k})
$$
（由德拉姆复形的函子性；见第 \ref{derham-section-de-rham-complex} 节），结论成立。
\end{proof}

\begin{lemma}
\label{derham-lemma-top-part-degenerates}
设 $k$ 是一个域。设 $X$ 是 $k$ 上维数为 $d$ 的等维光滑固有概形。映射
$$
\text{d} : H^d(X, \Omega^{d - 1}_{X/k}) \to H^d(X, \Omega^d_{X/k})
$$
为零。
\end{lemma}

\begin{proof}
很容易猜想这可由引理 \ref{derham-lemma-bottom-part-degenerates} 与
\ref{derham-lemma-duality-hodge} 结合推出。然而这并不可行，因为映射
$\mathcal{O}_X \to \Omega^1_{X/k}$ 和
$\Omega^{d - 1}_{X/k} \to \Omega^d_{X/k}$ 不是 $\mathcal{O}_X$-线性的，
所以不能利用概形对偶，备注
\ref{duality-remark-coherent-duality-proper-over-field}
中讨论的函子性来断定引理 \ref{derham-lemma-bottom-part-degenerates}
中的映射与本引理中的映射互为对偶。

\medskip\noindent
可以用 $X$ 的一个连通分支代替 $X$，因此可以假设 $X$ 不可约。由簇，引理
\ref{varieties-lemma-smooth-geometrically-normal} 和
\ref{varieties-lemma-proper-geometrically-reduced-global-sections}
可知 $k' = H^0(X, \mathcal{O}_X)$ 是有限可分扩张 $k'/k$。
由于 $\Omega_{k'/k} = 0$（例如见代数，引理
\ref{algebra-lemma-characterize-separable-algebraic-field-extensions}），
有 $\Omega_{X/k} = \Omega_{X/k'}$（见态射，引理
\ref{morphisms-lemma-triangle-differentials}）。
因此可以把 $k$ 替换为 $k'$，并假设 $H^0(X, \mathcal{O}_X) = k$。

\medskip\noindent
假设 $H^0(X, \mathcal{O}_X) = k$。由引理 \ref{derham-lemma-duality-hodge}，
有 $\dim H^d(X, \Omega^d_{X/k}) = 1$。
先假设 $k$ 的特征是素数 $p$。把相对 Frobenius 态射
$F_{X/k} : X \to X^{(p)}$（属于 $X$ 在 $k$ 上）记为上述符号；请记住引理
\ref{derham-lemma-Garel-map-frobenius-smooth-char-p} 中关于该态射证明的事实。
考虑交换图
$$
\xymatrix{
H^d(X, \Omega^{d - 1}_{X/k}) \ar[d] \ar[r] &
H^d(X^{(p)}, F_{X/k, *}\Omega^{d - 1}_{X/k}) \ar[d] \ar[r]_{\Theta^{d - 1}} &
H^d(X^{(p)}, \Omega^{d - 1}_{X^{(p)}/k}) \ar[d] \\
H^d(X, \Omega^d_{X/k}) \ar[r] &
H^d(X^{(p)}, F_{X/k, *}\Omega^d_{X/k}) \ar[r]^{\Theta^d} &
H^d(X^{(p)}, \Omega^d_{X^{(p)}/k})
}
$$
由于 $F_{X/k}$ 是有限的，左边两个水平箭头是同构；见概形上同调，引理
\ref{coherent-lemma-relative-affine-cohomology}。右方块交换，因为
$\Theta_{X^{(p)}/X}$ 是复形态射，且 $\Theta^{d - 1}$ 为零。
因此只需证明 $\Theta^d$ 非零（由上面的讨论，映射 $\Theta^d$
的源维数为 $1$）。然而我们知道
$$
\Theta^d : F_{X/k, *}\Omega^d_{X/k} \to \Omega^d_{X^{(p)}/k}
$$
是满射，因而应用右正合函子 $H^d(X^{(p)}, -)$ 后仍为满射
（右正合性来自次数超过 $d$ 的上同调消失；见上同调，命题
\ref{cohomology-proposition-vanishing-Noetherian}）。
最后，$H^d(X^{(d)}, \Omega^d_{X^{(d)}/k})$ 非零，例如因为它对偶于
$H^0(X^{(d)}, \mathcal{O}_{X^{(p)}})$；这是由引理
\ref{derham-lemma-duality-hodge} 应用于 $X^{(p)}$（在 $k$ 上）得出的。
这就完成了该情形的证明。

\medskip\noindent
最后，假设 $k$ 的特征为 $0$。可以把 $k$ 写成其有限型
$\mathbf{Z}$-子代数 $R$ 的滤过余极限。对其中某一个，可以找到概形的笛卡尔图
$$
\xymatrix{
X \ar[d] \ar[r] & Y \ar[d] \\
\Spec(k) \ar[r] & \Spec(R)
}
$$
使得 $Y \to \Spec(R)$ 是相对维数为 $d$ 的光滑固有态射。
见极限，引理 \ref{limits-lemma-descend-finite-presentation}、
\ref{limits-lemma-descend-smooth}、\ref{limits-lemma-descend-dimension-d} 和
\ref{limits-lemma-eventually-proper}。模
$M^{i, j} = H^j(Y, \Omega^i_{Y/R})$ 是有限 $R$-模；见概形上同调，引理
\ref{coherent-lemma-proper-over-affine-cohomology-finite}。
因此，用 $R$ 的一个局部化代替它之后，可以假设所有这些模都是有限自由的。由
$M^{i, j} \otimes_R k = H^j(X, \Omega^i_{X/k})$
；这是平坦基变换（概形上同调，引理
\ref{coherent-lemma-flat-base-change-cohomology}）的结果。
因此只需证明 $M^{d - 1, d} \to M^{d, d}$ 为零。这是整环上的有限自由模之间的映射，
所以只需找到稠密的素理想集合 $\mathfrak p \subset R$，使得与
$\kappa(\mathfrak p)$ 张量后所得映射为零。由于 $R$ 在 $\mathbf{Z}$
上有限型，可以取剩余域具有正特征的素理想 $\mathfrak p$
所组成的集合（略去细节）。注意，
$$
M^{d - 1, d} \otimes_R \kappa(\mathfrak p) =
H^d(Y_{\kappa(\mathfrak p)},
\Omega^{d - 1}_{Y_{\kappa(\mathfrak p)}/\kappa(\mathfrak p)})
$$
；例如这由极限，引理
\ref{limits-lemma-higher-direct-images-zero-above-dimension-fibre} 得出。
对 $M^{d, d}$ 也同样如此。因此
$M^{d - 1, d} \otimes_R \kappa(\mathfrak p) \to
M^{d, d} \otimes_R \kappa(\mathfrak p)$
由上面已经处理的正特征情形，此映射为零。
\end{proof}

\begin{proposition}
\label{derham-proposition-poincare-duality}
设 $k$ 是一个域。设 $X$ 是 $k$ 上非空、光滑、固有且等维的概形，
其维数为 $d$。存在 $k$-线性映射
$$
t : H^{2d}_{dR}(X/k) \longrightarrow k
$$
；它在先与 $H^0(X, \mathcal{O}_X)$ 中某个单位的乘法复合的意义下唯一，
并具有如下性质：对所有 $i$，配对
$$
H^i_{dR}(X/k) \times H_{dR}^{2d - i}(X/k)
\longrightarrow
k, \quad
(\xi, \xi') \longmapsto t(\xi \cup \xi')
$$
是完美配对。
\end{proposition}

\begin{proof}
由霍奇到德拉姆谱序列（第 \ref{derham-section-hodge-to-de-rham} 节）、
$\Omega^i_{X/k}$ 对 $i > d$ 的消失、上同调，命题
\ref{cohomology-proposition-vanishing-Noetherian} 中的消失，以及引理
\ref{derham-lemma-bottom-part-degenerates} 和
\ref{derham-lemma-top-part-degenerates} 的结果，可知
$H^0_{dR}(X/k) = H^0(X, \mathcal{O}_X)$，且
$H^d(X, \Omega^d_{X/k}) = H_{dR}^{2d}(X/k)$。
更确切地说，这些等同来自复形映射
$$
\Omega^\bullet_{X/k} \to \mathcal{O}_X[0]
\quad\text{且}\quad
\Omega^d_{X/k}[-d] \to \Omega^\bullet_{X/k}
$$
选择 $t : H_{dR}^{2d}(X/k) \to k$，使它经此等同对应于引理
\ref{derham-lemma-duality-hodge} 中那样的一个 $t$。于是无论如何，
该引理中所示的配对在 $i = 0$ 时是完美的。

\medskip\noindent
把 $\underline{k}$ 记为取值 $k$ 的常值层（在 $X$ 上）。
简记 $\Omega^\bullet = \Omega^\bullet_{X/k}$。
考虑映射 (\ref{derham-equation-wedge})，在当前情形下它写成
$$
\wedge :
\text{Tot}(\Omega^\bullet \otimes_{\underline{k}} \Omega^\bullet)
\longrightarrow
\Omega^\bullet
$$
对每个整数 $p = 0, 1, \ldots, d$，由于次数的原因，此映射在子复形
$\text{Tot}(\sigma_{> p} \Omega^\bullet \otimes_{\underline{k}}
\sigma_{\geq d - p} \Omega^\bullet)$ 上为零。因此，$\wedge$ 在子复形
$\text{Tot}(\Omega^\bullet \otimes_{\underline{k}}
\sigma_{\geq d - p}\Omega^\bullet)$ 上的限制通过复形映射
$$
\gamma_p :
\text{Tot}(\sigma_{\leq p} \Omega^\bullet \otimes_{\underline{k}}
\sigma_{\geq d - p} \Omega^\bullet)
\longrightarrow
\Omega^\bullet
$$
使用第 \ref{derham-section-cup-product} 节中的相同步骤，得到杯积
$$
H^i(X, \sigma_{\leq p} \Omega^\bullet) \times
H^{2d - i}(X, \sigma_{\geq d - p}\Omega^\bullet)
\longrightarrow
H_{dR}^{2d}(X, \Omega^\bullet)
$$
我们将对 $p$ 归纳证明：这些杯积经 $t$ 在
$H^i(X, \sigma_{\leq p} \Omega^\bullet)$ 与
$H^{2d - i}(X, \sigma_{\geq d - p}\Omega^\bullet)$ 之间诱导完美配对。
当 $p = d$ 时，这正是本命题的断言。

\medskip\noindent
基础情形是 $p = 0$。此时得到的恰是引理
\ref{derham-lemma-duality-hodge} 中 $H^i(X, \mathcal{O}_X)$ 与
$H^{d - i}(X, \Omega^d)$ 之间的配对，因而结论成立。

\medskip\noindent
归纳步骤。假设已知结论对 $p$ 成立。考虑可区别三角
$$
\Omega^{p + 1}[-p - 1] \to
\sigma_{\leq p + 1}\Omega^\bullet \to
\sigma_{\leq p}\Omega^\bullet \to
\Omega^{p + 1}[-p]
$$
以及可区别三角
$$
\sigma_{\geq d - p}\Omega^\bullet \to
\sigma_{\geq d - p - 1}\Omega^\bullet \to
\Omega^{d - p - 1}[-d + p + 1] \to
(\sigma_{\geq d - p}\Omega^\bullet)[1]
$$
注意，两者都是 $\underline{k}$-模层复形的同伦范畴中的可区别三角；
特别地，映射 $\sigma_{\leq p}\Omega^\bullet \to \Omega^{p + 1}[-p]$ 和
$\Omega^{d - p - 1}[-d + p + 1] \to (\sigma_{\geq d - p}\Omega^\bullet)[1]$
由实际的复形映射给出，即分别使用微分
$\Omega^p \to \Omega^{p + 1}$ 和
$\Omega^{d - p - 1} \to \Omega^{d - p}$。
考虑与这些可区别三角相伴的上同调长正合列
$$
\xymatrix{
H^{i - 1}(X, \sigma_{\leq p}\Omega^\bullet) \ar[d]_a \\
H^i(X, \Omega^{p + 1}[-p - 1]) \ar[d]_b \\
H^i(X, \sigma_{\leq p + 1}\Omega^\bullet) \ar[d]_c \\
H^i(X, \sigma_{\leq p}\Omega^\bullet) \ar[d]_d \\
H^{i + 1}(X, \Omega^{p + 1}[-p - 1])
}
\quad\quad
\xymatrix{
H^{2d - i  + 1}(X, \sigma_{\geq d - p}\Omega^\bullet) \\
H^{2d - i}(X, \Omega^{d - p - 1}[-d + p + 1]) \ar[u]_{a'} \\
H^{2d - i}(X, \sigma_{\geq d - p - 1}\Omega^\bullet) \ar[u]_{b'} \\
H^{2d - i}(X, \sigma_{\geq d - p}\Omega^\bullet) \ar[u]_{c'} \\
H^{2d - i - 1}(X, \Omega^{d - p - 1}[-d + p + 1]) \ar[u]_{d'}
}
$$
由归纳假设和引理 \ref{derham-lemma-duality-hodge}，我们知道上面在第一、第二、
第四和第五行的 $k$-向量空间之间构造的配对是完美的。由 $5$-引理，
为了证明中间行上同调群之间的配对是完美的，只需证明各对
$(a, a')$、$(b, b')$、$(c, c')$ 和 $(d, d')$
与给定配对相容（见下文）。

\medskip\noindent
先对 $(c, c')$ 证明这一点。这里只需注意到存在交换图
$$
\xymatrix{
\text{Tot}(\sigma_{\leq p} \Omega^\bullet \otimes_{\underline{k}}
\sigma_{\geq d - p} \Omega^\bullet) \ar[d]_{\gamma_p} &
\text{Tot}(\sigma_{\leq p + 1} \Omega^\bullet \otimes_{\underline{k}}
\sigma_{\geq d - p} \Omega^\bullet) \ar[l] \ar[d] \\
\Omega^\bullet &
\text{Tot}(\sigma_{\leq p + 1} \Omega^\bullet \otimes_{\underline{k}}
\sigma_{\geq d - p - 1} \Omega^\bullet) \ar[l]_-{\gamma_{p + 1}}
}
$$
因此，如果有 $\alpha \in H^i(X, \sigma_{\leq p + 1}\Omega^\bullet)$ 和
$\beta \in H^{2d - i}(X, \sigma_{\geq d - p}\Omega^\bullet)$，那么由杯积的函子性得到
$\gamma_p(\alpha \cup c'(\beta)) = \gamma_{p + 1}(c(\alpha) \cup \beta)$
。

\medskip\noindent
类似地，对 $(b, b')$ 使用交换图
$$
\xymatrix{
\text{Tot}(\sigma_{\leq p + 1} \Omega^\bullet \otimes_{\underline{k}}
\sigma_{\geq d - p - 1} \Omega^\bullet) \ar[d]_{\gamma_{p + 1}} &
\text{Tot}(\Omega^{p + 1}[-p - 1] \otimes_{\underline{k}}
\sigma_{\geq d - p - 1} \Omega^\bullet) \ar[l] \ar[d] \\
\Omega^\bullet &
\Omega^{p + 1}[-p - 1]
\otimes_{\underline{k}}
\Omega^{d - p - 1}[-d + p + 1] \ar[l]_-\wedge
}
$$
并作同样的论证。

\medskip\noindent
对 $(d, d')$ 使用交换图
$$
\xymatrix{
\Omega^{p + 1}[-p] \otimes_{\underline{k}}
\Omega^{d - p - 1}[-d + p] \ar[d] &
\text{Tot}(\sigma_{\leq p}\Omega^\bullet \otimes_{\underline{k}}
\Omega^{d - p - 1}[-d + p]) \ar[l] \ar[d] \\
\Omega^\bullet &
\text{Tot}(\sigma_{\leq p}\Omega^\bullet \otimes_{\underline{k}}
\sigma_{\geq d - p}\Omega^\bullet) \ar[l]
}
$$
并考察下列群中的上同调类：
$H^i(X, \sigma_{\leq p}\Omega^\bullet)$ 和
$H^{2d - i}(X, \Omega^{d - p - 1}[-d + p])$.
把 $i$ 改为 $i - 1$，就得到对 $(a, a')$ 的结论，
从而完成了这些配对是完美配对的证明。

\medskip\noindent
我们略去如下论证：$t$ 在先与 $H^0(X, \mathcal{O}_X)$
中某个单位的乘法复合的意义下唯一。
\end{proof}







\section{陈类}
\label{derham-section-chern-classes}

\noindent
目前已经证明的结果足以利用 Weil 上同调理论，第
\ref{weil-section-chern} 节中的讨论，在德拉姆上同调中构造陈类。

\begin{lemma}
\label{derham-lemma-chern-classes}
存在唯一的规则，它为每个拟紧拟分离概形 $X$ 指定全陈类
$$
c^{dR} :
K_0(\textit{Vect}(X))
\longrightarrow
\prod\nolimits_{i \geq 0} H^{2i}_{dR}(X/\mathbf{Z})
$$
并具有下列性质：
\begin{enumerate}
\item 有 $c^{dR}(\alpha + \beta) = c^{dR}(\alpha) c^{dR}(\beta)$，
其中 $\alpha, \beta \in K_0(\textit{Vect}(X))$；
\item 如果 $f : X \to X'$ 是拟紧拟分离概形之间的态射，那么
$c^{dR}(f^*\alpha) = f^*c^{dR}(\alpha)$；
\item 给定 $\mathcal{L} \in \Pic(X)$，有
$c^{dR}([\mathcal{L}]) = 1 + c_1^{dR}(\mathcal{L})$
\end{enumerate}
\end{lemma}

\noindent
该构造很容易推广到所有概形，但为此需要稍微加强 Weil 上同调理论，第
\ref{weil-section-chern} 节中的讨论。

\begin{proof}
应用 Weil 上同调理论，命题 \ref{weil-proposition-chern-class} 即可得到。

\medskip\noindent
令 $\mathcal{C}$ 为所有拟紧拟分离概形组成的范畴。它显然满足
Weil 上同调理论，第 \ref{weil-section-chern} 节的条件
(1)、(2) 以及 (3) 的 (a)、(b)、(c)。

\medskip\noindent
取逆变函子 $A$（从 $\mathcal{C}$ 到分次代数范畴），它把 $X$ 映到
$A(X) = \bigoplus_{i \geq 0} H_{dR}^{2i}(X/\mathbf{Z})$
，并赋予后者杯积。函子性在第 \ref{derham-section-de-rham-cohomology} 节讨论，
杯积在第 \ref{derham-section-cup-product} 节讨论。对于加性映射 $c_1^A$，
取第 \ref{derham-section-first-chern-class} 节构造的 $c_1^{dR}$。

\medskip\noindent
事实上，由引理 \ref{derham-lemma-cup-product-graded-commutative}
得到交换代数，这说明 $A$ 满足公理 (1)。

\medskip\noindent
为了检验 $A$ 的公理 (2)，只需检验
$H^*_{dR}(X \coprod Y/\mathbf{Z}) =  H^*_{dR}(X/\mathbf{Z}) \times
H^*_{dR}(Y/\mathbf{Z})$.
这是因为德拉姆上同调是取微分分次代数层（在扎里斯基拓扑中）的上同调构造的。

\medskip\noindent
对 $A$ 而言，公理 (3) 正是说取可逆模的第一陈类与拉回相容。
这由更一般的引理 \ref{derham-lemma-pullback-c1} 得出。

\medskip\noindent
对 $A$ 而言，公理 (4) 是我们在命题
\ref{derham-proposition-projective-space-bundle-formula} 中证明的射影空间丛公式。

\medskip\noindent
公理 (5)。设 $X$ 是拟紧拟分离概形，并设
$\mathcal{E} \to \mathcal{F}$ 是有限局部自由 $\mathcal{O}_X$-模之间的满射，
两者的秩分别为 $r + 1$ 和 $r$。把相应的包含态射记为
$i : P' = \mathbf{P}(\mathcal{F}) \to \mathbf{P}(\mathcal{E}) = P$。
这是 $X$ 上光滑射影概形之间的态射，它把 $P'$ 表示为 $P$
上的有效 Cartier 除子。因此由引理 \ref{derham-lemma-check-log-smooth}，
对 $P' \subset P$ 而言，其在 $\mathbf{Z}$ 上的对数极点复形有定义。
所以对于 $a \in A(P)$，若 $i^*a = 0$，则由引理
\ref{derham-lemma-log-complex-consequence}，有
$a \cup c_1^A(\mathcal{O}_P(P')) = 0$。证明完毕。
\end{proof}

\begin{remark}
\label{derham-remark-splitting-principle}
Weil 上同调理论，引理
\ref{weil-lemma-splitting-principle}（分裂原理）和
\ref{weil-lemma-chern-classes-E-tensor-L}（张量积的陈类）的类似结论，
对拟紧拟分离概形上的德拉姆陈类成立。这是显然的，因为我们已经在引理
\ref{derham-lemma-chern-classes} 的证明中说明，Weil 上同调理论，第
\ref{weil-section-chern} 节的所有公理均成立。
\end{remark}

\noindent
对 $\mathbf{Q}$ 上的概形，可以构造陈特征。

\begin{lemma}
\label{derham-lemma-chern-character}
存在唯一的规则，它为每个拟紧拟分离概形 $X$（在 $\mathbf{Q}$ 上）
指定陈特征
$$
ch^{dR} : K_0(\textit{Vect}(X)) \longrightarrow
\prod\nolimits_{i \geq 0} H_{dR}^{2i}(X/\mathbf{Q})
$$
并具有下列性质：
\begin{enumerate}
\item $ch^{dR}$ 对所有 $X$ 都是环同态；
\item 如果 $f : X' \to X$ 是 $\mathbf{Q}$ 上拟紧拟分离概形之间的态射，
那么 $f^* \circ ch^{dR} =  ch^{dR} \circ f^*$；
\item 给定 $\mathcal{L} \in \Pic(X)$，有
$ch^{dR}([\mathcal{L}]) = \exp(c_1^{dR}(\mathcal{L}))$。
\end{enumerate}
\end{lemma}

\noindent
该构造很容易推广到 $\mathbf{Q}$ 上的所有概形，但为此需要稍微加强
Weil 上同调理论，第 \ref{weil-section-chern} 节中的讨论。

\begin{proof}
完全仿照引理 \ref{derham-lemma-chern-classes} 的证明，可以说明
$\mathbf{Q}$ 上拟紧拟分离概形的范畴连同函子
$A^*(X) = \bigoplus_{i \geq 0} H_{dR}^{2i}(X/\mathbf{Q})$
满足 Weil 上同调理论，第 \ref{weil-section-chern} 节的公理。
此外，在这种情形下，$A(X)$ 是 $\mathbf{Q}$-代数（对所有 $X$）。
因此本引理由 Weil 上同调理论，命题
\ref{weil-proposition-chern-character} 得出。
\end{proof}







\section{一种 Weil 上同调理论}
\label{derham-section-weil}

\noindent
设 $k$ 是特征 $0$ 的域。本节证明函子
$$
X \longmapsto H^*_{dR}(X/k)
$$
定义了 $k$ 上以 $k$ 为系数的 Weil 上同调理论，其含义见
Weil 上同调理论，定义 \ref{weil-definition-weil-cohomology-theory}。
我们将检验本章前面的构造给出数据 (D0)、(D1)、(D2')，
并满足 Weil 上同调理论，第 \ref{weil-section-c1} 节的公理 (A1)--(A9)。

\medskip\noindent
在本节余下部分中，固定域 $k$（特征 $0$），并置 $F = k$。
取下列数据：
\begin{enumerate}
\item[(D0)] 对 $1$ 维 $F$-向量空间 $F(1)$，取
$F(1) = F = k$.
\item[(D1)] 对函子 $H^*$，取把光滑射影概形 $X$（在 $k$ 上）映到
$H^*_{dR}(X/k)$ 的函子。函子性在第
\ref{derham-section-de-rham-cohomology} 节讨论，杯积在第
\ref{derham-section-cup-product} 节讨论。由引理
\ref{derham-lemma-cup-product-graded-commutative}，得到分次交换 $F$-代数。
\item[(D2')] 对映射 $c_1^H : \Pic(X) \to H^2(X)(1)$，使用第
\ref{derham-section-first-chern-class} 节引入的德拉姆第一陈类。
\end{enumerate}
我们将证明公理 (A1)--(A9) 成立。

\medskip\noindent
本段使用 Weil 上同调理论，引理
\ref{weil-lemma-check-over-extension}，把公理的检验约化到 $k$
为代数闭域的情形。把 $k'$ 记为 $k$ 的代数闭包，置 $F' = k'$。
完全按照上述方式，得到 $k'$ 上系数域为 $F'$ 的数据 (D0)、(D1)、(D2')。
由引理 \ref{derham-lemma-proper-smooth-de-Rham}，存在函子性同构
$$
H_{dR}^{2d}(X/k) \otimes_k k'
\longrightarrow
H_{dR}^{2d}(X_{k'}/k')
$$
，其中 $X$ 是 $k$ 上光滑射影概形。此外，诸图
$$
\xymatrix{
\Pic(X) \ar[r]_{c^{dR}_1} \ar[d] & H_{dR}^2(X/k) \ar[d] \\
\Pic(X_{k'}) \ar[r]^{c^{dR}_1} & H_{dR}^2(X_{k'}/k')
}
$$
由引理 \ref{derham-lemma-pullback-c1} 交换。这就完成了约化的证明。

\medskip\noindent
假设 $k$ 是特征零的代数闭域。我们将对上面给出的数据
(D0)、(D1)、(D2') 证明公理 (A1)--(A9)。

\medskip\noindent
公理 (A1)。这里需要检验
$H^*_{dR}(X \coprod Y/k) =  H^*_{dR}(X/k) \times H^*_{dR}(Y/k)$.
这是因为德拉姆上同调是取微分分次代数层（在扎里斯基拓扑中）的上同调构造的。

\medskip\noindent
公理 (A2)。它正是说取可逆模的第一陈类与拉回相容。
这由更一般的引理 \ref{derham-lemma-pullback-c1} 得出。

\medskip\noindent
公理 (A3)。这由更一般的命题
\ref{derham-proposition-projective-space-bundle-formula} 得出。

\medskip\noindent
公理 (A4)。这由更一般的引理
\ref{derham-lemma-log-complex-consequence} 得出。

\medskip\noindent
至此，使用 Weil 上同调理论，引理
\ref{weil-lemma-chern-classes} 和
\ref{weil-lemma-cycle-classes}，已经得到陈特征和循环类映射
$$
\gamma :
\CH^*(X)
\longrightarrow
\bigoplus\nolimits_{i \geq 0} H^{2i}_{dR}(X/k)
$$
；对光滑射影概形 $X$（在 $k$ 上），它们是分次环同态，并与态射
$f : X \to Y$（在 $k$ 上光滑射影概形之间）的拉回相容。

\medskip\noindent
公理 (A5)。有 $H_{dR}^*(\Spec(k)/k) = k = F$（在次数 $0$ 上）。
由引理 \ref{derham-lemma-kunneth-de-rham}，
两个光滑射影 $k$-概形的积满足 K\"unneth 公式（注意，该陈述中的导出张量积
没有问题，因为是在域 $k$ 上张量）。

\medskip\noindent
公理 (A7)。这由命题 \ref{derham-proposition-blowup-split} 得出。

\medskip\noindent
公理 (A8)。设 $X$ 是 $k$ 上光滑射影概形。由 Weil 上同调理论，
第 \ref{weil-section-c1} 节对该公理的解释可知，
$k' = H^0(X, \mathcal{O}_X)$ 是有限可分 $k$-代数。于是
$H_{dR}^*(\Spec(k')/k) = k'$ 位于次数 $0$，因为
$\Omega_{k'/k} = 0$。由引理 \ref{derham-lemma-bottom-part-degenerates}，
还有 $H_{dR}^0(X, \mathcal{O}_X) = k'$，从而得到该公理。

\medskip\noindent
公理 (A6)。设 $X$ 是 $k$ 上非空光滑射影概形，等维且维数为 $d$。
把 $\Delta : X \to X \times_{\Spec(k)} X$
记为 $X$ 在 $k$ 上的对角态射。需要证明存在 $k$-线性映射
$$
\lambda : H_{dR}^{2d}(X/k) \longrightarrow k
$$
使得 $(1 \otimes \lambda)\gamma([\Delta]) = 1$ 在
$H^0_{dR}(X/k)$ 中成立。写成
$$
\gamma = \gamma([\Delta]) = \gamma_0 + \ldots + \gamma_{2d}
$$
，其中 $\gamma_i \in H_{dR}^i(X/k) \otimes_k H_{dR}^{2d - i}(X/k)$
是 K\"unneth 分量。问题在于证明存在线性映射
$\lambda : H_{dR}^{2d}(X/k) \to k$，使得
$(1 \otimes \lambda)\gamma_0 = 1$ 在 $H^0_{dR}(X/k)$ 中成立。

\medskip\noindent
令 $X = \coprod X_i$ 为 $X$ 分解成连通（因而不可约）分支的分解。
相应地，有 $\Delta = \coprod \Delta_i$，其中
$\Delta_i \subset X_i \times X_i$。于是
$$
\gamma([\Delta]) = \sum \gamma([\Delta_i])
$$
而且，$\gamma([\Delta_i])$ 对应于
$\Delta_i \subset X_i \times X_i$ 的类，这是经由分解
$$
H^*_{dR}(X \times X) = \prod\nolimits_{i, j} H^*_{dR}(X_i \times X_j)
$$
实现的。
略去细节；一种证明方法是利用 $\CH^0(X \times X)$ 中存在幂等元
$e_{i, j}$，它对应于开闭子概形 $X_i \times X_j$；并利用 $\gamma$ 是环同态，
它把 $e_{i, j}$ 映到所示上同调乘积分解中相应的幂等元。
如果能找到 $\lambda_i : H_{dR}^{2d}(X_i/k) \to k$，使得
$(1 \otimes \lambda_i)\gamma([\Delta_i]) = 1$ 在
$H^0_{dR}(X_i/k)$ 中成立，那么取 $\lambda = \sum \lambda_i$
就能对 $X$ 解决问题。因此可以并且确实假设 $X$ 不可约。

\medskip\noindent
证明不可约 $X$ 的公理 (A6)。由于 $k$ 代数闭，有
$H^0_{dR}(X/k) = k$，因为 $H^0(X, \mathcal{O}_X) = k$；
这里使用了 $X$ 是代数闭域上的射影簇（例如见簇，引理
\ref{varieties-lemma-proper-geometrically-reduced-global-sections}）。
令 $x \in X$ 为任意闭点。考虑笛卡尔图
$$
\xymatrix{
x \ar[d] \ar[r] & X \ar[d]^\Delta \\
X \ar[r]^-{x \times \text{id}} & X \times_{\Spec(k)} X
}
$$
$\gamma$ 与拉回的相容性表明，$\gamma([\Delta])$ 映到
$\gamma([x])$（在 $H_{dR}^{2d}(X/k)$ 中）；换言之，有
$\gamma_0 = 1 \otimes \gamma([x])$。由此得到两点：
(a) 类 $\gamma([x])$ 与 $x$ 无关；(b) 只需证明类
$\gamma([x])$ 非零；因而 (c) 只需找到任意零维循环
$\alpha$（在 $X$ 上），使得 $\gamma(\alpha) \not = 0$。为此，选择有限态射
$$
f : X \longrightarrow \mathbf{P}^d_k
$$
关于这种态射的存在性，见相交理论，第
\ref{intersection-section-projection} 节，特别是引理
\ref{intersection-lemma-projection-generically-finite}。注意，$f$
是有限 syntomic 态射（由更多态射，引理
\ref{more-morphisms-lemma-lci-permanence}，它是局部完全交态射；
由代数，引理 \ref{algebra-lemma-CM-over-regular-flat}，它是平坦的）。
由命题 \ref{derham-proposition-Garel}，有迹映射
$$
\Theta_f :
f_*\Omega^\bullet_{X/k}
\longrightarrow
\Omega^\bullet_{\mathbf{P}^d_k/k}
$$
，它与典范映射
$$
\Omega^\bullet_{\mathbf{P}^d_k/k}
\longrightarrow
f_*\Omega^\bullet_{X/k}
$$
的复合是乘以 $f$ 的次数。因此得到映射
$$
\Theta : H_{dR}^{2d}(X/k) \to H_{dR}^{2d}(\mathbf{P}^d_k/k)
$$
使得 $\Theta \circ f^*$ 是乘以某个正整数。因此，如果能找到
$\mathbf{P}^d_k$ 上类非零的零维循环，就可由 $\gamma$ 与拉回的相容性得出结论。
这由引理 \ref{derham-lemma-de-rham-cohomology-projective-space} 成立，
从而完成公理 (A6) 的证明。

\medskip\noindent
下文将不加说明地使用以下事实。第一，由 Weil 上同调理论，备注
\ref{weil-remark-trace}，映射
$\lambda_X : H^{2d}_{dR}(X/k) \to k$ 是唯一的。第二，
在公理 (A6) 的证明中已经看到
$\lambda_X(\gamma([x])) = 1$，只要 $X$ 不可约；也就是说，循环类映射
$\gamma : \CH^d(X) \to H_{dR}^{2d}(X/k)$ 与 $\lambda_X$
的复合是次数映射。

\medskip\noindent
公理 (A9)。设 $Y \subset X$ 是非空光滑等维射影概形
$X$（在 $k$ 上，维数为 $d$）上的非空光滑除子。需要证明图
$$
\xymatrix{
H_{dR}^{2d - 2}(X/k)
\ar[rrr]_{c^{dR}_1(\mathcal{O}_X(Y)) \cap -} \ar[d]_{restriction} & & &
H_{dR}^{2d}(X) \ar[d]^{\lambda_X} \\
H_{dR}^{2d - 2}(Y/k) \ar[rrr]^-{\lambda_Y} & & & k
}
$$
交换，其中 $\lambda_X$ 和 $\lambda_Y$ 如公理 (A6) 所述。
上面已经看到，如果把 $X = \coprod X_i$ 分解为连通（等价地，不可约）分支，
那么相应地 $\lambda_X = \sum \lambda_{X_i}$。
类似地，如果把 $Y = \coprod Y_j$ 分解为连通（等价地，不可约）分支，
那么 $\lambda_Y = \sum \lambda_{Y_j}$。此外，此时有
$\mathcal{O}_X(Y) = \otimes_j \mathcal{O}_X(Y_j)$，因而
$$
c_1^{dR}(\mathcal{O}_X(Y)) = \sum\nolimits_j
c^{dR}_1(\mathcal{O}_X(Y_j))
$$
在 $H_{dR}^2(X/k)$ 中成立。直接追图表明，只需在 $X$ 和 $Y$
都不可约时证明该图交换。此时 $H_{dR}^{2d - 2}(Y/k)$ 是 $1$ 维的，
因为由 Weil 上同调理论，引理 \ref{weil-lemma-poincare-duality}，
$Y$ 满足庞加莱对偶。由公理 (A4)，限制映射（左竖直箭头）的核
包含在与 $c^{dR}_1(\mathcal{O}_X(Y))$ 作杯积的映射之核中。
这意味着，只需找到一个上同调类
$a \in H_{dR}^{2d - 2}(X)$，其在 $Y$ 上的限制非零，并且对 $a$
该图交换。取任意丰沛可逆模 $\mathcal{L}$，并置
$$
a = c^{dR}_1(\mathcal{L})^{d - 1}
$$
于是 $a|_Y = c^{dR}_1(\mathcal{L}|_Y)^{d - 1}$，因而
$$
\lambda_Y(a|_Y) = \deg(c_1(\mathcal{L}|_Y)^{d - 1} \cap [Y])
$$
，这里使用了上面对 $\lambda_Y$ 的描述。由 Chow 同调，引理
\ref{chow-lemma-degrees-and-numerical-intersections} 结合簇，引理
\ref{varieties-lemma-ample-positive}，这是一个正整数。类似地，得到
$$
\lambda_X(c^{dR}_1(\mathcal{O}_X(Y)) \cap a) =
\deg(c_1(\mathcal{O}_X(Y)) \cap c_1(\mathcal{L})^{d - 1} \cap [X])
$$
由于几乎由定义就知道 $c_1(\mathcal{O}_X(Y)) \cap [X] = [Y]$，
因而有零维循环的等式
$$
(Y \to X)_*\left(c_1(\mathcal{L}|_Y)^{d - 1} \cap [Y]\right) =
c_1(\mathcal{O}_X(Y)) \cap c_1(\mathcal{L})^{d - 1} \cap [X]
$$
（在 $X$ 上）。因此这些循环具有相同次数，证明完毕。

\begin{proposition}
\label{derham-proposition-de-rham-is-weil}
设 $k$ 是特征零的域。把光滑射影概形 $X$（在 $k$ 上）映到
$H_{dR}^*(X/k)$ 的函子，是 Weil 上同调理论，含义见
Weil 上同调理论，定义 \ref{weil-definition-weil-cohomology-theory}。
\end{proposition}

\begin{proof}
在上面的讨论中，我们说明了数据 (D0)、(D1)、(D2') 满足
Weil 上同调理论，第 \ref{weil-section-c1} 节的公理 (A1)--(A9)。
因此结论由 Weil 上同调理论，命题 \ref{weil-proposition-get-weil} 得出。

\medskip\noindent
请不要阅读下文。在这些断言的证明中，我们还使用了引理
\ref{derham-lemma-proper-smooth-de-Rham}、
\ref{derham-lemma-pullback-c1},
\ref{derham-lemma-log-complex-consequence},
\ref{derham-lemma-kunneth-de-rham},
\ref{derham-lemma-bottom-part-degenerates} 和
\ref{derham-lemma-de-rham-cohomology-projective-space},
命题
\ref{derham-proposition-projective-space-bundle-formula},
\ref{derham-proposition-blowup-split} 和
\ref{derham-proposition-Garel},
Weil 上同调理论，引理
\ref{weil-lemma-check-over-extension},
\ref{weil-lemma-chern-classes},
\ref{weil-lemma-cycle-classes} 和
\ref{weil-lemma-poincare-duality},
Weil 上同调理论，备注 \ref{weil-remark-trace}，簇，引理
\ref{varieties-lemma-proper-geometrically-reduced-global-sections} 和
\ref{varieties-lemma-ample-positive},
相交理论，第 \ref{intersection-section-projection} 节及引理
\ref{intersection-lemma-projection-generically-finite}，
更多态射，引理 \ref{more-morphisms-lemma-lci-permanence}，
代数，引理 \ref{algebra-lemma-CM-over-regular-flat}，以及 Chow 同调，引理
\ref{chow-lemma-degrees-and-numerical-intersections}.
\end{proof}

\begin{remark}
\label{derham-remark-hodge-cohomology-is-weil}
完全按照上述方式，可以证明霍奇上同调
$X \mapsto H_{Hodge}^*(X/k)$ 配备 $c_1^{Hodge}$ 后确定一个 Weil 上同调理论。
如果将来需要这一点，我们会在这里精确表述并证明。这导致以下有趣结论：
如果 Weil 上同调理论的 Betti 数与所选 Weil 上同调理论无关
（在我们的域 $k$（特征 $0$）上），那么霍奇到德拉姆谱序列在
$E_1$ 页退化！当然，霍奇到德拉姆谱序列的退化是已知的
（例如见 \cite{Deligne-Illusie} 中精彩的代数证明），但它绝非容易的结果！
这提示证明 Betti 数的无关性同样困难；据我们所知，这仍是一个开放问题。
相关问题见 Weil 上同调理论，备注
\ref{weil-remark-betti-numbers-in-some-sense}。
\end{remark}







\section{闭浸入的 Gysin 映射}
\label{derham-section-gysin}

\noindent
本节定义闭浸入的 Gysin 映射。

\begin{remark}
\label{derham-remark-gysin-equations}
设 $X \to S$ 是概形的态射。设
$f_1, \ldots, f_c \in \Gamma(X, \mathcal{O}_X)$。设 $Z \subset X$
是由 $f_1, \ldots, f_c$ 截出的闭子概形。下面研究 {\it Gysin 映射}
\begin{equation}
\label{derham-equation-gysin}
\gamma^p_{f_1, \ldots, f_c} :
\Omega^p_{Z/S}
\longrightarrow
\mathcal{H}_Z^c(\Omega^{p + c}_{X/S})
\end{equation}
，其定义如下。给定局部截面 $\omega$（属于 $\Omega^p_{Z/S}$），
它是某个截面 $\tilde \omega$（属于 $\Omega^p_{X/S}$）的限制，置
$$
\gamma^p_{f_1, \ldots, f_c}(\omega) =
c_{f_1, \ldots, f_c}(\tilde \omega|_Z) \wedge
\text{d}f_1 \wedge \ldots \wedge \text{d}f_c
$$
，其中 $c_{f_1, \ldots, f_c} : \Omega^p_{X/S} \otimes \mathcal{O}_Z \to
\mathcal{H}_Z^c(\Omega^p_{X/S})$ 是概形的导出范畴，备注
\ref{perfect-remark-supported-map-c-equations} 中构造的映射。
此定义是良定义的：给定 $\omega$，可以把所选的
$\tilde \omega$ 改变如下形式的元素：
$\sum f_i \omega'_i + \sum \text{d}(f_i) \wedge \omega''_i$
，而该构造把这些元素映到零。
\end{remark}

\begin{lemma}
\label{derham-lemma-gysin-differential}
Gysin 映射 (\ref{derham-equation-gysin}) 与 $\Omega^\bullet_{X/S}$ 和
$\Omega^\bullet_{Z/S}$ 上的德拉姆微分相容。
\end{lemma}

\begin{proof}
一旦正确解释该陈述，它就由一个几乎平凡的计算得出。首先回忆，
函子 $\mathcal{H}^c_Z$ 在 $\mathcal{O}_X$-模范畴上计算时，
与 $X$ 上阿贝尔层范畴中类似定义的函子一致；见上同调，引理
\ref{cohomology-lemma-sections-support-abelian-unbounded}。
因此，微分 $\text{d} : \Omega^p_{X/S} \to \Omega^{p + 1}_{X/S}$
诱导映射
$\mathcal{H}^c_Z(\Omega^p_{X/S}) \to \mathcal{H}^c_Z(\Omega^{p + 1}_{X/S})$.
此外，概形的导出范畴，备注
\ref{perfect-remark-support-c-equations} 中扩张交错 {\v C}ech 复形的构造
也适用于阿贝尔层范畴。映射
$$
\Coker\left(\bigoplus \mathcal{F}_{1 \ldots \hat i \ldots c} \to
\mathcal{F}_{1 \ldots c}\right)
\longrightarrow
i_*\mathcal{H}^c_Z(\mathcal{F})
$$
用于概形的导出范畴，备注
\ref{perfect-remark-supported-map-c-equations} 中
$c_{f_1, \ldots, f_c}$ 的构造；它在 $X$ 上所有阿贝尔层的范畴中
是良定义且函子性的。因此本引理由等式
$$
\text{d}\left(
\frac{\tilde \omega \wedge \text{d}f_1 \wedge \ldots \wedge
\text{d}f_c}{f_1 \ldots f_c}\right) =
\frac{\text{d}(\tilde \omega) \wedge
\text{d}f_1 \wedge \ldots \wedge \text{d}f_c}{f_1 \ldots f_c}
$$
得出，而该等式显然成立。
\end{proof}

\begin{lemma}
\label{derham-lemma-gysin-global}
设 $X \to S$ 是概形的态射。设 $Z \to X$ 是有限表示的闭浸入，
其余法层 $\mathcal{C}_{Z/X}$ 是秩 $c$ 的局部自由模。于是存在典范映射
$$
\gamma^p : \Omega^p_{Z/S} \to \mathcal{H}^c_Z(\Omega^{p + c}_{X/S})
$$
；局部上，它由备注 \ref{derham-remark-gysin-equations} 中的映射
$\gamma^p_{f_1, \ldots, f_c}$ 给出。
\end{lemma}

\begin{proof}
这些假设表明，给定 $x \in Z \subset X$，存在开邻域 $U$（包含 $x$），
使得 $Z$ 由 $c$ 个元素 $f_1, \ldots, f_c \in \mathcal{O}_X(U)$ 截出。
因此只需证明：给定 $f_1, \ldots, f_c$ 和 $g_1, \ldots, g_c$
（它们位于 $\mathcal{O}_X(U)$ 中并截出 $Z \cap U$），映射
$\gamma^p_{f_1, \ldots, f_c}$ 与
$\gamma^p_{g_1, \ldots, g_c}$ 相同。为此，缩小 $U$ 后可以假设
$g_j = \sum a_{ji} f_i$，其中
$a_{ji} \in \mathcal{O}_X(U)$。于是由概形的导出范畴，引理
\ref{perfect-lemma-supported-map-determinant}，有
$c_{f_1, \ldots, f_c} = \det(a_{ji}) c_{g_1, \ldots, g_c}$。
另一方面，有
$$
\text{d}(g_1) \wedge \ldots \wedge \text{d}(g_c) \equiv
\det(a_{ji}) \text{d}(f_1) \wedge \ldots \wedge \text{d}(f_c)
\bmod (f_1, \ldots, f_c)\Omega^c_{X/S}
$$
合并这些关系，直接计算即得所需等式。
\end{proof}

\begin{lemma}
\label{derham-lemma-gysin-differential-global}
设 $X \to S$ 和 $i : Z \to X$ 如引理 \ref{derham-lemma-gysin-global} 所述。
Gysin 映射 $\gamma^p$ 与 $\Omega^\bullet_{X/S}$ 和
$\Omega^\bullet_{Z/S}$ 上的德拉姆微分相容。
\end{lemma}

\begin{proof}
可以局部地检验，而此时结论由引理 \ref{derham-lemma-gysin-differential} 得出。
\end{proof}

\begin{lemma}
\label{derham-lemma-gysin-projection}
设 $X \to S$ 和 $i : Z \to X$ 如引理 \ref{derham-lemma-gysin-global} 所述。
给定 $\alpha \in H^q(X, \Omega^p_{X/S})$，有
$\gamma^p(\alpha|_Z) = i^{-1}\alpha \wedge \gamma^0(1)$，这是在
$H^q(Z, \mathcal{H}^c_Z(\Omega^{p + c}_{X/S}))$.
其中的记号见证明。
\end{lemma}

\begin{proof}
限制 $\alpha|_Z$ 是由霍奇上同调的函子性给出的
$H^q(Z, \Omega^p_{Z/S})$ 中的元素。对上同调应用函子性并使用
$\gamma^p : \Omega^p_{Z/S} \to \mathcal{H}^c_Z(\Omega^{p + c}_{X/S})$
，得到 $\gamma^p(\alpha|_Z)$，它属于
$H^q(Z, \mathcal{H}^c_Z(\Omega^{p + c}_{X/S}))$。
这解释了公式的左边。

\medskip\noindent
为解释右边，先沿环化空间的映射
$i : (Z, i^{-1}\mathcal{O}_X) \to (X, \mathcal{O}_X)$ 拉回，
得到元素 $i^{-1}\alpha \in H^q(Z, i^{-1}\Omega^p_{X/S})$。
令 $\gamma^0(1) \in H^0(Z, \mathcal{H}_Z^c(\Omega^c_{X/S}))$
为 $1 \in H^0(Z, \mathcal{O}_Z) = H^0(Z, \Omega^0_{Z/S})$
在 $\gamma^0$ 下的像。使用杯积，得到元素
$$
i^{-1}\alpha \cup \gamma^0(1)
\in
H^{q + c}(Z, 
i^{-1}\Omega^p_{X/S} \otimes_{i^{-1}\mathcal{O}_X}
\mathcal{H}^c_Z(\Omega^c_{X/S}))
$$
使用上同调，备注 \ref{cohomology-remark-support-cup-product}
以及楔积，得到典范映射
$$
i^{-1}\Omega^p_{X/S} \otimes_{i^{-1}\mathcal{O}_X}^\mathbf{L}
R\mathcal{H}_Z(\Omega^c_{X/S}) \to
R\mathcal{H}_Z(\Omega^p_{X/S} \otimes_{\mathcal{O}_X}^\mathbf{L}
\Omega^c_{X/S}) \to
R\mathcal{H}_Z(\Omega^{p + c}_{X/S})
$$
由概形的导出范畴，引理
\ref{perfect-lemma-supported-trivial-vanishing}，对象
$R\mathcal{H}_Z(\Omega^j_{X/S})$ 的上同调层在次数 $> c$ 时消失。
因此在次数 $c$ 的上同调层上得到映射
$$
i^{-1}\Omega^p_{X/S} \otimes_{i^{-1}\mathcal{O}_X}
\mathcal{H}^c_Z(\Omega^c_{X/S}) \longrightarrow
\mathcal{H}^c_Z(\Omega^{p + c}_{X/S})
$$
由上同调的函子性，表达式 $i^{-1}\alpha \wedge \gamma^0(1)$
是杯积 $i^{-1}\alpha \cup \gamma^0(1)$ 的像。

\medskip\noindent
按这种方式解释公式的含义后，由杯积的一般性质
（上同调，第 \ref{cohomology-section-cup-product} 节），
现在只需证明图
$$
\xymatrix{
i^{-1}\Omega^p_X \otimes \Omega^0_Z \ar[rr]_{\text{id} \otimes \gamma^0}
\ar[d] & &
i^{-1}\Omega^p_X \otimes \mathcal{H}^c_Z(\Omega^c_X) \ar[d]^\wedge \\
\Omega^p_Z \otimes \Omega^0_Z \ar[r]^\wedge &
\Omega^p_Z \ar[r]^{\gamma^p} &
\mathcal{H}^c_Z(\Omega^{p + c}_X)
}
$$
在 $Z$ 上层的范畴中交换（这里显然略有滥用记号）。
这归结为对映射 $\gamma^j_{f_1, \ldots, f_c}$ 的简单计算，我们将其略去；
事实上，选择这些映射正是为了使上述交换性成立，并使 $1$ 映到
$\frac{\text{d}f_1 \wedge \ldots \wedge \text{d}f_c}{f_1 \ldots f_c}$.
\end{proof}

\begin{lemma}
\label{derham-lemma-gysin-transverse}
设 $c \geq 0$ 是整数，并设
$$
\xymatrix{
Z' \ar[d]_h \ar[r] & X' \ar[d]_g \ar[r] & S' \ar[d] \\
Z \ar[r] & X \ar[r] & S
}
$$
为概形的交换图。假设：
\begin{enumerate}
\item $Z \to X$ and $Z' \to X'$
满足引理 \ref{derham-lemma-gysin-global} 的假设；
\item 图中的左方块是笛卡尔方块；
\item $h^*\mathcal{C}_{Z/X} \to \mathcal{C}_{Z'/X'}$
（态射，引理 \ref{morphisms-lemma-conormal-functorial}）是同构。
\end{enumerate}
那么图
$$
\xymatrix{
h^*\Omega^p_{Z/S} \ar[rr]_-{h^{-1}\gamma^p} \ar[d] & &
\mathcal{O}_{X'}|_{Z'} \otimes_{h^{-1}\mathcal{O}_X|_Z}
h^{-1}\mathcal{H}^c_Z(\Omega^{p + c}_{X/S}) \ar[d] \\
\Omega^p_{Z'/S'} \ar[rr]^{\gamma^p} & &
\mathcal{H}^c_{Z'}(\Omega^{p + c}_{X'/S'})
}
$$
交换。左竖直箭头来自微分模的函子性，右竖直箭头使用上同调，备注
\ref{cohomology-remark-support-functorial}。
\end{lemma}

\begin{proof}
更确切地说，考虑复合
\begin{align*}
\mathcal{O}_{X'}|_{Z'} \otimes_{h^{-1}\mathcal{O}_X|_Z}^\mathbf{L}
h^{-1}R\mathcal{H}_Z(\Omega^{p + c}_{X/S})
& \to
R\mathcal{H}_{Z'}(Lg^*\Omega^{p + c}_{X/S}) \\
& \to
R\mathcal{H}_{Z'}(g^*\Omega^{p + c}_{X/S}) \\
& \to
R\mathcal{H}_{Z'}(\Omega^{p + c}_{X'/S'})
\end{align*}
其中第一个箭头由上同调，备注
\ref{cohomology-remark-support-functorial} 给出，最后一个箭头由微分的函子性给出。
由概形的导出范畴，引理
\ref{perfect-lemma-supported-trivial-vanishing}，上同调层在次数 $> c$
时消失，因此该复合诱导右竖直箭头。交换性可以局部检验。
因此可以假设 $Z$ 由
$f_1, \ldots, f_c \in \Gamma(X, \mathcal{O}_X)$ 截出。
那么 $Z'$ 由 $f'_i = g^\sharp(f_i)$ 截出。映射
$c_{f_1, \ldots, f_c}$ 和 $c_{f'_1, \ldots, f'_c}$
嵌入交换图
$$
\xymatrix{
h^*i^*\Omega^p_{X/S} \ar[rr]_-{h^{-1}c_{f_1, \ldots, f_c}} \ar[d] & &
\mathcal{O}_{X'}|_{Z'} \otimes_{h^{-1}\mathcal{O}_X|_Z}
h^{-1}\mathcal{H}^c_Z(\Omega^p_{X/S}) \ar[d] \\
(i')^*\Omega^p_{X'/S'} \ar[rr]^{c_{f'_1, \ldots, f'_c}} & &
\mathcal{H}^c_{Z'}(\Omega^p_{X'/S'})
}
$$
见概形的导出范畴，备注
\ref{perfect-remark-supported-functorial}。回忆，给定一个
$p$-形式 $\omega$（在 $Z$ 上），我们如下定义 $\gamma^p(\omega)$：
局部地在 $X$ 和 $Z$ 上，选择一个 $p$-形式 $\tilde \omega$
（在 $X$ 上）来提升 $\omega$，并取
$\gamma^p(\omega) =
c_{f_1, \ldots, f_c}(\tilde \omega) \wedge
\text{d}f_1 \wedge \ldots \wedge \text{d}f_c$.
由于形式
$\text{d}f_1 \wedge \ldots \wedge \text{d}f_c$ 拉回到
$\text{d}f'_1 \wedge \ldots \wedge \text{d}f'_c$，结论成立。
\end{proof}

\begin{remark}
\label{derham-remark-how-to-use}
设 $X \to S$、$i : Z \to X$ 和 $c \geq 0$ 如引理
\ref{derham-lemma-gysin-global} 所述。设 $p \geq 0$，并假设
$\mathcal{H}^i_Z(\Omega^{p + c}_{X/S}) = 0$ 对
$i = 0, \ldots, c - 1$ 成立。如果 $X \to S$ 光滑且
$Z \to X$ 是 Koszul 正则浸入，则此消失成立；见概形的导出范畴，引理
\ref{perfect-lemma-supported-vanishing}。于是得到映射
$$
\gamma^{p, q} :
H^q(Z, \Omega^p_{Z/S})
\longrightarrow
H^{q + c}(X, \Omega^{p + c}_{X/S})
$$
：先使用
$\gamma^p : \Omega^p_{Z/S} \to \mathcal{H}^c_Z(\Omega^{p + c}_{X/S})$
映入
$$
H^q(Z, \mathcal{H}^c_Z(\Omega^{p + c}_{X/S})) =
H^q(Z, R\mathcal{H}_Z(\Omega^{p + c}_{X/S})[c]) =
H^q(X, i_*R\mathcal{H}_Z(\Omega^{p + c}_{X/S})[c])
$$
，再使用伴随映射
$i_*R\mathcal{H}_Z(\Omega^{p + c}_{X/S}) \to \Omega^{p + c}_{X/S}$
继续映到所需的霍奇上同调模。
\end{remark}

\begin{lemma}
\label{derham-lemma-gysin-differential-hodge}
设 $X \to S$ 和 $i : Z \to X$ 如引理 \ref{derham-lemma-gysin-global} 所述。
假设 $X \to S$ 光滑且 $Z \to X$ 为 Koszul 正则浸入。
Gysin 映射 $\gamma^{p, q}$ 与 $\Omega^\bullet_{X/S}$ 和
$\Omega^\bullet_{Z/S}$ 上的德拉姆微分相容。
\end{lemma}

\begin{proof}
这立即由引理 \ref{derham-lemma-gysin-differential-global} 得出。
\end{proof}

\begin{lemma}
\label{derham-lemma-gysin-projection-global}
设 $X \to S$、$i : Z \to X$ 和 $c \geq 0$ 如引理
\ref{derham-lemma-gysin-global} 所述。假设 $X \to S$ 光滑且
$Z \to X$ 为 Koszul 正则浸入。给定
$\alpha \in H^q(X, \Omega^p_{X/S})$，有
$\gamma^{p, q}(\alpha|_Z) = \alpha \cup \gamma^{0, 0}(1)$，这是在
$H^{q + c}(X, \Omega^{p + c}_{X/S})$ 中成立，其中
$\gamma^{a, b}$ 如备注 \ref{derham-remark-how-to-use} 所述。
\end{lemma}

\begin{proof}
本引理由引理 \ref{derham-lemma-gysin-projection} 和上同调，引理
\ref{cohomology-lemma-support-cup-product} 得出。
建议读者跳过下面更详细的讨论。

\medskip\noindent
下文将不加说明地使用
$R\mathcal{H}_Z(\Omega^j_{X/S}) = \mathcal{H}^c_Z(\Omega^j_{X/S})[-c]$
对所有 $j$ 成立，正如备注 \ref{derham-remark-how-to-use} 所指出的。
还将不加说明地使用等同
$H^{q + c}_Z(X, \Omega^j_{X/S}) = H^{q + c}(Z, R\mathcal{H}_Z(\Omega^j_{X/S}) =
H^q(Z, \mathcal{H}^c_Z(\Omega^j_{X/S}))$；第一个等同见上同调，引理
\ref{cohomology-lemma-local-to-global-sections-with-support}。
在这些等同下：
\begin{enumerate}
\item $\gamma^0(1) \in H^c_Z(X, \Omega^c_{X/S})$ 映到 $\gamma^{0, 0}(1)$，
后者位于 $H^c(X, \Omega^c_{X/S})$ 中；
\item 等式的右边 $i^{-1}\alpha \wedge \gamma^0(1)$（在引理
\ref{derham-lemma-gysin-projection} 中）是杯积在楔积下的像；该杯积来自
上同调，备注 \ref{cohomology-remark-support-cup-product-global}，
作用于元素 $\alpha$ 和 $\gamma^0(1)$。换言之，引理
\ref{derham-lemma-gysin-projection} 的证明中的构造与上同调，备注
\ref{cohomology-remark-support-cup-product-global} 中的构造相符；
\item 由上同调，引理 \ref{cohomology-lemma-support-cup-product}，
它映到 $\alpha \cup \gamma^{0, 0}(1)$（在
$H^{q + c}(X, \Omega^p_{X/S} \otimes \Omega^c_{X/S})$ 中）；
\item 等式的左边 $\gamma^p(\alpha|_Z)$（在引理
\ref{derham-lemma-gysin-projection} 中）映到
$\gamma^{p, q}(\alpha|_Z)$。
\end{enumerate}
证明完毕。
\end{proof}

\begin{lemma}
\label{derham-lemma-gysin-transverse-global}
设 $c \geq 0$，并设
$$
\xymatrix{
Z' \ar[d]_h \ar[r] & X' \ar[d]_g \ar[r] & S' \ar[d] \\
Z \ar[r] & X \ar[r] & S
}
$$
满足引理 \ref{derham-lemma-gysin-transverse} 的假设；此外还假设
$X \to S$ 和 $X' \to S'$ 光滑，且 $Z \to X$ 和 $Z' \to X'$
是 Koszul 正则浸入。那么图
$$
\xymatrix{
H^q(Z, \Omega^p_{Z/S}) \ar[rr]_-{\gamma^{p, q}} \ar[d] & &
H^{q + c}(X, \Omega^{p + c}_{X/S}) \ar[d] \\
H^q(Z', \Omega^p_{Z'/S'}) \ar[rr]^{\gamma^{p, q}} & &
H^{q + c}(X', \Omega^{p + c}_{X'/S'})
}
$$
交换，其中 $\gamma^{p, q}$ 如备注 \ref{derham-remark-how-to-use} 所述。
\end{lemma}

\begin{proof}
这由引理 \ref{derham-lemma-gysin-transverse} 与上同调，引理
\ref{cohomology-lemma-support-functorial} 结合得出。
\end{proof}

\begin{lemma}
\label{derham-lemma-class-of-a-point}
设 $k$ 是一个域。设 $X$ 是 $k$ 上维数为 $d$ 的不可约光滑固有概形。
设 $Z \subset X$ 是仅由一个 $k$-有理点 $x$ 组成的既约闭子概形。那么
$1 \in k = H^0(Z, \mathcal{O}_Z) = H^0(Z, \Omega^0_{Z/k})$
在备注 \ref{derham-remark-how-to-use} 的映射
$H^0(Z, \Omega^0_{Z/k}) \to H^d(X, \Omega^d_{X/k})$ 下的像非零。
\end{lemma}

\begin{proof}
映射 $\gamma^0 : \mathcal{O}_Z \to
\mathcal{H}^d_Z(\Omega^d_{X/k}) = R\mathcal{H}_Z(\Omega^d_{X/k})[d]$
伴随于映射
$$
g^0 : i_*\mathcal{O}_Z \longrightarrow \Omega^d_{X/k}[d]
$$
（在 $D(\mathcal{O}_X)$ 中）。回忆，
$\Omega^d_{X/k} = \omega_X$ 是 $X/k$ 的对偶层；见概形对偶，引理
\ref{duality-lemma-duality-proper-over-field}。因此，本引理陈述中映射的
$k$-线性对偶是映射
$$
H^0(X, \mathcal{O}_X) \to \Ext^d_X(i_*\mathcal{O}_Z, \omega_X)
$$
，它把 $1$ 映到 $g^0$。因此只需证明 $g^0$ 非零。
这可以在任意邻域 $U$ 中完成，该邻域包含点 $x$。选择 $U$，使得存在
$f_1, \ldots, f_d \in \mathcal{O}_X(U)$，它们只在 $x$ 处消失，
并生成极大理想 $\mathfrak m_x \subset \mathcal{O}_{X, x}$。
可以假设 $U = \Spec(R)$ 是仿射的。回顾 $\gamma^0$ 的构造，
可见我们的扩张由下式给出：
$$
k \to
(R \to \bigoplus\nolimits_{i_0} R_{f_{i_0}} \to
\bigoplus\nolimits_{i_0 < i_1} R_{f_{i_0}f_{i_1}} \to
\ldots \to R_{f_1\ldots f_r})[d] \to R[d]
$$
其中 $1$ 在第一个映射下映到 $1/f_1 \ldots f_c$。
它非零，因为在此情形下 $1/f_1 \ldots f_c$ 是局部上同调群
$H^d_{(f_1, \ldots, f_d)}(R)$ 的非零元素，
\end{proof}






\section{相对庞加莱对偶}
\label{derham-section-relative-poincare-duality}

\noindent
本节证明基底上光滑固有概形的相对德拉姆上同调所满足的庞加莱对偶。
我们强烈建议读者先阅读第 \ref{derham-section-poincare-duality} 节。

\begin{situation}
\label{derham-situation-relative-duality}
这里 $S$ 是拟紧拟分离概形，而 $f : X \to S$ 是光滑固有的概形态射，
其所有纤维均非空且等维，维数为 $n$。
\end{situation}

\begin{lemma}
\label{derham-lemma-relative-bottom-part-degenerates}
在情形 \ref{derham-situation-relative-duality} 中，正向像
$f_*\mathcal{O}_X$ 是有限平展 $\mathcal{O}_S$-代数；并且在 $S$ 上局部地，
$Rf_*\mathcal{O}_X = f_*\mathcal{O}_X \oplus P$ 在 $D(\mathcal{O}_S)$ 中成立，
其中 $P$ 是 Tor 振幅落在 $[1, \infty)$ 中的完美对象。映射
$\text{d} : f_*\mathcal{O}_X \to f_*\Omega_{X/S}$ 为零。
\end{lemma}

\begin{proof}
陈述的第一部分由概形的导出范畴，引理
\ref{perfect-lemma-proper-flat-geom-red} 得出。置
$S' = \underline{\Spec}_S(f_*\mathcal{O}_X)$，便得到分解
$X \to S' \to S$（这是 Stein 分解；见态射进阶，第
\ref{more-morphisms-section-stein-factorization} 节，不过这里并不需要这一点）。
例如由态射，引理 \ref{morphisms-lemma-triangle-differentials} 和
\ref{morphisms-lemma-etale-at-point} 可见
$\Omega_{X/S} = \Omega_{X/S'}$。这当然蕴含
$\text{d} : f_*\mathcal{O}_X \to f_*\Omega_{X/S}$ 为零。
\end{proof}

\begin{lemma}
\label{derham-lemma-relative-duality-hodge}
在情形 \ref{derham-situation-relative-duality} 中，存在 $\mathcal{O}_S$-模映射
$$
t : Rf_*\Omega^n_{X/S}[n] \longrightarrow \mathcal{O}_S
$$
它在预复合以 $H^0(X, \mathcal{O}_X)$ 中某个单位的乘法这个意义下唯一，
并具有如下性质：对每个 $p$，配对
$$
Rf_*\Omega^p_{X/S}
\otimes_{\mathcal{O}_S}^\mathbf{L}
Rf_*\Omega^{n - p}_{X/S}[n]
\longrightarrow
\mathcal{O}_S
$$
由相对杯积再复合 $t$ 给出，并且是 $S$ 上完美复形之间的完美配对。
\end{lemma}

\begin{proof}
首先假设 $S$ 是诺特概形。令 $\omega^\bullet_{X/S}$ 为概形的对偶性，
备注 \ref{duality-remark-relative-dualizing-complex} 所述的
$X$ 在 $S$ 上的相对对偶复形，并令
$Rf_*\omega_{X/S}^\bullet \to \mathcal{O}_S$ 为其迹映射。由概形的对偶性，
引理 \ref{duality-lemma-smooth-proper}（此处使用了 $S$ 的诺特性），存在同构
$\omega^\bullet_{X/S} \cong \Omega^n_{X/S}[n]$；利用此同构便得到 $t$。
由引理 \ref{derham-lemma-proper-smooth-de-Rham}，复形 $Rf_*\Omega^p_{X/S}$ 是完美的。
由于 $\Omega^p_{X/S}$ 局部自由，并且
$\Omega^p_{X/S} \otimes_{\mathcal{O}_X} \Omega^{n - p}_{X/S} \to
\Omega^n_{X/S}$ 给出同构 $\Omega^p_{X/S} \cong
\SheafHom_{\mathcal{O}_X}(\Omega^{n - p}_{X/S}, \Omega^n_{X/S})$，
由概形的对偶性，备注
\ref{duality-remark-relative-dualizing-complex-relative-cup-product} 可知，
相对杯积诱导的配对是完美的。

\medskip\noindent
一般情形下存在性的证明。
由绝对诺特逼近，可以找到笛卡尔图
$$
\xymatrix{
X \ar[r] \ar[d]^f & X_0 \ar[d]^{f_0} \\
S \ar[r]^g & S_0
}
$$
其中 $S_0$ 是诺特概形，$f_0$ 是光滑固有态射，且其所有纤维均非空并且
等维，维数为 $n$。见极限，引理 \ref{limits-lemma-descend-finite-presentation}、
\ref{limits-lemma-descend-smooth}、\ref{limits-lemma-descend-dimension-d}、
\ref{limits-lemma-descend-surjective} 和
\ref{limits-lemma-eventually-proper}。选择陈述中那样的映射
$$
t_0 : Rf_{0, *}\Omega^n_{X_0/S_0}[n] \longrightarrow \mathcal{O}_{S_0}
$$
。由上同调与基变换，复形 $Rf_*(\Omega^p_{X/S})$ 等于
$Lg^*(Rf_{0, *}\Omega^p_{X_0/S_0})$，这对每个 $p \geq 0$ 均成立；见引理
\ref{derham-lemma-proper-smooth-de-Rham}。特别地，$t_0$ 沿 $g$ 的拉回给出
本引理陈述中的映射 $t$。具体地，由上同调，引理
\ref{cohomology-lemma-base-change-cup-product}，杯积映射
$$
Rf_*\Omega^p_{X/S}
\otimes_{\mathcal{O}_S}^\mathbf{L}
Rf_*\Omega^{n - p}_{X/S}[n]
\longrightarrow
Rf_*\Omega^n_{X/S}
$$
是沿 $g$ 拉回 $f_0$ 的相应杯积映射。因此，$t_0$ 给出完美配对这一事实
拉回后说明 $t$ 在 $S$ 上具有同一性质。

\medskip\noindent
映射 $t$ 的唯一性。选择可区别三角
$f_*\mathcal{O}_X \to Rf_*\mathcal{O}_X \to P \to f_*\mathcal{O}_X[1]$.
。由引理 \ref{derham-lemma-relative-bottom-part-degenerates}，对象 $P$ 是
Tor 振幅落在 $[1, \infty)$ 中的完美对象，并且该三角在 $S$ 上局部分裂。
因此 $R\SheafHom_{\mathcal{O}_X}(P, \mathcal{O}_X)$ 是 Tor 振幅落在
$(-\infty, -1]$ 中的完美对象。于是上述对偶性表明，在 $S$ 上局部地有
$$
Rf_*\Omega^n_{X/S}[n] \cong
R\SheafHom_{\mathcal{O}_S}(f_*\mathcal{O}_X, \mathcal{O}_S)
\oplus R\SheafHom_{\mathcal{O}_X}(P, \mathcal{O}_X)
$$
这说明 $R^nf_*\Omega^n_{X/S}$ 是有限局部自由的，并且得到完美的
$\mathcal{O}_S$-双线性配对
$$
f_*\mathcal{O}_X \times R^nf_*\Omega^n_{X/S} \longrightarrow \mathcal{O}_S
$$
；这里使用了 $t$。这蕴含任何 $\mathcal{O}_S$-线性映射
$t' : R^nf_*\Omega^n_{X/S} \to \mathcal{O}_S$ 都具有形式
$t' = t \circ g$，其中某个
$g \in \Gamma(S, f_*\mathcal{O}_X) = \Gamma(X, \mathcal{O}_X)$。
要使 $t'$ 仍确定完美配对，$g$ 必须是单位。证明完毕。
\end{proof}

\begin{lemma}
\label{derham-lemma-relative-top-part-degenerates}
在情形 \ref{derham-situation-relative-duality} 中，映射
$\text{d} : R^nf_*\Omega^{n - 1}_{X/S} \to R^nf_*\Omega^n_{X/S}$
为零。
\end{lemma}

\noindent
正如我们在引理 \ref{derham-lemma-top-part-degenerates} 的证明中提到的，
本引理并非引理
\ref{derham-lemma-relative-duality-hodge} 和
\ref{derham-lemma-relative-bottom-part-degenerates} 的简单推论。

\begin{proof}[$S$ 既约时的证明]
假设 $S$ 既约。注意
$\text{d} : R^nf_*\Omega^{n - 1}_{X/S} \to R^nf_*\Omega^n_{X/S}$
是（拟凝聚）$\mathcal{O}_S$-模之间的 $\mathcal{O}_S$-线性映射。
$\mathcal{O}_S$-模 $R^nf_*\Omega^n_{X/S}$ 是有限局部自由的
（由引理
\ref{derham-lemma-relative-duality-hodge} 和
\ref{derham-lemma-relative-bottom-part-degenerates}，它是有限局部自由
$\mathcal{O}_S$-模 $f_*\mathcal{O}_X$ 的对偶）。由于 $S$ 既约，
只需证明 $\text{d}$ 在每个一般点 $\eta \in S$ 处的茎为零；
这可通过考察仿射开集上的截面、利用 $\text{d}$ 的目标局部自由以及代数，
引理 \ref{algebra-lemma-reduced-ring-sub-product-fields} 的第 (2) 部分得出。
由于 $S$ 既约，有 $\mathcal{O}_{S, \eta} = \kappa(\eta)$；见代数，引理
\ref{algebra-lemma-minimal-prime-reduced-ring}。因此 $\text{d}_\eta$ 等同于映射
$$
\text{d} :
H^n(X_\eta, \Omega^{n - 1}_{X_\eta/\kappa(\eta)})
\longrightarrow
H^n(X_\eta, \Omega^n_{X_\eta/\kappa(\eta)})
$$
，它由引理 \ref{derham-lemma-top-part-degenerates} 为零。
\end{proof}

\begin{proof}[一般情形下的证明]
注意该问题在 $S$ 上对平坦拓扑是局部的：若 $S' \to S$ 是满的平坦概形态射，
并且该映射拉回到 $S'$ 后为零，则原映射为零。此外，由平坦基变换（概形的上同调，
引理 \ref{coherent-lemma-flat-base-change-cohomology}），该映射的形成与沿平坦态射的
基变换交换。

\medskip\noindent
考虑态射进阶，定理
\ref{more-morphisms-theorem-stein-factorization-general} 所述的 Stein 分解
$X \to S' \to S$。由引理 \ref{derham-lemma-relative-bottom-part-degenerates}，态射
$\pi : S' \to S$ 有限平展。态射 $f : X \to S'$ 固有（由该定理），
光滑（由态射进阶，引理
\ref{more-morphisms-lemma-smooth-etale-permanence}），并且由 Stein 分解定理，
其纤维几何连通。在引理 \ref{derham-lemma-relative-bottom-part-degenerates} 的证明中，
我们已看到 $\Omega_{X/S} = \Omega_{X/S'}$，因为 $S' \to S$ 平展。
因此 $\Omega^\bullet_{X/S} = \Omega^\bullet_{X/S'}$。我们有
$$
R^qf_*\Omega^p_{X/S} = \pi_*R^qf'_*\Omega^p_{X/S'}
$$
，对所有 $p, q$ 都成立。这由 Leray 谱序列（上同调，引理
\ref{cohomology-lemma-relative-Leray}）、$\pi$ 有限因而仿射，以及概形的上同调，
引理 \ref{coherent-lemma-relative-affine-vanishing} 得出（当然还使用了
$R^qf'_*\Omega^p_{X'/S}$ 拟凝聚）。因此，本引理中的映射是将 $\pi_*$ 作用于
$\text{d} : R^nf'_*\Omega^{n - 1}_{X/S'} \to R^nf'_*\Omega^n_{X/S'}$ 所得的映射。
换言之，为证明本引理，可以把 $f : X \to S$ 替换为 $f' : X \to S'$，
从而归约到下一段所讨论的情形。

\medskip\noindent
假设 $f$ 的纤维几何连通，并且 $f_*\mathcal{O}_X = \mathcal{O}_S$。
对每个 $s \in S$，可以选择平展邻域 $(S', s') \to (S, s)$，
使得该基变换 $X' \to S'$ 作为 $S$ 的基变换具有截面；见态射进阶，引理
\ref{more-morphisms-lemma-etale-nbhd-dominates-smooth}。
由证明开头的说明，这又归约到下一段所讨论的情形。

\medskip\noindent
假设 $f$ 的纤维几何连通，$f_*\mathcal{O}_X = \mathcal{O}_S$，并且有截面
$s : S \to X$，它是 $f$ 的截面。我们可以并且将会假设 $S = \Spec(A)$ 仿射。映射
$s^* : R\Gamma(X, \mathcal{O}_X) \to R\Gamma(S, \mathcal{O}_S) = A$
是映射 $A \to R\Gamma(X, \mathcal{O}_X)$ 的分裂。因此可写成
$$
R\Gamma(X, \mathcal{O}_X) = A \oplus P
$$
，其中 $P$ 是 $s^*$ 的“核”。由引理
\ref{derham-lemma-relative-bottom-part-degenerates}，对象 $P$（属于 $D(A)$）是
Tor 振幅落在 $[1, n]$ 中的完美对象。与引理
\ref{derham-lemma-relative-duality-hodge} 的证明相同，可见
$H^n(X, \Omega^n_{X/S})$ 是局部自由 $A$-模，其秩为 $1$
（事实上它与 $A$ 对偶，因而是秩为 $1$ 的自由模；我们很快会选择一个生成元，
但既不想核对它是否为同一个生成元，也无需如此）。

\medskip\noindent
用 $Z \subset X$ 表示 $s$ 的像；由概形，引理
\ref{schemes-lemma-section-immersion}，它是 $X$ 的闭子概形。注意，由除子，引理
\ref{divisors-lemma-regular-quasi-regular-immersion}，$Z \to X$ 是正则闭浸入
（从而更是 Koszul 正则闭浸入）；这由除子，引理
\ref{divisors-lemma-section-smooth-regular-immersion} 得出。当然，
$Z \to X$ 的余维数为 $n$。因此，由备注 \ref{derham-remark-how-to-use}，
可以考虑映射
$$
\gamma^{0, 0} : H^0(Z, \Omega^0_{Z/S}) \longrightarrow H^n(X, \Omega^n_{X/S})
$$
，并置 $\xi = \gamma^{0, 0}(1) \in H^n(X, \Omega^n_{X/S})$。

\medskip\noindent
我们断言 $\xi$ 是一个基元素。事实上，由于最高次数满足基变换（例如见极限，
引理 \ref{limits-lemma-higher-direct-images-zero-above-dimension-fibre}），可见
$H^n(X, \Omega^n_{X/S}) \otimes_A k = H^n(X_k, \Omega^n_{X_k/k})$
，这对任意环映射 $A \to k$ 均成立。由 $\xi$ 的构造与基变换相容（见引理
\ref{derham-lemma-gysin-transverse-global}），可见 $\xi$ 在
$H^n(X_k, \Omega^n_{X_k/k})$ 中的像由引理 \ref{derham-lemma-class-of-a-point} 非零，
只要 $k$ 是域。
因此 $\xi$ 是可逆模的处处不消失截面，从而是生成元。

\medskip\noindent
取 $\theta \in H^n(X, \Omega^{n - 1}_{X/S})$。我们必须证明
$\text{d}(\theta)$ 在 $H^n(X, \Omega^n_{X/S})$ 中为零。
可以写成 $\text{d}(\theta) = a \xi$，其中某个 $a \in A$，
因为 $\xi$ 是基元素。于是必须证明 $a = 0$。

\medskip\noindent
考虑闭浸入
$$
\Delta : X \to X \times_S X
$$
。它也是某个光滑态射（即任一投影）的截面，因而也是余维数为 $n$ 的正则且
Koszul 正则浸入。因此可以考虑备注 \ref{derham-remark-how-to-use} 中的映射
$$
\gamma^{p, q} :
H^q(X, \Omega^p_{X/S})
\longrightarrow
H^{q + n}(X \times_S X, \Omega^{p + n}_{X \times_S X/S})
$$
。考虑像
$$
\gamma^{n - 1, n}(\theta) \in
H^{2n}(X \times_S X, \Omega^{2n - 1}_{X \times_S X})
$$
。由引理 \ref{derham-lemma-de-rham-complex-product} 有
$$
\Omega^{2n - 1}_{X \times_S X} =
\Omega^{n - 1}_{X/S} \boxtimes \Omega^n_{X/S} \oplus
\Omega^n_{X/S} \boxtimes \Omega^{n - 1}_{X/S}
$$
由 K\"unneth 公式（可用概形的导出范畴，引理 \ref{perfect-lemma-kunneth}，
或概形的导出范畴，引理 \ref{perfect-lemma-kunneth-single-sheaf}）可见
$$
H^{2n}(X \times_S X, \Omega^{n - 1}_{X/S} \boxtimes \Omega^n_{X/S}) =
H^n(X, \Omega^{n - 1}_{X/S}) \otimes_A H^n(X, \Omega^n_{X/S})
$$
以及
$$
H^{2n}(X \times_S X, \Omega^n_{X/S} \boxtimes \Omega^{n - 1}_{X/S}) =
H^n(X, \Omega^n_{X/S}) \otimes_A H^n(X, \Omega^{n - 1}_{X/S})
$$
。事实上，由于我们考察的是最高次数，不会有更高的 Tor 群介入。再结合
$\xi$ 是生成元这一事实，可以写成
$$
\gamma^{n - 1, n}(\theta) = \theta_1 \otimes \xi + \xi \otimes \theta_2
$$
，其中 $\theta_1, \theta_2 \in H^n(X, \Omega^{n - 1}_{X/S})$。
完全相同的论证给出
$$
\gamma^{n, n}(\xi) = b \xi \otimes \xi
$$
，它位于
$H^{2n}(X \times_S X, \Omega^{2n}_{X \times_S X/S}) =
H^n(X, \Omega^n_{X/S}) \otimes_A H^n(X, \Omega^n_{X/S})$
中，其中某个 $b \in H^0(S, \mathcal{O}_S)$。

\medskip\noindent
{\bf 断言：} $\theta_1 = \theta$、$\theta_2 = \theta$ 且 $b = 1$。
先说明该断言蕴含所需结论 $a = 0$。事实上，由引理
\ref{derham-lemma-gysin-differential-hodge} 有
$$
\gamma^{n, n}(\text{d}(\theta)) = \text{d}(\gamma^{n - 1, n}(\theta))
$$
由上面的选择，这给出
$$
a \xi \otimes \xi =
\gamma^{n, n}(a\xi) =
\text{d}(\theta \otimes \xi + \xi \otimes \theta) =
a \xi \otimes \xi + (-1)^n a \xi \otimes \xi
$$
最右边的等式源于如下事实：映射
$\text{d} : \Omega^{2n - 1}_{X \otimes_S X/S} \to \Omega^{2n}_{X \times_S X/S}$
由引理 \ref{derham-lemma-de-rham-complex-product} 是微分
$\text{d} \boxtimes 1 : \Omega^{n - 1}_{X/S} \boxtimes \Omega^n_{X/S}
\to \Omega^n_{X/S} \boxtimes \Omega^n_{X/S}$
与微分
$(-1)^n 1 \boxtimes \text{d} : \Omega^n_{X/S} \boxtimes \Omega^{n - 1}_{X/S}
\to \Omega^n_{X/S} \boxtimes \Omega^n_{X/S}$ 之和。更多说明见第
\ref{derham-section-kunneth} 节和概形的导出范畴，第
\ref{perfect-section-kunneth-complexes} 节。由于 $\xi \otimes \xi$ 是秩为
$1$ 的自由 $A$-模
$H^n(X, \Omega^n_{X/S}) \otimes_A H^n(X, \Omega^n_{X/S})$
的一组基，故有
$$
a = a + (-1)^n a \Rightarrow a = 0
$$
，正是所需结论。

\medskip\noindent
在证明的其余部分，我们证明上述断言。记
$\eta = \gamma^{0, 0}(1) \in H^n(X \times_S X, \Omega^n_{X \times_S X/S})$。
由于 $\Omega^n_{X \times_S X/S} =
\bigoplus_{p + p' = n} \Omega^p_{X/S} \boxtimes \Omega^{p'}_{X/S}$，可写成
$$
\eta = \eta_0 + \eta_1 + \ldots + \eta_n
$$
，其中 $\eta_p$ 属于
$H^n(X \times_S X, \Omega^p_{X/S} \boxtimes \Omega^{n - p}_{X/S})$.
当 $p = 0$ 时，可写成
\begin{align*}
H^n(X \times_S X, \mathcal{O}_X \boxtimes \Omega^n_{X/S})
& =
H^n(R\Gamma(X, \mathcal{O}_X) \otimes_A^\mathbf{L}
R\Gamma(X, \Omega^n_{X/S})) \\
& =
A \otimes_A H^n(X, \Omega^n_{X/S}) \oplus
H^n(P \otimes_A^\mathbf{L} R\Gamma(X, \Omega^n_{X/S}))
\end{align*}
；这里使用了先前给出的分解 $R\Gamma(X, \mathcal{O}_X) = A \oplus P$。
考虑态射 $(s, \text{id}) : X \to X \times_S X$。于是概形论意义下有
$(s, \text{id})^{-1}(\Delta) = Z$。因此，由引理
\ref{derham-lemma-gysin-transverse-global} 可见 $(s, \text{id})^*\eta = \xi$。这意味着
$$
\xi = (s, \text{id})^*\eta = (s^* \otimes \text{id})(\eta_0)
$$
这恰好意味着 $\eta_0$ 在上述直和分解中的第一个分量是 $\xi$。换言之，可以写成
$$
\eta_0 = 1 \otimes \xi + \eta'_0
$$
，其中 $\eta'_0 \in H^n(P \otimes_A^\mathbf{L} R\Gamma(X, \Omega^n_{X/S}))$。
完全相同地，当 $p = n$ 时，可以写成
\begin{align*}
H^n(X \times_S X, \Omega^n_{X/S} \boxtimes \mathcal{O}_X)
& =
H^n(R\Gamma(X, \Omega^n_{X/S}) \otimes_A^\mathbf{L}
R\Gamma(X, \mathcal{O}_X)) \\
& =
H^n(X, \Omega^n_{X/S}) \otimes_A A \oplus
H^n(R\Gamma(X, \Omega^n_{X/S}) \otimes_A^\mathbf{L} P)
\end{align*}
并且可以写成
$$
\eta_n = \xi \otimes 1 + \eta'_n
$$
，其中 $\eta'_n \in H^n(R\Gamma(X, \Omega^n_{X/S}) \otimes_A^\mathbf{L} P)$。

\medskip\noindent
注意，$\text{pr}_1^*\theta = \theta \otimes 1$ 和
$\text{pr}_2^*\theta = 1 \otimes \theta$ 是 $X \times_S X$ 上的霍奇上同调类，
它们拉回为 $\theta$，这里的拉回沿 $\Delta$ 进行。因此，由引理
\ref{derham-lemma-gysin-projection-global} 有
$$
\theta_1 \otimes \xi + \xi \otimes \theta_2 =
\gamma^{n - 1, n}(\theta) =
(\theta \otimes 1) \cup \eta =
(1 \otimes \theta) \cup \eta
$$
，这是在 $X \times_S X$ 相对于 $S$ 的霍奇上同调环中成立。按照
$X \times_S X/S$ 的微分模上的直和分解，得到
$$
\theta_1 \otimes \xi =
(\theta \otimes 1) \cup \eta_0
\quad\text{且}\quad
\xi \otimes \theta_2 =
(1 \otimes \theta) \cup \eta_n
$$
考察上面所得公式 $\eta_0 = 1 \otimes \xi + \eta'_0$，可见要证明
$\theta_1 = \theta$，只需证明
$$
(\theta \otimes 1) \cup \eta'_0 = 0
$$
为此，注意与 $\theta \otimes 1$ 作杯积由如下映射在上同调上的作用给出：
$$
(P \otimes_A^\mathbf{L} R\Gamma(X, \Omega^n_{X/S}))[-n]
\xrightarrow{\theta \otimes 1}
R\Gamma(X, \Omega^{n - 1}_{X/S}) \otimes_A^\mathbf{L}
R\Gamma(X, \Omega^n_{X/S})
$$
，这是在导出范畴中；见上同调，备注
\ref{cohomology-remark-cup-with-element-map-total-cohomology}。该映射是以下两个映射的
导出张量积：
$$
\theta : P[-n] \to R\Gamma(X, \Omega^{n - 1}_{X/S})
\quad\text{且}\quad
1 : R\Gamma(X, \Omega^n_{X/S}) \to R\Gamma(X, \Omega^n_{X/S})
$$
；见概形的导出范畴，备注
\ref{perfect-remark-annoying-compatibility}。然而，其中第一个映射在 $D(A)$ 中为零，
因为它从 Tor 振幅落在 $[n + 1, 2n]$ 中的完美复形映到一个仅在次数
$0, 1, \ldots, n$ 具有上同调的复形；见代数进阶，引理
\ref{more-algebra-lemma-splitting-unique}。类似论证表明
$(1 \otimes \theta) \cup \eta'_n$ 也为零。最后，完全相同地得到
$$
b \xi \otimes \xi = \gamma^{n, n}(\xi) = (\xi \otimes 1) \cup \eta_0
$$
，并像前面一样，通过以同样方式证明
$(\xi \otimes 1) \cup \eta'_0 = 0$ 而得到结论。证明完毕。
\end{proof}

\begin{proposition}
\label{derham-proposition-relative-poincare-duality}
设 $S$ 是拟紧拟分离概形。设 $f : X \to S$ 是光滑固有的概形态射，
其所有纤维均非空且等维，维数为 $n$。存在 $\mathcal{O}_S$-模映射
$$
t : R^{2n}f_*\Omega^\bullet_{X/S} \longrightarrow \mathcal{O}_S
$$
它在预复合以 $H^0(X, \mathcal{O}_X)$ 中某个单位的乘法这个意义下唯一，
并具有如下性质：配对
$$
Rf_*\Omega^\bullet_{X/S}
\otimes_{\mathcal{O}_S}^\mathbf{L}
Rf_*\Omega^\bullet_{X/S}[2n]
\longrightarrow
\mathcal{O}_S, \quad
(\xi, \xi') \longmapsto t(\xi \cup \xi')
$$
是 $S$ 上完美复形之间的完美配对。
\end{proposition}

\begin{proof}
本证明与命题 \ref{derham-proposition-poincare-duality} 的证明完全相同。

\medskip\noindent
由相对霍奇--德拉姆谱序列
$$
E_1^{p, q} = R^qf_*\Omega^p_{X/S} \Rightarrow R^{p + q}f_*\Omega^\bullet_{X/S}
$$
（第 \ref{derham-section-hodge-to-de-rham} 节）、$\Omega^i_{X/S}$ 对 $i > n$ 的消失、
例如极限，引理 \ref{limits-lemma-higher-direct-images-zero-above-dimension-fibre}
中的消失，以及引理 \ref{derham-lemma-relative-bottom-part-degenerates} 和
\ref{derham-lemma-relative-top-part-degenerates} 的结论，可见
$R^0f_*\Omega_{X/S} = R^0f_*\mathcal{O}_X$ 且
$R^nf_*\Omega^n_{X/S} = R^{2n}f_*\Omega^\bullet_{X/S}$。
更确切地说，这些等同来自复形映射
$$
\Omega^\bullet_{X/S} \to \mathcal{O}_X[0]
\quad\text{且}\quad
\Omega^n_{X/S}[-n] \to \Omega^\bullet_{X/S}
$$
选择 $t : R^{2n}f_*\Omega_{X/S} \to \mathcal{O}_S$，使它经由此等同对应于
引理 \ref{derham-lemma-relative-duality-hodge} 中的一个 $t$。

\medskip\noindent
简记 $\Omega^\bullet = \Omega^\bullet_{X/S}$。考虑映射
(\ref{derham-equation-wedge})，它在当前情形中写作
$$
\wedge :
\text{Tot}(\Omega^\bullet \otimes_{f^{-1}\mathcal{O}_S} \Omega^\bullet)
\longrightarrow
\Omega^\bullet
$$
对每个整数 $p = 0, 1, \ldots, n$，由于次数的原因，该映射在子复形
$\text{Tot}(\sigma_{> p} \Omega^\bullet \otimes_{f^{-1}\mathcal{O}_S}
\sigma_{\geq n - p} \Omega^\bullet)$ 上为零。因此，$\wedge$ 在子复形
$\text{Tot}(\Omega^\bullet \otimes_{f^{-1}\mathcal{O}_S}
\geq_{n - p}\Omega^\bullet)$ 上的限制分解经过一个复形映射
$$
\gamma_p :
\text{Tot}(\sigma_{\leq p} \Omega^\bullet \otimes_{f^{-1}\mathcal{O}_S}
\sigma_{\geq n - p} \Omega^\bullet)
\longrightarrow
\Omega^\bullet
$$
采用第 \ref{derham-section-cup-product} 节中的相同步骤，得到相对杯积
$$
Rf_*\sigma_{\leq p} \Omega^\bullet
\otimes_{\mathcal{O}_S}^\mathbf{L}
Rf_*\sigma_{\geq n - p}\Omega^\bullet
\longrightarrow
Rf_*\Omega^\bullet
$$
我们将对 $p$ 归纳证明：这些杯积经由 $t$ 在
$Rf_*\sigma_{\leq p} \Omega^\bullet$ 与
$Rf_*\sigma_{\geq n - p}\Omega^\bullet[2n]$ 之间诱导完美配对。当 $p = n$ 时，
这正是命题的断言。

\medskip\noindent
基础情形是 $p = 0$。此时有
$$
Rf_*\sigma_{\leq p}\Omega^\bullet = Rf_*\mathcal{O}_X
\quad\text{且}\quad
Rf_*\sigma_{\geq n - p}\Omega^\bullet[2n] = Rf_*(\Omega^n[-n])[2n] =
Rf_*\Omega^n[n]
$$
此时所得恰好是引理 \ref{derham-lemma-relative-duality-hodge} 中
$Rf_*\mathcal{O}_X$ 与 $Rf_*\Omega^n[n]$ 之间的配对，故结论成立。

\medskip\noindent
归纳步骤。假设已知结论对 $p$ 成立。考虑可区别三角
$$
\Omega^{p + 1}[-p - 1] \to
\sigma_{\leq p + 1}\Omega^\bullet \to
\sigma_{\leq p}\Omega^\bullet \to
\Omega^{p + 1}[-p]
$$
以及可区别三角
$$
\sigma_{\geq n - p}\Omega^\bullet \to
\sigma_{\geq n - p - 1}\Omega^\bullet \to
\Omega^{n - p - 1}[-n + p + 1] \to
(\sigma_{\geq n - p}\Omega^\bullet)[1]
$$
注意，两者都是 $f^{-1}\mathcal{O}_S$-模层复形的同伦范畴中的可区别三角；
特别地，映射 $\sigma_{\leq p}\Omega^\bullet \to \Omega^{p + 1}[-p]$ 和
$\Omega^{n - p - 1}[-d + p + 1] \to (\sigma_{\geq n - p}\Omega^\bullet)[1]$
由实际的复形映射给出，具体使用微分 $\Omega^p \to \Omega^{p + 1}$ 和微分
$\Omega^{n - p - 1} \to \Omega^{n - p}$。考虑对这些可区别三角作用
$Rf_*$ 所得的相应可区别三角
$$
\xymatrix{
Rf_*\sigma_{\leq p}\Omega^\bullet \ar[d]_a \\
Rf_*\Omega^{p + 1}[-p - 1] \ar[d]_b \\
Rf_*\sigma_{\leq p + 1}\Omega^\bullet \ar[d]_c \\
Rf_*\sigma_{\leq p}\Omega^\bullet \ar[d]_d \\
Rf_*\Omega^{p + 1}[-p - 1]
}
\quad\quad
\xymatrix{
Rf_*\sigma_{\geq n - p}\Omega^\bullet \\
Rf_*\Omega^{n - p - 1}[-n + p + 1] \ar[u]_{a'} \\
Rf_*\sigma_{\geq n - p - 1}\Omega^\bullet \ar[u]_{b'} \\
Rf_*\sigma_{\geq n - p}\Omega^\bullet \ar[u]_{c'} \\
Rf_*\Omega^{n - p - 1}[-n + p + 1] \ar[u]_{d'}
}
$$
下文将证明各对 $(a, a')$、$(b, b')$、$(c, c')$ 和 $(d, d')$
与给定配对相容。这意味着我们得到一个从左侧可区别三角到如下可区别三角的映射：
对右侧可区别三角作用 $R\SheafHom(-, \mathcal{O}_S)$ 所得的可区别三角。
由归纳假设和引理 \ref{derham-lemma-duality-hodge}，上面在第一、第二、第四和第五行的
复形之间构造的配对是完美的；也就是说，对偶后它们确定同构。由导出范畴，
引理 \ref{derived-lemma-third-isomorphism-triangle}，中间一行复形之间的配对
也是完美的，正如所需。

\medskip\noindent
设 $e : K \to K'$ 和 $e' : M' \to M$ 是 $D(\mathcal{O}_S)$ 中对象的映射，
并设 $K \otimes_{\mathcal{O}_S}^\mathbf{L} M \to \mathcal{O}_S$ 和
$K' \otimes_{\mathcal{O}_S}^\mathbf{L} M' \to \mathcal{O}_S$ 是配对。
如果下图交换，就称这些配对相容：
$$
\xymatrix{
K' \otimes_{\mathcal{O}_S}^\mathbf{L} M' \ar[d] &
K \otimes_{\mathcal{O}_S}^\mathbf{L} M'
\ar[l]^{e \otimes 1} \ar[d]^{1 \otimes e'} \\
\mathcal{O}_S &
K \otimes_{\mathcal{O}_S}^\mathbf{L} M \ar[l]
}
$$
事实上，这意味着下图交换：
$$
\xymatrix{
K \ar[r] \ar[d]_e & R\SheafHom(M, \mathcal{O}_S) \ar[d]^{R\SheafHom(e', -)} \\
K' \ar[r] & R\SheafHom(M', \mathcal{O}_S)
}
$$
，所以足以满足我们的需要。

\medskip\noindent
先对 $(c, c')$ 证明这一点。这里只需注意有交换图
$$
\xymatrix{
\text{Tot}(\sigma_{\leq p} \Omega^\bullet \otimes_{f^{-1}\mathcal{O}_S}
\sigma_{\geq n - p} \Omega^\bullet) \ar[d]_{\gamma_p} &
\text{Tot}(\sigma_{\leq p + 1} \Omega^\bullet \otimes_{f^{-1}\mathcal{O}_S}
\sigma_{\geq n - p} \Omega^\bullet) \ar[l] \ar[d] \\
\Omega^\bullet &
\text{Tot}(\sigma_{\leq p + 1} \Omega^\bullet \otimes_{f^{-1}\mathcal{O}_S}
\sigma_{\geq n - p - 1} \Omega^\bullet) \ar[l]_-{\gamma_{p + 1}}
}
$$
由杯积的函子性，便得到所需图的交换性。

\medskip\noindent
类似地，对 $(b, b')$ 使用交换图
$$
\xymatrix{
\text{Tot}(\sigma_{\leq p + 1} \Omega^\bullet \otimes_{f^{-1}\mathcal{O}_S}
\sigma_{\geq n - p - 1} \Omega^\bullet) \ar[d]_{\gamma_{p + 1}} &
\text{Tot}(\Omega^{p + 1}[-p - 1] \otimes_{f^{-1}\mathcal{O}_S}
\sigma_{\geq n - p - 1} \Omega^\bullet) \ar[l] \ar[d] \\
\Omega^\bullet &
\Omega^{p + 1}[-p - 1]
\otimes_{f^{-1}\mathcal{O}_S}
\Omega^{n - p - 1}[-n + p + 1] \ar[l]_-\wedge
}
$$

\medskip\noindent
对 $(d, d')$ 和 $(a, a')$，使用交换图
$$
\xymatrix{
\Omega^{p + 1}[-p] \otimes_{f^{-1}\mathcal{O}_S}
\Omega^{n - p - 1}[-n + p] \ar[d] &
\text{Tot}(\sigma_{\leq p}\Omega^\bullet \otimes_{f^{-1}\mathcal{O}_S}
\Omega^{n - p - 1}[-n + p]) \ar[l] \ar[d] \\
\Omega^\bullet &
\text{Tot}(\sigma_{\leq p}\Omega^\bullet \otimes_{f^{-1}\mathcal{O}_S}
\sigma_{\geq n - p}\Omega^\bullet) \ar[l]
}
$$

\medskip\noindent
关于 $t$ 在预复合以 $H^0(X, \mathcal{O}_X)$ 中某个单位的乘法这个意义下
唯一的论证，从略。
\end{proof}
% Generated mechanically from the frozen local-cohomology.tex target.
% Do not edit this generated file; edit the bound locale source instead.

% OK, start here.
%

\setcounter{chapter}{50}
\chapter{局部上同调}



\phantomsection
\label{local-cohomology-section-phantom}


\section{引言}
\label{local-cohomology-section-introduction}

\noindent
本章继续研究局部上同调。参考文献之一是 \cite{SGA2}。
局部上同调的定义见《对偶复形》，第
\ref{dualizing-section-local-cohomology} 节。对诺特环而言，取局部上同调
等同于导出某个适当的挠函子；见《对偶复形》，第
\ref{dualizing-section-local-cohomology-noetherian} 节。
它与深度的关系见《对偶复形》，第
\ref{dualizing-section-depth} 节。

\medskip\noindent
我们将讨论局部上同调的有限性，从而证明 Grothendieck 有限性定理的一个相当一般的
版本；见定理 \ref{local-cohomology-theorem-finiteness} 和引理 \ref{local-cohomology-lemma-finiteness-Rjstar}
（凝聚模沿开浸入的高阶正像）。所用方法吸收了 Faltings 论文中的若干精妙论证；
见 \cite{Faltings-annulators} 和 \cite{Faltings-finiteness}。

\medskip\noindent
作为应用，我们将讨论 Hartshorne--Lichtenbaum 消失，还将讨论 Frobenius 与微分算子
在局部上同调上的作用。




\section{一般性质}
\label{local-cohomology-section-generalities}

\noindent
下一个引理说明函子 $R\Gamma_Z$ 与带支撑上同调之间的关系。

\begin{lemma}
\label{local-cohomology-lemma-local-cohomology-is-local-cohomology}
设 $A$ 是环，$I$ 是有限生成理想。置 $Z = V(I) \subset X = \Spec(A)$。
对 $K \in D(A)$，经由《概形的导出范畴》，引理
\ref{perfect-lemma-affine-compare-bounded}，它对应于
$\widetilde{K} \in D_\QCoh(\mathcal{O}_X)$；此时存在函子性同构
$$
R\Gamma_Z(K) = R\Gamma_Z(X, \widetilde{K})
$$
其中左边采用《对偶复形》中的等式
(\ref{dualizing-equation-local-cohomology})，右边采用《上同调》，第
\ref{cohomology-section-cohomology-support-bis} 节中的函子。
\end{lemma}

\begin{proof}
由《上同调》，引理
\ref{cohomology-lemma-triangle-sections-with-support}，存在可区别三角
$$
R\Gamma_Z(X, \widetilde{K})
\to R\Gamma(X, \widetilde{K})
\to R\Gamma(U, \widetilde{K})
\to R\Gamma_Z(X, \widetilde{K})[1]
$$
，其中 $U = X \setminus Z$。由《概形的导出范畴》，引理
\ref{perfect-lemma-affine-compare-bounded}，我们知道
$R\Gamma(X, \widetilde{K}) = K$。设 $I = (f_1, \ldots, f_r)$，便得到有限仿射开覆盖
$\mathcal{U} : U = D(f_1) \cup \ldots \cup D(f_r)$。由《概形的导出范畴》，
引理 \ref{perfect-lemma-alternating-cech-complex-complex-computes-cohomology}，
交错 {\v C}ech 复形
$\text{Tot}(\check{\mathcal{C}}_{alt}^\bullet(\mathcal{U},
\widetilde{K^\bullet}))$
计算 $R\Gamma(U, \widetilde{K})$，其中 $K^\bullet$ 是任意的 $A$-模复形，
它表示 $K$。逐一展开定义可得
$$
R\Gamma(U, \widetilde{K}) =
\text{Tot}\left(
K^\bullet \otimes_A
(\prod\nolimits_{i_0} A_{f_{i_0}} \to
\prod\nolimits_{i_0 < i_1} A_{f_{i_0}f_{i_1}} \to
\ldots \to A_{f_1\ldots f_r})\right)
$$
显然，
$K^\bullet = R\Gamma(X, \widetilde{K^\bullet}) \to
R\Gamma(U, \widetilde{K}^\bullet)$
由从 $A$ 到 $\prod A_{f_i}$ 的对角映射诱导。因此得到
$$
R\Gamma_Z(X, \widetilde{K}) =
\text{Tot}\left(
K^\bullet \otimes_A
(A \to \prod\nolimits_{i_0} A_{f_{i_0}} \to
\prod\nolimits_{i_0 < i_1} A_{f_{i_0}f_{i_1}} \to
\ldots \to A_{f_1\ldots f_r})\right)
$$
由《对偶复形》，引理 \ref{dualizing-lemma-local-cohomology-adjoint}，
右边的复形计算 $R\Gamma_Z(K)$，故本引理对给定的 $K$ 成立。

\medskip\noindent
最后证明可以使该同构关于 $K$ 具有函子性。考虑 $A$-模复形的映射
$$
M^\bullet = \text{Tot}\left(
K^\bullet \otimes_A
(A \to \prod\nolimits_{i_0} A_{f_{i_0}} \to
\prod\nolimits_{i_0 < i_1} A_{f_{i_0}f_{i_1}} \to
\ldots \to A_{f_1\ldots f_r})\right)
\longrightarrow
K^\bullet
$$
上述论证表明，对该箭头作用函子 $R\Gamma_Z(X, \widetilde{\ })$ 后得到同构，
并且 $R\Gamma(U, \widetilde{M}^\bullet) = 0$（细节从略）。于是有
$$
R\Gamma_Z(X, \widetilde{K}^\bullet)
\leftarrow
R\Gamma_Z(X, \widetilde{M}^\bullet)
\rightarrow
R\Gamma(X, \widetilde{M}^\bullet) = M^\bullet = R\Gamma_Z(K^\bullet)
$$
，其中所有映射都是 $D(A)$ 中的同构，并且关于 $K^\bullet$ 具有函子性，
正如所需。
\end{proof}

\begin{lemma}
\label{local-cohomology-lemma-local-cohomology}
设 $A$ 是环，$I \subset A$ 是有限生成理想。置
$X = \Spec(A)$、$Z = V(I)$、$U = X \setminus Z$，并令 $j : U \to X$
为包含态射。设 $\mathcal{F}$ 是拟凝聚 $\mathcal{O}_U$-模。则：
\begin{enumerate}
\item 存在 $A$-模 $M$，使得 $\mathcal{F}$ 是 $\widetilde{M}$ 在 $U$ 上的限制；
\item 给定 $M$，有正合列
$$
0 \to H^0_Z(M) \to M \to H^0(U, \mathcal{F}) \to H^1_Z(M) \to 0
$$
以及同构 $H^p(U, \mathcal{F}) = H^{p + 1}_Z(M)$，其中 $p \geq 1$；
\item 可以取 $M = H^0(U, \mathcal{F})$，此时
$H^0_Z(M) = H^1_Z(M) = 0$。
\end{enumerate}
\end{lemma}

\begin{proof}
$M$ 的存在性由《性质》，引理 \ref{properties-lemma-extend-trivial}
以及 $X$ 上的拟凝聚层与 $A$-模对应这一事实得出（《概形》，引理
\ref{schemes-lemma-equivalence-quasi-coherent}）。接着考察可区别三角
$$
R\Gamma_Z(X, \widetilde{M}) \to R\Gamma(X, \widetilde{M}) \to
R\Gamma(U, \widetilde{M}|_U) \to R\Gamma_Z(X, \widetilde{M})[1]
$$
，它来自《上同调》，引理
\ref{cohomology-lemma-triangle-sections-with-support}。由于 $X$ 仿射，
由《概形的上同调》，引理
\ref{coherent-lemma-quasi-coherent-affine-cohomology-zero}，有
$R\Gamma(X, \widetilde{M}) = M$。按 $M$ 的选择，
$\mathcal{F} = \widetilde{M}|_U$，因而得到正合列
$$
0 \to H^0_Z(X, \widetilde{M}) \to M \to H^0(U, \mathcal{F}) \to
H^1_Z(X, \widetilde{M}) \to 0
$$
以及同构 $H^p(U, \mathcal{F}) = H^{p + 1}_Z(X, \widetilde{M})$，其中
$p \geq 1$。由引理 \ref{local-cohomology-lemma-local-cohomology-is-local-cohomology}，
有 $H^i_Z(M) = H^i_Z(X, \widetilde{M})$，对所有 $i$ 均成立。故 (1) 与 (2) 成立。
最后，置 $M' = H^0(U, \mathcal{F})$，可见 $M \to M'$ 的核与余核均为
$I$-幂挠模。因此 $\widetilde{M}|_U \to \widetilde{M'}|_U$ 是同构，
确实可以像 (3) 所述那样使用 $M'$。不言而喻，$H^0_Z(M')$ 与
$H^0_Z(M')$ 二者均为零。
\end{proof}

\begin{lemma}
\label{local-cohomology-lemma-already-torsion}
设 $I, J \subset A$ 是环 $A$ 的有限生成理想。若 $M$ 是 $I$-幂挠模，则典范映射
$$
H^i_{V(I) \cap V(J)}(M) \to H^i_{V(J)}(M)
$$
对所有 $i$ 都是同构。
\end{lemma}

\begin{proof}
利用《对偶复形》，引理 \ref{dualizing-lemma-local-cohomology-ss} 的谱序列，
可归约为陈述 $R\Gamma_I(M) = M$；它立即由《对偶复形》，第
\ref{dualizing-section-local-cohomology} 节中的局部上同调构造得出。
\end{proof}

\begin{lemma}
\label{local-cohomology-lemma-multiplicative}
设 $S \subset A$ 是环 $A$ 的乘法集。设 $M$ 是 $A$-模，并满足
$S^{-1}M = 0$。则 $\colim_{f \in S} H^0_{V(f)}(M) = M$ 且
$\colim_{f \in S} H^1_{V(f)}(M) = 0$。
\end{lemma}

\begin{proof}
关于 $H^0$ 的陈述直接由定义得出。为证明关于 $H^1$ 的陈述，注意
$R\Gamma_{V(f)}$ 与 $H^1_{V(f)}$ 都和余极限交换。因此可假设 $M$
被某个 $f \in S$ 零化。于是 $H^1_{V(ff')}(M) = 0$，对所有
$f' \in S$ 均成立（例如由引理 \ref{local-cohomology-lemma-already-torsion}）。
\end{proof}

\begin{lemma}
\label{local-cohomology-lemma-elements-come-from-bigger}
设 $I \subset A$ 是环 $A$ 的有限生成理想。设 $\mathfrak p$ 是素理想，
$M$ 是 $A$-模。设 $i \geq 0$ 是整数，并考虑映射
$$
\Psi :
\colim_{f \in A, f \not \in \mathfrak p} H^i_{V((I, f))}(M)
\longrightarrow
H^i_{V(I)}(M)
$$
则：
\begin{enumerate}
\item $\Im(\Psi)$ 是在 $H^i_{V(I)}(M)_\mathfrak p$ 中映为零的元素全体；
\item 若 $H^{i - 1}_{V(I)}(M)_\mathfrak p = 0$，则 $\Psi$ 是单射；
\item 若 $H^{i - 1}_{V(I)}(M)_\mathfrak p = H^i_{V(I)}(M)_\mathfrak p = 0$，
则 $\Psi$ 是同构。
\end{enumerate}
\end{lemma}

\begin{proof}
对 $f \in A$ 且 $f \not \in \mathfrak p$，《对偶复形》，引理
\ref{dualizing-lemma-local-cohomology-ss} 的谱序列退化并给出短正合列
$$
0 \to H^1_{V(f)}(H^{i - 1}_{V(I)}(M)) \to
H^i_{V((I, f))}(M) \to H^0_{V(f)}(H^i_{V(I)}(M)) \to 0
$$
这证明了 (1)，而 (2) 由此及引理 \ref{local-cohomology-lemma-multiplicative} 得出。
(3) 是形式推论。
\end{proof}

\begin{lemma}
\label{local-cohomology-lemma-isomorphism}
设 $I \subset I' \subset A$ 是诺特环 $A$ 的有限生成理想。设 $M$ 是
$A$-模，$i \geq 0$ 是整数。考虑映射
$$
\Psi : H^i_{V(I')}(M) \to H^i_{V(I)}(M)
$$
下列陈述成立：
\begin{enumerate}
\item 若 $H^i_{\mathfrak pA_\mathfrak p}(M_\mathfrak p) = 0$ 对所有
$\mathfrak p \in V(I) \setminus V(I')$ 均成立，则 $\Psi$ 是满射；
\item 若 $H^{i - 1}_{\mathfrak pA_\mathfrak p}(M_\mathfrak p) = 0$ 对所有
$\mathfrak p \in V(I) \setminus V(I')$ 均成立，则 $\Psi$ 是单射；
\item 若 $H^i_{\mathfrak pA_\mathfrak p}(M_\mathfrak p) =
H^{i - 1}_{\mathfrak pA_\mathfrak p}(M_\mathfrak p) = 0$ 对所有
$\mathfrak p \in V(I) \setminus V(I')$ 均成立，则 $\Psi$ 是同构。
\end{enumerate}
\end{lemma}

\begin{proof}
(1) 的证明。取 $\xi \in H^i_{V(I)}(M)$。由于 $A$ 诺特，存在最大的理想
$I \subset I'' \subset I'$，使得 $\xi$ 是某个
$\xi'' \in H^i_{V(I'')}(M)$ 的像。若 $V(I'') = V(I')$，则证明完毕。
否则，选择一般点 $\mathfrak p \in V(I'')$，它不属于 $V(I')$。
由假设，有 $H^i_{V(I'')}(M)_\mathfrak p =
H^i_{\mathfrak pA_\mathfrak p}(M_\mathfrak p) = 0$。由引理
\ref{local-cohomology-lemma-elements-come-from-bigger} 可以增大 $I''$，这与其极大性矛盾。

\medskip\noindent
(2) 的证明。设 $\xi' \in H^i_{V(I')}(M)$ 属于 $\Psi$ 的核。
由于 $A$ 诺特，存在最大的理想 $I \subset I'' \subset I'$，使得 $\xi'$
在 $H^i_{V(I'')}(M)$ 中映为零。若 $V(I'') = V(I')$，则证明完毕。
否则，选择一般点 $\mathfrak p  \in V(I'')$，它不属于 $V(I')$。
由假设，有 $H^{i - 1}_{V(I'')}(M)_\mathfrak p =
H^{i - 1}_{\mathfrak pA_\mathfrak p}(M_\mathfrak p) = 0$。由引理
\ref{local-cohomology-lemma-elements-come-from-bigger} 可以增大 $I''$，这与其极大性矛盾。

\medskip\noindent
(3) 由 (1) 和 (2) 形式地得出。
\end{proof}





\section{Hartshorne 连通性引理}
\label{local-cohomology-section-hartshorne-connectedness}

\noindent
本节标题指的是下面这个结果。

\begin{lemma}
\label{local-cohomology-lemma-depth-2-connected-punctured-spectrum}
\begin{reference}
\cite[Proposition 2.1]{Hartshorne-connectedness}
\end{reference}
\begin{slogan}
Hartshorne 连通性
\end{slogan}
设 $A$ 是深度 $\geq 2$ 的诺特局部环。则 $A$、$A^h$ 和 $A^{sh}$ 的
穿孔谱均连通。
\end{lemma}

\begin{proof}
设 $U$ 是 $A$ 的穿孔谱。若 $U$ 不连通，则
$\Gamma(U, \mathcal{O}_U)$ 有非平凡幂等元。但 $A$ 是局部环，因而没有非平凡
幂等元。因此 $A \to \Gamma(U, \mathcal{O}_U)$ 不是同构。由引理
\ref{local-cohomology-lemma-local-cohomology}，可知 $H^0_\mathfrak m(A)$ 或
$H^1_\mathfrak m(A)$ 中至少一个非零。于是由《对偶复形》，引理
\ref{dualizing-lemma-depth}，有 $\text{depth}(A) \leq 1$。
关于 $A^h$ 与 $A^{sh}$ 的结论使用《代数进阶》，引理
\ref{more-algebra-lemma-henselization-depth}。
\end{proof}

\begin{lemma}
\label{local-cohomology-lemma-catenary-S2-equidimensional}
\begin{reference}
\cite[IV Corollary 5.10.9]{EGA}
\end{reference}
设 $A$ 是链状且满足 $(S_2)$ 的诺特局部环。则 $\Spec(A)$ 等维。
\end{lemma}

\begin{proof}
置 $X = \Spec(A)$，并记 $d = \dim(A) = \dim(X)$。在 $X$ 中，考虑不可约分支之并
$X_1$，其中各分支的维数为 $d$；以及不可约分支之并 $X_2$，其中各分支的维数
$< d$。
当然 $X = X_1 \cup X_2$。若 $X_2 = \emptyset$，则本引理成立。否则，
$Z = X_1 \cap X_2$ 是 $X$ 的非空闭子集，因为它至少包含 $X$ 的闭点。
因此可以选取一般点 $z \in Z$，它属于 $Z$ 的某个不可约分支。回忆，
$\mathcal{O}_{Z, z}$ 的谱是 $X$ 中特化到 $z$ 的点所成的集合。由于 $z$
既属于一个维数为 $d$ 的不可约分支，也属于一个维数 $< d$ 的不可约分支，
便得到非平凡特化 $x_1 \leadsto z$ 和 $x_2 \leadsto z$，使得 $x_1$ 与 $x_2$
的闭包具有不同维数。由于 $X$ 链状，只有当特化 $x_1 \leadsto z$ 与
$x_2 \leadsto z$ 中至少一个不是直接特化时，这种情形才可能发生！因此
$\dim(\mathcal{O}_{Z, z}) \geq 2$。因此
$\text{depth}(\mathcal{O}_{Z, z}) \geq 2$，因为 $A$ 满足 $(S_2)$。然而，穿孔谱
$U$（属于 $\mathcal{O}_{Z, z}$）不连通，因为闭子集
$U \cap X_1$ 与 $U \cap X_2$ 不相交
（由 $z$ 的选择），且二者覆盖 $U$。这与引理
\ref{local-cohomology-lemma-depth-2-connected-punctured-spectrum} 矛盾，证明完毕。
\end{proof}



\section{上同调维数}
\label{local-cohomology-section-cd}

\noindent
本节简要讨论上同调维数。

\begin{lemma}
\label{local-cohomology-lemma-cd}
设 $I \subset A$ 是环 $A$ 的一个有限生成理想。
置 $Y = V(I) \subset X = \Spec(A)$。设 $d \geq -1$ 是整数。
下列条件等价：
\begin{enumerate}
\item $H^i_Y(A) = 0$ 对所有 $i > d$ 成立；
\item $H^i_Y(M) = 0$ 对所有 $i > d$ 以及每个 $A$-模 $M$ 成立；以及
\item 若 $d = -1$，则 $Y = \emptyset$；若 $d = 0$，则
$Y$ 在 $X$ 中既开又闭；若 $d > 0$，则
$H^i(X \setminus Y, \mathcal{F}) = 0$ 对所有 $i \geq d$ 以及每个拟凝聚
$\mathcal{O}_{X \setminus Y}$-模 $\mathcal{F}$ 成立。
\end{enumerate}
\end{lemma}

\begin{proof}
注意到 $R\Gamma_Y(-)$ 具有有限的上同调维数，例如可由《对偶复形》，引理
\ref{dualizing-lemma-local-cohomology-adjoint} 看出。因此存在整数 $i_0$，
使得 $H^i_Y(M) = 0$ 对所有 $A$-模 $M$ 以及所有 $i \geq i_0$ 成立。

\medskip\noindent
先证明 (1) 与 (2) 等价。(2) 蕴含 (1) 是显然的。假设 (1) 成立。
我们对 $i > d$ 作降归纳，证明 $H^i_Y(M) = 0$ 对所有 $A$-模 $M$ 成立。
当 $i \geq i_0$ 时，上面已经证明了这一点。
为完成归纳步骤，设 $i_0 > i > d$。任取一个 $A$-模 $M$，并将其嵌入短正合列
$0 \to N \to F \to M \to 0$，其中 $F$ 是自由 $A$-模。由于
$R\Gamma_Y$ 是右伴随函子，$H^i_Y(-)$ 与直和交换。因此由假设 (1) 有
$H^i_Y(F) = 0$，这是因为 $i > d$。于是如所需，有
$H^i_Y(M) = H^{i + 1}_Y(N) = 0$。

\medskip\noindent
假设 $d = -1$ 且 (2) 成立。于是 $0 = H^0_Y(A/I) = A/I \Rightarrow A = I
\Rightarrow Y = \emptyset$，故 (3) 成立。反向蕴含的证明从略。

\medskip\noindent
假设 $d = 0$ 且 (2) 成立。置
$J = H^0_I(A) = \{x \in A \mid I^nx = 0 \text{ 对某个 }n > 0\}$.
则
$$
H^1_Y(A) = \Coker(A \to \Gamma(X \setminus Y, \mathcal{O}_{X \setminus Y}))
\quad\text{且}\quad
H^1_Y(I) = \Coker(I \to \Gamma(X \setminus Y, \mathcal{O}_{X \setminus Y}))
$$
且第一个映射的核等于 $J$；见引理 \ref{local-cohomology-lemma-local-cohomology}。
由 (2) 可得 $I(A/J) = A/J$。因此可以选取 $f \in I$，使其像为 $1$；
这里的像位于 $A/J$ 中。于是 $1 - f \in J$，故
$I^n(1 - f) = 0$ 对某个 $n > 0$ 成立。
从而 $f^n = f^{n + 1}$，所以 $e = f^n \in I$ 是幂等元。再考虑互补幂等元
$e' = 1 - f^n \in J$。对任意元素 $g \in I$，等式 $g^m e' = 0$
对某个 $m > 0$ 成立。因此 $I$ 包含于理想 $(e) \subset I$ 的根中。这说明
$Y = V(I) = V(e)$ 在 $X$ 中既开又闭，正如 (3) 所述。反过来，若
$Y = V(I)$ 既开又闭，则函子 $H^0_Y(-)$ 正合，且其高阶导出函子均消失。

\medskip\noindent
若 $d > 0$，则由引理 \ref{local-cohomology-lemma-local-cohomology} 立即可见 (2) 与 (3) 等价。
\end{proof}

\begin{definition}
\label{local-cohomology-definition-cd}
设 $I \subset A$ 是环 $A$ 的一个有限生成理想。满足上述等价条件的最小整数
$d \geq -1$（这些条件见引理 \ref{local-cohomology-lemma-cd}），称为
{\it $I$ 在 $A$ 中的上同调维数}，记为 $\text{cd}(A, I)$。
\end{definition}

\noindent
因此，$\text{cd}(A, I) = -1$ 当 $I = A$；而 $\text{cd}(A, I) = 0$
当 $I$ 局部幂零或由一个幂等元生成。下一个引理说明
$\text{cd}(A, I)$ 总是存在。

\begin{lemma}
\label{local-cohomology-lemma-bound-cd}
设 $I \subset A$ 是环 $A$ 的一个有限生成理想。则
\begin{enumerate}
\item $\text{cd}(A, I)$ 不超过 $I$ 的生成元个数；
\item $\text{cd}(A, I) \leq r$，只要存在 $f_1, \ldots, f_r \in A$
使得 $V(f_1, \ldots, f_r) = V(I)$；
\item $\text{cd}(A, I) \leq c$，只要 $\Spec(A) \setminus V(I)$
可由 $c$ 个仿射开集覆盖。
\end{enumerate}
\end{lemma}

\begin{proof}
关于 $R\Gamma_Y(-)$ 的显式描述（见《对偶复形》，引理
\ref{dualizing-lemma-local-cohomology-adjoint}）表明 (1) 成立。为由 (1) 推出
(2)，使用如下事实：$R\Gamma_Z$ 只依赖于闭子集 $Z$，而不依赖于有限生成理想
$I \subset A$ 的选取，其中 $V(I) = Z$。这个事实或者由《对偶复形》，第
\ref{dualizing-section-local-cohomology} 节中的局部上同调构造与《代数进阶》，
引理 \ref{more-algebra-lemma-local-cohomology-closed} 结合而得，或者由引理
\ref{local-cohomology-lemma-local-cohomology-is-local-cohomology} 得出。为证明 (3)，使用引理
\ref{local-cohomology-lemma-cd} 以及《概形的上同调》，引理
\ref{coherent-lemma-vanishing-nr-affines} 的消失结果。
\end{proof}

\begin{lemma}
\label{local-cohomology-lemma-cd-sum}
设 $I, J \subset A$ 是环 $A$ 的有限生成理想。则
$\text{cd}(A, I + J) \leq \text{cd}(A, I) + \text{cd}(A, J)$。
\end{lemma}

\begin{proof}
使用定义以及《对偶复形》，引理
\ref{dualizing-lemma-local-cohomology-ss}。
\end{proof}

\begin{lemma}
\label{local-cohomology-lemma-cd-change-rings}
设 $A \to B$ 是环映射，$I \subset A$ 是有限生成理想。则
$\text{cd}(B, IB) \leq \text{cd}(A, I)$。若 $A \to B$ 忠实平坦，则等号成立。
\end{lemma}

\begin{proof}
使用定义以及《对偶复形》，引理
\ref{dualizing-lemma-torsion-change-rings}。
\end{proof}

\begin{lemma}
\label{local-cohomology-lemma-cd-local}
设 $I \subset A$ 是环 $A$ 的有限生成理想。则
$\text{cd}(A, I) = \max \text{cd}(A_\mathfrak p, I_\mathfrak p)$。
\end{lemma}

\begin{proof}
置 $Y = V(I)$，并置 $Y' = V(I_\mathfrak p) \subset \Spec(A_\mathfrak p)$。
回忆，
$R\Gamma_Y(A) \otimes_A A_\mathfrak p = R\Gamma_{Y'}(A_\mathfrak p)$
由《对偶复形》，引理 \ref{dualizing-lemma-torsion-change-rings} 成立。
因此由《代数》，引理 \ref{algebra-lemma-characterize-zero-local} 得出结论。
\end{proof}

\begin{lemma}
\label{local-cohomology-lemma-cd-dimension}
设 $I \subset A$ 是环 $A$ 的有限生成理想。若 $M$ 是有限 $A$-模，则
$H^i_{V(I)}(M) = 0$ 对所有 $i > \dim(\text{Supp}(M))$ 成立。
特别地，有 $\text{cd}(A, I) \leq \dim(A)$。
\end{lemma}

\begin{proof}
先证明第二个断言。回忆，$\dim(A)$ 表示 Krull 维数。由引理
\ref{local-cohomology-lemma-cd-local}，可以假设 $A$ 是局部环。若 $V(I) = \emptyset$，则结论成立。
若 $V(I) \not = \emptyset$，则由于缺少闭点，有
$\dim(\Spec(A) \setminus V(I)) < \dim(A)$。注意到
$U = \Spec(A) \setminus V(I)$ 是谱空间 $\Spec(A)$ 的拟紧开子集，因而其本身也是
谱空间；见《代数》，引理 \ref{algebra-lemma-spec-spectral} 与《拓扑》，引理
\ref{topology-lemma-spectral-sub}。因此《上同调》，命题
\ref{cohomology-proposition-cohomological-dimension-spectral} 蕴含
$H^i(U, \mathcal{F}) = 0$ 对所有 $i \geq \dim(A)$ 成立，再由引理
\ref{local-cohomology-lemma-cd} 即得所需结论。在诺特情形，读者也可使用《Grothendieck 上同调》，
命题 \ref{cohomology-proposition-vanishing-Noetherian}。

\medskip\noindent
下面由第二个断言推出第一个断言。设 $\mathfrak a$ 是有限 $A$-模 $M$ 的零化子。
置 $B = A/\mathfrak a$。回忆，$\Spec(B) = \text{Supp}(M)$；见《代数》，引理
\ref{algebra-lemma-support-closed}。置 $J = IB$。则 $M$ 是 $B$-模，且
$H^i_{V(I)}(M) = H^i_{V(J)}(M)$；见《对偶复形》，引理
\ref{dualizing-lemma-local-cohomology-and-restriction}。由于第一部分给出
$\text{cd}(B, J) \leq \dim(B) = \dim(\text{Supp}(M))$，结论成立。
\end{proof}

\begin{lemma}
\label{local-cohomology-lemma-cd-is-one}
设 $I \subset A$ 是环 $A$ 的有限生成理想。若
$\text{cd}(A, I) = 1$，则 $\Spec(A) \setminus V(I)$ 是非空仿射概形。
\end{lemma}

\begin{proof}
这由引理 \ref{local-cohomology-lemma-cd} 与《概形的上同调》，引理
\ref{coherent-lemma-quasi-compact-h1-zero-covering} 得出。
\end{proof}

\begin{lemma}
\label{local-cohomology-lemma-cd-maximal}
设 $(A, \mathfrak m)$ 是维数为 $d$ 的诺特局部环。则
$H^d_\mathfrak m(A)$ 非零，且 $\text{cd}(A, \mathfrak m) = d$。
\end{lemma}

\begin{proof}
由维数的一种刻画，$A$ 存在一个由 $d$ 个元素生成的定义理想；见《代数》，命题
\ref{algebra-proposition-dimension}。因此，由引理 \ref{local-cohomology-lemma-bound-cd}，
$\text{cd}(A, \mathfrak m) \leq d$。所以 $H^d_\mathfrak m(A)$ 非零，当且仅当
$\text{cd}(A, \mathfrak m) = d$，又当且仅当
$\text{cd}(A, \mathfrak m) \geq d$。

\medskip\noindent
设 $A \to A^\wedge$ 是从 $A$ 到其完备化的映射。注意到 $A^\wedge$
是与 $A$ 同维的诺特局部环，其极大理想为 $\mathfrak m A^\wedge$。见《代数》，
各引理
\ref{algebra-lemma-completion-Noetherian-Noetherian}、
\ref{algebra-lemma-completion-complete}、
\ref{algebra-lemma-completion-faithfully-flat}，以及
《代数进阶》，引理 \ref{more-algebra-lemma-completion-dimension}。
由引理 \ref{local-cohomology-lemma-cd-change-rings}，只需对 $A^\wedge$ 证明本引理。

\medskip\noindent
由上一段，可以假设 $A$ 是完备局部环。于是 $A$ 有一个规范化对偶复形
$\omega_A^\bullet$（《对偶复形》，引理
\ref{dualizing-lemma-ubiquity-dualizing}）。局部对偶定理（取《对偶复形》，引理
\ref{dualizing-lemma-special-case-local-duality} 的形式）说明
$H^d_\mathfrak m(A)$ 与
$\text{Ext}^{-d}(A, \omega_A^\bullet) = H^{-d}(\omega_A^\bullet)$
Matlis 对偶；后者非零，例如可由《对偶复形》，引理
\ref{dualizing-lemma-nonvanishing-generically-local} 看出。
\end{proof}

\begin{lemma}
\label{local-cohomology-lemma-cd-bound-dim-local}
设 $(A, \mathfrak m)$ 是诺特局部环，$I \subset A$ 是真理想。设
$\mathfrak p \subset A$ 是满足
$V(\mathfrak p) \cap V(I) = \{\mathfrak m\}$ 的素理想。则
$\dim(A/\mathfrak p) \leq \text{cd}(A, I)$。
\end{lemma}

\begin{proof}
由引理 \ref{local-cohomology-lemma-cd-change-rings}，有
$\text{cd}(A, I) \geq \text{cd}(A/\mathfrak p, I(A/\mathfrak p))$.
由于 $V(I) \cap V(\mathfrak p) = \{\mathfrak m\}$，有
$\text{cd}(A/\mathfrak p, I(A/\mathfrak p)) =
\text{cd}(A/\mathfrak p, \mathfrak m/\mathfrak p)$.
由引理 \ref{local-cohomology-lemma-cd-maximal}，这等于 $\dim(A/\mathfrak p)$。
\end{proof}

\begin{lemma}
\label{local-cohomology-lemma-cd-blowup}
设 $A$ 是诺特环，$I \subset A$ 是理想。设
$b : X' \to X = \Spec(A)$ 是沿 $I$ 的吹起。若 $b$ 的纤维维数
$\leq d - 1$，则 $\text{cd}(A, I) \leq d$。
\end{lemma}

\begin{proof}
置 $U = X \setminus V(I)$。记 $j : U \to X'$ 为典范开浸入；见《除子》，第
\ref{divisors-section-blowing-up} 节。由于例外除子是有效 Cartier 除子
（《除子》，引理
\ref{divisors-lemma-blowing-up-gives-effective-Cartier-divisor}），可知 $j$
是仿射态射；见《除子》，引理
\ref{divisors-lemma-complement-locally-principal-closed-subscheme}。设
$\mathcal{F}$ 是拟凝聚 $\mathcal{O}_U$-模。则
$R^pj_*\mathcal{F} = 0$ 对所有 $p > 0$ 成立；见《概形的上同调》，引理
\ref{coherent-lemma-relative-affine-vanishing}。另一方面，由《极限》，引理
\ref{limits-lemma-higher-direct-images-zero-above-dimension-fibre}，有
$R^qb_*(j_*\mathcal{F}) = 0$ 对所有 $q \geq d$ 成立。因此，由 Leray 谱序列
（《上同调》，引理 \ref{cohomology-lemma-relative-Leray}），可得当
$R^n(b \circ j)_*\mathcal{F} = 0$ 对所有 $n \geq d$ 成立。于是
$H^n(U, \mathcal{F}) = 0$ 对所有 $n \geq d$ 成立（由《上同调》，引理
\ref{cohomology-lemma-apply-Leray}）。按照定义，这意味着
$\text{cd}(A, I) \leq d$。
\end{proof}







\section{更一般的支撑}
\label{local-cohomology-section-supports}

\noindent
设 $A$ 是诺特环，$M$ 是 $A$-模。设 $T \subset \Spec(A)$ 是在特化下稳定的
子集（《拓扑》，定义 \ref{topology-definition-specialization}）。定义
$$
H^0_T(M) = \colim_{Z \subset T} H^0_Z(M)
$$
其中余极限取遍由闭子集 $Z$ 构成的有向偏序集；这些闭子集位于
$\Spec(A)$ 中并包含于 $T$\footnote{由于 $T$ 在特化下稳定，有
$T = \bigcup_{Z \subset T} Z$；见《拓扑》，引理
\ref{topology-lemma-stable-specialization}。}。换言之，元素 $m$（属于 $M$）
属于 $H^0_T(M) \subset M$，当且仅当支撑
$V(\text{Ann}_R(m))$（即元素 $m$ 的支撑）包含于 $T$。

\begin{lemma}
\label{local-cohomology-lemma-support}
设 $A$ 是诺特环，$T \subset \Spec(A)$ 是在特化下稳定的子集。对 $A$-模 $M$，
下列条件等价：
\begin{enumerate}
\item $H^0_T(M) = M$；以及
\item $\text{Supp}(M) \subset T$.
\end{enumerate}
这类 $A$-模组成的范畴是 $A$-模范畴的 Serre 子范畴，并且在直和下封闭。
\end{lemma}

\begin{proof}
上述等价性成立，因为 $M$ 中每个元素的支撑都包含于 $M$ 的支撑；反过来，
$M$ 的支撑是其各元素支撑之并。由《同调》，定义
\ref{homology-definition-serre-subcategory} 与《代数》，引理
\ref{algebra-lemma-support-quotient}，这些模所成的范畴是
$\text{Mod}_A$ 的 Serre 子范畴。关于直和的断言，证明从略。
\end{proof}

\noindent
设 $A$ 是诺特环，$T \subset \Spec(A)$ 是在特化下稳定的子集。用
$\text{Mod}_{A, T} \subset \text{Mod}_A$ 表示引理 \ref{local-cohomology-lemma-support}
所述的 Serre 子范畴。用 $D_T(A) \subset D(A)$ 表示 $D(A)$ 的严格全饱和
三角子范畴（《导出范畴》，引理
\ref{derived-lemma-cohomology-in-serre-subcategory}）；它由如下复形组成：
这些复形由 $A$-模构成，且其上同调模属于 $\text{Mod}_{A, T}$。于是得到函子
$$
D(\text{Mod}_{A, T}) \to D_T(A) \to D(A)
$$
见《导出范畴》，第 \ref{derived-section-triangulated-sub} 节中的讨论。
用 $RH^0_T : D(A) \to D(\text{Mod}_{A, T})$ 表示 $H^0_T$ 的右导出延拓。用
$$
R\Gamma_T : D^+(A) \to D^+_T(A),
$$
表示 $RH^0_T : D^+(A) \to D^+(\text{Mod}_{A, T})$ 与
$D^+(\text{Mod}_{A, T}) \to D^+_T(A)$ 的复合。若 $A$ 的维数有限
\footnote{若 $\dim(A) = \infty$，这个构造在无界复形上可能具有预料之外的性质。}，
则用
$$
R\Gamma_T : D(A) \to D_T(A)
$$
表示 $RH^0_T$ 与 $D(\text{Mod}_{A, T}) \to D_T(A)$ 的复合。

\begin{lemma}
\label{local-cohomology-lemma-adjoint}
设 $A$ 是诺特环，$T \subset \Spec(A)$ 是在特化下稳定的子集。函子
$RH^0_T$ 是函子 $D(\text{Mod}_{A, T}) \to D(A)$ 的右伴随。
\end{lemma}

\begin{proof}
这是因为函子 $H^0_T(-)$ 是嵌入函子
$\text{Mod}_{A, T} \to \text{Mod}_A$ 的右伴随；见《导出范畴》，引理
\ref{derived-lemma-derived-adjoint-functors}。
\end{proof}

\begin{lemma}
\label{local-cohomology-lemma-adjoint-ext}
设 $A$ 是诺特环，$T \subset \Spec(A)$ 是在特化下稳定的子集。对任意对象
$K$（属于 $D(A)$），有
$$
H^i(RH^0_T(K)) = \colim_{Z \subset T\text{ 闭}} H^i_Z(K)
$$
\end{lemma}

\begin{proof}
设 $J^\bullet$ 是表示 $K$ 的 K-内射复形。按照定义，$RH^0_T$ 由复形
$$
H^0_T(J^\bullet) = \colim H^0_Z(J^\bullet)
$$
表示，其中等号来自 $H^0_T$ 的定义。由于滤过余极限正合，此复形的 $i$ 次上同调为
$\colim H^i(H^0_Z(J^\bullet)) = \colim H^i_Z(K)$
，正是所需结论。
\end{proof}

\begin{lemma}
\label{local-cohomology-lemma-equal-plus}
设 $A$ 是诺特环，$T \subset \Spec(A)$ 是在特化下稳定的子集。函子
$D^+(\text{Mod}_{A, T}) \to D^+_T(A)$ 是等价。
\end{lemma}

\begin{proof}
设 $M$ 是 $\text{Mod}_{A, T}$ 的对象。选取嵌入 $M \to J$，其目标是内射
$A$-模。由《对偶复形》，命题
\ref{dualizing-proposition-structure-injectives-noetherian}，模 $J$ 是剩余域
内射包的直和。设 $E$ 是 $\mathfrak p$ 的剩余域的内射包。由于 $E$ 是
$\mathfrak p$-幂挠模，可知 $H^0_T(E) = 0$ 当 $\mathfrak p \not \in T$，而
$H^0_T(E) = E$ 当 $\mathfrak p \in T$。因此，作为内射包的
直和，$H^0_T(J)$ 是内射的（由该命题），并且有嵌入
$M \to H^0_T(J)$。所以每个对象 $M$（属于 $\text{Mod}_{A, T}$）都有内射消解
$M \to J^\bullet$，其中 $J^n$ 也属于 $\text{Mod}_{A, T}$。由此可得
$RH^0_T(M) = M$。

\medskip\noindent
接着，假设 $K \in D_T^+(A)$。则谱序列
$$
R^qH^0_T(H^p(K)) \Rightarrow R^{p + q}H^0_T(K)
$$
收敛（《导出范畴》，引理
\ref{derived-lemma-two-ss-complex-functor}），而上面已经看到，只有 $q = 0$
的项可能非零。因此 $RH^0_T(K) \to K$ 是同构。所以函子
$D^+(\text{Mod}_{A, T}) \to D^+_T(A)$ 是等价，其拟逆由 $RH^0_T$ 给出。
\end{proof}

\begin{lemma}
\label{local-cohomology-lemma-equal-full}
设 $A$ 是诺特环，$T \subset \Spec(A)$ 是在特化下稳定的子集。若
$\dim(A) < \infty$，则函子 $D(\text{Mod}_{A, T}) \to D_T(A)$ 是等价。
\end{lemma}

\begin{proof}
记 $\dim(A) = d$。于是 $H^i_Z(M) = 0$ 对所有 $i > d$ 以及每个闭子集
$Z$（位于 $\Spec(A)$）成立；见引理
\ref{local-cohomology-lemma-cd-dimension}。由引理 \ref{local-cohomology-lemma-adjoint-ext} 可知，$H^0_T$
具有有界的上同调维数。

\medskip\noindent
设 $K \in D_T(A)$。我们断言 $RH^0_T(K) \to K$ 是同构。当 $K$ 下有界时，
这一点已经成立；见引理 \ref{local-cohomology-lemma-equal-plus}。不过，由于 $H^0_T$ 具有有界的
上同调维数，第 $i$ 个上同调（属于 $RH_T^0(K)$）只依赖于
$\tau_{\geq -d + i}K$，故结论成立。因此
$D(\text{Mod}_{A, T}) \to D_T(A)$ 是等价，其拟逆为 $RH^0_T$。
\end{proof}

\begin{remark}
\label{local-cohomology-remark-upshot}
设 $A$ 是诺特环，$T \subset \Spec(A)$ 是在特化下稳定的子集。上述讨论的结论是：
$R\Gamma_T : D^+(A) \to D_T^+(A)$ 是嵌入函子
$D_T^+(A) \to D^+(A)$ 的右伴随。若 $\dim(A) < \infty$，则
$R\Gamma_T : D(A) \to D_T(A)$ 是嵌入函子 $D_T(A) \to D(A)$ 的右伴随。
在两种情形中都有
$$
H^i_T(K) = H^i(R\Gamma_T(K)) = R^iH^0_T(K) =
\colim_{Z \subset T\text{ 闭}} H^i_Z(K)
$$
这由引理 \ref{local-cohomology-lemma-adjoint}、\ref{local-cohomology-lemma-adjoint-ext}、
\ref{local-cohomology-lemma-equal-plus} 及 \ref{local-cohomology-lemma-equal-full} 结合得出。
\end{remark}

\begin{lemma}
\label{local-cohomology-lemma-torsion-change-rings}
设 $A \to B$ 是诺特环之间的平坦同态。设 $T \subset \Spec(A)$ 是在特化下稳定的
子集，并设 $T' \subset \Spec(B)$ 是 $T$ 的逆像。则典范映射
$$
R\Gamma_T(K) \otimes_A^\mathbf{L} B
\longrightarrow
R\Gamma_{T'}(K \otimes_A^\mathbf{L} B)
$$
对 $K \in D^+(A)$ 是同构。若 $A$ 与 $B$ 都具有有限维数，则该结论对
$K \in D(A)$ 也成立。
\end{lemma}

\begin{proof}
从映射 $R\Gamma_T(K) \to K$ 得到映射
$R\Gamma_T(K) \otimes_A^\mathbf{L} B \to K \otimes_A^\mathbf{L} B$。
$R\Gamma_T(K) \otimes_A^\mathbf{L} B$ 的各上同调模支撑于 $T'$，因而得到本引理中的
箭头。若 $T$ 是 $\Spec(A)$ 的闭子集，则由《对偶复形》，引理
\ref{dualizing-lemma-torsion-change-rings}，该箭头是同构。回忆，
$H^i_T(K)$ 是 $H^i_Z(K)$ 的余极限，其中 $Z$ 取遍 $T$ 的闭子集所成的有向集；
见引理 \ref{local-cohomology-lemma-adjoint-ext}。相应地，
$H^i_{T'}(K \otimes_A^\mathbf{L} B) =
\colim H^i_{Z'}(K \otimes_A^\mathbf{L} B)$，其中 $Z'$ 是 $Z$ 的逆像。由于
$\otimes_A B$ 与滤过余极限交换，且不存在高阶 Tor，结论成立。
\end{proof}

\begin{lemma}
\label{local-cohomology-lemma-local-cohomology-ss}
设 $A$ 是环，并设 $T, T' \subset \Spec(A)$ 是在特化下稳定的子集。对
$K \in D^+(A)$，存在谱序列
$$
E_2^{p, q} = H^p_T(H^p_{T'}(K)) \Rightarrow H^{p + q}_{T \cap T'}(K)
$$
，如《导出范畴》，引理
\ref{derived-lemma-grothendieck-spectral-sequence} 所述。
\end{lemma}

\begin{proof}
设 $E$ 是 $D_{T \cap T'}(A)$ 的对象。则有
$$
\Hom(E, R\Gamma_T(R\Gamma_{T'}(K))) =
\Hom(E, R\Gamma_{T'}(K)) =
\Hom(E, K)
$$
第一个等号来自 $R\Gamma_T$ 的伴随性质，第二个等号来自 $R\Gamma_{T'}$ 的伴随性质。
另一方面，若 $J^\bullet$ 是表示 $K$ 的下有界内射复形，则
$H^0_{T'}(J^\bullet)$ 是由内射 $A$-模构成并表示 $R\Gamma_{T'}(K)$ 的复形，因而
$H^0_T(H^0_{T'}(J^\bullet))$ 是表示 $R\Gamma_T(R\Gamma_{T'}(K))$ 的复形。所以
$R\Gamma_T(R\Gamma_{T'}(K))$ 是 $D^+_{T \cap T'}(A)$ 的对象。结合这两个事实，
得到 $R\Gamma_{T \cap T'} = R\Gamma_T \circ R\Gamma_{T'}$。由陈述中引用的引理，
这便产生上述谱序列。
\end{proof}

\begin{lemma}
\label{local-cohomology-lemma-torsion-tensor-product}
设 $A$ 是诺特环，$T \subset \Spec(A)$ 是在特化下稳定的子集。假设 $A$ 具有有限维数。
则
$$
R\Gamma_T(K) = R\Gamma_T(A) \otimes_A^\mathbf{L} K
$$
对 $K \in D(A)$ 成立。对 $K, L \in D(A)$，有
$$
R\Gamma_T(K \otimes_A^\mathbf{L} L) =
K \otimes_A^\mathbf{L} R\Gamma_T(L) =
R\Gamma_T(K) \otimes_A^\mathbf{L} L =
R\Gamma_T(K) \otimes_A^\mathbf{L} R\Gamma_T(L)
$$
若 $K$ 或 $L$ 属于 $D_T(A)$，则 $K \otimes_A^\mathbf{L} L$ 也属于其中。
\end{lemma}

\begin{proof}
按照构造，对象 $R\Gamma_T(A)$ 可以由复形 $J^\bullet$ 表示，该复形位于
$\text{Mod}_{A, T}$ 中。因此，若对象 $K$ 由 K-平坦复形 $K^\bullet$ 表示，则
$R\Gamma_T(A) \otimes_A^\mathbf{L} K$ 由复形
$\text{Tot}(J^\bullet \otimes_A K^\bullet)$ 表示，该复形位于
$\text{Mod}_{A, T}$ 中。利用映射
$R\Gamma_T(A) \to A$，得到映射
$R\Gamma_T(A) \otimes_A^\mathbf{L} K\to K$。于是由 $R\Gamma_T$ 的伴随性质，
得到典范映射
$$
R\Gamma_T(A) \otimes_A^\mathbf{L} K \longrightarrow R\Gamma_T(K)
$$
，它分解刚刚构造的映射。注意到 $R\Gamma_T$ 与 $D(A)$ 中的直和交换；例如，
这可由引理 \ref{local-cohomology-lemma-adjoint-ext}、有向余极限与直和交换这一事实，以及通常的
局部上同调与直和交换这一事实得出（例如由《对偶复形》，引理
\ref{dualizing-lemma-local-cohomology-adjoint}）。因此，由《代数进阶》，注记
\ref{more-algebra-remark-P-resolution}，只需检验当 $K = A[k]$ 且
$k \in \mathbf{Z}$ 时该映射为同构。这是显然的。

\medskip\noindent
最后几个断言由刚刚证明的结果以及导出张量积的传递性得出。
\end{proof}





\section{局部上同调的滤过}
\label{local-cohomology-section-filter-local-cohomology}

\noindent
本节介绍一些与引理 \ref{local-cohomology-lemma-local-cohomology-ss} 的谱序列有关的技巧。

\begin{lemma}
\label{local-cohomology-lemma-filter-local-cohomology}
设 $A$ 是诺特环，$T \subset \Spec(A)$ 是在特化下稳定的子集。设
$T' \subset T$ 是 $T$ 中非极小素理想所成的集合。则 $T'$ 是
$\Spec(A)$ 的在特化下稳定的子集，并且对每个 $A$-模 $M$ 都有正合列
$$
0 \to
\colim_{Z, f} H^1_f(H^{i - 1}_Z(M)) \to
H^i_{T'}(M) \to H^i_T(M) \to
\bigoplus\nolimits_{\mathfrak p \in T \setminus T'}
H^i_{\mathfrak p A_\mathfrak p}(M_\mathfrak p)
$$
其中余极限取遍闭子集 $Z \subset T$ 以及 $f \in A$，并要求
$V(f) \cap Z \subset T'$。
\end{lemma}

\begin{proof}
对每个 $Z$ 和 $f$，《对偶复形》，引理
\ref{dualizing-lemma-local-cohomology-ss} 的谱序列退化，从而给出短正合列
$$
0 \to H^1_f(H^{i - 1}_Z(M)) \to
H^i_{Z \cap V(f)}(M) \to H^0_f(H^i_Z(M)) \to 0
$$
下文将直接使用这一点，不再另行说明。

\medskip\noindent
设 $\xi \in H^i_T(M)$ 在该直和中映为零。先把 $\xi$ 写成某个
$\xi' \in H^i_Z(M)$ 的像，其中 $Z \subset T$ 是闭子集；见引理
\ref{local-cohomology-lemma-adjoint-ext}。于是 $\xi'$ 在
$H^i_{\mathfrak p A_\mathfrak p}(M_\mathfrak p)$ 中映为零，对每个
$\mathfrak p \in Z$ 且 $\mathfrak p \not \in T'$ 都如此。这样的素理想只有有限个，
所以可以选取 $f \in A$，使其不属于其中任何一个，并且 $f$ 零化 $\xi'$。于是
$\xi'$ 是某个 $\xi'' \in H^i_{Z'}(M)$ 的像，其中
$Z' = Z \cap V(f)$。由 $f$ 的选取，有 $Z' \subset T'$，从而在倒数第二项处正合。

\medskip\noindent
设 $\xi \in H^i_{T'}(M)$ 在 $H^i_T(M)$ 中映为零。选取闭子集
$Z' \subset Z$，使得 $Z' \subset T'$ 且 $Z \subset T$，并使 $\xi$ 来自
$\xi' \in H^i_{Z'}(M)$，而后者在 $H^i_Z(M)$ 中映为零。于是可以找到
$f \in A$，使得 $V(f) \cap Z = Z'$，结论随即得出。
\end{proof}

\begin{lemma}
\label{local-cohomology-lemma-zero}
设 $A$ 是有限维诺特环，$T \subset \Spec(A)$ 是在特化下稳定的子集。设
$\{M_n\}_{n \geq 0}$ 是 $A$-模的逆系，并设 $i \geq 0$ 是整数。假设对每个
$m$ 都存在整数 $m'(m) \geq m$，使得对所有 $\mathfrak p \in T$，诱导映射
$$
H^i_{\mathfrak p A_\mathfrak p}(M_{k, \mathfrak p})
\longrightarrow
H^i_{\mathfrak p A_\mathfrak p}(M_{m, \mathfrak p})
$$
对 $k \geq m'(m)$ 为零。设 $m'' : \mathbf{N} \to \mathbf{N}$ 是
$2^{\dim(T)}$ 次自复合，其中参与复合的映射为 $m'$。则映射
$H^i_T(M_k) \to H^i_T(M_m)$ 对所有
$k \geq m''(m)$ 为零。
\end{lemma}

\begin{proof}
先作一个一般性的说明。假设有逆阿贝尔群系的正合列
$$
(A_n) \to (B_n) \to (C_n)
$$
，并假设对每个 $m$ 都存在整数 $m'(m) \geq m$，使得
$$
A_k \to A_m
\quad\text{且}\quad
C_k \to C_m
$$
对 $k \geq m'(m)$ 为零。则当 $k \geq m'(m'(m))$ 时，映射
$B_k \to B_m$ 为零。

\medskip\noindent
对 $\dim(T)$ 作归纳来证明本引理；由于 $\dim(A)$ 有限，它也是有限的。设
$T' \subset T$ 是 $T$ 中非极小素理想所成的集合。则 $T'$ 是
$\Spec(A)$ 的在特化下稳定的子集，而且本引理的假设适用于 $T'$。由于
$\dim(T') < \dim(T)$，已知本引理对 $T'$ 成立。对每个 $A$-模 $M$，都有正合列
$$
H^i_{T'}(M) \to H^i_T(M) \to
\bigoplus\nolimits_{\mathfrak p \in T \setminus T'}
H^i_{\mathfrak p A_\mathfrak p}(M_\mathfrak p)
$$
，这是引理 \ref{local-cohomology-lemma-filter-local-cohomology} 的结论。因此由证明开头的一般性说明
即可得出结论。
\end{proof}

\begin{lemma}
\label{local-cohomology-lemma-essential-image}
设 $A$ 是诺特环，$T \subset \Spec(A)$ 是在特化下稳定的子集。设
$\{M_n\}_{n \geq 0}$ 是 $A$-模的逆系，并设 $i \geq 0$ 是整数。假设 $A$ 的维数
有限，而且对每个 $m$ 都存在整数 $m'(m) \geq m$，使得对所有
$\mathfrak p \in T$ 有
\begin{enumerate}
\item $H^{i - 1}_{\mathfrak p A_\mathfrak p}(M_{k, \mathfrak p})
\to H^{i - 1}_{\mathfrak p A_\mathfrak p}(M_{m, \mathfrak p})$
对 $k \geq m'(m)$ 为零；以及
\item $ H^i_{\mathfrak p A_\mathfrak p}(M_{k, \mathfrak p}) \to
H^i_{\mathfrak p A_\mathfrak p}(M_{m, \mathfrak p})$
的像 $G(\mathfrak p, m)$ 与 $k \geq m'(m)$ 的选取无关，并且
$G(\mathfrak p, m)$ 单射地映入
$H^i_{\mathfrak p A_\mathfrak p}(M_{0, \mathfrak p})$。
\end{enumerate}
则存在整数 $m_0$，使得对每个 $m \geq m_0$，都存在整数
$m''(m) \geq m$，且对 $k \geq m''(m)$，映射
$H^i_T(M_k) \to H^i_T(M_m)$ 的像单射地映入 $H^i_T(M_{m_0})$。
\end{lemma}

\begin{proof}
先作一个一般性的说明。假设有逆阿贝尔群系的正合列
$$
(A_n) \to (B_n) \to (C_n) \to (D_n)
$$
。假设存在整数 $m_0$，使得对每个 $m \geq m_0$，都存在整数
$m'(m) \geq m$，使得映射
$$
\Im(B_k \to B_m) \longrightarrow B_{m_0}
\quad\text{且}\quad
\Im(D_k \to D_m) \longrightarrow D_{m_0}
$$
对 $k \geq m'(m)$ 是单射，并且 $A_k \to A_m$ 对 $k \geq m'(m)$ 为零。
那么，对 $m \geq m'(m_0)$ 和 $k \geq m'(m'(m))$，映射
$$
\Im(C_k \to C_m) \to C_{m'(m_0)}
$$
是单射。事实上，设 $c_0 \in C_m$ 是 $c_3 \in C_k$ 的像，并假设 $c_0$
在 $C_{m'(m_0)}$ 中映为零。图示为
$$
C_k \to C_{m'(m'(m))} \to C_{m'(m)} \to C_m \to C_{m'(m_0)},\quad
c_3 \mapsto c_2 \mapsto c_1 \mapsto c_0 \mapsto 0
$$
需要证明 $c_0 = 0$。像 $d_3$（它是 $c_3$ 的像）在 $C_{m_0}$ 中映为零，因而其像
$d_1 \in D_{m'(m)}$ 为零。因此可以选取 $b_1 \in B_{m'(m)}$，使其映为
$c_1$ 的像。由于 $c_3$ 在 $C_{m'(m_0)}$ 中映为零，可以找到元素
$a_{-1} \in A_{m'(m_0)}$，它映为像
$b_{-1} \in B_{m'(m_0)}$；后者是 $b_1$ 的像。由于 $a_{-1}$ 在
$A_{m_0}$ 中映为零，可得
$b_1$ 在 $B_{m_0}$ 中映为零。因此像 $b_0 \in B_m$ 为零，这当然蕴含所需的
$c_0 = 0$。

\medskip\noindent
对 $\dim(T)$ 作归纳来证明本引理；由于 $\dim(A)$ 有限，它也是有限的。设
$T' \subset T$ 是 $T$ 中非极小素理想所成的集合。则 $T'$ 是 $\Spec(A)$ 的
在特化下稳定的子集，而且本引理的假设适用于 $T'$。由于
$\dim(T') < \dim(T)$，已知本引理对 $T'$ 成立。对每个 $A$-模 $M$，都有正合列
$$
0 \to \colim_{Z, f} H^1_f(H^{i - 1}_Z(M)) \to
H^i_{T'}(M) \to H^i_T(M) \to
\bigoplus\nolimits_{\mathfrak p \in T \setminus T'}
H^i_{\mathfrak p A_\mathfrak p}(M_\mathfrak p)
$$
，这是引理 \ref{local-cohomology-lemma-filter-local-cohomology} 的结论。因此由证明开头的一般性说明
以及下面这个已知事实即可得出结论：群系
$$
\left\{\colim_{Z, f} H^1_f(H^{i - 1}_Z(M_n))\right\}_{n \geq 0}
$$
在引理 \ref{local-cohomology-lemma-zero} 中是本质为零的逆系；这里使用了函数
$m''(m)$（该引理中针对 $H^{i - 1}_Z(M)$ 的函数）与 $Z$ 无关这一点。
\end{proof}










\section{局部上同调的有限性，I}
\label{local-cohomology-section-finiteness}

\noindent
我们将沿用 Faltings 处理局部上同调模有限性的方法；见
\cite{Faltings-annulators} 与 \cite{Faltings-finiteness}。下面的引理表明，
为证明局部上同调模是有限模，只需证明它们具有零化理想。

\begin{lemma}
\label{local-cohomology-lemma-check-finiteness-local-cohomology-by-annihilator}
\begin{reference}
\cite[Lemma 3]{Faltings-annulators}
\end{reference}
设 $A$ 是诺特环，$T \subset \Spec(A)$ 是在特化下稳定的子集，$M$ 是有限
$A$-模，并设 $n \geq 0$。下列条件等价：
\begin{enumerate}
\item $H^i_T(M)$ 对 $i \leq n$ 是有限的；
\item 存在理想 $J \subset A$，满足 $V(J) \subset T$，使得 $J$ 对
$H^i_T(M)$ 起零化作用，对所有 $i \leq n$ 都如此。
\end{enumerate}
若 $T = V(I) = Z$ 对某个理想 $I \subset A$ 成立，则上述条件还等价于
\begin{enumerate}
\item[(3)] 存在 $e \geq 0$，使得 $I^e$ 零化 $H^i_Z(M)$，对所有
$i \leq n$ 都如此。
\end{enumerate}
\end{lemma}

\begin{proof}
对 $n$ 作归纳，证明 (1) 与 (2) 等价。当 $n = 0$ 时，
$H^0_T(M) \subset M$ 是有限的，故 (1) 成立。由于
$H^0_T(M) = \colim H^0_{V(J)}(M)$，其中 $J$ 如 (2) 所述，(2) 也成立。
假设 $n > 0$。

\medskip\noindent
假设 (1) 成立。回忆，$H^i_J(M) = H^i_{V(J)}(M)$；见《对偶复形》，引理
\ref{dualizing-lemma-local-cohomology-noetherian}。因此
$H^i_T(M) = \colim H^i_J(M)$，其中余极限取遍理想 $J \subset A$，并要求
$V(J) \subset T$；见引理 \ref{local-cohomology-lemma-adjoint-ext}。由于 $H^i_T(M)$ 对
$i \leq n$ 是有限生成的，可以找到 (2) 中的某个 $J \subset A$，使得
$H^i_J(M) \to H^i_T(M)$ 对 $i \leq n$ 满射。因此有限个生成元都是 $J$-幂挠元，
从而把 $J$ 替换成某个幂后，(2) 成立。

\medskip\noindent
假设有 (2) 中的 $J$。置 $N = H^0_T(M)$ 以及 $M' = M/N$。由 $R\Gamma_T$ 的构造，
可知 $H^i_T(N) = 0$ 对 $i > 0$ 成立，且 $H^0_T(N) = N$；见注记
\ref{local-cohomology-remark-upshot}。因此 $H^0_T(M') = 0$，并且
$H^i_T(M') = H^i_T(M)$ 对 $i > 0$ 成立。由此可以把 $M$ 替换为 $M'$。

所以可以假设 $H^0_T(M) = 0$。这意味着 $M$ 的有限个伴随素理想均不属于 $T$。
由素避（《代数》，引理 \ref{algebra-lemma-silly}），可以找到 $f \in J$，它不属于
$M$ 的任何伴随素理想。于是与短正合列
$$
0 \to M \to M \to M/fM \to 0
$$
相伴的局部上同调长正合列化为短正合列
$$
0 \to H^i_T(M) \to H^i_T(M/fM) \to H^{i + 1}_T(M) \to 0
$$
，其中 $i < n$。由此可得 $J^2$ 零化 $H^i_T(M/fM)$，对所有
$i < n$ 都如此。
由归纳假设，$H^i_T(M/fM)$ 对 $i < n$ 是有限的。再次使用该短正合列，可知
$H^{i + 1}_T(M)$ 对 $i < n$ 是有限的，正是所需结论。

\medskip\noindent
当 $T = V(I)$ 时，(2) 与 (3) 等价的证明从略。
\end{proof}

\noindent
下面这个 Faltings 的结果使我们能够在局部环层面证明局部上同调的有限性。

\begin{lemma}
\label{local-cohomology-lemma-check-finiteness-local-cohomology-locally}
\begin{reference}
这是 \cite[Satz 1]{Faltings-finiteness} 的一个特殊情形。
\end{reference}
设 $A$ 是诺特环，$I \subset A$ 是理想，$M$ 是有限 $A$-模，并且 $n \geq 0$
是整数。置 $Z = V(I)$。下列条件等价：
\begin{enumerate}
\item 模 $H^i_Z(M)$ 对 $i \leq n$ 是有限的；以及
\item 对所有 $\mathfrak p \in \Spec(A)$，模 $H^i_Z(M)_\mathfrak p$
在 $i \leq n$ 时是有限 $A_\mathfrak p$-模。
\end{enumerate}
\end{lemma}

\begin{proof}
(1) $\Rightarrow$ (2) 是显然的。对 $n$ 作归纳来证明反向蕴含。$n = 0$ 的情形
很清楚，因为此时 (1) 与 (2) 总是成立。

\medskip\noindent
假设 $n > 0$ 且 (2) 成立。置 $N = H^0_Z(M)$ 以及 $M' = M/N$。由《对偶复形》，
引理 \ref{dualizing-lemma-divide-by-torsion}，可以把 $M$ 替换为 $M'$。
因此可以假设 $H^0_Z(M) = 0$。这意味着 $\text{depth}_I(M) > 0$
（《对偶复形》，引理 \ref{dualizing-lemma-depth}）。选取 $f \in I$，使其为
$M$ 上的非零因子，并考虑短正合列
$$
0 \to M \to M \to M/fM \to 0
$$
，它产生长正合列
$$
0 \to H^0_Z(M/fM) \to H^1_Z(M) \to H^1_Z(M) \to H^1_Z(M/fM) \to
H^2_Z(M) \to \ldots
$$
，局部化后亦然。因此假设 (2) 蕴含模 $H^i_Z(M/fM)_\mathfrak p$ 对
$i < n$ 是有限的。于是由归纳假设，$H^i_Z(M/fM)$ 对 $i < n$ 是有限的。

\medskip\noindent
设 $\mathfrak p$ 是 $A$ 的素理想，并且它是 $H^i_Z(M)$ 的伴随素理想，其中
$i \leq n$。设 $\mathfrak p$ 是元素 $x \in H^i_Z(M)$ 的零化子。
则 $\mathfrak p \in Z$，因而 $f \in \mathfrak p$。所以 $fx = 0$，进而 $x$
来自 $H^{i - 1}_Z(M/fM)$ 的某个元素；所用映射是上述长正合列中的边界映射
$\delta$。由此可知，$\mathfrak p$ 是有限模
$\Im(\delta)$ 的伴随素理想。因此 $\text{Ass}(H^i_Z(M))$ 对 $i \leq n$ 是有限的；
见《代数》，引理 \ref{algebra-lemma-finite-ass}。

\medskip\noindent
回忆，
$$
H^i_Z(M) \subset
\prod\nolimits_{\mathfrak p \in \text{Ass}(H^i_Z(M))}
H^i_Z(M)_\mathfrak p
$$
，这是《代数》，引理 \ref{algebra-lemma-zero-at-ass-zero} 的结论。按照假设，
右端各模均为有限的 $I$-幂挠模，故可以找到整数
$e_{\mathfrak p, i} \geq 0$，其中 $i \leq n$ 且
$\mathfrak p \in \text{Ass}(H^i_Z(M))$，使得 $I^{e_{\mathfrak p, i}}$ 零化
$H^i_Z(M)_\mathfrak p$。由此可得 $I^e$（其中
$e = \max\{e_{\mathfrak p, i}\}$）零化 $H^i_Z(M)$，对所有 $i \leq n$ 都如此。
由引理 \ref{local-cohomology-lemma-check-finiteness-local-cohomology-by-annihilator} 可知，
$H^i_Z(M)$ 对 $i \leq n$ 是有限的。
\end{proof}

\begin{lemma}
\label{local-cohomology-lemma-annihilate-local-cohomology}
设 $A$ 是环，$J \subset I \subset A$ 是有限生成理想，并设 $i \geq 0$ 是整数。
置 $Z = V(I)$。若 $H^i_Z(A)$ 被某个 $J^n$ 零化，其中 $n$ 是某个整数，则
$H^i_Z(M)$ 被某个 $J^m$ 零化，其中 $m = m(M)$；这里的结论适用于每个有限呈示
$A$-模 $M$；这里要求 $M_f$ 是有限局部自由 $A_f$-模，对所有 $f \in I$
都如此。
\end{lemma}

\begin{proof}
考虑零化子 $\mathfrak a$，它是 $H^i_Z(M)$ 的零化子。设 $\mathfrak p \subset A$ 且
$\mathfrak p \not \in Z$。按照假设，存在 $f \in I$，满足
$f \not \in \mathfrak p$，并存在同构
$\varphi : A_f^{\oplus r} \to M_f$；这是 $A_f$-模的同构。消去分母
（并使用 $M$ 为有限呈示这一事实），
得到映射
$$
a : A^{\oplus r} \longrightarrow M
\quad\text{且}\quad
b : M \longrightarrow A^{\oplus r}
$$
，使得 $a_f = f^N \varphi$ 且 $b_f = f^N \varphi^{-1}$，其中 $N$ 是某个整数。
此外，可以假设 $a \circ b$ 与 $b \circ a$ 都等于乘以 $f^{2N}$ 的映射。因此
$H^i_Z(M)$ 被 $f^{2N}J^n$ 零化，即 $f^{2N}J^n \subset \mathfrak a$。

\medskip\noindent
由于 $U = \Spec(A) \setminus Z$ 拟紧，可以找到有限多个
$f_1, \ldots, f_t$ 与 $N_1, \ldots, N_t$，使得 $U = \bigcup D(f_j)$ 且
$f_j^{2N_j}J^n \subset \mathfrak a$。于是 $V(I) = V(f_1, \ldots, f_t)$。
由于 $I$ 有限生成，可得 $I^M \subset (f_1, \ldots, f_t)$ 对某个 $M$ 成立。
综上可知，$J^m \subset \mathfrak a$ 对 $m \gg 0$ 成立；例如取
$m = M (2N_1 + \ldots + 2N_t) n$ 即可。
\end{proof}

\begin{lemma}
\label{local-cohomology-lemma-local-finiteness-for-finite-locally-free}
设 $A$ 是诺特环，$I \subset A$ 是理想。置 $Z = V(I)$。设 $n \geq 0$ 是整数。
若 $H^i_Z(A)$ 对 $0 \leq i \leq n$ 是有限的，则同样的结论对
$H^i_Z(M)$ 以及 $0 \leq i \leq n$ 成立，其中涉及任意有限 $A$-模 $M$；
这里要求 $M_f$ 是有限局部自由 $A_f$-模，对所有 $f \in I$ 都如此。
\end{lemma}

\begin{proof}
假设 $H^i_Z(A)$ 对 $0 \leq i \leq n$ 是有限的。由此存在 $e \geq 0$，使得
$I^e$ 零化 $H^i_Z(A)$，对所有 $0 \leq i \leq n$ 都如此；见引理
\ref{local-cohomology-lemma-check-finiteness-local-cohomology-by-annihilator}。于是引理
\ref{local-cohomology-lemma-annihilate-local-cohomology} 蕴含 $H^i_Z(M)$ 在
$0 \leq i \leq n$ 时被某个 $I^m$ 零化，其中 $m = m(M, i)$。可以取同一个
$m$，并让它对所有 $0 \leq i \leq n$ 都适用。再由引理
\ref{local-cohomology-lemma-check-finiteness-local-cohomology-by-annihilator}，可知
$H^i_Z(M)$ 对 $0 \leq i \leq n$ 是有限的，正是所需结论。
\end{proof}





\section{推前的有限性，I}
\label{local-cohomology-section-finiteness-pushforward}

\noindent
本节讨论有限性定理中最简单的非平凡情形，即第一局部上同调的有限性；等价地，
讨论 $j_*\mathcal{F}$ 的有限性，其中 $j : U \to X$ 是开浸入，$X$ 局部诺特，
而 $\mathcal{F}$ 是 $U$ 上的凝聚层。沿用 Koll\'ar 的方法
（\cite{Kollar-variants} 与 \cite{Kollar-local-global-hulls}），我们得到一个
充要条件；见命题 \ref{local-cohomology-proposition-kollar}。对高阶直接像或高阶局部上同调群感兴趣的
读者，可直接跳到第 \ref{local-cohomology-section-finiteness-pushforward-II} 节或第
\ref{local-cohomology-section-finiteness-II} 节（它们独立于本节其余内容展开）。

\begin{lemma}
\label{local-cohomology-lemma-check-finiteness-pushforward-on-associated-points}
设 $X$ 是局部诺特概形，$j : U \to X$ 是开子概形的嵌入，其补集为 $Z$。对
$x \in U$，设 $i_x : W_x \to U$ 是以 $x$ 为一般点的整闭子概形。设
$\mathcal{F}$ 是凝聚 $\mathcal{O}_U$-模。下列条件等价：
\begin{enumerate}
\item 对所有 $x \in \text{Ass}(\mathcal{F})$，$\mathcal{O}_X$-模
$j_*i_{x, *}\mathcal{O}_{W_x}$ 都是凝聚的；
\item $j_*\mathcal{F}$ 是凝聚的。
\end{enumerate}
\end{lemma}

\begin{proof}
先证明 (1) 蕴含 (2)。假设 (1) 成立。该断言在 $X$ 上是局部的，故可假设 $X$
仿射。于是 $U$ 拟紧，因而 $\text{Ass}(\mathcal{F})$ 有限（《除子》，引理
\ref{divisors-lemma-finite-ass}）。所以可以对伴随点的个数作归纳。设 $x \in U$
是 $\mathcal{F}$ 的支撑的某个不可约分支的一般点。由《除子》，引理
\ref{divisors-lemma-finite-ass}，有 $x \in \text{Ass}(\mathcal{F})$。由 $x$ 的
选取，$\dim(\mathcal{F}_x) = 0$，这里将其视为 $\mathcal{O}_{X, x}$-模。因此
$\mathcal{F}_x$ 作为 $\mathcal{O}_{X, x}$-模具有有限长度（《代数》，引理
\ref{algebra-lemma-support-point}）。所以可以对这个长度作归纳。

\medskip\noindent
置 $\mathcal{G} = j_*i_{x, *}\mathcal{O}_{W_x}$。按照假设，这是凝聚
$\mathcal{O}_X$-模。我们有 $\mathcal{G}_x = \kappa(x)$。选取非零映射
$\varphi_x : \mathcal{F}_x \to \kappa(x) = \mathcal{G}_x$。由《概形的上同调》，
引理 \ref{coherent-lemma-map-stalks-local-map}，存在开集 $x \in V \subset U$
以及映射 $\varphi_V : \mathcal{F}|_V \to \mathcal{G}|_V$，其在 $x$ 处的茎为
$\varphi_x$。选取 $f \in \Gamma(X, \mathcal{O}_X)$，使其在 $x$ 处不消失且
$D(f) \subset V$。例如由《概形的上同调》，引理
\ref{coherent-lemma-homs-over-open}，可知 $\varphi_V$ 延拓为
$f^n\mathcal{F} \to \mathcal{G}|_U$，其中 $n$ 是某个整数。与乘以 $f^n$ 前复合，
得到映射 $\mathcal{F} \to \mathcal{G}|_U$，其在 $x$ 处的茎非零。设
$\mathcal{F}' \subset \mathcal{F}$ 为其核。注意到
$\text{Ass}(\mathcal{F}') \subset \text{Ass}(\mathcal{F})$；见《除子》，引理
\ref{divisors-lemma-ses-ass}。由于
$\text{length}_{\mathcal{O}_{X, x}}(\mathcal{F}'_x) =
\text{length}_{\mathcal{O}_{X, x}}(\mathcal{F}_x) - 1$
，可使用归纳假设，得出 $j_*\mathcal{F}'$ 凝聚。由于
$\mathcal{G} = j_*(\mathcal{G}|_U) = j_*i_{x, *}\mathcal{O}_{W_x}$ 凝聚，
可以考虑正合列
$$
0 \to j_*\mathcal{F}' \to j_*\mathcal{F} \to \mathcal{G}
$$
由《概形》，引理 \ref{schemes-lemma-push-forward-quasi-coherent}，层
$j_*\mathcal{F}$ 拟凝聚。因此，由《概形的上同调》，引理
\ref{coherent-lemma-coherent-Noetherian-quasi-coherent-sub-quotient}，
$j_*\mathcal{F}$ 在 $j_*(\mathcal{G}|_U)$ 中的像凝聚。最后，由《概形的上同调》，
引理 \ref{coherent-lemma-coherent-abelian-Noetherian}，$j_*\mathcal{F}$ 凝聚。

\medskip\noindent
假设 (2) 成立。完全按照上述方式，可归约到 $X$ 仿射的情形。选取
$x \in \text{Ass}(\mathcal{F})$，并置
$\mathcal{G} = j_*i_{x, *}\mathcal{O}_{W_x}$。然后选取非零映射
$\varphi_x : \mathcal{G}_x = \kappa(x) \to \mathcal{F}_x$；它之所以存在，
恰是因为 $x$ 是 $\mathcal{F}$ 的伴随点。完全按照上述论证，可以假设
$\varphi_x$ 延拓为 $\mathcal{O}_U$-模映射
$\varphi : \mathcal{G}|_U \to \mathcal{F}$。于是 $\varphi$ 是单射（例如由
《除子》，引理 \ref{divisors-lemma-check-injective-on-ass}），从而得到单射
$\mathcal{G} = j_*(\mathcal{G}|_V) \to j_*\mathcal{F}$。因此 (1) 成立。
\end{proof}

\begin{lemma}
\label{local-cohomology-lemma-finiteness-pushforwards-and-H1-local}
设 $A$ 是诺特环，$I \subset A$ 是理想。置 $X = \Spec(A)$、$Z = V(I)$、
$U = X \setminus Z$，并以 $j : U \to X$ 表示嵌入态射。设 $\mathcal{F}$ 是凝聚
$\mathcal{O}_U$-模。则
\begin{enumerate}
\item 存在有限 $A$-模 $M$，使得 $\mathcal{F}$ 是 $\widetilde{M}$ 在 $U$ 上的限制；
\item 给定 $M$ 后，有正合列
$$
0 \to H^0_Z(M) \to M \to H^0(U, \mathcal{F}) \to H^1_Z(M) \to 0
$$
以及同构 $H^p(U, \mathcal{F}) = H^{p + 1}_Z(M)$，其中 $p \geq 1$；
\item 给定 $M$ 与 $p \geq 0$ 后，下列条件等价：
\begin{enumerate}
\item $R^pj_*\mathcal{F}$ 凝聚；
\item $H^p(U, \mathcal{F})$ 是有限 $A$-模；
\item $H^{p + 1}_Z(M)$ 是有限 $A$-模；
\end{enumerate}
\item 若 (3) 中的等价条件对 $p = 0$ 成立，则可以取
$M = \Gamma(U, \mathcal{F})$；此时 $H^0_Z(M) = H^1_Z(M) = 0$。
\end{enumerate}
\end{lemma}

\begin{proof}
由《性质》，引理 \ref{properties-lemma-lift-finite-presentation}，存在凝聚
$\mathcal{O}_X$-模 $\mathcal{F}'$，其在 $U$ 上的限制同构于 $\mathcal{F}$。
设 $\mathcal{F}'$ 对应于 (1) 中的有限 $A$-模 $M$。注意到
$R^pj_*\mathcal{F}$ 拟凝聚（《概形的上同调》，引理
\ref{coherent-lemma-quasi-coherence-higher-direct-images}），并对应于 $A$-模
$H^p(U, \mathcal{F})$。由引理
\ref{local-cohomology-lemma-local-cohomology-is-local-cohomology} 以及《上同调》第
\ref{cohomology-section-cohomology-support} 节和第
\ref{cohomology-section-cohomology-support-bis} 节中的讨论，得到正合列
$$
0 \to H^0_Z(M) \to M \to H^0(U, \mathcal{F}) \to H^1_Z(M) \to 0
$$
以及同构 $H^p(U, \mathcal{F}) = H^{p + 1}_Z(M)$，其中 $p \geq 1$。这里使用了
$H^j(X, \mathcal{F}') = 0$ 对 $j > 0$ 成立这一事实，因为 $X$ 仿射且
$\mathcal{F}'$ 拟凝聚（《概形的上同调》，引理
\ref{coherent-lemma-quasi-coherent-affine-cohomology-zero}）。这就证明了 (2)。
(3) 和 (4) 直接由 (2) 得出；另见引理 \ref{local-cohomology-lemma-local-cohomology}。
\end{proof}

\begin{lemma}
\label{local-cohomology-lemma-finiteness-pushforward}
设 $X$ 是局部诺特概形，$j : U \to X$ 是开子概形的嵌入，其补集为 $Z$。设
$\mathcal{F}$ 是凝聚 $\mathcal{O}_U$-模。假设
\begin{enumerate}
\item $X$ 是 Nagata 概形；
\item $X$ 是普遍链状的；以及
\item 对 $x \in \text{Ass}(\mathcal{F})$ 和
$z \in Z \cap \overline{\{x\}}$，有
$\dim(\mathcal{O}_{\overline{\{x\}}, z}) \geq 2$.
\end{enumerate}
则 $j_*\mathcal{F}$ 凝聚。
\end{lemma}

\begin{proof}
由引理 \ref{local-cohomology-lemma-check-finiteness-pushforward-on-associated-points}，只需证明
$j_*i_{x, *}\mathcal{O}_{W_x}$ 对 $x \in \text{Ass}(\mathcal{F})$ 凝聚。设
$\pi : Y \to X$ 是 $X$ 在 $\Spec(\kappa(x))$ 中的正规化；见《态射》，第
\ref{morphisms-section-normalization} 节。由《态射》，引理
\ref{morphisms-lemma-nagata-normalization-finite-general}，态射 $\pi$ 有限。由于
$\pi$ 有限，由《概形的上同调》，引理
\ref{coherent-lemma-finite-pushforward-coherent}，可知
$\mathcal{G} = \pi_*\mathcal{O}_Y$ 是凝聚 $\mathcal{O}_X$-模。注意到
$W_x = U \cap \pi(Y)$。因此
$\pi|_{\pi^{-1}(U)} : \pi^{-1}(U) \to U$ 通过 $i_x : W_x \to U$ 分解，
从而得到典范映射
$$
i_{x, *}\mathcal{O}_{W_x}
\longrightarrow
(\pi|_{\pi^{-1}(U)})_*(\mathcal{O}_{\pi^{-1}(U)}) =
(\pi_*\mathcal{O}_Y)|_U = \mathcal{G}|_U
$$
该映射是单射（例如由《除子》，引理
\ref{divisors-lemma-check-injective-on-ass}）。因此
$j_*i_{x, *}\mathcal{O}_{W_x} \subset j_*\mathcal{G}|_U$
，而只需证明 $j_*\mathcal{G}|_U$ 凝聚。

\medskip\noindent
还需证明 $j_*(\mathcal{G}|_U)$ 凝聚。我们断言《除子》，引理
\ref{divisors-lemma-depth-2-hartog} 适用于
$$
\mathcal{G} \longrightarrow j_*(\mathcal{G}|_U)
$$
，这便完成证明。只需证明 $\text{depth}(\mathcal{G}_z) \geq 2$ 对
$z \in Z$ 成立。设 $y_1, \ldots, y_n \in Y$ 是映到 $z$ 的各点。由《代数》，引理
\ref{algebra-lemma-depth-goes-down-finite}，只需证明
$\text{depth}(\mathcal{O}_{Y, y_i}) \geq 2$ 对 $i = 1, \ldots, n$ 成立。
若非如此，则由《性质》，引理 \ref{properties-lemma-criterion-normal}，有
$\dim(\mathcal{O}_{Y, y_i}) = 1$ 对某个 $i$ 成立。这与态射
$\pi : Y \to \overline{\{x\}}$ 的维数公式（《态射》，引理
\ref{morphisms-lemma-dimension-formula}）以及假设 (3) 矛盾。
\end{proof}

\begin{lemma}
\label{local-cohomology-lemma-sharp-finiteness-pushforward}
设 $X$ 是整的局部诺特概形，$j : U \to X$ 是非空开子概形的嵌入，其补集为 $Z$。
假设对所有 $z \in Z$ 以及任意伴随素理想 $\mathfrak p$（属于完备化
$\mathcal{O}_{X, z}^\wedge$），都有
$\dim(\mathcal{O}_{X, z}^\wedge/\mathfrak p) \geq 2$。
则 $j_*\mathcal{O}_U$ 凝聚。
\end{lemma}

\begin{proof}
可以假设 $X$ 仿射。使用引理
\ref{local-cohomology-lemma-check-finiteness-local-cohomology-locally} 与
\ref{local-cohomology-lemma-finiteness-pushforwards-and-H1-local}，可归约到
$X = \Spec(A)$，其中 $(A, \mathfrak m)$ 是诺特局部整环且
$\mathfrak m \in Z$。然后可以对 $d = \dim(A)$ 作归纳。（基始情形为
$d = 0, 1$；按照我们对局部环的假设，这些情形不会发生。）置
$V = \Spec(A) \setminus \{\mathfrak m\}$。注意到 $V$ 的各局部环维数严格小于
$d$。对 $j' : U \to V$ 重复上述论证并使用归纳，可得
$j'_*\mathcal{O}_U$ 是凝聚 $\mathcal{O}_V$-模。选取非零元素
$f \in A$，它在 $Z$ 上消失。由于 $D(f) \cap V \subset U$，可以找到 $n$，使得
乘以 $f^n$ 的映射（位于 $U$ 上）延拓为映射
$f^n : j'_*\mathcal{O}_U \to \mathcal{O}_V$；该映射位于 $V$ 上（例如由《概形的上同调》，引理
\ref{coherent-lemma-homs-over-open}）。该映射是单射，因而存在单射
$$
j_*\mathcal{O}_U = j''_* j'_* \mathcal{O}_U \to j''_*\mathcal{O}_V
$$
位于 $X$ 上，其中 $j'' : V \to X$ 是嵌入态射。因此只需证明
$j''_*\mathcal{O}_V$ 凝聚。换言之，可以假设 $X$ 是某个诺特局部整环的谱，
而 $Z$ 仅由闭点组成。

\medskip\noindent
假设 $X = \Spec(A)$，其中 $(A, \mathfrak m)$ 是局部环且
$Z = \{\mathfrak m\}$。设 $A^\wedge$ 为 $A$ 的完备化。置
$X^\wedge = \Spec(A^\wedge)$、$Z^\wedge = \{\mathfrak m^\wedge\}$、
$U^\wedge = X^\wedge \setminus Z^\wedge$，并置
$\mathcal{F}^\wedge = \mathcal{O}_{U^\wedge}$。环 $A^\wedge$ 普遍链状且为 Nagata 环
（《代数》，注记
\ref{algebra-remark-Noetherian-complete-local-ring-universally-catenary} 与引理
\ref{algebra-lemma-Noetherian-complete-local-Nagata}）。此外，按照假设，引理
\ref{local-cohomology-lemma-finiteness-pushforward} 的条件 (3) 对
$X^\wedge, Z^\wedge, U^\wedge, \mathcal{F}^\wedge$ 成立！因此
$(U^\wedge \to X^\wedge)_*\mathcal{O}_{U^\wedge}$ 凝聚。由于态射
$c : X^\wedge \to X$ 平坦，可得 $j_*\mathcal{O}_U$ 的拉回为
$(U^\wedge \to X^\wedge)_*\mathcal{O}_{U^\wedge}$（《概形的上同调》，引理
\ref{coherent-lemma-flat-base-change-cohomology}）。最后，由于 $c$ 忠实平坦，
由《下降》，引理 \ref{descent-lemma-finite-type-descends}，可得
$j_*\mathcal{O}_U$ 凝聚。
\end{proof}

\begin{remark}
\label{local-cohomology-remark-closure}
设 $j : U \to X$ 是局部诺特概形之间的开浸入，且 $x \in U$。设
$i_x : W_x \to U$ 是以 $x$ 为一般点的整闭子概形，并设 $\overline{\{x\}}$
是其在 $X$ 中的闭包。则有交换图
$$
\xymatrix{
W_x \ar[d]_{i_x} \ar[r]_{j'} & \overline{\{x\}} \ar[d]^i \\
U \ar[r]^j & X
}
$$
我们有 $j_*i_{x, *}\mathcal{O}_{W_x} = i_*j'_*\mathcal{O}_{W_x}$。由于左侧竖直
箭头是闭浸入，可知 $j_*i_{x, *}\mathcal{O}_{W_x}$ 凝聚，当且仅当
$j'_*\mathcal{O}_{W_x}$ 凝聚。
\end{remark}

\begin{remark}
\label{local-cohomology-remark-no-finiteness-pushforward}
设 $X$ 是局部诺特概形，$j : U \to X$ 是开子概形的嵌入，其补集为 $Z$。设
$\mathcal{F}$ 是凝聚 $\mathcal{O}_U$-模。若存在
$x \in \text{Ass}(\mathcal{F})$ 以及 $z \in Z \cap \overline{\{x\}}$，使得
$\dim(\mathcal{O}_{\overline{\{x\}}, z}) \leq 1$，则 $j_*\mathcal{F}$ 不凝聚。
为证明这一点，可以作平坦基变换到 $\mathcal{O}_{X, z}$ 的谱。置
$X' = \overline{\{x\}}$。假设蕴含
$\mathcal{O}_{X' \cap U} \subset \mathcal{F}$。因此只需说明
$j_*\mathcal{O}_{X' \cap U}$ 不凝聚。这是显然的，因为 $X' = \{x, z\}$，所以
$j_*\mathcal{O}_{X' \cap U}$ 对应于 $\kappa(x)$，这里将其视为
$\mathcal{O}_{X, z}$-模；
由于 $x$ 不是闭点，它不可能是有限模。

\medskip\noindent
事实上，引理 \ref{local-cohomology-lemma-sharp-finiteness-pushforward} 的逆命题也成立：给定整诺特概形
之间的开浸入 $j : U \to X$，若存在 $z \in X \setminus U$ 以及伴随素理想
$\mathfrak p$（属于完备化 $\mathcal{O}_{X, z}^\wedge$），满足
$\dim(\mathcal{O}_{X, z}^\wedge/\mathfrak p) = 1$，则
$j_*\mathcal{O}_U$ 不凝聚。事实上，可以转到局部环，把 $U$ 扩大为穿孔谱，再转到
完备化；随后上述论证便给出非有限性。
\end{remark}

\begin{proposition}[Koll\'ar]
\label{local-cohomology-proposition-kollar}
\begin{reference}
见 \cite{k-coherent}；特殊情形见 \cite[IV, Proposition 7.2.2]{EGA}。
\end{reference}
\begin{slogan}
Hartogs 定理的弱类似：在诺特概形上，若一个开集的补集在层的支撑中余维为 2，
则凝聚层在该开集上的限制凝聚。
\end{slogan}
设 $j : U \to X$ 是局部诺特概形之间的开浸入，其补集为 $Z$。设 $\mathcal{F}$
是凝聚 $\mathcal{O}_U$-模。下列条件等价：
\begin{enumerate}
\item $j_*\mathcal{F}$ 凝聚；
\item 对 $x \in \text{Ass}(\mathcal{F})$、
$z \in Z \cap \overline{\{x\}}$，以及任意伴随素理想
$\mathfrak p$（属于完备化 $\mathcal{O}_{\overline{\{x\}}, z}^\wedge$），都有
$\dim(\mathcal{O}_{\overline{\{x\}}, z}^\wedge/\mathfrak p) \geq 2$。
\end{enumerate}
\end{proposition}

\begin{proof}
若 (2) 成立，则结合引理
\ref{local-cohomology-lemma-check-finiteness-pushforward-on-associated-points}、注记
\ref{local-cohomology-remark-closure} 与引理 \ref{local-cohomology-lemma-sharp-finiteness-pushforward}，可得 (1)。
若 (2) 不成立，则由注记 \ref{local-cohomology-remark-no-finiteness-pushforward} 中的讨论
（以及注记 \ref{local-cohomology-remark-closure}），$j_*i_{x, *}\mathcal{O}_{W_x}$ 对某个
$x \in \text{Ass}(\mathcal{F})$ 不是有限的。因此由引理
\ref{local-cohomology-lemma-check-finiteness-pushforward-on-associated-points}，
$j_*\mathcal{F}$ 不凝聚。
\end{proof}

\begin{lemma}
\label{local-cohomology-lemma-kollar-finiteness-H1-local}
设 $A$ 是诺特环，$I \subset A$ 是理想。置 $Z = V(I)$。设 $M$ 是有限 $A$-模。
下列条件等价：
\begin{enumerate}
\item $H^1_Z(M)$ 是有限 $A$-模；以及
\item 对所有 $\mathfrak p \in \text{Ass}(M)$，若 $\mathfrak p \not \in Z$，
则对所有 $\mathfrak q \in V(\mathfrak p + I)$，$(A/\mathfrak p)_\mathfrak q$
的完备化均不具有维数为 $1$ 的伴随素理想。
\end{enumerate}
\end{lemma}

\begin{proof}
这立即由命题 \ref{local-cohomology-proposition-kollar} 经引理
\ref{local-cohomology-lemma-finiteness-pushforwards-and-H1-local} 得出。
\end{proof}

\noindent
下一个引理的表述有一个优点：条件 (1) 和 (2) 可由 $X$ 上的有限型概形继承。
此外，这正是我们将在第 \ref{local-cohomology-section-finiteness-pushforward-II} 节推广到高阶直接像的
有限性形式。

\begin{lemma}
\label{local-cohomology-lemma-finiteness-pushforward-general}
设 $X$ 是局部诺特概形，$j : U \to X$ 是开子概形的嵌入，其补集为 $Z$。设
$\mathcal{F}$ 是凝聚 $\mathcal{O}_U$-模。假设
\begin{enumerate}
\item $X$ 普遍链状；
\item 对每个 $z \in Z$，$\mathcal{O}_{X, z}$ 的形式纤维满足 $(S_1)$。
\end{enumerate}
在这种情形下，下列条件等价：
\begin{enumerate}
\item[(a)] 对 $x \in \text{Ass}(\mathcal{F})$ 以及
$z \in Z \cap \overline{\{x\}}$，有
$\dim(\mathcal{O}_{\overline{\{x\}}, z}) \geq 2$；以及
\item[(b)] $j_*\mathcal{F}$ 凝聚。
\end{enumerate}
\end{lemma}

\begin{proof}
设 $x \in \text{Ass}(\mathcal{F})$。由命题 \ref{local-cohomology-proposition-kollar}，只需检验
$A = \mathcal{O}_{\overline{\{x\}}, z}$ 满足该命题中关于其完备化的伴随素理想的
条件，当且仅当 $\dim(A) \geq 2$。注意到 $A$ 普遍链状（这是显然的），并且其形式
纤维满足 $(S_1)$；这由《代数进阶》，引理
\ref{more-algebra-lemma-formal-fibres-normal} 与命题
\ref{more-algebra-proposition-finite-type-over-P-ring} 得出。设
$\mathfrak p' \subset A^\wedge$ 是伴随素理想。由于 $A \to A^\wedge$ 平坦，
由《代数》，引理 \ref{algebra-lemma-bourbaki}，可知 $\mathfrak p'$ 位于
$(0) \subset A$ 之上。形式纤维 $A^\wedge \otimes_A F$ 满足 $(S_1)$，其中 $F$
是 $A$ 的分式域。因此 $\mathfrak p'$ 是极小素理想；见《代数》，引理
\ref{algebra-lemma-criterion-no-embedded-primes}。由于 $A$ 普遍链状，由《代数进阶》，
命题 \ref{more-algebra-proposition-ratliff}，它形式链状。因此
$\dim(A^\wedge/\mathfrak p') = \dim(A)$，这就证明了所述等价性。
\end{proof}






\section{深度与维数}
\label{local-cohomology-section-dept-dimension}

\noindent
先给出若干辅助引理。

\begin{lemma}
\label{local-cohomology-lemma-ideal-depth-function}
设 $A$ 是诺特环，$I \subset A$ 是理想，$M$ 是有限 $A$-模。设
$\mathfrak p \in V(I)$ 是素理想，并假设
$e = \text{depth}_{IA_\mathfrak p}(M_\mathfrak p) < \infty$.
则存在非空开子集 $U \subset V(\mathfrak p)$，使得
$\text{depth}_{IA_\mathfrak q}(M_\mathfrak q) \geq e$
对所有 $\mathfrak q \in U$ 都成立。
\end{lemma}

\begin{proof}
由深度的定义，$IM_\mathfrak p \not = M_\mathfrak p$，且存在
$M_\mathfrak p$-正则序列
$f_1, \ldots, f_e \in IA_\mathfrak p$。将 $A$ 换成某个主局部化后，可以假设
$f_1, \ldots, f_e \in I$ 构成 $M$-正则序列；见
《代数》，引理 \ref{algebra-lemma-regular-sequence-in-neighbourhood}。
考虑模 $M' = M/IM$。由于 $\mathfrak p \in \text{Supp}(M')$，且有限模的支撑是闭集，
可得 $V(\mathfrak p) \subset \text{Supp}(M')$。因此，对
$\mathfrak q \in V(\mathfrak p)$，有 $IM_\mathfrak q \not = M_\mathfrak q$。
再利用局部化的正合性，由深度的定义可知
$\text{depth}_{IA_\mathfrak q}(M_\mathfrak q) \geq e$
对任意 $\mathfrak q \in V(I)$ 都成立。
\end{proof}

\begin{lemma}
\label{local-cohomology-lemma-depth-function}
设 $A$ 是诺特环，$M$ 是有限 $A$-模，$\mathfrak p$ 是素理想。假设
$e = \text{depth}_{A_\mathfrak p}(M_\mathfrak p) < \infty$.
则存在非空开子集 $U \subset V(\mathfrak p)$，使得
$\text{depth}_{A_\mathfrak q}(M_\mathfrak q) \geq e$
对所有 $\mathfrak q \in U$ 都成立，且除有限多个 $\mathfrak q \in U$ 外均有
$\text{depth}_{A_\mathfrak q}(M_\mathfrak q) > e$。
\end{lemma}

\begin{proof}
由深度的定义，$\mathfrak p M_\mathfrak p \not = M_\mathfrak p$，且存在
$M_\mathfrak p$-正则序列
$f_1, \ldots, f_e \in \mathfrak p A_\mathfrak p$。将 $A$ 换成某个主局部化后，
可以假设 $f_1, \ldots, f_e \in \mathfrak p$ 构成 $M$-正则序列；见
《代数》，引理 \ref{algebra-lemma-regular-sequence-in-neighbourhood}。
考虑模 $M' = M/(f_1, \ldots, f_e)M$。由于
$\mathfrak p \in \text{Supp}(M')$，且有限模的支撑是闭集，
可得 $V(\mathfrak p) \subset \text{Supp}(M')$。因此，对
$\mathfrak q \in V(\mathfrak p)$，有
$\mathfrak q M_\mathfrak q \not = M_\mathfrak q$。再利用局部化的正合性，
由深度的定义可知
$\text{depth}_{A_\mathfrak q}(M_\mathfrak q) \geq e$
对任意 $\mathfrak q \in V(I)$ 都成立。此外，只要 $\mathfrak q$ 不是模 $M'$ 的伴随素理想，深度就会增大。
于是最后一个断言由《代数》，引理 \ref{algebra-lemma-finite-ass} 得出。
\end{proof}

\begin{lemma}
\label{local-cohomology-lemma-finite-nr-points-next-S}
设 $X$ 是具有对偶复形 $\omega_X^\bullet$ 的诺特概形，
$\mathcal{F}$ 是凝聚 $\mathcal{O}_X$-模，$k \geq 0$ 是整数。假设
$\mathcal{F}$ 满足 $(S_k)$。则满足下式的点 $x \in X$ 只有有限多个：
$$
\text{depth}(\mathcal{F}_x) = k
\quad\text{且}\quad
\dim(\text{Supp}(\mathcal{F}_x)) > k
$$
\end{lemma}

\begin{proof}
对 $k$ 作归纳来证明。基础情形 $k = 0$ 断言 $\mathcal{F}$
只有有限多个嵌入伴随点；这由《除子》，引理
\ref{divisors-lemma-finite-ass} 得出。

\medskip\noindent
假设 $k > 0$，且结论对所有更小的 $k$ 都成立。
可以用有限多个仿射开集覆盖 $X$，故可假设 $X = \Spec(A)$ 是仿射的。
于是 $\mathcal{F}$ 是凝聚 $\mathcal{O}_X$-模；它对应于有限 $A$-模 $M$，
且该模满足 $(S_k)$。下文将直接使用《代数》，引理
\ref{algebra-lemma-one-equation-module} 和
\ref{algebra-lemma-depth-drops-by-one}，不再另行说明。

\medskip\noindent
设 $f \in A$ 是 $M$ 上的非零因子。则 $M/fM$ 满足 $(S_{k - 1})$。
由归纳假设，满足下列条件的素理想 $\mathfrak p \in V(f)$ 只有有限多个：
$\text{depth}((M/fM)_\mathfrak p) = k - 1$，且
$\dim(\text{Supp}((M/fM)_\mathfrak p)) > k - 1$.
这些恰好是满足下列条件的素理想 $\mathfrak p \in V(f)$：
$\text{depth}(M_\mathfrak p) = k$，且
$\dim(\text{Supp}(M_\mathfrak p)) > k$.
因此，在证明该有限性断言时，可以将 $A$ 换成 $A_f$，将 $M$ 换成 $M_f$。

\medskip\noindent
由于 $M$ 满足 $(S_k)$ 且 $k > 0$，可知 $M$ 没有嵌入伴随素理想
（《代数》，引理 \ref{algebra-lemma-criterion-no-embedded-primes}）。
因此 $\text{Ass}(M)$ 是 $M$ 的支撑的泛点集。于是《对偶复形》，引理
\ref{dualizing-lemma-CM-open} 表明集合
$U = \{\mathfrak q \mid M_\mathfrak q\text{ 为 Cohen--Macaulay}\}$
是包含 $\text{Ass}(M)$ 的开集。由素避（《代数》，引理
\ref{algebra-lemma-silly}），可以选取 $f \in A$，使得
$f \not \in \mathfrak p$ 对所有 $\mathfrak p \in \text{Ass}(M)$ 成立，并且
$D(f) \subset U$。于是 $f$ 是 $M$ 上的非零因子
（《代数》，引理 \ref{algebra-lemma-ass-zero-divisors}）。
如上将 $A$ 换成 $A_f$、将 $M$ 换成 $M_f$ 后，可知 $M$ 是
Cohen--Macaulay 模。因此，对所有 $\mathfrak q \subset A$，有
$\dim(M_\mathfrak q) = \text{depth}(M_\mathfrak q)$
，故引理所述集合为空，因而当然是有限的。
\end{proof}

\begin{lemma}
\label{local-cohomology-lemma-sitting-in-degrees}
设 $(A, \mathfrak m)$ 是具有规范化对偶复形
$\omega_A^\bullet$ 的诺特局部环，$M$ 是有限 $A$-模。
记 $E^i = \text{Ext}_A^{-i}(M, \omega_A^\bullet)$。则
\begin{enumerate}
\item $E^i$ 是有限 $A$-模，并且仅在
$0 \leq i \leq \dim(\text{Supp}(M))$,
时可能非零；
\item $\dim(\text{Supp}(E^i)) \leq i$；
\item $\text{depth}(M)$ 是满足下列条件的最小整数 $\delta \geq 0$：
$E^\delta \not = 0$；
\item $\mathfrak p \in \text{Supp}(E^0 \oplus \ldots \oplus E^i)
\Leftrightarrow
\text{depth}_{A_\mathfrak p}(M_\mathfrak p) + \dim(A/\mathfrak p) \leq i$；
\item $E^i$ 的零化理想等于 $H^i_\mathfrak m(M)$ 的零化理想。
\end{enumerate}
\end{lemma}

\begin{proof}
第 (1)、(2)、(3) 项即《对偶复形》，引理
\ref{dualizing-lemma-sitting-in-degrees} 中的断言。
对素理想 $\mathfrak p$（属于 $A$），复形
$(\omega_A^\bullet)_\mathfrak p[-\dim(A/\mathfrak p)]$
是 $A_\mathfrak p$ 的规范化对偶复形；见《对偶复形》，引理
\ref{dualizing-lemma-dimension-function}。因此
$$
E^i_\mathfrak p =
\text{Ext}^{-i}_A(M, \omega_A^\bullet)_\mathfrak p =
\text{Ext}^{-i + \dim(A/\mathfrak p)}_{A_\mathfrak p}
(M_\mathfrak p, (\omega_A^\bullet)_\mathfrak p[-\dim(A/\mathfrak p)])
$$
在
$i - \dim(A/\mathfrak p) < \text{depth}_{A_\mathfrak p}(M_\mathfrak p)$
时为零，而在
$i = \dim(A/\mathfrak p) + \text{depth}_{A_\mathfrak p}(M_\mathfrak p)$
时非零；这是把第 (3) 项用于 $A_\mathfrak p$ 所得。因此第 (4) 项成立。
若 $E$ 是 $A$ 的剩余域的内射包，则
$$
\Hom_A(H^i_\mathfrak m(M), E) =
\text{Ext}^{-i}_A(M, \omega_A^\bullet)^\wedge =
(E^i)^\wedge =
E^i \otimes_A A^\wedge
$$
；这由局部对偶定理得出（采用《对偶复形》，引理
\ref{dualizing-lemma-special-case-local-duality} 的形式）。
由于 $A \to A^\wedge$ 忠实平坦，第 (5) 项由 Matlis 对偶性得出
（《对偶复形》，命题 \ref{dualizing-proposition-matlis}）。
\end{proof}






\section{局部上同调的零化理想，I}
\label{local-cohomology-section-annihilators}

\noindent
本节讨论 Faltings 的一个结果；见 \cite{Faltings-annulators}。

\begin{proposition}
\label{local-cohomology-proposition-annihilator}
\begin{reference}
\cite{Faltings-annulators}.
\end{reference}
设 $A$ 是具有对偶复形的诺特环。设
$T \subset T' \subset \Spec(A)$ 是在特化下稳定的子集。
设 $s \geq 0$ 是整数，$M$ 是有限 $A$-模。下列条件等价：
\begin{enumerate}
\item 存在理想 $J \subset A$，满足 $V(J) \subset T'$，且
$J$ 将 $H^i_T(M)$ 对所有 $i \leq s$ 零化；
\item 对所有满足 $\mathfrak p \not \in T'$、
$\mathfrak q \in T$ 以及 $\mathfrak p \subset \mathfrak q$ 的情形，都有
$$
\text{depth}_{A_\mathfrak p}(M_\mathfrak p) +
\dim((A/\mathfrak p)_\mathfrak q) > s
$$
\end{enumerate}
\end{proposition}

\begin{proof}
设 $\omega_A^\bullet$ 是对偶复形，$\delta$ 是其维数函数；见《对偶复形》，
第 \ref{dualizing-section-dimension-function} 节。下列有限 $A$-模将起重要作用：
$$
E^i = \Ext_A^i(M, \omega_A^\bullet)
$$
对 $\mathfrak p \subset A$，记 $H^i_\mathfrak p$ 为 $A_\mathfrak p$-模关于
$\mathfrak pA_\mathfrak p$ 的局部上同调。于是下式的
$\mathfrak pA_\mathfrak p$-进完备化
$$
(E^i)_\mathfrak p =
\Ext^{\delta(\mathfrak p) + i}_{A_\mathfrak p}(M_\mathfrak p,
(\omega_A^\bullet)_\mathfrak p[-\delta(\mathfrak p)])
$$
与下式 Matlis 对偶：
$$
H^{-\delta(\mathfrak p) - i}_{\mathfrak p}(M_\mathfrak p)
$$
这由《对偶复形》，引理 \ref{dualizing-lemma-special-case-local-duality} 得出。
特别地，由此得到如下事实：理想 $J \subset A$ 将
$(E^i)_\mathfrak p$ 零化，当且仅当 $J$ 将
$H^{-\delta(\mathfrak p) - i}_{\mathfrak p}(M_\mathfrak p)$ 零化。

\medskip\noindent
置 $T_n = \{\mathfrak p \in T \mid \delta(\mathfrak p) \leq n\}$。
由于 $\delta$ 是有界函数，可知 $T_a = \emptyset$ 在 $a \ll 0$ 时成立，
而 $T_b = T$ 在 $b \gg 0$ 时成立。

\medskip\noindent
假设 (2) 成立。来证明 (1) 中 $J$ 的存在性。为此使用双重归纳。
对 $i \leq s$，考虑归纳假设 $IH_i$：$H^a_T(M)$ 被某个满足
$J \subset A$ 与 $V(J) \subset T'$ 的理想零化，其中
$0 \leq a \leq i$。情形 $IH_0$ 是显然的，
因为 $H^0_T(M)$ 是 $M$ 的子模，从而是有限的；因此它被某个满足
$J$ 及 $V(J) \subset T$ 的理想零化。

\medskip\noindent
归纳步骤。假设 $IH_{i - 1}$ 对某个 $0 < i \leq s$ 成立。
选取 $J'$，满足 $V(J') \subset T'$，使其将 $H^a_T(M)$
在 $0 \leq a \leq i - 1$ 时零化（归纳假设保证可以如此选取）。
将对 $n$ 作降归纳，证明存在理想 $J$，满足 $V(J) \subset T'$，
使 $J H^i_T(M)$ 的伴随素理想都属于 $T_n$。
当 $n \ll 0$ 时，这蕴含 $JH^i_T(M) = 0$
（《代数》，引理 \ref{algebra-lemma-ass-zero}），因而 $IH_i$ 成立。
基础情形 $n \gg 0$ 是显然的，因为此时 $T = T_n$，而
$H^i_T(M)$ 的所有伴随素理想都属于 $T$。

\medskip\noindent
于是设已给定 $J$，它对 $n$ 具有上述性质。取 $\mathfrak q \in T_n$。
令 $T_\mathfrak q \subset \Spec(A_\mathfrak q)$ 为 $T$ 的逆像。我们有
$H^j_T(M)_\mathfrak q = H^j_{T_\mathfrak q}(M_\mathfrak q)$
；见引理 \ref{local-cohomology-lemma-torsion-change-rings}。考虑谱序列
$$
H_\mathfrak q^p(H^q_{T_\mathfrak q}(M_\mathfrak q))
\Rightarrow
H^{p + q}_\mathfrak q(M_\mathfrak q)
$$
，它来自引理 \ref{local-cohomology-lemma-local-cohomology-ss}。下文将找到理想
$J'' \subset A$，满足 $V(J'') \subset T'$，使得
$H^i_\mathfrak q(M_\mathfrak q)$ 被 $J''$ 零化；这对所有
$\mathfrak q \in T_n \setminus T_{n - 1}$ 都成立。
断言：$J (J')^i J''$ 对 $n - 1$ 可行。事实上，取
$\mathfrak q \in T_n \setminus T_{n - 1}$。上述谱序列给出滤过
$$
E_\infty^{0, i} = E_{i + 2}^{0, i} \subset \ldots \subset E_3^{0, i} \subset
E_2^{0, i} = H^0_\mathfrak q(H^i_{T_\mathfrak q}(M_\mathfrak q))
$$
。模 $E_\infty^{0, i}$ 被 $J''$ 零化。子商
$E_j^{0, i}/E_{j + 1}^{0, i}$ 在 $i + 1 \geq j \geq 2$ 时被 $J'$ 零化，
因为 $d_j^{0, i}$ 的靶是下式的一个子商：
$$
H^j_\mathfrak q(H^{i - j + 1}_{T_\mathfrak q}(M_\mathfrak q)) =
H^j_\mathfrak q(H^{i - j + 1}_T(M)_\mathfrak q)
$$
，而 $H^{i - j + 1}_T(M)_\mathfrak q$ 被 $J'$ 零化；这是由 $J'$ 的选取得出的。
最后，由 $J$ 的选取，有
$J H^i_T(M)_\mathfrak q \subset H^0_\mathfrak q(H^i_T(M)_\mathfrak q)$
，因为 $\Spec(A_\mathfrak q)$ 的非闭点具有更大的 $\delta$ 值。因此
$\mathfrak q$ 不可能是 $J(J')^iJ'' H^i_T(M)$ 的伴随素理想，正如所需。

\medskip\noindent
由开头的讨论可知，$J''$ 应将
$$
(E^{-\delta(\mathfrak q) - i})_\mathfrak q =
(E^{-n - i})_\mathfrak q
$$
对所有 $\mathfrak q \in T_n \setminus T_{n - 1}$ 零化。
但若 $J''$ 对某个 $\mathfrak q$ 可行，则它对 $\mathfrak q$ 的一个开邻域中的所有
$\mathfrak q$ 都可行，因为模 $E^{-n - i}$ 是有限的。
由于 $\Spec(A)$ 的每个子集在诱导拓扑下都是诺特的（《拓扑》，引理
\ref{topology-lemma-Noetherian}），可知只需证明 $J''$ 对某个
$\mathfrak q$ 的存在性。

\medskip\noindent
由于这些 Ext 模是有限的，$J''$ 的存在性等价于
$$
\text{Supp}(E^{-n - i}) \cap \Spec(A_\mathfrak q) \subset T'.
$$
。这又等价于证明 $E^{-n - i}$ 在每个满足
$\mathfrak p \subset \mathfrak q$、$\mathfrak p \not \in T'$ 的点处的局部化为零。
在 $A_\mathfrak p$ 上使用局部对偶性，可知需要证明
$$
H^{i + n - \delta(\mathfrak p)}_\mathfrak p(M_\mathfrak p) =
H^{i - \dim((A/\mathfrak p)_\mathfrak q)}_\mathfrak p(M_\mathfrak p)
$$
为零（这里用到 $\delta$ 是维数函数）。由命题中的假设、$i \leq s$ 以及
《对偶复形》，引理 \ref{dualizing-lemma-depth}，它确实为零。

\medskip\noindent
为证明逆向蕴含，假设 (2) 不成立，并反向使用上述论证。首先选取
$\mathfrak q \in T$、$\mathfrak p \subset \mathfrak q$，满足
$\mathfrak p \not \in T'$，使得
$$
i = \text{depth}_{A_\mathfrak p}(M_\mathfrak p) +
\dim((A/\mathfrak p)_\mathfrak q) \leq s
$$
取最小值。于是
$H^{i - \dim((A/\mathfrak p)_\mathfrak q)}_\mathfrak p(M_\mathfrak p)$
非零；这由《对偶复形》，引理 \ref{dualizing-lemma-depth} 中的非消失性得出。
置 $n = \delta(\mathfrak q)$。此时不存在满足
$J \subset A$ 与 $V(J) \subset T'$ 的理想，使得
$J(E^{-n - i})_\mathfrak q = 0$。因此 $H^i_\mathfrak q(M_\mathfrak q)$
不被任何满足 $J \subset A$ 与 $V(J) \subset T'$ 的理想零化。
由 $i$ 的极小性和上面展示的谱序列，模 $H^i_T(M)_\mathfrak q$
不被任何满足 $J \subset A$ 与 $V(J) \subset T'$ 的理想零化。
因此 $H^i_T(M)$ 不被任何满足 $J \subset A$ 与 $V(J) \subset T'$
的理想零化。命题得证。
\end{proof}

\begin{lemma}
\label{local-cohomology-lemma-kill-local-cohomology-at-prime}
设 $I$ 是诺特环 $A$ 的理想，$M$ 是有限 $A$-模，
$\mathfrak p \subset A$ 是素理想，$s \geq 0$ 是整数。假设
\begin{enumerate}
\item $A$ 具有对偶复形；
\item $\mathfrak p \not \in V(I)$；
\item 对所有素理想 $\mathfrak p' \subset \mathfrak p$ 以及
$\mathfrak q \in V(I)$，若 $\mathfrak p' \subset \mathfrak q$，则有
$$
\text{depth}_{A_{\mathfrak p'}}(M_{\mathfrak p'}) +
\dim((A/\mathfrak p')_\mathfrak q) > s
$$
\end{enumerate}
则存在 $f \in A$，满足 $f \not \in \mathfrak p$，并将
$H^i_{V(I)}(M)$ 对所有 $i \leq s$ 零化。
\end{lemma}

\begin{proof}
考虑集合
$$
T = V(I)
\quad\text{且}\quad
T' = \bigcup\nolimits_{f \in A, f \not \in \mathfrak p} V(f)
$$
它们是在特化下稳定的 $\Spec(A)$ 的子集。注意
$T \subset T'$ 且 $\mathfrak p \not \in T'$。假设 (3) 表明命题
\ref{local-cohomology-proposition-annihilator} 的条件 (2) 成立。因此可找到
$J \subset A$，满足 $V(J) \subset T'$，使得
$J H^i_{V(I)}(M) = 0$ 对所有 $i \leq s$ 成立。选取
$f \in A$，满足 $f \not \in \mathfrak p$ 与 $V(J) \subset V(f)$。
$f$ 的某个幂将 $H^i_{V(I)}(M)$ 对所有 $i \leq s$ 零化。
\end{proof}





\section{局部上同调的有限性，II}
\label{local-cohomology-section-finiteness-II}

\noindent
继续第 \ref{local-cohomology-section-finiteness} 节开始的局部上同调有限性讨论。
利用 Faltings 零化理想定理，很容易证明下列基本结果。

\begin{proposition}
\label{local-cohomology-proposition-finiteness}
\begin{reference}
\cite{Faltings-annulators}.
\end{reference}
设 $A$ 是具有对偶复形的诺特环，
$T \subset \Spec(A)$ 是在特化下稳定的子集。
设 $s \geq 0$ 是整数，$M$ 是有限 $A$-模。下列条件等价：
\begin{enumerate}
\item $H^i_T(M)$ 是有限 $A$-模，其中 $i \leq s$；
\item 对所有满足 $\mathfrak p \not \in T$、$\mathfrak q \in T$ 以及
$\mathfrak p \subset \mathfrak q$ 的情形，都有
$$
\text{depth}_{A_\mathfrak p}(M_\mathfrak p) +
\dim((A/\mathfrak p)_\mathfrak q) > s
$$
\end{enumerate}
\end{proposition}

\begin{proof}
这是命题 \ref{local-cohomology-proposition-annihilator} 和引理
\ref{local-cohomology-lemma-check-finiteness-local-cohomology-by-annihilator} 的形式推论。
\end{proof}

\noindent
除给出一些供后文使用的引理外，本节余下部分研究命题
\ref{local-cohomology-proposition-finiteness} 中 $A$ 具有对偶复形这一条件能够弱化到何种程度。
粗略地说，答案是必须假设 $A$ 的形式纤维满足 $(S_n)$，其中 $n$ 充分大。

\medskip\noindent
设 $A$ 是诺特环，$I \subset A$ 是理想。置
$X = \Spec(A)$ 以及 $Z = V(I) \subset X$，并设 $M$ 是有限 $A$-模。定义
\begin{equation}
\label{local-cohomology-equation-cutoff}
s_{A, I}(M) =
\min \{
\text{depth}_{A_\mathfrak p}(M_\mathfrak p) + \dim((A/\mathfrak p)_\mathfrak q)
\mid
\mathfrak p \in X \setminus Z, \mathfrak q \in Z,
\mathfrak p \subset \mathfrak q
\}
\end{equation}
关于深度，我们约定 $0$ 的深度为 $\infty$；因此只需考虑素理想
$\mathfrak p$，且它属于 $M$ 的支撑。可以看出，$s_{A, I}(M)$
是这一情形的一个重要不变量。

\begin{lemma}
\label{local-cohomology-lemma-cutoff}
设 $A \to B$ 是诺特环之间的有限同态，$I \subset A$ 是理想，并置
$J = IB$。设 $M$ 是有限 $B$-模。若 $A$ 普遍链状，则
$s_{B, J}(M) = s_{A, I}(M)$.
\end{lemma}

\begin{proof}
设 $\mathfrak p \subset \mathfrak q \subset A$ 是满足
$I \subset \mathfrak q$ 和 $I \not \subset \mathfrak p$ 的素理想。
由于 $A \to B$ 是有限的，素理想 $\mathfrak p_i$（位于
$\mathfrak p$ 上方）只有有限多个。由《代数》，引理
\ref{algebra-lemma-depth-goes-down-finite}，有
$$
\text{depth}(M_\mathfrak p) = \min \text{depth}(M_{\mathfrak p_i})
$$
设 $\mathfrak p_i \subset \mathfrak q_{ij}$ 是位于 $\mathfrak q$ 上方的素理想。
由 $A \to B$ 的上升定理（《代数》，引理
\ref{algebra-lemma-integral-going-up}），至少存在一个 $\mathfrak q_{ij}$；
这对每个 $i$ 都成立。于是
$$
\dim((B/\mathfrak p_i)_{\mathfrak q_{ij}}) =
\dim((A/\mathfrak p)_\mathfrak q)
$$
；这由维数公式得出，见《代数》，引理
\ref{algebra-lemma-dimension-formula}。因此，用于定义
$s_{B, J}(M)$ 的量在各对 $(\mathfrak p_i, \mathfrak q_{ij})$ 上的最小值，
等于该量在数对 $(\mathfrak p, \mathfrak q)$ 上的值。引理得证。
\end{proof}

\begin{lemma}
\label{local-cohomology-lemma-change-completion}
设 $A$ 是具有对偶复形的诺特环，$I \subset A$ 是理想，
$M$ 是有限 $A$-模。设 $A', M'$ 是 $I$-进完备化，分别对应于 $A, M$。
设 $\mathfrak p' \subset \mathfrak q'$ 是 $A'$ 的素理想，其中
$\mathfrak q' \in V(IA')$，并且它们位于
$\mathfrak p \subset \mathfrak q$（在 $A$ 中）上方。则
$$
\text{depth}_{A_{\mathfrak p'}}(M'_{\mathfrak p'})
\geq
\text{depth}_{A_\mathfrak p}(M_\mathfrak p)
$$
并且
$$
\text{depth}_{A_{\mathfrak p'}}(M'_{\mathfrak p'}) +
\dim((A'/\mathfrak p')_{\mathfrak q'}) =
\text{depth}_{A_\mathfrak p}(M_\mathfrak p) +
\dim((A/\mathfrak p)_\mathfrak q)
$$
\end{lemma}

\begin{proof}
我们有
$$
\text{depth}(M'_{\mathfrak p'}) =
\text{depth}(M_\mathfrak p) +
\text{depth}(A'_{\mathfrak p'}/\mathfrak p A'_{\mathfrak p'})
\geq \text{depth}(M_\mathfrak p)
$$
；这由 $A \to A'$ 的平坦性得出，见《代数》，引理
\ref{algebra-lemma-apply-grothendieck-module}。由于 $A \to A'$ 的纤维是
Cohen--Macaulay 的（《对偶复形》，引理
\ref{dualizing-lemma-dualizing-gorenstein-formal-fibres} 以及
《代数进阶》，第 \ref{more-algebra-section-properties-formal-fibres} 节），可知
$\text{depth}(A'_{\mathfrak p'}/\mathfrak p A'_{\mathfrak p'}) =
\dim(A'_{\mathfrak p'}/\mathfrak p A'_{\mathfrak p'})$.
因此得到
\begin{align*}
\text{depth}(M'_{\mathfrak p'}) +
\dim((A'/\mathfrak p')_{\mathfrak q'})
& =
\text{depth}(M_\mathfrak p) +
\dim(A'_{\mathfrak p'}/\mathfrak p A'_{\mathfrak p'}) +
\dim((A'/\mathfrak p')_{\mathfrak q'}) \\
& =
\text{depth}(M_\mathfrak p) +
\dim((A'/\mathfrak p A')_{\mathfrak q'}) \\
& =
\text{depth}(M_\mathfrak p) +
\dim((A/\mathfrak p)_\mathfrak q)
\end{align*}
第二个等号成立是因为 $A'$ 链状；第三个等号由《代数进阶》，引理
\ref{more-algebra-lemma-completion-dimension} 得出，因为
$(A/\mathfrak p)_\mathfrak q$ 与 $(A'/\mathfrak p A')_{\mathfrak q'}$
具有相同的 $I$-进完备化。
\end{proof}





\begin{lemma}
\label{local-cohomology-lemma-cutoff-completion}
设 $A$ 是普遍链状诺特局部环，$I \subset A$ 是理想，
$M$ 是有限 $A$-模。则
$$
s_{A, I}(M) \geq s_{A^\wedge, I^\wedge}(M^\wedge)
$$
若 $A$ 的形式纤维满足 $(S_n)$，则
$\min(n + 1, s_{A, I}(M)) \leq s_{A^\wedge, I^\wedge}(M^\wedge)$.
\end{lemma}

\begin{proof}
记 $X = \Spec(A)$、$X^\wedge = \Spec(A^\wedge)$、$Z = V(I) \subset X$ 以及
$Z^\wedge = V(I^\wedge)$。设
$\mathfrak p' \subset \mathfrak q' \subset A^\wedge$ 是满足
$\mathfrak p' \not \in Z^\wedge$ 与 $\mathfrak q' \in Z^\wedge$ 的素理想。
设 $\mathfrak p \subset \mathfrak q$ 是 $A$ 中相应的素理想。于是
$\mathfrak p \not \in Z$ 且 $\mathfrak q \in Z$。图示为
$$
\xymatrix{
\mathfrak p' \ar[r] & \mathfrak q' \ar[r] & A^\wedge \\
\mathfrak p \ar[r] \ar@{-}[u] &
\mathfrak q \ar[r] \ar@{-}[u] & A \ar[u]
}
$$
记
\begin{align*}
a & = \dim(A/\mathfrak p) = \dim(A^\wedge/\mathfrak pA^\wedge),\\
b & = \dim(A/\mathfrak q) = \dim(A^\wedge/\mathfrak qA^\wedge),\\
a' & = \dim(A^\wedge/\mathfrak p'),\\
b' & = \dim(A^\wedge/\mathfrak q')
\end{align*}
这些等式由《代数进阶》，引理
\ref{more-algebra-lemma-completion-dimension} 得出。再记
\begin{align*}
p & = \dim(A^\wedge_{\mathfrak p'}/\mathfrak p A^\wedge_{\mathfrak p'}) =
\dim((A^\wedge/\mathfrak p A^\wedge)_{\mathfrak p'}) \\
q & = \dim(A^\wedge_{\mathfrak q'}/\mathfrak p A^\wedge_{\mathfrak q'}) =
\dim((A^\wedge/\mathfrak q A^\wedge)_{\mathfrak q'})
\end{align*}
由于 $A$ 普遍链状，可知
$A^\wedge/\mathfrak pA^\wedge = (A/\mathfrak p)^\wedge$
是维数为 $a$ 的等维环（《代数进阶》，命题
\ref{more-algebra-proposition-ratliff}）。因此 $a = a' + p$；类似地，
$b = b' + q$。将《代数》，引理
\ref{algebra-lemma-apply-grothendieck-module} 用于平坦局部环同态
$A_\mathfrak p \to A^\wedge_{\mathfrak p'}$
，可得
$$
\text{depth}(M^\wedge_{\mathfrak p'})
=
\text{depth}(M_\mathfrak p) +
\text{depth}(A^\wedge_{\mathfrak p'} / \mathfrak p A^\wedge_{\mathfrak p'})
$$
定义 $s_{A, I}(M)$ 时所取最小值的量是
$$
s(\mathfrak p, \mathfrak q) =
\text{depth}(M_\mathfrak p) + \dim((A/\mathfrak p)_\mathfrak q) =
\text{depth}(M_\mathfrak p) + a - b
$$
（最后一个等号成立是因为 $A$ 链状）。定义
$s_{A^\wedge, I^\wedge}(M^\wedge)$ 时所取最小值的量是
$$
s(\mathfrak p', \mathfrak q') =
\text{depth}(M^\wedge_{\mathfrak p'}) +
\dim((A^\wedge/\mathfrak p')_{\mathfrak q'}) =
\text{depth}(M^\wedge_{\mathfrak p'}) + a' - b'
$$
（最后一个等号成立是因为 $A^\wedge$ 链状）。现在记号已经足够，可以开始证明。

\medskip\noindent
设 $\mathfrak p \subset \mathfrak q \subset A$ 是满足
$\mathfrak p \not \in Z$ 与 $\mathfrak q \in Z$ 的素理想，并且
$s_{A, I}(M) = s(\mathfrak p, \mathfrak q)$.
可以选取 $\mathfrak q'$，使其在 $\mathfrak q A^\wedge$ 上方极小；再选取
$\mathfrak p' \subset \mathfrak q'$，使其在 $\mathfrak p A^\wedge$ 上方极小
（这里使用 $A \to A^\wedge$ 的下降定理）。于是得到如上的四个素理想，并且
$p = 0$、$q = 0$。此外，
$\text{depth}(A^\wedge_{\mathfrak p'} / \mathfrak p A^\wedge_{\mathfrak p'})=0$
也成立，因为 $p = 0$。这意味着
$s(\mathfrak p', \mathfrak q') = s(\mathfrak p, \mathfrak q)$.
从而得到第一个不等式。

\medskip\noindent
假设 $A$ 的形式纤维满足 $(S_n)$。则
$\text{depth}(A^\wedge_{\mathfrak p'} / \mathfrak p A^\wedge_{\mathfrak p'})
\geq \min(n, p)$.
因此
$$
s(\mathfrak p', \mathfrak q') \geq
s(\mathfrak p, \mathfrak q) + q + \min(n, p) - p \geq
s_{A, I}(M) + q + \min(n, p) - p
$$
因此，唯一可能出现问题的情形是 $p > n$。若果真如此，则
\begin{align*}
s(\mathfrak p', \mathfrak q')
& =
\text{depth}(M^\wedge_{\mathfrak p'}) +
\dim((A^\wedge/\mathfrak p')_{\mathfrak q'}) \\
& =
\text{depth}(M_\mathfrak p) +
\text{depth}(A^\wedge_{\mathfrak p'} / \mathfrak p A^\wedge_{\mathfrak p'}) +
\dim((A^\wedge/\mathfrak p')_{\mathfrak q'}) \\
& \geq
0 + n + 1
\end{align*}
，因为 $(A^\wedge/\mathfrak p')_{\mathfrak q'}$ 至少含有两个素理想。
这就证明了第二个不等式。
\end{proof}

\noindent
下述引理的证明方法适用于更一般的情形，不过所得更强结果将包含在后面的定理
\ref{local-cohomology-theorem-finiteness} 中。

\begin{lemma}
\label{local-cohomology-lemma-local-annihilator}
\begin{reference}
这是 \cite[Satz 1]{Faltings-annulators} 的一个特殊情形。
\end{reference}
设 $A$ 是 Gorenstein 诺特局部环，$I \subset A$ 是理想，并置
$Z = V(I) \subset \Spec(A)$。设 $M$ 是有限 $A$-模。令
$s = s_{A, I}(M)$ 如式 (\ref{local-cohomology-equation-cutoff}) 所示。则
$H^i_Z(M)$ 在 $i < s$ 时是有限的，但 $H^s_Z(M)$ 不是有限的。
\end{lemma}

\begin{proof}
由于 Gorenstein 局部环具有对偶复形，这是命题
\ref{local-cohomology-proposition-finiteness} 的一个特殊情形。若能给出该特殊情形的简短证明将很有帮助，
因为后面一般有限性定理的证明会用到它。
\end{proof}

\noindent
注意，下述定理的假设对优良诺特环成立（由定义），对具有对偶复形的诺特环成立
（《对偶复形》，引理 \ref{dualizing-lemma-universally-catenary} 以及
《对偶复形》，引理 \ref{dualizing-lemma-dualizing-gorenstein-formal-fibres}），
也对正则诺特环的商成立。

\begin{theorem}
\label{local-cohomology-theorem-finiteness}
\begin{reference}
这是 \cite[Satz 2]{Faltings-finiteness} 的一个特殊情形。
\end{reference}
设 $A$ 是诺特环，$I \subset A$ 是理想。置
$Z = V(I) \subset \Spec(A)$。设 $M$ 是有限 $A$-模。
令 $s = s_{A, I}(M)$ 如式 (\ref{local-cohomology-equation-cutoff}) 所示。假设
\begin{enumerate}
\item $A$ 普遍链状；
\item $A$ 的各局部环的形式纤维都是 Cohen--Macaulay 的。
\end{enumerate}
则 $H^i_Z(M)$ 在 $0 \leq i < s$ 时是有限的，而
$H^s_Z(M)$ 不是有限的。
\end{theorem}

\begin{proof}
由引理 \ref{local-cohomology-lemma-check-finiteness-local-cohomology-locally}，可以假设
$A$ 是局部环。

\medskip\noindent
若 $A$ 是诺特完备局部环，则由 Cohen 结构定理，可将 $A$ 写成正则完备局部环
$B$ 的商（《代数》，定理 \ref{algebra-theorem-cohen-structure-theorem}）。
使用引理 \ref{local-cohomology-lemma-cutoff} 和《对偶复形》，引理
\ref{dualizing-lemma-local-cohomology-and-restriction}，可归约到正则局部环的情形。
这一情形由引理 \ref{local-cohomology-lemma-local-annihilator} 得出，因为正则局部环是 Gorenstein 的
（《对偶复形》，引理 \ref{dualizing-lemma-regular-gorenstein}）。

\medskip\noindent
设 $A$ 是诺特局部环，$\mathfrak m$ 是其极大理想。可以假设
$I \subset \mathfrak m$，否则结论显然。设 $A^\wedge$ 是 $A$ 的完备化，
$Z^\wedge = V(IA^\wedge)$，并设 $M^\wedge = M \otimes_A A^\wedge$
是 $M$ 的完备化（《代数》，引理 \ref{algebra-lemma-completion-tensor}）。则
$H^i_Z(M) \otimes_A A^\wedge = H^i_{Z^\wedge}(M^\wedge)$；这由《对偶复形》，引理
\ref{dualizing-lemma-torsion-change-rings} 和 $A \to A^\wedge$ 的平坦性得出
（《代数》，引理 \ref{algebra-lemma-completion-flat}）。因此只需证明
$H^i_{Z^\wedge}(M^\wedge)$ 在 $i < s$ 时有限、在 $i = s$ 时不有限；见
《代数》，引理 \ref{algebra-lemma-descend-properties-modules}。既然已知结论对
$A^\wedge$ 成立，只需证明
$s_{A, I}(M) = s_{A^\wedge, I^\wedge}(M^\wedge)$。这由引理
\ref{local-cohomology-lemma-cutoff-completion} 得出。
\end{proof}

\begin{remark}
\label{local-cohomology-remark-astute-reader}
细心的读者会注意到，可以对 $A$ 的各局部环的形式纤维施加稍弱的条件。
具体而言，在定理 \ref{local-cohomology-theorem-finiteness} 的情形下，假设 $A$ 普遍链状，
但不对形式纤维作任何假设。设给定 $n$，并希望证明 $H^i_Z(M)$
在 $i \leq n$ 时有限。完全相同的证明表明，只需
$s_{A, I}(M) > n$，且 $A$ 的各局部环的形式纤维满足 $(S_n)$。
另一方面，若希望证明 $H^s_Z(M)$ 不有限，其中 $s = s_{A, I}(M)$，
则上述论证表明，只要形式纤维满足 $(S_{s - 1})$ 即可。
\end{remark}







\section{推前的有限性，II}
\label{local-cohomology-section-finiteness-pushforward-II}

\noindent
本节延续第 \ref{local-cohomology-section-finiteness-pushforward} 节的内容，并利用
第 \ref{local-cohomology-section-finiteness-II} 节中工作的成果。

\begin{lemma}
\label{local-cohomology-lemma-finiteness-Rjstar}
设 $X$ 是局部诺特概形，$j : U \to X$ 是开子概形的嵌入，其补集为 $Z$。
设 $\mathcal{F}$ 是凝聚 $\mathcal{O}_U$-模，$n \geq 0$ 是整数。假设
\begin{enumerate}
\item $X$ 普遍链状；
\item 对每个 $z \in Z$，$\mathcal{O}_{X, z}$ 的形式纤维满足 $(S_n)$。
\end{enumerate}
在这种情形下，下列条件等价：
\begin{enumerate}
\item[(a)] 对 $x \in \text{Supp}(\mathcal{F})$ 和
$z \in Z \cap \overline{\{x\}}$，有
$\text{depth}_{\mathcal{O}_{X, x}}(\mathcal{F}_x) +
\dim(\mathcal{O}_{\overline{\{x\}}, z}) > n$,
\item[(b)] $R^pj_*\mathcal{F}$ 在 $0 \leq p < n$ 时凝聚。
\end{enumerate}
\end{lemma}

\begin{proof}
该断言在 $X$ 上是局部的，故可假设 $X$ 仿射。记
$X = \Spec(A)$ 与 $Z = V(I)$。设 $M$ 是有限 $A$-模，使得与它对应的凝聚
$\mathcal{O}_X$-模限制为 $\mathcal{F}$（在 $U$ 上）；见引理
\ref{local-cohomology-lemma-finiteness-pushforwards-and-H1-local}。该引理还表明，
$R^pj_*\mathcal{F}$ 凝聚，当且仅当 $H^{p + 1}_Z(M)$ 是有限 $A$-模。
注意，下列表达式的最小值
$\text{depth}_{\mathcal{O}_{X, x}}(\mathcal{F}_x) +
\dim(\mathcal{O}_{\overline{\{x\}}, z})$
就是式 (\ref{local-cohomology-equation-cutoff}) 中的数 $s_{A, I}(M)$。于是引理由定理
\ref{local-cohomology-theorem-finiteness} 得出，正如注记 \ref{local-cohomology-remark-astute-reader} 所阐明的那样。
\end{proof}

\begin{lemma}
\label{local-cohomology-lemma-finiteness-for-finite-locally-free}
设 $X$ 是局部诺特概形，$j : U \to X$ 是开子概形的嵌入，其补集为 $Z$。
设 $n \geq 0$ 是整数。若 $R^pj_*\mathcal{O}_U$ 在
$0 \leq p < n$ 时凝聚，则 $R^pj_*\mathcal{F}$ 在
$0 \leq p < n$ 时也凝聚；这里取任意有限局部自由 $\mathcal{O}_U$-模
$\mathcal{F}$。
\end{lemma}

\begin{proof}
该问题在 $X$ 上是局部的，故可假设 $X$ 仿射。记
$X = \Spec(A)$ 与 $Z = V(I)$。由引理
\ref{local-cohomology-lemma-finiteness-pushforwards-and-H1-local}，本引理归结为引理
\ref{local-cohomology-lemma-local-finiteness-for-finite-locally-free}。
\end{proof}

\begin{lemma}
\label{local-cohomology-lemma-annihilate-Hp}
\begin{reference}
\cite[Lemma 1.9]{Bhatt-local}
\end{reference}
设 $A$ 是环，$J \subset I \subset A$ 是有限生成理想，$p \geq 0$ 是整数。
置 $U = \Spec(A) \setminus V(I)$。若 $H^p(U, \mathcal{O}_U)$
被 $J^n$ 零化（对某个 $n$），则 $H^p(U, \mathcal{F})$ 被
$J^m$ 零化（对某个 $m = m(\mathcal{F})$）；这对每个有限局部自由
$\mathcal{O}_U$-模 $\mathcal{F}$ 都成立。
\end{lemma}

\begin{proof}
考虑零化理想 $\mathfrak a$，它对应于 $H^p(U, \mathcal{F})$。取 $u \in U$。
存在开邻域 $u \in U' \subset U$ 以及同构
$\varphi : \mathcal{O}_{U'}^{\oplus r} \to \mathcal{F}|_{U'}$。
选取 $f \in A$，使得 $u \in D(f) \subset U'$。存在映射
$$
a : \mathcal{O}_U^{\oplus r} \longrightarrow \mathcal{F}
\quad\text{且}\quad
b : \mathcal{F} \longrightarrow \mathcal{O}_U^{\oplus r}
$$
，它们在 $D(f)$ 上的限制分别等于 $f^N \varphi$ 和
$f^N \varphi^{-1}$，其中 $N$ 适当。还可以假设
$a \circ b$ 与 $b \circ a$ 都等于乘以 $f^{2N}$ 的映射。
这由《性质》，引理 \ref{properties-lemma-section-maps-backwards} 得出，因为
$U$ 拟紧（$I$ 有限生成）且可分离，并且 $\mathcal{F}$ 与
$\mathcal{O}_U^{\oplus r}$ 都是有限呈示的。因此 $H^p(U, \mathcal{F})$
被 $f^{2N}J^n$ 零化，即 $f^{2N}J^n \subset \mathfrak a$。

\medskip\noindent
由于 $U$ 拟紧，可以找到有限多个 $f_1, \ldots, f_t$ 以及
$N_1, \ldots, N_t$，使得 $U = \bigcup D(f_i)$ 且
$f_i^{2N_i}J^n \subset \mathfrak a$。于是 $V(I) = V(f_1, \ldots, f_t)$；
由于 $I$ 有限生成，可知 $I^M \subset (f_1, \ldots, f_t)$
对某个 $M$ 成立。综上，$J^m \subset \mathfrak a$ 对
$m \gg 0$ 成立；例如取 $m = M (2N_1 + \ldots + 2N_t) n$ 即可。
\end{proof}











\section{局部上同调的零化理想，II}
\label{local-cohomology-section-annihilators-II}

\noindent
把第 \ref{local-cohomology-section-annihilators} 节关于局部上同调零化理想的讨论，
推广到上同调模有限的下有界复形。

\begin{definition}
\label{local-cohomology-definition-depth-complex}
设 $I$ 是诺特环 $A$ 的理想，并设 $K \in D^+_{\textit{Coh}}(A)$。
定义所谓{\it $I$-深度}（针对 $K$），记作 $\text{depth}_I(K)$，为最大
$m \in \mathbf{Z} \cup \{\infty\}$，使得 $H^i_I(K) = 0$
对所有 $i < m$ 成立。若 $A$ 是极大理想为 $\mathfrak m$
的局部环，则简称 $\text{depth}_\mathfrak m(K)$ 为 $K$ 的{\it 深度}。
\end{definition}

\noindent
该定义与《代数》，定义 \ref{algebra-definition-depth} 不冲突；这由《对偶复形》，引理
\ref{dualizing-lemma-depth} 得出。

\begin{proposition}
\label{local-cohomology-proposition-annihilator-complex}
设 $A$ 是具有对偶复形的诺特环，
$T \subset T' \subset \Spec(A)$ 是在特化下稳定的子集。
设 $s \in \mathbf{Z}$，并设 $K$ 是 $D_{\textit{Coh}}^+(A)$ 的对象。
下列条件等价：
\begin{enumerate}
\item 存在理想 $J \subset A$，满足 $V(J) \subset T'$，且
$J$ 将 $H^i_T(K)$ 对所有 $i \leq s$ 零化；
\item 对所有满足 $\mathfrak p \not \in T'$、
$\mathfrak q \in T$ 以及 $\mathfrak p \subset \mathfrak q$ 的情形，都有
$$
\text{depth}_{A_\mathfrak p}(K_\mathfrak p) +
\dim((A/\mathfrak p)_\mathfrak q) > s
$$
\end{enumerate}
\end{proposition}

\begin{proof}
这是命题 \ref{local-cohomology-proposition-annihilator} 的自然推广；读者应先阅读该命题的证明。
设 $\omega_A^\bullet$ 是对偶复形，$\delta$ 是其维数函数；见《对偶复形》，
第 \ref{dualizing-section-dimension-function} 节。下列有限 $A$-模将起重要作用：
$$
E^i = \Ext_A^i(K, \omega_A^\bullet)
$$
对 $\mathfrak p \subset A$，记 $H^i_\mathfrak p$ 为 $D(A_\mathfrak p)$
中对象关于 $\mathfrak pA_\mathfrak p$ 的局部上同调。于是下式的
$\mathfrak pA_\mathfrak p$-进完备化
$$
(E^i)_\mathfrak p =
\Ext^{\delta(\mathfrak p) + i}_{A_\mathfrak p}(K_\mathfrak p,
(\omega_A^\bullet)_\mathfrak p[-\delta(\mathfrak p)])
$$
与下式 Matlis 对偶：
$$
H^{-\delta(\mathfrak p) - i}_{\mathfrak p}(K_\mathfrak p)
$$
这由《对偶复形》，引理 \ref{dualizing-lemma-special-case-local-duality} 得出。
特别地，由此得到如下事实：理想 $J \subset A$ 将
$(E^i)_\mathfrak p$ 零化，当且仅当 $J$ 将
$H^{-\delta(\mathfrak p) - i}_{\mathfrak p}(K_\mathfrak p)$ 零化。

\medskip\noindent
置 $T_n = \{\mathfrak p \in T \mid \delta(\mathfrak p) \leq n\}$。
由于 $\delta$ 是有界函数，可知 $T_a = \emptyset$ 在 $a \ll 0$ 时成立，
而 $T_b = T$ 在 $b \gg 0$ 时成立。

\medskip\noindent
假设 (2) 成立。来证明 (1) 中 $J$ 的存在性。为此使用双重归纳。
对 $i \leq s$，考虑归纳假设 $IH_i$：$H^a_T(K)$ 被某个满足
$J \subset A$ 与 $V(J) \subset T'$ 的理想零化，其中 $a \leq i$。
情形 $IH_i$ 在 $i$ 充分小时是显然的，因为 $K$ 下有界。

\medskip\noindent
归纳步骤。假设 $IH_{i - 1}$ 对某个 $i \leq s$ 成立。选取 $J'$，满足
$V(J') \subset T'$，使其将 $H^a_T(K)$ 在 $a \leq i - 1$ 时零化
（归纳假设保证可以如此选取）。将对 $n$ 作降归纳，证明存在理想 $J$，
满足 $V(J) \subset T'$，使 $J H^i_T(K)$ 的伴随素理想都属于 $T_n$。
当 $n \ll 0$ 时，这蕴含 $JH^i_T(K) = 0$
（《代数》，引理 \ref{algebra-lemma-ass-zero}），因而 $IH_i$ 成立。
基础情形 $n \gg 0$ 是显然的，因为此时 $T = T_n$，而
$H^i_T(K)$ 的所有伴随素理想都属于 $T$。

\medskip\noindent
于是设已给定 $J$，它对 $n$ 具有上述性质。取 $\mathfrak q \in T_n$。
令 $T_\mathfrak q \subset \Spec(A_\mathfrak q)$ 为 $T$ 的逆像。我们有
$H^j_T(K)_\mathfrak q = H^j_{T_\mathfrak q}(K_\mathfrak q)$
；见引理 \ref{local-cohomology-lemma-torsion-change-rings}。考虑谱序列
$$
H_\mathfrak q^p(H^q_{T_\mathfrak q}(K_\mathfrak q))
\Rightarrow
H^{p + q}_\mathfrak q(K_\mathfrak q)
$$
，它来自引理 \ref{local-cohomology-lemma-local-cohomology-ss}。下文将找到理想
$J'' \subset A$，满足 $V(J'') \subset T'$，使得
$H^i_\mathfrak q(K_\mathfrak q)$ 被 $J''$ 零化；这对所有
$\mathfrak q \in T_n \setminus T_{n - 1}$ 都成立。
断言：$J (J')^i J''$ 对 $n - 1$ 可行。事实上，取
$\mathfrak q \in T_n \setminus T_{n - 1}$。上述谱序列给出滤过
$$
E_\infty^{0, i} = E_{i + 2}^{0, i} \subset \ldots \subset E_3^{0, i} \subset
E_2^{0, i} = H^0_\mathfrak q(H^i_{T_\mathfrak q}(K_\mathfrak q))
$$
。模 $E_\infty^{0, i}$ 被 $J''$ 零化。子商
$E_j^{0, i}/E_{j + 1}^{0, i}$ 在 $i + 1 \geq j \geq 2$ 时被 $J'$ 零化，
因为 $d_j^{0, i}$ 的靶是下式的一个子商：
$$
H^j_\mathfrak q(H^{i - j + 1}_{T_\mathfrak q}(K_\mathfrak q)) =
H^j_\mathfrak q(H^{i - j + 1}_T(K)_\mathfrak q)
$$
，而 $H^{i - j + 1}_T(K)_\mathfrak q$ 被 $J'$ 零化；这是由 $J'$ 的选取得出的。
最后，由 $J$ 的选取，有
$J H^i_T(K)_\mathfrak q \subset H^0_\mathfrak q(H^i_T(K)_\mathfrak q)$
，因为 $\Spec(A_\mathfrak q)$ 的非闭点具有更大的 $\delta$ 值。因此
$\mathfrak q$ 不可能是 $J(J')^iJ'' H^i_T(K)$ 的伴随素理想，正如所需。

\medskip\noindent
由开头的讨论可知，$J''$ 应将
$$
(E^{-\delta(\mathfrak q) - i})_\mathfrak q =
(E^{-n - i})_\mathfrak q
$$
对所有 $\mathfrak q \in T_n \setminus T_{n - 1}$ 零化。
但若 $J''$ 对某个 $\mathfrak q$ 可行，则它对 $\mathfrak q$ 的一个开邻域中的所有
$\mathfrak q$ 都可行，因为模 $E^{-n - i}$ 是有限的。
由于 $\Spec(A)$ 的每个子集在诱导拓扑下都是诺特的（《拓扑》，引理
\ref{topology-lemma-Noetherian}），可知只需证明 $J''$ 对某个
$\mathfrak q$ 的存在性。

\medskip\noindent
由于这些 Ext 模是有限的，$J''$ 的存在性等价于
$$
\text{Supp}(E^{-n - i}) \cap \Spec(A_\mathfrak q) \subset T'.
$$
。这又等价于证明 $E^{-n - i}$ 在每个满足
$\mathfrak p \subset \mathfrak q$、$\mathfrak p \not \in T'$ 的点处的局部化为零。
在 $A_\mathfrak p$ 上使用局部对偶性，可知需要证明
$$
H^{i + n - \delta(\mathfrak p)}_\mathfrak p(K_\mathfrak p) =
H^{i - \dim((A/\mathfrak p)_\mathfrak q)}_\mathfrak p(K_\mathfrak p)
$$
为零（这里用到 $\delta$ 是维数函数）。由命题中的假设、$i \leq s$ 以及
定义 \ref{local-cohomology-definition-depth-complex} 中的深度定义，它确实为零。

\medskip\noindent
为证明逆向蕴含，假设 (2) 不成立，并反向使用上述论证。首先选取
$\mathfrak q \in T$、$\mathfrak p \subset \mathfrak q$，满足
$\mathfrak p \not \in T'$，使得
$$
i = \text{depth}_{A_\mathfrak p}(K_\mathfrak p) +
\dim((A/\mathfrak p)_\mathfrak q) \leq s
$$
取最小值。于是
$H^{i - \dim((A/\mathfrak p)_\mathfrak q)}_\mathfrak p(K_\mathfrak p)$
非零；这由定义 \ref{local-cohomology-definition-depth-complex} 中的深度定义得出。
置 $n = \delta(\mathfrak q)$。此时不存在满足
$J \subset A$ 与 $V(J) \subset T'$ 的理想，使得
$J(E^{-n - i})_\mathfrak q = 0$。因此 $H^i_\mathfrak q(K_\mathfrak q)$
不被任何满足 $J \subset A$ 与 $V(J) \subset T'$ 的理想零化。
由 $i$ 的极小性和上面展示的谱序列，模 $H^i_T(K)_\mathfrak q$
不被任何满足 $J \subset A$ 与 $V(J) \subset T'$ 的理想零化。
因此 $H^i_T(K)$ 不被任何满足 $J \subset A$ 与 $V(J) \subset T'$
的理想零化。命题得证。
\end{proof}







\section{局部上同调的有限性，III}
\label{local-cohomology-section-finiteness-III}

\noindent
把第 \ref{local-cohomology-section-finiteness} 节和第 \ref{local-cohomology-section-finiteness-II} 节中
关于局部上同调有限性的讨论，推广到上同调模有限的下有界复形。

\begin{lemma}
\label{local-cohomology-lemma-check-finiteness-local-cohomology-by-annihilator-complex}
设 $A$ 是诺特环，$T \subset \Spec(A)$ 是在特化下稳定的子集，
$K$ 是 $D_{\textit{Coh}}^+(A)$ 的对象。设 $n \in \mathbf{Z}$。
下列条件等价：
\begin{enumerate}
\item $H^i_T(K)$ 在 $i \leq n$ 时有限；
\item 存在理想 $J \subset A$，满足 $V(J) \subset T$，且
$J$ 将 $H^i_T(K)$ 对所有 $i \leq n$ 零化。
\end{enumerate}
若 $T = V(I) = Z$，其中 $I \subset A$ 是理想，则它们还等价于
\begin{enumerate}
\item[(3)] 存在 $e \geq 0$，使得 $I^e$ 将 $H^i_Z(K)$
对所有 $i \leq n$ 零化。
\end{enumerate}
\end{lemma}

\begin{proof}
本引理是引理 \ref{local-cohomology-lemma-check-finiteness-local-cohomology-by-annihilator}
的自然推广；读者应先阅读后者的证明。假设 (1) 成立。回忆
$H^i_J(K) = H^i_{V(J)}(K)$；见《对偶复形》，引理
\ref{dualizing-lemma-local-cohomology-noetherian}。因此
$H^i_T(K) = \colim H^i_J(K)$，其中余极限遍历满足
$J \subset A$ 与 $V(J) \subset T$ 的理想；见引理
\ref{local-cohomology-lemma-adjoint-ext}。由于 $H^i_T(K)$ 在 $i \leq n$ 时有限生成，
可以找到 (2) 中的 $J \subset A$，使得
$H^i_J(K) \to H^i_T(K)$ 在 $i \leq n$ 时满射。
因此这有限组生成元都是 $J$-幂挠元，从而将 $J$ 换成某个幂后 (2) 成立。

\medskip\noindent
取整数 $a \in \mathbf{Z}$，使得 $H^i(K) = 0$ 对 $i < a$ 成立。
证明 (2) $\Rightarrow$ (1)，方法是对 $a$ 作降归纳。若 $a > n$，则
$H^i_T(K) = 0$ 对 $i \leq n$ 成立，因而 (1) 与 (2) 都成立，无须再证。

\medskip\noindent
假设已有 (2) 中的 $J$。注意 $N = H^a_T(K) = H^0_T(H^a(K))$ 是有限
$A$-模 $H^a(K)$ 的子模，因而有限。若 $n = a$，则已经完成；故以下假设
$a < n$。由 $R\Gamma_T$ 的构造，有 $H^i_T(N) = 0$ 对 $i > 0$ 成立，且
$H^0_T(N) = N$；见注记 \ref{local-cohomology-remark-upshot}。选取特异三角
$$
N[-a] \to K \to K' \to N[-a + 1]
$$
。于是 $H^a_T(K') = 0$，且 $H^i_T(K) = H^i_T(K')$ 对 $i > a$ 成立。
因此可以将 $K$ 换成 $K'$，从而可假设 $H^a_T(K) = 0$。
这意味着 $H^a(K)$ 的有限伴随素理想集与 $T$ 不交。由素避
（《代数》，引理 \ref{algebra-lemma-silly}），可找到 $f \in J$，
它不属于 $H^a(K)$ 的任何伴随素理想。选取特异三角
$$
L \to K \xrightarrow{f} K \to L[1]
$$
。由构造可知 $H^i(L) = 0$ 对 $i \leq a$ 成立。另一方面，有上同调长正合列
$$
0 \to H^{a + 1}_T(L) \to H^{a + 1}_T(K) \xrightarrow{f}
H^{a + 1}_T(K) \to H^{a + 2}_T(L) \to H^{a + 2}_T(K) \xrightarrow{f} \ldots
$$
，它分解为恒等式 $H^{a + 1}_T(L) = H^{a + 1}_T(K)$ 和短正合列
$$
0 \to H^{i - 1}_T(K) \to H^i_T(L) \to H^i_T(K) \to 0
$$
；这里 $i \leq n$，且使用了 $f \in J$。由此可知 $J^2$ 将
$H^i_T(L)$ 对所有 $i \leq n$ 零化。把归纳假设用于 $L$，可知
$H^i_T(L)$ 在 $i \leq n$ 时有限。再次使用短正合列，可知
$H^i_T(K)$ 在 $i \leq n$ 时有限，正如所需。

\medskip\noindent
在 $T = V(I)$ 的情形下，省略 (2) 与 (3) 等价性的证明。
\end{proof}

\begin{proposition}
\label{local-cohomology-proposition-finiteness-complex}
设 $A$ 是具有对偶复形的诺特环，
$T \subset \Spec(A)$ 是在特化下稳定的子集。
设 $s \in \mathbf{Z}$，$K \in D_{\textit{Coh}}^+(A)$。下列条件等价：
\begin{enumerate}
\item $H^i_T(K)$ 是有限 $A$-模，其中 $i \leq s$；
\item 对所有满足 $\mathfrak p \not \in T$、$\mathfrak q \in T$ 以及
$\mathfrak p \subset \mathfrak q$ 的情形，都有
$$
\text{depth}_{A_\mathfrak p}(K_\mathfrak p) +
\dim((A/\mathfrak p)_\mathfrak q) > s
$$
\end{enumerate}
\end{proposition}

\begin{proof}
这是命题 \ref{local-cohomology-proposition-annihilator-complex} 和引理
\ref{local-cohomology-lemma-check-finiteness-local-cohomology-by-annihilator-complex} 的形式推论。
\end{proof}






\section{改进凝聚模}
\label{local-cohomology-section-improve}

\noindent
类似的构造可见于 \cite{EGA}，以及较近期的
\cite{Kollar-local-global-hulls} 和 \cite{Kollar-variants}。

\begin{lemma}
\label{local-cohomology-lemma-get-depth-1-along-Z}
设 $X$ 是诺特概形，$T \subset X$ 是在特化下稳定的子集，
$\mathcal{F}$ 是凝聚 $\mathcal{O}_X$-模。则存在唯一同态
$\mathcal{F} \to \mathcal{F}'$，它是凝聚 $\mathcal{O}_X$-模之间的同态，并且
\begin{enumerate}
\item $\mathcal{F} \to \mathcal{F}'$ 是满射；
\item $\mathcal{F}_x \to \mathcal{F}'_x$ 对 $x \not \in T$ 是同构；
\item $\text{depth}_{\mathcal{O}_{X, x}}(\mathcal{F}'_x) \geq 1$ 对 $x \in T$ 成立。
\end{enumerate}
若 $f : Y \to X$ 是平坦态射且 $Y$ 诺特，则
$f^*\mathcal{F} \to f^*\mathcal{F}'$ 是相对于
$f^{-1}(T) \subset Y$ 和 $f^*\mathcal{F}$ 的相应商。
\end{lemma}

\begin{proof}
条件 (3) 恰好意味着 $\text{Ass}(\mathcal{F}') \cap T = \emptyset$。
因此，$\mathcal{F} \to \mathcal{F}'$ 是 $\mathcal{F}$ 对支集包含于
$T$ 的截面所成子层取商而得的商。这证明了唯一性。关于拉回的陈述由
《除子》，引理 \ref{divisors-lemma-bourbaki} 及唯一性得出。

\medskip\noindent
$\mathcal{F} \to \mathcal{F}'$ 的存在性。
由唯一性，只需在 $X$ 上局部证明存在性与唯一性；略去一个小细节。
因此可以假设 $X = \Spec(A)$ 是仿射的，且 $\mathcal{F}$ 是由有限
$A$-模 $M$ 所伴随的凝聚模。令 $M' = M / H^0_T(M)$，其中
$H^0_T(M)$ 如第 \ref{local-cohomology-section-supports} 节所定义。于是
$M_\mathfrak p = M'_\mathfrak p$ 对 $\mathfrak p \not \in T$ 成立，
这证明了 (1)。另一方面，$H^0_T(M) = \colim H^0_Z(M)$，其中 $Z$
遍历 $X$ 的、包含于 $T$ 的闭子集。因此，由《对偶复形》，引理
\ref{dualizing-lemma-divide-by-torsion}，有 $H^0_T(M') = 0$；也就是说，
$M'$ 的任何伴随素理想均不在 $T$ 中。因此，
$\text{depth}(M'_\mathfrak p) \geq 1$ 对 $\mathfrak p \in T$ 成立。
\end{proof}

\begin{lemma}
\label{local-cohomology-lemma-get-depth-2-along-Z}
设 $j : U \to X$ 是诺特概形间的开浸入，$\mathcal{F}$ 是凝聚
$\mathcal{O}_X$-模。假设 $\mathcal{F}' = j_*(\mathcal{F}|_U)$ 是凝聚的。
则 $\mathcal{F} \to \mathcal{F}'$ 是唯一满足下列条件的凝聚
$\mathcal{O}_X$-模同态：
\begin{enumerate}
\item $\mathcal{F}|_U \to \mathcal{F}'|_U$ 是同构；
\item $\text{depth}_{\mathcal{O}_{X, x}}(\mathcal{F}'_x) \geq 2$
对 $x \in X$ 且 $x \not \in U$ 成立。
\end{enumerate}
若 $f : Y \to X$ 是平坦态射且 $Y$ 诺特，则
$f^*\mathcal{F} \to f^*\mathcal{F}'$ 是相对于
$f^{-1}(U) \subset Y$ 的相应同态。
\end{lemma}

\begin{proof}
由《除子》，引理 \ref{divisors-lemma-depth-pushforward} 的第 (3) 部分，
有 $\text{depth}_{\mathcal{O}_{X, x}}(\mathcal{F}'_x) \geq 2$。
$\mathcal{F} \to \mathcal{F}'$ 的唯一性来自《除子》，引理
\ref{divisors-lemma-depth-2-hartog}。与平坦拉回的相容性来自平坦基变换；
见《概形的上同调》，引理
\ref{coherent-lemma-flat-base-change-cohomology}。
\end{proof}

\begin{lemma}
\label{local-cohomology-lemma-make-S2-along-Z}
设 $X$ 是诺特概形，$Z \subset X$ 是闭子概形，$\mathcal{F}$ 是凝聚
$\mathcal{O}_X$-模。假设 $X$ 普遍链状，且其局部环的形式纤维满足 $(S_1)$。
则存在唯一同态 $\mathcal{F} \to \mathcal{F}''$，它是凝聚
$\mathcal{O}_X$-模之间的同态，并且
\begin{enumerate}
\item $\mathcal{F}_x \to \mathcal{F}''_x$ 对 $x \in X \setminus Z$ 是同构；
\item $\mathcal{F}_x \to \mathcal{F}''_x$ 是满射，且
$\text{depth}_{\mathcal{O}_{X, x}}(\mathcal{F}''_x) = 1$，其中 $x \in Z$
满足存在直接特化 $x' \leadsto x$，并且 $x' \not \in Z$、
$x' \in \text{Ass}(\mathcal{F})$；
\item $\text{depth}_{\mathcal{O}_{X, x}}(\mathcal{F}''_x) \geq 2$
对其余 $x \in Z$ 成立。
\end{enumerate}
若 $f : Y \to X$ 是 Cohen--Macaulay 态射且 $Y$ 诺特，则
$f^*\mathcal{F} \to f^*\mathcal{F}''$ 对
$f^{-1}(Z) \subset Y$ 满足同样的性质。
\end{lemma}

\begin{proof}
令 $\mathcal{F} \to \mathcal{F}'$ 为引理
\ref{local-cohomology-lemma-get-depth-1-along-Z} 针对 $Z$（$X$ 的子集）所构造的同态。
回顾 $\mathcal{F}'$ 是 $\mathcal{F}$ 对支集在 $Z$ 上的截面所成子层取商而得的商。

\medskip\noindent
先证明唯一性。设 $\mathcal{F} \to \mathcal{F}''$ 如本引理所述。
得到分解 $\mathcal{F} \to \mathcal{F}' \to \mathcal{F}''$，因为由条件
(2) 与 (3)，有 $\text{Ass}(\mathcal{F}'') \cap Z = \emptyset$。
令 $U \subset X$ 为使得
$\mathcal{F}'|_U \to \mathcal{F}''|_U$ 是同构的极大开子概形。
可以看出，$U$ 包含 (2) 中所述的所有点。于是由《除子》，引理
\ref{divisors-lemma-depth-2-hartog}，可知
$\mathcal{F}'' = j_*(\mathcal{F}'|_U)$。这便给出唯一性
（一个小细节是：若有两个这样的 $\mathcal{F}''$，则取分别所得开集 $U$ 的交）。

\medskip\noindent
证明存在性。回顾
$\text{Ass}(\mathcal{F}') = \{x_1, \ldots, x_n\}$ 是有限集，且
$x_i \not \in Z$。令 $Y_i$ 为 $\{x_i\}$ 的闭包，并令 $Z_{i, j}$ 为
$Z \cap Y_i$ 的不可约分支。注意
$\text{Supp}(\mathcal{F}') \cap Z = \bigcup Z_{i, j}$。
令 $z_{i, j} \in Z_{i, j}$ 为泛点，并令
$$
d_{i, j} = \dim(\mathcal{O}_{\overline{\{x_i\}}, z_{i, j}})
$$
若 $d_{i, j} = 1$，则 $z_{i, j}$ 是 (2) 中所述的点之一，故无须在这些点处
修改 $\mathcal{F}'$。此外，仍假设 $d_{i, j} = 1$，利用引理
\ref{local-cohomology-lemma-depth-function}，可以找到开邻域
$z_{i, j} \in V_{i, j} \subset X$，使得
$\text{depth}_{\mathcal{O}_{X, z}}(\mathcal{F}'_z) \geq 2$
对 $z \in Z_{i, j} \cap V_{i, j}$ 且 $z \not = z_{i, j}$ 成立。
令
$$
Z' = X \setminus
\left(
X \setminus Z \cup \bigcup\nolimits_{d_{i, j} = 1} V_{i, j})
\right)
$$
记 $j' : X \setminus Z' \to X$。由我们对 $Z'$ 的选择，引理
\ref{local-cohomology-lemma-finiteness-pushforward-general} 的假设均得到满足。
令 $\mathcal{F}'' = j'_*(\mathcal{F}'|_{X \setminus Z'})$，并应用引理
\ref{local-cohomology-lemma-get-depth-2-along-Z}，即可得出结论。

\medskip\noindent
最后一个陈述来自深度沿平坦局部同态的变化公式（见《代数》，引理
\ref{algebra-lemma-apply-grothendieck-module}），以及 $f$ 为
Cohen--Macaulay 态射这一条件中所包含的对 $f$ 的纤维的假设。略去细节。
\end{proof}

\begin{lemma}
\label{local-cohomology-lemma-make-S2-along-T-simple}
设 $X$ 是局部具有对偶复形的诺特概形，$T' \subset X$ 是在特化下稳定的
子集，$\mathcal{F}$ 是凝聚 $\mathcal{O}_X$-模。假设只要
$x \leadsto x'$ 是 $X$ 中点的直接特化，且 $x' \in T'$、
$x \not \in T'$，就有 $\text{depth}(\mathcal{F}_x) \geq 1$。
则存在唯一同态 $\mathcal{F} \to \mathcal{F}''$，它是凝聚
$\mathcal{O}_X$-模之间的同态，并且
\begin{enumerate}
\item $\mathcal{F}_x \to \mathcal{F}''_x$ 对 $x \not \in T'$ 是同构；
\item $\text{depth}_{\mathcal{O}_{X, x}}(\mathcal{F}''_x) \geq 2$
对 $x \in T'$ 成立。
\end{enumerate}
若 $f : Y \to X$ 是 Cohen--Macaulay 态射且 $Y$ 诺特，则
$f^*\mathcal{F} \to f^*\mathcal{F}''$ 对
$f^{-1}(T') \subset Y$ 满足同样的性质。
\end{lemma}

\begin{proof}
证明唯一性。设 $\mathcal{F} \to \mathcal{F}''$ 如本引理所述。
令 $U \subset X$ 为使得 $\mathcal{F}|_U \to \mathcal{F}''|_U$
是同构的极大开子概形。由 (1)，可见 $U$ 包含集合 $X \setminus T'$。
由 (2) 以及《除子》，引理 \ref{divisors-lemma-depth-2-hartog}，可知
$\mathcal{F}'' = j_*(\mathcal{F}|_U)$，其中 $j : U \to X$ 是包含态射。
这便给出唯一性（一个小细节是：若有两个这样的 $\mathcal{F}''$，
则取分别所得开集 $U$ 的交）。

\medskip\noindent
证明存在性。令 $\mathcal{F} \to \mathcal{F}'$ 为 $\mathcal{F}$ 的商，
它由引理 \ref{local-cohomology-lemma-get-depth-1-along-Z} 利用 $T'$ 构造。
回顾 $\mathcal{F}'$ 是 $\mathcal{F}$ 对支集在 $T'$ 上的截面所成子层取商而得的商。
定义
$$
\mathcal{F}'' = \colim j_*(\mathcal{F}'|_V)
$$
其中 $j : V \to X$ 遍历满足 $X \setminus V \subset T'$ 的开子概形。
由于 $T'$ 在特化下稳定，注意此余极限是滤过的。每个同态
$\mathcal{F}' \to j_*(\mathcal{F}'|_V)$ 都是单射，因为
$\text{Ass}(\mathcal{F}')$ 与 $T'$ 不交。因此
$\mathcal{F}' \to \mathcal{F}''$ 是单射。

\medskip\noindent
假设 $X = \Spec(A)$ 是仿射的，且 $\mathcal{F}$ 对应于有限
$A$-模 $M$。则 $\mathcal{F}'$ 对应于
$M' = M / H^0_{T'}(M)$；见引理 \ref{local-cohomology-lemma-get-depth-1-along-Z} 的证明。
应用引理 \ref{local-cohomology-lemma-local-cohomology} 与 \ref{local-cohomology-lemma-adjoint-ext}，
可见 $\mathcal{F}''$ 对应于一个 $A$-模 $M''$，它位于短正合列
$$
0 \to M' \to M'' \to H^1_{T'}(M') \to 0
$$
中。由命题 \ref{local-cohomology-proposition-finiteness} 以及本引理陈述中对直接特化的条件，
可见 $M''$ 是有限 $A$-模。由此可见 $\mathcal{F}''$ 是凝聚的。

\medskip\noindent
由 $\mathcal{F}''$ 的凝聚性以及上述滤过余极限中转移同态的单射性，
可知 $\mathcal{F}'' = j_*(\mathcal{F}'|_V)$ 对任意如上足够小的 $V$
都成立。于是条件 (1) 与 (2) 由引理
\ref{local-cohomology-lemma-get-depth-2-along-Z} 得出。

\medskip\noindent
最后一个陈述来自深度沿平坦局部同态的变化公式（见《代数》，引理
\ref{algebra-lemma-apply-grothendieck-module}），以及 $f$ 为
Cohen--Macaulay 态射这一条件中所包含的对 $f$ 的纤维的假设。略去细节。
\end{proof}

\begin{lemma}
\label{local-cohomology-lemma-make-S2-along-T}
设 $X$ 是局部具有对偶复形的诺特概形。设
$T' \subset T \subset X$ 是在特化下稳定的子集，并且只要
$x \leadsto x'$ 是 $X$ 中点的直接特化且 $x' \in T'$，就有
$x \in T$。设 $\mathcal{F}$ 是凝聚 $\mathcal{O}_X$-模。
则存在唯一同态 $\mathcal{F} \to \mathcal{F}''$，它是凝聚
$\mathcal{O}_X$-模之间的同态，并且
\begin{enumerate}
\item $\mathcal{F}_x \to \mathcal{F}''_x$ 对 $x \not \in T$ 是同构；
\item $\mathcal{F}_x \to \mathcal{F}''_x$ 是满射，并且
$\text{depth}_{\mathcal{O}_{X, x}}(\mathcal{F}''_x) \geq 1$
对 $x \in T$ 且 $x \not \in T'$ 成立；
\item $\text{depth}_{\mathcal{O}_{X, x}}(\mathcal{F}''_x) \geq 2$
对 $x \in T'$ 成立。
\end{enumerate}
若 $f : Y \to X$ 是 Cohen--Macaulay 态射且 $Y$ 诺特，则
$f^*\mathcal{F} \to f^*\mathcal{F}''$ 对
$f^{-1}(T') \subset f^{-1}(T) \subset Y$ 满足同样的性质。
\end{lemma}

\begin{proof}
首先，令 $\mathcal{F} \to \mathcal{F}'$ 为 $\mathcal{F}$ 的商，
它由引理 \ref{local-cohomology-lemma-get-depth-1-along-Z} 利用 $T$ 构造。
其次，令 $\mathcal{F}' \to \mathcal{F}''$ 为引理
\ref{local-cohomology-lemma-make-S2-along-T-simple} 利用 $T'$ 所构造的唯一凝聚模同态。
于是 $\mathcal{F} \to \mathcal{F}''$ 正是所求。
\end{proof}










\section{Hartshorne--Lichtenbaum 消失定理}
\label{local-cohomology-section-Hartshorne-Lichtenbaum-vanishing}

\noindent
这一消失结果是 Lichtenbaum 定理的局部对应；读者可在《概形的对偶性》，
第 \ref{duality-section-lichtenbaum} 节中找到该定理。
这一结果以及其他许多内容可见于 \cite{CD}。

\begin{lemma}
\label{local-cohomology-lemma-cd-top-vanishing}
设 $A$ 是维数为 $d$ 的诺特环，$I \subset I' \subset A$ 是理想。
若 $I'$ 包含于 $A$ 的 Jacobson 根，并且
$\text{cd}(A, I') < d$，则 $\text{cd}(A, I) < d$。
\end{lemma}

\begin{proof}
由引理 \ref{local-cohomology-lemma-cd-dimension}，已知 $\text{cd}(A, I) \leq d$。
我们将利用引理 \ref{local-cohomology-lemma-isomorphism} 证明
$$
H^d_{V(I')}(A) \to H^d_{V(I)}(A)
$$
是满射，这将完成证明。取
$\mathfrak p \in V(I) \setminus V(I')$。由对 $I'$ 的假设，可见
$\mathfrak p$ 不是 $A$ 的极大理想。因此
$\dim(A_\mathfrak p) < d$。于是由引理 \ref{local-cohomology-lemma-cd-dimension}，
$H^d_{\mathfrak pA_\mathfrak p}(A_\mathfrak p) = 0$。
\end{proof}

\begin{lemma}
\label{local-cohomology-lemma-cd-top-vanishing-some-module}
设 $A$ 是维数为 $d$ 的诺特环，$I \subset A$ 是理想。若
$H^d_{V(I)}(M) = 0$ 对某个有限 $A$-模成立，且该模的支集包含所有维数为
$d$ 的不可约分支，则 $\text{cd}(A, I) < d$。
\end{lemma}

\begin{proof}
由引理 \ref{local-cohomology-lemma-cd-dimension}，已知 $\text{cd}(A, I) \leq d$。
因此，对任意有限 $A$-模 $N$，都有 $H^i_{V(I)}(N) = 0$，其中
$i > d$。我们约定：性质 $\mathcal{P}$ 对有限 $A$-模 $N$ 成立，是指
$H^d_{V(I)}(N) = 0$。
我们的一个假设是 $\mathcal{P}(M)$ 成立。注意
$\mathcal{P}(N_1 \oplus N_2)
\Leftrightarrow (\mathcal{P}(N_1) \wedge \mathcal{P}(N_2))$。
还要注意，若 $N \to N'$ 是满射，则
$\mathcal{P}(N) \Rightarrow \mathcal{P}(N')$，因为
$H^{d + 1}_{V(I)}$ 消失（见上文）。令
$\mathfrak p_1, \ldots, \mathfrak p_n$ 为 $A$ 的满足
$\dim(A/\mathfrak p_i) = d$ 的极小素理想。注意，由维数原因，
$\mathcal{P}(N)$ 在 $N$ 的支集与
$\{\mathfrak p_1, \ldots, \mathfrak p_n\}$ 不交时成立；见引理
\ref{local-cohomology-lemma-cd-dimension}。
对每个 $i$，令 $M_i = M/\mathfrak p_i M$。这是一个有限 $A$-模，
被 $\mathfrak p_i$ 零化，其支集等于
$V(\mathfrak p_i)$（这里使用了对 $M$ 的支集的假设）。最后，若
$J \subset A$ 是理想，则 $\mathcal{P}(JM_i)$ 成立，因为 $JM_i$
是若干个 $M$ 的直和的商。因此，由《概形的上同调》，引理
\ref{coherent-lemma-property-higher-rank-cohomological}，性质
$\mathcal{P}$ 对每个有限 $A$-模均成立。
\end{proof}

\begin{lemma}
\label{local-cohomology-lemma-top-coh-divisible}
设 $A$ 是维数为 $d$ 的诺特局部环。设 $f \in A$ 不包含于任何维数为
$d$ 的极小素理想。则同态
$f : H^d_{V(I)}(M) \to H^d_{V(I)}(M)$ 对任意有限 $A$-模 $M$
以及任意理想 $I \subset A$ 都是满射。
\end{lemma}

\begin{proof}
$M/fM$ 的支集维数 $< d$，这是由我们对 $f$ 的假设得出的。因此，由引理
\ref{local-cohomology-lemma-cd-dimension}，$H^d_{V(I)}(M/fM) = 0$。从而
$H^d_{V(I)}(fM) \to H^d_{V(I)}(M)$ 是满射。又因由引理
\ref{local-cohomology-lemma-cd-dimension} 已知 $\text{cd}(A, I) \leq d$，还可看出满射
$M \to fM$，$x \mapsto fx$ 诱导出满射
$H^d_{V(I)}(M) \to H^d_{V(I)}(fM)$。
\end{proof}

\begin{lemma}
\label{local-cohomology-lemma-cd-bound-dualizing}
设 $A$ 是具有规范化对偶复形 $\omega_A^\bullet$ 的诺特局部环。
设 $I \subset A$ 是理想。若 $H^0_{V(I)}(\omega_A^\bullet) = 0$，则
$\text{cd}(A, I) < \dim(A)$。
\end{lemma}

\begin{proof}
令 $d = \dim(A)$。设 $\mathfrak p_1, \ldots, \mathfrak p_n \subset A$
为维数为 $d$ 的极小素理想。回顾有限 $A$-模
$H^{-i}(\omega_A^\bullet)$ 仅当 $i \in \{0, \ldots, d\}$ 时可能非零，
而且 $H^{-i}(\omega_A^\bullet)$ 的支集维数 $\leq i$；见引理
\ref{local-cohomology-lemma-sitting-in-degrees}。令 $\omega_A = H^{-d}(\omega_A^\bullet)$。
由素避（《代数》，引理 \ref{algebra-lemma-silly}），可以找到
$f \in A$，$f \not \in \mathfrak p_i$，它零化
$H^{-i}(\omega_A^\bullet)$，其中 $i < d$。考虑特异三角
$$
\omega_A[d] \to \omega_A^\bullet \to
\tau_{\geq -d + 1}\omega_A^\bullet \to \omega_A[d + 1]
$$
见《导出范畴》，评注
\ref{derived-remark-truncation-distinguished-triangle}。
由《导出范畴》，引理 \ref{derived-lemma-trick-vanishing-composition}，
可见 $f^d$ 在 $\tau_{\geq -d + 1}\omega_A^\bullet$ 上诱导零自同态。
利用三角范畴的公理，可以找到一个同态
$$
\omega_A^\bullet \to \omega_A[d]
$$
它与 $\omega_A[d] \to \omega_A^\bullet$ 的复合是乘以 $f^d$，
这一乘法作用在 $\omega_A[d]$ 上。因此可知 $f^d$ 零化
$H^d_{V(I)}(\omega_A)$。由引理 \ref{local-cohomology-lemma-top-coh-divisible}，可知
$H^d_{V(I)}(\omega_A) = 0$。再结合引理
\ref{local-cohomology-lemma-cd-top-vanishing-some-module} 以及
$(\omega_A)_{\mathfrak p_i}$ 非零这一事实（例如见《对偶复形》，引理
\ref{dualizing-lemma-nonvanishing-generically-local}），便得到结论。
\end{proof}

\begin{lemma}
\label{local-cohomology-lemma-inverse-system-symbolic-powers}
设 $(A, \mathfrak m)$ 是完备诺特局部整环，
$\mathfrak p \subset A$ 是维数为 $1$ 的素理想。对每个 $n \geq 1$，
都存在 $m \geq n$，使得 $\mathfrak p^{(m)} \subset \mathfrak p^n$。
\end{lemma}

\begin{proof}
回顾符号幂 $\mathfrak p^{(m)}$ 定义为
$A \to A_\mathfrak p/\mathfrak p^mA_\mathfrak p$ 的核。
由于局部化是正合的，可知在短正合列
$$
0 \to \mathfrak a_n \to A/\mathfrak p^n \to A/\mathfrak p^{(n)} \to 0
$$
中，$\mathfrak a_n$ 的支集包含于 $\{\mathfrak m\}$。特别地，逆系统
$(\mathfrak a_n)$ 满足 Mittag--Leffler 条件，因为每个 $\mathfrak a_n$
都是阿廷 $A$-模。由此可知，本引理等价于要求
$\lim \mathfrak a_n = 0$。设 $f \in \lim \mathfrak a_n$。则 $f$ 是
$A = \lim A/\mathfrak p^n$ 的元素（这里使用了 $A$ 的完备性），
它在完备化 $A_\mathfrak p^\wedge$（即 $A_\mathfrak p$ 的完备化）中映为零。
由于 $A_\mathfrak p \to A_\mathfrak p^\wedge$ 忠实平坦，可见 $f$
在 $A_\mathfrak p$ 中映为零。由于 $A$ 是整环，可见 $f$ 如所需为零。
\end{proof}

\begin{proposition}
\label{local-cohomology-proposition-Hartshorne-Lichtenbaum-vanishing}
\begin{reference}
\cite[Theorem 3.1]{CD}
\end{reference}
设 $A$ 是诺特局部环，其完备化为 $A^\wedge$。设 $I \subset A$ 是理想，
并且
$$
\dim V(IA^\wedge + \mathfrak p) \geq 1
$$
对每个极小素理想 $\mathfrak p \subset A^\wedge$ 都成立，其中该素理想的维数为
$\dim(A)$。则 $\text{cd}(A, I) < \dim(A)$。
\end{proposition}

\begin{proof}
由于 $A \to A^\wedge$ 忠实平坦，由《对偶复形》，引理
\ref{dualizing-lemma-torsion-change-rings}，有
$H^d_{V(I)}(A) \otimes_A A^\wedge = H^d_{V(IA^\wedge)}(A^\wedge)$
。因此可以假设 $A$ 是完备的。

\medskip\noindent
假设 $A$ 完备。设 $\mathfrak p_1, \ldots, \mathfrak p_n \subset A$
为维数为 $d$ 的极小素理想。考虑完备局部环 $A_i = A/\mathfrak p_i$。
由《对偶复形》，引理
\ref{dualizing-lemma-local-cohomology-and-restriction}，有
$H^d_{V(I)}(A_i) = H^d_{V(IA_i)}(A_i)$。由引理
\ref{local-cohomology-lemma-cd-top-vanishing-some-module}，只需对 $(A_i, IA_i)$ 证明本引理。
因此可以假设 $A$ 是完备局部整环。

\medskip\noindent
假设 $A$ 是完备局部整环。可以选取素理想 $\mathfrak p \supset I$，使得
$\dim(A/\mathfrak p) = 1$。由引理 \ref{local-cohomology-lemma-cd-top-vanishing}，
只需对 $\mathfrak p$ 证明本引理。

\medskip\noindent
由引理 \ref{local-cohomology-lemma-cd-bound-dualizing}，只需证明
$H^0_{V(\mathfrak p)}(\omega_A^\bullet) = 0$。回顾
$$
H^0_{V(\mathfrak p)}(\omega_A^\bullet) =
\colim \text{Ext}^0_A(A/\mathfrak p^n, \omega_A^\bullet)
$$
由引理 \ref{local-cohomology-lemma-inverse-system-symbolic-powers}，可见此余极限等于
$$
\colim \text{Ext}^0_A(A/\mathfrak p^{(n)}, \omega_A^\bullet)
$$
由于 $\text{depth}(A/\mathfrak p^{(n)}) = 1$，由引理
\ref{local-cohomology-lemma-sitting-in-degrees}，可见这些 Ext 群为零，正如所需。
\end{proof}

\begin{lemma}
\label{local-cohomology-lemma-affine-complement}
设 $(A, \mathfrak m)$ 是诺特局部环，$I \subset A$ 是理想。假设 $A$
优良且正规，并且 $\dim V(I) \geq 1$。则
$\text{cd}(A, I) < \dim(A)$。特别地，若 $\dim(A) = 2$，则
$\Spec(A) \setminus V(I)$ 是仿射的。
\end{lemma}

\begin{proof}
由《代数进阶》，引理
\ref{more-algebra-lemma-completion-normal-local-ring}，完备化
$A^\wedge$ 是正规的，因而是整环。因此，命题
\ref{local-cohomology-proposition-Hartshorne-Lichtenbaum-vanishing} 的假设成立，结论随之得出。
关于仿射性的陈述来自引理 \ref{local-cohomology-lemma-cd-is-one}。
\end{proof}
















\section{Frobenius 作用}
\label{local-cohomology-section-frobenius}

\noindent
设 $p$ 是素数，$A$ 是使得 $p = 0$ 在 $A$ 中成立的环。
$A$ 的{\it Frobenius 自同态}是同态
$$
F : A \longrightarrow A,
\quad
a \longmapsto a^p
$$
本节证明若干关于具有 Frobenius 作用的模的引理。

\begin{lemma}
\label{local-cohomology-lemma-annihilator-frobenius-module}
设 $p$ 是素数。设 $(A, \mathfrak m, \kappa)$ 是诺特局部环，
且 $p = 0$ 在 $A$ 中成立。设 $M$ 是有限 $A$-模，并且
$M \otimes_{A, F} A \cong M$。则 $M$ 是有限自由模。
\end{lemma}

\begin{proof}
选取一个展示 $A^{\oplus m} \to A^{\oplus n} \to M$，它诱导同构
$\kappa^{\oplus n} \to M/\mathfrak m M$。令 $T = (a_{ij})$ 为同态
$A^{\oplus m} \to A^{\oplus n}$ 的矩阵。注意 $a_{ij} \in \mathfrak m$。
沿 $F$ 作基变换，并利用基变换的右正合性，得到一个展示
$A^{\oplus m} \to A^{\oplus n} \to M$，其中矩阵为
$T = (a_{ij}^p)$。因此，我们得到一个满足
$a_{ij} \in \mathfrak m^p$ 的展示。重复这一构造，可知对每个
$e \geq 1$，都存在满足 $a_{ij} \in \mathfrak m^e$ 的展示。这意味着
Fitting 理想（《代数进阶》，定义
\ref{more-algebra-definition-fitting-ideal}）$\text{Fit}_k(M)$ 在 $k < n$ 时包含于
$\bigcap_{e \geq 1} \mathfrak m^e$。由 Krull 交定理
（《代数》，引理 \ref{algebra-lemma-intersect-powers-ideal-module-zero}），
后者为零，因此由《代数进阶》，引理
\ref{more-algebra-lemma-fitting-ideal-finite-locally-free}，可知
$M$ 是秩为 $n$ 的自由模。
\end{proof}

\noindent
在本节中，称元素 $f_1, \ldots, f_r$（属于环 $A$）是{\it 独立的}，如果
$\sum a_if_i = 0$ 蕴含 $a_i \in (f_1, \ldots, f_r)$。换言之，令
$I = (f_1, \ldots, f_r)$，则 $I/I^2$ 是自由 $A/I$-模，其基为
$f_1, \ldots, f_r$。

\begin{lemma}
\label{local-cohomology-lemma-1}
\begin{reference}
见 \cite{Lech-inequalities} 以及 \cite[Lemma 1 page 299]{MatCA}。
\end{reference}
设 $A$ 是环。若 $f_1, \ldots, f_{r - 1}, f_rg_r$ 独立，则
$f_1, \ldots, f_r$ 独立。
\end{lemma}

\begin{proof}
设 $\sum a_if_i = 0$。则 $\sum a_ig_rf_i = 0$。因此
$a_r \in (f_1, \ldots, f_{r - 1}, f_rg_r)$。写成
$a_r = \sum_{i < r} b_i f_i + b f_rg_r$。于是
$0 = \sum_{i < r} (a_i + b_if_r)f_i + bf_r^2g_r$。从而
$a_i + b_i f_r \in (f_1, \ldots, f_{r - 1}, f_rg_r)$，这蕴含
$a_i \in (f_1, \ldots, f_r)$，正如所需。
\end{proof}

\begin{lemma}
\label{local-cohomology-lemma-2}
\begin{reference}
见 \cite{Lech-inequalities} 以及 \cite[Lemma 2 page 300]{MatCA}。
\end{reference}
设 $A$ 是环。若 $f_1, \ldots, f_{r - 1}, f_rg_r$ 独立，且 $A$-模
$A/(f_1, \ldots, f_{r - 1}, f_rg_r)$ 具有有限长度，则
\begin{align*}
& \text{length}_A(A/(f_1, \ldots, f_{r - 1}, f_rg_r)) \\
& =
\text{length}_A(A/(f_1, \ldots, f_{r - 1}, f_r)) +
\text{length}_A(A/(f_1, \ldots, f_{r - 1}, g_r))
\end{align*}
\end{lemma}

\begin{proof}
我们断言存在正合列
$$
0 \to
A/(f_1, \ldots, f_{r - 1}, g_r) \xrightarrow{f_r}
A/(f_1, \ldots, f_{r - 1}, f_rg_r) \to
A/(f_1, \ldots, f_{r - 1}, f_r) \to 0
$$
事实上，若 $a f_r \in (f_1, \ldots, f_{r - 1}, f_rg_r)$，则
$\sum_{i < r} a_i f_i + (a + bg_r)f_r = 0$ 对某些 $b, a_i \in A$ 成立。因此
$\sum_{i < r} a_i g_r f_i + (a + bg_r)g_rf_r = 0$，这蕴含
$a + bg_r \in (f_1, \ldots, f_{r - 1}, f_rg_r)$，也就是说，$a$ 在
$A/(f_1, \ldots, f_{r - 1}, g_r)$ 中映为零。这证明了断言。
最后使用长度的可加性（《代数》，引理
\ref{algebra-lemma-length-additive}）。
\end{proof}

\begin{lemma}
\label{local-cohomology-lemma-3}
\begin{reference}
见 \cite{Lech-inequalities} 以及 \cite[Lemma 3 page 300]{MatCA}。
\end{reference}
设 $(A, \mathfrak m)$ 是局部环。若 $\mathfrak m = (x_1, \ldots, x_r)$，
并且 $x_1^{e_1}, \ldots, x_r^{e_r}$ 对某些 $e_i > 0$ 独立，则
$\text{length}_A(A/(x_1^{e_1}, \ldots, x_r^{e_r})) = e_1\ldots e_r$。
\end{lemma}

\begin{proof}
使用引理 \ref{local-cohomology-lemma-1}、引理 \ref{local-cohomology-lemma-2} 以及归纳法。
\end{proof}

\begin{lemma}
\label{local-cohomology-lemma-flat-extension-independent}
设 $\varphi : A \to B$ 是平坦环同态。若
$f_1, \ldots, f_r \in A$ 独立，则
$\varphi(f_1), \ldots, \varphi(f_r) \in B$ 独立。
\end{lemma}

\begin{proof}
令 $I = (f_1, \ldots, f_r)$ 且 $J = \varphi(I)B$。由平坦性，有
$I/I^2 \otimes_A B = J/J^2$。因此，$I/I^2$ 作为 $A/I$-模的自由性
蕴含 $J/J^2$ 作为 $B/J$-模的自由性。
\end{proof}

\begin{lemma}[Kunz]
\label{local-cohomology-lemma-frobenius-flat-regular}
\begin{reference}
\cite{Kunz-flat}
\end{reference}
设 $p$ 是素数。设 $A$ 是满足 $p = 0$ 的诺特环。下列条件等价：
\begin{enumerate}
\item $A$ 是正则的；
\item $F : A \to A$，$a \mapsto a^p$，是平坦的。
\end{enumerate}
\end{lemma}

\begin{proof}
注意 $\Spec(F) : \Spec(A) \to \Spec(A)$ 是恒等映射。
正则性由局部环定义，而平坦性也是局部环的性质；见《代数》，引理
\ref{algebra-lemma-flat-localization}。因此可以并且确实假设 $A$
是以 $\mathfrak m$ 为极大理想的诺特局部环。

\medskip\noindent
假设 $A$ 正则。设 $x_1, \ldots, x_d$ 是 $A$ 的一组参数系。
应用 $F$，得到
$F(x_1), \ldots, F(x_d) = x_1^p, \ldots, x_d^p$，
它也是 $A$ 的一组参数系。因此 $F$ 平坦；见《代数》，引理
\ref{algebra-lemma-CM-over-regular-flat} 与
\ref{algebra-lemma-regular-ring-CM}。

\medskip\noindent
反之，假设 $F$ 平坦。写成 $\mathfrak m = (x_1, \ldots, x_r)$，其中
$r$ 极小。则 $x_1, \ldots, x_r$ 在上述意义下独立。由于 $F$ 平坦，
可见 $x_1^p, \ldots, x_r^p$ 独立；见引理
\ref{local-cohomology-lemma-flat-extension-independent}。因此，由引理 \ref{local-cohomology-lemma-3}，
$\text{length}_A(A/(x_1^p, \ldots, x_r^p)) = p^r$。
令 $\chi(n) = \text{length}_A(A/\mathfrak m^n)$，并回顾这是一个次数为
$\dim(A)$ 的数值多项式；见《代数》，命题
\ref{algebra-proposition-dimension}。选取 $n \gg 0$。注意
$$
\mathfrak m^{pn + pr} \subset F(\mathfrak m^n)A \subset \mathfrak m^{pn}
$$
这一点可通过考察 $x_1, \ldots, x_r$ 中的单项式看出。我们有
$$
A/F(\mathfrak m^n)A = A/\mathfrak m^n \otimes_{A, F} A
$$
由 $F$ 的平坦性，它的长度为
$\chi(n) \text{length}_A(A/F(\mathfrak m)A)$
（《代数》，引理 \ref{algebra-lemma-pullback-module}），
而由上文，这等于 $p^r\chi(n)$。所以
$$
\chi(pn + pr) \geq p^r\chi(n) \geq \chi(pn)
$$
考察首项可知这蕴含 $r = \dim(A)$，也就是说，$A$ 是正则的。
\end{proof}




\section{某些模的结构}
\label{local-cohomology-section-structure}

\noindent
这里给出若干关于正则局部环上某些类型模的结构的结果。这类结果以及更多内容
可见于 \cite{Huneke-Sharp}、\cite{Lyubeznik}、\cite{Lyubeznik2}。

\begin{lemma}
\label{local-cohomology-lemma-structure-torsion-D-module-regular}
\begin{reference}
\cite[Theorem 2.4]{Lyubeznik} 的特殊情形
\end{reference}
设 $k$ 是特征为 $0$ 的域，$d \geq 1$。设
$A = k[[x_1, \ldots, x_d]]$，其极大理想为 $\mathfrak m$。
设 $M$ 是 $\mathfrak m$-幂挠 $A$-模，并赋有满足 Leibniz 法则的加性算子
$D_1, \ldots, D_d$：
$$
D_i(fz) = \partial_i(f) z + f D_i(z)
$$
其中 $f \in A$ 且 $z \in M$。这里，$\partial_i$ 是关于 $x_i$ 的微分。
则 $M$ 同构于若干个单射包 $E$（即 $k$ 的单射包）的直和。
\end{lemma}

\begin{proof}
选取集合 $J$ 以及同构 $M[\mathfrak m] \to \bigoplus_{j \in J} k$。
由于 $\bigoplus_{j \in J} E$ 是单射的
（《对偶复形》，引理 \ref{dualizing-lemma-sum-injective-modules}），
可以将此同构延拓为 $A$-模同态
$\varphi : M \to \bigoplus_{j \in J} E$。
我们断言 $\varphi$ 是同构，也就是说，是双射。

\medskip\noindent
单射性。设 $z \in M$ 非零。由于 $M$ 是 $\mathfrak m$-幂挠的，
可以选取元素 $f \in A$，使得 $fz \in M[\mathfrak m]$ 且
$fz \not = 0$。于是 $\varphi(fz) = f\varphi(z)$ 非零，故
$\varphi(z)$ 非零。

\medskip\noindent
满射性。设 $z \in M$。则 $x_1^n z = 0$ 对某个 $n \geq 0$ 成立。
我们将证明 $z \in x_1M$，方法是对 $n$ 归纳。若 $n = 0$，则 $z = 0$，结论成立。
若 $n > 0$，则应用 $D_1$，得到
$0 = n x_1^{n - 1} z + x_1^nD_1(z)$。因此
$x_1^{n - 1}(nz + x_1D_1(z)) = 0$。由归纳假设，
$nz + x_1D_1(z) \in x_1M$。由于 $n$ 可逆，可知 $z \in x_1M$。
因此 $M$ 是 $x_1$-可除的。若 $\varphi$ 不是满射，则可选取
$e \in \bigoplus_{j \in J} E$，使其不在 $M$ 中。与上面类似地论证，
可以假设 $\mathfrak m e \subset M$；特别地，$x_1 e \in M$。
存在元素 $z_1 \in M$，满足 $x_1 z_1 = x_1 e$。因此
$x_1(z_1 - e) = 0$。把 $e$ 替换为 $e - z_1$，可以假设 $e$ 被 $x_1$ 零化。
所以只需证明
$$
\varphi[x_1] :
M[x_1]
\longrightarrow
\left(\bigoplus\nolimits_{j \in J} E\right)[x_1] =
\bigoplus\nolimits_{j \in J} E[x_1]
$$
是满射。若 $d = 1$，则由 $\varphi$ 的构造，这一点成立。
若 $d > 1$，则注意 $E[x_1]$ 是 $k[[x_2, \ldots, x_d]]$
的剩余域的单射包；见《对偶复形》，引理
\ref{dualizing-lemma-quotient}。注意，$M[x_1]$ 作为
$k[[x_2, \ldots, x_d]]$-模是 $\mathfrak m/(x_1)$-幂挠的，并赋有满足上述
Leibniz 法则的算子 $D_2, \ldots, D_d$。因此，对 $d$ 作归纳，可知
$\varphi[x_1]$ 如所需为满射。
\end{proof}

\begin{lemma}
\label{local-cohomology-lemma-structure-torsion-Frobenius-regular}
\begin{reference}
经过少量工作可由 \cite[Corollary 3.6]{Huneke-Sharp} 推出，
也可直接由 \cite[Theorem 1.4]{Lyubeznik2} 推出。
\end{reference}
设 $p$ 是素数。设 $(A, \mathfrak m, k)$ 是满足 $p = 0$ 的正则局部环。
记 $F : A \to A$，$a \mapsto a^p$，为 Frobenius 自同态。设 $M$ 是
$\mathfrak m$-幂挠模，并且 $M \otimes_{A, F} A \cong M$。则 $M$
同构于若干个单射包 $E$（即 $k$ 的单射包）的直和。
\end{lemma}

\begin{proof}
选取集合 $J$ 以及 $A$-模同态
$\varphi : M \to \bigoplus_{j \in J} E$，它将
$M[\mathfrak m]$ 同构地映到
$(\bigoplus_{j \in J} E)[\mathfrak m] = \bigoplus_{j \in J} k$ 上。
我们断言 $\varphi$ 是同构，也就是说，是双射。

\medskip\noindent
单射性。设 $z \in M$ 非零。由于 $M$ 是 $\mathfrak m$-幂挠的，
可以选取元素 $f \in A$，使得 $fz \in M[\mathfrak m]$ 且
$fz \not = 0$。于是 $\varphi(fz) = f\varphi(z)$ 非零，故
$\varphi(z)$ 非零。

\medskip\noindent
满射性。回顾 $F$ 是平坦的；见引理
\ref{local-cohomology-lemma-frobenius-flat-regular}。设 $x_1, \ldots, x_d$ 是
$\mathfrak m$ 的一个极小生成系。记
$$
M_n = M[x_1^{p^n}, \ldots, x_d^{p^n}]
$$
为 $M$ 的由被 $x_1^{p^n}, \ldots, x_d^{p^n}$ 零化的元素所成的子模。
因此 $M_0 = M[\mathfrak m]$ 是 $k$ 上的向量空间。此外，
$M = \bigcup M_n$；这是由 $M$ 是 $\mathfrak m$-幂挠的这一假设得出的。由于 $F^n$ 平坦，且
$F^n(x_i) = x_i^{p^n}$，有
$$
M_n \cong (M \otimes_{A, F^n} A)[x_1^{p^n}, \ldots, x_d^{p^n}] =
M[x_1, \ldots, x_d] \otimes_{A, F^n} A =
M_0 \otimes_k A/(x_1^{p^n}, \ldots, x_d^{p^n})
$$
因此 $M_n$ 是自由 $A/(x_1^{p^n}, \ldots, x_d^{p^n})$-模。
计算表明，$A/(x_1^{p^n}, \ldots, x_d^{p^n})$ 中每个被
$x_1^{p^n - 1}$ 零化的元素都可被 $x_1$ 整除；例如，可使用由《代数》，
引理 \ref{algebra-lemma-regular-complete-containing-coefficient-field} 得到的同构
$A/(x_1^{p^n}, \ldots, x_d^{p^n}) \cong
k[x_1, \ldots, x_d]/(x_1^{p^n}, \ldots, x_d^{p^n})$。
因此，同一结论对 $M_n$ 的每个元素成立。由于 $M$ 的每个元素都属于
$M_n$，而这对所有 $n \gg 0$ 成立，而且 $M$ 的每个元素都被 $x_1$ 的某个幂零化，
可知 $M$ 是 $x_1$-可除的。

\medskip\noindent
令 $x = x_1$。上面已经看到 $M$ 是 $x$-可除的。若 $\varphi$ 不是满射，
则可选取 $e \in \bigoplus_{j \in J} E$，使其不在 $M$ 中。
与上面类似地论证，可以假设 $\mathfrak m e \subset M$；特别地，
$x e \in M$。存在元素 $z_1 \in M$，满足 $x z_1 = x e$。因此
$x(z_1 - e) = 0$。把 $e$ 替换为 $e - z_1$，可以假设 $e$ 被 $x$ 零化。
所以只需证明
$$
\varphi[x] :
M[x]
\longrightarrow
\left(\bigoplus\nolimits_{j \in J} E\right)[x] =
\bigoplus\nolimits_{j \in J} E[x]
$$
是满射。若 $d = 1$，则由 $\varphi$ 的构造，这一点成立。若 $d > 1$，
则注意 $E[x]$ 是正则环 $A/xA$ 的剩余域的单射包；见《对偶复形》，引理
\ref{dualizing-lemma-quotient}。注意，$M[x]$ 作为 $A/xA$-模是
$\mathfrak m/(x)$-幂挠的，并且有
\begin{align*}
M[x] \otimes_{A/xA, F} A/xA
& = M[x] \otimes_{A, F} A \otimes_A A/xA \\
& = (M \otimes_{A, F} A)[x^p] \otimes_A A/xA \\
& \cong M[x^p] \otimes_A A/xA
\end{align*}
如前一样利用 $F$ 的平坦性论证。我们断言
$M[x^p] \otimes_A A/xA \to M[x]$，
$z \otimes 1 \mapsto x^{p - 1}z$，是同构。只需对上面定义的每个模
$M_n$，$n > 0$，证明这一点；而在这种情形下，它由
$A/(x_1^{p^n}, \ldots, x_d^{p^n})$ 与 $x = x_1$ 的相同结论得出。
因此，对 $\dim(A)$ 作归纳，可知 $\varphi[x]$ 如所需为满射。
\end{proof}





\section{局部上同调上的附加结构}
\label{local-cohomology-section-additional}

\noindent
下面给出一个示例性结果。

\begin{lemma}
\label{local-cohomology-lemma-derivation}
设 $A$ 是环，$I \subset A$ 是有限生成理想。令 $Z = V(I)$。
对每个导子 $\theta : A \to A$，都存在一个典范加性算子 $D$，它作用于
局部上同调模 $H^i_Z(A)$，并满足关于 $\theta$ 的 Leibniz 法则。
\end{lemma}

\begin{proof}
设 $f_1, \ldots, f_r$ 是生成 $I$ 的元素。回顾 $R\Gamma_Z(A)$ 由复形
$$
A \to \prod\nolimits_{i_0} A_{f_{i_0}} \to
\prod\nolimits_{i_0 < i_1} A_{f_{i_0}f_{i_1}}
\to \ldots \to A_{f_1\ldots f_r}
$$
计算；见《对偶复形》，引理 \ref{dualizing-lemma-local-cohomology-adjoint}。
由于 $\theta$ 唯一延拓为 $A$ 的任意局部化上的加性算子，并且该算子满足
关于 $\theta$ 的 Leibniz 法则，本引理显然成立。
\end{proof}

\begin{lemma}
\label{local-cohomology-lemma-frobenius}
设 $p$ 是素数。设 $A$ 是满足 $p = 0$ 的环。记
$F : A \to A$，$a \mapsto a^p$，为 Frobenius 自同态。
设 $I \subset A$ 是有限生成理想。令 $Z = V(I)$。存在同构
$R\Gamma_Z(A) \otimes_{A, F}^\mathbf{L} A \cong R\Gamma_Z(A)$.
\end{lemma}

\begin{proof}
这由《对偶复形》，引理
\ref{dualizing-lemma-torsion-change-rings}，以及
$Z = V(f_1^p, \ldots, f_r^p)$ 这一事实得出，其中
$I = (f_1, \ldots, f_r)$。
\end{proof}

\begin{lemma}
\label{local-cohomology-lemma-etale-derivation}
设 $A$ 是环。设 $V \to \Spec(A)$ 是拟紧、拟分离且 \'etale 的。
对每个导子 $\theta : A \to A$，都存在一个典范加性算子 $D$，它作用于
$H^i(V, \mathcal{O}_V)$，并满足关于 $\theta$ 的 Leibniz 法则。
\end{lemma}

\begin{proof}
若 $V$ 分离，则可使用仿射开覆盖
$V = \bigcup_{j = 1, \ldots m} V_j$ 来论证。事实上，由于 $V$ 分离，
可以写成 $V_{j_0 \ldots j_p} = \Spec(B_{j_0 \ldots j_p})$；见《概形》，
引理 \ref{schemes-lemma-characterize-separated}。于是可知 $A$-模
$H^i(V, \mathcal{O}_V)$ 是 {\v C}ech 复形的第 $i$ 个上同调群：
$$
\prod B_{j_0} \to
\prod B_{j_0j_1} \to
\prod B_{j_0j_1j_2} \to \ldots
$$
见《概形的上同调》，引理
\ref{coherent-lemma-cech-cohomology-quasi-coherent}。
每个 $B = B_{j_0 \ldots j_p}$ 都是 \'etale $A$-代数。因此
$\Omega_B = \Omega_A \otimes_A B$，由此可知 $\theta$ 唯一延拓为导子
$\theta_B : B \to B$。这些同态定义了 {\v C}ech 复形的一个自同态，
并在上同调群上定义了所需算子。

\medskip\noindent
在一般情形下，使用 $V$ 的由仿射开集组成的超覆盖，恰如《概形的上同调》，
引理 \ref{coherent-lemma-hypercoverings} 的证明第一部分所做的那样。
略去细节。
\end{proof}

\begin{remark}
\label{local-cohomology-remark-higher-order-operators}
可以加强引理 \ref{local-cohomology-lemma-derivation} 和引理 \ref{local-cohomology-lemma-etale-derivation}，
使其包含高阶微分算子；为此使用《代数》，引理
\ref{algebra-lemma-formally-etale-principal-parts} 与评注
\ref{algebra-remark-formally-etale-differential-operators}。
若以后需要这一点，我们将在此陈述并证明一个精确的引理。
\end{remark}

\begin{lemma}
\label{local-cohomology-lemma-etale-frobenius}
设 $p$ 是素数。设 $A$ 是满足 $p = 0$ 的环。记
$F : A \to A$，$a \mapsto a^p$，为 Frobenius 自同态。
若 $V \to \Spec(A)$ 是拟紧、拟分离且 \'etale 的，则存在同构
$R\Gamma(V, \mathcal{O}_V) \otimes_{A, F}^\mathbf{L} A \cong
R\Gamma(V, \mathcal{O}_V)$。
\end{lemma}

\begin{proof}
注意，相对 Frobenius 态射
$$
V \longrightarrow V \times_{\Spec(A), \Spec(F)} \Spec(A)
$$
是 $V$ 在 $A$ 上的一个同构；见《\'Etale 态射》，引理
\ref{etale-lemma-relative-frobenius-etale}。因此，本引理由上同调与基变换得出；
见《概形的导出范畴》，引理
\ref{perfect-lemma-compare-base-change}。注意，由于 $V$ 在 $A$ 上 \'etale，
它在 $A$ 上平坦。
\end{proof}





\section{某种一致性，I}
\label{local-cohomology-section-uniform-vanishing}

\noindent
本节的主要任务是陈述并证明引理
\ref{local-cohomology-lemma-characterize-vanishing-tor-ext-above-e}。

\begin{lemma}
\label{local-cohomology-lemma-map-tor-1-zero}
设 $R$ 是环。设 $M \to M'$ 是 $R$-模同态，其中 $M$ 有限展示，并且同态
$\text{Tor}_1^R(M, N) \to \text{Tor}_1^R(M', N)$ 对所有 $R$-模 $N$
均为零。则 $M \to M'$ 经由一个自由 $R$-模分解。
\end{lemma}

\begin{proof}
可以选取短正合列之间的同态
$$
\xymatrix{
0 \ar[r] &
K \ar[r] \ar[d] &
R^{\oplus r} \ar[r] \ar[d] &
M \ar[r] \ar[d] &
0 \\
0 \ar[r] &
K' \ar[r] &
\bigoplus_{i \in I} R \ar[r] &
M' \ar[r] &
0
}
$$
使其右侧竖直箭头为给定同态。可以让此同态经由短正合列
\begin{equation}
\label{local-cohomology-equation-pushout}
0 \to K' \to E \to M \to 0
\end{equation}
分解；该列是第一个短正合列沿 $K \to K'$ 的推出。追图可见，本引理的假设
蕴含由 (\ref{local-cohomology-equation-pushout}) 诱导的连接同态
$\text{Tor}_1^R(M, N) \to K' \otimes_R N$ 为零；也就是说，序列
(\ref{local-cohomology-equation-pushout}) 普遍正合。由《代数》，引理
\ref{algebra-lemma-universally-exact-split}，这蕴含
(\ref{local-cohomology-equation-pushout}) 分裂（这里用到了 $M$ 有限展示）。因此，同态
$M \to M'$ 经由 $\bigoplus_{i \in I} R$ 分解，结论得证。
\end{proof}

\begin{lemma}
\label{local-cohomology-lemma-characterize-vanishing-tor-ext-above-e}
设 $R$ 是环。设 $\alpha : M \to M'$ 是 $R$-模同态。设
$P_\bullet \to M$ 与 $P'_\bullet \to M'$ 是由投射 $R$-模组成的解消。
设 $e \geq 0$ 是整数。考虑下列条件：
\begin{enumerate}
\item 可以找到复形同态 $a_\bullet : P_\bullet \to P'_\bullet$，
它在上同调上诱导 $\alpha$，且 $a_i = 0$ 对 $i > e$ 成立。
\item 可以找到复形同态 $a_\bullet : P_\bullet \to P'_\bullet$，
它在上同调上诱导 $\alpha$，且 $a_{e + 1} = 0$。
\item 同态 $\Ext^i_R(M', N) \to \Ext^i_R(M, N)$ 对所有
$R$-模 $N$ 以及 $i > e$ 均为零。
\item 同态 $\Ext^{e + 1}_R(M', N) \to \Ext^{e + 1}_R(M, N)$
对所有 $R$-模 $N$ 均为零。
\item 令 $N = \Im(P'_{e + 1} \to P'_e)$，并以
$\xi \in \Ext^{e + 1}_R(M', N)$ 表示典范元素（见证明）。则 $\xi$
在 $\Ext^{e + 1}_R(M, N)$ 中映为零。
\item 同态 $\text{Tor}_i^R(M, N) \to \text{Tor}_i^R(M', N)$
对所有 $R$-模 $N$ 以及 $i > e$ 均为零。
\item 同态 $\text{Tor}_{e + 1}^R(M, N) \to \text{Tor}_{e + 1}^R(M', N)$
对所有 $R$-模 $N$ 均为零。
\end{enumerate}
则始终有蕴含关系
$$
(1) \Leftrightarrow (2) \Leftrightarrow (3) \Leftrightarrow (4)
\Leftrightarrow (5) \Rightarrow (6) \Leftrightarrow (7)
$$
若 $M$ 是 $(-e - 1)$-伪凝聚的（例如，若 $R$ 诺特且 $M$ 是有限
$R$-模），则所有条件等价。
\end{lemma}

\begin{proof}
显然，(2) 蕴含 (1)。若 $a_\bullet$ 如 (1) 所述，则可考虑复形同态
$a'_\bullet : P_\bullet \to P'_\bullet$，其中 $a'_i = a_i$ 对
$i \leq e + 1$ 成立，而 $a'_i = 0$ 对 $i \geq e + 1$ 成立，从而得到如 (2)
所述的复形同态。因此，(1) 等价于 (2)。

\medskip\noindent
由使用解消构造 $\Ext$ 与 $\text{Tor}$ 函子的方法（《代数》，
第 \ref{algebra-section-ext} 节与第 \ref{algebra-section-tor} 节），
可见 (1) 和 (2) 蕴含所有其他条件。

\medskip\noindent
显然，(3) 蕴含 (4)，而 (4) 蕴含 (5)。设 $N$ 如 (5) 所述。
典范同态 $\tilde \xi : P'_{e + 1} \to N$ 与
$P'_{e + 2} \to P'_{e + 1}$ 的预复合为零。因此，可以在下式中考虑类
$\xi$，即 $\tilde \xi$ 的类：
$$
\Ext^{e + 1}_R(M', N) =
\frac{\Ker(\Hom(P'_{e + 1}, N \to \Hom(P'_{e + 2}, N)}{
\Im(\Hom(P'_e, N \to \Hom(P'_{e + 1}, N)}
$$
选取复形同态 $a_\bullet : P_\bullet \to P'_\bullet$，它提升 $\alpha$；
见《导出范畴》，引理 \ref{derived-lemma-morphisms-lift-projective}。
若 $\xi$ 在 $\Ext^{e + 1}_R(M', N)$ 中映为零，则可找到同态
$\varphi : P_e \to N$，使得
$\tilde \xi  \circ a_{e + 1} = \varphi \circ d$。从而得到复形同态
$$
\xymatrix{
\ldots \ar[r] &
P_{e + 1} \ar[r] \ar[d]^0 &
P_e \ar[r] \ar[d]^{a_e - \varphi} &
P_{e - 1} \ar[r] \ar[d]^{a_{e - 1}} &
\ldots \\
\ldots \ar[r] &
P'_{e + 1} \ar[r] &
P'_e \ar[r] &
P'_{e - 1} \ar[r] &
\ldots
}
$$
如 (2) 所述。因此，(1)--(5) 等价。

\medskip\noindent
(6) 与 (7) 的等价性来自维数平移；略去细节。

\medskip\noindent
假设 $M$ 是 $(-e - 1)$-伪凝聚的。（本引理括号内的陈述来自
《代数进阶》，引理
\ref{more-algebra-lemma-Noetherian-pseudo-coherent}。）
我们将证明 (7) 蕴含 (4)，这将完成证明。对 $e$ 作归纳。基础情形为
$e = 0$。此时，由《代数进阶》，引理
\ref{more-algebra-lemma-n-pseudo-module}，$M$ 有限展示；再由引理
\ref{local-cohomology-lemma-map-tor-1-zero}，可知 $M \to M'$ 经由一个自由模分解。
当然，若 $M \to M'$ 经由一个自由模分解，则
$\Ext^i_R(M', N) \to \Ext^i_R(M, N)$ 对所有 $i > 0$ 均为零，正如所需。
假设 $e > 0$。可以选取短正合列之间的同态
$$
\xymatrix{
0 \ar[r] &
K \ar[r] \ar[d] &
R^{\oplus r} \ar[r] \ar[d] &
M \ar[r] \ar[d] &
0 \\
0 \ar[r] &
K' \ar[r] &
\bigoplus_{i \in I} R \ar[r] &
M' \ar[r] &
0
}
$$
使其右侧竖直箭头为给定同态。我们得到
$\text{Tor}_{i + 1}^R(M, N) = \text{Tor}^R_i(K, N)$
以及 $\Ext^{i + 1}_R(M, N) = \Ext^i_R(K, N)$，其中 $i \geq 1$，
且对所有 $R$-模 $N$ 都成立；对 $M', K'$ 也有类似结论。因此可见，同态
$\text{Tor}_e^R(K, N) \to \text{Tor}_e^R(K', N)$ 对所有 $R$-模
$N$ 均为零。由《代数进阶》，引理
\ref{more-algebra-lemma-cone-pseudo-coherent}，可见 $K$ 是
$(-e)$-伪凝聚的。由归纳假设，可知同态
$\Ext^e(K', N) \to \Ext^e(K, N)$ 对所有 $R$-模 $N$ 均为零，
这便给出所需结论。
\end{proof}

\begin{lemma}
\label{local-cohomology-lemma-cd-sequence-Koszul}
设 $I$ 是诺特环 $A$ 的理想。对所有 $n \geq 1$，都存在 $m > n$，
使得同态 $A/I^m \to A/I^n$ 满足引理
\ref{local-cohomology-lemma-characterize-vanishing-tor-ext-above-e} 中的等价条件，其中
$e = \text{cd}(A, I)$。
\end{lemma}

\begin{proof}
设 $\xi \in \Ext^{e + 1}_A(A/I^n, N)$ 为引理
\ref{local-cohomology-lemma-characterize-vanishing-tor-ext-above-e} 第 (5) 部分所构造的元素。
由于 $e = \text{cd}(A, I)$，有
$0 = H^{e + 1}_Z(N) = H^{e + 1}_I(N) = \colim \Ext^{e + 1}(A/I^m, N)$
；这是由《对偶复形》，引理
\ref{dualizing-lemma-local-cohomology-noetherian} 与
\ref{dualizing-lemma-local-cohomology-ext} 得出的。因此，可以选取
$m \geq n$，使得 $\xi$ 在 $\Ext^{e + 1}_A(A/I^m, N)$ 中映为零，正如所需。
\end{proof}






\section{某种一致性，II}
\label{local-cohomology-section-uniform-vanishing-bis}

\noindent
设 $I$ 是诺特环 $A$ 的理想，$M$ 是有限 $A$-模，且 $i > 0$。
由《代数进阶》，引理
\ref{more-algebra-lemma-tor-strictly-pro-zero}，存在
$c = c(A, I, M, i)$，使得
$\text{Tor}^A_i(M, A/I^n) \to \text{Tor}^A_i(M, A/I^{n - c})$
对所有 $n \geq c$ 均为零。本节讨论若干结果，它们表明有时可以选取一个常数
$c$，使其同时适用于所有 $A$-模 $M$（以及某个范围内的指标 $i$）。
这些材料与 \cite{Huneke-uniform} 和 \cite{AHS} 中讨论的一致
Artin--Rees 性质有关。

\medskip\noindent
在评注 \ref{local-cohomology-remark-strict-pro-isomorphism} 中，我们将应用这一点来说明
与导出完备化有关的若干 pro-系统严格 pro-同构（或者并非如此）。

\medskip\noindent
下面的引理可以显著加强。

\begin{lemma}
\label{local-cohomology-lemma-maps-zero-fixed-torsion}
设 $I$ 是诺特环 $A$ 的理想。对每个 $m \geq 0$ 和 $i > 0$，都存在
$c = c(A, I, m, i) \geq 0$，使得对每个 $A$-模 $M$（被 $I^m$ 零化），同态
$$
\text{Tor}^A_i(M, A/I^n) \to \text{Tor}^A_i(M, A/I^{n - c})
$$
对所有 $n \geq c$ 均为零。
\end{lemma}

\begin{proof}
对 $i$ 作归纳。基础情形是 $i = 1$。短正合列
$0 \to I^n \to A \to A/I^n \to 0$ 给出单射
$\text{Tor}_1^A(M, A/I^n) \subset I^n \otimes_A M$；见《代数》，评注
\ref{algebra-remark-Tor-ring-mod-ideal}。由于 $M$ 被 $I^m$ 零化，可见同态
$I^n \otimes_A M \to I^{n - m} \otimes_A M$ 对 $n \geq m$ 为零。
因此，取 $c = m$ 时结论成立。

\medskip\noindent
归纳步骤。设 $i > 1$，并假设 $c$ 对 $i - 1$ 有效。将《代数进阶》，引理
\ref{more-algebra-lemma-tor-strictly-pro-zero} 应用于 $M = A/I^m$，
可以选取 $c' \geq 0$，使得
$\text{Tor}_i(A/I^m, A/I^n) \to \text{Tor}_i(A/I^m, A/I^{n - c'})$
对 $n \geq c'$ 为零。设 $M$ 被 $I^m$ 零化。选取短正合列
$$
0 \to S \to \bigoplus\nolimits_{i \in I} A/I^m \to M \to 0
$$
相应的 Tor 长正合列给出正合列
$$
\text{Tor}_i^A(\bigoplus\nolimits_{i \in I} A/I^m, A/I^n) \to
\text{Tor}_i^A(M, A/I^n) \to
\text{Tor}_{i - 1}^A(S, A/I^n)
$$
对所有整数 $n \geq 0$ 均成立。若 $n \geq c + c'$，则同态
$\text{Tor}_{i - 1}^A(S, A/I^n) \to \text{Tor}_{i - 1}^A(S, A/I^{n - c})$
为零，并且同态 $\text{Tor}_i^A(A/I^m, A/I^{n - c}) \to 
\text{Tor}_i^A(A/I^m, A/I^{n - c - c'})$ 为零。结合这些短正合列，
可知结论对 $i$ 成立，其常数为 $c + c'$。
\end{proof}

\begin{lemma}
\label{local-cohomology-lemma-annihilates-affine}
设 $I = (a_1, \ldots, a_t)$ 是诺特环 $A$ 的理想。令 $a = a_1$，
并以 $B = A[\frac{I}{a}]$ 表示仿射爆破代数。存在 $c > 0$，使得
$\text{Tor}_i^A(B, M)$ 被 $I^c$ 零化；这对所有 $A$-模 $M$ 以及
$i \geq t$ 均成立。
\end{lemma}

\begin{proof}
回顾 $B$ 是
$A[x_2, \ldots, x_t]/(a_1x_2 - a_2, \ldots, a_1x_t - a_t)$
对其 $a_1$-挠子模取商所得的商；见《代数》，引理
\ref{algebra-lemma-affine-blowup-quotient-description}。令
$$
B_\bullet = \text{关于下列元素的 Koszul 复形： }a_1x_2 - a_2, \ldots, a_1x_t - a_t
\text{ 相对于 }A[x_2, \ldots, x_t]
$$
并将其视为位于次数 $(t - 1), \ldots, 0$ 的链复形。复形
$B_\bullet[1/a_1]$ 同构于由
$x_2 - a_2/a_1, \ldots, x_t - a_t/a_1$ 构成的 Koszul 复形；该序列是
$A[1/a_1][x_2, \ldots, x_t]$ 中的正则序列。由于正则序列是 Koszul 正则的，
可知增广
$$
\epsilon : B_\bullet \longrightarrow B
$$
在逆转 $a_1$ 后是拟同构。由于锥 $C_\bullet$（$\epsilon$ 的锥）的同调模是有限
$A[x_2, \ldots, x_n]$-模，且 $C_\bullet$ 有界，可知存在 $c \geq 0$，
使得 $a_1^c$ 将它们全部零化。由《导出范畴》，引理
\ref{derived-lemma-trick-vanishing-composition}，这蕴含在必要时将 $c$
替换为更大的整数后，$a_1^c$ 作用于 $C_\bullet$ 的同态在 $D(A)$ 中为零。
考察特异三角即可完成证明：
$$
B_\bullet \otimes_A^\mathbf{L} M \to
B \otimes_A^\mathbf{L} M \to
C_\bullet \otimes_A^\mathbf{L} M
$$
事实上，第一项由 $B_\bullet \otimes_A M$ 表示，它位于同调次数
$(t - 1), \ldots, 0$，因为 Koszul 复形 $B_\bullet$ 的各项都是自由
（因而平坦）$A$-模。因此
$\text{Tor}_i^A(B, M) = H_i(C_\bullet \otimes_A^\mathbf{L} M)$
对 $i > t - 1$ 成立，并且它被 $a_1^c$ 零化。由于
$a_1^cB = I^cB$，且该 Tor 模是 $B$-模，结论随之得出。
\end{proof}

\noindent
在本节余下的讨论中，固定诺特环 $A$ 和理想 $I \subset A$。以
$$
p : X \to \Spec(A)
$$
表示 $\Spec(A)$ 沿理想 $I$ 的爆破。换言之，$X$ 是 $\text{Proj}$；
这里使用的 Rees 代数为 $\bigoplus_{n \geq 0} I^n$。由《概形的上同调》，引理
\ref{coherent-lemma-coherent-on-proj} 与
\ref{coherent-lemma-recover-tail-graded-module}，可以选取整数
$q(A, I) \geq 0$，使得对所有 $q \geq q(A, I)$，都有
$H^i(X, \mathcal{O}_X(q)) = 0$（当 $i > 0$），并且
$H^0(X, \mathcal{O}_X(q)) = I^q$。

\begin{lemma}
\label{local-cohomology-lemma-compute-tor-Iq}
在上述情形下，对 $q \geq q(A, I)$ 以及任意 $A$-模 $M$，都有
$$
R\Gamma(X, Lp^*\widetilde{M}(q)) \cong M \otimes_A^\mathbf{L} I^q
$$
该同构成立于 $D(A)$ 中。
\end{lemma}

\begin{proof}
选取自由解消 $F_\bullet \to M$。则 $\widetilde{F}_\bullet$ 是
$\widetilde{M}$ 的平坦解消。因此，$Lp^*\widetilde{M}$ 由复形
$p^*\widetilde{F}_\bullet$ 给出，从而 $Lp^*\widetilde{M}(q)$ 由复形
$p^*\widetilde{F}_\bullet(q)$ 给出。$p^*\widetilde{F}_i(q)$ 关于
$\Gamma(X, -)$ 右无环，这是由我们对 $q \geq q(A, I)$ 的选择得出的；
并且我们有 $\Gamma(X, p^*\widetilde{F}_i(q)) = I^qF_i$，这也由我们对
$q \geq q(A, I)$ 的选择得出。于是，由《概形的导出范畴》，引理
\ref{perfect-lemma-acyclicity-lemma-global}，
$R\Gamma(X, Lp^*\widetilde{M}(q))$ 由各项为 $I^qF_i$ 的复形给出。
依定义，复形 $I^qF_\bullet$ 计算 $M \otimes_A^\mathbf{L} I^q$，故结论成立。
\end{proof}

\begin{lemma}
\label{local-cohomology-lemma-annihilates}
在上述情形下，设 $t$ 是 $I$ 的生成元个数的一个上界。存在整数
$c = c(A, I) \geq 0$，使得对任意 $A$-模 $M$，上同调层
$H^j(Lp^*\widetilde{M})$ 被 $I^c$ 零化，其中 $j \leq -t$。
\end{lemma}

\begin{proof}
写成 $I = (a_1, \ldots, a_t)$。该问题在 $X$ 上是仿射局部的。
对 $1 \leq i \leq t$，令 $B_i = A[\frac{I}{a_i}]$ 为仿射爆破代数。
则 $X$ 有一个仿射开覆盖，它由环 $B_i$ 的谱构成；见《除子》，引理
\ref{divisors-lemma-blowing-up-affine}。由《概形的导出范畴》，引理
\ref{perfect-lemma-quasi-coherence-pullback} 所给出的导出拉回描述，
只需证明对每个 $i$，都存在 $c \geq 0$，使得
$$
\text{Tor}_j^A(B_i, M)
$$
被 $I^c$ 零化，其中 $j \geq t$。这正是引理
\ref{local-cohomology-lemma-annihilates-affine}。
\end{proof}

\begin{lemma}
\label{local-cohomology-lemma-annihilates-tors}
在上述情形下，设 $t$ 是 $I$ 的生成元个数的一个上界。存在整数
$c = c(A, I) \geq 0$，使得对任意 $A$-模 $M$，Tor 模
$\text{Tor}_i^A(M, A/I^q)$ 被 $I^c$ 零化，其中 $i > t$，且
$q \geq 0$ 任意。
\end{lemma}

\begin{proof}
令 $q(A, I)$ 如上。对 $q \geq q(A, I)$，有
$$
R\Gamma(X, Lp^*\widetilde{M}(q)) = M \otimes_A^\mathbf{L} I^q
$$
，这是由引理 \ref{local-cohomology-lemma-compute-tor-Iq} 得出的。我们有有界收敛谱序列
$$
H^a(X, H^b(Lp^*\widetilde{M}(q))) \Rightarrow
\text{Tor}_{-a - b}^A(M, I^q)
$$
；见《概形的导出范畴》，引理 \ref{perfect-lemma-spectral-sequence}。
设 $d$ 是《概形的上同调》，引理
\ref{coherent-lemma-vanishing-nr-affines-quasi-separated}
中所述的整数（事实上可以取 $d = t$；见《概形的上同调》，引理
\ref{coherent-lemma-vanishing-nr-affines}）。于是可见
$H^{-i}(X, Lp^*\widetilde{M}(q)) = \text{Tor}_i^A(M, I^q)$
具有至多 $d$ 步的有限滤过，其分次对象是下列模的子商：
$$
H^a(X, H^{- i - a}(Lp^*\widetilde{M})(q)),\quad
a = 0, 1, \ldots, d - 1
$$
若 $i \geq t$，则所有这些模都被 $I^c$ 零化，其中
$c = c(A, I)$ 如引理 \ref{local-cohomology-lemma-annihilates} 所述；这是因为由该引理，
上同调层 $H^{- i - a}(Lp^*\widetilde{M})$ 均被 $I^c$ 零化。因此可见，
$\text{Tor}_i^A(M, I^q)$ 被 $I^{dc}$ 零化，其中 $q \geq q(A, I)$ 且
$i \geq t$。利用短正合列 $0 \to I^q \to A \to A/I^q \to 0$，
可知 $\text{Tor}_i(M, A/I^q)$ 被 $I^{dc}$ 零化，其中
$q \geq q(A, I)$ 且 $i > t$。于是可知 $I^m$（其中
$m = \max(dc, q(A, I) - 1)$）零化 $\text{Tor}_i^A(M, A/I^q)$，
这对所有 $q \geq 0$ 且 $i > t$ 成立，正如所需。
\end{proof}

\begin{lemma}
\label{local-cohomology-lemma-tor-maps-vanish}
设 $I$ 是诺特环 $A$ 的理想。设 $t \geq 0$ 是 $I$ 的生成元个数的一个上界。
存在 $N, c \geq 0$，使得同态
$$
\text{Tor}_{t + 1}^A(M, A/I^n) \to \text{Tor}_{t + 1}^A(M, A/I^{n - c})
$$
对任意 $A$-模 $M$ 以及所有 $n \geq N$ 均为零。
\end{lemma}

\begin{proof}
令 $c_1$ 为引理 \ref{local-cohomology-lemma-annihilates-tors} 中得到的常数。
请注意，此常数 $c_1$ 对 $\text{Tor}_i$ 同时有效，其中 $i > t$ 任意。

\medskip\noindent
写成 $I = (a_1, \ldots, a_t)$。对 $A$-模 $M$，令
$$
\ell(M) =
\#\{i \mid 1 \leq i \leq t,\ a_i^{c_1}\text{ 在下列对象上为零： }M\}
$$
这是 $\{0, 1, \ldots, t\}$ 的一个元素。我们将对 $0 \leq s \leq t$
作降归纳，证明下列陈述 $H_s$：存在 $N, c \geq 0$，使得对每个模
$M$（满足 $\ell(M) \geq s$），同态
$$
\text{Tor}_{t + 1 + i}^A(M, A/I^n) \to \text{Tor}_{t + 1 + i}^A(M, A/I^{n - c})
$$
对 $i = 0, \ldots, s$ 以及所有 $n \geq N$ 均为零。

\medskip\noindent
基础情形：$s = t$。若 $\ell(M) = t$，则 $M$ 被
$(a_1^{c_1}, \ldots, a_t^{c_1}\}$ 零化，因而被
$I^{t(c_1 - 1) + 1}$ 零化。由引理
\ref{local-cohomology-lemma-maps-zero-fixed-torsion}，可知 $H_t$ 成立：令 $c = N$ 为该引理中
得到的整数 $c(A, I, t(c_1 - 1) + 1, t + 1), \ldots, c(A, I, t(c_1 - 1) + 1, 2t + 1)$
的最大值即可。

\medskip\noindent
归纳步骤。设 $0 \leq s < t$，并设已有 $N, c$，如 $H_{s + 1}$ 所述。
考虑模 $M$，它满足 $\ell(M) = s$。可以选取 $i$，使得 $a_i^{c_1}$
在 $M$ 上不为零。于是 $\ell(M[a_i^c]) \geq s + 1$ 且
$\ell(M/a_i^{c_1}M) \geq s + 1$，归纳假设适用于这两个模。考虑正合列
$$
0 \to M[a_i^{c_1}] \to M \xrightarrow{a_i^{c_1}} M \to M/a_i^{c_1}M \to 0
$$
以 $E \subset M$ 表示中间箭头的像。考虑相应的 Tor 模图表
$$
\xymatrix{
&
&
\text{Tor}_{i + 1}(M/a_i^{c_1}M, A/I^q) \ar[d] \\
\text{Tor}_i(M[a_i^{c_1}], A/I^q) \ar[r] &
\text{Tor}_i(M, A/I^q) \ar[r] \ar[rd]^0 &
\text{Tor}_i(E, A/I^q) \ar[d] \\
&
&
\text{Tor}_i(M, A/I^q)
}
$$
其行与列均正合（对每个 $q$）。由我们对 $c_1$ 的选择，东南方向的箭头为零。
由此可知，模 $\text{Tor}_i(M, A/I^q)$ 夹在
$\text{Tor}_i(M[a_i^{c_1}], A/I^q)$ 的一个商模与
$\text{Tor}_{i + 1}(M/a_i^{c_1}M, A/I^q)$ 的一个子模之间。因此可知，
$H_s$ 成立，其中将 $N$ 替换为 $N + c$，并将 $c$ 替换为 $2c$。
略去一些细节。
\end{proof}

\begin{proposition}
\label{local-cohomology-proposition-uniform-artin-rees}
设 $I$ 是诺特环 $A$ 的理想。设 $t \geq 0$ 是 $I$ 的生成元个数的一个上界。
存在 $N, c \geq 0$，使得当 $n \geq N$ 时，同态
$$
A/I^n \to A/I^{n - c}
$$
满足引理 \ref{local-cohomology-lemma-characterize-vanishing-tor-ext-above-e} 中的等价条件，
其中 $e = t$。
\end{proposition}

\begin{proof}
这是引理
\ref{local-cohomology-lemma-tor-maps-vanish} 与
\ref{local-cohomology-lemma-characterize-vanishing-tor-ext-above-e} 的直接推论。
\end{proof}

\begin{remark}
\label{local-cohomology-remark-better-bound}
论文 \cite{AHS} 除了许多其他结果外，还证明了：若 $A$ 是局部的，
则命题 \ref{local-cohomology-proposition-uniform-artin-rees} 在将 $e = t$ 替换为
$e = \dim(A)$ 后仍成立。考察引理 \ref{local-cohomology-lemma-cd-sequence-Koszul}，
自然会问命题 \ref{local-cohomology-proposition-uniform-artin-rees} 在将 $e = t$ 替换为
$e = \text{cd}(A, I)$ 后是否成立。我们不知道答案。
\end{remark}

\begin{remark}
\label{local-cohomology-remark-strict-pro-isomorphism}
设 $I$ 是诺特环 $A$ 的理想。写成 $I = (f_1, \ldots, f_r)$。
以 $K_n^\bullet$ 表示《代数进阶》，情形
\ref{more-algebra-situation-koszul} 中由 $f_1^n, \ldots, f_r^n$
构成的 Koszul 复形，并以 $K_n \in D(A)$ 表示相应对象。
设 $M^\bullet$ 是由有限 $A$-模组成的有界复形，并以 $M \in D(A)$
表示相应对象。考虑 $D(A)$ 中的下列逆系统：
\begin{enumerate}
\item $M^\bullet/I^nM^\bullet$，即各项为 $M^i/I^nM^i$ 的复形；
\item $M \otimes_A^\mathbf{L} A/I^n$；
\item $M \otimes_A^\mathbf{L} K_n$；
\item $M \otimes_P^\mathbf{L} P/J^n$（见下文）。
\end{enumerate}
所有这些逆系统作为 pro-对象同构：(2) 与 (3) 之间的同构来自《代数进阶》，
引理 \ref{more-algebra-lemma-sequence-Koszul-complexes}；(1) 与 (2)
之间的同构由《代数进阶》，引理
\ref{more-algebra-lemma-derived-completion-plain-completion} 给出。
最后一个见下文。

\medskip\noindent
不过，可以询问这些 pro-系统的同构是否“严格”；这一术语与问题同
\cite[pages 61, 62]{quillenhomology} 中的讨论有关。具体而言，给定范畴
$\mathcal{C}$，可以定义一个“严格 pro-范畴”：其对象是逆系统 $(X_n)$，
而态射 $(X_n) \to (Y_n)$ 由二元组 $(c, \varphi_n)$ 给出；该二元组由
$c \geq 0$ 以及态射 $\varphi_n : X_n \to Y_{n - c}$ 组成，后者对所有
$n \geq c$ 给定；它们满足显然的相容条件，并且还要模去
某种等价关系（本质上由增大 $c$ 给出）。于是可以询问上述逆系统在这个严格
pro-范畴中是否同构。

\medskip\noindent
显然，即使 $M = A[0]$，(1) 与 (3) 也不可能如此。事实上，一般而言，
系统 $H^0(K_n) = A/(f_1^n, \ldots, f_r^n)$ 在模范畴中与系统 $A/I^n$
并非严格 pro-同构。例如，若取 $A = \mathbf{Z}[x_1, \ldots, x_r]$ 且
$f_i = x_i$，则 $H^0(K_n)$ 不被 $I^{r(n - 1)}$ 零化。\footnote{当然，
可以询问这些 pro-系统在如下范畴中是否同构：其对象为逆系统，而态射由三元组
$(r, c, \varphi_n)$ 给出；该三元组由 $r \geq 1$、$c \geq 0$ 以及态射
$\varphi_n : X_{rn} \to Y_{n - c}$ 组成，后者对 $n \geq c$ 给定。}

\medskip\noindent
事实上，上述结果表明，《代数进阶》，引理
\ref{more-algebra-lemma-derived-completion-plain-completion} 中讨论的从 (2)
到 (1) 的自然同态是严格 pro-同构。下面概述证明。先利用涉及朴素截断的标准论证，
约化到 $M^\bullet$ 由单个有限 $A$-模 $M$ 给出的情形，其中该模置于次数
$0$。选取命题 \ref{local-cohomology-proposition-uniform-artin-rees} 中的
$N, c \geq 0$。该命题蕴含：对 $n \geq N$，存在分解
$$
M \otimes_A^\mathbf{L} A/I^n
\to
\tau_{\geq -t}(M \otimes_A^\mathbf{L} A/I^n)
\to
M \otimes_A^\mathbf{L} A/I^{n - c}
$$
，它们是系统 (2) 中转移同态的分解。另一方面，由《代数进阶》，引理
\ref{more-algebra-lemma-tor-strictly-pro-zero}，可以找到另一个常数
$c' = c'(M) \geq 0$，使得同态
$\text{Tor}_i^A(M, A/I^{n'}) \to \text{Tor}_i(M, A/I^{n' - c'})$
对 $i = 1, 2, \ldots, t$ 以及 $n' \geq c'$ 均为零。于是由《导出范畴》，
引理 \ref{derived-lemma-trick-vanishing-composition}，同态
$$
\tau_{\geq -t}(M \otimes_A^\mathbf{L} A/I^{n + tc'})
\to
\tau_{\geq -t}(M \otimes_A^\mathbf{L} A/I^n)
$$
经由 $M \otimes_A^\mathbf{L}A/I^{n + tc'} \to M/I^{n + tc'}M$ 分解。
结合前一结果，得到分解
$$
M \otimes_A^\mathbf{L}A/I^{n + tc'} \to M/I^{n + tc'}M
\to M \otimes_A^\mathbf{L} A/I^{n - c}
$$
，这便给出所需结论。若以后需要这一结果，我们将仔细陈述并给出详细证明。

\medskip\noindent
对第 (4) 项，假设有诺特环 $P$、环同态 $P \to A$ 以及理想
$J \subset P$，并且 $I = JA$。由《代数进阶》，第
\ref{more-algebra-section-derived-base-change} 节，得到函子
$M \otimes_P^\mathbf{L} - : D(P) \to D(A)$，并得到如 (4) 中的逆系统
$M \otimes_P^\mathbf{L} P/J^n$，它位于 $D(A)$ 中。若 $P$ 诺特，
则 (4) 中的系统与 (1) 中的系统 pro-同构，因为可以用 Koszul 复形作比较。
若 $P \to A$ 有限，则系统 (4) 与系统 (2) 严格 pro-同构，因为逆系统
$A \otimes_P^\mathbf{L} P/J^n$ 与逆系统 $A/I^n$ 严格 pro-同构
（由上面的讨论），并且有
$$
M \otimes_P^\mathbf{L} P/J^n = M \otimes_A^\mathbf{L}
(A \otimes_P^\mathbf{L} P/J^n)
$$
；见《代数进阶》，引理 \ref{more-algebra-lemma-derived-base-change}。

\medskip\noindent
(4) 中的一个标准例子是取 $P = \mathbf{Z}[x_1, \ldots, x_r]$，
取同态 $P \to A$，它将 $x_i$ 映到 $f_i$，并取 $J = (x_1, \ldots, x_r)$。
在此情形下，可以证明
$$
M \otimes_P^\mathbf{L} P/J^n =
M \otimes_{A[x_1, \ldots, x_r]}^\mathbf{L}
A[x_1, \ldots, x_r]/(x_1, \ldots, x_r)^n
$$
并约化到上面讨论过的情形之一（不过，此情形严格更容易，因为
$A[x_1, \ldots, x_r]/(x_1, \ldots, x_r)^n$ 的 Tor 维数至多为
$r$，这对所有 $n$ 成立，因而可以省去使用命题
\ref{local-cohomology-proposition-uniform-artin-rees} 的步骤）。
\cite[Proposition 3.5.1]{BS} 的证明讨论了此情形。
\end{remark}








\section{某种一致性，III}
\label{local-cohomology-section-uniform}

\noindent
本节中，我们固定一个诺特环 $A$ 和一个理想 $I \subset A$。
我们的目标是证明引理 \ref{local-cohomology-lemma-bound-two-term-complex}；在后面的章节中，
我们将用它解决一个提升问题，见
形式空间的代数化，引理
\ref{restricted-lemma-get-morphism-general-better}.

\medskip\noindent
在本节中，我们始终以
$$
p : X \to \Spec(A)
$$
表示 $\Spec(A)$ 沿理想 $I$ 的爆破。换言之，$X$ 是 $\text{Proj}$，
所用的 Rees 代数为 $\bigoplus_{n \geq 0} I^n$。我们还考虑纤维积
$$
\xymatrix{
Y \ar[r] \ar[d] & X \ar[d]^p \\
\Spec(A/I) \ar[r] & \Spec(A)
}
$$
于是 $Y$ 是该爆破的例外除子，从而是 $X$ 上的一个有效 Cartier 除子，且
$\mathcal{O}_X(-1) = \mathcal{O}_X(Y)$。由于取 $\text{Proj}$ 与基变换可交换，
我们有
$$
Y = \text{Proj}(\bigoplus\nolimits_{n \geq 0} I^n/I^{n + 1}) = \text{Proj}(S)
$$
其中 $S = \text{Gr}_I(A) = \bigoplus_{n \geq 0} I^n/I^{n + 1}$。

\medskip\noindent
我们以
$d = d(S) = d(\text{Gr}_I(A)) = d(\bigoplus_{n \geq 0} I^n/I^{n + 1})$
表示 $p$ 的各纤维维数的最大值
（并把它定义为 $0$，若 $X = \emptyset$）。
这是良定义的。事实上，
\begin{enumerate}
\item $d \leq t - 1$，如果 $I = (a_1, \ldots, a_t)$，因为此时
$X \subset \mathbf{P}^{t - 1}_A$；
\item $d$ 也是 $\text{Proj}(S) \to \Spec(S_0)$ 的各纤维维数的最大值，
只要 $Y$ 非空；且 $d = 0$，若 $Y = \emptyset$
（等价地，$S = 0$；等价地，$I = A$）。
\end{enumerate}
因此 $d$ 仅依赖于 $S = \text{Gr}_I(A)$ 的同构类。注意，
$H^i(X, \mathcal{F}) = 0$ 对每个凝聚 $\mathcal{O}_X$-模
$\mathcal{F}$ 以及 $i > d$ 都成立；这是由概形上同调，引理
\ref{coherent-lemma-higher-direct-images-zero-above-dimension-fibre} 和
\ref{coherent-lemma-quasi-coherence-higher-direct-images-application} 得到的。
当然，对 $Y$ 上的凝聚模也有同样的结论。

\medskip\noindent
我们以
$q = q(S) = q(\text{Gr}_I(A)) = q(\bigoplus_{n \geq 0} I^n/I^{n + 1})$
表示如下定义的整数。注意，代数
$S = \bigoplus_{n \geq 0} I^n/I^{n + 1}$
是由次数 $1$ 部分在次数 $0$ 部分上生成的诺特分次环。因此由
概形上同调，引理 \ref{coherent-lemma-coherent-on-proj} 和
\ref{coherent-lemma-recover-tail-graded-module}，我们可以把 $q(S)$ 定义为满足
如下性质的最小整数 $q(S) \geq 0$：对所有 $q \geq q(S)$，有
$H^i(Y, \mathcal{O}_Y(q)) = 0$，其中 $1 \leq i \leq d$，并且
$H^0(Y, \mathcal{O}_Y(q)) = I^q/I^{q + 1}$。
（若 $S = 0$，则 $q(S) = 0$。）

\medskip\noindent
对 $n \geq 1$，我们可以考虑有效 Cartier 除子 $nY$，并将其记为 $Y_n$。

\begin{lemma}
\label{local-cohomology-lemma-bound-q-and-d}
取如上的 $q_0 = q(S)$ 和 $d = d(S)$。则
\begin{enumerate}
\item 对 $n \geq 1$、$q \geq q_0$ 和 $i > 0$，有
$H^i(X, \mathcal{O}_{Y_n}(q)) = 0$,
\item 对 $n \geq 1$ 和 $q \geq q_0$，有
$H^0(X, \mathcal{O}_{Y_n}(q)) = I^q/I^{q + n}$,
\item 对 $q \geq q_0$ 和 $i > 0$，有
$H^i(X, \mathcal{O}_X(q)) = 0$,
\item 对 $q \geq q_0$，有
$H^0(X, \mathcal{O}_X(q)) = I^q$.
\end{enumerate}
\end{lemma}

\begin{proof}
若 $I = A$，则 $X$ 是仿射的，所有断言都是平凡的。因此我们可以并且确实假设
$I \not = A$。于是 $Y$ 和 $X$ 都是非空概形。

\medskip\noindent
我们对 $n$ 归纳证明 (1) 和 (2)。基础情形 $n = 1$ 正是我们对
$q_0$ 的定义，因为 $Y_1 = Y$。回忆 $\mathcal{O}_X(1) = \mathcal{O}_X(-Y)$。
因此有短正合列
$$
0 \to \mathcal{O}_{Y_n}(1) \to \mathcal{O}_{Y_{n + 1}} \to \mathcal{O}_Y \to 0
$$
故对 $i > 0$，得到
$$
H^i(X, \mathcal{O}_{Y_n}(q + 1)) \to
H^i(X, \mathcal{O}_{Y_{n + 1}}(q)) \to
H^i(X, \mathcal{O}_{Y}(q))
$$
由两端项的已知消没，即得中间项所需的消没。对 $i = 0$，得到交换图
$$
\xymatrix{
0 \ar[r] &
I^{q + 1}/I^{q + 1 + n} \ar[d] \ar[r] &
I^q/I^{q + 1 + n} \ar[d] \ar[r] &
I^q/I^{q + 1} \ar[d] \ar[r] &
0 \\
0 \ar[r] &
H^0(X, \mathcal{O}_{Y_n}(q + 1)) \ar[r] &
H^0(X, \mathcal{O}_{Y_{n + 1}}(q)) \ar[r] &
H^0(Y, \mathcal{O}_Y(q)) \ar[r] &
0
}
$$
当 $q \geq q_0$ 时两行均正合（对下行，只需注意长上同调正合列的下一项在
$q \geq q_0$ 时消没）。由于 $q \geq q_0$，左、右两个竖直箭头都是同构，
因而中间的竖直箭头也是同构。

\medskip\noindent
(3) 和 (4) 的证明类似，我们略去。事实上，可以利用形式函子定理从 (1) 和 (2)
推出 (3) 和 (4)（但这未免大材小用）。
\end{proof}

\noindent
引入如下记号：给定 $n \geq c \geq 0$，所谓
{\it 一个 $(A, n, c)$-模}，是指一个有限 $A$-模 $M$，它被 $I^n$ 零化，
且将它视为 $A/I^n$-模时，它是 $I^c/I^n$-投射的；见
更多代数，\ref{more-algebra-section-near-projective} 节。

\medskip\noindent
我们将采用如下记号滥用：给定 $A$-模 $M$，以 $p^*M$ 表示沿 $p$
拉回拟凝聚模 $\widetilde{M}$ 所得的拟凝聚模；这里后者是在 $\Spec(A)$
上与 $M$ 关联的模。例如，$\mathcal{O}_{Y_n} = p^*(A/I^n)$。
对短正合列 $0 \to K \to L \to M \to 0$，若其各项都是 $A$-模，
我们得到正合列
$$
p^*K \to p^*L \to p^*M \to 0
$$
因为 $\widetilde{\ }$ 是正合函子，而 $p^*$ 是右正合函子。

\begin{lemma}
\label{local-cohomology-lemma-almost-exactness}
设 $0 \to K \to L \to M \to 0$ 是 $A$-模的短正合列，其中 $K$ 和 $L$
都被 $I^n$ 零化，且 $M$ 是一个 $(A, n, c)$-模。则
$p^*K \to p^*L$ 的核在概形意义下支撑于 $Y_c$。
\end{lemma}

\begin{proof}
设 $\Spec(B) \subset X$ 是一个仿射开集。该正合列在 $\Spec(B)$ 上的限制
对应于 $B$-模序列
$$
K \otimes_A B \to L \otimes_A B \to M \otimes_A B \to 0
$$
它同构于序列
$$
K \otimes_{A/I^n} B/I^nB \to
L \otimes_{A/I^n} B/I^nB \to
M \otimes_{A/I^n} B/I^nB \to 0
$$
因此第一个映射的核是模 $\text{Tor}_1^{A/I^n}(M, B/I^nB)$ 的像。
回忆例外除子 $Y$ 由 $I\mathcal{O}_X$ 切出。因此，只需证明
$\text{Tor}_1^{A/I^n}(M, B/I^nB)$ 被 $I^c$ 零化。由于任意
$a \in I^c$ 在 $M$ 上的乘法都经由某个有限自由 $A/I^n$-模分解，这是显然的。
\end{proof}

\noindent
我们有典范映射 $\mathcal{O}_X \to \mathcal{O}_X(1)$，它恰在 $Y$ 上消没。
因此，对每个凝聚 $\mathcal{O}_X$-模 $\mathcal{F}$，总有典范映射
$\mathcal{F}(q) \to \mathcal{F}(q + n)$，其中 $q \in \mathbf{Z}$ 任意且
$n \geq 0$。

\begin{lemma}
\label{local-cohomology-lemma-annihilated}
设 $\mathcal{F}$ 是凝聚 $\mathcal{O}_X$-模。则 $\mathcal{F}$
在概形意义下支撑于 $Y_c$，当且仅当典范映射
$\mathcal{F} \to \mathcal{F}(c)$ 为零。
\end{lemma}

\begin{proof}
这是因为 $\mathcal{O}_X \to \mathcal{O}_X(1)$ 恰在 $Y$ 上消没。
\end{proof}

\begin{lemma}
\label{local-cohomology-lemma-vanishing-coh-almost-projective}
取如上的 $q_0 = q(S)$ 和 $d = d(S)$。设有整数
$n \geq c \geq 0$、一个 $(A, n, c)$-模 $M$、一个指标
$i \in \{0, 1, \ldots, d\}$ 以及一个整数 $q$。则分为如下情形：
\begin{enumerate}
\item 在 $i = d \geq 1$ 且 $q \geq q_0$ 的情形，有
$H^d(X, p^*M(q)) = 0$.
\item 在 $i = d - 1 \geq 1$ 且 $q \geq q_0$ 的情形，有
$H^{d - 1}(X, p^*M(q)) = 0$.
\item 在 $d - 1 > i > 0$ 且 $q \geq q_0 + (d - 1 - i)c$ 的情形，映射
$H^i(X, p^*M(q)) \to H^i(X, p^*M(q - (d - 1 - i)c))$
为零。
\item 在 $i = 0$、$d \in \{0, 1\}$ 且 $q \geq q_0$ 的情形，有满射
$$
I^qM \longrightarrow H^0(X, p^*M(q))
$$
\item 在 $i = 0$、$d > 1$ 且 $q \geq q_0 + (d - 1)c$ 的情形，映射
$$
H^0(X, p^*M(q)) \to H^0(X, p^*M(q - (d - 1)c))
$$
的像包含于典范映射
$I^{q - (d - 1)c}M \to H^0(X, p^*M(q - (d - 1)c))$.
的像中。
\end{enumerate}
\end{lemma}

\begin{proof}
设 $M$ 是一个 $(A, n, c)$-模。选择短正合列
$$
0 \to K \to (A/I^n)^{\oplus r} \to M \to 0
$$
下文将用到 $K$ 是一个 $(A, n, c)$-模这一事实，见更多代数，
引理 \ref{more-algebra-lemma-ses-near-projective}。考虑相应的正合列
$$
p^*K \to (\mathcal{O}_{Y_n})^{\oplus r} \to p^*M \to 0
$$
将其拆成两个短正合列
$$
0 \to \mathcal{F} \to p^*K \to \mathcal{G} \to 0
\quad\text{且}\quad
0 \to \mathcal{G} \to (\mathcal{O}_{Y_n})^{\oplus r} \to p^*M \to 0
$$
由引理 \ref{local-cohomology-lemma-almost-exactness}，凝聚模 $\mathcal{F}$ 在概形意义下
支撑于 $Y_c$。

\medskip\noindent
(1) 的证明。设 $d > 0$。我们要证明 $H^d(X, p^*M(q)) = 0$，其中
$q \geq q_0$。利用 $\mathcal{G}$ 的扭曲在次数 $> d$ 中的
上同调消没，以及与上面第二个短正合列相伴的长上同调正合列，只需证明
$H^d(X, \mathcal{O}_{Y_n}(q)) = 0$。这由引理
\ref{local-cohomology-lemma-bound-q-and-d} 成立。

\medskip\noindent
(2) 的证明。设 $d > 1$。我们要证明 $H^{d - 1}(X, p^*M(q)) = 0$，其中
$q \geq q_0$。如上一段论证可知，只需证明
$H^d(X, \mathcal{G}(q)) = 0$。利用第一个短正合列以及
$\mathcal{F}$ 的扭曲在次数 $> d$ 中的上同调消没，又只需证明
$H^d(X, p^*K(q))$ 为零。由 (1) 以及 $K$ 是一个 $(A, n, c)$-模
这一事实（见上）即得此结论。

\medskip\noindent
(3) 的证明。设 $0 < i < d - 1$，并假设该断言对 $i + 1$ 成立；
在 $i = d - 2$ 的情形，则使用断言 (2)。利用与上面第二个短正合列相伴的
长上同调正合列，我们得到单射
$$
H^i(X, p^*M(q - (d - 1 - i)c)) \subset
H^{i + 1}(X, \mathcal{G}(q - (d - 1 - i)c))
$$
因为 $q - (d - 1 - i)c \geq q_0$ 给出
$H^i(X, \mathcal{O}_{Y_n}(q - (d - 1 - i)c))$
的消没（见上）。因此，只需证明映射
$H^{i + 1}(X, \mathcal{G}(q)) \to H^{i + 1}(X, \mathcal{G}(q - (d - 1 - i)c))$
为零。为研究这一点，考虑如下正合列之间的映射：
$$
\xymatrix{
H^{i + 1}(X, p^*K(q)) \ar[r] \ar[d] &
H^{i + 1}(X, \mathcal{G}(q)) \ar[r] \ar[d] \ar@{..>}[ld] &
H^{i + 2}(X, \mathcal{F}(q)) \ar[d] \\
H^{i + 1}(X, p^*K(q - c)) \ar[r] \ar[d] &
H^{i + 1}(X, \mathcal{G}(q - c)) \ar[r] \ar[d] &
H^{i + 2}(X, \mathcal{F}(q - c)) \\
H^{i + 1}(X, p^*K(q - (d - 1 - i)c)) \ar[r] &
H^{i + 1}(X, \mathcal{G}(q - (d - 1 - i)c))
}
$$
由于 $\mathcal{F}$ 在概形意义下支撑于 $Y_c$，由引理
\ref{local-cohomology-lemma-annihilated} 可知，典范映射
$\mathcal{G}(q) \to \mathcal{G}(q - c)$ 经由
$p^*K(q - c)$ 分解。这给出图中的虚线箭头。（事实上，为了证明，只需注意
最右端的竖直箭头为零，即可得到作为集合映射的虚线箭头。）因此，只需证明
$H^{i + 1}(X, p^*K(q - c)) \to
H^{i + 1}(X, p^*K(q - (d - 1 - i)c))$
为零。若 $i = d - 2$，则由于 $q - c \geq q_0$ 且 $K$ 是一个
$(A, n, c)$-模，由 (2) 可知此箭头的源为零。若 $i < d - 2$，则由于
$K$ 是一个 $(A, n, c)$-模，由归纳假设可知该映射确实为零，因为
$q - c - (q - (d - 1 - i)c) = (d - 2 - i)c = (d - 1 - (i + 1))c$，并且
$q - c \geq q_0 + (d - 1 - (i + 1))c$。由此完成 (3) 的证明。

\medskip\noindent
(4) 的证明。设 $d \in \{0, 1\}$ 且 $q \geq q_0$。
于是第一个短正合列给出满射
$H^1(X, p^*K(q)) \to H^1(X, \mathcal{G}(q))$
而由情形 (1)，此箭头的源为零。因此，对所有 $q \in \mathbf{Z}$，映射
$$
H^0(X, (\mathcal{O}_{Y_n})^{\oplus r}(q))
\longrightarrow
H^0(X, p^*M(q))
$$
都是满射。当 $q \geq q_0$ 时，由引理 \ref{local-cohomology-lemma-bound-q-and-d}，
其源等于 $(I^q/I^{q + n})^{\oplus r}$，这就容易证明该断言。

\medskip\noindent
(5) 的证明。设 $d > 1$。如 (4) 的证明中那样论证可知，只需证明
$$
H^0(X, p^*M(q))
\longrightarrow
H^0(X, p^*M(q - (d - 1)c))
$$
的像包含于
$$
H^0(X, (\mathcal{O}_{Y_n})^{\oplus r}(q - (d - 1)c))
\longrightarrow
H^0(X, p^*M(q - (d - 1)c))
$$
的像中。为证明上述包含关系，只需证明：对
$\sigma \in H^0(X, p^*M(q))$，若其边界为
$\xi \in H^1(X, \mathcal{G}(q))$，则 $\xi$ 在
$H^1(X, \mathcal{G}(q - (d - 1)c))$ 中的像为零。这由 (3) 的证明中
完全相同的论证得到。
\end{proof}

\begin{remark}
\label{local-cohomology-remark-duals}
给定一对 $(M, n)$，其中 $n \geq 0$ 是整数，并有一个有限
$A/I^n$-模 $M$，我们置 $M^\vee = \Hom_{A/I^n}(M, A/I^n)$。
给定一对 $(\mathcal{F}, n)$，其中 $n$ 是整数，并有一个凝聚
$\mathcal{O}_{Y_n}$-模 $\mathcal{F}$，我们置
$$
\mathcal{F}^\vee =
\SheafHom_{\mathcal{O}_{Y_n}}(\mathcal{F}, \mathcal{O}_{Y_n})
$$
对如上的 $(M, n)$，有典范映射
$$
can : p^*(M^\vee) \longrightarrow (p^*M)^\vee
$$
事实上，若选择一个展示
$(A/I^n)^{\oplus s} \to (A/I^n)^{\oplus r} \to M \to 0$
则得到一个展示
$\mathcal{O}_{Y_n}^{\oplus s} \to \mathcal{O}_{Y_n}^{\oplus r} \to
p^*M \to 0$。取对偶后，得到正合列
$$
0 \to M^\vee \to (A/I^n)^{\oplus r} \to (A/I^n)^{\oplus s}
$$
以及
$$
0 \to (p^*M)^\vee \to
\mathcal{O}_{Y_n}^{\oplus r} \to
\mathcal{O}_{Y_n}^{\oplus s}
$$
沿 $p$ 拉回第一个序列，就得到所需的映射 $can$。这个映射的构造关于有限
$A/I^n$-模 $M$ 是函子性的。$can$ 的核与余核在概形意义下支撑于
$Y_c$，如果 $M$ 是一个 $(A, n, c)$-模。事实上，在这种情况下，对
$a \in I^c$，映射 $a : M \to M$ 经由某个有限自由 $A/I^n$-模分解，
而对后者 $can$ 是同构。因此 $a$ 零化 $can$ 的核与余核。
\end{remark}

\begin{lemma}
\label{local-cohomology-lemma-factor-hom}
取如上的 $q_0 = q(S)$ 和 $d = d(S)$。设 $M$ 是一个 $(A, n, c)$-模，
并设 $\varphi : M \to I^n/I^{2n}$ 是一个 $A$-线性映射。假设
$n \geq \max(q_0 + (1 + d)c, (2 + d)c)$；若 $d = 0$，还假设
$n \geq q_0 + 2c$。则复合映射
$$
M \xrightarrow{\varphi} I^n/I^{2n} \to
I^{n - (1 + d)c}/I^{2n - (1 + d)c}
$$
具有形式 $\sum a_i \psi_i$，其中 $a_i \in I^c$ 且
$\psi_i : M \to I^{n - (2 + d)c}/I^{2n - (2 + d)c}$。
\end{lemma}

\begin{proof}
$d > 1$ 的情形。由于我们有相容的映射系统
$p^*(I^q) \to \mathcal{O}_X(q)$，其中 $q \geq 0$，故有典范映射
$p^*(I^q/I^{q + \nu}) \to \mathcal{O}_{Y_\nu}(q)$，其中 $\nu \geq 0$。
利用这一点并拉回
$\varphi$，得到映射
$$
\chi : p^*M \longrightarrow \mathcal{O}_{Y_n}(n)
$$
使得复合映射
$M \to H^0(X, p^*M) \to H^0(X, \mathcal{O}_{Y_n}(n))$
等于给定同态 $\varphi$ 与映射
$I^n/I^{2n} \to H^0(X, \mathcal{O}_{Y_n}(n))$ 的复合。由于
$\mathcal{O}_{Y_n}(n)$ 在 $Y_n$ 上可逆，线性映射 $\chi$ 决定一个截面
$$
\sigma \in \Gamma(X, (p^*M)^\vee(n))
$$
这里采用注记 \ref{local-cohomology-remark-duals} 中的记号。注记 \ref{local-cohomology-remark-duals} 中的讨论说明，
$can : p^*(M^\vee) \to (p^*M)^\vee$ 的余核和核在概形意义下支撑于 $Y_c$。
由引理 \ref{local-cohomology-lemma-annihilated}，映射
$(p^*M)^\vee(n) \to (p^*M)^\vee(n - 2c)$ 经由
$p^*(M^\vee)(n - 2c)$ 分解；略去一个小细节。因此，$\sigma$ 在
$\Gamma(X, (p^*M)^\vee(n - 2c))$ 中的像来自一个元素
$$
\sigma' \in \Gamma(X, p^*(M^\vee)(n - 2c))
$$
由引理 \ref{local-cohomology-lemma-vanishing-coh-almost-projective} 的 (5)，再利用
$M^\vee$ 是一个 $(A, n, c)$-模这一事实（见更多代数，引理
\ref{more-algebra-lemma-dual-near-projective}），以及
$n \geq q_0 + (1 + d)c$ 从而 $n - 2c \geq q_0 + (d - 1)c$ 这一事实，
可知 $\sigma'$ 在 $H^0(X, p^*M^\vee(n - (1 + d)c))$ 中的像，是某个元素
$\tau$ 在 $I^{n - (1 + d)c}M^\vee$ 中的像。写成
$\tau = \sum a_i \tau_i$，其中 $\tau_i \in I^{n - (2 + d)c}M^\vee$；
这是有意义的，因为 $n - (2 + d)c \geq 0$。于是 $\tau_i$ 决定模同态
$\psi_i : M \to I^{n - (2 + d)c}/I^{2n - (2 + d)c}$
其中使用了求值映射 $M \otimes M^\vee \to A/I^n$。

\medskip\noindent
下面证明这确实可行\footnote{我们希望有读者能为此提出一个不那么笨拙的证明。}。
取 $z \in M$，并证明 $\varphi(z)$ 与 $\sum a_i \psi_i(z)$ 在
$I^{n - (1 + d)c}/I^{2n - (1 + d)c}$ 中有相同的像。首先，元素 $z$ 决定映射
$p^*z : \mathcal{O}_{Y_n} \to p^*M$；它与 $\chi$ 的复合等于映射
$\mathcal{O}_{Y_n} \to \mathcal{O}_{Y_n}(n)$，后者对应于 $\varphi(z)$，
所经的映射为 $I^n/I^{2n} \to \Gamma(\mathcal{O}_{Y_n}(n))$。
其次，$z$ 与 $p^*z$ 决定求值映射
$e_z : M^\vee \to A/I^n$ 和 $e_{p^*z} : (p^*M)^\vee \to \mathcal{O}_{Y_n}$。
由于 $\chi(p^*z)$ 是对应于 $\varphi(z)$ 的截面，可知
$e_{p^*z}(\sigma)$ 是对应于 $\varphi(z)$ 的截面。这里以及下文都滥用记号：
对于模的映射 $a : \mathcal{F} \to \mathcal{G}$，若它定义在 $X$ 上，也以
$a : \mathcal{F}(t) \to \mathcal{F}(t)$ 表示相应的扭曲模映射。图
$$
\xymatrix{
p^*(M^\vee) \ar[d]_{can} \ar[r]_{p^*e_z} & \mathcal{O}_{Y_n} \ar@{=}[d] \\
(p^*M)^\vee \ar[r]^{e_{p^*z}} & \mathcal{O}_{Y_n}
}
$$
由于构造 $can$ 的函子性而交换。因此，$(p^*e_z)(\sigma')$ 作为
$\Gamma(Y_n, \mathcal{O}_{Y_n}(n - 2c))$ 中的元素，是对应于 $\varphi(z)$ 在
$I^{n - 2c}/I^{2n - 2c}$ 中之像的截面。下一步，$\sigma'$ 映到
$\sum a_i \tau_i$ 在 $H^0(X, p^*M^\vee(n - (1 + d)c))$ 中的像。这说明
$(p^*e_z)(\sum a_i \tau_i) = \sum a_i p^*e_z(\tau_i)$ 在
$\Gamma(Y_n, \mathcal{O}_{Y_n}(n - (1 + d)c))$ 中，是对应于
$\varphi(z)$ 在 $I^{n - (1 + d)c}/I^{2n - (1 + d)c}$ 中之像的截面。
回忆 $\psi_i$ 是利用一个求值映射从 $\tau_i$ 定义的。因此，若以
$$
\chi_i : p^*M \longrightarrow \mathcal{O}_{Y_n}(n - (2 + d)c)
$$
表示从 $\psi_i$ 得到的映射，则由与上面相同的论证可知，对应于
$\psi_i(z)$ 的截面是
$\chi_i(p^*z) = e_{p^*z}(\chi_i) = p^*e_z(\tau_i)$。因此可知，
$\varphi(z)$ 在 $\Gamma(Y_n, \mathcal{O}_{Y_n}(n - (1 + d)c))$ 中的像，
等于 $\sum a_i\psi_i(z)$ 的像。由于 $n - (1 + d)c \geq q_0$，由引理
\ref{local-cohomology-lemma-bound-q-and-d} 有
$\Gamma(Y_n, \mathcal{O}_{Y_n}(n - (1 + d)c)) =
I^{n - (1 + d)c}/I^{2n - (1 + d)c}$，故所需的相容性成立。

\medskip\noindent
$d = 1$ 的情形。像上面那样论证，得到
$$
\chi : p^*M \longrightarrow \mathcal{O}_{Y_n}(n),\quad
\sigma \in \Gamma(X, (p^*M)^\vee(n)),\quad
\sigma' \in \Gamma(X, p^*(M^\vee)(n - 2c)),
$$
然后利用引理 \ref{local-cohomology-lemma-vanishing-coh-almost-projective} 的 (4)，可知
$\sigma'$ 是某个元素 $\tau \in I^{n - 2c}M^\vee$ 的像。其余论证相同。

\medskip\noindent
$d = 0$ 的情形。论证与 $d = 1$ 的情形完全相同。
\end{proof}

\begin{lemma}
\label{local-cohomology-lemma-bound-two-term-complex}
取如上的 $d = d(S)$ 和 $q_0 = q(S)$。则
\begin{enumerate}
\item 对满足
$n \geq c \geq 0$ 且
$n \geq \max(q_0 + (1 + d)c, (2 + d)c)$,
的整数；
\item 对 $K$，它是 $D(A/I^n)$ 中的对象，且 $H^i(K) = 0$ 当
$i \not = -1, 0$，而 $H^i(K)$ 在 $i = -1, 0$ 时有限，并且满足
$\Ext^1_{A/I^c}(K, N)$ 被 $I^c$ 零化，对所有有限 $A/I^n$-模 $N$，
\end{enumerate}
映射
$$
\Ext^1_{A/I^n}(K, I^n/I^{2n})
\longrightarrow
\Ext^1_{A/I^n}(K, I^{n - (1 + d)c}/I^{2n - 2(1 + d)c})
$$
为零。
\end{lemma}

\begin{proof}
$d > 0$ 的情形。令 $K^{-1} \to K^0$ 是一个表示 $K$ 的复形；它是按照
更多代数，引理 \ref{more-algebra-lemma-ext-1-annihilated-definite} 的 (5)，
相对于理想 $I^c/I^n$ 选取的；该理想位于环 $A/I^n$ 中。特别地，$K^{-1}$ 是
$I^c/I^n$-投射的，因为 $I^c/I^n$ 中元素的乘法甚至经由 $K^0$ 分解。
由更多代数，引理 \ref{more-algebra-lemma-map-out-of-almost-free} 的 (1)，有
$$
\Ext^1_{A/I^n}(K, I^n/I^{2n}) =
\Coker(\Hom_{A/I^n}(K^0, I^n/I^{2n}) \to \Hom_{A/I^n}(K^{-1}, I^n/I^{2n}))
$$
其他 Ext 群也同样。因此，任何类 $\xi$，若它属于
$\Ext^1_{A/I^n}(K, I^n/I^{2n})$
都来自一个元素 $\varphi \in \Hom_{A/I^n}(K^{-1}, I^n/I^{2n})$。
以 $\varphi'$ 表示 $\varphi$ 在
$\Hom_{A/I^n}(K^{-1}, I^{n - (1 + d)c}/I^{2n - (1 + d)c})$.
中的像。由引理 \ref{local-cohomology-lemma-factor-hom}，可写成
$\varphi' = \sum a_i \psi_i$，其中 $a_i \in I^c$ 且
$\psi_i \in \Hom_{A/I^n}(M, I^{n - (2 + d)c}/I^{2n - (2 + d)c})$。
选择 $h_i : K^0 \to K^{-1}$，使得
$a_i \text{id}_{K^{-1}} = h_i \circ d_K^{-1}$。置
$\psi = \sum \psi_i \circ h_i : K^0 \to I^{n - (2 + d)c}/I^{2n - (2 + d)c}$.
则 $\varphi' = \psi \circ \text{d}_K^{-1}$，从而 $\xi$ 已经在
$\Ext^1_{A/I^n}(K, I^{n - (1 + d)c}/I^{2n - (1 + d)c})$
中映为零，因而更在
$\Ext^1_{A/I^n}(K, I^{n - (1 + d)c}/I^{2n - 2(1 + d)c})$.
中映为零。

\medskip\noindent
$d = 0$ 的情形\footnote{对 $d > 0$ 给出的论证也适用，但所得结果稍弱。}。
取如上的 $\xi$ 和 $\varphi$。考虑图
$$
\xymatrix{
K^0 \\
K^{-1} \ar[u] \ar[r]^\varphi & I^n/I^{2n} \ar[r] & I^{n - c}/I^{2n - c}
}
$$
拉回到 $X$，并利用映射
$p^*(I^n/I^{2n}) \to \mathcal{O}_{Y_n}(n)$
得到实线图
$$
\xymatrix{
p^*K^0 \ar@{..>}[rrd] \\
p^*K^{-1} \ar[u] \ar[r] &
\mathcal{O}_{Y_n}(n) \ar[r] &
\mathcal{O}_{Y_n}(n - c)
}
$$
可以覆盖 $X$，所用仿射开集 $U = \Spec(B)$ 具有如下性质：存在一个 $a \in I$，
满足如下性质：$IB = aB$，且 $a$ 是 $B$ 上的非零因子。事实上，可以用
仿射爆破代数的谱覆盖 $X$；见除子，引理
\ref{divisors-lemma-blowing-up-affine}。映射
$\mathcal{O}_{Y_n}(n) \to \mathcal{O}_{Y_n}(n - c)$ 在 $U$ 上的限制，
同构于拟凝聚 $\mathcal{O}_U$-模的映射；后者对应于 $B$-模映射
$a^c : B/a^nB \to B/a^nB$。由于 $a^c : K^{-1} \to K^{-1}$ 经由
$K^0$ 分解，可知虚线箭头在 $U$ 上存在。换言之，在 $X$ 上局部地，
我们可以找到虚线箭头！现在，与该图适配的虚线箭头所成的层，是在下述层作用下
主齐性的：
$$
\mathcal{F} = \SheafHom_{\mathcal{O}_X}(
\Coker(p^*K^{-1} \to p^*K^0), \mathcal{O}_{Y_n}(n - c))
$$
这是一个凝聚 $\mathcal{O}_X$-模。因此，寻找虚线箭头的障碍是
$H^1(X, \mathcal{F})$ 中的一个元素。由于 $1 > d = 0$，该上同调群为零；
见 $d = d(S)$ 的定义之后的讨论。这就证明了可以找到虚线箭头
$\psi : p^*K^0 \to \mathcal{O}_{Y_n}(n - c)$
与该图适配。由于 $n - c \geq q_0$，可知 $\psi$ 诱导映射
$K^0 \to I^{n - c}/I^{2n - c}$。追图得到
$\varphi' = \psi \circ \text{d}_K^{-1}$，证明如前完成。
\end{proof}
% Generated mechanically from the frozen algebraization.tex target.
% Do not edit this generated file; edit the bound locale source instead.

% OK, start here.
%

\setcounter{chapter}{51}
\chapter{代数几何与形式几何}



\phantomsection
\label{algebraization-section-phantom}


\section{引言}
\label{algebraization-section-introduction}

\noindent
本章继续研究形式代数几何，尤其研究形式对象是否为某个代数对象的完备化这一问题。
一部奠基性参考文献是 \cite{SGA2}。下面列出我们已经在 Stacks 项目中讨论过的结果：
\begin{enumerate}
\item 形式函数定理，见概形上同调，\ref{coherent-section-theorem-formal-functions} 节。
\item 凝聚形式模，见概形上同调，\ref{coherent-section-coherent-formal} 节。
\item Grothendieck 存在性定理，见概形上同调，
\ref{coherent-section-existence}、\ref{coherent-section-existence-proper} 和
\ref{coherent-section-existence-proper-support}.
\item Grothendieck 代数化定理，见概形上同调，\ref{coherent-section-algebraization} 节。
\item 更一般的 Grothendieck 存在性定理，见更多平坦性，
\ref{flat-section-existence} 和
\ref{flat-section-existence-derived}.
\end{enumerate}
下面概述本章内容。

\medskip\noindent
设 $X$ 是一个概形，$\mathcal{I} \subset \mathcal{O}_X$ 是有限型拟凝聚理想层。
本章的许多问题都涉及逆系统 $(\mathcal{F}_n)$，其各项是拟凝聚
$\mathcal{O}_X$-模，并且它满足
$\mathcal{F}_n = \mathcal{F}_{n + 1}/\mathcal{I}^n\mathcal{F}_{n + 1}$。
一个重要的特殊情形是：$X$ 是诺特环 $A$ 上的概形，并且
$\mathcal{I} = I \mathcal{O}_X$，其中 $I \subset A$ 是某个理想。在上同调，
\ref{cohomology-section-inverse-systems}、
\ref{cohomology-section-inverse-systems-bis} 和
\ref{cohomology-section-inverse-systems-tri} 节中，我们给出了一些一般结果。
在本章中，\ref{algebraization-section-ML-degree-zero} 和
\ref{algebraization-section-formal-functions-principal} 节包含概形与拟凝聚模所特有的结果。
在 \ref{algebraization-section-formal-sections-cd-one} 节中，我们证明：
$\lim H^p(X, \mathcal{F}_n)$ 上的极限拓扑是 $I$-进拓扑，如果
$\text{cd}(A, I) = 1$。本章的一个主题是说明：在主理想情形 $I = (f)$ 下证明的结果，
只假设 $\text{cd}(A, I) = 1$ 时仍然成立。

\medskip\noindent
在 \ref{algebraization-section-derived-completion} 节中，我们讨论环化位点
$(\mathcal{C}, \mathcal{O})$ 上的模关于有限型理想层 $\mathcal{I}$ 的导出完备化。
这一节自然延续更多代数，\ref{more-algebra-section-derived-completion} 节中
所述的交换代数导出完备化理论。第一个主要结果是导出完备化存在。
第二个主要结果是：对环化位点的态射 $f$，导出完备化与导出推出可交换：
$$
(Rf_*K)^\wedge = Rf_*(K^\wedge)
$$
前提是上方的理想层局部地由来自下方理想的截面生成；见引理
\ref{algebraization-lemma-pushforward-commutes-with-derived-completion}。需要强调的是，若所讨论的
理想都是全局有限生成的，则这两个主要结果都非常初等；本章中这一理论的所有应用
都属于这种情形。上面展示的等式是形式函数定理的“正确”版本，见
\ref{algebraization-section-formal-functions} 节中的讨论。

\medskip\noindent
设 $A$ 是诺特环，$I, J$ 是 $A$ 的两个理想。设 $M$ 是有限 $A$-模。
本章接下来的主题是映射
$$
R\Gamma_J(M) \longrightarrow R\Gamma_J(M)^\wedge
$$
它从 $M$ 的局部上同调映入其导出 $I$-进完备化。事实证明，若施加合适的深度条件，
这个映射在一定次数范围内的上同调上成为同构。在
\ref{algebraization-section-algebraization-sections-general} 节中，我们基本上在刚才所述的一般性下工作。
在 \ref{algebraization-section-algebraization-punctured} 节中，假设 $A$ 是局部环，且
$J = \mathfrak m$ 是极大理想。我们建议读者先读这一节，再读本章这一部分的另外两节。
最后，在 \ref{algebraization-section-bootstrap} 节中，我们从局部情形逐步推导，在一般情形中得到更强的结果。

\medskip\noindent
在本章的下一部分，我们利用局部上同调完备化的结果，给出关于凝聚模完备化之
上同调的一组并非穷尽的结果。更准确地说，设 $A$ 是诺特环，$I \subset A$
是理想，$U \subset \Spec(A)$ 是开子概形。若 $\mathcal{F}$ 是凝聚
$\mathcal{O}_U$-模，则可以考虑映射
$$
H^i(U, \mathcal{F}) \longrightarrow \lim H^i(U, \mathcal{F}/I^n\mathcal{F})
$$
并问在一定次数范围内是否得到同构。在 \ref{algebraization-section-algebraization-sections} 节中，
我们处理一些 $U$ 是局部环穿孔谱的例子。在
\ref{algebraization-section-algebraization-sections-coherent} 节中，我们讨论一般情形。
在 \ref{algebraization-section-connected} 节中，我们把所得的一些结果应用于代数几何中的连通性问题。

\medskip\noindent
本章余下各节专门讨论凝聚形式模的代数化。换言之，给定如上的凝聚模逆系统
$(\mathcal{F}_n)$，它位于 $U$ 上并满足
$\mathcal{F}_n = \mathcal{F}_{n + 1}/I^n\mathcal{F}_{n + 1}$
我们要问：是否存在一个凝聚 $\mathcal{O}_U$-模 $\mathcal{F}$，使得
$\mathcal{F}_n = \mathcal{F}/I^n\mathcal{F}$
对所有 $n$ 都成立。我们建议读者阅读 \ref{algebraization-section-algebraization-modules} 节，
其中给出了这一问题的精确表述、一个有用的一般结果
（引理 \ref{algebraization-lemma-when-done}）以及一个非平凡的应用
（引理 \ref{algebraization-lemma-algebraization-principal-variant}）。要证明实质上超出这一情形的结果，
还需要发展相当多的理论。本章所得此类结果中最强的那些，见
\ref{algebraization-section-algebraization-modules-conclusion} 节。



\section{形式截面，I}
\label{algebraization-section-ML-degree-zero}

\noindent
我们建议先阅读《上同调》，第 \ref{cohomology-section-inverse-systems} 节。

\begin{lemma}
\label{algebraization-lemma-properties-system}
设 $X$ 是概形。设 $\mathcal{I} \subset \mathcal{O}_X$
是拟凝聚理想层。设
$$
\ldots \to \mathcal{F}_3 \to \mathcal{F}_2 \to \mathcal{F}_1
$$
是拟凝聚 $\mathcal{O}_X$-模的一个逆系统，并且
$\mathcal{F}_n = \mathcal{F}_{n + 1}/\mathcal{I}^n\mathcal{F}_{n + 1}$.
置 $\mathcal{F} = \lim \mathcal{F}_n$。则
\begin{enumerate}
\item $\mathcal{F} = R\lim \mathcal{F}_n$，
\item 对任意仿射开集 $U \subset X$，有
$H^p(U, \mathcal{F}) = 0$，其中 $p > 0$，并且
\item 对每个 $p$ 都有短正合列
$0 \to R^1\lim H^{p - 1}(X, \mathcal{F}_n) \to
H^p(X, \mathcal{F}) \to \lim H^p(X, \mathcal{F}_n) \to 0$。
\end{enumerate}
此外，若 $\mathcal{I}$ 是有限型的，则
\begin{enumerate}
\item[(4)]
$\mathcal{F}_n = \mathcal{F}/\mathcal{I}^n\mathcal{F}$，并且
\item[(5)]
$\mathcal{I}^n \mathcal{F} = \lim_{m \geq n} \mathcal{I}^n\mathcal{F}_m$。
\end{enumerate}
\end{lemma}

\begin{proof}
(1)、(2) 与 (3) 是具有满射过渡映射的拟凝聚模逆系统的一般事实；见
《概形的导出范畴》，引理 \ref{perfect-lemma-Rlim-quasi-coherent}
与《上同调》，引理 \ref{cohomology-lemma-RGamma-commutes-with-Rlim}。
其次，设 $\mathcal{I}$ 是有限型的。
设 $U \subset X$ 是仿射开集。写 $U = \Spec(A)$，并设 $\mathcal{I}|_U$
对应于 $I \subset A$。注意 $I$ 是有限生成理想。
利用拟凝聚 $\mathcal{O}_U$-模与 $A$-模之间的范畴等价
（《概形》，引理 \ref{schemes-lemma-equivalence-quasi-coherent}），
可知 $M_n = \mathcal{F}_n(U)$ 是 $A$-模的一个逆系统，且
$M_n = M_{n + 1}/I^nM_{n + 1}$。因此
$$
M = \mathcal{F}(U) = \lim \mathcal{F}_n(U) = \lim M_n
$$
是一个 $I$-进完备模，并且由《代数》，引理
\ref{algebra-lemma-limit-complete}，有 $M/I^nM = M_n$。这证明了 (4)。
(5) 转化为命题
$\lim_{m \geq n} I^nM/I^mM = I^nM$。
由于 $I^mM = I^{m - n} \cdot I^nM$，这正是说
$I^mM$ 是 $I$-进完备的。这由《代数》，引理
\ref{algebra-lemma-hathat-finitely-generated}
以及 $M$ 的完备性得出。
\end{proof}





\section{形式截面，II}
\label{algebraization-section-formal-functions-principal}

\noindent
我们建议先阅读《上同调》，第
\ref{cohomology-section-inverse-systems-bis} 节与第
\ref{cohomology-section-inverse-systems-tri} 节。

\begin{lemma}
\label{algebraization-lemma-equivalent-f-good}
设 $X$ 是概形。设 $f \in \Gamma(X, \mathcal{O}_X)$。设
$$
\ldots \to \mathcal{F}_3 \to \mathcal{F}_2 \to \mathcal{F}_1
$$
是拟凝聚 $\mathcal{O}_X$-模的一个逆系统。下列条件等价：
\begin{enumerate}
\item 对所有 $n \geq 1$，映射
$f : \mathcal{F}_{n + 1} \to \mathcal{F}_{n + 1}$ 通过
$\mathcal{F}_{n + 1} \to \mathcal{F}_n$ 分解，并给出短正合列
$0 \to \mathcal{F}_n \to \mathcal{F}_{n + 1} \to \mathcal{F}_1 \to 0$；
\item 对所有 $n \geq 1$，映射
$f^n : \mathcal{F}_{n + 1} \to \mathcal{F}_{n + 1}$
通过 $\mathcal{F}_{n + 1} \to \mathcal{F}_1$ 分解，并给出短正合列
$0 \to \mathcal{F}_1 \to \mathcal{F}_{n + 1} \to \mathcal{F}_n \to 0$；
\item 存在一个 $\mathcal{O}_X$-模 $\mathcal{G}$，它是
$f$-可除的，并且 $\mathcal{F}_n = \mathcal{G}[f^n]$；
\item 存在一个 $\mathcal{O}_X$-模 $\mathcal{F}$，它无 $f$-挠，并且
$\mathcal{F}_n = \mathcal{F}/f^n\mathcal{F}$。
\end{enumerate}
\end{lemma}

\begin{proof}
(1)、(2)、(3) 的等价性以及蕴涵 (4) $\Rightarrow$ (1)，已在《上同调》，
引理 \ref{cohomology-lemma-equivalent-f-good} 中证明。
设 (1) 成立。置
$\mathcal{F} = \lim \mathcal{F}_n$。由引理 \ref{algebraization-lemma-properties-system}
的 (4)，有 $\mathcal{F}_n = \mathcal{F}/f^n\mathcal{F}$。
设 $U \subset X$ 是开集，且
$s = (s_n) \in \mathcal{F}(U) = \lim \mathcal{F}_n(U)$。
选取 $n \geq 1$。若 $fs = 0$，则由条件 (1)，$s_{n + 1}$ 属于
$\mathcal{F}_{n + 1} \to \mathcal{F}_n$ 的核。
故 $s_n = 0$。由于 $n$ 是任意的，可知 $s = 0$。
因此 $\mathcal{F}$ 无 $f$-挠。
\end{proof}

\begin{lemma}
\label{algebraization-lemma-formal-functions-principal}
\begin{reference}
这是 \cite[Lemma 1.6]{Bhatt-local} 的略加改进版本。
\end{reference}
设 $A$ 是环且 $f \in A$。设 $X$ 是 $A$ 上的概形。
设 $\mathcal{F}$ 是拟凝聚 $\mathcal{O}_X$-模。
假设 $\mathcal{F}[f^n] = \Ker(f^n : \mathcal{F} \to \mathcal{F})$
稳定。则
$$
R\Gamma(X, \lim \mathcal{F}/f^n\mathcal{F}) =
R\Gamma(X, \mathcal{F})^\wedge
$$
其中右边表示关于理想 $(f) \subset A$ 的导出完备化。因此，对
$p \in \mathbf{Z}$，得到交换图
$$
\xymatrix{
& 0 & 0 \\
0 \ar[r] &
\widehat{H^p(X, \mathcal{F})} \ar[r] \ar[u] &
\lim H^p(X, \mathcal{F}/f^n\mathcal{F}) \ar[r] \ar[u] &
T_f(H^{p + 1}(X, \mathcal{F})) \ar[r] &
0 \\
0 \ar[r] &
H^0(H^p(X, \mathcal{F})^\wedge) \ar[r] \ar[u] &
H^p(X, \lim \mathcal{F}/f^n\mathcal{F}) \ar[r] \ar[u] &
T_f(H^{p + 1}(X, \mathcal{F})) \ar[r] \ar@{=}[u] &
0 \\
&
R^1\lim H^p(X, \mathcal{F})[f^n] \ar[u] \ar[r]^\cong &
R^1\lim H^{p - 1}(X, \mathcal{F}/f^n\mathcal{F}) \ar[u] \\
& 0 \ar[u] & 0 \ar[u]
}
$$
其各行各列均正合，其中
$\widehat{H^p(X, \mathcal{F})} =
\lim H^p(X, \mathcal{F})/f^n H^p(X, \mathcal{F})$
是通常的 $f$-进完备化，而 $T_f(-)$ 表示 $f$-进 Tate 模，
如《代数进阶》，例
\ref{more-algebra-example-spectral-sequence-principal} 中所定义。
\end{lemma}

\begin{proof}
由引理 \ref{algebraization-lemma-properties-system}，有
$\lim \mathcal{F}/f^n\mathcal{F} = R\lim \mathcal{F}/f^n \mathcal{F}$.
其余结论由《上同调》，例
\ref{cohomology-example-formal-functions-principal} 得出。
\end{proof}









\section{形式截面，III}
\label{algebraization-section-formal-sections-cd-one}

\noindent
本节证明引理 \ref{algebraization-lemma-topology-I-adic}。在诺特概形与凝聚模的情形下，
当不假设理想 $I$ 是主理想，而是假设
$\text{cd}(A, I) = 1$ 时，该引理是《上同调》，引理
\ref{cohomology-lemma-topology-I-adic-f} 的类似结果。

\begin{lemma}
\label{algebraization-lemma-cd-one}
设 $I = (f_1, \ldots, f_r)$ 是诺特环 $A$ 的理想。
若 $\text{cd}(A, I) = 1$，则存在 $c \geq 1$ 与映射
$\varphi_j : I^c \to A$，使得 $\sum f_j \varphi_j : I^c \to I$
是包含映射。
\end{lemma}

\begin{proof}
由于 $\text{cd}(A, I) = 1$，补集 $U = \Spec(A) \setminus V(I)$
是仿射的（《局部上同调》，引理 \ref{local-cohomology-lemma-cd-is-one}）。
写 $U = \Spec(B)$。于是 $IB = B$，并且可以写成
$1 = \sum_{j = 1, \ldots, r} f_j b_j$，
其中 $b_j \in B$。由《概形上同调》，引理
\ref{coherent-lemma-homs-over-open}，可将 $b_j$ 表示为映射
$\varphi_j : I^c \to A$，其中 $c \geq 0$。
再由同一引理，$\sum f_j \varphi_j : I^c \to I \subset A$
是典范嵌入；这里必要时把 $c$ 换成更大的整数。
\end{proof}

\begin{lemma}
\label{algebraization-lemma-cd-one-extend}
设 $I = (f_1, \ldots, f_r)$ 是诺特环 $A$ 的理想，且
$\text{cd}(A, I) = 1$。设 $c \geq 1$ 以及 $\varphi_j : I^c \to A$，
$j = 1, \ldots, r$，如引理 \ref{algebraization-lemma-cd-one} 中所述。
则存在唯一的分次 $A$-代数映射
$$
\Phi : \bigoplus\nolimits_{n \geq 0} I^{nc} \to A[T_1, \ldots, T_r]
$$
使得 $\Phi(g) = \sum \varphi_j(g) T_j$，其中 $g \in I^c$。
此外，$\Phi$ 与映射
$A[T_1, \ldots, T_r] \to \bigoplus_{n \geq 0} I^n$,
$T_j \mapsto f_j$ 的复合是包含映射
$\bigoplus_{n \geq 0} I^{nc} \to \bigoplus_{n \geq 0} I^n$.
\end{lemma}

\begin{proof}
对每个 $j$ 与 $m \geq c$，$\varphi_j$ 在
$I^m$ 上的限制是映射 $\varphi_j : I^m \to I^{m - c}$。
给定 $j_1, \ldots, j_n \in \{1, \ldots, r\}$，我们断言复合
$$
\varphi_{j_1} \ldots \varphi_{j_n} :
I^{nc} \to I^{(n - 1)c} \to \ldots \to I^c \to A
$$
与指标 $j_1, \ldots, j_n$ 的次序无关。事实上，若
$g = g_1 \ldots g_n$ 且 $g_i \in I^c$，则可见
$$
(\varphi_{j_1} \ldots \varphi_{j_n})(g) =
\varphi_{j_1}(g_1) \ldots \varphi_{j_n}(g_n)
$$
与次序无关，因为 $A$ 中的乘法可交换。因此，可以如下定义 $\Phi$：
把 $g \in I^{nc}$ 映到
$$
\Phi(g) = \sum\nolimits_{e_1 + \ldots + e_r = n}
(\varphi_1^{e_1} \circ \ldots \circ \varphi_r^{e_r})(g)
T_1^{e_1} \ldots T_r^{e_r}
$$
直接验证可知，这是具有所需性质的分次 $A$-代数同态。
唯一性与最后一个性质都是立即可见的。这就证明了引理。
\end{proof}

\begin{lemma}
\label{algebraization-lemma-cd-one-extend-to-module}
设 $I = (f_1, \ldots, f_r)$ 是诺特环 $A$ 的理想，且
$\text{cd}(A, I) = 1$。设 $c \geq 1$ 以及 $\varphi_j : I^c \to A$，
$j = 1, \ldots, r$，如引理 \ref{algebraization-lemma-cd-one} 中所述。
设 $A \to B$ 是环映射，其中 $B$ 诺特，并设 $N$ 是有限
$B$-模。必要时增大 $c$ 并相应调整 $\varphi_j$ 之后，
存在唯一的分次 $B$-模映射
$$
\Phi_N : \bigoplus\nolimits_{n \geq 0} I^{nc}N \to N[T_1, \ldots, T_r]
$$
使得 $\Phi_N(g x) = \Phi(g) x$，其中 $g \in I^{nc}$ 且 $x \in N$，
这里 $\Phi$ 如引理
\ref{algebraization-lemma-cd-one-extend} 中所述。$\Phi_N$ 与映射
$N[T_1, \ldots, T_r] \to \bigoplus_{n \geq 0} I^nN$,
$T_j \mapsto f_j$ 的复合是包含映射
$\bigoplus_{n \geq 0} I^{nc}N \to \bigoplus_{n \geq 0} I^nN$.
\end{lemma}

\begin{proof}
由该公式以及引理 \ref{algebraization-lemma-cd-one-extend} 中 $\Phi$ 的唯一性，
唯一性是清楚的。考虑诺特 $A$-代数
$B' = B \oplus N$，其中 $N$ 是平方为零的理想。为了证明
$\Phi_N$ 存在，由引理 \ref{algebraization-lemma-cd-one}，只需证明 $\varphi_j$
可延拓为映射 $\varphi'_j : I^cB' \to B'$；这里必要时把 $c$
增大到某个 $c'$（并把 $\varphi_j$ 换成包含映射
$I^{c'} \to I^c$ 与 $\varphi_j$ 的复合）。回忆，$\varphi_j$ 对应于截面
$$
h_j \in \Gamma(\Spec(A) \setminus V(I), \mathcal{O}_{\Spec(A)})
$$
；见《概形上同调》，引理 \ref{coherent-lemma-homs-over-open}。
（事实上，在引理 \ref{algebraization-lemma-cd-one} 的证明中，我们正是这样选取
$\varphi_j$ 的。）利用同一引理，把拉回
$$
h'_j \in \Gamma(\Spec(B') \setminus V(IB'), \mathcal{O}_{\Spec(B')})
$$
（即 $h_j$ 的拉回）表示为 $B'$-线性映射
$\varphi'_j : I^{c'}B' \to B'$，其中 $c' \geq c$。
再次应用该引理可知，它将与 $\varphi_j$ 相符，只要 $c'$
充分大：事实上，只需在 $I^{c'}$
的一组有限生成元上检验二者相符。略去这个小细节。
\end{proof}

\begin{lemma}
\label{algebraization-lemma-cd-is-one-for-system}
设 $I = (f_1, \ldots, f_r)$ 是诺特环 $A$ 的理想，且
$\text{cd}(A, I) = 1$。设 $c \geq 1$ 以及 $\varphi_j : I^c \to A$，
$j = 1, \ldots, r$，如引理 \ref{algebraization-lemma-cd-one} 中所述。
设 $X$ 是 $\Spec(A)$ 上的诺特概形。设
$$
\ldots \to \mathcal{F}_3 \to \mathcal{F}_2 \to \mathcal{F}_1
$$
是凝聚 $\mathcal{O}_X$-模的一个逆系统，并且
$\mathcal{F}_n = \mathcal{F}_{n + 1}/I^n\mathcal{F}_{n + 1}$。
置 $\mathcal{F} = \lim \mathcal{F}_n$。
必要时增大 $c$ 并相应调整 $\varphi_j$ 之后，
存在唯一的分次 $\mathcal{O}_X$-模映射
$$
\Phi_\mathcal{F} :
\bigoplus\nolimits_{n \geq 0} I^{nc}\mathcal{F}
\longrightarrow
\mathcal{F}[T_1, \ldots, T_r]
$$
使得 $\Phi_\mathcal{F}(g s) = \Phi(g) s$，其中 $g \in I^{nc}$，
$s$ 是 $\mathcal{F}$ 的局部截面，而 $\Phi$ 如引理
\ref{algebraization-lemma-cd-one-extend} 中所述。$\Phi_\mathcal{F}$ 与映射
$\mathcal{F}[T_1, \ldots, T_r] \to \bigoplus_{n \geq 0} I^n\mathcal{F}$,
$T_j \mapsto f_j$
的复合是典范包含
$\bigoplus_{n \geq 0} I^{nc}\mathcal{F} \to
\bigoplus_{n \geq 0} I^n\mathcal{F}$.
\end{lemma}

\begin{proof}
由 $\mathcal{O}_X$-线性以及要求
$\Phi_\mathcal{F}(g s) = \Phi(g) s$ 对 $g \in I^{nc}$ 和
$s$ 是 $\mathcal{F}$ 的局部截面这一条件成立，唯一性立即可见。
因此可设 $X = \Spec(B)$ 是仿射的。注意 $(\mathcal{F}_n)$ 是范畴
$\textit{Coh}(X, I\mathcal{O}_X)$ 的对象；该范畴在《概形上同调》，
第 \ref{coherent-section-coherent-formal} 节中引入。
设 $B' = B^\wedge$ 是相对于 $I$ 而取的 $B$ 的完备化。
由《概形上同调》，引理 \ref{coherent-lemma-inverse-systems-affine}，
对象 $(\mathcal{F}_n)$ 对应于有限 $B'$-模 $N$；其含义是，
$\mathcal{F}_n$ 是与有限 $B$-模 $N/I^n N$ 相伴的凝聚模。
把引理 \ref{algebraization-lemma-cd-one-extend-to-module} 应用于
$I \subset A \to B'$ 与 $N$，可知必要时增大 $c$ 并相应调整
$\varphi_j$ 之后，得到唯一的映射
$$
\Phi_N : \bigoplus\nolimits_{n \geq 0} I^{nc}N \to N[T_1, \ldots, T_r]
$$
并具有相应性质。注意，在次数 $n$，得到映射
$N/I^mN \to \bigoplus_{e_1 + \ldots + e_r = n}
N/I^{m - nc}N \cdot T_1^{e_1} \ldots T_r^{e_r}$ 的逆系统，其中
$m \geq nc$。翻译回凝聚层，可知 $\Phi_N$ 对应于映射系统
$$
\Phi^n_m :
I^{nc}\mathcal{F}_m
\longrightarrow
\bigoplus\nolimits_{e_1 + \ldots + e_r = n}
\mathcal{F}_{m - nc} \cdot T_1^{e_1} \ldots T_r^{e_r}
$$
，其中 $m \geq nc$ 与 $n \geq 1$ 变化。对这些映射关于 $m$
取逆极限，得到 $\Phi_\mathcal{F} = \bigoplus_n \lim_m \Phi^n_m$。
注意 $\lim_m I^t\mathcal{F}_m = I^t \mathcal{F}$；例如在仿射开集上求值
即可看出。不过事实上并不需要这一点，因为显然有映射
$I^t\mathcal{F} \to \lim_m I^t\mathcal{F}_m$。
\end{proof}

\begin{lemma}
\label{algebraization-lemma-topology-I-adic}
设 $I$ 是诺特环 $A$ 的理想。设 $X$ 是 $\Spec(A)$ 上的诺特概形。设
$$
\ldots \to \mathcal{F}_3 \to \mathcal{F}_2 \to \mathcal{F}_1
$$
是凝聚 $\mathcal{O}_X$-模的一个逆系统，并且
$\mathcal{F}_n = \mathcal{F}_{n + 1}/I^n\mathcal{F}_{n + 1}$。
若 $\text{cd}(A, I) = 1$，则对所有 $p \in \mathbf{Z}$，
$\lim H^p(X, \mathcal{F}_n)$ 上的极限拓扑是 $I$-进拓扑。
\end{lemma}

\begin{proof}
首先，显然 $I^t \lim H^p(X, \mathcal{F}_n)$ 在
$H^p(X, \mathcal{F}_t)$ 中映为零。因此，$I$-进拓扑比极限拓扑更细。
反过来，置 $\mathcal{F} = \lim \mathcal{F}_n$，选取生成元
$f_1, \ldots, f_r$，它们生成 $I$；选取 $c \geq 1$，并选取引理
\ref{algebraization-lemma-cd-is-one-for-system} 中的 $\Phi_\mathcal{F}$。
下文将直接使用引理 \ref{algebraization-lemma-properties-system} 的结果而不再说明。
特别地，有短正合列
$$
0 \to R^1\lim H^{p - 1}(X, \mathcal{F}_n) \to H^p(X, \mathcal{F})
\to \lim H^p(X, \mathcal{F}_n) \to 0
$$
因此，可以把任意元素 $\xi$（它属于 $\lim H^p(X, \mathcal{F}_n)$）
提升为元素 $\xi' \in H^p(X, \mathcal{F})$。假设 $\xi$ 映为
$H^p(X, \mathcal{F}_{nc})$ 中的零，其中 $n$ 是某个整数；换言之，假设 $\xi$
在极限拓扑中是``小的''。有短正合列
$$
0 \to I^{nc}\mathcal{F} \to \mathcal{F} \to \mathcal{F}_{nc} \to 0
$$
因而该假设意味着可以把 $\xi'$ 提升为元素
$\xi'' \in H^p(X, I^{nc}\mathcal{F})$。应用 $\Phi_\mathcal{F}$，得到
$$
\Phi_\mathcal{F}(\xi'') = \sum\nolimits_{e_1 + \ldots + e_r = n}
\xi'_{e_1, \ldots, e_r} \cdot T_1^{e_1} \ldots T_r^{e_r}
$$
，其中某些 $\xi'_{e_1, \ldots, e_r} \in H^p(X, \mathcal{F})$。
令 $\xi_{e_1, \ldots, e_r} \in \lim H^p(X, \mathcal{F}_n)$ 为它们的像，
再利用引理 \ref{algebraization-lemma-cd-is-one-for-system} 的最后一个断言，可得
$$
\xi = \sum f_1^{e_1} \ldots f_r^{e_r} \xi_{e_1, \ldots, e_r}
$$
属于 $I^n \lim H^p(X, \mathcal{F}_n)$，如所需。
\end{proof}

\begin{example}
\label{algebraization-example-not-I-adic}
设 $k$ 是域。设 $A = k[x, y][[s, t]]/(xs - yt)$。
设 $I = (s, t)$ 且 $\mathfrak a = (x, y, s, t)$。
设 $X = \Spec(A) - V(\mathfrak a)$ 且
$\mathcal{F}_n = \mathcal{O}_X/I^n\mathcal{O}_X$。注意有理函数
$$
g = \frac{t}{x} = \frac{s}{y}
$$
在一个开邻域 $V \subset X$ 上正则，该邻域包含
$V(I\mathcal{O}_X)$。
因此，每个幂 $g^e$ 都确定一个截面
$g^e \in M = \lim H^0(X, \mathcal{F}_n)$。注意，$g^e \to 0$，
这里 $e \to \infty$，且收敛是相对于 $M$ 上的极限拓扑而言，因为 $g^e$ 在
$\mathcal{F}_e$ 中映为零。另一方面，$g^e \not \in IM$ 对任意 $e$ 成立；
读者可通过计算 $H^0(U, \mathcal{F}_n)$ 看出这一点，
此处略去计算。注意 $\text{cd}(A, I) = 2$。
因此，引理 \ref{algebraization-lemma-topology-I-adic} 的结果是最佳的。
\end{example}








\section{Mittag--Leffler 条件}
\label{algebraization-section-ML}

\noindent
对诺特局部环的极大理想取局部上同调时，我们往往自动得到
Mittag--Leffler 条件。这蕴涵：在穿孔谱上，对具有满射过渡映射的
凝聚层逆系统，其高阶上同调群也具有同一性质。

\begin{lemma}
\label{algebraization-lemma-descending-chain}
设 $(A, \mathfrak m)$ 是诺特局部环。
\begin{enumerate}
\item 设 $M$ 是有限 $A$-模。则 $A$-模
$H^i_\mathfrak m(M)$ 对任意 $i$ 都满足降链条件。
\item 设 $U = \Spec(A) \setminus \{\mathfrak m\}$ 是
$A$ 的穿孔谱。设 $\mathcal{F}$ 是凝聚 $\mathcal{O}_U$-模。
则 $A$-模 $H^i(U, \mathcal{F})$ 在 $i > 0$ 时满足降链条件。
\end{enumerate}
\end{lemma}

\begin{proof}
对 $M$ 的支撑维数作归纳来证明 (1)。当 $M = 0$ 时命题成立，
因而可以并且确实假设 $M$ 非零。

\medskip\noindent
归纳基。若 $\dim(\text{Supp}(M)) = 0$，则 $M$ 的支撑是
$\{\mathfrak m\}$，并且由局部上同调的构造可知
$H^0_\mathfrak m(M) = M$，且 $H^i_\mathfrak m(M) = 0$ 对
$i > 0$ 成立；见《对偶复形》，第
\ref{dualizing-section-local-cohomology} 节。
由于 $M$ 具有有限长度（《代数》，引理 \ref{algebra-lemma-length-finite}），
它满足降链条件。

\medskip\noindent
归纳步骤。假设 $\dim(\text{Supp}(M)) > 0$。
由归纳基，有限模 $H^0_\mathfrak m(M) \subset M$ 满足降链条件。
由《对偶复形》，引理 \ref{dualizing-lemma-divide-by-torsion}，
可以把 $M$ 换成 $M/H^0_\mathfrak m(M)$。
于是 $H^0_\mathfrak m(M) = 0$，即 $M$ 的深度 $\geq 1$；见
《对偶复形》，引理 \ref{dualizing-lemma-depth}。
选取 $x \in \mathfrak m$，使得 $x : M \to M$ 是单射。
由《代数》，引理 \ref{algebra-lemma-one-equation-module}，有
$\dim(\text{Supp}(M/xM)) = \dim(\text{Supp}(M)) - 1$，故可应用归纳假设。
选取指标 $i$ 并考虑正合列
$$
H^{i - 1}_\mathfrak m(M/xM) \to H^i_\mathfrak m(M) \xrightarrow{x}
H^i_\mathfrak m(M)
$$
，它来自短正合列 $0 \to M \xrightarrow{x} M \to M/xM \to 0$。
由此，$x$-挠子模 $H^i_\mathfrak m(M)[x]$ 是一个满足降链条件的模的商，
因而自身也满足降链条件。于是 $\mathfrak m$-挠子模
$H^i_\mathfrak m(M)[\mathfrak m]$ 满足降链条件（因而在
$A/\mathfrak m$ 上有限维）。所以由《对偶复形》，引理
\ref{dualizing-lemma-describe-categories}.
可知 $\mathfrak m$-幂挠模 $H^i_\mathfrak m(M)$ 满足降链条件。

\medskip\noindent
(2) 由 (1) 与《局部上同调》，引理
\ref{local-cohomology-lemma-finiteness-pushforwards-and-H1-local} 得出。
\end{proof}

\begin{lemma}
\label{algebraization-lemma-ML-local}
设 $(A, \mathfrak m)$ 是诺特局部环。
\begin{enumerate}
\item 设 $(M_n)$ 是有限 $A$-模的一个逆系统。则逆系统
$H^i_\mathfrak m(M_n)$ 对任意 $i$ 都满足 Mittag--Leffler 条件。
\item 设 $U = \Spec(A) \setminus \{\mathfrak m\}$ 是
$A$ 的穿孔谱。设 $\mathcal{F}_n$ 是凝聚 $\mathcal{O}_U$-模的一个逆系统。
则逆系统 $H^i(U, \mathcal{F}_n)$ 在 $i > 0$ 时
满足 Mittag--Leffler 条件。
\end{enumerate}
\end{lemma}

\begin{proof}
由引理 \ref{algebraization-lemma-descending-chain} 立即得出。
\end{proof}

\begin{lemma}
\label{algebraization-lemma-terrific}
设 $(A, \mathfrak m)$ 是诺特局部环。
设 $(M_n)$ 是有限 $A$-模的一个逆系统。
设 $M \to \lim M_n$ 是映射，其中 $M$ 是有限 $A$-模，
并且对某个 $i$，映射
$H^i_\mathfrak m(M) \to \lim H^i_\mathfrak m(M_n)$
是同构。则逆系统 $H^i_\mathfrak m(M_n)$
本质常值，其值为 $H^i_\mathfrak m(M)$。
\end{lemma}

\begin{proof}
由引理 \ref{algebraization-lemma-ML-local}，逆系统 $H^i_\mathfrak m(M_n)$
满足 Mittag--Leffler 条件。设 $E_n \subset H^i_\mathfrak m(M_n)$
是 $H^i_\mathfrak m(M_{n'})$ 在 $n' \gg n$ 时的像。
则 $(E_n)$ 是具有满射过渡映射的逆系统，并且
$H^i_\mathfrak m(M) = \lim E_n$。由引理 \ref{algebraization-lemma-descending-chain}，
$H^i_\mathfrak m(M)$ 满足降链条件，故满射
$H^i_\mathfrak m(M) \to E_n$ 中只有有限多个具有非平凡核。
因此，$E_n \to E_{n - 1}$ 对所有 $n \gg 0$ 都是同构，如所需。
\end{proof}

\begin{lemma}
\label{algebraization-lemma-local-cohomology-derived-completion}
设 $(A, \mathfrak m)$ 是诺特局部环。
设 $I \subset A$ 是理想。设 $M$ 是有限 $A$-模。则
$$
H^i(R\Gamma_\mathfrak m(M)^\wedge) = \lim H^i_\mathfrak m(M/I^nM)
$$
对所有 $i$ 成立，其中 $R\Gamma_\mathfrak m(M)^\wedge$ 表示导出
$I$-进完备化。
\end{lemma}

\begin{proof}
应用《对偶复形》，引理 \ref{dualizing-lemma-completion-local}
与引理 \ref{algebraization-lemma-ML-local}，可知各 $R^1\lim$ 项消失。
\end{proof}








\section{环化位点上的导出完备化}
\label{algebraization-section-derived-completion}

\noindent
我们强烈建议读者初次阅读时跳过本节。

\medskip\noindent
这些内容的代数版本见《代数进阶》，第
\ref{more-algebra-section-derived-completion} 节。
设 $\mathcal{O}$ 是位点 $\mathcal{C}$ 上的环层。
设 $f$ 是 $\mathcal{O}$ 的整体截面。用 $\mathcal{O}_f$ 表示与局部化预层
$U \mapsto \mathcal{O}(U)_f$ 相伴的层。

\begin{lemma}
\label{algebraization-lemma-map-twice-localize}
设 $(\mathcal{C}, \mathcal{O})$ 是环化位点。设 $f$ 是 $\mathcal{O}$ 的整体截面。
\begin{enumerate}
\item 对 $L, N \in D(\mathcal{O}_f)$，有
$R\SheafHom_\mathcal{O}(L, N) = R\SheafHom_{\mathcal{O}_f}(L, N)$.
特别地，两个 $\mathcal{O}_f$-结构在
$R\SheafHom_\mathcal{O}(L, N)$ 上相同。
\item 对 $K \in D(\mathcal{O})$ 与
$L \in D(\mathcal{O}_f)$，有
$$
R\SheafHom_\mathcal{O}(L, K) =
R\SheafHom_{\mathcal{O}_f}(L, R\SheafHom_\mathcal{O}(\mathcal{O}_f, K))
$$
特别地，
$R\SheafHom_\mathcal{O}(\mathcal{O}_f,
R\SheafHom_\mathcal{O}(\mathcal{O}_f, K)) =
R\SheafHom_\mathcal{O}(\mathcal{O}_f, K)$.
\item 若 $g$ 是 $\mathcal{O}$ 的另一个整体截面，则
$$
R\SheafHom_\mathcal{O}(\mathcal{O}_f, R\SheafHom_\mathcal{O}(\mathcal{O}_g, K))
= R\SheafHom_\mathcal{O}(\mathcal{O}_{gf}, K).
$$
\end{enumerate}
\end{lemma}

\begin{proof}
(1) 的证明。设 $\mathcal{J}^\bullet$ 是 $\mathcal{O}_f$-模 K-内射复形，
并且表示 $N$。由《位点上的上同调》，引理
\ref{sites-cohomology-lemma-K-injective-flat}，可知 $\mathcal{J}^\bullet$
也是 $\mathcal{O}$-模 K-内射复形。设 $\mathcal{F}^\bullet$ 是
$\mathcal{O}_f$-模复形，并且表示 $L$。则
$$
R\SheafHom_\mathcal{O}(L, N) =
R\SheafHom_\mathcal{O}(\mathcal{F}^\bullet, \mathcal{J}^\bullet) =
R\SheafHom_{\mathcal{O}_f}(\mathcal{F}^\bullet, \mathcal{J}^\bullet)
$$
；这是由《位点上的模》，引理
\ref{sites-modules-lemma-epimorphism-modules} 得出的，因为
$\mathcal{J}^\bullet$ 既是 $\mathcal{O}$-模 K-内射复形，
也是 $\mathcal{O}_f$-模 K-内射复形。

\medskip\noindent
(2) 的证明。设 $\mathcal{I}^\bullet$ 是 $\mathcal{O}$-模 K-内射复形，
并且表示 $K$。于是
$R\SheafHom_\mathcal{O}(\mathcal{O}_f, K)$ 由
$\SheafHom_\mathcal{O}(\mathcal{O}_f, \mathcal{I}^\bullet)$ 表示；由
《位点上的上同调》，引理 \ref{sites-cohomology-lemma-hom-K-injective} 与
\ref{sites-cohomology-lemma-K-injective-flat}，后者既是
$\mathcal{O}_f$-模 K-内射复形，也是 $\mathcal{O}$-模 K-内射复形。
设 $\mathcal{F}^\bullet$ 是 $\mathcal{O}_f$-模复形，并且表示 $L$。则
$$
R\SheafHom_\mathcal{O}(L, K) =
R\SheafHom_\mathcal{O}(\mathcal{F}^\bullet, \mathcal{I}^\bullet) =
R\SheafHom_{\mathcal{O}_f}(\mathcal{F}^\bullet,
\SheafHom_\mathcal{O}(\mathcal{O}_f, \mathcal{I}^\bullet))
$$
；这是由《位点上的模》，引理
\ref{sites-modules-lemma-adjoint-hom-restrict} 以及
$\SheafHom_\mathcal{O}(\mathcal{O}_f, \mathcal{I}^\bullet)$ 是
$\mathcal{O}_f$-模 K-内射复形得出的。

\medskip\noindent
(3) 的证明。这由以下两个事实得出：
$R\SheafHom_\mathcal{O}(\mathcal{O}_g, \mathcal{I}^\bullet)$
作为 $\mathcal{O}$-模复形是 K-内射的；并且
$\SheafHom_\mathcal{O}(\mathcal{O}_f,
\SheafHom_\mathcal{O}(\mathcal{O}_g, \mathcal{H})) = 
\SheafHom_\mathcal{O}(\mathcal{O}_{gf}, \mathcal{H})$
对所有 $\mathcal{O}$-模层 $\mathcal{H}$ 成立。
\end{proof}

\noindent
设 $K \in D(\mathcal{O})$。用 $T(K, f)$ 表示逆系统
$$
\ldots \to K \xrightarrow{f} K \xrightarrow{f} K
$$
在 $D(\mathcal{O})$ 中的一个导出极限（《导出范畴》，定义
\ref{derived-definition-derived-limit}）。

\begin{lemma}
\label{algebraization-lemma-hom-from-Af}
设 $(\mathcal{C}, \mathcal{O})$ 是环化位点。设 $f$ 是 $\mathcal{O}$
的整体截面。设 $K \in D(\mathcal{O})$。下列条件等价：
\begin{enumerate}
\item $R\SheafHom_\mathcal{O}(\mathcal{O}_f, K) = 0$,
\item $R\SheafHom_\mathcal{O}(L, K) = 0$ 对所有 $L$ 成立，
其中该对象位于 $D(\mathcal{O}_f)$ 中，
\item $T(K, f) = 0$.
\end{enumerate}
\end{lemma}

\begin{proof}
显然 (2) 蕴涵 (1)。蕴涵 (1) $\Rightarrow$ (2) 由引理
\ref{algebraization-lemma-map-twice-localize} 得出。$\mathcal{O}$-模
$\mathcal{O}_f$ 的一个自由解消为
$$
0 \to \bigoplus\nolimits_{n \in \mathbf{N}} \mathcal{O} \to
\bigoplus\nolimits_{n \in \mathbf{N}} \mathcal{O}
\to \mathcal{O}_f \to 0
$$
，其中第一个映射把局部截面 $(x_0, x_1, \ldots)$ 映到
$(x_0, x_1 - fx_0, x_2 - fx_1, \ldots)$，第二个映射把
$(x_0, x_1, \ldots)$ 映到 $x_0 + x_1/f + x_2/f^2 + \ldots$。
应用 $\SheafHom_\mathcal{O}(-, \mathcal{I}^\bullet)$，其中
$\mathcal{I}^\bullet$ 是 $\mathcal{O}$-模 K-内射复形，并且表示 $K$，
得到复形的短正合列
$$
0 \to \SheafHom_\mathcal{O}(\mathcal{O}_f, \mathcal{I}^\bullet) \to
\prod \mathcal{I}^\bullet \to \prod \mathcal{I}^\bullet \to 0
$$
，因为 $\mathcal{I}^n$ 是内射 $\mathcal{O}$-模。这些积是
$D(\mathcal{O})$ 中的积；见《内射对象》，引理
\ref{injectives-lemma-derived-products}。这意味着对象 $T(K, f)$
表示 $R\SheafHom_\mathcal{O}(\mathcal{O}_f, K)$，并且是在
$D(\mathcal{O})$ 中表示。
因此 (1) 与 (3) 等价。
\end{proof}

\begin{lemma}
\label{algebraization-lemma-ideal-of-elements-complete-wrt}
设 $(\mathcal{C}, \mathcal{O})$ 是环化位点。设 $K \in D(\mathcal{O})$。
对 $U$，令 $\mathcal{I}(U)$ 为截面 $f \in \mathcal{O}(U)$ 所成的集合，
这些截面满足 $T(K|_U, f) = 0$；这一规则给出 $\mathcal{O}$ 中的理想层。
\end{lemma}

\begin{proof}
下文直接使用引理 \ref{algebraization-lemma-hom-from-Af} 的结果而不再说明。
若 $f \in \mathcal{I}(U)$ 且 $g \in \mathcal{O}(U)$，则
$\mathcal{O}_{U, gf}$ 是 $\mathcal{O}_{U, f}$-模，因而
$R\SheafHom_\mathcal{O}(\mathcal{O}_{U, gf}, K|_U) = 0$，故
$gf \in \mathcal{I}(U)$。设 $f, g \in \mathcal{O}(U)$。则有短正合列
$$
0 \to \mathcal{O}_{U, f + g} \to
\mathcal{O}_{U, f(f + g)} \oplus \mathcal{O}_{U, g(f + g)} \to
\mathcal{O}_{U, gf(f + g)} \to 0
$$
，因为 $f, g$ 在 $\mathcal{O}(U)_{f + g}$ 中生成单位理想。
这由《代数》，引理 \ref{algebra-lemma-standard-covering}
以及最后一个箭头是满射这一简单事实得出。
因为 $R\SheafHom_\mathcal{O}( - , K|_U)$ 是三角范畴之间的正合函子，
下列三个对象的消失：
$R\SheafHom_{\mathcal{O}_U}(\mathcal{O}_{U, f(f + g)}, K|_U)$,
$R\SheafHom_{\mathcal{O}_U}(\mathcal{O}_{U, g(f + g)}, K|_U)$，以及
$R\SheafHom_{\mathcal{O}_U}(\mathcal{O}_{U, gf(f + g)}, K|_U)$,
蕴涵
$R\SheafHom_{\mathcal{O}_U}(\mathcal{O}_{U, f + g}, K|_U)$.
层条件的验证从略。
\end{proof}

\noindent
对任意环化位点，都可以作如下定义。

\begin{definition}
\label{algebraization-definition-derived-complete}
设 $(\mathcal{C}, \mathcal{O})$ 是环化位点。
设 $\mathcal{I} \subset \mathcal{O}$ 是理想层。
设 $K \in D(\mathcal{O})$。若 $K$ {\it 关于 $\mathcal{I}$ 导出完备}，
是指对每个对象 $U$（属于 $\mathcal{C}$）以及 $f \in \mathcal{I}(U)$，
对象 $T(K|_U, f)$ 在 $D(\mathcal{O}_U)$ 中都为零。
\end{definition}

\noindent
显然，由导出完备对象组成的全子范畴
$D_{comp}(\mathcal{O}) = D_{comp}(\mathcal{O}, \mathcal{I}) \subset
D(\mathcal{O})$
是饱和三角子范畴；见《导出范畴》，定义
\ref{derived-definition-triangulated-subcategory} 与
\ref{derived-definition-saturated}。该子范畴在 $D(\mathcal{O})$ 中对积与
同伦极限封闭，但一般不对可数直和封闭。

\begin{lemma}
\label{algebraization-lemma-derived-complete-internal-hom}
设 $(\mathcal{C}, \mathcal{O})$ 是环化位点。
设 $\mathcal{I} \subset \mathcal{O}$ 是理想层。
若 $K \in D(\mathcal{O})$ 且 $L \in D_{comp}(\mathcal{O})$，则
$R\SheafHom_\mathcal{O}(K, L) \in D_{comp}(\mathcal{O})$.
\end{lemma}

\begin{proof}
设 $U$ 是 $\mathcal{C}$ 的对象，并设 $f \in \mathcal{I}(U)$。回忆，
$$
\Hom_{D(\mathcal{O}_U)}(\mathcal{O}_{U, f}, R\SheafHom_\mathcal{O}(K, L)|_U)
=
\Hom_{D(\mathcal{O}_U)}(
K|_U \otimes_{\mathcal{O}_U}^\mathbf{L} \mathcal{O}_{U, f}, L|_U)
$$
；这是《位点上的上同调》，引理
\ref{sites-cohomology-lemma-internal-hom}。右边由引理
\ref{algebraization-lemma-hom-from-Af} 以及内部 Hom 与实际 Hom 之间的关系而为零；见
《位点上的上同调》，引理
\ref{sites-cohomology-lemma-section-RHom-over-U}。
对所有 $U'/U$，同样的消失成立。因此，对象
$R\SheafHom_{\mathcal{O}_U}(\mathcal{O}_{U, f},
R\SheafHom_\mathcal{O}(K, L)|_U)$ 作为 $D(\mathcal{O}_U)$ 的对象，
其第 $0$ 个上同调层消失（由上述引文）。其他上同调层亦然；也就是说，
$R\SheafHom_{\mathcal{O}_U}(\mathcal{O}_{U, f},
R\SheafHom_\mathcal{O}(K, L)|_U)$ 在 $D(\mathcal{O}_U)$ 中为零。
由引理 \ref{algebraization-lemma-hom-from-Af} 即得结论。
\end{proof}

\begin{lemma}
\label{algebraization-lemma-restriction-derived-complete}
设 $\mathcal{C}$ 是位点。设 $\mathcal{O} \to \mathcal{O}'$
是环层同态。设 $\mathcal{I} \subset \mathcal{O}$ 是理想层。
$D_{comp}(\mathcal{O}, \mathcal{I})$ 在限制函子
$D(\mathcal{O}') \to D(\mathcal{O})$ 下的原像是
$D_{comp}(\mathcal{O}', \mathcal{I}\mathcal{O}')$.
\end{lemma}

\begin{proof}
利用引理 \ref{algebraization-lemma-ideal-of-elements-complete-wrt} 可知，
$K' \in D(\mathcal{O}')$ 属于
$D_{comp}(\mathcal{O}', \mathcal{I}\mathcal{O}')$
当且仅当 $T(K'|_U, f)$ 对每个局部截面 $f \in \mathcal{I}(U)$ 都为零。
注意，$T(K'|_U, f)$ 的上同调层是在阿贝尔层范畴中计算的，
所以把 $f$ 看作 $\mathcal{O}$ 的截面，或者取 $f$ 在
$\mathcal{O}'$ 中的像，并无区别。由此及导出完备对象的定义，立即得出引理。
\end{proof}

\begin{lemma}
\label{algebraization-lemma-pushforward-derived-complete}
设 $f : (\Sh(\mathcal{D}), \mathcal{O}') \to (\Sh(\mathcal{C}), \mathcal{O})$
是环化拓扑斯态射。设 $\mathcal{I} \subset \mathcal{O}$ 与
$\mathcal{I}' \subset \mathcal{O}'$ 是理想层，并且 $f^\sharp$ 把
$f^{-1}\mathcal{I}$ 映入 $\mathcal{I}'$。则 $Rf_*$ 把
$D_{comp}(\mathcal{O}', \mathcal{I}')$ 映入
$D_{comp}(\mathcal{O}, \mathcal{I})$。
\end{lemma}

\begin{proof}
可以假设 $f$ 由环化位点态射给出，该态射对应于连续函子
$\mathcal{C} \to \mathcal{D}$（《位点上的模》，引理
\ref{sites-modules-lemma-morphism-ringed-topoi-comes-from-morphism-ringed-sites}
)。设 $U$ 是 $\mathcal{C}$ 的对象，并设 $g$ 是 $\mathcal{I}$ 在
$U$ 上的截面。必须证明
$\Hom_{D(\mathcal{O}_U)}(\mathcal{O}_{U, g}, Rf_*K|_U) = 0$
，只要 $K$ 关于 $\mathcal{I}'$ 导出完备。事实上，由《位点上的上同调》，引理
\ref{sites-cohomology-lemma-section-RHom-over-U}
，把这一点应用于 $U$ 上的所有对象及 $K$ 的所有平移，可推出
$R\SheafHom_{\mathcal{O}_U}(\mathcal{O}_{U, g}, Rf_*K|_U)$ 为零；
这又推出 $T(Rf_*K|_U, g)$ 为零（引理
\ref{algebraization-lemma-hom-from-Af}），而这正是需要证明的（定义
\ref{algebraization-definition-derived-complete}）。
设 $V$ 是 $\mathcal{D}$ 中 $U$ 的像。则
$$
\Hom_{D(\mathcal{O}_U)}(\mathcal{O}_{U, g}, Rf_*K|_U) =
\Hom_{D(\mathcal{O}'_V)}(\mathcal{O}'_{V, g'}, K|_V) = 0
$$
，其中 $g' = f^\sharp(g) \in \mathcal{I}'(V)$。第二个等号成立是因为
$K$ 导出完备；第一个等号成立是因为 $\mathcal{O}_{U, g}$ 的导出拉回是
$\mathcal{O}'_{V, g'}$，并应用《位点上的上同调》，引理
\ref{sites-cohomology-lemma-adjoint}。
\end{proof}

\noindent
下一个引理是存在导出完备化的最简单情形。

\begin{lemma}
\label{algebraization-lemma-derived-completion}
设 $(\mathcal{C}, \mathcal{O})$ 是环化位点。设 $f_1, \ldots, f_r$
是 $\mathcal{O}$ 的整体截面。设 $\mathcal{I} \subset \mathcal{O}$ 是由
$f_1, \ldots, f_r$ 生成的理想层。则包含函子
$D_{comp}(\mathcal{O}) \to D(\mathcal{O})$ 有左伴随；也就是说，
给定任意对象 $K$（它属于 $D(\mathcal{O})$），存在映射
$K \to K^\wedge$，其中 $K^\wedge$ 位于 $D_{comp}(\mathcal{O})$ 中，
使得映射
$$
\Hom_{D(\mathcal{O})}(K^\wedge, E) \longrightarrow \Hom_{D(\mathcal{O})}(K, E)
$$
在 $E$ 位于 $D_{comp}(\mathcal{O})$ 中时是双射。事实上，有
$$
K^\wedge =
R\SheafHom_\mathcal{O}
(\mathcal{O} \to \prod\nolimits_{i_0} \mathcal{O}_{f_{i_0}} \to
\prod\nolimits_{i_0 < i_1} \mathcal{O}_{f_{i_0}f_{i_1}} \to
\ldots \to \mathcal{O}_{f_1\ldots f_r}, K)
$$
，并且关于 $K$ 具有函子性。
\end{lemma}

\begin{proof}
用引理中最后一个展示公式定义 $K^\wedge$。有复形映射
$$
(\mathcal{O} \to \prod\nolimits_{i_0} \mathcal{O}_{f_{i_0}} \to
\prod\nolimits_{i_0 < i_1} \mathcal{O}_{f_{i_0}f_{i_1}} \to
\ldots \to \mathcal{O}_{f_1\ldots f_r}) \longrightarrow \mathcal{O}
$$
，它诱导映射 $K \to K^\wedge$。只需证明 $K^\wedge$ 导出完备，
并且 $K \to K^\wedge$ 在 $K$ 导出完备时是同构。

\medskip\noindent
设 $f$ 是 $\mathcal{O}$ 的整体截面。由引理
\ref{algebraization-lemma-map-twice-localize}，对象
$R\SheafHom_\mathcal{O}(\mathcal{O}_f, K^\wedge)$
等于
$$
R\SheafHom_\mathcal{O}(
(\mathcal{O}_f \to \prod\nolimits_{i_0} \mathcal{O}_{ff_{i_0}} \to
\prod\nolimits_{i_0 < i_1} \mathcal{O}_{ff_{i_0}f_{i_1}} \to
\ldots \to \mathcal{O}_{ff_1\ldots f_r}), K)
$$
若 $f = f_i$ 对某个 $i$ 成立，则 $f_1, \ldots, f_r$ 在
$\mathcal{O}_f$ 中生成单位理想，因而增广交错 {\v C}ech 复形
$$
\mathcal{O}_f \to \prod\nolimits_{i_0} \mathcal{O}_{ff_{i_0}} \to
\prod\nolimits_{i_0 < i_1} \mathcal{O}_{ff_{i_0}f_{i_1}} \to
\ldots \to \mathcal{O}_{ff_1\ldots f_r}
$$
为零（甚至同伦于零）。由此可见 $K^\wedge$ 导出完备。

\medskip\noindent
若 $K$ 导出完备，则 $R\SheafHom_\mathcal{O}(\mathcal{O}_f, K)$
对所有 $f = f_{i_0} \ldots f_{i_p}$、$p \geq 0$ 都为零。因此，
$K \to K^\wedge$ 在 $D(\mathcal{O})$ 中是同构。
\end{proof}

\noindent
下面说明导出完备化为何确实是一种完备化。

\begin{lemma}
\label{algebraization-lemma-derived-completion-koszul}
设 $(\mathcal{C}, \mathcal{O})$ 是环化位点。设 $f_1, \ldots, f_r$
是 $\mathcal{O}$ 的整体截面。设 $\mathcal{I} \subset \mathcal{O}$ 是由
$f_1, \ldots, f_r$ 生成的理想层。设 $K \in D(\mathcal{O})$。
引理 \ref{algebraization-lemma-derived-completion} 的导出完备化 $K^\wedge$ 由公式
$$
K^\wedge = R\lim K \otimes^\mathbf{L}_\mathcal{O} K_n
$$
给出，其中 $K_n = K(\mathcal{O}, f_1^n, \ldots, f_r^n)$ 是
关于 $f_1^n, \ldots, f_r^n$ 的 Koszul 复形，定义在 $\mathcal{O}$ 上。
\end{lemma}

\begin{proof}
在《代数进阶》，引理
\ref{more-algebra-lemma-extended-alternating-Cech-is-colimit-koszul}
中已经看到，增广交错 {\v C}ech 复形
$$
\mathcal{O} \to \prod\nolimits_{i_0} \mathcal{O}_{f_{i_0}} \to
\prod\nolimits_{i_0 < i_1} \mathcal{O}_{f_{i_0}f_{i_1}} \to
\ldots \to \mathcal{O}_{f_1\ldots f_r}
$$
是 Koszul 复形 $K^n = K(\mathcal{O}, f_1^n, \ldots, f_r^n)$ 的余极限，
这些复形位于次数 $0, \ldots, r$。注意，$K^n$ 是有限自由
$\mathcal{O}$-模的有限链复形，其对偶为
$\SheafHom_\mathcal{O}(K^n, \mathcal{O}) = K_n$，其中 $K_n$ 是位于次数
$-r, \ldots, 0$ 的 Koszul 上链复形（如通常约定）。由引理
\ref{algebraization-lemma-derived-completion}，函子 $E \mapsto E^\wedge$ 是从增广交错
{\v C}ech 复形取 $R\SheafHom$ 而得到的，其目标为 $E$：
$$
E^\wedge = R\SheafHom(\colim K^n, E)
$$
由《位点上的上同调》，引理
\ref{sites-cohomology-lemma-colim-and-lim-of-duals}，这等于
$R\lim (E \otimes_\mathcal{O}^\mathbf{L} K_n)$。
\end{proof}

\begin{lemma}
\label{algebraization-lemma-all-rings}
存在一种如下的构造方式：
\begin{enumerate}
\item 对每个二元组 $(A, I)$，其中 $A$ 是环而 $I \subset A$ 是有限生成理想，
构造复形 $K(A, I)$，它由 $A$-模组成；
\item 构造映射 $K(A, I) \to A$，它是 $A$-模复形之间的映射；
\item 对每个环映射 $A \to B$ 与有限生成理想 $I \subset A$，
构造复形映射 $K(A, I) \to K(B, IB)$；
\end{enumerate}
并满足：
\begin{enumerate}
\item[(a)] 对 $A \to B$ 与有限生成的 $I \subset A$，图
$$
\xymatrix{
K(A, I) \ar[r] \ar[d] & A \ar[d] \\
K(B, IB) \ar[r] & B
}
$$
交换；
\item[(b)] 对 $A \to B \to C$ 与有限生成的 $I \subset A$，映射的复合
$K(A, I) \to K(B, IB) \to K(C, IC)$ 就是映射
$K(A, I) \to K(C, IC)$；
\item[(c)] 对 $A \to B$ 与有限生成理想 $I \subset A$，诱导映射
$K(A, I) \otimes_A^\mathbf{L} B \to K(B, IB)$
在 $D(B)$ 中是同构；并且
\item[(d)] 若 $I = (f_1, \ldots, f_r) \subset A$，则有交换图
$$
\xymatrix{
(A \to \prod\nolimits_{i_0} A_{f_{i_0}} \to
\prod\nolimits_{i_0 < i_1} A_{f_{i_0}f_{i_1}} \to
\ldots \to A_{f_1\ldots f_r}) \ar[r] \ar[d] &  K(A, I) \ar[d] \\
A \ar[r]^1 & A
}
$$
，它位于 $D(A)$ 中，且横向箭头都是同构。
\end{enumerate}
\end{lemma}

\begin{proof}
设 $S$ 是形如 $A_0$ 的环所成的集合，其中
$A_0 = \mathbf{Z}[x_1, \ldots, x_n]/J$。
每个有限型 $\mathbf{Z}$-代数都同构于 $S$ 的一个元素。
设 $\mathcal{A}_0$ 为如下范畴：其对象是二元组 $(A_0, I_0)$，
其中 $A_0 \in S$ 且 $I_0 \subset A_0$ 是理想；从
$(A_0, I_0) \to (B_0, J_0)$ 的态射是环映射
$\varphi : A_0 \to B_0$，它满足 $J_0 = \varphi(I_0)B_0$。

\medskip\noindent
假设能以函子方式构造 $K(A_0, I_0) \to A_0$，这是对
$\mathcal{A}_0$ 的对象而言，并满足 (a)、(b)、(c)、(d)。此时取
$$
K(A, I) = \colim_{\varphi : (A_0, I_0) \to (A, I)} K(A_0, I_0)
$$
，其中余极限遍历环映射 $\varphi : A_0 \to A$，它们满足
$\varphi(I_0)A = I$，且 $(A_0, I_0)$ 位于 $\mathcal{A}_0$ 中。
在 $(A_0, I_0) \to (A, I)$ 与
$(A_0', I_0') \to (A, I)$ 之间的态射，由映射
$(A_0, I_0) \to (A_0', I_0')$ 给出；该映射位于 $\mathcal{A}_0$ 中，
并与到 $A$ 的映射可交换。
这些 $(A_0, I_0) \to (A, I)$ 所成的范畴是滤过的（略去细节）。
此外，$\colim_{\varphi : (A_0, I_0) \to (A, I)} A_0 = A$，
所以 $K(A, I)$ 是 $A$-模复形。
最后，给定 $\varphi : A \to B$ 与 $I \subset A$，对余极限中的每个
$(A_0, I_0) \to (A, I)$，复合
$(A_0, I_0) \to (B, IB)$ 位于 $(B, IB)$ 所对应的余极限中。
这样便得到余极限之间的映射。性质 (a)、(b)、(c)、(d) 立即由此及复形
$K(A_0, I_0)$ 的相应性质得出。

\medskip\noindent
在 $\mathcal{C}_0 = \mathcal{A}_0^{opp}$ 上赋予混沌拓扑。
在 $\mathcal{C}_0$ 上赋予环层
$\mathcal{O} : (A, I) \mapsto A$。各理想 $I$ 拼合成理想层
$\mathcal{I} \subset \mathcal{O}$。选取内射解消
$\mathcal{O} \to \mathcal{J}^\bullet$。考虑对象
$$
\mathcal{F}^\bullet = \bigcup\nolimits_n \mathcal{J}^\bullet[\mathcal{I}^n]
$$
设 $U = (A, I) \in \Ob(\mathcal{C}_0)$。
由于 $\mathcal{C}_0$ 上的拓扑是混沌的，$\mathcal{J}^\bullet(U)$
是由内射 $A$-模组成的 $A$ 的解消。因此，$\mathcal{F}^\bullet(U)$
是 $D(A)$ 的对象，表示 $R\Gamma_I(A)$ 在 $D(A)$ 中的像；见
《对偶复形》，第 \ref{dualizing-section-local-cohomology} 节。
选取 $\mathcal{O}$-模复形 $\mathcal{K}^\bullet$ 以及交换图
$$
\xymatrix{
\mathcal{O} \ar[r] & \mathcal{J}^\bullet \\
\mathcal{K}^\bullet \ar[r] \ar[u] & \mathcal{F}^\bullet \ar[u]
}
$$
，其中横向箭头是拟同构。由导出范畴 $D(\mathcal{O})$ 的构造，
这可以做到。置 $K(A, I) = \mathcal{K}^\bullet(U)$，其中
$U = (A, I)$。性质 (a)、(b) 是清楚的，性质 (c)、(d) 由
《对偶复形》，引理
\ref{dualizing-lemma-compute-local-cohomology-noetherian} 与
\ref{dualizing-lemma-local-cohomology-change-rings}.
\end{proof}

\begin{lemma}
\label{algebraization-lemma-global-extended-cech-complex}
设 $(\mathcal{C}, \mathcal{O})$ 是环化位点。设
$\mathcal{I} \subset \mathcal{O}$ 是有限型理想层。
存在映射 $K \to \mathcal{O}$，它位于 $D(\mathcal{O})$ 中，使得对每个
$U \in \Ob(\mathcal{C})$，若 $\mathcal{I}|_U$ 由
$f_1, \ldots, f_r \in \mathcal{I}(U)$ 生成，则有同构
$$
(\mathcal{O}_U \to \prod\nolimits_{i_0} \mathcal{O}_{U, f_{i_0}} \to
\prod\nolimits_{i_0 < i_1} \mathcal{O}_{U, f_{i_0}f_{i_1}} \to
\ldots \to \mathcal{O}_{U, f_1\ldots f_r}) \longrightarrow K|_U
$$
，且与到 $\mathcal{O}_U$ 的映射相容。
\end{lemma}

\begin{proof}
设 $\mathcal{C}' \subset \mathcal{C}$ 是由如下对象 $U$ 组成的全子范畴：
$\mathcal{I}|_U$ 由有限多个截面生成。则 $\mathcal{C}' \to \mathcal{C}$
是特殊余连续函子（《位点》，定义
\ref{sites-definition-special-cocontinuous-functor}）。
因此，只需在 $\mathcal{C}'$ 上工作；见《位点》，引理
\ref{sites-lemma-equivalence}。换言之，可以假设对每个对象
$U$（属于 $\mathcal{C}$），存在有限生成理想 $I \subset \mathcal{I}(U)$，使得
$\mathcal{I}|_U = \Im(I \otimes \mathcal{O}_U \to \mathcal{O}_U)$。
我们称 $I$ 生成 $\mathcal{I}|_U$。
警告：并不知道 $\mathcal{I}(U)$ 是 $\mathcal{O}(U)$ 中的有限生成理想。

\medskip\noindent
设 $U$ 是对象，并设 $I \subset \mathcal{O}(U)$ 是生成
$\mathcal{I}|_U$ 的有限生成理想。在范畴 $\mathcal{C}/U$ 上考虑预层复形
$$
K_{U, I}^\bullet : U'/U \longmapsto K(\mathcal{O}(U'), I\mathcal{O}(U'))
$$
，其中 $K(-, -)$ 如引理 \ref{algebraization-lemma-all-rings} 中所述。
我们断言其层化与 $I$ 的选择无关。事实上，若
$I' \subset \mathcal{O}(U)$ 是另一个也生成 $\mathcal{I}|_U$ 的有限生成理想，
则存在覆盖 $\{U_j \to U\}$，使得
$I\mathcal{O}(U_j) = I'\mathcal{O}(U_j)$。（提示：这是因为 $I$ 与 $I'$
均有限生成且都生成 $\mathcal{I}|_U$。）因此，
$K_{U, I}^\bullet$ 与 $K_{U, I'}^\bullet$ 在任意对象上
{\it 相同}，只要该对象位于某个 $U_j$ 之上。关于层化的断言随即得出。
用 $K_U^\bullet$ 表示这一共同值。

\medskip\noindent
对 $I$ 的选择无关还表明
$K_U^\bullet|_{\mathcal{C}/U'} = K_{U'}^\bullet$
，只要给定态射 $U' \to U$，从而给定局部化态射
$\mathcal{C}/U' \to \mathcal{C}/U$。因此，各复形 $K_U^\bullet$
可粘合成一个良定义的复形 $K^\bullet$，它由 $\mathcal{O}$-模组成。
映射 $K^\bullet \to \mathcal{O}$ 的存在性以及引理中的拟同构，
立即由引理 \ref{algebraization-lemma-all-rings} 中复形 $K(-, -)$ 的相应性质得出。
\end{proof}

\begin{proposition}
\label{algebraization-proposition-derived-completion}
设 $(\mathcal{C}, \mathcal{O})$ 是环化位点。
设 $\mathcal{I} \subset \mathcal{O}$ 是有限型理想层。
包含函子 $D_{comp}(\mathcal{O}) \to D(\mathcal{O})$ 有左伴随。
\end{proposition}

\begin{proof}
设 $K \to \mathcal{O}$ 位于 $D(\mathcal{O})$ 中，并如引理
\ref{algebraization-lemma-global-extended-cech-complex} 中所构造。设 $E \in D(\mathcal{O})$。
那么，$E^\wedge = R\SheafHom(K, E)$ 连同映射 $E \to E^\wedge$
即可满足要求。事实上，在位点 $\mathcal{C}$ 上局部地，
这恢复了引理 \ref{algebraization-lemma-derived-completion} 的伴随。
由此可知，$E^\wedge$ 总是导出完备的，并且
$E \to E^\wedge$ 在 $E$ 导出完备时是同构。
\end{proof}

\begin{remark}[与完备化的比较]
\label{algebraization-remark-compare-with-completion}
设 $(\mathcal{C}, \mathcal{O})$ 是一个环化位点。
设 $\mathcal{I} \subset \mathcal{O}$ 是一个有限型理想层。
设 $K \mapsto K^\wedge$ 为命题
\ref{algebraization-proposition-derived-completion} 的导出完备化函子。
对任意 $n \geq 1$，对象
$K \otimes_\mathcal{O}^\mathbf{L} \mathcal{O}/\mathcal{I}^n$
是导出完备的，因为它被 $\mathcal{I}$ 的局部截面的幂所零化。因此存在典范分解
$$
K \to K^\wedge \to K \otimes_\mathcal{O}^\mathbf{L} \mathcal{O}/\mathcal{I}^n
$$
，它分解了典范映射
$K \to K \otimes_\mathcal{O}^\mathbf{L} \mathcal{O}/\mathcal{I}^n$。
当 $n$ 变化时，这些映射彼此相容，于是得到比较映射
$$
K^\wedge
\longrightarrow
R\lim \left(K \otimes_\mathcal{O}^\mathbf{L} \mathcal{O}/\mathcal{I}^n\right)
$$
右端更容易被辨认为某种完备化。一般而言，这个比较映射并不是同构。
\end{remark}

\begin{remark}[局部化与导出完备化]
\label{algebraization-remark-localization-and-completion}
设 $(\mathcal{C}, \mathcal{O})$ 是一个环化位点。
设 $\mathcal{I} \subset \mathcal{O}$ 是一个有限型理想层。
设 $K \mapsto K^\wedge$ 为命题
\ref{algebraization-proposition-derived-completion} 的导出完备化函子。
由该命题证明中的构造可知，$K^\wedge|_U$
是 $K|_U$ 的导出完备化，其中 $U \in \Ob(\mathcal{C})$ 任意。
也可以如下证明这一点。由导出完备对象的定义可知，
$K^\wedge|_U$ 是导出完备的。因此得到典范映射
$a : (K|_U)^\wedge \to K^\wedge|_U$。
另一方面，若 $E$ 是 $D(\mathcal{O}_U)$ 的导出完备对象，则
$Rj_*E$ 是 $D(\mathcal{O})$ 的导出完备对象，这是引理
\ref{algebraization-lemma-pushforward-derived-complete} 的结论。这里 $j$ 是局部化态射
（《位点上的模》，\ref{sites-modules-section-localize} 节）。
因此还得到典范映射
$b : K^\wedge \to Rj_*((K|_U)^\wedge)$。
我们略去对 $b$ 的伴随映射为 $a$ 的逆映射这一（形式）事实的验证。
\end{remark}

\begin{remark}[完备张量积]
\label{algebraization-remark-completed-tensor-product}
设 $(\mathcal{C}, \mathcal{O})$ 是一个环化位点，并设
$\mathcal{I} \subset \mathcal{O}$ 是一个有限型理想层。
以 $K \mapsto K^\wedge$ 表示命题
\ref{algebraization-proposition-derived-completion} 中的伴随函子。置
$$
K \otimes^\wedge_\mathcal{O} L = (K \otimes_\mathcal{O}^\mathbf{L} L)^\wedge
$$
这个{\it 完备张量积}定义了函子
$D_{comp}(\mathcal{O}) \times D_{comp}(\mathcal{O}) \to D_{comp}(\mathcal{O})$
，使得
$$
\Hom_{D_{comp}(\mathcal{O})}(K, R\SheafHom_\mathcal{O}(L, M))
=
\Hom_{D_{comp}(\mathcal{O})}(K \otimes_\mathcal{O}^\wedge L, M)
$$
对 $K, L, M \in D_{comp}(\mathcal{O})$ 成立。注意，
$R\SheafHom_\mathcal{O}(L, M) \in D_{comp}(\mathcal{O})$；
这是引理 \ref{algebraization-lemma-derived-complete-internal-hom} 的结论。
\end{remark}

\begin{lemma}
\label{algebraization-lemma-map-identifies-koszul-and-cech-complexes}
设 $\mathcal{C}$ 是一个位点。
假设 $\varphi : \mathcal{O} \to \mathcal{O}'$ 是环层的平坦同态。
设 $f_1, \ldots, f_r$ 是 $\mathcal{O}$ 的全局截面，并满足
$\mathcal{O}/(f_1, \ldots, f_r) \cong
\mathcal{O}'/(f_1, \ldots, f_r)\mathcal{O}'$。
则增广交错 {\v C}ech 复形之间的映射
$$
\xymatrix{
\mathcal{O} \to
\prod_{i_0} \mathcal{O}_{f_{i_0}} \to
\prod_{i_0 < i_1} \mathcal{O}_{f_{i_0}f_{i_1}} \to \ldots \to
\mathcal{O}_{f_1\ldots f_r} \ar[d] \\
\mathcal{O}' \to
\prod_{i_0} \mathcal{O}'_{f_{i_0}} \to
\prod_{i_0 < i_1} \mathcal{O}'_{f_{i_0}f_{i_1}} \to \ldots \to
\mathcal{O}'_{f_1\ldots f_r}
}
$$
是拟同构。
\end{lemma}

\begin{proof}
注意，第二个复形是第一个复形与 $\mathcal{O}'$ 的张量积。
第一个增广交错 {\v C}ech 复形可以写成 Koszul 复形
$K_n = K(\mathcal{O}, f_1^n, \ldots, f_r^n)$ 的余极限；见《更多代数》，
引理 \ref{more-algebra-lemma-extended-alternating-Cech-is-colimit-koszul}。
因此，只需证明 $K_n \to K_n \otimes_\mathcal{O} \mathcal{O}'$
是拟同构。由于 $\mathcal{O} \to \mathcal{O}'$ 平坦，只需证明
$H^i \to H^i \otimes_\mathcal{O} \mathcal{O}'$ 是同构，其中
$H^i$ 是第 $i$ 个上同调层，即 $H^i = H^i(K_n)$。
这些层被 $f_1^n, \ldots, f_r^n$ 零化；见《更多代数》，
引理 \ref{more-algebra-lemma-homotopy-koszul}。
因此这些层被 $(f_1, \ldots, f_r)^m$ 零化，其中某个 $m \gg 0$。
于是 $H^i \to H^i \otimes_\mathcal{O} \mathcal{O}'$ 是同构；
这是《位点上的模》，引理
\ref{sites-modules-lemma-neighbourhood-isomorphism} 的结论。
\end{proof}

\begin{lemma}
\label{algebraization-lemma-restriction-derived-complete-equivalence}
设 $\mathcal{C}$ 是一个位点。设 $\mathcal{O} \to \mathcal{O}'$
是环层的同态，并设 $\mathcal{I} \subset \mathcal{O}$
是一个有限型理想层。
若 $\mathcal{O} \to \mathcal{O}'$ 平坦且
$\mathcal{O}/\mathcal{I} \cong \mathcal{O}'/\mathcal{I}\mathcal{O}'$,
则限制函子 $D(\mathcal{O}') \to D(\mathcal{O})$
诱导等价
$D_{comp}(\mathcal{O}', \mathcal{I}\mathcal{O}') \to
D_{comp}(\mathcal{O}, \mathcal{I})$.
\end{lemma}

\begin{proof}
由引理 \ref{algebraization-lemma-pushforward-derived-complete}，限制函子
$r : D(\mathcal{O}') \to D(\mathcal{O})$
把 $D_{comp}(\mathcal{O}', \mathcal{I}\mathcal{O}')$
映入 $D_{comp}(\mathcal{O}, \mathcal{I})$。我们将构造一个拟逆
$E \mapsto E'$。

\medskip\noindent
设 $K \to \mathcal{O}$ 是 $D(\mathcal{O})$ 中由引理
\ref{algebraization-lemma-global-extended-cech-complex} 构造的态射。置
$K' = K \otimes_\mathcal{O}^\mathbf{L} \mathcal{O}'$；这是
$D(\mathcal{O}')$ 中的对象。
则 $K' \to \mathcal{O}'$ 是 $D(\mathcal{O}')$ 中的映射；
它满足引理 \ref{algebraization-lemma-global-extended-cech-complex}
相对于 $\mathcal{I}' = \mathcal{I}\mathcal{O}'$ 的结论。
映射 $K \to r(K')$ 是拟同构，这是引理
\ref{algebraization-lemma-map-identifies-koszul-and-cech-complexes} 的结论。现在，对
$E \in D_{comp}(\mathcal{O}, \mathcal{I})$，置
$$
E' = R\SheafHom_\mathcal{O}(r(K'), E)
$$
把它视为 $D(\mathcal{O}')$ 中的对象，其中使用
$\mathcal{O}'$ 在 $K'$ 上的模结构。由于 $E$ 导出完备，我们有
$E = R\SheafHom_\mathcal{O}(K, E)$；见命题
\ref{algebraization-proposition-derived-completion} 的证明。
另一方面，由于 $K \to r(K')$ 是同构，可见
$E \to r(E')$ 是 $D(\mathcal{O})$ 中的同构。为了完成证明，
还须说明：若 $E = r(M')$，其中 $M'$ 是
$D_{comp}(\mathcal{O}', \mathcal{I}')$ 的对象，则
$E' \cong M'$。为得到一个映射，我们使用
$$
M' = R\SheafHom_{\mathcal{O}'}(\mathcal{O}', M') \to
R\SheafHom_\mathcal{O}(r(\mathcal{O}'), r(M')) \to
R\SheafHom_\mathcal{O}(r(K'), r(M')) = E'
$$
，其中第二个箭头使用映射 $K' \to \mathcal{O}'$。
要看出这是同构，只需说明对该箭头应用 $r$ 后，所得映射就是上面的同构
$E \to r(E')$。细节从略。
\end{proof}

\begin{lemma}
\label{algebraization-lemma-pushforward-derived-complete-adjoint}
设 $f : (\Sh(\mathcal{D}), \mathcal{O}') \to (\Sh(\mathcal{C}), \mathcal{O})$
是环化拓扑斯的态射。设 $\mathcal{I} \subset \mathcal{O}$
和 $\mathcal{I}' \subset \mathcal{O}'$ 是有限型理想层，并且
$f^\sharp$ 把 $f^{-1}\mathcal{I}$ 映入 $\mathcal{I}'$。
则 $Rf_*$ 把 $D_{comp}(\mathcal{O}', \mathcal{I}')$
映入 $D_{comp}(\mathcal{O}, \mathcal{I})$，并且有左伴随
$Lf_{comp}^*$；该左伴随由先应用 $Lf^*$ 再作导出完备化得到。
\end{lemma}

\begin{proof}
第一条陈述已见于引理 \ref{algebraization-lemma-pushforward-derived-complete}。
第二条陈述是有意义的，因为我们有导出完备化函子
$D(\mathcal{O}') \to D_{comp}(\mathcal{O}', \mathcal{I}')$；
这是命题 \ref{algebraization-proposition-derived-completion} 的结论。
现在设 $K \in D_{comp}(\mathcal{O}, \mathcal{I})$ 且
$M \in D_{comp}(\mathcal{O}', \mathcal{I}')$。则
$$
\Hom(K, Rf_*M) = \Hom(Lf^*K, M) = \Hom(Lf_{comp}^*K, M)
$$
，其中等式来自导出完备化的泛性质。
\end{proof}

\begin{lemma}
\label{algebraization-lemma-pushforward-commutes-with-derived-completion}
\begin{reference}
这是 \cite[Lemma 6.5.9 (2)]{BS} 的推广。在拟凝聚模与（导出）代数栈
的态射这一背景下，可与 \cite[Theorem 6.5]{HL-P} 比较。
\end{reference}
设 $f : (\Sh(\mathcal{D}), \mathcal{O}') \to (\Sh(\mathcal{C}), \mathcal{O})$
是环化拓扑斯的态射。设 $\mathcal{I} \subset \mathcal{O}$
是一个有限型理想层，并设 $\mathcal{I}' \subset \mathcal{O}'$
为由 $f^\sharp(f^{-1}\mathcal{I})$ 生成的理想。则
$Rf_*$ 与导出完备化可交换，即
$Rf_*(K^\wedge) = (Rf_*K)^\wedge$.
\end{lemma}

\begin{proof}
由命题 \ref{algebraization-proposition-derived-completion}，导出完备化函子存在。
由引理 \ref{algebraization-lemma-pushforward-derived-complete}，对象
$Rf_*(K^\wedge)$ 是导出完备的，因此由导出完备化的泛性质得到典范映射
$(Rf_*K)^\wedge \to Rf_*(K^\wedge)$。可以在 $\mathcal{C}$
上局部地检验这个映射是同构。由于导出完备化与局部化可交换
（备注 \ref{algebraization-remark-localization-and-completion}），可以假设
$\mathcal{I}$ 由全局截面 $f_1, \ldots, f_r$ 生成。
于是 $\mathcal{I}'$ 由 $g_i = f^\sharp(f_i)$ 生成。由引理
\ref{algebraization-lemma-derived-completion-koszul}，只需证明
$$
R\lim \left(
Rf_*K \otimes^\mathbf{L}_\mathcal{O} K(\mathcal{O}, f_1^n, \ldots, f_r^n)
\right)
=
Rf_*\left(
R\lim
K \otimes^\mathbf{L}_{\mathcal{O}'} K(\mathcal{O}', g_1^n, \ldots, g_r^n)
\right)
$$
由于 $Rf_*$ 与 $R\lim$ 可交换
（《位点上的上同调》，引理
\ref{sites-cohomology-lemma-Rf-commutes-with-Rlim}），只需证明
$$
Rf_*K \otimes^\mathbf{L}_\mathcal{O} K(\mathcal{O}, f_1^n, \ldots, f_r^n) =
Rf_*\left(
K \otimes^\mathbf{L}_{\mathcal{O}'} K(\mathcal{O}', g_1^n, \ldots, g_r^n)
\right)
$$
这由投影公式（《位点上的上同调》，引理
\ref{sites-cohomology-lemma-projection-formula}）以及等式
$Lf^*K(\mathcal{O}, f_1^n, \ldots, f_r^n) =
K(\mathcal{O}', g_1^n, \ldots, g_r^n)$ 得出。
\end{proof}

\begin{lemma}
\label{algebraization-lemma-formal-functions-general}
设 $A$ 是一个环，$I \subset A$ 是有限生成理想。
设 $\mathcal{C}$ 是一个位点，$\mathcal{O}$ 是 $A$-代数层。
设 $\mathcal{F}$ 是一个 $\mathcal{O}$-模层。则
$$
R\Gamma(\mathcal{C}, \mathcal{F})^\wedge =
R\Gamma(\mathcal{C}, \mathcal{F}^\wedge)
$$
在 $D(A)$ 中成立，其中 $\mathcal{F}^\wedge$ 是 $\mathcal{F}$
关于 $I\mathcal{O}$ 的导出完备化，而左端取关于 $I$ 的导出完备化。
这给出两个谱序列
$$
E_2^{i, j} = H^i(H^j(\mathcal{C}, \mathcal{F})^\wedge)
\quad\text{且}\quad
E_2^{p, q} = H^p(\mathcal{C}, H^q(\mathcal{F}^\wedge))
$$
，二者都收敛到
$H^*(R\Gamma(\mathcal{C}, \mathcal{F})^\wedge) =
H^*(\mathcal{C}, \mathcal{F}^\wedge)$
\end{lemma}

\begin{proof}
把引理 \ref{algebraization-lemma-pushforward-commutes-with-derived-completion}
应用于环化拓扑斯的态射 $(\mathcal{C}, \mathcal{O}) \to (pt, A)$，
再取上同调，即得第一条陈述。第二个谱序列是《导出范畴》，引理
\ref{derived-lemma-two-ss-complex-functor} 的第二个谱序列。
第一个谱序列则是把《更多代数》，例
\ref{more-algebra-example-derived-completion-spectral-sequence}
中的谱序列应用于 $R\Gamma(\mathcal{C}, \mathcal{F})^\wedge$ 所得。
\end{proof}

\begin{remark}
\label{algebraization-remark-local-calculation-derived-completion}
设 $(\mathcal{C}, \mathcal{O})$ 是一个环化位点。
设 $\mathcal{I} \subset \mathcal{O}$ 是一个有限型理想层。
设 $K \mapsto K^\wedge$ 为命题
\ref{algebraization-proposition-derived-completion} 的导出完备化。
设 $U \in \Ob(\mathcal{C})$ 是一个对象，并且作为理想层，
$\mathcal{I}$ 由 $f_1, \ldots, f_r \in \mathcal{I}(U)$ 生成。
置 $A = \mathcal{O}(U)$ 和 $I = (f_1, \ldots, f_r) \subset A$。
警告：未必有 $I = \mathcal{I}(U)$。于是
$$
R\Gamma(U, K^\wedge) = R\Gamma(U, K)^\wedge
$$
，其中右端是对象 $R\Gamma(U, K)$（它属于 $D(A)$）关于 $I$
的导出完备化。这是因为导出完备化与局部化可交换
（备注 \ref{algebraization-remark-localization-and-completion}），并且可应用引理
\ref{algebraization-lemma-formal-functions-general}。
\end{remark}








\section{形式函数定理}
\label{algebraization-section-formal-functions}

\noindent
我们暂时中断正文的推进，讨论概形上拟凝聚模背景中的导出完备化，
并用它给出形式函数定理的一个略有不同的证明。备注
\ref{algebraization-remark-references} 中列出了一些文献线索。

\medskip\noindent
引理 \ref{algebraization-lemma-pushforward-commutes-with-derived-completion}
是形式函数定理的一个（非常形式化的）导出版本
（《概形上同调》，定理 \ref{coherent-theorem-formal-functions}）。
更明确地说，设 $f : X \to S$ 是概形的态射，
$\mathcal{I} \subset \mathcal{O}_S$ 是有限型拟凝聚理想层，
并且 $\mathcal{F}$ 是 $X$ 上的拟凝聚层。则该引理断言
\begin{equation}
\label{algebraization-equation-formal-functions}
Rf_*(\mathcal{F}^\wedge) = (Rf_*\mathcal{F})^\wedge
\end{equation}
，其中 $\mathcal{F}^\wedge$ 是 $\mathcal{F}$ 关于
$f^{-1}\mathcal{I} \cdot \mathcal{O}_X$ 的导出完备化，右端则是
$Rf_*\mathcal{F}$ 关于 $\mathcal{I}$ 的导出完备化。要看出这确实还原了
形式函数定理，还需要做一些工作。

\begin{lemma}
\label{algebraization-lemma-sections-derived-completion-pseudo-coherent}
设 $X$ 是局部诺特概形。设 $\mathcal{I} \subset \mathcal{O}_X$
是拟凝聚理想层。设 $K$ 是 $D(\mathcal{O}_X)$ 的伪凝聚对象，
其导出完备化为 $K^\wedge$。则
$$
H^p(U, K^\wedge) = \lim H^p(U, K)/I^nH^p(U, K) =
H^p(U, K)^\wedge
$$
对任意仿射开集 $U \subset X$ 成立，其中 $I = \mathcal{I}(U)$，
而右端取关于 $I$ 的导出完备化。
\end{lemma}

\begin{proof}
写成 $U = \Spec(A)$。环 $A$ 是诺特的，因此 $I \subset A$
是有限生成理想。于是
$$
R\Gamma(U, K^\wedge) = R\Gamma(U, K)^\wedge
$$
，这是备注 \ref{algebraization-remark-local-calculation-derived-completion} 的结论。
现在，$R\Gamma(U, K)$ 是一个伪凝聚的 $A$-模复形
（《概形的导出范畴》，引理
\ref{perfect-lemma-pseudo-coherent-affine}）。
由《更多代数》，引理
\ref{more-algebra-lemma-derived-completion-pseudo-coherent}，
其第 $p$ 个上同调模（这里的复形是 $R\Gamma(U, K^\wedge)$）
等于相应模的 $I$-进完备化，也就是 $H^p(U, K)$ 的完备化。
这证明了第一个等式。
第二个（较不重要的）等式则由再次应用刚才的引理立即得到。
\end{proof}

\begin{lemma}
\label{algebraization-lemma-derived-completion-pseudo-coherent}
设 $X$ 是局部诺特概形。设 $\mathcal{I} \subset \mathcal{O}_X$
是拟凝聚理想层。设 $K$ 是 $D(\mathcal{O}_X)$ 的对象。则
\begin{enumerate}
\item 导出完备化 $K^\wedge$ 等于
$R\lim (K \otimes_{\mathcal{O}_X}^\mathbf{L} \mathcal{O}_X/\mathcal{I}^n)$.
\end{enumerate}
设 $K$ 是 $D(\mathcal{O}_X)$ 的伪凝聚对象。则
\begin{enumerate}
\item[(2)] 上同调层 $H^q(K^\wedge)$ 等于
$\lim H^q(K)/\mathcal{I}^nH^q(K)$.
\end{enumerate}
设 $\mathcal{F}$ 是凝聚 $\mathcal{O}_X$-模\footnote{例如可取
$H^q(K)$，其中 $K$ 是所讨论局部诺特概形 $X$ 上的伪凝聚对象。}。则
\begin{enumerate}
\item[(3)] 导出完备化 $\mathcal{F}^\wedge$ 等于
$\lim \mathcal{F}/\mathcal{I}^n\mathcal{F}$,
\item[(4)]
$\lim \mathcal{F}/\mathcal{I}^n \mathcal{F} =
R\lim \mathcal{F}/\mathcal{I}^n \mathcal{F}$,
\item[(5)] $H^p(U, \mathcal{F}^\wedge) = 0$ 对每个 $p \not = 0$
以及每个仿射开集 $U \subset X$ 都成立。
\end{enumerate}
\end{lemma}

\begin{proof}
(1) 的证明。存在典范映射
$$
K \longrightarrow
R\lim (K \otimes_{\mathcal{O}_X}^\mathbf{L} \mathcal{O}_X/\mathcal{I}^n),
$$
；见备注 \ref{algebraization-remark-compare-with-completion}。
导出完备化与取开子概形可交换
（备注 \ref{algebraization-remark-localization-and-completion}）。
$R\lim$ 的构造也与取开子概形可交换。因此，为检验上述映射是同构，
可以局部地工作。于是可假设 $X = U = \Spec(A)$。
写 $I = (f_1, \ldots, f_r)$，并设
$K_n = K(A, f_1^n, \ldots, f_r^n)$ 为 Koszul 复形。
由《更多代数》，引理 \ref{more-algebra-lemma-sequence-Koszul-complexes}，
前系统 $\{K_n\}$ 与 $\{A/I^n\}$ 在 $D(A)$ 中同构。
利用等价 $D(A) = D_{\QCoh}(\mathcal{O}_X)$
（《概形的导出范畴》，引理
\ref{perfect-lemma-affine-compare-bounded}），可知前系统
$\{K(\mathcal{O}_X, f_1^n, \ldots, f_r^n)\}$
与 $\{\mathcal{O}_X/\mathcal{I}^n\}$ 在
$D(\mathcal{O}_X)$ 中同构。这证明了下式中的第二个等式：
$$
K^\wedge = R\lim \left(
K \otimes_{\mathcal{O}_X}^\mathbf{L} K(\mathcal{O}_X, f_1^n, \ldots, f_r^n)
\right) =
R\lim (K \otimes_{\mathcal{O}_X}^\mathbf{L} \mathcal{O}_X/\mathcal{I}^n)
$$
第一个等式则是引理 \ref{algebraization-lemma-derived-completion-koszul}。

\medskip\noindent
假设 $K$ 伪凝聚。对仿射开集 $U \subset X$，由
$H^q(U, K^\wedge) = \lim H^q(U, K)/\mathcal{I}^n(U)H^q(U, K)$
和引理 \ref{algebraization-lemma-sections-derived-completion-pseudo-coherent} 可知该等式成立。
由于这对每个 $U$ 都成立，作为层有
$H^q(K^\wedge) = \lim H^q(K)/\mathcal{I}^nH^q(K)$。这证明了 (2)。

\medskip\noindent
(3) 是 (2) 的特殊情形。(4) 与 (5) 由《概形的导出范畴》，引理
\ref{perfect-lemma-Rlim-quasi-coherent} 得出。
\end{proof}

\begin{lemma}
\label{algebraization-lemma-formal-functions}
设 $A$ 是诺特环，$I \subset A$ 是理想。设 $X$ 是
$A$ 上的诺特概形，$\mathcal{F}$ 是凝聚 $\mathcal{O}_X$-模。
假设 $H^p(X, \mathcal{F})$ 是有限 $A$-模，并且对所有 $p$ 都如此。
则有短正合列
$$
0 \to R^1\lim H^{p - 1}(X, \mathcal{F}/I^n\mathcal{F}) \to
H^p(X, \mathcal{F})^\wedge \to \lim H^p(X, \mathcal{F}/I^n\mathcal{F}) \to 0
$$
，它们是 $A$-模的短正合列，其中 $H^p(X, \mathcal{F})^\wedge$
是通常的 $I$-进完备化。若 $f$ 是固有的，则 $R^1\lim$ 项为零。
\end{lemma}

\begin{proof}
考察引理 \ref{algebraization-lemma-formal-functions-general} 的两个谱序列。
由《更多代数》，引理
\ref{more-algebra-lemma-derived-completion-pseudo-coherent}，第一个谱序列退化。
得到 $H^p(X, \mathcal{F})^\wedge$，它位于次数 $p$。
这里用到了 $H^p(X, \mathcal{F})$ 是有限 $A$-模这一假设。
第二个谱序列退化，因为
$$
\mathcal{F}^\wedge = \lim \mathcal{F}/I^n\mathcal{F} =
R\lim \mathcal{F}/I^n\mathcal{F}
$$
由引理 \ref{algebraization-lemma-derived-completion-pseudo-coherent} 可知这是一个层。
得到 $H^p(X, \lim \mathcal{F}/I^n\mathcal{F})$，它位于次数 $p$。
由于 $R\Gamma(X, -)$ 与导出极限可交换
（《内射对象》，引理 \ref{injectives-lemma-RF-commutes-with-Rlim}），还有
$$
R\Gamma(X, \lim \mathcal{F}/I^n\mathcal{F}) =
R\Gamma(X, R\lim \mathcal{F}/I^n\mathcal{F}) =
R\lim R\Gamma(X, \mathcal{F}/I^n\mathcal{F})
$$
由《更多代数》，备注
\ref{more-algebra-remark-how-unique}，得到正合列
$$
0 \to
R^1\lim H^{p - 1}(X, \mathcal{F}/I^n\mathcal{F}) \to
H^p(X, \lim \mathcal{F}/I^n\mathcal{F}) \to
\lim H^p(X, \mathcal{F}/I^n\mathcal{F}) \to 0
$$
，它们是 $A$-模的正合列。结合以上结果，得到引理的第一条陈述。
$R^1\lim$ 项的消失由《概形上同调》，引理
\ref{coherent-lemma-ML-cohomology-powers-ideal} 得出。
\end{proof}

\begin{remark}
\label{algebraization-remark-references}
下面列出文献中讨论相关材料的一些参考文献。
对于固有映射，“导出形式函数定理”似乎可追溯到
\cite[Theorem 6.3.1]{lurie-thesis}。
\cite{dag12} 中也有讨论，尤其是第 4 章对仿射情形的讨论；另见《更多代数》，
\ref{more-algebra-section-derived-completion} 节。
\cite[Section 2.9]{G-R} 使用（ind）凝聚模讨论了固有基变换与导出完备化。
拟凝聚模复形的一个与（\ref{algebraization-equation-formal-functions}）类似的结果见
\cite[Theorem 6.5]{HL-P}。
\end{remark}







\section{局部上同调的代数化，I}
\label{algebraization-section-algebraization-sections-general}

\noindent
设 $A$ 是诺特环，$I$ 与 $J$ 是 $A$ 的两个理想。
设 $M$ 是有限 $A$-模。本节研究对象
$$
R\Gamma_J(M)^\wedge
\quad\text{属于}\quad
D(A)
$$
的上同调群，其中 ${}^\wedge$ 表示导出 $I$-进完备化。注意，在
《对偶化复形》，引理 \ref{dualizing-lemma-completion-local-H0} 中，
我们已经证明：若 $A$ 关于 $I$ 完备，则存在同构
$$
\colim H^0_Z(M) \longrightarrow H^0(R\Gamma_J(M)^\wedge)
$$
，其中（有向）余极限遍历闭子集 $Z = V(J')$，它们满足
$J' \subset J$ 和 $V(J') \cap V(I) = V(J) \cap V(I)$。
这些闭子集的并为
\begin{equation}
\label{algebraization-equation-associated-subset}
T = \{\mathfrak p \in \Spec(A) :
V(\mathfrak p) \cap V(I) \subset V(J) \cap V(I)\}
\end{equation}
这是 $\Spec(A)$ 的一个对特殊化稳定的子集。上述结果于是变为：
$$
H^0_T(M) \longrightarrow H^0(R\Gamma_J(M)^\wedge)
$$
若 $A$ 关于 $I$ 完备，则这个映射是同构；见《局部上同调》，引理
\ref{local-cohomology-lemma-adjoint-ext} 与备注
\ref{local-cohomology-remark-upshot}。
我们把这一同构推广到更高上同调群的方法基于下面的引理。

\begin{lemma}
\label{algebraization-lemma-kill-completion-general}
设 $I, J$ 是诺特环 $A$ 的理想。设 $M$ 是有限 $A$-模。
设 $\mathfrak p \subset A$ 是素理想，$s$ 与 $d$ 是整数。假设
\begin{enumerate}
\item $A$ 有对偶化复形，
\item $\mathfrak p \not \in V(J) \cap V(I)$,
\item $\text{cd}(A, I) \leq d$，且
\item 对所有素理想 $\mathfrak p' \subset \mathfrak p$，都有
$\text{depth}_{A_{\mathfrak p'}}(M_{\mathfrak p'}) +
\dim((A/\mathfrak p')_\mathfrak q) > d + s$
，其中 $\mathfrak q \in V(\mathfrak p') \cap V(J) \cap V(I)$ 任意。
\end{enumerate}
则存在 $f \in A$，$f \not \in \mathfrak p$，它零化
$H^i(R\Gamma_J(M)^\wedge)$，其中 $i \leq s$；这里 ${}^\wedge$
表示 $I$-进完备化。
\end{lemma}

\begin{proof}
我们将使用 $R\Gamma_J = R\Gamma_{V(J)}$，并对 $I + J$ 作同样处理；
见《对偶化复形》，引理
\ref{dualizing-lemma-local-cohomology-noetherian}。注意
$R\Gamma_J(M)^\wedge = R\Gamma_I(R\Gamma_J(M))^\wedge =
R\Gamma_{I + J}(M)^\wedge$；见
《对偶化复形》，引理
\ref{dualizing-lemma-complete-and-local} 与
\ref{dualizing-lemma-local-cohomology-ss}.
因此可以把 $J$ 替换为 $I + J$，并假设 $I \subset J$
以及 $\mathfrak p \not \in V(J)$。回忆
$$
R\Gamma_J(M)^\wedge = R\Hom_A(R\Gamma_I(A), R\Gamma_J(M))
$$
，这是把《更多代数》，引理
\ref{more-algebra-lemma-derived-completion} 中对导出完备化的描述
与《对偶化复形》，引理
\ref{dualizing-lemma-compute-local-cohomology-noetherian}
中对局部上同调的描述结合起来所得。假设 (3) 意味着
$R\Gamma_I(A)$ 仅在次数 $\leq d$ 上有非零上同调。
利用 $R\Gamma_I(A)$ 的典范截断，只需证明
$$
\text{Ext}^i(N, R\Gamma_J(M))
$$
被某个 $f \in A$，$f \not \in \mathfrak p$ 零化，其中
$i \leq s + d$，而所述结论对任意 $A$-模 $N$ 都成立。
再利用 $R\Gamma_J(M)$ 的典范截断，只需证明
$H^i_J(M)$ 被某个 $f \in A$，$f \not \in \mathfrak p$ 零化，
其中 $i \leq s + d$。这由《局部上同调》，引理
\ref{local-cohomology-lemma-kill-local-cohomology-at-prime} 得出。
\end{proof}

\begin{lemma}
\label{algebraization-lemma-kill-colimit-weak-general}
设 $I, J$ 是某个诺特环的理想。设 $M$ 是有限 $A$-模。
设 $s$ 与 $d$ 是整数。令 $T$ 如
（\ref{algebraization-equation-associated-subset}）中所定义，并假设
\begin{enumerate}
\item $A$ 有对偶化复形，
\item 若 $\mathfrak p \in V(I)$，则不加条件，
\item 若 $\mathfrak p \not \in V(I)$ 且 $\mathfrak p \in T$，则
$\dim((A/\mathfrak p)_\mathfrak q) \leq d$，其中存在
某个 $\mathfrak q \in V(\mathfrak p) \cap V(J) \cap V(I)$，
\item 若 $\mathfrak p \not \in V(I)$ 且 $\mathfrak p \not \in T$，则
$$
\text{depth}_{A_\mathfrak p}(M_\mathfrak p) \geq s
\quad\text{或}\quad
\text{depth}_{A_\mathfrak p}(M_\mathfrak p) +
\dim((A/\mathfrak p)_\mathfrak q) > d + s
$$
对所有 $\mathfrak q \in V(\mathfrak p) \cap V(J) \cap V(I)$ 成立。
\end{enumerate}
则存在理想 $J_0 \subset J$，满足
$V(J_0) \cap V(I) = V(J) \cap V(I)$，使得对任意 $J' \subset J_0$，若
$V(J') \cap V(I) = V(J) \cap V(I)$，则映射
$$
R\Gamma_{J'}(M) \longrightarrow R\Gamma_{J_0}(M)
$$
在次数 $\leq s$ 的上同调上诱导同构，并且这些模被 $J_0I$ 的某个幂零化。
\end{lemma}

\begin{proof}
考虑集合
$$
B = \{\mathfrak p \not \in V(I),\ \mathfrak p \in T,\text{ 且 }
\text{depth}(M_\mathfrak p) \leq s\}
$$
选择 $J_0 \subset J$，使 $V(J_0)$ 是 $B \cup V(J)$ 的闭包。

\medskip\noindent
断言 I：$V(J_0) \cap V(I) = V(J) \cap V(I)$。

\medskip\noindent
断言 I 的证明。包含关系 $\supset$ 由构造成立。
设 $\mathfrak p$ 是 $V(J_0)$ 的极小素理想。
若 $\mathfrak p \in B \cup V(J)$，则或者 $\mathfrak p \in T$，
或者 $\mathfrak p \in V(J)$；两种情形下都如所需有
$V(\mathfrak p) \cap V(I) \subset V(J) \cap V(I)$。
若 $\mathfrak p \not \in B \cup V(J)$，则
$V(\mathfrak p) \cap B$ 稠密，因而无限，可知
$\text{depth}(M_\mathfrak p) < s$；见《局部上同调》，引理
\ref{local-cohomology-lemma-depth-function}。
事实上，设
$V(\mathfrak p) \cap B = \{\mathfrak p_\lambda\}_{\lambda \in \Lambda}$.
依照 (3)，选择
$\mathfrak q_\lambda \in V(\mathfrak p_\lambda) \cap V(J) \cap V(I)$。
设 $\delta : \Spec(A) \to \mathbf{Z}$ 是与对偶化复形
$\omega_A^\bullet$（定义在 $A$ 上）相伴的维数函数。由于 $\Lambda$ 无限而 $\delta$ 有界，
存在无限子集 $\Lambda' \subset \Lambda$，使得
$\delta(\mathfrak q_\lambda)$ 在其上为常数。对
$\lambda \in \Lambda'$，有
$$
\text{depth}(M_{\mathfrak p_\lambda}) +
\delta(\mathfrak p_\lambda) - \delta(\mathfrak q_\lambda) =
\text{depth}(M_{\mathfrak p_\lambda}) +
\dim((A/\mathfrak p_\lambda)_{\mathfrak q_\lambda})
\leq d + s
$$
，这是 (3) 与 $B$ 的定义的结论。利用函数
$\text{depth} + \delta$ 的半连续性（见《概形的对偶性》，引理
\ref{duality-lemma-sitting-in-degrees}），可得
$$
\text{depth}(M_\mathfrak p) +
\dim((A/\mathfrak p)_{\mathfrak q_\lambda}) =
\text{depth}(M_\mathfrak p) + \delta(\mathfrak p) - \delta(\mathfrak q_\lambda)
\leq d + s
$$
又因为 $\mathfrak p \not \in V(I)$，由 (4) 可知
$\mathfrak p \in T$，即
$V(\mathfrak p) \cap V(I) \subset V(J) \cap V(I)$。断言 I 得证。

\medskip\noindent
断言 II：映射 $H^i_{J_0}(M) \to H^i_J(M)$ 是同构，其中
$i \leq s$，并且 $J' \subset J_0$ 满足
$V(J') \cap V(I) = V(J) \cap V(I)$。

\medskip\noindent
断言 II 的证明。选择 $\mathfrak p \in V(J')$，使其不属于 $V(J_0)$。
只需证明 $H^i_{\mathfrak pA_\mathfrak p}(M_\mathfrak p) = 0$
对 $i \leq s$ 成立；见《局部上同调》，引理
\ref{local-cohomology-lemma-isomorphism}。
注意 $\mathfrak p \in T$。因此，由于 $\mathfrak p$ 不属于 $B$，
可知 $\text{depth}(M_\mathfrak p) > s$，而这些群由《对偶化复形》，引理
\ref{dualizing-lemma-depth} 消失。

\medskip\noindent
断言 III：引理的最后一条陈述成立。

\medskip\noindent
由断言 II，对 $i \leq s$ 有
$$
H^i_T(M) = H^i_{J_0}(M) = H^i_{J'}(M)
$$
，其中 $J' \subset J_0$ 是任意满足
$V(J')  \cap V(I) = V(J) \cap V(I)$ 的理想。见《局部上同调》，引理
\ref{local-cohomology-lemma-adjoint-ext}。
现在检验《局部上同调》，命题
\ref{local-cohomology-proposition-annihilator} 对子集
$T \subset T \cup V(I)$、模 $M$ 和整数 $s$ 的假设。
须证明：给定 $\mathfrak p \subset \mathfrak q$，若
$\mathfrak p \not \in T \cup V(I)$ 且 $\mathfrak q \in T$，则
$$
\text{depth}_{A_\mathfrak p}(M_\mathfrak p) +
\dim((A/\mathfrak p)_\mathfrak q) > s
$$
若 $\text{depth}(M_\mathfrak p) \geq s$，则结论成立，因为
$(A/\mathfrak p)_\mathfrak q$ 的维数至少为 $1$。
因此可假设 $\text{depth}(M_\mathfrak p) < s$。
若 $\mathfrak q \in V(I)$，则 $\mathfrak q \in V(J) \cap V(I)$，
不等式由 (4) 成立。若 $\mathfrak q \not \in V(I)$，则由 (3)
可以选择 $\mathfrak q' \in V(\mathfrak q) \cap V(J) \cap V(I)$，使
$\dim((A/\mathfrak q)_{\mathfrak q'}) \leq d$。于是由假设 (4) 得
$$
\text{depth}_{A_\mathfrak p}(M_\mathfrak p) +
\dim((A/\mathfrak p)_{\mathfrak q'}) > s + d
$$
由于 $A$ 是链环，这蕴含所需不等式。应用《局部上同调》，命题
\ref{local-cohomology-proposition-annihilator}，得到
$J'' \subset A$，满足 $V(J'') \subset T \cup V(I)$，并且
$J''$ 零化 $H^i_T(M)$，其中 $i \leq s$。
于是可以写成 $V(J'') \cup V(J_0) \cup V(I) = V(J'I)$，
其中某个 $J' \subset J_0$ 满足
$V(J') \cap V(I) = V(J) \cap V(I)$。
把 $J_0$ 替换为 $J'$，证明完成。
\end{proof}

\begin{lemma}
\label{algebraization-lemma-kill-colimit-general}
在引理 \ref{algebraization-lemma-kill-colimit-weak-general} 中，若把空条件 (2)
替换为假设
\begin{enumerate}
\item[(2')] 若 $\mathfrak p \in V(I)$ 且
$\mathfrak p \not \in V(J) \cap V(I)$，则
$\text{depth}_{A_\mathfrak p}(M_\mathfrak p) +
\dim((A/\mathfrak p)_\mathfrak q) > s$
对所有 $\mathfrak q \in V(\mathfrak p) \cap V(J) \cap V(I)$ 成立，
\end{enumerate}
则这些条件还蕴含 $H^i_{J_0}(M)$ 是有限 $A$-模，其中 $i \leq s$。
\end{lemma}

\begin{proof}
回忆 $H^i_{J_0}(M) = H^i_T(M)$；见引理
\ref{algebraization-lemma-kill-colimit-weak-general} 的证明。因此只需检验：若
$\mathfrak p \not \in T$、$\mathfrak q \in T$ 且
$\mathfrak p \subset \mathfrak q$，则
$\text{depth}_{A_\mathfrak p}(M_\mathfrak p) +
\dim((A/\mathfrak p)_\mathfrak q) > s$；见《局部上同调》，命题
\ref{local-cohomology-proposition-finiteness}。
条件 (2') 告诉我们，这对 $\mathfrak p \in V(I)$ 成立。
由于 $H^i_T(M)$ 被 $IJ_0$ 的某个幂零化，当
$\mathfrak p \not \in V(IJ_0)$ 时该条件也成立；这是《局部上同调》，命题
\ref{local-cohomology-proposition-annihilator} 的结论。
这涵盖所有情形，证明完成。
\end{proof}

\begin{lemma}
\label{algebraization-lemma-kill-colimit-support-general}
若在引理 \ref{algebraization-lemma-kill-colimit-weak-general} 中还假设
\begin{enumerate}
\item[(6)] 若 $\mathfrak p \not \in V(I)$ 且 $\mathfrak p \in T$，则
$\text{depth}_{A_\mathfrak p}(M_\mathfrak p) > s$,
\end{enumerate}
则 $H^i_{J_0}(M) = H^i_J(M) = H^i_{J + I}(M)$ 对 $i \leq s$ 成立，
并且这些模被 $I$ 的某个幂零化。
\end{lemma}

\begin{proof}
选择 $\mathfrak p \in V(J)$ 或 $\mathfrak p \in V(J_0)$，但要求
$\mathfrak p \not \in V(J + I) = V(J_0 + I)$。
只需证明 $H^i_{\mathfrak pA_\mathfrak p}(M_\mathfrak p) = 0$
对 $i \leq s$ 成立；见《局部上同调》，引理
\ref{local-cohomology-lemma-isomorphism}。
这些群由条件 (6) 和《对偶化复形》，引理
\ref{dualizing-lemma-depth} 消失。最后一条陈述由《局部上同调》，命题
\ref{local-cohomology-proposition-annihilator} 得出。
\end{proof}

\begin{lemma}
\label{algebraization-lemma-algebraize-local-cohomology-general}
设 $I, J$ 是诺特环 $A$ 的理想。设 $M$ 是有限 $A$-模。
设 $s$ 与 $d$ 是整数。令 $T$ 如
（\ref{algebraization-equation-associated-subset}）中所定义，并假设
\begin{enumerate}
\item $A$ 是 $I$-进完备的，并且有对偶化复形，
\item 若 $\mathfrak p \in V(I)$，则不加条件，
\item $\text{cd}(A, I) \leq d$,
\item 若 $\mathfrak p \not \in V(I)$ 且 $\mathfrak p \not \in T$，则
$$
\text{depth}_{A_\mathfrak p}(M_\mathfrak p) \geq s
\quad\text{或}\quad
\text{depth}_{A_\mathfrak p}(M_\mathfrak p) +
\dim((A/\mathfrak p)_\mathfrak q) > d + s
$$
对所有 $\mathfrak q \in V(\mathfrak p) \cap V(J) \cap V(I)$ 成立，
\item 若 $\mathfrak p \not \in V(I)$、$\mathfrak p \not \in T$、
$V(\mathfrak p) \cap V(J) \cap V(I) \not = \emptyset$，并且
$\text{depth}(M_\mathfrak p) < s$，则下列条件之一成立
\footnote{我们的方法迫使我们加入这个额外条件。稍后将回到这一点
（插入未来的引用）。}：
\begin{enumerate}
\item $\dim(\text{Supp}(M_\mathfrak p)) < s + 2$\footnote{例如，若
$M$ 满足 Serre 条件 $(S_s)$，并且这一条件是在
$V(I) \cup T$ 的补集上要求的，则如此。}，或者
\item  $\delta(\mathfrak p) > d + \delta_{max} - 1$
，其中 $\delta$ 是维数函数，而 $\delta_{max}$
是 $\delta$ 在 $V(J) \cap V(I)$ 上的最大值，或者
\item $\text{depth}_{A_\mathfrak p}(M_\mathfrak p) +
\dim((A/\mathfrak p)_\mathfrak q) > d + s + \delta_{max} - \delta_{min} - 2$
对所有 $\mathfrak q \in V(\mathfrak p) \cap V(J) \cap V(I)$ 成立。
\end{enumerate}
\end{enumerate}
则存在理想 $J_0 \subset J$，满足
$V(J_0) \cap V(I) = V(J) \cap V(I)$
，使得对任意 $J' \subset J_0$，若
$V(J') \cap V(I) = V(J) \cap V(I)$，则映射
$$
R\Gamma_{J'}(M) \longrightarrow R\Gamma_J(M)^\wedge
$$
在次数 $\leq s$ 的上同调上诱导同构。
这里 ${}^\wedge$ 表示导出 $I$-进完备化。
\end{lemma}

\noindent
我们建议读者先阅读局部情形的证明
（引理 \ref{algebraization-lemma-algebraize-local-cohomology}），因为它无需处理所有这些
不等式便能解释证明的结构。

\begin{proof}
对理想 $\mathfrak a \subset A$，有
$R\Gamma_\mathfrak a = R\Gamma_{V(\mathfrak a)}$；见《对偶化复形》，引理
\ref{dualizing-lemma-local-cohomology-noetherian}。其次，注意
$$
R\Gamma_J(M)^\wedge =
R\Gamma_I(R\Gamma_J(M))^\wedge =
R\Gamma_{I + J}(M)^\wedge =
R\Gamma_{I + J'}(M)^\wedge =
R\Gamma_I(R\Gamma_{J'}(M))^\wedge =
R\Gamma_{J'}(M)^\wedge
$$
，这是《对偶化复形》，引理
\ref{dualizing-lemma-local-cohomology-ss} 与
\ref{dualizing-lemma-complete-and-local} 的结论。
这说明了如何定义引理陈述中的箭头。

\medskip\noindent
我们断言：引理 \ref{algebraization-lemma-kill-colimit-weak-general} 的假设
由当前对 $M$ 的假设蕴含。唯一需要检验的是假设 (3)。
因此设 $\mathfrak p \not \in V(I)$ 且 $\mathfrak p \in T$。
那么 $V(\mathfrak p) \cap V(I)$ 非空，因为 $I$ 包含在
$A$ 的 Jacobson 根中
（《代数》，引理 \ref{algebra-lemma-radical-completion}）。
由于 $\mathfrak p \in T$，有
$V(\mathfrak p) \cap V(I) = V(\mathfrak p) \cap V(J) \cap V(I)$。
设 $\mathfrak q \in V(\mathfrak p) \cap V(I)$ 是某个不可约分支的通用点。
有 $\text{cd}(A_\mathfrak q, I_\mathfrak q) \leq d$；这是《局部上同调》，引理
\ref{local-cohomology-lemma-cd-local} 的结论。
又有 $V(\mathfrak pA_\mathfrak q) \cap V(I_\mathfrak q) =
\{\mathfrak qA_\mathfrak q\}$，这是由 $\mathfrak q$ 的选择得到的；进而
$\dim((A/\mathfrak p)_\mathfrak q) \leq d$，见《局部上同调》，引理
\ref{local-cohomology-lemma-cd-bound-dim-local}。

\medskip\noindent
注意，引理对 $s < 0$ 成立。这并非平凡情形，因为不能先验地看出
$H^i(R\Gamma_J(M)^\wedge)$ 在 $i < 0$ 时为零。不过，这一消失已在
《对偶化复形》，引理 \ref{dualizing-lemma-completion-local} 中证明。
对于 $s \geq 0$，我们将用归纳法证明引理。

\medskip\noindent
$s = 0$ 时的引理由引理 \ref{algebraization-lemma-kill-colimit-weak-general}
的结论和《对偶化复形》，引理
\ref{dualizing-lemma-completion-local-H0} 立即得出。

\medskip\noindent
假设 $s > 0$，且引理对更小的 $s$ 已经证明。
设 $M' \subset M$ 是支集包含于 $V(I) \cup T$ 的极大子模。
则 $M'$ 是有限 $A$-模，并且其支集包含于
$V(J') \cup V(I)$，其中某个理想 $J' \subset J$ 满足
$V(J') \cap V(I) = V(J) \cap V(I)$。我们断言
$$
R\Gamma_{J'}(M') \to R\Gamma_J(M')^\wedge
$$
对 $J'$ 的任意选择都是同构。事实上，可以选择短正合列
$0 \to M_1 \oplus M_2 \to M' \to N \to 0$
，其中 $M_1$ 被 $J'$ 的某个幂零化，$M_2$ 被 $I$ 的某个幂零化，
而 $N$ 被 $I + J'$ 的某个幂零化。因此，只需对
$M_1$、$M_2$ 与 $N$ 证明该断言。对 $M_1$，有
$R\Gamma_{J'}(M_1) = M_1$；又因为 $M_1$ 是有限 $A$-模且
$I$-进完备，所以 $M_1^\wedge = M_1$。结合证明开头的说明，
这就对 $M_1$ 证明了断言。对 $M_2$，有
$H^i_J(M_2) = H^i_{I + J}(M) = H^i_{I + J'}(M) = H^i_{J'}(M_2)$
被 $I$ 的某个幂零化，因而导出完备。所以该情形中的断言也成立。
对 $N$，可以使用刚才两个论证中的任一个。考察短正合列
$0 \to M' \to M \to M/M' \to 0$
可知，只需对 $M/M'$ 证明引理。因此可假设
$\text{Ass}(M) \cap (V(I) \cup T) = \emptyset$。

\medskip\noindent
设 $\mathfrak p \in \text{Ass}(M)$ 满足
$V(\mathfrak p) \cap V(J) \cap V(I) = \emptyset$.
由于 $I$ 包含在 $A$ 的 Jacobson 根中，这蕴含
$V(\mathfrak p) \cap V(J') = \emptyset$ 对任意
$J' \subset J$ 都成立，只要 $V(J') \cap V(I) = V(J) \cap V(I)$。
置 $N = H^0_\mathfrak p(M)$，则
$R\Gamma_J(N) = R\Gamma_{J'}(N) = 0$ 对所有
$J' \subset J$ 都成立，只要 $V(J') \cap V(I) = V(J) \cap V(I)$。
特别地，$R\Gamma_J(N)^\wedge = 0$。
因此可以把 $M$ 替换为 $M/N$，因为这只在条件 (4) 或 (5)
不起作用的素理想处改变 $M$ 的结构。重复此过程，可假设
$V(\mathfrak p) \cap V(J) \cap V(I) \not = \emptyset$
对所有 $\mathfrak p \in \text{Ass}(M)$ 成立。

\medskip\noindent
假设 $\text{Ass}(M) \cap (V(I) \cup T) = \emptyset$，并且
$V(\mathfrak p) \cap V(J) \cap V(I) \not = \emptyset$
对所有 $\mathfrak p \in \text{Ass}(M)$ 成立。
设 $\mathfrak p \in \text{Ass}(M)$。我们要说明可以应用引理
\ref{algebraization-lemma-kill-completion-general}。正是在这一验证中，我们将使用补充条件 (5)。
选择 $\mathfrak p' \subset \mathfrak p$ 以及
$\mathfrak q' \subset V(\mathfrak p) \cap V(J) \cap V(I)$。
\begin{enumerate}
\item 若 $M_{\mathfrak p'} = 0$，则
$\text{depth}(M_{\mathfrak p'}) = \infty$，并且
$\text{depth}(M_{\mathfrak p'}) +
\dim((A/\mathfrak p')_{\mathfrak q'}) > d + s$.
\item 若 $\text{depth}(M_{\mathfrak p'}) < s$，则
$\text{depth}(M_{\mathfrak p'}) +
\dim((A/\mathfrak p')_{\mathfrak q'}) > d + s$，这是 (4) 的结论。
\end{enumerate}
在其余情形中，有 $M_{\mathfrak p'} \not = 0$ 且
$\text{depth}(M_{\mathfrak p'}) \geq s$。特别地，
$\mathfrak p'$ 属于 $M$ 的支集，并且可以选择
$\mathfrak p'' \subset \mathfrak p'$，使 $\mathfrak p'' \in \text{Ass}(M)$。
\begin{enumerate}
\item[(a)] 注意
$\dim((A/\mathfrak p'')_{\mathfrak p'}) \geq \text{depth}(M_{\mathfrak p'})$
，这是《代数》，引理 \ref{algebra-lemma-depth-dim-associated-primes} 的结论。
若等号成立，则
$$
\text{depth}(M_{\mathfrak p'}) + \dim((A/\mathfrak p')_{\mathfrak q'}) =
\text{depth}(M_{\mathfrak p''}) + \dim((A/\mathfrak p'')_{\mathfrak q'})
> s + d
$$
；这是把 (4) 应用于 $\mathfrak p''$ 所得，故这一情形完成。这意味着，
只有当
$\dim((A/\mathfrak p'')_{\mathfrak p'}) > \text{depth}(M_{\mathfrak p'})$.
时才有困难。这蕴含 $\dim(M_\mathfrak p) \geq s + 2$。
因此，若 (5)(a) 成立，就不会出现这种情形。
\item[(b)] 若 (5)(b) 成立，则
$$
\text{depth}(M_{\mathfrak p'}) + \dim((A/\mathfrak p')_{\mathfrak q'})
\geq s + \delta(\mathfrak p') - \delta(\mathfrak q')
\geq s + 1 + \delta(\mathfrak p) - \delta_{max}
> s + d
$$
，如所需。
\item[(c)] 若 (5)(c) 成立，则
\begin{align*}
\text{depth}(M_{\mathfrak p'}) + \dim((A/\mathfrak p')_{\mathfrak q'})
& \geq
s + \delta(\mathfrak p') - \delta(\mathfrak q') \\
& \geq
s + 1 + \delta(\mathfrak p) - \delta(\mathfrak q') \\
& =
s + 1 + \delta(\mathfrak p) - \delta(\mathfrak q) +
\delta(\mathfrak q) - \delta(\mathfrak q') \\
& >
s + 1 + (s + d + \delta_{max} - \delta_{min} - 2) +
\delta(\mathfrak q) - \delta(\mathfrak q') \\
& \geq 
2s + d - 1 \geq s + d
\end{align*}
，如所需。注意，这一论证成立是因为我们知道存在素理想
$\mathfrak q \in V(\mathfrak p) \cap V(J) \cap V(I)$。
\end{enumerate}
现在可以进行归纳步骤。

\medskip\noindent
选择理想 $J_0$，如引理 \ref{algebraization-lemma-kill-colimit-weak-general} 中所述，
并选择整数 $t > 0$，使 $(J_0I)^t$ 零化 $H^s_J(M)$。
引理 \ref{algebraization-lemma-kill-completion-general} 的假设对每个
$\mathfrak p \in \text{Ass}(M)$ 都成立（见上一段）。
因此，零化理想 $\mathfrak a \subset A$（它零化
$H^s(R\Gamma_J(M)^\wedge)$）不包含于 $\mathfrak p$，其中
$\mathfrak p \in \text{Ass}(M)$。于是可以找到
$f \in \mathfrak a(J_0I)^t$，它不属于 $M$ 的任何相伴素理想，
并且同时零化 $H^s(R\Gamma_J(M)^\wedge)$ 与 $H^s_J(M)$。
于是 $f$ 是 $M$ 上的非零因子，可以考察短正合列
$$
0 \to M \xrightarrow{f} M \to M/fM \to 0
$$
由 $f$ 的选择，得到
$$
\xymatrix{
H^{s - 1}_{J'}(M) \ar[d] \ar[r] &
H^{s - 1}_{J'}(M/fM) \ar[d] \ar[r] &
H^s_{J'}(M) \ar[d] \ar[r] & 0 \\
H^{s - 1}(R\Gamma_J(M)^\wedge) \ar[r] &
H^{s - 1}(R\Gamma_J(M/fM)^\wedge) \ar[r] &
H^s(R\Gamma_J(M)^\wedge) \ar[r] & 0
}
$$
，其中 $J' \subset J_0$ 任意，并满足
$V(J') \cap V(I) = V(J) \cap V(I)$。
因此，若选择 $J'$，使它对 $M$、$M/fM$ 和 $s - 1$ 都适用
（由归纳假设可以做到——见下一段），便可得出引理成立。

\medskip\noindent
为完成证明，须说明模 $M/fM$ 对 $s - 1$ 满足假设 (4) 与 (5)。
因此，设 $\mathfrak p$ 是 $M/fM$ 支集中的素理想，满足
$\text{depth}((M/fM)_\mathfrak p) < s - 1$，并且
$V(\mathfrak p) \cap V(J) \cap V(I)$ 非空。于是
$\dim(M_\mathfrak p) = \dim((M/fM)_\mathfrak p) + 1$，并且
$\text{depth}(M_\mathfrak p) = \text{depth}((M/fM)_\mathfrak p) + 1$。
特别地，我们知道 (4) 与 (5) 对 $\mathfrak p$ 和 $M$
在原来的值 $s$ 下成立。逐项检验即可得到所需不等式。
\end{proof}

\begin{example}
\label{algebraization-example-no-ML}
在引理 \ref{algebraization-lemma-algebraize-local-cohomology-general} 中，
我们不知道逆系统 $H^i_J(M/I^nM)$ 是否满足 Mittag--Leffler 条件。
例如，假设 $A = \mathbf{Z}_p[[x, y]]$、$I = (p)$、
$J = (p, x)$ 且 $M = A/(xy - p)$。则映射
$H^0_J(M/p^nM) \to H^0_J(M/pM)$
的像是由 $y^n$ 生成的理想；该理想位于
$M/pM = A/(p, xy)$ 中。
\end{example}





\section{局部上同调的代数化，II}
\label{algebraization-section-algebraization-punctured}

\noindent
本节重新给出 \ref{algebraization-section-algebraization-sections-general} 节的论证；
此时 $(A, \mathfrak m)$ 是局部环，并取局部上同调
$R\Gamma_\mathfrak m$，它关于 $\mathfrak m$ 定义。
与前面一样，主要工具是下面的引理。

\begin{lemma}
\label{algebraization-lemma-kill-completion}
设 $(A, \mathfrak m)$ 是诺特局部环。设 $I \subset A$ 是理想。
设 $M$ 是有限 $A$-模，$\mathfrak p \subset A$ 是素理想。
设 $s$ 与 $d$ 是整数。假设
\begin{enumerate}
\item $A$ 有对偶化复形，
\item $\text{cd}(A, I) \leq d$，并且
\item
$\text{depth}_{A_\mathfrak p}(M_\mathfrak p) + \dim(A/\mathfrak p) > d + s$.
\end{enumerate}
则存在 $f \in A \setminus \mathfrak p$，它零化
$H^i(R\Gamma_\mathfrak m(M)^\wedge)$，其中 $i \leq s$；这里 ${}^\wedge$
表示 $I$-进完备化。
\end{lemma}

\begin{proof}
根据《局部上同调》，引理
\ref{local-cohomology-lemma-sitting-in-degrees}，函数
$$
\mathfrak p' \longmapsto
\text{depth}_{A_{\mathfrak p'}}(M_{\mathfrak p'}) + \dim(A/\mathfrak p')
$$
在 $\Spec(A)$ 上下半连续。因此，该函数在
$\mathfrak p' \subset \mathfrak p$ 上的值为 $> s + d$。
所以，本引理是引理 \ref{algebraization-lemma-kill-completion-general} 的特殊情形，
只要 $\mathfrak p \not = \mathfrak m$。
若 $\mathfrak p = \mathfrak m$，则由深度与局部上同调的关系，
$H^i_\mathfrak m(M) = 0$ 对 $i \leq s + d$ 成立
（《对偶化复形》，引理 \ref{dualizing-lemma-depth}）。
因此，引理 \ref{algebraization-lemma-kill-completion-general} 证明中的论证表明，
在这个（退化）情形中 $H^i(R\Gamma_\mathfrak m(M)^\wedge) = 0$
对 $i \leq s$ 成立。
\end{proof}

\begin{lemma}
\label{algebraization-lemma-kill-colimit-weak}
设 $(A, \mathfrak m)$ 是诺特局部环。设 $I \subset A$ 是理想。
设 $M$ 是有限 $A$-模。设 $s$ 与 $d$ 是整数。假设
\begin{enumerate}
\item $A$ 有对偶化复形，
\item 若 $\mathfrak p \in V(I)$，则不加条件，
\item 若 $\mathfrak p \not \in V(I)$ 且
$V(\mathfrak p) \cap V(I) = \{\mathfrak m\}$，则
$\dim(A/\mathfrak p) \leq d$,
\item 若 $\mathfrak p \not \in V(I)$ 且
$V(\mathfrak p) \cap V(I) \not = \{\mathfrak m\}$，则
$$
\text{depth}_{A_\mathfrak p}(M_\mathfrak p) \geq s
\quad\text{或}\quad
\text{depth}_{A_\mathfrak p}(M_\mathfrak p) + \dim(A/\mathfrak p) > d + s
$$
\end{enumerate}
则存在理想 $J_0 \subset A$，满足
$V(J_0) \cap V(I) = \{\mathfrak m\}$，使得对任意 $J \subset J_0$，若
$V(J) \cap V(I) = \{\mathfrak m\}$，则映射
$$
R\Gamma_J(M) \longrightarrow R\Gamma_{J_0}(M)
$$
在次数 $\leq s$ 的上同调上诱导同构，并且这些模被 $J_0I$ 的某个幂零化。
\end{lemma}

\begin{proof}
这是引理 \ref{algebraization-lemma-kill-colimit-weak-general} 的特殊情形。
\end{proof}

\begin{lemma}
\label{algebraization-lemma-kill-colimit}
在引理 \ref{algebraization-lemma-kill-colimit-weak} 中，若把空条件 (2) 替换为假设
\begin{enumerate}
\item[(2')] 若 $\mathfrak p \in V(I)$ 且 $\mathfrak p \not = \mathfrak m$，
则 $\text{depth}_{A_\mathfrak p}(M_\mathfrak p) + \dim(A/\mathfrak p) > s$，
\end{enumerate}
则这些条件还蕴含 $H^i_{J_0}(M)$ 是有限 $A$-模，其中 $i \leq s$。
\end{lemma}

\begin{proof}
这是引理 \ref{algebraization-lemma-kill-colimit-general} 的特殊情形。
\end{proof}

\begin{lemma}
\label{algebraization-lemma-kill-colimit-support}
若在引理 \ref{algebraization-lemma-kill-colimit-weak} 中还假设
\begin{enumerate}
\item[(6)] 若 $\mathfrak p \not \in V(I)$ 且
$V(\mathfrak p) \cap V(I) = \{\mathfrak m\}$，则
$\text{depth}_{A_\mathfrak p}(M_\mathfrak p) > s$,
\end{enumerate}
则 $H^i_{J_0}(M) = H^i_J(M) = H^i_\mathfrak m(M)$ 对 $i \leq s$ 成立，
并且这些模被 $I$ 的某个幂零化。
\end{lemma}

\begin{proof}
这是引理 \ref{algebraization-lemma-kill-colimit-support-general} 的特殊情形。
\end{proof}

\begin{lemma}
\label{algebraization-lemma-algebraize-local-cohomology}
设 $(A, \mathfrak m)$ 是诺特局部环。设 $I \subset A$ 是理想。
设 $M$ 是有限 $A$-模。设 $s$ 与 $d$ 是整数。假设
\begin{enumerate}
\item $A$ 是 $I$-进完备的，并且有对偶化复形，
\item 若 $\mathfrak p \in V(I)$，则不加条件，
\item $\text{cd}(A, I) \leq d$,
\item 若 $\mathfrak p \not \in V(I)$ 且
$V(\mathfrak p) \cap V(I) \not = \{\mathfrak m\}$，则
$$
\text{depth}_{A_\mathfrak p}(M_\mathfrak p) \geq s
\quad\text{或}\quad
\text{depth}_{A_\mathfrak p}(M_\mathfrak p) + \dim(A/\mathfrak p) > d + s
$$
\end{enumerate}
则存在理想 $J_0 \subset A$，满足
$V(J_0) \cap V(I) = \{\mathfrak m\}$，使得对任意 $J \subset J_0$，若
$V(J) \cap V(I) = \{\mathfrak m\}$，则映射
$$
R\Gamma_J(M) \longrightarrow
R\Gamma_J(M)^\wedge = R\Gamma_\mathfrak m(M)^\wedge
$$
在次数 $\leq s$ 的上同调上诱导同构。
这里 ${}^\wedge$ 表示导出 $I$-进完备化。
\end{lemma}

\begin{proof}
本引理是引理 \ref{algebraization-lemma-algebraize-local-cohomology-general} 的特殊情形，
因为在 $\delta_{max} = \delta_{min} = \delta(\mathfrak m)$ 时，
条件 (4) 蕴含条件 (5)(c)。我们仍将给出这个重要特殊情形的证明，
因为它略为简单（需要检验的事项更少）。

\medskip\noindent
在当前情形中，$R\Gamma_\mathfrak a$ 与
$R\Gamma_{V(\mathfrak a)}$ 没有区别；见《对偶化复形》，引理
\ref{dualizing-lemma-local-cohomology-noetherian}。其次，注意
$$
R\Gamma_\mathfrak m(M)^\wedge =
R\Gamma_I(R\Gamma_J(M))^\wedge =
R\Gamma_J(M)^\wedge
$$
，这是《对偶化复形》，引理
\ref{dualizing-lemma-local-cohomology-ss} 与
\ref{dualizing-lemma-complete-and-local} 的结论，
也解释了引理陈述中的等号。

\medskip\noindent
注意，引理对 $s < 0$ 成立。这并非平凡情形，因为不能先验地看出
$H^s(R\Gamma_\mathfrak m(M)^\wedge)$ 在负的 $s$ 上为零。
不过，这一消失已在引理
\ref{algebraization-lemma-local-cohomology-derived-completion} 中证明。
对于 $s \geq 0$，我们将用归纳法证明引理。

\medskip\noindent
引理 \ref{algebraization-lemma-kill-colimit-weak} 的假设由《局部上同调》，引理
\ref{local-cohomology-lemma-cd-bound-dim-local} 保证。
$s = 0$ 时的引理由引理 \ref{algebraization-lemma-kill-colimit-weak} 和《对偶化复形》，
引理 \ref{dualizing-lemma-completion-local-H0} 得出。

\medskip\noindent
假设 $s > 0$，且引理对更小的 $s$ 成立。设 $M' \subset M$
是由下述元素组成的子模：其支集包含于 $V(I) \cup V(J)$，
其中某个理想 $J$ 满足 $V(J) \cap V(I) = \{\mathfrak m\}$。
则 $M'$ 是有限 $A$-模。我们断言
$$
R\Gamma_J(M') \to R\Gamma_\mathfrak m(M')^\wedge
$$
对 $J$ 的任意选择都是同构。事实上，对任意这样的模，都有短正合列
$0 \to M_1 \oplus M_2 \to M' \to N \to 0$，其中
$M_1$ 被 $J$ 的某个幂零化，$M_2$ 被 $I$ 的某个幂零化，
而 $N$ 被 $\mathfrak m$ 的某个幂零化。对 $M_1$，有
$R\Gamma_J(M_1) = M_1$；又因为 $M_1$ 是有限 $A$-模且
$I$-进完备，所以 $M_1^\wedge = M_1$。因此断言对 $M_1$ 成立。
对 $M_2$，$H^i_J(M_2)$ 被 $I$ 的某个幂零化，因而导出完备；
所以断言对 $M_2$ 也成立。对 $N$ 可用相同论证，故断言成立。
考察短正合列 $0 \to M' \to M \to M/M' \to 0$，
可知只需对 $M/M'$ 证明引理。因此可假设
$\mathfrak p \in \text{Ass}(M)$ 蕴含
$V(\mathfrak p) \cap V(I) \not = \{\mathfrak m\}$，即
$\mathfrak p$ 是 (4) 中所述的素理想。

\medskip\noindent
选择理想 $J_0$，如引理 \ref{algebraization-lemma-kill-colimit-weak} 中所述；
再选择整数 $t > 0$，使 $(J_0I)^t$ 零化 $H^s_J(M)$。
这里 $J$ 表示任意理想 $J \subset J_0$，并满足
$V(J) \cap V(I) = \{\mathfrak m\}$。
引理 \ref{algebraization-lemma-kill-completion} 的假设对每个
$\mathfrak p \in \text{Ass}(M)$ 都成立（见上一段）。因此，
零化理想 $\mathfrak a \subset A$（它零化
$H^s(R\Gamma_\mathfrak m(M)^\wedge)$）不包含于 $\mathfrak p$，
其中 $\mathfrak p \in \text{Ass}(M)$。于是可以找到
$f \in \mathfrak a(J_0I)^t$，它不属于 $M$ 的任何相伴素理想，
并且同时零化 $H^s(R\Gamma_\mathfrak m(M)^\wedge)$ 与 $H^s_J(M)$。
于是 $f$ 是 $M$ 上的非零因子，可以考察短正合列
$$
0 \to M \xrightarrow{f} M \to M/fM \to 0
$$
由 $f$ 的选择，得到
$$
\xymatrix{
H^{s - 1}_J(M) \ar[d] \ar[r] &
H^{s - 1}_J(M/fM) \ar[d] \ar[r] &
H^s_J(M) \ar[d] \ar[r] & 0 \\
H^{s - 1}(R\Gamma_\mathfrak m(M)^\wedge) \ar[r] &
H^{s - 1}(R\Gamma_\mathfrak m(M/fM)^\wedge) \ar[r] &
H^s(R\Gamma_\mathfrak m(M)^\wedge) \ar[r] & 0
}
$$
，其中 $J \subset J_0$ 任意，并满足
$V(J) \cap V(I) = \{\mathfrak m\}$。
因此，若选择 $J$，使它对 $M$、$M/fM$ 与 $s - 1$
都适用（由归纳假设可以做到），便可得出引理成立。
\end{proof}








\section{局部上同调的代数化，III}
\label{algebraization-section-bootstrap}

\noindent
本节利用第 \ref{algebraization-section-algebraization-sections-general} 节和
第 \ref{algebraization-section-algebraization-punctured} 节中的材料作进一步提升，
在如下情形中得到一个更强的结果。

\begin{situation}
\label{algebraization-situation-bootstrap}
这里 $A$ 是诺特环。给定理想 $I \subset A$、有限 $A$-模 $M$，
以及对特殊化稳定的子集 $T \subset V(I)$。另给定整数 $s$ 和 $d$。
我们假设
\begin{enumerate}
\item[(1)] $A$ 有对偶化复形，
\item[(3)] $\text{cd}(A, I) \leq d$,
\item[(4)] 对任意素理想 $\mathfrak p \subset \mathfrak r \subset \mathfrak q$，
若 $\mathfrak p \not \in V(I)$、
$\mathfrak r \in V(I) \setminus T$ 且
$\mathfrak q \in T$，则
$$
\text{depth}_{A_\mathfrak p}(M_\mathfrak p) \geq s
\quad\text{或}\quad
\text{depth}_{A_\mathfrak p}(M_\mathfrak p) +
\dim((A/\mathfrak p)_\mathfrak q) > d + s
$$
\item[(6)] 对任意 $\mathfrak q \in T$，以
$A', \mathfrak m', I', M'$ 分别表示关于 $I$ 的通常的进完备化，
它们来自 $A_\mathfrak q, \mathfrak qA_\mathfrak q, I_\mathfrak q, M_\mathfrak q$，
则
$$
\text{depth}(M'_{\mathfrak p'}) > s
$$
对所有 $\mathfrak p' \in \Spec(A') \setminus V(I')$ 中满足
$V(\mathfrak p') \cap V(I') = \{\mathfrak m'\}$ 者成立。
\end{enumerate}
\end{situation}

\noindent
下面的引理说明，在情形 \ref{algebraization-situation-bootstrap} 中，
在 (4) 中考察素理想三元组
$\mathfrak p \subset \mathfrak r \subset \mathfrak q$ 已经足够，
尽管实际的假设只涉及 $\mathfrak p$ 和 $\mathfrak q$。

\begin{lemma}
\label{algebraization-lemma-helper-bootstrap}
在情形 \ref{algebraization-situation-bootstrap} 中，设 $\mathfrak p \subset \mathfrak q$
是 $A$ 的素理想，且 $\mathfrak p \not \in V(I)$、
$\mathfrak q \in T$。若不存在满足
$\mathfrak r \in V(I) \setminus T$ 且
$\mathfrak p \subset \mathfrak r \subset \mathfrak q$ 的素理想，
则 $\text{depth}(M_\mathfrak p) > s$。
\end{lemma}

\begin{proof}
选取 $\mathfrak q' \in T$，使得
$\mathfrak p \subset \mathfrak q' \subset \mathfrak q$，
并且 $T$ 中不存在严格位于 $\mathfrak p$ 与 $\mathfrak q'$ 之间的
素理想。为证明该引理，我们可以并且确实把 $\mathfrak q$
替换为 $\mathfrak q'$。接着，令 $\mathfrak p' \subset A_\mathfrak q$
为与 $\mathfrak p$ 对应的素理想。由于不存在引理中所述的
素理想，这样便得到
$V(\mathfrak p') \cap V(IA_\mathfrak q) = \{\mathfrak q A_\mathfrak q\}$；
这里所用的正是这样的 $\mathfrak r$ 不存在这一事实。
令 $A', I', \mathfrak m', M'$ 为关于 $I$ 的进完备化，
它们分别来自
$A_\mathfrak q, I_\mathfrak q, \mathfrak qA_\mathfrak q, M_\mathfrak q$。
由于 $A_\mathfrak q \to A'$ 忠实平坦
（代数，引理 \ref{algebra-lemma-completion-faithfully-flat}），
我们可以选取 $\mathfrak p'' \subset A'$，它位于 $\mathfrak p'$ 之上，并使得
$\dim(A'_{\mathfrak p''}/\mathfrak p' A'_{\mathfrak p''}) = 0$。
于是可见
$$
\text{depth}(M'_{\mathfrak p''}) =
\text{depth}((M_\mathfrak q \otimes_{A_\mathfrak q} A')_{\mathfrak p''}) =
\text{depth}(M_\mathfrak p \otimes_{A_\mathfrak p} A'_{\mathfrak p''}) =
\text{depth}(M_\mathfrak p)
$$
这是由 $A \to A'$ 的平坦性以及对 $\mathfrak p''$ 的选取所得；参见
代数，引理 \ref{algebra-lemma-apply-grothendieck-module}。
由于 $\mathfrak p''$ 位于 $\mathfrak p'$ 之上，我们有
$V(\mathfrak p'') \cap V(I') = \{\mathfrak m'\}$。因此，
情形 \ref{algebraization-situation-bootstrap} 中的条件 (6) 蕴含
$\text{depth}(M'_{\mathfrak p''}) > s$，证明完毕。
\end{proof}

\noindent
下面这个冗长的引理说明可能施加的各组条件之间的关系。

\begin{lemma}
\label{algebraization-lemma-bootstrap-inherited}
在情形 \ref{algebraization-situation-bootstrap} 中，有
\begin{enumerate}
\item[(E)] 若 $T' \subset T$ 是更小的、对特殊化稳定的子集，则
$A, I, T', M$ 满足情形 \ref{algebraization-situation-bootstrap} 的假设；
\item[(F)] 若 $S \subset A$ 是乘法子集，则
$S^{-1}A, S^{-1}I, T', S^{-1}M$
满足情形 \ref{algebraization-situation-bootstrap} 的假设，
其中 $T' \subset V(S^{-1}I)$ 是 $T$ 的逆像；
\item[(G)] 四元组 $A', I', T', M'$
满足情形 \ref{algebraization-situation-bootstrap} 的假设，
其中 $A', I', M'$ 是通常的 $I$-进完备化，
它们分别来自 $A, I, M$，
而 $T' \subset V(I')$ 是 $T$ 的逆像。
\end{enumerate}
设 $I \subset \mathfrak a \subset A$ 是满足
$V(\mathfrak a) \subset T$ 的理想。则
\begin{enumerate}
\item[(A)] 若 $I$ 包含于 $A$ 的 Jacobson 根，则对于
$A, I, \mathfrak a, M$，引理 \ref{algebraization-lemma-kill-colimit-weak-general} 和
\ref{algebraization-lemma-kill-colimit-support-general} 的全部假设均成立；
\item[(B)] 若 $A$ 关于 $I$ 完备，则对于
$A, I, \mathfrak a, M$，引理
\ref{algebraization-lemma-algebraize-local-cohomology-general}
除可能的 (5) 外全部假设均成立；
\item[(C)] 若 $A$ 是以 $\mathfrak m = \mathfrak a$ 为极大理想的局部环，
则对于 $A, \mathfrak m, I, M$，
引理 \ref{algebraization-lemma-kill-colimit-weak} 和 \ref{algebraization-lemma-kill-colimit-support}
的全部假设均成立；
\item[(D)] 若 $A$ 是以 $\mathfrak m = \mathfrak a$ 为极大理想的局部环，
且 $I$-进完备，则对于 $A, \mathfrak m, I, M$，
引理 \ref{algebraization-lemma-algebraize-local-cohomology}
的全部假设均成立。
\end{enumerate}
\end{lemma}

\begin{proof}
证明 (E)。我们须证明情形 \ref{algebraization-situation-bootstrap} 的假设
(1)、(3)、(4)、(6) 对 $A, I, T, M$ 成立。把 $T$ 缩小为 $T'$
会减弱假设 (6) 而加强假设 (4)。不过，若有假设 (4) 中所述的
$\mathfrak p \subset \mathfrak r \subset \mathfrak q$，其中
$\mathfrak p \not \in V(I)$、$\mathfrak r \in V(I) \setminus T'$、
$\mathfrak q \in T'$，这正是 $A, I, T', M$ 的假设 (4)；那么要么可以选取
$\mathfrak r \in V(I) \setminus T$，此时可用 $A, I, T, M$ 的条件 (4)；
要么找不到这样的 $\mathfrak r$，此时由引理
\ref{algebraization-lemma-helper-bootstrap} 得到
$\text{depth}(M_\mathfrak p) > s$。
这便按要求证明了 (4) 对 $A, I, T', M$ 成立。

\medskip\noindent
证明 (F)。这是直接的，细节从略。

\medskip\noindent
证明 (G)。我们须证明情形 \ref{algebraization-situation-bootstrap} 的假设
(1)、(3)、(4)、(6) 对 $I$-进完备化
$A', I', T', M'$ 成立。请注意，
$\Spec(A') \to \Spec(A)$ 诱导同构 $V(I') \to V(I)$。

\medskip\noindent
假设 (1)：环 $A'$ 有对偶化复形；参见
对偶化复形，引理 \ref{dualizing-lemma-ubiquity-dualizing}。

\medskip\noindent
假设 (3)：由于 $I' = IA'$，这可由局部上同调，
引理 \ref{local-cohomology-lemma-cd-change-rings} 得出。

\medskip\noindent
假设 (4)：若有素理想
$\mathfrak p' \subset \mathfrak r' \subset \mathfrak q'$ 属于 $A'$，并满足
$\mathfrak p' \not \in V(I')$、
$\mathfrak r' \in V(I') \setminus T'$、
$\mathfrak q' \in T'$，则它们的像
$\mathfrak p \subset \mathfrak r \subset \mathfrak q$ 位于 $A$ 的谱中，并
满足
$\mathfrak p \not \in V(I)$、$\mathfrak r \in V(I) \setminus T$、
$\mathfrak q \in T$。
于是
$$
\text{depth}_{A_\mathfrak p}(M_\mathfrak p) \geq s
\quad\text{或}\quad
\text{depth}_{A_\mathfrak p}(M_\mathfrak p) +
\dim((A/\mathfrak p)_\mathfrak q) > d + s
$$
这是由 $A, I, T, M$ 的假设 (4) 得到的。又有
$\text{depth}(M'_{\mathfrak p'}) \geq \text{depth}(M_\mathfrak p)$，且
$\text{depth}(M'_{\mathfrak p'}) +
\dim((A'/\mathfrak p')_{\mathfrak q'}) =
\text{depth}(M_\mathfrak p) +
\dim((A/\mathfrak p)_\mathfrak q)$
由局部上同调，引理 \ref{local-cohomology-lemma-change-completion} 可得。
因此假设 (4) 对 $A', I', T', M'$ 成立。

\medskip\noindent
假设 (6)：令 $\mathfrak q' \in T'$ 位于素理想
$\mathfrak q \in T$ 之上。则 $A'_{\mathfrak q'}$
与 $A_\mathfrak q$ 有同构的 $I$-进完备化，
$M_\mathfrak q$ 与 $M'_{\mathfrak q'}$ 也同样如此。
因此，$A', I', T', M'$ 的假设 (6) 等价于
$A, I, T, M$ 的假设 (6)。

\medskip\noindent
证明 (A)。我们须对 $(A, I, \mathfrak a, M)$ 检验
引理 \ref{algebraization-lemma-kill-colimit-weak-general} 和
\ref{algebraization-lemma-kill-colimit-support-general} 的条件 (1)、(2)、(3)、(4)、(6)。
注意：这些引理陈述中的集合 $T$ 与上面的集合 $T$ 不同。

\medskip\noindent
条件 (1)：成立，因为我们在情形 \ref{algebraization-situation-bootstrap} 中假设了
$A$ 有对偶化复形。

\medskip\noindent
条件 (2)：此条件无内容。

\medskip\noindent
条件 (3)：设 $\mathfrak p \subset A$ 满足
$V(\mathfrak p) \cap V(I) \subset V(\mathfrak a)$。
由于 $I$ 包含于 $A$ 的 Jacobson 根，可见
$V(\mathfrak p) \cap V(I) \not = \emptyset$。
令 $\mathfrak q \in V(\mathfrak p) \cap V(I)$ 为一个一般点。
由于 $\text{cd}(A_\mathfrak q, I_\mathfrak q) \leq d$
（局部上同调，引理 \ref{local-cohomology-lemma-cd-local}），并且
$V(\mathfrak p A_\mathfrak q) \cap V(I_\mathfrak q) =
\{\mathfrak q A_\mathfrak q\}$，由局部上同调，
引理 \ref{local-cohomology-lemma-cd-bound-dim-local} 得到
$\dim((A/\mathfrak p)_\mathfrak q) \leq d$，从而证明 (3)。

\medskip\noindent
条件 (4)：设 $\mathfrak p \not \in V(I)$ 且
$\mathfrak q \in V(\mathfrak p) \cap V(\mathfrak a)$。
只须证明
$$
\text{depth}_{A_\mathfrak p}(M_\mathfrak p) \geq s
\quad\text{或}\quad
\text{depth}_{A_\mathfrak p}(M_\mathfrak p) +
\dim((A/\mathfrak p)_\mathfrak q) > d + s
$$
若存在素理想 $\mathfrak p \subset \mathfrak r \subset \mathfrak q$，
其中 $\mathfrak r \in V(I) \setminus T$，则这立即由情形
\ref{algebraization-situation-bootstrap} 中的假设 (4) 得出。
若不存在，则由引理 \ref{algebraization-lemma-helper-bootstrap} 有
$\text{depth}(M_\mathfrak p) > s$。

\medskip\noindent
条件 (6)：设 $\mathfrak p \not \in V(I)$ 满足
$V(\mathfrak p) \cap V(I) \subset V(\mathfrak a)$。
由于 $I$ 包含于 $A$ 的 Jacobson 根，可见
$V(\mathfrak p) \cap V(I) \not = \emptyset$。
选取 $\mathfrak q \in V(\mathfrak p) \cap V(I) \subset V(\mathfrak a)$。
显然不存在素理想
$\mathfrak p \subset \mathfrak r \subset \mathfrak q$，
使得 $\mathfrak r \in V(I) \setminus T$。
由引理 \ref{algebraization-lemma-helper-bootstrap} 有
$\text{depth}(M_\mathfrak p) > s$，这证明了 (6)。

\medskip\noindent
证明 (B)。我们须检验引理
\ref{algebraization-lemma-algebraize-local-cohomology-general} 的条件 (1)、(2)、(3)、(4)。
注意：该引理陈述中的集合 $T$ 与上面的集合 $T$ 不同。

\medskip\noindent
条件 (1)：成立，因为 $A$ 完备且有对偶化复形。

\medskip\noindent
条件 (2)：此条件无内容。

\medskip\noindent
条件 (3)：这与情形 \ref{algebraization-situation-bootstrap} 中的假设 (3) 相同。

\medskip\noindent
条件 (4)：这与引理 \ref{algebraization-lemma-kill-colimit-weak-general} 中的假设 (4)
相同，而我们已在 (A) 中证明了后者。

\medskip\noindent
证明 (C)。这是因为引理 \ref{algebraization-lemma-kill-colimit-weak} 和
\ref{algebraization-lemma-kill-colimit-support} 中的假设，与引理
\ref{algebraization-lemma-kill-colimit-weak-general} 和
\ref{algebraization-lemma-kill-colimit-support-general} 在局部情形下的假设相同，
而我们已在 (A) 中证明了这些假设成立。

\medskip\noindent
证明 (D)。这是因为引理
\ref{algebraization-lemma-algebraize-local-cohomology} 中的假设，与引理
\ref{algebraization-lemma-algebraize-local-cohomology-general}
中的假设 (1)、(2)、(3)、(4) 相同，而我们已在 (B) 中证明了这些假设成立。
\end{proof}

\begin{lemma}
\label{algebraization-lemma-algebraize-local-cohomology-bis}
在情形 \ref{algebraization-situation-bootstrap} 中，假设 $A$ 是以
$\mathfrak m$ 为极大理想的局部环，并且 $T = \{\mathfrak m\}$。则
$H^i_\mathfrak m(M) \to \lim H^i_\mathfrak m(M/I^nM)$
在 $i \leq s$ 时为同构，并且这些模都被 $I$ 的某个幂零化。
\end{lemma}

\begin{proof}
令 $A', I', \mathfrak m', M'$ 为通常的 $I$-进完备化，
分别来自 $A, I, \mathfrak m, M$。回忆我们有
$H^i_\mathfrak m(M) \otimes_A A' = H^i_{\mathfrak m'}(M')$
这是由 $A \to A'$ 的平坦性以及对偶化复形，引理
\ref{dualizing-lemma-torsion-change-rings} 得到的。
由于 $H^i_\mathfrak m(M)$ 是 $\mathfrak m$-幂挠模，故有
$H^i_\mathfrak m(M) = H^i_\mathfrak m(M) \otimes_A A'$；参见
代数进阶，引理 \ref{more-algebra-lemma-neighbourhood-equivalence}。
由此得到 $H^i_\mathfrak m(M) = H^i_{\mathfrak m'}(M')$。
完全相同的论证表明
$H^i_\mathfrak m(M/I^nM) =  H^i_{\mathfrak m'}(M'/(I')^nM')$
对所有 $n$ 和 $i$ 成立。

\medskip\noindent
由引理 \ref{algebraization-lemma-algebraize-local-cohomology}、
\ref{algebraization-lemma-kill-colimit-weak} 和
\ref{algebraization-lemma-kill-colimit-support}
以及引理 \ref{algebraization-lemma-bootstrap-inherited} 的 (C)、(D) 部分，
上述引理适用于 $A', \mathfrak m', I', M'$。
因此得到同构
$$
H^i_{\mathfrak m'}(M') \longrightarrow H^i(R\Gamma_{\mathfrak m'}(M')^\wedge)
$$
其中 $i \leq s$，而 ${}^\wedge$ 表示导出 $I'$-进完备化；
这些模都被 $I'$ 的某个幂零化。
由引理 \ref{algebraization-lemma-local-cohomology-derived-completion}
得到同构
$$
H^i_{\mathfrak m'}(M') \longrightarrow
\lim H^i_{\mathfrak m'}(M'/(I')^nM'))
$$
其中 $i \leq s$。结合已经建立的与 $A$ 上局部上同调的比较，
便得出该引理。
\end{proof}

\begin{lemma}
\label{algebraization-lemma-bootstrap-bis-bis}
设 $I \subset \mathfrak a$ 是诺特环 $A$ 的理想。
设 $M$ 是有限 $A$-模，并设 $s$ 和 $d$ 为整数。
若假设
\begin{enumerate}
\item[(a)] $A$ 有对偶化复形；
\item[(b)] $\text{cd}(A, I) \leq d$,
\item[(c)] 若 $\mathfrak p \not \in V(I)$ 且
$\mathfrak q \in V(\mathfrak p) \cap V(\mathfrak a)$，则
$\text{depth}_{A_\mathfrak p}(M_\mathfrak p) > s$，或
$\text{depth}_{A_\mathfrak p}(M_\mathfrak p) +
\dim((A/\mathfrak p)_\mathfrak q) > d + s$.
\end{enumerate}
则 $A, I, V(\mathfrak a), M, s, d$ 满足
情形 \ref{algebraization-situation-bootstrap} 的条件。
\end{lemma}

\begin{proof}
我们须证明情形 \ref{algebraization-situation-bootstrap} 的假设 (1)、(3)、(4)、(6) 成立。
显然 (a) $\Rightarrow$ (1)、
(b) $\Rightarrow$ (3)，且 (c) $\Rightarrow$ (4)。
为完成证明，下一段证明 (6) 成立。

\medskip\noindent
设 $\mathfrak q \in V(\mathfrak a)$。
用 $A', I', \mathfrak m', M'$ 表示关于 $I$ 的进完备化，
它们分别来自
$A_\mathfrak q, I_\mathfrak q, \mathfrak qA_\mathfrak q, M_\mathfrak q$。
设 $\mathfrak p' \subset A'$ 是满足
$V(\mathfrak p') \cap V(I') = \{\mathfrak m'\}$ 的非极大素理想。
由局部上同调，引理 \ref{local-cohomology-lemma-cd-bound-dim-local}，
注意到这蕴含 $\dim(A'/\mathfrak p') \leq d$。
记 $\mathfrak p \subset A$ 为 $\mathfrak p'$ 的像。我们有
$\text{depth}(M'_{\mathfrak p'}) \geq \text{depth}(M_\mathfrak p)$，且
$\text{depth}(M'_{\mathfrak p'}) +
\dim(A'/\mathfrak p') =
\text{depth}(M_\mathfrak p) +
\dim((A/\mathfrak p)_\mathfrak q)$
这是由局部上同调，引理 \ref{local-cohomology-lemma-change-completion} 得到的。
由假设 (c)，要么有
$\text{depth}(M'_{\mathfrak p'}) \geq \text{depth}(M_\mathfrak p) > s$
而结论已经得到；要么有
$\text{depth}(M'_{\mathfrak p'}) +
\dim(A'/\mathfrak p') > s + d$，结合已经证明的不等式，便蕴含
$\text{depth}(M'_{\mathfrak p'}) > s$，因为
$\dim(A'/\mathfrak p') \leq d$。两种情形都得到所需结论。
\end{proof}

\begin{lemma}
\label{algebraization-lemma-bootstrap}
在情形 \ref{algebraization-situation-bootstrap} 中，逆系统
$\{H^i_T(I^nM)\}_{n \geq 0}$ 在 $i \leq s$ 时为 pro-零。
此外，存在整数 $m_0$，使得对所有
$m \geq m_0$，都存在整数 $m'(m) \geq m$，满足：当
$k \geq m'(m)$ 时，映射
$H^{s + 1}_T(I^kM) \to H^{s + 1}_T(I^mM)$
的像单射映入 $H^{s + 1}_T(I^{m_0}M)$。
\end{lemma}

\begin{proof}
固定 $m$。设 $\mathfrak q \in T$。
由引理 \ref{algebraization-lemma-bootstrap-inherited} 和
\ref{algebraization-lemma-algebraize-local-cohomology-bis}
可知
$$
H^i_\mathfrak q(M_\mathfrak q)
\longrightarrow
\lim H^i_\mathfrak q(M_\mathfrak q/I^nM_\mathfrak q)
$$
在 $i \leq s$ 时为同构。逆系统
$\{H^i_\mathfrak q(I^nM_\mathfrak q)\}_{n \geq 0}$ 和
$\{H^i_\mathfrak q(M/I^nM)\}_{n \geq 0}$
对所有 $i$ 都满足 Mittag-Leffler 条件；参见
引理 \ref{algebraization-lemma-ML-local}。因此，考察如下长正合序列组成的逆系统：
$$
0 \to H^0_\mathfrak q(I^nM_\mathfrak q) \to
H^0_\mathfrak q(M_\mathfrak q) \to
H^0_\mathfrak q(M_\mathfrak q/I^nM_\mathfrak q) \to
H^1_\mathfrak q(I^nM_\mathfrak q) \to
H^1_\mathfrak q(M_\mathfrak q) \to \ldots
$$
可得（略去一些细节）：存在整数
$m'(m, \mathfrak q) \geq m$，使得对所有 $k \geq m'(m, \mathfrak q)$，
映射
$H^i_\mathfrak q(I^kM_\mathfrak q) \to H^i_\mathfrak q(I^mM_\mathfrak q)$
在 $i \leq s$ 时为零，并且映射
$H^{s + 1}_\mathfrak q(I^kM_\mathfrak q) \to
H^{s + 1}_\mathfrak q(I^mM_\mathfrak q)$
的像与 $k \geq m'(m, \mathfrak q)$ 的选取无关，并单射映入
$H^{s + 1}_\mathfrak q(M_\mathfrak q)$。

\medskip\noindent
假设能够证明 $m'(m, \mathfrak q)$ 的选取
与 $\mathfrak q \in T$ 无关。
那么该引理立即由局部上同调，引理 \ref{local-cohomology-lemma-zero} 和
\ref{local-cohomology-lemma-essential-image}.

\medskip\noindent
设 $\omega_A^\bullet$ 为对偶化复形。令
$\delta : \Spec(A) \to \mathbf{Z}$ 为相应的维数函数。
回忆 $\delta$ 只取有限多个值；参见
对偶化复形，引理 \ref{dualizing-lemma-universally-catenary}。
断言：对每个 $d \in \mathbf{Z}$，整数
$m'(m, \mathfrak q)$ 可以选得与满足
$\mathfrak q \in T$ 且 $\delta(\mathfrak q) = d$ 的素理想无关。
由上面的论述，该断言显然蕴含引理。

\medskip\noindent
选取 $\mathfrak q \in T$，使得 $\delta(\mathfrak q) = d$。
考察 Ext 模
$$
E(n, j) = \text{Ext}^j_A(I^nM, \omega_A^\bullet)
$$
我们将用到的一个关键性质是：这些都是有限 $A$-模。
由与对偶化复形相伴的维数函数之定义，回忆
$(\omega_A^\bullet)_\mathfrak q[-d]$ 是 $A_\mathfrak q$ 的规范化
对偶化复形；参见对偶化复形，第 \ref{dualizing-section-dimension-function} 节。
局部对偶定理（对偶化复形，引理
\ref{dualizing-lemma-special-case-local-duality}）告诉我们，
$\mathfrak qA_\mathfrak q$-进完备化后的
$E(n, -d - i)_\mathfrak q$ 与
$H^i_\mathfrak q(I^nM_\mathfrak q)$ Matlis 对偶。因此，第一段中
对 $m'(m, \mathfrak q)$ 的选取（其中 $i \leq s$）说明：
当 $k \geq m'(m, \mathfrak q)$ 且 $j \geq -d - s$ 时，映射
$$
E(m, j)_\mathfrak q \to E(k, j)_\mathfrak q
$$
为零。由于这些模是有限的，并且只对有限多个可能的 $j$ 非零
（略去一个小细节），我们可以找到开邻域
$W \subset \Spec(A)$，它包含 $\mathfrak q$，并使得
$$
E(m, j)_{\mathfrak q'} \to E(m'(m, \mathfrak q), j)_{\mathfrak q'}
$$
在 $j \geq -d - s$ 时，对所有 $\mathfrak q' \in W$ 均为零。
于是当然，映射 $E(m, j)_{\mathfrak q'} \to E(k, j)_{\mathfrak q'}$
对 $k \geq m'(m, \mathfrak q)$ 也为零。

\medskip\noindent
对于 $i = s + 1$，它对应于 $j = - d - s - 1$，
由局部对偶和第一段的结果可得
$$
K_{k, \mathfrak q} =
\Ker(E(m, -d - s - 1)_\mathfrak q \to E(k, -d - s - 1)_\mathfrak q)
$$
与 $k \geq m'(m, \mathfrak q)$ 无关，并且
$$
E(0, -d - s - 1)_\mathfrak q \to
E(m, -d - s - 1)_\mathfrak q/K_{m'(m, \mathfrak q), \mathfrak q}
$$
为满射。对 $k \geq m'(m, \mathfrak q)$，令
$$
K_k = \Ker(E(m, -d - s - 1) \to E(k, -d - s - 1))
$$
由于 $K_k$ 是有限模 $E(m, -d - s - 1)$ 的递增子模序列，
可见在稍微增大 $m'(m, \mathfrak q)$ 后，我们可以假设
$K_{m'(m, \mathfrak q)} = K_k$ 对 $k \geq m'(m, \mathfrak q)$ 成立。
必要时进一步缩小 $W$，还可假设
$$
E(0, -d - s - 1)_{\mathfrak q'} \to
E(m, -d - s - 1)_{\mathfrak q'}/K_{m'(m, \mathfrak q), \mathfrak q'}
$$
对所有 $\mathfrak q' \in W$ 都为满射（与前面一样，使用这些模的有限性，
以及该映射在 $\mathfrak q$ 处局部化后为满射）。

\medskip\noindent
任意子集，特别是
$T_d = \{\mathfrak q \in T \text{ 其中 }\delta(\mathfrak q) = d\}$,
作为赋予子空间拓扑的诺特拓扑空间 $\Spec(A)$ 的子集，
仍是诺特的，因而是拟紧的。
上面已经看到，对每个 $\mathfrak q \in T_d$，
都存在开邻域 $W$，使得 $m'(m, \mathfrak q)$
对所有 $\mathfrak q' \in T_d \cap W$ 都适用。
由此可知，可以找到整数 $m'(m, d)$，使得对所有
$\mathfrak q \in T_d$ 都有
$$
E(m, j)_\mathfrak q \to E(m'(m, d), j)_\mathfrak q
$$
在 $j \geq -d - s$ 时为零；并且令
$K_{m'(m, d)} = \Ker(E(m, -d - s - 1) \to E(m'(m, d), -d - s - 1))$
后，有
$$
K_{m'(m, d), \mathfrak q} =
\Ker(E(m, -d - s - 1)_{\mathfrak q} \to E(k, -d - s - 1)_{\mathfrak q})
$$
对所有 $k \geq m'(m, d)$ 成立，并且映射
$$
E(0, -d - s - 1)_\mathfrak q \to
E(m, -d - s - 1)_\mathfrak q/K_{m'(m, d), \mathfrak q}
$$
为满射。再次反向使用局部对偶定理，可得该断言成立。证明完毕。
\end{proof}

\begin{lemma}
\label{algebraization-lemma-final-bootstrap}
在情形 \ref{algebraization-situation-bootstrap} 中，存在整数 $m_0 \geq 0$，使得
\begin{enumerate}
\item $H^i_T(M) \to H^i_T(M/I^nM)$ 对所有 $i \leq s$
及 $n \geq m_0$ 都是单射；
\item $\{H^i_T(M/I^nM)\}_{n \geq 0}$
在 $i < s$ 时满足 Mittag-Leffler 条件；
\item $\{H^i_T(I^{m_0}M/I^nM)\}_{n \geq m_0}$
在 $i \leq s$ 时满足 Mittag-Leffler 条件；
\item $H^i_T(M) \to \lim H^i_T(M/I^nM)$
在 $i < s$ 时为同构；
\item $H^s_T(I^{m_0}M) \to \lim H^s_T(I^{m_0}M/I^nM)$
在 $i \leq s$ 时为同构；
\item $H^s_T(M) \to \lim H^s_T(M/I^nM)$ 为单射，
且其余核被 $I^{m_0}$ 零化；
\item $R^1\lim H^s_T(M/I^nM)$ 被 $I^{m_0}$ 零化。
\end{enumerate}
\end{lemma}

\begin{proof}
为证明 (1)，设 $0 \leq i \leq s$，并选取整数 $m > n > 0$，使得
$H^i_T(I^mM) \to H^i_T(I^nM)$ 为零；由引理
\ref{algebraization-lemma-bootstrap} 可以这样选取。考察各行正合的交换图
$$
\xymatrix{
H^i_T(I^nM) \ar[r] &
H^i_T(M) \ar[r] &
H^i_T(M/I^nM) \\
H^i_T(I^mM) \ar[r] \ar[u]^0 &
H^i_T(M) \ar[r] \ar[u]^1 &
H^i_T(M/I^mM) \ar[u] \\
}
$$
该图表明映射 $H^i_T(M) \to H^i_T(M/I^mM)$ 为单射。因此，
对任意充分大的 $m_0$，(1) 成立。

\medskip\noindent
考察长正合序列
$$
0 \to H^0_T(I^nM) \to H^0_T(M) \to
H^0_T(M/I^nM) \to H^1_T(I^nM) \to
H^1_T(M) \to \ldots
$$
由此及引理 \ref{algebraization-lemma-bootstrap} 得到 (2) 和 (4)。

\medskip\noindent
令 $m_0$ 和 $m'(-)$ 如引理 \ref{algebraization-lemma-bootstrap} 中所述。
对 $m \geq m_0$，考察长正合序列
$$
H^s_T(I^mM) \to H^s_T(I^{m_0}M) \to
H^s_T(I^{m_0}M/I^mM) \to H^{s + 1}_T(I^mM) \to
H^1_T(I^{m_0}M)
$$
于是，对 $k \geq m'(m)$，映射
$H^{s + 1}_T(I^kM) \to H^{s + 1}_T(I^mM)$
的像单射映入 $H^{s + 1}_T(I^{m_0}M)$。
因此，映射
$H^s_T(I^{m_0}M/I^kM) \to H^s_T(I^{m_0}M/I^mM)$
的像映为 $H^{s + 1}_T(I^mM)$ 中的零；这对所有 $k \geq m'(m)$ 成立。
由此得到 (3) 和 (5)。

\medskip\noindent
考察短正合序列
$0 \to I^{m_0}M \to M \to M/I^{m_0} M \to 0$ 和
$0 \to I^{m_0}M/I^nM \to M/I^nM \to M/I^{m_0} M \to 0$.
由此得到图表
$$
\xymatrix{
H^{s - 1}_T(M/I^{m_0}M) \ar[r] &
\lim H^s_T(I^{m_0}M/I^nM) \ar[r] &
\lim H^s_T(M/I^nM) \ar[r] &
H^s_T(M/I^{m_0}M) \\
H^{s - 1}_T(M/I^{m_0}M) \ar[r] \ar@{=}[u] &
H^s_T(I^{m_0}M) \ar[r] \ar[u]_{\cong} &
H^s_T(M) \ar[r] \ar[u] &
H^s_T(M/I^{m_0}M) \ar@{=}[u]
}
$$
其下行正合。由同调，引理
\ref{homology-lemma-apply-Mittag-Leffler}，上行在中间两处也正合。
由此得到 (6)。

\medskip\noindent
记 $B_n = H^s_T(M/I^nM)$。令 $A_n \subset B_n$ 为映射
$H^s_T(I^{m_0}M/I^nM) \to H^s_T(M/I^nM)$ 的像。
由 (3) 和同调，引理 \ref{homology-lemma-Mittag-Leffler}，
$(A_n)$ 满足 Mittag-Leffler 条件。
又有 $C_n = B_n/A_n$ 被 $I^{m_0}$ 零化。因此
$R^1\lim B_n \cong R^1\lim C_n$ 被 $I^{m_0}$ 零化，从而得到 (7)。
\end{proof}

\begin{theorem}
\label{algebraization-theorem-final-bootstrap}
在情形 \ref{algebraization-situation-bootstrap} 中，对 $i \leq s$，逆系统
$\{H^i_T(M/I^nM)\}_{n \geq 0}$ 本质上常值，其值为 $H^i_T(M)$。
特别地，$\{H^i_T(M/I^nM)\}_{n \geq 0}$ 满足 Mittag-Leffler 条件，
$H^i_T(M)$ 被 $I$ 的某个幂零化，并且
$H^i_T(M) = \lim H^i_T(M/I^nM)$.
\end{theorem}

\begin{proof}
设 $0 \leq i \leq s$。
若能证明 $\{H^i_T(M/I^nM)\}_{n \geq 0}$ 满足 Mittag-Leffler 条件并且
$H^i_T(M) = \lim H^i_T(M/I^nM)$，那么引理
\ref{algebraization-lemma-final-bootstrap} 中证明的映射
$H^i_T(M) \to H^i_T(M/I^nM)$ 在 $n \gg 0$ 时的单射性将蕴含：
$\{H^i_T(M/I^nM)\}_{n \geq 0}$ 本质上常值，其值为 $H^i_T(M)$。
略去一个小细节。

\medskip\noindent
由引理 \ref{algebraization-lemma-final-bootstrap}，我们已经知道
$\{H^i_T(M/I^nM)\}_{n \geq 0}$ 满足 Mittag-Leffler 条件，并且
$H^i_T(M) = \lim H^i_T(M/I^nM)$ 对 $i < s$ 成立。
我们还知道，对 $i = s$ 以及模 $I^{m_0}M$，这些陈述成立，
其中某个 $m_0 > 0$。为完成证明，我们将说明事实上
这些关于 $i = s$ 的断言对 $M$ 也成立。

\medskip\noindent
令 $M' = H^0_I(M)$ 且 $M'' = M/M'$，从而有短正合序列
$$
0 \to M' \to M \to M'' \to 0
$$
并且由对偶化复形，引理 \ref{dualizing-lemma-divide-by-torsion}，
$M''$ 满足 $H^0_I(M') = 0$。
由 Artin--Rees（代数，引理 \ref{algebra-lemma-Artin-Rees}），得到短正合序列
$$
0 \to M' \to M/I^n M \to M''/I^n M'' \to 0
$$
对充分大的 $n$ 成立。考察长正合序列
$$
H^s_T(M') \to
H^s_T(M/I^nM) \to
H^s_T(M''/I^nM'') \to
H^{s + 1}_T(M')
$$
现在容易看出，若逆系统 $\{H^s_T(M''/I^nM'')\}_{n \geq 0}$
满足 Mittag-Leffler 条件，则逆系统
$\{H^s_T(M/I^nM)\}_{n \geq 0}$ 也满足 Mittag-Leffler 条件。
（注意，群的逆系统 $G_n$ 是否满足 ML 条件，
只取决于该逆系统在充分大的 $n$ 处的值。）
此外，序列
$$
H^s_T(M') \to
\lim H^s_T(M/I^nM) \to
\lim H^s_T(M''/I^nM'') \to
H^{s + 1}_T(M')
$$
是正合的，因为在所需位置满足 ML 条件；参见
同调，引理 \ref{homology-lemma-apply-Mittag-Leffler}。
因此，若 $H^s_T(M'') \to \lim H^s_T(M''/I^nM'')$
为同构，则
$H^s_T(M) \to \lim H^s_T(M/I^nM)$
由五引理也为同构
（同调，引理 \ref{homology-lemma-five-lemma}）。
于是归结为下一段讨论的情形。

\medskip\noindent
假设 $H^0_I(M) = 0$。选取生成元
$f_1, \ldots, f_r$ 来生成 $I^{m_0}$，其中 $m_0$ 是引理
\ref{algebraization-lemma-final-bootstrap} 中为 $M$ 找到的整数。
考察正合序列
$$
0 \to M \xrightarrow{f_1, \ldots, f_r}
(I^{m_0}M)^{\oplus r} \to Q \to 0
$$
并以此定义 $Q$。作几点观察：第一个映射为单射，
恰恰是因为 $H^0_I(M) = 0$。该单射的余核 $Q$
是有限 $A$-模，并且对每个 $1 \leq j \leq r$ 都有
$Q_{f_j} \cong (M_{f_j})^{\oplus r - 1}$。
特别地，对素理想 $\mathfrak p \subset A$，若
$\mathfrak p \not \in V(I)$，则有
$Q_\mathfrak p \cong (M_\mathfrak p)^{\oplus r - 1}$。
类似地，给定 $\mathfrak q \in T$ 和
$\mathfrak p' \subset A' = (A_\mathfrak q)^\wedge$，
其不包含于 $V(IA')$，则有
$Q'_{\mathfrak p'} \cong (M'_{\mathfrak p'})^{\oplus r - 1}$，
其中 $Q' = (Q_\mathfrak q)^\wedge$ 且 $M' = (M_\mathfrak q)^\wedge$。
因此，情形 \ref{algebraization-situation-bootstrap} 中的条件对
$A, I, T, Q$ 成立。（注意，$Q$ 可能有非零的
$H^0_I(Q)$，但这无关紧要。）

\medskip\noindent
对任意 $n \geq 0$，令 $F^nM = M \cap I^n(I^{m_0}M)^{\oplus r}$，
从而得到短正合序列
$$
0 \to F^nM \to I^n(I^{m_0}M)^{\oplus r} \to I^nQ \to 0
$$
由 Artin--Rees（代数，引理 \ref{algebra-lemma-Artin-Rees}），
存在 $c \geq 0$，使得
$I^n M \subset F^nM \subset I^{n - c}M$ 对所有 $n \geq c$ 成立。
令 $m_0$ 为整数，并令 $m'(m)$ 为对 $m \geq m_0$ 定义的函数；
它们来自将引理 \ref{algebraization-lemma-bootstrap} 应用于 $M$ 的结果。
注意，整数 $m_0$ 与我们上面选取的整数 $m_0$ 相同
（无须核验这一点：也可以直接取两个整数的最大值）。最后，
将引理 \ref{algebraization-lemma-bootstrap} 应用于 $Q$ 可知，对每个整数 $m$，
都存在整数 $m''(m) \geq m$，使得
$H^s_T(I^kQ) \to H^s_T(I^mQ)$ 对所有 $k \geq m''(m)$ 都为零。

\medskip\noindent
固定 $m \geq m_0$。选取 $k \geq m'(m''(m + c))$。
选取 $\xi \in H^{s + 1}_T(I^kM)$，
使其在 $H^{s + 1}_T(M)$ 中的像为零。
我们要证明 $\xi$ 在 $H^{s + 1}_T(I^mM)$ 中的像为零。
这将完全如引理 \ref{algebraization-lemma-final-bootstrap} 的证明那样说明
$\{H^s_T(M/I^nM)\}_{n \geq 0}$ 满足 Mittag-Leffler 条件。
下面的图有助于直观理解该论证：
$$
\xymatrix{
&
H^{s + 1}_T(I^kM) \ar[r] \ar[d] &
H^{s + 1}_T(I^k(I^{m_0}M)^{\oplus r}) \ar[d] &
\\
H^s_T(I^{m''(m + c)}Q) \ar[r]_-\delta \ar[d] &
H^{s + 1}_T(F^{m''(m + c)}M) \ar[r] \ar[d] &
H^{s + 1}_T(I^{m''(m + c)}(I^{m_0}M)^{\oplus r}) \\
H^s_T(I^{m + c}Q) \ar[r] &
H^{s + 1}_T(F^{m + c}M) \ar[d] &
\\
&
H^{s + 1}_T(I^mM)
}
$$
$\xi$ 在 $H^{s + 1}_T(I^k(I^{m_0}M)^{\oplus r})$ 中的像
映为 $H^{s + 1}_T((I^{m_0}M)^{\oplus r})$ 中的零，
因而它也映为
$H^{s + 1}_T(I^{m''(m + c)}(I^{m_0}M)^{\oplus r})$ 中的零，
这是由 $m'(-)$ 的选取得到的。
因此，像 $\xi' \in H^{s + 1}_T(F^{m''(m + c)}M)$
映为 $H^{s + 1}_T(I^{m''(m + c)}(I^{m_0}M)^{\oplus r})$ 中的零，
从而有 $\xi' = \delta(\eta)$，其中某个
$\eta \in H^s_T(I^{m''(m + c)}Q)$。
由对 $m''(-)$ 的选取可知，$\eta$ 在
$H^s_T(I^{m + c}Q)$ 中的像为零。
这又意味着 $\xi'$ 在
$H^{s + 1}_T(F^{m + c}M)$ 中的像为零。
由于 $F^{m + c}M \subset I^mM$，结论成立。

\medskip\noindent
最后证明关于极限的陈述。考察短正合序列
$$
0 \to M/F^nM \to (I^{m_0}M)^{\oplus r}/I^n (I^{m_0}M)^{\oplus r}
\to Q/I^nQ \to 0
$$
由于这些逆系统 pro-同构，我们有
$\lim H^s_T(M/I^nM) = \lim H^s_T(M/F^nM)$。得到交换图
$$
\xymatrix{
H^{s - 1}_T(Q) \ar[r] \ar[d] &
\lim H^{s - 1}_T(Q/I^nQ) \ar[d] \\
H^s_T(M) \ar[r] \ar[d] &
\lim H^s_T(M/I^nM) \ar[d] \\
H^s_T((I^{m_0}M)^{\oplus r}) \ar[r] \ar[d] &
\lim H^s_T((I^{m_0}M)^{\oplus r}/I^n(I^{m_0}M)^{\oplus r}) \ar[d] \\
H^s_T(Q) \ar[r] &
\lim H^s_T(Q/I^nQ)
}
$$
右列正合，因为在所需位置满足 ML 条件；参见
同调，引理 \ref{homology-lemma-apply-Mittag-Leffler}。
由引理 \ref{algebraization-lemma-final-bootstrap} 的 (6)，最下面的水平箭头为单射（！）。
由引理 \ref{algebraization-lemma-final-bootstrap} 的 (5)，其上方的水平箭头为双射。
上同调次数 $\leq s - 1$ 中的箭头均为同构。
因此，由五引理（同调，引理 \ref{homology-lemma-five-lemma}），可知
$H^s_T(M) \to \lim H^s_T(M/I^nM)$ 为同构。
定理得证。
\end{proof}

\begin{lemma}
\label{algebraization-lemma-combine-two}
设 $I \subset \mathfrak a \subset A$ 是诺特环 $A$ 的理想，
并设 $M$ 是有限 $A$-模。设 $s$ 和 $d$ 为整数。假设
\begin{enumerate}
\item $A, I, V(\mathfrak a), M$ 对 $s$ 和 $d$ 满足
情形 \ref{algebraization-situation-bootstrap} 的假设；并且
\item $A, I, \mathfrak a, M$ 对 $s + 1$ 和 $d$ 满足
引理 \ref{algebraization-lemma-algebraize-local-cohomology-general}
的条件，其中 $J = \mathfrak a$。
\end{enumerate}
则存在理想
$J_0 \subset \mathfrak a$，满足 $V(J_0) \cap V(I) = V(\mathfrak a)$，
使得：对任意 $J \subset J_0$，若 $V(J) \cap V(I) = V(\mathfrak a)$，
映射
$$
H^{s + 1}_J(M) \longrightarrow \lim H^{s + 1}_\mathfrak a(M/I^nM)
$$
为同构。
\end{lemma}

\begin{proof}
具体而言，由引理 \ref{algebraization-lemma-algebraize-local-cohomology-general}，
存在 $J_0$，并且有同构
$H^{s + 1}_J(M) = H^{s + 1}(R\Gamma_\mathfrak a(M)^\wedge)$
；由对偶化复形，引理 \ref{dualizing-lemma-completion-local}，有短正合序列
$$
0 \to R^1\lim H^s_\mathfrak a(M/I^nM) \to
H^{s + 1}(R\Gamma_\mathfrak a(M)^\wedge) \to
\lim H^{s + 1}_\mathfrak a(M/I^nM) \to 0
$$
而模 $R^1\lim H^s_\mathfrak a(M/I^nM)$ 为零，因为由定理
\ref{algebraization-theorem-final-bootstrap}，
$\{H^s_\mathfrak a(M/I^nM)\}_{n \geq 0}$ 满足 Mittag-Leffler 条件。
\end{proof}








\section{形式截面的代数化，I}
\label{algebraization-section-algebraization-sections}

\noindent
本节研究局部情形下形式截面的代数化问题。
设 $(A, \mathfrak m)$ 为诺特局部环。
设 $I \subset A$ 为理想。令
$$
X = \Spec(A) \supset U = \Spec(A) \setminus \{\mathfrak m\}
$$
并用 $Y = V(I)$ 表示与 $I$ 对应的闭子概形。
设 $\mathcal{F}$ 为凝聚 $\mathcal{O}_U$-模。
本节考察极限
$$
\lim_n H^i(U, \mathcal{F}/I^n\mathcal{F})
$$
这与把 $\mathcal{F}$ 拉回到 $U$ 沿 $Y$ 的形式完备化所得之上同调密切相关；
不过，由于我们尚未引入形式概形，此处不能使用这种术语。

\begin{lemma}
\label{algebraization-lemma-compare-with-derived-completion}
设 $U$ 为诺特局部环 $A$ 的穿孔谱。
设 $\mathcal{F}$ 为凝聚 $\mathcal{O}_U$-模。
设 $I \subset A$ 为理想。则
$$
H^i(R\Gamma(U, \mathcal{F})^\wedge) =
\lim H^i(U, \mathcal{F}/I^n\mathcal{F})
$$
对所有 $i$ 成立，其中 $R\Gamma(U, \mathcal{F})^\wedge$ 表示
导出 $I$-进完备化。
\end{lemma}

\begin{proof}
由引理 \ref{algebraization-lemma-formal-functions-general} 和
\ref{algebraization-lemma-derived-completion-pseudo-coherent}，有
$$
R\Gamma(U, \mathcal{F})^\wedge =
R\Gamma(U, \mathcal{F}^\wedge) =
R\Gamma(U, R\lim \mathcal{F}/I^n\mathcal{F})
$$
因此得到短正合序列
$$
0 \to R^1\lim H^{i - 1}(U, \mathcal{F}/I^n\mathcal{F}) \to
H^i(R\Gamma(U, \mathcal{F})^\wedge) \to
\lim H^i(U, \mathcal{F}/I^n\mathcal{F}) \to 0
$$
这是由上同调，引理 \ref{cohomology-lemma-RGamma-commutes-with-Rlim} 得到的。
$R^1\lim$ 项为零，因为由引理 \ref{algebraization-lemma-ML-local}，群的逆系统
$H^i(U, \mathcal{F}/I^n\mathcal{F})$ 满足 Mittag-Leffler 条件。
\end{proof}

\begin{theorem}
\label{algebraization-theorem-algebraization-formal-sections}
\begin{reference}
本证明方法大体沿用 \cite[Theorem 1]{Faltings-algebraisation}
和 \cite[Satz 2]{Faltings-uber} 的证明方法。
所得结果几乎等同于
\cite[Theorem 1.1]{MRaynaud-paper}（补集为仿射的情形）和
\cite[Theorem 3.9]{MRaynaud-book}（补集为少数仿射开集之并的情形）。
\end{reference}
设 $(A, \mathfrak m)$ 为有对偶化复形且关于理想 $I$ 完备的诺特局部环。
令 $X = \Spec(A)$、$Y = V(I)$，且 $U = X \setminus \{\mathfrak m\}$。
设 $\mathcal{F}$ 为 $U$ 上的凝聚层。假设
\begin{enumerate}
\item $\text{cd}(A, I) \leq d$, i.e.,
$H^i(X \setminus Y, \mathcal{G}) = 0$ 对 $i \geq d$ 以及
拟凝聚 $\mathcal{G}$（在 $X$ 上）成立；
\item 对任意 $x \in X \setminus Y$，若其闭包
$\overline{\{x\}}$（在 $X$ 中）与 $U \cap Y$ 相交，则有
$$
\text{depth}_{\mathcal{O}_{X, x}}(\mathcal{F}_x) \geq s
\quad\text{或}\quad
\text{depth}_{\mathcal{O}_{X, x}}(\mathcal{F}_x)
+ \dim(\overline{\{x\}}) > d + s
$$
\end{enumerate}
则存在开集 $V_0 \subset U$，它包含 $U \cap Y$，
使得对任意开集 $V \subset V_0$，若它包含 $U \cap Y$，则映射
$$
H^i(V, \mathcal{F}) \to \lim H^i(U, \mathcal{F}/I^n\mathcal{F})
$$
在 $i < s$ 时为同构。若此外
$
\text{depth}_{\mathcal{O}_{X, x}}(\mathcal{F}_x) +
\dim(\overline{\{x\}}) > s
$
对所有 $x \in U \cap Y$ 成立，则这些上同调群是有限 $A$-模。
\end{theorem}

\begin{proof}
选取有限 $A$-模 $M$，使得 $\mathcal{F}$ 是在 $U$ 上限制所得的凝聚
$\mathcal{O}_X$-模，该模与 $M$ 相伴；参见局部上同调，
引理 \ref{local-cohomology-lemma-finiteness-pushforwards-and-H1-local}。
于是引理 \ref{algebraization-lemma-algebraize-local-cohomology}
的假设成立。
按该引理选取 $J_0$，并令 $V_0 = X \setminus V(J_0)$。
于是，开集 $V \subset V_0$ 若包含 $U \cap Y$，
便与理想以 $1$-对-$1$ 的方式对应；这里 $J \subset J_0$ 且
$V(J) \cap V(I) = \{\mathfrak m\}$。
此外，对这样的选取，有优三角
$$
R\Gamma_J(M) \to M \to R\Gamma(V, \mathcal{F}) \to
R\Gamma_J(M)[1]
$$
类似地，有优三角
$$
R\Gamma_\mathfrak m(M)^\wedge \to
M \to
R\Gamma(U, \mathcal{F})^\wedge \to
R\Gamma_\mathfrak m(M)^\wedge[1]
$$
其中涉及导出 $I$-进完备化。
由引理 \ref{algebraization-lemma-compare-with-derived-completion}，
$R\Gamma(U, \mathcal{F})^\wedge$ 的上同调群等于定理陈述中的极限。
这两个三角之间的典范映射和一些容易的论证表明，本定理可由主要引理
\ref{algebraization-lemma-algebraize-local-cohomology} 得出
（注意，此处为 $i < s$，而引理中为
$i \leq s$；这是由移位造成的）。
在附加假设下，上同调群的有限性由引理
\ref{algebraization-lemma-kill-colimit} 得出。
\end{proof}

\begin{lemma}
\label{algebraization-lemma-application-theorem}
设 $(A, \mathfrak m)$ 为有对偶化复形且关于理想 $I$ 完备的诺特局部环。
令 $X = \Spec(A)$、$Y = V(I)$，且 $U = X \setminus \{\mathfrak m\}$。
设 $\mathcal{F}$ 为 $U$ 上的凝聚层。
假设对任意相伴点 $x \in U$（相对于 $\mathcal{F}$）都有
$\dim(\overline{\{x\}}) > \text{cd}(A, I) + 1$，
其中 $\overline{\{x\}}$ 是其在 $X$ 中的闭包。则映射
$$
\colim H^0(V, \mathcal{F})
\longrightarrow
\lim H^0(U, \mathcal{F}/I^n\mathcal{F})
$$
是有限 $A$-模之间的同构，
其中余极限遍历开集 $V \subset U$，这些开集包含 $U \cap Y$。
\end{lemma}

\begin{proof}
在 $s = 1$ 时应用定理 \ref{algebraization-theorem-algebraization-formal-sections}
（同时也得到有限性）。
\end{proof}




\section{形式截面的代数化，II}
\label{algebraization-section-algebraization-sections-coherent}

\noindent
要简洁陈述第 \ref{algebraization-section-algebraization-sections-general} 节和
第 \ref{algebraization-section-bootstrap} 节中的结果对于拟仿射概形上凝聚层的上同调
及其关于理想的完备化所能给出的全部推论，有些困难。
本节给出关于 $H^0$ 的一份不穷尽的应用清单。
下一节讨论对高阶上同调的应用。

\begin{lemma}
\label{algebraization-lemma-ses-H0}
设 $I \subset \mathfrak a$ 是诺特环 $A$ 的理想。
设 $0 \to \mathcal{F}' \to \mathcal{F} \to \mathcal{F}'' \to 0$
为 $U = \Spec(A) \setminus V(\mathfrak a)$ 上凝聚模的短正合序列。
令 $\mathcal{V}$ 为开子概形 $V \subset U$ 的集合，其中各开子概形都包含
$U \cap V(I)$，并按反包含排序。考察交换图
$$
\xymatrix{
\colim_\mathcal{V} H^0(V, \mathcal{F}') \ar[d] \ar[r] &
\colim_\mathcal{V} H^0(V, \mathcal{F}) \ar[d] \ar[r] &
\colim_\mathcal{V} H^0(V, \mathcal{F}'') \ar[d] \\
\lim H^0(U, \mathcal{F}'/I^n\mathcal{F}') \ar[r] &
\lim H^0(U, \mathcal{F}'/I^n\mathcal{F}) \ar[r] &
\lim H^0(U, \mathcal{F}'/I^n\mathcal{F}'')
}
$$
若左边和右边的竖直箭头都是同构，则中间的也是同构。
若中间和左边的竖直箭头都是同构，则左边的也是同构。
\end{lemma}

\begin{proof}
图中的序列在中间正合，且第一个箭头为单射。
因此最后一个陈述由简单的追图得到。在证明的其余部分，
假设左边和右边的竖直箭头均为同构。追图表明中间的竖直箭头为单射。
只须再证明它是满射。

\medskip\noindent
可以选取有限 $A$-模 $M$ 和 $M'$，使得 $\mathcal{F}$
和 $\mathcal{F}'$ 分别是 $\widetilde{M}$ 和 $\widetilde{M}'$ 在 $U$ 上的限制；
参见局部上同调，引理
\ref{local-cohomology-lemma-finiteness-pushforwards-and-H1-local}。
把 $M'$ 替换为 $\mathfrak a^n M'$（其中某个 $n \geq 0$）后，
可以假设 $\mathcal{F}' \to \mathcal{F}$ 对应于模映射
$M' \to M$；参见概形上同调，引理 \ref{coherent-lemma-homs-over-open}。
把 $M'$ 替换为 $M' \to M$ 的像，并令 $M'' = M/M'$ 后，
可见原短正合序列对应于与短正合序列
$0 \to M' \to M \to M'' \to 0$ 相伴之凝聚模短正合序列的限制；
这里所述的是 $A$-模。

\medskip\noindent
令 $\hat s \in \lim H^0(U, \mathcal{F}/I^n\mathcal{F})$，其像为
$\hat s'' \in \lim H^0(U, \mathcal{F}''/I^n\mathcal{F}'')$。
由假设，找到 $V \in \mathcal{V}$ 以及截面
$s'' \in \mathcal{F}''(V)$，后者映为 $\hat s''$。
令 $J \subset A$ 为满足 $V(J) = \Spec(A) \setminus V$ 的理想。
由概形上同调，引理 \ref{coherent-lemma-homs-over-open}，
用 $J$ 的某个幂替换它之后，可以假设存在一个
$A$-线性映射 $\varphi : J \to M''$；它对应于 $s''$。
固定这个 $J$ 的选择；在证明的其余部分，将用 $V$ 的一个更小的
$V$ 来替换它，该开集属于 $\mathcal{V}$，即将有 $V \cap V(J) = \emptyset$。

\medskip\noindent
选取一个表现 $A^{\oplus m} \to A^{\oplus n} \to J \to 0$。
用 $g_1, \ldots, g_n \in J$ 表示 $A^{\oplus n}$ 的基向量的像，
从而 $J = (g_1, \ldots, g_n)$。设
$A^{\oplus m} \to A^{\oplus n}$ 由矩阵 $(a_{ji})$ 给出，
于是 $\sum a_{ji} g_i = 0$，其中 $j = 1, \ldots, m$。
由于 $M \to M''$ 为满射，对每个 $i$ 都可选取 $m_i \in M$，
使其映为 $\varphi(g_i) \in M''$。于是
元素 $g_i \hat s - m_i$ 属于
$\lim H^0(U, \mathcal{F}/I^n\mathcal{F})$，并落在子模
$\lim H^0(U, \mathcal{F}'/I^n\mathcal{F}')$ 中。
由假设，缩小 $V$ 后可假设存在
$s'_i \in \mathcal{F}'(V)$，$i = 1, \ldots, n$，
且 $s'_i$ 映为 $g_i \hat s - m_i$。
令 $s_i = s'_i + m_i$，这是在 $\mathcal{F}(V)$ 中。
注意，$\sum a_{ji} s_i$ 映为 $\sum a_{ji}g_i\hat s = 0$；
所用映射为
$$
\colim_\mathcal{V} \mathcal{F}(V')
\longrightarrow
\lim H^0(U, \mathcal{F}/I^n\mathcal{F})
$$
由于该映射为单射（见上文），缩小 $V$ 后可以假设
$\sum a_{ji}s_i = 0$ 在 $\mathcal{F}(V)$ 中对所有 $j = 1, \ldots, m$ 成立。
于是得到 $A$-模映射 $J \to \mathcal{F}(V)$，它把 $g_i$ 映为 $s_i$。
由 $\widetilde{J}$ 的泛性质，该 $A$-模映射对应于
$\mathcal{O}_V$-模映射 $\widetilde{J}|_V \to \mathcal{F}$。
不过，由于 $V(J) \cap V = \emptyset$，有
$\widetilde{J}|_V = \mathcal{O}_V$。因此构造出了截面
$s \in \mathcal{F}(V)$。略去验证 $s$ 经上面所示映射映为 $\hat s$ 的计算。
\end{proof}

\noindent
下面的引理将被命题 \ref{algebraization-proposition-application-H0} 取代。

\begin{lemma}
\label{algebraization-lemma-application-H0-pre}
设 $I \subset \mathfrak a$ 是诺特环 $A$ 的理想。
设 $\mathcal{F}$ 为 $U = \Spec(A) \setminus V(\mathfrak a)$ 上的凝聚模。
假设
\begin{enumerate}
\item $A$ 是 $I$-进完备的，并且有对偶化复形；
\item 若 $x \in \text{Ass}(\mathcal{F})$、$x \not \in V(I)$、
$\overline{\{x\}} \cap V(I) \not \subset V(\mathfrak a)$,
且 $z \in \overline{\{x\}} \cap V(\mathfrak a)$，则
$\dim(\mathcal{O}_{\overline{\{x\}}, z}) > \text{cd}(A, I) + 1$；
\item 下列条件之一成立：
\begin{enumerate}
\item $\mathcal{F}$ 在 $U \setminus V(I)$ 上的限制满足 $(S_1)$；
\item $V(\mathfrak a)$ 的维数至多为 $2$\footnote{这里的含义是：
在 $V(\mathfrak a)$ 上，来自 $\Spec(A)$ 上一个维数函数的最大值与
最小值之差至多为 $2$。}。
\end{enumerate}
\end{enumerate}
则得到同构
$$
\colim H^0(V, \mathcal{F})
\longrightarrow
\lim H^0(U, \mathcal{F}/I^n\mathcal{F})
$$
其中余极限遍历开集 $V \subset U$，这些开集包含 $U \cap V(I)$。
\end{lemma}

\begin{proof}
选取有限 $A$-模 $M$，使得 $\mathcal{F}$ 是在 $U$ 上的限制，
所限制的凝聚模与 $M$ 相伴；参见局部上同调，引理
\ref{local-cohomology-lemma-finiteness-pushforwards-and-H1-local}。
令 $d = \text{cd}(A, I)$。
设 $\mathfrak p$ 为 $A$ 的素理想，且不包含于 $V(I)$，
并设 $\mathfrak q \in V(\mathfrak p) \cap V(\mathfrak a)$。
那么要么 $\mathfrak p$ 不是 $M$ 的相伴素理想，因而
$\text{depth}(M_\mathfrak p) \geq 1$；
要么由 (2) 有 $\dim((A/\mathfrak p)_\mathfrak q) > d + 1$。
因此，引理 \ref{algebraization-lemma-algebraize-local-cohomology-general}
的假设对 $s = 1$ 和 $d$ 成立；这里使用了条件 (3)。
由此存在理想
$J_0 \subset \mathfrak a$，满足 $V(J_0) \cap V(I) = V(\mathfrak a)$，
使得对任意 $J \subset J_0$，若 $V(J) \cap V(I) = V(\mathfrak a)$，
则映射
$$
H^i_J(M) \longrightarrow H^i(R\Gamma_\mathfrak a(M)^\wedge)
$$
在 $i = 0, 1$ 时为同构。考察正合三角之间的态射
$$
\xymatrix{
R\Gamma_J(M)  \ar[d] \ar[r] &
M \ar[r] \ar[d] &
R\Gamma(V, \mathcal{F}) \ar[d] \ar[r] &
R\Gamma_J(M)[1]  \ar[d] \\
R\Gamma_J(M)^\wedge \ar[r] &
M \ar[r] &
R\Gamma(V, \mathcal{F})^\wedge \ar[r] &
R\Gamma_J(M)^\wedge[1] \\
R\Gamma_\mathfrak a(M)^\wedge \ar[r] \ar[u] &
M \ar[r] \ar[u] &
R\Gamma(U, \mathcal{F})^\wedge \ar[r] \ar[u] &
R\Gamma_\mathfrak a(M)^\wedge[1] \ar[u]
}
$$
其中 $V = \Spec(A) \setminus V(J)$。回忆
$R\Gamma_\mathfrak a(M)^\wedge \to R\Gamma_J(M)^\wedge$
为同构（因为 $\mathfrak a$、$\mathfrak a + I$ 和 $J + I$
截出同一个闭子概形；例如参见引理
\ref{algebraization-lemma-algebraize-local-cohomology-general} 的证明）。因此
$R\Gamma(U, \mathcal{F})^\wedge = R\Gamma(V, \mathcal{F})^\wedge$.
由此得到交换图
$$
\xymatrix{
0 \ar[r] &
H^0_J(M) \ar[r] \ar[d] &
M \ar[r] \ar[d] \ar[r] &
\Gamma(V, \mathcal{F}) \ar[d] \ar[r] &
H^1_J(M) \ar[d] \ar[r] &
0 \\
0 \ar[r] &
H^0(R\Gamma_J(M)^\wedge) \ar[r] &
M \ar[r] &
H^0(R\Gamma(V, \mathcal{F})^\wedge) \ar[r] &
H^1(R\Gamma_J(M)^\wedge) \ar[r] &
0 \\
0 \ar[r] &
H^0(R\Gamma_\mathfrak a(M)^\wedge) \ar[r] \ar[u] &
M \ar[r] \ar[u] &
H^0(R\Gamma(U, \mathcal{F})^\wedge) \ar[r] \ar[u] &
H^1(R\Gamma_\mathfrak a(M)^\wedge) \ar[r] \ar[u] &
0
}
$$
其各行正合，且下方的竖直箭头均为同构。因此得到同构
$\Gamma(V, \mathcal{F}) \to H^0(R\Gamma(U, \mathcal{F})^\wedge)$.
由引理 \ref{algebraization-lemma-formal-functions-general}
和 \ref{algebraization-lemma-derived-completion-pseudo-coherent}，有
$$
R\Gamma(U, \mathcal{F})^\wedge =
R\Gamma(U, \mathcal{F}^\wedge) =
R\Gamma(U, R\lim \mathcal{F}/I^n\mathcal{F})
$$
并且由上同调，引理 \ref{cohomology-lemma-RGamma-commutes-with-Rlim}，得到
$H^0(R\Gamma(U, \mathcal{F})^\wedge) =
\lim H^0(U, \mathcal{F}/I^n\mathcal{F})$。
\end{proof}

\noindent
现在提升前一引理，以去掉条件 (3)。

\begin{proposition}
\label{algebraization-proposition-application-H0}
设 $I \subset \mathfrak a$ 是诺特环 $A$ 的理想。
设 $\mathcal{F}$ 为 $U = \Spec(A) \setminus V(\mathfrak a)$ 上的凝聚模。
假设
\begin{enumerate}
\item $A$ 是 $I$-进完备的，并且有对偶化复形；
\item 若 $x \in \text{Ass}(\mathcal{F})$、$x \not \in V(I)$、
$\overline{\{x\}} \cap V(I) \not \subset V(\mathfrak a)$,
且 $z \in \overline{\{x\}} \cap V(\mathfrak a)$，则
$\dim(\mathcal{O}_{\overline{\{x\}}, z}) > \text{cd}(A, I) + 1$.
\end{enumerate}
则得到同构
$$
\colim H^0(V, \mathcal{F})
\longrightarrow
\lim H^0(U, \mathcal{F}/I^n\mathcal{F})
$$
其中余极限遍历开集 $V \subset U$，这些开集包含 $U \cap V(I)$。
\end{proposition}

\begin{proof}
令 $T \subset U$ 为点 $x$ 所成的集合，这些点满足
$\overline{\{x\}} \cap V(I) \subset V(\mathfrak a)$。
令 $\mathcal{F} \to \mathcal{F}'$ 为局部上同调，引理
\ref{local-cohomology-lemma-get-depth-1-along-Z} 中在 $U$ 上构造的
凝聚模满射。由于 $\mathcal{F} \to \mathcal{F}'$ 在开集
$V \subset U$ 上为同构，而该开集包含 $U \cap V(I)$，
只须证明把 $\mathcal{F}$ 替换为 $\mathcal{F}'$ 后的引理。
因此可以并且确实假设：对 $x \in U$，若
$\overline{\{x\}} \cap V(I) \subset V(\mathfrak a)$，
有 $\text{depth}(\mathcal{F}_x) \geq 1$。

\medskip\noindent
令 $\mathcal{V}$ 为开子概形 $V \subset U$ 所成的集合，
其中各开子概形包含 $U \cap V(I)$，并按反包含排序。
这是一个有向集。我们先断言
$$
\mathcal{F}(V)
\longrightarrow
\lim H^0(U, \mathcal{F}/I^n\mathcal{F})
$$
对任意 $V \in \mathcal{F}$ 都为单射（特别地，引理中的映射为单射）。
事实上，由上一段，相伴点 $x$（属于 $\mathcal{F}$）必须满足
$\overline{\{x\}} \cap U \cap Y \not = \emptyset$。
若 $y \in \overline{\{x\}} \cap U \cap Y$，则
$\mathcal{F}_x$ 是 $\mathcal{F}_y$ 的局部化，并且由 Krull 交定理
（代数，引理 \ref{algebra-lemma-intersect-powers-ideal-module-zero}），有
$\mathcal{F}_y \subset \lim \mathcal{F}_y/I^n \mathcal{F}_y$。
这证明了断言，因为核中的截面 $s \in \mathcal{F}(V)$
必有空支集，因而必为零。

\medskip\noindent
选取有限 $A$-模 $M$，使得 $\mathcal{F}$ 是 $\widetilde{M}$ 在 $U$ 上的限制；
参见局部上同调，引理
\ref{local-cohomology-lemma-finiteness-pushforwards-and-H1-local}。
可以并且确实假设 $H^0_\mathfrak a(M) = 0$。
令 $\text{Ass}(M) \setminus V(I) = \{\mathfrak p_1, \ldots, \mathfrak p_n\}$。
对 $n$ 作归纳来证明该引理。重新排序后，可以假设 $\mathfrak p_n$
是集合 $\{\mathfrak p_1, \ldots, \mathfrak p_n\}$ 关于包含关系的极小元，
即 $\mathfrak p_n$ 是 $M$ 的支集的一般点。令
$$
M' = H^0_{\mathfrak p_1 \ldots \mathfrak p_{n - 1} I}(M)
$$
并令 $M'' = M/M'$。设 $\mathcal{F}'$ 和 $\mathcal{F}''$ 为凝聚
$\mathcal{O}_U$-模，它们分别对应于 $M'$ 和 $M''$。
对偶化复形，引理 \ref{dualizing-lemma-divide-by-torsion}
蕴含 $M''$ 只有一个相伴素理想，即 $\mathfrak p_n$。
因此 $\mathcal{F}''$ 只有一个相伴点，并且可见引理
\ref{algebraization-lemma-application-H0-pre} 的条件 (3)(a) 成立；
从而映射 $\colim H^0(V, \mathcal{F}'')
\to \lim H^0(U, \mathcal{F}''/I^n\mathcal{F}'')$ 为同构。
另一方面，由于
$\mathfrak p_n \not \in V(\mathfrak p_1 \ldots \mathfrak p_{n - 1} I)$
可知 $\mathfrak p_n$ 不是 $M'$ 的相伴素理想。
因此归纳假设适用于 $M'$；注意，由于 $\mathcal{F}' \subset \mathcal{F}$，
条件 $\text{depth}(\mathcal{F}'_x) \geq 1$ 对点 $x$ 成立，
其中 $\overline{\{x\}} \cap V(I) \subset V(\mathfrak a)$；
参见代数，引理 \ref{algebra-lemma-depth-in-ses}。
因此映射
$\colim H^0(V, \mathcal{F}')
\to \lim H^0(U, \mathcal{F}'/I^n\mathcal{F}')$ 也为同构。
由引理 \ref{algebraization-lemma-ses-H0} 得出结论。
\end{proof}

\begin{lemma}
\label{algebraization-lemma-application-H0}
设 $I \subset \mathfrak a$ 是诺特环 $A$ 的理想。
设 $\mathcal{F}$ 为 $U = \Spec(A) \setminus V(\mathfrak a)$ 上的凝聚模。
假设
\begin{enumerate}
\item $A$ 是 $I$-进完备的，并且有对偶化复形；
\item 若 $x \in \text{Ass}(\mathcal{F})$、$x \not \in V(I)$、
$\overline{\{x\}} \cap V(I) \not \subset V(\mathfrak a)$，且
$z \in V(\mathfrak a) \cap \overline{\{x\}}$，则
$\dim(\mathcal{O}_{\overline{\{x\}}, z}) > \text{cd}(A, I) + 1$；
\item 对 $x \in U$，若 $\overline{\{x\}} \cap V(I) \subset V(\mathfrak a)$，
则 $\text{depth}(\mathcal{F}_x) \geq 2$。
\end{enumerate}
则得到同构
$$
H^0(U, \mathcal{F})
\longrightarrow
\lim H^0(U, \mathcal{F}/I^n\mathcal{F})
$$
\end{lemma}

\begin{proof}
令 $\hat s \in \lim H^0(U, \mathcal{F}/I^n\mathcal{F})$。
由命题 \ref{algebraization-proposition-application-H0}，可知 $\hat s$ 是某个元素
$s \in \mathcal{F}(V)$ 的像，其中 $V \subset U$ 是包含
$U \cap V(I)$ 的开集。不过，条件 (3) 表明
$\text{depth}(\mathcal{F}_x) \geq 2$ 对所有 $x \in U \setminus V$ 成立，
因而由除子，引理 \ref{divisors-lemma-depth-2-hartog} 得到
$\mathcal{F}(V) = \mathcal{F}(U)$，证明完毕。
\end{proof}

\begin{lemma}
\label{algebraization-lemma-alternative-colim-H0}
设 $A$ 为诺特环。设 $f \in \mathfrak a \subset A$
为 $A$ 的一个理想中的元素。设 $M$ 为有限 $A$-模。假设
\begin{enumerate}
\item $A$ 是 $f$-进完备的；
\item $f$ 在 $M$ 上是非零因子；
\item $H^1_\mathfrak a(M/fM)$ 是有限 $A$-模。
\end{enumerate}
则令 $U = \Spec(A) \setminus V(\mathfrak a)$，映射
$$
\colim_V \Gamma(V, \widetilde{M})
\longrightarrow
\lim \Gamma(U, \widetilde{M/f^nM})
$$
为同构，其中余极限遍历开集 $V \subset U$，这些开集包含
$U \cap V(f)$。
\end{lemma}

\begin{proof}
令 $\mathcal{F} = \widetilde{M}|_U$。
$H^1_\mathfrak a(M/fM)$ 的有限性蕴含
$H^0(U, \mathcal{F}/f\mathcal{F})$ 有限；参见局部上同调，引理
\ref{local-cohomology-lemma-finiteness-pushforwards-and-H1-local}。
由上同调，引理 \ref{cohomology-lemma-limit-finite}
（该引理适用，因为 $f$ 在 $\mathcal{F}$ 上是非零因子），可知
$N = \lim H^0(U, \mathcal{F}/f^n\mathcal{F})$
是有限 $A$-模、无 $f$-挠，并且
$N/fN \subset H^0(U, \mathcal{F}/f\mathcal{F})$。
另一方面，有映射 $M \to N$ 以及相容映射
$$
M/fM \longrightarrow H^0(U, \mathcal{F}/f\mathcal{F})
$$
对 $g \in \mathfrak a$，可见 $(M/fM)_g$ 同构地映至
$H^0(U \cap D(f), \mathcal{F}/f\mathcal{F})$，因为
$\mathcal{F}/f\mathcal{F}$ 是 $\widetilde{M/fM}$ 在 $U$ 上的限制。
由此可知 $M_g \to N_g$ 诱导同构
$$
M_g/fM_g = (M/fM)_g \to (N/fN)_g = N_g/fN_g
$$
由于 $f$ 在 $N$ 和 $M$ 上都是非零因子，可知
$M_g \to N_g$ 在 $f$-进完备化上诱导同构，
进而蕴含 $M_g \to N_g$ 在 $V(f) \cap D(g)$ 的某个开邻域上为同构。
由于 $g \in \mathfrak a$ 是任意的，可知 $M$ 和 $N$
在引理陈述中所述的某个开集 $V$ 上确定同构的凝聚模。证明完毕。
\end{proof}

\begin{proposition}
\label{algebraization-proposition-alternative-colim-H0}
设 $A$ 为诺特环。设 $f \in \mathfrak a \subset A$
为 $A$ 的一个理想中的元素。设 $\mathcal{F}$ 为
$U = \Spec(A) \setminus V(\mathfrak a)$ 上的凝聚模。假设
\begin{enumerate}
\item $A$ 是 $f$-进完备的，并且有对偶化复形；
\item 若 $x \in \text{Ass}(\mathcal{F})$、$x \not \in V(f)$、
$\overline{\{x\}} \cap V(f) \not \subset V(\mathfrak a)$,
且 $z \in \overline{\{x\}} \cap V(\mathfrak a)$，则
$\dim(\mathcal{O}_{\overline{\{x\}}, z}) > 2$.
\end{enumerate}
则映射
$$
\colim_V \Gamma(V, \mathcal{F})
\longrightarrow
\lim \Gamma(U, \mathcal{F}/f^n\mathcal{F})
$$
为同构，其中余极限遍历开集 $V \subset U$，这些开集包含
$U \cap V(f)$。
\end{proposition}

\begin{proof}[第一种证明]
回忆 $A$ 是泛链状的，并且有 Gorenstein 形式纤维；参见对偶化复形，引理
\ref{dualizing-lemma-dualizing-gorenstein-formal-fibres} 和
\ref{dualizing-lemma-universally-catenary}。因此可以考察映射
$\mathcal{F} \to \mathcal{F}'$，它由局部上同调，引理
\ref{local-cohomology-lemma-make-S2-along-Z} 针对闭子集
$V(f) \cap U$（属于 $U$）构造。注意
\begin{enumerate}
\item $\mathcal{F} \to \mathcal{F}'$ 的核和余核支撑在 $V(f) \cap U$ 上。
\item 模 $\mathcal{F}'$ 无 $f$-挠，因为其茎的深度 $\geq 1$，
这对 $V(f) \cap U$ 的所有点成立；换言之，$\mathcal{F}'$
在 $V(f) \cap U$ 中没有相伴点。
\item 若 $y \in V(f) \cap U$ 是
$\mathcal{F}'/f\mathcal{F}'$ 的相伴点，则
$\text{depth}(\mathcal{F}'_y) = 1$
并且由 $\mathcal{F}'$ 的构造，存在立即特殊化 $x \leadsto y$，
其中 $x \not \in V(f)$ 是 $\mathcal{F}$ 的相伴点。
由假设 (2) 可知，$y$ 不可能在 $\Spec(A)$ 中立即特殊化到点
$z \in V(\mathfrak a)$。
\item 由 (3) 可知 $H^0(U, \mathcal{F}'/f\mathcal{F}')$
是有限 $A$-模；参见局部上同调，引理
\ref{local-cohomology-lemma-finiteness-Rjstar}。
\end{enumerate}
这些观察将使我们能够完成证明。

\medskip\noindent
首先，断言该引理对 $\mathcal{F}'$ 成立。
具体而言，选取有限 $A$-模 $M'$，使得 $\mathcal{F}'$ 是在 $U$ 上的限制，
所限制的凝聚模与 $M'$ 相伴；参见局部上同调，引理
\ref{local-cohomology-lemma-finiteness-pushforwards-and-H1-local}。
由于 $\mathcal{F}'$ 无 $f$-挠，可以假设 $M'$ 也无 $f$-挠。
上面的观察 (4) 表明 $H^1_\mathfrak a(M')$ 是有限 $A$-模；
参见局部上同调，引理
\ref{local-cohomology-lemma-finiteness-pushforwards-and-H1-local}。
因此该断言由引理 \ref{algebraization-lemma-alternative-colim-H0} 得到。

\medskip\noindent
其次，注意该引理对任意凝聚 $\mathcal{O}_U$-模都显然成立，
只要它支撑在 $V(f) \cap U$ 上。
令 $\mathcal{K}$、$\mathcal{G}$、$\mathcal{Q}$ 分别为映射
$\mathcal{F} \to \mathcal{F}'$ 的核、像、余核。短正合序列
$0 \to \mathcal{G} \to \mathcal{F}' \to \mathcal{Q} \to 0$
和引理 \ref{algebraization-lemma-ses-H0} 表明结论对 $\mathcal{G}$ 成立。
然后对短正合序列再次作同样处理：
$0 \to \mathcal{K} \to \mathcal{F} \to \mathcal{G} \to 0$
从而完成证明。
\end{proof}

\begin{proof}[第二种证明]
该命题是命题 \ref{algebraization-proposition-application-H0} 的特殊情形。
\end{proof}

\begin{lemma}
\label{algebraization-lemma-alternative-H0}
设 $A$ 为诺特环。设 $f \in \mathfrak a \subset A$
为 $A$ 的一个理想中的元素。设 $M$ 为有限 $A$-模。假设
\begin{enumerate}
\item $A$ 是 $f$-进完备的；
\item $H^1_\mathfrak a(M)$ 和 $H^2_\mathfrak a(M)$ 被 $f$ 的某个幂零化。
\end{enumerate}
则令 $U = \Spec(A) \setminus V(\mathfrak a)$，映射
$$
\Gamma(U, \widetilde{M})
\longrightarrow
\lim \Gamma(U, \widetilde{M/f^nM})
$$
为同构。
\end{lemma}

\begin{proof}
可以把引理 \ref{algebraization-lemma-formal-functions-principal} 应用于 $U$ 以及
$\mathcal{F} = \widetilde{M}|_U$，因为 $\mathcal{F}$ 是凝聚
$\mathcal{O}_U$-模范畴中的诺特对象。
由于 $H^1(U, \mathcal{F}) = H^2_\mathfrak a(M)$
（局部上同调，引理
\ref{local-cohomology-lemma-finiteness-pushforwards-and-H1-local}）
被 $f$ 的某个幂零化，可知其 $f$-进 Tate 模为零。
因此该引理表明 $\lim H^0(U, \mathcal{F}/f^n \mathcal{F})$
等于通常的 $f$-进完备化，所完备化的是 $H^0(U, \mathcal{F})$。
考察短正合序列
$$
0 \to M/H^0_\mathfrak a(M) \to H^0(U, \mathcal{F}) \to H^1_\mathfrak a(M) \to 0
$$
它来自局部上同调，引理
\ref{local-cohomology-lemma-finiteness-pushforwards-and-H1-local}。
由于 $M/H^0_\mathfrak a(M)$ 是有限 $A$-模，它是完备的；
参见代数，引理 \ref{algebra-lemma-completion-tensor}。
由于 $H^1_\mathfrak a(M)$ 被 $f$ 的某个幂零化，
由代数，引理 \ref{algebra-lemma-completion-differ-by-torsion} 可知
$H^0(U, \mathcal{F})$ 也完备。证明完毕。
\end{proof}







\section{形式截面的代数化，III}
\label{algebraization-section-algebraization-sections-coherent-III}

\noindent
下一节给出了一个非穷尽的应用列表：将局部上同调完备化方面的材料
应用于拟仿射概形上凝聚模的高阶上同调，以及这些上同调关于某个理想的完备化。

\begin{proposition}
\label{algebraization-proposition-application-higher}
设 $I \subset \mathfrak a$ 为诺特环 $A$ 的理想。
设 $\mathcal{F}$ 是
$U = \Spec(A) \setminus V(\mathfrak a)$ 上的凝聚模。
设 $s \geq 0$。假设
\begin{enumerate}
\item $A$ 是 $I$-进完备的，并且有对偶化复形；
\item 若 $x \in U \setminus V(I)$，则
$\text{depth}(\mathcal{F}_x) > s$，或者
$$
\text{depth}(\mathcal{F}_x) +
\dim(\mathcal{O}_{\overline{\{x\}}, z}) > \text{cd}(A, I) + s + 1
$$
对所有 $z \in V(\mathfrak a) \cap \overline{\{x\}}$ 成立；
\item 下列条件之一成立：
\begin{enumerate}
\item $\mathcal{F}$ 在 $U \setminus V(I)$ 上的限制满足 $(S_{s + 1})$；或者
\item $V(\mathfrak a)$ 的维数至多为 $2$\footnote{这里的含义是：
该函数在 $V(\mathfrak a)$ 上所取最大值与最小值之差（这里该维数函数定义在
$\Spec(A)$ 上）至多为 $2$。}。
\end{enumerate}
\end{enumerate}
则映射
$$
H^i(U, \mathcal{F})
\longrightarrow
\lim H^i(U, \mathcal{F}/I^n\mathcal{F})
$$
在 $i < s$ 时为同构。此外，还有同构
$$
\colim H^s(V, \mathcal{F})
\longrightarrow
\lim H^s(U, \mathcal{F}/I^n\mathcal{F})
$$
其中余极限取遍开子概形 $V \subset U$，这些开子概形包含 $U \cap V(I)$。
\end{proposition}

\begin{proof}
可以假设 $s > 0$，因为 $s = 0$ 的情形已在
命题 \ref{algebraization-proposition-application-H0} 中处理。

\medskip\noindent
取一个有限 $A$-模 $M$，使得 $\mathcal{F}$ 是某个凝聚模
在 $U$ 上的限制，而该凝聚模与 $M$ 相伴；参见局部上同调，引理
\ref{local-cohomology-lemma-finiteness-pushforwards-and-H1-local}。
令 $d = \text{cd}(A, I)$。
设 $\mathfrak p$ 是 $A$ 的一个不包含在 $V(I)$ 中的素理想，
并设 $\mathfrak q \in V(\mathfrak p) \cap V(\mathfrak a)$。
于是或者 $\text{depth}(M_\mathfrak p) \geq s + 1 > s$，
或者由 (2) 有 $\dim((A/\mathfrak p)_\mathfrak q) > d + s + 1$。
由引理 \ref{algebraization-lemma-bootstrap-bis-bis} 可知，
对于 $A, I, V(\mathfrak a), M, s, d$，情形 \ref{algebraization-situation-bootstrap}
的假设均成立。另一方面，引理
\ref{algebraization-lemma-algebraize-local-cohomology-general}
对于 $s + 1$ 和 $d$ 的假设也成立；条件 (3) 正是在这里使用的。

\medskip\noindent
应用引理 \ref{algebraization-lemma-algebraize-local-cohomology-general}，
可得存在理想
$J_0 \subset \mathfrak a$，满足 $V(J_0) \cap V(I) = V(\mathfrak a)$，
使得对于任何 $J \subset J_0$，若 $V(J) \cap V(I) = V(\mathfrak a)$，
映射
$$
H^i_J(M) \longrightarrow H^i(R\Gamma_\mathfrak a(M)^\wedge)
$$
在 $i \leq s + 1$ 时为同构。

\medskip\noindent
当 $i \leq s$ 时，由引理 \ref{algebraization-lemma-bootstrap-inherited} 和
\ref{algebraization-lemma-kill-colimit-support-general}，映射
$H^i_\mathfrak a(M) \to H^i_J(M)$ 为同构。
利用上同调与局部上同调之间的比较
（局部上同调，引理 \ref{local-cohomology-lemma-local-cohomology}），
可推出映射 $H^i(U, \mathcal{F}) \to H^i(V,\mathcal{F})$
在 $V = \Spec(A) \setminus V(J)$ 且 $i < s$ 时为同构。

\medskip\noindent
由定理 \ref{algebraization-theorem-final-bootstrap}，等式
$H^i_\mathfrak a(M) = \lim H^i_\mathfrak a(M/I^nM)$
对 $i \leq s$ 成立。
由引理 \ref{algebraization-lemma-combine-two}，有
$H^{s + 1}_\mathfrak a(M) = \lim H^{s + 1}_\mathfrak a(M/I^nM)$。

\medskip\noindent
由以上结论和命题 \ref{algebraization-proposition-application-H0}，得到同构
$H^0(U, \mathcal{F}) = H^0(V, \mathcal{F}) =
\lim H^0(U, \mathcal{F}/I^n\mathcal{F})$。
当 $0 < i < s$ 时，利用局部上同调与上同调之间的关系，
以相同方式得到所需同构
$H^i(U, \mathcal{F}) = H^i(V, \mathcal{F}) =
\lim H^i(U, \mathcal{F}/I^n\mathcal{F})$；
这比 $i = 0$ 的情形更容易，因为当 $i > 0$ 时有
$$
H^i(U, \mathcal{F}) = H^{i + 1}_\mathfrak a(M),
\quad
H^i(V, \mathcal{F}) = H^{i + 1}_J(M),
\quad
H^i(R\Gamma(U, \mathcal{F})^\wedge) = 
H^{i + 1}(R\Gamma_\mathfrak a(M)^\wedge)
$$
最后一个断言的证明也类似。
\end{proof}

\begin{lemma}
\label{algebraization-lemma-alternative-higher}
设 $A$ 为诺特环。设 $f \in \mathfrak a \subset A$
是 $A$ 的一个理想中的元素。设 $M$ 为有限 $A$-模。
设 $s \geq 0$。假设
\begin{enumerate}
\item $A$ 是 $f$-进完备的；
\item $H^i_\mathfrak a(M)$ 被 $f$ 的某个幂零化，其中 $i \leq s + 1$。
\end{enumerate}
则令 $U = \Spec(A) \setminus V(\mathfrak a)$，映射
$$
H^i(U, \widetilde{M})
\longrightarrow
\lim H^i(U, \widetilde{M/f^nM})
$$
在 $i < s$ 时为同构。
\end{lemma}

\begin{proof}
对 $s$ 作归纳。若 $s = 0$，则断言为空。若 $s = 1$，则结论即为
引理 \ref{algebraization-lemma-alternative-H0}。假设 $s > 1$。由归纳法，
只需对 $i = s - 1 \geq 1$ 证明结论。
可以把引理 \ref{algebraization-lemma-formal-functions-principal}
应用于 $U$ 和 $\mathcal{F} = \widetilde{M}|_U$，
因为 $\mathcal{F}$ 是凝聚 $\mathcal{O}_U$-模范畴中的诺特对象。
由局部上同调，引理
\ref{local-cohomology-lemma-finiteness-pushforwards-and-H1-local}，
有 $H^j(U, \mathcal{F}) = H^{j + 1}_\mathfrak a(M)$，其中 $j$ 任意。
因此，当 $j = s = (s - 1) + 1$ 时，按假设该模被 $f$ 的某个幂零化。
于是由引理 \ref{algebraization-lemma-formal-functions-principal}，
$\lim H^{s - 1}(U, \mathcal{F}/f^n\mathcal{F})$
是关于 $f$ 的通常进完备化，所完备化的模为
$H^{s - 1}(U, \mathcal{F})$。
再次利用该模被 $f$ 的某个幂零化，可见此完备化恰好等于
$H^{s - 1}(U, \mathcal{F})$，这正是所需结论。
\end{proof}





\section{连通性的应用}
\label{algebraization-section-connected}

\noindent
本节讨论 Grothendieck 连通性定理及其变体；原始版本见
\cite[Exposee XIII, Theorem 2.1]{SGA2}。文献中有一个称为
Faltings 连通性定理的版本；我们猜测这指的是
\cite[Theorem 6]{Faltings-some}。下面陈述并证明
\cite[Theorem 1.6]{Varbaro} 中关于完备局部环的最优版本。

\begin{lemma}
\label{algebraization-lemma-punctured-still-connected}
\begin{reference}
\cite[Theorem 1.6]{Varbaro}
\end{reference}
设 $(A, \mathfrak m)$ 为诺特完备局部环。
设 $I$ 为 $A$ 的真理想。
令 $X = \Spec(A)$ 且 $Y = V(I)$。记
\begin{enumerate}
\item $d$ 为 $X$ 的不可约分支的维数最小值；
\item $c$ 为闭子集 $Z \subset X$ 的维数最小值，其中
$X \setminus Z$ 不连通。
\end{enumerate}
则对于闭子集 $Z \subset Y$，$Y \setminus Z$ 在
$\dim(Z) < \min(c, d - 1) - \text{cd}(A, I)$ 时连通。
特别地，$A/I$ 的穿孔谱在 $\text{cd}(A, I) < \min(c, d - 1)$ 时连通。
\end{lemma}

\begin{proof}
先证明最后一个断言。首先，若 $A/I$ 的穿孔谱为空，则由
局部上同调，引理 \ref{local-cohomology-lemma-cd-bound-dim-local}，
$X$ 的每个不可约分支的维数均
$\leq \text{cd}(A, I)$，从而 $\min(c, d - 1) - \text{cd}(A, I) < 0$；
这说明该引理在此情形下成立。因此可以假设
$U \cap Y$ 非空，其中 $U = X \setminus \{\mathfrak m\}$
是 $A$ 的穿孔谱。可以将 $A$ 替换为其既约化。
注意 $A$ 有对偶化复形
（对偶化复形，引理 \ref{dualizing-lemma-ubiquity-dualizing}），
并且 $A$ 关于 $I$ 完备
（代数，引理 \ref{algebra-lemma-complete-by-sub}）。
若假设 $d - 1 > \text{cd}(A, I)$，则应用
引理 \ref{algebraization-lemma-application-theorem} 可得
$$
\colim H^0(V, \mathcal{O}_V)
\longrightarrow
\lim H^0(U, \mathcal{O}_U/I^n\mathcal{O}_U)
$$
为同构，其中余极限取遍开子概形 $V \subset U$，且这些开子概形包含
$U \cap Y$。若 $U \cap Y$ 不连通，则它的第 $n$ 个无穷小邻域在
$U$ 中对所有 $n$ 都不连通，因而右端有一个非平凡幂等元
（这里使用了 $U \cap Y$ 非空这一点）。
因此可以找到不连通的 $V$。
令 $Z = X \setminus V$。由局部上同调，引理
\ref{local-cohomology-lemma-cd-bound-dim-local}，
可知 $Z$ 的每个不可约分支的维数均
$\leq \text{cd}(A, I)$。因此 $c \leq \text{cd}(A, I)$，
这确实证明了最后一个断言。

\medskip\noindent
现在可以从刚刚证明的结论推出引理的断言。设 $Z \subset Y$ 为闭子集，
$Y \setminus Z$ 不连通，且 $\dim(Z) = e$。回忆按照约定，连通空间非空。
因此，或者 (a) $Y = Z$，或者 (b)
$Y \setminus Z = W_1 \amalg W_2$，其中 $W_i$ 在 $Y \setminus Z$ 中
非空且既开又闭。在情形 (b) 中，可以选取 $w_i \in W_i$，
使其在 $U$ 中为闭点；参见态射，引理
\ref{morphisms-lemma-ubiquity-Jacobson-schemes}。
于是可以找到 $f_1, \ldots, f_e \in \mathfrak m$，使得
$V(f_1, \ldots, f_e) \cap Z = \{\mathfrak m\}$，
并且在情形 (b) 中可以假设 $w_i \in V(f_1, \ldots, f_e)$。
具体而言，利用素理想回避，可以归纳地选取 $f_i$，使得
$\dim V(f_1, \ldots, f_i) \cap Z = e - i$，
并使得在情形 (b) 中有 $w_1, w_2 \in V(f_i)$。
由此可知 $A/I + (f_1, \ldots, f_e)$ 的穿孔谱不连通
（略去一个小细节）。由于由局部上同调，引理
\ref{local-cohomology-lemma-cd-sum} 和
\ref{local-cohomology-lemma-bound-cd} 有
$\text{cd}(A, I + (f_1, \ldots, f_e)) \leq \text{cd}(A, I) + e$，故
$$
\text{cd}(A, I) + e \geq \min(c, d - 1)
$$
由证明的第一部分可得此式。这推出
$e \geq \min(c, d - 1) - \text{cd}(A, I)$，这正是所需结论。
\end{proof}

\begin{lemma}
\label{algebraization-lemma-connected}
设 $I \subset \mathfrak a$ 为诺特环 $A$ 的理想。假设
\begin{enumerate}
\item $A$ 是 $I$-进完备的，并且有对偶化复形；
\item 若 $\mathfrak p \subset A$ 是一个不包含在
$V(I)$ 中的极小素理想，并且 $\mathfrak q \in V(\mathfrak p) \cap V(\mathfrak a)$，则
$\dim((A/\mathfrak p)_\mathfrak q) > \text{cd}(A, I) + 1$；
\item 任何非空开子集 $V \subset \Spec(A)$，只要其包含
$V(I) \setminus V(\mathfrak a)$，就是连通的\footnote{例如，当 $A$ 为整环时。}。
\end{enumerate}
则 $V(I) \setminus V(\mathfrak a)$ 或为空，或为连通的。
\end{lemma}

\begin{proof}
可以将 $A$ 替换为其既约化。于是对于 $A$ 的所有相伴素理想，
(2) 中的不等式均成立。由命题 \ref{algebraization-proposition-application-H0} 可知
$$
\colim H^0(V, \mathcal{O}_V) = \lim H^0(T_n, \mathcal{O}_{T_n})
$$
其中余极限取遍 (3) 中的开子集 $V$，而 $T_n$ 表示第 $n$ 个
无穷小邻域；这是 $T = V(I) \setminus V(\mathfrak a)$ 在
$U = \Spec(A) \setminus V(\mathfrak a)$ 中的无穷小邻域。
因此 $T$ 或为空，或为连通的；否则右端将有非平凡幂等元，
而我们已假设左端没有。略去若干细节。
\end{proof}

\begin{lemma}
\label{algebraization-lemma-H0}
设 $A$ 为有对偶化复形的诺特整环，并且关于某个非零元
$f \in A$ 完备。设 $f \in \mathfrak a \subset A$，并把这里的中间项取为理想。
假设 $Z = V(\mathfrak a)$ 的每个不可约分支的
余维均 $> 2$，这里是在 $X = \Spec(A)$ 中计算；也就是说，假设 $Z$ 的每个
不可约分支的余维均 $> 1$，这里是在 $Y = V(f)$ 中计算。则 $Y \setminus Z$ 连通。
\end{lemma}

\begin{proof}
这是引理 \ref{algebraization-lemma-connected} 的特殊情形（该引理的证明依赖
命题 \ref{algebraization-proposition-application-H0}）。下面用较容易的命题
\ref{algebraization-proposition-alternative-colim-H0} 来证明它。

\medskip\noindent
令 $U = X \setminus Z$。由命题 \ref{algebraization-proposition-alternative-colim-H0}，
有同构
$$
\colim \Gamma(V, \mathcal{O}_V) \to
\lim_n \Gamma(U, \mathcal{O}_U/f^n \mathcal{O}_U)
$$
其中余极限取遍开子集 $V \subset U$，这些开子集包含 $U \cap Y$。
因此，若 $U \cap Y$ 不连通，则对于某个 $V$，
$\Gamma(V, \mathcal{O}_V)$ 中存在非平凡幂等元。这不可能，
因为 $V$ 是整概形，而 $X$ 是整环的谱。
\end{proof}






\section{完备化函子}
\label{algebraization-section-completion}

\noindent
设 $X$ 为诺特概形。设 $Y \subset X$ 为闭子概形，
其拟凝聚理想层为 $\mathcal{I} \subset \mathcal{O}_X$。
本节考虑凝聚 $\mathcal{O}_X$-模 $(\mathcal{F}_n)$ 的逆系统，其中
$\mathcal{F}_n$ 被 $I^n$ 零化，并且转移映射诱导同构
$\mathcal{F}_{n + 1}/I^n\mathcal{F}_{n + 1} \to \mathcal{F}_n$。
这些逆系统所成的范畴在概形上同调，第
\ref{coherent-section-coherent-formal} 节中记为
$$
\textit{Coh}(X, \mathcal{I})
$$
。这个范畴等价于 $X$ 沿 $Y$ 的形式完备化上的凝聚模范畴；
但是，我们还没有引入形式概形或其上的凝聚模，
所以这里不能使用这一术语。我们特别关注完备化函子
$$
\textit{Coh}(\mathcal{O}_X)
\longrightarrow
\textit{Coh}(X, \mathcal{I}),\quad
\mathcal{F} \longmapsto \mathcal{F}^\wedge
$$
；参见概形上同调，等式 (\ref{coherent-equation-completion-functor})。

\begin{lemma}
\label{algebraization-lemma-completion-fully-faithful}
设 $X$ 为诺特概形，$Y \subset X$ 为闭子概形。
设 $Y_n \subset X$ 为第 $n$ 个无穷小邻域；这是 $Y$ 在 $X$ 中的邻域。
考虑下列条件：
\begin{enumerate}
\item $X$ 是拟仿射的，并且
$\Gamma(X, \mathcal{O}_X) \to \lim \Gamma(Y_n, \mathcal{O}_{Y_n})$
为同构；
\item $X$ 有充足可逆模 $\mathcal{L}$，并且
$\Gamma(X, \mathcal{L}^{\otimes m}) \to
\lim \Gamma(Y_n, \mathcal{L}^{\otimes m}|_{Y_n})$
对所有 $m \gg 0$ 都是同构；
\item 对每个有限局部自由 $\mathcal{O}_X$-模
$\mathcal{E}$，映射
$\Gamma(X, \mathcal{E}) \to \lim \Gamma(Y_n, \mathcal{E}|_{Y_n})$
为同构；
\item 完备化函子
$\textit{Coh}(\mathcal{O}_X) \to \textit{Coh}(X, \mathcal{I})$
在有限局部自由对象所成的全子范畴上是全忠实的。
\end{enumerate}
则 (1) $\Rightarrow$ (2) $\Rightarrow$ (3) $\Rightarrow$ (4)，
并且 (4) $\Rightarrow$ (3)。
\end{lemma}

\begin{proof}
证明 (3) $\Rightarrow$ (4)。若 $\mathcal{F}$ 和 $\mathcal{G}$
在 $X$ 上有限局部自由，则考虑
$\mathcal{H} = \SheafHom_{\mathcal{O}_X}(\mathcal{G}, \mathcal{F})$，
并使用概形上同调，引理
\ref{coherent-lemma-completion-internal-hom}，
可知 (3) 推出 (4)。

\medskip\noindent
证明 (2) $\rightarrow$ (3)。设 $\mathcal{L}$ 在 $X$ 上充足，
并设 $\mathcal{E}$ 为有限局部自由 $\mathcal{O}_X$-模。
我们断言可以找到泛正合序列
$$
0 \to \mathcal{E} \to
(\mathcal{L}^{\otimes p})^{\oplus r} \to
(\mathcal{L}^{\otimes q})^{\oplus s}
$$
，其中某些 $r, s \geq 0$ 且 $0 \ll p \ll q$。若此断言成立，则使用正合序列
$$
0 \to \lim \Gamma(\mathcal{E}|_{Y_n}) \to
\lim \Gamma((\mathcal{L}^{\otimes p})^{\oplus r}|_{Y_n}) \to
\lim \Gamma((\mathcal{L}^{\otimes q})^{\oplus s}|_{Y_n})
$$
以及 (2) 中的同构，就得到 (3) 中的同构。
为证明该断言，考虑对偶局部自由模
$\SheafHom_{\mathcal{O}_X}(\mathcal{E}, \mathcal{O}_X)$
，并应用性质，命题 \ref{properties-proposition-characterize-ample}，
得到满射
$$
(\mathcal{L}^{\otimes -p})^{\oplus r}
\longrightarrow
\SheafHom_{\mathcal{O}_X}(\mathcal{E}, \mathcal{O}_X)
$$
取对偶便得到正合序列中的第一个映射
（它是泛单射，因为满射性是泛的）。对余核重复这一过程便得到第二个映射。
略去若干细节。

\medskip\noindent
证明 (1) $\Rightarrow$ (2)。这是因为若 $X$ 是拟仿射的，
则 $\mathcal{O}_X$ 是充足可逆模；参见性质，引理
\ref{properties-lemma-quasi-affine-O-ample}。

\medskip\noindent
略去 (4) $\Rightarrow$ (3) 的证明。
\end{proof}

\noindent
给定诺特概形和拟凝聚理想层
$\mathcal{I} \subset \mathcal{O}_X$，若对象 $(\mathcal{F}_n)$ 属于
$\textit{Coh}(X, \mathcal{I})$，且每个 $\mathcal{F}_n$ 都是有限局部自由
$\mathcal{O}_X/\mathcal{I}^n$-模，则称该对象为{\it 有限局部自由的}。

\begin{lemma}
\label{algebraization-lemma-completion-fully-faithful-general}
设 $X$ 为诺特概形，$Y \subset X$ 为闭子概形，
其理想层为 $\mathcal{I} \subset \mathcal{O}_X$。
设 $Y_n \subset X$ 为第 $n$ 个无穷小邻域；这是 $Y$ 在 $X$ 中的邻域。
设 $\mathcal{V}$ 为开子概形 $V \subset X$ 所成的集合，这些开子概形包含
$Y$；该集合按反包含排序。
\begin{enumerate}
\item $X$ 是拟仿射的，并且
$$
\colim_\mathcal{V} \Gamma(V, \mathcal{O}_V)
\longrightarrow
\lim \Gamma(Y_n, \mathcal{O}_{Y_n})
$$
为同构；
\item $X$ 有充足可逆模 $\mathcal{L}$，并且
$$
\colim_\mathcal{V} \Gamma(V, \mathcal{L}^{\otimes m})
\longrightarrow
\lim \Gamma(Y_n, \mathcal{L}^{\otimes m}|_{Y_n})
$$
对所有 $m \gg 0$ 都是同构；
\item 对每个 $V \in \mathcal{V}$ 和每个有限局部自由
$\mathcal{O}_V$-模 $\mathcal{E}$，映射
$$
\colim_{V' \geq V} \Gamma(V', \mathcal{E}|_{V'})
\longrightarrow
\lim \Gamma(Y_n, \mathcal{E}|_{Y_n})
$$
为同构；
\item 完备化函子
$$
\colim_\mathcal{V} \textit{Coh}(\mathcal{O}_V)
\longrightarrow
\textit{Coh}(X, \mathcal{I}),
\quad
\mathcal{F} \longmapsto \mathcal{F}^\wedge
$$
在有限局部自由对象所成的全子范畴上是全忠实的
（参见上面的说明）。
\end{enumerate}
则 (1) $\Rightarrow$ (2) $\Rightarrow$ (3) $\Rightarrow$ (4)，
并且 (4) $\Rightarrow$ (3)。
\end{lemma}

\begin{proof}
注意 $\mathcal{V}$ 是有向集，所以这些余极限就是范畴，第
\ref{categories-section-directed-colimits} 节中的余极限。
其余论证几乎与引理 \ref{algebraization-lemma-completion-fully-faithful} 的证明
中的论证完全相同；我们建议读者跳过它。

\medskip\noindent
证明 (3) $\Rightarrow$ (4)。若 $\mathcal{F}$ 和 $\mathcal{G}$
在 $V \in \mathcal{V}$ 上有限局部自由，则考虑
$\mathcal{H} = \SheafHom_{\mathcal{O}_V}(\mathcal{G}, \mathcal{F})$，
并使用概形上同调，引理
\ref{coherent-lemma-completion-internal-hom}，
可知 (3) 推出 (4)。

\medskip\noindent
证明 (2) $\Rightarrow$ (3)。设 $\mathcal{L}$ 在 $X$ 上充足，
并设 $\mathcal{E}$ 是有限局部自由 $\mathcal{O}_V$-模，
其中某个 $V \in \mathcal{V}$。
我们断言可以找到泛正合序列
$$
0 \to \mathcal{E} \to
(\mathcal{L}^{\otimes p})^{\oplus r}|_{V} \to
(\mathcal{L}^{\otimes q})^{\oplus s}|_{V}
$$
，其中某些 $r, s \geq 0$ 且 $0 \ll p \ll q$。若此断言成立，
则 (2) 中的同构将推出 (3) 中的同构。
为证明该断言，回忆 $\mathcal{L}|_V$ 是充足的；参见性质，引理
\ref{properties-lemma-ample-on-locally-closed}。
考虑对偶局部自由模
$\SheafHom_{\mathcal{O}_V}(\mathcal{E}, \mathcal{O}_V)$
，并应用性质，命题 \ref{properties-proposition-characterize-ample}，
得到满射
$$
(\mathcal{L}^{\otimes -p})^{\oplus r}|_V \longrightarrow
\SheafHom_{\mathcal{O}_V}(\mathcal{E}, \mathcal{O}_V)
$$
（它是泛单射，因为满射性是泛的）。
取对偶便得到正合序列中的第一个映射。
对余核重复这一过程便得到第二个映射。略去若干细节。

\medskip\noindent
证明 (1) $\Rightarrow$ (2)。这是因为若 $X$ 是拟仿射的，
则 $\mathcal{O}_X$ 是充足可逆模；参见性质，引理
\ref{properties-lemma-quasi-affine-O-ample}。

\medskip\noindent
略去 (4) $\Rightarrow$ (3) 的证明。
\end{proof}

\begin{lemma}
\label{algebraization-lemma-recognize-formal-coherent-modules}
设 $X$ 为诺特概形。设 $\mathcal{I} \subset \mathcal{O}_X$
为拟凝聚理想层。函子
$$
\textit{Coh}(X, \mathcal{I}) \longrightarrow \text{Pro-}\QCoh(\mathcal{O}_X)
$$
是全忠实的；参见范畴，注记 \ref{categories-remark-pro-category}。
\end{lemma}

\begin{proof}
设 $(\mathcal{F}_n)$ 和 $(\mathcal{G}_n)$ 为
$\textit{Coh}(X, \mathcal{I})$ 的对象。pro-对象态射 $\alpha$ 从
$(\mathcal{F}_n)$ 到 $(\mathcal{G}_n)$，它由一族与转移映射相容的映射
$\alpha_n : \mathcal{F}_{m(n)} \to \mathcal{G}_n$ 给出，其中
$\mathbf{N} \to \mathbf{N}$，$n \mapsto m(n)$ 是递增函数
（特别地，$m(n) \geq n$）。由于
$\mathcal{F}_n = \mathcal{F}_{m(n)}/\mathcal{I}^n\mathcal{F}_{m(n)}$，
并且 $\mathcal{G}_n$ 被 $\mathcal{I}^n$ 零化，
可知 $\alpha_n$ 诱导映射 $\mathcal{F}_n \to \mathcal{G}_n$。
\end{proof}

\noindent
下面给出若干例子，说明利用本章前面完成的工作能够证明何种全忠实性结果。

\begin{lemma}
\label{algebraization-lemma-fully-faithful}
设 $I \subset \mathfrak a$ 为诺特环 $A$ 的理想。
令 $U = \Spec(A) \setminus V(\mathfrak a)$。假设
\begin{enumerate}
\item $A$ 是 $I$-进完备的，并且有对偶化复形；
\item 对任意满足下列条件的相伴素理想 $\mathfrak p \subset A$：
$\mathfrak p \not \in V(I)$、
$V(\mathfrak p) \cap V(I) \not \subset V(\mathfrak a)$，且
$\mathfrak q \in V(\mathfrak p) \cap V(\mathfrak a)$，都有
$\dim((A/\mathfrak p)_\mathfrak q) > \text{cd}(A, I) + 1$；
\item 对 $\mathfrak p \subset A$，若 $\mathfrak p \not \in V(I)$ 且
$V(\mathfrak p) \cap V(I) \subset V(\mathfrak a)$，
有 $\text{depth}(A_\mathfrak p) \geq 2$。
\end{enumerate}
则完备化函子
$$
\textit{Coh}(\mathcal{O}_U)
\longrightarrow
\textit{Coh}(U, I\mathcal{O}_U),
\quad
\mathcal{F} \longmapsto \mathcal{F}^\wedge
$$
在有限局部自由对象所成的全子范畴上是全忠实的。
\end{lemma}

\begin{proof}
由引理 \ref{algebraization-lemma-completion-fully-faithful}，
只需证明
$$
\Gamma(U, \mathcal{O}_U) =
\lim \Gamma(U, \mathcal{O}_U/I^n\mathcal{O}_U)
$$
。这立即由引理 \ref{algebraization-lemma-application-H0} 得到。
\end{proof}

\begin{lemma}
\label{algebraization-lemma-fully-faithful-alternative}
设 $A$ 为诺特环。设 $f \in \mathfrak a \subset A$
为 $A$ 的一个理想中的元素。令 $U = \Spec(A) \setminus V(\mathfrak a)$。
假设
\begin{enumerate}
\item $A$ 是 $f$-进完备的；
\item $H^1_\mathfrak a(A)$ 和 $H^2_\mathfrak a(A)$ 被 $f$ 的某个幂零化。
\end{enumerate}
则完备化函子
$$
\textit{Coh}(\mathcal{O}_U)
\longrightarrow
\textit{Coh}(U, I\mathcal{O}_U),
\quad
\mathcal{F} \longmapsto \mathcal{F}^\wedge
$$
在有限局部自由对象所成的全子范畴上是全忠实的。
\end{lemma}

\begin{proof}
由引理 \ref{algebraization-lemma-completion-fully-faithful}，
只需证明
$$
\Gamma(U, \mathcal{O}_U) =
\lim \Gamma(U, \mathcal{O}_U/I^n\mathcal{O}_U)
$$
。这立即由引理 \ref{algebraization-lemma-alternative-H0} 得到。
\end{proof}

\begin{lemma}
\label{algebraization-lemma-fully-faithful-simple-two}
设 $A$ 为诺特环。设 $f \in \mathfrak a$ 为 $A$ 的一个理想中的元素。
令 $U = \Spec(A) \setminus V(\mathfrak a)$。假设
\begin{enumerate}
\item $A$ 有对偶化复形，并且关于 $f$ 完备；
\item 对每个素理想 $\mathfrak p \subset A$，若 $f \not \in \mathfrak p$，
且 $\mathfrak q \in V(\mathfrak p) \cap V(\mathfrak a)$，则
$\text{depth}(A_\mathfrak p) + \dim((A/\mathfrak p)_\mathfrak q) > 2$。
\end{enumerate}
则完备化函子
$$
\textit{Coh}(\mathcal{O}_U)
\longrightarrow
\textit{Coh}(U, I\mathcal{O}_U),
\quad
\mathcal{F} \longmapsto \mathcal{F}^\wedge
$$
在有限局部自由对象所成的全子范畴上是全忠实的。
\end{lemma}

\begin{proof}
结论由引理 \ref{algebraization-lemma-fully-faithful-alternative} 和
局部上同调，命题 \ref{local-cohomology-proposition-annihilator} 得到。
\end{proof}

\begin{lemma}
\label{algebraization-lemma-fully-faithful-general}
设 $I \subset \mathfrak a \subset A$ 为诺特环 $A$ 的理想。
令 $U = \Spec(A) \setminus V(\mathfrak a)$。设 $\mathcal{V}$ 为
$U$ 的开子概形所成的集合，这些开子概形包含 $U \cap V(I)$，
并按反包含排序。假设
\begin{enumerate}
\item $A$ 是 $I$-进完备的，并且有对偶化复形；
\item 对任意满足下列条件的相伴素理想
$\mathfrak p \subset A$：
$I \not \subset \mathfrak p$，
$V(\mathfrak p) \cap V(I) \not \subset V(\mathfrak a)$，
且 $\mathfrak q \in V(\mathfrak p) \cap V(\mathfrak a)$，都有
$\dim((A/\mathfrak p)_\mathfrak q) > \text{cd}(A, I) + 1$。
\end{enumerate}
则完备化函子
$$
\colim_\mathcal{V} \textit{Coh}(\mathcal{O}_V)
\longrightarrow
\textit{Coh}(U, I\mathcal{O}_U),
\quad
\mathcal{F} \longmapsto \mathcal{F}^\wedge
$$
在有限局部自由对象所成的全子范畴上是全忠实的。
\end{lemma}

\begin{proof}
由引理 \ref{algebraization-lemma-completion-fully-faithful-general}，
只需证明
$$
\colim_\mathcal{V} \Gamma(V, \mathcal{O}_V) =
\lim \Gamma(U, \mathcal{O}_U/I^n\mathcal{O}_U)
$$
。这立即由命题 \ref{algebraization-proposition-application-H0} 得到。
\end{proof}

\begin{lemma}
\label{algebraization-lemma-fully-faithful-general-alternative}
设 $A$ 为诺特环。设 $f \in \mathfrak a \subset A$
为 $A$ 的一个理想中的元素。令 $U = \Spec(A) \setminus V(\mathfrak a)$。
设 $\mathcal{V}$ 为 $U$ 的开子概形所成的集合，这些开子概形包含
$U \cap V(f)$，并按反包含排序。假设
\begin{enumerate}
\item $A$ 是 $f$-进完备的；
\item $f$ 是非零因子；
\item $H^1_\mathfrak a(A/fA)$ 是有限 $A$-模。
\end{enumerate}
则完备化函子
$$
\colim_\mathcal{V} \textit{Coh}(\mathcal{O}_V)
\longrightarrow
\textit{Coh}(U, f\mathcal{O}_U),
\quad
\mathcal{F} \longmapsto \mathcal{F}^\wedge
$$
在有限局部自由对象所成的全子范畴上是全忠实的。
\end{lemma}

\begin{proof}
由引理 \ref{algebraization-lemma-completion-fully-faithful-general}，
只需证明
$$
\colim_\mathcal{V} \Gamma(V, \mathcal{O}_V) =
\lim \Gamma(U, \mathcal{O}_U/I^n\mathcal{O}_U)
$$
。这立即由引理 \ref{algebraization-lemma-alternative-colim-H0} 得到。
\end{proof}

\begin{lemma}
\label{algebraization-lemma-fully-faithful-very-general}
设 $I \subset \mathfrak a \subset A$ 为诺特环 $A$ 的理想。
令 $U = \Spec(A) \setminus V(\mathfrak a)$。设 $\mathcal{V}$ 为
$U$ 的开子概形所成的集合，这些开子概形包含 $U \cap V(I)$，
并按反包含排序。设 $\mathcal{F}$ 和
$\mathcal{G}$ 是凝聚 $\mathcal{O}_V$-模，其中某个 $V \in \mathcal{V}$。
映射
$$
\colim_{V' \geq V} \Hom_V(\mathcal{G}|_{V'}, \mathcal{F}|_{V'})
\longrightarrow
\Hom_{\textit{Coh}(U, I\mathcal{O}_U)}(\mathcal{G}^\wedge, \mathcal{F}^\wedge)
$$
在下列假设成立时是双射：
\begin{enumerate}
\item $A$ 是 $I$-进完备的，并且有对偶化复形；
\item 若 $x \in \text{Ass}(\mathcal{F})$，$x \not \in V(I)$，
$\overline{\{x\}} \cap V(I) \not \subset V(\mathfrak a)$，
且 $z \in \overline{\{x\}} \cap V(\mathfrak a)$，则
$\dim(\mathcal{O}_{\overline{\{x\}}, z}) > \text{cd}(A, I) + 1$。
\end{enumerate}
\end{lemma}

\begin{proof}
可以选取凝聚 $\mathcal{O}_U$-模
$\mathcal{F}'$ 和 $\mathcal{G}'$，它们在 $V$ 上的限制
分别为 $\mathcal{F}$ 和 $\mathcal{G}$；参见性质，引理
\ref{properties-lemma-lift-finite-presentation}。
可以修改对 $\mathcal{F}'$ 的选择，以保证
$\text{Ass}(\mathcal{F}') \subset V$；例如参见局部上同调，引理
\ref{local-cohomology-lemma-get-depth-1-along-Z}。
因此可以而且确实把 $V$ 换成 $U$，并把 $\mathcal{F}$ 和 $\mathcal{G}$
换成 $\mathcal{F}'$ 和 $\mathcal{G}'$。
令 $\mathcal{H} = \SheafHom_{\mathcal{O}_U}(\mathcal{G}, \mathcal{F})$。
这是凝聚 $\mathcal{O}_U$-模。我们有
$$
\Hom_V(\mathcal{G}|_V, \mathcal{F}|_V) =
H^0(V, \mathcal{H})
\quad\text{且}\quad
\lim H^0(U, \mathcal{H}/\mathcal{I}^n\mathcal{H}) =
\Mor_{\textit{Coh}(U, I\mathcal{O}_U)}
(\mathcal{G}^\wedge, \mathcal{F}^\wedge)
$$
；参见概形上同调，引理 \ref{coherent-lemma-completion-internal-hom}。
因此，如果能够证明命题 \ref{algebraization-proposition-application-H0}
的假设对 $\mathcal{H}$ 成立，那么证明就完成了。
这是因为
$\text{Ass}(\mathcal{H}) \subset \text{Ass}(\mathcal{F})$；
参见概形上同调，引理
\ref{coherent-lemma-hom-into-depth}。
\end{proof}








\section{凝聚形式模的代数化，I}
\label{algebraization-section-algebraization-modules}

\noindent
完备化函子的本质满性（见下文）在
\cite{SGA2}、\cite{MRaynaud-book} 和 \cite{MRaynaud-paper}
中得到了系统研究。我们在下列仿射情形中工作。

\begin{situation}
\label{algebraization-situation-algebraize}
这里 $A$ 是诺特环，$I \subset \mathfrak a \subset A$ 是理想。
令 $X = \Spec(A)$，$Y = V(I) = \Spec(A/I)$，并令
$Z = V(\mathfrak a) = \Spec(A/\mathfrak a)$。此外，$U = X \setminus Z$。
\end{situation}

\noindent
本节试图寻找条件，以保证 $\textit{Coh}(U, I\mathcal{O}_U)$ 中的对象
位于完备化函子
$\textit{Coh}(\mathcal{O}_U) \to \textit{Coh}(U, I\mathcal{O}_U)$
的像中。参见概形上同调，第 \ref{coherent-section-coherent-formal} 节和
第 \ref{algebraization-section-completion} 节。

\begin{lemma}
\label{algebraization-lemma-system-of-modules}
在情形 \ref{algebraization-situation-algebraize} 中，考虑满足下列条件的
逆系统 $(M_n)$，它由 $A$-模组成：
\begin{enumerate}
\item $M_n$ 是有限 $A$-模；
\item $M_n$ 被 $I^n$ 零化；
\item $M_{n + 1}/I^nM_{n + 1} \to M_n$ 的核与余核
都是 $\mathfrak a$-幂挠的。
\end{enumerate}
则 $(\widetilde{M}_n|_U)$ 属于 $\textit{Coh}(U, I\mathcal{O}_U)$。
反之，$\textit{Coh}(U, I\mathcal{O}_U)$ 的每个对象都以这种方式得到。
\end{lemma}

\begin{proof}
略去对 $(\widetilde{M}_n|_U)$ 属于
$\textit{Coh}(U, I\mathcal{O}_U)$ 的验证。设 $(\mathcal{F}_n)$
为 $\textit{Coh}(U, I\mathcal{O}_U)$ 的对象。
由局部上同调，引理
\ref{local-cohomology-lemma-finiteness-pushforwards-and-H1-local}
可知 $\mathcal{F}_n = \widetilde{M_n}$，其中 $A/I^n$-模 $M_n$ 有限。
将 $M_n$ 除以 $H^0_\mathfrak a(M_n)$ 后，
可以假设 $M_n \subset H^0(U, \mathcal{F}_n)$；参见对偶化复形，引理
\ref{dualizing-lemma-divide-by-torsion} 以及已经引用的引理。
归纳地把 $M_{n + 1}$ 替换为 $M_n$ 在映射
$M_{n + 1} \to H^0(U, \mathcal{F}_{n + 1})
\to H^0(U, \mathcal{F}_n)$ 下的逆像之后，可以假设 $M_{n + 1}$ 映入
$M_n$。这给出满足 (1) 和 (2) 的逆系统 $(M_n)$，使得
$\mathcal{F}_n = \widetilde{M_n}$。为说明 (3) 成立，使用如下事实：
$M_{n + 1}/I^nM_{n + 1} \to M_n$ 是有限 $A$-模之间的映射，
在应用 $\widetilde{\ }$ 并限制到 $U$ 后诱导同构
（这里再次使用第一个所引引理）。
\end{proof}

\noindent
在情形 \ref{algebraization-situation-algebraize} 中，可以研究概形上同调，
等式 (\ref{coherent-equation-completion-functor}) 的完备化函子
\begin{equation}
\label{algebraization-equation-completion}
\textit{Coh}(\mathcal{O}_U)
\longrightarrow
\textit{Coh}(U, I\mathcal{O}_U),\quad
\mathcal{F} \longmapsto \mathcal{F}^\wedge
\end{equation}
若 $A$ 是 $I$-进完备的，则由前面对形式截面代数化的研究，
这个函子在适当的子范畴上全忠实；一些示例性结论见第
\ref{algebraization-section-completion} 节和引理 \ref{algebraization-lemma-fully-faithful-inequalities}。
接着，设 $(\mathcal{F}_n)$ 为 $\textit{Coh}(U, I\mathcal{O}_U)$ 的对象。
仍假设 $A$ 是 $I$-进完备的，可以问：$(\mathcal{F}_n)$ 何时位于
上面所示完备化函子的本质像中？

\begin{lemma}
\label{algebraization-lemma-essential-image-completion}
在情形 \ref{algebraization-situation-algebraize} 中，设 $(\mathcal{F}_n)$
为 $\textit{Coh}(U, I\mathcal{O}_U)$ 的对象。考虑下列条件：
\begin{enumerate}
\item $(\mathcal{F}_n)$ 位于函子 (\ref{algebraization-equation-completion}) 的本质像中；
\item $(\mathcal{F}_n)$ 是某个凝聚 $\mathcal{O}_U$-模的完备化；
\item $(\mathcal{F}_n)$ 是某个凝聚 $\mathcal{O}_V$-模的完备化，
其中 $U \cap Y \subset V \subset U$ 为开子概形；
\item $(\mathcal{F}_n)$ 是某个模在 $U$ 上的限制的完备化，而该模是凝聚
$\mathcal{O}_X$-模；
\item $(\mathcal{F}_n)$ 是某个完备化在 $U$ 上的限制，而被完备化者是凝聚
$\mathcal{O}_X$-模；
\item 存在对象 $(\mathcal{G}_n)$，它属于 $\textit{Coh}(X, I\mathcal{O}_X)$，
其在 $U$ 上的限制为 $(\mathcal{F}_n)$。
\end{enumerate}
则条件 (1)、(2)、(3)、(4) 和 (5) 等价，并且蕴含 (6)。
若 $A$ 是 $I$-进完备的，则条件 (6) 蕴含其余条件。
\end{lemma}

\begin{proof}
条件 (1) 和 (2) 等价，因为凝聚 $\mathcal{O}_U$-模 $\mathcal{F}$ 的完备化，
按定义就是 $\mathcal{F}$ 在函子 (\ref{algebraization-equation-completion}) 下的像。
若 $V \subset U$ 是包含 $U \cap Y$ 的开子概形，则有
$$
\textit{Coh}(V, I\mathcal{O}_V) =
\textit{Coh}(U, I\mathcal{O}_U)
$$
，因为凝聚 $\mathcal{O}_V$-模的范畴（这些模支撑在 $V \cap Y$ 上），
与凝聚 $\mathcal{O}_U$-模的范畴（这些模支撑在 $U \cap Y$ 上）相同。
因此，凝聚 $\mathcal{O}_V$-模的完备化是
$\textit{Coh}(U, I\mathcal{O}_U)$ 的对象。由此，(2)、(3)、(4) 和 (5)
等价，因为函子
$\textit{Coh}(\mathcal{O}_X) \to \textit{Coh}(\mathcal{O}_U) \to
\textit{Coh}(\mathcal{O}_V)$ 是本质满的。
参见性质，引理 \ref{properties-lemma-lift-finite-presentation}。

\medskip\noindent
(5) 总是蕴含 (6)。假设 $A$ 是 $I$-进完备的。
于是由概形上同调，引理 \ref{coherent-lemma-inverse-systems-affine}，
$\textit{Coh}(X, I\mathcal{O}_X)$ 的任意对象都对应一个有限 $A$-模。
所以在此情形下，(6) 蕴含 (5)。
\end{proof}

\begin{example}
\label{algebraization-example-not-algebraizable}
设 $k$ 为域。令 $A = k[x, y][[t]]$，其中 $I = (t)$ 且
$\mathfrak a = (x, y, t)$。使用情形 \ref{algebraization-situation-algebraize} 中的记号。
注意
$U \cap Y = (D(x) \cap Y) \cup (D(y) \cap Y)$ 是仿射开覆盖。
对 $n \geq 1$，考虑可逆模 $\mathcal{L}_n$；它是
$\mathcal{O}_U/t^n\mathcal{O}_U$-模，由 $A_x/t^nA_x$ 与 $A_y/t^nA_y$ 粘合而成，
粘合使用 $A_{xy}/t^nA_{xy}$ 中的可逆元，该元是任意如下形式幂级数的像：
$$
u = 1 + \frac{t}{xy} + \sum_{n \geq 2} a_n \frac{t^n}{(xy)^{\varphi(n)}}
$$
其中 $a_n \in k[x, y]$ 且 $\varphi(n) \in \mathbf{N}$。
于是 $(\mathcal{L}_n)$ 是 $\textit{Coh}(U, I\mathcal{O}_U)$ 的可逆对象，
但它不是某个凝聚 $\mathcal{O}_U$-模 $\mathcal{L}$ 的完备化。
这里只概述论证，并略去大部分细节。
设 $y \in U \cap Y$。那么茎 $\mathcal{L}_y$ 的完备化将是可逆模，
因而 $\mathcal{L}_y$ 可逆。于是将存在开子概形
$V \subset U$，它包含 $U \cap Y$，并且 $\mathcal{L}|_V$ 可逆。
由除子，引理 \ref{divisors-lemma-extend-invertible-module}，
可得一个可逆 $A$-模 $M$，满足
$\widetilde{M}|_V \cong \mathcal{L}|_V$。然而环 $A$ 是唯一分解整环，
所以 $M \cong A$，这将推出
$\mathcal{L}_n \cong \mathcal{O}_U/I^n\mathcal{O}_U$。
但由构造有 $\mathcal{L}_2 \not \cong \mathcal{O}_U/I^2\mathcal{O}_U$，
于是得到所需矛盾。

\medskip\noindent
注意，若取 $a_n = 0$，其中 $n \geq 2$，则可见
$\lim H^0(U, \mathcal{L}_n)$ 非零：在这种情形下，函数 $x$
（在 $D(x)$ 上）与函数 $x + t/y$（在 $D(y)$ 上）可以粘合。
另一方面，若取 $a_n = 1$ 且
$\varphi(n) = 2^n$，甚至取 $\varphi(n) = n^2$，则读者可以证明
$\lim H^0(U, \mathcal{L}_n)$ 为零；这给出了在此情形下
$(\mathcal{L}_n)$ 不可代数化的另一证明。
\end{example}

\noindent
若在情形 \ref{algebraization-situation-algebraize} 中环 $A$ 不是
$I$-进完备的，则引理 \ref{algebraization-lemma-essential-image-completion} 表明，
正确的问题是 $(\mathcal{F}_n)$ 是否位于限制函子的本质像中：
$$
\textit{Coh}(X, I\mathcal{O}_X)
\longrightarrow
\textit{Coh}(U, I\mathcal{O}_U)
$$
然而，这时不能再说这意味着 $(\mathcal{F}_n)$ 可代数化。
因此引入下列术语。

\begin{definition}
\label{algebraization-definition-algebraizable}
在情形 \ref{algebraization-situation-algebraize} 中，设 $(\mathcal{F}_n)$ 为
$\textit{Coh}(U, I\mathcal{O}_U)$ 的对象。称
{\it $(\mathcal{F}_n)$ 可延拓到 $X$}，如果存在对象
$(\mathcal{G}_n)$ 属于 $\textit{Coh}(X, I\mathcal{O}_X)$，且其在
$U$ 上的限制同构于 $(\mathcal{F}_n)$。
\end{definition}

\noindent
这个概念等价于在完备化上可代数化。

\begin{lemma}
\label{algebraization-lemma-algebraizable}
在情形 \ref{algebraization-situation-algebraize} 中，设 $(\mathcal{F}_n)$ 为
$\textit{Coh}(U, I\mathcal{O}_U)$ 的对象。设 $A', I', \mathfrak a'$
为 $I$-进完备化，所完备化的对象分别是 $A, I, \mathfrak a$。令 $X' = \Spec(A')$，
且令 $U' = X' \setminus V(\mathfrak a')$。下列条件等价：
\begin{enumerate}
\item $(\mathcal{F}_n)$ 可延拓到 $X$；
\item $(\mathcal{F}_n)$ 到 $U'$ 的拉回是某个凝聚
$\mathcal{O}_{U'}$-模的完备化。
\end{enumerate}
\end{lemma}

\begin{proof}
回忆 $A \to A'$ 是平坦环映射，并诱导同构
$A/I \to A'/I'$。参见代数，引理 \ref{algebra-lemma-completion-flat} 和
\ref{algebra-lemma-completion-complete}。
因此 $X' \to X$ 是平坦态射，并诱导同构 $Y' \to Y$。
从而 $U' \to U$ 是平坦态射，并诱导同构
$U' \cap Y' \to U \cap Y$。这说明在交换图
$$
\xymatrix{
\textit{Coh}(X', I\mathcal{O}_{X'}) \ar[r] &
\textit{Coh}(U', I\mathcal{O}_{U'}) \\
\textit{Coh}(X, I\mathcal{O}_X) \ar[u] \ar[r] &
\textit{Coh}(U, I\mathcal{O}_U) \ar[u]
}
$$
中，竖直函子是等价。参见概形上同调，引理
\ref{coherent-lemma-inverse-systems-pullback-equivalence}。
该引理由此以及引理 \ref{algebraization-lemma-essential-image-completion} 的结论形式地推出。
\end{proof}

\noindent
在情形 \ref{algebraization-situation-algebraize} 中，设 $(\mathcal{F}_n)$ 为
$\textit{Coh}(U, I\mathcal{O}_U)$ 的对象。为了判断
$(\mathcal{F}_n)$ 是否可延拓到 $X$，考察下列 $A$-模是自然的：
\begin{equation}
\label{algebraization-equation-guess}
M = \lim H^0(U, \mathcal{F}_n)
\end{equation}
注意 $M$ 有极限拓扑；它（先验地）比 $I$-进拓扑更粗，因为
$M \to H^0(U, \mathcal{F}_n)$ 零化 $I^nM$。有典范映射
$$
\widetilde{M}|_U \to \widetilde{M/I^nM}|_U \to
\widetilde{H^0(U, \mathcal{F}_n)}|_U \to
\mathcal{F}_n
$$
可以希望 $\widetilde{M}$ 在 $U$ 上限制为凝聚模，且
$(\mathcal{F}_n)$ 是该模的完备化。如果 $A$ 不完备，
这几乎不可能成立，因此这种想法过于天真。但即使 $A$ 是 $I$-进完备的，
这个概念也很难使用。一种不那么天真的做法是考虑下列要求。

\begin{definition}
\label{algebraization-definition-canonically-algebraizable}
在情形 \ref{algebraization-situation-algebraize} 中，设 $(\mathcal{F}_n)$ 为
$\textit{Coh}(U, I\mathcal{O}_U)$ 的对象。称
{\it $(\mathcal{F}_n)$ 可典范地延拓到 $X$}，如果逆系统
$$
\{\widetilde{H^0(U, \mathcal{F}_n)}\}_{n \geq 1}
$$
在 $\QCoh(\mathcal{O}_X)$ 中 pro-同构于某个对象
$(\mathcal{G}_n)$，而该对象属于 $\textit{Coh}(X, I\mathcal{O}_X)$。
\end{definition}

\noindent
引理 \ref{algebraization-lemma-canonically-algebraizable} 将说明，
定义 \ref{algebraization-definition-canonically-algebraizable} 中的条件强于
定义 \ref{algebraization-definition-algebraizable} 中的条件。

\begin{lemma}
\label{algebraization-lemma-canonically-algebraizable}
在情形 \ref{algebraization-situation-algebraize} 中，设 $(\mathcal{F}_n)$ 为
$\textit{Coh}(U, I\mathcal{O}_U)$ 的对象。若 $(\mathcal{F}_n)$
可典范地延拓到 $X$，则
\begin{enumerate}
\item $(\widetilde{H^0(U, \mathcal{F}_n)})$ pro-同构于
一个对象 $(\mathcal{G}_n)$，该对象属于 $\textit{Coh}(X, I \mathcal{O}_X)$，
该对象在唯一同构意义下唯一；
\item $(\mathcal{G}_n)$ 在 $U$ 上的限制同构于
$(\mathcal{F}_n)$；也就是说，$(\mathcal{F}_n)$ 可延拓到 $X$；
\item 逆系统 $\{H^0(U, \mathcal{F}_n)\}$ 满足 Mittag--Leffler 条件；
\item (\ref{algebraization-equation-guess}) 中的模 $M$ 在 $I$-进完备化上有限，
这里所完备化的环是 $A$，
且 $M$ 上的极限拓扑是 $I$-进拓扑。
\end{enumerate}
\end{lemma}

\begin{proof}
由定义 \ref{algebraization-definition-canonically-algebraizable}，(1) 中的
$(\mathcal{G}_n)$ 存在。由引理
\ref{algebraization-lemma-recognize-formal-coherent-modules}，(1) 中的
$(\mathcal{G}_n)$ 唯一。写成 $\mathcal{G}_n = \widetilde{M_n}$。
则 $\{M_n\}$ 是有限 $A$-模的逆系统，且
$M_n = M_{n + 1}/I^n M_{n + 1}$。
由定义 \ref{algebraization-definition-canonically-algebraizable}，逆系统
$\{H^0(U, \mathcal{F}_n)\}$ pro-同构于 $\{M_n\}$。
因此，逆系统 $\{H^0(U, \mathcal{F}_n)\}$ 满足 Mittag--Leffler 条件，且
$M = \lim M_n$（作为拓扑模）。所以，(4) 中 $M$ 的性质由代数，引理
\ref{algebra-lemma-limit-complete}、
\ref{algebra-lemma-finite-over-complete-ring}、
\ref{algebra-lemma-hathat-finitely-generated} 得到。
由于 $U$ 是拟仿射的，典范映射
$$
\widetilde{H^0(U, \mathcal{F}_n)}|_U \to \mathcal{F}_n
$$
是同构（性质，引理
\ref{properties-lemma-quasi-coherent-quasi-affine}）。
由此可知 $(\mathcal{G}_n|_U)$ 和 $(\mathcal{F}_n)$ pro-同构，
因而由引理 \ref{algebraization-lemma-recognize-formal-coherent-modules} 可知它们同构。
\end{proof}

\begin{lemma}
\label{algebraization-lemma-canonically-extend-base-change}
在情形 \ref{algebraization-situation-algebraize} 中，设 $(\mathcal{F}_n)$ 为
$\textit{Coh}(U, I\mathcal{O}_U)$ 的对象。设 $A \to A'$ 为平坦环映射。
令 $X' = \Spec(A')$，设 $U' \subset X'$ 为 $U$ 的逆像，
并记诱导态射为 $g : U' \to U$。令
$(\mathcal{F}'_n) = (g^*\mathcal{F}_n)$；参见概形上同调，引理
\ref{coherent-lemma-inverse-systems-pullback}。
若 $(\mathcal{F}_n)$ 可典范地延拓到 $X$，则
$(\mathcal{F}'_n)$ 可典范地延拓到 $X'$。
此外，引理 \ref{algebraization-lemma-canonically-algebraizable} 为
$(\mathcal{F}_n)$ 给出的延拓拉回为 $(\mathcal{F}'_n)$ 的延拓。
\end{lemma}

\begin{proof}
设 $f : X' \to X$ 为诱导态射。由平坦基变换，有
$H^0(U', \mathcal{F}'_n) = H^0(U, \mathcal{F}_n) \otimes_A A'$；
参见概形上同调，引理
\ref{coherent-lemma-flat-base-change-cohomology}。
因此，若 $(\mathcal{G}_n)$ 属于 $\textit{Coh}(X, I\mathcal{O}_X)$，并且
pro-同构于 $(\widetilde{H^0(U, \mathcal{F}_n)})$，则
$(f^*\mathcal{G}_n)$ pro-同构于
$$
(f^*\widetilde{H^0(U, \mathcal{F}_n)}) =
(\widetilde{H^0(U, \mathcal{F}_n) \otimes_A A'}) =
(\widetilde{H^0(U', \mathcal{F}'_n)})
$$
。证明完成。
\end{proof}

\begin{lemma}
\label{algebraization-lemma-when-done}
在情形 \ref{algebraization-situation-algebraize} 中，设 $(\mathcal{F}_n)$ 为
$\textit{Coh}(U, I\mathcal{O}_U)$ 的对象。设 $M$ 如 (\ref{algebraization-equation-guess})。
假设
\begin{enumerate}
\item[(a)] 逆系统 $H^0(U, \mathcal{F}_n)$ 具有 Mittag--Leffler 性质；
\item[(b)] $M$ 上的极限拓扑与 $I$-进拓扑一致；
\item[(c)] $M \to H^0(U, \mathcal{F}_n)$ 的像是有限 $A$-模，
对所有 $n$ 均如此。
\end{enumerate}
则 $(\mathcal{F}_n)$ 可典范地延拓到 $X$。
特别地，若 $A$ 是 $I$-进完备的，则
$(\mathcal{F}_n)$ 是某个凝聚 $\mathcal{O}_U$-模的完备化。
\end{lemma}

\begin{proof}
由于 $H^0(U, \mathcal{F}_n)$ 具有 Mittag--Leffler 性质，
且 $M$ 上的极限拓扑是 $I$-进拓扑，可知
$\{M/I^nM\}$ 与 $\{H^0(U, \mathcal{F}_n)\}$
是 pro-同构的 $A$-模逆系统。因此，若令
$$
\mathcal{G}_n = \widetilde{M/I^n M}
$$
则为了验证定义 \ref{algebraization-definition-canonically-algebraizable} 中的条件，
只需证明 $M$ 在 $I$-进完备化上为有限模，其中被完备化的环是 $A$。
这由上述结论、条件 (c) 给出的 $M/I^n M$ 的有限性，以及
代数，引理 \ref{algebra-lemma-finite-over-complete-ring} 得到。
\end{proof}

\noindent
下述结论在某种意义上是上面引理 \ref{algebraization-lemma-when-done}
最直接的可能应用。

\begin{lemma}
\label{algebraization-lemma-algebraization-principal-variant}
在情形 \ref{algebraization-situation-algebraize} 中，设
$(\mathcal{F}_n)$ 为 $\textit{Coh}(U, I\mathcal{O}_U)$ 的对象。假设
\begin{enumerate}
\item $I = (f)$ 是由非零因子 $f \in \mathfrak a$ 生成的主理想；
\item $\mathcal{F}_n$ 是有限局部自由
$\mathcal{O}_U/f^n\mathcal{O}_U$-模；
\item $H^1_\mathfrak a(A/fA)$ 和 $H^2_\mathfrak a(A/fA)$
是有限 $A$-模。
\end{enumerate}
则 $(\mathcal{F}_n)$ 可典范地延拓到 $X$。特别地，若 $A$
完备，则 $(\mathcal{F}_n)$ 是某个凝聚 $\mathcal{O}_U$-模的完备化。
\end{lemma}

\begin{proof}
我们将通过验证引理 \ref{algebraization-lemma-when-done} 的假设 (a)、(b) 和 (c)
来证明此结论。

\medskip\noindent
由于 $\mathcal{F}_n$ 在 $\mathcal{O}_U/f^n\mathcal{O}_U$ 上局部自由，
所以有短正合序列
$0 \to \mathcal{F}_n \to \mathcal{F}_{n + 1} \to \mathcal{F}_1 \to 0$
，对所有 $n$ 均如此。因此，由上同调，引理
\ref{cohomology-lemma-topology-I-adic-f}，条件 (b) 成立。

\medskip\noindent
由于 $f$ 是非零因子，得到短正合序列
$$
0 \to A/f^nA \xrightarrow{f} A/f^{n + 1}A \to A/fA \to 0
$$
以及相应的短正合序列
$0 \to \mathcal{F}_n \to \mathcal{F}_{n + 1} \to \mathcal{F}_1 \to 0$.
下文将不再另行说明地使用局部上同调，引理
\ref{local-cohomology-lemma-finiteness-pushforwards-and-H1-local}
。我们的假设蕴含
$H^0(U, \mathcal{O}_U/f\mathcal{O}_U)$ 和
$H^1(U, \mathcal{O}_U/f\mathcal{O}_U)$
都是有限 $A$-模。因此，对 $\mathcal{F}_1$ 也有相同结论；参见局部上同调，引理
\ref{local-cohomology-lemma-finiteness-for-finite-locally-free}。
利用归纳法和这些短正合序列，可知
$H^0(U, \mathcal{F}_n)$ 是有限 $A$-模，对所有 $n$ 均如此。
由此可见假设 (c) 成立。

\medskip\noindent
最后，由于 $H^1(U, \mathcal{F}_1)$ 是有限 $A$-模，
可以应用上同调，引理 \ref{cohomology-lemma-ML}，
从而看出假设 (a) 成立。
\end{proof}

\begin{remark}
\label{algebraization-remark-interesting-case-variant}
在引理 \ref{algebraization-lemma-algebraization-principal-variant} 中，
若 $A$ 是具有 Cohen--Macaulay 形式纤维的泛链状环
（例如，若 $A$ 有对偶化复形），则条件
$H^1_\mathfrak a(A/fA)$ 和 $H^2_\mathfrak a(A/fA)$
是有限 $A$-模，等价于
$$
\text{depth}((A/f)_\mathfrak p) + \dim((A/\mathfrak p)_\mathfrak q) > 2
$$
对所有 $\mathfrak p \in V(f) \setminus V(\mathfrak a)$
和 $\mathfrak q \in V(\mathfrak p) \cap V(\mathfrak a)$ 成立；
这是局部上同调，定理 \ref{local-cohomology-theorem-finiteness} 的结论。

\medskip\noindent
例如，若 $A/fA$ 满足 $(S_2)$，且 $Z = V(\mathfrak a)$ 的每个不可约分支
的余维均 $\geq 3$（余维在 $Y = \Spec(A/fA)$ 中取），则可得
$H^1_\mathfrak a(A/fA)$ 和 $H^2_\mathfrak a(A/fA)$ 的有限性。
应将此与引理 \ref{algebraization-lemma-algebraization-principal} 中稍弱的条件作比较
（另见注记 \ref{algebraization-remark-interesting-case}）。
\end{remark}






\section{凝聚形式模的代数化，II}
\label{algebraization-section-algebraization-modules-general}

\noindent
我们继续第 \ref{algebraization-section-algebraization-modules} 节中开始的讨论。
初次阅读时可以跳过本节。

\begin{lemma}
\label{algebraization-lemma-map-kernel-cokernel-on-closed}
在情形 \ref{algebraization-situation-algebraize} 中，设
$(\mathcal{F}_n) \to (\mathcal{F}'_n)$ 是
$\textit{Coh}(U, I\mathcal{O}_U)$ 中的一个态射，
其核与余核均被 $I$ 的某个幂零化。则
\begin{enumerate}
\item $(\mathcal{F}_n)$ 可延拓到 $X$，当且仅当
$(\mathcal{F}'_n)$ 可延拓到 $X$；
\item $(\mathcal{F}_n)$ 是某个凝聚 $\mathcal{O}_U$-模的完备化，
当且仅当 $(\mathcal{F}'_n)$ 也是如此。
\end{enumerate}
\end{lemma}

\begin{proof}
第 (2) 项直接由概形上同调，引理
\ref{coherent-lemma-existence-easy} 推出。
为证明第 (1) 项，先用引理 \ref{algebraization-lemma-algebraizable}
约化到 $A$ 为 $I$-进完备的情形。
但在这种情形下，引理 \ref{algebraization-lemma-essential-image-completion}
把第 (1) 项约化为第 (2) 项。
\end{proof}

\noindent
下面两个引理原本用于引理 \ref{algebraization-lemma-when-done} 的证明。
我们将其保留在这里，供有兴趣了解可得到哪些中间结果的读者参考。

\begin{lemma}
\label{algebraization-lemma-when-ML}
在情形 \ref{algebraization-situation-algebraize} 中，设 $(\mathcal{F}_n)$ 是
$\textit{Coh}(U, I\mathcal{O}_U)$ 的一个对象。若逆系统
$H^0(U, \mathcal{F}_n)$ 满足 Mittag--Leffler 条件，则典范映射
$$
\widetilde{M/I^nM}|_U \to \mathcal{F}_n
$$
对所有 $n$ 均为满射，其中 $M$ 如 (\ref{algebraization-equation-guess}) 所示。
\end{lemma}

\begin{proof}
满性可在某点 $y \in Y \setminus Z$ 的茎上检验。
若 $y$ 对应于素理想 $\mathfrak q \subset A$，则可以选取
$f \in \mathfrak a$，使得 $f \not \in \mathfrak q$。于是只需证明
$$
M_f \longrightarrow H^0(U, \mathcal{F}_n)_f = H^0(D(f), \mathcal{F}_n)
$$
为满射，因为 $D(f)$ 是仿射的（等号由性质，引理
\ref{properties-lemma-invert-f-sections} 成立）。由于满足
Mittag--Leffler 条件，可知
$$
\Im(M \to H^0(U, \mathcal{F}_n)) =
\Im(H^0(U, \mathcal{F}_m) \to H^0(U, \mathcal{F}_n))
$$
对某个 $m \geq n$ 成立。利用上同调长正合序列可见
$$
\Coker(H^0(U, \mathcal{F}_m) \to H^0(U, \mathcal{F}_n))
\subset
H^1(U, \Ker(\mathcal{F}_m \to \mathcal{F}_n))
$$
由于 $U = X \setminus V(\mathfrak a)$，此 $H^1$ 是
$\mathfrak a$-幂挠的。因此将 $f$ 可逆化后，余核变为零。
\end{proof}

\begin{lemma}
\label{algebraization-lemma-when-topology}
在情形 \ref{algebraization-situation-algebraize} 中，设 $(\mathcal{F}_n)$ 是
$\textit{Coh}(U, I\mathcal{O}_U)$ 的一个对象，并设 $M$
如 (\ref{algebraization-equation-guess}) 所示。令
$$
\mathcal{G}_n = \widetilde{M/I^nM}.
$$
若 $M$ 上的极限拓扑与 $I$-进拓扑一致，则
$\mathcal{G}_n|_U$ 是凝聚
$\mathcal{O}_U$-模，并且逆系统的映射
$$
(\mathcal{G}_n|_U) \longrightarrow (\mathcal{F}_n)
$$
在阿贝尔范畴 $\textit{Coh}(U, I\mathcal{O}_U)$ 中是单射。
\end{lemma}

\begin{proof}
注意，$\mathcal{G}_n$ 是拟凝聚 $\mathcal{O}_X$-模，
并且被 $I^n$ 零化，且
$\mathcal{G}_{n + 1}/I^n\mathcal{G}_{n + 1} = \mathcal{G}_n$.
考虑
$$
M_n = \Im(M \longrightarrow H^0(U, \mathcal{F}_n))
$$
该假设说明逆系统 $(M_n)$ 与
$(M/I^nM)$ 作为 $\text{Mod}_A$ 的 pro-对象同构。
选取 $f \in \mathfrak a$，使得 $D(f) \subset U$ 是仿射开集。于是有
$$
(M_n)_f \subset H^0(U, \mathcal{F}_n)_f = H^0(D(f), \mathcal{F}_n)
$$
由性质，引理 \ref{properties-lemma-invert-f-sections}，这里等号成立。
因此 $\widetilde{M_n}|_U \to \mathcal{F}_n$ 是单射。
由此推出 $\widetilde{M_n}|_U$ 是凝聚模
（概形上同调，引理
\ref{coherent-lemma-coherent-Noetherian-quasi-coherent-sub-quotient}）。
由于 $M \to M/I^nM$ 是满射，并且它分解为
$M_{n'} \to M/I^nM$，其中某个 $n' \geq n$，所以 $\mathcal{G}_n|_U$
作为凝聚模的商是凝聚的。
结合证明开头的观察，可知
$(\mathcal{G}_n|_U)$ 确实构成
$\textit{Coh}(U, I\mathcal{O}_U)$ 的一个对象。
最后，为证明该映射的单性，只需证明
$$
\lim (M/I^nM)_f = \lim H^0(D(f), \mathcal{G}_n) \to
\lim H^0(D(f), \mathcal{F}_n)
$$
为单射；参见概形上同调，引理
\ref{coherent-lemma-inverse-systems-abelian} 和
\ref{coherent-lemma-inverse-systems-affine}。
$\lim (M_n)_f \to \lim H^0(D(f), \mathcal{F}_n)$ 的单性是清楚的
（见上文），而由我们关于 pro-系统的说明，有
$\lim (M_n)_f = \lim (M/I^nM)_f$。证明完成。
\end{proof}





\section{距离函数}
\label{algebraization-section-distance}

\noindent
设 $Y$ 是诺特概形，$Z \subset Y$ 是闭子集。我们定义函数
\begin{equation}
\label{algebraization-equation-delta-Z}
\delta^Y_Z = \delta_Z : Y \longrightarrow \mathbf{Z}_{\geq 0} \cup \{\infty\}
\end{equation}
来度量 $Y$ 的一点到 $Z$ 的“距离”。
非正式讨论见注记 \ref{algebraization-remark-discussion}。
设 $y \in Y$。规定 $\delta_Z(y) = \infty$，如果 $y$ 位于
$Y$ 的一个与 $Z$ 不相交的连通分支中。
若 $y$ 位于 $Y$ 的一个与 $Z$ 相交的连通分支中，
则可找到 $k \geq 0$ 以及系统
$$
V_0 \subset W_0 \supset V_1 \subset W_1 \supset \ldots \supset
V_k \subset W_k
$$
，其中各项都是 $Y$ 的整的闭子概形，$V_0 \subset Z$，
且 $y \in W_k$ 是泛点。令
$c_i = \text{codim}(V_i, W_i)$，其中 $i = 0, \ldots, k$；并令
$b_i = \text{codim}(V_{i + 1}, W_i)$，其中 $i = 0, \ldots, k - 1$。
对这样的系统，令
$$
\delta(V_0, W_0, V_1, \ldots, W_k) =
k +
\max_{i = 0, 1, \ldots, k}
(c_i + c_{i + 1} + \ldots + c_k - b_i - b_{i + 1} - \ldots - b_{k - 1})
$$
这个数 $\geq k$，因为可以取 $i = k$，且 $c_k \geq 0$。
最后令
$$
\delta_Z(y) = \min \delta(V_0, W_0, V_1, \ldots, W_k)
$$
，其中最小值遍历上述所有由 $Y$ 的整的闭子概形组成的系统。

\begin{lemma}
\label{algebraization-lemma-discussion}
设 $Y$ 是诺特概形，$Z \subset Y$ 是闭子集。
\begin{enumerate}
\item 对 $y \in Y$，有 $\delta_Z(y) = 0 \Leftrightarrow y \in Z$。
\item 子集 $\{y \in Y \mid \delta_Z(y) \leq k\}$
对特化稳定。
\item 对 $y \in Y$ 和 $z \in \overline{\{y\}} \cap Z$，有
$\dim(\mathcal{O}_{\overline{\{y\}}, z}) \geq \delta_Z(y)$。
\item 若 $\delta$ 是 $Y$ 上的维数函数，则
$\delta_Z(y) \geq \delta(y) - \delta_{max}$，其中 $\delta_{max}$
是 $\delta$ 在 $Z$ 上的最大值。
\item 若 $Y = \Spec(A)$ 是链状诺特局部环的谱，
其极大理想为 $\mathfrak m$，且 $Z = \{\mathfrak m\}$，则
$\delta_Z(y) = \dim(\overline{\{y\}})$。
\item 若 $Y' \subset Y$ 是开子概形，则
$\delta^{Y'}_{Y' \cap Z}(y') \geq \delta^Y_Z(y')$ 对 $y' \in Y'$ 成立。
\end{enumerate}
假设 $Y$ 是链状的。则
\begin{enumerate}
\setcounter{enumi}{6}
\item 设 $y' \leadsto y$ 是 $Y$ 中点的直接特化。
若 $Y$ 是链状的，则 $\delta_Z(y') \leq \delta_Z(y) + 1$。
\item 给定如下特化图式
$$
\xymatrix{
& y'_0 \ar@{~>}[ld] \ar@{~>}[rd] &
& y'_1 \ar@{~>}[ld] & \ldots
& y'_{k - 1} \ar@{~>}[rd] &
\\
y_0 & &
y_1 & &
\ldots & &
y_k = y
}
$$
，其中各点都在 $Y$ 中，$y_0 \in Z$，并且
$y_i' \leadsto y_i$ 是直接特化，则 $\delta_Z(y_k) \leq k$。
\end{enumerate}
\end{lemma}

\begin{proof}
第 (1) 项的证明。若 $y \in Z$，则可取 $k = 0$ 以及
$V_0 = W_0 = \overline{\{y\}}$，从而得到 $\delta(V_0, W_0) = 0$，故
$\delta_Z(y) = 0$。若 $y \not \in Z$，则对每个系统
$V_0 \subset W_0 \supset V_1 \subset W_1 \supset \ldots \subset W_k$
，只要它是 $y$ 的系统，就或者有 $k = 0$ 且
$V_0 \not = W_0$，或者有 $k > 0$。
在两种情形下都有 $\delta(V_0, W_0, \ldots, W_k) > 0$。因此
$\delta_Z(y) > 0$。

\medskip\noindent
第 (2) 项的证明。设 $y \leadsto y'$ 是非平凡特化，并设
$V_0 \subset W_0 \supset V_1 \subset W_1 \supset \ldots \subset W_k$
是 $y$ 的一个系统。分两种情形。
情形 I：$V_k = W_k$，即 $c_k = 0$。在此情形下可令
$V'_k = W'_k = \overline{\{y'\}}$。
直接计算表明
$\delta(V_0, W_0, \ldots, V'_k, W'_k) \leq
\delta(V_0, W_0, \ldots, V_k, W_k)$
，因为只有 $b_{k - 1}$ 被换成了一个更大的整数。
情形 II：$V_k \not = W_k$，即 $c_k > 0$。
在此情形下，令 $V_{k + 1} = W_{k + 1} = \overline{\{y'\}}$，
可见 $V_0 \subset W_0 \supset \ldots \subset W_k \supset V_{k + 1}
\subset W_{k + 1}$ 是 $y'$ 的一个系统。此时 $c_{k + 1} = 0$，且
$b_k > 0$，所以得到
\begin{align*}
& \delta(V_0, \ldots, W_{k + 1}) \\
& =
k + 1 +
\max_{i = 0, 1, \ldots, k + 1}
(c_i + c_{i + 1} + \ldots + c_k + c_{k + 1}
- b_i - b_{i + 1} - \ldots - b_{k - 1} - b_k) \\
& =
k + 1 +
\max_{i = 0, 1, \ldots, k + 1}
(c_i + c_{i + 1} + \ldots + c_k
- b_i - b_{i + 1} - \ldots - b_{k - 1} - b_k) \\
& \leq
k + \max_{i = 0, 1, \ldots, k}
(c_i + c_{i + 1} + \ldots + c_k - b_i - b_{i + 1} - \ldots - b_{k - 1}) \\
& = \delta(V_0, \ldots, W_k)
\end{align*}
该不等式成立，是因为 $c_k > 0$ 且 $b_k > 0$；这尤其蕴含
$\delta(V_0, \ldots, W_k) \geq k + c_k \geq k + 1$。

\medskip\noindent
第 (3) 项的证明。给定 $y \in Y$ 和 $z \in \overline{\{y\}} \cap Z$，
得到系统
$$
V_0 = \overline{\{z\}} \subset W_0 = \overline{\{y\}}
$$
，并且由性质，引理 \ref{properties-lemma-codimension-local-ring}，有
$c_0 = \text{codim}(V_0, W_0) = \dim(\mathcal{O}_{\overline{\{y\}}, z})$。
因此 $\delta(V_0, W_0) = 0 + c_0 = c_0$，这就证明了所需结论。

\medskip\noindent
第 (4) 项的证明。设 $\delta$ 是 $Y$ 上的维数函数。
设 $V_0 \subset W_0 \supset V_1 \subset W_1 \supset \ldots \subset W_k$
是 $y$ 的一个系统。令 $y'_i \in W_i$ 和 $y_i \in V_i$ 为相应泛点，
于是 $y_0 \in Z$ 且 $y_k = y$。这时
$$
\delta(y_i) - \delta(y_{i - 1}) =
\delta(y'_{i - 1}) - \delta(y_{i - 1}) - \delta(y'_{i - 1}) + \delta(y_i) =
c_{i - 1} - b_{i - 1}
$$
最后有 $\delta(y'_k) - \delta(y_{k - 1}) = c_k$。因此
$$
\delta(y) - \delta(y_0) =
c_0 + \ldots + c_k - b_0 - \ldots - b_{k - 1}
$$
由此得出
$\delta(V_0, W_0, \ldots, W_k) \geq k + \delta(y) - \delta(y_0)$
，这就证明了所需结论。

\medskip\noindent
第 (5) 项的证明。函数 $\delta(y) = \dim(\overline{\{y\}})$
是维数函数。因此由第 (4) 项有 $\delta(y) \leq \delta_Z(y)$。
由第 (3) 项有 $\delta_Z(y) \leq \delta(y)$，证明完成。

\medskip\noindent
第 (6) 项的证明。这是清楚的，因为计算
$\delta^{Y'}_{Y' \cap Z}$ 时需要考虑的系统更少。

\medskip\noindent
第 (7) 项的证明。设
$V_0 \subset W_0 \supset V_1 \subset W_1 \supset \ldots \subset W_k$
是 $y$ 的一个系统。令 $W'_k = \overline{\{y'\}}$。
由于 $Y$ 是链状的，可知
$\text{codim}(V_k, W'_k) = \text{codim}(V_k, W_k) + 1$.
容易推出 $\delta(V_0, \ldots, V_k, W'_k) \leq
\delta(V_0, \ldots, V_k, W_k) + 1$，这就证明了所需结论。

\medskip\noindent
第 (8) 项的证明：结合第 (7) 项与第 (2) 项即可。
\end{proof}

\begin{lemma}
\label{algebraization-lemma-change-distance-function}
设 $Y$ 是泛链状诺特概形，$Z \subset Y$ 是闭子概形。
设 $f : Y' \to Y$ 是有限型态射，其所有纤维的维数均
$\leq e$。令 $Z' = f^{-1}(Z)$。则
$$
\delta_Z(y) \leq \delta_{Z'}(y') + e - \text{trdeg}_{\kappa(y)}(\kappa(y'))
$$
对 $y' \in Y'$ 成立，其中其像为 $y \in Y$。
\end{lemma}

\begin{proof}
若 $\delta_{Z'}(y') = \infty$，则无须证明。
若 $\delta_{Z'}(y') < \infty$，则选取整的闭子概形系统
$$
V'_0 \subset W'_0 \supset V'_1 \subset W'_1 \supset \ldots \subset W'_k
$$
，其中各项在 $Y'$ 中，$V'_0 \subset Z'$，且 $y'$ 是 $W'_k$ 的泛点，
并使得 $\delta_{Z'}(y') = \delta(V'_0, W'_0, \ldots, W'_k)$。
记
$$
V_0 \subset W_0 \supset V_1 \subset W_1 \supset \ldots \subset W_k
$$
为上述概形在 $Y$ 中的概形论像。注意，$y$ 是 $W_k$ 的泛点，
并且 $V_0 \subset Z$。对每个 $i$，考察图表
$$
\xymatrix{
V'_i \ar[r] \ar[d] & W'_i \ar[d] & V'_{i + 1} \ar[l] \ar[d] \\
V_i \ar[r] & W_i & V_{i + 1} \ar[l]
}
$$
记 $n_i$ 为 $V'_i/V_i$ 的相对维数，记
$m_i$ 为 $W'_i/W_i$ 的相对维数；更确切地说，
它们是相应函数域扩张的超越次数。令
$c_i = \text{codim}(V_i, W_i)$,
$c'_i = \text{codim}(V'_i, W'_i)$,
$b_i = \text{codim}(V_{i + 1}, W_i)$，以及
$b'_i = \text{codim}(V'_{i + 1}, W'_i)$。
由维数公式有
$$
c_i = c'_i + n_i - m_i
\quad\text{且}\quad
b_i = b'_i + n_{i + 1} - m_i
$$
参见态射，引理 \ref{morphisms-lemma-dimension-formula}。
因此 $c_i - b_i = c'_i - b'_i + n_i - n_{i + 1}$。于是
\begin{align*}
& c_i + c_{i + 1} + \ldots + c_k - b_i - b_{i + 1} - \ldots - b_{k - 1} \\
& =
c'_i + c'_{i + 1} + \ldots + c'_k - b'_i - b'_{i + 1} - \ldots - b'_{k - 1}
+ n_i - n_k + c_k - c'_k \\
& =
c'_i + c'_{i + 1} + \ldots + c'_k - b'_i - b'_{i + 1} - \ldots - b'_{k - 1}
+ n_i - m_k
\end{align*}
由此可见
\begin{align*}
\max_{i = 0, \ldots, k}
& (c_i + c_{i + 1} + \ldots + c_k - b_i - b_{i + 1} - \ldots - b_{k - 1}) \\
& =
\max_{i = 0, \ldots, k}
(c'_i + c'_{i + 1} + \ldots + c'_k - b'_i - b'_{i + 1} - \ldots - b'_{k - 1}
+ n_i - m_k) \\
& =
\max_{i = 0, \ldots, k}
(c'_i + c'_{i + 1} + \ldots + c'_k - b'_i - b'_{i + 1} - \ldots - b'_{k - 1}
+ n_i) - m_k \\
& \leq
\max_{i = 0, \ldots, k}
(c'_i + c'_{i + 1} + \ldots + c'_k - b'_i - b'_{i + 1} - \ldots - b'_{k - 1})
+ e - m_k
\end{align*}
由于 $m_k = \text{trdeg}_{\kappa(y)}(\kappa(y'))$，可得
$$
\delta(V_0, W_0, \ldots, W_k) \leq
\delta(V'_0, W'_0, \ldots, W'_k) + e - \text{trdeg}_{\kappa(y)}(\kappa(y'))
$$
，即为所需结论。
\end{proof}

\begin{remark}
\label{algebraization-remark-discussion}
设 $Y$ 是链状诺特概形，$Z \subset Y$ 是闭子集。
由引理 \ref{algebraization-lemma-discussion} 有
$$
\delta_Z(y) \leq \min
\left\{ k \middle|
\begin{matrix}
\text{ there exist specializations in }Y \\
y_0 \leftarrow y'_0 \rightarrow y_1 \leftarrow y'_1 \rightarrow \ldots
\leftarrow y'_{k - 1} \rightarrow y_k = y \\
\text{ 其中 }y_0 \in Z\text{ 且 }y_i' \leadsto y_i
\text{ immediate}
\end{matrix}
\right\}
$$
我们断言，若 $Y$ 在某个域上是有限型的，则等号成立。
如果以后需要这一结果，我们会在这里给出精确表述并加以证明。
然而，一般而言，若用该不等式的右端定义 $\delta_Z$，
则我们不知道引理 \ref{algebraization-lemma-change-distance-function} 是否仍然成立。
\end{remark}

\begin{example}
\label{algebraization-example-distance}
设 $k$ 是域，且 $Y = \mathbf{A}^n_k$。记
$\delta : Y \to \mathbf{Z}_{\geq 0}$ 为通常的维数函数。
\begin{enumerate}
\item 若 $Z = \{z\}$，其中 $z$ 是某个闭点，则
\begin{enumerate}
\item $\delta_Z(y) = \delta(y)$ 在 $y \leadsto z$ 时成立；
\item $\delta_Z(y) = \delta(y) + 1$ 在 $y \not \leadsto z$ 时成立。
\end{enumerate}
\item 若 $Z$ 是闭子簇，且 $W = \overline{\{y\}}$，则
\begin{enumerate}
\item $\delta_Z(y) = 0$ 在 $W \subset Z$ 时成立；
\item $\delta_Z(y) = \dim(W) - \dim(Z)$ 在 $Z$ 包含于 $W$ 时成立；
\item $\delta_Z(y) = 1$ 在 $\dim(W) \leq \dim(Z)$ 且
$W \not \subset Z$ 时成立；
\item $\delta_Z(y) = \dim(W) - \dim(Z) + 1$ 在
$\dim(W) > \dim(Z)$ 且 $Z \not \subset W$ 时成立。
\end{enumerate}
\end{enumerate}
情形 (1) 的推广如下：设 $Y$ 在某个域上是有限型的，
且 $Z = \{z\}$ 是闭点。则 $\delta_Z(y) = \delta(y) + t$，
其中 $t$ 是连接 $z$ 与 $\overline{\{y\}}$ 的某个闭点的曲线链的最小长度。
\end{example}







\section{凝聚形式模的代数化，III}
\label{algebraization-section-uniqueness}

\noindent
我们继续第 \ref{algebraization-section-algebraization-modules} 节和
第 \ref{algebraization-section-algebraization-modules-general} 节中开始的讨论。
我们将使用第 \ref{algebraization-section-distance} 节的距离函数，
表述情形 \ref{algebraization-situation-algebraize} 中凝聚形式模所满足的一些自然条件。

\medskip\noindent
在情形 \ref{algebraization-situation-algebraize} 中，给定一点 $y \in U \cap Y$，
可以考虑 $I$-进完备化
$$
\mathcal{O}_{X, y}^\wedge = \lim \mathcal{O}_{X, y}/I^n\mathcal{O}_{X, y}
$$
。这是关于 $I\mathcal{O}_{X, y}^\wedge$ 完备的诺特局部环，
其极大理想为 $\mathfrak m_y^\wedge$；参见代数，
第 \ref{algebra-section-completion-noetherian} 节。
设 $(\mathcal{F}_n)$ 是 $\textit{Coh}(U, I\mathcal{O}_U)$ 的一个对象。
将 $(\mathcal{F}_n)$ 在 $y$ 处的“茎”定义为
$$
\mathcal{F}_y^\wedge = \lim \mathcal{F}_{n, y}
$$
。这是 $\mathcal{O}_{X, y}^\wedge$ 上的有限模。参见代数，引理
\ref{algebra-lemma-limit-complete} 和
\ref{algebra-lemma-finite-over-complete-ring}。

\begin{definition}
\label{algebraization-definition-s-d-inequalities}
在情形 \ref{algebraization-situation-algebraize} 中，设 $(\mathcal{F}_n)$ 是
$\textit{Coh}(U, I\mathcal{O}_U)$ 的一个对象。设 $a, b$ 是整数，
并设 $\delta^Y_Z$ 如 (\ref{algebraization-equation-delta-Z}) 所示。
称 {\it $(\mathcal{F}_n)$ 满足 $(a, b)$-不等式组}，如果对
$y \in U \cap Y$ 以及素理想
$\mathfrak p \subset \mathcal{O}_{X, y}^\wedge$，只要
$\mathfrak p \not \in V(I\mathcal{O}_{X, y}^\wedge)$，就有
\begin{enumerate}
\item 若 $V(\mathfrak p) \cap V(I\mathcal{O}_{X, y}^\wedge) \not =
\{\mathfrak m_y^\wedge\}$，则
$$
\text{depth}((\mathcal{F}^\wedge_y)_\mathfrak p) + \delta^Y_Z(y) \geq a
\quad\text{或}\quad
\text{depth}((\mathcal{F}^\wedge_y)_\mathfrak p) +
\dim(\mathcal{O}_{X, y}^\wedge/\mathfrak p) + \delta^Y_Z(y) > b
$$
\item 若 $V(\mathfrak p) \cap V(I\mathcal{O}_{X, y}^\wedge) =
\{\mathfrak m_y^\wedge\}$，则
$$
\text{depth}((\mathcal{F}^\wedge_y)_\mathfrak p) + \delta^Y_Z(y) > a
$$
\end{enumerate}
称 {\it $(\mathcal{F}_n)$ 满足严格 $(a, b)$-不等式组}，
如果对 $y \in U \cap Y$ 以及素理想
$\mathfrak p \subset \mathcal{O}_{X, y}^\wedge$，只要
$\mathfrak p \not \in V(I\mathcal{O}_{X, y}^\wedge)$，就有
$$
\text{depth}((\mathcal{F}^\wedge_y)_\mathfrak p) + \delta^Y_Z(y) > a
\quad\text{或}\quad
\text{depth}((\mathcal{F}^\wedge_y)_\mathfrak p) +
\dim(\mathcal{O}_{X, y}^\wedge/\mathfrak p) + \delta^Y_Z(y) > b
$$
\end{definition}

\noindent
下面是一些初等观察。

\begin{lemma}
\label{algebraization-lemma-elementary}
在情形 \ref{algebraization-situation-algebraize} 中，设 $(\mathcal{F}_n)$ 是
$\textit{Coh}(U, I\mathcal{O}_U)$ 的一个对象，并设 $a, b$ 是整数。
\begin{enumerate}
\item 若 $(\mathcal{F}_n)$ 被 $I$ 的某个幂零化，则
$(\mathcal{F}_n)$ 满足 $(a, b)$-不等式组，对任意 $a, b$ 都如此。
\item 若 $(\mathcal{F}_n)$ 满足 $(a + 1, b)$-不等式组，则
$(\mathcal{F}_n)$ 满足严格 $(a, b)$-不等式组。
\end{enumerate}
若 $\text{cd}(A, I) \leq d$，且 $A$ 有对偶化复形，则
\begin{enumerate}
\item[(3)] $(\mathcal{F}_n)$ 满足 $(s, s + d)$-不等式组，
当且仅当对所有 $y \in U \cap Y$，元组
$\mathcal{O}_{X, y}^\wedge, I\mathcal{O}_{X, y}^\wedge,
\{\mathfrak m_y^\wedge\}, \mathcal{F}_y^\wedge, s - \delta^Y_Z(y), d$
如情形 \ref{algebraization-situation-bootstrap} 所示。
\item[(4)]
若 $(\mathcal{F}_n)$ 满足严格 $(s, s + d)$-不等式组，则
$(\mathcal{F}_n)$ 满足 $(s, s + d)$-不等式组。
\end{enumerate}
\end{lemma}

\begin{proof}
除第 (4) 项外均为立即的；第 (4) 项由引理
\ref{algebraization-lemma-bootstrap-bis-bis} 以及第 (3) 项中的对应关系推出。
\end{proof}

\begin{lemma}
\label{algebraization-lemma-explain-2-3-cd-1}
在情形 \ref{algebraization-situation-algebraize} 中，设 $(\mathcal{F}_n)$ 是
$\textit{Coh}(U, I\mathcal{O}_U)$ 的一个对象。若
$\text{cd}(A, I) = 1$，则 $\mathcal{F}$ 满足 $(2, 3)$-不等式组，
当且仅当
$$
\text{depth}((\mathcal{F}^\wedge_y)_\mathfrak p) +
\dim(\mathcal{O}_{X, y}^\wedge/\mathfrak p) + \delta^Y_Z(y) > 3
$$
对所有 $y \in U \cap Y$ 和
$\mathfrak p \subset \mathcal{O}_{X, y}^\wedge$，只要
$\mathfrak p \not \in V(I\mathcal{O}_{X, y}^\wedge)$ 就成立。
\end{lemma}

\begin{proof}
注意，对素理想 $\mathfrak p \subset \mathcal{O}_{X, y}^\wedge$，
若 $\mathfrak p \not \in V(I\mathcal{O}_{X, y}^\wedge)$，则有
$V(\mathfrak p) \cap V(I\mathcal{O}_{X, y}^\wedge) =
\{\mathfrak m_y^\wedge\}
\Leftrightarrow \dim(\mathcal{O}_{X, y}^\wedge/\mathfrak p) = 1$
，因为 $\text{cd}(A, I) = 1$。参见局部上同调，引理
\ref{local-cohomology-lemma-cd-change-rings} 和
\ref{local-cohomology-lemma-cd-bound-dim-local}。
现在考虑三个数
$\alpha = \text{depth}((\mathcal{F}^\wedge_y)_\mathfrak p) \geq 0$、
$\beta = \dim(\mathcal{O}_{X, y}^\wedge/\mathfrak p) \geq 1$ 和
$\gamma = \delta^Y_Z(y) \geq 1$。
由定义 \ref{algebraization-definition-s-d-inequalities}，所需条件为
\begin{enumerate}
\item 若 $\beta > 1$，则
$\alpha + \gamma \geq 2$ 或 $\alpha + \beta + \gamma > 3$；
\item 若 $\beta = 1$，则 $\alpha + \gamma > 2$。
\end{enumerate}
显然，这等价于 $\alpha + \beta + \gamma > 3$。
\end{proof}

\noindent
在本节余下部分中（建议读者初次阅读时跳过），我们将证明：
当 $A$ 是 $I$-进完备的时，由 $(\mathcal{F}_n)$ 构成如下范畴：
它们属于 $\textit{Coh}(U, I\mathcal{O}_U)$、可延拓到 $X$，并满足严格
$(1, 1 + \text{cd}(A, I))$-不等式组。该范畴等价于凝聚
$\mathcal{O}_U$-模范畴的一个满子范畴。

\begin{lemma}
\label{algebraization-lemma-sanity}
在情形 \ref{algebraization-situation-algebraize} 中，设 $\mathcal{F}$ 是凝聚
$\mathcal{O}_U$-模，且 $d \geq 1$。假设
\begin{enumerate}
\item $A$ 是 $I$-进完备的，有对偶化复形，并且
$\text{cd}(A, I) \leq d$；
\item 完备化 $\mathcal{F}^\wedge$（它来自 $\mathcal{F}$）
满足严格 $(1, 1 + d)$-不等式组。
\end{enumerate}
设 $x \in X$ 是一点，并令 $W = \overline{\{x\}}$。
若 $W \cap Y$ 有一个包含于 $Z$ 的不可约分支，
并且还有一个不包含于其中的不可约分支，则
$\text{depth}(\mathcal{F}_x) \geq 1$。
\end{lemma}

\begin{proof}
设 $W \cap Y = W_1 \cup \ldots \cup W_n$ 是不可约分支分解。
由假设，重新编号后可找到 $0 < m < n$，使得
$W_1, \ldots, W_m  \subset Z$ 且
$W_{m + 1}, \ldots, W_n \not \subset Z$。于是
$$
W \cap Y \setminus
\left((W_1 \cup \ldots \cup W_m) \cap (W_{m + 1} \cup \ldots \cup W_n)\right)
$$
不连通。由引理 \ref{algebraization-lemma-connected}，可找到
$1 \leq i \leq m < j \leq n$ 以及
$z \in W_i \cap W_j$，使得 $\dim(\mathcal{O}_{W, z}) \leq d + 1$。
选取直接特化 $y \leadsto z$，其中
$y \in W_j$ 且 $y \not \in Z$；$y$ 的存在性由性质，引理
\ref{properties-lemma-complement-closed-point-Jacobson} 得到。
注意，$\delta^Y_Z(y) = 1$ 且 $\dim(\mathcal{O}_{W, y}) \leq d$。
设 $\mathfrak p \subset \mathcal{O}_{X, y}$ 是对应于 $x$ 的素理想。
设 $\mathfrak p' \subset \mathcal{O}_{X, y}^\wedge$ 是
$\mathfrak p\mathcal{O}_{X, y}^\wedge$ 上方的一个极小素理想。则
$$
\text{depth}(\mathcal{F}_x) =
\text{depth}((\mathcal{F}^\wedge_y)_{\mathfrak p'})
\quad\text{且}\quad
\dim(\mathcal{O}_{W, y}) = \dim(\mathcal{O}_{X, y}^\wedge/\mathfrak p')
$$
参见代数，引理 \ref{algebra-lemma-apply-grothendieck-module}，以及
局部上同调，引理 \ref{local-cohomology-lemma-change-completion}。
现在直接从所给不等式组读出结论。
\end{proof}

\begin{lemma}
\label{algebraization-lemma-recover}
在情形 \ref{algebraization-situation-algebraize} 中，设 $\mathcal{F}$ 是凝聚
$\mathcal{O}_U$-模，且 $d \geq 1$。假设
\begin{enumerate}
\item $A$ 是 $I$-进完备的，有对偶化复形，并且
$\text{cd}(A, I) \leq d$；
\item 完备化 $\mathcal{F}^\wedge$（它来自 $\mathcal{F}$）
满足严格 $(1, 1+ d)$-不等式组；
\item 对 $x \in U$，若 $\overline{\{x\}} \cap Y \subset Z$，
则有 $\text{depth}(\mathcal{F}_x) \geq 2$。
\end{enumerate}
则 $H^0(U, \mathcal{F}) \to \lim H^0(U, \mathcal{F}/I^n\mathcal{F})$
是同构。
\end{lemma}

\begin{proof}
我们将通过证明引理 \ref{algebraization-lemma-application-H0} 可以应用来得到结论。
为此，取 $x \in \text{Ass}(\mathcal{F})$，其中 $x \not \in Y$，并令
$W = \overline{\{x\}}$。
由条件 (3)，有 $W \cap Y \not \subset Z$。
由引理 \ref{algebraization-lemma-sanity}，$W \cap Y$ 的任何不可约分支都不包含于 $Z$。
因此，若 $z \in W \cap Z$，则存在直接特化
$y \leadsto z$，其中 $y \in W \cap Y$ 且 $y \not \in Z$。
$y$ 的存在性使用性质，引理
\ref{properties-lemma-complement-closed-point-Jacobson}。
这时 $\delta^Y_Z(y) = 1$，且假设蕴含
$\dim(\mathcal{O}_{W, y}) > d$。
因此 $\dim(\mathcal{O}_{W, z}) > 1 + d$，结论得证。
\end{proof}

\begin{lemma}
\label{algebraization-lemma-fully-faithful-inequalities}
在情形 \ref{algebraization-situation-algebraize} 中，设 $\mathcal{F}$ 是凝聚
$\mathcal{O}_U$-模，且 $d \geq 1$。假设
\begin{enumerate}
\item $A$ 是 $I$-进完备的，有对偶化复形，并且
$\text{cd}(A, I) \leq d$；
\item 完备化 $\mathcal{F}^\wedge$（它来自 $\mathcal{F}$）
满足严格 $(1, 1 + d)$-不等式组；
\item 对 $x \in U$，若 $\overline{\{x\}} \cap Y \subset Z$，
则有 $\text{depth}(\mathcal{F}_x) \geq 2$。
\end{enumerate}
则映射
$$
\Hom_U(\mathcal{G}, \mathcal{F})
\longrightarrow
\Hom_{\textit{Coh}(U, I\mathcal{O}_U)}(\mathcal{G}^\wedge, \mathcal{F}^\wedge)
$$
对每个凝聚 $\mathcal{O}_U$-模 $\mathcal{G}$ 都是双射。
\end{lemma}

\begin{proof}
令 $\mathcal{H} = \SheafHom_{\mathcal{O}_U}(\mathcal{G}, \mathcal{F})$。
利用概形上同调，引理
\ref{coherent-lemma-hom-into-depth}，或更多代数，引理
\ref{more-algebra-lemma-hom-into-depth}，可知 $\mathcal{H}$ 的完备化
满足严格 $(1, 1 + d)$-不等式组，并且对
$x \in U$，若 $\overline{\{x\}} \cap Y \subset Z$，则有
$\text{depth}(\mathcal{H}_x) \geq 2$。细节从略。
因此由引理 \ref{algebraization-lemma-recover}，有
$$
\Hom_U(\mathcal{G}, \mathcal{F}) =
H^0(U, \mathcal{H}) =
\lim H^0(U, \mathcal{H}/\mathcal{I}^n\mathcal{H}) =
\Mor_{\textit{Coh}(U, I\mathcal{O}_U)}
(\mathcal{G}^\wedge, \mathcal{F}^\wedge)
$$
最后一个等号见概形上同调，引理
\ref{coherent-lemma-completion-internal-hom}。
\end{proof}

\begin{lemma}
\label{algebraization-lemma-construct-unique}
在情形 \ref{algebraization-situation-algebraize} 中，设 $(\mathcal{F}_n)$ 是
$\textit{Coh}(U, I\mathcal{O}_U)$ 的一个对象，且 $d \geq 1$。假设
\begin{enumerate}
\item $A$ 是 $I$-进完备的，有对偶化复形，并且
$\text{cd}(A, I) \leq d$；
\item $(\mathcal{F}_n)$ 是某个凝聚 $\mathcal{O}_U$-模的完备化；
\item $(\mathcal{F}_n)$ 满足严格 $(1, 1 + d)$-不等式组。
\end{enumerate}
则存在唯一的凝聚 $\mathcal{O}_U$-模 $\mathcal{F}$，
其完备化是 $(\mathcal{F}_n)$，并且对
$x \in U$，若 $\overline{\{x\}} \cap Y \subset Z$，则有
$\text{depth}(\mathcal{F}_x) \geq 2$。
\end{lemma}

\begin{proof}
选取一个凝聚 $\mathcal{O}_U$-模 $\mathcal{F}$，其完备化为
$(\mathcal{F}_n)$。令
$T = \{x \in U \mid \overline{\{x\}} \cap Y \subset Z\}$.
我们将构造 $\mathcal{F}$：把局部上同调，引理
\ref{local-cohomology-lemma-make-S2-along-T-simple}
应用于 $\mathcal{F}$ 和 $T$。
唯一性则由引理 \ref{algebraization-lemma-fully-faithful-inequalities} 的映射性质推出。

\medskip\noindent
由于 $T$ 在 $U$ 中对特化稳定，只需检验下述条件。
若 $x' \leadsto x$ 是 $U$ 中点的直接特化，且
$x \in T$、$x' \not \in T$，则
$\text{depth}(\mathcal{F}_{x'}) \geq 1$。
令 $W = \overline{\{x\}}$，并令 $W' = \overline{\{x'\}}$。
由于 $x' \not \in T$，可知 $W' \cap Y$ 不包含于 $Z$。
若 $W' \cap Y$ 含有一个包含于 $Z$ 的不可约分支，
则由引理 \ref{algebraization-lemma-sanity} 得证。
否则，选取一个不可约分支 $W_1$（它属于 $W \cap Y$），以及
一个不可约分支 $W'_1$（它属于 $W' \cap Y$），使得
$W_1 \subset W'_1$。
设 $z \in W_1$ 是泛点。设 $y \leadsto z$、$y \in W'_1$
构成直接特化，且 $y \not \in Z$；$y$ 的存在性由
$W'_1 \not \subset Z$（见上文）以及性质，引理
\ref{properties-lemma-complement-closed-point-Jacobson} 推出。
于是有 $z \in Z$、$x \leadsto z$、
$x' \leadsto y \leadsto z$、$y \in Y \setminus Z$，以及
$\delta^Y_Z(y) = 1$。
由局部上同调，引理 \ref{local-cohomology-lemma-cd-bound-dim-local}
以及 $z$ 是 $W \cap Y$ 的泛点这一事实，有
$\dim(\mathcal{O}_{W, z}) \leq d$。
由于 $x' \leadsto x$ 是直接特化，有
$\dim(\mathcal{O}_{W', z}) \leq d + 1$。
由于 $y \not = z$，可得
$\dim(\mathcal{O}_{W', y}) \leq d$。
若 $\text{depth}(\mathcal{F}_{x'}) = 0$，则与假设 (3) 矛盾；
从 $\mathcal{O}_{X, y}$ 过渡到其完备化的细节从略。
证明完成。
\end{proof}








\section{凝聚形式模的代数化，IV}
\label{algebraization-section-algebraization-modules-local}

\noindent
在本节中，我们证明引理 \ref{algebraization-lemma-algebraization-principal-variant}
在局部情形下的两个更强版本，即引理
\ref{algebraization-lemma-algebraization-principal} 和
\ref{algebraization-lemma-algebraization-principal-bis}。
虽然更一般的命题 \ref{algebraization-proposition-cd-1} 将使这些引理不再需要，
但它们的证明要容易得多。

\begin{lemma}
\label{algebraization-lemma-algebraization-principal}
在情形 \ref{algebraization-situation-algebraize} 中，设 $(\mathcal{F}_n)$ 是
$\textit{Coh}(U, I\mathcal{O}_U)$ 的一个对象。假设
\begin{enumerate}
\item $A$ 是局部环，且 $\mathfrak a = \mathfrak m$ 是极大理想；
\item $A$ 有对偶化复形；
\item $I = (f)$ 是由非零因子 $f \in \mathfrak m$ 生成的主理想；
\item $\mathcal{F}_n$ 是有限局部自由
$\mathcal{O}_U/f^n\mathcal{O}_U$-模；
\item 若 $\mathfrak p \in V(f) \setminus \{\mathfrak m\}$，则
$\text{depth}((A/f)_\mathfrak p) + \dim(A/\mathfrak p) > 1$；
\item 若 $\mathfrak p \not \in V(f)$ 且
$V(\mathfrak p) \cap V(f) \not = \{\mathfrak m\}$，则
$\text{depth}(A_\mathfrak p) + \dim(A/\mathfrak p) > 3$。
\end{enumerate}
则 $(\mathcal{F}_n)$ 可典范地延拓到 $X$。特别地，若 $A$ 完备，
则 $(\mathcal{F}_n)$ 是某个凝聚 $\mathcal{O}_U$-模的完备化。
\end{lemma}

\begin{proof}
我们将通过验证引理 \ref{algebraization-lemma-when-done} 的假设 (a)、(b) 和 (c)
来证明该结论。

\medskip\noindent
由于 $\mathcal{F}_n$ 在 $\mathcal{O}_U/f^n\mathcal{O}_U$ 上局部自由，
所以有短正合序列
$0 \to \mathcal{F}_n \to \mathcal{F}_{n + 1} \to \mathcal{F}_1 \to 0$
，对所有 $n$ 均成立。因此由上同调，引理
\ref{cohomology-lemma-topology-I-adic-f}，条件 (b) 成立。

\medskip\noindent
对 $n$ 归纳，并使用短正合序列
$0 \to A/f^n \to A/f^{n + 1} \to A/f \to 0$，可知
$A/f^nA$ 的相伴素理想与 $A/fA$ 的相伴素理想相同。
由于假设 (3)，$\mathcal{F}_n$ 的相伴点对应于
$A/f^nA$ 的不等于 $\mathfrak m$ 的相伴素理想。因此由 (5) 以及
局部上同调，命题 \ref{local-cohomology-proposition-kollar}，可知
$M_n = H^0(U, \mathcal{F}_n)$ 是有限 $A$-模。
故假设 (c) 成立。

\medskip\noindent
为完成证明，只需证明存在 $n > 1$，使映射
$$
H^1(U, \mathcal{F}_n) \longrightarrow H^1(U, \mathcal{F}_1)
$$
的像作为 $A$-模具有有限长度。事实上，由上同调，引理
\ref{cohomology-lemma-ML-better}，这将蕴含假设 (a)。由引理
\ref{algebraization-lemma-ML-local}，当 $n$ 充分大时，该像与 $n$ 无关。
设 $\omega_A^\bullet$ 是 $A$ 的归一化对偶化复形。
由局部对偶定理和 Matlis 对偶
（对偶化复形，引理 \ref{dualizing-lemma-special-case-local-duality}
以及命题 \ref{dualizing-proposition-matlis}），我们的断言等价于：
映射
$$
\text{Ext}^{-2}_A(M_1, \omega_A^\bullet) \to
\text{Ext}^{-2}_A(M_n, \omega_A^\bullet)
$$
的像对 $n \gg 1$ 具有有限长度。所涉及的模是有限 $A$-模，
其支集位于 $V(f)$。因此，只需证明在素理想 $\mathfrak q$
处局部化后该映射为零，其中它包含 $f$ 且不等于 $\mathfrak m$。
设 $\omega_{A_\mathfrak q}^\bullet$ 是 $A_\mathfrak q$ 上的
归一化对偶化复形，并回忆
$\omega_{A_\mathfrak q}^\bullet =
(\omega_A^\bullet)_\mathfrak q[\dim(A/\mathfrak q)]$；这是由
对偶化复形，引理 \ref{dualizing-lemma-dimension-function} 得到的。
利用 (4) 中给出的 $\mathcal{F}_n$ 的局部结构，只需证明
$$
\text{Ext}^{-2 + \dim(A/\mathfrak q)}_{A_\mathfrak q}(
A_\mathfrak q/f, \omega_{A_\mathfrak q}^\bullet)
\to
\text{Ext}^{-2 + \dim(A/\mathfrak q)}_{A_\mathfrak q}(
A_\mathfrak q/f^n, \omega_{A_\mathfrak q}^\bullet)
$$
在 $n$ 充分大时为零。若 $\dim(A/\mathfrak q) > 3$，则由局部上同调，
引理 \ref{local-cohomology-lemma-sitting-in-degrees} 立即得到。
在其他情形下，使用长正合序列
$$
\ldots
\xrightarrow{f^n}
H^{-1}(\omega_{A_\mathfrak q}^\bullet)
\to
\text{Ext}^{-1}_{A_\mathfrak q}(
A_\mathfrak q/f^n, \omega_{A_\mathfrak q}^\bullet) \to
H^0(\omega_{A_\mathfrak q}^\bullet)
\xrightarrow{f^n}
H^0(\omega_{A_\mathfrak q}^\bullet)
\to
\text{Ext}^0_{A_\mathfrak q}(
A_\mathfrak q/f^n, \omega_{A_\mathfrak q}^\bullet) \to 0
$$
若 $\dim(A/\mathfrak q) = 2$，则
$H^0(\omega_{A_\mathfrak q}^\bullet) = 0$
，因为 $\text{depth}(A_\mathfrak q) \geq 1$，这是由于
$f$ 是非零因子。
因此长正合序列表明，所需条件是映射
$$
f^{n - 1} :
H^{-1}(\omega_{A_\mathfrak q}^\bullet)/f \to
H^{-1}(\omega_{A_\mathfrak q}^\bullet)/f^n
$$
为零。现在 $H^{-1}(\omega^\bullet_\mathfrak q)$ 是有限模，
其支集包含在满足下式的素理想 $\mathfrak p \subset A_\mathfrak q$ 中：
$\text{depth}(A_\mathfrak p) + \dim((A/\mathfrak p)_\mathfrak q) \leq 1$.
由于 $\dim((A/\mathfrak p)_\mathfrak q) = \dim(A/\mathfrak p) - 2$，
条件 (6) 表明这些素理想包含于 $V(f)$。因此，当 $n$ 充分大时，
得到所需的零化。
最后，若 $\dim(A/\mathfrak q) = 1$，则条件 (5) 与 $f$ 为非零因子
共同保证 $A_\mathfrak q$ 的深度至少为 $2$。因此
$H^0(\omega_{A_\mathfrak q}^\bullet) =
H^{-1}(\omega_{A_\mathfrak q}^\bullet) = 0$
，且长正合序列表明，该断言等价于下述映射为零：
$$
f^{n - 1} :
H^{-2}(\omega_{A_\mathfrak q}^\bullet)/f \to
H^{-2}(\omega_{A_\mathfrak q}^\bullet)/f^n
$$
现在 $H^{-2}(\omega^\bullet_\mathfrak q)$ 是有限模，
其支集包含在满足 $\mathfrak p \subset A_\mathfrak q$ 以及
$\text{depth}(A_\mathfrak p) + \dim((A/\mathfrak p)_\mathfrak q)
\leq 2$ 的素理想中。由条件 (6)，所有这些素理想都包含于 $V(f)$。
因此，当 $n$ 充分大时，得到所需的零化。
\end{proof}

\begin{remark}
\label{algebraization-remark-interesting-case}
设 $(A, \mathfrak m)$ 是维数 $\geq 4$ 的完备诺特正规局部整环，
并设 $f \in \mathfrak m$ 非零。此时引理
\ref{algebraization-lemma-algebraization-principal} 的假设 (1)、(2)、(3)、(5) 和 (6)
都成立。因此，$U$ 沿 $U \cap V(f)$ 的形式完备化上的向量丛
可以代数化。在引理 \ref{algebraization-lemma-algebraization-principal-bis} 中，
我们会将此推广到更一般的凝聚形式模；另请比较注记
\ref{algebraization-remark-interesting-case-bis}。
\end{remark}

\begin{lemma}
\label{algebraization-lemma-helper-algebraize}
在情形 \ref{algebraization-situation-algebraize} 中，设 $(M_n)$ 是如引理
\ref{algebraization-lemma-system-of-modules} 所示的 $A$-模逆系统，并设
$(\mathcal{F}_n)$ 是 $\textit{Coh}(U, I\mathcal{O}_U)$ 中相应的对象。
设 $d \geq \text{cd}(A, I)$ 和 $s \geq 0$ 为整数。
沿用上述记号，假设
\begin{enumerate}
\item $A$ 是局部环，其极大理想为 $\mathfrak m = \mathfrak a$；
\item $A$ 有对偶化复形；
\item $(\mathcal{F}_n)$ 满足 $(s, s + d)$-不等式组
（定义 \ref{algebraization-definition-s-d-inequalities}）。
\end{enumerate}
设 $E$ 是 $A$ 的剩余域的内射包。则对 $i \leq s$，
存在有限 $A$-模 $N$，它被 $I$ 的某个幂零化，并且对 $n \gg 0$
存在相容映射
$$
H^i_\mathfrak m(M_n) \to \Hom_A(N, E)
$$
，其余核是有限长度 $A$-模，其核 $K_n$ 构成逆系统，并且
$\Im(K_{n''} \to K_{n'})$ 对 $n'' \gg n' \gg 0$ 具有有限长度。
\end{lemma}

\begin{proof}
设 $\omega_A^\bullet$ 是归一化对偶化复形。则
$\delta^Y_Z = \delta$ 是与该对偶化复形相伴的维数函数。
注意，$\Ext^{-i}_A(M_n, \omega_A^\bullet)$ 是有限 $A$-模，
并被 $I^n$ 零化。固定 $0 \leq i \leq s$。
下面将找到 $n_1 > n_0 > 0$，使得若令
$$
N = \Im(\Ext^{-i}_A(M_{n_0}, \omega_A^\bullet) \to
\Ext^{-i}_A(M_{n_1}, \omega_A^\bullet))
$$
，则映射
$$
N \to \Ext^{-i}_A(M_n, \omega_A^\bullet),\quad n \geq n_1
$$
的核是有限长度 $A$-模，而余核 $Q_n$ 构成一个系统，使得
$\Im(Q_{n'} \to Q_{n''})$ 对 $n'' \gg n' \gg n_1$ 具有有限长度。
这等价于系统
$\{\Ext^{-i}_A(M_n, \omega_A^\bullet)\}_{n \geq 1}$
在有限 $A$-模范畴模去有限长度 $A$-模的 Serre 子范畴所得的商中
本质上为常值。由局部对偶定理
（对偶化复形，引理 \ref{dualizing-lemma-special-case-local-duality}）
以及 Matlis 对偶
（对偶化复形，命题 \ref{dualizing-proposition-matlis}），可知存在映射
$$
H^i_\mathfrak m(M_n) \to \Hom_A(N, E),\quad n \geq n_1
$$
，具有引理陈述中的性质。

\medskip\noindent
选取 $f \in \mathfrak m$。令 $B = A_f^\wedge$ 为 $I$-进完备化，
它取自局部化 $A_f$。回忆
$\omega_{A_f}^\bullet = \omega_A^\bullet \otimes_A A_f$
和 $\omega_B^\bullet = \omega_A^\bullet \otimes_A B$ 都是对偶化复形
（对偶化复形，引理 \ref{dualizing-lemma-dualizing-localize} 和
\ref{dualizing-lemma-completion-henselization-dualizing}）。
令 $M$ 为有限 $B$-模 $\lim M_{n, f}$（比较概形上同调，引理
\ref{coherent-lemma-inverse-systems-affine} 中的讨论）。则
$$
\Ext^{-i}_A(M_n, \omega_A^\bullet)_f =
\Ext^{-i}_{A_f}(M_{n, f}, \omega_{A_f}^\bullet) =
\Ext^{-i}_B(M/I^n M, \omega_B^\bullet)
$$
由于 $\mathfrak m$ 可以由有限多个 $f \in \mathfrak m$ 生成，
只需证明对每个 $f$，系统
$$
\{\Ext^{-i}_B(M/I^n M, \omega_B^\bullet)\}_{n \geq 1}
$$
本质上为常值。若干细节从略。

\medskip\noindent
设 $\mathfrak q \subset IB$ 是素理想。则 $\mathfrak q$ 对应于一点
$y \in U \cap Y$。注意，
$\delta(\mathfrak q) = \dim(\overline{\{y\}})$
也是与 $\omega_B^\bullet$ 相伴的维数函数之值
（细节从略；使用 $\omega_B^\bullet$ 是由
$\omega_A^\bullet$ 与 $B$ 作张量积得到的这一事实）。
假设 (3) 通过引理 \ref{algebraization-lemma-elementary} 保证引理
\ref{algebraization-lemma-algebraize-local-cohomology-bis} 可应用于
$B_\mathfrak q, IB_\mathfrak q, \mathfrak qB_\mathfrak q, M_\mathfrak q$
，其中将 $s$ 换为 $s - \delta(y)$。于是得到
$$
H^{i - \delta(\mathfrak q)}_{\mathfrak qB_\mathfrak q}(M_\mathfrak q) =
\lim H^{i - \delta(\mathfrak q)}_{\mathfrak qB_\mathfrak q}(
(M/I^nM)_\mathfrak q)
$$
，且该模被 $I$ 的某个幂零化。
由引理 \ref{algebraization-lemma-terrific}，逆系统
$H^{i - \delta(\mathfrak q)}_{\mathfrak qB_\mathfrak q}((M/I^nM)_\mathfrak q)$
本质上为常值，其值为
$H^{i - \delta(\mathfrak q)}_{\mathfrak qB_\mathfrak q}(M_\mathfrak q)$.
由于 $(\omega_B^\bullet)_\mathfrak q[-\delta(\mathfrak q)]$ 是
$B_\mathfrak q$ 上的归一化对偶化复形，局部对偶定理表明系统
$$
\Ext^{-i}_B(M/I^n M, \omega_B^\bullet)_\mathfrak q
$$
本质上为常值，其值为
$\Ext^{-i}_B(M, \omega_B^\bullet)_\mathfrak q$.

\medskip\noindent
为完成证明，我们像引理 \ref{algebraization-lemma-bootstrap} 的证明那样作整体化；
这里的论证更容易，因为已经知道该系统的值。具体而言，考虑映射
$$
\alpha_n :
\Ext^{-i}_B(M/I^n M, \omega_B^\bullet)
\longrightarrow
\Ext^{-i}_B(M, \omega_B^\bullet)
$$
，其中 $n$ 变化。由上文，对每个 $\mathfrak q$，可找到一个
$n$，使得 $\alpha_n$ 在 $\mathfrak q$ 处局部化后为满射。
由于 $B$ 是诺特的，且 $\Ext^{-i}_B(M, \omega_B^\bullet)$
是有限模，可找到一个 $n$，使得 $\alpha_n$ 为满射。
对任意 $n$，若 $\alpha_n$ 为满射，则给定素理想
$\mathfrak q \in V(IB)$，可找到 $n' > n$，使得
$\Ker(\alpha_n)$ 到 $\Ext^{-i}(M/I^{n'}M, \omega_B^\bullet)$ 的映射为零，
至少在 $\mathfrak q$ 处局部化后如此。
由于 $\Ker(\alpha_n)$ 是有限 $A$-模，且截面的支集是拟紧的，
可找到一个 $n'$，使得 $\Ker(\alpha_n)$ 到
$\Ext^{-i}(M/I^{n'}M, \omega_B^\bullet)$ 的映射为零。
由此可见，$\Ext^{-i}(M/I^n M, \omega_B^\bullet)$ 本质上为常值，
其值为 $\Ext^{-i}(M, \omega_B^\bullet)$。证明完成。
\end{proof}

\noindent
下面是引理 \ref{algebraization-lemma-algebraization-principal} 的一个更一般版本。

\begin{lemma}
\label{algebraization-lemma-algebraization-principal-bis}
在情形 \ref{algebraization-situation-algebraize} 中，设 $(\mathcal{F}_n)$ 是
$\textit{Coh}(U, I\mathcal{O}_U)$ 的一个对象。假设
\begin{enumerate}
\item $A$ 是局部环，且 $\mathfrak a = \mathfrak m$ 是极大理想；
\item $A$ 有对偶化复形；
\item $I = (f)$ 是主理想；
\item $(\mathcal{F}_n)$ 满足 $(2, 3)$-不等式组。
\end{enumerate}
则 $(\mathcal{F}_n)$ 可延拓到 $X$。特别地，若 $A$ 是
$I$-进完备的，则 $(\mathcal{F}_n)$ 是某个凝聚
$\mathcal{O}_U$-模的完备化。
\end{lemma}

\begin{proof}
回忆 $\textit{Coh}(U, I\mathcal{O}_U)$ 是阿贝尔范畴；参见概形上同调，
引理 \ref{coherent-lemma-inverse-systems-abelian}。
在 $U$ 的仿射开集上，对象 $(\mathcal{F}_n)$ 对应于诺特环上的有限模
（概形上同调，引理 \ref{coherent-lemma-inverse-systems-affine}）。
因此，映射 $f^N : (\mathcal{F}_n) \to (\mathcal{F}_n)$ 的核在
$N$ 充分大时稳定。由引理
\ref{algebraization-lemma-map-kernel-cokernel-on-closed} 和
\ref{algebraization-lemma-essential-image-completion}，为证明本引理，
可以用这样一个映射的像替换 $(\mathcal{F}_n)$。
因此可假设 $f$ 在 $(\mathcal{F}_n)$ 上是单射。
替换后，引理 \ref{algebraization-lemma-equivalent-f-good} 的等价条件对逆系统
$(\mathcal{F}_n)$ 在 $U$ 上成立。证明余下部分将不再另行说明地使用这一点。

\medskip\noindent
我们将检验引理 \ref{algebraization-lemma-when-done} 的假设 (a)、(b) 和 (c)。
由上同调，引理 \ref{cohomology-lemma-topology-I-adic-f}，假设 (b) 成立。

\medskip\noindent
选取如引理 \ref{algebraization-lemma-system-of-modules} 所示的模逆系统 $\{M_n\}$。
可以假设 $H^0_\mathfrak m(M_n) = 0$：必要时把 $M_n$ 替换为
$M_n/H^0_\mathfrak m(M_n)$ 即可。于是得到短正合序列
$$
0 \to M_n \to H^0(U, \mathcal{F}_n) \to H^1_\mathfrak m(M_n) \to 0
$$
，对所有 $n$ 均成立。设 $E$ 是 $A$ 的剩余域的内射包。
由引理 \ref{algebraization-lemma-helper-algebraize} 和当前假设 (4)，可以选取
一个整数 $m \geq 0$、有限 $A$-模 $N_1$ 和 $N_2$；后二者被
$f^c$ 零化，其中某个 $c \geq 0$，并且存在相容映射系统
$$
H^i_\mathfrak m(M_n) \to \Hom_A(N_i, E), \quad i = 1, 2
$$
，对 $n \geq m$ 成立，并具有引理中陈述的性质。

\medskip\noindent
已知 $M = \lim H^0(U, \mathcal{F}_n)$ 是 $A$-模，
其极限拓扑为 $f$-进拓扑。因此，给定 $n$，模
$M/f^nM$ 是 $H^0(U, \mathcal{F}_N)$ 的子商，其中 $N \gg n$。
由上面得到的信息可知，$f^cM/f^nM$ 是有限 $A$-模。
由于 $f$ 在 $M$ 上是非零因子，可知 $M/f^{n - c}M$ 是有限 $A$-模。
由此可见，引理 \ref{algebraization-lemma-when-done} 的假设 (c) 成立。

\medskip\noindent
接下来研究模
$$
Ob = \lim H^1(U, \mathcal{F}_n) = \lim H^2_\mathfrak m(M_n)
$$
对 $n \geq m$，令 $K_n$ 为映射
$H^2_\mathfrak m(M_n) \to \Hom_A(N_2, E)$.
的核。令 $K = \lim K_n$。得到正合序列
$$
0 \to K \to Ob \to \Hom_A(N_2, E)
$$
由上文，$Ob = \lim H^2_\mathfrak m(M_n)$ 上的极限拓扑是
$f$-进拓扑。由于 $N_2$ 被 $f^c$ 零化，可知
$K = \lim K_n$ 上的极限拓扑也同样如此。
因此 $K/fK$ 是 $K_n$ 的子商，其中 $n \gg 1$。
然而，由于 $\{K_n\}$ 与有限长度 $A$-模的某个逆系统 pro-同构
（由引理 \ref{algebraization-lemma-helper-algebraize} 的结论），可知
$K/fK$ 是某个有限长度 $A$-模的子商。因此 $K$ 是有限 $A$-模；
参见代数，引理 \ref{algebra-lemma-finite-over-complete-ring}。
（事实上，甚至可见 $\dim(\text{Supp}(K)) = 1$，但这里不需要这一点。）

\medskip\noindent
给定 $n \geq 1$，考虑边缘映射
$$
\delta_n :
H^0(U, \mathcal{F}_n)
\longrightarrow
\lim_N H^1(U, f^n\mathcal{F}_N) \xrightarrow{f^{-n}} Ob
$$
（第二个映射是同构），它来自短正合序列
$$
0 \to f^n\mathcal{F}_N \to \mathcal{F}_N \to \mathcal{F}_n \to 0
$$
对每个 $n$，令
$$
P_n = \Im(H^0(U, \mathcal{F}_{n + m}) \to H^0(U, \mathcal{F}_n))
$$
，其中 $m$ 如上。注意，$\{P_n\}$ 是逆系统，并且映射
$f : \mathcal{F}_n \to \mathcal{F}_{n + 1}$ 在整体截面上把
$P_n$ 映入 $P_{n + 1}$。
若 $p \in P_n$，则 $\delta_n(p) \in K \subset Ob$，
因为 $\delta_n(p)$ 在
$H^1(U, f^n\mathcal{F}_{n + m}) = H^2_\mathfrak m(M_m)$
中的像为零，并且 $\delta_n$ 与 $Ob \to \Hom_A(N_2, E)$ 的复合
通过 $H^2_\mathfrak m(M_m)$ 分解；这是由 $m$ 的选取保证的。
因此
$$
\bigoplus\nolimits_{n \geq 0} \Im(P_n \to Ob)
$$
是有限生成分次 $A[T]$-模，其中 $T$ 通过乘以 $f$ 起作用。
事实上，它是 $K[T]$ 的分次子模，而 $K$ 在 $A$ 上有限。
仿照上同调，引理 \ref{cohomology-lemma-ML-general} 的证明进行论证
\footnote{为所显示的模选取形如 $\delta_{n_j}(p_j)$ 的齐次生成元。
若 $k = \max(n_j)$，则对 $n \geq k$ 和任意 $p \in P_n$，
可找到 $a_j \in A$，使得
$p - \sum a_j f^{n - n_j} p_j$ 位于 $\delta_n$ 的核中，
从而位于 $P_{n'}$ 的像中；这对所有 $n' \geq n$ 成立。
因此有
$\Im(P_n \to P_{n - k}) = \Im(P_{n'} \to P_{n - k})$，
对所有 $n' \geq n$ 均成立。}
可知逆系统 $\{P_n\}$ 满足 ML。
由于 $\{P_n\}$ 与 $\{H^0(U, \mathcal{F}_n)\}$ pro-同构，
可知 $\{H^0(U, \mathcal{F}_n)\}$ 满足 ML。
因此，引理 \ref{algebraization-lemma-when-done} 的假设 (a) 成立，证明完成。
\end{proof}

\noindent
引理 \ref{algebraization-lemma-algebraization-principal-bis} 的条件可以具体展开如下。

\begin{lemma}
\label{algebraization-lemma-unwinding-conditions}
在情形 \ref{algebraization-situation-algebraize} 中，设 $(\mathcal{F}_n)$ 为
$\textit{Coh}(U, I\mathcal{O}_U)$ 的一个对象。假设
\begin{enumerate}
\item $A$ 是局部环，其极大理想为 $\mathfrak a = \mathfrak m$；
\item $\text{cd}(A, I) = 1$.
\end{enumerate}
则 $(\mathcal{F}_n)$ 满足 $(2, 3)$-不等式，当且仅当对每个
$y \in U \cap Y$，若 $\dim(\{y\}) = 1$，以及对每个素理想
$\mathfrak p \subset \mathcal{O}_{X, y}^\wedge$，若
$\mathfrak p \not \in V(I\mathcal{O}_{X, y}^\wedge)$，便有
$$
\text{depth}((\mathcal{F}_y^\wedge)_\mathfrak p) +
\dim(\mathcal{O}_{X, y}^\wedge/\mathfrak p) > 2
$$
\end{lemma}

\begin{proof}
以下将不再另行说明地使用引理 \ref{algebraization-lemma-explain-2-3-cd-1}。
它立即表明该条件是必要的。反过来，假设该条件成立。
由引理 \ref{algebraization-lemma-discussion}，有
$\delta^Y_Z(y) = \dim(\overline{\{y\}})$；把这个函数简记为 $\delta$。
设 $y \in U \cap Y$。若 $\delta(y) > 2$，则引理
\ref{algebraization-lemma-explain-2-3-cd-1} 的不等式成立。
最后假设 $\delta(y) = 2$。我们必须证明
$$
\text{depth}((\mathcal{F}_y^\wedge)_\mathfrak p) +
\dim(\mathcal{O}_{X, y}^\wedge/\mathfrak p) > 1
$$
选取一个特化 $y \leadsto y'$，使得 $\delta(y') = 1$。此时存在环映射
$\mathcal{O}_{X, y'}^\wedge \to \mathcal{O}_{X, y}^\wedge$，
它把靶环识别为 $\mathcal{O}_{X, y'}^\wedge$ 在某个素理想
$\mathfrak q$ 处局部化之后的完备化，并且
$\dim(\mathcal{O}_{X, y'}^\wedge/\mathfrak q) = 1$。此外有
$$
\mathcal{F}_y^\wedge =
\mathcal{F}_{y'}^\wedge
\otimes_{\mathcal{O}_{X, y'}^\wedge}
\mathcal{O}_{X, y}^\wedge
$$
令 $\mathfrak p' \subset \mathcal{O}_{X, y'}^\wedge$ 为
$\mathfrak p$ 的像。由《局部上同调》，引理
\ref{local-cohomology-lemma-change-completion}，有
\begin{align*}
\text{depth}((\mathcal{F}_y^\wedge)_\mathfrak p) +
\dim(\mathcal{O}_{X, y}^\wedge/\mathfrak p)
& =
\text{depth}((\mathcal{F}_{y'}^\wedge)_{\mathfrak p'}) +
\dim((\mathcal{O}_{X, y}^\wedge/\mathfrak p)_{\mathfrak p'}) \\
& =
\text{depth}((\mathcal{F}_{y'}^\wedge)_{\mathfrak p'}) +
\dim(\mathcal{O}_{X, y}^\wedge/\mathfrak p') - 1
\end{align*}
最后一个等号成立，是因为这个特化是直接特化。
现在对 $y', \mathfrak p'$ 使用所假设的不等式，便证明了本引理。
\end{proof}

\begin{lemma}
\label{algebraization-lemma-unwinding-conditions-bis}
在情形 \ref{algebraization-situation-algebraize} 中，设 $(\mathcal{F}_n)$ 为
$\textit{Coh}(U, I\mathcal{O}_U)$ 的一个对象。假设
\begin{enumerate}
\item $A$ 是局部环，其极大理想为 $\mathfrak a = \mathfrak m$；
\item $A$ 有对偶化复形；
\item $\text{cd}(A, I) = 1$,
\item 对 $y \in U \cap Y$，模 $\mathcal{F}_y^\wedge$ 在
$V(I\mathcal{O}_{X, y}^\wedge)$ 外有限局部自由；例如，
若 $\mathcal{F}_n$ 是有限局部自由
$\mathcal{O}_U/I^n\mathcal{O}_U$-模，则此条件成立；并且
\item 下列条件之一成立：
\begin{enumerate}
\item $A_f$ 满足 $(S_2)$，且 $X$ 的每个不包含在 $Y$ 中的
不可约分支之维数都 $\geq 4$；或者
\item 若 $\mathfrak p \not \in V(f)$ 且
$V(\mathfrak p) \cap V(f) \not = \{\mathfrak m\}$，则
$\text{depth}(A_\mathfrak p) + \dim(A/\mathfrak p) > 3$。
\end{enumerate}
\end{enumerate}
则 $(\mathcal{F}_n)$ 满足 $(2, 3)$-不等式。
\end{lemma}

\begin{proof}
我们使用引理 \ref{algebraization-lemma-unwinding-conditions} 的判据。
取 $y \in U \cap Y$，它满足 $\dim(\overline{\{y\}} = 1$；并取素理想
$\mathfrak p$，它满足 $\mathfrak p \subset \mathcal{O}_{X, y}^\wedge$ 且
$\mathfrak p \not \in V(I\mathcal{O}_{X, y}^\wedge)$。
条件 (4) 表明
$\text{depth}((\mathcal{F}_y^\wedge)_\mathfrak p) =
\text{depth}((\mathcal{O}_{X, y}^\wedge)_\mathfrak p)$.
因此只需证明
$$
\text{depth}((\mathcal{O}_{X, y}^\wedge)_\mathfrak p) +
\dim(\mathcal{O}_{X, y}^\wedge/\mathfrak p) > 2
$$
令 $\mathfrak p_0 \subset A$ 为 $\mathfrak p$ 的像。
令 $\mathfrak q \subset A$ 为与 $y$ 对应的素理想。
由《局部上同调》，引理
\ref{local-cohomology-lemma-change-completion}，有
\begin{align*}
\text{depth}((\mathcal{O}_{X, y}^\wedge)_\mathfrak p) +
\dim(\mathcal{O}_{X, y}^\wedge/\mathfrak p)
& =
\text{depth}(A_{\mathfrak p_0}) + \dim((A/\mathfrak p_0)_\mathfrak q) \\
& =
\text{depth}(A_{\mathfrak p_0}) + \dim(A/\mathfrak p_0) - 1
\end{align*}
若 (5)(a) 成立，则这个量
$$
\geq \min(2, \dim(A_{\mathfrak p_0})) + \dim(A/\mathfrak p_0) - 1
$$
注意，无论如何都有 $\dim(A/\mathfrak p_0) \geq 2$。因此，若该最小值为
$2$，结论已经成立。否则，我们得到
$$
\dim(A_{\mathfrak p_0}) + \dim(A/\mathfrak p_0) - 1 \geq 4 - 1
$$
，因为 $\Spec(A)$ 的每个经过 $\mathfrak p_0$ 的分支之维数都
$\geq 4$。若 (5)(b) 成立，则立即得到结论。
\end{proof}

\begin{remark}
\label{algebraization-remark-interesting-case-bis}
设 $(A, \mathfrak m)$ 为诺特局部环，它有对偶化复形，并且关于
$f \in \mathfrak m$ 完备。设 $(\mathcal{F}_n)$ 为
$\textit{Coh}(U, f\mathcal{O}_U)$ 的一个对象，其中
$U$ 是 $A$ 的穿孔谱。置 $Y = V(f) \subset X = \Spec(A)$。
对 $y \in U \cap V(f)$，若它在 $U$ 中闭，也就是满足
$\dim(\overline{\{y\}}) = 1$ 的点，假设
$\mathcal{O}_{X, y}^\wedge$-模 $\mathcal{F}_y^\wedge$
满足下列两个条件：
\begin{enumerate}
\item $\mathcal{F}_y^\wedge[1/f]$ 满足 $(S_2)$；这里把它视为
$\mathcal{O}_{X, y}^\wedge[1/f]$-模；并且
\item 对 $\mathfrak p \in \text{Ass}(\mathcal{F}_y^\wedge[1/f])$，有
$\dim(\mathcal{O}_{X, y}^\wedge/\mathfrak p) \geq 3$。
\end{enumerate}
则 $(\mathcal{F}_n)$ 是 $U$ 上某个凝聚模的完备化。
这由引理 \ref{algebraization-lemma-algebraization-principal-bis}
和 \ref{algebraization-lemma-unwinding-conditions} 得出。
\end{remark}
































\section{改进凝聚形式模}
\label{algebraization-section-improving-formal-modules}

\noindent
设 $X$ 为诺特概形。设 $Y \subset X$ 为闭子概形，
其拟凝聚理想层为 $\mathcal{I} \subset \mathcal{O}_X$。
设 $(\mathcal{F}_n)$ 为 $\textit{Coh}(X, \mathcal{I})$ 的一个对象。
本节构造映射
$(\mathcal{F}_n) \to (\mathcal{F}'_n)$
，它们类似于《局部上同调》，第
\ref{local-cohomology-section-improve} 节中为凝聚模构造的映射。
对一点 $y \in Y$，置
$$
\mathcal{O}_{X, y}^\wedge = \lim \mathcal{O}_{X, y}/\mathcal{I}^n_y, \quad
\mathcal{I}_y^\wedge = \lim \mathcal{I}_y/\mathcal{I}^n_y
\quad\text{且}\quad
\mathfrak m_y^\wedge = \lim \mathfrak m_y/\mathcal{I}_y^n
$$
则 $\mathcal{O}_{X, y}^\wedge$ 是诺特局部环，其极大理想为
$\mathfrak m_y^\wedge$，并且关于
$\mathcal{I}_y^\wedge = \mathcal{I}_y\mathcal{O}_{X, y}^\wedge$ 完备。
还置
$$
\mathcal{F}_y^\wedge = \lim \mathcal{F}_{n, y}
$$
则 $\mathcal{F}_y^\wedge$ 是
$\mathcal{O}_{X, y}^\wedge$ 上的有限模，并且
$\mathcal{F}_y^\wedge/(\mathcal{I}_y^\wedge)^n\mathcal{F}_y^\wedge =
\mathcal{F}_{n, y}$ 对所有 $n$ 成立；参见《代数》，引理
\ref{algebra-lemma-limit-complete} 和
\ref{algebra-lemma-finite-over-complete-ring}.

\begin{lemma}
\label{algebraization-lemma-divide-torsion-formal-coherent-module}
在上述情形中，假设 $X$ 局部有对偶化复形。
设 $T \subset Y$ 为对特化稳定的子集。
假设对 $y \in T$ 以及满足
$\mathfrak p \subset \mathcal{O}_{X, y}^\wedge$ 的非极大素理想，
若 $V(\mathfrak p) \cap V(\mathcal{I}^\wedge_y) = \{\mathfrak m_y^\wedge\}$，
则
$$
\text{depth}_{(\mathcal{O}_{X, y})_\mathfrak p}
((\mathcal{F}^\wedge_y)_\mathfrak p) > 0
$$
于是存在典范映射
$(\mathcal{F}_n) \to (\mathcal{F}_n')$
，它是凝聚 $\mathcal{O}_X$-模逆系统之间的映射，并满足下列性质：
\begin{enumerate}
\item 对 $y \in T$，有 $\text{depth}(\mathcal{F}'_{n, y}) \geq 1$；
\item $(\mathcal{F}'_n)$ 作为 pro-系统同构于一个对象
$(\mathcal{G}_n)$，后者属于 $\textit{Coh}(X, \mathcal{I})$；
\item 所诱导的态射
$(\mathcal{F}_n) \to (\mathcal{G}_n)$ 在
$\textit{Coh}(X, \mathcal{I})$ 中是满射，其核被
$\mathcal{I}$ 的某个幂零化。
\end{enumerate}
\end{lemma}

\begin{proof}
对每个 $n$，令 $\mathcal{F}_n \to \mathcal{F}'_n$ 为
《局部上同调》，引理
\ref{local-cohomology-lemma-get-depth-1-along-Z} 中构造的满射。
这是 $\mathcal{F}_n$ 对支撑于 $T$ 的截面子层取商所得，故得到典范映射
$\mathcal{F}'_{n + 1} \to \mathcal{F}'_n$，从而得到映射
$(\mathcal{F}_n) \to (\mathcal{F}_n')$；它是凝聚
$\mathcal{O}_X$-模逆系统之间的映射。
性质 (1) 由构造成立。

\medskip\noindent
为证明性质 (2) 和 (3)，可以假设 $X = \Spec(A_0)$ 是仿射的，
且 $A_0$ 有对偶化复形。令 $I_0 \subset A_0$ 为与 $Y$ 对应的理想。
令 $A, I$ 为关于 $I_0$ 的进完备化，分别对应 $A_0, I_0$。
供后文使用，注意 $A$ 有对偶化复形
（《对偶化复形》，引理 \ref{dualizing-lemma-ubiquity-dualizing}）。
令 $M$ 为有限 $A$-模，它对应于 $(\mathcal{F}_n)$；参见
《概形上同调》，引理 \ref{coherent-lemma-inverse-systems-affine}。
于是 $\mathcal{F}_n$ 对应于 $M_n = M/I^nM$。回忆
$\mathcal{F}'_n$ 对应于商模 $M'_n = M_n / H^0_T(M_n)$；参见
《局部上同调》，引理 \ref{local-cohomology-lemma-get-depth-1-along-Z}
及其证明。

\medskip\noindent
置 $s = 0$ 和 $d = \text{cd}(A, I)$。
我们断言 $A, I, T, M, s, d$ 满足情形
\ref{algebraization-situation-bootstrap} 的假设 (1)、(3)、(4)、(6)。
事实上，(1) 和 (3) 由上文立即成立；由于 $s = 0$，(4) 是空条件；
而 (6) 正是本引理陈述中的假设。

\medskip\noindent
由定理 \ref{algebraization-theorem-final-bootstrap} 可知，$\{H^0_T(M_n)\}$
本质常值，其值为 $H^0_T(M)$。因此，短正合序列
$0 \to H^0_T(M_n) \to M_n \to M'_n \to 0$ 的极限是短正合序列
$$
0 \to H^0_T(M) \to M \to \lim M'_n \to 0
$$
置 $M' = \lim M'_n = M/H^0_T(M)$。这是有限 $A$-模。
令 $(\mathcal{G}_n)$ 为 $\textit{Coh}(X, \mathcal{I})$ 中与
$M'$ 对应的对象。为完成证明，必须证明典范映射
$\{M'/I^nM'\} \to \{M'_n\}$ 是 pro-同构。
这等价于断言
$\{H^0_T(M) + I^nM\} \to \{\Ker(M \to M'_n)\}$ 是 pro-同构。
除以 $I^nM$ 后，只需证明
$\{H^0_T(M)/H^0_T(M) \cap I^nM\} \to \{H^0_T(M_n)\}$
是 pro-同构。由于 $H^0_T(M)$ 被 $I$ 的某个幂零化，
由 Artin--Rees 引理可知，$H^0_T(M) \cap I^nM = 0$，
对所有 $n \gg 0$ 均成立。而上面已经看到
$\{H^0_T(M_n)\}$ 本质常值，其值为 $H^0_T(M)$；
故得到所需的 pro-同构。
\end{proof}

\begin{lemma}
\label{algebraization-lemma-improvement-formal-coherent-module-better}
在上述情形中，假设 $X$ 局部有对偶化复形。
设 $T' \subset T \subset Y$ 为对特化稳定的子集。
设 $d \geq 0$ 为整数。假设
\begin{enumerate}
\item[(a)] 在仿射局部有 $X = \Spec(A_0)$ 和 $Y = V(I_0)$，
并且 $\text{cd}(A_0, I_0) \leq d$；
\item[(b)] 对 $y \in T$ 以及满足
$\mathfrak p \subset \mathcal{O}_{X, y}^\wedge$ 的非极大素理想，若
$V(\mathfrak p) \cap V(\mathcal{I}_y^\wedge) = \{\mathfrak m_y^\wedge\}$，
则
$$
\text{depth}_{(\mathcal{O}_{X, y})_\mathfrak p}
((\mathcal{F}^\wedge_y)_\mathfrak p) > 0
$$
\item[(c)] 对 $y \in T'$ 以及满足
$\mathfrak p \subset \mathcal{O}_{X, y}^\wedge$ 的素理想，若
$\mathfrak p \not \in V(\mathcal{I}_y^\wedge)$ 且
$V(\mathfrak p) \cap V(\mathcal{I}_y^\wedge) \not =
\{\mathfrak m_y^\wedge\}$，则
$$
\text{depth}_{(\mathcal{O}_{X, y})_\mathfrak p}
((\mathcal{F}^\wedge_y)_\mathfrak p) \geq 1
\quad\text{或}\quad
\text{depth}_{(\mathcal{O}_{X, y})_\mathfrak p}
((\mathcal{F}^\wedge_y)_\mathfrak p) +
\dim(\mathcal{O}_{X, y}^\wedge/\mathfrak p) > 1 + d
$$
\item[(d)] 对 $y \in T'$ 以及满足
$\mathfrak p \subset \mathcal{O}_{X, y}^\wedge$ 的非极大素理想，若
$V(\mathfrak p) \cap V(\mathcal{I}_y^\wedge) = \{\mathfrak m_y^\wedge\}$，
则
$$
\text{depth}_{(\mathcal{O}_{X, y})_\mathfrak p}
((\mathcal{F}^\wedge_y)_\mathfrak p) > 1
$$
\item[(e)] 若 $y \leadsto y'$ 是直接特化且
$y' \in T'$，则 $y \in T$。
\end{enumerate}
于是存在典范映射 $(\mathcal{F}_n) \to (\mathcal{F}_n'')$，
它是凝聚 $\mathcal{O}_X$-模逆系统之间的映射，并满足下列性质：
\begin{enumerate}
\item 对 $y \in T$，有 $\text{depth}(\mathcal{F}''_{n, y}) \geq 1$；
\item 对 $y' \in T'$，有 $\text{depth}(\mathcal{F}''_{n, y'}) \geq 2$；
\item $(\mathcal{F}''_n)$ 作为 pro-系统同构于一个对象
$(\mathcal{H}_n)$，它属于 $\textit{Coh}(X, \mathcal{I})$；
\item 所诱导的态射 $(\mathcal{F}_n) \to (\mathcal{H}_n)$ 在
$\textit{Coh}(X, \mathcal{I})$ 中的核与余核被
$\mathcal{I}$ 的某个幂零化。
\end{enumerate}
\end{lemma}

\begin{proof}
仿照引理 \ref{algebraization-lemma-divide-torsion-formal-coherent-module} 及其证明，
对每个 $n$，令 $\mathcal{F}_n \to \mathcal{F}'_n$ 为
《局部上同调》，引理
\ref{local-cohomology-lemma-get-depth-1-along-Z} 中利用 $T$ 构造的满射。
接着，令 $\mathcal{F}'_n \to \mathcal{F}''_n$ 为
《局部上同调》，引理
\ref{local-cohomology-lemma-make-S2-along-T} 及其证明中构造的单射。
这些构造表明，得到典范映射
$\mathcal{F}''_{n + 1} \to \mathcal{F}''_n$，从而得到映射
$$
(\mathcal{F}_n) \longrightarrow (\mathcal{F}_n') \longrightarrow
(\mathcal{F}''_n)
$$
，它们是凝聚 $\mathcal{O}_X$-模逆系统之间的映射。
性质 (1) 和 (2) 由构造成立。

\medskip\noindent
为证明性质 (3) 和 (4)，可以假设 $X = \Spec(A_0)$ 是仿射的，
且 $A_0$ 有对偶化复形。令 $I_0 \subset A_0$ 为与 $Y$ 对应的理想。
令 $A, I$ 为关于 $I_0$ 的进完备化，分别对应 $A_0, I_0$。
供后文使用，注意 $A$ 有对偶化复形
（《对偶化复形》，引理 \ref{dualizing-lemma-ubiquity-dualizing}）。
令 $M$ 为有限 $A$-模，它对应于 $(\mathcal{F}_n)$；参见
《概形上同调》，引理 \ref{coherent-lemma-inverse-systems-affine}。
于是 $\mathcal{F}_n$ 对应于 $M_n = M/I^nM$。回忆
$\mathcal{F}'_n$ 对应于商模 $M'_n = M_n / H^0_T(M_n)$。
还回忆，$M' = \lim M'_n$ 是 $M$ 对 $H^0_T(M)$ 取的商，
$H^0_T(M)$ 是 $I$-幂挠模，并且
$\{M'_n\}$ 与 $\{M'/I^nM'\}$ 作为 pro-系统同构。
最后，$\mathcal{F}''_n$ 对应于扩张
$$
0 \to M'_n \to M''_n \to H^1_{T'}(M'_n) \to 0
$$
；参见《局部上同调》，引理
\ref{local-cohomology-lemma-make-S2-along-T} 的证明。

\medskip\noindent
置 $s = 1$。我们断言 $A, I, T', M', s, d$ 满足情形
\ref{algebraization-situation-bootstrap} 的假设 (1)、(3)、(4)、(6)。事实上，
(1) 和 (3) 立即成立，(c) 蕴含 (4)，而 (d) 蕴含 (6)。
我们省略 (c) $\Rightarrow$ (4) 和 (d) $\Rightarrow$ (6) 的证明。

\medskip\noindent
由定理 \ref{algebraization-theorem-final-bootstrap} 可知，$\{H^1_{T'}(M'/I^nM')\}$
本质常值，其值为 $H^1_{T'}(M')$。由此推出
$\{H^1_{T'}(M'_n)\}$ 本质常值，其值为 $H^1_{T'}(M')$。
因此，上面所显示的短正合序列之极限是短正合序列
$$
0 \to M' \to \lim M''_n \to H^1_{T'}(M') \to 0
$$
置 $M'' = \lim M''_n$。由《局部上同调》，命题
\ref{local-cohomology-proposition-finiteness} 可知，
$H^1_{T'}(M')$ 以及因而 $M''$ 都是有限 $A$-模
（验证条件时，使用 $M'$ 的深度至少为 $1$，这是对
$T - T'$ 中各素理想而言）。
令 $(\mathcal{H}_n)$ 为 $\textit{Coh}(X, \mathcal{I})$ 中与有限
$A$-模 $M''$ 对应的对象。为完成证明，必须证明典范映射
$\{M''/I^nM''\} \to \{M''_n\}$ 是 pro-同构。
已知 $\{M'/I^nM'\}$ 与 $\{M'_n\}$ pro-同构；读者可验证
（此处省略），这等价于要求
$\{H^1_{T'}(M')/I^nH^1_{T'}(M')\} \to \{H^1_{T'}(M'_n)\}$
是 pro-同构。由于
$H^1_{T'}(M')/I^nH^1_{T'}(M') = H^1_{T'}(M')$
在 $n$ 充分大时成立，
而上面已经看到 $\{H^1_{T'}(M'_n)\}$
本质常值，其值为 $H^1_{T'}(M')$，故结论成立。
\end{proof}

\begin{lemma}
\label{algebraization-lemma-improvement-application}
在情形 \ref{algebraization-situation-algebraize} 中，假设 $A$ 有对偶化复形。
设 $d \geq \text{cd}(A, I)$。设 $(\mathcal{F}_n)$ 为
$\textit{Coh}(U, I\mathcal{O}_U)$ 的一个对象。假设
$(\mathcal{F}_n)$ 满足 $(2, 2 + d)$-不等式；参见定义
\ref{algebraization-definition-s-d-inequalities}。
于是存在典范映射 $(\mathcal{F}_n) \to (\mathcal{F}_n'')$，
它是凝聚 $\mathcal{O}_U$-模逆系统之间的映射，并满足下列性质：
\begin{enumerate}
\item $\text{depth}(\mathcal{F}''_{n, y}) + \delta^Y_Z(y) \geq 3$
对所有 $y \in U \cap Y$ 成立；
\item $(\mathcal{F}''_n)$ 作为 pro-系统同构于一个对象
$(\mathcal{H}_n)$，后者属于 $\textit{Coh}(U, I\mathcal{O}_U)$；
\item 所诱导的态射 $(\mathcal{F}_n) \to (\mathcal{H}_n)$ 在
$\textit{Coh}(U, I\mathcal{O}_U)$ 中的核与余核被
$I$ 的某个幂零化；
\item 模 $H^0(U, \mathcal{F}''_n)$ 和 $H^1(U, \mathcal{F}''_n)$
都是有限 $A$-模，对所有 $n$ 均成立。
\end{enumerate}
\end{lemma}

\begin{proof}
存在性以及性质 (1)、(2)、(3) 由引理
\ref{algebraization-lemma-improvement-formal-coherent-module-better} 得出：把该引理应用于
$U$、$U \cap Y$、
$T = \{y \in U \cap Y : \delta^Y_Z(y) \leq 2\}$、
$T' = \{y \in U \cap Y : \delta^Y_Z(y) \leq 1\}$ 和
$(\mathcal{F}_n)$ 即可。模 $H^0(U, \mathcal{F}''_n)$ 和
$H^1(U, \mathcal{F}''_n)$ 的有限性由《局部上同调》，引理
\ref{local-cohomology-lemma-finiteness-Rjstar}，以及函数
$\delta^Y_Z(-)$ 在引理 \ref{algebraization-lemma-discussion} 中证明的初等性质得出。
\end{proof}











\section{凝聚形式模的代数化，V}
\label{algebraization-section-algebraization-modules-conclusion}

\noindent
本节证明关于凝聚形式模代数化的最一般结果。先在理想的上同调维数为
$1$ 时证明，再把所得结果应用于一次吹起，以证明更一般的结论。

\begin{lemma}
\label{algebraization-lemma-cd-1-canonical}
在情形 \ref{algebraization-situation-algebraize} 中，设 $(\mathcal{F}_n)$ 为
$\textit{Coh}(U, I\mathcal{O}_U)$ 的一个对象。假设
\begin{enumerate}
\item $A$ 有对偶化复形，且 $\text{cd}(A, I) = 1$；
\item $(\mathcal{F}_n)$ pro-同构于逆系统
$(\mathcal{F}_n'')$，后者由凝聚 $\mathcal{O}_U$-模组成，并且
$\text{depth}(\mathcal{F}''_{n, y}) + \delta^Y_Z(y) \geq 3$
对所有 $y \in U \cap Y$ 成立。
\end{enumerate}
则 $(\mathcal{F}_n)$ 典范地延拓到 $X$；参见定义
\ref{algebraization-definition-canonically-algebraizable}。
\end{lemma}

\begin{proof}
我们检验引理 \ref{algebraization-lemma-when-done} 的假设 (a)、(b) 和 (c)。
开始之前，先指出模
$H^0(U, \mathcal{F}''_n)$ 和 $H^1(U, \mathcal{F}''_n)$
都是有限 $A$-模，对所有 $n$ 均如此；这是由《局部上同调》，引理
\ref{local-cohomology-lemma-finiteness-Rjstar} 得出的。

\medskip\noindent
注意，对每个 $p \geq 0$，由引理 \ref{algebraization-lemma-topology-I-adic}，
$\lim H^p(U, \mathcal{F}_n)$ 上的极限拓扑是 $I$-进拓扑。
特别地，假设 (b) 成立。

\medskip\noindent
已知 $M = \lim H^0(U, \mathcal{F}_n)$ 是 $A$-模，
其极限拓扑为 $I$-进拓扑。因此，给定 $n$，模
$M/I^nM$ 是 $H^0(U, \mathcal{F}_N)$ 的子商，其中 $N \gg n$。
逆系统 $\{H^0(U, \mathcal{F}_N)\}$ pro-同构于有限 $A$-模的逆系统，
即 $\{H^0(U, \mathcal{F}''_N)\}$，故 $M/I^nM$ 有限。
由此 $M$ 有限；参见《代数》，引理
\ref{algebra-lemma-finite-over-complete-ring}。
特别地，假设 (c) 成立。

\medskip\noindent
对每个 $n \geq 0$，记 $Ob_n = \lim_N H^1(U, I^n\mathcal{F}_N)$。
一个特例是 $Ob = Ob_0 = \lim_N H^1(U, \mathcal{F}_N)$。
与上一段完全相同的论证表明，$Ob$ 是有限 $A$-模。
（事实上，还知道 $Ob/I Ob$ 被 $\mathfrak a$ 的某个幂零化，
但这似乎不太容易使用。）

\medskip\noindent
置 $\mathcal{F} = \lim \mathcal{F}_n$，选取生成元
$f_1, \ldots, f_r$，它们生成 $I$；再选取 $c \geq 1$，并像引理
\ref{algebraization-lemma-cd-is-one-for-system} 中那样选取 $\Phi_\mathcal{F}$。
以下将不再另行说明地使用引理 \ref{algebraization-lemma-properties-system} 的结论。
特别地，对每个 $n \geq 1$ 都有映射
$$
\delta_n :
H^0(U, \mathcal{F}_n)
\longrightarrow
H^1(U, I^n\mathcal{F})
\longrightarrow
Ob_n
$$
第一个映射来自短正合序列
$0 \to I^n\mathcal{F} \to \mathcal{F} \to \mathcal{F}_n \to 0$
，第二个映射来自 $I^n\mathcal{F} = \lim I^n\mathcal{F}_N$。
稍后将使用如下事实：若 $\delta_n(s) = 0$，其中
$s \in H^0(U, \mathcal{F}_n)$，则对每个 $n' \geq n$，
都可找到 $s' \in H^0(U, \mathcal{F}_{n'})$ 映到 $s$。
注意有交换图
$$
\xymatrix{
H^0(U, \mathcal{F}_{nc}) \ar[r] \ar[dd] &
H^1(U, I^{nc}\mathcal{F}) \ar[dd] \ar[rd]^{\Phi_\mathcal{F}} \\
& &
\bigoplus_{e_1 + \ldots + e_r = n}
H^1(U, \mathcal{F}) \cdot T_1^{e_1} \ldots T_r^{e_r} \ar[ld] \\
H^0(U, \mathcal{F}_n) \ar[r] &
H^1(U, I^n\mathcal{F})
}
$$
由此可知，障碍映射
$H^0(U, \mathcal{F}_n) \to Ob_n$
把映射
$H^0(U, \mathcal{F}_{nc}) \to H^0(U, \mathcal{F}_n)$
的像送入子模
$$
Ob'_n =
\Im\left(
\bigoplus\nolimits_{e_1 + \ldots + e_r = n}
Ob \cdot T_1^{e_1} \ldots T_r^{e_r} \to Ob_n
\right)
$$
；在直和项 $Ob \cdot T_1^{e_1} \ldots T_r^{e_r}$ 上，
这里使用的上同调映射来自对 $I$ 的各次幂取模后的约化，
所约化的乘法映射为
$f_1^{e_1} \ldots f_r^{e_r} : \mathcal{F} \to I^n\mathcal{F}$。
由构造，
$$
\bigoplus\nolimits_{n \geq 0} Ob'_n
$$
是 Rees 代数 $\bigoplus_{n \geq 0} I^n$ 上的有限分次模。
对每个 $n$，置
$$
M_n = \{s \in H^0(U, \mathcal{F}_n) \mid \delta_n(s) \in Ob'_n\}
$$
注意，$\{M_n\}$ 是逆系统，并且
$f_j : \mathcal{F}_n \to \mathcal{F}_{n + 1}$ 在整体截面上把
$M_n$ 映入 $M_{n + 1}$。
仿照《上同调》，引理 \ref{cohomology-lemma-ML-general} 的证明作完全相同的论证，
可知 $\{M_n\}$ 满足 ML。具体而言，由于 Rees 代数是诺特的，
可以为分次子模选取有限多个形如 $\delta_{n_j}(z_j)$ 的齐次生成元，
其中 $z_j \in M_{n_j}$；该分次子模为
$\bigoplus_{n \geq 0} \Im(M_n \to Ob'_n)$.
若 $k = \max(n_j)$，则对 $n \geq k$ 和任意 $z \in M_n$，
可找到 $a_j \in I^{n - n_j}$，使得
$z - \sum a_j z_j$ 位于 $\delta_n$ 的核中，因而位于
$M_{n'}$ 的像中，对所有 $n' \geq n$ 均成立
（因为 $\delta_n$ 的消失意味着，可以把
$z - \sum a_j z_j$ 提升为元素
$z' \in H^0(U, \mathcal{F}_{n'c})$，对所有 $n' \ge n$ 均可如此；
而后 $z'$ 在 $H^0(U, \mathcal{F}_{n'})$ 中的像由上文可知属于
$M_{n'}$）。因此
$\Im(M_n \to M_{n - k}) = \Im(M_{n'} \to M_{n - k})$
对所有 $n' \geq n$ 均成立。

\medskip\noindent
选取 $n$。由刚刚证明的 $\{M_n\}$ 的 Mittag--Leffler 性质，
可以找到 $n' \geq n$，使得 $M_{n'} \to M_n$ 的像
等于 $M' \to M_n$ 的像。由上文可知，
$M' \to M_n$ 的像包含
$H^0(U, \mathcal{F}_{n'c}) \to H^0(U, \mathcal{F}_n)$.
的像。因此，$\{M_n\}$ 与 $\{H^0(U, \mathcal{F}_n)\}$
pro-同构。于是 $\{H^0(U, \mathcal{F}_n)\}$ 满足 ML，
最终得到假设 (a) 成立。证明完成。
\end{proof}

\begin{proposition}[上同调维数为 1 时的代数化]
\label{algebraization-proposition-cd-1}
\begin{reference}
本结果的局部情形见 \cite[IV Corollaire 2.9]{MRaynaud-book}。
\end{reference}
在情形 \ref{algebraization-situation-algebraize} 中，设 $(\mathcal{F}_n)$ 为
$\textit{Coh}(U, I\mathcal{O}_U)$ 的一个对象。假设
\begin{enumerate}
\item $A$ 有对偶化复形，且 $\text{cd}(A, I) = 1$；
\item $(\mathcal{F}_n)$ 满足 $(2, 3)$-不等式；参见定义
\ref{algebraization-definition-s-d-inequalities}。
\end{enumerate}
则 $(\mathcal{F}_n)$ 延拓到 $X$。特别地，若 $A$ 是
$I$-进完备的，则 $(\mathcal{F}_n)$ 是某个凝聚
$\mathcal{O}_U$-模的完备化。
\end{proposition}

\begin{proof}
由引理 \ref{algebraization-lemma-map-kernel-cokernel-on-closed}，可以把
$(\mathcal{F}_n)$ 替换为对象 $(\mathcal{H}_n)$，它属于
$\textit{Coh}(U, I\mathcal{O}_U)$，并且是在引理
\ref{algebraization-lemma-improvement-application} 中找到的。因此，可以假设 $(\mathcal{F}_n)$ pro-同构于
逆系统 $(\mathcal{F}_n'')$，后者具有引理
\ref{algebraization-lemma-improvement-application} 中所述的性质。
引理 \ref{algebraization-lemma-cd-1-canonical} 已经证明
$(\mathcal{F}_n)$ 典范地延拓到 $X$。
最后一句由引理 \ref{algebraization-lemma-canonically-algebraizable} 得出。
\end{proof}

\begin{lemma}
\label{algebraization-lemma-blowup}
在情形 \ref{algebraization-situation-algebraize} 中，设 $(\mathcal{F}_n)$ 为
$\textit{Coh}(U, I\mathcal{O}_U)$ 的一个对象。假设
\begin{enumerate}
\item $A$ 有对偶化复形；
\item 吹起 $b : X' \to X$（沿着 $I$）的所有纤维之维数都
$\leq d - 1$；
\item 下列条件之一成立：
\begin{enumerate}
\item $(\mathcal{F}_n)$ 满足 $(d + 1, d + 2)$-不等式
（定义 \ref{algebraization-definition-s-d-inequalities}）；或者
\item 对 $y \in U \cap Y$ 以及素理想
$\mathfrak p \subset \mathcal{O}_{X, y}^\wedge$，若
$\mathfrak p \not \in V(I\mathcal{O}_{X, y}^\wedge)$
，则
$$
\text{depth}((\mathcal{F}^\wedge_y)_\mathfrak p) +
\dim(\mathcal{O}_{X, y}^\wedge/\mathfrak p) + \delta^Y_Z(y) > d + 2
$$
\end{enumerate}
\end{enumerate}
则 $(\mathcal{F}_n)$ 延拓到 $X$。
\end{lemma}

\begin{proof}
令 $Y' \subset X'$ 为例外除子。
令 $Z' \subset Y'$ 为 $Z \subset Y$ 的逆像。
则 $U' = X' \setminus Z'$ 是 $U$ 的逆像。
使用 (\ref{algebraization-equation-delta-Z}) 中的 $\delta^{Y'}_{Z'}$，置
$$
T' = \{y' \in Y' \mid \delta^{Y'}_{Z'}(y') = 1\}
\subset
T = \{y' \in Y' \mid \delta^{Y'}_{Z'}(y') = 1\text{ 或 }2\}
$$
由引理 \ref{algebraization-lemma-discussion} 的 (1) 和 (2)，它们是
$U' \cap Y' = Y' \setminus Z'$ 中对特化稳定的子集。
考虑对象 $(b|_{U'}^*\mathcal{F}_n)$；它属于
$\textit{Coh}(U', I\mathcal{O}_{U'})$。参见《概形上同调》，引理
\ref{coherent-lemma-inverse-systems-pullback}。
对 $y' \in U' \cap Y'$，记
$$
\mathcal{F}_{y'}^\wedge = \lim (b|_{U'}^*\mathcal{F}_n)_{y'}
$$
为该拉回在 $y'$ 处的“茎”。我们断言，引理
\ref{algebraization-lemma-improvement-formal-coherent-module-better} 的条件
(a)、(b)、(c)、(d)、(e) 对对象
$(b|_{U'}^*\mathcal{F}_n)$ 成立；该对象定义在 $U'$ 上，其中把 $d$ 换成 $1$，
并取子集 $T' \subset T \subset U' \cap Y'$。
条件 (a) 成立，因为 $Y'$ 是有效 Cartier 除子，故局部由
$1$ 个方程截出。条件 (e) 由引理 \ref{algebraization-lemma-discussion} 的 (7) 成立。
为证明 (b)、(c)、(d)，需要作一些准备。

\medskip\noindent
设 $y' \in U' \cap Y'$，并设
$\mathfrak p' \subset \mathcal{O}_{X', y'}^\wedge$
为不包含在 $V(I\mathcal{O}_{X', y'}^\wedge)$ 中的素理想。
记 $y = b(y') \in U \cap Y$。选取 $f \in I$，使得
$y'$ 包含在仿射吹起代数 $A[\frac{I}{f}]$ 的谱中；参见《除子》，引理
\ref{divisors-lemma-blowing-up-affine}。
对任意 $A$-代数 $B$，记 $B' = B[\frac{IB}{f}]$ 为相应的仿射吹起代数。
把 $I$-进完备化记作 ${\ }^\wedge$。由 $f$ 的选取，得到环映射
$(\mathcal{O}_{X, y}^\wedge)' \to \mathcal{O}_{X', y'}^\wedge$.
若令 $\mathfrak q' \subset (\mathcal{O}_{X, y}^\wedge)'$
为 $\mathfrak m_{y'}^\wedge$ 的逆像，则
$((\mathcal{O}_{X, y}^\wedge)'_{\mathfrak q'})^\wedge =
\mathcal{O}_{X', y'}^\wedge$.
令 $\mathfrak p \subset \mathcal{O}_{X, y}^\wedge$ 为相应的素理想。
此时有交换图
$$
\xymatrix{
\mathcal{O}_{X, y}^\wedge \ar[d] \ar[r] &
(\mathcal{O}_{X, y}^\wedge)' \ar[d]_\alpha \ar[r] &
(\mathcal{O}_{X, y}^\wedge)'_{\mathfrak q'} \ar[d] \ar[r]_\beta &
\mathcal{O}_{X', y'}^\wedge \ar[d] \\
\mathcal{O}_{X, y}^\wedge/\mathfrak p \ar[r] &
(\mathcal{O}_{X, y}^\wedge/\mathfrak p)' \ar[r] &
(\mathcal{O}_{X, y}^\wedge/\mathfrak p)'_{\mathfrak q'} \ar[r]^\gamma &
((\mathcal{O}_{X, y}^\wedge/\mathfrak p)'_{\mathfrak q'})^\wedge \ar[d] \\
 & & &
\mathcal{O}_{X', y'}^\wedge/\mathfrak p'
}
$$
，其中竖直箭头都是满射。由《代数续篇》，引理
\ref{more-algebra-lemma-completion-dimension} 和维数公式
（《代数》，引理 \ref{algebra-lemma-dimension-formula}），有
$$
\dim(((\mathcal{O}_{X, y}^\wedge/\mathfrak p)'_{\mathfrak q'})^\wedge) =
\dim((\mathcal{O}_{X, y}^\wedge/\mathfrak p)'_{\mathfrak q'}) =
\dim(\mathcal{O}_{X, y}^\wedge/\mathfrak p)
- \text{trdeg}(\kappa(y')/\kappa(y))
$$
沿着拉回、茎、局部化和完备化的定义逐一追踪，可得
$$
(\mathcal{F}_y^\wedge)_{\mathfrak p}
\otimes_{(\mathcal{O}_{X, y}^\wedge)_\mathfrak p}
(\mathcal{O}_{X', y'}^\wedge)_{\mathfrak p'}
=
(\mathcal{F}_{y'}^\wedge)_{\mathfrak p'}
$$
细节从略。图中的环映射 $\beta$ 和 $\gamma$ 都是平坦的，
其纤维为 Gorenstein（从而 Cohen--Macaulay）环，因为这些映射是
对具有对偶化复形的环取完备化所得。参见《对偶化复形》，引理
\ref{dualizing-lemma-formal-fibres-gorenstein} 和
\ref{dualizing-lemma-dualizing-gorenstein-formal-fibres}，以及
《代数续篇》，第 \ref{more-algebra-section-properties-formal-fibres} 节中的讨论。
注意
$(\mathcal{O}_{X, y}^\wedge)_\mathfrak p =
(\mathcal{O}_{X, y}^\wedge)'_{\tilde{\mathfrak p}}$
，其中 $\tilde{\mathfrak p}$ 是图中 $\alpha$ 的核。另一方面，
$(\mathcal{O}_{X, y}^\wedge)'_{\tilde{\mathfrak p}}
\to (\mathcal{O}_{X', y'}^\wedge)_{\mathfrak p'}$
由上文可知是纤维为 CM 环的平坦映射。因此
$(\mathcal{O}_{X, y}^\wedge)_\mathfrak p  \to
(\mathcal{O}_{X', y'}^\wedge)_{\mathfrak p'}$ 也是纤维为 CM 环的平坦映射。
使用《代数》，引理 \ref{algebra-lemma-apply-grothendieck-module}，可得
$$
\text{depth}((\mathcal{F}_{y'}^\wedge)_{\mathfrak p'}) =
\text{depth}((\mathcal{F}_y^\wedge)_{\mathfrak p}) +
\dim(F_\mathfrak r)
$$
，其中 $F$ 是
$(\mathcal{O}_{X, y}^\wedge/\mathfrak p)'_{\mathfrak q'}$
的泛形式纤维，而 $\mathfrak r$ 是与 $\mathfrak p'$ 对应的素理想。
由于 $(\mathcal{O}_{X, y}^\wedge/\mathfrak p)'_{\mathfrak q'}$
是泛链状局部整环，由 Ratliff 定理，其 $I$-进完备化
是等维且（泛）链状的
（《代数续篇》，命题 \ref{more-algebra-proposition-ratliff}）。
从而
$$
\dim(((\mathcal{O}_{X, y}^\wedge/\mathfrak p)'_{\mathfrak q'})^\wedge) =
\dim(F_\mathfrak r) + \dim(\mathcal{O}_{X', y'}^\wedge/\mathfrak p')
$$
结合引理 \ref{algebraization-lemma-change-distance-function}，得到
\begin{equation}
\label{algebraization-equation-one}
\begin{aligned}
&
\text{depth}((\mathcal{F}_{y'}^\wedge)_{\mathfrak p'}) +
\delta^{Y'}_{Z'}(y') \\
& =
\text{depth}((\mathcal{F}_y^\wedge)_{\mathfrak p}) +
\dim(F_\mathfrak r) + \delta^{Y'}_{Z'}(y') \\
& \geq
\text{depth}((\mathcal{F}_y^\wedge)_{\mathfrak p}) + \delta^Y_Z(y) +
\dim(F_\mathfrak r) + \text{trdeg}(\kappa(y')/\kappa(y)) - (d - 1) \\
& =
\text{depth}((\mathcal{F}_y^\wedge)_{\mathfrak p}) + \delta^Y_Z(y) - (d - 1)
+ \dim(\mathcal{O}_{X, y}^\wedge/\mathfrak p) -
\dim(\mathcal{O}_{X', y'}^\wedge/\mathfrak p')
\end{aligned}
\end{equation}
请注意
$\dim(\mathcal{O}_{X, y}^\wedge/\mathfrak p) \geq
\dim(\mathcal{O}_{X', y'}^\wedge/\mathfrak p')$。将其改写，得到
\begin{equation}
\label{algebraization-equation-two}
\begin{aligned}
&
\text{depth}((\mathcal{F}_{y'}^\wedge)_{\mathfrak p'}) +
\dim(\mathcal{O}_{X', y'}^\wedge/\mathfrak p') +
\delta^{Y'}_{Z'}(y') \\
& \geq
\text{depth}((\mathcal{F}_y^\wedge)_{\mathfrak p}) +
\dim(\mathcal{O}_{X, y}^\wedge/\mathfrak p) +
\delta^Y_Z(y) - (d - 1)
\end{aligned}
\end{equation}
这个不等式将使我们能够检验余下条件。

\medskip\noindent
检验引理 \ref{algebraization-lemma-improvement-formal-coherent-module-better} 的条件 (b) 和 (d)。
假设
$V(\mathfrak p') \cap V(I\mathcal{O}_{X', y'}^\wedge) =
\{\mathfrak m_{y'}^\wedge\}$.
这蕴含 $\dim(\mathcal{O}_{X', y'}^\wedge/\mathfrak p') = 1$，
因为 $Z'$ 是有效 Cartier 除子。
(b) 与 (d) 合起来等价于
$$
\text{depth}((\mathcal{F}_{y'}^\wedge)_{\mathfrak p'}) + \delta^{Y'}_{Z'}(y')
> 2
$$
若 $(\mathcal{F}_n)$ 满足 (3)(b) 中的不等式，
应用 (\ref{algebraization-equation-two}) 便立即得到这个结论。
若 $(\mathcal{F}_n)$ 满足 (3)(a)，即
$(d + 1, d + 2)$-不等式，则无论如何都有
$$
\text{depth}((\mathcal{F}_y^\wedge)_{\mathfrak p}) + \delta^Y_Z(y)
\geq d + 1
\quad\text{或}\quad
\text{depth}((\mathcal{F}_y^\wedge)_{\mathfrak p}) +
\dim(\mathcal{O}_{X, y}^\wedge/\mathfrak p) +
\delta^Y_Z(y) > d + 2
$$
考察上面的 (\ref{algebraization-equation-one}) 和 (\ref{algebraization-equation-two})，便得到所需结论，
唯一可能的例外是
$\dim(\mathcal{O}_{X, y}^\wedge/\mathfrak p) = 1$.
然而，若 $\dim(\mathcal{O}_{X, y}^\wedge/\mathfrak p) = 1$，则
$V(\mathfrak p) \cap V(I\mathcal{O}_{X, y}^\wedge) = \{\mathfrak m_y^\wedge\}$
，而且事实上
$$
\text{depth}((\mathcal{F}_y^\wedge)_{\mathfrak p}) + \delta^Y_Z(y) > d + 1
$$
，因为 $(\mathcal{F}_n)$ 满足 $(d + 1, d + 2)$-不等式；故仍得到结论。

\medskip\noindent
检验引理 \ref{algebraization-lemma-improvement-formal-coherent-module-better} 的条件 (c)。
假设
$V(\mathfrak p') \cap V(I\mathcal{O}_{X', y'}^\wedge) \not =
\{\mathfrak m_{y'}^\wedge\}$。则条件 (c) 等价于
$$
\text{depth}((\mathcal{F}_{y'}^\wedge)_{\mathfrak p'}) + \delta^{Y'}_{Z'}(y')
\geq 2
\quad\text{或}\quad
\text{depth}((\mathcal{F}_{y'}^\wedge)_{\mathfrak p'}) +
\dim(\mathcal{O}_{X', y'}^\wedge/\mathfrak p') +
\delta^{Y'}_{Z'}(y') > 3
$$
若 $(\mathcal{F}_n)$ 满足 (3)(b) 中的不等式，
应用 (\ref{algebraization-equation-two}) 可知所显示的两个不等式中的第二个成立。
若 $(\mathcal{F}_n)$ 满足 (3)(a)，即 $(d + 1, d + 2)$-不等式，
则由 (\ref{algebraization-equation-one}) 和 (\ref{algebraization-equation-two}) 立即得到结论。
这就证明了我们的断言。

\medskip\noindent
$(b|_{U'}^*\mathcal{F}_n) \to (\mathcal{F}_n'')$
以及 $(\mathcal{H}_n)$，后者在 $\textit{Coh}(U', I\mathcal{O}_{U'})$ 中；
像引理 \ref{algebraization-lemma-improvement-formal-coherent-module-better} 中那样选取它们。
对任意仿射开集 $W \subset X'$，注意到
$\delta^{W \cap Y'}_{W \cap Z'}(y') \geq \delta^{Y'}_{Z'}(y')$；
这是由引理 \ref{algebraization-lemma-discussion} 的 (6) 得出的。因此，
$(\mathcal{H}_n|_W)$ 满足引理 \ref{algebraization-lemma-cd-1-canonical} 的假设。
故 $(\mathcal{H}_n|_W)$ 典范地延拓到 $W$。
令 $(\mathcal{G}_{W, n})$ 为 $\textit{Coh}(W, I\mathcal{O}_W)$ 中
如引理 \ref{algebraization-lemma-canonically-algebraizable} 所述的典范延拓。
由引理 \ref{algebraization-lemma-canonically-extend-base-change}，可知对
$W' \subset W$ 存在唯一同构
$$
(\mathcal{G}_{W, n}|_{W'}) \longrightarrow
(\mathcal{G}_{W', n})
$$
，它与给定同构
$(\mathcal{G}_{W, n}|_{W \cap U}) \cong (\mathcal{H}_n|_{W \cap U})$.
相容。由此得到一个对象
$(\mathcal{G}_n)$，它属于 $\textit{Coh}(X', I\mathcal{O}_{X'})$，
其在 $U$ 上的限制同构于 $(\mathcal{H}_n)$。

\medskip\noindent
若 $A$ 是 $I$-进完备的，则可如下完成证明。
由 Grothendieck 存在性定理
（《概形上同调》，引理 \ref{coherent-lemma-existence-projective}），
可知 $(\mathcal{G}_n)$ 是某个凝聚 $\mathcal{O}_{X'}$-模的完备化。
再由《概形上同调》，引理 \ref{coherent-lemma-existence-easy}，
可知 $(b|_{U'}^*\mathcal{F}_n)$ 是某个凝聚
$\mathcal{O}_{U'}$-模 $\mathcal{F}'$ 的完备化。
由《概形上同调》，引理 \ref{coherent-lemma-inverse-systems-push-pull}，
存在映射
$$
(\mathcal{F}_n) \longrightarrow ((b|_{U'})_*\mathcal{F}')^\wedge
$$
，其核与余核被 $I$ 的某个幂零化。
最后应用引理 \ref{algebraization-lemma-map-kernel-cokernel-on-closed} 即得结论。

\medskip\noindent
若 $A$ 不完备，则在开始证明之前可以把 $A$ 换成其完备化；
参见引理 \ref{algebraization-lemma-algebraizable}。
完备化后各假设仍成立：条件 (3) 是立即的；条件 (1) 由
《对偶化复形》，引理 \ref{dualizing-lemma-ubiquity-dualizing} 得出；
条件 (2) 由《除子》，引理
\ref{divisors-lemma-flat-base-change-blowing-up} 得出。
因此，完备情形蕴含一般情形。
\end{proof}

\begin{proposition}[生成元较少的理想之代数化]
\label{algebraization-proposition-d-generators}
在情形 \ref{algebraization-situation-algebraize} 中，设 $(\mathcal{F}_n)$ 为
$\textit{Coh}(U, I\mathcal{O}_U)$ 的一个对象。假设
\begin{enumerate}
\item $A$ 有对偶化复形；
\item $V(I) = V(f_1, \ldots, f_d)$，其中某个 $d \geq 1$，并且
$f_1, \ldots, f_d \in A$；
\item 下列条件之一成立：
\begin{enumerate}
\item $(\mathcal{F}_n)$ 满足 $(d + 1, d + 2)$-不等式
（定义 \ref{algebraization-definition-s-d-inequalities}）；或者
\item 对 $y \in U \cap Y$ 以及素理想
$\mathfrak p \subset \mathcal{O}_{X, y}^\wedge$，若
$\mathfrak p \not \in V(I\mathcal{O}_{X, y}^\wedge)$
，则
$$
\text{depth}((\mathcal{F}^\wedge_y)_\mathfrak p) +
\dim(\mathcal{O}_{X, y}^\wedge/\mathfrak p) + \delta^Y_Z(y) > d + 2
$$
\end{enumerate}
\end{enumerate}
则 $(\mathcal{F}_n)$ 延拓到 $X$。特别地，若 $A$ 是
$I$-进完备的，则 $(\mathcal{F}_n)$ 是某个凝聚
$\mathcal{O}_U$-模的完备化。
\end{proposition}

\begin{proof}
可以假设 $I = (f_1, \ldots, f_d)$；参见《概形上同调》，引理
\ref{coherent-lemma-inverse-systems-ideals-equivalence}。
于是，$X$ 沿 $I$ 的吹起之所有纤维的维数至多为 $d - 1$。
故由引理 \ref{algebraization-lemma-blowup} 得到延拓。
最后一句由引理 \ref{algebraization-lemma-essential-image-completion} 得出。
\end{proof}

\noindent
请把下一个引理与注
\ref{algebraization-remark-interesting-case-variant}、
\ref{algebraization-remark-interesting-case},
\ref{algebraization-remark-interesting-case-bis} 和
\ref{algebraization-remark-interesting-case-ter} 比较。

\begin{lemma}
\label{algebraization-lemma-interesting-case-final}
在情形 \ref{algebraization-situation-algebraize} 中，设 $(\mathcal{F}_n)$ 为
$\textit{Coh}(U, I\mathcal{O}_U)$ 的一个对象。假设
\begin{enumerate}
\item $A$ 是有对偶化复形的局部环；
\item $X$ 的所有不可约分支具有相同维数；
\item 概形 $X \setminus Y$ 是 Cohen--Macaulay 的；
\item $I$ 由 $d$ 个元素生成；
\item $\dim(X) - \dim(Z) > d + 2$；并且
\item 对 $y \in U \cap Y$，模 $\mathcal{F}_y^\wedge$ 在
$V(I\mathcal{O}_{X, y}^\wedge)$ 外有限局部自由；例如，若
$\mathcal{F}_n$ 是有限局部自由
$\mathcal{O}_U/I^n\mathcal{O}_U$-模，则此条件成立。
\end{enumerate}
则 $(\mathcal{F}_n)$ 延拓到 $X$。特别地，若 $A$ 是 $I$-进
完备的，则 $(\mathcal{F}_n)$ 是某个凝聚 $\mathcal{O}_U$-模的完备化。
\end{lemma}

\begin{proof}
我们证明命题 \ref{algebraization-proposition-d-generators} 的假设 (1)、(2)、(3)(b)
均成立。其中 (1) 和 (2) 是显然的。

\medskip\noindent
设 $y \in U \cap Y$，并设 $\mathfrak p$ 为素理想
$\mathfrak p \subset \mathcal{O}_{X, y}^\wedge$，且
$\mathfrak p \not \in V(I\mathcal{O}_{X, y}^\wedge)$.
最后一个条件表明
$\text{depth}((\mathcal{F}_y^\wedge)_\mathfrak p) =
\text{depth}((\mathcal{O}_{X, y}^\wedge)_\mathfrak p)$.
由于 $X \setminus Y$ 是 Cohen--Macaulay 的，可知
$(\mathcal{O}_{X, y}^\wedge)_\mathfrak p$ 是 Cohen--Macaulay 的。
因此
\begin{align*}
& \text{depth}((\mathcal{F}^\wedge_y)_\mathfrak p) +
\dim(\mathcal{O}_{X, y}^\wedge/\mathfrak p) + \delta^Y_Z(y) \\
& =
\dim((\mathcal{O}_{X, y}^\wedge)_\mathfrak p) +
\dim(\mathcal{O}_{X, y}^\wedge/\mathfrak p) + \delta^Y_Z(y) \\
& =
\dim(\mathcal{O}_{X, y}^\wedge) + \delta^Y_Z(y)
\end{align*}
最后一个等号成立，是因为第二个条件表明
$\mathcal{O}_{X, y}$ 等维。
令 $\delta(y) = \dim(\overline{\{y\}})$。由于 $A$ 是链状局部环，
这是维数函数。由引理 \ref{algebraization-lemma-discussion}，有
$\delta^Y_Z(y) \geq \delta(y) - \dim(Z)$。由于 $X$ 等维，得到
$$
\dim(\mathcal{O}_{X, y}^\wedge) + \delta^Y_Z(y)
\geq \dim(\mathcal{O}_{X, y}^\wedge) + \delta(y) - \dim(Z)
= \dim(X) - \dim(Z)
$$
这就得到所需不等式，从而结论成立。
\end{proof}

\begin{remark}
\label{algebraization-remark-question}
我们无法证明或否证命题 \ref{algebraization-proposition-d-generators} 的如下类似命题：
把 $I$ 有 $d$ 个生成元这一假设替换为
$\text{cd}(A, I) \leq d$。
如果您知道证明或有反例，请发送电子邮件至
\href{mailto:stacks.project@gmail.com}{stacks.project@gmail.com}.
另一个显然的问题是，命题 \ref{algebraization-proposition-d-generators} 中的条件
在何种程度上是必要的。
\end{remark}




\section{凝聚形式模的代数化，VI}
\label{algebraization-section-algebraization-modules-yet-more}

\noindent
本节再给出几个较易证明的情形。

\begin{proposition}
\label{algebraization-proposition-algebraization-regular-sequence}
在情形 \ref{algebraization-situation-algebraize} 中，令
$(\mathcal{F}_n)$ 为 $\textit{Coh}(U, I\mathcal{O}_U)$ 的一个对象。
假设
\begin{enumerate}
\item 存在 $f_1, \ldots, f_d \in I$，使得对
$y \in U \cap Y$，理想 $I\mathcal{O}_{X, y}$
由 $f_1, \ldots, f_d$ 生成，并且
$f_1, \ldots, f_d$ 构成一个 $\mathcal{F}_y^\wedge$-正则序列；
\item $H^0(U, \mathcal{F}_1)$ 和 $H^1(U, \mathcal{F}_1)$
是有限 $A$-模。
\end{enumerate}
那么 $(\mathcal{F}_n)$ 可典范地延拓到 $X$。特别地，若 $A$
完备，则 $(\mathcal{F}_n)$ 是某个凝聚
$\mathcal{O}_U$-模的完备化。
\end{proposition}

\begin{proof}
我们通过验证引理 \ref{algebraization-lemma-when-done} 的假设 (a)、(b) 和 (c)
来证明这一点。对每个 $n$，都有短正合列
$$
0 \to I^n\mathcal{F}_{n + 1} \to \mathcal{F}_{n + 1} \to \mathcal{F}_n \to 0
$$
由于 $f_1, \ldots, f_d$ 在每个“茎”
$\mathcal{F}_y^\wedge$ 上构成正则序列（因而是拟正则序列，见
代数，引理 \ref{algebra-lemma-regular-quasi-regular}），并且我们有
$I\mathcal{F}_n = (f_1, \ldots, f_d)\mathcal{F}_n$ 对所有 $n$ 成立，
逐茎检验便得到
$$
I^n\mathcal{F}_{n + 1} =
\bigoplus\nolimits_{e_1 + \ldots + e_d = n} \mathcal{F}_1 \cdot
f_1^{e_1} \ldots f_d^{e_d}
$$
利用 $H^0(U, \mathcal{F}_1)$ 的有限性假设并作归纳，我们先得出
$M_n = H^0(U, \mathcal{F}_n)$ 是有限 $A$-模，且这对所有 $n$ 成立。
由此可见引理 \ref{algebraization-lemma-when-done} 的条件 (c) 成立。
我们还看到
$$
\bigoplus\nolimits_{n \geq 0} H^1(U, I^n\mathcal{F}_{n + 1})
$$
是有限分次 $R = \bigoplus I^n/I^{n +1}$-模。
由上同调，引理 \ref{cohomology-lemma-ML-general}，
可知引理 \ref{algebraization-lemma-when-done} 的条件 (a) 成立。最后，引理
\ref{algebraization-lemma-when-done} 的条件 (b) 也成立，因为
$\bigoplus H^0(U, I^n\mathcal{F}_{n + 1})$ 是有限分次 $R$-模，
而且可以应用上同调，引理
\ref{cohomology-lemma-topology-I-adic-general}。
\end{proof}

\begin{remark}
\label{algebraization-remark-interesting-case-ter}
在命题 \ref{algebraization-proposition-algebraization-regular-sequence} 的情形中，
若假设 $A$ 有对偶化复形，则
$H^0(U, \mathcal{F}_1)$ 和
$H^1(U, \mathcal{F}_1)$ 有限这一条件等价于
$$
\text{depth}(\mathcal{F}_{1, y}) +
\dim(\mathcal{O}_{\overline{\{y\}}, z}) > 2
$$
对所有 $y \in U \cap Y$ 以及 $z \in Z \cap \overline{\{y\}}$ 均成立。
见局部上同调，引理 \ref{local-cohomology-lemma-finiteness-Rjstar}。
例如，若 $\mathcal{F}_1$ 是有限局部自由
$\mathcal{O}_{U \cap Y}$-模，$Y$ 满足 $(S_2)$，并且
$\text{codim}(Z', Y') \geq 3$ 对如下每一对不可约分支都成立：
$Y'$ 是 $Y$ 的不可约分支，$Z'$ 是 $Z$ 的不可约分支，且
$Z' \subset Y'$，则上述条件成立。
\end{remark}

\begin{proposition}
\label{algebraization-proposition-algebraization-flat}
在情形 \ref{algebraization-situation-algebraize} 中，令
$(\mathcal{F}_n)$ 为 $\textit{Coh}(U, I\mathcal{O}_U)$ 的一个对象。
假设存在一个 Noether 局部环 $(R, \mathfrak m)$ 以及环同态
$R \to A$，使得
\begin{enumerate}
\item $I = \mathfrak m A$,
\item 对 $y \in U \cap Y$，茎 $\mathcal{F}_y^\wedge$ 是 $R$-平坦的；
\item $H^0(U, \mathcal{F}_1)$ 和 $H^1(U, \mathcal{F}_1)$ 是有限
$A$-模。
\end{enumerate}
那么 $(\mathcal{F}_n)$ 可典范地延拓到 $X$。特别地，若 $A$
完备，则 $(\mathcal{F}_n)$ 是某个凝聚
$\mathcal{O}_U$-模的完备化。
\end{proposition}

\begin{proof}
证明与命题 \ref{algebraization-proposition-algebraization-regular-sequence} 的证明
完全相同。具体而言，若 $\kappa = R/\mathfrak m$，则对 $n \geq 0$
有同构
$$
I^n \mathcal{F}_{n + 1} \cong
\mathcal{F}_1 \otimes_\kappa \mathfrak m^n/\mathfrak m^{n + 1}
$$
而右端是 $\mathcal{F}_1$ 的有限个副本的直和。这一点可以逐茎检验。
其余部分完全相同。
\end{proof}

\begin{remark}
\label{algebraization-remark-interesting-case-quater}
命题 \ref{algebraization-proposition-algebraization-flat} 是
\cite[Theorem 2.10 (i)]{Baranovsky} 的一个局部版本。从局部结果推出
整体结果很直接；下面简述论证。设 $(R, \mathfrak m)$ 是完备
Noether 局部环，且 $X \to \Spec(R)$ 是固有态射。对 $n \geq 1$，令
$X_n = X \times_{\Spec(R)} \Spec(R/\mathfrak m^n)$。
令 $Z \subset X_1$ 为特殊纤维的一个闭子集。置
$U = X \setminus Z$，并以 $j : U \to X$ 表示嵌入态射。给定一个对象
$$
(\mathcal{F}_n) \text{ 属于 } \textit{Coh}(U, \mathfrak m\mathcal{O}_U)
$$
它在 $R$ 上平坦，其含义是 $\mathcal{F}_n$ 在
$R/\mathfrak m^n$ 上平坦，且这对所有 $n$ 成立。假设
$j_*\mathcal{F}_1$ 和
$R^1j_*\mathcal{F}_1$ 是凝聚模。那么，由命题
\ref{algebraization-proposition-algebraization-flat}，仿射局部地在 $X$ 上可得到
$(\mathcal{F}_n)$ 的典范延拓；而该延拓的构造与局部化可交换
（由引理 \ref{algebraization-lemma-algebraization-principal-variant}）。
于是得到一个典范的整体对象 $(\mathcal{G}_n)$，它属于
$\textit{Coh}(X, \mathfrak m\mathcal{O}_X)$
且其在 $U$ 上的限制是 $(\mathcal{F}_n)$。由 Grothendieck
存在性定理（概形上同调，命题
\ref{coherent-proposition-existence-proper}），可知存在凝聚
$\mathcal{O}_X$-模 $\mathcal{G}$，其完备化为 $(\mathcal{G}_n)$。
由此可见 $(\mathcal{F}_n)$ 可代数化，也就是说，它是某个凝聚
$\mathcal{O}_U$-模的完备化。

\medskip\noindent
再补充一点：$j_*\mathcal{F}_1$ 和 $R^1j_*\mathcal{F}_1$
的凝聚性是特殊纤维上的条件。具体而言，若以
$j_1 : U_1 \to X_1$ 表示 $j : U \to X$ 的特殊纤维，则可以把
$\mathcal{F}_1$ 视为 $U_1$ 上的凝聚层，并且有
$j_*\mathcal{F}_1 = j_{1, *}\mathcal{F}_1$，且
$R^1j_*\mathcal{F}_1 = R^1j_{1, *}\mathcal{F}_1$。
因此，例如，若 $X_1$ 满足 $(S_2)$ 且不可约，并且
$\dim(X_1) - \dim(Z) \geq 3$，而 $\mathcal{F}_1$ 是局部自由
$\mathcal{O}_{U_1}$-模，则 $j_{1, *}\mathcal{F}_1$ 和
$R^1j_{1, *}\mathcal{F}_1$ 都是凝聚模。
\end{remark}








\section{完备化函子的应用}
\label{algebraization-section-completion-application}

\noindent
本节只是把若干已得结果组合起来，以便引用。这里还可以陈述许多
（更强的）结果。

\begin{lemma}
\label{algebraization-lemma-equivalence-better}
在情形 \ref{algebraization-situation-algebraize} 中，假设
\begin{enumerate}
\item $A$ 有对偶化复形，并且是 $I$-进完备的；
\item $I = (f)$ 由一个元素生成；
\item $A$ 是局部环，其极大理想为 $\mathfrak a = \mathfrak m$；
\item 下列条件之一成立：
\begin{enumerate}
\item $A_f$ 满足 $(S_2)$，并且对满足 $\mathfrak p \subset A$、
$f \not \in \mathfrak p$ 的极小素理想，有
$\dim(A/\mathfrak p) \geq 4$；或者
\item 若 $\mathfrak p \not \in V(f)$ 且
$V(\mathfrak p) \cap V(f) \not = \{\mathfrak m\}$，则
$\text{depth}(A_\mathfrak p) + \dim(A/\mathfrak p) > 3$。
\end{enumerate}
\end{enumerate}
那么，置 $U_0 = U \cap V(f)$，完备化函子
$$
\colim_{U_0 \subset U' \subset U\text{ 开}}
\textit{Coh}(\mathcal{O}_{U'})
\longrightarrow
\textit{Coh}(U, f\mathcal{O}_U)
$$
在有限局部自由对象所成的满子范畴上是等价。
\end{lemma}

\begin{proof}
由引理 \ref{algebraization-lemma-fully-faithful-general} 可知该函子充分忠实
（略去细节）。下面证明本质满。令 $(\mathcal{F}_n)$ 为
$\textit{Coh}(U, f\mathcal{O}_U)$ 中的有限局部自由对象。由引理
\ref{algebraization-lemma-algebraization-principal-bis} 或命题
\ref{algebraization-proposition-cd-1}，存在凝聚 $\mathcal{O}_U$-模 $\mathcal{F}$，
使得 $(\mathcal{F}_n)$ 是 $\mathcal{F}$ 的完备化。事实上，要应用
任一结果，唯一需要检验的是 $(\mathcal{F}_n)$ 满足 $(2, 3)$-不等式。
这由引理 \ref{algebraization-lemma-unwinding-conditions-bis} 给出。若 $y \in U_0$，
则 $f$-进完备化后的茎 $\mathcal{F}_y$ 同构于某个有限自由模，该模
定义在 $f$-进完备化后的 $\mathcal{O}_{U, y}$ 上。因此
$\mathcal{F}$ 在某个开邻域 $U'$ 上有限局部自由，而该邻域包含
$U_0$。证毕。
\end{proof}

\begin{lemma}
\label{algebraization-lemma-equivalence}
在情形 \ref{algebraization-situation-algebraize} 中，假设
\begin{enumerate}
\item $I = (f)$ 是主理想；
\item $A$ 是 $f$-进完备的；
\item $f$ 是非零因子；
\item $H^1_\mathfrak a(A/fA)$ 和 $H^2_\mathfrak a(A/fA)$
是有限 $A$-模。
\end{enumerate}
那么，置 $U_0 = U \cap V(f)$，完备化函子
$$
\colim_{U_0 \subset U' \subset U\text{ 开}}
\textit{Coh}(\mathcal{O}_{U'})
\longrightarrow
\textit{Coh}(U, f\mathcal{O}_U)
$$
在有限局部自由对象所成的满子范畴上是等价。
\end{lemma}

\begin{proof}
由引理 \ref{algebraization-lemma-fully-faithful-general-alternative}，该函子充分忠实。
本质满性由引理 \ref{algebraization-lemma-algebraization-principal-variant} 得到。
\end{proof}








\section{凝聚三元组}
\label{algebraization-section-coherent-triples}

\noindent
令 $(A, \mathfrak m)$ 为 Noether 局部环，令
$f \in \mathfrak m$ 为非零因子。置
$X = \Spec(A)$、$X_0 = \Spec(A/fA)$、$U = X \setminus V(\mathfrak m)$，以及
$U_0 = U \cap X_0$。若满足下列条件，就称
$(\mathcal{F}, \mathcal{F}_0, \alpha)$ 为一个{\it 凝聚三元组}：
\begin{enumerate}
\item $\mathcal{F}$ 是凝聚 $\mathcal{O}_U$-模，且
$f : \mathcal{F} \to \mathcal{F}$ 是单射；
\item $\mathcal{F}_0$ 是凝聚 $\mathcal{O}_{X_0}$-模；
\item $\alpha : \mathcal{F}/f\mathcal{F} \to \mathcal{F}_0|_{U_0}$
是同构。
\end{enumerate}
凝聚三元组之间有显然的{\it 态射}概念，它使所有凝聚三元组组成一个范畴。

\medskip\noindent
凝聚三元组范畴是加性范畴，但不是 Abel 范畴。不过，凝聚三元组的
短正合列应当是什么意思是清楚的。

\medskip\noindent
给定两个凝聚三元组 $(\mathcal{F}, \mathcal{F}_0, \alpha)$ 和
$(\mathcal{G}, \mathcal{G}_0, \beta)$，则
$(\mathcal{F} \otimes_{\mathcal{O}_U} \mathcal{G},
\mathcal{F}_0  \otimes_{\mathcal{O}_{X_0}} \mathcal{G}_0,
\alpha \otimes \beta)$ 未必是凝聚三元组\footnote{也就是说，$f$
在 $\mathcal{F} \otimes_{\mathcal{O}_U} \mathcal{G}$ 上未必是单射。}。
不过，若茎 $\mathcal{G}_x$ 对所有 $x \in U_0$ 都自由，则它确实是。

\medskip\noindent
称凝聚三元组 $(\mathcal{G}, \mathcal{G}_0, \beta)$ 为
{\it 局部自由的}（相应地，{\it 可逆的}），如果
$\mathcal{G}$ 和 $\mathcal{G}_0$ 分别是局部自由模（相应地，可逆模）。
在这种情形下，与
$(\mathcal{G}, \mathcal{G}_0, \beta)$ 作张量积是有意义的（见上文），
并且它把凝聚三元组的短正合列变为凝聚三元组的短正合列。

\begin{lemma}
\label{algebraization-lemma-prepare-chi-triple}
对任意凝聚三元组 $(\mathcal{F}, \mathcal{F}_0, \alpha)$，存在凝聚
$\mathcal{O}_X$-模 $\mathcal{F}'$，使得
$f : \mathcal{F}' \to \mathcal{F}'$ 是单射；还存在同构
$\alpha' : \mathcal{F}'|_U \to \mathcal{F}$ 以及映射
$\alpha'_0 : \mathcal{F}'/f\mathcal{F}' \to \mathcal{F}_0$
使得 $\alpha \circ (\alpha' \bmod f) = \alpha'_0|_{U_0}$。
\end{lemma}

\begin{proof}
选取有限 $A$-模 $M$。要求 $\mathcal{F}$ 是某个层在 $U$ 上的限制；
这个层是凝聚 $\mathcal{O}_X$-模，即由 $M$ 给出的相伴层，见
局部上同调，引理
\ref{local-cohomology-lemma-finiteness-pushforwards-and-H1-local}。
由于 $\mathcal{F}$ 无 $f$-挠，可以用 $M$ 对其
$f$-幂挠子模的商来代替它。另一方面，令
$M_0 = \Gamma(X_0, \mathcal{F}_0)$，于是 $\mathcal{F}_0$ 是凝聚
$\mathcal{O}_{X_0}$-模，并且它与有限 $A/fA$-模 $M_0$ 相伴。
由概形上同调，引理 \ref{coherent-lemma-homs-over-open}，存在 $n$，
使同构 $\alpha_0$ 对应于一个 $A/fA$-模同态
$\mathfrak m^n M/fM \to M_0$（其核与余核均被 $\mathfrak m$ 的某个幂
零化，但这里不需要这一点）。因此，若取 $M' = \mathfrak m^n M$，
并令 $\mathcal{F}'$ 为凝聚 $\mathcal{O}_X$-模，即与 $M'$ 相伴的模，
则引理显然成立。
\end{proof}

\noindent
令 $(\mathcal{F}, \mathcal{F}_0, \alpha)$ 为凝聚三元组。按照
引理 \ref{algebraization-lemma-prepare-chi-triple} 选取
$\mathcal{F}', \alpha', \alpha'_0$。置
\begin{equation}
\label{algebraization-equation-chi-triple}
\chi(\mathcal{F}, \mathcal{F}_0, \alpha) =
\text{length}_A(\Coker(\alpha'_0)) - 
\text{length}_A(\Ker(\alpha'_0))
\end{equation}
右端表达式有意义，因为 $\alpha'_0$ 在 $U_0$ 上是同构，故其核与余核
都是支撑在 $\{\mathfrak m\}$ 上的凝聚模，从而具有有限长度
（代数，引理 \ref{algebra-lemma-support-point}）。

\begin{lemma}
\label{algebraization-lemma-well-defined-chi-triple}
式 (\ref{algebraization-equation-chi-triple}) 中的量
$\chi(\mathcal{F}, \mathcal{F}_0, \alpha)$ 与按照引理
\ref{algebraization-lemma-prepare-chi-triple} 对
$\mathcal{F}', \alpha', \alpha'_0$ 的选择无关。
\end{lemma}

\begin{proof}
令 $\mathcal{F}', \alpha', \alpha'_0$ 和
$\mathcal{F}'', \alpha'', \alpha''_0$ 为两组这样的选择。对 $n > 0$，置
$\mathcal{F}'_n = \mathfrak m^n \mathcal{F}'$。由概形上同调，引理
\ref{coherent-lemma-homs-over-open}，存在某个 $n$ 以及
$\mathcal{O}_X$-模映射 $\mathcal{F}'_n \to \mathcal{F}''$，它与等同
$\mathcal{F}''|_U = \mathcal{F}'|_U$ 相符；这个等同由 $\alpha'$ 和
$\alpha''$ 确定。于是图
$$
\xymatrix{
\mathcal{F}'_n/f\mathcal{F}'_n \ar[r] \ar[d] &
\mathcal{F}'/f\mathcal{F}' \ar[d]^{\alpha_0'} \\
\mathcal{F}''/f\mathcal{F}'' \ar[r]^{\alpha_0''} &
\mathcal{F}_0
}
$$
限制到 $U_0$ 后交换。因此，由概形上同调，引理
\ref{coherent-lemma-homs-over-open}，该图限制到
$\mathfrak m^l(\mathcal{F}'_n/f\mathcal{F}'_n)$ 后交换，其中
$l > 0$ 是某个整数。由于
$\mathcal{F}'_{n + l}/f\mathcal{F}'_{n + l} \to \mathcal{F}'_n/f\mathcal{F}'_n$
通过 $\mathfrak m^l(\mathcal{F}'_n/f\mathcal{F}'_n)$ 分解，可知把
$n$ 替换为 $n + l$ 后该图交换。换言之，我们找到了第三组选择
$\mathcal{F}''', \alpha''', \alpha'''_0$
以及映射 $\mathcal{F}''' \to \mathcal{F}''$ 和
$\mathcal{F}''' \to \mathcal{F}'$；它们定义在 $X$ 上，并与 $U$ 和
$X_0$ 上的映射相容。
于是问题化归为下一段讨论的情形。

\medskip\noindent
假设有映射 $\mathcal{F}'' \to \mathcal{F}'$，它定义在 $X$ 上，并分别
与 $\alpha', \alpha''$ 在 $U$ 上以及与 $\alpha'_0, \alpha''_0$
在 $X_0$ 上相容。注意 $\mathcal{F}'' \to \mathcal{F}'$ 是单射，
因为它在 $U$ 上是同构，而且
$f : \mathcal{F}'' \to \mathcal{F}''$ 是单射。显然
$\mathcal{F}'/\mathcal{F}''$ 支撑在 $\{\mathfrak m\}$ 上，因而长度有限。
我们有如下凝聚 $\mathcal{O}_{X_0}$-模的映射：
$$
\mathcal{F}''/f\mathcal{F}'' \to
\mathcal{F}'/f\mathcal{F}' \xrightarrow{\alpha'_0}
\mathcal{F}_0
$$
其复合为 $\alpha''_0$，并且这些映射在 $U_0$ 上都是同构。初等同调代数
给出 $6$-项正合列
$$
\begin{matrix}
0 \to
\Ker(\mathcal{F}''/f\mathcal{F}'' \to \mathcal{F}'/f\mathcal{F}') \to
\Ker(\alpha''_0) \to
\Ker(\alpha'_0) \to \\
\Coker(\mathcal{F}''/f\mathcal{F}'' \to \mathcal{F}'/f\mathcal{F}') \to
\Coker(\alpha''_0) \to
\Coker(\alpha'_0) \to 0
\end{matrix}
$$
由长度的可加性（代数，引理 \ref{algebra-lemma-length-additive}），
只需证明
$$
\text{length}_A(
\Coker(\mathcal{F}''/f\mathcal{F}'' \to \mathcal{F}'/f\mathcal{F}')) -
\text{length}_A(
\Ker(\mathcal{F}''/f\mathcal{F}'' \to \mathcal{F}'/f\mathcal{F}')) = 0
$$
对下图应用蛇形引理即可得到这一点：
$$
\xymatrix{
0 \ar[r] &
\mathcal{F}'' \ar[r]_f \ar[d] &
\mathcal{F}'' \ar[r] \ar[d] &
\mathcal{F}''/f\mathcal{F}'' \ar[r] \ar[d] &
0 \\
0 \ar[r] &
\mathcal{F}' \ar[r]^f &
\mathcal{F}' \ar[r] &
\mathcal{F}'/f\mathcal{F}' \ar[r] &
0
}
$$
这里还使用了 $\mathcal{F}'/\mathcal{F}''$ 长度有限这一事实。
\end{proof}

\begin{lemma}
\label{algebraization-lemma-ses-chi-triple}
下式成立：
$\chi(\mathcal{G}, \mathcal{G}_0, \beta) =
\chi(\mathcal{F}, \mathcal{F}_0, \alpha) +
\chi(\mathcal{H}, \mathcal{H}_0, \gamma)$，其条件是
$$
0 \to
(\mathcal{F}, \mathcal{F}_0, \alpha) \to
(\mathcal{G}, \mathcal{G}_0, \beta) \to
(\mathcal{H}, \mathcal{H}_0, \gamma)
\to 0
$$
是凝聚三元组的短正合列。
\end{lemma}

\begin{proof}
选取 $\mathcal{G}', \beta', \beta'_0$，如引理
\ref{algebraization-lemma-prepare-chi-triple} 所述；这是针对三元组
$(\mathcal{G}, \mathcal{G}_0, \beta)$ 的选择。以
$j : U \to X$ 表示嵌入态射。令 $\mathcal{F}' \subset \mathcal{G}'$
为如下复合的核：
$$
\mathcal{G}' \xrightarrow{\beta'} j_*\mathcal{G} \to j_*\mathcal{H}
$$
注意 $\mathcal{H}' = \mathcal{G}'/\mathcal{F}'$ 是
$j_*\mathcal{H}$ 的凝聚子层，因而
$f : \mathcal{H}' \to \mathcal{H}'$ 是单射。因此，由蛇形引理得到短正合列
$$
0 \to \mathcal{F}'/f\mathcal{F}' \to
\mathcal{G}'/f\mathcal{G}' \to
\mathcal{H}'/f\mathcal{H}' \to 0
$$
由构造可得同构
$\alpha' : \mathcal{F}'|_U \to \mathcal{F}$,
$\beta' : \mathcal{G}'|_U \to \mathcal{G}$，以及
$\gamma' : \mathcal{H}'|_U \to \mathcal{H}$。为完成证明，需要构造映射
$\alpha'_0 : \mathcal{F}'/f\mathcal{F}' \to \mathcal{F}_0$ 和
$\gamma'_0 : \mathcal{H}'/f\mathcal{H}' \to \mathcal{H}_0$，使其如
引理 \ref{algebraization-lemma-prepare-chi-triple} 所述，并嵌入下列交换图：
$$
\xymatrix{
0 \ar[r] &
\mathcal{F}'/f\mathcal{F}' \ar[r] \ar@{..>}[d]^{\alpha'_0} &
\mathcal{G}'/f\mathcal{G}' \ar[r] \ar[d]^{\beta'_0} &
\mathcal{H}'/f\mathcal{H}' \ar[r] \ar@{..>}[d]^{\gamma'_0} &
0 \\
0 \ar[r] &
\mathcal{F}_0 \ar[r] &
\mathcal{G}_0 \ar[r] &
\mathcal{H}_0 \ar[r] &
0
}
$$
不过，对 $\mathcal{G}'$ 的初始选择，未必能做到这一点。从上图可见，
障碍恰好是复合
$$
\delta :
\mathcal{F}'/f\mathcal{F}' \to
\mathcal{G}'/f\mathcal{G}' \xrightarrow{\beta'_0}
\mathcal{G}_0 \to
\mathcal{H}_0
$$
注意，$\delta$ 限制到 $U_0$ 后为零，这是由我们对 $\mathcal{F}'$ 和
$\mathcal{H}'$ 的选择所致。因此，由概形上同调，引理
\ref{coherent-lemma-homs-over-open}，存在 $k > 0$，使得
$\delta$ 在 $\mathfrak m^k \cdot (\mathcal{F}'/f\mathcal{F}')$ 上为零。
对 $n > k$，置 $\mathcal{G}'_n = \mathfrak m^n \mathcal{G}'$、
$\mathcal{F}'_n = \mathcal{G}'_n \cap \mathcal{F}'$，以及
$\mathcal{H}'_n = \mathcal{G}'_n/\mathcal{F}'_n$。注意，$\beta'_0$
可与
$\mathcal{G}'_n/f\mathcal{G}'_n \to \mathcal{G}'/f\mathcal{G}'$
复合，从而给出映射
$\beta'_{n, 0} : \mathcal{G}'_n/f\mathcal{G}'_n \to \mathcal{G}_0$
，如引理 \ref{algebraization-lemma-prepare-chi-triple} 所述。由 Artin--Rees
（代数，引理 \ref{algebra-lemma-Artin-Rees}），可以选取 $n$，使得
$\mathcal{F}'_n \subset \mathfrak m^k \mathcal{F}'$。如上，映射
$f : \mathcal{F}'_n \to \mathcal{F}'_n$,
$f : \mathcal{G}'_n \to \mathcal{G}'_n$，以及
$f : \mathcal{H}'_n \to \mathcal{H}'_n$ 都是单射；再次使用蛇形引理，
得到短正合列
$$
0 \to \mathcal{F}'_n/f\mathcal{F}'_n \to
\mathcal{G}'_n/f\mathcal{G}'_n \to
\mathcal{H}'_n/f\mathcal{H}'_n \to 0
$$
如上，有同构
$\alpha'_n : \mathcal{F}'_n|_U \to \mathcal{F}$,
$\beta'_n : \mathcal{G}'_n|_U \to \mathcal{G}$，以及
$\gamma'_n : \mathcal{H}'_n|_U \to \mathcal{H}$。与前面一样，考虑障碍
$$
\delta_n :
\mathcal{F}'_n/f\mathcal{F}'_n \to
\mathcal{G}'_n/f\mathcal{G}'_n
\xrightarrow{\beta'_{n, 0}}
\mathcal{G}_0 \to
\mathcal{H}_0
$$
不过，交换图
$$
\xymatrix{
\mathcal{F}'_n/f\mathcal{F}'_n \ar[r] \ar[d] &
\mathcal{G}'_n/f\mathcal{G}'_n \ar[r]_{\beta'_{n, 0}} \ar[d] &
\mathcal{G}_0 \ar[r] \ar[d] &
\mathcal{H}_0 \ar[d] \\
\mathcal{F}'/f\mathcal{F}' \ar[r] &
\mathcal{G}'/f\mathcal{G}' \ar[r]^{\beta'_0} &
\mathcal{G}_0 \ar[r] &
\mathcal{H}_0
}
$$
结合对 $n$ 的选择和关于 $\delta$ 的观察，表明 $\delta_n = 0$。
这就产生了所需映射
$\alpha'_{n, 0} : \mathcal{F}'_n/f\mathcal{F}'_n \to \mathcal{F}_0$ 和
$\gamma'_{n, 0} : \mathcal{H}'_n/f\mathcal{H}'_n \to \mathcal{H}_0$.
因此，可以使用
$\mathcal{F}'_n, \alpha'_n, \alpha'_{n, 0}$,
$\mathcal{G}'_n, \beta'_n, \beta'_{n, 0}$，以及
$\mathcal{H}'_n, \gamma'_n, \gamma'_{n, 0}$
来计算
$\chi(\mathcal{F}, \mathcal{F}_0, \alpha)$,
$\chi(\mathcal{G}, \mathcal{G}_0, \beta)$，以及
$\chi(\mathcal{H}, \mathcal{H}_0, \gamma)$。最后，对下图应用蛇形引理
$$
\xymatrix{
0 \ar[r] &
\mathcal{F}'_n/f\mathcal{F}'_n \ar[r] \ar[d] &
\mathcal{G}'_n/f\mathcal{G}'_n \ar[r] \ar[d] &
\mathcal{H}'_n/f\mathcal{H}'_n \ar[r] \ar[d] &
0 \\
0 \ar[r] &
\mathcal{F}_0 \ar[r] &
\mathcal{G}_0 \ar[r] &
\mathcal{H}_0 \ar[r] &
0
}
$$
并使用长度的可加性（代数，引理
\ref{algebra-lemma-length-additive}），即可得到本引理。
\end{proof}

\begin{proposition}
\label{algebraization-proposition-hilbert-triple}
令 $(\mathcal{F}, \mathcal{F}_0, \alpha)$ 为凝聚三元组，并令
$(\mathcal{L}, \mathcal{L}_0, \lambda)$ 为可逆凝聚三元组。那么函数
$$
\mathbf{Z} \longrightarrow \mathbf{Z},\quad
n \longmapsto
\chi((\mathcal{F}, \mathcal{F}_0, \alpha) \otimes
(\mathcal{L}, \mathcal{L}_0, \lambda)^{\otimes n})
$$
是次数 $\leq \dim(\text{Supp}(\mathcal{F}))$ 的多项式。
\end{proposition}

\noindent
更确切地说，若 $\mathcal{F} = 0$，则该函数为常数。若
$\mathcal{F}$ 在 $U$ 中的支撑有限，则该函数为常数。若
$\mathcal{F}$ 在 $U$ 中的支撑维数为 $1$，即 $\mathcal{F}$ 的支撑
在 $X$ 中的闭包维数为 $2$，则该函数为线性函数，依此类推。

\begin{proof}
我们对 $\mathcal{F}$ 的支撑维数作归纳来证明。

\medskip\noindent
基本情形是 $\mathcal{F} = 0$。此时 $\mathcal{F}_0$ 或者为零，
或者其支撑为 $\{\mathfrak m\}$。在这种情形下，有
$$
(\mathcal{F}, \mathcal{F}_0, \alpha) \otimes
(\mathcal{L}, \mathcal{L}_0, \lambda)^{\otimes n} =
(0, \mathcal{F}_0 \otimes \mathcal{L}_0^{\otimes n}, 0) \cong
(0, \mathcal{F}_0, 0)
$$
因此命题中的函数为常数，其值等于 $\mathcal{F}_0$ 的长度。

\medskip\noindent
归纳步骤。假设 $\mathcal{F}$ 的支撑非空。以
$\mathcal{G}_0 \subset \mathcal{F}_0$ 表示由支撑在
$\{\mathfrak m\}$ 上的截面组成的子模。于是得到短正合列
$$
0 \to (0, \mathcal{G}_0, 0) \to
(\mathcal{F}, \mathcal{F}_0, \alpha) \to
(\mathcal{F}, \mathcal{F}_0/\mathcal{G}_0, \alpha) \to 0
$$
与可逆凝聚三元组 $(\mathcal{L}, \mathcal{L}_0, \lambda)$ 作张量积后，
此列仍然正合，见上文讨论。因此，由 $\chi$ 的可加性
（引理 \ref{algebraization-lemma-ses-chi-triple}）以及上述基本情形，只需对
$(\mathcal{F}, \mathcal{F}_0/\mathcal{G}_0, \alpha)$ 证明归纳步骤。
由此可假设 $\mathfrak m$ 不是 $\mathcal{F}_0$ 的相伴点。

\medskip\noindent
令 $T = \text{Ass}(\mathcal{F}) \cup \text{Ass}(\mathcal{F}/f\mathcal{F})$。
由于 $U$ 拟仿射，可以找到 $s \in \Gamma(U, \mathcal{L})$，它在每个
$u \in T$ 处均不消失，见性质，引理
\ref{properties-lemma-quasi-affine-invertible-nonvanishing-section}。
将 $s$ 乘以 $\mathfrak m$ 的适当元素后，可以假设
$\lambda(s \bmod f) = s_0|_{U_0}$，其中
$s_0 \in \Gamma(X_0, \mathcal{L}_0)$；略去细节。于是得到态射
$$
(s, s_0) :
(\mathcal{O}_U, \mathcal{O}_{X_0}, 1)
\longrightarrow
(\mathcal{L}, \mathcal{L}_0, \lambda)
$$
它属于凝聚三元组范畴。令
$\mathcal{G} = \Coker(s : \mathcal{F} \to \mathcal{F} \otimes \mathcal{L})$
并令
$\mathcal{G}_0 = \Coker(s_0 : \mathcal{F}_0 \to
\mathcal{F}_0 \otimes \mathcal{L}_0)$。注意，
$s_0 : \mathcal{F}_0 \to
\mathcal{F}_0 \otimes \mathcal{L}_0$ 是单射：它在 $U_0$ 上单射，
这是由 $s$ 的选择得到的；而且 $\mathfrak m$ 不是
$\mathcal{F}_0$ 的相伴点。因此存在同构
$\beta : \mathcal{G}/f\mathcal{G} \to \mathcal{G}_0|_{U_0}$，使得得到
短正合列
$$
0 \to
(\mathcal{F}, \mathcal{F}_0, \alpha) \to
(\mathcal{F}, \mathcal{F}_0, \alpha) \otimes
(\mathcal{L}, \mathcal{L}_0, \lambda) \to
(\mathcal{G}, \mathcal{G}_0, \beta) \to 0
$$
由支撑维数的归纳假设，命题对凝聚三元组
$(\mathcal{G}, \mathcal{G}_0, \beta)$ 成立。利用引理
\ref{algebraization-lemma-ses-chi-triple} 的可加性，可知
$$
n \longmapsto 
\chi((\mathcal{F}, \mathcal{F}_0, \alpha) \otimes
(\mathcal{L}, \mathcal{L}_0, \lambda)^{\otimes n + 1})
-
\chi((\mathcal{F}, \mathcal{F}_0, \alpha) \otimes
(\mathcal{L}, \mathcal{L}_0, \lambda)^{\otimes n})
$$
是多项式。对定义在全体整数上的函数应用代数，引理
\ref{algebra-lemma-numerical-polynomial} 的一个变体，即得结论
（略去细节）。
\end{proof}

\begin{lemma}
\label{algebraization-lemma-nonnegative-chi-triple}
假设 $\text{depth}(A) \geq 3$，等价地，
$\text{depth}(A/fA) \geq 2$。令 $(\mathcal{L}, \mathcal{L}_0, \lambda)$
为可逆凝聚三元组。那么
$$
\chi(\mathcal{L}, \mathcal{L}_0, \lambda) =
\text{length}_A \Coker(\Gamma(U, \mathcal{L}) \to \Gamma(U_0, \mathcal{L}_0))
$$
特别地，该数 $\geq 0$。此外，
$\chi(\mathcal{L}, \mathcal{L}_0, \lambda) = 0$ 当且仅当
$\mathcal{L} \cong \mathcal{O}_U$。
\end{lemma}

\begin{proof}
深度条件的等价性由代数，引理
\ref{algebra-lemma-depth-drops-by-one} 得到。由深度条件可知
$\Gamma(U, \mathcal{O}_U) = A$ 且
$\Gamma(U_0, \mathcal{O}_{U_0}) = A/fA$，见对偶化复形，引理
\ref{dualizing-lemma-depth} 以及局部上同调，引理
\ref{local-cohomology-lemma-finiteness-pushforwards-and-H1-local}。
利用局部上同调，引理
\ref{local-cohomology-lemma-finiteness-for-finite-locally-free}，可知
$M = \Gamma(U, \mathcal{L})$ 是有限 $A$-模。进而有
$\text{depth}(M) \geq 2$；例如，可用局部上同调，引理
\ref{local-cohomology-lemma-finiteness-pushforwards-and-H1-local} 的第 (4) 部分，
或除子，引理 \ref{divisors-lemma-depth-pushforward}。此外，
$\mathcal{L}_0 \cong \mathcal{O}_{X_0}$，因为 $X_0$ 是局部概形。因此还有
$M_0 = \Gamma(X_0, \mathcal{L}_0) = \Gamma(U_0, \mathcal{L}_0|_{U_0})$
，且该模同构于 $A/fA$。

\medskip\noindent
由上可知，$\mathcal{F}' = \widetilde{M}$ 是凝聚
$\mathcal{O}_X$-模，其在 $U$ 上的限制同构于 $\mathcal{L}$。同构
$\lambda : \mathcal{L}/f\mathcal{L} \to \mathcal{L}_0|_{U_0}$
在整体截面上确定映射 $M/fM \to M_0$，该映射在 $U_0$ 上是同构。
由于 $\text{depth}(M) \geq 2$，可知
$H^0_\mathfrak m(M/fM) = 0$，从而
$M/fM \to M_0$ 是单射。因此，由定义
$$
\chi(\mathcal{L}, \mathcal{L}_0, \lambda) =
\text{length}_A \Coker(M/fM \to M_0)
$$
这就给出引理的第一个断言。

\medskip\noindent
最后，若此长度为 $0$，则 $M \to M_0$ 是满射。因此可以找到
$s \in M = \Gamma(U, \mathcal{L})$，它映到 $\mathcal{L}_0$
的一个平凡化截面。考虑由下列正合列定义的有限 $A$-模 $K$、$Q$：
$$
0 \to K \to A \xrightarrow{s} M \to Q \to 0
$$
$K$ 和 $Q$ 的支撑均不与 $U_0$ 相交，因为 $s$ 在 $U_0$ 的各点
都非零。利用代数，引理 \ref{algebra-lemma-depth-in-ses}，可知
$\text{depth}(K) \geq 2$（注意，$As \subset M$ 有
$\text{depth} \geq 1$，因为它是 $M$ 的子模）。因此，由代数，引理
\ref{algebra-lemma-bound-depth}，$K$ 的支撑若非空，则其维数
$\geq 2$。这与
$\text{Supp}(M) \cap V(f) \subset \{\mathfrak m\}$ 矛盾，除非
$K = 0$。当 $K = 0$ 时，可知
$\text{depth}(Q) \geq 2$，并且如前得出 $Q = 0$。因此
$A \cong M$，且 $\mathcal{L}$ 平凡。
\end{proof}







\section{穿孔谱上的可逆模，I}
\label{algebraization-section-local-lefschetz-for-pic}

\noindent
本节证明关于 Picard 群的若干局部 Lefschetz 定理。其中一些思路取自
\cite{Kollar-pic}、\cite{Bhatt-local} 和 \cite{Kollar-map-pic}。

\begin{lemma}
\label{algebraization-lemma-injective-torsion-in-pic}
令 $(A, \mathfrak m)$ 为 Noether 局部环。令 $f \in \mathfrak m$
为非零因子，并假设 $\text{depth}(A/fA) \geq 2$，等价地，
$\text{depth}(A) \geq 3$。令 $U$、相应地 $U_0$ 分别为
$A$、相应地 $A/fA$ 的穿孔谱。映射
$$
\Pic(U) \to \Pic(U_0)
$$
在挠部分上是单射。
\end{lemma}

\begin{proof}
令 $\mathcal{L}$ 为可逆 $\mathcal{O}_U$-模。注意，$\mathcal{L}$
映到 $0$（在 $\Pic(U_0)$ 中），当且仅当能把 $\mathcal{L}$ 延拓为
可逆凝聚三元组 $(\mathcal{L}, \mathcal{L}_0, \lambda)$，如
第 \ref{algebraization-section-coherent-triples} 节所述。由命题
\ref{algebraization-proposition-hilbert-triple}，函数
$$
n \longmapsto \chi((\mathcal{L}, \mathcal{L}_0, \lambda)^{\otimes n})
$$
是多项式。由引理 \ref{algebraization-lemma-nonnegative-chi-triple}，该多项式的值为零
当且仅当 $\mathcal{L}^{\otimes n}$ 平凡。因此，若 $\mathcal{L}$ 为挠元，
则该多项式有无穷多个零点，从而恒为零，于是 $\mathcal{L}$ 平凡。
\end{proof}

\begin{proposition}[Koll\'ar]
\label{algebraization-proposition-injective-pic}
\begin{reference}
\cite[Theorem 1.9]{Kollar-map-pic}
\end{reference}
令 $(A, \mathfrak m)$ 为 Noether 局部环，并令 $f \in \mathfrak m$。假设
\begin{enumerate}
\item $A$ 有对偶化复形；
\item $f$ 是非零因子；
\item $\text{depth}(A/fA) \geq 2$，等价地，
$\text{depth}(A) \geq 3$；
\item 若 $f \in \mathfrak p \subset A$ 是素理想且
$\dim(A/\mathfrak p) = 2$，则 $\text{depth}(A_\mathfrak p) \geq 2$。
\end{enumerate}
令 $U$、相应地 $U_0$ 分别为 $A$、相应地 $A/fA$ 的穿孔谱。映射
$$
\Pic(U) \to \Pic(U_0)
$$
是单射。最后，若 (1)、(2)、(3) 成立，$A$ 满足 $(S_2)$，且
$\dim(A) \geq 4$，则 (4) 成立。
\end{proposition}

\begin{proof}
令 $\mathcal{L}$ 为可逆 $\mathcal{O}_U$-模。注意，$\mathcal{L}$
映到 $0$（在 $\Pic(U_0)$ 中），当且仅当能把 $\mathcal{L}$ 延拓为
可逆凝聚三元组 $(\mathcal{L}, \mathcal{L}_0, \lambda)$，如
第 \ref{algebraization-section-coherent-triples} 节所述。由命题
\ref{algebraization-proposition-hilbert-triple}，函数
$$
n \longmapsto \chi((\mathcal{L}, \mathcal{L}_0, \lambda)^{\otimes n})
$$
是多项式 $P$。由引理 \ref{algebraization-lemma-nonnegative-chi-triple}，不等式
$P(n) \geq 0$ 对所有 $n \in \mathbf{Z}$ 都成立，且等号成立当且仅当
$\mathcal{L}^{\otimes n}$ 平凡。特别地，$P(0) = 0$，而 $P$
或者恒为零，此时结论成立；或者 $P$ 的次数为偶数且 $\geq 2$。

\medskip\noindent
置 $M = \Gamma(U, \mathcal{L})$ 以及
$M_0 = \Gamma(X_0, \mathcal{L}_0) = \Gamma(U_0, \mathcal{L}_0)$。
于是 $M$ 是有限 $A$-模，其深度 $\geq 2$，并且
$M_0 \cong A/fA$，见引理 \ref{algebraization-lemma-nonnegative-chi-triple} 的证明。
注意，$H^2_\mathfrak m(M)$ 是有限 $A$-模；这是由局部上同调，引理
\ref{local-cohomology-lemma-local-finiteness-for-finite-locally-free}
以及下述事实得到的：$H^i_\mathfrak m(A) = 0$ 对 $i = 0, 1, 2$ 成立，
因为 $\text{depth}(A) \geq 3$。考虑短正合列
$$
0 \to M/fM \to M_0 \to Q \to 0
$$
引理 \ref{algebraization-lemma-nonnegative-chi-triple} 告诉我们，$Q$ 具有有限长度，
且该长度等于 $\chi(\mathcal{L}, \mathcal{L}_0, \lambda)$。由上式的
局部上同调长正合列，得到 $Q = H^1_\mathfrak m(M/fM)$ 以及
$H^i_\mathfrak m(M/fM) = H^i_\mathfrak m(M_0) \cong H^i_\mathfrak m(A/fA)$
对 $i > 1$ 成立。再考虑与序列
$0 \to M \to M \to M/fM \to 0$ 相伴的局部上同调长正合列，其开头为
$$
0 \to Q \to H^2_\mathfrak m(M) \to H^2_\mathfrak m(M) \to
H^2_\mathfrak m(A/fA)
$$
利用长度的可加性，可知
$\chi(\mathcal{L}, \mathcal{L}_0, \lambda)$
等于映射 $H^2_\mathfrak m(M) \to H^2_\mathfrak m(A/fA)$ 的像的长度。

\medskip\noindent
先在一个特殊情形中证明命题，以阐明余下证明。暂且假设
$H^2_\mathfrak m(A/fA)$ 是有限长度模。此时有
$P(1) \leq \text{length}_A H^2_\mathfrak m(A/fA)$。把完全相同的论证
应用于 $\mathcal{L}^{\otimes n}$，可知
$P(n) \leq \text{length}_A H^2_\mathfrak m(A/fA)$ 对所有 $n$ 成立。
因此 $P$ 不可能有正次数，结论成立。在余下证明中，我们将修改该论证，
给出 $P(n)$ 的线性上界，这已足够。

\medskip\noindent
研究映射
$H^2_\mathfrak m(M) \to H^2_\mathfrak m(M_0) \cong H^2_\mathfrak m(A/fA)$.
以 $\omega_A^\bullet$ 表示为 $A$ 选取的正规化对偶化复形。由局部对偶
（对偶化复形，引理
\ref{dualizing-lemma-special-case-local-duality}），此映射是下列映射的
Matlis 对偶：
$$
\text{Ext}^{-2}_A(M, \omega_A^\bullet)
\longleftarrow
\text{Ext}^{-2}_A(M_0, \omega_A^\bullet)
$$
因而其像具有相同的有限长度。有限 $A$-模
$\text{Ext}^{-2}_A(M_0, \omega_A^\bullet)$
的支撑若非空，则由 $\mathfrak m$ 以及有限多个素理想
$\mathfrak p_1, \ldots, \mathfrak p_r$ 组成；这些素理想包含 $f$，并且
$\dim(A/\mathfrak p_i) = 1$。事实上，由局部上同调，引理
\ref{local-cohomology-lemma-sitting-in-degrees}，该支撑包含于满足
$\mathfrak p \subset A$ 以及
$\text{depth}_{A_\mathfrak p}(M_{0, \mathfrak p}) + \dim(A/\mathfrak p) \leq 2$.
的素理想之集。因此，只需证明不存在素理想 $\mathfrak p$，它包含 $f$，
使得 $\dim(A/\mathfrak p) = 2$ 且
$\text{depth}_{A_\mathfrak p}(M_{0, \mathfrak p}) = 0$。然而，由
$M_{0, \mathfrak p} \cong (A/fA)_\mathfrak p$，这会推出
$\text{depth}(A_\mathfrak p) = 1$，与假设 (4) 矛盾。
选择截面 $t \in \Gamma(U, \mathcal{L}^{\otimes -1})$，使其在各点
$\mathfrak p_1, \ldots, \mathfrak p_r$ 均不消失，见性质，引理
\ref{properties-lemma-quasi-affine-invertible-nonvanishing-section}。
整体截面上乘以 $t$ 确定映射 $t : M \to A$；它给出同构
$M_{\mathfrak p_i} \to A_{\mathfrak p_i}$，这对
$i = 1, \ldots, r$ 成立。以
$t_0 = t|_{U_0}$ 表示对应截面
$\Gamma(U_0, \mathcal{L}_0^{\otimes -1})$；它同样确定映射
$t_0 : M_0 \to A/fA$，且与 $t$ 相容。于是有交换图
$$
\xymatrix{
\text{Ext}^{-2}_A(M, \omega_A^\bullet) &
\text{Ext}^{-2}_A(M_0, \omega_A^\bullet) \ar[l] \\
\text{Ext}^{-2}_A(A, \omega_A^\bullet) \ar[u]^t &
\text{Ext}^{-2}_A(A/fA, \omega_A^\bullet) \ar[l] \ar[u]_{t_0}
}
$$
由此可知，顶部水平映射之像的长度至多等于
$\text{Ext}^{-2}_A(A/fA, \omega_A^\bullet)$ 的长度加上 $t_0$ 的余核长度。

\medskip\noindent
不过，若把 $\mathcal{L}$ 替换为 $\mathcal{L}^n$，其中 $n > 1$，
则可以使用
$$
t^n :
M_n = \Gamma(U, \mathcal{L}^{\otimes n})
\longrightarrow
\Gamma(U, \mathcal{O}_U) = A
$$
来代替 $t$。这会把 $t_0$ 替换为它的 $n$ 次幂。因此，映射
$\text{Ext}^{-2}_A(M_n, \omega_A^\bullet) \leftarrow
\text{Ext}^{-2}_A(M_{n, 0}, \omega_A^\bullet)$
的像的长度至多等于 $\text{Ext}^{-2}_A(A/fA, \omega_A^\bullet)$ 的长度
加上下列映射的余核长度：
$$
t_0^n : 
\text{Ext}^{-2}_A(A/fA, \omega_A^\bullet)
\longrightarrow
\text{Ext}^{-2}_A(M_{n, 0}, \omega_A^\bullet)
$$
经由同构 $M_0 \cong A/fA$，映射 $t_0$ 变为
$g : A/fA \to A/fA$，其中 $g \in A/fA$；经由相应同构
$M_{n, 0} \cong A/fA$，映射 $t_0^n$ 变为
$g^n : A/fA \to A/fA$。因此，上述余核的长度就是
$\text{Ext}^{-2}_A(A/fA, \omega_A^\bullet)$ 对 $g^n$ 的商的长度。
由于 $\text{Ext}^{-2}_A(A/fA, \omega_A^\bullet)$ 是有限 $A$-模，
其支撑 $T$ 的维数为 $1$，并且 $V(g) \cap T$ 只含闭点，
这是由对 $t$ 的选择得到的。所以由代数，引理
\ref{algebra-lemma-support-dimension-d}，该长度关于 $n$ 线性增长。

\medskip\noindent
最后证明末尾断言。假设 $f \in \mathfrak m \subset A$ 满足
(1)、(2)、(3)，$A$ 满足 $(S_2)$，并且 $\dim(A) \geq 4$。
条件 (1) 蕴含 $A$ 是链状环，见对偶化复形，引理
\ref{dualizing-lemma-universally-catenary}。于是
$\Spec(A)$ 等维，见局部上同调，引理
\ref{local-cohomology-lemma-catenary-S2-equidimensional}。因此
$\dim(A_\mathfrak p) + \dim(A/\mathfrak p) \geq 4$ 对每个素理想
$\mathfrak p$ 都成立，而它是 $A$ 的素理想。于是
$\text{depth}(A_\mathfrak p) \geq \min(2, \dim(A_\mathfrak p))
\geq \min(2, 4 - \dim(A/\mathfrak p))$，从而 (4) 成立。
\end{proof}

\begin{remark}
\label{algebraization-remark-compare-SGA2}
在 SGA2 中可找到下列结果。令 $(A, \mathfrak m)$ 为 Noether 局部环，
并令 $f \in \mathfrak m$。假设 $A$ 是正则环的商，元素
$f$ 是非零因子，并且
\begin{enumerate}
\item[(a)] 若 $\mathfrak p \subset A$ 是素理想且
$\dim(A/\mathfrak p) = 1$，则 $\text{depth}(A_\mathfrak p) \geq 2$；
\item[(b)] $\text{depth}(A/fA) \geq 3$，等价地，
$\text{depth}(A) \geq 4$。
\end{enumerate}
令 $U$、相应地 $U_0$ 分别为 $A$、相应地 $A/fA$ 的穿孔谱。那么映射
$$
\Pic(U) \to \Pic(U_0)
$$
是单射。这就是 \cite[Exposee XI, Lemma 3.16]{SGA2}\footnote{条件
(a) 由条件 (b) 推出，见代数，引理
\ref{algebra-lemma-depth-localization}。}。SGA2 的这一结果由命题
\ref{algebraization-proposition-injective-pic} 推出，因为
\begin{enumerate}
\item 正则环的商有对偶化复形（见对偶化复形，引理
\ref{dualizing-lemma-regular-gorenstein} 和命题
\ref{dualizing-proposition-dualizing-essentially-finite-type}）；
\item 若 $\text{depth}(A) \geq 4$，则
$\text{depth}(A_\mathfrak p) \geq 2$ 对所有素理想 $\mathfrak p$ 都成立，
只要 $\dim(A/\mathfrak p) = 2$，见代数，引理
\ref{algebra-lemma-depth-localization}。
\end{enumerate}
\end{remark}






\section{穿孔谱上的可逆模，II}
\label{algebraization-section-local-lefschetz-for-pic-surjective}

\noindent
下面转向 Picard 群局部 Lefschetz 定理中的满性。首先，为了把 $U_0$
上的可逆模延拓到一个开邻域，有如下简单判据。

\begin{lemma}
\label{algebraization-lemma-surjective-Pic-first}
令 $(A, \mathfrak m)$ 为 Noether 局部环，且 $f \in \mathfrak m$。假设
\begin{enumerate}
\item $A$ 是 $f$-进完备的；
\item $f$ 是非零因子；
\item $H^1_\mathfrak m(A/fA)$ 和 $H^2_\mathfrak m(A/fA)$
是有限 $A$-模；
\item $H^3_\mathfrak m(A/fA) = 0$\footnote{注意，若
$\text{depth}(A/fA) \geq 4$，等价地
$\text{depth}(A) \geq 5$，则 (3) 和 (4) 成立。}。
\end{enumerate}
令 $U$、相应地 $U_0$ 分别为 $A$、相应地 $A/fA$ 的穿孔谱。那么
$$
\colim_{U_0 \subset U' \subset U\text{ 开}} \Pic(U')
\longrightarrow
\Pic(U_0)
$$
是满射。
\end{lemma}

\begin{proof}
令 $U_0 \subset U_n \subset U$ 为第 $n$ 个无穷小邻域，即 $U_0$
的无穷小邻域。注意，
$U_n$ 在 $U_{n + 1}$ 中的理想层同构于 $\mathcal{O}_{U_0}$，因为
$U_0 \subset U$ 是由非零因子 $f$ 截出的主闭子概形。因此有 Abel 群正合列
$$
\Pic(U_{n + 1}) \to \Pic(U_n) \to
H^2(U_0, \mathcal{O}_{U_0}) = H^3_\mathfrak m(A/fA) = 0
$$
见态射进阶，引理
\ref{more-morphisms-lemma-picard-group-first-order-thickening}。
因此，每个可逆 $\mathcal{O}_{U_0}$-模都是某个可逆凝聚形式模的限制，
也就是说，它是 $\textit{Coh}(U, f\mathcal{O}_U)$ 的可逆对象。
应用引理 \ref{algebraization-lemma-equivalence} 即得结论。
\end{proof}

\begin{remark}
\label{algebraization-remark-surjective-Pic-second}
令 $(A, \mathfrak m)$ 为 Noether 局部环，且 $f \in \mathfrak m$。
若假设下列条件，则引理 \ref{algebraization-lemma-surjective-Pic-first} 的结论成立：
\begin{enumerate}
\item $A$ 有对偶化复形；
\item $A$ 是 $f$-进完备的；
\item $f$ 是非零因子；
\item 下列条件之一成立：
\begin{enumerate}
\item $A_f$ 满足 $(S_2)$，并且对满足 $\mathfrak p \subset A$、
$f \not \in \mathfrak p$ 的极小素理想，有
$\dim(A/\mathfrak p) \geq 4$；或者
\item 若 $\mathfrak p \not \in V(f)$ 且
$V(\mathfrak p) \cap V(f) \not = \{\mathfrak m\}$，则
$\text{depth}(A_\mathfrak p) + \dim(A/\mathfrak p) > 3$。
\end{enumerate}
\item $H^3_{\mathfrak m}(A/fA) = 0$.
\end{enumerate}
证明与引理 \ref{algebraization-lemma-surjective-Pic-first} 的证明完全相同，只是用引理
\ref{algebraization-lemma-equivalence-better} 代替引理 \ref{algebraization-lemma-equivalence}。
这里需要说明两点：(a) 似乎很难找到能够知道
$H^3_{\mathfrak m}(A/fA) = 0$ 的例子，而又不假设
$\text{depth}(A/fA) \geq 4$；(b) 引理
\ref{algebraization-lemma-equivalence-better} 的证明比引理
\ref{algebraization-lemma-equivalence} 的证明难得多。
\end{remark}

\begin{lemma}
\label{algebraization-lemma-surjective-Pic-first-better}
令 $(A, \mathfrak m)$ 为 Noether 局部环，且 $f \in \mathfrak m$。假设
\begin{enumerate}
\item 引理 \ref{algebraization-lemma-surjective-Pic-first} 的条件成立；
\item 对每个极大理想 $\mathfrak p \subset A_f$，局部环
$(A_f)_\mathfrak p$ 的穿孔谱具有平凡 Picard 群。
\end{enumerate}
令 $U$、相应地 $U_0$ 分别为 $A$、相应地 $A/fA$ 的穿孔谱。那么
$$
\Pic(U) \longrightarrow \Pic(U_0)
$$
是满射。
\end{lemma}

\begin{proof}
令 $\mathcal{L}_0 \in \Pic(U_0)$。由引理
\ref{algebraization-lemma-surjective-Pic-first}，存在开集 $U_0 \subset U' \subset U$
以及 $\mathcal{L}' \in \Pic(U')$，其在 $U_0$ 上的限制为
$\mathcal{L}_0$。由于 $U' \supset U_0$，可知 $U \setminus U'$ 由与
(2) 中素理想 $\mathfrak p_1, \ldots, \mathfrak p_n$ 对应的点组成。
由假设，可以找到可逆模 $\mathcal{L}'_i$，它定义在
$\Spec(A_{\mathfrak p_i})$ 上，并与 $\mathcal{L}'$ 在穿孔谱
$U' \times_U \Spec(A_{\mathfrak p_i})$ 上相符，
因为平凡可逆模总能延拓。把极限，引理
\ref{limits-lemma-glueing-near-closed-point-modules} 应用 $n$ 次，
可知 $\mathcal{L}'$ 延拓为 $U$ 上的可逆模。
\end{proof}

\begin{lemma}
\label{algebraization-lemma-local-pic-to-completion}
令 $(A, \mathfrak m)$ 为深度 $\geq 2$ 的 Noether 局部环，并令
$A^\wedge$ 为其完备化。令 $U$、相应地 $U^\wedge$ 分别为
$A$、相应地 $A^\wedge$ 的穿孔谱。那么
$\Pic(U) \to \Pic(U^\wedge)$ 是单射。
\end{lemma}

\begin{proof}
令 $\mathcal{L}$ 为可逆 $\mathcal{O}_U$-模，其拉回为
$\mathcal{L}^\wedge$，定义在 $U^\wedge$ 上。由深度假设、对偶化复形，引理
\ref{dualizing-lemma-depth} 以及局部上同调，引理
\ref{local-cohomology-lemma-finiteness-pushforwards-and-H1-local}，有
$H^0(U, \mathcal{O}_U) = A$。因此，$\mathcal{L}$ 平凡当且仅当
$M = H^0(U, \mathcal{L})$ 作为 $A$-模同构于 $A$（略去细节）。
由于 $A \to A^\wedge$ 平坦，由平坦基变换有
$M \otimes_A A^\wedge = \Gamma(U^\wedge, \mathcal{L}^\wedge)$，见
概形上同调，引理 \ref{coherent-lemma-flat-base-change-cohomology}。
最后，容易看出 $M \cong A$ 当且仅当
$M \otimes_A A^\wedge \cong A^\wedge$。
\end{proof}

\begin{lemma}
\label{algebraization-lemma-trivial-local-pic-regular}
令 $(A, \mathfrak m)$ 为正则局部环。那么 $A$ 的穿孔谱的 Picard 群平凡。
\end{lemma}

\begin{proof}
结合除子，引理 \ref{divisors-lemma-extend-invertible-module} 与
代数进阶，引理 \ref{more-algebra-lemma-regular-local-UFD}。
\end{proof}

\noindent
现在可以自举前述结果，证明维数 $\geq 4$ 的完全交的穿孔谱具有平凡
Picard 群。回忆，若一个 Noether 局部环的完备化是某个正则局部环
关于一个由正则序列生成的理想的商，则称该环为完全交。见除幂代数，
第 \ref{dpa-section-lci} 节的讨论。

\begin{proposition}[Grothendieck]
\label{algebraization-proposition-trivial-local-pic-complete-intersection}
令 $(A, \mathfrak m)$ 为 Noether 局部环。若 $A$ 是维数
$\geq 4$ 的完全交，则 $A$ 的穿孔谱的 Picard 群平凡。
\end{proposition}

\begin{proof}
由引理 \ref{algebraization-lemma-local-pic-to-completion}，可以假设 $A$ 是完备局部环。
由假设可写成 $A = B/(f_1, \ldots, f_r)$，其中 $B$ 是完备正则局部环，
而 $f_1, \ldots, f_r$ 是正则序列。我们对 $r$ 作归纳以完成证明。
基本情形 $r = 0$ 由引理 \ref{algebraization-lemma-trivial-local-pic-regular} 得到。

\medskip\noindent
假设 $A = B/(f_1, \ldots, f_r)$，且命题对 $r - 1$ 成立。置
$A' = B/(f_1, \ldots, f_{r - 1})$，并把引理
\ref{algebraization-lemma-surjective-Pic-first-better} 应用于 $f_r \in A'$。
这是允许的：
\begin{enumerate}
\item 引理 \ref{algebraization-lemma-surjective-Pic-first} 的条件 (1) 成立，因为所用局部环
均完备；
\item 引理 \ref{algebraization-lemma-surjective-Pic-first} 的条件 (2) 成立，因为
$f_1, \ldots, f_r$ 是正则序列；
\item 引理 \ref{algebraization-lemma-surjective-Pic-first} 的条件 (3) 和 (4) 成立，因为
$A = A'/f_r A'$ 是维数 $\dim(A) \geq 4$ 的 Cohen--Macaulay 环；
\item 引理 \ref{algebraization-lemma-surjective-Pic-first-better} 的条件 (2) 由归纳假设
成立，因为 $\dim((A'_{f_r})_\mathfrak p) \geq 4$ 对极大素理想
$\mathfrak p$ 成立，而该素理想属于 $A'_{f_r}$；并且
$(A'_{f_r})_\mathfrak p = B_\mathfrak q/(f_1, \ldots, f_{r - 1})$
对某个素理想 $\mathfrak q \subset B$ 成立，且 $B_\mathfrak q$ 正则。
\end{enumerate}
证毕。
\end{proof}

\begin{example}
\label{algebraization-example-grothendieck-sharp}
命题 \ref{algebraization-proposition-trivial-local-pic-complete-intersection} 中的维数界是
尖锐的。例如，$A = k[[x, y, z, w]]/(xy - zw)$ 的穿孔谱的 Picard 群
非平凡。事实上，理想 $I = (x, z)$ 在穿孔谱上截出一个有效 Cartier
除子 $D$；该穿孔谱记为 $U$，它是 $A$ 的穿孔谱。容易看出
$I_x, I_y, I_z, I_w$ 分别是 $A_x, A_y, A_z, A_w$ 中的可逆理想。
另一方面，$A/I$ 的深度 $\geq 1$（事实上为 $2$），故 $I$ 的深度
$\geq 2$（事实上为 $3$），从而
$I = \Gamma(U, \mathcal{O}_U(-D))$。因此，若
$\mathcal{O}_U(-D)$ 平凡，则会有
$I \cong \Gamma(U, \mathcal{O}_U) = A$，但这是不成立的，因为 $I$
不是由 $1$ 个元素生成的。
\end{example}

\begin{example}
\label{algebraization-example-grothendieck-sharp-bis}
命题 \ref{algebraization-proposition-trivial-local-pic-complete-intersection} 不能推广到如下商：
$$
A = B/(f_1, \ldots, f_r)
$$
其中 $B$ 正则，且 $\dim(B) - r \geq 4$。换言之，
$f_1, \ldots, f_r$ 构成正则序列这一条件，一般是 $A$ 的穿孔谱的
Picard 群消失所必需的。
具体而言，令 $k$ 为域，并置
$$
A = k[[a, b, x, y, z, u, v, w]]/(a^3, b^3, xa^2 + yab + zb^2, w^2)
$$
注意 $A = A_0[w]/(w^2)$，其中
$A_0 = k[[a, b, x, y, z, u, v]]/(a^3, b^3, xa^2 + yab + zb^2)$。
下面将证明 $A_0$ 的深度为 $2$。以 $U$ 表示 $A$ 的穿孔谱，以
$U_0$ 表示 $A_0$ 的穿孔谱。注意有短正合列
$0 \to A_0 \to A \to A_0 \to 0$，其中第一个箭头由乘以 $w$ 给出。
由态射进阶，引理
\ref{more-morphisms-lemma-picard-group-first-order-thickening}，得到正合列
$$
H^0(U, \mathcal{O}_U^*) \to
H^0(U_0, \mathcal{O}_{U_0}^*) \to
H^1(U_0, \mathcal{O}_{U_0}) \to
\Pic(U)
$$
由于 $A_0$ 的深度、从而 $A$ 的深度为 $2$，可知
$H^0(U_0, \mathcal{O}_{U_0}) = A_0$ 且 $H^0(U, \mathcal{O}_U) = A$，
并且 $H^1(U_0, \mathcal{O}_{U_0})$ 非零，见对偶化复形，引理
\ref{dualizing-lemma-depth} 以及局部上同调，引理
\ref{local-cohomology-lemma-local-cohomology}。因此，上述最后一个箭头非零，
从而 $\Pic(U)$ 非零。

\medskip\noindent
要证明 $A_0$ 的深度为 $2$，只需证明
$A_1 = k[[a, b, x, y, z]]/(a^3, b^3, xa^2 + yab + zb^2)$
的深度为 $0$。这是因为 $a^2b^2$ 映到 $A_1$ 中的一个非零元素，
且该元素被每个变量 $a, b, x, y, z$ 零化。例如
$ya^2b^2 = (yab)(ab) = - (xa^2 + zb^2)(ab) = -xa^3b - yab^3 = 0$ 在 $A_1$ 中。
其他情形类似。
\end{example}










\section{Lefschetz 定理的应用}
\label{algebraization-section-lefschetz}

\noindent
本节讨论投射概形 $P$ 上的凝聚层与沿丰沛除子 $Q$
的形式完备化上的凝聚模之间的关系。

\medskip\noindent
设 $k$ 为域，$P$ 为 $k$ 上的固有概形。
设 $\mathcal{L}$ 为丰沛的可逆 $\mathcal{O}_P$-模。
取截面 $s \in \Gamma(P, \mathcal{L})$\footnote{我们不要求
$s$ 是正则截面。因此，$Q$ 仅为 $P$ 的局部主闭子概形，
而未必为有效 Cartier 除子。}，并令
$Q = Z(s)$ 为其零点概形，见
《除子》，定义 \ref{divisors-definition-zero-scheme-s}。
对所有 $n \geq 1$，记 $Q_n = Z(s^n)$
为第 $n$ 个无穷小邻域，其中心为 $Q$。
若 $\mathcal{F}$ 是凝聚 $\mathcal{O}_P$-模，则记
$\mathcal{F}_n = \mathcal{F}|_{Q_n}$ 为其限制；换言之，
它是 $\mathcal{F}$ 沿闭浸入 $Q_n \to P$ 的拉回。

\begin{proposition}
\label{algebraization-proposition-lefschetz}
在上述情形中，假设存在整数 $\sigma$，使得对每一点
$p \in P \setminus Q$ 都有
$$
\text{depth}(\mathcal{F}_p) + \dim(\overline{\{p\}}) > \sigma
$$
则映射
$$
H^i(P, \mathcal{F}) \longrightarrow \lim H^i(Q_n, \mathcal{F}_n)
$$
对 $0 \leq i < \sigma$ 为同构。
\end{proposition}

\begin{proof}
我们将使用《态射进阶》，引理 \ref{more-morphisms-lemma-apply-proj-spec}，
并在不再另行说明的情况下使用
《态射进阶》，第 \ref{more-morphisms-section-proj-spec} 节
中的记号和结果；若至少不了解该引理的陈述，本证明便无法理解。
注意，在我们的情形中
$A = \bigoplus_{m \geq 0} \Gamma(P, \mathcal{L}^{\otimes m})$
是有限型 $k$-代数，且它的每个分次部分都是有限维
$k$-向量空间，见《概形上同调》，
引理 \ref{coherent-lemma-coherent-proper-ample}。

\medskip\noindent
我们可以并且确实把 $s$ 看成元素 $f \in A_1 \subset A$，
即次数为 $1$ 的一个 $A$-齐次元素。记
$Y = V(f) \subset X$ 为由 $f$ 定义的闭子概形。
于是从概形论意义看，$U \cap Y = (\pi|_U)^{-1}(Q)$。
回顾记号
$\mathcal{F}_U = \pi^*\mathcal{F}|_U = (\pi|_U)^*\mathcal{F}$.
这是凝聚 $\mathcal{O}_U$-模。选取有限 $A$-模 $M$，使得
$\mathcal{F}_U = \widetilde{M}|_U$
（其存在性见《局部上同调》，引理
\ref{local-cohomology-lemma-finiteness-pushforwards-and-H1-local}）。
我们断言：$H^i_Z(M)$ 被 $f$ 的某个幂零化，只要
$i \leq \sigma + 1$。

\medskip\noindent
为证明该断言，我们应用《局部上同调》，
命题 \ref{local-cohomology-proposition-annihilator}。
将条件译成几何语言可知，只须证明：对
$u \in U$、$u \not \in Y$ 以及
$z \in \overline{\{u\}} \cap Z$，有
$$
\text{depth}(\mathcal{F}_{U, u}) +
\dim(\mathcal{O}_{\overline{\{u\}}, z}) > \sigma + 1
$$
这只需稍加思考。

\medskip\noindent
注意，$Z = \Spec(A_0)$ 是 $X$ 的有限个闭点所成的集合，
因为 $A_0$ 是有限维 $k$-代数。
（若读者希望 $Z$ 只有一个点，可以把有限 $k$-代数
$A_0$ 换成 $k$；这不会影响证明的其余部分。）

\medskip\noindent
态射 $\pi : L \to P$ 及其限制 $\pi|_U : U \to P$
都是相对维数为 $1$ 的光滑态射。
取 $u \in U$、$u \not \in Y$ 以及
$z \in \overline{\{u\}} \cap Z$，并令
$p = \pi(u) \in P \setminus Q$ 为其像。
于是 $u$ 或为 $\pi$ 在 $p$ 上纤维的泛点，或为该纤维的闭点。
若 $u$ 是纤维的泛点，则
$\text{depth}(\mathcal{F}_{U, u}) = \text{depth}(\mathcal{F}_p)$
且
$\dim(\overline{\{u\}}) = \dim(\overline{\{p\}}) + 1$.
若 $u$ 是纤维的闭点，则
$\text{depth}(\mathcal{F}_{U, u}) = \text{depth}(\mathcal{F}_p) + 1$
且
$\dim(\overline{\{u\}}) = \dim(\overline{\{p\}})$.
在这两种情形下都有
$\dim(\overline{\{u\}}) = \dim(\mathcal{O}_{\overline{\{u\}}, z})$
，因为 $Z$ 的每一点都是闭点。因此所需不等式由命题陈述中的
假设推出。

\medskip\noindent
令 $A'$ 为由 $f$ 所定义的 $A$ 的进完备化。由《代数》，
引理 \ref{algebra-lemma-completion-flat}，$A \to A'$ 是平坦的。
记 $U' \subset X' = \Spec(A')$ 为 $U$ 的逆像，并类似地定义
$Y'$ 与 $Z'$。令 $\mathcal{F}'$ 为 $U'$ 上由
$\mathcal{F}_U$ 拉回所得的模，并令
$M' = M \otimes_A A'$.
由局部上同调的平坦基变换
（《局部上同调》，引理
\ref{local-cohomology-lemma-torsion-change-rings}），有
$$
H^i_{Z'}(M') = H^i_Z(M) \otimes_A A'
$$
，故对 $i \leq \sigma + 1$，这些模均被 $f$ 的某个幂零化。
考虑交换图
$$
\xymatrix{
& H^i(U, \mathcal{F}_U) \ar[ld] \ar[d] \ar[r] &
\lim H^i(U, \mathcal{F}_U/f^n\mathcal{F}_U) \ar@{=}[d] \\
H^i(U, \mathcal{F}_U) \otimes_A A' \ar@{=}[r] &
H^i(U', \mathcal{F}') \ar[r] &
\lim H^i(U', \mathcal{F}'/f^n\mathcal{F}')
}
$$
由引理 \ref{algebraization-lemma-alternative-higher} 以及刚证明的挠性，
当 $i < \sigma$ 时，下方水平箭头为同构。
水平等号来自平坦基变换
（《概形上同调》，引理
\ref{coherent-lemma-flat-base-change-cohomology}）；
竖直等号则是因为 $U \cap Y$ 与 $U' \cap Y'$，
连同它们的第 $n$ 个无穷小邻域，都彼此同构地映射
（这是因为我们关于 $f$ 作完备化）。

\medskip\noindent
应用《态射进阶》，等式
(\ref{more-morphisms-equation-cohomology-torsor})
，得到相容的直和分解
$$
\lim H^i(U, \mathcal{F}_U/f^n\mathcal{F}_U) =
\lim
\left(
\bigoplus\nolimits_{m \in \mathbf{Z}}
H^i(Q_n, \mathcal{F}_n \otimes \mathcal{L}^{\otimes m})
\right)
$$
以及
$$
H^i(U, \mathcal{F}_U) =
\bigoplus\nolimits_{m \in \mathbf{Z}}
H^i(P, \mathcal{F} \otimes \mathcal{L}^{\otimes m})
$$
因此结论由《代数》，引理 \ref{algebra-lemma-daniel-litt} 得出。
\end{proof}

\begin{lemma}
\label{algebraization-lemma-lefschetz-addendum}
设 $k$ 为域，$X$ 为 $k$ 上的固有概形。
设 $\mathcal{L}$ 为丰沛的可逆 $\mathcal{O}_X$-模，
$s \in \Gamma(X, \mathcal{L})$。令 $Y = Z(s)$ 为 $s$ 的零点概形，
其第 $n$ 个无穷小邻域为 $Y_n = Z(s^n)$。
设 $\mathcal{F}$ 为凝聚 $\mathcal{O}_X$-模。
假设对所有 $x \in X \setminus Y$ 都有
$$
\text{depth}(\mathcal{F}_x) + \dim(\overline{\{x\}}) > 1
$$
则映射
$\Gamma(V, \mathcal{F}) \to \lim \Gamma(Y_n, \mathcal{F}|_{Y_n})$
对任意开子概形 $V \subset X$，只要它包含 $Y$，都是同构。
\end{lemma}

\begin{proof}
由命题 \ref{algebraization-proposition-lefschetz}，当 $V = X$ 时结论成立。
因此只须证明映射
$\Gamma(V, \mathcal{F}) \to \lim \Gamma(Y_n, \mathcal{F}|_{Y_n})$
是单射。若 $\sigma \in \Gamma(V, \mathcal{F})$
映为零，则其支撑与 $Y$ 不交（略去细节；提示：
使用 Krull 交定理）。于是，令 $T \subset X$ 为
$\text{Supp}(\sigma)$ 的闭包；它与 $Y$ 不交。
因此 $T$ 在 $k$ 上固有（因为它在 $X$ 中闭），
又是仿射的（因为它在仿射概形 $X \setminus Y$ 中闭，见
《态射》，引理 \ref{morphisms-lemma-proper-ample-delete-affine}），
从而在 $k$ 上有限
（《态射》，引理 \ref{morphisms-lemma-finite-proper}）。
故 $T$ 是 $X$ 的有限个闭点所成的集合。
由假设，$\text{depth}(\mathcal{F}_x) \geq 2$，因而深度至少为
$1$；这里 $x \in T$。
于是
$\Gamma(V, \mathcal{F}) \to \Gamma(V \setminus T, \mathcal{F})$
为单射，故如所需 $\sigma = 0$。
\end{proof}


\begin{example}
\label{algebraization-example-lefschetz}
设 $k$ 为域，$X$ 为 $k$ 上的固有簇。
设 $Y \subset X$ 为有效 Cartier 除子，且
$\mathcal{O}_X(Y)$ 丰沛；记 $Y_n$ 为它的第 $n$ 个无穷小邻域。
设 $\mathcal{E}$ 为有限局部自由 $\mathcal{O}_X$-模。
下面列出命题 \ref{algebraization-proposition-lefschetz} 的若干特殊情形。
\begin{enumerate}
\item 若 $X$ 是曲线，则得不到任何新结论。
\item 若 $X$ 是 Cohen-Macaulay 曲面（例如正规曲面），则
$$
H^0(X, \mathcal{E}) \to \lim H^0(Y_n, \mathcal{E}|_{Y_n})
$$
为同构。
\item 若 $X$ 是 Cohen-Macaulay 三维簇，则
$$
H^0(X, \mathcal{E}) \to \lim H^0(Y_n, \mathcal{E}|_{Y_n})
\quad\text{且}\quad
H^1(X, \mathcal{E}) \to \lim H^1(Y_n, \mathcal{E}|_{Y_n})
$$
均为同构。
\end{enumerate}
其规律应当已经清楚。若 $X$ 是正规三维簇，则可以对 $H^0$
得出结论，但不能对 $H^1$ 得出结论。
\end{example}

\noindent
在证明下一个主要结果之前，我们需要一个引理。

\begin{lemma}
\label{algebraization-lemma-Gm-equivariant-extend-canonically}
在情形 \ref{algebraization-situation-algebraize} 中，设 $(\mathcal{F}_n)$ 是
$\textit{Coh}(U, I\mathcal{O}_U)$ 的一个对象。假设
\begin{enumerate}
\item $A$ 是分次环，$\mathfrak a = A_+$，且 $I$ 是齐次理想；
\item $(\mathcal{F}_n) = (\widetilde{M_n}|_U)$，其中 $(M_n)$
是分次 $A$-模的逆系统；以及
\item $(\mathcal{F}_n)$ 可典范延拓到 $X$。
\end{enumerate}
则存在有限分次 $A$-模 $N$，使得
\begin{enumerate}
\item[(a)] 逆系统 $(N/I^nN)$ 与 $(M_n)$ 在分次 $A$-模
模去 $A_+$-幂挠模所得的范畴中为 pro-同构；并且
\item[(b)] $(\mathcal{F}_n)$ 是与 $N$ 相伴的凝聚模的完备化。
\end{enumerate}
\end{lemma}

\begin{proof}
令 $(\mathcal{G}_n)$ 为引理
\ref{algebraization-lemma-canonically-algebraizable} 中的典范延拓。
$A$ 与 $M_n$ 上的分次决定如下作用
$$
a : \mathbf{G}_m \times X \longrightarrow X
$$
；这是群概形 $\mathbf{G}_m$ 在 $X$ 上的作用，使
$(\widetilde{M_n})$ 成为
$\mathbf{G}_m$-等变拟凝聚 $\mathcal{O}_X$-模的逆系统，见
《群胚》，例 \ref{groupoids-example-Gm-on-affine}。
由于 $\mathfrak a$ 与 $I$ 是齐次理想，闭子概形 $Z$、$Y$
以及开子概形 $U$ 都是 $\mathbf{G}_m$-不变的。
$(\mathcal{F}_n)$ 是 $(\widetilde{M_n})$ 的限制，
并且是 $\mathbf{G}_m$-等变凝聚 $\mathcal{O}_U$-模的逆系统。
换言之，$(\mathcal{F}_n)$ 是 $\mathbf{G}_m$-等变凝聚形式模；
这意味着存在同构
$$
\alpha : (a^*\mathcal{F}_n) \longrightarrow (p^*\mathcal{F}_n)
$$
，它定义在 $\mathbf{G}_m \times U$ 上，并满足适当的余循环条件。
由于 $a$ 与 $p$ 是仿射概形之间的平坦态射，由引理
\ref{algebraization-lemma-canonically-extend-base-change} 可知，
存在唯一同构
$$
\beta : (a^*\mathcal{G}_n) \longrightarrow (p^*\mathcal{G}_n)
$$
，它定义在 $\mathbf{G}_m \times X$ 上，并将 $\alpha$
作为在 $\mathbf{G}_m \times U$ 上的限制。
唯一性保证 $\beta$ 满足相应的余循环条件。
由此，每个 $\mathcal{G}_n$ 都以与转移映射相容的方式
成为 $\mathbf{G}_m$-等变凝聚 $\mathcal{O}_X$-模。

\medskip\noindent
由《群胚》，引理 \ref{groupoids-lemma-Gm-equivariant-module}，
$\mathcal{G}_n$ 连同其 $\mathbf{G}_m$-等变结构
对应于分次 $A$-模 $N_n$。转移映射
$N_{n + 1} \to N_n$ 是分次模同态。注意，$N_n$ 是有限
$A$-模，并且 $N_n = N_{n + 1}/I^n N_{n + 1}$，因为
$(\mathcal{G}_n)$ 是 $\textit{Coh}(X, I\mathcal{O}_X)$ 的对象。
令 $N$ 为《代数》，引理
\ref{algebra-lemma-finiteness-graded} 中得到的有限分次 $A$-模。
于是 $N_n = N/I^nN$，故 $(\mathcal{G}_n)$
是与 $N$ 相伴的凝聚模的完备化，尤其 (b) 成立。

\medskip\noindent
为证明 (a)，还需更细致地展开上述情形。
首先，注意映射 $M_n \to H^0(U, \mathcal{F}_n)$ 的核与余核
都是 $A_+$-幂挠的（《局部上同调》，引理
\ref{local-cohomology-lemma-finiteness-pushforwards-and-H1-local}）。
注意，$H^0(U, \mathcal{F}_n)$ 带有自然分次，使这些映射及系统的
转移映射都是分次 $A$-模同态；例如，可以利用
$(U \to X)_*\mathcal{F}_n$ 是 $\mathbf{G}_m$-等变模，
并且定义在 $X$ 上，再应用《群胚》，
引理 \ref{groupoids-lemma-Gm-equivariant-module}。
其次，回顾定义 \ref{algebraization-definition-canonically-algebraizable}
和引理 \ref{algebraization-lemma-canonically-algebraizable}：
$(N_n)$ 与 $(H^0(U, \mathcal{F}_n))$ 为 pro-同构。
我们略去验证定义此 pro-同构的映射均为分次模同态。
因此，$(N_n)$ 与 $(M_n)$ 在分次 $A$-模模去
$A_+$-幂挠模所得的范畴中为 pro-同构。
\end{proof}

\noindent
设 $k$ 为域，$P$ 为 $k$ 上的固有概形。
设 $\mathcal{L}$ 为丰沛的可逆 $\mathcal{O}_P$-模。
取截面 $s \in \Gamma(P, \mathcal{L})$，并令
$Q = Z(s)$ 为其零点概形，见《除子》，
定义 \ref{divisors-definition-zero-scheme-s}。
令 $\mathcal{I} \subset \mathcal{O}_P$ 为 $Q$ 的理想层。
我们用 $\textit{Coh}(P, \mathcal{I})$ 表示
《概形上同调》，第 \ref{coherent-section-coherent-formal} 节
中引入的凝聚形式模范畴。

\begin{proposition}
\label{algebraization-proposition-lefschetz-existence}
在上述情形中，设 $(\mathcal{F}_n)$ 是
$\textit{Coh}(P, \mathcal{I})$ 的一个对象。假设对所有
$q \in Q$ 以及所有素理想
$\mathfrak p \in \mathcal{O}_{P, q}^\wedge$，若
$\mathfrak p \not \in V(\mathcal{I}_q^\wedge)$，都有
$$
\text{depth}((\mathcal{F}_q^\wedge)_\mathfrak p) +
\dim(\mathcal{O}_{P, q}^\wedge/\mathfrak p) +
\dim(\overline{\{q\}}) > 2
$$
则 $(\mathcal{F}_n)$ 是某个凝聚 $\mathcal{O}_P$-模的完备化。
\end{proposition}

\begin{proof}
由《概形上同调》，引理 \ref{coherent-lemma-existence-easy}，
为证明本命题，可以把 $(\mathcal{F}_n)$ 换成一个与它仅相差
$\mathcal{I}$-挠的对象（更精确的说明见下文）。
令 $T' = \{q \in Q \mid \dim(\overline{\{q\}}) = 0\}$，
并令 $T = \{q \in Q \mid \dim(\overline{\{q\}}) \leq 1\}$。
命题中的假设恰好说明
$Q \subset P$、$(\mathcal{F}_n)$ 以及
$T' \subset T \subset Q$ 在 $d = 1$ 时满足
引理 \ref{algebraization-lemma-improvement-formal-coherent-module-better}
的条件；除了直接整理不等式外，还要用到
$V(\mathfrak p) \cap V(\mathcal{I}^\wedge_y) = \{\mathfrak m^\wedge_y\}
\Leftrightarrow \dim(\mathcal{O}_{P, q}^\wedge/\mathfrak p) = 1$
，因为 $\mathcal{I}_y^\wedge$ 由 $1$ 个元素生成。
结合这两点，可以把 $(\mathcal{F}_n)$ 换成对象
$(\mathcal{H}_n)$；它属于 $\textit{Coh}(P, \mathcal{I})$，
并由引理 \ref{algebraization-lemma-improvement-formal-coherent-module-better}
给出。因此可以并且确实假设
$(\mathcal{F}_n)$ 与逆系统 $(\mathcal{F}_n'')$ 为
pro-同构；后者由凝聚 $\mathcal{O}_P$-模组成，并且满足
$\text{depth}(\mathcal{F}''_{n, q}) + \dim(\overline{\{q\}}) \geq 2$，
对所有 $q \in Q$ 都如此。

\medskip\noindent
我们将使用《态射进阶》，引理
\ref{more-morphisms-lemma-apply-proj-spec}，
并在不再另行说明的情况下使用
《态射进阶》，第 \ref{more-morphisms-section-proj-spec} 节
中的记号和结果；若至少不了解该引理的陈述，本证明便无法理解。
注意，在我们的情形中
$A = \bigoplus_{m \geq 0} \Gamma(P, \mathcal{L}^{\otimes m})$
是有限型 $k$-代数，且它的每个分次部分都是有限维
$k$-向量空间，见《概形上同调》，
引理 \ref{coherent-lemma-coherent-proper-ample}。

\medskip\noindent
由《概形上同调》，引理
\ref{coherent-lemma-inverse-systems-pullback}，
沿 $\pi|_U : U \to P$ 的拉回
$(\pi|_U^*\mathcal{F}_n)$ 是
$\textit{Coh}(U, f\mathcal{O}_U)$ 的对象，并且它与逆系统
$(\pi|_U^*\mathcal{F}_n'')$ 为 pro-同构；后者由凝聚
$\mathcal{O}_U$-模组成。
我们断言
$$
\text{depth}(\pi|_U^*\mathcal{F}''_{n, y}) + \delta_Z^Y(y) \geq 3
$$
，对所有 $y \in U \cap Y$ 都成立。由于 $Z$ 的所有点都是闭点，由引理
\ref{algebraization-lemma-discussion} 可知，
$\delta_Z^Y(y) \geq \dim(\overline{\{y\}})$ 对所有
$y \in U \cap Y$ 都成立。
令 $q \in Q$ 为 $y$ 的像。因为态射 $\pi : U \to P$
是相对维数为 $1$ 的光滑态射，所以 $y$ 或是 $\pi$
某条纤维的闭点，或是其泛点。因此
$$
\text{depth}(\pi^*\mathcal{F}''_{n, y}) + \delta_Z^Y(y)
\geq
\text{depth}(\pi^*\mathcal{F}''_{n, y}) + \dim(\overline{\{y\}}) =
\text{depth}(\mathcal{F}''_{n, q}) + \dim(\overline{\{q\}}) + 1
$$
，因为深度或维数二者之一增加 $1$。这就证明了断言。

\medskip\noindent
由引理 \ref{algebraization-lemma-cd-1-canonical} 可知，
$(\pi|_U^*\mathcal{F}_n)$ 可典范延拓到 $X$。注意
$$
M_n = \Gamma(U, \pi|_U^*\mathcal{F}_n) =
\bigoplus\nolimits_{m \in \mathbf{Z}}
\Gamma(P, \mathcal{F}_n \otimes_{\mathcal{O}_P} \mathcal{L}^{\otimes m})
$$
典范地成为分次 $A$-模，见《态射进阶》，等式
(\ref{more-morphisms-equation-cohomology-torsor})。
由《性质》，引理 \ref{properties-lemma-quasi-coherent-quasi-affine}，
有 $\pi|_U^*\mathcal{F}_n = \widetilde{M_n}|_U$。
因此可以应用引理
\ref{algebraization-lemma-Gm-equivariant-extend-canonically}，
得到有限分次 $A$-模 $N$，使 $(M_n)$ 与
$(N/I^nN)$ 在分次 $A$-模模去 $A_+$-挠模所得的范畴中
为 pro-同构。令 $\mathcal{F}$ 为凝聚 $\mathcal{O}_P$-模，
并要求它与 $N$ 相伴，见《概形上同调》，命题
\ref{coherent-proposition-coherent-modules-on-proj-general}。
同一命题还说明 $(\mathcal{F}/\mathcal{I}^n\mathcal{F})$
与 $(\mathcal{F}_n)$ 为 pro-同构。由于二者都是
$\textit{Coh}(P, \mathcal{I})$ 的对象，结论由引理
\ref{algebraization-lemma-recognize-formal-coherent-modules} 得出。
\end{proof}

\begin{example}
\label{algebraization-example-lefschetz-existence}
设 $k$ 为域，$X$ 为 $k$ 上的固有簇。
设 $Y \subset X$ 为有效 Cartier 除子，且
$\mathcal{O}_X(Y)$ 丰沛；记
$\mathcal{I} \subset \mathcal{O}_X$ 为相应的理想层。
设 $(\mathcal{E}_n)$ 是 $\textit{Coh}(X, \mathcal{I})$
的一个对象，且 $\mathcal{E}_n$ 有限局部自由。
下面列出命题 \ref{algebraization-proposition-lefschetz-existence} 的若干特殊情形。
\begin{enumerate}
\item 若 $X$ 是曲线或曲面，则得不到任何新结论。
\item 若 $X$ 是 Cohen-Macaulay 三维簇，则
$(\mathcal{E}_n)$ 是某个凝聚 $\mathcal{O}_X$-模
$\mathcal{E}$ 的完备化。
\item 更一般地，若 $\dim(X) \geq 3$ 且 $X$ 满足 $(S_3)$，
则 $(\mathcal{E}_n)$ 是某个凝聚 $\mathcal{O}_X$-模
$\mathcal{E}$ 的完备化。
\end{enumerate}
当然，若 $\mathcal{E}$ 存在，则 $\mathcal{E}$ 在 $Y$ 的某个开邻域上
有限局部自由。
\end{example}

\begin{proposition}
\label{algebraization-proposition-lefschetz-equivalence}
设 $k$ 为域，$X$ 为 $k$ 上的固有概形。
设 $\mathcal{L}$ 为丰沛的可逆 $\mathcal{O}_X$-模，并取
$s \in \Gamma(X, \mathcal{L})$。令 $Y = Z(s)$
为 $s$ 的零点概形，记 $\mathcal{I} \subset \mathcal{O}_X$
为相应的理想层。令 $\mathcal{V}$ 为 $X$ 中所有包含 $Y$
的开子概形所成的集合，并按反包含关系排序。
假设对所有 $x \in X \setminus Y$ 都有
$$
\text{depth}(\mathcal{O}_{X, x}) + \dim(\overline{\{x\}}) > 2
$$
则完备化函子
$$
\colim_\mathcal{V}
\textit{Coh}(\mathcal{O}_V)
\longrightarrow
\textit{Coh}(X, \mathcal{I})
$$
在有限局部自由对象所成的满子范畴之间给出等价。
\end{proposition}

\begin{proof}
为证明充分忠实，只须证明映射
$$
\colim_\mathcal{V} \Gamma(V, \mathcal{L}^{\otimes m})
\longrightarrow
\lim \Gamma(Y_n, \mathcal{L}^{\otimes m}|_{Y_n})
$$
对所有 $m$ 为同构，见引理
\ref{algebraization-lemma-completion-fully-faithful-general}。
这由引理 \ref{algebraization-lemma-lefschetz-addendum} 推出。

\medskip\noindent
本质满性。设 $(\mathcal{F}_n)$ 是
$\textit{Coh}(X, \mathcal{I})$ 的有限局部自由对象。
则对 $y \in Y$，有
$\mathcal{F}_y^\wedge = \lim \mathcal{F}_{n, y}$，且它是有限自由
$\mathcal{O}_{X, y}^\wedge$-模。
设 $\mathfrak p \subset \mathcal{O}_{X, y}^\wedge$
为满足 $\mathfrak p \not \in V(\mathcal{I}_y^\wedge)$ 的素理想。
那么 $\mathfrak p$ 位于某个素理想
$\mathfrak p_0 \subset \mathcal{O}_{X, y}$ 之上，而后者对应于
特化 $x \leadsto y$，其中 $x \not \in Y$。
由《局部上同调》，引理
\ref{local-cohomology-lemma-change-completion} 以及一些维数理论
（见《簇》，第 \ref{varieties-section-algebraic-schemes} 节），有
$$
\text{depth}((\mathcal{O}_{X, y}^\wedge)_\mathfrak p) +
\dim(\mathcal{O}_{X, y}^\wedge/\mathfrak p) =
\text{depth}(\mathcal{O}_{X, x}) +
\dim(\overline{\{x\}}) - \dim(\overline{\{y\}})
$$
因此，我们的假设蕴含命题
\ref{algebraization-proposition-lefschetz-existence} 的假设均成立，
从而 $(\mathcal{F}_n)$ 是某个凝聚
$\mathcal{O}_X$-模 $\mathcal{F}$ 的完备化。
于是 $\mathcal{F}_y$ 对所有 $y \in Y$ 都是有限自由的；
因此 $\mathcal{F}$ 在某个开邻域 $V$ 上有限局部自由，
且该邻域包含 $Y$。
证明完毕。
\end{proof}
% Generated mechanically from the frozen curves.tex target.
% Do not edit this generated file; edit the bound locale source instead.

% OK, start here.
%

\setcounter{chapter}{52}
\chapter{代数曲线}



\phantomsection
\label{curves-section-phantom}


\section{引言}
\label{curves-section-introduction}

\noindent
本章发展代数曲线理论的一些内容。
关于复数域上代数曲线的一部参考书是 \cite{ACGH}。

\medskip\noindent
已有知识。除一般的代数几何之外，我们还已经证明了
若干专门关于代数曲线的结果。现列举如下。
\begin{enumerate}
\item 我们在《簇》，第 \ref{varieties-section-dimension-one} 节中讨论了
$1$ 维 Noether 概形的仿射开集及其上的丰沛可逆层。
\item 我们在《簇》，第 \ref{varieties-section-curves} 节中看到，
一条曲线要么仿射，要么投射。
\item 我们在《簇》，第 \ref{varieties-section-divisors-curves} 节中讨论了
固有曲线上局部自由模的次数。
\item 我们在《曲线的 Picard 概形》，
第 \ref{pic-section-introduction} 节中讨论了代数闭域上
非奇异投射曲线的 Picard 概形。
\end{enumerate}





\section{曲线与函数域}
\label{curves-section-curves-function-fields}

\noindent
本节针对曲线的情形，详细阐述《簇》，
第 \ref{varieties-section-varieties-rational-maps} 节中的结果。

\begin{lemma}
\label{curves-lemma-extend-over-dvr}
设 $k$ 为域，$X$ 为曲线，$Y$ 为固有簇。
设 $U \subset X$ 为非空开集，$f : U \to Y$ 为态射。
若 $x \in X$ 是闭点，且 $\mathcal{O}_{X, x}$
为离散赋值环，则存在开集 $U \subset U' \subset X$，
它包含 $x$，并且存在簇态射 $f' : U' \to Y$，后者延拓 $f$。
\end{lemma}

\begin{proof}
这是《态射》，引理 \ref{morphisms-lemma-extend-across}
的一个特殊情形。
\end{proof}

\begin{lemma}
\label{curves-lemma-extend-over-normal-curve}
设 $k$ 为域，$X$ 为正规曲线，$Y$ 为固有簇。
从 $X$ 到 $Y$ 的有理映射所成的集合，等于态射
$X \to Y$ 所成的集合。
\end{lemma}

\begin{proof}
从 $X$ 到 $Y$ 的有理映射可由引理
\ref{curves-lemma-extend-over-dvr} 延拓为态射 $X \to Y$，
因为每个局部环都是离散赋值环
（例如见《簇》，引理 \ref{varieties-lemma-regular-point-on-curve}）。
反之，若两个态射 $f,g: X \to Y$ 作为有理映射等价，
则由《态射》，引理 \ref{morphisms-lemma-equality-of-morphisms}
可知 $f = g$。
\end{proof}

\begin{lemma}
\label{curves-lemma-flat}
设 $k$ 为域，$f : X \to Y$ 为 $k$ 上曲线之间的非常值态射。
若 $Y$ 正规，则 $f$ 平坦。
\end{lemma}

\begin{proof}
取 $x \in X$，其像为 $y \in Y$。那么 $\mathcal{O}_{Y, y}$
或为域，或为离散赋值环
（《簇》，引理 \ref{varieties-lemma-regular-point-on-curve}）。
由于 $f$ 非常值，它是支配的（因为它必将 $X$ 的泛点映到
$Y$ 的泛点）。这说明
$\mathcal{O}_{Y, y} \to \mathcal{O}_{X, x}$ 为单射
（《态射》，引理 \ref{morphisms-lemma-dominant-between-integral}）。
因此，$\mathcal{O}_{X, x}$ 作为 $\mathcal{O}_{Y, y}$-模无挠；
进而由《代数进阶》，引理
\ref{more-algebra-lemma-valuation-ring-torsion-free-flat}，
$\mathcal{O}_{X, x}$ 作为 $\mathcal{O}_{Y, y}$-模是平坦的。
\end{proof}

\begin{lemma}
\label{curves-lemma-finite}
设 $k$ 为域，$f : X \to Y$ 为 $k$ 上概形之间的态射。假设
\begin{enumerate}
\item $Y$ 在 $k$ 上分离；
\item $X$ 是维数 $\leq 1$ 的 $k$ 上固有概形；
\item $f(Z)$ 至少含有两个点，其中 $Z \subset X$
是维数为 $1$ 的任一不可约分支。
\end{enumerate}
则 $f$ 有限。
\end{lemma}

\begin{proof}
由《态射》，引理
\ref{morphisms-lemma-image-proper-scheme-closed}，态射 $f$ 固有。
因此 $f(X)$ 是闭集，闭点的像也是闭点。
设 $y \in Y$ 是 $X$ 中某个闭点的像。
那么由条件 (3)，$f^{-1}(\{y\})$ 是 $X$ 的闭子集，
并且不包含任何维数为 $1$ 的不可约分支的泛点。因此
$f^{-1}(\{y\})$ 有限。进而由《态射进阶》，引理
\ref{more-morphisms-lemma-proper-finite-fibre-finite-in-neighbourhood}，
$f$ 在 $y$ 的某个开邻域上有限
（若 $Y$ 为 Noether 概形，则可以使用较容易的
《概形上同调》，引理
\ref{coherent-lemma-proper-finite-fibre-finite-in-neighbourhood}）。
上面已经看到这样的点 $y$ 足够多，故证明完毕。
\end{proof}

\begin{lemma}
\label{curves-lemma-extend-to-completion}
设 $k$ 为域，$X \to Y$ 为簇之间的态射，其中
$Y$ 固有而 $X$ 为曲线。存在分解
$X \to \overline{X} \to Y$，其中
$X \to \overline{X}$ 是开浸入，且
$\overline{X}$ 是投射曲线。
\end{lemma}

\begin{proof}
这由引理 \ref{curves-lemma-extend-over-dvr} 以及《簇》，
引理 \ref{varieties-lemma-reduced-dim-1-projective-completion} 立即可得。
\end{proof}

\noindent
下面是本节的主要定理。若簇之间的态射
$f : X \to Y$ 的像 $f(X)$ 只由一个点 $y$ 组成，且该点属于 $Y$，
则称它为{\it 常值态射}。在这种情形下，$y$ 是 $Y$ 的闭点
（因为 $X$ 的闭点之像是 $Y$ 的闭点）。

\begin{theorem}
\label{curves-theorem-curves-rational-maps}
设 $k$ 为域。下列范畴典范等价：
\begin{enumerate}
\item 有限生成域扩张 $K/k$ 所成的范畴，其中扩张的超越次数为 $1$。
\item 以曲线为对象、以支配有理映射为态射的范畴。
\item 以正规投射曲线为对象、以非常值态射为态射的范畴。
\item 以非奇异投射曲线为对象、以非常值态射为态射的范畴。
\item 以正则投射曲线为对象、以非常值态射为态射的范畴。
\item 以正规固有曲线为对象、以非常值态射为态射的范畴。
\end{enumerate}
\end{theorem}

\begin{proof}
范畴 (1) 与 (2) 之间的等价，是《簇》，
定理 \ref{varieties-theorem-varieties-rational-maps}
中等价的限制。具体而言，一个簇为曲线，当且仅当其函数域的
超越次数为 $1$；例如见《簇》，
引理 \ref{varieties-lemma-dimension-locally-algebraic}。

\medskip\noindent
范畴 (3)、(4)、(5) 与 (6) 相同。首先，“正则”与“非奇异”
是同义词，见《性质》，定义 \ref{properties-definition-regular}。
对 Noether 的 $1$ 维概形而言，正规与正则是同一回事
（《性质》，引理 \ref{properties-lemma-regular-normal} 及
\ref{properties-lemma-normal-dimension-1-regular}）。
曲线的情形见《簇》，引理
\ref{varieties-lemma-regular-point-on-curve}。
因此 (3) 与 (5) 相同。最后，由《簇》，引理
\ref{varieties-lemma-dim-1-proper-projective}，(6) 与 (3) 相同。

\medskip\noindent
若 $f : X \to Y$ 是非奇异投射曲线之间的非常值态射，
则 $f$ 把泛点 $\eta$（即 $X$ 的泛点）映到泛点
$\xi$（即 $Y$ 的泛点）。
因此得到范畴 (1) 中的态射
$k(Y) = \mathcal{O}_{Y, \xi} \to \mathcal{O}_{X, \eta} = k(X)$
。若两个态射 $f,g: X \to Y$ 给出同一态射
$k(Y) \to k(X)$，则由 (1) 与 (2) 之间的等价，
$f$ 与 $g$ 作为有理映射等价；进而由引理
\ref{curves-lemma-extend-over-normal-curve}，有 $f=g$。
反之，设范畴 (1) 中给定映射
$k(Y) \to k(X)$。则得到态射 $U \to Y$，其中
$U \subset X$ 是某个非空开集。由引理 \ref{curves-lemma-extend-over-dvr}，
它可延拓到整个 $X$，从而得到范畴 (5) 中的态射。
所以存在充分忠实函子 (5)$\to$(1)。

\medskip\noindent
为完成证明，还须说明 (1) 中的每个 $K/k$
都是某条正规投射曲线的函数域。
我们已经知道 $K = k(X)$，其中 $X$ 是某条曲线。
将 $X$ 换成它的正规化
（这是一个与 $X$ 双有理的簇）后，可以假设 $X$ 正规
（《簇》，引理
\ref{varieties-lemma-normalization-locally-algebraic}）。
然后依照《簇》，引理
\ref{varieties-lemma-reduced-dim-1-projective-completion}，
选取 $X \to \overline{X}$，使
$\overline{X} \setminus X = \{x_1, \ldots, x_n\}$。
由于 $X$ 正规，而且每个局部环
$\mathcal{O}_{\overline{X}, x_i}$ 都正规，可知
$\overline{X}$ 如所需是一条正规投射曲线。
（备注：也可以先用《簇》，引理
\ref{varieties-lemma-dim-1-projective-completion} 作紧化，
再用《簇》，引理
\ref{varieties-lemma-normalization-locally-algebraic} 作正规化。
这样便可避免使用稍显微妙的《态射》，引理
\ref{morphisms-lemma-relative-normalization-normal-codim-1}。）
\end{proof}

\begin{definition}
\label{curves-definition-normal-projective-model}
设 $k$ 为域，$X$ 为曲线。
所谓{\it $X$ 的非奇异投射模型}，是指一个二元组
$(Y, \varphi)$，其中 $Y$ 是非奇异投射曲线，而
$\varphi : k(X) \to k(Y)$ 是函数域的同构。
\end{definition}

\noindent
由定理 \ref{curves-theorem-curves-rational-maps}，非奇异投射模型
在唯一同构意义下确定。因此我们常说``该非奇异投射模型''。
记号中通常省略 $\varphi$。
警告：$Y$ 未必在 $k$ 上光滑；但引理
\ref{curves-lemma-nonsingular-model-smooth} 表明，这种情形只能发生在正特征中。

\begin{lemma}
\label{curves-lemma-nonsingular-model-smooth}
设 $k$ 为域，$X$ 为曲线，且 $Y$ 是 $X$ 的非奇异投射模型。
若 $k$ 完美，则 $Y$ 是光滑投射曲线。
\end{lemma}

\begin{proof}
例如见《簇》，引理 \ref{varieties-lemma-regular-point-on-curve}。
\end{proof}

\begin{lemma}
\label{curves-lemma-smooth-models}
设 $k$ 为域，$X$ 为 $k$ 上的几何不可约曲线。
对域扩张 $K/k$，以 $Y_K$ 表示 $(X_K)_{red}$ 的一个非奇异投射模型。
\begin{enumerate}
\item 若 $X$ 固有，则 $Y_K$ 是 $X_K$ 的正规化。
\item 存在有限纯不可分扩张 $K/k$，使得 $Y_K$ 光滑。
\item 每当 $Y_K$ 光滑\footnote{甚至只需几何约化。}时，
都有 $H^0(Y_K, \mathcal{O}_{Y_K}) = K$。
\item 给定域的交换图
$$
\xymatrix{
\Omega & K' \ar[l] \\
K \ar[u] & k \ar[l] \ar[u]
}
$$
使得 $Y_K$ 与 $Y_{K'}$ 光滑，则
$Y_\Omega = (Y_K)_\Omega = (Y_{K'})_\Omega$。
\end{enumerate}
\end{lemma}

\begin{proof}
设 $X'$ 是 $X$ 的一个非奇异投射模型。则 $X'$ 与 $X$
有彼此同构的非空开子概形。特别地，$X'$ 几何不可约，
正如 $X$ 一样（省略若干细节）。所以可以假设 $X$ 投射。

\medskip\noindent
假设 $X$ 固有。则 $X_K$ 固有，因而其正规化
$(X_K)^\nu$ 也固有，因为它在一个固有概形上有限
（《簇》，引理 \ref{varieties-lemma-normalization-locally-algebraic}，
以及《态射》，引理 \ref{morphisms-lemma-finite-proper} 与
\ref{morphisms-lemma-composition-proper}）。
另一方面，$X_K$ 不可约，因为 $X$ 几何不可约。
因此 $X_K^\nu$ 固有、正规且不可约，并与 $(X_K)_{red}$ 双有理。
由于固有曲线是投射的（《簇》，引理
\ref{varieties-lemma-dim-1-proper-projective}），这就证明了 (1)。

\medskip\noindent
证明 (2)。由于 $X$ 固有并且已有 (1)，可应用《簇》，引理
\ref{varieties-lemma-finite-extension-geometrically-normal}，
找到有限纯不可分扩张 $K/k$，使得 $Y_K$ 几何正规。
由于对曲线而言正规与正则等价（《性质》，引理
\ref{properties-lemma-normal-dimension-1-regular}），
$Y_K$ 于是几何正则。由《簇》，引理
\ref{varieties-lemma-geometrically-regular-smooth}，$Y$ 是光滑簇。

\medskip\noindent
若 $Y_K$ 几何约化，则 $Y_K$ 几何整
（《簇》，引理 \ref{varieties-lemma-geometrically-integral}），
且由《簇》，引理
\ref{varieties-lemma-regular-functions-proper-variety} 可知
$H^0(Y_K, \mathcal{O}_{Y_K}) = K$。
因为光滑簇几何约化（甚至几何正则，见《簇》，引理
\ref{varieties-lemma-geometrically-regular-smooth}），这就证明了 (3)。

\medskip\noindent
若 $Y_K$ 光滑，则对每个扩张 $\Omega/K$，基变换
$(Y_K)_\Omega$ 在 $\Omega$ 上光滑
（《态射》，引理 \ref{morphisms-lemma-base-change-smooth}）。
因此显然有 $Y_\Omega = (Y_K)_\Omega$。这就证明了 (4)。
\end{proof}






\section{线性系}
\label{curves-section-linear-series}

\noindent
我们采用如下方式定义线性系，这与经典叙述略有不同
（见备注 \ref{curves-remark-classical-linear-series}）。

\begin{definition}
\label{curves-definition-linear-series}
设 $k$ 为域，$X$ 为维数 $\leq 1$ 的 $k$ 上固有概形。
设 $d \geq 0$ 且 $r \geq 0$。
所谓{\it 次数为 $d$、维数为 $r$ 的线性系}，是指二元组
$(\mathcal{L}, V)$，其中 $\mathcal{L}$ 是可逆
$\mathcal{O}_X$-模，次数为 $d$
（《簇》，定义 \ref{varieties-definition-degree-invertible-sheaf}），
而 $V \subset H^0(X, \mathcal{L})$ 是 $k$-向量子空间，
其维数为 $r + 1$。简言之，我们称 $(\mathcal{L}, V)$
是一个{\it $\mathfrak g^r_d$}，定义在 $X$ 上。
\end{definition}

\noindent
我们主要在 $X$ 为非奇异固有曲线时使用这个定义。
事实上，上述定义只是推广经典 $\mathfrak g^r_d$ 定义的一种方式。
例如，若 $X$ 是固有曲线，可以允许 $\mathcal{L}$ 是无挠凝聚
$\mathcal{O}_X$-模，秩为 $1$，从而推广线性系的概念。
在非奇异曲线上，每个无挠凝聚模都是局部自由的，
所以这与我们对非奇异固有曲线所用的概念一致。

\medskip\noindent
下述引理说明固有非奇异曲线上线性系的几何意义。

\begin{lemma}
\label{curves-lemma-linear-series}
设 $k$ 为域，$X$ 为 $k$ 上的非奇异固有曲线。
设 $(\mathcal{L}, V)$ 是一个 $\mathfrak g^r_d$，定义在 $X$ 上。
则存在态射
$$
\varphi : X \longrightarrow \mathbf{P}^r_k = \text{Proj}(k[T_0, \ldots, T_r])
$$
，它是 $k$ 上簇的态射；并存在映射
$\alpha : \varphi^*\mathcal{O}_{\mathbf{P}^r_k}(1) \to \mathcal{L}$
，使得 $\varphi^*T_0, \ldots, \varphi^*T_r$
被映为 $V$ 的一组基，所用映射为 $\alpha$。
\end{lemma}

\begin{proof}
设 $s_0, \ldots, s_r \in V$ 是一组 $k$-基。由于 $X$ 非奇异，
记该映射的像为 $\mathcal{L}' \subset \mathcal{L}$；该映射为
$s_0, \ldots, s_r : \mathcal{O}_X^{\oplus r + 1} \to \mathcal{L}$
。于是该像是可逆 $\mathcal{O}_X$-模；例如这可由《除子》，引理
\ref{divisors-lemma-torsion-free-over-regular-dim-1} 得到。
然后应用《构造》，引理 \ref{constructions-lemma-projective-space}，
得到态射
$$
\varphi = \varphi_{(\mathcal{L}', (s_0, \ldots, s_r))} :
X \longrightarrow \mathbf{P}^r_k
$$
，正如引理所述。
\end{proof}

\begin{lemma}
\label{curves-lemma-linear-series-trivial-existence}
设 $k$ 为域，$X$ 为维数 $\leq 1$ 的 $k$ 上固有概形。
若 $X$ 有一个 $\mathfrak g^r_d$，则 $X$ 有一个
$\mathfrak g^s_d$，对所有 $0 \leq s \leq r$ 都如此。
\end{lemma}

\begin{proof}
这是因为向量空间 $V$ 的维数为 $r + 1$，它定义在 $k$ 上；
该空间有维数为 $s + 1$ 的线性子空间，对所有
$0 \leq s \leq r$ 都如此。
\end{proof}

\begin{lemma}
\label{curves-lemma-g1d}
设 $k$ 为域，$X$ 为 $k$ 上的非奇异固有曲线。
设 $(\mathcal{L}, V)$ 是一个 $\mathfrak g^1_d$，定义在 $X$ 上。
则态射 $\varphi : X \to \mathbf{P}^1_k$
（见引理 \ref{curves-lemma-linear-series}）满足下列二者之一：
\begin{enumerate}
\item 它非常值，且次数 $\leq d$；或
\item 它分解通过 $\mathbf{P}^1_k$ 的一个闭点，并且在这种情形下
$H^0(X, \mathcal{O}_X) \not = k$。
\end{enumerate}
\end{lemma}

\begin{proof}
由引理 \ref{curves-lemma-linear-series} 可知
$\mathcal{L}' = \varphi^*\mathcal{O}_{\mathbf{P}^1_k}(1)$
有一个非零映射 $\mathcal{L}' \to \mathcal{L}$。
因此由《簇》，引理 \ref{varieties-lemma-check-invertible-sheaf-trivial}
可知 $0 \leq \deg(\mathcal{L}') \leq d$。
若 $\deg(\mathcal{L}') = 0$，则同一引理给出
$\mathcal{L}' \cong \mathcal{O}_X$；又因为有两个线性无关截面，
所以处于情形 (2)。
若 $\deg(\mathcal{L}') > 0$，则 $\varphi$ 非常值
（因为可逆模沿常值态射的拉回是平凡的）。因此
$$
\deg(\mathcal{L}') =
\deg(X/\mathbf{P}^1_k) \deg(\mathcal{O}_{\mathbf{P}^1_k}(1))
$$
，这由《簇》，引理
\ref{varieties-lemma-degree-pullback-map-proper-curves} 得到。
由于 $\mathcal{O}_{\mathbf{P}^1_k}(1)$ 的次数为 $1$，证明完成。
\end{proof}

\begin{lemma}
\label{curves-lemma-grd-inequalities}
设 $k$ 为域，$X$ 为 $k$ 上满足
$H^0(X, \mathcal{O}_X) = k$ 的固有曲线。若 $X$ 有一个
$\mathfrak g^r_d$，则 $r \leq d$。若等号成立，则
$H^1(X, \mathcal{O}_X) = 0$；换言之，$X$ 的亏格
（定义 \ref{curves-definition-genus}）为 $0$。
\end{lemma}

\begin{proof}
设 $(\mathcal{L}, V)$ 是一个 $\mathfrak g^r_d$。由于这样只会增大
$r$，可以假设 $V = H^0(X, \mathcal{L})$。选取非零元
$s \in V$。于是 $s$ 的零点概形是一个有效 Cartier 除子
$D \subset X$，并且 $\mathcal{L} = \mathcal{O}_X(D)$；还有短正合列
$$
0 \to \mathcal{O}_X \to \mathcal{L} \to \mathcal{L}|_D \to 0
$$
；见《除子》，引理 \ref{divisors-lemma-characterize-OD} 及
备注 \ref{divisors-remark-ses-regular-section}。
由《簇》，引理 \ref{varieties-lemma-degree-effective-Cartier-divisor}，
有 $\deg(D) = \deg(\mathcal{L}) = d$。
由于 $D$ 是 Artin 概形，有
$\mathcal{L}|_D \cong \mathcal{O}_D$\footnote{在我们的情形下，
这由《除子》，引理
\ref{divisors-lemma-finite-trivialize-invertible-upstairs}
得到，因为 $D \to \Spec(k)$ 有限。}。因此
$$
\dim_k H^0(D, \mathcal{L}|_D) = \dim_k H^0(D, \mathcal{O}_D) = \deg(D) = d
$$
另一方面，依假设
$\dim_k H^0(X, \mathcal{O}_X) = 1$，且
$\dim H^0(X, \mathcal{L}) = r + 1$。
因此 $r + 1 \leq 1 + d$，亦即引理所述的 $r \leq d$。

\medskip\noindent
假设等号成立。则
$H^0(X, \mathcal{L}) \to H^0(X, \mathcal{L}|_D)$ 满射。
若已知 $H^1(X, \mathcal{L})$ 为零，则由上同调长正合列可推出
$H^1(X, \mathcal{O}_X)$ 为零，证明便告完成。
我们的策略是把 $\mathcal{L}$ 换成一个上同调消没的高次幂。
由于 $\mathcal{L}|_D$ 是平凡可逆模（见上），可以找到截面
$t$，它取值于 $\mathcal{L}$，且其在 $D$ 上的限制生成
$\mathcal{L}|_D$。考虑乘法映射
$$
\mu :
H^0(X, \mathcal{L}) \otimes_k H^0(X, \mathcal{L})
\longrightarrow
H^0(X, \mathcal{L}^{\otimes 2})
$$
以及短正合列
$$
0 \to \mathcal{L} \xrightarrow{s}
\mathcal{L}^{\otimes 2} \to \mathcal{L}^{\otimes 2}|_D \to 0
$$
由于 $H^0(\mathcal{L}) \to H^0(\mathcal{L}|_D)$ 满射，且
$t$ 映到 $\mathcal{L}|_D$ 的一个平凡化，可知
$\mu(H^0(X, \mathcal{L}) \otimes t)$ 给出
$H^0(X, \mathcal{L}^{\otimes 2})$ 的一个子空间，并满射到
$\mathcal{L}^{\otimes 2}|_D$ 的整体截面。因此
$$
\dim H^0(X, \mathcal{L}^{\otimes 2}) = r + 1 + d = 2r + 1 =
\deg(\mathcal{L}^{\otimes 2}) + 1
$$
。所以 $\mathcal{L}^{\otimes 2}$ 与 $\mathcal{L}$ 具有同一性质，
亦即整体截面空间的维数等于次数加一。因为 $\mathcal{L}$ 丰沛
（《簇》，引理 \ref{varieties-lemma-ample-curve}），
存在某个 $n_0$，使 $\mathcal{L}^{\otimes n}$ 的 $H^1$
对所有 $n \geq n_0$ 都消没
（《概形的上同调》，引理 \ref{coherent-lemma-coherent-proper-ample}）。
因此，把上述论证应用于 $\mathcal{L}^{\otimes n}$，
其中 $n = 2^m$ 且 $m$ 足够大，便得引理成立。
\end{proof}

\begin{remark}[经典定义]
\label{curves-remark-classical-linear-series}
设 $X$ 为代数闭域 $k$ 上的光滑投射曲线。
称两个有效 Cartier 除子 $D, D' \subset X$ {\it 线性等价}，
当且仅当有 $\mathcal{O}_X(D) \cong \mathcal{O}_X(D')$，
这里把二者视为 $\mathcal{O}_X$-模。
由于 $\Pic(X) = \text{Cl}(X)$
（《除子》，引理 \ref{divisors-lemma-local-rings-UFD-c1-bijective}），
可知 $D$ 与 $D'$ 线性等价，当且仅当同 $D$ 与 $D'$
相伴的 Weil 除子定义了 $\text{Cl}(X)$ 的同一个元素。
给定有效 Cartier 除子 $D \subset X$，其次数为 $d$；
所谓它的{\it 完备线性系统}或{\it 完备线性系} $|D|$，
即 $D$ 的完备线性系，是由所有有效 Cartier 除子
$E \subset X$ 组成的集合，这些除子与 $D$ 线性等价。
换言之，$|D|$ 是下述态射某条纤维的闭点集合：
$$
\gamma_d :
\underline{\Hilbfunctor}^d_{X/k}
\longrightarrow
\underline{\Picardfunctor}^d_{X/k}
$$
这里取曲线上 Picard 概形，引理 \ref{pic-lemma-picard-pieces}
中对应于 $\mathcal{O}_X(D)$ 的闭点之上的纤维。
这赋予 $|D|$ 一个自然概形结构，并且事实证明
$|D| \cong \mathbf{P}^m_k$，其中
$m + 1 = h^0(\mathcal{O}_X(D))$。事实上，更典范地有
$$
|D| = \mathbf{P}(H^0(X, \mathcal{O}_X(D))^\vee)
$$
，其中 $(-)^\vee$ 表示 $k$-线性对偶，而 $\mathbf{P}$ 如《构造》，
例 \ref{constructions-example-projective-space} 中所定义。
用这种语言，$X$ 上的{\it 线性系统}或{\it 线性系}，
是可由线性方程截出的闭子簇 $L \subset |D|$。
若 $L$ 的维数为 $r$，则 $L = \mathbf{P}(V^\vee)$，其中
$V \subset H^0(X, \mathcal{O}_X(D))$ 是维数为 $r + 1$
的线性子空间。因此，经典线性系 $L \subset |D|$
对应于上述定义的线性系 $(\mathcal{O}_X(D), V)$。
\end{remark}







\section{对偶性}
\label{curves-section-duality}

\noindent
本节推导有关域上固有 $1$ 维概形的对偶化复形与对偶性的一般理论
所带来的结论。有关域上更高维固有概形的类似讨论，见《概形的对偶性》，
第 \ref{duality-section-duality-proper-over-field} 节。

\begin{lemma}
\label{curves-lemma-duality-dim-1}
设 $X$ 为维数 $\leq 1$ 的域 $k$ 上固有概形。
存在具有下列性质的对偶化复形 $\omega_X^\bullet$：
\begin{enumerate}
\item $H^i(\omega_X^\bullet)$ 仅在 $i = -1, 0$ 时非零；
\item $\omega_X = H^{-1}(\omega_X^\bullet)$
是凝聚 Cohen--Macaulay 模，其支撑是所有维数为 $1$ 的不可约分支；
\item 对闭点 $x \in X$，模 $H^0(\omega_{X, x}^\bullet)$ 非零，
当且仅当下列二者之一成立：
\begin{enumerate}
\item $\dim(\mathcal{O}_{X, x}) = 0$；或
\item $\dim(\mathcal{O}_{X, x}) = 1$
且 $\mathcal{O}_{X, x}$ 不是 Cohen--Macaulay 的；
\end{enumerate}
\item 对 $K \in D_\QCoh(\mathcal{O}_X)$ 有函子性同构\footnote{由
Yoneda 引理，这一性质在唯一同构意义下刻画
$\omega_X^\bullet$ 作为 $D_\QCoh(\mathcal{O}_X)$ 中的对象。由于
$\omega_X^\bullet$ 属于 $D^b_{\textit{Coh}}(\mathcal{O}_X)$，
事实上只需考虑 $K \in D^b_{\textit{Coh}}(\mathcal{O}_X)$。}
$$
\Ext^i_X(K, \omega_X^\bullet) = \Hom_k(H^{-i}(X, K), k)
$$
，并与平移及特异三角形相容；
\item 有函子性同构
$\Hom(\mathcal{F}, \omega_X) = \Hom_k(H^1(X, \mathcal{F}), k)$
，其中 $\mathcal{F}$ 是 $X$ 上的拟凝聚模；
\item 若 $X \to \Spec(k)$ 光滑且相对维数为 $1$，则
$\omega_X \cong \Omega_{X/k}$。
\end{enumerate}
\end{lemma}

\begin{proof}
以 $f : X \to \Spec(k)$ 表示结构态射。先取相对对偶化复形
$$
\omega_X^\bullet = \omega_{X/k}^\bullet = a(\mathcal{O}_{\Spec(k)})
$$
，如《概形的对偶性》，备注
\ref{duality-remark-relative-dualizing-complex} 中所述。
性质 (4) 由构造即得，因为 $a$ 是
$f_* : D_\QCoh(\mathcal{O}_X) \to D(\mathcal{O}_{\Spec(k)})$ 的右伴随。
由于 $f$ 固有，依定义有
$f^!(\mathcal{O}_{\Spec(k)}) = a(\mathcal{O}_{\Spec(k)})$；
见《概形的对偶性》，第 \ref{duality-section-upper-shriek} 节。
因此 $\omega_X^\bullet$ 与 $\omega_X$ 正如《概形的对偶性》，
例 \ref{duality-example-proper-over-local} 以及
例 \ref{duality-example-equidimensional-over-field} 中所述。
(1) 与 (2) 由《概形的对偶性》，引理
\ref{duality-lemma-vanishing-good-dualizing} 得到。
对闭点 $x \in X$，$\omega_{X, x}^\bullet$ 是
$\mathcal{O}_{X, x}$ 上的正规化对偶化复形；见《概形的对偶性》，
引理 \ref{duality-lemma-good-dualizing-normalized}。
于是 (3) 由《对偶化复形》，引理 \ref{dualizing-lemma-apply-CM} 得到。
对凝聚 $\mathcal{F}$，(5) 由《概形的对偶性》，引理
\ref{duality-lemma-dualizing-module-proper-over-A} 得到；一般情形则把
(4) 展开，并取 $K = \mathcal{F}[0]$ 与 $i = -1$。
(6) 由《概形的对偶性》，引理 \ref{duality-lemma-smooth-proper} 得到。
\end{proof}

\begin{lemma}
\label{curves-lemma-duality-dim-1-CM}
设 $X$ 为域 $k$ 上的固有概形，它是 Cohen--Macaulay 的，
并且等维且维数为 $1$。引理 \ref{curves-lemma-duality-dim-1} 中的模
$\omega_X$ 具有下列性质：
\begin{enumerate}
\item $\omega_X$ 是 $X$ 上的对偶化模
（《概形的对偶性》，第 \ref{duality-section-dualizing-module} 节）；
\item $\omega_X$ 是凝聚 Cohen--Macaulay 模，其支撑为 $X$；
\item 有函子性同构
$\Ext^i_X(K, \omega_X[1]) = \Hom_k(H^{-i}(X, K), k)$
，它们对 $K \in D_\QCoh(X)$ 与平移相容；
\item 有函子性同构
$\Ext^{1 + i}(\mathcal{F}, \omega_X) = \Hom_k(H^{-i}(X, \mathcal{F}), k)$
，其中 $\mathcal{F}$ 是 $X$ 上的拟凝聚模。
\end{enumerate}
\end{lemma}

\begin{proof}
回顾引理 \ref{curves-lemma-duality-dim-1} 的证明，其中
$\omega_X$ 正如《概形的对偶性》，例
\ref{duality-example-proper-over-local} 中所述，因而是对偶化模。
其余断言由引理 \ref{curves-lemma-duality-dim-1} 以及等式
$\omega_X^\bullet = \omega_X[1]$ 得到；后者成立是因为
$X$ 为 Cohen--Macaulay 概形（《概形的对偶性》，引理
\ref{duality-lemma-dualizing-module-CM-scheme}）。
\end{proof}

\begin{remark}
\label{curves-remark-rework-duality-locally-free}
设 $X$ 为维数 $\leq 1$ 的域 $k$ 上固有概形。
设 $\omega_X^\bullet$ 与 $\omega_X$ 如引理
\ref{curves-lemma-duality-dim-1} 中所定义。若 $\mathcal{E}$ 是有限局部自由
$\mathcal{O}_X$-模，其对偶为 $\mathcal{E}^\vee$，则有典范同构
$$
\Hom_k(H^{-i}(X, \mathcal{E}), k) =
H^i(X, \mathcal{E}^\vee \otimes_{\mathcal{O}_X}^\mathbf{L} \omega_X^\bullet)
$$
。这由上述引理及《上同调》，引理
\ref{cohomology-lemma-dual-perfect-complex} 得到。
若 $X$ 是 Cohen--Macaulay 的，并且等维且维数为 $1$，则有典范同构
$$
\Hom_k(H^{-i}(X, \mathcal{E}), k) =
H^{1 + i}(X, \mathcal{E}^\vee \otimes_{\mathcal{O}_X} \omega_X)
$$
；这由引理 \ref{curves-lemma-duality-dim-1-CM} 得到。特别地，若
$\mathcal{L}$ 是可逆 $\mathcal{O}_X$-模，则有
$$
\dim_k H^0(X, \mathcal{L}) =
\dim_k H^1(X, \mathcal{L}^{\otimes -1} \otimes_{\mathcal{O}_X} \omega_X)
$$
以及
$$
\dim_k H^1(X, \mathcal{L}) =
\dim_k H^0(X, \mathcal{L}^{\otimes -1} \otimes_{\mathcal{O}_X} \omega_X)
$$
\end{remark}

\noindent
下面对对偶化复形作一次合理性检验。

\begin{lemma}
\label{curves-lemma-sanity-check-duality}
设 $X$ 为维数 $\leq 1$ 的域 $k$ 上固有概形。
设 $\omega_X^\bullet$ 与 $\omega_X$ 如引理
\ref{curves-lemma-duality-dim-1} 中所定义。
\begin{enumerate}
\item 若 $X \to \Spec(k)$ 分解为
$X \to \Spec(k') \to \Spec(k)$，其中 $k'$ 是某个域，则 $\omega_X^\bullet$ 与
$\omega_X$ 满足性质 (4)、(5)、(6)，其中把 $k$ 换成 $k'$。
\item 若 $K/k$ 是域扩张，则 $\omega_X^\bullet$ 与 $\omega_X$
在基变换 $X_K$ 上的拉回，正如引理
\ref{curves-lemma-duality-dim-1} 对态射 $X_K \to \Spec(K)$ 所述。
\end{enumerate}
\end{lemma}

\begin{proof}
以 $f : X \to \Spec(k)$ 表示结构态射。断言 (1) 的确切含义是：
$\omega_X^\bullet$ 与 $\omega_X$ 正如引理
\ref{curves-lemma-duality-dim-1} 对态射 $f' : X \to \Spec(k')$ 所定义。
在引理 \ref{curves-lemma-duality-dim-1} 的证明中，我们取
$\omega_X^\bullet = a(\mathcal{O}_{\Spec(k)})$，其中 $a$ 是
《概形的对偶性》，引理
\ref{duality-lemma-twisted-inverse-image} 对 $f$ 所给函子的右伴随。
因此须证明
$a(\mathcal{O}_{\Spec(k)}) \cong a'(\mathcal{O}_{\Spec(k)})$
，其中 $a'$ 是《概形的对偶性》，引理
\ref{duality-lemma-twisted-inverse-image} 对 $f'$ 所给函子的右伴随。
由于 $k' \subset H^0(X, \mathcal{O}_X)$，可知 $k'/k$ 是有限扩张
（《概形的上同调》，引理
\ref{coherent-lemma-proper-over-affine-cohomology-finite}）。
由伴随的唯一性，有 $a = a' \circ b$，其中 $b$ 是《概形的对偶性》，
引理 \ref{duality-lemma-twisted-inverse-image} 对
$g : \Spec(k') \to \Spec(k)$ 所给函子的右伴随。
换言之，有 $f^! = (f')^! \circ g^!$。
因此只需证明 $\Hom_k(k', k) \cong k'$ 是一个 $k'$-模同构；
见《概形的对偶性》，例
\ref{duality-example-affine-twisted-inverse-image}。
这是因为二者都是 $k'$-向量空间，且维数相同（即维数 $1$）。

\medskip\noindent
证明 (2)。这是因为由《概形的对偶性》，引理
\ref{duality-lemma-more-base-change}，$a$ 满足基变换。
参见《概形的对偶性》，备注
\ref{duality-remark-relative-dualizing-complex} 中的讨论。
\end{proof}

\begin{lemma}
\label{curves-lemma-closed-immersion-dim-1-CM}
设 $X$ 为维数 $\leq 1$ 的域 $k$ 上固有概形。
设 $i : Y \to X$ 为闭浸入，并设
$\omega_X^\bullet$、$\omega_X$、$\omega_Y^\bullet$、$\omega_Y$
如引理 \ref{curves-lemma-duality-dim-1} 中所定义。则
\begin{enumerate}
\item $\omega_Y^\bullet = R\SheafHom(\mathcal{O}_Y, \omega_X^\bullet)$；
\item $\omega_Y = \SheafHom(\mathcal{O}_Y, \omega_X)$，且
$i_*\omega_Y = \SheafHom_{\mathcal{O}_X}(i_*\mathcal{O}_Y, \omega_X)$。
\end{enumerate}
\end{lemma}

\begin{proof}
以 $g : Y \to \Spec(k)$ 与 $f : X \to \Spec(k)$ 表示结构态射。
于是 $g = f \circ i$。以 $a, b, c$ 分别表示《概形的对偶性》，
引理 \ref{duality-lemma-twisted-inverse-image} 对 $f, g, i$
所给函子的右伴随。由右伴随的唯一性，有 $b = c \circ a$；
这里也用到 $Rg_* = Rf_* \circ Ri_*$。
在引理 \ref{curves-lemma-duality-dim-1} 的证明中，我们置
$\omega_X^\bullet = a(\mathcal{O}_{\Spec(k)})$ 且
$\omega_Y^\bullet = b(\mathcal{O}_{\Spec(k)})$。
因此 $\omega_Y^\bullet = c(\omega_X^\bullet)$；由《概形的对偶性》，
引理 \ref{duality-lemma-twisted-inverse-image-closed}，这蕴含 (1)。
由于 $\omega_X = H^{-1}(\omega_X^\bullet)$ 且
$\omega_Y = H^{-1}(\omega_Y^\bullet)$，可得
$\omega_Y = \SheafHom(\mathcal{O}_Y, \omega_X)$。
这又蕴含
$i_*\omega_Y = \SheafHom_{\mathcal{O}_X}(i_*\mathcal{O}_Y, \omega_X)$
；见《概形的对偶性》，引理
\ref{duality-lemma-sheaf-with-exact-support-ext}。
\end{proof}

\begin{lemma}
\label{curves-lemma-closed-subscheme-reduced-gorenstein}
设 $X$ 为域 $k$ 上的固有概形，它是 Gorenstein、约化的，
并且等维且维数为 $1$。设 $i : Y \to X$ 是等维且维数为 $1$
的约化闭子概形。设 $j : Z \to X$ 是 $X \setminus Y$
的概形论闭包。则
\begin{enumerate}
\item $Y$ 与 $Z$ 都是 Cohen--Macaulay 的；
\item 若 $\mathcal{I} \subset \mathcal{O}_X$、分别地
$\mathcal{J} \subset \mathcal{O}_X$ 是 $Y$、分别地 $Z$
在 $X$ 中的理想层，则
$$
\mathcal{I} = i_*\mathcal{I}'
\quad\text{且}\quad
\mathcal{J} = j_*\mathcal{J}'
$$
，其中 $\mathcal{I}' \subset \mathcal{O}_Z$、分别地
$\mathcal{J}' \subset \mathcal{O}_Y$ 是 $Y \cap Z$ 在
$Z$、分别地 $Y$ 中的理想层；
\item $\omega_Y = \mathcal{J}'(i^*\omega_X)$，且
$i_*(\omega_Y) = \mathcal{J}\omega_X$；
\item $\omega_Z = \mathcal{I}'(i^*\omega_X)$，且
$i_*(\omega_Z) = \mathcal{I}\omega_X$；
\item 有下列短正合列：
\begin{align*}
0 \to \omega_X \to i_*i^*\omega_X \oplus j_*j^*\omega_X \to
\mathcal{O}_{Y \cap Z} \to 0 \\
0 \to i_*\omega_Y \to \omega_X \to j_*j^*\omega_X \to 0 \\
0 \to j_*\omega_Z \to \omega_X \to i_*i^*\omega_X \to 0 \\
0 \to i_*\omega_Y \oplus j_*\omega_Z \to \omega_X \to
\mathcal{O}_{Y \cap Z} \to 0 \\
0 \to \omega_Y \to i^*\omega_X \to \mathcal{O}_{Y \cap Z} \to 0 \\
0 \to \omega_Z \to j^*\omega_X \to \mathcal{O}_{Y \cap Z} \to 0
\end{align*}
\end{enumerate}
这里 $\omega_X$、$\omega_Y$、$\omega_Z$ 如引理
\ref{curves-lemma-duality-dim-1} 中所定义。
\end{lemma}

\begin{proof}
约化的 $1$ 维 Noether 概形是 Cohen--Macaulay 的，所以 (1) 成立。
由于 $X$ 约化，在概形论意义下有 $X = Y \cup Z$。
采用《态射》，引理 \ref{morphisms-lemma-scheme-theoretic-union}
中的记号，并由该引理的断言，有短正合列
$$
0 \to \mathcal{O}_X \to
\mathcal{O}_Y \oplus \mathcal{O}_Z \to \mathcal{O}_{Y \cap Z} \to 0
$$
。由于 $\mathcal{J} = \Ker(\mathcal{O}_X \to \mathcal{O}_Z)$、
$\mathcal{J}' = \Ker(\mathcal{O}_Y \to \mathcal{O}_{Y \cap Z})$、
$\mathcal{I} = \Ker(\mathcal{O}_X \to \mathcal{O}_Y)$，以及
$\mathcal{I}' = \Ker(\mathcal{O}_Z \to \mathcal{O}_{Y \cap Z})$，
图追踪即得 (2)。
注意 $\mathcal{I} + \mathcal{J}$ 是 $Y \cap Z$ 的理想层，
并且 $\mathcal{I} \cap \mathcal{J} = 0$。因此有下列正合列：
\begin{align*}
0 \to \mathcal{O}_X \to \mathcal{O}_Y \oplus \mathcal{O}_Z \to
\mathcal{O}_{Y \cap Z} \to 0 \\
0 \to \mathcal{J} \to \mathcal{O}_X \to \mathcal{O}_Z \to 0 \\
0 \to \mathcal{I} \to \mathcal{O}_X \to \mathcal{O}_Y \to 0 \\
0 \to \mathcal{J} \oplus \mathcal{I} \to \mathcal{O}_X \to
\mathcal{O}_{Y \cap Z} \to 0 \\
0 \to \mathcal{J}' \to \mathcal{O}_Y \to \mathcal{O}_{Y \cap Z} \to 0 \\
0 \to \mathcal{I}' \to \mathcal{O}_Z \to \mathcal{O}_{Y \cap Z} \to 0
\end{align*}
由于 $X$ 是 Gorenstein 的，$\omega_X$ 是可逆
$\mathcal{O}_X$-模（《概形的对偶性》，引理
\ref{duality-lemma-gorenstein}）。
由于 $Y \cap Z$ 的维数为 $0$，有
$\omega_X|_{Y \cap Z} \cong \mathcal{O}_{Y \cap Z}$。
因此，只要证明 (3) 与 (4)，把上述短正合列与可逆模
$\omega_X$ 作张量积，便得到引理中的短正合列。
由对称性，只需证明 (3)；再由 (2)，只需证明
$i_*(\omega_Y) = \mathcal{J}\omega_X$。

\medskip\noindent
我们有
$i_*\omega_Y = \SheafHom_{\mathcal{O}_X}(i_*\mathcal{O}_Y, \omega_X)$
，这由引理 \ref{curves-lemma-closed-immersion-dim-1-CM} 得到。
再次利用 $\omega_X$ 可逆，最终只需证明
$\SheafHom_{\mathcal{O}_X}(\mathcal{O}_X/\mathcal{I}, \mathcal{O}_X)$
同构地映到 $\mathcal{J}$，其映射是在 $1$ 处取值。
换言之，$\mathcal{J}$ 是 $\mathcal{I}$ 的零化子。这由上文得到。
\end{proof}







\section{Riemann--Roch 定理}
\label{curves-section-Riemann-Roch}

\noindent
设 $k$ 为域，$X$ 为维数 $\leq 1$ 的 $k$ 上固有概形。
在《簇》，第 \ref{varieties-section-divisors-curves} 节中，
我们用下式定义常秩局部自由 $\mathcal{O}_X$-模
$\mathcal{E}$ 的次数：
\begin{equation}
\label{curves-equation-degree}
\deg(\mathcal{E}) =
\chi(X, \mathcal{E}) - \text{rank}(\mathcal{E})\chi(X, \mathcal{O}_X)
\end{equation}
；见《簇》，定义 \ref{varieties-definition-degree-invertible-sheaf}。
在《Chow 同调》一章中，我们把 $\mathcal{E}$ 的第一 Chern 类定义为
闭链上的运算（《Chow 同调》，第
\ref{chow-section-intersecting-chern-classes} 节），并证明
\begin{equation}
\label{curves-equation-degree-c1}
\deg(\mathcal{E}) = \deg(c_1(\mathcal{E}) \cap [X]_1)
\end{equation}
；见《Chow 同调》，引理 \ref{chow-lemma-degree-vector-bundle}。
结合 (\ref{curves-equation-degree}) 与 (\ref{curves-equation-degree-c1})，
得到第一种 Riemann--Roch 公式：
\begin{equation}
\label{curves-equation-rr}
\chi(X, \mathcal{E}) =
\deg(c_1(\mathcal{E}) \cap [X]_1) +
\text{rank}(\mathcal{E})\chi(X, \mathcal{O}_X)
\end{equation}
若 $\mathcal{L}$ 是可逆 $\mathcal{O}_X$-模，还可以考虑数值交数
$(\mathcal{L} \cdot X)$，其定义见《簇》，定义
\ref{varieties-definition-intersection-number}。不过，这并未给出新内容，因为
\begin{equation}
\label{curves-equation-numerical-degree}
(\mathcal{L} \cdot X) = \deg(\mathcal{L})
\end{equation}
；这由《簇》，引理
\ref{varieties-lemma-intersection-numbers-and-degrees-on-curves} 得到。
若 $\mathcal{L}$ 丰沛，则这个整数为正，并称下式
\begin{equation}
\label{curves-equation-degree-X}
\deg_\mathcal{L}(X) = (\mathcal{L} \cdot X) = \deg(\mathcal{L})
\end{equation}
为 $X$ 关于 $\mathcal{L}$ 的次数；见《簇》，定义
\ref{varieties-definition-degree}。

\medskip\noindent
为了得到真正的 Riemann--Roch 定理，我们希望把
$\chi(X, \mathcal{O}_X)$ 写成 $X$ 上某个典范零维闭链的次数。
完全一般的版本见 \cite{F}。我们将利用对偶性，在 $X$ 为
Gorenstein 时得到一个公式；不过，从某种意义上说这是一种取巧
（例如，因为这种方法在更高维中不能奏效）。

\medskip\noindent
我们先利用引理 \ref{curves-lemma-duality-dim-1} 与
\ref{curves-lemma-duality-dim-1-CM}，得到 $\mathcal{O}_X$ 的 Euler
示性数同对偶化复形或对偶化模的 Euler 示性数之间的关系。

\begin{lemma}
\label{curves-lemma-euler}
设 $X$ 为维数 $\leq 1$ 的域 $k$ 上固有概形。
令 $\omega_X^\bullet$ 与 $\omega_X$ 如引理
\ref{curves-lemma-duality-dim-1} 中所定义，则
$$
\chi(X, \mathcal{O}_X) = \chi(X, \omega_X^\bullet)
$$
。若 $X$ 是 Cohen--Macaulay 的，并且等维且维数为 $1$，则
$$
\chi(X, \mathcal{O}_X) = - \chi(X, \omega_X)
$$
\end{lemma}

\begin{proof}
第一条公式的右端定义如下：
$$
\chi(X, \omega_X^\bullet) =
\sum\nolimits_{i \in \mathbf{Z}} (-1)^i\dim_k H^i(X, \omega_X^\bullet)
$$
这是良定义的，因为 $\omega_X^\bullet$ 属于
$D^b_{\textit{Coh}}(\mathcal{O}_X)$；也可以由下式看出：
$$
H^i(X, \omega_X^\bullet) =
\Ext^i(\mathcal{O}_X, \omega_X^\bullet) =
H^{-i}(X, \mathcal{O}_X)
$$
它总是有限维的，并且仅在 $i = 0, -1$ 时非零。
这当然也证明了第一条公式。在 CM 情形下，
$\omega_X^\bullet = \omega_X[1]$，所以第二条公式由第一条公式推出；
见引理 \ref{curves-lemma-duality-dim-1-CM}。
\end{proof}

\noindent
我们将利用引理 \ref{curves-lemma-euler} 得到所需的
$\chi(X, \mathcal{O}_X)$ 公式，此时 $\omega_X$ 可逆，
亦即 $X$ 为 Gorenstein 的。其断言是：取 $-1/2$ 倍的
$\omega_X$ 的第一 Chern 类，再同闭链 $[X]_1$ 作帽积；
该闭链与 $X$ 相伴。所得闭链是 $X$ 上一个具有半整数系数的
自然零维闭链，其次数为
$\chi(X, \mathcal{O}_X)$。Riemann--Roch 定理的陈述中无法避免分数。

\begin{lemma}[Riemann--Roch]
\label{curves-lemma-rr}
设 $X$ 为域 $k$ 上的固有概形，它是 Gorenstein 的，
并且等维且维数为 $1$。设 $\omega_X$ 如引理
\ref{curves-lemma-duality-dim-1} 中所定义。则
\begin{enumerate}
\item $\omega_X$ 是可逆 $\mathcal{O}_X$-模；
\item $\deg(\omega_X) = -2\chi(X, \mathcal{O}_X)$；
\item 对常秩局部自由 $\mathcal{O}_X$-模 $\mathcal{E}$，有
$$
\chi(X, \mathcal{E}) = \deg(\mathcal{E}) -
\textstyle{\frac{1}{2}} \text{rank}(\mathcal{E}) \deg(\omega_X)
$$
，且 $\dim_k(H^i(X, \mathcal{E})) =
\dim_k(H^{1 - i}(X, \mathcal{E}^\vee \otimes_{\mathcal{O}_X} \omega_X))$
对所有 $i \in \mathbf{Z}$ 成立。
\end{enumerate}
\end{lemma}

\noindent
非奇异（正规）曲线是 Gorenstein 的；见《概形的对偶性》，引理
\ref{duality-lemma-regular-gorenstein}。

\begin{proof}
回顾 Gorenstein 概形是 Cohen--Macaulay 的
（《概形的对偶性》，引理 \ref{duality-lemma-gorenstein-CM}），
因而 $\omega_X$ 是 $X$ 上的对偶化模；见引理
\ref{curves-lemma-duality-dim-1-CM}。
从 Gorenstein 性质的定义几乎立即可知，对偶化层可逆；见《概形的对偶性》，
第 \ref{duality-section-gorenstein} 节。
把 (\ref{curves-equation-rr}) 应用于 $\omega_X$，得到
$$
\chi(X, \omega_X) = 
\deg(c_1(\omega_X) \cap [X]_1) + \chi(X, \mathcal{O}_X)
$$
。结合引理 \ref{curves-lemma-euler}，得到
$$
2\chi(X, \mathcal{O}_X) = - \deg(c_1(\omega_X) \cap [X]_1) = - \deg(\omega_X)
$$
；第二个等号由 (\ref{curves-equation-degree-c1}) 得到。
把它代回关于 $\mathcal{E}$ 的 (\ref{curves-equation-rr})，即得引理中的展示公式。
维数的对称性是 $X$ 的对偶性的推论；见备注
\ref{curves-remark-rework-duality-locally-free}。
\end{proof}





\section{一些消没结果}
\label{curves-section-vanishing}

\noindent
本节包含一些很弱的消没结果。更多且更有趣的若干结果见
第 \ref{curves-section-more-vanishing} 节。

\begin{lemma}
\label{curves-lemma-automatic}
设 $k$ 为域，$X$ 为 $k$ 上维数为 $1$ 的固有概形，且
$H^0(X, \mathcal{O}_X) = k$。则 $X$ 连通、Cohen--Macaulay，
并且等维且维数为 $1$。
\end{lemma}

\begin{proof}
由于 $\Gamma(X, \mathcal{O}_X) = k$ 没有非平凡幂等元，
$X$ 连通。这已经说明 $X$ 等维且维数为 $1$
（任何维数为 $0$ 的不可约分支都会是一个连通分支）。
设 $\mathcal{I} \subset \mathcal{O}_X$ 为支撑于闭点的最大凝聚子模。
这个 $\mathcal{I}$ 存在
（《除子》，引理 \ref{divisors-lemma-remove-embedded-points}），
并由整体截面生成（《簇》，引理
\ref{varieties-lemma-chi-tensor-finite}）。
由于 $1 \in \Gamma(X, \mathcal{O}_X)$ 不是 $\mathcal{I}$ 的截面，
可知 $\mathcal{I} = 0$。所以 $X$ 没有嵌入点
（《除子》，引理 \ref{divisors-lemma-remove-embedded-points}）。
再由《除子》，引理 \ref{divisors-lemma-S1-no-embedded}，
$X$ 满足 $(S_1)$。因此 $X$ 是 Cohen--Macaulay 的。
\end{proof}

\noindent
本节始终处于下述情形。

\begin{situation}
\label{curves-situation-Cohen-Macaulay-curve}
这里 $k$ 是域，$X$ 是 $k$ 上维数为 $1$ 的固有概形，并且
$H^0(X, \mathcal{O}_X) = k$。
\end{situation}

\noindent
由引理 \ref{curves-lemma-automatic}，概形 $X$ 是 Cohen--Macaulay 的，
并且等维且维数为 $1$。引理 \ref{curves-lemma-duality-dim-1} 与
\ref{curves-lemma-duality-dim-1-CM} 中讨论的对偶化模 $\omega_X$
具有不消没的 $H^1$，因为事实上
$\dim_k H^1(X, \omega_X) = \dim_k H^0(X, \mathcal{O}_X) = 1$。
事实证明，任何比 $\omega_X$ 稍微更“正”的对象都有消没的 $H^1$。

\begin{lemma}
\label{curves-lemma-vanishing}
在情形 \ref{curves-situation-Cohen-Macaulay-curve} 中，给定正合列
$$
\omega_X \to \mathcal{F} \to \mathcal{Q} \to 0
$$
，其中各项是凝聚 $\mathcal{O}_X$-模，且
$H^1(X, \mathcal{Q}) = 0$
（例如 $\dim(\text{Supp}(\mathcal{Q})) = 0$ 时如此），则或者
$H^1(X, \mathcal{F}) = 0$，或者
$\mathcal{F} = \omega_X \oplus \mathcal{Q}$。
\end{lemma}

\begin{proof}
（括号中的断言由《概形的上同调》，引理
\ref{coherent-lemma-coherent-support-dimension-0} 得到。）
由于 $H^0(X, \mathcal{O}_X) = k$ 与 $H^1(X, \omega_X)$ 对偶
（见第 \ref{curves-section-Riemann-Roch} 节），可知
$\dim H^1(X, \omega_X) = 1$。层 $\omega_X$ 表示函子
$\mathcal{F} \mapsto \Hom_k(H^1(X, \mathcal{F}), k)$
，这里范畴为凝聚 $\mathcal{O}_X$-模所成的范畴
（《概形的对偶性》，引理
\ref{duality-lemma-dualizing-module-proper-over-A}）。
考虑引理陈述中的正合列，并假设
$H^1(X, \mathcal{F}) \not = 0$。由于
$H^1(X, \mathcal{Q}) = 0$，映射
$H^1(X, \omega_X) \to H^1(X, \mathcal{F})$ 是同构。
由上述 $\omega_X$ 的泛性质，存在映射
$\mathcal{F} \to \omega_X$，其在 $H^1$ 上的作用是该同构的逆。
复合 $\omega_X \to \mathcal{F} \to \omega_X$
是恒等映射（由泛性质），引理得证。
\end{proof}

\begin{lemma}
\label{curves-lemma-vanishing-twist}
在情形 \ref{curves-situation-Cohen-Macaulay-curve} 中，设
$\mathcal{L}$ 为可逆 $\mathcal{O}_X$-模，它由整体截面生成，
且不同构于 $\mathcal{O}_X$。则
$H^1(X, \omega_X \otimes \mathcal{L}) = 0$。
\end{lemma}

\begin{proof}
由第 \ref{curves-section-Riemann-Roch} 节所讨论的对偶性，只需证明
$H^0(X, \mathcal{L}^{\otimes - 1}) = 0$。否则，可以选取整体截面
$t$，它属于 $\mathcal{L}^{\otimes - 1}$；再选取整体截面
$s$，它属于 $\mathcal{L}$，使得 $st \not = 0$。
但由假设，$st$ 是 $1$ 的常数倍，因为
$H^0(X, \mathcal{O}_X) = k$。于是
$\mathcal{L} \cong \mathcal{O}_X$，矛盾。
\end{proof}

\begin{lemma}
\label{curves-lemma-globally-generated}
在情形 \ref{curves-situation-Cohen-Macaulay-curve} 中，给定正合列
$$
\omega_X \to \mathcal{F} \to \mathcal{Q} \to 0
$$
，其中各项是凝聚 $\mathcal{O}_X$-模，
$\dim(\text{Supp}(\mathcal{Q})) = 0$ 且
$\dim_k H^0(X, \mathcal{Q}) \geq 2$；并假设不存在非零子模
$\mathcal{Q}' \subset \mathcal{F}$，使
$\mathcal{Q}' \to \mathcal{Q}$ 单射。
则 $\mathcal{F}$ 中由整体截面生成的子模满射到 $\mathcal{Q}$。
\end{lemma}

\begin{proof}
设 $\mathcal{F}' \subset \mathcal{F}$ 为由整体截面以及
$\omega_X \to \mathcal{F}$ 的像生成的子模。由于
$\dim_k H^0(X, \mathcal{Q}) \geq 2$，且
$\dim_k H^1(X, \omega_X) = \dim_k H^0(X, \mathcal{O}_X) = 1$，
可知 $\mathcal{F}' \to \mathcal{Q}$ 非零，并且
$\omega_X \to \mathcal{F}'$ 不是同构。
因此 $H^1(X, \mathcal{F}') = 0$；这由引理
\ref{curves-lemma-vanishing} 以及对 $\mathcal{F}$ 的假设得到。考虑短正合列
$$
0 \to \mathcal{F}' \to \mathcal{F} \to
\mathcal{Q}/\Im(\mathcal{F}' \to \mathcal{Q}) \to 0
$$
。若右端商非零，便产生矛盾，因为此时
$H^0(X, \mathcal{F})$ 比 $H^0(X, \mathcal{F}')$ 大。
\end{proof}

\noindent
下面给出一个整体生成的例子。

\begin{lemma}
\label{curves-lemma-globally-generated-curve}
在情形 \ref{curves-situation-Cohen-Macaulay-curve} 中，假设 $X$ 是整的。
设 $0 \to \omega_X \to \mathcal{F} \to \mathcal{Q} \to 0$
是凝聚 $\mathcal{O}_X$-模的短正合列，其中 $\mathcal{F}$ 无挠，
$\dim(\text{Supp}(\mathcal{Q})) = 0$，且
$\dim_k H^0(X, \mathcal{Q}) \geq 2$。则 $\mathcal{F}$ 由整体截面生成。
\end{lemma}

\begin{proof}
考虑由整体截面生成的子模 $\mathcal{F}'$。由引理
\ref{curves-lemma-globally-generated}，$\mathcal{F}' \to \mathcal{Q}$ 满射；
特别地，$\mathcal{F}' \not = 0$。由于 $X$ 是曲线，
$\mathcal{F}' \subset \mathcal{F}$ 是秩为 $1$ 的层之间的包含，
所以 $\mathcal{Q}' = \mathcal{F}/\mathcal{F}'$ 支撑于有限多个点。
为导出矛盾，假设 $\mathcal{Q}'$ 非零。则
$H^1(X, \mathcal{F}') \not = 0$。于是由泛性质得到非零映射
$\mathcal{F}' \to \omega_X$（《概形的对偶性》，引理
\ref{duality-lemma-dualizing-module-proper-over-A}）。
复合 $\mathcal{F}' \to \omega_X \to \mathcal{F}$ 的像由整体截面生成，
所以包含于 $\mathcal{F}'$。于是得到非零自映射
$\mathcal{F}' \to \mathcal{F}'$。由于 $\mathcal{F}'$ 是固有曲线上
秩为 $1$ 的无挠层，该映射必为自同构（省略细节）。
但这就蕴含 $\mathcal{F}'$ 包含于 $\omega_X \subset \mathcal{F}$，
同 $\mathcal{F}' \to \mathcal{Q}$ 的满射性矛盾。
\end{proof}

\begin{lemma}
\label{curves-lemma-tensor-omega-with-globally-generated-invertible}
在情形 \ref{curves-situation-Cohen-Macaulay-curve} 中，设 $\mathcal{L}$
为非常丰沛的可逆 $\mathcal{O}_X$-模，且
$\deg(\mathcal{L}) \geq 2$。则
$\omega_X \otimes_{\mathcal{O}_X} \mathcal{L}$ 由整体截面生成。
\end{lemma}

\begin{proof}
先假设 $k$ 代数闭。设 $x \in X$ 为闭点。设
$C_i \subset X$ 为各不可约分支，并对每个 $i$，以
$x_i \in C_i$ 表示其泛点。由《簇》，引理
\ref{varieties-lemma-very-ample-vanish-at-point}，可以选取截面
$s \in H^0(X, \mathcal{L})$，使 $s$ 在 $x$ 处消没，
但不在 $x_i$ 处消没，对所有 $i$ 都如此。相应的模映射
$s : \mathcal{O}_X \to \mathcal{L}$ 是单射，其余核
$\mathcal{Q}$ 支撑于有限多个点，并且
$H^0(X, \mathcal{Q}) \geq 2$。考虑相应的正合列
$$
0 \to \omega_X \to \omega_X \otimes \mathcal{L} \to
\omega_X \otimes \mathcal{Q} \to 0
$$
。由引理 \ref{curves-lemma-globally-generated}，由整体截面生成的模满射到
$\omega_X \otimes \mathcal{Q}$。由于 $x$ 任意，引理得证。
省略若干细节。

\medskip\noindent
下面把 $k$ 非代数闭的情形归约到代数闭域的情形。建议读者跳过
证明的其余部分。选取代数闭包 $\overline{k}$，其基域为 $k$，并考虑基变换
$X_{\overline{k}}$。我们来检验
$X_{\overline{k}} \to \Spec(\overline{k})$ 是情形
\ref{curves-situation-Cohen-Macaulay-curve} 的一个例子。由平坦基变换
（《概形的上同调》，引理
\ref{coherent-lemma-flat-base-change-cohomology}），有
$H^0(X_{\overline{k}}, \mathcal{O}) = \overline{k}$。
概形 $X_{\overline{k}}$ 在 $\overline{k}$ 上固有
（《态射》，引理 \ref{morphisms-lemma-base-change-proper}），
并且等维且维数为 $1$
（《态射》，引理 \ref{morphisms-lemma-dimension-fibre-after-base-change}）。
$\omega_X$ 在 $X_{\overline{k}}$ 上的拉回是
$X_{\overline{k}}$ 的对偶化模；这由引理
\ref{curves-lemma-sanity-check-duality} 得到。
$\mathcal{L}$ 在 $X_{\overline{k}}$ 上的拉回非常丰沛
（《态射》，引理 \ref{morphisms-lemma-very-ample-base-change}）。
$\mathcal{L}$ 在 $X_{\overline{k}}$ 上的拉回之次数，等于
$\mathcal{L}$ 在 $X$ 上的次数（《簇》，引理
\ref{varieties-lemma-degree-base-change}）。最后，
$\omega_X \otimes \mathcal{L}$ 由整体截面生成，当且仅当
其基变换也如此（《簇》，引理
\ref{varieties-lemma-globally-generated-base-change}）。
因此该结果由代数闭基域的情形推出。
\end{proof}




\section{非常丰沛可逆层}
\label{curves-section-very-ample}

\noindent
对于可逆模 $\mathcal{L}$，若它定义在代数闭域上有限型概形
$X$ 上，
常用的一条非常丰沛判据是：$\mathcal{L}$ 的截面分离点与切向量
（《簇》，第 \ref{varieties-section-separating-points-tangent-vectors} 节）。
下面给出曲线的另一条判据；请同《簇》，
第 \ref{varieties-subsection-regularity} 小节比较。

\begin{lemma}
\label{curves-lemma-criterion-very-ample}
设 $k$ 为域，$X$ 为 $k$ 上维数为 $1$ 的固有概形，且
$H^0(X, \mathcal{O}_X) = k$。设 $\mathcal{L}$ 为可逆
$\mathcal{O}_X$-模。假设
\begin{enumerate}
\item $\mathcal{L}$ 有一个正则整体截面；
\item $H^1(X, \mathcal{L}) = 0$；且
\item $\mathcal{L}$ 丰沛。
\end{enumerate}
则 $\mathcal{L}^{\otimes 6}$ 在 $X$ 上相对于 $k$ 非常丰沛。
\end{lemma}

\begin{proof}
设 $s$ 为 $\mathcal{L}$ 的正则整体截面。设
$i : Z = Z(s) \to X$ 为 $s$ 的零点概形；见《除子》，
第 \ref{divisors-section-effective-Cartier-invertible} 节。
由条件 (3)，有 $Z \not = \emptyset$（省略一个小细节）。
考虑短正合列
$$
0 \to \mathcal{O}_X \xrightarrow{s} \mathcal{L} \to
i_*(\mathcal{L}|_Z) \to 0
$$
。与 $\mathcal{L}$ 作张量积，得到
$$
0 \to \mathcal{L} \to \mathcal{L}^{\otimes 2} \to
i_*(\mathcal{L}^{\otimes 2}|_Z) \to 0
$$
。注意 $Z$ 的维数为 $0$
（《除子》，引理
\ref{divisors-lemma-effective-Cartier-makes-dimension-drop}），
所以它是 Artin 环的谱
（《簇》，引理 \ref{varieties-lemma-algebraic-scheme-dim-0}），
进而 $\mathcal{L}|_Z \cong \mathcal{O}_Z$
（《代数》，引理 \ref{algebra-lemma-locally-free-semi-local-free}）。
该短正合列还表明 $H^1(X, \mathcal{L}^{\otimes 2}) = 0$
（例如可用《簇》，引理 \ref{varieties-lemma-chi-tensor-finite}
看出右端相应位置消没）。对 $n \geq 1$ 作归纳，并使用正合列
$$
0 \to \mathcal{L}^{\otimes n} \xrightarrow{s}
\mathcal{L}^{\otimes n + 1} \to
i_*(\mathcal{L}^{\otimes n + 1}|_Z) \to 0
$$
，可知 $H^1(X, \mathcal{L}^{\otimes n}) = 0$ 对 $n > 0$ 成立，
并且存在整体截面 $t_{n + 1}$，它属于
$\mathcal{L}^{\otimes n + 1}$，并给出
$\mathcal{L}^{\otimes n + 1}|_Z \cong \mathcal{O}_Z$ 的平凡化。

\medskip\noindent
考虑乘法映射
$$
\mu_n :
H^0(X, \mathcal{L}) \otimes_k H^0(X, \mathcal{L}^{\otimes n})
\oplus
H^0(X, \mathcal{L}^{\otimes 2}) \otimes_k H^0(X, \mathcal{L}^{\otimes n - 1})
\longrightarrow
H^0(X, \mathcal{L}^{\otimes n + 1})
$$
我们断言它对 $n \geq 3$ 满射。为此考虑短正合列
$$
0 \to \mathcal{L}^{\otimes n} \xrightarrow{s}
\mathcal{L}^{\otimes n + 1} \to i_*(\mathcal{L}^{\otimes n + 1}|_Z) \to 0
$$
。该正合列中从左端而来的 $\mathcal{L}^{\otimes n + 1}$ 的截面，
属于 $\mu_n$ 的像。另一方面，由于
$H^0(\mathcal{L}^{\otimes 2}) \to H^0(\mathcal{L}^{\otimes 2}|_Z)$
满射（见上），且 $t_{n - 1}$ 映到
$\mathcal{L}^{\otimes n - 1}|_Z$ 的一个平凡化，可知
$\mu_n(H^0(X, \mathcal{L}^{\otimes 2}) \otimes t_{n - 1})$
给出 $H^0(X, \mathcal{L}^{\otimes n + 1})$ 的一个子空间，
该子空间满射到 $\mathcal{L}^{\otimes n + 1}|_Z$ 的整体截面。
断言得证。

\medskip\noindent
由上一段的断言可知，分次 $k$-代数
$$
S = \bigoplus\nolimits_{n \geq 0} H^0(X, \mathcal{L}^{\otimes n})
$$
由次数 $0, 1, 2, 3$ 的元素在 $k$ 上生成。
回顾 $X = \text{Proj}(S)$；见《态射》，引理
\ref{morphisms-lemma-proper-ample-is-proj}。
因此 $S^{(6)} = \bigoplus_{n} S_{6n}$ 由次数 $1$ 的元素生成。
这意味着 $\mathcal{L}^{\otimes 6}$ 如所需是非常丰沛的。
\end{proof}

\begin{lemma}
\label{curves-lemma-criterion-very-ample-bis}
设 $k$ 为域，$X$ 为 $k$ 上维数为 $1$ 的固有概形，且
$H^0(X, \mathcal{O}_X) = k$。设 $\mathcal{L}$ 为可逆
$\mathcal{O}_X$-模。假设
\begin{enumerate}
\item $\mathcal{L}$ 由整体截面生成；
\item $H^1(X, \mathcal{L}) = 0$；且
\item $\mathcal{L}$ 丰沛。
\end{enumerate}
则 $\mathcal{L}^{\otimes 2}$ 在 $X$ 上相对于 $k$ 非常丰沛。
\end{lemma}

\begin{proof}
选取基 $s_0, \ldots, s_n$；这是
$H^0(X, \mathcal{L}^{\otimes 2})$ 在 $k$ 上的一组基。由性质 (1)，
$\mathcal{L}^{\otimes 2}$ 由整体截面生成，并得到态射
$$
\varphi_{\mathcal{L}^{\otimes 2}, (s_0, \ldots, s_n)} :
X \longrightarrow \mathbf{P}^n_k
$$
；见《构造》，第 \ref{constructions-section-projective-space} 节。
引理断言此态射是闭浸入。为检验这一点，可以把 $k$ 换成其代数闭包；
见《下降》，引理
\ref{descent-lemma-descending-property-closed-immersion}。
所以可以假设 $k$ 代数闭。

\medskip\noindent
假设 $k$ 代数闭。对每个泛点 $\eta_i \in X$，设
$V_i \subset H^0(X, \mathcal{L})$ 为一个 $k$-向量子空间，
由在 $\eta_i$ 处消没的截面组成。由于 $\mathcal{L}$ 由整体截面生成，
$V_i \not = H^0(X, \mathcal{L})$。由于 $X$ 只有有限多个不可约分支，
且 $k$ 无限，可以找到 $s \in H^0(X, \mathcal{L})$，
它在 $\eta_i$ 处不消没，对所有 $i$ 都如此。
于是 $s$ 是 $\mathcal{L}$ 的正则截面
（因为由引理 \ref{curves-lemma-automatic}，$X$ 是 Cohen--Macaulay 的，
所以 $\mathcal{L}$ 没有嵌入相伴点）。

\medskip\noindent
特别地，引理 \ref{curves-lemma-criterion-very-ample} 的证明中给出的
所有断言对这个 $s$ 都成立。此外，因为 $\mathcal{L}$
由整体截面生成，可以找到整体截面
$t \in H^0(X, \mathcal{L})$，使 $t|_Z$ 不消没
（如上利用 $Z$ 只有有限多个点来论证）。
于是，在引理 \ref{curves-lemma-criterion-very-ample} 的证明中，
可用 $t$ 进一步看出乘法映射
$$
\mu_n :
H^0(X, \mathcal{L}) \otimes_k H^0(X, \mathcal{L}^{\otimes 2})
\longrightarrow
H^0(X, \mathcal{L}^{\otimes 3})
$$
满射。因此
$$
S = \bigoplus\nolimits_{n \geq 0} H^0(X, \mathcal{L}^{\otimes n})
$$
由次数 $0, 1, 2$ 的元素在 $k$ 上生成。仿照引理
\ref{curves-lemma-criterion-very-ample} 的证明论证，得到
$S^{(2)} = \bigoplus_{n} S_{2n}$ 由次数 $1$ 的元素生成。
这意味着 $\mathcal{L}^{\otimes 2}$ 如所需是非常丰沛的。
省略若干细节。
\end{proof}









\section{曲线的亏格}
\label{curves-section-genus}

\noindent
若 $X$ 是域 $k$ 上光滑、投射且几何不可约的曲线，我们先前已把
$X$ 的亏格定义为 $H^1(X, \mathcal{O}_X)$ 的维数；见《曲线的
Picard 概形》，定义 \ref{pic-definition-genus}。注意在此情形下
$H^0(X, \mathcal{O}_X) = k$；见《簇》，引理
\ref{varieties-lemma-regular-functions-proper-variety}。
现作如下推广。

\begin{definition}
\label{curves-definition-genus}
设 $k$ 为域，$X$ 为 $k$ 上维数为 $1$ 的固有概形，且
$H^0(X, \mathcal{O}_X) = k$。则 $X$ 的{\it 亏格}是
$g = \dim_k H^1(X, \mathcal{O}_X)$。
\end{definition}

\noindent
它有时称为 $X$ 的{\it 算术亏格}。文献中，固有曲线
$X$（其基域为 $k$）的算术亏格有时定义为
$$
p_a(X) = 1 - \chi(X, \mathcal{O}_X) =
1 - \dim_k H^0(X, \mathcal{O}_X) + \dim_k H^1(X, \mathcal{O}_X)
$$
。当我们的定义适用时，两者一致，因为我们假设
$H^0(X, \mathcal{O}_X) = k$。但须注意：
\begin{enumerate}
\item $p_a(X)$ 可以为负；且
\item $p_a(X)$ 依赖于基域 $k$，应写作 $p_a(X/k)$。
\end{enumerate}
例如，若 $k = \mathbf{Q}$ 且
$X = \mathbf{P}^1_{\mathbf{Q}(i)}$，则
$p_a(X/\mathbf{Q}) = -1$ 且 $p_a(X/\mathbf{Q}(i)) = 0$。

\medskip\noindent
我们定义中的假设 $H^0(X, \mathcal{O}_X) = k$ 有两个后果。
一方面，它意味着基域没有歧义。另一方面，它蕴含概形 $X$
是 Cohen--Macaulay 的，并且等维且维数为 $1$
（引理 \ref{curves-lemma-automatic}）。若 $\omega_X$ 表示引理
\ref{curves-lemma-duality-dim-1} 与 \ref{curves-lemma-duality-dim-1-CM}
中的对偶化模，则
\begin{equation}
\label{curves-equation-genus}
g = \dim_k H^1(X, \mathcal{O}_X) = \dim_k H^0(X, \omega_X)
\end{equation}
；这由对偶性得到，见备注
\ref{curves-remark-rework-duality-locally-free}。

\medskip\noindent
若 $X$ 在 $k$ 上固有、维数 $\leq 1$，且
$H^0(X, \mathcal{O}_X)$ 不等于基域 $k$，我们不用上面展示公式给出的
算术亏格 $p_a(X)$，而使用不变量 $\chi(X, \mathcal{O}_X)$。
事实上，\cite[page 276]{FAC} 与 \cite[Introduction]{Hirzebruch}
主张应把 $\chi(X, \mathcal{O}_X)$ 称为算术亏格。

\begin{lemma}
\label{curves-lemma-genus-base-change}
设 $k'/k$ 为域扩张，$X$ 为 $k$ 上维数为 $1$ 的固有概形，且
$H^0(X, \mathcal{O}_X) = k$。则 $X_{k'}$ 是 $k'$ 上维数为
$1$ 的固有概形，并且
$H^0(X_{k'}, \mathcal{O}_{X_{k'}}) = k'$。
此外，$X_{k'}$ 的亏格等于 $X$ 的亏格。
\end{lemma}

\begin{proof}
例如由《态射》，引理
\ref{morphisms-lemma-dimension-fibre-after-base-change}，
$X_{k'}$ 的维数为 $1$。态射 $X_{k'} \to \Spec(k')$ 固有，
这由《态射》，引理 \ref{morphisms-lemma-base-change-proper} 得到。
等式 $H^0(X_{k'}, \mathcal{O}_{X_{k'}}) = k'$ 由《概形的上同调》，
引理 \ref{coherent-lemma-flat-base-change-cohomology} 得到。
亏格相等也由同一引理得到。
\end{proof}

\begin{lemma}
\label{curves-lemma-genus-gorenstein}
设 $k$ 为域，$X$ 为 $k$ 上维数为 $1$ 的固有概形，且
$H^0(X, \mathcal{O}_X) = k$。若 $X$ 是 Gorenstein 的，则
$$
\deg(\omega_X) = 2g - 2
$$
，其中 $g$ 是 $X$ 的亏格，而 $\omega_X$ 如引理
\ref{curves-lemma-duality-dim-1} 中所定义。
\end{lemma}

\begin{proof}
由引理 \ref{curves-lemma-rr} 立即得到。
\end{proof}

\begin{lemma}
\label{curves-lemma-genus-smooth}
设 $X$ 为域 $k$ 上的光滑固有曲线，且
$H^0(X, \mathcal{O}_X) = k$。则
$$
\dim_k H^0(X, \Omega_{X/k}) = g
\quad\text{且}\quad
\deg(\Omega_{X/k}) = 2g - 2
$$
，其中 $g$ 是 $X$ 的亏格。
\end{lemma}

\begin{proof}
由引理 \ref{curves-lemma-duality-dim-1}，有 $\Omega_{X/k} = \omega_X$。
所以这些公式由 (\ref{curves-equation-genus}) 及引理
\ref{curves-lemma-genus-gorenstein} 得到。
\end{proof}






\section{平面曲线}
\label{curves-section-plane-curves}

\noindent
设 $k$ 为域。所谓{\it 平面曲线}，是指曲线 $X$，它同
$\mathbf{P}^2_k$ 的某个闭子概形同构。
我们常把嵌入 $X \to \mathbf{P}^2_k$ 视为已给定。
由《除子》，例 \ref{divisors-example-closed-subscheme-of-proj}，
曲线由相应的齐次理想确定：
$$
I(X) =
\Ker\left(
k[T_0, T_2, T_2] \longrightarrow \bigoplus \Gamma(X, \mathcal{O}_X(n))
\right)
$$
。回顾在这种情形下有
$$
X = \text{Proj}(k[T_0, T_2, T_2]/I)
$$
，这是 $\mathbf{P}^2_k$ 的闭子概形。
这些构造的一般资料见《除子》，例
\ref{divisors-example-closed-subscheme-of-proj} 及其中的参考文献。
事实证明，对某个齐次多项式有 $I(X) = (F)$，其中
$F \in k[T_0, T_1, T_2]$；见引理 \ref{curves-lemma-equation-plane-curve}。
由于 $X$ 不可约，$F$ 也不可约；见引理
\ref{curves-lemma-plane-curve}。此外，考察短正合列
$$
0 \to \mathcal{O}_{\mathbf{P}^2_k}(-d) \xrightarrow{F}
\mathcal{O}_{\mathbf{P}^2_k} \to \mathcal{O}_X \to 0
$$
，其中 $d = \deg(F)$，可知 $H^0(X, \mathcal{O}_X) = k$，且 $X$
的亏格为 $(d - 1)(d - 2)/2$；见引理
\ref{curves-lemma-genus-plane-curve} 的证明。

\medskip\noindent
寻找光滑平面曲线时，最容易的方法是写出显式方程。以 $p$
表示 $k$ 的特征。若 $p$ 不整除 $d$，可以取
$$
F = T_0^d + T_1^d + T_2^d
$$
。相应曲线 $X = V_+(F)$ 称为次数为 $d$ 的
{\it Fermat 曲线}。它光滑，因为在每个标准仿射片
$D_+(T_i)$ 上，所得曲线同下述仿射曲线同构：
$$
\Spec(k[x, y]/(x^d + y^d + 1))
$$
。环映射 $k \to k[x, y]/(x^d + y^d + 1)$ 光滑，这由《代数》，
引理 \ref{algebra-lemma-relative-global-complete-intersection-smooth}
得到，因为 $d x^{d - 1}$ 与 $d y^{d - 1}$ 在
$k[x, y]/(x^d + y^d + 1)$ 中生成单位理想。
若 $p | d$ 但 $p \not = 3$，则可用方程
$$
F = T_0^{d - 1}T_1 + T_1^{d - 1}T_2 + T_2^{d - 1}T_0
$$
。确切地说，在各仿射片上得到 $x + x^{d - 1}y + y^{d - 1}$，
其导数为 $1 - x^{d - 2}y$ 与 $x^{d - 1} - y^{d - 2}$；
三者的公共零点集为空\footnote{确实，由
$x^{d - 1} = y^{d - 2}$，有 $0 = x + x^{d - 1}y + y^{d - 1} =
x + 2 x^{d - 1} y$。由于 $x \not = 0$（这是因为
$1 = x^{d - 2}y$），故
$0 = 1 + 2x^{d - 2}y = 3$；除非 $3 = 0$，这不可能。}。
特征为 $3$ 的例子留给读者构造。

\medskip\noindent
更一般地，对任意域 $k$ 以及任意 $n$ 与 $d$，都存在次数为
$d$ 的光滑超曲面，它位于 $\mathbf{P}^n_k$ 中；
例如见 \cite{Poonen}。

\medskip\noindent
当然，这种方法只能得到亏格为三角数的光滑曲线。为了得到任意亏格的
光滑曲线，可以寻找位于
$\mathbf{P}^1 \times \mathbf{P}^1$ 上的光滑曲线
（以后在此插入参考文献）。

\begin{lemma}
\label{curves-lemma-equation-plane-curve}
设 $Z \subset \mathbf{P}^2_k$ 为闭子概形，它等维且维数为 $1$，
并且没有嵌入点（等价地，$Z$ 是 Cohen--Macaulay 的）。
则理想 $I(Z) \subset k[T_0, T_1, T_2]$ 是主理想；
这个理想同 $Z$ 相应。
\end{lemma}

\begin{proof}
这是《除子》，引理
\ref{divisors-lemma-equation-codim-1-in-projective-space} 的特殊情形
（另见《簇》，引理
\ref{varieties-lemma-equation-codim-1-in-projective-space}）。
括号中的断言来自如下事实：$1$ 维 Noether 概形是
Cohen--Macaulay 的，当且仅当它没有嵌入点；见《除子》，引理
\ref{divisors-lemma-noetherian-dim-1-CM-no-embedded-points}。
\end{proof}

\begin{lemma}
\label{curves-lemma-plane-curve}
设 $Z \subset \mathbf{P}^2_k$ 如引理
\ref{curves-lemma-equation-plane-curve} 中所述，并设
$I(Z) = (F)$，其中 $F \in k[T_0, T_1, T_2]$。
则 $Z$ 是曲线，当且仅当 $F$ 不可约。
\end{lemma}

\begin{proof}
若 $F$ 可约，比如 $F = F' F''$，令 $Z'$ 为
$\mathbf{P}^2_k$ 中由 $F'$ 定义的闭子概形。显然
$Z' \subset Z$ 且 $Z' \not = Z$。由于 $Z'$ 也具有维数 $1$，
可知 $Z$ 或者不约化，或者 $Z$ 不可约。
反之，写成 $Z = \sum a_i D_i$，其中 $D_i$ 是 $Z$
的不可约分支；见《除子》，引理
\ref{divisors-lemma-codim-1-part} 与
\ref{divisors-lemma-codimension-1-is-effective-Cartier}。
设 $F_i \in k[T_0, T_1, T_2]$ 为生成 $D_i$ 之理想的齐次多项式。
显然 $F$ 与 $\prod F_i^{a_i}$ 截出 $\mathbf{P}^2_k$
中的同一个闭子概形。因此有 $F = \lambda \prod F_i^{a_i}$，
其中 $\lambda \in k^*$；这是因为二者都生成 $Z$ 的理想。
所以若 $F$ 不可约，则 $Z$ 是素除子，亦即曲线。
\end{proof}

\begin{lemma}
\label{curves-lemma-genus-plane-curve}
设 $Z \subset \mathbf{P}^2_k$ 如引理
\ref{curves-lemma-equation-plane-curve} 中所述，并设
$I(Z) = (F)$，其中 $F \in k[T_0, T_1, T_2]$。
则 $H^0(Z, \mathcal{O}_Z) = k$，且 $Z$ 的亏格为
$(d - 1)(d - 2)/2$，其中 $d = \deg(F)$。
\end{lemma}

\begin{proof}
设 $S = k[T_0, T_1, T_2]$。有分次模的正合列
$$
0 \to S(-d) \xrightarrow{F} S \to S/(F) \to 0
$$
。以 $i : Z \to \mathbf{P}^2_k$ 表示给定的闭浸入。
应用正合函子 $\widetilde{\ }$
（《构造》，引理 \ref{constructions-lemma-proj-sheaves}），得到
$$
0 \to \mathcal{O}_{\mathbf{P}^2_k}(-d) \to
\mathcal{O}_{\mathbf{P}^2_k} \to i_*\mathcal{O}_Z \to 0
$$
，因为 $F$ 生成 $Z$ 的理想。注意
$\mathcal{O}_{\mathbf{P}^2_k}(-d)$ 与
$\mathcal{O}_{\mathbf{P}^2_k}$ 的上同调群由《概形的上同调》，引理
\ref{coherent-lemma-cohomology-projective-space-over-ring} 给出。
另一方面，有
$H^q(Z, \mathcal{O}_Z) = H^q(\mathbf{P}^2_k, i_*\mathcal{O}_Z)$；这由
《概形的上同调》，引理
\ref{coherent-lemma-relative-affine-cohomology}。
应用上同调长正合列，首先得到
$$
k = H^0(\mathbf{P}^2_k, \mathcal{O}_{\mathbf{P}^2_k}) \longrightarrow
H^0(Z, \mathcal{O}_Z)
$$
是同构；其次得到边界映射
$$
H^1(Z, \mathcal{O}_Z) \longrightarrow
H^2(\mathbf{P}^2_k, \mathcal{O}_{\mathbf{P}^2_k}(-d)) \cong
k[T_0, T_1, T_2]_{d - 3}
$$
是同构。容易看出后者的维数为 $(d - 1)(d - 2)/2$，证明完成。
\end{proof}

\begin{lemma}
\label{curves-lemma-smooth-plane-curve-point-over-separable}
设 $Z \subset \mathbf{P}^2_k$ 如引理
\ref{curves-lemma-equation-plane-curve} 中所述，并设
$I(Z) = (F)$，其中 $F \in k[T_0, T_1, T_2]$。
若 $Z \to \Spec(k)$ 至少在一点光滑，且 $k$ 无限，
则光滑轨迹中存在闭点 $z \in Z$，使得 $\kappa(z)/k$
是次数至多为 $d$ 的有限可分扩张。
\end{lemma}

\begin{proof}
设 $z' \in Z$ 是 $Z \to \Spec(k)$ 光滑的一个点。
必要时重新编号坐标，可以假设 $z'$ 属于 $D_+(T_0)$。
置 $f = F(1, x, y) \in k[x, y]$。则
$Z \cap D_+(X_0)$ 同 $k[x, y]/(f)$ 的谱同构。
设 $f_x, f_y$ 为 $f$ 对 $x, y$ 的偏导数。
由于 $z'$ 是 $Z/k$ 的光滑点，$f_x$ 或 $f_y$
在 $z'$ 处非零（见《代数》，第 \ref{algebra-section-smooth} 节中的讨论）。
重新编号坐标后，可以假设 $f_y$ 在 $z'$ 处非零。
因此存在非空开子概形 $V \subset Z \cap D_{+}(X_0)$，使投影
$$
p : V \longrightarrow \Spec(k[x])
$$
是 \'etale 的。由于把 $f$ 看作 $y$ 的多项式时次数至多为 $d$，
投影 $p$ 的各纤维之次数至多为 $d$
（见《态射》，第 \ref{morphisms-section-universally-bounded} 节中的讨论）。
此外，因为 $p$ 是 \'etale 的，$p$ 的像是开集
$U \subset \Spec(k[x])$。最后，由于 $k$ 无限，$k$-有理点集合
$U(k)$（它属于 $U$）无限，特别地非空。
任选 $t \in U(k)$，并令 $z \in V$ 为映到 $t$ 的一个点。
此 $z$ 即满足要求。
\end{proof}





\section{亏格为零的曲线}
\label{curves-section-genus-zero}

\noindent
以后我们需要知道固有的亏格为零曲线具有怎样的形状。
事实证明，Gorenstein 的固有亏格零曲线是次数为 $2$ 的平面曲线，
也就是一条二次曲线；见引理 \ref{curves-lemma-genus-zero}。
一般的固有亏格零曲线可由一条非奇异曲线（定义在某个更大的域上）
通过推出过程得到；见引理 \ref{curves-lemma-genus-zero-singular}。
由于非奇异曲线是 Gorenstein 的，这两个结果涵盖了所有可能情形。

\begin{lemma}
\label{curves-lemma-genus-zero-pic}
设 $X$ 是域 $k$ 上的固有曲线，且 $H^0(X, \mathcal{O}_X) = k$。
若 $X$ 的亏格为 $0$，则每个可逆 $\mathcal{O}_X$-模
$\mathcal{L}$，只要次数为 $0$，都是平凡的。
\end{lemma}

\begin{proof}
确实，由 Riemann--Roch（引理 \ref{curves-lemma-rr}）有
$\dim_k H^0(X, \mathcal{L}) \geq 0 + 1 - 0 = 1$，
故 $\mathcal{L}$ 有非零截面，进而由《簇》，引理
\ref{varieties-lemma-check-invertible-sheaf-trivial}，
有 $\mathcal{L} \cong \mathcal{O}_X$。
\end{proof}

\begin{lemma}
\label{curves-lemma-genus-zero-positive-degree}
设 $X$ 是域 $k$ 上的固有曲线，且 $H^0(X, \mathcal{O}_X) = k$。
假设 $X$ 的亏格为 $0$。设 $\mathcal{L}$ 为可逆
$\mathcal{O}_X$-模，其次数为 $d > 0$。则
\begin{enumerate}
\item $\dim_k H^0(X, \mathcal{L}) = d + 1$ 且 $\dim_k H^1(X, \mathcal{L}) = 0$；
\item $\mathcal{L}$ 非常丰沛，并定义了到 $\mathbf{P}^d_k$ 的闭浸入。
\end{enumerate}
\end{lemma}

\begin{proof}
由次数与亏格的定义，有
$$
\dim_k H^0(X, \mathcal{L}) - \dim_k H^1(X, \mathcal{L}) = d + 1
$$
。设 $s$ 为 $\mathcal{L}$ 的非零截面。
则 $s$ 的零点概形是有效 Cartier 除子
$D \subset X$，有 $\mathcal{L} = \mathcal{O}_X(D)$，并且有短正合列
$$
0 \to \mathcal{O}_X \to \mathcal{L} \to \mathcal{L}|_D \to 0
$$
；见《除子》，引理 \ref{divisors-lemma-characterize-OD} 与
注 \ref{divisors-remark-ses-regular-section}。
由假设 $H^1(X, \mathcal{O}_X) = 0$，故
$H^0(X, \mathcal{L}) \to H^0(X, \mathcal{L}|_D)$ 满射。
又因 $\mathcal{L}|_D$ 由整体截面生成
（因为 $\dim(D) = 0$；见《簇》，引理
\ref{varieties-lemma-chi-tensor-finite}），
所以可逆模 $\mathcal{L}$ 由整体截面生成。
事实上，由于 $D$ 是 Artin 概形，有
$\mathcal{L}|_D \cong \mathcal{O}_D$\footnote{在本情形中，这也由《除子》，
引理 \ref{divisors-lemma-finite-trivialize-invertible-upstairs} 得到，
因为 $D \to \Spec(k)$ 是有限的。}；
因此可以找到一个截面 $t$，它属于 $\mathcal{L}$，其在
$D$ 上的限制生成 $\mathcal{L}|_D$。
该短正合列还表明 $H^1(X, \mathcal{L}) = 0$。

\medskip\noindent
对 $n \geq 1$，考察乘法映射
$$
\mu_n :
H^0(X, \mathcal{L}) \otimes_k H^0(X, \mathcal{L}^{\otimes n})
\longrightarrow
H^0(X, \mathcal{L}^{\otimes n + 1})
$$
。我们断言它是满射。为此，考察短正合列
$$
0 \to \mathcal{L}^{\otimes n} \xrightarrow{s}
\mathcal{L}^{\otimes n + 1} \to \mathcal{L}^{\otimes n + 1}|_D \to 0
$$
。此列中由左端给出的 $\mathcal{L}^{\otimes n + 1}$ 的截面
属于 $\mu_n$ 的像。另一方面，由于
$H^0(\mathcal{L}) \to H^0(\mathcal{L}|_D)$ 满射，且
$t^n$ 映为 $\mathcal{L}^{\otimes n}|_D$ 的一个平凡化，
可见 $\mu_n(H^0(X, \mathcal{L}) \otimes t^n)$ 给出
$H^0(X, \mathcal{L}^{\otimes n + 1})$ 的一个子空间，它满射到
$\mathcal{L}^{\otimes n + 1}|_D$ 的整体截面空间。这就证明了断言。

\medskip\noindent
注意，由《簇》，引理 \ref{varieties-lemma-ample-curve}，
$\mathcal{L}$ 是丰沛的。因此《态射》，引理
\ref{morphisms-lemma-proper-ample-is-proj} 给出同构
$$
X \longrightarrow
\text{Proj}\left(
\bigoplus\nolimits_{n \geq 0} H^0(X, \mathcal{L}^{\otimes n})\right)
$$
。由于映射 $\mu_n$ 对所有 $n \geq 1$ 都是满射，
右端的分次代数是 $H^0(X, \mathcal{L})$ 上对称代数的商。
选取一个 $k$-基 $s_0, \ldots, s_d$，作为
$H^0(X, \mathcal{L})$ 的基，可见它是一个含 $d + 1$ 个变量的
多项式代数的商。分次环的商对应于 $\text{Proj}$ 的闭浸入
（《构造》，引理
\ref{constructions-lemma-surjective-graded-rings-generated-degree-1-map-proj}），
故得到闭浸入 $X \to \mathbf{P}^d_k$。
我们略去如下验证：该态射正是《构造》，引理
\ref{constructions-lemma-projective-space} 中同截面
$s_0, \ldots, s_d$（它们属于 $\mathcal{L}$）相联系的态射。
\end{proof}

\begin{lemma}
\label{curves-lemma-genus-zero}
设 $X$ 是域 $k$ 上的固有曲线，且 $H^0(X, \mathcal{O}_X) = k$。
若 $X$ 是 Gorenstein 的且亏格为 $0$，则 $X$
同次数为 $2$ 的平面曲线同构。
\end{lemma}

\begin{proof}
考察可逆层 $\mathcal{L} = \omega_X^{\otimes -1}$，其中
$\omega_X$ 如引理 \ref{curves-lemma-duality-dim-1} 所述。
由引理 \ref{curves-lemma-genus-gorenstein}，
$\deg(\omega_X) = -2$，故 $\deg(\mathcal{L}) = 2$。
由引理 \ref{curves-lemma-genus-zero-positive-degree}，
选取一个基 $s_0, s_1, s_2$，它是如下 $k$-向量空间的基：
$\mathcal{L}$ 的整体截面空间。由此得到闭浸入
$$
\varphi_{(\mathcal{L}, (s_0, s_1, s_2))} :
X \longrightarrow \mathbf{P}^2_k
$$
。因此 $X$ 是某个次数为 $d$ 的平面曲线。设
$F \in k[T_0, T_1, T_2]_d$ 为其方程
（引理 \ref{curves-lemma-equation-plane-curve}）。
由于 $X$ 的亏格为 $0$，可知 $d$ 是 $1$ 或 $2$
（引理 \ref{curves-lemma-genus-plane-curve}）。
注意，$F$ 在 $\varphi(X)$ 上的限制是零截面，因此
$F(s_0, s_1, s_2)$ 是 $\mathcal{L}^{\otimes 2}$ 的零截面。
由于 $s_0, s_1, s_2$ 线性无关，$F$ 不可能是线性的，
即 $d = \deg(F) \geq 2$。所以 $d = 2$，证明完成。
\end{proof}

\begin{proposition}[射影直线的刻画]
\label{curves-proposition-projective-line}
设 $k$ 为域，设 $X$ 为 $k$ 上的固有曲线。下列条件等价：
\begin{enumerate}
\item $X \cong \mathbf{P}^1_k$；
\item $X$ 在 $k$ 上光滑且几何不可约，$X$ 的亏格为 $0$，
并且 $X$ 有次数为奇数的可逆模；
\item $X$ 在 $k$ 上几何整，$X$ 的亏格为 $0$，$X$ 是 Gorenstein 的，
并且 $X$ 有次数为奇数的可逆层；
\item $H^0(X, \mathcal{O}_X) = k$，$X$ 的亏格为 $0$，$X$ 是 Gorenstein 的，
并且 $X$ 有次数为奇数的可逆层；
\item $X$ 在 $k$ 上几何整，$X$ 的亏格为 $0$，
并且 $X$ 有可逆 $\mathcal{O}_X$-模，其次数为 $1$；
\item $H^0(X, \mathcal{O}_X) = k$，$X$ 的亏格为 $0$，
并且 $X$ 有可逆 $\mathcal{O}_X$-模，其次数为 $1$；
\item $H^1(X, \mathcal{O}_X) = 0$，并且 $X$ 有可逆
$\mathcal{O}_X$-模，其次数为 $1$；
\item $H^1(X, \mathcal{O}_X) = 0$，且 $X$ 有闭点
$x_1, \ldots, x_n$，使得 $\mathcal{O}_{X, x_i}$ 正规且
$\gcd([\kappa(x_i) : k]) = 1$；以及
\item 以后在此添加更多条件。
\end{enumerate}
\end{proposition}

\begin{proof}
我们将证明条件 (2) -- (8) 中的每一个都推出 (1)，并略去
(1) 推出 (2) -- (8) 的验证。

\medskip\noindent
假设 (2)。$k$ 上的光滑概形是几何约化的
（《簇》，引理 \ref{varieties-lemma-smooth-geometrically-normal}）
且正则（《簇》，引理 \ref{varieties-lemma-smooth-regular}）。
因此 $X$ 是 Gorenstein 的（《概形的对偶性》，引理
\ref{duality-lemma-regular-gorenstein}）。于是归结为 (3)。

\medskip\noindent
假设 (3)。由于 $X$ 在 $k$ 上几何整，由《簇》，引理
\ref{varieties-lemma-regular-functions-proper-variety}，
有 $H^0(X, \mathcal{O}_X) = k$。于是归结为 (4)。

\medskip\noindent
假设 (4)。由于 $X$ 是 Gorenstein 的，引理
\ref{curves-lemma-duality-dim-1} 中的对偶化模 $\omega_X$
由引理 \ref{curves-lemma-genus-gorenstein} 具有次数
$\deg(\omega_X) = -2$。结合存在奇数次数可逆模的假设，
可知存在次数为 $1$ 的可逆模。这样便归结为 (6)。

\medskip\noindent
假设 (5)。由于 $X$ 在 $k$ 上几何整，由《簇》，引理
\ref{varieties-lemma-regular-functions-proper-variety}，
有 $H^0(X, \mathcal{O}_X) = k$。于是归结为 (6)。

\medskip\noindent
假设 (6)。由引理 \ref{curves-lemma-genus-zero-positive-degree}，
有 $X \cong \mathbf{P}^1_k$。

\medskip\noindent
假设 (7)。注意，由《簇》，引理
\ref{varieties-lemma-regular-functions-proper-variety}，
$\kappa = H^0(X, \mathcal{O}_X)$ 是 $k$ 的有限域扩张。
若 $d = [\kappa : k] > 1$，则每个可逆层的次数都被 $d$ 整除，
因而不可能存在次数为 $1$ 的可逆层。所以 $d = 1$，归结为情形 (6)。

\medskip\noindent
假设 (8)。注意，由《簇》，引理
\ref{varieties-lemma-regular-functions-proper-variety}，
$\kappa = H^0(X, \mathcal{O}_X)$ 是 $k$ 的有限域扩张。
由于 $\kappa \subset \kappa(x_i)$，次数最大公因数的假设给出
$k = \kappa$。同一条件使我们可以找到整数 $a_i$，使得
$1 = \sum a_i[\kappa(x_i) : k]$。
由《簇》，引理 \ref{varieties-lemma-regular-point-on-curve}，
$x_i$ 在 $X$ 上定义一个有效 Cartier 除子，故可考察可逆模
$\mathcal{O}_X(\sum a_i x_i)$。按我们对 $a_i$ 的选取，
$\mathcal{L}$ 的次数为 $1$。由引理
\ref{curves-lemma-genus-zero-positive-degree}，有
$X \cong \mathbf{P}^1_k$。
\end{proof}

\begin{lemma}
\label{curves-lemma-genus-zero-singular}
设 $X$ 是域 $k$ 上的固有曲线，且 $H^0(X, \mathcal{O}_X) = k$。
假设 $X$ 奇异且亏格为 $0$。则存在交换图
$$
\xymatrix{
x' \ar[d] \ar[r] & X' \ar[d]^\nu \ar[r] & \Spec(k') \ar[d] \\
x \ar[r] & X \ar[r] & \Spec(k)
}
$$
，其中
\begin{enumerate}
\item $k'/k$ 是非平凡有限扩张；
\item $X' \cong \mathbf{P}^1_{k'}$；
\item $x'$ 是一个 $k'$-有理点，它属于 $X'$；
\item $x$ 是一个 $k$-有理点，它属于 $X$；
\item $X' \setminus \{x'\} \to X \setminus \{x\}$ 是同构；
\item $0 \to \mathcal{O}_X \to \nu_*\mathcal{O}_{X'} \to k'/k \to 0$
是短正合列，其中
$k'/k = \kappa(x')/\kappa(x)$ 表示支撑在点 $x$ 上的摩天层。
\end{enumerate}
\end{lemma}

\begin{proof}
设 $\nu : X' \to X$ 为 $X$ 的正规化；见《簇》，第
\ref{varieties-section-normalization} 节与第
\ref{varieties-section-normalization-one-dimensional} 节。
由于 $X$ 奇异，$\nu$ 不是同构。于是
$k' = H^0(X', \mathcal{O}_{X'})$ 是 $k$ 的有限扩张
（《簇》，引理 \ref{varieties-lemma-regular-functions-proper-variety}）。
短正合列
$$
0 \to \mathcal{O}_X \to \nu_*\mathcal{O}_{X'} \to \mathcal{Q} \to 0
$$
以及 $\mathcal{Q}$ 支撑在有限多个闭点上的事实告诉我们：
\begin{enumerate}
\item $H^1(X', \mathcal{O}_{X'}) = 0$，即 $X'$ 的亏格为 $0$，
这里把它视作 $k'$ 上的曲线；
\item 有短正合列
$0 \to k \to k' \to H^0(X, \mathcal{Q}) \to 0$。
\end{enumerate}
特别地，$k'/k$ 是非平凡扩张。

\medskip\noindent
下面考察通常称为{\it 导子理想}的对象
$$
\mathcal{I} = \SheafHom_{\mathcal{O}_X}(\nu_*\mathcal{O}_{X'}, \mathcal{O}_X)
$$
。这是拟凝聚 $\mathcal{O}_X$-模。把 $\mathcal{I}$ 看成
$\mathcal{O}_X$ 中的理想，所用映射为
$\varphi \mapsto \varphi(1)$。于是 $\mathcal{I}(U)$ 是所有
$f \in \mathcal{O}_X(U)$ 的集合，它们满足
$f \left(\nu_*\mathcal{O}_{X'}(U)\right) \subset \mathcal{O}_X(U)$。
换言之，条件是
$f$ 零化 $\mathcal{Q}$。再换言之，有定义性正合列
$$
0 \to \mathcal{I} \to \mathcal{O}_X \to
\SheafHom_{\mathcal{O}_X}(\mathcal{Q}, \mathcal{Q})
$$
。设 $U \subset X$ 为包含 $\mathcal{Q}$ 之支撑的仿射开集。
则 $V = \mathcal{Q}(U) = H^0(X, \mathcal{Q})$ 是
$k$-向量空间，维数为 $n - 1$。映射
$\mathcal{O}_X(U) \to \Hom_k(V, V)$ 的像是交换子代数，
故其维数 $\leq n - 1$；这里维数取自 $k$
（这是矩阵代数的交换子代数所具有的性质；略去细节）。
由此得到短正合列
$$
0 \to \mathcal{I} \to \mathcal{O}_X \to \mathcal{A} \to 0
$$
，其中 $\text{Supp}(\mathcal{A}) = \text{Supp}(\mathcal{Q})$
且 $\dim_k H^0(X, \mathcal{A}) \leq n - 1$。
另一方面，描述
$\mathcal{I} = \SheafHom_{\mathcal{O}_X}(\nu_*\mathcal{O}_{X'}, \mathcal{O}_X)$
赋予 $\mathcal{I}$ 一个 $\nu_*\mathcal{O}_{X'}$-模结构，
使包含映射 $\mathcal{I} \to \nu_*\mathcal{O}_{X'}$
成为 $\nu_*\mathcal{O}_{X'}$-模映射。
因此有 $\mathcal{I} = \nu_*\mathcal{I}'$，其中
$\mathcal{I}' \subset \mathcal{O}_{X'}$ 是某个拟凝聚理想层；见《态射》，引理
\ref{morphisms-lemma-affine-equivalence-modules}。
把 $\mathcal{A}'$ 定义为余核：
$$
0 \to \mathcal{I}' \to \mathcal{O}_{X'} \to \mathcal{A}' \to 0
$$
。合并以上正合列，得到短正合列
$0 \to \mathcal{A} \to \nu_*\mathcal{A}' \to \mathcal{Q} \to 0$。
使用上述估计，并结合
$\dim_k H^0(X, \mathcal{Q}) = n - 1$，得到
$$
\dim_k H^0(X', \mathcal{A}') =
\dim_k H^0(X, \mathcal{A}) + \dim_k H^0(X, \mathcal{Q}) \leq 2 n - 2
$$
。然而，由于 $X'$ 是 $k'$ 上的曲线，左端被 $n$ 整除
（《簇》，引理 \ref{varieties-lemma-divisible}）。
由于 $\mathcal{A}$ 与 $\mathcal{A}'$ 都不可能为零，
可知 $\dim_k H^0(X', \mathcal{A}') = n$；
这意味着 $\mathcal{I}'$ 是某个 $k'$-有理点 $x'$ 的理想层。
由命题 \ref{curves-proposition-projective-line}，得到
$X' \cong \mathbf{P}^1_{k'}$。
回到上述等式，可知 $\dim_k H^0(X, \mathcal{A}) = 1$。
这意味着 $\mathcal{I}$ 是某个 $k$-有理点 $x$ 的理想层。
于是，作为摩天层，有
$\mathcal{A} = \kappa(x) = k$ 且
$\mathcal{A}' = \kappa(x') = k'$。
比较上述正合列，立刻得到引理所述的结构层结论。
\end{proof}

\begin{example}
\label{curves-example-squish-on-P1}
事实上，对任意非平凡有限扩张 $k'/k$，引理
\ref{curves-lemma-genus-zero-singular} 所描述的情形都会发生。
具体地，可以考察
$$
A = \{f \in k'[x] \mid f(0) \in k \}
$$
。$A$ 的谱是一条仿射曲线；可把它同
$B = k'[y]$ 的谱粘合，所用同构为
$A_x \cong B_y$，它把 $x^{-1}$ 映到 $y$。
所得结果是固有曲线 $X$，满足 $H^0(X, \mathcal{O}_X) = k$，
并有奇异点 $x$，它同极大理想 $A \cap (x)$ 相应。
$X$ 的正规化恰如引理中所述，是 $\mathbf{P}^1_{k'}$。
\end{example}







\section{几何亏格}
\label{curves-section-geometric-genus}

\noindent
若 $X$ 是 $k$ 上的固有且{\bf 光滑}的曲线，并且
$H^0(X, \mathcal{O}_X) = k$，则
$$
p_g(X) = \dim_k H^0(X, \Omega_{X/k})
$$
称为 $X$ 的{\it 几何亏格}。由引理 \ref{curves-lemma-genus-smooth}，
$X$ 的几何亏格同（算术）亏格一致。然而在更高维时，
几何亏格与算术亏格并不相同；见注 \ref{curves-remark-genus-higher-dimension}。

\medskip\noindent
对奇异曲线，几何亏格定义如下。

\begin{definition}
\label{curves-definition-geometric-genus}
设 $k$ 为域，设 $X$ 为 $k$ 上的几何不可约曲线。
$X$ 的{\it 几何亏格}，是 $X$ 的某个光滑射影模型的亏格；
该模型可以定义在 $k$ 的某个扩域上，如引理
\ref{curves-lemma-smooth-models} 中所述。
\end{definition}

\noindent
若 $k$ 完美，则非奇异射影模型 $Y$，也就是 $X$ 的非奇异射影模型，是光滑的
（引理 \ref{curves-lemma-nonsingular-model-smooth}），而 $X$
的几何亏格就是 $Y$ 的亏格。但若 $k$ 不完美，这可能不成立。
在这种情形下，选取扩张 $K/k$，使非奇异射影模型
$Y_K$（它是 $(X_K)_{red}$ 的模型）成为光滑射影曲线，并把 $X$ 的几何亏格定义为
$Y_K$ 的亏格。由引理 \ref{curves-lemma-smooth-models} 与
\ref{curves-lemma-genus-base-change}，这个定义是良定的。

\begin{remark}
\label{curves-remark-genus-higher-dimension}
设 $X$ 是 $d$ 维固有光滑簇，定义在代数闭域 $k$ 上。
此时{\it 算术亏格}通常定义为
$p_a(X) = (-1)^d(\chi(X, \mathcal{O}_X) - 1)$，而{\it 几何亏格}
定义为 $p_g(X) = \dim_k H^0(X, \Omega^d_{X/k})$。
在这种情形下，算术亏格与几何亏格不再一致，尽管仍有
$\omega_X \cong \Omega_{X/k}^d$。例如，若 $d = 2$，则
\begin{align*}
p_a(X) - p_g(X) & =
h^0(X, \mathcal{O}_X) - h^1(X, \mathcal{O}_X) + h^2(X, \mathcal{O}_X) - 1
- h^0(X, \Omega^2_{X/k}) \\
& =
- h^1(X, \mathcal{O}_X) + h^2(X, \mathcal{O}_X) - h^0(X, \omega_X) \\
& =
- h^1(X, \mathcal{O}_X)
\end{align*}
，其中 $h^i(X, \mathcal{F}) = \dim_k H^i(X, \mathcal{F})$，
最后一个等式来自对偶性。因此对曲面而言，差
$p_g(X) - p_a(X)$ 总是非负的；它有时称为曲面的不规则度。
若 $X = C_1 \times C_2$ 是亏格分别为 $g_1$ 与 $g_2$
的光滑射影曲线之积，则不规则度为 $g_1 + g_2$。
\end{remark}






\section{Riemann--Hurwitz 公式}
\label{curves-section-riemann-hurewitz}

\noindent
设 $k$ 为域。设 $f : X \to Y$ 为 $k$ 上光滑曲线之间的态射。
由《态射》，引理 \ref{morphisms-lemma-triangle-differentials}，
得到典范正合列
$$
f^*\Omega_{Y/k} \xrightarrow{\text{d}f} \Omega_{X/k}
\longrightarrow \Omega_{X/Y} \longrightarrow 0
$$
。由于 $X$ 与 $Y$ 光滑，层 $\Omega_{X/k}$ 与
$\Omega_{Y/k}$ 是可逆模；见《态射》，引理
\ref{morphisms-lemma-smooth-omega-finite-locally-free}。
假设第一个映射非零，也就是说，假设 $f$ 泛 \'etale；
见引理 \ref{curves-lemma-generically-etale}。设 $R \subset X$ 为
不同理想 $\mathfrak{D}_f$ 所截出的闭子概形；这是 $f$ 的不同理想。
由《判别式》，引理
\ref{discriminant-lemma-discriminant-quasi-finite-morphism-smooth}，
它也就是 $\text{d}f$ 的消没轨迹；它是有效 Cartier 除子，并且有
$$
f^*\Omega_{Y/k} \otimes_{\mathcal{O}_X} \mathcal{O}_X(R) = \Omega_{X/k}
$$
。特别地，若 $X$, $Y$ 是射影的，且
$k = H^0(Y, \mathcal{O}_Y) = H^0(X, \mathcal{O}_X)$，
并且 $X$, $Y$ 的亏格分别为 $g_X$, $g_Y$，则得到
Riemann--Hurwitz 公式
\begin{align*}
2g_X - 2 & =
\deg(\Omega_{X/k}) \\
& =
\deg(f^*\Omega_{Y/k} \otimes_{\mathcal{O}_X} \mathcal{O}_X(R)) \\
& =
\deg(f) \deg(\Omega_{Y/k}) + \deg(R) \\
& =
\deg(f) (2g_Y - 2) + \deg(R)
\end{align*}
。第一个与最后一个等式来自引理 \ref{curves-lemma-genus-smooth}。
第二个等式来自上面给出的可逆层同构。
第三个等式来自次数的可加性
（《簇》，引理 \ref{varieties-lemma-degree-tensor-product}）、
拉回次数的公式
（《簇》，引理 \ref{varieties-lemma-degree-pullback-map-proper-curves}），
以及最后关于 $\mathcal{O}_X(R)$ 次数的公式
（《簇》，引理 \ref{varieties-lemma-degree-effective-Cartier-divisor}）。

\medskip\noindent
为了使用 Riemann--Hurwitz 公式，需要计算
$\deg(R) = \dim_k \Gamma(R, \mathcal{O}_R)$。由 $k$ 上零维概形的结构
（例如见《簇》，引理 \ref{varieties-lemma-algebraic-scheme-dim-0}），
可知 $R$ 是 Artin 局部环之谱的有限不交并：
$R = \coprod_{x \in R} \Spec(\mathcal{O}_{R, x})$，
其中每个 $\mathcal{O}_{R, x}$ 在 $k$ 上都是有限维的。因此
$$
\deg(R) = \sum\nolimits_{x \in R} \dim_k \mathcal{O}_{R, x} =
\sum\nolimits_{x \in R} d_x [\kappa(x) : k]
$$
，其中
$$
d_x = \text{length}_{\mathcal{O}_{R, x}} \mathcal{O}_{R, x} =
\text{length}_{\mathcal{O}_{X, x}} \mathcal{O}_{R, x}
$$
是 $x$ 在 $R$ 中的重数
（见《代数》，引理 \ref{algebra-lemma-pushdown-module}）。
设 $x \in X$ 为闭点，其像为 $y \in Y$。
考察茎，得到正合列
$$
\Omega_{Y/k, y} \to \Omega_{X/k, x} \to \Omega_{X/Y, x} \to 0
$$
。选取局部生成元 $\eta_x$ 与 $\eta_y$；它们分别生成秩为
$1$ 的自由模 $\Omega_{X/k, x}$ 与 $\Omega_{Y/k, y}$。
可知
$
\eta_y \mapsto h \eta_x
$
，其中 $h \in \mathcal{O}_{X, x}$ 非零。由正合列可知，
$\Omega_{X/Y, x} \cong \mathcal{O}_{X, x}/h\mathcal{O}_{X, x}$；
这是作为 $\mathcal{O}_{X, x}$-模的同构。
由于除子 $R$ 由 $h$ 截出（见上），有
$\mathcal{O}_{R, x} = \mathcal{O}_{X, x}/h\mathcal{O}_{X, x}$。
于是得到下列等式：
\begin{align*}
d_x
& =
\text{length}_{\mathcal{O}_{X, x}}(\mathcal{O}_{R, x}) \\
& =
\text{length}_{\mathcal{O}_{X, x}}(\mathcal{O}_{X, x}/h\mathcal{O}_{X, x}) \\
& =
\text{length}_{\mathcal{O}_{X, x}}(\Omega_{X/Y, x}) \\
& =
\text{ord}_{\mathcal{O}_{X, x}}(h) \\
& =
\text{ord}_{\mathcal{O}_{X, x}}(``\eta_y/\eta_x")
\end{align*}
。第一个等式来自 $d_x$ 的定义。第二个与第三个等式已在上面看到。
第四个等式是 $\text{ord}$ 的定义；见《代数》，定义
\ref{algebra-definition-ord}。注意，由于 $\mathcal{O}_{X, x}$
是离散赋值环，整数 $\text{ord}_{\mathcal{O}_{X, x}}(h)$
就是 $h$ 的赋值。第五个等式是一种记忆法。

\medskip\noindent
下面给出一种可以用其他不变量来“计算”重数 $d_x$ 的情形。
确切地说，若 $\kappa(x)$ 在 $k$ 上可分，则可取
$\eta_x = \text{d}s$ 与 $\eta_y = \text{d}t$，其中 $s$ 与 $t$
分别是 $\mathcal{O}_{X, x}$ 与 $\mathcal{O}_{Y, y}$
中的一致化元（引理 \ref{curves-lemma-uniformizer-works}）。
于是有 $t \mapsto u s^{e_x}$，其中
$u \in \mathcal{O}_{X, x}$ 是某个单位，而 $e_x$ 是扩张
$\mathcal{O}_{Y, y} \subset \mathcal{O}_{X, x}$ 的分歧指数。因此
$$
\eta_y = \text{d}t = \text{d}(u s^{e_x}) =
e s^{e_x - 1} u \text{d}s + s^{e_x} \text{d}u
$$
。将 $\text{d}u = w \text{d}s$ 写成此形式，其中
$w \in \mathcal{O}_{X, x}$，可见
$$
``\eta_y/\eta_x" = e s^{e_x - 1} u + s^{e_x} w = (e_x u + s w)s^{e_x - 1}
$$
。因此其消没阶为 $e_x - 1$；例外是 $\kappa(x)$
的特征为 $p > 0$ 且 $p$ 整除 $e_x$ 的情形，此时消没阶
$> e_x - 1$。

\medskip\noindent
合并以上全部结果可知，若 $k$ 的特征为零，则
$$
2g_X - 2 = (2g_Y - 2)\deg(f) +
\sum\nolimits_{x \in X} (e_x - 1)[\kappa(x) : k]
$$
，其中 $e_x$ 是 $\mathcal{O}_{X, x}$ 在
$\mathcal{O}_{Y, f(x)}$ 上的分歧指数。当且仅当所有分歧都是温和的，
也就是剩余域扩张 $\kappa(x)/\kappa(y)$ 可分且
$e_x$ 与 $k$ 的特征互素时，这个精确公式才成立；
不过以上论证不足以证明这一点。参见引理 \ref{curves-lemma-rhe} 及其证明。

\begin{lemma}
\label{curves-lemma-generically-etale}
\begin{slogan}
光滑曲线之间的态射可分，当且仅当它几乎处处 \'etale
\end{slogan}
设 $k$ 为域。设 $f : X \to Y$ 为 $k$ 上光滑曲线之间的态射。
下列条件等价：
\begin{enumerate}
\item $\text{d}f : f^*\Omega_{Y/k} \to \Omega_{X/k}$ 非零；
\item $\Omega_{X/Y}$ 支撑在 $X$ 的某个真闭子集上；
\item 存在非空开集 $U \subset X$，使
$f|_U : U \to Y$ 非分歧；
\item 存在非空开集 $U \subset X$，使
$f|_U : U \to Y$ 为 \'etale；
\item 函数域扩张 $k(X)/k(Y)$ 是有限可分扩张。
\end{enumerate}
\end{lemma}

\begin{proof}
由于 $X$ 与 $Y$ 光滑，层 $\Omega_{X/k}$ 与
$\Omega_{Y/k}$ 是可逆模；见《态射》，引理
\ref{morphisms-lemma-smooth-omega-finite-locally-free}。
利用《态射》，引理 \ref{morphisms-lemma-triangle-differentials}
的正合列
$$
f^*\Omega_{Y/k} \longrightarrow \Omega_{X/k}
\longrightarrow \Omega_{X/Y} \longrightarrow 0
$$
，可知 (1) 与 (2) 等价，并且二者等价于
$f^*\Omega_{Y/k} \to \Omega_{X/k}$ 在泛点处非零这一条件。
(2) 与 (3) 的等价性来自《态射》，引理
\ref{morphisms-lemma-unramified-omega-zero}。
(3) 与 (4) 的等价性来自《态射》，引理
\ref{morphisms-lemma-flat-unramified-etale}，以及平坦性自动成立这一事实
（引理 \ref{curves-lemma-flat}）。
为看出 (5) 与 (4) 等价，可用《代数》，引理
\ref{algebra-lemma-smooth-at-generic-point}。略去若干细节。
\end{proof}

\begin{lemma}
\label{curves-lemma-rh}
设 $f : X \to Y$ 为域 $k$ 上光滑固有曲线之间的态射，
它满足引理 \ref{curves-lemma-generically-etale} 中的等价条件。
若 $k = H^0(Y, \mathcal{O}_Y) = H^0(X, \mathcal{O}_X)$，
并且 $X$ 与 $Y$ 的亏格分别为 $g_X$ 与 $g_Y$，则
$$
2g_X - 2 = (2g_Y - 2) \deg(f) + \deg(R)
$$
，其中 $R \subset X$ 是 $f$ 的不同理想所截出的有效 Cartier 除子。
\end{lemma}

\begin{proof}
见以上讨论；其中用到了《判别式》，引理
\ref{discriminant-lemma-discriminant-quasi-finite-morphism-smooth}、
引理 \ref{curves-lemma-genus-smooth}，以及《簇》，引理
\ref{varieties-lemma-degree-tensor-product} 与
\ref{varieties-lemma-degree-pullback-map-proper-curves}。
\end{proof}

\begin{lemma}
\label{curves-lemma-uniformizer-works}
设 $X \to \Spec(k)$ 光滑且相对维数为 $1$；这里是在闭点
$x \in X$ 处而言。
若 $\kappa(x)$ 在 $k$ 上可分，则对任意一致化元 $s$，
它位于离散赋值环 $\mathcal{O}_{X, x}$ 中，元素
$\text{d}s$ 都自由生成 $\Omega_{X/k, x}$；这里生成是在
$\mathcal{O}_{X, x}$ 上而言。
\end{lemma}

\begin{proof}
由《代数》，引理 \ref{algebra-lemma-characterize-smooth-over-field}，
环 $\mathcal{O}_{X, x}$ 是离散赋值环。由于 $x$ 是闭点，
$\kappa(x)$ 是 $k$ 的有限扩张。因此，若 $\kappa(x)/k$ 可分，
则由《代数》，引理 \ref{algebra-lemma-computation-differential}，
任意一致化元 $s$ 都映为
$\Omega_{X/k, x} \otimes_{\mathcal{O}_{X, x}} \kappa(x)$
中的非零元素。由于 $\Omega_{X/k, x}$ 是秩为 $1$
的自由模，而该模定义在 $\mathcal{O}_{X, x}$ 上，结论成立。
\end{proof}

\begin{lemma}
\label{curves-lemma-rhe}
沿用引理 \ref{curves-lemma-rh} 的记号与假设。对闭点
$x \in X$，设 $d_x$ 为 $x$ 在 $R$ 中的重数。则
$$
2g_X - 2 = (2g_Y - 2) \deg(f) + \sum\nolimits d_x [\kappa(x) : k]
$$
。此外，有下列结论：
\begin{enumerate}
\item $d_x = \text{length}_{\mathcal{O}_{X, x}}(\Omega_{X/Y, x})$；
\item $d_x \geq e_x - 1$，其中 $e_x$ 是
$\mathcal{O}_{X, x}$ 在 $\mathcal{O}_{Y, y}$ 上的分歧指数；
\item $d_x = e_x - 1$，当且仅当 $\mathcal{O}_{X, x}$
在 $\mathcal{O}_{Y, y}$ 上温和分歧。
\end{enumerate}
\end{lemma}

\begin{proof}
由引理 \ref{curves-lemma-rh} 与以上讨论
（其中用到了《簇》，引理
\ref{varieties-lemma-algebraic-scheme-dim-0} 以及《代数》，引理
\ref{algebra-lemma-pushdown-module}），只需证明重数 $d_x$
所满足的结论；这是 $x$ 在 $R$ 中的重数。第 (1) 部分已在以上讨论中证明。
以上讨论只在 $\kappa(x)$ 在 $k$ 上可分的情形证明了 (2) 与 (3)。
在证明的其余部分，我们用关于拟有限 Gorenstein 态射之不同理想的材料，
统一处理 (2) 与 (3)。

\medskip\noindent
首先注意，$f$ 是拟有限 Gorenstein 态射。
例如，这是因为 $f$ 是平坦的拟有限态射，且 $X$ 是 Gorenstein 的
（见《概形的对偶性》，引理
\ref{duality-lemma-flat-morphism-from-gorenstein-scheme}）；
也可以使用《判别式》，引理
\ref{discriminant-lemma-discriminant-quasi-finite-morphism-smooth}
的证明中所得的结果（上面已经使用过该引理）。
因此，由《判别式》，引理
\ref{discriminant-lemma-gorenstein-quasi-finite}，
$\omega_{X/Y}$ 可逆；在把 $X$ 换成开集以及沿某个
$Y' \to Y$ 作基变换之后，这一点仍然成立。以下将不再说明而直接使用。

\medskip\noindent
选取仿射开集 $U \subset X$ 与 $V \subset Y$，使得
$x \in U$、$y \in V$、$f(U) \subset V$，且 $x$ 是
$U$ 中位于 $y$ 上方的唯一一点。写成
$U = \Spec(A)$ 与 $V = \Spec(B)$。
于是 $R \cap U$ 是 $f|_U : U \to V$ 的不同理想。
由《判别式》，引理 \ref{discriminant-lemma-base-change-different}，
在本情形中，不同理想的形成同任意基变换交换。由 $U$ 与 $V$ 的选取，
有
$$
A \otimes_B \kappa(y) =
\mathcal{O}_{X, x} \otimes_{\mathcal{O}_{Y, y}} \kappa(y) =
\mathcal{O}_{X, x}/(s^{e_x})
$$
，其中 $e_x$ 是引理陈述中的分歧指数。令
$C = \mathcal{O}_{X, x}/(s^{e_x})$，把它视作
$\kappa(y)$ 上的有限代数。设 $\mathfrak{D}_{C/\kappa(y)}$
为 $C$ 在 $\kappa(y)$ 上的不同理想，其含义见《判别式》，定义
\ref{discriminant-definition-different}。
只需证明：$\mathfrak{D}_{C/\kappa(y)}$ 非零，当且仅当扩张
$\mathcal{O}_{Y, y} \subset \mathcal{O}_{X, x}$ 温和分歧；
并且在温和分歧情形，$\mathfrak{D}_{C/\kappa(y)}$
等于由 $s^{e_x - 1}$ 生成的理想；这是 $C$ 中的理想。
回顾温和分歧恰好意味着 $\kappa(x)/\kappa(y)$ 可分，且
$\kappa(y)$ 的特征不整除 $e_x$。另一方面，
$C/\kappa(y)$ 的不同理想非零，当且仅当
$\tau_{C/\kappa(y)} \in \omega_{C/\kappa(y)}$ 非零。
确实，由于 $\omega_{C/\kappa(y)}$ 是可逆 $C$-模
（它是 $\omega_{A/B}$ 的基变换），它是秩为 $1$ 的自由模；
设生成元为 $\lambda$。写成
$\tau_{C/\kappa(y)} = h\lambda$，其中 $h \in C$。
于是 $\mathfrak{D}_{C/\kappa(y)} = (h) \subset C$，从而得到该断言。
由《判别式》，引理 \ref{discriminant-lemma-tau-nonzero}，
有 $\tau_{C/\kappa(y)} \not = 0$，当且仅当
$\kappa(x)/\kappa(y)$ 可分且 $e_x$ 与该特征互素。
最后，即使 $\tau_{C/\kappa(y)}$ 非零，仍有
$s \tau_{C/\kappa(y)} = 0$，因为
$s\tau_{C/\kappa(y)} : C \to \kappa(y)$ 把 $c$ 映为
幂零算子 $sc$ 的迹，而这个迹为零。因此 $sh = 0$，从而
$h \in (s^{e_x - 1})$；这证明始终有
$\mathfrak{D}_{C/\kappa(y)} \subset (s^{e_x - 1})$。
由于 $(s^{e_x - 1}) \subset C$ 是最小的非零理想，最终断言得证。
\end{proof}






\section{不可分态射}
\label{curves-section-inseparable}

\noindent
下面对不可分态射下亏格的行为作一些说明。

\begin{lemma}
\label{curves-lemma-dominated-by-smooth}
设 $k$ 为域。设 $f : X \to Y$ 为 $k$ 上曲线之间的满态射。
若 $X$ 在 $k$ 上光滑，且 $Y$ 正规，则 $Y$ 在 $k$ 上光滑。
\end{lemma}

\begin{proof}
设 $y \in Y$。选取 $x \in X$，使它映到 $y$。
由《簇》，引理 \ref{varieties-lemma-flat-under-smooth}，
只需证明 $f$ 在 $x$ 处平坦。这由引理 \ref{curves-lemma-flat} 得到。
\end{proof}

\begin{lemma}
\label{curves-lemma-purely-inseparable}
设 $k$ 为特征 $p > 0$ 的域。设 $f : X \to Y$ 为
$k$ 上固有非奇异曲线之间的非常值态射。
若函数域扩张 $k(X)/k(Y)$ 纯不可分，则存在分解
$$
X = X_0 \to X_1 \to \ldots \to X_n = Y
$$
，使每个 $X_i$ 都是固有非奇异曲线，并且
$X_i \to X_{i + 1}$ 是次数为 $p$ 的态射，而
$k(X_{i + 1}) \subset k(X_i)$ 不可分。
\end{lemma}

\begin{proof}
这由定理 \ref{curves-theorem-curves-rational-maps} 以及如下事实得到：
域的有限纯不可分扩张总能写成一列次数为 $p$ 的（不可分）扩张；
见《域》，引理 \ref{fields-lemma-finite-purely-inseparable}。
\end{proof}

\begin{lemma}
\label{curves-lemma-inseparable-deg-p-smooth}
设 $k$ 为特征 $p > 0$ 的域。设 $f : X \to Y$ 为
$k$ 上固有非奇异曲线之间的非常值态射。
若 $X$ 光滑，且 $k(Y) \subset k(X)$ 是次数为 $p$ 的不可分扩张，
则存在唯一同构 $Y = X^{(p)}$，使 $f$ 就是 $F_{X/k}$。
\end{lemma}

\begin{proof}
相对 Frobenius 态射 $F_{X/k} : X \to X^{(p)}$
构造于《簇》，第 \ref{varieties-section-frobenius} 节。
注意，$X^{(p)}$ 是 $k$ 上的光滑固有曲线；这是 $X$ 的基变换。
由《簇》，引理 \ref{varieties-lemma-inseparable-deg-p-smooth}，
态射 $F_{X/k}$ 的次数为 $p$。因此
$k(X^{(p)})$ 与 $k(Y)$ 都是 $k(X)$ 的子域，并且
$[k(X) : k(Y)] = [k(X) : k(X^{(p)})] = p$。
为了证明引理，只需证明 $k(Y) = k(X^{(p)})$ 在
$k(X)$ 中成立；见定理 \ref{curves-theorem-curves-rational-maps}。

\medskip\noindent
写成 $K = k(X)$。考察映射 $\text{d} : K \to \Omega_{K/k}$。
由引理 \ref{curves-lemma-generically-etale} 可知，
$k(Y)$ 包含于 $\text{d}$ 的核中。由《簇》，引理
\ref{varieties-lemma-relative-frobenius-omega} 可知，
$k(X^{(p)})$ 也包含于 $\text{d}$ 的核中。
由于 $X$ 是光滑曲线，$\Omega_{K/k}$ 是维数为 $1$
的向量空间，定义在 $K$ 上。于是《代数进阶》，引理
\ref{more-algebra-lemma-p-basis} 蕴含
$\Ker(\text{d}) = kK^p$ 且 $[K : kK^p] = p$。
因次数相同，便有 $k(Y) = kK^p = k(X^{(p)})$。
\end{proof}

\begin{lemma}
\label{curves-lemma-purely-inseparable-smooth}
设 $k$ 为特征 $p > 0$ 的域。设 $f : X \to Y$ 为
$k$ 上固有非奇异曲线之间的非常值态射。
若 $X$ 光滑，且 $k(Y) \subset k(X)$ 纯不可分，则存在唯一的
$n \geq 0$ 与唯一同构 $Y = X^{(p^n)}$，使 $f$
是 $n$ 重相对 Frobenius 态射；这是 $X/k$ 的相对 Frobenius。
\end{lemma}

\begin{proof}
$n$ 重相对 Frobenius 态射（对 $X/k$ 而言）定义于《簇》，注
\ref{varieties-remark-n-fold-relative-frobenius}。
合并引理 \ref{curves-lemma-inseparable-deg-p-smooth} 与
\ref{curves-lemma-purely-inseparable} 即得结论。
\end{proof}

\begin{lemma}
\label{curves-lemma-purely-inseparable-smooth-genus}
设 $k$ 为特征 $p > 0$ 的域。设 $f : X \to Y$ 为
$k$ 上固有非奇异曲线之间的非常值态射。假设
\begin{enumerate}
\item $X$ 光滑；
\item $H^0(X, \mathcal{O}_X) = k$；
\item $k(X)/k(Y)$ 纯不可分。
\end{enumerate}
则 $Y$ 光滑，$H^0(Y, \mathcal{O}_Y) = k$，并且 $Y$
的亏格等于 $X$ 的亏格。
\end{lemma}

\begin{proof}
由引理 \ref{curves-lemma-purely-inseparable-smooth}，可知
$Y = X^{(p^n)}$ 是 $X$ 沿 $F_{\Spec(k)}^n$ 的基变换。
因此 $Y$ 光滑，而关于上同调与亏格的结论来自引理
\ref{curves-lemma-genus-base-change}。
\end{proof}

\begin{example}
\label{curves-example-inseparable}
本例说明，在非奇异射影曲线之间的纯不可分态射下，亏格可能改变。
设 $k$ 为特征 $3$ 的域。假设存在元素 $a \in k$，它不是
$3$ 次幂；例如可以取 $k = \mathbf{F}_3(a)$。
设 $X$ 为齐次方程
$$
F = T_1^2T_0 - T_2^3 + aT_0^3
$$
所定义的平面曲线，如第 \ref{curves-section-plane-curves} 节所述。
在仿射片 $D_+(T_0)$ 上使用坐标 $x = T_1/T_0$
与 $y = T_2/T_0$，得到 $x^2 - y^3 + a = 0$，
它定义一条非奇异仿射曲线。此外，无穷远点
$(0 : 1: 0)$ 是光滑点。因此 $X$ 是亏格为 $1$
的非奇异射影曲线（引理 \ref{curves-lemma-genus-plane-curve}）。
另一方面，考察态射 $f : X \to \mathbf{P}^1_k$；它在
$D_+(T_0)$ 上把 $(x, y)$ 映到
$x \in \mathbf{A}^1_k \subset \mathbf{P}^1_k$。
则 $f$ 是 $k$ 上固有非奇异曲线之间的态射，诱导次数为
$p = 3$ 的不可分函数域扩张；但 $X$ 的亏格为 $1$，
而 $\mathbf{P}^1_k$ 的亏格为 $0$。
\end{example}

\begin{proposition}
\label{curves-proposition-unwind-morphism-smooth}
设 $k$ 为特征 $p > 0$ 的域。设 $f : X \to Y$ 为
$k$ 上固有光滑曲线之间的非常值态射。则可把 $f$ 分解为
$$
X \longrightarrow X^{(p^n)} \longrightarrow Y
$$
，其中 $X^{(p^n)} \to Y$ 是固有光滑曲线之间的非常值态射，
诱导可分域扩张 $k(X^{(p^n)})/k(Y)$，且
$$
X^{(p^n)} = X \times_{\Spec(k), F_{\Spec(k)}^n} \Spec(k),
$$
而 $X \to X^{(p^n)}$ 是 $n$ 重相对 Frobenius 态射，作用于 $X$。
\end{proposition}

\begin{proof}
由《域》，引理 \ref{fields-lemma-separable-first}，存在子扩张
$k(X)/E/k(Y)$，使 $k(X)/E$ 纯不可分而 $E/k(Y)$ 可分。
由定理 \ref{curves-theorem-curves-rational-maps}，这对应于分解
$X \to Z \to Y$；它是 $f$ 的分解，其中 $Z$ 是非奇异固有曲线。
把引理 \ref{curves-lemma-purely-inseparable-smooth} 应用于态射
$X \to Z$ 即得结论。
\end{proof}

\begin{lemma}
\label{curves-lemma-inseparable-linear-system}
设 $k$ 为特征 $p > 0$ 的域。设 $X$ 为 $k$ 上的光滑固有曲线。
设 $(\mathcal{L}, V)$ 为 $\mathfrak g^r_d$，其中
$r \geq 1$。则下列两种情形之一成立：
\begin{enumerate}
\item 存在一个 $\mathfrak g^1_d$，其相应态射
$X \to \mathbf{P}^1_k$（引理 \ref{curves-lemma-linear-series}）
泛 \'etale（即如引理 \ref{curves-lemma-generically-etale} 中所述）；或者
\item 存在 $\mathfrak g^r_{d'}$，它在 $X^{(p)}$ 上，其中
$d' \leq d/p$。
\end{enumerate}
\end{lemma}

\begin{proof}
选取两个 $k$-线性无关元素 $s, t \in V$。
则 $f = s/t$ 是定义态射 $X \to \mathbf{P}^1_k$ 的有理函数；
该态射对应于线性系 $(\mathcal{L}, ks + kt)$。
若此态射并非泛 \'etale，则由命题
\ref{curves-proposition-unwind-morphism-smooth}，有
$f \in k(X^{(p)})$。现在选取一个基
$s_0, \ldots, s_r$，作为 $V$ 的基，并令 $\mathcal{L}' \subset \mathcal{L}$
为由 $s_0, \ldots, s_r$ 生成的可逆层。置
$f_i = s_i/s_0$，把它看作 $k(X)$ 中的元素。若对每一对 $(s_0, s_i)$ 都有
$f_i \in k(X^{(p)})$，则态射
$$
\varphi = \varphi_{(\mathcal{L}', (s_0, \ldots, s_r)} :
X
\longrightarrow
\mathbf{P}^r_k = \text{Proj}(k[T_0, \ldots, T_r])
$$
经由 $X^{(p)}$ 分解，因为这在仿射开集 $D_+(T_0)$ 上成立，
而由引理 \ref{curves-lemma-extend-over-normal-curve}，可以把仿射部分上的态射
延拓到整个光滑曲线 $X^{(p)}$ 上。引入记号，设有如下分解：
$$
X \xrightarrow{F_{X/k}} X^{(p)} \xrightarrow{\psi} \mathbf{P}^r_k
$$
；这是 $\varphi$ 的分解。则
$\mathcal{N} = \psi^*\mathcal{O}_{\mathbf{P}^1_k}(1)$
是可逆 $\mathcal{O}_{X^{(p)}}$-模，满足
$\mathcal{L}' = F_{X/k}^*\mathcal{N}$，并且
$\psi^*T_0, \ldots, \psi^*T_r$ 在 $k$ 上线性无关
（因为它们拉回为 $s_0, \ldots, s_r$，这是在 $X$ 上的拉回）。
最后有
$$
d = \deg(\mathcal{L}) \geq \deg(\mathcal{L}') =
\deg(F_{X/k}) \deg(\mathcal{N}) = p \deg(\mathcal{N})
$$
，正如所需。这里使用了《簇》，引理
\ref{varieties-lemma-check-invertible-sheaf-trivial}、
\ref{varieties-lemma-degree-pullback-map-proper-curves} 与
\ref{varieties-lemma-inseparable-deg-p-smooth}。
\end{proof}

\begin{lemma}
\label{curves-lemma-point-over-separable-extension}
设 $k$ 为域。设 $X$ 为 $k$ 上光滑固有曲线，满足
$H^0(X, \mathcal{O}_X) = k$，且亏格 $g \geq 2$。
则存在闭点 $x \in X$，使 $\kappa(x)/k$ 可分且次数
$\leq 2g - 2$。
\end{lemma}

\begin{proof}
置 $\omega = \Omega_{X/k}$。由引理 \ref{curves-lemma-genus-smooth}，
它的次数为 $2g - 2$，并有 $g$ 个整体截面。
因此有一个 $\mathfrak g^{g - 1}_{2g - 2}$。
由平凡的引理 \ref{curves-lemma-linear-series-trivial-existence}，
存在一个 $\mathfrak g^1_{2g - 2}$；再由引理
\ref{curves-lemma-g1d}，得到态射
$$
\varphi : X \longrightarrow \mathbf{P}^1_k
$$
，其次数为 $d \leq 2g - 2$。由于 $\varphi$ 平坦（引理 \ref{curves-lemma-flat}）且有限
（引理 \ref{curves-lemma-finite}），它是次数为 $d$ 的有限局部自由态射
（《态射》，引理 \ref{morphisms-lemma-finite-flat}）。
任取有理点 $t \in \mathbf{P}^1_k$，并任取点 $x \in X$，
使 $\varphi(x) = t$。则
$$
d \geq [\kappa(x) : \kappa(t)] = [\kappa(x) : k]
$$
，例如由《态射》，引理
\ref{morphisms-lemma-finite-locally-free-universally-bounded} 与
\ref{morphisms-lemma-characterize-universally-bounded}。
因此若 $k$ 完美（例如特征为零或有限），引理得证。
于是归结为下一段讨论的情形。

\medskip\noindent
假设 $k$ 是特征 $p > 0$ 的无限域。如上，将使用 $X$
有一个 $\mathfrak g^{g - 1}_{2g - 2}$ 这一事实。
光滑固有曲线 $X^{(p)}$ 与 $X$ 亏格相同，故其亏格 $> 0$。
由引理 \ref{curves-lemma-grd-inequalities}，可知 $X^{(p)}$
没有 $\mathfrak g^{g - 1}_d$，这里 $d \leq g - 1$ 任意。
把引理 \ref{curves-lemma-inseparable-linear-system} 应用于我们的
$\mathfrak g^{g - 1}_{2g - 2}$（并注意 $2g - 2/p \leq g - 1$），
可知情形 (2) 不会发生。因此得到态射
$$
\varphi : X \longrightarrow \mathbf{P}^1_k
$$
，它泛 \'etale（按该引理的含义），且次数 $\leq 2g - 2$。
设 $U \subset X$ 为 $\varphi$ 是 \'etale 的非空开子概形。
则 $\varphi(U) \subset \mathbf{P}^1_k$ 是非空 Zariski 开集；
可以选取 $k$-有理点 $t \in \varphi(U)$，因为 $k$ 无限。
令 $u \in U$ 为满足 $\varphi(u) = t$ 的一点。
则 $\kappa(u)/\kappa(t)$ 可分
（《态射》，引理 \ref{morphisms-lemma-etale-over-field}），
$\kappa(t) = k$，并且如前有
$[\kappa(u) : k] \leq 2g - 2$。
\end{proof}

\noindent
下面的引理其实不属于本节，但我们不知道还能把它放在何处。

\begin{lemma}
\label{curves-lemma-ramification-to-algebraic-closure}
设 $X$ 为域 $k$ 上的光滑曲线。设
$\overline{x} \in X_{\overline{k}}$ 为闭点，其像为 $x \in X$。
则 $\mathcal{O}_{X, x} \subset \mathcal{O}_{X_{\overline{k}}, \overline{x}}$
的分歧指数等于 $\kappa(x)/k$ 的不可分次数。
\end{lemma}

\begin{proof}
缩小 $X$ 后，可以假设存在 \'etale 态射
$\pi : X \to \mathbf{A}^1_k$；见《态射》，引理
\ref{morphisms-lemma-smooth-etale-over-affine-space}。
于是可以考察局部环的交换图
$$
\xymatrix{
\mathcal{O}_{X_{\overline{k}}, \overline{x}} &
\mathcal{O}_{\mathbf{A}^1_{\overline{k}}, \pi(\overline{x})} \ar[l] \\
\mathcal{O}_{X, x} \ar[u] &
\mathcal{O}_{\mathbf{A}^1_k, \pi(x)} \ar[l] \ar[u]
}
$$
。水平箭头的分歧指数为 $1$，因为它们对应于
\'etale 态射。此外，扩张 $\kappa(x)/\kappa(\pi(x))$ 可分，
故 $\kappa(x)$ 与 $\kappa(\pi(x))$ 在 $k$ 上具有相同的不可分次数。
由分歧指数的乘法性，只需在 $x$ 是仿射直线上的点时证明结论。

\medskip\noindent
假设 $X = \mathbf{A}^1_k$。此时 $X$ 在 $x$ 处的局部环形如
$$
\mathcal{O}_{X, x} = k[t]_{(P)}
$$
，其中 $P$ 是 $k$ 上不可约的首一多项式。
于是有 $P(t) = Q(t^q)$，其中 $Q \in k[t]$ 是某个可分多项式；
见《域》，引理
\ref{fields-lemma-irreducible-polynomials}。
注意，$\kappa(x) = k[t]/(P)$ 的不可分次数为 $q$，这是在 $k$ 上而言。
另一方面，在 $\overline{k}$ 上可以分解
$Q(t) = \prod (t - \alpha_i)$，其中 $\alpha_i$ 两两不同。
写成 $\alpha_i = \beta_i^q$，其中存在唯一的
$\beta_i \in \overline{k}$。于是点 $\overline{x}$
对应于其中一个 $\beta_i$。现在可以得出结论，因为
$$
k[t]_{(P)} \longrightarrow \overline{k}[t]_{(t - \beta_i)}
$$
的分歧指数确实等于 $q$：一致化元 $P$ 映为
$(t - \beta_i)^q$ 乘以一个单位。
\end{proof}





\section{推出}
\label{curves-section-pushouts}

\noindent
设 $k$ 为域。考察实线图
$$
\xymatrix{
Z' \ar[d] \ar[r]_{i'} & X' \ar@{..>}[d]^a \\
Z \ar@{..>}[r]^i & X
}
$$
；这是 $k$ 上概形的图，并满足
\begin{enumerate}
\item[(a)] $X'$ 是 $k$ 上分离的有限型概形，且维数 $\leq 1$；
\item[(b)] $i : Z' \to X'$ 是闭浸入；
\item[(c)] $Z'$ 与 $Z$ 在 $\Spec(k)$ 上有限；并且
\item[(d)] $Z' \to Z$ 满射。
\end{enumerate}
在这种情形下，$X'$ 的每个有限点集都包含于某个仿射开集中；
见《簇》，命题
\ref{varieties-proposition-finite-set-of-points-of-codim-1-in-affine}。
因此满足《态射进阶》，命题
\ref{more-morphisms-proposition-pushout-along-closed-immersion-and-integral}
的假设，并得到如下结论：
\begin{enumerate}
\item 推出 $X = Z \amalg_{Z'} X'$ 在概形范畴中存在；
\item $i : Z \to X$ 是闭浸入；
\item $a : X' \to X$ 是整满态射；
\item 由《态射进阶》，引理
\ref{more-morphisms-lemma-pushout-separated}，
$X \to \Spec(k)$ 是分离的；
\item 由《态射进阶》，引理
\ref{more-morphisms-lemma-pushout-finite-type}，
$X \to \Spec(k)$ 是有限型的；
\item 因此，由《态射》，引理
\ref{morphisms-lemma-finite-integral} 与
\ref{morphisms-lemma-permanence-finite-type}，
$a : X' \to X$ 是有限的；
\item 若 $X' \to \Spec(k)$ 固有，则由《态射》，引理
\ref{morphisms-lemma-image-proper-is-proper}，
$X \to \Spec(k)$ 固有。
\end{enumerate}

\noindent
下面的引理可以作显著推广。

\begin{lemma}
\label{curves-lemma-complete-local-ring-pushout}
在上述情形中，设 $Z = \Spec(k')$，其中 $k'$ 为域；并设
$Z' = \Spec(k'_1 \times \ldots \times k'_n)$，其中 $k'_i/k'$
是有限域扩张。设 $x \in X$ 为 $Z \to X$ 的像，并设
$x'_i \in X'$ 为 $\Spec(k'_i) \to X'$ 的像。则有纤维积图
$$
\xymatrix{
\prod\nolimits_{i = 1, \ldots, n} k'_i &
\prod\nolimits_{i = 1, \ldots, n} \mathcal{O}_{X', x'_i}^\wedge \ar[l] \\
k' \ar[u] &
\mathcal{O}_{X, x}^\wedge \ar[u] \ar[l]
}
$$
，其中水平箭头由到剩余域的映射给出。
\end{lemma}

\begin{proof}
选取仿射开邻域 $\Spec(A)$，它是 $x$ 在 $X$ 中的邻域。
设 $\Spec(A') \subset X'$ 为其逆像。由构造，有纤维积图
$$
\xymatrix{
\prod\nolimits_{i = 1, \ldots, n} k'_i &
A' \ar[l] \\
k' \ar[u] &
A \ar[u] \ar[l]
}
$$
。由于所有对象在 $A$ 上都有限，考察极大理想
$\mathfrak m \subset A$，它同 $x$ 相应。对它作完备化后，该图仍为纤维积图
（《代数》，引理 \ref{algebra-lemma-completion-flat}）。
最后应用《代数》，引理
\ref{algebra-lemma-completion-finite-extension}，
以识别 $A'$ 的完备化。
\end{proof}




\section{粘合与压缩}
\label{curves-section-glueing-squishing}

\noindent
以下用 $k[\epsilon]$ 表示 $k$ 上的对偶数代数，其定义见《簇》，定义
\ref{varieties-definition-dual-numbers}。

\begin{lemma}
\label{curves-lemma-no-in-between-over-k}
设 $k$ 为代数闭域。设 $k \subset A$ 为环扩张，并且 $A$
恰有两个 $k$-子代数。则 $A = k \times k$ 或
$A = k[\epsilon]$。
\end{lemma}

\begin{proof}
假设意味着 $k \not = A$，并且任意子环
$k \subset C \subset A$ 都等于 $k$ 或 $A$。
取 $t \in A$ 且 $t \not \in k$。则 $A$ 由 $t$ 在 $k$ 上生成，
故有 $A = k[x]/I$，其中 $I$ 是某个理想。
若 $I = (0)$，则会有不允许出现的子代数 $k[x^2]$。
否则，$I$ 由某个首一多项式 $P$ 生成。写成
$P = \prod_{i = 1}^d (t - a_i)$。若 $d > 2$，则由
$(t - a_1)(t - a_2)$ 生成的子代数给出矛盾。
所以 $d = 2$。若 $a_1 \not = a_2$，则 $A = k \times k$；
若 $a_1 = a_2$，则 $A = k[\epsilon]$。
\end{proof}

\begin{example}[粘合点]
\label{curves-example-glue-points}
设 $k$ 为代数闭域。设 $f : X' \to X$ 为代数
$k$-概形之间的态射。若下列条件成立，就称 $X$
由粘合 $a$ 与 $b$ 得到；这两点位于 $X'$ 中：
\begin{enumerate}
\item $a, b \in X'(k)$ 是不同的点，它们映到同一点
$x \in X(k)$；
\item $f$ 有限，并且
$f^{-1}(X \setminus \{x\}) \to X \setminus \{x\}$ 是同构；
\item 有短正合列
$$
0 \to \mathcal{O}_X \to f_*\mathcal{O}_{X'} \xrightarrow{a - b} x_*k \to 0
$$
，其中右端箭头把局部截面 $h$（它属于 $f_*\mathcal{O}_{X'}$）
映到差 $h(a) - h(b) \in k$。
\end{enumerate}
若属于这种情形，则还有短正合列
$$
0 \to \mathcal{O}_X^* \to f_*\mathcal{O}_{X'}^*
\xrightarrow{ab^{-1}} x_*k^* \to 0
$$
，其中右端箭头把局部截面 $h$（它属于 $f_*\mathcal{O}_{X'}^*$）
映到乘法差 $h(a)h(b)^{-1} \in k^*$。
\end{example}

\begin{example}[压缩切向量]
\label{curves-example-squish-tangent-vector}
设 $k$ 为代数闭域。设 $f : X' \to X$ 为代数
$k$-概形之间的态射。若下列条件成立，就称 $X$
由压缩切向量 $\vartheta$ 得到；该切向量位于 $X'$ 中：
\begin{enumerate}
\item $\vartheta : \Spec(k[\epsilon]) \to X'$ 是 $k$ 上的闭浸入，
且 $f \circ \vartheta$ 经由一点 $x \in X(k)$ 分解；
\item $f$ 有限，并且
$f^{-1}(X \setminus \{x\}) \to X \setminus \{x\}$ 是同构；
\item 有短正合列
$$
0 \to \mathcal{O}_X \to f_*\mathcal{O}_{X'} \xrightarrow{\vartheta} x_*k \to 0
$$
，其中右端箭头把局部截面 $h$（它属于 $f_*\mathcal{O}_{X'}$）
映到 $\epsilon$ 的系数；该系数取自
$\vartheta^\sharp(h) \in k[\epsilon]$。
\end{enumerate}
若属于这种情形，则还有短正合列
$$
0 \to \mathcal{O}_X^* \to f_*\mathcal{O}_{X'}^*
\xrightarrow{\vartheta} x_*k \to 0
$$
，其中右端箭头把局部截面 $h$（它属于 $f_*\mathcal{O}_{X'}^*$）
映到 $\text{d}\log(\vartheta^\sharp(h))$，这里
$\text{d}\log : k[\epsilon]^* \to k$
是把 $a + b\epsilon$ 映到 $b/a \in k$ 的交换群同态。
\end{example}

\begin{lemma}
\label{curves-lemma-factor-almost-isomorphism}
设 $k$ 为代数闭域。设 $f : X' \to X$ 为代数 $k$-概形之间的
有限态射，满足 $\mathcal{O}_X \subset f_*\mathcal{O}_{X'}$，
并且 $f$ 在某个有限点集之外是同构。则存在分解
$$
X' = X_n \to X_{n - 1} \to \ldots \to X_1 \to X_0 = X
$$
，使每个 $X_i \to X_{i - 1}$ 或者是两个点的粘合，
或者是一个切向量的压缩
（见例 \ref{curves-example-glue-points} 与
\ref{curves-example-squish-tangent-vector}）。
\end{lemma}

\begin{proof}
设 $U \subset X$ 为使 $f$ 成为同构的最大开集。
则 $X \setminus U = \{x_1, \ldots, x_n\}$，其中
$x_i \in X(k)$。我们将考察分解 $X' \to Y \to X$；
这是 $f$ 的分解，要求两个态射都有限，并且
$$
\mathcal{O}_X \subset g_*\mathcal{O}_Y \subset f_*\mathcal{O}_{X'}
$$
，其中 $g : Y \to X$ 为给定态射。由假设，
$\mathcal{O}_{X, x} \to (f_*\mathcal{O}_{X'})_x$
是同构，除非有 $x = x_i$，其中 $i$ 为某个指标。因此余核
$$
f_*\mathcal{O}_{X'}/\mathcal{O}_X = \bigoplus \mathcal{Q}_i
$$
是摩天层 $\mathcal{Q}_i$ 之直和；它们支撑在
$x_1, \ldots, x_n$ 上。由于所显示的商是凝聚
$\mathcal{O}_X$-模，可知 $\mathcal{Q}_i$ 在
$\mathcal{O}_{X, x_i}$ 上具有有限长度。因此可以对这些长度之和，
也就是整个余核的长度，作归纳。

\medskip\noindent
若 $n > 1$，则可以定义 $\mathcal{O}_X$-子代数
$\mathcal{A} \subset f_*\mathcal{O}_{X'}$，方法是取
$\mathcal{Q}_1$ 的逆像。
这将给出非平凡分解，于是可用归纳。

\medskip\noindent
假设 $n = 1$。简记 $x = x_1$。考察有限 $k$-代数扩张
$$
A = \mathcal{O}_{X, x} \subset (f_*\mathcal{O}_{X'})_x = B
$$
。注意，$\mathcal{Q} = \mathcal{Q}_1$ 是值为 $B/A$ 的摩天层。
有 $k$-子代数 $A \subset A + \mathfrak m_A B \subset B$。
若两个包含都严格，则得到非平凡分解，并如上用归纳。
若 $A + \mathfrak m_A B = B$，则由 Nakayama 有 $A = B$；
此时 $f$ 是同构，无事可证。由此可假设
$B = A + \mathfrak m_A B$。置 $C = B/\mathfrak m_A B$。
若 $C$ 有多于 $2$ 个 $k$-子代数，则得到介于 $A$ 与 $B$
之间的子代数，方法是在 $B$ 中取逆像。因此可假设
$C$ 恰有 $2$ 个 $k$-子代数。由引理
\ref{curves-lemma-no-in-between-over-k}，有 $C = k \times k$
或 $C = k[\epsilon]$。在这种情形下，$f$ 相应地是两个点的粘合，
或一个切向量的压缩。
\end{proof}

\begin{lemma}
\label{curves-lemma-glue-points}
设 $k$ 为代数闭域。若 $f : X' \to X$ 是如例
\ref{curves-example-glue-points} 中由粘合两点 $a, b$ 得到的态射，则有正合列
$$
k^* \to \Pic(X) \to \Pic(X') \to 0
$$
。若 $a$ 与 $b$ 位于 $X'$ 的不同连通分支上，则第一个映射为零；
若 $X'$ 固有且 $a$ 与 $b$ 位于 $X'$ 的同一连通分支上，
则第一个映射为单射。
\end{lemma}

\begin{proof}
由《簇》，引理 \ref{varieties-lemma-surjective-pic-birational-finite}，
映射 $\Pic(X) \to \Pic(X')$ 满射。使用短正合列
$$
0 \to \mathcal{O}_X^* \to f_*\mathcal{O}_{X'}^*
\xrightarrow{ab^{-1}} x_*k^* \to 0
$$
，得到
$$
H^0(X', \mathcal{O}_{X'}^*) \xrightarrow{ab^{-1}} k^* \to
H^1(X, \mathcal{O}_X^*) \to H^1(X, f_*\mathcal{O}_{X'}^*)
$$
。有 $H^1(X, f_*\mathcal{O}_{X'}^*) \subset H^1(X', \mathcal{O}_{X'}^*)$
（例如由 Leray 谱序列；见《上同调》，引理
\ref{cohomology-lemma-Leray}）。
因此 $\Pic(X) \to \Pic(X')$ 的核是
$ab^{-1} : H^0(X', \mathcal{O}_{X'}^*) \to k^*$ 的余核。
若 $a$ 与 $b$ 位于 $X'$ 的不同连通分支上，则
$ab^{-1}$ 满射。由于 $k$ 代数闭，$k$ 上约化连通固有概形的
任意正则函数都来自 $k$ 的一个元素；见《簇》，引理
\ref{varieties-lemma-proper-geometrically-reduced-global-sections}。
因此 $ab^{-1}$ 为零，如果 $X'$ 固有且 $a$ 与 $b$
位于同一连通分支上。
\end{proof}

\begin{lemma}
\label{curves-lemma-squish-tangent-vector}
设 $k$ 为代数闭域。若 $f : X' \to X$ 是如例
\ref{curves-example-squish-tangent-vector} 中对切向量
$\vartheta$ 的压缩，则有正合列
$$
(k, +) \to \Pic(X) \to \Pic(X') \to 0
$$
，并且若 $X'$ 固有且约化，则第一个映射为单射。
\end{lemma}

\begin{proof}
由《簇》，引理 \ref{varieties-lemma-surjective-pic-birational-finite}，
映射 $\Pic(X) \to \Pic(X')$ 满射。使用例
\ref{curves-example-squish-tangent-vector} 的短正合列
$$
0 \to \mathcal{O}_X^* \to f_*\mathcal{O}_{X'}^*
\xrightarrow{\vartheta} x_*k \to 0
$$
，得到
$$
H^0(X', \mathcal{O}_{X'}^*) \xrightarrow{\vartheta} k \to
H^1(X, \mathcal{O}_X^*) \to H^1(X, f_*\mathcal{O}_{X'}^*)
$$
。有 $H^1(X, f_*\mathcal{O}_{X'}^*) \subset H^1(X', \mathcal{O}_{X'}^*)$
（例如由 Leray 谱序列；见《上同调》，引理
\ref{cohomology-lemma-Leray}）。
因此 $\Pic(X) \to \Pic(X')$ 的核是映射
$\vartheta : H^0(X', \mathcal{O}_{X'}^*) \to k$ 的余核。
由于 $k$ 代数闭，$k$ 上约化连通固有概形的任意正则函数
都来自 $k$ 的一个元素；见《簇》，引理
\ref{varieties-lemma-proper-geometrically-reduced-global-sections}。
由此得到引理的最后一个断言。
\end{proof}


\section{多重交叉与结点奇点}
\label{curves-section-multicross}

\noindent
本节讨论最简单的一类曲线奇点。

\medskip\noindent
设 $k$ 为域。考察完备局部 $k$-代数
\begin{equation}
\label{curves-equation-multicross}
A = \{(f_1, \ldots, f_n) \in k[[t]] \times \ldots \times k[[t]] \mid
f_1(0) = \ldots = f_n(0)\}
\end{equation}
采用《簇》，定义 \ref{varieties-definition-wedge}
引入的术语，$A$ 是 $n$ 个 $1$ 元幂级数环（系数域为 $k$）
的楔和。注意，
$k[[t]] \times \ldots \times k[[t]]$
是 $A$ 在其全分式环中的整闭包。因此 $\delta$-不变量
在 $A$ 上的值为 $n - 1$。存在同构
$$
k[[x_1, \ldots, x_n]]/(\{x_ix_j\}_{i \not = j}) \longrightarrow A
$$
，它把 $x_i$ 映到 $(0, \ldots, 0, t, 0, \ldots, 0)$；
后者位于 $A$ 中。由此
$\dim(A) = 1$，且
$\dim_k \mathfrak m/\mathfrak m^2 = n$.
特别地，$A$ 正则当且仅当 $n = 1$。

\begin{lemma}
\label{curves-lemma-multicross-algebra}
设 $k$ 为可分闭域。设 $A$ 为 $1$ 维约化 Nagata 局部
$k$-代数，其剩余域为 $k$。则
$$
\delta\text{-invariant }A \geq \text{number of branches of }A - 1
$$
若等号成立，则 $A^\wedge$ 如式 (\ref{curves-equation-multicross}) 所示。
\end{lemma}

\begin{proof}
由于 $A$ 的剩余域可分闭，$A$ 的分支数等于 $A$ 的几何分支数；
见《代数补篇》，定义 \ref{more-algebra-definition-number-of-branches}。
不等式由《簇》，引理
\ref{varieties-lemma-delta-number-branches-inequality} 给出。
假设等号成立。可以把 $A$ 替换为 $A$ 的完备化；
这不改变分支数或 $\delta$-不变量；见《代数补篇》，引理
\ref{more-algebra-lemma-one-dimensional-number-of-branches}
以及《簇》，引理 \ref{varieties-lemma-delta-same-after-completion}。
此时 $A$ 严格 Hensel；见《代数》，引理
\ref{algebra-lemma-complete-henselian}。由《簇》，引理
\ref{varieties-lemma-delta-number-branches-inequality-sh}，
$A$ 是若干完备离散赋值环的楔和。由《代数》，引理
\ref{algebra-lemma-regular-complete-containing-coefficient-field}，
其中每一个都同构于 $k[[t]]$。所以 $A$ 如式
(\ref{curves-equation-multicross}) 所示。
\end{proof}

\begin{definition}
\label{curves-definition-multicross}
设 $k$ 为代数闭域。设 $X$ 为代数 $1$ 维 $k$-概形，并设
$x \in X$ 为闭点。称 $x$ 定义了一个{\it 多重交叉奇点}，
如果完备化
$\mathcal{O}_{X, x}^\wedge$
对某个 $n \geq 2$ 同构于式 (\ref{curves-equation-multicross})。
称 $x$ 是一个{\it 结点}或{\it 普通二重点}，也称它
{\it 定义了一个结点奇点}，如果 $n = 2$。
\end{definition}

\noindent
在某种意义上，这些是代数闭域上曲线所能具有的最简单奇点。

\begin{lemma}
\label{curves-lemma-multicross}
设 $k$ 为代数闭域。设 $X$ 为约化代数 $1$ 维 $k$-概形，
并设 $x \in X$。下列条件等价：
\begin{enumerate}
\item $x$ 定义一个多重交叉奇点；
\item $\delta$-不变量（对 $X$ 在 $x$ 处而言）等于
$X$ 在 $x$ 处的分支数减 $1$；
\item 存在态射序列
$U_n \to U_{n - 1} \to \ldots \to U_0 = U \subset X$
，其中 $U$ 是 $x$ 的开邻域，$U_n$ 非奇异，并且每个
$U_i \to U_{i - 1}$ 都是例 \ref{curves-example-glue-points}
所述的两点粘合。
\end{enumerate}
\end{lemma}

\begin{proof}
条件 (1) 与 (2) 的等价性就是引理 \ref{curves-lemma-multicross-algebra}。

\medskip\noindent
假设 (3) 成立。我们对 $i$ 作降序归纳，证明 $U_i$ 的所有奇点
都是多重交叉奇点。对 $U_n$ 这显然成立，因为 $U_n$ 没有奇点。
若 $U_i$ 是把 $U_{i + 1}$ 中的 $a, b \in U_{i + 1}$
粘合为一点 $c \in U_i$ 而得到的，则
$$
\mathcal{O}_{U_i, c}^\wedge \subset
\mathcal{O}_{U_{i + 1}, a}^\wedge \times \mathcal{O}_{U_{i + 1}, b}^\wedge
$$
恰由在 $k$ 中剩余类相同的元素组成。因此 $c$ 处的分支数等于
$a$ 与 $b$ 处分支数之和，而 $\delta$-不变量在 $c$ 处等于
$\delta$-不变量在 $a$ 与 $b$ 处的值之和再加 $1$
（因为所显示的包含之余维为 $1$）。这就按要求证明了 (2)。

\medskip\noindent
假设等价条件 (1) 与 (2) 成立。可以选取开集
$U \subset X$，使得 $x$ 是 $U$ 的唯一奇点。对正规化态射
应用引理 \ref{curves-lemma-factor-almost-isomorphism}：
$$
U^\nu = U_n \to U_{n - 1} \to \ldots \to U_1 \to U_0 = U
$$
只需证明每一步都没有压缩切向量。假设
$U_{i + 1} \to U_i$ 是第一次出现这种情形的步骤，即 $i$
最小，并且该步骤压缩切向量 $\theta$，其所在点为
$u' \in U_{i + 1}$，位于 $u \in U_i$ 上方。按上述论证可知，
$u_i$ 是多重交叉奇点（因为各映射
$U_i \to \ldots \to U_0$ 都是成对粘合点）。然而此时
$u'$ 与 $u$ 处的分支数相同，而 $\delta$-不变量在
$U_i$ 的 $u$ 处之值为 $1$ 加上 $\delta$-不变量在
$U_{i + 1}$ 的 $u'$ 处之值。由引理
\ref{curves-lemma-multicross-algebra}，
这说明 $u$ 不可能是多重交叉奇点，矛盾。
\end{proof}

\begin{lemma}
\label{curves-lemma-multicross-gorenstein-is-nodal}
设 $k$ 为代数闭域。设 $X$ 为约化代数 $1$ 维 $k$-概形。
设 $x \in X$ 为多重交叉奇点（定义
\ref{curves-definition-multicross}）。若 $X$ 为 Gorenstein，则 $x$ 为结点。
\end{lemma}

\begin{proof}
映射 $\mathcal{O}_{X, x} \to \mathcal{O}_{X, x}^\wedge$
平坦，并且在如下意义下非分歧：
$\kappa(x) = \mathcal{O}_{X, x}^\wedge/\mathfrak m_x \mathcal{O}_{X, x}^\wedge$.
（见《代数补篇》，节
\ref{more-algebra-section-permanence-completion}。）因此，$X$
为 Gorenstein 蕴含 $\mathcal{O}_{X, x}$ 为 Gorenstein；再由
《对偶复形》，引理 \ref{dualizing-lemma-flat-under-gorenstein}，
可推出 $\mathcal{O}_{X, x}^\wedge$ 为 Gorenstein。
所以只需证明式 (\ref{curves-equation-multicross}) 中的环 $A$
在 $n \geq 2$ 时为 Gorenstein 当且仅当 $n = 2$。

\medskip\noindent
若 $n = 2$，则 $A = k[[x, y]]/(xy)$ 是完全交，从而为
Gorenstein。例如，把《概形的对偶性》，引理
\ref{duality-lemma-gorenstein-lci} 应用于 $k[[x, y]] \to A$，
并使用正则局部环 $k[[x, y]]$ 由《对偶复形》，引理
\ref{dualizing-lemma-regular-gorenstein} 为 Gorenstein，即得此结论。

\medskip\noindent
假设 $n > 2$。若 $A$ 为 Gorenstein，则 $A$ 将是 $A$ 上的
对偶复形（《概形的对偶性》，定义
\ref{duality-definition-gorenstein}）。于是
$R\Hom(k, A)$ 将等于 $k[n]$，其中 $n \in \mathbf{Z}$；
见《对偶复形》，引理 \ref{dualizing-lemma-find-function}。
由此可推出 $\Ext^1_A(k, A) \cong k$ 或
$\Ext^1_A(k, A) = 0$（取决于 $n$ 的值；事实上
$n$ 必须为 $-1$，但这一点在此无关紧要）。利用正合列
$$
0 \to \mathfrak m_A \to A \to k \to 0
$$
可得
$$
\Ext^1_A(k, A) = \Hom_A(\mathfrak m_A, A)/A
$$
，其中 $A \to \Hom_A(\mathfrak m_A, A)$ 由
$a \mapsto (a' \mapsto aa')$ 给出。设
$e_i \in \Hom_A(\mathfrak m_A, A)$ 是把
$(f_1, \ldots, f_n) \in \mathfrak m_A$ 映到
$(0, \ldots, 0, f_i, 0, \ldots, 0)$ 的元素。容易验证，
$e_1, \ldots, e_{n - 1}$ 关于 $k$ 线性无关，其像位于
$\Hom_A(\mathfrak m_A, A)/A$ 中。因此
$\dim_k \Ext^1_A(k, A) \geq n - 1 \geq 2$，证毕。
（注意，$e_1 + \ldots + e_n$ 是 $1$ 在映射
$A \to \Hom_A(\mathfrak m_A, A)$ 下的像。）
\end{proof}




\section{Picard 群中的挠元}
\label{curves-section-torsion-in-pic}

\noindent
本节给出域上 $1$ 维固有概形的 Picard 群中挠元的界。
这将在研究曲线的半稳定约化时使用。

\medskip\noindent
似乎没有初等方法可得引理
\ref{curves-lemma-torsion-picard-smooth-projective} 的结论。
分析其证明可见两个关键要素：(1) 存在一个分类光滑射影曲线上
零次可逆层的阿贝尔簇；(2) 可以确定阿贝尔簇上挠点的结构。

\begin{lemma}
\label{curves-lemma-torsion-picard-smooth-projective}
设 $k$ 为代数闭域。设 $X$ 为亏格为 $g$ 的光滑射影曲线，
其底域为 $k$。
\begin{enumerate}
\item 若 $n \geq 1$ 在 $k$ 中可逆，则
$\Pic(X)[n] \cong (\mathbf{Z}/n\mathbf{Z})^{\oplus 2g}$.
\item 若 $k$ 的特征为 $p > 0$，则存在整数
$0 \leq f \leq g$，使得
$\Pic(X)[p^m] \cong (\mathbf{Z}/p^m\mathbf{Z})^{\oplus f}$ 对
所有 $m \geq 1$ 都成立。
\end{enumerate}
\end{lemma}

\begin{proof}
用 $\Pic^0(X) \subset \Pic(X)$ 表示由次数为 $0$ 的可逆层
组成的子群。换言之，有短正合列
$$
0 \to \Pic^0(X) \to \Pic(X) \xrightarrow{\deg} \mathbf{Z} \to 0.
$$
群 $\Pic^0(X)$ 是 $k$-点集；它取自群概形
$\underline{\Picardfunctor}^0_{X/k}$。见
《曲线的 Picard 概形》，引理 \ref{pic-lemma-picard-pieces}。
同一引理还说明 $\underline{\Picardfunctor}^0_{X/k}$
是 $g$ 维阿贝尔簇，定义在 $k$ 上；其定义见《群胚》，定义
\ref{groupoids-definition-abelian-variety}。于是结论来自
《群胚》，命题 \ref{groupoids-proposition-review-abelian-varieties}。
\end{proof}

\begin{lemma}
\label{curves-lemma-torsion-picard-becomes-visible}
设 $k$ 为域。设 $n$ 与 $k$ 的特征互素。设 $X$ 为 $k$ 上
光滑固有曲线，满足 $H^0(X, \mathcal{O}_X) = k$，且亏格为 $g$。
\begin{enumerate}
\item 若 $g = 1$，则存在有限可分扩张 $k'/k$，使得
$X_{k'}$ 有一个 $k'$-有理点，且
$\Pic(X_{k'})[n] \cong (\mathbf{Z}/n\mathbf{Z})^{\oplus 2}$.
\item 若 $g \geq 2$，则存在有限可分扩张 $k'/k$，满足
$[k' : k] \leq (2g - 2)(n^{2g})!$，使得 $X_{k'}$
有一个 $k'$-有理点，且
$\Pic(X_{k'})[n] \cong (\mathbf{Z}/n\mathbf{Z})^{\oplus 2g}$.
\end{enumerate}
\end{lemma}

\begin{proof}
假设 $g \geq 2$。首先可选取次数至多为 $2g - 2$ 的有限可分扩张，
使 $X$ 获得一个有理点；见引理
\ref{curves-lemma-point-over-separable-extension}。因此可假设 $X$
有一个 $k$-有理点 $x \in X(k)$，但此时必须在如下界下证明引理：
$[k' : k] \leq (n^{2g})!$.
设 $k \subset k^{sep} \subset \overline{k}$ 是代数闭包中的
一个可分代数闭包。由引理
\ref{curves-lemma-torsion-picard-smooth-projective} 有
$$
\Pic(X_{\overline{k}})[n] \cong (\mathbf{Z}/n\mathbf{Z})^{\oplus 2g}
$$
由《曲线的 Picard 概形》，引理 \ref{pic-lemma-torsion-descends}，
可得
$$
\Pic(X_{k^{sep}})[n] \cong (\mathbf{Z}/n\mathbf{Z})^{\oplus 2g}
$$
再由《曲线的 Picard 概形》，引理
\ref{pic-lemma-torsion-descends}，有连续作用
$$
\text{Gal}(k^{sep}/k)
\longrightarrow
\text{Aut}(\Pic(X_{k^{sep}})[n]
$$
，而对该映射核的固定域 $k'$，引理成立。由于作用连续，其核为开，
所以 $k'/k$ 有限。由 Galois 理论，
$\text{Gal}(k'/k)$ 是所显示箭头的像。基数为 $n^{2g}$ 的集合
之置换群的基数为 $(n^{2g})!$，故再由 Galois 理论可得
$[k' : k] \leq (n^{2g})!$。
（当然，这实际上以 $|\text{GL}_{2g}(\mathbf{Z}/n\mathbf{Z})|$
为界证明了引理，但这里所需的只是存在某个界。）

\medskip\noindent
若亏格为 $1$，则使 $X$ 获得有理点的有限可分域扩张之次数
不存在统一上界（略去细节）。不过，这样的扩张确实存在，例如由
《簇》，引理
\ref{varieties-lemma-smooth-separable-closed-points-dense} 可知。
其余证明与 $g \geq 2$ 的情形相同。
\end{proof}

\begin{proposition}
\label{curves-proposition-torsion-picard-reduced-proper}
设 $k$ 为代数闭域。设 $X$ 为 $k$ 上约化、连通且维数为 $1$
的固有概形。设 $g$ 为 $X$ 的亏格，并设 $g_{geom}$ 为 $X$
各不可约分支之几何亏格的和。对任意素数 $\ell$，若它不同于
$k$ 的特征，则有
$$
\dim_{\mathbf{F}_\ell} \Pic(X)[\ell]
\leq g + g_{geom}
$$
，且等号成立当且仅当 $X$ 的所有奇点都是多重交叉奇点。
\end{proposition}

\begin{proof}
设 $\nu : X^\nu \to X$ 为正规化（《簇》，引理
\ref{varieties-lemma-prepare-delta-invariant}）。选取分解
$$
X^\nu = X_n \to X_{n - 1} \to \ldots \to X_1 \to X_0 = X
$$
，如引理 \ref{curves-lemma-factor-almost-isomorphism} 所述。记
$h^0_i = \dim_k H^0(X_i, \mathcal{O}_{X_i})$ 以及
$h^1_i = \dim_k H^1(X_i, \mathcal{O}_{X_i})$。
由引理 \ref{curves-lemma-glue-points} 与
\ref{curves-lemma-squish-tangent-vector}，对每个 $n > i \geq 0$，
出现下列三种情形之一：
\begin{enumerate}
\item $X_i$ 由粘合位于不同连通分支上的
$a, b \in X_{i + 1}$ 得到；此时
$\Pic(X_i) = \Pic(X_{i + 1})$,
$h^0_{i + 1} = h^0_i + 1$, $h^1_{i + 1} = h^1_i$,
\item $X_i$ 由粘合位于同一连通分支上的
$a, b \in X_{i + 1}$ 得到；此时有短正合列
$$
0 \to k^* \to \Pic(X_i) \to \Pic(X_{i + 1}) \to 0,
$$
，且 $h^0_{i + 1} = h^0_i$，$h^1_{i + 1} = h^1_i - 1$；
\item $X_i$ 由压缩 $X_{i + 1}$ 中的一个切向量得到；
此时有短正合列
$$
0 \to (k, +) \to \Pic(X_i) \to \Pic(X_{i + 1}) \to 0,
$$
，且 $h^0_{i + 1} = h^0_i$，$h^1_{i + 1} = h^1_i - 1$。
\end{enumerate}
关于结构层上同调群维数的断言，使用例
\ref{curves-example-glue-points} 与
\ref{curves-example-squish-tangent-vector} 中的正合列即可证明。
由于 $k$ 代数闭且其特征与 $\ell$ 互素，可知
$(k, +)$ 与 $k^*$ 都是 $\ell$-可除群，并且其
$\ell$-挠子群满足 $(k, +)[\ell] = 0$ 以及
$k^*[\ell] \cong \mathbf{F}_\ell$。因此
$$
\dim_{\mathbf{F}_\ell} \Pic(X_{i + 1})[\ell] -
\dim_{\mathbf{F}_\ell}\Pic(X_i)[\ell]
$$
除情形 (2) 中等于 $-1$ 外，该量均为零。此过程结束时得到正规化
$X^\nu = X_n$，它是 $k$ 上光滑射影曲线的不交并。因此有
\begin{enumerate}
\item $h^1_n = g_{geom}$；并且
\item $\dim_{\mathbf{F}_\ell} \Pic(X_n)[\ell] = 2g_{geom}$.
\end{enumerate}
最后一个等式来自引理
\ref{curves-lemma-torsion-picard-smooth-projective}。由于 $g = h^1_0$，
可知类型 (2) 与 (3) 的步骤数至多为
$h^1_0 - h^1_n = g - g_{geom}$。由上述秩差的计算可得
$$
\dim_{\mathbf{F}_\ell} \Pic(X)[\ell] \leq
g - g_{geom} + 2g_{geom} = g + g_{geom}
$$
，且等号成立当且仅当不出现类型 (3) 的步骤。由引理
\ref{curves-lemma-multicross}，这恰好意味着 $X$ 的所有奇点都是
多重交叉奇点。反过来，若所有奇点都是多重交叉奇点，则引理
\ref{curves-lemma-multicross} 保证可以找到如上的序列
$X^\nu = X_n \to \ldots \to X_0 = X$，其中不出现类型 (3)
的步骤，于是引理中的等号成立（只需对每个点粘合局部序列，
即可得到一个同时适用于 $x$ 的所有奇点的序列；略去若干细节）。
\end{proof}





\section{亏格与几何亏格的比较}
\label{curves-section-genus-geometric-genus}

\noindent
设 $k$ 为域，其代数闭包为 $\overline{k}$。设 $X$ 为
维数 $\leq 1$ 的固有概形，其底域为 $k$。定义 $g_{geom}(X/k)$ 为
$X_{\overline{k}}$ 的所有 $1$ 维不可约分支之几何亏格的和。

\begin{lemma}
\label{curves-lemma-bound-geometric-genus}
设 $k$ 为域。设 $X$ 为维数 $\leq 1$ 的固有概形，其底域为 $k$。则
$$
g_{geom}(X/k) = \sum\nolimits_{C \subset X} g_{geom}(C/k)
$$
，其中求和遍历不可约分支 $C \subset X$，其维数为 $1$。
\end{lemma}

\begin{proof}
这直接来自定义以及如下事实：不可约分支 $\overline{Z}$
（属于 $X_{\overline{k}}$）满射到不可约分支 $Z$（属于 $X$）
（《簇》，引理 \ref{varieties-lemma-image-irreducible}），
且两者维数相同（把《态射》，引理
\ref{morphisms-lemma-dimension-fibre-after-base-change}
用于 $\overline{Z}$ 的一般点）。
\end{proof}

\begin{lemma}
\label{curves-lemma-geometric-genus-normalization}
设 $k$ 为域。设 $X$ 为维数 $\leq 1$ 的固有概形，其底域为 $k$。则
\begin{enumerate}
\item 有 $g_{geom}(X/k) = g_{geom}(X_{red}/k)$。
\item 若 $X' \to X$ 是双有理固有态射，则
$g_{geom}(X'/k) = g_{geom}(X/k)$.
\item 若 $X^\nu \to X$ 是正规化态射，则
$g_{geom}(X^\nu/k) = g_{geom}(X/k)$.
\end{enumerate}
\end{lemma}

\begin{proof}
第 (1) 点直接来自引理 \ref{curves-lemma-bound-geometric-genus}。
若 $X' \to X$ 双有理且固有，则它是有限态射，并在某个稠密开集上
为同构（见《簇》，引理
\ref{varieties-lemma-finite-in-codim-1} 与
\ref{varieties-lemma-modification-normal-iso-over-codimension-1}）。
因此 $X'_{\overline{k}} \to X_{\overline{k}}$ 在某个稠密开集上
为同构。所以 $X'_{\overline{k}}$ 与 $X_{\overline{k}}$
的不可约分支一一对应，而且对应分支的函数域同构。特别地，
这些分支的非奇异射影模型同构，因而具有相同的几何亏格。
这证明了 (2)。第 (3) 点来自 (1)、(2) 以及
$X^\nu \to X_{red}$ 双有理这一事实（《态射》，引理
\ref{morphisms-lemma-normalization-birational}）。
\end{proof}

\begin{lemma}
\label{curves-lemma-genus-goes-down}
设 $k$ 为域。设 $X$ 为维数 $\leq 1$ 的固有概形，其底域为 $k$。
设 $f : Y \to X$ 为有限态射，并且存在稠密开集 $U \subset X$，
使得 $f$ 在其上为闭浸入。则
$$
\dim_k H^1(X, \mathcal{O}_X) \geq \dim_k H^1(Y, \mathcal{O}_Y)
$$
\end{lemma}

\begin{proof}
考察正合列
$$
0 \to \mathcal{G} \to \mathcal{O}_X \to f_*\mathcal{O}_Y \to \mathcal{F} \to 0
$$
，它由 $X$ 上的凝聚层组成。由假设，$\mathcal{F}$ 支撑在有限多个
闭点上，因而高次上同调消失（《簇》，引理
\ref{varieties-lemma-chi-tensor-finite}）。另一方面，由《上同调》，
命题 \ref{cohomology-proposition-vanishing-Noetherian} 有
$H^2(X, \mathcal{G}) = 0$。于是形式地可知诱导映射
$H^1(X, \mathcal{O}_X) \to H^1(X, f_*\mathcal{O}_Y)$
满射。由于 $H^1(X, f_*\mathcal{O}_Y) = H^1(Y, \mathcal{O}_Y)$
（《概形的上同调》，引理
\ref{coherent-lemma-relative-affine-cohomology}），引理得证。
\end{proof}

\begin{lemma}
\label{curves-lemma-genus-normalization}
设 $k$ 为域。设 $X$ 为维数 $\leq 1$ 的固有概形，其底域为 $k$。
若 $X' \to X$ 是双有理固有态射，则
$$
\dim_k H^1(X, \mathcal{O}_X) \geq \dim_k H^1(X', \mathcal{O}_{X'})
$$
若 $X$ 约化，映射
$H^0(X, \mathcal{O}_X) \to H^0(X', \mathcal{O}_{X'})$
满射，并且等号成立，则 $X' = X$。
\end{lemma}

\begin{proof}
若 $f : X' \to X$ 双有理且固有，则它是有限态射，并在某个稠密
开集上为同构（见《簇》，引理
\ref{varieties-lemma-finite-in-codim-1} 与
\ref{varieties-lemma-modification-normal-iso-over-codimension-1}）。
因此不等式来自引理 \ref{curves-lemma-genus-goes-down}。假设 $X$ 约化。
则 $\mathcal{O}_X \to f_*\mathcal{O}_{X'}$ 为单射，并得到短正合列
$$
0 \to \mathcal{O}_X \to f_*\mathcal{O}_{X'} \to \mathcal{F} \to 0
$$
在第二个断言的假设下，由上同调长正合列可得
$H^0(X, \mathcal{F}) = 0$。于是 $\mathcal{F} = 0$，因为
$\mathcal{F}$ 由整体截面生成（《簇》，引理
\ref{varieties-lemma-chi-tensor-finite}）。并且
$\mathcal{O}_X = f_*\mathcal{O}_{X'}$。由于 $f$ 仿射，
这蕴含 $X = X'$。
\end{proof}

\begin{lemma}
\label{curves-lemma-bound-geometric-genus-curve}
设 $k$ 为域。设 $C$ 为 $k$ 上的固有曲线。置
$\kappa = H^0(C, \mathcal{O}_C)$。则
$$
[\kappa : k]_s \dim_\kappa H^1(C, \mathcal{O}_C) \geq g_{geom}(C/k)
$$
\end{lemma}

\begin{proof}
《簇》，引理 \ref{varieties-lemma-regular-functions-proper-variety}
说明 $\kappa$ 是域，并且是 $k$ 的有限扩张。由《域》，引理
\ref{fields-lemma-separable-degree}，有
$[\kappa : k]_s = |\Mor_k(\kappa, \overline{k})|$；从而
$\Spec(\kappa \otimes_k \overline{k})$ 有
$[\kappa : k]_s$ 个点，每个点的剩余域都是 $\overline{k}$。因此
$$
C_{\overline{k}} =
\bigcup\nolimits_{t \in \Spec(\kappa \otimes_k \overline{k})} C_t
$$
（集合论意义下的并）。这里
$C_t = C \times_{\Spec(\kappa), t} \Spec(\overline{k})$，而
$t$ 被视为 $k$-代数映射 $t : \kappa \to \overline{k}$。
由此得到 $g_{geom}(C/k) = \sum_t g_{geom}(C_t/\overline{k})$，
且求和指标集的大小为 $[\kappa : k]_s$。有
$$
H^0(C_t, \mathcal{O}_{C_t}) = \overline{k}
\quad\text{且}\quad
\dim_{\overline{k}} H^1(C_t, \mathcal{O}_{C_t}) =
\dim_\kappa H^1(C, \mathcal{O}_C)
$$
，这是上同调与基变换的结果（《概形的上同调》，引理
\ref{coherent-lemma-flat-base-change-cohomology}）。注意，正规化
$C_t^\nu$ 是 $C_t$ 各不可约分支之非奇异射影模型的不交并
（《态射》，引理
\ref{morphisms-lemma-normalization-in-terms-of-components}）。因此
$\dim_{\overline{k}} H^1(C_t^\nu, \mathcal{O}_{C_t^\nu})$
等于 $g_{geom}(C_t/\overline{k})$。由引理
\ref{curves-lemma-genus-goes-down} 有
$$
\dim_{\overline{k}} H^1(C_t, \mathcal{O}_{C_t}) \geq
\dim_{\overline{k}} H^1(C_t^\nu, \mathcal{O}_{C_t^\nu})
$$
，证毕。
\end{proof}

\begin{lemma}
\label{curves-lemma-bound-torsion-simple}
设 $k$ 为域。设 $X$ 为维数 $\leq 1$ 的固有概形，其底域为 $k$。
设 $\ell$ 为在 $k$ 中可逆的素数。则
$$
\dim_{\mathbf{F}_\ell} \Pic(X)[\ell] \leq
\dim_k H^1(X, \mathcal{O}_X) + g_{geom}(X/k)
$$
，其中 $g_{geom}(X/k)$ 如上定义。
\end{lemma}

\begin{proof}
由《簇》，引理 \ref{varieties-lemma-change-fields-pic}，映射
$\Pic(X) \to \Pic(X_{\overline{k}})$ 为单射。由《概形的上同调》，
引理 \ref{coherent-lemma-flat-base-change-cohomology}，
$\dim_k H^1(X, \mathcal{O}_X)$ 等于
$\dim_{\overline{k}} H^1(X_{\overline{k}}, \mathcal{O}_{X_{\overline{k}}})$。
因此可假设 $k$ 代数闭。

\medskip\noindent
设 $X_{red}$ 为 $X$ 的约化。则满射
$\mathcal{O}_X \to \mathcal{O}_{X_{red}}$ 诱导满射
$H^1(X, \mathcal{O}_X) \to H^1(X, \mathcal{O}_{X_{red}})$
，因为由《上同调》，命题
\ref{cohomology-proposition-vanishing-Noetherian}，
拟凝聚层的上同调在次数 $\geq 2$ 时消失。由于
$X_{red} \to X$ 在 $\overline{k}$ 上的不可约分支之间诱导同构，
并在 Picard 群的 $\ell$-挠子群上诱导同构
（《曲线的 Picard 概形》，引理
\ref{pic-lemma-torsion-descends}），可以把 $X$ 替换为 $X_{red}$。
这样就归约到命题
\ref{curves-proposition-torsion-picard-reduced-proper}。
\end{proof}






\section{结点曲线}
\label{curves-section-nodal}

\noindent
我们已经定义了代数闭域上的普通二重点；见定义
\ref{curves-definition-multicross}。具体地，若 $x \in X$ 是某个
$1$ 维代数概形的闭点，而该概形定义在代数闭域 $k$ 上，
则 $x$ 为普通二重点当且仅当
$$
\mathcal{O}_{X, x}^\wedge \cong k[[x, y]]/(xy)
$$
见式 (\ref{curves-equation-multicross}) 及节
\ref{curves-section-multicross} 中其后的讨论。

\begin{definition}
\label{curves-definition-nodal}
设 $k$ 为域。设 $X$ 为 $1$ 维局部代数 $k$-概形。
\begin{enumerate}
\item 称闭点 $x \in X$ 是{\it 结点}或{\it 普通二重点}，
也称它{\it 定义了结点奇点}，如果存在普通二重点
$\overline{x} \in X_{\overline{k}}$ 映到 $x$。
\item 称{\it $X$ 的奇点至多为结点型}，如果 $X$ 的每个闭点
或者位于结构态射 $X \to \Spec(k)$ 的光滑轨迹中，
或者是普通二重点。
\end{enumerate}
\end{definition}

\noindent
若 $1$ 维代数概形 $X$ 的奇点至多为结点型，通常就称 $X$ 为
{\it 结点曲线}。有时还要求结点曲线固有。按此方式定义的结点曲线
未必不可约，因此这与我们先前把曲线定义为维数为 $1$ 的簇相冲突。

\begin{lemma}
\label{curves-lemma-reduced-quotient-regular-ring-dim-2}
设 $(A, \mathfrak m)$ 为维数 $2$ 的正则局部环，并设
$I \subset \mathfrak m$ 为理想。
\begin{enumerate}
\item 若 $A/I$ 约化，则 $I = (0)$、$I = \mathfrak m$，或者
$I = (f)$，其中 $f \in \mathfrak m$ 非零。
\item 若 $A/I$ 的深度为 $1$，则 $I = (f)$，其中非零
$f \in \mathfrak m$。
\end{enumerate}
\end{lemma}

\begin{proof}
假设 $I \not = 0$。写成 $I = (f_1, \ldots, f_r)$。由于 $A$ 是
唯一分解整环（《代数补篇》，引理
\ref{more-algebra-lemma-regular-local-UFD}），可写成
$f_i = fg_i$，其中 $f$ 是 $f_1, \ldots, f_r$ 的最大公因子。
所以 $g_1, \ldots, g_r$ 的最大公因子为 $1$，这意味着不存在
高度 $1$ 的素理想包含 $g_1, \ldots, g_r$。于是或者
$(g_1, \ldots, g_r) = A$，从而 $I = (f)$；或者不是如此，
此时 $\dim(A) = 2$ 蕴含
$V(g_1, \ldots, g_r) = \{\mathfrak m\}$, i.e.,
$\mathfrak m = \sqrt{(g_1, \ldots, g_r)}$.

\medskip\noindent
假设 $A/I$ 约化，即 $I$ 为根理想。若 $f$ 是单位，则由 $I$
为根理想可知 $I = \mathfrak m$。若 $f \in \mathfrak m$，则
$f^n$ 在 $A/I$ 中的像为零。由约化性得到 $f \in I$，故
$I = (f)$。

\medskip\noindent
假设 $A/I$ 的深度为 $1$。则 $\mathfrak m$ 不是 $A/I$ 的
伴随素理想。由于 $f$ 模 $I$ 的类被 $g_1, \ldots, g_r$ 零化，
可知 $f$ 的类在 $A/I$ 中为零。因此按要求有 $I = (f)$。
\end{proof}

\noindent
设 $\kappa$ 为域，$V$ 为 $\kappa$ 上的向量空间。称
$q \in \text{Sym}^2_\kappa(V)$ {\it 非退化}，如果诱导的
$\kappa$-线性映射 $V^\vee \to V$ 为同构。若
$q = \sum_{i \leq j} a_{ij} x_i x_j$ 对某个 $\kappa$-基
$x_1, \ldots, x_n$ 成立，而该基属于 $V$，则这意味着
下列矩阵的行列式
$$
\left(
\begin{matrix}
2a_{11} & a_{12} & \ldots \\
a_{12} & 2a_{22} & \ldots \\
\ldots & \ldots & \ldots
\end{matrix}
\right)
$$
非零。这等价于如下条件：$q$ 关于 $x_i$ 的各偏导数在概形论意义下
截出 $0$。

\begin{lemma}
\label{curves-lemma-nodal-algebraic}
设 $k$ 为域。设 $(A, \mathfrak m, \kappa)$ 为 Noether 局部
$k$-代数。下列条件等价：
\begin{enumerate}
\item $\kappa/k$ 可分，$A$ 约化，
$\dim_\kappa(\mathfrak m/\mathfrak m^2) = 2$，并且存在非退化元素
$q \in \text{Sym}^2_\kappa(\mathfrak m/\mathfrak m^2)$
，它在 $\mathfrak m^2/\mathfrak m^3$ 中映为零；
\item $\kappa/k$ 可分，$\text{depth}(A) = 1$，
$\dim_\kappa(\mathfrak m/\mathfrak m^2) = 2$，并且存在非退化元素
$q \in \text{Sym}^2_\kappa(\mathfrak m/\mathfrak m^2)$
，它在 $\mathfrak m^2/\mathfrak m^3$ 中映为零；
\item $\kappa/k$ 可分，并且
$A^\wedge \cong \kappa[[x, y]]/(ax^2 + bxy + cy^2)$
为 $k$-代数同构，其中 $ax^2 + bxy + cy^2$ 是 $\kappa$ 上
的非退化二次型。
\end{enumerate}
\end{lemma}

\begin{proof}
假设 (3) 成立。则 $A^\wedge$ 约化，因为
$ax^2 + bxy + cy^2$ 或者不可约，或者是两个不相伴素元的乘积。因此
$A \subset A^\wedge$ 约化，从而 (1) 成立。

\medskip\noindent
假设 (1) 成立。$A$ 不可能是 Artin 环，否则由
$\mathfrak m \not = (0)$ 可知它不约化。因此
$\dim(A) \geq 1$；再由《代数》，引理
\ref{algebra-lemma-criterion-reduced} 得到
$\text{depth}(A) \geq 1$。另一方面，$\dim(A) = 2$ 蕴含
$A$ 正则；由《代数》，引理
\ref{algebra-lemma-regular-graded}，这与 $q$ 的存在矛盾。
所以 $\dim(A) \leq 1$；再由《代数》，引理
\ref{algebra-lemma-bound-depth} 得到
$\text{depth}(A) = 1$。故 (2) 成立。

\medskip\noindent
假设 (2) 成立。$A$ 的深度与 $A^\wedge$ 的深度相同
（《代数补篇》，引理 \ref{more-algebra-lemma-completion-depth}），
而其余条件不受完备化影响，因此可假设 $A$ 完备。按《代数补篇》，
引理 \ref{more-algebra-lemma-lift-residue-field} 选取
$\kappa \to A$。由于
$\dim_\kappa(\mathfrak m/\mathfrak m^2) = 2$，可以选取
$x_0, y_0 \in \mathfrak m$，使其像构成一组基。得到连续
$\kappa$-代数映射
$$
\kappa[[x, y]] \longrightarrow A
$$
，其规则为 $x \mapsto x_0$ 与 $y \mapsto y_0$。设
$q$ 是 $ax_0^2 + bx_0y_0 + cy_0^2$ 在
$\text{Sym}^2_\kappa(\mathfrak m/\mathfrak m^2)$ 中的类。
把 $Q(x, y) = ax^2 + bxy + cy^2$ 视为二元多项式。于是
$$
Q(x_0, y_0) = ax_0^2 + bx_0y_0 + cy_0^2 =
\sum\nolimits_{i + j = 3} a_{ij} x_0^iy_0^j
$$
，其中某些 $a_{ij}$ 属于 $A$。我们要证明，可以通过改变
$x_0$、$y_0$ 的选取来提高消没阶。假设
$x_1, y_1 \in \mathfrak m^2$。则
$$
Q(x_0 + x_1, y_0 + y_1) = Q(x_0, y_0) +
(2ax_0 + by_0)x_1 + (bx_0 + 2cy_0)y_1 \bmod \mathfrak m^4
$$
由引理之前的讨论，$Q$ 非退化恰好意味着
$2ax_0 + by_0$ 与 $bx_0 + 2cy_0$ 是
一组 $\kappa$-基；它是 $\mathfrak m/\mathfrak m^2$ 的基。
因此当然可以选取
$x_1, y_1 \in \mathfrak m^2$，使得
$Q(x_0 + x_1, y_0 + y_1) \in \mathfrak m^4$.
如此归纳下去，可以找到 $x_i, y_i \in \mathfrak m^{i + 1}$，使得
$$
Q(x_0 + x_1 + \ldots + x_n, y_0 + y_1 + \ldots + y_n) \in \mathfrak m^{n + 3}
$$
由于 $A$ 完备，可以置 $x_\infty = \sum x_i$ 以及
$y_\infty = \sum y_i$，并考察映射
$\kappa[[x, y]] \longrightarrow A$，它把 $x$ 映到 $x_\infty$，
把 $y$ 映到 $y_\infty$。由《代数》，引理
\ref{algebra-lemma-completion-generalities}，该映射诱导满射
$\kappa[[x, y]]/(Q) \longrightarrow A$。由引理
\ref{curves-lemma-reduced-quotient-regular-ring-dim-2}，
$k[[x, y]] \to A$ 的核为主理想。但该核不能包含 $Q$ 的真因子，
因为这样的因子的次数将为 $1$，所用变量为 $x, y$；这与
$\dim(\mathfrak m/\mathfrak m^2) = 2$ 矛盾。因此按要求，$Q$
生成该核。
\end{proof}

\begin{lemma}
\label{curves-lemma-2-branches-delta-1}
设 $k$ 为域。设 $(A, \mathfrak m, \kappa)$ 为 Nagata 局部
$k$-代数。下列条件等价：
\begin{enumerate}
\item $k \to A$ 如引理 \ref{curves-lemma-nodal-algebraic} 所述；
\item $\kappa/k$ 可分，$A$ 约化且维数为 $1$，
$\delta$-不变量在 $A$ 上为 $1$，并且 $A$ 有 $2$ 个几何分支。
\end{enumerate}
若这些条件成立，则整闭包 $A'$（即 $A$ 在其全分式环中的整闭包）
有 $1$ 个或 $2$ 个极大理想 $\mathfrak m'$，而扩张
$\kappa(\mathfrak m')/k$ 可分。
\end{lemma}

\begin{proof}
在两种条件下，$A$ 与 $A^\wedge$ 都约化。在情形 (2) 中，
这是因为约化 Nagata 局部环的完备化约化（《代数补篇》，引理
\ref{more-algebra-lemma-completion-reduced}）。在两种条件下，
$A$ 与 $A^\wedge$ 的维数都是 $1$（《代数补篇》，引理
\ref{more-algebra-lemma-completion-dimension}）。由《簇》，引理
\ref{varieties-lemma-delta-same-after-completion} 以及《代数补篇》，
引理 \ref{more-algebra-lemma-one-dimensional-number-of-branches}，
$\delta$-不变量以及 $A$ 与 $A^\wedge$ 的几何分支数都相同。
按《簇》，引理 \ref{varieties-lemma-pre-delta-invariant}，设
$A'$ 为 $A$ 在其全分式环中的整闭包。由《簇》，引理
\ref{varieties-lemma-normalization-same-after-completion}，
$A' \otimes_A A^\wedge$ 对 $A^\wedge$ 扮演同一角色。因此可以把
$A$ 替换为 $A^\wedge$，并假设 $A$ 完备。

\medskip\noindent
假设 (1) 成立。只需证明 $A$ 有两个几何分支且
$\delta$-不变量为 $1$。可以假设
$A = \kappa[[x, y]]/(ax^2 + bxy + cy^2)$，其中
$q = ax^2 + bxy + cy^2$ 非退化。分两种情形。

\medskip\noindent
情形 I：$q$ 在 $\kappa$ 上分裂。改变坐标后，可假设
$q = xy$。于是
$$
A' = \kappa[[x, y]]/(x) \times \kappa[[x, y]]/(y)
$$

\medskip\noindent
情形 II：$q$ 不分裂。此时 $c \not = 0$，而非退化意味着
$b^2 - 4ac \not = 0$。因此
$\kappa' = \kappa[t]/(a + bt + ct^2)$ 是次数为 $2$ 的可分扩张，
其底域为 $\kappa$。于是 $t = y/x$ 在 $A$ 上整，并得到
$$
A' = \kappa'[[x]]
$$
，其中 $y$ 在右端映到 $tx$。

\medskip\noindent
在两种情形下，直接验证可知 $\delta$-不变量为 $1$，
几何分支数为 $2$。由此可见 (1) 蕴含 (2)。
此外也得到引理最后一个断言。

\medskip\noindent
假设 (2) 成立。由《代数补篇》，引理
\ref{more-algebra-lemma-number-of-branches-1}，$A'$ 或者有两个
极大理想，或者 $A'$ 有一个极大理想，且
$[\kappa(\mathfrak m') : \kappa]_s = 2$。

\medskip\noindent
情形 I：$A'$ 有两个极大理想 $\mathfrak m'_1$、$\mathfrak m'_2$，
其剩余域分别为 $\kappa_1$、$\kappa_2$。由于 $\delta$-不变量
等于 $A'/A$ 的长度，并且存在满射
$A'/A \to (\kappa_1 \times \kappa_2)/\kappa$，可知
$\kappa = \kappa_1 = \kappa_2$。由于 $A$ 完备（并且由《代数》，
引理 \ref{algebra-lemma-complete-henselian} 可知它 Hensel），
而 $A'$ 在 $A$ 上有限，可知 $A' = A_1 \times A_2$
（由《代数》，引理 \ref{algebra-lemma-finite-over-henselian}）。
由于 $A'$ 为正规环，$A_1$ 与 $A_2$ 都是离散赋值环。因此由
《代数补篇》，引理
\ref{more-algebra-lemma-power-series-over-residue-field}，
$A_1$ 与 $A_2$ 都同构于 $\kappa[[t]]$（作为 $k$-代数）。
由于 $\delta$-不变量为 $1$，可知 $A$ 是 $A_1$ 与 $A_2$ 的楔和
（《簇》，定义 \ref{varieties-definition-wedge}）。容易得到
$A \cong \kappa[[x, y]]/(xy)$。

\medskip\noindent
情形 II：$A'$ 只有一个极大理想 $\mathfrak m'$，其剩余域为
$\kappa'$，且 $[\kappa' : \kappa]_s = 2$。完全照情形 I 的论证，
可知 $[\kappa' : \kappa] = 2$，并且 $\kappa'$ 在 $\kappa$ 上
可分。由于 $A'$ 正规，可知 $A'$ 同构于 $\kappa'[[t]]$
（见上面的引文）。由于 $A'/A$ 的长度为 $1$，得到
$$
A = \{f \in \kappa'[[t]] \mid f(0) \in \kappa\}
$$
于是简单计算表明 $A$ 如情形 (1) 所述。
\end{proof}

\begin{lemma}
\label{curves-lemma-fitting-ideal-well-defined}
设 $k$ 为域。设 $A = k[[x_1, \ldots, x_n]]$，并设
$I = (f_1, \ldots, f_m) \subset A$ 为理想。对任意 $r \geq 0$，
$A/I$ 中由 $r \times r$ 子式生成的理想（这些子式来自矩阵
$(\partial f_j/\partial x_i)$），与 $I$ 的生成元以及正则参数系
$x_1, \ldots, x_n$ 的选取无关；该参数系属于 $A$。
\end{lemma}

\begin{proof}
此引理的“正确”证明是说明该理想为一个连续微分模的
第 $(n - r)$ 个 Fitting 理想；该模属于 $A/I$ 相对于 $k$。
下面给出直接证明。若 $g_1, \ldots g_l$ 是 $I$ 的另一组生成元，
则可写成 $g_s = \sum a_{sj}f_j$，并有矩阵等式
$$
(\partial g_s/\partial x_i) = (a_{sj}) (\partial f_j/\partial x_i)
+ (\partial a_{sj}/\partial x_i f_j)
$$
最后一项在 $A/I$ 中为零。由 Cauchy--Binet 公式，$g_s$
所对应的子式理想包含于 $f_j$ 所对应的理想。由对称性，
两理想相同。若 $y_1, \ldots, y_n \in \mathfrak m_A$
是另一组正则参数系，则矩阵
$(\partial y_j/\partial x_i)$
可逆，此时可使用微分的链式法则。略去若干细节。
\end{proof}

\begin{lemma}
\label{curves-lemma-fitting-ideal}
设 $k$ 为域。设 $A = k[[x_1, \ldots, x_n]]$，并设
$I = (f_1, \ldots, f_m) \subset \mathfrak m_A$ 为理想。下列条件等价：
\begin{enumerate}
\item $k \to A/I$ 如引理 \ref{curves-lemma-nodal-algebraic} 所述；
\item $A/I$ 约化，并且各 $(n - 1) \times (n - 1)$ 子式
（来自矩阵 $(\partial f_j/\partial x_i)$）生成
$I + \mathfrak m_A$；
\item $\text{depth}(A/I) = 1$，并且各
$(n - 1) \times (n - 1)$ 子式（来自矩阵
$(\partial f_j/\partial x_i)$）生成 $I + \mathfrak m_A$。
\end{enumerate}
\end{lemma}

\begin{proof}
由引理 \ref{curves-lemma-fitting-ideal-well-defined}，证明中可以改变坐标系
以及生成元的选取。

\medskip\noindent
若 (1) 成立，则可改变坐标，使得 $x_1, \ldots, x_{n - 2}$
在 $A/I$ 中映为零，并且
$A/I = k[[x_{n - 1}, x_n]]/(a x_{n - 1}^2 + b x_{n - 1}x_n + c x_n^2)$
，其中 $a x_{n - 1}^2 + b x_{n - 1}x_n + c x_n^2$ 是某个
非退化二次型。直接计算即可证明 (2) 与 (3) 均成立。

\medskip\noindent
假设矩阵的各 $(n - 1) \times (n - 1)$ 子式（该矩阵为
$(\partial f_j/\partial x_i)$）生成 $I + \mathfrak m_A$。
假设对某个 $i$ 与 $j$，偏导数
$\partial f_j/\partial x_i$ 是 $A$ 中的单位。此时可使用参数系
$f_j, x_1, \ldots, x_{i - 1}, \hat x_i, x_{i + 1}, \ldots, x_n$
以及生成元
$f_j, f_1, \ldots, f_{j - 1}, \hat f_j, f_{j + 1}, \ldots, f_m$
来生成 $I$。于是得到正则参数系 $x_1, \ldots, x_n$，以及生成元
$x_1, f_2, \ldots, f_m$；它们生成 $I$。接着寻找满足
$i \geq 2$ 与 $j \geq 2$ 的指标，使
$\partial f_j/\partial x_i$ 是 $A$ 中的单位。若存在这样的指标对，
则可作上述替换，并假设有正则参数系 $x_1, \ldots, x_n$，
以及生成元 $x_1, x_2, f_3, \ldots, f_m$；它们生成 $I$。
继续进行，有限步后达到如下情形：存在正则参数系
$x_1, \ldots, x_n$，以及生成元
$x_1, \ldots, x_t, f_{t + 1}, \ldots, f_m$；它们生成 $I$，
并且 $\partial f_j/\partial x_i \in \mathfrak m_A$ 对所有
$i, j \geq t + 1$ 都成立。

\medskip\noindent
此时偏导数组成的矩阵具有如下分块形状：
$$
\left(
\begin{matrix}
I_{t \times t} & * \\
0 & \mathfrak m_A
\end{matrix}
\right)
$$
因此每个 $(n - 1) \times (n - 1)$ 子式都属于
$\mathfrak m_A^{n - 1 - t}$。注意，$I \not = \mathfrak m_A$，
否则子式理想将包含 $1$。由此 $n - 1 - t \leq 1$，因为存在
$\mathfrak m_A \setminus \mathfrak m_A^2 + I$ 中的元素
（否则由 Nakayama 有 $I = \mathfrak m_A$）。因此
$t \geq n - 2$。上面已经看到
$t \not = n$；类似地，若 $t = n - 1$，则存在可逆的
$(n - 1) \times (n - 1)$ 子式，这同样不允许。因此
$t = n - 2$。于是 $A/I$ 是 $k[[x_{n - 1}, x_n]]$ 的商。
引理 \ref{curves-lemma-reduced-quotient-regular-ring-dim-2} 说明，
在情形 (2) 与 (3) 中，$I$ 都由
$x_1, \ldots, x_{n - 2}, f$ 生成，其中
$f = f(x_{n - 1}, x_n)$。此时关于子式的条件恰好说明
$f$ 的二次项非退化，即 $A/I$ 如引理
\ref{curves-lemma-nodal-algebraic} 所述。
\end{proof}

\begin{lemma}
\label{curves-lemma-nodal}
设 $k$ 为域。设 $X$ 为 $1$ 维代数 $k$-概形，并设
$x \in X$ 为闭点。下列条件等价：
\begin{enumerate}
\item $x$ 是结点；
\item $k \to \mathcal{O}_{X, x}$ 如引理
\ref{curves-lemma-nodal-algebraic} 所述；
\item 任意 $\overline{x} \in X_{\overline{k}}$，若它映到 $x$，
就定义结点奇点；
\item $\kappa(x)/k$ 可分，$\mathcal{O}_{X, x}$ 约化，并且
$\Omega_{X/k}$ 的第一 Fitting 理想生成 $\mathfrak m_x$；
这一生成发生在 $\mathcal{O}_{X, x}$ 中；
\item $\kappa(x)/k$ 可分，$\text{depth}(\mathcal{O}_{X, x}) = 1$，
并且 $\Omega_{X/k}$ 的第一 Fitting 理想生成 $\mathfrak m_x$；
这一生成发生在 $\mathcal{O}_{X, x}$ 中；
\item $\kappa(x)/k$ 可分，并且 $\mathcal{O}_{X, x}$ 约化，
其 $\delta$-不变量为 $1$，且有 $2$ 个几何分支。
\end{enumerate}
\end{lemma}

\begin{proof}
先假设 $k$ 代数闭。此时 (1) 与 (3) 的等价性显然。
(1)、(3) 与 (2) 等价，因为二元非退化二次型在坐标变换下
只有 $xy$ 一种。(1) 与 (6) 的等价性就是引理
\ref{curves-lemma-multicross-algebra}。把 $X$ 替换为 $x$ 的一个仿射邻域后，
可假设存在闭浸入 $X \to \mathbf{A}^n_k$，它把 $x$ 映到 $0$。
设 $f_1, \ldots, f_m \in k[x_1, \ldots, x_n]$ 是理想 $I$
的生成元；该理想定义 $X$ 在 $\mathbf{A}^n_k$ 中的嵌入。此时
$\Omega_{X/k}$ 对应于 $R = k[x_1, \ldots, x_n]/I$-模
$\Omega_{R/k}$，后者有展示
$$
R^{\oplus m} \xrightarrow{(\partial f_j/\partial x_i)}
R^{\oplus n} \to \Omega_{R/k} \to 0
$$
（见《代数》，节 \ref{algebra-section-differentials} 与
\ref{algebra-section-netherlander}.)
因此 $\Omega_{R/k}$ 的第一 Fitting 理想就是由
$(n - 1) \times (n - 1)$ 子式生成的理想；这些子式来自矩阵
$(\partial f_j/\partial x_i)$。把引理
\ref{curves-lemma-fitting-ideal} 用于 $k[x_1, \ldots, x_n] \to R$
在极大理想 $(x_1, \ldots, x_n)$ 处的完备化，即知
(2)、(4)、(5) 等价。

\medskip\noindent
现在假设 $k$ 为任意域。在情形 (2)、(4)、(5)、(6) 中，
剩余域 $\kappa(x)$ 在 $k$ 上可分。我们来证明情形 (1)、(3)
也具有这一性质。设 $Z \subset X$ 为 $X$ 的闭子概形，
它由 $\Omega_{X/k}$ 的第一 Fitting 理想定义。$Z$ 的形成
与域扩张交换（《除子》，引理
\ref{divisors-lemma-base-change-and-fitting-ideal-omega}）。
若 (1) 或 (3) 成立，则存在一点 $\overline{x}$；它属于
$X_{\overline{k}}$，并且 $\overline{x}$ 是重数为 $1$ 的孤立点，
所在概形为 $Z_{\overline{k}}$（因为在 $\overline{k}$ 上，
引理的各条件等价）。
特别地，$Z_{\overline{x}}$ 在 $\overline{x}$ 处几何约化
（因为 $\overline{k}$ 代数闭）。所以 $Z$ 在 $x$ 处几何约化
（《簇》，引理
\ref{varieties-lemma-geometrically-reduced-upstairs}）。特别地，
$Z$ 在 $x$ 处约化，因而 $Z = \Spec(\kappa(x))$ 在 $x$ 的一个
邻域中成立，并且 $\kappa(x)$ 在 $k$ 上几何约化。
这意味着 $\kappa(x)/k$ 可分（《代数》，引理
\ref{algebra-lemma-characterize-separable-field-extensions}）。

\medskip\noindent
上一段的论证表明，若 (1) 或 (3) 成立，则
$\Omega_{X/k}$ 的第一 Fitting 理想生成 $\mathfrak m_x$。
由于 $\mathcal{O}_{X, x} \to \mathcal{O}_{X_{\overline{k}}, \overline{x}}$
平坦，并且
$\mathcal{O}_{X_{\overline{k}}, \overline{x}}$ 约化且深度为 $1$，
可知 (4) 与 (5) 成立（使用《代数》，引理
\ref{algebra-lemma-descent-reduced} 与
\ref{algebra-lemma-apply-grothendieck}).
反过来，由《代数》，引理 \ref{algebra-lemma-criterion-reduced}，
(4) 蕴含 (5)。若 (5) 成立，则 $Z$ 在 $x$ 处几何约化
（因为 $\kappa(x)/k$ 可分，并且在一个邻域中 $Z$ 就是 $x$）。
所以 $Z_{\overline{k}}$ 在任意点 $\overline{x}$ 处约化；
该点属于 $X_{\overline{k}}$ 且位于 $x$ 上方。换言之，
$\Omega_{X_{\overline{k}}/\overline{k}}$ 的第一 Fitting 理想生成
$\mathfrak m_{\overline{x}}$；这一生成发生在
$\mathcal{O}_{X_{\overline{k}, \overline{x}}}$ 中。此外，由于
$\mathcal{O}_{X, x} \to \mathcal{O}_{X_{\overline{k}}, \overline{x}}$
平坦，可知
$\text{depth}(\mathcal{O}_{X_{\overline{k}}, \overline{x}}) = 1$
（见上面的引文）。因此 (5) 对
$\overline{x} \in X_{\overline{k}}$ 成立，从而 (3) 成立
（使用代数闭域上各条件的等价性）。由此可见
(1)、(3)、(4)、(5) 等价。

\medskip\noindent
(2) 与 (6) 的等价性来自引理 \ref{curves-lemma-2-branches-delta-1}。

\medskip\noindent
最后证明 (2) $=$ (6) 与 (1) $=$ (3) $=$ (4) $=$ (5) 等价。
首先，由《簇》，引理
\ref{varieties-lemma-geometric-branches-and-change-of-fields}，
$X$ 在 $x$ 处的几何分支数等于
$X_{\overline{k}}$ 在 $\overline{x}$ 处的几何分支数。
根据目前已有的信息可知，在所有情形中这个数都等于 $2$。
另一方面，在情形 (1) 中，$X$ 显然在 $x$ 处几何约化，因而
$$
\delta\text{-invariant of }X\text{ 在 }x \leq
\delta\text{-invariant of }X_{\overline{k}}\text{ 在 }\overline{x}
$$
，这来自《簇》，引理
\ref{varieties-lemma-delta-invariant-and-change-of-fields}。
在情形 (1) 中，右端为 $1$，这迫使 $\delta$-不变量
（对 $X$ 在 $x$ 处而言）为 $1$（因为若其为零，则由《簇》，引理
\ref{varieties-lemma-delta-invariant-is-zero}，
$\mathcal{O}_{X, x}$ 将是离散赋值环，因而单分支，矛盾）。
所以 (5) 成立。反过来，若 (2) $=$ (5) 成立，则《簇》，引理
\ref{varieties-lemma-geometrically-normal-in-codim-1} 的假设
(a)、(b)、(c) 由引理 \ref{curves-lemma-2-branches-delta-1} 可知
对 $x \in X$ 成立。因此《簇》，引理
\ref{varieties-lemma-delta-invariant-and-change-of-fields-better}
适用，并说明上述不等式中等号成立。由此可知 (5) 对
$\overline{x} \in X_{\overline{k}}$ 成立；再由代数闭域上各条件的
等价性，回到情形 (1)。
\end{proof}

\begin{remark}[与结点相伴的二次扩张]
\label{curves-remark-quadratic-extension}
设 $k$ 为域。设 $(A, \mathfrak m, \kappa)$ 为 Noether 局部
$k$-代数。假设 $(A, \mathfrak m, \kappa)$ 如引理
\ref{curves-lemma-nodal-algebraic} 所述，或者 $A$ 为引理
\ref{curves-lemma-2-branches-delta-1} 所述的 Nagata 环，或者 $A$ 完备且
如引理 \ref{curves-lemma-fitting-ideal} 所述。则 $A$ 按如下方式典范地
定义次数为 $2$ 的可分 $\kappa$-代数 $\kappa'$：
\begin{enumerate}
\item 设 $q = ax^2 + bxy + cy^2$ 为引理
\ref{curves-lemma-nodal-algebraic} 所述的非退化二次型，并选取坐标
$x, y$ 使 $a \not = 0$，再置
$\kappa' = \kappa[x]/(ax^2 + bx + c)$,
\item 设 $A' \supset A$ 为 $A$ 在其全分式环中的整闭包，并置
$\kappa' = A'/\mathfrak m A'$；或者
\item 令 $\kappa'$ 为满足下式的 $\kappa$-代数：
$\text{Proj}(\bigoplus_{n \geq 0} \mathfrak m^n/\mathfrak m^{n + 1}) =
\Spec(\kappa')$.
\end{enumerate}
引理 \ref{curves-lemma-2-branches-delta-1} 的证明已经给出 (1) 与 (2)
的等价性。略去它们与 (3) 的等价性。若 $X$ 是局部 Noether
$k$-概形，而 $x \in X$ 是满足 $\mathcal{O}_{X, x} = A$ 的点，
则 (3) 表明
$\Spec(\kappa') = X^\nu \times_X \Spec(\kappa)$，其中
$\nu : X^\nu \to X$ 是正规化态射。
\end{remark}

\begin{remark}[平凡二次扩张]
\label{curves-remark-trivial-quadratic-extension}
设 $k$ 为域。设 $(A, \mathfrak m, \kappa)$ 如评注
\ref{curves-remark-quadratic-extension} 所述，并设 $\kappa'/\kappa$
为相伴的次数为 $2$ 的可分代数。下列条件等价：
\begin{enumerate}
\item $\kappa' \cong \kappa \times \kappa$ 为 $\kappa$-代数同构；
\item 引理 \ref{curves-lemma-nodal-algebraic} 中的形式 $q$
可选为 $xy$；
\item $A$ 有两个分支；
\item 引理 \ref{curves-lemma-2-branches-delta-1} 中的扩张 $A'/A$
有两个极大理想；并且
\item $A^\wedge \cong \kappa[[x, y]]/(xy)$ 为 $k$-代数同构。
\end{enumerate}
这些条件的等价性已在引理
\ref{curves-lemma-2-branches-delta-1} 的证明中给出。若 $X$ 是局部
Noether $k$-概形，而 $x \in X$ 满足 $\mathcal{O}_{X, x} = A$，
则这恰好意味着有两个点 $x_1, x_2$；它们属于正规化 $X^\nu$，
位于 $x$ 上方，并且 $\kappa(x) = \kappa(x_1) = \kappa(x_2)$。
\end{remark}

\begin{definition}
\label{curves-definition-split-node}
设 $k$ 为域。设 $X$ 为 $1$ 维代数 $k$-概形，并设
$x \in X$ 为闭点。称 $x$ 为{\it 分裂结点}，如果 $x$ 是结点，
$\kappa(x) = k$，并且评注
\ref{curves-remark-trivial-quadratic-extension} 中的等价断言对
$A = \mathcal{O}_{X, x}$ 成立。
\end{definition}

\noindent
把关于这一概念已经知道的内容表述为如下必备引理。

\begin{lemma}
\label{curves-lemma-split-node}
设 $k$ 为域。设 $X$ 为 $1$ 维代数 $k$-概形，并设
$x \in X$ 为闭点。下列条件等价：
\begin{enumerate}
\item $x$ 是分裂结点；
\item $x$ 是结点，并且恰有两个点 $x_1, x_2$；它们属于正规化
$X^\nu$，位于 $x$ 上方，且
$k = \kappa(x_1) = \kappa(x_2)$,
\item $\mathcal{O}_{X, x}^\wedge \cong k[[x, y]]/(xy)$ 为 $k$-代数同构；
并且
\item 此处增补更多内容。
\end{enumerate}
\end{lemma}

\begin{proof}
这来自评注 \ref{curves-remark-trivial-quadratic-extension} 中的讨论以及
引理 \ref{curves-lemma-nodal}。
\end{proof}

\begin{lemma}
\label{curves-lemma-node-field-extension}
设 $K/k$ 为域扩张。设 $X$ 为局部代数
$k$-概形，其维数为 $1$。设 $y \in X_K$ 为一个点，其像为
$x \in X$。以下条件等价：
\begin{enumerate}
\item $x$ 是 $X$ 的闭点且为结点；
\item $y$ 是 $Y$ 的闭点且为结点。
\end{enumerate}
\end{lemma}

\begin{proof}
若 $x$ 是 $X$ 的闭点，则 $y$ 也是闭点（考察剩余域即可）。
但逆命题未必成立；确实，$Y$ 的一个闭点可能映到 $X$ 的非闭点。
不过在这种情形下，$y$ 不可能是结点。事实上，此时 $X$ 在 $x$
处几何单分支（因为 $x$ 将是 $X$ 的泛点，且
$\mathcal{O}_{X, x}$ 为 Artin 环，而任意 Artin 局部环都几何单分支），
所以 $Y$ 在 $y$ 处几何单分支（《簇》，引理
\ref{varieties-lemma-geometrically-unibranch-and-change-of-fields}）。
由引理 \ref{curves-lemma-nodal}，这意味着 $y$ 不可能是结点。
因此我们可以并且确实假设 $x$ 与 $y$ 都是闭点。

\medskip\noindent
选取代数闭包 $\overline{k}$、$\overline{K}$，以及映射
$\overline{k} \to \overline{K}$，它延拓给定映射 $k \to K$。
利用引理 \ref{curves-lemma-nodal} 中 (1) 与 (3) 的等价性，
可归约到 $k$ 与 $K$ 都是代数闭域的情形。在这种情形下，
可以像引理 \ref{curves-lemma-nodal} 的证明那样论证：几何分支数和
$\delta$-不变量，对 $X$ 在 $x$ 处与对 $Y$ 在 $y$ 处相同。
另一种论证是选取一个同构
$k[[x_1, \ldots, x_n]]/(g_1, \ldots, g_m) \to \mathcal{O}_{X, x}^\wedge$
；按照《簇》，引理
\ref{varieties-lemma-complete-local-ring-structure-as-algebra}，
这是一个 $k$-代数同构。由《簇》，引理
\ref{varieties-lemma-base-change-complete-local-ring}，它给出
$K[[x_1, \ldots, x_n]]/(g_1, \ldots, g_m) \to \mathcal{O}_{Y, y}^\wedge$
这一 $K$-代数同构。根据定义，我们必须证明
$$
k[[x_1, \ldots, x_n]]/(g_1, \ldots, g_m) \cong k[[s, t]]/(st)
$$
当且仅当
$$
K[[x_1, \ldots, x_n]]/(g_1, \ldots, g_m) \cong K[[s, t]]/(st)
$$
。我们鼓励读者自行证明这一点。由于 $k$ 与 $K$ 都是代数闭域，
这等价于要求这些环满足引理 \ref{curves-lemma-nodal-algebraic} 中的条件。
借助引理 \ref{curves-lemma-fitting-ideal}，它转化为关于
$(n - 1) \times (n - 1)$ 阶子式的陈述；这些子式来自矩阵
$(\partial g_j/\partial x_i)$，且该陈述显然与所用的域无关。
细节从略。
\end{proof}

\begin{lemma}
\label{curves-lemma-node-etale-local}
设 $k$ 为域。设 $X$ 为局部代数
$k$-概形，其维数为 $1$。设 $Y \to X$ 为 \'etale 态射，且
$y \in Y$ 为一个像等于 $x \in X$ 的点。以下条件等价：
\begin{enumerate}
\item $x$ 是 $X$ 的闭点且为结点；
\item $y$ 是 $Y$ 的闭点且为结点。
\end{enumerate}
\end{lemma}

\begin{proof}
由引理 \ref{curves-lemma-node-field-extension}，我们可以基变换到 $k$
的代数闭包。此时 $x$ 与 $y$ 的剩余域都是 $k$。
所以映射 $\mathcal{O}_{X, x}^\wedge \to \mathcal{O}_{Y, y}^\wedge$
是同构（例如见《\'Etale 态射》，引理
\ref{etale-lemma-characterize-etale-completions}，或《代数进阶》，引理
\ref{more-algebra-lemma-flat-unramified}）。故结论显然。
\end{proof}

\begin{lemma}
\label{curves-lemma-node-over-separable-extension}
设 $k'/k$ 为有限可分域扩张。设 $X$ 为局部代数
$k'$-概形，其维数为 $1$。并设 $x \in X$ 为闭点。以下条件等价：
\begin{enumerate}
\item $x$ 是结点；
\item $x$ 是结点，其中将 $X$ 视为局部代数 $k$-概形。
\end{enumerate}
\end{lemma}

\begin{proof}
这立即来自引理 \ref{curves-lemma-nodal} 对结点的刻画。
\end{proof}

\begin{lemma}
\label{curves-lemma-nodal-lci}
设 $k$ 为域。设 $X$ 为局部代数 $k$-概形，等维且维数为 $1$。
以下条件等价：
\begin{enumerate}
\item $X$ 的奇点至多为结点；
\item $X$ 在 $k$ 上为局部完全交，并且闭子概形
$Z \subset X$ 由 $\Omega_{X/k}$ 的第一 Fitting 理想截出，
且在 $k$ 上非分歧。
\end{enumerate}
\end{lemma}

\begin{proof}
我们强烈建议读者自行寻找本引理的证明；下面只是把先前的结果
拼合起来，反而可能掩盖真正发生的事情。

\medskip\noindent
假设 (2)。由于 $Z \to \Spec(k)$ 是拟有限的
（《态射》，引理 \ref{morphisms-lemma-unramified-quasi-finite}），
由《态射》，引理 \ref{morphisms-lemma-residue-field-quasi-finite}，
点 $x \in Z$ 的剩余域是 $k$ 的有限扩张（并且可分）。
因此由《态射》，引理
\ref{morphisms-lemma-algebraic-residue-field-extension-closed-point-fibre}，
每个 $x \in Z$ 都是 $X$ 的闭点。由《代数》，引理
\ref{algebra-lemma-lci-CM}，局部环 $\mathcal{O}_{X, x}$ 是
Cohen--Macaulay 环。根据维数理论，
$\dim(\mathcal{O}_{X, x}) = 1$（《簇》，节
\ref{varieties-section-algebraic-schemes}），故
$\text{depth}(\mathcal{O}_{X, x})) = 1$。由引理
\ref{curves-lemma-nodal}，$x$ 是结点。若 $x \in X$ 且 $x \not \in Z$，
则由《除子》，引理
\ref{divisors-lemma-d-fitting-ideal-omega-smooth}，
$X \to \Spec(k)$ 在 $x$ 处光滑。

\medskip\noindent
假设 (1)。在此假设下，$X$ 在每个闭点处几何约化
（见《簇》，引理
\ref{varieties-lemma-geometrically-reduced-upstairs}）。
因此由《簇》，引理
\ref{varieties-lemma-geometrically-reduced-dense-smooth-open}，
$X \to \Spec(k)$ 在一个稠密开集上光滑。于是 $Z$ 是闭的，
并且仅由闭点组成。由《除子》，引理
\ref{divisors-lemma-d-fitting-ideal-omega-smooth}，态射
$X \setminus Z \to \Spec(k)$ 光滑。因此由《态射》，引理
\ref{morphisms-lemma-smooth-syntomic} 以及《态射》，定义
\ref{morphisms-definition-syntomic} 中对局部完全交的定义，
$X \setminus Z$ 是局部完全交。由引理 \ref{curves-lemma-nodal}，
对每个点 $x \in Z$，局部环 $\mathcal{O}_{Z, x}$ 等于
$\kappa(x)$，且 $\kappa(x)$ 在 $k$ 上可分。所以
$Z \to \Spec(k)$ 非分歧（《态射》，引理
\ref{morphisms-lemma-unramified-over-field}）。最后，引理
\ref{curves-lemma-nodal} 结合引理 \ref{curves-lemma-nodal-algebraic} 的 (3)
说明 $\mathcal{O}_{X, x}$ 是《除幂代数》，定义
\ref{dpa-definition-lci} 意义下的完全交。不过，《除幂代数》，引理
\ref{dpa-lemma-check-lci-agrees} 与《态射》，引理
\ref{morphisms-lemma-local-complete-intersection} 说明，这与定义
域上局部完全交概形时所用的概念一致。证明完成。
\end{proof}

\begin{lemma}
\label{curves-lemma-facts-about-nodal-curves}
设 $k$ 为域。设 $X$ 为局部代数 $k$-概形，等维且维数为
$1$，其奇点至多为结点。则 $X$ 为 Gorenstein 且几何约化。
\end{lemma}

\begin{proof}
Gorenstein 断言来自引理 \ref{curves-lemma-nodal-lci} 和
《概形的对偶性》，引理 \ref{duality-lemma-gorenstein-lci}。
也可以使用如下事实：只需在传到代数闭包后检验
（《概形的对偶性》，引理
\ref{duality-lemma-gorenstein-base-change}）；再利用 Noether
局部环为 Gorenstein 当且仅当其完备化为 Gorenstein
（《对偶复形》，引理
\ref{dualizing-lemma-flat-under-gorenstein}）；最后直接证明局部环
$k[[t]]$ 与 $k[[x, y]]/(xy)$ 为 Gorenstein。

\medskip\noindent
为证明 $X$ 几何约化，只需证明 $X_{\overline{k}}$ 约化
（《簇》，引理 \ref{varieties-lemma-perfect-reduced} 和
\ref{varieties-lemma-geometrically-reduced}）。但
$X_{\overline{k}}$ 是代数闭域上的结点曲线。因此
$X_{\overline{k}}$ 的完备局部环同构于
$\overline{k}[[t]]$ 或 $\overline{k}[[x, y]]/(xy)$；
这些环都约化，正合所需。
\end{proof}

\begin{lemma}
\label{curves-lemma-closed-subscheme-nodal-curve}
设 $k$ 为域。设 $X$ 为局部代数 $k$-概形，等维且维数为
$1$，其奇点至多为结点。若 $Y \subset X$ 是等维且维数为
$1$ 的约化闭子概形，则
\begin{enumerate}
\item $Y$ 的奇点至多为结点；
\item 若 $Z \subset X$ 是 $X \setminus Y$ 的概形论闭包，则
\begin{enumerate}
\item 概形论交 $Y \cap Z$ 是 $k$ 的若干有限可分扩张的谱的不交并；
\item $Y \cap Z$ 的每个点都是 $X$ 的结点；
\item $Y \to \Spec(k)$ 在 $Y \cap Z$ 的每个点处光滑。
\end{enumerate}
\end{enumerate}
\end{lemma}

\begin{proof}
由于 $X$ 与 $Y$ 都约化、等维且维数为 $1$，可知 $Y$
是 $X$ 的一部分不可约分支的概形论并（在约化环中，$(0)$
是诸极小素理想之交）。设 $y \in Y$ 为闭点。若 $y$ 位于
$X \to \Spec(k)$ 的光滑轨迹中，则 $y$ 位于 $X$ 的唯一一个
不可约分支上，并且 $Y$ 与 $X$ 在 $y$ 的某个开邻域中相同。
故 $Y \to \Spec(k)$ 在 $y$ 处光滑。若 $y$ 是 $X$ 的结点，
但仍只位于 $X$ 的一个不可约分支上，则由同一论证，$y$
是 $Y$ 的结点。假设 $y$ 位于多于 $1$ 个不可约分支上，
而这些分支属于 $X$。由于 $X$ 在 $y$ 处的几何分支数为
$2$（引理 \ref{curves-lemma-nodal}），由《性质》，引理
\ref{properties-lemma-number-of-branches-irreducible-components}，
不可约分支的数目至多为 $2$，并且这些分支通过 $y$。若 $Y$ 同时包含这两个
分支，则 $Y = X$ 在 $y$ 的某个开邻域中再次成立，且 $y$
是 $Y$ 的结点。最后，假设 $Y$ 只包含其中一个不可约分支。
把 $X$ 替换为 $x$ 的一个开邻域后，可以假设 $Y$ 是这两个
不可约分支之一，而 $Z$ 是另一个。再次由《性质》，引理
\ref{properties-lemma-number-of-branches-irreducible-components}，
可知 $X$ 在 $y$ 处有两个分支；也就是说，局部环
$\mathcal{O}_{X, y}$ 有两个分支，它们分别来自
$\mathcal{O}_{Y, y}$ 与 $\mathcal{O}_{Z, y}$。写成
$\mathcal{O}_{X, y}^\wedge \cong \kappa(y)[[u, v]]/(uv)$
，如评注 \ref{curves-remark-trivial-quadratic-extension} 所示。例如由引理
\ref{curves-lemma-nodal}，域 $\kappa(y)$ 是 $k$ 的有限可分扩张。
因此，必要时交换 $u$ 与 $v$ 的作用后，映射
$\mathcal{O}_{X, y} \to \mathcal{O}_{Y, Y}$
的完备化对应于 $\kappa(y)[[u, v]]/(uv) \to \kappa(y)[[u]]$，
而映射
$\mathcal{O}_{X, y} \to \mathcal{O}_{Y, Y}$
的完备化对应于 $\kappa(y)[[u, v]]/(uv) \to \kappa(y)[[v]]$。
$Y \cap Z$ 的概形论交由它们的理想之和截出；该理想之和在完备化中为
$(u, v)$，即极大理想。因此 (2)(a) 与 (2)(b) 显然。
最后，(2)(c) 成立：$\mathcal{O}_{Y, y}$ 的完备化正则，
所以 $\mathcal{O}_{Y, y}$ 正则（《代数进阶》，引理
\ref{more-algebra-lemma-completion-regular}）；又因
$\kappa(y)/k$ 可分，由《代数》，引理
\ref{algebra-lemma-separable-smooth}，它在某个开邻域中光滑。
\end{proof}




\section{结点曲线族}
\label{curves-section-families-nodal}

\noindent
在 Stacks 项目中，曲线是维数 $1$ 的不可约簇；但在文献中，
“半稳定曲线”或“结点曲线”通常并非不可约，并且常被假设为固有的，
尤其是在“半稳定曲线族”“结点曲线族”或“结点族”等说法中。
因此，要选取一个既与文献一致又在内部自洽的术语颇为困难。
此外，我们实际上想先研究下述引理中引入的概念
（该概念在源上是局部的）。

\begin{lemma}
\label{curves-lemma-nodal-family}
设 $f : X \to S$ 为概形态射。以下条件等价：
\begin{enumerate}
\item $f$ 平坦且局部有限表示，每个非空纤维
$X_s$ 都等维且维数为 $1$，并且 $X_s$ 的奇点至多为结点；
\item $f$ 是相对维数为 $1$ 的 syntomic 态射，并且闭子概形
$\text{Sing}(f) \subset X$ 由 $\Omega_{X/S}$ 的第一
Fitting 理想定义，且在 $S$ 上非分歧。
\end{enumerate}
\end{lemma}

\begin{proof}
回忆 $\text{Sing}(f)$ 的形成与基变换交换，见《除子》，引理
\ref{divisors-lemma-base-change-and-fitting-ideal-omega}。
于是本引理来自引理 \ref{curves-lemma-nodal-lci}、《态射》，引理
\ref{morphisms-lemma-syntomic-flat-fibres}，以及《态射》，引理
\ref{morphisms-lemma-unramified-etale-fibres}。（我们还使用显然的
《态射》，引理
\ref{morphisms-lemma-syntomic-locally-finite-presentation} 和
\ref{morphisms-lemma-syntomic-flat}。）
\end{proof}

\begin{definition}
\label{curves-definition-nodal-family}
设 $f : X \to S$ 为概形态射。若 $f$
{\it 相对维数为 $1$ 且至多为结点}，是指 $f$ 满足引理
\ref{curves-lemma-nodal-family} 中的等价条件。
\end{definition}

\noindent
之所以采用这个繁复术语，有以下几个原因\footnote{不过，如果您有更好的建议，
请给 Stacks 项目的维护者发送电子邮件。}。第一，我们要确保这个概念
不与文献中的其他概念混淆（见本节引言）。第二，可以设想该概念对
较高相对维数态射的若干推广（例如，可以要求态射在 \'etale 局部是
若干至多为结点的态射之复合；也可以要求态射的纤维维数更高，
但至多具有普通二重点）。

\begin{lemma}
\label{curves-lemma-smooth-relative-dimension-1}
相对维数为 $1$ 的光滑态射相对维数为 $1$ 且至多为结点。
\end{lemma}

\begin{proof}
略。
\end{proof}

\begin{lemma}
\label{curves-lemma-base-change-nodal-family}
设 $f : X \to S$ 相对维数为 $1$ 且至多为结点。
则 $f$ 的任意基变换也具有同一性质。
\end{lemma}

\begin{proof}
这是因为 syntomic 态射的基变换仍为 syntomic
（《态射》，引理 \ref{morphisms-lemma-base-change-syntomic}）；
相对维数为 $1$ 的态射之基变换仍有相对维数 $1$
（《态射》，引理
\ref{morphisms-lemma-base-change-relative-dimension-d}）；
$\text{Sing}(f)$ 的形成与基变换交换（《除子》，引理
\ref{divisors-lemma-base-change-and-fitting-ideal-omega}）；
并且非分歧态射的基变换仍非分歧（《态射》，引理
\ref{morphisms-lemma-base-change-unramified}）。
\end{proof}

\noindent
下面的引理说明，可以在纤维上检验一个态射是否相对维数为
$1$ 且至多为结点。

\begin{lemma}
\label{curves-lemma-locus-where-nodal}
设 $f : X \to S$ 为平坦且局部有限表示的概形态射。
则存在极大开子概形 $U \subset X$，使得
$f|_U : U \to S$ 相对维数为 $1$ 且至多为结点。
此外，$U$ 的形成与任意基变换交换。
\end{lemma}

\begin{proof}
由《态射》，引理
\ref{morphisms-lemma-set-points-where-fibres-lci}，存在这样的开集，
$f$ 在其上为 syntomic。因此可以假设 $f$ 是 syntomic 态射。
特别地，$f$ 是 Cohen--Macaulay 态射（《概形的对偶性》，引理
\ref{duality-lemma-lci-gorenstein} 和
\ref{duality-lemma-gorenstein-CM-morphism}）。因此 $X$ 是若干
开闭子概形的不交并，而 $f$ 在每个子概形上具有给定的相对维数；
见《态射》，引理
\ref{morphisms-lemma-flat-finite-presentation-CM-fibres-relative-dimension}。
这一分解被任意基变换保持；见《态射》，引理
\ref{morphisms-lemma-base-change-relative-dimension-d}。舍去除一个
部分之外的所有部分后，可以假设 $f$ 是相对维数为 $1$ 的
syntomic 态射。设 $\text{Sing}(f) \subset X$ 为
$\Omega_{X/S}$ 的第一 Fitting 理想定义的闭子概形。
存在极大开子概形 $W \subset \text{Sing}(f)$，使得
$W \to S$ 非分歧，并且它的形成与基变换交换
（《态射》，引理
\ref{morphisms-lemma-set-points-where-fibres-unramified}）。
又因为 $\text{Sing}(f)$ 的形成与基变换交换（《除子》，引理
\ref{divisors-lemma-base-change-and-fitting-ideal-omega}），可知
$$
U = (X \setminus \text{Sing}(f)) \cup W
$$
是 $X$ 的极大开子概形，使得
$f|_U : U \to S$ 相对维数为 $1$ 且至多为结点；
并且 $U$ 的形成与基变换交换。
\end{proof}

\begin{lemma}
\label{curves-lemma-nodal-family-precompose-etale}
设 $f : X \to S$ 相对维数为 $1$ 且至多为结点。
若 $Y \to X$ 是 \'etale 态射，则复合 $g : Y \to S$
相对维数为 $1$ 且至多为结点。
\end{lemma}

\begin{proof}
注意，$g$ 是平坦且局部有限表示的，因为它是两个平坦且局部
有限表示态射的复合（使用《态射》，引理
\ref{morphisms-lemma-etale-locally-finite-presentation}、
\ref{morphisms-lemma-etale-flat}、
\ref{morphisms-lemma-composition-finite-presentation} 和
\ref{morphisms-lemma-composition-flat}）。
因此只需证明纤维的奇点至多为结点。这来自引理
\ref{curves-lemma-node-etale-local}（以及如下事实：\'etale 态射与
光滑态射的复合仍光滑；见《态射》，引理
\ref{morphisms-lemma-etale-smooth-unramified} 和
\ref{morphisms-lemma-composition-smooth}）。
\end{proof}

\begin{lemma}
\label{curves-lemma-nodal-family-postcompose-etale}
设 $S' \to S$ 为概形的 \'etale 态射。
设 $f : X \to S'$ 相对维数为 $1$ 且至多为结点。
则复合 $g : X \to S$ 相对维数为 $1$ 且至多为结点。
\end{lemma}

\begin{proof}
注意，$g$ 是平坦且局部有限表示的，因为它是两个平坦且局部
有限表示态射的复合（使用《态射》，引理
\ref{morphisms-lemma-etale-locally-finite-presentation}、
\ref{morphisms-lemma-etale-flat}、
\ref{morphisms-lemma-composition-finite-presentation} 和
\ref{morphisms-lemma-composition-flat}）。因此只需证明 $g$
的纤维之奇点至多为结点。这来自引理
\ref{curves-lemma-node-over-separable-extension} 以及对光滑点的类似结果。
\end{proof}

\begin{lemma}
\label{curves-lemma-nodal-family-etale-local-source}
设 $f : X \to S$ 为概形态射。设 $\{U_i \to X\}$ 为
\'etale 覆盖。以下条件等价：
\begin{enumerate}
\item $f$ 相对维数为 $1$ 且至多为结点；
\item 每个 $U_i \to S$ 都相对维数为 $1$ 且至多为结点。
\end{enumerate}
换言之，相对维数为 $1$ 且至多为结点这一性质在源上是
\'etale 局部的。
\end{lemma}

\begin{proof}
一个方向已见于引理
\ref{curves-lemma-nodal-family-precompose-etale}。对另一个方向，
注意局部有限表示、平坦，或具有相对维数 $1$，都是在源上
\'etale 局部的性质（《下降》，引理
\ref{descent-lemma-locally-finite-presentation-fppf-local-source}、
\ref{descent-lemma-flat-fpqc-local-source} 和
\ref{descent-lemma-dimension-at-point}）。取纤维后，归约到
$S$ 为某个域的谱的情形。在这种情形下，结论来自引理
\ref{curves-lemma-node-etale-local}（以及光滑性在源上 \'etale 局部这一
事实；见《下降》，引理
\ref{descent-lemma-smooth-smooth-local-source}）。
\end{proof}

\begin{lemma}
\label{curves-lemma-nodal-family-fpqc-local-target}
设 $f : X \to S$ 为概形态射。设 $\{U_i \to S\}$ 为 fpqc 覆盖。
以下条件等价：
\begin{enumerate}
\item $f$ 相对维数为 $1$ 且至多为结点；
\item 每个 $X \times_S U_i \to U_i$ 都相对维数为 $1$
且至多为结点。
\end{enumerate}
换言之，相对维数为 $1$ 且至多为结点这一性质在靶上是 fpqc 局部的。
\end{lemma}

\begin{proof}
一个方向已见于引理 \ref{curves-lemma-base-change-nodal-family}。
对另一个方向，注意局部有限表示、平坦，或具有相对维数 $1$，
都是在靶上 fpqc 局部的性质（《下降》，引理
\ref{descent-lemma-descending-property-locally-finite-presentation}、
\ref{descent-lemma-descending-property-flat}，以及《态射》，引理
\ref{morphisms-lemma-dimension-fibre-after-base-change}）。
取纤维后，归约到 $S$ 为某个域的谱的情形。在这种情形下，
结论来自引理 \ref{curves-lemma-node-field-extension}（以及光滑性在靶上
fpqc 局部这一事实；见《下降》，引理
\ref{descent-lemma-descending-property-smooth}）。
\end{proof}

\begin{lemma}
\label{curves-lemma-descend-nodal-family}
设 $S = \lim S_i$ 为一族概形的有向系统之极限，转移态射均为仿射。
设 $0 \in I$，并设 $f_0 : X_0 \to Y_0$ 为 $S_0$ 上的概形态射。
假设 $S_0$、$X_0$、$Y_0$ 都拟紧且拟分离。设
$f_i : X_i \to Y_i$ 为 $f_0$ 到 $S_i$ 的基变换，并设
$f : X \to Y$ 为 $f_0$ 到 $S$ 的基变换。若
\begin{enumerate}
\item $f$ 相对维数为 $1$ 且至多为结点；
\item $f_0$ 局部有限表示，
\end{enumerate}
则存在 $i \geq 0$，使得 $f_i$ 相对维数为 $1$ 且至多为结点。
\end{lemma}

\begin{proof}
由《极限》，引理 \ref{limits-lemma-descend-syntomic}，存在 $i$，
使得 $f_i$ 为 syntomic。于是
$X_i = \coprod_{d \geq 0} X_{i, d}$ 是若干开闭子概形的不交并，
且 $X_{i, d} \to Y_i$ 具有相对维数 $d$；见《态射》，引理
\ref{morphisms-lemma-syntomic-relative-dimension}。根据《态射》，
引理 \ref{morphisms-lemma-dimension-fibre-after-base-change}
所给出的纤维维数在基变换下的行为，可知 $X \to X_i$
映入 $X_{i, 1}$。于是存在 $i' \geq i$，使得
$X_{i'} \to X_i$ 映入 $X_{i, 1}$；见《极限》，引理
\ref{limits-lemma-limit-contained-in-constructible}。因此
$f_{i'} : X_{i'} \to Y_{i'}$ 是相对维数为 $1$ 的 syntomic 态射
（再次使用《态射》，引理
\ref{morphisms-lemma-dimension-fibre-after-base-change}）。
考虑态射 $\text{Sing}(f_{i'}) \to Y_{i'}$。我们知道它到 $Y$
的基变换非分歧。因此由《极限》，引理
\ref{limits-lemma-descend-unramified}，增大 $i'$ 后，态射
$\text{Sing}(f_{i'}) \to Y_{i'}$ 变为非分歧。证明完成。
\end{proof}

\begin{lemma}
\label{curves-lemma-formal-local-structure-nodal-family}
设 $f : T \to S$ 为概形态射。设 $t \in T$ 的像为 $s \in S$。
假设
\begin{enumerate}
\item $f$ 在 $t$ 处平坦；
\item $\mathcal{O}_{S, s}$ 是 Noether 环；
\item $f$ 局部有限型；
\item $t$ 是纤维 $T_s$ 的分裂结点。
\end{enumerate}
则存在 $h \in \mathfrak m_s^\wedge$ 以及同构
$$
\mathcal{O}_{T, t}^\wedge \cong
\mathcal{O}_{S, s}^\wedge[[x, y]]/(xy - h)
$$
，它是 $\mathcal{O}_{S, s}^\wedge$-代数同构。
\end{lemma}

\begin{proof}
把 $S$ 替换为 $\Spec(\mathcal{O}_{S, s})$，并把 $T$ 替换为
到 $\Spec(\mathcal{O}_{S, s})$ 的基变换。此时 $T$ 局部
Noether，因而 $\mathcal{O}_{T, t}$ 是 Noether 环。置
$A = \mathcal{O}_{S, s}^\wedge$、$\mathfrak m = \mathfrak m_A$，
并置 $B = \mathcal{O}_{T, t}^\wedge$。由《代数进阶》，引理
\ref{more-algebra-lemma-flat-completion}，$A \to B$ 平坦。由于
$\mathcal{O}_{T, t}/\mathfrak m_s \mathcal{O}_{T, t} = \mathcal{O}_{T_s, t}$
，可知 $B/\mathfrak m B = \mathcal{O}_{T_s, t}^\wedge$。
由假设 (4) 和引理 \ref{curves-lemma-split-node}，存在
$\overline{u}, \overline{v} \in B/\mathfrak m B$，使得映射
$$
(A/\mathfrak m)[[x, y]] \longrightarrow B/\mathfrak m B,\quad
x \longmapsto \overline{u},
x \longmapsto \overline{v}
$$
为满射，且核为 $(xy)$。

\medskip\noindent
假设已有 $n \geq 1$ 以及 $u, v \in B$，它们映到
$\overline{u}, \overline{v}$，并且
$$
u v = h + \delta
$$
，其中 $h \in A$ 且 $\delta \in \mathfrak m^nB$。我们断言，
存在 $u', v' \in B$，满足
$u - u', v - v' \in \mathfrak m^n B$，并且
$$
u' v' = h' + \delta'
$$
，其中 $h' \in A$ 且 $\delta' \in \mathfrak m^{n + 1}B$。
为此，写成 $\delta = \sum f_i b_i$，其中
$f_i \in \mathfrak m^n$ 且 $b_i \in B$。再写成
$b_i = a_i + u b_{i, 1} + v b_{i, 2} + \delta_i$
，其中 $a_i \in A$、$b_{i, 1}, b_{i, 2} \in B$，且
$\delta_i \in \mathfrak m B$。这是可行的，因为 $B$ 的剩余域
等于 $A$ 的剩余域，并且 $u$ 与 $v$ 在 $B/\mathfrak m B$
中的像生成极大理想。然后置
$$
u' = u - \sum b_{i, 2}f_i,\quad
v' = v - \sum b_{i, 1}f_i
$$
，得到
$$
u'v' = h + \delta - \sum (b_{i, 1}u + b_{i, 2}v)f_i + \sum
c_{ij}f_if_j =
h + \sum a_if_i + \sum f_i \delta_i + \sum c_{ij}f_if_j
$$
，其中某些 $c_{i, j} \in B$。因此得到上述形式的等式，其中
$h' = h + \sum a_if_i$，且
$\delta' = \sum f_i \delta_i + \sum c_{ij}f_if_j$。

\medskip\noindent
作归纳论证，并从任意提升 $u_1, v_1 \in B$ 开始；它们提升
$\overline{u}, \overline{v}$。上一段的结果给出元素序列
$u_n, v_n \in B$ 以及 $h_n \in A$，使得
$u_n - u_{n + 1} \in \mathfrak m^n B$,
$v_n - v_{n + 1} \in \mathfrak m^n B$,
$h_n - h_{n + 1} \in \mathfrak m^n$,
并且 $u_n v_n - h_n \in \mathfrak m^n B$。由于 $A$ 与 $B$
完备，可以置
$u_\infty = \lim u_n$、$v_\infty = \lim v_n$，以及
$h_\infty = \lim h_n$，于是得到
$u_\infty v_\infty = h_\infty$ 于 $B$ 中成立。因此有 $A$-代数映射
$$
A[[x, y]]/(xy - h_\infty) \longrightarrow B
$$
，它把 $x$ 映到 $u_\infty$，把 $v$ 映到 $v_\infty$。
这是平坦 $A$-代数之间的映射，模去 $\mathfrak m$ 后是同构。
它模 $\mathfrak m$ 为满射，故由完备性和《代数》，引理
\ref{algebra-lemma-completion-generalities} 可知它是满射。
再应用《代数》，引理 \ref{algebra-lemma-mod-injective}，
可知它是同构。
\end{proof}

\begin{lemma}
\label{curves-lemma-genus-in-nodal-family-of-curves}
设 $f : X \to S$ 为概形态射。假设
\begin{enumerate}
\item $f$ 固有；
\item $f$ 相对维数为 $1$ 且至多为结点；
\item $f$ 的几何纤维连通。
\end{enumerate}
则 (a) $f_*\mathcal{O}_X = \mathcal{O}_S$，且此等式在任意
基变换后仍成立；(b) $R^1f_*\mathcal{O}_X$ 是有限局部自由
$\mathcal{O}_S$-模，其形成与任意基变换交换；(c)
$R^qf_*\mathcal{O}_X = 0$ 对 $q \geq 2$ 成立。
\end{lemma}

\begin{proof}
(a) 来自《概形的导出范畴》，引理
\ref{perfect-lemma-proper-flat-geom-red-connected}。由
《概形的导出范畴》，引理
\ref{perfect-lemma-proper-flat-h0}，在 $S$ 上局部可以写成
$Rf_*\mathcal{O}_X = \mathcal{O}_S \oplus P$，其中 $P$
是 Tor 振幅包含于 $[1, \infty)$ 的完美复形。回忆
$Rf_*\mathcal{O}_X$ 的形成与任意基变换交换
（《概形的导出范畴》，引理
\ref{perfect-lemma-flat-proper-perfect-direct-image-general}）。
所以对 $s \in S$ 有
$$
H^i(P \otimes_{\mathcal{O}_S}^\mathbf{L} \kappa(s)) =
H^i(X_s, \mathcal{O}_{X_s})
\text{ 对 }i \geq 1
$$
。除非 $i = 1$，否则该式为零，因为 $X_s$ 是维数 $1$
的 Noether 概形；见《上同调》，命题
\ref{cohomology-proposition-vanishing-Noetherian}。于是
$P = H^1(P)[-1]$，并且 $H^1(P)$ 有限局部自由，例如由
《代数进阶》，引理
\ref{more-algebra-lemma-lift-perfect-from-residue-field} 可知。
所有构造都与基变换相容，结论随即成立。
\end{proof}





\section{结点曲线族的 \'Etale 局部结构}
\label{curves-section-etale-local-nodal}


\noindent
考虑概形态射
$$
\Spec(\mathbf{Z}[u, v, a]/(uv - a))
\longrightarrow
\Spec(\mathbf{Z}[a])
$$
下一个引理说明，这个态射是结点曲线族的 \'etale 局部结构模型。

\begin{lemma}
\label{curves-lemma-etale-local-structure-nodal-family}
设 $f : X \to S$ 为概形态射。假设 $f$ 相对维数为 $1$ 且至多为结点。
设 $x \in X$ 是纤维 $X_s$ 的一个奇点。则存在概形的交换图
$$
\xymatrix{
X \ar[d] &
U \ar[rr] \ar[l] \ar[rd] & &
W \ar[r] \ar[ld] &
\Spec(\mathbf{Z}[u, v, a]/(uv - a)) \ar[d] \\
S & &
V \ar[ll] \ar[rr] & & \Spec(\mathbf{Z}[a])
}
$$
其中 $X \leftarrow U$、$S \leftarrow V$ 与 $U \to W$ 是 \'etale
态射，右侧方块是笛卡尔方块，并且存在一个点 $u \in U$ 映到
$x$（在 $X$ 中）。
\end{lemma}

\begin{proof}[第一种证明：使用 Artin 逼近]
首先使用绝对 Noether 逼近，把情形归约到有限型
$\mathbf{Z}$-概形。问题在 $X$ 与 $S$ 上是局部的，因此可假设
$X$ 与 $S$ 都是仿射的。于是可写 $S = \Spec(R)$，并把
$R$ 写成滤过余极限 $R = \colim R_i$，其中
$\mathbf{Z}$-代数均为有限型。由《极限》，引理
\ref{limits-lemma-descend-finite-presentation}，可找到 $i$ 以及
$f_i : X_i \to \Spec(R_i)$，其到 $S$ 的基变换为 $f$。
增大 $i$ 后可假设 $f_i$ 相对维数为 $1$ 且至多为结点，见引理
\ref{curves-lemma-descend-nodal-family}。$x_i \in X_i$ 是 $x$ 的像，
并且是其纤维的奇点，例如这是因为
$\text{Sing}(f)$ 的形成与基变换交换（《除子》，引理
\ref{divisors-lemma-base-change-and-fitting-ideal-omega}）。
若能对 $f_i : X_i \to S_i$ 与 $x_i$ 证明本引理，则通过基变换
即可得到 $f : X \to S$ 的结论。因此归约到下一段研究的情形。

\medskip\noindent
假设 $S$ 在 $\mathbf{Z}$ 上有限型。设 $s \in S$ 为 $x$ 的像。
回忆 $\kappa(x)$ 是 $\kappa(s)$ 的有限可分扩张，例如因为
$\text{Sing}(f) \to S$ 非分歧，或者因为 $x$ 是纤维 $X_s$
的结点，从而可应用引理 \ref{curves-lemma-nodal}。此外，设
$\kappa'/\kappa(x)$ 为次数为 $2$ 的、评注
\ref{curves-remark-quadratic-extension} 中与 $\mathcal{O}_{X_s, x}$
相伴的二次可分代数。由《态射进阶》，引理
\ref{more-morphisms-lemma-realize-prescribed-residue-field-extension-etale}，
可以选取一个 \'etale 邻域 $(V, v) \to (S, s)$，使扩张
$\kappa(v)/\kappa(s)$ 在 $\kappa(x)/\kappa(s)$ 的情形（即
$\kappa' \cong \kappa(x) \times \kappa(x)$）中实现前一扩张，
而在扩张 $\kappa'/\kappa(s)$ 中，若 $\kappa'$ 为域则实现之。
把 $X$ 替换为 $X \times_S V$，把 $S$
替换为 $V$，就归约到下一段的情形。

\medskip\noindent
假设 $S$ 在 $\mathbf{Z}$ 上有限型，且 $x \in X_s$ 是分裂结点，
见定义 \ref{curves-definition-split-node}。由引理
\ref{curves-lemma-formal-local-structure-nodal-family}，存在
$\mathcal{O}_{S, s}$-代数同构
$$
\mathcal{O}_{X, x}^\wedge \cong
\mathcal{O}_{S, s}^\wedge[[s, t]]/(st - h)
$$
，其中 $h \in \mathfrak m_s^\wedge \subset \mathcal{O}_{S, s}^\wedge$。
换言之，考虑同态
$$
\sigma : \mathbf{Z}[a] \longrightarrow \mathcal{O}_{S, s}^\wedge
$$
将 $a$ 映到 $h$，则存在 $\mathcal{O}_{S, s}$-代数同构
$$
\mathcal{O}_{X, x}^\wedge
\longrightarrow
\mathcal{O}_{Y_\sigma, y_\sigma}^\wedge
$$
其中
$$
Y_\sigma = \Spec(\mathbf{Z}[u, v, t]/(uv - a))
\times_{\Spec(\mathbf{Z}[a]), \sigma} \Spec(\mathcal{O}_{S, s}^\wedge)
$$
且 $y_\sigma$ 是 $Y_\sigma$ 中位于
$\Spec(\mathcal{O}_{S, s}^\wedge)$ 闭点之上的点，其坐标
$u, v$ 均为零。由于 $\mathcal{O}_{S, s}$ 是 G-环（《代数进阶》，
命题 \ref{more-algebra-proposition-ubiquity-G-ring}），可应用
《态射进阶》，引理
\ref{more-morphisms-lemma-relative-map-approximation-pre}，结论得证。
\end{proof}

\begin{proof}[不使用 Artin 逼近的证明]
我们只略述此证明；它由 Mohan Swaminathan 提供，并仿照
\cite[Proposition X.2.1]{ACG} 中的论证。与前一个证明不同，
本证明使用实际方程。

\medskip\noindent
以与第一个证明完全相同的方式，我们归约到 $S$ 在
$\mathbf{Z}$ 上有限型且 $x \in X_s$ 为分裂结点的情形；
见定义 \ref{curves-definition-split-node}。由引理
\ref{curves-lemma-split-node}，存在 $\kappa(s)$-代数同构
$$
\mathcal{O}_{X_s, x}^\wedge \cong \kappa(s)[[u, v]]/(uv)
$$
注意 $X_s$ 在 $x$ 处的切空间维数为 $2$。因此由《簇》，引理
\ref{varieties-lemma-immersion-into-affine}，经 Zariski 缩小后可假设
存在闭浸入 $X \to Y$，该浸入在 $S$ 上，且 $Y \to S$ 光滑、相对维数
为 $2$。由于 $X$ 按定义是相对维数为 $1$ 的 syntomic
概形，可知 $X$ 是 $Y$ 中的有效 Cartier 除子（见《除子》，引理
\ref{divisors-lemma-immersion-lci-into-smooth-regular-immersion}）。
于是缩小 $Y$ 与 $X$ 后，可假设 $Y = \Spec(B)$ 为仿射的，
并存在非零因子 $g \in B$，使得 $X = \Spec(B/gB)$。

\medskip\noindent
写 $S = \Spec(R)$，记 $\mathfrak p \subset R$ 为对应于 $s$ 的素理想。
令 $\mathfrak q \subset B$ 为对应于 $x$ 的素理想（视为 $Y$
中的点）。有映射
$$
B_\mathfrak q / \mathfrak p B_\mathfrak q
\longrightarrow
B_\mathfrak q / g B_\mathfrak q + \mathfrak p B_\mathfrak q
\longrightarrow
\kappa(s)[[u, v]]/(uv)
$$
其中第二个箭头在完备化后成为同构。把 $B$ 替换为适当的主元局部化后，
可假设存在 $U, V \in B$，它们分别映到 $u$ 与 $v$，模去
$(u, v)^2$ 意义下相等。于是
$B_\mathfrak q/\mathfrak p B_\mathfrak q$ 的完备化是
$\kappa(s)[[U, V]]$。此外，缩小 $Y$ 并以单位乘 $g$ 后，可假设
$g$ 映到 $UV$ 加上 $\kappa(s)[[U, V]]$ 中的高阶项。再缩小 $Y$，
还可假设
\begin{enumerate}
\item $\mathfrak q = \mathfrak pB + (U, V)$
\item $\Omega_{B/R}$ 在 $\text{d}U$ 与 $\text{d}V$ 上自由
\end{enumerate}
写 $\text{d}(g) = g_2 \text{d}U + g_1 \text{d}V$，其中 $g_1, g_2 \in B$。
注意，理想 $J = (g_1, g_2) \subset B$ 与坐标 $U$、$V$ 的选择无关
（$g$ 固定）。模去 $\mathfrak p B_\mathfrak q + (U, V)^3B_\mathfrak q$
考察，可知 $g_1$ 与 $g_2$ 分别与 $U$、$V$ 同余，模去
$\mathfrak p B_\mathfrak q + (U, V)^2B_\mathfrak q$。
因此可把上面的论证以 $U$ 替换为 $g_1$、以 $V$ 替换为 $g_2$
重新进行。于是得到 $J = (U, V)$，从而
$g_1, g_2 \in (U, V) = J$（当然，这里的 $g_1$ 与 $g_2$ 与先前
坐标选择所得的同名元素不同）。

\medskip\noindent
注意 $\Spec(B/J) \to S = \Spec(R)$ 在 $x$ 处 \'etale。
因此把 $S$ 替换为一个 \'etale 邻域并再次缩小 $Y$ 后，可假设
$R \to B \to B/J$ 是同构。记 $r \in R$ 为其在 $B/J$ 中的像
等于 $-g$ 的像的元素。于是可写
$$
g + r \equiv a U + b V \bmod (U, V)^2
$$
其中 $a, b \in R$。利用 $\text{d}g \in J\Omega_{B/R}$ 求导，
可得 $a = b = 0$。因此可写
$$
g + r = \alpha U^2 + \beta UV + \gamma V^2
$$
其中 $\alpha, \beta, \gamma \in B$。于是 $\beta$ 映到 $1$ 于
$\kappa(x)$，而 $\alpha$ 与 $\gamma$ 映到 $0$ 于 $\kappa(x)$。
所以可假设 $\beta$ 与 $\alpha + \beta + \gamma$ 在 $B$ 中可逆。
考虑由下式给出的覆盖 $Y'$，它覆盖 $Y$：
$$
Y' = \Spec(B[t]/(t(1 - t) - \alpha \gamma / \beta^2))
$$
其中 $x'$ 是 $x$ 上满足 $t = 1$ 的唯一点。注意
$Y' \to Y$ 在 $x'$ 处 \'etale。考虑
$$
T_1 = (\alpha + t \beta)U + ((1 - t)\beta + \gamma)V
\quad\text{且}\quad
T_2 =
\frac{(\alpha + (1 - t) \beta)U + (t\beta + \gamma)V}{\alpha + \beta + \gamma}
$$
它们属于 $\Gamma(Y', \mathcal{O}_{Y'})$。于是
$$
g = T_1 T_2 - r
$$
由于 $T_1$ 与 $T_2$ 在上述意义下也是坐标，可知态射
$Y' \to \mathbf{A}^2_S$（由 $T_1$ 与 $T_2$ 给出）在 $x'$ 的某个邻域内
\'etale，且 $X' \subset Y'$ 是 $X$ 的拉回，
是 $\mathbf{A}^2_S = \Spec(R[x_1, x_2])$ 中由方程
$x_1x_2 - r = 0$ 给出的闭子概形的逆像（必要时再次缩小 $Y'$）。
我们把这些数据拼合成所需的交换图，细节留给读者。
\end{proof}




\section{更多消没结果}
\label{curves-section-more-vanishing}

\noindent
这是第 \ref{curves-section-vanishing} 节的继续。

\begin{lemma}
\label{curves-lemma-h1-nonzero-degree-leq-2g-2}
在情形 \ref{curves-situation-Cohen-Macaulay-curve} 中，假设 $X$ 不可约且
亏格为 $g$。令 $\mathcal{L}$ 为可逆的 $\mathcal{O}_X$-模。
令 $Z \subset X$ 为一个 $0$ 维闭子概形，其理想
层为 $\mathcal{I} \subset \mathcal{O}_X$。若 $H^1(X, \mathcal{I}\mathcal{L})$
非零，则
$$
\deg(\mathcal{L}) \leq 2g - 2 + \deg(Z)
$$
除非 $\mathcal{I}\mathcal{L} \cong \omega_X$，否则不等号严格成立。
\end{lemma}

\begin{proof}
任意曲线（例如 $X$）都是 Cohen--Macaulay 的。
若 $H^1(X, \mathcal{I}\mathcal{L})$ 非零，则存在非零态射
$\mathcal{I}\mathcal{L} \to \omega_X$，见引理
\ref{curves-lemma-duality-dim-1-CM}。
由于 $\mathcal{I}\mathcal{L}$ 无挠，该态射是单射。
由于域是 Gorenstein 且 $X$ 约化，我们得到 Gorenstein
轨迹 $U \subset X$ 在 $X$ 中非空，见《概形的对偶性》引理
\ref{duality-lemma-gorenstein}。
该引理还说明 $\omega_X|_U$ 可逆。
因此我们有短正合序列
$$
0 \to \mathcal{I}\mathcal{L} \to \omega_X \to \mathcal{Q} \to 0
$$
其中 $\mathcal{Q}$ 的支集是零维的。因此
\begin{align*}
0 & \leq \dim \Gamma(X, \mathcal{Q})\\
& =
\chi(\mathcal{Q}) \\
& =
\chi(\omega_X) - \chi(\mathcal{I}\mathcal{L}) \\
& =
\chi(\omega_X) - \deg(\mathcal{L}) - \chi(\mathcal{I}) \\
& =
2g - 2 - \deg(\mathcal{L}) + \deg(Z)
\end{align*}
由引理 \ref{curves-lemma-euler}、\ref{curves-lemma-rr}，公式 (\ref{curves-equation-genus})，
以及《簇论》引理 \ref{varieties-lemma-chi-tensor-finite}
和 \ref{varieties-lemma-degree-on-proper-curve} 得到。这里还用了
$\deg(Z) = \dim_k \Gamma(Z, \mathcal{O}_Z) = \chi(\mathcal{O}_Z)$
以及短正合序列
$0 \to \mathcal{I} \to \mathcal{O}_X \to \mathcal{O}_Z \to 0$.
引理得证。
\end{proof}

\begin{lemma}
\label{curves-lemma-degree-more-than-2g-2}
\begin{reference}
\cite[Lemma 2]{Jongmin}
\end{reference}
在情形 \ref{curves-situation-Cohen-Macaulay-curve} 中，
假设 $X$ 不可约且亏格为 $g$。
令 $\mathcal{L}$ 为可逆的 $\mathcal{O}_X$-模。
令 $Z \subset X$ 为一个 $0$ 维闭子概形，其理想
层为 $\mathcal{I} \subset \mathcal{O}_X$。
若 $\deg(\mathcal{L}) > 2g - 2 + \deg(Z)$，则
$H^1(X, \mathcal{I}\mathcal{L}) = 0$，且下列可能性之一
成立：
\begin{enumerate}
\item $H^0(X, \mathcal{I}\mathcal{L}) \not = 0$，或
\item $g = 0$ 且 $\deg(\mathcal{L}) = \deg(Z) - 1$。
\end{enumerate}
在情形 (2) 中，若 $Z = \emptyset$，则 $X \cong \mathbf{P}^1_k$，且 $\mathcal{L}$
对应于 $\mathcal{O}_{\mathbf{P}^1}(-1)$。
\end{lemma}

\begin{proof}
由引理 \ref{curves-lemma-h1-nonzero-degree-leq-2g-2}，得
$H^1(X, \mathcal{I}\mathcal{L})$ 消没。
若 $H^0(X, \mathcal{I}\mathcal{L}) = 0$，则
$\chi(\mathcal{I}\mathcal{L}) = 0$。由短正合
序列 $0 \to \mathcal{I}\mathcal{L} \to \mathcal{L} \to \mathcal{O}_Z \to 0$
可得 $\deg(\mathcal{L}) = g - 1 + \deg(Z)$。
于是 $g - 1 + \deg(Z) > 2g - 2 + \deg(Z)$，从而 $g = 0$，
故 (2) 成立。在情形 (2) 且 $Z = \emptyset$ 时，
$\mathcal{L}^{-1}$ 是次数为 $1$ 的可逆层。
这意味着存在同构 $X \to \mathbf{P}^1_k$，且由引理
\ref{curves-lemma-genus-zero-positive-degree}，$\mathcal{L}^{-1}$ 是
$\mathcal{O}_{\mathbf{P}^1}(1)$ 的拉回。
\end{proof}

\begin{lemma}
\label{curves-lemma-degree-more-than-2g}
\begin{reference}
\cite[Lemma 3]{Jongmin}
\end{reference}
在情形 \ref{curves-situation-Cohen-Macaulay-curve} 中，
假设 $X$ 不可约且亏格为 $g$。
令 $\mathcal{L}$ 为可逆的 $\mathcal{O}_X$-模。
若 $\deg(\mathcal{L}) \geq 2g$，则 $\mathcal{L}$
由整体截面生成。
\end{lemma}

\begin{proof}
令 $Z \subset X$ 为由 $\mathcal{L}$ 的整体截面
切出的闭子概形。由引理 \ref{curves-lemma-degree-more-than-2g-2}
可知 $Z \not = X$。令 $\mathcal{I} \subset \mathcal{O}_X$
为切出 $Z$ 的理想层。考虑短正合序列
$$
0 \to \mathcal{I}\mathcal{L}
\to \mathcal{L} \to \mathcal{O}_Z \to 0
$$
若 $Z \not = \emptyset$，则由上同调的长正合序列可知
$H^1(X, \mathcal{I}\mathcal{L})$ 非零。
由引理 \ref{curves-lemma-duality-dim-1-CM} 得到一个
非零因而为单射的态射
$$
\mathcal{I}\mathcal{L}
\longrightarrow
\omega_X
$$
特别地，我们得到单射
$H^0(X, \mathcal{L}) = H^0(X, \mathcal{I}\mathcal{L})
\to H^0(X, \omega_X)$。这不可能，因为
$$
\dim_k H^0(X, \mathcal{L}) = \dim_k H^1(X, \mathcal{L}) +
\deg(\mathcal{L}) + 1 - g \geq g + 1
$$
而由 (\ref{curves-equation-genus}) 有 $\dim H^0(X, \omega_X) = g$。
\end{proof}

\begin{lemma}
\label{curves-lemma-degree-more-than-2g-1-and-Z}
在情形 \ref{curves-situation-Cohen-Macaulay-curve} 中，
假设 $X$ 不可约且亏格为 $g$。
令 $\mathcal{L}$ 为可逆的 $\mathcal{O}_X$-模。
令 $Z \subset X$ 为非空的 $0$ 维闭子概形。
若 $\deg(\mathcal{L}) \geq 2g - 1 + \deg(Z)$，则 $\mathcal{L}$
由整体截面生成，且 $H^0(X, \mathcal{L}) \to H^0(X, \mathcal{L}|_Z)$
是满射。
\end{lemma}

\begin{proof}
由引理 \ref{curves-lemma-degree-more-than-2g}，可知它由整体截面生成。
若 $\mathcal{I} \subset \mathcal{O}_X$ 是 $Z$ 的理想层，则由
引理 \ref{curves-lemma-h1-nonzero-degree-leq-2g-2} 有
$H^1(X, \mathcal{I}\mathcal{L}) = 0$。故限制态射满射。
\end{proof}

\begin{lemma}
\label{curves-lemma-degree-at-least-2g+1}
在情形 \ref{curves-situation-Cohen-Macaulay-curve} 中，
假设 $X$ 在 $k$ 上几何不可约且亏格为 $g$。
令 $\mathcal{L}$ 为可逆的 $\mathcal{O}_X$-模。
若 $\deg(\mathcal{L}) \geq 2g + 1$，则 $\mathcal{L}$
非常丰沛。
\end{lemma}

\begin{proof}
由引理 \ref{curves-lemma-degree-more-than-2g}，$\mathcal{L}$
由整体截面生成，因而确定一个态射
$f : X \to \mathbf{P}^n_k$，其中 $n = h^0(X,\mathcal{L}) - 1$。要证明
$\mathcal{L}$ 非常丰沛，就是要证明
$f$ 是闭浸入。只需检查 $f$ 到代数闭包 $\overline{k}$（$k$ 的）
的基变换是闭浸入（《下降》引理
\ref{descent-lemma-descending-property-closed-immersion}）。
因此可假设 $k$ 代数闭；由假设 $X$ 仍不可约。
引理 \ref{curves-lemma-degree-more-than-2g-1-and-Z}
说明对每个 $0$ 维闭子概形
$Z\subset X$，
限制态射 $H^0(X, \mathcal{L}) \to H^0(X, \mathcal{L}|_Z)$
都是满射。由《簇论》引理
\ref{varieties-lemma-variant-separate-points-tangent-vectors}，
$\mathcal{L}$ 非常丰沛。
\end{proof}


\begin{lemma}
\label{curves-lemma-vanishing-on-gorenstein}
\begin{reference}
\cite[Lemma 4]{Jongmin} 的弱形式
\end{reference}
令 $k$ 为一个域。令 $X$ 为 $k$ 上的 proper 概形，
它是约化、连通且维数为 $1$。
令 $\mathcal{L}$ 为可逆的 $\mathcal{O}_X$-模。
令 $Z \subset X$ 为一个 $0$ 维闭子概形，其理想
层为 $\mathcal{I} \subset \mathcal{O}_X$。
若 $H^1(X, \mathcal{I}\mathcal{L}) \not = 0$，则存在
一个约化的连通闭子概形 $Y \subset X$，
维数为 $1$，满足
$$
\deg(\mathcal{L}|_Y) \leq -2\chi(Y, \mathcal{O}_Y) + \deg(Z \cap Y)
$$
其中 $Z \cap Y$ 是概形论交。
\end{lemma}

\begin{proof}
若 $H^1(X, \mathcal{I}\mathcal{L})$ 非零，则存在非零态射
$\varphi : \mathcal{I}\mathcal{L} \to \omega_X$，见引理
\ref{curves-lemma-duality-dim-1-CM}。令 $Y \subset X$
为不可约分支 $C$（属于 $X$）的并，其中
$\varphi$ 在 $C$ 的一般点非零。则 $Y$ 是约化闭子概形。
令 $\mathcal{J} \subset \mathcal{O}_X$ 为 $Y$ 的理想层。
由于 $\mathcal{J}\mathcal{I}\mathcal{L}$
（作为 $\mathcal{L}$ 的子模）没有嵌入关联点，且 $\varphi$ 在
$\mathcal{J}$ 的支集的一般点处为零
（由 $Y$ 的选择及 $X$ 约化所致），可知 $\varphi$ 分解为
$$
\mathcal{I}\mathcal{L} \to
\mathcal{I}\mathcal{L}/\mathcal{J}\mathcal{I}\mathcal{L} \to \omega_X
$$
我们可以把 $\mathcal{I}\mathcal{L}/\mathcal{J}\mathcal{I}\mathcal{L}$
视为 $Y$ 上某个凝聚层的推前，并滥用记号仍用同一符号表示。
由于 $\omega_Y = \SheafHom(\mathcal{O}_Y, \omega_X)$，
由引理 \ref{curves-lemma-closed-immersion-dim-1-CM}
可得到态射
$$
\mathcal{I}\mathcal{L}/
\mathcal{J}\mathcal{I}\mathcal{L} 
\to \omega_Y
$$
它是 $\mathcal{O}_Y$-模的态射，并在 $Y$ 的一般点处为单射。
令 $\mathcal{I}' \subset \mathcal{O}_Y$ 为 $Z \cap Y$ 的理想
层。存在态射
$\mathcal{I}\mathcal{L}/\mathcal{J}\mathcal{I}\mathcal{L} \to
\mathcal{I}'\mathcal{L}|_Y$，其核支集在闭点上。
由于 $\omega_Y$ 是 Cohen--Macaulay 模，以上态射
经由单射 $\mathcal{I}'\mathcal{L}|_Y \to
\omega_Y$ 分解。于是得到短正合序列
$$
0 \to \mathcal{I}'\mathcal{L}|_Y \to \omega_Y \to \mathcal{Q} \to 0
$$
这是 $Y$ 上凝聚层的正合序列，其中 $\mathcal{Q}$ 的支集维数为 $0$
（这里用了 $\omega_Y$ 在 $Y$ 的一般点处是可逆模这一事实）。因此
$$
0 \leq \dim \Gamma(Y, \mathcal{Q}) =
\chi(\mathcal{Q}) = \chi(\omega_Y) - \chi(\mathcal{I}'\mathcal{L}) =
-2\chi(\mathcal{O}_Y) - \deg(\mathcal{L}|_Y) + \deg(Z \cap Y)
$$
由引理 \ref{curves-lemma-euler} 及《簇论》引理
\ref{varieties-lemma-chi-tensor-finite} 得到。
若 $Y$ 连通，则引理得证。
否则，对 $Y$ 的某个连通分支重复论证的最后部分。
\end{proof}

\begin{lemma}
\label{curves-lemma-global-generation}
令 $k$ 为一个域。令 $X$ 为 $k$ 上的 proper 概形，
它是约化、连通且维数为 $1$。
令 $\mathcal{L}$ 为可逆的 $\mathcal{O}_X$-模。
假设对每个约化连通闭子概形 $Y \subset X$（维数为 $1$），都有
$$
\deg(\mathcal{L}|_Y) \geq 2\dim_k H^1(Y, \mathcal{O}_Y)
$$
则 $\mathcal{L}$ 由整体截面生成。
\end{lemma}

\begin{proof}
对 $X$ 的不可约分支数作归纳。
若 $X$ 不可约，则把 $X$ 视为域
$k' = H^0(X, \mathcal{O}_X)$ 上的概形，应用引理
\ref{curves-lemma-degree-more-than-2g} 即得本引理。假设 $X$ 不可约。
在继续之前，若 $k$ 有限，则把 $k$ 替换为纯超越扩张 $K$。
这由《簇论》的下列引理所允许：
\ref{varieties-lemma-globally-generated-base-change},
\ref{varieties-lemma-degree-base-change},
\ref{varieties-lemma-geometrically-reduced-any-base-change}，以及
\ref{varieties-lemma-bijection-irreducible-components},
《概形的上同调》引理 \ref{coherent-lemma-flat-base-change-cohomology}、
引理 \ref{curves-lemma-sanity-check-duality}，以及 $K$ 在 $k$ 上几何不可约这一
初等事实。

\medskip\noindent
为导出矛盾，假设 $\mathcal{L}$ 不由整体截面生成。
于是可取凝聚理想层
$\mathcal{I} \subset \mathcal{O}_X$，使得
$H^0(X, \mathcal{I}\mathcal{L}) = H^0(X, \mathcal{L})$，
并且 $\mathcal{O}_X/\mathcal{I}$ 非零且
支集维数为 $0$。例如，可取 $\mathcal{I}$ 为 $\mathcal{L}$ 的整体截面
公共消没轨迹中任意闭点的理想层。
考虑短正合序列
$$
0 \to \mathcal{I}\mathcal{L} \to \mathcal{L} \to 
\mathcal{L}/\mathcal{I}\mathcal{L} \to 0
$$
由于 $\mathcal{L}/\mathcal{I}\mathcal{L}$ 的支集
维数为 $0$，可知 $\mathcal{L}/\mathcal{I}\mathcal{L}$
由整体截面生成（《簇论》引理
\ref{varieties-lemma-chi-tensor-finite}）。
由短正合序列以及
$H^0(X, \mathcal{I}\mathcal{L}) = H^0(X, \mathcal{L})$ 这一事实，得到单射
$H^0(X, \mathcal{L}/\mathcal{I}\mathcal{L}) \to H^1(X, \mathcal{I}\mathcal{L})$.

\medskip\noindent
回忆 $k$-向量空间 $H^1(X, \mathcal{I}\mathcal{L})$
是 $\Hom(\mathcal{I}\mathcal{L}, \omega_X)$ 的对偶。
取 $\varphi : \mathcal{I}\mathcal{L} \to \omega_X$。
由引理 \ref{curves-lemma-vanishing-on-gorenstein} 有
$H^1(X, \mathcal{L}) = 0$。因此
$$
\dim_k H^0(X, \mathcal{I}\mathcal{L}) = \dim_k H^0(X, \mathcal{L}) =
\deg(\mathcal{L}) + \chi(\mathcal{O}_X) > \dim_k H^1(X, \mathcal{O}_X) =
\dim_k H^0(X, \omega_X)
$$
我们得到 $\varphi$ 在整体截面上不是单射，特别地
$\varphi$ 不是单射。对每个不可约分支的一般点 $\eta \in X$
（它属于 $X$ 的不可约分支），记
$V_\eta \subset \Hom(\mathcal{I}\mathcal{L}, \omega_X)$ 为由 $k$-子向量
空间中的那些 $\varphi$ 构成的空间，它们在 $\eta$ 处为零。
由于 $\mathcal{I}\mathcal{L}$ 的每个关联点
都是 $X$ 的一般点，上述说明
$\Hom(\mathcal{I}\mathcal{L}, \omega_X) = \bigcup V_\eta$。
因为 $X$ 只有有限个一般点且 $k$ 是无限域，故有
$\Hom(\mathcal{I}\mathcal{L}, \omega_X) = V_\eta$ 对某个 $\eta$ 成立。
令 $\eta \in C \subset X$ 为相应的不可约分支。
令 $Y \subset X$ 为 $X$ 的其他不可约分支之并。
则 $Y$ 是非空的约化闭子概形且不等于 $X$。
令 $\mathcal{J} \subset \mathcal{O}_X$ 为 $Y$ 的理想层。
请注意，$\mathcal{J}$ 的支集是 $C$。

\medskip\noindent
令 $\varphi : \mathcal{I}\mathcal{L} \to \omega_X$ 为任意态射。
由于 $\mathcal{J}\mathcal{I}\mathcal{L}$
没有嵌入关联点（作为 $\mathcal{L}$ 的子模），且 $\varphi$ 在
支集的一般点 $\eta$（属于 $\mathcal{J}$）处为零，可知
$\varphi$ 分解为
$$
\mathcal{I}\mathcal{L} \to
\mathcal{I}\mathcal{L}/\mathcal{J}\mathcal{I}\mathcal{L} \to \omega_X
$$
我们可以把 $\mathcal{I}\mathcal{L}/\mathcal{J}\mathcal{I}\mathcal{L}$
视为 $Y$ 上某个凝聚层的推前，并滥用记号仍用同一符号表示。
由于 $\omega_Y = \SheafHom(\mathcal{O}_Y, \omega_X)$，
由引理 \ref{curves-lemma-closed-immersion-dim-1-CM}
得到分解
$$
\mathcal{I}\mathcal{L} \to
\mathcal{I}\mathcal{L}/
\mathcal{J}\mathcal{I}\mathcal{L} 
\xrightarrow{\varphi'} \omega_Y \to \omega_X
$$
这是 $\varphi$ 的分解。令 $\mathcal{I}' \subset \mathcal{O}_Y$ 为
$\mathcal{I} \subset \mathcal{O}_X$ 的像。存在满射
$\mathcal{I}\mathcal{L}/\mathcal{J}\mathcal{I}\mathcal{L} \to
\mathcal{I}'\mathcal{L}|_Y$，其核支集在闭点上。
由于 $\omega_Y$ 是 $Y$ 上的 Cohen--Macaulay 模，态射 $\varphi'$
经由态射
$\varphi'' : \mathcal{I}'\mathcal{L}|_Y \to \omega_Y$.
于是有交换图
$$
\vcenter{
\xymatrix{
0 \ar[r] &
\mathcal{I}\mathcal{L} \ar[r] \ar[d] &
\mathcal{L} \ar[r] \ar[d] &
\mathcal{L}/\mathcal{I}\mathcal{L} \ar[r] \ar[d] &
0 \\
0 \ar[r] &
\mathcal{I}'\mathcal{L}|_Y \ar[r] &
\mathcal{L}|_Y \ar[r] &
\mathcal{L}|_Y/\mathcal{I}'\mathcal{L}|_Y \ar[r] &
0
}
}
\quad\text{且}\quad
\vcenter{
\xymatrix{
\mathcal{I}\mathcal{L} \ar[r]_\varphi \ar[d] & \omega_X \\
\mathcal{I}'\mathcal{L}|_Y \ar[r]^{\varphi''} & \omega_Y \ar[u]
}
}
$$
现在可以如下完成证明：
由于对每个 $\varphi$ 都有一个 $\varphi''$，且
$\omega_X \in \textit{Coh}(\mathcal{O}_X)$
表示函子 $\mathcal{F} \mapsto \Hom_k(H^1(X, \mathcal{F}), k)$，
可知
$H^1(X, \mathcal{I}\mathcal{L}) \to H^1(Y, \mathcal{I}'\mathcal{L}|_Y)$
是单射。由于边界态射
$H^0(X, \mathcal{L}/\mathcal{I}\mathcal{L}) \to H^1(X, \mathcal{I}\mathcal{L})$
是单射，故复合态射
$$
H^0(X, \mathcal{L}/\mathcal{I}\mathcal{L}) \to
H^0(X, \mathcal{L}|_Y/\mathcal{I}'\mathcal{L}|_Y) \to
H^1(X, \mathcal{I}'\mathcal{L}|_Y)
$$
是单射。由于
$\mathcal{L}/\mathcal{I}\mathcal{L} \to
\mathcal{L}|_Y/\mathcal{I}'\mathcal{L}|_Y$
是凝聚模之间的满射，且其支集维数为 $0$，可知第一个态射
$H^0(X, \mathcal{L}/\mathcal{I}\mathcal{L}) \to
H^0(X, \mathcal{L}|_Y/\mathcal{I}'\mathcal{L}|_Y)$
是满射（因而是双射）。但由归纳假设，$\mathcal{L}|_Y$ 由整体
截面生成（当然，即使 $Y$ 不连通这仍然成立），故边界态射
$$
H^0(X, \mathcal{L}|_Y/\mathcal{I}'\mathcal{L}|_Y) \to
H^1(X, \mathcal{I}'\mathcal{L}|_Y)
$$
不可能是单射。
矛盾证明了结论。
\end{proof}





\section{压缩有理尾}
\label{curves-section-contracting-rational-tails}

\noindent
本节讨论压缩概形以改善其典范层正性性质的最简单情形。

\begin{example}[压缩有理尾]
\label{curves-example-rational-tail}
令 $k$ 为一个域。令 $X$ 为 $k$ 上维数为 $1$ 的 proper 概形，且
$H^0(X, \mathcal{O}_X) = k$。假设 $X$ 的奇点至多为结点。
称一个{\it 有理尾}为不可约分支 $C \subset X$
（视为整闭子概形），满足以下性质：
\begin{enumerate}
\item $X' \not = \emptyset$，其中 $X' \subset X$ 是 $X \setminus C$ 的
概形论闭包；
\item 概形论交 $C \cap X'$ 是单个
约化点 $x$；
\item $H^0(C, \mathcal{O}_C)$ 同构地映到
$x$ 的剩余域；
\item $C$ 的亏格为零。
\end{enumerate}
由于至少有两个 $X$ 的不可约分支经过
$x$，故 $x$ 是结点。
置 $k' = H^0(C, \mathcal{O}_C) = \kappa(x)$。
则 $k'/k$ 是有限可分域扩张
（引理 \ref{curves-lemma-nodal}）。存在典范态射
$$
c : X \longrightarrow X'
$$
它在 $X'$ 上诱导恒等态射，并将 $C$ 映到 $x \in X'$；该映射通过
典范态射 $C \to \Spec(k') = x$ 给出。这由
《态射》引理 \ref{morphisms-lemma-scheme-theoretic-union} 得到，
因为 $X$ 是 $C$ 与 $X'$ 的概形论并（$X$ 约化）。
此外，我们断言
$$
c_*\mathcal{O}_X = \mathcal{O}_{X'}
\quad\text{且}\quad
R^1c_*\mathcal{O}_X = 0
$$
为此，记 $i_C : C \to X$、$i_{X'} : X' \to X$ 及 $i_x : x \to X$
为这些嵌入，并使用正合序列
$$
0 \to \mathcal{O}_X \to
i_{C, *}\mathcal{O}_C \oplus i_{X', *}\mathcal{O}_{X'} \to
i_{x, *}\kappa(x) \to 0
$$
（《态射》引理 \ref{morphisms-lemma-scheme-theoretic-union}）。
考察高阶直接像的长正合序列可知，只需证明
$H^0(C, \mathcal{O}_C) = k'$ 且 $H^1(C, \mathcal{O}_C) = 0$，
而这由假设得到。
注意 $X'$ 也是 $k$ 上维数为 $1$ 的 proper 概形，
其奇点至多为结点
（引理 \ref{curves-lemma-closed-subscheme-nodal-curve}），
满足 $H^0(X', \mathcal{O}_{X'}) = k$，且
$X'$ 与 $X$ 有相同亏格。
我们称 $c : X \to X'$ 为
{\it 有理尾的压缩}。
\end{example}

\begin{lemma}
\label{curves-lemma-rational-tail-negative}
令 $k$ 为域。令 $X$ 为 $k$ 上维数为 $1$ 的 proper 概形，且
$H^0(X, \mathcal{O}_X) = k$。假设 $X$ 的奇点至多为结点。
令 $C \subset X$ 为有理尾
（例 \ref{curves-example-rational-tail}）。则 $\deg(\omega_X|_C) < 0$。
\end{lemma}

\begin{proof}
令 $X' \subset X$ 如例中所述。则有短正合
序列
$$
0 \to \omega_C \to \omega_X|_C \to \mathcal{O}_{C \cap X'} \to 0
$$
见引理 \ref{curves-lemma-closed-subscheme-reduced-gorenstein}、
\ref{curves-lemma-facts-about-nodal-curves} 及
\ref{curves-lemma-closed-subscheme-nodal-curve}。
取例中的 $k'$，则 $\deg(\omega_C) = -2[k' : k]$，
因为由命题 \ref{curves-proposition-projective-line} 有
$C \cong \mathbf{P}^1_{k'}$，且 $\deg(C \cap X') = [k' : k]$。
所以 $\deg(\omega_X|_C) = -[k' : k]$，为负数。
\end{proof}

\begin{lemma}
\label{curves-lemma-rational-tail-field-extension}
令 $k$ 为域。令 $X$ 为 $k$ 上维数为 $1$ 的 proper 概形，且
$H^0(X, \mathcal{O}_X) = k$。假设 $X$ 的奇点至多为结点。
令 $C \subset X$ 为有理尾
（例 \ref{curves-example-rational-tail}）。
对任意域扩张 $K/k$，基变换 $C_K \subset X_K$
是有限个互不相交有理尾的并。
\end{lemma}

\begin{proof}
令 $x \in C$，并令 $k' = \kappa(x)$ 如例中所述。
由命题 \ref{curves-proposition-projective-line}，有
$C \cong \mathbf{P}^1_{k'}$。
由于 $k'/k$ 是有限可分扩张，有
$k' \otimes_k K = K'_1 \times \ldots \times K'_n$，
它是有限个有限可分扩张 $K'_i/K$ 的乘积。
置 $C_i = \mathbf{P}^1_{K'_i}$，并记 $x_i \in C_i$
为 $x$ 的逆像。则 $C_K = \coprod C_i$ 且
$X'_K \cap C_i = x_i$，如所需。
\end{proof}

\begin{lemma}
\label{curves-lemma-no-rational-tail}
令 $k$ 为域。令 $X$ 为 $k$ 上维数为 $1$ 的 proper 概形，且
$H^0(X, \mathcal{O}_X) = k$。假设 $X$ 的奇点至多为结点。
若 $X$ 没有有理尾
（例 \ref{curves-example-rational-tail}），
则对每个约化连通闭子概形
$Y \subset X$，若 $Y \not = X$ 且维数为 $1$，有
$\deg(\omega_X|_Y) \geq \dim_k H^1(Y, \mathcal{O}_Y)$.
\end{lemma}

\begin{proof}
令 $Y \subset X$ 如陈述所述。则 $k' = H^0(Y, \mathcal{O}_Y)$
是域且是 $k$ 的有限扩张，并且 $[k' : k]$
整除下面与 $Y$ 及 $Y$ 上凝聚层相关的所有数值不变量，见
《簇论》引理 \ref{varieties-lemma-divisible}。
令 $Z \subset X$ 如引理
\ref{curves-lemma-closed-subscheme-reduced-gorenstein} 所述。
以下不再特别指出地使用该引理以及引理
\ref{curves-lemma-facts-about-nodal-curves} 和
\ref{curves-lemma-closed-subscheme-nodal-curve} 的结果。
于是得到短正合序列
$$
0 \to \omega_Y \to \omega_X|_Y \to \mathcal{O}_{Y \cap Z} \to 0
$$
见引理 \ref{curves-lemma-closed-subscheme-reduced-gorenstein}。
从而
$$
\deg(\omega_X|_Y) = \deg(Y \cap Z) + \deg(\omega_Y) =
\deg(Y \cap Z) - 2\chi(Y, \mathcal{O}_Y)
$$
因此，若引理不成立，则
$$
2[k' : k] > \deg(Y \cap Z) + \dim_k H^1(Y, \mathcal{O}_Y)
$$
由于 $Y \cap Z$ 非空，且由上述整除性可知，
只有在 $Y \cap Z$ 是单个 $k'$-有理点，且该点位于 $Y$ 的光滑轨迹中，
并且 $H^1(Y, \mathcal{O}_Y) = 0$ 时才可能发生这种情况。
若 $Y$ 不可约，则这意味着 $Y$ 是有理尾。
若 $Y$ 可约，则由于 $\deg(\omega_X|_Y) = -[k' : k]$，
可知存在不可约分支 $C$，属于 $Y$，满足
$\deg(\omega_X|_C) < 0$，见
《簇论》引理 \ref{varieties-lemma-degree-in-terms-of-components}。
将上述分析应用于 $C$，可知 $C$ 是有理尾。
\end{proof}

\begin{lemma}
\label{curves-lemma-no-rational-tail-semiample-genus-geq-2}
令 $k$ 为域。令 $X$ 为 $k$ 上维数为 $1$ 的 proper 概形，且
$H^0(X, \mathcal{O}_X) = k$。假设 $X$ 的奇点至多为结点。
假设 $X$ 没有有理尾
（例 \ref{curves-example-rational-tail}）。若
\begin{enumerate}
\item $X$ 的亏格为 $0$，则 $X$ 同构于不可约平面圆锥，且
$\omega_X^{\otimes -1}$ 非常丰沛；
\item $X$ 的亏格为 $1$，则 $\omega_X \cong \mathcal{O}_X$；
\item $X$ 的亏格满足 $\geq 2$，则
$\omega_X^{\otimes m}$ 对 $m \geq 2$ 由整体截面生成。
\end{enumerate}
\end{lemma}

\begin{proof}
由引理 \ref{curves-lemma-facts-about-nodal-curves}，$X$ 是
Gorenstein 的，即 $\omega_X$ 是可逆的 $\mathcal{O}_X$-模。

\medskip\noindent
若 $X$ 的亏格为零，则 $\deg(\omega_X) < 0$；因此若
$X$ 有多于一个不可约分支，就与引理
\ref{curves-lemma-no-rational-tail} 矛盾。在不可约情形，
由引理 \ref{curves-lemma-genus-zero} 可知 $X$ 同构于不可约平面圆锥，且
$\omega_X^{\otimes -1}$ 非常丰沛。

\medskip\noindent
若 $X$ 的亏格为 $1$，则 $\omega_X$ 有整体截面，且
$\deg(\omega_X|_C) = 0$ 对每个不可约分支成立。
确切地说，由引理 \ref{curves-lemma-no-rational-tail}，有
$\deg(\omega_X|_C) \geq 0$ 对所有不可约分支 $C$ 成立；由引理
\ref{curves-lemma-genus-gorenstein}，这些数之和为
$0$，于是可应用《簇论》引理
\ref{varieties-lemma-degree-in-terms-of-components}。
再由《簇论》引理
\ref{varieties-lemma-no-sections-dual-nef}，有
$\omega_X \cong \mathcal{O}_X$。

\medskip\noindent
假设亏格为 $g$ 的 $X$ 大于或等于 $2$。
若 $X$ 不可约，则由引理
\ref{curves-lemma-degree-more-than-2g} 即得结论。
假设 $X$ 可约。
由引理 \ref{curves-lemma-no-rational-tail}，引理
\ref{curves-lemma-global-generation} 的不等式对陈述中的每个
$Y \subset X$ 都成立，但 $Y = X$ 除外。分析引理
\ref{curves-lemma-global-generation} 的证明可知（在可约情形）
对 $Y = X$ 所用的唯一不等式是
$$
\deg(\omega_X^{\otimes m}) > -2 \chi(\mathcal{O}_X)
\quad\text{且}\quad
\deg(\omega_X^{\otimes m}) + \chi(\mathcal{O}_X) > \dim_k H^1(X, \mathcal{O}_X)
$$
由于在假设 $g \geq 2$ 且 $m \geq 2$ 下二者均成立，故得证。
\end{proof}

\begin{lemma}
\label{curves-lemma-contracting-rational-tails}
令 $k$ 为域。令 $X$ 为 $k$ 上维数为 $1$ 的 proper 概形，且
$H^0(X, \mathcal{O}_X) = k$。假设 $X$ 的奇点至多为结点。
考虑序列
$$
X = X_0 \to X_1 \to \ldots \to X_n = X'
$$
它由有理尾的压缩（例 \ref{curves-example-rational-tail}）组成，
直到不再有有理尾。则
\begin{enumerate}
\item 若 $X$ 的亏格为 $0$，则 $X'$ 是不可约平面圆锥；
\item 若 $X$ 的亏格为 $1$，则 $\omega_{X'} \cong \mathcal{O}_X$；
\item 若 $X$ 的亏格为 $> 1$，则
$\omega_{X'}^{\otimes m}$ 对 $m \geq 2$ 由整体截面生成。
\end{enumerate}
若 $X$ 的亏格满足 $\geq 1$，则态射 $X \to X'$
与选择无关，且该态射的形成与基域扩张相交换。
\end{lemma}

\begin{proof}
我们不断压缩有理尾，直到不再有有理尾。
由引理 \ref{curves-lemma-no-rational-tail-semiample-genus-geq-2}，
可知 (1)、(2)、(3) 成立。

\medskip\noindent
唯一性。要证明 $f : X \to X'$ 与所作选择无关，
只需证明：任意有理尾 $C \subset X$ 都被 $X \to X'$
映到一点；细节略去。
若不然，则存在截面
$s \in \Gamma(X', \omega_{X'}^{\otimes 2})$，它在不可约分支
$f(C)$ 的一般点处不为零。
由于每个压缩 $X_i \to X_{i + 1}$
都有截面 $X_{i + 1} \to X_i$，故存在截面
$X' \to X$，它是 $f$ 的截面。于是有正合序列
$$
0 \to \omega_{X'} \to \omega_X \to \omega_X|_{X''} \to 0
$$
其中 $X'' \subset X$ 是被 $f$ 压缩的不可约分支之并。
见引理 \ref{curves-lemma-closed-subscheme-reduced-gorenstein}。
于是得到态射 $\omega_{X'}^{\otimes 2} \to \omega_X^{\otimes 2}$，
取 $s$ 的像得到 $\omega_X^{\otimes 2}$ 的一个截面，
它在 $C$ 的一般点处不为零。
这与 $\omega_X$ 限制到有理尾上的次数为负相矛盾
（引理 \ref{curves-lemma-rational-tail-negative}）。

\medskip\noindent
关于基域扩张的陈述由引理
\ref{curves-lemma-rational-tail-field-extension} 得到。细节略去。
\end{proof}





\section{压缩有理桥}
\label{curves-section-contracting-rational-bridges}

\noindent
本节讨论继第 \ref{curves-section-contracting-rational-tails} 节所述情形之后，
压缩概形以改善其典范层正性性质的下一个最简单情形。

\begin{example}[压缩有理桥]
\label{curves-example-rational-bridge}
令 $k$ 为域。令 $X$ 为 $k$ 上维数为 $1$ 的 proper 概形，且
$H^0(X, \mathcal{O}_X) = k$。假设 $X$ 的奇点至多为结点。
称一个{\it 有理桥}为不可约分支 $C \subset X$
（视为整闭子概形），满足以下性质：
\begin{enumerate}
\item $X' \not = \emptyset$，其中 $X' \subset X$ 是 $X \setminus C$ 的
概形论闭包；
\item 概形论交 $C \cap X'$ 的次数为 $2$（相对于
$H^0(C, \mathcal{O}_C)$）；
\item $C$ 的亏格为零。
\end{enumerate}
置 $k' = H^0(C, \mathcal{O}_C)$ 且
$k'' = H^0(C \cap X', \mathcal{O}_{C \cap X'})$。
则 $k'$ 是域（《簇论》引理
\ref{varieties-lemma-proper-geometrically-reduced-global-sections}），
且 $\dim_{k'}(k'') = 2$。由于 $X$ 的至少
两个不可约分支经过 $C \cap X'$ 的每个点，故这些点是 $X$ 的结点，
并且在 $C$ 与 $X'$ 上均为光滑点
（引理 \ref{curves-lemma-closed-subscheme-nodal-curve}）。
因此 $k'/k$ 是有限可分域扩张，
且 $k''/k'$ 要么是次数为 $2$ 的可分域扩张，
要么 $k'' = k' \times k'$（引理 \ref{curves-lemma-nodal}）。由
第 \ref{curves-section-pushouts} 节，存在推出
$$
\xymatrix{
C \cap X' \ar[r] \ar[d] &
X' \ar[d]^a \\
\Spec(k') \ar[r] &
Y
}
$$
它具有许多良好性质（以下使用时不再逐一说明）。令 $y \in Y$ 为
$\Spec(k') \to Y$ 的像。则
$$
\mathcal{O}_{Y, y}^\wedge \cong k'[[s, t]]/(st)
\quad\text{或}\quad
\mathcal{O}_{Y, y}^\wedge \cong
\{f \in k''[[s]] : f(0) \in k'\}
$$
取决于 $C \cap X'$ 有 $2$ 个还是 $1$ 个点。
这由引理 \ref{curves-lemma-complete-local-ring-pushout} 以及事实
$\mathcal{O}_{X', p} \cong \kappa(p)[[t]]$ 得到，其中
$p \in C \cap X'$，而该事实见《代数补遗》引理
\ref{more-algebra-lemma-power-series-over-residue-field}。
因此 $y \in Y$ 是结点，见引理 \ref{curves-lemma-nodal} 与
\ref{curves-lemma-2-branches-delta-1}，特别是引理
\ref{curves-lemma-2-branches-delta-1} 中 (2) $\Rightarrow$ (1) 的证明里
情形 II 的讨论。
故 $Y$ 的奇点至多为结点。

\medskip\noindent
可将上面的交换图扩展为图
$$
\xymatrix{
C \cap X' \ar[r] \ar[d] &
X' \ar[d]^a \ar[r] & X \ar[ld]^c &
C \ar[ld] \ar[l] \\
\Spec(k') \ar[r] &
Y &
\Spec(k') \ar[l]
}
$$
其中下方两条水平态射相同。确切地说，$X$ 是 $X'$ 与 $C$ 的
概形论并（因此由《态射》引理
\ref{morphisms-lemma-scheme-theoretic-union} 为推出），
且态射 $C \to Y$ 与 $X' \to Y$ 在 $C \cap X'$ 上相同。
最后，我们断言
$$
c_*\mathcal{O}_X = \mathcal{O}_Y
\quad\text{且}\quad
R^1c_*\mathcal{O}_X = 0
$$
为此使用正合序列
$$
0 \to \mathcal{O}_X \to
\mathcal{O}_C \oplus \mathcal{O}_{X'} \to
\mathcal{O}_{C \cap X'} \to 0
$$
（《态射》引理 \ref{morphisms-lemma-scheme-theoretic-union}）。
高阶直接像的长正合序列为
$$
0 \to c_*\mathcal{O}_X \to c_*\mathcal{O}_C \oplus c_*\mathcal{O}_{X'} \to
c_*\mathcal{O}_{C \cap X'} \to
R^1c_*\mathcal{O}_X \to
R^1c_*\mathcal{O}_C \oplus R^1c_*\mathcal{O}_{X'}
$$
由于 $c|_{X'} = a$ 是仿射态射，故
$R^1c_*\mathcal{O}_{X'} = 0$.
由于 $c|_C$ 分解为 $C \to \Spec(k') \to X$，且
$C$ 亏格为零，故 $R^1c_*\mathcal{O}_C = 0$。
由于 $\mathcal{O}_{X'} \to \mathcal{O}_{C \cap X'}$ 是
满射，且 $c|_{X'}$ 是仿射态射，可知
$c_*\mathcal{O}_{X'} \to c_*\mathcal{O}_{C \cap X'}$
是满射。这证明了 $R^1c_*\mathcal{O}_X = 0$。
最后有
$\mathcal{O}_Y = c_*\mathcal{O}_X$，这是由正合序列以及
《态射补遗》命题
\ref{more-morphisms-proposition-pushout-along-closed-immersion-and-integral}
中对推出的结构层的描述得到的。

\medskip\noindent
综上，$Y$ 也是 $k$ 上维数为 $1$ 的 proper 概形，且
$H^0(Y, \mathcal{O}_Y) = k$，其奇点至多为结点
（引理 \ref{curves-lemma-closed-subscheme-nodal-curve}），
并且 $Y$ 与 $X$ 有相同亏格。
我们称 $c : X \to Y$ 为
{\it 有理桥的压缩}。
\end{example}

\begin{lemma}
\label{curves-lemma-rational-bridge-zero}
令 $k$ 为域。令 $X$ 为 $k$ 上维数为 $1$ 的 proper 概形，且
$H^0(X, \mathcal{O}_X) = k$。假设 $X$ 的奇点至多为结点。
令 $C \subset X$ 为有理桥
（例 \ref{curves-example-rational-bridge}）。则 $\deg(\omega_X|_C) = 0$。
\end{lemma}

\begin{proof}
令 $X' \subset X$ 如例中所述。则有短正合
序列
$$
0 \to \omega_C \to \omega_X|_C \to \mathcal{O}_{C \cap X'} \to 0
$$
见引理 \ref{curves-lemma-closed-subscheme-reduced-gorenstein}、
\ref{curves-lemma-facts-about-nodal-curves} 及
\ref{curves-lemma-closed-subscheme-nodal-curve}。
取例中的 $k''/k'/k$，有
$\deg(\omega_C) = -2[k' : k]$，因为 $C$ 的亏格为 $0$
（引理 \ref{curves-lemma-rr}），且
$\deg(C \cap X') = [k'' : k] = 2[k' : k]$。
因此 $\deg(\omega_X|_C) = 0$。
\end{proof}

\begin{lemma}
\label{curves-lemma-rational-bridge-field-extension}
令 $k$ 为域。令 $X$ 为 $k$ 上维数为 $1$ 的 proper 概形，且
$H^0(X, \mathcal{O}_X) = k$。假设 $X$ 的奇点至多为结点。
令 $C \subset X$ 为有理桥
（例 \ref{curves-example-rational-bridge}）。
对任意域扩张 $K/k$，基变换 $C_K \subset X_K$
是有限个互不相交有理桥的并。
\end{lemma}

\begin{proof}
令 $k''/k'/k$ 如例中所述。
由于 $k'/k$ 是有限可分扩张，有
$k' \otimes_k K = K'_1 \times \ldots \times K'_n$，
它是有限个有限可分扩张 $K'_i/K$ 的乘积。
相应的乘积分解
$k'' \otimes_k K = \prod K''_i$ 给出次数为 $2$ 的
可分代数扩张 $K''_i/K'_i$。
置 $C_i = C_{K'_i}$。则 $C_K = \coprod C_i$，
因而每个 $C_i$ 的亏格为 $0$（视为 $K'_i$ 上的曲线），
因为由平坦基变换有 $H^1(C_K, \mathcal{O}_{C_K}) = 0$。
最后，$X'_K \cap C_i = \Spec(K''_i)$ 的次数为 $2$（相对于 $K'_i$），
如所需。
\end{proof}

\begin{lemma}
\label{curves-lemma-rational-bridge-canonical}
令 $c : X \to Y$ 为有理桥的压缩
（例 \ref{curves-example-rational-bridge}）。
则 $c^*\omega_Y \cong \omega_X$。
\end{lemma}

\begin{proof}
可以通过直接计算证明这一点，但我们更愿意使用如下特征：将
$\omega_X$ 看作拟凝聚 $\mathcal{O}_X$-模，它表示函子
$\textit{Coh}(\mathcal{O}_X) \to \textit{Sets}$,
$\mathcal{F} \mapsto \Hom_k(H^1(X, \mathcal{F}), k) =
H^1(X, \mathcal{F})^\vee$，见引理
\ref{curves-lemma-duality-dim-1-CM}，或《概形的对偶性》引理
\ref{duality-lemma-dualizing-module-proper-over-A}。

\medskip\noindent
严格地说，记 $\mathcal{C}_Y$ 为如下范畴：对象是可逆的
$\mathcal{O}_Y$-模，态射是
$\mathcal{O}_Y$-模同态。记 $\mathcal{C}_X$ 为如下范畴：对象是可逆的
$\mathcal{O}_X$-模 $\mathcal{L}$，满足
$\mathcal{L}|_C \cong \mathcal{O}_C$，态射是
$\mathcal{O}_Y$-模同态。我们断言函子
$$
c^* : \mathcal{C}_Y \to \mathcal{C}_X
$$
是范畴等价。确切地说，由《态射的更多性质》引理
\ref{more-morphisms-lemma-bijection-on-Pic}，它本质满射。再由投影公式
（《上同调》引理 \ref{cohomology-lemma-projection-formula}）
得到 $c_*c^*\mathcal{N} = \mathcal{N}$，因而 $c^*$
是范畴等价，其拟逆为 $c_*$。

\medskip\noindent
我们断言 $\omega_X$ 是 $\mathcal{C}_X$ 的对象。事实上，有短正合列
$$
0 \to \omega_C \to \omega_X|_C \to \mathcal{O}_{C \cap X'} \to 0
$$
见引理 \ref{curves-lemma-closed-subscheme-reduced-gorenstein}。
取次数可得 $\deg(\omega_X|_C) = 0$（略去一个小细节）。
因此由引理 \ref{curves-lemma-genus-zero-pic}，$\omega_X|_C$ 是平凡的，
从而 $\omega_X$ 是 $\mathcal{C}_X$ 的对象。

\medskip\noindent
由于 $R^1c_*\mathcal{O}_X = 0$，投影公式表明
$R^1c_*c^*\mathcal{N} = 0$，其中
$\mathcal{N} \in \Ob(\mathcal{C}_Y)$。
因此，Leray 谱序列
（《上同调》引理 \ref{cohomology-lemma-apply-Leray}）给出如下
$$
\xymatrix{
\mathcal{C}_Y \ar[rr]_{c^*} \ar[dr]_{H^1(Y, -)^\vee} & &
\mathcal{C}_X \ar[ld]^{H^1(X, -)^\vee} \\
& \textit{Sets}
}
$$
范畴与函子的图表是交换的。由于
$\omega_Y \in \Ob(\mathcal{C}_Y)$ 表示东南方向的函子，且
$\omega_X \in \Ob(\mathcal{C}_X)$ 表示东南方向的函子，
由 Yoneda 引理（《范畴》引理 \ref{categories-lemma-yoneda}）即得结论。
\end{proof}

\begin{lemma}
\label{curves-lemma-no-rational-bridge-ample-genus-geq-2}
令 $k$ 为域。令 $X$ 为 $k$ 上的真概形，维数为 $1$，且
$H^0(X, \mathcal{O}_X) = k$。假设
\begin{enumerate}
\item $X$ 的奇点至多为结点，
\item $X$ 没有有理尾
（例 \ref{curves-example-rational-tail}），
\item $X$ 没有有理桥
（例 \ref{curves-example-rational-bridge}），
\item $g$ 为 $X$ 的亏格，且 $\geq 2$。
\end{enumerate}
则 $\omega_X$ 是丰沛的。
\end{lemma}

\begin{proof}
只需证明 $\deg(\omega_X|_C) > 0$ 对每个不可约分支 $C$（它属于 $X$）都成立，见《簇》引理
\ref{varieties-lemma-ampleness-in-terms-of-degrees-components}。
若 $X = C$ 不可约，则由 $g \geq 2$
及引理 \ref{curves-lemma-genus-gorenstein} 得到结论。
否则，置 $k' = H^0(C, \mathcal{O}_C)$。这是
$k$ 的一个域有限扩张，并且 $[k' : k]$
整除下文与 $C$ 及在 $C$ 上的拟凝聚层相关的所有数值不变量，见
《簇》引理 \ref{varieties-lemma-divisible}。
令 $X' \subset X$ 为 $X \setminus C$ 的闭包，如
引理 \ref{curves-lemma-closed-subscheme-reduced-gorenstein} 中所述。
下文将直接使用该引理以及引理
\ref{curves-lemma-facts-about-nodal-curves} 和
\ref{curves-lemma-closed-subscheme-nodal-curve} 的结果。
于是得到短正合列
$$
0 \to \omega_C \to \omega_X|_C \to \mathcal{O}_{C \cap X'} \to 0
$$
见引理 \ref{curves-lemma-closed-subscheme-reduced-gorenstein}。因此
$$
\deg(\omega_X|_C) = \deg(C \cap X') + \deg(\omega_C) =
\deg(C \cap X') - 2\chi(C, \mathcal{O}_C)
$$
所以若引理不成立，则
$$
2[k' : k] \geq \deg(C \cap X') + 2\dim_k H^1(C, \mathcal{O}_C)
$$
由于 $C \cap X'$ 非空，且由上面提到的整除性，只有以下两种情形之一才可能发生：
\begin{enumerate}
\item[(a)] $C \cap X'$ 是一个 $k'$-有理点，且属于 $C$，
$H^1(C, \mathcal{O}_C) = 0$；
\item[(b)] $C \cap X'$ 的次数为 $2$（在 $k'$ 上），且
$H^1(C, \mathcal{O}_C) = 0$。
\end{enumerate}
第一种可能意味着 $C$ 是有理尾，第二种意味着 $C$ 是有理桥。
由于两者都被排除，证明完成。
\end{proof}

\begin{lemma}
\label{curves-lemma-contracting-rational-bridges}
令 $k$ 为域。令 $X$ 为 $k$ 上维数为 $1$ 的真概形，满足
$H^0(X, \mathcal{O}_X) = k$ 且亏格为 $g \geq 2$。
假设 $X$ 的奇点至多为结点，且
$X$ 没有有理尾。考虑一列
$$
X = X_0 \to X_1 \to \ldots \to X_n = X'
$$
有理桥的压缩
（例 \ref{curves-example-rational-bridge}），直至不再有有理桥。
则 $\omega_{X'}$ 是丰沛的。
态射 $X \to X'$ 与选择无关，且该态射的构造与基域扩张相容。
\end{lemma}

\begin{proof}
我们不断压缩有理桥，直至不再有有理桥。由引理
\ref{curves-lemma-no-rational-bridge-ample-genus-geq-2}，此时
$\omega_{X'}$ 是丰沛的。

\medskip\noindent
记 $f : X \to X'$ 为其复合。由
引理 \ref{curves-lemma-rational-bridge-canonical} 及归纳法可知
$f^*\omega_{X'} = \omega_X$。
有 $f_*\mathcal{O}_X = \mathcal{O}_{X'}$，
因为有理桥的压缩具有这一性质。因此投影公式给出
$f_*f^*\mathcal{L} = \mathcal{L}$，对所有可逆的
$\mathcal{O}_{X'}$-模 $\mathcal{L}$ 均成立。
于是
$$
\Gamma(X', \omega_{X'}^{\otimes m}) = \Gamma(X, \omega_X^{\otimes m})
$$
对所有 $m$ 成立。由于 $X'$ 是这些空间直和的 Proj，
由《态射》引理 \ref{morphisms-lemma-proper-ample-is-proj}，
我们得到态射 $X \to X'$ 完全是典范的。

\medskip\noindent
令 $K/k$ 为域扩张，则
$\omega_{X_K}$ 是 $\omega_X$ 的拉回
（引理 \ref{curves-lemma-sanity-check-duality}），并且由
《概形的上同调》引理
\ref{coherent-lemma-flat-base-change-cohomology}，有
$\Gamma(X, \omega_X^{\otimes m}) \otimes_k K$
等于
$\Gamma(X_K, \omega_{X_K}^{\otimes m})$。
因此由上述论证，$f : X \to X'$ 的构造与
$K/k$ 的基变换相容。细节略。
\end{proof}






\section{压缩为稳定曲线}
\label{curves-section-contracting-to-stable}

\noindent
本节结合在
节 \ref{curves-section-contracting-rational-tails} 与
\ref{curves-section-contracting-rational-bridges} 中得到的压缩态射。
具体地，设 $k$ 为域，令 $X$ 为 $k$ 上维数为 $1$ 的真概形，满足
$H^0(X, \mathcal{O}_X) = k$ 且亏格 $g \geq 2$。
假设 $X$ 的奇点至多为结点。将引理
\ref{curves-lemma-contracting-rational-tails} 的态射与
引理 \ref{curves-lemma-contracting-rational-bridges} 的态射复合，得到态射
$$
c : X \longrightarrow Y
$$
使得 $Y$ 也是 $k$ 上维数为 $1$ 的真概形，奇点至多为结点，且
$k = H^0(Y, \mathcal{O}_Y)$，亏格为 $g$，并满足
$\mathcal{O}_Y = c_*\mathcal{O}_X$ 及 $R^1c_*\mathcal{O}_X = 0$，
以及 $\omega_Y$ 在 $Y$ 上丰沛。
引理 \ref{curves-lemma-characterize-contraction-to-stable}
表明这些条件事实上刻画了该态射。

\begin{lemma}
\label{curves-lemma-contract-gorenstein-canonical}
令 $k$ 为域。令 $c : X \to Y$ 为 $k$ 上真概形之间的态射。
假设
\begin{enumerate}
\item $\mathcal{O}_Y = c_*\mathcal{O}_X$ 且 $R^1c_*\mathcal{O}_X = 0$，
\item $X$ 与 $Y$ 都约化、Gorenstein，且维数为 $1$，
\item 满足 $\exists\ m \in \mathbf{Z}$，使得
$H^1(X, \omega_X^{\otimes m}) = 0$ 且 $\omega_X^{\otimes m}$
由整体截面生成。
\end{enumerate}
则 $c^*\omega_Y \cong \omega_X$。
\end{lemma}

\begin{proof}
由《态射的更多性质》定理
\ref{more-morphisms-theorem-stein-factorization-Noetherian}，$c$ 的纤维几何连通。
特别地，$c$ 是满射。
在有限多个闭点 $y = y_1, \ldots, y_r$（它们属于 $Y$）处 $X_y$ 的维数为 $1$，而在
$Y \setminus \{y_1, \ldots, y_r\}$ 上态射 $c$ 是同构。
细节略；提示：在 $\{y_1, \ldots, y_r\}$ 之外，
态射 $c$ 是有限的，见《概形的上同调》引理
\ref{coherent-lemma-characterize-finite}。

\medskip\noindent
下面仔细构造映射 $b : c^*\omega_Y \to \omega_X$。
记 $f : X \to \Spec(k)$ 及 $g : Y \to \Spec(k)$ 为结构态射。
我们有 $f^!k = \omega_X[1]$ 及 $g^!k = \omega_Y[1]$，见
引理 \ref{curves-lemma-duality-dim-1} 及其证明。于是
$f^! = c^! \circ g^!$，从而 $c^!\omega_Y = \omega_X$。
因此有函子性同构
$$
\Hom_{D(\mathcal{O}_X)}(\mathcal{F}, \omega_X)
\longrightarrow
\Hom_{D(\mathcal{O}_Y)}(Rc_*\mathcal{F}, \omega_Y)
$$
对拟凝聚 $\mathcal{O}_X$-模 $\mathcal{F}$，这是按 $c^!$ 的定义得到的\footnote{它是《概形的对偶性》引理
\ref{duality-lemma-twisted-inverse-image} 中右伴随函子限制到
$D^+_\QCoh(\mathcal{O}_Y)$ 后的结果。}。
该同构由迹映射 $t : Rc_*\omega_X \to \omega_Y$
（伴随的余单位）诱导。由投影公式
（《上同调》引理 \ref{cohomology-lemma-projection-formula}），
典范映射 $a : \omega_Y \to Rc_*c^*\omega_Y$ 是同构。
结合上述内容，得到满足下式的典范映射
$b : c^*\omega_Y \to \omega_X$：
$$
t \circ Rc_*(b) = a^{-1}
$$
特别地，将 $b$ 限制到 $c^{-1}(Y \setminus \{y_1, \ldots, y_r\})$
后，它是同构（因为它是可逆模之间的映射，且与另一个映射的复合给出同构 $a^{-1}$）。

\medskip\noindent
取 (3) 中的 $m \in \mathbf{Z}$，考虑映射
$$
b^{\otimes m} :
\Gamma(Y, \omega_Y^{\otimes m})
\longrightarrow
\Gamma(X, \omega_X^{\otimes m})
$$
由于 $Y$ 约化，且由构造中提到的 $b$ 的最后一个性质，该映射是单射。
由 Riemann--Roch（引理 \ref{curves-lemma-rr}）有
$\chi(X, \omega_X^{\otimes m}) =\chi(Y, \omega_Y^{\otimes m})$.
因此
$$
\dim_k \Gamma(Y, \omega_Y^{\otimes m}) \geq
\dim_k \Gamma(X, \omega_X^{\otimes m}) = \chi(X, \omega_X^{\otimes m})
$$
从而 $b^{\otimes m}$ 在整体截面上诱导同构。因此
$b^{\otimes m} : c^*\omega_Y^{\otimes m} \to \omega_X^{\otimes m}$
是满射，因为 $\omega_X^{\otimes m}$ 的生成元都在像中。
于是 $b^{\otimes m}$ 是同构，因而 $b$ 是同构。
\end{proof}

\begin{lemma}
\label{curves-lemma-characterize-contraction-to-stable}
令 $k$ 为域。令 $X$ 为 $k$ 上维数为 $1$ 的真概形，满足
$H^0(X, \mathcal{O}_X) = k$ 且亏格为 $g \geq 2$。
假设 $X$ 的奇点至多为结点。
存在唯一的态射（在唯一同构意义下）
$$
c : X \longrightarrow Y
$$
它是 $k$ 上概形之间的态射，并具有下列性质：
\begin{enumerate}
\item $Y$ 在 $k$ 上真，$\dim(Y) = 1$，且 $Y$ 的奇点至多为结点；
\item $\mathcal{O}_Y = c_*\mathcal{O}_X$ 且 $R^1c_*\mathcal{O}_X = 0$；
\item $\omega_Y$ 在 $Y$ 上丰沛。
\end{enumerate}
\end{lemma}

\begin{proof}
存在性：由本节引言中所述，结合引理
\ref{curves-lemma-contracting-rational-tails} 与
\ref{curves-lemma-contracting-rational-bridges}，存在具有所列全部性质的态射。
此外，它可写成复合
$$
X \to X_1 \to X_2 \ldots \to X_n \to X_{n + 1} \to \ldots \to X_{n + n'}
$$
其中前 $n$ 个态射是有理尾的压缩，后 $n'$ 个态射是有理桥的压缩。
注意，性质 (2) 对每个有理尾的压缩（例 \ref{curves-example-rational-tail}）
和有理桥的压缩（例 \ref{curves-example-rational-bridge}）都成立。
容易看出，该性质在态射的复合下保持。

\medskip\noindent
唯一性：令 $c : X \to Y$ 为满足条件
(1)、(2)、(3) 的态射。我们将证明存在唯一同构
$X_{n + n'} \to Y$，使其与态射 $X \to X_{n + n'}$ 及 $c$ 相容。

\medskip\noindent
开始证明前，先对 $c$ 作一些观察。
首先由《态射的更多性质》定理
\ref{more-morphisms-theorem-stein-factorization-Noetherian}，$c$ 的纤维几何连通。
特别地，$c$ 是满射。
对闭点 $y \in Y$，纤维 $X_y$ 满足
$$
H^1(X_y, \mathcal{O}_{X_y}) = 0
\quad\text{且}\quad
H^0(X_y, \mathcal{O}_{X_y}) = \kappa(y)
$$
第一个等式由《态射的更多性质》引理
\ref{more-morphisms-lemma-check-h1-fibre-zero} 得到，第二个等式由《态射的更多性质》引理
\ref{more-morphisms-lemma-h1-fibre-zero-check-h0-kappa} 得到。
因此，要么 $X_y = x$，其中 $x$ 是 $X$ 中映到 $y$ 的唯一点，且与 $y$ 有相同的剩余域；
要么 $X_y$ 是 $\kappa(y)$ 上维数为 $1$ 的真概形。注意在第二种情形
$X_y$ 是 Cohen--Macaulay 的（引理 \ref{curves-lemma-automatic}）。
然而，由于 $X$ 约化，$X_y$ 在其所有一般点处必为约化的（细节略），因而由《性质》引理
\ref{properties-lemma-criterion-reduced}，$X_y$ 约化。
由此 $X_y$ 的奇点至多为结点
（引理 \ref{curves-lemma-closed-subscheme-nodal-curve}）。
注意 $X_y$ 的亏格为零（见上文）。
最后，只有有限多个点 $y$ 使得纤维 $X_y$ 的维数为 $1$，记为
$\{y_1, \ldots, y_r\}$，并且 $c^{-1}(Y \setminus \{y_1, \ldots, y_r\})$
经 $c$ 同构映到 $Y \setminus \{y_1, \ldots, y_r\}$。
细节略；提示：在 $\{y_1, \ldots, y_r\}$ 之外，
态射 $c$ 是有限的，见《概形的上同调》引理
\ref{coherent-lemma-characterize-finite}。

\medskip\noindent
令 $C \subset X$ 为有理尾。
我们断言 $c$ 将 $C$ 映到一点。假设并非如此以导出矛盾，则 $C$ 的像是
$D \subset Y$ 的不可约分支。
回忆 $H^0(C, \mathcal{O}_C) = k'$ 是 $k$ 的有限可分扩张，且 $C$ 有一个
$k'$-有理点 $x$，它也是 $C$ 与 $X$ 的“其余部分”的唯一交点。
由上面的总体讨论可知
$C \setminus \{x\} \subset c^{-1}(Y \setminus \{y_1, \ldots, y_r\})$
经同构映到 $D$ 的一个开子集 $V$。令 $y = c(x) \in D$。
注意 $y$ 是 $D$ 与 $Y$ 的“其余部分”相交的唯一点。若
$y \not \in \{y_1, \ldots, y_r\}$，则 $C \cong D$，显然 $D$ 是
$Y$ 的有理尾，这与 $\omega_Y$ 的丰沛性矛盾
（引理 \ref{curves-lemma-rational-tail-negative}）。
因此 $y \in \{y_1, \ldots, y_r\}$ 且 $\dim(X_y) = 1$。
于是 $x \in X_y \cap C$，并且 $x$ 是 $X_y$ 及 $C$ 的光滑点
（引理 \ref{curves-lemma-closed-subscheme-nodal-curve}）。
若 $y \in D$ 是 $D$ 的奇点，则 $y$ 是结点，且 $Y = D$
（因为由引理 \ref{curves-lemma-closed-subscheme-nodal-curve}，$Y$ 不可能有另一个分支经过 $y$）。
于是 $X = X_y \cup C$，这意味着 $g = 0$，因为按
例 \ref{curves-example-rational-tail} 中的讨论，它等于 $X_y$ 的亏格；矛盾。
若 $y \in D$ 是 $D$ 的光滑点，则
$C \to D$ 是同构（因为非奇异射影模型唯一，且 $C$ 与 $D$ 双有理，见
节 \ref{curves-section-curves-function-fields}）。于是 $D$ 是
$Y$ 的有理尾，这与 $\omega_Y$ 的丰沛性矛盾。

\medskip\noindent
假设 $n \geq 1$。若 $C \subset X$ 是被 $X \to X_1$ 压缩的有理尾，
则由上一段 $C$ 被映到 $Y$ 中的一点。因此 $c : X \to Y$ 经由 $X \to X_1$
分解（因为 $X$ 是 $C$ 与 $X_1$ 的推出，见例 \ref{curves-example-rational-tail} 中的讨论）。
将 $X$ 替换为 $X_1$ 后，$n$ 减少了。由归纳法可假设
$n = 0$，即 $X$ 没有有理尾。

\medskip\noindent
假设 $n = 0$，即 $X$ 没有任何有理尾。
由引理 \ref{curves-lemma-no-rational-tail-semiample-genus-geq-2}，
$\omega_X^{\otimes 2}$ 与 $\omega_X^{\otimes 3}$ 由整体截面生成。
由引理 \ref{curves-lemma-vanishing-twist} 得到
$H^1(X, \omega_X^{\otimes 3}) = 0$。
应用引理 \ref{curves-lemma-contract-gorenstein-canonical}（取 $m = 3$），
可得 $c^*\omega_Y \cong \omega_X$。
由引理 \ref{curves-lemma-rational-bridge-canonical} 及归纳法，还有
$\omega_X = (X \to X_{n'})^*\omega_{X_{n'}}$。
对 $c$ 和 $X \to X_{n'}$ 均应用投影公式，得到
$$
\Gamma(X_{n'}, \omega_{X_{n'}}^{\otimes m}) =
\Gamma(X, \omega_X^{\otimes m}) =
\Gamma(Y, \omega_Y^{\otimes m})
$$
对所有 $m$ 成立。
由于 $X_{n'}$ 与 $Y$ 都是这些空间直和的 Proj，
由《态射》引理 \ref{morphisms-lemma-proper-ample-is-proj}，
得到典范同构 $X_{n'} = Y$，如所需。使图表交换的该同构的唯一性验证略去。
\end{proof}

\begin{lemma}
\label{curves-lemma-tricanonical}
令 $k$ 为域。令 $X$ 为 $k$ 上维数为 $1$ 的真概形，满足
$H^0(X, \mathcal{O}_X) = k$ 且亏格为 $g \geq 2$。假设 $X$ 的奇点
至多为结点且 $\omega_X$ 丰沛。则
$\omega_X^{\otimes 3}$ 非常丰沛，且 $H^1(X, \omega_X^{\otimes 3}) = 0$。
\end{lemma}

\begin{proof}
结合《簇》引理
\ref{varieties-lemma-ampleness-in-terms-of-degrees-components} 及
引理 \ref{curves-lemma-rational-tail-negative} 与 \ref{curves-lemma-rational-bridge-zero}，
可知 $X$ 不含有理尾或有理桥。
于是由引理 \ref{curves-lemma-contracting-rational-tails}，
$\omega_X^{\otimes 3}$ 由整体截面生成。
取 $H^0(X, \omega_X^{\otimes 3})$ 的 $k$-基 $s_0, \ldots, s_n$。
得到态射
$$
\varphi_{\omega_X^{\otimes 3}, (s_0, \ldots, s_n)} :
X \longrightarrow \mathbf{P}^n_k
$$
见《构造》节 \ref{constructions-section-projective-space}。
引理断言该态射是闭浸入。
为检验这一点，可将 $k$ 替换为其代数闭包，见《下降》引理
\ref{descent-lemma-descending-property-closed-immersion}。
因此可假设 $k$ 代数闭。

\medskip\noindent
假设 $k$ 代数闭。
我们将使用《簇》引理
\ref{varieties-lemma-variant-separate-points-tangent-vectors}
证明本引理。
令 $Z \subset X$ 为次数 $2$ 的闭子概形（在 $Z$ 上），
其理想层为 $\mathcal{I} \subset \mathcal{O}_X$。
我们需要证明
$$
H^0(X, \mathcal{L}) \to H^0(Z, \mathcal{L}|_Z)
$$
是满射。因此只需证明
$H^1(X, \mathcal{I}\mathcal{L}) = 0$。
为此使用引理 \ref{curves-lemma-vanishing-on-gorenstein}。
于是只需证明
$$
3\deg(\omega_X|_Y) > -2\chi(Y, \mathcal{O}_Y) + \deg(Z \cap Y)
$$
对每个约化连通闭子概形 $Y \subset X$ 都成立。
由于 $k$ 代数闭且 $Y$ 连通约化，有
$H^0(Y, \mathcal{O}_Y) = k$（《簇》引理
\ref{varieties-lemma-proper-geometrically-reduced-global-sections}).
因此 $\chi(Y, \mathcal{O}_Y) = 1 - \dim H^1(Y, \mathcal{O}_Y)$。
所以需要证明
$$
3\deg(\omega_X|_Y) > -2 + 2\dim H^1(Y, \mathcal{O}_Y) + \deg(Z \cap Y)
$$
由引理 \ref{curves-lemma-no-rational-tail} 可得，除非
$Y = X$ 或 $\deg(\omega_X|_Y) = 0$，上述不等式成立。
由于 $\omega_X$ 丰沛，第二种可能不发生
（见本证明引用的第一个引理）。最后，若 $Y = X$，由 Riemann--Roch（引理 \ref{curves-lemma-rr}）
及 $g \geq 2$ 可知该不等式成立。
对 $Z = \emptyset$ 使用同样论证可得
$H^1(X, \omega_X^{\otimes 3}) = 0$.
\end{proof}







\section{向量场}
\label{curves-section-vector-fields}

\noindent
本节研究曲线上的向量场空间。
向量场对应于无穷小自同构，见《态射的更多性质》节
\ref{more-morphisms-section-action-by-derivations}，
因而在模空间理论中发挥重要作用。

\medskip\noindent
令 $k$ 为代数闭域。
令 $X$ 为 $k$ 上的有限型概形。
令 $x \in X$ 为闭点。
若元素 $D \in \text{Der}_k(\mathcal{O}_X, \mathcal{O}_X)$
{\it 固定 $x$}，是指 $D(\mathcal{I}) \subset \mathcal{I}$，其中
$\mathcal{I} \subset \mathcal{O}_X$ 是 $x$ 的理想层。

\begin{lemma}
\label{curves-lemma-smooth-vector-fields}
令 $k$ 为代数闭域。
令 $X$ 为 $k$ 上光滑、真且连通的曲线。
令 $g$ 为 $X$ 的亏格。
\begin{enumerate}
\item 若 $g \geq 2$，则 $\text{Der}_k(\mathcal{O}_X, \mathcal{O}_X)$
为零；
\item 若 $g = 1$ 且 $D \in \text{Der}_k(\mathcal{O}_X, \mathcal{O}_X)$
非零，则 $D$ 不固定 $X$ 的任何闭点；
\item 若 $g = 0$ 且 $D \in \text{Der}_k(\mathcal{O}_X, \mathcal{O}_X)$
非零，则 $D$ 至多固定 $X$ 的 $2$ 个闭点。
\end{enumerate}
\end{lemma}

\begin{proof}
回忆我们有一个泛 $k$-导数
$d : \mathcal{O}_X \to \Omega_{X/k}$，因而对某个
$\mathcal{O}_X$-线性映射 $\theta : \Omega_{X/k} \to \mathcal{O}_X$ 有 $D = \theta \circ d$。
回忆 $\Omega_{X/k} \cong \omega_X$，见
引理 \ref{curves-lemma-duality-dim-1}。
由 Riemann--Roch 有 $\deg(\omega_X) = 2g - 2$
（引理 \ref{curves-lemma-rr}）。
因此当 $g > 1$ 时，由《簇》引理
\ref{varieties-lemma-check-invertible-sheaf-trivial}.
必有 $\theta$ 为零。这证明了 (1)。
若 $g = 1$，则非零 $\theta$ 处处不消失；若
$g = 0$，则非零 $\theta$ 的零点构成次数为 $2$ 的因子。
因此，只需证明闭点 $x \in X$ 处 $\theta$ 消失
等价于 $D$ 固定 $x$（如上所定义），即可得到 (2)、(3)。
令 $z \in \mathcal{O}_{X, x}$ 为局部参数。
则 $dz$ 是 $\Omega_{X, x}$ 的一个基元素，见
引理 \ref{curves-lemma-uniformizer-works}。
由于 $D(z) = \theta(dz)$，结论成立。
\end{proof}

\begin{lemma}
\label{curves-lemma-nodal-vector-fields}
令 $k$ 为代数闭域。
令 $X$ 为 $k$ 上奇点至多为结点的、真且连通的
维数为 $1$ 的概形。令 $\nu : X^\nu \to X$ 为正规化。
令 $S \subset X^\nu$ 为 $\nu$ 不是同构处的点集。则
$$
\text{Der}_k(\mathcal{O}_X, \mathcal{O}_X) =
\{D' \in \text{Der}_k(\mathcal{O}_{X^\nu}, \mathcal{O}_{X^\nu}) \mid
D' \text{ fixes every }x^\nu \in S\}
$$
\end{lemma}

\begin{proof}
令 $x \in X$ 为结点。令 $x', x'' \in X^\nu$ 为 $x$ 的逆像。
（由于 $k$ 代数闭，每个结点都是分裂结点，见
定义 \ref{curves-definition-split-node} 及引理
\ref{curves-lemma-split-node}。）
令 $u \in \mathcal{O}_{X^\nu, x'}$ 及 $v \in \mathcal{O}_{X^\nu, x''}$
为局部参数。注意我们有正合列
$$
0 \to \mathcal{O}_{X, x} \to
\mathcal{O}_{X^\nu, x'} \times \mathcal{O}_{X^\nu, x''} \to k \to 0
$$
这由引理 \ref{curves-lemma-multicross} 得到。
因此可将 $u$ 与 $v$ 视为 $\mathcal{O}_{X, x}$ 中的元素，且 $uv = 0$。

\medskip\noindent
令 $D \in \text{Der}_k(\mathcal{O}_X, \mathcal{O}_X)$。则
$0 = D(uv) = vD(u) + uD(v)$。由于 $(u)$ 是
$\mathcal{O}_{X, x}$ 中 $v$ 的湮灭子，反之亦然，故 $D(u) \in (u)$ 且
$D(v) \in (v)$。由于 $\mathcal{O}_{X^\nu, x'} = k + (u)$，
可将 $D$ 延拓到 $\mathcal{O}_{X^\nu, x'}$，且该延拓固定 $x'$。
这给出等式右端的某个 $D'$。
反之，给定固定 $x'$ 与 $x''$ 的 $D'$，可知 $D'$ 保持子环
$\mathcal{O}_{X, x} \subset
\mathcal{O}_{X^\nu, x'} \times \mathcal{O}_{X^\nu, x''}$，
这就是等式中从右向左的构造。
\end{proof}

\begin{lemma}
\label{curves-lemma-stable-vector-fields}
令 $k$ 为代数闭域。
令 $X$ 为 $k$ 上奇点至多为结点的、真且连通的
维数为 $1$ 的概形。假设 $X$ 的亏格至少为
$2$，且 $X$ 没有有理尾或有理桥。则
$\text{Der}_k(\mathcal{O}_X, \mathcal{O}_X) = 0$.
\end{lemma}

\begin{proof}
令 $D \in \text{Der}_k(\mathcal{O}_X, \mathcal{O}_X)$。
令 $X^\nu$ 为 $X$ 的正规化。
令 $D' \in \text{Der}_k(\mathcal{O}_{X^\nu}, \mathcal{O}_{X^\nu})$
为经引理 \ref{curves-lemma-nodal-vector-fields} 与 $D$ 对应的元素。
令 $C \subset X^\nu$ 为不可约分支。
若 $C$ 的亏格为 $> 1$，则由引理 \ref{curves-lemma-smooth-vector-fields} 第 (1) 部分，
$D'|_{\mathcal{O}_C} = 0$。
若 $C$ 的亏格为 $1$，则 $C$ 至少有一个闭点 $c$ 映到 $X$ 的结点
（否则 $X \cong C$ 将有亏格 $1$）。由对应关系，这意味着
$D'|_{\mathcal{O}_C}$ 固定 $c$，故由引理
\ref{curves-lemma-smooth-vector-fields} 第 (2) 部分它为零。
最后，若 $C$ 的亏格为零，则至少有
$3$ 个两两不同的闭点 $c_1, c_2, c_3 \in C$
映到 $X$ 的结点；否则会出现以下情形之一：
$X$ 是将 $C$ 的两个点粘合而成（$C$ 的两个点映到同一结点），或
$C$ 是有理桥（两个点映到 $X$ 的不同结点），或
$C$ 是有理尾（一个点映到 $X$ 的结点）。
这三种可能都不允许，因为 $C$ 的亏格 $\geq 2$，且没有有理桥或有理尾。
因此 $D'|_{\mathcal{O}_C}$ 固定 $c_1, c_2, c_3$，故由引理
\ref{curves-lemma-smooth-vector-fields} 第 (3) 部分它为零。
\end{proof}
% Generated mechanically from the frozen resolve.tex target.
% Do not edit this generated file; edit the bound locale source instead.

% OK, start here.
%

\setcounter{chapter}{53}
\chapter{曲面奇点的消解}



\phantomsection
\label{resolve-section-phantom}


\section{引言}
\label{resolve-section-introduction}

\noindent
本章讨论曲面的奇点消解，遵循 Lipman \cite{Lipman}，
并主要依照 Artin 在 \cite{Artin-Lipman} 中的论述。主要结果
（定理 \ref{resolve-theorem-resolve}）说明：Noether 概形
$2$ 维的概形 $Y$ 有奇点消解，只要它有有限正规化
$Y^\nu \to Y$，其中有有限多个奇点 $y_i \in Y^\nu$，并且对每个 $i$，
完备化 $\mathcal{O}_{Y^\nu, y_i}^\wedge$ 是正规的。

\medskip\noindent
确切地说，若 $Y$ 是一个 $2$ 维概形，有限型于拟优基环 $R$
（例如域，或分式域特征为 $0$ 的 Dedekind 域，如 $\mathbf{Z}$），
则 $Y$ 的正规化是有限的，
奇点有限，且局部环的完备化是正规的。见《更多代数》各节
\ref{more-algebra-section-singular-locus}、
\ref{more-algebra-section-G-ring} 以及
\ref{more-algebra-section-excellent} 中的讨论，以及
《更多代数》引理 \ref{more-algebra-lemma-normal-goes-up}。
因此这样的 $Y$ 有奇点消解。

\medskip\noindent
证明的大致轮廓如下。令 $A$ 为维数 $2$ 的 Noether 局部整环。证明步骤为
\begin{enumerate}
\item[N] 将 $A$ 替换为其正规化；
\item[V] 证明 Grauert--Riemenschneider 定理；
\item[B] 证明存在最大值 $g$，它等于 $H^1(X, \mathcal{O}_X)$
在所有正规修改 $X \to \Spec(A)$ 上的长度最大值，
并归约到 $g = 0$ 的情形；
\item[R] 我们称 $A$ 在 $g = 0$ 时定义有理奇点；在此情形经过有限次
爆破后可假设 $A$ 为 Gorenstein 且 $g = 0$；
\item[D] 我们称 $A$ 在 $g = 0$ 且 $A$ 为 Gorenstein 时定义有理双重点，
并在此情形显式地消解奇点。
\end{enumerate}
每一步都需要对环 $A$ 作出假设。下面依次讨论。

\medskip\noindent
关于 N：这里需要假设 $A$ 有有限正规化（这并非自动成立）。
本章大部分地方，在需要某个正规化有限时，我们假设概形是 Nagata 的。
然而，Nagata 条件比定理陈述中给出的条件略强。
解决这个小问题的一种办法是利用环 $A$ 形式非分歧（即其完备化
是约化的），并使用引理 \ref{resolve-lemma-formally-unramified}。
不过，按我们的证明方式，使用引理 \ref{resolve-lemma-normalization-completion}
将完备化上的正规化有限性提升为 $A$ 上正规化的有限性，反而更容易。

\medskip\noindent
关于 V：这是命题 \ref{resolve-proposition-Grauert-Riemenschneider}，
大致说明对正规修改 $f : X \to \Spec(A)$，有
$R^1f_*\omega_X = 0$，其中 $\omega_X$ 是 $X/A$ 的对偶化模
（备注 \ref{resolve-remark-dualizing-setup}）。事实上，由对偶性，该结果等价于关于
$Rf_*\mathcal{O}_X$ 在{\it 导出范畴} $D(A)$ 中的对象的一个陈述
（引理 \ref{resolve-lemma-R1-injective}）。
话虽如此，证明使用了标准事实：特殊纤维的分支具有正的余法层
（引理 \ref{resolve-lemma-nontrivial-normal-bundle}）。

\medskip\noindent
关于 B：在某种意义上，这是证明中最微妙的部分。
最后，当 $A$ 是完备 Noether 局部环时我们才需要使用这一步的输出，
尽管本文写得稍微一般些。术语在定义
\ref{resolve-definition-reduce-to-rational} 中给出。
若上面定义的 $g$ 有界，则直接论证表明可找到正规修改
$X \to \Spec(A)$，使 $X$ 的所有奇点都是有理奇点，见引理
\ref{resolve-lemma-reduce-to-rational}。
我们证明：给定有限扩张 $A \subset B$，若 $g$ 对 $B$ 有界（当它对 $A$ 有界时），
则在下列两种情形中成立：
(1) 分式域扩张可分，见引理 \ref{resolve-lemma-reduce-to-rational}；
(2) 分式域扩张次数为 $p$，特征为 $p$，且 $A$ 正则完备，见
引理 \ref{resolve-lemma-go-up-degree-p}。

\medskip\noindent
关于 R：这里将 $g = 0$ 的情形归约到 Gorenstein 情形。
使一切奏效的奇妙事实是：有理曲面奇点的爆破是正规的，见
引理 \ref{resolve-lemma-blow-up-normal-rational}。

\medskip\noindent
关于 D：有理双重点的消解基本上是手工进行的，见
节 \ref{resolve-section-rational-double-points}。
有理双重点是超曲面奇点（这是真的，但我们不证明，因为不需要）。局部方程形如
$$
a_{11} x_1^2 + a_{12} x_1x_2 + a_{13}x_1x_3 + a_{22} x_2^2 +
a_{23} x_2x_3 + a_{33} x_3^2 =
\sum a_{ijk} x_ix_jx_k
$$
利用二次部分不能为零（因为重数为 $2$，且在任意爆破后仍为 $2$）以及每次爆破
都正规的事实，很快即可得到消解。一个微妙之处是，我们没有在每次爆破下递减的不变量，
但依靠节 \ref{resolve-section-arcs} 中关于形式弧的材料来证明该过程终止。

\medskip\noindent
为将一切拼合起来，还需要一些额外工作。主要难点是要将完备化
$A^\wedge$ 的消解提升为 $\Spec(A)$ 的消解。
为此我们首先证明：若消解存在，则存在由规范化爆破给出的消解
（引理 \ref{resolve-lemma-existence-implies-existence-by-normalized-blowing-ups}）。
由引理 \ref{resolve-lemma-normalized-blowup-completion}，可将一列规范化爆破从完备化提升上来。
即使在完备局部环 $A$ 的消解证明中我们也使用这一点，因为我们的策略对有限包含
$A_0 \subset A$（其中 $A_0$ 正则）的次数作归纳，见引理 \ref{resolve-lemma-resolve-complete}。
若 B 步骤中有更强的结果（例如 Lipman 论文中证明的结果），则可避免这一步。




\section{正特征下的迹映射}
\label{resolve-section-trace}

\noindent
本节的一些结果可以从《德拉姆上同调》第 \ref{derham-section-trace} 节
关于微分形式上迹的更一般讨论中推出。相关讨论见备注
\ref{resolve-remark-compare-Garel}。

\medskip\noindent
固定素数 $p$。令 $R$ 为一个 $\mathbf{F}_p$-代数。给定
$a \in R$，置 $S = R[x]/(x^p - a)$。定义 $R$-线性映射
$$
\text{Tr}_x : \Omega_{S/R} \longrightarrow \Omega_R
$$
其规则为
$$
x^i\text{d}x \longmapsto
\left\{
\begin{matrix}
0 & \text{若} & 0 \leq i \leq p - 2, \\
\text{d}a & \text{若} & i = p - 1
\end{matrix}
\right.
$$
这是有意义的，因为 $\Omega_{S/R}$ 是自由 $R$-模，其基为
$x^i\text{d}x$（$0 \leq i \leq p - 1$）。
下面的引理说明迹映射是良定义的，即与坐标 $x$ 的选择无关。

\begin{lemma}
\label{resolve-lemma-trace-well-defined}
令 $\varphi : R[x]/(x^p - a) \to R[y]/(y^p - b)$ 为一个 $R$-代数
同态。则 $\text{Tr}_x = \text{Tr}_y \circ \varphi$。
\end{lemma}

\begin{proof}
设 $\varphi(x) = \lambda_0 + \lambda_1 y + \ldots + \lambda_{p - 1}y^{p - 1}$，
其中 $\lambda_i \in R$。将 $x$ 映到
$\lambda_0 + \lambda_1 y + \ldots + \lambda_{p - 1}y^{p - 1}$
能诱导 $R$-代数同态 $R[x]/(x^p - a) \to R[y]/(y^p - b)$
的条件等价于
$$
a = \lambda_0^p + \lambda_1^p b + \ldots + \lambda_{p - 1}^pb^{p - 1}
$$
在环 $R$ 中成立。考虑多项式环
$$
R_{univ} = \mathbf{F}_p[b, \lambda_0, \ldots, \lambda_{p - 1}]
$$
以及其中的元素
$a = \lambda_0^p + \lambda_1^p b + \ldots + \lambda_{p - 1}^pb^{p - 1}$
考虑泛代数映射
$\varphi_{univ} : R_{univ}[x]/(x^p - a) \to R_{univ}[y]/(y^p - b)$
它将 $x$ 映到
$\lambda_0 + \lambda_1 y + \ldots + \lambda_{p - 1}y^{p - 1}$.
得到典范映射
$$
R_{univ} \longrightarrow R
$$
将 $b, \lambda_i$ 分别送到 $b, \lambda_i$。按构造得到交换图
$$
\xymatrix{
R_{univ}[x]/(x^p - a) \ar[r] \ar[d]_{\varphi_{univ}} &
R[x]/(x^p - a) \ar[d]^\varphi \\
R_{univ}[y]/(y^p - b) \ar[r] & R[y]/(y^p - b)
}
$$
且水平箭头与迹映射相容。因此只需对映射 $\varphi_{univ}$ 证明引理。
于是可设 $R = \mathbf{F}_p[b, \lambda_0, \ldots, \lambda_{p - 1}]$
为多项式环。我们通过计算
$\text{Tr}_y(\varphi(x)^i\text{d}\varphi(x))$（$i = 0 , \ldots, p - 1$）
来验证此时引理成立。

\medskip\noindent
情形 $0 \leq i \leq p - 2$。展开
$$
(\lambda_0 + \lambda_1 y + \ldots + \lambda_{p - 1}y^{p - 1})^i
(\lambda_1 + 2 \lambda_2 y + \ldots + (p - 1)\lambda_{p - 1}y^{p - 2})
$$
在环 $R[y]/(y^p - b)$ 中展开上式。我们必须证明 $y^{p - 1}$ 的系数为零。
为此，只需证明上式作为 $y$ 的多项式时，在幂 $y^{pk - 1}$
前的系数都为零。将该多项式写成
$$
\frac{\text{d}}{(i + 1)\text{d}y}
(\lambda_0 + \lambda_1 y + \ldots + \lambda_{p - 1}y^{p - 1})^{i + 1}
$$
确实，$y^{kp - 1}$ 的系数全部为零。

\medskip\noindent
情形 $i = p - 1$。展开
$$
(\lambda_0 + \lambda_1 y + \ldots + \lambda_{p - 1}y^{p - 1})^{p - 1}
(\lambda_1 + 2 \lambda_2 y + \ldots + (p - 1)\lambda_{p - 1}y^{p - 2})
$$
在环 $R[y]/(y^p - b)$ 中展开上式。为了完成证明，需要说明 $y^{p - 1}$ 的系数乘以
$\text{d}b$ 等于 $\text{d}a$。这里使用 $R$ 是 $S/pS$，其中
$S = \mathbf{Z}[b, \lambda_0, \ldots, \lambda_{p - 1}]$.
于是上式作为 $y$ 的多项式等于
$$
\frac{\text{d}}{p\text{d}y}
(\lambda_0 + \lambda_1 y + \ldots + \lambda_{p - 1}y^{p - 1})^p
$$
由于 $\frac{\text{d}}{\text{d}y}(y^{pk}) = pk y^{pk - 1}$，
只需考察 $y^{pk}$ 在多项式
$(\lambda_0 + \lambda_1 y + \ldots + \lambda_{p - 1}y^{p - 1})^p$
中的系数模 $p$ 的值。这些项之和为
$$
\lambda_0^p + \lambda_1^py^p + \ldots + \lambda_{p - 1}^py^{p(p - 1)}
\bmod p
$$
因此，施加算子 $\frac{\text{d}}{p\text{d}y}$，并模 $y^p - b$ 化简后，
我们得到
$$
\lambda_1^p + 2\lambda_2^pb + \ldots + (p - 1)\lambda_{p - 1}b^{p - 2}
$$
这就是所要计算的 $y^{p - 1}$ 的系数。由于
$a = \lambda_0^p + \lambda_1^p b + \ldots + \lambda_{p - 1}^pb^{p - 1}$
在 $R$ 中成立，可得
$$
\text{d}a = (\lambda_1^p  + 2 \lambda_2^p b + \ldots +
(p - 1) \lambda_{p - 1}^p b^{p - 2}) \text{d}b
$$
在 $R$ 中成立。这证明 $y^{p - 1}$ 的系数正如所需。
\end{proof}

\begin{lemma}
\label{resolve-lemma-trace-higher}
令 $\mathbf{F}_p \subset \Lambda \subset R \subset S$ 为环扩张，
并假设 $S$ 同构于 $R[x]/(x^p - a)$，其中 $a \in R$。
则存在典范的 $R$-线性映射
$$
\text{Tr} :
\Omega^{t + 1}_{S/\Lambda}
\longrightarrow
\Omega_{R/\Lambda}^{t + 1}
$$
其中 $t \geq 0$，且满足
$$
\eta_1 \wedge \ldots \wedge \eta_t \wedge x^i\text{d}x
\longmapsto
\left\{
\begin{matrix}
0 & \text{若} & 0 \leq i \leq p - 2, \\
\eta_1 \wedge \ldots \wedge \eta_t \wedge \text{d}a & \text{若} & i = p - 1
\end{matrix}
\right.
$$
其中 $\eta_i \in \Omega_{R/\Lambda}$，并且 $\text{Tr}$ 湮灭映射
$S \otimes_R \Omega_{R/\Lambda}^{t + 1} \to \Omega_{S/\Lambda}^{t + 1}$.
\end{lemma}

\begin{proof}
当 $t = 0$ 时，使用复合映射
$$
\Omega_{S/\Lambda} \to \Omega_{S/R} \to \Omega_R \to \Omega_{R/\Lambda}
$$
其中第二个映射由引理 \ref{resolve-lemma-trace-well-defined} 给出。
存在正合序列
$$
H_1(L_{S/R}) \xrightarrow{\delta} \Omega_{R/\Lambda} \otimes_R S \to
\Omega_{S/\Lambda} \to \Omega_{S/R} \to 0
$$
（代数，引理 \ref{algebra-lemma-exact-sequence-NL}）。
模 $\Omega_{S/R}$ 在 $S$ 上以 $\text{d}x$ 为基自由，
模 $H_1(L_{S/R})$ 在 $S$ 上以 $x^p - a$ 为基自由，后者经 $\delta$
映到 $-\text{d}a \otimes 1$（在 $\Omega_{R/\Lambda} \otimes_R S$ 中）。
特别地，若置
$$
M = \Coker(R \to \Omega_{R/\Lambda}, 1 \mapsto -\text{d}a)
$$
则可见 $\Coker(\delta) = M \otimes_R S$。于是得到典范映射
$$
\Omega^{t + 1}_{S/\Lambda} \to
\wedge_S^t(\Coker(\delta)) \otimes_S \Omega_{S/R} =
\wedge^t_R(M) \otimes_R \Omega_{S/R}
$$
现在，由于映射
$\text{Tr} : \Omega_{S/R} \to \Omega_{R/\Lambda}$
的像包含于 $R\text{d}a$（引理 \ref{resolve-lemma-trace-well-defined}），
可知与像中的元素作楔积会湮灭 $\text{d}a$。
因此存在典范映射
$$
\wedge^t_R(M) \otimes_R \Omega_{S/R} \to \Omega_{R/\Lambda}^{t + 1}
$$
它将
$\overline{\eta}_1 \wedge \ldots \wedge \overline{\eta}_t \wedge \omega$
映到 $\eta_1 \wedge \ldots \wedge \eta_t \wedge \text{Tr}(\omega)$。
\end{proof}

\begin{remark}
\label{resolve-remark-compare-Garel}
令 $\mathbf{F}_p \subset \Lambda \subset R \subset S$ 及 $\text{Tr}$
如引理 \ref{resolve-lemma-trace-higher} 中所设。由《德拉姆上同调》命题
\ref{derham-proposition-Garel}，存在复形的典范映射
$$
\Theta_{S/R} :
\Omega_{S/\Lambda}^\bullet
\longrightarrow
\Omega_{R/\Lambda}^\bullet
$$
《德拉姆上同调》例 \ref{derham-example-Garel} 中的计算表明
$\Theta_{S/R}(x^i \text{d}x) = \text{Tr}_x(x^i\text{d}x)$
对所有 $i$ 都成立。
由于 $\text{Trace}_{S/R} = \Theta^0_{S/R}$ 恒等于零，并且
$$
\Theta_{S/R}(a \wedge b) = a \wedge \Theta_{S/R}(b)
$$
对 $a \in \Omega^i_{R/\Lambda}$ 及 $b \in \Omega^j_{S/\Lambda}$，
由此得到 $\text{Tr} = \Theta_{S/R}$。使用 $\text{Tr}$ 的优点在于，
它的构造要初等得多。
\end{remark}

\begin{lemma}
\label{resolve-lemma-trace-extends}
令 $S$ 为 $\mathbf{F}_p$ 上的概形。令 $f : Y \to X$ 为 $S$ 上 Noether
正规整概形之间的有限态射。假设
\begin{enumerate}
\item 函数域扩张是次数为 $p$ 的纯不可分扩张；
\item $\Omega_{X/S}$ 是相干的 $\mathcal{O}_X$-模（例如 $X$ 在 $S$ 上有限型）。
\end{enumerate}
则对 $i \geq 1$ 存在典范映射
$$
\text{Tr} : f_*\Omega^i_{Y/S} \longrightarrow (\Omega_{X/S}^i)^{**}
$$
其在 $X$ 的一般点处的茎恢复引理 \ref{resolve-lemma-trace-higher} 的迹映射。
\end{lemma}

\begin{proof}
正合序列 $f^*\Omega_{X/S} \to \Omega_{Y/S} \to \Omega_{Y/X} \to 0$
表明 $\Omega_{Y/S}$，从而 $f_*\Omega_{Y/S}$，也是相干模。因此只需证明：
一般点处的迹映射可延拓到满足
$x \in X$ 且 $\dim(\mathcal{O}_{X, x}) = 1$ 的茎；见《除子》引理
\ref{divisors-lemma-describe-reflexive-hull}。于是归结为下一段讨论的情形。

\medskip\noindent
设 $X = \Spec(A)$、$Y = \Spec(B)$，其中 $A$ 是离散赋值环，
且 $B$ 在 $A$ 上有限。由于诱导的分式域扩张 $L/K$ 是纯不可分的，
可知 $B$ 也是局部环。因此 $B$ 也是离散赋值环。于是有以下两种情形之一：
\begin{enumerate}
\item $B/A$ 的分歧指数为 $p$，因而 $B = A[x]/(x^p - a)$，
其中 $a \in A$ 是素元；
\item $\mathfrak m_B = \mathfrak m_A B$，且剩余域
$B/\mathfrak m_A B$ 是次数为 $p$ 的纯不可分扩张，扩张基域为
$\kappa_A = A/\mathfrak m_A$。
取任意 $x \in B$，使其剩余类不属于 $\kappa_A$，则
$B = A[x]/(x^p - a)$，其中 $a \in A$ 为单位。
\end{enumerate}
令 $\Spec(\Lambda) \subset S$ 为仿射开集，使得 $X$ 映入
$\Spec(\Lambda)$。应用引理 \ref{resolve-lemma-trace-higher} 可知，
迹映射延拓为
$\Omega^i_{B/\Lambda} \to \Omega^i_{A/\Lambda}$，对所有 $i \geq 1$。
\end{proof}


















\section{二次变换}
\label{resolve-section-quadratic}

\noindent
本节研究曲面上爆破非奇异点时发生的情形。我们不急于将这种态射
形式化地称为{\it 二次变换}：一方面常常使用其他名称，另一方面
“二次变换”这一短语有时具有不同的含义。

\begin{lemma}
\label{resolve-lemma-blowup}
令 $(A, \mathfrak m, \kappa)$ 为维数 $2$ 的正则局部环。
令 $f : X \to S = \Spec(A)$ 为沿 $A$ 中的 $\mathfrak m$ 爆破的态射，
带有例外除子 $E$。存在闭浸入
$$
r : X \longrightarrow \mathbf{P}^1_S
$$
在 $S$ 上满足
\begin{enumerate}
\item $r|_E : E \to \mathbf{P}^1_\kappa$ 是同构；
\item $\mathcal{O}_X(E) = \mathcal{O}_X(-1) =
r^*\mathcal{O}_{\mathbf{P}^1}(-1)$；
\item $\mathcal{C}_{E/X} = (r|_E)^*\mathcal{O}_{\mathbf{P}^1}(1)$，且
$\mathcal{N}_{E/X} = (r|_E)^*\mathcal{O}_{\mathbf{P}^1}(-1)$。
\end{enumerate}
\end{lemma}

\begin{proof}
由于 $A$ 是维数 $2$ 的正则环，可写 $\mathfrak m = (x, y)$。
将 $x$、$y$ 置于次数 $1$，它们生成 Rees 代数
$\bigoplus_{n \geq 0} \mathfrak m^n$，这是 $A$ 上的代数。回忆
$X = \text{Proj}(\bigoplus_{n \geq 0} \mathfrak m^n)$；见《除子》引理
\ref{divisors-lemma-blowing-up-affine}。于是满射
$$
A[T_0, T_1] \longrightarrow \bigoplus\nolimits_{n \geq 0} \mathfrak m^n,
\quad
T_0 \mapsto x,\ T_1 \mapsto y
$$
作为分次 $A$-代数诱导闭浸入
$r : X \to \mathbf{P}^1_S = \text{Proj}(A[T_0, T_1])$，
满足 $\mathcal{O}_X(1) = r^*\mathcal{O}_{\mathbf{P}^1_S}(1)$；见《构造》引理
\ref{constructions-lemma-surjective-graded-rings-generated-degree-1-map-proj}。
这证明 (2)，因为由《除子》引理
\ref{divisors-lemma-blowing-up-gives-effective-Cartier-divisor}，
$\mathcal{O}_X(E) = \mathcal{O}_X(-1)$。

\medskip\noindent
为证明 (1)，注意
$$
\left(\bigoplus\nolimits_{n \geq 0} \mathfrak m^n\right) \otimes_A \kappa =
\bigoplus\nolimits_{n \geq 0} \mathfrak m^n/\mathfrak m^{n + 1} \cong
\kappa[\overline{x}, \overline{y}]
$$
这是多项式代数；见《代数》引理 \ref{algebra-lemma-regular-graded}。
这证明 $X \to S$ 在 $\Spec(\kappa)$ 上的纤维等于
$\text{Proj}(\kappa[\overline{x}, \overline{y}]) = \mathbf{P}^1_\kappa$；见《构造》引理
\ref{constructions-lemma-base-change-map-proj}。
回忆 $E$ 是 $X$ 的闭子概形，由 $\mathfrak m\mathcal{O}_X$ 定义，
即 $E = X_\kappa$。由态射 $r$ 的选取可知，$r|_E$ 实际上
将 $E = X_\kappa$ 识别为 $\mathbf{P}^1_S \to S$ 的特殊纤维。

\medskip\noindent
(3) 由 (1)、(2) 以及《除子》引理
\ref{divisors-lemma-conormal-effective-Cartier-divisor} 得出。
\end{proof}

\begin{lemma}
\label{resolve-lemma-blowup-regular}
令 $(A, \mathfrak m, \kappa)$ 为维数 $2$ 的正则局部环。
令 $f : X \to S = \Spec(A)$ 为沿 $A$ 中的 $\mathfrak m$ 爆破的态射。
则 $X$ 是不可约正则概形。
\end{lemma}

\begin{proof}
由《除子》引理 \ref{divisors-lemma-blow-up-integral-scheme}
以及《代数》引理 \ref{algebra-lemma-regular-domain}，可知 $X$ 是整概形。
要看出 $X$ 正则，只需检查闭点处的
$\mathcal{O}_{X, x}$（其中 $x \in X$）正则；见《性质》引理
\ref{properties-lemma-characterize-regular}。
令 $x \in X$ 为闭点。由于 $f$ proper，$x$ 映到
$\mathfrak m$，即 $x$ 是例外除子 $E$ 上的点。
于是 $E$ 是有效 Cartier 除子，且 $E \cong \mathbf{P}^1_\kappa$。
因此，若 $g \in \mathfrak m_x \subset \mathcal{O}_{X, x}$ 是 $E$ 的局部
方程，则
$\mathcal{O}_{X, x}/(g) \cong \mathcal{O}_{\mathbf{P}^1_\kappa, x}$。
由于 $\mathbf{P}^1_\kappa$ 可由两个仿射开集覆盖，而它们是
在 $\kappa$ 上的多项式环的谱，可知
$\mathcal{O}_{\mathbf{P}^1_\kappa, x}$ 正则；见《代数》引理
\ref{algebra-lemma-dim-affine-space}。
最后应用《代数》引理
\ref{algebra-lemma-regular-mod-x}。
\end{proof}

\begin{lemma}
\label{resolve-lemma-blowup-pic}
令 $(A, \mathfrak m, \kappa)$ 为维数 $2$ 的正则局部环。
令 $f : X \to S = \Spec(A)$ 为沿 $A$ 中的 $\mathfrak m$ 爆破的态射。
则 $\Pic(X) = \mathbf{Z}$，由 $\mathcal{O}_X(E)$ 生成。
\end{lemma}

\begin{proof}
回忆 $E = \mathbf{P}^1_\kappa$ 的 Picard 群为 $\mathbf{Z}$，
生成元为 $\mathcal{O}(1)$；见《除子》引理
\ref{divisors-lemma-Pic-projective-space-UFD}。
由引理 \ref{resolve-lemma-blowup}，可逆 $\mathcal{O}_X$-模
$\mathcal{O}_X(E)$ 限制为 $\mathcal{O}(-1)$。因此
$\mathcal{O}_X(E)$ 在 $\Pic(X)$ 中生成无限循环群。
由于 $A$ 正则，它是 UFD；见《更多代数》引理
\ref{more-algebra-lemma-regular-local-UFD}。
于是去掉闭点的谱 $U = S \setminus \{\mathfrak m\} = X \setminus E$
的 Picard 群平凡；见《除子》引理
\ref{divisors-lemma-open-subscheme-UFD}。
因此，对每个可逆 $\mathcal{O}_X$-模 $\mathcal{L}$，
存在同构 $s : \mathcal{O}_U \to \mathcal{L}|_U$。
于是 $s$ 是 $\mathcal{L}$ 的正则亚纯截面，且
$\text{div}_\mathcal{L}(s) = nE$，其中
$n \in \mathbf{Z}$
（《除子》定义 \ref{divisors-definition-divisor-invertible-sheaf}）。
由《除子》引理 \ref{divisors-lemma-normal-c1-injective}
（以及由引理 \ref{resolve-lemma-blowup-regular} 知 $X$ 正规），
得到 $\mathcal{L} = \mathcal{O}_X(nE)$。
\end{proof}

\begin{lemma}
\label{resolve-lemma-cohomology-of-blowup}
令 $(A, \mathfrak m, \kappa)$ 为维数 $2$ 的正则局部环。
令 $f : X \to S = \Spec(A)$ 为沿 $A$ 中的 $\mathfrak m$ 爆破的态射。
令 $\mathcal{F}$ 为拟相干 $\mathcal{O}_X$-模。
\begin{enumerate}
\item $H^p(X, \mathcal{F}) = 0$，当 $p \not \in \{0, 1\}$ 时；
\item $H^1(X, \mathcal{O}_X(n)) = 0$，当 $n \geq -1$ 时；
\item $H^1(X, \mathcal{F}) = 0$，只要 $\mathcal{F}$ 或 $\mathcal{F}(1)$
是整体生成的；
\item $H^0(X, \mathcal{O}_X(n)) = \mathfrak m^{\max(0, n)}$；
\item $\text{length}_A H^1(X, \mathcal{O}_X(n)) = -n(-n - 1)/2$，
当 $n < 0$ 时。
\end{enumerate}
\end{lemma}

\begin{proof}
若 $\mathfrak m = (x, y)$，则 $X$ 被仿射爆破代数的谱覆盖，
这些代数为 $A[\frac{\mathfrak m}{x}]$ 和
$A[\frac{\mathfrak m}{y}]$；这是因为将 $x$、$y$ 置于次数 $1$ 后，
它们生成 Rees 代数 $\bigoplus \mathfrak m^n$（这是 $A$ 上的代数）。
见《除子》引理 \ref{divisors-lemma-blowing-up-affine} 以及
《构造》引理 \ref{constructions-lemma-proj-quasi-compact}。
由《构造》引理 \ref{constructions-lemma-proj-separated}，$X$ 是分离的，
故由《概形上同调》引理 \ref{coherent-lemma-vanishing-nr-affines}，
拟相干层的上同调在次数 $\geq 2$ 时消失。

\medskip\noindent
令 $i : E \to X$ 为例外除子，见《除子》定义
\ref{divisors-definition-blow-up}。
回忆 $\mathcal{O}_X(-E) = \mathcal{O}_X(1)$ 是
$f$-相对充足的；见《除子》引理
\ref{divisors-lemma-blowing-up-gives-effective-Cartier-divisor}。
因此有 $H^1(X, \mathcal{O}_X(-nE)) = 0$，其中 $n > 0$；
见《概形上同调》引理 \ref{coherent-lemma-kill-by-twisting}。
考虑滤链
$$
\mathcal{O}_X(-nE) \subset \mathcal{O}_X(-(n - 1)E) \subset
\ldots \subset \mathcal{O}_X(-E) \subset \mathcal{O}_X \subset \mathcal{O}_X(E)
$$
其相邻商层为
$$
\mathcal{O}_X(-t E)/\mathcal{O}_X(-(t + 1)E) =
\mathcal{O}_X(t)/\mathcal{I}(t) =
i_*\mathcal{O}_E(t)
$$
其中 $\mathcal{I} = \mathcal{O}_X(-E)$ 是 $E$ 的理想层。
由引理 \ref{resolve-lemma-blowup}，有 $E = \mathbf{P}^1_\kappa$，且
$\mathcal{O}_E(1)$ 确实对应于 $\mathbf{P}^1$ 上结构层的通常 Serre 扭曲。
因此 $\mathcal{O}_E(t)$ 的上同调在次数 $1$ 上（当 $t \geq -1$ 时）消失；见
《概形上同调》引理 \ref{coherent-lemma-cohomology-projective-space-over-ring}。
由于这等于 $H^1(X, i_*\mathcal{O}_E(t))$（由《概形上同调》引理
\ref{coherent-lemma-relative-affine-cohomology}），可知
$H^1(X, \mathcal{O}_X(-(t + 1)E)) \to H^1(X, \mathcal{O}_X(-tE))$
在 $t \geq -1$ 时满射。因此
$$
0 = H^1(X, \mathcal{O}_X(-nE))
\longrightarrow
H^1(X, \mathcal{O}_X(-tE)) = H^1(X, \mathcal{O}_X(t))
$$
在 $t \geq -1$ 时满射，这就证明了 (2)。

\medskip\noindent
令 $\mathcal{F}$ 为整体生成的。这意味着存在短正合序列
$$
0 \to \mathcal{G} \to \bigoplus\nolimits_{i \in I} \mathcal{O}_X
\to \mathcal{F} \to 0
$$
注意由《上同调》引理
\ref{cohomology-lemma-quasi-separated-cohomology-colimit}，有
$H^1(X, \bigoplus_{i \in I} \mathcal{O}_X) =
\bigoplus_{i \in I} H^1(X, \mathcal{O}_X)$。
由 (2) 得 $H^1(X, \mathcal{O}_X) = 0$。
若 $\mathcal{F}(1)$ 整体生成，则可找到满射
$\bigoplus_{i \in I} \mathcal{O}_X(-1) \to \mathcal{F}$，
并以类似方式论证。换言之，(3) 由 (2) 推出。

\medskip\noindent
对于 (4)，注意当 $n$ 足够大时，有
$\Gamma(X, \mathcal{O}_X(n)) = \mathfrak m^n$；见《概形上同调》引理
\ref{coherent-lemma-recover-tail-graded-module}。
若 $n \geq 0$，则可使用短正合序列
$$
0 \to \mathcal{O}_X(n) \to \mathcal{O}_X(n - 1) \to
i_*\mathcal{O}_E(n - 1) \to 0
$$
并利用左侧层的 $H^1$ 消失得到交换图
$$
\xymatrix{
0 \ar[r] &
\mathfrak m^{\max(0, n)} \ar[r] \ar[d] &
\mathfrak m^{\max(0, n - 1)} \ar[r] \ar[d] &
\mathfrak m^{\max(0, n)}/\mathfrak m^{\max(0, n - 1)} \ar[r] \ar[d] & 0\\
0 \ar[r] &
\Gamma(X, \mathcal{O}_X(n)) \ar[r] &
\Gamma(X, \mathcal{O}_X(n - 1)) \ar[r] &
\Gamma(E, \mathcal{O}_E(n - 1)) \ar[r] & 0
}
$$
其行均正合。事实上，当 $n < 0$ 时行也正合，
因为此时右侧的群为零。
在引理 \ref{resolve-lemma-blowup} 的证明中我们已看到右侧竖直箭头是同构
（细节略去）。因此若左侧竖直箭头是同构，中间的也是同构。
由此对 $n$ 作递降归纳即可得到 (4)。

\medskip\noindent
最后，我们利用下列序列对 $n$ 作递降归纳来证明 (5)：
$$
0 \to \mathcal{O}_X(n) \to \mathcal{O}_X(n - 1) \to
i_*\mathcal{O}_E(n - 1) \to 0
$$
具体地，当 $n \geq -1$ 时，我们已经知道
$H^1(X, \mathcal{O}_X(n)) = 0$。由于
$$
H^1(X, i_*\mathcal{O}_E(-2)) =
H^1(E, \mathcal{O}_E(-2)) =
H^1(\mathbf{P}^1_\kappa, \mathcal{O}(-2)) \cong \kappa
$$
由《概形上同调》引理
\ref{coherent-lemma-cohomology-projective-space-over-ring}，上述模长度为
$1$ 的 $A$-模，故由上同调长正合序列可知 (5) 对 $n = -2$ 成立。
如此继续即可。
\end{proof}

\begin{lemma}
\label{resolve-lemma-blowup-improve}
令 $(A, \mathfrak m)$ 为维数 $2$ 的正则局部环。
令 $f : X \to S = \Spec(A)$ 为沿 $A$ 中的 $\mathfrak m$ 爆破的态射。
令 $\mathfrak m^n \subset I \subset \mathfrak m$ 为理想。
令 $d \geq 0$ 为满足下式的最大整数：
$$
I \mathcal{O}_X \subset \mathcal{O}_X(-dE)
$$
其中 $E$ 为例外除子。置
$\mathcal{I}' = I\mathcal{O}_X(dE) \subset \mathcal{O}_X$.
则 $d > 0$，层 $\mathcal{O}_X/\mathcal{I}'$ 支持在有限多个
闭点 $x_1, \ldots, x_r$（它们属于 $X$）上，并且
\begin{align*}
\text{length}_A(A/I)
& >
\text{length}_A \Gamma(X, \mathcal{O}_X/\mathcal{I}') \\
& \geq
\sum\nolimits_{i = 1, \ldots, r}
\text{length}_{\mathcal{O}_{X, x_i}}
(\mathcal{O}_{X, x_i}/\mathcal{I}'_{x_i})
\end{align*}
\end{lemma}

\begin{proof}
由于 $I \subset \mathfrak m$，$I$ 的每个元素都在 $E$ 上消失，
故 $d \geq 1$。另一方面，由于 $\mathfrak m^n \subset I$，可知
$d \leq n$。考虑短正合序列
$$
0 \to I\mathcal{O}_X \to \mathcal{O}_X \to \mathcal{O}_X/I\mathcal{O}_X \to 0
$$
由于 $I\mathcal{O}_X$ 整体生成，由引理
\ref{resolve-lemma-cohomology-of-blowup} 有 $H^1(X, I\mathcal{O}_X) = 0$。
因此得到满射
$A/I \to \Gamma(X, \mathcal{O}_X/I\mathcal{O}_X)$。考虑短正合序列
$$
0 \to
\mathcal{O}_X(-dE)/I\mathcal{O}_X \to
\mathcal{O}_X/I\mathcal{O}_X \to
\mathcal{O}_X/\mathcal{O}_X(-dE) \to 0
$$
由《除子》引理 \ref{divisors-lemma-codim-1-part}，可知
$\mathcal{O}_X(-dE)/I\mathcal{O}_X$ 支持在 $X$ 的有限多个闭点上。
特别地，该相干层的高阶上同调群消失（细节略去）。于是下图中
$$
\xymatrix{
& & A/I \ar[d] \\
0 \ar[r] &
\Gamma(X, \mathcal{O}_X(-dE)/I\mathcal{O}_X) \ar[r] &
\Gamma(X, \mathcal{O}_X/I\mathcal{O}_X) \ar[r] &
\Gamma(X, \mathcal{O}_X/\mathcal{O}_X(-dE)) \ar[r] & 0
}
$$
底行正合且竖直箭头满射。我们有
$$
\text{length}_A \Gamma(X, \mathcal{O}_X(-dE)/I\mathcal{O}_X) <
\text{length}_A(A/I)
$$
因为 $\Gamma(X, \mathcal{O}_X/\mathcal{O}_X(-dE))$ 非零。
事实上，$1 \in \Gamma(X, \mathcal{O}_X)$ 的像在 $d > 0$ 时非零。

\medskip\noindent
为完成证明，将上述结果改写为引理中的陈述。由于
$\mathcal{O}_X(dE)$ 可逆，有
$$
\mathcal{O}_X/\mathcal{I}' =
\mathcal{O}_X(-dE)/I\mathcal{O}_X \otimes_{\mathcal{O}_X} \mathcal{O}_X(dE).
$$
因此 $\mathcal{O}_X/\mathcal{I}'$ 与 $\mathcal{O}_X(-dE)/I\mathcal{O}_X$
支持在同一有限点集上，记为 $x_1, \ldots, x_r \in E \subset X$。
此外有
$$
\Gamma(X, \mathcal{O}_X(-dE)/I\mathcal{O}_X) =
\bigoplus \mathcal{O}_X(-dE)_{x_i}/I\mathcal{O}_{X, x_i}
\cong
\bigoplus \mathcal{O}_{X, x_i}/\mathcal{I}'_{x_i} =
\Gamma(X, \mathcal{O}_X/\mathcal{I}')
$$
因为局部环上的可逆模是平凡的。由此得到严格不等式。
又由于
$$
\text{length}_A(\mathcal{O}_{X, x_i}/\mathcal{I}'_{x_i}) \geq
\text{length}_{\mathcal{O}_{X, x_i}}(\mathcal{O}_{X, x_i}/\mathcal{I}'_{x_i})
$$
这直接由长度的定义得到。
\end{proof}

\begin{lemma}
\label{resolve-lemma-differentials-of-blowup}
令 $(A, \mathfrak m, \kappa)$ 为维数 $2$ 的正则局部环。
令 $f : X \to S = \Spec(A)$ 为沿 $A$ 中的 $\mathfrak m$ 爆破的态射。
则 $\Omega_{X/S} = i_*\Omega_{E/\kappa}$，其中 $i : E \to X$
是例外除子的浸入。
\end{lemma}

\begin{proof}
记 $\mathbf{P}^1 = \mathbf{P}^1_S$，令
$r : X \to \mathbf{P}^1$ 如引理 \ref{resolve-lemma-blowup} 中所述。
则有正合序列
$$
\mathcal{C}_{X/\mathbf{P}^1} \to r^*\Omega_{\mathbf{P}^1/S} \to
\Omega_{X/S} \to 0
$$
见《态射》引理 \ref{morphisms-lemma-differentials-relative-immersion}。
由《态射》引理 \ref{morphisms-lemma-base-change-differentials}，
$\Omega_{\mathbf{P}^1/S}|_E = \Omega_{E/\kappa}$，因此只需证明第一箭头
映到典范映射
$r^*\Omega_{\mathbf{P}^1/S} \to i_*\Omega_{E/\kappa}$
的核上是满射。这可局部完成。沿用引理 \ref{resolve-lemma-blowup} 证明中的记号，
在 $X$ 的仿射开集上，态射 $f$ 对应于环映射
$$
A \to A[t]/(xt - y)
$$
其中 $x, y \in \mathfrak m$ 为生成元。因此
$\text{d}(xt - y) = x\text{d}t$ 且 $y\text{d}t = t \cdot x \text{d}t$，
这就证明了所需结论。
\end{proof}



\section{由二次变换支配}
\label{resolve-section-dominating-by-quadratic}

\noindent
利用上述结果，可以证明正则 $2$ 维概形的任意修改都由点上的爆破支配。

\medskip\noindent
令 $X$ 为概形，令 $x \in X$ 为闭点。照例，将
$i : x = \Spec(\kappa(x)) \to X$ 看作闭子概形。
称 {\it $X' \to X$} 为沿 $X$ 中的点 $x$ 的 $X$ 的爆破，
即在闭子概形 $x \subset X$ 上进行爆破。注意，若 $X$ 局部 Noether，则由《除子》引理
\ref{divisors-lemma-blowing-up-projective}，$X' \to X$ 是射影的（特别是 proper）。

\begin{lemma}
\label{resolve-lemma-make-ideal-principal}
令 $X$ 为 Noether 概形。令 $T \subset X$ 为有限个闭点 $x$ 的集合，
使得 $\mathcal{O}_{X, x}$ 是维数 $2$ 的正则环（对 $x \in T$）。
令 $\mathcal{I} \subset \mathcal{O}_X$ 为拟相干理想层，且
$\mathcal{O}_X/\mathcal{I}$ 支持在 $T$ 上。
则存在序列
$$
X_n \to X_{n - 1} \to \ldots \to X_1 \to X_0 = X
$$
其中 $X_{i + 1} \to X_i$ 是在 $X_i$ 上沿位于 $T$ 中某点上方的闭点进行的爆破，
并且 $\mathcal{I}\mathcal{O}_{X_n}$ 是可逆理想层。
\end{lemma}

\begin{proof}
设 $T = \{x_1, \ldots, x_r\}$. 记 $I_i$ 为 $\mathcal{I}$ 在
$x_i$ 处的茎。置
$$
n_i = \text{length}_{\mathcal{O}_{X, x_i}}(\mathcal{O}_{X, x_i}/I_i)
$$
由于 $\mathcal{O}_X/\mathcal{I}$ 支持在 $T$ 上，该长度有限，
从而 $\mathcal{O}_{X, x_i}/I_i$ 的支集等于
$\{\mathfrak m_{x_i}\}$（见《代数》引理 \ref{algebra-lemma-support-point}）。
我们对 $\sum n_i$ 作归纳。若 $n_i = 0$ 对所有 $i$ 都成立，则
$\mathcal{I} = \mathcal{O}_X$，结论成立。

\medskip\noindent
设 $n_i > 0$。令 $X' \to X$ 为在 $X$ 中沿 $x_i$ 的爆破
（见引理上方的讨论）。
由于 $\Spec(\mathcal{O}_{X, x_i}) \to X$ 平坦，可知
$X' \times_X \Spec(\mathcal{O}_{X, x_i})$ 是在极大理想处对环
$\mathcal{O}_{X, x_i}$ 的爆破；见《除子》引理
\ref{divisors-lemma-flat-base-change-blowing-up}。
因此交换图中的方块
$$
\xymatrix{
\text{Proj}(\bigoplus\nolimits_{d \geq 0} \mathfrak m_{x_i}^d) \ar[r] \ar[d] &
X' \ar[d] \\
\Spec(\mathcal{O}_{X, x_i}) \ar[r] & X
}
$$
是笛卡尔的。令 $E \subset X'$ 以及
$E' \subset \text{Proj}(\bigoplus\nolimits_{d \geq 0} \mathfrak m_{x_i}^d)$
为例外除子。令 $d \geq 1$ 为引理
\ref{resolve-lemma-blowup-improve} 对理想
$\mathcal{I}_i \subset \mathcal{O}_{X, x_i}$ 所得到的整数。
由于图中的水平箭头平坦，且 $E' \to E$ 满射，且 $E'$ 是 $E$ 的拉回，
可得
$$
\mathcal{I}\mathcal{O}_{X'} \subset \mathcal{O}_{X'}(-dE)
$$
（细节略去）。置 $\mathcal{I}' = \mathcal{I}\mathcal{O}_{X'}(dE) \subset \mathcal{O}_{X'}$。
于是 $\mathcal{O}_{X'}/\mathcal{I}'$ 支持在有限多个闭点
$T' \subset |X'|$ 上，因为这在
$X \setminus \{x_i\}$ 上以及拉回到
$\text{Proj}(\bigoplus\nolimits_{d \geq 0} \mathfrak m_{x_i}^d)$ 时都成立。
引理 \ref{resolve-lemma-blowup-improve} 的最后断言告诉我们：茎
$\mathcal{O}_{X', x'}/\mathcal{I}'\mathcal{O}_{X', x'}$
（其中 $x'$ 位于 $x_i$ 上方）的长度之和 $< n_i$。
因此长度总和下降了。

\medskip\noindent
由归纳假设，存在序列
$$
X'_n \to \ldots \to X'_1 \to X'
$$
它由位于 $T'$ 上方闭点处的爆破组成，且
$\mathcal{I}'\mathcal{O}_{X'_n}$ 可逆。由于
$\mathcal{I}'\mathcal{O}_{X'}(-dE) = \mathcal{I}\mathcal{O}_{X'}$，
可得 $\mathcal{I}\mathcal{O}_{X'_n} =
\mathcal{I}'\mathcal{O}_{X'_n}(-d(f')^{-1}E)$
其中 $f' : X'_n \to X'$ 是复合态射。
注意由《除子》引理
\ref{divisors-lemma-blow-up-pullback-effective-Cartier}，
$(f')^{-1}E$ 是有效 Cartier 除子。因此由《除子》引理
\ref{divisors-lemma-sum-effective-Cartier-divisors} 即得结论。
\end{proof}

\begin{lemma}
\label{resolve-lemma-dominate-by-blowing-up-in-points}
令 $X$ 为 Noether 概形。令 $T \subset X$ 为有限个闭点 $x$ 的集合，
使得 $\mathcal{O}_{X, x}$ 是维数 $2$ 的正则局部环。令
$f : Y \to X$ 为概形之间的 proper 态射，并且在
$U = X \setminus T$ 上是同构。则存在序列
$$
X_n \to X_{n - 1} \to \ldots \to X_1 \to X_0 = X
$$
其中 $X_{i + 1} \to X_i$ 是在 $X_i$ 上沿闭点 $x_i$ 进行的爆破，
且该点位于 $T$ 中某点上方，并且复合态射有分解 $X_n \to Y \to X$。
\end{lemma}

\begin{proof}
由《更多平坦性》引理 \ref{flat-lemma-dominate-modification-by-blowup}，
存在 $U$-容许爆破 $X' \to X$，支配 $Y \to X$。因此可设存在理想层
$\mathcal{I} \subset \mathcal{O}_X$，使得
$\mathcal{O}_X/\mathcal{I}$ 支持在 $T$ 上，且
$Y$ 是 $X$ 沿 $\mathcal{I}$ 的爆破。
由引理 \ref{resolve-lemma-make-ideal-principal}，存在序列
$$
X_n \to X_{n - 1} \to \ldots \to X_1 \to X_0 = X
$$
其中 $X_{i + 1} \to X_i$ 是在 $X_i$ 上沿闭点 $x_i$ 进行的爆破，
且该点位于 $T$ 中某点上方，$\mathcal{I}\mathcal{O}_{X_n}$ 是可逆理想层。
由爆破的泛性质（《除子》引理
\ref{divisors-lemma-universal-property-blowing-up}），
得到所需分解。
\end{proof}

\begin{lemma}
\label{resolve-lemma-extend-rational-map-blowing-up}
令 $S$ 为概形。令 $X$ 为 $S$ 上的概形，正则且维数为 $2$。
令 $Y$ 为 $S$ 上的 proper 概形。给定 $S$-有理映射
$f : U \to Y$，从 $X$ 到 $Y$，存在序列
$$
X_n \to X_{n - 1} \to \ldots \to X_1 \to X_0 = X
$$
以及 $S$-态射 $f_n : X_n \to Y$，使得 $X_{i + 1} \to X_i$
是在 $X_i$ 上沿不位于 $U$ 上方的闭点进行的爆破，且 $f_n$ 与 $f$ 一致。
\end{lemma}

\begin{proof}
可设 $U$ 包含每个余维 $1$ 的点；见《态射》引理
\ref{morphisms-lemma-extend-across}。
于是补集 $T \subset X$（对应于 $U$）是有限个闭点，
其局部环是维数 $2$ 的正则环。
应用《除子》引理 \ref{divisors-lemma-extend-rational-map-after-modification}，
得到 proper 态射 $p : X' \to X$，它在 $U$ 上为同构，
并有态射 $f' : X' \to Y$，与 $f$ 在 $U$ 上一致。
对态射 $p : X' \to X$ 应用引理 \ref{resolve-lemma-dominate-by-blowing-up-in-points}。
复合 $X_n \to X' \to Y$ 即为所需态射。
\end{proof}






\section{由正规化爆破支配}
\label{resolve-section-normalized-blowups}

\noindent
本节证明曲面的修改可由一列沿点的正规化爆破支配。

\begin{definition}
\label{resolve-definition-normalized-blowup}
令 $X$ 为概形，使每个拟紧开集都有有限多个不可约分支。
令 $x \in X$ 为闭点。{\it $X$ 在 $x$ 处的正规化爆破}是复合
$X'' \to X' \to X$，其中 $X' \to X$ 是沿 $X$ 中 $x$ 的爆破，
而 $X'' \to X'$ 是 $X'$ 的正规化。
\end{definition}

\noindent
这里，正规化 $X'' \to X'$ 的定义使用《除子》引理
\ref{divisors-lemma-blow-up-and-irreducible-components}：概形 $X'$
有由不可约分支有限的开集组成的开覆盖。
正规化的定义见《态射》定义 \ref{morphisms-definition-normalization}。

\medskip\noindent
一般而言，即使 $X$ Noether，正规化爆破也未必 proper。
回忆，若概形有由 Nagata 环的谱组成的仿射开覆盖，
则称其为 Nagata（《性质》定义 \ref{properties-definition-nagata}）。

\begin{lemma}
\label{resolve-lemma-Nagata-normalized-blowup}
在定义 \ref{resolve-definition-normalized-blowup} 中，若 $X$ 是 Nagata，
则 $X$ 在 $x$ 处的正规化爆破是正规、Nagata 的，且在 $X$ 上 proper。
\end{lemma}

\begin{proof}
爆破态射 $X' \to X$ 是 proper 的（由于 $X$ 局部 Noether，
可应用《除子》引理 \ref{divisors-lemma-blowing-up-projective}）。
因此 $X'$ 是 Nagata 的（《态射》引理
\ref{morphisms-lemma-finite-type-nagata}）。
于是正规化 $X'' \to X'$ 是有限的（《态射》引理
\ref{morphisms-lemma-nagata-normalization}），
从而 $X'' \to X$ 也 proper（《态射》引理
\ref{morphisms-lemma-finite-proper} 和
\ref{morphisms-lemma-composition-proper}）。
因此正规化爆破是正规（《态射》引理
\ref{morphisms-lemma-normalization-normal}）的 Nagata 代数空间。
\end{proof}

\noindent
在下面的引理中，我们需要假设 $X$ 是 Noether，以保证它有有限多个不可约分支。
于是 proper 性质保证 $f : Y \to X$ 的 $Y$ 也有有限多个不可约分支，
并且要求 $f$ 为双有理态射是有意义的
（《态射》定义 \ref{morphisms-definition-birational}）。

\begin{lemma}
\label{resolve-lemma-dominate-by-normalized-blowing-up}
令 $X$ 为 Noether、Nagata 且维数为 $2$ 的概形。
令 $f : Y \to X$ 为 proper 双有理态射。
则存在交换图
$$
\xymatrix{
X_n \ar[r] \ar[d] &
X_{n - 1} \ar[r] &
\ldots \ar[r] &
X_1 \ar[r] &
X_0 \ar[d] \\
Y \ar[rrrr]  & & & & X
}
$$
其中 $X_0 \to X$ 是正规化，且 $X_{i + 1} \to X_i$ 是在闭点处
对 $X_i$ 的正规化爆破。
\end{lemma}

\begin{proof}
下面将不再另行提及《态射》各节的结果：
\ref{morphisms-section-nagata}、
\ref{morphisms-section-dimension-formula} 以及
\ref{morphisms-section-normalization}。
可将 $Y$ 替换为其正规化。令 $X_0 \to X$ 为正规化。
态射 $Y \to X$ 经由 $X_0$ 分解。因此可设 $X$ 与 $Y$ 均正规。

\medskip\noindent
设 $X$ 与 $Y$ 正规。态射 $f : Y \to X$ 在某个开集上是同构，
该开集包含每个余维 $0$、$1$ 的点（位于 $Y$ 中）以及纤维有限的 $Y$ 中各点；见《簇》引理
\ref{varieties-lemma-modification-normal-iso-over-codimension-1}。
因此存在有限个闭点 $T \subset X$，使 $f$ 在 $X \setminus T$ 上为同构。
对每个 $x \in T$，纤维 $Y_x$ 是维数 $1$ 的 proper 几何连通概形，定义在 $\kappa(x)$ 上；见
《更多态射》引理
\ref{more-morphisms-lemma-geometrically-connected-fibres-towards-normal}。
于是
$$
BadCurves(f) = \{C \subset Y\text{ 闭} \mid
\dim(C) = 1, f(C) = \text{a point}\}
$$
是有限集。我们对 $BadCurves(f)$ 的元素个数作归纳证明引理。
基例是 $BadCurves(f)$ 为空的情形；此时 $f$ 是同构。

\medskip\noindent
固定 $x \in T$。令 $X' \to X$ 为 $X$ 在 $x$ 处的正规化爆破，令
$Y'$ 为 $Y \times_X X'$ 的正规化。图示为
$$
\xymatrix{
Y' \ar[r]_{f'} \ar[d] & X' \ar[d] \\
Y \ar[r]^f & X
}
$$
令 $x' \in X'$ 为位于 $x$ 上方的闭点，使纤维
$Y'_{x'}$ 的维数 $\geq 1$。令 $C' \subset Y'$ 为 $Y'_{x'}$ 的不可约分支，
即 $C' \in BadCurves(f')$。由于 $Y' \to Y \times_X X'$ 有限，
可知 $C'$ 必映到某个不可约分支 $C \subset Y_x$。
显然 $C \in BadCurves(f)$。
由于 $Y' \to Y$ 是双有理的，因而在余维 $1$ 的点上（位于 $Y$ 中）为同构，
可得单射
$$
BadCurves(f') \longrightarrow BadCurves(f)
$$
因此只需证明：经过有限次这样的正规化爆破后，至少消去一条坏曲线，
即上面显示的映射不满射。

\medskip\noindent
我们使用 Zariski 的论证消去一条坏曲线。
取 $C \in BadCurves(f)$，它位于 $x$ 上方。记 $\mathcal{O}_{Y, C}$
为 $Y$ 在 $C$ 的一般点处的局部环。选取元素
$u \in \mathcal{O}_{X, C}$，其在剩余域 $R(C)$ 中的像
在 $\kappa(x)$ 上超越（这是可能的，因为由《簇》引理
\ref{varieties-lemma-dimension-locally-algebraic}，
$R(C)$ 的超越次数为 $1$（相对于 $\kappa(x)$）。
可写 $u = a/b$，其中 $a, b \in \mathcal{O}_{X, x}$，因为
$\mathcal{O}_{Y, C}$ 与 $\mathcal{O}_{X, x}$ 有相同的分式域。
由 $u$ 的选取，必有 $a, b \in \mathfrak m_x$。于是
$$
N_{u, a, b} = \min
\{\text{ord}_{\mathcal{O}_{Y, C}}(a), \text{ord}_{\mathcal{O}_{Y, C}}(b)\} > 0
$$
因此可对该整数作递降归纳。令 $X' \to X$ 为 $x$ 处的正规化爆破，
并令 $Y'$ 为如上的 $X' \times_X Y$ 的正规化。我们将证明：若 $C$ 是某条坏曲线
$C' \subset Y'$ 的像，且位于 $x' \in X'$ 上方，则
存在选择 $a', b' \mathcal{O}_{X', x'}$，使
$N_{u, a', b'} < N_{u, a, b}$。这就完成证明。
确切地说，由于 $X' \to X$ 经由爆破分解，存在非零元素
$d \in \mathfrak m_{x'}$，使得 $a = a' d$ 且 $b = b' d$
（取爆破的例外除子的局部方程作为 $d$ 即可）。由于 $Y' \to Y$
在包含 $C$ 一般点的开集上是同构（如上所见），可知
$\mathcal{O}_{Y', C'} = \mathcal{O}_{Y, C}$。于是
$$
\text{ord}_{\mathcal{O}_{Y, C}}(a) =
\text{ord}_{\mathcal{O}_{Y', C'}}(a' d) =
\text{ord}_{\mathcal{O}_{Y', C'}}(a') +
\text{ord}_{\mathcal{O}_{Y', C'}}(d) >
\text{ord}_{\mathcal{O}_{Y', C'}}(a')
$$
对 $b$ 同理，证明完成。
\end{proof}

\begin{lemma}
\label{resolve-lemma-extend-rational-map-normalized-blowing-up}
令 $S$ 为概形。令 $X$ 为概形，定义在 $S$ 上，Noether、Nagata 且维数为 $2$。
令 $Y$ 为 $S$ 上的 proper 概形。给定
$S$-有理映射 $f : U \to Y$，从 $X$ 到 $Y$，存在序列
$$
X_n \to X_{n - 1} \to \ldots \to X_1 \to X_0 \to X
$$
以及 $S$-态射 $f_n : X_n \to Y$，使得 $X_0 \to X$ 是正规化，
$X_{i + 1} \to X_i$ 是在闭点处对 $X_i$ 的正规化爆破，且 $f_n$ 与 $f$ 一致。
\end{lemma}

\begin{proof}
应用《除子》引理 \ref{divisors-lemma-extend-rational-map-after-modification}，
得到 proper 态射 $p : X' \to X$，它在 $U$ 上为同构，
并有态射 $f' : X' \to Y$，与 $f$ 在 $U$ 上一致。
对态射 $p : X' \to X$ 应用引理 \ref{resolve-lemma-dominate-by-normalized-blowing-up}。
复合 $X_n \to X' \to Y$ 即为所需态射。
\end{proof}





\section{在局部环上作修改}
\label{resolve-section-modifications}

\noindent
令 $S$ 为概形。令 $s_1, \ldots, s_n \in S$ 为两两不同的闭点。
假设开浸入
$$
U = S \setminus \{s_1, \ldots, s_n\} \longrightarrow S
$$
是拟紧的。记 $FP_{S, \{s_1, \ldots, s_n\}}$ 为有限表示态射
$f : X \to S$ 的范畴，其中态射诱导同构 $f^{-1}(U) \to U$。
态射均为 $S$ 上概形的态射。
对每个 $i$，置 $S_i = \Spec(\mathcal{O}_{S, s_i})$，
并令 $V_i = S_i \setminus \{s_i\}$。记 $FP_{S_i, s_i}$ 为有限表示态射
$g_i : Y_i \to S_i$ 的范畴，其中态射诱导同构
$g_i^{-1}(V_i) \to V_i$。
态射均为 $S_i$ 上的态射。基变换定义函子
\begin{equation}
\label{resolve-equation-equivalence}
F :
FP_{S, \{s_1, \ldots, s_n\}}
\longrightarrow
FP_{S_1, s_1} \times \ldots \times FP_{S_n, s_n}
\end{equation}
为将本章至少部分问题归约到局部环情形，有以下引理。

\begin{lemma}
\label{resolve-lemma-equivalence}
函子 $F$（\ref{resolve-equation-equivalence}）是等价。
\end{lemma}

\begin{proof}
当 $n = 1$ 时，这是《极限》引理 \ref{limits-lemma-modifications}。
当 $n > 1$ 时可完全类似证明，或由此推出。例如，设
$g_i : Y_i \to S_i$ 是 $FP_{S_i, s_i}$ 中的对象。
由 $n = 1$ 的情形，可找到有限表示的 $f'_i : X'_i \to S$，
它们在 $S \setminus \{s_i\}$ 上为同构，且到 $S_i$ 的基变换为 $g_i$。
于是置
$$
f : X = X'_1 \times_S \ldots \times_S X'_n \to S
$$
这是 $FP_{S, \{s_1, \ldots, s_n\}}$ 中的对象，
其沿 $S_i \to S$ 的基变换恢复 $g_i$。因此函子本质满射。
满忠实性的证明略去。
\end{proof}

\begin{lemma}
\label{resolve-lemma-equivalence-properties}
令 $S, s_i, S_i$ 如（\ref{resolve-equation-equivalence}）中所设。
若 $f : X \to S$ 对应 $g_i : Y_i \to S_i$，对应关系由 $F$ 给出，
则 $f$ 分离、proper、有限，当且仅当 $g_i$ 对 $i = 1, \ldots, n$ 具有相应性质。
\end{lemma}

\begin{proof}
由《极限》引理
\ref{limits-lemma-modifications-properties} 得出。
\end{proof}

\begin{lemma}
\label{resolve-lemma-equivalence-fibre}
令 $S, s_i, S_i$ 如（\ref{resolve-equation-equivalence}）中所设。
若 $f : X \to S$ 对应 $g_i : Y_i \to S_i$，对应关系由 $F$ 给出，
则 $X_{s_i} \cong (Y_i)_{s_i}$ 是在 $\kappa(s_i)$ 上的概形同构。
\end{lemma}

\begin{proof}
显然。
\end{proof}

\begin{lemma}
\label{resolve-lemma-equivalence-sequence-blowups}
令 $S, s_i, S_i$ 如（\ref{resolve-equation-equivalence}）中所设，
并假设 $f : X \to S$ 对应 $g_i : Y_i \to S_i$，其中对应关系由 $F$ 给出。
则存在分解
$$
X = Z_m \to Z_{m - 1} \to \ldots \to Z_1 \to Z_0 = S
$$
这是 $f$ 的分解，其中 $Z_{j + 1} \to Z_j$ 是在 $Z_j$ 上沿闭点 $z_j$ 进行的爆破，
且该点位于 $\{s_1, \ldots, s_n\}$ 上方；当且仅当对每个
$i$ 存在分解
$$
Y_i = Z_{i, m_i} \to Z_{i, m_i - 1} \to \ldots \to Z_{i, 1} \to Z_{i, 0} = S_i
$$
这是 $g_i$ 的分解，其中 $Z_{i, j + 1} \to Z_{i, j}$ 是在 $Z_{i, j}$ 上沿闭点 $z_{i, j}$ 进行的爆破，
且该点位于 $s_i$ 上方。
\end{lemma}

\begin{proof}
从一列爆破开始
$Z_m \to Z_{m - 1} \to \ldots \to Z_1 \to Z_0 = S$。
第一个态射 $Z_1 \to S$ 是沿某个 $s_i$（不妨设为 $s_1$）的爆破。
应用 $F$ 于 $Z_1 \to S$，得到 $Z_{1, 1} \to S_1$ 在 $s_1$ 处的爆破，
即在 $s_1$ 处爆破，而有 $Z_{i, 0} = S_i$（当 $i > 1$ 时）。
下一步中，我们或者沿某个 $s_i$（$i \geq 2$）爆破 $Z_1$，
或者选取闭点 $z_1$，它属于 $Z_1 \to S$ 在 $s_1$ 上方的纤维。
第一种情形显然；第二种情形利用
$(Z_1)_{s_1} \cong (Z_{1, 1})_{s_1}$
（引理 \ref{resolve-lemma-equivalence-fibre}），得到闭点 $z_{1, 1} \in Z_{1, 1}$，
它对应于 $z_1$。
然后令 $Z_{1, 2} \to Z_{1, 1}$ 为沿 $z_{1, 1}$ 的爆破。
如此继续即可构造每个 $g_i$ 的分解。

\medskip\noindent
反之，给定爆破序列
$Z_{i, m_i} \to Z_{i, m_i - 1} \to \ldots \to Z_{i, 1} \to Z_{i, 0} = S_i$，
以完全相同的方式构造 $S$ 的爆破序列。
\end{proof}

\noindent
下面是关于正规化爆破的
引理 \ref{resolve-lemma-equivalence-sequence-blowups} 的类似结果。

\begin{lemma}
\label{resolve-lemma-equivalence-sequence-normalized-blowups}
令 $S, s_i, S_i$ 如（\ref{resolve-equation-equivalence}）中所设，
并假设 $f : X \to S$ 对应于 $g_i : Y_i \to S_i$，对应关系由 $F$ 给出。
假设 $S$ 的每个拟紧开子集都有有限个不可约分支。
则存在一个分解
$$
X = Z_m \to Z_{m - 1} \to \ldots \to Z_1 \to Z_0 = S
$$
这是 $f$ 的分解，其中 $Z_{j + 1} \to Z_j$ 是在 $Z_j$ 上沿闭点
$z_j$ 作的正规化爆破，该点位于 $\{x_1, \ldots, x_n\}$ 上方；当且仅当对每个 $i$
都存在一个分解
$$
Y_i = Z_{i, m_i} \to Z_{i, m_i - 1} \to \ldots \to Z_{i, 1} \to Z_{i, 0} = S_i
$$
这是 $g_i$ 的分解，其中 $Z_{i, j + 1} \to Z_{i, j}$ 是在 $Z_{i, j}$
上沿闭点 $z_{i, j}$ 作的正规化爆破，该点位于 $s_i$ 上方。
\end{lemma}

\begin{proof}
关于 $S$ 的假设用于逐步保证我们所正规化的概形局部只有有限个不可约分支，
从而上述陈述有意义。既然如此，该引理就由证明
引理 \ref{resolve-lemma-equivalence-sequence-blowups} 时所用的完全相同论证得到。
\end{proof}








\section{消没}
\label{resolve-section-vanishing}

\noindent
本节中我们经常在以下情形下工作。
回忆一下，修改是整概形之间的 proper 双有理态射（《态射》，定义
\ref{morphisms-definition-modification}）。

\begin{situation}
\label{resolve-situation-vanishing}
这里 $(A, \mathfrak m, \kappa)$ 是维数为 $2$ 的局部诺特正规整环。
令 $s$ 为 $S = \Spec(A)$ 的闭点，并令 $U = S \setminus \{s\}$。
令 $f : X \to S$ 为修改。记 $C_1, \ldots, C_r$ 为特殊纤维
$X_s$ 的不可约分支；这里该特殊纤维来自 $f$。
\end{situation}

\noindent
由《簇》引理
\ref{varieties-lemma-modification-normal-iso-over-codimension-1}，
态射 $f$ 定义了同构 $f^{-1}(U) \to U$。
由《簇》引理 \ref{varieties-lemma-dimension-fibre-in-higher-codimension}，
特殊纤维 $X_s$ 在 $\Spec(\kappa)$ 上 proper，且维数至多为 $1$。
由 Stein 分解（《态射进阶》引理
\ref{more-morphisms-lemma-geometrically-connected-fibres-towards-normal}），
有 $f_*\mathcal{O}_X = \mathcal{O}_S$，且特殊纤维 $X_s$ 在 $\kappa$ 上几何连通。
若 $X_s$ 的维数为 $0$，则 $f$ 是有限态射（《态射进阶》引理
\ref{more-morphisms-lemma-proper-finite-fibre-finite-in-neighbourhood}），
因而是同构（《态射》引理 \ref{morphisms-lemma-finite-birational-over-normal}）。
我们舍弃这个无趣的情形，于是得到
$\dim(C_i) = 1$，其中 $i = 1, \ldots, r$。

\begin{lemma}
\label{resolve-lemma-dominate-by-scheme-modification}
在情形 \ref{resolve-situation-vanishing} 中，存在 $U$-容许爆破
$X' \to S$，它支配 $X$。
\end{lemma}

\begin{proof}
这是《平坦性进阶》引理
\ref{flat-lemma-dominate-modification-by-blowup} 的特殊情形。
\end{proof}

\begin{lemma}
\label{resolve-lemma-nice-meromorphic-function}
在情形 \ref{resolve-situation-vanishing} 中，存在非零的
$f \in \mathfrak m$，使得对每个 $i = 1, \ldots, r$，存在
\begin{enumerate}
\item 闭点 $x_i \in C_i$，满足 $x_i \not \in C_j$（对 $j \not = i$），
\item 分解 $f = g_i f_i$，它是 $f$ 在 $\mathcal{O}_{X, x_i}$ 中的分解，其中
$g_i \in \mathfrak m_{x_i}$ 映到 $\mathcal{O}_{C_i, x_i}$ 的非零元素。
\end{enumerate}
\end{lemma}

\begin{proof}
以下不再特别提及情形 \ref{resolve-situation-vanishing} 之后的观察。
选取闭点 $x_i \in C_i$，它不属于 $C_j$，其中 $j \not = i$。
选取 $g_i \in \mathfrak m_{x_i}$，使其映到
$\mathcal{O}_{C_i, x_i}$ 的非零元素。由于 $A$ 的分式域就是
$\mathcal{O}_{X_i, x_i}$ 的分式域，可写成
$g_i = a_i/b_i$，其中 $a_i, b_i \in A$。取 $f = \prod a_i$。
\end{proof}

\begin{lemma}
\label{resolve-lemma-nontrivial-normal-bundle}
在情形 \ref{resolve-situation-vanishing} 中，假设 $X$ 正规。
令 $Z \subset X$ 为非空有效 Cartier 除子，并假设作为集合有
$Z \subset X_s$。则 $Z$ 的余法层非平凡。
更确切地说，存在某个 $i$ 使得 $C_i \subset Z$ 且
$\deg(\mathcal{C}_{Z/X}|_{C_i}) > 0$。
\end{lemma}

\begin{proof}
以下不再特别提及情形 \ref{resolve-situation-vanishing} 之后的观察。
令 $f$ 为引理 \ref{resolve-lemma-nice-meromorphic-function} 中的函数。
令 $\xi_i \in C_i$ 为泛点，令 $\mathcal{O}_i$ 为 $X$ 在 $\xi_i$ 处的局部环。
则 $\mathcal{O}_i$ 是离散赋值环。令 $e_i$ 为 $f$ 在 $\mathcal{O}_i$ 中的赋值，
于是 $e_i > 0$。令 $h_i \in \mathcal{O}_i$ 为 $Z$ 的局部方程，令 $d_i$ 为其赋值。
于是 $d_i \geq 0$。选取并固定 $i$，使 $d_i/e_i$ 最大
（于是 $d_i > 0$，因为 $Z$ 非空）。
将 $f$ 替换为 $f^{d_i}$，将 $Z$ 替换为 $e_iZ$。
这是允许的，因为有关系
$\mathcal{O}_X(e_i Z) = \mathcal{O}_X(Z)^{\otimes e_i}$，以及余法层与
$\mathcal{O}_X(Z)$ 之间的关系（见《除子》引理
\ref{divisors-lemma-invertible-sheaf-sum-effective-Cartier-divisors}
和 \ref{divisors-lemma-conormal-effective-Cartier-divisor}；又因次数乘以 $e_i$，见
《簇》引理 \ref{varieties-lemma-degree-tensor-product}）。
令 $\mathcal{I}$ 为 $Z$ 的理想层，使得 $\mathcal{C}_{Z/X} = \mathcal{I}|_Z$。
考虑像 $\overline{f}$，即 $f$ 在 $\Gamma(Z, \mathcal{O}_Z)$ 中的像。
由上述选择可知，$\overline{f}$ 在 $Z$ 的不可约分支的泛点处消失
（这些泛点都是某个 $C_j$ 的泛点，因为 $Z$ 包含于特殊纤维中）。另一方面，
由《除子》引理 \ref{divisors-lemma-normal-effective-Cartier-divisor-S1}，
$Z$ 是 $(S_1)$ 的。因而概形 $Z$ 没有嵌入关联点，于是由《除子》引理
\ref{divisors-lemma-S1-no-embedded} 和
\ref{divisors-lemma-restriction-injective-open-contains-weakly-ass} 得到
$\overline{f} = 0$。因此 $f$ 是 $\mathcal{I}$ 的整体截面，按构造生成
$\mathcal{I}_{\xi_i}$。故像 $s_i$，即 $f$ 在
$\Gamma(C_i, \mathcal{I}|_{C_i})$ 中的像，是非零的。然而，$f$ 的选择保证
$s_i$ 在 $x_i$ 处有零点。因此由《簇》引理
\ref{varieties-lemma-check-invertible-sheaf-trivial}，$\mathcal{I}|_{C_i}$ 的次数为 $> 0$。
\end{proof}

\begin{lemma}
\label{resolve-lemma-H1-injective}
在情形 \ref{resolve-situation-vanishing} 中，假设 $X$ 正规且 $A$ 是 Nagata 环。映射
$$
H^1(X, \mathcal{O}_X) \longrightarrow H^1(f^{-1}(U), \mathcal{O}_X)
$$
是单射。
\end{lemma}

\begin{proof}
令 $0 \to \mathcal{O}_X \to \mathcal{E} \to \mathcal{O}_X \to 0$ 为对应于
非平凡元素 $\xi$ 的扩张，该元素属于 $H^1(X, \mathcal{O}_X)$
（《上同调》引理 \ref{cohomology-lemma-h1-extensions}）。
令 $\pi : P = \mathbf{P}(\mathcal{E}) \to X$ 为与 $\mathcal{E}$ 关联的射影丛。
满射 $\mathcal{E} \to \mathcal{O}_X$ 定义截面 $\sigma : X \to P$，
其余法层同构于 $\mathcal{O}_X$（《除子》引理
\ref{divisors-lemma-conormal-sheaf-section-projective-bundle}）。
若 $\xi$ 在 $f^{-1}(U)$ 上的限制平凡，则得到映射
$\mathcal{E}|_{f^{-1}(U)} \to \mathcal{O}_{f^{-1}(U)}$，它劈裂单射
$\mathcal{O}_X \to \mathcal{E}$。这定义了第二个截面
$\sigma' : f^{-1}(U) \to P$，它与 $\sigma$ 不交。由于 $\xi$ 非平凡，可知 $\sigma'$ 不可能延拓到整个
$X$ 且仍与 $\sigma$ 不交。令 $X' \subset P$ 为 $\sigma'$ 的概形论像
（《态射》定义 \ref{morphisms-definition-scheme-theoretic-image}）。图示为
$$
\xymatrix{
& X' \ar[rd]_g \ar[r] & P \ar[d]_\pi \\
f^{-1}(U) \ar[ru]_{\sigma'} \ar[rr] & & X \ar@/_/[u]_\sigma
}
$$
态射 $P \setminus \sigma(X) \to X$ 是仿射的。
若 $X' \cap \sigma(X) = \emptyset$，则 $X' \to X$ 既仿射又 proper，因而有限
（《态射》引理 \ref{morphisms-lemma-finite-proper}），进而是同构
（因为 $X$ 正规，见《态射》引理
\ref{morphisms-lemma-finite-birational-over-normal}）。这与上述结论矛盾。

\medskip\noindent
令 $X^\nu$ 为 $X'$ 的正规化。由于 $A$ 是 Nagata 环，
$X^\nu \to X'$ 是有限态射（《态射》引理
\ref{morphisms-lemma-nagata-normalization} 和
\ref{morphisms-lemma-ubiquity-nagata}）。令 $Z \subset X^\nu$ 为有效 Cartier 除子
$\sigma(X) \subset P$ 的拉回。由上可知 $Z$ 非空，且包含在
$X^\nu \to S$ 的闭纤维中。由于 $P \to X$ 光滑，可知 $\sigma(X)$ 是有效
Cartier 除子（《除子》引理
\ref{divisors-lemma-section-smooth-regular-immersion}）。
因此 $Z \subset X^\nu$ 也是有效 Cartier 除子。由于 $\sigma(X)$ 在 $P$ 中的余法层为
$\mathcal{O}_X$，由《态射》引理
\ref{morphisms-lemma-conormal-functorial-flat}，$Z$ 在 $X^\nu$ 中的余法层
（先验上可逆）为 $\mathcal{O}_Z$。这与引理
\ref{resolve-lemma-nontrivial-normal-bundle} 矛盾，证明完成。
\end{proof}

\begin{lemma}
\label{resolve-lemma-R1-injective}
在情形 \ref{resolve-situation-vanishing} 中，假设 $X$ 正规且 $A$ 是 Nagata 环。则
$$
\Hom_{D(A)}(\kappa[-1], Rf_*\mathcal{O}_X)
$$
为零。这里使用 $D(A) = D_\QCoh(\mathcal{O}_S)$，将
$Rf_*\mathcal{O}_X$ 看作 $D(A)$ 的对象。
\end{lemma}

\begin{proof}
由 $Rf_*$ 与 $Lf^*$ 的伴随性，这样的映射等同于映射
$\alpha : Lf^*\kappa[-1] \to \mathcal{O}_X$。注意
$$
H^i(Lf^*\kappa[-1]) =
\left\{
\begin{matrix}
0 & \text{若} & i > 1 \\
\mathcal{O}_{X_s} & \text{若} & i = 1 \\
\text{some }\mathcal{O}_{X_s}\text{-模} & \text{若} & i \leq 0
\end{matrix}
\right.
$$
由于 $\Hom(H^0(Lf^*\kappa[-1]), \mathcal{O}_X) = 0$
（因为 $\mathcal{O}_X$ 无挠），故 $\Ext$ 的谱序列
（《位点上的上同调》例
\ref{sites-cohomology-example-hom-complex-into-sheaf}）蕴含
$\Hom_{D(\mathcal{O}_X)}(Lf^*\kappa[-1], \mathcal{O}_X)$ 等于
$\Ext^1_{\mathcal{O}_X}(\mathcal{O}_{X_s}, \mathcal{O}_X)$。
于是可知 $\alpha : Lf^*\kappa[-1] \to \mathcal{O}_X$ 由扩张
$$
0 \to \mathcal{O}_X \to \mathcal{E} \to \mathcal{O}_{X_s} \to 0
$$
给出。由引理 \ref{resolve-lemma-H1-injective}，该扩张沿满射
$\mathcal{O}_X \to \mathcal{O}_{X_s}$ 的拉回为零
（因为该拉回在 $f^{-1}(U)$ 上显然劈裂）。因此
$1 \in \mathcal{O}_{X_s}$ 可提升为整体截面 $s$，该截面属于 $\mathcal{E}$。
将 $s$ 乘以理想层 $\mathcal{I}$（即 $X_s$ 的理想层），得到 $\mathcal{O}_X$-模映射
$c_s : \mathcal{I} \to \mathcal{O}_X$。应用 $f_*$ 得到 $A$-线性映射
$f_*c_s : \mathfrak m \to A$。由于 $A$ 是诺特正规局部整环，该映射由某个元素
$a \in A$ 的乘法给出。将 $s$ 改为 $s -  a$，可知 $s$ 被
$\mathcal{I}$ 零化，故该扩张如期平凡。
\end{proof}

\begin{remark}
\label{resolve-remark-dualizing-setup}
令 $X$ 为维数为 $2$ 的整诺特正规概形。此时以下条件等价：
\begin{enumerate}
\item $X$ 有对偶复形 $\omega_X^\bullet$；
\item 存在凝聚 $\mathcal{O}_X$-模 $\omega_X$，使得
$\omega_X[n]$ 是对偶复形，其中 $n$ 可以是任意整数。
\end{enumerate}
这由 $X$ 是 Cohen--Macaulay 概形（《性质》引理
\ref{properties-lemma-normal-dimension-2-Cohen-Macaulay}）以及《概形的对偶性》引理
\ref{duality-lemma-dualizing-module-CM-scheme} 得出。在这种情形下，依照
《概形的对偶性》节 \ref{duality-section-dualizing-module}，称 $\omega_X$ 为
{\it 对偶模}。特别地，当 $A$ 是维数为 $2$ 的诺特正规局部整环时，
若上述条件成立，就称 {\it $A$ 有对偶模 $\omega_A$}。在这种情形下，若
$X \to \Spec(A)$ 是正规修改，则 $X$ 也有对偶模；见《概形的对偶性》例
\ref{duality-example-proper-over-local}。此时总以 $\omega_X$ 表示关于
$\omega_A$ 归一化的对偶模，即使得 $\omega_X[2]$ 是相对于
$\omega_A[2]$ 归一化的对偶复形。见《概形的对偶性》节
\ref{duality-section-glue}。
\end{remark}

\noindent
下一命题中的 Grauert--Riemenschneider 消没是引理
\ref{resolve-lemma-R1-injective} 与一般对偶理论的形式推论。

\begin{proposition}[Grauert-Riemenschneider]
\label{resolve-proposition-Grauert-Riemenschneider}
在情形 \ref{resolve-situation-vanishing} 中，假设
\begin{enumerate}
\item $X$ 是正规概形；
\item $A$ 是 Nagata 环，并有对偶复形 $\omega_A^\bullet$。
\end{enumerate}
令 $\omega_X$ 为 $X$ 的对偶模
（评注 \ref{resolve-remark-dualizing-setup}）。则 $R^1f_*\omega_X = 0$。
\end{proposition}

\begin{proof}
本证明中使用等同 $D(A) = D_\QCoh(\mathcal{O}_S)$，将拟凝聚
$\mathcal{O}_S$-模与 $A$-模相等同。此外，可以假设 $\omega_A^\bullet$
已经归一化；见《对偶复形》节 \ref{dualizing-section-dualizing-local}。
由于 $X$ 是诺特正规 $2$ 维概形，它是 Cohen--Macaulay 概形
（《性质》引理 \ref{properties-lemma-normal-dimension-2-Cohen-Macaulay}）。
因此 $\omega_X^\bullet = \omega_X[2]$（《概形的对偶性》引理
\ref{duality-lemma-dualizing-module-CM-scheme}，以及《概形的对偶性》例
\ref{duality-example-proper-over-local} 中的归一化）。
若命题不成立，则可找到非零映射 $R^1f_*\omega_X \to \kappa$。
换言之，得到非零映射 $\alpha : Rf_*\omega_X^\bullet \to \kappa[1]$。
应用 $R\Hom_A(-, \omega_A^\bullet)$，得到非零映射
$$
\beta : \kappa[-1] \longrightarrow Rf_*\mathcal{O}_X
$$
这与引理 \ref{resolve-lemma-R1-injective} 矛盾。为说明
$R\Hom_A(-, \omega_A^\bullet)$ 的确具有上述作用，先注意到
$$
R\Hom_A(\kappa[1], \omega_A^\bullet) =
R\Hom_A(\kappa, \omega_A^\bullet)[-1] =
\kappa[-1]
$$
因为 $\omega_A^\bullet$ 已归一化，并且有
$$
R\Hom_A(Rf_*\omega_X^\bullet, \omega_A^\bullet) =
Rf_*R\SheafHom_{\mathcal{O}_X}(\omega_X^\bullet, \omega_X^\bullet) =
Rf_*\mathcal{O}_X
$$
第一个等号来自《概形的对偶性》例
\ref{duality-example-iso-on-RSheafHom-noetherian} 以及按构造成立的
$\omega_X^\bullet = f^!\omega_A^\bullet$；第二个等号成立是因为
$\omega_X^\bullet$ 是 $X$ 的对偶复形（这归结为《概形的对偶性》引理
\ref{duality-lemma-shriek-dualizing}）。
\end{proof}





\section{有界性}
\label{resolve-section-bounded}

\noindent
本节开始讨论如何将 $2$ 维概形的情形归约到有理奇点情形。

\begin{lemma}
\label{resolve-lemma-exact-sequence}
令 $(A, \mathfrak m, \kappa)$ 为维数为 $2$ 的诺特正规局部整环。
考虑交换图
$$
\xymatrix{
X' \ar[rd]_{f'} \ar[rr]_g & & X \ar[ld]^f \\
& \Spec(A)
}
$$
其中 $f$ 与 $f'$ 是情形 \ref{resolve-situation-vanishing} 中的修改，且 $X$ 正规。
于是有短正合列
$$
0 \to H^1(X, \mathcal{O}_X) \to H^1(X', \mathcal{O}_{X'}) \to
H^0(X, R^1g_*\mathcal{O}_{X'}) \to 0
$$
此外，$\dim(\text{Supp}(R^1g_*\mathcal{O}_{X'})) = 0$，且
$R^1g_*\mathcal{O}_{X'}$ 由整体截面生成。
\end{lemma}

\begin{proof}
以下不再特别提及情形 \ref{resolve-situation-vanishing} 之后的观察。
由于 $X$ 正规，而 $g$ 是支配双有理态射，有
$g_*\mathcal{O}_{X'} = \mathcal{O}_X$；例如见《态射进阶》引理
\ref{more-morphisms-lemma-geometrically-connected-fibres-towards-normal}。
由于 $g$ 的纤维维数为 $\leq 1$，有
$R^pg_*\mathcal{O}_{X'} = 0$（当 $p > 1$ 时）；例如见《概形的上同调》引理
\ref{coherent-lemma-higher-direct-images-zero-above-dimension-fibre}。
$R^1g_*\mathcal{O}_{X'}$ 的支撑包含于 $X$ 中满足 $g'$ 的纤维维数为
$\geq 1$ 的点集。因此，它包含于那些不可约分支 $C' \subset X'_s$ 的像集，
其中这些分支映到 $X_s$ 的点；这是有限个闭点组成的集合
（回忆 $X'_s \to X_s$ 是 proper $1$ 维概形之间的态射，这些概形在 $\kappa$ 上）。
于是由《概形的上同调》引理
\ref{coherent-lemma-coherent-support-dimension-0}，$R^1g_*\mathcal{O}_{X'}$
由整体截面生成。利用态射 $f : X \to S$ 以及上述参考文献可知，
$H^p(X, \mathcal{F}) = 0$（当 $p > 1$ 时），这里 $\mathcal{O}_X$-模
$\mathcal{F}$ 可以是任意凝聚模。因此，引理中的短正合列由 $g$ 和
$\mathcal{O}_{X'}$ 的 Leray 谱序列得到；见《上同调》引理
\ref{cohomology-lemma-Leray}。
\end{proof}

\begin{lemma}
\label{resolve-lemma-bound-a-torsion}
令 $(A, \mathfrak m, \kappa)$ 为维数为 $2$ 的局部正规 Nagata 整环。
令 $a \in A$ 非零。存在整数 $N$，使得对每个修改
$f : X \to \Spec(A)$，若 $X$ 正规，则 $A$-模
$$
M_{X, a} = \Coker(A \longrightarrow H^0(Z, \mathcal{O}_Z))
$$
其中 $Z \subset X$ 由 $a$ 截出；该模的长度以 $N$ 为界。
\end{lemma}

\begin{proof}
由短正合列
$
0 \to \mathcal{O}_X \xrightarrow{a} \mathcal{O}_X \to \mathcal{O}_Z \to 0
$
可知
\begin{equation}
\label{resolve-equation-a-torsion}
M_{X, a} = H^1(X, \mathcal{O}_X)[a]
\end{equation}
这里置 $N[a] = \{n \in N \mid an = 0\}$；对一个 $A$-模 $N$ 采用此记号。因此若
$a$ 整除 $b$，则 $M_{X, a} \subset M_{X, b}$。
假设对某个 $c \in A$，模 $M_{X, c}$ 的长度有界。
那么对每个 $X$，有正合列
$$
0 \to M_{X, c} \to M_{X, c^2} \to M_{X, c}
$$
其中第二个箭头由乘以 $c$ 给出。因此 $M_{X, c^2}$ 的长度也有界。
所以，只须找到使本引理成立的 $c \in A$，并满足 $a$ 整除某个
$c^n$，其中 $n > 0$。由《代数进阶》引理
\ref{more-algebra-lemma-divides-radical}，可以假设 $A/(a)$ 是既约环。

\medskip\noindent
假设 $A/(a)$ 既约。令 $A/(a) \subset B$ 为 $A/(a)$ 在其商环中的正规化。
由于 $A$ 是 Nagata 环，$\Coker(A \to B)$ 是有限模。我们断言这个有限模的长度
就是一个界。为此，考虑引理中的 $f : X \to \Spec(A)$，并令
$Z' \subset Z$ 为 $Z \cap f^{-1}(U)$ 的概形论闭包。
例如由《簇》引理 \ref{varieties-lemma-finite-in-codim-1}，
$Z' \to \Spec(A/(a))$ 是有限态射。因此 $Z' = \Spec(B')$，其中
$A/(a) \subset B' \subset B$。另一方面，我们断言映射
$$
H^0(Z, \mathcal{O}_Z) \to H^0(Z', \mathcal{O}_{Z'})
$$
是单射。事实上，若 $s \in H^0(Z, \mathcal{O}_Z)$ 属于其核，
则 $s$ 在 $f^{-1}(U) \cap Z$ 上的限制为零。因此 $s$ 在
$H^1(X, \mathcal{O}_X)$ 中的像在 $H^1(f^{-1}(U), \mathcal{O}_X)$ 中消失。
由引理 \ref{resolve-lemma-H1-injective} 可知，$s$ 来自某个元素 $\tilde s$，该元素属于 $A$。
但按假设，$\tilde s$ 在 $B'$ 中的像为零，这蕴含 $s = 0$。
综上，$M_{X, a}$ 是 $B'/A$ 的子商；也就是说，$B'$ 的元素未必都能延拓为
$\mathcal{O}_Z$ 的整体截面，但无论如何，$M_{X, a}$ 的长度以 $B/A$ 的长度为界。
\end{proof}

\noindent
在某些情形下，奇点解消可归约到有理奇点情形。

\begin{definition}
\label{resolve-definition-reduce-to-rational}
令 $(A, \mathfrak m, \kappa)$ 为维数为 $2$ 的局部正规 Nagata 整环。
\begin{enumerate}
\item 称 $A$ {\it 定义一个有理奇点}，如果对每个正规修改
$X \to \Spec(A)$ 都有 $H^1(X, \mathcal{O}_X) = 0$。
\item 称 {\it 对 $A$ 可以归约到有理奇点}，如果这些 $A$-模
$$
H^1(X, \mathcal{O}_X)
$$
的长度对所有修改 $X \to \Spec(A)$（其中 $X$ 正规）一致有界。
\end{enumerate}
\end{definition}

\noindent
第 (2) 项中的说法由引理 \ref{resolve-lemma-reduce-to-rational} 解释。
粗略地说，下面的引理断言：$\Spec(A)$ 的修改的局部环，在 $A$ 定义有理奇点时
也定义有理奇点。

\begin{lemma}
\label{resolve-lemma-rational-propagates}
令 $(A, \mathfrak m, \kappa)$ 为维数为 $2$、定义有理奇点的局部正规 Nagata 整环。
令 $A \subset B$ 为具有相同分式域、且本质有限型的整环局部扩张，
并假设 $\dim(B) = 2$ 且 $B$ 正规。则 $B$ 定义有理奇点。
\end{lemma}

\begin{proof}
选取有限型 $A$-代数 $C$，使得 $B = C_\mathfrak q$，其中
$\mathfrak q \subset C$ 是素理想。将 $C$ 替换为 $C$ 在 $B$ 中的像后，可假设 $C$ 是整环，
其分式域等于 $A$ 的分式域。然后可以选取闭浸入
$\Spec(C) \to \mathbf{A}^n_A$，并在 $\mathbf{P}^n_A$ 中取闭包，从而得到 $B$
同构于 $\mathcal{O}_{X, x}$；这里 $x \in X$ 是某个射影修改
$X \to \Spec(A)$ 的闭点。（《态射》引理
\ref{morphisms-lemma-dimension-formula} 表明 $\kappa(x)$ 是 $\kappa$ 的有限扩张，
于是《态射》引理
\ref{morphisms-lemma-algebraic-residue-field-extension-closed-point-fibre}
表明 $x$ 是闭点。）令 $\nu : X^\nu \to X$ 为正规化。
由于 $A$ 是 Nagata 环，态射 $\nu$ 有限（《态射》引理
\ref{morphisms-lemma-nagata-normalization}）。因此由《态射进阶》引理
\ref{more-morphisms-lemma-category-projective}，$X^\nu$ 在 $A$ 上是射影的。
由于 $B = \mathcal{O}_{X, x}$ 正规，有
$\mathcal{O}_{X, x} = (\nu_*\mathcal{O}_{X^\nu})_x$。
因此存在唯一点 $x^\nu \in X^\nu$，它位于 $x$ 上方，且
$\mathcal{O}_{X^\nu, x^\nu} = \mathcal{O}_{X, x}$。
故可假设 $X$ 正规且在 $A$ 上射影。
令 $Y \to \Spec(\mathcal{O}_{X, x}) = \Spec(B)$ 为修改，并假设 $Y$ 正规。
需要证明 $H^1(Y, \mathcal{O}_Y) = 0$。由《极限》引理
\ref{limits-lemma-modifications}，可以找到概形态射 $g : X' \to X$，
它在 $X \setminus \{x\}$ 上为同构，并使
$X' \times_X \Spec(\mathcal{O}_{X, x})$ 同构于 $Y$。由《极限》引理
\ref{limits-lemma-modifications-properties}，$g$ 是 proper 的，故为修改。
$X'$ 在点 $x'$ 处的局部环同构于 $X$ 在 $g(x')$ 处的局部环，前提是
$g(x') \not = x$；
在 $g(x') = x$ 时，$X'$ 在 $x'$ 处的局部环同构于 $Y$ 在对应点处的局部环。
可知 $X'$ 正规，因为 $X$ 与 $Y$ 都正规。因此
$H^1(X', \mathcal{O}_{X'}) = 0$，这来自关于 $A$ 的假设。由引理 \ref{resolve-lemma-exact-sequence} 有
$R^1g_*\mathcal{O}_{X'} = 0$。这显然意味着如期有
$H^1(Y, \mathcal{O}_Y) = 0$。
\end{proof}

\begin{lemma}
\label{resolve-lemma-reduce-to-rational}
令 $(A, \mathfrak m, \kappa)$ 为维数为 $2$ 的局部正规 Nagata 整环。
若对 $A$ 可以归约到有理奇点，则存在一列有限的正规化爆破
$$
X = X_n \to X_{n - 1} \to \ldots \to X_1 \to X_0 = \Spec(A)
$$
其中心均为闭点，并且对任意闭点 $x \in X$，局部环
$\mathcal{O}_{X, x}$ 定义有理奇点。特别地，$X \to \Spec(A)$ 是修改，
而 $X$ 是在 $A$ 上射影的正规概形。
\end{lemma}

\begin{proof}
选取修改 $X \to \Spec(A)$，其中 $X$ 正规，并使
$H^1(X, \mathcal{O}_X)$ 的长度最大。由引理 \ref{resolve-lemma-exact-sequence}，
对任意进一步的修改 $g : X' \to X$，若 $X'$ 正规，则有
$R^1g_*\mathcal{O}_{X'} = 0$ 以及
$H^1(X, \mathcal{O}_X) = H^1(X', \mathcal{O}_{X'})$。

\medskip\noindent
令 $x \in X$ 为闭点。我们将证明 $\mathcal{O}_{X, x}$ 定义有理奇点。
令 $Y \to \Spec(\mathcal{O}_{X, x})$ 为修改，并假设 $Y$ 正规。
需要证明 $H^1(Y, \mathcal{O}_Y) = 0$。由《极限》引理
\ref{limits-lemma-modifications}，可找到概形态射 $g : X' \to X$，
它在 $X \setminus \{x\}$ 上为同构，并使
$X' \times_X \Spec(\mathcal{O}_{X, x})$ 同构于 $Y$。由《极限》引理
\ref{limits-lemma-modifications-properties}，$g$ 是 proper 的，故为修改。
$X'$ 在点 $x'$ 处的局部环同构于 $X$ 在 $g(x')$ 处的局部环，前提是
$g(x') \not = x$；若 $g(x') = x$，则 $X'$ 在 $x'$ 处的局部环同构于
$Y$ 在对应点处的局部环。可知 $X'$ 正规，因为 $X$ 与 $Y$ 都正规。
由极大性有 $R^1g_*\mathcal{O}_{X'} = 0$（见第一段）。这显然意味着如期有
$H^1(Y, \mathcal{O}_Y) = 0$。

\medskip\noindent
结论是，我们找到了一个正规修改 $X$，它是 $\Spec(A)$ 的修改，使 $X$ 在闭点处的局部环
全都定义有理奇点。然后选取一列正规化爆破
$X_n \to \ldots \to X_1 \to \Spec(A)$，使 $X_n$ 支配 $X$；见引理
\ref{resolve-lemma-dominate-by-normalized-blowing-up}。对闭点 $x' \in X_n$，若它映到
$x \in X$，则可将引理 \ref{resolve-lemma-rational-propagates} 应用于环映射
$\mathcal{O}_{X, x} \to \mathcal{O}_{X_n, x'}$，从而看出
$\mathcal{O}_{X_n, x'}$ 定义有理奇点。
\end{proof}

\begin{lemma}
\label{resolve-lemma-go-up-separable}
令 $A \to B$ 为维数为 $2$ 的局部正规 Nagata 整环之间的有限单射局部环映射。
假设所诱导的分式域扩张可分。若对 $A$ 可以归约到有理奇点，
则对 $B$ 也可以。
\end{lemma}

\begin{proof}
令 $n$ 为分式域扩张 $L/K$ 的次数。令
$\text{Trace}_{L/K} : L \to K$ 为迹。由于该扩张有限可分，迹配对
$(h, g) \mapsto \text{Trace}_{L/K}(fg)$ 是 $L$ 在 $K$ 上的非退化双线性形式。
见《域》引理 \ref{fields-lemma-separable-trace-pairing}。
选取 $b_1, \ldots, b_n \in B$，使其构成 $L$ 在 $K$ 上的一组基。
由上可知 $d = \det(\text{Trace}_{L/K}(b_ib_j)) \in A$ 非零。

\medskip\noindent
令 $Y \to \Spec(B)$ 为修改，并假设 $Y$ 正规。可以找到 $U$-容许爆破
$X'$，它是 $\Spec(A)$ 的爆破，使严格变换 $Y'$（来自 $Y$）在 $X'$ 上有限；
见《平坦性进阶》引理 \ref{flat-lemma-finite-after-blowing-up}。图示为
$$
\xymatrix{
Y' \ar[d] \ar[r] & Y \ar[r] & \Spec(B) \ar[d] \\
X' \ar[rr] & & \Spec(A)
}
$$
将 $X'$ 与 $Y'$ 替换为各自的正规化后，可假设 $X'$ 与 $Y'$ 分别是
$\Spec(A)$ 与 $\Spec(B)$ 的正规修改。如此便归约到存在交换图的情形
$$
\xymatrix{
Y \ar[d]_\pi \ar[r]_-g & \Spec(B) \ar[d] \\
X \ar[r]^-f & \Spec(A)
}
$$
其中 $X$ 与 $Y$ 分别是 $\Spec(A)$ 与 $\Spec(B)$ 的正规修改，且 $\pi$ 有限。

\medskip\noindent
$L$ 在 $K$ 上的迹映射延拓为 $\mathcal{O}_X$-模映射
$\text{Trace} : \pi_*\mathcal{O}_Y \to \mathcal{O}_X$。考虑映射
$$
\Phi : \pi_*\mathcal{O}_Y \longrightarrow \mathcal{O}_X^{\oplus n},\quad
s \longmapsto (\text{Trace}(b_1s), \ldots, \text{Trace}(b_ns))
$$
该映射是单射（因为它在泛点处是单射），并且有映射
$$
\mathcal{O}_X^{\oplus n} \longrightarrow \pi_*\mathcal{O}_Y,\quad
(s_1, \ldots, s_n) \longmapsto \sum b_i s_i
$$
它与 $\Phi$ 的复合以 $\text{Trace}(b_ib_j)$ 为矩阵。
因此 $\Phi$ 的余核被 $d$ 零化。由此得到正合列
$$
H^0(X, \Coker(\Phi)) \to H^1(Y, \mathcal{O}_Y) \to
H^1(X, \mathcal{O}_X)^{\oplus n}
$$
右端按假设有界，所以只须证明 $d$-挠子有界，这些挠子位于
$H^1(Y, \mathcal{O}_Y)$ 中。
这正是引理 \ref{resolve-lemma-bound-a-torsion} 与（\ref{resolve-equation-a-torsion}）的内容。
\end{proof}

\begin{lemma}
\label{resolve-lemma-regular-rational}
令 $A$ 为维数为 $2$ 的 Nagata 正则局部环。则 $A$ 定义有理奇点。
\end{lemma}

\begin{proof}
（本证明并不需要 $A$ 是 Nagata 环的假设，但我们只在 Nagata 的
$2$ 维正规局部整环情形下定义了有理奇点。）令
$X \to \Spec(A)$ 为修改，并假设 $X$ 正规。由引理
\ref{resolve-lemma-dominate-by-blowing-up-in-points}，可支配 $X$；把这个支配概形记作
$X_n$，它是以下序列的末项：
$$
X_n \to X_{n - 1} \to \ldots \to X_1 \to X_0 = \Spec(A)
$$
其各步均为沿闭点的爆破。由引理 \ref{resolve-lemma-blowup-regular}，概形
$X_i$ 正则，特别地正规（《代数》引理
\ref{algebra-lemma-regular-normal}）。由引理 \ref{resolve-lemma-exact-sequence} 有
$H^1(X, \mathcal{O}_X) \subset H^1(X_n, \mathcal{O}_{X_n})$。
所以只须证明 $H^1(X_n, \mathcal{O}_{X_n}) = 0$。再次使用引理
\ref{resolve-lemma-exact-sequence}，只须证明
$R^1(X_i \to X_{i - 1})_*\mathcal{O}_{X_i} = 0$，其中
$i = 1, \ldots, n$。这由引理 \ref{resolve-lemma-cohomology-of-blowup} 得到。
\end{proof}

\begin{lemma}
\label{resolve-lemma-bound-dualizing-implies-bound}
令 $A$ 为维数为 $2$ 且具有对偶复形 $\omega_A^\bullet$ 的局部正规 Nagata 整环。
若存在非零的 $d \in A$，使得对所有正规修改 $X \to \Spec(A)$，迹映射
$$
\Gamma(X, \omega_X) \to \omega_A
$$
的余核被 $d$ 零化，则对 $A$ 可以归约到有理奇点。
\end{lemma}

\begin{proof}
对命题中的 $X \to \Spec(A)$，需要约束 $H^1(X, \mathcal{O}_X)$。
令 $\omega_X$ 为 $X$ 的对偶模，如 Grauert--Riemenschneider 陈述中所述
（命题 \ref{resolve-proposition-Grauert-Riemenschneider}）。迹映射是《概形的对偶性》节
\ref{duality-section-trace} 中描述的映射 $Rf_*\omega_X \to \omega_A$。
由 Grauert--Riemenschneider 消没有 $Rf_*\omega_X = f_*\omega_X$，
所以迹映射确实给出映射 $\Gamma(X, \omega_X) \to \omega_A$。
由对偶性有 $Rf_*\omega_X = R\Hom_A(Rf_*\mathcal{O}_X, \omega_A)$
（这里使用了 $\omega_X[2]$ 是 $X$ 上相对于 $\omega_A[2]$ 归一化的对偶复形；
见《概形的对偶性》引理 \ref{duality-lemma-duality-bootstrap}，
或更直接地见节 \ref{duality-section-duality}，乃至例
\ref{duality-example-iso-on-RSheafHom-noetherian}）。区分三角
$$
A \to Rf_*\mathcal{O}_X \to R^1f_*\mathcal{O}_X[-1] \to A[1]
$$
经 $R\Hom_A(-, \omega_A)$ 变为短正合列
$$
0 \to f_*\omega_X \to \omega_A \to
\Ext_A^2(R^1f_*\mathcal{O}_X, \omega_A) \to 0
$$
（并且 $\Ext_A^i(R^1f_*\mathcal{O}_X, \omega_A) = 0$（当 $i \not = 2$ 时）；
这也将由下文讨论得到）。
由于 $R^1f_*\mathcal{O}_X$ 支撑在 $\{\mathfrak m\}$ 中，局部对偶定理告诉我们
$$
\Ext_A^2(R^1f_*\mathcal{O}_X, \omega_A) =
\Ext_A^0(R^1f_*\mathcal{O}_X, \omega_A[2]) =
\Hom_A(R^1f_*\mathcal{O}_X, E)
$$
是 $R^1f_*\mathcal{O}_X$ 的 Matlis 对偶（其他 Ext 群为零）；见《对偶复形》引理
\ref{dualizing-lemma-special-case-local-duality}。由 Matlis 对偶所固有的范畴等价
（《对偶复形》命题 \ref{dualizing-proposition-matlis}），若
$R^1f_*\mathcal{O}_X$ 不被 $d$ 零化，则上面的 $\Ext^2$ 也不被零化。
因此 $H^1(X, \mathcal{O}_X)$ 被 $d$ 零化。所需有界性遂由引理
\ref{resolve-lemma-bound-a-torsion} 与（\ref{resolve-equation-a-torsion}）得出。
\end{proof}

\begin{lemma}
\label{resolve-lemma-compare-differentials-dualizing}
令 $p$ 为素数。令 $A$ 为维数为 $2$、特征为 $p$ 的正则局部环。
令 $A_0 \subset A$ 为子环，使得 $\Omega_{A/A_0}$ 是秩为
$r < \infty$ 的自由模。置 $\omega_A = \Omega^r_{A/A_0}$。
若 $X \to \Spec(A)$ 是一列沿闭点爆破的结果，则存在映射
$$
\varphi_X : (\Omega^r_{X/\Spec(A_0)})^{**} \longrightarrow \omega_X
$$
延拓泛点处给定的等同。
\end{lemma}

\begin{proof}
注意 $A$ 是 Gorenstein 环（《对偶复形》引理
\ref{dualizing-lemma-regular-gorenstein}），因此可逆模 $\omega_A$ 的确可作为对偶模。
此外，引理中的任意 $X$ 都有可逆对偶模 $\omega_X$，因为 $X$ 正则
（故为 Gorenstein），且在 $A$ 上 proper；见评注
\ref{resolve-remark-dualizing-setup} 与引理 \ref{resolve-lemma-blowup-regular}。
假设已经构造映射 $\varphi_X : (\Omega^r_{X/A_0})^{**} \to \omega_X$，
并假设 $b : X' \to X$ 是沿闭点的爆破。置
$\Omega^r_X = (\Omega^r_{X/A_0})^{**}$ 以及
$\Omega^r_{X'} = (\Omega^r_{X'/A_0})^{**}$。由于
$\omega_{X'} = b^!(\omega_X)$，映射 $\Omega^r_{X'} \to \omega_{X'}$
等同于映射 $Rb_*(\Omega^r_{X'}) \to \omega_X$。见评注
\ref{resolve-remark-dualizing-setup} 与《概形的对偶性》节
\ref{duality-section-duality} 中的讨论。因此只须给出映射
$$
Rb_*(\Omega^r_{X'}) \longrightarrow \Omega^r_X
$$
层 $\Omega^r_{X'}$ 与 $\Omega^r_X$ 都可逆；见《除子》引理
\ref{divisors-lemma-reflexive-over-regular-dim-2}。考虑正合列
$$
b^*\Omega_{X/A_0} \to \Omega_{X'/A_0} \to \Omega_{X'/X} \to 0
$$
局部计算表明 $\Omega_{X'/X}$ 同构于例外除子 $E$ 上的可逆模；见引理
\ref{resolve-lemma-differentials-of-blowup}。由此得到以下两种情形之一：
$$
\Omega^r_{X'} \cong (b^*\Omega^r_X)(E)
\quad\text{或}\quad
\Omega^r_{X'} \cong b^*\Omega^r_X
$$
见《除子》引理 \ref{divisors-lemma-wedge-product-ses}。
（第二种可能在特征零时从不发生，但在特征 $p$ 时可能发生。）
在两种情形下，由引理 \ref{resolve-lemma-cohomology-of-blowup} 都有
$R^1b_*(\Omega^r_{X'}) = 0$ 且 $b_*(\Omega^r_{X'}) = \Omega^r_X$。
\end{proof}

\begin{lemma}
\label{resolve-lemma-go-up-degree-p}
令 $p$ 为素数。令 $A$ 为维数为 $2$、特征为 $p$ 的完备正则局部环。
令 $L/K$ 为次数为 $p$ 的不可分扩张，其中 $K$ 是 $A$ 的分式域。
令 $B \subset L$ 为 $A$ 的整闭包。则对 $B$ 可以归约到有理奇点。
\end{lemma}

\begin{proof}
有 $A = k[[x, y]]$。写 $L = K[x]/(x^p - f)$，其中 $f \in A$，
并以 $g \in B$ 表示 $x$ 的同余类，
即满足 $g^p = f$ 的元素。由《代数》引理
\ref{algebra-lemma-derivative-zero-pth-power} 可知，$\text{d}f$ 在
$\Omega_{K/\mathbf{F}_p}$ 中非零。由《代数进阶》引理
\ref{more-algebra-lemma-power-series-ring-subfields}，存在子域
$k^p \subset k' \subset k$，满足 $p^e = [k : k'] < \infty$，使得
$\text{d}f$ 在 $\Omega_{K/K_0}$ 中非零；这里 $K_0$ 是
$A_0 = k'[[x^p, y^p]] \subset A$ 的分式域。于是
$$
\Omega_{A/A_0} =
A \otimes_k \Omega_{k/k'} \oplus A \text{d}x \oplus A \text{d}y
$$
是秩为 $e + 2$ 的有限自由模。置 $\omega_A = \Omega^{e + 2}_{A/A_0}$。
考虑典范映射
$$
\text{Tr} :
\Omega^{e + 2}_{B/A_0}
\longrightarrow
\Omega^{e + 2}_{A/A_0} = \omega_A
$$
它来自引理 \ref{resolve-lemma-trace-extends}。由对偶性，这确定了映射
$$
c : \Omega^{e + 2}_{B/A_0} \to \omega_B = \Hom_A(B, \omega_A)
$$
断言：$c$ 的余核被 $B$ 的某个非零元素零化。

\medskip\noindent
由于 $\text{d}f$ 在 $\Omega_{A/A_0}$ 中非零，可以找到
$\eta_1, \ldots, \eta_{e + 1} \in \Omega_{A/A_0}$，使得
$\theta = \eta_1 \wedge \ldots \wedge \eta_{e + 1} \wedge \text{d}f$ 在
$\omega_A = \Omega^{e + 2}_{A/A_0}$ 中非零。为证明该断言，我们将构造元素
$\omega_i$，它们属于 $\Omega^{e + 2}_{B/A_0}$，其中
$i = 0, \ldots, p - 1$，它们映到
$\varphi_i \in \omega_B = \Hom_A(B, \omega_A)$，并满足
$\varphi_i(g^j) = \delta_{ij}\theta$，其中 $j = 0, \ldots, p - 1$。
由于 $\{1, g, \ldots, g^{p - 1}\}$ 是 $L/K$ 的一组基，这就证明了断言。
置 $\eta = \eta_1 \wedge \ldots \wedge \eta_{e + 1}$，使得
$\theta = \eta \wedge \text{d}f$。置
$\omega_i = \eta \wedge g^{p - 1 - i}\text{d}g$。于是按构造有
$$
\varphi_i(g^j) = \text{Tr}(g^j \eta \wedge g^{p - 1 - i}\text{d}g) =
\text{Tr}(\eta \wedge g^{p - 1 - i + j}\text{d}g) = \delta_{ij} \theta
$$
这来自引理 \ref{resolve-lemma-trace-higher} 中迹映射的显式描述。

\medskip\noindent
令 $Y \to \Spec(B)$ 为正规修改。完全如引理
\ref{resolve-lemma-go-up-separable} 的证明中一样，可以归约到 $Y$ 在某个修改
$X$ 上有限的情形，该修改属于 $\Spec(A)$。由引理
\ref{resolve-lemma-dominate-by-blowing-up-in-points}，甚至可以假设
$X \to \Spec(A)$ 是一列沿闭点爆破的结果。图示如下：
$$
\xymatrix{
Y \ar[d]_\pi \ar[r]_-g & \Spec(B) \ar[d] \\
X \ar[r]^-f & \Spec(A)
}
$$
可将引理 \ref{resolve-lemma-trace-extends} 应用于 $\pi$，从而得到下式的第一个箭头：
$$
\pi_*(\Omega^{e + 2}_{Y/A_0})
\xrightarrow{\text{Tr}}
(\Omega^{e + 2}_{X/A_0})^{**}
\xrightarrow{\varphi_X}
\omega_X
$$
第二个箭头来自引理 \ref{resolve-lemma-compare-differentials-dualizing}
（因为 $f$ 是一列沿闭点爆破）。由有限态射 $\pi$ 的对偶性，
这对应于映射
$$
c_Y : \Omega^{e + 2}_{Y/A_0} \longrightarrow \omega_Y
$$
它延拓上面的映射 $c$。因此，$\Gamma(Y, \omega_Y) \to \omega_B$ 的像
包含 $c$ 的像。由该断言可知，其余核被 $B$ 的某个固定非零元素零化。
最后应用引理 \ref{resolve-lemma-bound-dualizing-implies-bound}。
\end{proof}






\section{有理奇点}
\label{resolve-section-rational-singularities}

\noindent
本节把有理奇点的情形归约到 Gorenstein 有理奇点的情形。
见 \cite{Lipman-rational} 与 \cite{Mattuck}。

\begin{situation}
\label{resolve-situation-rational}
这里 $(A, \mathfrak m, \kappa)$ 是维数为 $2$、定义有理奇点的局部正规 Nagata 整环。
令 $s$ 为 $S = \Spec(A)$ 的闭点，并令 $U = S \setminus \{s\}$。
令 $f : X \to S$ 为修改，并假设 $X$ 正规。
记 $C_1, \ldots, C_r$ 为特殊纤维 $X_s$ 的不可约分支；该纤维来自 $f$。
\end{situation}

\begin{lemma}
\label{resolve-lemma-globally-generated}
在情形 \ref{resolve-situation-rational} 中，令 $\mathcal{F}$ 为拟凝聚
$\mathcal{O}_X$-模。则
\begin{enumerate}
\item $H^p(X, \mathcal{F}) = 0$（当 $p \not \in \{0, 1\}$ 时）；
\item $H^1(X, \mathcal{F}) = 0$，如果 $\mathcal{F}$ 由整体截面生成。
\end{enumerate}
\end{lemma}

\begin{proof}
第 (1) 项由《概形的上同调》引理
\ref{coherent-lemma-higher-direct-images-zero-above-dimension-fibre} 得出。
若 $\mathcal{F}$ 由整体截面生成，则存在满射
$\bigoplus_{i \in I} \mathcal{O}_X \to \mathcal{F}$。由第 (1) 项与上同调长正合列，
这在 $H^1$ 上诱导满射。由于 $H^1(X, \mathcal{O}_X) = 0$
（因为 $S$ 有有理奇点），而且 $H^1(X, -)$ 与直和交换
（《上同调》引理 \ref{cohomology-lemma-quasi-separated-cohomology-colimit}），
结论成立。
\end{proof}

\begin{lemma}
\label{resolve-lemma-sections-powers-I-rational}
在情形 \ref{resolve-situation-rational} 中，假设 $E = X_s$ 是有效 Cartier 除子。
令 $\mathcal{I}$ 为 $E$ 的理想层。则
$H^0(X, \mathcal{I}^n) = \mathfrak m^n$ 且
$H^1(X, \mathcal{I}^n) = 0$。
\end{lemma}

\begin{proof}
有 $H^0(X, \mathcal{O}_X) = A$；见情形 \ref{resolve-situation-vanishing} 后的讨论。
于是 $\mathfrak m \subset H^0(X, \mathcal{I}) \subset H^0(X, \mathcal{O}_X)$。
由于 $X_s \not = \emptyset$，第二个包含并非等号。因此
$H^0(X, \mathcal{I}) = \mathfrak m$。由于
$\mathcal{I}^n = \mathfrak m^n\mathcal{O}_X$，引理
\ref{resolve-lemma-globally-generated} 表明 $H^1(X, \mathcal{I}^n) = 0$。

\medskip\noindent
选取生成元 $x_1, \ldots, x_{\mu + 1}$，用来生成 $\mathfrak m$。
它们定义 $\mathcal{I}$ 的整体截面并生成该层。因此有短正合列
$$
0 \to \mathcal{F} \to \mathcal{O}_X^{\oplus \mu + 1} \to \mathcal{I} \to 0
$$
于是 $\mathcal{F}$ 是有限局部自由 $\mathcal{O}_X$-模，秩为 $\mu$，
且由《构造》引理 \ref{constructions-lemma-globally-generated-omega-twist-1}，
$\mathcal{F} \otimes \mathcal{I}$ 由整体截面生成。因此
$\mathcal{F} \otimes \mathcal{I}^n$ 对所有 $n \geq 1$ 都由整体截面生成。
所以当 $n \geq 2$ 时，可考虑正合列
$$
0 \to \mathcal{F} \otimes \mathcal{I}^{n - 1} \to
(\mathcal{I}^{n - 1})^{\oplus \mu + 1} \to
\mathcal{I}^n \to 0
$$
应用上同调长正合列，并使用引理 \ref{resolve-lemma-globally-generated} 给出的
$H^1(X, \mathcal{F} \otimes \mathcal{I}^{n - 1}) = 0$，得到
$H^0(X, \mathcal{I}^n)$ 的每个元素都具有形式 $\sum x_i a_i$，
其中某些 $a_i \in H^0(X, \mathcal{I}^{n - 1})$。归纳可得
$H^0(X, \mathcal{I}^n) = \mathfrak m^n$。
\end{proof}

\begin{lemma}
\label{resolve-lemma-blow-up-normal-rational}
在情形 \ref{resolve-situation-rational} 中，$\Spec(A)$ 沿 $\mathfrak m$ 的爆破是正规的。
\end{lemma}

\begin{proof}
令 $X' \to \Spec(A)$ 为该爆破，换言之，
$$
X' = \text{Proj}(A \oplus \mathfrak m \oplus \mathfrak m^2 \oplus \ldots).
$$
是 Rees 代数的 Proj。这尤其表明 $X'$ 是整概形，且
$X' \to \Spec(A)$ 是射影修改。令 $X$ 为 $X'$ 的正规化。
由于 $A$ 是 Nagata 环，$\nu : X \to X'$ 是有限态射
（《态射》引理 \ref{morphisms-lemma-nagata-normalization}）。
令 $E' \subset X'$ 为例外除子，并令 $E \subset X$ 为其逆像。
令 $\mathcal{I}' \subset \mathcal{O}_{X'}$ 与
$\mathcal{I} \subset \mathcal{O}_X$ 为它们的理想层。
回忆 $\mathcal{I}' = \mathcal{O}_{X'}(1)$（《除子》引理
\ref{divisors-lemma-blowing-up-projective}）。注意
$\mathcal{I} = \nu^*\mathcal{I}'$，且 $E$ 是有效 Cartier 除子
（《除子》引理 \ref{divisors-lemma-pullback-effective-Cartier-defined}）。
我们要证明 $\nu$ 是同构。由于 $\nu$ 有限，只须证明
$\mathcal{O}_{X'} \to \nu_*\mathcal{O}_X$ 是同构。否则可找到 $n \geq 0$，使得
$$
H^0(X', (\mathcal{I}')^n) \not =
H^0(X', (\nu_*\mathcal{O}_X) \otimes (\mathcal{I}')^n)
$$
例如，这是因为可以从关联分次模恢复拟凝聚 $\mathcal{O}_{X'}$-模；
见《性质》引理 \ref{properties-lemma-ample-quasi-coherent}。
由投影公式有
$$
H^0(X', (\nu_*\mathcal{O}_X) \otimes (\mathcal{I}')^n) =
H^0(X, \nu^*(\mathcal{I}')^n) =
H^0(X, \mathcal{I}^n) = \mathfrak m^n
$$
最后一个等号来自引理 \ref{resolve-lemma-sections-powers-I-rational}。
另一方面，显然有单射 $\mathfrak m^n \to H^0(X', (\mathcal{I}')^n)$。
由于 $H^0(X', (\mathcal{I}')^n)$ 无挠，可知等号对所有 $n$ 都成立，
故 $X = X'$。
\end{proof}

\begin{lemma}
\label{resolve-lemma-cohomology-blow-up-rational}
在情形 \ref{resolve-situation-rational} 中，令 $X$ 为 $\Spec(A)$ 沿 $\mathfrak m$ 的爆破。
令 $E \subset X$ 为例外除子。照常置 $\mathcal{O}_X(1) = \mathcal{I}$，并置
$\mathcal{O}_E(1) = \mathcal{O}_X(1)|_E$。则有
\begin{enumerate}
\item $E$ 是 $\kappa$ 上的 proper Cohen--Macaulay 曲线；
\item $\mathcal{O}_E(1)$ 非常丰沛；
\item $\deg(\mathcal{O}_E(1)) \geq 1$，且等号仅当 $A$ 为正则局部环时成立；
\item $H^1(E, \mathcal{O}_E(n)) = 0$，其中 $n \geq 0$；
\item $H^0(E, \mathcal{O}_E(n)) = \mathfrak m^n/\mathfrak m^{n + 1}$，
其中 $n \geq 0$。
\end{enumerate}
\end{lemma}

\begin{proof}
由于按构造 $\mathcal{O}_X(1)$ 非常丰沛，可知其在特殊纤维 $E$ 上的限制也非常丰沛。
由引理 \ref{resolve-lemma-blow-up-normal-rational}，概形 $X$ 正规。于是由《除子》引理
\ref{divisors-lemma-normal-effective-Cartier-divisor-S1}，$E$ 是 Cohen--Macaulay 的。
引理 \ref{resolve-lemma-sections-powers-I-rational} 可用，并且由正合列
$$
0 \to \mathcal{I}^{n + 1} \to \mathcal{I}^n \to i_*\mathcal{O}_E(n) \to 0
$$
以及上同调长正合列得到第 (4)、(5) 项。特别地，有
$$
\deg(\mathcal{O}_E(1)) = \chi(E, \mathcal{O}_E(1)) - \chi(E, \mathcal{O}_E) =
\dim(\mathfrak m/\mathfrak m^2) - 1
$$
这来自《簇》定义 \ref{varieties-definition-degree-invertible-sheaf}。
因而第 (3) 项也成立。
\end{proof}

\begin{lemma}
\label{resolve-lemma-dualizing-rational}
在情形 \ref{resolve-situation-rational} 中，假设 $A$ 有对偶复形
$\omega_A^\bullet$。令 $\omega_X$ 为 $X$ 的对偶模，则迹映射
$H^0(X, \omega_X) \to \omega_A$ 是同构；因此存在典范映射
$f^*\omega_A \to \omega_X$。
\end{lemma}

\begin{proof}
由 Grauert--Riemenschneider 消没（命题
\ref{resolve-proposition-Grauert-Riemenschneider}）有
$Rf_*\omega_X = f_*\omega_X$。由对偶性得到短正合列
$$
0 \to f_*\omega_X \to \omega_A \to
\Ext^2_A(R^1f_*\mathcal{O}_X, \omega_A) \to 0
$$
（例如见引理 \ref{resolve-lemma-bound-dualizing-implies-bound} 的证明），
而由于 $A$ 定义有理奇点，得到 $f_*\omega_X = \omega_A$。
\end{proof}

\begin{lemma}
\label{resolve-lemma-dualizing-blow-up-rational}
在情形 \ref{resolve-situation-rational} 中，假设 $A$ 有对偶复形
$\omega_A^\bullet$ 且不正则。令 $X$ 为 $\Spec(A)$ 沿 $\mathfrak m$ 的爆破，
例外除子为 $E \subset X$。令 $\omega_X$ 为 $X$ 的对偶模。则
\begin{enumerate}
\item $\omega_E = \omega_X|_E \otimes \mathcal{O}_E(-1)$；
\item $H^1(X, \omega_X(n)) = 0$，其中 $n \geq 0$；
\item 引理 \ref{resolve-lemma-dualizing-rational} 中的映射
$f^*\omega_A \to \omega_X$ 是满射。
\end{enumerate}
\end{lemma}

\begin{proof}
以下不再特别提及引理 \ref{resolve-lemma-cohomology-blow-up-rational} 的结果。
由《概形的对偶性》引理
\ref{duality-lemma-sheaf-with-exact-support-effective-Cartier} 和
\ref{duality-lemma-twisted-inverse-image-closed}，注意到
$\omega_E = \omega_X|_E \otimes \mathcal{O}_E(-1)$。因此
$\omega_X|_E = \omega_E(1)$。考虑短正合列
$$
0 \to \omega_X(n + 1) \to \omega_X(n) \to i_*\omega_E(n + 1) \to 0
$$
由《代数曲线》引理 \ref{curves-lemma-vanishing-twist} 可知
$H^1(E, \omega_E(n + 1)) = 0$（当 $n \geq 0$ 时）。因此映射
$$
\ldots \to H^1(X, \omega_X(2)) \to H^1(X, \omega_X(1)) \to H^1(X, \omega_X)
$$
均为满射。由于 $H^1(X, \omega_X(n))$ 在 $n \gg 0$ 时为零
（《概形的上同调》引理 \ref{coherent-lemma-kill-by-twisting}），
可知第 (2) 项成立。

\medskip\noindent
由《代数曲线》引理
\ref{curves-lemma-tensor-omega-with-globally-generated-invertible} 可知
$\omega_X|_E = \omega_E \otimes \mathcal{O}_E(1)$ 由整体截面生成。
由于上面已经看到 $H^1(X, \omega_X(1)) = 0$，映射
$H^0(X, \omega_X) \to H^0(E, \omega_X|_E)$ 是满射。因此
$\omega_X$ 由整体截面生成，所以第 (3) 项成立，因为引理
\ref{resolve-lemma-dualizing-rational} 使用 $\Gamma(X, \omega_X) = \omega_A$ 定义该映射。
\end{proof}

\begin{lemma}
\label{resolve-lemma-rational-to-gorenstein}
令 $(A, \mathfrak m, \kappa)$ 为维数为 $2$、定义有理奇点的局部正规 Nagata 整环。
假设 $A$ 有对偶复形。则存在一列有限的沿奇异闭点爆破
$$
X = X_n \to X_{n - 1} \to \ldots \to X_1 \to X_0 = \Spec(A)
$$
使得 $X_i$ 对每个 $i$ 都正规，并且对偶层 $\omega_X$（即 $X$ 的对偶层）是可逆
$\mathcal{O}_X$-模。
\end{lemma}

\begin{proof}
对偶模 $\omega_A$ 是有限 $A$-模，其泛点处的茎可逆。事实上，
$\omega_A \otimes_A K$ 是分式域 $K$ 的对偶模，而该域是 $A$ 的分式域，故秩为 $1$。
因此存在爆破 $b : Y \to \Spec(A)$，使 $\omega_A$ 关于 $b$ 的严格变换为可逆
$\mathcal{O}_Y$-模；见《除子》引理
\ref{divisors-lemma-blowup-fitting-ideal}。由引理
\ref{resolve-lemma-dominate-by-normalized-blowing-up}，可以选取一列正规化爆破
$$
X_n \to X_{n - 1} \to \ldots \to X_1 \to \Spec(A)
$$
使得 $X_n$ 支配 $Y$。由引理 \ref{resolve-lemma-blow-up-normal-rational} 并归纳可知，
每个 $X_i \to X_{i - 1}$ 就是爆破。

\medskip\noindent
我们断言 $\omega_{X_n}$ 可逆。由于 $\omega_{X_n}$ 是凝聚
$\mathcal{O}_{X_n}$-模，只须说明它的茎是可逆模。若 $x \in X_n$ 是正则点，
则由正则概形为 Gorenstein 概形这一事实可知结论显然
（《对偶复形》引理 \ref{dualizing-lemma-regular-gorenstein}）。
若 $x$ 是 $X_n$ 的奇异点，则每个像 $x_i \in X_i$（即 $x$ 的像）都是奇异点
（因为由引理 \ref{resolve-lemma-blowup-regular}，正则点的爆破仍正则）。
考虑引理 \ref{resolve-lemma-dualizing-rational} 中的典范映射
$f_n^*\omega_A \to \omega_{X_n}$。对每个 $i$，态射
$X_{i + 1} \to X_i$ 或者是 $x_i$ 的爆破，或者在 $x_i$ 处是同构。
由于 $x_i$ 总是奇异点，由引理 \ref{resolve-lemma-dualizing-blow-up-rational} 与归纳可知，
映射 $f_i^*\omega_A \to \omega_{X_i}$ 在 $x_i$ 处的茎上总是满射。因此
$$
(f_n^*\omega_A)_x \longrightarrow \omega_{X_n, x}
$$
是满射。另一方面，由 $b$ 的选择，$f_n^*\omega_A$ 模去其挠子模所得的商是可逆模
$\mathcal{L}$。此外，对偶模无挠（《概形的对偶性》引理
\ref{duality-lemma-dualizing-module}）。因此
$\mathcal{L}_x \cong  \omega_{X_n, x}$，证明完成。
\end{proof}






\section{形式弧}
\label{resolve-section-arcs}

\noindent
令 $X$ 为局部诺特概形。本节称 $X$ 中的{\it 形式弧}为态射
$a : T \to X$，其中 $T$ 是完备离散赋值环 $R$ 的谱，并把其剩余域
$\kappa$ 与像 $p$ 的剩余域等同；这里该像是 $\Spec(R)$ 的闭点之像。在这种情形下，称形式弧
$a$ {\it 以 $p$ 为中心}。称形式弧 $T \to X$ {\it 非奇异}，如果所诱导的映射
$\mathfrak m_p/\mathfrak m_p^2 \to \mathfrak m_R/\mathfrak m_R^2$ 为满射。

\medskip\noindent
令 $a : T \to X$、$T = \Spec(R)$ 为以闭点 $p$ 为中心的非奇异形式弧，且该点属于 $X$。
假设 $X$ 局部诺特。令 $b : X_1 \to X$ 为 $X$ 沿 $x$ 的爆破。
由于 $a$ 非奇异，存在元素 $f \in \mathfrak m_p$，它映到 $R$ 中的一个
一致化参数。特别地，$T$ 的泛点映到 $X$ 中不等于 $p$ 的点。
换言之，令 $K$ 为 $R$ 的分式域，则 $a$ 的限制定义态射
$\Spec(K) \to X \setminus \{p\}$。由于态射 $b$ 为固有态射，且在
$X \setminus \{x\}$ 上为同构，可应用固有性的赋值判别准则，得到唯一态射
$a_1$，使下图交换：
$$
\xymatrix{
T \ar[r]_{a_1} \ar[rd]_a & X_1 \ar[d]^{b} \\
& X
}
$$
令 $p_1 \in X_1$ 为 $T$ 的闭点之像。注意 $p_1$ 是闭点，因为它是
$\kappa = \kappa(p)$-有理点，并位于 $X_1 \to X$ 在 $x$ 上方的纤维中。
由于有分解
$$
\mathcal{O}_{X, x} \to \mathcal{O}_{X_1, p_1} \to R
$$
可知 $a_1$ 也是非奇异形式弧。

\medskip\noindent
可以重复这一过程，得到一列爆破
$$
\xymatrix{
T \ar[d]_a \ar[rd]_{a_1} \ar[rrd]_{a_2} \ar[rrrd]^{a_3} \\
(X, p) & (X_1, p_1) \ar[l] & (X_2, p_2) \ar[l] &
(X_3, p_3) \ar[l] & \ldots \ar[l]
}
$$
这类爆破序列可以刻画如下。

\begin{lemma}
\label{resolve-lemma-sequence-blowups}
令 $X$ 为局部诺特概形。设
$$
(X, p) = (X_0, p_0) \leftarrow (X_1, p_1) \leftarrow (X_2, p_2) \leftarrow
(X_3, p_3) \leftarrow \ldots
$$
为一列爆破，并满足：
\begin{enumerate}
\item $p_i$ 是闭点，映到 $p_{i - 1}$，且
$\kappa(p_i) = \kappa(p_{i - 1})$；
\item 存在 $x_1 \in \mathfrak m_p$，其在 $\mathfrak m_{p_i}$ 中的像
（$i > 0$）定义例外除子 $E_i \subset X_i$。
\end{enumerate}
则该序列由上述非奇异弧 $a : T \to X$ 得到。
\end{lemma}

\begin{proof}
记 $\mathcal{O}_n = \mathcal{O}_{X_n, p_n}$ 以及
$\mathcal{O} = \mathcal{O}_{X, p}$。以 $\mathfrak m \subset \mathcal{O}$ 和
$\mathfrak m_n \subset \mathcal{O}_n$ 表示极大理想。

\medskip\noindent
我们断言 $x_1^t \not \in \mathfrak m_n^{t + 1}$。
事实上，若不然，则在局部环 $\mathcal{O}_{n + 1}$ 中，元素 $x_1^t$
属于 $(t + 1)E_{n + 1}$ 的理想。这与 $x_1$ 定义 $E_{n + 1}$ 的假设矛盾。

\medskip\noindent
对每个 $n$，选取生成元 $y_{n, 1}, \ldots, y_{n, t_n}$ 来生成
$\mathfrak m_n$。由假设 (2) 有
$\mathfrak m_n \mathcal{O}_{n + 1} = x_1\mathcal{O}_{n + 1}$，故可写成
$y_{n, i} = a_{n, i} x_1$，其中某个
$a_{n, i} \in \mathcal{O}_{n + 1}$。由于映射
$\mathcal{O}_n \to \mathcal{O}_{n + 1}$ 按 (1) 在剩余域上定义同构，
可选取 $c_{n, i} \in \mathcal{O}_n$，使其剩余类与 $a_{n, i}$ 相同。
于是
$$
\mathfrak m_n = (x_1, z_{n, 1}, \ldots, z_{n, t_n}),
\quad z_{n, i} = y_{n, i} - c_{n, i} x_1
$$
而元素 $z_{n, i}$ 映到 $\mathfrak m_{n + 1}^2$ 的元素，
这一陈述在 $\mathcal{O}_{n + 1}$ 中成立。

\medskip\noindent
考虑
$$
J_n = \Ker(\mathcal{O} \to \mathcal{O}_n/\mathfrak m_n^{n + 1})
$$
我们断言 $\mathcal{O}/J_n$ 的长度为 $n + 1$，并且
$\mathcal{O}/(x_1) + J_n$ 等于剩余域。当 $n = 0$ 时这是直接的。
假设该断言对 $n$ 成立。令 $f \in J_n$。于是在 $\mathcal{O}_n$ 中有
$$
f = a x_1^{n + 1} + x_1^n A_1(z_{n, i}) +
x_1^{n - 1} A_2(z_{n, i}) + \ldots + A_{n + 1}(z_{n, i})
$$
其中 $a \in \mathcal{O}_n$，而 $A_i$ 是次数为 $i$、系数属于
$\mathcal{O}_n$ 的齐次式。由于 $\mathcal{O} \to \mathcal{O}_n$ 认同剩余域，
可选取 $a \in \mathcal{O}$（仿照上面对 $z_{n, i}$ 的构造）。
在 $\mathcal{O}_{n + 1}$ 中取像可见，$f$ 与 $a x_1^{n + 1}$
模 $\mathfrak m_{n + 1}^{n + 2}$ 有相同的像。由于
$x_n^{n + 1} \not \in \mathfrak m_{n + 1}^{n + 2}$，可知
$J_n/J_{n + 1}$ 的长度为 $1$，断言得证。

\medskip\noindent
考虑 $R = \lim \mathcal{O}/J_n$。这是 $\mathfrak m$-进完备化的一个商，
而该完备化取自 $\mathcal{O}$，故它是完备诺特局部环。另一方面，它的长度
不是有限的，且 $x_1$ 生成其极大理想。因此 $R$ 是完备离散赋值环。
映射 $\mathcal{O} \to R$ 提升为局部同态 $\mathcal{O}_n \to R$，
对每个 $n$ 均如此。有两种方法证明这一点：(1) 对每个 $n$，可用类似过程
构造 $\mathcal{O}_n \to R_n$，然后证明 $\mathcal{O} \to \mathcal{O}_n \to R_n$
经由同构 $R \to R_n$ 分解；或者 (2) 可用除子章引理
\ref{divisors-lemma-characterize-affine-blowup} 证明 $\mathcal{O}_n$ 是反复仿射
爆破代数的一个局部化，从而显式构造映射 $\mathcal{O}_n \to R$。
由此显然，这列爆破来自非奇异弧 $a : T = \Spec(R) \to X$。
\end{proof}

\noindent
下述引理是一类内龙（N\'eron）消奇异化引理。

\begin{lemma}
\label{resolve-lemma-sequence-blowups-along-arc-becomes-nonsingular}
令 $(A, \mathfrak m, \kappa)$ 为维数 $2$ 的诺特局部整环。令
$A \to R$ 为到一个完备离散赋值环的满射。这定义了非奇异弧
$a : T = \Spec(R) \to \Spec(A)$。令
$$
\Spec(A) = X_0 \leftarrow X_1 \leftarrow X_2 \leftarrow X_3 \leftarrow \ldots
$$
为由 $a$ 构造的爆破序列。若 $A_\mathfrak p$ 是正则局部环，其中
$\mathfrak p = \Ker(A \to R)$，则对某个 $i$，概形 $X_i$ 在 $x_i$ 处正则。
\end{lemma}

\begin{proof}
令 $x_1 \in \mathfrak m$ 映到 $R$ 的一个一致化参数。注意
$\kappa(\mathfrak p) = K$ 是 $R$ 的分式域。把
$\mathfrak p = (x_2, \ldots, x_r)$ 写成生成元个数 $r$ 最少的形式。
若 $r = 2$，则 $\mathfrak m = (x_1, x_2)$，且 $A$ 正则，引理成立。
假设 $r > 2$。必要时重新编号，可设 $x_2$ 映到 $A_\mathfrak p$ 的一个
一致化参数。于是 $\mathfrak p/\mathfrak p^2 + (x_2)$ 被 $x_1$ 的某个幂湮灭。
对 $i > 2$，可找到 $n_i \geq 0$ 以及 $a_i \in A$，使得
$$
x_1^{n_i} x_i - a_i x_2 = \sum\nolimits_{2 \leq j \leq k} a_{jk} x_jx_k
$$
其中 $a_{jk} \in A$。若出现 $n_i = 0$，其指标为某个 $i$，则可去掉
$x_i$ 这个属于 $\mathfrak p$ 的生成元，并对 $r$ 作归纳而得证。
若对某个 $i$，元素 $a_i$ 是单位，则可去掉 $x_2$ 这个属于
$\mathfrak p$ 的生成元，并以同样方式得证。因此或者 $a_i \in \mathfrak p$，或者
$a_i = u_i x_1^{m_1} \bmod \mathfrak p$，其中某个 $m_1 > 0$，且
$u_i \in A$ 为单位。因此有二者之一：
$$
x_1^{n_i} x_i = \sum\nolimits_{2 \leq j \leq k} a_{jk} x_jx_k
\quad\text{或}\quad
x_1^{n_i} x_i - u_i x_1^{m_i} x_2 =
\sum\nolimits_{2 \leq j \leq k} a_{jk} x_jx_k
$$
我们将证明，爆破之后整数 $n_i$、$m_i$ 减小，这就完成证明。

\medskip\noindent
考察这些方程在仿射爆破代数 $A' = A[\mathfrak m/x_1]$ 上的情形。
由于 $\mathfrak m = (x_1, \ldots, x_r)$，可见 $A'$ 由 $R$ 上的元素
$y_i = x_i/x_1$（$i \geq 2$）生成。显然 $A \to R$ 延拓为
$A' \to R$，其核为 $(y_2, \ldots, y_r)$。于是二者之一成立：
$$
x_1^{n_i - 1} y_i = \sum\nolimits_{2 \leq j \leq k} a_{jk} y_jy_k
\quad\text{或}\quad
x_1^{n_i - 1} y_i - u_i x_1^{m_1 - 1} y_2 =
\sum\nolimits_{2 \leq j \leq k} a_{jk} y_jy_k
$$
证明完成。
\end{proof}




\section{到完备化的基变换}
\label{resolve-section-aux}

\noindent
下述简单引理将成为后文的一个有用工具。

\begin{lemma}
\label{resolve-lemma-iso-completions}
令 $(A, \mathfrak m, \kappa)$ 为局部环，其极大理想 $\mathfrak m$ 有限生成。
令 $X$ 为 $A$ 上的概形。令
$Y = X \times_{\Spec(A)} \Spec(A^\wedge)$，其中 $A^\wedge$ 是关于
$\mathfrak m$ 的进完备化，而被完备化的环为 $A$。对点 $q \in Y$，设其像 $p \in X$ 位于
$\Spec(A)$ 的闭点上方，则局部环映射
$\mathcal{O}_{X, p} \to \mathcal{O}_{Y, q}$ 在完备化上诱导同构。
\end{lemma}

\begin{proof}
可设 $X$ 仿射，因而可写成 $X = \Spec(B)$。令
$\mathfrak q \subset B' = B \otimes_A A^\wedge$ 为对应于 $q$ 的素理想，
并令 $\mathfrak p \subset B$ 为对应于 $p$ 的素理想。由代数章引理
\ref{algebra-lemma-hathat-finitely-generated}，有
$$
B'/(\mathfrak m^\wedge)^n B' =
A^\wedge/(\mathfrak m^\wedge)^n \otimes_A B =
A/\mathfrak m^n \otimes_A B = B/\mathfrak m^n B
$$
对所有 $n$ 成立。由于 $\mathfrak m B \subset \mathfrak p$ 且
$\mathfrak m^\wedge B' \subset \mathfrak q$，可见 $B/\mathfrak p^n$ 与
$B'/\mathfrak q^n$ 都是上面所示环对同一素理想的第 $n$ 次幂取商。
引理得证。
\end{proof}

\begin{lemma}
\label{resolve-lemma-port-regularity-to-completion}
令 $(A, \mathfrak m, \kappa)$ 为诺特局部环。令 $X \to \Spec(A)$ 为局部有限型态射。
置 $Y = X \times_{\Spec(A)} \Spec(A^\wedge)$。令 $y \in Y$ 的像为 $x \in X$。则
\begin{enumerate}
\item 若 $\mathcal{O}_{Y, y}$ 正则，则 $\mathcal{O}_{X, x}$ 正则；
\item 若 $y$ 位于闭纤维中，则 $\mathcal{O}_{Y, y}$ 正则
$\Leftrightarrow \mathcal{O}_{X, x}$ 正则；
\item 若 $X$ 在 $A$ 上固有，则 $X$ 正则当且仅当 $Y$ 正则。
\end{enumerate}
\end{lemma}

\begin{proof}
由于 $A \to A^\wedge$ 忠实平坦（代数章引理
\ref{algebra-lemma-completion-faithfully-flat}），可知 $Y \to X$ 平坦。
故 (1) 由代数章引理 \ref{algebra-lemma-descent-regular} 得到。
引理 \ref{resolve-lemma-iso-completions} 表明，态射 $Y \to X$ 在特殊纤维各点的
完备局部环上诱导同构。故 (2) 由进阶代数章引理
\ref{more-algebra-lemma-completion-regular} 得到。若 $X$ 在 $A$ 上固有，
则 $Y$ 在 $A^\wedge$ 上固有（态射章引理
\ref{morphisms-lemma-base-change-proper}），并且 $X$ 与 $Y$ 的每个闭点都位于
闭纤维中。因此由性质章引理 \ref{properties-lemma-characterize-regular}，
$Y$ 为正则概形当且仅当 $X$ 正则。
\end{proof}

\begin{lemma}
\label{resolve-lemma-descend-admissible-blowup}
令 $(A, \mathfrak m)$ 为诺特局部环，其完备化为 $A^\wedge$。令
$U \subset \Spec(A)$ 与 $U^\wedge \subset \Spec(A^\wedge)$ 为相应的去心谱。
若 $Y \to \Spec(A^\wedge)$ 是 $U^\wedge$-容许爆破，则存在
$U$-容许爆破 $X \to \Spec(A)$，使得
$Y = X \times_{\Spec(A)} \Spec(A^\wedge)$。
\end{lemma}

\begin{proof}
按定义，存在理想 $J \subset A^\wedge$，使得
$V(J) = \{\mathfrak m A^\wedge\}$，且 $Y$ 是 $S^\wedge$ 沿由 $J$ 定义的
闭子概形的爆破；见除子章定义 \ref{divisors-definition-admissible-blowup}。
由于 $A^\wedge$ 诺特，这说明
$\mathfrak m^n A^\wedge \subset J$，其中 $n$ 为某个整数。由
$A^\wedge/\mathfrak m^n A^\wedge = A/\mathfrak m^n$，可找到理想
$\mathfrak m^n \subset I \subset A$，使得 $J = I A^\wedge$。
令 $X \to S$ 为沿 $I$ 的爆破。由于 $A \to A^\wedge$ 平坦，
由除子章引理 \ref{divisors-lemma-flat-base-change-blowing-up} 可知
$X$ 的基变换即为 $Y$。
\end{proof}

\begin{lemma}
\label{resolve-lemma-blowup-still-good}
令 $(A, \mathfrak m, \kappa)$ 为维数 $2$ 的 Nagata 局部正规整环。
假设 $A$ 定义有理奇点，且完备化 $A^\wedge$（所完备化的环为 $A$）正规。则
\begin{enumerate}
\item $A^\wedge$ 定义有理奇点；
\item 若 $X \to \Spec(A)$ 是沿 $\mathfrak m$ 的爆破，则对闭点
$x \in X$，完备化 $\mathcal{O}_{X, x}$ 正规。
\end{enumerate}
\end{lemma}

\begin{proof}
令 $Y \to \Spec(A^\wedge)$ 为修改，且 $Y$ 正规。需要证明
$H^1(Y, \mathcal{O}_Y) = 0$。由簇章引理
\ref{varieties-lemma-modification-normal-iso-over-codimension-1}，
$Y \to \Spec(A^\wedge)$ 在去心谱
$U^\wedge = \Spec(A^\wedge) \setminus \{\mathfrak m^\wedge\}$ 上为同构。
由引理 \ref{resolve-lemma-dominate-by-scheme-modification}，存在
$U^\wedge$-容许爆破 $Y' \to \Spec(A^\wedge)$，它支配 $Y$。由引理
\ref{resolve-lemma-descend-admissible-blowup}，存在 $U$-容许爆破
$X \to \Spec(A)$，其到 $A^\wedge$ 的基变换支配 $Y$。
由于 $A$ 是 Nagata 环，可用其正规化替换 $X$；替换后
$X \to \Spec(A)$ 是正规修改（但可能不再是 $U$-容许爆破）。
于是 $H^1(X, \mathcal{O}_X) = 0$，因为 $A$ 定义有理奇点。因此
$H^1(X \times_{\Spec(A)} \Spec(A^\wedge),
\mathcal{O}_{X \times_{\Spec(A)} \Spec(A^\wedge)}) = 0$
由平坦基变换得到（概形上同调章引理
\ref{coherent-lemma-flat-base-change-cohomology}，并用代数章引理
\ref{algebra-lemma-completion-flat} 所给的 $A \to A^\wedge$ 平坦性）。
由引理 \ref{resolve-lemma-exact-sequence} 可知 $H^1(Y, \mathcal{O}_Y) = 0$。

\medskip\noindent
最后，令 $X \to \Spec(A)$ 为 $\Spec(A)$ 沿 $\mathfrak m$ 的爆破。
则 $Y = X \times_{\Spec(A)} \Spec(A^\wedge)$ 是 $\Spec(A^\wedge)$
沿 $\mathfrak m^\wedge$ 的爆破。由引理 \ref{resolve-lemma-blow-up-normal-rational}，
$Y$ 与 $X$ 都正规。另一方面，$A^\wedge$ 优良（进阶代数章命题
\ref{more-algebra-proposition-ubiquity-excellent}），故 $Y$ 的每个仿射开子集
都是某个优良正规整环的谱（进阶代数章引理
\ref{more-algebra-lemma-finite-type-over-excellent}）。因此对 $y \in Y$，环映射
$\mathcal{O}_{Y, y} \to \mathcal{O}_{Y, y}^\wedge$ 正则；由进阶代数章引理
\ref{more-algebra-lemma-normal-goes-up} 可知 $\mathcal{O}_{Y, y}^\wedge$ 正规。
若 $x \in X$ 是特殊纤维的闭点，则存在唯一闭点 $y \in Y$，它位于 $x$ 上方。
由于 $\mathcal{O}_{X, x} \to \mathcal{O}_{Y, y}$ 在完备化上诱导同构
（引理 \ref{resolve-lemma-iso-completions}），结论成立。
\end{proof}

\begin{lemma}
\label{resolve-lemma-formally-unramified}
令 $(A, \mathfrak m)$ 为诺特局部环，令 $X$ 为 $A$ 上的概形。假设
\begin{enumerate}
\item $A$ 解析非分歧（代数章定义
\ref{algebra-definition-analytically-unramified}）；
\item $X$ 在 $A$ 上局部有限型；
\item $X \to \Spec(A)$ 在 $X$ 的各不可约分支的泛点处为 \'etale 态射。
\end{enumerate}
则 $X$ 的正规化在 $X$ 上有限。
\end{lemma}

\begin{proof}
由于 $A$ 解析非分歧，由代数章引理
\ref{algebra-lemma-analytically-unramified-easy} 可知它既约。由于 $X$ 的
正规化只依赖于 $X$ 的既约化，可替换 $X$，替换物为 $X_{red}$；注意
$X_{red} \to X$ 在开集 $U$ 上为同构，其中 $X \to \Spec(A)$ 为 \'etale 态射，
因为 $U$ 既约（下降章引理 \ref{descent-lemma-reduced-local-smooth}）。
故替换后条件 (3) 仍成立。此外，可设仿射概形 $X = \Spec(B)$。

\medskip\noindent
映射
$$
K = \prod\nolimits_{\mathfrak p \subset A\text{ minimal}} \kappa(\mathfrak p)
\longrightarrow
K^\wedge = \prod\nolimits_{\mathfrak p^\wedge \subset A^\wedge\text{ minimal}}
\kappa(\mathfrak p^\wedge)
$$
是单射，因为 $A \to A^\wedge$ 忠实平坦（代数章引理
\ref{algebra-lemma-completion-faithfully-flat}），从而在极小素理想集合之间诱导满射
（对平坦环映射使用降链性质；见代数章第
\ref{algebra-section-going-up} 节）。由于这里的环诺特，两边都是有限个域的乘积。
令 $L = \prod_{\mathfrak q \subset B\text{ minimal}} \kappa(\mathfrak q)$。
假设 (3) 蕴含 $L = B \otimes_A K$，且 $K \to L$ 是有限 \'etale 环映射
（这是因为 $A \to B$ 泛有限；例如使用代数章引理
\ref{algebra-lemma-generically-finite}，或态射章第
\ref{morphisms-section-generically-finite} 节中更详细的结果）。
由于 $B$ 既约，可知 $B \subset L$。这蕴含
$$
C = B \otimes_A A^\wedge \subset
L \otimes_A A^\wedge = L \otimes_K K^\wedge = M
$$
于是 $M$ 是 $C$ 的全分式环；作为 $K^\wedge$ 上的有限可分代数，
它是有限个域的乘积。因此 $C$ 既约，且其正规化 $C'$ 是 $C$ 在 $M$ 中的
整闭包。正规化 $B'$（被正规化的环为 $B$）是 $B$ 在 $L$ 中的整闭包。由
$A \to A^\wedge$ 的平坦性，得到单射 $B' \otimes_A A^\wedge \to M$，
其像包含在 $C'$ 中。图示为
$$
B' \otimes_A A^\wedge \longrightarrow C'
$$
由于 $A^\wedge$ 是 Nagata 环（代数章引理
\ref{algebra-lemma-Noetherian-complete-local-Nagata}），可知 $C'$ 在
$C = B \otimes_A A^\wedge$ 上有限；见代数章引理
\ref{algebra-lemma-Noetherian-complete-local-Nagata} 与
\ref{algebra-lemma-nagata-in-reduced-finite-type-finite-integral-closure}。
由于 $C$ 诺特，可知 $B' \otimes_A A^\wedge$ 在
$C = B \otimes_A A^\wedge$ 上有限。因此由忠实平坦下降（代数章引理
\ref{algebra-lemma-descend-properties-modules}），$B'$ 在 $B$ 上有限，正是所需结论。
\end{proof}

\begin{lemma}
\label{resolve-lemma-normalization-completion}
令 $(A, \mathfrak m, \kappa)$ 为诺特局部环。令 $X \to \Spec(A)$ 为局部有限型态射。
置 $Y = X \times_{\Spec(A)} \Spec(A^\wedge)$。若 $Y$ 中特殊纤维的补集正规，
则正规化 $X^\nu \to X$ 有限，且 $X^\nu$ 到 $\Spec(A^\wedge)$ 的基变换
给出 $Y$ 的正规化。
\end{lemma}

\begin{proof}
立即归约到仿射情形 $X = \Spec(B)$，其中 $B$ 是有限型 $A$-代数。
置 $C = B \otimes_A A^\wedge$，从而 $Y = \Spec(C)$。由于
$A \to A^\wedge$ 忠实平坦，对任意素理想 $\mathfrak q \subset B$，
存在素理想 $\mathfrak r \subset C$，它位于 $\mathfrak q$ 上方。
于是 $B_\mathfrak q \to C_\mathfrak r$ 忠实平坦。因此，若
$\mathfrak q$ 不位于 $\mathfrak m$ 上方，则 $C_\mathfrak r$ 正规，
这是关于 $Y$ 的假设所给；再由代数章引理 \ref{algebra-lemma-descent-normal}，
$B_\mathfrak q$ 正规。由此可见 $X$ 在特殊纤维之外正规。

\medskip\noindent
回忆完备诺特局部环 $A^\wedge$ 是 Nagata 环（代数章引理
\ref{algebra-lemma-Noetherian-complete-local-Nagata}）。因此正规化
$Y^\nu \to Y$ 有限（态射章引理 \ref{morphisms-lemma-nagata-normalization}），
并在特殊纤维之外为同构。记 $Y^\nu = \Spec(C')$。于是 $C \to C'$ 有限，
并在 $V(\mathfrak m C)$ 之外为同构。由于 $B \to C$ 平坦且诱导同构
$B/\mathfrak m B \to C/\mathfrak m C$，存在有限环映射 $B \to B'$，
其到 $C$ 的基变换给出 $C \to C'$；见进阶代数章引理
\ref{more-algebra-lemma-application-formal-glueing} 与注
\ref{more-algebra-remark-formal-glueing-algebras}。
因此得到有限态射 $X' \to X$，它在特殊纤维之外为同构，且其基变换给出
$Y^\nu \to Y$。由第一段的讨论，$X'$ 在不属于特殊纤维的点处正规。
对特殊纤维上的点 $x \in X'$，有相应的点 $y \in Y^\nu$ 和平坦映射
$\mathcal{O}_{X', x} \to \mathcal{O}_{Y^\nu, y}$。由于
$\mathcal{O}_{Y^\nu, y}$ 正规，$\mathcal{O}_{X', x}$ 也正规；见代数章引理
\ref{algebra-lemma-descent-normal}。因此 $X'$ 正规，从而它是 $X$ 的正规化。
\end{proof}

\begin{lemma}
\label{resolve-lemma-normalized-blowup-completion}
令 $(A, \mathfrak m, \kappa)$ 为诺特局部整环，其完备化 $A^\wedge$ 正规。
则对任意正规化爆破序列
$$
Y_n \to Y_{n - 1} \to \ldots \to Y_1 \to \Spec(A^\wedge)
$$
存在一列（固有的）正规化爆破
$$
X_n \to X_{n - 1} \to \ldots \to X_1 \to \Spec(A)
$$
其到 $A^\wedge$ 的基变换给出原序列。
\end{lemma}

\begin{proof}
给定序列 $Y_n \to \ldots \to Y_1 \to Y_0 = \Spec(A^\wedge)$，归纳构造
$X_n \to \ldots \to X_1 \to X_0 = \Spec(A)$。基本情形为 $i = 0$。
给定 $X_i$，其基变换为 $Y_i$，令 $Y'_i \to Y_i$ 为沿闭点
$y_i \in Y_i$ 的爆破，使得 $Y_{i + 1}$ 是 $Y_i$ 的正规化。
由于 $Y_i$ 与 $X_i$ 的闭纤维同构，点 $y_i$ 对应于闭点 $x_i$，后者位于
$X_i$ 的特殊纤维上。令 $X'_i \to X_i$ 为 $X_i$ 沿 $x_i$ 的爆破。
于是 $X'_i$ 到 $\Spec(A^\wedge)$ 的基变换同构于 $Y'_i$。
由引理 \ref{resolve-lemma-normalization-completion}，正规化
$X_{i + 1} \to X'_i$ 有限，且其到 $\Spec(A^\wedge)$ 的基变换同构于
$Y_{i + 1}$。
\end{proof}














\section{有理双点}
\label{resolve-section-rational-double-points}

\noindent
在第 \ref{resolve-section-rational-singularities} 节中，我们论证了 $2$ 维有理奇点的
消解可归约到 Gorenstein 情形。Gorenstein 有理曲面奇点是有理双点。
我们将通过显式计算消解它们。

\medskip\noindent
根据例子章第 \ref{examples-section-bad} 节的讨论，存在正规诺特局部整环
$A$，其完备化同构于 $\mathbf{C}[[x, y, z]]/(z^2)$。在这种情形下，
可以说 $A$ 有一个有理双点奇点；但另一方面，$\Spec(A)$ 不存在奇点消解。
若 $A$ 是 Nagata 环，这种现象不会发生；见代数章引理
\ref{algebra-lemma-local-nagata-domain-analytically-unramified}。

\medskip\noindent
然而还有更糟的情形：存在局部正规 Nagata 整环 $A$，其完备化为
$\mathbf{C}[[x, y, z]]/(yz)$；另有一个这样的环，其完备化为
$\mathbf{C}[[x, y, z]]/(y^2 - z^3)$。这是 \cite{Nishimura-few} 的例 2.5。
因此，本节需要假设所论环的完备化正规。

\begin{situation}
\label{resolve-situation-rational-double-point}
此处 $(A, \mathfrak m, \kappa)$ 为维数 $2$ 的 Nagata 局部正规整环，
它定义有理奇点，其完备化正规，并且它是 Gorenstein 环。假设 $A$ 不正则。
\end{situation}

\noindent
本节的论证将表明，在这一情境下，反复爆破奇点即可消解 $\Spec(A)$。
证明过程中需要下述引理。

\begin{lemma}
\label{resolve-lemma-issquare}
令 $\kappa$ 为域，令 $I \subset \kappa[x, y]$ 为理想。设
$$
a + b x + c y + d x^2 + exy + f y^2 \in I^2
$$
其中 $a, b, c, d, e, f \in k$ 不全为零。若 $I$ 在 $\kappa[x, y]$ 中的
余长度为 $> 1$，则
$a + b x + c y + d x^2 + exy + f y^2 = j(g + hx + iy)^2$
，其中某些 $j, g, h, i \in \kappa$。
\end{lemma}

\begin{proof}
考虑偏导数 $b + 2dx + ey$ 与 $c + ex + 2fy$。由 Leibniz 法则，
它们属于 $I$。若其中一个非零，则作形如
$x \mapsto \alpha + \beta x + \gamma y$ 与
$y \mapsto \delta + \epsilon x + \zeta y$ 的线性坐标变换后，可设
$x \in I$。于是 $I = (x)$，或者 $I = (x, F)$，其中 $F$ 是次数
$\geq 2$ 的首一多项式，以 $y$ 为变量。第一种情形结论显然。第二种情形下，
注意 $I^2$ 中任意元素都可写成
$$
A(x, y) x^2 + B(y) x F + C(y) F^2
$$
的形式，其中 $A(x, y) \in \kappa[x, y]$ 且 $B, C \in \kappa[y]$。因此
$$
a + b x + c y + d x^2 + exy + f y^2 = A(x, y) x^2 + B(y) x F + C(y) F^2
$$
由次数可知 $B = C = 0$，且 $A$ 为常数。

\medskip\noindent
为完成证明，还需处理两个偏导数都为零的情形。这只能发生在特征 $2$ 时，
此时得到
$$
a + d x^2 + f y^2 \in I^2
$$
可设 $f$ 非零（否则交换 $x$ 与 $y$ 的角色）。除以 $f$ 后，归结为
$\kappa$ 的特征为 $2$ 且
$$
a + d x^2 + y^2 \in I^2
$$
若 $a$ 与 $d$ 都是 $\kappa$ 中的平方，则结论成立。否则，存在导子
$\theta : \kappa \to \kappa$，使得 $\theta(a) \not = 0$ 或
$\theta(d) \not = 0$；见代数章引理
\ref{algebra-lemma-derivative-zero-pth-power}。可将它延拓为
$\kappa[x, y]$ 的导子，所取延拓满足 $\theta(x) = \theta(y) = 0$。于是
$$
\theta(a) + \theta(d) x^2 \in I
$$
情形 $\theta(d) = 0$ 导致矛盾。因此可设
$\alpha + x^2 \in I$，其中某个 $\alpha \in \kappa$。结合上式可知
$a + \alpha d + y^2 \in I$。故
$$
J = (\alpha + x^2, a + \alpha d + y^2) \subset I
$$
其余维数至多为 $2$。注意 $J/J^2$ 是 $\kappa[x, y]/J$ 上的自由模，基为
$\alpha + x^2$ 与 $a + \alpha d + y^2$。因此
$a + d x^2 + y^2 =
1 \cdot (a + \alpha d + y^2) + d \cdot (\alpha + x^2) \in I^2$
蕴含包含关系 $J \subset I$ 是严格的。于是可找到形如
$g + hx + iy + jxy$ 的非零元素，它属于 $I$。若 $j = 0$，则 $I$ 包含一个线性型，
可如第一段那样得出结论。因此 $j \not = 0$，且
$\dim_\kappa(I/J) = 1$（否则可在 $I$ 中找到上述形式且 $j = 0$ 的元素）。
所以 $I$ 具有形式
$(\alpha + x^2, \beta + y^2, g + hx + iy + jxy)$
，其中 $j \not = 0$，且余长度为 $3$。在这种情形下，
$a + dx^2 + y^2 \in I^2$ 不可能成立。可直接计算证明这一点，但我们采用
如下论证。为证明该断言，可设 $\kappa$ 代数闭。于是作坐标变换
$x \mapsto \sqrt{\alpha} + x$ 与 $y \mapsto \sqrt{\beta} + y$，并可设
$I = (x^2, y^2, g' + h'x + i'y + jxy)$，其中 $j$ 不变。此时
$g' = h' = i' = 0$，否则 $I$ 的余长度不是 $3$。故得到
$I = (x^2, y^2, xy)$，结论显然。
\end{proof}

\noindent
令 $(A, \mathfrak m, \kappa)$ 如情境
\ref{resolve-situation-rational-double-point}。令 $X \to \Spec(A)$ 为沿
$\mathfrak m$ 的爆破，所爆破的概形是 $\Spec(A)$。由引理
\ref{resolve-lemma-blow-up-normal-rational}，$X$ 正规。由引理
\ref{resolve-lemma-rational-propagates}，$X$ 的所有奇点都是有理奇点。
由于 $\omega_A = A$，由引理 \ref{resolve-lemma-dualizing-blow-up-rational} 可知
$\omega_X \cong \mathcal{O}_X$（约定见注 \ref{resolve-remark-dualizing-setup} 中的讨论）。
因此 $X$ 的所有奇点都是 Gorenstein 奇点。此外，由引理
\ref{resolve-lemma-blowup-still-good}，$X$ 在闭点处的局部环具有正规完备化。
换言之，爆破 $\Spec(A)$ 后得到正规曲面 $X$，其奇点如情境
\ref{resolve-situation-rational-double-point} 所述。下文将不再说明而直接使用这一点。
（注意：在下述讨论过程中将看到，这些奇点只有有限多个。）

\medskip\noindent
令 $E \subset X$ 为例外除子。由引理 \ref{resolve-lemma-dualizing-blow-up-rational}，
有 $\omega_E = \mathcal{O}_E(-1)$。由引理
\ref{resolve-lemma-cohomology-blow-up-rational}，有
$\kappa = H^0(E, \mathcal{O}_E)$。因此 $E$ 是 Gorenstein 曲线；由代数曲线章
第 \ref{curves-section-Riemann-Roch} 节所讨论的 Riemann--Roch，有
$$
\chi(E, \mathcal{O}_E) = 1 - g = -(1/2) \deg(\omega_E) =
(1/2)\deg(\mathcal{O}_E(1))
$$
其中 $g = \dim_\kappa H^1(E, \mathcal{O}_E) \geq 0$。由簇章引理
\ref{varieties-lemma-ampleness-in-terms-of-degrees-components}，
$\deg(\mathcal{O}_E(1))$ 为正，故 $g = 0$ 且
$\deg(\mathcal{O}_E(1)) = 2$。因此有
$$
\dim_\kappa (\mathfrak m^n/\mathfrak m^{n + 1}) = 2n + 1
$$
；这是由引理 \ref{resolve-lemma-cohomology-blow-up-rational} 以及 $E$ 上的
Riemann--Roch 得到的。

\medskip\noindent
选取 $x_1, x_2, x_3 \in \mathfrak m$，使其映到
$\mathfrak m/\mathfrak m^2$ 的一组基。因为
$\dim_\kappa(\mathfrak m^2/\mathfrak m^3) = 5$，元素 $x_i x_j$ 的像满足
一个关系；这里 $i \geq j$，且所论空间是这个 $\kappa$-向量空间。
换言之，可以找到 $a_{ij} \in A$（$i \geq j$），它们不全属于
$\mathfrak m$，并使得
$$
a_{11} x_1^2 + a_{12} x_1x_2 + a_{13}x_1x_3 + a_{22} x_2^2 +
a_{23} x_2x_3 + a_{33} x_3^2 =
\sum a_{ijk} x_ix_jx_k
$$
，其中某些 $a_{ijk} \in A$，且 $i \leq j \leq k$。用
$a \mapsto \overline{a}$ 表示映射 $A \to \kappa$。二次型
$q = \sum \overline{a}_{ij} t_i t_j \in \kappa[t_1, t_2, t_3]$
由上述选择确定到相差 $\kappa^*$ 中元素的乘法为止。若在论证过程中发现
$\overline{a}_{ij} = 0$ 于 $\kappa$ 中，则可把项 $a_{ij} x_i x_j$
并入右端，并设 $a_{ij} = 0$；这一操作改变 $a_{ijk}$，但不改变其他
$a_{i'j'}$。

\medskip\noindent
该爆破由 $3$ 个仿射图覆盖，它们对应于“变量” $x_1, x_2, x_3$。
由对称性，只需研究其中一个图。为此令
$$
A' = A[\mathfrak m/x_1]
$$
为仿射爆破代数（如代数章第 \ref{algebra-section-blow-up} 节）。
由于 $x_1, x_2, x_3$ 生成 $\mathfrak m$，可知 $A'$ 由
$y_2 = x_2/x_1$ 与 $y_3 = x_3/x_1$ 生成，基环为 $A$。为简化公式，偶尔使用
$y_1 = 1$。此外，考察上述关系可得
$$
a_{11} + a_{12} y_2 + a_{13} y_3 + a_{22} y_2^2 +
a_{23} y_2y_3 + a_{33} y_3^2 =
x_1 (\sum a_{ijk} y_iy_jy_k)
$$
于 $A'$ 中成立。回忆 $x_1 \in A'$ 定义例外除子 $E$；这里是在某个仿射开子集上，
而该开子集属于 $X$。因而该除子在概形意义下由下式给出：
$$
\kappa[y_2, y_3]/
(\overline{a}_{11} + \overline{a}_{12} y_2 + \overline{a}_{13} y_3 +
\overline{a}_{22} y_2^2 + \overline{a}_{23} y_2y_3 + \overline{a}_{33} y_3^2)
$$
换言之，$E \subset \mathbf{P}^2_\kappa = \text{Proj}(\kappa[t_1, t_2, t_3])$
是上述二次型 $q$ 的零概形。

\medskip\noindent
二次型 $q$ 是由 $A$ 定义的奇点的一个重要不变量。若 $q$ 是某个线性型的
平方再乘以 $\kappa^*$ 中的元素，则称处于{\bf 情形 II}；否则称处于
{\bf 情形 I}。注意，恰好在改变 $x_1, x_2, x_3$ 的选择后有
$$
x_3^2 = \sum a_{ijk}x_ix_jx_k
$$
于局部环 $A$ 中成立时，我们处于情形 II。

\medskip\noindent
令 $\mathfrak m' \subset A'$ 为位于 $\mathfrak m$ 上方的极大理想，
其剩余域为 $\kappa'$。换言之，$\mathfrak m'$ 对应于例外除子上的闭点
$p \in E$。回忆满射
$$
\kappa[y_2, y_3] \to \kappa'
$$
的核由两个元素 $f_2, f_3 \in \kappa[y_2, y_3]$ 生成（例如见代数章例
\ref{algebra-example-spec-kxy}，或代数章引理
\ref{algebra-lemma-dim-affine-space} 的证明）。令 $z_2, z_3 \in A'$ 映到
$f_2, f_3$，这里的像在 $\kappa[y_2, y_3]$ 中考察。于是
$\mathfrak m' = (x_1, z_2, z_3)$，因为 $x_2$ 与 $x_3$ 都变得可被
$x_1$ 整除，而这一切发生在 $A'$ 中。

\medskip\noindent
{\bf 断言。} 若 $X$ 在 $p$ 处奇异，则 $\kappa' = \kappa$，或者我们处于
情形 II。事实上，若 $A'_{\mathfrak m'}$ 奇异，则
$\dim_{\kappa'} \mathfrak m'/(\mathfrak m')^2 = 3$，这蕴含
$\dim_{\kappa'} \overline{\mathfrak m}'/(\overline{\mathfrak m}')^2 = 2$
，其中 $\overline{m}'$ 是
$\mathcal{O}_{E, p} = \mathcal{O}_{X, p}/x_1\mathcal{O}_{X, p}$ 的极大理想。
这蕴含
$$
q(1, y_2, y_3) =
\overline{a}_{11} + \overline{a}_{12} y_2 + \overline{a}_{13} y_3 +
\overline{a}_{22} y_2^2 + \overline{a}_{23} y_2y_3 + \overline{a}_{33} y_3^2
\in (f_2, f_3)^2
$$
，否则 $z_2$ 与 $z_3$ 的类在
$\overline{\mathfrak m}'/(\overline{\mathfrak m}')^2$ 中会满足一个关系。
现在断言由引理 \ref{resolve-lemma-issquare} 得到。

\medskip\noindent
情形 I 中的消解。由断言，$X$ 的任意奇点都是 $\kappa$-有理点。
选取这样的奇点 $p$。可以选取 $x_1, x_2, x_3 \in \mathfrak m$，使得
$p$ 位于上述仿射图中，且坐标为 $y_2 = y_3 = 0$。由于它是奇点，
仿照断言的证明可得 $q(1, y_2, y_3) \in (y_2, y_3)^2$。
因此可选 $a_{11} = a_{12} = a_{13} = 0$，并且
$q(t_1, t_2, t_3) = q(t_2, t_3)$。于是
$$
E = V(q) \subset \mathbf{P}^1_\kappa
$$
或者是两条在 $p$ 相交的不同直线之并，或者是次数 $2$ 的曲线，且具有唯一
$\kappa$-有理点（略去一个小细节；使用 $q$ 在相差标量的意义下
不是线性型的平方）。在两种情形下，都可推出 $X$ 有唯一奇点 $p$，
且该点是 $\kappa$-有理点。在这种情形下还需要更多一点信息。
首先，考察上式中的高次项，可得 $\overline{a}_{111} = 0$；这是因为
$p$ 奇异。于是可写成
$a_{111} = b_{111} x_1 \bmod (x_2, x_3)$，其中某个 $b_{111} \in A$。
此时在 $p$ 处，相对于生成元 $x_1, y_2, y_3$，二次型如下；这些元素生成
$\mathfrak m'$：
$$
q' =
\overline{b}_{111} t_1^2 +
\overline{a}_{112} t_1 t_2 + \overline{a}_{113} t_1 t_3 +
\overline{a}_{22} t_2^2 +
\overline{a}_{23} t_2 t_3 +
\overline{a}_{33} t_3^2
$$
可见 $E' = V(q')$ 与直线 $t_1 = 0$ 相交于两个点，或相交于一个次数
$2$ 的点。因此 $p$ 处于情形 I。

\medskip\noindent
假设爆破 $X' \to X$；这里爆破的是 $X$，中心为 $p$，并且再次有奇点 $p'$。则 $p'$ 是
$\kappa$-有理点，并且可以继续爆破而得到 $X'' \to X'$。若这一过程不停止，
便得到爆破序列
$$
\Spec(A) \leftarrow X \leftarrow X' \leftarrow X'' \leftarrow \ldots
$$
我们要证明引理 \ref{resolve-lemma-sequence-blowups} 适用于此情形。为此必须说明
元素 $x_1$ 的选择；此元素属于 $\mathfrak m$。假设 $A$ 处于情形 I，且 $X$ 有奇点。
称 $x_1 \in \mathfrak m$ 为{\it 良坐标}，如果对 $x_2, x_3$ 的任意选择
（等价地，某个选择），二次型 $q(t_1, t_2, t_3)$ 满足
$q(0, t_2, t_3)$ 不是一个平方的标量倍。上面已经看到良坐标存在。
若 $x_1$ 是良坐标，则奇点 $p \in E$（它是 $X$ 的奇点）不在超曲面 $t_1 = 0$ 上；
这是因为该超曲面或者没有有理点，或者即使有，该点在 $X$ 上也不奇异。
注意，这等价于说 $x_1$ 在 $\mathcal{O}_{X, p}$ 中的像截出例外除子 $E$。
上述计算还表明，若 $x_1$ 是 $A$ 的良坐标，则
$x_1 \in \mathfrak m'\mathcal{O}_{X, p}$ 是 $p$ 处的良坐标。当然，这里用到了
良坐标的概念不依赖于计算中所用的 $x_2$, $x_3$ 的选择。因此 $x_1$ 在
$p'$、$p''$ 等处都映为良坐标。于是引理 \ref{resolve-lemma-sequence-blowups}
适用，并且我们的爆破序列来自非奇异弧 $A \to R$。此时映射
$A^\wedge \to R$ 是满射。由于 $A$ 的完备化正规，由引理
\ref{resolve-lemma-sequence-blowups-along-arc-becomes-nonsingular} 可知，经过有限次爆破后，
$$
\Spec(A^\wedge) \leftarrow X^\wedge \leftarrow (X')^\wedge \leftarrow \ldots
$$
所得概形 $(X^{(n)})^\wedge$ 正则。由于
$(X^{(n)})^\wedge \to X^{(n)}$ 在完备局部环上诱导同构（引理
\ref{resolve-lemma-iso-completions}），可知 $X^{(n)}$ 也正则。

\medskip\noindent
情形 II 中的消解。此时有
$$
x_3^2 = \sum a_{ijk}x_ix_jx_k
$$
于 $A$ 中成立，其中 $x_1, x_2, x_3$ 是 $\mathfrak m$ 的某组生成元。
于是 $q = t_3^2$ 且 $E = 2C$，其中 $C$ 是一条直线。回忆在 $A'$ 中有
$$
y_3^2 = x_1(\sum a_{ijk} y_iy_jy_k)
$$
由于已知 $X$ 正规，得到离散赋值环 $\mathcal{O}_{X, \xi}$；这里 $\xi$ 是
$C$ 的泛点。元素 $y_3 \in A'$ 映到 $\mathcal{O}_{X, \xi}$ 的一个
一致化参数。由于 $x_1$ 在概形意义下截出 $E$，而 $C$ 在其中具有重数 $2$，
可知 $x_1$ 是 $y_3^2$ 的一个单位倍，这一陈述在 $\mathcal{O}_{X, \xi}$ 中成立。
考察上式可推出
$$
h(y_2) = \overline{a}_{111} + \overline{a}_{112} y_2 +
\overline{a}_{122} y_2^2 +
\overline{a}_{222} y_2^3
$$
在 $\xi$ 的剩余域中必须非零。现在，假设 $p \in C$ 定义一个奇点。
则 $y_3$ 在 $p$ 处为零；由证明上述断言时所用的论证，$p$ 必对应于
$h$ 的一个零点。若 $h$ 在 $p$ 处没有二重零点，则二次型 $q'$ 在 $p$ 处
不是平方，故 $p$ 落入上面已经处理的情形 I\footnote{极大理想在 $p$ 处，
而所处环为 $A'$；它由 $y_3, x_1$ 和第三个元素 $g$ 生成。后者在
$\kappa[y_2]$ 中的像是 $h$ 的素因子，并对应于 $p$。若该素因子不整除
$h$ 两次，则 $p$ 处的二次型形如
$$
y_3^2 - x_1((something)x_1 + (something)y_3 + (unit)g)
$$
，它绝不可能在 $\kappa[y_3, x_1, g]$ 中为平方。}。由于 $h$ 的次数为
$3$，至多有一个落入情形 II 的奇点 $p \in C$；而且它是
$\kappa$-有理点。改变 $x_1, x_2, x_3$ 的选择后，可设该点为
$y_2 = y_3 = 0$。于是
$h = \overline{a}_{122} y_2^2 + \overline{a}_{222} y_2^3$。此外，仍须有
$\overline{a}_{113} = 0$，二次型 $q'$ 才具有所需形状。因此局部环
$\mathcal{O}_{X, p}$ 定义下一段所述的奇点。

\medskip\noindent
最后处理如下情形：可以选取生成元 $x_1, x_2, x_3$，它们生成 $\mathfrak m$，使得
$$
x_3^2 + x_1(a x_2^2 + b x_2x_3 + c x_3^2) \in \mathfrak m^4
$$
，其中某些 $a, b, c \in A$。这是情形 II 的一个子类。若
$\overline{a} = 0$，则可写成 $a = a_1 x_1 + a_2 x_2 + a_3 x_3$，
爆破后得到
$$
y_3^2 + x_1(a_1 x_1 y_2^2 + a_2 x_1 y_2^3 + a_3 x_1 y_2^2 y_3 +
b y_2 y_3 + c y_3^2) = x_1^2 (\sum a_{ijkl}y_iy_jy_ky_l)
$$
这意味着 $X$ 不正规\footnote{事实上，该方程表明，沿 $1$ 维轨迹
$x_1 = y_3 = 0$ 会出现奇异情形，而这不可能发生在正规曲面上。}，矛盾。
由上一段的结果，若爆破 $X$ 有落入情形 II 的奇点 $p$，则这样的点只有一个，
并且它是 $\kappa$-有理点。计算仿射爆破代数
$A[\frac{\mathfrak m}{x_2}]$ 和 $A[\frac{\mathfrak m}{x_3}]$
，读者容易看出 $p$ 不可能包含在 $X$ 的相应开子集中。因此 $p$ 位于
$A[\frac{\mathfrak m}{x_1}]$ 的谱中。像前面一样作爆破，可知 $p$ 必为坐标
$y_2 = y_3 = 0$ 的点，且新方程形如
$$
y_3^2 + x_1(a y_2^2 + b y_2 y_3 + c y_3^2) \in (\mathfrak m')^4
$$
，它与之前具有相同形状，并且 $x_1$ 定义例外除子。因此，若过程不停止，
便得到无限爆破序列，并且在每一步中，$x_1$ 都在奇点的局部环中定义例外除子。
于是可用引理 \ref{resolve-lemma-sequence-blowups}、
\ref{resolve-lemma-sequence-blowups-along-arc-becomes-nonsingular} 以及与此前相同的论证
完成证明。

\begin{lemma}
\label{resolve-lemma-resolve-rational-double-points}
令 $(A, \mathfrak m, \kappa)$ 为维数 $2$ 的局部正规 Nagata 整环，
它定义有理奇点，其完备化正规，并且它是 Gorenstein 环。
则存在一列有限次沿奇异闭点的爆破
$$
X_n \to X_{n - 1} \to \ldots \to X_1 \to X_0 = \Spec(A)
$$
，使得 $X_n$ 正则，并且每个中间概形 $X_i$ 正规，只有有限多个同类型奇点。
\end{lemma}

\begin{proof}
这正是上述讨论已经证明的内容。
\end{proof}




\section{由存在性蕴含的性质}
\label{resolve-section-existence-gives}

\noindent
本节证明，对于诺特整概形，正则改造的存在性会带来相当多的后果。
不关心“坏”诺特环的读者可略过本节。

\begin{lemma}
\label{resolve-lemma-regular-alteration-implies}
令 $Y$ 为诺特整概形。假设存在改造
$f : X \to Y$，且 $X$ 正则。则正规化 $Y^\nu \to Y$
是有限的，并且 $Y$ 有一个正则的稠密开子概形。
\end{lemma}

\begin{proof}
只需在 $Y = \Spec(A)$ 且 $A$ 为诺特整环的情形证明此结论。
令 $B$ 为 $A$ 在其分式域中的整闭包。
置 $C = \Gamma(X, \mathcal{O}_X)$。由概形上同调章引理
\ref{coherent-lemma-proper-over-affine-cohomology-finite}
可知 $C$ 是有限 $A$-模。由于 $X$ 正规
（性质章引理 \ref{properties-lemma-regular-normal}），
$C$ 是正规整环
（性质章引理 \ref{properties-lemma-normal-integral-sections}）。
因此 $B \subset C$，且可知 $B$ 在 $A$ 上有限，因为 $A$ 诺特。

\medskip\noindent
存在非空开子概形 $V \subset Y$，使得 $f^{-1}V \to V$
有限；见态射章定义 \ref{morphisms-definition-alteration}。
缩小 $V$ 后，可假设 $f^{-1}V \to V$ 平坦
（态射章命题 \ref{morphisms-proposition-generic-flatness}）。
因此 $f^{-1}V \to V$ 忠实平坦。由代数章引理
\ref{algebra-lemma-descent-regular}，$V$ 正则。
\end{proof}

\begin{lemma}
\label{resolve-lemma-algebra-helper}
令 $(A, \mathfrak m)$ 为诺特局部环，令 $B \subset C$
为有限 $A$-代数。假设 (a) $B$ 是正规环，并且
(b) $\mathfrak m$-进完备化 $C^\wedge$ 是正规环。
则 $B^\wedge$ 是正规环。
\end{lemma}

\begin{proof}
考虑交换图
$$
\xymatrix{
B \ar[r] \ar[d] & C \ar[d] \\
B^\wedge \ar[r] & C^\wedge
}
$$
回忆：$\mathfrak m$-进完备化在有限 $A$-模范畴上是正合的，
因为它由张量上平坦 $A$-代数 $A^\wedge$ 给出
（代数章引理 \ref{algebra-lemma-completion-flat}）。
我们用 Serre 判据
（代数章引理 \ref{algebra-lemma-criterion-normal}）
证明诺特环 $B^\wedge$ 正规。
令 $\mathfrak q \subset B^\wedge$ 为位于
$\mathfrak p \subset B$ 上方的素理想。若 $\dim(B_\mathfrak p) \geq 2$，则
$\text{depth}(B_\mathfrak p) \geq 2$；又因
$B_\mathfrak p \to B^\wedge_\mathfrak q$ 平坦，可知
$\text{depth}(B^\wedge_\mathfrak q) \geq 2$
（代数章引理 \ref{algebra-lemma-apply-grothendieck}）。
若 $\dim(B_\mathfrak p) \leq 1$，则 $B_\mathfrak p$
或者是离散赋值环，或者是域。
此时 $C_\mathfrak p$ 在 $B_\mathfrak p$ 上忠实平坦
（因为它有限且无挠）。
所以 $B^\wedge_\mathfrak p \to C^\wedge_\mathfrak p$
忠实平坦；在 $\mathfrak q$ 处局部化后仍然如此。
由于 $C^\wedge$ 及其任意局部化都满足 $(S_2)$，由代数章引理
\ref{algebra-lemma-descent-Sk} 可知 $B^\wedge_\mathfrak p$ 满足 $(S_2)$。
综上可知，$(S_2)$ 对 $B^\wedge$ 成立。为证明 $B^\wedge$
满足 $(R_1)$，只需考虑素理想 $\mathfrak q \subset B^\wedge$，且它满足
$\dim(B^\wedge_\mathfrak q) \leq 1$。由
$\dim(B^\wedge_\mathfrak q) = \dim(B_\mathfrak p) +
\dim(B^\wedge_\mathfrak q/\mathfrak p B^\wedge_\mathfrak q)$ 以及
代数章引理 \ref{algebra-lemma-dimension-base-fibre-total}，
可知 $\dim(B_\mathfrak p) \leq 1$；并且同前可知
$B^\wedge_\mathfrak q \to C^\wedge_\mathfrak q$
忠实平坦。结论由代数章引理 \ref{algebra-lemma-descent-Rk} 得到。
\end{proof}

\begin{lemma}
\label{resolve-lemma-regular-alteration-implies-local}
令 $(A, \mathfrak m, \kappa)$ 为诺特局部整环。
假设存在改造 $f : X \to \Spec(A)$，且 $X$ 正则。则
\begin{enumerate}
\item 存在非零元 $f \in A$，使得 $A_f$ 正则；
\item 整闭包 $B$，即 $A$ 在其分式域中的整闭包，在 $A$ 上有限；
\item $\mathfrak m$-进完备化后的 $B$ 是正规环；换言之，
$B$ 在其各极大理想处的完备化都是正规整环；
\item $A$ 的泛形式纤维正则。
\end{enumerate}
\end{lemma}

\begin{proof}
第 (1) 与 (2) 项由引理 \ref{resolve-lemma-regular-alteration-implies} 得到。
为给第 (3) 项的证明建立记号，需要重做该引理证明的一部分。
置 $C = \Gamma(X, \mathcal{O}_X)$。由概形上同调章引理
\ref{coherent-lemma-proper-over-affine-cohomology-finite}
可知 $C$ 是有限 $A$-模。由于 $X$ 正规
（性质章引理 \ref{properties-lemma-regular-normal}），
$C$ 是正规整环
（性质章引理 \ref{properties-lemma-normal-integral-sections}）。
因此 $B \subset C$，且可知 $B$ 在 $A$ 上有限，因为 $A$ 诺特。
由引理 \ref{resolve-lemma-algebra-helper}，为证明第 (3) 项，只需证明
$\mathfrak m$-进完备化 $C^\wedge$ 正规。

\medskip\noindent
由代数章引理 \ref{algebra-lemma-completion-finite-extension}，
完备化 $C^\wedge$ 是 $C$ 的如下各完备化之乘积：取 $C$ 中位于
$\mathfrak m$ 上方的素理想，并在这些素理想处完备化。这样的素理想只有有限个，可列为
$\mathfrak m_1, \ldots, \mathfrak m_r$，而它们正是 $C$ 的极大理想。
（对 $B$ 的相应结论说明了本引理最后一句。）
因此，将 $A$ 替换为 $C_{\mathfrak m_i}$，并将 $X$ 替换为
$X_i = X \times_{\Spec(C)} \Spec(C_{\mathfrak m_i})$，
即可归约到下一段讨论的情形。
（注意，由概形上同调章引理 \ref{coherent-lemma-flat-base-change-cohomology}，
$\Gamma(X_i, \mathcal{O}) = C_{\mathfrak m_i}$。）

\medskip\noindent
此时 $A$ 是诺特局部正规整环，$f : X \to \Spec(A)$ 是满足
$\Gamma(X, \mathcal{O}_X) = A$ 的正则改造。需要证明完备化
$A^\wedge$（所完备化的环为 $A$）是正规整环。由引理
\ref{resolve-lemma-port-regularity-to-completion}，
$Y = X \times_{\Spec(A)} \Spec(A^\wedge)$ 正则。
又由概形上同调章引理 \ref{coherent-lemma-flat-base-change-cohomology}，
$\Gamma(Y, \mathcal{O}_Y) = A^\wedge$，故与前面一样可知
$A^\wedge$ 正规。具体地，由性质章引理
\ref{properties-lemma-regular-normal}，$Y$ 正规。
因为 $\Gamma(Y, \mathcal{O}_Y) = A^\wedge$ 是局部环，所以它连通。
因而 $Y$ 正规且整（对诺特概形，连通且正规蕴含整）。于是由性质章引理
\ref{properties-lemma-normal-integral-sections}，
$\Gamma(Y, \mathcal{O}_Y) = A^\wedge$ 是正规整环。
这证明了第 (3) 项。

\medskip\noindent
证明第 (4) 项。令 $\eta \in \Spec(A)$ 表示泛点，并用下标
$\eta$ 表示到 $\eta$ 的基变换。由于 $f$ 是改造，概形 $X_\eta$
在 $\eta$ 上有限且忠实平坦。由引理
\ref{resolve-lemma-port-regularity-to-completion}，
$Y = X \times_{\Spec(A)} \Spec(A^\wedge)$ 正则，故 $Y_\eta$ 正则
（因为它是 $Y$ 中开子集的极限）。于是
$Y_\eta \to \Spec(A^\wedge \otimes_A \kappa(\eta))$
是到泛形式纤维的有限忠实平坦态射。结论由代数章引理
\ref{algebra-lemma-descent-regular} 得到。
\end{proof}











\section{奇点消解}
\label{resolve-section-resolution}



\noindent
先给出一个定义。

\begin{definition}
\label{resolve-definition-resolution}
令 $Y$ 为诺特整概形。$Y$ 的一个{\it 奇点消解}是修改
$f : X \to Y$，其中 $X$ 正则。
\end{definition}

\noindent
对于曲面，我们有时还需要稍多一些信息。

\begin{definition}
\label{resolve-definition-resolution-surface}
令 $Y$ 为 $2$ 维诺特整概形。称 $Y$ 有一个
{\it 由正规化爆破给出的奇点消解}，如果存在序列
$$
Y_n \to Y_{n - 1} \to \ldots \to Y_1 \to Y_0 \to Y
$$
满足
\begin{enumerate}
\item $Y_i$ 在 $Y$ 上固有，其中 $i = 0, \ldots, n$；
\item $Y_0 \to Y$ 是正规化；
\item $Y_i \to Y_{i - 1}$ 是正规化爆破，其中 $i = 1, \ldots, n$；
\item $Y_n$ 正则。
\end{enumerate}
\end{definition}

\noindent
注意，条件 (1) 蕴含：正规化 $Y_0$（即 $Y$ 的正规化）在 $Y$ 上有限，
并且正规化爆破中所用的正规化也都是有限的。

\begin{lemma}
\label{resolve-lemma-existence-implies-existence-by-normalized-blowing-ups}
令 $(A, \mathfrak m, \kappa)$ 为诺特局部环。
假设 $A$ 正规且维数为 $2$。若 $\Spec(A)$ 有奇点消解，
则 $\Spec(A)$ 有由正规化爆破给出的奇点消解。
\end{lemma}

\begin{proof}
由引理 \ref{resolve-lemma-regular-alteration-implies-local}，完备化
$A^\wedge$（所完备化的是 $A$）正规。由引理
\ref{resolve-lemma-port-regularity-to-completion} 可知
$\Spec(A^\wedge)$ 有奇点消解。由引理
\ref{resolve-lemma-normalized-blowup-completion}，任意正规化爆破序列
$Y_n \to Y_{n - 1} \to \ldots \to \Spec(A^\wedge)$
都来自正规化爆破序列 $X_n \to \ldots \to \Spec(A)$。此外，若 $Y_n$
正则，则由引理 \ref{resolve-lemma-port-regularity-to-completion}，$X_n$ 正则。
因此只需在 $A$ 完备的情形证明本引理。

\medskip\noindent
再假设 $A$ 完备。我们要用到：$A$ 是 Nagata 环
（代数章命题 \ref{algebra-proposition-ubiquity-nagata}）、优良环
（进阶代数章命题 \ref{more-algebra-proposition-ubiquity-excellent}），
且有对偶化复形
（对偶化复形章引理 \ref{dualizing-lemma-ubiquity-dualizing}）。
此外，对任意在 $A$ 上本质有限型的环，同样的结论成立。
若 $B$ 是优良局部正规整环，则完备化 $B^\wedge$ 正规
（因为 $B \to B^\wedge$ 正则，并可应用进阶代数章引理
\ref{more-algebra-lemma-normal-goes-up}）。在以下证明中，我们将使用此事实而不再说明。

\medskip\noindent
令 $X \to \Spec(A)$ 为奇点消解。选取正规化爆破序列
$$
Y_n \to Y_{n - 1} \to \ldots \to Y_1 \to \Spec(A)
$$
支配 $X$（引理 \ref{resolve-lemma-dominate-by-normalized-blowing-up}）。
态射 $Y_n \to X$ 在 $X$ 的有限多个点之外是同构。
因此可应用引理 \ref{resolve-lemma-dominate-by-blowing-up-in-points}，
找到沿闭点的爆破序列
$$
X_m \to X_{m - 1} \to \ldots \to X
$$
使得 $X_m$ 支配 $Y_n$。图示为
$$
\xymatrix{
& Y_n \ar[rd] \ar[rr] & & \Spec(A) \\
X_m \ar[rr] \ar[ru] & & X \ar[ru]
}
$$
为证明本引理，只需证明对 $Y_n$ 作有限次正规化爆破便可得到正则概形。
由上图可知 $Y_n$ 有奇点消解，即 $X_m$。由于 $Y_n$ 是正规曲面，
这说明 $Y_n$ 至多有有限多个奇点 $y_1, \ldots, y_t$
（因为 $X_m \to Y_n$ 在维数为 $1$ 的纤维之外是同构；见簇章引理
\ref{varieties-lemma-modification-normal-iso-over-codimension-1}）。

\medskip\noindent
令 $x_a \in X$ 为 $y_a$ 的像。则 $\mathcal{O}_{X, x_a}$ 正则，
因而定义有理奇点（引理 \ref{resolve-lemma-regular-rational}）。
将引理 \ref{resolve-lemma-rational-propagates} 应用于
$\mathcal{O}_{X, x_a} \to \mathcal{O}_{Y_n, y_a}$，
可知 $\mathcal{O}_{Y_n, y_a}$ 定义有理奇点。由引理
\ref{resolve-lemma-rational-to-gorenstein}，存在沿奇异闭点的有限爆破序列
$$
Y_{a, n_a} \to Y_{a, n_a - 1} \to \ldots \to \Spec(\mathcal{O}_{Y_n, y_a})
$$
使得 $Y_{a, n_a}$ 是 Gorenstein 的，即有可逆对偶化模。由（本质上平凡的）
引理 \ref{resolve-lemma-equivalence-sequence-blowups}，取 $n' = \sum n_a$，
这些序列对应于爆破序列
$$
Y_{n + n'} \to Y_{n + n' - 1} \to \ldots \to Y_n
$$
使得 $Y_{n + n'}$ 正规，且 $Y_{n + n'}$ 的局部环都是 Gorenstein 的。
利用上面给出的引理，可支配 $Y_{n + n'}$；更具体地，可取爆破序列
$X_{m + m'} \to \ldots \to X_m$ 支配 $Y_{n + n'}$，如下图所示：
$$
\xymatrix{
 & Y_{n + n'} \ar[rr] & & Y_n \ar[rd] \ar[rr] & & \Spec(A) \\
X_{m + m'} \ar[ru] \ar[rr] & & X_m \ar[rr] \ar[ru] & & X \ar[ru]
}
$$
因此 $Y_{n + n'}$ 同样只有有限多个奇点
$y'_1, \ldots, y'_s$；但此时这些奇点是有理双点。更准确地说，局部环
$\mathcal{O}_{Y_{n + n'}, y'_b}$ 满足引理
\ref{resolve-lemma-resolve-rational-double-points} 的情形。
完全按照上面的论证即可得到本引理。
\end{proof}

\begin{lemma}
\label{resolve-lemma-resolve-complete}
令 $(A, \mathfrak m, \kappa)$ 为诺特完备局部环。
假设 $A$ 是维数为 $2$ 的正规整环。则 $\Spec(A)$ 有奇点消解。
\end{lemma}

\begin{proof}
诺特完备局部环是 J-2 的
（进阶代数章命题 \ref{more-algebra-proposition-ubiquity-J-2}）、Nagata 的
（代数章命题 \ref{algebra-proposition-ubiquity-nagata}）、优良的
（进阶代数章命题 \ref{more-algebra-proposition-ubiquity-excellent}），
并且有对偶化复形
（对偶化复形章引理 \ref{dualizing-lemma-ubiquity-dualizing}）。
此外，对任意在 $A$ 上本质有限型的环，同样的结论成立。
若 $B$ 是优良局部正规整环，则完备化 $B^\wedge$ 正规
（因为 $B \to B^\wedge$ 正则，并可应用进阶代数章引理
\ref{more-algebra-lemma-normal-goes-up}）。换言之，在以下证明中遇到的局部环
都具有所需的“优良性”性质。

\medskip\noindent
选取 $A_0 \subset A$，使 $A_0$ 为正则完备局部环且
$A_0 \to A$ 有限；见代数章引理
\ref{algebra-lemma-complete-local-Noetherian-domain-finite-over-regular}。
这诱导分式域的有限扩张 $K/K_0$。我们对 $[K : K_0]$ 作归纳。
起始情形是次数为 $1$；此时 $A_0 = A$，结论成立。

\medskip\noindent
假设存在中间域 $K_0 \subset L \subset K$，且
$K_0 \not = L \not = K$。令 $B \subset A$ 为 $A_0$ 在 $L$ 中的整闭包。
由归纳假设，选取奇点消解 $Y \to \Spec(B)$。令 $X$ 为
$Y \times_{\Spec(B)} \Spec(A)$ 的正规化。图示为
$$
\xymatrix{
X \ar[r] \ar[d] & \Spec(A) \ar[d] \\
Y \ar[r] & \Spec(B)
}
$$
由于 $A$ 是 J-2 的，$X$ 的正则轨迹是开的。由于 $X$ 是正规曲面，
可知 $X$ 至多有有限多个奇点 $x_1, \ldots, x_n$；它们是满足
$\dim(\mathcal{O}_{X, x_i}) = 2$ 的闭点。
对每个 $i$，令 $y_i \in Y$ 为其像。由于
$\mathcal{O}_{Y, y_i}^\wedge \to \mathcal{O}_{X, x_i}^\wedge$
有限且次数小于原先的次数，由归纳假设可知
$\mathcal{O}_{X, x_i}^\wedge$ 有奇点消解。由引理
\ref{resolve-lemma-existence-implies-existence-by-normalized-blowing-ups}，存在正规化爆破序列
$$
Z^\wedge_{i, n_i} \to \ldots \to Z^\wedge_{i, 1} \to
\Spec(\mathcal{O}_{X, x_i}^\wedge)
$$
且 $Z^\wedge_{i, n_i}$ 正则。由引理
\ref{resolve-lemma-normalized-blowup-completion}，有相应的正规化爆破序列
$$
Z_{i, n_i} \to \ldots \to Z_{i, 1} \to \Spec(\mathcal{O}_{X, x_i})
$$
由引理 \ref{resolve-lemma-port-regularity-to-completion}，$Z_{i, n_i}$ 是正则概形。
由引理 \ref{resolve-lemma-equivalence-sequence-normalized-blowups}，
可将这些正规化爆破嵌入相应的序列
$$
Z_n \to Z_{n - 1} \to \ldots \to Z_1 \to X
$$
当然，考察局部环可知 $Z_n$ 也正则。这证明了归纳步骤。

\medskip\noindent
假设不存在满足 $K_0 \subset L \subset K$ 且
$K_0 \not = L \not = K$ 的中间域。则或者 $K/K_0$ 可分，
或者 $K$ 的特征为 $p$ 且 $[K : K_0] = p$。
于是，引理 \ref{resolve-lemma-go-up-separable} 或引理 \ref{resolve-lemma-go-up-degree-p}
蕴含可以归约到有理奇点。由引理 \ref{resolve-lemma-reduce-to-rational} 可知，
存在正规修改 $X \to \Spec(A)$，使得每个奇点 $x$（位于 $X$ 上）
所对应的局部环 $\mathcal{O}_{X, x}$ 定义有理奇点。由于 $A$ 是 J-2 的，
可知 $X$ 有有限多个奇点 $x_1, \ldots, x_n$。
由引理 \ref{resolve-lemma-rational-to-gorenstein}，存在沿奇异闭点的有限爆破序列
$$
X_{i, n_i} \to X_{i, n_i - 1} \to \ldots \to \Spec(\mathcal{O}_{X, x_i})
$$
使得 $X_{i, n_i}$ 是 Gorenstein 的，即有可逆对偶化模。由（本质上平凡的）
引理 \ref{resolve-lemma-equivalence-sequence-blowups}，取 $n = \sum n_a$，
这些序列对应于爆破序列
$$
X_n \to X_{n - 1} \to \ldots \to X
$$
使得 $X_n$ 正规，且 $X_n$ 的局部环都是 Gorenstein 的。
$X_n$ 同样只有有限多个奇点 $x'_1, \ldots, x'_s$，
但此时这些奇点是有理双点。更准确地说，局部环
$\mathcal{O}_{X_n, x'_i}$ 满足引理
\ref{resolve-lemma-resolve-rational-double-points} 的情形。
完全按照上面的论证即可得到本引理。
\end{proof}

\noindent
最后来到本章的主定理。

\begin{theorem}[Lipman]
\label{resolve-theorem-resolve}
\begin{reference}
\cite[Theorem on page 151]{Lipman}
\end{reference}
令 $Y$ 为二维诺特整概形。下列条件等价：
\begin{enumerate}
\item 存在改造 $X \to Y$，且 $X$ 正则；
\item $Y$ 存在奇点消解；
\item $Y$ 有由正规化爆破给出的奇点消解；
\item 正规化 $Y^\nu \to Y$ 有限，$Y^\nu$ 有有限多个奇点
$y_1, \ldots, y_m$，并且对每个 $y_i$，
$\mathcal{O}_{Y^\nu, y_i}$ 的完备化正规。
\end{enumerate}
\end{theorem}

\begin{proof}
蕴含关系 (3) $\Rightarrow$ (2) $\Rightarrow$ (1) 是立即的。

\medskip\noindent
令 $X \to Y$ 为改造，且 $X$ 正则。由引理
\ref{resolve-lemma-regular-alteration-implies}，$Y^\nu \to Y$ 有限。
考虑态射章引理 \ref{morphisms-lemma-normalization-normal} 给出的分解
$f : X \to Y^\nu$。由簇章引理 \ref{varieties-lemma-finite-in-codim-1}，
态射 $f$ 在开子概形 $V \subset Y^\nu$ 上有限；该开子概形包含每个余维
$\leq 1$ 的点，其中所说的余维是在 $Y^\nu$ 中计算的。由代数章引理
\ref{algebra-lemma-CM-over-regular-flat} 以及以下事实，$f$ 在 $V$ 上平坦：
由 Serre 判据，维数 $\leq 2$ 的正规局部环是 Cohen--Macaulay 的
（代数章引理 \ref{algebra-lemma-criterion-normal}）。
于是由代数章引理 \ref{algebra-lemma-descent-regular}，$V$ 正则。
由于 $Y^\nu$ 诺特，可知
$Y^\nu \setminus V = \{y_1, \ldots, y_m\}$ 是有限集。
由引理 \ref{resolve-lemma-regular-alteration-implies-local}，
$\mathcal{O}_{Y^\nu, y_i}$ 的完备化正规。
这样便得到 (1) $\Rightarrow$ (4)。

\medskip\noindent
假设 (4)。需要证明 (3)。可以立即用正规化替换 $Y$。
令 $y_1, \ldots, y_m \in Y$ 为奇点。应用引理
\ref{resolve-lemma-resolve-complete} 与
\ref{resolve-lemma-existence-implies-existence-by-normalized-blowing-ups}，
可知存在有限的正规化爆破序列
$$
Y_{i, n_i} \to Y_{i, n_i - 1} \to \ldots \to \Spec(\mathcal{O}^\wedge_{Y, y_i})
$$
使得 $Y_{i, n_i}$ 正则。由引理
\ref{resolve-lemma-normalized-blowup-completion}，有相应的正规化爆破序列
$$
X_{i, n_i} \to \ldots \to X_{i, 1} \to \Spec(\mathcal{O}_{Y, y_i})
$$
由引理 \ref{resolve-lemma-port-regularity-to-completion}，$X_{i, n_i}$ 是正则概形。
由引理 \ref{resolve-lemma-equivalence-sequence-normalized-blowups}，
可将这些正规化爆破嵌入相应的序列
$$
X_n \to X_{n - 1} \to \ldots \to X_1 \to Y
$$
当然，考察局部环可知 $X_n$ 也正则。证明完毕。
\end{proof}






\section{嵌入消解}
\label{resolve-section-embedded}

\noindent
给定曲面上的曲线，存在一次爆破，可将该曲线化为严格正规交叉除子。
本节将使用如下事实：一维局部诺特概形正规当且仅当它正则
（代数章引理 \ref{algebra-lemma-characterize-dvr}）。
还将使用如下事实：局部诺特概形上的任意点都特殊化到某个闭点
（性质章引理 \ref{properties-lemma-locally-Noetherian-closed-point}）。

\begin{lemma}
\label{resolve-lemma-resolve-curve}
令 $Y$ 为一维诺特整概形。下列条件等价：
\begin{enumerate}
\item 存在改造 $X \to Y$，且 $X$ 正则；
\item $Y$ 存在奇点消解；
\item 存在沿闭点的有限爆破序列
$Y_n \to Y_{n - 1} \to \ldots \to Y_1 \to Y$，且 $Y_n$ 正则；
\item 正规化 $Y^\nu \to Y$ 有限。
\end{enumerate}
\end{lemma}

\begin{proof}
蕴含关系 (3) $\Rightarrow$ (2) $\Rightarrow$ (1) 是立即的。
蕴含关系 (1) $\Rightarrow$ (4) 由引理
\ref{resolve-lemma-regular-alteration-implies} 得到。
注意，一维正规概形正则，因此蕴含关系 (4) $\Rightarrow$ (2) 也显然。
故只需证明等价条件 (1)、(2) 与 (4) 蕴含 (3)。

\medskip\noindent
令 $f : X \to Y$ 为奇点消解。由于 $Y$ 的维数为一，由簇章引理
\ref{varieties-lemma-finite-in-codim-1} 可知 $f$ 有限。
我们将构造分解
$$
X \to \ldots \to Y_2 \to Y_1 \to Y
$$
其中，$Y_i \to Y_{i - 1}$ 是某个闭点的爆破，且只要
$Y_{i - 1}$ 不正则，它就不是同构。这些态射中的每一个都有限（理由同上），
并得到相应的系统
$$
f_*\mathcal{O}_X \supset \ldots \supset
f_{2, *}\mathcal{O}_{Y_2} \supset
f_{1, *}\mathcal{O}_{Y_1} \supset \mathcal{O}_Y
$$
其中 $f_i : Y_i \to Y$ 是结构态射。由于 $Y$ 诺特，这个凝聚子模的递增序列
必定稳定（概形上同调章引理 \ref{coherent-lemma-acc-coherent}）。
这证明存在某个 $n$，使得概形 $Y_n$ 如所需那样正则。
为构造 $Y_i$，在给定 $Y_{i - 1}$ 后，选取奇异闭点
$y_{i - 1} \in Y_{i - 1}$，并令 $Y_i \to Y_{i - 1}$ 为相应的爆破。
由于 $X$ 正则且维数为 $1$（因而闭点处的局部环是离散赋值环，
特别地是主理想整环），理想层
$\mathfrak m_{y_{i - 1}} \cdot \mathcal{O}_X$ 可逆。
由爆破的泛性质（除子章引理
\ref{divisors-lemma-universal-property-blowing-up}），得到分解
$X \to Y_i$。最后，$Y_i \to Y_{i - 1}$ 不是同构，
因为 $\mathfrak m_{y_{i - 1}}$ 不是可逆理想。
\end{proof}

\begin{lemma}
\label{resolve-lemma-blowup-curve}
令 $X$ 为诺特概形。令 $Y \subset X$ 为维数 $1$ 的整闭子概形，
并假设它满足引理 \ref{resolve-lemma-resolve-curve} 的等价条件。则存在沿闭点的有限爆破序列
$$
X_n \to X_{n - 1} \to \ldots \to X_1 \to X
$$
使得 $Y$ 在 $X_n$ 中的严格变换是正则曲线。
\end{lemma}

\begin{proof}
令 $Y_n \to Y_{n - 1} \to \ldots \to Y_1 \to Y$ 为引理
\ref{resolve-lemma-resolve-curve} 给出的爆破序列。令
$X_n \to X_{n - 1} \to \ldots \to X_1 \to X$ 为 $X$ 的相应爆破序列。
这确实可行，因为由除子章引理 \ref{divisors-lemma-strict-transform}，
严格变换就是相应的爆破。
\end{proof}

\noindent
令 $X$ 为局部诺特概形，令 $Y, Z \subset X$ 为闭子概形。
令 $p \in Y \cap Z$ 为闭点。假设 $Y$ 是维数为 $1$ 的整概形，
且 $Y$ 的泛点不包含在 $Z$ 中。在此情形可考虑不变量
\begin{equation}
\label{resolve-equation-multiplicity}
m_p(Y \cap Z) =
\text{length}_{\mathcal{O}_{X, p}}(\mathcal{O}_{Y \cap Z, p})
\end{equation}
这是整数 $\geq 1$。事实上，若 $I, J \subset \mathcal{O}_{X, p}$
是对应于 $Y, Z$ 的理想，则
$\mathcal{O}_{Y \cap Z, p} = \mathcal{O}_{X, p}/I + J$
的支集等于 $\{\mathfrak m_p\}$，因为我们假设 $Y \cap Z$ 不包含
$Y$ 中特殊化到 $p$ 的唯一一点。因此由代数章引理
\ref{algebra-lemma-support-point}，其长度有限。

\begin{lemma}
\label{resolve-lemma-blowup-nonsingular-curves-meeting-at-point}
在上述情形中，令 $X' \to X$ 为 $X$ 在 $p$ 处的爆破。
令 $Y', Z' \subset X'$ 为 $Y, Z$ 的严格变换。
若 $\mathcal{O}_{Y, p}$ 正则，则
\begin{enumerate}
\item $Y' \to Y$ 是同构；
\item $Y'$ 与例外纤维 $E \subset X'$ 交于一点
$q$，且 $m_q(Y \cap E) = 1$；
\item 若还有 $q \in Z'$，则 $m_q(Y \cap Z') < m_p(Y \cap Z)$。
\end{enumerate}
\end{lemma}

\begin{proof}
由于 $\mathcal{O}_{X, p} \to \mathcal{O}_{Y, p}$ 满射，且
$\mathcal{O}_{Y, p}$ 是离散赋值环，可以选取元素
$x_1 \in \mathfrak m_p$，其在 $\mathcal{O}_{Y, p}$ 中的像是一致化参数。
选取仿射开集 $U = \Spec(A)$，它包含 $p$，并使得 $x_1 \in A$。
令 $\mathfrak m \subset A$ 为对应于 $p$ 的极大理想。令
$I, J \subset A$ 为定义 $Y, Z$ 的理想，这里所处的仿射概形为 $\Spec(A)$。
缩小 $U$ 后，可假设 $\mathfrak m = I + (x_1)$；换言之，
在概形论意义下 $V(x_1) \cap U \cap Y = \{p\}$。
由此可知 $p$ 是 $Y$ 上的有效 Cartier 除子；又因 $Y'$ 是 $Y$ 在 $p$ 处的爆破
（除子章引理 \ref{divisors-lemma-strict-transform}），
由除子章引理 \ref{divisors-lemma-blow-up-effective-Cartier-divisor} 可知
$Y' \to Y$ 是同构。关系 $\mathfrak m = I + (x_1)$ 蕴含
$\mathfrak m^n \subset I + (x_1^n)$，故可定义映射
$$
\psi : A[\textstyle{\frac{\mathfrak m}{x_1}}] \longrightarrow A/I
$$
如下定义：将 $y/x_1^n \in A[\frac{\mathfrak m}{x_1}]$ 映到
$a$ 在 $A/I$ 中的类，其中选取 $a$ 使得 $y \equiv ax_1^n \bmod I$。
于是 $\psi$ 对应于从 $Y \cap U$ 到 $X'$ 的态射；该态射在 $U$ 上，
并由 $Y' \cong Y$ 给出。
由于 $x_1$ 在 $A[\frac{\mathfrak m}{x_1}]$ 中的像截出例外除子，
可知 $m_q(Y', E) = 1$。最后，由 $J \subset \mathfrak m$ 可知理想
$J' \subset A[\frac{\mathfrak m}{x_1}]$ 必然包含所有元素
$f/x_1$，其中 $f \in J$。因此，若选取 $f \in J$，使其像
$\overline{f}$（该像取在 $A/I$ 中）的最小赋值等于 $m_p(Y \cap Z)$，则可知
$\psi(f/x_1) = \overline{f}/x_1$ 在 $A/I$ 中的赋值少一，
这证明了本引理的最后一项。
\end{proof}

\begin{lemma}
\label{resolve-lemma-blowup-curves}
令 $X$ 为诺特概形。令 $Y_i \subset X$，$i = 1, \ldots, n$，
为维数 $1$ 的整闭子概形，并各自满足引理 \ref{resolve-lemma-resolve-curve}
的等价条件。则存在沿闭点的有限爆破序列
$$
X_n \to X_{n - 1} \to \ldots \to X_1 \to X
$$
使得严格变换 $Y'_i \subset X_n$（即 $Y_i$ 在 $X_n$ 中的严格变换）
是两两不交的正则曲线。
\end{lemma}

\begin{proof}
由引理 \ref{resolve-lemma-blowup-curve}，可假设 $Y_i$ 是正则曲线，其中
$i = 1, \ldots, n$。对每个 $i \not = j$ 以及
$p \in Y_i \cap Y_j$，都有不变量
$m_p(Y_i \cap Y_j)$（\ref{resolve-equation-multiplicity}）。
若这些数的最大值 $> 1$，则在最大值达到的所有点 $p$ 处爆破，
可使它减小（引理 \ref{resolve-lemma-blowup-nonsingular-curves-meeting-at-point}）。
若最大值为 $1$，则在所有这些点 $p$ 处爆破，使用同一引理即可分离这些曲线。
\end{proof}

\noindent
当曲线包含在正则曲面上时，我们通常希望将它化为正规交叉除子。

\begin{lemma}
\label{resolve-lemma-turn-into-effective-Cartier}
令 $X$ 为维数 $2$ 的正则概形，令 $Z \subset X$ 为真闭子概形。
存在沿闭点的爆破序列
$$
X_n \to \ldots \to X_1 \to X
$$
使得逆像 $Z_n$（即 $Z$ 在 $X_n$ 中的逆像）是有效 Cartier 除子。
\end{lemma}

\begin{proof}
令 $D \subset Z$ 为包含在 $Z$ 中的最大有效 Cartier 除子。
则由除子章引理 \ref{divisors-lemma-codim-1-part}，
$\mathcal{I}_Z \subset \mathcal{I}_D$，且商支承在闭点上。
因此可写成 $\mathcal{I}_Z = \mathcal{I}_{Z'} \mathcal{I}_D$，
其中 $Z' \subset X$ 是闭子概形，并且在集合论意义下由有限多个闭点组成。
应用引理 \ref{resolve-lemma-make-ideal-principal}，可得到本引理所述的爆破序列，
使得 $\mathcal{I}_{Z'}\mathcal{O}_{X_n}$ 可逆。这证明了本引理。
\end{proof}

\begin{lemma}
\label{resolve-lemma-embedded-resolution}
令 $X$ 为维数 $2$ 的正则概形。令 $Z \subset X$ 为真闭子概形，
并假设每个不可约分支 $Y \subset Z$，若其维数为 $1$，就满足引理
\ref{resolve-lemma-resolve-curve} 的等价条件。则存在沿闭点的爆破序列
$$
X_n \to \ldots \to X_1 \to X
$$
使得逆像 $Z_n$（即 $Z$ 在 $X_n$ 中的逆像）是有效 Cartier 除子，
并支承在严格正规交叉除子上。
\end{lemma}

\begin{proof}
令 $X' \to X$ 为闭点 $p$ 处的爆破。则逆像 $Z' \subset X'$（即 $Z$ 的逆像）
支承在 $Z$ 的严格变换与例外除子上。例外除子是正则曲线
（引理 \ref{resolve-lemma-blowup}），而严格变换 $Y'$（所对应的不可约分支为 $Y$）
或者等于 $Y$，或者是 $Y$ 在 $p$ 处的爆破。因此在此过程中不会产生新的
维数 $1$ 的奇异分支。于是由引理 \ref{resolve-lemma-turn-into-effective-Cartier}
与引理 \ref{resolve-lemma-blowup-curves}，可假设 $Z$ 是有效 Cartier 除子，
且所有不可约分支 $Y$（它们属于 $Z$）都正则。
（当然，不能假设各不可约分支两两不交，因为每次在 $Z$ 的一点处爆破，
都会向 $Z$ 添加一个新的不可约分支，即例外除子。）

\medskip\noindent
假设 $Z$ 是有效 Cartier 除子，其不可约分支 $Y_i$ 都正则。
对每个 $i \not = j$ 以及 $p \in Y_i \cap Y_j$，都有不变量
$m_p(Y_i \cap Y_j)$（\ref{resolve-equation-multiplicity}）。
若这些数的最大值 $> 1$，则在最大值达到的所有点 $p$ 处爆破，
可使它减小（引理 \ref{resolve-lemma-blowup-nonsingular-curves-meeting-at-point}）
（注意，“新”的不变量 $m_{q_i}(Y'_i \cap E)$ 始终为 $1$）。
若最大值为 $1$，并且例如有 $p \in Y_1 \cap \ldots \cap Y_r$，
其中某个 $r > 2$，而该点不位于任何其他分支上，则在 $p$ 处爆破后，
可知 $Y'_1, \ldots, Y'_r$ 不会在 $p$ 上方的点相交；并且
$m_{q_i}(Y'_i, E) = 1$，其中 $Y'_i \cap E = \{q_i\}$。
因此，继续在多于 $3$ 个分支（这些分支属于 $Z$）相交的点处爆破，便达到如下情形：
对每个闭点 $p \in X$，或者
(a) 没有曲线 $Y_i$ 经过 $p$；
(b) 恰有一条曲线 $Y_i$ 经过 $p$，且 $\mathcal{O}_{Y_i, p}$ 正则；
或者 (c) 恰有两条曲线 $Y_i$、$Y_j$ 经过 $p$，局部环
$\mathcal{O}_{Y_i, p}$、$\mathcal{O}_{Y_j, p}$ 正则，且
$m_p(Y_i \cap Y_j) = 1$。
这意味着 $\sum Y_i$ 是正则曲面 $X$ 上的严格正规交叉除子；见
平展态射章引理 \ref{etale-lemma-strict-normal-crossings}。
\end{proof}






\section{收缩例外曲线}
\label{resolve-section-minus-one}

\noindent
令 $X$ 为诺特概形。令 $E \subset X$ 为满足下列性质的闭子概形：
\begin{enumerate}
\item $E$ 是 $X$ 上的有效 Cartier 除子；
\item 存在域 $k$ 以及概形同构 $\mathbf{P}^1_k \to E$；
\item 法层 $\mathcal{N}_{E/X}$ 拉回为
$\mathcal{O}_{\mathbf{P}^1}(-1)$。
\end{enumerate}
这样的闭子概形称为{\it 第一类例外曲线}。

\medskip\noindent
令 $X'$ 为诺特概形，令 $x \in X'$ 为闭点，并使
$\mathcal{O}_{X', x}$ 是维数 $2$ 的正则局部环。令
$b : X \to X'$ 为 $X'$ 在 $x$ 处的爆破。此时例外纤维
$E \subset X$ 是第一类例外曲线。这由引理 \ref{resolve-lemma-blowup} 得到。

\medskip\noindent
问题：每条第一类例外曲线是否都能作为上述爆破的纤维得到？
换言之，是否总存在概形的固有态射 $X \to X'$，使得
$E$ 映到闭点 $x \in X'$，使得 $\mathcal{O}_{X', x}$
是维数 $2$ 的正则局部环，并且 $X$ 是 $X'$ 在 $x$ 处的爆破？
若是如此，就称{\it 存在 $E$ 的收缩}。

\begin{lemma}
\label{resolve-lemma-factor-through-contraction}
令 $X$ 为诺特概形。令 $E \subset X$ 为第一类例外曲线。
若存在收缩 $X \to X'$，且它收缩 $E$，则它具有如下泛性质：
对每个态射 $\varphi : X \to Y$，若 $\varphi(E)$ 为一点，
就存在唯一分解 $X \to X' \to Y$，它给出 $\varphi$。
\end{lemma}

\begin{proof}
令 $b : X \to X'$ 为 $E$ 的收缩。作为拓扑空间，$X'$ 是 $X$
关于将 $E$ 的所有点等同为一点这一关系所得的商。事实上，$b$ 固有
（除子章引理 \ref{divisors-lemma-blowing-up-projective} 与
态射章引理 \ref{morphisms-lemma-locally-projective-proper}）且满，
因而给出拓扑空间的商映射（拓扑章引理
\ref{topology-lemma-closed-morphism-quotient-topology}）。
另一方面，典范映射 $\mathcal{O}_{X'} \to b_*\mathcal{O}_X$
是同构。确实，在像点 $x \in X'$（即 $E$ 的像）的补集上这是显然的，
而在 $x$ 处的茎上，该映射由引理 \ref{resolve-lemma-cohomology-of-blowup}
第 (4) 项可知是同构。因此，环空间对 $(X', \mathcal{O}_{X'})$
可由 $X$ 构造如下：先取拓扑空间的商，再赋予它
$b_*\mathcal{O}_X$ 作为结构层。

\medskip\noindent
给定 $\varphi$，令 $\varphi' : X' \to Y$ 为满足
$\varphi = \varphi' \circ b$ 的唯一拓扑空间映射。于是映射
$$
\varphi^\sharp : \varphi^{-1}\mathcal{O}_Y =
b^{-1}((\varphi')^{-1}\mathcal{O}_Y) \to \mathcal{O}_X
$$
伴随于映射
$$
(\varphi')^\sharp :
(\varphi')^{-1}\mathcal{O}_Y \to b_*\mathcal{O}_X = \mathcal{O}_{X'}
$$
于是 $(\varphi', (\varphi')^\sharp)$ 是从 $X'$ 到 $Y$ 的环空间态射，
并给出所需分解。由于 $\varphi$ 是局部环空间态射，$\varphi'$ 也是。
事实上，只需验证映射 $\mathcal{O}_{Y, y} \to \mathcal{O}_{X', x}$
是局部的，其中 $y \in Y$ 是 $E$ 在 $\varphi$ 下的像。确实，
元素 $f \in \mathfrak m_y$ 拉回为 $X$ 上在 $E$ 的每一点都为零的函数，
故 $f$ 到 $X'$ 的拉回是在 $x$ 于 $X'$ 中的一个邻域上定义且具有同一性质的函数。
显然，该函数必在 $x$ 处消失，正如所需。
\end{proof}

\begin{lemma}
\label{resolve-lemma-contraction-unique}
令 $X$ 为诺特概形。令 $E \subset X$ 为第一类例外曲线。
若存在 $E$ 的收缩，则它在唯一同构意义下唯一。
\end{lemma}

\begin{proof}
这由引理 \ref{resolve-lemma-factor-through-contraction} 的泛性质立即得到。
\end{proof}

\begin{lemma}
\label{resolve-lemma-exceptional-first-kind-local}
令 $X$ 为诺特概形。令 $E \subset X$ 为第一类例外曲线。
令 $E_n = nE$，并以 $\mathcal{O}_n$ 表示其结构层。则
$$
A = \lim H^0(E_n, \mathcal{O}_n)
$$
是维数 $2$ 的完备诺特正则局部环，并且
$\Ker(A \to H^0(E_n, \mathcal{O}_n))$ 是其极大理想的 $n$ 次幂。
\end{lemma}

\begin{proof}
回顾，存在同构 $\mathbf{P}^1_k \to E$，使得 $E$ 在 $X$ 中的法层
拉回为 $\mathcal{O}(-1)$。于是 $H^0(E, \mathcal{O}_E) = k$。
以 $\mathcal{O}_n(iE)$ 表示可逆层 $\mathcal{O}_X(iE)$ 到 $E_n$ 的限制，
其中 $n \geq 1$ 且 $i \in \mathbf{Z}$。
回顾，$\mathcal{O}_X(-nE)$ 是 $E_n$ 的理想层。因此，对 $d \geq 0$
得到短正合列
$$
0 \to \mathcal{O}_E(-(d + n)E) \to
\mathcal{O}_{n + 1}(-dE) \to
\mathcal{O}_n(-dE) \to 0
$$
由于 $\mathcal{O}_E(-(d + n)E) = \mathcal{O}_{\mathbf{P}^1_k}(d + n)$，
对所有 $d \geq 0$ 与 $n \geq 1$，其第一上同调群消失。
从而系统 $H^0(E_n, \mathcal{O}_n(-dE))$ 的转移映射都是满射。
当 $d = 0$ 时，得到由环的满射构成的逆系统，且每个转移映射的核都是幂零理想。
因此 $A = \lim H^0(E_n, \mathcal{O}_n)$ 是剩余域为 $k$ 的局部环，
其极大理想为
$$
\lim \Ker(H^0(E_n, \mathcal{O}_n) \to H^0(E, \mathcal{O}_E)) =
\lim H^0(E_n, \mathcal{O}_n(-E))
$$
在此核中选取 $x, y$，使其映到一组 $k$-基；该基属于
$H^0(E, \mathcal{O}_E(-E)) = H^0(\mathbf{P}^1_k, \mathcal{O}(1))$。
于是 $x^d, x^{d - 1}y, \ldots, y^d$ 是
$\lim H^0(E_n, \mathcal{O}_n(-dE))$ 的元素，并映到
$H^0(E, \mathcal{O}_E(-dE)) = H^0(\mathbf{P}^1_k, \mathcal{O}(d))$
的一组基。由此可见，$A$ 关于由诸核
$$
I_n = \Ker(A \longrightarrow H^0(E_n, \mathcal{O}_n))
$$
所定义的线性拓扑是分离且完备的。我们有
$x, y \in I_1$、$I_d I_{d'} \subset I_{d + d'}$，并且
$I_d/I_{d + 1}$ 是自由 $k$-模，其基为
$x^d, x^{d - 1}y, \ldots, y^d$。我们将证明 $I_d = (x, y)^d$。事实上，若
$z_e \in I_e$ 且 $e \geq d$，则可写成
$$
z_e = a_{e, 0} x^d + a_{e, 1} x^{d - 1}y + \ldots + a_{e, d}y^d + z_{e + 1}
$$
其中 $a_{e, j} \in (x, y)^{e - d}$，而 $z_{e + 1} \in I_{e + 1}$；
这是由我们对 $I_d/I_{d + 1}$ 的描述得到的。因此，从某个
$z = z_d \in I_d$ 出发，可递归地进行这一过程，得到
$$
z = \sum\nolimits_{e \geq d} \sum\nolimits_j a_{e, j} x^{d - j} y^j
$$
其中某些 $a_{e, j} \in (x, y)^{e - d}$。于是
$a_j = \sum_{e \geq d} a_{e, j}$ 存在（由完备性以及
$a_{e, j} \in I_{e - d}$），并且
$z = \sum a_{e, j} x^{d - j} y^j$。故 $I_d = (x, y)^d$。
于是 $A$ 关于 $(x, y)$-进拓扑完备。由代数章引理
\ref{algebra-lemma-completion-Noetherian}，$A$ 是诺特的。
显然其维数为 $2$，这是由 $(x, y)^d/(x, y)^{d + 1}$ 的描述以及代数章命题
\ref{algebra-proposition-dimension} 得到的。
由于极大理想由两个元素生成，该局部环是正则的。
\end{proof}

\begin{lemma}
\label{resolve-lemma-contraction}
令 $X$ 为诺特概形。令 $E \subset X$ 为第一类例外曲线。
若存在态射 $f : X \to Y$，使得
\begin{enumerate}
\item $Y$ 诺特；
\item $f$ 固有；
\item $f$ 将 $E$ 映到一点 $y$，该点属于 $Y$；
\item $f$ 在不属于 $E$ 的每一点处拟有限；
\end{enumerate}
则存在 $E$ 的收缩，并且它就是 $f$ 的 Stein 分解。
\end{lemma}

\begin{proof}
应用态射进阶章定理
\ref{more-morphisms-theorem-stein-factorization-Noetherian}，
得到 Stein 分解 $X \to X' \to Y$。于是 $X \to X'$ 满足本引理的所有假设
（略去一些细节）。因此，在让 $Y$ 由 $X'$ 取代后，还可假设
$f_*\mathcal{O}_X = \mathcal{O}_Y$，并且 $f$ 的纤维几何连通。

\medskip\noindent
假设 $f_*\mathcal{O}_X = \mathcal{O}_Y$，且 $f$ 的纤维几何连通。
注意到 $y \in Y$ 是闭点，因为 $f$ 是闭映射且 $E$ 闭。态射
$f^{-1}(Y \setminus \{y\}) \to Y \setminus \{y\}$（即 $f$ 的限制）是有限态射
（态射进阶章引理 \ref{more-morphisms-lemma-characterize-finite}）。
由于有限态射是仿射的，结合 $f_*\mathcal{O}_X = \mathcal{O}_Y$，
该限制是同构。为证明 $\mathcal{O}_{Y, y}$ 是维数 $2$ 的正则局部环，
考虑概形上同调章引理 \ref{coherent-lemma-formal-functions-stalk} 的同构
$$
\mathcal{O}_{Y, y}^\wedge \longrightarrow
\lim H^0(X \times_Y \Spec(\mathcal{O}_{Y, y}/\mathfrak m_y^n), \mathcal{O})
$$
令 $E_n = nE$，如引理 \ref{resolve-lemma-exceptional-first-kind-local} 所示。
注意
$$
E_n \subset X \times_Y \Spec(\mathcal{O}_{Y, y}/\mathfrak m_y^n)
$$
因为 $E \subset X_y = X \times_Y \Spec(\kappa(y))$。另一方面，
由于在集合论意义下 $E = f^{-1}(\{y\})$（因为 $f$ 的纤维几何连通），
可知概形论纤维 $X_y$ 在概形论意义下包含于 $E_n$，这里取某个 $n > 0$。
确实，对凝聚 $\mathcal{O}_X$-模 $\mathcal{F} = \mathcal{O}_{X_y}$
与理想层 $\mathcal{I}$（即 $E$ 的理想层）应用概形上同调章引理
\ref{coherent-lemma-power-ideal-kills-sheaf}，并使用
$\mathcal{I}^n$ 是 $E_n$ 的理想层。这说明
$$
X \times_Y \Spec(\mathcal{O}_{Y, y}/\mathfrak m_y^m) \subset E_{nm}
$$
因此，上面显示的逆极限等于 $\lim H^0(E_n, \mathcal{O}_n)$，
而由引理 \ref{resolve-lemma-exceptional-first-kind-local}，后者是二维
正则局部环。故 $\mathcal{O}_{Y, y}$ 是二维正则局部环，
因为其完备化具有此性质（代数进阶章引理
\ref{more-algebra-lemma-completion-regular} 与
\ref{more-algebra-lemma-completion-dimension}）。

\medskip\noindent
还需证明 $f : X \to Y$ 是爆破 $b : Y' \to Y$，即 $Y$ 在 $y$ 处的爆破。
我们鼓励读者自行寻找证明。首先注意，引理
\ref{resolve-lemma-exceptional-first-kind-local} 还蕴含在概形论意义下
$X_y = E$。由于 $E$ 的理想层可逆，这说明
$f^{-1}\mathfrak m_y \cdot \mathcal{O}_X$ 可逆。因此由爆破的泛性质得到分解
$$
X \to Y' \to Y
$$
这给出态射 $f$ 的分解；参见除子章引理
\ref{divisors-lemma-universal-property-blowing-up}。
回顾，由引理 \ref{resolve-lemma-blowup}，例外纤维 $E' \subset Y'$
是第一类例外曲线。令 $g : E \to E'$ 为诱导的态射。
由于 $E'$ 与 $E$ 的余法层都由 $a$、$b$（的拉回）生成，典范映射
$g^*\mathcal{C}_{E'/Y'} \to \mathcal{C}_{E/X}$
是满射（态射章引理 \ref{morphisms-lemma-conormal-functorial}）。
由于二者都可逆，该映射是同构。由于 $\mathcal{C}_{E/X}$ 的次数为正，
$g$ 不可能是常值态射。因此 $g$ 具有有限纤维，故 $g$ 是有限态射
（同上引用）。然而，$Y'$ 在 $E'$ 的所有点处正则（因而正规），
且 $X \to Y'$ 双有理并在 $E'$ 外是同构，所以由簇章引理
\ref{varieties-lemma-modification-normal-iso-over-codimension-1}，
$X \to Y'$ 是同构。
\end{proof}

\begin{lemma}
\label{resolve-lemma-pic-blowup}
令 $b : X \to X'$ 为第一类例外曲线 $E \subset X$ 的收缩。
则有短正合列
$$
0 \to \Pic(X') \to \Pic(X) \to \mathbf{Z} \to 0
$$
其中第一个映射是沿 $b$ 的拉回，第二个映射将 $\mathcal{L}$
映到 $\mathcal{L}$ 在例外曲线 $E$ 上的次数。映射
$n \mapsto \mathcal{O}_X(-nE)$ 给出该正合列的分裂。
\end{lemma}

\begin{proof}
由于 $E = \mathbf{P}^1_k$，$E$ 的 Picard 群是 $\mathbf{Z}$，见除子章引理
\ref{divisors-lemma-Pic-projective-space-UFD}。因此，可将最后一个映射看作
$\mathcal{L} \mapsto \mathcal{L}|_E$。由第一类例外曲线的定义，
$\mathcal{O}_X(E)$ 到 $E$ 的限制之次数为 $-1$。综合这些说明，
只需证明 $\Pic(X') \to \Pic(X)$ 是单射，且其像恰由这样的可逆层组成：
限制为 $\mathcal{O}_E$，而限制取在 $E$ 上。

\medskip\noindent
给定可逆 $\mathcal{O}_{X'}$-模 $\mathcal{L}'$，我们断言映射
$\mathcal{L}' \to b_*b^*\mathcal{L}'$ 是同构。除像点
$x \in X'$（即 $E$ 的像）外，这处处显然。为在 $x$ 处的茎上验证它是同构，
可用 $X'$ 的一个开邻域（含 $x$）代替它，并假设 $\mathcal{L}'$ 是
$\mathcal{O}_{X'}$。于是只需证明映射
$\mathcal{O}_{X'} \to b_*\mathcal{O}_X$ 是同构。这由引理
\ref{resolve-lemma-cohomology-of-blowup} 第 (4) 项得到。

\medskip\noindent
令 $\mathcal{L}$ 为可逆 $\mathcal{O}_X$-模，并满足
$\mathcal{L}|_E = \mathcal{O}_E$。我们断言：(1) $b_*\mathcal{L}$ 可逆；
(2) $b^*b_*\mathcal{L} \to \mathcal{L}$ 是同构。
在 $X' \setminus \{x\}$ 上，断言 (1)、(2) 显然。因此，只需在基变换到
$\Spec(\mathcal{O}_{X', x})$ 后证明 (1)、(2)。计算 $b_*$ 与平坦基变换交换
（概形上同调章引理 \ref{coherent-lemma-flat-base-change-cohomology}），
$b^*$ 及伴随映射的构造也同样如此。但若 $X'$ 是正则局部环的谱，
则由引理 \ref{resolve-lemma-blowup-pic} 中 Picard 群的描述，$\mathcal{L}$ 平凡。
这就证明了断言。

\medskip\noindent
综合前两段证明的断言，可知映射
$\mathcal{L} \mapsto b_*\mathcal{L}$ 是下列映射的逆：
$$
\Pic(X') \longrightarrow \Ker(\Pic(X) \to \Pic(E))
$$
本引理得证。
\end{proof}

\begin{remark}
\label{resolve-remark-pic-blowup}
令 $b : X \to X'$ 为第一类例外曲线 $E \subset X$ 的收缩。
由引理 \ref{resolve-lemma-pic-blowup}，得到等同
$$
\Pic(X) = \Pic(X') \oplus \mathbf{Z}
$$
其中，$\mathcal{L}$ 对应于偶 $(\mathcal{L}', n)$ 当且仅当
$\mathcal{L} = (b^*\mathcal{L}')(-nE)$，即
$\mathcal{L}(nE) = b^*\mathcal{L}'$。事实上，引理
\ref{resolve-lemma-pic-blowup} 的证明表明
$\mathcal{L}' = b_*\mathcal{L}(nE)$。当然，赋值
$\mathcal{L} \mapsto \mathcal{L}'$ 是群同态。
\end{remark}

\begin{lemma}
\label{resolve-lemma-lift-sections-and-h1}
令 $X$ 为诺特概形。令 $E \subset X$ 为第一类例外曲线。
令 $\mathcal{L}$ 为可逆 $\mathcal{O}_X$-模。令 $n$ 为如下整数：
$\mathcal{L}|_E$ 的次数为 $n$；这里将它看作 $\mathbf{P}^1$ 上的可逆模。则
\begin{enumerate}
\item 若 $H^1(X, \mathcal{L}) = 0$ 且 $n \geq 0$，则
$H^1(X, \mathcal{L}(iE)) = 0$ 对所有 $0 \leq i \leq n + 1$ 成立。
\item 若 $n \leq 0$，则
$H^1(X, \mathcal{L}) \subset H^1(X, \mathcal{L}(E))$。
\end{enumerate}
\end{lemma}

\begin{proof}
由除子章引理 \ref{divisors-lemma-Pic-projective-space-UFD}，
注意到 $\mathcal{L}|_E = \mathcal{O}(n)$。对短正合列
$$
0 \to \mathcal{L} \to \mathcal{L}(E) \to \mathcal{L}(E)|_E \to 0,
$$
应用归纳法及相应的上同调长正合列，并使用如下事实：
$H^1(\mathbf{P}^1, \mathcal{O}(d)) = 0$ 对 $d \geq -1$ 成立，而
$H^0(\mathbf{P}^1, \mathcal{O}(d)) = 0$ 对 $d \leq -1$ 成立。
略去一些细节。
\end{proof}

\begin{lemma}
\label{resolve-lemma-contract-ample}
令 $S = \Spec(R)$ 为仿射诺特概形。令 $X \to S$ 为固有态射。
令 $\mathcal{L}$ 为 $X$ 上的丰沛可逆层。令 $E \subset X$
为第一类例外曲线。则
\begin{enumerate}
\item 存在收缩 $b : X \to X'$，它收缩 $E$；
\item $X'$ 在 $S$ 上固有；
\item 可逆 $\mathcal{O}_{X'}$-模 $\mathcal{L}'$ 是丰沛的，
其中 $\mathcal{L}'$ 如注释 \ref{resolve-remark-pic-blowup} 所述。
\end{enumerate}
\end{lemma}

\begin{proof}
令 $n$ 为引理 \ref{resolve-lemma-lift-sections-and-h1} 中
$\mathcal{L}|_E$ 的次数。注意到 $n > 0$，因为 $\mathcal{L}$ 在 $E$ 上丰沛
（簇章引理 \ref{varieties-lemma-ample-curve} 与性质章引理
\ref{properties-lemma-ample-on-closed}）。以 $\mathcal{L}$ 的某个幂代替它后，
可假设 $H^i(X, \mathcal{L}^{\otimes e}) = 0$ 对所有
$i > 0$ 与 $e > 0$ 成立；见概形上同调章引理
\ref{coherent-lemma-vanshing-gives-ample}。最后，再以 $\mathcal{L}$ 的另一个幂代替它，
可假设存在整体截面 $t_0, \ldots, t_n$，它们属于 $\mathcal{L}$ 并定义闭浸入
$\psi : X \to \mathbf{P}^n_S$；见态射章引理
\ref{morphisms-lemma-finite-type-over-affine-ample-very-ample}。

\medskip\noindent
置 $\mathcal{M} = \mathcal{L}(nE)$。则
$\mathcal{M}|_E \cong \mathcal{O}_E$。由于有短正合列
$$
0 \to \mathcal{M}(-E) \to \mathcal{M} \to \mathcal{O}_E \to 0
$$
且 $H^1(X, \mathcal{M}(-E))$ 为零（由引理
\ref{resolve-lemma-lift-sections-and-h1} 以及 $n > 0$），可选取截面
$s_{n + 1}$，它属于 $\mathcal{M}$ 并生成 $\mathcal{M}|_E$。
最后，以 $s_0, \ldots, s_n$ 表示 $\mathcal{M}$ 的截面；它们由上面选取的
截面 $t_0, \ldots, t_n$（属于 $\mathcal{L}$）通过
$\mathcal{L} \subset \mathcal{L}(nE) = \mathcal{M}$ 得到。
合在一起，截面 $s_0, \ldots, s_n, s_{n + 1}$ 生成 $\mathcal{M}$，
并且在 $X$ 的每一点都如此，因而定义态射
$$
\varphi : X \longrightarrow \mathbf{P}^{n + 1}_S
$$
它在 $S$ 上；见构造章引理 \ref{constructions-lemma-projective-space}。

\medskip\noindent
下面将验证引理 \ref{resolve-lemma-contraction} 的条件。验证完成后可知，
Stein 分解 $X \to X' \to \mathbf{P}^{n + 1}_S$（即 $\varphi$ 的 Stein 分解）
就是所需收缩，这证明 (1)。此外，态射
$X' \to \mathbf{P}^{n + 1}_S$ 有限，故 $X'$ 在 $S$ 上固有
（态射章引理 \ref{morphisms-lemma-finite-proper} 与
\ref{morphisms-lemma-composition-proper}）。这证明 (2)。注意，$X'$ 上有丰沛可逆层。
事实上，$\mathcal{M}'$ 是 $\mathcal{O}_{\mathbf{P}^{n + 1}_S}(1)$ 的拉回，
由态射章引理
\ref{morphisms-lemma-pullback-ample-tensor-relatively-ample}，它是丰沛的。
注意，$\mathcal{M}'$ 的拉回为 $\mathcal{M}$，这里拉回到 $X$
（由构造章引理 \ref{constructions-lemma-projective-space}）。
最后，$\mathcal{M} = \mathcal{L}(nE)$。由于在上述论证中用原来
$\mathcal{L}$ 的一个正幂代替了它，由性质章引理
\ref{properties-lemma-ample-power-ample}，本引理 (3) 中的可逆
$\mathcal{O}_{X'}$-模 $\mathcal{L}'$ 在 $X'$ 上丰沛。

\medskip\noindent
容易看出：$\mathbf{P}^{n + 1}_S$ 诺特，且 $\varphi$ 固有。
略去细节。

\medskip\noindent
其次，注意到 $U = X \setminus E$ 的任一点都映到开子概形
$W$，它位于 $\mathbf{P}^{n + 1}_S$ 中；在该开子概形上，前
$n + 1$ 个齐次坐标中至少一个非零。另一方面，$E$ 的任一点都映到
前 $n + 1$ 个齐次坐标全为零的点，特别地映到 $W$ 的补集。
此外，显然有分解
$$
U = \varphi^{-1}(W) \xrightarrow{\varphi|_U}
W \xrightarrow{\text{pr}} \mathbf{P}^n_S
$$
它是 $\psi|_U$ 的分解，其中 $\text{pr}$ 是使用前 $n + 1$ 个坐标的投影，
而 $\psi : X \to \mathbf{P}^n_S$ 是上面选取的嵌入。因此
$\varphi|_U : U \to W$ 拟有限。

\medskip\noindent
最后，考虑映射 $\varphi|_E : E \to \mathbf{P}^{n + 1}_S$。
注意，对任意点 $x \in E$，像 $\varphi(x)$ 的前 $n + 1$ 个坐标全为零，
即态射 $\varphi|_E$ 通过闭子概形 $\mathbf{P}^0_S \cong S$ 分解。
由概形章引理 \ref{schemes-lemma-morphism-into-affine}，态射
$E \to S = \Spec(R)$ 分解为
$E \to \Spec(H^0(E, \mathcal{O}_E)) \to \Spec(R)$。
由于按假设 $H^0(E, \mathcal{O}_E)$ 是域，可知 $E$ 映到
$S \subset \mathbf{P}^{n + 1}_S$ 中的一点，证明完毕。
\end{proof}

\begin{lemma}
\label{resolve-lemma-contract-when-quasi-projective}
令 $S$ 为诺特概形。令 $f : X \to S$ 为有限型态射。
令 $E \subset X$ 为包含在 $f$ 的某个纤维中的第一类例外曲线。
\begin{enumerate}
\item 若 $X$ 在 $S$ 上射影，则存在收缩 $X \to X'$，它收缩 $E$，
并且 $X'$ 在 $S$ 上射影。
\item 若 $X$ 在 $S$ 上拟射影，则存在收缩 $X \to X'$，它收缩 $E$，
并且 $X'$ 在 $S$ 上拟射影。
\end{enumerate}
\end{lemma}

\begin{proof}
两种情形都由引理 \ref{resolve-lemma-contract-ample} 结合关于丰沛可逆模与
（拟）射影态射的标准结果得到。

\medskip\noindent
(1) 的证明。$f$ 的射影性意味着 $f$ 固有且存在一个
$f$-丰沛可逆模 $\mathcal{L}$，见态射章引理
\ref{morphisms-lemma-projective-is-quasi-projective-proper} 与定义
\ref{morphisms-definition-quasi-projective}。令 $U \subset S$ 为包含
$E$ 的像的仿射开集。由引理 \ref{resolve-lemma-contract-ample}，存在收缩
$c : f^{-1}(U) \to V'$，它收缩 $E$；另有丰沛可逆模
$\mathcal{N}'$，定义在 $V'$ 上，
其到 $f^{-1}(U)$ 的拉回等于
$\mathcal{L}(nE)|_{f^{-1}(U)}$。令 $v \in V'$ 为使 $c$ 是在 $v$ 处爆破的闭点。
粘合 $V'$ 与 $X \setminus E$，所沿开集为
$f^{-1}(U) \setminus E = V' \setminus \{v\}$，便得到概形 $X'$，它在 $S$ 上。态射 $c$ 与
$\text{id}_{X \setminus E}$ 粘合为态射 $b : X \to X'$，它就是 $E$ 的收缩。
$U$ 在 $X'$ 中的逆像在 $U$ 上固有。另一方面，$X' \to S$ 在
$v$ 于 $S$ 中的像之补集上的限制，同构于 $X \to S$ 在该开集上的限制。
因此 $X' \to S$ 固有（固有性在基底上是局部的，见态射章引理
\ref{morphisms-lemma-proper-local-on-the-base}）。最后，$\mathcal{N}'$ 与
$\mathcal{L}|_{X \setminus E}$ 在
$f^{-1}(U) \setminus E = V' \setminus \{v\}$ 上限制为同构的可逆模，
故它们粘合为可逆模 $\mathcal{L}'$；该模定义在 $X'$ 上。$\mathcal{L}'$ 在
$U$ 于 $X'$ 中的逆像上的限制是丰沛的，因为 $\mathcal{N}'$ 具有此性质。
对 $S$ 的仿射开集，若它避开 $v$ 的像，则由 $\mathcal{L}$ 具有此性质可知同样成立。
故由态射章引理 \ref{morphisms-lemma-characterize-relatively-ample}，
$\mathcal{L}'$ 是关于 $(X' \to S)$ 相对丰沛的，(1) 得证。

\medskip\noindent
(2) 的证明。由态射章引理
\ref{morphisms-lemma-quasi-projective-open-projective}，可将 $X$ 写成
概形 $\overline{X}$ 的开子概形，而该概形在 $S$ 上射影。由 (1)，存在收缩
$b : \overline{X} \to \overline{X}'$，且 $\overline{X}'$ 在 $S$ 上射影。
令 $X' \subset \overline{X}$ 为 $X \to \overline{X}'$ 的像；
由于 $b$ 在 $E$ 外是同构，这是开集。于是 $X \to X'$ 是所需收缩。
注意，$X'$ 在 $S$ 上拟射影，因为由第 (1) 项证明中的构造，
它具有一个 $S$-相对丰沛可逆模。
\end{proof}

\begin{lemma}
\label{resolve-lemma-regular-dim-2-quasi-projective}
令 $S$ 为诺特概形。令 $f : X \to S$ 为有限型分离态射，
并且 $X$ 是维数 $2$ 的正则概形。则 $X$ 在 $S$ 上拟射影。
\end{lemma}

\begin{proof}
由 Chow 引理（概形上同调章引理
\ref{coherent-lemma-chow-Noetherian}），存在固有态射
$\pi : X' \to X$，它在稠密开集 $U \subset X$ 上是同构，
并且 $X' \to S$ 是 H-拟射影的。由引理
\ref{resolve-lemma-extend-rational-map-blowing-up}，存在沿闭点的爆破序列
$$
X_n \to \ldots \to X_1 \to X_0 = X
$$
以及一个 $S$-态射 $X_n \to X'$，它延拓有理映射 $U \to X'$。
由除子章引理 \ref{divisors-lemma-blowing-up-projective} 与态射章引理
\ref{morphisms-lemma-composition-projective}，注意到 $X_n \to X$ 是射影的。
这由态射章引理 \ref{morphisms-lemma-projective-permanence} 蕴含
$X_n \to X'$ 是射影的。因此由态射章引理
\ref{morphisms-lemma-composition-quasi-projective}，$X_n \to S$ 是拟射影的
（也使用射影态射是拟射影的这一事实，见态射章引理
\ref{morphisms-lemma-projective-quasi-projective}）。由引理
\ref{resolve-lemma-contract-when-quasi-projective}（以及引理
\ref{resolve-lemma-contraction-unique} 所述收缩的唯一性），可知
$X_{n - 1}, \ldots, X_0 = X$ 都在 $S$ 上拟射影，正如所需。
\end{proof}

\begin{lemma}
\label{resolve-lemma-regular-dim-2-projective}
令 $S$ 为诺特概形。令 $f : X \to S$ 为固有态射，
并且 $X$ 是维数 $2$ 的正则概形。则 $X$ 在 $S$ 上射影。
\end{lemma}

\begin{proof}
这由引理 \ref{resolve-lemma-regular-dim-2-quasi-projective} 与态射章引理
\ref{morphisms-lemma-projective-is-quasi-projective-proper} 得到。
\end{proof}






\section{双有理映射的分解}
\label{resolve-section-factorizing}

\noindent
非奇异曲面之间的固有双有理态射由二次变换序列给出。

\begin{lemma}
\label{resolve-lemma-proper-birational-regular-surfaces}
令 $f : X \to Y$ 为整诺特概形之间的固有双有理态射，
这些概形是维数 $2$ 的正则概形。则 $f$ 是沿闭点的爆破序列。
\end{lemma}

\begin{proof}
令 $V \subset Y$ 为 $f$ 在其上是同构的最大开集。则 $V$ 包含所有
余维数 $1$ 的点，而这些点属于 $V$（簇章引理
\ref{varieties-lemma-modification-normal-iso-over-codimension-1}）。
令 $y \in Y$ 为不包含于 $V$ 的闭点。我们要证明，$f$ 通过爆破
$b : Y' \to Y$ 分解，该爆破是 $Y$ 在 $y$ 处的爆破。确实，若这成立，
则纤维 $f^{-1}(y)$ 的至少一个（事实上恰好一个）分支将同构地映到
$Y'$ 中的例外曲线，并且 $X \to Y'$ 的纤维中的曲线数将严格小于
$X \to Y$ 的纤维中的曲线数，故可用归纳法结束证明。略去一些细节。

\medskip\noindent
由引理 \ref{resolve-lemma-extend-rational-map-blowing-up}，存在爆破序列
$$
X' = X_n \to X_{n - 1} \to \ldots \to X_1 \to X_0 = X
$$
它们沿着位于纤维 $f^{-1}(y)$ 上方的闭点进行；并且存在态射
$X' \to Y'$，使得
$$
\xymatrix{
X' \ar[d]_{f'} \ar[r] & X \ar[d]^f \\
Y' \ar[r] & Y
}
$$
交换。我们要证明态射 $X' \to Y'$ 通过 $X$ 分解；于是可对 $n$ 使用归纳法，
归约到如下情形：$X' \to X$ 是 $X$ 在闭点 $x \in X$ 处的爆破，
而该点映到 $y$。

\medskip\noindent
令 $E \subset X'$ 为爆破 $X' \to X$ 的例外纤维。若 $E$ 映到 $Y'$ 中的一点，
则由引理 \ref{resolve-lemma-factor-through-contraction} 得到所需分解。
我们将证明，若非如此便会产生矛盾。确实，若 $f'(E)$ 不是一点，
则 $E' = f'(E)$ 必为 $Y'$ 中的例外曲线。图式为
$$
\xymatrix{
E \ar[r] \ar[d]_g & X' \ar[d]_{f'} \ar[r] & X \ar[d]^f \\
E' \ar[r] & Y' \ar[r] & Y
}
$$
与前面相同地论证，可知 $f'$ 在 $E'$ 的泛点的某个开邻域上是同构。
因此 $g : E \to E'$ 是有限双有理态射。由态射章引理
\ref{morphisms-lemma-extend-across}，$g$ 的逆（一个有理映射）处处有定义，
故 $g$ 是同构。考虑映射
$$
g^*\mathcal{C}_{E'/Y'} \longrightarrow \mathcal{C}_{E/X'}
$$
它来自态射章引理 \ref{morphisms-lemma-conormal-functorial}。
由于源与靶是次数为 $1$ 的可逆模，并且它们位于
$E = E' = \mathbf{P}^1_\kappa$ 上，
并且该映射非零（因为 $f'$ 在 $E$ 的泛点处是同构），可知它是同构。
由态射章引理 \ref{morphisms-lemma-two-immersions}，得到
$\Omega_{X'/Y'}|_E = 0$。这意味着 $f'$ 在 $E$ 的每一点处非分歧
（态射章引理 \ref{morphisms-lemma-unramified-at-point}）。因此 $f'$ 在 $E$
的每一点处拟有限（态射章引理 \ref{morphisms-lemma-unramified-quasi-finite}）。
因此，由簇章引理
\ref{varieties-lemma-modification-normal-iso-over-codimension-1}，最大开集
$V' \subset Y'$（$f'$ 在其上为同构）包含 $E'$。
这又蕴含 $y$ 在 $X'$ 中的逆像是 $E'$。故 $y$ 在 $X$ 中的逆像是 $x$。
因此，由簇章引理
\ref{varieties-lemma-modification-normal-iso-over-codimension-1}，
$x \in X$ 属于 $f$ 在其上为同构的最大开集。这与我们假设 $y$ 不在该开集中矛盾。
\end{proof}

\begin{lemma}
\label{resolve-lemma-birational-regular-surfaces}
令 $S$ 为诺特概形。令 $X$ 与 $Y$ 为 $S$ 上的固有整概形，
并且它们是维数 $2$ 的正则概形。则 $X$ 与 $Y$ 是 $S$-双有理的，
当且仅当存在由 $S$-态射组成的图式
$$
X = X_0 \leftarrow X_1 \leftarrow \ldots \leftarrow X_n = Y_m
\to \ldots \to Y_1 \to Y_0 = Y
$$
其中每个态射都是沿闭点的爆破。
\end{lemma}

\begin{proof}
令 $U \subset X$ 为开集，并令 $f : U \to Y$ 为给定的
$S$-有理映射（它作为 $S$-有理映射可逆）。由引理
\ref{resolve-lemma-extend-rational-map-blowing-up}，可将 $f$ 分解为
$X_n \to \ldots \to X_1 \to X_0 = X$ 与 $f_n : X_n \to Y$。
由于 $X_n$ 在 $S$ 上固有，且 $Y$ 在 $S$ 上分离，态射 $f_n$ 固有。
显然 $f_n$ 是双有理的。因此由引理
\ref{resolve-lemma-proper-birational-regular-surfaces}，$f_n$ 是若干收缩的复合。
略去逆命题的证明。
\end{proof}
% Generated mechanically from the frozen models.tex target.
% Do not edit this generated file; edit the bound locale source instead.

% OK, start here.
%

\setcounter{chapter}{54}
\chapter{半稳定约化}



\phantomsection
\label{models-section-phantom}


\section{引言}
\label{models-section-introduction}

\noindent
本章证明曲线的半稳定约化定理。我们将采用 Artin 与 Winters
在论文 \cite{Artin-Winters} 中的方法。

\medskip\noindent
事实上，无须使用任何奇点消解结果，也可以证明曲线的半稳定约化定理。
具体而言，可以运用 Mumford 所发展的几何不变量理论（GIT），参见
\cite{GIT}，来确立半稳定曲线模空间的存在性与射影性。Gieseker
大力推进了这一方法，并在其讲义 \cite{Gieseker}\footnote{Gieseker
的讲义以代数闭域为基域写成，但同一方法在 $\mathbf{Z}$ 上也适用。}
中证明了完整结果。这是一项相当惊人的成就：人们似乎很难凭直觉相信，
在从未真正研究正维基底上的曲线族的情况下，仍能证明这样的结果。

\medskip\noindent
从历史上看，曲线半稳定约化定理的首个证明见 Deligne 与 Mumford
的论文 \cite{DM}。该证明把问题约化到阿贝尔簇的情形；由于
Grothendieck 等人的工作，这一情形在当时已经解决，参见
\cite{SGA7-I} 与 \cite{SGA7-II}）。阿贝尔簇的半稳定约化定理使用
N\'eron 模型理论，而该理论又建立在对基底上双有理群律的处理之上。

\medskip\noindent
Artin 与 Winters 论文中的方法依靠曲面奇点的消解来获得正则模型。
一旦正则模型的存在性成立，证明便归结为分析特殊纤维的各种可能情形，
再利用域上 $1$ 维概形的 Picard 群中挠子的一个不等式得出结论。
Saito 的论文 \cite{Saito} 中可见类似论证；他直接运用 \'etale 上同调，
并得到更强的结果：可以用惯性群在 $\ell$-进 \'etale 上同调上的作用
来刻画半稳定约化。

\medskip\noindent
证明该定理的另一种途径是运用刚性解析几何技术。这里请读者参见
\cite{vanderPut} 与 \cite{Arzdorf-Wewers}。

\medskip\noindent
Temkin 的论文 \cite{Temkin} 运用赋值论技术（此外还证明了更多结果）；
该文附录 A 还对不同证明以及它们与 $2$ 维概形奇点消解之间的关系
作了很好的综述。

\medskip\noindent
读者或许还愿意参阅 Ahmed Abbes 撰写的另一篇综述论文
\cite{Abbes-ssr}。






\section{线性代数}
\label{models-section-linear-algebra}

\noindent
下面给出稍后将用到的若干引理。

\begin{lemma}
\label{models-lemma-recurring}
\begin{reference}
\cite[Theorem I]{Taussky}
\end{reference}
令 $A = (a_{ij})$ 为复 $n \times n$ 矩阵。
\begin{enumerate}
\item 若 $|a_{ii}| > \sum_{j \not = i} |a_{ij}|$ 对每个 $i$ 都成立，
则 $\det(A)$ 非零。
\item 若存在实向量 $m = (m_1, \ldots, m_n)$，满足 $m_i > 0$，
并且 $|a_{ii} m_i| > \sum_{j \not = i} |a_{ij}m_j|$ 对每个 $i$
都成立，则 $\det(A)$ 非零。
\end{enumerate}
\end{lemma}

\begin{proof}
若 $A$ 如 (1) 且 $\det(A) = 0$，则存在非零向量 $z$ 使得
$Az = 0$。选取 $r$ 使 $|z_r|$ 最大。于是
$$
|a_{rr} z_r| = |\sum\nolimits_{k \not = r} a_{rk}z_k| \leq
\sum\nolimits_{k \not = r} |a_{rk}||z_k| \leq
|z_r| \sum\nolimits_{k \not = r} |a_{rk}| < |a_{rr}||z_r|
$$
这给出矛盾。为证明 (2)，对矩阵 $(a_{ij}m_j)$ 应用 (1)；
其行列式为 $m_1 \ldots m_n \det(A)$。
\end{proof}

\begin{lemma}
\label{models-lemma-recurring-real}
令 $A = (a_{ij})$ 为实 $n \times n$ 矩阵，并且 $a_{ij} \geq 0$
对 $i \not = j$ 成立。令 $m = (m_1, \ldots, m_n)$ 为满足 $m_i > 0$
的实向量。对 $I \subset \{1, \ldots, n\}$，令
$x_I \in \mathbf{R}^n$ 为如下向量：其第 $i$ 个坐标为 $m_i$（若
$i \in I$），否则为 $0$。若
\begin{equation}
\label{models-equation-ineq}
-a_{ii}m_i \geq \sum\nolimits_{j \not = i} a_{ij}m_j
\end{equation}
对每个 $i$ 都成立，则 $\Ker(A)$ 是由满足下列条件的向量 $x_I$
张成的向量空间：
\begin{enumerate}
\item $a_{ij} = 0$ 对 $i \in I$、$j \not \in I$ 成立；
\item 对 $i \in I$，(\ref{models-equation-ineq}) 中等号成立。
\end{enumerate}
\end{lemma}

\begin{proof}
将 $a_{ij}$ 替换为 $a_{ij}m_j$ 后，可假设 $m_i = 1$ 对所有 $i$ 成立。
若 $I \subset \{1, \ldots, n\}$ 满足 (1) 与 (2)，则直接计算可知
$x_I$ 属于 $A$ 的核。反过来，令
$x = (x_1, \ldots, x_n) \in \mathbf{R}^n$ 为 $A$ 的核中的非零向量。
我们将对 $x$ 的非零坐标数作归纳，证明 $x$ 属于满足 (1)、(2) 的向量
$x_I$ 所张成的空间。令 $I \subset \{1, \ldots, n\}$ 为指标 $r$
的集合，其中 $|x_r|$ 取最大值。对 $r \in I$，有
$$
|a_{rr} x_r| = |\sum\nolimits_{k \not = r} a_{rk}x_k| \leq
\sum\nolimits_{k \not = r} a_{rk}|x_k| \leq
|x_r| \sum\nolimits_{k \not = r} a_{rk} \leq |a_{rr}||x_r|
$$
因此处处取等号。特别地，$a_{rk} = 0$ 对 $r \in I$、$k \not \in I$
成立，并且对 $r \in I$，(\ref{models-equation-ineq}) 中等号成立。
于是可将 $x_I$ 的适当倍数从 $x$ 中减去，从而减少非零坐标的数目。
\end{proof}

\begin{lemma}
\label{models-lemma-recurring-symmetric-real}
令 $A = (a_{ij})$ 为对称实 $n \times n$ 矩阵，并且 $a_{ij} \geq 0$
对 $i \not = j$ 成立。令 $m = (m_1, \ldots, m_n)$ 为满足 $m_i > 0$
的实向量。假设
\begin{enumerate}
\item $Am = 0$,
\item 不存在真非空子集 $I \subset \{1, \ldots, n\}$，使得
$a_{ij} = 0$ 对 $i \in I$、$j \not \in I$ 成立。
\end{enumerate}
则 $x^t A x \leq 0$，且等号成立当且仅当 $x = qm$，其中某个
$q \in \mathbf{R}$。
\end{lemma}

\begin{proof}[第一种证明]
将 $a_{ij}$ 替换为 $a_{ij}m_im_j$ 后，可假设 $m_i = 1$ 对所有
$i$ 都成立。条件 (1) 意味着 $-a_{ii} = \sum_{j \not = i} a_{ij}$
对所有 $i$ 都成立。回忆
$x^tAx = \sum_{i, j} x_ia_{ij}x_j$。于是
\begin{align*}
\sum\nolimits_{i \not = j} -a_{ij}(x_j - x_i)^2 & =
\sum\nolimits_{i \not = j} -a_{ij}x_j^2 + 2a_{ij}x_ix_i - a_{ij}x_i^2 \\
& =
\sum\nolimits_j a_{jj} x_j^2 +
\sum\nolimits_{i \not = j} 2a_{ij}x_ix_i +
\sum\nolimits_j a_{jj} x_i^2 \\
& = 2x^tAx
\end{align*}
这显然 $\leq 0$。若等号成立，令 $I$ 为指标 $i$ 的集合，其中
$x_i \not = x_1$。则 $a_{ij} = 0$ 对 $i \in I$、$j \not \in I$
成立。因此由条件 (2)，$I = \{1, \ldots, n\}$，而 $x$
是 $m = (1, \ldots, 1)$ 的倍数。
\end{proof}

\begin{proof}[第二种证明]
由谱定理，矩阵 $A$ 的特征值都是实数。我们断言所有特征值都
$\leq 0$。事实上，由于性质 (1) 意味着
$-a_{ii}m_i = \sum_{j \not = i} a_{ij}m_j$ 对所有 $i$ 都成立，
可知矩阵 $A' = A - \lambda I$（其中 $\lambda > 0$）满足
$|a'_{ii}m_i| > \sum a'_{ij}m_j = \sum |a'_{ij}m_j|$ 对所有 $i$
都成立。因此由引理 \ref{models-lemma-recurring}，$A'$ 可逆。
这蕴含对称双线性形式 $x^tAy$
是半负定的，即 $x^tAx \leq 0$ 对所有 $x$ 都成立。由此，$A$ 的核
等于向量 $x$ 的集合，其中 $x^tAx = 0$。引理
\ref{models-lemma-recurring-real} 对核的描述给出本引理的最后断言。
\end{proof}

\begin{lemma}
\label{models-lemma-orthogonal-direct-sum}
令 $L$ 为有限自由 $\mathbf{Z}$-模，并赋有整值对称正定双线性形式
$\langle\ ,\ \rangle : L \times L \to \mathbf{Z}$。
令 $A \subset L$ 为子模，且 $L/A$ 无挠。置
$B = \{b \in L \mid \langle a, b\rangle = 0,\ \forall a \in A\}$.
则有单射
$$
A^\#/A \leftarrow L/(A \oplus B) \rightarrow B^\#/B
$$
其余核均为 $L^\#/L$ 的商。这里
$A^\# = \{a' \in A \otimes \mathbf{Q} \mid
\langle a, a'\rangle \in \mathbf{Z},\ \forall a \in A\}$
，而 $B$ 与 $L$ 的相应记号作类似定义。
\end{lemma}

\begin{proof}
注意
$L \otimes \mathbf{Q} = A \otimes \mathbf{Q} \oplus B \otimes \mathbf{Q}$
，因为由正定性，该形式在 $A$ 上非退化。以
$\pi_B : L \otimes \mathbf{Q} \to B \otimes \mathbf{Q}$ 表示投影。
由于该形式为整值，$\pi_B(x) \in B^\#$ 对 $x \in L$ 成立。
这给出正合列
$$
0 \to A \to L \xrightarrow{\pi_B} B^\# \to Q \to 0
$$
其中 $Q$ 是 $L \to B^\#$ 的余核。注意 $Q$ 是 $L^\#/L$ 的商，
因为映射 $L^\# \to B^\#$ 是满射：它是 $\mathbf{Z}$-线性对偶于
$B \to L$ 的映射，而后者作为 $\mathbf{Z}$-模映射是分裂的。
模去 $A \oplus B$，得到短正合列
$$
0 \to L/(A \oplus B) \to B^\#/B \to Q \to 0
$$
这就证明了本引理。
\end{proof}

\begin{lemma}
\label{models-lemma-coker}
令 $L_0$、$L_1$ 为有限自由 $\mathbf{Z}$-模，并赋有整值对称正定
双线性形式 $\langle\ ,\ \rangle : L_i \times L_i \to \mathbf{Z}$。
令 $\text{d} : L_0 \to L_1$ 与 $\text{d}^* : L_1 \to L_0$
互为伴随。若 $\langle\ ,\ \rangle$ 在 $L_0$ 上是幺模的，则存在同构
$$
\Phi :
\Coker(\text{d}^*\text{d})_{torsion}
\longrightarrow
\Im(\text{d})^\#/\Im(\text{d})
$$
，其中记号同引理 \ref{models-lemma-orthogonal-direct-sum}。
\end{lemma}

\begin{proof}
令 $x \in L_0$ 表示 $\Coker(\text{d}^*\text{d})$ 中的一个挠类。
则对某个 $a > 0$，可以写成 $ax = \text{d}^*\text{d}(y)$。
对任意 $z \in \Im(\text{d})$，设 $z = \text{d}(y')$，则
$$
\langle (1/a)\text{d}(y), z \rangle =
\langle (1/a)\text{d}(y), \text{d}(y') \rangle =
\langle x, y' \rangle \in \mathbf{Z}
$$
因此 $(1/a)\text{d}(y) \in \Im(\text{d})^\#$。定义
$\Phi(x) = (1/a)\text{d}(y) \bmod \Im(\text{d})$.
略去 $\Phi$ 良定义、可加且为单射的证明。

\medskip\noindent
为证明 $\Phi$ 为满射，令 $z \in \Im(\text{d})^\#$。
则 $z$ 定义线性映射 $L_0 \to \mathbf{Z}$，其规则为
$x \mapsto \langle z, \text{d}(x)\rangle$。由于假设 $L_0$ 上的配对为幺模的，
可找到 $x' \in L_0$，使得
$\langle x', x \rangle = \langle z, \text{d}(x)\rangle$
对所有 $x \in L_0$ 成立。特别地，$x'$ 与 $\Ker(\text{d})$ 的配对为零。
由于 $\Im(\text{d}^*\text{d}) \otimes \mathbf{Q}$ 是
$\Ker(\text{d}) \otimes \mathbf{Q}$ 的正交补，这意味着 $x'$ 在
$\Coker(\text{d}^*\text{d})$ 中定义一个挠类。我们断言
$\Phi(x') = z$。确实，写成 $a x' = \text{d}^*\text{d}(y)$，其中某个
$y \in L_0$ 且 $a > 0$。对任意 $x \in L_0$，得到
$$
\langle z, \text{d}(x)\rangle =
\langle x', x \rangle =
\langle (1/a)\text{d}^*\text{d}(y), x \rangle =
\langle (1/a)\text{d}(y),\text{d}(x) \rangle
$$
因此 $z = \Phi(x')$，证明完成。
\end{proof}

\begin{lemma}
\label{models-lemma-recurring-symmetric-integer}
令 $A = (a_{ij})$ 为对称整数 $n \times n$ 矩阵，并且
$a_{ij} \geq 0$ 对 $i \not = j$ 成立。令 $m = (m_1, \ldots, m_n)$
为满足 $m_i > 0$ 的整数向量。假设
\begin{enumerate}
\item $Am = 0$,
\item 不存在真非空子集 $I \subset \{1, \ldots, n\}$，使得
$a_{ij} = 0$ 对 $i \in I$、$j \not \in I$ 成立。
\end{enumerate}
令 $e$ 为序对 $(i, j)$ 的数目，其中 $i < j$ 且 $a_{ij} > 0$。
则对素数 $\ell$，若它与所有 $a_{ij}$ 和 $m_i$ 互素，便有
$$
\dim_{\mathbf{F}_\ell}(\Coker(A)[\ell]) \leq 1 - n + e
$$
\end{lemma}

\begin{proof}
由引理 \ref{models-lemma-recurring-symmetric-real}，$A$ 的秩为 $n - 1$。
复合映射
$$
\mathbf{Z}^{\oplus n} \xrightarrow{\text{diag}(m_1, \ldots, m_n)}
\mathbf{Z}^{\oplus n} \xrightarrow{(a_{ij})}
\mathbf{Z}^{\oplus n} \xrightarrow{\text{diag}(m_1, \ldots, m_n)}
\mathbf{Z}^{\oplus n}
$$
的矩阵为 $a_{ij}m_im_j$。由对 $\ell$ 的限制，首尾两个映射的余核
都是阶与 $\ell$ 互素的挠群，因此只须对以 $a_{ij}m_im_j$ 为矩阵元的
矩阵证明本引理。故可假设 $m = (1, \ldots, 1)$。

\medskip\noindent
假设 $m = (1, \ldots, 1)$。置 $V = \{1, \ldots, n\}$ 以及
$E = \{(i, j) \mid i < j\text{ 且 }a_{ij} > 0\}$。对
$e = (i, j) \in E$，置 $a_e = a_{ij}$。定义映射 $s, t : E \to V$，
令 $s(i, j) = i$ 以及 $t(i, j) = j$。置
$\mathbf{Z}(V) = \bigoplus_{i \in V} \mathbf{Z}i$ 以及
$\mathbf{Z}(E) = \bigoplus_{e \in E} \mathbf{Z}e$.
在 $\mathbf{Z}(V)$ 与 $\mathbf{Z}(E)$ 上定义对称正定整值配对如下：
$$
\langle i, i \rangle = 1\text{ 对 }i \in V, \quad
\langle e, e \rangle = a_e\text{ 对 }e \in E
$$
并令其他所有配对均为零。考虑映射
$$
\text{d} : \mathbf{Z}(V) \to \mathbf{Z}(E), \quad
i \longmapsto
\sum\nolimits_{e \in E,\ s(e) = i} e - \sum\nolimits_{e \in E,\ t(e) = i} e
$$
以及
$$
\text{d}^*(e) = a_e(s(e) - t(e))
$$
计算表明
$$
\langle d(x), y\rangle = \langle x, \text{d}^*(y) \rangle
$$
，换言之，$\text{d}$ 与 $\text{d}^*$ 互为伴随。接下来计算
\begin{align*}
\text{d}^*\text{d}(i)
& = 
\text{d}^*(
\sum\nolimits_{e \in E,\ s(e) = i} e - \sum\nolimits_{e \in E,\ t(e) = i} e) \\
& =
\sum\nolimits_{e \in E,\ s(e) = i} a_e(s(e) - t(e)) -
\sum\nolimits_{e \in E,\ t(e) = i} a_e(s(e) - t(e))
\end{align*}
$i$ 在 $\text{d}^*\text{d}(i)$ 中的系数为
$$
\sum\nolimits_{e \in E,\ s(e) = i} a_e +
\sum\nolimits_{e \in E,\ t(e) = i} a_e = - a_{ii}
$$
，因为 $\sum_j a_{ij} = 0$，而 $j \not = i$ 在
$\text{d}^*\text{d}(i)$ 中的系数为 $-a_{ij}$。
因此 $\Coker(A) = \Coker(\text{d}^*\text{d})$。

\medskip\noindent
考虑包含关系
$$
\Im(\text{d}) \oplus \Ker(\text{d}^*) \subset \mathbf{Z}(E)
$$
左端是正交直和。显然 $\mathbf{Z}(E)/\Ker(\text{d}^*)$ 无挠。
我们断言 $\mathbf{Z}(E)/\Im(\text{d})$ 也无挠。事实上，设
$x = \sum x_e e \in \mathbf{Z}(E)$ 且 $a > 1$，并且
$ax = \text{d}y$ 对某个 $y = \sum y_i i \in \mathbf{Z}(V)$ 成立。
则 $a x_e = y_{s(e)} - y_{t(e)}$。由性质 (2)，所有 $y_i$ 模 $a$
具有同一个同余类。因此可以写成
$y = a y' + (y_1, y_1, \ldots, y_1)$。由于
$\text{d}(y_1, y_1, \ldots, y_1) = 0$，可知
$x = \text{d}(y')$，这正是所需证明的结论。

\medskip\noindent
因此可以应用引理 \ref{models-lemma-orthogonal-direct-sum}，得到单射
$$
\Im(\text{d})^\#/\Im(\text{d}) \leftarrow
\mathbf{Z}(E)/(\Im(\text{d}) \oplus \Ker(\text{d}^*)) \rightarrow
\Ker(\text{d}^*)^\#/\Ker(\text{d}^*)
$$
其余核被所有 $a_e$ 的乘积零化，而该乘积与 $\ell$ 互素。由于
$\Ker(\text{d}^*)$ 是秩为 $1 - n + e$ 的格，可知只要证明存在同构
$$
\Phi : M_{torsion} \longrightarrow \Im(\text{d})^\#/\Im(\text{d})
$$
即可完成证明。这已在引理 \ref{models-lemma-coker} 中证明。
\end{proof}









\section{数值型}
\label{models-section-numerical-types}

\noindent
部分论证将涉及以下数据结构的组合性质。

\begin{definition}
\label{models-definition-type}
一个{\it 数值型} $T$ 由下列数据给出：
$$
n, m_i, a_{ij}, w_i, g_i
$$
其中 $n \geq 1$ 为整数，而 $m_i$、$a_{ij}$、$w_i$、$g_i$
是在 $1 \leq i, j \leq n$ 范围内满足下列条件的整数：
\begin{enumerate}
\item $m_i > 0$, $w_i > 0$, $g_i \geq 0$,
\item 矩阵 $A = (a_{ij})$ 对称，且 $a_{ij} \geq 0$ 对
$i \not = j$ 成立；
\item 不存在真非空子集 $I \subset \{1, \ldots, n\}$，使得
$a_{ij} = 0$ 对 $i \in I$、$j \not \in I$ 成立；
\item 对每个 $i$，都有 $\sum_j a_{ij}m_j = 0$；
\item $w_i | a_{ij}$.
\end{enumerate}
\end{definition}

\noindent
显然，这类结构用起来多少有些麻烦；但它恰好出现在光滑几何连通曲线的
固有正则模型的特殊纤维中。当然，我们只在重排指标的意义下考察这些型。

\begin{definition}
\label{models-definition-type-equivalent}
称两个数值型 $n, m_i, a_{ij}, w_i, g_i$ 与
$n', m'_i, a'_{ij}, w'_i, g'_i$ 为{\it 等价型}，如果存在一个置换
$\sigma$，作用于 $\{1, \ldots, n\}$，使得
$m_i = m'_{\sigma(i)}$、$a_{ij} = a'_{\sigma(i)\sigma(j)}$、
$w_i = w'_{\sigma(i)}$，并且 $g_i = g'_{\sigma(i)}$。
\end{definition}

\noindent
一个数值型具有亏格。

\begin{lemma}
\label{models-lemma-genus}
令 $n, m_i, a_{ij}, w_i, g_i$ 为数值型。则表达式
$$
g = 1 + \sum m_i(w_i(g_i - 1) - \frac{1}{2} a_{ii})
$$
是整数。
\end{lemma}

\begin{proof}
为证明 $g$ 是整数，须证明 $\sum a_{ii}m_i$ 为偶数。如下模 $2$
计算即可看出这一点：
\begin{align*}
\sum\nolimits_i a_{ii} m_i 
& \equiv
\sum\nolimits_{i,\ m_i\text{ 为奇数}} a_{ii}m_i \\
& \equiv
\sum\nolimits_{i,\ m_i\text{ 为奇数}} \sum\nolimits_{j \not = i} a_{ij}m_j \\
& \equiv
\sum\nolimits_{i,\ m_i\text{ 为奇数}}
\sum\nolimits_{j \not = i,\ m_j\text{ 为奇数}} a_{ij}m_j \\
& \equiv
\sum\nolimits_{i < j,\ m_i\text{ 且 }m_j\text{ 为奇数}} a_{ij}(m_i + m_j) \\
& \equiv
0
\end{align*}
这里使用了 $a_{ij} = a_{ji}$，以及 $\sum_j a_{ij}m_j = 0$
对所有 $i$ 成立。
\end{proof}

\begin{definition}
\label{models-definition-genus}
称 $n, m_i, a_{ij}, w_i, g_i$ 为{\it 亏格 $g$ 的数值型}，如果
$g = 1 + \sum m_i(w_i(g_i - 1) - \frac{1}{2} a_{ii})$ 是引理
\ref{models-lemma-genus} 中的整数。
\end{definition}

\noindent
下面将证明（引理 \ref{models-lemma-genus-nonnegative}）亏格几乎总是
$\geq 0$。不过，确实存在负亏格的数值型。

\begin{lemma}
\label{models-lemma-irreducible}
令 $n, m_i, a_{ij}, w_i, g_i$ 为亏格 $g$ 的数值型。
若 $n = 1$，则 $a_{11} = 0$ 且 $g = 1 + m_1w_1(g_1 - 1)$。
此外，所有这样的数值型可分类如下：
\begin{enumerate}
\item 若 $g < 0$，则 $g_1 = 0$，而亏格为 $g$ 且 $n = 1$ 的可能数值型
只有有限多个，它们对应于分解式 $m_1w_1 = 1 - g$。
\item 若 $g = 0$，则 $m_1 = 1$、$w_1 = 1$、$g_1 = 0$，
如引理 \ref{models-lemma-genus-zero} 所述。
\item 若 $g = 1$，则推出 $g_1 = 1$，但 $m_1, w_1$ 可以是任意正整数；
这就是情形 (\ref{models-item-one})，见引理 \ref{models-lemma-genus-one}。
\item 若 $g > 1$，则 $g_1 > 1$，而亏格为 $g$ 且 $n = 1$ 的可能数值型
只有有限多个，它们对应于分解式 $m_1w_1(g_1 - 1) = g - 1$。
\end{enumerate}
\end{lemma}

\begin{proof}
该引理不证自明。
\end{proof}

\begin{lemma}
\label{models-lemma-diagonal-negative}
令 $n, m_i, a_{ij}, w_i, g_i$ 为亏格 $g$ 的数值型。
若 $n > 1$，则 $a_{ii} < 0$ 对所有 $i$ 成立。
\end{lemma}

\begin{proof}
对矩阵 $A$ 应用引理 \ref{models-lemma-recurring-symmetric-real}。
\end{proof}

\begin{lemma}
\label{models-lemma-minus-one}
令 $n, m_i, a_{ij}, w_i, g_i$ 为亏格 $g$ 的数值型。
假设 $n > 1$。若指标 $i$ 的贡献
$m_i(w_i(g_i - 1) - \frac{1}{2} a_{ii})$ 对亏格 $g$ 为 $< 0$，
则 $g_i = 0$ 且 $a_{ii} = -w_i$。
\end{lemma}

\begin{proof}
这立即由引理 \ref{models-lemma-diagonal-negative} 以及
$w_i > 0$、$g_i \geq 0$、$w_i | a_{ii}$ 得出。
\end{proof}

\begin{definition}
\label{models-definition-type-minus-one}
令 $n, m_i, a_{ij}, w_i, g_i$ 为数值型。称 $i$ 为一个
{\it $(-1)$-指标}，若 $g_i = 0$ 且 $a_{ii} = -w_i$。
\end{definition}

\noindent
我们可以“收缩”$(-1)$-指标。

\begin{lemma}
\label{models-lemma-contract}
令 $n, m_i, a_{ij}, w_i, g_i$ 为数值型 $T$。
假设 $n$ 是一个 $(-1)$-指标。则存在数值型 $T'$，它由
$n', m'_i, a'_{ij}, w'_i, g'_i$ 给出并满足
\begin{enumerate}
\item $n' = n - 1$,
\item $m'_i = m_i$,
\item $a'_{ij} = a_{ij} - a_{in}a_{jn}/a_{nn}$,
\item $w'_i = w_i/2$，若 $a_{in}/w_n$ 为偶数且
$a_{in}/w_i$ 为奇数；否则 $w'_i = w_i$；
\item $g'_i =
\frac{w_i}{w'_i}(g_i - 1) + 1 + \frac{a_{in}^2 - w_na_{in}}{2w'_iw_n}$.
\end{enumerate}
此外，$g = g'$。
\end{lemma}

\begin{proof}
注意，例如由引理 \ref{models-lemma-irreducible} 可知 $n > 1$，因而
$n' \geq 1$。我们对 $n', m'_i, a'_{ij}, w'_i, g'_i$ 检验定义
\ref{models-definition-type} 的条件 (1)--(5)。

\medskip\noindent
条件 (1) 是立即的。

\medskip\noindent
条件 (2)。$A' = (a'_{ij})$ 的对称性是立即的；并且由引理
\ref{models-lemma-diagonal-negative} 有 $a_{nn} < 0$，故
$a'_{ij} \geq a_{ij} \geq 0$ 对 $i \not = j$ 成立。

\medskip\noindent
条件 (3)。假设 $I \subset \{1, \ldots, n - 1\}$，并且
$a'_{ii'} = 0$ 对 $i \in I$ 与
$i' \in \{1, \ldots, n - 1\} \setminus I$ 成立。于是可见，对每个
$i \in I$ 与 $i' \in I'$ 都有 $a_{in}a_{i'n} = 0$。因此，或者
$a_{in} = 0$ 对所有 $i \in I$ 成立，而
$I \subset \{1, \ldots, n\}$ 与 $T$ 的性质 (3) 矛盾；或者
$a_{i'n} = 0$ 对所有 $i' \in \{1, \ldots, n - 1\} \setminus I$
成立，而 $I \cup \{n\} \subset \{1, \ldots, n\}$ 与 $T$ 的
性质 (3) 矛盾。因此 (3) 对 $T'$ 成立。

\medskip\noindent
条件 (4)。计算得
$$
\sum\nolimits_{j = 1}^{n - 1} a'_{ij}m_j =
\sum\nolimits_{j = 1}^{n - 1}
(a_{ij}m_j - \frac{a_{in}a_{jn}m_j}{a_{nn}}) =
- a_{in}m_n - \frac{a_{in}}{a_{nn}}(-a_{nn}m_n) = 0
$$
，正如所需。

\medskip\noindent
条件 (5)。须证明 $w'_i$ 整除 $a_{in}a_{jn}/a_{nn}$。
这是显然的，因为 $a_{nn} = -w_n$，且 $w_n | a_{jn}$、$w_i | a_{in}$。

\medskip\noindent
为证明 $g = g'$，先写成
\begin{align*}
g
& =
1 + \sum\nolimits_{i = 1}^n m_i(w_i(g_i - 1) - \frac{1}{2}a_{ii}) \\
& =
1 + \sum\nolimits_{i = 1}^{n - 1} m_i(w_i(g_i - 1) - \frac{1}{2}a_{ii})
-\frac{1}{2}m_nw_n \\
& =
1 +  \sum\nolimits_{i = 1}^{n - 1}
m_i(w_i(g_i - 1) - \frac{1}{2}a_{ii} - \frac{1}{2}a_{in})
\end{align*}
与 $g'$ 的表达式比较，可知只须
$$
w'_i(g'_i - 1) - \frac{1}{2}a'_{ii} =
w_i(g_i - 1) - \frac{1}{2}a_{in} - \frac{1}{2}a_{ii}
$$
对 $i \leq n - 1$ 成立。换言之，有
$$
g'_i = \frac{2w_i(g_i - 1) - a_{in} - a_{ii} + a'_{ii} + 2w'_i}{2w'_i} =
\frac{w_i}{w'_i}(g_i - 1) + 1 + \frac{a_{in}^2 - w_na_{in}}{2w'_iw_n}
$$
直接检验可知，这是一个整数 $\geq 0$，若按 (4) 选取 $w'_i$。
\end{proof}

\begin{lemma}
\label{models-lemma-top-genus}
令 $n, m_i, a_{ij}, w_i, g_i$ 为数值型。令 $e$ 为序对
$(i, j)$ 的数目，其中 $i < j$ 且 $a_{ij} > 0$。则表达式
$g_{top} = 1 - n + e$ 为 $\geq 0$。
\end{lemma}

\begin{proof}
若非如此，则 $e < n - 1$，这意味着存在一个 $i$，使得
$a_{ij} = 0$ 对所有 $j \not = i$ 成立。这与定义
\ref{models-definition-type} 的假设 (3) 矛盾。
\end{proof}

\begin{definition}
\label{models-definition-top-genus}
令 $n, m_i, a_{ij}, w_i, g_i$ 为数值型 $T$。{\it $T$ 的拓扑亏格}
是引理 \ref{models-lemma-top-genus} 中的非负整数 $g_{top} = 1 - n + e$。
\end{definition}

\noindent
我们想用拓扑亏格来控制亏格。然而，这并非总是成立，例如引理
\ref{models-lemma-irreducible} 中 $n = 1$ 的数值型便是如此。但它对极小数值型
成立；极小数值型定义如下。

\begin{definition}
\label{models-definition-type-minimal}
称数值型 $n, m_i, a_{ij}, w_i, g_i$（其亏格为 $g$）是{\it 极小的}，
如果不存在 $i$ 满足 $g_i = 0$ 与 $a_{ii} = -w_i$；换言之，
不存在 $(-1)$-指标。
\end{definition}

\noindent
我们将证明，极小型的亏格 $g$（其中 $n > 1$）大于或等于
$\max(1, g_{top})$。

\begin{lemma}
\label{models-lemma-non-irreducible-minimal-type-genus-at-least-one}
若 $n, m_i, a_{ij}, w_i, g_i$ 为满足 $n > 1$ 的极小数值型，
则 $g \geq 1$。
\end{lemma}

\begin{proof}
这是因为 $g = 1 + \sum \Phi_i$，而
$\Phi_i = m_i(w_i(g_i - 1) - \frac{1}{2} a_{ii})$ 由引理
\ref{models-lemma-minus-one} 与极小型的定义为非负。
\end{proof}

\begin{lemma}
\label{models-lemma-genus-nonnegative}
若 $n, m_i, a_{ij}, w_i, g_i$ 为满足 $n > 1$ 的极小数值型，
则 $g \geq g_{top}$。
\end{lemma}

\begin{proof}
只对与固有正则模型相伴的数值型感兴趣的读者可以跳过此证明，
因为稍后我们将在几何情形下重新证明这一结果。可以写成
$$
g_{top} = 1 - n + \frac{1}{2}\sum\nolimits_{a_{ij} > 0} 1 =
1 + \sum\nolimits_i (-1 +
\frac{1}{2}\sum\nolimits_{j \not = i,\ a_{ij} > 0} 1) 
$$
另一方面，有
\begin{align*}
g & =
1 + \sum m_i(w_i(g_i - 1) - \frac{1}{2} a_{ii}) \\
& =
1 + \sum m_iw_ig_i - \sum m_iw_i +
\frac{1}{2} \sum\nolimits_{i \not = j} a_{ij}m_j \\
& =
1 + \sum\nolimits_i
m_iw_i(-1 + g_i + \frac{1}{2} \sum\nolimits_{j \not = i} \frac{a_{ij}}{w_i})
\end{align*}
第一个等号来自定义，第二个等号使用 $\sum a_{ij}m_j = 0$，
最后一个等号使用 $a_{ij} = a_{ji}$ 并交换求和次序。与 $g_{top}$
的公式比较，可知若
$$
\Psi_i =
m_iw_i(-1 + g_i + \frac{1}{2} \sum\nolimits_{j \not = i} \frac{a_{ij}}{w_i})
- (-1 + \frac{1}{2}\sum\nolimits_{j \not = i,\ a_{ij} > 0} 1)
$$
为 $\geq 0$ 对每个 $i$ 都成立，本引理便成立。然而，事实未必如此。
下面分析哪些指标可能满足 $\Psi_i < 0$。首先注意
$$
(-1 + g_i + \frac{1}{2}\sum\nolimits_{j \not = i} \frac{a_{ij}}{w_i}) \geq
(-1 + \frac{1}{2}\sum\nolimits_{j \not = i,\ a_{ij} > 0} 1)
$$
，因为 $a_{ij}/w_i$ 是非负整数。由于 $m_iw_i$ 是正整数，
可知 $\Psi_i \geq 0$，只要 $m_iw_i = 1$，或该不等式的左端
为 $\geq 0$；后者发生在 $g_i > 0$ 时，或 $a_{ij} > 0$ 对至少两个指标
$j$ 成立时，或存在一个 $j$ 使得 $a_{ij} > w_i$ 时。因此
$$
P = \{i : \Psi_i < 0\}
$$
恰为满足下列条件的指标 $i$ 的集合：$m_iw_i > 1$、$g_i = 0$；
$a_{ij} > 0$ 对唯一的 $j$ 成立；并且 $a_{ij} = w_i$ 对这个 $j$ 成立。
此外，
$$
i \in P \Rightarrow \Psi_i = \frac{1}{2}(-m_iw_i + 1)
$$
证明策略是说明：给定 $i \in P$，可以从 $\Psi_j$ 借用一点，其中
$j$ 是 $i$ 的相邻指标，即 $a_{ij} > 0$。然而，这并不能直接奏效，
因为 $j$ 可能是满足 $\Psi_j = 0$ 的指标。

\medskip\noindent
考虑集合
$$
Z = \{j : g_j = 0\text{ 且 }
j\text{ has exactly two neighbours }i, k\text{ 其中 }
a_{ij} = w_j = a_{jk}\}
$$
对 $j \in Z$，有 $\Psi_j = 0$。我们考虑序列
$M = (i, j_1, \ldots, j_s)$，其中 $s \geq 0$、$i \in P$、
$j_1, \ldots, j_s \in Z$，并且
$a_{ij_1} > 0, a_{j_1j_2} > 0, \ldots, a_{j_{s - 1}j_s} > 0$。
若我们的数值型由两个属于 $P$ 的指标组成，或更一般地，由两个属于
$P$ 的指标及全都属于 $Z$ 的其余指标组成，则 $g_{top} = 0$，
由引理 \ref{models-lemma-non-irreducible-minimal-type-genus-at-least-one}
即可得证。我们可以并且确实排除这些情形。

\medskip\noindent
令 $M = (i, j_1, \ldots, j_s)$ 为极大序列，并令 $k$ 为 $j_s$
的第二个相邻指标。（若 $s = 0$，则 $k$ 是 $i$ 的唯一相邻指标。）
由极大性，$k \not \in Z$；由刚才所述，$k \not \in P$。注意
$w_i = a_{ij_1} = w_{j_1} = a_{j_1j_2} = \ldots = w_{j_s} =
a_{j_sk}$。考察数值型的定义，可知
\begin{align*}
m_ia_{ii} + m_{j_1}w_i & = 0,\\
m_iw_i + m_{j_1}a_{j_1j_1} + m_{j_2}w_i & = 0,\\
\ldots & \ldots \\
m_{j_{s - 1}}w_i + m_{j_s}a_{j_sj_s} + m_kw_i & = 0
\end{align*}
第一个等式蕴含 $m_{j_1} \geq 2m_i$，因为该数值型是极小的。
于是第二个等式蕴含 $m_{j_2} \geq 3m_i$，依此类推。无论如何，
可得 $m_k \geq 2m_i$（也包括 $s = 0$ 的情形）。

\medskip\noindent
令 $k$ 为如下指标：存在 $t > 0$ 以及两两不同的上述极大序列
$M_1, \ldots, M_t$，其中 $M_b = (i_b, j_{b, 1}, \ldots, j_{b, s_b})$，
使得 $k$ 是 $j_{b, s_b}$ 的相邻指标，对 $b = 1, \ldots, t$ 成立。
我们将证明 $\Phi_j + \sum_{b = 1, \ldots, t} \Phi_{i_b} \geq 0$。
由上述讨论，这将完成本引理的证明。令 $M$ 为出现在 $M_b$ 中的指标
之并，其中 $b = 1, \ldots, t$。写成
$$
\Psi_k =
-\sum\nolimits_{b = 1, \ldots, t} \Psi_{i_b} + \Psi_k'
$$
其中
\begin{align*}
\Psi_k' & =
m_kw_k\left(-1 + g_k +
\frac{1}{2} \sum\nolimits_{b = 1, \ldots t}
(\frac{a_{kj_{b, s_b}}}{w_k} - \frac{m_{i_b}w_{i_b}}{m_kw_k}) +
\frac{1}{2} \sum\nolimits_{l \not = k,\ l \not \in M}
\frac{a_{kl}}{w_k}
\right) \\
&
-\left(
-1 + \frac{1}{2}\sum\nolimits_{l \not = k,\ l \not \in M,\ a_{kl} > 0} 1
\right)
\end{align*}
假设 $\Psi_k' < 0$，以求矛盾。若集合
$\{l : l \not = k,\ l \not \in M,\ a_{kl} > 0\}$ 为空，
则 $\{1, \ldots, n\} = M \cup \{k\}$ 且 $g_{top} = 0$，
因为此时 $e = n - 1$；结果由引理
\ref{models-lemma-non-irreducible-minimal-type-genus-at-least-one} 成立。
因此可假设至少存在一个这样的 $l$，它对第一个括号内的和贡献
$(1/2)a_{kl}/w_k \geq 1/2$。对每个 $b = 1, \ldots, t$，有
$$
\frac{a_{kj_{b, s_b}}}{w_k} - \frac{m_{i_b}w_{i_b}}{m_kw_k} =
\frac{w_{i_b}}{w_k}(1 - \frac{m_{i_b}}{m_k})
$$
该表达式为 $\geq \frac{1}{2}$，因为由上一段有
$m_k \geq 2m_{i_b}$；并且它为 $\geq 1$，如果 $w_k < w_{i_b}$。
由此，$\Psi_k' < 0$ 蕴含 $g_k = 0$。若 $t \geq 2$，或
$t = 1$ 且 $w_k < w_{i_1}$，则 $\Psi_k' \geq 0$（这里使用了
上面所示的 $l$ 的存在性），这也给出矛盾。因此 $t = 1$ 且
$w_k = w_{i_1}$。若关于 $l$ 的和中至少有两个非零项，或有一个这样的
$k$ 且 $a_{kl} > w_k$，则同样有 $\Psi_k' \geq 0$。最后一种可能是
$t = 1$，并且存在一个 $l$ 使 $a_{kl} = w_k$。这是不允许的，
因为它将意味着 $k \in Z$，与 $M_1$ 的极大性矛盾。
\end{proof}

\begin{lemma}
\label{models-lemma-minus-two}
令 $n, m_i, a_{ij}, w_i, g_i$ 为亏格 $g$ 的数值型。
假设 $n > 1$。若 $i$ 使得贡献
$m_i(w_i(g_i - 1) - \frac{1}{2} a_{ii})$
对亏格 $g$ 为 $0$，则 $g_i = 0$ 且 $a_{ii} = -2w_i$。
\end{lemma}

\begin{proof}
这立即由引理 \ref{models-lemma-diagonal-negative} 以及
$w_i > 0$、$g_i \geq 0$、$w_i | a_{ii}$ 得出。
\end{proof}

\noindent
满足这一关系的指标在极小数值型的结构中起着重要作用。
因此我们为它们命名。

\begin{definition}
\label{models-definition-type-minus-two}
令 $n, m_i, a_{ij}, w_i, g_i$ 为亏格 $g$ 的数值型。
称 $i$ 为一个{\it $(-2)$-指标}，若 $g_i = 0$ 且
$a_{ii} = -2w_i$。
\end{definition}

\noindent
给定亏格 $g$ 的极小数值型，$(-2)$-指标恰好是在下列公式中
不对亏格贡献正数的指标：
$$
g = 1 + \sum m_i(w_i(g_i - 1) - \frac{1}{2} a_{ii})
$$
因此，正如稍后将看到的，要控制与 $(-2)$-指标相伴的量会有些棘手。

\begin{remark}
\label{models-remark-genus-equality}
令 $n, m_i, a_{ij}, w_i, g_i$ 为满足 $n > 1$ 的极小数值型。
引理 \ref{models-lemma-genus-nonnegative} 中可以有等式 $g = g_{top}$。
例如，若 $m_i = w_i = 1$，且 $g_i = 0$ 对所有 $i$ 成立，
并且 $a_{ij} \in \{0, 1\}$ 对 $i < j$ 成立。
\end{remark}


\section{数值型的 Picard 群}
\label{models-section-picard-group}

\noindent
定义如下。

\begin{definition}
\label{models-definition-picard-group}
令 $n, m_i, a_{ij}, w_i, g_i$ 为数值型 $T$。{\it $T$ 的 Picard 群}
是矩阵 $(a_{ij}/w_i)$ 的余核；更确切地，
$$
\Pic(T) =
\Coker\left(
\mathbf{Z}^{\oplus n} \to \mathbf{Z}^{\oplus n},\quad
e_i
\mapsto
\sum \frac{a_{ij}}{w_j}e_j
\right)
$$
其中 $e_i$ 表示第 $i$ 个标准基向量，它属于 $\mathbf{Z}^{\oplus n}$。
\end{definition}

\begin{lemma}
\label{models-lemma-picard-rank-1}
令 $n, m_i, a_{ij}, w_i, g_i$ 为数值型 $T$。
$T$ 的 Picard 群是秩为 $1$ 的有限生成阿贝尔群。
\end{lemma}

\begin{proof}
若 $n = 1$，则 $A = (a_{ij})$ 是零矩阵，结论显然。若 $n > 1$，
则矩阵 $A$ 的秩为 $n - 1$；这可由引理 \ref{models-lemma-recurring-real}
或引理 \ref{models-lemma-recurring-symmetric-real} 得出。当然，按 $1/w_i$
缩放各行不影响秩。这就证明了本引理。
\end{proof}

\begin{lemma}
\label{models-lemma-picard-T-and-A}
令 $n, m_i, a_{ij}, w_i, g_i$ 为数值型 $T$。
则 $\Pic(T) \subset \Coker(A)$，其中 $A = (a_{ij})$。
\end{lemma}

\begin{proof}
由于 $\Pic(T)$ 是 $(a_{ij}/w_i)$ 的余核，可知存在交换图
$$
\xymatrix{
0 \ar[r] &
\mathbf{Z}^{\oplus n} \ar[rr]_A & &
\mathbf{Z}^{\oplus n} \ar[rr] & &
\Coker(A) \ar[r] & 0 \\
0 \ar[r] &
\mathbf{Z}^{\oplus n} \ar[rr]^{(a_{ij}/w_i)} \ar[u]_{\text{id}} & &
\mathbf{Z}^{\oplus n} \ar[rr] \ar[u]_{\text{diag}(w_1, \ldots, w_n)} & &
\Pic(T) \ar[r] \ar[u] & 0
}
$$
，其两行均正合。由蛇形引理可得 $\Pic(T) \subset \Coker(A)$。
\end{proof}

\begin{lemma}
\label{models-lemma-contract-picard-group}
令 $n, m_i, a_{ij}, w_i, g_i$ 为数值型 $T$。
假设 $n$ 是一个 $(-1)$-指标。令 $T'$ 为引理
\ref{models-lemma-contract} 中构造的数值型。存在单射
$$
\Pic(T) \to \Pic(T')
$$
，其余核为初等阿贝尔 $2$-群。
\end{lemma}

\begin{proof}
回忆 $n' = n - 1$。令 $e_i$、相应地 $e'_i$ 为第 $i$ 个基向量，
分别属于 $\mathbf{Z}^{\oplus n}$、相应地 $\mathbf{Z}^{\oplus n - 1}$。
首先记
$$
q : \mathbf{Z}^{\oplus n} \to \mathbf{Z}^{\oplus n - 1},
\quad e_n \mapsto 0\text{ 且 }e_i \mapsto e'_i\text{ 对 }i \leq n - 1
$$
并置
$$
p : \mathbf{Z}^{\oplus n} \to \mathbf{Z}^{\oplus n - 1},\quad
e_n \mapsto \sum\nolimits_{j = 1}^{n - 1} \frac{a_{nj}}{w'_j} e'_j
\text{ 且 }
e_i \mapsto \frac{w_i}{w'_i} e'_i\text{ 对 }i \leq n - 1
$$
计算（略去）表明存在交换图
$$
\xymatrix{
\mathbf{Z}^{\oplus n} \ar[rr]_{(a_{ij}/w_i)} \ar[d]_q & &
\mathbf{Z}^{\oplus n} \ar[d]^p \\
\mathbf{Z}^{\oplus n'} \ar[rr]^{(a'_{ij}/w'_i)} & &
\mathbf{Z}^{\oplus n'}
}
$$
由于上方箭头的余核为 $\Pic(T)$，而下方箭头的余核为 $\Pic(T')$，
得到所需的 Picard 群同态。由于 $\frac{w_i}{w'_i} \in \{1, 2\}$，
可知 $\Pic(T) \to \Pic(T')$ 的余核被 $2$ 零化（因为 $2e'_i$
属于 $p$ 的像，对所有 $i \leq n - 1$ 都成立）。最后证明
$\Pic(T) \to \Pic(T')$ 是单射。令 $L = (l_1, \ldots, l_n)$
表示 $\Pic(T)$ 中映到 $\Pic(T')$ 的零元的元素。由于 $q$ 为满射，
追图可知可以假设 $L$ 属于 $p$ 的核。这意味着
$l_na_{ni}/w'_i + l_iw_i/w'_i = 0$，即 $l_i = - a_{ni}/w_i l_n$。
因此 $L$ 是 $-l_ne_n$ 在映射 $(a_{ij}/w_j)$ 下的像，引理得证。
\end{proof}

\begin{lemma}
\label{models-lemma-picard-group-genus-nonpositive}
令 $n, m_i, a_{ij}, w_i, g_i$ 为数值型 $T$。
若亏格 $g$（属于 $T$）为 $\leq 0$，则 $\Pic(T) = \mathbf{Z}$。
\end{lemma}

\begin{proof}
对 $n$ 作归纳。若 $n = 1$，断言显然。若 $n > 1$，则由引理
\ref{models-lemma-non-irreducible-minimal-type-genus-at-least-one}，$T$ 不是极小的。
用一个等价型替换 $T$ 后，可假设 $n$ 是一个 $(-1)$-指标。
由引理 \ref{models-lemma-contract-picard-group}，得到
$\Pic(T) \subset \Pic(T')$。由引理 \ref{models-lemma-contract} 可知
$T'$ 的亏格等于 $T$ 的亏格，归纳即得结论。
\end{proof}








\section{真子图的分类}
\label{models-section-classify-proper-subgraphs}

\noindent
本节固定一个数值型 $n, m_i, a_{ij}, w_i, g_i$，其亏格为 $g$。
我们将完整列出所有完全由
$(-2)$-指标（定义 \ref{models-definition-type-minus-two}）组成的可能“子图”，
同时分类所有亏格为 $1$ 的极小数值型。换言之，本节将证明
命题 \ref{models-proposition-classify-subgraphs} 和
引理 \ref{models-lemma-genus-one}。

\medskip\noindent
策略如下。令 $n, m_i, a_{ij}, w_i, g_i$ 为亏格为 $g$ 的数值型。
令 $I \subset \{1, \ldots, n\}$ 为一个由 $(-2)$-指标组成的子集，
并且不存在非空真子集 $J \subset I$，使得 $a_{jj'} = 0$，其中
$j \in J$、$j' \in I \setminus J$。我们按基数 $|I|$ 对上述 $I$
作归纳。若 $I = \{i\}$ 仅含 $1$ 个指标，则 $m_i$、
$a_{ii}$ 和 $w_i$ 只受到如下约束：定义 \ref{models-definition-type}
给出的 $w_i | a_{ii}$，以及引理 \ref{models-lemma-diagonal-negative}
给出的 $a_{ii} < 0$；这就是归纳起点。在归纳步骤中，先把归纳假设
用于子集 $I' \subset I$，其大小满足 $|I'| < |I|$。这会对可能的
$m_i, a_{ij}, w_i$（其中 $i, j \in I$）施加若干约束。
特别地，因为 $|I'| < |I| \leq n$，由 $\sum a_{ij}m_j = 0$ 和
引理 \ref{models-lemma-recurring-symmetric-real} 可知，子矩阵
$(a_{ij})_{i, j \in I'}$ 是负定的，且其行列式的符号为 $(-1)^m$。
对余下的每一种可能性，计算 $(a_{ij})_{i, j \in I}$ 的行列式。
若其符号为 $-(-1)^{|I|}$，则 Sylvester 定理表明矩阵
$(a_{ij})_{i, j \in I}$ 并非半负定，因而可以排除该情形。
若其符号为 $(-1)^{|I|}$，则 $|I| < n$，我们暂且断定该情形
可以作为真子图出现，并把它列入本节的某个引理。若行列式为 $0$，
则再次由引理 \ref{models-lemma-recurring-symmetric-real} 必有 $|I| = n$，
且 $g = 0$。在这些情形中，我们会实际求出所有可能的
$m_i, a_{ij}, w_i$（$i, j \in I$），并将其列入引理
\ref{models-lemma-genus-one}。完成这一论证后，我们也就得到所有亏格为
$1$ 且满足 $n > 1$ 的极小数值型：由亏格公式以及上一节的讨论，
它们必然完全由 $(-2)$-指标组成，因而必会出现在归纳过程中。
最后，读者可以逐一查看引理 \ref{models-lemma-genus-one} 中列出的可能性，
从而看出本节列为可能情形的所有真子图配置，实际上都已经出现在
亏格为 $1$ 的数值型中。

\medskip\noindent
设 $i$ 和 $j$ 是 $(-2)$-指标，并满足 $a_{ij} > 0$。
由引理 \ref{models-lemma-recurring-symmetric-real}，矩阵 $A = (a_{ij})$
是半负定的，故矩阵
$$
\left(
\begin{matrix}
-2w_i & a_{ij} \\
a_{ij} & -2w_j
\end{matrix}
\right)
$$
是负定的，除非 $n = 2$。情形 $n = 2$ 确实可能发生：此时行列式
$4w_1w_2 - a_{12}^2$ 为零。利用 $\text{lcm}(w_1, w_2)$ 整除
$a_{12}$，读者很容易看出只有如下可能：
$$
(w_1, w_2, a_{12}) = (w, w, 2w), (w, 4w, 4w), \text{ 或 }(4w, w, 4w)
$$
注意，情形 $(4w, w, 4w)$ 可由情形 $(w, 4w, 4w)$ 通过交换指标
$i, j$ 得到。在这些情形中 $g = 1$。这就得到情形
(\ref{models-item-two-cycle}) 和 (\ref{models-item-up4})，它们列于引理
\ref{models-lemma-genus-one}。现在假设 $n > 2$。上述矩阵的行列式
$4w_iw_j - a_{ij}^2$ 为 $> 0$，所以 $a_{ij}^2/w_iw_j < 4$。
另一方面，已知 $\text{lcm}(w_i, w_j) | a_{ij}$，因而
$a_{ij}^2/w_iw_j$ 是整数。因此
$a_{ij}^2/w_iw_j \in \{1, 2, 3\}$，并且 $w_i | w_j$，
或反之。于是有如下可能性：
$$
(w_1, w_2, a_{12}) = (w, w, w), (w, 2w, 2w), (w, 3w, 3w),
(2w, w, 2w), \text{ 或 }(3w, w, 3w)
$$
注意，情形 $(2w, w, 2w)$ 可由情形 $(w, 2w, 2w)$ 通过交换指标
$i, j$ 得到；情形 $(3w, w, 3w)$ 与 $(w, 3w, 3w)$ 也同样如此。
前三个解给出情形 (\ref{models-item-A2})、(\ref{models-item-B2}) 和
(\ref{models-item-G2})，它们列于引理 \ref{models-lemma-two-by-two}。
在该引理中，我们还写出了对整数 $m_i$ 和 $m_j$ 的推论；这里利用了
$\sum_l a_{kl}m_l = 0$ 对每个 $k$ 都成立，因而特别推出
$a_{ii}m_i + a_{ij}m_j \leq 0$（当 $k = i$）以及
$a_{ij}m_i + a_{jj}m_j \leq 0$（当 $k = j$）。

\begin{lemma}
\label{models-lemma-two-by-two}
如下形状的真子图的分类：
$$
\xymatrix{
\bullet \ar@{-}[r] & \bullet
}
$$
若 $n > 2$，给定一对指标 $i, j$；它们是 $(-2)$-指标且满足
$a_{ij} > 0$，则适当重排后，$m$、$a$、$w$ 数据如下：
\begin{enumerate}
\item
\label{models-item-A2}
分别为
$$
\left(
\begin{matrix}
m_1 \\
m_2
\end{matrix}
\right),
\quad
\left(
\begin{matrix}
-2w & w \\
w & -2w
\end{matrix}
\right),
\quad
\left(
\begin{matrix}
w \\
w
\end{matrix}
\right)
$$
其中 $w$ 任意，且 $2m_1 \geq m_2$、$2m_2 \geq m_1$；或
\item
\label{models-item-B2}
分别为
$$
\left(
\begin{matrix}
m_1 \\
m_2
\end{matrix}
\right),
\quad
\left(
\begin{matrix}
-2w & 2w \\
2w & -4w
\end{matrix}
\right),
\quad
\left(
\begin{matrix}
w \\
2w
\end{matrix}
\right)
$$
其中 $w$ 任意，且 $m_1 \geq m_2$、$2m_2 \geq m_1$；或
\item
\label{models-item-G2}
分别为
$$
\left(
\begin{matrix}
m_1 \\
m_2
\end{matrix}
\right),
\quad
\left(
\begin{matrix}
-2w & 3w \\
3w & -6w
\end{matrix}
\right),
\quad
\left(
\begin{matrix}
w \\
3w
\end{matrix}
\right)
$$
其中 $w$ 任意，且 $2m_1 \geq 3m_2$、$2m_2 \geq m_1$。
\end{enumerate}
\end{lemma}

\begin{proof}
见上文的讨论。
\end{proof}

\noindent
设 $i$、$j$、$k$ 是三个 $(-2)$-指标，并且 $a_{ij} > 0$、
$a_{jk} > 0$。换言之，指标 $i$ 与 $j$ “相交”，而 $j$ 与 $k$
“相交”。下文将不再特别说明：每一对 $(i, j)$、$(i, k)$、
$(j, k)$ 都是引理 \ref{models-lemma-two-by-two} 所列的情形之一。
矩阵 $A = (a_{ij})$ 由引理 \ref{models-lemma-recurring-symmetric-real}
可知是半负定的，故矩阵
$$
\left(
\begin{matrix}
-2w_i & a_{ij} & a_{ik} \\
a_{ij} & -2w_j & a_{jk} \\
a_{ik} & a_{jk} & -2w_k
\end{matrix}
\right)
$$
是负定的，除非 $n = 3$。情形 $n = 3$ 确实可能发生：此时该矩阵的
行列式\footnote{它是
$-8w_iw_jw_k + 2a_{ij}^2w_k + 2a_{jk}^2w_i + 2a_{ik}^2w_j +
2a_{ij}a_{jk}a_{ik}$.}为零，于是得到整数等式
$$
4 = \frac{a_{ij}^2}{w_iw_j} +
\frac{a_{jk}^2}{w_jw_k} +
\frac{a_{ik}^2}{w_iw_k} +
\frac{a_{ij}a_{ik}a_{jk}}{w_iw_jw_k}
$$
。该等式右端最后一项由其余各项决定，因为
$$
\left(\frac{a_{ij}a_{ik}a_{jk}}{w_iw_jw_k}\right)^2 =
\frac{a_{ij}^2}{w_iw_j} \frac{a_{jk}^2}{w_jw_k} \frac{a_{ik}^2}{w_iw_k}
$$
上文已经看到，$\frac{a_{ij}^2}{w_iw_j}, \frac{a_{jk}^2}{w_jw_k}$
属于 $\{1, 2, 3\}$，而 $\frac{a_{ik}^2}{w_iw_k}$ 属于
$\{0, 1, 2, 3\}$，故只有如下可能：
$$
(\frac{a_{ij}^2}{w_iw_j}, \frac{a_{jk}^2}{w_jw_k}, \frac{a_{ik}^2}{w_iw_k}) =
(1, 1, 1), (1, 3, 0), (2, 2, 0),\text{ 或 } (3, 1, 0)
$$
注意，情形 $(3, 1, 0)$ 可由情形 $(1, 3, 0)$ 通过颠倒指标
$i, j, k$ 的顺序得到。在每一种情形中都有 $g = 1$；读者可在情形
(\ref{models-item-three-cycle})、(\ref{models-item-equal-up3})、
(\ref{models-item-equal-down3})、(\ref{models-item-up-up})、(\ref{models-item-up-down})、
(\ref{models-item-down-up}) 中找到它们，这些情形列于引理
\ref{models-lemma-genus-one}：其中一种对应 $(1, 1, 1)$，两种对应
$(1, 3, 0)$，三种对应 $(2, 2, 0)$。
假设 $n > 3$，则得到整数不等式
$$
4 > \frac{a_{ij}^2}{w_iw_j} + \frac{a_{ik}^2}{w_iw_k} +
\frac{a_{jk}^2}{w_jw_k} + \frac{a_{ij}a_{ik}a_{jk}}{w_iw_jw_k}
$$
。利用上文对这些数的限制，可知只有如下可能：
$$
(\frac{a_{ij}^2}{w_iw_j}, \frac{a_{jk}^2}{w_jw_k}, \frac{a_{ik}^2}{w_iw_k}) =
(1, 1, 0), (1, 2, 0),\text{ 或 }(2, 1, 0)
$$
特别地，$a_{ik} = 0$（回忆我们假设 $a_{ij} > 0$ 且
$a_{jk} > 0$）。注意，情形 $(2, 1, 0)$ 可由情形 $(1, 2, 0)$
通过颠倒指标 $i, j, k$ 的顺序得到。前两个解给出情形
(\ref{models-item-A3})、(\ref{models-item-C3}) 和 (\ref{models-item-B3})，它们列于引理
\ref{models-lemma-three-by-three}；其中还写出了对整数 $m_i$、$m_j$、
$m_k$ 的推论。

\begin{lemma}
\label{models-lemma-three-by-three}
如下形状的真子图的分类：
$$
\xymatrix{
\bullet \ar@{-}[r] &
\bullet \ar@{-}[r] &
\bullet
}
$$
若 $n > 3$，给定三个指标 $i, j, k$，它们都是 $(-2)$-指标，
并且 $a_{ij}, a_{ik}, a_{jk}$ 中至少有两个非零，则适当重排后，
$m$、$a$、$w$ 数据如下：
\begin{enumerate}
\item
\label{models-item-A3}
分别为
$$
\left(
\begin{matrix}
m_1 \\
m_2 \\
m_3
\end{matrix}
\right),
\quad
\left(
\begin{matrix}
-2w & w & 0 \\
w & -2w & w \\
0 & w & -2w
\end{matrix}
\right),
\quad
\left(
\begin{matrix}
w \\
w \\
w
\end{matrix}
\right)
$$
其中 $2m_1 \geq m_2$、$2m_2 \geq m_1 + m_3$、$2m_3 \geq m_2$；或
\item
\label{models-item-C3}
分别为
$$
\left(
\begin{matrix}
m_1 \\
m_2 \\
m_3
\end{matrix}
\right),
\quad
\left(
\begin{matrix}
-2w & w & 0 \\
w & -2w & 2w \\
0 & 2w & -4w
\end{matrix}
\right),
\quad
\left(
\begin{matrix}
w \\
w \\
2w
\end{matrix}
\right)
$$
其中 $2m_1 \geq m_2$、$2m_2 \geq m_1 + 2m_3$、$2m_3 \geq m_2$；或
\item
\label{models-item-B3}
分别为
$$
\left(
\begin{matrix}
m_1 \\
m_2 \\
m_3
\end{matrix}
\right),
\quad
\left(
\begin{matrix}
-4w & 2w & 0 \\
2w & -4w & 2w \\
0 & 2w & -2w
\end{matrix}
\right),
\quad
\left(
\begin{matrix}
2w \\
2w \\
w
\end{matrix}
\right)
$$
其中 $2m_1 \geq m_2$、$2m_2 \geq m_1 + m_3$、$m_3 \geq m_2$。
\end{enumerate}
\end{lemma}

\begin{proof}
见上文的讨论。
\end{proof}

\noindent
设 $i$、$j$、$k$、$l$ 是四个 $(-2)$-指标，并且 $a_{ij} > 0$、
$a_{jk} > 0$、$a_{kl} > 0$。换言之，指标 $i$ 与 $j$ “相交”，
$j$ 与 $k$ “相交”，而 $k$ 与 $l$ “相交”。由引理
\ref{models-lemma-three-by-three} 可知 $a_{ik} = a_{jl} = 0$。
矩阵 $A = (a_{ij})$ 是半负定的，故矩阵
$$
\left(
\begin{matrix}
-2w_i & a_{ij} & 0 & a_{il} \\
a_{ij} & -2w_j & a_{jk} & 0 \\
0 & a_{jk} & -2w_k & a_{kl} \\
a_{il} & 0 & a_{kl} & -2w_l
\end{matrix}
\right)
$$
是负定的，除非 $n = 4$。情形 $n = 4$ 确实可能发生：此时该矩阵的
行列式\footnote{它是
$16w_iw_jw_kw_l - 4a_{ij}^2w_kw_l - 4a_{jk}^2w_iw_l - 4a_{kl}^2w_iw_j -
4a_{il}^2w_jw_k + a_{ij}^2a_{kl}^2 + a_{jk}^2a_{il}^2 -
2a_{ij}a_{il}a_{jk}a_{kl}$.}为零，于是得到非负整数等式
$$
16 +
\frac{a_{ij}^2}{w_iw_j}\frac{a_{kl}^2}{w_kw_l} +
\frac{a_{jk}^2}{w_jw_k}\frac{a_{il}^2}{w_iw_l} =
4\frac{a_{ij}^2}{w_iw_j} + 4\frac{a_{jk}^2}{w_jw_k} + 4\frac{a_{kl}^2}{w_kw_l}
+ 4\frac{a_{il}^2}{w_iw_l} + 2\frac{a_{ij}a_{il}a_{jk}a_{kl}}{w_iw_jw_kw_l}
$$
。该等式右端最后一项由其余各项决定，因为
$$
\left(\frac{a_{ij}a_{il}a_{jk}a_{kl}}{w_iw_jw_kw_l}\right)^2 =
\frac{a_{ij}^2}{w_iw_j} \frac{a_{jk}^2}{w_jw_k}
\frac{a_{kl}^2}{w_kw_l} \frac{a_{il}^2}{w_iw_l}
$$
上文已经看到，
$\frac{a_{ij}^2}{w_iw_j}, \frac{a_{jk}^2}{w_jw_k}, \frac{a_{kl}^2}{w_kw_l}$
属于 $\{1, 2\}$，而 $\frac{a_{il}^2}{w_iw_l}$ 属于 $\{0, 1, 2\}$，
故只有如下解：
$$
(\frac{a_{ij}^2}{w_iw_j}, \frac{a_{jk}^2}{w_jw_k}, \frac{a_{kl}^2}{w_kw_l},
\frac{a_{il}^2}{w_iw_l}) =
(1, 1, 1, 1) \text{ 或 } (2, 1, 2, 0)
$$
且此时 $g = 1$；读者可在情形 (\ref{models-item-four-cycle})、
(\ref{models-item-up-equal-up})、(\ref{models-item-up-equal-down}) 和
(\ref{models-item-down-equal-up}) 中找到它们，这些情形列于引理
\ref{models-lemma-genus-one}。假设 $n > 4$，则得到非负整数不等式
$$
16 +
\frac{a_{ij}^2}{w_iw_j}\frac{a_{kl}^2}{w_kw_l} +
\frac{a_{jk}^2}{w_jw_k}\frac{a_{il}^2}{w_iw_l} >
4\frac{a_{ij}^2}{w_iw_j} + 4\frac{a_{jk}^2}{w_jw_k} + 4\frac{a_{kl}^2}{w_kw_l}
+ 4\frac{a_{il}^2}{w_iw_l} + 2\frac{a_{ij}a_{il}a_{jk}a_{kl}}{w_iw_jw_kw_l}
$$
。利用上文对这些数的限制，可知只有如下可能：
$$
(\frac{a_{ij}^2}{w_iw_j}, \frac{a_{jk}^2}{w_jw_k}, \frac{a_{kl}^2}{w_kw_l},
\frac{a_{il}^2}{w_iw_l}) =
(1, 1, 1, 0), (1, 1, 2, 0), (1, 2, 1, 0), \text{ 或 }(2, 1, 1, 0)
$$
特别地，$a_{il} = 0$（回忆我们假设其余三项非零）。注意，情形
$(2, 1, 1, 0)$ 可由情形 $(1, 1, 2, 0)$ 通过颠倒指标
$i, j, k, l$ 的顺序得到。前三个解给出情形 (\ref{models-item-A4})、
(\ref{models-item-C4})、(\ref{models-item-B4}) 和 (\ref{models-item-F4})，它们列于引理
\ref{models-lemma-four-by-four}；其中还写出了对整数 $m_i$、$m_j$、$m_k$、
$m_l$ 的推论。

\begin{lemma}
\label{models-lemma-four-by-four}
如下形状的真子图的分类：
$$
\xymatrix{
\bullet \ar@{-}[r] &
\bullet \ar@{-}[r] &
\bullet \ar@{-}[r] &
\bullet
}
$$
若 $n > 4$，给定四个 $(-2)$-指标 $i, j, k, l$，并且
$a_{ij}, a_{jk}, a_{kl}$ 非零，则适当重排后，$m$、$a$、$w$ 数据如下：
\begin{enumerate}
\item
\label{models-item-A4}
分别为
$$
\left(
\begin{matrix}
m_1 \\
m_2 \\
m_3 \\
m_4
\end{matrix}
\right),
\quad
\left(
\begin{matrix}
-2w & w & 0 & 0 \\
w & -2w & w & 0 \\
0 & w & -2w & w \\
0 & 0 & w & -2w 
\end{matrix}
\right),
\quad
\left(
\begin{matrix}
w \\
w \\
w \\
w
\end{matrix}
\right)
$$
其中 $2m_1 \geq m_2$、$2m_2 \geq m_1 + m_3$、
$2m_3 \geq m_2 + m_4$，并且 $2m_4 \geq m_3$；或
\item
\label{models-item-C4}
分别为
$$
\left(
\begin{matrix}
m_1 \\
m_2 \\
m_3 \\
m_4
\end{matrix}
\right),
\quad
\left(
\begin{matrix}
-2w & w & 0 & 0 \\
w & -2w & w & 0 \\
0 & w & -2w & 2w \\
0 & 0 & 2w & -4w 
\end{matrix}
\right),
\quad
\left(
\begin{matrix}
w \\
w \\
w \\
2w
\end{matrix}
\right)
$$
其中 $2m_1 \geq m_2$、$2m_2 \geq m_1 + m_3$、
$2m_3 \geq m_2 + 2m_4$，并且 $2m_4 \geq m_3$；或
\item
\label{models-item-B4}
分别为
$$
\left(
\begin{matrix}
m_1 \\
m_2 \\
m_3 \\
m_4
\end{matrix}
\right),
\quad
\left(
\begin{matrix}
-4w & 2w & 0 & 0 \\
2w & -4w & 2w & 0 \\
0 & 2w & -4w & 2w \\
0 & 0 & 2w & -2w 
\end{matrix}
\right),
\quad
\left(
\begin{matrix}
2w \\
2w \\
2w \\
w
\end{matrix}
\right)
$$
其中 $2m_1 \geq m_2$、$2m_2 \geq m_1 + m_3$、
$2m_3 \geq m_2 + m_4$，并且 $m_4 \geq m_3$；或
\item
\label{models-item-F4}
分别为
$$
\left(
\begin{matrix}
m_1 \\
m_2 \\
m_3 \\
m_4
\end{matrix}
\right),
\quad
\left(
\begin{matrix}
-2w & w & 0 & 0 \\
w & -2w & 2w & 0 \\
0 & 2w & -4w & 2w \\
0 & 0 & 2w & -4w 
\end{matrix}
\right),
\quad
\left(
\begin{matrix}
w \\
w \\
2w \\
2w
\end{matrix}
\right)
$$
其中 $2m_1 \geq m_2$、$2m_2 \geq m_1 + 2m_3$、
$2m_3 \geq m_2 + m_4$，并且 $2m_4 \geq m_3$。
\end{enumerate}
\end{lemma}

\begin{proof}
见上文的讨论。
\end{proof}

\noindent
设 $i$、$j$、$k$、$l$ 是四个 $(-2)$-指标，并且 $a_{ij} > 0$、
$a_{ij} > 0$、$a_{il} > 0$。换言之，指标 $i$ 与指标
$j$、$k$、$l$ “相交”。由引理 \ref{models-lemma-three-by-three} 可知
$a_{jk} = a_{jl} = a_{kl} = 0$。矩阵 $A = (a_{ij})$ 是半负定的，
故矩阵
$$
\left(
\begin{matrix}
-2w_i & a_{ij} & a_{ik} & a_{il} \\
a_{ij} & -2w_j & 0 & 0 \\
a_{ik} & 0 & -2w_k & 0 \\
a_{il} & 0 & 0 & -2w_l
\end{matrix}
\right)
$$
是负定的，除非 $n = 4$。情形 $n = 4$ 确实可能发生：此时该矩阵的
行列式\footnote{它是
$16w_iw_jw_kw_l - 4a_{ij}^2w_kw_l - 4a_{ik}^2w_jw_l - 4a_{il}^2w_jw_k$.}
为零，于是得到非负整数等式
$$
4 = \frac{a_{ij}^2}{w_iw_j} + \frac{a_{ik}^2}{w_iw_k} + \frac{a_{il}^2}{w_jw_l}
$$
。上文已经看到，
$\frac{a_{ij}^2}{w_iw_j}, \frac{a_{ik}^2}{w_iw_k}, \frac{a_{il}^2}{w_iw_l}$
属于 $\{1, 2\}$，故适当重排后只有一种可能：$4 = 1 + 1 + 2$。
在每一种情形中都有 $g = 1$；读者可在情形
(\ref{models-item-triple-with-up}) 和 (\ref{models-item-triple-with-down}) 中找到它们，
这些情形列于引理 \ref{models-lemma-genus-one}。假设 $n > 4$，则得到
非负整数不等式
$$
4 > \frac{a_{ij}^2}{w_iw_j} + \frac{a_{ik}^2}{w_iw_k} + \frac{a_{il}^2}{w_jw_l}
$$
。这推出 $\frac{a_{ij}^2}{w_iw_j} =
\frac{a_{ik}^2}{w_iw_k} = \frac{a_{il}^2}{w_jw_l} = 1$，并且
$w_i = w_j = w_k = w_l$。这就得到情形 (\ref{models-item-D4})，它列于引理
\ref{models-lemma-D4}；其中还写出了对整数 $m_i$、$m_j$、$m_k$、$m_l$ 的推论。

\begin{lemma}
\label{models-lemma-D4}
如下形状的真子图的分类：
$$
\xymatrix{
\bullet \ar@{-}[r] & \bullet \ar@{-}[r] \ar@{-}[d] & \bullet \\
& \bullet
}
$$
若 $n > 4$，给定四个 $(-2)$-指标 $i, j, k, l$，并且
$a_{ij}, a_{ik}, a_{il}$ 非零，则适当重排后，$m$、$a$、$w$ 数据如下：
\begin{enumerate}
\item
\label{models-item-D4}
分别为
$$
\left(
\begin{matrix}
m_1 \\
m_2 \\
m_3 \\
m_4
\end{matrix}
\right),
\quad
\left(
\begin{matrix}
-2w & w & w & w \\
w & -2w & 0 & 0 \\
w & 0 & -2w & 0 \\
w & 0 & 0 & -2w 
\end{matrix}
\right),
\quad
\left(
\begin{matrix}
w \\
w \\
w \\
w
\end{matrix}
\right)
$$
其中 $2m_1 \geq m_2 + m_3 + m_4$、$2m_2 \geq m_1$、
$2m_3 \geq m_1$、$2m_4 \geq m_1$。注意，这推出
$m_1 \geq \max(m_2, m_3, m_4)$。
\end{enumerate}
\end{lemma}

\begin{proof}
见上文的讨论。
\end{proof}

\noindent
设 $h$、$i$、$j$、$k$、$l$ 是五个 $(-2)$-指标，并且
$a_{hi} > 0$、$a_{ij} > 0$、$a_{jk} > 0$、$a_{kl} > 0$。
换言之，指标 $h$ 与 $i$ “相交”，$i$ 与 $j$ “相交”，
$j$ 与 $k$ “相交”，而 $k$ 与 $l$ “相交”。应用引理
\ref{models-lemma-three-by-three} 和 \ref{models-lemma-four-by-four} 可知
$a_{hj} = a_{hk} = a_{ik} = a_{il} = a_{jl} = 0$，并且分数
$\frac{a_{hi}^2}{w_hw_i}, \frac{a_{ij}^2}{w_iw_j}, \frac{a_{jk}^2}{w_jw_k},
\frac{a_{kl}^2}{w_kw_l}$ 属于 $\{1, 2\}$，而分数
$\frac{a_{hl}^2}{w_hw_l} \in \{0, 1, 2\}$。
矩阵 $A = (a_{ij})$ 是半负定的，故矩阵
$$
\left(
\begin{matrix}
-2w_h & a_{hi} & 0 & 0 & a_{hl} \\
a_{hi} & -2w_i & a_{ij} & 0 & 0 \\
0 & a_{ij} & -2w_j & a_{jk} & 0 \\
0 & 0 & a_{jk} & -2w_k & a_{kl} \\
a_{hl} & 0 & 0 & a_{kl} & -2w_l
\end{matrix}
\right)
$$
是负定的，除非 $n = 5$。情形 $n = 5$ 确实可能发生：此时该矩阵的
行列式\footnote{它是
$-32w_hw_iw_jw_kw_l +
8a_{hi}^2w_jw_kw_l +
8a_{ij}^2w_hw_kw_l +
8a_{jk}^2w_hw_iw_l +
8a_{kl}^2w_hw_iw_j +
8a_{hl}^2w_iw_jw_k -
2a_{hi}^2a_{jk}^2w_l -
2a_{hi}^2a_{kl}^2w_j -
2a_{ij}^2a_{kl}^2w_h -
2a_{hl}^2a_{ij}^2w_k -
2a_{hl}^2a_{jk}^2w_i +
2a_{hi}a_{ij}a_{jk}a_{kl}a_{hl}
$
.}
为零，于是得到非负整数等式
\begin{align*}
16 + 
\frac{a_{hi}^2}{w_hw_i}\frac{a_{jk}^2}{w_jw_k} +
\frac{a_{hi}^2}{w_hw_i}\frac{a_{kl}^2}{w_kw_l} +
\frac{a_{ij}^2}{w_iw_j}\frac{a_{kl}^2}{w_kw_l} +
\frac{a_{hl}^2}{w_hw_l}\frac{a_{ij}^2}{w_iw_j} +
\frac{a_{hl}^2}{w_hw_l}\frac{a_{jk}^2}{w_jw_k} \\
= 4\frac{a_{hi}^2}{w_hw_i} +
4\frac{a_{ij}^2}{w_iw_j} +
4\frac{a_{jk}^2}{w_jw_k} +
4\frac{a_{kl}^2}{w_kw_l} +
4\frac{a_{hl}^2}{w_hw_l} +
\frac{a_{hi}a_{ij}a_{jk}a_{kl}a_{hl}}{w_hw_iw_jw_kw_l}
\end{align*}
。该等式右端最后一项由其余各项决定，因为
$$
\left(\frac{a_{hi}a_{ij}a_{jk}a_{kl}a_{hl}}{w_hw_iw_jw_kw_l} \right)^2 =
\frac{a_{hi}^2}{w_hw_i} \frac{a_{ij}^2}{w_iw_j}
\frac{a_{jk}^2}{w_jw_k} \frac{a_{kl}^2}{w_kw_l} \frac{a_{hl}^2}{w_hw_l}
$$
由此可知只有如下解：
$$
(\frac{a_{hi}^2}{w_hw_i}, \frac{a_{ij}^2}{w_iw_j},
\frac{a_{jk}^2}{w_jw_k}, \frac{a_{kl}^2}{w_kw_l}, \frac{a_{hl}^2}{w_hw_l}) =
(1, 1, 1, 1, 1), (1, 1, 2, 1, 0), (1, 2, 1, 1, 0), \text{ 或 }(2, 1, 1, 2, 0)
$$
注意，情形 $(1, 2, 1, 1, 0)$ 可由情形 $(1, 1, 2, 1, 0)$ 通过颠倒指标
$h, i, j, k, l$ 的顺序得到。在这些情形中 $g = 1$；读者可在情形
(\ref{models-item-five-cycle})、(\ref{models-item-equal-equal-up-equal})、
(\ref{models-item-equal-equal-down-equal})、(\ref{models-item-up-equal-equal-up})、
(\ref{models-item-up-equal-equal-down}) 和 (\ref{models-item-down-equal-equal-up})
中找到它们，这些情形列于引理 \ref{models-lemma-genus-one}：其中一种对应
$(1, 1, 1, 1, 1)$，两种对应 $(1, 1, 2, 1, 0)$，三种对应
$(2, 1, 1, 2, 0)$。假设 $n > 5$，则得到非负整数不等式
\begin{align*}
16 + 
\frac{a_{hi}^2}{w_hw_i}\frac{a_{jk}^2}{w_jw_k} +
\frac{a_{hi}^2}{w_hw_i}\frac{a_{kl}^2}{w_kw_l} +
\frac{a_{ij}^2}{w_iw_j}\frac{a_{kl}^2}{w_kw_l} +
\frac{a_{hl}^2}{w_hw_l}\frac{a_{ij}^2}{w_iw_j} +
\frac{a_{hl}^2}{w_hw_l}\frac{a_{jk}^2}{w_jw_k} \\
> 4\frac{a_{hi}^2}{w_hw_i} +
4\frac{a_{ij}^2}{w_iw_j} +
4\frac{a_{jk}^2}{w_jw_k} +
4\frac{a_{kl}^2}{w_kw_l} +
4\frac{a_{hl}^2}{w_hw_l} +
\frac{a_{hi}a_{ij}a_{jk}a_{kl}a_{hl}}{w_hw_iw_jw_kw_l}
\end{align*}
。利用上文对这些数的限制，可知只有如下可能：
$$
(\frac{a_{hi}^2}{w_hw_i}, \frac{a_{ij}^2}{w_iw_j},
\frac{a_{jk}^2}{w_jw_k}, \frac{a_{kl}^2}{w_kw_l}, \frac{a_{hl}^2}{w_hw_l}) =
(1, 1, 1, 1, 0), (1, 1, 1, 2, 0), \text{ 或 } (2, 1, 1, 1, 0)
$$
特别地，$a_{hl} = 0$（回忆我们假设其余四项非零）。注意，情形
$(1, 1, 1, 2, 0)$ 可由情形 $(2, 1, 1, 1, 0)$ 通过颠倒指标
$h, i, j, k, l$ 的顺序得到。前两个解给出情形 (\ref{models-item-A5})、
(\ref{models-item-C5}) 和 (\ref{models-item-B5})，它们列于引理
\ref{models-lemma-five-by-five}；其中还写出了对整数 $m_h$、$m_i$、$m_j$、
$m_k$、$m_l$ 的推论。

\begin{lemma}
\label{models-lemma-five-by-five}
如下形状的真子图的分类：
$$
\xymatrix{
\bullet \ar@{-}[r] &
\bullet \ar@{-}[r] &
\bullet \ar@{-}[r] &
\bullet \ar@{-}[r] &
\bullet
}
$$
若 $n > 5$，给定五个 $(-2)$-指标 $h, i, j, k, l$，并且
$a_{hi}, a_{ij}, a_{jk}, a_{kl}$ 非零，则适当重排后，
$m$、$a$、$w$ 数据如下：
\begin{enumerate}
\item
\label{models-item-A5}
分别为
$$
\left(
\begin{matrix}
m_1 \\
m_2 \\
m_3 \\
m_4 \\
m_5
\end{matrix}
\right),
\quad
\left(
\begin{matrix}
-2w & w & 0 & 0 & 0 \\
w & -2w & w & 0 & 0 \\
0 & w & -2w & w & 0 \\
0 & 0 & w & -2w & w \\
0 & 0 & 0 & w & -2w
\end{matrix}
\right),
\quad
\left(
\begin{matrix}
w \\
w \\
w \\
w \\
w
\end{matrix}
\right)
$$
其中 $2m_1 \geq m_2$、$2m_2 \geq m_1 + m_3$、
$2m_3 \geq m_2 + m_4$、$2m_4 \geq m_3 + m_5$，并且
$2m_5 \geq m_4$；或
\item
\label{models-item-C5}
分别为
$$
\left(
\begin{matrix}
m_1 \\
m_2 \\
m_3 \\
m_4 \\
m_5
\end{matrix}
\right),
\quad
\left(
\begin{matrix}
-2w & w & 0 & 0 & 0 \\
w & -2w & w & 0 & 0 \\
0 & w & -2w & w & 0 \\
0 & 0 & w & -2w & 2w \\
0 & 0 & 0 & 2w & -4w
\end{matrix}
\right),
\quad
\left(
\begin{matrix}
w \\
w \\
w \\
w \\
2w
\end{matrix}
\right)
$$
其中 $2m_1 \geq m_2$、$2m_2 \geq m_1 + m_3$、
$2m_3 \geq m_2 + 2m_4$、$2m_4 \geq m_3 + m_5$，并且
$2m_5 \geq m_4$；或
\item
\label{models-item-B5}
分别为
$$
\left(
\begin{matrix}
m_1 \\
m_2 \\
m_3 \\
m_4 \\
m_5
\end{matrix}
\right),
\quad
\left(
\begin{matrix}
-4w & 2w & 0 & 0 & 0 \\
2w & -4w & 2w & 0 & 0 \\
0 & 2w & -4w & 2w & 0 \\
0 & 0 & 2w & -4w & 2w \\
0 & 0 & 0 & 2w & -2w
\end{matrix}
\right),
\quad
\left(
\begin{matrix}
2w \\
2w \\
2w \\
2w \\
w
\end{matrix}
\right)
$$
其中 $2m_1 \geq m_2$、$2m_2 \geq m_1 + m_3$、
$2m_3 \geq m_2 + m_4$、$2m_4 \geq m_3 + m_5$，并且
$m_4 \geq m_3$。
\end{enumerate}
\end{lemma}

\begin{proof}
见上文的讨论。
\end{proof}

\noindent
设 $h$、$i$、$j$、$k$、$l$ 是五个 $(-2)$-指标，并且
$a_{hi} > 0$、$a_{hj} > 0$、$a_{hk} > 0$、$a_{hl} > 0$。
换言之，指标 $h$ 与指标 $i$、$j$、$k$、$l$ “相交”。由引理
\ref{models-lemma-three-by-three} 可知
$a_{ij} = a_{ik} = a_{il} = a_{jk} = a_{jl} = a_{kl} = 0$；再由引理
\ref{models-lemma-D4} 可知，有 $w_h = w_i = w_j = w_k = w_l = w$，其中
$w > 0$ 是整数，并且
$a_{hi} = a_{hj} = a_{hk} = a_{hl} = -2w$。相应的矩阵
$$
\left(
\begin{matrix}
-2w & w & w & w & w \\
w & -2w & 0 & 0 & 0 \\
w & 0 & -2w & 0 & 0 \\
w & 0 & 0 & -2w & 0 \\
w & 0 & 0 & 0 & -2w
\end{matrix}
\right)
$$
是奇异的。因此这种情形只能在 $n = 5$ 且 $g = 1$ 时发生。
读者可在情形 (\ref{models-item-quadruple}) 中找到它，该情形列于引理
\ref{models-lemma-genus-one}。

\begin{lemma}
\label{models-lemma-fourfold}
如下形状的真子图不存在：
$$
\xymatrix{
\bullet \ar@{-}[r] & \bullet \ar@{-}[ld] \ar@{-}[r] \ar@{-}[d] & \bullet \\
\bullet & \bullet
}
$$
若 $n > 5$，则{\bf 不}存在五个 $(-2)$-指标
$h$、$i$、$j$、$k$，使得 $a_{hi} > 0$、$a_{hj} > 0$、
$a_{hk} > 0$、$a_{hl} > 0$。
\end{lemma}

\begin{proof}
见上文的讨论。
\end{proof}

\noindent
设 $h$、$i$、$j$、$k$、$l$ 是五个 $(-2)$-指标，并且
$a_{hi} > 0$、$a_{ij} > 0$、$a_{jk} > 0$、$a_{jl} > 0$。
换言之，指标 $h$ 与 $i$ “相交”，而指标 $j$ 与指标
$i$、$k$、$l$ “相交”。由引理 \ref{models-lemma-D4} 可知
$a_{ik} = a_{il} = a_{kl} = 0$、$w_i = w_j = w_k = w_l = w$，
并且 $a_{ij} = a_{jk} = a_{jl} = w$，其中 $w$ 为某个整数。
把引理 \ref{models-lemma-four-by-four} 应用于两个四元组 $h, i, j, k$ 和
$h, i, j, l$，可知 $a_{hj} = a_{hk} = a_{hl} = 0$，且
$w_h = \frac{1}{2}w$、$w$ 或 $2w$；相应地，$a_{hi} = w$、$w$
或 $2w$。由于 $A$ 是半负定的，可知矩阵
$$
\left(
\begin{matrix}
-2w_h & a_{hi} & 0 & 0 & 0 \\
a_{hi} & -2w & w & 0 & 0 \\
0 & w & -2w & w & w \\
0 & 0 & w & -2w & 0 \\
0 & 0 & w & 0 & -2w
\end{matrix}
\right)
$$
是负定的，除非 $n = 5$。读者计算可得，该矩阵的行列式为 $0$ 的情形是
$w_h = \frac{1}{2}w$ 或 $2w$。
这就得到情形 (\ref{models-item-triple-extended-up}) 和
(\ref{models-item-triple-extended-down})，它们列于引理 \ref{models-lemma-genus-one}。
当 $w_h = w$ 时，得到情形 (\ref{models-item-D5})，它列于引理
\ref{models-lemma-D5}。

\begin{lemma}
\label{models-lemma-D5}
如下形状的真子图的分类：
$$
\xymatrix{
\bullet \ar@{-}[r] & \bullet \ar@{-}[r] &
\bullet \ar@{-}[r] \ar@{-}[d] & \bullet \\
& & \bullet
}
$$
若 $n > 5$，给定五个 $(-2)$-指标 $h, i, j, k, l$，并且
$a_{hi}, a_{ij}, a_{jk}, a_{jl}$ 非零，则适当重排后，
$m$、$a$、$w$ 数据如下：
\begin{enumerate}
\item
\label{models-item-D5}
分别为
$$
\left(
\begin{matrix}
m_1 \\
m_2 \\
m_3 \\
m_4 \\
m_5
\end{matrix}
\right),
\quad
\left(
\begin{matrix}
-2w & w & 0 & 0 & 0 \\
w & -2w & w & 0 & 0 \\
0 & w & -2w & w & w \\
0 & 0 & w & -2w & 0 \\
0 & 0 & w & 0 & -2w
\end{matrix}
\right),
\quad
\left(
\begin{matrix}
w \\
w \\
w \\
w \\
w
\end{matrix}
\right)
$$
其中 $2m_1 \geq m_2$、$2m_2 \geq m_1 + m_3$、
$2m_3 \geq m_2 + m_4 + m_5$、$2m_4 \geq m_3$，并且
$2m_5 \geq m_3$。
\end{enumerate}
\end{lemma}

\begin{proof}
见上文的讨论。
\end{proof}

\noindent
设 $t > 5$，并设 $i_1, \ldots, i_t$ 是 $t$ 个互异的 $(-2)$-指标，
使得 $a_{i_ji_{j + 1}}$ 非零，其中 $j = 1, \ldots, t - 1$。我们对
$t$ 作归纳，证明：若 $n = t$，则得到可能情形
(\ref{models-item-n-cycle})、(\ref{models-item-up-chain-equal-up})、
(\ref{models-item-up-chain-equal-down})、(\ref{models-item-down-chain-equal-up})，
它们列于引理 \ref{models-lemma-genus-one}；若
$n > t$，则得到情形 (\ref{models-item-An})、(\ref{models-item-Cn})、
(\ref{models-item-Bn})，它们列于引理 \ref{models-lemma-long}。
首先，若 $a_{i_1i_t}$ 非零，则由引理 \ref{models-lemma-five-by-five} 的结果
立即可知 $w_{i_1} = \ldots = w_{i_t} = w$，并且有
$a_{i_ji_{j + 1}} = w$（$j = 1, \ldots, t - 1$），且
$a_{i_1i_t} = w$。于是向量 $(1, \ldots, 1)$ 属于相应
$t \times t$ 矩阵的核。因此必有 $n = t$，亏格为 $1$，并且得到情形
(\ref{models-item-n-cycle})，它列于引理 \ref{models-lemma-genus-one}。
故可假设 $a_{i_1i_t} = 0$。由归纳假设（若 $t = 6$，则用引理
\ref{models-lemma-five-by-five}），有 $a_{i_ji_k} = 0$，其中 $k > j + 1$。
此外，有 $w_{i_1} = \ldots = w_{i_{t - 1}} = w$，其中 $w$ 为整数，
且 $w_{i_1}, w_{i_t} \in \{\frac{1}{2}w, w, 2w\}$。再者，
$w_{i_1}$、相应地 $w_{i_t}$ 的取值为 $\frac{1}{2}w$、$w$ 或 $2w$，
意味着 $a_{i_1i_2}$、相应地 $a_{i_{t - 1}i_t}$ 的取值为 $w$、$w$
或 $2w$。共有 $9$ 种可能。逐一判断都很容易：
\begin{enumerate}
\item 若 $(w_{i_1}, w_{i_t}) = (\frac{1}{2}w, \frac{1}{2}w)$，则得到情形
(\ref{models-item-up-chain-equal-down})，它列于引理 \ref{models-lemma-genus-one}；
\item 若 $(w_{i_1}, w_{i_t}) = (\frac{1}{2}w, w)$ 或
$(w, \frac{1}{2}w)$，则得到情形 (\ref{models-item-Bn})，它列于引理
\ref{models-lemma-long}；
\item 若 $(w_{i_1}, w_{i_t}) = (\frac{1}{2}w, 2w)$ 或
$(2w, \frac{1}{2}w)$，则得到情形 (\ref{models-item-up-chain-equal-up})，
它列于引理 \ref{models-lemma-genus-one}；
\item 若 $(w_{i_1}, w_{i_t}) = (w, w)$，则得到情形
(\ref{models-item-An})，它列于引理 \ref{models-lemma-long}；
\item 若 $(w_{i_1}, w_{i_t}) = (w, 2w)$ 或 $(2w, w)$，则得到情形
(\ref{models-item-Cn})，它列于引理 \ref{models-lemma-long}；
\item 若 $(w_{i_1}, w_{i_t}) = (2w, 2w)$，则得到情形
(\ref{models-item-down-chain-equal-up})，它列于引理 \ref{models-lemma-genus-one}。
\end{enumerate}

\begin{lemma}
\label{models-lemma-long}
如下形状的真子图的分类：
$$
\xymatrix{
\bullet \ar@{-}[r] &
\bullet \ar@{-}[r] &
\bullet \ar@{..}[r] &
\bullet \ar@{-}[r] &
\bullet \ar@{-}[r] &
\bullet
}
$$
设 $t > 5$ 且 $n > t$。给定 $t$ 个互异的 $(-2)$-指标
$i_1, \ldots, i_t$，使得 $a_{i_ji_{j + 1}}$ 非零，其中
$j = 1, \ldots, t - 1$，则在允许颠倒这些指标次序的意义下，
$a$、$w$ 数据如下：
\begin{enumerate}
\item
\label{models-item-An}
分别满足 $w_{i_1} = w_{i_2} = \ldots = w_{i_t} = w$、
$a_{i_ji_{j + 1}} = w$，且 $a_{i_ji_k} = 0$（当 $k > j + 1$）；或
\item
\label{models-item-Cn}
分别满足 $w_{i_1} = w_{i_2} = \ldots = w_{i_{t - 1}} = w$、
$w_{j_t} = 2w$、$a_{i_ji_{j + 1}} = w$（$j < t - 1$）、
$a_{i_{t - 1}i_t} = 2w$，且 $a_{i_ji_k} = 0$（当 $k > j + 1$）；或
\item
\label{models-item-Bn}
分别满足 $w_{i_1} = w_{i_2} = \ldots = w_{i_{t - 1}} = 2w$、
$w_{j_t} = w$、$a_{i_ji_{j + 1}} = 2w$、
$a_{i_{t - 1}i_t} = 2w$，且 $a_{i_ji_k} = 0$（当 $k > j + 1$）。
\end{enumerate}
\end{lemma}

\begin{proof}
见上文的讨论。
\end{proof}

\noindent
设 $t > 4$，并设 $i_1, \ldots, i_{t + 1}$ 是 $t + 1$ 个互异的
$(-2)$-指标，使得 $a_{i_ji_{j + 1}} > 0$（$j = 1, \ldots, t - 1$），
并且 $a_{j_{t - 1}j_{t + 1}} > 0$。参见引理 \ref{models-lemma-Dn} 中的图。
我们对 $t$ 作归纳，证明：若 $n = t + 1$，则得到可能情形
(\ref{models-item-Dn-extended-up}) 和 (\ref{models-item-Dn-extended-down})，
它们列于引理 \ref{models-lemma-genus-one}；若 $n > t + 1$，则得到情形
(\ref{models-item-Dn})，它列于引理 \ref{models-lemma-Dn}。
由归纳假设（当 $t = 5$ 时用引理 \ref{models-lemma-D5}），可知
$a_{i_ji_k}$ 在指定的非零项之外均为零（$j, k \geq 2$）。此外，有
$w_2 = \ldots = w_{t + 1} = w$，其中 $w$ 为整数，并且非零的
$a_{i_ji_k}$（$j, k \geq 2$）均等于 $w$。
把引理 \ref{models-lemma-long}（若 $t = 5$，则用引理
\ref{models-lemma-five-by-five}）应用于序列 $i_1, \ldots, i_t$ 以及序列
$i_1, \ldots, i_{t - 1}, i_{t + 1}$，可知 $a_{i_1 i_j} = 0$
（$j \geq 3$），且 $w_1$ 等于 $\frac{1}{2}w$、$w$ 或 $2w$；
相应地，$a_{i_1i_2}$ 等于 $w, w, 2w$。共有 $3$ 种可能，
逐一判断都很容易：
\begin{enumerate}
\item 若 $w_1 = \frac{1}{2}w$，则得到情形
(\ref{models-item-Dn-extended-down})，它列于引理 \ref{models-lemma-genus-one}。
\item 若 $w_1 = w$，则得到情形 (\ref{models-item-Dn})，它列于引理
\ref{models-lemma-Dn}。
\item 若 $w_1 = 2w$，则得到情形 (\ref{models-item-Dn-extended-up})，
它列于引理 \ref{models-lemma-genus-one}。
\end{enumerate}

\begin{lemma}
\label{models-lemma-Dn}
如下形状的真子图的分类：
$$
\xymatrix{
\bullet \ar@{-}[r] & \bullet \ar@{..}[r] & \bullet \ar@{-}[r] &
\bullet \ar@{-}[r] \ar@{-}[d] & \bullet \\
& & & \bullet
}
$$
设 $t > 4$ 且 $n > t + 1$。给定 $t + 1$ 个互异的 $(-2)$-指标
$i_1, \ldots, i_{t + 1}$，使得 $a_{i_ji_{j + 1}}$ 非零，其中
$j = 1, \ldots, t - 1$，并且 $a_{i_{t - 1}i_{t + 1}}$ 非零，
则 $a$、$w$ 数据如下：
\begin{enumerate}
\item
\label{models-item-Dn}
分别满足 $w_{i_1} = w_{i_2} = \ldots = w_{i_{t + 1}} = w$，
$a_{i_ji_{j + 1}} = w$（$j = 1, \ldots, t - 1$），
$a_{i_{t - 1}i_{t + 1}} = w$，并且有 $a_{i_ji_k} = 0$；这里指其余
指标对 $(j, k)$，且 $j > k$。
\end{enumerate}
\end{lemma}

\begin{proof}
见上文的讨论。
\end{proof}

\noindent
假设给定 $6$ 个互异的 $(-2)$-指标 $g, h, i, j, k, l$，使得
$a_{gh}, a_{hi}, a_{ij}, a_{jk}, a_{il}$ 非零。参见引理
\ref{models-lemma-E6} 中的图。应用引理 \ref{models-lemma-D5} 可知，必处于引理
\ref{models-lemma-E6} 所述的情形。由于行列式为 $3w^6 > 0$，可知此时绝不可能有
$n = 6$！

\begin{lemma}
\label{models-lemma-E6}
如下形状的真子图的分类：
$$
\xymatrix{
\bullet \ar@{-}[r] & \bullet \ar@{-}[r] & \bullet \ar@{-}[r] \ar@{-}[d] &
\bullet \ar@{-}[r] & \bullet \\
& & \bullet
}
$$
设 $n > 6$。给定 $6$ 个互异的 $(-2)$-指标 $i_1, \ldots, i_6$，
使得 $a_{12}, a_{23}, a_{34}, a_{45}, a_{36}$ 非零，则
$m$、$a$、$w$ 数据如下：
\begin{enumerate}
\item
\label{models-item-E6}
分别为
$$
\left(
\begin{matrix}
m_1 \\
m_2 \\
m_3 \\
m_4 \\
m_5 \\
m_6
\end{matrix}
\right),
\quad
\left(
\begin{matrix}
-2w & w & 0 & 0 & 0 & 0 \\
w & -2w & w & 0 & 0 & 0 \\
0 & w & -2w & w & 0 & w \\
0 & 0 & w & -2w & w & 0 \\
0 & 0 & 0 & w & -2w & 0 \\
0 & 0 & w & 0 & 0 & -2w
\end{matrix}
\right),
\quad
\left(
\begin{matrix}
w \\
w \\
w \\
w \\
w \\
w
\end{matrix}
\right)
$$
其中 $2m_1 \geq m_2$、$2m_2 \geq m_1 + m_3$、
$2m_3 \geq m_2 + m_4 + m_6$、$2m_4 \geq m_3 + m_5$、
$2m_5 \geq m_3$，并且 $2m_6 \geq m_3$。
\end{enumerate}
\end{lemma}

\begin{proof}
见上文的讨论。
\end{proof}

\noindent
设 $t \geq 4$，并设 $i_0, \ldots, i_{t + 1}$ 是 $t + 2$ 个互异的
$(-2)$-指标，使得 $a_{i_ji_{j + 1}} > 0$（$j = 1, \ldots, t - 1$），
并且 $a_{i_0i_2} > 0$、$a_{i_{t - 1}i_{t + 1}} > 0$。
参见引理 \ref{models-lemma-double-triple} 中的图。应用引理 \ref{models-lemma-D5}
和 \ref{models-lemma-Dn} 可知，其余所有 $a_{i_ji_k}$ 均为零（$j < k$）；
并且有 $w_{i_0} = \ldots = w_{i_{t + 1}} = w$，其中 $w$ 为整数，
而 $A$ 中指定的非零非对角元均等于 $w$。计算表明相应矩阵的行列式
为零。因此 $n = t + 2$，并得到情形 (\ref{models-item-double-triple})，
它列于引理 \ref{models-lemma-genus-one}。

\begin{lemma}
\label{models-lemma-double-triple}
如下形状的真子图不存在：
$$
\xymatrix{
\bullet \ar@{-}[r] & \bullet \ar@{..}[r] \ar@{-}[d] &
\bullet \ar@{-}[d] \ar@{-}[r] & \bullet \\
& \bullet & \bullet
}
$$
假设 $t \geq 4$ 且 $n > t + 2$。{\bf 不}存在 $t + 2$ 个互异的
$(-2)$-指标 $i_0, \ldots, i_{t + 1}$，使得
$a_{i_ji_{j + 1}} > 0$（$j = 1, \ldots, t - 1$），并且
$a_{i_0i_2} > 0$、$a_{i_{t - 1}i_{t + 1}} > 0$。
\end{lemma}

\begin{proof}
见上文的讨论。
\end{proof}

\noindent
假设给定 $7$ 个互异的 $(-2)$-指标 $f, g, h, i, j, k, l$，使得数
$a_{fg}, a_{gh}, a_{ij}, a_{jh}, a_{kl}, a_{lh}$ 非零。
参见引理 \ref{models-lemma-E6-completed} 中的图。应用引理 \ref{models-lemma-D5}
可知相应矩阵为
$$
\left(
\begin{matrix}
-2w & w & 0 & 0 & 0 & 0 & 0 \\
w & -2w & w & 0 & 0 & 0 & 0 \\
0 & w & -2w & 0 & w & 0 & w \\
0 & 0 & 0 & -2w & w & 0 & 0 \\
0 & 0 & w & w & -2w & 0 & 0 \\
0 & 0 & 0 & 0 & 0 & -2w & w \\
0 & 0 & w & 0 & 0 & w & -2w
\end{matrix}
\right)
$$
由于行列式为 $0$，可知必有 $n = 7$ 且 $g = 1$，并得到情形
(\ref{models-item-E6-completed})，它列于引理 \ref{models-lemma-genus-one}。

\begin{lemma}
\label{models-lemma-E6-completed}
如下形状的真子图不存在：
$$
\xymatrix{
\bullet \ar@{-}[r] & \bullet \ar@{-}[r] & \bullet \ar@{-}[r] \ar@{-}[d] &
\bullet \ar@{-}[r] & \bullet \\
& & \bullet \ar@{-}[d] \\
& & \bullet
}
$$
假设 $n > 7$。{\bf 不}存在 $7$ 个互异的 $(-2)$-指标
$f, g, h, i, j, k, l$，使得
$a_{fg}, a_{gh}, a_{ij}, a_{jh}, a_{kl}, a_{lh}$ 非零。
\end{lemma}

\begin{proof}
见上文的讨论。
\end{proof}

\noindent
假设给定 $7$ 个互异的 $(-2)$-指标 $f, g, h, i, j, k, l$，使得数
$a_{fg}, a_{gh}, a_{hi}, a_{ij}, a_{jk}, a_{il}$ 非零。
参见引理 \ref{models-lemma-E7} 中的图。应用引理 \ref{models-lemma-D5} 和
\ref{models-lemma-Dn} 可知，必处于引理 \ref{models-lemma-E7} 所述的情形。
由于行列式为 $-8w^7 > 0$，可知此时绝不可能有 $n = 7$！

\begin{lemma}
\label{models-lemma-E7}
如下形状的真子图的分类：
$$
\xymatrix{
\bullet \ar@{-}[r] & \bullet \ar@{-}[r] & \bullet \ar@{-}[r] &
\bullet \ar@{-}[r] \ar@{-}[d] & \bullet \ar@{-}[r] & \bullet \\
& & & \bullet
}
$$
设 $n > 7$。给定 $7$ 个互异的 $(-2)$-指标 $i_1, \ldots, i_7$，
使得 $a_{12}, a_{23}, a_{34}, a_{45}, a_{56}, a_{47}$ 非零，
则 $m$、$a$、$w$ 数据如下：
\begin{enumerate}
\item
\label{models-item-E7}
分别为
$$
\left(
\begin{matrix}
m_1 \\
m_2 \\
m_3 \\
m_4 \\
m_5 \\
m_6 \\
m_7
\end{matrix}
\right),
\quad
\left(
\begin{matrix}
-2w & w & 0 & 0 & 0 & 0 & 0 \\
w & -2w & w & 0 & 0 & 0 & 0 \\
0 & w & -2w & w & 0 & 0 & 0 \\
0 & 0 & w & -2w & w & 0 & w \\
0 & 0 & 0 & w & -2w & w & 0 \\
0 & 0 & 0 & 0 & w & -2w & 0 \\
0 & 0 & 0 & w & 0 & 0 & -2w
\end{matrix}
\right),
\quad
\left(
\begin{matrix}
w \\
w \\
w \\
w \\
w \\
w \\
w
\end{matrix}
\right)
$$
其中 $2m_1 \geq m_2$、$2m_2 \geq m_1 + m_3$、
$2m_3 \geq m_2 + m_4$、$2m_4 \geq m_3 + m_5 + m_7$、
$2m_5 \geq m_4 + m_6$、$2m_6 \geq m_5$，并且 $2m_7 \geq m_4$。
\end{enumerate}
\end{lemma}

\begin{proof}
见上文的讨论。
\end{proof}

\noindent
假设给定 $8$ 个互异的 $(-2)$-指标，相应非零元 $a_{ij}$ 在矩阵
$A$ 中的图形如下：
$$
\xymatrix{
\bullet \ar@{-}[r] & \bullet \ar@{-}[r] & \bullet \ar@{-}[r] &
\bullet \ar@{-}[r] & \bullet \ar@{-}[r] \ar@{-}[d] &
\bullet \ar@{-}[r] & \bullet \\
& & & & \bullet
}
$$
或如下：
$$
\xymatrix{
\bullet \ar@{-}[r] & \bullet \ar@{-}[r] & \bullet \ar@{-}[r] &
\bullet \ar@{-}[r] \ar@{-}[d] & \bullet \ar@{-}[r] &
\bullet \ar@{-}[r] & \bullet \\
& & & \bullet
}
$$
完全仿照引理 \ref{models-lemma-E7} 的证明，可知第一种图形得到情形
(\ref{models-item-E8})，它列于引理 \ref{models-lemma-E8}；而且不会在引理
\ref{models-lemma-genus-one} 中产生新情形。完全仿照引理
\ref{models-lemma-E6-completed} 的证明，可知若 $n > 8$，第二种图形不会出现；
但它会给出情形 (\ref{models-item-E7-completed})，该情形列于引理
\ref{models-lemma-genus-one}，此时 $n = 8$。

\begin{lemma}
\label{models-lemma-E8}
如下形状的真子图的分类：
$$
\xymatrix{
\bullet \ar@{-}[r] & \bullet \ar@{-}[r] & \bullet \ar@{-}[r] &
\bullet \ar@{-}[r] & \bullet \ar@{-}[r] \ar@{-}[d] &
\bullet \ar@{-}[r] & \bullet \\
& & & & \bullet
}
$$
设 $n > 8$。给定 $8$ 个互异的 $(-2)$-指标 $i_1, \ldots, i_8$，
使得 $a_{12}, a_{23}, a_{34}, a_{45}, a_{56}, a_{65}, a_{57}$ 非零，
则 $m$、$a$、$w$ 数据如下：
\begin{enumerate}
\item
\label{models-item-E8}
分别为
$$
\left(
\begin{matrix}
m_1 \\
m_2 \\
m_3 \\
m_4 \\
m_5 \\
m_6 \\
m_7 \\
m_8
\end{matrix}
\right),
\quad
\left(
\begin{matrix}
-2w & w & 0 & 0 & 0 & 0 & 0 & 0 \\
w & -2w & w & 0 & 0 & 0 & 0 & 0 \\
0 & w & -2w & w & 0 & 0 & 0 & 0 \\
0 & 0 & w & -2w & w & 0 & 0 & 0 \\
0 & 0 & 0 & w & -2w & w & 0 & w \\
0 & 0 & 0 & 0 & w & -2w & w & 0 \\
0 & 0 & 0 & 0 & 0 & w & -2w & 0 \\
0 & 0 & 0 & 0 & w & 0 & 0 & -2w
\end{matrix}
\right),
\quad
\left(
\begin{matrix}
w \\
w \\
w \\
w \\
w \\
w \\
w \\
w
\end{matrix}
\right)
$$
其中 $2m_1 \geq m_2$、$2m_2 \geq m_1 + m_3$、
$2m_3 \geq m_2 + m_4$、$2m_4 \geq m_3 + m_5$、
$2m_5 \geq m_4 + m_6 + m_8$、$2m_6 \geq m_5 + m_7$、
$2m_7 \geq m_6$，并且 $2m_8 \geq m_5$。
\end{enumerate}
\end{lemma}

\begin{proof}
见上文的讨论。
\end{proof}

\begin{lemma}
\label{models-lemma-E7-completed}
如下形状的真子图不存在：
$$
\xymatrix{
\bullet \ar@{-}[r] &
\bullet \ar@{-}[r] &
\bullet \ar@{-}[r] & \bullet \ar@{-}[r] \ar@{-}[d] &
\bullet \ar@{-}[r] & \bullet \ar@{-}[r] & \bullet \\
& & & \bullet
}
$$
假设 $n > 8$。{\bf 不}存在 $8$ 个互异的 $(-2)$-指标
$e, f, g, h, i, j, k, l$，使得
$a_{ef}, a_{fg}, a_{gh}, a_{hi}, a_{ij}, a_{jk}, a_{lh}$ 非零。
\end{lemma}

\begin{proof}
见上文的讨论。
\end{proof}

\noindent
假设给定 $9$ 个互异的 $(-2)$-指标，相应非零元 $a_{ij}$ 在矩阵
$A$ 中的图形如下：
$$
\xymatrix{
\bullet \ar@{-}[r] & \bullet \ar@{-}[r] &
\bullet \ar@{-}[r] & \bullet \ar@{-}[r] &
\bullet \ar@{-}[r] & \bullet \ar@{-}[r] \ar@{-}[d] &
\bullet \ar@{-}[r] & \bullet \\
& & & & & \bullet
}
$$
完全仿照引理 \ref{models-lemma-E6-completed} 的证明，可知若 $n > 9$，
这种图形不会出现；但当 $n = 9$ 时，它给出情形
(\ref{models-item-E8-completed})，该情形列于引理 \ref{models-lemma-genus-one}。

\begin{lemma}
\label{models-lemma-E8-completed}
如下形状的真子图不存在：
$$
\xymatrix{
\bullet \ar@{-}[r] & \bullet \ar@{-}[r] &
\bullet \ar@{-}[r] & \bullet \ar@{-}[r] &
\bullet \ar@{-}[r] & \bullet \ar@{-}[r] \ar@{-}[d] &
\bullet \ar@{-}[r] & \bullet \\
& & & & & \bullet
}
$$
假设 $n > 9$。{\bf 不}存在 $9$ 个互异的 $(-2)$-指标
$d, e, f, g, h, i, j, k, l$，使得
$a_{de}, a_{ef}, a_{fg}, a_{gh}, a_{hi}, a_{ij}, a_{jk}, a_{lh}$ 非零。
\end{lemma}

\begin{proof}
见上文的讨论。
\end{proof}

\noindent
汇总上述信息，得到如下命题。

\begin{proposition}
\label{models-proposition-classify-subgraphs}
令 $n, m_i, a_{ij}, w_i, g_i$ 为亏格为 $g$ 的数值型。
令 $I \subset \{1, \ldots, n\}$ 为基数 $\geq 2$ 的真子集，
它由 $(-2)$-指标组成，并且不存在非空真子集 $I' \subset I$，使得
$a_{i'i} = 0$ 对 $i' \in I$、$i \in I \setminus I'$ 成立。
则适当重排后，$m_i$、$a_{ij}$、$w_i$（$i, j \in I$）的数据
正是下列各引理所列者：引理 \ref{models-lemma-two-by-two}、
\ref{models-lemma-three-by-three},
\ref{models-lemma-four-by-four},
\ref{models-lemma-D4},
\ref{models-lemma-five-by-five},
\ref{models-lemma-D5},
\ref{models-lemma-long},
\ref{models-lemma-Dn},
\ref{models-lemma-E6},
\ref{models-lemma-E7} 或
\ref{models-lemma-E8}。
\end{proposition}

\begin{proof}
这由上文的讨论立即得到；另见第 \ref{models-section-classify-proper-subgraphs} 节
开头的讨论。
\end{proof}













\section{亏格零与一的极小型分类}
\label{models-section-classification-genus-one}

\noindent
本节标题已经说明全部内容。

\begin{lemma}[亏格零]
\label{models-lemma-genus-zero}
亏格为零的极小数值型只有
$n = 1$, $m_1 = 1$, $a_{11} = 0$, $w_1 = 1$, $g_1 = 0$.
\end{lemma}

\begin{proof}
这由引理 \ref{models-lemma-non-irreducible-minimal-type-genus-at-least-one}
和 \ref{models-lemma-irreducible} 得出。
\end{proof}

\begin{lemma}[亏格一]
\label{models-lemma-genus-one}
在等价意义下，亏格为一的极小数值型如下：
\begin{enumerate}
\item
\label{models-item-one}
$n = 1$、$a_{11} = 0$、$g_1 = 1$，而 $m_1, w_1 \geq 1$ 任意；
\item
\label{models-item-two-cycle}
$n = 2$，且 $m_i, a_{ij}, w_i, g_i$ 如下：
$$
\left(
\begin{matrix}
m \\
m
\end{matrix}
\right),
\quad
\left(
\begin{matrix}
-2w & 2w \\
2w & -2w
\end{matrix}
\right),
\quad
\left(
\begin{matrix}
w \\
w
\end{matrix}
\right),
\quad
\left(
\begin{matrix}
0 \\
0
\end{matrix}
\right)
$$
其中 $w$ 和 $m$ 任意；
\item
\label{models-item-up4}
$n = 2$，且 $m_i, a_{ij}, w_i, g_i$ 如下：
$$
\left(
\begin{matrix}
2m \\
m
\end{matrix}
\right),
\quad
\left(
\begin{matrix}
-2w & 4w \\
4w & -8w
\end{matrix}
\right),
\quad
\left(
\begin{matrix}
w \\
4w
\end{matrix}
\right),
\quad
\left(
\begin{matrix}
0 \\
0
\end{matrix}
\right)
$$
其中 $w$ 和 $m$ 任意；
\item
\label{models-item-three-cycle}
$n = 3$，且 $m_i, a_{ij}, w_i, g_i$ 如下：
$$
\left(
\begin{matrix}
m \\
m \\
m
\end{matrix}
\right),
\quad
\left(
\begin{matrix}
-2w & w & w \\
w & -2w & w \\
w & w & -2w
\end{matrix}
\right),
\quad
\left(
\begin{matrix}
w \\
w \\
w
\end{matrix}
\right),
\quad
\left(
\begin{matrix}
0 \\
0 \\
0
\end{matrix}
\right)
$$
其中 $w$ 和 $m$ 任意；
\item
\label{models-item-equal-up3}
$n = 3$，且 $m_i, a_{ij}, w_i, g_i$ 如下：
$$
\left(
\begin{matrix}
m \\
2m \\
m
\end{matrix}
\right),
\quad
\left(
\begin{matrix}
-2w & w & 0 \\
w & -2w & 3w \\
0 & 3w & -6w
\end{matrix}
\right),
\quad
\left(
\begin{matrix}
w \\
w \\
3w
\end{matrix}
\right),
\quad
\left(
\begin{matrix}
0 \\
0 \\
0
\end{matrix}
\right)
$$
其中 $w$ 和 $m$ 任意；
\item
\label{models-item-equal-down3}
$n = 3$，且 $m_i, a_{ij}, w_i, g_i$ 如下：
$$
\left(
\begin{matrix}
m \\
2m \\
3m
\end{matrix}
\right),
\quad
\left(
\begin{matrix}
-6w & 3w & 0 \\
3w & -6w & 3w \\
0 & 3w & -2w
\end{matrix}
\right),
\quad
\left(
\begin{matrix}
3w \\
3w \\
w
\end{matrix}
\right),
\quad
\left(
\begin{matrix}
0 \\
0 \\
0
\end{matrix}
\right)
$$
其中 $w$ 和 $m$ 任意；
\item
\label{models-item-up-up}
$n = 3$，且 $m_i, a_{ij}, w_i, g_i$ 如下：
$$
\left(
\begin{matrix}
2m \\
2m \\
m
\end{matrix}
\right),
\quad
\left(
\begin{matrix}
-2w & 2w & 0 \\
2w & -4w & 4w \\
0 & 4w & -8w
\end{matrix}
\right),
\quad
\left(
\begin{matrix}
w \\
2w \\
4w
\end{matrix}
\right),
\quad
\left(
\begin{matrix}
0 \\
0 \\
0
\end{matrix}
\right)
$$
其中 $w$ 和 $m$ 任意；
\item
\label{models-item-up-down}
$n = 3$，且 $m_i, a_{ij}, w_i, g_i$ 如下：
$$
\left(
\begin{matrix}
m \\
m \\
m
\end{matrix}
\right),
\quad
\left(
\begin{matrix}
-2w & 2w & 0 \\
2w & -4w & 2w \\
0 & 2w & -2w
\end{matrix}
\right),
\quad
\left(
\begin{matrix}
w \\
2w \\
w
\end{matrix}
\right),
\quad
\left(
\begin{matrix}
0 \\
0 \\
0
\end{matrix}
\right)
$$
其中 $w$ 和 $m$ 任意；
\item
\label{models-item-down-up}
$n = 3$，且 $m_i, a_{ij}, w_i, g_i$ 如下：
$$
\left(
\begin{matrix}
m \\
2m \\
m
\end{matrix}
\right),
\quad
\left(
\begin{matrix}
-4w & 2w & 0 \\
2w & -2w & 2w \\
0 & 2w & -4w
\end{matrix}
\right),
\quad
\left(
\begin{matrix}
2w \\
w \\
2w
\end{matrix}
\right),
\quad
\left(
\begin{matrix}
0 \\
0 \\
0
\end{matrix}
\right)
$$
其中 $w$ 和 $m$ 任意；
\item
\label{models-item-four-cycle}
$n = 4$，且 $m_i, a_{ij}, w_i, g_i$ 如下：
$$
\left(
\begin{matrix}
m \\
m \\
m \\
m
\end{matrix}
\right),
\quad
\left(
\begin{matrix}
-2w & w & 0 & w \\
w & -2w & w & 0 \\
0 & w & -2w & w \\
w & 0 & w & -2w
\end{matrix}
\right),
\quad
\left(
\begin{matrix}
w \\
w \\
w \\
w
\end{matrix}
\right),
\quad
\left(
\begin{matrix}
0 \\
0 \\
0 \\
0
\end{matrix}
\right)
$$
其中 $w$ 和 $m$ 任意；
\item
\label{models-item-up-equal-up}
$n = 4$，且 $m_i, a_{ij}, w_i, g_i$ 如下：
$$
\left(
\begin{matrix}
2m \\
2m \\
2m \\
m
\end{matrix}
\right),
\quad
\left(
\begin{matrix}
-2w & 2w & 0 & 0 \\
2w & -4w & 2w & 0 \\
0 & 2w & -4w & 4w \\
0 & 0 & 4w & -8w
\end{matrix}
\right),
\quad
\left(
\begin{matrix}
w \\
2w \\
2w \\
4w
\end{matrix}
\right),
\quad
\left(
\begin{matrix}
0 \\
0 \\
0 \\
0
\end{matrix}
\right)
$$
其中 $w$ 和 $m$ 任意；
\item
\label{models-item-up-equal-down}
$n = 4$，且 $m_i, a_{ij}, w_i, g_i$ 如下：
$$
\left(
\begin{matrix}
m \\
m \\
m \\
m
\end{matrix}
\right),
\quad
\left(
\begin{matrix}
-2w & 2w & 0 & 0 \\
2w & -4w & 2w & 0 \\
0 & 2w & -4w & 2w \\
0 & 0 & 2w & -2w
\end{matrix}
\right),
\quad
\left(
\begin{matrix}
w \\
2w \\
2w \\
w
\end{matrix}
\right),
\quad
\left(
\begin{matrix}
0 \\
0 \\
0 \\
0
\end{matrix}
\right)
$$
其中 $w$ 和 $m$ 任意；
\item
\label{models-item-down-equal-up}
$n = 4$，且 $m_i, a_{ij}, w_i, g_i$ 如下：
$$
\left(
\begin{matrix}
m \\
2m \\
2m \\
m
\end{matrix}
\right),
\quad
\left(
\begin{matrix}
-4w & 2w & 0 & 0 \\
2w & -2w & w & 0 \\
0 & w & -2w & 2w \\
0 & 0 & 2w & -4w
\end{matrix}
\right),
\quad
\left(
\begin{matrix}
2w \\
w \\
w \\
2w
\end{matrix}
\right),
\quad
\left(
\begin{matrix}
0 \\
0 \\
0 \\
0
\end{matrix}
\right)
$$
其中 $w$ 和 $m$ 任意；
\item
\label{models-item-triple-with-up}
$n = 4$，且 $m_i, a_{ij}, w_i, g_i$ 如下：
$$
\left(
\begin{matrix}
2m \\
m \\
m \\
m
\end{matrix}
\right),
\quad
\left(
\begin{matrix}
-2w & w & w & 2w \\
w & -2w & 0 & 0 \\
w & 0 & -2w & 0 \\
2w & 0 & 0 & -4w
\end{matrix}
\right),
\quad
\left(
\begin{matrix}
w \\
w \\
w \\
2w
\end{matrix}
\right),
\quad
\left(
\begin{matrix}
0 \\
0 \\
0 \\
0
\end{matrix}
\right)
$$
其中 $w$ 和 $m$ 任意；
\item
\label{models-item-triple-with-down}
$n = 4$，且 $m_i, a_{ij}, w_i, g_i$ 如下：
$$
\left(
\begin{matrix}
2m \\
m \\
m \\
2m
\end{matrix}
\right),
\quad
\left(
\begin{matrix}
-4w & 2w & 2w & 2w \\
2w & -4w & 0 & 0 \\
2w & 0 & -4w & 0 \\
2w & 0 & 0 & -2w
\end{matrix}
\right),
\quad
\left(
\begin{matrix}
2w \\
2w \\
2w \\
w
\end{matrix}
\right),
\quad
\left(
\begin{matrix}
0 \\
0 \\
0 \\
0
\end{matrix}
\right)
$$
其中 $w$ 和 $m$ 任意；
\item
\label{models-item-five-cycle}
$n = 5$，且 $m_i, a_{ij}, w_i, g_i$ 如下：
$$
\left(
\begin{matrix}
m \\
m \\
m \\
m \\
m
\end{matrix}
\right),
\quad
\left(
\begin{matrix}
-2w & w & 0 & 0 & w \\
w & -2w & w & 0 & 0 \\
0 & w & -2w & w & 0 \\
0 & 0 & w & -2w & w \\
w & 0 & 0 & w & -2w \\
\end{matrix}
\right),
\quad
\left(
\begin{matrix}
w \\
w \\
w \\
w \\
w
\end{matrix}
\right),
\quad
\left(
\begin{matrix}
0 \\
0 \\
0 \\
0 \\
0
\end{matrix}
\right)
$$
其中 $w$ 和 $m$ 任意；
\item
\label{models-item-equal-equal-up-equal}
$n = 5$，且 $m_i, a_{ij}, w_i, g_i$ 如下：
$$
\left(
\begin{matrix}
m \\
2m \\
3m \\
2m \\
m
\end{matrix}
\right),
\quad
\left(
\begin{matrix}
-2w & w & 0 & 0 & 0 \\
w & -2w & w & 0 & 0 \\
0 & w & -2w & 2w & 0 \\
0 & 0 & 2w & -4w & 2w \\
0 & 0 & 0 & 2w & -4w \\
\end{matrix}
\right),
\quad
\left(
\begin{matrix}
w \\
w \\
w \\
2w \\
2w
\end{matrix}
\right),
\quad
\left(
\begin{matrix}
0 \\
0 \\
0 \\
0 \\
0
\end{matrix}
\right)
$$
其中 $w$ 和 $m$ 任意；
\item
\label{models-item-equal-equal-down-equal}
$n = 5$，且 $m_i, a_{ij}, w_i, g_i$ 如下：
$$
\left(
\begin{matrix}
m \\
2m \\
3m \\
4m \\
2m
\end{matrix}
\right),
\quad
\left(
\begin{matrix}
-4w & 2w & 0 & 0 & 0 \\
2w & -4w & 2w & 0 & 0 \\
0 & 2w & -4w & 2w & 0 \\
0 & 0 & 2w & -2w & w \\
0 & 0 & 0 & w & -2w \\
\end{matrix}
\right),
\quad
\left(
\begin{matrix}
2w \\
2w \\
2w \\
w \\
w
\end{matrix}
\right),
\quad
\left(
\begin{matrix}
0 \\
0 \\
0 \\
0 \\
0
\end{matrix}
\right)
$$
其中 $w$ 和 $m$ 任意；
\item
\label{models-item-up-equal-equal-up}
$n = 5$，且 $m_i, a_{ij}, w_i, g_i$ 如下：
$$
\left(
\begin{matrix}
2m \\
2m \\
2m \\
2m \\
m
\end{matrix}
\right),
\quad
\left(
\begin{matrix}
-2w & 2w & 0 & 0 & 0 \\
2w & -4w & 2w & 0 & 0 \\
0 & 2w & -4w & 2w & 0 \\
0 & 0 & 2w & -4w & 4w \\
0 & 0 & 0 & 4w & -8w \\
\end{matrix}
\right),
\quad
\left(
\begin{matrix}
w \\
2w \\
2w \\
2w \\
4w
\end{matrix}
\right),
\quad
\left(
\begin{matrix}
0 \\
0 \\
0 \\
0 \\
0
\end{matrix}
\right)
$$
其中 $w$ 和 $m$ 任意；
\item
\label{models-item-up-equal-equal-down}
$n = 5$，且 $m_i, a_{ij}, w_i, g_i$ 如下：
$$
\left(
\begin{matrix}
m \\
m \\
m \\
m \\
m
\end{matrix}
\right),
\quad
\left(
\begin{matrix}
-2w & 2w & 0 & 0 & 0 \\
2w & -4w & 2w & 0 & 0 \\
0 & 2w & -4w & 2w & 0 \\
0 & 0 & 2w & -4w & 2w \\
0 & 0 & 0 & 2w & -2w \\
\end{matrix}
\right),
\quad
\left(
\begin{matrix}
w \\
2w \\
2w \\
2w \\
w
\end{matrix}
\right),
\quad
\left(
\begin{matrix}
0 \\
0 \\
0 \\
0 \\
0
\end{matrix}
\right)
$$
其中 $w$ 和 $m$ 任意；
\item
\label{models-item-down-equal-equal-up}
$n = 5$，且 $m_i, a_{ij}, w_i, g_i$ 如下：
$$
\left(
\begin{matrix}
m \\
2m \\
2m \\
2m \\
m
\end{matrix}
\right),
\quad
\left(
\begin{matrix}
-4w & 2w & 0 & 0 & 0 \\
2w & -2w & w & 0 & 0 \\
0 & w & -2w & w & 0 \\
0 & 0 & w & -2w & 2w \\
0 & 0 & 0 & 2w & -4w \\
\end{matrix}
\right),
\quad
\left(
\begin{matrix}
2w \\
w \\
w \\
w \\
2w
\end{matrix}
\right),
\quad
\left(
\begin{matrix}
0 \\
0 \\
0 \\
0 \\
0
\end{matrix}
\right)
$$
其中 $w$ 和 $m$ 任意；
\item
\label{models-item-quadruple}
$n = 5$，且 $m_i, a_{ij}, w_i, g_i$ 如下：
$$
\left(
\begin{matrix}
2m \\
m \\
m \\
m \\
m
\end{matrix}
\right),
\quad
\left(
\begin{matrix}
-2w & w & w & w & w \\
w & -2w & 0 & 0 & 0 \\
w & 0 & -2w & 0 & 0 \\
w & 0 & 0 & -2w & 0 \\
w & 0 & 0 & 0 & -2w \\
\end{matrix}
\right),
\quad
\left(
\begin{matrix}
w \\
w \\
w \\
w \\
w
\end{matrix}
\right),
\quad
\left(
\begin{matrix}
0 \\
0 \\
0 \\
0 \\
0
\end{matrix}
\right)
$$
其中 $w$ 和 $m$ 任意；
\item
\label{models-item-triple-extended-up}
$n = 5$，且 $m_i, a_{ij}, w_i, g_i$ 如下：
$$
\left(
\begin{matrix}
m \\
2m \\
2m \\
m \\
m
\end{matrix}
\right),
\quad
\left(
\begin{matrix}
-4w & 2w & 0 & 0 & 0 \\
2w & -2w & w & 0 & 0 \\
0 & w & -2w & w & w \\
0 & 0 & w & -2w & 0 \\
0 & 0 & w & 0 & -2w \\
\end{matrix}
\right),
\quad
\left(
\begin{matrix}
2w \\
w \\
w \\
w \\
w
\end{matrix}
\right),
\quad
\left(
\begin{matrix}
0 \\
0 \\
0 \\
0 \\
0
\end{matrix}
\right)
$$
其中 $w$ 和 $m$ 任意；
\item
\label{models-item-triple-extended-down}
$n = 5$，且 $m_i, a_{ij}, w_i, g_i$ 如下：
$$
\left(
\begin{matrix}
2m \\
2m \\
2m \\
m \\
m
\end{matrix}
\right),
\quad
\left(
\begin{matrix}
-2w & 2w & 0 & 0 & 0 \\
2w & -4w & 2w & 0 & 0 \\
0 & 2w & -4w & 2w & 2w \\
0 & 0 & 2w & -4w & 0 \\
0 & 0 & 2w & 0 & -4w \\
\end{matrix}
\right),
\quad
\left(
\begin{matrix}
w \\
2w \\
2w \\
2w \\
2w
\end{matrix}
\right),
\quad
\left(
\begin{matrix}
0 \\
0 \\
0 \\
0 \\
0
\end{matrix}
\right)
$$
其中 $w$ 和 $m$ 任意；
\item
\label{models-item-n-cycle}
$n \geq 6$，并有一个推广情形 (\ref{models-item-five-cycle}) 的 $n$-圈：
\begin{enumerate}
\item $m_1 = \ldots = m_n = m$,
\item $a_{12} = \ldots = a_{(n - 1) n} = w$, $a_{1n} = w$,
而对其余 $i < j$，有 $a_{ij} = 0$；
\item $w_1 = \ldots = w_n = w$
\end{enumerate}
其中 $w$ 和 $m$ 任意；
\item
\label{models-item-up-chain-equal-up}
$n \geq 6$，并有一条推广情形 (\ref{models-item-up-equal-equal-up}) 的链：
\begin{enumerate}
\item $m_1 = \ldots = m_{n - 1} = 2m$, $m_n = m$,
\item $a_{12} = \ldots = a_{(n - 2) (n - 1)} = 2w$, $a_{(n - 1) n} = 4w$,
而对其余 $i < j$，有 $a_{ij} = 0$；
\item $w_1 = w$, $w_2 = \ldots = w_{n - 1} = 2w$, $w_n = 4w$
\end{enumerate}
其中 $w$ 和 $m$ 任意；
\item
\label{models-item-up-chain-equal-down}
$n \geq 6$，并有一条推广情形 (\ref{models-item-up-equal-equal-down}) 的链：
\begin{enumerate}
\item $m_1 = \ldots = m_n = m$,
\item $a_{12} = \ldots = a_{(n - 1) n} = w$,
而对其余 $i < j$，有 $a_{ij} = 0$；
\item $w_1 = w$, $w_2 = \ldots = w_{n - 1} = 2w$, $w_n = w$
\end{enumerate}
其中 $w$ 和 $m$ 任意；
\item
\label{models-item-down-chain-equal-up}
$n \geq 6$，并有一条推广情形 (\ref{models-item-down-equal-equal-up}) 的链：
\begin{enumerate}
\item $m_1 = w$, $w_2 = \ldots = m_{n - 1} = 2m$, $m_n = m$,
\item $a_{12} = 2w$, $a_{23} = \ldots = a_{(n - 2) (n - 1)} = w$,
$a_{(n - 1) n} = 2w$，而对其余 $i < j$，有 $a_{ij} = 0$；
\item $w_1 = 2w$, $w_2 = \ldots = w_{n - 1} = w$, $w_n = 2w$
\end{enumerate}
其中 $w$ 和 $m$ 任意；
\item
\label{models-item-Dn-extended-up}
$n \geq 6$，并有一个推广情形 (\ref{models-item-triple-extended-up}) 的型：
\begin{enumerate}
\item $m_1 = m$, $m_2 = \ldots = m_{n - 3} = 2m$, $m_{n - 1} = m_n = m$,
\item $a_{12} = 2w$, $a_{23} = \ldots = a_{(n - 2) (n - 1)} = w$,
$a_{(n - 2) n} = w$，而对其余 $i < j$，有 $a_{ij} = 0$；
\item $w_1 = 2w$, $w_2 = \ldots = w_n = w$
\end{enumerate}
其中 $w$ 和 $m$ 任意；
\item
\label{models-item-Dn-extended-down}
$n \geq 6$，并有一个推广情形 (\ref{models-item-triple-extended-down}) 的型：
\begin{enumerate}
\item $m_1 = \ldots = m_{n - 3} = 2m$, $m_{n - 1} = m_n = m$,
\item $a_{12} = \ldots = a_{(n - 2) (n - 1)} = 2w$,
$a_{(n - 2) n} = 2w$，而对其余 $i < j$，有 $a_{ij} = 0$；
\item $w_1 = w$, $w_2 = \ldots = w_n = 2w$
\end{enumerate}
其中 $w$ 和 $m$ 任意；
\item
\label{models-item-double-triple}
$n \geq 6$，并有一个推广情形 (\ref{models-item-quadruple}) 的型：
\begin{enumerate}
\item $m_1 = m_2 = m$, $m_3 = \ldots = m_{n - 2} = 2m$, $m_{n - 1} = m_n = m$,
\item $a_{13} = w$, $a_{23} = \ldots = a_{(n - 2) (n - 1)} = w$,
$a_{(n - 2) n} = w$，而对其余 $i < j$，有 $a_{ij} = 0$；
\item $w_1 = \ldots = w_n = w$,
\end{enumerate}
其中 $w$ 和 $m$ 任意；
\item
\label{models-item-E6-completed}
$n = 7$，且 $m_i, a_{ij}, w_i, g_i$ 如下：
$$
\left(
\begin{matrix}
m \\
2m \\
3m \\
m \\
2m \\
m \\
2m
\end{matrix}
\right),
\quad
\left(
\begin{matrix}
-2w & w & 0 & 0 & 0 & 0 & 0 \\
w & -2w & w & 0 & 0 & 0 & 0 \\
0 & w & -2w & 0 & w & 0 & w \\
0 & 0 & 0 & -2w & w & 0 & 0 \\
0 & 0 & w & w & -2w & 0 & 0 \\
0 & 0 & 0 & 0 & 0 & -2w & w \\
0 & 0 & w & 0 & 0 & w & -2w
\end{matrix}
\right),
\quad
\left(
\begin{matrix}
w \\
w \\
w \\
w \\
w \\
w \\
w
\end{matrix}
\right),
\quad
\left(
\begin{matrix}
0 \\
0 \\
0 \\
0 \\
0 \\
0 \\
0
\end{matrix}
\right)
$$
其中 $w$ 和 $m$ 任意；
\item
\label{models-item-E7-completed}
$n = 8$，且 $m_i, a_{ij}, w_i, g_i$ 如下：
$$
\left(
\begin{matrix}
m \\
2m \\
3m \\
4m \\
3m \\
2m \\
m \\
2m
\end{matrix}
\right),
\quad
\left(
\begin{matrix}
-2w & w & 0 & 0 & 0 & 0 & 0 & 0 \\
w & -2w & w & 0 & 0 & 0 & 0 & 0 \\
0 & w & -2w & w & 0 & 0 & 0 & 0 \\
0 & 0 & w & -2w & w & 0 & 0 & w \\
0 & 0 & 0 & w & -2w & w & 0 & 0 \\
0 & 0 & 0 & 0 & w & -2w & w & 0 \\
0 & 0 & 0 & 0 & 0 & w & -2w & 0 \\
0 & 0 & 0 & w & 0 & 0 & 0 & -2w \\
\end{matrix}
\right),
\quad
\left(
\begin{matrix}
w \\
w \\
w \\
w \\
w \\
w \\
w \\
w
\end{matrix}
\right),
\quad
\left(
\begin{matrix}
0 \\
0 \\
0 \\
0 \\
0 \\
0 \\
0 \\
0
\end{matrix}
\right)
$$
其中 $w$ 和 $m$ 任意；
\item
\label{models-item-E8-completed}
$n = 9$，且 $m_i, a_{ij}, w_i, g_i$ 如下：
$$
\left(
\begin{matrix}
m \\
2m \\
3m \\
4m \\
5m \\
6m \\
4m \\
2m \\
3m
\end{matrix}
\right),
\quad
\left(
\begin{matrix}
-2w & w & 0 & 0 & 0 & 0 & 0 & 0 & 0 \\
w & -2w & w & 0 & 0 & 0 & 0 & 0 & 0 \\
0 & w & -2w & w & 0 & 0 & 0 & 0 & 0 \\
0 & 0 & w & -2w & w & 0 & 0 & 0 & 0 \\
0 & 0 & 0 & w & -2w & w & 0 & 0 & 0 \\
0 & 0 & 0 & 0 & w & -2w & w & 0 & w \\
0 & 0 & 0 & 0 & 0 & w & -2w & w & 0 \\
0 & 0 & 0 & 0 & 0 & 0 & w & -2w & 0 \\
0 & 0 & 0 & 0 & 0 & w & 0 & 0 & -2w \\
\end{matrix}
\right),
\quad
\left(
\begin{matrix}
w \\
w \\
w \\
w \\
w \\
w \\
w \\
w \\
w
\end{matrix}
\right),
\quad
\left(
\begin{matrix}
0 \\
0 \\
0 \\
0 \\
0 \\
0 \\
0 \\
0 \\
0
\end{matrix}
\right)
$$
其中 $w$ 和 $m$ 任意。
\end{enumerate}
\end{lemma}

\begin{proof}
这已经在第 \ref{models-section-classify-proper-subgraphs} 节中证明。
另见第 \ref{models-section-classify-proper-subgraphs} 节开头的讨论。
\end{proof}






\section{数值型不变量的界}
\label{models-section-towards-classification}

\noindent
在证明曲线的半稳定约化时，我们将用到亏格为 $g$ 的数值型之
Picard 群的一个界；本节将证明这个界。

\begin{lemma}
\label{models-lemma-bound-neighbours}
令 $n, m_i, a_{ij}, w_i, g_i$ 为亏格为 $g$ 的数值型。
给定 $i, j$，并假设 $a_{ij} > 0$，则有
$m_ia_{ij} \leq m_j|a_{jj}|$ 且 $m_iw_i \leq m_j|a_{jj}|$。
\end{lemma}

\begin{proof}
对每个指标 $j$，都有
$m_j a_{jj} + \sum_{i \not = j} m_ia_{ij} = 0$。因此，若已知
$|a_{jj}|$ 和 $m_j$ 的上界，则也得到非零（因而为正）的 $a_{ij}$
以及 $m_i$ 的上界。回忆 $w_i$ 整除 $a_{ij}$，读者很容易看出本引理成立。
\end{proof}

\begin{lemma}
\label{models-lemma-bound-heart}
固定 $g \geq 2$。对每个极小数值型 $n, m_i, a_{ij}, w_i, g_i$，
若其亏格为 $g$ 且 $n > 1$，则有：
\begin{enumerate}
\item 集合 $J \subset \{1, \ldots, n\}$ 由非 $(-2)$-指标构成，且
至多有 $2g - 2$ 个元素；
\item 对 $j \in J$，有 $g_j < g$；
\item 对 $j \in J$，有 $m_j|a_{jj}| \leq 6g - 6$；
\item 对 $j \in J$ 以及 $i \in \{1, \ldots, n\}$，有
$m_ia_{ij} \leq 6g - 6$。
\end{enumerate}
\end{lemma}

\begin{proof}
回忆 $g = 1 + \sum m_j(w_j(g_j - 1) - \frac{1}{2} a_{jj})$。
对 $j \in J$，量 $m_j(w_j(g_j - 1) - \frac{1}{2} a_{jj})$
对亏格 $g$ 的贡献为 $> 0$，因而
$\geq 1/2$。这里使用了引理 \ref{models-lemma-minus-one}、
定义 \ref{models-definition-type-minus-one}、定义 \ref{models-definition-type-minimal}、
引理 \ref{models-lemma-minus-two} 和定义 \ref{models-definition-type-minus-two}；
下文将不再特别说明这些结果。因此 $J$ 至多有 $2(g - 1)$ 个元素。
这证明了 (1)。

\medskip\noindent
回忆由引理 \ref{models-lemma-diagonal-negative}，$-a_{ii} > 0$ 对所有 $i$
成立。因此，对 $j \in J$，量
$m_j(w_j(g_j - 1) - \frac{1}{2} a_{jj})$ 对亏格 $g$ 的贡献为
$> m_jw_j(g_j - 1)$。故
$$
g - 1 > m_jw_j(g_j - 1) \Rightarrow g_j < (g - 1)/m_jw_j + 1
$$
这确实推出 $g_j < g$，从而证明了 (2)。

\medskip\noindent
对 $j \in J$，若 $g_j > 0$，则量
$m_j(w_j(g_j - 1) - \frac{1}{2} a_{jj})$ 对亏格 $g$ 的贡献为
$\geq -\frac{1}{2}m_ja_{jj}$，于是立即得到
$m_j|a_{jj}| \leq 2(g - 1)$。否则 $a_{jj} = -kw_j$，其中
$k \geq 3$ 为整数（因为 $j \in J$），并且有
$$
m_jw_j(-1 + \frac{k}{2}) \leq g - 1
\Rightarrow
m_jw_j \leq \frac{2(g - 1)}{k - 2}
$$
将它代回 $a_{jj} = -km_jw_j$，得到
$$
m_j|a_{jj}| \leq 2(g - 1) \frac{k}{k - 2} \leq 6(g - 1)
$$
这证明了 (3)。

\medskip\noindent
(4) 由引理 \ref{models-lemma-bound-neighbours} 和 (3) 得出。
\end{proof}

\begin{lemma}
\label{models-lemma-bound-wm}
固定 $g \geq 2$。对每个极小数值型 $n, m_i, a_{ij}, w_i, g_i$，
若其亏格为 $g$，则有 $m_i|a_{ij}| \leq 768g$。
\end{lemma}

\begin{proof}
由引理 \ref{models-lemma-bound-neighbours}，只需证明
$m_i|a_{ii}| \leq 768g$ 对所有 $i$ 成立。令
$J \subset \{1, \ldots, n\}$ 为引理
\ref{models-lemma-bound-heart} 中的非 $(-2)$-指标集。注意 $J$ 非空，因为
$g \geq 2$。此外，由该引理，有 $m_j|a_{jj}| \leq 6g$，其中
$j \in J$。

\medskip\noindent
设有 $j \in J$ 以及一个序列 $i_1, \ldots, i_7$，它由
$(-2)$-指标组成，并使得 $a_{ji_1}$ 以及 $a_{i_1i_2}$、
$a_{i_2i_3}$、$a_{i_3i_4}$、$a_{i_4i_5}$、$a_{i_5i_6}$、
$a_{i_6i_7}$ 非零。由引理 \ref{models-lemma-bound-neighbours} 可知
$m_{i_1}w_{i_1} \leq 6g$ 且 $m_{i_1}a_{ji_1} \leq 6g$。
因为 $i_1$ 是 $(-2)$-指标，所以 $a_{i_1i_1} = -2w_{i_1}$，
从而 $m_{i_1}|a_{i_1i_1}| \leq 12g$。重复这一论证，得到
$m_{i_2}w_{i_2} \leq 12g$ 且 $m_{i_2}a_{i_1i_2} \leq 12g$。
继而 $m_{i_2}|a_{i_2i_2}| \leq 24g$，依此类推。最终得到
$m_{i_k}|a_{i_ki_k}| \leq 2^k(6g) \leq 768g$，其中
$k = 1, \ldots, 7$。

\medskip\noindent
令 $I \subset \{1, \ldots, n\} \setminus J$ 为极大连通子集。
换言之，不存在非空真子集 $I' \subset I$，使得
$a_{i'i} = 0$（$i' \in I'$ 且 $i \in I \setminus I'$）；并且
$I$ 在此性质下极大。
特别地，由于数值型按定义是连通的，存在 $j \in J$ 和 $i \in I$，
使得 $a_{ij} > 0$。查看命题 \ref{models-proposition-classify-subgraphs} 中
对此类 $I$ 的分类，并使用上一段的结果，可知
$w_i|a_{ii}| \leq 768g$ 对所有 $i \in I$ 成立，除非 $I$ 是引理
\ref{models-lemma-long} 或引理
\ref{models-lemma-Dn} 所述的情形。因此可以假设 $a_{ii'}$ 的非零图形
（$i, i' \in I$）为
$$
\xymatrix{
\bullet \ar@{-}[r] &
\bullet \ar@{-}[r] &
\bullet \ar@{..}[r] &
\bullet \ar@{-}[r] &
\bullet \ar@{-}[r] &
\bullet
}
$$
（它有引理 \ref{models-lemma-long} 详述的 3 个子情形），或为
$$
\xymatrix{
\bullet \ar@{-}[r] &
\bullet \ar@{-}[r] &
\bullet \ar@{..}[r] &
\bullet \ar@{-}[r] &
\bullet \ar@{-}[r] \ar@{-}[d] &
\bullet \\
& & & & \bullet
}
$$
我们将证明该界对引理 \ref{models-lemma-long} 的第一个子情形成立；其余情形
留给读者，其论证几乎完全相同。

\medskip\noindent
重新编号后，可以假设 $I = \{1, \ldots, t\} \subset \{1, \ldots, n\}$，
并且存在整数 $w$，使得
$$
w = w_1 = \ldots = w_t =
a_{12} = \ldots = a_{(t - 1) t} =
-\frac{1}{2} a_{i_1i_2} = \ldots = -\frac{1}{2} a_{(t - 1) t}
$$
等式 $a_{ii}m_i + \sum_{j \not = i} a_{ij}m_j = 0$ 推出
$$
2m_2 \geq m_1 + m_3, \ldots, 2m_{t - 1} \geq m_{t - 2} + m_t
$$
不等式 $2m_i \geq m_{i - 1} + m_{i + 1}$ 中等号成立，当且仅当
$i$ 除 $i - 1$ 和 $i + 1$ 外不与任何指标“相交”。若 $i$ 确实与另一个
指标相交，则该指标属于 $J$（由 $I$ 的极大性）。特别地，映射
$\{1, \ldots, t\} \to \mathbf{Z}$，$i \mapsto m_i$ 是凹的。

\medskip\noindent
令 $m = \max(m_i, i \in \{1, \ldots, t\})$。则
$m_i|a_{ii}| \leq 2mw$ 对 $i \leq t$ 成立，我们的目标是证明
$2mw \leq 768g$。
令 $s$、相应地 $s'$ 为 $\{1, \ldots, t\}$ 中满足
$m_s = m = m_{s'}$ 的最小、相应地最大指标。由凹性可知
$m_i = m$ 对 $s \leq i \leq s'$ 成立。若 $s > 1$，则不等式
$2m_s \geq m_{s - 1} + m_{s + 1}$ 中等号不成立，因而 $s$ 与
$J$ 中的一个指标相交。此时由证明第二段的结果，$2mw \leq 12g$。
同样，若 $s' < t$，则 $s'$ 与 $J$ 中的一个指标相交，仍得到
$2mw \leq 12g$。但若 $s = 1$ 且 $s' = t$，则有 $a_{ij} = 0$，
其中 $j \in J$ 且 $i \in \{2, \ldots, t - 1\}$。然而我们已经看到，
必存在一对 $(i, j) \in I \times J$，使得 $a_{ij} > 0$；于是只能有
$i = 1$ 或 $i = t$。仍按前述方法得到 $2mw \leq 12g$（因为此时
$m_1 = m = m_t$）。
\end{proof}

\begin{proposition}
\label{models-proposition-bound-picard-group}
令 $g \geq 2$。对每个数值型 $T$（其亏格为 $g$）以及素数
$\ell > 768g$，有
$$
\dim_{\mathbf{F}_\ell} \Pic(T)[\ell] \leq g
$$
其中 $\Pic(T)$ 如定义 \ref{models-definition-picard-group} 所述。若 $T$ 极小，
则甚至有
$$
\dim_{\mathbf{F}_\ell} \Pic(T)[\ell] \leq g_{top} \leq g
$$
其中 $g_{top}$ 如定义 \ref{models-definition-top-genus} 所述。
\end{proposition}

\begin{proof}
设 $T$ 由 $n, m_i, a_{ij}, w_i, g_i$ 给出。若 $T$ 不极小，
则存在一个 $(-1)$-指标。用等价型替换 $T$ 后，可假设 $n$ 是
$(-1)$-指标。应用引理 \ref{models-lemma-contract-picard-group}，得到
$\Pic(T) \subset \Pic(T')$，其中 $T'$ 是亏格为 $g$ 的数值型
（引理 \ref{models-lemma-contract}），并且有 $n - 1$ 个指标。因此，只要对
极小数值型证明本引理，即可对 $n$ 作归纳而得到结论。

\medskip\noindent
假设 $T$ 是亏格 $\geq 2$ 的极小数值型。由引理
\ref{models-lemma-genus-nonnegative} 可知 $g_{top} \leq g$。若 $A = (a_{ij})$，
则由引理 \ref{models-lemma-picard-T-and-A} 有
$\Pic(T) \subset \Coker(A)$。因此，只需对 $\Coker(A)$ 证明本引理。
由引理 \ref{models-lemma-bound-wm} 可知，有 $m_i|a_{ij}| \leq 768g$，
其中 $i, j$ 任意。因此结论由引理
\ref{models-lemma-recurring-symmetric-integer} 得出。
\end{proof}







\section{模型}
\label{models-section-models}

\noindent
在本章中，$R$ 表示一个离散赋值环，$K$ 表示它的分式域。必要时，
我们用 $\pi \in R$ 表示一个一致化元，并用 $k = R/(\pi)$ 表示它的剩余域。

\medskip\noindent
设 $V$ 是一个代数 $K$-概形
（Varieties, Definition \ref{varieties-definition-algebraic-scheme}）。
$V$ 的一个{\it 模型}是指一个平坦有限型\footnote{有时允许模型在 $R$ 上局部有限型会很有用；
到需要时我们再处理这一情形。}
态射 $X \to \Spec(R)$，并配有一个同构
$V \to X_K = X \times_{\Spec(R)} \Spec(K)$。我们通常把 $V$ 与一般纤维
$X_K$（即 $X$ 的一般纤维）等同起来，并直接写作 $V = X_K$。
特殊纤维为 $X_k = X \times_{\Spec(R)} \Spec(k)$。
所谓{\it $X \to X'$ 这个 $V$ 的模型态射}，是指 $X \to X'$ 是一个
$R$ 上概形的态射，其在 $V$ 上诱导恒等态射。

\medskip\noindent
我们称{\it $X$ 是 $V$ 的一个固有模型}，如果 $X$ 是 $V$ 的一个模型
且结构态射 $X \to \Spec(R)$ 是固有的。
分离模型、光滑模型以及这里还应补充的其他模型也作类似定义。
我们称{\it $X$ 是 $V$ 的一个正则模型}，如果 $X$ 是 $V$ 的一个模型
且 $X$ 是正则概形。
正规模型、约化模型以及这里还应补充的其他模型也作类似定义。

\medskip\noindent
设 $R \subset R'$ 是离散赋值环的一个扩张
（More on Algebra, Definition
\ref{more-algebra-definition-extension-discrete-valuation-rings}）。
它诱导分式域的一个扩张 $K'/K$。
给定一个代数概形 $V$（定义在 $K$ 上），以 $V'$ 表示基变换
$V \times_{\Spec(K)} \Spec(K')$。于是有函子
$$
\text{models for }V\text{ 相对于 }R
\longrightarrow
\text{models for }V'\text{ 相对于 }R'
$$
把 $X$ 映到 $X \times_{\Spec(R)} \Spec(R')$。

\begin{lemma}
\label{models-lemma-closure-is-model}
设 $V_1 \to V_2$ 是 $K$ 上代数概形的闭浸入。
若 $X_2$ 是 $V_2$ 的一个模型，则 $V_1 \to X_2$ 的概形论像
是 $V_1$ 的一个模型。
\end{lemma}

\begin{proof}
由 Morphisms, Lemma \ref{morphisms-lemma-quasi-compact-scheme-theoretic-image}
与 Example \ref{morphisms-example-scheme-theoretic-image}，
问题归结为下面的代数陈述。设 $A_1$ 是一个在 $R$ 上平坦的有限型
$R$-代数。设 $A_1 \otimes_R K \to B_2$ 是满射。那么
$A_2 = A_1 / \Ker(A_1 \to B_2)$ 是一个在 $R$ 上平坦的有限型
$R$-代数，并且 $B_2 = A_2 \otimes_R K$。
我们略去详细证明；可用
More on Algebra, Lemma \ref{more-algebra-lemma-dedekind-torsion-free-flat}
证明 $A_2$ 平坦。
\end{proof}

\begin{lemma}
\label{models-lemma-normalization}
设 $X$ 是几何正规簇 $V$（定义在 $K$ 上）的一个模型。
则正规化 $\nu : X^\nu \to X$ 是有限的，而且 $X^\nu$ 到完备化
$R^\wedge$ 的基变换是 $X$ 的基变换的正规化。此外，对每个
$x \in X^\nu$，$\mathcal{O}_{X^\nu, x}$ 的完备化都是正规的。
\end{lemma}

\begin{proof}
注意 $R^\wedge$ 是一个离散赋值环
（More on Algebra, Lemma \ref{more-algebra-lemma-completion-dvr}）。
令 $Y = X \times_{\Spec(R)} \Spec(R^\wedge)$。
由于 $R^\wedge$ 是离散赋值环，我们有
$$
Y \setminus Y_k =
Y \times_{\Spec(R^\wedge)} \Spec(K^\wedge) =
V \times_{\Spec(K)} \Spec(K^\wedge)
$$
其中 $K^\wedge$ 是 $R^\wedge$ 的分式域。
因为 $V$ 几何正规，右端是正规概形。因此引理的第一部分由
Resolution of Surfaces, Lemma \ref{resolve-lemma-normalization-completion}
推出。

\medskip\noindent
为证明第二部分，由第一部分可设 $X$ 和 $Y$ 都正规。若 $x$ 位于一般纤维，
则 $\mathcal{O}_{X, x} = \mathcal{O}_{V, x}$ 是一个本质上有限型于某个域的
正规局部环。这种环是优环（More on Algebra, Proposition
\ref{more-algebra-proposition-ubiquity-excellent}）。
若 $x$ 是特殊纤维上的一点，其像为 $y \in Y$，则由
Resolution of Surfaces, Lemma \ref{resolve-lemma-iso-completions}，
$\mathcal{O}_{X, x}^\wedge = \mathcal{O}_{Y, y}^\wedge$。
在此情形下，$\mathcal{O}_{Y, y}$ 是优正规局部整环；这是因为由与上面相同的文献，
$R^\wedge$ 是优环。若 $B$ 是优正规局部整环，
则其完备化 $B^\wedge$ 正规（因为 $B \to B^\wedge$ 正则，且可应用
More on Algebra, Lemma \ref{more-algebra-lemma-normal-goes-up}）。
证明完毕。
\end{proof}

\begin{lemma}
\label{models-lemma-regular}
设 $X$ 是光滑曲线 $C$（定义在 $K$ 上）的一个模型。则 $X$ 存在奇点消解，
并且任意消解都是 $C$ 的一个模型。
\end{lemma}

\begin{proof}
我们验证 Lipman 定理的条件 (4)
（Resolution of Surfaces, Theorem \ref{resolve-theorem-resolve}）成立。
除 $X^\nu$ 仅有有限个奇点这一陈述外，这由
Lemma \ref{models-lemma-normalization} 立即可得。为证明该陈述，由
More on Algebra, Proposition \ref{more-algebra-proposition-ubiquity-J-2}
可知 $R$ 是 J-2 环，因而非奇异轨迹在 $X^\nu$ 中是开集。
由于 $X^\nu$ 正规且维数 $\leq 2$，奇点都是闭点；因此奇异轨迹是闭集
便意味着它只含有限多个点（因为 $X$ 拟紧）。注意 $X$ 的任意消解
都是 $X$ 的一个修改
（Resolution of Surfaces, Definition \ref{resolve-definition-resolution}）。
由 Varieties, Lemma
\ref{varieties-lemma-modification-normal-iso-over-codimension-1}，
它在 $X$ 的正规轨迹上是同构。由于正规点集包含 $C = X_K$，
故任意消解都是 $C$ 的一个模型。
\end{proof}

\begin{definition}
\label{models-definition-minimal-model}
设 $C$ 是 $K$ 上满足 $H^0(C, \mathcal{O}_C) = K$ 的光滑射影曲线。
所谓{\it 极小模型}，是指一个正则固有模型 $X$，它是 $C$ 的模型，
并且 $X$ 不包含第一类例外曲线
（Resolution of Surfaces, Section \ref{resolve-section-minus-one}）。
\end{definition}

\noindent
严格说来，这种对象应称为极小正则固有模型，甚至应称为相对极小正则射影模型。
不过，只要我们始终处理离散赋值环上的模型（本章正是如此），就不会产生混淆。

\medskip\noindent
极小模型总是存在
（Proposition \ref{models-proposition-exists-minimal-model}），并且在亏格 $> 0$
时唯一（Lemma \ref{models-lemma-minimal-model-unique}）。

\begin{lemma}
\label{models-lemma-pre-exists-minimal-model}
\begin{slogan}
曲线的正则固有模型可由极小模型经连续吹起得到
\end{slogan}
设 $C$ 是 $K$ 上满足 $H^0(C, \mathcal{O}_C) = K$ 的光滑射影曲线。
若 $X$ 是 $C$ 的一个正则固有模型，则存在态射序列
$$
X = X_m \to X_{m - 1} \to \ldots \to X_1 \to X_0
$$
其中各项均为 $C$ 的固有正则模型，每个态射都是一条第一类例外曲线的收缩，
且 $X_0$ 是一个极小模型。
\end{lemma}

\begin{proof}
由 Resolution of Surfaces, Lemma \ref{resolve-lemma-regular-dim-2-projective}
可知 $X$ 在 $R$ 上是射影的。因此，由
More on Morphisms, Lemma \ref{more-morphisms-lemma-projective}，
$X$ 有一个丰沛可逆层（下文将用到这一点）。
设 $E \subset X$ 是一条第一类例外曲线；参见
Resolution of Surfaces, Section \ref{resolve-section-minus-one}。
由 Resolution of Surfaces, Lemma \ref{resolve-lemma-contract-ample}，
可收缩 $E$；相应态射为 $X \to X'$，并且 $X'$ 正则且在 $R$ 上射影。
显然，$X'_k$ 的不可约分支数恰比 $X_k$ 的不可约分支数少一。
因此，这样的收缩只能执行有限次，最终便得到一个极小模型。
\end{proof}

\begin{proposition}
\label{models-proposition-exists-minimal-model}
设 $C$ 是 $K$ 上满足 $H^0(C, \mathcal{O}_C) = K$ 的光滑射影曲线。
则极小模型存在。
\end{proposition}

\begin{proof}
选取一个闭浸入 $C \to \mathbf{P}^n_K$。令 $X$ 为
$C \to \mathbf{P}^n_R$ 的概形论像。由
Lemma \ref{models-lemma-closure-is-model}，$X \to \Spec(R)$ 是 $C$ 的一个射影模型。
由 Lemma \ref{models-lemma-regular}，存在奇点消解 $X' \to X$，且 $X'$ 是
$C$ 的一个模型。于是，作为固有态射的复合，$X' \to \Spec(R)$ 是固有的。
现在应用 Lemma \ref{models-lemma-pre-exists-minimal-model} 即得一个极小模型。
\end{proof}





\section{正则模型的几何}
\label{models-section-special-fibre}

\noindent
本节描述一个固有正则模型 $X$ 的几何；它是光滑射影曲线 $C$ 的模型，
该曲线定义在 $K$ 上并满足 $H^0(C, \mathcal{O}_C) = K$。

\begin{lemma}
\label{models-lemma-divisor-special-fiber}
设 $X$ 是光滑曲线 $C$（定义在 $K$ 上）的一个正则模型。
\begin{enumerate}
\item 特殊纤维 $X_k$ 是 $X$ 上的一个有效 Cartier 除子；
\item 每个不可约分支 $C_i$（属于 $X_k$）都是 $X$ 上的一个有效 Cartier 除子；
\item $X_k = \sum m_i C_i$（有效 Cartier 除子之和），
其中 $m_i$ 是 $C_i$ 在 $X_k$ 中的重数；
\item $\mathcal{O}_X(X_k) \cong \mathcal{O}_X$。
\end{enumerate}
\end{lemma}

\begin{proof}
回忆 $R$ 是一个离散赋值环，其一致化元为 $\pi$，剩余域为
$k = R/(\pi)$。由于 $X \to \Spec(R)$ 平坦，元素 $\pi$ 在 $X$ 上仿射局部地
是非零因子（见 More on Algebra, Lemma
\ref{more-algebra-lemma-dedekind-torsion-free-flat}）。因此，若
$U = \Spec(A) \subset X$ 是仿射开集，则
$$
X_k \cap U = U_k = \Spec(A \otimes_R k) = \Spec(A/\pi A)
$$
并且 $\pi$ 是 $A$ 中的非零因子。因此，由 Divisors, Lemma
\ref{divisors-lemma-characterize-effective-Cartier-divisor}，
$X_k = V(\pi)$ 是一个有效 Cartier 除子。这就证明了 (1)。

\medskip\noindent
以上讨论表明，偶对 $(\mathcal{O}_X(X_k), 1)$ 与偶对
$(\mathcal{O}_X, \pi)$ 同构，从而证明了 (4)。

\medskip\noindent
由 Divisors, Lemma \ref{divisors-lemma-effective-Cartier-divisor-is-a-sum}，
存在两两不同的整有效 Cartier 除子 $D_i \subset X$ 和整数
$a_i \geq 0$，使得 $X_k = \sum a_i D_i$。可以删去所有除子
$D_i$ 中满足 $a_i = 0$ 的那些。此时，由有效 Cartier 除子加法的定义立即可知，
集合论上有 $X_k = \bigcup D_i$。因此 $C_i = D_i$ 是 $X_k$ 的不可约分支，
这就证明了 (2)。令 $\xi_i$ 为 $C_i$ 的一般点。则
$\mathcal{O}_{X, \xi_i}$ 是离散赋值环
（Divisors, Lemma \ref{divisors-lemma-integral-effective-Cartier-divisor-dvr}）。
一致化元 $\pi_i \in \mathcal{O}_{X, \xi_i}$ 是 $C_i$ 的一个局部方程，
而 $\pi$ 的像是 $X_k$ 的一个局部方程。由 $X_k = \sum a_i C_i$ 可知，
$\pi$ 与 $\pi_i^{a_i}$ 在 $\mathcal{O}_{X, \xi_i}$ 中生成同一个理想。
另一方面，$C_i$ 在 $X_k$ 中的重数为
$$
m_i = \text{length}_{\mathcal{O}_{C_i, \xi_i}} \mathcal{O}_{X_k, \xi_i} =
\text{length}_{\mathcal{O}_{C_i, \xi_i}} \mathcal{O}_{X, \xi_i}/(\pi) =
\text{length}_{\mathcal{O}_{C_i, \xi_i}} \mathcal{O}_{X, \xi_i}/(\pi_i^{a_i}) =
a_i
$$
参见 Chow Homology, Definition
\ref{chow-definition-cycle-associated-to-closed-subscheme}。
所以 $a_i = m_i$，从而 (3) 得证。
\end{proof}

\begin{lemma}
\label{models-lemma-gorenstein}
设 $X$ 是光滑曲线 $C$（定义在 $K$ 上）的一个正则模型。则
\begin{enumerate}
\item $X \to \Spec(R)$ 是相对维数为 $1$ 的 Gorenstein 态射；
\item 每个不可约分支 $C_i$（属于 $X_k$）都是 Gorenstein 的。
\end{enumerate}
\end{lemma}

\begin{proof}
由于 $X \to \Spec(R)$ 平坦，为证明 (1)，只需证明各纤维都是 Gorenstein 的
（Duality for Schemes, Lemma \ref{duality-lemma-gorenstein-morphism}）。
一般纤维是光滑曲线，因而正则，进而是 Gorenstein 的
（Duality for Schemes, Lemma \ref{duality-lemma-regular-gorenstein}）。
对于特殊纤维 $X_k$，它是一个正则概形（因而是 Gorenstein 概形）上的有效
Cartier 除子，所以也是 Gorenstein 的；例如可用 Dualizing Complexes, Lemma
\ref{dualizing-lemma-gorenstein-divide-by-nonzerodivisor}。
同一论证说明各曲线 $C_i$ 也是 Gorenstein 的。
\end{proof}

\begin{situation}
\label{models-situation-regular-model}
设 $R$ 是离散赋值环，其分式域为 $K$，剩余域为 $k$，一致化元为 $\pi$。
设 $C$ 是定义在 $K$ 上且满足 $H^0(C, \mathcal{O}_C) = K$ 的光滑射影曲线。
设 $X$ 是 $C$ 的一个正则固有模型。设 $C_1, \ldots, C_n$ 是特殊纤维
$X_k$ 的不可约分支。依照 Lemma \ref{models-lemma-divisor-special-fiber}，写作
$X_k = \sum m_i C_i$。
\end{situation}

\begin{lemma}
\label{models-lemma-regular-model-connected}
在 Situation \ref{models-situation-regular-model} 中，特殊纤维 $X_k$ 是连通的。
\end{lemma}

\begin{proof}
这是 More on Morphisms, Lemma
\ref{more-morphisms-lemma-geometrically-connected-fibres-towards-normal}
的推论。
\end{proof}

\begin{lemma}
\label{models-lemma-regular-model-pic}
在 Situation \ref{models-situation-regular-model} 中，有正合序列
$$
0 \to \mathbf{Z} \to \mathbf{Z}^{\oplus n} \to
\Pic(X) \to \Pic(C) \to 0
$$
其中第一个映射把 $1$ 送到 $(m_1, \ldots, m_n)$，第二个映射把第
$i$ 个基向量送到 $\mathcal{O}_X(C_i)$。
\end{lemma}

\begin{proof}
注意 $C \subset X$ 是开子概形。由
Divisors, Lemma \ref{divisors-lemma-extend-invertible-module}，限制映射
$\Pic(X) \to \Pic(C)$ 是满射。设 $\mathcal{L}$ 是一个可逆
$\mathcal{O}_X$-模，并且存在同构 $s : \mathcal{O}_C \to \mathcal{L}|_C$。
于是 $s$ 是 $\mathcal{L}$ 的一个正则亚纯截面，而且有
$\text{div}_\mathcal{L}(s) = \sum a_i C_i$，其中 $a_i \in \mathbf{Z}$
（Divisors, Definition \ref{divisors-definition-divisor-invertible-sheaf}）。
由 Divisors, Lemma \ref{divisors-lemma-normal-c1-injective}
（以及 $X$ 正规这一事实），可得 $\mathcal{L} = \mathcal{O}_X(\sum a_iC_i)$。
最后，假设 $\mathcal{O}_X(\sum a_i C_i) \cong \mathcal{O}_X$。
那么存在一个元素 $g$，它属于 $X$ 的函数域，并满足
$\text{div}_X(g) = \sum a_i C_i$。特别地，有理函数 $g$ 在一般纤维
$C$（即 $X$ 的一般纤维）上既无零点也无极点。由于 $C$ 是正规概形，这说明
$g \in H^0(C, \mathcal{O}_C) = K$。因此有 $g = \pi^a u$，其中
$a \in \mathbf{Z}$ 且 $u \in R^*$。于是
$\text{div}_X(g) = a \sum m_i C_i$，证明完毕。
\end{proof}

\noindent
在 Situation \ref{models-situation-regular-model} 中，对每个可逆
$\mathcal{O}_X$-模 $\mathcal{L}$ 和每个 $i$，都有一个整数
$$
\deg(\mathcal{L}|_{C_i}) =
\chi(C_i, \mathcal{L}|_{C_i}) - \chi(C_i, \mathcal{O}_{C_i})
$$
它是 $\mathcal{L}$ 到 $C_i$ 的限制相对于基域 $k$ 的次数\footnote{注意可能出现域
$\kappa_i = H^0(C_i, \mathcal{O}_{C_i})$ 严格大于 $k$ 的情形。此时 $C_i$
上的每个可逆模，其次数（按上式定义）都被 $[\kappa_i : k]$ 整除。}；
参见 Varieties, Section \ref{varieties-section-divisors-curves}。

\begin{lemma}
\label{models-lemma-intersection-pairing}
在 Situation \ref{models-situation-regular-model} 中，给定 $\mathcal{L}$（一个可逆
$\mathcal{O}_X$-模）以及
$a = (a_1, \ldots, a_n) \in \mathbf{Z}^{\oplus n}$，定义
$$
\langle a, \mathcal{L} \rangle = \sum a_i\deg(\mathcal{L}|_{C_i})
$$
则 $\langle , \rangle$ 是双线性的，而且对
$b = (b_1, \ldots, b_n) \in \mathbf{Z}^{\oplus n}$ 有
$$
\left\langle a, \mathcal{O}_X(\sum b_i C_i) \right\rangle =
\left\langle b, \mathcal{O}_X(\sum a_i C_i) \right\rangle
$$
\end{lemma}

\begin{proof}
双线性立即由定义和 Varieties, Lemma
\ref{varieties-lemma-degree-tensor-product} 得出。为证明对称性，只需设 $a$ 和
$b$ 是 $\mathbf{Z}^{\oplus n}$ 中的标准基向量。因此只需证明
$$
\deg(\mathcal{O}_X(C_j)|_{C_i}) = \deg(\mathcal{O}_X(C_i)|_{C_j})
$$
对所有 $1 \leq i, j \leq n$ 成立。若 $i = j$，则无需证明。
若 $i \not = j$，则典范截面 $1$（属于 $\mathcal{O}_X(C_j)$）限制为
$\mathcal{O}_X(C_j)|_{C_i}$ 的一个非零截面（因而是正则截面），
其零点概形恰为 $C_i \cap C_j$（概形论相交）。换言之，
$C_i \cap C_j$ 是 $C_i$ 上的一个有效 Cartier 除子，并且
$$
\deg(\mathcal{O}_X(C_j)|_{C_i}) = \deg(C_i \cap C_j)
$$
由 Varieties, Lemma \ref{varieties-lemma-degree-effective-Cartier-divisor}
可得上式。利用对称性，另一端也得到同一个（！）公式，证明完毕。
\end{proof}

\noindent
在 Situation \ref{models-situation-regular-model} 中，常常可以方便地把
$\mathbf{Z}^{\oplus n}$ 看成集合 $\{C_1, \ldots, C_n\}$ 上的自由阿贝尔群。
我们把这个群的元素记作 $\sum a_i C_i$；这里将它视为形式和，尽管等价地，
我们也可以（而且有时确实会）把这样的和看作 $X$ 上支撑于特殊纤维
$X_k$ 的 Weil 除子。现在，Lemma \ref{models-lemma-intersection-pairing}
使我们可以在这个自由阿贝尔群上依照下式定义一个对称双线性形式
$(\ \cdot\ )$：
\begin{equation}
\label{models-equation-form}
\left(\sum a_i C_i \cdot \sum b_j C_j\right) =
\left\langle a, \mathcal{O}_X(\sum b_j C_j) \right\rangle =
\left\langle b, \mathcal{O}_X(\sum a_i C_i) \right\rangle
\end{equation}
下面证明这个双线性形式的一些性质。

\begin{lemma}
\label{models-lemma-properties-form}
在 Situation \ref{models-situation-regular-model} 中，对称双线性形式
(\ref{models-equation-form}) 具有以下性质：
\begin{enumerate}
\item $(C_i \cdot C_j) \geq 0$ 若 $i \not = j$，且等号成立当且仅当
$C_i \cap C_j = \emptyset$；
\item $(\sum m_i C_i \cdot C_j) = 0$；
\item 不存在非空真子集 $I \subset \{1, \ldots, n\}$，使得
$(C_i \cdot C_j) = 0$ 对 $i \in I$、$j \not \in I$ 成立；
\item $(\sum a_i C_i \cdot \sum a_i C_i) \leq 0$，且等号成立当且仅当
存在 $q \in \mathbf{Q}$，使得 $a_i = qm_i$ 对 $i = 1, \ldots, n$ 成立。
\end{enumerate}
\end{lemma}

\begin{proof}
在 Lemma \ref{models-lemma-intersection-pairing} 的证明中，我们看到
$(C_i \cdot C_j) = \deg(C_i \cap C_j)$ 若 $i \not = j$。
它满足 $\geq 0$，且 $> 0 $ 当且仅当 $C_i \cap C_j \not = \emptyset$。
这证明了 (1)。

\medskip\noindent
(2) 的证明。由 Lemma \ref{models-lemma-divisor-special-fiber}，与
$\sum m_i C_i$ 相伴的可逆层是平凡的，而平凡层的次数为零，故结论成立。

\medskip\noindent
(3) 的证明。借助 (1) 的证明中对相交数的描述，这正表达了 $X_k$ 连通这一事实
（Lemma \ref{models-lemma-regular-model-connected}）。

\medskip\noindent
(4) 由 (1)、(2)、(3) 以及 Lemma \ref{models-lemma-recurring-symmetric-real} 推出。
\end{proof}

\begin{lemma}
\label{models-lemma-multiple-fibre-normal-bundle}
在 Situation \ref{models-situation-regular-model} 中，令
$d = \gcd(m_1, \ldots, m_n)$，并令 $D = \sum (m_i/d)C_i$ 为有效
Cartier 除子。则 $\mathcal{O}_X(D)$ 的阶整除 $d$（在 $\Pic(X)$ 中），
并且 $\mathcal{C}_{D/X}$ 是一个可逆 $\mathcal{O}_D$-模，
它的阶整除 $d$（在 $\Pic(D)$ 中）。
\end{lemma}

\begin{proof}
我们有
$$
\mathcal{O}_X(D)^{\otimes d} = \mathcal{O}_X(dD) =
\mathcal{O}_X(X_k) = \mathcal{O}_X
$$
这是 Lemma \ref{models-lemma-divisor-special-fiber} 的结论。又因为
$\mathcal{C}_{D/X}$ 是 $\mathcal{O}_X(-D)$ 的拉回，故第二个结论也成立。
\end{proof}

\begin{lemma}
\label{models-lemma-regular-model-field}
\begin{reference}
\cite[Lemma 2.6]{Artin-Winters}
\end{reference}
在 Situation \ref{models-situation-regular-model} 中，令
$d = \gcd(m_1, \ldots, m_n)$。令 $D = \sum (m_i/d) C_i$ 为有效
Cartier 除子。则存在有效 Cartier 除子序列
$$
(X_k)_{red} = Z_0 \subset Z_1 \subset \ldots \subset Z_m = D
$$
使得 $Z_j = Z_{j - 1} + C_{i_j}$，其中某个
$i_j \in \{1, \ldots, n\}$ 对 $j = 1, \ldots, m$ 取值；并且
$H^0(Z_j, \mathcal{O}_{Z_j})$ 都是 $k$ 的有限扩域，其中 $j = 0, \ldots m$。
\end{lemma}

\begin{proof}
约化概形 $D_{red} = (X_k)_{red} = \sum C_i$ 是连通的
（Lemma \ref{models-lemma-regular-model-connected}），并且在 $k$ 上固有。因此由
Varieties, Lemma
\ref{varieties-lemma-proper-geometrically-reduced-global-sections}，
$H^0(D_{red}, \mathcal{O})$ 是一个域，也是 $k$ 的有限扩张。
所以结论对 $Z_0 = D_{red} = (X_k)_{red}$ 成立。假设已经构造出
$$
(X_k)_{red} = Z_0 \subset Z_1 \subset \ldots \subset Z_t \subset D
$$
其中 $Z_j = Z_{j - 1} + C_{i_j}$，某个 $i_j \in \{1, \ldots, n\}$
对 $j = 1, \ldots, t$ 取值，并且 $H^0(Z_j, \mathcal{O}_{Z_j})$
都是 $k$ 的有限扩域，其中 $j = 0, \ldots, t$。写
$Z_t = \sum a_i C_i$，其中 $1 \leq a_i \leq m_i/d$。
若 $a_i = m_i/d$ 对所有 $i$ 成立，则 $Z_t = D$，引理得证。
否则，有 $a_i < m_i/d$ 对某个 $i$ 成立，由
Lemma \ref{models-lemma-properties-form} 可得 $(Z_t \cdot Z_t) < 0$。
这意味着 $(D - Z_t \cdot Z_t) > 0$，因为该引理还给出
$(D \cdot Z_t) = 0$。因此可以找到一个 $i$，满足
$a_i < m_i/d$ 且 $(C_i \cdot Z_t) > 0$。令
$Z_{t + 1} = Z_t + C_i$ 以及 $i_{t + 1} = i$。考虑短正合序列
$$
0 \to \mathcal{O}_X(-Z_t)|_{C_i} \to \mathcal{O}_{Z_{t + 1}} \to
\mathcal{O}_{Z_t} \to 0
$$
它来自 Divisors, Lemma \ref{divisors-lemma-ses-add-divisor}。
根据我们对 $i$ 的选择，$\mathcal{O}_X(-Z_t)|_{C_i}$ 是固有曲线
$C_i$ 上的负次数可逆层，因而没有非零整体截面
（Varieties, Lemma \ref{varieties-lemma-check-invertible-sheaf-trivial}）。
因此，$H^0(\mathcal{O}_{Z_{t + 1}}) \subset H^0(\mathcal{O}_{Z_t})$
是一个域（这一点很清楚，也可由 Algebra, Lemma
\ref{algebra-lemma-integral-under-field} 得出），并且是 $k$ 的有限扩张。
这样便延长了序列。由于过程必然停止，例如因为
$t \leq \sum (m_i/d - 1)$，证明完毕。
\end{proof}

\begin{lemma}
\label{models-lemma-regular-model-genus}
\begin{reference}
\cite[Lemma 2.6]{Artin-Winters}
\end{reference}
在 Situation \ref{models-situation-regular-model} 中，令
$d = \gcd(m_1, \ldots, m_n)$。令 $D = \sum (m_i/d) C_i$ 为 $X$ 上的
有效 Cartier 除子。则
$$
1 - g_C = d [\kappa : k] (1 - g_D)
$$
其中 $g_C$ 是 $C$ 的亏格，$g_D$ 是 $D$ 的亏格，并且
$\kappa = H^0(D, \mathcal{O}_D)$。
\end{lemma}

\begin{proof}
由 Lemma \ref{models-lemma-regular-model-field} 可知 $\kappa$ 是 $k$ 的一个有限扩域。
又因为 $H^0(C, \mathcal{O}_C) = K$，所以 $C$ 与 $D$ 的亏格均有定义
（见 Algebraic Curves, Definition \ref{curves-definition-genus}），并且
$g_C = \dim_K H^1(C, \mathcal{O}_C)$ 以及
$g_D = \dim_\kappa H^1(D, \mathcal{O}_D)$。由
Derived Categories of Schemes, Lemma
\ref{perfect-lemma-chi-locally-constant-geometric}，有
$$
1 - g_C = \chi(C, \mathcal{O}_C) =
\chi(X_k, \mathcal{O}_{X_k}) = \dim_k H^0(X_k, \mathcal{O}_{X_k})
- \dim_k H^1(X_k, \mathcal{O}_{X_k})
$$
我们断言
$$
\chi(X_k, \mathcal{O}_{X_k}) = d \chi(D, \mathcal{O}_D)
$$
这将证明引理，因为
$$
\chi(D, \mathcal{O}_D) =
\dim_k H^0(D, \mathcal{O}_D) - \dim_k H^1(D, \mathcal{O}_D) =
[\kappa : k](1 - g_D)
$$
注意，作为有效 Cartier 除子有 $X_k = dD$。为证明该断言，我们对
$1 \leq r \leq d$ 归纳证明
$\chi(rD, \mathcal{O}_{rD}) = r \chi(D, \mathcal{O}_D)$。
基础情形 $r = 1$ 显然成立。若 $1 \leq r < d$，则考虑短正合序列
$$
0 \to \mathcal{O}_X(rD)|_D \to \mathcal{O}_{(r + 1)D} \to
\mathcal{O}_{rD} \to 0
$$
它来自 Divisors, Lemma \ref{divisors-lemma-ses-add-divisor}。由欧拉示性数的
可加性（Varieties, Lemma \ref{varieties-lemma-euler-characteristic-additive}），
只需证明 $\chi(D, \mathcal{O}_X(rD)|_D) = \chi(D, \mathcal{O}_D)$。
这是因为 $\mathcal{O}_X(rD)|_D$ 是 $\Pic(D)$ 中的挠元
（Lemma \ref{models-lemma-multiple-fibre-normal-bundle}），并且线丛的次数具有可加性
（Varieties, Lemma \ref{varieties-lemma-degree-tensor-product}），
所以挠可逆层的次数为零。
\end{proof}

\begin{lemma}
\label{models-lemma-exceptional-curves-dont-meet}
在 Situation \ref{models-situation-regular-model} 中，给定一对指标 $i, j$，
若 $C_i$ 和 $C_j$ 都是第一类例外曲线，并且
$C_i \cap C_j \not = \emptyset$，则 $n = 2$、$m_1 = m_2 = 1$、
$C_1 \cong \mathbf{P}^1_k$、$C_2 \cong \mathbf{P}^1_k$，而且 $C_1$ 与
$C_2$ 相交于一个 $k$-有理点，且 $C$ 的亏格为 $0$。
\end{lemma}

\begin{proof}
选取同构 $C_i = \mathbf{P}^1_{\kappa_i}$ 和
$C_j = \mathbf{P}^1_{\kappa_j}$。概形 $C_i \cap C_j$ 在 $C_i$ 与 $C_j$
中都是非空有效 Cartier 除子。因此
$$
(C_i \cdot C_j) = \deg(C_i \cap C_j) \geq \max([\kappa_i: k], [\kappa_j : k])
$$
第一个等号已在 Lemma \ref{models-lemma-intersection-pairing} 的证明中给出。
另一方面，自交数 $(C_i \cdot C_i)$ 等于 $\mathcal{O}_X(C_i)$ 在 $C_i$
上的次数，即 $-[\kappa_i : k]$，因为 $C_i$ 是第一类例外曲线。
对 $C_j$ 也同样成立。由 Lemma \ref{models-lemma-properties-form}，
$$
0 \geq (C_i + C_j)^2 = -[\kappa_i : k] + 2(C_i \cdot C_j) - [\kappa_j : k]
$$
这说明 $[\kappa_i : k] = \deg(C_i \cap C_j) = [\kappa_j : k]$，
而且有 $(C_i + C_j)^2 = 0$。再次考察该引理可知 $n = 2$、
$\{1, 2\} = \{i, j\}$ 且 $m_1 = m_2$。此外，概形论相交
$C_i \cap C_j$ 只包含一个点 $p$，其剩余域为 $\kappa$，并且
$\kappa_i \to \kappa \leftarrow \kappa_j$ 都是同构。令
$D = C_1 + C_2$ 为 $X$ 上的有效 Cartier 除子。注意 $D$ 是 $C_1$ 与
$C_2$ 的概形论并
（Divisors, Lemma \ref{divisors-lemma-sum-effective-Cartier-divisors-union}），
故有短正合序列
$$
0 \to \mathcal{O}_D \to \mathcal{O}_{C_1} \oplus \mathcal{O}_{C_2} \to
\mathcal{O}_p \to 0
$$
它来自 Morphisms, Lemma \ref{morphisms-lemma-scheme-theoretic-union}。
我们已知 $C_i \cong \mathbf{P}^1_\kappa$ 的上同调
（Cohomology of Schemes, Lemma
\ref{coherent-lemma-cohomology-projective-space-over-ring}），
故由上同调长正合序列可得 $H^0(D, \mathcal{O}_D) = \kappa$ 以及
$H^1(D, \mathcal{O}_D) = 0$。再由 Lemma \ref{models-lemma-regular-model-genus} 得到
$$
1 - g_C = d[\kappa : k](1 - 0)
$$
其中 $d = m_1 = m_2$。由此可得 $g_C = 0$、$d = m_1 = m_2 = 1$，
并且 $\kappa = k$。
\end{proof}





\section{极小模型的唯一性}
\label{models-section-uniqueness}

\noindent
若一般纤维的亏格为正，则极小模型是唯一的
（Lemma \ref{models-lemma-minimal-model-unique}），因而具有合适的映射性质
（Lemma \ref{models-lemma-minimal-model-mapping-property}）。

\begin{lemma}
\label{models-lemma-minimal-model-unique}
设 $C$ 是定义在 $K$ 上的光滑射影曲线，满足
$H^0(C, \mathcal{O}_C) = K$ 且亏格 $> 0$。
则 $C$ 有唯一的极小模型。
\end{lemma}

\begin{proof}
我们已经证明了引理中困难的部分，即极小模型的存在性（其证明依赖于
曲面奇点的消解）；参见 Proposition \ref{models-proposition-exists-minimal-model}。
为证明唯一性，设 $X$ 和 $Y$ 是两个极小模型。由
Resolution of Surfaces, Lemma \ref{resolve-lemma-birational-regular-surfaces}，
存在如下 $S$-态射图：
$$
X = X_0 \leftarrow X_1 \leftarrow \ldots \leftarrow X_n = Y_m
\to \ldots \to Y_1 \to Y_0 = Y
$$
其中每个态射都是在一个闭点处的吹起。态射 $X_n \to X_{n - 1}$ 的例外纤维
是一条第一类例外曲线 $E$。我们断言 $E$ 在态射 $X_n = Y_m \to Y$
下被收缩到一点。若断言成立，则由 Resolution of Surfaces, Lemma
\ref{resolve-lemma-factor-through-contraction}，$X_n \to Y$ 经由
$X_{n - 1}$ 分解。在这种情形下，由 Resolution of Surfaces, Lemma
\ref{resolve-lemma-proper-birational-regular-surfaces}，态射
$X_{n - 1} \to Y$ 仍是一列例外曲线的收缩。因此对 $n$ 归纳即可得出结论。
（基础情形 $n = 0$ 意味着存在收缩序列
$X = Y_m \to \ldots \to Y_1 \to Y_0 = Y$
以 $Y$ 结束。但由于 $X$ 是极小模型，它不包含第一类例外曲线，
故 $m = 0$ 且 $X = Y$。）

\medskip\noindent
断言的证明。我们对 $m$ 归纳证明：任意第一类例外曲线 $E \subset Y_m$
都被态射 $Y_m \to Y$ 映到一点。若 $m = 0$，则这由 $Y$ 是极小模型立即可知。
若 $m > 0$，则或者 $Y_m \to Y_{m - 1}$ 收缩 $E$（此时已经完成），
或者例外纤维 $E' \subset Y_m$（属于 $Y_m \to Y_{m - 1}$）是另一条
第一类例外曲线。由于 $E$ 和 $E'$ 都是特殊纤维的不可约分支，
且按假设 $g_C > 0$，故由 Lemma \ref{models-lemma-exceptional-curves-dont-meet}
可得 $E \cap E' = \emptyset$。于是 $E$ 在 $Y_{m - 1}$ 中的像是
一条第一类例外曲线（这是因为态射 $Y_m \to Y_{m - 1}$ 在 $E$ 的某个邻域上
是同构）。归纳假设表明 $Y_{m - 1} \to Y$ 收缩这条曲线，证明完毕。
\end{proof}

\begin{lemma}
\label{models-lemma-minimal-model-mapping-property}
设 $C$ 是定义在 $K$ 上的光滑射影曲线，满足
$H^0(C, \mathcal{O}_C) = K$ 且亏格 $> 0$。设 $X$ 是 $C$ 的极小模型
（Lemma \ref{models-lemma-minimal-model-unique}）。设 $Y$ 是 $C$ 的一个正则固有模型。
则存在唯一的模型态射 $Y \to X$，它是一列第一类例外曲线的收缩。
\end{lemma}

\begin{proof}
态射 $X \to Y$ 的存在性和性质立即由
Lemma \ref{models-lemma-pre-exists-minimal-model} 以及极小模型的唯一性得出。
态射 $Y \to X$ 是唯一的，因为 $C \subset Y$ 在概形论意义下稠密，
而且 $X$ 是分离的（见 Morphisms, Lemma
\ref{morphisms-lemma-equality-of-morphisms}）。
\end{proof}

\begin{example}
\label{models-example-nonunique-in-genus-zero}
若 $C$ 的亏格为 $0$，则极小模型确实不唯一。具体地，考虑闭子概形
$$
X \subset \mathbf{P}^2_R
$$
它由 $T_1T_2 - \pi T_0^2 = 0$ 定义。更准确地说，$X$ 定义为
$\text{Proj}(R[T_0, T_1, T_2]/(T_1T_2 - \pi T_0^2))$。于是特殊纤维
$X_k$ 是两条例外曲线 $C_1$、$C_2$ 的并，二者都同构于
$\mathbf{P}^1_k$（恰如 Lemma \ref{models-lemma-exceptional-curves-dont-meet} 中那样）。
从 $(0 : 1 : 0)$ 投影定义了一个态射 $X \to \mathbf{P}^1_R$，它收缩
$C_2$，并诱导 $C_1$ 与 $\mathbf{P}^1_R$ 的特殊纤维之间的同构。
从 $(0 : 0 : 1)$ 投影定义了一个态射 $X \to \mathbf{P}^1_R$，它收缩
$C_1$，并诱导 $C_2$ 与 $\mathbf{P}^1_R$ 的特殊纤维之间的同构。
更准确地说，这些态射对应于分次 $R$-代数映射
$$
R[T_0, T_1] \longrightarrow
R[T_0, T_1, T_2]/(T_1T_2 - \pi T_0^2) \longleftarrow
R[T_0, T_2]
$$
我们将在 Lemma \ref{models-lemma-nonuniqueness} 中研究这一现象。
\end{example}








\section{亏格公式}
\label{models-section-genus-formula}

\noindent
来自固有正则模型的组合结构还满足一个限制。

\begin{lemma}
\label{models-lemma-add-component}
在 Situation \ref{models-situation-regular-model} 中，假设有有效 Cartier 除子
$D, D' \subset X$，使得 $D' = D + C_i$ 对某个
$i \in \{1, \ldots, n\}$ 成立，并且 $D' \subset X_k$。则
$$
\chi(X_k, \mathcal{O}_{D'}) - \chi(X_k, \mathcal{O}_D) =
\chi(X_k, \mathcal{O}_X(-D)|_{C_i}) =
-(D \cdot C_i) + \chi(C_i, \mathcal{O}_{C_i})
$$
\end{lemma}

\begin{proof}
第二个等号由 (\ref{models-equation-form}) 中双线性形式 $(\ \cdot\ )$ 的定义以及
Lemma \ref{models-lemma-intersection-pairing} 得出。为证明第一个等号，分两种情形。
若 $C_i \not \subset D$，则 $D'$ 是 $D$ 与 $C_i$ 的概形论并
（由 Divisors, Lemma
\ref{divisors-lemma-sum-effective-Cartier-divisors-union}），
因而得到短正合序列
$$
0 \to \mathcal{O}_{D'} \to
\mathcal{O}_D \times \mathcal{O}_{C_i} \to
\mathcal{O}_{D \cap C_i} \to 0
$$
它来自 Morphisms, Lemma \ref{morphisms-lemma-scheme-theoretic-union}。
又有正合序列
$$
0 \to \mathcal{O}_X(-D)|_{C_i} \to
\mathcal{O}_{C_i} \to \mathcal{O}_{D \cap C_i} \to 0
$$
（Divisors, Remark \ref{divisors-remark-ses-regular-section}），
故由欧拉示性数的可加性
（Varieties, Lemma \ref{varieties-lemma-euler-characteristic-additive}）
可得所需断言。另一方面，若 $C_i \subset D$，则有正合序列
$$
0 \to \mathcal{O}_X(-D)|_{C_i} \to \mathcal{O}_{D'} \to \mathcal{O}_D \to 0
$$
它来自 Divisors, Lemma \ref{divisors-lemma-ses-add-divisor}，
由此立即看出引理成立。
\end{proof}

\begin{lemma}
\label{models-lemma-genus-formula}
在 Situation \ref{models-situation-regular-model} 中，有
$$
g_C = 1 + \sum\nolimits_{i = 1, \ldots, n}
m_i\left([\kappa_i : k] (g_i - 1) - \frac{1}{2}(C_i \cdot C_i)\right)
$$
其中 $\kappa_i = H^0(C_i, \mathcal{O}_{C_i})$，$g_i$ 是 $C_i$ 的亏格，
$g_C$ 是 $C$ 的亏格。
\end{lemma}

\begin{proof}
基本工具是 Derived Categories of Schemes, Lemma
\ref{perfect-lemma-chi-locally-constant-geometric}，它给出
$$
1 - g_C = \chi(C, \mathcal{O}_C) =
\chi(X_k, \mathcal{O}_{X_k})
$$
选取有效 Cartier 除子序列
$$
X_k = D_m \supset D_{m - 1} \supset \ldots \supset D_1 \supset D_0 = \emptyset
$$
使得 $D_{j + 1} = D_j + C_{i_j}$ 对每个 $j$ 都成立。（逐次把
$D_{j + 1}$ 的一个非零重数减一，显然可以选出这样的序列。）从
$\chi(\mathcal{O}_{D_0}) = 0$ 开始应用 Lemma \ref{models-lemma-add-component}，得到
\begin{align*}
1 - g_C
& =
\chi(X_k, \mathcal{O}_{X_k}) \\
& =
\sum\nolimits_j
\left(-(D_j \cdot C_{i_j}) +  \chi(C_{i_j}, \mathcal{O}_{C_{i_j}})\right) \\
& =
- \sum\nolimits_j
(C_{i_1} + C_{i_2} + \ldots + C_{i_{j - 1}} \cdot C_{i_j}) +
\sum\nolimits_j \chi(C_{i_j}, \mathcal{O}_{C_{i_j}}) \\
& =
-\frac{1}{2}\sum\nolimits_{j \not = j'} (C_{i_{j'}} \cdot C_{i_j}) +
\sum m_i \chi(C_i, \mathcal{O}_{C_i}) \\
& =
\frac{1}{2} \sum m_i(C_i \cdot C_i) + \sum m_i \chi(C_i, \mathcal{O}_{C_i})
\end{align*}
最后一个等号或许需要解释。由于 $\sum_j C_{i_j} = \sum m_i C_i$，
由 Lemma \ref{models-lemma-properties-form} 有
$(\sum_j C_{i_j} \cdot \sum_j C_{i_j}) = 0$。因此
$$
0 = \sum\nolimits_{j \not = j'} (C_{i_{j'}} \cdot C_{i_j}) +
\sum m_i(C_i \cdot C_i)
$$
这是把该乘积分成“非对角”项与“对角”项所得。注意，由
Varieties, Lemma \ref{varieties-lemma-regular-functions-proper-variety}，
$\kappa_i$ 是 $k$ 的有限扩域。因此 $C_i$ 的亏格有定义，并且
$\chi(C_i, \mathcal{O}_{C_i}) = [\kappa_i : k](1 - g_i)$。
汇总以上各式并重新排列各项，得到
$$
g_C = - \frac{1}{2}\sum m_i(C_i \cdot C_i) +
\sum m_i[\kappa_i : k](g_i - 1) + 1
$$
这也正是引理的结论。
\end{proof}

\begin{lemma}
\label{models-lemma-numerical-type-of-model}
在 Situation \ref{models-situation-regular-model} 中，令
$\kappa_i = H^0(C_i, \mathcal{O}_{C_i})$，并令 $g_i$ 为 $C_i$ 的亏格，
则数据
$$
n, m_i, (C_i \cdot C_j), [\kappa_i : k], g_i
$$
构成一个数值型，其亏格等于 $C$ 的亏格。
\end{lemma}

\begin{proof}
（在 Lemma \ref{models-lemma-genus-formula} 的证明中，我们已经看到引理陈述所用的量
都有良好定义。）需要验证 Definition \ref{models-definition-type} 的条件 (1)--(5)。

\medskip\noindent
条件 (1) 立即成立。

\medskip\noindent
条件 (2)。矩阵 $(C_i \cdot C_j)$ 的对称性由 Equation
(\ref{models-equation-form}) 和 Lemma \ref{models-lemma-intersection-pairing} 得出。
$(C_i \cdot C_j)$ 在 $i \not = j$ 时的非负性是
Lemma \ref{models-lemma-properties-form} 的 (1)。

\medskip\noindent
条件 (3) 是 Lemma \ref{models-lemma-properties-form} 的 (3)。

\medskip\noindent
条件 (4) 是 Lemma \ref{models-lemma-properties-form} 的 (2)。

\medskip\noindent
条件 (5) 来自以下事实：$(C_i \cdot C_j)$ 是 $C_i$ 上某个可逆模的次数，
因此被 $[\kappa_i : k]$ 整除；参见 Varieties, Lemma
\ref{varieties-lemma-divisible}。

\medskip\noindent
Lemma \ref{models-lemma-genus-formula} 中证明的亏格公式说明，该数值型具有陈述中的亏格；
参见 Definition \ref{models-definition-genus}。
\end{proof}

\begin{definition}
\label{models-definition-numerical-type-model}
在 Situation \ref{models-situation-regular-model} 中，所谓{\it 与 $X$ 相伴的数值型}，
是指 Lemma \ref{models-lemma-numerical-type-of-model} 所描述的数值型。
\end{definition}

\noindent
现在把模型的极小性与型的极小性对应起来。

\begin{lemma}
\label{models-lemma-numerical-type-minimal-model}
在 Situation \ref{models-situation-regular-model} 中，以下条件等价：
\begin{enumerate}
\item $X$ 是极小模型；
\item 与 $X$ 相伴的数值型是极小的。
\end{enumerate}
\end{lemma}

\begin{proof}
若数值型极小，则不存在 $i$ 同时满足 $g_i = 0$ 和
$(C_i \cdot C_i) = -[\kappa_i: k]$；参见
Definition \ref{models-definition-type-minimal}。这显然说明曲线 $C_i$ 中没有
第一类例外曲线。

\medskip\noindent
反过来，假设数值型不极小。则存在 $i$，使得 $g_i = 0$ 且
$(C_i \cdot C_i) = -[\kappa_i: k]$。我们断言这说明 $C_i$ 是第一类例外曲线。
事实上，可逆层 $\mathcal{O}_X(-C_i)|_{C_i}$ 的次数为
$-(C_i \cdot C_i) = [\kappa_i : k]$；这是把 $C_i$ 看作 $k$ 上的固有曲线时
所得。因此它的次数为 $1$；这是把 $C_i$ 看作 $\kappa_i$ 上的固有曲线时所得。
应用 Algebraic Curves, Proposition \ref{curves-proposition-projective-line} 可知
$C_i \cong \mathbf{P}^1_{\kappa_i}$，这是作为 $\kappa_i$ 上的概形而言。
由于域上的 $\mathbf{P}^1$ 的 Picard 群为
$\mathbf{Z}$，可知 $C_i$ 在 $X$ 中的法层同构于
$\mathcal{O}_{\mathbf{P}_{\kappa_i}}(-1)$，证明完毕。
\end{proof}

\begin{remark}
\label{models-remark-numerical-type-not-from-model}
并非每个数值型都来自模型，一个简单的理由是存在亏格为负的数值型。
也存在亏格为正的极小数值型，它并不是某个模型（定义在某个离散赋值环上）
所相伴的数值型，而该模型的一般纤维是光滑射影几何不可约曲线
（定义在该离散赋值环的分式域上）。一个简单例子是
$n = 1$, $m_1 = 1$, $a_{11} = 0$, $w_1 = 6$, $g_1 = 1$.
在此情形下，特殊纤维 $X_k$ 不会几何连通，因为它定义在一个扩域
$\kappa$ 上，而该域是 $k$ 的 $6$ 次扩张。这与一般纤维几何连通矛盾（见
More on Morphisms, Lemma
\ref{more-morphisms-lemma-geometrically-connected-fibres-towards-normal}）。
类似地，$n = 2$、$m_1 = m_2 = 1$、
$-a_{11} = -a_{22} = a_{12} = a_{21} = 6$、$w_1 = w_2 = 6$、
$g_1 = g_2 = 1$ 也因同一理由给出一个例子（细节略去）。
但当 $w_i$ 的最大公因数为 $1$ 时，我们没有例子。
\end{remark}

\begin{lemma}
\label{models-lemma-numerical-type-rational-point}
在 Situation \ref{models-situation-regular-model} 中，假设 $C$ 有一个 $K$-有理点。则
\begin{enumerate}
\item $X_k$ 有一个 $k$-有理点 $x$，它是 $X_k$ 在 $k$ 上的光滑点；
\item 若 $x \in C_i$，则 $H^0(C_i, \mathcal{O}_{C_i}) = k$ 且 $m_i = 1$；
\item $H^0(X_k, \mathcal{O}_{X_k}) = k$，并且 $X_k$ 的亏格等于 $C$ 的亏格。
\end{enumerate}
\end{lemma}

\begin{proof}
由于 $X \to \Spec(R)$ 固有，由固有性的赋值判别准则，$K$-有理点延拓为态射
$a : \Spec(R) \to X$（Morphisms, Lemma
\ref{morphisms-lemma-characterize-proper}）。令 $x \in X$ 为 $a$ 下
$\Spec(R)$ 的闭点的像。于是 $a$ 对应于一个 $R$-代数同态
$\psi : \mathcal{O}_{X, x} \to R$（见 Schemes, Section
\ref{schemes-section-points}）。因此 $\pi \not \in \mathfrak m_x^2$（因为 $\pi$
在 $R$ 中的像不属于 $\mathfrak m_R^2$）。故
$\mathcal{O}_{X_k, x} = \mathcal{O}_{X, x}/\pi \mathcal{O}_{X, x}$ 是正则的
（Algebra, Lemma \ref{algebra-lemma-regular-ring-CM}）。进而由 Algebra, Lemma
\ref{algebra-lemma-separable-smooth}，$X_k \to \Spec(k)$ 在 $x$ 处光滑。
因此 $x$ 包含在唯一一个不可约分支 $C_i$（属于 $X_k$）中，且
$\mathcal{O}_{C_i, x} = \mathcal{O}_{X_k, x}$，并且 $m_i = 1$。
$C_i$ 有 $k$-有理点这一事实说明域
$\kappa_i = H^0(C_i, \mathcal{O}_{C_i})$
（Varieties, Lemma \ref{varieties-lemma-regular-functions-proper-variety}）
等于 $k$。这证明了 (1)。我们有 $H^0(X_k, \mathcal{O}_{X_k}) = k$，
因为 $H^0(X_k, \mathcal{O}_{X_k})$ 是 $k$ 的一个扩域
（Lemma \ref{models-lemma-regular-model-field}），并映入
$H^0(C_i, \mathcal{O}_{C_i}) = k$。亏格相等由
Lemma \ref{models-lemma-regular-model-genus} 得出。
\end{proof}

\begin{lemma}
\label{models-lemma-genus-reduction-smaller}
在 Situation \ref{models-situation-regular-model} 中，假设 $X$ 是极小模型，
$\gcd(m_1, \ldots, m_n) = 1$，并且
$H^0((X_k)_{red}, \mathcal{O}) = k$。则映射
$$
H^1(X_k, \mathcal{O}_{X_k}) \to H^1((X_k)_{red}, \mathcal{O}_{(X_k)_{red}})
$$
是满射，而且只要 $(X_k)_{red} \not = X_k$，其核就非平凡。
\end{lemma}

\begin{proof}
由于次数 $\geq 2$ 的上同调在 $X_k$ 上消失
（Cohomology, Proposition \ref{cohomology-proposition-vanishing-Noetherian}），
$X_k$ 上阿贝尔层的任意满射都会在 $H^1$ 上诱导满射。考虑
Lemma \ref{models-lemma-regular-model-field} 的序列
$$
(X_k)_{red} = Z_0 \subset Z_1 \subset \ldots \subset Z_m = X_k
$$
由于域同态
$H^0(Z_j, \mathcal{O}_{Z_j}) \to
H^0((X_k)_{red}, \mathcal{O}_{(X_k)_{red}}) = k$ 是单射，
可得 $H^0(Z_j, \mathcal{O}_{Z_j}) = k$ 对 $j = 0, \ldots, m$ 成立。
于是 $H^0(X_k, \mathcal{O}_{X_k}) \to H^0(Z_{m - 1}, \mathcal{O}_{Z_{m - 1}})$
是满射。令 $C = C_{i_m}$。
则 $X_k = Z_{m - 1} + C$。令
$\mathcal{L} = \mathcal{O}_X(-Z_{m - 1})|_C$。于是 $\mathcal{L}$ 是一个可逆
$\mathcal{O}_C$-模。与 Lemma \ref{models-lemma-regular-model-field} 的证明相同，有如下凝聚层正合序列：
$$
0 \to \mathcal{L} \to \mathcal{O}_{X_k} \to \mathcal{O}_{Z_{m - 1}} \to 0
$$
这是 $X_k$ 上的序列。因此得到短正合序列
$$
0 \to
H^1(C, \mathcal{L}) \to H^1(X_k, \mathcal{O}_{X_k}) \to
H^1(Z_{m - 1}, \mathcal{O}_{Z_{m - 1}}) \to 0
$$
$\mathcal{L}$ 在 $C$ 上相对于 $k$ 的次数为
$$
(C \cdot -Z_{m - 1}) = (C \cdot C - X_k) = (C \cdot C)
$$
令 $\kappa = H^0(C, \mathcal{O}_C)$ 且 $w = [\kappa : k]$。
由可逆层次数的定义可知
$$
\chi(C, \mathcal{L}) =
\chi(C, \mathcal{O}_C) + (C \cdot C) =
w(1 - g_C) + (C \cdot C)
$$
其中 $g_C$ 是 $C$ 的亏格。该表达式 $< 0$，因为 $X$ 极小，
因而 $C$ 不是第一类例外曲线
（见 Lemma \ref{models-lemma-numerical-type-minimal-model} 的证明）。
故 $\dim_k H^1(C, \mathcal{L}) > 0$，证明完毕。
\end{proof}

\begin{lemma}
\label{models-lemma-genus-reduction-bigger-than}
在 Situation \ref{models-situation-regular-model} 中，假设 $X_k$ 有一个 $k$-有理点
$x$，它是 $X_k \to \Spec(k)$ 的光滑点。则
$$
\dim_k H^1((X_k)_{red}, \mathcal{O}_{(X_k)_{red}}) \geq
g_{top} + g_{geom}(X_k/k)
$$
其中 $g_{geom}$ 如 Algebraic Curves, Section
\ref{curves-section-genus-geometric-genus} 中所定义，而 $g_{top}$ 是与 $X_k$
相伴的数值型的拓扑亏格（Definition \ref{models-definition-top-genus}；
Definition \ref{models-definition-numerical-type-model}）。
\end{lemma}

\begin{proof}
我们将证明不等式
$$
\dim_k H^1(D, \mathcal{O}_D) \geq g_{top}(D) + g_{geom}(D/k)
$$
对所有连通约化有效 Cartier 除子 $D \subset (X_k)_{red}$（包含 $x$）成立；
证明对 $D$ 的不可约分支数归纳。这里 $g_{top}(D) = 1 - m + e$，其中
$m$ 是 $D$ 的不可约分支数，$e$ 是 $D$ 中彼此相交的无序分支对的数目。

\medskip\noindent
基础情形：$D$ 只有一个不可约分支。此时 $D = C_i$ 是包含 $x$ 的唯一不可约分支。
在此情形下，$\dim_k H^1(D, \mathcal{O}_D) = g_i$ 且 $g_{top}(D) = 0$。
由于 $C_i$ 有一个 $k$-有理光滑点，它是几何整的
（Varieties, Lemma \ref{varieties-lemma-variety-with-smooth-rational-point}）。
因此 $g_i$ 是 $C_{i, \overline{k}}$ 的亏格
（Algebraic Curves, Lemma \ref{curves-lemma-genus-base-change}）。
同时，$g_{geom}(D/k)$ 是正规化 $C_{i, \overline{k}}^\nu$
（属于 $C_{i, \overline{k}}$）的亏格。把 Algebraic Curves, Lemma
\ref{curves-lemma-genus-normalization} 应用于正规化态射
$C_{i, \overline{k}}^\nu \to C_{i, \overline{k}}$，得到
\begin{equation}
\label{models-equation-genus-change-special-component}
\text{genus of }C_{i, \overline{k}} \geq
\text{genus of }C_{i, \overline{k}}^\nu
\end{equation}
综合以上结果可知，在这种情形下
$\dim_k H^1(D, \mathcal{O}_D) \geq g_{top}(D) + g_{geom}(D/k)$。

\medskip\noindent
归纳步骤。设 $D$ 有多于 $1$ 个不可约分支。可以写成
$D = C_i + D'$，其中 $x \in D'$ 且 $D'$ 仍然连通。这是一个留给读者的
图论练习（提示：取 $C_i$ 为 $D$ 中离 $x$ 最远的分支）。下面计算各不变量的变化。
由于 $x \in D'$，有
$H^0(D, \mathcal{O}_D) = H^0(D', \mathcal{O}_{D'}) = k$。考察层的短正合序列
$$
0 \to \mathcal{O}_D \to \mathcal{O}_{C_i} \oplus \mathcal{O}_{D'}
\to \mathcal{O}_{C_i \cap D'} \to 0
$$
（Morphisms, Lemma \ref{morphisms-lemma-scheme-theoretic-union}），
并利用欧拉示性数的可加性，得到
\begin{align*}
\dim_k H^1(D, \mathcal{O}_D) - \dim_k H^1(D', \mathcal{O}_{D'})
& =
-\chi(\mathcal{O}_{C_i}) + \chi(\mathcal{O}_{C_i \cap D'}) \\
& =
w_i(g_i - 1) + \sum\nolimits_{C_j \subset D'} a_{ij}
\end{align*}
这里与 Lemma \ref{models-lemma-numerical-type-of-model} 一样，令
$w_i = [\kappa_i : k]$、$\kappa_i = H^0(C_i, \mathcal{O}_{C_i})$，
$g_i$ 为 $C_i$ 的亏格，并令 $a_{ij} = (C_i \cdot C_j)$。我们有
$$
g_{top}(D) - g_{top}(D') = -1 +
\sum\nolimits_{C_j \subset D'\text{ meeting }C_i} 1
$$
以及
$$
g_{geom}(D/k) - g_{geom}(D'/k) = g_{geom}(C_i/k)
$$
这个等式是 Algebraic Curves, Lemma \ref{curves-lemma-bound-geometric-genus}
的结论。把这些等式与归纳假设结合，只需证明
$$
w_i g_i - g_{geom}(C_i/k) +
\sum\nolimits_{C_j \subset D'\text{ meets } C_i} (a_{ij} - 1) - (w_i - 1)
$$
非负。事实上，有
\begin{equation}
\label{models-equation-genus-change}
w_i g_i \geq [\kappa_i : k]_s g_i \geq g_{geom}(C_i/k)
\end{equation}
第二个不等式由 Algebraic Curves, Lemma
\ref{curves-lemma-bound-geometric-genus-curve} 给出。另一方面，由于 $w_i$ 整除
$a_{ij}$（Varieties, Lemma \ref{varieties-lemma-divisible}），显然
\begin{equation}
\label{models-equation-change-intersections}
\sum\nolimits_{C_j \subset D'\text{ meets } C_i} (a_{ij} - 1) - (w_i - 1)
\geq 0
\end{equation}
这是因为至少有一个 $C_j \subset D'$ 与 $C_i$ 相交。
\end{proof}

\begin{lemma}
\label{models-lemma-equality-genus-reduction-bigger-than}
若 Lemma \ref{models-lemma-genus-reduction-bigger-than} 中的等号成立，则
\begin{enumerate}
\item $X_k$ 中包含 $x$ 的唯一不可约分支是 $k$ 上光滑射影几何不可约曲线；
\item 若 $C \subset X_k$ 是另一个不可约分支，则
$\kappa = H^0(C, \mathcal{O}_C)$ 是 $k$ 的有限可分扩张，$C$ 有一个
$\kappa$-有理点，并且 $C$ 在 $\kappa$ 上光滑。
\end{enumerate}
\end{lemma}

\begin{proof}
回顾 Lemma \ref{models-lemma-genus-reduction-bigger-than} 的证明可知，要使等号成立，
不等式 (\ref{models-equation-genus-change-special-component})、
(\ref{models-equation-genus-change}) 和 (\ref{models-equation-change-intersections})
都必须是等式。

\medskip\noindent
令 $C_i$ 为包含 $x$ 的不可约分支。等式
(\ref{models-equation-genus-change-special-component}) 通过
Algebraic Curves, Lemma \ref{curves-lemma-genus-normalization} 表明
$C_{i, \overline{k}}^\nu \to C_{i, \overline{k}}$ 是同构。因此
$C_{i, \overline{k}}$ 光滑，(1) 成立。

\medskip\noindent
接着，令 $C_i \subset X_k$ 是另一个不可约分支。可以假设有
$D =  D' + C_i$，如 Lemma \ref{models-lemma-genus-reduction-bigger-than} 证明的
归纳步骤中那样。等式 (\ref{models-equation-genus-change}) 立即说明
$\kappa_i/k$ 是有限可分扩张。等式 (\ref{models-equation-change-intersections})
说明：或者 $a_{ij} = 1$ 对某个 $j$ 成立，或者存在唯一的
$C_j \subset D'$ 与 $C_i$ 相交，并且 $a_{ij} = w_i$。在两种情形下，
$C_i$ 都有一个 $\kappa_i$-有理点 $c$，且在概形论意义下
$c = C_i \cap C_j$。由于 $\mathcal{O}_{X, c}$ 是正则局部环，这说明
$C_i$ 与 $C_j$ 的局部方程构成局部环 $\mathcal{O}_{X, c}$ 的一组正则参数。
于是由 Algebra, Lemma \ref{algebra-lemma-regular-ring-CM}，
$\mathcal{O}_{C_i, c}$ 是正则的。我们得出 $C_i \to \Spec(\kappa_i)$ 在 $c$
处光滑（Algebra, Lemma \ref{algebra-lemma-separable-smooth}）。因此 $C_i$
在 $\kappa_i$ 上几何整
（Varieties, Lemma \ref{varieties-lemma-variety-with-smooth-rational-point}）。
最后还需证明 $C_i$ 在 $\kappa_i$ 上光滑。注意
$$
C_{i, \overline{k}} = C_i \times_{\Spec(k)} \Spec(\overline{k})
= \coprod\nolimits_{\kappa_i \to \overline{k}}
C_i \times_{\Spec(\kappa_i)} \Spec(\overline{k})
$$
其中有 $[\kappa_i : k]$ 个分支。因此，若 $C_i$ 在 $\kappa_i$ 上不光滑，
则这些曲线中的每一条都不光滑，进而都不正规；正规化态射会降低亏格
（Algebraic Curves, Lemma \ref{curves-lemma-genus-normalization}）。
这是不允许的，因为它会降低 $C_i/k$ 的几何亏格，与
$[\kappa_i : k] g_i = g_{geom}(C_i/k)$ 矛盾。
\end{proof}








\section{吹降例外曲线}
\label{models-section-contract}

\noindent
下面的引理说明，在正则固有模型中收缩一条第一类例外曲线时，相交数如何变化。
把它放在这里，主要是为了与 Lemma \ref{models-lemma-contract} 中引入的数值收缩比较。
我们将在 Remark \ref{models-remark-compare-contractions} 中比较几何收缩与数值收缩。

\begin{lemma}
\label{models-lemma-blowdown-regular-model}
在 Situation \ref{models-situation-regular-model} 中，假设 $C_n$ 是第一类例外曲线。
令 $f : X \to X'$ 为 $C_n$ 的收缩。令 $C'_i = f(C_i)$，并写作
$X'_k = \sum m'_i C'_i$。则 $X'$、$C'_i$、$i = 1, \ldots, n' = n - 1$
以及 $m'_i = m_i$ 如 Situation \ref{models-situation-regular-model} 中所述，并且
\begin{enumerate}
\item 对 $i, j < n$，有
$(C'_i \cdot C'_j) =
(C_i \cdot C_j) - (C_i \cdot C_n) (C_j \cdot C_n) /(C_n \cdot C_n)$,
\item 对 $i < n$，若 $C_i \cap C_n \not = \emptyset$，则有映射
$\kappa_i \leftarrow \kappa'_i \rightarrow \kappa_n$。
\end{enumerate}
这里 $\kappa_i = H^0(C_i, \mathcal{O}_{C_i})$ 且
$\kappa'_i = H^0(C'_i, \mathcal{O}_{C'_i})$。
\end{lemma}

\begin{proof}
由 Resolution of Surfaces, Lemma \ref{resolve-lemma-contract-ample}，
可以收缩 $C_n$；所用态射为 $f : X \to X'$，并且 $X'$ 正则且在 $R$ 上射影。
因此 $X'$ 如 Situation \ref{models-situation-regular-model} 中所述。令 $x \in X'$ 为
$C_n$ 的像。由于 $f$ 给出同构
$X \setminus C_n \to X' \setminus \{x\}$，显然 $m'_i = m_i$ 对 $i < n$ 成立。

\medskip\noindent
引理的 (2) 立即由态射 $C_i \to C'_i$ 和 $C_n \to x \to C'_i$ 的存在性得出。

\medskip\noindent
由 Divisors, Lemma \ref{divisors-lemma-blow-up-pullback-effective-Cartier}，
拉回 $f^{-1}C'_i$ 有定义。由 Divisors, Lemma
\ref{divisors-lemma-effective-Cartier-divisor-is-a-sum}，可知
$f^{-1}C'_i = C_i + e_i C_n$ 对某个 $e_i \geq 0$ 成立。由于
$\mathcal{O}_X(C_i + e_i C_n) = \mathcal{O}_X(f^{-1}C'_i) =
f^*\mathcal{O}_{X'}(C'_i)$
（Divisors, Lemma \ref{divisors-lemma-pullback-effective-Cartier-divisors}），
并且可逆层的拉回限制到 $C_n$ 上是平凡可逆层，故
$$
0 = \deg_{C_n}(\mathcal{O}_X(C_i + e_i C_n)) =
(C_i + e_i C_n \cdot C_n) = (C_i \cdot C_n) + e_i(C_n \cdot C_n)
$$
由于 $f_j = f|_{C_j} : C_j \to C_j$ 是 $k$ 上固有曲线之间的固有双有理态射，
由 Varieties, Lemma \ref{varieties-lemma-degree-birational-pullback} 可知，
$\deg_{C'_j}(\mathcal{O}_{X'}(C'_i)|_{C'_j})$ 等于
$\deg_{C_j}(f_j^*\mathcal{O}_{X'}(C'_i)|_{C'_j})$。考察交换图
$$
\xymatrix{
C_j \ar[r] \ar[d]_{f_j} & X \ar[d]^f \\
C'_j \ar[r] & X'
}
$$
并应用 Divisors, Lemma
\ref{divisors-lemma-pullback-effective-Cartier-divisors}，可得
$$
(C'_i \cdot C'_j) = \deg_{C'_j}(\mathcal{O}_{X'}(C'_i)|_{C'_j}) =
\deg_{C_j}(\mathcal{O}_X(C_i + e_i C_n)) = (C_i + e_i C_n \cdot C_j)
$$
代入上面求得的 $e_i$ 公式，便得到 (1)。
\end{proof}

\begin{remark}
\label{models-remark-genus-change}
在 Lemma \ref{models-lemma-blowdown-regular-model} 的情形下，还可以精确描述
亏格 $g_i$（属于 $C_i$）与亏格 $g'_i$（属于 $C'_i$）之间的关系。公式为
$$
g'_i = \frac{w_i}{w'_i}(g_i - 1) + 1 +
\frac{(C_i \cdot C_n)^2 - w_n(C_i \cdot C_n)}{2w'_iw_n}
$$
其中 $w_i = [\kappa_i : k]$、$w_n = [\kappa_n : k]$ 且
$w'_i = [\kappa'_i : k]$。为证明这一点，考虑短正合序列
$$
0 \to \mathcal{O}_{X'}(-C'_i) \to \mathcal{O}_{X'} \to
\mathcal{O}_{C'_i} \to 0
$$
及其到 $X$ 的拉回，即
$$
0 \to \mathcal{O}_X(-C'_i - e_iC_n) \to \mathcal{O}_X \to
\mathcal{O}_{C_i + e_i C_n} \to 0
$$
其中 $e_i$ 如 Lemma \ref{models-lemma-blowdown-regular-model} 的证明中所述。
由于 $Rf_*f^*\mathcal{L} = \mathcal{L}$ 对任意可逆模 $\mathcal{L}$
（在 $X'$ 上）都成立（细节略去），故
$$
Rf_*\mathcal{O}_{C_i + e_i C_n} = \mathcal{O}_{C'_i}
$$
作为 $X'_k$ 上的凝聚层复形成立。因此两端具有相同的欧拉示性数，
它也等于 $\mathcal{O}_{C_i + e_i C_n}$ 在 $X_k$ 上的欧拉示性数。
利用正合序列
$$
0 \to \mathcal{O}_{C_i + e_i C_n} \to
\mathcal{O}_{C_i} \oplus \mathcal{O}_{e_iC_n} \to
\mathcal{O}_{C_i \cap e_iC_n} \to 0
$$
并进一步过滤 $\mathcal{O}_{e_iC_n}$（细节略去），可得
$$
\chi(\mathcal{O}_{C'_i}) =
\chi(\mathcal{O}_{C_i}) - {e_i + 1 \choose 2}(C_n \cdot C_n)
- e_i(C_i \cdot C_n)
$$
由于 $e_i = -(C_i \cdot C_n)/(C_n \cdot C_n)$ 且
$(C_n \cdot C_n) = -w_n$，这就导出本注记开头的公式。
若以后需要这一结果，我们会把它表述为引理并给出详细证明。
\end{remark}

\begin{remark}
\label{models-remark-compare-contractions}
设 $f : X \to X'$ 如 Lemma \ref{models-lemma-blowdown-regular-model} 中所述。
设 $n, m_i, a_{ij}, w_i, g_i$ 是与 $X$ 相伴的数值型，
$n', m'_i, a'_{ij}, w'_i, g'_i$ 是与 $X'$ 相伴的数值型。
由 Lemma \ref{models-lemma-blowdown-regular-model} 和 Remark \ref{models-remark-genus-change}
立即可知，除 $w'_i$ 的取值外，这与 Lemma \ref{models-lemma-contract} 中数值型的收缩一致。
在几何情形下，$w'_i$ 是同时整除 $w_i$ 和 $w_n$ 的某个正整数。
在数值情形下，我们把 $w'_i$ 选为整除 $w_i$ 的尽可能大的整数，
使得 $g'_i$（由该公式给出）为整数。这个选择在数值设定中效果很好，
因为它有助于比较数值型的 Picard 群；参见
Lemma \ref{models-lemma-contract-picard-group}（尽管下文实际只用到单射性，
而对更小的 $w'_i$，这一单射性也成立）。
\end{remark}

\begin{lemma}
\label{models-lemma-nonuniqueness}
设 $C$ 是定义在 $K$ 上的光滑射影曲线，满足
$H^0(C, \mathcal{O}_C) = K$ 且亏格为 $0$。若 $C$ 有不止一个极小模型，
则每个极小模型的特殊纤维都同构于 $\mathbf{P}^1_k$。
\end{lemma}

\noindent
这个引理可以加强为：两个不同构极小模型之间的双有理变换，可以分解为一列如
Example \ref{models-example-nonunique-in-genus-zero} 中的初等变换。
若以后需要这一结论，我们会在这里精确表述并证明它。

\begin{proof}
设 $X$ 是 $C$ 的某个极小模型。与 $X$ 相伴的数值型亏格为 $0$ 且极小
（Definition \ref{models-definition-numerical-type-model} 和
Lemma \ref{models-lemma-numerical-type-minimal-model}）。因此由
Lemma \ref{models-lemma-genus-zero} 可知，$X_k$ 约化、不可约，满足
$H^0(X_k, \mathcal{O}_{X_k}) = k$，且亏格为 $0$。设 $Y$ 是 $C$ 的第二个
极小模型，它不同构于 $X$。由 Resolution of Surfaces, Lemma
\ref{resolve-lemma-birational-regular-surfaces}，存在如下 $S$-态射图：
$$
X = X_0 \leftarrow X_1 \leftarrow \ldots \leftarrow X_n = Y_m
\to \ldots \to Y_1 \to Y_0 = Y
$$
其中每个态射都是在一个闭点处的吹起。我们对 $m$ 归纳证明引理。
基础情形是 $m = 0$；此时，由于已假设 $Y$ 极小，这将意味着 $n = 0$，
但 $X$ 不同构于 $Y$，故这种情形不会发生，即无需验证任何事情。

\medskip\noindent
继续之前，注意 $n + 1 = m + 1$ 等于 $X_n = Y_m$ 的特殊纤维的不可约分支数，
因为 $X_k$ 和 $Y_k$ 都不可约。下文还会用到另一个观察：若
$X' \to X''$ 是 $C$ 的正则固有模型之间的态射，则 $X' \to X''$ 在
$X''$ 的某个开集上是同构，该开集的补是 $X''$ 的特殊纤维上有限个闭点；
参见 Varieties, Lemma
\ref{varieties-lemma-modification-normal-iso-over-codimension-1}。
事实上，任何这样的 $X' \to X''$ 都是一列闭点处的吹起
（Resolution of Surfaces, Lemma
\ref{resolve-lemma-proper-birational-regular-surfaces}），而吹起的数目等于
$X'$ 与 $X''$ 的特殊纤维的不可约分支数之差。

\medskip\noindent
令 $E_i \subset Y_i$（$m \geq i \geq 1$）为被态射
$Y_i \to Y_{i - 1}$ 收缩的曲线。令 $i$ 为满足如下条件的最大指标：
$E_i$ 的重数 $> 1$（在 $Y_i$ 的特殊纤维中）。于是后续吹起
$Y_m \to \ldots \to Y_{i + 1} \to Y_i$ 在 $E_i$ 上都是同构；否则，
$E_j$ 对某个 $j > i$ 的重数将 $> 1$。令 $E \subset Y_m$ 为 $E_i$ 的逆像。
由刚才的说明，$E \subset Y_m$ 是第一类例外曲线。令 $Y_m \to Y'$ 为 $E$
的收缩（其存在性由 Resolution of Surfaces, Lemma
\ref{resolve-lemma-contract-when-quasi-projective} 给出）。态射 $Y_m \to X$
必须收缩 $E$，因为 $X_k$ 约化。因此存在态射 $Y' \to Y$ 和 $Y' \to X$
（由 Resolution of Surfaces, Lemma
\ref{resolve-lemma-factor-through-contraction}），它们都是至多
$n - 1 = m - 1$ 个例外曲线收缩的复合（见以上讨论）。对 $m$ 应用归纳假设即可。
总之，可以假设所有曲线 $X_i$ 和 $Y_i$ 的特殊纤维都约化。

\medskip\noindent
由于 $X_i$ 与 $Y_i$ 的纤维约化，吹起 $X_i \to X_{i - 1}$ 和
$Y_i \to Y_{i - 1}$ 必须发生在特殊纤维的正则闭点处。事实上，若 $X''$
是 $C$ 的一个正则模型，$x \in X''$ 是特殊纤维的闭点，并且
$\pi \in \mathfrak m_x^2$，则例外纤维 $E$（属于吹起 $X' \to X''$，
中心为 $x$）的重数至少为 $2$（在 $X'$ 的特殊纤维中；局部计算略去）。因此，如所断言的，
$\mathcal{O}_{X''_k, x} = \mathcal{O}_{X'', x}/\pi$ 是正则的
（Algebra, Lemma \ref{algebra-lemma-regular-ring-CM}）。特别地，$x$ 是它所在的
唯一不可约分支 $Z'$（属于 $X''_k$）上的 Cartier 除子
（Varieties, Lemma \ref{varieties-lemma-regular-point-on-curve}）。
因此，严格变换 $Z \subset X'$（来自 $Z'$）同构地映到 $Z'$
（使用 Divisors, Lemmas \ref{divisors-lemma-strict-transform} 和
\ref{divisors-lemma-blow-up-effective-Cartier-divisor}）。换言之，若一个
不可约分支 $Z$（属于 $X_i$）在映射 $X_i \to X_j$（$i > j$）下不被收缩，
则它同构地映到其像。

\medskip\noindent
现在可以证明引理了。令 $E \subset Y_m$ 为被态射
$Y_m \to Y_{m - 1}$ 收缩的第一类例外曲线。若 $E$ 被态射
$Y_m = X_n \to X$ 收缩，则存在分解 $Y_{m - 1} \to X$
（Resolution of Surfaces, Lemma \ref{resolve-lemma-factor-through-contraction}），
而且 $Y_{m - 1} \to X$ 是一列闭点处的吹起
（Resolution of Surfaces, Lemma
\ref{resolve-lemma-proper-birational-regular-surfaces}）。在这种情形下，减小 $m$
并应用归纳假设即可。最后，假设 $E$ 不被态射 $Y_m \to X$ 收缩。
$E \to X_k$ 是满射，因为 $X_k$ 不可约；由上面的讨论，这意味着它是同构。
所以 $X_k$ 同构于射影直线，正是所需结论。
\end{proof}





\section{模型的 Picard 群}
\label{models-section-picard-groups-models}

\noindent
假设 $R, K, k, \pi, C, X, n, C_1, \ldots, C_n, m_1, \ldots, m_n$
如 Situation \ref{models-situation-regular-model} 中所述。
在 Lemma \ref{models-lemma-regular-model-pic} 中，我们得到了正合序列
$$
0 \to \mathbf{Z} \to \mathbf{Z}^{\oplus n} \to
\Pic(X) \to \Pic(C) \to 0
$$
我们希望用这个序列研究适当素数 $\ell$ 下这些 Picard 群中的 $\ell$-挠。

\begin{lemma}
\label{models-lemma-characterize-trivial}
在 Situation \ref{models-situation-regular-model} 中，令 $d = \gcd(m_1, \ldots, m_n)$。
若 $\mathcal{L}$ 是一个可逆 $\mathcal{O}_X$-模，并且
\begin{enumerate}
\item 它在 $C$ 上的限制是平凡可逆模，并且
\item 它的次数为 $0$（在每个 $C_i$ 上），
\end{enumerate}
则 $\mathcal{L}^{\otimes d} \cong \mathcal{O}_X$。
\end{lemma}

\begin{proof}
由 Lemma \ref{models-lemma-regular-model-pic}，有
$\mathcal{L} \cong \mathcal{O}_X(\sum a_i C_i)$，其中
$a_i \in \mathbf{Z}$。
$\mathcal{L}|_{C_j}$ 的次数是 $\sum_j a_j(C_i \cdot C_j)$。特别地，
$(\sum a_i C_i \cdot \sum a_i C_i) = 0$。
因此由 Lemma \ref{models-lemma-properties-form} 可知，
$(a_1, \ldots, a_n) = q(m_1, \ldots, m_n)$，其中
$q \in \mathbf{Q}$。于是 $\mathcal{L} = \mathcal{O}_X(lD)$，
其中 $l \in \mathbf{Z}$；
这里 $D = \sum (m_i/d) C_i$ 如
Lemma \ref{models-lemma-multiple-fibre-normal-bundle} 中所述，结论随即成立。
\end{proof}

\begin{lemma}
\label{models-lemma-canonical-map-of-pic}
在 Situation \ref{models-situation-regular-model} 中，
令 $T$ 为与 $X$ 相伴的数值型。存在典范映射
$$
\Pic(C) \to \Pic(T)
$$
其核恰由 $C$ 上如下可逆模组成：它们是可逆模 $\mathcal{L}$ 的限制，
其中该模定义在 $X$ 上，并且满足
$\deg_{C_i}(\mathcal{L}|_{C_i}) = 0$，其中 $i = 1, \ldots, n$。
\end{lemma}

\begin{proof}
回忆 $w_i = [\kappa_i : k]$，其中
$\kappa_i = H^0(C_i, \mathcal{O}_{C_i)})$；还回忆，$C_i$ 上任意可逆模的次数
均可被 $w_i$ 整除（Varieties, Lemma \ref{varieties-lemma-divisible}）。
因此可以考虑映射
$$
\frac{\deg}{w} : \Pic(X) \to \mathbf{Z}^{\oplus n}, \quad
\mathcal{L} \mapsto
(\frac{\deg(\mathcal{L}|_{C_1})}{w_1}, \ldots,
\frac{\deg(\mathcal{L}|_{C_n})}{w_n})
$$
在此映射下，$\mathcal{O}_X(C_j)$ 的像是
$$
((C_j \cdot C_1)/w_1, \ldots, (C_j \cdot C_n)/w_n) =
(a_{1j}/w_1, \ldots, a_{nj}/w_n)
$$
这恰是第 $j$ 个基向量在映射
$(a_{ij}/w_i) : \mathbf{Z}^{\oplus n} \to \mathbf{Z}^{\oplus n}$
下的像；该映射定义了 $T$ 的 Picard 群，参见
Definition \ref{models-definition-picard-group}。
因此，引理中的典范映射来自交换图
$$
\xymatrix{
\mathbf{Z}^{\oplus n} \ar[r] \ar[d]_{\text{id}} &
\Pic(X) \ar[r] \ar[d]^{\frac{\deg}{w}} &
\Pic(C) \ar[r] \ar[d] & 0 \\
\mathbf{Z}^{\oplus n} \ar[r]^{(a_{ij}/w_i)} &
\mathbf{Z}^{\oplus n} \ar[r] &
\Pic(T) \ar[r] & 0
}
$$
其两行均正合（上行由 Lemma \ref{models-lemma-regular-model-pic} 得到）。
关于核的描述是清楚的。
\end{proof}

\begin{lemma}
\label{models-lemma-sequence-torsion}
在 Situation \ref{models-situation-regular-model} 中，令
$d = \gcd(m_1, \ldots, m_n)$，并令 $T$ 为与 $X$ 相伴的数值型。
设 $h \geq 1$ 是与 $d$ 互素的整数。存在正合序列
$$
0 \to \Pic(X)[h] \to \Pic(C)[h] \to \Pic(T)[h]
$$
\end{lemma}

\begin{proof}
对 Lemma \ref{models-lemma-regular-model-pic} 的正合序列取 $h$-挠，
得到 $0 \to \Pic(X)[h] \to \Pic(C)[h]$ 的正合性，
因为 $h$ 与 $d$ 互素。
利用 Lemma \ref{models-lemma-canonical-map-of-pic} 的映射，可得映射
$\Pic(C)[h] \to \Pic(T)[h]$，它把 $\Pic(X)[h]$ 中的元素送到零。
反之，若 $\xi \in \Pic(C)[h]$ 在 $\Pic(T)[h]$ 中的像为零，
则可以找到一个可逆 $\mathcal{O}_X$-模 $\mathcal{L}$，
使得 $\deg(\mathcal{L}|_{C_i}) = 0$ 对所有 $i$ 都成立，且它在 $C$ 上的限制为
$\xi$。于是由 Lemma \ref{models-lemma-characterize-trivial}，
$\mathcal{L}^{\otimes h}$ 是 $d$-挠的。
取整数 $d'$ 使得 $dd' \equiv 1 \bmod h$。
这样的整数存在，因为 $h$ 与 $d$ 互素。
于是 $\mathcal{L}^{\otimes dd'}$ 是一个 $h$-挠可逆层（在 $X$ 上），
且它在 $C$ 上的限制为 $\xi$。
\end{proof}

\begin{lemma}
\label{models-lemma-torsion-embeds}
在 Situation \ref{models-situation-regular-model} 中，设 $h$ 是与 $k$ 的特征互素的整数。
则映射
$$
\Pic(X)[h] \longrightarrow \Pic((X_k)_{red})[h]
$$
是单射。
\end{lemma}

\begin{proof}
注意 $X \times_{\Spec(R)} \Spec(R/\pi^n)$ 是 $(X_k)_{red}$ 的有限阶增厚
（例如可由 Cohomology of Schemes, Lemma
\ref{coherent-lemma-power-ideal-kills-sheaf} 得出）。因此典范映射
$\Pic(X \times_{\Spec(R)} \Spec(R/\pi^n)) \to \Pic((X_k)_{red})$
在 $h$-挠上给出同构；这是由 More on Morphisms, Lemma
\ref{more-morphisms-lemma-torsion-pic-thickening} 以及关于 $h$ 的假设得到的。
因此，若 $\mathcal{L}$ 是一个 $h$-挠可逆层（在 $X$ 上），且其在
$(X_k)_{red}$ 上的限制是平凡层，则
$\mathcal{L}$ 在 $X \times_{\Spec(R)} \Spec(R/\pi^n)$ 上的限制对所有 $n$
都是平凡层。
于是
\begin{align*}
H^0(X, \mathcal{L})^\wedge
& =
\lim H^0(X \times_{\Spec(R)} \Spec(R/\pi^n),
\mathcal{L}|_{X \times_{\Spec(R)} \Spec(R/\pi^n)}) \\
& \cong
\lim H^0(X \times_{\Spec(R)} \Spec(R/\pi^n),
\mathcal{O}_{X \times_{\Spec(R)} \Spec(R/\pi^n)}) \\
& =
R^\wedge
\end{align*}
这里，第一个和最后一个等号使用形式函数定理
（Cohomology of Schemes, Theorem \ref{coherent-theorem-formal-functions}），
中间的同构例如使用 More on Algebra, Lemma
\ref{more-algebra-lemma-isomorphic-completions}。
由于 $H^0(X, \mathcal{L})$ 是有限 $R$-模，而 $R$ 是离散赋值环，
这说明 $H^0(X, \mathcal{L})$ 是秩为 $1$ 的自由 $R$-模。
取基元 $s \in H^0(X, \mathcal{L})$。
沿上述同构反向追踪可知，
$s|_{X \times_{\Spec(R)} \Spec(R/\pi^n)}$ 对所有 $n$ 都给出平凡化。
由于 $s$ 的零化轨迹在 $X$ 中闭，且 $X \to \Spec(R)$ 固有，
故 $s$ 的零化轨迹为空，正是所需结论。
\end{proof}






\section{半稳定约化}
\label{models-section-semistable-reduction}

\noindent
本节将仔细定义半稳定约化的含义。

\begin{example}
\label{models-example-blowup}
设 $R$ 是离散赋值环，其一致化元为 $\pi$。
给定 $n \geq 0$，考虑环映射
$$
R \longrightarrow A = R[x, y]/(xy - \pi^n)
$$
置 $X = \Spec(A)$ 且 $S = \Spec(R)$。
若 $n = 0$，则 $X \to S$ 光滑。
对所有 $n$，态射 $X \to S$ 都是 Algebraic Curves, Section
\ref{curves-section-families-nodal} 所定义的相对维数为 $1$ 的至多结点态射。
若 $n = 1$，则 $X$ 正则；但若 $n > 1$，则 $X$ 不正则，
因为 $(x, y)$ 不再生成极大理想 $\mathfrak m = (\pi, x, y)$。
为改善 $n > 1$ 时的情形，考虑吹起 $b : X' \to X$；
这是 $X$ 沿 $\mathfrak m$ 的吹起。参见 Divisors, Section
\ref{divisors-section-blowing-up}。
按构造，$X'$ 由三个仿射片覆盖，它们分别对应吹起代数
$A[\frac{\mathfrak m}{\pi}]$、$A[\frac{\mathfrak m}{x}]$ 和
$A[\frac{\mathfrak m}{y}]$。

\medskip\noindent
代数 $A[\frac{\mathfrak m}{\pi}]$ 有生成元
$x' = x/\pi$ 和 $y' = y/\pi$，并满足 $x'y' = \pi^{n - 2}$。
因此，$X'$ 的这一部分是 $R[x', y'](x'y' - \pi^{n - 2})$ 的谱。

\medskip\noindent
代数 $A[\frac{\mathfrak m}{x}]$ 有生成元 $x$、$u = \pi/x$，
并满足关系 $xu - \pi$。注意，此环包含
$y/x = \pi^n/x^2 = u^2\pi^{n - 2}$。因此，$X'$ 的这一部分正则。

\medskip\noindent
由对称性，代数 $A[\frac{\mathfrak m}{y}]$ 的情形与
$A[\frac{\mathfrak m}{x}]$ 的情形相同。

\medskip\noindent
因此可知，$X' \to S$ 是相对维数为 $1$ 的至多结点态射，并且
$X'$ 正则，至多除去一个点；该点有一个与上面完全相同的仿射开邻域，
只是把 $n$ 换成 $n - 2$。
对 $n$ 归纳，可得一列闭点处的吹起
$$
X_{\lfloor n/2 \rfloor} \to \ldots \to X_1 \to X_0 = X
$$
使得 $X_{\lfloor n/2 \rfloor} \to S$ 是相对维数为 $1$ 的至多结点态射，
且 $X_{\lfloor n/2 \rfloor}$ 正则。
\end{example}

\begin{lemma}
\label{models-lemma-etale-local-at-worst-nodal}
设 $R$ 是离散赋值环。设 $X$ 是相对维数为 $1$ 的至多结点概形
（在 $R$ 上）。
设 $x \in X$ 是 $X$ 在 $R$ 上的特殊纤维中的一点。
则存在交换图
$$
\xymatrix{
X \ar[d] &
U \ar[r] \ar[d] \ar[l] &
\Spec(A) \ar[dl] \\
\Spec(R) &
\Spec(R') \ar[l]
}
$$
其中 $R \subset R'$ 是离散赋值环的 \'etale 扩张，
态射 $U \to X$ 是 \'etale 的，态射 $U \to \Spec(A)$ 是 \'etale 的，
存在一点 $x' \in U$ 映到 $x$，并且
$$
A = R'[u, v]/(uv)
\quad\text{或}\quad
A = R'[u, v]/(uv - \pi^n)
$$
其中 $n \geq 0$，且 $\pi \in R'$ 是一致化元。
\end{lemma}

\begin{proof}
我们已经在远为一般的情形下证明了这个引理；参见
Algebraic Curves, Lemma \ref{curves-lemma-etale-local-structure-nodal-family}。
这里所要做的，只是把那里的陈述改写成上面的陈述。

\medskip\noindent
首先，若态射 $X \to \Spec(R)$ 在 $x$ 处光滑，则可以找到 \'etale 态射
$U \to \mathbf{A}^1_R = \Spec(R[u])$，其中 $U \subset X$ 是 $x$ 的某个
仿射开邻域。这就是
Morphisms, Lemma \ref{morphisms-lemma-smooth-etale-over-affine-space}。
必要时把坐标 $u$ 换成 $u + 1$，便可假设 $x$ 映到标准开集
$D(u) \subset \mathbf{A}^1_R$ 中的一点。此时 $D(u) = \Spec(A)$，其中
$A = R[u, v]/(uv - 1)$，因而结论在此情形下成立。

\medskip\noindent
其次，假设 $x$ 是该纤维的奇点。应用
Algebraic Curves, Lemma \ref{curves-lemma-etale-local-structure-nodal-family}
可得交换图
$$
\xymatrix{
X \ar[d] &
U \ar[rr] \ar[l] \ar[rd] & &
W \ar[r] \ar[ld] &
\Spec(\mathbf{Z}[u, v, a]/(uv - a)) \ar[d] \\
\Spec(R) & &
V \ar[ll] \ar[rr] & & \Spec(\mathbf{Z}[a])
}
$$
并具有所引引理陈述中的全部性质。
令 $x' \in U$ 为该引理保证的、映到 $x$ 的点。
先把 $V$ 缩小为 $x'$ 的像的一个仿射邻域，记作
$V = \Spec(R')$。于是 $R \to R'$ 是 \'etale 的。
由于 $R$ 是离散赋值环，$R'$ 是若干拟局部 Dedekind 整环的有限直积
（使用 More on Algebra, Lemma
\ref{more-algebra-lemma-Dedekind-etale-extension}）。
因此，例如使用素理想回避，可以找到标准开集
$D(f) \subset V = \Spec(R')$，它包含 $x'$ 的像，并使得 $R'_f$ 是离散赋值环。
把 $R'$ 替换为 $R'_f$，便达到如下情形：
$V = \Spec(R')$，且 $R \subset R'$ 是离散赋值环的 \'etale 扩张
（离散赋值环的扩张定义于 More on Algebra, Definition
\ref{more-algebra-definition-extension-discrete-valuation-rings}）。

\medskip\noindent
态射 $V \to \Spec(\mathbf{Z}[a])$ 由像 $h$ 决定；它是 $a$ 在
$R'$ 中的像。
于是 $W = \Spec(R'[u, v]/(uv - h))$。
因此，引理对 $A = R'[u, v]/(uv - h)$ 成立。
若 $h = 0$，便显然得到引理中所列的第一种情形。
若 $h \not = 0$，则可写成 $h = \epsilon \pi^n$，其中
$n \geq 0$，且 $\epsilon$ 是 $R'$ 的单位。
作坐标变换 $u_{new} = \epsilon u$ 和 $v_{new} = v$，
便得到引理所列的 $A$ 的第二种同构类型。
\end{proof}

\begin{lemma}
\label{models-lemma-blowup-at-worst-nodal}
设 $R$ 是离散赋值环。设 $X$ 是一个拟紧概形，它是
相对维数为 $1$ 的至多结点概形（在 $R$ 上），且一般纤维光滑。
则存在 $m \geq 0$ 以及一列态射
$$
X_m \to \ldots \to X_1 \to X_0 = X
$$
满足
\begin{enumerate}
\item $X_{i + 1} \to X_i$ 是闭点 $x_i$ 处的吹起，其中 $X_i$ 在该点奇异；
\item $X_i \to \Spec(R)$ 是相对维数为 $1$ 的至多结点态射；
\item $X_m$ 正则。
\end{enumerate}
\end{lemma}

\noindent
一个稍强而同样成立的陈述是：无论怎样在奇点处吹起，最终都会得到一个解消，
且所有中间吹起都是相对维数为 $1$ 的至多结点概形（在 $R$ 上）。

\begin{proof}
由于 $X$ 拟紧，其特殊纤维 $X_k$ 也拟紧。
由于 $X_k$ 的奇点至多为结点，$X_k$ 只有有限多个结点，并且在其余各点
都在 $k$ 上光滑。由于 $X \to \Spec(R)$ 平坦且一般纤维光滑，
$X$ 在 $R$ 上光滑，例外仅为 $X_k$ 的有限多个结点
（使用 Morphisms, Lemma \ref{morphisms-lemma-smooth-at-point}）。
由此，$X$ 在每一点都正则，可能的例外仅为其特殊纤维的结点
（参见 Algebra, Lemma \ref{algebra-lemma-regular-goes-up}）。
令 $x \in X$ 为这样的一个结点。选取交换图
$$
\xymatrix{
X \ar[d] &
U \ar[r] \ar[d] \ar[l] &
\Spec(A) \ar[dl] \\
\Spec(R) &
\Spec(R') \ar[l]
}
$$
如 Lemma \ref{models-lemma-etale-local-at-worst-nodal} 中所述。
注意，情形 $A = R'[u, v]/(uv)$ 不可能发生，因为这会意味着
$X/R$ 的一般纤维奇异（略去一个小细节）。
因此有 $A = R'[u, v]/(uv - \pi^n)$，其中 $n \geq 0$。
由于 $x$ 是奇点，故 $n \geq 2$；参见
Example \ref{models-example-blowup} 中的讨论。

\medskip\noindent
缩小 $U$ 后，可假设存在唯一一点 $u \in U$ 映到 $x$。
令 $w \in \Spec(A)$ 为 $u$ 的像。
还可假设 $u$ 是 $U$ 中映到 $w$ 的唯一一点。
由于两条水平箭头均为 \'etale 的，把 $u$ 视为 $U$ 的闭子概形，
便可知它既是 $x \in X$ 的概形论逆像，也是
$w \in \Spec(A)$ 的概形论逆像。
由于吹起与平坦基变换可交换
（Divisors, Lemma \ref{divisors-lemma-flat-base-change-blowing-up}），
得到交换图
$$
\xymatrix{
X' \ar[d] &
U' \ar[l] \ar[d] \ar[r] &
W' \ar[d] \\
X & U \ar[l] \ar[r] & \Spec(A)
}
$$
其中两个方块均为 Cartesian 方块，竖直箭头分别是沿
$x, u, w$ 的吹起（分别在 $X, U, \Spec(A)$ 中）。
概形 $W'$ 已在 Example \ref{models-example-blowup} 中描述过。
那里已经看到，$W'$ 是相对维数为 $1$ 的至多结点概形（在 $R'$ 上）。
因此 $W'$ 也是相对维数为 $1$ 的至多结点概形（在 $R$ 上）
（Algebraic Curves, Lemma
\ref{curves-lemma-nodal-family-postcompose-etale}）。
所以 $U'$ 是相对维数为 $1$ 的至多结点概形（在 $R$ 上）
（参见 Algebraic Curves, Lemma
\ref{curves-lemma-nodal-family-etale-local-source}）。
由于 $X' \to X$ 在 $x$ 的补集上是同构，再次使用
Algebraic Curves, Lemma
\ref{curves-lemma-nodal-family-etale-local-source}，
可知 $X'/R$ 也具有同样的性质。

\medskip\noindent
最后，需要说明，进行有限次这样的吹起后会得到正则模型 $X_m$。
这相当清楚，因为上述吹起使“不变量”$n$ 减少 $2$；参见
Example \ref{models-example-blowup} 中的计算。
不过，为避免精确定义这个不变量并证明其性质，我们改作如下论证。
若 $n = 2$，则 $W'$ 正则，因而 $X'$ 在位于 $x$ 上方的所有点处都正则，
并且 $X$ 的奇点数减少了 $1$。
若 $n > 2$，则 $w'$ 是 $W'$ 中位于 $w$ 上方的唯一奇点，并满足
$\kappa(w) = \kappa(w')$。因此 $U'$ 有唯一奇点 $u'$ 位于 $u$ 上方，
且 $\kappa(u) = \kappa(u')$。
显然，这意味着 $X'$ 有唯一奇点 $x'$ 位于 $x$ 上方，
即 $u'$ 的像。于是完全按照上面的论证，可得交换图
$$
\xymatrix{
X'' \ar[d] &
U'' \ar[l] \ar[d] \ar[r] &
W'' \ar[d] \\
X' & U' \ar[l] \ar[r] & W'
}
$$
其中两个方块均为 Cartesian 方块，竖直箭头分别是沿
$x', u', w'$ 的吹起（分别在 $X', U', W'$ 中）。
如此继续，便得到一列相容的吹起，并在
$\lfloor n/2 \rfloor$ 步之后停止。过程结束时，概形
$X^{(\lfloor n/2 \rfloor)}$ 的奇点数将比 $X$ 少一个。
对奇点数归纳即完成证明。
\end{proof}

\begin{lemma}
\label{models-lemma-blowdown-at-worst-nodal}
设 $R$ 是离散赋值环，其分式域为 $K$、剩余域为 $k$。
假设 $X \to \Spec(R)$ 是相对维数为 $1$ 的至多结点态射（在 $R$ 上）。
设 $X \to X'$ 是第一类例外曲线 $E \subset X$ 的收缩。
则 $X'$ 是相对维数为 $1$ 的至多结点概形（在 $R$ 上）。
\end{lemma}

\begin{proof}
具体地，令 $x' \in X'$ 为 $E$ 的像。
唯一需要说明的是，$X' \to \Spec(R)$ 是相对维数为 $1$ 的至多结点态射
（在 $x'$ 的某个邻域中）。
$X \to \Spec(R)$ 的闭纤维约化，所以 $\pi \in R$ 的消没阶为 $1$
（在 $E$ 上）。
这立即说明，把 $\pi$ 看作
$\mathfrak m_{x'} \subset \mathcal{O}_{X', x'}$ 的元素时，
它不属于 $\mathfrak m_{x'}^2$。
由于 $\mathcal{O}_{X', x'}$ 是维数为 $2$ 的正则环
（依 Resolution of Surfaces, Section \ref{resolve-section-minus-one}
中收缩的定义），可知 $\mathcal{O}_{X'_k, x'}$
是维数为 $1$ 的正则环
（Algebra, Lemma \ref{algebra-lemma-regular-ring-CM}）。
另一方面，曲线 $E$ 必须与至少另一个分支 $C$ 相交，其中该分支属于闭纤维
$X_k$。取 $x \in E \cap C$。
则 $x$ 是特殊纤维 $X_k$ 的结点，因而 $\kappa(x)/k$ 是有限可分扩张；
参见 Algebraic Curves, Lemma \ref{curves-lemma-nodal}。
由于 $x \mapsto x'$，可知 $\kappa(x')/k$ 是有限可分扩张。
由 Algebra, Lemma \ref{algebra-lemma-separable-smooth} 可知，
$X'_k \to \Spec(k)$ 在 $x'$ 的某个开邻域中光滑。
结合平坦性，这证明 $X' \to \Spec(R)$ 在 $x'$ 的某个邻域中光滑
（Morphisms, Lemma \ref{morphisms-lemma-smooth-at-point}）。
证明至此完成，因为相对维数为 $1$ 的光滑态射是相对维数为 $1$
的至多结点态射
（Algebraic Curves, Lemma
\ref{curves-lemma-smooth-relative-dimension-1}）。
\end{proof}

\begin{lemma}
\label{models-lemma-semistable}
设 $R$ 是分式域为 $K$ 的离散赋值环。
设 $C$ 是 $K$ 上满足 $H^0(C, \mathcal{O}_C) = K$ 的光滑射影曲线。
下列条件等价：
\begin{enumerate}
\item 存在 $C$ 的一个固有模型，它是相对维数为 $1$ 的至多结点概形（在 $R$ 上）；
\item 存在 $C$ 的一个极小模型，它是相对维数为 $1$ 的至多结点概形（在 $R$ 上）；
\item $C$ 的任意极小模型都是相对维数为 $1$ 的至多结点概形（在 $R$ 上）。
\end{enumerate}
\end{lemma}

\begin{proof}
为理解这个陈述，回忆：极小模型被定义为不含第一类例外曲线的正则固有模型
（Definition \ref{models-definition-minimal-model}）；极小模型存在
（Proposition \ref{models-proposition-exists-minimal-model}）；
而若 $C$ 的亏格 $> 0$，极小模型唯一
（Lemma \ref{models-lemma-minimal-model-unique}）。
据此，蕴含关系 (2) $\Rightarrow$ (1) 和 (3) $\Rightarrow$ (2) 是清楚的。

\medskip\noindent
假设 (1)。设 $X$ 是 $C$ 的一个固有模型，并且是
相对维数为 $1$ 的至多结点概形（在 $R$ 上）。
应用 Lemma \ref{models-lemma-blowup-at-worst-nodal}，还可假设 $X$ 正则。
令
$$
X = X_m \to X_{m - 1} \to \ldots \to X_1 \to X_0
$$
如 Lemma \ref{models-lemma-pre-exists-minimal-model} 中所述。
由 Lemma \ref{models-lemma-blowdown-at-worst-nodal} 和归纳法可知，
$X_0$ 是相对维数为 $1$ 的至多结点概形（在 $R$ 上）。

\medskip\noindent
要完成证明，只需说明 (2) 蕴含 (3)。
若 $C$ 的亏格 $> 0$，这一点是清楚的，因为此时极小模型唯一
（见以上讨论）。另一方面，若极小模型不唯一，则对任意极小模型，
态射 $X \to \Spec(R)$ 都光滑，因为由
Lemma \ref{models-lemma-nonuniqueness}，其特殊纤维同构于 $\mathbf{P}^1_k$。
\end{proof}

\begin{definition}
\label{models-definition-semistable}
设 $R$ 是分式域为 $K$ 的离散赋值环。
设 $C$ 是 $K$ 上满足 $H^0(C, \mathcal{O}_C) = K$ 的光滑射影曲线。
若 Lemma \ref{models-lemma-semistable} 中的等价条件成立，则称
$C$ 具有{\it 半稳定约化}。
\end{definition}

\begin{lemma}
\label{models-lemma-good}
设 $R$ 是分式域为 $K$ 的离散赋值环。
设 $C$ 是 $K$ 上满足 $H^0(C, \mathcal{O}_C) = K$ 的光滑射影曲线。
下列条件等价：
\begin{enumerate}
\item 存在 $C$ 的一个固有光滑模型；
\item 存在 $C$ 的一个极小模型，它在 $R$ 上光滑；
\item 任意极小模型都在 $R$ 上光滑。
\end{enumerate}
\end{lemma}

\begin{proof}
若 $X$ 是光滑固有模型，则其特殊纤维连通
（Lemma \ref{models-lemma-regular-model-connected}）且光滑，因而不可约。
这立即说明它是极小模型。因此 (1) 蕴含 (2)。
要完成证明，只需说明 (2) 蕴含 (3)。
若 $C$ 的亏格 $> 0$，这一点是清楚的，因为此时极小模型唯一
（Lemma \ref{models-lemma-minimal-model-unique}）。
另一方面，若极小模型不唯一，则对任意极小模型，
态射 $X \to \Spec(R)$ 都光滑，因为由
Lemma \ref{models-lemma-nonuniqueness}，其特殊纤维同构于 $\mathbf{P}^1_k$。
\end{proof}

\begin{definition}
\label{models-definition-good}
设 $R$ 是分式域为 $K$ 的离散赋值环。
设 $C$ 是 $K$ 上满足 $H^0(C, \mathcal{O}_C) = K$ 的光滑射影曲线。
若 Lemma \ref{models-lemma-good} 中的等价条件成立，则称
$C$ 具有{\it 良约化}。
\end{definition}








\section{亏格零情形的半稳定约化}
\label{models-section-semistable-reduction-genus-zero}

\noindent
本节对亏格零曲线证明半稳定约化定理
（Theorem \ref{models-theorem-semistable-reduction}）。

\medskip\noindent
设 $R$ 是分式域为 $K$ 的离散赋值环。
设 $C$ 是 $K$ 上满足 $H^0(C, \mathcal{O}_C) = K$ 的光滑射影曲线。
若 $C$ 的亏格为 $0$，则 $C$ 同构于一条二次曲线；
参见 Algebraic Curves, Lemma \ref{curves-lemma-genus-zero}。
因此存在有限可分扩张 $K'/K$，其次数至多为 $2$，使得
$C(K') \not = \emptyset$；参见 Algebraic Curves, Lemma
\ref{curves-lemma-smooth-plane-curve-point-over-separable}。
令 $R' \subset K'$ 为 $R$ 的整闭包；参见 More on Algebra, Remark
\ref{more-algebra-remark-finite-separable-extension} 中的讨论。
我们将证明，$C_{K'}$ 在 $R'_{\mathfrak m}$ 上具有半稳定约化，
其中 $\mathfrak m$ 遍历 $R'$ 的每个极大理想
（当然，在当前情形下至多有两个这样的理想）。
把 $R$ 替换为 $R'_{\mathfrak m}$，并把 $C$ 替换为 $C_{K'}$，
便约化到下一段所讨论的情形。

\medskip\noindent
本段中，$R$ 是分式域为 $K$ 的离散赋值环，$C$ 是 $K$ 上满足
$H^0(C, \mathcal{O}_C) = K$ 的光滑射影曲线，其亏格为 $0$，
并且 $C$ 有一个 $K$-有理点。在这种情形下，由
Algebraic Curves, Proposition \ref{curves-proposition-projective-line}，
有 $C \cong \mathbf{P}^1_K$。
因此可以用 $\mathbf{P}^1_R$ 作为模型，并可知 $C$ 同时具有良约化和半稳定约化。

\begin{example}
\label{models-example-extension-necessary-genus-zero}
令 $R = \mathbf{R}[[\pi]]$，并考虑概形
$$
X = V(T_1^2 + T_2^2 - \pi T_0^2) \subset \mathbf{P}^2_R
$$
$X$ 到 $\mathbf{C}[[\pi]]$ 的基变换同构于
Example \ref{models-example-nonunique-in-genus-zero} 中定义的概形，
因为有分解 $T_1^2 + T_2^2 = (T_1 + iT_2)(T_1 - iT_2)$
（在 $\mathbf{C}$ 上）。
因此 $X$ 正则，而其特殊纤维不可约但奇异；所以
$X$ 是其一般纤维的唯一极小模型
（使用 Lemma \ref{models-lemma-nonuniqueness}）。
由此可见，即使亏格为 $0$，扩张也是必要的。
\end{example}



\section{亏格一情形的半稳定约化}
\label{models-section-semistable-reduction-genus-one}

\noindent
本节对亏格一曲线证明半稳定约化定理
（Theorem \ref{models-theorem-semistable-reduction}）。
建议读者先阅读亏格 $\geq 2$ 情形的证明
（Section \ref{models-section-semistable-reduction-genus-at-least-two}）。
我们将尽可能使用 Lemma \ref{models-lemma-genus-one} 中给出的
亏格为 $1$ 的极小数值型分类。

\medskip\noindent
设 $R$ 是分式域为 $K$ 的离散赋值环。
设 $C$ 是 $K$ 上满足 $H^0(C, \mathcal{O}_C) = K$ 的光滑射影曲线。
假设 $C$ 的亏格为 $1$。
选取素数 $\ell \geq 7$，它不同于 $k$ 的特征。
选取有限可分扩张 $K'/K$，使得 $C(K') \not = \emptyset$ 且
$\Pic(C_{K'})[\ell] \cong (\mathbf{Z}/\ell \mathbf{Z})^{\oplus 2}$.
参见 Algebraic Curves, Lemma
\ref{curves-lemma-torsion-picard-becomes-visible}。
令 $R' \subset K'$ 为 $R$ 的整闭包；参见 More on Algebra, Remark
\ref{more-algebra-remark-finite-separable-extension}.
可以把 $R$ 替换为 $R'_{\mathfrak m}$，其中 $\mathfrak m$ 是
$R'$ 的某个极大理想，并把 $C$ 替换为 $C_{K'}$。
这就约化到下一段讨论的情形。

\medskip\noindent
在本节余下部分中，$R$ 是分式域为 $K$ 的离散赋值环，
$C$ 是 $K$ 上满足 $H^0(C, \mathcal{O}_C) = K$ 的光滑射影曲线，
其亏格为 $1$，具有一个 $K$-有理点，并满足
$\Pic(C)[\ell] \cong (\mathbf{Z}/\ell \mathbf{Z})^{\oplus 2}$
，其中 $\ell \geq 7$ 是不同于 $k$ 的特征的某个素数。
我们将证明 $C$ 具有半稳定约化。

\medskip\noindent
令 $X$ 为 $C$ 的一个极小模型；参见
Proposition \ref{models-proposition-exists-minimal-model}.
令 $T = (n, m_i, (a_{ij}), w_i, g_i)$ 为与 $X$ 相伴的数值型
（Definition \ref{models-definition-numerical-type-model}）。
则 $T$ 是极小数值型
（Lemma \ref{models-lemma-numerical-type-minimal-model}）。
由于 $C$ 有有理点，由 Lemma \ref{models-lemma-numerical-type-rational-point}，
存在 $i$ 使得 $m_i = w_i = 1$。
考察 Lemma \ref{models-lemma-genus-one} 中亏格为 $1$ 的极小数值型分类，
可知 $m = w = 1$，并且以下情形不允许：
(\ref{models-item-up4}),
(\ref{models-item-equal-down3}),
(\ref{models-item-up-up}),
(\ref{models-item-down-up}),
(\ref{models-item-up-equal-up}),
(\ref{models-item-down-equal-up}),
(\ref{models-item-triple-with-down}),
(\ref{models-item-equal-equal-down-equal}),
(\ref{models-item-up-equal-equal-up}),
(\ref{models-item-down-equal-equal-up}),
(\ref{models-item-triple-extended-down}),
(\ref{models-item-up-chain-equal-up}),
(\ref{models-item-down-chain-equal-up}),
(\ref{models-item-Dn-extended-down})（因为不存在同时满足 $w_i$ 和 $m_i$
等于 $1$ 的指标）。
令 $e$ 为指标对 $(i, j)$ 的数目，其中 $i < j$ 且 $a_{ij} > 0$。
对于其余情形，有
\begin{enumerate}
\item[(A)] 对以下情形，$e = n - 1$：
(\ref{models-item-one}),
(\ref{models-item-two-cycle}),
(\ref{models-item-equal-up3}),
(\ref{models-item-up-down}),
(\ref{models-item-up-equal-down}),
(\ref{models-item-triple-with-up}),
(\ref{models-item-equal-equal-up-equal}),
(\ref{models-item-up-equal-equal-down}),
(\ref{models-item-quadruple}),
(\ref{models-item-triple-extended-up}),
(\ref{models-item-up-chain-equal-down}),
(\ref{models-item-Dn-extended-up}),
(\ref{models-item-double-triple}),
(\ref{models-item-E6-completed}),
(\ref{models-item-E7-completed}) 以及
(\ref{models-item-E8-completed})；
\item[(B)] 对以下情形，$e = n$：
(\ref{models-item-three-cycle}),
(\ref{models-item-four-cycle}),
(\ref{models-item-five-cycle}) 以及
(\ref{models-item-n-cycle}).
\end{enumerate}
我们将分别讨论这两类情形。

\medskip\noindent
情形 (A)。在这种情形下，$\Pic(T)[\ell]$ 平凡
（数值型的 Picard 群定义于 Section \ref{models-section-picard-group}）。
消没性来自 $\Pic(T) \subset \Coker(A)$
（Lemma \ref{models-lemma-picard-T-and-A}）以及
$\Coker(A)[\ell] = 0$；后者由 Lemma
\ref{models-lemma-recurring-symmetric-integer} 和所选 $\ell$ 与 $a_{ij}$、$m_i$
互素这一事实得到。
由 Lemmas \ref{models-lemma-sequence-torsion} 和 \ref{models-lemma-torsion-embeds}，
可知存在嵌入
$$
(\mathbf{Z}/\ell \mathbf{Z})^{\oplus 2} \subset \Pic((X_k)_{red})[\ell].
$$
由 Algebraic Curves, Lemma \ref{curves-lemma-bound-torsion-simple}，得到
$$
2 \leq \dim_k H^1((X_k)_{red}, \mathcal{O}_{(X_k)_{red}}) +
g_{geom}((X_k)_{red}/k)
$$
由 Algebraic Curves, Lemmas \ref{curves-lemma-bound-geometric-genus} 和
\ref{curves-lemma-bound-geometric-genus-curve}
可知 $g_{geom}((X_k)_{red}/k) \leq \sum w_ig_i$。
Lemma \ref{models-lemma-genus-reduction-smaller} 的假设
由 Lemma \ref{models-lemma-numerical-type-rational-point} 得以满足，因而
$\dim_k H^1((X_k)_{red}, \mathcal{O}_{(X_k)_{red}}) \leq g = 1$。
合并这些不等式可得
$$
2 \leq 1 + \sum w_i g_i
$$
考察该列表可知，此数值型由
$n = 1$、$w_1 = m_1 = g_1 = 1$ 给出。
由于处处取等，可知 $g_{geom}(C_1/k) = 1$。
另一方面，已知 $C_1$ 有一个 $k$-有理点 $x$，使得
$C_1 \to \Spec(k)$ 在 $x$ 处光滑。
因此 $C_1$ 几何整
（Varieties, Lemma \ref{varieties-lemma-variety-with-smooth-rational-point}）。
于是 $g_{geom}(C_1/k) = 1$ 既等于 $C_{1, \overline{k}}$ 的正规化的亏格，
也等于 $C_{1, \overline{k}}$ 的亏格。
因此正规化态射
$C_{1, \overline{k}}^\nu \to C_{1, \overline{k}}$
是同构
（Algebraic Curves, Lemma \ref{curves-lemma-genus-normalization}）。
所以 $C_1$ 在 $k$ 上光滑，正是所需结论。

\medskip\noindent
情形 (B)。这里我们只能得出存在嵌入
$$
\mathbf{Z}/\ell \mathbf{Z} \subset \Pic(X_k)[\ell]
$$
由各型的分类可知，$m_i = w_i = 1$ 且 $g_i = 0$，这对每个 $i$ 成立。
因此每个 $C_i$ 都是 $k$ 上的亏格零曲线。
此外，对每个 $i$，存在 $j$，使得 $C_i \cap C_j$ 是一个 $k$-有理点。
于是由 Algebraic Curves, Proposition
\ref{curves-proposition-projective-line}，可得
$C_i \cong \mathbf{P}^1_k$。
特别地，由于 $X_k$ 是各 $C_i$ 的概形论并，故
$X_{\overline{k}}$ 是各 $C_{i, \overline{k}}$ 的概形论并。
因此 $X_{\overline{k}}$ 是维数为 $1$ 的约化连通固有概形（在 $\overline{k}$ 上），
且 $\dim_{\overline{k}}
H^1(X_{\overline{k}}, \mathcal{O}_{X_{\overline{k}}}) = 1$。
此外，由 Varieties, Lemma \ref{varieties-lemma-change-fields-pic}
以及上述结论，仍有
$$
\dim_{\mathbf{F}_\ell}(\Pic(X_{\overline{k}}) \geq 1
$$
由 Algebraic Curves, Proposition
\ref{curves-proposition-torsion-picard-reduced-proper}
可知 $X_{\overline{k}}$ 只有多重交叉奇点。
但由于 $X_k$ 是 Gorenstein 的（Lemma \ref{models-lemma-gorenstein}），
$X_{\overline{k}}$ 也是 Gorenstein 的（Duality for Schemes, Lemma
\ref{duality-lemma-gorenstein-base-change}）。
由 Algebraic Curves, Lemma
\ref{curves-lemma-multicross-gorenstein-is-nodal} 可知，
$X_{\overline{k}}$ 至多为结点。这就完成了情形 (B) 的证明。

\begin{example}
\label{models-example-extension-necessary-genus-one}
设 $k$ 是代数闭域。设 $Z$ 是 $k$ 上亏格为正数 $g$ 的光滑射影曲线。
设 $n \geq 1$ 是与 $k$ 的特征互素的整数。
设 $\mathcal{L}$ 是可逆 $\mathcal{O}_Z$-模，其阶为 $n$；参见
Algebraic Curves, Lemma
\ref{curves-lemma-torsion-picard-smooth-projective}。
选取同构 $\varphi : \mathcal{L}^{\otimes n} \to \mathcal{O}_Z$。
置 $R = k[[\pi]]$，其分式域为 $K = k((\pi))$。
用 $Z_R$ 表示 $Z$ 到 $R$ 的基变换。
令 $\mathcal{L}_R$ 为 $\mathcal{L}$ 到 $Z_R$ 的拉回。
考虑有限平坦态射
$$
p : X \longrightarrow Z_R
$$
使得
$$
p_*\mathcal{O}_X =
\text{Sym}^*_{\mathcal{O}_{Z_R}}(\mathcal{L}_R)/(\varphi - \pi) =
\mathcal{O}_{Z_R} \oplus \mathcal{L}_R \oplus
\mathcal{L}_R^{\otimes 2} \oplus \ldots \oplus \mathcal{L}_R^{\otimes n - 1}
$$
更精确地，若 $U = \Spec(A) \subset Z$ 是仿射开集，
使得 $\mathcal{L}|_U$ 由截面 $s$ 平凡化，并满足
$\varphi(s^{\otimes n}) = f$（其中 $f$ 是单位），则
$$
p^{-1}(U_R) = \Spec\left(
(A \otimes_R R[[\pi]])[x]/(x^n - \pi f)
\right)
$$
读者可验证，一般纤维之间的态射 $X_K \to Z_K$ 是有限 \'etale 的。
考察结构层的描述可知，$H^0(X, \mathcal{O}_X) = R$ 且
$H^0(X_K, \mathcal{O}_{X_K}) = K$。
由 Riemann--Hurwitz 公式
（Algebraic Curves, Lemma \ref{curves-lemma-rhe}），
$X_K$ 的亏格为 $n(g - 1) + 1$。特别地，$X_K$ 的亏格为 $1$，
若 $Z$ 的亏格为 $1$。
另一方面，由上述局部方程，概形 $X$ 正则；作为有效 Cartier 除子，
特殊纤维 $X_k$ 是约化特殊纤维的 $n$ 倍。
因此，任何有限扩张 $K'/K$，若在其上 $X_K$ 获得半稳定约化，
都必须分歧，且分歧指数至少为 $n$（略去一些细节）。
所以，不存在如下普适上界：使亏格为 $1$ 的曲线获得半稳定约化所需扩张的次数。
\end{example}






\section{亏格至少为二时的半稳定约化}
\label{models-section-semistable-reduction-genus-at-least-two}

\noindent
本节对亏格 $\geq 2$ 的曲线证明半稳定约化定理
（Theorem \ref{models-theorem-semistable-reduction}）。
固定 $g \geq 2$。

\medskip\noindent
设 $R$ 是分式域为 $K$ 的离散赋值环。
设 $C$ 是 $K$ 上满足 $H^0(C, \mathcal{O}_C) = K$ 的光滑射影曲线。
假设 $C$ 的亏格为 $g$。
选取素数 $\ell > 768g$，它不同于 $k$ 的特征。
选取有限可分扩张 $K'/K$，使得 $C(K') \not = \emptyset$ 且
$\Pic(C_{K'})[\ell] \cong (\mathbf{Z}/\ell \mathbf{Z})^{\oplus 2g}$.
参见 Algebraic Curves, Lemma
\ref{curves-lemma-torsion-picard-becomes-visible}。
令 $R' \subset K'$ 为 $R$ 的整闭包；参见 More on Algebra, Remark
\ref{more-algebra-remark-finite-separable-extension}.
可以把 $R$ 替换为 $R'_{\mathfrak m}$，其中 $\mathfrak m$ 是
$R'$ 的某个极大理想，并把 $C$ 替换为 $C_{K'}$。
这就约化到下一段讨论的情形。

\medskip\noindent
在本节余下部分中，$R$ 是分式域为 $K$ 的离散赋值环，
$C$ 是 $K$ 上满足 $H^0(C, \mathcal{O}_C) = K$ 的光滑射影曲线，
其亏格为 $g$，具有一个 $K$-有理点，并满足
$\Pic(C)[\ell] \cong (\mathbf{Z}/\ell \mathbf{Z})^{\oplus 2g}$
，其中 $\ell \geq 768g$ 是不同于 $k$ 的特征的某个素数。
我们将证明 $C$ 具有半稳定约化。

\medskip\noindent
在本节余下部分中，我们将不再另行说明地使用
Lemma \ref{models-lemma-numerical-type-rational-point} 的各项结论。

\medskip\noindent
令 $X$ 为 $C$ 的一个极小模型；参见
Proposition \ref{models-proposition-exists-minimal-model}.
令 $T = (n, m_i, (a_{ij}), w_i, g_i)$ 为与 $X$ 相伴的数值型
（Definition \ref{models-definition-numerical-type-model}）。
则 $T$ 是亏格为 $g$ 的极小数值型
（Lemma \ref{models-lemma-numerical-type-minimal-model}）。
由 Proposition \ref{models-proposition-bound-picard-group}，有
$$
\dim_{\mathbf{F}_\ell} \Pic(T)[\ell] \leq g_{top}
$$
由 Lemmas \ref{models-lemma-sequence-torsion} 和 \ref{models-lemma-torsion-embeds}，
可知存在嵌入
$$
(\mathbf{Z}/\ell \mathbf{Z})^{\oplus 2g - g_{top}}
\subset \Pic((X_k)_{red})[\ell].
$$
由 Algebraic Curves, Lemma \ref{curves-lemma-bound-torsion-simple}，得到
$$
2g - g_{top} \leq
\dim_k H^1((X_k)_{red}, \mathcal{O}_{(X_k)_{red}}) + g_{geom}(X_k/k)
$$
由 Lemmas \ref{models-lemma-genus-reduction-smaller} 和
\ref{models-lemma-genus-reduction-bigger-than}
，有
$$
g \geq \dim_k H^1((X_k)_{red}, \mathcal{O}_{(X_k)_{red}}) \geq
g_{top} + g_{geom}(X_k/k)
$$
初等数论告诉我们，这 $3$ 个不等式能够同时成立的唯一方式，是它们全都取等。
考察 Lemma \ref{models-lemma-genus-reduction-smaller}，可知
$m_i = 1$ 对所有 $i$ 都成立。
考察 Lemma \ref{models-lemma-equality-genus-reduction-bigger-than}，可知
$X_k$ 的每个不可约分支都在 $k$ 上光滑。

\medskip\noindent
特别地，由于 $X_k$ 是其不可约分支 $C_i$ 的概形论并，
可知 $X_{\overline{k}}$ 是各 $C_{i, \overline{k}}$ 的概形论并。
因此 $X_{\overline{k}}$ 是维数为 $1$ 的约化连通固有概形（在 $\overline{k}$ 上），
并且 $\dim_{\overline{k}}
H^1(X_{\overline{k}}, \mathcal{O}_{X_{\overline{k}}}) = g$。
此外，由 Varieties, Lemma \ref{varieties-lemma-change-fields-pic}
以及上述结论，仍有
$$
\dim_{\mathbf{F}_\ell}(\Pic(X_{\overline{k}})[\ell])
\geq 2g - g_{top} =
\dim_{\overline{k}} H^1(X_{\overline{k}}, \mathcal{O}_{X_{\overline{k}}})
+ g_{geom}(X_{\overline{k}})
$$
由 Algebraic Curves, Proposition
\ref{curves-proposition-torsion-picard-reduced-proper}
可知 $X_{\overline{k}}$ 只有多重交叉奇点。
但由于 $X_k$ 是 Gorenstein 的（Lemma \ref{models-lemma-gorenstein}），
$X_{\overline{k}}$ 也是 Gorenstein 的（Duality for Schemes, Lemma
\ref{duality-lemma-gorenstein-base-change}）。
由 Algebraic Curves, Lemma
\ref{curves-lemma-multicross-gorenstein-is-nodal} 可知，
$X_{\overline{k}}$ 至多为结点。证明完成。










\section{曲线的半稳定约化}
\label{models-section-semistable-reduction-theorem}

\noindent
本节完成该定理的证明。
对 $g \geq 2$，令 $768g < \ell' < \ell$ 为前两个素数 $> 768g$，并置
\begin{equation}
\label{models-equation-bound}
B_g = (2g - 2)(\ell^{2g})!
\end{equation}
$B_g$ 的精确形式并不重要；我们要强调的是，它只依赖于 $g$。

\begin{theorem}
\label{models-theorem-semistable-reduction}
\begin{reference}
\cite[Corollary 2.7]{DM}
\end{reference}
设 $R$ 是分式域为 $K$ 的离散赋值环。设 $C$ 是
$K$ 上满足 $H^0(C, \mathcal{O}_C) = K$ 的光滑射影曲线。
则存在离散赋值环的扩张 $R \subset R'$，它在分式域上诱导有限可分扩张
$K'/K$，使得 $C_{K'}$ 具有半稳定约化。
更精确地，有下列结论：
\begin{enumerate}
\item 若 $C$ 的亏格为零，则存在次数为 $2$ 的可分扩张
$K'/K$，使得 $C_{K'} \cong \mathbf{P}^1_{K'}$；因此
$C_{K'}$ 同构于光滑射影概形 $\mathbf{P}^1_{R'}$ 的一般纤维，
其中 $R'$ 是 $R$ 在 $K'$ 中的整闭包。
\item 若 $C$ 的亏格为一，则存在有限可分扩张 $K'/K$，使得
$C_{K'}$ 在 $R'_\mathfrak m$ 上具有半稳定约化；这对每个极大理想
$\mathfrak m$ 都成立，其中该理想属于整闭包 $R'$，即 $R$ 在 $K'$ 中的整闭包。
此外，$C_{K'}$ 在 $R'_\mathfrak m$ 上的（唯一）极小模型的特殊纤维，
要么是光滑亏格一曲线，要么是一圈有理曲线。
\item 若亏格 $g$（曲线 $C$ 的亏格）大于一，则存在有限可分扩张
$K'/K$，其次数至多为 $B_g$（\ref{models-equation-bound}），使得
$C_{K'}$ 在 $R'_\mathfrak m$ 上具有半稳定约化；这对每个极大理想
$\mathfrak m$ 都成立，其中该理想属于整闭包 $R'$，即 $R$ 在 $K'$ 中的整闭包。
\end{enumerate}
\end{theorem}

\begin{proof}
亏格零的情形见 Section \ref{models-section-semistable-reduction-genus-zero}。
亏格一的情形见 Section \ref{models-section-semistable-reduction-genus-one}。
亏格大于一的情形见
Section \ref{models-section-semistable-reduction-genus-at-least-two}。
为看出次数 $[K' : K]$ 有界，可以使用
Algebraic Curves, Lemma \ref{curves-lemma-torsion-picard-becomes-visible}
中证明的界：使所有 $\ell$-挠或 $\ell'$-挠可见所需扩张的次数有界。
（使用 $\ell$ 和 $\ell'$ 的原因是需要避开剩余域 $k$ 的特征。）
\end{proof}

\begin{remark}[改进该界]
\label{models-remark-improving-bound}
文献中的结果表明，Theorem \ref{models-theorem-semistable-reduction}
陈述中的界可以改进。例如，\cite{DM} 证明：若剩余域完美，
$C$ 的半稳定约化与其 Jacobian 的半稳定约化等价；这大概对一般剩余域也成立。
对阿贝尔簇而言，若 Galois 在 $\ell$-挠上的作用平凡，
其中 $\ell \geq 3$ 且不等于剩余特征，则它具有半稳定约化。
因此大概可以取 $\ell = 5$（在公式（\ref{models-equation-bound}）中用于 $B_g$）
（不过证明需要多得多的工作；若以后需要，我们会在这里给出精确陈述和证明）。
\end{remark}
% Generated mechanically from the frozen functors.tex target.
% Do not edit this generated file; edit the bound locale source instead.

% OK, start here.
%

\setcounter{chapter}{55}
\chapter{函子与态射}



\phantomsection
\label{functors-section-phantom}


\section{引言}
\label{functors-section-introduction}

\noindent
设 $X$ 和 $Y$ 为概形。本章围绕函子
$\QCoh(\mathcal{O}_Y) \to \QCoh(\mathcal{O}_X)$ 与概形态射
$X \to Y$ 之间的关系展开。更一般地，我们研究
$\QCoh(\mathcal{O}_X)$ 与 $X$ 之间的关系；若 $X$ 为 Noether 概形，
则研究 $\textit{Coh}(\mathcal{O}_X)$ 与 $X$ 之间的关系。
Gabriel 在 \cite{Gabriel} 中研究了这种关系。






\section{模范畴上的函子}
\label{functors-section-preliminary}

\noindent
对环 $A$，记 $\text{Mod}^{fp}_A$ 为有限表示 $A$-模所成的范畴。

\begin{lemma}
\label{functors-lemma-functor-on-fp-modules}
设 $A$ 为环，$\mathcal{B}$ 为具有滤过余极限的范畴，并设
$F : \text{Mod}^{fp}_A \to \mathcal{B}$ 为函子。则 $F$ 唯一地延拓为
与滤过余极限交换的函子 $F' : \text{Mod}_A \to \mathcal{B}$。
\end{lemma}

\begin{proof}
这由 Categories, Lemma
\ref{categories-lemma-extend-functor-by-colim} 得出。
为说明该引理适用，注意由 Algebra, Lemma
\ref{algebra-lemma-characterize-finitely-presented-module-hom}，
有限表示 $A$-模是 $\text{Mod}_A$ 中范畴意义下的紧对象。
此外，由 Algebra, Lemma \ref{algebra-lemma-module-colimit-fp}，
每个 $A$-模都是有限表示 $A$-模的滤过余极限。
\end{proof}

\noindent
若范畴 $\mathcal{B}$ 可加且具有滤过余极限，则 $\mathcal{B}$ 具有任意直和：
任一直和都可写成有限直和的滤过余极限。

\begin{lemma}
\label{functors-lemma-functor-on-fp-modules-additive}
设 $A$、$\mathcal{B}$、$F$ 如引理 \ref{functors-lemma-functor-on-fp-modules} 中所述。
假设 $\mathcal{B}$ 可加且 $F$ 可加。则 $F'$ 可加，并与任意直和交换。
\end{lemma}

\begin{proof}
为证明 $F'$ 可加，只需证明态射
$F'(M) \oplus F'(M') \to F'(M \oplus M')$ 对任意 $A$-模 $M$、$M'$
都是同构；参见
Homology, Lemma \ref{homology-lemma-additive-functor}。
把 $M = \colim_i M_i$ 和 $M' = \colim_j M'_j$ 写成有限表示
$A$-模 $M_i$ 的滤过余极限。于是
$F'(M) = \colim_i F(M_i)$、$F'(M') = \colim_j F(M'_j)$，以及
\begin{align*}
F'(M \oplus M')
& =
F'(\colim_{i, j} M_i \oplus M'_j) \\
& =
\colim_{i, j} F(M_i \oplus M'_j) \\
& =
\colim_{i, j} F(M_i) \oplus F(M'_j) \\
& =
F'(M) \oplus F'(M')
\end{align*}
正合所需。为证明 $F'$ 与直和交换，假设
$M = \bigoplus_{i \in I} M_i$。则
$M = \colim_{I' \subset I\text{ 有限}} \bigoplus_{i \in I'} M_i$
是一个滤过余极限。由此得到
\begin{align*}
F'(M)
& =
\colim_{I' \subset I\text{ 有限}}
F'(\bigoplus\nolimits_{i \in I'} M_i) \\
& =
\colim_{I' \subset I\text{ 有限}}
\bigoplus\nolimits_{i \in I'} F'(M_i) \\
& =
\bigoplus\nolimits_{i \in I} F'(M_i)
\end{align*}
第二个等号由已经证明的 $F'$ 的可加性成立。
\end{proof}

\noindent
若范畴 $\mathcal{B}$ 可加、具有滤过余极限且具有余核，则
$\mathcal{B}$ 具有任意余极限；参见上述讨论和 Categories, Lemma
\ref{categories-lemma-colimits-coproducts-coequalizers}。

\begin{lemma}
\label{functors-lemma-functor-on-fp-modules-right-exact}
设 $A$、$\mathcal{B}$、$F$ 如引理 \ref{functors-lemma-functor-on-fp-modules} 中所述。
假设 $\mathcal{B}$ 可加、具有余核，且 $F$ 右正合。则 $F'$ 可加、
右正合，并与任意直和交换。
\end{lemma}

\begin{proof}
由于 $F$ 右正合，$F$ 与二元余积交换，而这种余积由直和表示。
故由 Homology, Lemma \ref{homology-lemma-additive-functor}，$F$ 可加。
再由引理 \ref{functors-lemma-functor-on-fp-modules-additive}，$F'$ 可加并与直和交换。
我们建议读者自行证明 $F'$ 右正合，而不要阅读下面的证明。

\medskip\noindent
为证明 $F'$ 右正合，只需证明 $F'$ 与等化余子交换；参见
Categories, Lemma \ref{categories-lemma-characterize-right-exact}。
若 $a, b : K \to L$ 为 $A$-模态射，则 $a$ 与 $b$ 的等化余子是
$a - b : K \to L$ 的余核。因此，设
$K \to L \to M \to 0$ 为 $A$-模的正合序列。我们必须证明，在
$$
F'(K) \to F'(L) \to F'(M) \to 0
$$
第二个箭头是第一个箭头在 $\mathcal{B}$ 中的余核
（若 $\mathcal{B}$ 为 Abel 范畴，我们会说所示序列正合）。
由 Algebra, Lemma \ref{algebra-lemma-module-colimit-fp}，把
$M = \colim_{i \in I} M_i$ 写成有限表示 $A$-模的滤过余极限。
令 $L_i = L \times_M M_i$。
由此得到一族正合序列 $K \to L_i \to M_i \to 0$，它们以 $I$ 为参数。
由 Categories, Lemma \ref{categories-lemma-colimits-commute}，
余极限与余极限交换；又因为余核是一类等化余子，所以只需证明
$F'(L_i) \to F(M_i)$ 是 $F'(K) \to F'(L_i)$ 在 $\mathcal{B}$ 中的余核，
对所有 $i \in I$ 均成立。换言之，可以假设 $M$ 有限表示。
把 $L = \colim_{i \in I} L_i$ 写成有限表示 $A$-模的滤过余极限，
使每个 $L_i$ 都满射到 $M$。令 $K_i = K \times_L L_i$。
由此得到一族短正合序列 $K_i \to L_i \to M \to 0$，它们以 $I$ 为参数。
重复上述论证，可归约为证明：$F(L_i) \to F(M_i)$ 是
$F'(K) \to F(L_i)$ 在 $\mathcal{B}$ 中的余核，对所有 $i \in I$ 均成立。
换言之，可以假设 $L$ 和 $M$ 都是有限表示
$A$-模。此时模 $\Ker(L \to M)$ 是有限的
（Algebra, Lemma \ref{algebra-lemma-extension}）。因此可把
$K = \colim_{i \in I} K_i$ 写成有限表示 $A$-模的滤过余极限，
其中每一项都满射到 $\Ker(L \to M)$。由此得到一族短正合序列
$K_i \to L \to M \to 0$，它们以 $I$ 为参数。重复上述论证，
可归约为证明：$F(L) \to F(M)$ 是 $F(K_i) \to F(L)$
在 $\mathcal{B}$ 中的余核，对所有 $i \in I$ 均成立。换言之，可以假设
$K$、$L$、$M$ 都是有限表示 $A$-模。最后一种情形由 $F$ 右正合的假设得出。
\end{proof}

\noindent
若范畴 $\mathcal{B}$ 可加且具有核，则 $\mathcal{B}$ 具有有限极限。
事实上，有限积是存在的直和，而 $a, b : L \to M$ 的等化子是
存在的 $a - b : K \to L$ 的核。因此由 Categories, Lemma
\ref{categories-lemma-finite-limits-exist}，所有有限极限都存在。

\begin{lemma}
\label{functors-lemma-functor-on-fp-modules-left-exact}
设 $A$、$\mathcal{B}$、$F$ 如引理 \ref{functors-lemma-functor-on-fp-modules} 中所述。
假设 $A$ 是凝聚环（Algebra, Definition
\ref{algebra-definition-coherent}），$\mathcal{B}$ 可加且具有核，
滤过余极限与取核交换，并且 $F$ 左正合。则 $F'$ 可加、左正合，
并与任意直和交换。
\end{lemma}

\begin{proof}
由于 $A$ 凝聚，范畴 $\text{Mod}^{fp}_A$ 是 Abel 范畴，且其核与余核
同 $\text{Mod}_A$ 中的一样；参见 Algebra, Lemmas
\ref{algebra-lemma-coherent-ring} 和 \ref{algebra-lemma-coherent}。
因此 $\text{Mod}^{fp}_A$ 中所有有限极限都存在，并且 Categories,
Definition \ref{categories-definition-exact} 适用。由于 $F$ 左正合，
$F$ 与二元积交换，而这种积由直和表示。故由 Homology, Lemma
\ref{homology-lemma-additive-functor}，$F$ 可加。再由引理
\ref{functors-lemma-functor-on-fp-modules-additive}，$F'$ 可加并与直和交换。
我们建议读者自行证明 $F'$ 左正合，而不要阅读下面的证明。

\medskip\noindent
为证明 $F'$ 左正合，只需证明 $F'$ 与等化子交换；参见 Categories,
Lemma \ref{categories-lemma-characterize-left-exact}。若
$a, b : L \to M$ 为 $A$-模态射，则 $a$ 与 $b$ 的等化子是
$a - b : L \to M$ 的核。因此，设 $0 \to K \to L \to M$
为 $A$-模的正合序列。我们必须证明，在
$$
0 \to F'(K) \to F'(L) \to F'(M)
$$
箭头 $F'(K) \to F'(L)$ 是 $F'(L) \to F'(M)$ 在 $\mathcal{B}$ 中的核
（若 $\mathcal{B}$ 为 Abel 范畴，我们会说所示序列正合）。
由 Algebra, Lemma \ref{algebra-lemma-module-colimit-fp}，把
$M = \colim_{i \in I} M_i$ 写成有限表示 $A$-模的滤过余极限。
令 $L_i = L \times_M M_i$。由此得到一族正合序列
$0 \to K \to L_i \to M_i$，它们以 $I$ 为参数。根据假设，$\mathcal{B}$ 中的滤过余极限
与取核交换，故只需证明 $F'(K) \to F'(L_i)$ 是
$F'(L_i) \to F(M_i)$ 在 $\mathcal{B}$ 中的核，对所有 $i \in I$ 均成立。换言之，可以假设
$M$ 有限表示。把 $L = \colim_{i \in I} L_i$ 写成有限表示 $A$-模的
滤过余极限。令 $K_i = K \times_L L_i$。由此得到一族短正合序列
$0 \to K_i \to L_i \to M$，它们以 $I$ 为参数。重复上述论证，可归约为证明：
$F'(K_i) \to F(L_i)$ 是 $F(L_i) \to F(M)$ 在 $\mathcal{B}$ 中的核，
对所有 $i \in I$ 均成立。换言之，可以假设
$L$ 和 $M$ 都是有限表示 $A$-模。由于 $A$ 凝聚，$A$-模
$K = \Ker(L \to M)$ 具有有限表示，因为有限表示 $A$-模的范畴
是 Abel 范畴（见上述参考文献）。换言之，$K$、$L$、$M$ 三个模
都是有限表示 $A$-模。最后一种情形由 $F$ 左正合的假设得出。
\end{proof}

\noindent
若范畴 $\mathcal{B}$ 可加且具有余核，则 $\mathcal{B}$ 具有有限余极限。
事实上，有限余积是存在的直和，而 $a, b : K \to L$ 的等化余子是
存在的 $a - b : K \to L$ 的余核。因此由 Categories, Lemma
\ref{categories-lemma-colimits-exist}，所有有限余极限都存在。

\begin{lemma}
\label{functors-lemma-functor-on-modules-fp}
设 $A$ 为环，$\mathcal{B}$ 为具有余核的可加范畴。下列两个范畴等价：
\begin{enumerate}
\item 右正合函子 $F : \text{Mod}^{fp}_A \to \mathcal{B}$ 所成的范畴；
\item 二元组 $(K, \kappa)$ 所成的范畴，其中 $K \in \Ob(\mathcal{B})$，
且 $\kappa : A \to \text{End}_\mathcal{B}(K)$ 为环同态。
\end{enumerate}
此等价由把 $F$ 映到 $F(A)$ 并赋予后者自然 $A$-作用这一规则给出。
\end{lemma}

\begin{proof}
设 $(K, \kappa)$ 如 (2) 中所述。我们将构造函子
$F : \text{Mod}^{fp}_A \to \mathcal{B}$，使 $F(A) = K$ 带有所给的
$A$-作用 $\kappa$。具体地，对整数 $n \geq 0$，令
$$
F(A^{\oplus n}) = K^{\oplus n}
$$
给定 $A$-线性映射 $\varphi : A^{\oplus m} \to A^{\oplus n}$，
其矩阵为 $(a_{ij}) \in \text{Mat}(n \times m, A)$，定义
$$
F(\varphi) :
F(A^{\oplus m}) = K^{\oplus m}
\longrightarrow
K^{\oplus n} = F(A^{\oplus n})
$$
为矩阵为 $(\kappa(a_{ij}))$ 的映射。由此定义了可加函子 $F$，
它从 $\text{Mod}^{fp}_A$ 中以 $0$、$A$、$A^{\oplus 2}$、$\ldots$ 为对象的
全子范畴取值于 $\mathcal{B}$；验证从略。

\medskip\noindent
对每个对象 $M$（属于 $\text{Mod}^{fp}_A$），选取一个表示
$$
A^{\oplus m_M} \xrightarrow{\varphi_M} A^{\oplus n_M} \to M \to 0
$$
这是 $M$ 作为 $A$-模的表示。采用平凡表示
$0 \to A^{\oplus n} \xrightarrow{1} A^{\oplus n} \to 0$，若 $M = A^{\oplus n}$
（这并非必需，但可简化叙述）。对每个态射 $f : M \to N$
（属于 $\text{Mod}^{fp}_A$），可以选取交换图
\begin{equation}
\label{functors-equation-map}
\vcenter{
\xymatrix{
A^{\oplus m_M} \ar[r]_{\varphi_M} \ar[d]_{\psi_f} &
A^{\oplus n_M} \ar[r] \ar[d]_{\chi_f} &
M \ar[r] \ar[d]_f & 0 \\
A^{\oplus m_N} \ar[r]^{\varphi_N} &
A^{\oplus n_N} \ar[r] &
N \ar[r] & 0
}
}
\end{equation}
作出这些选择后，可以如下定义：对对象 $M$（属于 $\text{Mod}^{fp}_A$），令
$$
F(M) = \Coker(F(\varphi_M) : F(A^{\oplus m_M}) \to F(A^{\oplus n_M}))
$$
而对态射 $f : M \to N$（属于 $\text{Mod}^{fp}_A$），令
$$
F(f) = \text{映射 }F(M) \to F(N)\text{ 由下列对象诱导： }
F(\psi_f)\text{ 且 }F(\chi_f)\text{ on cokernels}
$$
注意，该规则延拓了已有的函子 $F$；该函子原定义于由自由模
$A^{\oplus n}$ 组成的全子范畴。还需证明 $F$ 是函子、$F$ 可加并且 $F$ 右正合。

\medskip\noindent
设 $f : M \to N$ 是 $\text{Mod}^{fp}_A$ 的态射。我们断言，上面定义的
$F(f)$ 与 (\ref{functors-equation-map}) 中 $\psi_f$ 和 $\chi_f$ 的选择无关。
事实上，设
$$
\xymatrix{
A^{\oplus m_M} \ar[r]_{\varphi_M} \ar[d]_\psi &
A^{\oplus n_M} \ar[r] \ar[d]_\chi &
M \ar[r] \ar[d]_f & 0 \\
A^{\oplus m_N} \ar[r]^{\varphi_N} &
A^{\oplus n_N} \ar[r] &
N \ar[r] & 0
}
$$
也交换。记 $F(f)' : F(M) \to F(N)$ 为由 $F(\psi)$ 和 $F(\chi)$
诱导的映射。考察这些交换图，由初等交换代数，存在映射
$\omega : A^{\oplus n_M} \to A^{\oplus m_N}$，使得
$\chi = \chi_f + \varphi_N \circ \omega$。应用 $F$，得到
$F(\chi) = F(\chi_f) + F(\varphi_N) \circ F(\omega)$。
由于 $F(N)$ 是 $F(\varphi_N)$ 的余核，映射
$F(A^{\oplus n_M}) \to F(M)$ 等化 $F(f)$ 与 $F(f)'$。
余核是满态射，故 $F(f) = F(f)'$。

\medskip\noindent
下面证明 $F$ 是函子。首先，$F(\text{id}_M) = \text{id}_{F(M)}$，
因为可以在上图中把 $\psi_f$ 和 $\chi_f$ 都取为恒等映射，
当 $f = \text{id}_M$ 时。其次，设有 $f : M \to N$ 和 $g : L \to M$。
此时 $\psi = \psi_f \circ \psi_g$ 和
$\chi = \chi_f \circ \chi_g$ 可用于 $f \circ g$ 的
(\ref{functors-equation-map})。因此它们诱导出正确的映射，恰好说明
$F(f) \circ F(g) = F(f \circ g)$。

\medskip\noindent
下面证明 $F$ 可加。设有 $f, g : M \to N$。此时
$\psi = \psi_f + \psi_g$ 和 $\chi = \chi_f + \chi_g$
可用于 $f + g$ 的 (\ref{functors-equation-map})。因此它们诱导出正确的映射，
恰好说明 $F(f) + F(g) = F(f + g)$。

\medskip\noindent
最后证明 $F$ 右正合。只需证明 $F$ 与等化余子交换；参见 Categories,
Lemma \ref{categories-lemma-characterize-right-exact}。为此，只需证明
$F$ 与余核交换。设 $K \to L \to M \to 0$ 为 $A$-模的正合序列，
且 $K$、$L$、$M$ 都有限表示。由于 $F$ 是可加函子，这当然给出复形
$$
F(K) \to F(L) \to F(M) \to 0
$$
而我们必须证明第二个箭头是第一个箭头在 $\mathcal{B}$ 中的余核。
无论如何，都得到一个映射 $\Coker(F(K) \to F(L)) \to F(M)$。
由初等交换代数，存在交换图
$$
\xymatrix{
A^{\oplus m_M} \ar[r]_{\varphi_M} \ar[d]_\psi &
A^{\oplus n_M} \ar[r] \ar[d]_\chi &
M \ar[r] \ar[d]_1 & 0 \\
K \ar[r] &
L \ar[r] &
M \ar[r] & 0
}
$$
对此图应用 $F$，并使用把 $F(M)$ 构造成 $F(\varphi_M)$ 的余核这一事实，
可知存在映射 $F(M) \to \Coker(F(K) \to F(L))$，它是映射
$\Coker(F(K) \to F(L)) \to F(M)$ 的右逆。这首先说明
$F(L) \to F(M)$ 总是满态射。其次，上述论证给出
$$
\Coker(F(K) \to F(L)) = F(M) \oplus E
$$
并且这一直接和分解同时与
$F(M) \to \Coker(F(K) \to F(L))$ 和
$\Coker(F(K) \to F(L)) \to F(M)$ 相容。然而，此时满态射
$p : F(L) \to E$ 与 $F(K) \to F(L)$ 复合后为零，
与 $F(A^{n_M}) \to F(L)$ 复合后也为零。又因为
$K \oplus A^{n_M} \to L$ 满射（省略代数论证），由上述结果可知
$F(K \oplus A^{n_M}) \to F(L)$ 是满态射，故 $E = 0$。证明完毕。
\end{proof}

\begin{lemma}
\label{functors-lemma-functor-on-modules}
设 $A$ 为环，$\mathcal{B}$ 为具有任意直和及余核的可加范畴。
下列两个范畴等价：
\begin{enumerate}
\item 右正合且与任意直和交换的函子
$F : \text{Mod}_A \to \mathcal{B}$ 所成的范畴；
\item 二元组 $(K, \kappa)$ 所成的范畴，其中 $K \in \Ob(\mathcal{B})$，
且 $\kappa : A \to \text{End}_\mathcal{B}(K)$ 为环同态。
\end{enumerate}
此等价由把 $F$ 映到 $F(A)$ 并赋予后者自然 $A$-作用这一规则给出。
\end{lemma}

\begin{proof}
合并引理 \ref{functors-lemma-functor-on-modules-fp} 和
\ref{functors-lemma-functor-on-fp-modules-right-exact} 即得。
\end{proof}






\section{模范畴之间的函子}
\label{functors-section-functors}

\noindent
下列引理是本章各结果的典型范例。

\begin{lemma}
\label{functors-lemma-functor}
设 $A$ 和 $B$ 为环，并设 $F : \text{Mod}_A \to \text{Mod}_B$ 为函子。
下列条件等价：
\begin{enumerate}
\item $F$ 同构于函子 $M \mapsto M \otimes_A K$，
这里所用的是某个 $A \otimes_\mathbf{Z} B$-模 $K$；
\item $F$ 右正合且与所有直和交换；
\item $F$ 与所有余极限交换；
\item $F$ 有右伴随函子 $G$。
\end{enumerate}
\end{lemma}

\begin{proof}
若 (1) 成立，则 (4) 成立，因为 $M \mapsto M \otimes_A K$ 的右伴随是
$N \mapsto \Hom_B(K, N)$；参见 Differential Graded Algebra, Lemma
\ref{dga-lemma-tensor-hom-adjunction}。若 (4) 成立，则由 Categories,
Lemma \ref{categories-lemma-adjoint-exact}，(3) 成立。
(3) $\Rightarrow$ (2) 由定义立即得到。

\medskip\noindent
假设 (2)。我们证明 (1)。由 Homology, Section
\ref{homology-section-functors} 中的讨论，函子 $F$ 可加。因此 $F$ 诱导环同态
$A \to \text{End}_B(F(M))$, $a \mapsto F(a \cdot \text{id}_M)$，
对每个 $A$-模 $M$ 都如此。
从而 $F(M)$ 成为 $A \otimes_\mathbf{Z} B$-模，并且关于 $M$ 具有函子性。
令 $K = F(A)$。定义
$$
M \otimes_A K = M \otimes_A F(A) \longrightarrow F(M),
\quad m \otimes k \longmapsto F(\varphi_m)(k)
$$
这里 $\varphi_m : A \to M$ 的映射规则为 $a \to am$。规则
$(m, k) \mapsto F(\varphi_m)(k)$ 是 $A$-双线性的（并且关于右侧为
$B$-线性），因而如所需给出上示 $A \otimes_\mathbf{Z} B$-线性映射。
此构造关于 $M$ 具有函子性，故定义函子变换
$- \otimes_A K \to F(-)$；在 $A$ 上取值时它是同构。
对每个 $A$-模 $M$，可以选取正合序列
$$
\bigoplus\nolimits_{j \in J} A \to
\bigoplus\nolimits_{i \in I} A \to
M \to 0
$$
利用上面构造的映射，得到交换图
$$
\xymatrix{
(\bigoplus\nolimits_{j \in J} A) \otimes_A K \ar[r] \ar[d] &
(\bigoplus\nolimits_{i \in I} A) \otimes_A K \ar[r] \ar[d] &
M \otimes_A K \ar[r] \ar[d] &
0 \\
F(\bigoplus\nolimits_{j \in J} A) \ar[r] &
F(\bigoplus\nolimits_{i \in I} A) \ar[r] &
F(M) \ar[r] &  0
}
$$
下行因 $F$ 右正合而正合。上行因与 $K$ 的张量积右正合而正合。
由于 $F$ 与直和交换，左边两个竖直箭头都是双射。结论随即成立。
\end{proof}

\begin{example}
\label{functors-example-functor-modules}
设 $R$ 为环，$A$ 和 $B$ 为 $R$-代数，并设 $K$ 为
$A \otimes_R B$-模。此时可以考虑函子
\begin{equation}
\label{functors-equation-FM-modules}
F : \text{Mod}_A \longrightarrow \text{Mod}_B,\quad
M \longmapsto M \otimes_A K
\end{equation}
此函子是 $R$-线性的、右正合的，与任意直和及所有余极限交换，
并且有右伴随（引理 \ref{functors-lemma-functor}）。
\end{example}

\begin{lemma}
\label{functors-lemma-functor-modules}
设 $R$ 为环，$A$ 和 $B$ 为 $R$-代数。下列两个范畴等价：
\begin{enumerate}
\item 右正合且与任意直和交换的 $R$-线性函子
$F : \text{Mod}_A \to \text{Mod}_B$ 所成的范畴；
\item 范畴 $\text{Mod}_{A \otimes_R B}$。
\end{enumerate}
此等价把 $K$ 映到 (\ref{functors-equation-FM-modules}) 中的函子 $F$。
\end{lemma}

\begin{proof}
设 $F$ 为第一个范畴的对象。由引理 \ref{functors-lemma-functor}，可以假设
$F(M) = M \otimes_A K$ 关于 $M$ 函子性地成立，其中所用的是某个
$A \otimes_\mathbf{Z} B$-模 $K$。$R$-线性这一 $F$ 的性质立即说明，
$A \otimes_\mathbf{Z} B$-模结构在 $K$ 上来自唯一的
$A \otimes_R B$-模结构在 $K$ 上。因此，(\ref{functors-equation-FM-modules}) 中把
$K$ 映到 $F$ 的函子本质满。

\medskip\noindent
为证明该函子全忠实，必须证明：给定 $A \otimes_R B$-模 $K$ 和 $K'$，
对应函子之间的任意变换 $t : F \to F'$ 都来自唯一的
$\varphi : K \to K'$。由于 $K = F(A)$ 且 $K' = F'(A)$，可把
$\varphi$ 取为值 $t_A : F(A) \to F'(A)$，即 $t$ 在 $A$ 处的值。
此映射是 $A \otimes_R B$-线性的；这是由引理 \ref{functors-lemma-functor} 证明中
$A \otimes B$-模结构在 $F(A)$ 和 $F'(A)$ 上的定义得出的。
\end{proof}

\begin{remark}
\label{functors-remark-composition}
设 $R$ 为环，$A$、$B$、$C$ 为 $R$-代数。设
$F : \text{Mod}_A \to \text{Mod}_B$ 和
$F' : \text{Mod}_B \to \text{Mod}_C$ 是与任意直和交换的
$R$-线性右正合函子。若在引理 \ref{functors-lemma-functor-modules} 的等价下，
对象 $K$（位于 $\text{Mod}_{A \otimes_R B}$）对应于 $F$，而对象
$K'$（位于 $\text{Mod}_{B \otimes_R C}$）对应于 $F'$，则把
$K \otimes_B K'$ 视为 $\text{Mod}_{A \otimes_R C}$ 的对象时，
它对应于 $F' \circ F$。
\end{remark}

\begin{remark}
\label{functors-remark-exact-flat}
在引理 \ref{functors-lemma-functor-modules} 的情形下，假设 $F$ 对应于 $K$。则
$F$ 正合 $\Leftrightarrow$ $K$ 在 $A$ 上平坦。
\end{remark}

\begin{remark}
\label{functors-remark-finite}
在引理 \ref{functors-lemma-functor-modules} 的情形下，假设 $F$ 对应于 $K$。则
$F$ 把有限 $A$-模映到有限 $B$-模
$\Leftrightarrow$ $K$ 作为 $B$-模是有限的。
\end{remark}

\begin{remark}
\label{functors-remark-finite-presentation}
在引理 \ref{functors-lemma-functor-modules} 的情形下，假设 $F$ 对应于 $K$。则
$F$ 把有限表示 $A$-模映到有限表示 $B$-模
$\Leftrightarrow$ $K$ 作为 $B$-模是有限表示的。
\end{remark}

\begin{lemma}
\label{functors-lemma-functor-equivalence}
设 $A$ 和 $B$ 为环。若
$$
F : \text{Mod}_A \longrightarrow \text{Mod}_B
$$
是范畴等价，则存在环同构 $A \to B$ 及可逆 $B$-模 $L$，使得
$F$ 同构于函子 $M \mapsto (M \otimes_A B) \otimes_B L$。
\end{lemma}

\begin{proof}
由于等价与所有余极限交换，引理 \ref{functors-lemma-functor} 适用。令 $K$ 为
$A \otimes_\mathbf{Z} B$-模，使 $F$ 同构于函子
$M \mapsto M \otimes_A K$。令 $K'$ 为 $B \otimes_\mathbf{Z} A$-模，
使 $F$ 的一个拟逆同构于函子 $N \mapsto N \otimes_B K'$。
由注 \ref{functors-remark-composition} 和引理 \ref{functors-lemma-functor-modules}，有同构
$$
\psi : K \otimes_B K' \longrightarrow A
$$
它是 $A \otimes_\mathbf{Z} A$-模同构。同样，有同构
$$
\psi' : K' \otimes_A K \longrightarrow B
$$
它是 $B \otimes_\mathbf{Z} B$-模同构。
选取元素 $\xi = \sum_{i = 1, \ldots, n} x_i \otimes y_i \in K \otimes_B K'$，
使 $\psi(\xi) = 1$。考察同构
$$
K \xrightarrow{\psi^{-1} \otimes \text{id}_K}
K \otimes_B K' \otimes_A K \xrightarrow{\text{id}_K \otimes \psi'} K
$$
其复合是同构，并由下式给出：
$$
k \longmapsto \sum x_i \psi'(y_i \otimes k)
$$
由此可知，该自同构分解为
$$
K \to B^{\oplus n} \to K
$$
这是 $B$-模映射的分解。所以 $K$ 作为 $B$-模是有限投射的。

\medskip\noindent
我们断言 $K$ 作为 $B$-模可逆。这等价于要求 $K$ 作为 $B$-模的秩
恒为 $1$；参见 More on Algebra, Lemma
\ref{more-algebra-lemma-invertible} 和 Algebra, Lemma
\ref{algebra-lemma-finite-projective}。否则，存在极大理想
$\mathfrak m \subset B$，使得或者 (a) $K \otimes_B B/\mathfrak m = 0$，
或者 (b) 存在满射 $K \to (B/\mathfrak m)^{\oplus 2}$，它是
$B$-模满射。情形 (a) 不可能，因为 $K' \otimes_A K \otimes_B N = N$
对所有 $B$-模 $N$ 都成立。
情形 (b) 将给出满射
$$
A = K \otimes_B K' \longrightarrow (B/\mathfrak m \otimes_B K')^{\oplus 2}
$$
这是（右）$A$-模满射。
这是不可能的，因为目标是至少需要两个生成元的 $A$-模：
$B/\mathfrak m \otimes_B K'$ 非零，因为它是非零模
$B/\mathfrak m$ 在 $F$ 的拟逆下的像。

\medskip\noindent
由于 $K$ 作为 $B$-模可逆，有 $\Hom_B(K, K) = B$。又因为 $K = F(A)$，
$A$ 在 $K$ 上的作用定义环同构 $A \to B$。引理得证。
\end{proof}

\begin{lemma}
\label{functors-lemma-functor-equivalence-linear}
设 $R$ 为环，$A$ 和 $B$ 为 $R$-代数。若
$$
F : \text{Mod}_A \longrightarrow \text{Mod}_B
$$
是 $R$-线性范畴等价，则存在同构 $A \to B$（作为 $R$-代数）及可逆 $B$-模
$L$，使 $F$ 同构于函子 $M \mapsto (M \otimes_A B) \otimes_B L$。
\end{lemma}

\begin{proof}
由引理 \ref{functors-lemma-functor-equivalence} 得到 $A \to B$ 和 $L$。
为完成证明，只需说明 $R$-线性这一 $F$ 的性质迫使 $A \to B$ 为
$R$-代数态射。细节从略。
\end{proof}

\begin{remark}
\label{functors-remark-monoidal}
设 $A$ 和 $B$ 为环。用模的张量积赋予 $\text{Mod}_A$ 和
$\text{Mod}_B$ 通常的幺半结构。设
$F : \text{Mod}_A \to \text{Mod}_B$ 为幺半范畴之间的函子；参见
Categories, Definition \ref{categories-definition-functor-monoidal-categories}。
有如下说明：
\begin{enumerate}
\item 由我们的定义，$F(A)$ 是单位对象，故 $F(A) = B$。
\item 得到乘法映射 $\varphi : A \to B$：它把 $a \in A$ 映到其在
$F(A) = B$ 上的作用。
\item 取 $A = B$ 且 $F(M) = M \otimes_A M$。此时 $\varphi(a) = a^2$。
\item 若 $F$ 可加，则 $\varphi$ 是环态射。
\item 取 $A = B = \mathbf{Z}$ 且 $F(M) = M/\text{torsion}$。此时
$\varphi = \text{id}_\mathbf{Z}$，但 $F$ 不是恒等函子。
\item 若 $F$ 右正合且与直和交换，则由引理 \ref{functors-lemma-functor}，
$F(M) = M \otimes_{A, \varphi} B$。
\end{enumerate}
换言之，环态射 $A \to B$ 与那些同所有余极限交换的幺半范畴函子
$\text{Mod}_A \to \text{Mod}_B$ 的同构类之间有双射。
\end{remark}




\section{模范畴上函子的延拓}
\label{functors-section-functors-extend}

\noindent
对环 $A$，用 $\text{Mod}^{fp}_A$ 表示有限表示 $A$-模所成的范畴。

\begin{lemma}
\label{functors-lemma-functor-fp-modules}
设 $A$ 和 $B$ 为环，并设
$F : \text{Mod}^{fp}_A \to \text{Mod}^{fp}_B$ 为函子。
则 $F$ 唯一延拓为与滤过余极限交换的函子
$F' : \text{Mod}_A \to \text{Mod}_B$。
\end{lemma}

\begin{proof}
这是引理 \ref{functors-lemma-functor-on-fp-modules} 的特殊情形。
\end{proof}

\begin{remark}
\label{functors-remark-monoidal-extension}
令 $A$、$B$、$F$ 和 $F'$ 如引理 \ref{functors-lemma-functor-fp-modules} 中所述。
注意，两个有限表示模的张量积仍是有限表示的；参见 Algebra, Lemma
\ref{algebra-lemma-tensor-finiteness}。因此可用模的张量积赋予
$\text{Mod}^{fp}_A$、$\text{Mod}^{fp}_B$、$\text{Mod}_A$ 和
$\text{Mod}_B$ 通常的幺半结构。此时，若 $F$ 是幺半范畴之间的函子，
则 $F'$ 也是。这立即来自模的张量积与滤过余极限交换这一事实。
\end{remark}

\begin{lemma}
\label{functors-lemma-functor-fp-modules-exact}
令 $A$、$B$、$F$ 和 $F'$ 如引理 \ref{functors-lemma-functor-fp-modules} 中所述。
\begin{enumerate}
\item 若 $F$ 可加，则 $F'$ 可加并与任意直和交换；
\item 若 $F$ 右正合，则 $F'$ 右正合。
\end{enumerate}
\end{lemma}

\begin{proof}
这由引理 \ref{functors-lemma-functor-on-fp-modules-additive} 和
\ref{functors-lemma-functor-on-fp-modules-right-exact} 得出。
\end{proof}

\begin{remark}
\label{functors-remark-monoidal-extension-exact}
结合注 \ref{functors-remark-monoidal}、注 \ref{functors-remark-monoidal-extension} 和引理
\ref{functors-lemma-functor-fp-modules-exact}，得到如下结论。给定环 $A$ 和 $B$，
环态射 $A \to B$ 所成的集合与右正合幺半范畴函子
$\text{Mod}^{fp}_A \to \text{Mod}^{fp}_B$ 的同构类集合之间有双射。
\end{remark}

\begin{lemma}
\label{functors-lemma-functor-fp-modules-left-exact}
令 $A$、$B$、$F$ 和 $F'$ 如引理 \ref{functors-lemma-functor-fp-modules} 中所述。
假设 $A$ 为相干环
（Algebra, Definition \ref{algebra-definition-coherent}）。
若 $F$ 左正合，则 $F'$ 左正合。
\end{lemma}

\begin{proof}
这是引理 \ref{functors-lemma-functor-on-fp-modules-left-exact} 的特殊情形。
\end{proof}

\noindent
对环 $A$，用 $\text{Mod}^{fg}_A$ 表示有限生成 $A$-模（亦称有限
$A$-模）所成的范畴。

\begin{lemma}
\label{functors-lemma-functor-finite-modules}
设 $A$ 和 $B$ 为诺特环，并设
$F : \text{Mod}^{fg}_A \to \text{Mod}^{fg}_B$ 为函子。
则 $F$ 唯一延拓为与滤过余极限交换的函子
$F' : \text{Mod}_A \to \text{Mod}_B$。若 $F$ 可加，则 $F'$ 可加并与
任意直和交换。若 $F$ 正合、左正合或右正合，则 $F'$ 也具有相应性质。
\end{lemma}

\begin{proof}
参见引理 \ref{functors-lemma-functor-fp-modules-exact} 和
\ref{functors-lemma-functor-fp-modules-left-exact}。还要用到有限 $A$-模是有限表示
$A$-模这一事实；参见 Algebra, Lemma
\ref{algebra-lemma-Noetherian-finite-type-is-finite-presentation}；以及诺特环
是相干环这一事实；参见 Algebra, Lemma
\ref{algebra-lemma-Noetherian-coherent}。
\end{proof}









\section{拟凝聚模范畴之间的函子}
\label{functors-section-functor-quasi-coherent}

\noindent
本节简要研究拟凝聚模范畴之间的函子。

\begin{example}
\label{functors-example-functor-quasi-coherent}
设 $R$ 为环，$X$ 和 $Y$ 为 $R$ 上的概形，且 $X$ 拟紧且拟分离。
设 $\mathcal{K}$ 为拟凝聚 $\mathcal{O}_{X \times_R Y}$-模。于是可考虑函子
\begin{equation}
\label{functors-equation-FM-QCoh}
F : \QCoh(\mathcal{O}_X) \longrightarrow \QCoh(\mathcal{O}_Y),\quad
\mathcal{F} \longmapsto
\text{pr}_{2, *}(\text{pr}_1^*\mathcal{F}
\otimes_{\mathcal{O}_{X \times_R Y}} \mathcal{K})
\end{equation}
态射 $\text{pr}_2$ 拟紧且拟分离
（Schemes, Lemmas \ref{schemes-lemma-quasi-compact-preserved-base-change}
和 \ref{schemes-lemma-separated-permanence}）。故沿此态射的推前保持拟凝聚模；
参见 Schemes, Lemma \ref{schemes-lemma-push-forward-quasi-coherent}。
此外，该函子是 $R$-线性的，并与任意直和交换；参见 Cohomology of
Schemes, Lemma \ref{coherent-lemma-colimit-cohomology}。
\end{example}

\noindent
下述引理是引理 \ref{functors-lemma-functor-modules} 的自然推广。

\begin{lemma}
\label{functors-lemma-functor-quasi-coherent-from-affine}
设 $R$ 为环，$X$ 和 $Y$ 为 $R$ 上的概形，且 $X$ 仿射。下列两个范畴等价：
\begin{enumerate}
\item $R$-线性函子
$F : \QCoh(\mathcal{O}_X) \to \QCoh(\mathcal{O}_Y)$
中右正合且与任意直和交换者所成的范畴；
\item 范畴 $\QCoh(\mathcal{O}_{X \times_R Y})$。
\end{enumerate}
此等价把 $\mathcal{K}$ 送到 (\ref{functors-equation-FM-QCoh}) 中的函子 $F$。
\end{lemma}

\begin{proof}
设 $\mathcal{K}$ 为 $\QCoh(\mathcal{O}_{X \times_R Y})$ 的对象，且
$F_\mathcal{K}$ 为函子 (\ref{functors-equation-FM-QCoh})。由例
\ref{functors-example-functor-quasi-coherent} 中的讨论，已经知道 $F$ 是 $R$-线性的，
并与任意直和交换。由于 $\text{pr}_2 : X \times_R Y \to Y$ 是仿射的
（Morphisms, Lemma \ref{morphisms-lemma-base-change-affine}），函子
$\text{pr}_{2, *}$ 正合；参见 Cohomology of Schemes, Lemma
\ref{coherent-lemma-relative-affine-vanishing}。故 $F$ 也右正合；换言之，
$F$ 如 (1) 中所述。

\medskip\noindent
设 $F$ 如 (1) 中所述。不妨设 $X = \Spec(A)$。考虑拟凝聚
$\mathcal{O}_Y$-模 $\mathcal{G} = F(\mathcal{O}_X)$。函子 $F$ 诱导
$R$-线性映射
$A \to \text{End}_{\mathcal{O}_Y}(\mathcal{G})$,
$a \mapsto F(a \cdot \text{id})$。因此 $\mathcal{G}$ 是下述层上的模层：
$$
A \otimes_R \mathcal{O}_Y = \text{pr}_{2, *}\mathcal{O}_{X \times_R Y}
$$
由 Morphisms, Lemma \ref{morphisms-lemma-affine-equivalence-modules}，存在唯一的
拟凝聚模 $\mathcal{K}$ 在 $X \times_R Y$ 上，使得
$F(\mathcal{O}_X) = \mathcal{G} =
\text{pr}_{2, *}\mathcal{K}$，且与 $A$ 和 $\mathcal{O}_Y$ 的作用相容。以
$F_\mathcal{K}$ 表示由 (\ref{functors-equation-FM-QCoh}) 给出的函子。存在等价
$\text{Mod}_A \to \QCoh(\mathcal{O}_X)$，它把 $A$ 送到 $\mathcal{O}_X$；参见
Schemes, Lemma
\ref{schemes-lemma-equivalence-quasi-coherent}。因此，由引理
\ref{functors-lemma-functor-on-modules} 得到同构 $F \cong F_\mathcal{K}$，因为按构造
已有同构 $F(\mathcal{O}_X) \cong F_\mathcal{K}(\mathcal{O}_X)$，且它与
$A$-作用相容。

\medskip\noindent
这说明把 $\mathcal{K}$ 送到 $F_\mathcal{K}$ 的函子本质满。其全忠实性的验证从略。
\end{proof}

\begin{remark}
\label{functors-remark-affine-morphism}
下文将用到：对仿射态射 $h : T \to S$，有
$h_*\mathcal{G} \otimes_{\mathcal{O}_S} \mathcal{H} =
h_*(\mathcal{G} \otimes_{\mathcal{O}_T} h^*\mathcal{H})$，其中
$\mathcal{G} \in \QCoh(\mathcal{O}_T)$ 且
$\mathcal{H} \in \QCoh(\mathcal{O}_S)$。将其翻译成代数后立即可得。
\end{remark}

\begin{lemma}
\label{functors-lemma-functor-quasi-coherent-from-affine-compose}
在引理 \ref{functors-lemma-functor-quasi-coherent-from-affine} 中，设 $F$ 对应于
$\mathcal{K}$，后者属于 $\QCoh(\mathcal{O}_{X \times_R Y})$。则：
\begin{enumerate}
\item 若 $f : X' \to X$ 为仿射态射，则 $F \circ f_*$ 对应于
$(f \times \text{id}_Y)^*\mathcal{K}$；
\item 若 $g : Y' \to Y$ 为平坦态射，则 $g^* \circ F$ 对应于
$(\text{id}_X \times g)^*\mathcal{K}$；
\item 若 $j : V \to Y$ 为开浸入，则 $j^* \circ F$ 对应于
$\mathcal{K}|_{X \times_R V}$。
\end{enumerate}
\end{lemma}

\begin{proof}
(1) 的证明。考虑交换图
$$
\xymatrix{
X' \times_R Y \ar[rrd]^{\text{pr}'_2} \ar[rd]_{f \times \text{id}_Y}
\ar[dd]_{\text{pr}'_1} \\
& X \times_R Y \ar[r]_{\text{pr}_2} \ar[d]_{\text{pr}_1} & Y \\
X' \ar[r]^f & X
}
$$
设 $\mathcal{F}'$ 为 $X'$ 上的拟凝聚模。则
\begin{align*}
\text{pr}_{2, *}(\text{pr}_1^*f_*\mathcal{F}'
\otimes_{\mathcal{O}_{X \times_R Y}} \mathcal{K})
& =
\text{pr}_{2, *}((f \times \text{id}_Y)_*
(\text{pr}'_1)^*\mathcal{F}'
\otimes_{\mathcal{O}_{X \times_R Y}} \mathcal{K}) \\
& =
\text{pr}_{2, *}(f \times \text{id}_Y)_*
\left((\text{pr}'_1)^*\mathcal{F}'
\otimes_{\mathcal{O}_{X' \times_R Y}}
(f \times \text{id}_Y)^*\mathcal{K})\right)  \\
& =
\text{pr}'_{2, *}((\text{pr}'_1)^*\mathcal{F}'
\otimes_{\mathcal{O}_{X' \times_R Y}} (f \times \text{id}_Y)^*\mathcal{K})
\end{align*}
这里，第一个等式是图中左侧方块的仿射基变换；参见 Cohomology of Schemes,
Lemma \ref{coherent-lemma-affine-base-change}。第二个等式由注
\ref{functors-remark-affine-morphism} 得到。第三个等式是模的推前的函子性。这证明了 (1)。

\medskip\noindent
(2) 的证明。考虑交换图
$$
\xymatrix{
X \times_R Y' \ar[rr]_-{\text{pr}'_2} \ar[rd]^{\text{id}_X \times g}
\ar[rdd]_{\text{pr}'_1} & & Y' \ar[d]^g \\
& X \times_R Y \ar[r]_-{\text{pr}_2} \ar[d]^{\text{pr}_1} & Y \\
& X
}
$$
有
\begin{align*}
g^*\text{pr}_{2, *}(\text{pr}_1^*\mathcal{F}
\otimes_{\mathcal{O}_{X \times_R Y}} \mathcal{K})
& =
\text{pr}'_{2, *}(
(\text{id}_X \times g)^*(
\text{pr}_1^*\mathcal{F} \otimes_{\mathcal{O}_{X \times_R Y}} \mathcal{K})) \\
& =
\text{pr}'_{2, *}((\text{pr}'_1)^*\mathcal{F}
\otimes_{\mathcal{O}_{X \times_R Y'}}
(\text{id}_X \times g)^*\mathcal{K})
\end{align*}
第一个等式由图中方块的平坦基变换得到；参见 Cohomology of Schemes, Lemma
\ref{coherent-lemma-flat-base-change-cohomology}。第二个等式由拉回的函子性以及
张量积的拉回是各拉回的张量积这一事实得到。

\medskip\noindent
(3) 是 (2) 的特殊情形。
\end{proof}

\begin{lemma}
\label{functors-lemma-functor-quasi-coherent-from-affine-diagonal-pre}
设 $R$ 为环，$X$ 和 $Y$ 为 $R$ 上的概形。假设 $X$ 拟紧且对角态射仿射。设
$F : \QCoh(\mathcal{O}_X) \to \QCoh(\mathcal{O}_Y)$
为与任意直和交换的 $R$-线性右正合函子。则可构造：
\begin{enumerate}
\item 拟凝聚模 $\mathcal{K}$ 在 $X \times_R Y$ 上；
\item 自然变换 $t : F \to F_\mathcal{K}$，其中 $F_\mathcal{K}$ 表示函子
(\ref{functors-equation-FM-QCoh})。
\end{enumerate}
它们使得 $t : F \circ f_* \to F_\mathcal{K} \circ f_*$ 是同构，
其中 $f : X' \to X$ 为源是仿射概形的任意态射。
\end{lemma}

\begin{proof}
考虑态射 $f' : X' \to X$，其中 $X'$ 仿射。由于 $X$ 的对角态射仿射，
$f'$ 是仿射态射（Morphisms, Lemma
\ref{morphisms-lemma-affine-permanence}）。因此
$f'_* : \QCoh(\mathcal{O}_{X'}) \to \QCoh(\mathcal{O}_X)$ 是 $R$-线性正合函子
（Cohomology of Schemes, Lemma
\ref{coherent-lemma-relative-affine-vanishing}），并与直和交换
（Cohomology of Schemes, Lemma \ref{coherent-lemma-colimit-cohomology}）。故
$F \circ f'_*$ 是与任意直和交换的 $R$-线性右正合函子。于是由引理
\ref{functors-lemma-functor-quasi-coherent-from-affine}，有
$F \circ f'_* = F_{\mathcal{K}'}$，其中 $\mathcal{K}'$ 在
$X' \times_R Y$ 上。此外，给定态射 $f'' : X'' \to X'$，其中 $X''$ 仿射，
由已给出的引用与引理 \ref{functors-lemma-functor-quasi-coherent-from-affine-compose}
可得典范识别
$(f'' \times \text{id}_Y)^*\mathcal{K}' = \mathcal{K}''$
。再给定态射 $f''' : X''' \to X''$ 时，这些识别满足余圈条件；其写法留给读者。

\medskip\noindent
选取仿射开覆盖 $X = \bigcup_{i = 1, \ldots, n} U_i$。由于 $X$ 的对角态射仿射，
交 $U_{i_0 \ldots i_p} = U_{i_0} \cap \ldots \cap U_{i_p}$ 是仿射的。如上，
包含态射 $j_{i_0 \ldots i_p} : U_{i_0 \ldots i_p} \to X$ 是仿射的。以
$\mathcal{K}_{i_0 \ldots i_p}$ 表示 $U_{i_0 \ldots i_p} \times_R Y$ 上对应于
$F \circ j_{i_0 \ldots i_p *}$ 的上述拟凝聚模。由上可得识别
$$
\mathcal{K}_{i_0 \ldots i_p} =
\mathcal{K}_{i_0 \ldots \hat i_j \ldots i_p}|_{U_{i_0 \ldots i_p} \times_R Y}
$$
它们满足粘合所需的通常相容性。换言之，得到唯一的拟凝聚模
$\mathcal{K}$ 在 $X \times_R Y$ 上，其在 $U_{i_0 \ldots i_p} \times_R Y$
上的限制为 $\mathcal{K}_{i_0 \ldots i_p}$，并与所示识别相容。

\medskip\noindent
接着构造变换 $t$。给定拟凝聚 $\mathcal{O}_X$-模 $\mathcal{F}$，以
$\mathcal{F}_{i_0 \ldots i_p}$ 表示 $\mathcal{F}$ 在
$U_{i_0 \ldots i_p}$ 上的限制，并以
$(\text{pr}_1^*\mathcal{F} \otimes \mathcal{K})_{i_0 \ldots i_p}$ 表示
$\text{pr}_1^*\mathcal{F} \otimes \mathcal{K}$ 在
$U_{i_0 \ldots i_p} \times_R Y$ 上的限制。注意
\begin{align*}
F(j_{i_0 \ldots i_p *}\mathcal{F}_{i_0 \ldots i_p})
& =
\text{pr}_{i_0 \ldots i_p, 2, *}(
\text{pr}_{i_0 \ldots i_p, 1}^*\mathcal{F}_{i_0 \ldots i_p}
\otimes \mathcal{K}_{i_0 \ldots i_p}) \\
& =
\text{pr}_{i_0 \ldots i_p, 2, *}
(\text{pr}_1^*\mathcal{F} \otimes \mathcal{K})_{i_0 \ldots i_p}
\end{align*}
其中 $\text{pr}_{i_0 \ldots i_p, 2} : U_{i_0 \ldots i_p} \times_R Y \to Y$
是投影，另一投影亦类似。此外，这些识别与上一段所示识别相容。回顾
Cohomology of Schemes, Lemma
\ref{coherent-lemma-separated-case-relative-cech}：相对 {\v C}ech 复形
$$
\bigoplus 
\text{pr}_{i_0, 2, *}
(\text{pr}_1^*\mathcal{F} \otimes \mathcal{K})_{i_0}
\to
\bigoplus 
\text{pr}_{i_0i_1, 2, *}
(\text{pr}_1^*\mathcal{F} \otimes \mathcal{K})_{i_0i_1}
\to
\bigoplus 
\text{pr}_{i_0i_1i_2, 2, *}
(\text{pr}_1^*\mathcal{F} \otimes \mathcal{K})_{i_0i_1i_2}
\to \ldots
$$
计算 $R\text{pr}_{2, *}(\text{pr}_1^*\mathcal{F} \otimes \mathcal{K})$。因此其
$0$ 次上同调层是 $F_\mathcal{K}(\mathcal{F})$。于是由考察下列交换图得到所需映射
$t : F(\mathcal{F}) \to F_\mathcal{K}(\mathcal{F})$
$$
\xymatrix{
&
F(\mathcal{F}) \ar[r] \ar@{..>}[d] &
\bigoplus F(j_{i_0*}\mathcal{F}_{i_0}) \ar[r] \ar[d] &
\bigoplus F(j_{i_0i_1*}\mathcal{F}_{i_0i_1}) \ar[d] \\
0 \ar[r] &
F_\mathcal{K}(\mathcal{F}) \ar[r] &
\bigoplus 
\text{pr}_{i_0, 2, *}
(\text{pr}_1^*\mathcal{F} \otimes \mathcal{K})_{i_0}
\ar[r] &
\bigoplus 
\text{pr}_{i_0i_1, 2, *}
(\text{pr}_1^*\mathcal{F} \otimes \mathcal{K})_{i_0i_1}
}
$$
上行由把 $F$ 作用于（正合）复形
$0 \to \mathcal{F} \to \bigoplus j_{i_0*}\mathcal{F}_{i_0} \to
\bigoplus j_{i_0i_1*}\mathcal{F}_{i_0i_1}$ 得到（但由于 $F$ 不正合，上行只是
复形而不一定正合）。实线竖直箭头是上述识别。这确实按要求定义了虚线箭头。
该箭头关于 $\mathcal{F}$ 函子性地变化；细节从略。

\medskip\noindent
还需证明最后一个断言。设 $f : X' \to X$ 如引理陈述所述，并设
$\mathcal{K}'$ 为证明第一段中在 $X' \times_R Y$ 上构造的拟凝聚模。若态射
$f : X' \to X$ 的像落在某个开集 $U_i$ 中，则结果由引理
\ref{functors-lemma-functor-quasi-coherent-from-affine-compose} 得出，因为此时已知
$\mathcal{K}_i = \mathcal{K}|_{U_i \times_R Y}$ 拉回为 $\mathcal{K}$。一般地，
得到仿射开覆盖 $X' = \bigcup U'_i$，其中 $U'_i = f^{-1}(U_i)$，以及同构
$\mathcal{K}'|_{U'_i} = f_i^*\mathcal{K}_i$，其中
$f_i : U'_i \to U_i$ 是诱导态射。这些态射满足粘合为同构
$\mathcal{K}' = f^*\mathcal{K}$ 所需的相容条件，结论得证。部分细节从略。
\end{proof}

\begin{lemma}
\label{functors-lemma-coh-noetherian-from-affine-flat}
在引理 \ref{functors-lemma-functor-quasi-coherent-from-affine} 或引理
\ref{functors-lemma-functor-quasi-coherent-from-affine-diagonal-pre} 中，若 $F$ 是正合函子，
则对应对象 $\mathcal{K}$ 属于 $\QCoh(\mathcal{O}_{X \times_R Y})$，并在 $X$ 上平坦。
\end{lemma}

\begin{proof}
可假设 $X$ 仿射，于是处于引理
\ref{functors-lemma-functor-quasi-coherent-from-affine} 的情形。由引理
\ref{functors-lemma-functor-quasi-coherent-from-affine-compose}，可假设 $Y$ 仿射。在仿射
情形，该断言翻译为注 \ref{functors-remark-exact-flat}。
\end{proof}

\begin{lemma}
\label{functors-lemma-functor-quasi-coherent-from-affine-diagonal}
设 $R$ 为环，$X$ 和 $Y$ 为 $R$ 上的概形。假设 $X$ 拟紧且对角态射仿射。
下列两个范畴等价：
\begin{enumerate}
\item $R$-线性正合函子
$F : \QCoh(\mathcal{O}_X) \to \QCoh(\mathcal{O}_Y)$
中与任意直和交换者所成的范畴；
\item $\QCoh(\mathcal{O}_{X \times_R Y})$ 的全子范畴，其对象为满足下列条件的
$\mathcal{K}$：
\begin{enumerate}
\item $\mathcal{K}$ 在 $X$ 上平坦；
\item 对 $\mathcal{F} \in \QCoh(\mathcal{O}_X)$，有
$R^q\text{pr}_{2, *}(\text{pr}_1^*\mathcal{F}
\otimes_{\mathcal{O}_{X \times_R Y}} \mathcal{K}) = 0$，只要 $q > 0$。
\end{enumerate}
\end{enumerate}
此等价把 $\mathcal{K}$ 送到 (\ref{functors-equation-FM-QCoh}) 中的函子 $F$。
\end{lemma}

\begin{proof}
设 $\mathcal{K}$ 如 (2) 中所述。(\ref{functors-equation-FM-QCoh}) 中的函子 $F$ 与直和
交换。由于由 (1)(a) 模 $\mathcal{K}$ 在 $X$ 上平坦，给定短正合列
$0 \to \mathcal{F}_1 \to \mathcal{F}_2 \to \mathcal{F}_3 \to 0$，得到短正合列
$$
0 \to
\text{pr}_1^*\mathcal{F}_1 \otimes_{\mathcal{O}_{X \times_R Y}} \mathcal{K} \to
\text{pr}_1^*\mathcal{F}_2 \otimes_{\mathcal{O}_{X \times_R Y}} \mathcal{K} \to
\text{pr}_1^*\mathcal{F}_3 \otimes_{\mathcal{O}_{X \times_R Y}} \mathcal{K} \to
0
$$
由于由 (2)(b)，第一项的高阶直接像 $R^1\text{pr}_{2, *}$ 为零，遂得
$0 \to F(\mathcal{F}_1) \to F(\mathcal{F}_2) \to F(\mathcal{F}_3) \to 0$
正合，故 $F$ 如 (1) 中所述。

\medskip\noindent
设 $F$ 如 (1) 中所述，并令 $\mathcal{K}$ 和 $t : F \to F_\mathcal{K}$ 如引理
\ref{functors-lemma-functor-quasi-coherent-from-affine-diagonal-pre} 中所述。由引理
\ref{functors-lemma-coh-noetherian-from-affine-flat}，$\mathcal{K}$ 在 $X$ 上平坦。为完成
证明，须说明 $t$ 是同构并证明有关高阶直接像的断言。两者都来自相对
{\v C}ech 复形
$$
\bigoplus 
\text{pr}_{i_0, 2, *}
(\text{pr}_1^*\mathcal{F} \otimes \mathcal{K})_{i_0}
\to
\bigoplus 
\text{pr}_{i_0i_1, 2, *}
(\text{pr}_1^*\mathcal{F} \otimes \mathcal{K})_{i_0i_1}
\to
\bigoplus 
\text{pr}_{i_0i_1i_2, 2, *}
(\text{pr}_1^*\mathcal{F} \otimes \mathcal{K})_{i_0i_1i_2}
\to \ldots
$$
计算 $R\text{pr}_{2, *}(\text{pr}_1^*\mathcal{F} \otimes \mathcal{K})$ 这一事实。
关于记号及其原因，参见引理
\ref{functors-lemma-functor-quasi-coherent-from-affine-diagonal-pre} 的证明。在引理
\ref{functors-lemma-functor-quasi-coherent-from-affine-diagonal-pre} 的证明中还已看到，
此复形等于 $F$ 作用于复形
$$
\bigoplus j_{i_0*}\mathcal{F}_{i_0} \to
\bigoplus j_{i_0i_1*}\mathcal{F}_{i_0i_1} \to
\bigoplus j_{i_0i_1i_2*}\mathcal{F}_{i_0i_1i_2} \to \ldots
$$
此复形除零次外正合，且零次上同调层等于 $\mathcal{F}$。因此，由于 $F$ 是
正合函子，可得 $F = F_\mathcal{K}$，并且 (2)(b) 成立。

\medskip\noindent
从略如下证明：把 $F$ 送到 $\mathcal{K}$ 的构造具有函子性，并且是把
$\mathcal{K}$ 送到由 (\ref{functors-equation-FM-QCoh}) 确定的函子 $F_\mathcal{K}$ 的
函子的拟逆。
\end{proof}

\begin{remark}
\label{functors-remark-characterize-FM-QCoh}
设 $R$ 为环，$X$ 和 $Y$ 为 $R$ 上的概形。假设 $X$ 拟紧且对角态射仿射。
引理 \ref{functors-lemma-functor-quasi-coherent-from-affine-diagonal} 可推广如下：与
$X \times_R Y$ 上拟凝聚模相关联的函子 (\ref{functors-equation-FM-QCoh})，恰为满足
下列性质的函子
$F : \QCoh(\mathcal{O}_X) \to \QCoh(\mathcal{O}_Y)$
：
\begin{enumerate}
\item $F$ 是 $R$-线性的，并与任意直和交换；
\item $F \circ j_*$ 右正合，只要 $j : U \to X$ 是仿射开集的包含；
\item 序列 $0 \to F(\mathcal{F}) \to F(\mathcal{G}) \to F(\mathcal{H})$
正合，只要 $0 \to \mathcal{F} \to \mathcal{G} \to \mathcal{H} \to 0$ 是正合列，
且对所有 $x \in X$，茎上的序列
$0 \to \mathcal{F}_x \to \mathcal{G}_x \to \mathcal{H}_x \to 0$
是分裂短正合列。
\end{enumerate}
确切地说，这些假设足以构造引理
\ref{functors-lemma-functor-quasi-coherent-from-affine-diagonal-pre} 中那样的变换
$t : F \to F_\mathcal{K}$ 并证明它是同构。此外，函子
(\ref{functors-equation-FM-QCoh}) 的确具有性质 (1)、(2)、(3)。若以后需要此结论，
我们将在此仔细陈述并证明它。
\end{remark}

\begin{lemma}
\label{functors-lemma-compose-FM-QCoh}
设 $R$ 为环，$X$、$Y$、$Z$ 为 $R$ 上的概形。假设 $X$ 和 $Y$ 拟紧且对角态射仿射。设
$$
F : \QCoh(\mathcal{O}_X) \to \QCoh(\mathcal{O}_Y)
\quad\text{且}\quad
G : \QCoh(\mathcal{O}_Y) \to \QCoh(\mathcal{O}_Z)
$$
为与任意直和交换的 $R$-线性正合函子。令 $\mathcal{K}$ 属于
$\QCoh(\mathcal{O}_{X \times_R Y})$，$\mathcal{L}$ 属于
$\QCoh(\mathcal{O}_{Y \times_R Z})$，它们是相应的“核”；参见引理
\ref{functors-lemma-functor-quasi-coherent-from-affine-diagonal}。则 $G \circ F$ 对应于
$\text{pr}_{13, *}(\text{pr}_{12}^*\mathcal{K}
\otimes_{\mathcal{O}_{X \times_R Y \times_R Z}}
\text{pr}_{23}^*\mathcal{L})$，它属于 $\QCoh(\mathcal{O}_{X \times_R Z})$。
\end{lemma}

\begin{proof}
由于 $G \circ F : \QCoh(\mathcal{O}_X) \to \QCoh(\mathcal{O}_Z)$ 是 $R$-线性
正合函子且与任意直和交换，由引理
\ref{functors-lemma-functor-quasi-coherent-from-affine-diagonal}，存在 $\mathcal{M}$ 属于
$\QCoh(\mathcal{O}_{X \times_R Z})$，它对应于 $G \circ F$。另一方面，记
$\mathcal{E} = \text{pr}_{13, *}(\text{pr}_{12}^*\mathcal{K}
\otimes \text{pr}_{23}^*\mathcal{L})$。这里及证明以下部分均省略张量积的下标。
设 $U \subset X$ 和 $W \subset Z$ 为仿射开子概形。为证明引理，将构造同构
$$
\Gamma(U \times_R W, \mathcal{E})
\cong
\Gamma(U \times_R W, \mathcal{M})
$$
并使其与随 $U$ 和 $W$ 变化的限制映射相容。

\medskip\noindent
首先注意
$$
\Gamma(U \times_R W, \mathcal{E}) =
\Gamma(U \times_R Y \times_R W,
\text{pr}_{12}^*\mathcal{K} \otimes \text{pr}_{23}^*\mathcal{L})
$$
这是由构造得到的。因此只需证明 $\mathcal{M}$ 也有同样的性质。

\medskip\noindent
写 $U = \Spec(A)$，并以 $j : U \to X$ 表示包含态射。回顾引理
\ref{functors-lemma-functor-quasi-coherent-from-affine} 的证明中 $\mathcal{M}$ 的构造：
$$
\Gamma(U \times_R W, \mathcal{M}) =
\Gamma(W, G(F(j_*\mathcal{O}_U)))
$$
其中右端的 $A$-模作用由 $A$ 在 $\mathcal{O}_U$ 上的作用给出。$F$ 与
$\mathcal{K}$ 之间的对应告诉我们
$F(j_*\mathcal{O}_U) = b_*(a^*j_*\mathcal{O}_U \otimes \mathcal{K})$
其中 $a : X \times_R Y \to X$ 和 $b : X \times_R Y \to Y$ 是投影态射。由于
$j$ 是仿射态射，由 Cohomology of Schemes, Lemma
\ref{coherent-lemma-affine-base-change}，有
$a^*j_*\mathcal{O}_U = (j \times \text{id}_Y)_*\mathcal{O}_{U \times_R Y}$
。接着，例如由注 \ref{functors-remark-affine-morphism}，有
$(j \times \text{id}_Y)_*\mathcal{O}_{U \times_R Y} \otimes \mathcal{K} =
(j \times \text{id}_Y)_*\mathcal{K}|_{U \times_R Y}$。
汇总所得结果，得到
$$
F(j_*\mathcal{O}_U) =
(U \times_R Y \to Y)_*\mathcal{K}|_{U \times_R Y}
$$
它带有显然的 $A$-作用。（此公式隐含在引理
\ref{functors-lemma-functor-quasi-coherent-from-affine} 的证明中。）应用函子 $G$ 得
$$
G(F(j_*\mathcal{O}_U)) =
t_*(s^*((U \times_R Y \to Y)_*\mathcal{K}|_{U \times_R Y})
\otimes \mathcal{L})
$$
其中 $s : Y \times_R Z \to Y$ 和 $t : Y \times_R Z \to Z$ 为投影态射。再次
使用仿射基变换（Cohomology of Schemes, Lemma
\ref{coherent-lemma-affine-base-change}），但这次用于方块
$$
\xymatrix{
U \times_R Y \times_R Z \ar[r] \ar[d] & U \times_R Y \ar[d] \\
Y \times_R Z \ar[r] & Y
}
$$
得到
$$
s^*((U \times_R Y \to Y)_*\mathcal{K}|_{U \times_R Y}) =
(U \times_R Y \times_R Z \to Y \times_R Z)_*
\text{pr}_{12}^*\mathcal{K}|_{U \times_R Y \times_R Z}
$$
再次使用注 \ref{functors-remark-affine-morphism}，得到
\begin{align*}
(U \times_R Y \times_R Z \to Y \times_R Z)_*
\text{pr}_{12}^*\mathcal{K}|_{U \times_R Y \times_R Z}
\otimes \mathcal{L} \\
=
(U \times_R Y \times_R Z \to Y \times_R Z)_*
\left(\text{pr}_{12}^*\mathcal{K} \otimes
\text{pr}_{23}^*\mathcal{L}\right)|_{U \times_R Y \times_R Z}
\end{align*}
对此应用函子 $\Gamma(W, t_*(-)) = \Gamma(Y \times_R W, -)$，得到
\begin{align*}
\Gamma(U \times_R W, \mathcal{M})
& =
\Gamma(W, G(F(j_*\mathcal{O}_U))) \\
& =
\Gamma(Y \times_R W, (U \times_R Y \times_R Z \to Y \times_R Z)_*
(\text{pr}_{12}^*\mathcal{K} \otimes
\text{pr}_{23}^*\mathcal{L})|_{U \times_R Y \times_R Z}) \\
& =
\Gamma(U \times_R Y \times_R W,
\text{pr}_{12}^*\mathcal{K} \otimes \text{pr}_{23}^*\mathcal{L})
\end{align*}
如所需。关于这些同构与限制映射相容的验证从略。
\end{proof}

\begin{lemma}
\label{functors-lemma-persistence-exactness}
令 $R$、$X$、$Y$ 和 $\mathcal{K}$ 如引理
\ref{functors-lemma-functor-quasi-coherent-from-affine-diagonal} 的 (2) 中所述。则对任意
概形 $T$ 在 $R$ 上，有
$$
R^q\text{pr}_{13, *}(\text{pr}_{12}^*\mathcal{F}
\otimes_{\mathcal{O}_{T \times_R X \times_R Y}}
\text{pr}_{23}^*\mathcal{K}) = 0
$$
其中 $\mathcal{F}$ 为 $T \times_R X$ 上的拟凝聚模，且 $q > 0$。
\end{lemma}

\begin{proof}
问题在 $T$ 上是局部的，故可假设 $T$ 仿射。此时可考虑图
$$
\xymatrix{
T \times_R X \ar[d] &
T \times_R X \times_R Y \ar[d] \ar[l] \ar[r] &
T \times_R Y \ar[d] \\
X &
X \times_R Y \ar[l] \ar[r] &
Y
}
$$
其竖直箭头均为仿射态射。特别地，沿 $T \times_R Y \to Y$ 的推前是忠实正合的
（Cohomology of Schemes, Lemma
\ref{coherent-lemma-relative-affine-vanishing} 和 Morphisms, Lemma
\ref{morphisms-lemma-affine-equivalence-modules}）。利用沿仿射态射的高阶直接像
消失这一事实（见上述引用）追踪此图，可知只需证明
$$
R^q\text{pr}_{2, *}(
\text{pr}_{23, *}(\text{pr}_{12}^*\mathcal{F}
\otimes_{\mathcal{O}_{T \times_R X \times_R Y}}
\text{pr}_{23}^*\mathcal{K})) =
R^q\text{pr}_{2, *}(
\text{pr}_{23, *}(\text{pr}_{12}^*\mathcal{F})
\otimes_{\mathcal{O}_{X \times_R Y}}
\mathcal{K}))
$$
为零；这由对 $\mathcal{K}$ 的假设成立。等号由注
\ref{functors-remark-affine-morphism} 得出。
\end{proof}

\begin{lemma}
\label{functors-lemma-functor-quasi-coherent-from-separated}
在引理 \ref{functors-lemma-functor-quasi-coherent-from-affine-diagonal} 中，令 $F$ 和
$\mathcal{K}$ 相互对应。若 $X$ 分离且在 $R$ 上平坦，则存在满射
$\mathcal{O}_X \boxtimes F(\mathcal{O}_X) \to \mathcal{K}$.
\end{lemma}

\begin{proof}
设 $\Delta : X \to X \times_R X$ 为对角态射，并置
$\mathcal{O}_\Delta = \Delta_*\mathcal{O}_X$。由于 $\Delta$ 是闭浸入，有短正合列
$$
0 \to \mathcal{I} \to 
\mathcal{O}_{X \times_R X} \to \mathcal{O}_\Delta \to 0
$$
由于 $\mathcal{K}$ 在 $X$ 上平坦，其拉回 $\text{pr}_{23}^*\mathcal{K}$ 到
$X \times_R X \times_R Y$ 后，在 $X \times_R X$ 上平坦。因此得到短正合列
$$
0 \to 
\text{pr}_{12}^*\mathcal{I}
\otimes
\text{pr}_{23}^*\mathcal{K} \to
\text{pr}_{23}^*\mathcal{K} \to
\text{pr}_{12}^*\mathcal{O}_\Delta
\otimes
\text{pr}_{23}^*\mathcal{K} \to 0
$$
在 $X \times_R X \times_R Y$ 上；参见 Modules, Lemma
\ref{modules-lemma-pullback-tensor-flat-module}。于是由引理
\ref{functors-lemma-persistence-exactness} 得到满射
$$
\text{pr}_{13, *}(\text{pr}_{23}^*\mathcal{K})
\to
\text{pr}_{13, *}(
\text{pr}_{12}^*\mathcal{O}_\Delta
\otimes
\text{pr}_{23}^*\mathcal{K})
$$
由平坦基变换（Cohomology of Schemes, Lemma
\ref{coherent-lemma-flat-base-change-cohomology}），此箭头的源等于
$\text{pr}_2^*\text{pr}_{2, *}\mathcal{K}
= \mathcal{O}_X \boxtimes F(\mathcal{O}_X)$。另一方面，其目标等于
$$
\text{pr}_{13, *}(
\text{pr}_{12}^*\mathcal{O}_\Delta
\otimes
\text{pr}_{23}^*\mathcal{K}) =
\text{pr}_{13, *} (\Delta \times \text{id}_Y)_* \mathcal{K} =
\mathcal{K}
$$
这就完成了证明。例如，第一个等式由 Cohomology, Lemma
\ref{cohomology-lemma-projection-formula-closed-immersion} 以及事实
$\text{pr}_{12}^*\mathcal{O}_\Delta =
(\Delta \times \text{id}_Y)_*\mathcal{O}_{X \times_R Y}$ 得出。
\end{proof}










\section{Gabriel--Rosenberg 重构}
\label{functors-section-gabriel}

\noindent
本节标题所指的是命题 \ref{functors-proposition-gabriel-rosenberg} 这一类结果。除 Gabriel
的原始论文 \cite{Gabriel} 外，也请参见 \cite{Brandenburg}；其中证明了拟分离
概形的相应结果，并讨论了相关文献。本节仅证明拟紧拟分离概形的
Gabriel--Rosenberg 重构。

\begin{lemma}
\label{functors-lemma-categorically-compact-QCoh}
设 $X$ 为拟紧拟分离概形，$\mathcal{F}$ 为拟凝聚 $\mathcal{O}_X$-模。则
$\mathcal{F}$ 是 $\QCoh(\mathcal{O}_X)$ 的范畴紧对象，当且仅当
$\mathcal{F}$ 是有限表示的。
\end{lemma}

\begin{proof}
关于范畴中范畴紧对象的概念，参见 Categories, Definition
\ref{categories-definition-compact-object}。若 $\mathcal{F}$ 是有限表示的，则由
Modules, Lemma \ref{modules-lemma-finite-presentation-quasi-compact-colimit}，
它是范畴紧的。反过来，任意拟凝聚模 $\mathcal{F}$ 都可写为滤过余极限
$\mathcal{F} = \colim \mathcal{F}_i$，其中各项是有限表示（因而拟凝聚）
$\mathcal{O}_X$-模；参见
Properties, Lemma \ref{properties-lemma-directed-colimit-finite-presentation}。
若 $\mathcal{F}$ 范畴紧，则存在某个 $i$ 及态射
$\mathcal{F} \to \mathcal{F}_i$，它是给定态射
$\mathcal{F}_i \to \mathcal{F}$ 的右逆。因此 $\mathcal{F}$ 是某有限表示模的
直和因子，故其自身也有限表示。
\end{proof}

\begin{lemma}
\label{functors-lemma-supported-on-support-pre}
设 $X$ 为仿射概形，$\mathcal{F}$ 为有限表示 $\mathcal{O}_X$-模，且
$\mathcal{E}$ 为非零拟凝聚 $\mathcal{O}_X$-模。若
$\text{Supp}(\mathcal{E}) \subset \text{Supp}(\mathcal{F})$,
则存在非零映射 $\mathcal{F} \to \mathcal{E}$。
\end{lemma}

\begin{proof}
把该断言翻译成代数。设 $A$ 为环，$M$ 为有限表示 $A$-模，$N$ 为非零
$A$-模。假设 $\text{Supp}(N) \subset \text{Supp}(M)$。要证
$\Hom_A(M, N)$ 非零。可假设 $N = A/I$ 是循环模（用 $N$ 的任意非零循环子模
代替它）。选取表示
$$
A^{\oplus m} \xrightarrow{T} A^{\oplus n} \to M \to 0
$$
回顾 $\text{Supp}(M)$ 由 $\text{Fit}_0(M)$ 截出；后者是由 $n \times n$ 阶
子式生成的理想，这里矩阵为 $T$。参见 More on Algebra, Lemma
\ref{more-algebra-lemma-fitting-ideal-basics}。假设
$\text{Supp}(N) \subset \text{Supp}(M)$ 现在意味着 $\text{Fit}_0(M)$ 的元素
在 $A/I$ 中幂零。考虑正合列
$$
0 \to \Hom_A(M, A/I) \to (A/I)^{\oplus n} \xrightarrow{T^t} (A/I)^{\oplus m}
$$
须证明 $T^t$ 不可能单射；建议读者利用 $\text{Fit}_0(M)$ 的元素在 $A/I$ 中
幂零这一事实自行寻找证明。下面给出我们的证明。由于 $\text{Fit}_0(M)$ 有限
生成，幂零性意味着零化子 $J \subset A/I$ 非零；它是 $\text{Fit}_0(M)$ 在
$A/I$ 中的零化子。为证明 $T^t$ 不单射，可在一个素理想处局部化。选取适当
素理想后，可假设 $A$ 是局部环且 $J$ 仍非零。于是由 More on Algebra, Lemma
\ref{more-algebra-lemma-exact-length-1}，$T^t$ 有非零核。
\end{proof}

\begin{lemma}
\label{functors-lemma-supported-on-support}
设 $X$ 为拟紧拟分离概形，$\mathcal{F}$ 为有限表示 $\mathcal{O}_X$-模。
$\QCoh(\mathcal{O}_X)$ 的下列两个子范畴相等：
\begin{enumerate}
\item 全子范畴 $\mathcal{A} \subset \QCoh(\mathcal{O}_X)$，其对象为支集在集合
意义下包含于 $\text{Supp}(\mathcal{F})$ 的拟凝聚模；
\item 最小 Serre 子范畴 $\mathcal{B} \subset \QCoh(\mathcal{O}_X)$，它包含
$\mathcal{F}$，并对扩张和任意直和封闭。
\end{enumerate}
\end{lemma}

\begin{proof}
注意，该断言有意义，因为有限表示 $\mathcal{O}_X$-模是拟凝聚的。由于
$\mathcal{A}$ 是对扩张和直和封闭的 Serre 子范畴，且 $\mathcal{F}$ 是
$\mathcal{A}$ 的对象，故 $\mathcal{B} \subset \mathcal{A}$。只需再证明
$\mathcal{A}$ 包含于 $\mathcal{B}$。

\medskip\noindent
设 $\mathcal{E}$ 为 $\mathcal{A}$ 的对象。存在极大子模
$\mathcal{E}' \subset \mathcal{E}$，它属于 $\mathcal{B}$。确切地说，设
$\mathcal{E}_i \subset \mathcal{E}$、$i \in I$ 是所有属于 $\mathcal{B}$ 的
子对象。则 $\bigoplus \mathcal{E}_i$ 属于 $\mathcal{B}$，故
$$
\mathcal{E}' = \Im(\bigoplus \mathcal{E}_i \longrightarrow \mathcal{E})
$$
显然这正是所求极大子模。

\medskip\noindent
现在假设有非零映射
$\mathcal{G} \to \mathcal{E}/\mathcal{E}'$
，其中 $\mathcal{G}$ 属于 $\mathcal{B}$。则
$\mathcal{G}' = \mathcal{E} \times_{\mathcal{E}/\mathcal{E}'} \mathcal{G}$
属于 $\mathcal{B}$，因为它是 $\mathcal{E}'$ 与 $\mathcal{G}$ 的扩张。于是
$\mathcal{G}' \to \mathcal{E}$ 的像会严格大于 $\mathcal{E}'$，与
$\mathcal{E}'$ 的极大性矛盾。因此只需证明下一段的断言。

\medskip\noindent
设 $\mathcal{E}$ 为 $\mathcal{A}$ 的非零对象。断言存在非零映射
$\mathcal{G} \to \mathcal{E}$，其中 $\mathcal{G}$ 属于 $\mathcal{B}$。对此按
最小个数 $n$ 归纳；这里 $U_i$ 是 $X$ 的仿射开集，且满足
$\text{Supp}(\mathcal{E}) \subset U_1 \cup \ldots \cup U_n$.
置 $U = U_n$，并以 $j : U \to X$ 表示包含态射。记
$\mathcal{E}' = \Im(\mathcal{E} \to j_*\mathcal{E}|_U)$。其核
$\mathcal{E}''$（这里所指的满射为 $\mathcal{E} \to \mathcal{E}'$）的支集包含于
$U_1 \cup \ldots \cup U_{n - 1}$。若 $\mathcal{E}''$ 非零，结论即得。因此可假设
$\mathcal{E} \subset j_*\mathcal{E}|_U$。特别地，$\mathcal{E}|_U$ 非零。由引理
\ref{functors-lemma-supported-on-support-pre}，存在非零映射
$\mathcal{F}|_U \to \mathcal{E}|_U$。它对应于映射
$$
\varphi : \mathcal{F} \longrightarrow j_*(\mathcal{E}|_U)
$$
其在 $U$ 上的限制非零。置 $\mathcal{G} = \varphi^{-1}(\mathcal{E})$ 即得结论。
\end{proof}

\begin{lemma}
\label{functors-lemma-quotient-supported-on-closed}
设 $X$ 为拟紧拟分离概形，$Z \subset X$ 为闭子集，并且
$U = X \setminus Z$ 拟紧。设 $\mathcal{A} \subset \QCoh(\mathcal{O}_X)$ 为全子
范畴，其对象是在 $Z$ 上有支集的拟凝聚模。则限制函子
$\QCoh(\mathcal{O}_X) \to \QCoh(\mathcal{O}_U)$ 诱导等价
$\QCoh(\mathcal{O}_X)/\mathcal{A} \cong \QCoh(\mathcal{O}_U)$。
\end{lemma}

\begin{proof}
由商构造的泛性质（Homology, Lemma
\ref{homology-lemma-serre-subcategory-is-kernel}），当然得到诱导函子
$\QCoh(\mathcal{O}_X)/\mathcal{A} \cong \QCoh(\mathcal{O}_U)$。以
$j : U \to X$ 表示包含态射。由于 $j$ 拟紧且拟分离，得到函子
$j_* : \QCoh(\mathcal{O}_U) \to \QCoh(\mathcal{O}_X)$。读者可验证它定义了拟逆；
细节从略。
\end{proof}

\begin{lemma}
\label{functors-lemma-characterize-affine}
设 $X$ 为拟紧拟分离概形。若 $\QCoh(\mathcal{O}_X)$ 等价于某个环上的模范畴，
则 $X$ 仿射。
\end{lemma}

\begin{proof}
设 $F : \text{Mod}_R \to \QCoh(\mathcal{O}_X)$ 为等价。则
$\mathcal{F} = F(R)$ 具有下列性质：
\begin{enumerate}
\item 它是有限表示 $\mathcal{O}_X$-模
（引理 \ref{functors-lemma-categorically-compact-QCoh}）；
\item $\Hom_X(\mathcal{F}, -)$ 正合；
\item $\Hom_X(\mathcal{F}, \mathcal{F})$ 是交换环；
\item $\QCoh(\mathcal{O}_X)$ 的每个对象都是若干份 $\mathcal{F}$ 的直和的商。
\end{enumerate}
设 $x \in X$ 为闭点。考虑满射
$$
\mathcal{O}_X \to i_*\kappa(x)
$$
其目标是 $\kappa(x)$ 沿包含态射 $i : x \to X$ 的推前。有
$$
\Hom_X(\mathcal{F}, i_*\kappa(x)) =
\Hom_{\mathcal{O}_{X, x}}(\mathcal{F}_x, \kappa(x))
$$
这首先由 (4) 蕴含 $\mathcal{F}_x$ 非零。由 (2)，每个映射
$\mathcal{F}_x \to \kappa(x)$ 都提升为映射
$\mathcal{F}_x \to \mathcal{O}_{X, x}$（事实上它甚至提升为整体映射
$\mathcal{F} \to \mathcal{O}_X$）。由于 $\mathcal{F}_x$ 是有限
$\mathcal{O}_{X, x}$-模，这蕴含 $\mathcal{F}_x$ 是非零有限自由
$\mathcal{O}_{X, x}$-模。又因 $\mathcal{F}$ 有限表示，故 $\mathcal{F}$ 在 $x$ 的某个开邻域
上是正秩有限自由的（Modules, Lemma
\ref{modules-lemma-finite-presentation-stalk-free}）。由于 $X$ 的每个闭子集都含有
闭点（Topology, Lemma \ref{topology-lemma-quasi-compact-closed-point}），这蕴含
$\mathcal{F}$ 是正秩有限局部自由的。类似地，映射
$$
\Hom_X(\mathcal{F}, \mathcal{F}) \to
\Hom_X(\mathcal{F}, i_*i^*\mathcal{F}) =
\Hom_{\kappa(x)}(\mathcal{F}_x/\mathfrak m_x \mathcal{F}_x,
\mathcal{F}_x/\mathfrak m_x \mathcal{F}_x)
$$
满射。由性质 (3)，可知秩 $\mathcal{F}_x$ 必为 $1$。因此 $\mathcal{F}$ 是可逆
$\mathcal{O}_X$-模。于是可知函子
$$
\mathcal{H}
\longmapsto
\Gamma(X, \mathcal{H}) = \Hom_X(\mathcal{O}_X, \mathcal{H}) =
\Hom_X(\mathcal{F}, \mathcal{H} \otimes_{\mathcal{O}_X} \mathcal{F})
$$
在 $\QCoh(\mathcal{O}_X)$ 上也正合。这蕴含第一 $\Ext$ 群
$$
\Ext^1_{\QCoh(\mathcal{O}_X)}(\mathcal{O}_X, \mathcal{H}) = 0
$$
在阿贝尔范畴 $\QCoh(\mathcal{O}_X)$ 中计算时，对所有 $\mathcal{H}$（属于
$\QCoh(\mathcal{O}_X)$）都消失。不过，由于
$\QCoh(\mathcal{O}_X) \subset \textit{Mod}(\mathcal{O}_X)$ 对扩张封闭
（Schemes, Section \ref{schemes-section-quasi-coherent}），拟凝聚模之间的
$\Ext^1$ 在 $\QCoh(\mathcal{O}_X)$ 中的计算与在
$\textit{Mod}(\mathcal{O}_X)$ 中的计算相同。因此
$$
H^1(X, \mathcal{H}) =
\Ext^1_{\textit{Mod}(\mathcal{O}_X)}(\mathcal{O}_X, \mathcal{H}) = 0
$$
对所有 $\mathcal{H}$（属于 $\QCoh(\mathcal{O}_X)$）都成立。例如由 Cohomology
of Schemes, Lemma \ref{coherent-lemma-quasi-compact-h1-zero-covering}，这蕴含
$X$ 仿射。
\end{proof}

\begin{proposition}
\label{functors-proposition-gabriel-rosenberg}
\begin{reference}
\cite[Theorem 1.2]{Brandenburg} 的特殊情形
\end{reference}
设 $X$ 和 $Y$ 为拟紧拟分离概形。若
$F : \QCoh(\mathcal{O}_X) \to \QCoh(\mathcal{O}_Y)$ 为等价，则存在概形同构
$f : Y \to X$ 以及可逆 $\mathcal{O}_Y$-模 $\mathcal{L}$，使得
$F(\mathcal{F}) = f^*\mathcal{F} \otimes \mathcal{L}$。
\end{proposition}

\begin{proof}
当然，$F$ 可加、正合，与所有极限、所有余极限、直和等交换。设
$U \subset X$ 为仿射开子概形。设 $\mathcal{I} \subset \mathcal{O}_X$ 为有限型
拟凝聚理想层，使得 $Z = V(\mathcal{I})$ 是 $U$ 在 $X$ 中的补集；参见
Properties, Lemma \ref{properties-lemma-quasi-coherent-finite-type-ideals}。
则 $\mathcal{O}_X/\mathcal{I}$ 是有限表示 $\mathcal{O}_X$-模。因此由引理
\ref{functors-lemma-categorically-compact-QCoh}，$\mathcal{G} = F(\mathcal{O}_X/\mathcal{I})$
是有限表示 $\mathcal{O}_Y$-模。以 $T \subset Y$ 表示 $\mathcal{G}$ 的支集，
并置 $V = Y \setminus T$。由于 $\mathcal{G}$ 有限表示，概形 $V$ 是 $Y$ 的
拟紧开集。由引理 \ref{functors-lemma-supported-on-support}，$F$ 在下列范畴之间诱导等价：
\begin{enumerate}
\item $\QCoh(\mathcal{O}_X)$ 中由支集在 $Z$ 上的模组成的全子范畴；
\item $\QCoh(\mathcal{O}_Y)$ 中由支集在 $T$ 上的模组成的全子范畴。
\end{enumerate}
由引理 \ref{functors-lemma-quotient-supported-on-closed} 得到交换图
$$
\xymatrix{
\QCoh(\mathcal{O}_X) \ar[r]_F \ar[d] &
\QCoh(\mathcal{O}_Y) \ar[d] \\
\QCoh(\mathcal{O}_U) \ar[r]^{F_U} &
\QCoh(\mathcal{O}_V)
}
$$
其中竖直箭头为限制函子，水平箭头为等价。由引理
\ref{functors-lemma-characterize-affine}，$V$ 仿射。仿射情形已有引理
\ref{functors-lemma-functor-equivalence}。因此存在同构 $f_U : V \to U$ 及可逆
$\mathcal{O}_V$-模 $\mathcal{L}_U$，使得 $F_U$ 是函子
$\mathcal{F} \mapsto f_U^*\mathcal{F} \otimes \mathcal{L}_U$.

\medskip\noindent
注意上述各图关于 $X$ 的仿射开子概形之间的包含满足显然的相容性，即可完成证明。
因此态射 $f_U$ 与可逆模 $\mathcal{L}_U$ 可粘合。细节从略。
\end{proof}










\section{相干模范畴之间的函子}
\label{functors-section-functor-coherent}

\noindent
下述引理保证：给定相干模范畴之间的函子时，可以使用拟凝聚模范畴之间函子的材料。

\begin{lemma}
\label{functors-lemma-functor-coherent}
设 $X$ 和 $Y$ 为诺特概形，并设
$F : \textit{Coh}(\mathcal{O}_X) \to \textit{Coh}(\mathcal{O}_Y)$
为函子。则 $F$ 唯一延拓为与滤过余极限交换的函子
$\QCoh(\mathcal{O}_X) \to \QCoh(\mathcal{O}_Y)$
。若 $F$ 可加，则其延拓与任意直和交换。若 $F$ 正合、左正合或右正合，
则其延拓也具有相应性质。
\end{lemma}

\begin{proof}
延拓的存在唯一性是一般事实；参见 Categories, Lemma
\ref{categories-lemma-extend-functor-by-colim}。为说明该引理适用，注意相干模是
有限表示的（Modules, Lemma \ref{modules-lemma-coherent-finite-presentation}），
因而由 Modules, Lemma
\ref{modules-lemma-finite-presentation-quasi-compact-colimit}，它们是
$\textit{Mod}(\mathcal{O}_X)$ 的范畴紧对象。最后，例如由 Properties, Lemma
\ref{properties-lemma-quasi-coherent-colimit-finite-type}，每个拟凝聚模都是相干模的
滤过余极限。

\medskip\noindent
假设 $F$ 可加。若 $\mathcal{F} = \bigoplus_{j \in J} \mathcal{H}_j$，其中
$\mathcal{H}_j$ 拟凝聚，则
$\mathcal{F} = \colim_{J' \subset J\text{ 有限}}
\bigoplus_{j \in J'} \mathcal{H}_j$.
仍把 $F$ 的延拓记作 $F$，得到
\begin{align*}
F(\mathcal{F})
& =
\colim_{J' \subset J\text{ 有限}}
F(\bigoplus\nolimits_{j \in J'} \mathcal{H}_j) \\
& =
\colim_{J' \subset J\text{ 有限}}
\bigoplus\nolimits_{j \in J'} F(\mathcal{H}_j) \\
& =
\bigoplus\nolimits_{j \in J} F(\mathcal{H}_j)
\end{align*}
故 $F$ 与任意直和交换。

\medskip\noindent
设 $0 \to \mathcal{F} \to \mathcal{F}' \to \mathcal{F}'' \to 0$ 为拟凝聚
$\mathcal{O}_X$-模的短正合列。把 $\mathcal{F}' = \bigcup \mathcal{F}'_i$ 写成
其相干子模的并；参见 Properties, Lemma
\ref{properties-lemma-quasi-coherent-colimit-finite-type}。以
$\mathcal{F}''_i \subset \mathcal{F}''$ 表示 $\mathcal{F}'_i$ 的像，并置
$\mathcal{F}_i = \mathcal{F} \cap \mathcal{F}'_i =
\Ker(\mathcal{F}'_i \to \mathcal{F}''_i)$。显然
$\mathcal{F} = \bigcup \mathcal{F}_i$ 且
$\mathcal{F}'' = \bigcup \mathcal{F}''_i$，并且有短正合列
$$
0 \to \mathcal{F}_i \to \mathcal{F}_i' \to \mathcal{F}_i'' \to 0
$$
由于该延拓与滤过余极限交换，有
$F(\mathcal{F}) = \colim_{i \in I} F(\mathcal{F}_i)$、
$F(\mathcal{F}') = \colim_{i \in I} F(\mathcal{F}'_i)$，以及
$F(\mathcal{F}'') = \colim_{i \in I} F(\mathcal{F}''_i)$。
由于滤过余极限正合（Modules, Lemma \ref{modules-lemma-limits-colimits}），
$F$ 的正合性性质由其延拓继承。
\end{proof}

\begin{lemma}
\label{functors-lemma-equivalence-coherent}
设 $X$ 和 $Y$ 为诺特概形，并设
$F : \textit{Coh}(\mathcal{O}_X) \to \textit{Coh}(\mathcal{O}_Y)$
为范畴等价。则存在同构 $f : Y \to X$ 及可逆 $\mathcal{O}_Y$-模
$\mathcal{L}$，使得 $F(\mathcal{F}) = f^*\mathcal{F} \otimes \mathcal{L}$。
\end{lemma}

\begin{proof}
由引理 \ref{functors-lemma-functor-coherent}，得到唯一的函子
$F' : \QCoh(\mathcal{O}_X) \to \QCoh(\mathcal{O}_Y)$ 延拓 $F$。$F$ 的拟逆亦然，
故由唯一性，$F'$ 是等价。由命题 \ref{functors-proposition-gabriel-rosenberg}，存在同构
$f : Y \to X$ 及可逆 $\mathcal{O}_Y$-模 $\mathcal{L}$，使得
$F'(\mathcal{F}) = f^*\mathcal{F} \otimes \mathcal{L}$。于是 $f$ 和
$\mathcal{L}$ 对 $F$ 也适用。
\end{proof}

\begin{remark}
\label{functors-remark-equivalence-coherent-linear}
在引理 \ref{functors-lemma-equivalence-coherent} 中，若 $X$ 和 $Y$ 定义在同一基环 $R$
上，且 $F$ 是 $R$-线性的，则同构 $f$ 是 $R$ 上的概形态射。
\end{remark}

\begin{lemma}
\label{functors-lemma-characterize-finite}
设 $f : V \to X$ 为诺特概形之间的拟有限分离态射。若存在相干
$\mathcal{O}_V$-模 $\mathcal{K}$，其支集为 $V$，且 $f_*\mathcal{K}$ 相干并有
$R^qf_*\mathcal{K} = 0$，则 $f$ 有限。
\end{lemma}

\begin{proof}
由 Zariski 主定理，可找到开浸入 $j : V \to Y$ 在 $X$ 上，其中
$\pi : Y \to X$ 有限；参见 More on Morphisms, Lemma
\ref{more-morphisms-lemma-quasi-finite-separated-pass-through-finite}。由于 $\pi$
仿射，函子 $\pi_*$ 在相干 $\mathcal{O}_X$-模范畴上正合且忠实。因此
$j_*\mathcal{K}$ 相干，且 $R^qj_*\mathcal{K}$ 在 $q > 0$ 时为零。换言之，
问题归约到下一段讨论的情形。

\medskip\noindent
假设 $f$ 为开浸入。可将 $X$ 替换为 $V$ 的概形论闭包。为导出矛盾，假设
$X \setminus V$ 非空。选取泛点 $\xi \in X \setminus V$，它属于
$X \setminus V$ 的某个不可约分支。利用平坦基变换，在沿
$\Spec(\mathcal{O}_{X, \xi}) \to X$ 基变换后考察情形，并使用 Local Cohomology,
Lemma \ref{local-cohomology-lemma-finiteness-pushforwards-and-H1-local}，归约到下一段
讨论的代数问题。

\medskip\noindent
设 $(A, \mathfrak m)$ 为诺特局部环，$M$ 为有限 $A$-模，且其支集为
$\Spec(A)$。则 $H^i_\mathfrak m(M) \not = 0$，其中 $i$ 为某个整数。这由 Dualizing
Complexes, Lemma \ref{dualizing-lemma-depth} 以及 $M$ 非零从而具有有限深度这一
事实得出。
\end{proof}

\noindent
下一个引理可推广到 $k$ 为诺特环且 $X$ 在 $k$ 上平坦的情形（其他假设不变）。

\begin{lemma}
\label{functors-lemma-functor-coherent-over-field}
设 $k$ 为域，$X$、$Y$ 为 $k$ 上的有限型概形，且 $X$ 分离。下列两个范畴等价：
\begin{enumerate}
\item $k$-线性正合函子
$F : \textit{Coh}(\mathcal{O}_X) \to \textit{Coh}(\mathcal{O}_Y)$
所成的范畴；
\item 相干 $\mathcal{O}_{X \times Y}$-模 $\mathcal{K}$ 所成的范畴，其中这些模
在 $X$ 上平坦，且其支集在 $Y$ 上有限。
\end{enumerate}
此等价把 $\mathcal{K}$ 送到函子 (\ref{functors-equation-FM-QCoh}) 在
$\textit{Coh}(\mathcal{O}_X)$ 上的限制。
\end{lemma}

\begin{proof}
设 $\mathcal{K}$ 如 (2) 中所述。由引理
\ref{functors-lemma-functor-quasi-coherent-from-affine-diagonal}，
(\ref{functors-equation-FM-QCoh}) 给出的函子 $F$ 正合且 $k$-线性。此外，例如由
Cohomology of Schemes, Lemma
\ref{coherent-lemma-support-proper-over-base-pushforward}，$F$ 把
$\textit{Coh}(\mathcal{O}_X)$ 送入 $\textit{Coh}(\mathcal{O}_Y)$。

\medskip\noindent
构造该构造的拟逆。设 $F$ 如 (1) 中所述。由引理
\ref{functors-lemma-functor-coherent}，可把 $F$ 延拓为拟凝聚模范畴上的 $k$-线性正合函子，
且它与任意直和交换。由引理
\ref{functors-lemma-functor-quasi-coherent-from-affine-diagonal}，该延拓对应于唯一的拟凝聚模
$\mathcal{K}$，它在 $X$ 上平坦，并满足
$R^q\text{pr}_{2, *}(\text{pr}_1^*\mathcal{F}
\otimes_{\mathcal{O}_{X \times Y}} \mathcal{K}) = 0$，其中 $q > 0$，并且上述
$\mathcal{O}_X$-模 $\mathcal{F}$ 可为任意拟凝聚模。由于
$F(\mathcal{O}_X)$ 是相干 $\mathcal{O}_Y$-模，由引理
\ref{functors-lemma-functor-quasi-coherent-from-separated} 可知 $\mathcal{K}$ 相干。

\medskip\noindent
对闭点 $x \in X$，以 $\mathcal{O}_x$ 表示在 $x$ 处取值为 $x$ 的剩余域的摩天层。
有
$$
F(\mathcal{O}_x) =
\text{pr}_{2, *}(\text{pr}_1^*\mathcal{O}_x \otimes \mathcal{K}) =
(x \times Y \to Y)_*(\mathcal{K}|_{x \times Y})
$$
由于 $x \times Y \to Y$ 有限，沿该态射的推前忠实。因此，若 $y \in Y$ 属于
$\mathcal{K}|_{x \times Y}$ 的支集之像，则 $y$ 属于 $F(\mathcal{O}_x)$ 的支集。

\medskip\noindent
设 $Z \subset X \times Y$ 为概形论支集 $Z$，即 $\mathcal{K}$ 的概形论支集；参见 Morphisms,
Definition \ref{morphisms-definition-scheme-theoretic-support}。先证明 $Z \to Y$
拟有限；为此证明其在闭点上的纤维有限。确切地说，若 $Z \to Y$ 在闭点
$y \in Y$ 上的纤维维数 $> 0$，则可找到无穷多个两两不同的闭点
$x_1, x_2, \ldots$，它们位于 $Z_y \to X$ 的像中。由于存在满射
$\mathcal{O}_X \to \bigoplus_{i = 1, \ldots, n} \mathcal{O}_{x_i}$
，得到满射
$$
F(\mathcal{O}_X) \to \bigoplus\nolimits_{i = 1, \ldots, n} F(\mathcal{O}_{x_i})
$$
由上述讨论，点 $y$ 属于每个相干模 $F(\mathcal{O}_{x_i})$ 的支集。由于
$F(\mathcal{O}_X)$ 是相干模，这会导出矛盾：$F(\mathcal{O}_X)$ 在 $y$ 处的茎
可由 $< n$ 个元素生成，只要 $n$ 足够大。因此 $Z \to Y$ 拟有限。由于
$\text{pr}_{2, *}\mathcal{K}$ 相干，且 $R^q\text{pr}_{2, *}\mathcal{K} = 0$
对 $q > 0$ 成立，由引理 \ref{functors-lemma-characterize-finite}，
$Z \to Y$ 有限。
\end{proof}

\begin{lemma}
\label{functors-lemma-pushforward-invertible-pre}
设 $f : X \to Y$ 为概形之间的有限型分离态射。设 $\mathcal{F}$ 为 $X$ 上的
有限型拟凝聚模，其支集在 $Y$ 上有限，并且
$\mathcal{L} = f_*\mathcal{F}$ 是可逆 $\mathcal{O}_X$-模。则存在截面
$s : Y \to X$，使得 $\mathcal{F} \cong s_*\mathcal{L}$。
\end{lemma}

\begin{proof}
在仿射局部考察，这翻译为下述代数问题。设 $A \to B$ 为环态射，$N$ 为
$B$-模，并且作为 $A$-模可逆。则零化子 $J$，即 $N$ 在 $B$ 中的零化子，具有如下性质：
$A \to B/J$ 是同构。细节从略。
\end{proof}

\begin{lemma}
\label{functors-lemma-pushforward-invertible}
设 $f : X \to Y$ 为概形之间的有限型分离态射，且有截面 $s : Y \to X$。设
$\mathcal{F}$ 为 $X$ 上的有限型拟凝聚模，在集合意义下支集于 $s(Y)$，并且
$\mathcal{L} = f_*\mathcal{F}$
为可逆 $\mathcal{O}_X$-模。若 $Y$ 既约，则 $\mathcal{F} \cong s_*\mathcal{L}$。
\end{lemma}

\begin{proof}
由引理 \ref{functors-lemma-pushforward-invertible-pre}，存在截面 $s' : Y  \to X$，使得
$\mathcal{F} = s'_*\mathcal{L}$。由于 $s'(Y)$ 与 $s(Y)$ 有相同的底闭子集，
并且二者都是 $X$ 的既约闭子概形，故必相等。因此 $s = s'$，引理得证。
\end{proof}

\begin{lemma}
\label{functors-lemma-equivalence-coherent-over-field}
\begin{reference}
\cite{Gabriel} 中“拟凝聚模范畴决定概形同构类”这一结果的弱版本。
\end{reference}
设 $k$ 为域，$X$、$Y$ 为 $k$ 上的有限型概形，且 $X$ 分离、$Y$ 既约。若存在
$k$-线性范畴等价
$F : \textit{Coh}(\mathcal{O}_X) \to \textit{Coh}(\mathcal{O}_Y)$
，则存在同构 $f : Y \to X$ 在 $k$ 上，以及可逆 $\mathcal{O}_Y$-模
$\mathcal{L}$，使得 $F(\mathcal{F}) = f^*\mathcal{F} \otimes \mathcal{L}$。
\end{lemma}

\begin{proof}[使用 Gabriel--Rosenberg 重构的证明]
本引理是引理 \ref{functors-lemma-equivalence-coherent} 与注
\ref{functors-remark-equivalence-coherent-linear} 中所讨论结果的弱形式。
\end{proof}

\begin{proof}[不依赖 Gabriel--Rosenberg 重构的证明]
由引理 \ref{functors-lemma-functor-coherent-over-field}，得到相干
$\mathcal{O}_{X \times Y}$-模 $\mathcal{K}$，它在 $X$ 上平坦、支集在 $Y$ 上有限，
并且 $F$ 由函子 (\ref{functors-equation-FM-QCoh}) 在
$\textit{Coh}(\mathcal{O}_X)$ 上的限制给出。若能证明 $F(\mathcal{O}_X)$ 是可逆
$\mathcal{O}_Y$-模，则由引理 \ref{functors-lemma-pushforward-invertible-pre}，有
$\mathcal{K} = s_*\mathcal{L}$，其中 $s : Y \to X \times Y$ 是
$\text{pr}_2$ 的某个截面，另有某个可逆 $\mathcal{O}_Y$-模 $\mathcal{L}$。这将说明
$F$ 具有所述形式，其中 $f = \text{pr}_1 \circ s$。部分细节从略。

\medskip\noindent
只需再证明 $F(\mathcal{O}_X)$ 可逆。这里只勾勒证明并省略一些细节。对闭点
$x \in X$，以 $\mathcal{O}_x$ 表示 $\textit{Coh}(\mathcal{O}_X)$ 中在 $x$ 处
取值为 $\kappa(x)$ 的摩天层。首先注意，范畴
$\textit{Coh}(\mathcal{O}_X)$ 的单对象恰为这些摩天层 $\mathcal{O}_x$。对 $Y$
也同样成立。因此，对每个闭点 $y \in Y$，存在闭点 $x \in X$，使得
$\mathcal{O}_y \cong F(\mathcal{O}_x)$。此外，考察自同态可知，有
$\kappa(x) \cong \kappa(y)$，它们是 $k$ 的有限扩张。于是
$$
\Hom_Y(F(\mathcal{O}_X), \mathcal{O}_y) \cong
\Hom_Y(F(\mathcal{O}_X), F(\mathcal{O}_x)) \cong
\Hom_X(\mathcal{O}_X, \mathcal{O}_x) \cong \kappa(x) \cong \kappa(y)
$$
这蕴含相干 $\mathcal{O}_Y$-模 $F(\mathcal{O}_X)$ 在 $y \in Y$ 处的茎可由
$1$ 个生成元生成（且不能更少），对每个闭点 $y \in Y$ 都如此。立即可知
$F(\mathcal{O}_X)$ 局部可由 $1$ 个元素生成（且不能更少）；又因 $Y$ 既约，
这确实说明它是可逆模。
\end{proof}
% Generated mechanically from the frozen equiv.tex target.
% Do not edit this generated file; edit the bound locale source instead.

% OK, start here.
%

\setcounter{chapter}{56}
\chapter{代数簇的导出范畴}



\phantomsection
\label{equiv-section-phantom}


\section{引言}
\label{equiv-section-introduction}

\noindent
本章继续 Derived Categories of Schemes, Section
\ref{perfect-section-introduction} 中开始的讨论。
我们将讨论 Fourier--Mukai 变换；Mukai 在 \cite{Mukai} 中最先研究了这种变换。
我们将证明关于导出等价的 Orlov 定理（\cite{Orlov-K3}）。
我们还将讨论 Anel 与 To\"en 在 \cite{AT} 中证明的导出等价类的可数性。

\medskip\noindent
Daniel Huybrechts 的著作 \cite{Huybrechts} 是本专题的良好入门。
其他推动这一专题普及的论文包括：
\begin{enumerate}
\item Bondal 与 Kapranov 的论文，参见 \cite{Bondal-Kapranov}；
\item Bondal 与 Orlov 的论文，参见 \cite{Bondal-Orlov}；
\item Bondal 与 Van den Bergh 的论文，参见 \cite{BvdB}；
\item Beilinson 的论文，参见
\cite{Beilinson} 与 \cite{Beilinson-derived}；
\item Orlov 的论文，参见 \cite{Orlov-AV}；
\item Orlov 的论文，参见 \cite{Orlov-motives}；
\item Rouquier 的论文，参见 \cite{Rouquier-dimensions}；
\item 这里还可以列举许多其他论文。
\end{enumerate}




\section{约定与记号}
\label{equiv-section-conventions}

\noindent
设 $k$ 为域。$k$-线性三角范畴 $\mathcal{T}$ 是一个三角范畴
（Derived Categories, Section
\ref{derived-section-triangulated-definitions}），其上赋有 $k$-线性结构
（Differential Graded Algebra, Section \ref{dga-section-linear}），
并且平移函子 $[n] : \mathcal{T} \to \mathcal{T}$ 是 $k$-线性的；
这对所有 $n \in \mathbf{Z}$ 都成立。

\medskip\noindent
设 $k$ 为域。用 $\text{Vect}_k$ 表示 $k$-向量空间范畴。
对 $k$-向量空间 $V$，用 $V^\vee$ 表示 $k$-线性意义下 $V$ 的对偶，
即 $V^\vee = \Hom_k(V, k)$。

\medskip\noindent
设 $X$ 为概形。用 $D_{perf}(\mathcal{O}_X)$ 表示
$D(\mathcal{O}_X)$ 中由完美复形组成的全子范畴
（Cohomology, Section \ref{cohomology-section-perfect}）。
若 $X$ 为 Noether 概形，则
$D_{perf}(\mathcal{O}_X) \subset D^b_{\textit{Coh}}(\mathcal{O}_X)$；参见
Derived Categories of Schemes, Lemma \ref{perfect-lemma-perfect-on-noetherian}。
若 $X$ 为正则 Noether 概形，则
$D_{perf}(\mathcal{O}_X) = D^b_{\textit{Coh}}(\mathcal{O}_X)$；参见
Derived Categories of Schemes, Lemma \ref{perfect-lemma-perfect-on-regular}。

\medskip\noindent
设 $k$ 为域，$X$ 和 $Y$ 为 $k$ 上的概形。在这种情形下，
用 $X \times Y$ 代替 $X \times_{\Spec(k)} Y$。


\medskip\noindent
设 $S$ 为概形，$X$、$Y$ 为 $S$ 上的概形。
设 $\mathcal{F}$ 为 $\mathcal{O}_X$-模，
$\mathcal{G}$ 为 $\mathcal{O}_Y$-模。定义
$$
\mathcal{F} \boxtimes \mathcal{G} =
\text{pr}_1^*\mathcal{F} \otimes_{\mathcal{O}_{X \times_S Y}}
\text{pr}_2^*\mathcal{G}
$$
作为 $\mathcal{O}_{X \times_S Y}$-模。若
$K \in D(\mathcal{O}_X)$ 且 $M \in D(\mathcal{O}_Y)$，则定义
$$
K \boxtimes M =
L\text{pr}_1^*K \otimes_{\mathcal{O}_{X \times_S Y}}^\mathbf{L} L\text{pr}_2^*M
$$
作为 $D(\mathcal{O}_{X \times_S Y})$ 的对象。
因此这一记号可能有歧义，但上下文应能说明所指的是二者中的哪一个。





\section{Serre 函子}
\label{equiv-section-Serre-functors}

\noindent
本节材料取自 \cite{Bondal-Kapranov}。

\begin{lemma}
\label{equiv-lemma-Serre-functor-exists}
设 $k$ 为域，$\mathcal{T}$ 为 $k$-线性三角范畴，并且满足
$\dim_k \Hom_\mathcal{T}(X, Y) < \infty$，其中
$X, Y \in \Ob(\mathcal{T})$ 任意。下列条件等价：
\begin{enumerate}
\item 存在 $k$-线性等价
$S : \mathcal{T} \to \mathcal{T}$ 以及 $k$-线性同构
$c_{X, Y} : \Hom_\mathcal{T}(X, Y) \to \Hom_\mathcal{T}(Y, S(X))^\vee$
，且该同构关于 $X, Y \in \Ob(\mathcal{T})$ 具有函子性；
\item 对每个 $X \in \Ob(\mathcal{T})$，函子
$Y \mapsto \Hom_\mathcal{T}(X, Y)^\vee$ 可表，而函子
$Y \mapsto \Hom_\mathcal{T}(Y, X)^\vee$ 余可表。
\end{enumerate}
\end{lemma}

\begin{proof}
条件 (1) 蕴含 (2)，因为给定 $(S, c)$ 和 $X \in \Ob(\mathcal{T})$ 后，
对象 $S(X)$ 表示函子
$Y \mapsto \Hom_\mathcal{T}(X, Y)^\vee$，而对象 $S^{-1}(X)$ 余表示函子
$Y \mapsto \Hom_\mathcal{T}(Y, X)^\vee$。

\medskip\noindent
假设 (2)。下面将反复使用 Yoneda 引理；参见
Categories, Lemma \ref{categories-lemma-yoneda}。
对每个 $X$，以 $S(X)$ 表示函子
$Y \mapsto \Hom_\mathcal{T}(X, Y)^\vee$ 的表示对象。给定
$\varphi : X \to X'$，得到唯一的箭头 $S(\varphi) : S(X) \to S(X')$；
它由相应的函子变换
$\Hom_\mathcal{T}(X, -)^\vee \to \Hom_\mathcal{T}(X', -)^\vee$ 确定。
这样，$S$ 是函子，并且按构造得到同构 $c_{X, Y}$。
还须证明 $S$ 是等价。对每个 $X$，以 $S'(X)$ 表示函子
$Y \mapsto \Hom_\mathcal{T}(Y, X)^\vee$ 的余表示对象。与上面相同的论证表明
$S'$ 是函子。我们断言 $S'$ 是 $S$ 的拟逆。为此，注意
$$
\Hom_\mathcal{T}(X, Y) = \Hom_\mathcal{T}(Y, S(X))^\vee =
\Hom_\mathcal{T}(S'(S(X)), Y)
$$
在双函子意义下成立；即得到 $S' \circ S \cong \text{id}_\mathcal{T}$。
类似地，有
$$
\Hom_\mathcal{T}(Y, X) = \Hom_\mathcal{T}(S'(X), Y)^\vee =
\Hom_\mathcal{T}(Y, S(S'(X)))
$$
从而得到 $S \circ S' \cong \text{id}_\mathcal{T}$。
\end{proof}

\begin{definition}
\label{equiv-definition-Serre-functor}
设 $k$ 为域，$\mathcal{T}$ 为 $k$-线性三角范畴，并且满足
$\dim_k \Hom_\mathcal{T}(X, Y) < \infty$，其中
$X, Y \in \Ob(\mathcal{T})$ 任意。若引理 \ref{equiv-lemma-Serre-functor-exists}
中的等价条件成立，则称{\it 存在 Serre 函子}。在这种情形下，
{\it Serre 函子}是一个 $k$-线性等价
$S : \mathcal{T} \to \mathcal{T}$，并配备 $k$-线性同构
$c_{X, Y} : \Hom_\mathcal{T}(X, Y) \to \Hom_\mathcal{T}(Y, S(X))^\vee$
；这些同构关于 $X, Y \in \Ob(\mathcal{T})$ 具有函子性。
\end{definition}

\begin{remark}
\label{equiv-remark-matress}
设 $X^0 \to X^1 \to X^2 \to X^0[1]$ 和
$Y^0 \to Y^1 \to Y^2 \to Y^0[1]$ 是某三角范畴中的可区别三角。
对 $p \in \mathbf{Z}$，写成 $p = 3n + i$，其中
$i \in \{0, 1, 2\}$，并定义 $X^p = X^i[n]$。对 $Y^q$ 作类似定义。
考虑项为
$$
K^{p, q} = \Hom(X^{-p}, Y^q)
$$
的双复形。微分 $d_1 : K^{p, q} \to K^{p + 1, q}$ 由映射
$X^{-p - 1} \to X^{-p}$ 给出（它等于第一个可区别三角中的相应映射，至多差一个平移）；
微分 $d_2 : K^{p, q} \to K^{p, q + 1}$ 类似地由映射
$Y^q \to Y^{q + 1}$ 给出。由 Derived Categories, Lemma
\ref{derived-lemma-representable-homological} 可知，这个双复形的各行各列都正合。
此外，$K^{p, q} = K^{p + 3, q - 3}$，且这些等式与微分相容。最后，公理 TR3
还蕴含一个性质：给定
$\alpha \in K^{p, q}$ 和 $\beta \in K^{p - 1, q + 1}$，若
$d_2 \alpha = d_1 \beta$，则存在 $\gamma \in K^{p - 2, q + 2}$，使得
$d_1 \gamma = d_2 \beta $ 在 $K^{p - 1, q + 2}$ 中成立，并且
$d_2 \gamma = d_1 \alpha $ 在 $K^{p - 2, q + 3} = K^{p + 1, q}$ 中成立。
（提示：当 $p = q = 0$ 时，这正是 TR3 的陈述；其他指标的情形由平移证明。）
具有这些性质的双复形称为 {\it matress}；参见 \cite{Bondal-Kapranov}。
\end{remark}

\begin{remark}
\label{equiv-remark-dual-matress}
设 $k$ 为域，$K^{\bullet, \bullet}$ 为有限维 $k$-向量空间的双复形，
并且它还是注 \ref{equiv-remark-matress} 中的 matress。我们断言，对偶双复形
$L^{p, q} = \Hom_k(K^{-q, -p}, k)$ 也是 matress。各行各列的正合性以及
3-周期性都是直接的。为验证附加条件，考虑线性映射
$$
\partial :
K^{p, q} \oplus K^{p - 1, q + 1} \oplus K^{p - 2, q + 2}
\longrightarrow
K^{p, q + 1} \oplus K^{p - 1, q + 2} \oplus K^{p - 2, q + 3}
$$
它把 $(\alpha, \beta, \gamma)$ 映到
$(d_2 \alpha - d_1 \beta, d_2 \gamma - d_1 \alpha, d_1 \gamma - d_2 \beta)$
；这里采用注 \ref{equiv-remark-matress} 中的识别。成为 matress 的条件是
$$
\Im(\partial) \cap
\left( 0 \oplus K^{p - 1, q + 2} \oplus K^{p - 2, q + 3} \right)
=
\Im(\partial|_{0 \oplus 0 \oplus K^{p - 2, q + 2}})
$$
以 ${}^\wedge$ 表示向量空间或映射的对偶。取对偶得到
$$
\Ker(\partial^\wedge) +
\left(L^{-p, -q - 1} \oplus 0 \oplus 0 \right)
=
\Ker(\text{pr}_{L^{-p + 2, -q - 2}} \circ \partial^\wedge)
$$
其中 $\text{pr}$ 项表示投影。这蕴含
$\Im(\partial^\wedge) \cap (L^{-p, -q} \oplus L^{-p + 1, -q - 1} \oplus 0)$
等于
$\Im(\partial^\wedge|_{L^{-p, -q - 1} \oplus 0 \oplus 0})$；这正转化为
$L^{\bullet, \bullet}$ 的 matress 条件。略去若干细节。
\end{remark}

\begin{lemma}
\label{equiv-lemma-Serre-functor}
在定义 \ref{equiv-definition-Serre-functor} 的情形下，若 Serre 函子存在，
则它在唯一同构意义下唯一，并且是三角范畴之间的正合函子。
\end{lemma}

\begin{proof}
给定 Serre 函子 $S$，对象 $S(X)$ 表示函子
$Y \mapsto \Hom_\mathcal{T}(X, Y)^\vee$。
因此，由 Yoneda 引理（Categories, Lemma \ref{categories-lemma-yoneda}），
对象 $S(X)$ 连同函子性识别
$\Hom_\mathcal{T}(X, Y)^\vee = \Hom_\mathcal{T}(Y, S(X))$
在唯一同构意义下确定。此外，对 $\varphi : X \to X'$，箭头
$S(\varphi) : S(X) \to S(X')$ 由相应的函子变换
$\Hom_\mathcal{T}(X, -)^\vee \to \Hom_\mathcal{T}(X', -)^\vee$ 唯一确定。

\medskip\noindent
对对象 $X, Y$（属于 $\mathcal{T}$），有
\begin{align*}
\Hom(Y, S(X)[1])^\vee
& =
\Hom(Y[-1], S(X))^\vee \\
& =
\Hom(X, Y[-1]) \\
& =
\Hom(X[1], Y) \\
& =
\Hom(Y, S(X[1]))^\vee
\end{align*}
由 Yoneda 引理，存在唯一同构
$S(X[1]) \to S(X)[1]$，它诱导上述从左上到右下的同构。
由于上面的每个同构都同时关于 $X$ 和 $Y$ 具有函子性，
这就定义了函子同构
$S \circ [1] \to [1] \circ S$。

\medskip\noindent
设 $(A, B, C, f, g, h)$ 是 $\mathcal{T}$ 中的可区别三角。
须证明三角 $(S(A), S(B), S(C), S(f), S(g), S(h))$ 是可区别的。
这里使用上面构造的典范同构 $S(A[1]) \to S(A)[1]$，把目标
$S(A[1])$（即 $S(h)$ 的目标）与 $S(A)[1]$ 识别。先注意，对任意 $X$（属于 $\mathcal{T}$），
三角 $(S(A), S(B), S(C), S(f), S(g), S(h))$ 诱导长正合列
$$
\ldots \to
\Hom(X, S(A)) \to
\Hom(X, S(B)) \to
\Hom(X, S(C)) \to
\Hom(X, S(A)[1]) \to \ldots
$$
，其中各项都是有限维 $k$-向量空间。事实上，这个序列是下列序列的
$k$-线性对偶：
$$
\ldots \leftarrow
\Hom(A, X) \leftarrow
\Hom(B, X) \leftarrow
\Hom(C, X) \leftarrow
\Hom(A[1], X) \leftarrow
\ldots
$$
后一个序列由 Derived Categories, Lemma
\ref{derived-lemma-representable-homological} 可知是正合的。
接着，选取可区别三角 $(S(A), E, S(C), i, p, S(h))$；
由公理 TR1 和 TR2，这样的选取是可能的。我们要构造虚线箭头，使下图交换：
$$
\xymatrix{
S(C)[-1] \ar[r]_-{S(h[-1])} &
S(A) \ar[r]_{S(f)} &
S(B) \ar[r]_{S(g)} &
S(C) \ar[r]_{S(h)} &
S(A)[1] \\
S(C)[-1] \ar[r]^-{S(h[-1])} \ar@{=}[u] &
S(A) \ar[r]^i \ar@{=}[u] &
E \ar[r]^p \ar@{..>}[u]^\varphi &
S(C) \ar[r]^{S(h)} \ar@{=}[u] &
S(A)[1] \ar@{=}[u]
}
$$
具体地，若已有 $\varphi$，则断言对任意 $X$，由此得到的映射
$\Hom(X, E) \to \Hom(X, S(B))$ 是 $k$-向量空间同构。事实上，将得到交换图
$$
\xymatrix{
\Hom(X, S(C)[-1]) \ar[r] &
\Hom(X, S(A)) \ar[r] &
\Hom(X, S(B)) \ar[r] &
\Hom(X, S(C)) \ar[r] &
\Hom(X, S(A)[1]) \\
\Hom(X, S(C)[-1]) \ar[r] \ar@{=}[u] &
\Hom(X, S(A)) \ar[r] \ar@{=}[u] &
\Hom(X, E) \ar[r] \ar[u]^\varphi &
\Hom(X, S(C)) \ar[r] \ar@{=}[u] &
\Hom(X, S(A)[1]) \ar@{=}[u]
}
$$
其两行都正合（见上文）。应用五引理
（Homology, Lemma \ref{homology-lemma-five-lemma}），可知中间箭头是同构。
再由 Yoneda 引理，得到 $\varphi$ 是同构。

\medskip\noindent
为求得 $\varphi$，由可区别三角
$A \to B \to C \to A[1]$ 和 $S(A) \to E \to S(C) \to S(A)[1]$
构造注 \ref{equiv-remark-matress} 中的 matress。由注
\ref{equiv-remark-dual-matress}，它的 $k$-线性对偶也是 matress。
利用 $S$ 是 Serre 函子这一事实，与上面相同可知，这个对偶 matress
正是由三角（其中第二个尚未知道是可区别的）
$S(A) \to E \to S(C) \to S(A)[1]$ 和
$S(A) \to S(B) \to S(C) \to S(A)[1]$ 构成的双复形。
而 matress 条件恰好说明映射
$S(A) \to S(A)$ 和 $S(C) \to S(C)$ 可以补成三角之间的态射；
即存在使上图交换的 $\varphi$。
\end{proof}






\section{Serre 函子的例}
\label{equiv-section-examples-Serre-functors}

\noindent
下述引理给出标准例。

\begin{lemma}
\label{equiv-lemma-Serre-functor-Gorenstein-proper}
设 $k$ 为域，$X$ 为 $k$ 上的 Gorenstein 固有概形。
考虑 Duality for Schemes, Lemmas
\ref{duality-lemma-duality-proper-over-field} 中的复形 $\omega_X^\bullet$。
则函子
$$
S : D_{perf}(\mathcal{O}_X) \longrightarrow D_{perf}(\mathcal{O}_X),\quad
K \longmapsto S(K) = \omega_X^\bullet \otimes_{\mathcal{O}_X}^\mathbf{L} K
$$
是 Serre 函子。
\end{lemma}

\begin{proof}
该陈述有意义，因为由 Derived Categories of Schemes, Lemma
\ref{perfect-lemma-ext-finite}，有 $\dim \Hom_X(K, L) < \infty$，其中
$K, L \in D_{perf}(\mathcal{O}_X)$。由于 $X$ 是 Gorenstein 的，对偶复形
$\omega_X^\bullet$ 是 $D(\mathcal{O}_X)$ 的可逆对象；参见
Duality for Schemes, Lemma \ref{duality-lemma-gorenstein}。
特别地，在 $X$ 上局部，复形 $\omega_X^\bullet$ 只有一个非零上同调层，
且它是可逆模；参见 Cohomology, Lemma
\ref{cohomology-lemma-invertible-derived}。因此 $S(K)$ 属于
$D_{perf}(\mathcal{O}_X)$。另一方面，$\omega_X^\bullet$ 的可逆性显然蕴含
$S$ 是 $D_{perf}(\mathcal{O}_X)$ 的自等价。最后，须找到同构
$$
c_{K, L} : \Hom_X(K, L) \longrightarrow
\Hom_X(L, \omega_X^\bullet \otimes_{\mathcal{O}_X}^\mathbf{L} K)^\vee
$$
，并且它关于 $K, L$ 具有双函子性。为此，使用典范同构
$$
\Hom_X(K, L) = H^0(X, L \otimes_{\mathcal{O}_X}^\mathbf{L} K^\vee)
$$
以及
$$
\Hom_X(L, \omega_X^\bullet \otimes_{\mathcal{O}_X}^\mathbf{L} K) =
H^0(X, 
\omega_X^\bullet \otimes_{\mathcal{O}_X}^\mathbf{L} K
\otimes_{\mathcal{O}_X}^\mathbf{L} L^\vee)
$$
；它们由 Cohomology, Lemma \ref{cohomology-lemma-dual-perfect-complex} 给出。
由于 $(L \otimes_{\mathcal{O}_X}^\mathbf{L} K^\vee)^\vee =
(K^\vee)^\vee \otimes_{\mathcal{O}_X}^\mathbf{L} L^\vee$
，并且存在典范同构 $K \to (K^\vee)^\vee$，由
Duality for Schemes, Lemma
\ref{duality-lemma-duality-proper-over-field-perfect} 可知，这些
$k$-向量空间典范对偶。这就给出同构 $c_{K, L}$。
略去这些同构具有函子性的证明。
\end{proof}





\section{相干模的刻画}
\label{equiv-section-coherent}

\noindent
本节在某种意义上延续了 Derived Categories of Schemes, Section
\ref{perfect-section-pseudo-coherent} 与 More on Morphisms, Section
\ref{more-morphisms-section-characterize-pseudo-coherent} 中的讨论。

\medskip\noindent
在陈述结果之前，需要先引入一些记号。设 $k$ 为域，设 $n \geq 0$
为整数。令 $S = k[X_0, \ldots, X_n]$。对于整数 $e$，以
$S_e \subset S$ 表示次数为 $e$ 的齐次多项式所成的子空间。考察
（非交换）$k$-代数
$$
R =
\left(
\begin{matrix}
S_0 & S_1 & S_2 & \ldots & \ldots \\
0 & S_0 & S_1 & \ldots & \ldots\\
0 & 0 & S_0 & \ldots & \ldots \\
\ldots & \ldots & \ldots & \ldots & \ldots \\
0 & \ldots & \ldots & \ldots & S_0
\end{matrix}
\right)
$$
（有 $n + 1$ 行和列），其乘法与加法按显然方式定义。

\begin{lemma}
\label{equiv-lemma-perfect-for-R}
沿用上述 $k$、$n$ 和 $R$，设 $K$ 为 $D(R)$ 的对象。下列条件等价：
\begin{enumerate}
\item $\sum_{i \in \mathbf{Z}} \dim_k H^i(K) < \infty$；
\item $K$ 是紧对象。
\end{enumerate}
\end{lemma}

\begin{proof}
若 $K$ 是紧对象，则 $K$ 可由复形 $M^\bullet$ 表示，并且后者作为分次
$R$-模是有限投射的；参见 Differential Graded Algebra, Lemma
\ref{dga-lemma-compact}。由于 $\dim_k R < \infty$，可得
$\sum \dim_k M^i < \infty$，从而当然有
$\sum \dim_k H^i(M^\bullet) < \infty$。（也可以很容易地由较初等的
Differential Graded Algebra, Proposition \ref{dga-proposition-compact}
推出这一蕴含。）

\medskip\noindent
假设 $K$ 满足 (1)。考虑截断的可区别三角
$\tau_{\leq m}K \to K \to \tau_{\geq m + 1}K$；参见
Derived Categories, Remark
\ref{derived-remark-truncation-distinguished-triangle}。显然，
$\tau_{\leq m}K$ 与 $\tau_{\geq m + 1} K$ 都满足 (1)。若能证明二者
都是紧对象，则 $K$ 也是紧对象；参见 Derived Categories, Lemma
\ref{derived-lemma-compact-objects-subcategory}。因此，对 $K$ 的非零
上同调模的数目作归纳，可以假设 $H^i(K)$ 仅对一个 $i$ 非零。
作平移后，可以假设 $K$ 由只含一个有限维 $R$-模 $M$ 的复形给出，
该模位于次数 $0$。

\medskip\noindent
由于 $\dim_k(M) < \infty$，可知 $M$ 作为 $R$-模是 Artin 的。因此，
只需证明每个单 $R$-模都表示 $D(R)$ 的紧对象。注意
$$
I =
\left(
\begin{matrix}
0 & S_1 & S_2 & \ldots & \ldots \\
0 & 0 & S_1 & \ldots & \ldots\\
0 & 0 & 0 & \ldots & \ldots \\
\ldots & \ldots & \ldots & \ldots & \ldots \\
0 & \ldots & \ldots & \ldots & 0
\end{matrix}
\right)
$$
是 $R$ 的幂零双边理想，而 $R/I$ 是交换 $k$-代数，并且同构于
$n + 1$ 个 $k$ 的乘积（各因子沿矩阵对角线放置；换言之，$R/I$
可提升为 $k$-子代数，并嵌入 $R$）。由此，$R$ 恰有 $n + 1$ 个单模
同构类 $M_0, \ldots, M_n$（沿对角线放置）。考虑由行向量组成的右
$R$-模 $P_i$
$$
P_i =
\left(
\begin{matrix}
0 &
\ldots &
0 &
S_0 &
\ldots &
S_{i - 1} &
S_i
\end{matrix}
\right)
$$
，其乘法 $P_i \times R \to P_i$ 按显然方式定义。于是
$R \cong P_0 \oplus \ldots \oplus P_n$，这是一个右 $R$-模同构。显然
$R$ 是 $D(R)$ 的紧对象，
故每个 $P_i$ 都是 $D(R)$ 的紧对象。（当然也可由此推出每个 $P_i$
作为 $R$-模都是投射的，但这并不是本证明所需的结论。）显然，
$P_0 = M_0$ 是上述单 $R$-模中的第一个。对于 $P_1$，有短正合序列
$$
0 \to P_0^{\oplus n + 1} \to P_1 \to M_1 \to 0
$$
这说明 $M_1$ 属于一个另外两项均为紧对象的可区别三角，因而 $M_1$
是 $D(R)$ 的紧对象。更一般地，存在短正合序列
$$
0 \to C_i \to P_i \to M_i \to 0
$$
，其中 $C_i$ 是有限维 $R$-模，其单组成因子都同构于 $M_j$，这里
$j < i$。由归纳假设，先得到 $C_i$ 确定 $D(R)$ 的紧对象，继而得到
$M_i$ 也如此，正如所需。
\end{proof}

\begin{lemma}
\label{equiv-lemma-coherent-on-projective-space}
设 $k$ 为域，设 $n \geq 0$，并设
$K \in D_\QCoh(\mathcal{O}_{\mathbf{P}^n_k})$。下列条件等价：
\begin{enumerate}
\item $K$ 属于 $D^b_{\textit{Coh}}(\mathcal{O}_{\mathbf{P}^n_k})$；
\item $\sum_{i \in \mathbf{Z}}
\dim_k H^i(\mathbf{P}^n_k, E \otimes^\mathbf{L} K) < \infty$
对于任意完美对象 $E$ 成立，该对象属于
$D(\mathcal{O}_{\mathbf{P}^n_k})$；
\item $\sum_{i \in \mathbf{Z}}
\dim_k \Ext^i_{\mathbf{P}^n_k}(E, K) < \infty$
对于任意完美对象 $E$ 成立，该对象属于
$D(\mathcal{O}_{\mathbf{P}^n_k})$；
\item $\sum_{i \in \mathbf{Z}} \dim_k H^i(\mathbf{P}^n_k,
K \otimes^\mathbf{L} \mathcal{O}_{\mathbf{P}^n_k}(d)) < \infty$
对于 $d = 0, 1, \ldots, n$ 成立。
\end{enumerate}
\end{lemma}

\begin{proof}
由 Cohomology, Lemma \ref{cohomology-lemma-dual-perfect-complex}，
(2) 与 (3) 等价。若 (1) 成立，则对于完美的 $E$，导出张量积
$E \otimes^\mathbf{L} K$ 属于
$D^b_{\textit{Coh}}(\mathcal{O}_{\mathbf{P}^n_k})$，再由
Derived Categories of Schemes, Lemma
\ref{perfect-lemma-direct-image-coherent} 可知 (2) 成立。显然 (2) 蕴含
(4)，因为 $\mathcal{O}_{\mathbf{P}^n_k}(d)$ 可视为 $\mathbf{P}^n_k$ 的
导出范畴中的完美对象。因此，只需证明 (4) 蕴含 (1)。

\medskip\noindent
假设 (4) 成立。令 $R$ 如 Lemma \ref{equiv-lemma-perfect-for-R} 中所定义，并令
$P = \bigoplus_{d = 0, \ldots, n} \mathcal{O}_{\mathbf{P}^n_k}(-d)$。
回顾 $R = \text{End}_{\mathbf{P}^n_k}(P)$，而 $P$ 的其余自 Ext 都为零，
并且由 Derived Categories of Schemes, Lemma
\ref{perfect-lemma-Pn-module-category}，$P$ 确定等价
$- \otimes^\mathbf{L} P : D(R) \to D_\QCoh(\mathcal{O}_{\mathbf{P}^n_k})$。
设 $K$ 所对应的对象为 $L$，它属于 $D(R)$。于是
\begin{align*}
H^i(L)
& =
\Ext^i_{D(R)}(R, L) \\
& =
\Ext^i_{\mathbf{P}^n_k}(P, K) \\
& =
H^i(\mathbf{P}^n_k, K \otimes P^\vee) \\
& =
\bigoplus\nolimits_{d = 0, \ldots, n}
H^i(\mathbf{P}^n_k, K \otimes \mathcal{O}(d))
\end{align*}
；这里使用了 Differential Graded Algebra, Lemma
\ref{dga-lemma-upgrade-tensor-with-complex-derived}
（以及 $- \otimes^\mathbf{L} P$ 是等价这一事实）和 Cohomology, Lemma
\ref{cohomology-lemma-dual-perfect-complex}。因此，假设 (4) 蕴含 $L$
满足 Lemma \ref{equiv-lemma-perfect-for-R} 的条件 (2)，从而是 $D(R)$ 的紧对象。
于是 $K$ 是 $D_\QCoh(\mathcal{O}_{\mathbf{P}^n_k})$ 的紧对象。再由
Derived Categories of Schemes, Proposition
\ref{perfect-proposition-compact-is-perfect}，$K$ 是完美的。由于
$D_{perf}(\mathcal{O}_{\mathbf{P}^n_k}) =
D^b_{\textit{Coh}}(\mathcal{O}_{\mathbf{P}^n_k})$
，其中用到了 Derived Categories of Schemes, Lemma
\ref{perfect-lemma-perfect-on-regular}，故 (1) 成立。
\end{proof}

\begin{lemma}
\label{equiv-lemma-finiteness}
设 $X$ 是域 $k$ 上的固有概形。设
$K \in D^b_{\textit{Coh}}(\mathcal{O}_X)$，并设 $E$ 属于
$D(\mathcal{O}_X)$ 且是完美的。则
$\sum_{i \in \mathbf{Z}} \dim_k \Ext^i_X(E, K) < \infty$.
\end{lemma}

\begin{proof}
例如，结合 Derived Categories of Schemes, Lemmas
\ref{perfect-lemma-ext-finite} 与
\ref{perfect-lemma-ext-from-perfect-into-bounded-QCoh} 即得此结论。
另一种证明是结合 Derived Categories of Schemes, Lemmas
\ref{perfect-lemma-perfect-on-noetherian} 与
\ref{perfect-lemma-direct-image-coherent}。
\end{proof}

\begin{lemma}
\label{equiv-lemma-characterize-dbcoh-projective}
\begin{reference}
在射影情形，这是 \cite[Lemma 7.46]{Rouquier-dimensions}，并且也隐含在
\cite[Theorem A.1]{BvdB} 中。
\end{reference}
设 $X$ 是域 $k$ 上的固有概形，并设
$K \in \Ob(D_\QCoh(\mathcal{O}_X))$。下列条件等价：
\begin{enumerate}
\item $K \in D^b_{\textit{Coh}}(\mathcal{O}_X)$；
\item $\sum_{i \in \mathbf{Z}} \dim_k \Ext^i_X(E, K) < \infty$
对于所有完美 $E$ 成立，这些对象属于 $D(\mathcal{O}_X)$。
\end{enumerate}
\end{lemma}

\begin{proof}
蕴含 (1) $\Rightarrow$ (2) 由 Lemma \ref{equiv-lemma-finiteness} 得到。
蕴含 (2) $\Rightarrow$ (1) 由 More on Morphisms, Lemma
\ref{more-morphisms-lemma-characterize-relatively-perfect}
得到（关于域上的相对完美对象的含义，参见 Derived Categories of Schemes, Example
\ref{perfect-example-relatively-perfect-field}
）；下一段给出射影情形下较容易的证明。

\medskip\noindent
假设 (2) 成立且 $X$ 在 $k$ 上射影。选取闭浸入
$i : X \to \mathbf{P}^n_k$。只需证明 $Ri_*K$ 属于
$D^b_{\textit{Coh}}(\mathbf{P}^n_k)$，因为拟相干模 $\mathcal{F}$ 定义在
$X$ 上时，它是相干的（相应地，为零），当且仅当 $i_*\mathcal{F}$
是相干的（相应地，为零）。对于完美对象 $E$（它属于
$D(\mathcal{O}_{\mathbf{P}^n_k})$），$Li^*E$ 是 $D(\mathcal{O}_X)$ 的
完美对象，并且
$$
\Ext^q_{\mathbf{P}^n_k}(E, Ri_*K) = \Ext^q_X(Li^*E, K)
$$
因此，由假设可知
$\sum_{q \in \mathbf{Z}} \dim_k \Ext^q_{\mathbf{P}^n_k}(E, Ri_*K) < \infty$.
由 Lemma \ref{equiv-lemma-coherent-on-projective-space} 即得结论。
\end{proof}





\section{一个可表性定理}
\label{equiv-section-bondal-van-den-bergh}

\noindent
本节材料取自 \cite{BvdB}。

\medskip\noindent
设 $\mathcal{T}$ 为 $k$-线性三角范畴。本节考察 $k$-线性上同调函子
$H$，其定义域为 $\mathcal{T}$，值域为 $k$-向量空间范畴。这意味着
$H$ 是函子
$$
H : \mathcal{T}^{opp} \longrightarrow \text{Vect}_k
$$
，它是 $k$-线性的，并且对于任意可区别三角 $X \to Y \to Z$（位于
$\mathcal{T}$ 中），序列 $H(Z) \to H(Y) \to H(X)$ 是 $k$-向量空间的
正合序列。参见 Derived Categories, Definition
\ref{derived-definition-homological} 与 Differential Graded Algebra,
Section \ref{dga-section-linear}。

\begin{lemma}
\label{equiv-lemma-maps-from-compact-filtered}
设 $\mathcal{D}$ 为三角范畴，设 $\mathcal{D}' \subset \mathcal{D}$ 为
全三角子范畴，并设 $X \in \Ob(\mathcal{D})$。以箭头 $E \to X$ 为对象且
$E \in \Ob(\mathcal{D}')$ 的范畴是滤过的。
\end{lemma}

\begin{proof}
检验 Categories, Definition \ref{categories-definition-directed} 中的条件。
该范畴非空，因为它包含 $0 \to X$。若 $E_i \to X$（$i = 1, 2$）是对象，
则 $E_1 \oplus E_2 \to X$ 是对象，并且存在态射
$(E_i \to X) \to (E_1 \oplus E_2 \to X)$。最后，设
$a, b : (E \to X) \to (E' \to X)$ 是态射。选取可区别三角
$E \xrightarrow{a - b} E' \to E''$，它位于 $\mathcal{D}'$ 中。由公理 TR3，
得到三角的态射
$$
\xymatrix{
E \ar[r]_{a - b} \ar[d] &
E' \ar[d] \ar[r] & E'' \ar[d] \\
0 \ar[r] &
X \ar[r] &
X
}
$$
，而所得箭头 $(E' \to X) \to (E'' \to X)$ 使 $a$ 与 $b$ 相等。
\end{proof}

\begin{lemma}
\label{equiv-lemma-van-den-bergh}
\begin{reference}
\cite[Lemma 2.14]{CKN}
\end{reference}
设 $k$ 为域。设 $\mathcal{D}$ 是具有直和且紧生成的 $k$-线性三角范畴。
以 $\mathcal{D}_c$ 表示紧对象所成的全子范畴。设
$H : \mathcal{D}_c^{opp} \to \text{Vect}_k$ 是 $k$-线性上同调函子，并满足
$\dim_k H(X) < \infty$，其中 $X \in \Ob(\mathcal{D}_c)$ 任取。则 $H$
同构于函子 $X \mapsto \Hom(X, Y)$，其中 $Y \in \Ob(\mathcal{D})$
是某个对象。
\end{lemma}

\begin{proof}
下文将反复使用 Derived Categories, Lemma
\ref{derived-lemma-compact-objects-subcategory} 而不再说明。记
$G : \mathcal{D}_c \to \text{Vect}_k$ 为 $k$-线性同调函子，它把 $X$ 送到
$H(X)^\vee$。对于任意对象 $Y$（属于 $\mathcal{D}$），置
$$
G'(Y) = \colim_{X \to Y, X \in \Ob(\mathcal{D}_c)} G(X)
$$
由 Lemma \ref{equiv-lemma-maps-from-compact-filtered}，该余极限是滤过的。我们断言
$G'$ 是 $k$-线性同调函子，$G'$ 在 $\mathcal{D}_c$ 上的限制是 $G$，并且
$G'$ 把直和送到直和。

\medskip\noindent
具体而言，设 $Y_1 \to Y_2 \to Y_3$ 是可区别三角。设
$\xi \in G'(Y_2)$ 在 $G'(Y_3)$ 中的像为零。由于该余极限是滤过的，
$\xi$ 可由某个 $X \to Y_2$ 以及 $X \in \Ob(\mathcal{D}_c)$、
$g \in G(X)$ 表示。$\xi$ 在 $G'(Y_3)$ 中的像为零，意味着复合
$X \to Y_2 \to Y_3$ 分解为 $X \to X' \to Y_3$，其中
$X' \in \mathcal{D}_c$，且 $g$ 在 $G(X')$ 中的像为零。选取可区别三角
$X'' \to X \to X'$。于是 $X'' \in \Ob(\mathcal{D}_c)$。由于 $G$ 是
同调函子，可知 $g$ 是某个 $g'' \in G'(X'')$ 的像。由公理 TR3，态射
$X \to Y_2$ 与 $X' \to Y_3$ 可嵌入可区别三角的态射
$(X'' \to X \to X') \to (Y_1 \to Y_2 \to Y_3)$。由此确实可见，
$\xi$ 是 $G'(Y_1)$ 中由 $X'' \to Y_1$ 及 $g'' \in G(X'')$ 表示的元素的像。

\medskip\noindent
若 $Y \in \Ob(\mathcal{D}_c)$，则 $\text{id} : Y \to Y$ 是以
$X \to Y$ 为对象且 $X \in \Ob(\mathcal{D}_c)$ 的箭头范畴中的终对象。
所以在此情形 $G'(Y) = G(Y)$，从而关于限制的断言成立。设
$Y = \bigoplus_{i \in I} Y_i$ 是一个直和。设 $a : X \to Y$，其中
$X \in \Ob(\mathcal{D}_c)$，并设 $g \in G(X)$ 表示元素 $\xi$，该元素属于
$G'(Y)$。态射 $a : X \to Y$ 可唯一写成态射 $a_i : X \to Y_i$ 之和，
且几乎所有项为零；这是因为 $X$ 是 $\mathcal{D}$ 的紧对象。令
$I' = \{i \in I \mid a_i \not = 0\}$。则 $a$ 可分解为复合
$$
X \xrightarrow{(1, \ldots, 1)}
\bigoplus\nolimits_{i \in I'} X
\xrightarrow{\bigoplus_{i \in I'} a_i}
\bigoplus\nolimits_{i \in I} Y_i = Y
$$
于是 $\xi = \sum_{i \in I'} \xi_i$ 是各元素 $\xi_i \in G'(Y_i)$ 的像之和，
其中这些元素对应于 $a_i : X \to Y_i$ 和 $g \in G(X)$。因此
$\bigoplus G'(Y_i) \to G'(Y)$ 是满射。其为单射的（平凡）验证从略。

\medskip\noindent
由此，函子 $Y \mapsto G'(Y)^\vee$ 是上同调函子，并把直和送到直积。
因此，由 Brown 可表性（参见 Derived Categories, Proposition
\ref{derived-proposition-brown}），存在 $Y \in \Ob(\mathcal{D})$ 以及同构
$G'(Z)^\vee = \Hom(Z, Y)$；该同构关于 $Z$ 是函子性的。对于
$X \in \Ob(\mathcal{D}_c)$，有
$G'(X)^\vee = G(X)^\vee = (H(X)^\vee)^\vee = H(X)$，因为
$\dim_k H(X) < \infty$。证明完成。
\end{proof}

\begin{theorem}
\label{equiv-theorem-bondal-van-den-bergh}
\begin{reference}
在射影情形，这是 \cite[Theorem A.1]{BvdB}。
\end{reference}
设 $X$ 是域 $k$ 上的固有概形。设
$F : D_{perf}(\mathcal{O}_X)^{opp} \to \text{Vect}_k$ 是 $k$-线性上同调函子，
并满足
$$
\sum\nolimits_{n \in \mathbf{Z}} \dim_k F(E[n]) < \infty
$$
对所有 $E \in D_{perf}(\mathcal{O}_X)$ 成立。则
$F$ 同构于形如 $E \mapsto \Hom_X(E, K)$ 的函子，其中
$K \in D^b_{\textit{Coh}}(\mathcal{O}_X)$ 为某个对象。
\end{theorem}

\begin{proof}
导出范畴 $D_\QCoh(\mathcal{O}_X)$ 具有直和且是紧生成的，而
$D_{perf}(\mathcal{O}_X)$ 是紧对象所成的全子范畴；参见
Derived Categories of Schemes, Lemma
\ref{perfect-lemma-quasi-coherence-direct-sums},
Theorem \ref{perfect-theorem-bondal-van-den-Bergh} 及
Proposition \ref{perfect-proposition-compact-is-perfect}.
由 Lemma \ref{equiv-lemma-van-den-bergh}，可以假设 $F(E) = \Hom_X(E, K)$，其中
$K \in \Ob(D_\QCoh(\mathcal{O}_X))$ 为某个对象。于是由
Lemma \ref{equiv-lemma-characterize-dbcoh-projective} 可知，$K$ 属于
$D^b_{\textit{Coh}}(\mathcal{O}_X)$。
\end{proof}

\begin{lemma}
\label{equiv-lemma-homological-representable}
设 $X$ 是域 $k$ 上的正则固有概形。设
$G : D_{perf}(\mathcal{O}_X) \to \text{Vect}_k$ 是 $k$-线性同调函子，并满足
$$
\sum\nolimits_{n \in \mathbf{Z}} \dim_k G(E[n]) < \infty
$$
对所有 $E \in D_{perf}(\mathcal{O}_X)$ 成立。则 $G$ 同构于形如
$E \mapsto \Hom_X(K, E)$ 的函子，其中 $K \in D_{perf}(\mathcal{O}_X)$
为某个对象。
\end{lemma}

\begin{proof}
考察反变函子 $E \mapsto E^\vee$，其定义在
$D_{perf}(\mathcal{O}_X)$ 上；参见 Cohomology, Lemma
\ref{cohomology-lemma-dual-perfect-complex}。该函子是
$D_{perf}(\mathcal{O}_X)$ 的正合反自等价。因此，可将 Theorem
\ref{equiv-theorem-bondal-van-den-bergh} 应用于函子 $F(E) = G(E^\vee)$，得到
$K \in D_{perf}(\mathcal{O}_X)$，使得 $G(E^\vee) = \Hom_X(E, K)$。于是
$G(E) = \Hom_X(E^\vee, K) = \Hom_X(K^\vee, E)$，故取 $K^\vee$ 即可。
\end{proof}






\section{伴随函子的存在性}
\label{equiv-section-adjoints}

\noindent
作为 Bondal 与 van den Bergh 论文中结果的推论，得到下述伴随函子自动存在性。

\begin{lemma}
\label{equiv-lemma-always-right-adjoints}
设 $k$ 为域，设 $X$ 与 $Y$ 是 $k$ 上的固有概形。若 $X$ 正则，则任意
$k$-线性正合函子
$F : D_{perf}(\mathcal{O}_X) \to D_{perf}(\mathcal{O}_Y)$
都有正合右伴随与正合左伴随。
\end{lemma}

\begin{proof}
若伴随函子存在，则由非常一般的 Derived Categories, Lemma
\ref{derived-lemma-adjoint-is-exact}，它是正合函子。

\medskip\noindent
先证明右伴随的存在性。为此，只需证明对于
$M \in D_{perf}(\mathcal{O}_Y)$，反变函子
$K \mapsto \Hom_Y(F(K), M)$ 可表。该函子是反变、$k$-线性且上同调的。
因此，由 Theorem \ref{equiv-theorem-bondal-van-den-bergh}，只需证明
$$
\sum\nolimits_{i \in \mathbf{Z}} \dim_k \Ext^i_Y(F(K), M) < \infty
$$
即可。这由 Lemma \ref{equiv-lemma-finiteness} 得到。

\medskip\noindent
对于左伴随的存在性，论证相同，只需使用 Lemma
\ref{equiv-lemma-homological-representable} 代替 Theorem
\ref{equiv-theorem-bondal-van-den-bergh}。
\end{proof}






\section{Fourier--Mukai 函子}
\label{equiv-section-fourier-mukai}

\noindent
这些函子最早在 \cite{Mukai} 中引入。

\begin{definition}
\label{equiv-definition-fourier-mukai-functor}
设 $S$ 为概形，设 $X$ 与 $Y$ 为 $S$ 上的概形，并设
$K \in D(\mathcal{O}_{X \times_S Y})$。三角范畴之间的正合函子
$$
\Phi_K : D(\mathcal{O}_X) \longrightarrow D(\mathcal{O}_Y),\quad
M \longmapsto R\text{pr}_{2, *}(
L\text{pr}_1^*M \otimes_{\mathcal{O}_{X \times_S Y}}^\mathbf{L} K)
$$
称为一个 {\it Fourier--Mukai 函子}，而 $K$ 称为该函子的一个
{\it Fourier--Mukai 核}。此外：
\begin{enumerate}
\item 若 $\Phi_K$ 把 $D_\QCoh(\mathcal{O}_X)$ 送入 $D_\QCoh(\mathcal{O}_Y)$，
则所得正合函子
$\Phi_K : D_\QCoh(\mathcal{O}_X) \to D_\QCoh(\mathcal{O}_Y)$
称为 Fourier--Mukai 函子；
\item 若 $\Phi_K$ 把 $D_{perf}(\mathcal{O}_X)$ 送入
$D_{perf}(\mathcal{O}_Y)$，则所得正合函子
$\Phi_K : D_{perf}(\mathcal{O}_X) \to D_{perf}(\mathcal{O}_Y)$
称为 Fourier--Mukai 函子；
\item 若 $X$ 与 $Y$ 都是 Noether 概形，且 $\Phi_K$ 把
$D^b_{\textit{Coh}}(\mathcal{O}_X)$ 送入 $D^b_{\textit{Coh}}(\mathcal{O}_Y)$，
则所得正合函子
$\Phi_K : D^b_{\textit{Coh}}(\mathcal{O}_X) \to
D^b_{\textit{Coh}}(\mathcal{O}_Y)$
称为 Fourier--Mukai 函子。对于 $D_{\textit{Coh}}$、
$D^+_{\textit{Coh}}$、$D^-_{\textit{Coh}}$ 也作类似定义。
\end{enumerate}
\end{definition}

\begin{lemma}
\label{equiv-lemma-fourier-Mukai-QCoh}
设 $S$ 为概形，设 $X$ 与 $Y$ 为 $S$ 上的概形，并设
$K \in D(\mathcal{O}_{X \times_S Y})$。相应的 Fourier--Mukai 函子 $\Phi_K$
把 $D_\QCoh(\mathcal{O}_X)$ 送入 $D_\QCoh(\mathcal{O}_Y)$，只要 $K$ 属于
$D_\QCoh(\mathcal{O}_{X \times_S Y})$ 且 $X \to S$ 拟紧、拟分离。
\end{lemma}

\begin{proof}
这是因为：导出拉回保持 $D_\QCoh$（Derived Categories of Schemes, Lemma
\ref{perfect-lemma-quasi-coherence-pullback}）；导出张量积保持 $D_\QCoh$
（Derived Categories of Schemes, Lemma
\ref{perfect-lemma-quasi-coherence-tensor-product}）；投影
$\text{pr}_2 : X \times_S Y \to Y$ 拟紧且拟分离（Schemes, Lemmas
\ref{schemes-lemma-quasi-compact-preserved-base-change} 与
\ref{schemes-lemma-separated-permanence}）；沿拟分离且拟紧态射的总直像保持
$D_\QCoh$（Derived Categories of Schemes, Lemma
\ref{perfect-lemma-quasi-coherence-direct-image}）。
\end{proof}

\begin{lemma}
\label{equiv-lemma-compose-fourier-mukai}
设 $S$ 为概形，设 $X, Y, Z$ 为 $S$ 上的概形。假设 $X \to S$、
$Y \to S$ 与 $Z \to S$ 都拟紧且拟分离。设
$K \in D_\QCoh(\mathcal{O}_{X \times_S Y})$，并设
$K' \in D_\QCoh(\mathcal{O}_{Y \times_S Z})$。考虑 Fourier--Mukai 函子
$\Phi_K : D_\QCoh(\mathcal{O}_X) \to D_\QCoh(\mathcal{O}_Y)$
与 $\Phi_{K'} : D_\QCoh(\mathcal{O}_Y) \to D_\QCoh(\mathcal{O}_Z)$。若
$X$ 与 $Z$ 在 $S$ 上 Tor 独立，且 $Y \to S$ 平坦，则
$$
\Phi_{K'} \circ \Phi_K = \Phi_{K''} :
D_\QCoh(\mathcal{O}_X)
\longrightarrow
D_\QCoh(\mathcal{O}_Z)
$$
，其中
$$
K'' = R\text{pr}_{13, *}(
L\text{pr}_{12}^*K
\otimes_{\mathcal{O}_{X \times_S Y \times_S Z}}^\mathbf{L}
L\text{pr}_{23}^*K')
$$
属于 $D_\QCoh(\mathcal{O}_{X \times_S Z})$。
\end{lemma}

\begin{proof}
该陈述由 Lemma \ref{equiv-lemma-fourier-Mukai-QCoh} 可知是有意义的。下文将使用
Derived Categories of Schemes, Lemmas
\ref{perfect-lemma-quasi-coherence-pullback},
\ref{perfect-lemma-quasi-coherence-tensor-product} 以及
\ref{perfect-lemma-quasi-coherence-direct-image}
以及 Schemes, Lemmas
\ref{schemes-lemma-quasi-compact-preserved-base-change} 与
\ref{schemes-lemma-separated-permanence}
，而不再说明。由 Derived Categories of Schemes, Lemma
\ref{perfect-lemma-flat-base-change-tor-independent}
可知 $X \times_S Y$ 与 $Y \times_S Z$ 在 $Y$ 上 Tor 独立。这意味着对于
下列笛卡尔图有基变换：
$$
\xymatrix{
X \times_S Y \times_S Z \ar[d] \ar[r] &
Y \times_S Z \ar[d]^{p^{YZ}_Y} \\
X \times_S Y \ar[r]^{p^{XY}_Y} & Y
}
$$
这里所指的是具有拟相干上同调层的复形；参见 Derived Categories of Schemes,
Lemma \ref{perfect-lemma-compare-base-change}。简记 $p^* = Lp^*$、
$p_* = Rp_*$ 及 $\otimes = \otimes^\mathbf{L}$，则对于
$M \in D_\QCoh(\mathcal{O}_X)$ 有等式链
\begin{align*}
\Phi_{K'}(\Phi_K(M))
& =
p^{YZ}_{Z, *}(p^{YZ, *}_Y p^{XY}_{Y, *}(p^{XY, *}_X M \otimes K) \otimes K') \\
& =
p^{YZ}_{Z, *}(\text{pr}_{23, *} \text{pr}_{12}^*(p^{XY, *}_X M \otimes K)
\otimes K') \\
& =
p^{YZ}_{Z, *}(\text{pr}_{23, *}(\text{pr}_1^*M \otimes \text{pr}_{12}^*K)
\otimes K') \\
& =
p^{YZ}_{Z, *}(\text{pr}_{23, *}(\text{pr}_1^*M \otimes \text{pr}_{12}^*K
\otimes \text{pr}_{23}^*K')) \\
& =
\text{pr}_{3, *}(\text{pr}_1^*M \otimes \text{pr}_{12}^*K
\otimes \text{pr}_{23}^*K') \\
& =
p^{XZ}_{Z, *}\text{pr}_{13, *}(\text{pr}_1^*M \otimes \text{pr}_{12}^*K
\otimes \text{pr}_{23}^*K') \\
& =
p^{XZ}_{Z, *} (p^{XZ, *}_X M \otimes \text{pr}_{13, *}(\text{pr}_{12}^*K
\otimes \text{pr}_{23}^*K'))
\end{align*}
，正如所需。第二个等式使用了上述基变换，而第 $4$ 个等式与最后一个等式
使用了 Derived Categories of Schemes, Lemma
\ref{perfect-lemma-cohomology-base-change}。
\end{proof}

\begin{lemma}
\label{equiv-lemma-fourier-mukai}
设 $S$ 为概形，设 $X$ 与 $Y$ 为 $S$ 上的概形，并设
$K \in D(\mathcal{O}_{X \times_S Y})$。若下列条件中至少一个成立，则相应的
Fourier--Mukai 函子 $\Phi_K$ 把 $D_{perf}(\mathcal{O}_X)$ 送入
$D_{perf}(\mathcal{O}_Y)$：
\begin{enumerate}
\item $S$ 是 Noether 概形，$X \to S$ 与 $Y \to S$ 是有限型的，
$K \in D^b_{\textit{Coh}}(\mathcal{O}_{X \times_S Y})$，$H^i(K)$ 的支撑在
$Y$ 上固有，其中 $i$ 任取，并且 $K$ 作为
$D(\text{pr}_2^{-1}\mathcal{O}_Y)$ 的对象具有有限 Tor 维数；
\item $X \to S$ 是有限表示的，并且 $K$ 可由有界复形
$\mathcal{K}^\bullet$ 表示；该复形由有限表示
$\mathcal{O}_{X \times_S Y}$-模组成，在 $Y$ 上平坦且其支撑在 $Y$ 上固有；
\item $X \to S$ 是有限表示的固有平坦态射，且 $K$ 完美；
\item $S$ 是 Noether 概形，$X \to S$ 平坦且固有，并且 $K$ 完美；
\item $X \to S$ 是有限表示的固有平坦态射，且 $K$ 是 $Y$-完美的；
\item $S$ 是 Noether 概形，$X \to S$ 平坦且固有，并且 $K$ 是
$Y$-完美的。
\end{enumerate}
\end{lemma}

\begin{proof}
若 $M$ 在 $X$ 上完美，则 $L\text{pr}_1^*M$ 在 $X \times_S Y$ 上完美；
参见 Cohomology, Lemma \ref{cohomology-lemma-perfect-pullback}。下文将使用
此事实而不再说明。还要使用如下事实：若 $X \to S$ 是有限型、固有、
平坦或有限表示的，则基变换 $\text{pr}_2 : X \times_S Y \to Y$ 也具有
相应性质；参见 Morphisms, Lemmas
\ref{morphisms-lemma-base-change-finite-type},
\ref{morphisms-lemma-base-change-proper},
\ref{morphisms-lemma-base-change-flat} 与
\ref{morphisms-lemma-base-change-finite-presentation}.

\medskip\noindent
情形 (1) 由
Derived Categories of Schemes, Lemma \ref{perfect-lemma-perfect-direct-image}
与 Derived Categories of Schemes, Lemma
\ref{perfect-lemma-perfect-on-noetherian} 结合得到。

\medskip\noindent
情形 (2) 由 Derived Categories of Schemes, Lemma
\ref{perfect-lemma-base-change-tensor-perfect} 得到。

\medskip\noindent
情形 (3) 由 Derived Categories of Schemes, Lemma
\ref{perfect-lemma-flat-proper-perfect-direct-image-general} 得到。

\medskip\noindent
情形 (4) 由情形 (3) 以及如下事实得到：由 Morphisms, Lemma
\ref{morphisms-lemma-noetherian-finite-type-finite-presentation}，Noether 概形之间的
有限型态射是有限表示的。

\medskip\noindent
情形 (5) 由 Derived Categories of Schemes, Lemma
\ref{perfect-lemma-derived-pushforward-rel-perfect} 与 Derived Categories of Schemes,
Lemma \ref{perfect-lemma-perfect-relatively-perfect} 结合得到。

\medskip\noindent
情形 (6) 由情形 (5) 得到，其方式与情形 (4) 由情形 (3) 得到的方式相同。
\end{proof}

\begin{lemma}
\label{equiv-lemma-fourier-mukai-Coh}
设 $S$ 为 Noether 概形，设 $X$ 与 $Y$ 为 $S$ 上的有限型概形，并设
$K \in D^b_{\textit{Coh}}(\mathcal{O}_{X \times_S Y})$。若下列条件中至少
一个成立，则相应的 Fourier--Mukai 函子 $\Phi_K$ 把
$D^b_{\textit{Coh}}(\mathcal{O}_X)$ 送入 $D^b_{\textit{Coh}}(\mathcal{O}_Y)$：
\begin{enumerate}
\item $H^i(K)$ 的支撑在 $Y$ 上固有，其中 $i$ 任取，并且 $K$ 作为
$D(\text{pr}_1^{-1}\mathcal{O}_X)$ 的对象具有有限 Tor 维数；
\item $K$ 可由有界复形 $\mathcal{K}^\bullet$ 表示；该复形由相干
$\mathcal{O}_{X \times_S Y}$-模组成，在 $X$ 上平坦且其支撑在 $Y$ 上固有；
\item $H^i(K)$ 的支撑在 $Y$ 上固有，其中 $i$ 任取，并且 $X$ 是正则概形；
\item $K$ 完美，$H^i(K)$ 的支撑在 $Y$ 上固有，其中 $i$ 任取，并且
$Y \to S$ 平坦。
\end{enumerate}
此外，在每种情形下，若 $X \to S$ 固有，则支撑条件自动成立。
\end{lemma}

\begin{proof}
设 $M$ 为 $D^b_{\textit{Coh}}(\mathcal{O}_X)$ 的对象。在每种情形下，都将使用
Derived Categories of Schemes, Lemma \ref{perfect-lemma-direct-image-coherent}
来证明
$$
\Phi_K(M) = R\text{pr}_{2, *}(
L\text{pr}_1^*M
\otimes_{\mathcal{O}_{X \times_S Y}}^\mathbf{L}
K)
$$
属于 $D^b_{\textit{Coh}}(\mathcal{O}_Y)$。导出张量积
$L\text{pr}_1^*M \otimes_{\mathcal{O}_{X \times_S Y}}^\mathbf{L} K$
是 $D(\mathcal{O}_{X \times_S Y})$ 的伪凝聚对象（由 Cohomology, Lemma
\ref{cohomology-lemma-pseudo-coherent-pullback}、Derived Categories of Schemes,
Lemma \ref{perfect-lemma-identify-pseudo-coherent-noetherian} 以及 Cohomology,
Lemma \ref{cohomology-lemma-tensor-pseudo-coherent}），因而具有相干上同调层
（再次由 Derived Categories of Schemes, Lemma
\ref{perfect-lemma-identify-pseudo-coherent-noetherian}）。在每种情形下，上同调层
$H^i(L\text{pr}_1^*M \otimes_{\mathcal{O}_{X \times_S Y}}^\mathbf{L} K)$
的支撑在 $Y$ 上固有，因为这些支撑包含于各 $H^i(K)$ 的支撑之并。因此，
在每种情形下，只需证明这个张量积下有界。

\medskip\noindent
情形 (1)。由 Cohomology, Lemma
\ref{cohomology-lemma-variant-derived-pullback}，有
$$
L\text{pr}_1^*M
\otimes_{\mathcal{O}_{X \times_S Y}}^\mathbf{L}
K
\cong
\text{pr}_1^{-1}M
\otimes_{\text{pr}_1^{-1}\mathcal{O}_X}^\mathbf{L}
K
$$
，其中记号含义显然。因此，Tor 维数的假设与 $M$ 仅有有限个非零上同调层
这一事实给出了所需的界。

\medskip\noindent
情形 (2) 成立，因为这里的假设蕴含 $K$ 作为
$D(\text{pr}_1^{-1}\mathcal{O}_X)$ 的对象具有有限 Tor 维数，故可应用上一段的论证。

\medskip\noindent
在情形 (3) 中，$K$ 作为 $D(\text{pr}_1^{-1}\mathcal{O}_X)$ 的对象也具有
有限 Tor 维数。事实上，选取仿射开集 $U = \Spec(A)$ 与
$V = \Spec(B)$，它们分别属于 $X$ 与 $Y$，并映入仿射开集
$W = \Spec(R)$，后者属于 $S$。于是
$K|_{U \times V}$ 由有限 $A \otimes_R B$-模组成的有界复形
$M^\bullet$ 给出。由于 $A$ 是有限维正则环，可知每个 $M^i$ 作为 $A$-模
具有有限投射维数（Algebra, Lemma
\ref{algebra-lemma-finite-gl-dim-finite-dim-regular}），因而作为 $A$-模具有
有限 Tor 维数。因此 $M^\bullet$ 作为 $A$-模复形具有有限 Tor 维数
（More on Algebra, Lemma
\ref{more-algebra-lemma-complex-finite-tor-dimension-modules}）。由于
$X \times Y$ 拟紧，存在 $[a, b]$，使得对于每个点 $z \in X \times Y$，茎
$K_z$ 的 Tor 振幅落在 $[a, b]$ 中（相对于
$\mathcal{O}_{X, \text{pr}_1(z)}$）。
这蕴含 $K$ 作为 $D(\text{pr}_1^{-1}\mathcal{O}_X)$ 的对象具有有界 Tor 维数；
参见 Cohomology, Lemma \ref{cohomology-lemma-tor-amplitude-stalk}。按前两段
的方式即可得到结论。

\medskip\noindent
情形 (4)。沿用上述记号，环映射 $R \to B$ 平坦。因此环映射
$A \to A \otimes_R B$ 平坦，故任意投射 $A \otimes_R B$-模都是
$A$-平坦的。于是，任意完美 $A \otimes_R B$-模复形作为 $A$-模复形
都具有有限 Tor 维数，结论同前。
\end{proof}

\begin{example}
\label{equiv-example-diagonal-fourier-mukai}
设 $X \to S$ 为概形之间的分离态射。则对角态射
$\Delta : X \to X \times_S X$ 是闭浸入，因而
$\mathcal{O}_\Delta = \Delta_*\mathcal{O}_X = R\Delta_*\mathcal{O}_X$
是有限型拟相干 $\mathcal{O}_{X \times_S X}$-模，并且在 $X$ 上平坦
（关于任一投影）。在此情形下，Fourier--Mukai 函子
$\Phi_{\mathcal{O}_\Delta}$ 等于恒等函子。事实上，对任意
$M \in D(\mathcal{O}_X)$，有
\begin{align*}
L\text{pr}_1^*M \otimes_{\mathcal{O}_{X \times_S X}}^\mathbf{L}
\mathcal{O}_\Delta
& =
L\text{pr}_1^*M \otimes_{\mathcal{O}_{X \times_S X}}^\mathbf{L}
R\Delta_*\mathcal{O}_X \\
& =
R\Delta_*(
L\Delta^*L\text{pr}_1^*M \otimes_{\mathcal{O}_X}^\mathbf{L} \mathcal{O}_X) \\
& =
R\Delta_*(M)
\end{align*}
第一个等式已在上文讨论。第二个等式是 Cohomology, Lemma
\ref{cohomology-lemma-projection-formula-closed-immersion}.
第三个等式成立是因为 $\text{pr}_1 \circ \Delta = \text{id}_X$，并且有
Cohomology, Lemma \ref{cohomology-lemma-derived-pullback-composition}.
若将其推前到 $X$，所用函子为 $R\text{pr}_{2, *}$，则得到 $M$；这是由
Cohomology, Lemma \ref{cohomology-lemma-derived-pushforward-composition}
以及 $\text{pr}_2 \circ \Delta = \text{id}_X$ 这一事实推出的。
\end{example}

\begin{lemma}
\label{equiv-lemma-fourier-mukai-right-adjoint}
\begin{reference}
参见 \cite{Rizzardo} 中的讨论。
\end{reference}
设 $X \to S$ 与 $Y \to S$ 为拟紧拟分离概形之间的态射。设
$\Phi : D_\QCoh(\mathcal{O}_X) \to D_\QCoh(\mathcal{O}_Y)$
为具有伪凝聚核的 Fourier--Mukai 函子，其中
$K \in D_\QCoh(\mathcal{O}_{X \times_S Y})$。设
$a : D_\QCoh(\mathcal{O}_Y) \to  D_\QCoh(\mathcal{O}_{X \times_S Y})$
为 $R\text{pr}_{2, *}$ 的右伴随；参见 Duality for Schemes, Lemma
\ref{duality-lemma-twisted-inverse-image}。记
$$
K' = (Y \times_S X \to X \times_S Y)^*
R\SheafHom_{\mathcal{O}_{X \times_S Y}}(K, a(\mathcal{O}_Y)) \in
D_\QCoh(\mathcal{O}_{Y \times_S X})
$$
并以 $\Phi' : D_\QCoh(\mathcal{O}_Y) \to D_\QCoh(\mathcal{O}_X)$
记相应的 Fourier--Mukai 变换。存在典范映射
$$
\Hom_X(M, \Phi'(N)) \longrightarrow \Hom_Y(\Phi(M), N)
$$
它关于 $M$（作为 $D_\QCoh(\mathcal{O}_X)$ 的对象）以及
$N$（作为 $D_\QCoh(\mathcal{O}_Y)$ 的对象）都具有函子性，并且在下列任一情形下为同构：
\begin{enumerate}
\item $N$ 完美；或
\item $K$ 完美，且 $X \to S$ 是有限表示的固有平坦态射。
\end{enumerate}
\end{lemma}

\begin{proof}
由 Lemma \ref{equiv-lemma-fourier-Mukai-QCoh} 得到陈述中的函子 $\Phi$。注意，由
Duality for Schemes, Lemma
\ref{duality-lemma-twisted-inverse-image-bounded-below}，
$a(\mathcal{O}_Y)$ 属于 $D^+_\QCoh(\mathcal{O}_{X \times_S Y})$。
因此，当 $K$ 伪凝聚时，由 Derived Categories of Schemes, Lemma
\ref{perfect-lemma-quasi-coherence-internal-hom} 有
$K' \in D_\QCoh(\mathcal{O}_{Y \times_S X})$，从而得到所述的 $\Phi'$。

\medskip\noindent
简记
$\otimes^\mathbf{L} = \otimes_{\mathcal{O}_{X \times_S Y}}^\mathbf{L}$
以及
$\SheafHom = R\SheafHom_{\mathcal{O}_{X \times_S Y}}$.
设 $M$ 属于 $D_\QCoh(\mathcal{O}_X)$，并设
$N$ 属于 $D_\QCoh(\mathcal{O}_Y)$。则
\begin{align*}
\Hom_Y(\Phi(M), N)
& =
\Hom_Y(R\text{pr}_{2, *}(L\text{pr}_1^*M \otimes^\mathbf{L} K), N) \\
& =
\Hom_{X \times_S Y}(L\text{pr}_1^*M \otimes^\mathbf{L} K, a(N)) \\
& =
\Hom_{X \times_S Y}(L\text{pr}_1^*M,
R\SheafHom(K, a(N))) \\
& =
\Hom_X(M, R\text{pr}_{1, *}R\SheafHom(K, a(N)))
\end{align*}
其中使用了 Cohomology, Lemmas \ref{cohomology-lemma-internal-hom}
与 \ref{cohomology-lemma-adjoint}。存在典范映射
$$
L\text{pr}_2^*N \otimes^\mathbf{L} R\SheafHom(K, a(\mathcal{O}_Y))
\xrightarrow{\alpha}
R\SheafHom(K, L\text{pr}_2^*N \otimes^\mathbf{L} a(\mathcal{O}_Y))
\xrightarrow{\beta}
R\SheafHom(K, a(N))
$$
这里 $\alpha$ 是 Cohomology, Lemma
\ref{cohomology-lemma-internal-hom-diagonal-better} 中的映射，而 $\beta$ 是
Duality for Schemes, Equation
(\ref{duality-equation-compare-with-pullback}) 中的映射。合并所有这些箭头，
便得到本引理陈述中所显示的函子性箭头。

\medskip\noindent
箭头 $\alpha$ 由 Derived Categories of Schemes, Lemma
\ref{perfect-lemma-internal-hom-evaluate-tensor-isomorphism} 可知为同构，只要
$K$ 或 $N$ 完美。箭头 $\beta$ 在 $N$ 完美时由 Duality for Schemes, Lemma
\ref{duality-lemma-compare-with-pullback-perfect} 可知为同构；更一般地，若
$X \to S$ 是有限表示的平坦固有态射，则由 Duality for Schemes,
Lemma \ref{duality-lemma-compare-with-pullback-flat-proper}，它也为同构。
\end{proof}

\begin{lemma}
\label{equiv-lemma-fourier-mukai-left-adjoint}
\begin{reference}
参见 \cite{Rizzardo} 中的讨论。
\end{reference}
设 $S$ 为 Noether 概形。设 $Y \to S$ 为平坦固有 Gorenstein 态射，并设
$X \to S$ 为有限型态射。记 $\omega^\bullet_{Y/S}$ 为 $Y$ 在 $S$ 上的
相对对偶化复形。设
$\Phi : D_\QCoh(\mathcal{O}_X) \to D_\QCoh(\mathcal{O}_Y)$
为具有完美核的 Fourier--Mukai 函子，其中
$K \in D_\QCoh(\mathcal{O}_{X \times_S Y})$。记
$$
K' = (Y \times_S X \to X \times_S Y)^*(K^\vee
\otimes_{\mathcal{O}_{X \times_S Y}}^\mathbf{L}
L\text{pr}_2^*\omega^\bullet_{Y/S})
\in
D_\QCoh(\mathcal{O}_{Y \times_S X})
$$
并以 $\Phi' : D_\QCoh(\mathcal{O}_Y) \to D_\QCoh(\mathcal{O}_X)$
记相应的 Fourier--Mukai 变换。存在典范同构
$$
\Hom_Y(N, \Phi(M)) \longrightarrow \Hom_X(\Phi'(N), M)
$$
它关于 $M$（作为 $D_\QCoh(\mathcal{O}_X)$ 的对象）以及
$N$（作为 $D_\QCoh(\mathcal{O}_Y)$ 的对象）都具有函子性。
\end{lemma}

\begin{proof}
由 Lemma \ref{equiv-lemma-fourier-Mukai-QCoh} 得到陈述中的函子 $\Phi$。

\medskip\noindent
注意，在当前设定下，相对对偶化复形的形成与基变换交换；参见
Duality for Schemes, Remark
\ref{duality-remark-relative-dualizing-complex}。因此
$L\text{pr}_2^*\omega^\bullet_{Y/S} = \omega^\bullet_{X \times_S Y/X}$。
此外，$\omega^\bullet_{Y/S}$ 是导出范畴中的可逆对象；参见
Duality for Schemes, Lemma
\ref{duality-lemma-affine-flat-Noetherian-gorenstein}；因而它尤其完美。

\medskip\noindent
为真正证明本引理，我们采用一个取巧办法。具体地，将 $X$ 与 $Y$ 的角色以及
$K$ 与 $K'$ 的角色互换后，所得数据符合
Lemma \ref{equiv-lemma-fourier-mukai-right-adjoint} 的情形，从而得到结论。显然，
$K'$ 作为完美对象的张量积仍是完美的，故 Lemma
\ref{equiv-lemma-fourier-mukai-right-adjoint} 中的讨论适用于它。要说明将
Lemma \ref{equiv-lemma-fourier-mukai-right-adjoint} 的过程应用于 $K'$（在
$Y \times_S X$ 上）会产生一个同构于 $K$ 的复形，只需说明下式
（略去细节）：
$$
R\SheafHom(R\SheafHom(K, \omega^\bullet_{X \times_S Y/X}),
\omega^\bullet_{X \times_S Y/X}) = K
$$
这是显然的，因为 $K$ 完美而 $\omega^\bullet_{X \times_S Y/X}$ 可逆；
细节从略。因此 Lemma \ref{equiv-lemma-fourier-mukai-right-adjoint} 给出映射
$$
\Hom_Y(N, \Phi(M)) \longrightarrow \Hom_X(\Phi'(N), M)
$$
它关于 $M$（作为 $D_\QCoh(\mathcal{O}_X)$ 的对象）以及
$N$（作为 $D_\QCoh(\mathcal{O}_Y)$ 的对象）都具有函子性；由于
$K'$ 完美，该映射为同构。证明完毕。
\end{proof}

\begin{lemma}
\label{equiv-lemma-fourier-mukai-flat-proper-over-noetherian}
设 $S$ 为 Noether 概形。
\begin{enumerate}
\item 对于 $X$、$Y$（二者在 $S$ 上固有且平坦），若 $K$ 是
$D_{perf}(\mathcal{O}_{X \times_S Y})$ 的对象，则得到 Fourier--Mukai 函子
$\Phi_K : D_{perf}(\mathcal{O}_X) \to D_{perf}(\mathcal{O}_Y)$。
\item 对于 $X$、$Y$、$Z$（三者在 $S$ 上固有且平坦），以及 $K \in
D_{perf}(\mathcal{O}_{X \times_S Y})$，$K' \in
D_{perf}(\mathcal{O}_{Y \times_S Z})$，复合
$\Phi_{K'} \circ \Phi_K : D_{perf}(\mathcal{O}_X) \to D_{perf}(\mathcal{O}_Z)$
等于 $\Phi_{K''}$，其中 $K'' \in D_{perf}(\mathcal{O}_{X \times_S Z})$
按 Lemma \ref{equiv-lemma-compose-fourier-mukai} 中的方式计算；
\item 对于如 (1) 中的 $X$、$Y$、$K$、$\Phi_K$，若 $X \to S$ 是 Gorenstein
态射，则
$\Phi_{K'} : D_{perf}(\mathcal{O}_Y) \to D_{perf}(\mathcal{O}_X)$ 是
$\Phi_K$ 的右伴随，其中 $K' \in D_{perf}(\mathcal{O}_{Y \times_S X})$
是 $L\text{pr}_1^*\omega_{X/S}^\bullet
\otimes_{\mathcal{O}_{X \times_S Y}}^\mathbf{L} K^\vee$ 沿
$Y \times_S X \to X \times_S Y$ 的拉回；
\item 对于如 (1) 中的 $X$、$Y$、$K$、$\Phi_K$，若 $Y \to S$ 是 Gorenstein
态射，则
$\Phi_{K''} : D_{perf}(\mathcal{O}_Y) \to D_{perf}(\mathcal{O}_X)$ 是
$\Phi_K$ 的左伴随，其中 $K'' \in D_{perf}(\mathcal{O}_{Y \times_S X})$
是 $L\text{pr}_2^*\omega_{Y/S}^\bullet
\otimes_{\mathcal{O}_{X \times_S Y}}^\mathbf{L} K^\vee$ 沿
$Y \times_S X \to X \times_S Y$ 的拉回。
\end{enumerate}
\end{lemma}

\begin{proof}
第 (1) 部分立即由 Lemma \ref{equiv-lemma-fourier-mukai} 的第 (4) 部分得到。

\medskip\noindent
第 (2) 部分由 Lemma \ref{equiv-lemma-compose-fourier-mukai} 以及下述事实得到：
$K'' = R\text{pr}_{13, *}(
L\text{pr}_{12}^*K
\otimes_{\mathcal{O}_{X \times_S Y \times_S Z}}^\mathbf{L}
L\text{pr}_{23}^*K')$ 是完美的；例如这可由 Derived Categories of Schemes,
Lemma
\ref{perfect-lemma-flat-proper-perfect-direct-image}.

\medskip\noindent
第 (3) 部分在所有具有拟相干上同调层的复形上的伴随性，由
Lemma \ref{equiv-lemma-fourier-mukai-right-adjoint} 得到；其中 $K'$ 等于
$R\SheafHom_{\mathcal{O}_{X \times_S Y}}(K, a(\mathcal{O}_Y))$
沿 $Y \times_S X \to X \times_S Y$ 的拉回，而 $a$ 是
$R\text{pr}_{2, *} : D_\QCoh(\mathcal{O}_{X \times_S Y}) \to
D_\QCoh(\mathcal{O}_Y)$ 的右伴随。以 $f : X \to S$ 记 $X$ 的结构态射。
由于 $f$ 固有，函子
$f^! : D_\QCoh^+(\mathcal{O}_S) \to D_\QCoh^+(\mathcal{O}_X)$
是下述右伴随在 $D_\QCoh^+(\mathcal{O}_S)$ 上的限制：
$Rf_* : D_\QCoh(\mathcal{O}_X) \to D_\QCoh(\mathcal{O}_S)$；参见
Duality for Schemes, Section \ref{duality-section-upper-shriek}。因此，
Duality for Schemes, Remark
\ref{duality-remark-relative-dualizing-complex}
中定义的相对对偶化复形 $\omega_{X/S}^\bullet$ 等于
$\omega_{X/S}^\bullet = f^!\mathcal{O}_S$。由于相对对偶化复形的形成与
基变换交换（参见 Duality for Schemes, Remark
\ref{duality-remark-relative-dualizing-complex}），可知
$a(\mathcal{O}_Y) = L\text{pr}_1^*\omega_{X/S}^\bullet$.
因此
$$
R\SheafHom_{\mathcal{O}_{X \times_S Y}}(K, a(\mathcal{O}_Y))
\cong
L\text{pr}_1^*\omega_{X/S}^\bullet
\otimes_{\mathcal{O}_{X \times_S Y}}^\mathbf{L} K^\vee
$$
，由 Cohomology, Lemma \ref{cohomology-lemma-dual-perfect-complex}。
最后，由于假设 $X \to S$ 是 Gorenstein 态射，相对对偶化复形可逆；这由
Duality for Schemes, Lemma
\ref{duality-lemma-affine-flat-Noetherian-gorenstein}.
于是 $\omega_{X/S}^\bullet$ 完美
(Cohomology, Lemma \ref{cohomology-lemma-invertible-derived})
，从而 $K'$ 完美。因此 $\Phi_{K'}$ 确实把
$D_{perf}(\mathcal{O}_Y)$ 送入 $D_{perf}(\mathcal{O}_X)$，第 (3) 部分得证。

\medskip\noindent
第 (4) 部分的证明与第 (3) 部分相同，只需使用
Lemma \ref{equiv-lemma-fourier-mukai-left-adjoint} 代替
Lemma \ref{equiv-lemma-fourier-mukai-right-adjoint}。
\end{proof}














\section{解消与界}
\label{equiv-section-tricks-smooth}

\noindent
光滑固有概形的对角具有良好的解消。

\begin{lemma}
\label{equiv-lemma-on-product}
设 $R$ 为 Noether 环。设 $X$、$Y$ 为 $R$ 上具有解消性质的有限型概形。
对于任意相干 $\mathcal{O}_{X \times_R Y}$-模 $\mathcal{F}$，存在满射
$\mathcal{E} \boxtimes \mathcal{G} \to \mathcal{F}$，其中
$\mathcal{E}$ 是有限局部自由 $\mathcal{O}_X$-模，而
$\mathcal{G}$ 是有限局部自由 $\mathcal{O}_Y$-模。
\end{lemma}

\begin{proof}
设 $U \subset X$ 与 $V \subset Y$ 为仿射开子概形。设
$\mathcal{I} \subset \mathcal{O}_X$ 为 $X \setminus U$ 上约化诱导闭子概形
结构的理想层。类似地，设 $\mathcal{I}' \subset \mathcal{O}_Y$ 为
$Y \setminus V$ 上约化诱导闭子概形结构的理想层。则理想层
$$
\mathcal{J} = \Im(\text{pr}_1^*\mathcal{I} \otimes_{\mathcal{O}_{X \times_R Y}}
\text{pr}_2^*\mathcal{I}' \to \mathcal{O}_{X \times_R Y})
$$
满足 $V(\mathcal{J}) = X \times_R Y \setminus U \times_R V$。对于任意截面
$s \in \mathcal{F}(U \times_R V)$，可找到整数 $n > 0$ 与映射
$\mathcal{J}^n \to \mathcal{F}$，其在 $U \times_R V$ 上的限制给出 $s$；
参见 Cohomology of Schemes, Lemma \ref{coherent-lemma-homs-over-open}。
由假设，可选取满射 $\mathcal{E} \to \mathcal{I}$ 与
$\mathcal{G} \to \mathcal{I}'$。它们相应地产生满射
$$
\mathcal{E} \boxtimes \mathcal{G} \to \mathcal{J}
\quad\text{且}\quad
\mathcal{E}^{\otimes n} \boxtimes \mathcal{G}^{\otimes n} \to \mathcal{J}^n
$$
，从而产生映射
$\mathcal{E}^{\otimes n} \boxtimes \mathcal{G}^{\otimes n} \to \mathcal{F}$
，其像包含截面 $s$（在 $U \times_R V$ 上）。由于 $X \times_R Y$
可由有限多个形如 $U \times_R V$ 的仿射开集覆盖，并且
$\mathcal{F}|_{U \times_R V}$ 由有限多个截面生成（Properties, Lemma
\ref{properties-lemma-finite-type-module}），故存在满射
$$
\bigoplus\nolimits_{j = 1, \ldots, N}
\mathcal{E}_j^{\otimes n_j} \boxtimes \mathcal{G}_j^{\otimes n_j}
\to \mathcal{F}
$$
其中 $\mathcal{E}_j$ 在 $X$ 上有限局部自由，而
$\mathcal{G}_j$ 在 $Y$ 上有限局部自由。令
$\mathcal{E} = \bigoplus \mathcal{E}_j^{\otimes n_j}$ 且
$\mathcal{G} = \bigoplus \mathcal{G}_j^{\otimes n_j}$，便得到本引理。
\end{proof}

\begin{lemma}
\label{equiv-lemma-on-product-general}
设 $R$ 为环。设 $X$、$Y$ 为 $R$ 上具有解消性质的拟紧拟分离概形。
对于任意有限型拟相干 $\mathcal{O}_{X \times_R Y}$-模 $\mathcal{F}$，
存在满射 $\mathcal{E} \boxtimes \mathcal{G} \to \mathcal{F}$，其中
$\mathcal{E}$ 是有限局部自由 $\mathcal{O}_X$-模，而
$\mathcal{G}$ 是有限局部自由 $\mathcal{O}_Y$-模。
\end{lemma}

\begin{proof}
由极限论证从 Lemma \ref{equiv-lemma-on-product} 得到。建议读者略过此证明。
由于 $X \times_R Y$ 是 $X \times_\mathbf{Z} Y$ 的闭子概形，把 $R$ 替换为
$\mathbf{Z}$ 并无妨碍。由 Properties, Lemma
\ref{properties-lemma-finite-directed-colimit-surjective-maps}，可将
$\mathcal{F}$ 写成某个有限表示 $\mathcal{O}_{X \times_R Y}$-模的商。因此可假设
$\mathcal{F}$ 是有限表示的。接着可写成 $X = \lim X_i$，其中 $X_i$ 在
$\mathbf{Z}$ 上有限表示；类似地，$Y = \lim Y_j$；参见 Limits, Proposition
\ref{limits-proposition-approximate}。于是 $\mathcal{F}$ 可下降为
$\mathcal{F}_{ij}$，后者位于某个 $X_i \times_R Y_j$ 上
（Limits, Lemma \ref{limits-lemma-descend-modules-finite-presentation}），
而具有解消性质这一性质也可下降（Derived Categories of Schemes, Lemma
\ref{perfect-lemma-resolution-property-descends}）。然后把
Lemma \ref{equiv-lemma-on-product} 应用于 $\mathcal{F}_{ij}$，再拉回即可。
\end{proof}

\begin{lemma}
\label{equiv-lemma-diagonal-resolution}
设 $R$ 为 Noether 环。设 $X$ 为 $R$ 上具有解消性质的分离有限型概形。令
$\mathcal{O}_\Delta = \Delta_*(\mathcal{O}_X)$，其中
$\Delta : X \to X \times_R X$ 是 $X/k$ 的对角。存在解消
$$
\ldots \to
\mathcal{E}_2 \boxtimes \mathcal{G}_2 \to
\mathcal{E}_1 \boxtimes \mathcal{G}_1 \to
\mathcal{E}_0 \boxtimes \mathcal{G}_0 \to
\mathcal{O}_\Delta \to 0
$$
其中每个 $\mathcal{E}_i$ 与 $\mathcal{G}_i$ 都是有限局部自由
$\mathcal{O}_X$-模。
\end{lemma}

\begin{proof}
由于 $X$ 分离，对角态射 $\Delta$ 是闭浸入，因而 $\mathcal{O}_\Delta$ 是
相干 $\mathcal{O}_{X \times_R X}$-模（Cohomology of Schemes, Lemma
\ref{coherent-lemma-i-star-equivalence}）。故本引理立即由
Lemma \ref{equiv-lemma-on-product} 得到。
\end{proof}

\begin{lemma}
\label{equiv-lemma-Ext-0-regular}
设 $X$ 为维数 $d < \infty$ 的正则 Noether 概形。则
\begin{enumerate}
\item 对于 $\mathcal{F}$、$\mathcal{G}$ 这两个相干 $\mathcal{O}_X$-模，
有 $\Ext^n_X(\mathcal{F}, \mathcal{G}) = 0$，只要 $n > d$；并且
\item 对于 $K, L \in D^b_{\textit{Coh}}(\mathcal{O}_X)$ 与
$a \in \mathbf{Z}$，若 $H^i(K) = 0$ 对所有 $i < a + d$ 成立，且
$H^i(L) = 0$ 对所有 $i \geq a$ 成立，则 $\Hom_X(K, L) = 0$。
\end{enumerate}
\end{lemma}

\begin{proof}
为证明 (1)，使用谱序列
$$
H^p(X, \SheafExt^q(\mathcal{F}, \mathcal{G})) \Rightarrow
\Ext^{p + q}_X(\mathcal{F}, \mathcal{G})
$$
；参见 Cohomology, Section \ref{cohomology-section-ext}。设 $x \in X$。有
$$
\SheafExt^q(\mathcal{F}, \mathcal{G})_x =
\SheafExt^q_{\mathcal{O}_{X, x}}(\mathcal{F}_x, \mathcal{G}_x)
$$
；参见 Cohomology, Lemma \ref{cohomology-lemma-stalk-internal-hom}
（这里还使用了 Derived Categories of Schemes, Lemma
\ref{perfect-lemma-identify-pseudo-coherent-noetherian} 所给出的
$\mathcal{F}$ 的伪凝聚性）。令 $d_x = \dim(\mathcal{O}_{X, x})$。
由于 $\mathcal{O}_{X, x}$ 正则，环 $\mathcal{O}_{X, x}$ 的整体维数为
$d_x$；参见 Algebra, Proposition
\ref{algebra-proposition-regular-finite-gl-dim}。因此
$\SheafExt^q_{\mathcal{O}_{X, x}}(\mathcal{F}_x, \mathcal{G}_x)$
在 $q > d_x$ 时为零。于是模 $\SheafExt^q(\mathcal{F}, \mathcal{G})$
的支撑维数至多为 $d - q$。故
$H^p(X, \SheafExt^q(\mathcal{F}, \mathcal{G})) = 0$，只要 $p > d - q$
；这是由 Cohomology, Proposition
\ref{cohomology-proposition-vanishing-Noetherian} 得到的。第 (1) 部分得证。

\medskip\noindent
证明 (2)。可对 $K$ 与 $L$ 的非零上同调层数目作归纳。当这些数目为
$0, 1$ 时，结论由 (1) 得到。若 $K$ 的非零上同调层数目为 $> 1$，则令
$i \in \mathbf{Z}$ 为使 $H^i(K)$ 非零的最小整数。得到特异三角
$$
H^i(K)[-i] \to K \to \tau_{\geq i + 1}K
$$
(Derived Categories, Remark
\ref{derived-remark-truncation-distinguished-triangle})
，并且 $\Hom(K, L)$ 的消失可由
$\Hom(H^i(K)[-i], L)$ 与 $\Hom(\tau_{\geq i + 1}K, L)$ 的消失通过
Derived Categories, Lemma \ref{derived-lemma-representable-homological} 推出。
若 $L$ 有多于一个非零上同调层，论证类似。
\end{proof}

\begin{lemma}
\label{equiv-lemma-split-complex-regular}
设 $X$ 为维数 $d < \infty$ 的正则 Noether 概形。设
$K \in D^b_{\textit{Coh}}(\mathcal{O}_X)$ 且 $a \in \mathbf{Z}$。若
$H^i(K) = 0$ 对所有 $a < i < a + d$ 成立，则
$K = \tau_{\leq a}K \oplus \tau_{\geq a + d}K$.
\end{lemma}

\begin{proof}
由假设的上同调层消失，有
$\tau_{\leq a}K = \tau_{\leq a + d - 1}K$。由 Derived Categories,
Remark \ref{derived-remark-truncation-distinguished-triangle}，有特异三角
$$
\tau_{\leq a}K \to K \to \tau_{\geq a + d}K \xrightarrow{\delta}
(\tau_{\leq a}K)[1]
$$
由 Derived Categories, Lemma \ref{derived-lemma-split}，只需证明态射
$\delta$ 为零。这由 Lemma \ref{equiv-lemma-Ext-0-regular} 得到。
\end{proof}

\begin{lemma}
\label{equiv-lemma-diagonal-trick}
设 $k$ 为域。设 $X$ 为 $k$ 上拟紧、分离且光滑的概形。存在有限局部自由
$\mathcal{O}_X$-模 $\mathcal{E}$ 与 $\mathcal{G}$，使得
$$
\mathcal{O}_\Delta \in \langle \mathcal{E} \boxtimes \mathcal{G} \rangle
$$
在 $D(\mathcal{O}_{X \times X})$ 中成立；记号同 Derived Categories,
Section \ref{derived-section-generators}。
\end{lemma}

\begin{proof}
由 Varieties, Lemma \ref{varieties-lemma-smooth-regular}，回忆 $X$ 正则。
因此由 Derived Categories of Schemes, Lemma
\ref{perfect-lemma-regular-resolution-property}，$X$ 具有解消性质。故可选取
Lemma \ref{equiv-lemma-diagonal-resolution} 中的解消。记 $\dim(X) = d$。由于
$X \times X$ 在 $k$ 上光滑，它是正则的。因此 $X \times X$ 是正则
Noether 概形，且 $\dim(X \times X) = 2d$。对象
$$
K = (\mathcal{E}_{2d} \boxtimes \mathcal{G}_{2d} \to
\ldots \to
\mathcal{E}_0 \boxtimes \mathcal{G}_0)
$$
属于 $D_{perf}(\mathcal{O}_{X \times X})$；其上同调层为
$\mathcal{O}_\Delta$（次数 $0$）以及
$\Ker(\mathcal{E}_{2d} \boxtimes \mathcal{G}_{2d} \to
\mathcal{E}_{2d-1} \boxtimes \mathcal{G}_{2d-1})$（次数 $-2d$），
其余次数均为零。
因此由 Lemma \ref{equiv-lemma-split-complex-regular} 可知，$\mathcal{O}_\Delta$
是 $K$ 在 $D_{perf}(\mathcal{O}_{X \times X})$ 中的直和项。显然，对象
$K$ 属于
$$
\left\langle
\bigoplus\nolimits_{i = 0, \ldots, 2d} \mathcal{E}_i \boxtimes \mathcal{G}_i
\right\rangle
\subset
\left\langle
\left(\bigoplus\nolimits_{i = 0, \ldots, 2d} \mathcal{E}_i\right)
\boxtimes
\left(\bigoplus\nolimits_{i = 0, \ldots, 2d} \mathcal{G}_i\right)
\right\rangle
$$
，证明完毕。（读者可参见 Derived Categories, Lemmas
\ref{derived-lemma-generated-by-E-explicit} 与
\ref{derived-lemma-in-cone-n}，以确认我们的对象属于此范畴。）
\end{proof}

\begin{lemma}
\label{equiv-lemma-smooth-proper-strong-generator}
设 $k$ 为域。设 $X$ 为 $k$ 上固有且光滑的概形。则
$D_{perf}(\mathcal{O}_X)$ 具有强生成元。
\end{lemma}

\begin{proof}
使用 Lemma \ref{equiv-lemma-diagonal-trick}，选取有限局部自由
$\mathcal{O}_X$-模 $\mathcal{E}$ 与 $\mathcal{G}$，使得
$\mathcal{O}_\Delta \in \langle \mathcal{E} \boxtimes \mathcal{G} \rangle$
在 $D(\mathcal{O}_{X \times X})$ 中成立。断言 $\mathcal{G}$ 是
$D_{perf}(\mathcal{O}_X)$ 的强生成元。沿用 Derived Categories, Section
\ref{derived-section-operate-on-full} 中的记号，选取 $m, n \geq 1$，使得
$$
\mathcal{O}_\Delta \in
smd(add(\mathcal{E} \boxtimes \mathcal{G}[-m, m])^{\star n})
$$
这可由 Derived Categories, Lemma
\ref{derived-lemma-find-smallest-containing-E} 做到。设 $K$ 为
$D_{perf}(\mathcal{O}_X)$ 的对象。由于
$L\text{pr}_1^*K \otimes_{\mathcal{O}_{X \times X}}^\mathbf{L} -$
是正合函子，并且
$$
L\text{pr}_1^*K \otimes_{\mathcal{O}_{X \times X}}^\mathbf{L}
(\mathcal{E} \boxtimes \mathcal{G}) =
(K \otimes_{\mathcal{O}_X}^\mathbf{L} \mathcal{E}) \boxtimes \mathcal{G}
$$
，由 Derived Categories, Remark \ref{derived-remark-operations-functor} 可知
$$
L\text{pr}_1^*K
\otimes_{\mathcal{O}_{X \times X}}^\mathbf{L}
\mathcal{O}_\Delta
\in
smd(add(
(K \otimes_{\mathcal{O}_X}^\mathbf{L} \mathcal{E})
\boxtimes \mathcal{G}[-m, m])^{\star n})
$$
应用正合函子 $R\text{pr}_{2, *}$，并注意到
$$
R\text{pr}_{2, *}
\left((K \otimes_{\mathcal{O}_X}^\mathbf{L} \mathcal{E}) \boxtimes
\mathcal{G}\right) =
R\Gamma(X, K \otimes_{\mathcal{O}_X}^\mathbf{L} \mathcal{E})
\otimes_k \mathcal{G}
$$
；由 Derived Categories of Schemes, Lemma
\ref{perfect-lemma-cohomology-base-change}，得到
$$
K = R\text{pr}_{2, *}(L\text{pr}_1^*K
\otimes_{\mathcal{O}_{X \times X}}^\mathbf{L} \mathcal{O}_\Delta)
\in
smd(add(R\Gamma(X, K \otimes_{\mathcal{O}_X}^\mathbf{L} \mathcal{E})
\otimes_k \mathcal{G}[-m, m])^{\star n})
$$
等号由 Example \ref{equiv-example-diagonal-fourier-mukai} 中的讨论得到。由于
$K$ 完美，存在 $a \leq b$，使得 $H^i(X, K)$ 仅在
$i \in [a, b]$ 时非零。由于 $X$ 固有，每个 $H^i(X, K)$ 都是有限维的。
因此右端包含于
$smd(add(\mathcal{G}[-m + a, m + b])^{\star n})$，而后者又由上述某个引用
包含于 $\langle \mathcal{G} \rangle_n$。证明完毕。
\end{proof}

\begin{lemma}
\label{equiv-lemma-diagonal-trick-proper}
设 $k$ 为域。设 $X$ 为 $k$ 上固有光滑的概形。存在整数 $m, n \geq 1$
与有限局部自由 $\mathcal{O}_X$-模 $\mathcal{G}$，使每个相干
$\mathcal{O}_X$-模都包含于 $smd(add(\mathcal{G}[-m, m])^{\star n})$；
记号同 Derived Categories, Section \ref{derived-section-operate-on-full}。
\end{lemma}

\begin{proof}
在 Lemma \ref{equiv-lemma-smooth-proper-strong-generator} 的证明中，已经说明存在
$m', n \geq 1$，使得对任意相干 $\mathcal{O}_X$-模 $\mathcal{F}$，有
$$
\mathcal{F} \in smd(add(\mathcal{G}[-m' + a, m' + b])^{\star n})
$$
；这里 $a \leq b$ 可任取，只要 $H^i(X, \mathcal{F})$ 仅在
$i \in [a, b]$ 时非零。因此可取 $a = 0$ 与 $b = \dim(X)$。再取
$m = \max(m', m' + b)$，证明完毕。
\end{proof}

\noindent
下一个引理就是本节标题所指的有界性结果。

\begin{lemma}
\label{equiv-lemma-boundedness}
设 $k$ 为域。设 $X$ 为 $k$ 上光滑固有的概形。设 $\mathcal{A}$ 为阿贝尔范畴。
设 $H : D_{perf}(\mathcal{O}_X) \to \mathcal{A}$ 为同调函子
（Derived Categories, Definition \ref{derived-definition-homological}），并且
对所有 $K$（作为 $D_{perf}(\mathcal{O}_X)$ 的对象），对象 $H^i(K)$ 仅对
有限多个 $i \in \mathbf{Z}$ 非零。则存在整数 $m \geq 1$，使得
$H^i(\mathcal{F}) = 0$ 对任意相干 $\mathcal{O}_X$-模 $\mathcal{F}$
以及 $i \not \in [-m, m]$ 成立。上同调函子的情形类似。
\end{lemma}

\begin{proof}
结合 Lemma \ref{equiv-lemma-diagonal-trick-proper} 与
Derived Categories, Lemma \ref{derived-lemma-forward-cone-n} 即得。
\end{proof}

\begin{lemma}
\label{equiv-lemma-bounded-fibres}
设 $k$ 为域。设 $X$、$Y$ 为 $k$ 上的有限型概形。设
$K_0 \to K_1 \to K_2 \to \ldots$ 为
$D_{perf}(\mathcal{O}_{X \times Y})$ 中的对象系统，并设整数 $m \geq 0$，
使得
\begin{enumerate}
\item $H^q(K_i)$ 仅在 $q \leq m$ 时非零；
\item 对于每个相干 $\mathcal{O}_X$-模 $\mathcal{F}$，若
$\dim(\text{Supp}(\mathcal{F})) = 0$，则对象
$$
R\text{pr}_{2, *}(
\text{pr}_1^*\mathcal{F} \otimes_{\mathcal{O}_{X \times Y}}^\mathbf{L}
K_n)
$$
在 $[-m, m] \cup [-m - n, m - n]$ 以外的次数上具有零上同调层，并且当
$n > 2m$ 时，转移映射在 $[-m, m]$ 中各次数的上同调层上诱导同构。
\end{enumerate}
则 $K_n$ 在 $[-m, m] \cup [-m - n, m - n]$ 以外的次数上具有零上同调层，
并且当 $n > 2m$ 时，转移映射在 $[-m, m]$ 中各次数的上同调层上诱导同构。
此外，若 $X$ 与 $Y$ 在 $k$ 上光滑，则对充分大的 $n$，有
$K_n = K \oplus C_n$，该等式在 $D_{perf}(\mathcal{O}_{X \times Y})$ 中成立；
其中 $K$ 的上同调仅位于次数 $[-m, m]$，而 $C_n$ 的上同调仅位于次数
$[-m - n, m - n]$，转移映射在各个 $K$ 的副本之间给出同构。
\end{lemma}

\begin{proof}
设 $Z$ 为 (2) 中某个 $\mathcal{F}$ 的概形论支撑。则
$Z \to \Spec(k)$ 有限，因而 $Z \times Y \to Y$ 有限。因此，对于对象
$M$（属于 $D_\QCoh(\mathcal{O}_{X \times Y})$），若其上同调层支撑于
$Z \times Y$，则有
$H^i(R\text{pr}_{2, *}(M)) = \text{pr}_{2, *}H^i(M)$，并且函子
$\text{pr}_{2, *}$ 在支撑于 $Z \times Y$ 的拟相干模上忠实；细节从略。
故对象
$$
\text{pr}_1^*\mathcal{F} \otimes_{\mathcal{O}_{X \times Y}}^\mathbf{L} K_n
$$
在 $D_{perf}(\mathcal{O}_{X \times Y})$ 中的上同调层在
$[-m, m] \cup [-m - n, m - n]$ 之外为零，并且当 $n > 2m$ 时，转移映射
在 $[-m, m]$ 中各次数的上同调层上诱导同构。设
$z \in X \times Y$ 为映到闭点 $x \in X$ 的闭点。则
$$
K_{n, z} \otimes_{\mathcal{O}_{X \times Y, z}}^\mathbf{L}
\mathcal{O}_{X \times Y, z}/\mathfrak m_x^t\mathcal{O}_{X \times Y, z}
$$
仅在区间 $[-m, m] \cup [-m - n, m - n]$ 中具有非零上同调。由
More on Algebra, Lemma
\ref{more-algebra-lemma-kollar-kovacs-pseudo-coherent} 可知，$K_{n, z}$
仅在次数 $[-m, m] \cup [-m - n, m - n]$ 中具有非零上同调。由于这对
$X \times Y$ 的所有闭点都成立，可知 $K_n$ 仅在次数
$[-m, m] \cup [-m - n, m - n]$ 中具有非零上同调层。完全相同的论证表明，
映射 $K_n \to K_{n + 1}$ 在 $[-m, m]$ 中各次数的上同调层上是同构，只要
$n > 2m$。

\medskip\noindent
若 $X$ 与 $Y$ 在 $k$ 上光滑，则 $X \times Y$ 在 $k$ 上光滑，因而由
Varieties, Lemma \ref{varieties-lemma-smooth-regular} 它是正则的。因此，由
Lemma \ref{equiv-lemma-split-complex-regular}，便得到 $K_n$ 的直和分解，只要
$n > 2m + \dim(X \times Y)$。最后一个陈述由此显然。
\end{proof}





\section{同胞函子}
\label{equiv-section-sibling}

\noindent
本节围绕下述概念证明若干范畴论结果。

\begin{definition}
\label{equiv-definition-siblings}
设 $\mathcal{A}$ 为阿贝尔范畴，设 $\mathcal{D}$ 为三角范畴。称三角范畴之间的
两个正合函子
$$
F, F' : D^b(\mathcal{A}) \longrightarrow \mathcal{D}
$$
为{\it 同胞函子}，或者称 $F'$ 是 $F$ 的一个{\it 同胞}，如果下列两个条件成立：
\begin{enumerate}
\item 函子 $F \circ i$ 与 $F' \circ i$ 同构，其中
$i : \mathcal{A} \to D^b(\mathcal{A})$ 为包含函子；并且
\item 有 $F(K) \cong F'(K)$，对任意 $K$（属于 $D^b(\mathcal{A})$）。
\end{enumerate}
\end{definition}

\noindent
有时第二个条件可由第一个条件推出。

\begin{lemma}
\label{equiv-lemma-sibling-fully-faithful}
设 $\mathcal{A}$ 为阿贝尔范畴，设 $\mathcal{D}$ 为三角范畴。设
$F, F' : D^b(\mathcal{A}) \longrightarrow \mathcal{D}$
为三角范畴之间的正合函子。假设
\begin{enumerate}
\item 函子 $F \circ i$ 与 $F' \circ i$ 同构，其中
$i : \mathcal{A} \to D^b(\mathcal{A})$ 为包含函子；并且
\item 对所有 $X, Y \in \Ob(\mathcal{A})$，有
$\Ext^q_\mathcal{D}(F(X), F(Y)) = 0$，只要 $q < 0$（例如当
$F$ 全忠实时）。
\end{enumerate}
则 $F$ 与 $F'$ 是同胞函子。
\end{lemma}

\begin{proof}
设 $K \in D^b(\mathcal{A})$。将证明 $F(K)$ 同构于 $F'(K)$。可将 $K$
表示为有界复形 $A^\bullet$，其项是 $\mathcal{A}$ 中的对象。将 $K$ 替换为
某个平移后，可假设 $A^i = 0$ 对所有 $i > 0$ 成立。选取 $n \geq 0$，使得
$A^{-i} = 0$ 对所有 $i > n$ 成立。对象
$$
M_i = (A^{-i} \to \ldots \to A^0)[-i],\quad i = 0, \ldots, n
$$
在 $D^b(\mathcal{A})$ 中组成复形
$A^\bullet = A^{-n} \to \ldots \to A^0$ 的一个 Postnikov 系统，而该复形
也属于 $D^b(\mathcal{A})$；参见 Derived Categories, Example
\ref{derived-example-key-postnikov}。由于 $F$ 与 $F'$ 都是三角范畴之间的正合函子，
$$
F(M_i)
\quad\text{且}\quad
F'(M_i)
$$
在 $\mathcal{D}$ 中都组成下述复形的 Postnikov 系统：
$$
F(A^{-n}) \to \ldots \to F(A^0) =
F'(A^{-n}) \to \ldots \to F'(A^0)
$$
由假设，这些对象之间的所有负次数 $\Ext$ 都消失，故由 Postnikov 系统的唯一性
（Derived Categories, Lemma \ref{derived-lemma-existence-postnikov-system}）得到
$F(K) = F(M_n[n]) \cong F'(M_n[n]) = F'(K)$。
\end{proof}

\begin{lemma}
\label{equiv-lemma-sibling-faithful}
设 $F$ 与 $F'$ 为 Definition \ref{equiv-definition-siblings} 中的同胞函子。则
\begin{enumerate}
\item 若 $F$ 本质满，则 $F'$ 本质满；
\item 若 $F$ 全忠实，则 $F'$ 全忠实。
\end{enumerate}
\end{lemma}

\begin{proof}
第 (1) 部分立即由同胞函子的性质 (2) 得到。

\medskip\noindent
假设 $F$ 全忠实。以 $\mathcal{D}' \subset \mathcal{D}$ 记 $F$ 的本质像，于是
$F : D^b(\mathcal{A}) \to \mathcal{D}'$ 是等价。由同胞函子的性质 (2)，函子
$F'$ 通过 $\mathcal{D}'$ 分解，故可考虑函子
$H = F^{-1} \circ F' : D^b(\mathcal{A}) \to D^b(\mathcal{A})$。注意 $H$
是恒等函子的同胞。只需证明 $H$ 全忠实，因而问题化归为下一段所讨论的情形。

\medskip\noindent
令 $\mathcal{D} = D^b(\mathcal{A})$。须证明恒等函子的同胞
$F : \mathcal{D} \to \mathcal{D}$ 全忠实。以 $a_X : X \to F(X)$ 记
Definition \ref{equiv-definition-siblings} 对 $X \in \Ob(\mathcal{A})$ 给出的函子性同构。
对于任意 $K$（属于 $\mathcal{D}$）以及特异三角
$K_1 \to K_2 \to K_3$（属于 $\mathcal{D}$），若映射
$$
F : \Hom(K, K_i[n]) \to \Hom(F(K), F(K_i[n]))
$$
对所有 $n \in \mathbf{Z}$ 以及 $i = 1, 3$ 都是同构，则对 $i = 2$ 以及所有
$n \in \mathbf{Z}$ 也成立。这里使用 $5$-引理 Homology, Lemma
\ref{homology-lemma-five-lemma} 与 Derived Categories, Lemma
\ref{derived-lemma-representable-homological}；细节从略。类似地，若映射
$$
F : \Hom(K_i[n], K) \to \Hom(F(K_i[n]), F(K))
$$
对所有 $n \in \mathbf{Z}$ 以及 $i = 1, 3$ 都是同构，则对 $i = 2$ 以及所有
$n \in \mathbf{Z}$ 也成立。使用典范截断并对非零上同调对象的数目作归纳，
可知只需证明
$$
F : \Ext^q(X, Y) \to \Ext^q(F(X), F(Y))
$$
对所有 $X, Y \in \Ob(\mathcal{A})$ 以及所有 $q \in \mathbf{Z}$ 都是双射。
由于 $F$ 是 $\text{id}$ 的同胞，有 $F(X) \cong X$ 与 $F(Y) \cong Y$，
故右端在 $q < 0$ 时为零。情形 $q = 0$ 由 $F$ 是恒等函子的同胞这一假设成立。
还需证明 $q > 0$ 的情形。

\medskip\noindent
情形 $q = 1$：单射性。元素 $\xi$（属于 $\Ext^1(X, Y)$）给出特异三角
$$
Y \to E \to X \xrightarrow{\xi} Y[1]
$$
注意 $E \in \Ob(\mathcal{A})$。由于 $F$ 是恒等函子的同胞，得到交换图
$$
\xymatrix{
E \ar[d] \ar[r] & X \ar[d] \\
F(E) \ar[r] & F(X)
}
$$
，其中竖直箭头是同构 $a_E$ 与 $a_X$。由 TR3，起初与 $\xi$ 相伴的特异三角
同构于特异三角
$$
F(Y) \to F(E) \to F(X) \xrightarrow{F(\xi)} F(Y[1]) = F(Y)[1]
$$
因此 $\xi = 0$ 当且仅当 $F(\xi)$ 为零，亦即
$F : \Ext^1(X, Y) \to \Ext^1(F(X), F(Y))$ 是单射。

\medskip\noindent
情形 $q = 1$：满射性。设 $\theta$ 为 $\Ext^1(F(X), F(Y))$ 的元素。
它定义了 $F(X)$ 被 $F(Y)$ 的一个扩张（位于 $\mathcal{A}$ 中）；可将其
中项写为 $F(E)$，因为 $F$ 是恒等函子的同胞。于是得到特异三角
$$
F(Y) \xrightarrow{F(\alpha)} F(E)
\xrightarrow{F(\beta)} F(X)
\xrightarrow{\theta} F(Y[1]) = F(Y)[1]
$$
，其中有态射 $\alpha : Y \to E$ 与 $\beta : E \to X$。由于 $F$ 是恒等
函子的同胞，序列 $0 \to Y \to E \to X \to 0$ 是 $\mathcal{A}$ 中的短正合列。
因此得到特异三角
$$
Y \xrightarrow{\alpha} E \xrightarrow{\beta} X \xrightarrow{\delta} Y[1]
$$
，其中有态射 $\delta : X \to Y[1]$。应用正合函子 $F$，得到特异三角
$$
F(Y) \xrightarrow{F(\alpha)} F(E) \xrightarrow{F(\beta)} F(X)
\xrightarrow{F(\delta)} F(Y)[1]
$$
按上述方式论证，可知这两个三角同构。因此存在交换图
$$
\xymatrix{
F(X) \ar[d]^\gamma \ar[r]_{F(\delta)} & F(Y[1]) \ar[d]_\epsilon \\
F(X) \ar[r]^\theta & F(Y[1])
}
$$
，其中 $\gamma$、$\epsilon$ 是某些同构（还可说得更多，但这里不需要）。
可写成 $\gamma = F(\gamma')$ 与 $\epsilon = F(\epsilon')$。于是
$\theta = F(\epsilon' \circ \delta \circ (\gamma')^{-1})$
，从而得到满射性。

\medskip\noindent
情形 $q > 1$：满射性。使用 Yoneda 扩张（参见 Derived Categories, Section
\ref{derived-section-ext}），对于任意元素 $\xi$（属于 $\Ext^q(F(X), F(Y))$），
可找到
$F(X) = B_0, B_1, \ldots, B_{q - 1}, B_q = F(Y) \in \Ob(\mathcal{A})$ 以及元素
$$
\xi_i \in \Ext^1(B_{i - 1}, B_i)
$$
，使得 $\xi$ 是复合 $\xi_q \circ \ldots \circ \xi_1$。写成
$B_i = F(A_i)$（当然有 $A_i = B_i$，但不需要使用这一点），于是
$$
\xi_i = F(\eta_i) \in \Ext^1(F(A_{i - 1}), F(A_i))
\quad\text{其中}\quad
\eta_i \in \Ext^1(A_{i - 1}, A_i)
$$
；这是由 $q = 1$ 时的满射性得到的。于是
$\eta = \eta_q \circ \ldots \circ \eta_1$ 是 $\Ext^q(X, Y)$ 的元素，
且 $F(\eta) = \xi$。

\medskip\noindent
情形 $q > 1$：单射性。元素 $\xi$（属于 $\Ext^q(X, Y)$）给出特异三角
$$
Y[q - 1] \to E \to X \xrightarrow{\xi} Y[q]
$$
应用 $F$，得到特异三角
$$
F(Y)[q - 1] \to F(E) \to F(X) \xrightarrow{F(\xi)} F(Y)[q]
$$
若 $F(\xi) = 0$，则 $F(E) \cong F(Y)[q - 1] \oplus F(X)$ 在
$\mathcal{D}$ 中成立；参见 Derived Categories, Lemma
\ref{derived-lemma-split}。由于 $F$ 是恒等函子的同胞，有
$E \cong F(E)$，因而
$$
E \cong F(E) \cong F(Y)[q - 1] \oplus F(X) \cong Y[q - 1] \oplus X
$$
换言之，$E$ 同构于其上同调对象的直和。这蕴含最初的特异三角分裂，
亦即 $\xi = 0$。
\end{proof}

\noindent
作如下非标准定义。设 $\mathcal{A}$ 为阿贝尔范畴。称 $\mathcal{A}$
{\it 有足够多的负对象}，如果对任意给定的 $X \in \Ob(\mathcal{A})$，
存在对象 $N$，使得
\begin{enumerate}
\item 存在满射 $N \to X$；并且
\item $\Hom(X, N) = 0$.
\end{enumerate}
下面证明关于这一概念的两个引理，以便用于 Proposition
\ref{equiv-proposition-siblings-isomorphic} 的证明。

\begin{lemma}
\label{equiv-lemma-good-map}
设 $\mathcal{A}$ 为有足够多负对象的阿贝尔范畴。设
$X \in D^b(\mathcal{A})$。设 $b \in \mathbf{Z}$ 满足
$H^i(X) = 0$ 对 $i > b$ 成立。则存在映射 $N[-b] \to X$，
使诱导映射 $N \to H^b(X)$ 为满射，且 $\Hom(H^b(X), N) = 0$。
\end{lemma}

\begin{proof}
利用截断函子，可将 $X$ 表示为复形
$A^a \to A^{a + 1} \to \ldots \to A^b$，其各项是 $\mathcal{A}$ 中的对象。
选取 $N$（属于 $\mathcal{A}$），使得存在满射 $t : N \to A^b$，并且
$\Hom(A^b, N) = 0$。于是满射 $t$ 定义所需的映射 $N[-b] \to X$。
\end{proof}

\begin{lemma}
\label{equiv-lemma-good-map-zero}
设 $\mathcal{A}$ 为有足够多负对象的阿贝尔范畴。设
$f : X \to X'$ 为 $D^b(\mathcal{A})$ 的态射。设 $b \in \mathbf{Z}$，
使得 $H^i(X) = 0$ 对 $i > b$ 成立，而 $H^i(X') = 0$ 对
$i \geq b$ 成立。则存在映射 $N[-b] \to X$，使诱导映射
$N \to H^b(X)$ 为满射，使得 $\Hom(H^b(X), N) = 0$，并且使复合
$N[-b] \to X \to X'$ 为零。
\end{lemma}

\begin{proof}
可将 $f$ 表示为有界复形之间的映射
$f^\bullet : A^\bullet \to B^\bullet$，其中复形的各项是 $\mathcal{A}$ 中的对象；
例如参见 Derived Categories, Lemma \ref{derived-lemma-bounded-derived}。考虑对象
$$
C = \Ker(A^b \to A^{b + 1}) \times_{\Ker(B^b \to B^{b + 1})} B^{b - 1}
$$
（属于 $\mathcal{A}$）。由于 $H^b(B^\bullet) = 0$，可知
$C \to H^b(A^\bullet)$ 为满射。
另一方面，映射 $C \to A^b \to B^b$ 与映射
$C \to B^{b - 1} \to B^b$ 相同，因而复合 $C[-b] \to X \to X'$
为零。由于 $\mathcal{A}$ 有足够多负对象，可找到对象 $N$，它带有满射
$N \to C \oplus H^b(X)$，并且 $\Hom(C \oplus H^b(X), N) = 0$。
于是 $N$ 连同映射 $N[-b] \to X$ 给出引理所要求的解。
\end{proof}

\noindent
我们建议读者阅读原始的 \cite[Proposition 2.16]{Orlov-K3}，以领会
下述命题证明中所蕴含的精彩思想。

\begin{proposition}
\label{equiv-proposition-siblings-isomorphic}
\begin{reference}
\cite[Proposition 2.16]{Orlov-K3}；在上述负对象的定义中，无须假设
$\Ext^q(N, X)$ 对 $q > 0$ 消失这一事实归功于 \cite{Canonaco-Stellari}。
\end{reference}
设 $F$ 与 $F'$ 为 Definition \ref{equiv-definition-siblings} 所述的同胞函子。
假设 $F$ 全忠实，并且 $\mathcal{A}$ 有足够多负对象（见上）。则
$F$ 与 $F'$ 是同构函子。
\end{proposition}

\begin{proof}
由 Definition \ref{equiv-definition-siblings} 的第 (2) 部分，函子 $F'$ 的像包含在
函子 $F$ 的本质像中。因此函子 $H = F^{-1} \circ F'$ 是恒等函子的同胞。
这将问题化归为下一段所述的情形。

\medskip\noindent
令 $\mathcal{D} = D^b(\mathcal{A})$。须证明恒等函子的同胞
$F : \mathcal{D} \to \mathcal{D}$ 与恒等函子同构。给定对象 $X$（属于
$\mathcal{D}$），若 $X$ 的{\it 宽度}为 $w = w(X)$，其含义是 $w \geq 0$
为满足下述条件的最小整数：存在整数 $a \in \mathbf{Z}$，使得
$H^i(X) = 0$ 对 $i \not \in [a, a + w - 1]$ 成立。由于 $F$ 是恒等函子的同胞，
并且 $F \circ [n] = [n] \circ F$，我们已经有同构
$$
c_X : X \to F(X)
$$
，它们对 $w(X) \leq 1$ 定义并与平移相容。此外，若
$X = A[-a]$ 且 $X' = A'[-a]$，其中 $A, A' \in \Ob(\mathcal{A})$，
则对任意态射 $f : X \to X'$，图
\begin{equation}
\label{equiv-equation-to-show}
\vcenter{
\xymatrix{
X \ar[d]_{c_X} \ar[r]_f &
X' \ar[d]^{c_{X'}} \\
F(X) \ar[r]^{F(f)} &
F(X')
}
}
\end{equation}
交换。

\medskip\noindent
下面证明：对任意态射 $f : X \to X'$，若 $w(X), w(X') \leq 1$，则图
(\ref{equiv-equation-to-show}) 交换。若 $X$ 或 $X'$ 为零，这一点显然。否则可唯一地写成
$X = A[-a]$ 与 $X' = A'[-a']$，其中 $A, A'$ 属于 $\mathcal{A}$，并且
$a, a' \in \mathbf{Z}$。情形 $a = a'$ 已在上面讨论。若 $a' > a$，则
$f = 0$（Derived Categories, Lemma \ref{derived-lemma-negative-exts}），结论显然。
若 $a' < a$，则 $f$ 对应于元素 $\xi \in \Ext^q(A, A')$，其中
$q = a - a'$。使用 Yoneda 扩张（参见 Derived Categories, Section
\ref{derived-section-ext}），可找到
$A = A_0, A_1, \ldots, A_{q - 1}, A_q = A' \in \Ob(\mathcal{A})$ 以及元素
$$
\xi_i \in \Ext^1(A_{i - 1}, A_i)
$$
，使得 $\xi$ 是复合 $\xi_q \circ \ldots \circ \xi_1$。换言之，令
$X_i = A_i[-a + i]$，便得到态射
$$
X = X_0 \xrightarrow{f_1} X_1 \to \ldots \to X_{q - 1}
\xrightarrow{f_q} X_q = X'
$$
，其复合为 $f$。由于图 (\ref{equiv-equation-to-show}) 对 $f_1, \ldots, f_q$ 的交换性
蕴含它对 $f$ 的交换性，故可化归到 $q = 1$ 的情形。在此情形中，作平移后
可假设有特异三角
$$
A' \to E \to A \xrightarrow{f} A'[1]
$$
。注意 $E$ 是 $\mathcal{A}$ 的对象。考虑下图
$$
\xymatrix{
E \ar[d]_{c_E} \ar[r] &
A \ar[d]_{c_A} \ar[r]_f &
A'[1] \ar[d]^{c_{A'}[1]}
\ar@{..>}@<-1ex>[d]_\gamma \ar@{..>}[ld]^\epsilon \ar[r] &
E[1] \ar[d]^{c_E[1]} \\
F(E) \ar[r] &
F(A) \ar[r]^{F(f)} &
F(A')[1] \ar[r] &
F(E)[1]
}
$$
，其两行都是特异三角。右侧方块已经交换，但尚不知道中间方块是否交换。
由三角范畴的公理，可找到态射 $\gamma$ 使整个图交换。于是
$\gamma - c_{A'}[1]$ 与 $F(A')[1] \to F(E)[1]$ 的复合为零，故可找到
$\epsilon : A'[1] \to F(A)$，使得
$\gamma - c_{A'}[1] = F(f) \circ \epsilon$。然而，任何箭头
$A'[1] \to F(A)$ 都为零，因为它是 $\mathcal{A}$ 中对象之间的负次数 Ext 类。
因此 $\gamma = c_{A'}[1]$，从而中间方块也交换，这正是所要证明的。

\medskip\noindent
为完成证明，我们对 $w$ 归纳，以证明存在同构 $c_X : X \to F(X)$，
对所有 $X$（满足 $w(X) \leq w$）定义，并与这些对象之间的所有态射相容。
基础情形 $w = 1$ 已在上面证明。假设对某个 $w \geq 1$ 已知结论。

\medskip\noindent
设 $X$ 为满足 $w(X) = w + 1$ 的对象。选取 $a \in \mathbf{Z}$，使得
$H^i(X) = 0$ 对 $i \not \in [a, a + w]$ 成立。令 $b = a + w$，于是
$H^b(X)$ 非零。依照 Lemma \ref{equiv-lemma-good-map} 选取 $N[-b] \to X$。
选取特异三角
$$
N[-b] \to X \to Y \to N[-b + 1]
$$
。计算上同调长正合列，得到 $w(Y) \leq w$。因此由归纳假设，得到下图中的
实线箭头
$$
\xymatrix{
N[-b] \ar[r] \ar[d]_{c_N[-b]} &
X \ar[r] \ar@{..>}[d]_{c_{N[-b] \to X}} &
Y \ar[r] \ar[d]^{c_Y} &
N[-b + 1] \ar[d]^{c_N[-b + 1]} \\
F(N)[-b] \ar[r] &
F(X) \ar[r] &
F(Y) \ar[r] &
F(N)[-b + 1]
}
$$
。由此得到虚线箭头 $c_{N[-b] \to X}$。根据 Derived Categories, Lemma
\ref{derived-lemma-uniqueness-third-arrow}，该虚线箭头是唯一的，因为
$\Hom(X, F(N)[-b]) \cong \Hom(X, N[-b]) = 0$；这里使用了我们对 $N$ 的选取。
事实上，
$c_{N[-b] \to X}$ 是使顶点为 $X, Y, F(X), F(Y)$ 的方块交换的唯一虚线箭头。

\medskip\noindent
设 $N'[-b] \to X$ 是 Lemma \ref{equiv-lemma-good-map} 所述的另一映射；我们证明
$c_{N[-b] \to X} = c_{N'[-b] \to X}$。注意映射
$(N \oplus N')[-b] \to X$ 也满足 Lemma \ref{equiv-lemma-good-map} 的条件。
因此可假设 $N'[-b] \to X$ 通过 $N'[-b] \to N[-b] \to X$ 分解，其中
有某个态射 $N' \to N$。选取特异三角
$N[-b] \to X \to Y \to N[-b + 1]$ 与
$N'[-b] \to X \to Y' \to N'[-b + 1]$。由公理 TR3，可找到态射
$g : Y' \to Y$，它与 $\text{id}_X$ 及 $N' \to N$ 一起构成三角之间的态射。
由于图 (\ref{equiv-equation-to-show}) 对 $g$ 交换，故有
$$
(F(X) \to F(Y)) \circ c_{N'[-b] \to X} = (F(X) \to F(Y)) \circ c_{N[-b] \to X}
$$
。由上述构造中指出的 $c_{N[-b] \to X}$ 的唯一性，现在得到
$c_{N'[-b] \to X} = c_{N[-b] \to X}$。

\medskip\noindent
因此现在可对 $X$（宽度为 $w + 1$）定义同构 $c_X : X \to F(X)$，
将它取为所有映射 $c_{N[-b] \to X}$ 的共同值，其中 $N[-b] \to X$ 如
Lemma \ref{equiv-lemma-good-map} 所述。为完成证明，须说明：对满足
上述条件的对象之间的所有态射 $f : X \to X'$，若
$w(X) \leq w + 1$ 且 $w(X') \leq w + 1$，图 (\ref{equiv-equation-to-show}) 都交换。选取
$a \leq b \leq a + w$，使得 $H^i(X) = 0$ 对 $i \not \in [a, b]$ 成立；
并选取 $a' \leq b' \leq a' + w$，使得 $H^i(X') = 0$ 对
$i \not \in [a', b']$ 成立。我们对 $(b' - a') + (b - a)$ 归纳以证明断言。
（基础情形是该数为零；由于 $w \geq 1$，这没有问题。）分两种情形。

\medskip\noindent
情形 I：$b' < b$。在此情形中，由 Lemma \ref{equiv-lemma-good-map-zero}，可依
Lemma \ref{equiv-lemma-good-map} 选取 $N[-b] \to X$，使得复合
$N[-b] \to X \to X'$ 为零。选取特异三角
$N[-b] \to X \to Y \to N[-b + 1]$。由于 $N[-b] \to X'$ 为零，可知
$f$ 通过 $X \to Y \to X'$ 分解。由于 $H^i(Y)$ 只可能对
$i \in [a, b - 1]$ 非零，由归纳可知图 (\ref{equiv-equation-to-show}) 对
$Y \to X'$ 交换。图 (\ref{equiv-equation-to-show}) 对 $X \to Y$ 的交换性：若
$w(X) = w + 1$，则来自构造；若 $w(X) \leq w$，则来自第一个归纳假设。
故图 (\ref{equiv-equation-to-show}) 对 $f$ 交换。

\medskip\noindent
情形 II：$b' \geq b$。在此情形中，依 Lemma \ref{equiv-lemma-good-map} 选取
$N'[-b'] \to X'$。还可假设 $\Hom(H^{b'}(X), N') = 0$（这只在
$b' = b$ 时相关）；例如，可将 $N'$ 替换为一个对象 $N''$，使后者满射到
$N' \oplus H^{b'}(X)$，并且 $\Hom(N' \oplus H^{b'}(X), N'') = 0$。
选取特异三角 $N'[-b'] \to X' \to Y' \to N'[-b' + 1]$。
$\Hom(X, X') \to \Hom(X, Y')$ 由我们对 $N'$ 的选取而为单射（细节从略）；于是
$\Hom(X, F(X')) \to \Hom(X, F(Y'))$ 也为单射。因此在此情形中，只需验证
图 (\ref{equiv-equation-to-show}) 对复合 $X \to Y'$ 交换；该复合由态射
$X \to X' \to Y'$ 构成。由于 $H^i(Y')$ 只可能对 $i \in [a', b' - 1]$ 非零，
结论由归纳假设得到。
\end{proof}










\section{推出全忠实性}
\label{equiv-section-get-fully-faithful}

\noindent
判定一个函子何时全忠实将很有用；为此给出
\cite[Lemma 2.15]{Orlov-K3} 的如下变体。

\begin{lemma}
\label{equiv-lemma-get-fully-faithful}
\begin{reference}
\cite[Lemma 2.15]{Orlov-K3} 的变体
\end{reference}
设 $F : \mathcal{D} \to \mathcal{D}'$ 为三角范畴之间的正合函子。设
$S \subset \Ob(\mathcal{D})$ 为对象的集合。假设
\begin{enumerate}
\item $F$ 同时有右伴随和左伴随；
\item 对 $K \in \mathcal{D}$，若 $\Hom(E, K[i]) = 0$ 对所有
$E \in S$ 及 $i \in \mathbf{Z}$ 成立，则 $K = 0$；
\item 对 $K \in \mathcal{D}$，若 $\Hom(K, E[i]) = 0$ 对所有
$E \in S$ 及 $i \in \mathbf{Z}$ 成立，则 $K = 0$；
\item 映射 $\Hom(E, E'[i]) \to \Hom(F(E), F(E')[i])$（由 $F$ 诱导）对所有
$E, E' \in S$ 及 $i \in \mathbf{Z}$ 都是双射。
\end{enumerate}
则 $F$ 全忠实。
\end{lemma}

\begin{proof}
分别以 $F_r$ 与 $F_l$ 表示 $F$ 的右伴随和左伴随。对 $E \in S$，
选取特异三角
$$
E \to F_r(F(E)) \to C \to E[1]
$$
，其中第一支箭头是伴随的单位。对 $E' \in S$，有
$$
\Hom(E', F_r(F(E))[i]) = \Hom(F(E'), F(E)[i]) = \Hom(E', E[i])
$$
。最后一个等号由假设 (4) 成立。因此把同调函子 $\Hom(E', -)$
（Derived Categories, Lemma \ref{derived-lemma-representable-homological}）
应用于上述特异三角，可知 $\Hom(E', C[i]) = 0$ 对所有
$i \in \mathbf{Z}$ 与 $E' \in S$ 成立。由假设 (2)，得到
$C = 0$ 与 $E = F_r(F(E))$。

\medskip\noindent
对 $K \in \Ob(\mathcal{D})$，选取特异三角
$$
F_l(F(K)) \to K \to C \to F_l(F(K))[1]
$$
，其中第一支箭头是伴随的余单位。对 $E \in S$，有
$$
\Hom(F_l(F(K)), E[i]) = \Hom(F(K), F(E)[i]) =
\Hom(K, F_r(F(E))[i]) = \Hom(K, E[i])
$$
，其中最后一个等号由第一段的结论成立。于是与前面一样得到
$\Hom(C, E[i]) = 0$ 对所有 $E \in S$ 及 $i \in \mathbf{Z}$ 成立。
故由假设 (3)，$C = 0$。再由 Categories, Lemma
\ref{categories-lemma-adjoint-fully-faithful}，$F$ 全忠实。
\end{proof}

\begin{lemma}
\label{equiv-lemma-duality-at-point}
设 $k$ 为域。设 $X$ 为有限型 $k$-概形，并且是正则的。设 $x \in X$
为闭点。给定凝聚 $\mathcal{O}_X$-模 $\mathcal{F}$，其支撑在 $x$ 上；
选取凝聚 $\mathcal{O}_X$-模 $\mathcal{F}'$，其支撑也在 $x$ 上，使得
$\mathcal{F}_x$ 与 $\mathcal{F}'_x$ 互为 Matlis 对偶。则存在同构
$$
\Hom_X(\mathcal{F}, M) =
H^0(X, M \otimes_{\mathcal{O}_X}^\mathbf{L} \mathcal{F}'[-d_x])
$$
，其中 $d_x = \dim(\mathcal{O}_{X, x})$；该同构对 $M$
（属于 $D_{perf}(\mathcal{O}_X)$）具有函子性。
\end{lemma}

\begin{proof}
由于 $\mathcal{F}$ 支撑在 $x$ 上，有
$$
\Hom_X(\mathcal{F}, M) =
\Hom_{\mathcal{O}_{X, x}}(\mathcal{F}_x, M_x)
$$
；类似地，有
$$
H^0(X, M \otimes_{\mathcal{O}_X}^\mathbf{L} \mathcal{F}'[-d_x]) =
\text{Tor}^{\mathcal{O}_{X, x}}_{d_x}(M_x, \mathcal{F}'_x)
$$
。因此只需证明：给定 Noether 正则局部环 $A$（维数为 $d$）与有限长度
$A$-模 $N$，若 $N'$ 是 $N$ 的 Matlis 对偶，则存在函子性同构
$$
\Hom_A(N, K) = \text{Tor}^A_d(K, N')
$$
，对 $K$（属于 $D_{perf}(A)$）成立。左端可写成
$H^0(R\Hom_A(N, A) \otimes_A^\mathbf{L} K)$；这是由 More on Algebra, Lemma
\ref{more-algebra-lemma-dual-perfect-complex} 以及 $N$ 确定 $D(A)$ 的完美对象这一事实得到的。
公式成立是因为
$$
R\Hom_A(N, A) = R\Hom_A(N, A[d])[-d] = N'[-d]
$$
；这里使用 Dualizing Complexes, Lemma
\ref{dualizing-lemma-dualizing-finite-length}，以及 $A[d]$ 是 $A$ 上的正规化
对偶化复形这一事实（由 Dualizing Complexes, Lemma
\ref{dualizing-lemma-regular-gorenstein}，$A$ 是 Gorenstein 环）。
\end{proof}

\begin{lemma}
\label{equiv-lemma-orthogonal-point-sheaf}
设 $k$ 为域。设 $X$ 为有限型 $k$-概形，并且是正则的。设 $x \in X$ 为闭点，
以 $\mathcal{O}_x$ 表示在 $x$ 处取值为 $\kappa(x)$ 的摩天层。设 $K$ 属于
$D_{perf}(\mathcal{O}_X)$。
\begin{enumerate}
\item 若 $\Ext^i_X(\mathcal{O}_x, K) = 0$，则存在开邻域 $U$（它是 $x$ 的邻域），使得
$H^{i - d_x}(K)|_U = 0$，其中 $d_x = \dim(\mathcal{O}_{X, x})$。
\item 若 $\Hom_X(\mathcal{O}_x, K[i]) = 0$ 对所有 $i \in \mathbf{Z}$ 成立，
则 $K$ 在 $x$ 的某个开邻域上为零。
\item 若 $\Ext^i_X(K, \mathcal{O}_x) = 0$，则存在开邻域 $U$（它是 $x$ 的邻域），使得
$H^i(K^\vee)|_U = 0$。
\item 若 $\Hom_X(K, \mathcal{O}_x[i]) = 0$ 对所有 $i \in \mathbf{Z}$ 成立，
则 $K$ 在 $x$ 的某个开邻域上为零。
\item 若 $H^i(X, K \otimes_{\mathcal{O}_X}^\mathbf{L} \mathcal{O}_x) = 0$，
则存在开邻域 $U$（它是 $x$ 的邻域），使得 $H^i(K)|_U = 0$。
\item 若 $H^i(X, K \otimes_{\mathcal{O}_X}^\mathbf{L} \mathcal{O}_x) = 0$
对 $i \in \mathbf{Z}$ 成立，则 $K$ 在 $x$ 的某个开邻域上为零。
\end{enumerate}
\end{lemma}

\begin{proof}
注意 $H^i(X, K \otimes_{\mathcal{O}_X}^\mathbf{L} \mathcal{O}_x)$ 等于
$K_x  \otimes_{\mathcal{O}_{X, x}}^\mathbf{L} \kappa(x)$。因此第 (5) 部分由
More on Algebra, Lemma \ref{more-algebra-lemma-cut-complex-in-two} 得到。
第 (6) 部分由第 (5) 部分得到。第 (1) 部分由第 (5) 部分、Lemma
\ref{equiv-lemma-duality-at-point}，以及 $\kappa(x)$ 的 Matlis 对偶是 $\kappa(x)$ 这一事实得到。
第 (2) 部分由第 (1) 部分得到。第 (3) 部分由第 (5) 部分以及下述事实得到：
$\Ext^i(K, \mathcal{O}_x) =
H^i(X, K^\vee \otimes_{\mathcal{O}_X}^\mathbf{L} \mathcal{O}_x)$
由 Cohomology, Lemma \ref{cohomology-lemma-dual-perfect-complex} 成立。
第 (4) 部分由第 (3) 部分以及刚引用的引理给出的 $K \cong (K^\vee)^\vee$ 得到。
\end{proof}

\begin{lemma}
\label{equiv-lemma-hom-into-point-sheaf}
设 $X$ 为 Noether 概形。设 $x \in X$ 为闭点，以 $\mathcal{O}_x$ 表示在
$x$ 处取值为 $\kappa(x)$ 的摩天层。设 $K$ 属于
$D^b_{\textit{Coh}}(\mathcal{O}_X)$。设 $b \in \mathbf{Z}$。下列条件等价：
\begin{enumerate}
\item $H^i(K)_x = 0$ 对所有 $i > b$ 成立；
\item $\Hom_X(K, \mathcal{O}_x[-i]) = 0$ 对所有 $i > b$ 成立。
\end{enumerate}
\end{lemma}

\begin{proof}
考虑复形 $K_x$（属于 $D^b_{\textit{Coh}}(\mathcal{O}_{X, x})$）。存在整数
$b_x \in \mathbf{Z}$，使得 $K_x$ 可由下列上有界复形表示：
$$
\ldots \to
\mathcal{O}_{X, x}^{\oplus n_{b_x - 2}} \to
\mathcal{O}_{X, x}^{\oplus n_{b_x - 1}} \to
\mathcal{O}_{X, x}^{\oplus n_{b_x}} \to 0 \to \ldots
$$
，其中 $\mathcal{O}_{X, x}^{\oplus n_i}$ 位于次数 $i$，而所有转移映射都由
系数属于 $\mathfrak m_x$ 的矩阵给出。参见 More on Algebra, Lemma
\ref{more-algebra-lemma-lift-pseudo-coherent-from-residue-field}。
结论由此容易得到（并且这些等价条件成立当且仅当 $b \geq b_x$）。
\end{proof}

\begin{lemma}
\label{equiv-lemma-get-fully-faithful-geometric}
设 $k$ 为域。设 $X$ 与 $Y$ 为 $k$ 上的固有概形。假设 $X$ 正则。则 $k$-线性正合函子
$F : D_{perf}(\mathcal{O}_X) \to D_{perf}(\mathcal{O}_Y)$
全忠实，当且仅当对任意闭点 $x, x' \in X$，映射
$$
F : \Ext^i_X(\mathcal{O}_x, \mathcal{O}_{x'})
\longrightarrow
\Ext^i_Y(F(\mathcal{O}_x), F(\mathcal{O}_{x'}))
$$
对所有 $i \in \mathbf{Z}$ 都是同构。这里 $\mathcal{O}_x$ 是在 $x$ 处取值为
$\kappa(x)$ 的摩天层。
\end{lemma}

\begin{proof}
由 Lemma \ref{equiv-lemma-always-right-adjoints}，函子 $F$ 同时有左伴随与右伴随。
因此可应用 Lemma \ref{equiv-lemma-get-fully-faithful} 的判据，因为该引理的假设 (2) 与 (3)
由 Lemma \ref{equiv-lemma-orthogonal-point-sheaf} 得到。
\end{proof}

\begin{lemma}
\label{equiv-lemma-noah-pre}
\begin{reference}
Noah Olander 于 2020 年 6 月 9 日的电子邮件
\end{reference}
设 $k$ 为域。设 $X$ 为正则的固有 $k$-概形。设
$F : D_{perf}(\mathcal{O}_X) \to D_{perf}(\mathcal{O}_X)$ 为 $k$-线性正合函子。
假设对每个凝聚 $\mathcal{O}_X$-模 $\mathcal{F}$，若
$\dim(\text{Supp}(\mathcal{F})) = 0$，则都有同构 $\mathcal{F} \cong F(\mathcal{F})$。
则 $F$ 全忠实。
\end{lemma}

\begin{proof}
由 Lemma \ref{equiv-lemma-get-fully-faithful-geometric}，只需证明映射
$$
F : \Ext^i_X(\mathcal{O}_x, \mathcal{O}_{x'})
\longrightarrow
\Ext^i_X(F(\mathcal{O}_x), F(\mathcal{O}_{x'}))
$$
对所有 $i \in \mathbf{Z}$ 以及所有闭点 $x, x' \in X$ 都是同构。由假设，
源与目标同构。若 $x \not = x'$，则两端都为零，结论成立。若 $x = x'$，
只需证明该映射是单射或满射之一。对 $i < 0$，两端都为零，结论成立。
对 $i = 0$，任意非零映射 $\alpha : \mathcal{O}_x \to \mathcal{O}_x$
（作为 $\mathcal{O}_X$-模态射）都是同构。因此 $F(\alpha)$ 也是同构，
故 $F(\alpha)$ 非零。这就证明了 $i = 0$ 的情形。对 $i = 1$，
非零元素 $\xi$（属于 $\Ext^1(\mathcal{O}_x, \mathcal{O}_x)$）对应于非分裂短正合列
$$
0 \to \mathcal{O}_x \to \mathcal{F} \to \mathcal{O}_x \to 0
$$
。由于 $F(\mathcal{F}) \cong \mathcal{F}$，可知 $F(\mathcal{F})$ 也是
$\mathcal{O}_x$ 被 $\mathcal{O}_x$ 的非分裂扩张。由于
$\mathcal{O}_x \cong F(\mathcal{O}_x)$ 是简单 $\mathcal{O}_X$-模，并且
$\mathcal{F} \cong F(\mathcal{F})$ 的长度为 $2$，可知在特异三角
$$
F(\mathcal{O}_x) \to F(\mathcal{F}) \to F(\mathcal{O}_x)
\xrightarrow{F(\xi)} F(\mathcal{O}_x)[1]
$$
中，前两支箭头必组成短正合列；该短正合列必与上述短正合列同构，因而非分裂。
由此 $F(\xi)$ 非零，得到 $i = 1$ 的结论。对 $i > 1$，Ext 类的复合定义满射
$$
\Ext^1(F(\mathcal{O}_x), F(\mathcal{O}_x)) \otimes \ldots \otimes
\Ext^1(F(\mathcal{O}_x), F(\mathcal{O}_x))
\longrightarrow
\Ext^i(F(\mathcal{O}_x), F(\mathcal{O}_x))
$$
。参见 Duality for Schemes, Lemma \ref{duality-lemma-regular-ideal-ext}。
因此次数 $1$ 的满射性蕴含对 $i > 0$ 的满射性。证明完毕。
\end{proof}










\section{特殊函子}
\label{equiv-section-special-functors}

\noindent
本节证明关于一类特殊函子的若干结果，它们将在本章后面使用。

\begin{definition}
\label{equiv-definition-siblings-geometric}
设 $k$ 为域。设 $X$、$Y$ 为有限型 $k$-概形。回忆
$D^b_{\textit{Coh}}(\mathcal{O}_X) = D^b(\textit{Coh}(\mathcal{O}_X))$
；这是 Derived Categories of Schemes, Proposition
\ref{perfect-proposition-DCoh} 的结论。称两个 $k$-线性正合函子
$$
F, F' :
D^b_{\textit{Coh}}(\mathcal{O}_X) = D^b(\textit{Coh}(\mathcal{O}_X))
\longrightarrow
D^b_{\textit{Coh}}(\mathcal{O}_Y)
$$
为{\it 同胞函子}，或者称 $F'$ 是 $F$ 的一个{\it 同胞}，如果 $F$ 与 $F'$
在 Definition \ref{equiv-definition-siblings} 的意义下为同胞函子，其中阿贝尔范畴取为
$\textit{Coh}(\mathcal{O}_X)$。若 $X$ 正则，则
$D_{perf}(\mathcal{O}_X) = D^b_{\textit{Coh}}(\mathcal{O}_X)$，这是
Derived Categories of Schemes, Lemma \ref{perfect-lemma-perfect-on-noetherian}；
并且对 $k$-线性正合函子
$F, F' : D_{perf}(\mathcal{O}_X) \to D_{perf}(\mathcal{O}_Y)$.
也使用相同术语。
\end{definition}

\begin{lemma}
\label{equiv-lemma-exact-functor-preserving-Coh}
设 $k$ 为域。设 $X$、$Y$ 为有限型 $k$-概形，并且 $X$ 分离。设
$F : D^b_{\textit{Coh}}(\mathcal{O}_X) \to D^b_{\textit{Coh}}(\mathcal{O}_Y)$
为 $k$-线性正合函子，它把
$\textit{Coh}(\mathcal{O}_X) \subset D^b_{\textit{Coh}}(\mathcal{O}_X)$
映入
$\textit{Coh}(\mathcal{O}_Y) \subset D^b_{\textit{Coh}}(\mathcal{O}_Y)$.
。则存在 Fourier--Mukai 函子
$F' : D^b_{\textit{Coh}}(\mathcal{O}_X) \to D^b_{\textit{Coh}}(\mathcal{O}_Y)$
，其核是凝聚 $\mathcal{O}_{X \times Y}$-模 $\mathcal{K}$；它在 $X$ 上平坦，
其支撑在 $Y$ 上有限，并且该函子是 $F$ 的同胞。
\end{lemma}

\begin{proof}
以 $H : \textit{Coh}(\mathcal{O}_X) \to \textit{Coh}(\mathcal{O}_Y)$ 表示
$F$ 的限制。由于 $F$ 是三角范畴之间的正合函子，可知 $H$ 是阿贝尔范畴之间的
正合函子。当然，$H$ 是 $k$-线性的，因为 $F$ 也是。由
Functors and Morphisms, Lemma \ref{functors-lemma-functor-coherent-over-field}
，得到凝聚 $\mathcal{O}_{X \times Y}$-模 $\mathcal{K}$，它在 $X$ 上平坦，
且支撑在 $Y$ 上有限。令 $F'$ 为以 $\mathcal{K}$ 定义的 Fourier--Mukai 函子，
使得 $F'$ 限制为 $H$（当限制到 $ \textit{Coh}(\mathcal{O}_X)$ 时）。由 Lemma
\ref{equiv-lemma-fourier-mukai-Coh}，函子 $F'$ 把
$D^b_{\textit{Coh}}(\mathcal{O}_X)$ 映入 $D^b_{\textit{Coh}}(\mathcal{O}_Y)$。
注意 $F$ 与 $F'$ 满足 Lemma \ref{equiv-lemma-sibling-fully-faithful} 的第一、第二个条件，
故它们是同胞函子。
\end{proof}

\begin{remark}
\label{equiv-remark-difficult}
若 $F, F' : D^b_{\textit{Coh}}(\mathcal{O}_X) \to \mathcal{D}$ 是同胞函子，$F$
全忠实，并且 $X$ 是约化的射影 $k$-概形，则 $F \cong F'$；这可由
Proposition \ref{equiv-proposition-siblings-isomorphic} 配合 Theorem
\ref{equiv-theorem-fully-faithful} 证明中的论证得到。然而一般而言，我们不知道同胞函子
是否同构。即便在 Lemma \ref{equiv-lemma-exact-functor-preserving-Coh} 的情形中，
证明同胞函子 $F$ 与 $F'$ 同构似乎也很困难。若 $X$ 是光滑的固有 $k$-概形，
并且 $F$ 全忠实，则由 \cite{Noah} 可知 $F \cong F'$。
若你在更一般的情形中有证明或反例，请发送电子邮件至
\href{mailto:stacks.project@gmail.com}{stacks.project@gmail.com}.
\end{remark}

\begin{lemma}
\label{equiv-lemma-two-functors-pre}
设 $k$ 为域。设 $X$、$Y$ 为 $k$ 上的固有概形。假设 $X$ 正则。设
$F, G : D_{perf}(\mathcal{O}_X) \to D_{perf}(\mathcal{O}_Y)$
为 $k$-线性正合函子，且满足
\begin{enumerate}
\item $F(\mathcal{F}) \cong G(\mathcal{F})$ 对任意凝聚
$\mathcal{O}_X$-模 $\mathcal{F}$（满足 $\dim(\text{Supp}(\mathcal{F})) = 0$）成立；
\item $F$ 全忠实。
\end{enumerate}
则 $G$ 的本质像包含于 $F$ 的本质像。
\end{lemma}

\begin{proof}
回忆 $F$ 与 $G$ 都有左右伴随；参见 Lemma \ref{equiv-lemma-always-right-adjoints}。
特别地，本质像 $\mathcal{A} \subset D_{perf}(\mathcal{O}_Y)$（即 $F$ 的本质像）满足
Derived Categories, Lemma \ref{derived-lemma-right-adjoint} 的等价条件。我们断言
$G$ 通过 $\mathcal{A}$ 分解。由 Derived Categories, Lemma
\ref{derived-lemma-right-adjoint}，$\mathcal{A} = {}^\perp(\mathcal{A}^\perp)$，
故只需证明 $\Hom_Y(G(M), N) = 0$ 对所有 $M$（属于
$D_{perf}(\mathcal{O}_X)$）以及 $N \in \mathcal{A}^\perp$ 成立。我们有
$$
\Hom_Y(G(M), N) = \Hom_X(M, G_r(N))
$$
，其中 $G_r$ 是 $G$ 的右伴随。因此只需证明 $G_r(N) = 0$。有
$G(\mathcal{F}) \cong F(\mathcal{F})$，对 (1) 中所述的 $\mathcal{F}$ 成立，故
$$
\Hom_X(\mathcal{F}, G_r(N)) =
\Hom_Y(G(\mathcal{F}), N) =
\Hom_Y(F(\mathcal{F}), N) = 0
$$
，因为 $N$ 属于本质像 $\mathcal{A}$（即 $F$ 的本质像）的右正交。当然，
$\Hom_X(\mathcal{F}, G_r(N)[i])$ 对任意 $i \in \mathbf{Z}$ 也同样消失。因此由
Lemma \ref{equiv-lemma-orthogonal-point-sheaf}，$G_r(N) = 0$，证毕。
\end{proof}

\begin{lemma}
\label{equiv-lemma-noah}
\begin{reference}
Noah Olander 于 2020 年 6 月 8 日的电子邮件
\end{reference}
设 $k$ 为域。设 $X$ 为正则的固有 $k$-概形。设
$F : D_{perf}(\mathcal{O}_X) \to D_{perf}(\mathcal{O}_X)$ 为 $k$-线性正合函子。
假设对每个凝聚 $\mathcal{O}_X$-模 $\mathcal{F}$，若
$\dim(\text{Supp}(\mathcal{F})) = 0$，则存在同构
$\mathcal{F} \cong F(\mathcal{F})$。那么存在自同构 $f : X \to X$（定义在 $k$ 上），
它在底拓扑空间上诱导恒等映射\footnote{这常常迫使 $f$ 本身为恒等映射；参见
Varieties, Lemma \ref{varieties-lemma-automorphism}.}，以及可逆
$\mathcal{O}_X$-模 $\mathcal{L}$，使得 $F$ 与
$F'(M) = f^*M \otimes_{\mathcal{O}_X}^\mathbf{L} \mathcal{L}$ 是同胞函子。
\end{lemma}

\begin{proof}
由 Lemma \ref{equiv-lemma-noah-pre}，函子 $F$ 全忠实。由 Lemma
\ref{equiv-lemma-two-functors-pre}，恒等函子的本质像包含于 $F$ 的本质像；亦即 $F$
本质满。因此 $F$ 是等价。注意拟逆 $F^{-1}$ 满足与 $F$ 相同的假设。

\medskip\noindent
设 $M \in D_{perf}(\mathcal{O}_X)$，并设 $H^i(M) = 0$ 对 $i > b$ 成立。
由于 $F$ 全忠实，有
$$
\Hom_X(M, \mathcal{O}_x[-i]) =
\Hom_X(F(M), F(\mathcal{O}_x)[-i]) \cong
\Hom_X(F(M), \mathcal{O}_x[-i])
$$
，对任意 $i \in \mathbf{Z}$ 以及任意闭点 $x$（属于 $X$）成立。因此由 Lemma
\ref{equiv-lemma-hom-into-point-sheaf}，可知 $F(M)$ 的上同调层在次数 $> b$ 中消失。

\medskip\noindent
设 $\mathcal{F}$ 为凝聚 $\mathcal{O}_X$-模。由上可知 $F(\mathcal{F})$ 的上同调层
只可能在次数 $\leq 0$ 非零。令 $\mathcal{G} = H^0(F(\mathcal{F}))$。选取特异三角
$$
K \to F(\mathcal{F}) \to \mathcal{G} \to K[1]
$$
。则 $K$ 的上同调层只可能在次数 $\leq -1$ 非零。应用 $F^{-1}$，得到特异三角
$$
F^{-1}(K) \to \mathcal{F} \to F^{-1}(\mathcal{G}) \to F^{-1}(K')[1]
$$
。由于 $F^{-1}(K)$ 的上同调层只可能在次数 $\leq -1$ 非零（把上一段应用于
$F^{-1}$），箭头 $F^{-1}(K) \to \mathcal{F}$ 为零（Derived Categories, Lemma
\ref{derived-lemma-negative-exts}）。因此 $K \to F(\mathcal{F})$ 为零；由第一个
特异三角的选取，这蕴含 $F(\mathcal{F}) = \mathcal{G}$。

\medskip\noindent
由上一段，$F$ 保持 $\textit{Coh}(\mathcal{O}_X)$，并且确实定义等价
$H : \textit{Coh}(\mathcal{O}_X) \to \textit{Coh}(\mathcal{O}_X)$。由
Functors and Morphisms, Lemma
\ref{functors-lemma-equivalence-coherent-over-field}
，得到自同构 $f : X \to X$（在 $k$ 上）以及可逆 $\mathcal{O}_X$-模
$\mathcal{L}$，使得 $H(\mathcal{F}) = f^*\mathcal{F} \otimes \mathcal{L}$。
令 $F'(M) = f^*M \otimes_{\mathcal{O}_X}^\mathbf{L} \mathcal{L}$。使用 Lemma
\ref{equiv-lemma-sibling-fully-faithful}，可知 $F$ 与 $F'$ 是同胞函子。为说明 $f$
在 $X$ 的底拓扑空间上为恒等映射，使用 $F(\mathcal{O}_x) \cong \mathcal{O}_x$
以及 $\mathcal{O}_x$ 的支撑为 $\{x\}$ 这一事实。证明完毕。
\end{proof}

\begin{lemma}
\label{equiv-lemma-two-functors}
设 $k$ 为域。设 $X$、$Y$ 为 $k$ 上的固有概形。假设 $X$ 正则。设
$F, G : D_{perf}(\mathcal{O}_X) \to D_{perf}(\mathcal{O}_Y)$ 为
$k$-线性正合函子，且满足
\begin{enumerate}
\item $F(\mathcal{F}) \cong G(\mathcal{F})$ 对任意凝聚
$\mathcal{O}_X$-模 $\mathcal{F}$（满足 $\dim(\text{Supp}(\mathcal{F})) = 0$）成立；
\item $F$ 全忠实；
\item $G$ 是 Fourier--Mukai 函子，其核属于 $D_{perf}(\mathcal{O}_{X \times Y})$。
\end{enumerate}
则存在 Fourier--Mukai 函子
$F' : D_{perf}(\mathcal{O}_X) \to D_{perf}(\mathcal{O}_Y)$
，其核属于 $D_{perf}(\mathcal{O}_{X \times Y})$，并且 $F$ 与 $F'$ 是同胞函子。
\end{lemma}

\begin{proof}
由 Lemma \ref{equiv-lemma-two-functors-pre}，$G$ 的本质像包含于 $F$ 的本质像。
考虑函子 $H = F^{-1} \circ G$；由于 $F$ 全忠实，该函子有意义。由 Lemma
\ref{equiv-lemma-noah}，得到自同构 $f : X \to X$ 与可逆 $\mathcal{O}_X$-模
$\mathcal{L}$，使得函子 $H' : K \mapsto f^*K \otimes \mathcal{L}$ 是 $H$ 的同胞。
特别地，由 Lemma \ref{equiv-lemma-sibling-faithful}，$H$ 是自等价；并且 $H$ 诱导
$\textit{Coh}(\mathcal{O}_X)$ 的自等价（因为其同胞函子 $H'$ 具有这一性质）。
因此拟逆 $H^{-1}$ 与 $(H')^{-1}$ 存在且互为同胞（略去一个小细节），并且
$(H')^{-1}$ 把 $M$ 映为
$(f^{-1})^*(M \otimes_{\mathcal{O}_X}^\mathbf{L} \mathcal{L}^{\otimes -1})$
；这是 Fourier--Mukai 函子（细节从略）。于是当然，$F = G \circ H^{-1}$ 是
$G \circ (H')^{-1}$ 的同胞。由 Lemma \ref{equiv-lemma-compose-fourier-mukai}，
Fourier--Mukai 函子的复合仍为 Fourier--Mukai 函子，结论得证。
\end{proof}





\section{全忠实函子}
\label{equiv-section-fully-faithful}

\noindent
依照 \cite{Orlov-K3} 与 \cite{Ballard}，我们的目标是证明导出范畴之间的
全忠实函子是 Fourier--Mukai 函子的同胞。

\begin{situation}
\label{equiv-situation-fully-faithful}
这里 $k$ 为域。设 $X$ 与 $Y$ 为 $k$ 上的光滑固有概形。给定 $k$-线性、
正合且全忠实的函子
$F : D_{perf}(\mathcal{O}_X) \to D_{perf}(\mathcal{O}_Y)$。
\end{situation}

\noindent
继续阅读之前，最好至少阅读 Derived Categories, Section
\ref{derived-section-postnikov} 的一部分。

\medskip\noindent
回忆 $X$ 正则，因而具有解消性质（Varieties, Lemma
\ref{varieties-lemma-smooth-regular} 与 Derived Categories of Schemes, Lemma
\ref{perfect-lemma-regular-resolution-property}）。因此可在 $X \times X$ 上选取解消
$$
\ldots \to
\mathcal{E}_2 \boxtimes \mathcal{G}_2 \to
\mathcal{E}_1 \boxtimes \mathcal{G}_1 \to
\mathcal{E}_0 \boxtimes \mathcal{G}_0 \to
\mathcal{O}_\Delta \to 0
$$
，其中每个 $\mathcal{E}_i$ 与 $\mathcal{G}_i$ 都是有限局部自由
$\mathcal{O}_X$-模；参见 Lemma \ref{equiv-lemma-diagonal-resolution}。使用复形
\begin{equation}
\label{equiv-equation-original-complex}
\ldots \to
\mathcal{E}_2 \boxtimes \mathcal{G}_2 \to
\mathcal{E}_1 \boxtimes \mathcal{G}_1 \to
\mathcal{E}_0 \boxtimes \mathcal{G}_0
\end{equation}
（属于 $D_{perf}(\mathcal{O}_{X \times X})$），如 Derived Categories, Example
\ref{derived-example-key-postnikov} 所述；若对每个 $n$ 记
$$
M_n = (\mathcal{E}_n \boxtimes \mathcal{G}_n \to \ldots \to
\mathcal{E}_0 \boxtimes \mathcal{G}_0)[-n]
$$
，便得到复形 (\ref{equiv-equation-original-complex}) 的无穷 Postnikov 系统。这意味着态射
$M_0 \to M_1[1] \to M_2[2] \to \ldots$、
$M_n \to \mathcal{E}_n \boxtimes \mathcal{G}_n$ 以及
$\mathcal{E}_n \boxtimes \mathcal{G}_n \to M_{n - 1}$ 满足 Derived Categories,
Definition \ref{derived-definition-postnikov-system} 中记载的若干条件。令
$$
\mathcal{F}_n = \Ker(\mathcal{E}_n \boxtimes \mathcal{G}_n \to
\mathcal{E}_{n - 1} \boxtimes \mathcal{G}_{n - 1})
$$
。注意，由于 $\mathcal{O}_\Delta$ 在 $X$ 上平坦（通过 $\text{pr}_1$），
$\mathcal{F}_n$ 对所有 $n$ 也具有同一性质（这是方便但非必要的观察）。有
$$
H^q(M_n[n]) = \left\{
\begin{matrix}
\mathcal{O}_\Delta & \text{若} & q = 0 \\
\mathcal{F}_n & \text{若} & q = -n \\
0 & \text{若} & q \not = 0, -n
\end{matrix}
\right.
$$
因此对 $n \geq \dim(X \times X)$，有
$$
M_n[n] \cong \mathcal{O}_\Delta \oplus \mathcal{F}_n[n]
$$
；该同构在 $D_{perf}(\mathcal{O}_{X \times X})$ 中成立，见
Lemma \ref{equiv-lemma-split-complex-regular}。

\medskip\noindent
我们关心复形
\begin{equation}
\label{equiv-equation-complex}
\ldots \to
\mathcal{E}_2 \boxtimes F(\mathcal{G}_2) \to
\mathcal{E}_1 \boxtimes F(\mathcal{G}_1) \to
\mathcal{E}_0 \boxtimes F(\mathcal{G}_0)
\end{equation}
（属于 $D_{perf}(\mathcal{O}_{X \times Y})$），因为该复形的“总体化”应当给出
我们试图构造的 Fourier--Mukai 函子的核。对所有 $i, j \geq 0$，有
\begin{align*}
\Ext^q_{X \times Y}(\mathcal{E}_i \boxtimes F(\mathcal{G}_i),
\mathcal{E}_j \boxtimes F(\mathcal{G}_j))
& =
\bigoplus\nolimits_p
\Ext^{q + p}_X(\mathcal{E}_i, \mathcal{E}_j) \otimes_k
\Ext^{-p}_Y(F(\mathcal{G}_i), F(\mathcal{G}_j)) \\
& =
\bigoplus\nolimits_p
\Ext^{q + p}_X(\mathcal{E}_i, \mathcal{E}_j) \otimes_k
\Ext^{-p}_X(\mathcal{G}_i, \mathcal{G}_j)
\end{align*}
。第二个等号成立是因为 $F$ 全忠实，第一个等号则由 Derived Categories of Schemes,
Lemma \ref{perfect-lemma-kunneth-Ext} 得到。可知这些 $\Ext^q$ 对 $q < 0$ 为零。
因此由 Derived Categories, Lemma \ref{derived-lemma-existence-postnikov-system}，
可为复形 (\ref{equiv-equation-complex}) 构造无穷 Postnikov 系统
$K_0, K_1, K_2, \ldots$（位于 $D_{perf}(\mathcal{O}_{X \times Y})$ 中）。与
$M_0, M_1, M_2, \ldots$ 的情形平行，这意味着得到态射
$K_0 \to K_1[1] \to K_2[2] \to \ldots$、
$K_n \to \mathcal{E}_n \boxtimes F(\mathcal{G}_n)$ 以及
$\mathcal{E}_n \boxtimes F(\mathcal{G}_n) \to K_{n - 1}$；它们属于
$D_{perf}(\mathcal{O}_{X \times Y})$，并满足 Derived Categories, Definition
\ref{derived-definition-postnikov-system} 中记载的若干条件。

\medskip\noindent
设 $\mathcal{F}$ 为凝聚 $\mathcal{O}_X$-模，其支撑只含有限多个点；亦即
$\dim(\text{Supp}(\mathcal{F})) = 0$。考虑三角范畴之间的正合函子
$$
D_{perf}(\mathcal{O}_{X \times Y})
\longrightarrow
D_{perf}(\mathcal{O}_Y),\quad
N \longmapsto R\text{pr}_{2, *}(\text{pr}_1^*\mathcal{F}
\otimes^\mathbf{L}_{\mathcal{O}_{X \times Y}} N)
$$
于是对象 $R\text{pr}_{2, *}(\text{pr}_1^*\mathcal{F}
\otimes^\mathbf{L}_{\mathcal{O}_{X \times Y}} K_i)$ 构成一个 Postnikov 系统，
对应于 $D_{perf}(\mathcal{O}_Y)$ 中以如下对象为项的复形：
$$
R\text{pr}_{2, *}(
(\mathcal{F} \otimes \mathcal{E}_i) \boxtimes F(\mathcal{G}_i)) =
\Gamma(X, \mathcal{F} \otimes \mathcal{E}_i) \otimes_k F(\mathcal{G}_i) =
F(\Gamma(X, \mathcal{F} \otimes \mathcal{E}_i) \otimes_k \mathcal{G}_i)
$$
。这里使用了 $\mathcal{F} \otimes \mathcal{E}_i$ 的高次上同调消失，因为其支撑
维数为 $0$。另一方面，应用正合函子
$$
D_{perf}(\mathcal{O}_{X \times X})
\longrightarrow
D_{perf}(\mathcal{O}_Y),\quad
N \longmapsto F(R\text{pr}_{2, *}(\text{pr}_1^*\mathcal{F}
\otimes^\mathbf{L}_{\mathcal{O}_{X \times X}} N))
$$
，可知对象
$F(R\text{pr}_{2, *}(\text{pr}_1^*\mathcal{F}
\otimes^\mathbf{L}_{\mathcal{O}_{X \times X}} M_n))$
构成第二个无穷 Postnikov 系统，对应于 $D_{perf}(\mathcal{O}_Y)$ 中以如下对象为项的复形：
$$
F(R\text{pr}_{2, *}(
(\mathcal{F} \otimes \mathcal{E}_i) \boxtimes \mathcal{G}_i)) =
F(\Gamma(X, \mathcal{F} \otimes \mathcal{E}_i) \otimes_k \mathcal{G}_i)
$$
。这与前面相同！Postnikov 系统的唯一性（Derived Categories, Lemma
\ref{derived-lemma-existence-postnikov-system}）可以应用，因为
$$
\Ext^q_Y(
F(\Gamma(X, \mathcal{F} \otimes \mathcal{E}_i) \otimes_k \mathcal{G}_i),
F(\Gamma(X, \mathcal{F} \otimes \mathcal{E}_j) \otimes_k \mathcal{G}_j)) = 0,
\quad q < 0
$$
；这里使用 $F$ 全忠实。因此得到一族同构
$$
F(R\text{pr}_{2, *}(\text{pr}_1^*\mathcal{F}
\otimes^\mathbf{L}_{\mathcal{O}_{X \times X}} M_n[n]))
\cong
R\text{pr}_{2, *}(\text{pr}_1^*\mathcal{F}
\otimes^\mathbf{L}_{\mathcal{O}_{X \times Y}} K_n[n])
$$
（位于 $D_{perf}(\mathcal{O}_Y)$ 中），它们与 $D_{perf}(\mathcal{O}_Y)$ 中由下列态射
诱导的态射相容：
$$
M_{n - 1}[n - 1] \to M_n[n]
\quad\text{且}\quad
K_{n - 1}[n - 1] \to K_n[n]
$$
$$
M_n \to \mathcal{E}_n \boxtimes \mathcal{G}_n
\quad\text{且}\quad
K_n \to \mathcal{E}_n \boxtimes F(\mathcal{G}_n)
$$
$$
\mathcal{E}_n \boxtimes \mathcal{G}_n \to M_{n - 1}
\quad\text{且}\quad
\mathcal{E}_n \boxtimes F(\mathcal{G}_n) \to K_{n - 1}
$$
；这些态射是 Postnikov 系统结构的一部分。对充分大的 $n$，得到直和分解
$$
F(R\text{pr}_{2, *}(\text{pr}_1^*\mathcal{F}
\otimes^\mathbf{L}_{\mathcal{O}_{X \times X}} M_n[n]))
=
F(\mathcal{F}) \oplus
F(R\text{pr}_{2, *}(
\text{pr}_1^*\mathcal{F} \otimes_{\mathcal{O}_{X \times Y}} \mathcal{F}_n
))[n]
$$
，它对应于上面构造的 $M_n$ 的直和分解（这里使用 $\mathcal{F}_n$ 在 $X$ 上的
平坦性，即通过 $\text{pr}_1$ 的平坦性，以便在上述公式中写普通张量积；但这对论证
并非必要）。由 Lemma \ref{equiv-lemma-boundedness}，可知存在整数 $m \geq 0$，使该直和
分解的第一个直和项只在区间 $[-m, m]$ 中可能有非零上同调层，而第二个直和项只在
区间 $[-m - n, m + \dim(X) - n]$ 中可能有非零上同调层。因此系统
$K_0 \to K_1[1] \to K_2[2] \to \ldots$（位于
$D_{perf}(\mathcal{O}_{X \times Y})$ 中）满足 Lemma
\ref{equiv-lemma-bounded-fibres} 的假设，必要时可把 $m$ 换成更大的整数。于是可写成
$$
K_n[n] = K \oplus C_n
$$
，对 $n \gg 0$ 成立，并与转移映射相容；且 $C_n$ 只在范围
$[-m - n, m - n]$ 中可能有非零上同调层。以 $G$ 表示对应于 $K$ 的
Fourier--Mukai 函子。综合以上结果，得到
$$
\begin{matrix}
G(\mathcal{F}) \oplus
R\text{pr}_{2, *}(
\text{pr}_1^*\mathcal{F} \otimes_{\mathcal{O}_{X \times Y}}^\mathbf{L} C_n)
\cong \\
R\text{pr}_{2, *}(\text{pr}_1^*\mathcal{F}
\otimes^\mathbf{L}_{\mathcal{O}_{X \times Y}} K_n[n]) \cong \\
F(R\text{pr}_{2, *}(\text{pr}_1^*\mathcal{F}
\otimes^\mathbf{L}_{\mathcal{O}_{X \times X}} M_n[n]))
\cong \\
F(\mathcal{F}) \oplus
F(R\text{pr}_{2, *}(
\text{pr}_1^*\mathcal{F} \otimes_{\mathcal{O}_{X \times Y}} \mathcal{F}_n
))[n]
\end{matrix}
$$
观察这些对象所处的次数，可知当 $n \gg m$ 时得到同构
$$
F(\mathcal{F}) \cong G(\mathcal{F})
$$
。此外，回忆这对每个凝聚 $\mathcal{F}$（定义在 $X$ 上且支撑维数为 $0$）都成立。

\begin{lemma}
\label{equiv-lemma-fully-faithful}
设 $k$ 为域。设 $X$ 与 $Y$ 为 $k$ 上的光滑固有概形。给定 $k$-线性、
正合且全忠实的函子
$F : D_{perf}(\mathcal{O}_X) \to D_{perf}(\mathcal{O}_Y)$
，存在 Fourier--Mukai 函子
$F' : D_{perf}(\mathcal{O}_X) \to D_{perf}(\mathcal{O}_Y)$，其核属于
$D_{perf}(\mathcal{O}_{X \times Y})$，并且它是 $F$ 的同胞。
\end{lemma}

\begin{proof}
把 Lemma \ref{equiv-lemma-two-functors} 应用于 $F$ 与上面构造的函子 $G$。
\end{proof}

\noindent
下述定理在不假设 $X$ 射影时也成立；参见 \cite{Noah}。

\begin{theorem}[Orlov]
\label{equiv-theorem-fully-faithful}
\begin{reference}
\cite[Theorem 2.2]{Orlov-K3}；\cite{Noah} 在不假设 $X$ 射影的情况下证明了这一点
\end{reference}
设 $k$ 为域。设 $X$ 与 $Y$ 为光滑固有 $k$-概形，并且 $X$ 在 $k$ 上射影。
任意 $k$-线性、全忠实的正合函子
$F : D_{perf}(\mathcal{O}_X) \to D_{perf}(\mathcal{O}_Y)$ 是 Fourier--Mukai 函子，
其某个核属于 $D_{perf}(\mathcal{O}_{X \times Y})$。
\end{theorem}

\begin{proof}
依 Lemma \ref{equiv-lemma-fully-faithful}，令 $F'$ 为 $F$ 的同胞 Fourier--Mukai 函子。
要由 Proposition \ref{equiv-proposition-siblings-isomorphic} 得到 $F \cong F'$，
只需证明 $\textit{Coh}(\mathcal{O}_X)$ 有足够多负对象。然而，例如当
$X = \Spec(k)$ 时，这并不成立。因此先把 $X = \coprod X_i$ 分解为其连通
（亦即不可约）分支，并论证只需对每个（全忠实的）复合函子证明结论：
$$
F_i :
D_{perf}(\mathcal{O}_{X_i}) \to
D_{perf}(\mathcal{O}_X) \to
D_{perf}(\mathcal{O}_Y)
$$
。细节从略。因此可假设 $X$ 不可约。

\medskip\noindent
情形 $\dim(X) = 0$。此时 $X$ 是有限（可分）扩张 $k'/k$ 的谱，因而
$D_{perf}(\mathcal{O}_X)$ 等价于分次 $k'$-向量空间的范畴；在此等价下，
$\mathcal{O}_X$ 对应于平凡的 $1$-维向量空间，位于次数 $0$。容易看出任意两个同胞函子
$F, F' : D_{perf}(\mathcal{O}_X) \to D_{perf}(\mathcal{O}_Y)$ 都同构。事实上，
我们有同构 $F(\mathcal{O}_X) \cong F'(\mathcal{O}_X)$，它与 $k$-代数
$k' = \text{End}_{D_{perf}(\mathcal{O}_X)}(\mathcal{O}_X)$ 的作用相容；该同构
典范地延拓到任意分次 $k'$-向量空间上的同构。

\medskip\noindent
情形 $\dim(X) > 0$。此时 $X$ 是维数 $> 1$ 的光滑射影簇。设
$\mathcal{F}$ 为凝聚 $\mathcal{O}_X$-模。须证明存在凝聚模 $\mathcal{N}$，使得
\begin{enumerate}
\item 存在满射 $\mathcal{N} \to \mathcal{F}$；并且
\item $\Hom(\mathcal{F}, \mathcal{N}) = 0$.
\end{enumerate}
选取丰沛可逆 $\mathcal{O}_X$-模 $\mathcal{L}$。我们断言
$\mathcal{N} = (\mathcal{L}^{\otimes n})^{\oplus r}$ 在 $n \ll 0$ 且 $r$ 足够大时
满足要求。条件 (1) 由 Properties, Proposition
\ref{properties-proposition-characterize-ample} 得到。最后，有
$$
\Hom(\mathcal{F}, \mathcal{L}^{\otimes n}) =
H^0(X, \SheafHom(\mathcal{F}, \mathcal{L}^{\otimes n})) =
H^0(X, \SheafHom(\mathcal{F}, \mathcal{O}_X) \otimes \mathcal{L}^{\otimes n})
$$
。由于对偶 $\SheafHom(\mathcal{F}, \mathcal{O}_X)$ 无挠，由 Varieties, Lemma
\ref{varieties-lemma-vanishin-h0-negative}，它对 $n \ll 0$ 消失。证明完毕。
\end{proof}

\begin{proposition}
\label{equiv-proposition-equivalence}
设 $k$ 为域。设 $X$ 与 $Y$ 为 $k$ 上的光滑固有概形。若
$F : D_{perf}(\mathcal{O}_X) \to D_{perf}(\mathcal{O}_Y)$ 是三角范畴之间的
$k$-线性正合等价，则存在 Fourier--Mukai 函子
$F' : D_{perf}(\mathcal{O}_X) \to D_{perf}(\mathcal{O}_Y)$，
其核属于 $D_{perf}(\mathcal{O}_{X \times Y})$，且该函子既是等价又是 $F$ 的同胞。
\end{proposition}

\begin{proof}
Lemma \ref{equiv-lemma-fully-faithful} 的函子 $F'$ 由 Lemma
\ref{equiv-lemma-sibling-faithful} 可知是等价。
\end{proof}

\begin{lemma}
\label{equiv-lemma-uniqueness}
设 $k$ 为域。设 $X$ 为 $k$ 上的光滑固有概形。设
$K \in D_{perf}(\mathcal{O}_{X \times X})$。若 Fourier--Mukai 函子
$\Phi_K : D_{perf}(\mathcal{O}_X) \to D_{perf}(\mathcal{O}_X)$ 与恒等函子同构，
则 $K \cong \Delta_*\mathcal{O}_X$ 在 $_{perf}(\mathcal{O}_{X \times X})$ 中成立。
\end{lemma}

\begin{proof}
令 $i$ 为使上同调层 $H^i(K)$ 非零的最小整数。设 $\mathcal{E}$ 与
$\mathcal{G}$ 为有限局部自由 $\mathcal{O}_X$-模。则
\begin{align*}
H^i(X \times X, K \otimes_{\mathcal{O}_{X \times X}}^\mathbf{L}
(\mathcal{E} \boxtimes \mathcal{G}))
& =
H^i(X, R\text{pr}_{2, *}(K \otimes_{\mathcal{O}_{X \times X}}^\mathbf{L}
(\mathcal{E} \boxtimes \mathcal{G}))) \\
& =
H^i(X, \Phi_K(\mathcal{E}) \otimes_{\mathcal{O}_X}^\mathbf{L} \mathcal{G}) \\
& \cong
H^i(X, \mathcal{E} \otimes \mathcal{G})
\end{align*}
；当 $i < 0$ 时它为零。另一方面，可选取 $\mathcal{E}$ 与 $\mathcal{G}$，使得存在满射
$\mathcal{E}^\vee \boxtimes \mathcal{G}^\vee \to H^i(K)$
；这由 Lemma \ref{equiv-lemma-on-product} 得到。在此情形中，上述等式的左端非零。
因此得到 $H^i(K) = 0$ 对 $i < 0$ 成立。

\medskip\noindent
令 $i$ 为使 $H^i(K)$ 非零的最大整数。同一论证令 $\mathcal{E}$ 与
$\mathcal{G}$ 的支撑维数为 $0$，得到 $i \leq 0$。因此可知 $K$ 由单个凝聚
$\mathcal{O}_{X \times X}$-模 $\mathcal{K}$ 给出，它位于次数 $0$。

\medskip\noindent
由于 $R\text{pr}_{2, *}(\text{pr}_1^*\mathcal{F} \otimes \mathcal{K})$
就是 $\mathcal{F}$，取支撑在闭点上的 $\mathcal{F}$，可知 $\mathcal{K}$ 的支撑
在 $X$ 上经 $\text{pr}_2$ 有限。由于
$R\text{pr}_{2, *}(\mathcal{K}) \cong \mathcal{O}_X$，由 Functors and Morphisms,
Lemma \ref{functors-lemma-pushforward-invertible-pre} 得到
$\mathcal{K} = s_*\mathcal{O}_X$，其中 $s : X \to X \times X$ 是第二投影的某个截面。
于是 $\Phi_K(M) = f^*M$，其中 $f = \text{pr}_1 \circ s$；这只有在 $s$ 是对角态射时
才可能发生，正是所求。
\end{proof}








\section{Fourier--Mukai 核的范畴}
\label{equiv-section-category-Fourier-Mukai-kernels}

\noindent
设 $S$ 为概形。我们断言存在如下范畴：
\begin{enumerate}
\item 对象为 $S$ 上的光滑固有概形。
\item 从 $X$ 到 $Y$ 的态射为 $D_{perf}(\mathcal{O}_{X \times_S Y})$
中对象的同构类。
\item $K \in D_{perf}(\mathcal{O}_{X \times_S Y})$ 的同构类与 $K'$ 在
$D_{perf}(\mathcal{O}_{Y \times_S Z})$ 中的同构类之复合，是下列对象的同构类：
$$
R\text{pr}_{13, *}(
L\text{pr}_{12}^*K
\otimes_{\mathcal{O}_{X \times_S Y \times_S Z}}^\mathbf{L}
L\text{pr}_{23}^*K')
$$
由 Derived Categories of Schemes, Lemma
\ref{perfect-lemma-flat-proper-perfect-direct-image-general}，它属于
$D_{perf}(\mathcal{O}_{X \times_S Z})$。
\item 从 $X$ 到 $X$ 的恒等态射是 $\Delta_{X/S, *}\mathcal{O}_X$ 的同构类。
由 More on Morphisms, Lemma
\ref{more-morphisms-lemma-perfect-closed-immersion-perfect-direct-image}
以及下述事实，它属于 $D_{perf}(\mathcal{O}_{X \times_S X})$：由 Divisors, Lemma
\ref{divisors-lemma-immersion-smooth-into-smooth-regular-immersion} 与
More on Morphisms, Lemma \ref{more-morphisms-lemma-regular-immersion-perfect}，
$\Delta_{X/S}$ 是完美态射。
\end{enumerate}
下面验证态射复合的结合律；略去恒等态射确为恒等元的验证。为此，设有
$X, Y, Z, W$，以及
$c \in D_{perf}(\mathcal{O}_{X \times_S Y})$,
$c' \in D_{perf}(\mathcal{O}_{Y \times_S Z})$ 和
$c'' \in D_{perf}(\mathcal{O}_{Z \times_S W})$。则
\begin{align*}
c'' \circ (c' \circ c)
& \cong
\text{pr}^{134}_{14, *}(
\text{pr}^{134, *}_{13}
\text{pr}^{123}_{13, *}(\text{pr}^{123, *}_{12}c \otimes
\text{pr}^{123, *}_{23}c')
\otimes \text{pr}^{134, *}_{34}c'') \\
& \cong
\text{pr}^{134}_{14, *}(
\text{pr}^{1234}_{134, *}
\text{pr}^{1234, *}_{123}(\text{pr}^{123, *}_{12}c \otimes
\text{pr}^{123, *}_{23}c')
\otimes \text{pr}^{134, *}_{34}c'') \\
& \cong
\text{pr}^{134}_{14, *}(
\text{pr}^{1234}_{134, *}
(\text{pr}^{1234, *}_{12}c \otimes
\text{pr}^{1234, *}_{23}c')
\otimes \text{pr}^{134, *}_{34}c'') \\
& \cong
\text{pr}^{134}_{14, *}
\text{pr}^{1234}_{134, *}
((\text{pr}^{1234, *}_{12}c \otimes
\text{pr}^{1234, *}_{23}c')
\otimes \text{pr}^{1234, *}_{34}c'') \\
& \cong
\text{pr}^{1234}_{14, *}(
(\text{pr}^{1234, *}_{12}c \otimes
\text{pr}^{1234, *}_{23}c') \otimes
\text{pr}^{1234, *}_{34}c'')
\end{align*}
这里采用记号
$$
p^{1234}_{134} : X \times_S Y \times_S Z \times_S W
\to X \times_S Z \times_S W
\quad\text{且}\quad
p^{134}_{14} : X \times_S Z \times_S W \to X \times_S W
$$
表示这些投影；其他指标也采用类似记号。我们还用 $\text{pr}_*$ 代替
$R\text{pr}_*$，用 $\text{pr}^*$ 代替 $L\text{pr}^*$，并省略
$\otimes$ 的所有上标与下标。第一个等式就是复合的定义。第二个等式成立，是因为
$\text{pr}^{134, *}_{13} \text{pr}^{123}_{13, *} =
\text{pr}^{1234}_{134, *} \text{pr}^{1234, *}_{123}$
由基变换成立（Derived Categories of Schemes, Lemma
\ref{perfect-lemma-compare-base-change}）。第三个等式成立，是因为拉回能正确复合并与张量积相容；
见 Cohomology, Lemmas \ref{cohomology-lemma-derived-pullback-composition} 与
\ref{cohomology-lemma-pullback-tensor-product}。第四个等式由 $p^{1234}_{134}$ 的“投影公式”得到；
见 Derived Categories of Schemes, Lemma
\ref{perfect-lemma-cohomology-base-change}。第五个等式说明固有推前与复合相容；见
Cohomology, Lemma \ref{cohomology-lemma-derived-pushforward-composition}。
再由张量积的结合律，即证得复合的结合律。

\begin{lemma}
\label{equiv-lemma-base-change-is-functor}
设 $S' \to S$ 为概形态射。下述规则
\begin{enumerate}
\item 把光滑固有概形 $X$（定义在 $S$ 上）映为 $X' =  S' \times_S X$；并且
\item 把对象 $K$（属于 $D_{perf}(\mathcal{O}_{X \times_S Y})$）的同构类映为
$L(X' \times_{S'} Y' \to X \times_S Y)^*K$
在 $D_{perf}(\mathcal{O}_{X' \times_{S'} Y'})$ 中的同构类，
\end{enumerate}
给出从为 $S$ 定义的范畴到为 $S'$ 定义的范畴的函子。
\end{lemma}

\begin{proof}
为验证此点，设有 $X, Y, Z$，以及
$K \in D_{perf}(\mathcal{O}_{X \times_S Y})$ 和
$M \in D_{perf}(\mathcal{O}_{Y \times_S Z})$。记
$K' \in D_{perf}(\mathcal{O}_{X' \times_{S'} Y'})$ 和
$M' \in D_{perf}(\mathcal{O}_{Y' \times_{S'} Z'})$
为该引理陈述中的相应拉回。图表
$$
\xymatrix{
X' \times_{S'} Y' \times_{S'} Z' \ar[r] \ar[d]_{\text{pr}'_{13}} &
X \times_S Y \times_S Z \ar[d]^{\text{pr}_{13}} \\
X' \times_{S'} Z' \ar[r] &
X \times_S Z
}
$$
是笛卡儿的，并且 $\text{pr}_{13}$ 光滑且固有。由 Derived Categories of Schemes, Lemma
\ref{perfect-lemma-flat-proper-perfect-direct-image-general}
可知，下方水平箭头对下列复合的导出拉回
$$
R\text{pr}_{13, *}(
L\text{pr}_{12}^*K
\otimes_{\mathcal{O}_{X \times_S Y \times_S Z}}^\mathbf{L}
L\text{pr}_{23}^*M)
$$
确实（典范地）同构于
$$
R\text{pr}'_{13, *}(
L(\text{pr}'_{12})^*K'
\otimes_{\mathcal{O}_{X' \times_{S'} Y' \times_{S'} Z'}}^\mathbf{L}
L(\text{pr}'_{23})^*M')
$$
，正如所需。略去若干细节。
\end{proof}





\section{相对等价}
\label{equiv-section-relative-equivalences}

\noindent
本节证明关于下述概念的若干引理。

\begin{definition}
\label{equiv-definition-relative-equivalence-kernel}
设 $S$ 为概形，并设 $X \to S$ 与 $Y \to S$ 为光滑固有态射。称对象
$K \in D_{perf}(\mathcal{O}_{X \times_S Y})$ 为{\it 从 $X$ 到 $Y$、定义在 $S$ 上
的相对等价的 Fourier--Mukai 核}，如果存在对象
$K' \in D_{perf}(\mathcal{O}_{X \times_S Y})$，使得
$$
\Delta_{X/S, *}\mathcal{O}_X \cong
R\text{pr}_{13, *}(L\text{pr}_{12}^*K
\otimes_{\mathcal{O}_{X \times_S Y \times_S X}}^\mathbf{L}
L\text{pr}_{23}^*K')
$$
在 $D(\mathcal{O}_{X \times_S X})$ 中成立，并且
$$
\Delta_{Y/S, *}\mathcal{O}_Y \cong
R\text{pr}_{13, *}(L\text{pr}_{12}^*K'
\otimes_{\mathcal{O}_{Y \times_S X \times_S Y}}^\mathbf{L}
L\text{pr}_{23}^*K)
$$
在 $D(\mathcal{O}_{Y \times_S Y})$ 中成立。换言之，$K$ 的同构类在
Section \ref{equiv-section-category-Fourier-Mukai-kernels} 所定义的范畴中给出可逆箭头。
\end{definition}

\noindent
这里刻意采用了较为繁复的说法。

\begin{lemma}
\label{equiv-lemma-equivalences-rek}
沿用 Definition \ref{equiv-definition-relative-equivalence-kernel} 的记号，设 $K$ 为从 $X$
到 $Y$、定义在 $S$ 上的相对等价的 Fourier--Mukai 核。则相应的 Fourier--Mukai 函子
$\Phi_K : D_\QCoh(\mathcal{O}_X) \to D_\QCoh(\mathcal{O}_Y)$
（Lemma \ref{equiv-lemma-fourier-Mukai-QCoh}）以及
$\Phi_K : D_{perf}(\mathcal{O}_X) \to D_{perf}(\mathcal{O}_Y)$
（Lemma \ref{equiv-lemma-fourier-mukai}）都是等价。
\end{lemma}

\begin{proof}
由 Lemma \ref{equiv-lemma-compose-fourier-mukai} 与
Example \ref{equiv-example-diagonal-fourier-mukai} 立即得到。
\end{proof}

\begin{lemma}
\label{equiv-lemma-base-change-rek}
沿用 Definition \ref{equiv-definition-relative-equivalence-kernel} 的记号，设 $K$ 为从 $X$
到 $Y$、定义在 $S$ 上的相对等价的 Fourier--Mukai 核。设 $S_1 \to S$ 为概形态射，并令
$X_1 = S_1 \times_S X$ 与 $Y_1 = S_1 \times_S Y$。则拉回
$K_1 = L(X_1 \times_{S_1} Y_1 \to X \times_S Y)^*K$ 是从 $X_1$
到 $Y_1$、定义在 $S_1$ 上的相对等价的 Fourier--Mukai 核。
\end{lemma}

\begin{proof}
设 $K' \in D_{perf}(\mathcal{O}_{Y \times_S X})$ 为 Definition
\ref{equiv-definition-relative-equivalence-kernel} 中假定存在的对象。记 $K'_1$ 为 $K'$ 沿
$Y_1 \times_{S_1} X_1 \to Y \times_S X$ 的拉回。只需证明
$$
\Delta_{X_1/S_1, *}\mathcal{O}_X \cong
R\text{pr}_{13, *}(L\text{pr}_{12}^*K_1
\otimes_{\mathcal{O}_{X_1 \times_{S_1} Y_1 \times_{S_1} X_1}}^\mathbf{L}
L\text{pr}_{23}^*K_1')
$$
在 $D(\mathcal{O}_{X_1 \times_{S_1} X_1})$ 中成立；另一个条件同理。由于
$$
\xymatrix{
X_1 \times_{S_1} Y_1 \times_{S_1} X_1 \ar[r] \ar[d]_{\text{pr}_{13}} &
X \times_S Y \times_S X \ar[d]^{\text{pr}_{13}} \\
X_1 \times_{S_1} X_1 \ar[r] &
X \times_S X
}
$$
是笛卡儿图，由 Derived Categories of Schemes, Lemma
\ref{perfect-lemma-flat-proper-perfect-direct-image-general}
可知，只需证明
$$
\Delta_{X_1/S_1, *}\mathcal{O}_{X_1}
\cong
L(X_1 \times_{S_1} X_1 \to X \times_S X)^*\Delta_{X/S, *}\mathcal{O}_X 
$$
而只要 $X$ 与 $X_1 \times_{S_1} X_1$ 在 $X \times_S X$ 上 Tor 无关，此式便成立；
见 Derived Categories of Schemes, Lemma \ref{perfect-lemma-compare-base-change}。
这种 Tor 无关性可直接看出，也可由更一般的 More on Morphisms, Lemma
\ref{more-morphisms-lemma-case-of-tor-independence} 得到：把它应用于四个顶点为
$X, X, X, S$ 的方块及其沿 $S_1 \to S$ 的基变换即可。
\end{proof}

\begin{lemma}
\label{equiv-lemma-descend-rek}
设 $S = \lim_{i \in I} S_i$ 是概形的有向系统的极限，其转移态射
$g_{i'i} : S_{i'} \to S_i$ 都是仿射的。假设 $S_i$ 对每个 $i \in I$ 都拟紧且拟分离。
取 $0 \in I$。设 $X_0 \to S_0$ 与 $Y_0 \to S_0$ 为光滑固有态射。令
$X_i = S_i \times_{S_0} X_0$（$i \geq 0$），并令 $X = S \times_{S_0} X_0$；对 $Y_0$ 也作类似定义。
若 $K$ 是从 $X$ 到 $Y$、定义在 $S$ 上的相对等价的 Fourier--Mukai 核，则对某个 $i \geq 0$，
存在从 $X_i$ 到 $Y_i$、定义在 $S_i$ 上的相对等价的 Fourier--Mukai 核。
\end{lemma}

\begin{proof}
设 $K' \in D_{perf}(\mathcal{O}_{Y \times_S X})$ 为 Definition
\ref{equiv-definition-relative-equivalence-kernel} 中假定存在的对象。由于
$X \times_S Y = \lim X_i \times_{S_i} Y_i$，存在指标 $i$ 以及对象 $K_i$ 与 $K'_i$，
它们属于 $D_{perf}(\mathcal{O}_{Y_i \times_{S_i} X_i})$，
使其拉回到 $Y \times_S X$ 后分别给出 $K$ 与 $K'$；见 Derived Categories of Schemes,
Lemma \ref{perfect-lemma-descend-perfect}。由 Derived Categories of Schemes, Lemma
\ref{perfect-lemma-flat-proper-perfect-direct-image-general}
，对象
$$
R\text{pr}_{13, *}(L\text{pr}_{12}^*K_i
\otimes_{\mathcal{O}_{X_i \times_{S_i} Y_i \times_{S_i} X_i}}^\mathbf{L}
L\text{pr}_{23}^*K_i')
$$
是完美的，且其到 $X \times_S X$ 的拉回等于
$$
R\text{pr}_{13, *}(L\text{pr}_{12}^*K
\otimes_{\mathcal{O}_{X \times_S Y \times_S X}}^\mathbf{L}
L\text{pr}_{23}^*K') \cong \Delta_{X/S, *}\mathcal{O}_X
$$
；参见 Lemma \ref{equiv-lemma-base-change-rek} 的证明。另一方面，由于 $X_i \to S$ 光滑且分离，对象
$$
\Delta_{i, *}\mathcal{O}_{X_i}
$$
（属于 $D(\mathcal{O}_{X_i \times_{S_i} X_i})$）也是完美的（由 More on Morphisms, Lemmas
\ref{more-morphisms-lemma-smooth-diagonal-perfect} 与
\ref{more-morphisms-lemma-perfect-proper-perfect-direct-image}），且其到
$X \times_S X$ 的拉回等于
$$
\Delta_{X/S, *}\mathcal{O}_X
$$
；参见 Lemma \ref{equiv-lemma-base-change-rek} 的证明。因此由 Derived Categories of Schemes,
Lemma \ref{perfect-lemma-descend-perfect}，增大 $i$ 后可假设
$$
\Delta_{i, *}\mathcal{O}_{X_i} \cong
R\text{pr}_{13, *}(L\text{pr}_{12}^*K_i
\otimes_{\mathcal{O}_{X_i \times_{S_i} Y_i \times_{S_i} X_i}}^\mathbf{L}
L\text{pr}_{23}^*K_i')
$$
，正如所需。交换 $K$ 与 $K'$ 的角色可作同样论证。
\end{proof}








\section{无形变}
\label{equiv-section-no-deformations}

\noindent
本节标题所指的是 Lemma \ref{equiv-lemma-no-deformations}。

\begin{lemma}
\label{equiv-lemma-deform-koszul}
设 $(R, \mathfrak m, \kappa) \to (A, \mathfrak n, \lambda)$ 为本质有限表示的
局部环平坦局部同态。设
$\overline{f}_1, \ldots, \overline{f}_r \in \mathfrak n/\mathfrak m A
\subset A/\mathfrak m A$ 为正则序列，并设 $K \in D(A)$。假设
\begin{enumerate}
\item $K$ 是完美的；
\item $K \otimes_A^\mathbf{L} A/\mathfrak m A$ 在 $D(A/\mathfrak m A)$ 中同构于
$\overline{f}_1, \ldots, \overline{f}_r$ 上的 Koszul 复形。
\end{enumerate}
则 $K$ 在 $D(A)$ 中同构于某个正则序列 $f_1, \ldots, f_r \in A$ 上的 Koszul 复形，
该序列提升给定元素 $\overline{f}_1, \ldots, \overline{f}_r$。此外，
$A/(f_1, \ldots, f_r)$ 在 $R$ 上平坦。
\end{lemma}

\begin{proof}
本引理的证明采用链复形。Koszul 复形
$K_\bullet(\overline{f}_1, \ldots, \overline{f}_r)$ 定义于 More on Algebra,
Definition \ref{more-algebra-definition-koszul-complex}。由 More on Algebra, Lemma
\ref{more-algebra-lemma-lift-complex-stably-frees}，可用下列复形表示 $K$：
$$
K_\bullet :
A \to A^{\oplus r} \to \ldots \to A^{\oplus r} \to A
$$
$A/\mathfrak mA$ 与它的张量积恰好等于（！）
$K_\bullet(\overline{f}_1, \ldots, \overline{f}_r)$。记
$f_1, \ldots, f_r \in A$ 为箭头 $A^{\oplus r} \to A$ 的各分量。这些 $f_i$
是 $\overline{f}_i$ 的提升。由 Algebra, Lemma
\ref{algebra-lemma-grothendieck-regular-sequence-general}
，$f_1, \ldots, f_r$ 在 $A$ 中构成正则序列，并且 $A/(f_1, \ldots, f_r)$
在 $R$ 上平坦。令 $J = (f_1, \ldots, f_r) \subset A$。考虑图表
$$
\xymatrix{
K_\bullet \ar[rd] \ar@{..>}[rr]_{\varphi_\bullet} & &
K_\bullet(f_1, \ldots, f_r) \ar[ld] \\
& A/J
}
$$
由于 $f_1, \ldots, f_r$ 是正则序列，西南方向的箭头为拟同构（见 More on Algebra,
Lemma \ref{more-algebra-lemma-regular-koszul-regular}）。因此，例如由 Algebra,
Lemma \ref{algebra-lemma-compare-resolutions}，可找到使图表交换的虚线箭头。
模去 $\mathfrak m$，得到交换图表
$$
\xymatrix{
K_\bullet(\overline{f}_1, \ldots, \overline{f}_r)
\ar[rd] \ar[rr]_{\overline{\varphi}_\bullet} & &
K_\bullet(\overline{f}_1, \ldots, \overline{f}_r) \ar[ld] \\
& (A/\mathfrak m A)/(\overline{f}_1, \ldots, \overline{f}_r)
}
$$
；这是由 $K_\bullet$ 的选取所得。因此 $\overline{\varphi}$ 在导出范畴
$D(A/\mathfrak m A)$ 中为同构，从而
$\overline{\varphi} \otimes_{A/\mathfrak m A}^\mathbf{L} \lambda$ 为同构。由于
$\overline{f}_i \in \mathfrak n / \mathfrak m A$，可知
$$
\text{Tor}_i^{A/\mathfrak m A}(
K_\bullet(\overline{f}_1, \ldots, \overline{f}_r), \lambda)
=
K_i(\overline{f}_1, \ldots, \overline{f}_r) \otimes_{A/\mathfrak m A} \lambda
$$
所以 $\varphi_i \bmod \mathfrak n$ 可逆。由于 $A$ 是局部环，这说明 $\varphi_i$
是同构，证明完毕。
\end{proof}

\begin{lemma}
\label{equiv-lemma-limit-arguments}
设 $R \to S$ 为 Noether 环之间的有限型平坦环同态。设 $\mathfrak q \subset S$
为位于 $\mathfrak p \subset R$ 上方的素理想，并设 $K \in D(S)$ 为完美对象。
设 $f_1, \ldots, f_r \in \mathfrak q S_\mathfrak q$ 为正则序列，使得
$S_\mathfrak q/(f_1, \ldots, f_r)$ 在 $R$ 上平坦，并且
$K \otimes_S^\mathbf{L} S_\mathfrak q$ 同构于 $f_1, \ldots, f_r$ 上的 Koszul 复形。
则存在 $g \in S$，$g \not \in \mathfrak q$，使得
\begin{enumerate}
\item $f_1, \ldots, f_r$ 是 $f'_1, \ldots, f'_r \in S_g$ 的像；
\item $f'_1, \ldots, f'_r$ 在 $S_g$ 中构成正则序列；
\item $S_g/(f'_1, \ldots, f'_r)$ 在 $R$ 上平坦；
\item $K \otimes_S^\mathbf{L} S_g$ 同构于 $f_1, \ldots, f_r$ 上的 Koszul 复形。
\end{enumerate}
\end{lemma}

\begin{proof}
由局部化的定义，可找到满足 (1) 的 $g \in S$，$g \not \in \mathfrak q$。把 $g$
替换为 $gg'$，其中 $g' \in S$ 且 $g' \not \in \mathfrak q$，便可假设 (2) 成立；
见 Algebra, Lemma \ref{algebra-lemma-regular-sequence-in-neighbourhood}。由 Algebra,
Theorem \ref{algebra-theorem-openness-flatness}，$S_g/(f'_1, \ldots, f'_r)$
在 $R$ 上的平坦性在 $\mathfrak q$ 的某个开邻域中成立。因此再次把 $g$ 替换为
$gg'$，其中 $g' \in S$ 且 $g' \not \in \mathfrak q$，便可同时假设 (3) 成立。
最后，由 More on Algebra, Lemma \ref{more-algebra-lemma-colimit-perfect-complexes}，
再作一次替换即可得到 (4)。
\end{proof}

\noindent
下述引理的推广见 More on Morphisms of Spaces, Lemma
\ref{spaces-more-morphisms-lemma-where-isomorphism}.

\begin{lemma}
\label{equiv-lemma-isomorphism-in-neighbourhood}
设 $S$ 为 Noether 概形，$s \in S$。设 $p : X \to Y$ 为 $S$ 上的概形态射。假设
\begin{enumerate}
\item $Y \to S$ 与 $X \to S$ 都是固有态射；
\item $X$ 在 $S$ 上平坦；
\item $X_s \to Y_s$ 是同构。
\end{enumerate}
则存在开子集 $U \subset S$，它是 $s$ 的邻域，并且基变换 $X_U \to Y_U$ 是同构。
\end{lemma}

\begin{proof}
由 Morphisms, Lemma \ref{morphisms-lemma-closed-immersion-proper}，态射 $p$ 是固有的。
由 Cohomology of Schemes, Lemma
\ref{coherent-lemma-proper-finite-fibre-finite-in-neighbourhood}
，存在开子集 $Y_s \subset V \subset Y$，使得
$p|_{p^{-1}(V)} : p^{-1}(V) \to V$ 是有限态射。由 More on Morphisms, Theorem
\ref{more-morphisms-theorem-criterion-flatness-fibre-Noetherian}
，存在开子集 $X_s \subset U \subset X$，使得 $p|_U : U \to Y$ 平坦。
去掉 $X \setminus U$ 与 $Y \setminus V$ 的像（它们是不含 $s$ 的闭子集）后，
可假设 $p$ 平坦且有限。于是 $p$ 是开映射（Morphisms, Lemma
\ref{morphisms-lemma-fppf-open}），并且 $Y_s \subset p(X) \subset Y$；故缩小 $S$ 后，
可假设 $p$ 满射。由于 $p_s : X_s \to Y_s$ 是同构，态射
$$
p^\sharp : \mathcal{O}_Y \longrightarrow p_*\mathcal{O}_X
$$
是凝聚 $\mathcal{O}_Y$-模之间的态射（$p$ 有限）；沿 $i : Y_s \to Y$ 拉回后，
它成为同构（例如由 Cohomology of Schemes, Lemma
\ref{coherent-lemma-affine-base-change}）。Nakayama 引理于是说明
$\mathcal{O}_{Y, y} \to (p_*\mathcal{O}_X)_y$ 对所有 $y \in Y_s$ 都是满射。
因此存在开子集 $Y_s \subset V \subset Y$，使得 $p^\sharp|_V$ 满射
（Modules, Lemma \ref{modules-lemma-finite-type-surjective-on-stalk}）。故再次缩小 $S$ 后，
可假设 $p^\sharp$ 满射；这说明 $p$ 是闭浸入（因为 $p$ 已经有限）。于是 $p$
是 Noether 概形之间满射、平坦的闭浸入，因而是同构；见 Morphisms, Section
\ref{morphisms-section-flat-closed-immersions}。
\end{proof}

\begin{lemma}
\label{equiv-lemma-no-deformations}
设 $k$ 为域。设 $S$ 为 $k$ 上具有 $k$-有理点 $s$ 的有限型概形。设 $Y \to S$
为光滑固有态射。设 $X = Y_s \times S \to S$ 为纤维等于 $Y_s$ 的常值族。
设 $K$ 为从 $X$ 到 $Y$、定义在 $S$ 上的相对等价的 Fourier--Mukai 核。假设限制
$$
L(Y_s \times_S Y_s \to X \times_S Y)^*K \cong 
\Delta_{Y_s/k, *} \mathcal{O}_{Y_s}
$$
在 $D(\mathcal{O}_{Y_s \times Y_s})$ 中成立。则存在开邻域 $s \in U \subset S$，
使得 $Y|_U$ 同构于 $Y_s \times U$，且该同构定义在 $U$ 上。
\end{lemma}

\begin{proof}
记 $i : Y_s \times Y_s = X_s \times Y_s \to X \times_S Y$ 为自然闭浸入。
（从现在起，我们把 $Y_s$ 而非 $X_s$ 用作 $X$ 在 $s$ 上纤维的记号。）设
$z \in Y_s \times Y_s = (X \times_S Y)_s \subset X \times_S Y$ 为闭点。
如上所示，我们既把 $z$ 看作 $Y_s \times Y_s$ 的闭点，也把它看作
$X \times_S Y$ 的闭点。

\medskip\noindent
情形 I：$z \not \in \Delta_{Y_s/k}(Y_s)$。记 $\mathcal{O}_z$ 为凝聚
$\mathcal{O}_{Y_s \times Y_s}$-模，其支撑在 $z$ 上，值为 $\kappa(z)$。则
$i_*\mathcal{O}_z$ 是凝聚 $\mathcal{O}_{X \times_S Y}$-模，其支撑在 $z$ 上，
值为 $\kappa(z)$。我们的假设意味着
$$
K \otimes_{\mathcal{O}_{X \times_S Y}}^\mathbf{L} i_*\mathcal{O}_z =
Li^*K \otimes_{\mathcal{O}_{Y_s \times Y_s}}^\mathbf{L} \mathcal{O}_z = 0
$$
因此由 Lemma \ref{equiv-lemma-orthogonal-point-sheaf}，存在开集
$U(z) \subset X \times_S Y$，它是 $z$ 的邻域，并使得 $K|_{U(z)} = 0$。在此情形，令
$Z(z) = \emptyset$，把它视为 $U(z)$ 的闭子概形。

\medskip\noindent
情形 II：$z \in \Delta_{Y_s/k}(Y_s)$。由于 $Y_s$ 在 $k$ 上光滑，
$\Delta_{Y_s/k} : Y_s \to Y_s \times Y_s$ 是正则浸入；见 More on Morphisms,
Lemma \ref{more-morphisms-lemma-smooth-diagonal-perfect}。选取正则序列
$\overline{f}_1, \ldots, \overline{f}_r \in
\mathcal{O}_{Y_s \times Y_s, z}$，使它生成 $\Delta_{Y_s/k}(Y_s)$ 的理想层。
由于正则序列是 Koszul-正则的（More on Algebra, Lemma
\ref{more-algebra-lemma-regular-koszul-regular}），我们的假设意味着
$$
K_z \otimes_{\mathcal{O}_{X \times_S Y, z}}^\mathbf{L}
\mathcal{O}_{Y_s \times Y_s, z}
\in D(\mathcal{O}_{Y_s \times Y_s, z})
$$
由 $\overline{f}_1, \ldots, \overline{f}_r$ 在
$\mathcal{O}_{Y_s \times Y_s, z}$ 上的 Koszul 复形表示。把 Lemma
\ref{equiv-lemma-deform-koszul} 应用于
$\mathcal{O}_{S, s} \to \mathcal{O}_{X \times_S Y, z}$，可知
$K_z \in D(\mathcal{O}_{X \times_S Y, z})$ 由某个正则序列
$f_1, \ldots, f_r \in \mathcal{O}_{X \times_S Y, z}$ 的 Koszul 复形表示；
该序列提升正则序列 $\overline{f}_1, \ldots, \overline{f}_r$，而且
$\mathcal{O}_{X \times_S Y}/(f_1, \ldots, f_r)$ 在 $\mathcal{O}_{S, s}$ 上平坦。
由若干极限论证（Lemma \ref{equiv-lemma-limit-arguments}），存在仿射开集
$U(z) \subset X \times_S Y$，它是 $z$ 的邻域；并存在闭子概形
$Z(z) \subset U(z)$，使得
\begin{enumerate}
\item $Z(z) \to U(z)$ 是正则闭浸入；
\item $K|_{U(z)}$ 拟同构于 $\mathcal{O}_{Z(z)}$；
\item $Z(z) \to S$ 平坦；
\item $Z(z)_s = \Delta_{Y_s/k}(Y_s) \cap U(z)_s$
作为 $U(z)_s$ 的闭子概形成立。
\end{enumerate}

\noindent
由性质 (2)，对 $z, z' \in Y_s \times Y_s$，有
$Z(z) \cap U(z') = Z(z') \cap U(z)$，等式在闭子概形意义下成立。因此得到开邻域
$$
U = \bigcup\nolimits_{z \in Y_s \times Y_s\text{ 闭}} U(z)
$$
，它是 $Y_s \times Y_s$ 在 $X \times_S Y$ 中的开邻域；并得到闭子概形 $Z \subset U$，
使得 (1) $Z \to U$ 是正则闭浸入，(2) $Z \to S$ 平坦，并且
(3) $Z_s = \Delta_{Y_s/k}(Y_s)$。由于 $X \times_S Y \to S$ 固有，把 $S$
替换为 $s$ 的某个开邻域后，可假设 $U = X \times_S Y$。由于投影
$Z_s \to Y_s$ 与 $Z_s \to X_s$ 都是同构，缩小 $S$ 后可假设
$Z \to Y$ 与 $Z \to X$ 都是同构；见 Lemma
\ref{equiv-lemma-isomorphism-in-neighbourhood}。证明完毕。
\end{proof}

\begin{lemma}
\label{equiv-lemma-no-deformations-better}
设 $k$ 为代数闭域，并设 $X$ 为 $k$ 上的光滑固有概形。设
$f : Y \to S$ 为光滑固有态射，其中 $S$ 在 $k$ 上有限型。设 $K$ 为从
$X \times S$ 到 $Y$、定义在 $S$ 上的相对等价的 Fourier--Mukai 核。则 $S$
可由开子概形 $U$ 覆盖，使得存在 $U$-同构
$f^{-1}(U) \cong Y_0 \times U$，其中某个 $Y_0$ 在 $k$ 上固有且光滑。
\end{lemma}

\begin{proof}
选取闭点 $s \in S$。由于 $k$ 代数闭，这是一个 $k$-有理点。令 $Y_0 = Y_s$。
记 $K_0$ 为 $K$ 到 $X \times Y_0$ 的限制。由 Lemma \ref{equiv-lemma-base-change-rek}，它
是从 $X$ 到 $Y_0$、定义在 $\Spec(k)$ 上的相对等价的 Fourier--Mukai 核。
设 $K'_0$ 为 $D_{perf}(\mathcal{O}_{Y_0 \times X})$ 中由 Definition
\ref{equiv-definition-relative-equivalence-kernel} 假定存在的对象。则由 Definition
\ref{equiv-definition-relative-equivalence-kernel} 所固有的对称性，$K'_0$ 是从 $Y_0$ 到 $X$、
定义在 $\Spec(k)$ 上的相对等价的 Fourier--Mukai 核。因此由 Lemma
\ref{equiv-lemma-base-change-rek}，拉回
$$
M = (Y_0 \times X \times S \to Y_0 \times X)^*K'_0
$$
（定义在 $(Y_0 \times S) \times_S (X \times S) = Y_0 \times X \times S$ 上）
是从 $Y_0 \times S$ 到 $X \times S$、定义在 $S$ 上的相对等价的 Fourier--Mukai 核。
现在考虑核
$$
K_{new} =
R\text{pr}_{13, *}(L\text{pr}_{12}^*M
\otimes_{\mathcal{O}_{(Y_0 \times S) \times_S (X \times S)
\times_S Y}}^\mathbf{L}
L\text{pr}_{23}^*K)
$$
，它定义在 $(Y_0 \times S) \times_S Y$ 上。这是从 $Y_0 \times S$ 到 $Y$、
定义在 $S$ 上的相对等价的 Fourier--Mukai 核，因为它是 Section
\ref{equiv-section-category-Fourier-Mukai-kernels} 所构造范畴中两个可逆箭头的复合。
此外，该复合与基变换相容（Lemma \ref{equiv-lemma-base-change-is-functor}）。因此可知，
$K_{new}$ 到 $((Y_0 \times S) \times_S Y)_s = Y_0 \times Y_0$ 的拉回等于
$K_0$ 与 $K'_0$ 的复合，因而等于此范畴中的恒等元。换言之，
$$
L(Y_0 \times Y_0 \to (Y_0 \times S) \times_S Y)^*K_{new}
\cong
\Delta_{Y_0/k, *}\mathcal{O}_{Y_0}
$$
因此由 Lemma \ref{equiv-lemma-no-deformations}，$Y \to S$ 同构于 $Y_0 \times S$，
该同构在 $s$ 的某个开邻域上成立。证明完毕。
\end{proof}






\section{可数性}
\label{equiv-section-countability}

\noindent
本节证明关于某些集合之可数性的若干初等引理。设 $\mathcal{C}$ 为范畴。
本节称 $\mathcal{C}$ 为{\it 可数的}，如果
\begin{enumerate}
\item 对任意 $X, Y \in \Ob(\mathcal{C})$，集合 $\Mor_\mathcal{C}(X, Y)$ 可数；并且
\item $\mathcal{C}$ 的对象同构类之集合可数。
\end{enumerate}

\begin{lemma}
\label{equiv-lemma-countable-finite-type}
设 $R$ 为可数 Noether 环。则 $R$ 上有限型概形的范畴可数。
\end{lemma}

\begin{proof}
略。
\end{proof}

\begin{lemma}
\label{equiv-lemma-countable-abelian}
设 $\mathcal{A}$ 为可数阿贝尔范畴。则 $D^b(\mathcal{A})$ 可数。
\end{lemma}

\begin{proof}
只需对 $D(\mathcal{A})$ 证明该陈述，因为其他范畴都是它的全子范畴。由于
$D(\mathcal{A})$ 中每个对象都是由 $\mathcal{A}$ 的对象组成的复形，立即可知
$D^b(\mathcal{A})$ 的对象同构类之集合可数。此外，对有界复形
$A^\bullet$ 与 $B^\bullet$（它们由 $\mathcal{A}$ 的对象组成），显然
$\Hom_{K^b(\mathcal{A})}(A^\bullet, B^\bullet)$ 可数。
有
$$
\Hom_{D^b(\mathcal{A})}(A^\bullet, B^\bullet) =
\colim_{s : (A')^\bullet \to A^\bullet
\text{ qis and }(A')^\bullet\text{ bounded}}
\Hom_{K^b(\mathcal{A})}((A')^\bullet, B^\bullet)
$$
，这是 Derived Categories, Lemma \ref{derived-lemma-bounded-derived}。因此它是可数集，
因为它是可数个可数集的余极限。
\end{proof}

\begin{lemma}
\label{equiv-lemma-countable-perfect}
设 $X$ 为某个可数 Noether 环上的有限型概形。则范畴
$D_{perf}(\mathcal{O}_X)$ 与 $D^b_{\textit{Coh}}(\mathcal{O}_X)$ 都可数。
\end{lemma}

\begin{proof}
由 Morphisms, Lemma \ref{morphisms-lemma-finite-type-noetherian}，$X$ 是 Noether 概形。
因而由 Derived Categories of Schemes, Lemma
\ref{perfect-lemma-perfect-on-noetherian}，$D_{perf}(\mathcal{O}_X)$ 是
$D^b_{\textit{Coh}}(\mathcal{O}_X)$ 的全子范畴。所以只需对
$D^b_{\textit{Coh}}(\mathcal{O}_X)$ 证明结论。回忆
$D^b_{\textit{Coh}}(\mathcal{O}_X) = D^b(\textit{Coh}(\mathcal{O}_X))$
，这是 Derived Categories of Schemes, Proposition \ref{perfect-proposition-DCoh}。
故由 Lemma \ref{equiv-lemma-countable-abelian}，只需证明
$\textit{Coh}(\mathcal{O}_X)$ 可数。略去此证明。
\end{proof}

\begin{lemma}
\label{equiv-lemma-countable-isos}
设 $K$ 为代数闭域，$S$ 为 $K$ 上的有限型概形，并设 $X \to S$ 与 $Y \to S$
为有限型态射。存在可数集 $I$，并且对 $i \in I$ 存在一对
$(S_i \to S, h_i)$，满足下列性质：
\begin{enumerate}
\item $S_i \to S$ 是有限型态射；令 $X_i = X \times_S S_i$ 与
$Y_i = Y \times_S S_i$；
\item $h_i : X_i \to Y_i$ 是 $S_i$ 上的同构；
\item 对任意闭点 $s \in S(K)$，若 $X_s \cong Y_s$ 在 $K = \kappa(s)$ 上成立，
则 $s$ 属于 $S_i \to S$ 的像，其中 $i$ 为某个指标。
\end{enumerate}
\end{lemma}

\begin{proof}
域 $K$ 是其可数子域的滤过并。对偶地，$\Spec(K)$ 是 $K$ 的这些可数子域之谱的
余滤过极限。因此 Limits, Lemma \ref{limits-lemma-descend-finite-presentation}
保证存在可数子域 $k$，以及态射 $X_0 \to S_0$ 与 $Y_0 \to S_0$；这些态射位于
$k$ 上的有限型概形之间，并使得 $X \to S$ 与 $Y \to S$ 是它们的基变换。

\medskip\noindent
由 Lemma \ref{equiv-lemma-countable-finite-type}，存在可数集 $I$ 与若干对
$(S_{0, i} \to S_0, h_{0, i})$，满足
\begin{enumerate}
\item $S_{0, i} \to S_0$ 是有限型态射；令
$X_{0, i} = X_0 \times_{S_0} S_{0, i}$ 与
$Y_{0, i} = Y_0 \times_{S_0} S_{0, i}$；
\item $h_{0, i} : X_{0, i} \to Y_{0, i}$ 是 $S_{0, i}$ 上的同构。
\end{enumerate}
而且每一对 $(T \to S_0, h_T)$，其中 $T \to S_0$ 有限型且
$h_T : X_0 \times_{S_0} T \to Y_0 \times_{S_0} T$ 为同构，都同构于上述某一对。
记 $(S_i \to S, h_i)$ 为 $(S_{0, i} \to S_0, h_{0, i})$ 沿
$\Spec(K) \to \Spec(k)$ 的基变换。我们断言这组数据满足要求。

\medskip\noindent
设 $s \in S(K)$，并设 $h_s : X_s \to Y_s$ 为 $K = \kappa(s)$ 上的同构。
可把 $K$ 写成其有限生成 $k$-子代数的滤过并。因此由 Limits, Proposition
\ref{limits-proposition-characterize-locally-finite-presentation} 与 Lemma
\ref{limits-lemma-descend-finite-presentation}，可找到这样的有限生成 $k$-子代数
$K \supset A \supset k$，使得
\begin{enumerate}
\item 存在交换图表
$$
\xymatrix{
\Spec(K) \ar[d]_s \ar[r] &
\Spec(A) \ar[d]^{s'} \\
S \ar[r] &
S_0}
$$
其中 $s' : \Spec(A) \to S_0$ 是 $k$ 上的某个态射；
\item $h_s$ 是下列同构的基变换：
$h_{s'} : X_0 \times_{S_0, s'} \Spec(A) \to
X_0 \times_{S_0, s'} \Spec(A)$，该同构定义在 $A$ 上。
\end{enumerate}
于是 $(s' : \Spec(A) \to S_0, h_{s'})$ 当然同构于
$(S_{0, i} \to S_0, h_{0, i})$，其中 $i \in I$。这便完成证明，因为 (1) 中的交换图表说明，
$s$ 属于 $s'$ 到 $\Spec(K)$ 的基变换之像。
\end{proof}

\begin{lemma}
\label{equiv-lemma-countable-equivs}
设 $K$ 为代数闭域。存在可数集 $I$，并且对 $i \in I$ 存在一组
$(S_i/K, X_i \to S_i, Y_i \to S_i, M_i)$，满足下列性质：
\begin{enumerate}
\item $S_i$ 是 $K$ 上的有限型概形；
\item $X_i \to S_i$ 与 $Y_i \to S_i$ 都是光滑固有概形态射；
\item $M_i \in D_{perf}(\mathcal{O}_{X_i \times_{S_i} Y_i})$
是从 $X_i$ 到 $Y_i$、定义在 $S_i$ 上的相对等价的 Fourier--Mukai 核；
\item 对任意光滑固有概形 $X$ 与 $Y$（它们定义在 $K$ 上），若存在 $K$-线性正合等价
$D_{perf}(\mathcal{O}_X) \to D_{perf}(\mathcal{O}_Y)$
，则存在 $i \in I$ 与 $s \in S_i(K)$，使得
$X \cong (X_i)_s$ 且 $Y \cong (Y_i)_s$。
\end{enumerate}
\end{lemma}

\begin{proof}
选取可数子域 $k \subset K$，例如素域。由 Lemmas
\ref{equiv-lemma-countable-finite-type} 与 \ref{equiv-lemma-countable-perfect}，满足本引理 (1)、(2)、(3)
的 $k$ 上系统只有可数多个同构类。因此可选取可数集 $I$，并对每个 $i \in I$
选取一个这样的系统
$$
(S_{0, i}/k, X_{0, i} \to S_{0, i}, Y_{0, i} \to S_{0, i}, M_{0, i})
$$
，它定义在 $k$ 上，并使每个同构类至少出现一次。记
$(S_i/K, X_i \to S_i, Y_i \to S_i, M_i)$ 为上述系统到 $K$ 的基变换。
由 Lemma \ref{equiv-lemma-base-change-rek}，该系统具有性质 (1)、(2)、(3)。下面证明性质 (4)。

\medskip\noindent
考虑光滑固有概形 $X$ 与 $Y$（它们定义在 $K$ 上），并设有 $K$-线性正合等价
$F : D_{perf}(\mathcal{O}_X) \to D_{perf}(\mathcal{O}_Y)$。
由 Proposition \ref{equiv-proposition-equivalence}，可假设存在对象
$M \in D_{perf}(\mathcal{O}_{X \times Y})$，使得 $F = \Phi_M$ 是相应的
Fourier--Mukai 函子。由 Lemma \ref{equiv-lemma-fourier-mukai-flat-proper-over-noetherian}，
存在 $M'$，它属于 $D_{perf}(\mathcal{O}_{Y \times X})$，使得 $\Phi_{M'}$ 是 $\Phi_M$
的右伴随。由于 $\Phi_M$ 是等价，这说明 $\Phi_{M'}$ 是 $\Phi_M$ 的拟逆。
由 Lemma \ref{equiv-lemma-fourier-mukai-flat-proper-over-noetherian}，下列对象所定义的
Fourier--Mukai 函子
$$
A = R\text{pr}_{13, *}(
L\text{pr}_{12}^*M
\otimes_{\mathcal{O}_{X \times Y \times X}}^\mathbf{L}
L\text{pr}_{23}^*M')
$$
（该对象属于 $D_{perf}(\mathcal{O}_{X \times X})$）以及
$$
B = R\text{pr}_{13, *}(
L\text{pr}_{12}^*M'
\otimes_{\mathcal{O}_{Y \times X \times Y}}^\mathbf{L}
L\text{pr}_{23}^*M)
$$
（该对象属于 $D_{perf}(\mathcal{O}_{Y \times Y})$）分别同构于
$\text{id} : D_{perf}(\mathcal{O}_X) \to D_{perf}(\mathcal{O}_X)$
与
$\text{id} : D_{perf}(\mathcal{O}_Y) \to D_{perf}(\mathcal{O}_Y)$
。因此
$A \cong \Delta_{X/K, *}\mathcal{O}_X$ 且
$B \cong \Delta_{Y/K, *}\mathcal{O}_Y$，这是 Lemma \ref{equiv-lemma-uniqueness}。
故由定义，$M$ 是从 $X$ 到 $Y$、定义在 $K$ 上的相对等价的 Fourier--Mukai 核。

\medskip\noindent
可把 $K$ 写成其有限型 $k$-子代数 $A \subset K$ 的滤过余极限。由 Limits, Lemma
\ref{limits-lemma-descend-finite-presentation}，可找到 $X_0, Y_0$，它们在 $A$ 上有限型，
它们到 $K$ 的基变换分别产生 $X$ 与 $Y$。由 Limits, Lemmas
\ref{limits-lemma-eventually-proper} 与 \ref{limits-lemma-descend-smooth}，增大 $A$ 后，
可假设 $X_0$ 与 $Y_0$ 在 $A$ 上光滑且固有。由 Lemma \ref{equiv-lemma-descend-rek}，
再次增大 $A$ 后，可假设 $M$ 是某个
$M_0  \in D_{perf}(\mathcal{O}_{X_0 \times_{\Spec(A)} Y_0})$ 的拉回；该对象是从
$X_0$ 到 $Y_0$、定义在 $\Spec(A)$ 上的相对等价的 Fourier--Mukai 核。因此
$(S_0/k, X_0 \to S_0, Y_0 \to S_0, M_0)$ 同构于
$(S_{0, i}/k, X_{0, i} \to S_{0, i}, Y_{0, i} \to S_{0, i}, M_{0, i})$
，其中 $i \in I$。由于 $S_i = S_{0, i} \times_{\Spec(k)} \Spec(K)$，
性质 (4) 对 $s : \Spec(K) \to S_i$ 成立；该态射由
$\Spec(K) \to \Spec(A) \cong S_{0, i}$ 诱导，而后者来自 $A \subset K$。
\end{proof}








\section{导出等价簇的可数性}
\label{equiv-section-countable-derived-equivalent}

\noindent
本节证明 Anel 和 To\"en 的一个结果，见 \cite{AT}。

\begin{definition}
\label{equiv-definition-derived-equivalent}
设 $k$ 为域，$X$ 和 $Y$ 为 $k$ 上的光滑射影概形。称 $X$ 和 $Y$
是{\it 导出等价的}，是指存在一个 $k$-线性的正合等价
$D_{perf}(\mathcal{O}_X) \to D_{perf}(\mathcal{O}_Y)$.
\end{definition}

\noindent
结果如下。

\begin{theorem}
\label{equiv-theorem-countable}
\begin{reference}
对 \cite{AT} 的轻微改进
\end{reference}
设 $K$ 为代数闭域。设 $\mathbf{X}$ 为 $K$ 上的光滑固有概形。光滑固有概形
$\mathbf{Y}$（定义在 $K$ 上且与 $\mathbf{X}$ 导出等价）的同构类至多可数。
\end{theorem}

\begin{proof}
取一个可数集 $I$，并对每个 $i \in I$ 取满足引理
\ref{equiv-lemma-countable-equivs} 的性质 (1)、(2)、(3) 和 (4) 的系统
$(S_i/K, X_i \to S_i, Y_i \to S_i, M_i)$。固定 $i \in I$，并置
$S = S_i$、$X = X_i$、$Y = Y_i$ 以及 $M = M_i$。显然，只需证明：
纤维 $Y_s$ 的同构类所成之集是可数的，其中 $s \in S(K)$ 且
$X_s \cong \mathbf{X}$。我们在下一段证明这一点。

\medskip\noindent
设 $S$ 为 $K$ 上的有限型概形，$X \to S$ 和 $Y \to S$ 为固有光滑态射，
并设 $M \in D_{perf}(\mathcal{O}_{X \times_S Y})$ 为一个从 $X$ 到 $Y$、
定义在 $S$ 上的相对等价的 Fourier--Mukai 核。我们将证明：纤维 $Y_s$ 的同构类
所成之集是可数的，其中 $s \in S(K)$ 且 $X_s \cong \mathbf{X}$。对族
$\mathbf{X} \times S \to S$ 和 $X \to S$ 应用引理
\ref{equiv-lemma-countable-isos}，可得一个可数集 $I$，并对每个 $i \in I$
得到一对 $(S_i \to S, h_i)$，满足下列性质：
\begin{enumerate}
\item $S_i \to S$ 是有限型态射，并置 $X_i = X \times_S S_i$；
\item $h_i : \mathbf{X} \times S_i \to X_i$ 是 $S_i$ 上的同构；
\item 对任意闭点 $s \in S(K)$，如果同构 $\mathbf{X} \cong X_s$ 是定义在
$K = \kappa(s)$ 上的，则点 $s$ 位于态射 $S_i \to S$（对某个 $i$）的像中。
\end{enumerate}
置 $Y_i = Y \times_S S_i$，并用
$M_i \in D_{perf}(\mathcal{O}_{X_i \times_{S_i} Y_i})$
表示 $M$ 的拉回。由引理 \ref{equiv-lemma-base-change-rek}，$M_i$ 是一个从 $X_i$ 到
$Y_i$、定义在 $S_i$ 上的相对等价的 Fourier--Mukai 核。由于 $I$ 可数，
由性质 (3)，只需证明纤维 $Y_{i, s}$ 的同构类所成之集是可数的，其中
$s \in S_i(K)$。事实上，由引理 \ref{equiv-lemma-no-deformations-better}，这个数是有限的，
证明完毕。
\end{proof}
% Generated mechanically from the frozen pione.tex target.
% Do not edit this generated file; edit the bound locale source instead.

% OK, start here.
%

\setcounter{chapter}{57}
\chapter{概形的基本群}



\phantomsection
\label{pione-section-phantom}


\section{引言}
\label{pione-section-introduction}

\noindent
本章讨论概形的 Grothendieck 基本群及其应用。基础性参考文献是 \cite{SGA1}。
一篇很好的导论是 \cite{Lenstra}。
其他参考文献包括 \cite{Murre-lectures} 和 \cite{Grothendieck-Murre}。










\section{点上的 \'etale 概形}
\label{pione-section-schemes-etale-point}

\noindent
本节描述域的谱上的 \'etale 概形。在陈述结果之前，我们先对拓扑群 $G$
引入 $G$-集范畴。

\begin{definition}
\label{pione-definition-G-set-continuous}
设 $G$ 为拓扑群。一个{\it $G$-集}，有时也称为{\it 离散 $G$-集}，
是一个集合 $X$ 连同一个左作用 $a : G \times X \to X$；要求 $a$ 连续，
其中给 $X$ 赋予离散拓扑，给 $G \times X$ 赋予乘积拓扑。
一个{\it $G$-集的态射} $f : X \to Y$ 就是任意的 $G$-等变映射，
其源为 $X$、靶为 $Y$。$G$-集范畴记作 {\it $G\textit{-Sets}$}。
\end{definition}

\noindent
条件 $a : G \times X \to X$ 连续，无非表示任意 $x \in X$ 的稳定子群在 $G$
中是开的。如果 $G$ 是抽象群 $G$（即它是群但不是拓扑群），只要给它赋予
那么上述定义与先前的定义一致（例如见 Sites, Example
\ref{sites-example-site-on-group}），前提是给 $G$ 赋予离散拓扑。

\medskip\noindent
回忆，如果 $L/K$ 是无限 Galois 扩张，则 Galois 群
$G = \text{Gal}(L/K)$ 带有典范拓扑；见 Fields, Section
\ref{fields-section-infinite-galois}。

\begin{lemma}
\label{pione-lemma-sheaves-point}
设 $K$ 为域。设 $K^{sep}$ 为 $K$ 的一个可分闭包。考虑 profinite 群
$G = \text{Gal}(K^{sep}/K)$。函子
$$
\begin{matrix}
\text{schemes \'etale over }K &
\longrightarrow &
G\textit{-Sets} \\
X/K & \longmapsto &
\Mor_{\Spec(K)}(\Spec(K^{sep}), X)
\end{matrix}
$$
是范畴等价。
\end{lemma}

\begin{proof}
概形 $X$ 定义在 $K$ 上；它在 $K$ 上为 \'etale，当且仅当
$X \cong \coprod_{i\in I} \Spec(K_i)$，其中每个 $K_i$ 都是 $K$ 的有限可分扩张
（Morphisms, Lemma \ref{morphisms-lemma-etale-over-field}）。该引理的函子将
$X$ 送到 $G$-集
$$
\coprod\nolimits_i \Hom_K(K_i, K^{sep})
$$
及其自然左 $G$-作用。由 $G$ 上拓扑的定义，每个元素都有开稳定子群。反之，
任意 $G$-集 $S$ 都是其轨道的不交并。设 $S = \coprod S_i$。选取
$s_i \in S_i$，并用 $G_i \subset G$ 表示它的开稳定子群。由 Galois 理论
（Fields, Theorem \ref{fields-theorem-inifinite-galois-theory}），域
$(K^{sep})^{G_i}$ 是 $K$ 的有限可分扩张，因而概形
$$
\coprod\nolimits_i \Spec((K^{sep})^{G_i})
$$
在 $K$ 上为 \'etale。这给出了该引理函子的逆。略去一些细节。
\end{proof}

\begin{remark}
\label{pione-remark-covering-surjective}
在引理 \ref{pione-lemma-sheaves-point} 的对应下，小 \'etale 位点
$\Spec(K)_\etale$，即 $K$ 的小 \'etale 位点，其中的覆盖对应于
$G\textit{-Sets}$ 中的满射映射族。
\end{remark}









\section{Galois 范畴}
\label{pione-section-galois}

\noindent
本节讨论读者可在 \cite[Expos\'e V, Sections 4, 5, and 6]{SGA1}
中找到的部分材料。

\medskip\noindent
设 $F : \mathcal{C} \to \textit{Sets}$ 为函子。回忆，依照我们的约定，
范畴的对象构成集合，并且任意一对对象之间的态射也构成集合。有典范单射
\begin{equation}
\label{pione-equation-embedding-product}
\text{Aut}(F)
\longrightarrow
\prod\nolimits_{X \in \Ob(\mathcal{C})} \text{Aut}(F(X))
\end{equation}
对集合 $E$，给 $\text{Aut}(E)$ 赋予紧开拓扑；见 Topology, Example
\ref{topology-example-automorphisms-of-a-set}。当然，当 $E$ 有限时，这就是
离散拓扑；本节所关心的正是这种情形\footnote{讨论 pro-\'etale 基本群时，
一般情形也将是我们所关心的。}。给 $\text{Aut}(F)$ 赋予由
(\ref{pione-equation-embedding-product}) 右端的乘积拓扑诱导的拓扑。特别地，作用映射
$$
\text{Aut}(F) \times F(X) \longrightarrow F(X)
$$
在给 $F(X)$ 赋予离散拓扑时是连续的，因为作用映射
$\text{Aut}(E) \times E \to E$ 对任意集合 $E$ 都具有这一性质。
$\text{Aut}(F)$ 上该拓扑的泛性质如下：设 $G$ 为拓扑群，而
$G \to \text{Aut}(F)$ 为群同态，并且诱导作用 $G \times F(X) \to F(X)$
对所有 $X \in \Ob(\mathcal{C})$ 都连续，其中给 $F(X)$ 赋予离散拓扑。
那么 $G \to \text{Aut}(F)$ 连续。

\medskip\noindent
下述引理说明，取值于有限集范畴的函子之自同构群自动成为 profinite 群。

\begin{lemma}
\label{pione-lemma-aut-inverse-limit}
设 $\mathcal{C}$ 为范畴，$F : \mathcal{C} \to \textit{Sets}$ 为函子。映射
(\ref{pione-equation-embedding-product}) 将 $\text{Aut}(F)$ 识别为
$\prod_{X \in \Ob(\mathcal{C})} \text{Aut}(F(X))$ 的闭子群。特别地，如果
$F(X)$ 对所有 $X$ 都有限，那么 $\text{Aut}(F)$ 是 profinite 群。
\end{lemma}

\begin{proof}
设 $\xi = (\gamma_X) \in \prod \text{Aut}(F(X))$ 是不属于
$\text{Aut}(F)$ 的元素。于是存在态射 $f : X \to X'$（属于 $\mathcal{C}$）
及元素 $x \in F(X)$，使得
$F(f)(\gamma_X(x)) \not = \gamma_{X'}(F(f)(x))$.
考虑开邻域
$U = \{\gamma \in \text{Aut}(F(X)) \mid \gamma(x) = \gamma_X(x)\}$，
它是 $\gamma_X$ 的邻域；以及开邻域
$U' = \{\gamma' \in \text{Aut}(F(X')) \mid \gamma'(F(f)(x)) =
\gamma_{X'}(F(f)(x))\}$.
于是
$U \times U' \times \prod_{X'' \not = X, X'} \text{Aut}(F(X''))$
是 $\xi$ 的开邻域，且不与 $\text{Aut}(F)$ 相交。最后一个断言来自如下事实：
$\prod \text{Aut}(F(X))$ 在每个 $F(X)$ 都有限时是 profinite 空间。
\end{proof}

\begin{example}
\label{pione-example-galois-category-G-sets}
设 $G$ 为拓扑群。一个重要例子是遗忘函子
\begin{equation}
\label{pione-equation-forgetful}
\textit{Finite-}G\textit{-Sets} \longrightarrow \textit{Sets}
\end{equation}
其中 $\textit{Finite-}G\textit{-Sets}$ 是 $G\textit{-Sets}$ 的全子范畴，
其对象是有限 $G$-集。$G\textit{-Sets}$（即 $G$-集范畴）定义于 Definition
\ref{pione-definition-G-set-continuous}。
\end{example}

\noindent
设 $G$ 为拓扑群。$G$ 的{\it profinite 完备化}是 profinite 群
$$
G^\wedge =
\lim_{U \subset G\text{ open, normal, finite index}} G/U
$$
及其 profinite 拓扑。注意，该极限是余滤过的，因为有限多个开、正规且
具有有限指数的子群之交仍具有这些性质。profinite 完备化的泛性质是：任意连续映射
$G \to H$，其中 $H$ 为 profinite 群，都典范地分解为
$G \to G^\wedge \to H$。

\begin{lemma}
\label{pione-lemma-single-out-profinite}
设 $G$ 为拓扑群。给函子 (\ref{pione-equation-forgetful}) 的自同构群赋予引理
\ref{pione-lemma-aut-inverse-limit} 所给的 profinite 拓扑，则该群是 $G$ 的
profinite 完备化。
\end{lemma}

\begin{proof}
用 $F_G$ 表示函子 (\ref{pione-equation-forgetful})。任意态射 $X \to Y$（属于
$\textit{Finite-}G\textit{-Sets}$）都与 $G$ 的作用可交换。
因此任意 $g \in G$ 都定义 $F_G$ 的一个自同构，从而得到典范群同态
$G \to \text{Aut}(F_G)$。注意，任意有限 $G$-集 $X$ 都是若干 $G$-集的
有限不交并，这些集合形如 $G/H_i$ 并带有典范 $G$-作用，其中 $H_i \subset G$
是有限指数开子群。于是 $U_i = \bigcap gH_ig^{-1}$ 是开、正规且具有有限指数的。
此外，$U_i$ 在 $G/H_i$ 上的作用平凡，故 $U = \bigcap U_i$ 在 $F_G(X)$ 上的
作用平凡。因此作用 $G \times F_G(X) \to F_G(X)$ 连续。由
$\text{Aut}(F_G)$ 上拓扑的泛性质，映射 $G \to \text{Aut}(F_G)$ 连续。
由引理 \ref{pione-lemma-aut-inverse-limit} 和 profinite 完备化的泛性质，诱导出连续群同态
$$
G^\wedge \longrightarrow \text{Aut}(F_G)
$$
此外，由于 $G/U$ 在 $G/U$ 上的作用忠实，该映射为单射。如果其像稠密，
则该映射为满射，因而由 Topology, Lemma \ref{topology-lemma-bijective-map}
可知它是同胚。

\medskip\noindent
设 $\gamma \in \text{Aut}(F_G)$ 且 $X \in \Ob(\mathcal{C})$。我们将证明存在
$g \in G$，使 $\gamma$ 与 $g$ 在 $F_G(X)$ 上诱导相同的作用。这将完成证明。
如前所述，$X$ 是若干 $G/H_i$ 的有限不交并。取上述 $U_i$ 和 $U$，则有限
$G$-集 $Y = G/U$ 满射到 $G/H_i$（对所有 $i$），故只需找到 $g \in G$，使 $\gamma$
与 $g$ 在 $F_G(G/U) = G/U$ 上诱导相同作用。设 $e \in G$ 为单位元，并设
$\gamma(eU) = g_0U$，其中 $g_0 \in G$。对任意 $g_1 \in G$，态射
$$
R_{g_1} : G/U \longrightarrow G/U,\quad gU \longmapsto gg_1U
$$
属于 $\textit{Finite-}G\textit{-Sets}$，并且与 $\gamma$ 的作用可交换。因此
$$
\gamma(g_1U) = \gamma(R_{g_1}(eU)) = R_{g_1}(\gamma(eU)) =
R_{g_1}(g_0U) = g_0g_1U
$$
所以 $g = g_0$ 符合要求。
\end{proof}

\noindent
回忆，正合函子是与所有有限极限和有限余极限可交换的函子。特别地，这种函子
与等化子、余等化子、纤维积、推出等构造可交换。

\begin{lemma}
\label{pione-lemma-second-fundamental-functor}
设 $G$ 为拓扑群。设
$F : \textit{Finite-}G\textit{-Sets} \to \textit{Sets}$
为正合函子，并且 $F(X)$ 对所有 $X$ 都有限。那么 $F$ 同构于函子
(\ref{pione-equation-forgetful})。
\end{lemma}

\begin{proof}
设 $X$ 为 $\textit{Finite-}G\textit{-Sets}$ 的非空对象。图
$$
\xymatrix{
X \ar[r] \ar[d] & \{*\} \ar[d] \\
\{*\} \ar[r] & \{*\}
}
$$
为余笛卡儿图。因此 $F(X)$ 非空。设 $U \subset G$ 为有限指数开正规子群。
注意
$$
G/U \times G/U = \coprod\nolimits_{gU \in G/U} G/U
$$
其中与 $gU$ 对应的直和项，对应于左端 $(eU, gU)$ 的轨道。于是
$$
F(G/U) \times F(G/U) = F(G/U \times G/U) = \coprod\nolimits_{gU \in G/U} F(G/U)
$$
故 $|F(G/U)| = |G/U|$，因为 $F(G/U)$ 非空。因此
$$
\lim_{U \subset G\text{ open, normal, finite idex}} F(G/U)
$$
非空（Categories, Lemma \ref{categories-lemma-nonempty-limit}）。选取该极限中的
元素 $\gamma = (\gamma_U)$。用 $F_G$ 表示函子 (\ref{pione-equation-forgetful})。
可将 $F_G$ 识别为函子
$$
X \longmapsto \colim_U \Mor(G/U, X)
$$
其中 $f : G/U \to X$ 对应于 $f(eU) \in X = F_G(X)$（略去细节）。因此元素
$\gamma$ 决定一个良定义的映射
$$
t : F_G \longrightarrow F
$$
具体而言，给定 $x \in X$，选取 $U$ 和 $f : G/U \to X$，使它把 $eU$
送到 $x$，再置 $t_X(x) = F(f)(\gamma_U)$。我们将证明，$t$ 诱导双射
$t_{G/U} : F_G(G/U) \to F(G/U)$（对任意 $U$）。由此容易推出 $t$ 是同构
（略去细节）。由于 $|F_G(G/U)| = |F(G/U)|$，只需证明 $t_{G/U}$ 为满射。
其像至少含有元素
$t_{G/U}(eU) = F(\text{id}_{G/U})(\gamma_U) = \gamma_U$.
对 $g \in G$，用 $R_g : G/U \to G/U$ 表示右乘。映射
$F(R_g) : F(G/U) \to F(G/U)$ 的不动点集等于 $F(\emptyset) = \emptyset$，
如果 $g \not \in U$；这是因为 $F$ 与等化子可交换。于是，如果
$g_1, \ldots, g_{|G/U|}$ 是 $G/U$
的一组代表元，则各元素 $F(R_{g_i})(\gamma_U)$ 两两不同，因而充满
$F(G/U)$。此时
$$
t_{G/U}(g_iU) = F(R_{g_i})(\gamma_U)
$$
证明完毕。
\end{proof}

\begin{example}
\label{pione-example-from-C-F-to-G-sets}
设 $\mathcal{C}$ 为范畴，而 $F : \mathcal{C} \to \textit{Sets}$ 为函子，
使得 $F(X)$ 对所有 $X \in \Ob(\mathcal{C})$ 都有限。由引理
\ref{pione-lemma-aut-inverse-limit}，$G = \text{Aut}(F)$ 典范地带有 profinite
拓扑群结构。得到函子
\begin{equation}
\label{pione-equation-remember}
\mathcal{C} \longrightarrow \textit{Finite-}G\textit{-Sets},\quad
X \longmapsto F(X)
\end{equation}
其中给 $F(X)$ 赋予所诱导的 $G$-作用。由我们对 $\text{Aut}(F)$ 上拓扑的构造，
该作用连续。
\end{example}

\noindent
定义 Galois 范畴的目的是刻画那些使函子 (\ref{pione-equation-remember}) 成为等价的
对 $(\mathcal{C}, F)$。我们如下定义 Galois 范畴。

\begin{definition}
\label{pione-definition-galois-category}
\begin{reference}
不同于 \cite[Expos\'e V, Definition 5.1]{SGA1} 中的定义。
可与 \cite[Definition 7.2.1]{BS} 比较。
\end{reference}
设 $\mathcal{C}$ 为范畴，而 $F : \mathcal{C} \to \textit{Sets}$ 为函子。如果
满足下列条件，则称对 $(\mathcal{C}, F)$ 为{\it Galois 范畴}：
\begin{enumerate}
\item $\mathcal{C}$ 有有限极限和有限余极限；
\item
\label{pione-item-connected-components}
$\mathcal{C}$ 的每个对象都是连通对象的有限（可以为空）余积；
\item $F(X)$ 对所有 $X \in \Ob(\mathcal{C})$ 都有限；
\item $F$ 反映同构\footnote{即，给定态射 $f$（属于 $\mathcal{C}$），如果 $F(f)$
是同构，则 $f$ 是同构。}，并且是正合的\footnote{这表示 $F$ 与有限极限和
有限余极限可交换；见 Categories, Section
\ref{categories-section-exact-functor}。}。
\end{enumerate}
这里称 $X \in \Ob(\mathcal{C})$ 连通，是指它不是始对象，并且对任意单态射
$Y \to X$，要么 $Y$ 是始对象，要么 $Y \to X$ 是同构。
\end{definition}

\noindent
{\bf 警告：}这个定义与大多数参考文献所给的定义不同（尽管最终会看出二者等价）。
具体而言，\cite[Expos\'e V, Definition 5.1]{SGA1} 把 Galois 范畴定义为
与 $\textit{Finite-}G\textit{-Sets}$ 等价的范畴，其中 $G$ 为某个 profinite 群。
随后 Grothendieck 用一组公理 (G1) -- (G6) 刻画 Galois 范畴；这些公理比上述
公理弱。我们作此选择，是为了强调有限极限和有限余极限的存在性以及函子 $F$
的正合性。其代价是稍后应用本节结果时必须多做一些工作。

\begin{lemma}
\label{pione-lemma-epi-mono}
设 $(\mathcal{C}, F)$ 为 Galois 范畴。设
$X \to Y \in \text{Arrows}(\mathcal{C})$。则
\begin{enumerate}
\item $F$ 忠实；
\item $X \to Y$ 是单态射
$\Leftrightarrow F(X) \to F(Y)$ 是单射；
\item $X \to Y$ 是满态射
$\Leftrightarrow F(X) \to F(Y)$ 是满射；
\item 对象 $A$（属于 $\mathcal{C}$）是始对象，当且仅当 $F(A) = \emptyset$；
\item 对象 $Z$（属于 $\mathcal{C}$）是终对象，当且仅当 $F(Z)$ 为单元集；
\item 如果 $X$ 和 $Y$ 连通，则 $X \to Y$ 是满态射；
\item
\label{pione-item-one-element}
如果 $X$ 连通，而 $a, b : X \to Y$ 为两个态射，则等式 $a = b$ 在如下条件下成立：
$F(a)$ 和 $F(b)$ 在 $F(X)$ 的某个元素上一致；
\item 如果 $X = \coprod_{i = 1, \ldots, n} X_i$ 且
$Y = \coprod_{j = 1, \ldots, m} Y_j$，其中 $X_i$、$Y_j$ 连通，则存在映射
$\alpha : \{1, \ldots, n\} \to \{1, \ldots, m\}$，使 $X \to Y$ 来自一族态射
$X_i \to Y_{\alpha(i)}$。
\end{enumerate}
\end{lemma}

\begin{proof}
(1) 的证明。设 $a, b : X \to Y$ 且 $F(a) = F(b)$。令 $E$ 为 $a$ 和 $b$
的等化子。于是 $F(E) = F(X)$；可知 $E = X$，因为 $F$ 反映同构。

\medskip\noindent
(2) 的证明。函子 $F$ 把态射 $X \to X \times_Y X$ 变为映射
$F(X) \to F(X) \times_{F(Y)} F(X)$，并且 $F$ 反映同构，故断言成立。

\medskip\noindent
(3) 的证明。函子 $F$ 把态射 $Y \amalg_X Y \to Y$ 变为映射
$F(Y) \amalg_{F(X)} F(Y) \to F(Y)$，并且 $F$ 反映同构，故断言成立。

\medskip\noindent
(4) 的证明。存在始对象 $A$，而且显然 $F(A) = \emptyset$。反之，如果对象
$X$ 满足 $F(X) = \emptyset$，则唯一映射 $A \to X$ 诱导双射
$F(A) \to F(X)$，故 $A \to X$ 是同构。

\medskip\noindent
(5) 的证明。存在终对象 $Z$，而且显然 $F(Z)$ 为单元集。反之，如果对象 $X$
的 $F(X)$ 为单元集，则唯一映射 $X \to Z$ 诱导双射 $F(X) \to F(Z)$，故
$X \to Z$ 是同构。

\medskip\noindent
(6) 的证明。等化子 $E$（对应两个映射 $Y \to Y \amalg_X Y$）不是
$\mathcal{C}$ 的始对象，因为 $X \to Y$ 经过 $E$ 分解，且
$F(X) \not = \emptyset$。因此 $E = Y$，断言得证。

\medskip\noindent
(\ref{pione-item-one-element}) 的证明。等化子 $E$ 是 $a$ 和 $b$ 的等化子，它带有单态射
$E \to X$，而 $F(E) \subset F(X)$ 是 $F(a)$ 和 $F(b)$ 取值相同的元素所成之集。
最后使用如下事实：要么 $E$ 是始对象，要么 $E = X$。

\medskip\noindent
(8) 的证明。对每个 $i, j$，$E_{ij} = X_i \times_Y Y_j$ 要么是始对象，
要么等于 $X_i$。选取 $s \in F(X_i)$，可知 $E_{ij} = X_i$ 当且仅当 $s$
映到 $F(Y_j) \subset F(Y)$ 的某个元素，故这种情形对唯一的
$j = \alpha(i)$ 成立。
\end{proof}

\noindent
由上述引理可知，给定 Galois 范畴中的连通对象 $X$；若该范畴记为
$(\mathcal{C}, F)$，则自同构群 $\text{Aut}(X)$ 的阶至多为 $|F(X)|$。事实上，给定
$s \in F(X)$ 和 $g \in \text{Aut}(X)$，由 (\ref{pione-item-one-element})，
$F(g)(s) = s$ 当且仅当 $g = \text{id}_X$。如果等号成立，就称 $X$ 是
{\it Galois 的}。等价地，$X$ 是 Galois 的，当且仅当它连通，并且
$\text{Aut}(X)$ 在 $F(X)$ 上传递作用。

\begin{lemma}
\label{pione-lemma-galois}
设 $(\mathcal{C}, F)$ 为 Galois 范畴。对任意连通对象 $X$（属于 $\mathcal{C}$），
存在 Galois 对象 $Y$ 和态射 $Y \to X$。
\end{lemma}

\begin{proof}
以下将不再明言地使用引理 \ref{pione-lemma-epi-mono} 的结果。令
$n = |F(X)|$。考虑 $X^n$ 及 $S_n$ 在其上的自然作用。设
$$
X^n = \coprod\nolimits_{t \in T} Z_t
$$
是连通对象分解。选取 $t$，使 $F(Z_t)$ 含有 $(s_1, \ldots, s_n)$，其中各
$s_i$ 两两不同。若 $(s'_1, \ldots, s'_n) \in F(Z_t)$ 是另一元素，则我们断言
各 $s'_i$ 也两两不同。否则，比如 $s'_i = s'_j$，则 $Z_t$ 是 $X^{n - 1}$
的某个连通分支在对角态射
$$
\Delta_{ij} : X^{n - 1} \longrightarrow X^n
$$
下的像。由于连通对象之间的态射是满态射，并且应用 $F$ 后诱导满射，
将推出 $s_i = s_j$，与假设矛盾。

\medskip\noindent
令 $G \subset S_n$ 为满足 $g(Z_t) = Z_t$ 的元素所成子群。考察 $S_n$ 在
$$
F(X)^n = F(X^n) = \coprod\nolimits_{t' \in T} F(Z_{t'})
$$
上的作用，可知
$G = \{g \in S_n \mid g(s_1, \ldots, s_n) \in F(Z_t)\}$。现选取第二个元素
$(s'_1, \ldots, s'_n) \in F(Z_t)$。上面已证明各 $s'_i$ 两两不同。因此可找到
$g \in S_n$，使 $g(s_1, \ldots, s_n) = (s'_1, \ldots, s'_n)$。换言之，
$G$ 在 $F(Z_t)$ 上的作用是传递的，证明完毕。
\end{proof}

\noindent
下面是一个关键引理。

\begin{lemma}
\label{pione-lemma-tame}
\begin{reference}
可与 \cite[Definition 7.2.4]{BS} 比较。
\end{reference}
设 $(\mathcal{C}, F)$ 为 Galois 范畴。设 $G = \text{Aut}(F)$ 如例
\ref{pione-example-from-C-F-to-G-sets} 所述。对任意连通对象 $X$（属于 $\mathcal{C}$），
$G$ 在 $F(X)$ 上的作用是传递的。
\end{lemma}

\begin{proof}
以下将不再明言地使用引理 \ref{pione-lemma-epi-mono} 的结果。令 $I$ 为
$\mathcal{C}$ 中 Galois 对象同构类所成之集。对每个 $i \in I$，令 $X_i$
为该同构类的代表。选取 $\gamma_i \in F(X_i)$（$i \in I$）。如下在 $I$ 上定义偏序：
$i \geq i'$ 当且仅当存在态射 $f_{ii'} : X_i \to X_{i'}$。给定这样的态射，
可以后复合一个自同构 $X_{i'} \to X_{i'}$，以保证
$F(f_{ii'})(\gamma_i) = \gamma_{i'}$。在这种规范化下，态射 $f_{ii'}$ 唯一。
注意 $I$ 是有向偏序集（Categories, Definition
\ref{categories-definition-directed-set}）：如果 $i_1, i_2 \in I$，则存在
Galois 对象 $Y$ 及态射 $Y \to X_{i_1} \times X_{i_2}$；这是对
$X_{i_1} \times X_{i_2}$ 的某个连通分支应用引理 \ref{pione-lemma-galois} 得到的。
于是 $Y \cong X_i$，其中某个 $i \in I$，且 $i \geq i_1$、$i \geq I_2$。

\medskip\noindent
我们断言函子 $F$ 同构于函子 $F'$，后者把 $X$ 送到
$$
F'(X) = \colim_I \Mor_\mathcal{C}(X_i, X)
$$
这是通过如下定义的函子变换 $t : F' \to F$ 实现的：给定 $f : X_i \to X$，
置 $t_X(f) = F(f)(\gamma_i)$。由 (\ref{pione-item-one-element}) 可知 $t_X$ 是单射。
为证满射，取 $\gamma \in F(X)$。由 Galois 范畴的定义，可立即归约到 $X$
连通的情形。再由引理 \ref{pione-lemma-galois}，可设 $X$ 是 Galois 的。此时
$X$ 同构于 $X_i$，其中 $i$ 是某个指标；并且可以选取同构 $X_i \to X$，使 $\gamma_i$ 映到
$\gamma$（由 Galois 对象的定义）。因此 $t$ 是同构。

\medskip\noindent
置 $A_i = \text{Aut}(X_i)$。我们断言，对 $i \geq i'$，存在典范映射
$h_{ii'} : A_i \to A_{i'}$，使得对所有 $a \in A_i$，图
$$
\xymatrix{
X_i \ar[d]_a \ar[r]_{f_{ii'}} & X_{i'} \ar[d]^{h_{ii'}(a)} \\
X_i \ar[r]^{f_{ii'}} & X_{i'}
}
$$
可交换。具体而言，令 $h_{ii'}(a) = a' : X_{i'} \to X_{i'}$ 为满足
$F(a')(\gamma_{i'}) = F(f_{ii'} \circ a)(\gamma_i)$.
的唯一自同构。如前，这使该图可交换，而且选择是唯一的。因此
$h_{i'i''} \circ h_{ii'} = h_{ii''}$
只要 $i \geq i' \geq i''$。由于 $F(X_i) \to F(X_{i'})$ 为满射，
可知 $A_i \to A_{i'}$ 为满射。取逆极限，得到群
$$
A = \lim_I A_i
$$
这是 profinite 群，因为各自同构群有限。由 Categories, Lemma
\ref{categories-lemma-nonempty-limit}，映射 $A \to A_i$ 对所有 $i$ 都是满射。

\medskip\noindent
由于 $A$ 的元素作用于逆系统 $X_i$，通过预复合得到 $A$ 在 $F'$ 上的右作用。
换言之，得到同态 $A^{opp} \to G$。由于 $A \to A_i$ 为满射，可知 $G$
在 $F(X_i)$ 上传递作用，对所有 $i$ 均如此。又由于每个连通对象都被某个 $X_i$ 支配，引理得证。
\end{proof}

\begin{proposition}
\label{pione-proposition-galois}
\begin{reference}
这是 \cite[Expos\'e V]{SGA1} 的弱版本。
证明取自 \cite[Theorem 7.2.5]{BS}。
\end{reference}
设 $(\mathcal{C}, F)$ 为 Galois 范畴。设 $G = \text{Aut}(F)$ 如例
\ref{pione-example-from-C-F-to-G-sets} 所述。函子
$F : \mathcal{C} \to \textit{Finite-}G\textit{-Sets}$
(\ref{pione-equation-remember}) 是等价。
\end{proposition}

\begin{proof}
以下将不再明言地使用引理 \ref{pione-lemma-epi-mono} 的结果。特别地，我们知道该函子
忠实。由引理 \ref{pione-lemma-tame}，对任意连通的 $X$，$G$ 在 $F(X)$ 上的作用
传递。因此，对一般的 $X$，函子 $F$ 把 $X$ 的连通分支（其存在性来自
Galois 范畴的一条公理）送到 $G$ 在 $F(X)$ 上的轨道。设 $X$ 和 $Y$ 为对象，
而 $s : F(X) \to F(Y)$ 为映射。于是
$\Gamma_s \subset F(X) \times F(Y)$（即 $s$ 的图像）是连通分支的并。因此存在
一个 $Z$，它是 $X \times Y$ 的若干连通分支之并，并带有单态射 $Z \to X \times Y$，
且 $F(Z) = \Gamma_s$。由于 $F(Z) \to F(X)$ 为双射，$Z \to X$ 是同构，
从而 $s = F(f)$，其中 $f : X \cong Z \to Y$ 为复合。因此 $F$ 全忠实。

\medskip\noindent
为完成证明，我们来证明 $F$ 本质满。只需证明 $G/H$ 位于本质像中，其中
$H \subset G$ 是任意有限指数开子群。由 $G$ 上拓扑的定义，存在有限族对象
$X_i$，使
$$
\Ker(G \longrightarrow \prod\nolimits_i \text{Aut}(F(X_i)))
$$
包含于 $H$。可设 $X_i$ 连通，对所有 $i$ 均如此。由引理 \ref{pione-lemma-galois}，
可以选取 Galois 对象 $Y$，使其映到 $\prod X_i$ 的某个连通分支。选取同构
$F(Y) = G/U$（位于 $G\textit{-sets}$ 中），其中 $U \subset G$ 为某个开子群。由于 $Y$
是 Galois 的，群 $\text{Aut}(Y) = \text{Aut}_{G\textit{-Sets}}(G/U)$
在 $F(Y) = G/U$ 上传递作用。这说明 $U$ 正规。由于 $F(Y)$ 满射到
$F(X_i)$，对每个 $i$ 均如此，可知 $U \subset H$。令 $M \subset \text{Aut}(Y)$ 为对应于
$$
(H/U)^{opp} \subset (G/U)^{opp} = \text{Aut}_{G\textit{-Sets}}(G/U)
= \text{Aut}(Y).
$$
的有限子群。置 $X = Y/M$，即 $X$ 是各箭头 $m : Y \to Y$（$m \in M$）
的余等化子。由于 $F$ 正合，可知 $F(X) = G/H$，证明完毕。
\end{proof}

\begin{lemma}
\label{pione-lemma-functoriality-galois}
设 $(\mathcal{C}, F)$ 和 $(\mathcal{C}', F')$ 为 Galois 范畴。设
$H : \mathcal{C} \to \mathcal{C}'$ 为正合函子。存在同构
$t : F' \circ H \to F$。$t$ 的选择决定连续同态
$h : G' = \text{Aut}(F') \to \text{Aut}(F) = G$ 以及 $2$-交换图
$$
\xymatrix{
\mathcal{C} \ar[r]_H \ar[d] & \mathcal{C}' \ar[d] \\
\textit{Finite-}G\textit{-Sets} \ar[r]^h &
\textit{Finite-}G'\textit{-Sets}
}
$$
映射 $h$ 与 $t$ 无关，仅允许相差 $G$ 的内自同构。反之，给定连续同态
$h : G' \to G$，存在正合函子 $H : \mathcal{C} \to \mathcal{C}'$ 和同构
$t$，按上述方式恢复 $h$。
\end{lemma}

\begin{proof}
由命题 \ref{pione-proposition-galois} 和引理 \ref{pione-lemma-single-out-profinite}，
可设 $\mathcal{C} = \textit{Finite-}G\textit{-Sets}$ 且 $F$ 为遗忘函子；
对 $\mathcal{C}'$ 同理。于是 $t$ 的存在性来自引理
\ref{pione-lemma-second-fundamental-functor}。映射 $h$ 来自经由 $t$ 的结构传递。
图的交换性显然。$h$ 在相差 $G$ 中某个元素的内共轭意义下唯一，这是因为
$t$ 的选择在相差 $G$ 的一个元素意义下唯一。最后一个断言是直接的。
\end{proof}





\section{函子与同态}
\label{pione-section-translation}

\noindent
设 $(\mathcal{C}, F)$、$(\mathcal{C}', F')$、$(\mathcal{C}'', F'')$
为 Galois 范畴。置 $G = \text{Aut}(F)$、$G' = \text{Aut}(F')$，以及
$G'' = \text{Aut}(F'')$。设 $H : \mathcal{C} \to \mathcal{C}'$
和 $H' : \mathcal{C}' \to \mathcal{C}''$ 为正合函子。
设 $h : G' \to G$ 和 $h' : G'' \to G'$ 为引理
\ref{pione-lemma-functoriality-galois} 所述的相应连续同态。本节考虑相应的 $2$-交换图
\begin{equation}
\label{pione-equation-translation}
\vcenter{
\xymatrix{
\mathcal{C} \ar[r]_H \ar[d] &
\mathcal{C}' \ar[r]_{H'} \ar[d] &
\mathcal{C}'' \ar[d] \\
\textit{Finite-}G\textit{-Sets} \ar[r]^h &
\textit{Finite-}G'\textit{-Sets} \ar[r]^{h'} &
\textit{Finite-}G''\textit{-Sets}
}
}
\end{equation}
并把序列 $1 \to G'' \to G' \to G \to 1$ 的正合性质与函子 $H$ 和 $H'$
的性质联系起来。

\begin{lemma}
\label{pione-lemma-functoriality-galois-surjective}
在图 (\ref{pione-equation-translation}) 中，下列条件等价：
\begin{enumerate}
\item $h : G' \to G$ 是满射；
\item $H : \mathcal{C} \to \mathcal{C}'$ 全忠实；
\item 如果 $X \in \Ob(\mathcal{C})$ 连通，则 $H(X)$ 连通；
\item 如果 $X \in \Ob(\mathcal{C})$ 连通，且存在态射
$*' \to H(X)$（位于 $\mathcal{C}'$ 中），则存在态射 $* \to X$；以及
\item 对任意对象 $X$（属于 $\mathcal{C}$），映射
$\Mor_\mathcal{C}(*, X) \to \Mor_{\mathcal{C}'}(*', H(X))$
为双射。
\end{enumerate}
这里 $*$ 和 $*'$ 分别是 $\mathcal{C}$ 和 $\mathcal{C}'$ 的终对象。
\end{lemma}

\begin{proof}
(5) $\Rightarrow$ (4) 和 (2) $\Rightarrow$ (5) 显然。

\medskip\noindent
假设 (3)。设 $X$ 为 $\mathcal{C}$ 的连通对象，并设 $*' \to H(X)$
为态射。由 (3)，$H(X)$ 连通，故 $*' \to H(X)$ 是同构。因此相应的
$G'$-集（对应于 $H(X)$）恰有一个元素，这意味着相应的 $G$-集
（对应于 $X$）有一个元素，进而意味着 $X$ 同构于 $\mathcal{C}$ 的终对象；
特别地，存在映射 $* \to X$。由此可见 (3) $\Rightarrow$ (4)。

\medskip\noindent
如果 (1) 成立，则函子
$\textit{Finite-}G\textit{-Sets} \to \textit{Finite-}G'\textit{-Sets}$
全忠实：在这种情形下，$G$-集之间的映射与 $G$ 的作用可交换，当且仅当它与
$G'$ 的作用可交换。因此 (1) $\Rightarrow$ (2)。

\medskip\noindent
如果 (1) 成立，则对 $G$-集 $X$，其 $G$-轨道与 $G'$-轨道一致。
因此 (1) $\Rightarrow$ (3)。

\medskip\noindent
为完成证明，只需证明 (4) 蕴含 (1)。如果 (1) 不成立，即 $h$ 不是满射，
则存在开子群 $U \subset G$，它包含 $h(G')$ 但不等于 $G$。于是有限
$G$-集 $M = G/U$ 上的作用是传递的，但 $G'$ 有一个不动点。对象
$X$（属于 $\mathcal{C}$，对应于 $M$）将与 (3) 矛盾。由此可见
(3) $\Rightarrow$ (1)，证明完毕。
\end{proof}

\begin{lemma}
\label{pione-lemma-composition-trivial}
在图 (\ref{pione-equation-translation}) 中，下列条件等价：
\begin{enumerate}
\item $h \circ h'$ 是平凡的；以及
\item $H' \circ H$ 的像由同构于终对象有限余积的对象组成。
\end{enumerate}
\end{lemma}

\begin{proof}
可以把 $H$ 和 $H'$ 替换为典范函子
$\textit{Finite-}G\textit{-Sets} \to \textit{Finite-}G'\textit{-Sets}
\to \textit{Finite-}G''\textit{-Sets}$；这些函子由 $h$ 和 $h'$ 决定。
于是所述命题变为：$G''$ 在每个
$G$-集上的作用平凡，当且仅当同态 $G'' \to G$ 平凡。这是显然的。
\end{proof}

\begin{lemma}
\label{pione-lemma-functoriality-galois-ses}
在图 (\ref{pione-equation-translation}) 中，下列条件等价：
\begin{enumerate}
\item 序列 $G'' \xrightarrow{h'} G' \xrightarrow{h} G \to 1$
在如下意义下正合：$h$ 是满射，$h \circ h'$ 平凡，并且 $\Ker(h)$ 是包含
$\Im(h')$ 的最小闭正规子群；
\item $H$ 全忠实，而且对象 $X'$（属于 $\mathcal{C}'$）位于 $H$ 的本质像中，
当且仅当 $H'(X')$ 同构于终对象的有限余积；以及
\item $H$ 全忠实，$H \circ H'$ 把每个对象送到终对象的有限余积，并且对
对象 $X'$（属于 $\mathcal{C}'$），如果 $H'(X')$ 是终对象的有限余积，则存在
对象 $X$（属于 $\mathcal{C}$）及满态射 $H(X) \to X'$。
\end{enumerate}
\end{lemma}

\begin{proof}
由引理 \ref{pione-lemma-functoriality-galois-surjective} 和
\ref{pione-lemma-composition-trivial}，可设 $H$ 全忠实，$h$ 是满射，
$H' \circ H$ 把对象送到终对象的不交并，并且 $h \circ h'$ 平凡。
令 $N \subset G'$ 为包含 $h'$ 的像的最小闭正规子群。显然
$N \subset \Ker(h)$。可设函子 $H$ 和 $H'$ 是典范函子
$\textit{Finite-}G\textit{-Sets} \to \textit{Finite-}G'\textit{-Sets}
\to \textit{Finite-}G''\textit{-Sets}$；这些函子由 $h$ 和 $h'$ 决定。

\medskip\noindent
假设 (2) 成立。这意味着，对有限 $G'$-集 $X'$，如果 $G''$ 的作用平凡，
则 $G'$ 的作用经过 $G$ 分解。把它应用于 $X' = G'/U'N$，其中 $U'$ 是
$G'$ 的充分小开子群。于是有 $\Ker(h) \subset U'N$，对所有 $U'$ 均如此。
由于 $N$ 闭，这蕴含 $\Ker(h) \subset N$，即 (1) 成立。

\medskip\noindent
假设 (1) 成立。这意味着 $N = \Ker(h)$。设 $X'$ 为有限 $G'$-集，且
$G''$ 的作用平凡。这意味着 $\Ker(G' \to \text{Aut}(X'))$ 是包含
$\Im(h')$ 的闭正规子群。因此它包含 $N = \Ker(h)$，而 $G'$ 在 $X'$
上的作用经过 $G$ 分解，即 (2) 成立。

\medskip\noindent
假设 (3) 成立。这意味着，对有限 $G'$-集 $X'$，如果 $G''$ 的作用平凡，
则存在 $G'$-集的满射 $X \to X'$，其中 $X$ 是 $G$-集。显然，这意味着
$G'$ 在 $X'$ 上的作用经过 $G$ 分解，即 (2) 成立。

\medskip\noindent
(2) $\Rightarrow$ (3) 立即成立。证明完毕。
\end{proof}

\begin{lemma}
\label{pione-lemma-functoriality-galois-injective}
在图 (\ref{pione-equation-translation}) 中，下列条件等价：
\begin{enumerate}
\item $h'$ 是单射；以及
\item 对每个连通对象 $X''$（属于 $\mathcal{C}''$），存在对象 $X'$
（属于 $\mathcal{C}'$）以及图
$$
X'' \leftarrow Y'' \rightarrow H(X')
$$
它位于 $\mathcal{C}''$ 中，其中 $Y'' \to X''$ 是满态射，而
$Y'' \to H(X')$ 是单态射。
\end{enumerate}
\end{lemma}

\begin{proof}
可以把 $H'$ 替换为有限 $G'$-集范畴与有限 $G''$-集范畴之间的相应函子。

\medskip\noindent
假设 $h' : G'' \to G'$ 是单射。设 $H'' \subset G''$ 为开子群。
由于 $G''$ 上的拓扑是由 $G'$ 诱导的拓扑，存在开子群 $H' \subset G'$，使
$(h')^{-1}(H') \subset H''$。于是所需图为
$$
G''/H'' \leftarrow G''/(h')^{-1}(H') \rightarrow G'/H'
$$
反之，假设 (2) 对函子
$\textit{Finite-}G'\textit{-Sets} \to \textit{Finite-}G''\textit{-Sets}$.
成立。设 $g'' \in \Ker(h')$。任取开子群 $H'' \subset G''$。
由假设，存在有限 $G'$-集 $X'$ 及图
$$
G''/H'' \leftarrow Y'' \rightarrow X'
$$
它由 $G''$-集构成，左箭头为满射，右箭头为单射。由于 $g''$ 位于 $h'$ 的核中，
可知 $g''$ 在 $X'$ 上作用平凡。因此 $g''$ 在 $Y''$ 上作用平凡，进而在
$G''/H''$ 上作用平凡。故 $g'' \in H''$。由于这对所有开子群都成立，可知
$g''$ 是单位元，正如所需。
\end{proof}

\begin{lemma}
\label{pione-lemma-functoriality-galois-normal}
在图 (\ref{pione-equation-translation}) 中，下列条件等价：
\begin{enumerate}
\item $h'$ 的像是正规的；以及
\item 对每个连通对象 $X'$（属于 $\mathcal{C}'$），如果存在从
$\mathcal{C}''$ 的终对象到 $H'(X')$ 的态射，则 $H'(X')$ 同构于终对象的有限余积。
\end{enumerate}
\end{lemma}

\begin{proof}
这转化为关于连续群同态 $h' : G'' \to G'$ 的如下陈述：$h'$ 的像正规，
当且仅当每个开子群 $U' \subset G'$，只要它包含 $h'(G'')$，也包含
$h'(G'')$ 的每个共轭。结论很容易由此推出；略去若干细节。
\end{proof}






\section{有限 \'etale 态射}
\label{pione-section-finite-etale}

\noindent
本节证明关于有限 \'etale 态射的足够多基本结果，以便构造 \'etale 基本群。

\medskip\noindent
设 $X$ 为概形。用记号 $\textit{F\'Et}_X$ 表示在 $X$ 上有限且 \'etale
的概形所成范畴。于是
\begin{enumerate}
\item $\textit{F\'Et}_X$ 的对象是有限 \'etale 态射
$Y \to X$，其目标为 $X$；以及
\item $\textit{F\'Et}_X$ 中从 $Y \to X$ 到 $Y' \to X$ 的态射，是使图
$Y \to Y'$
$$
\xymatrix{
Y \ar[rr] \ar[rd] & &  Y' \ar[ld] \\
& X
}
$$
可交换的态射。
\end{enumerate}
我们常把 $\textit{F\'Et}_X$ 的对象称为 $X$ 的{\it 有限 \'etale 覆盖}
（即使 $Y$ 为空）。事实上，在概形范畴上存在栈
$p : \textit{F\'Et} \to \Sch$，其在 $X$ 处的纤维正是刚定义的范畴
$\textit{F\'Et}_X$。见 Examples of Stacks, Section
\ref{examples-stacks-section-finite-etale}.

\begin{example}
\label{pione-example-finite-etale-geometric-point}
设 $k$ 为代数闭域，且 $X = \Spec(k)$。在这种情形下，
$\textit{F\'Et}_X$ 等价于有限集范畴。更一般地，当 $k$ 可分代数闭时仍然成立。
这是因为在 $k$ 上 \'etale 的概形是 $k$ 的有限可分扩域之谱的不交并；见
Morphisms, Lemma \ref{morphisms-lemma-etale-over-field}.
\end{example}

\begin{lemma}
\label{pione-lemma-finite-etale-covers-limits-colimits}
设 $X$ 为概形。范畴 $\textit{F\'Et}_X$ 有有限极限和有限余极限，并且对任意态射
$X' \to X$，基变换函子
$\textit{F\'Et}_X \to \textit{F\'Et}_{X'}$ 正合。
\end{lemma}

\begin{proof}
有限极限与左正合性。由 Categories, Lemma
\ref{categories-lemma-finite-limits-exist}，只需证明
$\textit{F\'Et}_X$ 有终对象和纤维积。这是显然的，因为 $X$ 上所有概形的范畴
有终对象（即 $X$）和纤维积。此外，在 $X$ 上有限 \'etale 的概形之纤维积仍
在 $X$ 上有限 \'etale。显然，基变换与这些运算可交换，因而基变换左正合
（Categories, Lemma
\ref{categories-lemma-characterize-left-exact}).

\medskip\noindent
有限余极限与右正合性。由 Categories, Lemma
\ref{categories-lemma-colimits-exist}，只需证明 $\textit{F\'Et}_X$
有有限余积和余等化子。有限余积由不交并给出（空余积是空概形）。
设 $a, b : Z \to Y$ 为 $\textit{F\'Et}_X$ 的两个态射。由于
$Z \to X$ 和 $Y \to X$ 有限 \'etale，可以写成
$Z = \underline{\Spec}(\mathcal{C})$ 和 $Y = \underline{\Spec}(\mathcal{B})$，
其中某些有限局部自由 $\mathcal{O}_X$-代数 $\mathcal{C}$ 和 $\mathcal{B}$
给出上述表示。态射 $a, b$ 诱导两个映射
$a^\sharp, b^\sharp : \mathcal{B} \to \mathcal{C}$。令
$\mathcal{A} = \text{Eq}(a^\sharp, b^\sharp)$ 为它们的等化子。如果
$$
\underline{\Spec}(\mathcal{A}) \longrightarrow X
$$
有限 \'etale，则显然这就是余等化子（毕竟，$\textit{F\'Et}_X$ 的任意对象
都可写成某个 $\mathcal{O}_X$-代数层的相对谱）。把 $X$ 替换为某个
\'etale 覆盖的成员后即可验证这一点（Descent, Lemmas
\ref{descent-lemma-descending-property-finite} 和
\ref{descent-lemma-descending-property-etale}）。因此由 \'Etale Morphisms,
Lemma \ref{etale-lemma-finite-etale-etale-local}，可以假设
$Y = \coprod_{i = 1, \ldots, n} X$ 和 $Z = \coprod_{j = 1, \ldots, m} X$。
此时
$$
\mathcal{C} = \prod\nolimits_{1 \leq j \leq m} \mathcal{O}_X
\quad\text{且}\quad
\mathcal{B} = \prod\nolimits_{1 \leq i \leq n} \mathcal{O}_X
$$
进一步替换为某个开覆盖的成员后，可设 $a, b$ 对应于映射
$a_s, b_s : \{1, \ldots, m\} \to \{1, \ldots, n\}$；也就是说，直和项
$X$（属于 $Z$ 并对应于指标 $j$）映到直和项 $X$（属于 $Y$ 并对应于指标
$a_s(j)$、相应地 $b_s(j)$），分别在态射 $a$、$b$ 下。
令 $\{1, \ldots, n\} \to T$ 为 $a_s, b_s$ 的余等化子。于是
$$
\mathcal{A} = \prod\nolimits_{t \in T} \mathcal{O}_X
$$
其谱显然在 $X$ 上有限 \'etale。略去它与基变换相容的验证。
因此基变换是右正合函子。
\end{proof}

\begin{remark}
\label{pione-remark-colimits-commute-forgetful}
设 $X$ 为概形。考虑自然函子
$F_1 : \textit{F\'Et}_X \to \Sch$ 和 $F_2 : \textit{F\'Et}_X \to \Sch/X$。
则
\begin{enumerate}
\item 函子 $F_1$ 和 $F_2$ 与有限余极限可交换。
\item 函子 $F_2$ 与有限极限可交换。
\item 函子 $F_1$ 与连通有限极限可交换，即与等化子和纤维积可交换。
\end{enumerate}
关于极限的结果立即来自引理
\ref{pione-lemma-finite-etale-covers-limits-colimits} 证明中的讨论以及
Categories, Lemma \ref{categories-lemma-connected-limit-over-X}。
显然，$F_1$ 和 $F_2$ 与有限余积可交换。由 Categories, Lemma
\ref{categories-lemma-characterize-left-exact}
的对偶，只需证明 $F_1$ 和 $F_2$ 与余等化子可交换。在引理
\ref{pione-lemma-finite-etale-covers-limits-colimits} 的证明中，我们看到
$\textit{F\'Et}_X$ 中的余等化子在 \'etale 局部具有如下形式：
$$
\xymatrix{
\coprod_{j \in J} U \ar@<1ex>[r]^a \ar@<-1ex>[r]_b &
\coprod_{i \in I} U \ar[r] &
\coprod_{t \in \text{Coeq}(a, b)} U
}
$$
这在概形范畴中显然是余等化子。因此，结论来自余等化子性质是 fpqc 局部的这一事实；
其精确表述见 Descent, Lemma \ref{descent-lemma-coequalizer-fpqc-local}。
\end{remark}

\begin{lemma}
\label{pione-lemma-internal-hom-finite-etale}
设 $X$ 为概形。给定 $U, V$ 在 $X$ 上有限 \'etale，存在概形 $W$ 在 $X$
上有限 \'etale，使
$$
\Mor_X(X, W) = \Mor_X(U, V)
$$
并且在任意基变换后同一结论仍然成立。
\end{lemma}

\begin{proof}
由 More on Morphisms, Lemma
\ref{more-morphisms-lemma-hom-from-finite-locally-free-separated-lqf}，
存在概形 $W$，它表示 $\mathit{Mor}_X(U, V)$。（这里使用 Morphisms, Lemmas
\ref{morphisms-lemma-etale-locally-quasi-finite} 所说明的 \'etale 态射局部拟有限，
以及有限态射是分离的。）这个概形在任意基变换后显然仍满足该公式。
为完成证明，必须证明 $W \to X$ 有限 \'etale。把 $X$ 替换为某个
\'etale 覆盖的成员后即可验证这一点（Descent, Lemmas
\ref{descent-lemma-descending-property-finite} 和
\ref{descent-lemma-descending-property-separated}）。因此由 \'Etale Morphisms,
Lemma \ref{etale-lemma-finite-etale-etale-local}，可设
$U = \coprod_{i = 1, \ldots, n} X$ 和
$V = \coprod_{j = 1, \ldots, m} X$。在这种情形下，直接检验可得（略去细节）
$W = \coprod_{\alpha : \{1, \ldots, n\} \to \{1, \ldots, m\}} X$
，证明完毕。
\end{proof}

\noindent
设 $X$ 为概形。$X$ 的一个{\it 几何点}是态射
$\Spec(k) \to X$，其中 $k$ 代数闭。这样的点通常记作 $\overline{x}$，
即用带上划线的小写字母表示。我们常用 $\overline{x}$ 同时表示概形
$\Spec(k)$ 和该态射，并用 $\kappa(\overline{x})$ 表示 $k$。称
$\overline{x}$ {\it 位于} $x$ {\it 之上}，是指 $x \in X$ 是
$\overline{x}$ 的像。进一步讨论见
\'Etale Cohomology, Section \ref{etale-cohomology-section-stalks}.
给定 $\overline{x}$ 和 \'etale 态射 $U \to X$，可以考虑
$$
|U_{\overline{x}}| : \text{the underlying set of points of the
scheme }U_{\overline{x}} = U \times_X \overline{x}
$$
由于 $U_{\overline{x}}$ 作为 $\overline{x}$ 上的概形，是若干个
$\overline{x}$ 的不交并（Morphisms, Lemma
\ref{morphisms-lemma-etale-over-field}），也可以把该集合描述为
$$
|U_{\overline{x}}| =
\left\{
\begin{matrix}
\text{commutative} \\
\text{diagrams}
\end{matrix}
\vcenter{
\xymatrix{
\overline{x} \ar[rd]_{\overline{x}} \ar[r]_{\overline{u}} & U \ar[d] \\
& X
}
}
\right\}
$$
对应 $U \mapsto |U_{\overline{x}}|$ 是一个函子，常记作
$F_{\overline{x}}$。

\begin{lemma}
\label{pione-lemma-finite-etale-connected-galois-category}
设 $X$ 为连通概形。设 $\overline{x}$ 为几何点。函子
$$
F_{\overline{x}} : \textit{F\'Et}_X \longrightarrow \textit{Sets},\quad
Y \longmapsto |Y_{\overline{x}}|
$$
定义了一个 Galois 范畴（Definition \ref{pione-definition-galois-category}）。
\end{lemma}

\begin{proof}
把 $\textit{F\'Et}_{\overline{x}}$ 与有限集范畴识别后（Example
\ref{pione-example-finite-etale-geometric-point}），可知函子
$F_{\overline{x}}$ 正是态射 $\overline{x} \to X$ 的基变换函子。因此，由引理
\ref{pione-lemma-finite-etale-covers-limits-colimits} 可知 $\textit{F\'Et}_X$
有有限极限和有限余极限，并且 $F_{\overline{x}}$ 正合。还将使用如下事实：
$\textit{F\'Et}_X$ 中的有限极限与 $X$ 上概形范畴中的相应有限极限一致；见
Remark \ref{pione-remark-colimits-commute-forgetful}。

\medskip\noindent
如果 $Y' \to Y$ 是 $\textit{F\'Et}_X$ 中的单态射，则
$Y' \to Y' \times_Y Y'$ 是同构，因而 $Y' \to Y$ 是概形的单态射。
于是 $Y' \to Y$ 是开浸入（\'Etale Morphisms, Theorem
\ref{etale-theorem-etale-radicial-open}）。由于 $Y'$ 在 $X$ 上有限，而
$Y$ 在 $X$ 上分离，态射 $Y' \to Y$ 有限（Morphisms, Lemma
\ref{morphisms-lemma-finite-permanence}），因而是闭的（Morphisms, Lemma
\ref{morphisms-lemma-finite-proper}），所以它是 $Y$ 的开闭子概形之包含。
由此，$Y$ 是范畴 $\textit{F\'Et}_X$ 的连通对象（如 Definition
\ref{pione-definition-galois-category}），当且仅当 $Y$ 作为概形连通。于是由
Topology, Lemma \ref{topology-lemma-finite-fibre-connected-components}，
$Y$ 无论作为概形还是在 Definition \ref{pione-definition-galois-category}
的意义下，都是其连通分支的有限余积。

\medskip\noindent
设 $Y \to Z$ 为 $\textit{F\'Et}_X$ 中的态射，它诱导双射
$F_{\overline{x}}(Y) \to F_{\overline{x}}(Z)$。必须证明 $Y \to Z$
是同构。由上可设 $Z$ 连通。由于 $Y \to Z$ 有限 \'etale，因而有限局部自由，
只需证明 $Y \to Z$ 是次数为 $1$ 的有限局部自由态射。这在 $Z$ 中位于
$\overline{x}$ 之上的任意点的某个邻域内成立；由于 $Z$ 连通且次数局部常值，
结论成立。
\end{proof}



\section{基本群}
\label{pione-section-fundamental-groups}

\noindent
本节定义 Grothendieck 的代数基本群。由于引理
\ref{pione-lemma-finite-etale-connected-galois-category}，下述定义是有意义的。

\begin{definition}
\label{pione-definition-fundamental-group}
设 $X$ 为连通概形。设 $\overline{x}$ 为 $X$ 的几何点。概形 $X$ 以
$\overline{x}$ 为{\it 基点}的{\it 基本群}是群
$$
\pi_1(X, \overline{x}) = \text{Aut}(F_{\overline{x}})
$$
即纤维函子
$F_{\overline{x}} : \textit{F\'Et}_X \to \textit{Sets}$
的自同构群，并赋予引理 \ref{pione-lemma-aut-inverse-limit} 所给的典范 profinite 拓扑。
\end{definition}

\noindent
把上述结果与 Section \ref{pione-section-galois} 的内容结合，得到下述定理。

\begin{theorem}
\label{pione-theorem-fundamental-group}
设 $X$ 为连通概形。设 $\overline{x}$ 为 $X$ 的几何点。
\begin{enumerate}
\item 纤维函子 $F_{\overline{x}}$ 定义范畴等价
$$
\textit{F\'Et}_X \longrightarrow
\textit{Finite-}\pi_1(X, \overline{x})\textit{-Sets}
$$
\item 给定第二个几何点 $\overline{x}'$（属于 $X$），存在同构
$t : F_{\overline{x}} \to F_{\overline{x}'}$。它给出与 (1) 中等价相容的同构
$\pi_1(X, \overline{x}) \to \pi_1(X, \overline{x}')$。在相差内共轭的意义下，
该同构与 $t$ 无关。
\item 给定连通概形之间的态射 $f : X \to Y$，记
$\overline{y} = f \circ \overline{x}$。存在典范连续同态
$$
f_* : \pi_1(X, \overline{x}) \to \pi_1(Y, \overline{y})
$$
使图
$$
\xymatrix{
\textit{F\'Et}_Y \ar[r]_{\text{base change}} \ar[d]_{F_{\overline{y}}} &
\textit{F\'Et}_X \ar[d]^{F_{\overline{x}}} \\
\textit{Finite-}\pi_1(Y, \overline{y})\textit{-Sets} \ar[r]^{f_*} &
\textit{Finite-}\pi_1(X, \overline{x})\textit{-Sets}
}
$$
可交换。
\end{enumerate}
\end{theorem}

\begin{proof}
(1) 来自引理 \ref{pione-lemma-finite-etale-connected-galois-category} 和命题
\ref{pione-proposition-galois}。(2) 是引理 \ref{pione-lemma-functoriality-galois}
的特殊情形。对于 (3)，注意图
$$
\xymatrix{
\textit{F\'Et}_Y \ar[r] \ar[d]_{F_{\overline{y}}} &
\textit{F\'Et}_X \ar[d]^{F_{\overline{x}}} \\
\textit{Sets} \ar@{=}[r] & \textit{Sets}
}
$$
可交换（实际上可交换，而不只是 $2$-交换），因为
$\overline{y} = f \circ \overline{x}$。因此可以把引理
\ref{pione-lemma-functoriality-galois} 与隐含的函子变换一起应用，从而得到 (3)。
\end{proof}

\begin{lemma}
\label{pione-lemma-fundamental-group-Galois-group}
设 $K$ 为域，并置 $X = \Spec(K)$。设 $\overline{K}$ 为代数闭包，并以
$\overline{x} : \Spec(\overline{K}) \to X$ 表示相应的几何点。设
$K^{sep} \subset \overline{K}$ 为可分闭包。
\begin{enumerate}
\item 引理 \ref{pione-lemma-sheaves-point} 的函子诱导等价
$$
\textit{F\'Et}_X \longrightarrow
\textit{Finite-}\text{Gal}(K^{sep}/K)\textit{-Sets}.
$$
它与 $F_{\overline{x}}$ 及函子
$\textit{Finite-}\text{Gal}(K^{sep}/K)\textit{-Sets} \to \textit{Sets}$.
相容。
\item 它诱导典范同构
$$
\text{Gal}(K^{sep}/K) \longrightarrow \pi_1(X, \overline{x})
$$
，这是 profinite 拓扑群的同构。
\end{enumerate}
\end{lemma}

\begin{proof}
引理 \ref{pione-lemma-sheaves-point} 的函子与函子 $F_{\overline{x}}$ 相同，
因为对任意 $Y$，若它在 $X$ 上 \'etale，则有
$$
\Mor_X(\Spec(\overline{K}), Y) = \Mor_X(\Spec(K^{sep}), Y)
$$
事实上，如引理 \ref{pione-lemma-sheaves-point} 的证明所示，
$Y = \coprod_{i \in I} \Spec(L_i)$，其中 $L_i/K$ 是 $K$ 的有限可分扩张。
因此任意 $K$-代数同态 $L_i \to \overline{K}$ 都经过 $K^{sep}$ 分解。
还要注意，$F_{\overline{x}}(Y)$ 有限，当且仅当 $I$ 有限，又当且仅当
$Y \to X$ 有限 \'etale。这证明了 (1)。

\medskip\noindent
(2) 是 (1)、引理 \ref{pione-lemma-functoriality-galois} 和引理
\ref{pione-lemma-single-out-profinite} 的形式推论。（另见下述评注。）
\end{proof}

\begin{remark}
\label{pione-remark-variance}
在引理 \ref{pione-lemma-fundamental-group-Galois-group} 的情形下，下面给出同构
$\text{Gal}(K^{sep}/K) \to
\pi_1(X, \overline{x}) = \text{Aut}(F_{\overline{x}})$ 的更明确构造。注意
$\text{Gal}(K^{sep}/K) = \text{Aut}(\overline{K}/K)$
，因为 $\overline{K}$ 是 $K^{sep}$ 的完美化。由于
$F_{\overline{x}}(Y) = \Mor_X(\Spec(\overline{K}), Y)$，可以考虑映射
$$
\text{Aut}(\overline{K}/K) \times F_{\overline{x}}(Y) \to F_{\overline{x}}(Y),
\quad
(\sigma, \overline{y}) \mapsto
\sigma \cdot \overline{y} = \overline{y} \circ \Spec(\sigma)
$$
这是一个作用，因为
$$
\sigma\tau \cdot \overline{y} =
\overline{y} \circ \Spec(\sigma\tau) =
\overline{y} \circ \Spec(\tau) \circ \Spec(\sigma) =
\sigma \cdot (\tau \cdot \overline{y})
$$
该作用关于 $Y \in \textit{F\'Et}_X$ 具有函子性，从而得到所需映射。
\end{remark}





\section{连通概形的 Galois 覆盖}
\label{pione-section-finite-etale-under-galois}

\noindent
设 $X$ 为带有几何点 $\overline{x}$ 的连通概形。由于
$F_{\overline{x}} : \textit{F\'Et}_X \to \textit{Sets}$ 是 Galois 范畴
（Lemma \ref{pione-lemma-finite-etale-connected-galois-category}），Section
\ref{pione-section-galois} 的内容适用。本节把若干术语和结果明确转移到概形与有限
\'etale 态射的情形。

\medskip\noindent
称有限 \'etale 态射 $Y \to X$ 为{\it Galois 覆盖}，如果 $Y$ 定义了
$\textit{F\'Et}_X$ 的 Galois 对象。对有限 \'etale 态射 $Y \to X$，置
$G = \text{Aut}_X(Y)$，则下列条件等价：
\begin{enumerate}
\item $Y$ 是 $X$ 的 Galois 覆盖；
\item $Y$ 连通，且 $|G|$ 等于 $Y \to X$ 的次数；
\item $Y$ 连通，且 $G$ 在 $F_{\overline{x}}(Y)$ 上传递作用；以及
\item $Y$ 连通，且 $G$ 在 $F_{\overline{x}}(Y)$ 上单传递作用。
\end{enumerate}
这立即来自 Section \ref{pione-section-galois} 中的讨论。

\medskip\noindent
对任意有限 \'etale 态射 $f : Y \to X$，如果 $Y$ 连通，则存在有限
\'etale Galois 覆盖 $Y' \to X$，它支配 $Y$（Lemma \ref{pione-lemma-galois}）。

\medskip\noindent
$\textit{F\'Et}_X$ 的 Galois 对象经由等价
$$
F_{\overline{x}} : \textit{F\'Et}_X \to
\textit{Finite-}\pi_1(X, \overline{x})\textit{-Sets}
$$
（Theorem \ref{pione-theorem-fundamental-group}），对应于有限
$\pi_1(X, \overline{x})\textit{-Sets}$，它们形如
$G = \pi_1(X, \overline{x})/H$，其中 $H$ 为正规开子群。等价地，
如果 $G$ 是有限群，且 $\pi_1(X, \overline{x}) \to G$ 是连续满射，则把
$G$ 看作 $\pi_1(X, \overline{x})$-集时，它对应于 Galois 覆盖。

\medskip\noindent
如果 $Y_i \to X$（$i = 1, 2$）是 Galois 群为 $G_i$ 的有限 \'etale
Galois 覆盖，则存在有限 \'etale Galois 覆盖 $Y \to X$，其 Galois 群是
$G_1 \times G_2$ 的子群。事实上，取相应的连续同态
$\pi_1(X, \overline{x}) \to G_i$，并令 $G$ 为诱导连续同态
$\pi_1(X, \overline{x}) \to G_1 \times G_2$ 的像。









\section{基本群的拓扑不变性}
\label{pione-section-topological-invariance}

\noindent
本节的主要结果是，连通概形的泛同胚会诱导基本群之间的同构。
见命题 \ref{pione-proposition-universal-homeomorphism}。

\medskip\noindent
我们通常不直接证明两个概形具有相同的基本群，而是证明它们的有限 \'etale
覆盖范畴相同。当然，只要它们连通，这就蕴含它们的基本群相同。

\begin{lemma}
\label{pione-lemma-what-equivalence-gives}
设 $f : X \to Y$ 是拟紧拟分离概形的态射，使得基变换函子
$\textit{F\'Et}_Y \to \textit{F\'Et}_X$ 是范畴等价。此时
\begin{enumerate}
\item $f$ 诱导同胚 $\pi_0(X) \to \pi_0(Y)$；
\item 若 $X$ 连通，等价地，若 $Y$ 连通，则
$\pi_1(X, \overline{x}) = \pi_1(Y, \overline{y})$。
\end{enumerate}
\end{lemma}

\begin{proof}
设 $Y = Y_0 \amalg Y_1$ 为到两个非空开闭子概形的不交并分解。我们断言
$f(X)$ 与两个 $Y_i$ 都相交。事实上，若非如此，不妨设
$f(X) \subset Y_1$，则可考虑有限 \'etale 态射 $V = Y_1 \to Y$。
它不是同构，但 $V \times_Y X \to X$ 是同构，矛盾。

\medskip\noindent
设 $X = X_0 \amalg X_1$ 为到两个开闭子概形的不交并分解。考虑有限
\'etale 态射 $U = X_1 \to X$。于是有 $U = X \times_Y V$，其中
$V \to Y$ 是某个有限 \'etale 态射。态射 $V \to Y$ 的次数局部常值，
故可得如下到开闭子概形的分解
$Y = \coprod_{d \geq 0} Y_d$，使得 $V \to Y$ 的次数为 $d$ 的部分恰为
$Y_d$。因为 $f^{-1}(Y_d) = \emptyset$ 在 $d > 1$ 时成立，由上面的论证
可知，$Y_d = \emptyset$ 在 $d > 1$ 时成立。进而可知，
$f^{-1}(Y_i) = X_i$ 对 $i = 0, 1$ 成立。

\medskip\noindent
于是 $f^{-1}$ 在 $Y$ 的开闭子集的集合与 $X$ 的开闭子集的集合之间诱导
双射。注意 $X$ 和 $Y$ 都是谱空间，见《性质》引理
\ref{properties-lemma-quasi-compact-quasi-separated-spectral}。
由《拓扑》引理 \ref{topology-lemma-connected-component-intersection}，
谱空间的开闭子集格决定其连通分支的集合。因此
$\pi_0(X) \to \pi_0(Y)$ 是双射。由于 $\pi_0(X)$ 和 $\pi_0(Y)$ 都是
profinite 空间（《拓扑》引理 \ref{topology-lemma-pi0-profinite}），
由《拓扑》引理 \ref{topology-lemma-bijective-map} 可知
$\pi_0(X) \to \pi_0(Y)$ 是同胚。这证明了 (1)。断言 (2) 是立即的。
\end{proof}

\noindent
下一个引理说明，Hensel 对的基本群就是相应闭子集的基本群。

\begin{lemma}
\label{pione-lemma-gabber}
设 $(A, I)$ 是 Hensel 对。置 $X = \Spec(A)$ 且 $Z = \Spec(A/I)$。函子
$$
\textit{F\'Et}_X \longrightarrow \textit{F\'Et}_Z,\quad
U \longmapsto U \times_X Z
$$
是范畴等价。
\end{lemma}

\begin{proof}
这只是《代数补篇》引理
\ref{more-algebra-lemma-finite-etale-equivalence} 的几何转述。
\end{proof}

\noindent
下一个引理说明，增厚的基本群与原概形的基本群相同。我们将在下面关于泛同胚
的较强命题的证明中使用这一点。

\begin{lemma}
\label{pione-lemma-thickening}
设 $X \subset X'$ 是概形的增厚。函子
$$
\textit{F\'Et}_{X'} \longrightarrow \textit{F\'Et}_X,\quad
U' \longmapsto U' \times_{X'} X
$$
是范畴等价。
\end{lemma}

\begin{proof}
关于增厚的讨论见《态射补篇》第
\ref{more-morphisms-section-thickenings} 节。设 $U' \to X'$ 是
\'etale 态射，并且 $U = U' \times_{X'} X \to X$ 是有限 \'etale
态射。则 $U' \to X'$ 也是有限 \'etale 态射。例如，这可由《态射补篇》
引理 \ref{more-morphisms-lemma-properties-that-extend-over-thickenings}
得到。现在，若 $X \subset X'$ 是有限阶增厚，则结合这一观察与《\'Etale
态射》定理 \ref{etale-theorem-remarkable-equivalence} 即可证明本引理。
下面我们将对一般的增厚证明本引理，不过建议读者跳过这个证明。

\medskip\noindent
设 $X' = \bigcup X_i'$ 是仿射开覆盖。置
$X_i = X \times_{X'} X_i'$，$X_{ij}' = X'_i \cap X'_j$，
$X_{ij} = X \times_{X'} X_{ij}'$, $X_{ijk}' = X'_i \cap X'_j \cap X'_k$,
$X_{ijk} = X \times_{X'} X_{ijk}'$。
假设对每个增厚
$X_i \subset X'_i$、$X_{ij} \subset X_{ij}'$ 和
$X_{ijk} \subset X_{ijk}'$ 都能证明该定理。则由概形的相对粘合，可得
$X \subset X'$ 的结论，见《构造》第
\ref{constructions-section-relative-glueing} 节。注意，概形
$X_i'$、$X_{ij}'$、$X_{ijk}'$ 作为仿射概形的开子概形，都是分离的。
把上述论证再重复一次，便可归约到概形
$X'_i$、$X_{ij}'$、$X_{ijk}'$ 都仿射的情形。

\medskip\noindent
在仿射情形下，有 $X' = \Spec(A')$ 和 $X = \Spec(A'/I')$，其中
$I'$ 是局部幂零理想。于是 $(A', I')$ 是 Hensel 对（《代数补篇》引理
\ref{more-algebra-lemma-locally-nilpotent-henselian}），结论便由引理
\ref{pione-lemma-gabber} 得出（在此情形下，该引理的证明容易得多）。
\end{proof}

\noindent
证明下一个命题的“正确”方式，是从 \'etale 景的不变性推出它，见《\'Etale
上同调》定理 \ref{etale-cohomology-theorem-topological-invariance}。

\begin{proposition}
\label{pione-proposition-universal-homeomorphism}
设 $f : X \to Y$ 是概形的泛同胚。则
$$
\textit{F\'Et}_Y \longrightarrow \textit{F\'Et}_X,\quad
V \longmapsto V \times_Y X
$$
是范畴等价。因此，若 $X$ 和 $Y$ 连通，则 $f$ 诱导基本群的同构
$\pi_1(X, \overline{x}) \to \pi_1(Y, \overline{y})$。
\end{proposition}

\begin{proof}
回忆，泛同胚恰好就是整、泛单射且满射的态射，见《态射》引理
\ref{morphisms-lemma-universal-homeomorphism}。特别地，由《态射》引理
\ref{morphisms-lemma-universally-injective}，对角态射
$\Delta : X \to X \times_Y X$ 是增厚。因此，由引理
\ref{pione-lemma-thickening} 可知，给定有限 \'etale 态射 $U \to X$，存在唯一的同构
$$
\varphi : U \times_Y X \to X \times_Y U
$$
；这是 $X \times_Y X$ 上两个有限 \'etale 概形之间的同构，并且沿 $\Delta$ 拉回后，
成为恒等态射 $\text{id} : U \to U$；该态射定义在 $X$ 上。由于
$X \to X \times_Y X \times_Y X$ 也是增厚（它是双射闭浸入），可知
$(U, \varphi)$ 是相对于 $X/Y$ 的下降资料。由《\'Etale 态射》命题
\ref{etale-proposition-effective} 可知，$U = X \times_Y V$，其中
$V \to Y$ 拟紧、分离且 \'etale。我们略去 $V \to Y$ 有限的证明（提示：
态射 $U \to V$ 是满射，而 $U \to Y$ 是整态射）。因此
$\textit{F\'Et}_Y \to \textit{F\'Et}_X$ 是本质满的。

\medskip\noindent
与上面同样论证可知，给定 $V_1 \to Y$ 和 $V_2 \to Y$ 这两个
$\textit{F\'Et}_Y$ 中的对象，任意态射
$a : X \times_Y V_1 \to X \times_Y V_2$（在 $X$ 上）都与典范下降资料相容。因此，由
《\'Etale 态射》引理 \ref{etale-lemma-fully-faithful-cases}，$a$ 下降为
态射 $V_1 \to V_2$；该态射定义在 $Y$ 上。
\end{proof}










\section{固有概形的有限 \'etale 覆盖}
\label{pione-section-finite-etale-over-proper}

\noindent
本节将证明，Hensel 局部环上的连通固有概形与其特殊纤维具有相同的基本群。
我们还将证明该结果在 Hensel 对情形下的一个变体。

\medskip\noindent
我们还将证明，代数闭域 $k$ 上的连通固有概形的基本群在把 $k$ 换成
一个代数闭域扩张后保持不变。

\medskip\noindent
我们不在连通情形下陈述并证明这些结果，而是在一般情形下证明它们，
并把利用引理 \ref{pione-lemma-what-equivalence-gives} 推出基本群结论的工作
留给读者。

\begin{lemma}
\label{pione-lemma-finite-etale-on-proper-over-henselian}
设 $A$ 是 Hensel 局部环，$X$ 是 $A$ 上的固有概形，其闭纤维为 $X_0$。
则函子
$$
\textit{F\'Et}_X \to \textit{F\'Et}_{X_0},\quad
U \longmapsto U_0 = U \times_X X_0
$$
是范畴等价。
\end{lemma}

\begin{proof}
这里给出的证明是应用代数化与逼近的一个例子。我们分若干步进行。

\medskip\noindent
当 $A$ 是完备 Noether 局部环时的本质满性。令
$X_n = X \times_{\Spec(A)} \Spec(A/\mathfrak m^{n + 1})$。
由《\'Etale 态射》定理 \ref{etale-theorem-remarkable-equivalence}，嵌入
$$
X_0 \to X_1 \to X_2 \to \ldots
$$
在 $X_0$ 上的 \'etale 概形范畴与在 $X_n$ 上的 \'etale 概形范畴之间
诱导范畴等价。此外，若 $U_n \to X_n$ 对应于有限 \'etale 态射
$U_0 \to X_0$，则 $U_n \to X_n$ 也是有限的；例如，这可由《态射补篇》
引理 \ref{more-morphisms-lemma-thicken-property-morphisms-cartesian} 得到。
在这种情形下，态射 $U_0 \to \Spec(A/\mathfrak m)$ 是固有的，因为
$X_0$ 在 $A/\mathfrak m$ 上固有。因此可应用 Grothendieck
代数化定理（采用《概形上同调》引理
\ref{coherent-lemma-algebraize-formal-scheme-finite-over-proper} 的形式），
得到有限态射 $U \to X$，其在 $X_0$ 上的限制还原为 $U_0$。由《态射补篇》
引理 \ref{more-morphisms-lemma-check-smoothness-on-infinitesimal-nbhds}
可知，$U \to X$ 在 $U_0$ 的每一点处都是 \'etale 的。然而，因为 $U$
的每一点都特化到 $U_0$ 的某一点（这是因为 $U$ 在 $A$ 上固有），
可知 $U \to X$ 是 \'etale 态射。由此，该函子本质满。

\medskip\noindent
当 $A$ 是完备 Noether 局部环时的全忠实性。设 $U \to X$ 和 $V \to X$
是有限 \'etale 态射，且 $\varphi_0 : U_0 \to V_0$ 是 $X_0$ 上的态射。
考虑态射
$$
\Gamma_{\varphi_0} : U_0 \longrightarrow U_0 \times_{X_0} V_0
$$
它既是有限 \'etale 态射，也是闭浸入。把本质满性应用于
$X = U \times_X V$，可得有限 \'etale 态射 $W \to U \times_X V$，
其特殊纤维同构于 $\Gamma_{\varphi_0}$。考虑投影 $W \to U$。依构造，
它是有限 \'etale 态射，并且在 $U_0$ 上是同构。由《\'Etale 态射》引理
\ref{etale-lemma-finite-etale-one-point}，$W \to U$ 在 $U_0$ 的一个开邻域
上是同构。因此它处处是同构，而复合 $\varphi : U \cong W \to V$
正是所需的 $\varphi_0$ 的提升。

\medskip\noindent
当 $A$ 是 Hensel Noether 局部 G-环时的本质满性。设 $U_0 \to X_0$
是有限 \'etale 态射。令 $A^\wedge$ 为 $A$ 关于极大理想的完备化，
并令 $X^\wedge$ 为 $X$ 到 $A^\wedge$ 的基变换。由上面的结果，存在
有限 \'etale 态射 $V \to X^\wedge$，其特殊纤维为 $U_0$。写成
$A^\wedge = \colim A_i$，其中 $A \to A_i$ 为有限型。由《极限》引理
\ref{limits-lemma-descend-finite-presentation}，存在某个 $i$ 以及有限表现态射
$U_i \to X_{A_i}$，其到 $X^\wedge$ 的基变换为 $V$。增大 $i$ 后，
可设 $U_i \to X_{A_i}$ 有限且 \'etale（《极限》引理
\ref{limits-lemma-descend-finite-finite-presentation} 和
\ref{limits-lemma-descend-etale}）。写
$$
A_i = A[x_1, \ldots, x_n]/(f_1, \ldots, f_m)
$$
则环同态 $A_i \to A^\wedge$ 可重新解释为一个解
$(a_1, \ldots, a_n)$，它是 $A^\wedge$ 中方程组 $f_j = 0$ 的解。
由《环同态的光滑化》定理
\ref{smoothing-theorem-approximation-property}，可把这个解逼近到例如
$11$ 阶；所得解为 $(b_1, \ldots, b_n)$，位于 $A$ 中。把它重新翻译成
环同态，便得到 $A$-代数同态 $A_i \to A$；它与原同态
$A_i \to A^\wedge$ 给出同一个闭点（因为 $11 > 1$）。因此，
$U \to X$，即沿这个环同态对 $V \to X_{A_i}$ 作的基变换，是有限 \'etale 态射，
其特殊纤维同构于 $U_0$。

\medskip\noindent
当 $A$ 是 Hensel Noether 局部 G-环时的全忠实性。它可由本质满性推出，
论证与 $A$ 是完备 Noether 环时完全相同。

\medskip\noindent
一般情形。设 $(A, \mathfrak m)$ 是 Hensel 局部环。置 $S = \Spec(A)$，
并以 $s \in S$ 表示闭点。由《极限》引理
\ref{limits-lemma-proper-limit-of-proper-finite-presentation-noetherian}，
可把 $X \to \Spec(A)$ 写成固有态射 $X_i \to S_i$ 的余滤过极限，其中
$S_i$ 在 $\mathbf{Z}$ 上为有限型。对每个 $i$，令 $s_i \in S_i$ 为 $s$
的像。由于 $S = \lim S_i$ 且 $A = \mathcal{O}_{S, s}$，有
$A = \colim \mathcal{O}_{S_i, s_i}$。环 $A_i = \mathcal{O}_{S_i, s_i}$
是 Noether 局部 G-环（《代数补篇》命题
\ref{more-algebra-proposition-ubiquity-G-ring}）。由《代数补篇》引理
\ref{more-algebra-lemma-henselization-colimit} 可知
$A = \colim A_i^h$。由《代数补篇》引理
\ref{more-algebra-lemma-henselization-G-ring}，各环 $A_i^h$ 都是 G-环。
因此 $A = \colim A_i^h$，并且作为概形有
$$
X = \lim (X_i \times_{S_i} \Spec(A_i^h))
$$
。$X$ 上的有限 \'etale 概形范畴，是
$X_i \times_{S_i} \Spec(A_i^h)$ 上的有限 \'etale 概形范畴的极限（由《极限》引理
\ref{limits-lemma-descend-finite-presentation}、
\ref{limits-lemma-descend-finite-finite-presentation} 和
\ref{limits-lemma-descend-etale})
。同样的结论也适用于
$X_0 = \lim (X_i \times_{S_i} s_i)$ 上的有限 \'etale 概形。因此，
$X / \Spec(A)$ 的结果可从上面已经处理过的
$(X_i \times_{S_i} \Spec(A_i^h)) / \Spec(A_i^h)$ 的结果形式地推出。
\end{proof}

\begin{lemma}
\label{pione-lemma-finite-etale-on-proper-over-henselian-pair}
设 $(A, I)$ 是 Hensel 对，$X$ 是 $A$ 上的固有概形。置
$X_0 = X \times_{\Spec(A)} \Spec(A/I)$。则函子
$$
\textit{F\'Et}_X \to \textit{F\'Et}_{X_0},\quad
U \longmapsto U_0 = U \times_X X_0
$$
是范畴等价。
\end{lemma}

\begin{proof}
本引理的证明与引理
\ref{pione-lemma-finite-etale-on-proper-over-henselian} 的证明完全相同。

\medskip\noindent
当 $A$ 是 Noether 环且 $I$-进完备时的本质满性。令
$X_n = X \times_{\Spec(A)} \Spec(A/I^{n + 1})$。由《\'Etale 态射》定理
\ref{etale-theorem-remarkable-equivalence}，嵌入
$$
X_0 \to X_1 \to X_2 \to \ldots
$$
在 $X_0$ 上的 \'etale 概形范畴与在 $X_n$ 上的 \'etale 概形范畴之间
诱导范畴等价。此外，若 $U_n \to X_n$ 对应于有限 \'etale 态射
$U_0 \to X_0$，则 $U_n \to X_n$ 也是有限的；例如，这可由《态射补篇》
引理 \ref{more-morphisms-lemma-thicken-property-morphisms-cartesian} 得到。
在这种情形下，态射 $U_0 \to \Spec(A/I)$ 是固有的，因为
$X_0$ 在 $A/I$ 上固有。因此可应用 Grothendieck 代数化定理
（采用《概形上同调》引理
\ref{coherent-lemma-algebraize-formal-scheme-finite-over-proper} 的形式），
得到有限态射 $U \to X$，其在 $X_0$ 上的限制还原为 $U_0$。由《态射补篇》
引理 \ref{more-morphisms-lemma-check-smoothness-on-infinitesimal-nbhds}
可知，$U \to X$ 在 $U_0$ 的每一点处都是 \'etale 的。然而，因为 $U$
的每一点都特化到 $U_0$ 的某一点（这是因为 $U$ 在 $A$ 上固有），
可知 $U \to X$ 是 \'etale 态射。由此，该函子本质满。

\medskip\noindent
当 $A$ 是 Noether 环且 $I$-进完备时的全忠实性。设 $U \to X$ 和
$V \to X$ 是有限 \'etale 态射，且 $\varphi_0 : U_0 \to V_0$ 是
$X_0$ 上的态射。考虑态射
$$
\Gamma_{\varphi_0} : U_0 \longrightarrow U_0 \times_{X_0} V_0
$$
它既是有限 \'etale 态射，也是闭浸入。把本质满性应用于
$X = U \times_X V$，可得有限 \'etale 态射 $W \to U \times_X V$，
其特殊纤维同构于 $\Gamma_{\varphi_0}$。考虑投影 $W \to U$。依构造，
它是有限 \'etale 态射，并且在 $U_0$ 上是同构。由《\'Etale 态射》引理
\ref{etale-lemma-finite-etale-one-point}，$W \to U$ 在 $U_0$ 的一个开邻域
上是同构。因此它处处是同构，而复合 $\varphi : U \cong W \to V$
正是所需的 $\varphi_0$ 的提升。

\medskip\noindent
当 $(A, I)$ 是 Hensel 对且 $A$ 是 Noether G-环时的本质满性。
设 $U_0 \to X_0$ 是有限 \'etale 态射。令 $A^\wedge$ 为 $A$ 关于 $I$
的完备化。注意，$A^\wedge$ 是 $IA^\wedge$-进完备的 Noether 环，
见《代数》引理 \ref{algebra-lemma-completion-complete} 和
\ref{algebra-lemma-completion-Noetherian-Noetherian}。令 $X^\wedge$
为 $X$ 到 $A^\wedge$ 的基变换。由上面的结果，存在有限 \'etale 态射
$V \to X^\wedge$，其特殊纤维为 $U_0$。写成
$A^\wedge = \colim A_i$，其中 $A \to A_i$ 为有限型。由《极限》引理
\ref{limits-lemma-descend-finite-presentation}，存在某个 $i$ 以及有限表现态射
$U_i \to X_{A_i}$，其到 $X^\wedge$ 的基变换为 $V$。增大 $i$ 后，
可设 $U_i \to X_{A_i}$ 有限且 \'etale（《极限》引理
\ref{limits-lemma-descend-finite-finite-presentation} 和
\ref{limits-lemma-descend-etale}）。写
$$
A_i = A[x_1, \ldots, x_n]/(f_1, \ldots, f_m)
$$
则环同态 $A_i \to A^\wedge$ 可重新解释为一个解
$(a_1, \ldots, a_n)$，它是 $A^\wedge$ 中方程组 $f_j = 0$ 的解。
由《环同态的光滑化》引理
\ref{smoothing-lemma-henselian-pair}，可把这个解逼近到例如 $11$ 阶；
所得解为 $(b_1, \ldots, b_n)$，位于 $A$ 中。把它重新翻译成
环同态，便得到 $A$-代数同态 $A_i \to A$；它与原同态
$A_i \to A^\wedge$ 给出同一个闭点（因为 $11 > 1$）。因此，
$U \to X$，即沿这个环同态对 $V \to X_{A_i}$ 作的基变换，是有限 \'etale 态射，
其特殊纤维同构于 $U_0$。

\medskip\noindent
当 $(A, I$ 是 Hensel 对且 $A$ 是 Noether G-环时的全忠实性。它可由
本质满性推出，论证与 $A$ 是完备 Noether 环时完全相同。

\medskip\noindent
一般情形。设 $(A, I)$ 是 Hensel 对。置 $S = \Spec(A)$，并记
$S_0 = \Spec(A/I)$。由《极限》引理
\ref{limits-lemma-proper-limit-of-proper-finite-presentation-noetherian}，
可把 $X \to \Spec(A)$ 写成固有态射 $X_i \to S_i$ 的余滤过极限，其中
$S_i$ 仿射且在 $\mathbf{Z}$ 上为有限型。写 $S_i = \Spec(A_i)$，并以
$I_i \subset A_i$ 表示理想 $I$ 沿同态 $A_i \to A$ 的逆像。置
$S_{i, 0} = \Spec(A_i/I_i)$。由于 $S = \lim S_i$，有
$A = \colim A_i$。因而还有 $I = \colim I_i$ 及 $A/I = \colim A_i/I_i$。
环 $A_i$ 是 Noether G-环（《代数补篇》命题
\ref{more-algebra-proposition-ubiquity-G-ring}）。以 $(A_i^h, I_i^h)$
表示对 $(A_i, I_i)$ 的 Hensel 化。由《代数补篇》引理
\ref{more-algebra-lemma-henselization-colimit} 可知
$A = \colim A_i^h$。由《代数补篇》引理
\ref{more-algebra-lemma-henselization-pair-G-ring}，各环 $A_i^h$ 都是
G-环。因此 $A = \colim A_i^h$，并且作为概形有
$$
X = \lim (X_i \times_{S_i} \Spec(A_i^h))
$$
。$X$ 上的有限 \'etale 概形范畴，是
$X_i \times_{S_i} \Spec(A_i^h)$ 上的有限 \'etale 概形范畴的极限（由《极限》引理
\ref{limits-lemma-descend-finite-presentation}、
\ref{limits-lemma-descend-finite-finite-presentation} 和
\ref{limits-lemma-descend-etale})
。同样的结论也适用于
$X_0 = \lim (X_i \times_{S_i} S_{i, 0})$ 上的有限 \'etale 概形。因此，
$X / \Spec(A)$ 的结果可从上面已经处理过的
$(X_i \times_{S_i} \Spec(A_i^h)) / \Spec(A_i^h)$ 的结果形式地推出。
\end{proof}

\begin{lemma}
\label{pione-lemma-finite-etale-invariant-over-proper}
设 $k'/k$ 是代数闭域的扩张，$X$ 是 $k$ 上的固有概形。则函子
$$
U \longmapsto U_{k'}
$$
在 $X$ 上的有限 \'etale 概形范畴与在 $X_{k'}$ 上的有限 \'etale
概形范畴之间给出范畴等价。
\end{lemma}

\begin{proof}
我们来证明该函子本质满。设 $U' \to X_{k'}$ 是有限 \'etale 态射。
把 $k' = \colim A_i$ 写成有限型 $k$-代数的滤过余极限。由《极限》引理
\ref{limits-lemma-descend-finite-presentation}，存在某个 $i$ 以及有限表现态射
$U_i \to X_{A_i}$，其到 $X_{k'}$ 的基变换为 $U'$。增大 $i$ 后，
可设 $U_i \to X_{A_i}$ 有限且 \'etale（《极限》引理
\ref{limits-lemma-descend-finite-finite-presentation} 和
\ref{limits-lemma-descend-etale}）。
由于 $k$ 代数闭，可找到一个 $k$-值点 $t$，它位于 $\Spec(A_i)$ 中。令
$U = (U_i)_t$ 为 $U_i$ 在 $t$ 上的纤维。令 $A_i^h$ 为
$(A_i)_{\mathfrak m}$ 的 Hensel 化，其中 $\mathfrak m$ 是与点 $t$
对应的极大理想。由引理
\ref{pione-lemma-finite-etale-on-proper-over-henselian} 可知，
$(U_i)_{A_i^h} = U \times \Spec(A_i^h)$；这是 $X_{A_i^h}$ 上的概形同构。
现在，由于 $A_i^h$ 在 $A_i$
上是代数的（例如见《环同态的光滑化》例
\ref{smoothing-example-describe-henselian} 中的讨论），而且 $k'$ 代数闭，
可以找到环同态 $A_i^h \to k'$，它扩张给定的包含
$A_i \subset k'$。因此 $U'$ 同构于 $U$ 的基变换。
全忠实性的证明完全相同。
\end{proof}








\section{局部连通性}
\label{pione-section-unibranch}

\noindent
本节研究映射 $\pi_1(U) \to \pi_1(X)$ 何时为满射，其中 $U$ 是概形
$X$ 的稠密开子概形。我们将看到，粗略地说，只要 $U \cap B$ 连通，
其中 $B$ 是点 $x \in X \setminus U$ 附近任意小的“球”，上述映射就是满射。

\begin{lemma}
\label{pione-lemma-dense-faithful}
设 $f : X \to Y$ 是概形态射。若 $f(X)$ 在 $Y$ 中稠密，则基变换函子
$\textit{F\'Et}_Y \to \textit{F\'Et}_X$ 忠实。
\end{lemma}

\begin{proof}
因为有限 \'etale 覆盖的范畴具有内部 Hom（引理
\ref{pione-lemma-internal-hom-finite-etale}），只需证明以下命题：给定 $W$
（它在 $Y$ 上有限 \'etale），以及态射 $s : X \to W$（在 $X$ 上），
截面 $t : Y \to W$ 若满足 $s = t \circ f$，则至多有一个。设有两个截面
$t_1, t_2 : Y \to W$，使得 $s = t_1 \circ f = t_2 \circ f$。
$t_1$ 与 $t_2$ 的等化子在 $Y$ 中闭（《概形》引理
\ref{schemes-lemma-where-are-they-equal}），而 $f(X)$ 在 $Y$ 中稠密，
故 $t_1$ 与 $t_2$ 在 $Y_{red}$ 上相同。于是 $t_1$ 与 $t_2$ 具有相同的像；
该像是 $W$ 的开闭子概形，并同构地映到 $Y$（《\'Etale 态射》命题
\ref{etale-proposition-properties-sections}）。因此二者相等。
\end{proof}

\noindent
下一个引理假设严格 Hensel 化的穿孔谱连通。例如，局部环几何单枝时，
这一条件成立，见《代数补篇》引理
\ref{more-algebra-lemma-geometrically-unibranch}。《性质》引理
\ref{properties-lemma-geometrically-unibranch} 给出了一个部分逆命题。

\begin{lemma}
\label{pione-lemma-same-etale-extensions}
设 $(A, \mathfrak m)$ 是局部环。置 $X = \Spec(A)$，并令
$U = X \setminus \{\mathfrak m\}$。若 $A$ 的严格 Hensel 化的穿孔谱连通，则
$$
\textit{F\'Et}_X \longrightarrow \textit{F\'Et}_U,\quad
Y \longmapsto Y \times_X U
$$
是全忠实函子。
\end{lemma}

\begin{proof}
先设 $A$ 严格 Hensel。在这种情形下，有限 \'etale 覆盖 $Y$（在 $X$
上）都同构于有限多个 $X$ 的不交并。因此，只需证明：态射
$U \to U \amalg \ldots \amalg U$（在 $U$ 上）唯一延拓为态射
$X \to X \amalg \ldots \amalg X$（在 $X$ 上）。若 $U$ 连通（特别地，非空），
这显然成立。

\medskip\noindent
一般情形。因为有限 \'etale 覆盖的范畴具有内部 Hom（引理
\ref{pione-lemma-internal-hom-finite-etale}），只需证明：给定 $Y$（它在 $X$ 上有限
\'etale），任意态射 $s : U \to Y$（在 $X$ 上）都延拓为态射
$t : X \to Y$（在 $X$ 上）。令 $A^{sh}$ 为 $A$ 的严格 Hensel 化，并记
$X^{sh} = \Spec(A^{sh})$、$U^{sh} = U \times_X X^{sh}$、
$Y^{sh} = Y \times_X X^{sh}$。由第一段及对 $A$ 的假设，可把基变换
$s^{sh} : U^{sh} \to Y^{sh}$（它来自 $s$）延拓为
$t^{sh} : X^{sh} \to Y^{sh}$。置 $A' = A^{sh} \otimes_A A^{sh}$。
于是，两个拉回 $t'_1, t'_2$（来自 $t^{sh}$，拉到 $X' = \Spec(A')$）
都是拉回 $s'$（来自 $s$，拉到 $U' = U \times_X X'$）的延拓。由于
$A \to A'$ 平坦，$U' \subset X'$ 在拓扑意义下稠密；这是由
$A \to A'$ 的降链性质（《代数》引理
\ref{algebra-lemma-flat-going-down}）得到的。因此，由引理
\ref{pione-lemma-dense-faithful}，$t'_1 = t'_2$。
所以 $t^{sh}$ 下降为态射 $t : X \to Y$；例如，这可由《下降》引理
\ref{descent-lemma-fpqc-universal-effective-epimorphisms} 得到。
\end{proof}

\noindent
鉴于引理 \ref{pione-lemma-same-etale-extensions}，自然要问：一个环（及其严格
Hensel 化）的穿孔谱何时连通。Hartshorne 的一个著名引理给出了充分条件，
见《局部上同调》引理
\ref{local-cohomology-lemma-depth-2-connected-punctured-spectrum}。

\begin{lemma}
\label{pione-lemma-quasi-compact-dense-open-connected-at-infinity-Noetherian}
设 $X$ 是概形，$U \subset X$ 是稠密开子概形。假设
\begin{enumerate}
\item $X$ 的底拓扑空间是 Noether 空间；
\item 对每个 $x \in X \setminus U$，$\mathcal{O}_{X, x}$ 的严格 Hensel
化的穿孔谱都连通。
\end{enumerate}
则 $\textit{F\'Et}_X \to \textit{F\'et}_U$ 全忠实。
\end{lemma}

\begin{proof}
设 $Y_1, Y_2$ 在 $X$ 上有限 \'etale，并设
$\varphi : (Y_1)_U \to (Y_2)_U$ 是 $U$ 上的态射。须证明 $\varphi$
唯一提升为态射 $Y_1 \to Y_2$（在 $X$ 上）。唯一性由引理
\ref{pione-lemma-dense-faithful} 得出。

\medskip\noindent
设 $x \in X \setminus U$ 是 $X \setminus U$ 的一个不可约分支的泛点。置
$V = U \times_X \Spec(\mathcal{O}_{X, x})$。由 $x$ 的选取方式，它就是
$\Spec(\mathcal{O}_{X, x})$ 的穿孔谱。由引理
\ref{pione-lemma-same-etale-extensions}，态射
$\varphi_V : (Y_1)_V \to (Y_2)_V$ 唯一延拓为态射
$(Y_1)_{\Spec(\mathcal{O}_{X, x})} \to (Y_2)_{\Spec(\mathcal{O}_{X, x})}$。
由《极限》引理 \ref{limits-lemma-glueing-near-point}，可找到开子概形
$U \subset U'$（它包含 $x$），以及态射
$\varphi' : (Y_1)_{U'} \to (Y_2)_{U'}$，它延拓 $\varphi$。由于 $X$ 的底拓扑空间是
Noether 空间，对 $\varphi$ 已定义的开子概形的补集作 Noether 归纳，
便完成证明。
\end{proof}

\begin{lemma}
\label{pione-lemma-retrocompact-dense-open-connected-at-infinity-closed}
设 $X$ 是概形，$U \subset X$ 是稠密开子概形。假设
\begin{enumerate}
\item $U \to X$ 拟紧；
\item $X \setminus U$ 的每一点都是闭点；
\item 对每个 $x \in X \setminus U$，$\mathcal{O}_{X, x}$ 的严格 Hensel
化的穿孔谱都连通。
\end{enumerate}
则 $\textit{F\'Et}_X \to \textit{F\'et}_U$ 全忠实。
\end{lemma}

\begin{proof}
设 $Y_1, Y_2$ 在 $X$ 上有限 \'etale，并设
$\varphi : (Y_1)_U \to (Y_2)_U$ 是 $U$ 上的态射。须证明 $\varphi$
唯一提升为态射 $Y_1 \to Y_2$（在 $X$ 上）。唯一性由引理
\ref{pione-lemma-dense-faithful} 得出。

\medskip\noindent
设 $x \in X \setminus U$。置 $V = U \times_X \Spec(\mathcal{O}_{X, x})$。
由于 $X \setminus U$ 的每一点都闭，$V$ 就是
$\Spec(\mathcal{O}_{X, x})$ 的穿孔谱。由引理
\ref{pione-lemma-same-etale-extensions}，态射
$\varphi_V : (Y_1)_V \to (Y_2)_V$ 唯一延拓为态射
$(Y_1)_{\Spec(\mathcal{O}_{X, x})} \to (Y_2)_{\Spec(\mathcal{O}_{X, x})}$。
由《极限》引理 \ref{limits-lemma-glueing-near-point}（这里用到 $U$
在 $X$ 中逆紧），可找到开子概形 $U \subset U'_x$（它包含 $x$），
以及态射 $\varphi'_x : (Y_1)_{U'_x} \to (Y_2)_{U'_x}$，它延拓
$\varphi$。注意，对任意两点
$x, x' \in X \setminus U$，态射 $\varphi'_x$ 与 $\varphi'_{x'}$
在 $U'_x \cap U'_{x'}$ 上相同，因为 $U$ 在该开子概形中稠密（引理
\ref{pione-lemma-dense-faithful}）。因此可按所需把 $\varphi$ 延拓到
$\bigcup U'_x = X$。
\end{proof}

\begin{lemma}
\label{pione-lemma-quasi-compact-dense-open-connected-at-infinity}
设 $X$ 是概形，$U \subset X$ 是稠密开子概形。假设
\begin{enumerate}
\item $X$ 的每个拟紧开子概形只有有限多个不可约分支；
\item 对每个 $x \in X \setminus U$，$\mathcal{O}_{X, x}$ 的严格 Hensel
化的穿孔谱都连通。
\end{enumerate}
则 $\textit{F\'Et}_X \to \textit{F\'et}_U$ 全忠实。
\end{lemma}

\begin{proof}
设 $Y_1, Y_2$ 在 $X$ 上有限 \'etale，并设
$\varphi : (Y_1)_U \to (Y_2)_U$ 是 $U$ 上的态射。须证明 $\varphi$
唯一提升为态射 $Y_1 \to Y_2$（在 $X$ 上）。唯一性由引理
\ref{pione-lemma-dense-faithful} 得出。为证明存在性，我们将说明可扩大 $U$
（若 $U \not = X$），再用 Zorn 引理完成证明。

\medskip\noindent
设 $x \in X \setminus U$ 是 $X \setminus U$ 的一个不可约分支的泛点。置
$V = U \times_X \Spec(\mathcal{O}_{X, x})$。由 $x$ 的选取方式，它就是
$\Spec(\mathcal{O}_{X, x})$ 的穿孔谱。由引理
\ref{pione-lemma-same-etale-extensions}，态射
$\varphi_V : (Y_1)_V \to (Y_2)_V$（唯一）延拓为态射
$(Y_1)_{\Spec(\mathcal{O}_{X, x})} \to (Y_2)_{\Spec(\mathcal{O}_{X, x})}$。
选取仿射邻域 $W \subset X$（它包含 $x$）。因为 $U \cap W$ 在 $W$ 中稠密，
它包含各泛点 $\eta_1, \ldots, \eta_n$，即 $W$ 的各泛点。选取仿射开子概形
$W' \subset W \cap U$，使其包含 $\eta_1, \ldots, \eta_n$。置
$V' = W' \times_X \Spec(\mathcal{O}_{X, x})$。把《极限》引理
\ref{limits-lemma-glueing-near-point} 应用于 $x \in W \supset W'$，
可找到开子概形 $W' \subset W'' \subset W$，其中 $x \in W''$，
以及态射 $\varphi'' : (Y_1)_{W''} \to (Y_2)_{W''}$，它与
$\varphi$ 在 $W'$ 上相同。由于 $W'$ 在 $W'' \cap U$ 中稠密，由引理
\ref{pione-lemma-dense-faithful} 可知，$\varphi$ 与 $\varphi''$ 在
$U \cap W'$ 上相同。因此 $\varphi$ 与 $\varphi''$ 粘合为态射
$\varphi'$，它定义在 $U' = U \cup W''$ 上，并与 $\varphi$ 在 $U$
相同。注意 $x \in U'$，所以我们已把 $\varphi$ 延拓到一个严格更大的开子概形。

\medskip\noindent
考虑集合 $\mathcal{S}$，其元素是二元组 $(U', \varphi')$，满足
$U \subset U'$，且 $\varphi'$ 是 $\varphi$ 的延拓。按显然的方式在
$\mathcal{S}$ 上赋予偏序。若 $(U'_i, \varphi'_i)$ 是全序子集，则它有
最大元 $(U', \varphi')$：只需取 $U' = \bigcup U'_i$，并令
$\varphi' : (Y_1)_{U'} \to (Y_2)_{U'}$ 为与 $\varphi'_i$ 在
$U'_i$ 上相同的态射。因此可应用 Zorn 引理，$\mathcal{S}$ 有极大元。
由上面的论证，这个极大元就是 $\varphi$ 在整个 $X$ 上的延拓。
\end{proof}

\begin{lemma}
\label{pione-lemma-local-exact-sequence}
设 $(A, \mathfrak m)$ 是局部环。置 $X = \Spec(A)$，并令
$U = X \setminus \{\mathfrak m\}$。令 $U^{sh}$ 为严格 Hensel 化
$A^{sh}$（即 $A$ 的严格 Hensel 化）的穿孔谱。假设 $U$ 拟紧且 $U^{sh}$
连通。则序列
$$
\pi_1(U^{sh}, \overline{u}) \to \pi_1(U, \overline{u}) \to
\pi_1(X, \overline{u}) \to 1
$$
按引理 \ref{pione-lemma-functoriality-galois-ses} 第 (1) 部分的意义是正合的。
\end{lemma}

\begin{proof}
映射 $\pi_1(U) \to \pi_1(X)$ 由引理
\ref{pione-lemma-same-etale-extensions} 和
\ref{pione-lemma-functoriality-galois-surjective} 可知是满射。

\medskip\noindent
记 $X^{sh} = \Spec(A^{sh})$。设 $Y \to X$ 是有限 \'etale 态射。则
$Y^{sh} = Y \times_X X^{sh} \to X^{sh}$ 是有限 \'etale 态射。由于
$A^{sh}$ 严格 Hensel，$Y^{sh}$ 同构于有限多个 $X^{sh}$ 的不交并。因此
$Y \times_X U^{sh}$ 亦然。于是复合
$\pi_1(U^{sh}) \to \pi_1(U) \to \pi_1(X)$ 平凡，见引理
\ref{pione-lemma-composition-trivial}。

\medskip\noindent
为完成证明，按照引理 \ref{pione-lemma-functoriality-galois-ses}，只需证明：
给定有限 \'etale 态射 $V \to U$，若 $V \times_U U^{sh}$ 是若干个
$U^{sh}$ 的不交并，则可以找到有限 \'etale 态射 $Y \to X$，使得
$V \cong Y \times_X U$ 在 $U$ 上成立。该假设蕴含存在有限 \'etale 态射
$Y^{sh} \to X^{sh}$ 以及同构
$V \times_U U^{sh} \cong Y^{sh} \times_{X^{sh}} U^{sh}$。考虑图表
$$
\xymatrix{
U \ar[d] & U^{sh} \ar[d] \ar[l] &
U^{sh} \times_U U^{sh} \ar[d] \ar@<1ex>[l] \ar@<-1ex>[l] &
U^{sh} \times_U U^{sh} \times_U U^{sh}
\ar[d] \ar@<1ex>[l] \ar[l] \ar@<-1ex>[l] \\
X & X^{sh} \ar[l] &
X^{sh} \times_X X^{sh} \ar@<1ex>[l] \ar@<-1ex>[l] &
X^{sh} \times_X X^{sh} \times_X X^{sh} \ar@<1ex>[l] \ar[l] \ar@<-1ex>[l]
}
$$
由于按假设 $U \subset X$ 拟紧，所有向下的箭头都是拟紧开浸入。设
$\xi \in X^{sh} \times_X X^{sh}$ 是不属于
$U^{sh} \times_U U^{sh}$ 的点。则 $\xi$ 位于闭点 $x^{sh}$（在
$X^{sh}$ 中）上方。考虑局部环同态
$$
A^{sh} = \mathcal{O}_{X^{sh}, x^{sh}} \to
\mathcal{O}_{X^{sh} \times_X X^{sh}, \xi}
$$
它由第一投影 $X^{sh} \times_X X^{sh}$ 决定。这个同态是一些局部同态的
滤过余极限，而这些局部同态由 \'etale 环映射局部化而得。由于 $A^{sh}$
严格 Hensel，可知该同态是同构。由于对补集中的每个 $\xi$ 都是如此，
这些点之间不存在特殊化关系，因而每个这样的 $\xi$ 都是闭点（这也可以
直接证明）。点 $\xi$ 处的局部环同构于 $A^{sh}$，故它严格 Hensel，且其
穿孔谱连通。对于点 $\xi$（它属于
$X^{sh} \times_X X^{sh} \times_X X^{sh}$ 而不属于
$U^{sh} \times_U U^{sh} \times_U U^{sh}$），同理可得这一结论。
由引理 \ref{pione-lemma-retrocompact-dense-open-connected-at-infinity-closed}，沿各
竖直箭头的拉回在有限 \'etale 概形范畴上诱导全忠实函子。因此，
$V \times_U U^{sh}$ 相对于 fpqc 覆盖 $\{U^{sh} \to U\}$ 的典范下降数据，
转化为 $Y^{sh}$ 相对于 fpqc 覆盖 $\{X^{sh} \to X\}$ 的下降数据。由于
$Y^{sh} \to X^{sh}$ 有限，因而仿射，该下降数据有效（《下降》引理
\ref{descent-lemma-affine}）。于是得到仿射态射 $Y \to X$，以及与下降数据
相容的同构 $Y \times_X X^{sh} \to Y^{sh}$。由下降数据的全忠实性
（如《下降》引理 \ref{descent-lemma-refine-coverings-fully-faithful}），得到
同构 $V \to U \times_X Y$。最后，$Y \to X$ 有限 \'etale，因为
$Y^{sh} \to X^{sh}$ 如此；见《下降》引理
\ref{descent-lemma-descending-property-etale} 和
\ref{descent-lemma-descending-property-finite}。
\end{proof}

\noindent
设 $X$ 是不可约概形，$\eta \in X$ 是其泛点。典范态射 $\eta \to X$
诱导典范映射
\begin{equation}
\label{pione-equation-inclusion-generic-point}
\text{Gal}(\kappa(\eta)^{sep}/\kappa(\eta)) = \pi_1(\eta, \overline{\eta})
\longrightarrow \pi_1(X, \overline{\eta})
\end{equation}
左侧的等同来自引理 \ref{pione-lemma-fundamental-group-Galois-group}。

\begin{lemma}
\label{pione-lemma-irreducible-geometrically-unibranch}
设 $X$ 是不可约几何单枝概形。对任意非空开子概形 $U \subset X$，典范映射
$$
\pi_1(U, \overline{u}) \longrightarrow \pi_1(X, \overline{u})
$$
是满射。映射 (\ref{pione-equation-inclusion-generic-point})
$\pi_1(\eta, \overline{\eta}) \to \pi_1(X, \overline{\eta})$
也是满射。
\end{lemma}

\begin{proof}
由引理 \ref{pione-lemma-thickening}，可以作约化后代替 $X$。所以可设
$X$ 是整概形。由引理 \ref{pione-lemma-functoriality-galois-surjective}，本引理的
断言等价于函子 $\textit{F\'Et}_X \to \textit{F\'Et}_U$ 和
$\textit{F\'Et}_X \to \textit{F\'Et}_\eta$ 全忠实。

\medskip\noindent
关于 $\textit{F\'Et}_X \to \textit{F\'Et}_U$ 的结论来自引理
\ref{pione-lemma-quasi-compact-dense-open-connected-at-infinity}，以及如下事实：若
局部环 $A$ 几何单枝，则其严格 Hensel 化的谱不可约。见《代数补篇》引理
\ref{more-algebra-lemma-geometrically-unibranch}。

\medskip\noindent
注意，剩余域 $\kappa(\eta) = \mathcal{O}_{X, \eta}$ 是
$\mathcal{O}_X(U)$ 的滤过余极限，其中 $U \subset X$ 遍历非空仿射开子概形。
因此，$\textit{F\'Et}_\eta$ 是相应范畴 $\textit{F\'Et}_U$ 的余极限，其中
$U$ 如上；见《极限》引理 \ref{limits-lemma-descend-finite-presentation}、
\ref{limits-lemma-descend-finite-finite-presentation} 和
\ref{limits-lemma-descend-etale}。
再作形式论证，便可知函子 $\textit{F\'Et}_X \to \textit{F\'Et}_\eta$ 的
全忠实性可由各函子 $\textit{F\'Et}_X \to \textit{F\'Et}_U$ 的全忠实性推出。
\end{proof}

\begin{lemma}
\label{pione-lemma-exact-sequence-finite-nr-closed-pts}
设 $X$ 是概形。设 $x_1, \ldots, x_n \in X$ 是有限多个闭点，满足
\begin{enumerate}
\item $U = X \setminus \{x_1, \ldots, x_n\}$ 连通，且是 $X$ 的逆紧开子概形；
\item 对每个 $i$，穿孔谱 $U_i^{sh}$（它属于 $\mathcal{O}_{X, x_i}$ 的
严格 Hensel 化）连通。
\end{enumerate}
则映射 $\pi_1(U) \to \pi_1(X)$ 是满射，且其核是 $\pi_1(U)$ 的最小闭正规
子群，它包含各映射 $\pi_1(U_i^{sh}) \to \pi_1(U)$ 的像，其中
$i = 1, \ldots, n$。
\end{lemma}

\begin{proof}
满射性由引理 \ref{pione-lemma-retrocompact-dense-open-connected-at-infinity-closed} 和
\ref{pione-lemma-functoriality-galois-surjective} 得出。可以考虑映射序列
$$
\pi_1(U)  \to \ldots \to
\pi_1(X \setminus \{x_1, x_2\}) \to \pi_1(X \setminus \{x_1\}) \to \pi_1(X)
$$
群论论证表明，只需在 $n = 1$ 时证明关于核的断言（略去细节）。记
$x = x_1$、$U^{sh} = U_1^{sh}$，置 $A = \mathcal{O}_{X, x}$，并令
$A^{sh}$ 为其严格 Hensel 化。考虑图表
$$
\xymatrix{
U \ar[d] &
\Spec(A) \setminus \{\mathfrak m\} \ar[l] \ar[d] &
U^{sh} \ar[d] \ar[l] \\
X & \Spec(A) \ar[l] & \Spec(A^{sh}) \ar[l]
}
$$
由引理 \ref{pione-lemma-functoriality-galois-ses}，须证明：若有限 \'etale 态射
$V \to U$ 拉回到 $U^{sh}$ 后成为平凡覆盖，则它延拓为 $X$ 上的有限
\'etale 概形。由引理 \ref{pione-lemma-local-exact-sequence}，我们已知 $A$ 的
穿孔谱上有限 \'etale 概形的相应断言。另一方面，由《极限》引理
\ref{limits-lemma-glueing-near-closed-point}，$X$ 上有限表示概形等同于把
$U$ 上和 $A$ 上的有限表示概形沿 $A$ 的穿孔谱粘合而得的概形。证明完毕。
\end{proof}









\section{正规概形的基本群}
\label{pione-section-normal}

\noindent
设 $X$ 是整的几何单枝概形。上一节已经看到，$X$ 的基本群是函数域
$K$（即 $X$ 的函数域）的 Galois 群之商。由于该映射连续，其核是 Galois 群的闭正规
子群。因此，由 Galois 理论，该核对应于 Galois 扩张 $M/K$（《域》定理
\ref{fields-theorem-inifinite-galois-theory}）。本节将确定 $M$；这里假设
$X$ 是正规整概形。

\medskip\noindent
设 $X$ 是函数域为 $K$ 的正规整概形。设 $L/K$ 是有限扩张。考虑正规化
$Y \to X$，即 $X$ 在态射 $\Spec(L) \to X$ 中的正规化，其定义见《态射》
第 \ref{morphisms-section-normalization-X-in-Y} 节。在此情形下，称
{\it $X$ 在 $L$ 中不分歧}，如果 $Y \to X$ 是概形的非分歧态射。我们将在
引理 \ref{pione-lemma-unramified} 中阐明这一条件。注意，$Y \to X$ 在
$\Spec(K)$ 上的概形论纤维是 $\Spec(L)$。因此，域扩张 $L/K$ 是可分的，
只要 $X$ 在 $L$ 中不分歧；见《态射》引理
\ref{morphisms-lemma-unramified-over-field}。

\begin{lemma}
\label{pione-lemma-unramified-in-L}
在上述情形下，下列条件等价：
\begin{enumerate}
\item $X$ 在 $L$ 中不分歧；
\item $Y \to X$ 是 \'etale 态射；
\item $Y \to X$ 是有限 \'etale 态射。
\end{enumerate}
\end{lemma}

\begin{proof}
注意，$Y \to X$ 是整态射。在每种情形下，按照定义，态射 $Y \to X$
局部有限型。因此，由《态射》引理 \ref{morphisms-lemma-finite-integral}，
在每种情形下 $Y \to X$ 都有限。特别地，(2) 与 (3) 等价。
\'Etale 态射非分歧，故 (2) 蕴含 (1)。

\medskip\noindent
反之，假设 $Y \to X$ 非分歧。由于正规概形几何单枝（《性质》引理
\ref{properties-lemma-normal-geometrically-unibranch}），由《态射补篇》引理
\ref{more-morphisms-lemma-global-unramified-dominant-unibranch-is-etale} 可知
态射 $Y \to X$ 是 \'etale 态射。下一段还将给出直接证明。

\medskip\noindent
设 $x \in X$。可以选取 \'etale 邻域 $(U, u) \to (X, x)$，使得
$$
Y \times_X U = \coprod V_j \longrightarrow U
$$
是若干闭浸入的不交并；见《\'Etale 态射》引理
\ref{etale-lemma-finite-unramified-etale-local}。通过缩小，可以假设 $U$
拟紧。于是 $U$ 只有有限多个不可约分支（《下降》引理
\ref{descent-lemma-locally-finite-nr-irred-local-fppf}）。由于 $U$ 正规
（《下降》引理 \ref{descent-lemma-normal-local-smooth}），$U$ 的不可约分支
都是开闭的（《性质》引理
\ref{properties-lemma-normal-locally-finite-nr-irreducibles}），故可以假设
$U$ 不可约。于是 $U$ 是整概形，其泛点 $\xi$ 映到 $X$ 的泛点。另一方面，
已知 $Y \times_X U$ 是 $U$ 在 $\Spec(L) \times_X U$ 中的正规化；见
《态射补篇》引理
\ref{more-morphisms-lemma-normalization-smooth-localization}。
$\Spec(L) \times_X U$ 的每一点都映到 $\xi$。因此，由《态射》引理
\ref{morphisms-lemma-normalization-generic}，每个 $V_j$ 都包含一个映到
$\xi$ 的点。态射 $V_j \to U$ 是同构，因为 $U = \overline{\{\xi\}}$。
所以 $Y \times_X U \to U$ 是 \'etale 态射。由《下降》引理
\ref{descent-lemma-descending-property-etale}，可知 $Y \to X$ 在
$U \to X$ 的像（即 $x$ 的一个开邻域）上是 \'etale 态射。
\end{proof}

\begin{lemma}
\label{pione-lemma-finite-etale-covering-normal-unramified}
设 $X$ 是函数域为 $K$ 的正规整概形。设 $Y \to X$ 是有限 \'etale 态射。
若 $Y$ 连通，则 $Y$ 是正规整概形，且 $Y$ 是 $X$ 在 $Y$ 的函数域中的
正规化。
\end{lemma}

\begin{proof}
由《下降》引理 \ref{descent-lemma-normal-local-smooth}，概形 $Y$ 正规。
由于 $Y \to X$ 平坦，由《态射》引理
\ref{morphisms-lemma-generalizations-lift-flat}，$Y$ 的每个泛点都映到 $X$
的泛点。由于 $Y \to X$ 有限，$Y$ 只有有限多个不可约分支。因此，由
《性质》引理 \ref{properties-lemma-normal-locally-finite-nr-irreducibles}，
$Y$ 是有限多个正规整概形的不交并。所以，若 $Y$ 连通，则 $Y$ 是正规
整概形。

\medskip\noindent
设 $L$ 是 $Y$ 的函数域，并令 $Y' \to X$ 为 $X$ 在 $L$ 中的正规化。
由《态射》引理 \ref{morphisms-lemma-characterize-normalization}，得到分解
$Y' \to Y \to X$，并且 $Y' \to Y$ 是 $Y$ 在 $L$ 中的正规化。由于 $Y$
正规，显然 $Y' = Y$（这也可由《态射》引理
\ref{morphisms-lemma-finite-birational-over-normal} 得出）。
\end{proof}

\begin{proposition}
\label{pione-proposition-normal}
设 $X$ 是函数域为 $K$ 的正规整概形。则典范映射
(\ref{pione-equation-inclusion-generic-point})
$$
\text{Gal}(K^{sep}/K) = \pi_1(\eta, \overline{\eta})
\longrightarrow \pi_1(X, \overline{\eta})
$$
等同于商映射 $\text{Gal}(K^{sep}/K) \to \text{Gal}(M/K)$，其中
$M \subset K^{sep}$ 是有限子扩张 $L$ 之并，其中 $X$ 在 $L$ 中不分歧。
\end{proposition}

\begin{proof}
正规概形 $X$ 几何单枝（《性质》引理
\ref{properties-lemma-normal-geometrically-unibranch}）。因此，引理
\ref{pione-lemma-irreducible-geometrically-unibranch} 可应用于 $X$。所以
$\pi_1(\eta, \overline{\eta}) \to \pi_1(X, \overline{\eta})$ 是满射，且
交换图表的上方水平箭头
$$
\xymatrix{
\textit{F\'Et}_X \ar[r] \ar[d] \ar[rd]_c & \textit{F\'Et}_\eta \ar[d] \\
\textit{Finite-}\pi_1(X, \overline{\eta})\textit{-sets} \ar[r] &
\textit{Finite-}\text{Gal}(K^{sep}/K)\textit{-sets}
}
$$
全忠实。左侧竖直箭头是定理 \ref{pione-theorem-fundamental-group} 的等价，
右侧竖直箭头是引理 \ref{pione-lemma-fundamental-group-Galois-group} 的等价。
下方水平箭头由本命题中的映射诱导。由引理
\ref{pione-lemma-unramified-in-L} 和
\ref{pione-lemma-finite-etale-covering-normal-unramified}，可知 $c$ 的本质像由
同构于下列形式集合的 $\text{Gal}(K^{sep}/K)\textit{-Sets}$ 构成：
$$
S = \Hom_K(\prod\nolimits_{i = 1, \ldots, n} L_i, K^{sep}) =
\coprod\nolimits_{i = 1, \ldots, n} \Hom_K(L_i, K^{sep})
$$
其中 $L_i/K$ 是有限可分扩张，且 $X$ 在 $L_i$ 中不分歧。因此，若
$M \subset K^{sep}$ 如该引理的陈述所述，则 $\text{Gal}(K^{sep}/M)$
恰为 $\text{Gal}(K^{sep}/K)$ 中在 $c$ 的本质像的每个对象上平凡作用的
子群。另一方面，$c$ 的本质像恰为这样的 $S$ 所成的范畴：
$\text{Gal}(K^{sep}/K)$ 的作用经由满射
$\text{Gal}(K^{sep}/K) \to \pi_1(X, \overline{\eta})$ 分解。因此，
$\text{Gal}(K^{sep}/M)$ 是其核。于是 $\text{Gal}(K^{sep}/M)$ 是正规子群，
$M/K$ 是 Galois 扩张，并且有短正合列
$$
1 \to \text{Gal}(K^{sep}/M) \to
\text{Gal}(K^{sep}/K) \to
\text{Gal}(M/K) \to 1
$$
这是由 Galois 理论得到的（《域》定理
\ref{fields-theorem-inifinite-galois-theory} 和引理
\ref{fields-lemma-ses-infinite-galois}）。证明完毕。
\end{proof}

\begin{lemma}
\label{pione-lemma-local-exact-sequence-normal}
设 $(A, \mathfrak m)$ 是正规局部环。置 $X = \Spec(A)$。令 $A^{sh}$ 为
$A$ 的严格 Hensel 化。令 $K$ 和 $K^{sh}$ 分别为 $A$ 和 $A^{sh}$ 的
分式域。则序列
$$
\pi_1(\Spec(K^{sh})) \to \pi_1(\Spec(K)) \to \pi_1(X) \to 1
$$
按引理 \ref{pione-lemma-functoriality-galois-ses} 第 (1) 部分的意义是正合的。
\end{lemma}

\begin{proof}
注意，$A^{sh}$ 是正规整环，见《代数补篇》引理
\ref{more-algebra-lemma-henselization-normal}。映射
$\pi_1(\Spec(K)) \to \pi_1(X)$ 由命题 \ref{pione-proposition-normal} 可知是满射。

\medskip\noindent
记 $X^{sh} = \Spec(A^{sh})$。设 $Y \to X$ 是有限 \'etale 态射。则
$Y^{sh} = Y \times_X X^{sh} \to X^{sh}$ 是有限 \'etale 态射。由于
$A^{sh}$ 严格 Hensel，$Y^{sh}$ 同构于若干个 $X^{sh}$ 的不交并。因此，
$Y \times_X \Spec(K^{sh})$ 亦然。于是复合
$\pi_1(\Spec(K^{sh})) \to \pi_1(X)$ 平凡，见引理
\ref{pione-lemma-composition-trivial}。

\medskip\noindent
为完成证明，按照引理 \ref{pione-lemma-functoriality-galois-ses}，只需证明：
给定有限 \'etale 态射 $V \to \Spec(K)$，若
$V \times_{\Spec(K)} \Spec(K^{sh})$ 是若干个 $\Spec(K^{sh})$ 的不交并，
则可以找到有限 \'etale 态射 $Y \to X$，使得
$V \cong Y \times_X \Spec(K)$ 在 $\Spec(K)$ 上成立。记
$V = \Spec(L)$，则 $L$ 是 $K$ 的若干有限可分扩张之有限直积。令
$B \subset L$ 为 $A$ 在 $L$ 中的整闭包。若 $A \to B$ 是 \'etale 态射，
则可以取 $Y = \Spec(B)$，证明即告完成。由《代数》引理
\ref{algebra-lemma-integral-closure-commutes-smooth}（以及一个略去的极限
论证），可知 $B \otimes_A A^{sh}$ 是 $A^{sh}$ 在
$L^{sh} = L \otimes_K K^{sh}$ 中的整闭包。我们的假设是 $L^{sh}$ 为若干个
$K^{sh}$ 的直积，因而 $B^{sh}$ 为若干个 $A^{sh}$ 的直积。所以
$A^{sh} \to B^{sh}$ 是 \'etale 态射。由于 $A \to A^{sh}$ 忠实平坦，
$A \to B$ 是 \'etale 态射（《下降》引理
\ref{descent-lemma-descending-property-etale}），正如所需。
\end{proof}







\section{群作用与整闭包}
\label{pione-section-group-actions-integral}

\noindent
本节继续《代数补篇》第
\ref{more-algebra-section-group-actions-integral} 节的讨论。
回顾正规局部环按定义是整环。

\begin{lemma}
\label{pione-lemma-get-algebraic-closure}
设 $A$ 是正规整环，其分式域 $K$ 可分代数闭。设
$\mathfrak p \subset A$ 是非零素理想。则剩余域
$\kappa(\mathfrak p)$ 代数闭。
\end{lemma}

\begin{proof}
反设该引理不成立以导出矛盾。于是存在首一不可约多项式
$P(T) \in \kappa(\mathfrak p)[T]$，其次数 $d > 1$。将 $P$ 替换为
$a^d P(a^{-1}T)$，其中 $a \in A$ 适当（以清除分母）之后，可以假设
$P$ 是 $Q$ 的像，其中后者是 $A[T]$ 中的首一多项式。注意，$Q$ 在
$K[T]$ 中不可约。
事实上，由《代数》引理 \ref{algebra-lemma-polynomials-divide}，$K$ 上的一个
因式分解会给出 $A$ 上的一个因式分解；将其模 $\mathfrak p$ 约化，就会得到
$P$ 的因式分解。由于 $K$ 可分闭，$Q$ 不是可分多项式
（《域》定义 \ref{fields-definition-separable}）。于是 $K$ 的特征为
$p > 0$，且 $Q$ 的一次项为零（《域》定义
\ref{fields-definition-separable}）。然而，此时可将 $Q$ 替换为
$Q + a T$，其中 $a \in \mathfrak p$ 非零，从而得到矛盾。
\end{proof}

\begin{lemma}
\label{pione-lemma-normal-local-domain-separablly-closed-fraction-field}
分式域可分闭的正规局部环是严格 Hensel 环。
\end{lemma}

\begin{proof}
设 $(A, \mathfrak m, \kappa)$ 是正规局部环，其分式域 $K$ 可分闭。
若 $A = K$，则结论成立。否则，由引理
\ref{pione-lemma-get-algebraic-closure}，剩余域 $\kappa$ 代数闭，因而只需验证
$A$ 是 Hensel 环。设 $f \in A[T]$ 首一，并设 $a_0 \in \kappa$ 是重数为
$1$ 的根，这里所约化的多项式为 $\overline{f} \in \kappa[T]$。设
$f = \prod f_i$ 是其在 $K[T]$ 中的因式分解。由《代数》引理
\ref{algebra-lemma-polynomials-divide}，有 $f_i \in A[T]$。因此，$a_0$ 是
某个 $f_i$ 的根，其中指标为 $i$。将 $f$ 替换为 $f_i$ 后，可以假设
$f$ 不可约。此时，导数 $f'$ 不可能在 $A[T]$ 中为零，因为 $a_0$ 是单根；
又可断定 $f$ 必为一次多项式，因为 $K$ 可分代数闭。因此 $A$ 是 Hensel 环，见《代数》
定义 \ref{algebra-definition-henselian}。
\end{proof}

\begin{lemma}
\label{pione-lemma-inertia-base-change}
设有限群 $G$ 作用于环 $R$。设 $R^G \to A$ 是环同态。设
$\mathfrak q' \subset A \otimes_{R^G} R$ 是位于素理想
$\mathfrak q \subset R$ 之上的素理想。则
$$
I_\mathfrak q = \{\sigma \in G \mid
\sigma(\mathfrak q) = \mathfrak q\text{ 且 }
\sigma \bmod \mathfrak q = \text{id}_{\kappa(\mathfrak q)}\}
$$
等于
$$
I_{\mathfrak q'} = \{\sigma \in G \mid
\sigma(\mathfrak q') = \mathfrak q'\text{ 且 }
\sigma \bmod \mathfrak q' = \text{id}_{\kappa(\mathfrak q')}\}
$$
\end{lemma}

\begin{proof}
由于 $\mathfrak q$ 是 $\mathfrak q'$ 的逆像，且
$\kappa(\mathfrak q) \subset \kappa(\mathfrak q')$，可得
$I_{\mathfrak q'} \subset I_\mathfrak q$。反过来，若
$\sigma \in I_\mathfrak q$，则该 $\sigma$ 在纤维环
$A \otimes_{R^G} \kappa(\mathfrak q)$ 上平凡作用。因此，$\sigma$ 固定所有
位于 $\mathfrak q$ 之上的素理想，并在它们的剩余域上诱导恒等映射。
\end{proof}

\begin{lemma}
\label{pione-lemma-inertia-invariants-etale}
设有限群 $G$ 作用于环 $R$。设 $\mathfrak q \subset R$ 是素理想。置
$$
I = \{\sigma \in G \mid \sigma(\mathfrak q) = \mathfrak q
\text{ 且 } \sigma \bmod \mathfrak q = \text{id}_\mathfrak q\}
$$
则 $R^G \to R^I$ 在 $R^I \cap \mathfrak q$ 处为 \'etale 的。
\end{lemma}

\begin{proof}
证明的策略是利用 \'etale 局部化，约化到 $R \to R^I$ 在
$R^I \cap \mathfrak p$ 处为局部同构的情形。设 $R^G \to A$ 是 \'etale
环同态。我们断言：如果对于 $G$ 在 $A \otimes_{R^G} R$ 上的作用以及
某个素理想 $\mathfrak q'$（它是 $A \otimes_{R^G} R$ 中位于
$\mathfrak q$ 之上的素理想），该结论成立，那么原结论亦成立。

\medskip\noindent
为验证这一点，注意，由于 $R^G \to A$ 平坦，有
$A = (A \otimes_{R^G} R)^G$，见《代数补篇》引理
\ref{more-algebra-lemma-base-change-invariants}。由引理
\ref{pione-lemma-inertia-base-change}，群 $I$ 不变。再次应用《代数补篇》引理
\ref{more-algebra-lemma-base-change-invariants} 可知
$A \otimes_{R^G} R^I = (A \otimes_{R^G} R)^I$
（因为 $R^I \to A \otimes_{R^G} R^I$ 平坦）。因此
$$
\xymatrix{
\Spec((A \otimes_{R^G} R)^I) \ar[d] \ar[r] & \Spec(R^I) \ar[d] \\
\Spec(A) \ar[r] & \Spec(R^G)
}
$$
是笛卡儿图，且水平箭头为 \'etale 的。因此，若左侧竖直箭头在某个开邻域
$W$（它是 $(A \otimes_{R^G} R)^I \cap \mathfrak q'$ 的邻域）上为
\'etale 的，则右侧竖直箭头在 $W$ 于 $\Spec(R^I)$ 中的（开）像的各点处
为 \'etale 的，见《下降》引理 \ref{descent-lemma-smooth-permanence}。
特别地，态射 $\Spec(R^I) \to \Spec(R^G)$ 在
$R^I \cap \mathfrak q$ 处为 \'etale 的。

\medskip\noindent
令 $\mathfrak p = R^G \cap \mathfrak q$。由《代数补篇》引理
\ref{more-algebra-lemma-one-orbit}，$\Spec(R) \to \Spec(R^G)$ 在
$\mathfrak p$ 上的纤维是有限的。再由《代数补篇》引理
\ref{more-algebra-lemma-one-orbit-geometric}，这些点处的剩余域扩张是
代数且正规的，并具有有限自同构群。因此，可以把《态射补篇》引理
\ref{more-morphisms-lemma-etale-makes-integral-split} 应用于整同态
$R^G \to R$ 和素理想 $\mathfrak p$。结合上述断言，可约化到如下情形：
$R = A_1 \times \ldots \times A_n$，每个 $A_i$ 恰有一个素理想
$\mathfrak q_i$ 位于 $\mathfrak p$ 之上，且剩余域扩张
$\kappa(\mathfrak q_i)/\kappa(\mathfrak p)$ 纯不可分。当然，
$\mathfrak q$ 是这些素理想之一，不妨设 $\mathfrak q = \mathfrak q_1$。

\medskip\noindent
$G$ 未必置换各因子 $A_i$（若 $A_i$ 的谱连通，则会如此；例如，若
$R^G$ 是局部环）。可按如下方式修正这一点；我们建议读者自行推演，或许可以
用幂等元代替拓扑来论证。回顾上述乘积分解给出 $\Spec(R)$ 相应的开闭子集
$U_i$ 的不交并分解。由于 $G$ 有限，可以把这个覆盖细化为有限不交并分解
$\Spec(R) = \coprod_{j \in J} W_j$，其中 $W_j$ 是开闭子集，使得对所有
$j \in J$，存在
$j' \in J$ 满足 $\sigma(W_j) = W_{j'}$。在这些 $W_j$ 中，不与
$\{\mathfrak q_1, \ldots, \mathfrak q_n\}$ 相交者的并是一个不与
$\mathfrak p$ 上纤维相交的闭子集。它因而映到 $\Spec(R^G)$ 中一个不含
$\mathfrak p$ 的闭子集，这是因为
$\Spec(R) \to \Spec(R^G)$ 是闭映射。因此，将 $R^G$ 替换为一个主局部化之后（上述断言
允许这样做），可以假设每个 $W_j$ 都与某个点 $\mathfrak q_i$ 相交。此时，
置 $U_i = W_j$，其中 $\mathfrak q_i \in W_j$。相应的乘积分解
$R = A_1 \times \ldots \times A_n$ 具有 $G$ 置换各因子 $A_i$ 的性质。

\medskip\noindent
所以，可以假设有一个乘积分解 $R = A_1 \times \ldots \times A_n$，
它与 $G$-作用相容；其中每个 $A_i$ 恰有一个素理想 $\mathfrak q_i$ 位于
$\mathfrak p$ 之上，且域扩张
$\kappa(\mathfrak q_i)/\kappa(\mathfrak p)$ 纯不可分。记
$A' = A_2 \times \ldots \times A_n$，于是
$$
R = A_1 \times A'
$$
由于 $\mathfrak q = \mathfrak q_1$，每个 $\sigma \in I$ 都保持上述乘积分解。
因此
$$
R^I = (A_1)^I \times (A')^I
$$
注意，有 $I = D = \{\sigma \in G \mid \sigma(\mathfrak q) = \mathfrak q\}$，
因为 $\kappa(\mathfrak q)/\kappa(\mathfrak p)$ 纯不可分。
由于 $G$ 在 $\mathfrak p$ 之上的素理想上的作用是传递的（《代数补篇》引理
\ref{more-algebra-lemma-one-orbit}），可知 $I$ 在 $G$ 中的指数为 $n$，并且可写成
$G = eI \amalg \sigma_2I \amalg \ldots \amalg \sigma_nI$，使得
$A_i = \sigma_i(A_1)$ 对 $i = 2, \ldots, n$ 成立。由此得到
$$
R^G = (A_1)^I.
$$
所以映射 $R^G \to R^I$ 在 $R^I \cap \mathfrak q$ 处为 \'etale 的，
证明完毕。
\end{proof}

\noindent
下一个引理推广了《代数补篇》引理
\ref{more-algebra-lemma-inertial-invariants-unramified}。

\begin{lemma}
\label{pione-lemma-inertial-invariants-unramified}
设 $A$ 是分式域为 $K$ 的正规整环。设 $L/K$ 是一个（可能无限的）Galois
扩张。令 $G = \text{Gal}(L/K)$，并令 $B$ 为 $A$ 在 $L$ 中的整闭包。
设 $\mathfrak q \subset B$。置
$$
I = \{\sigma \in G \mid
\sigma(\mathfrak q) = \mathfrak q \text{ 且 }
\sigma \bmod \mathfrak q = \text{id}_{\kappa(\mathfrak q)}\}
$$
则 $(B^I)_{B^I \cap \mathfrak q}$ 是 \'etale $A$-代数的滤过余极限。
\end{lemma}

\begin{proof}
可将 $L$ 写成 $K$ 的有限 Galois 扩张的滤过余极限。因此，只需在 $L/K$
为有限 Galois 扩张时证明该引理，见《代数》引理
\ref{algebra-lemma-colimit-colimit-etale}。因为 $A = B^G$（$A$ 在
$K = L^G$ 中整闭），故结论由引理 \ref{pione-lemma-inertia-invariants-etale} 得出。
\end{proof}







\section{分歧理论}
\label{pione-section-ramification}

\noindent
本节继续 More on Algebra, Section
\ref{more-algebra-section-ramification} 中的讨论，并将其与我们对概形基本群的讨论联系起来。

\medskip\noindent
设 $(A, \mathfrak m, \kappa)$ 是分式域为 $K$ 的正规局部环。选取一个可分代数闭包
$K^{sep}$。设 $A^{sep}$ 为 $A$ 在 $K^{sep}$ 中的整闭包，并选取极大理想
$\mathfrak m^{sep} \subset A^{sep}$。设 $A \subset A^h \subset A^{sh}$
分别为 Hensel 化和严格 Hensel 化。注意，$A^h$ 和 $A^{sh}$ 也都是正规环
（More on Algebra, Lemma \ref{more-algebra-lemma-henselization-normal}）。以
$K^h$ 和 $K^{sh}$ 表示它们的分式域。由 Lemma
\ref{pione-lemma-normal-local-domain-separablly-closed-fraction-field}，
$(A^{sep})_{\mathfrak m^{sep}}$ 是严格 Hensel 环，故可以选取一个 $A$-代数映射
$A^{sh} \to (A^{sep})_{\mathfrak m^{sep}}$。具体来说，先选取一个
$\kappa$-嵌入\footnote{这是可行的，因为 $\kappa(\mathfrak m^{sh})$ 是
$\kappa$ 的可分代数闭包，而由 Lemma \ref{pione-lemma-get-algebraic-closure}，
$\kappa(\mathfrak m^{sep})$ 是 $\kappa$ 的代数闭包。}
$\kappa(\mathfrak m^{sh}) \to \kappa(\mathfrak m^{sep})$，再由 Algebra,
Lemma \ref{algebra-lemma-strictly-henselian-functorial} 将其（唯一地）延拓为
一个 $A$-代数同态。由此得到交换图
$$
\xymatrix{
K^{sep} & K^{sh} \ar[l] & K^h \ar[l] & K \ar[l] \\
(A^{sep})_{\mathfrak m^{sep}} \ar[u] &
A^{sh} \ar[u] \ar[l] &
A^h \ar[u] \ar[l] &
A \ar[u] \ar[l]
}
$$
可以取这些环之谱的基本群。当然，因为 $K^{sep}$、
$(A^{sep})_{\mathfrak m^{sep}}$ 和 $A^{sh}$ 都是严格 Hensel 的，所得的群均为平凡群。
因此有意义的部分是
\begin{equation}
\label{pione-equation-inertia-diagram-pione}
\vcenter{
\xymatrix{
\pi_1(U^{sh}) \ar[r] \ar[rd]_1 & \pi_1(U^h) \ar[d] \ar[r] & \pi_1(U) \ar[d] \\
& \pi_1(X^h) \ar[r] & \pi_1(X)
}
}
\end{equation}
这里 $X^h$ 和 $X$ 分别是 $A^h$ 和 $A$ 的谱，而 $U^{sh}$、$U^h$、$U$
分别是 $K^{sh}$、$K^h$、$K$ 的谱。标号 $1$ 表示该映射是平凡的；这是因为它经过平凡群
$\pi_1(X^{sh})$ 分解。另一方面，profinite 群
$G = \text{Gal}(K^{sep}/K)$ 作用在 $A^{sep}$ 上，因而可以作如下定义：
$$
D = \{\sigma \in G \mid \sigma(\mathfrak m^{sep}) = \mathfrak m^{sep}\}
\supset
I = \{\sigma \in D \mid \sigma \bmod \mathfrak m^{sep} =
\text{id}_{\kappa(\mathfrak m^{sep})}\}
$$
这两个群有时分别称为{\it 分解群}和{\it 惯性群}，尤其是在 $A$ 为离散赋值环时。

\begin{lemma}
\label{pione-lemma-identify-inertia}
在上述情形中，通过同构 $\pi_1(U) = \text{Gal}(K^{sep}/K)$，交换图
(\ref{pione-equation-inertia-diagram-pione}) 转化为交换图
$$
\xymatrix{
I \ar[r] \ar[rd]_1 & D \ar[d] \ar[r] & \text{Gal}(K^{sep}/K) \ar[d] \\
& \text{Gal}(\kappa(\mathfrak m^{sh})/\kappa) \ar[r] & \text{Gal}(M/K)
}
$$
其中 $K^{sep}/M/K$ 是关于 $A$ 非分歧的最大子扩张。此外，竖直箭头均为满射，
左侧竖直箭头的核为 $I$，右侧竖直箭头的核是
$\text{Gal}(K^{sep}/K)$ 中包含 $I$ 的最小闭正规子群。
\end{lemma}

\begin{proof}
按构造，群 $D$ 作用于 $(A^{sep})_{\mathfrak m^{sep}}$，且该作用在 $A$ 上。
给定剩余域上的映射后，
$A^{sh} \to (A^{sep})_{\mathfrak m^{sep}}$ 具有唯一性（Algebra, Lemma
\ref{algebra-lemma-strictly-henselian-functorial}），故映射
$A^{sh} \to (A^{sep})_{\mathfrak m^{sep}}$ 的像包含于
$((A^{sep})_{\mathfrak m^{sep}})^I$。另一方面，Lemma
\ref{pione-lemma-inertial-invariants-unramified} 表明
$((A^{sep})_{\mathfrak m^{sep}})^I$ 是 $A$ 的 \'etale 扩张的滤过余极限。
由于 $A^{sh}$ 是最大的这种扩张，遂有
$A^{sh} = ((A^{sep})_{\mathfrak m^{sep}})^I$，从而 $K^{sh} = (K^{sep})^I$。

\medskip\noindent
回顾 $I$ 是一个满射 $D \to \text{Aut}(\kappa(\mathfrak m^{sep})/\kappa)$
的核，见 More on Algebra, Lemma
\ref{more-algebra-lemma-one-orbit-geometric-galois}。我们有
$\text{Aut}(\kappa(\mathfrak m^{sep})/\kappa) =
\text{Gal}(\kappa(\mathfrak m^{sh})/\kappa)$，因为由上文可知，这些域分别是
$\kappa$ 的代数闭包与可分代数闭包。另一方面，由 Algebra, Lemma
\ref{algebra-lemma-henselian-functorial} 中的唯一性，$A^{sh}$ 在 $A$ 上的任意自同构
也是 $A^{sh}$ 在 $A^h$ 上的自同构。此外，$A^{sh}$ 是有限 \'etale 扩张
$A^h \subset A'$ 的余极限，而这些扩张与有限可分扩张成 $1$ 对 $1$ 的对应，
后者形如 $\kappa'/\kappa$，
见 Algebra, Remark \ref{algebra-remark-construct-sh-from-h}。因此
$$
\text{Aut}(A^{sh}/A) = \text{Aut}(A^{sh}/A^h) =
\text{Gal}(\kappa(\mathfrak m^{sh})/\kappa)
$$
设 $\kappa'/\kappa$ 是 Galois 群为 $G$ 的有限 Galois 扩张。由 Algebra, Lemma
\ref{algebra-lemma-henselian-cat-finite-etale}，设 $A^h \subset A'$ 是与
$\kappa \subset \kappa'$ 对应的有限 \'etale 扩张。考察分式域与次数可得
$(A')^G = A^h$（略去一个小细节）。取余极限即得
$(A^{sh})^{\text{Gal}(\kappa(\mathfrak m^{sh})/\kappa)} = A^h$。
综合以上各点，可得 $A^h = ((A^{sep})_{\mathfrak m^{sep}})^D$，从而
$K^h = (K^{sep})^D$。

\medskip\noindent
因为 $U$、$U^h$、$U^{sh}$ 分别是域 $K$、$K^h$、$K^{sh}$ 的谱，所以由 Lemma
\ref{pione-lemma-fundamental-group-Galois-group} 可知两个交换图的上行彼此对应。由 Lemma
\ref{pione-lemma-gabber} 有
$\pi_1(X^h) = \text{Gal}(\kappa(\mathfrak m^{sh})/\kappa)$。上文已经指出正合列
$1 \to I \to D \to \text{Gal}(\kappa(\mathfrak m^{sh})/\kappa) \to 1$
的正合性。由 Proposition \ref{pione-proposition-normal} 可知
$\pi_1(X) = \text{Gal}(M/K)$。最后，关于
$\text{Gal}(K^{sep}/K) \to \text{Gal}(M/K) = \pi_1(X)$ 之核的断言由 Lemma
\ref{pione-lemma-local-exact-sequence-normal} 得出。证毕。
\end{proof}

\noindent
设 $X$ 是函数域为 $K$ 的正规整概形，$K^{sep}$ 是 $K$ 的一个可分代数闭包，
而 $X^{sep} \to X$ 是 $X$ 在 $K^{sep}$ 中的正规化。由于
$G = \text{Gal}(K^{sep}/K)$ 作用于 $K^{sep}$，便得到 $G$ 在 $X^{sep}$ 上的右作用。
对 $y \in X^{sep}$，定义
$$
D_y = \{\sigma \in G \mid \sigma(y) = y\} \supset
I_y = \{\sigma \in D \mid \sigma \bmod \mathfrak m_y =
\text{id}_{\kappa(y)} \}
$$
与上文类似。另一方面，对 $x \in X$，设 $\mathcal{O}_{X, x}^{sh}$ 是一个严格
Hensel 化，设 $K_x^{sh}$ 是 $\mathcal{O}_{X, x}^{sh}$ 的分式域，并选取一个
$K$-嵌入 $K_x^{sh} \to K^{sep}$。

\begin{lemma}
\label{pione-lemma-normal-pione-quotient-inertia}
设 $X$ 是函数域为 $K$ 的正规整概形。沿用上述记号，
$\text{Gal}(K^{sep}/K) = \pi_1(\Spec(K))$ 的下列三个子群相等：
\begin{enumerate}
\item 满射
$\text{Gal}(K^{sep}/K) \longrightarrow \pi_1(X)$,
的核；
\item 包含所有 $I_y$（其中 $y \in X^{sep}$）的最小闭正规子群；以及
\item 包含所有 $\text{Gal}(K^{sep}/K_x^{sh})$（其中 $x \in  X$）的最小闭正规子群。
\end{enumerate}
\end{lemma}

\begin{proof}
(2) 与 (3) 的等价性来自 Lemma \ref{pione-lemma-identify-inertia}：$I_y$ 与
$\text{Gal}(K^{sep}/K_x^{sh})$ 共轭，如果 $y$ 位于 $x$ 上方。由 Lemma
\ref{pione-lemma-local-exact-sequence-normal} 可知
$\text{Gal}(K^{sep}/K_x^{sh})$ 到 $\pi_1(\Spec(\mathcal{O}_{X, x}))$ 的映射是平凡的，
因此 (2) 与 (3) 中的子群 $N \subset G = \text{Gal}(K^{sep}/K)$ 包含于
$G \longrightarrow \pi_1(X)$ 的核。

\medskip\noindent
为证明反向包含关系，因 $N$ 正规，只需证明：若 $N \subset U \subset G$ 且 $U$ 是开正规子群，
则商映射 $G \to G/U$ 经过 $\pi_1(X)$ 分解。换言之，若 $L/K$ 是与 $U$ 对应的 Galois
扩张，则必须证明 $X$ 在 $L$ 中非分歧（Section \ref{pione-section-normal}，特别是
Proposition \ref{pione-proposition-normal}）。只需在 $X$ 仿射时证明这一点（如此便可在以下证明中引用代数结论）。
设 $Y \to X$ 是 $X$ 在 $L$ 中的正规化。包含映射 $L \subset K^{sep}$ 诱导态射
$\pi : X^{sep} \to Y$。对 $y \in X^{sep}$，$\pi(y)$ 在 $\text{Gal}(L/K)$ 中的惯性群
是 $I_y$ 在 $\text{Gal}(L/K)$ 中的像；这由 More on Algebra, Lemma
\ref{more-algebra-lemma-one-orbit-geometric-galois-compare} 得出。由于 $N \subset U$，
所有这些惯性群都平凡。应用 Lemma \ref{pione-lemma-inertia-invariants-etale} 可知
$Y \to X$ 是 \'etale 的。（另一种做法是使用 Lemma
\ref{pione-lemma-local-exact-sequence-normal}，看出 $Y$ 在
$\Spec(\mathcal{O}_{X, x})$ 上的拉回对所有 $x \in X$ 都是 \'etale 的，
再多做一些工作便可得出结论。）
\end{proof}

\begin{example}
\label{pione-example-bigger-codim}
设 $X$ 是函数域为 $K$ 的正规整 Noether 概形。分歧轨迹的纯性（见下文）告诉我们，
若 $X$ 正则，则在 Lemma \ref{pione-lemma-normal-pione-quotient-inertia} 中，只需考虑惯性群
$I = \pi_1(\Spec(K_x^{sh}))$，其中 $x$ 是余维为 $1$ 的 $X$ 中的点。
但一般而言这并不充分。事实上，令
$Y = \mathbf{A}_k^n = \Spec(k[t_1, \ldots, t_n])$，其中 $k$ 是特征不为 $2$ 的域。
令 $G = \{\pm 1\}$ 是阶为 $2$ 的群，它通过各坐标上的乘法作用于 $Y$。令
$$
X = \Spec(k[t_it_j, i, j \in \{1, \ldots, n\}])
$$
嵌入 $k[t_it_j] \subset k[t_1, \ldots, t_n]$ 定义了一个次数为 $2$ 的态射
$Y \to X$；除极大理想 $\mathfrak m = (t_it_j)$ 上方外，它处处非分歧，
而该理想给出一个余维为 $n$ 的 $X$ 中的点。
\end{example}

\begin{lemma}
\label{pione-lemma-unramified}
设 $X$ 是函数域为 $K$ 的整正规概形。设 $L/K$ 是有限扩张，并设 $Y \to X$
是 $X$ 在 $L$ 中的正规化。下列条件等价：
\begin{enumerate}
\item 按 Section \ref{pione-section-normal} 中的定义，$X$ 在 $L$ 中非分歧；
\item $Y \to X$ 是概形的非分歧态射；
\item $Y \to X$ 是概形的 \'etale 态射；
\item $Y \to X$ 是概形的有限 \'etale 态射；
\item 对 $x \in X$，投影
$Y \times_X \Spec(\mathcal{O}_{X, x}) \to \Spec(\mathcal{O}_{X, x})$
非分歧；
\item 将 (5) 中的局部环换成 $\mathcal{O}_{X, x}^h$ 后，结论仍成立；
\item 将 (5) 中的局部环换成 $\mathcal{O}_{X, x}^{sh}$ 后，结论仍成立；
\item 对 $x \in X$，概形论纤维 $Y_x$ 在 $x$ 上是次数 $\geq [L : K]$ 的 \'etale 概形。
\end{enumerate}
若 $L/K$ 是 Galois 群为 $G$ 的 Galois 扩张，则上述条件还等价于
\begin{enumerate}
\item[(9)] 对 $y \in Y$，群
$I_y = \{g \in G \mid g(y) = y\text{ 且 }
g \bmod \mathfrak m_y = \text{id}_{\kappa(y)}\}$ 平凡。
\end{enumerate}
\end{lemma}

\begin{proof}
(1) 与 (2) 的等价性就是 (1) 的定义。(2)、(3)、(4) 的等价性由 Lemma
\ref{pione-lemma-unramified-in-L} 给出。容易证明 (4) $\Rightarrow$ (5)、
(5) $\Rightarrow$ (6)、(6) $\Rightarrow$ (7)。

\medskip\noindent
假设 (7) 成立。注意 $\mathcal{O}_{X, x}^{sh}$ 是正规局部整环
（More on Algebra, Lemma \ref{more-algebra-lemma-henselization-normal}）。
令 $L^{sh} = L \otimes_K K_x^{sh}$，其中 $K_x^{sh}$ 是
$\mathcal{O}_{X, x}^{sh}$ 的分式域。于是
$L^{sh} = \prod_{i = 1, \ldots, n} L_i$，且 $L_i/K_x^{sh}$ 是有限可分扩张。
由 Algebra, Lemma \ref{algebra-lemma-integral-closure-commutes-smooth}
（以及一个略去的极限论证）可知，$Y \times_X \Spec(\mathcal{O}_{X, x}^{sh})$
是 $\Spec(\mathcal{O}_{X, x}^{sh})$ 在 $L^{sh}$ 中的整闭包。因此，由 Lemma
\ref{pione-lemma-unramified-in-L}（分别应用于各因子 $L_i$，它们来自 $L^{sh}$），可知
$Y \times_X \Spec(\mathcal{O}_{X, x}^{sh}) \to \Spec(\mathcal{O}_{X, x}^{sh})$
是有限 \'etale 的。考察一般点可知其次数等于 $[L : K]$，故 (8) 成立。

\medskip\noindent
假设 (8) 成立。设 $x \in X$，并设概形论纤维 $Y_x$ 在 $x$ 上是次数
$\geq [L : K]$ 的 \'etale 概形。这意味着 $Y$ 有
$\geq [L : K]$ 个位于 $x$ 上方的几何点。我们将证明 $Y \to X$ 在 $x$ 的某个邻域上是有限
\'etale 的，这便证明了 (1)。
为此可以假设 $X = \Spec(R)$，点 $x$ 对应于素理想 $\mathfrak p \subset R$，
并且 $Y = \Spec(S)$。应用 More on Morphisms, Lemma
\ref{more-morphisms-lemma-etale-makes-integral-split}，可找到一个 \'etale 邻域
$(U, u) \to (X, x)$，使得
$Y \times_X U = V_1 \amalg \ldots \amalg V_m$，并且 $V_i$ 有唯一一个点 $v_i$，
它位于 $u$ 上方，而 $\kappa(v_i)/\kappa(u)$ 是纯不可分扩张。必要时缩小 $U$，
可以假设 $U$ 是一般点为 $\xi$ 的正规整概形（使用 Descent, Lemmas
\ref{descent-lemma-locally-finite-nr-irred-local-fppf}、
\ref{descent-lemma-normal-local-smooth}，以及 Properties, Lemma
\ref{properties-lemma-normal-locally-finite-nr-irreducibles}）。
由上面对几何点的说明可知 $m \geq [L : K]$。另一方面，由 More on Morphisms, Lemma
\ref{more-morphisms-lemma-normalization-smooth-localization} 可知
$\coprod V_i \to U$ 是 $U$ 在 $\Spec(L) \times_X U$ 中的正规化。因为
$K \subset \kappa(\xi)$ 是有限可分扩张，可以写成
$\Spec(L) \times_X U = \Spec(\prod_{i = 1, \ldots, n} L_i)$，
其中 $L_i/\kappa(\xi)$ 有限，且 $[L : K] = \sum [L_i : \kappa(\xi)]$。
由于 $V_j$ 对每个 $j$ 都非空，并且 $m \geq [L : K]$，所以
$m = n$，且 $[L_i : \kappa(\xi)] = 1$ 对所有 $i$ 都成立。于是
$V_j \to U$ 是同构，特别地是 \'etale 的，故 $Y \times_X U \to U$ 是
\'etale 的。由 Descent, Lemma \ref{descent-lemma-descending-property-etale}，
可知 $Y \to X$ 在 $U \to X$ 的像（即 $x$ 的一个开邻域）上是 \'etale 的。

\medskip\noindent
假设 $L/K$ 是 Galois 扩张且 (9) 成立。由 Lemma
\ref{pione-lemma-inertial-invariants-unramified}，$Y \to X$ 是 \'etale 的。
略去 (1) 蕴含 (9) 的证明。
\end{proof}

\noindent
对于离散赋值环的无限 Galois 扩张，还可以多说一点。为此引入如下记号。
称整数子集 $S \subset \mathbf{N}$ {\it 乘法有向}，如果 $1 \in S$，且对
$n, m \in S$，存在 $k \in S$ 使得 $n | k$ 且 $m | k$。在 $S$ 上定义偏序：
$n \geq_S m$ 当且仅当 $m | n$。给定域 $\kappa$，得到有限群的逆系统
$\{\mu_n(\kappa)\}_{n \in S}$，其转移映射为
$$
\mu_n(\kappa) \longrightarrow \mu_m(\kappa),\quad
\zeta \longmapsto \zeta^{n/m}
$$
其中 $n \geq_S m$。于是可以构造 profinite 群
$$
\lim_{n \in S} \mu_n(\kappa)
$$
注意，该极限是余滤过的（因为 $S$ 有向）。此构造关于 $\kappa$ 具有函子性。
特别地，$\text{Aut}(\kappa)$ 作用于这个 profinite 群。例如，若
$S = \{1, n\}$，则得到 $\mu_n(\kappa)$。若
$S = \{1, \ell, \ell^2, \ell^3, \ldots\}$，其中 $\ell$ 是不同于 $\kappa$
之特征的素数，则得到 $\lim_n \mu_{\ell^n}(\kappa)$；它有时称为
$\ell$-进 Tate 模，即 $\kappa$ 的乘法群的 Tate 模（参见 More on Algebra, Example
\ref{more-algebra-example-spectral-sequence-principal}）。

\begin{lemma}
\label{pione-lemma-structure-decomposition}
设 $A$ 是分式域为 $K$ 的离散赋值环。设 $L/K$ 是一个（可能无限的）Galois 扩张，
$B$ 是 $A$ 在 $L$ 中的整闭包，$\mathfrak m$ 是 $B$ 的一个极大理想。令
$G = \text{Gal}(L/K)$、
$D = \{\sigma \in G \mid \sigma(\mathfrak m) = \mathfrak m\}$，以及
$I = \{\sigma \in D \mid \sigma \bmod \mathfrak m =
\text{id}_{\kappa(\mathfrak m)}\}$.
分解群 $D$ 属于如下典范正合列：
$$
1 \to I \to D \to \text{Aut}(\kappa(\mathfrak m)/\kappa_A) \to 1
$$
惯性群 $I$ 属于如下典范正合列：
$$
1 \to P \to I \to I_t \to 1
$$
其中
\begin{enumerate}
\item 野惯性群 $P$ 是 $D$ 的正规子群；
\item $P$ 是 pro-$p$-群，如果 $\kappa_A$ 的特征为 $p > 1$；并且
$P = \{1\}$，如果 $\kappa_A$ 的特征为零；
\item 存在一个乘法有向子集 $S \subset \mathbf{N}$，使得
$\kappa(\mathfrak m)$ 含有一个本原 $n$ 次单位根，对每个 $n \in S$
（$S$ 的元素均与 $p$ 互素）；
\item 存在一个典范满射
$$
\theta_{can} : I \to \lim_{n \in S} \mu_n(\kappa(\mathfrak m))
$$
其核为 $P$，且满足
$\theta_{can}(\tau \sigma \tau^{-1}) = \tau(\theta_{can}(\sigma))$
（其中 $\tau \in D$、$\sigma \in I$），并诱导同构
$I_t \to \lim_{n \in S} \mu_n(\kappa(\mathfrak m))$；这里的商群称为温和惯性群。
\end{enumerate}
\end{lemma}

\begin{proof}
这基本上只是重述 More on Algebra, Section
\ref{more-algebra-section-ramification} 中所证明的有限 Galois 扩张结果。
映射 $D \to \text{Aut}(\kappa(\mathfrak m)/\kappa)$ 的满射性就是
More on Algebra, Lemma \ref{more-algebra-lemma-one-orbit-geometric-galois}。
由此得到第一个正合列。

\medskip\noindent
为构造第二个短正合列，令 $\Lambda$ 为有限 Galois 子扩张之集；也就是说，
$\lambda \in \Lambda$ 对应于 $L/L_\lambda/K$。令
$G_\lambda = \text{Gal}(L_\lambda/K)$。回顾 $G_\lambda$ 构成转移映射均为满射的有限群逆系统，
并且 $G = \lim_{\lambda \in \Lambda} G_\lambda$，见 Fields, Lemma
\ref{fields-lemma-infinite-galois-limit}。令 $B_\lambda$ 为 $A$ 在 $L_\lambda$ 中的整闭包。
再令 $\mathfrak m_\lambda = \mathfrak m \cap B_\lambda$，并以
$P_\lambda, I_\lambda, D_\lambda$ 分别表示 $\mathfrak m_\lambda$ 的
野惯性群、惯性群和分解群，见 More on Algebra, Lemma
\ref{more-algebra-lemma-galois-inertia}。对 $\lambda \geq \lambda'$，限制映射给出交换图
$$
\xymatrix{
P_\lambda \ar[d] \ar[r] &
I_\lambda \ar[d] \ar[r] &
D_\lambda \ar[d] \ar[r] &
G_\lambda \ar[d] \\
P_{\lambda'} \ar[r] &
I_{\lambda'} \ar[r] &
D_{\lambda'} \ar[r] &
G_{\lambda'}
}
$$
其中竖直映射均为满射，见 More on Algebra, Lemma
\ref{more-algebra-lemma-compare-inertia}。

\medskip\noindent
由定义立刻可知，$I = \lim I_\lambda$ 且 $D = \lim D_\lambda$；
这里使用的是上述同构 $G = \lim G_\lambda$。因为 $L = \colim L_\lambda$，所以
$B = \colim B_\lambda$ 且
$\kappa(\mathfrak m) = \colim \kappa(\mathfrak m_\lambda)$。
系统 $D_\lambda$ 的转移映射与映射
$D_\lambda \to \text{Aut}(\kappa(\mathfrak m_\lambda)/\kappa)$ 相容
（见 More on Algebra, Lemma \ref{more-algebra-lemma-compare-inertia}），
故映射 $D \to \text{Aut}(\kappa(\mathfrak m)/\kappa)$ 是诸映射
$D_\lambda \to \text{Aut}(\kappa(\mathfrak m_\lambda)/\kappa)$ 的极限。

\medskip\noindent
存在典范映射
$$
\theta_{\lambda, can} :
I_\lambda
\longrightarrow
\mu_{n_\lambda}(\kappa(\mathfrak m_\lambda))
$$
其中 $n_\lambda = |I_\lambda|/|P_\lambda|$，群
$\mu_{n_\lambda}(\kappa(\mathfrak m_\lambda))$ 的阶为 $n_\lambda$，并且
$\theta_{\lambda, can}(\tau \sigma \tau^{-1}) =
\tau(\theta_{\lambda, can}(\sigma))$ 对 $\tau \in D_\lambda$ 和
$\sigma \in I_\lambda$ 成立，而且得到交换图
$$
\xymatrix{
I_\lambda \ar[r]_-{\theta_{\lambda, can}} \ar[d] &
\mu_{n_\lambda}(\kappa(\mathfrak m_\lambda))
\ar[d]^{(-)^{n_\lambda/n_{\lambda'}}} \\
I_{\lambda'} \ar[r]^-{\theta_{\lambda', can}} &
\mu_{n_{\lambda'}}(\kappa(\mathfrak m_{\lambda'}))
}
$$
见 More on Algebra, Remark
\ref{more-algebra-remark-canonical-inertia-character}。

\medskip\noindent
令 $S \subset \mathbf{N}$ 为整数 $n_\lambda$ 的集合。由于 $\Lambda$ 有向，
可知 $S$ 是乘法有向的。对上面交换图中的映射 $\theta_{\lambda, can}$ 取极限，得到
$$
\theta_{can} : I \to \lim_{n \in S} \mu_n(\kappa(\mathfrak m)).
$$
此映射连续（略去一个小细节）。由于系统 $I_\lambda$ 的转移映射都是满射，
且 $\Lambda$ 有向，投影 $I \to I_\lambda$ 均为满射。对每个 $\lambda$，交换图
$$
\xymatrix{
I \ar[d] \ar[r]_-{\theta_{can}} &
\lim_{n \in S} \mu_n(\kappa(\mathfrak m)) \ar[d] \\
I_{\lambda} \ar[r]^-{\theta_{\lambda, can}} &
\mu_{n_\lambda}(\kappa(\mathfrak m_\lambda))
}
$$
成立。因此 $\theta_{can}$ 的像满射到有限群
$\mu_{n_\lambda}(\kappa(\mathfrak m)) =
\mu_{n_\lambda}(\kappa(\mathfrak m_\lambda))$；该群的阶为 $n_\lambda$
（见上文）。于是 $\theta_{can}$ 的像稠密。另一方面，$\theta_{can}$ 连续且源是
profinite 群。因此，一个拓扑论证表明 $\theta_{can}$ 是满射。

\medskip\noindent
等式 $\theta_{can}(\tau \sigma \tau^{-1}) = \tau(\theta_{can}(\sigma))$
对 $\tau \in D$、$\sigma \in I$ 成立；这是由映射 $\theta_{\lambda, can}$
的相应性质，以及映射 $D \to \text{Aut}(\kappa(\mathfrak m))$ 同各映射
$D_\lambda \to \text{Aut}(\kappa(\mathfrak m_\lambda))$ 的相容性得出的。
令 $P = \Ker(\theta_{can})$，这表明 $P$ 是 $D$ 的正规子群。令 $I_t = I/P$；
得到同构 $I_t \to \lim_{n \in S} \mu_n(\kappa(\mathfrak m))$，
这是由 $\theta_{can}$ 的满射性推出的。

\medskip\noindent
为完成证明，我们证明 $P = \lim P_\lambda$，这说明 $P$ 是 pro-$p$-群。回顾温和惯性群
$I_{\lambda, t} = I_\lambda/P_\lambda$ 的阶为 $n_\lambda$。由于转移映射
$P_\lambda \to P_{\lambda'}$ 均为满射，且 $\Lambda$ 有向，得到短正合列
$$
1 \to \lim P_\lambda \to I \to \lim I_{\lambda, t} \to 1
$$
（略去细节）。因为对每个 $\lambda$，映射 $\theta_{\lambda, can}$ 诱导同构
$I_{\lambda, t} \cong \mu_{n_\lambda}(\kappa(\mathfrak m))$
，所需结论随即得出。
\end{proof}

\begin{lemma}
\label{pione-lemma-structure-decomposition-separable-closure}
设 $A$ 是分式域为 $K$ 的离散赋值环，$K^{sep}$ 是 $K$ 的一个可分闭包，
$A^{sep}$ 是 $A$ 在 $K^{sep}$ 中的整闭包，而 $\mathfrak m^{sep}$ 是
$A^{sep}$ 的一个极大理想。令 $\mathfrak m = \mathfrak m^{sep} \cap A$、
$\kappa = A/\mathfrak m$，以及
$\overline{\kappa} = A^{sep}/\mathfrak m^{sep}$。于是
$\overline{\kappa}$ 是 $\kappa$ 的代数闭包。令
$G = \text{Gal}(K^{sep}/K)$、
$D = \{\sigma \in G \mid \sigma(\mathfrak m^{sep}) = \mathfrak m^{sep}\}$，以及
$I = \{\sigma \in D \mid \sigma \bmod \mathfrak m^{sep} =
\text{id}_{\kappa(\mathfrak m^{sep})}\}$.
分解群 $D$ 属于如下典范正合列：
$$
1 \to I \to D \to \text{Gal}(\kappa^{sep}/\kappa) \to 1
$$
其中 $\kappa^{sep} \subset \overline{\kappa}$ 是 $\kappa$ 的可分闭包。
惯性群 $I$ 属于如下典范正合列：
$$
1 \to P \to I \to I_t \to 1
$$
其中
\begin{enumerate}
\item 野惯性群 $P$ 是 $D$ 的正规子群；
\item $P$ 是 pro-$p$-群，如果 $\kappa_A$ 的特征为 $p > 1$；并且
$P = \{1\}$，如果 $\kappa_A$ 的特征为零；
\item 存在一个典范满射
$$
\theta_{can} : I \to \lim_{n\text{ prime to }p} \mu_n(\kappa^{sep})
$$
其核为 $P$，且满足
$\theta_{can}(\tau \sigma \tau^{-1}) = \tau(\theta_{can}(\sigma))$
（其中 $\tau \in D$、$\sigma \in I$），并诱导同构
$I_t \to \lim_{n\text{ prime to }p} \mu_n(\kappa^{sep})$。
\end{enumerate}
\end{lemma}

\begin{proof}
由 Lemma \ref{pione-lemma-get-algebraic-closure}，域 $\overline{\kappa}$ 是
$\kappa$ 的代数闭包。大部分断言立即由 Lemma
\ref{pione-lemma-structure-decomposition} 的相应部分得出。例如，因为
$\text{Aut}(\overline{\kappa}/\kappa) = \text{Gal}(\kappa^{sep}/\kappa)$，
便得到第一个正合列。
因此，唯一还需证明的断言是：若 $S$ 如 Lemma
\ref{pione-lemma-structure-decomposition} 中所述，则极限
$\lim_{n \in S} \mu_n(\overline{\kappa})$ 等于
$\lim_{n\text{ prime to }p} \mu_n(\kappa^{sep})$。为此，只需证明每个整数 $n$
只要与 $p$ 互素，就整除 $S$ 的某个元素。
令 $\pi \in A$ 是一个素元，并考虑分裂域 $L$，即多项式 $X^n - \pi$ 的分裂域。
因为该多项式可分，$L$ 是 $K$ 的有限 Galois 扩张。选取嵌入
$L \to K^{sep}$。注意，若 $B$ 是 $A$ 在 $L$ 中的整闭包，则
$A \to B_{\mathfrak m^{sep} \cap B}$ 的分歧指数可被 $n$ 整除
（因为 $\pi$ 有一个 $n$ 次根位于 $B$ 中；事实上该分歧指数等于 $n$，但此处用不到）。
于是，由 Lemma \ref{pione-lemma-structure-decomposition} 的证明中 $S$ 的构造可知，
$n$ 整除 $S$ 的某个元素。
\end{proof}








\section{几何基本群与算术基本群}
\label{pione-section-galois-action}

\noindent
本节研究如下比较：概形 $X$ 在域 $k$ 上的基本群与 $X_{\overline{k}}$
的基本群之间有何关系，其中 $\overline{k}$ 是 $k$ 的代数闭包。

\begin{lemma}
\label{pione-lemma-limit}
设 $I$ 是有向集。设 $X_i$ 是以 $I$ 为指标、转移态射均为仿射态射的拟紧拟分离概形逆系统。
依照 Limits, Section \ref{limits-section-limits}，令 $X = \lim X_i$。
则有范畴等价
$$
\colim \textit{F\'Et}_{X_i} = \textit{F\'Et}_X
$$
若 $X_i$ 对所有充分大的 $i$ 都连通，并且 $\overline{x}$ 是 $X$ 的一个几何点，则
$$
\pi_1(X, \overline{x}) = \lim \pi_1(X_i, \overline{x})
$$
\end{lemma}

\begin{proof}
该范畴等价由 Limits, Lemmas
\ref{limits-lemma-descend-finite-presentation},
\ref{limits-lemma-descend-finite-finite-presentation} 和
\ref{limits-lemma-descend-etale} 得出。第二个断言是范畴断言的形式推论。
\end{proof}

\begin{lemma}
\label{pione-lemma-perfection}
设 $k$ 是完美闭包为 $k^{perf}$ 的域。设 $X$ 是 $k$ 上的连通概形。
则 $X_{k^{perf}}$ 连通，且
$\pi_1(X_{k^{perf}}) \to \pi_1(X)$ 是同构。
\end{lemma}

\begin{proof}
这是基本群拓扑不变性的一个特例，见 Proposition
\ref{pione-proposition-universal-homeomorphism}。要看出
$\Spec(k^{perf}) \to \Spec(k)$ 是泛同胚，可以使用 Algebra, Lemma
\ref{algebra-lemma-radicial-integral-bijective}。
\end{proof}

\begin{lemma}
\label{pione-lemma-ses-field}
设 $k$ 是代数闭包为 $\overline{k}$ 的域。设 $X$ 是 $k$ 上拟紧且拟分离的概形。
若基变换 $X_{\overline{k}}$ 连通，则有短正合列
$$
1 \to \pi_1(X_{\overline{k}}) \to \pi_1(X) \to \pi_1(\Spec(k)) \to 1
$$
其中各项均为 profinite 拓扑群。
\end{lemma}

\begin{proof}
$\textit{F\'Et}_{\Spec(k)}$ 的连通对象形如
$\Spec(k') \to \Spec(k)$，其中 $k'/k$ 是有限可分扩张。于是
$X_{\Spec{k'}}$ 连通，因为态射 $X_{\overline{k}} \to X_{\Spec(k')}$ 是满射，
且按假设 $X_{\overline{k}}$ 连通。因此，由 Lemma
\ref{pione-lemma-functoriality-galois-surjective}，
$\pi_1(X) \to \pi_1(\Spec(k))$ 是满射。

\medskip\noindent
继续之前，注意可以假设 $k$ 是完美域。事实上，由 Lemma \ref{pione-lemma-perfection} 有
$\pi_1(X_{k^{perf}}) = \pi_1(X)$ 且
$\pi_1(\Spec(k^{perf})) = \pi_1(\Spec(k))$。

\medskip\noindent
显然，函子的复合
$\textit{F\'Et}_{\Spec(k)} \to \textit{F\'Et}_X \to
\textit{F\'Et}_{X_{\overline{k}}}$ 将对象送到 $X_{\Spec(\overline{k})}$ 的若干副本的不交并。
因此复合
$\pi_1(X_{\overline{k}}) \to \pi_1(X) \to \pi_1(\Spec(k))$
由 Lemma \ref{pione-lemma-composition-trivial} 可知是平凡同态。

\medskip\noindent
设 $U \to X$ 是有限 \'etale 态射，且 $U$ 连通。注意
$U \times_X X_{\overline{k}} = U_{\overline{k}}$。假设
$U_{\overline{k}} \to X_{\overline{k}}$ 有一个截面
$s : X_{\overline{k}} \to U_{\overline{k}}$。于是 $s(X_{\overline{k}})$ 是
$U_{\overline{k}}$ 的一个开连通分支。对
$\sigma \in \text{Gal}(\overline{k}/k)$，以 $s^\sigma$ 表示 $s$ 沿
$\Spec(\sigma)$ 的基变换。由于 $U_{\overline{k}} \to X_{\overline{k}}$ 是有限
\'etale 态射，它只有有限多个截面。因此
$$
\overline{T} = \bigcup s^\sigma(X_{\overline{k}})
$$
是有限并，且 $\overline{T}$ 是一个在
$\text{Gal}(\overline{k}/k)$ 作用下稳定的开闭子集。由 Varieties, Lemma
\ref{varieties-lemma-closed-fixed-by-Galois} 可知，$\overline{T}$ 是某个闭子集
$T \subset U$ 的原像。由于 $U_{\overline{k}} \to U$ 是开映射
（Morphisms, Lemma \ref{morphisms-lemma-scheme-over-field-universally-open}），
$T$ 也开。又因 $U$ 连通，故 $T = U$。因此 $U_{\overline{k}}$ 是
$X_{\overline{k}}$ 的若干副本的（有限）不交并。由 Lemma
\ref{pione-lemma-functoriality-galois-normal} 可知
$\pi_1(X_{\overline{k}}) \to \pi_1(X)$ 的像是正规子群。

\medskip\noindent
设 $V \to X_{\overline{k}}$ 是有限 \'etale 覆盖。回顾 $\overline{k}$ 是 $k$
的有限可分扩张之并。由 Lemma \ref{pione-lemma-limit}，可找到有限可分扩张 $k'/k$
和有限 \'etale 态射 $U \to X_{k'}$，使得
$V = X_{\overline{k}} \times_{X_{k'}} U =
U \times_{\Spec(k')} \Spec(\overline{k})$.
于是复合 $U \to X_{k'} \to X$ 是有限 \'etale 态射，并且
$U \times_{\Spec(k)} \Spec(\overline{k})$ 包含
$V = U \times_{\Spec(k')} \Spec(\overline{k})$ 作为开闭子概形。（这是因为
$\Spec(\overline{k})$ 是
$\Spec(k') \times_{\Spec(k)} \Spec(\overline{k})$ 的一个开闭子概形，
它由乘法映射 $k' \otimes_k \overline{k} \to \overline{k}$ 给出。）
由 Lemma \ref{pione-lemma-functoriality-galois-injective} 可知
$\pi_1(X_{\overline{k}}) \to \pi_1(X)$ 是单射。

\medskip\noindent
最后，必须证明：对任意有限 \'etale 态射 $U \to X$，若
$U_{\overline{k}}$ 是 $X_{\overline{k}}$ 的若干副本的不交并，则存在有限
\'etale 态射 $V \to \Spec(k)$ 和满射 $V \times_{\Spec(k)} X \to U$。
见 Lemma \ref{pione-lemma-functoriality-galois-ses}。像上面一样应用 Lemma
\ref{pione-lemma-limit}，可找到有限可分扩张 $k'/k$，使得存在同构
$U_{k'} \cong \coprod_{i = 1, \ldots, n} X_{k'}$.
因此令 $V = \coprod_{i = 1, \ldots, n} \Spec(k')$ 即得结论。
\end{proof}






\section{同伦正合列}
\label{pione-section-homotopy-exact-sequence}

\noindent
本节讨论如下结果。设 $f : X \to S$ 是有限表示的平坦且固有的态射，
其几何纤维均连通且约化。假设 $S$ 连通，并设 $\overline{s}$ 是 $S$
的一个几何点。则有基本群的正合列
$$
\pi_1(X_{\overline{s}}) \to \pi_1(X) \to \pi_1(S) \to 1
$$
见 Proposition \ref{pione-proposition-first-homotopy-sequence}。

\begin{lemma}
\label{pione-lemma-stein-factorization-etale}
\begin{reference}
\cite[Expose X, Proposition 1.2, p. 262]{SGA1}.
\end{reference}
设 $f : X \to S$ 是概形的固有态射。设 $X \to S' \to S$ 是 $f$ 的 Stein 分解，
见 More on Morphisms, Theorem
\ref{more-morphisms-theorem-stein-factorization-general}。若 $f$ 有限表示且平坦，
并具有几何约化纤维，则 $S' \to S$ 是有限 \'etale 态射。
\end{lemma}

\begin{proof}
这由 Derived Categories of Schemes, Lemma
\ref{perfect-lemma-proper-flat-geom-red} 以及 More on Morphisms, Theorem
\ref{more-morphisms-theorem-stein-factorization-general} 中的信息得出。
\end{proof}

\begin{proposition}
\label{pione-proposition-first-homotopy-sequence}
设 $f : X \to S$ 是有限表示的平坦且固有的态射，其几何纤维均连通且约化。
假设 $S$ 连通，并设 $\overline{s}$ 是 $S$ 的一个几何点。则有基本群的正合列
$$
\pi_1(X_{\overline{s}}) \to \pi_1(X) \to \pi_1(S) \to 1
$$
\end{proposition}

\begin{proof}
设 $Y \to X$ 是有限 \'etale 态射。考虑 Stein 分解
$$
\xymatrix{
Y \ar[d] \ar[r] & X \ar[d] \\
T \ar[r] & S
}
$$
它是 $Y \to S$ 的 Stein 分解。由 Lemma
\ref{pione-lemma-stein-factorization-etale}，态射 $T \to S$ 是有限 \'etale 的。
由此得到函子 $\textit{F\'Et}_X \to \textit{F\'Et}_S$。对任意有限 \'etale 态射
$U \to S$，态射 $Y \to U \times_S X$ 在 $X$ 上给出，当且仅当态射
$Y \to U$ 在 $S$ 上给出；这样的态射唯一地经过 Stein 分解，亦即对应于唯一的态射
$T \to U$
（这是因为 More on Morphisms, Lemma
\ref{more-morphisms-lemma-stein-universally-closed} 将 Stein 分解构造成相对正规化，
再应用 Morphisms, Lemma
\ref{morphisms-lemma-characterize-normalization} 的分解性质）。因此函子
$\textit{F\'Et}_X \to \textit{F\'Et}_S$ 与
$\textit{F\'Et}_S \to \textit{F\'Et}_X$ 互为伴随。注意，
$U \times_S X \to S$ 的 Stein 分解是 $U$，因为 $U \times_S X \to U$
的各纤维都几何连通。

\medskip\noindent
由上述讨论和 Categories, Lemma
\ref{categories-lemma-adjoint-fully-faithful}，可知
$\textit{F\'Et}_S \to \textit{F\'Et}_X$ 全忠实；也就是说，
$\pi_1(X) \to \pi_1(S)$ 是满射（Lemma
\ref{pione-lemma-functoriality-galois-surjective}）。

\medskip\noindent
显然，复合
$\textit{F\'Et}_S \to \textit{F\'Et}_X \to \textit{F\'Et}_{X_{\overline{s}}}$
将任意 $U$ 送到 $X_{\overline{s}}$ 的若干副本的不交并。因此由 Lemma
\ref{pione-lemma-composition-trivial}，
$\pi_1(X_{\overline{s}}) \to \pi_1(X) \to \pi_1(S)$ 是平凡同态。

\medskip\noindent
设 $Y \to X$ 是有限 \'etale 态射，$Y$ 连通，并且
$Y \times_X X_{\overline{s}}$ 含有一个连通分支 $Z$，它同构于
$X_{\overline{s}}$。考虑上述 Stein 分解 $T$。设
$\overline{t} \in T_{\overline{s}}$ 是与纤维 $Z$ 对应的点。注意 $T$ 连通
（因为它是连通概形的像），而且由上面的满射性，$T \times_S X$ 连通。
现在考虑分解
$$
\pi : Y \longrightarrow T \times_S X
$$
设 $\overline{x} \in X_{\overline{s}}$ 是任意闭点。注意
$\kappa(\overline{t}) = \kappa(\overline{s}) = \kappa(\overline{x})$
是代数闭域。于是 $\pi$ 在 $(\overline{t}, \overline{x})$ 上的纤维仅含一个点，
即唯一的点 $\overline{z} \in Z$；它对应于
$\overline{x} \in X_{\overline{s}}$，这一对应通过同构
$Z \to X_{\overline{s}}$ 给出。因此有限 \'etale 态射 $\pi$ 的次数为 $1$，
这是在 $(\overline{t}, \overline{x})$ 的一个邻域上。由于
$T \times_S X$ 连通，该态射处处次数为 $1$，从而 $Y \cong T \times_S X$。
所以 $Y \times_X X_{\overline{s}}$ 完全分裂。综合以上各点，Lemmas
\ref{pione-lemma-functoriality-galois-ses} 和
\ref{pione-lemma-functoriality-galois-normal} 均可应用，证明完成。
\end{proof}




\section{特化映射}
\label{pione-section-specialization-map}

\noindent
本节构造特化映射。设 $f : X \to S$ 是具有几何连通纤维的概形固有态射。
设 $s' \leadsto s$ 是 $S$ 中点的特化。设 $\overline{s}$ 和 $\overline{s}'$
分别是位于 $s$ 和 $s'$ 上方的几何点。则有特化映射
$$
sp : \pi_1(X_{\overline{s}'}) \longrightarrow \pi_1(X_{\overline{s}})
$$
此映射的构造如下。令 $A$ 为 $\mathcal{O}_{S, s}$ 关于
$\kappa(s) \subset \kappa(s)^{sep} \subset \kappa(\overline{s})$ 的严格
Hensel 化，见 Algebra, Definition \ref{algebra-definition-henselization}。
因为 $s' \leadsto s$，点 $s'$ 对应于 $\Spec(\mathcal{O}_{S, s})$ 的一个点，
故 $\Spec(A)$ 中至少有一个（也可能有许多个）位于 $s'$ 上方的点，
其剩余域是 $\kappa(s')$ 的可分代数扩张。由于
$\kappa(\overline{s}')$ 代数闭，可以选取一个态射
$\varphi : \overline{s}' \to \Spec(A)$，从而得到交换图
$$
\xymatrix{
\overline{s}' \ar[r]_-\varphi \ar[rd] &
\Spec(A) \ar[d] &
\overline{s} \ar[l] \ar[ld] \\
& S
}
$$
特化映射是复合
$$
\pi_1(X_{\overline{s}'}) \longrightarrow
\pi_1(X_A) =
\pi_1(X_{\kappa(s)^{sep}}) =
\pi_1(X_{\overline{s}})
$$
其中第一个等号由 Lemma \ref{pione-lemma-finite-etale-on-proper-over-henselian} 给出，
第二个等号由 Lemmas \ref{pione-lemma-perfection} 和
\ref{pione-lemma-finite-etale-invariant-over-proper} 得出。按构造，特化映射属于交换图
$$
\xymatrix{
\pi_1(X_{\overline{s}'}) \ar[rr]_{sp} \ar[rd] & &
\pi_1(X_{\overline{s}}) \ar[ld] \\
& \pi_1(X)
}
$$
只要 $X$ 连通。特化映射依赖于上面对
$\varphi : \overline{s}' \to \Spec(A)$ 的选择；若要表明这一点，便记作 $sp_\varphi$。

\begin{lemma}
\label{pione-lemma-specialization-map-base-change}
考虑概形的交换图
$$
\xymatrix{
Y \ar[d]_g \ar[r] & X \ar[d]^f \\
T \ar[r] & S
}
$$
其中 $f$ 和 $g$ 都是具有几何连通纤维的固有态射。设 $t' \leadsto t$ 是 $T$
中点的特化，并考虑上述特化映射
$sp : \pi_1(Y_{\overline{t}'}) \to \pi_1(Y_{\overline{t}})$。
则有特化映射的交换图
$$
\xymatrix{
\pi_1(Y_{\overline{t}'}) \ar[r]_{sp} \ar[d] & \pi_1(Y_{\overline{t}}) \ar[d] \\
\pi_1(X_{\overline{s}'}) \ar[r]^{sp} & \pi_1(X_{\overline{s}})
}
$$
其中 $\overline{s}$ 和 $\overline{s}'$ 分别是 $\overline{t}$ 和
$\overline{t}'$ 的像。
\end{lemma}

\begin{proof}
令 $B$ 为 $\mathcal{O}_{T, t}$ 关于
$\kappa(t) \subset \kappa(t)^{sep} \subset \kappa(\overline{t})$ 的严格 Hensel 化。
像构造特化映射时那样，选取
$\psi : \overline{t}' \to \Spec(B)$，它提升 $\overline{t}' \to T$。
以 $s$ 和 $s'$ 表示 $t$ 和 $t'$ 在 $S$ 中的像。
令 $A$ 为 $\mathcal{O}_{S, s}$ 关于
$\kappa(s) \subset \kappa(s)^{sep} \subset \kappa(\overline{s})$ 的严格 Hensel 化。
由于 $\kappa(\overline{s}) = \kappa(\overline{t})$，由严格 Hensel 化的函子性
（Algebra, Lemma \ref{algebra-lemma-strictly-henselian-functorial}），
得到属于下述交换图的环同态 $A \to B$
$$
\xymatrix{
\overline{t}' \ar[r]_-\psi \ar[d] & \Spec(B) \ar[d] \ar[r] & T \ar[d] \\
\overline{s}' \ar[r]^-\varphi & \Spec(A) \ar[r] & S
}
$$
这里直接取态射 $\varphi : \overline{s}' \to \Spec(A)$ 为复合
$\overline{t}' \to \Spec(B) \to \Spec(A)$。施行基变换得到交换图
$$
\xymatrix{
Y_{\overline{t}'} \ar[r] \ar[d] & Y_B \ar[d] \\
X_{\overline{s}'} \ar[r] & X_A
}
$$
由特化映射的构造，此图的交换性蕴含引理中交换图的交换性。
\end{proof}

\begin{lemma}
\label{pione-lemma-specialization-map-composition}
设 $f : X \to S$ 是具有几何连通纤维的固有态射。设
$s'' \leadsto s' \leadsto s$ 是 $S$ 中点的特化。特化映射的复合
$\pi_1(X_{\overline{s}''}) \to \pi_1(X_{\overline{s}'}) \to
\pi_1(X_{\overline{s}})$ 是特化映射
$\pi_1(X_{\overline{s}''}) \to \pi_1(X_{\overline{s}})$.
\end{lemma}

\begin{proof}
设 $\mathcal{O}_{S, s} \to A$ 是利用
$\kappa(s) \to \kappa(\overline{s})$ 构造的严格 Hensel 化。设
$A \to \kappa(\overline{s}')$ 是构造第一个特化映射时所用的映射。设
$\mathcal{O}_{S, s'} \to A'$ 是利用
$\kappa(s') \subset \kappa(\overline{s}')$ 构造的严格 Hensel 化。
由严格 Hensel 化的函子性，存在映射 $A \to A'$，使它同
$A' \to \kappa(\overline{s}')$ 的复合就是给定映射
（Algebra, Lemma \ref{algebra-lemma-map-into-henselian-colimit}）。
接着，设 $A' \to \kappa(\overline{s}'')$ 是构造第二个特化映射时所用的映射。
显然，第一与第二个特化映射的复合就是特化映射
$\pi_1(X_{\overline{s}''}) \to \pi_1(X_{\overline{s}})$
，它利用 $A \to A' \to \kappa(\overline{s}'')$ 构造。
\end{proof}

\noindent
设 $X \to S$ 是具有几何连通纤维的固有态射。设 $R$ 是分式域代数闭的严格
Hensel 赋值环，并设 $\Spec(R) \to S$ 是态射。设
$\eta, s \in \Spec(R)$ 分别为一般点和闭点。于是可以考虑特化映射
$$
sp_R : \pi_1(X_\eta) \to \pi_1(X_s)
$$
它属于基变换 $X_R/\Spec(R)$。注意，这是有意义的，因为 $\eta$ 和 $s$
的剩余域均代数闭。

\begin{lemma}
\label{pione-lemma-specialization-map-valuation-ring}
设 $f : X \to S$ 是具有几何连通纤维的固有态射。设 $s' \leadsto s$ 是 $S$
中点的特化，并设
$sp : \pi_1(X_{\overline{s}'}) \to \pi_1(X_{\overline{s}})$
是特化映射。则存在分式域代数闭的严格 Hensel 赋值环 $R$，它位于 $S$ 上，
使 $sp$ 同构于上面定义的 $sp_R$。
\end{lemma}

\begin{proof}
设 $\mathcal{O}_{S, s} \to A$ 是利用
$\kappa(s) \to \kappa(\overline{s})$ 构造的严格 Hensel 化。设
$A \to \kappa(\overline{s}')$ 是构造 $sp$ 时所用的映射。设
$R \subset \kappa(\overline{s}')$ 是分式域为 $\kappa(\overline{s}')$、
且支配 $A$ 之像的赋值环。见 Algebra, Lemma \ref{algebra-lemma-dominate}。
例如，由 Lemma \ref{pione-lemma-normal-local-domain-separablly-closed-fraction-field}
和 Algebra, Lemma \ref{algebra-lemma-valuation-ring-normal} 可知 $R$ 是严格
Hensel 环。引理随即成立。
\end{proof}

\noindent
设 $X \to S$ 是具有几何连通纤维的固有态射。设 $R$ 是严格 Hensel 离散赋值环，
并设 $\Spec(R) \to S$ 是态射。设 $\eta, s \in \Spec(R)$ 分别为一般点和闭点。
于是可以考虑特化映射
$$
sp_R : \pi_1(X_{\overline{\eta}}) \to \pi_1(X_s)
$$
它属于基变换 $X_R/\Spec(R)$。注意，这是有意义的，因为 $s$ 的剩余域代数闭。

\begin{lemma}
\label{pione-lemma-specialization-map-discrete-valuation-ring}
设 $f : X \to S$ 是具有几何连通纤维的固有态射。设 $s' \leadsto s$ 是 $S$
中点的特化，并设
$sp : \pi_1(X_{\overline{s}'}) \to \pi_1(X_{\overline{s}})$
是特化映射。若 $S$ 是 Noether 概形，则存在严格 Hensel 离散赋值环 $R$，
它位于 $S$ 上，使得 $sp$ 同构于上面定义的 $sp_R$。
\end{lemma}

\begin{proof}
设 $\mathcal{O}_{S, s} \to A$ 是利用
$\kappa(s) \to \kappa(\overline{s})$ 构造的严格 Hensel 化。设
$A \to \kappa(\overline{s}')$ 是构造 $sp$ 时所用的映射。设
$R \subset \kappa(\overline{s}')$ 是支配 $A$ 之像的离散赋值环，见 Algebra, Lemma
\ref{algebra-lemma-exists-dvr}。选取域的交换图
$$
\xymatrix{
\kappa(\overline{s}) \ar[r] & k \\
A/\mathfrak m_A \ar[r] \ar[u] & R/\mathfrak m_R \ar[u]
}
$$
其中 $k$ 代数闭。设 $R^{sh}$ 是 $R$ 的严格 Hensel 化，它利用 $R \to k$ 构造。
由 Algebra, Lemma \ref{algebra-lemma-strictly-henselian-functorial}，
得到映射 $A \to R^{sh}$。由 More on Algebra, Lemma
\ref{more-algebra-lemma-henselization-dvr}，环 $R^{sh}$ 是离散赋值环。
以 $\eta, o$ 表示 $\Spec(R^{sh})$ 的一般点和闭点。因为概形图
$$
\xymatrix{
\overline{\eta} \ar[d] \ar[r] & \Spec(R^{sh}) \ar[d] &
\Spec(k) \ar[d] \ar[l] \\
\overline{s}' \ar[r] & \Spec(A) & \overline{s} \ar[l]
}
$$
交换，所以得到交换图
$$
\xymatrix{
\pi_1(X_{\overline{\eta}}) \ar[d] \ar[r]_{sp_{R^{sh}}} & \pi_1(X_o) \ar[d] \\
\pi_1(X_{\overline{s}'}) \ar[r]^{sp} & X_{\overline{s}}
}
$$
按这些映射的构造，这是特化映射的交换图。由于竖直箭头均为同构
（Lemma \ref{pione-lemma-finite-etale-invariant-over-proper}），引理得证。
\end{proof}








\section{限制到闭子概形}
\label{pione-section-lefschetz}

\noindent
本节证明关于限制函子
$$
\textit{F\'Et}_X \longrightarrow \textit{F\'Et}_Y,\quad
U \longmapsto V = U \times_X Y
$$
的一些结果，其中 $X$ 是概形，而 $Y$ 是闭子概形。利用基本群的拓扑不变性，
可以把对此函子的研究同有限局部自由模上的完备化函子联系起来。

\medskip\noindent
以下引理使用 Cohomology of Schemes, Section
\ref{coherent-section-coherent-formal} 中定义的凝聚形式模概念。给定 Noether 概形和拟凝聚理想层
$\mathcal{I} \subset \mathcal{O}_X$，若对象 $(\mathcal{F}_n)$ 属于
$\textit{Coh}(X, \mathcal{I})$，且每个 $\mathcal{F}_n$ 都是有限局部自由
$\mathcal{O}_X/\mathcal{I}^n$-模，就称它{\it 有限局部自由}。

\begin{lemma}
\label{pione-lemma-restriction-fully-faithful}
设 $X$ 是 Noether 概形，$Y \subset X$ 是理想层为
$\mathcal{I} \subset \mathcal{O}_X$ 的闭子概形。假设完备化函子
$$
\textit{Coh}(\mathcal{O}_X)
\longrightarrow 
\textit{Coh}(X, \mathcal{I}),\quad
\mathcal{F} \longmapsto \mathcal{F}^\wedge
$$
在有限局部自由对象的全子范畴上全忠实（见上文）。则限制函子
$\textit{F\'Et}_X \to \textit{F\'Et}_Y$ 全忠实。
\end{lemma}

\begin{proof}
由于有限 \'etale 覆盖的范畴具有内部 Hom（Lemma
\ref{pione-lemma-internal-hom-finite-etale}），只需证明如下断言：给定 $U$，它是 $X$
上的有限 \'etale 概形，并给定态射 $t : Y \to U$，它在 $X$ 上，存在唯一的截面
$s : X \to U$ 使得 $t = s|_Y$。图示为
$$
\xymatrix{
& U \ar[d]^f \\
Y \ar[r] \ar[ru] & X \ar@{..>}@/^1em/[u]
}
$$
寻找虚线箭头 $s$ 等价于寻找 $\mathcal{O}_X$-代数映射
$$
s^\sharp : f_*\mathcal{O}_U \longrightarrow \mathcal{O}_X
$$
它模掉 $Y$ 的理想层后化为给定的代数映射
$t^\sharp : f_*\mathcal{O}_U \to \mathcal{O}_Y$。由 Lemma
\ref{pione-lemma-thickening}，可以将 $t$ 唯一提升为相容映射系统
$t_n : Y_n \to U$，从而得到映射
$$
\lim t_n^\sharp : f_*\mathcal{O}_U \longrightarrow \lim \mathcal{O}_{Y_n}
$$
它是 $X$ 上代数层的映射。由于 $f_*\mathcal{O}_U$ 是有限局部自由
$\mathcal{O}_X$-模，可知存在唯一的 $\mathcal{O}_X$-模映射
$\sigma : f_*\mathcal{O}_U \to \mathcal{O}_X$，其完备化为
$\lim t_n^\sharp$。为证明 $\sigma$ 是代数同态，需要检验交换图
$$
\xymatrix{
f_*\mathcal{O}_U \otimes_{\mathcal{O}_X} f_*\mathcal{O}_U
\ar[r] \ar[d]_{\sigma \otimes \sigma} &
f_*\mathcal{O}_U \ar[d]^\sigma \\
\mathcal{O}_X \otimes_{\mathcal{O}_X} \mathcal{O}_X \ar[r] &
\mathcal{O}_X
}
$$
成立。对每个 $n$，已知此图限制到 $Y_n$ 后交换；也就是说，施行完备化函子后该图交换。
因此由完备化函子的忠实性得出结论。
\end{proof}

\begin{lemma}
\label{pione-lemma-restriction-equivalence}
设 $X$ 是 Noether 概形，$Y \subset X$ 是理想层为
$\mathcal{I} \subset \mathcal{O}_X$ 的闭子概形。假设完备化函子
$$
\textit{Coh}(\mathcal{O}_X)
\longrightarrow 
\textit{Coh}(X, \mathcal{I}),\quad
\mathcal{F} \longmapsto \mathcal{F}^\wedge
$$
在有限局部自由对象的全子范畴上给出等价（见上文）。则限制函子
$\textit{F\'Et}_X \to \textit{F\'Et}_Y$ 是等价。
\end{lemma}

\begin{proof}
由 Lemma \ref{pione-lemma-restriction-fully-faithful}，限制函子全忠实。

\medskip\noindent
设 $U_1 \to Y$ 是有限 \'etale 态射。为完成证明，我们将说明 $U_1$
属于限制函子的本质像。

\medskip\noindent
对 $n \geq 1$，设 $Y_n$ 为第 $n$ 个 $Y$ 的无穷小邻域。由 Lemma
\ref{pione-lemma-thickening}，存在唯一的有限 \'etale 态射
$\pi_n : U_n \to Y_n$，其到 $Y = Y_1$ 的基变换恢复 $U_1 \to Y_1$。
考虑层 $\mathcal{F}_n = \pi_{n, *}\mathcal{O}_{U_n}$。可以并且将
$\mathcal{F}_n$ 视为一个 $\mathcal{O}_X$-模，它位于 $X$ 上；该模局部同构于
$(\mathcal{O}_X/f^{n + 1}\mathcal{O}_X)^{\oplus r}$。这个
$(\mathcal{F}_n)$ 是 $\textit{Coh}(X, \mathcal{I})$ 中的有限局部自由对象。
按假设，存在有限局部自由 $\mathcal{O}_X$-模 $\mathcal{F}$ 和相容同构系统
$$
\mathcal{F}/\mathcal{I}^n\mathcal{F} \to \mathcal{F}_n
$$
，它们是 $\mathcal{O}_X$-模同构。

\medskip\noindent
为在 $\mathcal{F}$ 上构造代数结构，考虑乘法映射
$\mathcal{F}_n \otimes_{\mathcal{O}_X} \mathcal{F}_n \to \mathcal{F}_n$
；它们来自 $\mathcal{F}_n = \pi_{n, *}\mathcal{O}_{U_n}$ 是代数层这一事实。
这些映射定义
$$
(\mathcal{F}\otimes_{\mathcal{O}_X} \mathcal{F})^\wedge
\longrightarrow
\mathcal{F}^\wedge
$$
，它是范畴 $\textit{Coh}(X, \mathcal{I})$ 中的映射。因此按假设，可以设存在映射
$\mu : \mathcal{F}\otimes_{\mathcal{O}_X} \mathcal{F} \to \mathcal{F}$
，它限制到 $Y_n$ 后给出上述乘法映射。由引理所述函子的忠实性，可知
$\mu$ 定义了交换 $\mathcal{O}_X$-代数结构，其底层模为 $\mathcal{F}$，并与
$\mathcal{F}_n$ 上给定的代数结构相容。令
$$
U = \underline{\Spec}_X((\mathcal{F}, \mu))
$$
得到有限局部自由概形 $\pi : U \to X$，其到 $Y$ 的限制同构于 $U_1$。
$\pi$ 的判别式是截面
$$
\det(Q_\pi) :
\mathcal{O}_X
\longrightarrow
\wedge^{top}(\pi_*\mathcal{O}_U)^{\otimes -2}
$$
的零点集；该截面构造于 Discriminants, Section
\ref{discriminant-section-discriminant}。由 Discriminants, Lemma
\ref{discriminant-lemma-discriminant}，它到 $Y_n$ 的限制对所有 $n$
都是同构，故它自身也是同构。再由 Discriminants, Lemma
\ref{discriminant-lemma-discriminant}，$\pi$ 是 \'etale 的。
\end{proof}

\begin{lemma}
\label{pione-lemma-restriction-fully-faithful-general}
设 $X$ 是 Noether 概形，$Y \subset X$ 是理想层为
$\mathcal{I} \subset \mathcal{O}_X$ 的闭子概形。令 $\mathcal{V}$ 为开子概形
$V \subset X$ 所成的集合，其中每个都包含 $Y$；以反包含关系排序。假设完备化函子
$$
\colim_\mathcal{V} \textit{Coh}(\mathcal{O}_V)
\longrightarrow
\textit{Coh}(X, \mathcal{I}),
\quad
\mathcal{F} \longmapsto \mathcal{F}^\wedge
$$
在有限局部自由对象所成的全子范畴上全忠实（见上文）。则限制函子
$\colim_\mathcal{V} \textit{F\'Et}_V \to \textit{F\'Et}_Y$
全忠实。
\end{lemma}

\begin{proof}
注意 $\mathcal{V}$ 是有向集，故这里的余极限具有 Categories, Section
\ref{categories-section-directed-colimits} 中所述的形式。其余论证几乎与 Lemma
\ref{pione-lemma-restriction-fully-faithful} 的证明完全相同；建议读者略去。

\medskip\noindent
由于有限 \'etale 覆盖的范畴具有内部 Hom（Lemma
\ref{pione-lemma-internal-hom-finite-etale}），只需证明如下断言：给定 $U$，它在
$V \in \mathcal{V}$ 上有限 \'etale，并给定态射 $t : Y \to U$，它在 $V$ 上；
存在 $V' \geq V$ 以及态射 $s : V' \to U$，后者在 $V$ 上，并且 $t = s|_Y$。图示为
$$
\xymatrix{
& & U \ar[d]^f \\
Y \ar[r] \ar[rru] & V' \ar@{..>}[ru] \ar[r] & V
}
$$
寻找虚线箭头 $s$ 等价于寻找 $\mathcal{O}_{V'}$-代数映射
$$
s^\sharp : f_*\mathcal{O}_U|_{V'} \longrightarrow \mathcal{O}_{V'}
$$
，它模掉 $Y$ 的理想层后化为给定的代数映射
$t^\sharp : f_*\mathcal{O}_U \to \mathcal{O}_Y$。由 Lemma
\ref{pione-lemma-thickening}，可以将 $t$ 唯一提升为相容映射系统
$t_n : Y_n \to U$，从而得到映射
$$
\lim t_n^\sharp : f_*\mathcal{O}_U \longrightarrow \lim \mathcal{O}_{Y_n}
$$
，它是 $V$ 上代数层的映射。注意 $f_*\mathcal{O}_U$ 是有限局部自由
$\mathcal{O}_V$-模。因此可取 $V' \geq V$ 以及映射
$\sigma : f_*\mathcal{O}_U|_{V'} \to \mathcal{O}_{V'}$
，其完备化为 $\lim t_n^\sharp$。为证明 $\sigma$ 是代数同态，需要检验交换图
$$
\xymatrix{
(f_*\mathcal{O}_U \otimes_{\mathcal{O}_V} f_*\mathcal{O}_U)|_{V'}
\ar[r] \ar[d]_{\sigma \otimes \sigma} &
f_*\mathcal{O}_U|_{V'} \ar[d]^\sigma \\
\mathcal{O}_{V'} \otimes_{\mathcal{O}_{V'}} \mathcal{O}_{V'} \ar[r] &
\mathcal{O}_{V'}
}
$$
成立。对每个 $n$，已知此图限制到 $Y_n$ 后交换；也就是说，施行完备化函子后该图交换。
因此由完备化函子的忠实性可知，存在 $V'' \geq V'$，使得
$\sigma|_{V''}$ 如所需是代数同态。
\end{proof}

\begin{lemma}
\label{pione-lemma-restriction-equivalence-general}
设 $X$ 是 Noether 概形，$Y \subset X$ 是理想层为
$\mathcal{I} \subset \mathcal{O}_X$ 的闭子概形。令 $\mathcal{V}$ 为开子概形
$V \subset X$ 所成的集合，其中每个都包含 $Y$；以反包含关系排序。假设完备化函子
$$
\colim_\mathcal{V} \textit{Coh}(\mathcal{O}_V)
\longrightarrow
\textit{Coh}(X, \mathcal{I}),
\quad
\mathcal{F} \longmapsto \mathcal{F}^\wedge
$$
在有限局部自由对象所成的全子范畴之间给出等价（见上文说明）。则限制函子
$$
\colim_\mathcal{V} \textit{F\'Et}_V \to \textit{F\'Et}_Y
$$
是等价。
\end{lemma}

\begin{proof}
注意 $\mathcal{V}$ 是有向集，故这里的余极限具有 Categories, Section
\ref{categories-section-directed-colimits} 中所述的形式。其余论证几乎与 Lemma
\ref{pione-lemma-restriction-equivalence} 的证明完全相同；建议读者略去。

\medskip\noindent
由 Lemma \ref{pione-lemma-restriction-fully-faithful-general}，限制函子全忠实。

\medskip\noindent
设 $U_1 \to Y$ 是有限 \'etale 态射。为完成证明，我们将说明 $U_1$
属于限制函子的本质像。

\medskip\noindent
对 $n \geq 1$，设 $Y_n$ 为第 $n$ 个 $Y$ 的无穷小邻域。由 Lemma
\ref{pione-lemma-thickening}，存在唯一的有限 \'etale 态射
$\pi_n : U_n \to Y_n$，其到 $Y = Y_1$ 的基变换恢复 $U_1 \to Y_1$。
考虑层 $\mathcal{F}_n = \pi_{n, *}\mathcal{O}_{U_n}$。可以并且将
$\mathcal{F}_n$ 视为一个 $\mathcal{O}_X$-模，它位于 $X$ 上；该模局部同构于
$(\mathcal{O}_X/f^{n + 1}\mathcal{O}_X)^{\oplus r}$.
这个 $(\mathcal{F}_n)$ 是 $\textit{Coh}(X, \mathcal{I})$ 中的有限局部自由对象。
按假设，存在 $V \in \mathcal{V}$、有限局部自由 $\mathcal{O}_V$-模
$\mathcal{F}$，以及相容同构系统
$$
\mathcal{F}/\mathcal{I}^n\mathcal{F} \to \mathcal{F}_n
$$
，它们是 $\mathcal{O}_V$-模同构。

\medskip\noindent
为在 $\mathcal{F}$ 上构造代数结构，考虑乘法映射
$\mathcal{F}_n \otimes_{\mathcal{O}_V} \mathcal{F}_n \to \mathcal{F}_n$
；它们来自 $\mathcal{F}_n = \pi_{n, *}\mathcal{O}_{U_n}$ 是代数层这一事实。
这些映射定义
$$
(\mathcal{F}\otimes_{\mathcal{O}_V} \mathcal{F})^\wedge
\longrightarrow
\mathcal{F}^\wedge
$$
，它是范畴 $\textit{Coh}(X, \mathcal{I})$ 中的映射。因此按假设，缩小
$V$ 后可假设存在映射
$\mu : \mathcal{F}\otimes_{\mathcal{O}_V} \mathcal{F} \to \mathcal{F}$
，它限制到 $Y_n$ 后给出上述乘法映射。必要时进一步缩小后，可假设 $\mu$
定义了交换 $\mathcal{O}_V$-代数结构，其底层模为 $\mathcal{F}$，并与
$\mathcal{F}_n$ 上给定的代数结构相容。令
$$
U = \underline{\Spec}_V((\mathcal{F}, \mu))
$$
，得到 $V$ 上的有限局部自由概形，其到 $Y$ 的限制同构于 $U_1$。由此
$U \to V$ 在位于 $Y$ 上方的所有点处都是 \'etale 的；参见 More on Morphisms,
Lemma \ref{more-morphisms-lemma-check-smoothness-on-infinitesimal-nbhds}。
因此再次缩小 $V$ 后，可假设 $U \to V$ 有限 \'etale。证明完成。
\end{proof}

\begin{lemma}
\label{pione-lemma-restriction-faithful}
设 $X$ 是概形，$Y \subset X$ 是闭子概形。若 $X$ 的每个连通分支都与
$Y$ 相交，则限制函子 $\textit{F\'Et}_X \to \textit{F\'Et}_Y$ 忠实。
\end{lemma}

\begin{proof}
设 $a, b : U \to U'$ 是两个 $X$ 上有限 \'etale 概形之间的态射，并且它们到
$Y$ 的限制相同。$U$ 的一个连通分支的像是 $X$ 的一个连通分支；这是将
Topology, Lemma \ref{topology-lemma-finite-fibre-connected-components}
应用于 $U \to X$ 到 $X$ 的一个连通分支上的限制所得。因此按假设，$U$
每个连通分支的像都与 $Y$ 相交。由 \'Etale Morphisms, Proposition
\ref{etale-proposition-equality}，可知 $a = b$ 在 $U$ 的每个连通分支上的限制成立。
由于 $a$ 与 $b$ 的等化子是 $U$ 的开子概形（因为 $U'$ 在 $X$ 上的对角是开的），
结论随即成立。
\end{proof}

\begin{lemma}
\label{pione-lemma-restriction-fully-faithful-special}
设 $X$ 是 Noether 概形，$Y \subset X$ 是闭子概形。令 $Y_n \subset X$ 为第
$n$ 个无穷小邻域，即 $Y$ 在 $X$ 中的相应邻域。假设以下条件之一成立：
\begin{enumerate}
\item $X$ 拟仿射，且
$\Gamma(X, \mathcal{O}_X) \to \lim \Gamma(Y_n, \mathcal{O}_{Y_n})$
是同构；或
\item $X$ 有丰沛可逆模 $\mathcal{L}$，且
$\Gamma(X, \mathcal{L}^{\otimes m}) \to
\lim \Gamma(Y_n, \mathcal{L}^{\otimes m}|_{Y_n})$
对所有 $m \gg 0$ 都是同构；或
\item 对每个有限局部自由 $\mathcal{O}_X$-模 $\mathcal{E}$，映射
$\Gamma(X, \mathcal{E}) \to \lim \Gamma(Y_n, \mathcal{E}|_{Y_n})$
是同构。
\end{enumerate}
则限制函子 $\textit{F\'Et}_X \to \textit{F\'Et}_Y$ 全忠实。
\end{lemma}

\begin{proof}
本引理形式地由 Lemma \ref{pione-lemma-restriction-fully-faithful} 和
Algebraic and Formal Geometry, Lemma
\ref{algebraization-lemma-completion-fully-faithful} 推出。
\end{proof}

\begin{lemma}
\label{pione-lemma-restriction-fully-faithful-general-special}
设 $X$ 是 Noether 概形，$Y \subset X$ 是闭子概形。令 $Y_n \subset X$ 为第
$n$ 个无穷小邻域，即 $Y$ 在 $X$ 中的相应邻域。令 $\mathcal{V}$ 为开子概形
$V \subset X$ 所成的集合，其中每个都包含 $Y$；以反包含关系排序。假设以下条件之一成立：
\begin{enumerate}
\item $X$ 拟仿射，且
$$
\colim_\mathcal{V} \Gamma(V, \mathcal{O}_V)
\longrightarrow
\lim \Gamma(Y_n, \mathcal{O}_{Y_n})
$$
是同构；或
\item $X$ 有丰沛可逆模 $\mathcal{L}$，且
$$
\colim_\mathcal{V} \Gamma(V, \mathcal{L}^{\otimes m})
\longrightarrow
\lim \Gamma(Y_n, \mathcal{L}^{\otimes m}|_{Y_n})
$$
对所有 $m \gg 0$ 都是同构；或
\item 对每个 $V \in \mathcal{V}$ 以及每个有限局部自由
$\mathcal{O}_V$-模 $\mathcal{E}$，映射
$$
\colim_{V' \geq V} \Gamma(V', \mathcal{E}|_{V'})
\longrightarrow
\lim \Gamma(Y_n, \mathcal{E}|_{Y_n})
$$
是同构。
\end{enumerate}
则函子
$$
\colim_\mathcal{V} \textit{F\'Et}_V \to \textit{F\'Et}_Y
$$
全忠实。
\end{lemma}

\begin{proof}
本引理形式地由 Lemma \ref{pione-lemma-restriction-fully-faithful-general} 和
Algebraic and Formal Geometry, Lemma
\ref{algebraization-lemma-completion-fully-faithful-general} 推出。
\end{proof}






\section{推出与基本群}
\label{pione-section-pushouts}

\noindent
主要结果如下。

\begin{lemma}
\label{pione-lemma-pushout-along-closed-immersion-and-integral}
在 More on Morphisms, Situation
\ref{more-morphisms-situation-pushout-along-closed-immersion-and-integral} 中，
例如若 $Z \to Y$ 与 $Z \to X$ 都是概形的闭浸入，则有范畴等价
$$
\textit{F\'Et}_{Y \amalg_Z X}
\longrightarrow
\textit{F\'Et}_Y
\times_{\textit{F\'Et}_Z}
\textit{F\'Et}_X
$$
\end{lemma}

\begin{proof}
由 More on Morphisms, Proposition
\ref{more-morphisms-proposition-pushout-along-closed-immersion-and-integral}，
该推出存在。此函子把推出上的有限 \'etale 概形 $U$ 送到基变换
$Y' = U \times_{Y \amalg_Z X} Y$ 与 $X' = U \times_{Y \amalg_Z X} X$，
以及自然同构 $Y' \times_Y Z \to X' \times_X Z$；此同构位于 $Z$ 上。
为证明此函子是等价，使用 More on Morphisms, Lemma
\ref{more-morphisms-lemma-pushout-functor-equivalence-flat} 构造拟逆函子。
唯一剩下要证明的是：给定分离、拟有限且 \'etale 的态射
$U \to Y \amalg_Z X$，并且 $X' \to X$ 与 $Y' \to Y$ 都有限，则
$U \to Y \amalg_Z X$ 有限。这既可由相应的代数事实（More on Algebra,
Lemma \ref{more-algebra-lemma-finite-module-over-fibre-product}）推出，也可由下述事实看出：
$$
X' \amalg Y' \to U
$$
满射，而 $X'$ 与 $Y'$ 在 $Y \amalg_Z X$ 上都是固有的（这里使用 More on Morphisms,
Proposition \ref{more-morphisms-proposition-pushout-along-closed-immersion-and-integral}
对推出的描述）。于是可应用 Morphisms, Lemma
\ref{morphisms-lemma-scheme-theoretic-image-is-proper}，得出 $U$ 在
$Y \amalg_Z X$ 上固有。由于拟有限且固有的态射是有限的（More on Morphisms,
Lemma \ref{more-morphisms-lemma-characterize-finite}），结论成立。
\end{proof}





\section{穿孔谱的有限 \'etale 覆盖，I}
\label{pione-section-pi1-punctured-spec}

\noindent
我们先证明一些 Lefschetz 型结果。

\begin{situation}
\label{pione-situation-local-lefschetz}
设 $(A, \mathfrak m)$ 是 Noether 局部环，且 $f \in \mathfrak m$。
置 $X = \Spec(A)$ 与 $X_0 = \Spec(A/fA)$，并令
$U = X \setminus \{\mathfrak m\}$ 与
$U_0 = X_0 \setminus \{\mathfrak m\}$ 分别为
$A$ 与 $A/fA$ 的穿孔谱。
\end{situation}

\noindent
回顾：对概形 $X$，在 $X$ 上有限
\'etale 的概形所成的范畴记为 $\textit{F\'Et}_X$；见
第 \ref{pione-section-finite-etale} 节。
在情形 \ref{pione-situation-local-lefschetz} 中，
我们将研究基变换函子
$$
\xymatrix{
\textit{F\'Et}_X \ar[d] \ar[r] & \textit{F\'Et}_U \ar[d] \\
\textit{F\'Et}_{X_0} \ar[r] & \textit{F\'Et}_{U_0}
}
$$
在许多情形下，右侧竖直箭头是忠实的。

\begin{lemma}
\label{pione-lemma-faithful}
在情形 \ref{pione-situation-local-lefschetz} 中，假设下列条件之一成立：
\begin{enumerate}
\item 对每个极小素理想都有 $\dim(A/\mathfrak p) \geq 2$，其中
$\mathfrak p \subset A$ 且 $f \not \in \mathfrak p$；或
\item $U$ 的每个连通分支都与 $U_0$ 相交。
\end{enumerate}
则
$$
\textit{F\'Et}_U \longrightarrow \textit{F\'Et}_{U_0},\quad
V \longmapsto V_0 = V \times_U U_0
$$
是忠实函子。
\end{lemma}

\begin{proof}
情形 (2) 立即由引理 \ref{pione-lemma-restriction-faithful} 得到。
假设 (1) 蕴含 $U$ 的每个不可约分支都与 $U_0$ 相交；见
Algebra, Lemma \ref{algebra-lemma-one-equation}。
因此 (1) 由 (2) 推出。
\end{proof}

\noindent
在证明更有内容的结果之前，需要先给出几个引理。

\begin{lemma}
\label{pione-lemma-fill-in-missing}
在情形 \ref{pione-situation-local-lefschetz} 中，设 $V \to U$ 是有限
态射。令 $A^\wedge$ 为 $\mathfrak m$-进完备化（对环 $A$ 而言），
令 $X' = \Spec(A^\wedge)$，并令 $U'$ 与 $V'$ 分别为
$U$ 与 $V$ 到 $X'$ 的基变换。若 $Y' \to X'$ 是满足
$V' = Y' \times_{X'} U'$ 的有限态射，则存在有限态射 $Y \to X$，
使得 $V = Y \times_X U$ 且 $Y' = Y \times_X X'$。
\end{lemma}

\begin{proof}
这是 More on Algebra, Proposition
\ref{more-algebra-proposition-equivalence} 的直接应用。
具体地，选取生成元 $f_1, \ldots, f_t$，它们生成 $\mathfrak m$。
对每个 $i$，记 $V \times_U D(f_i) = \Spec(B_i)$。
对 $1 \leq i, j \leq n$，得到同构
$\alpha_{ij} : (B_i)_{f_j} \to (B_j)_{f_i}$；它是
$A_{f_if_j}$-代数同构，因为二者的谱都表示
$V \times_U D(f_if_j)$。
记 $Y' = \Spec(B')$。由于 $V \times_U U' = Y \times_{X'} U'$，
得到同构 $\alpha_i : B'_{f_i} \to B_i \otimes_A A^\wedge$。
直接验证可知，$(B', B_i, \alpha_i, \alpha_{ij})$ 是
$\text{Glue}(A \to A^\wedge, f_1, \ldots, f_t)$ 的对象；见
More on Algebra, Remark \ref{more-algebra-remark-glueing-data}。
应用上引命题（并使用 More on Algebra, Remark
\ref{more-algebra-remark-formal-glueing-algebras} 得到代数结构），
可找到 $A$-代数 $B$，使得 $\text{Can}(B)$ 同构于
$(B', B_i, \alpha_i, \alpha_{ij})$。置 $Y = \Spec(B)$，
可见 $Y \to X$ 是配备有相容同构
$V \cong Y \times_X U$ 与 $Y' = Y \times_X X'$ 的态射，正合所需。
\end{proof}

\begin{lemma}
\label{pione-lemma-fully-faithful-henselian-completion}
在情形 \ref{pione-situation-local-lefschetz} 中，假设 $A$ 是 Hensel 环，
或更一般地，$(A, (f))$ 是 Hensel 对。
令 $A^\wedge$ 为 $\mathfrak m$-进完备化（对环 $A$ 而言），
令 $X' = \Spec(A^\wedge)$，并令 $U'$ 与 $U'_0$ 分别为
$U$ 与 $U_0$ 到 $X'$ 的基变换。若
$\textit{F\'Et}_{U'} \to \textit{F\'Et}_{U'_0}$ 全忠实，则
$\textit{F\'Et}_U \to \textit{F\'Et}_{U_0}$ 全忠实。
\end{lemma}

\begin{proof}
假设 $\textit{F\'Et}_{U'} \longrightarrow \textit{F\'Et}_{U'_0}$
全忠实。由于 $X' \to X$ 忠实平坦，立即可知函子
$V \to V_0 = V \times_U U_0$ 忠实。
由于有限 \'etale 覆盖的范畴具有内部 Hom
（引理 \ref{pione-lemma-internal-hom-finite-etale}），
只需证明以下断言：给定 $V$，它在 $U$ 上有限且 \'etale，则有
$$
\Mor_U(U, V) = \Mor_{U_0}(U_0, V_0)
$$
于是，假设给定态射 $s_0 : U_0 \to V_0$（在 $U_0$ 上），
我们将构造态射 $s : U \to V$（在 $U$ 上）。

\medskip\noindent
由假设，确实存在态射 $s' : U' \to V'$，其在 $V'_0$ 上的限制
为 $s'_0$，后者是 $s_0$ 的基变换。由于 $V' \to U'$ 有限且 \'etale，
这意味着 $V' = s'(U') \amalg W'$，其中 $W' \to U'$ 有限且 \'etale。
选取有限态射 $Z' \to X'$，使得
$W' = Z' \times_{X'} U'$。这可由下述形式的 Zariski 主定理得到：
More on Morphisms, Lemma
\ref{more-morphisms-lemma-quasi-finite-separated-pass-through-finite}
（略去一个小细节）。于是
$$
V' = s'(U') \amalg W' \longrightarrow X' \amalg Z' = Y'
$$
是开浸入，并且 $V' = Y' \times_{X'} U'$。
由引理 \ref{pione-lemma-fill-in-missing}，可找到有限态射 $Y \to X$，
使得 $V = Y \times_X U$ 且 $Y' = Y \times_X X'$。
记 $Y = \Spec(B)$，于是 $Y' = \Spec(B \otimes_A A^\wedge)$。
此时 $B \otimes_A A^\wedge$ 有幂等元 $e'$，它对应于开闭子概形
$X'$；后者包含于 $Y' = X' \amalg Z'$。

\medskip\noindent
先考虑 $A$ 为 Hensel 环的情形（稍为容易）。记 $\overline{e}$
为 $e'$ 在 $B \otimes_A \kappa(\mathfrak m) = B/\mathfrak mB$ 中的像；
它可提升为幂等元 $e$（属于 $B$），因为 $A$ 是 Hensel 环
（由 Algebra, Lemma \ref{algebra-lemma-characterize-henselian}，
$B$ 是局部环的乘积）。再由幂等元提升的唯一性可知，$e$ 映到 $e'$
（这里使用 $B \otimes_A A^\wedge$ 是局部环的乘积）。
令 $Y_1 \subset Y$ 为对应于 $e$ 的开闭子概形。于是
$Y_1 \times_X X' = s'(X')$，这蕴含 $Y_1 \to X$ 是同构
（由忠实平坦下降），并给出所需截面。

\medskip\noindent
再考虑 $(A, (f))$ 为 Hensel 对的情形。这里使用 $s'$ 是
$s'_0$ 的提升。具体地，令 $Y_{0, 1} \subset Y_0 = Y \times_X X_0$
为 $s_0(U_0) \subset V_0 = Y_0 \times_{X_0} U_0$ 的闭包。
由于 $X' \to X$ 平坦，基变换 $Y'_{0, 1} \subset Y'_0$
是 $s'_0(U'_0)$ 的闭包，因而等于 $X'_0 \subset Y'_0$
（见 Morphisms, Lemma
\ref{morphisms-lemma-flat-base-change-scheme-theoretic-image}）。
由于 $Y'_0 \to Y_0$ 是下沉的
（Morphisms, Lemma \ref{morphisms-lemma-fpqc-quotient-topology}），
可知 $Y_{0, 1}$ 在 $Y_0$ 中是开闭的。令 $e_0 \in B/fB$
为相应的幂等元。由 More on Algebra, Lemma
\ref{more-algebra-lemma-characterize-henselian-pair}，
可将 $e_0$ 提升为幂等元 $e \in B$。随后同前得出结论。
\end{proof}

\noindent
在情形 \ref{pione-situation-local-lefschetz} 中，限制函子
$\textit{F\'Et}_U \longrightarrow \textit{F\'Et}_{U_0}$
的全忠实性在相当温和的假设下即成立。特别地，这些假设通常并不蕴含
$U$ 是连通概形，但结论保证 $U$ 与 $U_0$ 的连通分支数相同。

\begin{lemma}
\label{pione-lemma-fully-faithful-simple}
在情形 \ref{pione-situation-local-lefschetz} 中，假设
\begin{enumerate}
\item[(a)] $A$ 有对偶化复形；
\item[(b)] $(A, (f))$ 是 Hensel 对；
\item[(c)] 下列条件之一成立：
\begin{enumerate}
\item[(i)] $A_f$ 满足 $(S_2)$，且 $X$ 的每个不包含于
$X_0$ 的不可约分支的维数都 $\geq 3$；或
\item[(ii)] 对每个素理想 $\mathfrak p \subset A$，若
$f \not \in \mathfrak p$，则有
$\text{depth}(A_\mathfrak p) + \dim(A/\mathfrak p) > 2$。
\end{enumerate}
\end{enumerate}
则限制函子
$\textit{F\'Et}_U \longrightarrow \textit{F\'Et}_{U_0}$
全忠实。
\end{lemma}

\begin{proof}
令 $A'$ 为关于 $\mathfrak m$ 的环 $A$ 的进完备化。我们将证明这些假设对
$A'$ 仍成立。对条件 (a) 与 (b)，这是显然的。条件 (c)(ii) 由
Local Cohomology, Lemma \ref{local-cohomology-lemma-change-completion}
在完备化下保持。其次，假设 (c)(i) 成立。由于 $A$ 普遍链状
（Dualizing Complexes, Lemma \ref{dualizing-lemma-universally-catenary}），
可知 $\Spec(A')$ 的每个不包含于 $V(f)$ 的不可约分支的维数都
$\geq 3$；见
More on Algebra, Proposition \ref{more-algebra-proposition-ratliff}。
由于 $A \to A'$ 平坦且具有 Gorenstein 纤维，
$A_f$ 满足 $(S_2)$ 蕴含 $A'_f$ 满足 $(S_2)$。
这里使用的参考结果为：
Dualizing Complexes, Section \ref{dualizing-section-formal-fibres},
More on Algebra, Section \ref{more-algebra-section-properties-formal-fibres},
以及 Algebra, Lemma \ref{algebra-lemma-Sk-goes-up}。
因此，由引理 \ref{pione-lemma-fully-faithful-henselian-completion}，
可设 $A$ 是完备 Noether 局部环。

\medskip\noindent
除其他假设外，再假设 $A$ 是完备局部环。由引理
\ref{pione-lemma-restriction-fully-faithful}，结论可由
Algebraic and Formal Geometry, Lemma
\ref{algebraization-lemma-fully-faithful-simple-two} 得到。
\end{proof}

\begin{lemma}
\label{pione-lemma-fully-faithful-minimal}
\begin{reference}
\cite[Corollary 1.11]{Bhatt-local}
\end{reference}
在情形 \ref{pione-situation-local-lefschetz} 中，假设
\begin{enumerate}
\item $H^1_\mathfrak m(A)$ 与 $H^2_\mathfrak m(A)$ 均被 $f$ 的某个幂零化；并且
\item $A$ 是 Hensel 环，或更一般地，$(A, (f))$ 是 Hensel 对。
\end{enumerate}
则限制函子
$\textit{F\'Et}_U \longrightarrow \textit{F\'Et}_{U_0}$
全忠实。
\end{lemma}

\begin{proof}
由引理 \ref{pione-lemma-fully-faithful-henselian-completion}，
可设 $A$ 是完备 Noether 局部环。
（这些假设会保持；使用 Dualizing Complexes, Lemma
\ref{dualizing-lemma-torsion-change-rings}。）
由引理 \ref{pione-lemma-restriction-fully-faithful}，结论可由
Algebraic and Formal Geometry, Lemma
\ref{algebraization-lemma-fully-faithful-alternative} 得到。
\end{proof}

\begin{lemma}
\label{pione-lemma-fully-faithful}
在情形 \ref{pione-situation-local-lefschetz} 中，假设 $A$ 的深度 $\geq 3$，
并且 $A$ 是 Hensel 环，或更一般地，$(A, (f))$ 是 Hensel 对。则限制函子
$\textit{F\'Et}_U \to \textit{F\'Et}_{U_0}$
全忠实。
\end{lemma}

\begin{proof}
深度假设蕴含
$H^1_\mathfrak m(A) = H^2_\mathfrak m(A) = 0$；见
Dualizing Complexes, Lemma \ref{dualizing-lemma-depth}。
因此可应用引理 \ref{pione-lemma-fully-faithful-minimal}。
\end{proof}



























\section{局部情形中的纯性，I}
\label{pione-section-local-purity}

\noindent
设 $(A, \mathfrak m)$ 是 Noether 局部环。置 $X = \Spec(A)$，
并令 $U = X \setminus \{\mathfrak m\}$ 为穿孔谱。若对
$(A, \mathfrak m)$，限制函子
$$
\textit{F\'Et}_X \longrightarrow \textit{F\'Et}_U
$$
本质满，则称{\it 纯性成立}。本节试图理解：从 $X$ 过渡到
超曲面 $X_0$（位于 $X$ 中）时，这个问题会如何变化。
换言之，我们研究穿孔谱基本群的一种局部 Lefschetz 性质。
这些结果将用于在证明局部纯性的主要结果时对维数作归纳；这些结果是
引理 \ref{pione-lemma-local-purity}、
命题 \ref{pione-proposition-purity-complete-intersection} 以及
命题 \ref{pione-proposition-purity-smooth-over-depth2}。

\begin{lemma}
\label{pione-lemma-sections-over-punctured-spec}
设 $(A, \mathfrak m)$ 是 Noether 局部环。置 $X = \Spec(A)$，
并令 $U = X \setminus \{\mathfrak m\}$。设 $\pi : Y \to X$
是有限态射，并且对所有闭点都有
$\text{depth}(\mathcal{O}_{Y, y}) \geq 2$，其中 $y \in Y$。
则 $Y$ 是 $B = \mathcal{O}_Y(\pi^{-1}(U))$ 的谱。
\end{lemma}

\begin{proof}
置 $V = \pi^{-1}(U)$，并以 $\pi' : V \to U$ 表示 $\pi$ 的限制。
考虑 $\mathcal{O}_X$-模映射
$$
\pi_*\mathcal{O}_Y \longrightarrow j_*\pi'_*\mathcal{O}_V
$$
其中 $j : U \to X$ 是包含态射。我们断言 Divisors, Lemma
\ref{divisors-lemma-depth-2-hartog} 可应用于此映射。若然，则
$B = \Gamma(Y, \mathcal{O}_Y)$，从而引理成立。令 $x \in X$
为闭点。只需证明
$\text{depth}((\pi_*\mathcal{O}_Y)_x) \geq 2$。
令 $y_1, \ldots, y_n \in Y$ 为映到 $x$ 的各点。
由 Algebra, Lemma \ref{algebra-lemma-depth-goes-down-finite}，
只需证明 $\text{depth}(\mathcal{O}_{Y, y_i}) \geq 2$
对 $i = 1, \ldots, n$ 成立。
这正是引理的假设，证明完毕。
\end{proof}

\begin{lemma}
\label{pione-lemma-reformulate-purity}
设 $(A, \mathfrak m)$ 是 Noether 局部环。置 $X = \Spec(A)$，
并令 $U = X \setminus \{\mathfrak m\}$。
设 $V$ 在 $U$ 上有限且 \'etale。假设 $A$ 的深度 $\geq 2$。
则下列条件等价：
\begin{enumerate}
\item 有 $V = Y \times_X U$，其中 $Y \to X$ 有限且 \'etale；
\item $B = \Gamma(V, \mathcal{O}_V)$ 在 $A$ 上有限且 \'etale。
\end{enumerate}
\end{lemma}

\begin{proof}
以 $\pi : V \to U$ 表示给定的有限 \'etale 态射。
假设 (1) 中的 $Y$ 存在。令 $x \in X$ 为对应于 $\mathfrak m$ 的点。
令 $y \in Y$ 为映到 $x$ 的点。我们断言
$\text{depth}(\mathcal{O}_{Y, y}) \geq 2$。
这是因为 $Y \to X$ 是 \'etale 的，因而
$A = \mathcal{O}_{X, x}$ 与 $\mathcal{O}_{Y, y}$ 具有相同深度
（Algebra, Lemma \ref{algebra-lemma-apply-grothendieck}）。
因此可应用引理 \ref{pione-lemma-sections-over-punctured-spec}，并有
$Y = \Spec(B)$。

\medskip\noindent
蕴含 (2) $\Rightarrow$ (1) 较为容易，细节从略。
\end{proof}

\begin{lemma}
\label{pione-lemma-reformulate-purity-normal}
设 $(A, \mathfrak m)$ 是 Noether 局部环。置 $X = \Spec(A)$，
并令 $U = X \setminus \{\mathfrak m\}$。假设 $A$ 正规且维数
$\geq 2$。函子
$$
\textit{F\'Et}_U \longrightarrow
\left\{
\begin{matrix}
\text{finite normal }A\text{-algebras }B\text{ such} \\
\text{that }\Spec(B) \to X\text{ is \'etale over }U
\end{matrix}
\right\},
\quad
V \longmapsto \Gamma(V, \mathcal{O}_V)
$$
是等价。此外，有 $V = Y \times_X U$，其中 $Y \to X$
有限且 \'etale，当且仅当 $B = \Gamma(V, \mathcal{O}_V)$
在 $A$ 上有限且 \'etale。
\end{lemma}

\begin{proof}
注意 $\text{depth}(A) \geq 2$，因为 $A$ 正规
（Serre 正规性判据，Algebra, Lemma
\ref{algebra-lemma-criterion-normal}）。
因此最后的陈述由引理 \ref{pione-lemma-reformulate-purity} 得到。
给定有限 \'etale 态射 $\pi : V \to U$，置
$B = \Gamma(V, \mathcal{O}_V)$。若能证明 $B$ 正规且在 $A$ 上有限，
便得到所显示的函子。由于有显然的拟逆函子，这也就是全部所需证明的内容。

\medskip\noindent
由于 $A$ 正规，概形 $V$ 正规
（Descent, Lemma \ref{descent-lemma-normal-local-smooth}）。
因此 $V$ 是有限多个整概形的不交并
（Properties, Lemma \ref{properties-lemma-normal-Noetherian}）。
所以可设 $V$ 整。在这种情形下，函数域 $L$（属于 $V$）
（Morphisms, Section \ref{morphisms-section-rational-maps}）
是 $A$ 的分式域的有限可分扩张
（这是通过考察 $V \to U$ 的一般纤维并使用 Morphisms, Lemma
\ref{morphisms-lemma-etale-over-field} 得到的）。
由 Algebra, Lemma
\ref{algebra-lemma-Noetherian-normal-domain-finite-separable-extension}，
整闭包 $B' \subset L$（即 $A$ 在 $L$ 中的整闭包）在 $A$ 上有限。
由 More on Algebra, Lemma
\ref{more-algebra-lemma-integral-closure-reflexive}，$B'$ 是自反的
$A$-模；进而由 More on Algebra, Lemma
\ref{more-algebra-lemma-reflexive-over-normal} 可知
$\text{depth}_A(B') \geq 2$。

\medskip\noindent
设 $f \in \mathfrak m$。则
$B_f = \Gamma(V \times_U D(f), \mathcal{O}_V)$
（Properties, Lemma \ref{properties-lemma-invert-f-sections}）。
所以 $B'_f = B_f$，因为 $B_f$ 正规（见上文）、在 $A_f$ 上有限，
且分式域为 $L$。于是 $V = \Spec(B') \times_X U$。
由引理 \ref{pione-lemma-sections-over-punctured-spec} 可得 $B = B'$；
这里将该引理应用于 $\Spec(B') \to X$。该引理适用，是因为局部化
$B'_{\mathfrak m'}$（它是 $B'$ 在极大理想
$\mathfrak m' \subset B'$ 处的局部化，而这些理想位于
$\mathfrak m$ 上方）的深度
$\geq 2$；这由 Algebra, Lemma
\ref{algebra-lemma-depth-goes-down-finite} 以及上一段关于深度的说明得到。
\end{proof}

\begin{lemma}
\label{pione-lemma-purity-and-completion}
设 $(A, \mathfrak m)$ 是 Noether 局部环。置 $X = \Spec(A)$，
并令 $U = X \setminus \{\mathfrak m\}$。设 $V$ 在 $U$ 上有限且
\'etale。令 $A^\wedge$ 为 $\mathfrak m$-进完备化（对环 $A$ 而言），
令 $X' = \Spec(A^\wedge)$，并令 $U'$ 与 $V'$ 分别为
$U$ 与 $V$ 到 $X'$ 的基变换。下列条件等价：
\begin{enumerate}
\item 有 $V = Y \times_X U$，其中 $Y \to X$ 有限且 \'etale；并且
\item 有 $V' = Y' \times_{X'} U'$，其中 $Y' \to X'$
有限且 \'etale。
\end{enumerate}
\end{lemma}

\begin{proof}
蕴含 (1) $\Rightarrow$ (2) 可由对一个解 $Y \to X$ 作基变换得到。
设 $Y' \to X'$ 如 (2) 所述。由引理 \ref{pione-lemma-fill-in-missing}，
可找到有限态射 $Y \to X$，使得 $V = Y \times_X U$ 且
$Y' = Y \times_X X'$。由下降可知 $Y \to X$ 有限且 \'etale
（Algebra, Lemmas \ref{algebra-lemma-descend-properties-modules} 以及
\ref{algebra-lemma-etale}）。证明完毕。
\end{proof}

\noindent
以下两个引理的要点在于，这些假设并不强制 $A$ 的深度
$\geq 3$。例如，若 $A$ 是维数 $\geq 3$ 的完备正规局部整环，
且 $f \in \mathfrak m$ 非零，则这些假设成立。

\begin{lemma}
\label{pione-lemma-lift-simple}
在情形 \ref{pione-situation-local-lefschetz} 中，设 $V$ 在 $U$ 上有限且
\'etale。假设
\begin{enumerate}
\item[(a)] $A$ 有对偶化复形；
\item[(b)] $(A, (f))$ 是 Hensel 对；
\item[(c)] 下列条件之一成立：
\begin{enumerate}
\item[(i)] $A_f$ 满足 $(S_2)$，且 $X$ 的每个不包含于
$X_0$ 的不可约分支的维数都 $\geq 3$；或
\item[(ii)] 对每个素理想 $\mathfrak p \subset A$，若
$f \not \in \mathfrak p$，则
$\text{depth}(A_\mathfrak p) + \dim(A/\mathfrak p) > 2$。
\end{enumerate}
\item[(d)] $V_0 = V \times_U U_0$ 等于 $Y_0 \times_{X_0} U_0$，
其中 $Y_0 \to X_0$ 有限且 \'etale。
\end{enumerate}
则有 $V = Y \times_X U$，其中 $Y \to X$ 有限且 \'etale。
\end{lemma}

\begin{proof}
使用引理 \ref{pione-lemma-purity-and-completion}，归约到完备情形。
（这些假设会保持；见引理 \ref{pione-lemma-fully-faithful-simple} 的证明。）

\medskip\noindent
在完备情形下，由 More on Algebra, Lemma
\ref{more-algebra-lemma-finite-etale-equivalence}，可将
$Y_0 \to X_0$ 提升为有限 \'etale 态射 $Y \to X$；注意，由
More on Algebra, Lemma \ref{more-algebra-lemma-complete-henselian}，
$(A, fA)$ 是 Hensel 对。随后可使用引理
\ref{pione-lemma-fully-faithful-simple} 看出 $V$ 同构于
$Y \times_X U$，证明完毕。
\end{proof}

\begin{lemma}
\label{pione-lemma-lift-purity-general}
在情形 \ref{pione-situation-local-lefschetz} 中，设 $V$ 在 $U$ 上有限且
\'etale。假设
\begin{enumerate}
\item $H^1_\mathfrak m(A)$ 与 $H^2_\mathfrak m(A)$ 均被 $f$ 的某个幂零化；
\item $V_0 = V \times_U U_0$ 等于 $Y_0 \times_{X_0} U_0$，
其中 $Y_0 \to X_0$ 有限且 \'etale。
\end{enumerate}
则有 $V = Y \times_X U$，其中 $Y \to X$ 有限且 \'etale。
\end{lemma}

\begin{proof}
使用引理 \ref{pione-lemma-purity-and-completion}，归约到完备情形。
（这些假设会保持；使用 Dualizing Complexes, Lemma
\ref{dualizing-lemma-torsion-change-rings}。）

\medskip\noindent
在完备情形下，由 More on Algebra, Lemma
\ref{more-algebra-lemma-finite-etale-equivalence}，可将
$Y_0 \to X_0$ 提升为有限 \'etale 态射 $Y \to X$；注意，由
More on Algebra, Lemma \ref{more-algebra-lemma-complete-henselian}，
$(A, fA)$ 是 Hensel 对。随后可使用引理
\ref{pione-lemma-fully-faithful-minimal} 看出 $V$ 同构于
$Y \times_X U$，证明完毕。
\end{proof}

\begin{lemma}
\label{pione-lemma-lift-purity}
在情形 \ref{pione-situation-local-lefschetz} 中，设 $V$ 在 $U$ 上有限且
\'etale。假设
\begin{enumerate}
\item $A$ 的深度 $\geq 3$；
\item $V_0 = V \times_U U_0$ 等于 $Y_0 \times_{X_0} U_0$，
其中 $Y_0 \to X_0$ 有限且 \'etale。
\end{enumerate}
则有 $V = Y \times_X U$，其中 $Y \to X$ 有限且 \'etale。
\end{lemma}

\begin{proof}
深度假设蕴含
$H^1_\mathfrak m(A) = H^2_\mathfrak m(A) = 0$；见
Dualizing Complexes, Lemma \ref{dualizing-lemma-depth}。
因此可应用引理 \ref{pione-lemma-lift-purity-general}。
\end{proof}










\section{分歧轨迹的纯性}
\label{pione-section-purity}

\noindent
我们将使用有限局部自由态射的判别式。见《判别式》Section
\ref{discriminant-section-discriminant}。

\begin{lemma}
\label{pione-lemma-find-point-codim-1}
设 $(A, \mathfrak m)$ 是满足 $\dim(A) \geq 1$ 的 Noether 局部环。
设 $f \in \mathfrak m$。则存在 $\mathfrak p \in V(f)$，使得
$\dim(A_\mathfrak p) = 1$。
\end{lemma}

\begin{proof}
对 $\dim(A)$ 作归纳。若 $\dim(A) = 1$，则取 $\mathfrak p = \mathfrak m$
即可。若 $\dim(A) > 1$，则设 $Z \subset \Spec(A)$ 是维数 $> 1$ 的一个
不可约分支。于是 $V(f) \cap Z$ 的维数 $> 0$
（《代数》引理 \ref{algebra-lemma-one-equation}）。选取素理想
$\mathfrak q \in V(f) \cap Z$，$\mathfrak q \not = \mathfrak m$，
它对应于 $A$ 的穿孔谱的一个闭点；由《性质》引理
\ref{properties-lemma-complement-closed-point-Jacobson}，这是可以做到的。
于是 $\mathfrak q$ 不是 $Z$ 的一般点。因此
$0 < \dim(A_\mathfrak q) < \dim(A)$ 且 $f \in \mathfrak q A_\mathfrak q$。
由关于维数的归纳，可以找到
$f \in \mathfrak p \subset A_\mathfrak q$，使得
$\dim((A_\mathfrak q)_\mathfrak p) = 1$。
于是 $\mathfrak p \cap A$ 满足要求。
\end{proof}

\begin{lemma}
\label{pione-lemma-ramification-quasi-finite-flat}
设 $f : X \to Y$ 是局部 Noether 概形的态射。设 $x \in X$。假设
\begin{enumerate}
\item $f$ 平坦，
\item $f$ 在 $x$ 处拟有限，
\item $x$ 不是 $X$ 的某个不可约分支的一般点，
\item 对特化 $x' \leadsto x$，若
$\dim(\mathcal{O}_{X, x'}) = 1$，则态射 $f$ 在 $x'$ 处非分歧。
\end{enumerate}
则 $f$ 在 $x$ 处是 \'etale 的。
\end{lemma}

\begin{proof}
注意，$f$ 非分歧的点所成集合与 $f$ 为 \'etale 的点所成集合相同，并且
这个集合是开的。见《态射》定义 \ref{morphisms-definition-unramified}
和 \ref{morphisms-definition-etale} 以及引理
\ref{morphisms-lemma-flat-unramified-etale}。
为了检验 $f$ 在 $x$ 处是 \'etale 的，我们可以在基底和目标上 \'etale
局部地工作（《下降》引理 \ref{descent-lemma-descending-property-etale} 和
\ref{descent-lemma-etale-etale-local-source}）。
因此，可以应用《态射补篇》引理
\ref{more-morphisms-lemma-etale-makes-quasi-finite-finite-at-point}
并假设 $f : X \to Y$ 有限，且 $x$ 是 $X$ 中位于 $y = f(x)$ 上方的唯一一点。
于是 $f$ 有限局部自由（《态射》引理
\ref{morphisms-lemma-finite-flat}）。

\medskip\noindent
假设 $f$ 有限局部自由，且 $x$ 是 $X$ 中位于 $y = f(x)$ 上方的唯一一点。
由《判别式》引理 \ref{discriminant-lemma-discriminant}，可得局部主闭子概形
$D_\pi \subset Y$，使得 $y' \in D_\pi$ 当且仅当存在 $x' \in X$，满足
$f(x') = y'$ 且 $f$ 在 $x'$ 处分歧。因此必须证明 $y \not \in D_\pi$。
假设 $y \in D_\pi$，以导出矛盾。

\medskip\noindent
由条件 (3)，有 $\dim(\mathcal{O}_{X, x}) \geq 1$。由《代数》引理
\ref{algebra-lemma-dimension-base-fibre-equals-total}，有
$\dim(\mathcal{O}_{X, x}) = \dim(\mathcal{O}_{Y, y})$。
由引理 \ref{pione-lemma-find-point-codim-1}，可以找到 $y' \in D_\pi$，它特化到 $y$，
且满足 $\dim(\mathcal{O}_{Y, y'}) = 1$。
选取 $x' \in X$，满足 $f(x') = y'$ 且 $f$ 在该点分歧。由于 $f$
有限，故它是闭态射，从而 $x' \leadsto x$。同前可得
$\dim(\mathcal{O}_{X, x'}) = \dim(\mathcal{O}_{Y, y'}) = 1$。
这与性质 (4) 矛盾。
\end{proof}

\begin{lemma}
\label{pione-lemma-local-purity}
设 $(A, \mathfrak m)$ 是维数为 $d \geq 2$ 的正则局部环。置
$X = \Spec(A)$ 和 $U = X \setminus \{\mathfrak m\}$。则
\begin{enumerate}
\item 函子 $\textit{F\'Et}_X \to \textit{F\'Et}_U$ 本质满，即纯性对 $A$
成立；
\item 任意有限同态 $A \to B$，若 $B$ 正规且在穿孔谱上诱导有限 \'etale
态射，则它是 \'etale 的。
\end{enumerate}
\end{lemma}

\begin{proof}
回顾由《代数》引理 \ref{algebra-lemma-regular-normal}，正则局部环是正规环。
因此，由引理 \ref{pione-lemma-reformulate-purity-normal}，(1) 与 (2) 等价。
我们对 $d$ 作归纳来证明该引理。

\medskip\noindent
考虑 $d = 2$ 的情形。此时 $A \to B$ 平坦。事实上，由《代数》命题
\ref{algebra-proposition-going-down-normal-integral}，$A \to B$ 满足降链性质。
再由《代数》引理 \ref{algebra-lemma-dimension-base-fibre-equals-total}，有
$\dim(B_{\mathfrak m'}) = 2$ 对所有极大理想 $\mathfrak m' \subset B$ 成立。
由《代数》引理 \ref{algebra-lemma-criterion-normal}，$B_{\mathfrak m'}$ 是
Cohen--Macaulay 环。因此，关键的一步是应用《代数》引理
\ref{algebra-lemma-CM-over-regular-flat}，得到 $A \to B_{\mathfrak m'}$ 平坦。
再由《代数》引理 \ref{algebra-lemma-flat-localization}，$A \to B$ 平坦。
于是可以应用引理 \ref{pione-lemma-ramification-quasi-finite-flat}（也可以直接应用较简单的
《判别式》引理 \ref{discriminant-lemma-discriminant} 进行论证），看出
$A \to B$ 是 \'etale 的。

\medskip\noindent
考虑 $d \geq 3$ 的情形。设 $V \to U$ 有限 \'etale。取
$f \in \mathfrak m_A$，$f \not \in \mathfrak m_A^2$。于是 $A/fA$ 是维数为
$d - 1 \geq 2$ 的正则局部环，见《代数》引理
\ref{algebra-lemma-regular-ring-CM}。设 $U_0$ 是 $A/fA$ 的穿孔谱，并置
$V_0 = V \times_U U_0$。由引理 \ref{pione-lemma-lift-purity}，只需证明
$V_0$ 位于 $\textit{F\'Et}_{\Spec(A/fA)} \to \textit{F\'Et}_{U_0}$ 的本质像中。
这由归纳假设得出。
\end{proof}

\begin{lemma}[分歧轨迹的纯性]
\label{pione-lemma-purity}
\begin{reference}
\cite{Nagata-Purity} 和 \cite[Exp. X, Thm. 3.1]{SGA1}
\end{reference}
\begin{history}
Zariski 最早在 \cite{Zariski-Purity} 中以几何形式陈述并证明了这一结果。
Nagata 对非完美基域的推广，作为同卷下一篇文章发表于
《美国国家科学院院刊》，见 \cite{Nagata-Remarks-Purity}。次年，Nagata 在
\cite{Nagata-Purity} 中证明了 Noether 局部环的结果。他的证明使用了 Chow
的一个结果，即完备局部整环的 Bertini 定理，见 \cite{Chow-Bertini}；
Bertini 定理的历史见 Kleiman 的历史性文章 \cite{Kleiman-Bertini}。
几年后，Auslander 找到了一个完全不同的证明，见 \cite{Auslander-Purity}。
\end{history}
设 $f : X \to Y$ 是局部 Noether 概形的态射。设 $x \in X$，并置
$y = f(x)$。假设
\begin{enumerate}
\item $\mathcal{O}_{X, x}$ 正规；
\item $\mathcal{O}_{Y, y}$ 正则；
\item $f$ 在 $x$ 处拟有限；
\item $\dim(\mathcal{O}_{X, x}) = \dim(\mathcal{O}_{Y, y}) \geq 1$；
\item 对特化 $x' \leadsto x$，若
$\dim(\mathcal{O}_{X, x'}) = 1$，则态射 $f$ 在 $x'$ 处非分歧。
\end{enumerate}
则 $f$ 在 $x$ 处是 \'etale 的。
\end{lemma}

\begin{proof}
我们对 $d = \dim(\mathcal{O}_{X, x}) = \dim(\mathcal{O}_{Y, y})$ 作归纳来证明该引理。

\medskip\noindent
先考虑较平凡的情形 $d = 1$。此时已假设 $f$ 在 $x$ 处非分歧，且
$\mathcal{O}_{Y, y}$ 是离散赋值环（《代数》引理
\ref{algebra-lemma-characterize-dvr}）。于是 $\mathcal{O}_{X, x}$ 在
$\mathcal{O}_{Y, y}$ 上平坦（否则该态射在 $x$ 处不会拟有限），从而看出
$f$ 在 $x$ 处平坦。由于平坦 $+$ 非分歧即为 \'etale，结论成立（略去一些细节）。

\medskip\noindent
考虑 $d \geq 2$ 的情形。我们将对 $d$ 作归纳，把问题约化到引理
\ref{pione-lemma-local-purity} 中讨论的情形。为了检验 $f$ 在 $x$ 处是 \'etale 的，
可以在基底和目标上 \'etale 局部地工作（《下降》引理
\ref{descent-lemma-descending-property-etale} 和
\ref{descent-lemma-etale-etale-local-source}）。因此，可以应用《态射补篇》引理
\ref{more-morphisms-lemma-etale-makes-quasi-finite-finite-at-point}
并假设 $f : X \to Y$ 有限，且 $x$ 是 $X$ 中位于 $y$ 上方的唯一一点。这里使用了
局部环的 \'etale 扩张不改变维数、正规性和正则性；见《代数补篇》Section
\ref{more-algebra-section-permanence-etale} 和《\'Etale 态射》Section
\ref{etale-section-properties-permanence}。

\medskip\noindent
接着，可以由 $\Spec(\mathcal{O}_{Y, y})$ 作基变换，并假设 $Y$ 是一个正则局部环的谱。
于是 $X = \Spec(\mathcal{O}_{X, x})$，因为 $X$ 的每一点都必定特化到 $x$。

\medskip\noindent
环同态 $\mathcal{O}_{Y, y} \to \mathcal{O}_{X, x}$ 有限，并且必为单射
（由维数相等）。由《代数》命题
\ref{algebra-proposition-going-down-normal-integral}，可知
$\mathcal{O}_{Y, y} \to \mathcal{O}_{X, x}$ 满足降链性质（并使用了由《代数》引理
\ref{algebra-lemma-regular-normal} 得出的正则环是正规环这一事实）。选取
$x' \in X$，$x' \not = x$，并令其像为 $y' = f(x')$。由于正规整环的局部化仍
正规，$\mathcal{O}_{X, x'}$ 正规。类似地，$\mathcal{O}_{Y, y'}$ 正则（见《代数》引理
\ref{algebra-lemma-localization-of-regular-local-is-regular}）。由《代数》引理
\ref{algebra-lemma-dimension-base-fibre-equals-total}（上面已经验证降链性质），有
$\dim(\mathcal{O}_{X, x'}) = \dim(\mathcal{O}_{Y, y'})$。这些维数当然严格小于
$d$，因为 $x' \not = x$；再由关于 $d$ 的归纳，得出 $f$ 在 $x'$ 处是 \'etale 的。

\medskip\noindent
于是归结到如下情形：存在 Noether 局部环之间的有限局部同态 $A \to B$，
这些环的维数为 $d \geq 2$，其中 $A$ 正则，$B$ 正规，并且它在穿孔谱上诱导有限 \'etale 态射
$V \to U$。我们的目标是证明 $A \to B$ 为 \'etale。这由引理
\ref{pione-lemma-local-purity} 得出，证明完毕。
\end{proof}

\noindent
下面的引理有时可用于找出有限 \'etale 态射能够延拓到的最大开子集。

\begin{lemma}
\label{pione-lemma-extend-S2}
设 $j : U \to X$ 是局部 Noether 概形的开浸入，且
$\text{depth}(\mathcal{O}_{X, x}) \geq 2$ 对 $x \not \in U$ 成立。
设 $\pi : V \to U$ 有限 \'etale。则
\begin{enumerate}
\item $\mathcal{B} = j_*\pi_*\mathcal{O}_V$ 是自反凝聚 $\mathcal{O}_X$-代数；置
$Y = \underline{\Spec}_X(\mathcal{B})$；
\item $Y \to X$ 是满足如下条件的唯一有限态射：$V = Y \times_X U$，且
$\text{depth}(\mathcal{O}_{Y, y}) \geq 2$ 对 $y \in Y \setminus V$ 成立；
\item $Y \to X$ 在 $y$ 处为 \'etale 当且仅当 $Y \to X$ 在 $y$ 处平坦；
\item $Y \to X$ 为 \'etale 当且仅当 $\mathcal{B}$ 作为 $\mathcal{O}_X$-模有限局部自由。
\end{enumerate}
此外，(a) $\mathcal{B}$ 和 $Y \to X$ 的构造与局部 Noether 概形的平坦态射
$X' \to X$ 所给的基变换可交换；(b) 若 $V' \to U'$ 是有限 \'etale 态射，
$U \subset U' \subset X$ 为开子概形，并且它限制为 $V \to U$（在 $U$ 上），
则存在唯一同构 $Y' \times_X U' = V'$（在 $U'$ 上）。
\end{lemma}

\begin{proof}
注意，$\pi_*\mathcal{O}_V$ 是有限局部自由 $\mathcal{O}_U$-模，特别地是自反模。
由《除子》引理 \ref{divisors-lemma-reflexive-S2-extend}，模
$j_*\pi_*\mathcal{O}_V$ 是 $X$ 上以 $\pi_*\mathcal{O}_V$ 为其在 $U$ 上限制的
唯一自反凝聚模。这证明了 (1)。

\medskip\noindent
由构造，$Y \times_X U = V$。由于 $\mathcal{B}$ 凝聚，可知 $Y \to X$ 有限。
由《除子》引理 \ref{divisors-lemma-reflexive-S2}，
$\text{depth}(\mathcal{B}_x) \geq 2$ 对 $x \in X \setminus U$ 成立。因此，由
《代数》引理 \ref{algebra-lemma-depth-goes-down-finite}，
$\text{depth}(\mathcal{O}_{Y, y}) \geq 2$ 对 $y \in Y \setminus V$ 成立。反之，假设
$\pi' : Y' \to X$ 是有限态射，满足 $V = Y' \times_X U$，且
$\text{depth}(\mathcal{O}_{Y', y'}) \geq 2$ 对 $y' \in Y' \setminus V$ 成立。
则 $\pi'_*\mathcal{O}_{Y'}$ 限制为 $\pi_*\mathcal{O}_V$（在 $U$ 上），并且
$\text{depth}((\pi'_*\mathcal{O}_{Y'})_x) \geq 2$ 对
$x \in X \setminus U$ 成立；这由《代数》引理
\ref{algebra-lemma-depth-goes-down-finite} 得出。
例如，由《除子》引理 \ref{divisors-lemma-depth-2-hartog}，
$\pi'_*\mathcal{O}_{Y'}$ 与 $j_*\pi_*\mathcal{O}_V$ 典范同构。这证明了 (2)。

\medskip\noindent
若 $Y \to X$ 在 $y$ 处为 \'etale，则 $Y \to X$ 在 $y$ 处平坦。反之，假设
$Y \to X$ 在 $y$ 处平坦。若 $y \in V$，则 $Y \to X$ 在 $y$ 处为 \'etale。
若 $y \not \in V$，则检验引理 \ref{pione-lemma-ramification-quasi-finite-flat} 的
(1)、(2)、(3)、(4) 成立，即可看出 $Y \to X$ 在 $y$ 处为 \'etale。
(1) 和 (2) 显然成立；由于 $\text{depth}(\mathcal{O}_{Y, y}) \geq 2$，(3) 也显然成立。
若 $y' \leadsto y$ 是特化且 $\dim(\mathcal{O}_{Y, y'}) = 1$，则
$y' \in V$，因为否则由《代数》引理 \ref{algebra-lemma-bound-depth}，
这个局部环的深度将为 $2$，矛盾。因此，$Y \to X$ 在 $y'$ 处为 \'etale，
从而引理 \ref{pione-lemma-ramification-quasi-finite-flat} 的 (4) 也成立。
这就完成了 (3) 的证明。

\medskip\noindent
(4) 由 (3) 以及如下事实得出：$((Y \to X)_*\mathcal{O}_Y)_x$ 是平坦
$\mathcal{O}_{X, x}$-模，当且仅当 $\mathcal{O}_{Y, y}$ 是平坦
$\mathcal{O}_{X, x}$-模；这里取遍所有 $y \in Y$，且它们映到 $x$。见《代数》引理
\ref{algebra-lemma-flat-localization}。这里还使用了 Noether 环上的有限平坦模是
有限局部自由模这一事实，见《代数》引理 \ref{algebra-lemma-finite-projective}
（以及《代数》引理
\ref{algebra-lemma-Noetherian-finite-type-is-finite-presentation}）。

\medskip\noindent
关于引理最后的断言，(a) 由平坦基变换得出，见《概形的上同调》引理
\ref{coherent-lemma-flat-base-change-cohomology}；把 (2) 中的唯一性应用于限制
$Y \times_X U'$，便得到 (b)。
\end{proof}

\begin{lemma}
\label{pione-lemma-extend-pure}
设 $j : U \to X$ 是 Noether 概形的开浸入，并且纯性对
$\mathcal{O}_{X, x}$ 成立；这里取遍所有 $x \not \in U$。则
$$
\textit{F\'Et}_X \longrightarrow \textit{F\'Et}_U
$$
本质满。
\end{lemma}

\begin{proof}
设 $V \to U$ 是有限 \'etale 态射。由 Noether 归纳，只需把 $V \to U$ 延拓为
到 $X$ 的一个严格更大开子集的有限 \'etale 态射。设 $x \in X \setminus U$
是 $X \setminus U$ 的一个不可约分支的一般点。则逆像 $U_x$，即 $U$ 在
$\Spec(\mathcal{O}_{X, x})$ 中的逆像，是 $\mathcal{O}_{X, x}$ 的穿孔谱。
由假设，$V_x = V \times_U U_x$ 是有限 \'etale 态射
$Y_x \to \Spec(\mathcal{O}_{X, x})$ 到 $U_x$ 的限制。由《极限》引理
\ref{limits-lemma-glueing-near-point}，可以找到开子概形 $U \subset U' \subset X$，
它包含 $x$，以及有限表示态射 $V' \to U'$，使得它在 $U$ 上的限制
恢复 $V \to U$，并且它在 $\Spec(\mathcal{O}_{X, x})$ 上的限制恢复 $Y_x$。
最后，态射 $V' \to U'$ 在可能把 $U'$ 缩小为更小的开子概形后是有限 \'etale 的，
见《极限》引理 \ref{limits-lemma-glueing-near-point-properties}。
\end{proof}
















\section{穿孔谱的有限 \'etale 覆盖，II}
\label{pione-section-pi1-punctured-spec-II}

\noindent
本节证明 Section \ref{pione-section-pi1-punctured-spec} 中所讨论内容的一些变式。
设有 Noether 局部环 $(A, \mathfrak m)$ 和 $f \in \mathfrak m$。置
$X = \Spec(A)$ 和 $X_0 = \Spec(A/fA)$，并令
$U = X \setminus \{\mathfrak m\}$ 和 $U_0 = X_0 \setminus \{\mathfrak m\}$
分别为 $A$ 和 $A/fA$ 的穿孔谱。所有这些恰如情形
\ref{pione-situation-local-lefschetz}。不同之处在于，我们将考虑限制函子
$$
\colim_{U_0 \subset U' \subset U\text{ 开}} \textit{F\'Et}_{U'}
\longrightarrow
\textit{F\'Et}_{U_0}
$$
换言之，我们不会试图把 $U_0$ 的有限 \'etale 覆盖提升到整个 $U$，而只提升到
某个开邻域 $U'$，即 $U_0$ 在 $U$ 中的一个开邻域。

\begin{lemma}
\label{pione-lemma-faithful-general}
在情形 \ref{pione-situation-local-lefschetz} 下，设 $U' \subset U$ 是包含 $U_0$ 的开子集。
假设对素理想 $\mathfrak p \subset A$，若点 $\mathfrak p \in U'$ 是极小点，且
$\mathfrak p \not \in U_0$，则
$\dim(A/\mathfrak p) \geq 2$。那么
$$
\textit{F\'Et}_{U'} \longrightarrow \textit{F\'Et}_{U_0},\quad
V' \longmapsto V_0 = V' \times_{U'} U_0
$$
是忠实函子。此外，存在满足该假设的 $U'$；而任意更小开子集
$U'' \subset U'$，若包含 $U_0$，也满足该假设。特别地，限制函子
$$
\colim_{U_0 \subset U' \subset U\text{ 开}} \textit{F\'Et}_{U'}
\longrightarrow
\textit{F\'Et}_{U_0}
$$
是忠实的。
\end{lemma}

\begin{proof}
由《代数》引理 \ref{algebra-lemma-one-equation} 可知，$V(\mathfrak p)$ 与
$U_0$ 相交，其中 $\mathfrak p$ 是 $A$ 的任意满足
$\dim(A/\mathfrak p) \geq 2$ 的素理想。因此，由引理
\ref{pione-lemma-restriction-faithful}，陈述中的这种 $U$
所对应的展示函子是忠实的。为了看出这样的 $U'$ 存在，注意若
$\mathfrak p \subset A$ 满足 $\mathfrak p \in U$、$\mathfrak p \not \in U_0$ 和
$\dim(A/\mathfrak p) = 1$，则 $\mathfrak p$ 对应于 $U$ 的一个闭点，因而
$V(\mathfrak p) \cap U_0 = \emptyset$。因此，可以取 $U'$ 为 $X$ 的那些不与
$U_0$ 相交且维数为 $1$ 的不可约分支之补集。
\end{proof}

\begin{lemma}
\label{pione-lemma-fully-faithful-general-better}
在情形 \ref{pione-situation-local-lefschetz} 下，假设
\begin{enumerate}
\item $A$ 有对偶化复形，且关于 $f$-进拓扑完备；
\item $X$ 的每个不包含于 $X_0$ 的不可约分支的维数都 $\geq 3$。
\end{enumerate}
则限制函子
$$
\colim_{U_0 \subset U' \subset U\text{ 开}} \textit{F\'Et}_{U'}
\longrightarrow
\textit{F\'Et}_{U_0}
$$
全忠实。
\end{lemma}

\begin{proof}
由基本群的拓扑不变性，可以在证明中把 $A$ 替换为其约化；见引理
\ref{pione-lemma-thickening}。于是，结论由引理
\ref{pione-lemma-restriction-fully-faithful-general} 和《代数化与形式几何》引理
\ref{algebraization-lemma-fully-faithful-general} 得出。
\end{proof}

\begin{lemma}
\label{pione-lemma-fully-faithful-general}
在情形 \ref{pione-situation-local-lefschetz} 下，假设
\begin{enumerate}
\item $A$ 关于 $f$-进拓扑完备；
\item $f$ 是非零因子；
\item $H^1_\mathfrak m(A/fA)$ 是有限 $A$-模。
\end{enumerate}
则限制函子
$$
\colim_{U_0 \subset U' \subset U\text{ 开}} \textit{F\'Et}_{U'}
\longrightarrow
\textit{F\'Et}_{U_0}
$$
全忠实。
\end{lemma}

\begin{proof}
这由引理 \ref{pione-lemma-restriction-fully-faithful-general} 和《代数化与形式几何》引理
\ref{algebraization-lemma-fully-faithful-general-alternative} 得出。
\end{proof}








\section{穿孔谱的有限 \'etale 覆盖，III}
\label{pione-section-pi1-punctured-spec-III}

\noindent
本节研究在情形 \ref{pione-situation-local-lefschetz} 下，限制函子
$$
\colim_{U_0 \subset U' \subset U\text{ 开}} \textit{F\'Et}_{U'}
\longrightarrow
\textit{F\'Et}_{U_0}
$$
何时为范畴等价。

\begin{lemma}
\label{pione-lemma-essentially-surjective-general-better}
在情形 \ref{pione-situation-local-lefschetz} 下，假设
\begin{enumerate}
\item $A$ 有对偶化复形，且关于 $f$-进拓扑完备；
\item 下列条件之一成立：
\begin{enumerate}
\item $A_f$ 满足 $(S_2)$，且 $X$ 的每个不包含于 $X_0$ 的不可约分支的维数都
$\geq 4$；或
\item 若 $\mathfrak p \not \in V(f)$ 且
$V(\mathfrak p) \cap V(f) \not = \{\mathfrak m\}$，则
$\text{depth}(A_\mathfrak p) + \dim(A/\mathfrak p) > 3$。
\end{enumerate}
\end{enumerate}
则限制函子
$$
\colim_{U_0 \subset U' \subset U\text{ 开}} \textit{F\'Et}_{U'}
\longrightarrow
\textit{F\'Et}_{U_0}
$$
为等价。
\end{lemma}

\begin{proof}
这由引理 \ref{pione-lemma-restriction-equivalence-general} 和《代数化与形式几何》引理
\ref{algebraization-lemma-equivalence-better} 得出。
\end{proof}

\begin{lemma}
\label{pione-lemma-essentially-surjective-general}
在情形 \ref{pione-situation-local-lefschetz} 下，假设
\begin{enumerate}
\item $A$ 关于 $f$-进拓扑完备；
\item $f$ 是非零因子；
\item $H^1_\mathfrak m(A/fA)$ 和 $H^2_\mathfrak m(A/fA)$ 都是有限 $A$-模。
\end{enumerate}
则限制函子
$$
\colim_{U_0 \subset U' \subset U\text{ 开}} \textit{F\'Et}_{U'}
\longrightarrow
\textit{F\'Et}_{U_0}
$$
为等价。
\end{lemma}

\begin{proof}
这由引理 \ref{pione-lemma-restriction-equivalence-general} 和《代数化与形式几何》引理
\ref{algebraization-lemma-equivalence} 得出。
\end{proof}

\begin{remark}
\label{pione-remark-combine}
设 $(A, \mathfrak m)$ 是完备 Noether 局部环，并设 $f \in \mathfrak m$ 非零。
假设 $A_f$ 满足 $(S_2)$，且 $\Spec(A)$ 的每个不可约分支的维数都 $\geq 4$。
那么，引理 \ref{pione-lemma-essentially-surjective-general-better} 告诉我们，范畴
$$
\colim\nolimits_{U' \subset U\text{ open, }U_0 \subset U}
\text{ category of schemes finite \'etale over }U'
$$
等价于在 $U_0$ 上有限 \'etale 的概形范畴。例如，当 $A$ 是维数 $\geq 4$
的正规整环时，这一结论成立！
\end{remark}




\section{穿孔谱的有限 \'etale 覆盖，IV}
\label{pione-section-pi1-punctured-spec-IV}

\noindent
设 $X, X_0, U, U_0$ 如情形 \ref{pione-situation-local-lefschetz}。本节考察限制函子
$$
\textit{F\'Et}_U
\longrightarrow
\textit{F\'Et}_{U_0}
$$
何时本质满。我们将先取用 Section \ref{pione-section-pi1-punctured-spec-III} 的结果，
再用纯性填补其间的缺口。回顾：若对 Noether 局部环 $(A, \mathfrak m)$，
限制函子 $\textit{F\'Et}_X \to \textit{F\'Et}_U$ 本质满，其中
$X = \Spec(A)$ 且 $U = X \setminus \{\mathfrak m\}$，则称该环{\it 满足纯性}。

\begin{lemma}
\label{pione-lemma-equivalence-better}
在情形 \ref{pione-situation-local-lefschetz} 下，假设
\begin{enumerate}
\item $A$ 有对偶化复形，且关于 $f$-进拓扑完备；
\item 下列条件之一成立：
\begin{enumerate}
\item $A_f$ 满足 $(S_2)$，且 $X$ 的每个不包含于 $X_0$ 的不可约分支的维数都
$\geq 4$；或
\item 若 $\mathfrak p \not \in V(f)$ 且
$V(\mathfrak p) \cap V(f) \not = \{\mathfrak m\}$，则
$\text{depth}(A_\mathfrak p) + \dim(A/\mathfrak p) > 3$。
\end{enumerate}
\item 对每个极大理想 $\mathfrak p \subset A_f$，纯性对
$(A_f)_\mathfrak p$ 成立。
\end{enumerate}
则限制函子 $\textit{F\'Et}_U \to \textit{F\'Et}_{U_0}$ 本质满。
\end{lemma}

\begin{proof}
设 $V_0 \to U_0$ 是有限 \'etale 态射。由引理
\ref{pione-lemma-essentially-surjective-general-better}，存在开子集 $U' \subset U$，
它包含 $U_0$，以及有限 \'etale 态射 $V' \to U$，其到 $U_0$ 的基变换同构于
$V_0 \to U_0$。由于 $U' \supset U_0$，可见 $U \setminus U'$ 由与 (3) 中的素理想
$\mathfrak p_1, \ldots, \mathfrak p_n$ 对应的点组成。由假设，可以找到有限
\'etale 态射 $V'_i \to \Spec(A_{\mathfrak p_i})$，它与 $V' \to U'$ 一致
（在 $U' \times_U \Spec(A_{\mathfrak p_i})$ 上）。把《极限》引理
\ref{limits-lemma-glueing-near-closed-point} 应用 $n$ 次，可知 $V' \to U'$ 延拓为
有限 \'etale 态射 $V \to U$。
\end{proof}

\begin{lemma}
\label{pione-lemma-equivalence}
设 $(A, \mathfrak m)$ 是 Noether 局部环。设 $f \in \mathfrak m$。假设
\begin{enumerate}
\item $A$ 关于 $f$-进拓扑完备；
\item $f$ 是非零因子；
\item $H^1_\mathfrak m(A/fA)$ 和 $H^2_\mathfrak m(A/fA)$ 都是有限 $A$-模；
\item 对每个极大理想 $\mathfrak p \subset A_f$，纯性对
$(A_f)_\mathfrak p$ 成立。
\end{enumerate}
则限制函子 $\textit{F\'Et}_U \to \textit{F\'Et}_{U_0}$ 本质满。
\end{lemma}

\begin{proof}
证明与引理 \ref{pione-lemma-equivalence-better} 的证明完全相同，只需使用引理
\ref{pione-lemma-essentially-surjective-general} 代替引理
\ref{pione-lemma-essentially-surjective-general-better}。
\end{proof}



\section{局部情形中的纯性，II}
\label{pione-section-local-purity-II}

\noindent
本节是 Section \ref{pione-section-local-purity} 的续篇。回顾：若对 Noether 局部环
$(A, \mathfrak m)$，限制函子 $\textit{F\'Et}_X \to \textit{F\'Et}_U$
本质满，其中 $X = \Spec(A)$ 且 $U = X \setminus \{\mathfrak m\}$，
则称该环{\it 满足纯性}。

\begin{lemma}
\label{pione-lemma-purity-inherited-by-hypersurface-better}
设 $(A, \mathfrak m)$ 是 Noether 局部环。设 $f \in \mathfrak m$。假设
\begin{enumerate}
\item $A$ 有对偶化复形，且关于 $f$-进拓扑完备；
\item 下列条件之一成立：
\begin{enumerate}
\item $A_f$ 满足 $(S_2)$，且 $X$ 的每个不包含于 $X_0$ 的不可约分支的维数都
$\geq 4$；或
\item 若 $\mathfrak p \not \in V(f)$ 且
$V(\mathfrak p) \cap V(f) \not = \{\mathfrak m\}$，则
$\text{depth}(A_\mathfrak p) + \dim(A/\mathfrak p) > 3$。
\end{enumerate}
\item 对每个极大理想 $\mathfrak p \subset A_f$，纯性对
$(A_f)_\mathfrak p$ 成立；并且
\item 纯性对 $A$ 成立。
\end{enumerate}
则纯性对 $A/fA$ 成立。
\end{lemma}

\begin{proof}
置 $X = \Spec(A)$，并令 $U = X \setminus \{\mathfrak m\}$ 为穿孔谱。类似地，
有 $X_0 = \Spec(A/fA)$ 和 $U_0 = X_0 \setminus \{\mathfrak m\}$。设
$V_0 \to U_0$ 是有限 \'etale 态射。由引理 \ref{pione-lemma-equivalence-better}，
可得有限 \'etale 态射 $V \to U$，其到 $U_0$ 的基变换同构于
$V_0 \to U_0$。由假设 (5)，$V \to U$ 延拓为有限 \'etale 态射
$Y \to X$。于是，$Y$ 到 $X_0$ 的限制就是 $V_0 \to U_0$ 所需的延拓。
\end{proof}

\begin{lemma}
\label{pione-lemma-purity-inherited-by-hypersurface}
设 $(A, \mathfrak m)$ 是 Noether 局部环。设 $f \in \mathfrak m$。假设
\begin{enumerate}
\item $A$ 关于 $f$-进拓扑完备；
\item $f$ 是非零因子；
\item $H^1_\mathfrak m(A/fA)$ 和 $H^2_\mathfrak m(A/fA)$ 都是有限 $A$-模；
\item 对每个极大理想 $\mathfrak p \subset A_f$，纯性对
$(A_f)_\mathfrak p$ 成立；
\item 纯性对 $A$ 成立。
\end{enumerate}
则纯性对 $A/fA$ 成立。
\end{lemma}

\begin{proof}
证明与引理 \ref{pione-lemma-purity-inherited-by-hypersurface-better} 的证明完全相同，
只需使用引理 \ref{pione-lemma-equivalence} 代替引理 \ref{pione-lemma-equivalence-better}。
\end{proof}

\noindent
现在可以利用前面的结果，证明纯性对维数 $\geq 3$ 的完全交成立。回顾：若一个
Noether 局部环的完备化，是某个正则局部环模一条正则序列所生成理想的商，
则称该环为完全交。见《除幂代数》Section \ref{dpa-section-lci} 中的讨论。

\begin{proposition}
\label{pione-proposition-purity-complete-intersection}
设 $(A, \mathfrak m)$ 是 Noether 局部环。若 $A$ 是维数 $\geq 3$ 的完全交，
则纯性对 $A$ 成立；也就是说，穿孔谱的任意有限 \'etale 覆盖都可延拓。
\end{proposition}

\begin{proof}
由引理 \ref{pione-lemma-purity-and-completion}，可设 $A$ 是完备局部环。由假设，可以写成
$A = B/(f_1, \ldots, f_r)$，其中 $B$ 是完备正则局部环，且
$f_1, \ldots, f_r$ 是正则序列。我们对 $r$ 作归纳以完成证明。基本情形为
$r = 0$；这由引理 \ref{pione-lemma-local-purity} 得出，该引理适用于维数 $\geq 2$
的正则环。

\medskip\noindent
假设 $A = B/(f_1, \ldots, f_r)$，且命题对 $r - 1$ 成立。置
$A' = B/(f_1, \ldots, f_{r - 1})$，并把引理
\ref{pione-lemma-purity-inherited-by-hypersurface} 应用于 $f_r \in A'$。这是允许的：
条件 (1) 成立，因为 $f_1, \ldots, f_r$ 是正则序列；
条件 (2) 成立，因为 $B$ 完备，从而 $A'$ 完备；
条件 (3) 成立，因为 $A = A'/f_r A'$ 是维数为 $\dim(A) \geq 3$ 的
Cohen--Macaulay 环，见《对偶化复形》引理 \ref{dualizing-lemma-depth}；
条件 (4) 由归纳假设成立，因为 $\dim((A'_{f_r})_\mathfrak p) \geq 3$，其中
$\mathfrak p$ 是 $A'_{f_r}$ 的一个极大素理想，并且
$(A'_{f_r})_\mathfrak p = B_\mathfrak q/(f_1, \ldots, f_{r - 1})$ 对某个
$\mathfrak q \subset B$ 成立；
条件 (5) 由归纳假设成立。
\end{proof}




\section{局部情形中的纯性，III}
\label{pione-section-local-purity-III}

\noindent
本节继续 Sections \ref{pione-section-local-purity} 和
\ref{pione-section-local-purity-II} 中的讨论。

\begin{lemma}
\label{pione-lemma-fully-faithful-power-series-over-depth2}
设 $(A, \mathfrak m)$ 是深度 $\geq 2$ 的 Noether 局部环。设
$B = A[[x_1, \ldots, x_d]]$，其中 $d \geq 1$。置
$Y = \Spec(B)$ 和 $Y_0 = V(x_1, \ldots, x_d)$。对任意开子概形
$V \subset Y$，若 $V_0 = V \cap Y_0$ 等于
$Y_0 \setminus \{\mathfrak m_B\}$，则限制函子
$$
\textit{F\'Et}_V \longrightarrow \textit{F\'Et}_{V_0}
$$
全忠实。
\end{lemma}

\begin{proof}
置 $I = (x_1, \ldots, x_d)$。置 $X = \Spec(A)$。利用态射 $Y \to X$
把 $Y_0$ 与 $X$ 等同后，$V_0$ 就与穿孔谱 $U$ 等同，即 $A$ 的穿孔谱。
沿这个仿射态射
前推模，得到
\begin{align*}
\lim_n \Gamma(V_0, \mathcal{O}_V/I^n\mathcal{O}_V)
& =
\lim_n \Gamma(V_0, \mathcal{O}_Y/I^n\mathcal{O}_Y) \\
& =
\lim_n \Gamma(U, \mathcal{O}_U[x_1, \ldots, x_d]/(x_1, \ldots, x_d)^n) \\
& =
\lim_n A[x_1, \ldots, x_d]/(x_1, \ldots, x_d)^n \\
& =
B
\end{align*}
事实上，由于 $A$ 的深度 $\geq 2$，有 $\Gamma(U, \mathcal{O}_U) = A$；
见《局部上同调》引理
\ref{local-cohomology-lemma-finiteness-pushforwards-and-H1-local}。
因此，对引理中任意开子集 $V \subset Y$，有
$$
B = \Gamma(Y, \mathcal{O}_Y) \to \Gamma(V, \mathcal{O}_V) \to
\lim_n \Gamma(V_0, \mathcal{O}_Y/I^n\mathcal{O}_Y) = B
$$
这蕴含两个箭头都是同构（略去一个小细节）。由《代数化与形式几何》引理
\ref{algebraization-lemma-completion-fully-faithful}，可知
$\textit{Coh}(\mathcal{O}_V) \to \textit{Coh}(V, I\mathcal{O}_V)$
在有限局部自由对象所成的满子范畴上全忠实。因此，由引理
\ref{pione-lemma-restriction-fully-faithful} 即得结论。
\end{proof}

\begin{lemma}
\label{pione-lemma-purity-power-series-over-depth2}
\begin{slogan}
有限 \'etale 覆盖的 Ramanujam--Samuel 定理
\end{slogan}
设 $(A, \mathfrak m)$ 是深度 $\geq 2$ 的 Noether 局部环。设
$B = A[[x_1, \ldots, x_d]]$，其中 $d \geq 1$。对任意开子集
$V \subset Y = \Spec(B)$，若它包含
\begin{enumerate}
\item 任意素理想 $\mathfrak q \subset B$，只要
$\mathfrak q \cap A \not = \mathfrak m$；
\item 素理想 $\mathfrak m B$，
\end{enumerate}
则函子
$
\textit{F\'Et}_Y
\to
\textit{F\'Et}_V
$
为等价。特别地，纯性对 $B$ 成立。
\end{lemma}

\begin{proof}
素理想 $\mathfrak q \subset B$ 若不属于 $V$，则位于 $\mathfrak m$ 上方。在这种情形下，
$A \to B_\mathfrak q$ 是平坦局部同态，因而
$\text{depth}(B_\mathfrak q) \geq 2$（《代数》引理
\ref{algebra-lemma-apply-grothendieck}）。因此，结合引理
\ref{pione-lemma-quasi-compact-dense-open-connected-at-infinity-Noetherian} 与《局部上同调》引理
\ref{local-cohomology-lemma-depth-2-connected-punctured-spectrum}，该函子全忠实。

\medskip\noindent
设 $W \to V$ 是有限 \'etale 态射。设 $B \to C$ 是唯一的有限环同态，使得
$\Spec(C) \to Y$ 是引理 \ref{pione-lemma-extend-S2} 中构造的、延拓 $W \to V$
的有限态射。注意，$C = \Gamma(W, \mathcal{O}_W)$。

\medskip\noindent
置 $Y_0 = V(x_1, \ldots, x_d)$ 和 $V_0 = V \cap Y_0$。置 $X = \Spec(A)$。
利用态射 $Y \to X$ 把 $Y_0$ 与 $X$ 等同后，$V_0$ 就与穿孔谱 $U$ 等同，
即 $A$ 的穿孔谱。因此，可把 $W_0 = W \times_Y Y_0$ 看作 $U$ 上的有限 \'etale 概形。于是
$$
W_0 \times_U (U \times_X Y)
\quad\text{且}\quad
W \times_V (U \times_X Y)
$$
是 $U \times_X Y$ 上的有限 \'etale 概形，它们在 $V_0$ 上的限制同构。把引理
\ref{pione-lemma-fully-faithful-power-series-over-depth2} 应用于开子集 $U \times_X Y$，
得到同构
$$
W_0 \times_U (U \times_X Y) \longrightarrow W \times_V (U \times_X Y)
$$
（在 $U \times_X Y$ 上）。

\medskip\noindent
注意，$C_0 = \Gamma(W_0, \mathcal{O}_{W_0})$ 是有限 $A$-代数；这是把引理
\ref{pione-lemma-extend-S2} 应用于 $W_0 \to U \subset X$ 所得（与上面对
$B \to C$ 所做的完全相同）。由于引理
\ref{pione-lemma-extend-S2} 中的构造与平坦基变换及开集的改变相容，上面的同构诱导同构
$$
\Psi : C \longrightarrow C_0 \otimes_A B
$$
（作为有限 $B$-代数）。然而，我们知道，$\Spec(C) \to Y$ 在位于 $Y$
中至少一个点上方的所有点处是 \'etale 的，而该点位于 $\mathfrak m \in X$
上方。由于 $\Psi$ 是同构，可知 $\Spec(C_0) \to X$ 在 $\mathfrak m$
上方是 \'etale 的（略去一个小细节）。当然，这意味着 $A \to C_0$ 有限
\'etale，因而 $B \to C$ 有限 \'etale。
\end{proof}

\begin{lemma}
\label{pione-lemma-purity-smooth-over-depth2}
设 $f : X \to S$ 是概形的态射。设 $U \subset X$ 是开子概形。假设
\begin{enumerate}
\item $f$ 光滑；
\item $S$ 是 Noether 概形；
\item 对 $s \in S$，若 $\text{depth}(\mathcal{O}_{S, s}) \leq 1$，则有
$X_s = U_s$；
\item $U_s \subset X_s$ 对所有 $s \in S$ 都稠密。
\end{enumerate}
则 $\textit{F\'Et}_X \to \textit{F\'Et}_U$ 为等价。
\end{lemma}

\begin{proof}
结合引理 \ref{pione-lemma-quasi-compact-dense-open-connected-at-infinity-Noetherian} 与
《局部上同调》引理 \ref{local-cohomology-lemma-depth-2-connected-punctured-spectrum}，
该函子全忠实（此外应用《代数》引理
\ref{algebra-lemma-apply-grothendieck} 来检验深度条件）。

\medskip\noindent
设 $\pi : V \to U$ 是有限 \'etale 态射。设 $Y \to X$ 是引理
\ref{pione-lemma-extend-S2} 中构造的有限态射。必须证明 $Y \to X$ 有限 \'etale。
为了证明这一点对所有 $x \in X$ 成立，固定其像 $s \in S$；可以用 Noether
概形的平坦态射 $S' \to S$ 作基变换，只要 $s$ 位于其像中。这由引理
\ref{pione-lemma-extend-S2} 中的构造与平坦基变换相容得出。

\medskip\noindent
扩大 $U$ 后，可设 $U \subset X$ 是使 $Y \to X$ 有限 \'etale 的最大开集。
令 $Z \subset X$ 为 $U$ 的补集。为导出矛盾，假设 $Z \not = \emptyset$。
选取一点 $s \in S$，它位于 $Z \to S$ 的像中，并且 $s$ 的任何严格一般化都不在该像中。
对 $\Spec(\mathcal{O}_{S, s})$ 作基变换后，可知 $S = \Spec(A)$，其中
$(A, \mathfrak m, \kappa)$ 是深度 $\geq 2$ 的 Noether 局部环，且 $Z$
包含于闭纤维 $X_s$ 中并在 $X_s$ 中无处稠密。选取闭点 $z \in Z$。于是
$\kappa(z)/\kappa$ 有限（由 Hilbert 零点定理；见《代数》定理
\ref{algebra-theorem-nullstellensatz}）。选取局部概形的有限平坦态射
$(S', s') \to (S, s)$ 来实现剩余域扩张 $\kappa(z)/\kappa$；见《代数》引理
\ref{algebra-lemma-finite-free-given-residue-field-extension}。由 $S' \to S$
作基变换后，约化到 $\kappa(z) = \kappa$ 的情形。

\medskip\noindent
由《态射补篇》引理 \ref{more-morphisms-lemma-slice-smooth}，存在局部闭子概形
$S' \subset X$，它经过 $z$，使得 $S' \to S$ 在 $z$ 处为 \'etale。由
$S' \to S$ 作基变换后，可设存在截面 $\sigma : S \to X$，满足
$\sigma(s) = z$。选取仿射邻域 $\Spec(B) \subset X$，它包含 $s$。于是
$A \to B$ 是光滑环同态，并有截面 $\sigma : B \to A$。记
$I = \Ker(\sigma)$，并以 $B^\wedge$ 表示关于 $I$ 的 $B$ 之完备化。于是有
$B^\wedge \cong A[[x_1, \ldots, x_d]]$，其中 $d \geq 0$；见《代数》引理
\ref{algebra-lemma-section-smooth}。注意 $d > 0$，否则 $X \to S$ 在 $z$
处为 \'etale，这将蕴含 $z$ 是 $X_s$ 的一般点，继而由假设 (4) 有
$z \in U$。类似地，若 $d > 0$，则 $\mathfrak m B^\wedge$ 映入 $U$，这里所用态射为
$\Spec(B^\wedge) \to X$。只需证明 $Y \to X$ 在基变换到
$\Spec(B^\wedge)$ 后有限 \'etale。由于 $B \to B^\wedge$ 平坦（《代数》引理
\ref{algebra-lemma-completion-flat}），这由引理
\ref{pione-lemma-purity-power-series-over-depth2} 和 $Y \to X$ 构造中的唯一性得出。
\end{proof}

\begin{proposition}
\label{pione-proposition-purity-smooth-over-depth2}
设 $A \to B$ 是 Noether 局部环之间的局部同态。假设 $A$ 的深度 $\geq 2$，
$A \to B$ 关于 $\mathfrak m_B$-进拓扑形式光滑，且 $\dim(B) > \dim(A)$。
对任意开子集 $V \subset Y = \Spec(B)$，若它包含
\begin{enumerate}
\item 任意素理想 $\mathfrak q \subset B$，只要
$\mathfrak q \cap A \not = \mathfrak m_A$；
\item 素理想 $\mathfrak m_A B$，
\end{enumerate}
则函子 $\textit{F\'Et}_Y \to \textit{F\'Et}_V$ 为等价。特别地，纯性对 $B$ 成立。
\end{proposition}

\begin{proof}
素理想 $\mathfrak q \subset B$ 若不属于 $V$，则位于 $\mathfrak m_A$ 上方。
在这种情形下，$A \to B_\mathfrak q$ 是平坦局部同态，因而
$\text{depth}(B_\mathfrak q) \geq 2$（《代数》引理
\ref{algebra-lemma-apply-grothendieck}）。因此，结合引理
\ref{pione-lemma-quasi-compact-dense-open-connected-at-infinity-Noetherian} 与《局部上同调》引理
\ref{local-cohomology-lemma-depth-2-connected-punctured-spectrum}，该函子全忠实。

\medskip\noindent
以 $A^\wedge$ 和 $B^\wedge$ 表示 $A$ 和 $B$ 关于各自极大理想的完备化。注意，
命题的假设对 $A^\wedge \to B^\wedge$ 成立；见《代数补篇》引理
\ref{more-algebra-lemma-completion-dimension}、
\ref{more-algebra-lemma-completion-depth} 和
\ref{more-algebra-lemma-formally-smooth-completion}。由引理
\ref{pione-lemma-extend-S2} 中构造的唯一性及其与平坦基变换的相容性，只需对
$A^\wedge \to B^\wedge$ 和 $V$ 的逆像证明本质满性（略去细节；当 $V$
为穿孔谱时，可与引理 \ref{pione-lemma-purity-and-completion} 比较）。由《代数补篇》命题
\ref{more-algebra-proposition-fs-regular}，这意味着可设 $A \to B$ 正则。

\medskip\noindent
设 $W \to V$ 是有限 \'etale 态射。由 Popescu 定理（《环同态的光滑化》定理
\ref{smoothing-theorem-popescu}），可把 $B = \colim B_i$ 写成光滑 $A$-代数的
滤过余极限。可以选取某个 $i$ 和开子集 $V_i \subset \Spec(B_i)$，使其逆像为
$V$（《极限》引理 \ref{limits-lemma-descend-opens}）。增大 $i$ 后，可设存在有限
\'etale 态射 $W_i \to V_i$，其到 $V$ 的基变换为 $W \to V$；见《极限》引理
\ref{limits-lemma-descend-finite-presentation}、
\ref{limits-lemma-descend-finite-finite-presentation} 和
\ref{limits-lemma-descend-etale}。可以假设 $V_i$ 的补集包含于
$\Spec(B_i) \to \Spec(A)$ 的闭纤维，因为这一点对 $V$ 成立（或者如此选取
$V_i$，或者使用上面的引理证明它对充分大的 $i$ 成立）。设 $\eta$ 是
$\Spec(B) \to \Spec(A)$ 之闭纤维的一般点。由于 $\eta \in V$，$\eta$ 的像位于
$V_i$。因此，把 $V_i$ 替换为 $\Spec(B)$ 之闭点的像的一个仿射开邻域后，
可设 $\Spec(B_i) \to \Spec(A)$ 的闭纤维不可约，且其一般点包含于 $V_i$
（略去细节；使用域上光滑的概形是不可约概形的不交并这一事实）。此时可以应用引理
\ref{pione-lemma-purity-smooth-over-depth2}，看出 $W_i \to V_i$ 延拓为有限 \'etale 态射
$\Spec(C_i) \to \Spec(B_i)$。把它拉回到 $\Spec(B)$，便可得出 $W$ 位于函子
$\textit{F\'Et}_Y \to \textit{F\'Et}_V$ 的本质像中，正如所需。
\end{proof}





\section{基本群的 Lefschetz 定理}
\label{pione-section-global-lefschetz}

\noindent
当然，我们已经在局部情形中证明了许多这种类型的结果。本节讨论射影情形。

\begin{proposition}
\label{pione-proposition-lefschetz-fully-faithful}
设 $k$ 是域。设 $X$ 是 $k$ 上的固有概形。设 $\mathcal{L}$ 是丰沛可逆
$\mathcal{O}_X$-模。设 $s \in \Gamma(X, \mathcal{L})$。设 $Y = Z(s)$ 是
$s$ 的零点概形。假设对所有 $x \in X \setminus Y$，都有
$$
\text{depth}(\mathcal{O}_{X, x}) + \dim(\overline{\{x\}}) > 1
$$
则限制函子 $\textit{F\'Et}_X \to \textit{F\'Et}_Y$ 全忠实。事实上，
对任意开子概形 $V \subset X$，若它包含 $Y$，则限制函子
$\textit{F\'Et}_V \to \textit{F\'Et}_Y$ 全忠实。
\end{proposition}

\begin{proof}
第一个断言是引理 \ref{pione-lemma-restriction-fully-faithful-special} 和《代数化与形式几何》命题
\ref{algebraization-proposition-lefschetz} 的形式推论。第二个断言由引理
\ref{pione-lemma-restriction-fully-faithful-special} 和《代数化与形式几何》引理
\ref{algebraization-lemma-lefschetz-addendum} 得出。
\end{proof}

\begin{proposition}
\label{pione-proposition-lefschetz-equivalence-general}
设 $k$ 是域。设 $X$ 是 $k$ 上的固有概形。设 $\mathcal{L}$ 是丰沛可逆
$\mathcal{O}_X$-模。设 $s \in \Gamma(X, \mathcal{L})$。设 $Y = Z(s)$ 是
$s$ 的零点概形。令 $\mathcal{V}$ 为 $X$ 中包含 $Y$ 的开子概形所成集合，
按逆包含关系排序。假设对所有 $x \in X \setminus Y$，都有
$$
\text{depth}(\mathcal{O}_{X, x}) + \dim(\overline{\{x\}}) > 2
$$
则限制函子
$$
\colim_\mathcal{V} \textit{F\'Et}_V \to \textit{F\'Et}_Y
$$
为等价。
\end{proposition}

\begin{proof}
这是引理 \ref{pione-lemma-restriction-equivalence-general} 和《代数化与形式几何》命题
\ref{algebraization-proposition-lefschetz-equivalence} 的形式推论。
\end{proof}

\begin{proposition}
\label{pione-proposition-lefschetz-equivalence}
设 $k$ 是域。设 $X$ 是 $k$ 上的固有概形。设 $\mathcal{L}$ 是丰沛可逆
$\mathcal{O}_X$-模。设 $s \in \Gamma(X, \mathcal{L})$。设 $Y = Z(s)$ 是
$s$ 的零点概形。假设对所有 $x \in X \setminus Y$，都有
$$
\text{depth}(\mathcal{O}_{X, x}) + \dim(\overline{\{x\}}) > 2
$$
并且对闭点 $x \in X \setminus Y$，纯性对 $\mathcal{O}_{X, x}$ 成立。则限制函子
$\textit{F\'Et}_X \to \textit{F\'Et}_Y$ 为等价。若 $X$（等价地，$Y$）连通，则
$$
\pi_1(Y, \overline{y}) \to \pi_1(X, \overline{y})
$$
该态射对任意几何点 $\overline{y}$ 均为同构；这里该点取自 $Y$。
\end{proposition}

\begin{proof}
由命题 \ref{pione-proposition-lefschetz-fully-faithful}，全忠实性成立。由命题
\ref{pione-proposition-lefschetz-equivalence-general}，$\textit{F\'Et}_Y$ 的任意对象都同构于
纤维积 $U \times_V Y$，其中 $U \to V$ 是某个有限 \'etale 态射，且
$V \subset X$ 是包含 $Y$ 的开子概形。补集 $T = X \setminus V$
是\footnote{事实上，$T$ 在 $k$ 上固有（因为它在 $X$ 中闭），并且仿射
（因为它在仿射概形 $X \setminus Y$ 中闭；见《态射》引理
\ref{morphisms-lemma-proper-ample-delete-affine}），因此在 $k$ 上有限
（《态射》引理 \ref{morphisms-lemma-finite-proper}）。所以 $T$ 是闭点的有限集合。}
$X \setminus Y$ 中闭点的有限集合。写成 $T = \{x_1, \ldots, x_n\}$。
由假设，可以找到有限 \'etale 态射
$V'_i \to \Spec(\mathcal{O}_{X, x_i})$，它与 $U \to V$ 一致（在
$V \times_X \Spec(\mathcal{O}_{X, x_i})$ 上）。把《极限》引理
\ref{limits-lemma-glueing-near-closed-point} 应用 $n$ 次，可见 $U \to V$
按要求延拓为有限 \'etale 态射 $U' \to X$。最后一个断言见引理
\ref{pione-lemma-what-equivalence-gives}。
\end{proof}










\section{分歧轨迹的纯性}
\label{pione-section-purity-ramification}

\noindent
本节讨论泛有限态射情形下与分支轨迹纯性相对应的结果。该结果似乎归功于
Gabber。其一个特殊情形是 van der Waerden 的纯性定理：从正规簇到光滑簇的
双有理态射不为同构的轨迹具有纯性。

\begin{lemma}
\label{pione-lemma-characterize-rational-singularity}
设 $A$ 是维数为 $2$ 的 Noether 正规局部整环。假设 $A$ 是 Nagata 环，
具有对偶化模 $\omega_A$，并且存在奇点解消
$f : X \to \Spec(A)$。按《曲面奇点的消解》注记
\ref{resolve-remark-dualizing-setup} 取 $\omega_X$。若
$\omega_X \cong \mathcal{O}_X(E)$，其中 $E \subset X$ 是某个支撑在例外纤维上的
有效 Cartier 除子，则 $A$ 定义一个有理奇点。
若 $f$ 是极小解消，则 $E = 0$。
\end{lemma}

\begin{proof}
存在迹映射 $Rf_*\omega_X \to \omega_A$，见《概形的对偶性》节
\ref{duality-section-trace}。由 Grauert--Riemenschneider 定理
（《曲面奇点的消解》命题
\ref{resolve-proposition-Grauert-Riemenschneider}），有
$R^1f_*\omega_X = 0$。因此迹映射就是映射
$f_*\omega_X \to \omega_A$。于是可以考虑
$$
\mathcal{O}_{\Spec(A)} = f_*\mathcal{O}_X \to f_*\omega_X \to \omega_A
$$
其中第一个映射来自引理陈述中假设存在的映射
$\mathcal{O}_X \to \mathcal{O}_X(E) = \omega_X$。由《除子》引理
\ref{divisors-lemma-check-isomorphism-via-depth-and-ass}，该复合映射是同构：
它在 $A$ 的穿孔谱上是同构（由引理中的假设以及 $f$ 在穿孔谱上为同构）；
而 $A$ 与 $\omega_A$ 都是 $A$-模，且深度为 $2$（由《代数》引理
\ref{algebra-lemma-criterion-normal} 与《对偶化复形》引理
\ref{dualizing-lemma-depth-dualizing-module}）。因此
$f_*\omega_X \to \omega_A$ 满射，进而为同构。于是
$Rf_*\omega_X = \omega_A$；由对偶性，这蕴含
$Rf_*\mathcal{O}_X = \mathcal{O}_{\Spec(A)}$。从而
$H^1(X, \mathcal{O}_X) = 0$，这说明 $A$ 定义一个有理奇点（见《曲面奇点的消解》
节 \ref{resolve-section-bounded} 中的讨论，特别是引理
\ref{resolve-lemma-regular-rational} 与
\ref{resolve-lemma-exact-sequence}）。若 $f$ 极小，则 $E = 0$：反复应用
《曲面奇点的消解》引理
\ref{resolve-lemma-dualizing-blow-up-rational} 可知映射
$f^*\omega_A \to \omega_X$ 满射，而如上所见 $\omega_A \cong A$。
\end{proof}

\begin{lemma}
\label{pione-lemma-key-purity-ramification}
设 $f : X \to \Spec(A)$ 是有限型态射。设 $x \in X$ 是一点。假设
\begin{enumerate}
\item $A$ 是优良正则局部环；
\item $\mathcal{O}_{X, x}$ 正规且维数为 $2$；
\item $f$ 在 $\overline{\{x\}}$ 之外为 \'etale。
\end{enumerate}
则 $f$ 在 $x$ 处为 \'etale。
\end{lemma}

\begin{proof}
首先把 $X$ 换成 $x$ 的一个仿射开邻域。注意，
$\mathcal{O}_{X, x}$ 是优良局部环（《代数补篇》引理
\ref{more-algebra-lemma-finite-type-over-excellent}）。因此可以选取一个极小奇点解消
$W \to \Spec(\mathcal{O}_{X, x})$，见《曲面奇点的消解》定理
\ref{resolve-theorem-resolve}。必要时把 $X$ 换成 $x$ 的仿射开邻域，可以找到固有态射
$b : X' \to X$，使得
$X' \times_X \Spec(\mathcal{O}_{X, x}) = W$，见《极限》引理
\ref{limits-lemma-glueing-near-closed-point}。进一步缩小 $X$ 后，可设 $X'$ 正则。
事实上，已知 $W$ 正则且 $X'$ 优良，而优良环的谱的正则轨迹是开的。由于
$W \to \Spec(\mathcal{O}_{X, x})$ 是射影态射（它是一列正规化爆破），
缩小 $X$ 后可设 $b$ 射影（略去细节）。令
$U = X \setminus \overline{\{x\}}$。由于
$W \to \Spec(\mathcal{O}_{X, x})$ 在穿孔谱上为同构，可设
$b : X' \to X$ 在 $U$ 上为同构。因此也把 $U$ 看作 $X'$ 的开子概形。
置 $f' = f \circ b : X' \to \Spec(A)$。

\medskip\noindent
由于 $A$ 正则，可见 $\mathcal{O}_Y$ 是 $Y$ 的对偶化复形。因此
$f^!\mathcal{O}_Y$ 是 $X$ 上的对偶化复形（《概形的对偶性》引理
\ref{duality-lemma-shriek-dualizing}）。由《概形的对偶性》引理
\ref{duality-lemma-dualizing-module-CM-scheme}，$X$ 的 Cohen--Macaulay
轨迹是开的（也可以利用优良性证明这一点）。由于
$\mathcal{O}_{X, x}$ 是 Cohen--Macaulay 的，缩小 $X$ 后可设
$X$ 是 Cohen--Macaulay 的。注意，\'etale 态射是局部完全交。因此可以把
《概形的对偶性》引理
\ref{duality-lemma-fundamental-class-almost-lci} 用于 $r = 0$，得到映射
$$
\mathcal{O}_X \longrightarrow \omega_{X/Y} = H^0(f^!\mathcal{O}_Y)
$$
它在 $X \setminus \overline{\{x\}}$ 上为同构。由《概形的对偶性》引理
\ref{duality-lemma-shriek-over-CM}，$\omega_{X/Y}$ 满足 $(S_2)$，故由
《除子》引理 \ref{divisors-lemma-check-isomorphism-via-depth-and-ass}
可知该映射为同构。这已经说明 $X$ 是 Gorenstein 的，特别地，
$\mathcal{O}_{X, x}$ 是 Gorenstein 环。

\medskip\noindent
置 $\omega_{X'/Y} = H^0((f')^!\mathcal{O}_Y)$。与上面完全相同的论证表明
$(f')^!\mathcal{O}_Y = \omega_{X'/Y}[0]$ 是 $X'$ 的对偶化复形。由于
$X'$ 正则，态射 $X' \to Y$ 是局部完全交态射，见《态射补篇》引理
\ref{more-morphisms-lemma-morphism-regular-schemes-is-lci}。由《概形的对偶性》引理
\ref{duality-lemma-fundamental-class-lci}，存在映射
$$
\mathcal{O}_{X'} \longrightarrow \omega_{X'/Y}
$$
它在 $U$ 上为同构。于是
$\omega_{X'/Y} = \mathcal{O}_{X'}(E)$，其中 $E \subset X'$ 是某个与
$U$ 不交的有效 Cartier 除子。

\medskip\noindent
由于 $\omega_{X/Y} = \mathcal{O}_Y$，可见
$\omega_{X'/Y} = b^! f^!\mathcal{O}_Y = b^!\mathcal{O}_X$。回到
$W \to \Spec(\mathcal{O}_{X, x})$，可得
$\omega_W = \mathcal{O}_W(E|_W)$。由引理
\ref{pione-lemma-characterize-rational-singularity}，有 $E|_W = 0$。
于是由上面已经用过的《概形的对偶性》引理
\ref{duality-lemma-fundamental-class-lci}，$f' : X' \to Y$ 为 \'etale。
这立即完成证明，因为 $f'$ 为 \'etale 迫使 $b$ 是同构。
\end{proof}

\begin{lemma}[分歧轨迹的纯性]
\label{pione-lemma-purity-ramification}
\begin{reference}
复空间情形下的这一结果见 \cite{Fischer} 第 170 页。一般情形是
\cite[Theorem 2.4]{Zong} 中归于 Gabber 的结果。
\end{reference}
设 $f : X \to Y$ 是局部 Noether 概形之间的态射。设 $x \in X$，并置
$y = f(x)$。假设
\begin{enumerate}
\item $\mathcal{O}_{X, x}$ 正规且维数 $\geq 1$；
\item $\mathcal{O}_{Y, y}$ 正则；
\item $f$ 是局部有限型的；并且
\item 对特化 $x' \leadsto x$，若
$\dim(\mathcal{O}_{X, x'}) = 1$，则态射 $f$ 在 $x'$ 处为 \'etale。
\end{enumerate}
则 $f$ 在 $x$ 处为 \'etale。
\end{lemma}

\begin{proof}
对 $d = \dim(\mathcal{O}_{X, x})$ 作归纳来证明该引理。

\medskip\noindent
先看平凡情形 $d = 1$；此时由假设，态射 $f$ 在 $x$ 处为 \'etale。
以下假设 $d \geq 2$。

\medskip\noindent
沿 $\Spec(\mathcal{O}_{Y, y}) \to Y$ 作基变换不会影响引理的结论，见
《态射》引理 \ref{morphisms-lemma-set-points-where-fibres-etale}。因此可设
$Y = \Spec(A)$，其中 $A$ 是正则局部环，而 $y$ 所对应的是极大理想
$\mathfrak m$（在 $A$ 中）。

\medskip\noindent
设 $x' \leadsto x$ 是满足 $x' \not = x$ 的特化。
$\mathcal{O}_{X, x'}$ 正规，因为它是 $\mathcal{O}_{X, x}$ 的局部化。若 $x'$ 不是
$X$ 的一般点，则 $1 \leq \dim(\mathcal{O}_{X, x'}) < d$，由归纳假设可知
$f$ 在 $x'$ 处为 \'etale。因此可设 $f$ 在所有特化到 $x$ 的点处均为
\'etale。由于 $f$ 为 \'etale 的点所成集合在 $X$ 中是开的（依定义），
把 $X$ 换成 $x$ 的一个开邻域后，可设 $f$ 在
$\overline{\{x\}}$ 之外为 \'etale。特别地，可见 $f$ 仅可能在位于闭点
$y \in Y = \Spec(A)$ 上方的点处不是 \'etale。

\medskip\noindent
令 $X' = X \times_{\Spec(A)} \Spec(A^\wedge)$。设 $x' \in X'$ 是位于
$x$ 上方的唯一一点。由上可见，$X'$ 在闭纤维之外于
$\Spec(A^\wedge)$ 上为 \'etale，因而 $X'$ 在闭纤维之外正规。由于
$X$ 正规，由《曲面奇点的消解》引理
\ref{resolve-lemma-normalization-completion} 可知 $X'$ 正规。于是，只要能证明
$X' \to \Spec(A^\wedge)$ 在 $x'$ 处为 \'etale，就可推出 $f$ 在 $x$
处为 \'etale（由前述《态射》引理
\ref{morphisms-lemma-set-points-where-fibres-etale}）。故可以并且确实假设
$A$ 是正则完备局部环。

\medskip\noindent
此时 $d = 2$ 的情形由引理 \ref{pione-lemma-key-purity-ramification} 得出。

\medskip\noindent
假设 $d > 2$。取 $t \in \mathfrak m$，满足
$t \not \in \mathfrak m^2$。置 $Y_0 = \Spec(A/tA)$ 以及
$X_0 = X \times_Y Y_0$。则 $X_0 \to Y_0$ 在闭点上方的纤维之外为
\'etale。由于 $d > 2$，有
$\dim(\mathcal{O}_{X_0, x}) = d - 1$ 且该数 $\geq 2$。正规化
$X_0' \to X_0$ 是满射且有限（因为我们在完备局部环上工作，而这种环是
Nagata 环）。设 $x' \in X_0'$ 是映到 $x$ 的一点。由归纳假设，态射
$X'_0 \to Y$ 在 $x'$ 处为 \'etale。由包含关系
$\kappa(y) \subset \kappa(x) \subset \kappa(x')$
可知 $\kappa(x)$ 是 $\kappa(y)$ 的有限扩张。因此 $x$ 是
$X \to Y$ 在 $y$ 上方纤维的闭点。但 $x$ 也是该纤维的一般点，所以
$f$ 在 $x$ 处拟有限，从而归约到分支轨迹纯性的情形，见引理
\ref{pione-lemma-purity}。
\end{proof}




\section{分歧轨迹补集的仿射性}
\label{pione-section-stronger-purity}

\noindent
设 $f : X \to Y$ 是 Noether 概形之间的有限型态射，其中 $X$ 正规而
$Y$ 正则。设 $V \subset X$ 是 $f$ 为 \'etale 的最大开子概形。
\cite[Chapter IV, Section 21.12]{EGA} 中的讨论表明，
$V \to X$ 可能是仿射态射。注意，若 $V \to X$ 仿射，则可推出分歧轨迹的纯性
（引理 \ref{pione-lemma-purity-ramification}）；这里使用了《除子》引理
\ref{divisors-lemma-complement-affine-open-immersion}。因此，$V \to X$ 的仿射性是
分歧轨迹纯性的一种“强”形式。本节证明 $V \to X$ 仿射，条件是
$X$ 与 $Y$ 等特征且优良，见定理 \ref{pione-theorem-global}。有理由猜测，当 $X$ 与 $Y$
具有混合特征（但仍优良）时，该结果仍然成立。

\begin{lemma}
\label{pione-lemma-structure-cohomology}
设 $(A, \mathfrak m)$ 是包含一个域的正则局部环。设
$f : V \to \Spec(A)$ 为 \'etale 且准紧。假设
$\mathfrak m \not \in f(V)$，并且
$g : V \to \Spec(A) \setminus \{\mathfrak m\}$ 仿射。则对
$H^i(V, \mathcal{O}_V)$（$i > 0$）同构于 $A$ 的剩余域的单射包的若干副本的直和。
\end{lemma}

\begin{proof}
记 $U = \Spec(A) \setminus \{\mathfrak m\}$ 为穿孔谱。于是
$g : V \to U$ 仿射。由《概形的上同调》引理
\ref{coherent-lemma-relative-affine-cohomology}，有
$H^i(V, \mathcal{O}_V) = H^i(U, g_*\mathcal{O}_V)$。由《概形》引理
\ref{schemes-lemma-push-forward-quasi-coherent}，
$\mathcal{O}_U$-模 $g_*\mathcal{O}_V$ 是拟凝聚的。对任意拟凝聚
$\mathcal{O}_U$-模 $\mathcal{F}$，上同调 $H^i(U, \mathcal{F})$（$i > 0$）是
$\mathfrak m$-幂挠的；例如见《局部上同调》引理
\ref{local-cohomology-lemma-local-cohomology}。特别地，$A$-模
$H^i(V, \mathcal{O}_V)$（$i > 0$）是 $\mathfrak m$-幂挠的。

对任意平坦环映射 $A \to A'$，平坦基变换给出
$H^i(V, \mathcal{O}_V) \otimes_A A' = H^i(V', \mathcal{O}_{V'})$，
其中 $V' = V \times_{\Spec(A)} \Spec(A')$；见《概形的上同调》引理
\ref{coherent-lemma-flat-base-change-cohomology}。取 $A'$ 为 $A$ 的完备化
（由《代数补篇》节 \ref{more-algebra-section-permanence-completion}，它是平坦的），
则
$$
H^i(V, \mathcal{O}_V) = H^i(V, \mathcal{O}_V) \otimes_A A' =
H^i(V', \mathcal{O}_{V'}),\quad\text{对 } i > 0
$$
第一个等号由刚证明的挠性和《代数补篇》引理
\ref{more-algebra-lemma-neighbourhood-equivalence} 得出。此外，剩余域
$k$ 的单射包对 $A$ 与 $A'$ 相同，见《对偶化复形》引理
\ref{dualizing-lemma-compare}。这样便归约到
$A = k[[x_1, \ldots, x_d]]$ 的情形；见《代数》节
\ref{algebra-section-cohen-structure-theorem}。

\medskip\noindent
假设 $k$ 的特征为 $p > 0$。映射 $F : A \to A$，
$a \mapsto a^p$ 是平坦的（《局部上同调》引理
\ref{local-cohomology-lemma-frobenius-flat-regular}）；而由《\'Etale 态射》引理
\ref{etale-lemma-relative-frobenius-etale}，有
$V \times_{\Spec(A), \Spec(F)} \Spec(A) \cong V$，这是 $\Spec(A)$ 上概形的同构。
因此上面的结论给出
$H^i(V, \mathcal{O}_V) \otimes_{A, F} A \cong H^i(V, \mathcal{O}_V)$。
于是所求结论由《局部上同调》引理
\ref{local-cohomology-lemma-structure-torsion-Frobenius-regular} 得出。

\medskip\noindent
假设 $k$ 的特征为 $0$。由《局部上同调》引理
\ref{local-cohomology-lemma-etale-derivation}，存在加性算子
$D_j$，$j = 1, \ldots, d$，作用在 $H^i(V, \mathcal{O}_V)$ 上，
它们满足关于 $\partial_j = \partial/\partial x_j$ 的 Leibniz 法则。
于是所求结论由《局部上同调》引理
\ref{local-cohomology-lemma-structure-torsion-D-module-regular} 得出。
\end{proof}

\begin{lemma}
\label{pione-lemma-conclude}
在引理 \ref{pione-lemma-structure-cohomology} 的情形下，假设
$H^i(V, \mathcal{O}_V) = 0$ 对 $i \geq \dim(A) - 1$ 成立。则 $V$ 仿射。
\end{lemma}

\begin{proof}
令 $k = A/\mathfrak m$。由于
$V \times_{\Spec(A)} \Spec(k) = \emptyset$，由上同调与基变换可得
$$
R\Gamma(V, \mathcal{O}_V) \otimes_A^\mathbf{L} k = 0
$$
见《概形的导出范畴》引理 \ref{perfect-lemma-compare-base-change}。因此存在谱序列
（《代数补篇》例 \ref{more-algebra-example-tor}）
$$
E_2^{p, q} = \text{Tor}_{-p}(k, H^q(V, \mathcal{O}_V)),\quad
d_2^{p, q} : E_2^{p, q} \to E_2^{p + 2, q - 1}
$$
以及 $d_r^{p, q} : E_r^{p, q} \to E_r^{p + r, q - r + 1}$，
并且它收敛到零。由引理 \ref{pione-lemma-structure-cohomology}、《对偶化复形》引理
\ref{dualizing-lemma-tor-injective-hull} 以及假设
$H^i(V, \mathcal{O}_V) = 0$（$i \geq \dim(A) - 1$），可知没有非零微分进入或离开
$(p, q) = (0, 0)$ 项。因此
$H^0(V, \mathcal{O}_V) \otimes_A k = 0$。这意味着，若
$\mathfrak m = (x_1, \ldots, x_d)$，则有开覆盖
$V = \bigcup V \times_{\Spec(A)} \Spec(A_{x_i})$，它由仿射开子概形
$V \times_{\Spec(A)} \Spec(A_{x_i})$ 组成（因为 $V$ 在 $A$ 的穿孔谱上仿射），
并且 $x_1, \ldots, x_d$ 在 $\Gamma(V, \mathcal{O}_V)$ 中生成单位理想。
由《性质》引理 \ref{properties-lemma-characterize-affine}，这蕴含 $V$ 仿射。
\end{proof}

\begin{theorem}
\label{pione-theorem-global}
设 $Y$ 是域上的优良正则概形。设 $f : X \to Y$ 是概形的有限型态射，
且 $X$ 正规。设 $V \subset X$ 是 $f$ 为 \'etale 的最大开子概形。则含入态射
$V \to X$ 仿射。
\end{theorem}

\begin{proof}
设 $x \in X$，其像为 $y \in Y$。只须证明 $V \cap W$ 仿射，其中 $W$ 是
$x$ 的某个仿射开邻域。由于 $\Spec(\mathcal{O}_{X, x})$ 是这些概形 $W$ 的极限，
这当且仅当
$$
V_x = V \times_X \Spec(\mathcal{O}_{X, x})
$$
仿射（《极限》引理 \ref{limits-lemma-limit-affine}）。因此，只要该定理对态射
$X \times_Y \Spec(\mathcal{O}_{Y, y}) \to \Spec(\mathcal{O}_{Y, y})$ 成立，
原定理便成立。特别地，可以设 $Y$ 正则且维数有限，从而可对维数
$d = \dim(Y)$ 作归纳。再次结合相同论证，可以设 $Y$ 是以 $y$ 为闭点的局部概形，
并且 $V \cap (X \setminus f^{-1}(\{y\}) \to X \setminus f^{-1}(\{y\})$
仿射。

\medskip\noindent
设 $x \in X$ 是位于 $y$ 上方的一点。若 $x \in V$，则无须证明。注意，
$f^{-1}(\{y\}) \cap V$ 是由有限个闭点组成的集合（\'etale 态射的纤维是离散的）。
因此把 $X$ 换成 $x$ 的一个仿射开邻域后，可设
$y \not \in f(V)$。我们要证明 $V$ 仿射。

\medskip\noindent
令 $e(V)$ 为最大整数 $i$，使 $H^i(V, \mathcal{O}_V) \not = 0$。
由于 $X$ 仿射，整数 $e(V)$ 等于诸数 $e(V_x)$ 的最大值，其中
$x \in X \setminus V$；见《局部上同调》引理
\ref{local-cohomology-lemma-cd-local} 以及《局部上同调》引理
\ref{local-cohomology-lemma-cd} 中对上同调维数的刻画。由《局部上同调》引理
\ref{local-cohomology-lemma-cd-dimension}，有
$e(V_x) \leq \dim(\mathcal{O}_{X, x}) - 1$。若
$\dim(\mathcal{O}_{X, x}) \geq 2$，则分歧轨迹的纯性（引理
\ref{pione-lemma-purity-ramification}）表明 $V_x$ 是
$\mathcal{O}_{X, x}$ 的穿孔谱的真开子概形。由于 $\mathcal{O}_{X, x}$ 正规且优良，
Hartshorne--Lichtenbaum 消失定理（《局部上同调》引理
\ref{local-cohomology-lemma-affine-complement}）蕴含
$e(V_x) \leq \dim(\mathcal{O}_{X, x}) - 2$。另一方面，由于
$X \to Y$ 是有限型的，并且 $V \subset X$ 稠密（必要时把 $X$ 换成
$V$ 的闭包），维数公式（《态射》引理
\ref{morphisms-lemma-dimension-formula}）表明
$\dim(\mathcal{O}_{X, x}) \leq d$。因而
$e(V) \leq \max(0, d -  2)$。于是由引理
\ref{pione-lemma-conclude} 可知 $V$ 仿射，只要 $d \geq 2$。若 $d = 1$ 或 $d = 0$，则
$\mathcal{O}_{Y, y}$ 的穿孔谱仿射，故 $V$ 仿射。
\end{proof}




\section{光滑固有情形下的特化映射}
\label{pione-section-specialization-smooth-proper}

\noindent
本节讨论如下结果。设 $f : X \to S$ 是概形的固有光滑态射。设
$s \leadsto s'$ 是 $S$ 中点的一个特化。则特化映射
$$
sp : \pi_1(X_{\overline{s}}) \longrightarrow \pi_1(X_{\overline{s}'})
$$
（见节 \ref{pione-section-specialization-map}）是满射，并且
\begin{enumerate}
\item 若 $\kappa(s')$ 的特征为零，则它是同构；或者
\item 若 $\kappa(s')$ 的特征为 $p > 0$，则它在最大与 $p$ 互素商上诱导同构。
\end{enumerate}

\begin{lemma}
\label{pione-lemma-specialization-map-surjective}
设 $f : X \to S$ 是具有几何连通纤维的平坦固有态射。设
$s' \leadsto s$ 是一个特化。若 $X_s$ 几何既约，则特化映射
$sp : \pi_1(X_{\overline{s}'}) \to \pi_1(X_{\overline{s}})$ 是满射。
\end{lemma}

\begin{proof}
由于 $X_s$ 几何既约，必要时缩小 $S$ 后，可以假设所有纤维都几何既约；见
《态射补篇》引理 \ref{more-morphisms-lemma-geometrically-reduced-open}。令
$\mathcal{O}_{S, s} \to A \to \kappa(\overline{s}')$ 如特化映射的构造中那样；
见节 \ref{pione-section-specialization-map}。因此只须证明
$$
\pi_1(X_{\overline{s}'}) \to \pi_1(X_A)
$$
是满射。这由命题 \ref{pione-proposition-first-homotopy-sequence} 和
$\pi_1(\Spec(A)) = \{1\}$ 得出。
\end{proof}

\begin{proposition}
\label{pione-proposition-specialization-map-isomorphism}
设 $f : X \to S$ 是具有几何连通纤维的光滑固有态射。设
$s' \leadsto s$ 是一个特化。若 $\kappa(s)$ 的特征为零，则特化映射
$$
sp : \pi_1(X_{\overline{s}'}) \to \pi_1(X_{\overline{s}})
$$
是同构。
\end{proposition}

\begin{proof}
由引理 \ref{pione-lemma-specialization-map-surjective}，该映射是满射。因此只须证明它是单射。

\medskip\noindent
可设 $S$ 仿射。于是 $S$ 是一族在 $\mathbf{Z}$ 上有限型的仿射概形的余滤过极限。
因此可以假设 $X \to S$ 是 $X_0 \to S_0$ 的基变换，其中 $S_0$ 是一个有限型
$\mathbf{Z}$-代数的谱，而 $X_0 \to S_0$ 光滑且固有。见《极限》引理
\ref{limits-lemma-descend-finite-presentation}、
\ref{limits-lemma-descend-smooth} 和
\ref{limits-lemma-eventually-proper}。由引理
\ref{pione-lemma-specialization-map-base-change}，归约到基为 Noether 的情形。

\medskip\noindent
应用引理 \ref{pione-lemma-specialization-map-discrete-valuation-ring}，归约到如下情形：
基 $S$ 是严格 Hensel 离散赋值环 $A$ 的谱，并且所考察的是 $A$ 上的特化映射。
令 $K$ 为 $A$ 的分式域。选取代数闭包 $\overline{K}$，它对应于几何一般点
$\overline{\eta}$（在 $\Spec(A)$ 上）。对有限可分的
$\overline{K}/L/K$，令 $B \subset L$ 为 $A$ 在 $L$ 中的整闭包。由《代数补篇》注
\ref{more-algebra-remark-finite-separable-extension}，这是离散赋值环。

\medskip\noindent
设 $X \to \Spec(A)$ 如上一段。为了证明特化映射的单射性，只须证明每个有限
\'etale 覆盖 $V$（在 $X_{\overline{\eta}}$ 上）都是某个有限 \'etale 覆盖
$Y \to X$ 的基变换。事实上，此时由引理
\ref{pione-lemma-functoriality-galois-injective}，
$\pi_1(X_{\overline{\eta}}) \to \pi_1(X) = \pi_1(X_s)$ 是单射。

\medskip\noindent
给定 $V$，先由引理 \ref{pione-lemma-perfection} 把 $V$ 下降为
$V' \to X_{K^{sep}}$，再由引理 \ref{pione-lemma-limit} 下降为
$V'' \to X_L$。令 $Z \to X_B$ 为 $X_B$ 在 $V''$ 中的正规化。注意，
$Z$ 正规，并且有 $Z_L = V''$（作为 $X_L$ 上的概形）。因此
$Z \to X_B$ 在一般纤维上有限 \'etale。问题在于尚不知道
$Z \to X_B$ 是否处处 \'etale。由于 $X \to \Spec(A)$ 具有几何连通的光滑纤维，
可见特殊纤维 $X_s$ 几何不可约。因此
$X_B \to \Spec(B)$ 的特殊纤维不可约；令 $\xi_B$ 为其一般点。令
$\xi_1, \ldots, \xi_r$ 为 $Z$ 中映到 $\xi_B$ 的点。现在第一个问题
（事实将表明也是唯一的问题）是离散赋值环的扩张
$$
\mathcal{O}_{X_B, \xi_B} \subset \mathcal{O}_{Z, \xi_i}
$$
可能分歧。令 $e_i$ 为该扩张的分歧指数。注意，由于
$\kappa(s)$ 的特征为零，该分歧是温和的！

\medskip\noindent
为消除分歧，选取更进一步的有限可分扩张 $K^{sep}/L'/L/K$，使分歧指数
$e$（对应诱导扩张 $B'/B$）可被 $e_i$ 整除。考虑正规化基变换
$Z'$，它由 $Z$ 沿 $\Spec(B') \to \Spec(B)$ 得到；见《态射补篇》节
\ref{more-morphisms-section-reduced-fibre-theorem} 中的讨论。令
$\xi_{i, j}$ 为 $Z'$ 中既映到 $\xi_{B'}$ 又映到 $\xi_i$（在 $Z$ 中）的点。则局部环
$$
\mathcal{O}_{Z', \xi_{i, j}}
$$
是 $\mathcal{O}_{Z, \xi_i}$ 在 $L' \otimes_L F_i$ 中的整闭包的局部化，其中
$F_i$ 是 $\mathcal{O}_{Z, \xi_i}$ 的分式域；略去细节。因此 Abhyankar 引理
（《代数补篇》引理 \ref{more-algebra-lemma-abhyankar}）表明
$$
\mathcal{O}_{X_{B'}, \xi_{B'}} \subset \mathcal{O}_{Z', \xi_{i, j}}
$$
不分歧。由此可知态射 $Z' \to X_{B'}$ 在所有余维数 $1$ 的点处为 \'etale。
于是由分支轨迹的纯性（引理 \ref{pione-lemma-purity}），可见
$Z' \to X_{B'}$ 是有限 \'etale 的！

\medskip\noindent
不过，$A \to B'$ 诱导的剩余域扩张是平凡的（因为 $A$ 的剩余域可分闭且特征为零，
所以代数闭）。两次应用引理
\ref{pione-lemma-finite-etale-on-proper-over-henselian}，可知 $Z'$ 是某个有限
\'etale 覆盖 $Y \to X$ 的基变换（第一次得到 $Y$（在 $A$ 上），第二次证明拉回到
$B$ 后与 $Z'$ 同构）。这就完成了证明。
\end{proof}

\noindent
设 $G$ 是 profinite 群。设 $p$ 是素数。{\it 最大与 $p$ 互素商}按定义为
$$
G' = \lim_{U \subset G\text{ open, normal, index prime to }p} G/U
$$
若 $X$ 是连通概形且给定了 $p$，则最大与 $p$ 互素商（即 $\pi_1(X)$ 的相应商）
记为 $\pi'_1(X)$。

\begin{theorem}
\label{pione-theorem-specialization-map-isomorphism-prime-to-p}
设 $f : X \to S$ 是具有几何连通纤维的光滑固有态射。设
$s' \leadsto s$ 是一个特化。若 $\kappa(s)$ 的特征为 $p$，则特化映射
$$
sp : \pi_1(X_{\overline{s}'}) \to \pi_1(X_{\overline{s}})
$$
是满射，并诱导同构
$$
\pi'_1(X_{\overline{s}'}) \cong \pi'_1(X_{\overline{s}})
$$
于最大与 p 互素商之间。
\end{theorem}

\begin{proof}
证明与命题 \ref{pione-proposition-specialization-map-isomorphism} 完全相同，但有以下区别：
\begin{enumerate}
\item 给定 $X/A$，不再证明函子
$\textit{F\'Et}_X \to \textit{F\'Et}_{X_{\overline{\eta}}}$
本质满。只证明 Galois 群的阶与 $p$ 互素的 Galois 对象位于本质像中。
这足以用完全相同的论证推出
$\pi'_1(X_{\overline{s}'}) \to \pi'_1(X_{\overline{s}})$ 是单射。
\item 扩张
$\mathcal{O}_{X_B, \xi_B} \subset \mathcal{O}_{Z, \xi_i}$
温和分歧，因为相应的分式域扩张是 Galois 扩张，且 Galois 群的阶与 $p$ 互素。
见《代数补篇》引理 \ref{more-algebra-lemma-galois-conclusion}。
\item 扩张 $\kappa_B/\kappa_A$ 未必仍然平凡，但它是纯不可分的。因此态射
$X_{\kappa_B} \to X_{\kappa_A}$ 是泛同胚，并由命题
\ref{pione-proposition-universal-homeomorphism} 在基本群上诱导同构。
\end{enumerate}
\end{proof}













\section{温和分歧}
\label{pione-section-tame}

\noindent
设 $X \to Y$ 是 $\mathbf{Z}$ 上有限型概形之间的有限 \'etale 态射。
关于 $f$ 在 $\infty$ 处温和分歧的含义，有多种定义。文献
\cite{Kerz-Schmidt} 讨论了这些概念在何种程度上相互等价。

\medskip\noindent
本节讨论另一个更初等的问题，它先于无穷远处温和性的概念。
请与文献 \cite{Grothendieck-Murre} 中略有不同的讨论比较。
假设给定
\begin{enumerate}
\item 局部 Noether 概形 $X$；
\item 稠密开子概形 $U \subset X$；
\item 有限 \'etale 态射 $f : Y \to U$，
\end{enumerate}
使得对每个素除子 $Z \subset X$，若 $Z \cap U = \emptyset$，则局部环
$\mathcal{O}_{X, \xi}$，即 $X$ 在一般点 $\xi$（$Z$ 的一般点）处的局部环，
都是离散赋值环。令 $K_\xi$ 为 $\mathcal{O}_{X, \xi}$ 的分式域，
便得到概形的 Cartesian 方形
$$
\xymatrix{
\Spec(K_\xi) \ar[r] \ar[d] & U \ar[d] \\
\Spec(\mathcal{O}_{X, \xi}) \ar[r] & X
}
$$
特别地，$Y \times_U \Spec(K_\xi)$ 是某个有限可分代数
$L_\xi/K_\xi$ 的谱。
此时称 {\it $Y$ 在 $X$ 上于余维数 $1$ 处不分歧}，分别地称
{\it $Y$ 在 $X$ 上于余维数 $1$ 处温和分歧}，若
$L_\xi/K_\xi$ 关于 $\mathcal{O}_{X, \xi}$ 对每个上述
$(Z, \xi)$ 不分歧，分别地温和分歧；见
《代数补篇》定义 \ref{more-algebra-definition-types-of-extensions}。
更确切地说，将 $L_\xi$ 分解为 $K_\xi$ 的有限可分域扩张之积，
并要求每个因子关于 $\mathcal{O}_{X, \xi}$ 不分歧，分别地温和分歧。

\begin{lemma}
\label{pione-lemma-pullback-tame-codim1}
设 $X' \to X$ 是局部 Noether 概形的态射，$U \subset X$ 是稠密开子概形。假设
\begin{enumerate}
\item $U' = f^{-1}(U)$ 是 $X'$ 中的稠密开子概形；
\item 对每个素除子 $Z \subset X$，若 $Z \cap U = \emptyset$，则局部环
$\mathcal{O}_{X, \xi}$，即 $X$ 在一般点 $\xi$（$Z$ 的一般点）处的局部环，是离散赋值环；
\item 对每个素除子 $Z' \subset X'$，若 $Z' \cap U' = \emptyset$，则局部环
$\mathcal{O}_{X', \xi'}$，即 $X'$ 在一般点 $\xi'$（$Z'$ 的一般点）处的局部环，是离散赋值环；
\item 若 $\xi' \in X'$ 如 (3) 所述，则 $\xi = f(\xi')$ 如 (2) 所述。
\end{enumerate}
若 $f : Y \to U$ 有限 \'etale，且 $Y$ 在 $X$ 上于余维数 $1$ 处不分歧，
分别地温和分歧，则 $Y' = Y \times_X X' \to U'$ 有限 \'etale，且
$Y'$ 在 $X'$ 上于余维数 $1$ 处不分歧，分别地温和分歧。
\end{lemma}

\begin{proof}
本引理中唯一值得说明的事实是《代数补篇》引理
\ref{more-algebra-lemma-tame-goes-up} 给出的交换代数结果。
\end{proof}

\noindent
利用刚引入的术语，可以把先前得到的纯性结果重新表述为如下简洁形式。

\begin{lemma}
\label{pione-lemma-purity-one-divisor-general}
设 $X$ 是局部 Noether 概形，$U \subset X$ 开且稠密，并设
$Y \to U$ 是有限 \'etale 态射。假设
\begin{enumerate}
\item $Y$ 在 $X$ 上于余维数 $1$ 处不分歧；
\item 局部环 $\mathcal{O}_{X, x}$ 对所有 $x \in X \setminus U$ 都正则。
\end{enumerate}
则存在有限 \'etale 态射 $Y' \to X$，其在 $U$ 上的限制为 $Y$。
\end{lemma}

\begin{proof}
令 $\xi \in X \setminus U$ 为 $X \setminus U$ 的某个余维数 $1$ 不可约分支的一般点。
则 $\mathcal{O}_{X, \xi}$ 是离散赋值环。如上文讨论，记
$Y \times_U \Spec(K_\xi) = \Spec(L_\xi)$.
以 $B_\xi$ 表示 $\mathcal{O}_{X, \xi}$ 在 $L_\xi$ 中的整闭包。
$Y$ 在 $X$ 上于余维数 $1$ 处不分歧这一假设意味着
$\mathcal{O}_{X, \xi} \to B_\xi$ 有限 \'etale。因此得到有限 \'etale 态射
$Y_\xi \to \Spec(\mathcal{O}_{X, \xi})$ 以及同构
$$
Y \times_U \Spec(K_\xi) \cong
Y_\xi \times_{\Spec(\mathcal{O}_{X, \xi})} \Spec(K_\xi)
$$
（在 $\Spec(K_\xi)$ 上）。由《极限》引理
\ref{limits-lemma-glueing-near-point}，可以找到开子概形
$U \subset U' \subset X$（它包含 $\xi$）以及有限表现态射 $Y' \to U'$，使其在 $U$ 上的限制恢复
$Y$，而在 $\Spec(\mathcal{O}_{X, \xi})$ 上的限制恢复 $Y_\xi$。
再由《极限》引理 \ref{limits-lemma-glueing-near-point-properties}，态射
$Y' \to U'$ 在必要时把 $U'$ 缩小为更小的开集后有限 \'etale。
对 $X \setminus U$ 的其他余维数 $1$ 分支的一般点重复此论证，可以假设已有
有限 \'etale 态射 $Y' \to U'$，它把 $Y \to U$ 延拓到某个开子概形
$U' \subset X$，该开子概形包含 $U$ 以及所有余维数 $1$ 的 $X \setminus U$ 中的点。
最后对 $Y' \to U'$ 应用引理 \ref{pione-lemma-extend-pure}。事实上，所有局部环
$\mathcal{O}_{X, x}$（其中 $x \in X \setminus U'$）都正则且满足
$\dim(\mathcal{O}_{X, x}) \geq 2$。故由引理 \ref{pione-lemma-local-purity}，
$\mathcal{O}_{X, x}$ 具有纯性。
\end{proof}

\begin{lemma}
\label{pione-lemma-purity-one-divisor}
设 $X$ 是局部 Noether 概形，$D \subset X$ 是有效 Cartier 除子，且
$D$ 为正则概形。设 $Y \to X \setminus D$ 是有限 \'etale 态射。
若 $Y$ 在 $X$ 上于余维数 $1$ 处不分歧，则存在有限 \'etale 态射
$Y' \to X$，其在 $X \setminus D$ 上的限制为 $Y$。
\end{lemma}

\begin{proof}
这是引理 \ref{pione-lemma-purity-one-divisor-general} 的特殊情形。首先，
$D$ 在 $X$ 中处处不稠密（见《除子》节
\ref{divisors-section-effective-Cartier-divisors} 中的讨论），故
$X \setminus D$ 在 $X$ 中稠密。其次，环 $\mathcal{O}_{X, x}$ 对所有
$x \in D$ 都是正则局部环；这是由《代数》引理
\ref{algebra-lemma-regular-mod-x} 以及 $\mathcal{O}_{D, x}$ 正则这一假设得出的。
\end{proof}

\begin{example}[标准温和分歧态射]
\label{pione-example-tamely-ramified}
设 $A$ 是 Noether 环，$f \in A$ 是非零因子，且 $A/fA$ 既约。
这意味着 $A_\mathfrak p$ 是以 $f$ 为一致化元的离散赋值环，其中
$\mathfrak p$ 是 $f$ 上方的任意极小素理想。设 $e \geq 1$ 是在 $A$ 中可逆的整数。令
$$
C = A[x]/(x^e - f)
$$
则 $\Spec(C) \to \Spec(A)$ 是有限局部自由态射，并在 $A_f$ 的谱上 \'etale。
有限 \'etale 态射
$$
\Spec(C_f) \longrightarrow \Spec(A_f)
$$
在 $\Spec(A)$ 上于余维数 $1$ 处温和分歧。温和性立即由《代数补篇》引理
\ref{more-algebra-lemma-characterize-tame} 对温和分歧扩张的刻画得出。
\end{example}

\noindent
下面给出适用于正则除子的 Abhyankar 引理版本。

\begin{lemma}[正则除子的 Abhyankar 引理]
\label{pione-lemma-abhyankar-one-divisor}
设 $X$ 是局部 Noether 概形，$D \subset X$ 是有效 Cartier 除子，且
$D$ 为正则概形。设 $Y \to X \setminus D$ 是有限 \'etale 态射。
若 $Y$ 在 $X$ 上于余维数 $1$ 处温和分歧，则在 $X$ 上 \'etale 局部，
态射 $Y \to X$ 是例 \ref{pione-example-tamely-ramified} 所述标准温和分歧态射的有限不交并。
\end{lemma}

\begin{proof}
对每个 $x \in X$，将找出一个 \'etale 邻域 $(U, u) \to (X, x)$，其中
$U = \Spec(A)$，使得基变换 $Y \times_X U \to U$ 是例
\ref{pione-example-tamely-ramified} 中标准温和分歧态射的有限不交并。
可设 $x \in D$；当 $x \not \in D$ 时，结论由《\'Etale 态射》引理
\ref{etale-lemma-finite-etale-etale-local} 以及在例
\ref{pione-example-tamely-ramified} 中取 $e = 1$、$f = 1$ 得出。

\medskip\noindent
本段把问题归约到 $X$ 为严格 Hensel 局部环之谱且 $x$ 为闭点的情形。
具体地，缩小 $X$ 后可以假设 $X = \Spec(A)$，而 $D \subset X$ 由非零因子
$f \in A$ 给出。此时 $Y$ 作为 $\Spec(A_f)$ 的有限 \'etale 覆盖是仿射的；记
$Y = \Spec(B)$。令 $A^{sh}$ 为 $\mathcal{O}_{X, x}$ 的严格 Hensel 化。
由于 $A \to A^{sh}$ 平坦且 $x \in D$，可见 $f$ 映为 $A^{sh}$ 的极大理想中的非零因子。
注意，$A^{sh}/fA^{sh} = (A/fA)^{sh}$（《代数》引理
\ref{algebra-lemma-quotient-strict-henselization}）是 $\mathcal{O}_{D, x}$
的严格 Hensel 化，故为正则局部环；见《代数补篇》引理
\ref{more-algebra-lemma-henselization-regular}。由引理
\ref{pione-lemma-pullback-tame-codim1}，$Y$ 到 $\Spec(A^{sh})$ 的基变换在余维数
$1$ 处温和分歧。假设已经对 $Y$ 到 $\Spec(A^{sh})$ 的基变换证明了断言。由于 $A^{sh}$ 严格 Hensel，
每个 \'etale 邻域都有截面，从而得到同构
$$
A^{sh} \otimes_A B \cong
\prod\nolimits_{i \in I} A^{sh}_f[x]/(x^{e_i} - f)
$$
其中 $I$ 是有限集，每个 $e_i$ 都是在 $A^{sh}$ 中可逆的整数。把
$A^{sh} = \colim A_\lambda$ 写成 \'etale $A$-代数 $A_\lambda$ 的滤过余极限。
上示同构下降为同构
$$
A_\lambda \otimes_A B \cong \prod\nolimits_{i \in I}
(A_\lambda)_f[x]/(x^{e_i} - f)
$$
（对充分大的 $\lambda$）；例如见《代数》引理
\ref{algebra-lemma-colimit-category-fp-algebras}。再稍微增大 $\lambda$，
可以假设 $e_i$ 在 $A_\lambda$ 中可逆。于是
$\Spec(A_\lambda) \to \Spec(A)$ 就是所求的 $x$ 的 \'etale 邻域。

\medskip\noindent
现在设 $X = \Spec(A)$，其中 $A$ 是严格 Hensel 局部环，而 $x \in X$
对应于极大理想 $\mathfrak m$，这是 $A$ 的极大理想。令 $f \in \mathfrak m$ 为切出 $D$ 的非零因子。
则 $A/fA$ 正则，由《代数》引理 \ref{algebra-lemma-regular-mod-x}，
$A$ 也正则。将使用正则局部环的一些性质，例如它们是正规整环；见《代数》引理
\ref{algebra-lemma-regular-domain} 和 \ref{algebra-lemma-regular-normal}。
特别地，$A$ 是整环。如上可见 $Y = \Spec(B)$ 仿射，且 $A_f \to B$ 有限
\'etale。故由《代数》引理 \ref{algebra-lemma-normal-goes-up}，$B$ 也正规；
再由《代数》引理 \ref{algebra-lemma-characterize-reduced-ring-normal}，
$Y$ 是若干正规整环之谱的有限不交并。因此也可假设 $B$ 是整环。
由于 $A$ 和 $A/fA$ 都是整环，$f$ 是 $A$ 的素元，并生成高度为 $1$ 的素理想
$\mathfrak p = fA$。注意，$A_\mathfrak p$ 是以 $f$ 为一致化元的离散赋值环。

\medskip\noindent
令 $K$ 为 $A$ 的分式域，$L$ 为 $B$ 的分式域。温和分歧假设意味着
$L$ 关于 $A_\mathfrak p$ 温和分歧。按《代数补篇》注
\ref{more-algebra-remark-finite-separable-extension}，令 $e \geq 1$ 为
$L$ 在 $A_\mathfrak p$ 上各分歧指数之积。于是 $e$ 在
$\kappa(\mathfrak p)$ 中可逆，但此时尚不知道 $e$ 在 $A/fA$ 或 $A$ 中是否可逆。

\medskip\noindent
考虑有限自由 $A$-代数
$$
A' = A[x]/(x^e - f)
$$
它是严格 Hensel 局部环的局部有限扩张，因而仍严格 Hensel。
注意，$f' = x$ 是 $A'$ 中的非零因子，而 $A'/f'A' \cong A/fA$ 是正则环。
因此如上，$A'$ 正则，进而是整环。令
$B' = B \otimes_A A' = B \otimes_{A_f} A'_{f'}$.
由 Abhyankar 引理（《代数补篇》引理 \ref{more-algebra-lemma-abhyankar}），
$\Spec(B')$ 在 $\Spec(A')$ 上于余维数 $1$ 处不分歧。事实上，由引理
\ref{pione-lemma-pullback-tame-codim1}，$\Spec(B')$ 在 $\Spec(A')$ 上于余维数
$1$ 处至少仍温和分歧，而 Abhyankar 引理说明所有分歧指数都已变为 $1$。
由引理 \ref{pione-lemma-purity-one-divisor}，$\Spec(B') \to \Spec(A'_{f'})$
延拓为有限 \'etale 态射 $\Spec(C) \to \Spec(A')$。然而 $A'$ 严格 Hensel，
所以 $\Spec(C)$ 是若干份 $\Spec(A')$ 的有限不交并。因而至少存在一个态射
$\Spec(A'_{f'}) \to \Spec(B')$，从而存在包含
$B \subset A'_{f'}$，它是 $A_f$-代数的包含。

\medskip\noindent
于是有包含 $K \subset L \subset K[f^{1/e}]$。由《代数补篇》引理
\ref{more-algebra-lemma-pull-root-uniformizer}，$L$ 等于
$K[f^{1/n}]$，其中 $n$ 是 $e$ 的某个因子。$B$ 的正规性随即推出
$B \cong A_f[y]/(y^n - f)$；略去细节。不过，环映射
$A_f \to A_f[y]/(y^n - f)$ 的分歧轨迹由 $ny^{n - 1}$ 切出，
这必然意味着 $n$ 在 $A_f$ 中为单位。由于 $f$ 是素元，这意味着
$n = u f^s$，其中 $u$ 是 $A$ 的某个单位，$s \in \mathbf{Z}$。
又因 $n$ 在 $\kappa(\mathfrak p)$ 中映为单位，故 $s = 0$，证明完成。
\end{proof}

\begin{lemma}
\label{pione-lemma-extend-tame-covering-normal}
在引理 \ref{pione-lemma-abhyankar-one-divisor} 的情形中，$X$ 在 $Y$ 中的正规化是
有限局部自由态射 $\pi : Y' \to X$，并且
\begin{enumerate}
\item $Y'$ 在 $X \setminus D$ 上的限制同构于 $Y$；
\item $D' = \pi^{-1}(D)_{red}$ 是 $Y'$ 上的有效 Cartier 除子；
\item $D'$ 是正则概形。
\end{enumerate}
此外，在 $X$ 上 \'etale 局部，态射 $Y' \to X$ 是下列态射的有限不交并：
$$
\Spec(A[x]/(x^e - f)) \to \Spec(A)
$$
其中 $A$ 是 Noether 环，$f \in A$ 是非零因子，$A/fA$ 正则，并且
$e \geq 1$ 在 $A$ 中可逆。
\end{lemma}

\begin{proof}
这只是引理 \ref{pione-lemma-abhyankar-one-divisor} 的补充；事实上，读过该引理的证明后，
本引理几乎立即成立。也可以直接从引理 \ref{pione-lemma-abhyankar-one-divisor} 的结论推出。
具体地，取 $X$ 在 $Y$ 中的正规化与 \'etale 基变换可交换；见《态射补篇》引理
\ref{more-morphisms-lemma-normalization-smooth-localization}。所以关于
$Y' \to X$ 局部结构的断言可以在 $X$ 上 \'etale 局部证明。由引理
\ref{pione-lemma-abhyankar-one-divisor}，可以假设 $X = \Spec(A)$，其中 $A$ 是
Noether 环；存在非零因子 $f\in A$ 使 $A/fA$ 正则；并且 $Y$ 是若干环
$A_f[x]/(x^e - f)$ 之谱的有限不交并，其中 $e$ 在 $A$ 中可逆。
略去如下验证：$A$ 在 $A_f[x]/(x^e - f)$ 中的整闭包等于
$A' = A[x]/(x^e - f)$。（为此，只须证明 $A'$ 在位于 $(f)$ 上方的素理想处的
局部化正则。）其余细节从略。
\end{proof}

\begin{lemma}
\label{pione-lemma-tame-covering-split}
在引理 \ref{pione-lemma-abhyankar-one-divisor} 的情形中，令 $Y' \to X$ 如引理
\ref{pione-lemma-extend-tame-covering-normal} 所述。设 $R$ 是分式域为 $K$ 的离散赋值环。令
$$
t : \Spec(R) \to X
$$
为一个态射，使得概形论逆像 $t^{-1}D$ 是 $\Spec(R)$ 的约化闭点。
\begin{enumerate}
\item 若 $t|_{\Spec(K)}$ 可提升为 $Y$ 的一点，则得到提升
$t' : \Spec(R) \to Y'$，并且 $Y' \to X$ 沿 $t'(\Spec(R))$ 为 \'etale；
\item 若 $\Spec(K) \times_X Y$ 同构于若干份 $\Spec(K)$ 的不交并，
则 $Y' \to X$ 在 $t(\Spec(R))$ 的某个开邻域上有限 \'etale。
\end{enumerate}
\end{lemma}

\begin{proof}
把固有性的赋值判据应用于有限态射 $Y' \to X$，可知
$\Spec(K)$-值点（在 $Y$ 中）若与 $t|_{\Spec(K)}$ 作为到 $X$ 的态射相符，便唯一提升为态射
$t' : \Spec(R) \to Y'$。因而陈述 (1) 有意义。

\medskip\noindent
选取 \'etale 邻域 $(U, u) \to (X, t(\mathfrak m_R))$，使
$U = \Spec(A)$，并使 $Y' \times_X U \to U$ 对某个 $f \in A$ 具有引理
\ref{pione-lemma-extend-tame-covering-normal} 中的描述。于是
$\Spec(R) \times_X U \to \Spec(R)$ \'etale 且满。若 $R'$ 表示
$\Spec(R) \times_X U$ 在 $\Spec(R)$ 闭点上方的局部环，则 $R'$ 是离散赋值环，
且 $R \subset R'$ 是离散赋值环的不分歧扩张（《代数补篇》引理
\ref{more-algebra-lemma-Dedekind-etale-extension}）。关于 $t$ 的假设意味着映射
$A \to R'$（它对应于
$$
\Spec(R') \to \Spec(R) \times_X U \to U
$$
）把 $f$ 映为一致化元 $\pi \in R'$。现设
$$
Y' \times_X U =
\coprod\nolimits_{i \in I} \Spec(A[x]/(x^{e_i} - f))
$$
其中 $e_i \geq 1$。于是
$$
\Spec(R') \times_U (Y' \times_X U) =
\coprod\nolimits_{i \in I} \Spec(R'[x]/(x^{e_i} - \pi))
$$
环 $R'[x]/(x^{e_i} - f)$ 是离散赋值环（《代数补篇》引理
\ref{more-algebra-lemma-pull-root-uniformizer}），所以它们不存在到 $R'$ 的分式域的映射，
除非 $e_i = 1$。

\medskip\noindent
(1) 的证明。在此情形中，映射 $t' : \Spec(R) \to Y'$ 经基变换给出相应映射
$t'' : \Spec(R') \to Y' \times_X U$。由上文讨论，它必映入某个对应于
$i \in I$ 且 $e_i = 1$ 的分支。因此显然 $Y' \times_X U \to U$ 沿
$t''$ 的像为 \'etale。\'etale 性可在 \'etale 基变换后检验，故 (1) 得证。

\medskip\noindent
(2) 的证明。在此情形中，假设蕴含 $e_i = 1$，对所有 $i \in I$ 均如此。
因此 $Y' \times_X U \to U$ 有限 \'etale，结论同上。
\end{proof}

\begin{lemma}
\label{pione-lemma-extend-covering}
设 $S$ 是一般点为 $\eta$ 的正规整 Noether 概形。设 $f : X \to S$ 是具有几何连通
纤维的光滑态射。设 $\sigma : S \to X$ 是 $f$ 的截面。设 $Z \to X_\eta$ 是有限
\'etale Galois 覆盖（节 \ref{pione-section-finite-etale-under-galois}），其群 $G$ 的阶在
$S$ 上可逆，并且 $Z$ 有一个 $\kappa(\eta)$-理性点映到 $\sigma(\eta)$。
则存在有限 \'etale Galois 覆盖 $Y \to X$，其群为 $G$，且其在
$X_\eta$ 上的限制为 $Z$。
\end{lemma}

\begin{proof}
先设 $S = \Spec(R)$ 是离散赋值环 $R$ 的谱，闭点为 $s \in S$。于是
$X_s$ 是 $X$ 中的有效 Cartier 除子，并且作为域上的光滑概形，$X_s$ 正则。
此外，一般纤维 $X_\eta$ 就是开子概形 $X \setminus X_s$。由《代数补篇》引理
\ref{more-algebra-lemma-galois-conclusion} 及关于 $G$ 的假设，$Z$ 在 $X$ 上于
余维数 $1$ 处温和分歧。令 $Z' \to X$ 如引理
\ref{pione-lemma-extend-tame-covering-normal} 所述。注意，$G$ 在 $Z$ 上的作用延拓为
$G$ 在 $Z'$ 上的作用。由引理 \ref{pione-lemma-tame-covering-split}，
$Z' \to X$ 在 $\sigma(y)$ 的某个开邻域上有限 \'etale。由于 $X_s$ 不可约，
这蕴含 $Z \to X_\eta$ 在 $X$ 上于余维数 $1$ 处不分歧。于是得到有限 \'etale 态射
$Y \to X$，其在 $X_\eta$ 上的限制为 $Z$；这是由引理
\ref{pione-lemma-purity-one-divisor} 得出的。当然 $Y \cong Z'$（略去细节；提示：在 \'etale
局部计算），故 $Y$ 是群为 $G$ 的 Galois 覆盖。

\medskip\noindent
一般情形。令 $U \subset S$ 为满足如下条件的极大开子概形：存在有限
\'etale Galois 覆盖 $Y \to X \times_S U$，其群为 $G$，且其在 $X_\eta$ 上的限制同构于 $Z$。
为导出矛盾，假设 $U \not = S$。令 $s \in S \setminus U$ 为
$S \setminus U$ 某个不可约分支的一般点。则逆像 $U_s$（即 $U$ 在
$\Spec(\mathcal{O}_{S, s})$ 中的逆像）是
$\mathcal{O}_{S, s}$ 的穿孔谱。我们断言
$Y \times_S U_s \to X \times_S U_s$ 是有限 \'etale Galois 覆盖
$Y'_s \to X \times_S \Spec(\mathcal{O}_{S, s})$ 的限制，该覆盖的群为 $G$。

\medskip\noindent
先说明该断言如何产生所需矛盾。由《极限》引理
\ref{limits-lemma-glueing-near-point}，可找到开子概形
$U \subset U' \subset S$（它包含 $s$）以及有限表现态射 $Y'' \to U'$，使其在 $U$ 上的限制
恢复 $Y' \to U$，而在 $\Spec(\mathcal{O}_{S, s})$ 上的限制恢复 $Y'_s$。
此外，由该引理给出的范畴等价，缩小 $U'$ 后可假设存在态射
$Y'' \to U' \times_S X$，并且 $G$ 在 $Y''$ 上、在 $U' \times_S X$ 上方作用；
这些态射与作用经基变换到 $U$ 和 $\Spec(\mathcal{O}_{S, s})$ 后与给定者相容。
必要时进一步缩小 $U'$，可假设 $Y'' \to U \times_S X$ 有限 \'etale；见
《极限》引理 \ref{limits-lemma-glueing-near-point-properties}。这说明找到了
$S$ 的一个严格更大的开集，使 $Y$ 在其上延拓为群为 $G$ 的有限 \'etale Galois
覆盖，从而得到所需矛盾。

\medskip\noindent
断言的证明。替换 $S$，所用者为 $\Spec(\mathcal{O}_{S, s})$。于是
$S = \Spec(A)$，其中 $(A, \mathfrak m)$ 是正规局部整环；$U \subset S$
为穿孔谱，并已有有限 \'etale Galois 覆盖
$Y \to X \times_S U$，其群为 $G$。若 $\dim(A) = 1$，则由证明第一段可把 $Y$ 延拓为
$X$ 的 Galois 覆盖。因此可设 $\dim(A) \geq 2$，从而有
$\text{depth}(A) \geq 2$，因为 $S$ 正规；见《代数》引理
\ref{algebra-lemma-criterion-normal}。由于 $X \to S$ 平坦，可知
$\text{depth}(\mathcal{O}_{X, x}) \geq 2$，对每个点 $x \in X$（它映到 $s$）均如此；
见《代数》引理 \ref{algebra-lemma-apply-grothendieck}。令
$$
Y' \longrightarrow X
$$
为用 $Y \to X \times_S U$ 在引理 \ref{pione-lemma-extend-S2} 中构造的有限态射。
注意，得到 $G$ 在 $Y$ 上的典范作用。因此只剩证明 $Y'$ 在 $X$ 上 \'etale。
事实上，例如由引理 \ref{pione-lemma-purity-smooth-over-depth2}，只须证明
$Y' \to X$ 在纤维 $X_s$ 的（唯一）一般点上方为 \'etale。为此，在
$\sigma(s)$ 的一个（形式）邻域中作局部计算。

\medskip\noindent
选取仿射开子概形 $\Spec(B) \subset X$，它包含 $\sigma(s)$。则
$A \to B$ 是具有截面 $\sigma : B \to A$ 的光滑环映射。记
$I = \Ker(\sigma)$，并以 $B^\wedge$ 表示 $I$-进的 $B$ 的完备化。由《代数》引理
\ref{algebra-lemma-section-smooth}，有
$B^\wedge \cong A[[x_1, \ldots, x_d]]$，其中 $d \geq 0$。当然 $B \to B^\wedge$ 平坦
（《代数》引理 \ref{algebra-lemma-completion-flat}），且
$\Spec(B^\wedge) \to X$ 的像包含 $X_s$ 的一般点。令
$V \subset \Spec(B^\wedge)$ 为 $U$ 的逆像。考虑有限 \'etale 态射
$$
W = Y \times_{(X \times_S U)} V \longrightarrow V
$$
引理 \ref{pione-lemma-extend-S2} 中 $Y'$ 的构造与平坦基变换相容，因此基变换
$Y' \times_X \Spec(B^\wedge) \to \Spec(B^\wedge)$ 是从 $W \to V$
在 $\Spec(B^\wedge)$ 上按引理 \ref{pione-lemma-extend-S2} 的步骤构造的。
令 $V_0 = V \cap V(x_1, \ldots, x_d) \subset V$，并令
$W_0 = W \times_V V_0$。这是正规整概形，它映入 $\sigma(S)$，所用态射为
$\Spec(B^\wedge) \to X$，且事实上与 $\sigma(U)$ 等同。
因此 $W_0 \to V_0 = U$ 完全分裂：其一般纤维具有此性质，因为假设
$Z \to X_\eta$ 有一个 $\kappa(\eta)$-理性点位于 $\sigma(\eta)$ 上方
（而 $G$ 的作用当然随即蕴含整个纤维 $Z_{\sigma(\eta)}$ 是若干份概形
$\eta = \Spec(\kappa(\eta))$ 的不交并）。最后，由引理
\ref{pione-lemma-fully-faithful-power-series-over-depth2} 有
$$
W_0 \times_U V \cong W
$$
这说明 $W$ 是若干份 $V$ 的不交并，因而
$Y' \times_X \Spec(B^\wedge)$ 是若干份 $\Spec(B^\wedge)$ 的不交并，证明完成。
\end{proof}

\begin{lemma}
\label{pione-lemma-extend-covering-general}
设 $S$ 是一般点为 $\eta$ 的拟紧拟分离正规整概形。设 $f : X \to S$ 是具有
几何连通纤维的拟紧拟分离光滑态射。设 $\sigma : S \to X$ 是 $f$ 的截面。
设 $Z \to X_\eta$ 是有限 \'etale Galois 覆盖（节
\ref{pione-section-finite-etale-under-galois}），其群 $G$ 的阶在 $S$ 上可逆，
并且 $Z$ 有一个 $\kappa(\eta)$-理性点映到 $\sigma(\eta)$。
则存在有限 \'etale Galois 覆盖 $Y \to X$，其群为 $G$，且其在 $X_\eta$ 上的限制为 $Z$。
\end{lemma}

\begin{proof}
若 $S$ 为 Noether 概形，这就是引理 \ref{pione-lemma-extend-covering} 的结论。
一般情形由标准极限论证推出。我们强烈建议读者略过本证明。

\medskip\noindent
可把 $S = \lim S_i$ 写成一族概形的有向极限，其中过渡态射仿射，而每个
$S_i$ 都在 $\mathbf{Z}$ 上有限型；见《极限》命题
\ref{limits-proposition-approximate}。对每个 $i$，令
$S \to S'_i \to S_i$ 为 $S_i$ 在 $S$ 中的正规化；见《态射》节
\ref{morphisms-section-normalization-X-in-Y}。结合《代数》命题
\ref{algebra-proposition-ubiquity-nagata} 以及《态射》引理
\ref{morphisms-lemma-nagata-normalization-finite} 和
\ref{morphisms-lemma-normal-normalization}，可知 $S'_i$ 在 $\mathbf{Z}$ 上有限型，
在 $S_i$ 上有限，并且 $S'_i$ 是正规整概形，而 $S \to S'_i$ 占优。
由《态射》引理 \ref{morphisms-lemma-functoriality-normalization}，得到过渡态射
$S'_{i'} \to S'_i$；它们与过渡态射 $S_{i'} \to S_i$ 及源为 $S$ 的态射相容。
断言 $S = \lim S'_i$。略去证明（提示：在 $S_i$ 的一个选定仿射开集上方，
对某个 $i$，考察仿射开集，
把问题转化为直接的代数问题）。因此可把
$S = \lim S_i$ 写成一族正规整概形 $S_i$ 的有向极限，其中过渡态射仿射，且
$S_i$ 在 $\mathbf{Z}$ 上有限型。

\medskip\noindent
对某个 $i$，可找到有限表现光滑态射 $X_i \to S_i$，其到 $S$ 的基变换为
$X \to S$；见《极限》引理 \ref{limits-lemma-descend-finite-presentation} 和
\ref{limits-lemma-descend-smooth}。增大 $i$ 后，可假设截面 $\sigma$ 提升为截面
$\sigma_i : S_i \to X_i$（由《极限》引理
\ref{limits-lemma-descend-finite-presentation} 中的范畴等价）。可把 $X_i$ 替换为
《态射补篇》节 \ref{more-morphisms-section-connected-components} 所研究的开子概形
$X_i^0$，因为 $X \to X_i$ 的像显然映入其中（其开放性见《态射补篇》引理
\ref{more-morphisms-lemma-connected-along-section-open}）。于是可假设
$X_i \to S_i$ 的纤维几何连通。再增大 $i$，可假设 $|G|$ 在 $S_i$ 上可逆。
令 $\eta_i \in S_i$ 为一般点。由于 $X_\eta$ 是概形 $X_{i, \eta_i}$ 的极限，
可用完全相同的论证，把 $Z \to X_\eta$ 下降为某个有限 \'etale Galois 覆盖
$Z_i \to X_{i, \eta_i}$，必要时增大 $i$；见引理 \ref{pione-lemma-limit}。
必要时再次增大 $i$，可假设 $Z_i$ 有一个 $\kappa(\eta_i)$-理性点，
它映到 $\sigma_i(\eta_i)$。此时应用 Noether 情形的引理，再拉回到 $X$，即得结论。
\end{proof}











\section{正特征中的技巧}
\label{pione-section-achinger}

\noindent
Piotr Achinger 的论文 \cite{Achinger} 证明了正特征中的仿射概形总是
$K(\pi, 1)$。本节说明 \cite{Achinger} 中较初等的部分。具体地，证明：
对正特征域 $k$，一个在 $\mathbf{A}^n_k$ 上 \'etale 的仿射概形实际上通过另一个态射
在 $\mathbf{A}^n_k$ 上有限 \'etale。还将证明，正特征中连通仿射概形的闭浸入在
\'etale 基本群上诱导单射。

\medskip\noindent
设 $k$ 是特征为 $p > 0$ 的域。设
$$
k[x_1, \ldots, x_n] \longrightarrow A
$$
是有限型 $k$-代数的满射，其源是关于 $x_1, \ldots, x_n$ 的多项式代数。
以 $I \subset k[x_1, \ldots, x_n]$ 表示核，于是
$A = k[x_1, \ldots, x_n]/I$。不假设 $A$ 非零（换言之，允许 $A$ 为零环且
$I = k[x_1, \ldots, x_n]$）。最后，假设给定有限 \'etale 环映射
$\pi : A \to B$。

\medskip\noindent
假设给定 $k, n, k[x_1, \ldots, x_n] \to A, I, \pi : A \to B$。
设 $C$ 是 $k$-代数。考虑交换图
$$
\xymatrix{
& B \\
C \ar[r] & C/\varphi(I)C \ar[u]^\tau \\
k[x_1, \ldots, x_n] \ar[u]^\varphi \ar[r] &
A \ar[u] \ar@/_3em/[uu]_\pi
}
$$
其中 $\varphi$ 是 \'etale $k$-代数映射，$\tau$ 是满的 $k$-代数映射。
给定 $C, \varphi, \tau$。对任意 $r \geq 0$ 以及元素
$y_1, \ldots, y_r \in C$，若它们生成代数 $C$（在 $\Im(\varphi)$ 上），令
$s = s(r, y_1, \ldots, y_r) \in \{0, \ldots, r\}$ 为满足下列条件的最大元素：
元素 $y_i$ 在 $\Im(\varphi)$ 上整，其中 $1 \leq i \leq s$。定义
$NF(C, \varphi, \tau)$ 为
$r - s = r - s(r, y_1, \ldots, y_r)$ 的最小值，其中遍历所有上述
$r$ 和 $y_1, \ldots, y_r$ 的选择。注意，
$NF(C, \varphi, \tau)$ 为 $0$ 当且仅当 $\varphi$ 有限。

\begin{lemma}
\label{pione-lemma-lower-invariant}
在上述情形中，若 $NF(C, \varphi, \tau) > 0$，则存在 \'etale
$k$-代数映射 $\varphi'$ 和满的 $k$-代数映射 $\tau'$，它们嵌入交换图
$$
\xymatrix{
& B \\
C \ar[r] & C/\varphi'(I)C \ar[u]_{\tau'} \\
k[x_1, \ldots, x_n] \ar[u]^{\varphi'} \ar[r] &
A \ar[u] \ar@/_3em/[uu]_\pi
}
$$
并满足 $NF(C, \varphi', \tau') < NF(C, \varphi, \tau)$。
\end{lemma}

\begin{proof}
选取 $r \geq 0$ 和 $y_1, \ldots, y_r \in C$，它们生成
$C$，作为 $\Im(\varphi)$ 上的代数；再选取 $0 \leq s \leq r$，使得
$y_1, \ldots, y_s$ 在 $\Im(\varphi)$ 上整，并且
$r - s = NF(C, \varphi, \tau) > 0$。由于 $B$ 在 $A$ 上有限，
$y_{s + 1}$ 在 $B$ 中的像满足一个以 $A$ 中元素为系数的首一多项式。
因此可以找到 $d \geq 1$ 和
$f_1, \ldots, f_d \in k[x_1, \ldots, x_n]$，使得
$$
z = y_{s + 1}^d + \varphi(f_1) y_{s + 1}^{d - 1} + \ldots + \varphi(f_d) \in
J = \Ker(C \to C/\varphi(I)C \xrightarrow{\tau} B)
$$
由于 $\varphi : k[x_1, \ldots, x_n] \to C$ 为 \'etale，可以找到非零非常数多项式
$g \in k[T_1, \ldots, T_{n + 1}]$，使得
$$
g(\varphi(x_1), \ldots, \varphi(x_n), z) = 0
\quad\text{于}\quad
C
$$
为说明这一点，例如可用如下事实：
$C \otimes_{\varphi, k[x_1, \ldots, x_n]} k(x_1, \ldots, x_n)$
是 $k(x_1, \ldots, x_n)$ 的有限可分域扩张的有限积
（见《代数》引理 \ref{algebra-lemma-etale-over-field}），因此 $z$ 满足一个
以 $k(x_1, \ldots, x_n)$ 中元素为系数的首一多项式。清除分母即得 $g$。

\medskip\noindent
$g$ 的存在性与《代数》引理 \ref{algebra-lemma-helper-polynomial} 给出整数
$e_1, e_2, \ldots, e_n \geq 1$，使得 $z$ 在子环 $C'$ 上整，而该子环属于 $C$ 并由
$t_1 = \varphi(x_1) + z^{pe_1}, \ldots, t_n = \varphi(x_n) + z^{pe_n}$ 生成。
当然，元素 $\varphi(x_1), \ldots, \varphi(x_n)$ 也在 $C'$ 上整，元素
$y_1, \ldots, y_s$ 亦然。最后，由 $z$ 的选取，元素
$y_{s + 1}$ 也在 $C'$ 上整。

\medskip\noindent
考虑环映射
$$
\varphi' : k[x_1, \ldots, x_n] \longrightarrow C, \quad
x_i \longmapsto t_i
$$
其像为 $C'$。由于
$\text{d}(\varphi(x_i)) = \text{d}(t_i) = \text{d}(\varphi'(x_i))$
在 $\Omega_{C/k}$ 中成立（这里用到 $k$ 的特征为 $p > 0$），所以
$\varphi'$ 为 \'etale，因为 $\varphi$ 为 \'etale；见《代数》引理
\ref{algebra-lemma-characterize-etale-over-polynomial-ring}。注意，
$\varphi'(x_i) - \varphi(x_i) = t_i - \varphi(x_i) = z^{pe_i}$ 属于核
$J$，它是映射 $C \to C/\varphi(I)C \to B$ 的核；这是因为选取的 $z$ 是 $J$ 的元素。
因此，对 $f \in I$，元素
$$
\varphi'(f) =
f(t_1, \ldots, t_n) =
f(\varphi(x_1) + z^{pe_1}, \ldots, \varphi(x_n) + z^{pe_n}) =
\varphi(f) + \text{element of }(z)
$$
也属于 $J$。换言之，$\varphi'(I)C \subset J$，并得到满射
$$
\tau' : C/\varphi'(I)C \longrightarrow C/J \cong B
$$
（在 $A$ 上 \'etale 的代数之间）。最后，代数 $C$ 由元素
$\varphi(x_1), \ldots, \varphi(x_n), y_1, \ldots, y_r$ 生成（在
$C' = \Im(\varphi')$ 上），而
$\varphi(x_1), \ldots, \varphi(x_n), y_1, \ldots, y_{s + 1}$
在 $C' = \Im(\varphi')$ 上整。因此
$NF(C, \varphi', \tau') < r - s = NF(C, \varphi, \tau)$。证明完成。
\end{proof}

\begin{lemma}
\label{pione-lemma-affine-etale-over-affine-space}
设 $k$ 是特征为 $p > 0$ 的域。设 $X \to \mathbf{A}^n_k$ 是 \'etale 态射，
且 $X$ 仿射。则存在有限 \'etale 态射 $X \to \mathbf{A}^n_k$。
\end{lemma}

\begin{proof}
记 $X = \Spec(C)$。令 $A = 0$，并记 $I = k[x_1, \ldots, x_n]$。
由假设，存在某个 \'etale $k$-代数映射
$\varphi : k[x_1, \ldots, x_n] \to C$。以 $\tau : C/\varphi(I)C \to 0$
表示唯一满射。可以选取 $\varphi$ 和 $\tau$，使 $N(C, \varphi, \tau)$ 最小。
由引理 \ref{pione-lemma-lower-invariant}，得到 $N(C, \varphi, \tau) = 0$。
因此 $\varphi$ 有限 \'etale。
\end{proof}

\begin{lemma}
\label{pione-lemma-dominate-affine-space}
设 $k$ 是特征为 $p > 0$ 的域。设 $Z \subset \mathbf{A}^n_k$ 是闭子概形。
设 $Y \to Z$ 有限 \'etale。则存在有限 \'etale 态射
$f : U \to \mathbf{A}^n_k$，使得存在开且闭的浸入
$Y \to f^{-1}(Z)$（在 $Z$ 上）。
\end{lemma}

\begin{proof}
把问题转化为代数。记
$\mathbf{A}^n_k = \Spec(k[x_1, \ldots, x_n])$.
则 $Z = \Spec(A)$，其中 $A = k[x_1, \ldots, x_n]/I$，这里
$I \subset k[x_1, \ldots, x_n]$ 是某个理想。记 $Y = \Spec(B)$，
于是 $Y \to Z$ 对应于有限 \'etale $k$-代数映射 $A \to B$。

\medskip\noindent
由《代数》引理 \ref{algebra-lemma-lift-etale}，存在 \'etale 环映射
$$
\varphi : k[x_1, \ldots, x_n] \to C
$$
以及满的 $A$-代数映射 $\tau : C/\varphi(I)C \to B$。
（甚至可以选取 $C, \varphi, \tau$ 使 $\tau$ 为同构，但这里不用这一点。）
可以选取 $\varphi$ 和 $\tau$，使 $N(C, \varphi, \tau)$ 最小。
由引理 \ref{pione-lemma-lower-invariant}，得到 $N(C, \varphi, \tau) = 0$。
因此 $\varphi$ 有限 \'etale。

\medskip\noindent
令 $f : U = \Spec(C) \to \mathbf{A}^n_k$ 为对应于 $\varphi$ 的有限 \'etale
态射。态射 $Y \to f^{-1}(Z) = \Spec(C/\varphi(I)C)$（由 $\tau$ 诱导）
是闭浸入，因为 $\tau$ 满；
它也是开的，因为由《态射》引理
\ref{morphisms-lemma-etale-permanence}，它是 \'etale 态射。证明完成。
\end{proof}

\noindent
下面是主要结果。

\begin{proposition}
\label{pione-proposition-injective}
设 $p$ 是素数。设 $i : Z \to X$ 是 $\mathbf{F}_p$ 上连通仿射概形的闭浸入。
对任意几何点 $\overline{z}$（取自 $Z$），映射
$$
\pi_1(Z, \overline{z}) \to \pi_1(X, \overline{z})
$$
是单射。
\end{proposition}

\begin{proof}
设 $Y \to Z$ 是有限 \'etale 态射。只须构造有限 \'etale 态射
$f : U \to X$，使 $Y$ 同构于 $f^{-1}(Z)$ 的一个开且闭子概形；见引理
\ref{pione-lemma-functoriality-galois-injective}。记 $Y = \Spec(A)$ 且
$X = \Spec(R)$，于是闭浸入 $Y \to X$ 由满射 $R \to A$ 给出。
可把 $A = \colim A_i$ 写成它的有限型 $\mathbf{F}_p$-子代数的滤过余极限。
由引理 \ref{pione-lemma-limit}，可以找到一个 $i$ 和有限 \'etale 态射
$Y_i \to Z_i = \Spec(A_i)$，使得 $Y = Z \times_{Z_i} Y_i$。

\medskip\noindent
选取满射 $\mathbf{F}_p[x_1, \ldots, x_n] \to A_i$。它确定闭浸入
$$
Z_i = \Spec(A_i)
\longrightarrow
X_i = \mathbf{A}^n_{\mathbf{F}_p} = \Spec(\mathbf{F}_p[x_1, \ldots, x_n])
$$
由多项式代数的泛性质以及 $R \to A$ 为满射这一事实，可以找到交换图
$$
\xymatrix{
\mathbf{F}_p[x_1, \ldots, x_n] \ar[r] \ar[d] &
A_i \ar[d] \\
R \ar[r] &
A
}
$$
（在 $\mathbf{F}_p$-代数范畴中）。因此有交换图
$$
\xymatrix{
Y_i \ar[r] &
Z_i \ar[r] &
X_i \\
Y \ar[u] \ar[r] &
Z \ar[u] \ar[r] &
X \ar[u]
}
$$
其右方形为 Cartesian。显然，若能找到有限 \'etale 的
$f_i : U_i \to X_i$，使 $Y_i$ 同构于 $f_i^{-1}(Z_i)$ 的一个开且闭子概形，
则基变换 $f : U \to X$（即 $f_i$ 沿 $X \to X_i$ 的基变换）就是本问题的解。
因此，应用引理 \ref{pione-lemma-dominate-affine-space} 于
$Y_i \to Z_i \to X_i = \mathbf{A}^n_{\mathbf{F}_p}$ 即得结论。
\end{proof}
% Generated mechanically from the frozen etale-cohomology.tex target.
% Do not edit this generated file; edit the bound locale source instead.

% OK, start here.
%

\setcounter{chapter}{58}
\chapter{\'Etale 上同调}



\phantomsection
\label{etale-cohomology-section-phantom}



\section{引言}
\label{etale-cohomology-section-introduction}

\noindent
本章是关于概形 \'etale 上同调系列章节的第一章。本章讨论 \'etale
拓扑以及该拓扑中阿贝尔层上同调的最基本内容。所讨论的许多主题在首次阅读时
可以放心略过；关于如何决定略过哪些内容，请参阅下一节的建议。

\medskip\noindent
本章的初版源自 Johan de Jong 于 2009 年秋季在哥伦比亚大学讲授的
一门 \'etale 上同调课程第一部分的讲义。最初记录讲义的是
Thibaut Pugin、Zachary Maddock 和 Min Lee。课程第二部分见迹公式一章，参见
The Trace Formula，第 \ref{trace-section-introduction} 节。




\section{首次阅读时可略过哪些小节？}
\label{etale-cohomology-section-skip}

\noindent
我们希望将本章材料用于发展与代数空间、Deligne--Mumford 叠、代数叠等
有关的理论。因此，我们在原本关于概形 \'etale 上同调的阐述中加入了一些
相当技术性的材料。读者可以通过 ``topos'' 一词出现的频繁程度、与集合论
有关的讨论，或处理概形态射非常一般性质的证明来辨认这些材料。其中一些
讨论在首次阅读时可以略过。

\medskip\noindent
特别地，我们建议读者略过下列小节：
\begin{enumerate}
\item 比较大拓扑斯与小拓扑斯，
第 \ref{etale-cohomology-section-compare} 节。
\item 恢复态射，
第 \ref{etale-cohomology-section-morphisms} 节。
\item 推与拉，
第 \ref{etale-cohomology-section-monomorphisms} 节。
\item 性质 (A)，
第 \ref{etale-cohomology-section-A} 节。
\item 性质 (B)，
第 \ref{etale-cohomology-section-B} 节。
\item 性质 (C)，
第 \ref{etale-cohomology-section-C} 节。
\item 小 \'etale 网站的拓扑不变性，
第 \ref{etale-cohomology-section-topological-invariance} 节。
\item 整的处处单射态射，
第 \ref{etale-cohomology-section-integral-universally-injective} 节。
\item 大网站与前推，
第 \ref{etale-cohomology-section-big} 节。
\item 大下撇函子的正合性，
第 \ref{etale-cohomology-section-exactness-lower-shriek} 节。
\end{enumerate}
除这些小节外，还有一些零散结果也可以略过；读者可根据上述关键词辨认它们。



%9.08.09
\section{序言}
\label{etale-cohomology-section-prologue}

\noindent
本讲义讨论另一种上同调理论。首先要指出，Zariski 拓扑并不完全令人满意。
它无法给出我们所希望结果的主要原因之一是：若 $X$ 是复簇且
$\mathcal{F}$ 是常值层，则
$$
H^i(X, \mathcal{F}) = 0, \quad \text{ 对所有 } i > 0.
$$
原因如下。在不可约概形（特别是簇）中，任意两个非空开子集都有交，因此
常值层的限制映射是满射的。我们称这样的层为 {\it flasque}（柔软层）。
此时所有高阶 {\v C}ech 上同调群都消失，所有高阶 Zariski 上同调群也都消失。
换言之，Zariski 拓扑中的开集 ``不够多''，无法检测这种高阶上同调。

\medskip\noindent
另一方面，若 $X$ 是光滑射影复簇，则
$$
H_{Betti}^{2 \dim X}(X (\mathbf{C}), \Lambda) = \Lambda \quad \text{ 对于 }
\Lambda = \mathbf{Z}, \ \mathbf{Z}/n\mathbf{Z},
$$
其中 $X(\mathbf{C})$ 表示 $X$ 的复点集。这是代数几何中值得复现的性质，
尤其是在正特征情形下。




\section{\'Etale 拓扑}
\label{etale-cohomology-section-etale-topology}

\noindent
要仅仅通过 ``增加'' 额外开集来细化 Zariski 拓扑是很困难的。定义拓扑的一种
有效方法是不只考虑开集，还考虑位于开集之上的某些概形。为了定义 \'etale
拓扑，我们考虑所有 \'etale 的态射 $\varphi : U \to X$。若 $X$ 是
$\mathbf{C}$ 上的光滑射影簇，这意味着
\begin{enumerate}
\item $U$ 是光滑簇的不交并，且
\item $\varphi$（从解析意义上说）局部为同构。
\end{enumerate}
这里 ``解析地'' 指的是 $\mathbf{C}$ 上通常的（超越的）拓扑。因此第二个
条件意味着 $\varphi$ 的导数处处满秩（特别地，$U$ 的每个连通分支都与
$X$ 具有相同维数）。

\medskip\noindent
例如，一个二重覆盖——粗略地说，即簇之间次数为 $2$ 的有限映射——
$$
\Spec(\mathbf{C}[t])
\longrightarrow
\Spec(\mathbf{C}[t]),
\quad t \longmapsto t^2
$$
若其纤维含有单个点，就不是 \'etale 态射。在该例中，这发生在 $t = 0$ 时。
要使 $\mathbf{C}$ 上的簇之间有限映射为 \'etale，所有纤维都应具有相同的点数。
从上述映射的源中去掉点 $t = 0$ 会使态射变为 \'etale。但也可以从源中去掉
其他点，所得态射仍然是 \'etale。为了考虑 \'etale 拓扑，我们必须考察所有
这类态射。与 Zariski 拓扑不同，它们不必只是 $X$ 的开子集，尽管它们的像
总是开子集。

\begin{definition}
\label{etale-cohomology-definition-etale-covering-initial}
若态射族 $\{ \varphi_i : U_i \to X\}_{i \in I}$ 中每个 $\varphi_i$ 都是
\'etale 态射，且它们的像覆盖 $X$，即
$X = \bigcup_{i \in I} \varphi_i(U_i)$.
\end{definition}

\noindent
这就 ``定义'' 了 \'etale 拓扑。换言之，现在可以说明层是什么。一个
{\it \'etale 层} $\mathcal{F}$（集合层，分别地，阿贝尔群层、向量空间层等）
在 $X$ 上是如下数据：
\begin{enumerate}
\item 对每个 \'etale 态射 $\varphi : U \to X$，给出一个集合
（分别地，阿贝尔群、向量空间等）$\mathcal{F}(U)$；
\item 对每一对 \'etale 概形 $U, \ U'$（它们位于 $X$ 上），以及每个态射
$U \to U'$（它位于 $X$ 上且自动是 \'etale 的），给出限制映射
$\rho^{U'}_U : \mathcal{F}(U') \to \mathcal{F}(U)$
\end{enumerate}
这些数据必须满足：对于恒等态射，有 $\rho^U_U = \text{id}$，其中恒等态射为
$U \to U$；有 $\rho^{U'}_U \circ \rho^{U''}_{U'} = \rho^{U''}_U$，当有态射
$U \to U' \to U''$（它们是位于 $X$ 上的 \'etale 概形之间的态射）时；此外还满足如下
{\it 层公理}：
\begin{itemize}
\item[(*)] 对每个 \'etale 覆盖 $\{ \varphi_i : U_i \to U\}_{i \in
I}$，图
$$
\xymatrix{
\mathcal{F} (U) \ar[r] &
\Pi_{i \in I} \mathcal{F} (U_i) \ar@<1ex>[r] \ar@<-1ex>[r] &
\Pi_{i, j \in I} \mathcal{F} (U_i \times_U U_j)
}
$$
在集合（分别地，阿贝尔群、向量空间等）范畴中是等化子图。
\end{itemize}

\begin{remark}
\label{etale-cohomology-remark-i-is-j}
在最后一个陈述中，不能遗漏 $i = j$ 的情形；这通常是一个高度非平凡的
条件（不同于 Zariski 拓扑）。事实上，许多重要覆盖只有一个元素。
\end{remark}

\noindent
由于恒等态射是 \'etale 态射，我们可以计算 \'etale 层的整体截面，而上同调
就只是相应的右导出函子。换言之，待进一步发展理论并使陈述精确化之后，
定义上同调将不再有障碍。




\section{\'Etale 拓扑的性质}
\label{etale-cohomology-section-feats}

\noindent
对自然数 $n \in \mathbf{N} = \{1, 2, 3, 4, \ldots\}$，有
$$
H_\etale^2 (\mathbf{P}^1_\mathbf{C}, \mathbf{Z}/n\mathbf{Z}) =
\mathbf{Z}/n\mathbf{Z}.
$$
更一般地，若 $X$ 是复簇，则其取有限域系数的 \'etale Betti 数与
$X(\mathbf{C})$ 的通常 Betti 数相同，即
$$
\dim_{\mathbf{F}_q} H_\etale^{2i} (X, \mathbf{F}_q) =
\dim_{\mathbf{F}_q} H_{Betti}^{2i} (X(\mathbf{C}), \mathbf{F}_q).
$$
这令人非常满意。然而，这些等式只对挠系数成立，并非一般成立。对整数系数，有
$$
H_\etale^2 (\mathbf{P}^1_\mathbf{C}, \mathbf{Z}) = 0.
$$
相反，$H_{Betti}^2(\mathbf{P}^1(\mathbf{C}), \mathbf{Z}) = \mathbf{Z}$，
因为拓扑空间 $\mathbf{P}^1(\mathbf{C})$ 同胚于 $2$-球面。
通过极限过程，可以从挠系数回到非挠系数；我们很快就会讨论这种方法。




\section{一个计算}
\label{etale-cohomology-section-computation}

\noindent
如何计算 $\mathbf{P}^1_\mathbf{C}$ 取系数
$\Lambda = \mathbf{Z}/n\mathbf{Z}$ 时的上同调？
我们使用 {\v C}ech 上同调。$\mathbf{P}^1_\mathbf{C}$ 的一个覆盖由两个标准开集
$U_0, U_1$ 给出；二者都同构于 $\mathbf{A}^1_\mathbf{C}$，而它们的交同构于
$\mathbf{A}^1_\mathbf{C} \setminus \{0\} = \mathbf{G}_{m, \mathbf{C}}$。
事实证明，Mayer--Vietoris 序列在 \'etale 上同调中成立，从而得到正合序列
$$
H_\etale^{i-1}(U_0\cap U_1, \Lambda) \to
H_\etale^i(\mathbf{P}^1_\mathbf{C}, \Lambda) \to
H_\etale^i(U_0, \Lambda) \oplus
H_\etale^i(U_1, \Lambda) \to H_\etale^i(U_0\cap U_1,
\Lambda).
$$
为了得到预期答案，我们需要证明第三项中的直和消失。事实上，与通常拓扑一样，有
$$
H_\etale^q (\mathbf{A}^1_\mathbf{C}, \Lambda) = 0
\quad \text{ 当 } q \geq 1,
$$
并且
$$
H_\etale^q (\mathbf{A}^1_\mathbf{C} \setminus \{0\}, \Lambda) = \left\{
\begin{matrix}
\Lambda & \text{ 当 }q = 1\text{ 时；} \\
0 & \text{ 当 }q \geq 2\text{ 时。}
\end{matrix}
\right.
$$
这些结果已经相当困难（是否存在初等证明？）。下面说明一旦掌握 \'etale
上同调的工具，我们将如何进行这个计算。

\medskip\noindent
{\bf 高阶上同调。}下面的一般事实处理了这一点：若 $X$ 是 $\mathbf{C}$ 上的
仿射曲线，则
$$
H_\etale^q (X, \mathbf{Z}/n\mathbf{Z}) = 0 \quad \text{ 当 } q \geq 2.
$$
这可通过考察曲线的泛点并进行一些 Galois 上同调来证明。因此只需考虑一阶上同调。

\medskip\noindent
{\bf 一阶上同调。}我们使用如下识别：
\begin{eqnarray*}
H_\etale^1 (X, \mathbf{Z}/n\mathbf{Z}) = \left\{
\begin{matrix}
\text{集合层 }\mathcal{F}\text{，它位于 \'etale 网站 }X_\etale
\text{ 上并带有} \\
\text{作用 }\mathbf{Z}/n\mathbf{Z} \times \mathcal{F} \to \mathcal{F}
\text{，使得 }\mathcal{F}\text{ 是一个 }\mathbf{Z}/n\mathbf{Z}\text{-挠子。}
\end{matrix}
\right\}
\Big/ \cong
\\
 = \left\{
\begin{matrix}
\text{态射 }Y \to X\text{，它们是有限 \'etale 的，并带有} \\
\text{自由的 }\mathbf{Z}/n\mathbf{Z}\text{ 作用，使得 }
X = Y/(\mathbf{Z}/n\mathbf{Z}).
\end{matrix}
\right\}
\Big/ \cong.
\end{eqnarray*}
第一个识别非常一般（对网站上的任意上同调理论都成立），与 \'etale 拓扑无关。
第二个识别是下降理论的结果。最后一个集合描述了一族我们可以实际操作的几何对象。

\medskip\noindent
曲线 $\mathbf{A}^1_\mathbf{C}$ 没有非平凡有限 \'etale 覆盖，因此
$H_\etale^1 (\mathbf{A}^1_\mathbf{C}, \mathbf{Z}/n\mathbf{Z}) = 0$.
这既可从拓扑上看出，也可使用下一段中的论证。

\medskip\noindent
下面描述有限 \'etale 覆盖
$\varphi : Y \to \mathbf{A}^1_\mathbf{C} \setminus \{0\}$。
只需考虑 $Y$ 连通的情形，以下即作此假设。我们将应用 Riemann--Hurwitz
公式来确定 $Y$ 的可能形式（当然这有些迂回；如果愿意，也可以直接略过下一节）。
设该态射是 $n$ 对 1，并考虑射影紧化
$$
\xymatrix{
{Y\ } \ar@{^{(}->}[r] \ar[d]^\varphi &
{\bar Y} \ar[d]^{\bar\varphi} \\
{\mathbf{A}^1_\mathbf{C} \setminus \{0\}} \ar@{^{(}->}[r] &
{\mathbf{P}^1_\mathbf{C}}
}
$$
尽管 $\varphi$ 是 \'etale 且不分歧，$\bar{\varphi}$ 可能在 0 和 $\infty$ 处分歧。
设 0 的原像是点 $y_1,
\ldots, y_r$，分歧指数为 $e_1, \ldots e_r$；
$\infty$ 的原像是点 $y_1', \ldots, y_s'$，分歧指数为 $d_1, \ldots d_s$。
特别地，$\sum e_i = n = \sum d_j$。应用 Riemann--Hurwitz 公式得到
$$
2 g_Y - 2 = -2n + \sum (e_i - 1) + \sum (d_j - 1)
$$
因此 $g_Y = 0$、$r = s = 1$ 且 $e_1 = d_1 = n$。
于是 $Y \cong {\mathbf{A}^1_\mathbf{C} \setminus \{0\}}$，容易看出
$\varphi(z) = \lambda z^n$，其中 $\lambda \in \mathbf{C}^*$。
重新参数化 $Y$ 后可假设 $\lambda = 1$。因此该覆盖就是取 $n$ 次根，
作用在 $\mathbf{A}^1_{\mathbf{C}} \setminus \{0\}$ 上的坐标。

\medskip\noindent
请记住，我们需要分类
${\mathbf{A}^1_\mathbf{C} \setminus \{0\}}$ 的覆盖，并同时考虑其上的自由
$\mathbf{Z}/n\mathbf{Z}$-作用。在本例中，任意这样的作用都对应于 $Y$ 的一个
自同构，将 $z$ 送到 $\zeta_n z$，其中 $\zeta_n$ 是本原 $n$ 次单位根。
这样的作用有 $\phi(n)$ 个（这里 $\phi(n)$ 表示 Euler 函数）。因此，连通有限 \'etale
覆盖恰有 $\phi(n)$ 个；给定自由 $\mathbf{Z}/n\mathbf{Z}$-作用时，每个对应一个本原
$n$ 次单位根。留给读者验证：非连通有限 \'etale 次数为 $n$ 的覆盖
$\mathbf{A}^1_{\mathbf{C}} \setminus \{0\}$，带有给定自由
$\mathbf{Z}/n\mathbf{Z}$-作用，与非本原的 $n$ 次单位根（它们是 $1$ 的根）一一对应。
换言之，这个计算表明
$$
H_\etale^1 (\mathbf{A}^1_\mathbf{C} \setminus \{0\},
\mathbf{Z}/n\mathbf{Z}) =
\Hom(\mu_n(\mathbf{C}), \mathbf{Z}/n\mathbf{Z}) \cong \mathbf{Z}/n\mathbf{Z}.
$$
第一个识别是典范的，第二个不是；见备注 \ref{etale-cohomology-remark-normalize-H1-Gm}。
由于 Riemann--Hurwitz 公式的证明不使用上同调计算，上述论证实际上构成一个证明
（只要补充消失等事项的细节）。




\section{非挠系数}
\label{etale-cohomology-section-nontorsion}

\noindent
为研究非挠系数，作如下定义：
$$
H_\etale^i (X, \mathbf{Q}_\ell) :=
\left( \lim_n H_\etale^i(X, \mathbf{Z}/\ell^n\mathbf{Z}) \right)
\otimes_{\mathbf{Z}_\ell} \mathbf{Q}_\ell.
$$
符号 $\lim_n$ 表示上同调群系统
$H_\etale^i(X, \mathbf{Z}/\ell^n\mathbf{Z})$ 按 $n$ 编号的 {\it 极限}，见
Categories, Section \ref{categories-section-posets-limits}。
因此，我们需要研究满足某些相容条件的层系统。




\section{层理论}
\label{etale-cohomology-section-sheaf-theory}
%9.10.09

\noindent
从这里开始，我们正式讨论网站与层。接下来的几节可以容纳数量惊人的有用
抽象材料。其中一部分已经在前面的章节中展开，例如网站、网站上的模以及
网站上的上同调等章节。这里尽量不加入过多材料，只保留足以使本章后续内容
能够理解的部分。




\section{预层}
\label{etale-cohomology-section-presheaves}

\noindent
本节的参考资料是
Sites，第 \ref{sites-section-presheaves} 节。

\begin{definition}
\label{etale-cohomology-definition-presheaf}
设 $\mathcal{C}$ 是一个范畴。$\mathcal{C}$ 上的{\it 集合预层}（分别地，
{\it 阿贝尔预层}）是函子 $\mathcal{C}^{opp} \to
\textit{Sets}$（分别地，$\textit{Ab}$）。
\end{definition}

\noindent
{\bf 术语。}若 $U \in \Ob(\mathcal{C})$，则 $\mathcal{F}(U)$ 的元素称为
$\mathcal{F}$ 在 $U$ 上的{\it 截面}。对 $\varphi : V \to U$（其中它是
$\mathcal{C}$ 中的态射），映射 $\mathcal{F}(\varphi) : \mathcal{F}(U) \to \mathcal{F}(V)$
称为{\it 限制映射}，通常记作 $s \mapsto s|_V$，有时也记作
$s \mapsto \varphi^*s$。记号 $s|_V$ 有歧义，因为限制映射取决于 $\varphi$，
但这是约定俗成的记号滥用。我们还使用记号
$\Gamma(U, \mathcal{F}) = \mathcal{F}(U)$。

\medskip\noindent
说 $\mathcal{F}$ 是函子，意味着如果 $W \to V \to U$ 是
$\mathcal{C}$ 中的态射，且 $s \in \Gamma(U, \mathcal{F})$，那么
$(s|_V)|_W = s |_W$，这里仍使用上述记号滥用。此外，与恒等态射
$\text{id}_U : U \to U$ 对应的限制映射就是恒等映射。

\medskip\noindent
在 $\mathcal{C}$ 上的集合预层（分别地，阿贝尔预层）的范畴记作
$\textit{PSh} (\mathcal{C})$（分别地，$\textit{PAb}
(\mathcal{C})$）。它是从 $\mathcal{C}^{opp}$ 到
$\textit{Sets}$（分别地，$\textit{Ab}$）的函子范畴，也就是说，预层之间的
态射是函子的自然变换。我们只考虑
$\textit{PSh}(\mathcal{C})$ 和 $\textit{PAb}(\mathcal{C})$，并且此处范畴
$\mathcal{C}$ 为小范畴。
（我们的约定是：除非另有说明，范畴都是小的；若范畴不是小的，应列在
Categories，备注 \ref{categories-remark-big-categories} 中。）

\begin{example}
\label{etale-cohomology-example-representable-presheaf}
给定对象 $X \in \Ob(\mathcal{C})$，考虑如下规则：
$$
\begin{matrix}
U & \longmapsto & h_X(U) = \Mor_\mathcal{C}(U, X) \\
V \xrightarrow{\varphi} U & \longmapsto &
(\psi \mapsto \psi \circ \varphi) : h_X(U) \to h_X(V).
\end{matrix}
$$
这定义了函子 $h_X : \mathcal{C}^{opp} \to \textit{Sets}$，因而定义了一个预层。
它称为与 $X$ 关联的{\it 可表预层}。并非在每个网站的每个拓扑中，
可表预层都是层。
\end{example}

\begin{lemma}[Yoneda]
\label{etale-cohomology-lemma-yoneda}
\begin{slogan}
对象之间的态射与它们所表示的函子之间的自然变换构成一一对应。
\end{slogan}
设 $\mathcal{C}$ 是一个范畴，且 $X, Y \in
\Ob(\mathcal{C})$。存在自然双射
$$
\begin{matrix}
\Mor_\mathcal{C}(X, Y) &
\longrightarrow &
\Mor_{\textit{PSh}(\mathcal{C})} (h_X, h_Y) \\
\psi &
\longmapsto &
h_\psi = \psi \circ - : h_X \to h_Y.
\end{matrix}
$$
\end{lemma}

\begin{proof}
见
Categories，引理 \ref{categories-lemma-yoneda}。
\end{proof}




\section{网站}
\label{etale-cohomology-section-sites}


\begin{definition}
\label{etale-cohomology-definition-family-morphisms-fixed-target}
设 $\mathcal{C}$ 是一个范畴。一个以 U 为固定目标的{\it 态射族}
$\mathcal{U} = \{\varphi_i : U_i \to U\}_{i\in I}$ 由以下数据组成：
\begin{enumerate}
\item 对象 $U \in \mathcal{C}$；
\item 集合 $I$（可以为空）；以及
\item 对每个 $i\in I$，给出态射 $\varphi_i : U_i \to U$，它属于
$\mathcal{C}$ 且以 $U$ 为目标。
\end{enumerate}
\end{definition}

\noindent
我们有以固定目标的态射族之间的{\it 态射}这一概念；其特殊情形就是
{\it 加细}。这部分材料的参考资料是
Sites，第 \ref{sites-section-refinements} 节。

\begin{definition}
\label{etale-cohomology-definition-site}
一个{\it 网站}\footnote{我们称为 site 的对象，在
\cite[Expos\'e II, D\'efinition 1.3]{SGA4} 中称为带有预拓扑的范畴；
在 \cite{ArtinTopologies} 中称为带有 Grothendieck 拓扑的范畴。}
由范畴 $\mathcal{C}$ 及集合 $\text{Cov}(\mathcal{C})$ 组成；后者由以固定目标的
态射族构成，称为{\it 覆盖}，并满足：
\begin{enumerate}
\item （同构）若 $\varphi : V \to U$ 是 $\mathcal{C}$ 中的同构，
则 $\{\varphi : V \to U\}$ 是覆盖；
\item （局部性）若 $\{\varphi_i : U_i \to U\}_{i\in I}$ 是覆盖，并且
对每个 $i \in I$ 给定覆盖
$\{\psi_{ij} : U_{ij} \to U_i \}_{j\in I_i}$，则
$$
\{
\varphi_i \circ \psi_{ij} : U_{ij} \to U
\}_{(i, j)\in \prod_{i\in I} \{i\} \times I_i}
$$
也是覆盖；以及
\item （基变换）若 $\{U_i \to U\}_{i\in I}$
是覆盖，且 $V \to U$ 是 $\mathcal{C}$ 中的态射，则
\begin{enumerate}
\item 对每个 $i \in I$，纤维积 $U_i \times_U V$ 在 $\mathcal{C}$ 中存在；并且
\item $\{U_i \times_U V \to V\}_{i\in I}$ 是覆盖。
\end{enumerate}
\end{enumerate}
\end{definition}

\noindent
对我们而言，网站的底层范畴总是“``小''”的，即其对象的集合构成一个集合，
网站的覆盖集合也同样是一个集合（如上面的定义所述）。在本章中，我们通常省略
那些把对象族和覆盖族限制为集合的论证。进一步讨论见
Sites，备注 \ref{sites-remark-no-big-sites}。

\begin{example}
\label{etale-cohomology-example-site-topological-space}
若 $X$ 是拓扑空间，则有如下定义的关联网站 $X_{Zar}$：$X_{Zar}$ 的对象是
$X$ 的开子集，它们之间的态射是包含映射，覆盖是通常的拓扑（满射）覆盖。
注意，若 $U, V \subset W \subset X$ 是开子集，则 $U \times_W V = U \cap V$
存在；也就是说，这个范畴有纤维积。所有验证都是平凡的，结果正如预期。
\end{example}




\section{层}
\label{etale-cohomology-section-sheaves}

\begin{definition}
\label{etale-cohomology-definition-sheaf}
集合预层 $\mathcal{F}$（分别地，阿贝尔预层）在网站 $\mathcal{C}$ 上称为
{\it 分离预层}，如果对每个覆盖
$\{\varphi_i : U_i \to U\}_{i\in I} \in \text{Cov} (\mathcal{C})$，映射
$$
\mathcal{F}(U) \longrightarrow \prod\nolimits_{i\in I} \mathcal{F}(U_i)
$$
是单射。这里该映射为 $s \mapsto (s|_{U_i})_{i\in I}$。
预层 $\mathcal{F}$ 若对每个覆盖 $\{\varphi_i : U_i \to U\}_{i\in I} \in \text{Cov} (\mathcal{C})$，图
\begin{equation}
\label{etale-cohomology-equation-sheaf-axiom}
\xymatrix{
\mathcal{F}(U) \ar[r] &
\prod_{i\in I} \mathcal{F}(U_i) \ar@<1ex>[r] \ar@<-1ex>[r] &
\prod_{i, j \in I} \mathcal{F}(U_i \times_U U_j),
}
\end{equation}
其中第一个映射是 $s \mapsto (s|_{U_i})_{i\in I}$，右侧两个映射分别为
$(s_i)_{i\in I} \mapsto (s_i |_{U_i \times_U U_j})$ 和
$(s_i)_{i\in I} \mapsto (s_j |_{U_i \times_U U_j})$；该图在集合范畴
（分别地，阿贝尔群范畴）中是等化子图。
\end{definition}

\begin{remark}
\label{etale-cohomology-remark-empty-covering}
对于空覆盖（即 $I = \emptyset$），这意味着 $\mathcal{F}(\emptyset)$ 是空积，
它是相应范畴中的终对象（对 $\textit{Sets}$ 和 $\textit{Ab}$ 都是单元素集）。
\end{remark}

\begin{example}
\label{etale-cohomology-example-sheaf-site-space}
对拓扑空间关联的网站 $X_{Zar}$ 作上述展开（见例 \ref{etale-cohomology-example-site-topological-space}），
得到通常的层概念。
\end{example}

\begin{definition}
\label{etale-cohomology-definition-category-sheaves}
我们用 $\Sh(\mathcal{C})$（分别地，$\textit{Ab}(\mathcal{C})$）表示
$\textit{PSh}(\mathcal{C})$（分别地，$\textit{PAb}(\mathcal{C})$）的满子范畴，
其对象是层。这就是 $\mathcal{C}$ 上的{\it 集合层范畴}
（分别地，{\it 阿贝尔层范畴}）。
\end{definition}




\section{G-集的例子}
\label{etale-cohomology-section-G-sets}

\noindent
设 $G$ 是群，并如下定义网站 $\mathcal{T}_G$：其底层范畴是 $G$-集范畴，
即对象是带有左 $G$-作用的集合，态射是等变映射；$\mathcal{T}_G$ 的覆盖是态射族
$\{\varphi_i : U_i \to U\}_{i\in I}$，满足
$U = \bigcup_{i\in I} \varphi_i(U_i)$。

\medskip\noindent
网站 $\mathcal{T}_G$ 中有一个特殊对象，即带有左平移自然作用的 $G$-集 $G$。
记作 ${}_G G$。注意存在自然群同构
$$
\begin{matrix}
\rho : & G^{opp} & \longrightarrow & \text{Aut}_{G\textit{-Sets}}({}_G G) \\
& g & \longmapsto & (h \mapsto hg).
\end{matrix}
$$
特别地，对任意预层 $\mathcal{F}$，集合 $\mathcal{F}({}_G G)$ 继承一个
$G$-作用（通过 $\rho$）。（注意，由于 $\mathcal{F}$ 是反变的，集合 $\mathcal{F}({}_G G)$
仍是左 $G$-集。）事实上，函子
$$
\begin{matrix}
\Sh(\mathcal{T}_G) & \longrightarrow & G\textit{-Sets} \\
\mathcal{F} & \longmapsto & \mathcal{F}({}_G G)
\end{matrix}
$$
是范畴等价。其拟逆函子是 $X \mapsto
h_X$。这里不提供完整证明（可见
Sites，第 \ref{sites-section-example-sheaf-G-sets} 节），而尝试解释其原因。
\begin{enumerate}
\item
若 $S$ 是 $G$-集，可以将其分解为轨道 $S = \coprod_{i\in I}
O_i$。
覆盖 $\{O_i \to S\}_{i\in I}$ 的层公理说明
$$
\xymatrix{
\mathcal{F}(S) \ar[r] &
\prod_{i\in I} \mathcal{F}(O_i) \ar@<1ex>[r] \ar@<-1ex>[r] &
\prod_{i, j \in I} \mathcal{F}(O_i \times_S O_j)
}
$$
是等化子。注意 $G\textit{-Sets}$ 中的纤维积由 $\textit{Sets}$ 中的纤维积诱导，
并利用 $\mathcal{F}(\emptyset)$ 是单元素 $G$-集这一事实，可得
$$
\prod_{i, j \in I} \mathcal{F}(O_i \times_S O_j) = \prod_{i \in I}
\mathcal{F}(O_i)
$$
且上面的两个映射实际上相同。因此层公理仅仅说明
$\mathcal{F}(S) = \prod_{i\in I} \mathcal{F}(O_i)$。
\item
若 $S$ 是 $G$-集 $S= G/H$，且 $\mathcal{F}$ 是 $\mathcal{T}_G$ 上的层，
则有
$$
\mathcal{F}(G/H) = \mathcal{F}({}_G G)^H
$$
特别地，$\mathcal{F}(\{*\}) = \mathcal{F}({}_G G)^G$。为此，使用覆盖
$\{ {}_G G \to G/H \}$ 的层公理（该覆盖覆盖 $S$）。我们有
\begin{eqnarray*}
{}_G G \times_{G/H} {}_G G & \cong & G \times H \\
(g_1, g_2) & \longmapsto & (g_1, g_1^{-1} g_2)
\end{eqnarray*}
是 ${}_G G$ 的不交并（作为 $G$-集）。因此层公理变为
$$
\xymatrix{
\mathcal{F} (G/H) \ar[r] &
\mathcal{F}({}_G G) \ar@<1ex>[r] \ar@<-1ex>[r] &
\prod_{h\in H} \mathcal{F}({}_G G)
}
$$
其中右侧两个映射为 $s \mapsto (s)_{h \in H}$ 和 $s \mapsto
(hs)_{h \in H}$。因此 $\mathcal{F}(G/H) = \mathcal{F}({}_G G)^H$，如所声称。
\end{enumerate}
这还不能完全证明所声称的范畴等价，但至少说明层 $\mathcal{F}$ 完全由其在
${}_G G$ 上的截面决定。细节（以及集合论方面的说明）见
Sites，第 \ref{sites-section-example-sheaf-G-sets} 节。




\section{层化}
\label{etale-cohomology-section-sheafification}

\begin{definition}
\label{etale-cohomology-definition-0-cech}
设 $\mathcal{F}$ 是网站 $\mathcal{C}$ 上的预层，且
$\mathcal{U} = \{U_i \to U\} \in \text{Cov} (\mathcal{C})$。
定义 $\mathcal{F}$ 关于 $\mathcal{U}$ 的{\it 零阶 {\v C}ech 上同调群}为
$$
\check H^0 (\mathcal{U}, \mathcal{F}) =
\left\{
(s_i)_{i\in I} \in \prod\nolimits_{i\in I }\mathcal{F}(U_i)
\text{ 使得 }
s_i|_{U_i \times_U U_j} = s_j |_{U_i \times_U U_j}
\right\}.
$$
\end{definition}

\noindent
存在典范映射
$\mathcal{F}(U) \to \check H^0 (\mathcal{U}, \mathcal{F})$，
$s \mapsto (s |_{U_i})_{i\in I}$。
从覆盖 $\mathcal{V} = \{V_j \to V\}_{j \in J}$ 到 $\mathcal{U}$ 的
{\it 覆盖态射}定义为三元组 $(\chi, \alpha, \chi_j)$，其中
$\chi : V \to U$ 是态射，$\alpha : J \to I$ 是集合映射，并且对每个
$j \in J$，态射 $\chi_j$ 使下图交换：
$$
\xymatrix{
V_j \ar[rr]_{\chi_j} \ar[d] & & U_{\alpha(j)} \ar[d] \\
V \ar[rr]^\chi & & U.
}
$$
给定数据 $\chi, \alpha, \{\chi_j\}_{j \in J}$，定义
\begin{eqnarray*}
\check H^0(\mathcal{U}, \mathcal{F}) & \longrightarrow &
\check H^0(\mathcal{V}, \mathcal{F}) \\
(s_i)_{i\in I} & \longmapsto &
\left(\chi_j^*\left(s_{\alpha(j)}\right)\right)_{j\in J}.
\end{eqnarray*}
我们断言：
\begin{enumerate}
\item 该映射定义良好；并且
\item 它只依赖于 $\chi$，与 $\alpha, \{\chi_j\}_{j \in J}$ 的选择无关。
\end{enumerate}
第一点的证明从略。为说明第 (2) 点，考虑另一个三元组
$(\psi, \beta, \psi_j)$，满足 $\chi = \psi$。于是有交换图
$$
\xymatrix{
V_j \ar[rrr]_{(\chi_j, \psi_j)} \ar[dd] & & &
U_{\alpha(j)} \times_U U_{\beta(j)} \ar[dl] \ar[dr] \\
& & U_{\alpha(j)} \ar[dr] & &
U_{\beta(j)} \ar[dl] \\
V \ar[rrr]^{\chi = \psi} & & & U.
}
$$
给定截面 $s \in \mathcal{F}(\mathcal{U})$，其在 $\mathcal{F}(V_j)$ 中、由
$(\chi, \alpha, \{\chi_j\}_{j \in J})$ 给出的映射下的像为
$\chi_j^*s_{\alpha(j)}$；在由
$(\psi, \beta, \{\psi_j\}_{j \in J})$ 给出的映射下，其像为
$\psi_j^*s_{\beta(j)}$。二者相等，因为按假设
$s \in \check H^0(\mathcal{U}, \mathcal{F})$，所以二者都等于共同值
$$
s_{\alpha(j)}|_{U_{\alpha(j)} \times_U U_{\beta(j)}} =
s_{\beta(j)}|_{U_{\alpha(j)} \times_U U_{\beta(j)}}
$$
沿图中的映射 $(\chi_j, \psi_j)$ 的拉回。

\begin{theorem}
\label{etale-cohomology-theorem-sheafification}
设 $\mathcal{C}$ 是网站，$\mathcal{F}$ 是 $\mathcal{C}$ 上的预层。
\begin{enumerate}
\item 规则
$$
U \mapsto \mathcal{F}^+(U) :=
\colim_{\mathcal{U} \text{ 为 }U\text{ 的覆盖}}
\check H^0(\mathcal{U}, \mathcal{F})
$$
是预层，并且该余极限是有向余极限。
\item 存在预层的典范映射 $\mathcal{F} \to \mathcal{F}^+$。
\item 若 $\mathcal{F}$ 是分离预层，则 $\mathcal{F}^+$ 是层，且 (2) 中的映射是单射。
\item $\mathcal{F}^+$ 是分离预层。
\item $\mathcal{F}^\# = (\mathcal{F}^+)^+$ 是层，并且典范映射诱导函子同构
$$
\Hom_{\textit{PSh}(\mathcal{C})}(\mathcal{F}, \mathcal{G}) =
\Hom_{\Sh(\mathcal{C})}(\mathcal{F}^\#, \mathcal{G})
$$
对任意 $\mathcal{G} \in \Sh(\mathcal{C})$ 都成立。
\end{enumerate}
\end{theorem}

\begin{proof}
见 Sites，定理 \ref{sites-theorem-plus}。
\end{proof}

\noindent
换言之，这意味着自然映射 $\mathcal{F} \to \mathcal{F}^\#$ 是遗忘函子
$\Sh(\mathcal{C}) \to \textit{PSh}(\mathcal{C})$ 的左伴随。




\section{上同调}
\label{etale-cohomology-section-cohomology}

\noindent
下面的基本结果使我们能够定义网站上阿贝尔层的上同调。

\begin{theorem}
\label{etale-cohomology-theorem-enough-injectives}
网站上的阿贝尔层范畴是一个具有足够内射对象的阿贝尔范畴。
\end{theorem}

\begin{proof}
见
Modules on Sites，引理 \ref{sites-modules-lemma-abelian-abelian}，以及
Injectives，定理 \ref{injectives-theorem-sheaves-injectives}。
\end{proof}

\noindent
因此，可以将上同调定义为截面函子的右导出函子：若 $U \in \Ob(\mathcal{C})$ 且
$\mathcal{F} \in \textit{Ab}(\mathcal{C})$，则
$$
H^p(U, \mathcal{F}) :=
R^p\Gamma(U, \mathcal{F}) =
H^p(\Gamma(U, \mathcal{I}^\bullet))
$$
其中 $\mathcal{F} \to \mathcal{I}^\bullet$ 是内射分解。为此需要检查函子
$\Gamma(U, -)$ 是左正合的。这一点成立，也是范畴 $\textit{Ab}(\mathcal{C})$
成为阿贝尔范畴的原因之一，见
Modules on Sites，引理 \ref{sites-modules-lemma-abelian-abelian}。
关于网站上上同调的更一般讨论（包括整体截面函子及其右导出函子），见
Cohomology on Sites，第 \ref{sites-cohomology-section-cohomology-sheaves} 节。



\section{fpqc 拓扑}
\label{etale-cohomology-section-fpqc}
%9.15.09

\noindent
在讨论 \'etale 上同调之前，我们先研究一些 fpqc 拓扑，因为它对拟凝聚层很有效。

\begin{definition}
\label{etale-cohomology-definition-fpqc-covering}
设 $T$ 是概形。$T$ 的{\it fpqc 覆盖}是满足以下条件的族
$\{ \varphi_i : T_i \to T\}_{i \in I}$：
\begin{enumerate}
\item 每个 $\varphi_i$ 都是平坦态射，且
$\bigcup_{i\in I} \varphi_i(T_i) = T$，并且
\item 对每个仿射开集 $U \subset T$，存在有限集合 $K$、映射
$\mathbf{i} : K \to I$ 以及仿射开集
$U_{\mathbf{i}(k)} \subset T_{\mathbf{i}(k)}$，使得
$U = \bigcup_{k \in K} \varphi_{\mathbf{i}(k)}(U_{\mathbf{i}(k)})$.
\end{enumerate}
\end{definition}

\begin{remark}
\label{etale-cohomology-remark-fpqc}
第一个条件对应 fp，它代表法语 {\it fid\`element plat}（忠实平坦）；
第二个条件对应 qc，即{\it quasi-compact}（拟紧）。当第二个条件成立时，
第一个条件的后半部分是不必要的。
\end{remark}

\begin{example}
\label{etale-cohomology-example-fpqc-coverings}
fpqc 覆盖的例子。
\begin{enumerate}
\item $T$ 的任意 Zariski 开覆盖都是 fpqc 覆盖。
\item 族 $\{\Spec(B) \to \Spec(A)\}$ 是 fpqc 覆盖，当且仅当 $A \to B$
是忠实平坦的环映射。
\item 若 $f: X \to Y$ 平坦、满射且拟紧，则 $\{ f: X\to
Y\}$ 是 fpqc 覆盖。
\item 态射
$\varphi :
\coprod_{x \in \mathbf{A}^1_k} \Spec(\mathcal{O}_{\mathbf{A}^1_k, x})
\to \mathbf{A}^1_k$,
其中 $k$ 是域；该态射平坦且满射，但不是拟紧的，事实上族 $\{\varphi\}$
不是 fpqc 覆盖。
\item 记
$\mathbf{A}^2_k = \Spec(k[x, y])$。记 $i_x : D(x) \to \mathbf{A}^2_k$
和 $i_y : D(y) \to \mathbf{A}^2_k$ 为标准开子集。
那么族
$\{i_x, i_y, \Spec(k[[x, y]]) \to \mathbf{A}^2_k\}$
以及
$\{i_x, i_y, \Spec(\mathcal{O}_{\mathbf{A}^2_k, 0}) \to \mathbf{A}^2_k\}$
都是 fpqc 覆盖。
\end{enumerate}
\end{example}

\begin{lemma}
\label{etale-cohomology-lemma-site-fpqc}
概形范畴上的 fpqc 覆盖族满足网站公理。
\end{lemma}

\begin{proof}
见 Topologies，引理 \ref{topologies-lemma-fpqc}。
\end{proof}

\noindent
看起来，这个引理允许我们定义概形范畴的 fpqc 网站。然而，考虑 fpqc 拓扑时会
出现集合论问题，见 Topologies，第 \ref{topologies-section-fpqc} 节。
问题源于我们要求网站是“``小''”的，但没有任何小的概形范畴能够包含给定非空概形的
fpqc 覆盖的余最终系统。严格说来，这并不妨碍定义“部分”的 fpqc 网站，
但这样做似乎并不明智。变通办法是允许使用 fpqc 拓扑的层概念（见下文），
但禁止考虑所有 fpqc 层所成的范畴。

\begin{definition}
\label{etale-cohomology-definition-sheaf-property-fpqc}
设 $S$ 是概形。$S$ 上概形的范畴记作
$\Sch/S$。考虑函子
$\mathcal{F} : (\Sch/S)^{opp} \to \textit{Sets}$，即集合预层。称 $\mathcal{F}$
{\it 满足 fpqc 拓扑的层性质}，如果对每个 fpqc 覆盖
$\{U_i \to U\}_{i \in I}$（其覆盖 $S$ 上的概形），图 (\ref{etale-cohomology-equation-sheaf-axiom}) 是等化子图。
\end{definition}

\noindent
类似地，称 $\mathcal{F}$ {\it 满足 Zariski 拓扑的层性质}，如果对每个开覆盖
$U = \bigcup_{i \in I} U_i$，图 (\ref{etale-cohomology-equation-sheaf-axiom}) 是等化子图。见
Schemes，定义 \ref{schemes-definition-representable-by-open-immersions}。
显然，这等价于：对每个概形 $T$（位于 $S$ 上），将 $\mathcal{F}$ 限制到 $T$ 的开集上
所得的是一个（通常意义下的）层。

\begin{lemma}
\label{etale-cohomology-lemma-fpqc-sheaves}
设 $\mathcal{F}$ 是 $\Sch/S$ 上的预层。则 $\mathcal{F}$ 满足 fpqc 拓扑的层性质
当且仅当
\begin{enumerate}
\item $\mathcal{F}$ 满足 Zariski 拓扑的层性质；并且
\item 对每个忠实平坦态射 $\Spec(B) \to \Spec(A)$（它是 $S$ 上仿射概形之间的态射），
覆盖 $\{\Spec(B) \to \Spec(A)\}$ 满足层公理。也就是说，
$$
\xymatrix{
\mathcal{F}(\Spec(A)) \ar[r] &
\mathcal{F}(\Spec(B)) \ar@<1ex>[r] \ar@<-1ex>[r] &
\mathcal{F}(\Spec(B \otimes_A B))
}
$$
是等化子图。
\end{enumerate}
\end{lemma}

\begin{proof}
见 Topologies，引理 \ref{topologies-lemma-sheaf-property-fpqc}。
\end{proof}

\noindent
对于满足 fpqc 拓扑层条件的预层 $\mathcal{F}$（它定义在 $\Sch/S$ 上），
另一种理解方式是以下数据：
\begin{enumerate}
\item 对每个 $T/S$，给出通常意义（即 Zariski）的层 $\mathcal{F}_T$（位于 $T_{Zar}$ 上）；
\item 对每个映射 $f : T' \to T$（位于 $S$ 上），给出限制映射
$f^{-1}\mathcal{F}_T \to \mathcal{F}_{T'}$；
\end{enumerate}
满足
\begin{enumerate}
\item[(a)] 限制映射满足函子性；
\item[(b)] 若 $f : T' \to T$ 是开浸入，则限制映射
$f^{-1}\mathcal{F}_T \to \mathcal{F}_{T'}$ 是同构；并且
\item[(c)] 对每个忠实平坦态射
$\Spec(B) \to \Spec(A)$（位于 $S$ 上），图
$$
\xymatrix{
\mathcal{F}_{\Spec(A)}(\Spec(A)) \ar[r] &
\mathcal{F}_{\Spec(B)}(\Spec(B)) \ar@<1ex>[r] \ar@<-1ex>[r] &
\mathcal{F}_{\Spec(B \otimes_A B)}(\Spec(B \otimes_A B))
}
$$
是等化子。
\end{enumerate}
数据 (1)、(2) 以及条件 (a)、(b) 给出 $\Sch/S$ 上满足 Zariski 拓扑层条件的预层。
由引理 \ref{etale-cohomology-lemma-fpqc-sheaves}，条件 (c) 即足以得到 fpqc 拓扑的层条件。

\begin{example}
\label{etale-cohomology-example-quasi-coherent}
考虑预层
$$
\begin{matrix}
\mathcal{F} : & (\Sch/S)^{opp} & \longrightarrow & \textit{Ab} \\
& T/S & \longmapsto & \Gamma(T, \Omega_{T/S}).
\end{matrix}
$$
微分与局部化的相容性说明，$\mathcal{F}$ 是 Zariski 网站上的层。
然而，它不满足 fpqc 拓扑的层条件。具体地，考虑
$S = \Spec(\mathbf{F}_p)$ 以及态射
$$
\varphi :
V = \Spec(\mathbf{F}_p[v])
\to
U = \Spec(\mathbf{F}_p[u])
$$
它把 $u$ 映到 $v^p$。族 $\{\varphi\}$ 是 fpqc 覆盖，
但限制映射
$\mathcal{F}(U) \to \mathcal{F}(V)$
把生成元 $\text{d}u$ 映到 $\text{d}(v^p) = 0$，所以它是零映射，而图
$$
\xymatrix{
\mathcal{F}(U) \ar[r]^{0} &
\mathcal{F}(V) \ar@<1ex>[r] \ar@<-1ex>[r] &
\mathcal{F}(V \times_U V)
}
$$
不是等化子。稍后将看到，$\mathcal{F}$ 的确在 \'etale 网站与光滑网站上
给出一个层。
\end{example}

\begin{lemma}
\label{etale-cohomology-lemma-representable-sheaf-fpqc}
$\Sch/S$ 上任意可表预层都满足 fpqc 拓扑的层条件。
\end{lemma}

\begin{proof}
见 Descent，引理 \ref{descent-lemma-fpqc-universal-effective-epimorphisms}。
\end{proof}

\noindent
稍后还会回到这一点，因为该事实的证明要用到拟凝聚层的下降，而这将在下一节
讨论。用一种较为精致的说法，该引理断言
{\it fpqc 拓扑弱于典范拓扑}，或者说 fpqc 拓扑是{\it 次典范的}。
在网站的框架下，这一点见 Sites，第 \ref{sites-section-representable-sheaves} 节。

\begin{remark}
\label{etale-cohomology-remark-fpqc-finest}
fpqc 拓扑比 Zariski、\'etale、光滑、syntomic 与 fppf 拓扑都更细。因此，
任意满足 fpqc 拓扑层条件的预层，在 Zariski、\'etale、光滑、syntomic 与
fppf 网站上都是层。特别地，例如可表预层在概形的 \'etale 网站上是层。
\end{remark}

\begin{example}
\label{etale-cohomology-example-additive-group-sheaf}
设 $S$ 是概形。考虑加法群概形
$\mathbf{G}_{a, S} = \mathbf{A}^1_S$（在 $S$ 上），见
Groupoids，例 \ref{groupoids-example-additive-group}。
关联的可表预层由下式给出：
$$
h_{\mathbf{G}_{a, S}}(T) =
\Mor_S(T, \mathbf{G}_{a, S}) =
\Gamma(T, \mathcal{O}_T).
$$
由上文可知，这是满足 fpqc 拓扑层条件的集合预层。另一方面，它显然也是环预层。
因此可以将其看作函子
$$
\begin{matrix}
\mathcal{O} : &
(\Sch/S)^{opp} &
\longrightarrow &
\textit{Rings} \\
&
T/S &
\longmapsto &
\Gamma(T, \mathcal{O}_T)
\end{matrix}
$$
且满足 fpqc 拓扑的层条件。相应地，还可以定义 $\mathcal{O}$-模，等等。
\end{example}




\section{忠实平坦下降}
\label{etale-cohomology-section-fpqc-descent}

\noindent
本节讨论拟凝聚模的忠实平坦下降。更准确地说，我们将证明：对于 fpqc 覆盖，
拟凝聚模满足有效下降。

\begin{definition}
\label{etale-cohomology-definition-descent-datum}
设 $\mathcal{U} = \{ t_i : T_i \to T\}_{i \in I}$ 是以固定目标为目标的概形态射族。
关于 $\mathcal{U}$ 的拟凝聚层{\it 下降数据}是集合
$((\mathcal{F}_i)_{i \in I}, (\varphi_{ij})_{i, j \in I})$，其中
\begin{enumerate}
\item $\mathcal{F}_i$ 是 $T_i$ 上的拟凝聚层；
\item $\varphi_{ij} : \text{pr}_0^* \mathcal{F}_i \to
\text{pr}_1^* \mathcal{F}_j$ 是 $T_i \times_T T_j$ 上的模同构；
\end{enumerate}
并满足{\it 循环条件}：下图
$$
\xymatrix{
\text{pr}_0^*\mathcal{F}_i \ar[dr]_{\text{pr}_{02}^*\varphi_{ik}}
\ar[rr]^{\text{pr}_{01}^*\varphi_{ij}} & &
\text{pr}_1^*\mathcal{F}_j \ar[dl]^{\text{pr}_{12}^*\varphi_{jk}} \\
& \text{pr}_2^*\mathcal{F}_k
}
$$
在 $T_i \times_T T_j \times_T T_k$ 上交换。
若存在拟凝聚层 $\mathcal{F}$ 位于 $T$ 上，以及
$\mathcal{O}_{T_i}$-模同构 $\varphi_i : t_i^* \mathcal{F} \cong \mathcal{F}_i$，
且它们与映射 $\varphi_{ij}$ 相容，即
$$
\varphi_{ij} = \text{pr}_1^* (\varphi_j) \circ \text{pr}_0^* (\varphi_i)^{-1}.
$$
\end{definition}

\noindent
本节及下一节讨论下述定理证明中的一些成分，以及相关材料。

\begin{theorem}
\label{etale-cohomology-theorem-descent-quasi-coherent}
若 $\mathcal{V} = \{T_i \to T\}_{i\in I}$ 是 fpqc 覆盖，则关于 $\mathcal{V}$ 的
所有拟凝聚层下降数据都是有效的。
\end{theorem}

\begin{proof}
见
Descent，命题 \ref{descent-proposition-fpqc-descent-quasi-coherent}。
\end{proof}

\noindent
换言之，拟凝聚层的纤维范畴是 fpqc 网站上的一个叠。
定理的证明分两步。第一步是注意到对 Zariski 覆盖，这一点利用层的标准粘合（见
Sheaves，第 \ref{sheaves-section-glueing-sheaves} 节）以及拟凝聚性的局部性，
很容易（或众所周知）。第二步是处理形如 $\{\Spec(B) \to \Spec(A)\}$ 的 fpqc 覆盖，
其中 $A \to B$ 是忠实平坦的环映射。这是一个代数引理，现予以陈述。

\medskip\noindent
{\bf 模的下降。}若 $A \to B$ 是环映射，考虑复形
$$
(B/A)_\bullet : B \to B \otimes_A B \to B \otimes_A B \otimes_A B \to \ldots
$$
其中 $B$ 位于次数 0，$B \otimes_A B$ 位于次数 1，依此类推；映射由下式给出：
\begin{eqnarray*}
b & \mapsto & 1 \otimes b - b \otimes 1, \\
b_0 \otimes b_1 & \mapsto & 1 \otimes b_0 \otimes b_1 - b_0 \otimes 1 \otimes
b_1 + b_0 \otimes b_1 \otimes 1, \\
& \text{etc.}
\end{eqnarray*}

\begin{lemma}
\label{etale-cohomology-lemma-algebra-descent}
若 $A \to B$ 忠实平坦，则复形 $(B/A)_\bullet$ 在正次数上正合，且
$H^0((B/A)_\bullet) = A$。
\end{lemma}

\begin{proof}
见 Descent，引理 \ref{descent-lemma-ff-exact}。
\end{proof}

\noindent
Grothendieck 分三步证明这一点。首先，他假设映射 $A
\to B$ 有截面，并构造到如下复形的显式同伦：只有次数 0 的 $A$ 项非零。其次，
他指出，只需在忠实平坦基变换 $A \to
A'$ 后证明结果，此时将 $B$ 替换为 $B' = B \otimes_A A'$。第三，他取忠实平坦基变换
$A \to A' = B$，并注意映射
$A' = B \to B' = B \otimes_A B$ 有自然截面。

\medskip\noindent
同样的策略证明以下引理。

\begin{lemma}
\label{etale-cohomology-lemma-descent-modules}
若 $A \to B$ 忠实平坦且 $M$ 是 $A$-模，则复形
$(B/A)_\bullet \otimes_A M$ 在正次数上正合，且
$H^0((B/A)_\bullet \otimes_A M) = M$。
\end{lemma}

\begin{proof}
见 Descent，引理 \ref{descent-lemma-ff-exact}。
\end{proof}

\begin{definition}
\label{etale-cohomology-definition-descent-datum-modules}
设 $A \to B$ 是环映射，$N$ 是 $B$-模。$N$ 关于 $A \to B$ 的
一个{\it 下降数据}，是 $B \otimes_A B$-模同构
$\varphi : N \otimes_A B \cong B \otimes_A N$，使得下列
$B \otimes_A B \otimes_A B$-模图
$$
\xymatrix{
{N \otimes_A B \otimes_A B} \ar[dr]_{\varphi_{02}} \ar[rr]^{\varphi_{01}} & &
{B \otimes_A N \otimes_A B} \ar[dl]^{\varphi_{12}} \\
& {B \otimes_A B \otimes_A N}
}
$$
交换，其中 $\varphi_{01} = \varphi \otimes \text{id}_B$，而
$\varphi_{12}$ 与 $\varphi_{02}$ 以类似方式定义。
\end{definition}

\noindent
若 $N' = B \otimes_A M$，其中 M 是某个 $A$-模，则该模具有由下列映射
给出的典范下降数据：
$$
\begin{matrix}
\varphi_\text{can}: & N' \otimes_A B & \to & B \otimes_A N' \\
& b_0 \otimes m \otimes b_1 & \mapsto & b_0 \otimes b_1 \otimes m.
\end{matrix}
$$

\begin{definition}
\label{etale-cohomology-definition-effective-modules}
若存在 $A$-模 $M$，使得
$(N, \varphi) \cong (B \otimes_A M, \varphi_\text{can})$，则称下降数据
$(N, \varphi)$ 是{\it 有效的}；这里采用显然的下降数据同构概念。
\end{definition}

\noindent
定理 \ref{etale-cohomology-theorem-descent-quasi-coherent} 是下述结果的推论。

\begin{theorem}
\label{etale-cohomology-theorem-descent-modules}
若 $A \to B$ 忠实平坦，则关于 $A\to B$ 的下降数据都是有效的。
\end{theorem}

\begin{proof}
见 Descent，命题 \ref{descent-proposition-descent-module}。
另见 Descent，评注
\ref{descent-remark-homotopy-equivalent-cosimplicial-algebras}，其中给出证明的另一种看法。
\end{proof}

\begin{remarks}
\label{etale-cohomology-remarks-theorem-modules-exactness}
模的下降结果有若干应用：
\begin{enumerate}
\item 对覆盖 $\{\Spec(B) \to \Spec(A)\}$，其中 $A \to B$ 忠实平坦，
{\v C}ech 复形在正次数上正合。这将给出若干上同调消失结果。
\item 若 $(N, \varphi)$ 是关于忠实平坦映射 $A \to B$ 的下降数据，则相应的
$A$-模由下式给出：
$$
M = \Ker \left(
\begin{matrix}
N & \longrightarrow & B \otimes_A N \\
n & \longmapsto & 1 \otimes n - \varphi(n \otimes 1)
\end{matrix}
\right).
$$
见 Descent，命题 \ref{descent-proposition-descent-module}。
\end{enumerate}
\end{remarks}




%9.17.09
\section{拟凝聚层}
\label{etale-cohomology-section-quasi-coherent}

\noindent
我们可以利用模的下降来研究拟凝聚层。

\begin{proposition}
\label{etale-cohomology-proposition-quasi-coherent-sheaf-fpqc}
对任意拟凝聚层 $\mathcal{F}$（定义在 $S$ 上），预层
$$
\begin{matrix}
\mathcal{F}^a : & \Sch/S & \to & \textit{Ab}\\
& (f: T \to S) & \mapsto & \Gamma(T, f^*\mathcal{F})
\end{matrix}
$$
是满足 fpqc 拓扑层条件的 $\mathcal{O}$-模。
\end{proposition}

\begin{proof}
这在 Descent，引理 \ref{descent-lemma-sheaf-condition-holds} 中证明。
这里给出证明梗概。由引理 \ref{etale-cohomology-lemma-fpqc-sheaves} 可知，
只需检查仿射概形的 Zariski 覆盖和忠实平坦态射上的层性质。
Zariski 覆盖的层性质是标准概形理论，因为
$\Gamma(U, i^\ast \mathcal{F}) = \mathcal{F}(U)$，其中
$i : U \hookrightarrow S$ 是开浸入。

\medskip\noindent
对于 $\left\{\Spec(B)\to \Spec(A)\right\}$，若 $A\to B$ 忠实
平坦且 $\mathcal{F}|_{\Spec(A)} = \widetilde{M}$，这对应于
$M = H^0\left((B/A)_\bullet \otimes_A M \right)$，也就是说
\begin{align*}
0 \to M \to B \otimes_A M \to B \otimes_A B \otimes_A M
\end{align*}
由引理 \ref{etale-cohomology-lemma-descent-modules} 可知是正合的。
\end{proof}

\noindent
在环化网站上还有拟凝聚层的抽象概念。这里简要介绍它；
更多信息见 Modules on Sites，第 \ref{sites-modules-section-local} 节。
设 $\mathcal{C}$ 是一个范畴，$U$ 是 $\mathcal{C}$ 的对象。
记 $\mathcal{C}/U$ 为位于 $U$ 之上的对象所成的范畴，见
Categories，例 \ref{categories-example-category-over-X}。
若 $\mathcal{C}$ 是网站，则 $\mathcal{C}/U$ 也是网站；具体地，
$V/U$ 的覆盖是形如 $\{V_i/U \to V/U\}$ 的
$\mathcal{C}/U$ 态射族，并满足 $\{V_i \to V\}$ 是
$\mathcal{C}$ 的覆盖。此外，给定任意层 $\mathcal{F}$（定义在
$\mathcal{C}$ 上），其{\it 限制}
$\mathcal{F}|_{\mathcal{C}/U}$（按显然方式定义）也是层。详见
Sites，第 \ref{sites-section-localize} 节。

\begin{definition}
\label{etale-cohomology-definition-ringed-site}
设 $\mathcal{C}$ 是一个{\it 环化网站}，即带有环层
$\mathcal{O}$ 的网站。称 $\mathcal{O}$-模层 $\mathcal{F}$
在 $\mathcal{C}$ 上是{\it 拟凝聚的}，如果对所有
$U \in \Ob(\mathcal{C})$，都存在由
$\{U_i \to U\}_{i\in I}$ 给出的 $\mathcal{C}$ 覆盖，使得限制
$\mathcal{F}|_{\mathcal{C}/U_i}$ 同构于自由 $\mathcal{O}$-模之间某个
$\mathcal{O}$-线性映射的余核
$$
\bigoplus\nolimits_{k \in K} \mathcal{O}|_{\mathcal{C}/U_i}
\longrightarrow
\bigoplus\nolimits_{l \in L} \mathcal{O}|_{\mathcal{C}/U_i}.
$$
对 $K$ 的直和是与预层
$V \mapsto \bigoplus_{k \in K} \mathcal{O}(V)$ 关联的层，另一侧同理。
\end{definition}

\noindent
虽然能够给出上述一般定义很有用，但这个概念一般而言性质并不好。

\begin{remark}
\label{etale-cohomology-remark-final-object}
若 $\mathcal{C}$ 有终对象（例如 $S$），则上述定义中的条件只需
对 $U = S$ 检查即可。见 Modules on Sites，引理
\ref{sites-modules-lemma-local-final-object}。
\end{remark}

\begin{theorem}[拟凝聚层元定理]
\label{etale-cohomology-theorem-quasi-coherent}
设 $S$ 是概形。设 $\mathcal{C}$ 是一个网站，并假设
\begin{enumerate}
\item 底层范畴 $\mathcal{C}$ 是 $\Sch/S$ 的全子范畴；
\item $T \in \Ob(\mathcal{C})$ 的任意 Zariski 覆盖都可由
$\mathcal{C}$ 的某个覆盖细化；
\item $S/S$ 是 $\mathcal{C}$ 的对象；
\item $\mathcal{C}$ 的每个覆盖都是概形的 fpqc 覆盖。
\end{enumerate}
则预层 $\mathcal{O}$ 是 $\mathcal{C}$ 上的层，并且
任意拟凝聚 $\mathcal{O}$-模在 $(\mathcal{C}, \mathcal{O})$ 上都具有
$\mathcal{F}^a$ 的形式，其中 $\mathcal{F}$ 是定义在 $S$ 上的某个
拟凝聚层。
\end{theorem}

\begin{proof}
经过一些形式论证，这正是定理 \ref{etale-cohomology-theorem-descent-quasi-coherent}。
细节略去。在 Descent，命题
\ref{descent-proposition-equivalence-quasi-coherent} 中，我们对 $S$ 的
大 Zariski、fppf、\'etale、光滑和 syntomic 网站，以及 $S$ 的小 Zariski
和 \'etale 网站证明了该定理的更精确版本。
\end{proof}

\noindent
换言之，$S$ 上的拟凝聚模与定理中这类网站上的拟凝聚
$\mathcal{O}$-模（网站记为 $\mathcal{C}$）之间没有区别。关于大、小网站
$(\Sch/S)_{fppf}$、$S_\etale$ 等的更精确陈述，可见 Descent 的
以下各节：
\ref{descent-section-quasi-coherent-sheaves},
\ref{descent-section-quasi-coherent-cohomology}，和
\ref{descent-section-quasi-coherent-sheaves-bis}。
在本章中，为了方便陈述适用于上述任一情形的结果，
我们有时将其称为“定理 \ref{etale-cohomology-theorem-quasi-coherent} 中的这类网站”。






\section{{\v C}ech 上同调}
\label{etale-cohomology-section-cech-cohomology}

\noindent
下一目标是用下降理论证明等式
$H^i(\mathcal{C}, \mathcal{F}^a) = H_{Zar}^i(S, \mathcal{F})$
对所有拟凝聚层 $\mathcal{F}$（定义在 $S$ 上）以及定理
\ref{etale-cohomology-theorem-quasi-coherent} 中的任意网站 $\mathcal{C}$ 都成立。
为此，我们引入网站上的 {\v C}ech 上同调。更多细节见
\cite{ArtinTopologies}，以及 Cohomology on Sites 的各节
\ref{sites-cohomology-section-cech}、
\ref{sites-cohomology-section-cech-functor}
和 \ref{sites-cohomology-section-cech-cohomology-cohomology}。

\begin{definition}
\label{etale-cohomology-definition-cech-complex}
设 $\mathcal{C}$ 是一个范畴。设
$\mathcal{U} = \{U_i \to U\}_{i \in I}$ 是 $\mathcal{C}$ 中具有共同靶
的态射族。假设所有纤维积
$U_{i_0} \times_U \ldots \times_U U_{i_p}$ 都存在于 $\mathcal{C}$ 中。
设 $\mathcal{F} \in \textit{PAb}(\mathcal{C})$ 是阿贝尔预层。定义
{\it {\v C}ech 复形} $\check{\mathcal{C}}^\bullet(\mathcal{U}, \mathcal{F})$ 为
$$
\prod_{i_0 \in I} \mathcal{F}(U_{i_0}) \to
\prod_{i_0, i_1 \in I} \mathcal{F}(U_{i_0} \times_U U_{i_1}) \to
\prod_{i_0, i_1, i_2 \in I}
\mathcal{F}(U_{i_0} \times_U U_{i_1} \times_U U_{i_2}) \to \ldots
$$
其中第一项位于次数 0，各映射取通常的映射。定义
{\it {\v C}ech 上同调群}为
$$
\check{H}^p(\mathcal{U}, \mathcal{F}) =
H^p(\check{\mathcal{C}}^\bullet(\mathcal{U}, \mathcal{F})).
$$
\end{definition}

\noindent
在上述定义中，允许指标 $(i_0, \ldots, i_p)$ 重复至关重要。

\begin{lemma}
\label{etale-cohomology-lemma-cech-presheaves}
记号和假设同定义 \ref{etale-cohomology-definition-cech-complex}。函子
$\check{\mathcal{C}}^\bullet(\mathcal{U}, -)$ 在范畴
$\textit{PAb}(\mathcal{C})$ 上正合。
\end{lemma}

\noindent
换言之，若 $0\to \mathcal{F}_1\to \mathcal{F}_2\to \mathcal{F}_3\to 0$
是阿贝尔群预层的短正合列，则
$$
0 \to \check{\mathcal{C}}^\bullet\left(\mathcal{U}, \mathcal{F}_1\right)
\to\check{\mathcal{C}}^\bullet(\mathcal{U}, \mathcal{F}_2) \to
\check{\mathcal{C}}^\bullet(\mathcal{U}, \mathcal{F}_3)\to 0
$$
是复形的短正合列。

\begin{proof}
这由预层短正合列的定义立即得到。事实上，阿贝尔预层范畴是某个范畴到
$\textit{Ab}$ 的函子范畴，因而自动成为阿贝尔范畴：列
$\mathcal{F}_1\to \mathcal{F}_2\to \mathcal{F}_3$ 在
$\textit{PAb}$ 中正合，当且仅当对所有 $U \in \Ob(\mathcal{C})$，列
$\mathcal{F}_1(U) \to \mathcal{F}_2(U) \to \mathcal{F}_3(U)$ 在
$\textit{Ab}$ 中正合。因此，上述复形在每个次数上不过是短正合列的积。
另见 Cohomology on Sites，引理
\ref{sites-cohomology-lemma-cech-exact-presheaves}。
\end{proof}

\noindent
这说明 $\check{H}^\bullet(\mathcal{U}, -)$ 是 $\delta$-函子。
下面证明它是泛 $\delta$-函子。为此，需要证明它是{\it 可抹的}函子。
我们先回顾 Yoneda 引理。

\begin{lemma}[Yoneda 引理]
\label{etale-cohomology-lemma-yoneda-presheaf}
对任意预层 $\mathcal{F}$（定义在范畴 $\mathcal{C}$ 上）以及
$U \in \Ob(\mathcal{C})$，有函子性同构
$$
\Hom_{\textit{PSh}(\mathcal{C})}(h_U, \mathcal{F}) =
\mathcal{F}(U).
$$
\end{lemma}

\begin{proof}
见 Categories，引理 \ref{categories-lemma-yoneda}。
\end{proof}

\noindent
给定集合 $E$，本节记 $\mathbf{Z}[E]$ 为由 $E$ 生成的自由阿贝尔群。
用公式表示即为 $\mathbf{Z}[E] = \bigoplus_{e \in E} \mathbf{Z}$；也就是说，
$\mathbf{Z}[E]$ 是一个自由 $\mathbf{Z}$-模，其基由 $E$ 的元素组成。
利用这一记号，我们引入集合预层上的自由阿贝尔预层。

\begin{definition}
\label{etale-cohomology-definition-free-abelian-presheaf}
设 $\mathcal{C}$ 是一个范畴。给定集合预层 $\mathcal{G}$，定义
{\it $\mathcal{G}$ 上的自由阿贝尔预层}，记作
$\mathbf{Z}_\mathcal{G}$，规则为
$$
\mathbf{Z}_\mathcal{G}(U)
=
\mathbf{Z}[\mathcal{G}(U)]
$$
其中 $U \in \Ob(\mathcal{C})$，限制映射由 $\mathcal{G}$ 的限制映射诱导。
在特殊情形 $\mathcal{G} = h_U$ 下，简记为
$\mathbf{Z}_U = \mathbf{Z}_{h_U}$。
\end{definition}

\noindent
函子 $\mathcal{G} \mapsto \mathbf{Z}_\mathcal{G}$ 是遗忘函子
$\textit{PAb}(\mathcal{C}) \to \textit{PSh}(\mathcal{C})$ 的左伴随。
因此，对任意预层 $\mathcal{F}$，有典范同构
$$
\Hom_{\textit{PAb}(\mathcal{C})}(\mathbf{Z}_U, \mathcal{F})
=
\Hom_{\textit{PSh}(\mathcal{C})}(h_U, \mathcal{F})
=
\mathcal{F}(U)
$$
最后一个等式由 Yoneda 引理得到。特别地，有以下结果。

\begin{lemma}
\label{etale-cohomology-lemma-cech-complex-describe}
记号和假设同定义 \ref{etale-cohomology-definition-cech-complex}。{\v C}ech 复形
$\check{\mathcal{C}}^\bullet(\mathcal{U}, \mathcal{F})$ 可显式描述如下：
\begin{eqnarray*}
\check{\mathcal{C}}^\bullet(\mathcal{U}, \mathcal{F})
& = &
\left(
\prod_{i_0 \in I}
\Hom_{\textit{PAb}(\mathcal{C})}(\mathbf{Z}_{U_{i_0}}, \mathcal{F}) \to
\prod_{i_0, i_1 \in I}
\Hom_{\textit{PAb}(\mathcal{C})}(
\mathbf{Z}_{U_{i_0} \times_U U_{i_1}}, \mathcal{F}) \to \ldots
\right) \\
& = &
\Hom_{\textit{PAb}(\mathcal{C})}\left(
\left(
\bigoplus_{i_0 \in I} \mathbf{Z}_{U_{i_0}} \leftarrow
\bigoplus_{i_0, i_1 \in I} \mathbf{Z}_{U_{i_0} \times_U U_{i_1}} \leftarrow
\ldots
\right), \mathcal{F}\right)
\end{eqnarray*}
\end{lemma}

\begin{proof}
这由上式得到。见 Cohomology on Sites，引理
\ref{sites-cohomology-lemma-cech-map-into}。
\end{proof}

\noindent
这使我们只需研究最后一个 $\Hom$ 的第一变元中的复形。

\begin{lemma}
\label{etale-cohomology-lemma-exact}
记号和假设同定义 \ref{etale-cohomology-definition-cech-complex}。阿贝尔预层复形
\begin{align*}
\mathbf{Z}_\mathcal{U}^\bullet \quad : \quad
\bigoplus_{i_0 \in I} \mathbf{Z}_{U_{i_0}} \leftarrow
\bigoplus_{i_0, i_1 \in I} \mathbf{Z}_{U_{i_0} \times_U U_{i_1}} \leftarrow
\bigoplus_{i_0, i_1, i_2 \in I}
\mathbf{Z}_{U_{i_0} \times_U U_{i_1} \times_U U_{i_2}} \leftarrow
\ldots
\end{align*}
除次数 $0$ 外，该复形在 $\textit{PAb}(\mathcal{C})$ 中的所有次数上正合。
\end{lemma}

\begin{proof}
对任意 $V\in \Ob(\mathcal{C})$，阿贝尔群复形
$\mathbf{Z}_\mathcal{U}^\bullet(V)$ 为
$$
\begin{matrix}
\mathbf{Z}\left[
\coprod_{i_0\in I} \Mor_\mathcal{C}(V, U_{i_0})\right]
\leftarrow
\mathbf{Z}\left[
\coprod_{i_0, i_1 \in I}
\Mor_\mathcal{C}(V, U_{i_0} \times_U U_{i_1})\right]
\leftarrow \ldots = \\
\bigoplus_{\varphi : V \to U}
\left(
\mathbf{Z}\left[\coprod_{i_0 \in I} \Mor_\varphi(V, U_{i_0})\right]
\leftarrow
\mathbf{Z}\left[\coprod_{i_0, i_1\in I} \Mor_\varphi(V, U_{i_0}) \times
\Mor_\varphi(V, U_{i_1})\right]
\leftarrow
\ldots
\right)
\end{matrix}
$$
其中
$$
\Mor_{\varphi}(V, U_i)
=
\{ V \to U_i \text{，使得 } V \to U_i \to U \text{ 等于 } \varphi \}.
$$
令 $S_\varphi = \coprod_{i\in I} \Mor_\varphi(V, U_i)$，于是
$$
\mathbf{Z}_\mathcal{U}^\bullet(V)
=
\bigoplus_{\varphi : V \to U}
\left(
\mathbf{Z}[S_\varphi] \leftarrow
\mathbf{Z}[S_\varphi \times S_\varphi] \leftarrow
\mathbf{Z}[S_\varphi \times S_\varphi \times S_\varphi] \leftarrow
\ldots \right).
$$
因此，只需证明对每个 $S = S_\varphi$，复形
\begin{align*}
\mathbf{Z}[S] \leftarrow
\mathbf{Z}[S \times S] \leftarrow
\mathbf{Z}[S \times S \times S] \leftarrow \ldots
\end{align*}
在负次数上正合。为此，可以给出一个显式同伦。固定 $s\in S$，并定义
$K: n_{(s_0, \ldots, s_p)} \mapsto n_{(s, s_0,
\ldots, s_p)}.$ 容易验证，$K$ 是下述算子的零同伦：
$$
\delta :
\eta_{(s_0, \ldots, s_p)}
\mapsto
\sum\nolimits_{i = 0}^p (-1)^p \eta_{(s_0, \ldots, \hat s_i, \ldots, s_p)}.
$$
更多细节见 Cohomology on Sites，引理
\ref{sites-cohomology-lemma-homology-complex}。
\end{proof}

\begin{lemma}
\label{etale-cohomology-lemma-hom-injective}
记号和假设同定义 \ref{etale-cohomology-definition-cech-complex}。若 $\mathcal{I}$ 是
$\textit{PAb}(\mathcal{C})$ 的内射对象，则
$\check H^p(\mathcal{U}, \mathcal{I}) = 0$ 对所有 $p > 0$ 成立。
\end{lemma}

\begin{proof}
{\v C}ech 复形是将函子
$\Hom_{\textit{PAb}(\mathcal{C})}(-, \mathcal{I}) $ 作用于复形 $
\mathbf{Z}^\bullet_\mathcal{U} $ 所得，也就是说
$$
\check H^p(\mathcal{U}, \mathcal{I}) = H^p
(\Hom_{\textit{PAb}(\mathcal{C})} (\mathbf{Z}^\bullet_\mathcal{U},
\mathcal{I})).
$$
但我们刚刚看到 $\mathbf{Z}^\bullet_\mathcal{U}$ 在负次数上正合，而函子
$\Hom_{\textit{PAb}(\mathcal{C})}(-,
\mathcal{I})$ 正合，因此 $\Hom_{\textit{PAb}(\mathcal{C})}
(\mathbf{Z}^\bullet_\mathcal{U}, \mathcal{I})$ 在正次数上正合。
\end{proof}

\begin{theorem}
\label{etale-cohomology-theorem-cech-derived}
记号和假设同定义 \ref{etale-cohomology-definition-cech-complex}。在
$\textit{PAb}(\mathcal{C})$ 上，函子 $\check{H}^p(\mathcal{U}, -)$ 是
$\check{H}^0(\mathcal{U}, -)$ 的右导出函子。
\end{theorem}

\begin{proof}
由引理 \ref{etale-cohomology-lemma-hom-injective}，函子
$\check H^p(\mathcal{U}, -)$ 因为可抹而成为泛 $\delta$-函子。
$\check H^0(\mathcal{U}, -)$ 的右导出函子也是如此。由于二者在次数
$0$ 上一致，泛 $\delta$-函子的泛性质说明它们一致。更多细节见
Cohomology on Sites，引理
\ref{sites-cohomology-lemma-cech-cohomology-derived-presheaves}。
\end{proof}

\begin{remark}
\label{etale-cohomology-remark-presheaves-no-topology}
注意，前面所有陈述都只涉及预层，因此尚未使用拓扑。
\end{remark}




\section{{\v C}ech 到上同调谱序列}
\label{etale-cohomology-section-cech-ss}

\noindent
这个谱序列是证明层上同调基本结果的基础工具。

\begin{lemma}
\label{etale-cohomology-lemma-forget-injectives}
遗忘函子 $\textit{Ab}(\mathcal{C})\to \textit{PAb}(\mathcal{C})$
将内射对象映为内射对象。
\end{lemma}

\begin{proof}
这是形式上的：遗忘函子有左伴随，即层化，而层化是正合函子。
更多细节见 Cohomology on Sites，引理
\ref{sites-cohomology-lemma-injective-abelian-sheaf-injective-presheaf}。
\end{proof}

\begin{theorem}
\label{etale-cohomology-theorem-cech-ss}
设 $\mathcal{C}$ 是网站。对任意覆盖
$\mathcal{U} = \{U_i \to U\}_{i \in I}$（它覆盖
$U \in \Ob(\mathcal{C})$）以及任意阿贝尔层 $\mathcal{F}$（在
$\mathcal{C}$ 上），存在谱序列
$$
E_2^{p, q}
=
\check H^p(\mathcal{U}, \underline{H}^q(\mathcal{F}))
\Rightarrow
H^{p+q}(U, \mathcal{F}),
$$
其中 $\underline{H}^q(\mathcal{F})$ 是阿贝尔预层
$V \mapsto H^q(V, \mathcal{F})$。
\end{theorem}

\begin{proof}
选取内射消解 $\mathcal{F}\to \mathcal{I}^\bullet$（位于
$\textit{Ab}(\mathcal{C})$ 中），并考虑双复形
$\check{\mathcal{C}}^\bullet(\mathcal{U}, \mathcal{I}^\bullet)$ 及映射
$$
\xymatrix{
\Gamma(U, I^\bullet) \ar[r] &
\check{\mathcal{C}}^\bullet(\mathcal{U}, \mathcal{I}^\bullet) \\
& \check{\mathcal{C}}^\bullet(\mathcal{U}, \mathcal{F}) \ar[u]
}
$$
这里，水平映射是到最左列的自然映射
$\Gamma(U, I^\bullet) \to
\check{\mathcal{C}}^0(\mathcal{U}, \mathcal{I}^\bullet)$，垂直映射由
$\mathcal{F}\to \mathcal{I}^0$ 诱导，并落在最下行。根据假设，
$\mathcal{I}^\bullet$ 是 $\textit{Ab}(\mathcal{C})$ 中的内射对象复形，
故由引理 \ref{etale-cohomology-lemma-forget-injectives}，它也是
$\textit{PAb}(\mathcal{C})$ 中的内射对象复形。因此，双复形的各行在
正次数上正合（引理 \ref{etale-cohomology-lemma-hom-injective}），并且
$\check{\mathcal{C}}^0(\mathcal{U}, \mathcal{I}^\bullet)
\to \check{\mathcal{C}}^1(\mathcal{U}, \mathcal{I}^\bullet)$ 的核等于
$\Gamma(U, \mathcal{I}^\bullet)$，因为 $\mathcal{I}^\bullet$
是层的复形。特别地，总复形的上同调就是全局截面函子
$H^0(U, \mathcal{F})$ 的标准上同调。

\medskip\noindent
在垂直方向上，第 $q$ 个上同调群（取自第 $p$ 列）为
$$
\prod_{i_0, \ldots, i_p}
H^q(U_{i_0} \times_U \ldots \times_U U_{i_p}, \mathcal{F})
=
\prod_{i_0, \ldots, i_p}
\underline{H}^q(\mathcal{F})(U_{i_0} \times_U \ldots \times_U U_{i_p})
$$
它位于 $E_1^{p, q}$ 项。因此，这是标准的双复形谱序列，且
$E_2$-页如所述。更多细节见 Cohomology on Sites，引理
\ref{sites-cohomology-lemma-cech-spectral-sequence}。
\end{proof}

\begin{remark}
\label{etale-cohomology-remark-grothendieck-ss}
这是下述函子复合的 Grothendieck 谱序列：
$$
\textit{Ab}(\mathcal{C}) \longrightarrow
\textit{PAb}(\mathcal{C}) \xrightarrow{\check H^0} \textit{Ab}.
$$
\end{remark}








\section{概形的大网站与小网站}
\label{etale-cohomology-section-big-small}

\noindent
设 $S$ 是概形。设 $\tau$ 是我们将讨论的拓扑之一。因此
$\tau \in \{fppf, syntomic, smooth, \etale, Zariski\}$。
当然，如果只关心 \'etale 拓扑，则可始终假设 $\tau = \etale$。
此外，从第 \ref{etale-cohomology-section-etale-site} 节起，我们将更详细地讨论
\'etale 态射、\'etale 覆盖和 \'etale 网站。为了继续讨论拟凝聚层的
上同调，引入大 $\tau$-网站很方便；当
$\tau \in \{\etale, Zariski\}$ 时，还要引入小 $\tau$-网站，即 $S$ 的相应小网站。
为此，先引入 $\tau$-覆盖的概念。

\begin{definition}
\label{etale-cohomology-definition-tau-covering}
（见 Topologies，定义
\ref{topologies-definition-fppf-covering},
\ref{topologies-definition-syntomic-covering},
\ref{topologies-definition-smooth-covering},
\ref{topologies-definition-etale-covering}，和
\ref{topologies-definition-zariski-covering}。）
设 $\tau \in \{fppf, syntomic, smooth, \etale, Zariski\}$。
称具有共同靶的概形态射族 $\{f_i : T_i \to T\}_{i \in I}$ 是
{\it $\tau$-覆盖}，当且仅当每个 $f_i$ 分别是有限表示平坦态射、
syntomic 态射、光滑态射、\'etale 态射或开浸入，并且
$\bigcup f_i(T_i) = T$。
\end{definition}

\noindent
所有 $\tau$-覆盖所成的类满足定义 \ref{etale-cohomology-definition-site}（我们对网站的定义）
中的公理 (1)、(2) 和 (3)，见 Topologies，引理
\ref{topologies-lemma-fppf},
\ref{topologies-lemma-syntomic},
\ref{topologies-lemma-smooth},
\ref{topologies-lemma-etale}，和
\ref{topologies-lemma-zariski}。

\medskip\noindent
下面引入我们将使用的网站。与 \cite{SGA4} 中的做法不同，我们不想选取一个
宇宙，而是选取一个“部分宇宙”（即 Sets，第
\ref{sets-section-sets-hierarchy} 节所述的一个适当大的集合），并考虑包含在
这个集合中的所有概形。当然，我们确保所选基概形 $S$ 包含在部分宇宙中。
选定底层范畴后，再选取一个适当大的 $\tau$-覆盖集合，使之成为网站。
细节见概形拓扑一章；这些选择有很大自由度，但最终的具体选择不会影响
$S$ 的 \'etale（或其他）上同调（正如在 \cite{SGA4} 中，宇宙的具体选择
最终并不重要）。此外，本材料的写法也允许愿意使用强不可达基数（即宇宙）
的读者以之替代部分宇宙。

\begin{definition}
\label{etale-cohomology-definition-tau-site}
设 $S$ 是概形。
设 $\tau \in \{fppf, syntomic, smooth, \etale, \linebreak[0] Zariski\}$。
\begin{enumerate}
\item {\it 大 $\tau$-网站（基为 $S$）}是按上述方式构造的任一网站
$(\Sch/S)_\tau$；更详细的构造见 Topologies，定义
\ref{topologies-definition-big-small-fppf},
\ref{topologies-definition-big-small-syntomic},
\ref{topologies-definition-big-small-smooth},
\ref{topologies-definition-big-small-etale}，和
\ref{topologies-definition-big-small-Zariski}.
\item 若 $\tau \in \{\etale, Zariski\}$，则{\it 小 $\tau$-网站（基为 $S$）}
是全子范畴 $S_\tau$，位于 $(\Sch/S)_\tau$ 中；其对象是概形 $T$（位于 $S$ 上），
其结构态射 $T \to S$ 分别为 \'etale 态射或开浸入。$S_\tau$ 中的覆盖是族
$\{U_i \to U\}$，它在 $(\Sch/S)_\tau$ 中构成覆盖，其中 $U$ 是 $S_\tau$ 的对象。
\end{enumerate}
\end{definition}

\noindent
网站 $(\Sch/S)_\tau$ 的底层范畴具有合理的“封闭性”。例如，给定其中的概形
$T$，$T$ 的任意局部闭子概形都同构于 $(\Sch/S)_\tau$ 的某个对象。其他这类
封闭性包括：对概形的纤维积封闭；可取可数不交并；可取给定概形上的有限型
概形；给定仿射概形 $\Spec(R)$，可对 $R$ 作完备化、局部化或关于某个理想取商，
而仍留在该范畴中；等等。另一方面，例如对 $(\Sch/S)_\tau$ 中的概形取任意
不交并，则可能离开这个范畴。还要注意，给定对象 $T$（位于 $(\Sch/S)_\tau$ 中），可能
存在 $\tau$-覆盖 $\{T_i \to T\}_{i \in I}$（如定义
\ref{etale-cohomology-definition-tau-covering} 所述），但它不是 $(\Sch/S)_\tau$ 中的覆盖；例如，
概形 $T_i$ 可能不是范畴 $(\Sch/S)_\tau$ 的对象。不过，我们对网站
$(\Sch/S)_\tau$ 的选取保证，总存在覆盖 $\{U_j \to T\}_{j \in J}$
（属于 $(\Sch/S)_\tau$），它加细覆盖 $\{T_i \to T\}_{i \in I}$，见
Topologies，引理
\ref{topologies-lemma-fppf-induced},
\ref{topologies-lemma-syntomic-induced},
\ref{topologies-lemma-smooth-induced},
\ref{topologies-lemma-etale-induced}，和
\ref{topologies-lemma-zariski-induced}.
本章大多不再理会这些问题。

\medskip\noindent
若 $\mathcal{F}$ 是 $(\Sch/S)_\tau$ 或 $S_\tau$ 上的层，则以
$$
H^p_\tau(U, \mathcal{F}), \text{ 特别地 }
H^p_\tau(S, \mathcal{F})
$$
表示 $\mathcal{F}$ 在网站对象 $U$ 上的上同调群，见第
\ref{etale-cohomology-section-cohomology} 节。因此有
$H^p_{fppf}(S, \mathcal{F})$,
$H^p_{syntomic}(S, \mathcal{F})$,
$H^p_{smooth}(S, \mathcal{F})$,
$H^p_\etale(S, \mathcal{F})$，和
$H^p_{Zar}(S, \mathcal{F})$。后两种记号可能有歧义，因为它们既可能指大
\'etale 或 Zariski 网站，也可能指相应的小网站。不过，由下述引理可知，
这种歧义无害。

\begin{lemma}
\label{etale-cohomology-lemma-compare-cohomology-big-small}
设 $\tau \in \{\etale, Zariski\}$。若 $\mathcal{F}$ 是定义在
$(\Sch/S)_\tau$ 上的阿贝尔层，则 $\mathcal{F}$ 在 $S$ 上的上同调群与
$\mathcal{F}|_{S_\tau}$ 在 $S$ 上的上同调群相同。
\end{lemma}

\begin{proof}
由 Topologies，引理 \ref{topologies-lemma-at-the-bottom} 和
\ref{topologies-lemma-at-the-bottom-etale}
可知，函子 $S_\tau \to (\Sch/S)_\tau$ 满足 Sites，引理
\ref{sites-lemma-bigger-site} 的假设。因此，本引理由 Cohomology on Sites，
引理 \ref{sites-cohomology-lemma-cohomology-bigger-site} 得出。
\end{proof}

\noindent
$S$ 的大或小 \'etale 网站上的层范畴，只依赖于 $(\Sch/S)_\etale$ 或
$S_\etale$ 中由仿射对象组成的全子范畴；并且只需考虑这些对象之间的标准
\'etale 覆盖（定义见下文）。由此得到网站 $(\textit{Aff}/S)_\etale$ 和
$S_{affine, \etale}$，见 Topologies，定义
\ref{topologies-definition-big-small-etale}。比较结果见 Topologies，引理
\ref{topologies-lemma-affine-big-site-etale} 和
\ref{topologies-lemma-alternative}。下面定义本章将考虑的若干拓扑中的标准覆盖。

\begin{definition}
\label{etale-cohomology-definition-standard-tau}
（见 Topologies，定义
\ref{topologies-definition-standard-fppf},
\ref{topologies-definition-standard-syntomic},
\ref{topologies-definition-standard-smooth},
\ref{topologies-definition-standard-etale}，和
\ref{topologies-definition-standard-Zariski}.)
设 $\tau \in \{fppf, syntomic, smooth, \etale, Zariski\}$，并设 $T$ 是仿射概形。
一个{\it 标准 $\tau$-覆盖（靶为 $T$）}是一个族
$\{f_j : U_j \to T\}_{j = 1, \ldots, m}$，其中每个 $U_j$ 都是仿射的，
每个 $f_j$ 分别是有限表示平坦态射、标准 syntomic 态射、标准光滑态射、
\'etale 态射或 $T$ 中标准主开集的浸入，并且 $T = \bigcup f_j(U_j)$。
\end{definition}

\begin{lemma}
\label{etale-cohomology-lemma-tau-affine}
设 $\tau \in \{fppf, syntomic, smooth, \etale, Zariski\}$。
仿射概形的任意 $\tau$-覆盖都可由某个标准 $\tau$-覆盖加细。
\end{lemma}

\begin{proof}
见 Topologies，引理
\ref{topologies-lemma-fppf-affine},
\ref{topologies-lemma-syntomic-affine},
\ref{topologies-lemma-smooth-affine},
\ref{topologies-lemma-etale-affine}，和
\ref{topologies-lemma-zariski-affine}.
\end{proof}

\noindent
为完整起见，我们陈述并证明这里所考虑的上同调群不依赖于部分宇宙的选取。
下文将在引理 \ref{etale-cohomology-lemma-etale-topos-independent-partial-universe} 中证明小
\'etale 拓扑斯的不变性。本引理使用的记号和术语见 Topologies，第
\ref{topologies-section-change-alpha} 节。

\begin{lemma}
\label{etale-cohomology-lemma-cohomology-enlarge-partial-universe}
设 $\tau \in \{fppf, syntomic, smooth, \etale, Zariski\}$。设 $S$ 是概形。
设 $(\Sch/S)_\tau$ 和 $(\Sch'/S)_\tau$ 是两个大 $\tau$-网站（基为 $S$），
并假设前者包含于后者。此时：
\begin{enumerate}
\item 对任意阿贝尔层 $\mathcal{F}'$（定义在 $(\Sch'/S)_\tau$ 上），以及
任意对象 $U$（属于 $(\Sch/S)_\tau$），有
$$
H^p_\tau(U, \mathcal{F}'|_{(\Sch/S)_\tau}) =
H^p_\tau(U, \mathcal{F}')
$$
换言之，在较大网站中计算的 $\mathcal{F}'$ 在 $U$ 上的上同调，与把
$\mathcal{F}'$ 限制到较小网站后在 $U$ 上的上同调相同。
\item 对任意阿贝尔层 $\mathcal{F}$（在 $(\Sch/S)_\tau$ 上），存在阿贝尔层
$\mathcal{F}'$（在 $(\Sch/S)_\tau'$ 上），其到 $(\Sch/S)_\tau$ 的限制
同构于 $\mathcal{F}$。
\end{enumerate}
\end{lemma}

\begin{proof}
由 Topologies，引理 \ref{topologies-lemma-change-alpha} 可知，包含函子
$(\Sch/S)_\tau \to (\Sch'/S)_\tau$ 满足 Sites，引理
\ref{sites-lemma-bigger-site} 的假设。这推出 (2)，而 (1) 由
Cohomology on Sites，引理 \ref{sites-cohomology-lemma-cohomology-bigger-site}
得出。
\end{proof}




\section{\'Etale 拓扑斯}
\label{etale-cohomology-section-etale-topos}

\noindent
{\it 拓扑斯}是某个网站上的集合层所成的范畴，见 Sites，定义
\ref{sites-definition-topos}。因此，通常用“概形的 \'etale 拓扑斯”这一说法
指概形的小 \'etale 网站上的层范畴。正式定义如下。

\begin{definition}
\label{etale-cohomology-definition-etale-topos}
设 $S$ 是概形。
\begin{enumerate}
\item $S$ 的{\it \'etale 拓扑斯}，或称{\it 小 \'etale 拓扑斯}，是集合层范畴
$\Sh(S_\etale)$，这些层位于 $S$ 的小 \'etale 网站上。
\item $S$ 的{\it Zariski 拓扑斯}，或称{\it 小 Zariski 拓扑斯}，是集合层范畴
$\Sh(S_{Zar})$，这些层位于 $S$ 的小 Zariski 网站上。
\item 对 $\tau \in \{fppf, syntomic, smooth, \etale, Zariski\}$，
{\it 大 $\tau$-拓扑斯}是某个大 $\tau$-拓扑斯（基为 $S$）上的集合层范畴。
\end{enumerate}
\end{definition}

\noindent
注意，$S$ 的小 Zariski 拓扑斯就是 $S$ 的底层拓扑空间上的集合层范畴，见
Topologies，引理 \ref{topologies-lemma-Zariski-usual}。小 \'etale 拓扑斯
不依赖于构造小 \'etale 网站时所作的选择，但一般而言，大拓扑斯确实依赖于
这些选择。

\medskip\noindent
事实上，大或小 \'etale 拓扑斯只依赖于 $(\Sch/S)_\etale$ 或 $S_\etale$
中由仿射对象组成的全子范畴，见 Topologies，引理
\ref{topologies-lemma-affine-big-site-etale} 和
\ref{topologies-lemma-alternative}。例如，在下述引理的证明中将用到这一点。

\begin{lemma}
\label{etale-cohomology-lemma-etale-topos-independent-partial-universe}
设 $S$ 是概形。$S$ 的 \'etale 拓扑斯不依赖于定义
\ref{etale-cohomology-definition-tau-site} 中小 \'etale 网站的构造（在典范等价意义下）。
\end{lemma}

\begin{proof}
需要证明：给定两个大 \'etale 网站 $\Sch_\etale$ 和 $\Sch_\etale'$，
它们都包含 $S$；采用显然的记号，有
$\Sh(S_\etale) \cong \Sh(S_\etale')$。由 Topologies，引理
\ref{topologies-lemma-contained-in}，可假设
$\Sch_\etale \subset \Sch_\etale'$。由 Sets，引理
\ref{sets-lemma-what-is-in-it}，$S$ 上任意 \'etale 的仿射概形都同构于
$\Sch_\etale$ 和 $\Sch_\etale'$ 二者中的对象。因此，诱导函子
$S_{affine, \etale} \to S_{affine, \etale}'$
是等价。此外，显然该函子及其某个拟逆都把标准 \'etale 覆盖变为标准
\'etale 覆盖。因此，结论由 Topologies，引理
\ref{topologies-lemma-alternative} 得出。
\end{proof}





\section{拟凝聚层的上同调}
\label{etale-cohomology-section-cohomology-quasi-coherent}
%9.22.09

\noindent
先从一个简单引理开始（它在比所述情形更一般的范围内成立）。该引理说明，
标准覆盖的 {\v C}ech 复形等于形如 $\{\Spec(B) \to \Spec(A)\}$ 的 fpqc
覆盖的 {\v C}ech 复形，其中 $A \to B$ 忠实平坦。

\begin{lemma}
\label{etale-cohomology-lemma-cech-complex}
设 $\tau \in \{fppf, syntomic, smooth, \etale, Zariski\}$。设 $S$ 是概形。
设 $\mathcal{F}$ 是 $(\Sch/S)_\tau$ 上的阿贝尔层；也可设它是
$S_\tau$ 上的阿贝尔层（当 $\tau = \etale$ 时）。再设
$\mathcal{U} = \{U_i \to U\}_{i \in I}$
是该网站的标准 $\tau$-覆盖。令 $V = \coprod_{i \in I} U_i$。则：
\begin{enumerate}
\item $V$ 是仿射概形；
\item $\mathcal{V} = \{V \to U\}$ 是 fpqc 覆盖；它也是 $\tau$-覆盖，
除非 $\tau = Zariski$；
\item {\v C}ech 复形 $\check{\mathcal{C}}^\bullet (\mathcal{U}, \mathcal{F})$
与 $\check{\mathcal{C}}^\bullet (\mathcal{V}, \mathcal{F})$ 相同。
\end{enumerate}
\end{lemma}

\begin{proof}
标准 $\tau$-覆盖的定义见 Topologies，定义
\ref{topologies-definition-standard-Zariski},
\ref{topologies-definition-standard-etale},
\ref{topologies-definition-standard-smooth},
\ref{topologies-definition-standard-syntomic}，和
\ref{topologies-definition-standard-fppf}.
由定义，每个概形 $U_i$ 都是仿射的，且 $I$ 是有限集。因此 $V$ 是仿射概形。
显然 $V \to U$ 平坦且满射，故 $\mathcal{V}$ 是 fpqc 覆盖，见例
\ref{etale-cohomology-example-fpqc-coverings}。除 Zariski 情形外，覆盖 $\mathcal{V}$ 也是
$\tau$-覆盖，见 Topologies，定义
\ref{topologies-definition-etale-covering},
\ref{topologies-definition-smooth-covering},
\ref{topologies-definition-syntomic-covering}，和
\ref{topologies-definition-fppf-covering}.

\medskip\noindent
注意，$\mathcal{U}$ 是 $\mathcal{V}$ 的加细，因而有 {\v C}ech 复形的映射
$\check{\mathcal{C}}^\bullet (\mathcal{V}, \mathcal{F}) \to
\check{\mathcal{C}}^\bullet (\mathcal{U}, \mathcal{F})$，见
Cohomology on Sites，等式
(\ref{sites-cohomology-equation-map-cech-complexes})。接着注意，若
$T = \coprod_{j \in J} T_j$ 是 $\mathcal{F}$ 所定义的网站中若干概形的不交并，
则具有共同靶的态射族 $\{T_j \to T\}_{j \in J}$ 是 Zariski 覆盖，故
\begin{equation}
\label{etale-cohomology-equation-sheaf-coprod}
\mathcal{F}(T) =
\mathcal{F}(\coprod\nolimits_{j \in J} T_j) =
\prod\nolimits_{j \in J} \mathcal{F}(T_j)
\end{equation}
由 $\mathcal{F}$ 的层条件成立。这说明上面的 {\v C}ech 复形映射在每一次数上
都是同构，因为作为概形有
$$
V \times_U \ldots \times_U V
=
\coprod\nolimits_{i_0, \ldots i_p} U_{i_0} \times_U \ldots \times_U U_{i_p}
$$
\end{proof}

\noindent
注意，等式 (\ref{etale-cohomology-equation-sheaf-coprod}) 对一般预层不成立。即使对层，它也不在
任意网站上都成立，因为余积未必产生覆盖，也未必是不交的。不过，对所有常用
网站（至少对我们将研究的全部网站），该等式成立。

\begin{remark}
\label{etale-cohomology-remark-refinement}
在引理 \ref{etale-cohomology-lemma-cech-complex} 的陈述中，覆盖 $\mathcal{U}$ 是
$\mathcal{V}$ 的加细，反向却不成立。形如 $\{V \to U\}$ 的覆盖并不构成
$U$ 的所有覆盖所成范畴的始子范畴。不过，仍可借助这些覆盖计算 {\v C}ech
上同调 $\check H^n(U, \mathcal{F})$；这里它定义为：在覆盖 $\mathcal{U}$
（其靶为 $U$）所成范畴的反范畴上，对 $\mathcal{F}$ 关于 $\mathcal{U}$ 的
{\v C}ech 上同调群取余极限，而计算可借助覆盖 $\{V \to U\}$。若以后需要，
我们将陈述一个精确引理（它只对层成立）并补在此处。
\end{remark}

\begin{lemma}[上同调的局部性]
\label{etale-cohomology-lemma-locality-cohomology}
设 $\mathcal{C}$ 是网站，$\mathcal{F}$ 是 $\mathcal{C}$ 上的阿贝尔层，
$U$ 是 $\mathcal{C}$ 的对象，$p > 0$ 是整数，并且 $\xi \in
H^p(U, \mathcal{F})$。则存在覆盖
$\mathcal{U} = \{U_i \to U\}_{i \in I}$；它是 $U$ 在 $\mathcal{C}$ 中的
覆盖，并且 $\xi |_{U_i} = 0$ 对所有 $i \in I$ 都成立。
\end{lemma}

\begin{proof}
选取内射消解 $\mathcal{F} \to \mathcal{I}^\bullet$。于是，$\xi$ 由上闭链
$\tilde{\xi} \in \mathcal{I}^p(U)$ 表示，并且
$d^p(\tilde{\xi}) = 0$。根据假设，序列
$\mathcal{I}^{p - 1} \to \mathcal{I}^p \to \mathcal{I}^{p + 1}$ 在
$\textit{Ab}(\mathcal{C})$ 中正合。这意味着存在覆盖
$\mathcal{U} = \{U_i \to U\}_{i \in I}$，使得
$\tilde{\xi}|_{U_i} = d^{p - 1}(\xi_i)$，其中
$\xi_i \in \mathcal{I}^{p-1}(U_i)$。上同调类 $\xi|_{U_i}$ 由上闭链
$\tilde{\xi}|_{U_i}$ 表示，而后者是上边界，故该上同调类为零。更多细节见
Cohomology on Sites，引理
\ref{sites-cohomology-lemma-kill-cohomology-class-on-covering}。
\end{proof}

\begin{theorem}
\label{etale-cohomology-theorem-zariski-fpqc-quasi-coherent}
设 $S$ 是概形，$\mathcal{F}$ 是拟凝聚 $\mathcal{O}_S$-模。设
$\mathcal{C}$ 或者是 $(\Sch/S)_\tau$，其中
$\tau \in \{fppf, \allowbreak syntomic, \allowbreak smooth, \allowbreak \etale, \allowbreak Zariski\}$；或者是
$S_\etale$。则
$$
H^p(S, \mathcal{F}) = H^p_\tau(S, \mathcal{F}^a)
$$
对所有 $p \geq 0$ 都成立，其中：
\begin{enumerate}
\item 左端表示层 $\mathcal{F}$ 在概形 $S$ 的底层拓扑空间上的通常上同调；
\item 右端表示阿贝尔层 $\mathcal{F}^a$（见命题
\ref{etale-cohomology-proposition-quasi-coherent-sheaf-fpqc}）在网站 $\mathcal{C}$ 上的上同调。
\end{enumerate}
\end{theorem}

\begin{proof}
我们将证明
$H^p(U, f^*\mathcal{F}) = H^p_\tau(U, \mathcal{F}^a)$；这里
$f : U \to S$ 是网站 $\mathcal{C}$ 的任意对象。由层性质，结论在
$p = 0$ 时成立。

\medskip\noindent
假设 $U$ 是仿射的。我们要证明
$H^p_\tau(U, \mathcal{F}^a) = 0$ 对所有 $p > 0$ 都成立。对 $p$ 作归纳。
\begin{enumerate}
\item[p = 1]
取 $\xi \in H^1_\tau(U, \mathcal{F}^a)$。由引理
\ref{etale-cohomology-lemma-locality-cohomology}，存在 fpqc 覆盖
$\mathcal{U} = \{U_i \to U\}_{i \in I}$，使得 $\xi|_{U_i} = 0$ 对所有
$i \in I$ 都成立。加细 $\mathcal{U}$ 后，可假设 $\mathcal{U}$ 是标准
$\tau$-覆盖。应用定理 \ref{etale-cohomology-theorem-cech-ss} 的谱序列可知，$\xi$ 来自上同调类
$\check \xi \in \check H^1(\mathcal{U}, \mathcal{F}^a)$。考虑覆盖
$\mathcal{V} = \{\coprod_{i\in I} U_i \to U\}$。由引理
\ref{etale-cohomology-lemma-cech-complex}，
$\check H^\bullet(\mathcal{U}, \mathcal{F}^a) =
\check H^\bullet(\mathcal{V}, \mathcal{F}^a)$.
另一方面，因为 $\mathcal{V}$ 是形如 $\{\Spec(B) \to \Spec(A)\}$ 的覆盖，
且 $f^*\mathcal{F} = \widetilde{M}$；这里存在 $A$-模 $M$。所以 {\v C}ech 复形
$\check{\mathcal{C}}^\bullet(\mathcal{V}, \mathcal{F})$
就是复形 $(B/A)_\bullet \otimes_A M$。由引理
\ref{etale-cohomology-lemma-descent-modules}，有
$H^p((B/A)_\bullet \otimes_A M) = 0$（当 $p > 0$ 时），故
$\check \xi = 0$，进而
$\xi = 0$。
\item[p > 1]
取 $\xi \in H^p_\tau(U, \mathcal{F}^a)$。由引理
\ref{etale-cohomology-lemma-locality-cohomology}，存在 fpqc 覆盖
$\mathcal{U} = \{U_i \to U\}_{i \in I}$，使得 $\xi|_{U_i} = 0$ 对所有
$i \in I$ 都成立。加细 $\mathcal{U}$ 后，可假设 $\mathcal{U}$ 是标准
$\tau$-覆盖。应用定理 \ref{etale-cohomology-theorem-cech-ss} 的谱序列。注意，各个交
$U_{i_0} \times_U \ldots \times_U U_{i_p}$ 都是仿射的，故由归纳假设，上同调群
$$
E_2^{p, q} = \check H^p(\mathcal{U}, \underline{H}^q(\mathcal{F}^a))
$$
对所有 $0 < q < p$ 都消失。由此可见，$\xi$ 必来自某个
$\check \xi \in \check H^p(\mathcal{U}, \mathcal{F}^a)$。把
$\mathcal{U}$ 换成只含一个态射的覆盖 $\mathcal{V}$，并再次使用引理
\ref{etale-cohomology-lemma-descent-modules}，便知 {\v C}ech 上同调类 $\check \xi$ 必为零，
故 $\xi = 0$。
\end{enumerate}
接着假设 $U$ 是分离的。选取仿射开覆盖
$U = \bigcup_{i \in I} U_i$（它覆盖 $U$）。于是族
$\mathcal{U} = \{U_i \to U\}_{i \in I}$ 是 fpqc 覆盖；并且所有交
$U_{i_0} \times_U \ldots \times_U U_{i_p}$ 都是仿射的，因为 $U$ 分离。因此，定理
\ref{etale-cohomology-theorem-cech-ss} 的谱序列除第零行外各行均为零。所以
$$
H^p_\tau(U, \mathcal{F}^a) =
\check H^p(\mathcal{U}, \mathcal{F}^a) =
\check H^p(\mathcal{U}, \mathcal{F}) = H^p(U, \mathcal{F})
$$
；最后一个等式来自标准概形理论，见 Cohomology of Schemes，引理
\ref{coherent-lemma-cech-cohomology-quasi-coherent}。

\medskip\noindent
一般情形较为技术性；若要延伸这里给出的证明，还需讨论谱序列之间的映射，
故我们不作处理。它由 Descent，命题
\ref{descent-proposition-same-cohomology-quasi-coherent}（其证明采用稍有不同的
方法）与 Cohomology on Sites，引理
\ref{sites-cohomology-lemma-cohomology-of-open} 结合得出。
\end{proof}

\begin{remark}
\label{etale-cohomology-remark-right-derived-global-sections}
对定理 \ref{etale-cohomology-theorem-zariski-fpqc-quasi-coherent} 作一点说明。由于 $S$ 是范畴
$\mathcal{C}$ 的终对象，右端的上同调群就是整体截面函子的右导出函子。事实上，
证明表明
$H^p(U, f^*\mathcal{F}) = H^p_\tau(U, \mathcal{F}^a)$
对任意对象 $f : U \to S$（属于网站 $\mathcal{C}$）都成立。
\end{remark}





\section{层的例子}
\label{etale-cohomology-section-examples-sheaves}

\noindent
设 $S$ 和 $\tau$ 如第 \ref{etale-cohomology-section-big-small} 节所述。前面已经看到，任意可表
预层都是 $(\Sch/S)_\tau$ 或 $S_\tau$ 上的层，见引理
\ref{etale-cohomology-lemma-representable-sheaf-fpqc} 和注
\ref{etale-cohomology-remark-fpqc-finest}。下面给出一些特殊情形。

\begin{definition}
\label{etale-cohomology-definition-additive-sheaf}
在第 \ref{etale-cohomology-section-big-small} 节的任一网站 $(\Sch/S)_\tau$ 或 $S_\tau$ 上：
\begin{enumerate}
\item 层 $T \mapsto \Gamma(T, \mathcal{O}_T)$ 记作 $\mathcal{O}_S$ 或
$\mathbf{G}_a$；若要标明基概形，则记作 $\mathbf{G}_{a, S}$。
\item 类似地，层 $T \mapsto \Gamma(T, \mathcal{O}^*_T)$ 记作
$\mathcal{O}_S^*$ 或 $\mathbf{G}_m$；若要标明基概形，则记作
$\mathbf{G}_{m, S}$。
\item 任意网站上的{\it 常值层} $\underline{\mathbf{Z}/n\mathbf{Z}}$，是常值预层
$U \mapsto \mathbf{Z}/n\mathbf{Z}$ 的层化。
\end{enumerate}
\end{definition}

\noindent
例如，由定理 \ref{etale-cohomology-theorem-quasi-coherent} 可知第一个是层。第二个是第一个的
子预层，而且容易看出它本身也是层。第三个按定义即为层。注意，这些层都可表。
前两个分别由概形 $\mathbf{G}_{a, S}$ 和 $\mathbf{G}_{m, S}$ 表示，见
Groupoids，第 \ref{groupoids-section-group-schemes} 节。第三个由有限 \'etale
群概形 $\mathbf{Z}/n\mathbf{Z}_S$ 表示，它有时记作
$(\mathbf{Z}/n\mathbf{Z})_S$；这就是 $n$ 个 $S$ 的副本，并赋予 $S$ 上显然的
群概形结构，见 Groupoids，例 \ref{groupoids-example-constant-group} 及下述注。

\begin{remark}
\label{etale-cohomology-remark-constant-locally-constant-maps}
设 $G$ 是抽象群。在第 \ref{etale-cohomology-section-big-small} 节的任一网站
$(\Sch/S)_\tau$ 或 $S_\tau$ 上，层化 $\underline{G}$（即与 $G$ 相伴的
常值预层在该网站的{\it Zariski 拓扑}中的层化）已经给出
$$
\Gamma(U, \underline{G}) =
\{\text{Zariski 局部常值映射 }U \to G\}
$$
。由 Groupoids，例 \ref{groupoids-example-constant-group}，这个 Zariski 层
由群概形 $G_S$ 表示。由引理 \ref{etale-cohomology-lemma-representable-sheaf-fpqc}，任意可表预层
也满足 $\tau$-拓扑的层条件，故上述 Zariski 层化 $\underline{G}$ 同时也是
$\tau$-层化。
\end{remark}

\begin{definition}
\label{etale-cohomology-definition-structure-sheaf}
设 $S$ 是概形。$S$ 的{\it 结构层}是环层 $\mathcal{O}_S$，它位于上述任一网站
$S_{Zar}$、$S_\etale$ 或 $(\Sch/S)_\tau$ 上。
\end{definition}

\noindent
若所使用的网站可能产生混淆，我们就用下标加以标明。例如，可用
$\mathcal{O}_{S_\etale}$ 强调我们是在 $S$ 的小 \'etale 网站上工作。

\begin{remark}
\label{etale-cohomology-remark-special-case-fpqc-cohomology-quasi-coherent}
用上面引入的术语，定理 \ref{etale-cohomology-theorem-zariski-fpqc-quasi-coherent} 的一个
特殊情形是
$$
H_{fppf}^p(X, \mathbf{G}_a) =
H_\etale^p(X, \mathbf{G}_a) =
H_{Zar}^p(X, \mathbf{G}_a) =
H^p(X, \mathcal{O}_X)
$$
，对所有 $p \geq 0$ 成立。此外，可以用记号
$H^p_{fppf}(X, \mathcal{O}_X)$ 表示结构层在 $X$ 的大 fppf 网站上的上同调。
\end{remark}




\section{Picard 群}
\label{etale-cohomology-section-picard-groups}

\noindent
下述定理有时称为“Hilbert 定理 90”。

\begin{theorem}
\label{etale-cohomology-theorem-picard-group}
对任意概形 $X$，有典范等同
\begin{align*}
H_{fppf}^1(X, \mathbf{G}_m) & = H^1_{syntomic}(X, \mathbf{G}_m) \\
& = H^1_{smooth}(X, \mathbf{G}_m) \\
& = H_\etale^1(X, \mathbf{G}_m) \\
& = H^1_{Zar}(X, \mathbf{G}_m) \\
& = \Pic(X) \\
& = H^1(X, \mathcal{O}_X^*)
\end{align*}
\end{theorem}

\begin{proof}
设 $\tau$ 是第 \ref{etale-cohomology-section-big-small} 节所考虑的拓扑之一。由
Cohomology on Sites，引理
\ref{sites-cohomology-lemma-h1-invertible} 可知
$H^1_\tau(X, \mathbf{G}_m) =
H^1_\tau(X, \mathcal{O}_\tau^*) =
\Pic(\mathcal{O}_\tau)$
，其中 $\mathcal{O}_\tau$ 是网站 $(\Sch/X)_\tau$ 的结构层。于是，可逆
$\mathcal{O}_\tau$-模是拟凝聚 $\mathcal{O}_\tau$-模。由定理
\ref{etale-cohomology-theorem-quasi-coherent}，或更精确的 Descent，命题
\ref{descent-proposition-equivalence-quasi-coherent}，可知
$\Pic(\mathcal{O}_\tau) = \Pic(X)$。最后一个等式以同样方式证明。
\end{proof}






\section{\'Etale 网站}
\label{etale-cohomology-section-etale-site}

\noindent
从这里开始，我们更详细地研究概形的 \'etale 网站。第一步是简要讨论
\'etale 态射的概念。





\section{\'Etale 态射}
\label{etale-cohomology-section-etale-morphism}

\noindent
正式定义见 Morphisms，第 \ref{morphisms-section-etale} 节；\'etale 态射的
若干有趣性质综述见 \'Etale Morphisms，诸节
\ref{etale-section-etale-morphisms},
\ref{etale-section-structure-etale-map},
\ref{etale-section-etale-smooth},
\ref{etale-section-topological-etale},
\ref{etale-section-functorial-etale}，和
\ref{etale-section-properties-permanence}。

\medskip\noindent
回忆：代数 $A$（在代数闭域 $k$ 上）称为{\it 光滑的}，如果它是有限型的，且微分模
$\Omega_{A/k}$ 是秩等于维数的有限局部自由模。概形 $X$（在 $k$ 上）称为在 $k$
上{\it 光滑}，如果它局部为有限型，并且每个仿射开集都是某个光滑 $k$-代数的
谱。若 $k$ 不是代数闭的，则称 $k$-代数 $A$ 是光滑 $k$-代数，如果
$A \otimes_k \overline{k}$ 是光滑 $\overline{k}$-代数。环映射 $A \to B$
称为光滑的，如果它平坦、有限表示，并且对所有素理想
$\mathfrak p \subset A$，纤维环 $\kappa(\mathfrak p) \otimes_A B$ 在剩余域
$\kappa(\mathfrak p)$ 上光滑。更一般地，概形态射称为{\it 光滑的}，如果它
平坦、局部有限表示，并且各几何纤维光滑。

\medskip\noindent
关于这些事实，见 Morphisms，第 \ref{morphisms-section-smooth} 节。由此可如下
定义 \'etale 态射。

\begin{definition}
\label{etale-cohomology-definition-etale-morphism}
若概形态射是相对维数为 0 的光滑态射，则称它是{\it \'etale 的}。
\end{definition}

\noindent
特别地，若概形态射 $X \to S$ 光滑且 $\Omega_{X/S} = 0$，则它是 \'etale 的。

\begin{proposition}
\label{etale-cohomology-proposition-etale-morphisms}
关于 \'etale 态射的若干事实。
\begin{enumerate}
\item 设 $k$ 是域。概形态射 $U \to \Spec(k)$ 是 \'etale 的，当且仅当
$U \cong \coprod_{i \in I} \Spec(k_i)$，其中对每个 $i \in I$，环 $k_i$ 是域，
并且是 $k$ 的有限可分扩张。
\item 设 $\varphi : U \to S$ 是概形态射。下列条件等价：
\begin{enumerate}
\item $\varphi$ 是 \'etale 的；
\item $\varphi$ 局部有限表示且平坦，并且它的所有纤维都是 \'etale 的；
\item $\varphi$ 平坦、非分歧且局部有限表示。
\end{enumerate}
\item 环映射 $A \to B$ 是 \'etale 的，当且仅当
$B \cong A[x_1, \ldots, x_n]/(f_1, \ldots, f_n)$
，并且 $\Delta = \det \left( \frac{\partial f_i}{\partial x_j} \right)$
在 $B$ 中可逆。
\item \'Etale 态射的基变换仍是 \'etale 态射。
\item \'Etale 态射的复合仍是 \'etale 态射。
\item \'Etale 态射的纤维积和乘积仍是 \'etale 态射。
\item \'Etale 态射的相对维数为 0。
\item 设 $Y \to X$ 是 \'etale 态射。若 $X$ 是既约的（相应地，正则的），
则 $Y$ 也是如此。
\item \'Etale 态射是开的。
\item 若 $X \to S$ 和 $Y \to S$ 都是 \'etale 的，则任意 $S$-态射
$X \to Y$ 也是 \'etale 的。
\end{enumerate}
\end{proposition}

\begin{proof}
我们已在前面各章证明了这些事实以及更多结果。参考文献如下：
(1) Morphisms，引理 \ref{morphisms-lemma-etale-over-field}。
(2) Morphisms，引理 \ref{morphisms-lemma-etale-flat-etale-fibres} 和
\ref{morphisms-lemma-flat-unramified-etale}。
(3) Algebra，引理 \ref{algebra-lemma-etale-standard-smooth}。
(4) Morphisms，引理 \ref{morphisms-lemma-base-change-etale}。
(5) Morphisms，引理 \ref{morphisms-lemma-composition-etale}。
(6) 由 (4) 和 (5) 形式地得出。
(7) Morphisms，引理 \ref{morphisms-lemma-etale-locally-quasi-finite} 和
\ref{morphisms-lemma-locally-quasi-finite-rel-dimension-0}。
(8) 见 Algebra，引理 \ref{algebra-lemma-reduced-goes-up} 和
\ref{algebra-lemma-Rk-goes-up}；此类更多结果还见 \'Etale Morphisms，第
\ref{etale-section-properties-permanence} 节。
(9) 见 Morphisms，引理 \ref{morphisms-lemma-fppf-open} 和
\ref{morphisms-lemma-etale-flat}。
(10) 见 Morphisms，引理 \ref{morphisms-lemma-etale-permanence}。
\end{proof}

\begin{definition}
\label{etale-cohomology-definition-standard-etale}
若环映射 $A \to B$ 满足
$B \cong \left(A[t]/(f)\right)_g$，其中 $f, g \in A[t]$、$f$ 首一，且
$\text{d}f/\text{d}t$ 在 $B$ 中可逆，则称它是{\it 标准 \'etale 的}。
\end{definition}

\noindent
标准 \'etale 环映射确实是 \'etale 的。事实上，假设
$B = \left(A[t]/(f)\right)_g$，其中 $f, g \in A[t]$、$f$ 首一，且
$\text{d}f/\text{d}t$ 在 $B$ 中可逆。则 $A[t]/(f)$ 是有限自由 $A$-模，
其秩等于首一多项式 $f$ 的次数。因此，$B$ 作为这个自由代数的局部化，
在 $A$ 上有限表示且平坦。要完成 $B$ 为 \'etale 的证明，只需说明纤维环
$$
\kappa(\mathfrak p) \otimes_A B
\cong
\kappa(\mathfrak p) \otimes_A (A[t]/(f))_g
\cong
\kappa(\mathfrak p)[t, 1/\overline{g}]/(\overline{f})
$$
是有限可分域扩张的有限乘积。这里
$\overline{f}, \overline{g} \in \kappa(\mathfrak p)[t]$ 分别是 $f$ 和 $g$ 的像。
令
$$
\overline{f} = \overline{f}_1 \ldots \overline{f}_a
\overline{f}_{a + 1}^{e_1} \ldots \overline{f}_{a + b}^{e_b}
$$
是 $\overline{f}$ 分解成两两不同的不可约首一因子 $\overline{f}_i$ 的幂，
其中 $e_1, \ldots, e_b > 1$。根据假设，
$\text{d}\overline{f}/\text{d}t$ 在
$\kappa(\mathfrak p)[t, 1/\overline{g}]$ 中可逆。因此至少所有
$\overline{f}_i$（其中 $i > a$）都可逆。由此得出
$$
\kappa(\mathfrak p)[t, 1/\overline{g}]/(\overline{f})
\cong
\prod\nolimits_{i \in I} \kappa(\mathfrak p)[t]/(\overline{f}_i)
$$
，其中 $I \subset \{1, \ldots, a\}$ 是指标 $i$ 的子集，所取指标满足
$\overline{f}_i$ 不整除 $\overline{g}$。此外，
$\text{d}\overline{f}/\text{d}t$ 在因子
$\kappa(\mathfrak p)[t]/(\overline{f}_i)$ 中的像，显然等于某个单位乘以
$\text{d}\overline{f}_i/\text{d}t$。因此，
$\kappa_i = \kappa(\mathfrak p)[t]/(\overline{f}_i)$ 是
$\kappa(\mathfrak p)$ 的有限域扩张，由一个极小多项式可分的元素生成；也就是说，
域扩张 $\kappa_i/\kappa(\mathfrak p)$ 如所需是有限可分的。

\medskip\noindent
事实上，任意 \'etale 环映射都局部标准 \'etale。为陈述这一点，引入下述记号。
称环映射 $A \to B${\it 在素理想 $\mathfrak q$（它是 $B$ 的素理想）处
\'etale}，如果存在 $h \in B$、$h \not \in \mathfrak q$，使得
$A \to B_h$ 是 \'etale 的。结果如下。

\begin{theorem}
\label{etale-cohomology-theorem-standard-etale}
环映射 $A \to B$ 在素理想 $\mathfrak q$ 处 \'etale，当且仅当存在
$g \in B$、$g \not \in \mathfrak q$，使得 $B_g$ 在 $A$ 上标准 \'etale。
\end{theorem}

\begin{proof}
见 Algebra，命题 \ref{algebra-proposition-etale-locally-standard}。
\end{proof}





\section{\'Etale 覆盖}
\label{etale-cohomology-section-etale-covering}

\noindent
回忆其定义。

\begin{definition}
\label{etale-cohomology-definition-etale-covering}
概形 $U$ 的一个{\it \'etale 覆盖}是概形态射族
$\{\varphi_i : U_i \to U\}_{i \in I}$，满足：
\begin{enumerate}
\item 每个 $\varphi_i$ 都是 \'etale 态射；
\item 诸 $U_i$ 覆盖 $U$，即 $U = \bigcup_{i\in I}\varphi_i(U_i)$。
\end{enumerate}
\end{definition}

\begin{lemma}
\label{etale-cohomology-lemma-etale-fpqc}
任意 \'etale 覆盖都是 fpqc 覆盖。
\end{lemma}

\begin{proof}
（另见 Topologies，引理
\ref{topologies-lemma-zariski-etale-smooth-syntomic-fppf-fpqc}。）设
$\{\varphi_i : U_i \to U\}_{i \in I}$ 是 \'etale 覆盖。由于 \'etale 态射平坦，
且覆盖的成员覆盖其靶，性质 fp（忠实平坦）成立。为验证性质 qc（拟紧），设
$V \subset U$ 是仿射开集，并写成
$\varphi_i^{-1}(V) = \bigcup_{j \in J_i} V_{ij}$，其中
$V_{ij} \subset U_i$ 是某些仿射开集。因为 $\varphi_i$ 是开态射
（\'etale 态射是开的），所以
$V = \bigcup_{i\in I} \bigcup_{j \in J_i} \varphi_i(V_{ij})$
是 $V$ 的开覆盖。又因 $V$ 拟紧，该覆盖有有限加细。
\end{proof}

\noindent
因此，对 fpqc 覆盖成立的任意陈述，{\it 更加}对 \'etale 覆盖成立。例如，
\'etale 网站是次典范的。

\begin{definition}
\label{etale-cohomology-definition-big-etale-site}
（更多细节见第 \ref{etale-cohomology-section-big-small} 节，或 Topologies，第
\ref{topologies-section-etale} 节。）设 $S$ 是概形。{\it $S$ 上的大 \'etale
网站}是网站 $(\Sch/S)_\etale$，见定义 \ref{etale-cohomology-definition-tau-site}。
{\it $S$ 上的小 \'etale 网站}是网站 $S_\etale$，见定义
\ref{etale-cohomology-definition-tau-site}。类似地，定义 $S$ 上的{\it 大}和{\it 小 Zariski
网站}，分别记作 $(\Sch/S)_{Zar}$ 和 $S_{Zar}$。
\end{definition}

\noindent
粗略地说，$S$ 的大 \'etale 网站由 $S$ 上的概形组成，其覆盖为 \'etale 覆盖；
$S$ 的小 \'etale 网站由在 $S$ 上 \'etale 的概形组成，其覆盖仍为 \'etale 覆盖。
事实上，由命题 \ref{etale-cohomology-proposition-etale-morphisms}，$S_\etale$ 的对象之间的任意态射
都是 \'etale 的。因此，要验证 $\{U_i \to U\}_{i \in I}$（位于 $S_\etale$ 中）
是覆盖，只需验证 $\coprod U_i \to U$ 满射。

\medskip\noindent
小 \'etale 网站的对象少于大 \'etale 网站；它只包含 $S$ 的 \'etale 拓扑中的
“开集”。它是大 \'etale 网站的全子范畴，其拓扑由大网站上的拓扑诱导。因此，
从大 \'etale 网站到小网站的限制函子正合，并把内射对象映为内射对象。由此得到
下述结论。

\begin{proposition}
\label{etale-cohomology-proposition-cohomology-restrict-small-site}
设 $S$ 是概形，$\mathcal{F}$ 是 $(\Sch/S)_\etale$ 上的阿贝尔层。则
$\mathcal{F}|_{S_\etale}$ 是 $S_\etale$ 上的层，并且
$$
H^p_\etale(S, \mathcal{F}|_{S_\etale}) =
H^p_\etale(S, \mathcal{F})
$$
对所有 $p \geq 0$ 都成立。
\end{proposition}

\begin{proof}
这是引理 \ref{etale-cohomology-lemma-compare-cohomology-big-small} 的特殊情形。
\end{proof}

\noindent
按照第 \ref{etale-cohomology-section-big-small} 节引入的一般记号，把上述上同调群记作
$H_\etale^p(S, \mathcal{F})$。





%9.24.09
\section{Kummer 理论}
\label{etale-cohomology-section-kummer}

\noindent
设 $n \in \mathbf{N}$，并考虑如下定义的函子 $\mu_n$：
$$
\begin{matrix}
\Sch^{opp} & \longrightarrow & \textit{Ab} \\
S & \longmapsto &
\mu_n(S)
=
\{t \in \Gamma(S, \mathcal{O}_S^*) \mid t^n = 1 \}.
\end{matrix}
$$
由 Groupoids，例 \ref{groupoids-example-roots-of-unity}，这是可表函子，表示它的
概形也记作 $\mu_n$。由引理 \ref{etale-cohomology-lemma-representable-sheaf-fpqc}，该函子满足
fpqc 拓扑的层条件（特别地，它也满足 \'etale、Zariski 等拓扑的层条件）。

\begin{lemma}
\label{etale-cohomology-lemma-kummer-sequence}
若 $n\in \mathcal{O}_S^*$，则
$$
0 \to
\mu_{n, S} \to
\mathbf{G}_{m, S} \xrightarrow{(\cdot)^n}
\mathbf{G}_{m, S} \to 0
$$
在 $S$ 的小 \'etale 网站和大 \'etale 网站上都是层的短正合列。
\end{lemma}

\begin{proof}
按定义，层 $\mu_{n, S}$ 是映射 $(\cdot)^n$ 的核。因此，只需证明最后一个映射
满射。设 $U$ 是 $S$ 上的概形。设
$f \in \mathbf{G}_m(U) = \Gamma(U, \mathcal{O}_U^*)$。需要证明：可找到 $U$ 的
一个 \'etale 覆盖，使得 $f$ 限制到每个覆盖成员后都是 $n$ 次幂。令
$$
U' =
\underline{\Spec}_U(\mathcal{O}_U[T]/(T^n-f))
\xrightarrow{\pi}
U.
$$
（关于相对谱，见 Constructions，第
\ref{constructions-section-spec-via-glueing} 或
\ref{constructions-section-spec} 节。）设 $\Spec(A) \subset U$ 是仿射开集，
并设 $f|_{\Spec(A)}$ 对应于单位 $a \in A^*$。则
$\pi^{-1}(\Spec(A)) = \Spec(B)$，其中 $B = A[T]/(T^n - a)$。环映射
$A \to B$ 是秩为 $n$ 的有限自由映射，因而忠实平坦，所以
$\Spec(B) \to \Spec(A)$ 满射。该结论对 $U$ 中每个仿射开集都成立，故 $\pi$
满射。此外，$n$ 和 $T^{n - 1}$ 在 $B$ 中可逆，所以
$nT^{n-1} \in B^*$，并且环映射 $A \to B$ 标准 \'etale，特别地是 \'etale 的。
该结论对 $U$ 中每个仿射开集都成立，故 $\pi$ 是 \'etale 的。因此
$\mathcal{U} = \{\pi : U' \to U\}$ 是 \'etale 覆盖。此外，
$f|_{U'} = (f')^n$，其中 $f'$ 是 $T$ 在
$\Gamma(U', \mathcal{O}_{U'}^*)$ 中的类，故 $\mathcal{U}$ 具有所需性质。
\end{proof}

\begin{remark}
\label{etale-cohomology-remark-no-kummer-sequence-zariski}
若把引理 \ref{etale-cohomology-lemma-kummer-sequence} 中的“\'etale”换成“Zariski”，则引理不成立。
由于 \'etale 拓扑比光滑拓扑粗，见 Topologies，引理
\ref{topologies-lemma-zariski-etale-smooth}，可知该序列在光滑拓扑中也正合。
\end{remark}

\noindent
由定理 \ref{etale-cohomology-theorem-picard-group}、引理 \ref{etale-cohomology-lemma-kummer-sequence} 以及上同调的
一般性质，得到上同调长正合列
$$
\xymatrix{
0 \ar[r] &
H_\etale^0(S, \mu_{n, S}) \ar[r] &
\Gamma(S, \mathcal{O}_S^*) \ar[r]^{(\cdot)^n} &
\Gamma(S, \mathcal{O}_S^*) \ar@(rd, ul)[rdllllr]
\\
& H_\etale^1(S, \mu_{n, S}) \ar[r] &
\Pic(S) \ar[r]^{(\cdot)^n} &
\Pic(S) \ar@(rd, ul)[rdllllr] \\
& H_\etale^2(S, \mu_{n, S}) \ar[r] &
\ldots
}
$$
；至少当 $n$ 在 $S$ 上可逆时如此。当 $n$ 在 $S$ 上不可逆时，可应用下述引理。

\begin{lemma}
\label{etale-cohomology-lemma-kummer-sequence-syntomic}
对任意 $n \in \mathbf{N}$，序列
$$
0 \to
\mu_{n, S} \to
\mathbf{G}_{m, S} \xrightarrow{(\cdot)^n}
\mathbf{G}_{m, S} \to 0
$$
在网站 $(\Sch/S)_{fppf}$ 和 $(\Sch/S)_{syntomic}$ 上都是层的短正合列。
\end{lemma}

\begin{proof}
按定义，层 $\mu_{n, S}$ 是映射 $(\cdot)^n$ 的核。因此，只需证明最后一个映射
满射。由于 syntomic 拓扑弱于 fppf 拓扑，见 Topologies，引理
\ref{topologies-lemma-zariski-etale-smooth-syntomic-fppf}，只需对 syntomic 拓扑
证明此事。设 $U$ 是 $S$ 上的概形。设
$f \in \mathbf{G}_m(U) = \Gamma(U, \mathcal{O}_U^*)$。需要证明：可找到 $U$ 的
一个 syntomic 覆盖，使得 $f$ 限制到每个覆盖成员后都是 $n$ 次幂。令
$$
U' =
\underline{\Spec}_U(\mathcal{O}_U[T]/(T^n-f))
\xrightarrow{\pi}
U.
$$
（关于相对谱，见 Constructions，第
\ref{constructions-section-spec-via-glueing} 或
\ref{constructions-section-spec} 节。）设 $\Spec(A) \subset U$ 是仿射开集，
并设 $f|_{\Spec(A)}$ 对应于单位 $a \in A^*$。则
$\pi^{-1}(\Spec(A)) = \Spec(B)$，其中 $B = A[T]/(T^n - a)$。环映射
$A \to B$ 是秩为 $n$ 的有限自由映射，因而忠实平坦，所以
$\Spec(B) \to \Spec(A)$ 满射。该结论对 $U$ 中每个仿射开集都成立，故 $\pi$
满射。此外，$B$ 在 $A$ 上是整体相对完全交，所以环映射 $A \to B$ 标准
syntomic，特别地是 syntomic 的。该结论对 $U$ 中每个仿射开集都成立，故 $\pi$
是 syntomic 的。因此 $\mathcal{U} = \{\pi : U' \to U\}$ 是 syntomic 覆盖。
此外，$f|_{U'} = (f')^n$，其中 $f'$ 是 $T$ 在
$\Gamma(U', \mathcal{O}_{U'}^*)$ 中的类，故 $\mathcal{U}$ 具有所需性质。
\end{proof}

\begin{remark}
\label{etale-cohomology-remark-no-kummer-sequence-smooth-etale-zariski}
引理 \ref{etale-cohomology-lemma-kummer-sequence-syntomic} 对光滑、\'etale 或 Zariski 拓扑不成立。
\end{remark}

\noindent
由定理 \ref{etale-cohomology-theorem-picard-group}、引理
\ref{etale-cohomology-lemma-kummer-sequence-syntomic} 以及上同调的一般性质，得到上同调长正合列
$$
\xymatrix{
0 \ar[r] &
H_{fppf}^0(S, \mu_{n, S}) \ar[r] &
\Gamma(S, \mathcal{O}_S^*) \ar[r]^{(\cdot)^n} &
\Gamma(S, \mathcal{O}_S^*) \ar@(rd, ul)[rdllllr]
\\
& H_{fppf}^1(S, \mu_{n, S}) \ar[r] &
\Pic(S) \ar[r]^{(\cdot)^n} &
\Pic(S) \ar@(rd, ul)[rdllllr] \\
& H_{fppf}^2(S, \mu_{n, S}) \ar[r] &
\ldots
}
$$
，它对任意概形 $S$ 和任意整数 $n$ 成立。当然，对 syntomic 上同调也有类似序列。

\medskip\noindent
设 $n \in \mathbf{N}$，并设 $S$ 是任意概形。还有一种更直接的方法描述取值于
$\mu_n$ 的第一上同调群。考虑序对 $(\mathcal{L}, \alpha)$，其中
$\mathcal{L}$ 是 $S$ 上的可逆层，而
$\alpha : \mathcal{L}^{\otimes n} \to \mathcal{O}_S$ 是第 $n$ 次张量幂的平凡化，
这一张量幂来自 $\mathcal{L}$。设 $(\mathcal{L}', \alpha')$ 是另一这样的序对。同构
$\varphi : (\mathcal{L}, \alpha) \to (\mathcal{L}', \alpha')$ 是可逆层的同构
$\varphi : \mathcal{L} \to \mathcal{L}'$，并使下图
$$
\xymatrix{
\mathcal{L}^{\otimes n} \ar[d]_{\varphi^{\otimes n}} \ar[r]_\alpha &
\mathcal{O}_S \ar[d]^1 \\
(\mathcal{L}')^{\otimes n} \ar[r]^{\alpha'} &
\mathcal{O}_S \\
}
$$
交换。因此有
\begin{equation}
\label{etale-cohomology-equation-isomorphisms-pairs}
\mathit{Isom}_S((\mathcal{L}, \alpha), (\mathcal{L}', \alpha'))
=
\left\{
\begin{matrix}
\emptyset & \text{若} & \text{二者不同构} \\
H^0(S, \mu_{n, S})\cdot \varphi & \text{若} &
\varphi \text{ 是序对同构}
\end{matrix}
\right.
\end{equation}
此外，给定两个序对 $(\mathcal{L}, \alpha)$、$(\mathcal{L}', \alpha')$，其张量积
$$
(\mathcal{L}, \alpha) \otimes (\mathcal{L}', \alpha')
=
(\mathcal{L} \otimes \mathcal{L}', \alpha \otimes \alpha')
$$
仍是一个序对。序对 $(\mathcal{O}_S, 1)$ 是该张量积运算的单位元，而逆元由
$$
(\mathcal{L}, \alpha)^{-1} = (\mathcal{L}^{\otimes -1}, \alpha^{\otimes -1}).
$$
给出。因此，序对的同构类集合构成阿贝尔群。注意，
$$
(\mathcal{L}, \alpha)^{\otimes n}
=
(\mathcal{L}^{\otimes n}, \alpha^{\otimes n})
\xrightarrow{\alpha}
(\mathcal{O}_S, 1)
$$
是同构，故该群中每个元素的阶都整除 $n$。提醒读者：一般而言，这个群{\bf 并非}
$n$-挠子群（位于 $\Pic(S)$ 中）。

\begin{lemma}
\label{etale-cohomology-lemma-describe-h1-mun}
设 $S$ 是概形。有典范等同
$$
H_\etale^1(S, \mu_n) =
\text{上述序对 }(\mathcal{L}, \alpha)\text{ 的同构类群}
$$
，条件是 $n$ 在 $S$ 上可逆。一般而言，有
$$
H_{fppf}^1(S, \mu_n) =
\text{上述序对 }(\mathcal{L}, \alpha)\text{ 的同构类群}.
$$
把 fppf 换成 syntomic 后，同一结果仍成立。
\end{lemma}

\begin{proof}
先证明第二个同构。设 $(\mathcal{L}, \alpha)$ 是上述序对。选取仿射开覆盖
$S = \bigcup U_i$，使得 $\mathcal{L}|_{U_i} \cong \mathcal{O}_{U_i}$。设
$s_i \in \mathcal{L}(U_i)$ 是生成元。于是
$\alpha(s_i^{\otimes n}) = f_i \in \mathcal{O}_S^*(U_i)$。写
$U_i = \Spec(A_i)$，则存在整体相对完全交
$A_i \to B_i = A_i[T]/(T^n - f_i)$，使 $f_i$ 的像是 $n$ 次幂（在 $B_i$ 中）。
换言之，令 $V_i = \Spec(B_i)$，便得到 syntomic 覆盖
$\mathcal{V} = \{V_i \to S\}_{i \in I}$ 和平凡化
$\varphi_i : (\mathcal{L}, \alpha)|_{V_i} \to (\mathcal{O}_{V_i}, 1)$。

\medskip\noindent
我们用这个结果（覆盖 $\mathcal{V}$ 的存在性）把
$H^1_{syntomic}(S, \mu_{n, S})$ 中的一个上同调类与该序对相联系。给出两种
（等价的）构造。

\medskip\noindent
第一种构造：使用 {\v C}ech 上同调。在二重交 $V_i \times_S V_j$ 上，有序对的同构
$$
(\mathcal{O}_{V_i \times_S V_j}, 1)
\xrightarrow{\text{pr}_0^*\varphi_i^{-1}}
(\mathcal{L}|_{V_i \times_S V_j}, \alpha|_{V_i \times_S V_j})
\xrightarrow{\text{pr}_1^*\varphi_j}
(\mathcal{O}_{V_i \times_S V_j}, 1)
$$
。由 (\ref{etale-cohomology-equation-isomorphisms-pairs})，它由元素
$\zeta_{ij} \in \mu_n(V_i \times_S V_j)$ 给出。略去如下验证：这些
$\zeta_{ij}$ 给出一个 $1$-上闭链，也就是说，给出元素
$(\zeta_{i_0i_1}) \in \check C(\mathcal{V}, \mu_n)$，并且
$d(\zeta_{i_0i_1}) = 0$。因此，它的类是
$\check H^1(\mathcal{V}, \mu_n)$ 中的元素，并且由定理
\ref{etale-cohomology-theorem-cech-ss}，它映到 $H^1_{syntomic}(S, \mu_{n, S})$ 中的上同调类。

\medskip\noindent
第二种构造：使用挠子。考虑预层
$$
\mu_n(\mathcal{L}, \alpha) :
U
\longmapsto
\mathit{Isom}_U((\mathcal{O}_U, 1), (\mathcal{L}, \alpha)|_U)
$$
，它位于 $(\Sch/S)_{syntomic}$ 上。可把它看成
$\SheafHom_\mathcal{O}(\mathcal{O}, \mathcal{L})$ 的子预层（这是内 Hom 层，
见 Modules on Sites，第 \ref{sites-modules-section-internal-hom} 节）。由于定义
该子预层的条件是局部的，它是层。由
(\ref{etale-cohomology-equation-isomorphisms-pairs})，这个层带有层 $\mu_{n, S}$ 的自由作用。
因此，只需验证它局部有截面。平凡化覆盖 $\mathcal{V}$ 的存在性说明了这一点。
所以 $\mu_n(\mathcal{L}, \alpha)$ 是 $\mu_{n, S}$-挠子；由
Cohomology on Sites，引理 \ref{sites-cohomology-lemma-torsors-h1}，得到
$H^1_{syntomic}(S, \mu_{n, S})$ 中相应的元素。

\medskip\noindent
现在还需证明下列各点：
\begin{enumerate}
\item 两种构造给出相同的上同调类。
\item 同构的序对给出相同的上同调类。
\item
$(\mathcal{L}, \alpha) \otimes (\mathcal{L}', \alpha')$
的上同调类，是 $(\mathcal{L}, \alpha)$ 与 $(\mathcal{L}', \alpha')$ 的上同调类之和。
\item 若上同调类平凡，则序对平凡。
\item $H^1_{syntomic}(S, \mu_{n, S})$ 的任意元素都是某个序对的上同调类。
\end{enumerate}
略去 (1) 的证明。(2) 由第二种构造显然成立，因为同构挠子给出相同的上同调类。
(3) 由第一种构造显然成立，因为所得 {\v C}ech 类相加。(4) 由第二种构造显然
成立，因为挠子平凡当且仅当它有整体截面，见 Cohomology on Sites，引理
\ref{sites-cohomology-lemma-trivial-torsor}。

\medskip\noindent
(5) 可如下看出（虽然直接证明会更好）。假设
$\xi \in H^1_{syntomic}(S, \mu_{n, S})$。则 $\xi$ 映到元素
$\overline{\xi} \in H^1_{syntomic}(S, \mathbf{G}_{m, S})$，且
$n \overline{\xi} = 0$。由定理 \ref{etale-cohomology-theorem-picard-group} 可知，
$\overline{\xi}$ 对应于可逆层 $\mathcal{L}$，其第 $n$ 次张量幂同构于
$\mathcal{O}_S$。因此，存在序对 $(\mathcal{L}, \alpha')$，其上同调类 $\xi'$
有同样的像 $\overline{\xi'}$，该像位于
$H^1_{syntomic}(S, \mathbf{G}_{m, S})$ 中。所以只需证明 $\xi - \xi'$ 是某个
序对的类。由构造及上面的上同调长正合列可知，
$\xi - \xi' = \partial(f)$，其中某个 $f \in H^0(S, \mathcal{O}_S^*)$。
考虑序对 $(\mathcal{O}_S, f)$。略去如下验证：
该序对的上同调类是 $\partial(f)$。这就完成第一个等同（把 fppf 换成 syntomic）
的证明。

\medskip\noindent
为看出第一个等同，注意：若 $n$ 在 $S$ 上可逆，则证明第一部分构造的覆盖
$\mathcal{V}$ 实际上是 \'etale 覆盖（比较引理
\ref{etale-cohomology-lemma-kummer-sequence} 的证明）。证明的其余部分与拓扑无关，只有最后一个
论证使用了 Kummer 序列的正合性，即使用了引理
\ref{etale-cohomology-lemma-kummer-sequence}。
\end{proof}






\section{邻域、茎与点}
\label{etale-cohomology-section-stalks}

\noindent
可为 $S$ 的每个几何点配上一个正合的茎函子。$S_\etale$ 上的层态射是同构，
当且仅当它在所有这些茎上都是同构。阿贝尔层复形正合，当且仅当其茎复形在
每个几何点处都正合。合起来说，这意味着概形 $S$ 的小 \'etale 网站有足够多的点。
而且，$S$ 的小 \'etale 拓扑斯的每个点（这是一个抽象概念）都由一个几何点给出。
因此，在某种意义上，可用几何点及其邻域来理解 $S$ 的小 \'etale 拓扑斯。

\begin{definition}
\label{etale-cohomology-definition-geometric-point}
设 $S$ 是概形。
\begin{enumerate}
\item $S$ 的一个{\it 几何点}是态射
$\Spec(k) \to S$，其中 $k$ 是代数闭域。这种点通常记作 $\overline{s}$，
即用带上划线的小写字母表示。我们常用 $\overline{s}$ 同时表示概形
$\Spec(k)$ 和该态射，并用 $\kappa(\overline{s})$ 表示 $k$。
\item 说 $\overline{s}${\it 位于} $s${\it 之上}，意指
$s \in S$ 是 $\overline{s}$ 的像。
\item 几何点 $\overline{s}$ 在 $S$ 中的一个{\it \'etale 邻域}是交换图
$$
\xymatrix{
& U \ar[d]^\varphi \\
{\overline{s}} \ar[r]^{\overline{s}} \ar[ur]^{\bar u} & S
}
$$
，其中 $\varphi$ 是概形的 \'etale 态射。把它记作
$(U, \overline{u}) \to (S, \overline{s})$。
\item {\it \'Etale 邻域的态射}
$(U, \overline{u}) \to (U', \overline{u}')$
是一个 $S$-态射 $h: U \to U'$，满足
$\overline{u}' = h \circ \overline{u}$。
\end{enumerate}
\end{definition}

\begin{remark}
\label{etale-cohomology-remark-etale-between-etale}
由于 $U$ 和 $U'$ 在 $S$ 上 \'etale，它们之间的任意 $S$-态射也都是
\'etale 的，见命题 \ref{etale-cohomology-proposition-etale-morphisms}。特别地，
\'etale 邻域的所有态射都是 \'etale 的。
\end{remark}

\begin{remark}
\label{etale-cohomology-remark-etale-neighbourhoods}
设 $S$ 是概形，$s \in S$ 是一点。在 More on Morphisms，定义
\ref{more-morphisms-definition-etale-neighbourhood} 中，我们定义了
\'etale 邻域 $(U, u) \to (S, s)$ 作为 $(S, s)$ 的邻域这一概念。若
$\overline{s}$ 是 $S$ 中位于 $s$ 之上的几何点，则任意 \'etale 邻域
$(U, \overline{u}) \to (S, \overline{s})$ 都给出 $(U, u)$，即
$(S, s)$ 的一个 \'etale 邻域：取 $u \in U$ 为 $U$ 中使
$\overline{u}$ 位于 $u$ 之上的唯一一点。反过来，给定
$(U, u)$，即 $(S, s)$ 的一个 \'etale 邻域，剩余域扩张
$\kappa(u)/\kappa(s)$ 有限可分（见命题
\ref{etale-cohomology-proposition-etale-morphisms}），因而可找到嵌入
$\kappa(u) \subset \kappa(\overline{s})$，它位于 $\kappa(s)$ 上。
换言之，可找到几何点 $\overline{u}$，它是 $U$ 中位于 $u$ 之上的点，
使得 $(U, \overline{u})$ 是 $(S, \overline{s})$ 的 \'etale 邻域。
我们将用这些观察在两类
\'etale 邻域之间转换。
\end{remark}

\begin{lemma}
\label{etale-cohomology-lemma-cofinal-etale}
设 $S$ 是概形，$\overline{s}$ 是 $S$ 的几何点。\'Etale 邻域的范畴是余滤过的。
更确切地说：
\begin{enumerate}
\item 设 $(U_i, \overline{u}_i)_{i = 1, 2}$ 是 $\overline{s}$ 在 $S$
中的两个 \'etale 邻域。则存在第三个 \'etale 邻域 $(U, \overline{u})$
以及态射
$(U, \overline{u}) \to (U_i, \overline{u}_i)$，$i = 1, 2$。
\item 设
$h_1, h_2: (U, \overline{u}) \to (U', \overline{u}')$
是 $\overline{s}$ 的 \'etale 邻域之间的两个态射。则存在 \'etale 邻域
$(U'', \overline{u}'')$ 和态射
$h : (U'', \overline{u}'') \to (U, \overline{u})$，使 $h_1$ 与 $h_2$
相等，即 $h_1 \circ h = h_2 \circ h$。
\end{enumerate}
\end{lemma}

\begin{proof}
对于 (1)，考虑纤维积 $U = U_1 \times_S U_2$。由于 \'etale 态射在基变换下
保持，见命题 \ref{etale-cohomology-proposition-etale-morphisms}，它在 $U_1$ 和 $U_2$ 上都
\'etale。映射 $\overline{s} \to U$（由
$(\overline{u}_1, \overline{u}_2)$ 定义）赋予它一个映到 $U_1$ 和
$U_2$ 的 \'etale 邻域结构。
对于 (2)，把 $U''$ 定义为纤维积
$$
\xymatrix{
U'' \ar[r] \ar[d] & U \ar[d]^{(h_1, h_2)} \\
U' \ar[r]^-\Delta & U' \times_S U'.
}
$$
由于 $\overline{u}$ 和 $\overline{u}'$ 在 $S$ 上都与 $\overline{s}$
相容，可知 $\overline{u}'' = (\overline{u}, \overline{u}')$ 是 $U''$
的几何点。特别地，$U'' \not = \emptyset$。此外，由于 $U'$ 在 $S$ 上
\'etale，纤维积 $U'\times_S U'$ 也如此（见命题
\ref{etale-cohomology-proposition-etale-morphisms}）。故由上面的注
\ref{etale-cohomology-remark-etale-between-etale}，竖直箭头 $(h_1, h_2)$ 是 \'etale 的。
于是经基变换，$U''$ 在 $U'$ 上 \'etale，因而也在 $S$ 上 \'etale
（因为 \'etale 态射的复合仍是 \'etale 的）。所以
$(U'', \overline{u}'')$ 给出了所求解答。
\end{proof}

\begin{lemma}
\label{etale-cohomology-lemma-geometric-lift-to-cover}
设 $S$ 是概形，$\overline{s}$ 是 $S$ 的几何点，
$(U, \overline{u})$ 是 $\overline{s}$ 的 \'etale 邻域，并设
$\mathcal{U} = \{\varphi_i : U_i \to U \}_{i\in I}$ 是 \'etale 覆盖。
则存在 $i \in I$ 和 $\overline{u}_i : \overline{s} \to U_i$，使得
$\varphi_i : (U_i, \overline{u}_i) \to (U, \overline{u})$
是 \'etale 邻域的态射。
\end{lemma}

\begin{proof}
因为 $U = \bigcup_{i\in I} \varphi_i(U_i)$，纤维积
$\overline{s} \times_{\overline{u}, U, \varphi_i} U_i$ 对某个 $i$ 非空。
考察笛卡尔图
$$
\xymatrix{
\overline{s} \times_{\overline{u}, U, \varphi_i} U_i
\ar[d]^{\text{pr}_1} \ar[r]_-{\text{pr}_2} & U_i
\ar[d]^{\varphi_i} \\
\Spec(k) = \overline{s} \ar@/^1pc/[u]^\sigma
\ar[r]^-{\overline{u}} & U
}
$$
。投影 $\text{pr}_1$ 是 \'etale 态射的基变换，故为 \'etale 的，见命题
\ref{etale-cohomology-proposition-etale-morphisms}。因此，由命题
\ref{etale-cohomology-proposition-etale-morphisms}，
$\overline{s} \times_{\overline{u}, U, \varphi_i} U_i$
是 $k$ 的有限可分扩张的谱的不交并；这里
$\overline{s} = \Spec(k)$。但 $k$ 代数闭，所以这些扩张全都平凡，
从而存在截面 $\sigma$，它是 $\text{pr}_1$ 的截面。复合
$\text{pr}_2 \circ \sigma$ 给出与 $\overline{u}$ 相容的映射。
\end{proof}

\begin{definition}
\label{etale-cohomology-definition-stalk}
设 $S$ 是概形，$\mathcal{F}$ 是 $S_\etale$ 上的预层，而
$\overline{s}$ 是 $S$ 的几何点。$\mathcal{F}$ 在 $\overline{s}$ 处的
{\it 茎}是
$$
\mathcal{F}_{\overline{s}}
=
\colim_{(U, \overline{u})} \mathcal{F}(U)
$$
，其中 $(U, \overline{u})$ 遍历 $\overline{s}$ 在 $S$ 中的所有
\'etale 邻域。
\end{definition}

\noindent
由引理 \ref{etale-cohomology-lemma-cofinal-etale}，该余极限取遍一个滤过指标范畴，即
\'etale 邻域范畴的反范畴。换言之，可把
$\mathcal{F}_{\overline{s}}$ 的元素看作三元组
$(U, \overline{u}, \sigma)$，其中 $\sigma \in \mathcal{F}(U)$。
两个三元组 $(U, \overline{u}, \sigma)$、
$(U', \overline{u}', \sigma')$ 定义茎中的同一元素，当且仅当存在第三个
\'etale 邻域 $(U'', \overline{u}'')$ 以及 \'etale 邻域态射
$h : (U'', \overline{u}'') \to (U, \overline{u})$、
$h' : (U'', \overline{u}'') \to (U', \overline{u}')$，使得
$h^*\sigma = (h')^*\sigma'$ 在 $\mathcal{F}(U'')$ 中成立。见
Categories，第 \ref{categories-section-directed-colimits} 节。

\begin{lemma}
\label{etale-cohomology-lemma-stalk-gives-point}
设 $S$ 是概形，$\overline{s}$ 是 $S$ 的几何点。考虑函子
\begin{align*}
u : S_\etale & \longrightarrow \textit{Sets}, \\
U & \longmapsto
|U_{\overline{s}}|
=
\{\overline{u} \text{，使得 }(U, \overline{u})
\text{ 是 }\overline{s}\text{ 的一个 \'etale 邻域}\}.
\end{align*}
这里 $|U_{\overline{s}}|$ 表示几何纤维的底集。则 $u$ 定义一个点
$p$，它是网站 $S_\etale$ 的点（Sites，定义
\ref{sites-definition-point}），
而与之相伴的茎函子 $\mathcal{F} \mapsto \mathcal{F}_p$
（Sites，等式 \ref{sites-equation-stalk}）就是上面定义的函子
$\mathcal{F} \mapsto \mathcal{F}_{\overline{s}}$。
\end{lemma}

\begin{proof}
在引理 \ref{etale-cohomology-lemma-geometric-lift-to-cover} 的证明中已经看到，概形
$U_{\overline{s}}$ 是与 $\overline{s}$ 同构的概形的不交并。因此，也可把
$|U_{\overline{s}}|$ 看作 $U$ 中位于 $\overline{s}$ 之上的几何点集合，
即所有能嵌入定义 \ref{etale-cohomology-definition-geometric-point} 中图式的态射
$\overline{u} : \overline{s} \to U$ 所成的集合。由此可知 $u(S)$ 是单点集，并且
$u(U \times_V W) = u(U) \times_{u(V)} u(W)$
，只要 $U \to V$ 和 $W \to V$ 是 $S_\etale$ 中的态射。又给定覆盖
$\{U_i \to U\}_{i \in I}$，它位于 $S_\etale$ 中，由引理
\ref{etale-cohomology-lemma-geometric-lift-to-cover} 可知
$\coprod u(U_i) \to u(U)$ 满射。因此可应用 Sites，命题
\ref{sites-proposition-point-limits}，故 $p$ 是网站 $S_\etale$ 的一个点。
最后，我们的函子 $\mathcal{F} \mapsto \mathcal{F}_{\overline{s}}$
与 Sites，等式 \ref{sites-equation-stalk} 中的函子
$\mathcal{F} \mapsto \mathcal{F}_p$（它和 $p$ 相伴）由完全相同的
余极限给出，这就证明了
最后一个断言。
\end{proof}

\begin{remark}
\label{etale-cohomology-remark-map-stalks}
设 $S$ 是概形，并设 $\overline{s} : \Spec(k) \to S$ 和
$\overline{s}' : \Spec(k') \to S$ 是 $S$ 的两个几何点。几何点的一个
{\it 态射 $a : \overline{s} \to \overline{s}'$} 就是一个态射
$a : \Spec(k) \to \Spec(k')$，它满足
$\overline{s}' \circ a = \overline{s}$。给定这样的态射，得到从
$\overline{s}'$ 的 \'etale 邻域范畴到 $\overline{s}$ 的 \'etale 邻域范畴的
函子，其规则为
$(U, \overline{u}') \mapsto (U, \overline{u}' \circ a)$。因此得到典范映射
$$
\mathcal{F}_{\overline{s}'}
=
\colim_{(U, \overline{u}')} \mathcal{F}(U)
\longrightarrow
\colim_{(U, \overline{u})} \mathcal{F}(U)
=
\mathcal{F}_{\overline{s}}
$$
，见 Categories，引理 \ref{categories-lemma-functorial-colimit}。利用把茎中元素
描述为三元组的方式，该映射把 $\mathcal{F}_{\overline{s}'}$ 中由三元组
$(U, \overline{u}', \sigma)$ 表示的元素，映到
$\mathcal{F}_{\overline{s}}$ 中由三元组
$(U, \overline{u}' \circ a, \sigma)$ 表示的元素。由于上述函子显然是等价，
可知这个典范映射是茎函子的同构。

\medskip\noindent
我们来确认，与 $a$ 对应的茎映射方向无误。注意，按 Sites，定义
\ref{sites-definition-morphism-points}，上述内容意味着 $a$ 定义态射
$a : p \to p'$；这是点 $p, p'$ 之间的态射，它们是网站
$S_\etale$ 的点，而这些点由
引理 \ref{etale-cohomology-lemma-stalk-gives-point} 分别与
$\overline{s}, \overline{s}'$ 相伴。还有更一般的点态射（对应于 $S$ 中点的
特殊化），我们稍后会描述；这些态射将不是同构，见第
\ref{etale-cohomology-section-specialization} 节。
\end{remark}

\begin{lemma}
\label{etale-cohomology-lemma-stalk-exact}
设 $S$ 是概形，$\overline{s}$ 是 $S$ 的几何点。
\begin{enumerate}
\item 茎函子
$\textit{PAb}(S_\etale) \to \textit{Ab}$,
$\mathcal{F}  \mapsto  \mathcal{F}_{\overline{s}}$
是正合的。
\item 有 $(\mathcal{F}^\#)_{\overline{s}} = \mathcal{F}_{\overline{s}}$，
其中 $\mathcal{F}$ 是 $S_\etale$ 上任意集合预层。
\item 函子
$\textit{Ab}(S_\etale) \to \textit{Ab}$,
$\mathcal{F} \mapsto \mathcal{F}_{\overline{s}}$ 是正合的。
\item 类似地，函子
$\textit{PSh}(S_\etale) \to \textit{Sets}$ 以及
$\Sh(S_\etale) \to \textit{Sets}$（由茎函子
$\mathcal{F} \mapsto \mathcal{F}_{\overline{s}}$ 给出）也是正合的（见
Categories，定义
\ref{categories-definition-exact}），并与任意余极限交换。
\end{enumerate}
\end{lemma}

\begin{proof}
在说明如何直接证明之前，先注意该结果可由 Modules on Sites，第
\ref{sites-modules-section-stalks} 节的一般材料得出。之所以如此，是因为
$\mathcal{F} \mapsto \mathcal{F}_{\overline{s}}$ 来自 $S$ 的小 \'etale
网站的一个点，见引理 \ref{etale-cohomology-lemma-stalk-gives-point}。我们只直接证明
(1)、(2) 和 (3)，略去 (4) 的直接证明。

\medskip\noindent
它作为 $\textit{PAb}(S_\etale)$ 上函子的正合性，形式地来自如下事实：
有向余极限与所有余极限及有限极限交换。(2) 中茎的等同由映射
$$
\kappa :
\mathcal{F}_{\overline{s}}
\longrightarrow
(\mathcal{F}^\#)_{\overline{s}}
$$
给出；该映射由自然态射 $\mathcal{F}\to \mathcal{F}^\#$ 诱导，见定理
\ref{etale-cohomology-theorem-sheafification}。我们断言这个映射是阿贝尔群的同构。
下面证明单射性，并略去满射性的证明。

\medskip\noindent
设 $\sigma\in \mathcal{F}_{\overline{s}}$。存在 \'etale 邻域
$(U, \overline{u})\to (S, \overline{s})$，使 $\sigma$ 是某个截面
$s \in \mathcal{F}(U)$ 的像。若 $\kappa(\sigma) = 0$ 在
$(\mathcal{F}^\#)_{\overline{s}}$ 中成立，则存在 \'etale 邻域态射
$(U', \overline{u}')\to (U, \overline{u})$，使 $s|_{U'}$ 在
$\mathcal{F}^\#(U')$ 中为零。由此存在 \'etale 覆盖
$\{U_i'\to U'\}_{i\in I}$，使得 $s|_{U_i'}=0$ 在
$\mathcal{F}(U_i')$ 中成立，对所有 $i$ 都如此。由引理
\ref{etale-cohomology-lemma-geometric-lift-to-cover}，
存在 $i \in I$ 和态射 $\overline{u}_i': \overline{s} \to U_i'$，使
$(U_i', \overline{u}_i') \to (U', \overline{u}')\to (U, \overline{u})$
是 \'etale 邻域的态射。于是 $\sigma = 0$，因为
$(U_i', \overline{u}_i') \to (U, \overline{u})$ 是 \'etale 邻域态射，
并且 $s|_{U'_i}=0$。这证明了 $\kappa$ 是单射。

\medskip\noindent
为证明函子 $\textit{Ab}(S_\etale) \to \textit{Ab}$ 正合，考虑
$\textit{Ab}(S_\etale)$ 中任意短正合列：
$
0\to \mathcal{F}\to \mathcal{G}\to \mathcal H \to 0.
$
它给出预层的正合列
$$
0 \to \mathcal{F} \to \mathcal{G} \to \mathcal H \to
\mathcal H/^p\mathcal{G} \to 0,
$$
，其中 $/^p$ 表示 $\textit{PAb}(S_\etale)$ 中的商。在
$\overline{s}$ 处取茎，可知
$(\mathcal H /^p\mathcal{G})_{\bar{s}} =
(\mathcal H /\mathcal{G})_{\bar{s}} = 0$，因为
$\mathcal H/^p\mathcal{G}$ 的层化是 $0$。因此，
$$
0\to \mathcal{F}_{\overline{s}} \to \mathcal{G}_{\overline{s}} \to
\mathcal{H}_{\overline{s}} \to 0 = (\mathcal H/^p\mathcal{G})_{\overline{s}}
$$
正合，因为取茎作为预层上的函子是正合的。
\end{proof}

\begin{theorem}
\label{etale-cohomology-theorem-exactness-stalks}
设 $S$ 是概形。集合层的态射 $a : \mathcal{F} \to \mathcal{G}$ 是单射
（相应地，满射），当且仅当茎上的映射
$a_{\overline{s}} : \mathcal{F}_{\overline{s}} \to \mathcal{G}_{\overline{s}}$
对 $S$ 的每个几何点都是单射（相应地，满射）。$S_\etale$ 上的阿贝尔层序列正合，当且仅当它在
$S$ 的几何点处的所有茎上都正合。
\end{theorem}

\begin{proof}
茎上正合的必要性由引理 \ref{etale-cohomology-lemma-stalk-exact} 得出。反过来，只需证明：
层态射满射（相应地，单射），当且仅当它在所有茎上满射（相应地，单射）。
我们证明满射情形，略去单射情形的证明。

\medskip\noindent
设 $\alpha : \mathcal{F} \to \mathcal{G}$ 是层态射，并且对所有几何点，
$\mathcal{F}_{\overline{s}} \to \mathcal{G}_{\overline{s}}$ 都满射。固定
$U \in \Ob(S_\etale)$ 和 $s \in \mathcal{G}(U)$。对每个 $u \in U$，
选取某个 $\overline{u} \to U$（它位于 $u$ 之上）以及 \'etale 邻域
$(V_u , \overline{v}_u) \to (U, \overline{u})$，使得
$s|_{V_u} = \alpha(s_{V_u})$ 对某个
$s_{V_u} \in \mathcal{F}(V_u)$ 成立。这是可能的，因为 $\alpha$ 在茎上满射。
于是 $\{V_u \to U\}_{u \in U}$ 是一个 \'etale 覆盖，在其上 $s$ 的限制属于
映射 $\alpha$ 的像。因此 $\alpha$ 满射，见 Sites，第
\ref{sites-section-sheaves-injective} 节。
\end{proof}

\begin{remarks}
\label{etale-cohomology-remarks-enough-points}
关于几何网站的点。
\begin{enumerate}
\item 定理 \ref{etale-cohomology-theorem-exactness-stalks} 说明：$S_\etale$ 的点族
（它由 $S$ 的几何点给出，见引理 \ref{etale-cohomology-lemma-stalk-gives-point}）是保守的，见
Sites，定义 \ref{sites-definition-enough-points}。特别地，$S_\etale$
有足够多的点。
\item 假设 $\mathcal{F}$ 是 $S$ 的大 \'etale 网站上的层
\label{etale-cohomology-item-stalks-big}
。设 $T \to S$ 是 $S$ 的大 \'etale 网站的对象，并设
$\overline{t}$ 是 $T$ 的几何点。把 $\mathcal{F}_{\overline{t}}$ 定义为
限制 $\mathcal{F}|_{T_\etale}$ 的茎；这是把 $\mathcal{F}$ 限制到
$T$ 的小 \'etale 网站所得的限制。换言之，可定义 $\mathcal{F}$ 在任意概形
$T/S \in \Ob((\Sch/S)_\etale)$ 的任意几何点处的茎。
\item 考虑该网站所有对象的所有几何点，可知 $S$ 的大 \'etale 网站也有足够多的点，
见 (\ref{etale-cohomology-item-stalks-big})。
\end{enumerate}
\end{remarks}

\noindent
初次阅读时应跳过下述引理。

\begin{lemma}
\label{etale-cohomology-lemma-points-small-etale-site}
设 $S$ 是概形。
\begin{enumerate}
\item 设 $p$ 是小 \'etale 网站
$S_\etale$ 的一个点；该网站属于概形 $S$，且该点由函子
$u : S_\etale \to \textit{Sets}$ 给出。则存在一个几何点 $\overline{s}$，
它是 $S$ 的几何点，使得 $p$ 与 $S_\etale$ 的如下点同构：引理
\ref{etale-cohomology-lemma-stalk-gives-point} 中与 $\overline{s}$ 相伴的点。
\item 设 $p : \Sh(pt) \to \Sh(S_\etale)$ 是 $S$ 的小 \'etale 拓扑斯的一个点。
则 $p$ 来自 $S$ 的某个几何点；也就是说，茎函子
$\mathcal{F} \mapsto \mathcal{F}_p$ 与定义
\ref{etale-cohomology-definition-stalk} 所定义的某个茎函子同构。
\end{enumerate}
\end{lemma}

\begin{proof}
由 Sites，引理 \ref{sites-lemma-point-site-topos}，网站的点与相伴拓扑斯的点
一一对应，故只需证明 (1)。由 Sites，命题
\ref{sites-proposition-point-limits}，函子 $u$ 具有下列性质：
(a) $u(S) = \{*\}$；(b) $u(U \times_V W) = u(U) \times_{u(V)} u(W)$；
并且 (c) 若 $\{U_i \to U\}$ 是一个 \'etale 覆盖，则
$\coprod u(U_i) \to u(U)$ 满射。
特别地，若 $U' \subset U$ 是开子概形，则 $u(U') \subset u(U)$。此外，由
Sites，引理 \ref{sites-lemma-point-site-topos}，可写成
$u(U) = p^{-1}(h_U^\#)$；换言之，$u(U)$ 是可表层 $h_U$ 的茎。若
$U = V \amalg W$，则 $h_U = (h_V \amalg h_W)^\#$，并且得到
$u(U) = u(V) \amalg u(W)$，因为 $p^{-1}$ 正合。

\medskip\noindent
考虑 $u$ 到 $S_{Zar}$ 的限制。由 Sites，例
\ref{sites-example-point-topological} 和
\ref{sites-example-point-topology}，存在唯一一点 $s \in S$，使得对开集
$S' \subset S$，$u(S') = \{*\}$ 当且仅当 $s \in S'$，而
$u(S') = \emptyset$ 当且仅当 $s \not \in S'$。注意，若
$\varphi : U \to S$ 是 $S_\etale$ 的对象，则
$\varphi(U) \subset S$ 是开集（见命题
\ref{etale-cohomology-proposition-etale-morphisms}），而
$\{U \to \varphi(U)\}$ 是 \'etale 覆盖。因此可知
$u(U) = \emptyset \Leftrightarrow s \not \in \varphi(U)$。

\medskip\noindent
取一个几何点 $\overline{s} : \overline{s} \to S$，它位于 $s$ 上方；关于记号的惯常滥用，
见定义 \ref{etale-cohomology-definition-geometric-point}。设 $\varphi : U \to S$ 是
$S_\etale$ 的对象，且 $U$ 仿射。注意，$\varphi$ 是分离的，而纤维
$U_s$（即 $\varphi$ 在 $s$ 上的纤维）是 $\Spec(\kappa(s))$ 上的仿射概形；
它是一个有限乘积之谱，其中各因子 $k_i$ 是 $\kappa(s)$ 的有限可分扩张。因此可以应用
\'Etale Morphisms，引理 \ref{etale-lemma-etale-etale-local-technical}，
得到一个 \'etale 邻域 $(V, \overline{v})$，它属于 $(S, \overline{s})$，并使得
$$
U \times_S V = U_1 \amalg \ldots \amalg U_n \amalg W
$$
其中 $U_i \to V$ 是同构，且 $W$ 没有位于 $\overline{v}$ 上方的点。因此可得
$$
u(U) \times u(V) =
u(U \times_S V) =
u(U_1) \amalg \ldots \amalg u(U_n) \amalg u(W)
$$
当然还有 $u(U_i) = u(V)$。将 $V$ 略微缩小后，可假定 $V$ 恰有一个位于
$s$ 上方的点，因而 $W$ 没有位于 $s$ 上方的点。由上述结论，这给出
$u(W) = \emptyset$。于是得到
$$
u(U) \times u(V) =
u(U_1) \amalg \ldots \amalg u(U_n) =
\coprod\nolimits_{i = 1, \ldots, n} u(V)
$$
其中各等号都与到集合 $u(V)$ 的映射相容。注意 $u(V) \not = \emptyset$，
因为 $s$ 属于 $V \to S$ 的像。特别地，在这种情形下 $u(U)$ 是一个含
$n$ 个元素的有限集：左边的纤维为 $u(U)$，它位于 $u(V)$ 的任一元素上方，
而右边的纤维基数为 $n$。

\medskip\noindent
考虑极限
$$
\lim_{(V, \overline{v})} u(V)
$$
这里极限取遍 \'etale 邻域 $(V, \overline{v})$ 所成的范畴；这些邻域属于
$\overline{s}$。
显然，只在 $(V, \overline{v})$ 所成、其中 $V$ 仿射的子范畴上取极限，所得值相同。
由上一段（交换 $V$ 与 $U$ 的角色）可知，在这种情形下 $u(V)$ 总是有限非空集。
此外，该极限是余滤过的，见引理 \ref{etale-cohomology-lemma-cofinal-etale}。因此由 Categories，第
\ref{categories-section-codirected-limits} 节，该极限非空。从中取一个元素 $x$。
这意味着我们得到一个 $x_{V, \overline{v}} \in u(V)$；这里
$(V, \overline{v})$ 可取为 $(S, \overline{s})$ 的任一 \'etale 邻域，并且对
\'etale 邻域的每个态射
$\varphi : (V', \overline{v}') \to (V, \overline{v})$，都有
$u(\varphi)(x_{V', \overline{v}'}) = x_{V, \overline{v}}$。

\medskip\noindent
我们将利用 $x$ 的选取来构造函子性双射
$$
c : |U_{\overline{s}}| \longrightarrow u(U)
$$
（其中 $U \in \Ob(S_\etale)$），这将完成证明。关于 $|U_{\overline{s}}|$ 的描述，
见引理 \ref{etale-cohomology-lemma-stalk-gives-point} 及其证明。首先声称，只需在 $U$ 仿射时构造该映射；
我们略去此断言的证明。设 $U \to S$ 是 $S_\etale$ 中的对象，且 $U$ 仿射；
并设 $\overline{u} : \overline{s} \to U$ 是 $|U_{\overline{s}}|$ 的一个元素。
选取 $(V, \overline{v})$，使 $U \times_S V$ 按证明第三段所述分解。
于是序对 $(\overline{u}, \overline{v})$ 给出 $U \times_S V$ 的一个几何点，
该点位于 $\overline{v}$ 上方，并确定一个分支 $U_i$，它是 $U \times_S V$ 的分支。
更确切地，存在一个截面 $\sigma : V \to U \times_S V$，它是投影
$\text{pr}_U$ 的截面，并使得
$(\overline{u}, \overline{v}) = \sigma \circ \overline{v}$。令
$c(\overline{u}) = u(\text{pr}_U)(u(\sigma)(x_{V, \overline{v}})) \in u(U)$。
需要验证它与 $(V, \overline{v})$ 的选取无关。由引理
\ref{etale-cohomology-lemma-cofinal-etale}，\'etale 邻域所成的范畴是余滤过的。因此只需证明：
给定 \'etale 邻域的态射
$\varphi : (V', \overline{v}') \to (V, \overline{v})$，并选取投影的截面
$\sigma' : V' \to U \times_S V'$，使得
$(\overline{u}, \overline{v'}) = \sigma' \circ \overline{v}'$，则有
$u(\sigma')(x_{V', \overline{v}'}) = u(\sigma)(x_{V, \overline{v}})$。
考虑图表
$$
\xymatrix{
V' \ar[d]^{\sigma'} \ar[r]_\varphi & V \ar[d]^\sigma \\
U \times_S V' \ar[r]^{1 \times \varphi} &
U \times_S V
}
$$
该图表未必交换。这是因为概形 $V'$ 和 $V$ 可能不连通，因而上面用来构造
$\sigma'$ 与 $\sigma$ 的分解可能不唯一。不过由构造确有
$\sigma \circ \varphi \circ \overline{v}' =
(1 \times \varphi) \circ \sigma' \circ \overline{v}'$
。因此，由于 $U \times_S V$ 在 $S$ 上是 \'etale 的，存在
$V'' \subset V'$；这是 $\overline{v'}$ 的一个开邻域，并使得把该图表限制到 $V''$
后它交换，见 Morphisms，引理 \ref{morphisms-lemma-value-at-one-point}。
这意味着可将上图扩充为
$$
\xymatrix{
V'' \ar[r] \ar[d]^{\sigma'|_{V''}} &
V' \ar[d]^{\sigma'} \ar[r]_\varphi &
V \ar[d]^\sigma \\
U \times_S V'' \ar[r] &
U \times_S V' \ar[r]^{1 \times \varphi} &
U \times_S V
}
$$
使得左方块与外围矩形都交换。由于 $u$ 是函子，这蕴含
$x_{V'', \overline{v}'}$ 沿图表中的任一条路径，都映到
$u(U \times_S V)$ 中的同一个元素。另一方面，它分别映到元素
$x_{V', \overline{v}'}$ 与 $x_{V, \overline{v}}$，它们分别属于
$u(V')$ 与 $u(V)$。这就蕴含所需等式
$u(\sigma')(x_{V', \overline{v}'}) = u(\sigma)(x_{V, \overline{v}})$。

\medskip\noindent
类似地可证明，构造 $c : |U_{\overline{s}}| \to u(U)$ 关于 $U$ 是函子性的；
细节略去。最后，由第三段的结果，映射 $c$ 显然为双射。这就完成引理的证明。
\end{proof}







\section{其他拓扑中的点}
\label{etale-cohomology-section-points-topologies}

\noindent
本节简要讨论概形的 \'etale 网站以外若干网站的点是否存在。关于本节所用术语，参见
Sites，第 \ref{sites-section-sites-enough-points} 节，以及
Topologies，第 \ref{topologies-section-procedure} 节及后续各节。
所有几何网站都有足够多的点。

\begin{lemma}
\label{etale-cohomology-lemma-points-fppf}
设 $S$ 是概形。下列所有网站都有足够多的点：
$S_{affine, Zar}$, $S_{Zar}$, $S_{affine, \etale}$, $S_\etale$,
$(\Sch/S)_{Zar}$, $(\textit{Aff}/S)_{Zar}$,
$(\Sch/S)_\etale$, $(\textit{Aff}/S)_\etale$,
$(\Sch/S)_{smooth}$, $(\textit{Aff}/S)_{smooth}$,
$(\Sch/S)_{syntomic}$, $(\textit{Aff}/S)_{syntomic}$,
$(\Sch/S)_{fppf}$，以及 $(\textit{Aff}/S)_{fppf}$。
\end{lemma}

\begin{proof}
对每个大网站，其相伴拓扑斯等价于网站 $(\textit{Aff}/S)_\tau$ 所定义的拓扑斯，见
Topologies，引理 \ref{topologies-lemma-affine-big-site-Zariski}、
\ref{topologies-lemma-affine-big-site-etale},
\ref{topologies-lemma-affine-big-site-smooth},
\ref{topologies-lemma-affine-big-site-syntomic} 和
\ref{topologies-lemma-affine-big-site-fppf}。网站 $(\textit{Aff}/S)_\tau$
的结论立即由 Deligne 的结果 Sites，引理 \ref{sites-lemma-criterion-points} 得出。

\medskip\noindent
关于 $S_{Zar}$ 的结论是显然的。关于 $S_{affine, Zar}$ 的结论由 Deligne 的结果得出。
关于 $S_\etale$ 的结论，或者由定理 \ref{etale-cohomology-theorem-exactness-stalks}（的证明）得出，
或者由 Topologies，引理 \ref{topologies-lemma-alternative} 以及将 Deligne 的结果应用于
$S_{affine, \etale}$ 得出。
\end{proof}

\noindent
上述引理保证点的存在性，却没有说明这些点的形貌。可以如下显式构造{\it 一些}点。
设 $\overline{s} : \Spec(k) \to S$ 是几何点，且 $k$ 代数闭。考虑函子
$$
u : (\Sch/S)_{fppf} \longrightarrow \textit{Sets},
\quad
u(U) = U(k) = \Mor_S(\Spec(k), U).
$$
注意 $U \mapsto U(k)$ 与有限极限交换，因为 $S(k) = \{\overline{s}\}$，且
$(U_1 \times_U U_2)(k) = U_1(k) \times_{U(k)} U_2(k)$.
此外，若 $\{U_i \to U\}$ 是 fppf 覆盖，则 $\coprod U_i(k) \to U(k)$ 满射。
由 Sites，命题 \ref{sites-proposition-point-limits}，可知 $u$ 定义一个点 $p$；
这是 $(\Sch/S)_{fppf}$ 的点，其茎为
$$
\mathcal{F}_p = \colim_{(U, x)} \mathcal{F}(U)
$$
其中照例对序对 $U \to S$、$x \in U(k)$ 取余极限。然而，这个范畴有一个始对象，
即 $(\Spec(k), \text{id})$，因此可知
$$
\mathcal{F}_p = \mathcal{F}(\Spec(k))
$$
，这并没有多大意思！事实上，一般而言，这些点不会组成保守点族。下述评注描述一类更有意思的点。

\begin{remark}
\label{etale-cohomology-remark-points-fppf-site}
\begin{reference}
这一点在 \cite{Schroeer} 中讨论。
\end{reference}
设 $S = \Spec(A)$ 是仿射概形。设 $(p, u)$ 是网站
$(\textit{Aff}/S)_{fppf}$ 的一个点，见 Sites，第 \ref{sites-section-points} 节和
第 \ref{sites-section-construct-points} 节。令 $B = \mathcal{O}_p$ 为结构层在点
$p$ 处的茎。回顾
$$
B = \colim_{(U, x)} \mathcal{O}(U) =
\colim_{(\Spec(C), x_C)} C
$$
，其中 $x_C \in u(\Spec(C))$。可能发生如下情形：$\Spec(B)$ 是
$(\textit{Aff}/S)_{fppf}$ 的对象，并且存在元素 $x_B \in u(\Spec(B))$，
它映到相容族 $x_C$。此时邻域系有始对象，因而有
$\mathcal{F}_p = \mathcal{F}(\Spec(B))$；其中 $\mathcal{F}$ 是
$(\textit{Aff}/S)_{fppf}$ 上的任意层。不难看出，若
$\mathcal{F} \mapsto \mathcal{F}(\Spec(B))$ 定义
$\Sh((\textit{Aff}/S)_{fppf})$ 的一个点，则 $B$ 必须是局部 $A$-代数，
使得每个忠实平坦、有限表示的环映射 $B \to B'$ 都有一个截面 $B' \to B$。
反之，对任意这样的 $A$-代数 $B$，函子
$\mathcal{F} \mapsto \mathcal{F}(\Spec(B))$ 是某个点的茎函子。细节略去。
网站 $(\textit{Aff}/S)_{fppf}$ 的一般点是什么样的，尚不清楚。
\end{remark}











\section{阿贝尔层的支集}
\label{etale-cohomology-section-support}

\noindent
先讨论局部截面的支集。

\begin{lemma}
\label{etale-cohomology-lemma-support-subsheaf-final}
设 $S$ 是概形。设 $\mathcal{F}$ 是 $S$ 的 \'etale 拓扑斯之终对象的子层（见
Sites，例 \ref{sites-example-singleton-sheaf}）。则存在唯一开集
$W \subset S$，使得 $\mathcal{F} = h_W$。
\end{lemma}

\begin{proof}
该条件意味着 $\mathcal{F}(U)$ 是单点集或空集；这里
$\varphi : U \to S$ 遍历 $\Ob(S_\etale)$。特别地，局部截面总能黏合。若
$\mathcal{F}(U) \not = \emptyset$，则 $\mathcal{F}(\varphi(U)) \not = \emptyset$，
因为 $\{\varphi : U \to \varphi(U)\}$ 是覆盖。因此可取
$W = \bigcup_{\varphi : U \to S, \mathcal{F}(U) \not = \emptyset} \varphi(U)$.
\end{proof}

\begin{lemma}
\label{etale-cohomology-lemma-zero-over-image}
设 $S$ 是概形。设 $\mathcal{F}$ 是 $S_\etale$ 上的阿贝尔层。
设 $\sigma \in \mathcal{F}(U)$ 是局部截面。存在开子集 $W \subset U$，使得
\begin{enumerate}
\item $W \subset U$ 是 $U$ 中使 $\sigma|_W = 0$ 的最大 Zariski 开子集；
\item 对每个 $\varphi : V \to U$（它属于 $S_\etale$），有
$$
\sigma|_V = 0 \Leftrightarrow \varphi(V) \subset W,
$$
\item 对每个几何点 $\overline{u}$（它是 $U$ 的几何点），有
$$
(U, \overline{u}, \sigma) = 0\text{ 于 }\mathcal{F}_{\overline{s}}
\Leftrightarrow
\overline{u} \in W
$$
，其中 $\overline{s} = (U \to S) \circ \overline{u}$。
\end{enumerate}
\end{lemma}

\begin{proof}
由于 $\mathcal{F}$ 是 \'etale 拓扑中的层，$\mathcal{F}$ 在 $U_{Zar}$ 上的限制是 $U$ 的
Zariski 拓扑中的层。因此存在满足 (1) 的 Zariski 开集 $W$，见 Modules，引理
\ref{modules-lemma-support-section-closed}。设 $\varphi : V \to U$ 是
$S_\etale$ 中的箭头。注意，$\varphi(V) \subset U$ 是开子集，并且
$\{V \to \varphi(V)\}$ 是 \'etale 覆盖。因此，若 $\sigma|_V = 0$，则由
$\mathcal{F}$ 的层条件可知 $\sigma|_{\varphi(V)} = 0$。这证明了 (2)。
为证明 (3)，需要证明：若 $(U, \overline{u}, \sigma)$ 定义
$\mathcal{F}_{\overline{s}}$ 的零元素，则 $\overline{u} \in W$。这是因为该假设意味着
存在 \'etale 邻域的一个态射 $(V, \overline{v}) \to (U, \overline{u})$，使得
$\sigma|_V = 0$。故由 (2) 可知 $V \to U$ 映入 $W$，从而
$\overline{u} \in W$。
\end{proof}

\noindent
设 $S$ 是概形，并设 $s \in S$。设 $\mathcal{F}$ 是 $S_\etale$ 上的层。
由评注 \ref{etale-cohomology-remark-map-stalks}，层 $\mathcal{F}$ 在位于 $s$ 上方的几何点处之茎的
同构类是良定义的。

\begin{definition}
\label{etale-cohomology-definition-support}
设 $S$ 是概形。设 $\mathcal{F}$ 是 $S_\etale$ 上的阿贝尔层。
\begin{enumerate}
\item {\it $\mathcal{F}$ 的支集}是所有点 $s \in S$ 组成的集合；它们满足
$\mathcal{F}_{\overline{s}} \not = 0$，其中 $\overline{s}$ 是位于 $s$ 上方的
任一（等价地，某个）几何点。
\item 设 $\sigma \in \mathcal{F}(U)$ 是截面。{\it $\sigma$ 的支集}是闭子集
$U \setminus W$，其中 $W \subset U$ 是 $U$ 中使 $\sigma$ 限制为零的最大开子集
（见引理 \ref{etale-cohomology-lemma-zero-over-image}）。
\end{enumerate}
\end{definition}

\noindent
一般而言，阿贝尔层的支集不是闭集。例如，设 $S = \mathbf{A}^1_{\mathbf{C}}$。
令 $i_t : \Spec(\mathbf{C}) \to S$ 是点 $t \in \mathbf{C}$ 的嵌入。
稍后将看到，$\mathcal{F}_t = i_{t, *}(\mathbf{Z}/2\mathbf{Z})$ 是支集恰为
$\{t\}$ 的阿贝尔层，见第 \ref{etale-cohomology-section-closed-immersions} 节。于是
$$
\bigoplus\nolimits_{n \in \mathbf{N}} \mathcal{F}_n
$$
是支集为 $\{1, 2, 3, \ldots\} \subset S$ 的阿贝尔层。这是因为取茎与余极限交换，
见引理 \ref{etale-cohomology-lemma-stalk-exact}。这就给出了支集不闭的阿贝尔层之例。
下面是关于层与截面的支集的一些基本事实。

\begin{lemma}
\label{etale-cohomology-lemma-support-section-closed}
设 $S$ 是概形。设 $\mathcal{F}$ 是 $S_\etale$ 上的阿贝尔层。
设 $U \in \Ob(S_\etale)$ 且 $\sigma \in \mathcal{F}(U)$。
\begin{enumerate}
\item $\sigma$ 的支集在 $U$ 中是闭的。
\item $\sigma + \sigma'$ 的支集包含于 $\sigma, \sigma' \in \mathcal{F}(U)$
二者支集的并。
\item 若 $\varphi : \mathcal{F} \to \mathcal{G}$ 是 $S_\etale$ 上阿贝尔层的映射，
则 $\varphi(\sigma)$ 的支集包含于 $\sigma \in \mathcal{F}(U)$ 的支集。
\item $\mathcal{F}$ 的支集是 $\mathcal{F}$ 的所有局部截面之支集的像的并。
\item 若 $\mathcal{F} \to \mathcal{G}$ 满射，则 $\mathcal{G}$ 的支集是
$\mathcal{F}$ 的支集的子集。
\item 若 $\mathcal{F} \to \mathcal{G}$ 单射，则 $\mathcal{F}$ 的支集是
$\mathcal{G}$ 的支集的子集。
\end{enumerate}
\end{lemma}

\begin{proof}
(1) 由定义成立。(2) 与 (3) 成立，因为对 $\mathcal{F}$ 与 $\mathcal{G}$ 在
$U_{Zar}$ 上的限制，相应结论成立，见 Modules，引理
\ref{modules-lemma-support-section-closed}。(4) 是引理
\ref{etale-cohomology-lemma-zero-over-image} 的 (3) 的直接推论。(5) 与 (6) 由其他各项得出。
\end{proof}

\begin{lemma}
\label{etale-cohomology-lemma-support-sheaf-rings-closed}
$S_\etale$ 上环层的支集是闭集。
\end{lemma}

\begin{proof}
这是因为（按照我们的约定）环是 $0$ 当且仅当 $1 = 0$，所以环层的支集就是单位截面的支集。
\end{proof}




\section{Hensel 环}
\label{etale-cohomology-section-henselian-ring}

\noindent
先陈述一个在 Stacks project 中已经多次使用的定理。该结果有许多版本；这里仅陈述其代数版本。

\begin{theorem}
\label{etale-cohomology-theorem-quasi-finite-etale-locally}
设 $A\to B$ 是有限型环映射，且 $\mathfrak p \subset A$ 是素理想。则存在
\'etale 环映射 $A \to A'$ 和素理想 $\mathfrak p' \subset A'$；后者位于
$\mathfrak p$ 上方，并且
\begin{enumerate}
\item
$\kappa(\mathfrak p) = \kappa(\mathfrak p')$,
\item
$ B \otimes_A A' = B_1\times \ldots \times B_r \times C$,
\item
$ A'\to B_i$ 是有限的，并且存在唯一素理想 $q_i\subset B_i$ 位于
$\mathfrak p'$ 上方；以及
\item 纤维 $\Spec(C \otimes_{A'} \kappa(\mathfrak p'))$ 的所有不可约分支的维数都至少为 1；
这是 $C$ 在 $\mathfrak p'$ 上的纤维。
\end{enumerate}
\end{theorem}

\begin{proof}
见 Algebra，引理 \ref{algebra-lemma-etale-makes-quasi-finite-finite}，或见
\cite[Th\'eor\`eme 18.12.1]{EGA4}。关于用概形态射表述的一系列版本，见
More on Morphisms，第 \ref{more-morphisms-section-etale-localization} 节。
\end{proof}

\noindent
回顾 Hensel 引理。该引理有许多版本。下面列出两个：
\begin{enumerate}
\item[(f)] 若 $f\in \mathbf{Z}_p[T]$ 首一，且
$f \bmod p = g_0 h_0$，其中 $gcd(g_0, h_0) = 1$，则 $f$ 可分解为
$f = gh$，其中 $\bar g = g_0$ 且 $\bar h = h_0$；
\item[(r)] 若 $f \in \mathbf{Z}_p[T]$ 首一，且 $a_0 \in \mathbf{F}_p$ 满足
$\bar f(a_0) =0$ 而 $\bar f'(a_0) \neq 0$，则存在 $a \in \mathbf{Z}_p$，使得
$f(a) = 0$ 且 $\bar a = a_0$。
\end{enumerate}
两个版本都成立（稍后会看到这一点）。第一个版本要求把互素因子分解提升，
第二个版本要求把模极大理想的单根提升。事实证明，对一般局部环要求这两个条件是等价的，
并且还等价于许多其他条件。我们用根提升性质来定义 Hensel 局部环，因为它通常最容易检验。

%10.01.09
\begin{definition}
\label{etale-cohomology-definition-henselian}
（见 Algebra，定义 \ref{algebra-definition-henselian}。）称局部环
$(R, \mathfrak m, \kappa)$ 是{\it Hensel 的}，如果对所有首一
$f \in R[T]$ 以及所有 $a_0 \in \kappa$，只要
$\bar f(a_0) = 0$ 且 $\bar f'(a_0) \neq 0$，就存在 $a \in R$，使得
$f(a) = 0$ 且 $a \bmod \mathfrak m = a_0$。
\end{definition}

\noindent
应当记住的 Hensel 局部环的一个好例子是完备局部环。回顾（Algebra，定义
\ref{algebra-definition-complete-local-ring}），完备局部环是局部环
$(R, \mathfrak m)$，它满足 $R \cong \lim_n R/\mathfrak m^n$；换言之，
它关于 $\mathfrak m$-进拓扑是完备且分离的。

\begin{theorem}
\label{etale-cohomology-theorem-hensel}
完备局部环是 Hensel 的。
\end{theorem}

\begin{proof}
使用 Newton 方法。见 Algebra，引理 \ref{algebra-lemma-complete-henselian}。
\end{proof}

\begin{theorem}
\label{etale-cohomology-theorem-henselian}
设 $(R, \mathfrak m, \kappa)$ 是局部环。下列条件等价：
\begin{enumerate}
\item $R$ 是 Hensel 的；
\item 对任意 $f\in R[T]$ 以及任意分解 $\bar f = g_0 h_0$（它在
$\kappa[T]$ 中且满足 $\gcd(g_0, h_0)=1$），存在分解 $f = gh$（它在
$R[T]$ 中），并满足 $\bar g = g_0$ 且 $\bar h = h_0$；
\item 任意有限 $R$-代数 $S$ 都同构于若干在 $R$ 上有限的局部环的有限乘积；
\item 任意有限型 $R$-代数 $A$ 都同构于一个乘积
$A \cong A' \times C$，其中 $A' \cong A_1 \times \ldots \times A_r$
是有限局部 $R$-代数的乘积，并且 $C \otimes_R \kappa$ 的所有不可约分支的维数至少为 1；
\item 若 $A$ 是 \'etale $R$-代数，且 $\mathfrak n$ 是 $A$ 的一个位于
$\mathfrak m$ 上方的极大理想，满足 $\kappa \cong A/\mathfrak n$，则存在同构
$\varphi : A \cong R \times A'$，使得
$\varphi(\mathfrak n) = \mathfrak m \times A' \subset R \times A'$。
\end{enumerate}
\end{theorem}

\begin{proof}
这只是 Algebra，引理 \ref{algebra-lemma-characterize-henselian} 中若干结果的子集。
注意，上面的 (5) 对应于 Algebra，引理
\ref{algebra-lemma-characterize-henselian} 的 (8)，但表述略有不同。
\end{proof}

\begin{lemma}
\label{etale-cohomology-lemma-finite-over-henselian}
若 $R$ 是 Hensel 的，且 $A$ 是有限 $R$-代数，则 $A$ 是 Hensel 局部环的有限乘积。
\end{lemma}

\begin{proof}
见 Algebra，引理 \ref{algebra-lemma-finite-over-henselian}。
\end{proof}

\begin{definition}
\label{etale-cohomology-definition-strictly-henselian}
若局部环 $R$ 是 Hensel 的，且其剩余域可分闭，则称其为{\it 严格 Hensel 的}。
\end{definition}

\begin{example}
\label{etale-cohomology-example-powerseries}
在 $R = \mathbf{C}[[t]]$ 的情形，\'etale $R$-代数是平凡扩张 $R \to R$ 与扩张
$R \to R[X, X^{-1}]/(X^n-t)$ 的有限乘积。后一类扩张经开集
$D(t) \subset \Spec(R)$ 分解，故任意 \'etale 覆盖都可由覆盖
$\{\text{id} : \Spec(R) \to \Spec(R)\}$ 加细。下面将看到，对 Hensel 环之谱的
\'etale 覆盖而言，这是一个颇为一般的事实。由此将说明，严格 Hensel 环之谱的高阶
\'etale 上同调为零。
\end{example}

\begin{theorem}
\label{etale-cohomology-theorem-henselization}
设 $(R, \mathfrak m, \kappa)$ 是局部环，且
$\kappa\subset\kappa^{sep}$ 是可分代数闭包。存在典范平坦局部环映射
$R \to R^h \to R^{sh}$，其中
\begin{enumerate}
\item $R^h$、$R^{sh}$ 是 \'etale $R$-代数的滤过余极限；
\item $R^h$ 是 Hensel 的，$R^{sh}$ 是严格 Hensel 的；
\item $\mathfrak m R^h$（相应地 $\mathfrak m R^{sh}$）是 $R^h$（相应地
$R^{sh}$）的极大理想；以及
\item $\kappa = R^h/\mathfrak m R^h$，且
$\kappa^{sep} = R^{sh}/\mathfrak m R^{sh}$；这些等式按 $\kappa$ 的扩张来理解。
\end{enumerate}
\end{theorem}

\begin{proof}
$R^h$ 与 $R^{sh}$ 的结构见 Algebra，引理 \ref{algebra-lemma-henselization} 和
\ref{algebra-lemma-strict-henselization}。
\end{proof}

\noindent
定理 \ref{etale-cohomology-theorem-henselization} 中构造的环分别称为局部环 $R$ 的
{\it Hensel 化}和{\it 严格 Hensel 化}，见 Algebra，定义
\ref{algebra-definition-henselization}。$R$ 的许多性质会反映在其（严格）Hensel 化中，
见 More on Algebra，第 \ref{more-algebra-section-permanence-henselization} 节。




\section{结构层的茎}
\label{etale-cohomology-section-stalks-structure-sheaf}

\noindent
本节把结构层在几何点处的茎与相应“通常”点处局部环的严格 Hensel 化等同起来。

\begin{lemma}
\label{etale-cohomology-lemma-describe-etale-local-ring}
\begin{slogan}
概形结构层在 etale 拓扑中的茎是严格 Hensel 化。
\end{slogan}
设 $S$ 是概形。设 $\overline{s}$ 是 $S$ 的几何点，位于 $s \in S$ 上方。
令 $\kappa = \kappa(s)$，并令
$\kappa \subset \kappa^{sep} \subset \kappa(\overline{s})$ 表示
$\kappa$ 在 $\kappa(\overline{s})$ 中的可分代数闭包。则有典范等同
$$
(\mathcal{O}_{S, s})^{sh}
\cong
(\mathcal{O}_S)_{\overline{s}}
$$
，其中左边是定理 \ref{etale-cohomology-theorem-henselization} 所述局部环
$\mathcal{O}_{S, s}$ 的严格 Hensel 化，右边是结构层 $\mathcal{O}_S$
（定义在 $S_\etale$ 上）在几何点 $\overline{s}$ 处的茎。
\end{lemma}

\begin{proof}
设 $\Spec(A) \subset S$ 是 $s$ 的仿射邻域。设 $\mathfrak p \subset A$
是对应于 $s$ 的素理想。作出这些选择后，有典范同构
$\mathcal{O}_{S, s} = A_{\mathfrak p}$ 和 $\kappa(s) = \kappa(\mathfrak p)$。
因而有
$\kappa(\mathfrak p) \subset \kappa^{sep} \subset \kappa(\overline{s})$.
回顾
$$
(\mathcal{O}_S)_{\overline{s}} =
\colim_{(U, \overline{u})} \mathcal{O}(U)
$$
，其中极限取遍 $(S, \overline{s})$ 的 \'etale 邻域。一个余共尾系由如下
\'etale 邻域 $(U, \overline{u})$ 给出：$U$ 仿射，且 $U \to S$ 经
$\Spec(A)$ 分解。换言之，可知
$$
(\mathcal{O}_S)_{\overline{s}} = \colim_{(B, \mathfrak q, \phi)} B
$$
，其中余极限取遍 \'etale $A$-代数 $B$；它们配备素理想 $\mathfrak q$，
该理想位于 $\mathfrak p$ 上方，并配备一个 $\kappa(\mathfrak p)$-代数映射
$\phi : \kappa(\mathfrak q) \to \kappa(\overline{s})$。注意，由于
$\kappa(\mathfrak q)$ 在 $\kappa(\mathfrak p)$ 上有限可分，$\phi$ 的像包含于
$\kappa^{sep}$。经这些转译，本引理的结论等价于 Algebra，引理
\ref{algebra-lemma-strict-henselization-different} 的结论。
\end{proof}

\begin{definition}
\label{etale-cohomology-definition-etale-local-rings}
设 $S$ 是概形。设 $\overline{s}$ 是 $S$ 的几何点，位于点 $s \in S$ 上方。
\begin{enumerate}
\item {\it $S$ 在 $\overline{s}$ 处的 \'etale 局部环}是结构层 $\mathcal{O}_S$
（定义在 $S_\etale$ 上）在 $\overline{s}$ 处的茎。有时称它为
{\it $\mathcal{O}_{S, s}$ 的严格 Hensel 化}，这里是相对于几何点
$\overline{s}$ 而言。
所用记号为 $\mathcal{O}_{S, \overline{s}}^{sh}$。
\item {\it $\mathcal{O}_{S, s}$ 的 Hensel 化}是 $S$ 在 $s$ 处的局部环的 Hensel 化。
见 Algebra，定义 \ref{algebra-definition-henselization}，以及定理
\ref{etale-cohomology-theorem-henselization}。记为 $\mathcal{O}_{S, s}^h$。
\item {\it $S$ 在 $\overline{s}$ 处的严格 Hensel 化}是概形
$\Spec(\mathcal{O}_{S, \overline{s}}^{sh})$。
\item {\it $S$ 在 $s$ 处的 Hensel 化}是概形 $\Spec(\mathcal{O}_{S, s}^h)$。
\end{enumerate}
\end{definition}

\noindent
设 $f : T \to S$ 是概形态射。设 $\overline{t}$ 是 $T$ 的几何点，其像为
$\overline{s}$，后者位于 $S$ 中。设 $t \in T$ 和 $s \in S$ 为它们的像。则得到典范交换图
$$
\xymatrix{
\Spec(\mathcal{O}^{sh}_{T, \overline{t}}) \ar[r] \ar[d] &
\Spec(\mathcal{O}^h_{T, t}) \ar[r] \ar[d] &
T \ar[d]^f \\
\Spec(\mathcal{O}^{sh}_{S, \overline{s}}) \ar[r] &
\Spec(\mathcal{O}^h_{S, s}) \ar[r] &
S
}
$$
；这是 $T$ 与 $S$ 的 Hensel 化和严格 Hensel 化所成的图。可通过选取 $t$ 和 $s$
的仿射邻域，并使用 Algebra，引理
\ref{algebra-lemma-henselian-functorial-improve} 和
\ref{algebra-lemma-strictly-henselian-functorial-improve} 给出的（严格）Hensel 化的函子性来证明。

\begin{lemma}
\label{etale-cohomology-lemma-describe-henselization}
设 $S$ 是概形。设 $s \in S$。则有
$$
\mathcal{O}_{S, s}^h =
\colim_{(U, u)} \mathcal{O}(U)
$$
，其中余极限取遍滤过范畴中的 \'etale 邻域 $(U, u)$；它们是
$(S, s)$ 的邻域，并满足 $\kappa(s) = \kappa(u)$。
\end{lemma}

\begin{proof}
本引理是 More on Morphisms，引理
\ref{more-morphisms-lemma-describe-henselization} 的一个副本。
\end{proof}

\begin{remark}
\label{etale-cohomology-remark-henselization-Noetherian}
设 $S$ 是概形。设 $s \in S$。若 $S$ 局部 Noether，则
$\mathcal{O}_{S, s}^h$ 也 Noether，并且二者有相同的完备化：
$$
\widehat{\mathcal{O}_{S, s}} \cong \widehat{\mathcal{O}_{S, s}^h}.
$$
特别地，
$\mathcal{O}_{S, s} \subset
\mathcal{O}_{S, s}^h \subset
\widehat{\mathcal{O}_{S, s}}$.
$\mathcal{O}_{S, s}$ 的 Hensel 化一般远小于它的完备化，并继承它的许多性质。
例如，若 $\mathcal{O}_{S, s}$ 既约，则 $\mathcal{O}_{S, s}^h$ 也既约；
但对完备化而言，一般并非如此。此处待补充未来的参考文献。
\end{remark}

\begin{lemma}
\label{etale-cohomology-lemma-etale-site-locally-ringed}
设 $S$ 是概形。小 \'etale 网站 $S_\etale$ 配备其结构层 $\mathcal{O}_S$ 后，
是局部环化网站，见 Modules on Sites，定义
\ref{sites-modules-definition-locally-ringed}。
\end{lemma}

\begin{proof}
这是因为各茎 $(\mathcal{O}_S)_{\overline{s}} = \mathcal{O}^{sh}_{S, \overline{s}}$
都是局部环，且 $S_\etale$ 有足够多的点；见引理
\ref{etale-cohomology-lemma-describe-etale-local-ring}、定理 \ref{etale-cohomology-theorem-exactness-stalks} 和评注
\ref{etale-cohomology-remarks-enough-points}。关于这些事实蕴含小 \'etale 网站是局部环化网站，见
Modules on Sites，引理 \ref{sites-modules-lemma-locally-ringed-stalk} 和
\ref{sites-modules-lemma-ringed-stalk-not-zero}。
\end{proof}






\section{小 \'etale 拓扑斯的函子性}
\label{etale-cohomology-section-functoriality}

\noindent
到目前为止，我们尚未讨论 \'etale 网站的函子性，换言之，给定概形态射时会发生什么。
精确的形式化讨论见 Topologies，第 \ref{topologies-section-etale} 节。
本节和随后几节将在小 \'etale 网站的框架下简要讨论这些内容。

\medskip\noindent
设 $f : X \to Y$ 是概形态射。得到函子
\begin{equation}
\label{etale-cohomology-equation-functorial}
u : Y_\etale \longrightarrow X_\etale, \quad
V/Y \longmapsto X \times_Y V/X.
\end{equation}
该函子具有下列重要性质：
\begin{enumerate}
\item $u(\text{终对象}) = \text{终对象}$；
\item $u$ 保持纤维积；
\item 若 $\{V_j \to V\}$ 是 $Y_\etale$ 中的覆盖，则
$\{u(V_j) \to u(V)\}$ 是 $X_\etale$ 中的覆盖。
\end{enumerate}
每一项都容易检验（略去）。因此得到所谓的{\it 网站态射}
$$
f_{small} : X_\etale \longrightarrow Y_\etale,
$$
，见 Sites，定义 \ref{sites-definition-morphism-sites} 和 Sites，命题
\ref{sites-proposition-get-morphism}。处理 \'etale 层与 \'etale 上同调时，
不必详细掌握这个抽象概念。通常只需知道 \'etale 层上存在函子
$f_{small, *}$（推前）和 $f_{small}^{-1}$（拉回），并了解它们的一些简单性质。
随后几节将讨论这些性质；不过有时我们会引用较抽象的材料来作证明，因为那里往往是证明它们的自然框架。


\section{正像}
\label{etale-cohomology-section-direct-image}

\noindent
定义预层的推前。

\begin{definition}
\label{etale-cohomology-definition-direct-image-presheaf}
设 $f: X\to Y$ 是概形态射。设 $\mathcal{F} $ 是 $X_\etale$ 上的集合预层。
$\mathcal{F}$（在 $f$ 下）的{\it 正像}，或称{\it 推前}，是
$$
f_*\mathcal{F} : Y_\etale^{opp} \longrightarrow \textit{Sets}, \quad
(V/Y) \longmapsto \mathcal{F}(X \times_Y V/X).
$$
有时写作 $f_* = f_{small, *}$，以区别于其他正像函子（例如通常的 Zariski 推前或
$f_{big, *}$）。
\end{definition}

\noindent
这是良定义的 \'etale 预层，因为 \'etale 态射的基变换仍是 \'etale 的。
更范畴化地说，$f_*\mathcal{F}$ 是函子的复合 $\mathcal{F} \circ u$，其中
$u$ 如方程 (\ref{etale-cohomology-equation-functorial}) 所示。这清楚地说明，该构造关于预层
$\mathcal{F}$ 是函子性的，故得到函子
$$
f_* = f_{small, *} :
\textit{PSh}(X_\etale)
\longrightarrow
\textit{PSh}(Y_\etale)
$$
。注意，若 $\mathcal{F}$ 是阿贝尔群预层，则 $f_*\mathcal{F}$ 也是阿贝尔群预层，
并得到
$$
f_* = f_{small, *} :
\textit{PAb}(X_\etale)
\longrightarrow
\textit{PAb}(Y_\etale)
$$
，与前面相同（即由完全相同的规则定义）。

\begin{remark}
\label{etale-cohomology-remark-direct-image-sheaf}
我们断言，层的正像仍是层。事实上，若 $\{V_j \to V\}$ 是 $Y_\etale$ 中的
\'etale 覆盖，则 $\{X \times_Y V_j \to X \times_Y V\}$ 是
$X_\etale$ 中的 \'etale 覆盖。因此，$\mathcal{F}$ 关于
$\{X \times_Y V_i \to X \times_Y V\}$ 的层条件，等价于
$f_*\mathcal{F}$ 关于 $\{V_i \to V\}$ 的层条件。故若 $\mathcal{F}$ 是层，
则 $f_*\mathcal{F}$ 也是层。
\end{remark}

\begin{definition}
\label{etale-cohomology-definition-direct-image-sheaf}
设 $f: X\to Y$ 是概形态射。设 $\mathcal{F} $ 是 $X_\etale$ 上的集合层。
$\mathcal{F}$（在 $f$ 下）的{\it 正像}，或称{\it 推前}，是
$$
f_*\mathcal{F} : Y_\etale^{opp} \longrightarrow \textit{Sets}, \quad
(V/Y) \longmapsto \mathcal{F}(X \times_Y V/X)
$$
；由评注 \ref{etale-cohomology-remark-direct-image-sheaf}，这是层。有时写作
$f_* = f_{small, *}$，以区别于其他正像函子（例如通常的 Zariski 推前或
$f_{big, *}$）。
\end{definition}

\noindent
完全相同的讨论仍然适用，并得到函子
$$
f_* = f_{small, *} :
\Sh(X_\etale)
\longrightarrow
\Sh(Y_\etale)
$$
以及
$$
f_* = f_{small, *} :
\textit{Ab}(X_\etale)
\longrightarrow
\textit{Ab}(Y_\etale)
$$
，仍称为{\it 正像}。

\medskip\noindent
阿贝尔层上的函子 $f_*$ 是左正合的。（关于阿贝尔范畴之间的函子为左正合是什么意思，
见 Homology，第 \ref{homology-section-functors} 节。）事实上，若
$0 \to \mathcal{F}_1 \to \mathcal{F}_2 \to \mathcal{F}_3$
在 $X_\etale$ 上正合，则对每个
$U/X \in \Ob(X_\etale)$
，阿贝尔群序列
$0 \to \mathcal{F}_1(U) \to \mathcal{F}_2(U) \to \mathcal{F}_3(U)$
正合。因此，对每个 $V/Y \in \Ob(Y_\etale)$，阿贝尔群序列
$0 \to f_*\mathcal{F}_1(V) \to f_*\mathcal{F}_2(V) \to f_*\mathcal{F}_3(V)$
正合，因为这就是前一个序列取 $U = X \times_Y V$ 所得的序列。

\begin{definition}
\label{etale-cohomology-definition-higher-direct-images}
设 $f: X \to Y$ 是概形态射。右导出函子 $\{R^pf_*\}_{p \geq 1}$（它们来自函子
$f_* : \textit{Ab}(X_\etale) \to \textit{Ab}(Y_\etale)$）称为{\it 高阶正像}。
\end{definition}

\noindent
关于高阶正像及其导出范畴版本，将在（此处待补充未来的参考文献）中作更详细的讨论。



\section{逆像}
\label{etale-cohomology-section-inverse-image}

\noindent
本节简要讨论小 \'etale 网站上层的拉回。其精确构造见 Topologies，第
\ref{topologies-section-etale} 节。

\begin{definition}
\label{etale-cohomology-definition-inverse-image}
设 $f: X\to Y$ 是概形态射。{\it 逆像}，或称{\it 拉回}\footnote{对于集合层或
阿贝尔群层的拉回，我们使用记号 $f^{-1}$；而把 $f^*$ 留给通过环化网站/拓扑斯态射所得的
模层拉回。}函子是
$$
f^{-1} = f_{small}^{-1} :
\Sh(Y_\etale)
\longrightarrow
\Sh(X_\etale)
$$
以及
$$
f^{-1} = f_{small}^{-1} :
\textit{Ab}(Y_\etale)
\longrightarrow
\textit{Ab}(X_\etale)
$$
；它们左伴随于 $f_* = f_{small, *}$。因此，$f^{-1}$ 由下列事实刻画：
$$
\Hom_{{\Sh(X_\etale)}} (f^{-1}\mathcal{G}, \mathcal{F})
=
\Hom_{\Sh(Y_\etale)} (\mathcal{G}, f_*\mathcal{F})
$$
对任意 $\mathcal{F} \in \Sh(X_\etale)$ 和
$\mathcal{G} \in \Sh(Y_\etale)$ 函子性地成立。类似地，有
$$
\Hom_{{\textit{Ab}(X_\etale)}} (f^{-1}\mathcal{G}, \mathcal{F})
=
\Hom_{\textit{Ab}(Y_\etale)} (\mathcal{G}, f_*\mathcal{F})
$$
，其中 $\mathcal{F} \in \textit{Ab}(X_\etale)$ 且
$\mathcal{G} \in \textit{Ab}(Y_\etale)$。
\end{definition}

\noindent
这种伴随是否存在并不显然。另一方面，它在相当一般的框架中存在，见下面的评注
\ref{etale-cohomology-remark-functoriality-general}。一般机制说明，$f^{-1}\mathcal{G}$ 是下述预层的相伴层：
\begin{equation}
\label{etale-cohomology-equation-pullback}
U/X
\longmapsto
\colim_{U \to X \times_Y V} \mathcal{G}(V/Y)
\end{equation}
，其中余极限取遍序对
$(V/Y, \varphi : U/X \to X \times_Y V/X)$ 所成的范畴。为看出这一点，
将 Sites，命题 \ref{sites-proposition-get-morphism} 应用于方程
(\ref{etale-cohomology-equation-functorial}) 的函子 $u$，并使用 Sites，第
\ref{sites-section-continuous-functors} 节和第
\ref{sites-section-functoriality-PSh} 节对 $u_s = (u_p\ )^\#$ 的描述。
我们偶尔会使用这个拉回公式来证明它的一些基本性质。

\begin{lemma}
\label{etale-cohomology-lemma-stalk-pullback}
设 $f : X \to Y$ 是概形态射。
\begin{enumerate}
\item 函子
$f^{-1} : \textit{Ab}(Y_\etale) \to \textit{Ab}(X_\etale)$
是正合的。
\item 函子
$f^{-1} : \Sh(Y_\etale) \to \Sh(X_\etale)$
是正合的；换言之，它与有限极限和余极限交换，见 Categories，定义
\ref{categories-definition-exact}。
\item 设 $\overline{x} \to X$ 是几何点。设 $\mathcal{G}$ 是 $Y_\etale$ 上的层。
则有典范等同
$$
(f^{-1}\mathcal{G})_{\overline{x}} = \mathcal{G}_{\overline{y}}.
$$
，其中 $\overline{y} = f \circ \overline{x}$。
\item 对任意 \'etale 态射 $V \to Y$，有 $f^{-1}h_V = h_{X \times_Y V}$。
\end{enumerate}
\end{lemma}

\begin{proof}
$f^{-1}$ 在集合层上的正合性，是把 Sites，命题
\ref{sites-proposition-get-morphism} 应用于方程
(\ref{etale-cohomology-equation-functorial}) 中的函子 $u$ 所得的推论。
事实上，拉回的正合性是拓扑斯态射（也可说网站态射）定义的一部分。因此 (2) 成立。
它蕴含 (1)，因为给定阿贝尔层 $\mathcal{G}$（在 $Y_\etale$ 上），
$f^{-1}\mathcal{F}$ 的底层集合层等于 $f^{-1}$ 应用于 $\mathcal{F}$ 的底层集合层，
见 Sites，第 \ref{sites-section-sheaves-algebraic-structures} 节。另见 Modules on Sites，
引理 \ref{sites-modules-lemma-flat-pullback-exact}。文献中有时通过定理
\ref{etale-cohomology-theorem-exactness-stalks} 从 (3) 推出 (1) 与 (2)。

\medskip\noindent
(3) 是关于拉回之茎的一般事实，见 Sites，引理
\ref{sites-lemma-point-morphism-sites}。下面也直接证明 (3)。注意，由引理
\ref{etale-cohomology-lemma-stalk-exact}，取茎与层化交换。回顾 $f^{-1}\mathcal{G}$ 是下述预层的相伴层：
$$
U \longrightarrow \colim_{U \to X \times_Y V} \mathcal{G}(V),
$$
见方程 (\ref{etale-cohomology-equation-pullback})。因此有
\begin{align*}
(f^{-1}\mathcal{G})_{\overline{x}}
& = \colim_{(U, \overline{u})} f^{-1}\mathcal{G}(U) \\
& = \colim_{(U, \overline{u})}
\colim_{a : U \to X \times_Y V} \mathcal{G}(V) \\
& = \colim_{(V, \overline{v})} \mathcal{G}(V) \\
& = \mathcal{G}_{\overline{y}}
\end{align*}
；在第三个等号中，序对 $(U, \overline{u})$ 和映射
$a : U \to X \times_Y V$ 对应于序对 $(V, a \circ \overline{u})$。

\medskip\noindent
(4) 可用类似方法证明，只需等同定义 $f^{-1}h_V$ 的余极限。也可以使用 Yoneda 引理
（Categories，引理 \ref{categories-lemma-yoneda}）和函子性等式
$$
\Mor_{\Sh(X_\etale)}(f^{-1}h_V, \mathcal{F}) =
\Mor_{\Sh(Y_\etale)}(h_V, f_*\mathcal{F}) =
f_*\mathcal{F}(V) = \mathcal{F}(X \times_Y V)
$$
，再结合可表预层是层这一事实。完全一般的结果另见 Sites，引理
\ref{sites-lemma-pullback-representable-sheaf}。
\end{proof}

\noindent
函子对 $(f_*, f^{-1})$ 定义一个小 \'etale 拓扑斯态射
$$
f_{small} :
\Sh(X_\etale)
\longrightarrow
\Sh(Y_\etale)
$$
。关于层上同调的许多一般结论对拓扑斯与拓扑斯态射都成立。我们将尽量指出哪些结果是一般的，
哪些是 \'etale 拓扑斯所特有的。

\begin{remark}
\label{etale-cohomology-remark-functoriality-general}
更一般地，设 $\mathcal{C}_1, \mathcal{C}_2$ 是网站，并假定它们有终对象和纤维积。
设 $u: \mathcal{C}_2 \to \mathcal{C}_1$ 是满足下列条件的函子：
\begin{enumerate}
\item 若 $\{V_i \to V\}$ 是 $\mathcal{C}_2$ 的覆盖，则
$\{u(V_i) \to u(V)\}$ 是 $\mathcal{C}_1$ 的覆盖（此时称 $u$ 是{\it 连续的}）；以及
\item $u$ 与有限极限交换（即 $u$ 左正合；也即 $u$ 保持纤维积与终对象）。
\end{enumerate}
于是可以定义
$f_*: \Sh(\mathcal{C}_1) \to \Sh(\mathcal{C}_2)$
，令 $ f_* \mathcal{F}(V) = \mathcal{F}(u(V))$。此外，存在正合函子
$f^{-1}$，它左伴随于 $f_*$，见 Sites，定义
\ref{sites-definition-morphism-sites} 和命题
\ref{sites-proposition-get-morphism}。警告：仅要求 $u$ 连续且与纤维积交换，
不足以得到拓扑斯态射。
\end{remark}




\section{大拓扑斯的函子性}
\label{etale-cohomology-section-functoriality-big-topoi}

\noindent
给定概形态射 $f : X \to Y$，有大量与 $f$ 相伴的拓扑斯态射；其列表见
Topologies，第 \ref{topologies-section-change-topologies} 节。也许最常用的是拓扑斯态射
$$
f_{big} = f_{big, \tau} :
\Sh((\Sch/X)_\tau)
\longrightarrow
\Sh((\Sch/Y)_\tau)
$$
，其中 $\tau \in \{Zariski, \etale, smooth, syntomic, fppf\}$。
它们分别对应于连续函子
$$
(\Sch/Y)_\tau \longrightarrow (\Sch/X)_\tau, \quad
V/Y \longmapsto X \times_Y V/X
$$
；该函子保持终对象、纤维积与覆盖，因而定义网站态射
$$
f_{big} : (\Sch/X)_\tau \longrightarrow (\Sch/Y)_\tau.
$$
见 Topologies，第 \ref{topologies-section-zariski} 节、
\ref{topologies-section-etale},
\ref{topologies-section-smooth},
\ref{topologies-section-syntomic} 和
\ref{topologies-section-fppf}。特别地，沿 $f_{big}$ 的推前由规则
$$
(f_{big, *}\mathcal{F})(V/Y) = \mathcal{F}(X \times_Y V/X)
$$
给出。事实证明，这些拓扑斯态射有一个容易描述的逆像函子 $f_{big}^{-1}$。即有
$$
(f_{big}^{-1}\mathcal{G})(U/X) = \mathcal{G}(U/Y)
$$
，其中 $U/Y$ 的结构态射是结构态射 $U \to X$ 与 $f$ 的复合，见 Topologies，引理
\ref{topologies-lemma-morphism-big}、
\ref{topologies-lemma-morphism-big-etale},
\ref{topologies-lemma-morphism-big-smooth},
\ref{topologies-lemma-morphism-big-syntomic} 和
\ref{topologies-lemma-morphism-big-fppf}.







\section{函子性与模层}
\label{etale-cohomology-section-morphisms-modules}

\noindent
本节将在 \'etale 拓扑的框架下，重新表述 Descent，第
\ref{descent-section-quasi-coherent-sheaves} 节、
\ref{descent-section-quasi-coherent-cohomology} 节和
\ref{descent-section-quasi-coherent-sheaves-bis}
所说明的一些材料。设 $f : X \to Y$ 是概形态射。前面已经看到（见第
\ref{etale-cohomology-section-functoriality}、\ref{etale-cohomology-section-direct-image} 和
\ref{etale-cohomology-section-inverse-image} 节），它诱导小 \'etale 网站的态射 $f_{small}$。
在 Descent，评注 \ref{descent-remark-change-topologies-ringed} 中已经看到，
$f$ 还诱导自然映射
$$
f_{small}^\sharp :
\mathcal{O}_{Y_\etale}
\longrightarrow
f_{small, *}\mathcal{O}_{X_\etale}
$$
；这是 $Y_\etale$ 上环层的映射，并且 $(f_{small}, f_{small}^\sharp)$
是环化网站的态射。环化网站态射的定义见 Modules on Sites，定义
\ref{sites-modules-definition-ringed-site}。这里只回顾，$f_{small}^\sharp$ 由相容映射系
$$
\text{pr}_V^\sharp : \mathcal{O}(V) \longrightarrow \mathcal{O}(X \times_Y V)
$$
定义，其中 $V$ 遍历 $Y_\etale$ 的对象。

\medskip\noindent
显然，该构造与概形态射的复合相容。更确切地，若 $f : X \to Y$ 和
$g : Y \to Z$ 是概形态射，则有
$$
(g_{small}, g_{small}^\sharp)
\circ
(f_{small}, f_{small}^\sharp)
=
((g \circ f)_{small}, (g \circ f)_{small}^\sharp)
$$
，这是环化拓扑斯态射的等式。此外，由 Modules on Sites，定义
\ref{sites-modules-definition-pushforward} 可知，给定概形态射 $f : X \to Y$，
可得到良定义的拉回与正像函子
\begin{align*}
f_{small}^* :
\textit{Mod}(\mathcal{O}_{Y_\etale})
\longrightarrow
\textit{Mod}(\mathcal{O}_{X_\etale}), \\
f_{small, *} :
\textit{Mod}(\mathcal{O}_{X_\etale})
\longrightarrow
\textit{Mod}(\mathcal{O}_{Y_\etale})
\end{align*}
，它们按通常方式互为伴随。若 $g : Y \to Z$ 是另一个概形态射，则有
$(g \circ f)_{small}^* = f_{small}^* \circ g_{small}^*$
以及 $(g \circ f)_{small, *} = g_{small, *} \circ f_{small, *}$，
这是由上面关于复合的讨论得出的。

\medskip\noindent
由所有 $\mathcal{O}_X$-模（定义在 $X$ 上）组成的范畴，与由所有
$\mathcal{O}_{X_\etale}$-模（定义在 $X_\etale$ 上）组成的范畴之间，有相当大的差别。
但是 Descent，第
\ref{descent-section-quasi-coherent-sheaves} 节、
\ref{descent-section-quasi-coherent-cohomology} 节和
\ref{descent-section-quasi-coherent-sheaves-bis}
的结果告诉我们，考虑 $S$ 上的拟凝聚模与考虑 $S_\etale$ 上的拟凝聚模之间，
并没有太大差别。（例如，已经在定理 \ref{etale-cohomology-theorem-quasi-coherent} 中看到这一点。）
特别地，若 $f : X \to Y$ 是任意概形态射，则对拟凝聚层而言，拉回函子
$f_{small}^*$ 与 $f^*$ 相符，见 Descent，命题
\ref{descent-proposition-equivalence-quasi-coherent-functorial}。
此外，只要 $f$ 拟紧且拟分离，对推前也有相同结论，见 Descent，引理
\ref{descent-lemma-higher-direct-images-small-etale}。

\medskip\noindent
最后简述大网站上结构层的函子性。设 $f : X \to Y$ 是概形态射。任取拓扑
$\tau \in \{Zariski,\linebreak[0]
\etale,\linebreak[0] smooth,\linebreak[0] syntomic,\linebreak[0]
fppf\}$。于是态射
$f_{big} : (\Sch/X)_\tau \to (\Sch/Y)_\tau$
通过映射
$$
f_{big}^\sharp : \mathcal{O}_Y \longrightarrow f_{big, *}\mathcal{O}_X
$$
成为环化网站态射，见 Descent，评注
\ref{descent-remark-change-topologies-ringed}。事实上，它由上面小网站情形中的相同构造给出。












\section{比较拓扑}
\label{etale-cohomology-section-compare-topologies}

\noindent
本节开始研究比较不同拓扑下的层时会发生什么。

\begin{lemma}
\label{etale-cohomology-lemma-where-sections-are-equal}
设 $S$ 是概形。设 $\mathcal{F}$ 是 $S_\etale$ 上的集合层。设
$s, t \in \mathcal{F}(S)$。则存在开集 $W \subset S$，它由下列性质刻画：
态射 $f : T \to S$ 经 $W$ 分解，当且仅当 $s|_T = t|_T$（限制是通过
$f_{small}$ 的拉回）。
\end{lemma}

\begin{proof}
考虑如下预层：对 $U \in \Ob(S_\etale)$，若 $s|_U \not = t|_U$，则赋予空集；
否则赋予单点集。显然，这是 $\Sh(S_\etale)$ 的终对象的子层。由引理
\ref{etale-cohomology-lemma-support-subsheaf-final}，找到表示该预层的开集 $W \subset S$。
对几何点 $\overline{x}$（它位于 $S$ 上），可知 $\overline{x} \in W$ 当且仅当
$s$ 与 $t$ 在 $\overline{x}$ 处的茎相同。由引理
\ref{etale-cohomology-lemma-stalk-pullback} 对拉回之茎的描述，可知 $W$ 具有所需性质。
\end{proof}

\begin{lemma}
\label{etale-cohomology-lemma-describe-pullback}
设 $S$ 是概形。设 $\tau \in \{Zariski, \etale\}$。考虑态射
$$
\pi_S : (\Sch/S)_\tau \longrightarrow S_\tau
$$
，它见于 Topologies，引理 \ref{topologies-lemma-at-the-bottom} 或
\ref{topologies-lemma-at-the-bottom-etale}。设 $\mathcal{F}$ 是 $S_\tau$ 上的层。
则 $\pi_S^{-1}\mathcal{F}$ 由规则
$$
(\pi_S^{-1}\mathcal{F})(T) = \Gamma(T_\tau, f_{small}^{-1}\mathcal{F})
$$
给出，其中 $f : T \to S$。此外，$\pi_S^{-1}\mathcal{F}$ 关于 fpqc 覆盖满足层条件。
\end{lemma}

\begin{proof}
注意，有态射 $i_f : \Sh(T_\tau) \to \Sh(\Sch/S)_\tau)$，使得
$\pi_S \circ i_f = f_{small}$；这里把它们视为态射 $T_\tau \to S_\tau$。见 Topologies，引理
\ref{topologies-lemma-put-in-T}、
\ref{topologies-lemma-morphism-big-small},
\ref{topologies-lemma-put-in-T-etale} 和
\ref{topologies-lemma-morphism-big-small-etale}.
由于拉回具有传递性，可知
$i_f^{-1} \pi_S^{-1}\mathcal{F} = f_{small}^{-1}\mathcal{F}$，正如所需。

\medskip\noindent
设 $\{g_i : T_i \to T\}_{i \in I}$ 是 fpqc 覆盖。最后一个断言的含义如下：
给定层 $\mathcal{G}$（定义在 $T_\tau$ 上），并给定截面
$s_i \in \Gamma(T_i, g_{i, small}^{-1}\mathcal{G})$，它们到
$T_i \times_T T_j$ 的拉回相同，则存在唯一截面 $s$；它是 $\mathcal{G}$ 在
$T$ 上的截面，且其到 $T_i$ 的拉回与 $s_i$ 相同。

\medskip\noindent
设 $V \to T$ 是 $T_\tau$ 的对象，并设 $t \in \mathcal{G}(V)$。
对每个 $i$，存在最大的开集 $W_i \subset T_i \times_T V$，使得
$s_i$ 与 $t$ 的拉回作为 $\mathcal{G}$ 到 $W_i \subset T_i \times_T V$
之拉回的截面相同；见引理 \ref{etale-cohomology-lemma-where-sections-are-equal}。
由于 $s_i$ 与 $s_j$ 在 $T_i \times_T T_j$ 上相同，可知
$W_i$ 与 $W_j$ 拉回到 $T_i \times_T T_j \times_T V$ 后给出同一个开集。
由 Descent，引理 \ref{descent-lemma-open-fpqc-covering}，得到开集
$W \subset V$，它在 $T_i \times_T V$ 中的逆像恢复 $W_i$。

\medskip\noindent
根据 $g_{i, small}^{-1}\mathcal{G}$ 的构造，存在一个
$\tau$-覆盖 $\{T_{ij} \to T_i\}_{j \in J_i}$；对每个 $j$，存在开浸入或
\'etale 态射 $V_{ij} \to T$、截面 $t_{ij} \in \mathcal{G}(V_{ij})$，以及交换图
$$
\xymatrix{
T_{ij} \ar[r] \ar[d] & V_{ij} \ar[d] \\
T_i \ar[r] & T
}
$$
使得 $s_i|_{T_{ij}}$ 是 $t_{ij}$ 的拉回。换言之，在将覆盖
$\{T_i \to T\}$ 替换为 $\{T_{ij} \to T\}$ 后，可以假设存在分解
$T_i \to V_i \to T$，其中 $V_i \in \Ob(T_\tau)$，并且截面
$t_i \in \mathcal{G}(V_i)$ 拉回后等于 $s_i$（在 $T_i$ 上）。
由上一段的结果，得到开集 $W_i \subset V_i$，使得 $t_i|_{W_i}$ 与每个
$s_j$ 在 $T_j \times_T W_i$ 上``相同''。注意 $T_i \to V_i$ 经
$W_i$ 分解。因此 $\{W_i \to T\}$ 是一个 $\tau$-覆盖，引理得证。
\end{proof}

\begin{lemma}
\label{etale-cohomology-lemma-sections-upstairs}
设 $S$ 是概形。设 $f : T \to S$ 是满足下列条件的态射：
\begin{enumerate}
\item $f$ 平坦且拟紧；
\item $f$ 的几何纤维连通。
\end{enumerate}
设 $\mathcal{F}$ 是 $S_\etale$ 上的层。则
$\Gamma(S, \mathcal{F}) = \Gamma(T, f^{-1}_{small}\mathcal{F})$。
\end{lemma}

\begin{proof}
存在典范映射
$\Gamma(S, \mathcal{F}) \to \Gamma(T, f_{small}^{-1}\mathcal{F})$.
由于 $f$ 是满射（因为它的纤维连通），可知此映射是单射。

\medskip\noindent
为证明该映射是满射，设
$\alpha \in \Gamma(T, f_{small}^{-1}\mathcal{F})$。
由于 $\{T \to S\}$ 是 fpqc 覆盖，可以用引理 \ref{etale-cohomology-lemma-describe-pullback}
看出，只需证明 $\alpha$ 通过两个投影拉回到 $T \times_S T$ 后给出同一个截面。
设 $\overline{s} \to S$ 是几何点。只需证明这种相同性在
$(T \times_S T)_{\overline{s}}$ 上成立，因为 $T \times_S T$ 的每个几何点
都包含在这些几何纤维之一中。换言之，我们要证明
$\alpha|_{T_{\overline{s}}}$ 拉回到
$$
(T \times_S T)_{\overline{s}} =
T_{\overline{s}} \times_{\overline{s}} T_{\overline{s}}
$$
后，经到 $T_{\overline{s}}$ 的两个投影所得的截面相同。
然而，$\mathcal{F}|_{T_{\overline{s}}}$ 是 $\mathcal{F}|_{\overline{s}}$ 的拉回，
故它是取值于 $\mathcal{F}_{\overline{s}}$ 的常值层。根据假设，
$T_{\overline{s}}$ 连通，所以常值层的任一截面都是常值的。因此
$\alpha|_{T_{\overline{s}}}$ 对应于 $\mathcal{F}_{\overline{s}}$ 的一个元素。
于是，到 $(T \times_S T)_{\overline{s}}$ 的两个拉回均对应于这个相同元素，结论得证。
\end{proof}

\noindent
下面给出引理 \ref{etale-cohomology-lemma-sections-upstairs} 的一个版本，其中不假设该态射平坦。

\begin{lemma}
\label{etale-cohomology-lemma-sections-upstairs-submersive}
设 $S$ 是概形。设 $f : X \to S$ 是满足下列条件的态射：
\begin{enumerate}
\item $f$ 是商映射；
\item $f$ 的几何纤维连通。
\end{enumerate}
设 $\mathcal{F}$ 是 $S_\etale$ 上的层。则
$\Gamma(S, \mathcal{F}) = \Gamma(X, f^{-1}_{small}\mathcal{F})$。
\end{lemma}

\begin{proof}
存在典范映射
$\Gamma(S, \mathcal{F}) \to \Gamma(X, f_{small}^{-1}\mathcal{F})$.
由于 $f$ 是满射（因为它的纤维连通），可知此映射是单射。

\medskip\noindent
为证明该映射是满射，设
$\tau \in \Gamma(X, f_{small}^{-1}\mathcal{F})$。
只需找到一个 \'etale 覆盖 $\{U_i \to S\}$ 和截面
$\sigma_i \in \mathcal{F}(U_i)$，使得 $\sigma_i$ 拉回后等于
$\tau|_{X \times_S U_i}$。事实上，上面证明的单射性保证 $\sigma_i$ 与
$\sigma_j$ 在 $\mathcal{F}$ 到 $U_i \times_S U_j$ 上的限制相同。
因此得到唯一截面 $\sigma \in \mathcal{F}(S)$，其限制为 $\sigma_i$
（在 $U_i$ 上）。于是 $\sigma$ 到 $X$ 的拉回就是 $\tau$，因为这在局部成立。

\medskip\noindent
设 $\overline{x}$ 是 $X$ 的几何点，像为 $\overline{s}$（在 $S$ 中）。
考虑 $\tau$ 在下列茎中的像：
$$
(f_{small}^{-1}\mathcal{F})_{\overline{x}} = \mathcal{F}_{\overline{s}}
$$
见引理 \ref{etale-cohomology-lemma-stalk-pullback}。可以找到一个 \'etale 邻域
$U \to S$（它是 $\overline{s}$ 的邻域）和截面 $\sigma \in \mathcal{F}(U)$，
它映到茎中的这个像。因此，在对 $S$ 以 $U$ 作替换，并对 $X$ 以
$X \times_S U$ 作替换后，可以假设存在截面 $\sigma$，它是 $\mathcal{F}$
在 $S$ 上的截面，并且它在
$(f_{small}^{-1}\mathcal{F})_{\overline{x}}$ 中的像与 $\tau$ 相同。

\medskip\noindent
由引理 \ref{etale-cohomology-lemma-where-sections-are-equal}，存在最大的开集 $W \subset X$，
使得 $f_{small}^{-1}\sigma$ 与 $\tau$ 在 $W$ 上相同，且 $W$ 的构造与进一步的拉回可交换。
注意，$\mathcal{F}$ 到几何纤维 $X_{\overline{s}}$ 的拉回，是把
$\mathcal{F}_{\overline{s}}$ 看作 $\overline{s}$ 上的层后经
$X_{\overline{s}} \to \overline{s}$ 所作的拉回。因此，$\tau$ 与 $\sigma$
给出取值于 $\mathcal{F}_{\overline{s}}$ 的常值层的截面（该层在 $X_{\overline{s}}$ 上），
并且它们在一点处相同。根据假设 $X_{\overline{s}}$ 连通，所以 $W$ 包含
$X_s$。对不同几何纤维作同样的论证可知，$W$ 包含它所遇到的每条纤维。
由于 $f$ 是商映射，可知 $W$ 是 $s$ 的某个开邻域（在 $S$ 中）的逆像。
证明完成。
\end{proof}

\begin{lemma}
\label{etale-cohomology-lemma-sections-base-field-extension}
设 $K/k$ 是域扩张，且 $k$ 是可分代数闭域。设 $S$ 是 $k$ 上的概形。记
$p : S_K = S \times_{\Spec(k)} \Spec(K) \to S$ 为投影。设 $\mathcal{F}$ 是
$S_\etale$ 上的层。则
$\Gamma(S, \mathcal{F}) = \Gamma(S_K, p^{-1}_{small}\mathcal{F})$。
\end{lemma}

\begin{proof}
这由引理 \ref{etale-cohomology-lemma-sections-upstairs} 得出。事实上，$p$ 作为
$\Spec(K) \to \Spec(k)$ 的基变换，显然平坦且拟紧。另一方面，若
$\overline{s} : \Spec(L) \to S$ 是几何点，则 $p$ 在 $\overline{s}$ 上的纤维
是 $K \otimes_k L$ 的谱；由 Algebra，引理
\ref{algebra-lemma-separably-closed-irreducible}，它不可约，因而连通。
\end{proof}









\section{恢复态射}
\label{etale-cohomology-section-morphisms}

\noindent
本节证明：把概形对应到其局部环化小 \'etale 拓扑斯的规则，在适当意义下是全忠实的；
见定理 \ref{etale-cohomology-theorem-fully-faithful}。

\begin{lemma}
\label{etale-cohomology-lemma-morphism-locally-ringed}
设 $f : X \to Y$ 是概形态射。环化网站态射
$(f_{small}, f_{small}^\sharp)$（它与 $f$ 对应）是局部环化网站的态射；见 Modules on Sites，
定义 \ref{sites-modules-definition-morphism-locally-ringed-topoi}。
\end{lemma}

\begin{proof}
注意，该断言有意义，因为我们已经看到
$(X_\etale, \mathcal{O}_{X_\etale})$ 和
$(Y_\etale, \mathcal{O}_{Y_\etale})$ 是局部环化网站；见引理
\ref{etale-cohomology-lemma-etale-site-locally-ringed}。此外，我们知道 $X_\etale$ 有足够多的点；
见定理 \ref{etale-cohomology-theorem-exactness-stalks} 和注记 \ref{etale-cohomology-remarks-enough-points}。
因此，只需证明 $(f_{small}, f_{small}^\sharp)$ 满足 Modules on Sites，
引理 \ref{sites-modules-lemma-locally-ringed-morphism} 的条件 (3)。
为此取一点 $p$（它属于 $X_\etale$）。由引理
\ref{etale-cohomology-lemma-points-small-etale-site}，$p$ 对应于几何点 $\overline{x}$（它位于 $X$ 上）。
由引理 \ref{etale-cohomology-lemma-stalk-pullback}，点 $q = f_{small} \circ p$ 对应于几何点
$\overline{y} = f \circ \overline{x}$（它位于 $Y$ 上）。因此，我们要证明的断言是：
所诱导的茎映射
$$
(\mathcal{O}_Y)_{\overline{y}} \longrightarrow (\mathcal{O}_X)_{\overline{x}}
$$
是局部环同态。设 $a \in (\mathcal{O}_Y)_{\overline{y}}$ 是左侧的一个元素，
它映到右侧的极大理想中的元素。设 $a$ 是三元组 $(V, \overline{v}, a)$ 的等价类，
其中 $V \to Y$ 是 \'etale 态射，$\overline{v} : \overline{x} \to V$ 是 $Y$ 上的态射，
且 $a \in \mathcal{O}(V)$。它映到下列三元组的等价类：
$(X \times_Y V, \overline{x} \times \overline{v}, \text{pr}_V^\sharp(a))$
（它位于局部环 $(\mathcal{O}_X)_{\overline{x}}$ 中）。显然，属于极大理想意味着：
把 $\text{pr}_V^\sharp(a)$ 拉回为 $\kappa(\overline{x})$ 的一个元素后得到零。
因而，把 $a$ 拉回到 $\kappa(\overline{x})$ 也得到零。这意味着 $a$ 属于
$(\mathcal{O}_Y)_{\overline{y}}$ 的极大理想。
\end{proof}

\begin{lemma}
\label{etale-cohomology-lemma-2-morphism}
设 $X$、$Y$ 是概形。设 $f : X \to Y$ 是概形态射。设 $t$ 是一个
$2$-态射，其源和目标均为 $(f_{small}, f_{small}^\sharp)$；见 Modules on Sites，
定义 \ref{sites-modules-definition-2-morphism-ringed-topoi}。则 $t = \text{id}$。
\end{lemma}

\begin{proof}
这意味着 $t : f^{-1}_{small} \to f^{-1}_{small}$ 是一个函子变换，并且下图交换：
$$
\xymatrix{
f_{small}^{-1}\mathcal{O}_Y
\ar[rd]_{f_{small}^\sharp}  & &
f_{small}^{-1}\mathcal{O}_Y \ar[ll]^t \ar[ld]^{f_{small}^\sharp} \\
& \mathcal{O}_X
}
$$
设 $V \to Y$ 是 \'etale 态射，且 $V$ 仿射。由 Morphisms，引理
\ref{morphisms-lemma-quasi-affine-finite-type-over-S}，可以选择浸入
$i : V \to \mathbf{A}^n_Y$（在 $Y$ 上）。用层来表述，这意味着 $i$ 诱导层的单射
$h_i : h_V \to \prod_{j = 1, \ldots, n} \mathcal{O}_Y$。
基变换 $i'$（它把 $i$ 基变换到 $X$ 得到）是浸入
（Schemes，引理 \ref{schemes-lemma-base-change-immersion}）。因此
$i' : X \times_Y V \to \mathbf{A}^n_X$ 是浸入，这又意味着
$h_{i'} : h_{X \times_Y V} \to \prod_{j = 1, \ldots, n} \mathcal{O}_X$
是层的单射。借助引理 \ref{etale-cohomology-lemma-stalk-pullback} 的等同
$f_{small}^{-1}h_V = h_{X \times_Y V}$，映射 $h_{i'}$ 等于
$$
\xymatrix{
f_{small}^{-1}h_V \ar[r]^-{f^{-1}h_i} &
\prod_{j = 1, \ldots, n} f_{small}^{-1}\mathcal{O}_Y
\ar[r]^{\prod f^\sharp} &
\prod_{j = 1, \ldots, n} \mathcal{O}_X
}
$$
（验证略去）。这意味着映射 $t : f_{small}^{-1}h_V \to f_{small}^{-1}h_V$
可置于下列交换图中：
$$
\xymatrix{
f_{small}^{-1}h_V \ar[r]^-{f^{-1}h_i} \ar[d]^t &
\prod_{j = 1, \ldots, n} f_{small}^{-1}\mathcal{O}_Y
\ar[r]^-{\prod f^\sharp} \ar[d]^{\prod t} &
\prod_{j = 1, \ldots, n} \mathcal{O}_X \ar[d]^{\text{id}}\\
f_{small}^{-1}h_V \ar[r]^-{f^{-1}h_i} &
\prod_{j = 1, \ldots, n} f_{small}^{-1}\mathcal{O}_Y
\ar[r]^-{\prod f^\sharp} &
\prod_{j = 1, \ldots, n} \mathcal{O}_X
}
$$
右方块的交换性由上面说明的关于 $t$ 的假设得出。
根据前面的讨论，横箭头的复合是单射，故左边的竖箭头也是恒等映射。
$Y_\etale$ 上的任何集合层都接受一个满射，其源是形如 $h_V$ 的层的（巨大）余积，
其中 $V$ 仿射（合用 Topologies，引理 \ref{topologies-lemma-alternative} 与
Sites，引理 \ref{sites-lemma-sheaf-coequalizer-representable}）。因此
$t : f_{small}^{-1} \to f_{small}^{-1}$ 是所需的恒等变换。
\end{proof}

\begin{lemma}
\label{etale-cohomology-lemma-faithful}
设 $X$、$Y$ 是概形。任何两个概形态射 $a, b : X \to Y$，如果存在一个 $2$-同构
$(a_{small}, a_{small}^\sharp) \cong (b_{small}, b_{small}^\sharp)$
（它位于环化拓扑斯的 $2$-范畴中），则二者相等。
\end{lemma}

\begin{proof}
让我们仔细论证，因为这里有些容易混淆。设
$t : a_{small}^{-1} \to b_{small}^{-1}$ 为该 $2$-同构。考虑任意开集
$V \subset Y$。注意 $h_V$ 是终层 $*$ 的子层。因此
$a_{small}^{-1}h_V = h_{a^{-1}(V)}$ 和
$b_{small}^{-1}h_V = h_{b^{-1}(V)}$ 都是终层的子层。故同构
$$
t : a_{small}^{-1}h_V = h_{a^{-1}(V)} \to b_{small}^{-1}h_V = h_{b^{-1}(V)}
$$
必为恒等，并且 $a^{-1}(V) = b^{-1}(V)$。因此 $a$ 与 $b$ 在底拓扑空间上相等。
接着取截面 $f \in \mathcal{O}_Y(V)$。它决定且被下列集合层映射决定：
$f : h_V \to \mathcal{O}_Y$。将其拉回并作用 $t$，得到交换图
$$
\xymatrix{
h_{b^{-1}(V)} \ar@{=}[r] &
b_{small}^{-1}h_V \ar[d]^{b_{small}^{-1}(f)} & &
a_{small}^{-1}h_V \ar[d]^{a_{small}^{-1}(f)} \ar[ll]^t &
h_{a^{-1}(V)} \ar@{=}[l]
\\
& b_{small}^{-1}\mathcal{O}_Y
\ar[rd]_{b^\sharp}  & &
a_{small}^{-1}\mathcal{O}_Y \ar[ll]^t \ar[ld]^{a^\sharp} \\
& & \mathcal{O}_X
}
$$
其中三角形的交换性来自 Modules on Sites，第
\ref{sites-modules-section-2-category} 节中 $2$-同构的定义。
上面已看到，顶部水平箭头的复合来自恒等式 $a^{-1}(V) = b^{-1}(V)$。
因此图的交换性给出
$a_{small}^\sharp(f) = b_{small}^\sharp(f)$ 在
$\mathcal{O}_X(a^{-1}(V)) = \mathcal{O}_X(b^{-1}(V))$ 中成立。
由于这对每个开集 $V$ 和每个 $f \in \mathcal{O}_Y(V)$ 都成立，可知概形态射
$a = b$。
\end{proof}

\begin{lemma}
\label{etale-cohomology-lemma-morphism-ringed-etale-topoi-affines}
设 $X$、$Y$ 是仿射概形。设
$$
(g, g^\#) :
(\Sh(X_\etale), \mathcal{O}_X)
\longrightarrow
(\Sh(Y_\etale), \mathcal{O}_Y)
$$
是局部环化拓扑斯的态射。则存在唯一的概形态射 $f : X \to Y$，使得
$(g, g^\#)$ $2$-同构于 $(f_{small}, f_{small}^\sharp)$；见 Modules on Sites，
定义 \ref{sites-modules-definition-2-morphism-ringed-topoi}。
\end{lemma}

\begin{proof}
本证明中用 $\mathcal{O}_X$ 表示小 \'etale 网站 $X_\etale$ 的结构层，
$\mathcal{O}_Y$ 亦然。记 $Y = \Spec(B)$ 且 $X = \Spec(A)$。由于
$B = \Gamma(Y_\etale, \mathcal{O}_Y)$、
$A = \Gamma(X_\etale, \mathcal{O}_X)$，可知 $g^\sharp$ 诱导环同态
$\varphi : B \to A$。令 $f = \Spec(\varphi) : X \to Y$ 为相应的仿射概形态射。
我们将证明这个 $f$ 满足要求。

\medskip\noindent
设 $V \to Y$ 是仿射概形，并且在 $Y$ 上是 \'etale 的。因此可以写成
$V = \Spec(C)$，其中 $C$ 是 \'etale $B$-代数。可以写成
$$
C = B[x_1, \ldots, x_n]/(P_1, \ldots, P_n)
$$
其中 $P_i$ 是多项式，并且 $\Delta = \det(\partial P_i/ \partial x_j)$ 在 $C$ 中可逆；
例如见 Algebra，引理 \ref{algebra-lemma-etale-standard-smooth}。
若 $T$ 是 $Y$ 上的概形，则一个 $T$-值点（取值于 $V$）由 $n$ 个截面
（属于 $\Gamma(T, \mathcal{O}_T)$）给出，它们满足多项式方程
$P_1 = 0, \ldots, P_n = 0$。换言之，层 $h_V$（在 $Y_\etale$ 上）是下列两个映射的等化子：
$$
\xymatrix{
\prod\nolimits_{i = 1, \ldots, n} \mathcal{O}_Y
\ar@<1ex>[r]^a \ar@<-1ex>[r]_b &
\prod\nolimits_{j = 1, \ldots, n} \mathcal{O}_Y
}
$$
其中 $b(h_1, \ldots, h_n) = 0$，而
$a(h_1, \ldots, h_n) =
(P_1(h_1, \ldots, h_n), \ldots, P_n(h_1, \ldots, h_n))$。
由于 $g^{-1}$ 正合，可知下列实线交换图的顶行也是等化子图：
$$
\xymatrix{
g^{-1}h_V \ar[r] \ar@{..>}[d] &
\prod\nolimits_{i = 1, \ldots, n} g^{-1}\mathcal{O}_Y
\ar@<1ex>[r]^{g^{-1}a} \ar@<-1ex>[r]_{g^{-1}b} \ar[d]^{\prod g^\sharp} &
\prod\nolimits_{j = 1, \ldots, n} g^{-1}\mathcal{O}_Y \ar[d]^{\prod g^\sharp}\\
h_{X \times_Y V} \ar[r] &
\prod\nolimits_{i = 1, \ldots, n} \mathcal{O}_X
\ar@<1ex>[r]^{a'} \ar@<-1ex>[r]_{b'} &
\prod\nolimits_{j = 1, \ldots, n} \mathcal{O}_X  \\
}
$$
这里 $b'$ 是零映射，而 $a'$ 是由诸像
$P'_i = \varphi(P_i) \in A[x_1, \ldots, x_n]$ 按同一规则定义的映射：
$a'(h_1, \ldots, h_n) =
(P'_1(h_1, \ldots, h_n), \ldots, P'_n(h_1, \ldots, h_n))$；这就是定义 $a$ 时所用的规则。
图的交换性来自整体截面上的等式 $\varphi = g^\sharp$。底行同样是等化子图；
理由与前面完全相同，因为 $X \times_Y V$ 是仿射概形
$\Spec(A \otimes_B C)$，并且
$A \otimes_B C = A[x_1, \ldots, x_n]/(P'_1, \ldots, P'_n)$。
因此，得到唯一的虚线箭头 $g^{-1}h_V \to h_{X \times_Y V}$，使之嵌入该图。

\medskip\noindent
我们断言层映射 $g^{-1}h_V \to h_{X \times_Y V}$ 是同构。由于 $X$ 的小
\'etale 网站有足够多的点（定理 \ref{etale-cohomology-theorem-exactness-stalks}），只需在茎上证明。
因此，设 $\overline{x}$ 是 $X$ 的几何点，并以 $p$ 表示 $X$ 的小 \'etale 拓扑斯中
相应的点。置 $q = g \circ p$。这是 $Y$ 的小 \'etale 拓扑斯的一个点。
由引理 \ref{etale-cohomology-lemma-points-small-etale-site}，可知 $q$ 对应于几何点
$\overline{y}$（它位于 $Y$ 上）。考虑茎映射
$$
(g^\sharp)_p :
(\mathcal{O}_Y)_{\overline{y}} =
\mathcal{O}_{Y, q} =
(g^{-1}\mathcal{O}_Y)_p
\longrightarrow
\mathcal{O}_{X, p} =
(\mathcal{O}_X)_{\overline{x}}
$$
由于 $(g, g^\sharp)$ 是{\it 局部}环化拓扑斯的态射，$(g^\sharp)_p$ 是严格 Hensel
局部环之间的局部环同态。对上面的大交换图取局部化，并应用 Algebra，引理
\ref{algebra-lemma-strictly-henselian-solutions}，可知
$(g^{-1}h_V)_p \to (h_{X \times_Y V})_p$ 是所需的同构。

\medskip\noindent
我们断言同构 $g^{-1}h_V \to h_{X \times_Y V}$ 是函子性的。具体地，设
$V_1 \to V_2$ 是仿射概形间的态射，且这些概形在 $Y$ 上是 \'etale 的。写成
$V_i = \Spec(C_i)$，其中
$$
C_i = B[x_{i, 1}, \ldots, x_{i, n_i}]/(P_{i, 1}, \ldots, P_{i, n_i})
$$
态射 $V_1 \to V_2$ 由 $B$-代数同态 $C_2 \to C_1$ 给出，后者又由一些多项式
$Q_j \in B[x_{1, 1}, \ldots, x_{1, n_1}]$（其中 $j = 1, \ldots, n_2$）给出。
于是容易证明下列层的图交换：
$$
\xymatrix{
h_{V_1} \ar[d] \ar[r] & \prod_{i = 1, \ldots, n_1} \mathcal{O}_Y
\ar[d]^{Q_1, \ldots, Q_{n_2}}\\
h_{V_2} \ar[r] & \prod_{i = 1, \ldots, n_2} \mathcal{O}_Y
}
$$
将它拉回到 $X_\etale$，得到下列实线交换图：
$$
\xymatrix{
g^{-1}h_{V_1} \ar@{..>}[dd] \ar[rrd] \ar[r] &
\prod_{i = 1, \ldots, n_1} g^{-1}\mathcal{O}_Y
\ar[dd]^{g^\sharp}
\ar[rrd]^{Q_1, \ldots, Q_{n_2}} \\
& & g^{-1}h_{V_2} \ar@{..>}[dd] \ar[r] &
\prod_{i = 1, \ldots, n_2} g^{-1}\mathcal{O}_Y
\ar[dd]^{g^\sharp} \\
h_{X \times_Y V_1} \ar[r] \ar[rrd] &
\prod\nolimits_{i = 1, \ldots, n_1} \mathcal{O}_X
\ar[rrd]^{Q'_1, \ldots, Q'_{n_2}} \\
& & h_{X \times_Y V_2} \ar[r] &
\prod\nolimits_{i = 1, \ldots, n_2} \mathcal{O}_X
}
$$
其中 $Q'_j \in A[x_{1, 1}, \ldots, x_{1, n_1}]$ 是 $Q_j$ 经 $\varphi$ 得到的像。
由于虚线箭头存在、使两个方块交换，而且水平箭头是单射，所以整个图交换。
这证明了函子性（也证明了 $g^{-1}h_V \to h_{X \times_Y V}$ 的构造不依赖于
表示的选择，尽管严格说来我们不需要证明这一点）。

\medskip\noindent
此时可以证明 $f_{small, *} \cong g_*$。事实上，设 $\mathcal{F}$ 是 $X_\etale$
上的层。对每个仿射的 $V \in \Ob(X_\etale)$，有
\begin{align*}
(g_*\mathcal{F})(V)
& =
\Mor_{\Sh(Y_\etale)}(h_V, g_*\mathcal{F}) \\
& =
\Mor_{\Sh(X_\etale)}(g^{-1}h_V, \mathcal{F}) \\
& =
\Mor_{\Sh(X_\etale)}(h_{X \times_Y V}, \mathcal{F}) \\
& =
\mathcal{F}(X \times_Y V) \\
& =
f_{small, *}\mathcal{F}(V)
\end{align*}
其中第三个等式使用上面构造的同构 $g^{-1}h_V \cong h_{X \times_Y V}$。
这些同构关于 $\mathcal{F}$ 显然是函子性的，关于 $V$ 也具有函子性，因为同构
$g^{-1}h_V \cong h_{X \times_Y V}$ 本身是函子性的。$Y_\etale$ 上的任意层都由其
到仿射概形子范畴的限制所决定（Topologies，引理
\ref{topologies-lemma-alternative}），所以得到所需的函子同构
$f_{small, *} \cong g_*$。

\medskip\noindent
最后，需要检验：借助上面的同构 $f_{small, *} \cong g_*$，映射
$f_{small}^\sharp$ 与 $g^\sharp$ 相同。根据构造，这对 $\mathcal{O}_Y$ 的整体截面，
即对 $B$ 的元素已经成立。只需对仿射概形 $V$ 上的截面检验，其中它在
$Y$ 上是 \'etale 的（再次使用 Topologies，引理
\ref{topologies-lemma-alternative}）。像前面一样写
$V = \Spec(C)$、$C = B[x_i]/(P_j)$，只需检验坐标函数 $x_i$ 映到
$\mathcal{O}_X$ 在 $X \times_Y V$ 上的同一截面。这恰好就是下图交换的含义：
$$
\xymatrix{
g^{-1}h_V \ar[r] \ar@{..>}[d] &
\prod\nolimits_{i = 1, \ldots, n} g^{-1}\mathcal{O}_Y
\ar[d]^{\prod g^\sharp} \\
h_{X \times_Y V} \ar[r] &
\prod\nolimits_{i = 1, \ldots, n} \mathcal{O}_X
}
$$
因此引理得证。
\end{proof}

\noindent
下面给出一般概形的版本。

\begin{theorem}
\label{etale-cohomology-theorem-fully-faithful}
设 $X$、$Y$ 是概形。设
$$
(g, g^\#) :
(\Sh(X_\etale), \mathcal{O}_X)
\longrightarrow
(\Sh(Y_\etale), \mathcal{O}_Y)
$$
是局部环化拓扑斯的态射。则存在唯一的概形态射 $f : X \to Y$，使得
$(g, g^\#)$ 同构于 $(f_{small}, f_{small}^\sharp)$。换言之，构造
$$
\Sch \longrightarrow \textit{Locally ringed topoi},
\quad
X \longrightarrow (X_\etale, \mathcal{O}_X)
$$
是全忠实的（右侧的态射取到 $2$-同构为止）。
\end{theorem}

\begin{proof}
可以仔细调整引理 \ref{etale-cohomology-lemma-morphism-ringed-etale-topoi-affines} 的证明论证，
使之适用于整体情形，从而证明本定理。不过，我们希望说明如何利用引理
\ref{etale-cohomology-lemma-2-morphism} 的结论所提供的刚性，把前一引理的结果粘合成整体态射。
遗憾的是，这有些繁琐。

\medskip\noindent
先在 $Y$ 仿射时证明存在性。此时选取仿射开覆盖 $X = \bigcup U_i$。对每个 $i$，
包含态射 $j_i : U_i \to X$ 诱导局部环化拓扑斯的态射
$(j_{i, small}, j_{i, small}^\sharp) :
(\Sh(U_{i, \etale}), \mathcal{O}_{U_i})
\to
(\Sh(X_\etale), \mathcal{O}_X)$
，这是引理 \ref{etale-cohomology-lemma-morphism-locally-ringed} 的结论。把它与 $(g, g^\sharp)$
复合，得到局部环化拓扑斯的态射
$$
(g, g^\sharp) \circ (j_{i, small}, j_{i, small}^\sharp) :
(\Sh(U_{i, \etale}), \mathcal{O}_{U_i})
\to
(\Sh(Y_\etale), \mathcal{O}_Y)
$$
；见 Modules on Sites，引理
\ref{sites-modules-lemma-composition-morphisms-locally-ringed-topoi}。
由引理 \ref{etale-cohomology-lemma-morphism-ringed-etale-topoi-affines}，存在唯一的概形态射
$f_i : U_i \to Y$ 和一个 $2$-同构
$$
t_i :
(f_{i, small}, f_{i, small}^\sharp)
\longrightarrow
(g, g^\sharp) \circ (j_{i, small}, j_{i, small}^\sharp).
$$
置 $U_{i, i'} = U_i \cap U_{i'}$，并以 $j_{i, i'} : U_{i, i'} \to U_i$ 表示包含态射。
由于有
$j_i \circ j_{i, i'} = j_{i'} \circ j_{i', i}$
，可知
\begin{align*}
(g, g^\sharp) \circ
(j_{i, small}, j_{i, small}^\sharp) \circ
(j_{i, i', small}, j_{i, i', small}^\sharp)
= \\
(g, g^\sharp) \circ
(j_{i', small}, j_{i', small}^\sharp) \circ
(j_{i', i, small}, j_{i', i, small}^\sharp)
\end{align*}
因此由唯一性（见引理 \ref{etale-cohomology-lemma-faithful}），得到
$f_i \circ j_{i, i'} = f_{i'} \circ j_{i', i}$。换言之，概形态射
$f_i = f \circ j_i$ 是某个整体概形态射 $f : X \to Y$ 的限制。考虑下列
$2$-同构的图（为简化记号，我们省略分量 ${}^\sharp$）：
$$
\xymatrix{
g \circ j_{i, small} \circ j_{i, i', small}
\ar[rr]^{t_i \star \text{id}_{j_{i, i', small}}}
\ar@{=}[d] & &
f_{small} \circ j_{i, small} \circ j_{i, i', small} \ar@{=}[d] \\
g \circ j_{i', small} \circ j_{i', i, small}
\ar[rr]^{t_{i'} \star \text{id}_{j_{i', i, small}}} & &
f_{small} \circ j_{i', small} \circ j_{i', i, small}
}
$$
记号 $\star$ 表示水平复合；一般情形见 Categories，定义
\ref{categories-definition-2-category}，本情形见 Sites，第
\ref{sites-section-2-category} 节。由引理 \ref{etale-cohomology-lemma-2-morphism}，该图交换。
因此，对任意层 $\mathcal{G}$（在 $Y_\etale$ 上），同构
$t_i : f_{small}^{-1}\mathcal{G}|_{U_i} \to g^{-1}\mathcal{G}|_{U_i}$
在 $U_{i, i'}$ 上相同，从而得到整体同构
$t : f_{small}^{-1}\mathcal{G} \to g^{-1}\mathcal{G}$。显然，该同构关于
$\mathcal{G}$ 是函子性的，并且与映射 $f_{small}^\sharp$ 和 $g^\sharp$ 相容
（因为它在局部与这些映射相容）。这证明了 $Y$ 仿射时的定理。

\medskip\noindent
在一般情形，设 $V \subset Y$ 是仿射开集。那么 $h_V$ 是终层 $*$ 的子层
（该终层位于 $Y_\etale$ 上）。由于 $g$ 正合，可知 $g^{-1}h_V$ 是
$X_\etale$ 上终层的子层。因此，由引理 \ref{etale-cohomology-lemma-support-subsheaf-final}，
存在开子概形 $W \subset X$，使得 $g^{-1}h_V = h_W$。由 Modules on Sites，
引理 \ref{sites-modules-lemma-localize-morphism-locally-ringed-topoi}，存在下列
局部环化拓扑斯态射的交换图：
$$
\xymatrix{
(\Sh(W_\etale), \mathcal{O}_W) \ar[r] \ar[d]_{g'} &
(\Sh(X_\etale), \mathcal{O}_X) \ar[d]^g \\
(\Sh(V_\etale), \mathcal{O}_V) \ar[r] &
(\Sh(Y_\etale), \mathcal{O}_Y)
}
$$
其中水平箭头是局部化态射（由包含态射 $V \to Y$ 和 $W \to X$ 诱导），
而 $g'$ 由 $g$ 诱导。由上一段的结论，得到概形态射 $f' : W \to V$ 和一个
$2$-同构
$t : (f'_{small}, (f'_{small})^\sharp) \to (g', (g')^\sharp)$.
与前面完全一样，这些态射 $f'$（当仿射开集 $V \subset Y$ 变化时）由唯一性
在交叠处相同，故得到态射 $f : X \to Y$。此外，这些 $2$-同构 $t$ 由引理
\ref{etale-cohomology-lemma-2-morphism} 再次可知在交叠处相容，因而得到整体 $2$-同构
$(f_{small}, (f_{small})^\sharp) \to (g, (g)^\sharp)$.
这正是所需结论。略去了一些细节。
\end{proof}










\section{推前与拉回}
\label{etale-cohomology-section-monomorphisms}

\noindent
设 $f : X \to Y$ 是概形态射。下面列出我们将要考虑的条件：
\begin{enumerate}
\item[(A)] 对每个 \'etale 态射 $U \to X$ 和 $u \in U$，存在 \'etale 态射
$V \to Y$、不交并分解 $X \times_Y V = W \amalg W'$，以及态射
$h : W \to U$（在 $X$ 上），使得 $u$ 属于 $h$ 的像。
\item[(B)] 对每个 \'etale 态射 $V \to Y$ 和每个 \'etale 覆盖
$\{U_i \to X \times_Y V\}$，存在 \'etale 覆盖 $\{V_j \to V\}$，使得对每个
$j$ 都有 $X \times_Y V_j = \coprod W_{ij}$，其中 $W_{ij} \to X \times_Y V$
经 $U_i \to X \times_Y V$ 分解（对某个 $i$）。
\item[(C)] 对每个 \'etale 态射 $U \to X$，存在 \'etale 态射 $V \to Y$ 和满态射
$X \times_Y V \to U$（在 $X$ 上）。
\end{enumerate}
这些性质中的每一个都可以用函子 $f_{small, *}$ 的行为来表述。我们将在随后几节
阐明这一点。



\section{性质 (A)}
\label{etale-cohomology-section-A}

\noindent
性质 (A) 的定义见第 \ref{etale-cohomology-section-monomorphisms} 节。

\begin{lemma}
\label{etale-cohomology-lemma-property-A-implies}
设 $f : X \to Y$ 是概形态射。假设 (A) 成立。
\begin{enumerate}
\item
$f_{small, *} :
\textit{Ab}(X_\etale)
\to
\textit{Ab}(Y_\etale)$
 保持对单射和满射的反映，
\item 映射 $f_{small}^{-1}f_{small, *}\mathcal{F} \to \mathcal{F}$
对任意 Abel 层 $\mathcal{F}$（在 $X_\etale$ 上）都是满射，
\item
$f_{small, *} :
\textit{Ab}(X_\etale)
\to
\textit{Ab}(Y_\etale)$
是忠实的。
\end{enumerate}
\end{lemma}

\begin{proof}
设 $\mathcal{F}$ 是 Abel 层（在 $X_\etale$ 上）。设 $U$ 是 $X_\etale$ 的对象。
根据假设，可以找到覆盖 $\{W_i \to U\}$（它位于 $X_\etale$ 中），使每个 $W_i$
都是 $X \times_Y V_i$ 的开闭子概形，其中 $V_i$ 是 $Y_\etale$ 的某个对象。
层条件给出
$$
\mathcal{F}(U) \subset \prod \mathcal{F}(W_i)
$$
，且 $\mathcal{F}(W_i)$ 是
$\mathcal{F}(X \times_Y V_i) = f_{small, *}\mathcal{F}(V_i)$ 的直和项。
因此，$f_{small, *}$ 显然反映单射。

\medskip\noindent
接着，设 $a : \mathcal{G} \to \mathcal{F}$ 是 Abel 层的映射，且
$f_{small, *}a$ 是满射。设 $s \in \mathcal{F}(U)$，其中 $U$ 如上。
用上面的 $W_i$、$V_i$，可知只需证明 $s|_{W_i}$ 在 \'etale 局部是
$\mathcal{G}$ 的某个截面在 $a$ 下的像。由于 $\mathcal{F}(W_i)$ 是
$\mathcal{F}(X \times_Y V_i)$ 的直和项，只需证明：对任意
$V \in \Ob(Y_\etale)$，任意元素 $s \in \mathcal{F}(X \times_Y V)$
在 $X \times_Y V$ 的 \'etale 局部是 $\mathcal{G}$ 的某个截面在 $a$ 下的像。
由于
$\mathcal{F}(X \times_Y V) = f_{small, *}\mathcal{F}(V)$
，根据假设存在覆盖 $\{V_j \to V\}$，使得 $s$ 是
$s_j \in f_{small, *}\mathcal{G}(V_j) = \mathcal{G}(X \times_Y V_j)$ 的像。
这证明了 $f_{small, *}$ 反映满射。

\medskip\noindent
(2)、(3) 由 (1) 形式地得出；见 Modules on Sites，引理
\ref{sites-modules-lemma-reflect-surjections}。
\end{proof}

\begin{lemma}
\label{etale-cohomology-lemma-locally-quasi-finite-A}
设 $f : X \to Y$ 是分离且局部拟有限的概形态射。则上面的性质 (A) 成立。
\end{lemma}

\begin{proof}
设 $U \to X$ 是 \'etale 态射，且 $u \in U$。几何断言 (A) 直接归结为
$U$ 与 $Y$ 都是仿射概形的情形。以 $x \in X$ 和 $y \in Y$ 表示 $u$ 的像。
由于 $X \to Y$ 局部拟有限，而 $U \to X$ 局部拟有限（见 Morphisms，引理
\ref{morphisms-lemma-etale-locally-quasi-finite}），可知 $U \to Y$ 局部拟有限
（见 Morphisms，引理 \ref{morphisms-lemma-composition-quasi-finite}）。此外，
$X \to Y$ 与 $U \to Y$ 都分离。因此 More on Morphisms，引理
\ref{more-morphisms-lemma-etale-splits-off-quasi-finite-part-technical-variant}
可应用于这两个态射。这意味着可以选择 \'etale 邻域
$(V, v) \to (Y, y)$，使得
$$
X \times_Y V = W \amalg R, \quad
U \times_Y V = W' \amalg R'
$$
以及点 $w \in W$、$w' \in W'$，使得
\begin{enumerate}
\item $W$、$R$ 在 $X \times_Y V$ 中开且闭；
\item $W'$、$R'$ 在 $U \times_Y V$ 中开且闭；
\item $W \to V$ 与 $W' \to V$ 有限；
\item $w$、$w'$ 映到 $v$；
\item $\kappa(v) \subset \kappa(w)$ 和 $\kappa(v) \subset \kappa(w')$
是纯不可分扩张；
\item $W$ 或 $W'$ 中没有其他点映到 $v$。
\end{enumerate}
有如下交换图：
$$
\xymatrix{
U \ar[d] & U \times_Y V \ar[l] \ar[d] & W' \amalg R' \ar[d] \ar[l] \\
X \ar[d] & X \times_Y V \ar[l] \ar[d] & W \amalg R \ar[l] \\
Y & V \ar[l]
}
$$
缩小 $V$ 后，可以假设 $W'$ 映入 $W$：只需去掉 $R$ 的逆像（它在 $W'$ 中）的像；
这是闭集（因为 $W' \to V$ 有限），且不含 $v$。于是 $W' \to W$ 有限，
因为 $W \to V$ 与 $W' \to V$ 都有限。因此 $W' \to W$ 是有限 \'etale 态射，
并且在 $w$ 上的纤维中恰有一个点满足 $\kappa(w) = \kappa(w')$。故
$W' \to W$ 在某个开邻域 $W^\circ$（它包含 $w$）上是同构；见 \'Etale Morphisms，
引理 \ref{etale-lemma-finite-etale-one-point}。由于 $W \to V$ 有限，
$W \setminus W^\circ$ 的像是闭子集 $T$（在 $V$ 中），且不含 $v$。因此，把
$V$ 替换成 $V \setminus T$ 后，可以假设 $W' \to W$ 是同构。此时分解
$X \times_Y V = W \amalg R$ 和态射 $W \to U$ 正是所需，证明完成。
\end{proof}

\begin{lemma}
\label{etale-cohomology-lemma-integral-A}
设 $f : X \to Y$ 是整概形态射。则性质 (A) 成立。
\end{lemma}

\begin{proof}
设 $U \to X$ 是 \'etale 态射，且 $u \in U$ 是一点。需要找到 \'etale 态射
$V \to Y$、不交并分解 $X \times_Y V = W \amalg W'$，以及 $X$-态射
$W \to U$，使得 $u$ 属于其像。可以缩小 $U$ 和 $Y$，并假设 $U$ 与 $Y$
仿射。此时 $X$ 也仿射，因为整态射按定义是仿射的。写成
$Y = \Spec(A)$、$X = \Spec(B)$ 和 $U = \Spec(C)$。那么 $A \to B$ 是整环同态，
而 $B \to C$ 是 \'etale 环同态。由 Algebra，引理 \ref{algebra-lemma-etale}，
可以找到有限 $A$-子代数 $B' \subset B$ 和 \'etale 环同态 $B' \to C'$，使得
$C = B \otimes_{B'} C'$。因此，问题归结为 \'etale 态射
$U' = \Spec(C') \to X' = \Spec(B')$，它位于有限态射 $X' \to Y$ 上。在这种情形，结论由引理
\ref{etale-cohomology-lemma-locally-quasi-finite-A} 得出。
\end{proof}

\begin{lemma}
\label{etale-cohomology-lemma-when-push-pull-surjective}
设 $f : X \to Y$ 是概形态射。以
$f_{small} :
\Sh(X_\etale)
\to
\Sh(Y_\etale)$
表示相应的小 \'etale 拓扑斯态射。假设下列条件至少有一个成立：
\begin{enumerate}
\item $f$ 是整态射；
\item $f$ 分离且局部拟有限。
\end{enumerate}
则函子
$f_{small, *} :
\textit{Ab}(X_\etale)
\to
\textit{Ab}(Y_\etale)$
具有下列性质：
\begin{enumerate}
\item 映射
$f_{small}^{-1}f_{small, *}\mathcal{F} \to \mathcal{F}$
总是满射；
\item $f_{small, *}$ 是忠实的；
\item $f_{small, *}$ 反映单射和满射。
\end{enumerate}
\end{lemma}

\begin{proof}
合用引理 \ref{etale-cohomology-lemma-locally-quasi-finite-A}、
\ref{etale-cohomology-lemma-integral-A} 和 \ref{etale-cohomology-lemma-property-A-implies}。
\end{proof}



\section{性质 (B)}
\label{etale-cohomology-section-B}

\noindent
性质 (B) 的定义见第 \ref{etale-cohomology-section-monomorphisms} 节。

\begin{lemma}
\label{etale-cohomology-lemma-property-B-implies}
设 $f : X \to Y$ 是概形态射。假设 (B) 成立。则函子
$f_{small, *} :
\Sh(X_\etale)
\to
\Sh(Y_\etale)$
把满射变为满射。
\end{lemma}

\begin{proof}
这由 Sites，引理 \ref{sites-lemma-weaker} 得出。
\end{proof}

\begin{lemma}
\label{etale-cohomology-lemma-simplify-B}
设 $f : X \to Y$ 是概形态射。假设
\begin{enumerate}
\item $V \to Y$ 是概形的 \'etale 态射；
\item $\{U_i \to X \times_Y V\}$ 是 \'etale 覆盖；
\item $v \in V$ 是一点。
\end{enumerate}
假设对任意这样的数据，都存在 \'etale 邻域 $(V', v') \to (V, v)$、不交并分解
$X \times_Y V' = \coprod W'_i$，以及态射 $W'_i \to U_i$
（在 $X \times_Y V$ 上）。则性质 (B) 成立。
\end{lemma}

\begin{proof}
略。
\end{proof}

\begin{lemma}
\label{etale-cohomology-lemma-finite-B}
设 $f : X \to Y$ 是有限概形态射。则性质 (B) 成立。
\end{lemma}

\begin{proof}
考虑 \'etale 态射 $V \to Y$、\'etale 覆盖 $\{U_i \to X \times_Y V\}$，以及
$v \in V$。需要找到引理 \ref{etale-cohomology-lemma-simplify-B} 中的 $V' \to V$、分解和映射。
可以缩小 $V$ 和 $Y$，因而可以假设 $V$ 与 $Y$ 仿射。由于 $X$ 在 $Y$ 上有限，
这也蕴含 $X$ 仿射。在证明过程中，可以（有限多次）把 $(V, v)$ 替换为
\'etale 邻域 $(V', v')$，并相应地把覆盖 $\{U_i \to X \times_Y V\}$ 替换为
$\{V' \times_V U_i \to X \times_Y V'\}$。

\medskip\noindent
由于 $X \times_Y V \to V$ 有限，存在有限多个（两两不同的）点
$x_1, \ldots, x_n \in X \times_Y V$ 映到 $v$。把 More on Morphisms，引理
\ref{more-morphisms-lemma-etale-splits-off-quasi-finite-part-technical-variant}
应用于 $X \times_Y V \to V$ 及诸点 $x_1, \ldots, x_n$（它们位于 $v$ 上），得到
\'etale 邻域 $(V', v') \to (V, v)$，使得
$$
X \times_Y V' = R \amalg \coprod T_a
$$
，其中 $T_a \to V'$ 有限，并且恰有一个点 $p_a$ 位于 $v'$ 上；此外
$\kappa(v') \subset \kappa(p_a)$ 是纯不可分扩张，而 $R \to V'$ 在 $v'$ 上的纤维为空。
由于 $X \to Y$ 有限，$R \to V'$ 也有限。因此，缩小 $V'$ 后可以假设
$R = \emptyset$。于是可以假设
$X \times_Y V = X_1 \amalg \ldots \amalg X_n$，其中恰有一个点
$x_l \in X_l$ 位于 $v$ 上，并且 $\kappa(v) \subset \kappa(x_l)$ 是纯不可分扩张。
注意，这个性质在加细 \'etale 邻域 $(V, v)$ 时保持不变。

\medskip\noindent
对每个 $l$，选择一个 $i_l$ 和一点 $u_l \in U_{i_l}$，它映到 $x_l$。现在，把性质 (A)
应用于有限态射 $X \times_Y V \to V$、\'etale 态射
$U_{i_l} \to X \times_Y V$ 以及诸点 $u_l$。引理 \ref{etale-cohomology-lemma-integral-A}
保证可以这样做。由此得到 \'etale 邻域 $(V', v') \to (V, v)$ 和分解
$$
X \times_Y V' = W_l \amalg R_l
$$
，以及 $X$-态射 $a_l : W_l \to U_{i_l}$，其像包含 $u_{i_l}$。图示如下：
$$
\xymatrix{
& & & U_{i_l} \ar[d] & \\
W_l \ar[rrru] \ar[r] & W_l \amalg R_l \ar@{=}[r] &
X \times_Y V' \ar[r] \ar[d] &
X \times_Y V \ar[r] \ar[d] & X \ar[d] \\
& & V' \ar[r] & V \ar[r] & Y
}
$$
把 $(V, v)$ 替换为 $(V', v')$ 后，可知每个 $x_l$ 都包含在开闭邻域 $W_l$ 中，
使得包含态射 $W_l \to X \times_Y V$ 经 $U_i \to X \times_Y V$ 分解
（对某个 $i$）。把 $W_l$ 替换为 $W_l \cap X_l$ 后，可知这些开闭集彼此不交，
而且 $\{x_1, \ldots, x_n\} \subset W_1 \cup \ldots \cup W_n$。
由于 $X \times_Y V \to V$ 有限，可以缩小 $V$ 并假设
$X \times_Y V = W_1 \amalg \ldots \amalg W_n$，正是所需。
\end{proof}

\begin{lemma}
\label{etale-cohomology-lemma-integral-B}
设 $f : X \to Y$ 是整概形态射。则性质 (B) 成立。
\end{lemma}

\begin{proof}
考虑 \'etale 态射 $V \to Y$、\'etale 覆盖 $\{U_i \to X \times_Y V\}$，以及
$v \in V$。需要找到引理 \ref{etale-cohomology-lemma-simplify-B} 中的 $V' \to V$、分解和映射。
可以缩小 $V$ 与 $Y$，因而可以假设 $V$ 与 $Y$ 仿射。由于 $X$ 在 $Y$ 上整，
这也蕴含 $X$ 与 $X \times_Y V$ 仿射。可以加细覆盖
$\{U_i \to X \times_Y V\}$，因而可以假设
$\{U_i \to X \times_Y V\}_{i = 1, \ldots, n}$ 是标准 \'etale 覆盖。写成
$Y = \Spec(A)$、$X = \Spec(B)$、$V = \Spec(C)$ 和 $U_i = \Spec(B_i)$。
那么 $A \to B$ 是整环同态，而 $B \otimes_A C \to B_i$ 是 \'etale 环同态。
由 Algebra，引理 \ref{algebra-lemma-etale}，可以找到有限 $A$-子代数
$B' \subset B$ 和 \'etale 环同态 $B' \otimes_A C \to B'_i$
（其中 $i = 1, \ldots, n$），使得 $B_i = B \otimes_{B'} B'_i$。
因此，问题归结为 \'etale 覆盖
$\{\Spec(B'_i) \to X' \times_Y V\}_{i = 1, \ldots, n}$，其中
$X' = \Spec(B')$ 在 $Y$ 上有限。在这种情形，结论由引理 \ref{etale-cohomology-lemma-finite-B} 得出。
\end{proof}

\begin{lemma}
\label{etale-cohomology-lemma-what-integral}
设 $f : X \to Y$ 是概形态射。假设 $f$ 是整态射（例如有限态射）。则
\begin{enumerate}
\item $f_{small, *}$ 把满射变为满射（对集合层和 Abel 层都是如此）；
\item $f_{small}^{-1}f_{small, *}\mathcal{F} \to \mathcal{F}$
对任意 Abel 层 $\mathcal{F}$（在 $X_\etale$ 上）都是满射；
\item
$f_{small, *} :
\textit{Ab}(X_\etale)
\to
\textit{Ab}(Y_\etale)$
是忠实的，并反映单射和满射；
\item
$f_{small, *} :
\textit{Ab}(X_\etale)
\to
\textit{Ab}(Y_\etale)$
正合。
\end{enumerate}
\end{lemma}

\begin{proof}
(2)、(3) 已在引理 \ref{etale-cohomology-lemma-when-push-pull-surjective} 中看到。
(1) 由引理 \ref{etale-cohomology-lemma-integral-B} 和 \ref{etale-cohomology-lemma-property-B-implies} 得出。
(4) 是 (1) 的推论；见 Modules on Sites，引理
\ref{sites-modules-lemma-exactness}。
\end{proof}








\section{性质 (C)}
\label{etale-cohomology-section-C}

\noindent
性质 (C) 的定义见第 \ref{etale-cohomology-section-monomorphisms} 节。

\begin{lemma}
\label{etale-cohomology-lemma-property-C-implies}
设 $f : X \to Y$ 是概形态射。假设 (C) 成立。则函子
$f_{small, *} :
\Sh(X_\etale)
\to
\Sh(Y_\etale)$
反映单射和满射。
\end{lemma}

\begin{proof}
这由 Sites，引理 \ref{sites-lemma-cover-from-below} 得出。略去如下验证：
性质 (C) 蕴含函子
$Y_\etale \to X_\etale$, $V \mapsto X \times_Y V$
满足 Sites，引理 \ref{sites-lemma-cover-from-below} 的假设。
\end{proof}

\begin{remark}
\label{etale-cohomology-remark-property-C-strong}
若 $f : X \to Y$ 是开浸入，则性质 (C) 成立。事实上，若
$U \in \Ob(X_\etale)$，则也可以把 $U$ 看作 $Y_\etale$ 的对象，并且
$U \times_Y X = U$。因此，性质 (C) 并不蕴含 $f_{small, *}$ 正合，
因为开浸入一般不具有该正合性。
\end{remark}

\begin{lemma}
\label{etale-cohomology-lemma-property-C-closed-implies}
设 $f : X \to Y$ 是概形态射。假设对任意 \'etale 态射 $V \to Y$，都有
\begin{enumerate}
\item $X \times_Y V \to V$ 具有性质 (C)；
\item $X \times_Y V \to V$ 是闭映射。
\end{enumerate}
则函子
$Y_\etale \to X_\etale$, $V \mapsto X \times_Y V$
几乎余连续；见 Sites，定义 \ref{sites-definition-almost-cocontinuous}。
\end{lemma}

\begin{proof}
设 $V \to Y$ 是 $Y_\etale$ 的对象，并设
$\{U_i \to X \times_Y V\}_{i \in I}$ 是 $X_\etale$ 的覆盖。由假设 (1)，
对每个 $i$，可以找到 \'etale 态射 $h_i : V_i \to V$ 和满态射
$X \times_Y V_i \to U_i$（在 $X \times_Y V$ 上）。注意
$\bigcup h_i(V_i) \subset V$ 是包含闭集
$Z = \Im(X \times_Y V \to V)$ 的开集。令
$h_0 : V_0 = V \setminus Z \to V$ 为开浸入。显然，
$\{V_i \to V\}_{i \in I \cup \{0\}}$ 是 \'etale 覆盖，使得对每个
$i \in I \cup \{0\}$，或者 $V_i \times_Y X = \emptyset$（即 $i = 0$），
或者 $V_i \times_Y X \to V \times_Y X$ 经 $U_i \to X \times_Y V$ 分解
（若 $i \not = 0$）。因此函子 $Y_\etale \to X_\etale$ 几乎余连续。
\end{proof}

\begin{lemma}
\label{etale-cohomology-lemma-integral-homeo-onto-image-C}
设 $f : X \to Y$ 是整概形态射，它给出 $X$ 与 $Y$ 的一个闭子集之间的同胚。
则性质 (C) 成立。
\end{lemma}

\begin{proof}
设 $g : U \to X$ 是 \'etale 态射。需要找到一个对象 $V \to Y$（属于
$Y_\etale$）和满态射 $X \times_Y V \to U$（在 $X$ 上）。假设对每个
$u \in U$，都能找到一个对象 $V_u \to Y$（属于 $Y_\etale$）和态射
$h_u : X \times_Y V_u \to U$（在 $X$ 上），且 $u \in \Im(h_u)$。
那么可以取 $V = \coprod V_u$ 和 $h = \coprod h_u$，从而得证。因此，给定一点
$u \in U$，我们来寻找如上的一对 $(V_u, h_u)$。为此可以缩小 $U$ 并假设
$U$ 仿射。在这种情形，$g : U \to X$ 局部拟有限。设
$g^{-1}(g(\{u\})) = \{u, u_2, \ldots, u_n\}$。由于不存在特殊化
$u_i \leadsto u$，可以用一个仿射邻域替换 $U$，使得
$g^{-1}(g(\{u\})) = \{u\}$。

\medskip\noindent
像 $g(U) \subset X$ 是开集，故 $f(g(U))$ 在 $Y$ 中局部闭。选择开集
$V \subset Y$，使得 $f(g(U)) = f(X) \cap V$。于是 $g$ 经
$X \times_Y V$ 分解，并且所得的 $\{U \to X \times_Y V\}$ 是 \'etale 覆盖。
由于 $f$ 具有性质 (B)，见引理 \ref{etale-cohomology-lemma-integral-B}，存在 \'etale 覆盖
$\{V_j \to V\}$，使得 $X \times_Y V_j \to X \times_Y V$ 经 $U$ 分解。
这蕴含 $V' = \coprod V_j$ 在 $Y$ 上是 \'etale 的，并且存在态射
$h : X \times_Y V' \to U$，其像满射到 $g(U)$。由于 $u$ 是其纤维中唯一的点，
它必属于 $h$ 的像，从而得证。
\end{proof}

\noindent
我们建议读者把下面的引理看作通往整且普遍单射态射的 $f_{small, *}$ 之最终真相的
旅途中的驿站\footnote{驿站是人们在长途旅行中停下来进食和休息的地方。}。

\begin{lemma}
\label{etale-cohomology-lemma-integral-universally-injective}
设 $f : X \to Y$ 是概形态射。假设 $f$ 普遍单射且整（例如闭浸入）。则
\begin{enumerate}
\item
$f_{small, *} :
\Sh(X_\etale)
\to
\Sh(Y_\etale)$
反映单射和满射；
\item
$f_{small, *} :
\Sh(X_\etale)
\to
\Sh(Y_\etale)$
与推出和余等化子可交换（更一般地，与有限连通余极限可交换）；
\item $f_{small, *}$ 把满射变为满射（对集合层和 Abel 层都是如此）；
\item 映射
$f_{small}^{-1}f_{small, *}\mathcal{F} \to \mathcal{F}$
对任意层 $\mathcal{F}$（集合层或 Abel 群层，在 $X_\etale$ 上）都是满射；
\item 函子 $f_{small, *}$ 是忠实的（对集合层和 Abel 层都是如此）；
\item
$f_{small, *} :
\textit{Ab}(X_\etale)
\to
\textit{Ab}(Y_\etale)$
正合；
\item 函子 $Y_\etale \to X_\etale$、$V \mapsto X \times_Y V$ 几乎余连续。
\end{enumerate}
\end{lemma}

\begin{proof}
由引理 \ref{etale-cohomology-lemma-integral-A}、\ref{etale-cohomology-lemma-integral-B} 和
\ref{etale-cohomology-lemma-integral-homeo-onto-image-C}，可知态射 $f$ 具有性质 (A)、(B) 和 (C)。
此外，由引理 \ref{etale-cohomology-lemma-property-C-closed-implies}，可知函子
$Y_\etale \to X_\etale$ 几乎余连续。现在有
\begin{enumerate}
\item 由引理 \ref{etale-cohomology-lemma-property-C-implies}，性质 (C) 蕴含 (1)；
\item 由 Sites，引理 \ref{sites-lemma-morphism-of-sites-almost-cocontinuous}，
几乎余连续蕴含 (2)；
\item 由引理 \ref{etale-cohomology-lemma-property-B-implies}，性质 (B) 蕴含 (3)。
\end{enumerate}
性质 (4)、(5)、(6) 从前三项形式地得出；见 Sites，引理
\ref{sites-lemma-exactness-properties} 和 Modules on Sites，引理
\ref{sites-modules-lemma-exactness}。性质 (7) 已在上面看到。
\end{proof}




\section{小 \'etale 网站的拓扑不变性}
\label{etale-cohomology-section-topological-invariance}

\noindent
下一个定理表明，小 \'etale 网站在如下意义下是拓扑不变量：若
$f : X \to Y$ 是作为普遍同胚的概形态射，则网站
$X_\etale \cong Y_\etale$。这改进了 \'Etale Morphisms，定理
\ref{etale-theorem-remarkable-equivalence} 的结果。我们先对态射证明该结论，
然后陈述范畴层面的结论。

\begin{theorem}
\label{etale-cohomology-theorem-etale-topological}
设 $X$ 和 $Y$ 是基概形 $S$ 上的两个概形。设 $S' \to S$ 是普遍同胚。
以 $X'$（相应地以 $Y'$）表示到 $S'$ 的基变换。若 $X$ 在 $S$ 上是 \'etale 的，
则映射
$$
\Mor_S(Y, X) \longrightarrow \Mor_{S'}(Y', X')
$$
是双射。
\end{theorem}

\begin{proof}
沿 $Y \to S$ 作基变换后，可以假设 $Y = S$。因此，可以并且确实假设
$X$ 与 $Y$ 都在 $S$ 上是 \'etale 的。换言之，定理断言：基变换函子从
$S$ 上 \'etale 概形的范畴到 $S'$ 上 \'etale 概形的范畴是全忠实函子。

\medskip\noindent
考虑遗忘函子
\begin{equation}
\label{etale-cohomology-equation-descent-etale-forget}
\begin{matrix}
\text{关于 }S'/S\text{ 的下降数据 }(X', \varphi') \\
\text{，其中 }X'\text{ 在 }S'\text{ 上 \'etale}
\end{matrix}
\longrightarrow
\text{在 }S'\text{ 上 \'etale 的概形 }X'
\end{equation}
我们断言这个函子是等价。另一方面，函子
\begin{equation}
\label{etale-cohomology-equation-descent-etale}
\text{在 }S\text{ 上 \'etale 的概形 }X \longrightarrow
\begin{matrix}
\text{关于 }S'/S\text{ 的下降数据 }(X', \varphi') \\
\text{，其中 }X'\text{ 在 }S'\text{ 上 \'etale}
\end{matrix}
\end{equation}
由 \'Etale Morphisms，引理 \ref{etale-lemma-fully-faithful-cases} 是全忠实的。
因此，该断言蕴含本定理。

\medskip\noindent
证明该断言。回忆，普遍同胚等价于整、普遍单射且满的态射；见 Morphisms，引理
\ref{morphisms-lemma-universal-homeomorphism}。特别地，由 Morphisms，引理
\ref{morphisms-lemma-universally-injective}，对角态射
$\Delta : S' \to S' \times_S S'$ 是加厚。因此，由 \'Etale Morphisms，定理
\ref{etale-theorem-etale-topological}，可知给定 \'etale 态射 $X' \to S'$，
存在唯一同构
$$
\varphi' : X' \times_S S' \to S' \times_S X'
$$
，它是 $S' \times_S S'$ 上 \'etale 概形的同构，并且沿 $\Delta$ 拉回后成为
$\text{id} : X' \to X'$（在 $S'$ 上）。由于
$S' \to S' \times_S S' \times_S S'$ 也是加厚（它是双射且为闭浸入），
可知 $(X', \varphi')$ 是相对于 $S'/S$ 的下降数据。$\varphi'$ 构造的典范性表明，
它与 $S'$ 上 \'etale 概形之间的态射相容。换言之，得到函子
(\ref{etale-cohomology-equation-descent-etale-forget}) 的拟逆 $X' \mapsto (X', \varphi')$。
这证明了断言，并完成了定理的证明。
\end{proof}

\begin{theorem}
\label{etale-cohomology-theorem-topological-invariance}
\begin{reference}
\cite[IV Theorem 18.1.2]{EGA}
\end{reference}
设 $f : X \to Y$ 是概形态射。假设 $f$ 整、普遍单射且满
（即 $f$ 是普遍同胚；见 Morphisms，引理
\ref{morphisms-lemma-universal-homeomorphism}）。函子
$$
V \longmapsto V_X = X \times_Y V
$$
给出范畴等价
$$
\{
\text{在 }Y\text{ 上 \'etale 的概形 }V
\}
\leftrightarrow
\{
\text{在 }X\text{ 上 \'etale 的概形 }U
\}
$$
\end{theorem}

\noindent
我们给出两个证明。第一个证明使用拟紧、分离、\'etale 态射相对于满整态射的下降
有效性。第二个证明使用本章前面讨论的性质 (A)、(B)、(C)。

\begin{proof}[第一个证明]
由定理 \ref{etale-cohomology-theorem-etale-topological}，可知该函子全忠实。还需证明它本质满。
设 $U \to X$ 是概形的 \'etale 态射。

\medskip\noindent
假设 $U$ 与 $Y$ 仿射时结论成立。在这种情形，选择仿射开覆盖
$U = \bigcup U_i$，使每个 $U_i$ 映入 $Y$ 的一个仿射开集。根据假设
（仿射情形），可以找到 \'etale 态射 $V_i \to Y$，使得
$X \times_Y V_i \cong U_i$（作为 $X$ 上的概形）。令 $V_{i, i'} \subset V_i$ 是开子概形，
其底拓扑空间对应于 $U_i \cap U_{i'}$。由于有同构
$$
X \times_Y V_{i, i'} \cong U_i \cap U_{i'} \cong X \times_Y V_{i', i}
$$
（作为 $X$ 上的概形），由全忠实性得到同构
$\theta_{i, i'} : V_{i, i'} \to V_{i', i}$（它是 $Y$ 上概形的同构）。略去这些同构满足 Schemes，第
\ref{schemes-section-glueing-schemes} 节的余圈条件的验证。应用 Schemes，引理
\ref{schemes-lemma-glue-schemes}，得到概形 $V \to Y$；它是把概形 $V_i$ 沿同一
$\theta_{i, i'}$ 粘合得到的。根据构造，$V \to Y$ 显然是 \'etale 的，且
$X \times_Y V \cong U$。

\medskip\noindent
因此，只需在 $U$ 与 $Y$ 仿射时证明该引理。回忆，在定理
\ref{etale-cohomology-theorem-etale-topological} 的证明中，我们已证明 $U$ 带有唯一下降数据
$(U, \varphi)$（相对于 $X/Y$）。由 \'Etale Morphisms，命题
\ref{etale-proposition-effective}（它适用是因为通过归约到仿射情形，
$U \to X$ 不仅是 \'etale 的，而且拟紧且分离），存在 \'etale 态射
$V \to Y$，使得 $X \times_Y V \cong U$，证明完成。
\end{proof}

\begin{proof}[第二个证明]
由定理 \ref{etale-cohomology-theorem-etale-topological}，可知该函子全忠实。还需证明它本质满。
设 $U \to X$ 是概形的 \'etale 态射。

\medskip\noindent
假设 $U$ 与 $Y$ 仿射时结论成立。在这种情形，选择仿射开覆盖
$U = \bigcup U_i$，使每个 $U_i$ 映入 $Y$ 的一个仿射开集。根据假设
（仿射情形），可以找到 \'etale 态射 $V_i \to Y$，使得
$X \times_Y V_i \cong U_i$（作为 $X$ 上的概形）。令
$V_{i, i'} \subset V_i$ 是开子概形，其底拓扑空间对应于
$U_i \cap U_{i'}$。由于有同构
$$
X \times_Y V_{i, i'} \cong U_i \cap U_{i'} \cong X \times_Y V_{i', i}
$$
（作为 $X$ 上的概形），由全忠实性得到同构
$\theta_{i, i'} : V_{i, i'} \to V_{i', i}$（它是 $Y$ 上概形的同构）。
略去这些同构满足 Schemes，第 \ref{schemes-section-glueing-schemes} 节的余圈条件
的验证。应用 Schemes，引理 \ref{schemes-lemma-glue-schemes}，得到概形
$V \to Y$；它是把概形 $V_i$ 沿同一 $\theta_{i, i'}$ 粘合得到的。根据构造，
$V \to Y$ 显然是 \'etale 的，且 $X \times_Y V \cong U$。

\medskip\noindent
因此，只需在 $X$ 与 $Y$ 仿射时证明函子
\begin{equation}
\label{etale-cohomology-equation-affine-etale}
\{
\text{在 }Y\text{ 上 \'etale 的仿射概形 }V
\}
\leftrightarrow
\{
\text{在 }X\text{ 上 \'etale 的仿射概形 }U
\}
\end{equation}
本质满。

\medskip\noindent
设 $U \to X$ 是在 $X$ 上 \'etale 的仿射概形。需要找到 \'etale（且仿射）的
$V \to Y$，使得 $X \times_Y V$ 同构于 $U$（在 $X$ 上）。注意，仿射概形间的
\'etale 态射具有普遍有界的纤维；见 Morphisms，引理
\ref{morphisms-lemma-etale-locally-quasi-finite} 和
\ref{morphisms-lemma-locally-quasi-finite-qc-source-universally-bounded}。
因此，可以对整数 $n$ 作归纳，它限制 $U \to X$ 之纤维的次数。\'etale 态射情形下
这个整数的描述见 Morphisms，引理
\ref{morphisms-lemma-etale-universally-bounded}。若 $n = 1$，则
$U \to X$ 是开浸入（见 \'Etale Morphisms，定理
\ref{etale-theorem-etale-radicial-open}），结论显然。假设 $n > 1$。

\medskip\noindent
由引理 \ref{etale-cohomology-lemma-integral-homeo-onto-image-C}，存在概形的 \'etale 态射
$W \to Y$ 和满态射 $W_X \to U$（在 $X$ 上）。由于 $U$ 拟紧，
可以用 $W$ 的有限多个仿射开集的不交并替换 $W$，因而也可以假设 $W$ 仿射。
有如下图：
$$
\xymatrix{
U \ar[d] & U \times_Y W \ar[l] \ar[d] & W_X \amalg R \ar@{=}[l]\\
X \ar[d] & W_X \ar[l] \ar[d] \\
Y & W \ar[l]
}
$$
这个不交并分解之所以出现，是因为按构造，仿射概形的 \'etale 态射
$U \times_Y W \to W_X$ 有截面。现在可知，态射 $R \to X \times_Y W$ 是仿射概形的
\'etale 态射，其纤维次数普遍不超过 $n - 1$。因此，由归纳假设，存在 \'etale 概形
$V' \to W$，使得 $R \cong W_X \times_W V'$。取 $V'' = W \amalg V'$，得到概形
$V''$（它在 $W$ 上是 \'etale 的），其到 $W_X$ 的基变换同构于 $U \times_Y W$
（在 $X \times_Y W$ 上）。

\medskip\noindent
此时可以使用下降找到 $V$（在 $Y$ 上），使其到 $X$ 的基变换同构于 $U$
（在 $X$ 上）。
具体地，由与普遍同胚
$X \times_Y (W \times_Y W) \to (W \times_Y W)$ 对应的函子
(\ref{etale-cohomology-equation-affine-etale}) 的全忠实性，存在唯一同构
$\varphi : V'' \times_Y W \to W \times_Y V''$，它到
$X \times_Y (W \times_Y W)$ 的基变换是 $U \times_Y W$ 在
$X \times_Y W$ 上的典范下降数据。特别地，$\varphi$ 满足余圈条件。因此，由
Descent，引理 \ref{descent-lemma-affine}，可知 $\varphi$ 是有效的（回忆上面的概形
全都仿射）。于是得到 $V \to Y$ 和同构 $V'' \cong W \times_Y V$，使得
$W \times_Y V/W/Y$ 上的典范下降数据与 $\varphi$ 相同。注意，由 Descent，引理
\ref{descent-lemma-descending-property-etale}，$V \to Y$ 是 \'etale 的。此外，
存在同构 $V_X \cong U$，它来自下列同构的下降：
$$
V_X \times_X W_X =
X \times_Y V \times_Y W =
(X \times_Y W) \times_W (W \times_Y V) \cong
W_X  \times_W V'' \cong U \times_Y W
$$
；这个同构由构造给出。略去一些细节。
\end{proof}

\begin{remark}
\label{etale-cohomology-remark-affine-inside-equivalence}
在定理 \ref{etale-cohomology-theorem-topological-invariance} 的情形下，$V \mapsto V_X$ 还诱导
两类对象之间的等价：一类是 \'etale 态射 $V \to Y$ 且 $V$ 仿射，另一类是
\'etale 态射 $U \to X$ 且 $U$ 仿射。例如，这由 Limits，命题
\ref{limits-proposition-affine} 得出。
\end{remark}

\begin{proposition}[\'Etale 上同调的拓扑不变性]
\label{etale-cohomology-proposition-topological-invariance}
设 $X_0 \to X$ 是概形的普遍同胚（例如由幂零理想层定义的闭浸入）。则
\begin{enumerate}
\item \'etale 网站 $X_\etale$ 与 $(X_0)_\etale$ 同构；
\item \'etale 拓扑斯 $\Sh(X_\etale)$ 与 $\Sh((X_0)_\etale)$ 等价；
\item $H^q_\etale(X, \mathcal{F}) = H^q_\etale(X_0, \mathcal{F}|_{X_0})$
对所有 $q$ 以及任意 Abel 层 $\mathcal{F}$（在 $X_\etale$ 上）成立。
\end{enumerate}
\end{proposition}

\begin{proof}
范畴等价 $X_\etale \to (X_0)_\etale$ 由定理
\ref{etale-cohomology-theorem-topological-invariance} 给出。略去在此等价下 \'etale 覆盖相互对应的
证明。因此 (1) 成立，而 (2)、(3) 从 (1) 形式地得出。
\end{proof}






\section{闭浸入与推前}
\label{etale-cohomology-section-closed-immersions}

\noindent
在以正确的一般性陈述并证明命题
\ref{etale-cohomology-proposition-closed-immersion-pushforward} 之前，我们先对闭浸入简要陈述并证明它。
事实上，前面的一些论证在闭浸入情形下更容易理解，所以在此以简化形式重复它们。

\medskip\noindent
本节余下部分中，$i : Z \to X$ 是闭浸入。函子
$$
\Sch/X \longrightarrow \Sch/Z, \quad
U \longmapsto U_Z = Z \times_X U
$$
如上记为 $U \mapsto U_Z$。由于闭浸入在任意基变换下保持，概形 $U_Z$ 是
$U$ 的闭子概形。

\begin{lemma}
\label{etale-cohomology-lemma-closed-immersion-almost-full}
设 $i : Z \to X$ 是概形的闭浸入。设 $U, U'$ 是 $X$ 上的 \'etale 概形。
设 $h : U_Z \to U'_Z$ 是 $Z$ 上的态射。则存在下图
$$
\xymatrix{
U & W \ar[l]_a \ar[r]^b & U'
}
$$
，它位于 $X_\etale$ 中，使得 $a_Z : W_Z \to U_Z$ 是同构，且
$h = b_Z \circ (a_Z)^{-1}$。
\end{lemma}

\begin{proof}
考虑概形 $M = U \times_X U'$。$\Gamma_h \subset M_Z$（即 $h$ 的图像）是开集。
例如，这是因为 $\Gamma_h$ 是 \'etale 态射
$\text{pr}_{1, Z} : M_Z \to U_Z$ 的某个截面的像；见 \'Etale Morphisms，命题
\ref{etale-proposition-properties-sections}。因此，存在开子概形 $W \subset M$，
它与闭子集 $M_Z$ 的交是 $\Gamma_h$。置 $a = \text{pr}_1|_W$ 和
$b = \text{pr}_2|_W$。
\end{proof}

\begin{lemma}
\label{etale-cohomology-lemma-closed-immersion-almost-essentially-surjective}
设 $i : Z \to X$ 是概形的闭浸入。设 $V \to Z$ 是概形的 \'etale 态射。
存在 \'etale 态射 $U_i \to X$ 和态射 $U_{i, Z} \to V$，使得
$\{U_{i, Z} \to V\}$ 是 $V$ 的 Zariski 覆盖。
\end{lemma}

\begin{proof}
由于只需找到 $V$ 的 Zariski 覆盖，其成员是形如 $U_Z$ 的概形，且 $U$ 在
$X$ 上是 \'etale 的，因此可以在 $X$ 和 $V$ 上作 Zariski 局部化。因此可以假设
$X$ 与 $V$ 仿射。在仿射情形，这就是 Algebra，引理
\ref{algebra-lemma-lift-etale}。
\end{proof}

\noindent
若 $\overline{x} : \Spec(k) \to X$ 是 $X$ 的几何点，则或者 $\overline{x}$
（唯一地）经闭子概形 $Z$ 分解，或者 $Z_{\overline{x}} = \emptyset$。若
$\overline{x}$ 经 $Z$ 分解，就称 $\overline{x}$ 为 $Z$ 的几何点（它确实如此），
并用记号``$\overline{x} \in Z$''表示这一点。

\begin{lemma}
\label{etale-cohomology-lemma-stalk-pushforward-closed-immersion}
设 $i : Z \to X$ 是概形的闭浸入。设 $\mathcal{G}$ 是 $Z_\etale$ 上的集合层。
设 $\overline{x}$ 是 $X$ 的几何点。则
$$
(i_{small, *}\mathcal{G})_{\overline{x}} =
\left\{
\begin{matrix}
* & \text{若} & \overline{x} \not \in Z \\
\mathcal{G}_{\overline{x}} & \text{若} & \overline{x} \in Z
\end{matrix}
\right.
$$
其中 $*$ 表示单点集。
\end{lemma}

\begin{proof}
令 $U = X \setminus Z$。注意，$i_{small, *}\mathcal{G}|_{U_\etale} = *$ 是
$U$ 上 \'etale 层范畴中的终对象，即给每个在 $U$ 上 \'etale 的概形赋予单点集的层。
这说明了 $(i_{small, *}\mathcal{G})_{\overline{x}}$ 的取值
（当 $\overline{x} \not \in Z$ 时）。

\medskip\noindent
接着，假设 $\overline{x} \in Z$。注意
$$
(i_{small, *}\mathcal{G})_{\overline{x}}
=
\colim_{(U, \overline{u})} \mathcal{G}(U_Z)
$$
，另一方面
$$
\mathcal{G}_{\overline{x}}
=
\colim_{(V, \overline{v})} \mathcal{G}(V).
$$
令 $\mathcal{C}_1 = \{(U, \overline{u})\}^{opp}$ 为 $\overline{x}$ 在 $X$ 中的
\'etale 邻域范畴的反范畴，并令
$\mathcal{C}_2 = \{(V, \overline{v})\}^{opp}$ 为 $\overline{x}$ 在 $Z$ 中的
\'etale 邻域范畴的反范畴。典范映射
$$
\mathcal{G}_{\overline{x}}
\longrightarrow
(i_{small, *}\mathcal{G})_{\overline{x}}
$$
对应于函子 $F : \mathcal{C}_1 \to \mathcal{C}_2$，其中
$F(U, \overline{u}) = (U_Z, \overline{x})$。现在，引理
\ref{etale-cohomology-lemma-closed-immersion-almost-essentially-surjective} 和
\ref{etale-cohomology-lemma-closed-immersion-almost-full} 蕴含 $\mathcal{C}_1$ 在
$\mathcal{C}_2$ 中余终；见 Categories，定义 \ref{categories-definition-cofinal}。
因此，上面的箭头是同构；见 Categories，引理 \ref{categories-lemma-cofinal}。
\end{proof}

\begin{proposition}
\label{etale-cohomology-proposition-closed-immersion-pushforward}
设 $i : Z \to X$ 是概形的闭浸入。
\begin{enumerate}
\item 函子
$$
i_{small, *} :
\Sh(Z_\etale)
\longrightarrow
\Sh(X_\etale)
$$
是全忠实的，其本质像由如下集合层 $\mathcal{F}$ 构成：它在 $X_\etale$ 上，
且到 $X \setminus Z$ 的限制同构于 $*$；
\item 函子
$$
i_{small, *} :
\textit{Ab}(Z_\etale)
\longrightarrow
\textit{Ab}(X_\etale)
$$
是全忠实的，其本质像由 $X_\etale$ 上且支集包含于 $Z$ 的 Abel 层构成。
\end{enumerate}
在两种情形下，$i_{small}^{-1}$ 都是函子 $i_{small, *}$ 的左逆。
\end{proposition}

\begin{proof}
先讨论集合层的情形。对任意层 $\mathcal{G}$（定义于 $Z$ 上），态射
$i_{small}^{-1}i_{small, *}\mathcal{G} \to \mathcal{G}$
由引理 \ref{etale-cohomology-lemma-stalk-pushforward-closed-immersion}（以及定理
\ref{etale-cohomology-theorem-exactness-stalks}）是同构。这形式地蕴含 $i_{small, *}$ 全忠实；见
Sites，引理 \ref{sites-lemma-exactness-properties}。显然，
$i_{small, *}\mathcal{G}|_{U_\etale} \cong *$，其中 $U = X \setminus Z$。
反之，设 $\mathcal{F}$ 是 $X$ 上的集合层，且
$\mathcal{F}|_{U_\etale} \cong *$。考虑伴随映射
$$
\mathcal{F} \longrightarrow i_{small, *}i_{small}^{-1}\mathcal{F}
$$
合用引理 \ref{etale-cohomology-lemma-stalk-pushforward-closed-immersion} 和
\ref{etale-cohomology-lemma-stalk-pullback}，可知它是同构。这完成了 (1) 的证明。
(2) 的证明相同。
\end{proof}





\section{整普遍单射态射}
\label{etale-cohomology-section-integral-universally-injective}

\noindent
下面给出命题 \ref{etale-cohomology-proposition-closed-immersion-pushforward} 的一般版本。

\begin{proposition}
\label{etale-cohomology-proposition-integral-universally-injective-pushforward}
设 $f : X \to Y$ 是整且普遍单射的概形态射。
\begin{enumerate}
\item 函子
$$
f_{small, *} :
\Sh(X_\etale)
\longrightarrow
\Sh(Y_\etale)
$$
是全忠实的，其本质像由集合层 $\mathcal{F}$ 构成：它定义于 $Y_\etale$ 上，
且到 $Y \setminus f(X)$ 的限制同构于 $*$；
\item 函子
$$
f_{small, *} :
\textit{Ab}(X_\etale)
\longrightarrow
\textit{Ab}(Y_\etale)
$$
是全忠实的，其本质像由 $Y_\etale$ 上且支集包含于 $f(X)$ 的 Abel 层构成。
\end{enumerate}
在两种情形下，$f_{small}^{-1}$ 都是函子 $f_{small, *}$ 的左逆。
\end{proposition}

\begin{proof}
可以把 $f$ 分解为
$$
\xymatrix{
X \ar[r]^h & Z \ar[r]^i & Y
}
$$
，其中 $h$ 整、普遍单射且满，而 $i : Z \to Y$ 是闭浸入。对 $i$ 应用命题
\ref{etale-cohomology-proposition-closed-immersion-pushforward}，再对 $h$ 应用定理
\ref{etale-cohomology-theorem-topological-invariance}。
\end{proof}









\section{大网站与推前}
\label{etale-cohomology-section-big}

\noindent
本节对若干类概形态射证明关于 $f_{big, *}$ 的一些技术性结果。

\begin{lemma}
\label{etale-cohomology-lemma-monomorphism-big-push-pull}
设 $\tau \in \{Zariski, \etale, smooth, syntomic, fppf\}$。
设 $f : X \to Y$ 是概形的单态射。则典范映射
$f_{big}^{-1}f_{big, *}\mathcal{F} \to \mathcal{F}$ 是同构；这里
$\mathcal{F}$ 是定义于 $(\Sch/X)_\tau$ 上的任意层。
\end{lemma}

\begin{proof}
在这种情形下，函子 $(\Sch/X)_\tau \to (\Sch/Y)_\tau$ 连续、余连续且全忠实。
因此，结论由 Sites，引理 \ref{sites-lemma-back-and-forth} 得到。
\end{proof}

\begin{remark}
\label{etale-cohomology-remark-push-pull-shriek}
在引理 \ref{etale-cohomology-lemma-monomorphism-big-push-pull} 的情形下，典范映射
$\mathcal{F} \to f_{big}^{-1}f_{big!}\mathcal{F}$ 是同构；这里
$\mathcal{F}$ 是定义于 $(\Sch/X)_\tau$ 上的任意集合层。证明相同。
这对 Abel 群层也成立。不过要注意，Abel 群层的函子 $f_{big!}$ 定义于
Modules on Sites，第 \ref{sites-modules-section-exactness-lower-shriek} 节，
且一般不同于集合层上的 $f_{big!}$。Abel 群层的结论由
Modules on Sites，引理 \ref{sites-modules-lemma-back-and-forth} 得到。
\end{remark}

\begin{lemma}
\label{etale-cohomology-lemma-closed-immersion-cover-from-below}
设 $f : X \to Y$ 是概形的闭浸入。设 $U \to X$ 是 syntomic（分别为光滑、
\'etale）态射。则存在 syntomic（分别为光滑、\'etale）态射 $V_i \to Y$ 以及态射
$V_i \times_Y X \to U$，使得 $\{V_i \times_Y X \to U\}$ 是 $U$ 的 Zariski 覆盖。
\end{lemma}

\begin{proof}
先证明 $\tau = syntomic$ 时的引理。问题在 $U$ 上是局部的。因此，可以假设
$U$ 是映入 $Y$ 的某个仿射开子概形的仿射概形。于是，归结为证明下列情形：
$Y = \Spec(A)$、$X = \Spec(A/I)$ 且 $U = \Spec(\overline{B})$，其中
$A/I \to \overline{B}$ 是 syntomic 环同态。由 Algebra，引理
\ref{algebra-lemma-lift-syntomic}，可以找到元素
$\overline{g}_i \in \overline{B}$，使得 $\overline{B}_{\overline{g}_i} = A_i/IA_i$
对某些 syntomic 环同态 $A \to A_i$ 成立。这证明了 syntomic 情形。
光滑情形的证明相同，只是改用 Algebra，引理 \ref{algebra-lemma-lift-smooth}。
\'etale 情形使用 Algebra，引理 \ref{algebra-lemma-lift-etale}。
\end{proof}

\begin{lemma}
\label{etale-cohomology-lemma-prepare-closed-immersion-almost-cocontinuous}
设 $f : X \to Y$ 是概形的闭浸入。设 $\{U_i \to X\}$ 是 syntomic（分别为光滑、
\'etale）覆盖。则存在 syntomic（分别为光滑、\'etale）覆盖 $\{V_j \to Y\}$，
使得对每个 $j$，或者 $V_j \times_Y X = \emptyset$，或者态射
$V_j \times_Y X \to X$ 通过 $U_i$ 分解，其中 $i$ 可以依赖于所给覆盖成员。
\end{lemma}

\begin{proof}
对每个 $i$，可以选取 syntomic（分别为光滑、\'etale）态射
$g_{ij} : V_{ij} \to Y$ 以及态射 $V_{ij} \times_Y X \to U_i$（在 $X$ 上），
使得 $\{V_{ij} \times_Y X \to U_i\}$ 是 Zariski 覆盖；见引理
\ref{etale-cohomology-lemma-closed-immersion-cover-from-below}。这特别蕴含
$\bigcup_{ij} g_{ij}(V_{ij})$ 包含闭子集 $f(X)$。因此，syntomic（分别为光滑、
\'etale）映射族 $g_{ij}$ 与开浸入 $Y \setminus f(X) \to Y$ 一起构成所需的
$Y$ 的 syntomic（分别为光滑、\'etale）覆盖。
\end{proof}

\begin{lemma}
\label{etale-cohomology-lemma-closed-immersion-almost-cocontinuous}
设 $f : X \to Y$ 是概形的闭浸入。设 $\tau \in \{syntomic, smooth, \etale\}$。
函子 $V \mapsto X \times_Y V$ 给出几乎余连续函子
（见 Sites，定义 \ref{sites-definition-almost-cocontinuous}）
$(\Sch/Y)_\tau \to (\Sch/X)_\tau$；这是大 $\tau$ 网站之间的函子。
\end{lemma}

\begin{proof}
需要证明如下事实：给定态射 $V \to Y$ 和任意 syntomic（分别为光滑、\'etale）
覆盖 $\{U_i \to X \times_Y V\}$，存在光滑（分别为光滑、\'etale）覆盖
$\{V_j \to V\}$，使得对每个 $j$，或者 $X \times_Y V_j$ 为空，或者
$X \times_Y V_j \to Z \times_Y V$ 通过某个 $U_i$ 分解。这可由对闭浸入
$X \times_Y V \to V$ 应用上面的引理
\ref{etale-cohomology-lemma-prepare-closed-immersion-almost-cocontinuous} 得到。
\end{proof}

\begin{lemma}
\label{etale-cohomology-lemma-closed-immersion-pushforward-exact}
设 $f : X \to Y$ 是概形的闭浸入。设 $\tau \in \{syntomic, smooth, \etale\}$。
\begin{enumerate}
\item 推前
$f_{big, *} :
\Sh((\Sch/X)_\tau)
\to
\Sh((\Sch/Y)_\tau)$
与余等化子及推出交换。
\item 推前
$f_{big, *} :
\textit{Ab}((\Sch/X)_\tau)
\to
\textit{Ab}((\Sch/Y)_\tau)$
是正合的。
\end{enumerate}
\end{lemma}

\begin{proof}
这由 Sites，引理 \ref{sites-lemma-morphism-of-sites-almost-cocontinuous}，
Modules on Sites，引理
\ref{sites-modules-lemma-morphism-ringed-sites-almost-cocontinuous}，以及上面的引理
\ref{etale-cohomology-lemma-closed-immersion-almost-cocontinuous} 得到。
\end{proof}

\begin{remark}
\label{etale-cohomology-remark-fppf-closed-immersion-not-closed}
引理 \ref{etale-cohomology-lemma-closed-immersion-pushforward-exact} 中缺少 $\tau = fppf$ 的情形。
原因是，给定环 $A$、理想 $I$ 以及忠实平坦有限表示的环同态
$A/I \to \overline{B}$，没有理由认为可以找到{\it 任意}平坦有限表示环同态
$A \to B$，它满足 $B/IB \not = 0$，并使得 $A/I \to B/IB$ 通过
$\overline{B}$ 分解。
因此，引理 \ref{etale-cohomology-lemma-closed-immersion-almost-cocontinuous} 的证明不适用于
fppf 拓扑。事实上，
$f_{big, *} : \textit{Ab}((\Sch/X)_{fppf})
\to \textit{Ab}((\Sch/Y)_{fppf})$
在 $f$ 是闭浸入时为正合的说法很可能是假的。如果知道一个例子，请发送电子邮件至
\href{mailto:stacks.project@gmail.com}{stacks.project@gmail.com}。
\end{remark}












\section{大网站下叹号函子的正合性}
\label{etale-cohomology-section-exactness-lower-shriek}

\noindent
这里只需下列技术性结果。注意，一般而言，函子 $f_{big!}$ 与紧支集上同调毫无关系。

\begin{lemma}
\label{etale-cohomology-lemma-exactness-lower-shriek}
设 $\tau \in \{Zariski, \etale, smooth, syntomic, fppf\}$。
设 $f : X \to Y$ 是概形态射。令
$$
f_{big} :
\Sh((\Sch/X)_\tau)
\longrightarrow
\Sh((\Sch/Y)_\tau)
$$
为 Topologies，引理 \ref{topologies-lemma-morphism-big}、
\ref{topologies-lemma-morphism-big-etale},
\ref{topologies-lemma-morphism-big-smooth},
\ref{topologies-lemma-morphism-big-syntomic} 或
\ref{topologies-lemma-morphism-big-fppf} 中相应的拓扑斯态射。
\begin{enumerate}
\item 函子
$f_{big}^{-1} : \textit{Ab}((\Sch/Y)_\tau) \to \textit{Ab}((\Sch/X)_\tau)$
有左伴随
$$
f_{big!} : \textit{Ab}((\Sch/X)_\tau) \to \textit{Ab}((\Sch/Y)_\tau)
$$
，并且后者正合。
\item 函子
$f_{big}^* :
\textit{Mod}((\Sch/Y)_\tau, \mathcal{O})
\to
\textit{Mod}((\Sch/X)_\tau, \mathcal{O})$
有左伴随
$$
f_{big!} :
\textit{Mod}((\Sch/X)_\tau, \mathcal{O})
\to
\textit{Mod}((\Sch/Y)_\tau, \mathcal{O})
$$
，并且后者正合。
\end{enumerate}
此外，这两个函子 $f_{big!}$ 在底层 Abel 群层上相同。
\end{lemma}

\begin{proof}
回忆，$f_{big}$ 是由连续且余连续的函子
$u : (\Sch/X)_\tau \to (\Sch/Y)_\tau$、$U/X \mapsto U/Y$ 所联系的拓扑斯态射。
此外，有 $f_{big}^{-1}\mathcal{O} = \mathcal{O}$。因此，$f_{big!}$ 的存在性分别由
Modules on Sites，引理 \ref{sites-modules-lemma-g-shriek-adjoint} 和
Modules on Sites，引理 \ref{sites-modules-lemma-lower-shriek-modules} 得到。
注意，若 $U$ 是 $(\Sch/X)_\tau$ 的对象，则函子 $u$ 诱导范畴等价
$$
u' :
(\Sch/X)_\tau/U
\longrightarrow
(\Sch/Y)_\tau/U
$$
，因为箭头两边都等于 $(\Sch/U)_\tau$。因此，$f_{big!}$ 在底层 Abel 层上相同由
Modules on Sites，注 \ref{sites-modules-remark-when-shriek-equal} 中的讨论得到。
$f_{big!}$ 的正合性由 Modules on Sites，引理
\ref{sites-modules-lemma-exactness-lower-shriek} 得到，因为上面的函子 $u$
与纤维积和等化子交换。
\end{proof}

\noindent
下面证明一个技术性引理，稍后比较代数栈所联系的不同网站上的模层时将会用到它。

\begin{lemma}
\label{etale-cohomology-lemma-compare-structure-sheaves}
设 $X$ 是概形。设 $\tau \in \{Zariski, \etale, smooth, syntomic, fppf\}$。
设 $\mathcal{C}_1 \subset \mathcal{C}_2 \subset (\Sch/X)_\tau$ 是满足下列性质的
全子范畴：
\begin{enumerate}
\item 对对象 $U/X$（它属于 $\mathcal{C}_t$），
\begin{enumerate}
\item 若 $\{U_i \to U\}$ 是 $(\Sch/X)_\tau$ 的覆盖，则 $U_i/X$ 是
$\mathcal{C}_t$ 的对象；
\item $U \times \mathbf{A}^1/X$ 是 $\mathcal{C}_t$ 的对象。
\end{enumerate}
\item $X/X$ 是 $\mathcal{C}_t$ 的对象。
\end{enumerate}
赋予 $\mathcal{C}_t$ 网站结构；其覆盖恰为那些覆盖 $\{U_i \to U\}$，它们是
$(\Sch/X)_\tau$ 的覆盖并满足 $U \in \Ob(\mathcal{C}_t)$。则
\begin{enumerate}
\item[(a)] 函子 $\mathcal{C}_1 \to \mathcal{C}_2$ 全忠实、连续且余连续。
\end{enumerate}
以 $g : \Sh(\mathcal{C}_1) \to \Sh(\mathcal{C}_2)$ 表示相应的拓扑斯态射。
以 $\mathcal{O}_t$ 表示 $\mathcal{O}$ 到 $\mathcal{C}_t$ 的限制。以 $g_!$ 表示
Modules on Sites，定义 \ref{sites-modules-definition-g-shriek} 中的函子。
\begin{enumerate}
\item[(b)] 典范映射 $g_!\mathcal{O}_1 \to \mathcal{O}_2$ 是同构。
\end{enumerate}
\end{lemma}

\begin{proof}
断言 (a) 由定义立即可得。在本证明中，所有概形都是 $X$ 上的概形，所有概形态射
都是 $X$ 上的概形态射。注意，$g^{-1}$ 由限制给出，所以对对象 $U$
（它属于 $\mathcal{C}_1$），有 $\mathcal{O}_1(U) = \mathcal{O}_2(U) = \mathcal{O}(U)$。
回忆，$g_!\mathcal{O}_1$ 是预层 $g_{p!}\mathcal{O}_1$ 所联系的层；该预层把
对象 $V$（在 $\mathcal{C}_2$ 中）映到群
$$
\colim_{V \to U} \mathcal{O}(U)
$$
，其中 $U$ 遍历 $\mathcal{C}_1$ 的对象，余极限取在 Abel 群范畴中。下面将反复使用
如下事实：若
$$
V \to U \to U'
$$
是满足 $U, U' \in \Ob(\mathcal{C}_1)$ 的态射，且
$f' \in \mathcal{O}(U')$ 限制为 $f \in \mathcal{O}(U)$，则
$(V \to U, f)$ 与 $(V \to U', f')$ 定义余极限的同一个元素。此外，
$g_!\mathcal{O}_1 \to \mathcal{O}_2$ 只是把元素 $(V \to U, f)$ 映到 $f$ 在
$V$ 上的拉回。

\medskip\noindent
满性。设 $V$ 是概形，并设 $h \in \mathcal{O}(V)$。于是得到态射
$V \to X \times \mathbf{A}^1$；它由 $h$ 和结构态射 $V \to X$ 诱导。写成
$\mathbf{A}^1 = \Spec(\mathbf{Z}[x])$，可见元素
$x \in \mathcal{O}(X \times \mathbf{A}^1)$ 拉回到 $h$。由假设 (1)(b) 和 (2)，
$X \times \mathbf{A}^1$ 是 $\mathcal{C}_1$ 的对象，于是得到所需的满性。

\medskip\noindent
单性。设 $V$ 是概形。设
$s = \sum_{i = 1, \ldots, n} (V \to U_i, f_i)$ 是上面所示余极限的一个元素。
对任意 $i$，可以使用态射 $f_i : U_i \to X \times \mathbf{A}^1$，看出
$(V \to U_i, f_i)$ 与 $(f_i : V \to X \times \mathbf{A}^1, x)$ 定义余极限的
同一个元素。然后考虑
$$
f_1 \times \ldots \times f_n : V \to X \times \mathbf{A}^n
$$
，可见 $s$ 在余极限中等价于
$$
\sum\nolimits_{i = 1, \ldots, n}
(f_1 \times \ldots \times f_n : V \to X \times \mathbf{A}^n, x_i) =
(f_1 \times \ldots \times f_n : V \to X \times \mathbf{A}^n,
x_1 + \ldots + x_n)
$$
现在，若 $x_1 + \ldots + x_n$ 在 $V$ 上限制为零，则可见
$f_1 \times \ldots \times f_n$ 通过
$X \times \mathbf{A}^{n - 1} = V(x_1 + \ldots + x_n)$ 分解。因此，$s$
在余极限中等价于零。
\end{proof}











%9.29.09
\section{\'Etale 上同调}
\label{etale-cohomology-section-etale-cohomology}

\noindent
以下各节将证明一些关于 \'etale 上同调的基本结果。下面举出一个对拓扑空间上同调
成立、并且对 \'etale 上同调也成立的例子。

\begin{lemma}[\'Etale 上同调的 Mayer--Vietoris 序列]
\label{etale-cohomology-lemma-mayer-vietoris}
设 $X$ 是概形。假设 $X = U \cup V$ 是两个开子概形之并。对任意 Abel 层
$\mathcal{F}$（定义于 $X_\etale$ 上），存在上同调长正合序列
$$
\begin{matrix}
0 \to
H^0_\etale(X, \mathcal{F}) \to
H^0_\etale(U, \mathcal{F}) \oplus H^0_\etale(V, \mathcal{F}) \to
H^0_\etale(U \cap V, \mathcal{F}) \phantom{\to \ldots} \\
\phantom{0} \to H^1_\etale(X, \mathcal{F}) \to
H^1_\etale(U, \mathcal{F}) \oplus H^1_\etale(V, \mathcal{F}) \to
H^1_\etale(U \cap V, \mathcal{F}) \to \ldots
\end{matrix}
$$
这个长正合序列关于 $\mathcal{F}$ 是函子性的。
\end{lemma}

\begin{proof}
注意，若 $\mathcal{I}$ 是内射 Abel 层，则
$$
0 \to \mathcal{I}(X) \to \mathcal{I}(U) \oplus \mathcal{I}(V) \to
\mathcal{I}(U \cap V) \to 0
$$
正合。由于 $\mathcal{I}$ 是层，首项和中间项处的正合性成立。右端的正合性成立，
因为由 Cohomology on Sites，引理
\ref{sites-cohomology-lemma-restriction-along-monomorphism-surjective}，
$\mathcal{I}(U) \to \mathcal{I}(U \cap V)$ 是满射。另一种证明方法是指出映射
$\mathcal{I}(U) \oplus \mathcal{I}(V) \to
\mathcal{I}(U \cap V)$ 的余核就是
$\mathcal{I}$ 关于覆盖 $X = U \cup V$ 的第一 {\v C}ech 上同调群；它由引理
\ref{etale-cohomology-lemma-hom-injective} 和 \ref{etale-cohomology-lemma-forget-injectives} 消失。因此，若
$\mathcal{F} \to \mathcal{I}^\bullet$ 是内射解消，则
$$
0 \to \mathcal{I}^\bullet(X) \to
\mathcal{I}^\bullet(U) \oplus \mathcal{I}^\bullet(V) \to
\mathcal{I}^\bullet(U \cap V) \to 0
$$
是复形的短正合序列，而相联系的上同调长正合序列正是引理陈述中的序列。
\end{proof}

\begin{lemma}[相对 Mayer--Vietoris 序列]
\label{etale-cohomology-lemma-relative-mayer-vietoris}
设 $f : X \to Y$ 是概形态射。假设 $X = U \cup V$ 是两个开子概形之并。记
$a = f|_U : U \to Y$、$b = f|_V : V \to Y$ 以及
$c = f|_{U \cap V} : U \cap V \to Y$。对每个 Abel 层 $\mathcal{F}$
（定义于 $X_\etale$ 上），存在长正合序列
$$
0 \to
f_*\mathcal{F} \to
a_*(\mathcal{F}|_U) \oplus b_*(\mathcal{F}|_V) \to
c_*(\mathcal{F}|_{U \cap V}) \to
R^1f_*\mathcal{F} \to \ldots
$$
（定义于 $Y_\etale$ 上）。这个长正合序列关于 $\mathcal{F}$ 是函子性的。
\end{lemma}

\begin{proof}
设 $\mathcal{F} \to \mathcal{I}^\bullet$ 是 $\mathcal{F}$ 在 $X_\etale$ 上的
内射解消。我们断言得到复形的短正合序列
$$
0 \to
f_*\mathcal{I}^\bullet \to
a_*\mathcal{I}^\bullet|_U \oplus b_*\mathcal{I}^\bullet|_V \to
c_*\mathcal{I}^\bullet|_{U \cap V} \to
0.
$$
事实上，对任意 $W$（属于 $Y_\etale$）以及任意 $n \geq 0$，在 $W$ 上相应的
截面群序列
$$
0 \to
\mathcal{I}^n(W \times_Y X) \to
\mathcal{I}^n(W \times_Y U)
\oplus \mathcal{I}^n(W \times_Y V) \to
\mathcal{I}^n(W \times_Y (U \cap V)) \to
0
$$
已在引理 \ref{etale-cohomology-lemma-mayer-vietoris} 的证明中证明为短正合。取上同调层，并使用
$\mathcal{I}^\bullet|_U$ 是 $\mathcal{F}|_U$ 的内射解消这一事实，即得引理；
$\mathcal{I}^\bullet|_V$、$\mathcal{I}^\bullet|_{U \cap V}$ 的情形相同。
\end{proof}







\section{余极限}
\label{etale-cohomology-section-colimit}

\noindent
回忆，若 $(\mathcal{F}_i, \varphi_{ii'})$ 是网站 $\mathcal{C}$ 上的层图，则它在层范畴中
的余极限是预层 $U \mapsto \colim_i \mathcal{F}_i(U)$ 的层化；见 Sites，引理
\ref{sites-lemma-colimit-sheaves}。若系统是定向的，$U$ 是 $\mathcal{C}$ 的拟紧对象，
并且它有由拟紧对象组成的余终覆盖系，则
$\mathcal{F}(U) = \colim \mathcal{F}_i(U)$；见 Sites，引理
\ref{sites-lemma-directed-colimits-sections}。关于 Abel 层余极限之高阶上同调群的结果，
见 Cohomology on Sites，引理
\ref{sites-cohomology-lemma-colim-works-over-collection}。

\medskip\noindent
在 Cohomology on Sites，引理 \ref{sites-cohomology-lemma-colimit} 中，我们把这个结果
推广到网站逆系统上的层系统。下面给出概形逆系统上的 \'etale 层系统这一相应概念。

\begin{definition}
\label{etale-cohomology-definition-inverse-system-sheaves}
设 $I$ 是预序集。设 $(X_i, f_{i'i})$ 是以 $I$ 为指标的概形逆系统。
一个{\it 层系统 $(\mathcal{F}_i, \varphi_{i'i})$}（定义于 $(X_i, f_{i'i})$ 上）
由下列数据给出：
\begin{enumerate}
\item 一个层 $\mathcal{F}_i$（定义于 $(X_i)_\etale$ 上），对所有 $i \in I$；
\item 对 $i' \geq i$，层的映射
$\varphi_{i'i} : f_{i'i}^{-1}\mathcal{F}_i \to \mathcal{F}_{i'}$
（定义于 $(X_{i'})_\etale$ 上）。
\end{enumerate}
这些数据满足等式
$\varphi_{i''i} = \varphi_{i''i'} \circ f_{i'' i'}^{-1}\varphi_{i'i}$，
只要 $i'' \geq i' \geq i$。
\end{definition}

\noindent
在定义 \ref{etale-cohomology-definition-inverse-system-sheaves} 的情形下，假设 $I$ 是定向集，并且转移态射
$f_{i'i}$ 仿射。设 $X = \lim X_i$ 是概形范畴中的极限；见 Limits，第
\ref{limits-section-limits} 节。以 $f_i : X \to X_i$ 表示投影态射，并考虑映射
$$
f_i^{-1}\mathcal{F}_i = f_{i'}^{-1}f_{i'i}^{-1}\mathcal{F}_i
\xrightarrow{f_{i'}^{-1}\varphi_{i'i}}
f_{i'}^{-1}\mathcal{F}_{i'}
$$
这使 $f_i^{-1}\mathcal{F}_i$ 成为 $X_\etale$ 上以 $I$ 为指标的层系统
（验证这一点是个很好的练习）。我们经常希望知道是否存在同构
$$
H^q_\etale(X, \colim f_i^{-1}\mathcal{F}_i) =
\colim H^q_\etale(X_i, \mathcal{F}_i)
$$
后面会证明：若 $X_i$ 对所有 $i$ 都拟紧且拟分离，则此式成立；见定理
\ref{etale-cohomology-theorem-colimit}。

\begin{lemma}
\label{etale-cohomology-lemma-colimit-affine-sites}
设 $I$ 是定向集。设 $(X_i, f_{i'i})$ 是以 $I$ 为指标、具有仿射转移态射的概形逆系统。
设 $X = \lim_{i \in I} X_i$。采用 Topologies，引理
\ref{topologies-lemma-alternative} 中的记号，有
$$
X_{affine, \etale} = \colim (X_i)_{affine, \etale}
$$
；这里的网站等式取 Sites，引理 \ref{sites-lemma-colimit-sites} 的意义。
\end{lemma}

\begin{proof}
先在 $X$ 与 $X_i$ 都拟紧且拟分离（对所有 $i$）时证明此式；所有我们关心的情形
都是如此。在这种情形下，$X_{affine, \etale}$ 的任意对象，以及相应地
$(X_i)_{affine, \etale}$ 的任意对象，都在 $X$ 上有限表示。此外，$X$ 上有限表示概形
的范畴是 $X_i$ 上有限表示概形范畴的余极限；见 Limits，引理
\ref{limits-lemma-descend-finite-presentation}。由 Limits，引理
\ref{limits-lemma-limit-affine} 和 \ref{limits-lemma-descend-etale}，对 $X$ 上
\'etale 的仿射对象所成子范畴，同一结论也成立。最后，若 $\{U^j \to U\}$ 是
$X_{affine, \etale}$ 的覆盖，并且 $U_i^j \to U_i$ 是 $X_i$ 上 \'etale 的仿射概形
间态射，其到 $X$ 的基变换为 $U^j \to U$，则
$\{U^j_i \to U_i\}$ 到某个 $X_{i'}$ 的基变换是覆盖（当 $i'$ 足够大）；见 Limits，引理
\ref{limits-lemma-descend-surjective}。

\medskip\noindent
在一般情形下，设 $U$ 是 $X_{affine, \etale}$ 的对象。则 $U \to X$ 是 \'etale
且分离的（因为 $U$ 分离），但一般不拟紧。不过，$U \to X$ 局部有限表示。因此，
由 Limits，引理 \ref{limits-lemma-descend-finite-presentation-variant}，存在某个 $i$、
拟紧拟分离概形 $U_i$ 以及局部有限表示态射 $U_i \to X_i$，其到 $X$ 的基变换
为 $U \to X$。于是 $U = \lim_{i' \geq i} U_{i'}$，其中
$U_{i'} = U_i \times_{X_i} X_{i'}$。增大 $i$ 后，可以假设 $U_i$ 仿射；见
Limits，引理 \ref{limits-lemma-limit-affine}。为验证 $U_i \to X_i$ 对充分大的 $i$
是 \'etale 的，选取有限仿射开覆盖
$U_i = U_{i, 1} \cup \ldots \cup U_{i, m}$，使得
$U_{i, j} \to U_i \to X_i$ 映入仿射开集 $W_{i, j} \subset X_i$。应用 Limits，引理
\ref{limits-lemma-descend-etale}，可知 $U_{i, j} \to W_{i, j}$ 是 \'etale 的
（必要时再次增大 $i$）。由此可见函子
$\colim (X_i)_{affine, \etale} \to X_{affine, \etale}$ 本质满。全忠实性直接由
已经使用过的 Limits，引理
\ref{limits-lemma-descend-finite-presentation-variant} 得到。关于覆盖的断言与证明第一段
完全同样地证明。
\end{proof}

\noindent
利用上面的结果，得到下列关于余极限和上同调的一般结论。

\begin{theorem}
\label{etale-cohomology-theorem-colimit}
设 $X = \lim_{i \in I} X_i$ 是具有仿射转移态射
$f_{i'i} : X_{i'} \to X_i$ 的概形定向系统的极限。假设 $X_i$ 对所有 $i \in I$
都拟紧且拟分离。设 $(\mathcal{F}_i, \varphi_{i'i})$ 是 $(X_i, f_{i'i})$ 上的
Abel 层系统。以 $f_i : X \to X_i$ 表示投影，并置
$\mathcal{F} = \colim f_i^{-1}\mathcal{F}_i$。则
$$
\colim_{i\in I} H_\etale^p(X_i, \mathcal{F}_i) = H_\etale^p(X, \mathcal{F}).
$$
对所有 $p \geq 0$ 成立。
\end{theorem}

\begin{proof}
由 Topologies，引理 \ref{topologies-lemma-alternative}，可以计算
$\mathcal{F}$ 在 $X_{affine, \etale}$ 上的上同调。因此，结论由引理
\ref{etale-cohomology-lemma-colimit-affine-sites} 和 Cohomology on Sites，引理
\ref{sites-cohomology-lemma-colimit} 合并得到。
\end{proof}

\noindent
下面两个结果是上述定理的特殊情形。

\begin{lemma}
\label{etale-cohomology-lemma-colimit}
设 $X$ 是拟紧拟分离概形。设 $I$ 是定向集。设
$(\mathcal{F}_i, \varphi_{ij})$ 是 $X_\etale$ 上以 $I$ 为指标的 Abel 层系统。则
$$
\colim_{i\in I} H_\etale^p(X, \mathcal{F}_i) = H_\etale^p(X,
\colim_{i\in I} \mathcal{F}_i).
$$
\end{lemma}

\begin{proof}
这是定理 \ref{etale-cohomology-theorem-colimit} 的特殊情形。再概述一个直接证明。同时对所有 $X$
证明结论，并对 $p$ 作归纳。
\begin{enumerate}
\item
对任意拟紧拟分离概形 $X$ 以及任意 \'etale 覆盖 $\mathcal{U}$（它覆盖 $X$），证明存在加细
$\mathcal{V} = \{V_j \to X\}_{j\in J}$，其中 $J$ 有限，每个 $V_j$ 拟紧且拟分离，
并且所有 $V_{j_0} \times_X \ldots \times_X V_{j_p}$ 也拟紧且拟分离。
\item
利用上一步和层范畴中余极限的定义，证明定理对 $p = 0$ 以及所有 $X$ 成立。
\item
利用上同调的局部性（引理 \ref{etale-cohomology-lemma-locality-cohomology}）、从 {\v C}ech 到上同调的
谱序列（定理 \ref{etale-cohomology-theorem-cech-ss}），以及归纳假设适用于上述情形中的所有
$V_{j_0} \times_X \ldots \times_X V_{j_p}$ 这一事实，证明归纳步骤 $p \to p + 1$。
\end{enumerate}
\end{proof}

\begin{lemma}
\label{etale-cohomology-lemma-directed-colimit-cohomology}
设 $A$ 是环，$(I, \leq)$ 是定向集，$(B_i, \varphi_{ij})$ 是 $A$-代数系统。
置 $B = \colim_{i\in I} B_i$。设 $X \to \Spec(A)$ 是拟紧拟分离的概形态射。
设 $\mathcal{F}$ 是 $X_\etale$ 上的 Abel 层。记
$Y_i = X \times_{\Spec(A)} \Spec(B_i)$、
$Y = X \times_{\Spec(A)} \Spec(B)$、
$\mathcal{G}_i = (Y_i \to X)^{-1}\mathcal{F}$ 以及
$\mathcal{G} = (Y \to X)^{-1}\mathcal{F}$。则
$$
H_\etale^p(Y, \mathcal{G}) =
\colim_{i\in I} H_\etale^p (Y_i, \mathcal{G}_i).
$$
\end{lemma}

\begin{proof}
这是定理 \ref{etale-cohomology-theorem-colimit} 的特殊情形。再如下概述一个直接证明。
\begin{enumerate}
\item 给定 \'etale 态射 $V \to Y$，其中 $V$ 拟紧且拟分离，则存在 $i\in I$ 和
$V_i \to Y_i$，使得 $V = V_i \times_{Y_i} Y$。若所考虑的概形全都仿射，
这对应于下列代数断言：若 $B = \colim B_i$ 且 $B \to C$ 是 \'etale 的，
则存在 $i \in I$ 以及 \'etale 同态 $B_i \to C_i$，使得
$C \cong B \otimes_{B_i} C_i$。这在 Algebra，引理 \ref{algebra-lemma-etale} 中证明。
\item 在 (1) 的情形下，证明
$\mathcal{G}(V) = \colim_{i' \geq i} \mathcal{G}_{i'}(V_{i'})$
，其中 $V_{i'}$ 是 $V_i$ 到 $Y_{i'}$ 的基变换。
\item 由 (1) 可见，对每个 \'etale 覆盖
$\mathcal{V} = \{V_j \to Y\}_{j\in J}$，其中 $J$ 有限且诸 $V_j$ 拟紧拟分离，
存在 $i \in I$ 以及 \'etale 覆盖
$\mathcal{V}_i = \{V_{ij} \to Y_i\}_{j \in J}$，使得
$\mathcal{V} \cong \mathcal{V}_i \times_{Y_i} Y$。
\item 证明 (2) 和 (3) 蕴含
$$
\check H^*(\mathcal{V}, \mathcal{G})=
\colim_{i\in I} \check H^*(\mathcal{V}_i, \mathcal{G}_i).
$$
\item 巧妙地使用从 {\v C}ech 到上同调的谱序列（定理 \ref{etale-cohomology-theorem-cech-ss}）。
\end{enumerate}
\end{proof}

\begin{lemma}
\label{etale-cohomology-lemma-higher-direct-images}
设 $f: X\to Y$ 是概形态射，并且 $\mathcal{F}\in
\textit{Ab}(X_\etale)$。则 $R^pf_*\mathcal{F}$ 是下列预层所联系的层：
$$
(V \to Y) \longmapsto H_\etale^p(X \times_Y V, \mathcal{F}|_{X \times_Y V}).
$$
更一般地，对 $K \in D(X_\etale)$，$R^pf_*K$ 是下列预层所联系的层：
$$
(V \to Y) \longmapsto H_\etale^p(X \times_Y V, K|_{X \times_Y V}).
$$
\end{lemma}

\begin{proof}
这个引理对拓扑空间也成立，并且该情形的证明相同。层的情形见 Cohomology on Sites，
引理 \ref{sites-cohomology-lemma-higher-direct-images}；Abel 层复形的情形见
Cohomology on Sites，引理
\ref{sites-cohomology-lemma-sheafification-cohomology}。
\end{proof}

\begin{lemma}
\label{etale-cohomology-lemma-relative-colimit}
设 $S$ 是概形。设 $X = \lim_{i \in I} X_i$ 是 $S$ 上概形定向系统的极限，
其转移态射 $f_{i'i} : X_{i'} \to X_i$ 仿射。假设结构态射
$g_i : X_i \to S$ 与 $g : X \to S$ 拟紧且拟分离。设
$(\mathcal{F}_i, \varphi_{i'i})$ 是 $(X_i, f_{i'i})$ 上的 Abel 层系统。
以 $f_i : X \to X_i$ 表示投影，并置
$\mathcal{F} = \colim f_i^{-1}\mathcal{F}_i$。则
$$
\colim_{i\in I} R^p g_{i, *} \mathcal{F}_i = R^p g_* \mathcal{F}
$$
对所有 $p \geq 0$ 成立。
\end{lemma}

\begin{proof}
回忆（引理 \ref{etale-cohomology-lemma-higher-direct-images}），$R^p g_{i, *} \mathcal{F}_i$ 是预层
$U \mapsto H^p_\etale(U \times_S X_i, \mathcal{F}_i)$ 所联系的层；
$R^pg_*\mathcal{F}$ 也类似。此外，层系统的余极限是预层层次上余极限的层化。
注意，$S_\etale$ 的每个对象都有由拟紧拟分离对象（例如仿射概形）组成的覆盖。
此外，若 $U$ 是拟紧拟分离对象，则有
$$
\colim H^p_\etale(U \times_S X_i, \mathcal{F}_i) =
H^p_\etale(U \times_S X, \mathcal{F})
$$
；这由定理 \ref{etale-cohomology-theorem-colimit} 得到。因此引理成立。
\end{proof}

\begin{lemma}
\label{etale-cohomology-lemma-relative-colimit-general}
设 $I$ 是定向集。设 $g_i : X_i \to S_i$ 是以 $I$ 为指标的概形态射逆系统。
假设 $g_i$ 拟紧且拟分离，并且对 $i' \geq i$，转移态射
$f_{i'i} : X_{i'} \to X_i$ 与 $h_{i'i} : S_{i'} \to S_i$ 仿射。设
$g : X \to S$ 是态射 $g_i$ 的极限；见 Limits，第 \ref{limits-section-limits} 节。
以 $f_i : X \to X_i$ 和 $h_i : S \to S_i$ 表示投影。设
$(\mathcal{F}_i, \varphi_{i'i})$ 是 $(X_i, f_{i'i})$ 上的层系统。置
$\mathcal{F} = \colim f_i^{-1}\mathcal{F}_i$。则
$$
R^p g_* \mathcal{F} =
\colim_{i \in I} h_i^{-1}R^p g_{i, *} \mathcal{F}_i
$$
对所有 $p \geq 0$ 成立。
\end{lemma}

\begin{proof}
引理中的映射如何构造？对 $i' \geq i$，有交换图
$$
\xymatrix{
X \ar[r]_{f_{i'}} \ar[d]_g &
X_{i'} \ar[r]_{f_{i'i}} \ar[d]_{g_{i'}} &
X_i \ar[d]^{g_i} \\
S \ar[r]^{h_{i'}} &
S_{i'} \ar[r]^{h_{i'i}} &
S_i
}
$$
把基变换映射
$h_{i'i}^{-1}Rg_{i, *}\mathcal{F}_i \to Rg_{i', *}f_{i'i}^{-1}\mathcal{F}_i$
（Cohomology on Sites，引理
\ref{sites-cohomology-lemma-base-change-map-flat-case} 或
注 \ref{sites-cohomology-remark-base-change}）
与映射 $Rg_{i', *}\varphi_{i'i}$ 合并，得到
$\psi_{i'i} : h_{i' i}^{-1} R^p g_{i, *} \mathcal{F}_i \to
R^pg_{i', *} \mathcal{F}_{i'}$。类似地，利用图中的左方块，得到映射
$\psi_i : h_i^{-1}R^pg_{i, *}\mathcal{F}_i \to R^pg_*\mathcal{F}$.
映射 $h_{i'}^{-1}\psi_{i'i}$ 与 $\psi_i$ 就是引理陈述中使用的映射。为使其有意义，
必须验证
$\psi_{i''i} = \psi_{i''i'} \circ h_{i''i'}^{-1}\psi_{i'i}$ 且
$\psi_{i'} \circ h_{i'}^{-1}\psi_{i'i} = \psi_i$；这由 Cohomology on Sites，注
\ref{sites-cohomology-remark-compose-base-change-horizontal} 得到。

\medskip\noindent
证明该等式。第一种证明使用维数移动\footnote{也可以用这种方法构造引理中的映射。}。
对任意 $U$（它在 $X$ 上仿射且 \'etale），由定理 \ref{etale-cohomology-theorem-colimit} 有
$$
g_*\mathcal{F}(U) =
H^0(U \times_S X, \mathcal{F}) =
\colim H^0(U_i \times_{S_i} X_i, \mathcal{F}_i) =
\colim g_{i, *}\mathcal{F}_i(U_i)
$$
，其中余极限取遍充分大的 $i$，从而存在某个 $i$ 和某个 $U_i$
（它在 $S_i$ 上仿射且 \'etale），其基变换为 $U$（在 $S$ 上；见引理
\ref{etale-cohomology-lemma-colimit-affine-sites}）。
由 Sites，引理 \ref{sites-lemma-colimit}，右边等于
$(\colim h_i^{-1}g_{i, *}\mathcal{F}_i)(U)$。这证明了 $p = 0$ 时的引理。
若 $(\mathcal{G}_i, \varphi_{i'i})$ 是满足
$\mathcal{G} = \colim f_i^{-1}\mathcal{G}_i$ 的系统，使得 $\mathcal{G}_i$ 是
$X_i$ 上的内射 Abel 层（对所有 $i$），则对任意 $U$（它在 $X$ 上仿射且 \'etale），由定理
\ref{etale-cohomology-theorem-colimit} 有
$$
H^p(U \times_S X, \mathcal{G}) =
\colim H^p(U_i \times_{S_i} X_i, \mathcal{G}_i) = 0
$$
，对 $p > 0$ 成立（余极限与前相同）。因此 $R^pg_*\mathcal{G} = 0$，
并且对这样的系统得到 $p > 0$ 时的结论。一般而言，可以选取系统的短正合序列
$$
0 \to (\mathcal{F}_i, \varphi_{i'i}) \to
(\mathcal{G}_i, \varphi_{i'i}) \to
(\mathcal{Q}_i, \varphi_{i'i}) \to 0
$$
，其中 $(\mathcal{G}_i, \varphi_{i'i})$ 如上；见 Cohomology on Sites，引理
\ref{sites-cohomology-lemma-colim-sites-injective}。由归纳，引理对 $p - 1$ 成立；
由上文，对 $p$ 和 $(\mathcal{G}_i, \varphi_{i'i})$ 有消失。因此，上同调长正合序列
给出 $p$ 和 $(\mathcal{F}_i, \varphi_{i'i})$ 的结论。

\medskip\noindent
第二种证明。回忆，$S_{affine, \etale} = \colim (S_i)_{affine, \etale}$；见引理
\ref{etale-cohomology-lemma-colimit-affine-sites}。因此，若 $U$ 是 $S_{affine, \etale}$ 的对象，
则可写成 $U = U_i \times_{S_i} S$，其中有某个 $i$ 以及某个 $U_i$
（它属于 $(S_i)_{affine, \etale}$），并且
$$
(\colim_{i \in I} h_i^{-1}R^p g_{i, *} \mathcal{F}_i)(U) =
\colim_{i' \geq i} (R^p g_{i', *}\mathcal{F}_{i'})(U_i \times_{S_i} S_{i'})
$$
；这由 Sites，引理 \ref{sites-lemma-colimit} 以及上面所述系统中转移映射的构造得到。
由于 $R^pg_{i', *}\mathcal{F}_{i'}$ 是预层
$U_{i'} \mapsto H^p(U_{i'} \times_{S_{i'}} X_{i'}, \mathcal{F}_{i'})$ 所联系的层，
而 $R^pg_*\mathcal{F}$ 是预层 $U \mapsto H^p(U \times_S X, \mathcal{F})$
所联系的层（引理 \ref{etale-cohomology-lemma-higher-direct-images}），得到典范交换图
$$
\xymatrix{
\colim_{i' \geq i}
H^p(U_i \times_{S_i} X_{i'}, \mathcal{F}_{i'}) \ar[r] \ar[d] &
\colim_{i' \geq i}
(R^p g_{i', *}\mathcal{F}_{i'})(U_i \times_{S_i} S_{i'}) \ar[d] \\
H^p(U \times_S X, \mathcal{F}) \ar[r] &
R^pg_*\mathcal{F}(U)
}
$$
注意，由定理 \ref{etale-cohomology-theorem-colimit}，左侧竖直箭头是同构。我们要证明右侧竖直箭头
是同构。不过，已经知道该箭头的源和靶都是 $S_{affine, \etale}$ 上的层。因此，只需
证明：(1) 靶中的元素局部来自源中的元素；(2) 源中映到靶中零的元素局部消失。
(1) 立即由上文和下方水平箭头来自预层映射这一事实得到；该预层映射在层化后成为同构。
对于 (2)，设
$\xi \in  \colim_{i' \geq i}
(R^p g_{i', *}\mathcal{F}_{i'})(U_i \times_{S_i} S_{i'})$ 属于核。选取
$i' \geq i$ 以及元素
$\xi_{i'} \in (R^p g_{i', *}\mathcal{F}_{i'})(U_i \times_{S_i} S_{i'})$，
它表示 $\xi$。
选取标准 \'etale 覆盖
$\{U_{i', k} \to U_i \times_{S_i} S_{i'}\}_{k = 1, \ldots, m}$，
使得 $\xi_{i'}|_{U_{i', k}}$ 来自
$\xi_{i', k} \in H^p(U_{i', k} \times_{S_{i'}} X_{i'}, \mathcal{F}_{i'})$。
由于只需证明 $\xi$ 局部消失，可以替换 $U$，使用 \'etale 覆盖
$\{U_{i', k} \times_{S_{i'}} S \to U = U_i \times_{S_i} S\}$ 的成员。
替换后，可见 $\xi$ 是某个元素 $\xi'$ 的像；该元素属于上图中的群
$\colim_{i' \geq i} H^p(U_i \times_{S_i} X_{i'}, \mathcal{F}_{i'})$。
由于 $\xi'$ 在 $R^pg_*\mathcal{F}(U)$ 中映到零，可以再作一次替换，
并假设 $\xi'$ 在 $H^p(U \times_S X, \mathcal{F})$ 中映到零。不过，由于左侧竖直箭头
是同构，于是得到 $\xi' = 0$，从而如愿有 $\xi = 0$。
\end{proof}

\begin{lemma}
\label{etale-cohomology-lemma-linus-hamann}
设 $X = \lim_{i \in I} X_i$ 是具有仿射转移态射 $f_{i'i}$ 和投影态射
$f_i : X \to X_i$ 的概形定向极限。设 $\mathcal{F}$ 是 $X_\etale$ 上的层。则
\begin{enumerate}
\item 存在典范映射
$\varphi_{i'i} : f_{i'i}^{-1}f_{i, *}\mathcal{F} \to f_{i', *}\mathcal{F}$
，使得 $(f_{i, *}\mathcal{F}, \varphi_{i'i})$ 是定义
\ref{etale-cohomology-definition-inverse-system-sheaves} 意义下 $(X_i, f_{i'i})$ 上的层系统；
\item $\mathcal{F} = \colim f_i^{-1}f_{i, *}\mathcal{F}$.
\end{enumerate}
\end{lemma}

\begin{proof}
通过 Topologies，引理 \ref{topologies-lemma-alternative} 和引理
\ref{etale-cohomology-lemma-colimit-affine-sites}，这是 Sites，引理
\ref{sites-lemma-colimit-push-pull} 的特殊情形。
\end{proof}

\begin{lemma}
\label{etale-cohomology-lemma-compute-strangely}
设 $I$ 是定向集。设 $g_i : X_i \to S_i$ 是以 $I$ 为指标的概形态射逆系统。
假设 $g_i$ 拟紧且拟分离，并且对 $i' \geq i$，转移态射
$X_{i'} \to X_i$ 和 $S_{i'} \to S_i$ 仿射。设 $g : X \to S$ 是态射 $g_i$ 的极限；
见 Limits，第 \ref{limits-section-limits} 节。以 $f_i : X \to X_i$ 和
$h_i : S \to S_i$ 表示投影。设 $\mathcal{F}$ 是 $X$ 上的 Abel 层。则有
$$
R^pg_*\mathcal{F} = \colim_{i \in I} h_i^{-1}R^pg_{i, *}(f_{i, *}\mathcal{F})
$$
\end{lemma}

\begin{proof}
这是引理 \ref{etale-cohomology-lemma-relative-colimit-general} 和
\ref{etale-cohomology-lemma-linus-hamann} 的形式组合。
\end{proof}











\section{余极限与复形}
\label{etale-cohomology-section-colimit-complexes}

\noindent
本节讨论如何在各种情形下对复形系统取上同调，继续第
\ref{etale-cohomology-section-colimit} 节中关于层的讨论。
除非绝对必要，我们强烈建议读者不要阅读本节。

\begin{lemma}
\label{etale-cohomology-lemma-colimit-variant-complexes}
设 $X = \lim_{i \in I} X_i$ 是概形定向系统的极限，其仿射转移态射为
$f_{i'i} : X_{i'} \to X_i$。假设概形 $X_i$ 对所有 $i \in I$
均拟紧且拟分离。设 $\mathcal{F}_i^\bullet$ 是
$X_{i, \etale}$ 上的 Abel 层复形。设
$\varphi_{i'i} : f_{i'i}^{-1}\mathcal{F}_i^\bullet \to
\mathcal{F}_{i'}^\bullet$ 是 $X_{i, \etale}$ 上的复形映射，
并且等式
$\varphi_{i''i} = \varphi_{i''i'} \circ f_{i'' i'}^{-1}\varphi_{i'i}$
在每当 $i'' \geq i' \geq i$ 时成立。假设存在整数 $a$，使得
$\mathcal{F}_i^n = 0$ 对 $n < a$ 以及所有 $i \in I$ 成立。则
$$
H^p_\etale(X, \colim f_i^{-1}\mathcal{F}_i^\bullet) =
\colim H^p_\etale(X_i, \mathcal{F}^\bullet_i)
$$
，其中 $f_i : X \to X_i$ 是投影。
\end{lemma}

\begin{proof}
这是定理 \ref{etale-cohomology-theorem-colimit} 的推论。置
$\mathcal{F}^\bullet = \colim f_i^{-1}\mathcal{F}_i^\bullet$。
该定理告诉我们
$$
\colim_{i \in I} H_\etale^p(X_i, \mathcal{F}_i^n) =
H_\etale^p(X, \mathcal{F}^n)
$$
对所有 $n, p \in \mathbf{Z}$ 成立。使用谱序列
$$
E_{1, i}^{s, t} = H_\etale^t(X_i, \mathcal{F}_i^s) \Rightarrow
H_\etale^{s + t}(X_i, \mathcal{F}_i^\bullet)
$$
以及
$$
E_1^{s, t} = H_\etale^t(X, \mathcal{F}^s) \Rightarrow
H_\etale^{s + t}(X, \mathcal{F}^\bullet)
$$
；见 Derived Categories，引理 \ref{derived-lemma-two-ss-complex-functor}。
由于 $\mathcal{F}_i^n = 0$ 对 $n < a$ 成立（其中 $a$ 与 $i$ 无关），
可见只有固定的有限多个项 $E_{1, i}^{s, t}$（与 $i$ 无关）和
$E_1^{s, t}$ 对
$H^q_\etale(X_i, \mathcal{F}_i^\bullet)$ 与
$H^q_\etale(X, \mathcal{F}^\bullet)$ 有贡献，并且
$E_1^{s, t} = \colim E_{i, i}^{s, t}$。
这就给出了所需结论；略去一些细节。
（也可使用复形的“笨截断”给出另一论证，从而避免使用谱序列。）
\end{proof}

\begin{lemma}
\label{etale-cohomology-lemma-direct-sum-bounded-below-cohomology}
设 $X$ 是拟紧且拟分离概形。设
$K_i \in D(X_\etale)$，$i \in I$，是一族对象。
假设给定 $a \in \mathbf{Z}$，使得 $H^n(K_i) = 0$ 对
$n < a$ 和 $i \in I$ 成立。则
$R\Gamma(X, \bigoplus_i K_i) =
\bigoplus_i R\Gamma(X, K_i)$。
\end{lemma}

\begin{proof}
必须证明
$H^p(X, \bigoplus_i K_i) =
\bigoplus_i H^p(X, K_i)$ 对所有 $p \in \mathbf{Z}$ 成立。
选取复形 $\mathcal{F}_i^\bullet$ 来表示 $K_i$，使得
$\mathcal{F}_i^n = 0$ 对 $n < a$ 成立。
由 Injectives，引理 \ref{injectives-lemma-derived-products}，
复形 $\mathcal{F}_i^\bullet$ 的直和表示对象 $\bigoplus K_i$。
由于 $\bigoplus \mathcal{F}^\bullet$ 是有限直和的滤过余极限，
结论由引理 \ref{etale-cohomology-lemma-colimit-variant-complexes} 得到。
\end{proof}

\begin{lemma}
\label{etale-cohomology-lemma-colimit-complexes}
设 $S$ 是概形。设 $X = \lim_{i \in I} X_i$ 是 $S$ 上概形定向系统的极限，
其仿射转移态射为 $f_{i'i} : X_{i'} \to X_i$。假设概形 $X_i$ 对所有
$i \in I$ 均拟紧且拟分离。设
$K \in D^+(S_\etale)$。则
$$
\colim_{i \in I} H_\etale^p(X_i, K|_{X_i}) = H_\etale^p(X, K|_X).
$$
对所有 $p \in \mathbf{Z}$ 成立，其中 $K|_{X_i}$ 与 $K|_X$
分别是 $K$ 到 $X_i$ 与 $X$ 的拉回。
\end{lemma}

\begin{proof}
可以将 $K$ 表示为向下有界复形 $\mathcal{G}^\bullet$；
该复形由 $S_\etale$ 上的 Abel 层组成。设 $\mathcal{G}^n = 0$ 对
$n < a$ 成立。分别以 $\mathcal{F}^\bullet_i$ 和
$\mathcal{F}^\bullet$ 表示这个复形到 $X_i$ 和 $X$ 的拉回。
这些复形分别表示对象 $K|_{X_i}$ 和 $K|_X$，而且逐项有
$\mathcal{F}^\bullet = \colim f_i^{-1}\mathcal{F}_i^\bullet$。
因此，引理由引理 \ref{etale-cohomology-lemma-colimit-variant-complexes} 得到。
\end{proof}

\begin{lemma}
\label{etale-cohomology-lemma-relative-colimit-general-complexes}
令 $I$、$g_i : X_i \to S_i$、$g : X \to S$、$f_i$、$g_i$、$h_i$
如引理 \ref{etale-cohomology-lemma-relative-colimit-general} 中所述。
设 $0 \in I$ 且 $K_0 \in D^+(X_{0, \etale})$。
对 $i \geq 0$，以 $K_i$ 表示 $K_0$ 到 $X_i$ 的拉回；
以 $K$ 表示 $K_0$ 到 $X$ 的拉回。则
$$
R^pg_*K = \colim_{i \geq 0} h_i^{-1}R^pg_{i, *}K_i
$$
对所有 $p \in \mathbf{Z}$ 成立。
\end{lemma}

\begin{proof}
固定整数 $p_0 \in \mathbf{Z}$。
设 $a$ 是整数，使得 $H^j(K_0) = 0$ 对 $j < a$ 成立。
我们将证明该公式对所有 $p \leq p_0$ 成立，并对 $a$ 作下降归纳。
若 $a > p_0$，则由平凡消失可见，对 $p \leq p_0$，公式两边均为零；
见 Derived Categories，引理 \ref{derived-lemma-negative-vanishing}。
假设 $a \leq p_0$。考虑可区别三角
$$
H^a(K_0)[-a] \to K_0 \to \tau_{\geq a + 1}K_0
$$
将这个可区别三角拉回到 $X_i$ 和 $X$，得到关于 $K_i$ 和 $K$
的相容可区别三角。对 $p \leq p_0$，考虑交换图
$$
\xymatrix{
\colim_{i \geq 0} h_i^{-1}R^{p - 1}g_{i, *}(\tau_{\geq a + 1}K_i)
\ar[r]_-\alpha \ar[d] &
R^{p - 1}g_*(\tau_{\geq a + 1}K) \ar[d] \\
\colim_{i \geq 0} h_i^{-1}R^pg_{i, *}(H^a(K_i)[-a])
\ar[r]_-\beta \ar[d] &
R^pg_*(H^a(K)[-a]) \ar[d] \\
\colim_{i \geq 0} h_i^{-1}R^pg_{i, *}K_i
\ar[r]_-\gamma \ar[d] &
R^pg_*K \ar[d] \\
\colim_{i \geq 0} R^pg_{i, *}\tau_{\geq a + 1}K_i
\ar[r]_-\delta \ar[d] &
R^pg_*\tau_{\geq a + 1}K \ar[d] \\
\colim_{i \geq 0} R^{p + 1}g_{i, *}(H^a(K_i)[-a])
\ar[r]^-\epsilon &
R^{p + 1}g_*(H^a(K)[-a])
}
$$
，其各列正合。由引理 \ref{etale-cohomology-lemma-relative-colimit-general}，
箭头 $\beta$ 和 $\epsilon$ 是同构。
由归纳假设，箭头 $\alpha$ 和 $\delta$ 是同构。
因此 $\gamma$ 是同构，正合所求。
\end{proof}

\begin{lemma}
\label{etale-cohomology-lemma-relative-colimit-general-really-complexes}
令 $I$、$g_i : X_i \to S_i$、$g : X \to S$、$f_{ii'}$、$f_i$、$g_i$、$h_i$
如引理 \ref{etale-cohomology-lemma-relative-colimit-general} 中所述。
设 $\mathcal{F}_i^\bullet$ 是 $X_{i, \etale}$ 上的 Abel 层复形。
设 $\varphi_{i'i} : f_{i'i}^{-1}\mathcal{F}_i^\bullet \to
\mathcal{F}_{i'}^\bullet$ 是 $X_{i, \etale}$ 上的复形映射，
并且等式
$\varphi_{i''i} = \varphi_{i''i'} \circ f_{i'' i'}^{-1}\varphi_{i'i}$
在每当 $i'' \geq i' \geq i$ 时成立。假设存在整数 $a$，使得
$\mathcal{F}_i^n = 0$ 对 $n < a$ 以及所有 $i \in I$ 成立。则
$$
R^pg_*(\colim f_i^{-1}\mathcal{F}_i^\bullet) =
\colim_{i \geq 0} h_i^{-1}R^pg_{i, *}\mathcal{F}_i^\bullet
$$
对所有 $p \in \mathbf{Z}$ 成立。
\end{lemma}

\begin{proof}
这是引理 \ref{etale-cohomology-lemma-relative-colimit-general} 的推论。置
$\mathcal{F}^\bullet = \colim f_i^{-1}\mathcal{F}_i^\bullet$。
该引理告诉我们
$$
\colim_{i \in I} h_i^{-1}R^pg_{i, *}\mathcal{F}_i^n = R^pg_*\mathcal{F}^n
$$
对所有 $n, p \in \mathbf{Z}$ 成立。使用谱序列
$$
E_{1, i}^{s, t} = R^tg_{i, *}\mathcal{F}_i^s \Rightarrow
R^{s + t}g_{i, *}\mathcal{F}_i^\bullet
$$
以及
$$
E_1^{s, t} = R^tg_*\mathcal{F}^s \Rightarrow
R^{s + t}g_*\mathcal{F}^\bullet
$$
；见 Derived Categories，引理 \ref{derived-lemma-two-ss-complex-functor}。
由于 $\mathcal{F}_i^n = 0$ 对 $n < a$ 成立（其中 $a$ 与 $i$ 无关），
可见只有固定的有限多个项 $E_{1, i}^{s, t}$（与 $i$ 无关）和
$E_1^{s, t}$ 有贡献，并且
$E_1^{s, t} = \colim E_{i, i}^{s, t}$。
这就给出了所需结论；略去一些细节。
（也可使用复形的“笨截断”给出另一论证，从而避免使用谱序列。）
\end{proof}

\begin{lemma}
\label{etale-cohomology-lemma-direct-sum-bounded-below-pushforward}
设 $f : X \to Y$ 是拟紧且拟分离的概形态射。设
$K_i \in D(X_\etale)$，$i \in I$，是一族对象。
假设给定 $a \in \mathbf{Z}$，使得 $H^n(K_i) = 0$ 对
$n < a$ 和 $i \in I$ 成立。则
$Rf_*(\bigoplus_i K_i) = \bigoplus_i Rf_*K_i$。
\end{lemma}

\begin{proof}
必须证明
$R^pf_*(\bigoplus_i K_i) =
\bigoplus_i R^pf_*K_i$ 对所有 $p \in \mathbf{Z}$ 成立。
选取复形 $\mathcal{F}_i^\bullet$ 来表示 $K_i$，使得
$\mathcal{F}_i^n = 0$ 对 $n < a$ 成立。
由 Injectives，引理 \ref{injectives-lemma-derived-products}，
复形 $\mathcal{F}_i^\bullet$ 的直和表示对象 $\bigoplus K_i$。
由于 $\bigoplus \mathcal{F}^\bullet$ 是有限直和的滤过余极限，
结论由引理 \ref{etale-cohomology-lemma-relative-colimit-general-really-complexes} 得到。
\end{proof}














\section{高阶直接像的茎}
\label{etale-cohomology-section-stalks-direct-image}

\noindent
高阶直接像的茎常可按如下方式计算。

\begin{theorem}
\label{etale-cohomology-theorem-higher-direct-images}
设 $f: X \to S$ 是拟紧且拟分离的概形态射，
$\mathcal{F}$ 是 $X_\etale$ 上的 Abel 层，并且 $\overline{s}$ 是
$S$ 的一个位于 $s \in S$ 上方的几何点。则
$$
\left(R^nf_* \mathcal{F}\right)_{\overline{s}} =
H_\etale^n( X \times_S \Spec(\mathcal{O}_{S, \overline{s}}^{sh}),
p^{-1}\mathcal{F})
$$
，其中 $p : X \times_S \Spec(\mathcal{O}_{S, \overline{s}}^{sh}) \to X$
是投影。对 $K \in D^+(X_\etale)$ 和 $n \in \mathbf{Z}$，有
$$
\left(R^nf_*K\right)_{\overline{s}} =
H_\etale^n(X \times_S \Spec(\mathcal{O}_{S, \overline{s}}^{sh}), p^{-1}K)
$$
事实上，有
$$
\left(Rf_*K\right)_{\overline{s}}
=
R\Gamma_\etale(X \times_S \Spec(\mathcal{O}_{S, \overline{s}}^{sh}), p^{-1}K)
$$
作为 $D^+(\textit{Ab})$ 中的等式。
\end{theorem}

\begin{proof}
设 $\mathcal{I}$ 是 $\overline{s}$ 在 $S$ 上的 \'etale 邻域所成的范畴。
由引理 \ref{etale-cohomology-lemma-higher-direct-images}，有
$$
(R^nf_*\mathcal{F})_{\overline{s}} =
\colim_{(V, \overline{v}) \in \mathcal{I}^{opp}}
H_\etale^n(X \times_S V, \mathcal{F}|_{X \times_S V}).
$$
可以将 $\mathcal{I}$ 替换为由 $\overline{s}$ 的仿射 \'etale 邻域
组成的始子范畴。注意
$$
\Spec(\mathcal{O}_{S, \overline{s}}^{sh}) =
\lim_{(V, \overline{v}) \in \mathcal{I}} V
$$
；这是引理 \ref{etale-cohomology-lemma-describe-etale-local-ring} 以及
Limits，引理
\ref{limits-lemma-directed-inverse-system-affine-schemes-has-limit}
的结论。由于纤维积与极限交换，还得到
$$
X \times_S \Spec(\mathcal{O}_{S, \overline{s}}^{sh}) =
\lim_{(V, \overline{v}) \in \mathcal{I}} X \times_S V
$$
结论由引理 \ref{etale-cohomology-lemma-directed-colimit-cohomology} 得到。
对于第二种形式，使用相同论证，但以引理
\ref{etale-cohomology-lemma-colimit-complexes} 代替引理
\ref{etale-cohomology-lemma-directed-colimit-cohomology}。

\medskip\noindent
为证明最后一个陈述，只需构造映射
$\left(Rf_*K\right)_{\overline{s}} \to
R\Gamma_\etale(X \times_S \Spec(\mathcal{O}_{S, \overline{s}}^{sh}), p^{-1}K)$，
它属于 $D^+(\textit{Ab})$，并在上述每个次数 $n$
的上同调群上实现同构（对所有 $n$）。
为此，选取向下有界复形 $\mathcal{J}^\bullet$，它由
$X_\etale$ 上的内射 Abel 层组成，并用它表示 $K$。复形
$f_*\mathcal{J}^\bullet$ 表示 $Rf_*K$。因此，复形
$$
(f_*\mathcal{J}^\bullet)_{\overline{s}} =
\colim_{(V, \overline{v}) \in \mathcal{I}^{opp}}
(f_*\mathcal{J}^\bullet)(V)
$$
表示 $(Rf_*K)_{\overline{s}}$。对每个 $V$，有映射
$$
(f_*\mathcal{J}^\bullet)(V) =
\Gamma(X \times_S V, \mathcal{J}^\bullet)
\longrightarrow
\Gamma(X \times_S \Spec(\mathcal{O}_{S, \overline{s}}^{sh}),
p^{-1}\mathcal{J}^\bullet)
$$
，而靶复形表示
$R\Gamma_\etale(X \times_S \Spec(\mathcal{O}_{S, \overline{s}}^{sh}), p^{-1}K)$，
这是 $D^+(\textit{Ab})$ 中的对象。
对这些映射取余极限，即得结论。
\end{proof}

\begin{remark}
\label{etale-cohomology-remark-stalk-fibre}
设 $f : X \to S$ 是概形态射。
设 $K \in D(X_\etale)$。
设 $\overline{s}$ 是 $S$ 的几何点。
总有典范映射
$$
(Rf_*K)_{\overline{s}}
\longrightarrow
R\Gamma(X \times_S \Spec(\mathcal{O}_{S, \overline{s}}^{sh}), p^{-1}K)
\longrightarrow
R\Gamma(X_{\overline{s}}, K|_{X_{\overline{s}}})
$$
，其中 $p : X \times_S \Spec(\mathcal{O}_{S, \overline{s}}^{sh}) \to X$
是投影。具体而言，考虑交换图
$$
\xymatrix{
X_{\overline{s}} \ar[r] \ar[d]^{f_{\overline{s}}} &
X \times_S \Spec(\mathcal{O}_{S, \overline{s}}^{sh}) \ar[r]_-p \ar[d]^{f'} &
X \ar[d]^f \\
\overline{s} \ar[r]^-i &
\Spec(\mathcal{O}_{S, \overline{s}}^{sh}) \ar[r]^-j &
S
}
$$
对于该图中的两个方块，有基变换映射
$$
i^{-1}Rf'_*(p^{-1}K) \to Rf_{\overline{s}, *}(K|_{X_{\overline{s}}})
\quad\text{且}\quad
j^{-1}Rf_*K \to Rf'_*(p^{-1}K)
$$
（Cohomology on Sites，注
\ref{sites-cohomology-remark-base-change}）。
取整体截面，即得所需映射。
由 Cohomology on Sites，注
\ref{sites-cohomology-remark-compose-base-change-horizontal}，
这两个映射的复合就是通常的（基变换）映射
$(Rf_*K)_{\overline{s}} \to R\Gamma(X_{\overline{s}}, K|_{X_{\overline{s}}})$。
\end{remark}






\section{Leray 谱序列}
\label{etale-cohomology-section-leray}

\begin{lemma}
\label{etale-cohomology-lemma-prepare-leray}
设 $f: X \to Y$ 是态射，而 $\mathcal{I}$ 是
$\textit{Ab}(X_\etale)$ 的内射对象。设 $V \in \Ob(Y_\etale)$。则
\begin{enumerate}
\item 对任意覆盖 $\mathcal{V} = \{V_j\to V\}_{j \in J}$，有
$\check H^p(\mathcal{V}, f_*\mathcal{I}) = 0$ 对所有 $p > 0$ 成立；
\item $f_*\mathcal{I}$ 关于函子 $\Gamma(V, -)$ 是无环的；并且
\item 若 $g : Y \to Z$，则 $f_*\mathcal{I}$ 关于 $g_*$ 是无环的。
\end{enumerate}
\end{lemma}

\begin{proof}
注意
$\check{\mathcal{C}}^\bullet(\mathcal{V}, f_*\mathcal{I}) =
\check{\mathcal{C}}^\bullet(\mathcal{V} \times_Y X, \mathcal{I})$，
由引理 \ref{etale-cohomology-lemma-hom-injective}，它的高阶上同调群消失。
这证明了 (1)。第二个陈述如下得到：若一个层对任意覆盖的高阶
{\v C}ech 上同调群均消失，则它的高阶上同调群也消失。
用 {\v C}ech 到上同调的谱序列证明这一点是一道很精彩的练习；
细节以及更精确、更一般的陈述见 Cohomology on Sites，引理
\ref{sites-cohomology-lemma-cech-vanish-collection}。
(3) 是 (2) 以及引理 \ref{etale-cohomology-lemma-higher-direct-images} 中
$R^pg_*$ 的描述的推论。
\end{proof}

\noindent
应用 Grothendieck 谱序列的形式体系，得到如下结论。

\begin{proposition}[Leray 谱序列]
\label{etale-cohomology-proposition-leray}
设 $f: X \to Y$ 是概形态射，而 $\mathcal{F}$ 是 $X$ 上的 \'etale 层。
则存在谱序列
$$
E_2^{p, q} = H_\etale^p(Y, R^qf_*\mathcal{F}) \Rightarrow
H_\etale^{p+q}(X, \mathcal{F}).
$$
\end{proposition}

\begin{proof}
见引理 \ref{etale-cohomology-lemma-prepare-leray} 以及 Derived Categories，第
\ref{derived-section-composition-right-derived-functors} 节。
\end{proof}









\section{有限态射高阶直接像的消失}
\label{etale-cohomology-section-vanishing-finite-morphism}

\noindent
下一个目标是证明，概形有限态射的高阶直接像消失。

\begin{lemma}
\label{etale-cohomology-lemma-vanishing-etale-cohomology-strictly-henselian}
设 $R$ 是严格 Hensel 局部环。置 $S = \Spec(R)$，并设
$\overline{s}$ 是它的闭点。则整体截面函子
$\Gamma(S, -) : \textit{Ab}(S_\etale) \to \textit{Ab}$
是正合的。事实上，有
$\Gamma(S, \mathcal{F}) = \mathcal{F}_{\overline{s}}$，
对任意集合层 $\mathcal{F}$ 成立。特别地，
$$
\forall p\geq 1, \quad H_\etale^p(S, \mathcal{F})=0
$$
对所有 $\mathcal{F}\in \textit{Ab}(S_\etale)$ 成立。
\end{lemma}

\begin{proof}
若证明了 $\Gamma(S, \mathcal{F}) = \mathcal{F}_{\overline{s}}$，
则由于茎函子正合，$\Gamma(S, -)$ 也正合。
设 $(U, \overline{u})$ 是 $\overline{s}$ 的一个 \'etale 邻域。
选取仿射开邻域 $\Spec(A)$；它是 $\overline{u}$ 在 $U$ 中的邻域。
于是 $R \to A$ 是 \'etale 的，并且
$\kappa(\overline{s}) = \kappa(\overline{u})$。
由定理 \ref{etale-cohomology-theorem-henselian}，可见存在同构
$A \cong R \times A'$；它是 $R$-代数同构，并与到
$\kappa(\overline{s}) = \kappa(\overline{u})$ 的映射相容。因此得到截面
$$
\xymatrix{
\Spec(A) \ar[r] & U \ar[d]\\
& S \ar[ul]
}
$$
于是，在 $\overline{s}$ 的 \'etale 邻域系统中，恒等映射
$(S, \overline{s}) \to (S, \overline{s})$ 是余终的。
因此 $\Gamma(S, \mathcal{F}) = \mathcal{F}_{\overline{s}}$。
引理的最后一个陈述是因为正合函子的高阶导出函子为零；见
Derived Categories，引理 \ref{derived-lemma-right-derived-exact-functor}。
\end{proof}

\begin{proposition}
\label{etale-cohomology-proposition-finite-higher-direct-image-zero}
设 $f : X \to Y$ 是概形的有限态射。
\begin{enumerate}
\item 对任意几何点 $\overline{y} : \Spec(k) \to Y$，有
$$
(f_*\mathcal{F})_{\overline{y}} =
\prod\nolimits_{\overline{x} : \Spec(k) \to X,\ f(\overline{x}) =
\overline{y}} \mathcal{F}_{\overline{x}}.
$$
对 $\mathcal{F}$ 属于 $\Sh(X_\etale)$ 的情形成立，并且
$$
(f_*\mathcal{F})_{\overline{y}} =
\bigoplus\nolimits_{\overline{x} : \Spec(k) \to X,\ f(\overline{x}) =
\overline{y}} \mathcal{F}_{\overline{x}}.
$$
对 $\mathcal{F}$ 属于 $\textit{Ab}(X_\etale)$ 的情形成立。
\item 对任意 $q \geq 1$，有 $R^q f_*\mathcal{F} = 0$，
其中 $\mathcal{F}$ 属于 $\textit{Ab}(X_\etale)$。
\end{enumerate}
\end{proposition}

\begin{proof}
以 $X_{\overline{y}}^{sh}$ 表示纤维积
$X \times_Y \Spec(\mathcal{O}_{Y, \overline{y}}^{sh})$。
由定理 \ref{etale-cohomology-theorem-higher-direct-images}，
$R^qf_*\mathcal{F}$ 在 $\overline{y}$ 处的茎由
$H_\etale^q(X_{\overline{y}}^{sh}, \mathcal{F})$ 计算。
由于 $f$ 有限，$X_{\bar y}^{sh}$ 在
$\Spec(\mathcal{O}_{Y, \overline{y}}^{sh})$ 上有限，因而有
$X_{\bar y}^{sh} = \Spec(A)$，其中某个环 $A$ 在
$\mathcal{O}_{Y, \bar y}^{sh}$ 上有限。
由于后一个环严格 Hensel，引理 \ref{etale-cohomology-lemma-finite-over-henselian}
说明 $A$ 是 Hensel 局部环的有限乘积
$A = A_1 \times \ldots \times A_r$。由于
$\mathcal{O}_{Y, \overline{y}}^{sh}$ 的剩余域可分闭，每个
$A_i$ 也具有该性质。因此 $A_i$ 严格 Hensel。
这说明 $X_{\overline{y}}^{sh} = \coprod_{i = 1}^r \Spec(A_i)$。
引理 \ref{etale-cohomology-lemma-vanishing-etale-cohomology-strictly-henselian}
中的消失给出 $(R^qf_*\mathcal{F})_{\overline{y}} = 0$ 对
$q > 0$ 成立，再由定理 \ref{etale-cohomology-theorem-exactness-stalks} 得到 (2)。
(1) 由引理 \ref{etale-cohomology-lemma-vanishing-etale-cohomology-strictly-henselian}
中的相应陈述得到。
\end{proof}

\begin{lemma}
\label{etale-cohomology-lemma-finite-pushforward-commutes-with-base-change}
考虑笛卡尔方块
$$
\xymatrix{
X' \ar[r]_{g'} \ar[d]_{f'} & X \ar[d]^f \\
Y' \ar[r]^g & Y
}
$$
，其中 $f$ 是有限态射。
对任意集合层 $\mathcal{F}$（定义在 $X_\etale$ 上），有
$f'_*(g')^{-1}\mathcal{F} = g^{-1}f_*\mathcal{F}$。
\end{lemma}

\begin{proof}
在极大的一般性下，存在拉回映射
$g^{-1}f_*\mathcal{F} \to f'_*(g')^{-1}\mathcal{F}$；见
Sites，第 \ref{sites-section-pullback} 节。
只需在茎上验证（定理 \ref{etale-cohomology-theorem-exactness-stalks}）。
设 $\overline{y}' : \Spec(k) \to Y'$ 是几何点。有
\begin{align*}
(f'_*(g')^{-1}\mathcal{F})_{\overline{y}'}
& =
\prod\nolimits_{\overline{x}' : \Spec(k) \to X',\ f' \circ \overline{x}' =
\overline{y}'}
((g')^{-1}\mathcal{F})_{\overline{x}'} \\
& =
\prod\nolimits_{\overline{x}' : \Spec(k) \to X',\ f' \circ \overline{x}' =
\overline{y}'} \mathcal{F}_{g' \circ \overline{x}'} \\
& =
\prod\nolimits_{\overline{x} : \Spec(k) \to X,\ f \circ \overline{x} =
g \circ \overline{y}'} \mathcal{F}_{\overline{x}} \\
& =
(f_*\mathcal{F})_{g \circ \overline{y}'} \\
& =
(g^{-1}f_*\mathcal{F})_{\overline{y}'}
\end{align*}
第一个等式由命题 \ref{etale-cohomology-proposition-finite-higher-direct-image-zero} 得到。
第二个等式由引理 \ref{etale-cohomology-lemma-stalk-pullback} 得到。
第三个等式成立，是因为该图是笛卡尔方块，故映射
$$
\{\overline{x}' : \Spec(k) \to X',\ f' \circ \overline{x}' =
\overline{y}'\}
\longrightarrow
\{\overline{x} : \Spec(k) \to X,\ f \circ \overline{x} =
g \circ \overline{y}'\}
$$
把 $\overline{x}'$ 映到 $g' \circ \overline{x}'$，且是双射。
第四个等式由命题 \ref{etale-cohomology-proposition-finite-higher-direct-image-zero} 得到。
第五个等式由引理 \ref{etale-cohomology-lemma-stalk-pullback} 得到。
\end{proof}

\begin{lemma}
\label{etale-cohomology-lemma-integral-pushforward-commutes-with-base-change}
考虑笛卡尔方块
$$
\xymatrix{
X' \ar[r]_{g'} \ar[d]_{f'} & X \ar[d]^f \\
Y' \ar[r]^g & Y
}
$$
，其中 $f$ 是整态射。
对任意集合层 $\mathcal{F}$（定义在 $X_\etale$ 上），有
$f'_*(g')^{-1}\mathcal{F} = g^{-1}f_*\mathcal{F}$。
\end{lemma}

\begin{proof}
问题在 $Y$ 上是局部的，因而可以假设 $Y$ 仿射。
于是可写成 $X = \lim X_i$，其中 $f_i : X_i \to Y$ 有限
（在仿射情形这很容易，不过可参见
Limits，引理 \ref{limits-lemma-integral-limit-finite-and-finite-presentation}）。
以 $p_{i'i} : X_{i'} \to X_i$ 表示转移态射，并以
$p_i : X \to X_i$ 表示投影。置
$\mathcal{F}_i = p_{i, *}\mathcal{F}$，则由引理
\ref{etale-cohomology-lemma-linus-hamann} 得到系统
$(\mathcal{F}_i, \varphi_{i'i})$，它满足
$\mathcal{F} = \colim p_i^{-1}\mathcal{F}_i$。
由引理 \ref{etale-cohomology-lemma-relative-colimit} 得到
$f_*\mathcal{F} = \colim f_{i, *}\mathcal{F}_i$。
置 $X'_i = Y' \times_Y X_i$，其投影为 $f'_i$ 与 $g'_i$。
由于极限与极限交换，有 $X' = \lim X'_i$。
以 $p'_i : X' \to X'_i$ 表示投影。有
\begin{align*}
g^{-1}f_*\mathcal{F}
& =
g^{-1} \colim f_{i, *}\mathcal{F}_i \\
& =
\colim g^{-1}f_{i, *}\mathcal{F}_i \\
& =
\colim f'_{i, *}(g'_i)^{-1}\mathcal{F}_i \\
& =
f'_*(\colim (p'_i)^{-1}(g'_i)^{-1}\mathcal{F}_i) \\
& =
f'_*(\colim (g')^{-1}p_i^{-1}\mathcal{F}_i) \\
& =
f'_*(g')^{-1} \colim p_i^{-1}\mathcal{F}_i \\
& =
f'_*(g')^{-1}\mathcal{F}
\end{align*}
正合所求。第一个等式见上文。
第二个等式使用拉回与余极限交换。
第三个等式使用有限情形；见引理
\ref{etale-cohomology-lemma-finite-pushforward-commutes-with-base-change}。
第四个等式使用引理 \ref{etale-cohomology-lemma-relative-colimit}。
第五个等式使用 $g'_i \circ p'_i = p_i \circ g'$。
第六个等式使用拉回与余极限交换。
第七个等式使用 $\mathcal{F} = \colim p_i^{-1}\mathcal{F}_i$。
\end{proof}

\noindent
下一个引理是处理 \'etale 层与满有限态射的一种上同调下降。
证明正常基变换定理后，我们将大幅推广这一结果。

\begin{lemma}
\label{etale-cohomology-lemma-cohomological-descent-finite}
设 $f : X \to Y$ 是概形的满有限态射。
置 $f_n : X_n \to Y$ 为 $(n + 1)$ 重纤维积；该纤维积由
$X$ 在 $Y$ 上取得。
对 $\mathcal{F} \in \textit{Ab}(Y_\etale)$，置
$\mathcal{F}_n = f_{n, *}f_n^{-1}\mathcal{F}$。存在正合序列
$$
0 \to \mathcal{F} \to \mathcal{F}_0 \to \mathcal{F}_1 \to
\mathcal{F}_2 \to \ldots
$$
于 $Y_\etale$ 上。此外，存在谱序列
$$
E_1^{p, q} = H^q_\etale(X_p, f_p^{-1}\mathcal{F})
$$
收敛到 $H^{p + q}(Y_\etale, \mathcal{F})$。
该谱序列关于 $\mathcal{F}$ 函子性。
\end{lemma}

\begin{proof}
若证明了引理的第一个陈述，则得到以
$E_1^{p, q} = H^q_\etale(Y, \mathcal{F}_p)$ 为首项并收敛到
$H^{p + q}(Y_\etale, \mathcal{F})$ 的谱序列；见
Derived Categories，引理 \ref{derived-lemma-two-ss-complex-functor}。
另一方面，由于
$R^if_{p, *}f_p^{-1}\mathcal{F} = 0$ 对 $i > 0$ 成立
（命题 \ref{etale-cohomology-proposition-finite-higher-direct-image-zero}），得到
$$
H^q_\etale(X_p, f_p^{-1}\mathcal{F}) =
H^q_\etale(Y, f_{p, *}f_p^{-1} \mathcal{F}) =
H^q_\etale(Y, \mathcal{F}_p)
$$
；这是命题 \ref{etale-cohomology-proposition-leray} 的结论，从而得到引理中的谱序列。

\medskip\noindent
为证明引理的第一个陈述，注意 $X_n$ 构成 $Y$ 上的单纯概形；
见 Simplicial，例 \ref{simplicial-example-fibre-products-simplicial-object}。
还要注意，对每个投影 $d_j : X_{n + 1} \to X_n$，存在映射
$d_j^{-1} f_n^{-1}\mathcal{F} \to f_{n + 1}^{-1}\mathcal{F}$。
这些映射诱导映射
$$
\delta_j : \mathcal{F}_n \to \mathcal{F}_{n + 1}
$$
，其中 $j = 0, \ldots, n + 1$。我们用这些映射的交错和来定义微分
$\mathcal{F}_n \to \mathcal{F}_{n + 1}$。
类似地，有典范增广 $\mathcal{F} \to \mathcal{F}_0$；
它就是典范映射 $\mathcal{F} \to f_*f^{-1}\mathcal{F}$。
要验证这个层序列是正合复形，只需在几何点的茎上验证
（定理 \ref{etale-cohomology-theorem-exactness-stalks}）。
因此，设 $\overline{y} : \Spec(k) \to Y$ 是几何点。置
$E = \{\overline{x} : \Spec(k) \to X \mid f(\overline{x}) = \overline{y}\}$。
于是 $E$ 是有限非空集合，并且
$$
(\mathcal{F}_n)_{\overline{y}} =
\bigoplus\nolimits_{e \in E^{n + 1}} \mathcal{F}_{\overline{y}}
$$
；这是命题 \ref{etale-cohomology-proposition-finite-higher-direct-image-zero}
和引理 \ref{etale-cohomology-lemma-stalk-pullback} 的结论。
因此，必须验证：给定 Abel 群 $M$，序列
$$
0 \to M \to \bigoplus\nolimits_{e \in E} M \to
\bigoplus\nolimits_{e \in E^2} M \to \ldots
$$
正合。这里第一个映射是对角映射，而映射
$\bigoplus_{e \in E^{n + 1}} M  \to \bigoplus_{e \in E^{n + 2}} M$
是由 $(n + 2)$ 个投影 $E^{n + 2} \to E^{n + 1}$ 所诱导映射的交错和。
这可以直接证明，也可将 Simplicial，引理
\ref{simplicial-lemma-fibre-products-simplicial-object-w-section}
应用于映射 $E \to \{*\}$ 得到。
\end{proof}

\begin{remark}
\label{etale-cohomology-remark-cohomological-descent-finite}
在引理 \ref{etale-cohomology-lemma-cohomological-descent-finite} 的情形中，
若 $\mathcal{G}$ 是 $Y_\etale$ 上的集合层，则
$$
\Gamma(Y, \mathcal{G}) =
\text{等化子}(
\xymatrix{
\Gamma(X_0, f_0^{-1}\mathcal{G})
\ar@<1ex>[r] \ar@<-1ex>[r] &
\Gamma(X_1, f_1^{-1}\mathcal{G})
}
)
$$
证明完全相同：说明层 $\mathcal{G}$ 是两个映射
$f_{0, *}f_0^{-1}\mathcal{G} \to f_{1, *}f_1^{-1}\mathcal{G}$
的等化子即可。
\end{remark}










%10.06.09
\section{茎上的 Galois 作用}
\label{etale-cohomology-section-galois-action-stalks}

\noindent
本节定义如下作用：点 $s$（其所在概形为 $S$）的剩余域之绝对 Galois 群，
作用于任意位于 $s$ 上方的几何点处的茎函子。

\medskip\noindent
茎上的 Galois 作用。
设 $S$ 是概形。
设 $\overline{s}$ 是 $S$ 的几何点。
设 $\sigma \in \text{Aut}(\kappa(\overline{s})/\kappa(s))$。
如下定义 $\sigma$ 在茎 $\mathcal{F}_{\overline{s}}$
上的作用，其中 $\mathcal{F}$ 是层：
\begin{equation}
\label{etale-cohomology-equation-galois-action}
\begin{matrix}
\mathcal{F}_{\overline{s}} &
\longrightarrow &
\mathcal{F}_{\overline{s}} \\
(U, \overline{u}, t) &
\longmapsto &
(U, \overline{u} \circ \Spec(\sigma), t).
\end{matrix}
\end{equation}
这里使用定义 \ref{etale-cohomology-definition-stalk} 后讨论中以三元组描述茎元素的方法。
这是左作用，因为若
$\sigma_i \in \text{Aut}(\kappa(\overline{s})/\kappa(s))$，则
\begin{align*}
\sigma_1 \cdot (\sigma_2 \cdot (U, \overline{u}, t))
& =
\sigma_1 \cdot (U, \overline{u} \circ \Spec(\sigma_2), t) \\
& =
(U, \overline{u} \circ \Spec(\sigma_2) \circ \Spec(\sigma_1), t) \\
& =
(U, \overline{u} \circ \Spec(\sigma_1 \circ \sigma_2), t) \\
& =
(\sigma_1 \circ \sigma_2) \cdot (U, \overline{u}, t)
\end{align*}
显然，这个作用关于层 $\mathcal{F}$ 是函子性的。
也可以直接引用注 \ref{etale-cohomology-remark-map-stalks} 来定义这个作用。

\begin{definition}
\label{etale-cohomology-definition-algebraic-geometric-point}
设 $S$ 是概形。
设 $\overline{s}$ 是位于点 $s$（属于 $S$）上方的几何点。
以
$\kappa(s) \subset \kappa(s)^{sep} \subset \kappa(\overline{s})$
表示 $\kappa(s)$ 在代数闭域 $\kappa(\overline{s})$ 中的可分代数闭包。
\begin{enumerate}
\item 在此情形，$\kappa(s)$ 的{\it 绝对 Galois 群}是
$\text{Gal}(\kappa(s)^{sep}/\kappa(s))$。有时记作
$\text{Gal}_{\kappa(s)}$。
\item 若几何点 $\overline{s}$ 所对应的扩张
$\kappa(s) \subset \kappa(\overline{s})$ 是 $\kappa(s)$ 的代数闭包，
则称它是{\it 代数的}。
\end{enumerate}
\end{definition}

\begin{example}
\label{etale-cohomology-example-stupid}
几何点
$\Spec(\mathbf{C}) \to \Spec(\mathbf{Q})$
不是代数的。
\end{example}

\noindent
令
$\kappa(s) \subset \kappa(s)^{sep} \subset \kappa(\overline{s})$
如定义中所述。注意，由于 $\kappa(\overline{s})$ 代数闭，映射
$$
\text{Aut}(\kappa(\overline{s})/\kappa(s))
\longrightarrow
\text{Gal}(\kappa(s)^{sep}/\kappa(s)) = \text{Gal}_{\kappa(s)}
$$
是满射。设 $(U, \overline{u})$ 是 $\overline{s}$ 的
\'etale 邻域，并设 $\overline{u}$ 位于点 $u$（属于 $U$）上方。
由于 $U \to S$ 是 \'etale 的，剩余域扩张
$\kappa(u)/\kappa(s)$ 有限可分。
由此得到
\begin{enumerate}
\item 若 $\sigma \in \text{Aut}(\kappa(\overline{s})/\kappa(s)^{sep})$，
则 $\sigma$ 在 $\mathcal{F}_{\overline{s}}$ 上的作用平凡。
\item 更精确地，
$\text{Aut}(\kappa(\overline{s})/\kappa(s))$ 的作用决定绝对 Galois 群
$\text{Gal}_{\kappa(s)}$ 在 $\mathcal{F}_{\overline{s}}$
上的作用，并且由后者决定。
\item 给定 $(U, \overline{u}, t)$，它表示元素 $\xi$，而该元素属于
$\mathcal{F}_{\overline{s}}$；群
$\text{Gal}(\kappa(s)^{sep}/K)$ 的任意元素作用平凡，其中
$\kappa(s) \subset K \subset \kappa(s)^{sep}$ 是
$\overline{u}^\sharp : \kappa(u) \to \kappa(\overline{s})$ 的像。
\end{enumerate}
综上，$\mathcal{F}_{\overline{s}}$ 成为
$\text{Gal}_{\kappa(s)}$-集合（见 Fundamental Groups，定义
\ref{pione-definition-G-set-continuous}）。
因此，可以把茎函子看作函子
$$
\Sh(S_\etale) \longrightarrow
\text{Gal}_{\kappa(s)}\textit{-Sets},
\quad
\mathcal{F} \longmapsto \mathcal{F}_{\overline{s}}
$$
；此后通常都这样看待茎函子。

\begin{theorem}
\label{etale-cohomology-theorem-equivalence-sheaves-point}
设 $S = \Spec(K)$，其中 $K$ 是域。
设 $\overline{s}$ 是 $S$ 的几何点。
以 $G = \text{Gal}_{\kappa(s)}$ 表示绝对 Galois 群。
取茎诱导范畴等价
$$
\Sh(S_\etale) \longrightarrow G\textit{-Sets},
\quad
\mathcal{F} \longmapsto \mathcal{F}_{\overline{s}}.
$$
\end{theorem}

\begin{proof}
构造这个函子的逆。在 Fundamental Groups，引理
\ref{pione-lemma-sheaves-point} 中已经看到：给定 $G$-集合 $M$，
存在 \'etale 态射
$X \to \Spec(K)$，
使得 $\Mor_K(\Spec(K^{sep}), X)$ 同构于 $M$，且这是 $G$-集合的同构。
考虑层 $\mathcal{F}$，它定义在 $\Spec(K)_\etale$ 上，并由规则
$U \mapsto \Mor_K(U, X)$ 定义。由于 \'etale 拓扑是次典范的，
这确实是层。于是可见
$\mathcal{F}_{\overline{s}} = \Mor_K(\Spec(K^{sep}), X) = M$
作为 $G$-集合成立（略去细节）。这给出了该函子的逆，结论得证。
\end{proof}

\begin{remark}
\label{etale-cohomology-remark-every-sheaf-representable}
定理 \ref{etale-cohomology-theorem-equivalence-sheaves-point} 与 Fundamental Groups，引理
\ref{pione-lemma-sheaves-point} 的结论也可以表述为：
$\Spec(K)_\etale$ 上的每个层都可由概形 $X$ 表示，且它
在 $\Spec(K)$ 上是 \'etale 的。
这并不意味着每个层在 Sites，定义
\ref{sites-definition-representable-sheaf} 的意义下都可表示。
原因是，在构造 $\Spec(K)_\etale$ 时，我们选取了一个充分大的、在
$\Spec(K)$ 上 \'etale 的概形集合，而
$\Spec(K)_\etale$ 上的层构成真类。
\end{remark}

\begin{lemma}
\label{etale-cohomology-lemma-global-sections-point}
假设与记号同定理 \ref{etale-cohomology-theorem-equivalence-sheaves-point}。
存在函子性双射
$$
\Gamma(S, \mathcal{F}) = (\mathcal{F}_{\overline{s}})^G
$$
\end{lemma}

\begin{proof}
可以使用形式论证和定理 \ref{etale-cohomology-theorem-equivalence-sheaves-point}
的结果来证明。给定层 $\mathcal{F}$，它对应于
$G$-集合 $M = \mathcal{F}_{\overline{s}}$，有
\begin{eqnarray*}
\Gamma(S, \mathcal{F}) & = &
\Mor_{\Sh(S_\etale)}(h_{\Spec(K)},  \mathcal{F})
\\
& = & \Mor_{G\textit{-Sets}}(\{*\}, M) \\
& = & M^G
\end{eqnarray*}
这里第一个等同由 Sites，第 \ref{sites-section-presheaves} 节和
第 \ref{sites-section-representable-sheaves} 节解释；
第二个由定理 \ref{etale-cohomology-theorem-equivalence-sheaves-point} 得到；
第三个显然。还给出一个直接证明\footnote{献给仍心存疑虑的读者。}。

\medskip\noindent
设 $t \in \Gamma(S, \mathcal{F})$ 是整体截面。
则三元组 $(S, \overline{s}, t)$ 定义
$\mathcal{F}_{\overline{s}}$ 的一个元素；它显然在 $G$ 的作用下不变。
反之，设 $(U, \overline{u}, t)$ 定义
$\mathcal{F}_{\overline{s}}$ 的一个不变元素。
可以缩小 $U$，并假设 $U = \Spec(L)$，其中
$K$ 有一个有限可分域扩张；见命题 \ref{etale-cohomology-proposition-etale-morphisms}。
在这种情形，映射 $\mathcal{F}(U) \to \mathcal{F}_{\overline{s}}$
是单射，因为对 \'etale 邻域的任意态射
$(U', \overline{u}') \to (U, \overline{u})$，限制映射
$\mathcal{F}(U) \to \mathcal{F}(U')$ 是单射；这是因为
$U' \to U$ 是 $S_\etale$ 的覆盖。
稍微扩大 $L$ 后，可以假设 $K \subset L$ 是有限 Galois 扩张。
此时使用
$$
\Spec(L) \times_{\Spec(K)} \Spec(L)
=
\coprod\nolimits_{\sigma \in \text{Gal}(L/K)} \Spec(L)
$$
，其中映射 $\Spec(L) \to \Spec(L \otimes_K L)$
来自环映射 $a \otimes b \mapsto a\sigma(b)$。因此可见，
$(U, \overline{u}, t)$ 在整个 $G$ 下不变这一条件说明
$t \in \mathcal{F}(\Spec(L))$ 通过任一投影作限制后，
都映到
$\mathcal{F}(\Spec(L) \times_{\Spec(K)} \Spec(L))$
的同一个元素（这里使用了上述单射性；略去细节）。
于是，层 $\mathcal{F}$ 关于 \'etale 覆盖
$\{\Spec(L) \to \Spec(K)\}$ 的层条件发挥作用，得到
$t$ 来自 $\mathcal{F}$ 在 $\Spec(K)$ 上的唯一截面。
\end{proof}

\begin{remark}
\label{etale-cohomology-remark-stalk-pullback}
设 $S$ 是概形，并设 $\overline{s} : \Spec(k) \to S$
是 $S$ 的几何点。按定义，这意味着 $k$ 代数闭。
特别地，$k$ 的绝对 Galois 群平凡。因此，由定理
\ref{etale-cohomology-theorem-equivalence-sheaves-point}，
$\Spec(k)_\etale$ 上的层范畴等价于集合范畴。
该等价由取 $\Spec(k)$ 上的截面给出。
这最终给出了茎函子的另一种定义。具体而言，函子
$$
\Sh(S_\etale) \longrightarrow \textit{Sets}, \quad
\mathcal{F} \longmapsto \mathcal{F}_{\overline{s}}
$$
同构于函子
$$
\Sh(S_\etale)
\longrightarrow
\Sh(\Spec(k)_\etale) = \textit{Sets},
\quad
\mathcal{F} \longmapsto \overline{s}^*\mathcal{F}
$$
要严格证明这一点，可以使用引理 \ref{etale-cohomology-lemma-stalk-pullback} 的 (3)，
其中 $f = \overline{s}$。此外，由此，引理
\ref{etale-cohomology-lemma-stalk-pullback} 的 (3) 的一般情形
也可从拉回的函子性得到。
\end{remark}






\section{群上同调}
\label{etale-cohomology-section-group-cohomology}

\noindent
下文若写 $H^i(G, M)$，其含义是：$G$ 是拓扑群，
$M$ 是具有连续 $G$-作用的离散 $G$-模，而 $H^i(G, -)$ 是
第 $i$ 个右导出函子，定义在范畴 $\text{Mod}_G$ 上；
该范畴由这样的 $G$-模组成；
见定义 \ref{etale-cohomology-definition-G-module-continuous} 和
\ref{etale-cohomology-definition-galois-cohomology}。这包括抽象群 $G$ 的情形，
其含义就是把 $G$ 看成具有离散拓扑的拓扑群。

\medskip\noindent
当模具有非离散拓扑时，将使用记号
$H^i_{cont}(G, M)$ 表示 \cite{Tate} 中引入的连续上同调群；
见第 \ref{etale-cohomology-section-continuous-group-cohomology} 节。

\begin{definition}
\label{etale-cohomology-definition-G-module-continuous}
设 $G$ 是拓扑群。
\begin{enumerate}
\item {\it $G$-模}（有时称为{\it 离散 $G$-模}）是 Abel 群
$M$，它带有由群同态给出的左作用 $a : G \times M \to M$，
并且 $a$ 连续（此时给 $M$ 配备离散拓扑）。
\item {\it $G$-模的态射} $f : M \to N$ 是
$G$-等变同态，其源为 $M$、靶为 $N$。
\item $G$-模的范畴记作 $\text{Mod}_G$。
\end{enumerate}
设 $R$ 是环。
\begin{enumerate}
\item {\it $R\text{-}G$-模}是 $R$-模 $M$，它带有左作用
$a : G \times M \to M$，该作用由 $R$-线性映射给出；并且
$a$ 连续（此时给 $M$ 配备离散拓扑）。
\item {\it $R\text{-}G$-模的态射} $f : M \to N$ 是
$G$-等变 $R$-模映射，其源为 $M$、靶为 $N$。
\item $R\text{-}G$-模的范畴记作 $\text{Mod}_{R, G}$。
\end{enumerate}
\end{definition}

\noindent
条件 $a : G \times M \to M$ 连续，等价于任意 $x \in M$
的稳定子在 $G$ 中开。若 $G$ 是抽象群，则在给 $G$ 配备离散拓扑后，
这就对应于带 $G$-作用的 Abel 群这一概念。注意
$\text{Mod}_{\mathbf{Z}, G} = \text{Mod}_G$。

\medskip\noindent
范畴 $\text{Mod}_G$ 有足够多内射对象；见 Injectives，引理
\ref{injectives-lemma-G-modules}。考虑左正合函子
$$
\text{Mod}_G \longrightarrow \textit{Ab},
\quad
M \longmapsto M^G =
\{x \in M \mid g \cdot x = x\ \forall g \in G\}
$$
有时记 $M^G = H^0(G, M)$，有时写作
$M^G = \Gamma_G(M)$。该函子有全右导出函子
$R\Gamma_G(M)$，并且其第 $i$ 个右导出函子为
$R^i\Gamma_G(M) = H^i(G, M)$，对任意 $i \geq 0$ 成立。

\medskip\noindent
同一构造也适用于
$H^0(G, -) : \text{Mod}_{R, G} \to \text{Mod}_R$。
引理 \ref{etale-cohomology-lemma-modules-abelian} 将说明，这与底层 $G$-模的上同调一致。

\begin{definition}
\label{etale-cohomology-definition-galois-cohomology}
设 $G$ 是拓扑群。设 $M$ 是具有连续 $G$-作用的离散 $G$-模。
换言之，$M$ 是定义 \ref{etale-cohomology-definition-G-module-continuous}
中引入的范畴 $\text{Mod}_G$ 的对象。
\begin{enumerate}
\item 右导出函子 $H^i(G, M)$（它导自 $H^0(G, M)$，并定义在范畴
$\text{Mod}_G$ 上）称为 $M$ 的{\it 连续群上同调群}。
\item 若 $G$ 是配备离散拓扑的抽象群，则
$H^i(G, M)$ 称为 $M$ 的{\it 群上同调群}。
\item 若 $G$ 是 Galois 群，则群 $H^i(G, M)$ 称为
$M$ 的{\it Galois 上同调群}。
\item 若 $G$ 是域 $K$ 的绝对 Galois 群，则群
$H^i(G, M)$ 有时称为{\it $K$ 的、以 $M$ 为系数的 Galois 上同调群}。
在此情形，有时写
$H^i(K, M)$ 来代替 $H^i(G, M)$。
\end{enumerate}
\end{definition}

\begin{lemma}
\label{etale-cohomology-lemma-modules-abelian}
设 $G$ 是拓扑群。设 $R$ 是环。
对每个 $i \geq 0$，图
$$
\xymatrix{
\text{Mod}_{R, G} \ar[rr]_{H^i(G, -)} \ar[d] & &
\text{Mod}_R \ar[d] \\
\text{Mod}_G \ar[rr]^{H^i(G, -)} & &
\textit{Ab}
}
$$
交换，其中竖直箭头是遗忘函子。
\end{lemma}

\begin{proof}
以 $F : \text{Mod}_{R, G} \to \text{Mod}_G$ 表示遗忘函子。
则 $F$ 有左伴随
$H : \text{Mod}_G \to \text{Mod}_{R, G}$，由
$H(M) = M \otimes_\mathbf{Z} R$ 给出。注意，
$\text{Mod}_G$ 的每个对象都是形如
$\mathbf{Z}[G/U]$ 的模的直和之商，其中 $U \subset G$ 是开子群。
这里 $\mathbf{Z}[G/U]$ 表示 $G$-模；它由有限
$\mathbf{Z}$-线性组合组成，所组合的是右 $U$ 同余类（位于 $G$ 中），
并带有左 $G$-作用。
因此，$\text{Mod}_G$ 中每个向上有界复形都拟同构于
$\text{Mod}_G$ 中的一个向上有界复形，后者的底层各项是平坦
$\mathbf{Z}$-模
（Derived Categories，引理 \ref{derived-lemma-subcategory-left-resolution}）。
于是显然，$LH$ 存在于 $D^-(\text{Mod}_G)$ 上，并可在任意复形上逐项
计算 $H$，只要其各项是平坦 $\mathbf{Z}$-模；这由 Derived Categories，引理
\ref{derived-lemma-subcategory-left-acyclics} 和命题
\ref{derived-proposition-enough-acyclics} 得到。
由 Derived Categories，引理
\ref{derived-lemma-pre-derived-adjoint-functors} 可知
$$
\text{Ext}^i(\mathbf{Z}, F(M)) = \text{Ext}^i(R, M)
$$
对 $M$（它属于 $\text{Mod}_{R, G}$）成立。
注意，有
$H^0(G, -) = \Hom(\mathbf{Z}, -)$ 于 $\text{Mod}_G$ 上，其中
$\mathbf{Z}$ 表示具有平凡作用的 $G$-模。因此，有
$H^i(G, -) = \text{Ext}^i(\mathbf{Z}, -)$ 于 $\text{Mod}_G$ 上。
类似地，有 $H^i(G, -) = \text{Ext}^i(R, -)$ 于
$\text{Mod}_{R, G}$ 上。综合以上即得引理。
\end{proof}

\begin{lemma}
\label{etale-cohomology-lemma-ext-modules-hom}
设 $G$ 是拓扑群。设 $R$ 是环。
设 $M$、$N$ 是 $R\text{-}G$-模。若 $M$ 作为 $R$-模有限投射，则
$\text{Ext}^i(M, N) = H^i(G, M^\vee \otimes_R N)$
（记号见证明）。
\end{lemma}

\begin{proof}
模 $M^\vee = \Hom_R(M, R)$ 带有 $G$ 的逆步作用。
具体而言，有 $(g \cdot \lambda)(m) = \lambda(g^{-1} \cdot m)$，
其中 $g \in G$、$\lambda \in M^\vee$、$m \in M$。
$G$ 在 $M^\vee \otimes_R N$ 上的作用是对角作用，即
$g \cdot (\lambda \otimes n) = g \cdot \lambda \otimes g \cdot n$。
注意，对第三个 $R\text{-}G$-模 $E$，有
$\Hom(E, M^\vee \otimes_R N) = \Hom(M \otimes_R E, N)$。
事实上，由 Algebra，引理 \ref{algebra-lemma-hom-from-tensor-product}
和 \ref{algebra-lemma-evaluation-map-iso-finite-projective}，
这在 $R$-模层面成立；而 $G$-作用的定义正是为了使它对
$R\text{-}G$-模仍然成立。因此，$M^\vee \otimes_R N$
是内射 $R\text{-}G$-模，只要 $N$ 是内射 $R\text{-}G$-模。
所以，若 $N \to N^\bullet$ 是内射分解，则
$M^\vee \otimes_R N \to M^\vee \otimes_R N^\bullet$
也是内射分解。于是
$$
\Hom(M, N^\bullet) = \Hom(R, M^\vee \otimes_R N^\bullet) =
(M^\vee \otimes_R N^\bullet)^G
$$
由于左边计算 $\text{Ext}^i(M, N)$，右边计算
$H^i(G, M^\vee \otimes_R N)$，证明完成。
\end{proof}

\begin{lemma}
\label{etale-cohomology-lemma-finite-dim-group-cohomology}
设 $G$ 是拓扑群。设 $k$ 是域。
设 $V$ 是 $k\text{-}G$-模。
若 $G$ 拓扑有限生成且
$\dim_k(V) < \infty$，则 $\dim_k H^1(G, V) < \infty$。
\end{lemma}

\begin{proof}
设 $g_1, \ldots, g_r \in G$ 是拓扑生成 $G$ 的元素；
也就是说，由 $g_1, \ldots, g_r$ 生成的子群稠密。
由引理 \ref{etale-cohomology-lemma-ext-modules-hom} 可见，
$H^1(G, V)$ 是 $k$-向量空间，其元素是扩张
$$
0 \to V \to E \to k \to 0
$$
；这些扩张属于 $k\text{-}G$-模。选取 $e \in E$，它映到 $1 \in k$。
写成
$$
g_i \cdot e = v_i + e
$$
，其中某个 $v_i \in V$。这是可能的，因为 $g_i \cdot 1 = 1$。
断言元素列表 $v_1, \ldots, v_r \in V$ 决定扩张 $E$ 的同构类。
一旦证明这一点，引理即得，因为这意味着我们的 Ext 向量空间同构于
$k$-向量空间 $V^{\oplus r}$ 的一个子商；略去一些细节。
由于 $E$ 是定义 \ref{etale-cohomology-definition-G-module-continuous} 中定义的范畴的对象，
知道存在开子群 $U$，使得 $u \cdot e = e$
对所有 $u \in U$ 成立。
现任取 $g \in G$。则 $gU$ 包含一个字 $w$，它由元素
$g_1, \ldots, g_r$ 组成。
设 $gu = w$。由于元素 $w \cdot e$ 由
$v_1, \ldots, v_r$ 决定，可见
$g \cdot e = (gu) \cdot e = w \cdot e$ 也由其决定。
\end{proof}

\begin{lemma}
\label{etale-cohomology-lemma-profinite-group-cohomology-torsion}
设 $G$ 是预有限拓扑群。则
\begin{enumerate}
\item $H^i(G, M)$ 是挠群，只要 $i > 0$，且任取
$G$-模 $M$；并且
\item $H^i(G, M) = 0$，若 $M$ 是 $\mathbf{Q}$-向量空间。
\end{enumerate}
\end{lemma}

\begin{proof}
证明 (1)。由维数移动，只需证明 $H^1(G, M)$ 是挠群，
其中任取 $G$-模 $M$。
选取正合序列 $0 \to M \to I \to N \to 0$，其中 $I$
是 $G$-模范畴的内射对象。
则 $H^1(G, M)$ 的任意元素都是某个元素
$y \in N^G$ 的像。选取 $x \in I$，它映到 $y$。
稳定子 $U \subset G$（关于 $x$）是开的，因此具有有限指数 $r$。
设 $g_1, \ldots, g_r \in G$ 是 $G/U$ 的一组代表元。
于是 $\sum g_i(x)$ 是 $I$ 的不变元素，映到 $ry$。
因此 $r$ 消去最初选取的 $H^1(G, M)$ 的元素。
(2) 得证，因为此时 $H^i(G, M)$ 既是 $\mathbf{Q}$-向量空间又是挠群。
\end{proof}





\section{Tate 连续上同调}
\label{etale-cohomology-section-continuous-group-cohomology}

\noindent
Tate 连续上同调（\cite{Tate}）由连续非齐次上链复形定义。
当 $M$ 是带有连续 $G$-作用的任意拓扑 Abel 群时，可以作此定义。
具体而言，考虑复形
$$
C^\bullet_{cont}(G, M) :
M \to \text{Maps}_{cont}(G, M) \to
\text{Maps}_{cont}(G \times G, M) \to \ldots
$$
其中对 $n \geq 1$，边界映射由下式定义：
\begin{align*}
\text{d}(f)(g_1, \ldots, g_{n + 1})
& = g_1(f(g_2, \ldots, g_{n + 1})) \\
&
+ \sum\nolimits_{j = 1, \ldots, n}
(-1)^jf(g_1, \ldots, g_jg_{j + 1}, \ldots, g_{n + 1}) \\
&
+ (-1)^{n + 1}f(g_1, \ldots, g_n)
\end{align*}
而当 $n = 0$ 时，它把 $m \in M$ 映到映射
$g \mapsto g(m) - m$。定义
$$
H^i_{cont}(G, M) = H^i(C^\bullet_{cont}(G, M))
$$
由于该复形的各项涉及从 $G$ 以及 $G$ 的自身乘积到拓扑模
$M$ 的连续映射，并不清楚它是否把拓扑模的短正合序列
变为上同调长正合序列。另一个困难是拓扑 Abel 群范畴并非 Abel 范畴！

\medskip\noindent
不过，离散 $G$-模的短正合序列确实给出连续上链复形的短正合序列，
从而给出连续上同调群 $H^i_{cont}(G, -)$ 的上同调长正合序列。
因此，在定义 \ref{etale-cohomology-definition-G-module-continuous} 的范畴
$\text{Mod}_G$ 上，函子 $H^i_{cont}(G, M)$ 构成上同调
$\delta$-函子，其定义见 Homology，第
\ref{homology-section-cohomological-delta-functor} 节。
由于定义 \ref{etale-cohomology-definition-galois-cohomology} 的上同调
$H^i(G, M)$ 是泛 $\delta$-函子
（Derived Categories，引理 \ref{derived-lemma-right-derived-delta-functor}），
得到典范映射
$$
H^i(G, M) \longrightarrow H^i_{cont}(G, M)
$$
，其中 $M \in \text{Mod}_G$。已知当 $G$ 是抽象群
（即 $G$ 具有离散拓扑），或者当 $G$ 是预有限群时，
这些映射是同构（此处插入未来的参考文献）。
若你知道某个例子，说明对拓扑群 $G$ 和
$M \in \Ob(\text{Mod}_G)$，此映射不是同构，请发送电子邮件至
\href{mailto:stacks.project@gmail.com}{stacks.project@gmail.com}。




\section{点的上同调}
\label{etale-cohomology-section-cohomology-point}

\noindent
作为前几节讨论的推论，得到域的谱的 \'etale 上同调与 Galois 上同调的等价。

\begin{lemma}
\label{etale-cohomology-lemma-equivalence-abelian-sheaves-point}
设 $S = \Spec(K)$，其中 $K$ 是域。
设 $\overline{s}$ 是 $S$ 的几何点。
以 $G = \text{Gal}_{\kappa(s)}$ 表示绝对 Galois 群。
茎函子诱导范畴等价
$$
\textit{Ab}(S_\etale) \longrightarrow \text{Mod}_G,
\quad
\mathcal{F} \longmapsto \mathcal{F}_{\overline{s}}.
$$
\end{lemma}

\begin{proof}
在定理 \ref{etale-cohomology-theorem-equivalence-sheaves-point} 中，
已经看到集合层与 $G$-集合之间的等价。
本引理由此形式地得到，因为 Abel 层无非是带有交换群律的集合层，
而 $G$-模无非是带有交换群律的 $G$-集合。
\end{proof}

\begin{lemma}
\label{etale-cohomology-lemma-compare-cohomology-point}
记号与假设同引理 \ref{etale-cohomology-lemma-equivalence-abelian-sheaves-point}。
设 $\mathcal{F}$ 是 $\Spec(K)_\etale$ 上的 Abel 层，
对应于 $G$-模 $M$。则
\begin{enumerate}
\item 在 $D(\textit{Ab})$ 中有典范同构
$R\Gamma(S, \mathcal{F}) = R\Gamma_G(M)$；
\item $H_\etale^0(S, \mathcal{F}) = M^G$；并且
\item $H_\etale^q(S, \mathcal{F}) = H^q(G, M)$。
\end{enumerate}
\end{lemma}

\begin{proof}
结合引理 \ref{etale-cohomology-lemma-equivalence-abelian-sheaves-point}
与引理 \ref{etale-cohomology-lemma-global-sections-point}。
\end{proof}

\begin{example}
\label{etale-cohomology-example-sheaves-point}
$\Spec(K)_\etale$ 上的层。
设 $G = \text{Gal}(K^{sep}/K)$ 是 $K$ 的绝对 Galois 群。
\begin{enumerate}
\item 常值层 $\underline{\mathbf{Z}/n\mathbf{Z}}$ 对应于模
$\mathbf{Z}/n\mathbf{Z}$，其 $G$-作用平凡；
\item 层 $\mathbf{G}_m|_{\Spec(K)_\etale}$ 对应于
$(K^{sep})^*$，带有其 $G$-作用；
\item 层 $\mathbf{G}_a|_{\Spec(K^{sep})}$ 对应于
$(K^{sep}, +)$，带有其 $G$-作用；并且
\item 层 $\mu_n|_{\Spec(K^{sep})}$ 对应于
$\mu_n(K^{sep})$，带有其 $G$-作用。
\end{enumerate}
由注 \ref{etale-cohomology-remark-special-case-fpqc-cohomology-quasi-coherent}
和定理 \ref{etale-cohomology-theorem-picard-group}，有如下上同调群等同：
\begin{align*}
H_\etale^0(S_\etale, \mathbf{G}_m) & =
\Gamma(S, \mathcal{O}_S^*) \\
H_\etale^1(S_\etale, \mathbf{G}_m) & =
H_{Zar}^1(S, \mathcal{O}_S^*) = \Pic(S) \\
H_\etale^i(S_\etale, \mathbf{G}_a) & =
H_{Zar}^i(S, \mathcal{O}_S)
\end{align*}
此外，对任意拟凝聚层 $\mathcal{F}$（定义在 $S_\etale$ 上），有
$$
H^i(S_\etale, \mathcal{F}) = H_{Zar}^i(S, \mathcal{F}),
$$
见定理 \ref{etale-cohomology-theorem-zariski-fpqc-quasi-coherent}。
特别地，这给出等式序列
$$
0 =
\Pic(\Spec(K)) =
H_\etale^1(\Spec(K)_\etale, \mathbf{G}_m) =
H^1(G, (K^{sep})^*)
$$
，它正是 Hilbert 定理 90。类似地，对 $i \geq 1$，
$$
0 = H^i(\Spec(K), \mathcal{O})
= H_\etale^i(\Spec(K)_\etale, \mathbf{G}_a)
= H^i(G, K^{sep})
$$
其中 $K^{sep}$ 表示作为 Galois 模的 $K^{sep}$，其群律为加法。
这样，可以把迄今所做的工作看作计算 Galois 上同调群的一种复杂方法。
\end{example}

\noindent
下一个结果是一则趣闻，初读时应跳过。

\begin{lemma}
\label{etale-cohomology-lemma-all-modules-quasi-coherent}
设 $R$ 是维数 $0$ 的局部环。设 $S = \Spec(R)$。
则每个 $\mathcal{O}_S$-模（定义在 $S_\etale$ 上）都是拟凝聚的。
\end{lemma}

\begin{proof}
设 $\mathcal{F}$ 是 $\mathcal{O}_S$-模，定义在 $S_\etale$ 上。
必须证明 $\mathcal{F}$ 由 $R$-模
$M = \Gamma(S, \mathcal{F})$ 决定。
更精确地，若 $\pi : X \to S$ 是 \'etale 的，必须证明
$\Gamma(X, \mathcal{F}) = \Gamma(X, \pi^*\widetilde{M})$。

\medskip\noindent
设 $\mathfrak m \subset R$ 是极大理想，并设
$\kappa$ 是剩余域。由 Algebra，引理
\ref{algebra-lemma-local-dimension-zero-henselian}，
局部环 $R$ 是 Hensel 的。若 $X \to S$ 是 \'etale 的，
则由 Morphisms，引理 \ref{morphisms-lemma-etale-over-field}，
$X$ 的底层拓扑空间离散，因而 $X$ 是仿射概形的不交并，
且每个分支只有一个点。此外，若 $X = \Spec(A)$ 仿射且只有一个点，
则由 Algebra，引理 \ref{algebra-lemma-mop-up}，
$R \to A$ 是有限 \'etale 的。
在此情形，必须证明 $\Gamma(X, \mathcal{F}) = M \otimes_R A$。

\medskip\noindent
函子 $A \mapsto A/\mathfrak m A$ 定义如下范畴等价：
有限 \'etale $R$-代数的范畴等价于有限可分 $\kappa$-代数的范畴；
见 Algebra，引理 \ref{algebra-lemma-henselian-cat-finite-etale}。
先考虑 $A/\mathfrak m A$ 是 $\kappa$ 的 Galois 扩张且
Galois 群为 $G$ 的情形。
对每个 $\sigma \in G$，以 $\sigma : A \to A$ 表示
$A$ 在 $R$ 上的相应自同构。
置 $N = \Gamma(X, \mathcal{F})$。
于是 $\Spec(\sigma) : X \to X$ 是 $S$ 上的自同构，
故由其拉回定义映射 $\sigma : N \to N$；
这是 $\sigma$-线性映射：
$\sigma(an) = \sigma(a) \sigma(n)$，其中 $a \in A$ 且 $n \in N$。
将 Galois 下降应用于拟凝聚模
$\widetilde{N}$（定义在 $X$ 上），它配备来自 $\sigma$ 在 $N$ 上作用的同构。
见 Descent，引理 \ref{descent-lemma-galois-descent-more-general}。
该引理说明存在同构 $N = N^G \otimes_R A$。
另一方面，显然 $N^G = M$；这是 $\mathcal{F}$ 的层条件。
因此所需同构成立。

\medskip\noindent
一般情形（$A$ 局部且在 $R$ 上有限 \'etale）如下归结为
Galois 情形。选取有限 \'etale 的 $A \to B$，使得 $B$ 局部，
其剩余域在 $\kappa$ 上 Galois。置
$G = \text{Aut}(B/R) = \text{Gal}(\kappa_B/\kappa)$。
设 $H \subset G$ 是对应于 Galois 扩张
$\kappa_B/\kappa_A$ 的 Galois 群。于是如上可证
$\Gamma(X, \mathcal{F}) = \Gamma(\Spec(B), \mathcal{F})^H$。
将 Galois 扩张的结果使用两次，得到
$$
\Gamma(X, \mathcal{F}) = (M \otimes_R B)^H = M \otimes_R A
$$
正合所求。
\end{proof}








\section{曲线的上同调}
\label{etale-cohomology-section-cohomology-curves}

\noindent
眼下的下一个任务是计算代数闭域上光滑曲线取挠系数的 \'etale 上同调，
特别要证明它在次数至少为 3 时消失。为此，将计算一般点处的上同调，
这归结为某些 Galois 上同调。





\section{布劳尔群}
\label{etale-cohomology-section-brauer-groups}

\noindent
域的布劳尔群用有限维中心单代数定义。
本节回顾布劳尔群的相关事实，其中大部分在 Brauer Groups，第
\ref{brauer-section-introduction} 节中讨论。
其他参考文献见 \cite{SerreCorpsLocaux}、
\cite{SerreGaloisCohomology} 或 \cite{Weil}。

\begin{theorem}
\label{etale-cohomology-theorem-central-simple-algebra}
设 $K$ 是域。对含幺、结合（不一定交换）的 $K$-代数 $A$，
以下条件等价：
\begin{enumerate}
\item $A$ 是有限维中心单 $K$-代数；
\item $A$ 是有限维 $K$-向量空间，$K$ 是 $A$ 的中心，并且
$A$ 没有非平凡双边理想；
\item 存在 $d \geq 1$，使得
$A \otimes_K \bar K \cong \text{Mat}(d \times d, \bar K)$；
\item 存在 $d \geq 1$，使得
$A \otimes_K K^{sep} \cong \text{Mat}(d \times d, K^{sep})$；
\item 存在 $d \geq 1$ 以及有限 Galois 扩张 $K'/K$，使得
$A \otimes_K K' \cong \text{Mat}(d \times d, K')$；
\item 存在 $n \geq 1$ 以及有限维中心斜域 $D$（定义在 $K$ 上），使得
$A \cong \text{Mat}(n \times n, D)$。
\end{enumerate}
整数 $d$ 称为 $A$ 的{\it 次数}。
\end{theorem}

\begin{proof}
这是 Brauer Groups，引理
\ref{brauer-lemma-finite-central-simple-algebra} 的复述。
\end{proof}

\begin{lemma}
\label{etale-cohomology-lemma-brauer-inverse}
设 $A$ 是 $K$ 上的有限维中心单代数。则
$$
\begin{matrix}
A \otimes_K A^{opp} & \longrightarrow & \text{End}_K(A) \\
\ a \otimes a' & \longmapsto & (x \mapsto a x a')
\end{matrix}
$$
是 $K$ 上代数的同构。
\end{lemma}

\begin{proof}
见 Brauer Groups，引理 \ref{brauer-lemma-inverse}。
\end{proof}

\begin{definition}
\label{etale-cohomology-definition-brauer-equivalent}
若两个有限维中心单代数 $A_1$ 与 $A_2$（均定义在 $K$ 上）满足：
存在 $m, n \geq 1$，使得
$\text{Mat}(n \times n, A_1)
\cong \text{Mat}(m \times m, A_2)$，
则称它们{\it 相似}或{\it 等价}。记作 $A_1 \sim A_2$。
\end{definition}

\noindent
由 Brauer Groups，引理 \ref{brauer-lemma-similar}，这是等价关系。

\begin{definition}
\label{etale-cohomology-definition-brauer-group}
设 $K$ 是域。$K$ 的{\it 布劳尔群}是相似类集合
$\text{Br} (K)$，其中各类来自 $K$ 上的有限维中心单代数，
其群律由张量积（在 $K$ 上）诱导。
$A$ 在 $\text{Br}(K)$ 中的类记作 $[A]$。单位元为
$[K] = [\text{Mat}(d \times d, K)]$，其中任意 $d \geq 1$。
\end{definition}

\noindent
前一引理说明逆元存在，且 $-[A] = [A^{opp}]$。
域的布劳尔群总是挠群。
事实上，将看到 $[A]$ 的阶整除 $\deg(A)$，其中 $A$
是任意有限维中心单代数（见引理 \ref{etale-cohomology-lemma-annihilated-by-degree}）。
一般而言，布劳尔群不是有限生成的；例如，非 Archimedes 局部域的布劳尔群是
$\mathbf{Q}/\mathbf{Z}$。
$\mathbf{C}(x, y)$ 的布劳尔群不可数。

\begin{lemma}
\label{etale-cohomology-lemma-central-simple-algebra-pgln}
\begin{slogan}
中心单代数由 PGL 的 Galois 上同调分类。
\end{slogan}
设 $K$ 是域，并设 $K^{sep}$ 是可分代数闭包。
则次数为 $d$、定义在 $K$ 上的中心单代数之同构类集合，与非 Abel 上同调
$H^1(\text{Gal}(K^{sep}/K), \text{PGL}_d(K^{sep}))$
之间存在双射。
\end{lemma}

\begin{proof}[证明概要。]
Skolem--Noether 定理（见 Brauer Groups，定理
\ref{brauer-theorem-skolem-noether}）说明，对任意域 $L$，群
$\text{Aut}_{L\text{-Algebras}}(\text{Mat}_d(L))$
等于 $\text{PGL}_d(L)$。由定理 \ref{etale-cohomology-theorem-central-simple-algebra}，
可见次数为 $d$ 的中心单代数对应于 $K$-代数
$\text{Mat}_d(K)$ 的形式。
结合二者可见，次数为 $d$ 的中心单代数之同构类对应于
$H^1(\text{Gal}(K^{sep}/K), \text{PGL}_d(K^{sep}))$
的元素。关于扭曲的更多细节，例如见 \cite{SilvermanEllipticCurves}。
\end{proof}

\noindent
若 $A$ 是次数为 $d$ 的有限维中心单代数，定义在域 $K$ 上，
以 $\xi_A$ 表示相应的上同调类
$H^1(\text{Gal}(K^{sep}/K), \text{PGL}_d(K^{sep}))$。
考虑短正合序列
$$
1 \to (K^{sep})^* \to \text{GL}_d(K^{sep}) \to \text{PGL}_d(K^{sep}) \to 1,
$$
它给出上同调长正合序列（直到次数 2），其连接映射为
$$
\delta_d :
H ^1(\text{Gal}(K^{sep}/K), \text{PGL}_d(K^{sep}))
\longrightarrow
H^2(\text{Gal}(K^{sep}/K), (K^{sep})^*).
$$
具体如下：若 $\xi$ 是由 1-上闭链 $(g_\sigma)$ 表示的上同调类，
则 $\delta_d(\xi)$ 是 2-上闭链之类
\begin{equation}
\label{etale-cohomology-equation-two-cocycle}
(\sigma, \tau)
\longmapsto
\tilde g_\sigma^{-1} \tilde g_{\sigma \tau} \sigma(\tilde g_\tau^{-1})
\in (K^{sep})^*
\end{equation}
，其中 $\tilde g_\sigma \in \text{GL}_d(K^{sep})$ 是
$g_\sigma$ 的一个提升。
利用这一点，可以将映射
$$
\delta : \text{Br}(K) \longrightarrow H^2(\text{Gal}(K^{sep}/K), (K^{sep})^*),
\quad
[A] \longmapsto \delta_{\deg A} (\xi_A)
$$
具体化如下。假设 $A$ 次数为 $d$，且定义在 $K$ 上。选取同构
$\varphi : \text{Mat}_d(K^{sep}) \to A \otimes_K K^{sep}$。
对 $\sigma \in \text{Gal}(K^{sep}/K)$，选取元素
$\tilde g_\sigma \in \text{GL}_d(K^{sep})$，使得
$\varphi^{-1} \circ \sigma(\varphi)$ 等于映射
$x \mapsto \tilde g_\sigma x \tilde g_\sigma^{-1}$。
$H^2$ 中的类由二上闭链（\ref{etale-cohomology-equation-two-cocycle}）定义。

\begin{theorem}
\label{etale-cohomology-theorem-brauer-delta}
设 $K$ 是具有可分代数闭包 $K^{sep}$ 的域。上面定义的映射
$\delta : \text{Br}(K) \to H^2(\text{Gal}(K^{sep}/K), (K^{sep})^*)$
是群同构。
\end{theorem}

\begin{proof}[证明概要]
为证明 $\delta$ 定义群同态，即
$\delta(A \otimes_K B) = \delta(A) + \delta(B)$，
直接用上闭链计算。

\medskip\noindent
证明 $\delta$ 的单射性。在 Abel 情形（$d = 1$）中，有等同
$$
H^1(\text{Gal}(K^{sep}/K), \text{GL}_d(K^{sep})) =
H_\etale^1(\Spec(K), \text{GL}_d(\mathcal{O}))
$$
后者由 fpqc 下降是平凡的。若这在非 Abel 情形也成立，
便立即推出 $\delta$ 的单射性。（见 \cite{SGA4.5}。）
更确切地说，为证明这一点，可以把 $\delta([A])$ 重新解释为
如下存在性的障碍：存在 $K$-向量空间 $V$，它带左 $A$-模结构，
并满足 $\dim_K V = \deg A$。当 $V$ 存在时，有
$A \cong \text{End}_K(V)$。

\medskip\noindent
为证明满射性，选取上同调类
$\xi \in H^2(\text{Gal}(K^{sep}/K), (K^{sep})^*)$，
则存在有限 Galois 扩张 $K^{sep}/K'/K$，使得 $\xi$ 是某个
$\xi' \in H^2(\text{Gal}(K'|K), (K')^*)$ 的像。
然后写出 $K$ 上的一个显式中心单代数，使用数据 $K', \xi'$。
\end{proof}










\section{概形的布劳尔群}
\label{etale-cohomology-section-brauer-scheme}

\noindent
设 $S$ 为概形。若一个 $\mathcal{O}_S$-代数
$\mathcal{A}$ 在 \'etale 局部是矩阵代数，则称其为 {\it Azumaya 代数}；
换言之，若存在一个 \'etale 覆盖
$\mathcal{U} = \{ \varphi_i : U_i \to S\}_{i \in I}$，使得
$\varphi_i^*\mathcal{A} \cong \text{Mat}_{d_i}(\mathcal{O}_{U_i})$
对某些 $d_i \geq 1$ 成立。对于两个这样的
$\mathcal{A}$ 与 $\mathcal{B}$，若存在有限局部自由
$\mathcal{O}_S$-模 $\mathcal{F}$ 与 $\mathcal{G}$，
它们在每个 $s \in S$ 处的秩都为正，并且
$$
\mathcal{A} \otimes_{\mathcal{O}_S}
\SheafHom_{\mathcal{O}_S}(\mathcal{F}, \mathcal{F})
\cong
\mathcal{B} \otimes_{\mathcal{O}_S}
\SheafHom_{\mathcal{O}_S}(\mathcal{G}, \mathcal{G})
$$
作为 $\mathcal{O}_S$-代数同构，则称它们{\it 等价}。$S$ 的{\it 布劳尔群}
是集合 $\text{Br}(S)$，其元素为 Azumaya
$\mathcal{O}_S$-代数的等价类，
其运算由（在 $\mathcal{O}_S$ 上的）张量积诱导。

\begin{lemma}
\label{etale-cohomology-lemma-end-unique-up-to-invertible}
设 $S$ 为概形。设 $\mathcal{F}$ 与 $\mathcal{G}$ 为正秩的有限局部自由
$\mathcal{O}_S$-模层。若存在同构
$\SheafHom_{\mathcal{O}_S}(\mathcal{F}, \mathcal{F}) \cong
\SheafHom_{\mathcal{O}_S}(\mathcal{G}, \mathcal{G})$，且它是
$\mathcal{O}_S$-代数同构，则存在可逆层
$\mathcal{L}$，它定义在 $S$ 上，并且
$\mathcal{F} \otimes_{\mathcal{O}_S} \mathcal{L} \cong \mathcal{G}$，
并且此同构诱导给定的自同态代数同构。
\end{lemma}

\begin{proof}
固定一个同构
$\SheafHom_{\mathcal{O}_S}(\mathcal{F}, \mathcal{F}) \to
\SheafHom_{\mathcal{O}_S}(\mathcal{G}, \mathcal{G})$。
考虑层 $\mathcal{L} \subset \SheafHom(\mathcal{F}, \mathcal{G})$：
作为 $\mathcal{O}_S$-模，它由所有局部同构
$\varphi : \mathcal{F} \to \mathcal{G}$ 生成，其中由
$\varphi$ 作共轭得到的正是给定的自同态代数同构。
局部计算（归结到 $\mathcal{F}$ 与 $\mathcal{G}$ 有限自由且
$S$ 仿射的情形）表明 $\mathcal{L}$ 可逆。另一项局部计算表明取值映射
$$
\mathcal{F} \otimes_{\mathcal{O}_S} \mathcal{L} \longrightarrow \mathcal{G}
$$
是同构。
\end{proof}

\noindent
下述引理证明中的论证可见
\cite{Saltman-torsion}。

\begin{lemma}
\label{etale-cohomology-lemma-annihilated-by-degree}
\begin{reference}
论证取自 \cite{Saltman-torsion}。
\end{reference}
设 $S$ 为概形。设 $\mathcal{A}$ 为 Azumaya 代数，并且是秩为
$d^2$ 的局部自由代数（相对于 $S$）。则
$\mathcal{A}$ 在 $S$ 的布劳尔群中的类被 $d$ 零化。
\end{lemma}

\begin{proof}
选取 \'etale 覆盖 $\{U_i \to S\}$，并选取同构
$\mathcal{A}|_{U_i} \to \SheafHom(\mathcal{F}_i, \mathcal{F}_i)$，
其中某个局部自由 $\mathcal{O}_{U_i}$-模 $\mathcal{F}_i$
秩为 $d$。（可以假设 $\mathcal{F}_i$ 自由。）考虑复合
$$
p_i : \mathcal{F}_i^{\otimes d} \to
\wedge^d(\mathcal{F}_i) \to \mathcal{F}_i^{\otimes d}
$$
第一个箭头是通常的投影；第二个箭头把 $\mathcal{F}_i$ 的最高外幂
同构到 $\mathcal{F}_i^{\otimes d}$ 的如下截面子模：其中的截面在
$d$ 个字母上的对称群作用下依符号特征变换。于是
$p_i^2 = d! p_i$，并且 $p_i$ 的秩为 $1$。利用给定同构
$\mathcal{A}|_{U_i} \to \SheafHom(\mathcal{F}_i, \mathcal{F}_i)$
以及典范同构
$$
\SheafHom(\mathcal{F}_i, \mathcal{F}_i)^{\otimes d} =
\SheafHom(\mathcal{F}_i^{\otimes d}, \mathcal{F}_i^{\otimes d})
$$
，可把 $p_i$ 看作 $\mathcal{A}^{\otimes d}$ 在
$U_i$ 上的截面。我们断言，有
$p_i|_{U_i \times_S U_j} = p_j|_{U_i \times_S U_j}$，
且这是作为 $\mathcal{A}^{\otimes d}$ 的截面的等式。事实上，应用
引理 \ref{etale-cohomology-lemma-end-unique-up-to-invertible}，得到可逆层
$\mathcal{L}_{ij}$ 以及典范同构
$$
\mathcal{F}_i|_{U_i \times_S U_j} \otimes \mathcal{L}_{ij}
\longrightarrow
\mathcal{F}_j|_{U_i \times_S U_j}.
$$
由此同构可见 $p_i$ 映到 $p_j$。由于 $\mathcal{A}^{\otimes d}$
是 $S_\etale$ 上的层（命题
\ref{etale-cohomology-proposition-quasi-coherent-sheaf-fpqc}），我们得到典范整体截面
$p \in \Gamma(S, \mathcal{A}^{\otimes d})$。局部计算表明
$$
\mathcal{H} =
\Im(\mathcal{A}^{\otimes d} \to \mathcal{A}^{\otimes d}, f \mapsto fp)
$$
是秩为 $d^d$ 的局部自由模，并且由 $\mathcal{A}^{\otimes d}$ 作（左）乘
诱导同构
$\mathcal{A}^{\otimes d} \to \SheafHom(\mathcal{H}, \mathcal{H})$。
换言之，$\mathcal{A}^{\otimes d}$ 是 $S$ 的布劳尔群中的平凡元素，
正如所需。
\end{proof}

\noindent
在此情形下，定理 \ref{etale-cohomology-theorem-brauer-delta} 中同构 $\delta$ 的类似物
是映射
$$
\delta_S: \text{Br}(S) \to H_\etale^2(S, \mathbf{G}_m).
$$
映射 $\delta_S$ 确实是单射。若 $S$ 拟紧或连通，则
$\text{Br}(S)$ 是挠群，因而在此情形下 $\delta_S$ 的像包含在 $S$
的{\it 上同调布劳尔群}
$$
\text{Br}'(S) := H_\etale^2(S, \mathbf{G}_m)_\text{torsion}.
$$
中。因此，若 $S$ 拟紧或连通，就有包含关系 $\text{Br}(S)
\subset \text{Br}'(S)$。它并非总是等式：存在非分离奇异曲面 $S$，
使得 $\text{Br}(S) \subset
\text{Br}'(S)$ 是严格包含。若 $S$ 拟射影，则
$\text{Br}(S) = \text{Br}'(S)$。不过，举例来说，
尚不知这一点对 $\mathbf{C}$ 上的光滑固有簇是否成立。



\section{Artin--Schreier 序列}
\label{etale-cohomology-section-artin-schreier}

\noindent
设 $p$ 为素数。设 $S$ 为特征 $p$ 的概形。
{\it Artin--Schreier} 序列是短正合列
$$
0 \longrightarrow \underline{\mathbf{Z}/p\mathbf{Z}}_S \longrightarrow
\mathbf{G}_{a, S} \xrightarrow{F-1} \mathbf{G}_{a, S} \longrightarrow 0
$$
，其中 $F - 1$ 是映射 $x \mapsto x^p - x$。

\begin{lemma}
\label{etale-cohomology-lemma-vanishing-affine-char-p-p}
设 $p$ 为素数。设 $S$ 为特征 $p$ 的概形。
\begin{enumerate}
\item 若 $S$ 仿射，则
$H_\etale^q(S, \underline{\mathbf{Z}/p\mathbf{Z}}) = 0$ 对所有
$q \geq 2$ 都成立。
\item 若 $S$ 是维数为 $d$ 的拟紧拟分离概形，则
$H_\etale^q(S, \underline{\mathbf{Z}/p\mathbf{Z}}) = 0$
对所有 $q \geq 2 + d$ 都成立。
\end{enumerate}
\end{lemma}

\begin{proof}
回忆结构层的 \'etale 上同调等于它在底拓扑空间上的上同调
（定理 \ref{etale-cohomology-theorem-zariski-fpqc-quasi-coherent}）。
第一条由 Artin--Schreier 正合列以及仿射概形上结构层上同调的消失得到
（《概形上同调》引理
\ref{coherent-lemma-quasi-coherent-affine-cohomology-zero}）。
第二条由同一论证、《上同调》命题
\ref{cohomology-proposition-cohomological-dimension-spectral}
中的消失，以及 $S$ 是谱空间这一事实得到
（《性质》引理
\ref{properties-lemma-quasi-compact-quasi-separated-spectral}）。
\end{proof}

\begin{lemma}
\label{etale-cohomology-lemma-F-1}
设 $k$ 是特征 $p > 0$ 的代数闭域。
设 $V$ 是有限维 $k$-向量空间。设 $F : V \to V$
为 Frobenius 线性映射，即加法映射，并且满足
$F(\lambda v) = \lambda^p F(v)$，其中任取 $\lambda \in k$
与 $v \in V$。则 $F - 1 : V \to V$ 是满射，其核是有限维
$\mathbf{F}_p$-向量空间，维数 $\leq \dim_k(V)$。
\end{lemma}

\begin{proof}
若 $F = 0$，则结论成立。若 $V$ 有一个由
$F$-稳定子向量空间组成的滤过，并且结论对每个分次片都成立，
则它对 $(V, F)$ 成立。结合这两点，可以假设 $F$ 的核为零。

\medskip\noindent
选取基 $v_1, \ldots, v_n$，它们构成 $V$ 的一组基，并写成
$F(v_i) = \sum a_{ij} v_j$。注意，$v = \sum \lambda_i v_i$
属于该核，当且仅当 $\sum \lambda_i^p a_{ij} v_j = 0$。
由于 $k$ 代数闭，这蕴含矩阵 $(a_{ij})$ 可逆。设
$(b_{ij})$ 为其逆矩阵。为了说明 $F - 1$ 为满射，选取
$w = \sum \mu_i v_i \in V$，并尝试求解
$$
(F - 1)(\sum \lambda_iv_i) =
\sum \lambda_i^p a_{ij} v_j - \sum \lambda_j v_j = \sum \mu_j v_j
$$
这等价于
$$
\sum \lambda_j^p v_j - \sum b_{ij} \lambda_i v_j = \sum b_{ij} \mu_i v_j
$$
，换言之，
$$
\lambda_j^p - \sum b_{ij} \lambda_i = \sum b_{ij} \mu_i,
\quad j = 1, \ldots, \dim(V).
$$
代数
$$
A = k[x_1, \ldots, x_n]/
(x_j^p - \sum b_{ij} x_i - \sum b_{ij} \mu_i)
$$
在 $k$ 上标准光滑（《代数》定义
\ref{algebra-definition-standard-smooth}），因为矩阵 $(b_{ij})$
可逆，而 $x_j^p$ 的偏导数都为零。$A$ 在 $k$ 上的一组基是所有单项式
$x_1^{e_1} \ldots x_n^{e_n}$，其中 $e_i < p$，故
$\dim_k(A) = p^n$。由于 $k$ 代数闭，可见 $\Spec(A)$ 恰有
$p^n$ 个点。因此 $F - 1$ 为满射，并且每条纤维都有
$p^n$ 个点；也就是说，$F - 1$ 的核是有 $p^n$ 个元素的群。
\end{proof}

\begin{lemma}
\label{etale-cohomology-lemma-F-2}
在引理 \ref{etale-cohomology-lemma-F-1} 的情形下，设 $k'/k$ 是代数闭域扩张。
令 $V' = V \otimes_k k'$，并令
$F' : V' \to V'$ 为 $F$ 诱导的 Frobenius 线性映射。
则 $\Ker(F - 1)$ 同构地映到 $\Ker(F' - 1)$ 上。
\end{lemma}

\begin{proof}
设 $v_1, \ldots, v_n$ 是 $V$ 的一组基，并设
$F(v_i) = \sum a_{ij}v_j$。则 $v = \sum \lambda_i v_i$
属于 $\Ker(F - 1)$，当且仅当有
$\sum \lambda_i^pa_{ij} = \lambda_j$，其中 $j = 1, \ldots, n$。
由引理 \ref{etale-cohomology-lemma-F-1}，
这些方程截出一个闭子概形 $Z$；它位于
$\mathbf{A}^n_k = \Spec(k[\lambda_1, \ldots, \lambda_n])$ 中，
并且在 $\Spec(k)$ 上有限。
于是 $Z(k) = Z(k')$，结论随之成立。
\end{proof}

\begin{lemma}
\label{etale-cohomology-lemma-top-cohomology-coherent}
设 $X$ 是域 $k$ 上分离的有限型概形。设 $\mathcal{F}$
为 $\mathcal{O}_X$-模的凝聚层。则
$\dim_k H^d(X, \mathcal{F}) < \infty$，其中 $d = \dim(X)$。
\end{lemma}

\begin{proof}
对 $d$ 作归纳。情形 $d = 0$ 成立，因为此时
$X$ 是有限维 $k$-代数 $A$ 的谱
（《簇》引理 \ref{varieties-lemma-algebraic-scheme-dim-0}），
而每个凝聚层 $\mathcal{F}$ 都对应于有限 $A$-模
$M = H^0(X, \mathcal{F})$，它满足 $\dim_k M < \infty$。

\medskip\noindent
假设 $d > 0$，且已对维数 $< d$ 的分离有限型概形证明结论。
概形 $X$ 是诺特的。考虑性质
$\mathcal{P}$；它定义在 $X$ 上的凝聚层上，规则为
$$
\mathcal{P}(\mathcal{F}) \Leftrightarrow
\dim_k H^d(X, \mathcal{F}) < \infty
$$
我们将使用《概形上同调》引理
\ref{coherent-lemma-property-initial} 的结果，证明
$\mathcal{P}$ 对 $X$ 上每个凝聚层都成立。

\medskip\noindent
设
$$
0 \to \mathcal{F}_1 \to \mathcal{F} \to \mathcal{F}_2 \to 0
$$
是 $X$ 上凝聚层的短正合列。考虑上同调长正合列中的一段
$$
H^d(X, \mathcal{F}_1) \to
H^d(X, \mathcal{F}) \to
H^d(X, \mathcal{F}_2)
$$
因此，若 $\mathcal{P}$ 对 $\mathcal{F}_1$ 与 $\mathcal{F}_2$
成立，则它对 $\mathcal{F}$ 也成立。

\medskip\noindent
设 $Z \subset X$ 为整闭子概形。设 $\mathcal{I}$
为 $Z$ 上的凝聚理想层。为完成证明，须说明
$H^d(X, i_*\mathcal{I}) = H^d(Z, \mathcal{I})$ 是有限维的。
若 $\dim(Z) < d$，则结论成立，因为该上同调群为零
（《上同调》命题
\ref{cohomology-proposition-vanishing-Noetherian}）。
这样便归结到下一段讨论的情形。

\medskip\noindent
假设 $X$ 是维数为 $d$ 的簇，并且
$\mathcal{F} = \mathcal{I}$ 是凝聚理想层。在此情形下有短正合列
$$
0 \to \mathcal{I} \to \mathcal{O}_X \to i_*\mathcal{O}_Z \to 0
$$
，其中 $i : Z \to X$ 是由 $\mathcal{I}$ 定义的闭子概形。
由归纳假设可见
$H^{d - 1}(Z, \mathcal{O}_Z) = H^{d - 1}(X, i_*\mathcal{O}_Z)$
是有限维的。因此，只需对结构层证明结论。

\medskip\noindent
可以把 Chow 引理
（《概形上同调》引理 \ref{coherent-lemma-chow-Noetherian}）
应用于态射 $X \to \Spec(k)$。由此得到 Chow 引理陈述中的图
$$
\xymatrix{
X \ar[rd]_g & X' \ar[d]^-{g'} \ar[l]^\pi \ar[r]_i & \mathbf{P}^n_k \ar[dl] \\
& \Spec(k) &
}
$$
再设 $U \subset X$ 是满足 $\pi^{-1}(U) \to U$ 为同构的稠密开子概形。
还可以假设 $X'$ 也是簇，见
《概形上同调》评注 \ref{coherent-remark-chow-Noetherian}。
态射 $i' = (i, \pi) : X' \to \mathbf{P}^n_X$
是闭浸入（同上）。因此
$$
\mathcal{L} = i^*\mathcal{O}_{\mathbf{P}^n_k}(1) \cong
(i')^*\mathcal{O}_{\mathbf{P}^n_X}(1)
$$
是 $\pi$-相对丰沛的（例如由《态射》引理
\ref{morphisms-lemma-characterize-ample-on-finite-type}）。
所以由《概形上同调》引理
\ref{coherent-lemma-kill-by-twisting}，存在 $n \geq 0$，使得
$R^p\pi_*\mathcal{L}^{\otimes n} = 0$ 对所有 $p > 0$ 都成立。
置 $\mathcal{G} = \pi_*\mathcal{L}^{\otimes n}$。
任取非零整体截面 $s$；它是 $\mathcal{L}^{\otimes n}$ 的截面。
由于 $\mathcal{G} = \pi_*\mathcal{L}^{\otimes n}$，截面 $s$
对应于 $\mathcal{G}$ 的一个截面，即一个映射
$\mathcal{O}_X \to \mathcal{G}$。
有 $s|_U \not = 0$，因为 $X'$ 是簇且 $\mathcal{L}$ 可逆，
所以 $\mathcal{O}_X|_U \to \mathcal{G}|_U$ 非零。
由于 $\mathcal{G}|_U = \mathcal{L}^{\otimes n}|_{\pi^{-1}(U)}$
可逆，得到短正合列
$$
0 \to \mathcal{O}_X \to \mathcal{G} \to \mathcal{Q} \to 0
$$
，其中 $\mathcal{Q}$ 凝聚且支撑在 $X$ 的真闭子概形上。
像先前一样使用归纳假设，可见只需证明
$\dim H^d(X, \mathcal{G}) < \infty$。

\medskip\noindent
由 Leray 谱序列
（《上同调》引理 \ref{cohomology-lemma-apply-Leray}）可见
$H^d(X, \mathcal{G}) = H^d(X', \mathcal{L}^{\otimes n})$。
设 $\overline{X}' \subset \mathbf{P}^n_k$ 为 $X'$ 的闭包。
则 $\overline{X}'$ 是维数为 $d$、定义在 $k$ 上的射影簇，而
$X' \subset \overline{X}'$ 是稠密开子概形。
可逆层 $\mathcal{L}^{\otimes n}$ 是
$\mathcal{O}_{\overline{X}'}(n)$ 在 $X'$ 上的限制。由
《上同调》命题
\ref{cohomology-proposition-cohomological-dimension-spectral}，映射
$$
H^d(\overline{X}', \mathcal{O}_{\overline{X}'}(n))
\longrightarrow
H^d(X', \mathcal{L}^{\otimes n})
$$
是满射。左侧上同调群由《概形上同调》引理
\ref{coherent-lemma-coherent-projective} 是有限维的，证明完成。
\end{proof}

\begin{lemma}
\label{etale-cohomology-lemma-vanishing-variety-char-p-p}
设 $X$ 是代数闭域 $k$ 上分离的有限型概形，且该域的特征为
$p > 0$。则
$H_\etale^q(X, \underline{\mathbf{Z}/p\mathbf{Z}}) = 0$
对 $q \geq dim(X) + 1$ 成立。
\end{lemma}

\begin{proof}
设 $d = \dim(X)$。由引理
\ref{etale-cohomology-lemma-vanishing-affine-char-p-p} 中建立的消失，只需说明
$H_\etale^{d + 1}(X, \underline{\mathbf{Z}/p\mathbf{Z}}) = 0$。
由引理 \ref{etale-cohomology-lemma-top-cohomology-coherent} 可见
$H^d(X, \mathcal{O}_X)$ 是有限维 $k$-向量空间。
因此，与 Artin--Schreier 序列相伴的上同调长正合列以
$$
H^d(X, \mathcal{O}_X) \xrightarrow{F - 1}
H^d(X, \mathcal{O}_X) \to H^{d + 1}_\etale(X, \mathbf{Z}/p\mathbf{Z}) \to 0
$$
结尾。由引理 \ref{etale-cohomology-lemma-F-1}，此序列中的映射 $F - 1$ 是满射。
这就证明了该引理。
\end{proof}

\begin{lemma}
\label{etale-cohomology-lemma-finiteness-proper-variety-char-p-p}
设 $X$ 是代数闭域 $k$ 上的固有概形，且该域的特征为 $p > 0$。则
\begin{enumerate}
\item $H_\etale^q(X, \underline{\mathbf{Z}/p\mathbf{Z}})$
是有限 $\mathbf{Z}/p\mathbf{Z}$-模，对所有 $q$ 都如此；并且
\item 映射
$H^q_\etale(X, \underline{\mathbf{Z}/p\mathbf{Z}}) \to
H^q_\etale(X_{k'}, \underline{\mathbf{Z}/p\mathbf{Z}})$
是同构，只要 $k'/k$ 是代数闭域扩张。
\end{enumerate}
\end{lemma}

\begin{proof}
由《概形上同调》引理
\ref{coherent-lemma-proper-over-affine-cohomology-finite}
以及定理 \ref{etale-cohomology-theorem-zariski-fpqc-quasi-coherent} 中的上同调比较，
上同调群
$H^q_\etale(X, \mathbf{G}_a) = H^q(X, \mathcal{O}_X)$
是有限维 $k$-向量空间。因此，由引理 \ref{etale-cohomology-lemma-F-1}，
与 Artin--Schreier 序列相伴的上同调长正合列分解为短正合列
$$
0 \to H_\etale^q(X, \underline{\mathbf{Z}/p\mathbf{Z}}) \to
H^q(X, \mathcal{O}_X) \xrightarrow{F - 1} H^q(X, \mathcal{O}_X) \to 0
$$
而且 $\mathbf{F}_p$-维数（对上同调群
$H_\etale^q(X, \underline{\mathbf{Z}/p\mathbf{Z}})$ 而言）
小于或等于 $k$-维数（对向量空间
$H^q(X, \mathcal{O}_X)$ 而言）。
这证明了第一条。第二条由引理 \ref{etale-cohomology-lemma-F-2} 得到，因为
$H^q(X, \mathcal{O}_X) \otimes_k k' \to H^q(X_{k'}, \mathcal{O}_{X_{k'}})$
由平坦基变换是同构
（《概形上同调》引理
\ref{coherent-lemma-flat-base-change-cohomology}）。
\end{proof}








\section{局部常值层}
\label{etale-cohomology-section-locally-constant}

\noindent
本节是《网站上的模》第
\ref{sites-modules-section-locally-constant} 节针对 \'etale 网站的类似版本。

\begin{definition}
\label{etale-cohomology-definition-finite-locally-constant}
设 $X$ 为概形。
设 $\mathcal{F}$ 为 $X_\etale$ 上的集合层。
\begin{enumerate}
\item 设 $E$ 为集合。称 $\mathcal{F}$ 为{\it 取值为 $E$ 的常值层}，
如果 $\mathcal{F}$ 是预层 $U \mapsto E$ 的层化。
记作 $\underline{E}_X$ 或 $\underline{E}$。
\item 若 $\mathcal{F}$ 同构于第 (1) 项中的层，则称其为{\it 常值层}。
\item 称 $\mathcal{F}$ {\it 局部常值}，如果存在覆盖
$\{U_i \to X\}$，使得 $\mathcal{F}|_{U_i}$ 为常值层。
\item 若 $\mathcal{F}$ 局部常值，而且各常值取值都是有限集合，
则称其{\it 有限局部常值}。
\end{enumerate}
设 $\mathcal{F}$ 为 $X_\etale$ 上的阿贝尔群层。
\begin{enumerate}
\item 设 $A$ 为阿贝尔群。
称 $\mathcal{F}$ 为{\it 取值为 $A$ 的常值层}，
如果 $\mathcal{F}$ 是预层 $U \mapsto A$ 的层化。
记作 $\underline{A}_X$ 或 $\underline{A}$。
\item 若 $\mathcal{F}$ 作为阿贝尔群层同构于第 (1) 项中的层，
则称其为{\it 常值层}。
\item 称 $\mathcal{F}$ {\it 局部常值}，如果存在覆盖
$\{U_i \to X\}$，使得 $\mathcal{F}|_{U_i}$ 为常值层。
\item 若 $\mathcal{F}$ 局部常值，而且各常值取值都是有限阿贝尔群，
则称其{\it 有限局部常值}。
\end{enumerate}
设 $\Lambda$ 为环。设 $\mathcal{F}$ 为 $\Lambda$-模层，
它定义在 $X_\etale$ 上。
\begin{enumerate}
\item 设 $M$ 为 $\Lambda$-模。
称 $\mathcal{F}$ 为{\it 取值为 $M$ 的常值层}，
如果 $\mathcal{F}$ 是预层 $U \mapsto M$ 的层化。
记作 $\underline{M}_X$ 或 $\underline{M}$。
\item 若 $\mathcal{F}$ 作为 $\Lambda$-模层同构于第 (1) 项中的层，
则称其为{\it 常值层}。
\item 称 $\mathcal{F}$ {\it 局部常值}，如果存在覆盖
$\{U_i \to X\}$，使得 $\mathcal{F}|_{U_i}$ 为常值层。
\end{enumerate}
\end{definition}

\begin{lemma}
\label{etale-cohomology-lemma-pullback-locally-constant}
设 $f : X \to Y$ 为概形态射。若 $\mathcal{G}$ 是局部常值的
集合层、阿贝尔群层或 $\Lambda$-模层，并且定义在 $Y_\etale$ 上，
则 $f^{-1}\mathcal{G}$ 在 $X_\etale$ 上也具有同样的性质。
\end{lemma}

\begin{proof}
这对任意拓扑斯态射都成立，见《网站上的模》引理
\ref{sites-modules-lemma-pullback-locally-constant}。
\end{proof}

\begin{lemma}
\label{etale-cohomology-lemma-pushforward-locally-constant}
设 $f : X \to Y$ 为有限 \'etale 概形态射。
若 $\mathcal{F}$ 是（有限）局部常值的集合层、
（有限）局部常值的阿贝尔群层或（有限型）局部常值的
$\Lambda$-模层，并且定义在 $X_\etale$ 上，则
$f_*\mathcal{F}$ 在 $Y_\etale$ 上也具有同样的性质。
\end{lemma}

\begin{proof}
$f_*$ 的构造与 \'etale 局部化可交换。
有限 \'etale 态射局部同构于若干同构的不交并，见
《\'Etale 态射》引理 \ref{etale-lemma-finite-etale-etale-local}。
因此，该引理是说：若 $\mathcal{F}_i$（$i = 1, \ldots, n$）
都是（有限）局部常值的集合层，则
$\prod_{i = 1, \ldots, n} \mathcal{F}_i$ 也是。这是显然的。
阿贝尔群层和模层的情形类似。
\end{proof}

\begin{lemma}
\label{etale-cohomology-lemma-characterize-finite-locally-constant}
设 $X$ 为概形，$\mathcal{F}$ 为 $X_\etale$ 上的集合层。
则下列条件等价：
\begin{enumerate}
\item $\mathcal{F}$ 有限局部常值；并且
\item 有 $\mathcal{F} = h_U$，其中 $U \to X$ 是某个有限 \'etale 态射。
\end{enumerate}
\end{lemma}

\begin{proof}
有限 \'etale 态射局部同构于若干同构的不交并，见
《\'Etale 态射》引理 \ref{etale-lemma-finite-etale-etale-local}。
因此 (2) 蕴含 (1)。反过来，若 $\mathcal{F}$ 有限局部常值，
则存在 \'etale 覆盖 $\{X_i \to X\}$，使得
$\mathcal{F}|_{X_i}$ 可由有限 \'etale 态射 $U_i \to X_i$ 表示。
完全仿照《下降》引理
\ref{descent-lemma-descent-data-sheaves} 的证明进行论证，
可得到概形下降数据 $(U_i, \varphi_{ij})$，它相对于
$\{X_i \to X\}$（细节从略）。例如由《下降》引理
\ref{descent-lemma-affine}，此下降数据是有效的；所得概形态射
$U \to X$ 由《下降》引理
\ref{descent-lemma-descending-property-finite} 和
\ref{descent-lemma-descending-property-etale} 是有限 \'etale 的。
\end{proof}

\begin{lemma}
\label{etale-cohomology-lemma-morphism-locally-constant}
设 $X$ 为概形。
\begin{enumerate}
\item 设 $\varphi : \mathcal{F} \to \mathcal{G}$ 是
$X_\etale$ 上局部常值集合层的映射。
若 $\mathcal{F}$ 有限局部常值，则存在 \'etale 覆盖
$\{U_i \to X\}$，使得 $\varphi|_{U_i}$
是由集合映射所伴随的常值层映射。
\item 设 $\varphi : \mathcal{F} \to \mathcal{G}$ 是
$X_\etale$ 上局部常值阿贝尔群层的映射。
若 $\mathcal{F}$ 有限局部常值，则存在 \'etale 覆盖
$\{U_i \to X\}$，使得 $\varphi|_{U_i}$
是由阿贝尔群映射所伴随的常值阿贝尔群层映射。
\item 设 $\Lambda$ 为环。
设 $\varphi : \mathcal{F} \to \mathcal{G}$ 是
局部常值 $\Lambda$-模层的映射，并且这些层定义在 $X_\etale$ 上。
若 $\mathcal{F}$ 有限型，则存在 \'etale 覆盖
$\{U_i \to X\}$，使得 $\varphi|_{U_i}$
是由 $\Lambda$-模映射所伴随的常值 $\Lambda$-模层映射。
\end{enumerate}
\end{lemma}

\begin{proof}
这对任意网站都成立，见《网站上的模》引理
\ref{sites-modules-lemma-morphism-locally-constant}。
\end{proof}

\begin{lemma}
\label{etale-cohomology-lemma-kernel-finite-locally-constant}
设 $X$ 为概形。
\begin{enumerate}
\item 有限局部常值集合层的范畴在 $\Sh(X_\etale)$ 内
对有限极限与有限余极限封闭。
\item 有限局部常值阿贝尔群层的范畴是
$\textit{Ab}(X_\etale)$ 的弱 Serre 子范畴。
\item 设 $\Lambda$ 为诺特环。有限型局部常值
$\Lambda$-模层（定义在 $X_\etale$ 上）的范畴是
$\textit{Mod}(X_\etale, \Lambda)$ 的弱 Serre 子范畴。
\end{enumerate}
\end{lemma}

\begin{proof}
这对任意网站都成立，见《网站上的模》引理
\ref{sites-modules-lemma-kernel-finite-locally-constant}。
\end{proof}

\begin{lemma}
\label{etale-cohomology-lemma-tensor-product-locally-constant}
设 $X$ 为概形。设 $\Lambda$ 为环。
两个局部常值 $\Lambda$-模层（定义在 $X_\etale$ 上）的张量积仍是
局部常值 $\Lambda$-模层。
\end{lemma}

\begin{proof}
这对任意网站都成立，见《网站上的模》引理
\ref{sites-modules-lemma-tensor-product-locally-constant}。
\end{proof}

\begin{lemma}
\label{etale-cohomology-lemma-connected-locally-constant}
设 $X$ 为连通概形。设 $\Lambda$ 为环，并设
$\mathcal{F}$ 为局部常值 $\Lambda$-模层。
则存在 $\Lambda$-模 $M$ 以及 \'etale 覆盖
$\{U_i \to X\}$，使得 $\mathcal{F}|_{U_i} \cong \underline{M}|_{U_i}$。
\end{lemma}

\begin{proof}
选取 \'etale 覆盖
$\{U_i \to X\}$，使得 $\mathcal{F}|_{U_i}$ 为常值层；设
$\mathcal{F}|_{U_i} \cong \underline{M_i}_{U_i}$。
注意，$U_i \times_X U_j$ 为空，只要 $M_i$ 不同构于 $M_j$。
对每个 $\Lambda$-模 $M$，令 $I_M = \{i \in I \mid M_i \cong M\}$。
由于 \'etale 态射是开映射，可见
$U_M = \bigcup_{i \in I_M} \Im(U_i \to X)$
是 $X$ 的开子集。于是 $X = \coprod U_M$ 是 $X$ 的不交开覆盖。
由于 $X$ 连通，只有一个 $U_M$ 非空，引理得证。
\end{proof}







\section{局部常值层与基本群}
\label{etale-cohomology-section-pione}

\noindent
在某些情形下，可以把局部常值层与概形的基本群联系起来。

\begin{lemma}
\label{etale-cohomology-lemma-locally-constant-on-connected}
设 $X$ 为连通概形。设 $\overline{x}$ 为 $X$ 的几何点。
\begin{enumerate}
\item 存在范畴等价
$$
\left\{
\begin{matrix}
\text{有限局部常值}\\
X_\etale\text{ 上的集合层}
\end{matrix}
\right\}
\longleftrightarrow
\left\{
\begin{matrix}
\text{有限 }\pi_1(X, \overline{x})\text{-集合}
\end{matrix}
\right\}
$$
\item 存在范畴等价
$$
\left\{
\begin{matrix}
\text{有限局部常值}\\
X_\etale\text{ 上的阿贝尔群层}
\end{matrix}
\right\}
\longleftrightarrow
\left\{
\begin{matrix}
\text{有限 }\pi_1(X, \overline{x})\text{-模}
\end{matrix}
\right\}
$$
\item 设 $\Lambda$ 为有限环。存在范畴等价
$$
\left\{
\begin{matrix}
\text{有限型局部常值}\\
X_\etale\text{ 上的 }\Lambda\text{-模层}
\end{matrix}
\right\}
\longleftrightarrow
\left\{
\begin{matrix}
\text{带有相容结构的有限 }\pi_1(X, \overline{x})\text{-模，}\\
\text{其 }\Lambda\text{-模结构与作用可交换}
\end{matrix}
\right\}
$$
\end{enumerate}
\end{lemma}

\begin{proof}
注意，$\pi_1(X, \overline{x})$ 是预有限拓扑群，见
《基本群》定义 \ref{pione-definition-fundamental-group}。
左侧各范畴定义于第 \ref{etale-cohomology-section-locally-constant} 节。
右侧范畴所用的记号，对集合取自《基本群》定义
\ref{pione-definition-G-set-continuous}，对阿贝尔群取自定义
\ref{etale-cohomology-definition-G-module-continuous}。这就解释了这些记号。

\medskip\noindent
断言 (1) 由引理 \ref{etale-cohomology-lemma-characterize-finite-locally-constant}
和《基本群》定理 \ref{pione-theorem-fundamental-group} 得到。
第 (2)、(3) 条由此立即得到，只须赋予底层（层的）集合额外结构。
例如，$X_\etale$ 上有限局部常值的阿贝尔群层，
等同于一个有限局部常值集合层 $\mathcal{F}$，连同一个满足通常公理的映射
$+ : \mathcal{F} \times \mathcal{F} \to \mathcal{F}$。
(1) 中的等价把积映到积，因而把 $+$ 映到相应有限
$\pi_1(X, \overline{x})$-集合上的加法。由于
$\pi_1(X, \overline{x})$-模等同于带相容阿贝尔群结构的
$\pi_1(X, \overline{x})$-集合，便得到 (2)。第 (3) 条完全同理。
\end{proof}

\begin{lemma}
\label{etale-cohomology-lemma-locally-constant-on-connected-geom-unibranch}
设 $X$ 为不可约几何单分支概形。
设 $\overline{x}$ 为 $X$ 的几何点。
设 $\Lambda$ 为环。存在范畴等价
$$
\left\{
\begin{matrix}
\text{有限型局部常值}\\
X_\etale\text{ 上的 }\Lambda\text{-模层}
\end{matrix}
\right\}
\longleftrightarrow
\left\{
\begin{matrix}
\text{有限 }\Lambda\text{-模 }M\text{，带有连续的}\\
\pi_1(X, \overline{x})\text{-作用}
\end{matrix}
\right\}
$$
\end{lemma}

\begin{proof}
引理 \ref{etale-cohomology-lemma-locally-constant-on-connected} 中给出的证明不适用，
因为有限 $\Lambda$-模 $M$ 的底层集合未必有限。

\medskip\noindent
设 $\nu : X^\nu \to X$ 为正规化态射。由《态射》引理
\ref{morphisms-lemma-normalization-geom-unibranch-univ-homeo}，
这是万有同胚。由《基本群》命题
\ref{pione-proposition-universal-homeomorphism}，它诱导同构
$\pi_1(X^\nu, \overline{x}) \to \pi_1(X, \overline{x})$；
由定理 \ref{etale-cohomology-theorem-topological-invariance}，可得有限型局部常值
$\Lambda$-模层的两个范畴之间的等价；这两个范畴分别定义在
$X_\etale$ 与 $X^\nu_\etale$ 上。
这把问题归结到 $X$ 为整正规概形的情形。

\medskip\noindent
假设 $X$ 是整正规概形。设 $\eta \in X$ 为泛点。
设 $\overline{\eta}$ 为位于 $\eta$ 上方的几何点。
由《基本群》命题 \ref{pione-proposition-normal}，有连续满射
$$
\text{Gal}(\kappa(\eta)^{sep}/\kappa(\eta)) = \pi_1(\eta, \overline{\eta})
\longrightarrow
\pi_1(X, \overline{\eta})
$$
，其核描述于《基本群》引理
\ref{pione-lemma-normal-pione-quotient-inertia}。
设 $\mathcal{F}$ 为有限型局部常值的
$\Lambda$-模层，它定义在 $X_\etale$ 上。设 $M = \mathcal{F}_{\overline{\eta}}$
为 $\mathcal{F}$ 在 $\overline{\eta}$ 处的茎。
由第 \ref{etale-cohomology-section-galois-action-stalks} 节，得到
$\text{Gal}(\kappa(\eta)^{sep}/\kappa(\eta))$ 在 $M$ 上的连续作用。
目标是说明此作用经由上面显示的满射分解。
由于 $\mathcal{F}$ 有限型，$M$ 是有限 $\Lambda$-模。
由于 $\mathcal{F}$ 局部常值，对每个 $x \in X$，
$\mathcal{F}$ 到 $\Spec(\mathcal{O}_{X, x}^{sh})$ 的限制都是常值的。
因此，$\text{Gal}(K^{sep}/K_x^{sh})$ 在 $M$ 上的作用平凡
（记号如《基本群》引理
\ref{pione-lemma-normal-pione-quotient-inertia}）。
由此得到所需分解。

\medskip\noindent
反过来，假设有一个有限 $\Lambda$-模 $M$，
带有 $\pi_1(X, \overline{\eta})$ 的连续作用。
我们将构造 $\mathcal{F}$，使得
$M \cong \mathcal{F}_{\overline{\eta}}$ 作为
$\Lambda[\pi_1(X, \overline{\eta})]$-模。
选取生成元 $m_1, \ldots, m_r \in M$。
由于 $\pi_1(X, \overline{\eta})$ 在 $M$ 上的作用连续，
对每个 $i$ 都存在开子群 $N_i$，它是预有限群
$\pi_1(X, \overline{\eta})$ 的开子群，
使得每个 $\gamma \in H_i$ 都固定 $m_i$。因此开子群
$H = \bigcap_{i = 1, \ldots, r} H_i$ 的每个元素都固定
$M$ 的每个元素。缩小 $H$ 后，可以假设 $H$ 是
$\pi_1(X, \overline{\eta})$ 的开正规子群。
置 $G = \pi_1(X, \overline{\eta})/H$。设 $f : Y \to X$
为相应的 Galois 有限 \'etale $G$-覆盖。
可把 $f_*\underline{\mathbf{Z}}$ 看作
$\mathbf{Z}[G]$-模层，它定义在 $X_\etale$ 上。然后只须取
$$
\mathcal{F} =
f_*\underline{\mathbf{Z}} \otimes_{\underline{\mathbf{Z}[G]}} \underline{M}
$$
把 $\mathcal{F}_{\overline{\eta}}$ 的计算留给读者。
我们也略去此构造是上一段构造之逆的验证。
\end{proof}

\begin{remark}
\label{etale-cohomology-remark-functorial-locally-constant-on-connected}
引理 \ref{etale-cohomology-lemma-locally-constant-on-connected} 与
\ref{etale-cohomology-lemma-locally-constant-on-connected-geom-unibranch} 中的等价
与拉回相容。例如，假设 $f : Y \to X$ 是连通概形之间的态射。
设 $\overline{y}$ 为 $Y$ 的几何点，并置
$\overline{x} = f(\overline{y})$。则图
$$
\xymatrix{
Y_\etale\text{ 上有限局部常值的集合层}
\ar[r] &
\text{有限 }\pi_1(Y, \overline{y})\text{-集合} \\
X_\etale\text{ 上有限局部常值的集合层}
\ar[r] \ar[u]_{f^{-1}} &
\text{有限 }\pi_1(X, \overline{x})\text{-集合} \ar[u]
}
$$
可交换，其中右侧竖直箭头来自连续同态
$\pi_1(Y, \overline{y}) \to \pi_1(X, \overline{x})$；它由 $f$ 诱导。
这立即由《基本群》定理
\ref{pione-theorem-fundamental-group} 中的交换图得到。
其余情形有类似结果。
\end{remark}









\section{迹方法}
\label{etale-cohomology-section-trace-method}

\noindent
本节的参考文献是 \cite[Expos\'e IX, \S 5]{SGA4}。
这里的材料将在下文引理 \ref{etale-cohomology-lemma-vanishing-easier} 的证明中使用。

\medskip\noindent
设 $f : Y \to X$ 为概形的 \'etale 态射。有一列伴随函子
$$
f_!, f^{-1}, f_*
$$
，它们位于 $\textit{Ab}(X_\etale)$ 与 $\textit{Ab}(Y_\etale)$ 之间。
函子 $f_!$ 在第 \ref{etale-cohomology-section-extension-by-zero} 节讨论。
伴随映射 $\text{id} \to f_* f^{-1}$ 称为{\it 限制}。
伴随映射 $f_! f^{-1} \to \text{id}$ 通常称为{\it 迹映射}。
若 $f$ 有限 \'etale，则 $f_* = f_!$
（引理 \ref{etale-cohomology-lemma-shriek-equals-star-finite-etale}），因而可把它看作映射
$f_*f^{-1} \to \text{id}$。

\begin{definition}
\label{etale-cohomology-definition-trace-map}
设 $f : Y \to X$ 为有限 \'etale 概形态射。
上面描述并在下面明确给出的映射 $f_* f^{-1} \to \text{id}$
称为{\it 迹}。
\end{definition}

\noindent
设 $f : Y \to X$ 为有限 \'etale 概形态射。迹映射由以下两个性质刻画：
\begin{enumerate}
\item 它与 $X$ 上的 \'etale 局部化可交换；并且
\item 若 $Y = \coprod_{i = 1}^d X$，则迹映射是求和映射
$f_*f^{-1} \mathcal{F} = \mathcal{F}^{\oplus d} \to \mathcal{F}$。
\end{enumerate}
由《\'Etale 态射》引理 \ref{etale-lemma-finite-etale-etale-local}，
每个有限 \'etale 态射 $f : Y \to X$ 在 $X$ 的 \'etale 局部，
都具有 (2) 中给出的形式，其中某个整数 $d \geq 0$。
因此可以用上述刻画定义迹映射；特别地，为构造迹映射，
无须知道 $f_!$ 的存在性以及 $f_!$ 与 $f_*$ 的相等性。
此描述表明，若 $f$ 的常次数为 $d$，则复合
$$
\mathcal{F} \xrightarrow{res}
f_* f^{-1} \mathcal{F} \xrightarrow{trace}
\mathcal{F}
$$
是乘以 $d$。所谓“迹方法”是如下观察：若 $\mathcal{F}$
是 $X_\etale$ 上的阿贝尔群层，且乘以 $d$ 在 $\mathcal{F}$ 上
是同构，则映射
$$
H^n_\etale(X, \mathcal{F}) \longrightarrow H^n_\etale(Y, f^{-1}\mathcal{F})
$$
是单射。事实上，由高阶直像的消失
（命题 \ref{etale-cohomology-proposition-finite-higher-direct-image-zero}）
和 Leray 谱序列（命题 \ref{etale-cohomology-proposition-leray}），有
$$
H^n_\etale(Y, f^{-1}\mathcal{F}) = H^n_\etale(X, f_*f^{-1}\mathcal{F})
$$
。因此可以考虑映射
$$
H^n_\etale(X, \mathcal{F}) \to
H^n_\etale(Y, f^{-1}\mathcal{F})= H^n_\etale(X, f_*f^{-1}\mathcal{F})
\xrightarrow{trace}
H^n_\etale(X, \mathcal{F})
$$
，而复合是同构（在我们对 $\mathcal{F}$ 与 $f$ 的假设下）。
特别地，若 $H_\etale^q(Y, f^{-1}\mathcal{F}) = 0$，则
$H_\etale^q(X, \mathcal{F}) = 0$ 也成立。
确实，乘以 $d$ 在 $H_\etale^q(X, \mathcal{F})$ 上诱导同构，
而此同构经由 $H_\etale^q(Y, f^{-1}\mathcal{F})= 0$ 分解。

\medskip\noindent
这常与下面的引理结合使用。

\begin{lemma}
\label{etale-cohomology-lemma-pullback-filtered}
设 $S$ 为连通概形。设 $\ell$ 为素数。设
$\mathcal{F}$ 为有限型局部常值的
$\mathbf{F}_\ell$-向量空间层，它定义在 $S_\etale$ 上。
则存在有限 \'etale 态射 $f : T \to S$，其次数与
$\ell$ 互素，并且 $f^{-1}\mathcal{F}$
有有限滤过，其逐次商均为
$\underline{\mathbf{Z}/\ell\mathbf{Z}}_T$。
\end{lemma}

\begin{proof}
选取几何点 $\overline{s}$，它属于 $S$。
通过引理 \ref{etale-cohomology-lemma-locally-constant-on-connected} 的等价，
层 $\mathcal{F}$ 对应于有限维
$\mathbf{F}_\ell$-向量空间 $V$，它带有
$\pi_1(S, \overline{s})$ 的连续作用。
设 $G \subset \text{Aut}(V)$ 为给出该作用的同态
$\rho : \pi_1(S, \overline{s}) \to \text{Aut}(V)$ 的像。
注意 $G$ 有限。连续满同态
$\overline{\rho} : \pi_1(S, \overline{s}) \to G$
对应于 Galois 对象 $Y \to S$，它属于 $\textit{F\'Et}_S$，
其自同构群为 $G = \text{Aut}(Y/S)$，见《基本群》第
\ref{pione-section-finite-etale-under-galois} 节。
设 $H \subset G$ 为 $\ell$-Sylow 子群。
我们断言 $T = Y/H \to S$ 符合要求。事实上，设
$\overline{t} \in T$ 为位于 $\overline{s}$ 上方的几何点。
由基本群的函子性，
$\pi_1(T, \overline{t}) \to \pi_1(S, \overline{s})$ 的像
是 $(\overline{\rho})^{-1}(H)$。因此，
$\pi_1(T, \overline{t})$ 在 $V$ 上的作用（对应于
$f^{-1}\mathcal{F}$）经由映射
$\pi_1(T, \overline{t}) \to H$ 给出，见评注
\ref{etale-cohomology-remark-functorial-locally-constant-on-connected}。
由于 $H$ 是有限 $\ell$-群，表示
$\rho|_{\pi_1(T, \overline{t})}$ 的每个不可约成分
都是秩为 $1$ 的平凡表示（这是有限群表示论中的简单引理；
此处插入将来的参考文献）。
通过引理 \ref{etale-cohomology-lemma-locally-constant-on-connected} 的等价，
这意味着 $f^{-1}\mathcal{F}$ 是一列逐次扩张，
其中各常值层的取值为 $\underline{\mathbf{Z}/\ell\mathbf{Z}}_T$。
此外，$T = Y/H \to S$ 的次数与 $\ell$ 互素，
因为它等于 $H$ 在 $G$ 中的指数。
\end{proof}

\begin{lemma}
\label{etale-cohomology-lemma-action-l-group}
设 $\Lambda$ 为诺特环。
设 $\ell$ 为素数且 $n \geq 1$。
设 $H$ 为有限 $\ell$-群。
设 $M$ 为有限 $\Lambda[H]$-模，并被 $\ell^n$ 零化。
则存在有限滤过
$0 = M_0 \subset M_1 \subset \ldots \subset M_t = M$，
它由 $\Lambda[H]$-子模组成，使得 $H$ 在
$M_{i + 1}/M_i$ 上的作用平凡，对所有
$i = 0, \ldots, t - 1$ 都如此。
\end{lemma}

\begin{proof}
从略。提示：证明增广理想 $\mathfrak m$ 是幂零的；它属于非交换环
$\mathbf{Z}/\ell^n\mathbf{Z}[H]$。
\end{proof}

\begin{lemma}
\label{etale-cohomology-lemma-pullback-filtered-modules}
设 $S$ 为不可约几何单分支概形。
设 $\ell$ 为素数且 $n \geq 1$。设 $\Lambda$
为诺特环。设 $\mathcal{F}$ 为有限型局部常值的
$\Lambda$-模层，它定义在 $S_\etale$ 上，并且被 $\ell^n$ 零化。
则存在有限 \'etale 态射 $f : T \to S$，其次数与 $\ell$
互素，并且
$f^{-1}\mathcal{F}$ 有有限滤过，其逐次商形如
$\underline{M}_T$，其中取某个有限 $\Lambda$-模 $M$。
\end{lemma}

\begin{proof}
选取几何点 $\overline{s}$，它属于 $S$。
通过引理 \ref{etale-cohomology-lemma-locally-constant-on-connected-geom-unibranch}
的等价，层 $\mathcal{F}$ 对应于有限 $\Lambda$-模
$M$，它带有 $\pi_1(S, \overline{s})$ 的连续作用。
设 $G \subset \text{Aut}(V)$ 为给出该作用的同态
$\rho : \pi_1(S, \overline{s}) \to \text{Aut}(M)$ 的像。
注意，$G$ 有限，因为 $M$ 是有限 $\Lambda$-模
（见引理 \ref{etale-cohomology-lemma-locally-constant-on-connected-geom-unibranch} 的证明）。
连续满同态
$\overline{\rho} : \pi_1(S, \overline{s}) \to G$
对应于 Galois 对象 $Y \to S$，它属于 $\textit{F\'Et}_S$，
其自同构群为 $G = \text{Aut}(Y/S)$，见《基本群》第
\ref{pione-section-finite-etale-under-galois} 节。
设 $H \subset G$ 为 $\ell$-Sylow 子群。
我们断言 $T = Y/H \to S$ 符合要求。事实上，设
$\overline{t} \in T$ 为位于 $\overline{s}$ 上方的几何点。
由基本群的函子性，
$\pi_1(T, \overline{t}) \to \pi_1(S, \overline{s})$ 的像
是 $(\overline{\rho})^{-1}(H)$。因此，
$\pi_1(T, \overline{t})$ 在 $M$ 上的作用（对应于
$f^{-1}\mathcal{F}$）经由映射
$\pi_1(T, \overline{t}) \to H$ 给出，见评注
\ref{etale-cohomology-remark-functorial-locally-constant-on-connected}。
取引理 \ref{etale-cohomology-lemma-action-l-group} 中的滤过
$0 = M_0 \subset M_1 \subset \ldots \subset M_t = M$。
它诱导滤过
$0 = \mathcal{F}_0 \subset \mathcal{F}_1 \subset \ldots
\subset \mathcal{F}_t = f^{-1}\mathcal{F}$，
其逐次商是取值为 $M_{i + 1}/M_i$ 的常值层。
最后，$T = Y/H \to S$ 的次数与 $\ell$ 互素，
因为它等于 $H$ 在 $G$ 中的指数。
\end{proof}







\section{Galois 上同调}
\label{etale-cohomology-section-galois-cohomology}

\noindent
本节证明一个关于 Galois 上同调的结果
（命题 \ref{etale-cohomology-proposition-serre-galois}），所用工具是 \'etale 上同调以及
第 \ref{etale-cohomology-section-trace-method} 节中的技巧。
这将使我们能够证明域的谱上的高阶 \'etale 上同调群消失。

\begin{lemma}
\label{etale-cohomology-lemma-nonvanishing-inherited}
设 $\ell$ 是素数，$n$ 是整数且 $> 0$。
设 $S$ 是拟紧拟分离概形。
设 $X = \lim_{i \in I} X_i$ 是一个 $S$-概形有向系的极限，
其中每个 $X_i \to S$ 都是常次数的有限 \'etale 态射，且该次数与 $\ell$ 互素。
下列条件等价：
\begin{enumerate}
\item 存在一个 $\ell$-幂挠层 $\mathcal{G}$ 于 $S$ 上，使得
$H_\etale^n(S, \mathcal{G}) \neq 0$；
\item 存在一个 $\ell$-幂挠层 $\mathcal{F}$ 于 $X$ 上，使得
$H_\etale^n(X, \mathcal{F}) \neq 0$。
\end{enumerate}
事实上，给定 $\mathcal{G}$，可以取
$\mathcal{F} = g^{-1}\mathcal{F}$；给定 $\mathcal{F}$，可以取
$\mathcal{G} = g_*\mathcal{F}$。
\end{lemma}

\begin{proof}
以 $g : X \to S$ 和 $g_i : X_i \to S$ 表示结构态射。
固定一个 $\ell$-幂挠层 $\mathcal{G}$ 于 $S$ 上，使得
$H^n_\etale(S, \mathcal{G}) \not = 0$。
由 $\mathcal{G}_i = g_i^{-1}\mathcal{G}$ 给出的系统满足定理
\ref{etale-cohomology-theorem-colimit} 的条件，其余极限层为 $g^{-1}\mathcal{G}$。因此
$$
\colim_{i\in I} H^n_\etale(X_i, g_i^{-1}\mathcal{G}) =
H^n_\etale(X, \mathcal{G})
$$
由于 $g_i$ 是次数与 $\ell$ 互素的有限 \'etale 态射，可以应用“迹方法”，得到映射
$$
H^n_\etale(S, \mathcal{G}) \to H^n_\etale(X_i, g_i^{-1}\mathcal{G})
$$
全都是单射（并且与转移映射相容）。见第 \ref{etale-cohomology-section-trace-method} 节。
因此余极限非零，即 $H^n(X,g^{-1}\mathcal{G}) \neq 0$，
取 $\mathcal{F} = g^{-1}\mathcal{G}$ 即得所需结论。

\medskip\noindent
反之，设给定一个 $\ell$-幂挠层 $\mathcal{F}$ 于 $X$ 上，使得
$H^n_\etale(X, \mathcal{F}) \not = 0$。注意，由于 $g_i$ 是有限态射，
其高阶正像消失（命题 \ref{etale-cohomology-proposition-finite-higher-direct-image-zero}）。
于是应用引理 \ref{etale-cohomology-lemma-relative-colimit}，可知 $g$ 也具有同样的性质。
高阶正像的消失以及 Leray（命题 \ref{etale-cohomology-proposition-leray}）给出
$H^n_\etale(X, \mathcal{F}) = H^n(S, g_*\mathcal{F}) \neq 0$；
取 $\mathcal{G} = g_*\mathcal{F}$ 即得所需结论。
\end{proof}

\begin{lemma}
\label{etale-cohomology-lemma-reduce-to-l-group}
设 $\ell$ 是素数，$n$ 是整数且 $> 0$。设 $K$ 是域，
$G = Gal(K^{sep}/K)$，并设 $H \subset G$ 是极大 pro-$\ell$ 子群，
$L/K$ 是相应的域扩张。则
$H^n_\etale(\Spec(K), \mathcal{F}) = 0$ 对所有 $\ell$-幂挠层
$\mathcal{F}$ 成立，当且仅当
$H^n_\etale(\Spec(L), \underline{\mathbf{Z}/\ell\mathbf{Z}}) = 0$。
\end{lemma}

\begin{proof}
将 $L = \bigcup L_i$ 写成其在 $K$ 上有限子扩张的并。
由 $H$ 的选取，$[L_i : K]$ 与 $\ell$ 互素。因此如引理
\ref{etale-cohomology-lemma-nonvanishing-inherited} 中那样，
$\Spec(L) = \lim_{i \in I} \Spec(L_i)$。
故可将 $K$ 替换为 $L$，并假设 $G$，即 $K$ 的绝对 Galois 群，
是预有限 pro-$\ell$ 群。

\medskip\noindent
假设 $H^n(\Spec(K), \underline{\mathbf{Z}/\ell\mathbf{Z}}) = 0$。
设 $\mathcal{F}$ 是一个 $\ell$-幂挠层，定义在 $\Spec(K)_\etale$ 上。
我们要证明 $H^n_\etale(\Spec(K), \mathcal{F}) = 0$。
由引理 \ref{etale-cohomology-lemma-equivalence-abelian-sheaves-point} 中给出的对应，
层 $\mathcal{F}$ 对应于一个 $\ell$-幂挠 $G$-模 $M$。
由连续性，任意有限个元素 $x_1, \ldots, x_m \in M$ 都被某个开子群 $U$ 固定。
设 $M'$ 是由 $x_1, \ldots, x_m$ 的轨道生成的模。
这是有限交换 $\ell$-群，因为每个 $x_i$ 都被 $\ell$ 的某个幂零化，
并且轨道有限。由于 $M$ 是这些子模 $M'$ 的滤过余极限，
可见 $\mathcal{F}$ 是相应子层 $\mathcal{F}' \subset \mathcal{F}$ 的滤过余极限。
对此余极限应用定理 \ref{etale-cohomology-theorem-colimit}，便归结到
$\mathcal{F}$ 是有限局部常值层的情形。

\medskip\noindent
设 $M$ 是有限交换 $\ell$-群，带有预有限 pro-$\ell$ 群 $G$ 的连续作用。
则存在 $G$-不变滤过
$$
0 = M_0 \subset M_1 \subset \ldots \subset M_r = M
$$
使得 $M_{i + 1}/M_i \cong \mathbf{Z}/\ell \mathbf{Z}$，且其上的
$G$-作用平凡（这是有限群表示论中的一个简单引理；此处以后补引文）。
因此相应的层 $\mathcal{F}$ 有滤过
$$
0 = \mathcal{F}_0 \subset \mathcal{F}_1 \subset \ldots \subset
\mathcal{F}_r = \mathcal{F}
$$
其逐次商同构于 $\underline{\mathbf{Z}/\ell \mathbf{Z}}$。
于是由归纳法和上同调长正合列得到结论。
\end{proof}

\begin{lemma}
\label{etale-cohomology-lemma-reduce-to-l-group-higher}
设 $\ell$ 是素数，$n$ 是整数且 $> 0$。设 $K$ 是域，
$G = Gal(K^{sep}/K)$，并设 $H \subset G$ 是极大 pro-$\ell$ 子群，
$L/K$ 是相应的域扩张。则
$H^q_\etale(\Spec(K),\mathcal{F}) = 0$ 对 $q \geq n$ 以及所有
$\ell$-挠层 $\mathcal{F}$ 成立，当且仅当
$H^n_\etale(\Spec(L), \underline{\mathbf{Z}/\ell\mathbf{Z}}) = 0$。
\end{lemma}

\begin{proof}
正向是显然的，故只需证明反向。对 $q$ 作归纳。
$q = n$ 的情形就是引理 \ref{etale-cohomology-lemma-reduce-to-l-group}。
现设 $\mathcal{F}$ 是一个 $\ell$-幂挠层，定义在 $\Spec(K)$ 上。
设 $f : \Spec(K^{sep}) \rightarrow \Spec(K)$ 是一个几何点的包含。
考虑正合列
$$
0 \rightarrow \mathcal{F} \xrightarrow{res}
f_* f^{-1} \mathcal{F} \rightarrow f_* f^{-1} \mathcal{F}/\mathcal{F} 
\rightarrow 0
$$
注意，$K^{sep}$ 可写成有限可分扩张的滤过余极限。因此 $f$
是一个有限 \'etale 态射有向系的极限。如引理
\ref{etale-cohomology-lemma-nonvanishing-inherited} 的证明所示，可知 $f$ 的高阶正像消失。
因而 $f_* f^{-1} \mathcal{F}$ 的高阶上同调可表示为几何点上的高阶上同调，
后者显然消失。于是，由于这里的一切仍是 $\ell$-挠的，
结合归纳假设与上同调长正合列，即可得到 $q + 1$ 的结论。
\end{proof}

\begin{proposition}
\label{etale-cohomology-proposition-serre-galois}
\begin{reference}
\cite[Chapter II, Section 3, Proposition 5]{SerreGaloisCohomology}
\end{reference}
设 $K$ 是域，其可分代数闭包为 $K^{sep}$。假设对任意有限扩张
$K'$（扩张自 $K$）都有 $\text{Br}(K') = 0$。则
\begin{enumerate}
\item $H^q(\text{Gal}(K^{sep}/K), (K^{sep})^*) = 0$
对所有 $q \geq 1$ 成立；
\item $H^q(\text{Gal}(K^{sep}/K), M) = 0$
对任意挠 $\text{Gal}(K^{sep}/K)$-模 $M$ 和任意 $q \geq 2$ 成立。
\end{enumerate}
\end{proposition}

\begin{proof}
置 $p = \text{char}(K)$。由引理 \ref{etale-cohomology-lemma-compare-cohomology-point}、
定理 \ref{etale-cohomology-theorem-brauer-delta} 和例 \ref{etale-cohomology-example-sheaves-point}，
本命题等价于证明：若
$H^2(\Spec(K'),\mathbf{G}_m|_{\Spec(K')_\etale}) = 0$
对所有有限扩张 $K'/K$ 都成立，则
\begin{itemize}
\item $H^q(\Spec(K),\mathbf{G}_m|_{\Spec(K)_\etale}) = 0$
对所有 $q \geq 1$ 成立；
\item $H^q(\Spec(K),\mathcal{F}) = 0$
对任意挠层 $\mathcal{F}$ 和任意 $q \geq 2$ 成立。
\end{itemize}
先证明第二部分。由于 $\mathcal{F}$ 是挠层，可以利用 $\ell$-准素分解，
以及上同调与余极限（即直和，见定理 \ref{etale-cohomology-theorem-colimit}）的相容性，
将问题归结为证明 $H^q(\Spec(K),\mathcal{F}) = 0$, $q \geq 2$，
其中对所有 $\ell$-幂挠层并对每个素数 $\ell$ 都要成立。
这样便可逐个分析各素数。

\medskip\noindent
假设 $\ell \neq p$。对任意扩张 $K'/K$，考虑 Kummer 序列
（引理 \ref{etale-cohomology-lemma-kummer-sequence}）
$$
0 \to
\mu_{\ell, \Spec{K'}} \to
\mathbf{G}_{m, \Spec{K'}} \xrightarrow{(\cdot)^{\ell}}
\mathbf{G}_{m, \Spec{K'}} \to 0
$$
按假设，$H^q(\Spec{K'},\mathbf{G}_m|_{\Spec(K')_\etale}) = 0$
在 $q = 2$ 时成立；由定理 \ref{etale-cohomology-theorem-picard-group}，在 $q = 1$ 时，
结合 $\Pic(K) = (0)$ 也成立。故由上同调长正合列可知，
$H^2(\Spec{K'}, \mu_\ell) = 0$ 对任意可分扩张 $K'/K$ 都成立。
现设 $H$ 是极大 pro-$\ell$ 子群，位于 $K$ 的绝对 Galois 群中，
$L$ 是相应扩张。可以把 $L$ 写成有限扩张的余极限；
对此余极限应用定理 \ref{etale-cohomology-theorem-colimit}，得到
$H^2(\Spec(L), \mu_\ell) = 0$。
此时 $\mu_\ell$ 必为常值层。否则就意味着存在一个次数与 $\ell$ 互素的
$L$ 的 Galois 扩张，这与 $L$ 的定义不符（该扩张即由添加
$\ell$ 次单位根到 $L$ 而得到）。
因此由引理 \ref{etale-cohomology-lemma-reduce-to-l-group-higher}，得到 $\ell \neq p$ 时的结论。

\medskip\noindent
现假设 $\ell = p$。考虑 Artin--Schrier 正合列
（第 \ref{etale-cohomology-section-artin-schreier} 节）
$$
0 \longrightarrow \underline{\mathbf{Z}/p\mathbf{Z}}_{\Spec{K}} \longrightarrow
\mathbf{G}_{a, \Spec{K}} \xrightarrow{F-1} \mathbf{G}_{a, \Spec{K}} 
\longrightarrow 0
$$
其中 $F - 1$ 是映射 $x \mapsto x^p - x$。由注
\ref{etale-cohomology-remark-special-case-fpqc-cohomology-quasi-coherent} 以及仿射概形结构层的
高阶上同调消失（Cohomology of Schemes，引理
\ref{coherent-lemma-quasi-coherent-affine-cohomology-zero}），
$\mathbf{G}_{a, \Spec{K}}$ 的高阶上同调消失。注意，这可应用于任意特征为
$p$ 的域；特别地，可应用于域扩张 $L$，该扩张由极大 pro-$p$ 子群 $H$
确定。由此得到
$H^n(\Spec{L}, \underline{\mathbf{Z}/p\mathbf{Z}}_{\Spec{L}}) = 0$
对 $n \geq 2$ 成立；
再由引理 \ref{etale-cohomology-lemma-reduce-to-l-group-higher} 即得 $\ell = p$ 时的结论。

\medskip\noindent
为完成证明，还需证明
$H^q(\text{Gal}(K^{sep}/K), (K^{sep})^*) = 0$
对所有 $q \geq 1$ 成立。
置 $G = \text{Gal}(K^{sep}/K)$，并将 $M = (K^{sep})^*$ 视为 $G$-模。
上文已经证明 $H^1(G, M) = 0$ 和 $H^2(G, M) = 0$。考虑正合列
$$
0 \to A \to M \to M \otimes \mathbf{Q} \to B \to 0
$$
，它是 $G$-模的正合列。由上文有 $H^i(G, A) = 0$ 和
$H^i(G, B) = 0$（$i > 1$），因为 $A$ 和 $B$ 是挠 $G$-模。
由引理 \ref{etale-cohomology-lemma-profinite-group-cohomology-torsion}，有
$H^i(G, M \otimes \mathbf{Q}) = 0$（$i > 0$）。不难验证，这蕴含
$H^i(G, M) = 0$ 对 $i \geq 3$ 也成立。
\end{proof}

%10.08.09

\begin{definition}
\label{etale-cohomology-definition-Cr}
称域 $K$ 为 {\it $C_r$ 域}，如果对任意 $0 < d^r < n$ 以及任意
$f \in K[T_1,
\ldots, T_n]$，若它是次数为 $d$ 的齐次多项式，则都存在
$\alpha = (\alpha_1,
\ldots, \alpha_n)$，其中 $\alpha_i \in K$ 不全为零，使得
$f(\alpha) = 0$。
这样的 $\alpha$ 称为 $f$ 的一个{\it 非平凡解}。
\end{definition}

\begin{example}
\label{etale-cohomology-example-algebraically-closed-field-Cr}
代数闭域是 $C_r$ 域。
\end{example}

\noindent
事实上有如下简单引理。

\begin{lemma}
\label{etale-cohomology-lemma-algebraically-closed-find-solutions}
设 $k$ 是代数闭域。设 $f_1, \ldots, f_s \in k[T_1, \ldots, T_n]$
是次数分别为 $d_1, \ldots, d_s$ 的齐次多项式，其中 $d_i
> 0$。
若 $s < n$，则 $f_1 = \ldots = f_s = 0$ 有一个公共非平凡解。
\end{lemma}

\begin{proof}
这由维数理论得出，例如将 Varieties，引理
\ref{varieties-lemma-intersection-in-affine-space} 应用 $s - 1$ 次。
\end{proof}

\noindent
下面的结果计算 $C_1$ 域的布劳尔群。

\begin{theorem}
\label{etale-cohomology-theorem-C1-brauer-group-zero}
设 $K$ 是 $C_1$ 域。则 $\text{Br}(K) = 0$。
\end{theorem}

\begin{proof}
设 $D$ 是中心为 $K$ 的有限维 $K$-除代数。前面已见
$$
D \otimes_K K^{sep} \cong \text{Mat}_d(K^{sep})
$$
在内同构的意义下唯一。因此行列式 $\det :
\text{Mat}_d(K^{sep}) \to K^{sep}$ 是 Galois 不变的，并下降为次数 $d$ 的齐次映射
$$
\det = N_\text{red} : D \longrightarrow K
$$
，称为{\it 约化范数}。由于 $K$ 是 $C_1$ 域，若 $d > 1$，则存在非零
$x \in D$，使得 $N_\text{red}(x) = 0$。这显然蕴含 $x$ 不可逆，产生矛盾。
故 $\text{Br}(K) = 0$。
\end{proof}

\begin{definition}
\label{etale-cohomology-definition-variety}
设 $k$ 是域。{\it 簇}是 $k$ 上有限型的分离整概形。
{\it 曲线}是维数为 $1$ 的簇。
\end{definition}

\begin{theorem}[Tsen's theorem]
\label{etale-cohomology-theorem-tsen}
维数为 $r$ 的、代数闭域 $k$ 上的簇的函数域是 $C_r$ 域。
\end{theorem}

\begin{proof}
对射影空间，可直接证明域 $k(x_1, \ldots, x_r)$ 是 $C_r$ 域（练习）。

\medskip\noindent
一般情形。不失一般性，可以假设 $X$ 是射影的。设
$f \in k(X)[T_1, \ldots, T_n]_d$ 且 $0 < d^r < n$。
设 $f$ 的系数属于 $\Gamma(X, \mathcal{O}_X(H))$，其中
$H \subset X$ 是某个丰沛对象。取
$\mathbf{\alpha} = (\alpha_1, \ldots, \alpha_n)$，其中
$\alpha_i \in \Gamma(X,
\mathcal{O}_X(eH))$。则 $f(\mathbf{\alpha}) \in \Gamma(X,
\mathcal{O}_X((de + 1)H))$。考虑方程组 $f(\mathbf{\alpha})
=0$。由渐近 Riemann--Roch
（Varieties，命题 \ref{varieties-proposition-asymptotic-riemann-roch}），
存在 $c > 0$，使得
\begin{itemize}
\item 变量数为 $n\dim_k \Gamma(X, \mathcal{O}_X(eH)) \sim n e^r c$；
\item 方程数为
$\dim_k \Gamma(X, \mathcal{O}_X((de + 1)H)) \sim (de + 1)^r c$。
\end{itemize}
由于 $n > d^r$，变量多于方程。方程是齐次的，故由引理
\ref{etale-cohomology-lemma-algebraically-closed-find-solutions} 存在一个解。
\end{proof}

\begin{lemma}
\label{etale-cohomology-lemma-curve-brauer-zero}
设 $C$ 是代数闭域 $k$ 上的曲线。则 $C$ 的函数域的布劳尔群为零：
$\text{Br}(k(C)) = 0$。
\end{lemma}

\begin{proof}
这直接由 Tsen 定理（定理 \ref{etale-cohomology-theorem-tsen}）和定理
\ref{etale-cohomology-theorem-C1-brauer-group-zero} 得出。
\end{proof}

\begin{lemma}
\label{etale-cohomology-lemma-cohomology-Gm-function-field-curve}
设 $k$ 是代数闭域，$K/k$ 是超越次数为 1 的域扩张。则对所有
$q \geq 1$，都有 $H_\etale^q(\Spec(K), \mathbf{G}_m) = 0$。
\end{lemma}

\begin{proof}
回顾一下，由引理 \ref{etale-cohomology-lemma-compare-cohomology-point}，
$H_\etale^q(\Spec(K), \mathbf{G}_m) = H^q(\text{Gal}(K^{sep}/K), (K^{sep})^*)$。
因此由命题 \ref{etale-cohomology-proposition-serre-galois}，只需证明：若 $K'/K$ 是有限域扩张，
则 $\text{Br}(K') = 0$。注意，$K' = \colim K''$，其中 $K''$ 遍历
$k$ 的有限生成子扩张，这些子扩张包含在 $K'$ 中且超越次数为 $1$。
又有 $\text{Br}(K') = \colim \text{Br}(K'')$，故问题归结到有限生成域扩张
$K''/k$，它的超越次数为 $1$。这样的域是 $k$ 上某条曲线的函数域，
所以由引理 \ref{etale-cohomology-lemma-curve-brauer-zero}，其布劳尔群平凡。
\end{proof}






\section{乘法群的高阶消失}
\label{etale-cohomology-section-higher-Gm}

\noindent
本节固定一个代数闭域 $k$ 和一条光滑曲线 $X$，后者定义在 $k$ 上。
以 $i_x : x \hookrightarrow X$ 表示 $X$ 的闭点的包含，以
$j : \eta \hookrightarrow X$ 表示泛点的包含。又以 $X_0$ 表示 $X$ 的闭点集。

\begin{theorem}[基本正合列]
\label{etale-cohomology-theorem-fundamental-exact-sequence}
$X$ 上有一个 \'etale 层的短正合列
$$
0 \longrightarrow
\mathbf{G}_{m, X} \longrightarrow
j_* \mathbf{G}_{m, \eta} \longrightarrow
\bigoplus\nolimits_{x \in X_0} {i_x}_* \underline{\mathbf{Z}}
\longrightarrow 0.
$$
\end{theorem}

\begin{proof}
设 $\varphi : U \to X$ 是 \'etale 态射。由 \'etale 态射的性质
（命题 \ref{etale-cohomology-proposition-etale-morphisms}），有 $U = \coprod_i U_i$，
其中每个 $U_i$ 都是一条映到 $X$ 的光滑曲线。$U$ 的上述序列是每个
$U_i$ 的相应序列之积，故只需处理 $U$ 连通、从而不可约的情形。
在此情形，有熟知的正合列
$$
1 \longrightarrow
\Gamma(U, \mathcal{O}_U^*) \longrightarrow
k(U)^* \longrightarrow \bigoplus\nolimits_{y \in U_0} \mathbf{Z}_y.
$$
这等价于序列
$$
0 \longrightarrow \Gamma(U, \mathcal{O}_U^*) \longrightarrow
\Gamma(\eta \times_X U, \mathcal{O}_{\eta \times_X U}^*) \longrightarrow
\bigoplus\nolimits_{x \in X_0}
\Gamma(x \times_X U, \underline{\mathbf{Z}})
$$
展开定义，这正是序列
$$
0 \longrightarrow \mathbf{G}_m(U) \longrightarrow
j_* \mathbf{G}_{m, \eta}(U) \longrightarrow
\left(\bigoplus\nolimits_{x \in X_0} {i_x}_* \underline{\mathbf{Z}}\right)(U).
$$
这定义了基本正合列中的映射，并表明该列除最后一步可能尚未验证外是正合的。
为看出满射性，回顾：若 $U$ 是非奇异曲线，$D$ 是 $U$ 上的除子，
则存在一个 Zariski 开覆盖 $\{U_j \to U\}$ 于 $U$ 上，使得
$D|_{U_j} = \text{div}(f_j)$，其中某个 $f_j \in k(U)^*$ 具有此性质。
\end{proof}

\begin{lemma}
\label{etale-cohomology-lemma-higher-direct-jstar-Gm}
对任意 $q \geq 1$，都有 $R^q j_*\mathbf{G}_{m, \eta} = 0$。
\end{lemma}

\begin{proof}
需证明 $(R^q j_*\mathbf{G}_{m, \eta})_{\bar x} = 0$ 对每个几何点
$\bar x$（它是 $X$ 的点）都成立。

\medskip\noindent
假设 $\bar x$ 位于闭点 $x$ 之上，该点属于 $X$。设 $\Spec(A)$ 是 $x$ 在 $X$
中的一个仿射开邻域，$K$ 是 $A$ 的分式域。则
$$
\Spec(\mathcal{O}^{sh}_{X, \bar x}) \times_X \eta =
\Spec(\mathcal{O}^{sh}_{X, \bar x} \otimes_A K).
$$
环 $\mathcal{O}^{sh}_{X, \bar x} \otimes_A K$ 是离散赋值环
$\mathcal{O}^{sh}_{X, \bar x}$ 的一个局部化，故它要么仍是
$\mathcal{O}^{sh}_{X, \bar x}$，要么是其分式域 $K^{sh}_{\bar x}$。
但由于某个局部一致化元被取逆，它必为后者。因此
$$
(R^q j_*\mathbf{G}_{m, \eta})_{(X, \bar x)} = H_\etale^q(\Spec
K^{sh}_{\bar x}, \mathbf{G}_m).
$$
现回顾 $\mathcal{O}^{sh}_{X, \bar x} =
\colim_{(U, \bar u) \to \bar x} \mathcal{O}(U) = \colim_{A \subset B} B$
，其中 $A \to B$ 是 \'etale 的。因此 $K^{sh}_{\bar x}$ 是
$K = k(X)$ 的代数扩张，可以应用引理
\ref{etale-cohomology-lemma-cohomology-Gm-function-field-curve} 得到消失结论。

\medskip\noindent
假设 $\bar x = \bar \eta$ 位于泛点 $\eta$ 之上，而后者属于 $X$
（事实上这一情形是多余的）。此时
$\mathcal{O}^{sh}_{X, \bar \eta} = \kappa(\eta)^{sep}$，故
\begin{eqnarray*}
(R^q j_*\mathbf{G}_{m, \eta})_{\bar \eta}
& = &
H_\etale^q(\Spec(\kappa(\eta)^{sep}) \times_X \eta, \mathbf{G}_m) \\
& = & H_\etale^q (\Spec(\kappa(\eta)^{sep}), \mathbf{G}_m) \\
& = & 0 \ \ \text{ 对 } q \geq 1
\end{eqnarray*}
，因为相应的 Galois 群是平凡的。
\end{proof}

\begin{lemma}
\label{etale-cohomology-lemma-cohomology-jstar-Gm}
对所有 $p \geq 1$，都有 $H_\etale^p(X, j_*\mathbf{G}_{m, \eta}) = 0$。
\end{lemma}

\begin{proof}
Leray 谱序列为
$$
E_2^{p, q} = H_\etale^p(X, R^qj_*\mathbf{G}_{m, \eta}) \Rightarrow
H_\etale^{p+q}(\eta, \mathbf{G}_{m, \eta}),
$$
由引理 \ref{etale-cohomology-lemma-cohomology-Gm-function-field-curve}，右端在
$p + q \geq 1$ 时消失。由引理 \ref{etale-cohomology-lemma-higher-direct-jstar-Gm}，
左端在 $q \geq 1$ 时消失。故得到所需的消失结论。
\end{proof}

\begin{lemma}
\label{etale-cohomology-lemma-cohomology-istar-Z}
对所有 $q \geq 1$，都有 $H_\etale^q(X, \bigoplus_{x \in X_0} {i_x}_*
\underline{\mathbf{Z}}) = 0$。
\end{lemma}

\begin{proof}
由于 $X$ 拟紧拟分离，上同调与余极限可交换，故只需证明
$H_\etale^q(X, {i_x}_*
\underline{\mathbf{Z}})$ 消失。但闭点的包含
$i_x$ 是有限的，故由命题 \ref{etale-cohomology-proposition-finite-higher-direct-image-zero}，
有 $R^p {i_x}_* \underline{\mathbf{Z}} = 0$ 对所有 $p \geq 1$ 成立。
应用 Leray 谱序列可见
$H_\etale^q(X, {i_x}_* \underline{\mathbf{Z}}) =
H_\etale^q(x, \underline{\mathbf{Z}})$.
最后，由于 $x$ 是一个代数闭域的谱，$x$ 上所有高阶上同调都消失。
\end{proof}

\noindent
由以上一系列引理，得到如下结果。

\begin{theorem}
\label{etale-cohomology-theorem-vanishing-cohomology-Gm-curve}
设 $X$ 是代数闭域上的光滑曲线。则
$$
H_\etale^q(X, \mathbf{G}_m) = 0 \ \ \text{ 对所有 } q \geq 2.
$$
\end{theorem}

\begin{proof}
见上面的讨论。
\end{proof}

\noindent
还得到上同调长正合列
$$
0 \to
H_\etale^0(X, \mathbf{G}_m) \to
H_\etale^0(X, j_*\mathbf{G}_{m\eta}) \to
H_\etale^0(X, \bigoplus {i_x}_*\underline{\mathbf{Z}}) \to
H_\etale^1(X, \mathbf{G}_m) \to 0
$$
不过这就是熟知的序列
$$
0 \to H_{Zar}^0(X, \mathcal{O}_X^*) \to k(X)^* \to \text{Div}(X)
\to \Pic(X) \to 0.
$$





\section{曲线的 Picard 群}
\label{etale-cohomology-section-pic-curves}

\noindent
下一步利用 Kummer 序列推出曲线以有限群为系数的上同调群的一些信息。
为在长正合列中得到消失结论，先回顾关于 Picard 群的一些事实。

\medskip\noindent
设 $X$ 是代数闭域 $k$ 上的光滑射影曲线。以
$g = \dim_k H^1(X, \mathcal{O}_X)$ 表示 $X$ 的亏格。存在短正合列
$$
0 \to \Pic^0(X) \to \Pic(X) \xrightarrow{\deg} \mathbf{Z} \to 0.
$$
交换群 $\Pic^0(X)$ 可与 $\Pic^0(X) = \underline{\Picardfunctor}^0_{X/k}(k)$
等同，即它是 $k$-值点，所对应的阿贝尔簇
$\underline{\Picardfunctor}^0_{X/k}$ 定义在 $k$ 上且维数为 $g$。
因此，若 $n \in k^*$，则
$\Pic^0(X)[n] \cong (\mathbf{Z}/n\mathbf{Z})^{2g}$
作为交换群同构。见 Picard Schemes of Curves，第
\ref{pic-section-picard-curve} 节，以及 Groupoids，第
\ref{groupoids-section-abelian-varieties} 节。这个关键事实，即代数闭域上
光滑射影曲线的 Picard 群中挠元的描述，看来没有初等证明。

\begin{lemma}
\label{etale-cohomology-lemma-cohomology-smooth-projective-curve}
设 $X$ 是亏格为 $g$ 的光滑射影曲线，定义在代数闭域 $k$ 上，并设
$n \geq 1$ 在 $k$ 中可逆。则有典范等同
$$
H_\etale^q(X, \mu_n) =
\left\{
\begin{matrix}
\mu_n(k) & \text{ 若 }q = 0, \\
\Pic^0(X)[n] & \text{ 若 }q = 1, \\
\mathbf{Z}/n\mathbf{Z} & \text{ 若 }q = 2, \\
0 & \text{ 若 }q \geq 3.
\end{matrix}
\right.
$$
由于 $\mu_n \cong \underline{\mathbf{Z}/n\mathbf{Z}}$，这给出
（非典范的）等同
$$
H_\etale^q(X, \underline{\mathbf{Z}/n\mathbf{Z}}) \cong
\left\{
\begin{matrix}
\mathbf{Z}/n\mathbf{Z} & \text{ 若 }q = 0, \\
(\mathbf{Z}/n\mathbf{Z})^{2g} & \text{ 若 }q = 1, \\
\mathbf{Z}/n\mathbf{Z} & \text{ 若 }q = 2, \\
0 & \text{ 若 }q \geq 3.
\end{matrix}
\right.
$$
\end{lemma}

\begin{proof}
定理 \ref{etale-cohomology-theorem-picard-group} 和
\ref{etale-cohomology-theorem-vanishing-cohomology-Gm-curve} 确定了
$\mathbf{G}_m$ 在 $X$ 上的 \'etale 上同调，其表述使用 $X$ 的 Picard 群。
于是 Kummer 序列
$0\to \mu_{n, X} \to \mathbf{G}_{m, X}
\to \mathbf{G}_{m, X}\to 0$（引理 \ref{etale-cohomology-lemma-kummer-sequence}）
给出上同调长正合列
$$
\xymatrix{
0 \ar[r] & \mu_n(k) \ar[r] &
k^* \ar[r]^{(\cdot)^n} &
k^* \ar@(rd, ul)[rdllllr] \\
& H_\etale^1(X, \mu_n) \ar[r] &
\Pic(X) \ar[r]^{(\cdot)^n} &
\Pic(X) \ar@(rd, ul)[rdllllr] \\
& H_\etale^2(X, \mu_n) \ar[r] & 0 \ar[r] & 0 \ldots
}
$$
这里的 $n$ 次幂映射 $k^* \to k^*$ 是满射，因为 $k$ 是代数闭的。
因此需要计算映射 $n : \Pic(X) \to \Pic(X)$ 的核与余核。
考虑如下各行正合的交换图
$$
\xymatrix{
0 \ar[r] &
\Pic^0(X) \ar[r] \ar@{>>}[d]^{(\cdot)^n} &
\Pic(X) \ar[r]_-\deg \ar[d]^{(\cdot)^n} &
\mathbf{Z} \ar[r] \ar@{^{(}->}[d]^n & 0 \\
0 \ar[r] &
\Pic^0(X) \ar[r] &
\Pic(X) \ar[r]^-\deg &
\mathbf{Z} \ar[r] & 0
}
$$
群 $\Pic^0(X)$ 是 $k$-点，属于群概形
$\underline{\Picardfunctor}^0_{X/k}$，
见 Picard Schemes of Curves，引理 \ref{pic-lemma-picard-pieces}。
同一引理还说明 $\underline{\Picardfunctor}^0_{X/k}$ 是一个
$g$ 维阿贝尔簇，定义在 $k$ 上；这里的定义见 Groupoids，定义
\ref{groupoids-definition-abelian-variety}。故由 Groupoids，命题
\ref{groupoids-proposition-review-abelian-varieties}，左边的竖直映射是满射。
应用蛇引理便得到本引理所述的典范等同。

\medskip\noindent
为得到本引理中的非典范等同，需要证明
$n : \Pic^0(X) \to \Pic^0(X)$ 的核同构于
$(\mathbf{Z}/n\mathbf{Z})^{\oplus 2g}$。这也是 Groupoids，命题
\ref{groupoids-proposition-review-abelian-varieties} 的一部分。
\end{proof}

\begin{lemma}
\label{etale-cohomology-lemma-pullback-on-h2-curve}
设 $\pi : X \to Y$ 是代数闭域 $k$ 上光滑射影曲线之间的非常值态射，
并设 $n \geq 1$ 在 $k$ 中可逆。映射
$$
\pi^* : H^2_\etale(Y, \mu_n) \longrightarrow H^2_\etale(X, \mu_n)
$$
由乘以 $\pi$ 的次数给出。
\end{lemma}

\begin{proof}
注意，该陈述有意义，因为在引理
\ref{etale-cohomology-lemma-cohomology-smooth-projective-curve} 中，我们已经把两个上同调群
$ H^2_\etale(Y, \mu_n)$ 和 $H^2_\etale(X, \mu_n)$ 都与
$\mathbf{Z}/n\mathbf{Z}$ 等同。事实上，若 $\mathcal{L}$ 是次数为
$1$ 的线丛，定义在 $Y$ 上，其类为
$[\mathcal{L}] \in H^1_\etale(Y, \mathbf{G}_m)$，
则 $[\mathcal{L}]$ 的余边界是 $H^2_\etale(Y, \mu_n)$ 的生成元。
这里的余边界是与 Kummer 序列相伴的上同调长正合列中的余边界。
因此，本引理的结果由如下事实得出：线丛 $\pi^*\mathcal{L}$ 在
$X$ 上的次数为 $\deg(\pi)$。省略若干细节。
\end{proof}

\begin{lemma}
\label{etale-cohomology-lemma-vanishing-cohomology-mu-smooth-curve}
设 $X$ 是代数闭域 $k$ 上的仿射光滑曲线，且 $n \in k^*$。
设 $X \subset \overline{X}$ 是一个光滑射影紧化
（Varieties，注 \ref{varieties-remark-smooth-projective-compactification}）。
设 $g$ 是 $\overline{X}$ 的亏格，$r$ 是 $\overline{X} \setminus X$
中的点数。则
\begin{enumerate}
\item $H_\etale^0(X, \mu_n) = \mu_n(k)$;
\item $H_\etale^1(X, \mu_n) \cong (\mathbf{Z}/n\mathbf{Z})^{2g+r-1}$；
\item $H_\etale^q(X, \mu_n) = 0$ 对所有 $q \geq 2$ 成立。
\end{enumerate}
\end{lemma}

\begin{proof}
写成 $X = \overline{X} - \{ x_1, \ldots, x_r\}$。则
$\Pic(X) = \Pic(\overline{X})/ R$，其中 $R$ 是由
$\mathcal{O}_{\overline{X}}(x_i)$（$1 \leq i \leq r$）生成的子群。
由于 $r \geq 1$，可见 $\Pic^0(\overline{X}) \to \Pic(X)$ 是满射，
故 $\Pic(X)$ 可除（见引理
\ref{etale-cohomology-lemma-cohomology-smooth-projective-curve} 的证明中的讨论）。
应用 Kummer 序列得到 (1) 和 (3)。对 (2)，使用
\begin{align*}
H_\etale^1(X, \mu_n)
& =
\{(\mathcal L, \alpha)\}/\cong \\
& =
(\{(\bar{\mathcal L}, D, \bar \alpha)\}/\cong)/\tilde{R}
\end{align*}
第一个等号和记号见引理 \ref{etale-cohomology-lemma-describe-h1-mun}。
第二行中的三元组 $(\bar{\mathcal L}, D, \bar \alpha)$ 由以下数据组成：
可逆模 $\bar{\mathcal L}$，定义在 $\overline{X}$ 上且次数为 $0$；
除子 $D$，定义在 $\overline{X}$ 上且次数为 $0$，其支集包含于
$\left\{x_1, \ldots, x_r\right\}$；以及同构
$\bar{\alpha}: \bar{\mathcal L}^{\otimes n} \cong \mathcal{O}_{\bar{X}}(D)$。
三元组的同构按引理 \ref{etale-cohomology-lemma-describe-h1-mun} 之前讨论中的方式定义。
此外，$\tilde{R}$ 是由如下形式的三元组生成的子群：
$(\mathcal{O}_{\overline{X}}(D'), n D', 1^{\otimes n})$
，其中 $D'$ 的支集包含于 $\left\{x_1, \ldots, x_r\right\}$ 且次数为 0。
上述讨论表明，第二个等号由如下映射给出：它把三元组
$(\bar{\mathcal L},\ D,\ \bar \alpha)$ 送到对
$(\bar{\mathcal L}|_X, \bar \alpha|_X)$。于是得到正合列
$$
0 \longrightarrow
H_\etale^1(\overline{X}, \mu_n) \longrightarrow
H_\etale^1(X, \mu_n) \longrightarrow
\bigoplus_{i = 1}^r \mathbf{Z}/n\mathbf{Z}
\xrightarrow{\ \sum\ }
\mathbf{Z}/n\mathbf{Z} \longrightarrow 0
$$
其中中间的映射把三元组 $(\bar{ \mathcal L}, D, \bar \alpha)$ 的类
（这里 $D = \sum_{i = 1}^r a_i (x_i)$）送到 $r$-元组
$(a_i)_{i = 1}^r$。现只需利用引理
\ref{etale-cohomology-lemma-cohomology-smooth-projective-curve} 计算秩。
\end{proof}

\begin{remark}
\label{etale-cohomology-remark-natural-proof}
证明前一个推论的“自然”方法是把 $X$ 从 $\bar
X$ 中切除。
这是可行的，只是我们尚未建立该理论。
\end{remark}

\begin{remark}
\label{etale-cohomology-remark-normalize-H1-Gm}
设 $k$ 是代数闭域。设 $n$ 是与 $k$ 的特征互素的整数。回顾
$$
\mathbf{G}_{m, k} = \mathbf{A}^1_k \setminus \{0\} =
\mathbf{P}^1_k \setminus \{0, \infty\}
$$
我们断言存在典范同构
$$
H^1_\etale(\mathbf{G}_{m, k}, \mu_n) = \mathbf{Z}/n\mathbf{Z}
$$
这是什么意思？这意味着存在元素 $1_k$ 于
$H^1_\etale(\mathbf{G}_{m, k}, \mu_n)$ 中，使得对每个态射
$\Spec(k') \to \Spec(k)$，映射 $\mathbf{G}_{m, k'} \to \mathbf{G}_{m, k}$
在 \'etale 上同调上诱导的拉回映射把 $1_k$ 送到 $1_{k'}$。
（特别地，此元素被 $k$ 的所有自同构固定。）为看出这一点，考虑
$\mu_{n, \mathbf{Z}}$-挠子
$\mathbf{G}_{m, \mathbf{Z}} \to \mathbf{G}_{m, \mathbf{Z}}$，
$x \mapsto x^n$。由挠子与一阶上同调的等同，拉回此挠子便得到我们的
典范元素 $1_k$。反向扭曲可见存在典范等同
$$
H^1_\etale(\mathbf{G}_{m, k}, \mathbf{Z}/n\mathbf{Z}) =
\Hom(\mu_n(k), \mathbf{Z}/n\mathbf{Z}),
$$
即这些同构关于代数闭域之间的映射相容，特别地，关于 $k$ 的自同构相容。
\end{remark}










\section{零延拓}
\label{etale-cohomology-section-extension-by-zero}

\noindent
Modules on Sites，第 \ref{sites-modules-section-localize} 节中的一般材料
使我们可以作如下定义。

\begin{definition}
\label{etale-cohomology-definition-extension-zero}
设 $j : U \to X$ 是概形的 \'etale 态射。
\begin{enumerate}
\item 限制函子
$j^{-1} : \Sh(X_\etale) \to \Sh(U_\etale)$
有左伴随
$j_!^{Sh} : \Sh(U_\etale) \to \Sh(X_\etale)$.
\item 限制函子
$j^{-1} : \textit{Ab}(X_\etale) \to \textit{Ab}(U_\etale)$
有左伴随，记为
$j_! : \textit{Ab}(U_\etale) \to \textit{Ab}(X_\etale)$
，称为{\it 零延拓}。
\item 设 $\Lambda$ 是环。限制函子
$j^{-1} : \textit{Mod}(X_\etale, \Lambda) \to
\textit{Mod}(U_\etale, \Lambda)$
有左伴随，记为
$j_! : \textit{Mod}(U_\etale, \Lambda) \to
\textit{Mod}(X_\etale, \Lambda)$
，称为{\it 零延拓}。
\end{enumerate}
\end{definition}

\noindent
若 $\mathcal{F}$ 是 $X_\etale$ 上的交换群层，则一般有
$j_!\mathcal{F} \not = j_!^{Sh}\mathcal{F}$。另一方面，函子 $j_!$ 对
$\Lambda$-模层的作用与底层交换群层上的 $j_!$ 一致
（Modules on Sites，注 \ref{sites-modules-remark-localize-shriek-equal}）。
函子 $j_!$ 由函子性同构
$$
\Hom_X(j_!\mathcal{F}, \mathcal{G}) = \Hom_U(\mathcal{F}, j^{-1}\mathcal{G})
$$
刻画，其中对所有 $\mathcal{F} \in \textit{Ab}(U_\etale)$ 和
$\mathcal{G} \in \textit{Ab}(X_\etale)$ 均成立。$\Lambda$-模层的情形类似。

\medskip\noindent
为更明确地描述定义 \ref{etale-cohomology-definition-extension-zero} 中的函子，回顾：
$j^{-1}$ 就是沿函子 $U_\etale \to X_\etale$ 的限制。换言之，
$j^{-1}\mathcal{G}(U') = \mathcal{G}(U')$ 对 $U'$ 成立，其中它在
$U$ 上是 \'etale 的。另一方面，对
$\mathcal{F} \in \textit{Ab}(U_\etale)$，考虑预层
\begin{equation}
\label{etale-cohomology-equation-j-p-shriek}
j_{p!}\mathcal{F} : X_\etale \longrightarrow \textit{Ab},
\quad
V \longmapsto \bigoplus\nolimits_{V \to U} \mathcal{F}(V \to U)
\end{equation}
则 $j_!\mathcal{F}$ 是 $j_{p!}\mathcal{F}$ 的层化。
这一点在 Modules on Sites，引理
\ref{sites-modules-lemma-extension-by-zero} 中证明；更一般地，见
Modules on Sites，第 \ref{sites-modules-section-localize} 节和第
\ref{sites-modules-section-exactness-lower-shriek} 节中的讨论。

\begin{exercise}
\label{etale-cohomology-exercise-jshriek-direct}
直接证明：函子 $j_!$（它定义为 (\ref{etale-cohomology-equation-j-p-shriek}) 中函子
$j_{p!}$ 的层化）是 $j^{-1}$ 的左伴随。
\end{exercise}

\begin{proposition}
\label{etale-cohomology-proposition-describe-jshriek}
设 $j : U \to X$ 是概形的 \'etale 态射。设
$\mathcal{F}$ 属于 $\textit{Ab}(U_\etale)$。若
$\overline{x} : \Spec(k) \to X$ 是 $X$ 的几何点，则
$$
(j_!\mathcal{F})_{\overline{x}} =
\bigoplus\nolimits_{\overline{u} : \Spec(k) \to U,\ j(\overline{u}) =
\overline{x}} \mathcal{F}_{\bar{u}}.
$$
特别地，$j_!$ 是正合函子。
\end{proposition}

\begin{proof}
$j_!$ 的正合性是非常一般的，见 Modules on Sites，引理
\ref{sites-modules-lemma-extension-by-zero-exact}。当然，它也可由茎的描述得出。
茎的公式由 Modules on Sites，引理
\ref{sites-modules-lemma-stalk-j-shriek}，以及用几何点描述小 \'etale
网站的点这一结果得出；后者见引理 \ref{etale-cohomology-lemma-points-small-etale-site}。

\medskip\noindent
为以后使用，注意同构
\begin{align*}
(j_!\mathcal{F})_{\overline{x}}
& =
(j_{p!}\mathcal{F})_{\overline{x}} \\
& =
\colim_{(V, \overline{v})} j_{p!}\mathcal{F}(V) \\
& =
\colim_{(V, \overline{v})}
\bigoplus\nolimits_{\varphi : V \to U}
\mathcal{F}(V \xrightarrow{\varphi} U) \\
& \to
\bigoplus\nolimits_{\overline{u} : \Spec(k) \to U,\ j(\overline{u}) =
\overline{x}} \mathcal{F}_{\bar{u}}.
\end{align*}
在 Modules on Sites，引理 \ref{sites-modules-lemma-stalk-j-shriek}
中构造；它把 $(V, \overline{v}, \varphi, s)$ 送到 $s$ 在
$\mathcal{F}$ 于 $\overline{u} = \varphi(\overline{v})$ 处的茎中的类。
\end{proof}

\begin{lemma}
\label{etale-cohomology-lemma-jshriek-open}
设 $j : U \to X$ 是概形的开浸入。对任意交换群层 $\mathcal{F}$ 于
$U_\etale$ 上，伴随映射 $j^{-1}j_*\mathcal{F} \to \mathcal{F}$ 和
$\mathcal{F} \to j^{-1}j_!\mathcal{F}$ 都是同构。事实上，
$j_!\mathcal{F}$ 是 $X_\etale$ 上唯一的交换群层，其到 $U$ 的限制为
$\mathcal{F}$，并且它在 $X \setminus U$ 的几何点处的茎为零。
\end{lemma}

\begin{proof}
我们鼓励读者展开定义证明第一项陈述，但这里仅使用它是非常一般的
Modules on Sites，引理 \ref{sites-modules-lemma-restrict-back} 的特殊情形。
对第二项陈述，注意：若 $\mathcal{G}$ 是 $X_\etale$ 上的交换群层，
其到 $U$ 的限制为 $\mathcal{F}$，则由伴随性得到映射
$j_!\mathcal{F} \to \mathcal{G}$。由命题
\ref{etale-cohomology-proposition-describe-jshriek}，此映射在 $U$ 的几何点的茎上是同构。
因此，若 $\mathcal{G}$ 在 $X \setminus U$ 的几何点处的茎消失，
则由定理 \ref{etale-cohomology-theorem-exactness-stalks}，
$j_!\mathcal{F} \to \mathcal{G}$ 是同构。
\end{proof}

\begin{lemma}[零延拓与基变换可交换]
\label{etale-cohomology-lemma-shriek-base-change}
设 $f: Y \to X$ 是概形态射。设 $j: V \to X$ 是 \'etale 态射。
考虑纤维积
$$
\xymatrix{
V' = Y \times_X V \ar[d]_{f'} \ar[r]_-{j'} & Y \ar[d]^f \\
V \ar[r]^j & X
}
$$
则在交换群层和模层上都有 $j'_! f'^{-1} = f^{-1} j_!$。
\end{lemma}

\begin{proof}
这是因为 $j'_! f'^{-1}$ 是 $f'_* (j')^{-1}$ 的左伴随，而
$f^{-1} j_!$ 是 $j^{-1}f_*$ 的左伴随。又有
$f'_* (j')^{-1} = j^{-1}f_*$，因为 $f_*$ 按构造与 \'etale 局部化可交换。
事实上，本引理在网站态射的框架下非常一般地成立，见
Modules on Sites，引理
\ref{sites-modules-lemma-localize-morphism-ringed-sites}。
\end{proof}

\begin{lemma}
\label{etale-cohomology-lemma-shriek-into-star-separated-etale}
设 $j : U \to X$ 是分离 \'etale 态射。则有函子性单射
$j_!\mathcal{F} \to j_*\mathcal{F}$，定义在交换群层和 $\Lambda$-模层上。
\end{lemma}

\begin{proof}
我们在交换群层的情形证明此结论。对任意交换群层，构造典范映射
$$
j_{p!}\mathcal{F} \to j_*\mathcal{F}
$$
，它是 $X_\etale$ 上交换群预层的映射；这里 $\mathcal{F}$ 定义在
$U_\etale$ 上，$j_{p!}$ 如 (\ref{etale-cohomology-equation-j-p-shriek}) 中所定义。
将此映射层化便得到所需映射 $j_!\mathcal{F} \to j_*\mathcal{F}$。
在 \'etale 态射 $V \to X$ 上计算两端，得到
$$
j_{p!}\mathcal{F}(V) =
\bigoplus\nolimits_{\varphi : V \to U} \mathcal{F}(V \xrightarrow{\varphi} U)
\quad\text{且}\quad
j_*\mathcal{F}(V) = \mathcal{F}(V \times_X U)
$$
对每个 $\varphi$，有既开又闭的浸入
$$
\Gamma_\varphi = (1, \varphi) : V \longrightarrow V \times_X U
$$
，它定义在 $U$ 上。它是开的，因为它是 $U$ 上 \'etale 概形之间的态射；
它是闭的，因为它是一个 $V$ 上分离概形的截面（Schemes，引理
\ref{schemes-lemma-section-immersion}）。因此，对截面
$s_\varphi \in \mathcal{F}(V \xrightarrow{\varphi} U)$，存在唯一截面
$s'_\varphi$ 于 $\mathcal{F}(V \times_X U)$ 中，使其拉回为
$s_\varphi$，所沿映射为 $\Gamma_\varphi$，且其到 $\Gamma_\varphi$
的像之补集的限制为零。

\medskip\noindent
为证明映射是单射，假设上述公式中的元素
$\sum_{i = 1, \ldots, n} s_{\varphi_i}$ 属于 $j_{p!}\mathcal{F}(V)$，
且它在 $j_*\mathcal{F}(V)$ 中的像为零。任务是证明
$\sum_{i = 1, \ldots, n} s_{\varphi_i}$ 到 $V$ 的某个 \'etale 覆盖的
各成员上的限制为零。考察所有两两等化子（它们在 $V$ 中既开又闭），
这些等化子来自态射 $\varphi_i : V \to U$；在 $V$ 上局部工作，可以假设态射
$\Gamma_{\varphi_1}, \ldots, \Gamma_{\varphi_n}$ 的像两两不交。
由于假设 $\sum_{i = 1, \ldots, n} s'_{\varphi_i} = 0$，立即得到
$s'_{\varphi_i} = 0$ 对每个 $i$ 都成立（因为这些截面的支集两两不交），
进而 $s_{\varphi_i} = 0$ 对所有 $i$ 都成立，如所需。
\end{proof}

\begin{lemma}
\label{etale-cohomology-lemma-shriek-equals-star-finite-etale}
设 $j : U \to X$ 是有限 \'etale 态射。则映射
$j_! \to j_*$（见引理 \ref{etale-cohomology-lemma-shriek-into-star-separated-etale}）
在交换群层和 $\Lambda$-模层上都是同构。
\end{lemma}

\begin{proof}
只需验证 $j_!\mathcal{F} \to j_*\mathcal{F}$ 是同构，且可在
$X$ 上 \'etale 局部地验证。因此可以假设
$U \to X$ 是有限个同构的不交并，见 \'Etale Morphisms，引理
\ref{etale-lemma-finite-etale-etale-local}。省略这一情形的证明。
\end{proof}

\begin{lemma}
\label{etale-cohomology-lemma-ses-associated-to-open}
设 $X$ 是概形。设 $Z \subset X$ 是闭子概形，$U \subset X$ 是其补集。
以 $i : Z \to X$ 和 $j : U \to X$ 表示包含态射。对每个交换群层
$\mathcal{F}$ 于 $X_\etale$ 上，有典范短正合列
$$
0 \to j_!j^{-1}\mathcal{F} \to \mathcal{F} \to i_*i^{-1}\mathcal{F} \to 0
$$
于 $X_\etale$ 上。
\end{lemma}

\begin{proof}
这些映射由所涉函子的伴随性质得到。对几何点 $\overline{x}$ 于
$X$ 中，或者 $\overline{x} \in U$；此时左侧映射在茎上是同构，
而 $i_*i^{-1}\mathcal{F}$ 的茎为零。或者 $\overline{x} \in Z$；
此时右侧映射在茎上是同构，而 $j_!j^{-1}\mathcal{F}$ 的茎为零。
这里使用了引理 \ref{etale-cohomology-lemma-stalk-pushforward-closed-immersion} 和命题
\ref{etale-cohomology-proposition-describe-jshriek} 中对茎的描述。
\end{proof}

\begin{lemma}
\label{etale-cohomology-lemma-compatible-shriek-push-finite}
考虑概形的笛卡儿图
$$
\xymatrix{
U \ar[d]_g \ar[r]_{j'} & X \ar[d]^f \\
V \ar[r]^j & Y
}
$$
，其中 $f$ 是有限态射，$j$ 是开浸入。则有
$f_* \circ j'_! = j_! \circ g_*$，作为函子
$\textit{Ab}(U_\etale) \to \textit{Ab}(Y_\etale)$。
\end{lemma}

\begin{proof}
设 $\mathcal{F}$ 是 $\textit{Ab}(U_\etale)$ 的对象。设
$\overline{y}$ 是 $Y$ 的几何点，且不包含在开集 $V$ 中。则
$$
(f_*j'_!\mathcal{F})_{\overline{y}} =
\bigoplus\nolimits_{\overline{x},\ f(\overline{x}) = \overline{y}}
(j'_!\mathcal{F})_{\overline{x}} = 0
$$
由命题 \ref{etale-cohomology-proposition-finite-higher-direct-image-zero}，并且由引理
\ref{etale-cohomology-lemma-jshriek-open}，$j'_!\mathcal{F}$ 在
$\overline{x} \not \in U$ 处的茎为零。另一方面，有
$$
j^{-1}f_*j'_!\mathcal{F} =
g_*(j')^{-1}j'_!\mathcal{F} =
g_*\mathcal{F}
$$
，这由引理 \ref{etale-cohomology-lemma-finite-pushforward-commutes-with-base-change} 和
\ref{etale-cohomology-lemma-jshriek-open} 得出。因此由引理 \ref{etale-cohomology-lemma-jshriek-open}
中对 $j_!$ 的刻画，可知 $f_*j'_!\mathcal{F} = j_!g_*\mathcal{F}$。
省略验证此等同关于 $\mathcal{F}$ 具有函子性。
\end{proof}







\section{可构造层}
\label{etale-cohomology-section-constructible}

\noindent
设 $X$ 为概形。$X$ 的一个{\it 可构造局部闭子概形}，是指一个局部闭子概形
$T \subset X$，使得 $T$ 的底拓扑空间是 $X$ 的可构造子集。
若 $T, T' \subset X$ 是具有同一底拓扑空间的局部闭子概形，则由 \'etale
位点的拓扑不变性（定理 \ref{etale-cohomology-theorem-topological-invariance}），有
$T_\etale \cong T'_\etale$。因此，在以下定义中可以假定所考虑的局部闭
子概形都是既约的。

\begin{definition}
\label{etale-cohomology-definition-constructible}
设 $X$ 为概形。
\begin{enumerate}
\item 若 $X_\etale$ 上的集合层满足如下条件，则称它是{\it 可构造的}：
对每个仿射开集 $U \subset X$，都存在将 $U$ 有限分解为可构造局部闭
子概形的分解 $U = \coprod_i U_i$，使得 $\mathcal{F}|_{U_i}$ 对所有 $i$
都是有限局部常值的。
\item 若 $X_\etale$ 上的阿贝尔群层满足如下条件，则称它是{\it 可构造的}：
对每个仿射开集 $U \subset X$，都存在将 $U$ 有限分解为可构造局部闭
子概形的分解 $U = \coprod_i U_i$，使得 $\mathcal{F}|_{U_i}$ 对所有 $i$
都是有限局部常值的。
\item 设 $\Lambda$ 为诺特环。若一个 $\Lambda$-模层在 $X_\etale$ 上满足如下条件，
则称它是{\it 可构造的}：对每个仿射开集 $U \subset X$，都存在将 $U$
有限分解为可构造局部闭子概形的分解 $U = \coprod_i U_i$，使得
$\mathcal{F}|_{U_i}$ 对所有 $i$ 都是有限型且局部常值的。
\end{enumerate}
\end{definition}

\noindent
这看来是公认的定义。另有一种更便于推广到概形 \'etale 位点以外情形的
定义方式：从
$h_U$, $j_{U!}\mathbf{Z}/n\mathbf{Z}$ 以及 $j_{U!}\underline{\Lambda}$
出发，其中 $U$ 遍历 $X_\etale$ 的所有拟紧拟分离对象，再取
$\Sh(X_\etale)$、$\textit{Ab}(X_\etale)$ 以及
$\textit{Mod}(X_\etale, \underline{\Lambda})$ 中包含这些对象且对有限极限和
有限余极限封闭的最小满子范畴。由引理 \ref{etale-cohomology-lemma-constructible-abelian}
以及引理 \ref{etale-cohomology-lemma-category-constructible-sets},
\ref{etale-cohomology-lemma-category-constructible-abelian} 和
\ref{etale-cohomology-lemma-category-constructible-modules}
可知，当 $X$ 拟紧拟分离时，这种做法给出同一范畴；但一般而言，它们给出的
范畴并不相同。

\medskip\noindent
概形的一个由局部闭子概形组成的不交并分解 $U = \coprod U_i$，称为 $U$ 的一个
{\it 分划}（参见拓扑学，第 \ref{topology-section-stratifications} 节）。

\begin{lemma}
\label{etale-cohomology-lemma-constructible-quasi-compact-quasi-separated}
设 $X$ 为拟紧拟分离概形，$\mathcal{F}$ 为 $X_\etale$ 上的集合层。以下条件等价：
\begin{enumerate}
\item $\mathcal{F}$ 是可构造的；
\item 存在开覆盖 $X = \bigcup U_i$，使得 $\mathcal{F}|_{U_i}$ 是可构造的；
\item 存在由可构造局部闭子概形组成的分划 $X = \bigcup X_i$，使得
$\mathcal{F}|_{X_i}$ 是有限局部常值的。
\end{enumerate}
若 $\Lambda$ 为诺特环，则对阿贝尔层和 $\Lambda$-模层也有类似结论。
\end{lemma}

\begin{proof}
(1) 蕴含 (2) 是显然的。

\medskip\noindent
假设 (2) 成立。对每个 $x \in X$，可以找到某个 $i$ 以及一个仿射开邻域
$V_x \subset U_i$，后者是 $x$ 的邻域。因而可以找到有限仿射开覆盖
$X = \bigcup V_j$，使得对每个
$j$，都有一个由局部闭可构造子集组成的有限分解
$V_j = \coprod V_{j, k}$，并且 $\mathcal{F}|_{V_{j, k}}$ 是有限局部常值的。
由拓扑学，引理
\ref{topology-lemma-quasi-compact-open-immersion-constructible-image}
可知，每个 $V_{j, k}$ 作为 $X$ 的子集都是可构造的。由拓扑学，引理
\ref{topology-lemma-constructible-partition-refined-by-stratification}
可找到一个具有可构造局部闭层的有限分层 $X = \coprod X_l$，使得每个
$V_{j, k}$ 都是若干 $X_l$ 的并。因此 (3) 成立。

\medskip\noindent
假设 (3) 成立。设 $U \subset X$ 为仿射开集。则 $U \cap X_i$ 是 $U$ 的可构造
局部闭子集（例如由性质，引理 \ref{properties-lemma-locally-constructible}），
而 $U = \coprod U \cap X_i$ 是定义 \ref{etale-cohomology-definition-constructible} 中所述的
$U$ 的分划。因此 (1) 成立。
\end{proof}

\begin{lemma}
\label{etale-cohomology-lemma-constructible-constructible}
设 $X$ 为拟紧拟分离概形。设 $\mathcal{F}$ 为集合层、阿贝尔群层或
$\Lambda$-模层（其中 $\Lambda$ 为诺特环），定义在 $X_\etale$ 上。若存在可构造局部闭子概形
$T_i \subset X$，使得 (a) $X = \bigcup T_j$ 且 (b) $\mathcal{F}|_{T_j}$
是可构造的，则 $\mathcal{F}$ 是可构造的。
\end{lemma}

\begin{proof}
首先，可以假定该覆盖是有限的，因为 $X$ 在谱拓扑下是拟紧的
（拓扑学，引理 \ref{topology-lemma-constructible-hausdorff-quasi-compact}
以及性质，引理
\ref{properties-lemma-quasi-compact-quasi-separated-spectral}）。
注意到每个 $T_i$ 本身都是拟紧拟分离概形（因为它在 $X$ 中可构造；细节略去）。
因此，可以找到一个有限分划 $T_i = \coprod T_{i, j}$，其各部分是
$T_i$ 的局部闭可构造子集，并且 $\mathcal{F}|_{T_{i, j}}$ 是有限局部常值的
（引理 \ref{etale-cohomology-lemma-constructible-quasi-compact-quasi-separated}）。由拓扑学，
引理 \ref{topology-lemma-constructible-in-constructible} 可知，$T_{i, j}$ 是
$X$ 的可构造局部闭子概形。于是可以对 $X = \bigcup T_{i, j}$ 应用拓扑学，引理
\ref{topology-lemma-constructible-partition-refined-by-stratification}
以得到所需的 $X$ 的分划。
\end{proof}

\begin{lemma}
\label{etale-cohomology-lemma-constructible-local}
设 $X$ 为概形。集合层、阿贝尔群层或 $\Lambda$-模层（其中 $\Lambda$ 为诺特环）
的可构造性可以在 $X$ 上按 Zariski 局部检验。
\end{lemma}

\begin{proof}
该陈述的含义是：若 $X = \bigcup U_i$ 是开覆盖，且 $\mathcal{F}|_{U_i}$
是可构造的，则 $\mathcal{F}$ 是可构造的。若 $U \subset X$ 为仿射开集，则
$U = \bigcup U \cap U_i$，且 $\mathcal{F}|_{U \cap U_i}$ 是可构造的
（可构造层限制到开集仍可构造是显然的）。由引理
\ref{etale-cohomology-lemma-constructible-quasi-compact-quasi-separated} 可知，
$\mathcal{F}|_U$ 是可构造的，也就是说，$U$ 存在一个合适的分划。
\end{proof}

\begin{lemma}
\label{etale-cohomology-lemma-pullback-constructible}
设 $f : X \to Y$ 为概形态射。若 $\mathcal{F}$ 是可构造集合层、阿贝尔群层或
$\Lambda$-模层（其中 $\Lambda$ 为诺特环），定义在 $Y_\etale$ 上，则
$f^{-1}\mathcal{F}$ 在 $X_\etale$ 上也具有同样的性质。
\end{lemma}

\begin{proof}
由引理 \ref{etale-cohomology-lemma-constructible-local}，可归约到 $X$ 和 $Y$ 都是仿射的情形。
由引理 \ref{etale-cohomology-lemma-constructible-quasi-compact-quasi-separated}，只需找到一个
由可构造局部闭子概形组成的 $X$ 的有限分划，使得 $f^{-1}\mathcal{F}$ 在每一部分
上都是有限局部常值的。为此，只须拉回 $Y$ 的适配于 $\mathcal{F}$ 的分划，并应用
引理 \ref{etale-cohomology-lemma-pullback-locally-constant}。
\end{proof}

\begin{lemma}
\label{etale-cohomology-lemma-constructible-abelian}
设 $X$ 为概形。
\begin{enumerate}
\item 可构造集合层范畴在 $\Sh(X_\etale)$ 中对有限极限和有限余极限封闭。
\item 可构造阿贝尔层范畴是 $\textit{Ab}(X_\etale)$ 的弱 Serre 子范畴。
\item 设 $\Lambda$ 为诺特环。可构造 $\Lambda$-模层在 $X_\etale$ 上的范畴
是 $\textit{Mod}(X_\etale, \Lambda)$ 的弱 Serre 子范畴。
\end{enumerate}
\end{lemma}

\begin{proof}
证明 (3)。我们将使用同调，引理
\ref{homology-lemma-characterize-weak-serre-subcategory} 的判据。假设
$\varphi : \mathcal{F} \to \mathcal{G}$ 是可构造 $\Lambda$-模层之间的态射。
需要证明 $\mathcal{K} = \Ker(\varphi)$ 和 $\mathcal{Q} = \Coker(\varphi)$
都是可构造的。类似地，假设
$0 \to \mathcal{F} \to \mathcal{E} \to \mathcal{G} \to 0$
是 $\Lambda$-模层的短正合列，其中 $\mathcal{F}$、$\mathcal{G}$ 可构造。
需要证明 $\mathcal{E}$ 可构造。在这两种情形中，都可以用某个仿射开覆盖的成员
替换 $X$，因而可以假定 $X$ 是仿射的。进而，可以用 $X$ 的一个有限分划的各个
成员替换 $X$；该分划由可构造局部闭子概形组成，并且 $\mathcal{F}$ 和
$\mathcal{G}$ 在其上都是有限型且局部常值的。于是应用引理
\ref{etale-cohomology-lemma-kernel-finite-locally-constant} 即得结论。

\medskip\noindent
(1) 和 (2) 的证明非常类似，故略去。
\end{proof}

\begin{lemma}
\label{etale-cohomology-lemma-support-constructible}
设 $X$ 为拟紧拟分离概形。
\begin{enumerate}
\item 设 $\mathcal{F} \to \mathcal{G}$ 为 $X_\etale$ 上可构造集合层之间的态射。
则满足如下条件的点 $x \in X$ 所组成的集合，即
$\mathcal{F}_{\overline{x}} \to \mathcal{G}_{\overline{x}}$ 为满射、分别地为单射、
分别地同构于某个给定集合映射的点集，在 $X$ 中可构造。
\item 设 $\mathcal{F}$ 为 $X_\etale$ 上的可构造阿贝尔层。则 $\mathcal{F}$ 的
支集可构造。
\item 设 $\Lambda$ 为诺特环。设 $\mathcal{F}$ 为可构造 $\Lambda$-模层，
定义在 $X_\etale$ 上。则 $\mathcal{F}$ 的支集可构造。
\end{enumerate}
\end{lemma}

\begin{proof}
证明 (1)。设 $X = \coprod X_i$ 是由局部闭可构造子概形组成的 $X$ 的分划，
使得 $\mathcal{F}$ 和 $\mathcal{G}$ 在各部分上都是有限局部常值的（对
$\mathcal{F}$ 和 $\mathcal{G}$ 分别应用引理
\ref{etale-cohomology-lemma-constructible-quasi-compact-quasi-separated}，再取一个共同加细）。
然后，对该态射在每一部分上的限制应用引理
\ref{etale-cohomology-lemma-morphism-locally-constant}。

\medskip\noindent
(2) 和 (3) 的证明略去。
\end{proof}

\noindent
以下引理稍后会非常有用。粗略地说，它表明可构造层范畴具有某种弱
``诺特'' 性质。

\begin{lemma}
\label{etale-cohomology-lemma-colimit-constructible}
设 $X$ 为拟紧拟分离概形。设
$\mathcal{F} = \colim_{i \in I} \mathcal{F}_i$ 是集合层、阿贝尔层或模层的
滤过余极限。
\begin{enumerate}
\item 若 $\mathcal{F}$ 和 $\mathcal{F}_i$ 都是可构造集合层，则 ind-对象
$\mathcal{F}_i$ 本质上为常值，取值为 $\mathcal{F}$。
\item 若 $\mathcal{F}$ 和 $\mathcal{F}_i$ 都是可构造阿贝尔群层，则 ind-对象
$\mathcal{F}_i$ 本质上为常值，取值为 $\mathcal{F}$。
\item 设 $\Lambda$ 为诺特环。若 $\mathcal{F}$ 和 $\mathcal{F}_i$ 都是可构造
$\Lambda$-模层，则 ind-对象 $\mathcal{F}_i$ 本质上为常值，取值为
$\mathcal{F}$。
\end{enumerate}
\end{lemma}

\begin{proof}
证明 (1)。下文不再另行说明地使用如下事实：可构造层的有限极限和有限余极限
仍然可构造（引理 \ref{etale-cohomology-lemma-kernel-finite-locally-constant}）。对每个 $i$，
令 $T_i \subset X$ 为所有点 $x \in X$ 所组成的集合，其中
$\mathcal{F}_{i, \overline{x}} \to \mathcal{F}_{\overline{x}}$ 不是满射。
由于 $\mathcal{F}_i$ 和 $\mathcal{F}$ 可构造，
$T_i$ 是 $X$ 的可构造子集（引理 \ref{etale-cohomology-lemma-support-constructible}）。因为
$\mathcal{F}$ 的茎都是有限的，且
$\mathcal{F} = \colim_{i \in I} \mathcal{F}_i$，所以对每个 $x \in X$，
$x \not \in T_i$ 对充分大的 $i$ 成立。由性质，引理
\ref{properties-lemma-quasi-compact-quasi-separated-spectral}
可知 $X$ 是谱空间；由拓扑学，引理
\ref{topology-lemma-constructible-hausdorff-quasi-compact}，$X$ 上的可构造拓扑
是拟紧的。因此，$T_i = \emptyset$ 对充分大的 $i$ 成立。从而，
$\mathcal{F}_i \to \mathcal{F}$ 对充分大的 $i$ 是满射。现在假设
$\mathcal{F}_i \to \mathcal{F}$ 对所有 $i$ 都是满射。取 $i \in I$。对 $i' \geq i$，
记 $S_{i'} \subset X$ 为如下点 $x$ 组成的集合：
$\Im(\mathcal{F}_{i, \overline{x}} \to \mathcal{F}_{\overline{x}})$
中的元素个数不等于
$\Im(\mathcal{F}_{i, \overline{x}} \to \mathcal{F}_{i', \overline{x}})$
中的元素个数。由于 $\mathcal{F}_i$、$\mathcal{F}_{i'}$ 和 $\mathcal{F}$
可构造，$S_{i'}$ 是 $X$ 的可构造子集（细节略去；提示：使用引理
\ref{etale-cohomology-lemma-support-constructible}）。由于 $\mathcal{F}_i$ 和 $\mathcal{F}$
的茎都是有限的，且 $\mathcal{F} = \colim_{i' \geq i} \mathcal{F}_{i'}$，
所以对每个 $x \in X$，$x \not \in S_{i'}$ 对充分大的 $i'$ 成立。
由与上面相同的论证，可找到充分大的 $i'$，使得 $S_{i'} = \emptyset$。
因此，$\mathcal{F}_i \to \mathcal{F}_{i'}$ 如所需地通过 $\mathcal{F}$ 分解。

\medskip\noindent
证明 (2)。注意到可构造阿贝尔层也是可构造集合层。因此 (2) 由 (1) 得出。

\medskip\noindent
证明 (3)。下文不再另行说明地使用如下事实：可构造 $\Lambda$-模层范畴是阿贝尔
范畴（引理 \ref{etale-cohomology-lemma-kernel-finite-locally-constant}）。对每个 $i$，令
$\mathcal{Q}_i$ 为态射 $\mathcal{F}_i \to \mathcal{F}$ 的余核。令 $T_i$ 为
$\mathcal{Q}_i$ 的支集；它是 $X$ 的可构造子集，因为 $\mathcal{Q}_i$ 可构造
（引理 \ref{etale-cohomology-lemma-support-constructible}）。由于 $\mathcal{F}$ 的茎都是有限
$\Lambda$-模，且 $\mathcal{F} = \colim_{i \in I} \mathcal{F}_i$，所以对每个
$x \in X$，$x \not \in T_i$ 对充分大的 $i$ 成立。由性质，引理
\ref{properties-lemma-quasi-compact-quasi-separated-spectral}
可知 $X$ 是谱空间；由拓扑学，引理
\ref{topology-lemma-constructible-hausdorff-quasi-compact}，$X$ 上的可构造拓扑
是拟紧的。因此，$T_i = \emptyset$ 对充分大的 $i$ 成立。这就证明了第一个断言。
对于第二个断言，现在假设 $\mathcal{F}_i \to \mathcal{F}$ 对所有 $i$ 都是
满射。取 $i \in I$。对 $i' \geq i$，记 $\mathcal{K}_{i'}$ 为
$\Ker(\mathcal{F}_i \to \mathcal{F})$ 在 $\mathcal{F}_{i'}$ 中的像。令
$S_{i'}$ 为 $\mathcal{K}_{i'}$ 的支集；它是 $X$ 的可构造子集，因为
$\mathcal{K}_{i'}$ 可构造。由于 $\Ker(\mathcal{F}_i \to \mathcal{F})$
的茎都是有限 $\Lambda$-模，且
$\mathcal{F} = \colim_{i' \geq i} \mathcal{F}_{i'}$，所以对每个 $x \in X$，
$x \not \in S_{i'}$ 对充分大的 $i'$ 成立。由与上面相同的论证，可找到充分大的
$i'$，使得 $S_{i'} = \emptyset$。因此，
$\mathcal{F}_i \to \mathcal{F}_{i'}$ 如所需地通过 $\mathcal{F}$ 分解。
\end{proof}

\begin{lemma}
\label{etale-cohomology-lemma-tensor-product-constructible}
设 $X$ 为概形，$\Lambda$ 为诺特环。两个可构造 $\Lambda$-模层在
$X_\etale$ 上的张量积仍是可构造 $\Lambda$-模层。
\end{lemma}

\begin{proof}
该问题立即归约到 $X$ 为仿射的情形。由于 $X$ 的任意两个由可构造局部闭层组成的
分划都有同类的共同加细，并且拉回与张量积可交换，所以可归约到引理
\ref{etale-cohomology-lemma-tensor-product-locally-constant}。
\end{proof}

\begin{lemma}
\label{etale-cohomology-lemma-tensor-constructible}
设 $\Lambda \to \Lambda'$ 为诺特环同态。设 $X$ 为概形，$\mathcal{F}$ 为
可构造 $\Lambda$-模层，定义在 $X_\etale$ 上。则
$\mathcal{F} \otimes_{\underline{\Lambda}} \underline{\Lambda'}$
是可构造 $\Lambda'$-模层。
\end{lemma}

\begin{proof}
略。提示：在仿射局部，可以使用同一个分层。
\end{proof}










\section{关于态射的辅助引理}
\label{etale-cohomology-section-stratify-morphisms}

\noindent
下面给出一些在证明可构造层的函子性性质时有用的引理。

\begin{lemma}
\label{etale-cohomology-lemma-etale-stratified-finite}
设 $U \to X$ 为拟紧拟分离概形之间的 \'etale 态射（例如诺特概形之间的
\'etale 态射）。则存在一个由可构造局部闭子概形组成的分划
$X = \coprod_i X_i$，使得 $X_i \times_X U \to X_i$ 对所有 $i$ 都是有限
\'etale 态射。
\end{lemma}

\begin{proof}
若 $U \to X$ 是分离的，则这就是态射进阶，引理
\ref{more-morphisms-lemma-stratify-flat-fp-qf}。一般而言，可以假定 $X$ 是仿射的。
选择一个有限仿射开覆盖 $U = \bigcup U_j$。对所有态射 $U_j \to X$ 和
$U_j \cap U_{j'} \to X$ 应用前一种情形，并选取所得各分划的共同加细
$X = \coprod X_i$。进一步加细该分划后，还可以假定 $X_i$ 是仿射的。
固定 $i$，并令 $V = U \times_X X_i$。态射
$V_j = U_j \times_X X_i \to X_i$ 以及
$V_{jj'} = (U_j \cap U_{j'}) \times_X X_i \to X_i$ 都是有限 \'etale 的。
因此，$V_j$ 和 $V_{jj'}$ 都是仿射概形，且 $V_{jj'} \subset V_j$ 既闭又开
（因为 $V_{jj'} \to X_i$ 是固有的，故可应用态射，引理
\ref{morphisms-lemma-image-proper-scheme-closed}）。于是 $V = \bigcup V_j$ 是分离的，
因为 $\mathcal{O}(V_j) \to \mathcal{O}(V_{jj'})$ 是满射，见概形，引理
\ref{schemes-lemma-characterize-separated}。所以前一种情形适用于 $V \to X_i$，
而且如果需要，可以进一步加细分划（事实上不需要，但我们不必用到这一点）。
\end{proof}

\noindent
在诺特情形中，可以用诺特归纳法和下面这个有趣的引理来证明前一个引理。

\begin{lemma}
\label{etale-cohomology-lemma-generically-finite}
设 $f: X \to Y$ 为拟紧、拟分离且局部有限型的概形态射。若 $\eta$ 是 $Y$
的某个不可约分支的泛点，且 $f^{-1}(\eta)$ 是有限的，则存在开集
$V \subset Y$，它包含 $\eta$，并且 $f^{-1}(V) \to V$ 是有限态射。
\end{lemma}

\begin{proof}
这就是态射，引理 \ref{morphisms-lemma-generically-finite}。
\end{proof}

\noindent
下面引理的陈述还可以略作加强。

\begin{lemma}
\label{etale-cohomology-lemma-decompose-quasi-finite-morphism}
设 $f : Y \to X$ 为仿射概形之间的拟有限且有限表现态射。
\begin{enumerate}
\item 存在仿射概形的满态射 $X' \to X$ 和闭子概形
$Z' \subset Y' = X' \times_X Y$，使得
\begin{enumerate}
\item $Z' \subset Y'$ 是一个加厚；
\item $Z' \to X'$ 是有限 \'etale 态射。
\end{enumerate}
\item 存在一个由局部闭、可构造、仿射的层组成的有限分划
$X = \coprod X_i$，以及有限局部自由满态射 $X'_i \to X_i$，使得
$Y'_i = X'_i \times_X Y \to X'_i$ 的既约化同构于
$\coprod_{j = 1}^{n_i} (X'_i)_{red} \to (X'_i)_{red}$，其中 $n_i$ 为某个整数。
\end{enumerate}
\end{lemma}

\begin{proof}
令 $X' = \coprod X'_i$，可见 (2) 蕴含 (1)。写成 $X = \Spec(A)$ 和
$Y = \Spec(B)$。把 $A$ 写成有限型 $\mathbf{Z}$-代数 $A_i$ 的滤过余极限。
由于 $B$ 是有限表现 $A$-代数，可知存在 $0 \in I$ 以及有限型环同态
$A_0 \to B_0$，使得 $B = \colim B_i$，其中
$B_i = A_i \otimes_{A_0} B_0$，见代数，引理
\ref{algebra-lemma-colimit-category-fp-algebras}。当 $i$ 充分大时，
$A_i \to B_i$ 是拟有限的，见极限，引理
\ref{limits-lemma-descend-quasi-finite}。因此，可归约到 $\mathbf{Z}$ 上有限型代数
的情形，特别地，可归约到诺特情形。（细节略去。）

\medskip\noindent
假设 $X$ 和 $Y$ 都是诺特的。在此情形中，$X$ 的任意局部闭子集都是可构造的。
由引理 \ref{etale-cohomology-lemma-generically-finite} 和诺特归纳法可知，存在有限分划
$X = \coprod X_i$，它是 $X$ 的由局部闭层组成的分划，并使得
$Y \times_X X_i \to X_i$ 是有限的。可以加细该分划，使其各层仿射。因此，
把 $X$ 替换为 $X' = \coprod X_i$ 后，
可以假定 $Y \to X$ 是有限的。

\medskip\noindent
假设 $X$ 和 $Y$ 都是诺特的，且 $Y \to X$ 有限。假设在沿某个满、平坦、
拟有限态射 $U \to X$ 作基变换后可以证明 (2)。于是有分划
$U = \coprod U_i$ 和有限局部自由态射 $U'_i \to U_i$，使得
$U'_i \times_X Y \to U'_i$ 同构于
$\coprod_{j = 1}^{n_i} (U'_i)_{red} \to (U'_i)_{red}$，其中 $n_i$ 为某个整数。
由前一段的论证，可以找到一个具有局部闭仿射层的分划 $X = \coprod X_j$，使得
$X_j \times_X U_i \to X_j$ 对所有 $i, j$ 都是有限的。由态射，引理
\ref{morphisms-lemma-finite-flat}，每个 $X_j \times_X U_i \to X_j$ 都是有限
局部自由的。因此，$X_j \times_X U'_i \to X_j$ 是有限局部自由的（态射，引理
\ref{morphisms-lemma-composition-finite-locally-free}）。由此，
$X = \coprod X_j$ 和 $X_j' = \coprod_i X_j \times_X U'_i$ 给出了
$Y \to X$ 的一个解。因此，只需在作满、平坦、拟有限基变换后证明该结论
（在诺特情形中）。

\medskip\noindent
应用态射，引理 \ref{morphisms-lemma-massage-finite}，可知可以假定 $Y$ 是某个
仿射概形 $Z$ 的闭子概形，而后者在集合意义下是闭子概形的有限并
$Z = \bigcup_{i \in I} Z_i$，其中这些闭子概形同构地映到 $X$。在这种情形下，
我们将找到一个可行的、具有仿射局部闭层的有限分划 $X = \coprod X_j$
（换言之，$X'_j = X_j$）。令 $T_i = Y \cap Z_i$。这是 $X$ 的闭子概形。
由于 $X$ 是诺特的，可以找到一个由仿射局部闭子概形组成的有限分划
$X = \coprod X_j$，使得每个 $X_j \times_X T_i$ 在集合意义下都是若干层
$X_j \times_X Z_i$ 的并。把 $X$ 替换为 $X_j$，可知可以假定
$I = I_1 \amalg I_2$，其中 $Z_i \subset Y$ 对 $i \in I_1$ 成立，而
$Z_i \cap Y = \emptyset$ 对 $i \in I_2$ 成立。把 $Z$ 替换为
$\bigcup_{i \in I_1} Z_i$，可知可以假定 $Y = Z$。
最后，可以再次用上述分划的各成员替换 $X$，使得对每个 $i, i' \subset I$，交
$Z_i \cap Z_{i'}$ 或者为空，或者在集合意义下同时等于 $Z_i$ 和 $Z_{i'}$。
这显然意味着 $Y$ 在集合意义下等于诸 $Z_i$ 的不交并，而这正是所要证明的。
\end{proof}








\section{关于可构造层的进一步结果}
\label{etale-cohomology-section-more-constructible}

\noindent
设 $\Lambda$ 为诺特环，$X$ 为概形。我们经常把 $X_\etale$ 看成环化位点，
其环层为 $\underline{\Lambda}$。在阿贝尔层的情形中，我们经常取
$\Lambda = \mathbf{Z}/n\mathbf{Z}$，其中 $n$ 为某个合适的整数。

\begin{lemma}
\label{etale-cohomology-lemma-jshriek-constructible}
设 $j : U \to X$ 为拟紧拟分离概形之间的 \'etale 态射。
\begin{enumerate}
\item 层 $h_U$ 是可构造集合层。
\item 层 $j_!\underline{M}$ 对有限阿贝尔群 $M$ 是可构造阿贝尔层。
\item 若 $\Lambda$ 为诺特环，且 $M$ 是有限 $\Lambda$-模，则
$j_!\underline{M}$ 是可构造 $\Lambda$-模层，定义在 $X_\etale$ 上。
\end{enumerate}
\end{lemma}

\begin{proof}
由引理 \ref{etale-cohomology-lemma-etale-stratified-finite}，存在分划 $\coprod_i X_i$，使得
$\pi_i : j^{-1}(X_i) \to X_i$ 是有限 \'etale 态射。$h_U$ 在 $X_i$ 上的限制是
$h_{j^{-1}(X_i)}$；由引理 \ref{etale-cohomology-lemma-characterize-finite-locally-constant}，
它是有限局部常值的。对于 (2) 和 (3)，注意到
$$
j_!(\underline{M})|_{X_i} =
\pi_{i!}(\underline{M}) =
\pi_{i*}(\underline{M})
$$
由引理 \ref{etale-cohomology-lemma-shriek-base-change} 和
\ref{etale-cohomology-lemma-shriek-equals-star-finite-etale} 成立。因此，只需对有限 \'etale 态射
$\pi : Y \to X$ 证明该引理。这就是引理
\ref{etale-cohomology-lemma-pushforward-locally-constant}。
\end{proof}

\begin{lemma}
\label{etale-cohomology-lemma-torsion-colimit-constructible}
设 $X$ 为拟紧拟分离概形。
\begin{enumerate}
\item 设 $\mathcal{F}$ 为 $X_\etale$ 上的集合层。则 $\mathcal{F}$ 是可构造
集合层的滤过余极限。
\item 设 $\mathcal{F}$ 为 $X_\etale$ 上的挠阿贝尔层。则 $\mathcal{F}$ 是
可构造阿贝尔层的滤过余极限。
\item 设 $\Lambda$ 为诺特环，$\mathcal{F}$ 为 $\Lambda$-模层，定义在
$X_\etale$ 上。则 $\mathcal{F}$ 是可构造 $\Lambda$-模层的滤过余极限。
\end{enumerate}
\end{lemma}

\begin{proof}
令 $\mathcal{B}$ 为 $X_\etale$ 的拟紧拟分离对象所组成的集合。由位点上的模，
引理 \ref{sites-modules-lemma-module-filtered-colimit-constructibles}，
任意集合层都是如下形式的层的滤过余极限：
$$
\text{Coequalizer}\left(
\xymatrix{
\coprod\nolimits_{j = 1, \ldots, m} h_{V_j}
\ar@<1ex>[r] \ar@<-1ex>[r] &
\coprod\nolimits_{i = 1, \ldots, n} h_{U_i}
}
\right)
$$
其中 $V_j$ 和 $U_i$ 是 $X_\etale$ 的拟紧拟分离对象。由引理
\ref{etale-cohomology-lemma-jshriek-constructible} 和 \ref{etale-cohomology-lemma-constructible-abelian}，
这些余等化子是可构造的。这就证明了 (1)。

\medskip\noindent
设 $\Lambda$ 为诺特环。由位点上的模，引理
\ref{sites-modules-lemma-module-filtered-colimit-constructibles}，
$\Lambda$-模层 $\mathcal{F}$ 是如下形式的模层的滤过余极限：
$$
\Coker\left(
\bigoplus\nolimits_{j = 1, \ldots, m} j_{V_j!}\underline{\Lambda}_{V_j}
\longrightarrow
\bigoplus\nolimits_{i = 1, \ldots, n} j_{U_i!}\underline{\Lambda}_{U_i}
\right)
$$
其中 $V_j$ 和 $U_i$ 是 $X_\etale$ 的拟紧拟分离对象。由引理
\ref{etale-cohomology-lemma-jshriek-constructible} 和 \ref{etale-cohomology-lemma-constructible-abelian}，
这些余核是可构造的。这就证明了 (3)。

\medskip\noindent
证明 (2)。首先写成 $\mathcal{F} = \bigcup \mathcal{F}[n]$，其中
$\mathcal{F}[n]$ 是 $n$-挠子层。于是可以把 $\mathcal{F}[n]$ 看成
$\mathbf{Z}/n\mathbf{Z}$-模层，并应用 (3)。
\end{proof}

\begin{lemma}
\label{etale-cohomology-lemma-check-constructible}
设 $f : X \to Y$ 为拟紧拟分离概形之间的满态射。
\begin{enumerate}
\item 设 $\mathcal{F}$ 为 $Y_\etale$ 上的集合层。则 $\mathcal{F}$ 可构造，
当且仅当 $f^{-1}\mathcal{F}$ 可构造。
\item 设 $\mathcal{F}$ 为 $Y_\etale$ 上的阿贝尔层。则 $\mathcal{F}$ 可构造，
当且仅当 $f^{-1}\mathcal{F}$ 可构造。
\item 设 $\Lambda$ 为诺特环。设 $\mathcal{F}$ 为 $\Lambda$-模层，定义在
$Y_\etale$ 上。则 $\mathcal{F}$ 可构造，当且仅当 $f^{-1}\mathcal{F}$ 可构造。
\end{enumerate}
\end{lemma}

\begin{proof}
一个方向由引理 \ref{etale-cohomology-lemma-pullback-constructible} 得出。反过来，假设
$f^{-1}\mathcal{F}$ 可构造。使用引理
\ref{etale-cohomology-lemma-torsion-colimit-constructible}，把
$\mathcal{F} = \colim \mathcal{F}_i$ 写成可构造层（集合层、阿贝尔群层或模层）的
滤过余极限。由于 $f^{-1}$ 是左伴随，它与余极限可交换（范畴，引理
\ref{categories-lemma-adjoint-exact}），从而
$f^{-1}\mathcal{F} = \colim f^{-1}\mathcal{F}_i$.
由引理 \ref{etale-cohomology-lemma-colimit-constructible} 可知，
$f^{-1}\mathcal{F}_i \to f^{-1}\mathcal{F}$
对所有充分大的 $i$ 都是满射。由于 $f$ 是满射，通过使用引理
\ref{etale-cohomology-lemma-stalk-pullback} 和定理 \ref{etale-cohomology-theorem-exactness-stalks} 考察茎可得，
$\mathcal{F}_i \to \mathcal{F}$ 对所有充分大的 $i$ 都是满射。因此，
$\mathcal{F}$ 是某个可构造层 $\mathcal{G}$ 的商。再把该论证应用于
$\mathcal{G} \times_\mathcal{F} \mathcal{G}$ 或
$\mathcal{G} \to \mathcal{F}$ 的核，并使用 $f^{-1}$ 的正合性以及可构造层
（集合层、阿贝尔群层或模层）范畴在下列范畴内部对有限（余）极限或（余）核封闭：
$\Sh(Y_\etale)$, $\Sh(X_\etale)$, $\textit{Ab}(Y_\etale)$,
$\textit{Ab}(X_\etale)$, $\textit{Mod}(Y_\etale, \Lambda)$ 以及
$\textit{Mod}(X_\etale, \Lambda)$，见引理
\ref{etale-cohomology-lemma-constructible-abelian}，即可得出结论。
\end{proof}

\begin{lemma}
\label{etale-cohomology-lemma-pushforward-constructible}
设 $f : X \to Y$ 为概形的有限 \'etale 态射，$\Lambda$ 为诺特环。若
$\mathcal{F}$ 是可构造集合层、可构造阿贝尔群层或可构造 $\Lambda$-模层，
定义在 $X_\etale$ 上，则 $f_*\mathcal{F}$ 在 $Y_\etale$ 上也具有同样的性质。
\end{lemma}

\begin{proof}
由引理 \ref{etale-cohomology-lemma-constructible-local}，只需在 $Y$ 上按 Zariski 局部检验；
由引理 \ref{etale-cohomology-lemma-check-constructible}，可以用某个 \'etale 覆盖替换 $Y$
（$f_*$ 的构造与 \'etale 局部化可交换）。有限 \'etale 态射在 \'etale 局部
同构于若干同构的不交并，见 \'Etale 态射，引理
\ref{etale-lemma-finite-etale-etale-local}。因此，在集合层的情形中，该引理是说：
若 $\mathcal{F}_i$（$i = 1, \ldots, n$）是可构造集合层，则
$\prod_{i = 1, \ldots, n} \mathcal{F}_i$ 也是可构造的。这是显然的。
对阿贝尔群层和模层也类似。
\end{proof}

\begin{lemma}
\label{etale-cohomology-lemma-category-constructible-sets}
设 $X$ 为拟紧拟分离概形。可构造集合层范畴正是 $\Sh(X_\etale)$ 的如下满子范畴：
其对象为可写成下列余等化子的层 $\mathcal{F}$：
$$
\xymatrix{
\mathcal{F}_1
\ar@<1ex>[r] \ar@<-1ex>[r]
&
\mathcal{F}_0 \ar[r]
&
\mathcal{F}}
$$
其中 $\mathcal{F}_i$（$i = 0, 1$）是形如 $h_U$ 的层的有限余积，而 $U$ 是
$X_\etale$ 的拟紧拟分离对象。
\end{lemma}

\begin{proof}
在引理 \ref{etale-cohomology-lemma-torsion-colimit-constructible} 的证明中，我们已经看到这种形式的
层是可构造的。反过来，假设对每个可构造集合层 $\mathcal{F}$，都能找到一个满射
$\mathcal{F}_0 \to \mathcal{F}$，其中 $\mathcal{F}_0$ 如引理中所述。那么就能找到
所需的满射
$\mathcal{F}_1 \to \mathcal{F}_0 \times_\mathcal{F} \mathcal{F}_0$
，因为由引理 \ref{etale-cohomology-lemma-constructible-abelian}，后一个层是可构造的。

\medskip\noindent
由拓扑学，引理
\ref{topology-lemma-constructible-partition-refined-by-stratification}
，可以选择有限分层 $X = \coprod_{i \in I} X_i$，使得 $\mathcal{F}$ 在每一层
上都是有限局部常值的。我们对 $I$ 的基数作归纳来证明该结论。设 $i \in I$
是 $I$ 的偏序中的极小元。则 $X_i \subset X$ 是闭的。由归纳假设，存在有限多个
拟紧拟分离对象 $U_\alpha$，它们属于 $(X \setminus X_i)_\etale$，并且存在满射
$\coprod h_{U_\alpha} \to \mathcal{F}|_{X \setminus X_i}$。它们确定一个态射
$$
\coprod h_{U_\alpha} \to \mathcal{F}
$$
，它限制到 $X \setminus X_i$ 后是满射。由引理
\ref{etale-cohomology-lemma-characterize-finite-locally-constant}，可知
$\mathcal{F}|_{X_i} = h_V$，其中某个概形 $V$ 在 $X_i$ 上有限 \'etale。设
$\overline{v}$ 为 $V$ 的位于 $\overline{x} \in X_i$ 上方的几何点。可以把
$\overline{v}$ 看成茎 $\mathcal{F}_{\overline{x}} = V_{\overline{x}}$ 的元素。
因此，可以找到一个 \'etale 邻域 $(U, \overline{u})$，其中心为 $\overline{x}$，以及截面
$s \in \mathcal{F}(U)$，其在 $\overline{x}$ 处的茎给出 $\overline{v}$。把 $s$
看成态射 $s : h_U \to \mathcal{F}$，并限制到 $X_i$，得到态射
$s|_{X_i} : U \times_X X_i \to V$；它在 $X_i$ 上，并把 $\overline{u}$ 映到
$\overline{v}$。由于 $V$ 是拟紧的（它在闭子概形 $X_i$ 上有限，而后者属于拟紧概形
$X$），有限多个 $s^{(1)}, \ldots, s^{(m)}$（它们是 $\mathcal{F}$ 在
$U^{(1)}, \ldots, U^{(m)}$ 上的这些截面）
上，将确定一个联合满态射
$$
\coprod s^{(j)}|_{X_i} : \coprod U^{(j)} \times_X X_i \longrightarrow V
$$
于是得到所需的满射
$$
\coprod h_{U_\alpha} \amalg \coprod h_{U^{(j)}} \to \mathcal{F}
$$
。
\end{proof}

\begin{lemma}
\label{etale-cohomology-lemma-category-constructible-modules}
设 $X$ 为拟紧拟分离概形，$\Lambda$ 为诺特环。可构造 $\Lambda$-模层范畴
正是由如下形式的模层组成的范畴：
$$
\Coker\left(
\bigoplus\nolimits_{j = 1, \ldots, m} j_{V_j!}\underline{\Lambda}_{V_j}
\longrightarrow
\bigoplus\nolimits_{i = 1, \ldots, n} j_{U_i!}\underline{\Lambda}_{U_i}
\right)
$$
其中 $V_j$ 和 $U_i$ 是 $X_\etale$ 的拟紧拟分离对象。事实上，甚至可以假定
$U_i$ 和 $V_j$ 都是仿射的。
\end{lemma}

\begin{proof}
在引理 \ref{etale-cohomology-lemma-torsion-colimit-constructible} 的证明中，我们已经看到这种形式的
模层是可构造的。由于可构造模层范畴是阿贝尔范畴（引理
\ref{etale-cohomology-lemma-constructible-abelian}），只需证明：给定可构造模层 $\mathcal{F}$，
存在满射
$$
\bigoplus\nolimits_{i = 1, \ldots, n} j_{U_i!}\underline{\Lambda}_{U_i}
\longrightarrow \mathcal{F}
$$
，其中 $U_i$ 是 $X_\etale$ 的某些仿射对象。由位点上的模，引理
\ref{sites-modules-lemma-module-filtered-colimit-constructibles}
，存在满射
$$
\Psi :
\bigoplus\nolimits_{i \in I} j_{U_i!}\underline{\Lambda}_{U_i}
\longrightarrow
\mathcal{F}
$$
，其中 $U_i$ 是仿射的，且直和取遍一个可能无限的指标集 $I$。对每个有限子集
$I' \subset I$，令
$$
T_{I'} = \text{Supp}(\Coker(
\bigoplus\nolimits_{i \in I'} j_{U_i!}\underline{\Lambda}_{U_i}
\longrightarrow \mathcal{F}))
$$
由可构造层的定义，集合 $T_{I'}$ 是 $X$ 的可构造子集。我们要证明
$T_{I'} = \emptyset$ 对某个 $I'$ 成立。由于每个茎 $\mathcal{F}_{\overline{x}}$ 都是有限型
$\Lambda$-模，且 $\Psi$ 是满射，所以对每个 $x \in X$，存在某个 $I'$，使得
$x \not \in T_{I'}$。换言之，有
$\emptyset = \bigcap_{I' \subset I\text{ 有限}} T_{I'}$。由于
$X$ 由性质，引理
\ref{properties-lemma-quasi-compact-quasi-separated-spectral}
是谱空间，而由拓扑学，引理
\ref{topology-lemma-constructible-hausdorff-quasi-compact}，$X$ 上的可构造拓扑
是拟紧的。因此，$T_{I'} = \emptyset$ 对某个有限的 $I' \subset I$ 成立，
如所需。
\end{proof}

\begin{lemma}
\label{etale-cohomology-lemma-category-constructible-abelian}
设 $X$ 为拟紧拟分离概形。可构造阿贝尔层范畴正是由如下形式的阿贝尔层组成的范畴：
$$
\Coker\left(
\bigoplus\nolimits_{j = 1, \ldots, m}
j_{V_j!}\underline{\mathbf{Z}/m_j\mathbf{Z}}_{V_j}
\longrightarrow
\bigoplus\nolimits_{i = 1, \ldots, n}
j_{U_i!}\underline{\mathbf{Z}/n_i\mathbf{Z}}_{U_i}
\right)
$$
其中 $V_j$ 和 $U_i$ 是 $X_\etale$ 的拟紧拟分离对象，且 $m_j$、$n_i$ 为正整数。
事实上，甚至可以假定 $U_i$ 和 $V_j$ 都是仿射的。
\end{lemma}

\begin{proof}
这由对 $\Lambda = \mathbf{Z}/n\mathbf{Z}$ 应用引理
\ref{etale-cohomology-lemma-category-constructible-modules}，以及如下事实得出：由于 $X$ 拟紧，
每个可构造阿贝尔层都被某个正整数 $n$ 零化（细节略去）。
\end{proof}

\begin{lemma}
\label{etale-cohomology-lemma-constructible-is-compact}
设 $X$ 为拟紧拟分离概形，$\Lambda$ 为诺特环。设 $\mathcal{F}$ 为可构造集合层、
阿贝尔群层或 $\Lambda$-模层，定义在 $X_\etale$ 上。设
$\mathcal{G} = \colim \mathcal{G}_i$ 为集合层、阿贝尔群层或 $\Lambda$-模层的
滤过余极限。则
$$
\Mor(\mathcal{F}, \mathcal{G}) = \colim \Mor(\mathcal{F}, \mathcal{G}_i)
$$
在集合层、阿贝尔群层或 $\Lambda$-模层的范畴中成立，该范畴定义在 $X_\etale$ 上。
\end{lemma}

\begin{proof}
先考虑集合层的情形。由引理 \ref{etale-cohomology-lemma-category-constructible-sets}，只需对
$h_U$ 证明该引理，其中 $U$ 是 $X_\etale$ 的拟紧拟分离对象。回忆
$\Mor(h_U, \mathcal{G}) = \mathcal{G}(U)$。因此，结论由位点，引理
\ref{sites-lemma-directed-colimits-sections} 得出。

\medskip\noindent
在阿贝尔层或模层的情形中，使用引理
\ref{etale-cohomology-lemma-category-constructible-abelian} 和
\ref{etale-cohomology-lemma-category-constructible-modules}，以同样方式得到结论。对于阿贝尔层，
还要补充注意到
$\Mor(j_{U!}\underline{\mathbf{Z}/n\mathbf{Z}}, \mathcal{G})$
等于所有 $n$-挠元素所组成的集合，这些元素属于 $\mathcal{G}(U)$。
\end{proof}

\begin{lemma}
\label{etale-cohomology-lemma-finite-pushforward-constructible}
设 $f : X \to Y$ 为概形的有限且有限表现态射，$\Lambda$ 为诺特环。若
$\mathcal{F}$ 是可构造集合层、阿贝尔群层或 $\Lambda$-模层，定义在 $X_\etale$ 上，
则 $f_*\mathcal{F}$ 也是可构造的。
\end{lemma}

\begin{proof}
由引理 \ref{etale-cohomology-lemma-constructible-local}，只需在 $X$ 和 $Y$ 都是仿射时证明该结论。
由引理 \ref{etale-cohomology-lemma-finite-pushforward-commutes-with-base-change} 和
\ref{etale-cohomology-lemma-check-constructible}，可以基变换到任意满射到 $X$ 的仿射概形。由引理
\ref{etale-cohomology-lemma-decompose-quasi-finite-morphism}，这归约到有限 \'etale 态射的情形
（因为加厚给出 \'etale 拓扑斯、甚至小 \'etale 位点之间的等价，见定理
\ref{etale-cohomology-theorem-topological-invariance}）。有限 \'etale 情形就是引理
\ref{etale-cohomology-lemma-pushforward-constructible}。
\end{proof}

\begin{lemma}
\label{etale-cohomology-lemma-category-is-colimit}
设 $X = \lim_{i \in I} X_i$ 为一个具有仿射转移态射的概形有向系的极限。
假设 $X_i$ 对所有 $i \in I$ 都拟紧拟分离。
\begin{enumerate}
\item $X_\etale$ 上的可构造集合层范畴是各 $(X_i)_\etale$ 上可构造集合层范畴
的余极限。
\item $X_\etale$ 上的可构造阿贝尔层范畴是各 $(X_i)_\etale$ 上可构造阿贝尔层
范畴的余极限。
\item 设 $\Lambda$ 为诺特环。可构造 $\Lambda$-模层在 $X_\etale$ 上的范畴是
可构造 $\Lambda$-模层在各 $(X_i)_\etale$ 上的范畴的余极限。
\end{enumerate}
\end{lemma}

\begin{proof}
证明 (1)。用 $f_i : X \to X_i$ 表示投影态射。证明分为三部分，分别对应于
``忠实''、``全忠实'' 和 ``本质满射''。

\medskip\noindent
忠实性。选取 $0 \in I$，并设 $\mathcal{F}_0$、$\mathcal{G}_0$ 为 $X_0$ 上的
可构造层。假设 $a, b : \mathcal{F}_0 \to \mathcal{G}_0$ 是满足
$f_0^{-1}a = f_0^{-1}b$ 的态射。令 $E \subset X_0$ 为如下点
$x \in X_0$ 所组成的集合：$a_{\overline{x}} = b_{\overline{x}}$。由引理
\ref{etale-cohomology-lemma-support-constructible}，子集 $E \subset X_0$ 可构造。根据假设，
$X \to X_0$ 映入 $E$。由极限，引理
\ref{limits-lemma-limit-contained-in-constructible}，可找到 $i \geq 0$，使得
$X_i \to X_0$ 映入 $E$。因此 $f_{i0}^{-1}a = f_{i0}^{-1}b$。

\medskip\noindent
全忠实性。选取 $0 \in I$，并设 $\mathcal{F}_0$、$\mathcal{G}_0$ 为 $X_0$ 上的
可构造层。假设 $a : f_0^{-1}\mathcal{F}_0 \to f_0^{-1}\mathcal{G}_0$ 是一个态射。
我们断言，存在某个 $i$ 和态射
$a_i : f_{i0}^{-1}\mathcal{F}_0 \to f_{i0}^{-1}\mathcal{G}_0$
，它的拉回即为 $a$，这是 $X$ 上的态射。由引理
\ref{etale-cohomology-lemma-category-constructible-sets}，可以替换 $\mathcal{F}_0$，所用对象是若干
由 $(X_0)_\etale$ 的拟紧拟分离对象表示的层的有限余积。因此，需要证明：
若 $U_0 \to X_0$ 是 $(X_0)_\etale$ 的这样一个
对象，则
$$
f_0^{-1}\mathcal{G}(U) = \colim_{i \geq 0} f_{i0}^{-1}\mathcal{G}(U_i)
$$
，其中 $U = X \times_{X_0} U_0$ 且 $U_i = X_i \times_{X_0} U_0$。这是定理
\ref{etale-cohomology-theorem-colimit} 的特殊情形。

\medskip\noindent
本质满射性。需要证明，每个可构造层 $\mathcal{F}$（定义在 $X$ 上）都同构于
$f_i^{-1}\mathcal{F}$，其中某个可构造层 $\mathcal{F}_i$ 定义在 $X_i$ 上。
应用引理 \ref{etale-cohomology-lemma-category-constructible-sets} 并使用前两段的结论，可知只需对
$h_U$ 证明这一点，其中 $U$ 是 $X_\etale$ 的某个拟紧拟分离对象。在这种情形中，
需要证明 $U$ 是某个在 $X_i$ 上 \'etale 的拟紧拟分离概形的基变换，其中 $i$
是某个指标。这由极限，引理 \ref{limits-lemma-descend-finite-presentation} 和
\ref{limits-lemma-descend-etale} 得出。

\medskip\noindent
证明 (3)。论证与集合层的论证非常类似，只是用引理
\ref{etale-cohomology-lemma-category-constructible-modules} 代替引理
\ref{etale-cohomology-lemma-category-constructible-sets}。细节略去。(2) 由 (3) 得出，因为拟紧概形上的
每个可构造阿贝尔层都是可构造 $\mathbf{Z}/n\mathbf{Z}$-模层，其中 $n$ 为某个整数。
\end{proof}

\begin{lemma}
\label{etale-cohomology-lemma-category-loc-cst-is-colimit}
设 $X = \lim_{i \in I} X_i$ 为一个具有仿射转移态射的概形有向系的极限。
假设 $X_i$ 对所有 $i \in I$ 都拟紧拟分离。
\begin{enumerate}
\item $X_\etale$ 上的有限局部常值层范畴是各 $(X_i)_\etale$ 上有限局部常值层
范畴的余极限。
\item $X_\etale$ 上的有限局部常值阿贝尔层范畴是各 $(X_i)_\etale$ 上有限局部
常值阿贝尔层范畴的余极限。
\item 设 $\Lambda$ 为诺特环。有限型局部常值 $\Lambda$-模层在 $X_\etale$ 上的
范畴是有限型局部常值 $\Lambda$-模层在各 $(X_i)_\etale$ 上的范畴的余极限。
\end{enumerate}
\end{lemma}

\begin{proof}
由引理 \ref{etale-cohomology-lemma-category-is-colimit}，每种情形中的函子都是全忠实的。由同一引理，
为完成情形 (1) 的证明，只需证明如下结论：给定可构造层 $\mathcal{F}_i$，定义在
$X_i$ 上；若其拉回 $\mathcal{F}$ 定义在 $X$ 上，并且是有限局部常值的，则存在
$i' \geq i$，使得拉回 $\mathcal{F}_{i'}$（它来自 $\mathcal{F}_i$，并定义在
$X_{i'}$ 上）
是有限局部常值的。根据假设，存在 \'etale 覆盖
$\mathcal{U} = \{U_j \to X\}_{j \in J}$
，使得有 $\mathcal{F}|_{U_j} \cong \underline{S_j}$，其中 $S_j$ 为某个有限集。
可以假定 $U_j$ 对所有 $j \in J$ 都是仿射的。由于 $X$ 拟紧，可以假定 $J$
有限。由引理 \ref{etale-cohomology-lemma-colimit-affine-sites}，可找到 $i' \geq i$ 和 \'etale 覆盖
$\mathcal{U}_{i'} = \{U_{i', j} \to X_{i'}\}_{j \in J}$，其到 $X$ 的基变换为
$\mathcal{U}$。于是，可构造层 $\mathcal{F}_{i'}|_{U_{i', j}}$ 和
$\underline{S_j}$ 定义在 $(U_{i', j})_\etale$ 上，并且拉回到 $U_j$ 后同构。因此，
增大 $i'$ 后可使 $\mathcal{F}_{i'}|_{U_{i', j}}$ 和 $\underline{S_j}$ 同构。
所以 $\mathcal{F}_{i'}$ 是有限局部常值的。情形 (2) 和 (3) 的证明完全相同。
\end{proof}

\begin{lemma}
\label{etale-cohomology-lemma-irreducible-subsheaf-constant-zero}
设 $X$ 为具有泛点 $\eta$ 的不可约概形。
\begin{enumerate}
\item 设 $S' \subset S$ 为集合的包含。若在
$\underline{S'} \subset \mathcal{G} \subset \underline{S}$
在 $\Sh(X_\etale)$ 中成立，且 $S' = \mathcal{G}_{\overline{\eta}}$，则
$\mathcal{G} = \underline{S'}$。
\item 设 $A' \subset A$ 为阿贝尔群的包含。若在
$\underline{A'} \subset \mathcal{G} \subset \underline{A}$
在 $\textit{Ab}(X_\etale)$ 中成立，且 $A' = \mathcal{G}_{\overline{\eta}}$，则
$\mathcal{G} = \underline{A'}$。
\item 设 $M' \subset M$ 为环 $\Lambda$ 上模的包含。若
$\underline{M'} \subset \mathcal{G} \subset \underline{M}$ 在
$\textit{Mod}(X_\etale, \underline{\Lambda})$ 中成立，且
$M' = \mathcal{G}_{\overline{\eta}}$，则 $\mathcal{G} = \underline{M'}$。
\end{enumerate}
\end{lemma}

\begin{proof}
这是因为，对每个 \'etale 态射 $U \to X$，只要 $U \not = \emptyset$，点
$\eta$ 都在其像中。
\end{proof}

\begin{lemma}
\label{etale-cohomology-lemma-push-constant-sheaf-from-open}
设 $X$ 为函数域为 $K$ 的整正规概形。设 $E$ 为集合。
\begin{enumerate}
\item 设 $g : \Spec(K) \to X$ 为泛点的包含。则
$g_*\underline{E} = \underline{E}$。
\item 设 $j : U \to X$ 为非空开集的包含。则
$j_*\underline{E} = \underline{E}$。
\end{enumerate}
\end{lemma}

\begin{proof}
证明 (1)。设 $x \in X$ 为一点。令
$\mathcal{O}^{sh}_{X, \overline{x}}$
为 $\mathcal{O}_{X, x}$ 的一个严格 Hensel 化。由代数进阶，引理
\ref{more-algebra-lemma-henselization-normal}，可知
$\mathcal{O}^{sh}_{X, \overline{x}}$ 是正规整环。因此
$\Spec(K) \times_X \Spec(\mathcal{O}^{sh}_{X, \overline{x}})$ 不可约。
由此，茎 $(g_*\underline{E})_{\overline{x}}$ 等于 $E$，见定理
\ref{etale-cohomology-theorem-higher-direct-images}。

\medskip\noindent
证明 (2)。由于 $g$ 通过 $j$ 分解，存在态射
$j_*\underline{E} \to g_*\underline{E}$。该态射是单射，因为对于每个概形 $V$
（它在 $X$ 上 \'etale），集合 $\Spec(K) \times_X V$ 在 $U \times_X V$ 中稠密。
另一方面，有态射 $\underline{E} \to j_*\underline{E}$，因而得出结论。
\end{proof}

\begin{lemma}
\label{etale-cohomology-lemma-zero-in-generic-point}
设 $X$ 为拟紧拟分离概形。设 $\eta \in X$ 为 $X$ 的某个不可约分支的泛点。
\begin{enumerate}
\item 设 $\mathcal{F}$ 为 $X_\etale$ 上的挠阿贝尔层，其茎
$\mathcal{F}_{\overline{\eta}}$ 为零。则 $\mathcal{F} = \colim \mathcal{F}_i$
是可构造阿贝尔层 $\mathcal{F}_i$ 的滤过余极限，使得对每个 $i$，
$\mathcal{F}_i$ 的支集都包含在某个不含 $\eta$ 的闭子概形中。
\item 设 $\Lambda$ 为诺特环，$\mathcal{F}$ 为 $\Lambda$-模层，定义在 $X_\etale$ 上，
其茎 $\mathcal{F}_{\overline{\eta}}$ 为零。则
$\mathcal{F} = \colim \mathcal{F}_i$
是可构造 $\Lambda$-模层 $\mathcal{F}_i$ 的滤过余极限，使得对每个 $i$，
$\mathcal{F}_i$ 的支集都包含在某个不含 $\eta$ 的闭子概形中。
\end{enumerate}
\end{lemma}

\begin{proof}
证明 (1)。由引理 \ref{etale-cohomology-lemma-torsion-colimit-constructible}，可以写成
$\mathcal{F} = \colim_{i \in I} \mathcal{F}_i$，其中 $\mathcal{F}_i$ 是可构造
阿贝尔层。选取 $i \in I$。由于根据假设 $\mathcal{F}|_\eta$ 为零，可知存在
$i'(i) \geq i$，使得 $\mathcal{F}_i|_\eta \to \mathcal{F}_{i'(i)}|_\eta$ 为零，
见引理 \ref{etale-cohomology-lemma-colimit-constructible}。于是
$\mathcal{G}_i = \Im(\mathcal{F}_i \to \mathcal{F}_{i'(i)})$ 是可构造阿贝尔层
（引理 \ref{etale-cohomology-lemma-constructible-abelian}），且其在 $\eta$ 处的茎为零。因此，令
$E_i$ 为 $\mathcal{G}_i$ 的支集；它是 $X$ 的不含 $\eta$ 的可构造子集。由于 $\eta$
是 $X$ 的某个不可约分支的泛点，由拓扑学，引理
\ref{topology-lemma-generic-point-in-constructible} 可知
$\eta \not \in Z_i = \overline{E_i}$。定义新的有向集 $I'$，其底集取为 $I$，
并用如下规则定义序：$i_1$ 大于或等于 $i_2$，当且仅当
$i_1 \geq i'(i_2)$。于是，层 $\mathcal{G}_i$ 构成 $I'$ 上的一个系，其余极限为
$\mathcal{F}$，证明完成。

\medskip\noindent
情形 (2) 的证明完全相同，故略去。
\end{proof}








\section{诺特概形上的可构造层}
\label{etale-cohomology-section-noetherian-constructible}

\noindent
若 $X$ 是诺特概形，则任意局部闭子集都是可构造局部闭子集
（拓扑学，引理 \ref{topology-lemma-constructible-Noetherian-space}）。
因此，阿贝尔层 $\mathcal{F}$ 定义在 $X_\etale$ 上时，它是可构造的，当且仅当存在有限分划
$X = \coprod X_i$，使得 $\mathcal{F}|_{X_i}$ 是有限局部常值的。
（按照约定，拓扑空间的分划以局部闭子集为各部分；见拓扑学，
第 \ref{topology-section-stratifications} 节。）换言之，在定义
\ref{etale-cohomology-definition-constructible} 中可以省略形容词“可构造”。事实上，诺特概形上的
可构造层范畴还具有一些附加性质，我们将在本节中逐一列出。

\begin{proposition}
\label{etale-cohomology-proposition-constructible-over-noetherian}
设 $X$ 为诺特概形，$\Lambda$ 为诺特环。
\begin{enumerate}
\item 可构造集合层的任意子层或商层都是可构造的。
\item $X_\etale$ 上的可构造阿贝尔层范畴是
$\textit{Ab}(X_\etale)$ 的（强）Serre 子范畴。特别地，$X_\etale$ 上可构造
阿贝尔层的任意子层和商层都是可构造的。
\item 可构造 $\Lambda$-模层在 $X_\etale$ 上的范畴是
$\textit{Mod}(X_\etale, \Lambda)$ 的（强）Serre 子范畴。特别地，
可构造 $\Lambda$-模层在 $X_\etale$ 上的任意子模层和商模层都是可构造的。
\end{enumerate}
\end{proposition}

\begin{proof}
证明 (1)。设 $\mathcal{G} \subset \mathcal{F}$，其中 $\mathcal{F}$ 是
$X_\etale$ 上的可构造集合层。设 $\eta \in X$ 为 $X$ 的某个不可约分支的泛点。
由诺特归纳法，只需找到一个开邻域 $U$，它包含 $\eta$，且使 $\mathcal{G}|_U$ 局部常值。
为此，可以替换 $X$，代之以 $\eta$ 的一个 \'etale 邻域。因而可以假定
$\mathcal{F}$ 是常值层且 $X$ 不可约。

\medskip\noindent
设 $\mathcal{F} = \underline{S}$，其中 $S$ 为某个有限集。于是
$S' = \mathcal{G}_{\overline{\eta}} \subset S$；记
$S' = \{s_1, \ldots, s_t\}$。选取一个 \'etale 邻域
$(U, \overline{u})$，它是 $\overline{\eta}$ 的邻域；再选取截面
$\sigma_1, \ldots, \sigma_t \in \mathcal{G}(U)$，使其映到 $s_i$，后者视为
$\mathcal{G}_{\overline{\eta}} \subset S$ 中的元素。由于 $\sigma_i$ 映到元素
$s_i \in S' \subset S = \Gamma(X, \mathcal{F})$，故 $\sigma_i$ 在
$U \times_X U$ 上的两个拉回作为 $\mathcal{G}$ 的截面相同。由 $\mathcal{G}$ 的
层条件可知，$\sigma_i$ 来自 $\mathcal{G}$ 在开集
$\Im(U \to X)$ 上的一个截面；这是 $X$ 的开集。缩小 $X$ 后，可以假定
$\underline{S'} \subset \mathcal{G} \subset \underline{S}$。于是由引理
\ref{etale-cohomology-lemma-irreducible-subsheaf-constant-zero} 可知
$\underline{S'} = \mathcal{G}$。

\medskip\noindent
设 $\mathcal{F} \to \mathcal{Q}$ 为满射，其中 $\mathcal{F}$ 是 $X_\etale$ 上的
可构造集合层。令
$\mathcal{G} = \mathcal{F} \times_\mathcal{Q} \mathcal{F}$。
由证明的第一部分，$\mathcal{G}$ 作为 $\mathcal{F} \times \mathcal{F}$ 的子层是
可构造的。这又蕴含 $\mathcal{Q}$ 可构造；见引理
\ref{etale-cohomology-lemma-constructible-abelian}。

\medskip\noindent
证明 (3)。我们已经知道，可构造模层构成一个弱 Serre 子范畴；见引理
\ref{etale-cohomology-lemma-constructible-abelian}。因此，只需证明关于子模层的断言。

\medskip\noindent
设 $\mathcal{G} \subset \mathcal{F}$ 是某个可构造 $\Lambda$-模层的子模层，
该层定义在 $X_\etale$ 上。设 $\eta \in X$ 为 $X$ 的某个不可约分支的泛点。
由诺特归纳法，只需找到一个开邻域 $U$，它包含 $\eta$，且使 $\mathcal{G}|_U$ 局部常值。
为此，可以替换 $X$，代之以 $\eta$ 的一个 \'etale 邻域。因而可以假定
$\mathcal{F}$ 是常值层且 $X$ 不可约。

\medskip\noindent
设 $\mathcal{F} = \underline{M}$，其中所取的有限 $\Lambda$-模记作 $M$。于是
$M' = \mathcal{G}_{\overline{\eta}} \subset M$。选取有限多个元素
$s_1, \ldots, s_t$，使其生成 $M'$ 这个 $\Lambda$-模。（这是可能的，因为
$\Lambda$ 是诺特的且 $M$ 有限。）选取一个 \'etale 邻域
$(U, \overline{u})$，它是 $\overline{\eta}$ 的邻域；再选取截面
$\sigma_1, \ldots, \sigma_t \in \mathcal{G}(U)$，使其映到 $s_i$，后者视为
$\mathcal{G}_{\overline{\eta}} \subset M$ 中的元素。由于 $\sigma_i$ 映到元素
$s_i \in M' \subset M = \Gamma(X, \mathcal{F})$，故 $\sigma_i$ 在
$U \times_X U$ 上的两个拉回作为 $\mathcal{G}$ 的截面相同。由 $\mathcal{G}$ 的
层条件可知，$\sigma_i$ 来自 $\mathcal{G}$ 在开集
$\Im(U \to X)$ 上的一个截面；这是 $X$ 的开集。缩小 $X$ 后，可以假定
$\underline{M'} \subset \mathcal{G} \subset \underline{M}$。于是由引理
\ref{etale-cohomology-lemma-irreducible-subsheaf-constant-zero} 可知
$\underline{M'} = \mathcal{G}$。

\medskip\noindent
证明 (2)。这以通常方式由 (3) 得出。细节略去。
\end{proof}

\noindent
下面的引理说明，$X$ 上可构造层的阿贝尔范畴中的每个对象都是“诺特的”，
即对子对象满足升链条件（a.c.c.）。

\begin{lemma}
\label{etale-cohomology-lemma-constructible-over-noetherian-noetherian}
设 $X$ 为诺特概形，$\Lambda$ 为诺特环。考虑包含关系
$$
\mathcal{F}_1 \subset \mathcal{F}_2 \subset \mathcal{F}_3 \subset \ldots
\subset \mathcal{F}
$$
，它们位于集合层、阿贝尔群层或 $\Lambda$-模层的范畴中。若 $\mathcal{F}$
可构造，则对某个 $n$ 有
$\mathcal{F}_n = \mathcal{F}_{n + 1} = \mathcal{F}_{n + 2} = \ldots$。
\end{lemma}

\begin{proof}
由命题 \ref{etale-cohomology-proposition-constructible-over-noetherian} 可知，
$\mathcal{F}_i$ 和 $\colim \mathcal{F}_i$ 都可构造。于是该引理由引理
\ref{etale-cohomology-lemma-colimit-constructible} 得出。
\end{proof}

\begin{lemma}
\label{etale-cohomology-lemma-constructible-maps-into-constant}
设 $X$ 为诺特概形。
\begin{enumerate}
\item 设 $\mathcal{F}$ 为 $X_\etale$ 上的可构造集合层。存在层的单射
$$
\mathcal{F} \longrightarrow
\prod\nolimits_{i = 1, \ldots, n} f_{i, *}\underline{E_i}
$$
，其中 $f_i : Y_i \to X$ 是有限态射，$E_i$ 是有限集。
\item 设 $\mathcal{F}$ 为 $X_\etale$ 上的可构造阿贝尔层。存在阿贝尔层的单射
$$
\mathcal{F} \longrightarrow
\bigoplus\nolimits_{i = 1, \ldots, n} f_{i, *}\underline{M_i}
$$
，其中 $f_i : Y_i \to X$ 是有限态射，$M_i$ 是有限阿贝尔群。
\item 设 $\Lambda$ 为诺特环，设 $\mathcal{F}$ 为可构造 $\Lambda$-模层，
定义在 $X_\etale$ 上。存在模层的单射
$$
\mathcal{F} \longrightarrow
\bigoplus\nolimits_{i = 1, \ldots, n} f_{i, *}\underline{M_i}
$$
，其中 $f_i : Y_i \to X$ 是有限态射，$M_i$ 是有限 $\Lambda$-模。
\end{enumerate}
此外，可以假定每个 $Y_i$ 都不可约且既约，并映满某个不可约既约闭子概形
$Z_i \subset X$，且 $Y_i \to Z_i$ 在 $Z_i$ 的某个非空开集上有限 \'etale。
\end{lemma}

\begin{proof}
证明 (1)。由于 $\mathcal{F}$ 的子层满足升链条件（引理
\ref{etale-cohomology-lemma-constructible-over-noetherian-noetherian}），只需证明：对每个点
$x \in X$，可以找到态射 $\varphi : \mathcal{F} \to f_*\underline{E}$，其中
$f : Y \to X$ 有限且 $E$ 为有限集，并使
$\varphi_{\overline{x}} : \mathcal{F}_{\overline{x}} \to
(f_*S)_{\overline{x}}$ 为单射。（也可以避开这一论证：选取引理
\ref{etale-cohomology-lemma-constructible-quasi-compact-quasi-separated} 中那样的 $X$ 的一个分划，
并对每个部分的每个不可约分支构造一个 $Y_i \to X$。）令 $Z \subset X$ 为
$\overline{\{x\}}$ 上的诱导既约概形结构（概形，定义
\ref{schemes-definition-reduced-induced-scheme}）。由于 $\mathcal{F}$ 可构造，
存在有限可分扩张 $K/\kappa(x)$，使得 $\mathcal{F}|_{\Spec(K)}$ 是以
$E$ 为值的常值层，其中 $E$ 为某个有限集。令 $Y \to Z$ 为 $Z$ 在
$\Spec(K)$ 中的正规化。由态射，引理
\ref{morphisms-lemma-normal-normalization}，$Y$ 是正规整概形。由于
$K/\kappa(x)$ 是有限扩张，显然 $K$ 是 $Y$ 的函数域。用
$g : \Spec(K) \to Y$ 表示该包含。态射
$\mathcal{F}|_{\Spec(K)} \to \underline{E}$ 与态射
$\mathcal{F}|_Y \to g_*\underline{E} = \underline{E}$ 伴随（引理
\ref{etale-cohomology-lemma-push-constant-sheaf-from-open}）。后者又与态射
$\varphi : \mathcal{F} \to f_*\underline{E}$ 伴随。注意，$\varphi$ 在几何点
$\overline{x}$ 处的茎是单射：可以选取一个提升 $\overline{y} \in Y$，
它是 $\overline{x}$ 的提升，而交换图
$$
\xymatrix{
\mathcal{F}_{\overline{x}} \ar@{=}[r] \ar[d] &
(\mathcal{F}|_Y)_{\overline{y}} \ar@{=}[d] \\
(f_*\underline{E})_{\overline{x}} \ar[r] &
\underline{E}_{\overline{y}}
}
$$
证明了单射性。不过，证明尚未完成，因为态射 $f : Y \to Z$ 是整的，
但一般并非有限的\footnote{若 $X$ 是 Nagata 概形，例如域上的有限型概形，
则 $Y \to Z$ 是有限的。}。

\medskip\noindent
为解决前一段末句所述的问题，写成 $Y = \lim_{i \in I} Y_i$，其中 $Y_i$
不可约、整，且在 $Z$ 上有限。具体而言，把性质，引理
\ref{properties-lemma-integral-algebra-directed-colimit-finite} 应用于
$f_*\mathcal{O}_Y$，把它视为 $\mathcal{O}_Z$-代数层，再应用函子
$\underline{\Spec}_Z$。于是由引理 \ref{etale-cohomology-lemma-relative-colimit} 有
$f_*\underline{E} = \colim f_{i, *}\underline{E}$。由引理
\ref{etale-cohomology-lemma-constructible-is-compact}，态射 $\mathcal{F} \to f_*\underline{E}$
通过 $f_{i, *}\underline{E}$ 分解，其中 $i$ 是某个指标。由于 $Y_i \to Z$ 是整概形之间的
有限态射，且该态射诱导的函数域扩张有限可分，故该态射在 $Z$ 的某个非空开集上
有限 \'etale（使用代数，引理 \ref{algebra-lemma-smooth-at-generic-point}；
细节略去）。这就完成了 (1) 的证明。

\medskip\noindent
(2) 和 (3) 的证明与 (1) 的证明完全相同。
\end{proof}

\noindent
在下面的引理中，我们用一个标准技巧把一个非常一般的陈述归约到诺特情形。

\begin{lemma}
\label{etale-cohomology-lemma-constructible-maps-into-constant-general}
\begin{reference}
\cite[Exposee IX, Proposition 2.14]{SGA4}
\end{reference}
设 $X$ 为拟紧拟分离概形。
\begin{enumerate}
\item 设 $\mathcal{F}$ 为 $X_\etale$ 上的可构造集合层。存在层的单射
$$
\mathcal{F} \longrightarrow
\prod\nolimits_{i = 1, \ldots, n} f_{i, *}\underline{E_i}
$$
，其中 $f_i : Y_i \to X$ 是有限且有限表现的态射，$E_i$ 是有限集。
\item 设 $\mathcal{F}$ 为 $X_\etale$ 上的可构造阿贝尔层。存在阿贝尔层的单射
$$
\mathcal{F} \longrightarrow
\bigoplus\nolimits_{i = 1, \ldots, n} f_{i, *}\underline{M_i}
$$
，其中 $f_i : Y_i \to X$ 是有限且有限表现的态射，$M_i$ 是有限阿贝尔群。
\item 设 $\Lambda$ 为诺特环，设 $\mathcal{F}$ 为可构造 $\Lambda$-模层，
定义在 $X_\etale$ 上。存在模层的单射
$$
\mathcal{F} \longrightarrow
\bigoplus\nolimits_{i = 1, \ldots, n} f_{i, *}\underline{M_i}
$$
，其中 $f_i : Y_i \to X$ 是有限且有限表现的态射，$M_i$ 是有限 $\Lambda$-模。
\end{enumerate}
\end{lemma}

\begin{proof}
我们通过绝对诺特逼近把该引理归约到诺特情形。具体而言，由极限，命题
\ref{limits-proposition-approximate}，可以写成
$X = \lim_{t \in T} X_t$，其中每个 $X_t$ 在 $\Spec(\mathbf{Z})$ 上有限型，
且转移态射为仿射态射。由引理 \ref{etale-cohomology-lemma-category-is-colimit}，可构造层
（集合层、阿贝尔群层或 $\Lambda$-模层）在 $X_\etale$ 上的范畴是各
$X_t$ 上相应范畴的余极限。因此，我们的可构造层 $\mathcal{F}$ 是同类可构造层
$\mathcal{F}_t$ 的拉回；后者定义在 $X_t$ 上，其中 $t$ 是某个指标。随后对诺特情形应用引理
\ref{etale-cohomology-lemma-constructible-maps-into-constant}，得到单射
$$
\mathcal{F}_t \longrightarrow
\prod\nolimits_{i = 1, \ldots, n} f_{i, *}\underline{E_i}
\quad\text{或}\quad
\mathcal{F}_t \longrightarrow
\bigoplus\nolimits_{i = 1, \ldots, n} f_{i, *}\underline{M_i}
$$
，它位于 $X_t$ 上，其中 $f_i : Y_i \to X_t$ 是某些有限态射。由于 $X_t$
是诺特的，态射 $f_i$ 是有限表现的。由于拉回正合，而且 $f_{i, *}$ 的形成与
基变换可交换（引理 \ref{etale-cohomology-lemma-finite-pushforward-commutes-with-base-change}），
结论得证。
\end{proof}

\begin{lemma}
\label{etale-cohomology-lemma-support-in-subset}
设 $X$ 为诺特概形，$E \subset X$ 为在特殊化下封闭的子集。
\begin{enumerate}
\item 设 $\mathcal{F}$ 为 $X_\etale$ 上的挠阿贝尔层，其支集包含于 $E$。则
$\mathcal{F} = \colim \mathcal{F}_i$ 是可构造阿贝尔层 $\mathcal{F}_i$ 的滤过
余极限，且对每个 $i$，$\mathcal{F}_i$ 的支集包含在某个包含于 $E$ 的闭子集中。
\item 设 $\Lambda$ 为诺特环，$\mathcal{F}$ 为 $\Lambda$-模层，定义在 $X_\etale$ 上，
其支集包含于 $E$。则 $\mathcal{F} = \colim \mathcal{F}_i$ 是可构造
$\Lambda$-模层 $\mathcal{F}_i$ 的滤过余极限，且对每个 $i$，$\mathcal{F}_i$
的支集包含在某个包含于 $E$ 的闭子集中。
\end{enumerate}
\end{lemma}

\begin{proof}
证明 (1)。由引理 \ref{etale-cohomology-lemma-torsion-colimit-constructible}，可以写成
$\mathcal{F} = \colim_{i \in I} \mathcal{F}_i$，其中 $\mathcal{F}_i$ 是可构造
阿贝尔层。由命题 \ref{etale-cohomology-proposition-constructible-over-noetherian}，像
$\mathcal{F}'_i \subset \mathcal{F}$（它来自态射
$\mathcal{F}_i \to \mathcal{F}$）是可构造的。因此
$\mathcal{F} = \colim \mathcal{F}'_i$，而 $\mathcal{F}'_i$ 的支集包含在 $E$ 中。
由于 $\mathcal{F}'_i$ 的支集是可构造的（根据可构造层的定义），故其闭包也包含于
$E$；例如见拓扑学，引理
\ref{topology-lemma-constructible-stable-specialization-closed}。

\medskip\noindent
情形 (2) 的证明完全相同，故略去。
\end{proof}








\section{特殊化与 \'etale 层}
\label{etale-cohomology-section-specialization}

\noindent
拓扑图景如下：设 $X$ 为拓扑空间，$x' \leadsto x$ 为 $X$ 中点的一个特殊化。
则 $x$ 的每个开邻域都包含 $x'$。因此，对任意层 $\mathcal{F}$（定义在 $X$ 上），
存在茎的{\it 特殊化映射}
$$
sp : \mathcal{F}_x \longrightarrow \mathcal{F}_{x'}
$$
，它把二元组 $(U, s)$ 在 $\mathcal{F}_x$ 中的等价类映到
二元组 $(U, s)$ 在 $\mathcal{F}_{x'}$ 中的等价类。关于用二元组的等价类描述茎，
见层，第 \ref{sheaves-section-stalks} 节。这个映射当然关于 $\mathcal{F}$ 具有
函子性；换言之，$sp$ 是函子之间的变换。

\medskip\noindent
对 \'etale 拓扑中的层，可以仿照这一构造；见
\cite[Exposee VIII, 7.7, page 397]{SGA4}。为此，设有概形 $S$、几何点
$\overline{s}$（它属于 $S$），以及几何点 $\overline{t}$（它属于
$\Spec(\mathcal{O}^{sh}_{S, \overline{s}})$）。对任意层 $\mathcal{F}$（定义在
$S_\etale$ 上），我们将构造
{\it 特殊化映射}
$$
sp : \mathcal{F}_{\overline{s}} \longrightarrow \mathcal{F}_{\overline{t}}
$$
这里略有语言上的滥用：不应写成 $\mathcal{F}_{\overline{t}}$，而应写成
$\mathcal{F}_{p(\overline{t})}$，其中
$p : \Spec(\mathcal{O}^{sh}_{S, \overline{s}}) \to S$ 是典范态射。回忆
$$
\mathcal{F}_{\overline{s}} = \colim_{(U, \overline{u})} \mathcal{F}(U)
$$
，其中余极限取遍所有 \'etale 邻域 $(U, \overline{u})$，它们是
$(S, \overline{s})$ 的邻域；见第 \ref{etale-cohomology-section-stalks} 节。由于
$\mathcal{O}^{sh}_{S, \overline{s}}$ 是结构层的茎，对每个 \'etale 邻域
$(U, \overline{u})$（它是 $(S, \overline{s})$ 的邻域），都有典范映射
$\mathcal{O}_{U, u} \to \mathcal{O}^{sh}_{S, \overline{s}}$。因此得到唯一分解
$$
\Spec(\mathcal{O}^{sh}_{S, \overline{s}}) \to U \to S
$$
若 $\overline{v}$ 表示 $\overline{t}$ 在 $U$ 中的像，则
$(U, \overline{v})$ 是 $(S, \overline{t})$ 的一个 \'etale 邻域。该构造定义了一个
函子，它从 $(S, \overline{s})$ 的 \'etale 邻域范畴到
$(S, \overline{t})$ 的 \'etale 邻域范畴。因此，可以定义映射
$sp : \mathcal{F}_{\overline{s}} \to \mathcal{F}_{\overline{t}}$
：它把 $(U, \overline{u}, \sigma)$ 的等价类（其中
$\sigma \in \mathcal{F}(U)$）映到 $(U, \overline{v}, \sigma)$ 的等价类。

\medskip\noindent
设 $K \in D(S_\etale)$。对如上的 $\overline{s}$ 和 $\overline{t}$，有
{\it 特殊化映射}
$$
sp : K_{\overline{s}} \longrightarrow K_{\overline{t}}
\quad\text{于}\quad D(\textit{Ab})
$$
具体而言，若 $K$ 由阿贝尔层复形 $\mathcal{F}^\bullet$ 表示，则直接取映射
$$
K_{\overline{s}} = \mathcal{F}^\bullet_{\overline{s}}
\longrightarrow
\mathcal{F}^\bullet_{\overline{t}} = K_{\overline{t}}
$$
，其各项由上面构造的层的特殊化映射给出。由于茎函子正合，这一映射与表示
$K$ 的复形之选择无关（即逐项取复形的茎在导出范畴上有良好定义）。

\medskip\noindent
显然，该构造关于层 $\mathcal{F}$ 具有函子性，该层定义在 $S_\etale$ 上。
若把茎函子看成拓扑斯的态射
$\overline{s}, \overline{t} : \textit{Sets} \to \Sh(S_\etale)$，则可以把
$sp$ 看成如下 $2$-态射：
$$
\xymatrix{
\textit{Sets}
\rrtwocell^{\overline{t}}_{\overline{s}}{\ sp}
&
&
\Sh(S_\etale)
}
$$
。

\begin{remark}[sp 的另一种描述]
\label{etale-cohomology-remark-alternative-sp}
设 $S$、$\overline{s}$ 和 $\overline{t}$ 如上。描述特殊化映射的另一种方式是使用
$$
\mathcal{F}_{\overline{s}} =
\Gamma(\Spec(\mathcal{O}^{sh}_{S, \overline{s}}), p^{-1}\mathcal{F})
\quad\text{且}\quad
\mathcal{F}_{\overline{t}} = \Gamma(\overline{t},
\overline{t}^{-1}p^{-1}\mathcal{F})
$$
。第一个等式由把定理 \ref{etale-cohomology-theorem-higher-direct-images} 应用于
$\text{id}_S : S \to S$ 得出，第二个等式由引理
\ref{etale-cohomology-lemma-stalk-pullback} 得出。于是，可以把 $sp$ 看成映射
$$
sp :
\mathcal{F}_{\overline{s}} =
\Gamma(\Spec(\mathcal{O}^{sh}_{S, \overline{s}}), p^{-1}\mathcal{F})
\xrightarrow{\text{pullback by }\overline{t}}
\Gamma(\overline{t}, \overline{t}^{-1}p^{-1}\mathcal{F}) =
\mathcal{F}_{\overline{t}}
$$
\end{remark}

\begin{remark}[sp 的又一种描述]
\label{etale-cohomology-remark-another-sp}
设 $S$、$\overline{s}$ 和 $\overline{t}$ 如上。还可以使用唯一态射
$$
c :
\Spec(\mathcal{O}^{sh}_{S, \overline{t}})
\longrightarrow
\Spec(\mathcal{O}^{sh}_{S, \overline{s}})
$$
，它位于 $S$ 上，并且与给定态射
$\overline{t} \to \Spec(\mathcal{O}^{sh}_{S, \overline{s}})$ 及态射
$\overline{t} \to \Spec(\mathcal{O}^{sh}_{t, \overline{t}})$ 相容。把代数，引理
\ref{algebra-lemma-map-into-henselian-colimit} 应用于 $\mathcal{O}_{S, s}$、
$\mathcal{O}^{sh}_{S, \overline{t}}$ 和
$\mathcal{O}^{sh}_{S, \overline{s}} \to \kappa(\overline{t})$，即可得到所示箭头的
唯一性与存在性。于是有
$$
sp :
\mathcal{F}_{\overline{s}} =
\Gamma(\Spec(\mathcal{O}^{sh}_{S, \overline{s}}), \mathcal{F})
\xrightarrow{\text{pullback by }c}
\Gamma(\Spec(\mathcal{O}^{sh}_{S, \overline{t}}), \mathcal{F}) =
\mathcal{F}_{\overline{t}}
$$
（采用显然的记号约定）。事实上，该过程也适用于对象 $K$，该对象属于
$D(S_\etale)$：
$K$ 的特殊化映射等于映射
$$
sp :
K_{\overline{s}} =
R\Gamma(\Spec(\mathcal{O}^{sh}_{S, \overline{s}}), K)
\xrightarrow{\text{pullback by }c}
R\Gamma(\Spec(\mathcal{O}^{sh}_{S, \overline{t}}), K) =
K_{\overline{t}}
$$
。这些等号成立，是因为在严格 Hensel 概形
$\Spec(\mathcal{O}^{sh}_{S, \overline{s}})$ 和
$\Spec(\mathcal{O}^{sh}_{S, \overline{t}})$ 上取全局截面是正合的（并且分别等同于
在 $\overline{s}$ 和 $\overline{t}$ 处取茎）；因此，不会出现与 $K$ 可能无界这一
事实有关的微妙问题。
\end{remark}

\begin{remark}[提升特殊化]
\label{etale-cohomology-remark-can-lift}
设 $S$ 为概形，$t \leadsto s$ 为 $S$ 上点的一个特殊化。选取几何点
$\overline{t}$ 和 $\overline{s}$，它们分别位于 $t$ 和 $s$ 上方。由概形，引理
\ref{schemes-lemma-specialize-points}，$t$ 对应于
$\Spec(\mathcal{O}_{S, s})$ 的一个点；又因为
$\mathcal{O}_{S, s} \to \mathcal{O}^{sh}_{S, \overline{s}}$ 忠实平坦，故可找到点
$t' \in \Spec(\mathcal{O}^{sh}_{S, \overline{s}})$，它映到 $t$。由于
$\Spec(\mathcal{O}^{sh}_{S, \overline{s}})$ 是 $S$ 上 \'etale 概形的极限，
$\kappa(t')/\kappa(t)$ 是可分代数扩张（当然通常不是有限的）。由于
$\kappa(\overline{t})$ 是代数闭的，可以选取嵌入
$\kappa(t') \to \kappa(\overline{t})$，它是 $\kappa(t)$ 上扩张之间的嵌入。
这一选择给出交换图
$$
\xymatrix{
\overline{t} \ar[d] \ar[r] &
\Spec(\mathcal{O}^{sh}_{S, \overline{s}}) \ar[d] &
\overline{s} \ar[l] \ar[d] \\
t \ar[r] & S & s \ar[l]
}
$$
，它由点和几何点组成。因此，若 $t \leadsto s$，总可以把 $\overline{t}$“提升”为
$S$ 在 $\overline{s}$ 处的严格 Hensel 化的一个几何点，并得到上述特殊化映射。
\end{remark}

\begin{lemma}
\label{etale-cohomology-lemma-specialization-map-pullback}
设 $g : S' \to S$ 为概形态射，$\mathcal{F}$ 为 $S_\etale$ 上的层。设
$\overline{s}'$ 为 $S'$ 的几何点，$\overline{t}'$ 为
$\Spec(\mathcal{O}^{sh}_{S', \overline{s}'})$ 的几何点。记
$\overline{s} = g(\overline{s}')$、$\overline{t} = h(\overline{t}')$，其中典范态射为
$h : \Spec(\mathcal{O}^{sh}_{S', \overline{s}'}) \to
\Spec(\mathcal{O}^{sh}_{S, \overline{s}})$。对任意层 $\mathcal{F}$（定义在
$S_\etale$ 上），特殊化映射
$$
sp :
(g^{-1}\mathcal{F})_{\overline{s}'}
\longrightarrow
(g^{-1}\mathcal{F})_{\overline{t}'}
$$
在下列等同下等于特殊化映射
$sp : \mathcal{F}_{\overline{s}} \to \mathcal{F}_{\overline{t}}$：
$(g^{-1}\mathcal{F})_{\overline{s}'} = \mathcal{F}_{\overline{s}}$ 以及
$(g^{-1}\mathcal{F})_{\overline{t}'} = \mathcal{F}_{\overline{t}}$。
这些等同来自引理 \ref{etale-cohomology-lemma-stalk-pullback}。
\end{lemma}

\begin{proof}
略。
\end{proof}

\begin{lemma}
\label{etale-cohomology-lemma-characterize-locally-constant}
设 $S$ 为概形，并且 $S$ 的每个拟紧开集都只有有限个不可约分支（例如，$S$ 的
底拓扑空间为诺特空间，或 $S$ 局部诺特）。设 $\mathcal{F}$ 为 $S_\etale$ 上的集合层。
下列条件等价：
\begin{enumerate}
\item $\mathcal{F}$ 是有限局部常值的；
\item $\mathcal{F}$ 的所有茎都是有限集，并且所有特殊化映射
$sp : \mathcal{F}_{\overline{s}} \to \mathcal{F}_{\overline{t}}$
都是双射。
\end{enumerate}
\end{lemma}

\begin{proof}
假设 (2)。设 $\overline{s}$ 为 $S$ 的几何点，位于 $s \in S$ 上方。为证明 (1)，
必须找到一个 \'etale 邻域 $(U, \overline{u})$，它是 $(S, \overline{s})$ 的邻域，并使得
$\mathcal{F}|_U$ 是常值层。可以且确实假定 $S$ 是仿射的。

\medskip\noindent
由于 $\mathcal{F}_{\overline{s}}$ 有限，可以选取 $(U, \overline{u})$、
$n \geq 0$ 以及两两不同的元素
$\sigma_1, \ldots, \sigma_n \in \mathcal{F}(U)$，使得
$\{\sigma_1, \ldots, \sigma_n\} \subset \mathcal{F}(U)$ 双射地映到
$\mathcal{F}_{\overline{s}}$，所用映射是
$\mathcal{F}(U) \to \mathcal{F}_{\overline{s}}$。考虑映射
$$
\varphi : \underline{\{1, \ldots, n\}} \longrightarrow \mathcal{F}|_U
$$
，它定义在 $U_\etale$ 上，并由 $\sigma_1, \ldots, \sigma_n$ 给出。按照构造，
该映射在 $\overline{u}$ 处的茎上是双射。考虑子集
$$
E = \{u' \in U \mid \varphi_{\overline{u}'}\text{ 是双射}\} \subset U
$$
。这里，$\overline{u}'$ 是 $U$ 的位于 $u'$ 上方的任意几何点（由注
\ref{etale-cohomology-remark-map-stalks}，该条件与选择无关）。点 $u \in U$ 是 $\overline{u}$ 的像，并
属于 $E$。由关于 $\mathcal{F}$ 的特殊化映射的假设、注
\ref{etale-cohomology-remark-can-lift} 以及引理 \ref{etale-cohomology-lemma-specialization-map-pullback}，
$E$ 在拓扑空间 $U$ 中对特殊化和一般化封闭。

\medskip\noindent
缩小 $U$ 后，还可以假定 $U$ 仿射。由下降，引理
\ref{descent-lemma-locally-finite-nr-irred-local-fppf}，$U$ 只有有限个不可约分支。
删去不经过 $u$ 的不可约分支后，可以假定 $U$ 的每个不可约分支都经过 $u$。
由于 $U$ 是 Sober 拓扑空间，故 $E = U$；再由定理
\ref{etale-cohomology-theorem-exactness-stalks} 可知 $\varphi$ 是同构。因此 (1) 成立。

\medskip\noindent
省略 (1) 蕴含 (2) 的证明。
\end{proof}

\begin{lemma}
\label{etale-cohomology-lemma-characterize-locally-constant-module}
设 $S$ 为概形，并且 $S$ 的每个拟紧开集都只有有限个不可约分支（例如，$S$ 的
底拓扑空间为诺特空间，或 $S$ 局部诺特）。设 $\Lambda$ 为诺特环。设
$\mathcal{F}$ 为 $\Lambda$-模层，定义在 $S_\etale$ 上。下列条件等价：
\begin{enumerate}
\item $\mathcal{F}$ 是有限型局部常值 $\Lambda$-模层；
\item $\mathcal{F}$ 的所有茎都是有限 $\Lambda$-模，并且所有特殊化映射
$sp : \mathcal{F}_{\overline{s}} \to \mathcal{F}_{\overline{t}}$
都是双射。
\end{enumerate}
\end{lemma}

\begin{proof}
该引理的证明与引理 \ref{etale-cohomology-lemma-characterize-locally-constant} 的证明相同。
假设 (2)。设 $\overline{s}$ 为 $S$ 的几何点，位于 $s \in S$ 上方。为证明 (1)，
必须找到一个 \'etale 邻域 $(U, \overline{u})$，它是 $(S, \overline{s})$ 的邻域，并使得
$\mathcal{F}|_U$ 是常值层。可以且确实假定 $S$ 是仿射的。

\medskip\noindent
由于 $M = \mathcal{F}_{\overline{s}}$ 是有限 $\Lambda$-模，且 $\Lambda$ 诺特，
可以选取一个表现
$$
\Lambda^{\oplus m} \xrightarrow{A} \Lambda^{\oplus n} \to M \to 0
$$
，其中 $A = (a_{ji})$ 是一个系数位于 $\Lambda$ 中的矩阵。可以选取
$(U, \overline{u})$ 以及元素
$\sigma_1, \ldots, \sigma_n \in \mathcal{F}(U)$，使得
$\sum a_{ji}\sigma_i = 0$ 在 $\mathcal{F}(U)$ 中成立，并且 $\sigma_i$ 在
$\mathcal{F}_{\overline{s}} = M$ 中的像，是 $\Lambda^n$ 的标准基元素在上述 $M$ 的表现中的
像。考虑映射
$$
\varphi : \underline{M} \longrightarrow \mathcal{F}|_U
$$
，它定义在 $U_\etale$ 上，并由 $\sigma_1, \ldots, \sigma_n$ 给出。按照构造，
该映射在 $\overline{u}$ 处的茎上是双射。考虑子集
$$
E = \{u' \in U \mid \varphi_{\overline{u}'}\text{ 是双射}\} \subset U
$$
。这里，$\overline{u}'$ 是 $U$ 的位于 $u'$ 上方的任意几何点（由注
\ref{etale-cohomology-remark-map-stalks}，该条件与选择无关）。点 $u \in U$ 是 $\overline{u}$ 的像，并
属于 $E$。由关于 $\mathcal{F}$ 的特殊化映射的假设、注
\ref{etale-cohomology-remark-can-lift} 以及引理 \ref{etale-cohomology-lemma-specialization-map-pullback}，
$E$ 在拓扑空间 $U$ 中对特殊化和一般化封闭。

\medskip\noindent
缩小 $U$ 后，还可以假定 $U$ 仿射。由下降，引理
\ref{descent-lemma-locally-finite-nr-irred-local-fppf}，$U$ 只有有限个不可约分支。
删去不经过 $u$ 的不可约分支后，可以假定 $U$ 的每个不可约分支都经过 $u$。
由于 $U$ 是 Sober 拓扑空间，故 $E = U$；再由定理
\ref{etale-cohomology-theorem-exactness-stalks} 可知 $\varphi$ 是同构。因此 (1) 成立。

\medskip\noindent
省略 (1) 蕴含 (2) 的证明。
\end{proof}

\begin{lemma}
\label{etale-cohomology-lemma-specialization-map-pushforward}
设 $f : X \to S$ 为概形的拟紧拟分离态射。设 $K \in D^+(X_\etale)$。
设 $\overline{s}$ 为 $S$ 的几何点，$\overline{t}$ 为
$\Spec(\mathcal{O}^{sh}_{S, \overline{s}})$ 的几何点。则有交换图
$$
\xymatrix{
(Rf_*K)_{\overline{s}} \ar[r]_{sp} \ar@{=}[d] &
(Rf_*K)_{\overline{t}} \ar@{=}[d] \\
R\Gamma(X \times_S \Spec(\mathcal{O}^{sh}_{S, \overline{s}}), K)
\ar[r] &
R\Gamma(X \times_S \Spec(\mathcal{O}^{sh}_{S, \overline{t}}), K)
}
$$
，其中下方水平箭头来自沿态射 $\text{id}_X \times c$ 的拉回，而
$c : \Spec(\mathcal{O}^{sh}_{S, \overline{t}}) \to
\Spec(\mathcal{O}^{sh}_{S, \overline{s}})$ 是注 \ref{etale-cohomology-remark-another-sp}
中引入的态射。竖直箭头由定理 \ref{etale-cohomology-theorem-higher-direct-images} 给出。
\end{lemma}

\begin{proof}
这立即由注 \ref{etale-cohomology-remark-another-sp} 中对 $sp$ 的描述得出。
\end{proof}

\begin{remark}
\label{etale-cohomology-remark-specialization-map-and-fibres}
设 $f : X \to S$ 为概形态射，$K \in D(X_\etale)$。设 $\overline{s}$ 为 $S$
的几何点，$\overline{t}$ 为 $\Spec(\mathcal{O}^{sh}_{S, \overline{s}})$ 的几何点。
设 $c$ 如注 \ref{etale-cohomology-remark-another-sp} 中所述。总可以作出交换图
$$
\xymatrix{
(Rf_*K)_{\overline{s}} \ar[r] \ar[d]_{sp} &
R\Gamma(X \times_S \Spec(\mathcal{O}_{S, \overline{s}}^{sh}), K)
\ar[r] \ar[d]_{(\text{id}_X \times c)^{-1}} &
R\Gamma(X_{\overline{s}}, K) \\
(Rf_*K)_{\overline{t}} \ar[r] &
R\Gamma(X \times_S \Spec(\mathcal{O}_{S, \overline{t}}^{sh}), K) \ar[r] &
R\Gamma(X_{\overline{t}}, K)
}
$$
，其中水平箭头是注 \ref{etale-cohomology-remark-stalk-fibre} 中的箭头。一般而言，右侧并不存在一个
夹在此图中的竖直映射，它连接 $K$ 在诸纤维上的上同调；即使在引理
\ref{etale-cohomology-lemma-specialization-map-pushforward} 的情形中也是如此。
\end{remark}















\section{具有可构造上同调的复形}
\label{etale-cohomology-section-Dc}

\noindent
设 $\Lambda$ 为环。记 $D(X_\etale, \Lambda)$ 为 $\Lambda$-模层（定义在
$X_\etale$ 上）的导出范畴。记 $D^b(X_\etale, \Lambda)$（分别记
$D^+$、$D^-$）为 $D(X_\etale, \Lambda)$ 中有界（分别为上有界、下有界）
复形所成的充满子范畴。

\begin{definition}
\label{etale-cohomology-definition-c}
设 $X$ 为概形，$\Lambda$ 为诺特环。记
{\it $D_c(X_\etale, \Lambda)$} 为 $D(X_\etale, \Lambda)$ 的充满子范畴，
其对象是上同调层均为可构造 $\Lambda$-模层的复形。
\end{definition}

\noindent
由引理 \ref{etale-cohomology-lemma-constructible-abelian} 及导出范畴，第
\ref{derived-section-triangulated-sub} 节，此定义是合理的。因此，
$D_c(X_\etale, \Lambda)$ 是 $D(X_\etale, \Lambda)$ 的严格充满、饱和三角子范畴。

\begin{lemma}
\label{etale-cohomology-lemma-restrict-and-shriek-from-etale-c}
设 $\Lambda$ 为诺特环。若 $j : U \to X$ 为概形的 \'etale 态射，则
\begin{enumerate}
\item 结论 $K|_U \in D_c(U_\etale, \Lambda)$ 在
$K \in D_c(X_\etale, \Lambda)$ 时成立；
\item 结论 $j_!M \in D_c(X_\etale, \Lambda)$ 在
$M \in D_c(U_\etale, \Lambda)$ 且态射 $j$ 拟紧拟分离时成立。
\end{enumerate}
\end{lemma}

\begin{proof}
第一条是显然的。第二条由 $j_!$ 正合及引理
\ref{etale-cohomology-lemma-jshriek-constructible} 得出。
\end{proof}

\begin{lemma}
\label{etale-cohomology-lemma-pullback-c}
设 $\Lambda$ 为诺特环，$f : X \to Y$ 为概形态射。若
$K \in D_c(Y_\etale, \Lambda)$，则 $Lf^*K \in D_c(X_\etale, \Lambda)$。
\end{lemma}

\begin{proof}
这是因为 $f^{-1} = f^*$ 正合，再由引理
\ref{etale-cohomology-lemma-pullback-constructible} 即得。
\end{proof}

\begin{lemma}
\label{etale-cohomology-lemma-one-constructible}
设 $X$ 为拟紧拟分离概形，$\Lambda$ 为诺特环。设
$K \in D(X_\etale, \Lambda)$、$b \in \mathbf{Z}$，且 $H^b(K)$ 可构造。
则存在层 $\mathcal{F}$，它是有限多个 $j_{U!}\underline{\Lambda}$ 的直和，
其中 $U \in \Ob(X_\etale)$ 为仿射对象；并且存在映射
$\mathcal{F}[-b] \to K$，它位于 $D(X_\etale, \Lambda)$ 中，并诱导满射
$\mathcal{F} \to H^b(K)$。
\end{lemma}

\begin{proof}
将 $K$ 表示为复形 $\mathcal{K}^\bullet$，该复形由 $\Lambda$-模层组成。考虑满射
$$
\Ker(\mathcal{K}^b \to \mathcal{K}^{b + 1})
\longrightarrow
H^b(K)
$$
。由位点上的模，引理
\ref{sites-modules-lemma-module-quotient-direct-sum}
，可以选取满射
$\bigoplus_{i \in I} j_{U_i!} \underline{\Lambda} \to
\Ker(\mathcal{K}^b \to \mathcal{K}^{b + 1})$
，其中 $U_i$ 仿射。对有限子集 $I' \subset I$，记
$\mathcal{H}_{I'} \subset H^b(K)$ 为
$\bigoplus_{i \in I'} j_{U_i!} \underline{\Lambda}$ 的像。由引理
\ref{etale-cohomology-lemma-colimit-constructible}，可知 $\mathcal{H}_{I'} = H^b(K)$ 对某个有限子集
$I' \subset I$ 成立。取
$\mathcal{F} = \bigoplus_{i \in I'} j_{U_i!} \underline{\Lambda}$ 即得结论。
\end{proof}

\begin{lemma}
\label{etale-cohomology-lemma-bounded-above-c}
设 $X$ 为拟紧拟分离概形，$\Lambda$ 为诺特环，并设
$K \in D^-(X_\etale, \Lambda)$。下列条件等价：
\begin{enumerate}
\item $K$ 属于 $D_c(X_\etale, \Lambda)$；
\item $K$ 可由一个上有界复形表示，该复形的每一项都是有限多个
$j_{U!}\underline{\Lambda}$ 的直和，其中 $U \in \Ob(X_\etale)$ 仿射；
\item $K$ 可由一个上有界的平坦可构造 $\Lambda$-模层复形表示。
\end{enumerate}
\end{lemma}

\begin{proof}
(2) 蕴含 (3)、(3) 蕴含 (1) 都是显然的。假设
$K$ 属于 $D_c^-(X_\etale, \Lambda)$。设 $H^i(K) = 0$ 对所有 $i > b$ 成立。
我们对 $a$ 归纳构造复形 $\mathcal{F}^a \to \ldots \to \mathcal{F}^b$，
使得每个 $\mathcal{F}^i$ 都是有限多个 $j_{U!}\underline{\Lambda}$ 的直和，
其中 $U \in \Ob(X_\etale)$ 仿射；还构造映射
$\mathcal{F}^\bullet \to K$，它诱导同构
$H^i(\mathcal{F}^\bullet) \to H^i(K)$（对 $i > a$），并诱导满射
$H^a(\mathcal{F}^\bullet) \to H^a(K)$。当 $a = b$ 时，由引理
\ref{etale-cohomology-lemma-one-constructible} 可以完成此构造。给定这样的数据，选取优三角
$$
\mathcal{F}^\bullet \to K \to L \to \mathcal{F}^\bullet[1]
$$
。于是 $H^i(L) = 0$ 对所有 $i \geq a$ 成立。按引理
\ref{etale-cohomology-lemma-one-constructible}，选取
$\mathcal{F}^{a - 1}[-a +1] \to L$。复合
$\mathcal{F}^{a - 1}[-a +1] \to L \to \mathcal{F}^\bullet$
对应于映射 $\mathcal{F}^{a - 1} \to \mathcal{F}^a$，且它与
$\mathcal{F}^a \to \mathcal{F}^{a + 1}$ 的复合为零。由 TR4 得到映射
$$
(\mathcal{F}^{a - 1} \to \ldots \to \mathcal{F}^b) \to K
$$
，它位于 $D(X_\etale, \Lambda)$ 中。这就完成了归纳步骤及引理的证明。
\end{proof}

\begin{lemma}
\label{etale-cohomology-lemma-tensor-c}
设 $X$ 为概形，$\Lambda$ 为诺特环，并设
$K, L \in D_c^-(X_\etale, \Lambda)$。则
$K \otimes_\Lambda^\mathbf{L} L$ 属于 $D_c^-(X_\etale, \Lambda)$。
\end{lemma}

\begin{proof}
这由引理 \ref{etale-cohomology-lemma-bounded-above-c} 和
\ref{etale-cohomology-lemma-tensor-product-constructible} 得出。
\end{proof}




\section{具有可构造上同调的有限 Tor 维数复形}
\label{etale-cohomology-section-ctf}

\noindent
设 $X$ 为概形，$\Lambda$ 为诺特环。导出范畴 $D_c(X_\etale, \Lambda)$
（定义见第 \ref{etale-cohomology-section-Dc} 节）的一个常用子范畴，是由（局部）具有有限
Tor 维数的对象组成的充满子范畴。下面给出正式定义。

\begin{definition}
\label{etale-cohomology-definition-ctf}
设 $X$ 为概形，$\Lambda$ 为诺特环。记
{\it $D_{ctf}(X_\etale, \Lambda)$} 为 $D_c(X_\etale, \Lambda)$ 的充满子范畴，
其对象局部具有有限 Tor 维数。
\end{definition}

\noindent
这是 $D_c(X_\etale, \Lambda)$ 和 $D(X_\etale, \Lambda)$ 的严格充满、饱和
三角子范畴。按照我们的约定，见位点上的上同调，定义
\ref{sites-cohomology-definition-tor-amplitude}，有
$$
D_{ctf}(X_\etale, \Lambda) \subset D^b_c(X_\etale, \Lambda)
\subset D^b(X_\etale, \Lambda)
$$
，只要 $X$ 拟紧。关于 $D_{ctf}(X_\etale, \Lambda)$ 中的对象，引理
\ref{etale-cohomology-lemma-when-ctf} 给出了一种很好的理解方式。

\begin{remark}
\label{etale-cohomology-remark-different}
从某种意义上说，导出范畴 $D_{ctf}(X_\etale, \Lambda)$ 中的对象，比
$D(\mathcal{O}_X)$ 中的完美对象具有更好的整体性质。具体而言，一个
$\mathcal{O}_X$-模复形可能局部拟同构于有限局部自由
$\mathcal{O}_X$-模的有限复形，却不整体拟同构于局部自由
$\mathcal{O}_X$-模的有界复形。下面的引理说明，在诺特概形上的
$D_{ctf}$ 中不会出现这种现象。
\end{remark}

\begin{lemma}
\label{etale-cohomology-lemma-when-ctf}
设 $\Lambda$ 为诺特环，$X$ 为拟紧拟分离概形，并设
$K \in D(X_\etale, \Lambda)$。下列条件等价：
\begin{enumerate}
\item $K \in D_{ctf}(X_\etale, \Lambda)$；
\item $K$ 可由可构造平坦 $\Lambda$-模层的有限复形表示。
\end{enumerate}
事实上，若 $K$ 的 Tor 振幅落在 $[a, b]$ 中，则可将 $K$ 表示为复形
$\mathcal{F}^a \to \ldots \to \mathcal{F}^b$，其中
$\mathcal{F}^p$ 是可构造平坦 $\Lambda$-模层。
\end{lemma}

\begin{proof}
可构造平坦 $\Lambda$-模层的有限复形具有有限 Tor 维数，这是显然的。
它属于 $D_c(X_\etale, \Lambda)$ 也是显然的。因此 (2) 蕴含 (1)。

\medskip\noindent
假设 (1)。选取 $a, b \in \mathbf{Z}$，使得
$H^i(K \otimes_\Lambda^\mathbf{L} \mathcal{G}) = 0$ if
$i \not \in [a, b]$，并且这对所有 $\Lambda$-模层 $\mathcal{G}$ 都成立。
我们对 $b - a$ 归纳证明最后一个断言。若
$a = b$，则 $K = H^a(K)[-a]$ 是平坦可构造层，结论成立。接着假设
$b > a$。将 $K$ 表示为复形 $\mathcal{K}^\bullet$，其项为 $\Lambda$-模层。
考虑满射
$$
\Ker(\mathcal{K}^b \to \mathcal{K}^{b + 1})
\longrightarrow
H^b(K)
$$
。由引理 \ref{etale-cohomology-lemma-category-constructible-modules}，可找到有限多个仿射概形
$U_i$，它们在 $X$ 上 \'etale；还可找到满射
$\bigoplus j_{U_i!}\underline{\Lambda}_{U_i} \to H^b(K)$。把 $U_i$ 换成标准
\'etale 覆盖 $\{U_{ij} \to U_i\}$ 后，可以假定该满射提升为映射
$\mathcal{F} = \bigoplus j_{U_i!}\underline{\Lambda}_{U_i} \to
\Ker(\mathcal{K}^b \to \mathcal{K}^{b + 1})$.
该映射决定优三角
$$
\mathcal{F}[-b] \to K \to L \to \mathcal{F}[-b + 1]
$$
，它位于 $D(X_\etale, \Lambda)$ 中。由于 $D_{ctf}(X_\etale, \Lambda)$ 是三角
子范畴，$L$ 也属于其中。事实上，$L$ 的 Tor 振幅落在 $[a, b - 1]$ 中，因为
$\mathcal{F}$ 满射到 $H^b(K)$（细节略）。由归纳假设，可以找到有限复形
$\mathcal{F}^a \to \ldots \to \mathcal{F}^{b - 1}$，它由可构造平坦
$\Lambda$-模层组成，并表示 $L$。映射
$L \to \mathcal{F}[-b + 1]$ 对应于映射
$\mathcal{F}^b \to \mathcal{F}$，后者在
$\mathcal{F}^{b - 1} \to \mathcal{F}^b$ 的像上为零。于是由公理 TR3，$K$ 可由复形
$$
\mathcal{F}^a \to \ldots \to \mathcal{F}^{b - 1} \to \mathcal{F}^b
$$
表示，从而完成证明。
\end{proof}

\begin{remark}
\label{etale-cohomology-remark-projective-each-degree}
设 $\Lambda$ 为诺特环，$X$ 为概形。对有界复形 $K^\bullet$，它由可构造平坦
$\Lambda$-模层组成并定义在 $X_\etale$ 上，每个茎
$K^p_{\overline{x}}$ 都是有限投射 $\Lambda$-模。因此，该复形的各个茎都是
$\Lambda$-模的完美复形。
\end{remark}

\begin{lemma}
\label{etale-cohomology-lemma-restrict-and-shriek-from-etale-ctf}
设 $\Lambda$ 为诺特环。若 $j : U \to X$ 为概形的 \'etale 态射，则
\begin{enumerate}
\item 结论 $K|_U \in D_{ctf}(U_\etale, \Lambda)$ 在
$K \in D_{ctf}(X_\etale, \Lambda)$ 时成立；
\item 结论 $j_!M \in D_{ctf}(X_\etale, \Lambda)$ 在
$M \in D_{ctf}(U_\etale, \Lambda)$ 且态射 $j$ 拟紧拟分离时成立。
\end{enumerate}
\end{lemma}

\begin{proof}
证明本引理或许最容易的方法，是先约化到 $X$ 仿射的情形，再应用引理
\ref{etale-cohomology-lemma-when-ctf}，把它转化为关于可构造平坦 $\Lambda$-模层有限复形的断言；
此时结论由引理 \ref{etale-cohomology-lemma-jshriek-constructible} 得出。
\end{proof}

\begin{lemma}
\label{etale-cohomology-lemma-pullback-ctf}
设 $\Lambda$ 为诺特环，$f : X \to Y$ 为概形态射。若
$K \in D_{ctf}(Y_\etale, \Lambda)$，则
$Lf^*K \in D_{ctf}(X_\etale, \Lambda)$。
\end{lemma}

\begin{proof}
应用引理 \ref{etale-cohomology-lemma-when-ctf}，把问题约化为关于可构造平坦
$\Lambda$-模层有限复形的问题。于是，由 $f^{-1} = f^*$ 正合及引理
\ref{etale-cohomology-lemma-pullback-constructible} 即得结论。
\end{proof}

\begin{lemma}
\label{etale-cohomology-lemma-connected-ctf-locally-constant}
设 $X$ 为连通概形，$\Lambda$ 为诺特环。设
$K \in D_{ctf}(X_\etale, \Lambda)$ 具有局部常值上同调层。则存在有限投射
$\Lambda$-模的有限复形 $M^\bullet$ 以及 \'etale 覆盖 $\{U_i \to X\}$，使得
$K|_{U_i} \cong \underline{M^\bullet}|_{U_i}$ 在 $D(U_{i, \etale}, \Lambda)$ 中成立。
\end{lemma}

\begin{proof}
选取 \'etale 覆盖 $\{U_i \to X\}$，使得 $K|_{U_i}$ 为常值复形，比如
$K|_{U_i} \cong \underline{M_i^\bullet}_{U_i}$；这里所用有限复形由有限
$\Lambda$-模组成，记为 $M_i^\bullet$。见位点上的上同调，引理
\ref{sites-cohomology-lemma-locally-constant}.
注意，$U_i \times_X U_j$ 在 $M_i^\bullet$ 与 $M_j^\bullet$ 于
$D(\Lambda)$ 中不同构时为空。对每个 $\Lambda$-模复形 $M^\bullet$，令
$I_{M^\bullet} =
\{i \in I \mid M_i^\bullet \cong M^\bullet\text{ 于 }D(\Lambda)\}$.
由于 \'etale 态射是开映射，故
$U_{M^\bullet} = \bigcup_{i \in I_{M^\bullet}} \Im(U_i \to X)$
是 $X$ 的开子集。于是 $X = \coprod U_{M^\bullet}$ 是 $X$ 的不交开覆盖。
由于 $X$ 连通，只有一个 $U_{M^\bullet}$ 非空。又因为 $K$ 属于
$D_{ctf}(X_\etale, \Lambda)$，可知 $M^\bullet$ 是 $\Lambda$-模的完美复形；
见代数的进一步结果，引理 \ref{more-algebra-lemma-perfect}。因此，可以假定
$M^\bullet$ 是有限投射 $\Lambda$-模的有限复形。
\end{proof}








\section{挠层}
\label{etale-cohomology-section-torsion}

\noindent
本节简要讨论挠阿贝尔层及其 \'etale 上同调。设 $\mathcal{C}$ 为位点。我们已经在
位点上的上同调，引理 \ref{sites-cohomology-lemma-torsion} 中证明：
$D(\mathcal{C})$ 中上同调层均为挠层的任何对象，都可由每一项皆为挠层的复形表示。

\begin{lemma}
\label{etale-cohomology-lemma-torsion-cohomology}
设 $X$ 为拟紧拟分离概形。
\begin{enumerate}
\item 若 $\mathcal{F}$ 是 $X_\etale$ 上的挠阿贝尔层，则
$H^n_\etale(X, \mathcal{F})$ 对所有 $n$ 都是挠阿贝尔群。
\item 若 $K$ 属于 $D^+(X_\etale)$ 且具有挠上同调层，则
$H^n_\etale(X, K)$ 对所有 $n$ 都是挠阿贝尔群。
\end{enumerate}
\end{lemma}

\begin{proof}
为证明 (1)，写成 $\mathcal{F} = \bigcup \mathcal{F}[n]$，其中
$\mathcal{F}[d]$ 是 $d$-挠子层。由引理 \ref{etale-cohomology-lemma-colimit}，有
$H^n_\etale(X, \mathcal{F}) = \colim H^n_\etale(X, \mathcal{F}[d])$.
由于 $H^n_\etale(X, \mathcal{F}[d])$ 被 $d$ 零化，这就证明了 (1)。

\medskip\noindent
为证明 (2)，可以使用谱序列
$E_2^{p, q} = H^p_\etale(X, H^q(K))$，它收敛到 $H^n_\etale(X, K)$
（导出范畴，引理 \ref{derived-lemma-two-ss-complex-functor}），以及对层所得的结论。
\end{proof}

\begin{lemma}
\label{etale-cohomology-lemma-torsion-direct-image}
设 $f : X \to Y$ 为概形的拟紧拟分离态射。
\begin{enumerate}
\item 若 $\mathcal{F}$ 是 $X_\etale$ 上的挠阿贝尔层，则
$R^nf_*\mathcal{F}$ 是 $Y_\etale$ 上的挠阿贝尔层，对所有 $n$ 都成立。
\item 若 $K$ 属于 $D^+(X_\etale)$ 且具有挠上同调层，则
$Rf_*K$ 是 $D^+(Y_\etale)$ 的对象，并且其上同调层均为挠阿贝尔层。
\end{enumerate}
\end{lemma}

\begin{proof}
(1) 的证明。回忆 $R^nf_*\mathcal{F}$ 是与预层
$V \mapsto H^n_\etale(X \times_Y V, \mathcal{F})$ 相伴的层，该预层定义在
$Y_\etale$ 上。见位点上的上同调，引理
\ref{sites-cohomology-lemma-higher-direct-images}。若选取仿射的 $V$，则
$X \times_Y V$ 因 $f$ 拟紧拟分离而拟紧拟分离，故可应用引理
\ref{etale-cohomology-lemma-torsion-cohomology}，看出
$H^n_\etale(X \times_Y V, \mathcal{F})$ 是挠群。

\medskip\noindent
(2) 的证明。回忆 $R^nf_*K$ 是与预层
$V \mapsto H^n_\etale(X \times_Y V, K)$ 相伴的层，该预层定义在
$Y_\etale$ 上。见位点上的上同调，引理
\ref{sites-cohomology-lemma-unbounded-describe-higher-direct-images}。若选取仿射的
$V$，则 $X \times_Y V$ 因 $f$ 拟紧拟分离而拟紧拟分离，故可应用引理
\ref{etale-cohomology-lemma-torsion-cohomology}，看出
$H^n_\etale(X \times_Y V, K)$ 是挠群。
\end{proof}









\section{支撑在闭子概形上的上同调}
\label{etale-cohomology-section-cohomology-support}

\noindent
设 $X$ 是一个概形，且 $Z \subset X$ 是一个闭子概形。
设 $\mathcal{F}$ 是 $X_\etale$ 上的一个阿贝尔层。令
$$
\Gamma_Z(X, \mathcal{F}) =
\{s \in \mathcal{F}(X) \mid \text{Supp}(s) \subset Z\}
$$
为支撑在 $Z$ 上的截面（定义 \ref{etale-cohomology-definition-support}）。
这是一个左正合函子，但一般并不正合。
因此得到导出函子
$$
R\Gamma_Z(X, -) : D(X_\etale) \longrightarrow D(\textit{Ab})
$$
并把支撑在 $Z$ 上的上同调群定义为
$H^q_Z(X, \mathcal{F}) = R^q\Gamma_Z(X, \mathcal{F})$。

\medskip\noindent
设 $\mathcal{I}$ 是 $X_\etale$ 上的一个内射阿贝尔层。
令 $U = X \setminus Z$。
则限制映射 $\mathcal{I}(X) \to \mathcal{I}(U)$ 是满射
（位点上的上同调，引理
\ref{sites-cohomology-lemma-restriction-along-monomorphism-surjective})
，其核为 $\Gamma_Z(X, \mathcal{I})$。立即可知，对
$K \in D(X_\etale)$ 有优三角
$$
R\Gamma_Z(X, K) \to R\Gamma(X, K) \to R\Gamma(U, K) \to R\Gamma_Z(X, K)[1]
$$
，它位于 $D(\textit{Ab})$ 中。由此得到上同调长正合列
$$
\ldots \to H^i_Z(X, K) \to H^i(X, K) \to H^i(U, K) \to
H^{i + 1}_Z(X, K) \to \ldots
$$
，对任意 $K$（作为 $D(X_\etale)$ 中的对象）都成立。

\medskip\noindent
对阿贝尔层 $\mathcal{F}$（定义在 $X_\etale$ 上），可以考虑{\it 支撑在
$Z$ 上的截面子层}，记作 $\mathcal{H}_Z(\mathcal{F})$，其定义规则为
$$
\mathcal{H}_Z(\mathcal{F})(U) =
\{s \in \mathcal{F}(U) \mid \text{Supp}(s) \subset U \times_X Z\}
$$
这里使用定义 \ref{etale-cohomology-definition-support} 中截面的支撑。利用命题
\ref{etale-cohomology-proposition-closed-immersion-pushforward} 的等价，可以把
$\mathcal{H}_Z(\mathcal{F})$ 看作 $Z_\etale$ 上的一个阿贝尔层。
于是得到函子
$$
\textit{Ab}(X_\etale) \longrightarrow \textit{Ab}(Z_\etale),\quad
\mathcal{F} \longmapsto \mathcal{H}_Z(\mathcal{F})
$$
，它是左正合的，但一般并不正合。

\begin{lemma}
\label{etale-cohomology-lemma-sections-with-support-acyclic}
设 $i : Z \to X$ 是概形的一个闭浸入。
设 $\mathcal{I}$ 是 $X_\etale$ 上的一个内射阿贝尔层。
则 $\mathcal{H}_Z(\mathcal{I})$ 是 $Z_\etale$ 上的一个内射阿贝尔层。
\end{lemma}

\begin{proof}
注意，对任意阿贝尔层 $\mathcal{G}$（定义在 $Z_\etale$ 上），有
$$
\Hom_Z(\mathcal{G}, \mathcal{H}_Z(\mathcal{F})) =
\Hom_X(i_*\mathcal{G}, \mathcal{F})
$$
，因为 $i_*\mathcal{G}$ 的任意截面都支撑在 $Z$ 上。由于
$i_*$ 是正合的（第 \ref{etale-cohomology-section-closed-immersions} 节），而
$\mathcal{I}$ 在 $X_\etale$ 上是内射的，所以
$\mathcal{H}_Z(\mathcal{I})$ 在 $Z_\etale$ 上是内射的。
\end{proof}

\noindent
将导出函子记作
$$
R\mathcal{H}_Z : D(X_\etale) \longrightarrow D(Z_\etale)
$$
。置
$\mathcal{H}^q_Z(\mathcal{F}) = R^q\mathcal{H}_Z(\mathcal{F})$，于是
$\mathcal{H}^0_Z(\mathcal{F}) = \mathcal{H}_Z(\mathcal{F})$。
由上面的引理，有 Grothendieck 谱序列
$$
E_2^{p, q} = H^p(Z, \mathcal{H}^q_Z(\mathcal{F}))
\Rightarrow H^{p + q}_Z(X, \mathcal{F})
$$

\begin{lemma}
\label{etale-cohomology-lemma-cohomology-with-support-sheaf-on-support}
设 $i : Z \to X$ 是概形的一个闭浸入。
设 $\mathcal{G}$ 是 $Z_\etale$ 上的一个内射阿贝尔层。
则 $\mathcal{H}^p_Z(i_*\mathcal{G}) = 0$ 对 $p > 0$ 成立。
\end{lemma}

\begin{proof}
这是因为函子 $i_*$ 正合，并且把内射阿贝尔层变为内射阿贝尔层
（位点上的上同调，引理
\ref{sites-cohomology-lemma-pushforward-injective-flat}）。
\end{proof}

\begin{lemma}
\label{etale-cohomology-lemma-cohomology-with-support-triangle}
设 $i : Z \to X$ 是概形的一个闭浸入。
设 $j : U \to X$ 是 $Z$ 的补集的包含映射。
设 $\mathcal{F}$ 是 $X_\etale$ 上的一个阿贝尔层。
有优三角
$$
i_*R\mathcal{H}_Z(\mathcal{F}) \to \mathcal{F} \to Rj_*(\mathcal{F}|_U) \to
i_*R\mathcal{H}_Z(\mathcal{F})[1]
$$
，它位于 $D(X_\etale)$ 中。由此得到正合列
$$
0 \to i_*\mathcal{H}_Z(\mathcal{F}) \to \mathcal{F} \to
j_*(\mathcal{F}|_U) \to i_*\mathcal{H}^1_Z(\mathcal{F}) \to 0
$$
以及同构
$R^pj_*(\mathcal{F}|_U) \cong i_*\mathcal{H}^{p + 1}_Z(\mathcal{F})$
，其中 $p \geq 1$。
\end{lemma}

\begin{proof}
为得到这个优三角，选取内射分解
$\mathcal{F} \to \mathcal{I}^\bullet$。于是由上面的讨论得到复形的短正合列
$$
0 \to
i_*\mathcal{H}_Z(\mathcal{I}^\bullet) \to \mathcal{I}^\bullet
\to j_*(\mathcal{I}^\bullet|_U) \to 0
$$
。因此由导出范畴，第
\ref{derived-section-canonical-delta-functor} 节，得到上述优三角。
\end{proof}

\noindent
设 $X$ 是一个概形，且 $Z \subset X$ 是一个闭子概形。用
$D_Z(X_\etale)$ 表示 $D(X_\etale)$ 的严格充满、饱和三角子范畴，
其对象是上同调层支撑在 $Z$ 上的复形。注意，
$D_Z(X_\etale)$ 仅依赖于 $X$ 的底闭子集。

\begin{lemma}
\label{etale-cohomology-lemma-complexes-with-support-on-closed}
设 $i : Z \to X$ 是概形的一个闭浸入。函子
$Ri_{small, *} = i_{small, *} : D(Z_\etale) \to D(X_\etale)$
诱导等价 $D(Z_\etale) \to D_Z(X_\etale)$，其拟逆为
$$
i_{small}^{-1}|_{D_Z(X_\etale)} = R\mathcal{H}_Z|_{D_Z(X_\etale)}
$$
\end{lemma}

\begin{proof}
回忆 $i_{small}^{-1}$ 与 $i_{small, *}$ 构成一对正合函子的伴随，
并且 $i_{small}^{-1}i_{small, *}$ 同构于阿贝尔层范畴上的恒等函子。
见命题 \ref{etale-cohomology-proposition-closed-immersion-pushforward} 和引理
\ref{etale-cohomology-lemma-stalk-pullback}。因此
$i_{small, *} : D(Z_\etale) \to D_Z(X_\etale)$ 充分忠实，而
$i_{small}^{-1}$ 给出一个左逆。另一方面，设 $K$ 是
$D_Z(X_\etale)$ 的一个对象，并考虑伴随映射
$K \to i_{small, *}i_{small}^{-1}K$。利用 $i_{small, *}$ 和
$i_{small}^{-1}$ 的正合性，它在上同调层上诱导伴随映射
$H^n(K) \to i_{small, *}i_{small}^{-1}H^n(K)$。由于这些上同调层
支撑在 $Z$ 上，可知这些伴随映射都是同构，因而
$D(Z_\etale) \to D_Z(X_\etale)$ 是一个等价。

\medskip\noindent
为完成证明，还须说明 $R\mathcal{H}_Z(K) = i_{small}^{-1}K$，
其中 $K$ 是 $D_Z(X_\etale)$ 的对象。为此可以使用刚刚证明的等式
$K = i_{small, *}i_{small}^{-1}K$。于是可以选取一个 K-内射代表
$\mathcal{I}^\bullet$ 来表示 $i_{small}^{-1}K$。
由于 $i_{small, *}$ 是正合函子 $i_{small}^{-1}$ 的右伴随，复形
$i_{small, *}\mathcal{I}^\bullet$ 是 K-内射的
（导出范畴，引理 \ref{derived-lemma-adjoint-preserve-K-injectives}）。
可见 $R\mathcal{H}_Z(K)$ 由
$\mathcal{H}_Z(i_{small, *}\mathcal{I}^\bullet) = \mathcal{I}^\bullet$
计算，这正是所需结论。
\end{proof}

\begin{lemma}
\label{etale-cohomology-lemma-cohomology-with-support-quasi-coherent}
设 $X$ 是一个概形，$Z \subset X$ 是一个闭子概形。
设 $\mathcal{F}$ 是一个拟凝聚 $\mathcal{O}_X$-模，并用
$\mathcal{F}^a$ 表示 $X$ 的小 \'etale 网站上的伴随拟凝聚层
（命题 \ref{etale-cohomology-proposition-quasi-coherent-sheaf-fpqc}）。则
\begin{enumerate}
\item $H^q_Z(X, \mathcal{F})$ 与 $H^q_Z(X_\etale, \mathcal{F}^a)$ 一致；
\item 若 $Z$ 的补集在 $X$ 中是逆紧致的，则
$i_*\mathcal{H}^q_Z(\mathcal{F}^a)$ 是一个拟凝聚
$\mathcal{O}_X$-模层，并且等于 $(i_*\mathcal{H}^q_Z(\mathcal{F}))^a$。
\end{enumerate}
\end{lemma}

\begin{proof}
设 $j : U \to X$ 是 $Z$ 的补集的包含映射。关于上同调群的陈述 (1)
由带支撑上同调的长正合列以及等式
$H^q(X_\etale, \mathcal{F}^a) = H^q(X, \mathcal{F})$ 和
$H^q(U_\etale, \mathcal{F}^a) = H^q(U, \mathcal{F})$ 得出，见定理
\ref{etale-cohomology-theorem-zariski-fpqc-quasi-coherent}。若
$j : U \to X$ 是拟紧态射，即若 $U \subset X$ 是逆紧致的，则
$R^qj_*$ 把拟凝聚层变为拟凝聚层
（概形的上同调，引理
\ref{coherent-lemma-quasi-coherence-higher-direct-images})
，并且与在 \'etale 网站上取伴随层可交换
（下降，引理 \ref{descent-lemma-higher-direct-images-small-etale}）。
应用引理 \ref{etale-cohomology-lemma-cohomology-with-support-triangle} 即得结论。
\end{proof}







\section{具有严格 Hensel 局部环的概形}
\label{etale-cohomology-section-strictly-henselian-local-rings}

\noindent
本节汇集一些关于局部环严格 Hensel 的概形之 \'etale 上同调的结果。
例如，下面是引理
\ref{etale-cohomology-lemma-vanishing-etale-cohomology-strictly-henselian}
的一个颇有趣的推广。

\begin{lemma}
\label{etale-cohomology-lemma-local-rings-strictly-henselian}
设 $S$ 是一个所有局部环都严格 Hensel 的概形。
则对任意阿贝尔层 $\mathcal{F}$（定义在 $S_\etale$ 上），有
$H^i(S_\etale, \mathcal{F}) = H^i(S_{Zar}, \mathcal{F})$。
\end{lemma}

\begin{proof}
设 $\epsilon : S_\etale \to S_{Zar}$ 是由包含函子给出的网站态射。
Zariski 层 $R^p\epsilon_*\mathcal{F}$ 是预层
$U \mapsto H^p_\etale(U, \mathcal{F})$ 的伴随层。因此，它在
$x \in X$ 处的茎为
$\colim H^p_\etale(U, \mathcal{F}) =
H^p_\etale(\Spec(\mathcal{O}_{X, x}), \mathcal{G}_x)$
，其中 $\mathcal{G}_x$ 表示 $\mathcal{F}$ 到
$\Spec(\mathcal{O}_{X, x})$ 的拉回，见引理
\ref{etale-cohomology-lemma-directed-colimit-cohomology}。因此，由引理
\ref{etale-cohomology-lemma-vanishing-etale-cohomology-strictly-henselian}，
高阶正像 $R^p\epsilon_*\mathcal{F}$ 为零；再由 Leray 谱序列即得结论。
\end{proof}

\begin{lemma}
\label{etale-cohomology-lemma-gabber-criterion}
设 $R$ 是一个所有局部环都严格 Hensel 的环。
设 $\mathcal{F}$ 是 $\Spec(R)_\etale$ 上的一个层。
假设对所有 $f, g \in R$，映射
$$
H^1_\etale(D(f + g), \mathcal{F})
\longrightarrow
H^1_\etale(D(f(f + g)), \mathcal{F}) \oplus
H^1_\etale(D(g(f + g)), \mathcal{F})
$$
的核为零。则 $H^q_\etale(\Spec(R), \mathcal{F}) = 0$ 对 $q > 0$ 成立。
\end{lemma}

\begin{proof}
由引理 \ref{etale-cohomology-lemma-local-rings-strictly-henselian} 可知，
$\mathcal{F}$ 的 \'etale 上同调在 $\Spec(R)$ 的任意开子集上
都与 Zariski 上同调一致。我们对 $i$ 归纳证明如下陈述：若
$h \in R$，则 $H^q_\etale(D(h), \mathcal{F}) = 0$ 对
$1 \leq q \leq i$ 成立。基础情形 $i = 0$ 是平凡的。假设 $i \geq 1$。

\medskip\noindent
取 $\xi \in H^q_\etale(D(h), \mathcal{F})$，其中
$1 \leq q \leq i$ 且 $h \in R$。若 $q < i$，则由归纳已经完成，
故假设 $q = i$。把 $R$ 替换为 $R_h$ 后，可以假设
$\xi \in H^i_\etale(\Spec(R), \mathcal{F})$；这里略去一些细节。
令 $I \subset R$ 为元素 $f \in R$ 所成的集合，这些元素满足
$\xi|_{D(f)} = 0$。由于 $\xi$ 在 Zariski 局部平凡，
对每个素理想 $\mathfrak p$（属于 $R$），存在 $f \in I$ 使得
$f \not \in \mathfrak p$。因此，只要证明 $I$ 是理想，就有
$1 \in I$，从而完成证明。显然，$f \in I$、$r \in R$ 蕴含
$rf \in I$。于是设 $f, g \in I$，并证明 $f + g \in I$。
若 $q = i = 1$，这恰好就是引理的假设！因而得到 $i = 1$ 的情形。
当 $q = i > 1$ 时，注意
$$
D(f + g) = D(f(f + g)) \cup D(g(f + g))
$$
由 Mayer--Vietoris（上同调，引理
\ref{cohomology-lemma-mayer-vietoris}；它适用是因为
$\Spec(R)$ 的开子概形上的 \'etale 上同调等于 Zariski 上同调），
有正合列
$$
\xymatrix{
H^{i - 1}_\etale(D(fg(f + g)), \mathcal{F}) \ar[d] \\
H^i_\etale(D(f + g), \mathcal{F}) \ar[d] \\
H^i_\etale(D(f(f + g)), \mathcal{F}) \oplus
H^i_\etale(D(g(f + g)), \mathcal{F})
}
$$
而第一个群由归纳为零，故结论成立。
\end{proof}

\begin{lemma}
\label{etale-cohomology-lemma-affine-only-closed-points}
设 $S$ 是一个仿射概形，满足：(1) 所有点都是闭点；(2) 所有剩余域
都是可分代数闭的。则对任意阿贝尔层 $\mathcal{F}$（定义在
$S_\etale$ 上），有 $H^i(S_\etale, \mathcal{F}) = 0$ 对 $i > 0$ 成立。
\end{lemma}

\begin{proof}
条件 (1) 蕴含 $S$ 的底拓扑空间是预有限的，见代数，引理
\ref{algebra-lemma-ring-with-only-minimal-primes}。因此，拓扑空间
$S$ 上阿贝尔层的高阶上同调群（即 Zariski 上同调）为零，见上同调，
引理 \ref{cohomology-lemma-vanishing-for-profinite}。由代数，引理
\ref{algebra-lemma-local-dimension-zero-henselian}，其局部环严格 Hensel。
于是由引理 \ref{etale-cohomology-lemma-local-rings-strictly-henselian}，$S$ 的
\'etale 上同调由 Zariski 上同调计算，证明完成。
\end{proof}

\noindent
绝对整闭环的谱是所有局部环都严格 Hensel 的概形的一个例子，见代数进阶，
引理
\ref{more-algebra-lemma-absolutely-integrally-closed-strictly-henselian}。
事实证明，具有可分闭分式域的正规整环具有更强的性质，如下面的引理所述。

\begin{lemma}
\label{etale-cohomology-lemma-normal-scheme-with-alg-closed-function-field}
设 $X$ 是一个函数域可分闭的整正规概形。
\begin{enumerate}
\item 分离 \'etale 态射 $U \to X$ 是开浸入的不交并。
\item $X$ 的所有局部环都严格 Hensel。
\end{enumerate}
\end{lemma}

\begin{proof}
设 $R$ 是一个分式域可分代数闭的正规整环。设
$R \to A$ 是 \'etale 环映射。则 $A \otimes_R K$ 作为
$K$-代数是有限乘积 $\prod_{i = 1, \ldots, n} K$，其各因子都是
$K$。设 $e_i$（$i = 1, \ldots, n$）
是 $A \otimes_R K$ 中相应的幂等元。由于 $A$ 正规
（代数，引理 \ref{algebra-lemma-normal-goes-up}），这些幂等元
$e_i$ 属于 $A$（代数，引理
\ref{algebra-lemma-normal-ring-integrally-closed}）。因此
$A = \prod Ae_i$，可以假设 $A \otimes_R K = K$。由于
$A \subset A \otimes_R K = K$（利用 $R \to A$ 的平坦性以及
$R \subset K$），可知 $A$ 是整环。由同一论证可知
$A \otimes_R A \subset (A \otimes_R A) \otimes_R K = K$。
于是映射 $A \otimes_R A \to A$ 既是单射也是满射。因此，由态射，
引理 \ref{morphisms-lemma-universally-injective} 和 \'Etale 态射，
定理 \ref{etale-theorem-etale-radicial-open}，$R \to A$ 定义一个开浸入。

\medskip\noindent
设 $f : U \to X$ 是一个分离 \'etale 态射。设 $\eta \in X$
为一般点，并令 $f^{-1}(\{\eta\}) = \{\xi_i\}_{i \in I}$。
上一段的结果说明：对任意仿射开集 $U' \subset U$，只要其在
$X$ 中的像包含于一个仿射开集，就有
$U' = \coprod_{i \in I} U'_i$，其中 $U'_i$ 是 $U'$ 中所有
$\xi_i$ 的特殊化点所成的集合。此外，态射 $U'_i \to X$
是开浸入。由此可知，$U_i = \overline{\{\xi_i\}}$ 是 $U$ 的
开闭子概形，并且 $U_i \to X$ 在源上局部是同构。由态射，引理
\ref{morphisms-lemma-distinct-local-rings}，$U_i \to X$ 的分离性蕴含
$U_i \to X$ 是单射，故 $U_i \to X$ 是开浸入，即 (1) 成立。

\medskip\noindent
由 (1) 以及 $\mathcal{O}_{X, x}$ 的严格 Hensel 化可描述为
$\overline{x}$ 处局部环（位于 $X$ 的 \'etale 网站上）这一事实
（引理 \ref{etale-cohomology-lemma-describe-etale-local-ring}），可得 (2)。
也可以直接证明，见基本群，引理
\ref{pione-lemma-normal-local-domain-separablly-closed-fraction-field}。
\end{proof}

\begin{lemma}
\label{etale-cohomology-lemma-Rf-star-zero-normal-with-alg-closed-function-field}
设 $f : X \to Y$ 是概形态射，其中 $X$ 是函数域可分闭的整正规概形。
则 $R^qf_*\underline{M} = 0$ 对 $q > 0$ 以及任意阿贝尔群 $M$ 成立。
\end{lemma}

\begin{proof}
回忆，$R^qf_*\underline{M}$ 是预层
$V \mapsto H^q_\etale(V \times_Y X, M)$ 在 $Y_\etale$ 上的伴随层，见引理
\ref{etale-cohomology-lemma-higher-direct-images}。若 $V$ 仿射，则
$V \times_Y X \to X$ 是分离 \'etale 态射。因此
$V \times_Y X = \coprod U_i$ 是开子概形 $U_i$（属于 $X$）的不交并，
见引理 \ref{etale-cohomology-lemma-normal-scheme-with-alg-closed-function-field}。
由引理 \ref{etale-cohomology-lemma-local-rings-strictly-henselian}，
$H^q_\etale(U_i, M)$ 等于 $H^q_{Zar}(U_i, M)$。由上同调，引理
\ref{cohomology-lemma-irreducible-constant-cohomology-zero}，后者为零。
\end{proof}

\begin{lemma}
\label{etale-cohomology-lemma-closed-of-affine-normal-scheme-with-alg-closed-function-field}
设 $X$ 是函数域可分闭的仿射整正规概形。设
$Z \subset X$ 是一个闭子概形。设 $V \to Z$ 是 \'etale 态射，
且 $V$ 仿射。则 $V$ 是 $Z$ 的有限多个开子概形的不交并。
若 $V \to Z$ 是满射且有限 \'etale，则 $V \to Z$ 有一个截面。
\end{lemma}

\begin{proof}
由代数，引理 \ref{algebra-lemma-lift-etale}，可以把 $V$
提升为仿射概形 $U$，它在 $X$ 上是 \'etale 的。对 $U \to X$ 应用引理
\ref{etale-cohomology-lemma-normal-scheme-with-alg-closed-function-field}，得到第一个陈述。

\medskip\noindent
最后一个陈述是第一个陈述的推论。设
$V = \coprod_{i = 1, \ldots, n} V_i$ 是分解为开闭子概形的一个有限分解，
并且 $V_i \to Z$ 是开浸入。由于 $V \to Z$ 有限，可知
$V_i \to Z$ 也是闭的。设 $U_i \subset Z$ 是其像。于是有开闭子概形分解
$$
Z =
\coprod\nolimits_{(A, B)}
\bigcap\nolimits_{i \in A} U_i \cap
\bigcap\nolimits_{i \in B} U_i^c
$$
，其中不交并遍历所有分解 $\{1, \ldots, n\} = A \amalg B$，
且 $A$ 至少有一个元素。每个层都包含在某个 $U_i$ 中，
于是得到所需截面。
\end{proof}

\begin{lemma}
\label{etale-cohomology-lemma-gabber-for-h1-absolutely-algebraically-closed}
设 $X$ 是函数域可分闭的正规整仿射概形。设
$Z \subset X$ 是一个闭子概形。对任意有限阿贝尔群
$M$，有 $H^1_\etale(Z, \underline{M}) = 0$。
\end{lemma}

\begin{proof}
由位点上的上同调，引理 \ref{sites-cohomology-lemma-torsors-h1}，
$H^1_\etale(Z, \underline{M})$ 的一个元素对应于一个
$\underline{M}$-旋子 $\mathcal{F}$，它定义在 $Z_\etale$ 上。
这样的旋子显然是有限局部常值层。
因此，由引理 \ref{etale-cohomology-lemma-characterize-finite-locally-constant}，
$\mathcal{F}$ 可由一个概形 $V$ 表示，它在 $Z$ 上是有限 \'etale 的。
旋子局部平凡，所以 $V \to Z$ 当然是满射。由引理
\ref{etale-cohomology-lemma-closed-of-affine-normal-scheme-with-alg-closed-function-field}，
$V \to Z$ 有截面，故得结论。
\end{proof}

\begin{lemma}
\label{etale-cohomology-lemma-gabber-for-absolutely-algebraically-closed}
设 $X$ 是函数域可分闭的正规整仿射概形。设
$Z \subset X$ 是一个闭子概形。对任意有限阿贝尔群 $M$，有
$H^q_\etale(Z, \underline{M}) = 0$ 对 $q \geq 1$ 成立。
\end{lemma}

\begin{proof}
写 $X = \Spec(R)$ 且 $Z = \Spec(R')$，于是有环满射
$R \to R'$。由引理
\ref{etale-cohomology-lemma-normal-scheme-with-alg-closed-function-field} 和代数，引理
\ref{algebra-lemma-quotient-strict-henselization}，$R'$ 的所有局部环
都严格 Hensel。此外可知，对任意 $f' \in R'$，都有满射
$R_f \to R'_{f'}$，其中 $f \in R$ 是 $f'$ 的一个提升。
由于 $R_f$ 是分式域可分闭的正规整环，由引理
\ref{etale-cohomology-lemma-gabber-for-h1-absolutely-algebraically-closed} 可知
$H^1_\etale(D(f'), \underline{M}) = 0$。因此可以对
$Z = \Spec(R')$ 应用引理 \ref{etale-cohomology-lemma-gabber-criterion} 得到结论。
\end{proof}

\begin{lemma}
\label{etale-cohomology-lemma-integral-cover-trivial-cohomology}
设 $X$ 是一个仿射概形。
\begin{enumerate}
\item 存在整满射 $X' \to X$，使得对每个闭子概形
$Z' \subset X'$、每个有限阿贝尔群 $M$ 以及每个 $q \geq 1$，都有
$H^q_\etale(Z', \underline{M}) = 0$。
\item 对任意闭子概形 $Z \subset X$、有限阿贝尔群 $M$、
$q \geq 1$ 以及 $\xi \in H^q_\etale(Z, \underline{M})$，存在有限表现的
有限满射 $X' \to X$，使得 $\xi$ 在
$H^q_\etale(X' \times_X Z, \underline{M})$ 中的拉回为零。
\end{enumerate}
\end{lemma}

\begin{proof}
写 $X = \Spec(A)$。写 $A = \mathbf{Z}[x_i]/J$，其中 $J$ 是某个理想。
设 $R$ 是 $\mathbf{Z}[x_i]$ 的分式域的
某个代数闭包中 $\mathbf{Z}[x_i]$ 的整闭包。令
$A' = R/JR$，并置 $X' = \Spec(A')$。由引理
\ref{etale-cohomology-lemma-gabber-for-absolutely-algebraically-closed}，
这给出 (1) 中的一个例子。

\medskip\noindent
证明 (2)。设 $X' \to X$ 是上面找到的整满射。显然，$\xi$ 在
$H^q_\etale(X' \times_X Z, \underline{M})$ 中的像为零。可以把
$X'$ 写成极限 $X' = \lim X'_i$，其中各概形在 $X$ 上有限且有限表现；
在当前仿射情形下这很容易做到，而一般地，这是极限，
引理 \ref{limits-lemma-integral-limit-finite-and-finite-presentation}
的一个特殊情形。由引理 \ref{etale-cohomology-lemma-directed-colimit-cohomology}，
可知 $\xi$ 在 $H^q_\etale(X'_i \times_X Z, \underline{M})$
中的像为零，其中 $i$ 充分大。
\end{proof}








\section{绝对整闭情形的消失}
\label{etale-cohomology-section-aic-vanishing}

\noindent
回忆，若环 $R$ 上的每个首一多项式都在 $R$ 中有根，则称 $R$
是绝对整闭的（代数进阶，定义
\ref{more-algebra-definition-absolutely-integrally-closed}）。
本节证明 $\Spec(R)$ 以有限挠群为系数的 \'etale 上同调在正次数消失
（命题 \ref{etale-cohomology-proposition-aic-vanishing}），从而略微改进先前的引理
\ref{etale-cohomology-lemma-gabber-for-absolutely-algebraically-closed}。
建议读者跳过本节。

\begin{lemma}
\label{etale-cohomology-lemma-find-extension}
设 $A$ 是一个环。设 $a, b \in A$，满足
$aA + bA = A$，且 $a \bmod bA$ 是一个单位根。
则存在单生成扩张 $A \subset B$ 和元素 $y \in B$，
使得 $u = a - by$ 是一个单位。
\end{lemma}

\begin{proof}
设 $a^n \equiv 1 \bmod bA$。特别地，$a^i$ 模
$b^mA$ 对所有 $i, m \geq 1$ 都是单位。我们断言，存在
$a_1, \ldots, a_n \in A$，使得
$$
1 = a^n + a_1 a^{n - 1}b + a_2 a^{n - 2}b^2 + \ldots + a_n b^n
$$
事实上，由于 $1 - a^n \in bA$，可以找到元素 $a_1 \in A$，使得
$1 - a^n - a_1 a^{n - 1} b \in b^2A$；这里使用了
$a^{n - 1}$ 模 $bA$ 为单位这一事实。接着可以找到元素
$a_2 \in A$，使得
$1 - a^n - a_1 a^{n - 1} b - a_2 a^{n - 2} b^2 \in b^3A$。
如此继续。最终找到 $a_1, \ldots, a_{n - 1} \in A$，使得
$1 - (a^n + a_1 a^{n - 1}b + a_2 a^{n - 2}b^2 +
\ldots + a_{n - 1} ab^{n - 1}) \in b^nA$。于是可以找到
$a_n \in A$，使上述等式成立。

\medskip\noindent
对上述 $a_1, \ldots, a_n$，我们断言取
$$
B = A[y]/(y^n + a_1 y^{n - 1} + a_2 y^{n - 2} + \ldots + a_n)
$$
即可。事实上，设 $\mathfrak q \subset B$ 是位于
$\mathfrak p \subset A$ 之上的素理想。为得到矛盾，假设
$u = a - by$ 属于 $\mathfrak q$。若 $b \in \mathfrak p$，
则 $a \not \in \mathfrak p$，因为 $aA + bA = A$，从而
$u$ 不属于 $\mathfrak q$。因此可以假设 $b \not \in \mathfrak p$，
即 $b \not \in \mathfrak q$。这蕴含 $y \bmod \mathfrak q$
等于 $a/b \bmod \mathfrak q$。然而此时得到
$$
0 = y^n + a_1 y^{n - 1} + a_2 y^{n - 2} + \ldots + a_n =
b^{-n}(a^n + a_1 a^{n - 1}b + a_2 a^{n - 2}b^2 + \ldots + a_nb^n) =
b^{-n}
$$
，矛盾。证明完成。
\end{proof}

\noindent
为说明证明，需要引入一些群概形。固定一个素数 $\ell$。令
$$
A = \mathbf{Z}[\zeta] =
\mathbf{Z}[x]/(x^{\ell - 1} + x^{\ell - 2} + \ldots + 1)
$$
换言之，$A$ 是 $\mathbf{Z}$ 的单生成扩张，它由一个本原
$\ell$ 次单位根 $\zeta$ 生成。置
$$
\pi = \zeta - 1
$$
一个略去的计算表明，$\ell$ 可被 $\pi^{\ell - 1}$ 整除；
这里的计算在 $A$ 中进行。
第一个 $A$ 上群概形为
$$
G = \Spec(A[s, \frac{1}{\pi s + 1}])
$$
，其群律由余乘法
$$
\mu :
A[s, \frac{1}{\pi s + 1}]
\longrightarrow
A[s, \frac{1}{\pi s + 1}] \otimes_A A[s, \frac{1}{\pi s + 1}],\quad
s \longmapsto \pi s \otimes s + s \otimes 1 + 1 \otimes s
$$
给出。由此有
$$
\mu(\pi s + 1) = (\pi s + 1) \otimes (\pi s + 1)
$$
，所以确实得到所示的 $A$-代数映射。略去它确实定义群律的验证。
第二个 $A$ 上群概形为
$$
H = \Spec(A[t, \frac{1}{\pi^\ell t + 1}])
$$
，其群律由余乘法
$$
\mu : A[t, \frac{1}{\pi^\ell t + 1}]
\longrightarrow
A[t, \frac{1}{\pi^\ell t + 1}] \otimes_A A[t, \frac{1}{\pi^\ell t + 1}],\quad
t \longmapsto \pi^\ell t \otimes t + t \otimes 1 + 1 \otimes t
$$
给出。与前面相同的验证表明，这定义了一个群律。接着注意多项式
$$
\Phi(s) = \frac{(\pi s + 1)^\ell - 1}{\pi^\ell}
$$
属于 $A[s]$，次数为 $\ell$，并且关于 $s$ 首一。事实上，
$s^i$ 的系数在 $0 < i < \ell$ 时等于
${\ell \choose i}\pi^{i - \ell}$；由于
$\pi^{\ell - 1}$ 整除 $\ell$（在 $A$ 中），这是 $A$ 的元素。于是得到环映射
$$
A[t, \frac{1}{\pi^\ell t + 1}]
\longrightarrow
A[s, \frac{1}{\pi s + 1}],\quad
t \longmapsto \Phi(s)
$$
；读者容易验证它与余乘法相容。因此得到群概形态射
$$
f : G \to H
$$
下面的引理特别说明，该态射忠实平坦（事实上将看到它是有限
\'etale 满射）。

\begin{lemma}
\label{etale-cohomology-lemma-monogenic-one}
有
$$
A[s, \frac{1}{\pi s + 1}] =
\left(A[t, \frac{1}{\pi^\ell t + 1}]\right)[s]/(\Phi(s) - t)
$$
特别地，$G$ 的 Hopf 代数是 $H$ 的 Hopf 代数的单生成扩张。
\end{lemma}

\begin{proof}
这由上面的讨论和 $\Phi(s)$ 的形式得出。特别注意，利用
$\Phi(s) = t$，元素 $\frac{1}{\pi^\ell t + 1}$ 变为元素
$\frac{1}{(\pi s + 1)^\ell}$。
\end{proof}

\noindent
接下来计算 $f$ 的核。由于 $H$ 的单位元
由 $t = 0$ 在 $H$ 中给出，可知 $f$ 的核由 $\Phi(s) = 0$ 给出。现在注意，
下列 $A$-值点
$\sigma_0, \ldots, \sigma_{\ell - 1}$
属于 $G$，其定义为
$$
\sigma_i :
s = \frac{\zeta^i - 1}{\pi} = \frac{\zeta^i - 1}{\zeta - 1} =
\zeta^{i - 1} + \zeta^{i - 2} + \ldots + 1,\quad
i = 0, 1, \ldots, \ell - 1
$$
。它们显然包含在 $\Ker(f)$ 中。此外，它们在 $G \to \Spec(A)$
的{\bf 每个}纤维中都两两不同。读者还可计算出
$\sigma_i +_G \sigma_j = \sigma_{i + j \bmod \ell}$。因此得到群概形的闭浸入
$$
\underline{\mathbf{Z}/\ell \mathbf{Z}}_A \longrightarrow \Ker(f)
$$
，它把 $i$ 送到 $\sigma_i$。然而，按构造，$\Ker(f)$ 在
$\Spec(A)$ 上是次数为 $\ell$ 的有限平坦概形。因此该映射是同构。
总之，得到 $A$ 上群概形的短正合列
\begin{equation}
\label{etale-cohomology-equation-ses}
0 \to
\underline{\mathbf{Z}/\ell \mathbf{Z}}_A
\to G
\to H
\to 0
\end{equation}
。

\begin{lemma}
\label{etale-cohomology-lemma-lift-points-H-to-G}
设 $R$ 是一个绝对整闭的 $A$-代数。则 $G(R) \to H(R)$ 是满射。
\end{lemma}

\begin{proof}
设 $h \in H(R)$ 对应于 $A$-代数映射
$A[t, \frac{1}{\pi^\ell t + 1}] \to R$，它把 $t$ 送到 $a \in A$。
由于 $\Phi(s)$ 首一，可以找到 $b \in A$，使得 $\Phi(b) = a$。
由引理 \ref{etale-cohomology-lemma-monogenic-one}，把 $s$ 送到 $b$，得到唯一的
$A$-代数映射 $A[s, \frac{1}{\pi s + 1}] \to R$，它与上面的映射
$A[t, \frac{1}{\pi^\ell t + 1}] \to R$ 相容。后者又对应于元素
$g \in G(R)$，它映到 $h \in H(R)$。
\end{proof}

\begin{lemma}
\label{etale-cohomology-lemma-interpolate}
设 $R$ 是一个绝对整闭的 $A$-代数。设 $I, J \subset R$
是满足 $I + J = R$ 的理想。存在 $g \in G(R)$，使得
$g \bmod I = \sigma_0$ 且 $g \bmod J = \sigma_1$。
\end{lemma}

\begin{proof}
选取 $x \in I$，使得 $x \equiv 1 \bmod J$。可以并且确实把
$I$ 替换为 $xR$，把 $J$ 替换为 $(x - 1)R$。于是要找
$s \in R$，使得
\begin{enumerate}
\item $1 + \pi s$ 是单位；
\item $s \equiv 0 \bmod xR$；以及
\item $s \equiv 1 \bmod (x - 1)R$。
\end{enumerate}
后两个条件说明，$s = x + x(x - 1)y$，其中某个 $y \in R$。
第一个条件说明
$1 + \pi s = 1 + \pi x + \pi x (x - 1) y$ 需要是 $R$ 的单位。
但注意，$1 + \pi x$ 与 $\pi x (x - 1)$ 生成 $R$ 的单位理想，
而且 $1 + \pi x$ 是 $\ell$ 次方等于 $1$ 的元素（模
$\pi x (x - 1)$）
\footnote{因为 $1 + \pi x$ 同余于 $1$（模 $\pi$），同余于
$1$（模 $x$），并且同余于 $1 + \pi = \zeta$（模 $x - 1$）；
此外，$(\pi) \cap (x) \cap (x - 1) = (\pi x (x - 1))$ 在
$A[x]$ 中成立。}。
由引理 \ref{etale-cohomology-lemma-find-extension} 以及 $R$ 绝对整闭这一事实即得结论。
\end{proof}

\begin{proposition}
\label{etale-cohomology-proposition-aic-vanishing}
设 $R$ 是一个绝对整闭环。设 $M$ 是有限阿贝尔群。
则 $H^i_\etale(\Spec(R), \underline{M}) = 0$ 对 $i > 0$ 成立。
\end{proposition}

\begin{proof}
任意有限阿贝尔群都有一个有限滤过，其子商是素数阶循环群，
故可以假设 $M = \mathbf{Z}/\ell\mathbf{Z}$，其中 $\ell$ 是素数。

\medskip\noindent
注意，$R$ 的所有局部环都严格 Hensel，见代数进阶，引理
\ref{more-algebra-lemma-absolutely-integrally-closed-strictly-henselian}。
此外，由代数进阶，引理
\ref{more-algebra-lemma-absolutely-integrally-closed-quotient-localization}，
$R$ 的任意局部化也绝对整闭。因此，引理
\ref{etale-cohomology-lemma-gabber-criterion} 告诉我们，只需证明映射
$$
H^1_\etale(D(f + g), \mathbf{Z}/\ell\mathbf{Z})
\longrightarrow
H^1_\etale(D(f(f + g)), \mathbf{Z}/\ell\mathbf{Z}) \oplus
H^1_\etale(D(g(f + g)), \mathbf{Z}/\ell\mathbf{Z})
$$
的核对任意 $f, g \in R$ 都为零。把 $R$ 替换为 $R_{f + g}$ 后，
归结为如下断言：给定
$\xi \in H^1_\etale(\Spec(R), \mathbf{Z}/\ell\mathbf{Z})$
以及仿射开覆盖 $\Spec(R) = U \cup V$，若 $\xi|_U$ 和
$\xi|_V$ 平凡，则 $\xi = 0$。

\medskip\noindent
令 $A = \mathbf{Z}[\zeta]$ 如上。由于 $\mathbf{Z} \subset A$
是单生成的，可以找到环映射 $A \to R$。从现在起，把 $R$
看作 $A$-代数，并把 $\Spec(R)$ 看作 $\Spec(A)$ 上的概形。
把短正合列 (\ref{etale-cohomology-equation-ses}) 基变换到 $\Spec(R)$ 并取
\'etale 上同调，得到
$$
G(R) \to H(R) \to
H^1_\etale(\Spec(R), \mathbf{Z}/\ell\mathbf{Z}) \to
H^1_\etale(\Spec(R), G)
$$
请在证明余下部分始终记住这一点。

\medskip\noindent
设 $\tau \in \Gamma(U \cap V, \mathbf{Z}/\ell\mathbf{Z})$ 是一个截面，
它在 Mayer--Vietoris 正合列（引理 \ref{etale-cohomology-lemma-mayer-vietoris}）中的
边界给出 $\xi$。对 $i = 0, 1, \ldots, \ell - 1$，令
$A_i \subset U \cap V$ 为 $\tau$ 取值 $i \bmod \ell$ 的开闭子集。
于是有有限不交并分解
$$
U \cap V = A_0 \amalg \ldots \amalg A_{\ell - 1}
$$
，使得 $\tau$ 在每个 $A_i$ 上常值。对
$i = 0, 1, \ldots, \ell - 1$，用
$\tau_i \in H^0(U \cap V, \mathbf{Z}/\ell\mathbf{Z})$ 表示如下元素：
它取值 $1$ 的地方是 $A_i$，取值 $0$ 的地方是 $A_j$，其中
$j \not = i$。
于是 $\tau$ 是各 $\tau_i$ 的倍数之和\footnote{不计计算错误，
有 $\tau = \sum i \tau_i$。}。因此只需证明对应于 $\tau_i$
的上同调类平凡。这把问题归结到 $\tau$ 只取两个不同值
$1$ 和 $0$ 的情形。

\medskip\noindent
假设 $\tau$ 只取值 $1$ 和 $0$。写作
$$
U \cap V = A \amalg B
$$
其中 $A$ 是 $\tau = 0$ 的点集，$B$ 是
$\tau = 1$ 的点集。于是 $A$ 与 $B$ 是互不相交的闭子集。
用 $\overline{A}$ 和 $\overline{B}$ 表示 $A$ 与 $B$ 在
$\Spec(R)$ 中的闭包。于是得到一个“香蕉形”：即有
$$
\overline{A} \cap \overline{B} = Z_1 \amalg Z_2
$$
其中 $Z_1 \subset U$ 与 $Z_2 \subset V$ 是互不相交的闭子集。
置 $T_1 = \Spec(R) \setminus V$ 和 $T_2 = \Spec(R) \setminus U$。
注意 $Z_1 \subset T_1 \subset U$、$Z_2 \subset T_2 \subset V$，且
$T_1 \cap T_2 = \emptyset$。从拓扑上可写作
$$
\Spec(R) = \overline{A} \cup \overline{B} \cup T_1 \cup T_2
$$
建议画图以便直观理解。
为证明 $\xi$ 为零，我们可以并且确实把 $R$ 替换为其约化
（命题 \ref{etale-cohomology-proposition-topological-invariance}）。
下文把 $A$、$\overline{A}$、$B$、$\overline{B}$、$T_1$、$T_2$
看作 $\Spec(R)$ 的约化闭子概形。接着，在 $Z_1$
和 $Z_2$ 上采用如下概形结构：
$$
Z_1 = \overline{A} \cap (\overline{B} \cup T_1)
\quad\text{且}\quad
Z_2 = \overline{A} \cap (\overline{B} \cup T_2)
$$
（即态射，定义
\ref{morphisms-definition-scheme-theoretic-intersection-union}
中的概形论并与交）。

\medskip\noindent
用 $X$ 表示一个 $G$-旋子；它定义在 $\Spec(R)$ 上，并对应于
$\xi$ 在 $H^1(\Spec(R), G)$ 中的像。若 $X$ 平凡，则 $\xi$
来自一个元素 $h \in H(R)$（见上面的上同调正合列）。
然而，由引理 \ref{etale-cohomology-lemma-lift-points-H-to-G}，元素 $h$
可提升为 $G(R)$ 的一个元素，因而如所愿得到
$\xi = 0$。所以我们的目标是证明 $X$ 平凡。

\medskip\noindent
回忆，嵌入 $\mathbf{Z}/\ell \mathbf{Z} \to G(R)$
把 $i \bmod \ell$ 送到 $\sigma_i \in G(R)$。注意 $\overline{A}$
是一个绝对整闭环的谱（即 $R$ 的一个商）。由
引理 \ref{etale-cohomology-lemma-interpolate}，可以找到 $g \in G(\overline{A})$，使得
$g|_{\overline{A} \cap Z_1} = \sigma_0$ 且
$g|_{\overline{A} \cap Z_2} = \sigma_1$（按概形论意义）。
于是可以定义
\begin{enumerate}
\item $g_1 \in G(U)$：它等于 $g$ 于 $\overline{A} \cap U$，
等于 $\sigma_0$ 于 $\overline{B} \cap U$，并且等于 $\sigma_0$ 于 $T_1$；以及
\item $g_2 \in G(V)$：它等于 $g$ 于 $\overline{A} \cap V$，
等于 $\sigma_1$ 于 $\overline{B} \cap V$，并且等于 $\sigma_1$ 于 $T_2$。
\end{enumerate}
具体而言，为找到 (1) 中的 $g_1$，把截面
$\sigma_0$ 在 $\Omega = (\overline{B} \cup T_1) \cap U$ 上的取值
与截面 $g$ 在 $\Omega' = \overline{A} \cap U$ 上的限制黏合。
注意 $U = \Omega \cup \Omega'$（按概形论意义），
因为 $U$ 是约化的，并且
$\Omega \cap \Omega' = Z_1$（按概形论意义）；这里使用了我们对
$Z_1$ 的选取。因此由
态射，引理 \ref{morphisms-lemma-scheme-theoretic-union}，
$U$ 是 $\Omega$ 与 $\Omega'$ 沿 $Z_1$ 的推出。
所以可以找到 $g_1$。
(2) 中 $g_2$ 的存在性同理。
于是有
$$
\tau = g_2|_{A \cup B} - g_1|_{A \cup B} \quad(\text{addition in group law})
$$
从而看出 $X$ 平凡，证明完成。
\end{proof}




















\section{真基变换的仿射类似}
\label{etale-cohomology-section-gabber-affine-proper}

\noindent
本节讨论 Ofer Gabber 的一个结果，见
\cite{gabber-affine-proper}。Roland Huber 也证明了这一结果，见
\cite{Huber-henselian}。Gabber 证明所需的一部分工作已经在
第 \ref{etale-cohomology-section-strictly-henselian-local-rings} 节中完成。

\begin{lemma}
\label{etale-cohomology-lemma-efface-cohomology-on-closed-by-finite-cover}
设 $X$ 是仿射概形。设 $\mathcal{F}$ 是 $X_\etale$ 上的挠阿贝尔层。
设 $Z \subset X$ 是闭子概形。设
$\xi \in H^q_\etale(Z, \mathcal{F}|_Z)$，其中 $q > 0$。
则存在挠阿贝尔层的单射 $\mathcal{F} \to \mathcal{F}'$
（在 $X_\etale$ 上），使得 $\xi$ 在
$H^q_\etale(Z, \mathcal{F}'|_Z)$ 中的像为零。
\end{lemma}

\begin{proof}
由引理 \ref{etale-cohomology-lemma-torsion-colimit-constructible} 和 \ref{etale-cohomology-lemma-colimit}，
可以找到态射 $\mathcal{G} \to \mathcal{F}$，其中 $\mathcal{G}$
是可构造阿贝尔层，且 $\xi$ 来自某个元素 $\zeta$，该元素属于
$H^q_\etale(Z, \mathcal{G}|_Z)$。假设可以找到挠阿贝尔层的单射
$\mathcal{G} \to \mathcal{G}'$（在 $X_\etale$ 上），使得 $\zeta$ 在
$H^q_\etale(Z, \mathcal{G}'|_Z)$ 中的像为零。
此时可以取 $\mathcal{F}'$ 为推出
$$
\mathcal{F}' = \mathcal{G}' \amalg_{\mathcal{G}} \mathcal{F}
$$
，从而得到引理的结论。（注意，限制到 $Z$ 是正合的，故与有限极限
和余极限交换；此外，它作为推前的左伴随，也与任意余极限交换。）
因此可以假设 $\mathcal{F}$ 可构造。

\medskip\noindent
假设 $\mathcal{F}$ 可构造。由
引理 \ref{etale-cohomology-lemma-constructible-maps-into-constant-general}，
只需在 $\mathcal{F}$ 形如 $f_*\underline{M}$ 时证明结论，
其中 $M$ 是有限阿贝尔群，而 $f : Y \to X$ 是有限呈示的有限态射
（由引理 \ref{etale-cohomology-lemma-finite-pushforward-constructible}，
这种层仍然可构造，但这里不需要这一点）。
由于 $f_*$ 的形成与任意基变换交换
（引理 \ref{etale-cohomology-lemma-finite-pushforward-commutes-with-base-change}），
$f_*\underline{M}$ 限制到 $Z$ 等于 $\underline{M}$ 沿
$Y \times_X Z \to Z$ 的推前。由 Leray 谱序列
（命题 \ref{etale-cohomology-proposition-leray}）
和高阶正像的消失
（命题 \ref{etale-cohomology-proposition-finite-higher-direct-image-zero}），得到
$$
H^q_\etale(Z, f_*\underline{M}|_Z) = H^q_\etale(Y \times_X Z, \underline{M}).
$$
由引理 \ref{etale-cohomology-lemma-integral-cover-trivial-cohomology}，
可以找到有限呈示的有限满态射 $Y' \to Y$，使得 $\xi$ 在
$H^q(Y' \times_X Z, \underline{M})$ 中的像为零。
用 $f' : Y' \to X$ 表示复合 $Y' \to Y \to X$。我们断言映射
$$
f_*\underline{M} \longrightarrow f'_*\underline{M}
$$
是单射；结合上述讨论，这就完成证明。为验证所需的单射性，
可以考察茎。具体地，若 $\overline{x} : \Spec(k) \to X$ 是几何点，则
$$
(f_*\underline{M})_{\overline{x}} =
\bigoplus\nolimits_{f(\overline{y}) = \overline{x}} M
$$
，这由命题 \ref{etale-cohomology-proposition-finite-higher-direct-image-zero} 得到；
另一个层也同样。由于 $Y' \to Y$ 是有限满态射，
它在位于 $\overline{x}$ 上方的几何点之间诱导的映射也满，
故得结论。
\end{proof}

\noindent
上面的引理将处理 Gabber 结果中的高阶上同调群。
下面的引理用于处理整体截面。

\begin{lemma}
\label{etale-cohomology-lemma-gabber-h0}
设 $X$ 是拟紧拟分离概形。设 $i : Z \to X$ 是闭浸入。假设
\begin{enumerate}
\item 对任意层 $\mathcal{F}$（在 $X_{Zar}$ 上），映射
$\Gamma(X, \mathcal{F}) \to \Gamma(Z, i^{-1}\mathcal{F})$
是双射；并且
\item 对任意有限态射 $X' \to X$，假设 (1) 对
$Z \times_X X' \to X'$ 成立。
\end{enumerate}
则对任意层 $\mathcal{F}$（在 $X_\etale$ 上），有
$\Gamma(X, \mathcal{F}) = \Gamma(Z, i^{-1}_{small}\mathcal{F})$。
\end{lemma}

\begin{proof}
设 $\mathcal{F}$ 是 $X_\etale$ 上的层。有典范的（基变换）映射
$$
i^{-1}(\mathcal{F}|_{X_{Zar}})
\longrightarrow
(i_{small}^{-1}\mathcal{F})|_{Z_{Zar}}
$$
，这是 $Z_{Zar}$ 上层的映射。通过考察茎来证明该映射为单射。
左边在 $z \in Z$ 处的茎是 $\mathcal{F}|_{X_{Zar}}$ 在 $z$ 处的茎。
右边的茎是对所有初等 \'etale 邻域
$(U, u) \to (X, z)$ 取余极限，其中 $U \times_X Z \to Z$
在 $z$ 的某个邻域上有截面。由于 \'etale 态射是开映射，
$U \to X$ 的像是开邻域 $U_0$，它是 $z$ 在 $X$ 中的邻域。
映射 $\mathcal{F}(U_0) \to \mathcal{F}(U)$ 是单射；
这是由层 $\mathcal{F}$ 关于 \'etale 覆盖 $U \to U_0$ 的层条件得到的。
对所有 $U$ 和 $U_0$ 取余极限，便得到茎上的单射性。

\medskip\noindent
由此和假设 (1)，映射
$\Gamma(X, \mathcal{F}) \to \Gamma(Z, i^{-1}_{small}\mathcal{F})$
是单射。由 (2)，对所有 $X'$（在 $X$ 上有限）也有同一结论。

\medskip\noindent
在证明满射性之前先引入一些记号。对任意对象 $U$（属于 $X_\etale$）
和 $s \in \mathcal{F}(U)$，用
$i^{-1}s \in (i_{small}^{-1}\mathcal{F})(U \times_X Z)$
表示拉回（更准确地说，这是 $s$ 在位点，引理
\ref{sites-lemma-recover} 的映射与从预层拉回到层拉回的映射之复合下的像）。
这一构造与 $X_\etale$ 中态射的拉回以及（有限）态射
$X' \to X$ 的拉回相容，具体细节留给读者说明。

\medskip\noindent
设 $t \in \Gamma(Z, i^{-1}_{small}\mathcal{F})$。由
$i^{-1}_{small}\mathcal{F}$ 的构造，存在 \'etale 覆盖
$\{V_j \to Z\}_{j \in J}$、\'etale 态射 $U_j \to X$、截面
$s_j \in \mathcal{F}(U_j)$ 以及态射 $V_j \to U_j$（在 $X$ 上），
使得 $t|_{V_j}$ 是 $i^{-1}s_j$ 沿
$V_j \to U_j \times_X Z$ 的拉回。由于 $V_j \to U_j \times_X Z$
是 \'etale 的，其像是开子概形。由于
$U_j \times_X Z \subset U_j$ 是闭的，把 $U_j$ 替换为一个开子概形后，
可以假设 $V_j \to U_j \times_X Z$ 满。于是
$\{V_j \to U_j \times_X Z\}$ 是 \'etale 覆盖，从而
$t|_{U_j \times_X Z} = i^{-1}s_j$。
注意，每个非空闭子概形 $T \subset X$ 都与 $Z$ 相交；例如，把假设 (1)
应用于层 $(T \to X)_*\underline{\mathbf{Z}}$ 即得。
因此 $\coprod U_j \to X$ 满。由态射进阶，引理
\ref{more-morphisms-lemma-there-is-a-scheme-integral-over}
，可以找到有限满态射 $X' \to X$，使得 $X' \to X$ 在 Zariski 局部
经 $\coprod U_j \to X$ 分解。设
$X' = \bigcup_{\alpha \in A} W_\alpha$ 是开覆盖，$j : A \to J$ 是映射，
并且存在态射 $g_\alpha : W_\alpha \to U_{j(\alpha)}$（在 $X$ 上）。
用 $\mathcal{F}'$ 表示 $\mathcal{F}$ 到 $X'_\etale$ 的拉回。
用 $s'_\alpha \in \mathcal{F}'(W_\alpha)$ 表示
$s_{j(\alpha)}$ 沿 $g_\alpha$ 的拉回。
置 $Z' = X' \times_X Z$。用 $i' : Z' \to X'$ 表示 $i$ 的基变换。
于是 $(i'_{small})^{-1}\mathcal{F}' =
(Z' \to Z)_{small}^{-1}i_{small}^{-1}\mathcal{F}$。
在这个等同下，用
$t' \in
\Gamma(Z', (i')^{-1}_{small}\mathcal{F}')$
表示 $t$ 到 $Z'$ 的拉回。
按构造，$t'|_{W_\alpha \times_X Z}$ 等于
$(i')^{-1}s'_\alpha$。因此 $t'$ 是如下子层的截面：
$$
(i')^{-1}(\mathcal{F}'|_{X'_{Zar}})
\subset
(i'_{small})^{-1}\mathcal{F}')|_{Z'_{Zar}}
$$
由假设 (2)，截面 $t'$ 来自截面
$s' \in \Gamma(X', \mathcal{F}'|_{X'_{Zar}}) = \Gamma(X', \mathcal{F}')$。
由第二段证明的单射性，$s'$ 到 $X' \times_X X'$ 的两个拉回相同
（毕竟 $t'$ 到 $Z' \times_Z Z'$ 的拉回具有这一性质）。
因此由注 \ref{etale-cohomology-remark-cohomological-descent-finite}，
$s'$ 来自 $\mathcal{F}$ 在 $X$ 上的一个截面。
\end{proof}

\begin{lemma}
\label{etale-cohomology-lemma-connected-topological}
设 $Z \subset X$ 是拓扑空间 $X$ 的闭子集。假设
\begin{enumerate}
\item $X$ 是谱空间
（拓扑，定义 \ref{topology-definition-spectral-space}）；并且
\item 对 $x \in X$，交 $Z \cap \overline{\{x\}}$
连通（特别地非空）。
\end{enumerate}
若 $Z = Z_1 \amalg Z_2$，且 $Z_i$ 在 $Z$ 中闭，
则存在分解 $X = X_1 \amalg X_2$，其中 $X_i$ 在 $X$ 中闭，
并且 $Z_i = Z \cap X_i$。
\end{lemma}

\begin{proof}
注意 $Z_i$ 拟紧。因此，由拓扑，引理
\ref{topology-lemma-make-spectral-space}，点集
$W_i$（其点特化到 $Z_i$）在可构造拓扑中闭。假设 (2) 蕴含
$X = W_1 \amalg W_2$。
设 $x \in \overline{W_1}$。由拓扑，引理
\ref{topology-lemma-constructible-stable-specialization-closed} 的 (1)，
存在特化 $x_1 \leadsto x$，其中 $x_1 \in W_1$。
于是 $\overline{\{x\}} \subset \overline{\{x_1\}}$，从而
$x \in W_1$。换言之，取 $X_i = W_i$ 即可。
\end{proof}

\begin{lemma}
\label{etale-cohomology-lemma-h0-topological}
设 $Z \subset X$ 是拓扑空间 $X$ 的闭子集。假设
\begin{enumerate}
\item $X$ 是谱空间
（拓扑，定义 \ref{topology-definition-spectral-space}）；并且
\item 对 $x \in X$，交 $Z \cap \overline{\{x\}}$
连通（特别地非空）。
\end{enumerate}
则对任意层 $\mathcal{F}$（在 $X$ 上），有
$\Gamma(X, \mathcal{F}) = \Gamma(Z, \mathcal{F}|_Z)$。
\end{lemma}

\begin{proof}
若 $x \leadsto x'$ 是点的特化，则有典范映射
$\mathcal{F}_{x'} \to \mathcal{F}_x$；它与开集上的截面相容，
并且关于 $\mathcal{F}$ 函子性。由于 $X$ 的每个点都特化到 $Z$ 中的点，
可知 $\Gamma(X, \mathcal{F}) \to \Gamma(Z, \mathcal{F}|_Z)$ 为单射。
困难之处是证明它满。

\medskip\noindent
用 $\mathcal{B}$ 表示 $X$ 的全体拟紧开集。把 $\mathcal{F}$ 写成滤过余极限
$\mathcal{F} = \colim \mathcal{F}_i$，其中每个 $\mathcal{F}_i$
都如模，等式 (\ref{modules-equation-towards-constructible-sets}) 所述。
见模，引理 \ref{modules-lemma-filtered-colimit-constructibles}。
于是有 $\mathcal{F}|_Z = \colim \mathcal{F}_i|_Z$；
这是因为限制到 $Z$ 是左伴随（范畴，引理
\ref{categories-lemma-adjoint-exact} 和层，引理
\ref{sheaves-lemma-f-map}）。
由层，引理 \ref{sheaves-lemma-directed-colimits-sections}，函子
$\Gamma(X, -)$ 与 $\Gamma(Z, -)$ 和滤过余极限交换。
因此可以假设层 $\mathcal{F}$ 如模，等式
(\ref{modules-equation-towards-constructible-sets}) 所述。

\medskip\noindent
假设有嵌入 $\mathcal{F} \subset \mathcal{G}$。则
$$
\Gamma(X, \mathcal{F}) =
\Gamma(Z, \mathcal{F}|_Z) \cap \Gamma(X, \mathcal{G})
$$
，其中交取在 $\Gamma(Z, \mathcal{G}|_Z)$ 内。
这由证明开头的观察得出：通过考察茎，可以检验 $\mathcal{G}$
的整体截面是否属于 $\mathcal{F}$，而 $X$ 的每个点都特化到 $Z$ 中的点。

\medskip\noindent
由模，引理 \ref{modules-lemma-constructible-in-constant}，存在单射
$\mathcal{F} \to \prod (Z_i \to X)_*\underline{S_i}$，
其中乘积有限，$Z_i \subset X$ 闭，且 $S_i$ 有限。
因此只需对层 $(Z_i \to X)_*\underline{S_i}$ 证明满射性。注意
$$
\Gamma(X, (Z_i \to X)_*\underline{S_i}) = \Gamma(Z_i, \underline{S_i})
\quad\text{且}\quad
\Gamma(X, (Z_i \to X)_*\underline{S_i}|_Z) =
\Gamma(Z \cap Z_i, \underline{S_i})
$$
此外，条件 (1) 和 (2) 由 $Z_i$ 继承；对 (2) 这是显然的，
对 (1) 则由拓扑，引理 \ref{topology-lemma-spectral-sub} 得出。
因此只需在（有限）常值层的情形证明引理。
这一情形正是引理 \ref{etale-cohomology-lemma-connected-topological} 的重新表述，
证明完成。
\end{proof}

\begin{example}
\label{etale-cohomology-example-quinard}
若 $X$ 不是谱空间，则引理 \ref{etale-cohomology-lemma-h0-topological} 不成立。
下面给出一个例子。设 $Y$ 是 $T_1$ 拓扑空间，$y \in Y$ 是非开点。
令 $X = Y \amalg \{ x \}$，并赋予如下拓扑：其闭集为
$\emptyset$、$\{y\}$ 以及所有 $F \amalg \{ x \}$，
其中 $F$ 是 $Y$ 的闭子集。于是 $Z = \{x, y\}$ 是 $X$ 的闭子集，
并满足引理 \ref{etale-cohomology-lemma-h0-topological} 的假设 (2)。
但 $X$ 连通而 $Z$ 不连通。因此，对取值于 $\{0, 1\}$ 的常值层
（在 $X$ 上），引理的结论不成立。
\end{example}

\begin{lemma}
\label{etale-cohomology-lemma-h0-henselian-pair}
设 $(A, I)$ 是 Hensel 对。置 $X = \Spec(A)$ 和
$Z = \Spec(A/I)$。对任意层 $\mathcal{F}$（在 $X_\etale$ 上），有
$\Gamma(X, \mathcal{F}) = \Gamma(Z, \mathcal{F}|_Z)$。
\end{lemma}

\begin{proof}
回忆，任意环的谱都是谱空间，见代数，引理
\ref{algebra-lemma-spec-spectral}。由代数进阶，引理
\ref{more-algebra-lemma-irreducible-henselian-pair-connected}
可知，$\overline{\{x\}} \cap Z$ 连通，其中 $x \in X$ 任意。
由引理 \ref{etale-cohomology-lemma-h0-topological}，该断言对 $X_{Zar}$ 上的层成立。
对任意有限态射 $X' \to X$，有 $X' = \Spec(A')$ 和
$Z \times_X X' = \Spec(A'/IA')$，且 $(A', IA')$ 是 Hensel 对，
见代数进阶，引理
\ref{more-algebra-lemma-integral-over-henselian-pair}
；故同一断言对 $(X')_{Zar}$ 上的层成立。
因此应用引理 \ref{etale-cohomology-lemma-gabber-h0} 即得结论。
\end{proof}

\noindent
最后，可以陈述并证明 Gabber 定理。

\begin{theorem}[Gabber]
\label{etale-cohomology-theorem-gabber}
设 $(A, I)$ 是 Hensel 对。置 $X = \Spec(A)$ 和
$Z = \Spec(A/I)$。对任意挠阿贝尔层 $\mathcal{F}$（在 $X_\etale$ 上），有
$H^q_\etale(X, \mathcal{F}) = H^q_\etale(Z, \mathcal{F}|_Z)$。
\end{theorem}

\begin{proof}
由引理 \ref{etale-cohomology-lemma-h0-henselian-pair}，结论对 $q = 0$ 成立。
设 $q \geq 1$。假设结论已在所有次数 $< q$ 中证明。
设 $\mathcal{F}$ 是挠阿贝尔层。令
$\mathcal{F} \to \mathcal{F}'$
为挠阿贝尔层的单射（稍后选取），其余核为 $\mathcal{Q}$，
从而有短正合列
$$
0 \to \mathcal{F} \to \mathcal{F}' \to \mathcal{Q} \to 0
$$
，这是 $X_\etale$ 上挠阿贝尔层的短正合列。它给出 $X$ 和 $Z$ 上
上同调长正合列之间的映射，其中一部分如下：
$$
\xymatrix{
H^{q - 1}_\etale(X, \mathcal{F}') \ar[d] \ar[r] &
H^{q - 1}_\etale(X, \mathcal{Q}) \ar[d] \ar[r] &
H^q_\etale(X, \mathcal{F}) \ar[d] \ar[r] &
H^q_\etale(X, \mathcal{F}') \ar[d] \\
H^{q - 1}_\etale(Z, \mathcal{F}'|_Z) \ar[r] &
H^{q - 1}_\etale(Z, \mathcal{Q}|_Z) \ar[r] &
H^q_\etale(Z, \mathcal{F}|_Z) \ar[r] &
H^q_\etale(Z, \mathcal{F}'|_Z)
}
$$
利用这个行正合的阿贝尔群交换图来完成证明。

\medskip\noindent
先证关于 $\mathcal{F}$ 的单射性。设 $\xi$ 是
$H^q_\etale(X, \mathcal{F})$ 的非零元素。把引理
\ref{etale-cohomology-lemma-efface-cohomology-on-closed-by-finite-cover} 应用于
$Z = X$（！），可以找到 $\mathcal{F} \subset \mathcal{F}'$，
使得 $\xi$ 向右的像为零。于是 $\xi$ 是
$H^{q - 1}_\etale(X, \mathcal{Q})$ 中某个元素的像，而对
$q - 1$ 的双射性说明 $\xi$ 在
$H^q_\etale(Z, \mathcal{F}|_Z)$ 中的像不为零。

\medskip\noindent
再证关于 $\mathcal{F}$ 的满射性。设 $\xi$ 是
$H^q_\etale(Z, \mathcal{F}|_Z)$ 的一个元素。把引理
\ref{etale-cohomology-lemma-efface-cohomology-on-closed-by-finite-cover} 应用于
$Z = Z$，可以找到 $\mathcal{F} \subset \mathcal{F}'$，
使得 $\xi$ 向右的像为零。于是 $\xi$ 是
$H^{q - 1}_\etale(Z, \mathcal{Q}|_Z)$ 中某个元素的像，而对
$q - 1$ 的双射性说明 $\xi$ 属于竖直映射的像。
\end{proof}

\begin{lemma}
\label{etale-cohomology-lemma-vanishing-restriction-injective}
设 $X$ 是具有仿射对角的概形，并可由 $n + 1$ 个仿射开集覆盖。
设 $Z \subset X$ 是闭子概形。设 $\mathcal{A}$ 是 $X_\etale$
上的挠环层，而 $\mathcal{I}$ 是 $\mathcal{A}$-模的内射层
（在 $X_\etale$ 上）。则
$H^q_\etale(Z, \mathcal{I}|_Z) = 0$ 对 $q > n$ 成立。
\end{lemma}

\begin{proof}
对 $n$ 作归纳。若 $n = 0$，则 $X$ 仿射。设
$X = \Spec(A)$ 且 $Z = \Spec(A/I)$。令 $A^h$ 为所有 \'etale
$A$-代数 $B$ 的滤过余极限，其中 $A/I \to B/IB$ 是同构。
于是 $(A^h, IA^h)$ 是 Hensel 对，且 $A/I = A^h/IA^h$；
见代数进阶，引理 \ref{more-algebra-lemma-henselization} 及其证明。
置 $X^h = \Spec(A^h)$。由定理 \ref{etale-cohomology-theorem-gabber}，有
$$
H^q_\etale(Z, \mathcal{I}|_Z) = H^q_\etale(X^h, \mathcal{I}|_{X^h})
$$
由定理 \ref{etale-cohomology-theorem-colimit}，有
$$
H^q_\etale(X^h, \mathcal{I}|_{X^h}) =
\colim_{A \to B} H^q_\etale(\Spec(B), \mathcal{I}|_{\Spec(B)})
$$
，其中余极限取遍上述 $A$-代数 $B$。
由于态射 $\Spec(B) \to \Spec(A)$ 是 \'etale 的，限制
$\mathcal{I}|_{\Spec(B)}$ 是 $\mathcal{A}|_{\Spec(B)}$-模的内射层
（位点上的上同调，引理
\ref{sites-cohomology-lemma-cohomology-of-open}）。
因此右边的上同调群为零，从而得到这一情形的结论。

\medskip\noindent
归纳步骤可用 Mayer--Vietoris 完成。具体地，假设
$X = U \cup V$，其中 $U$ 是 $n$ 个仿射开集之并，而 $V$ 仿射。
利用 $X$ 的对角是仿射的，可知 $U \cap V$ 是 $n$ 个仿射开集之并。
Mayer--Vietoris 给出正合列
$$
H^{q - 1}_\etale(U \cap V \cap Z, \mathcal{I}|_Z) \to
H^q_\etale(Z, \mathcal{I}|_Z) \to
H^q_\etale(U \cap Z, \mathcal{I}|_Z) \oplus
H^q_\etale(V \cap Z, \mathcal{I}|_Z)
$$
，由归纳假设，当 $q > n$ 时得到所需的消失。
\end{proof}





\section{曲线上挠层的上同调}
\label{etale-cohomology-section-vanishing-torsion}

\noindent
本节的目标是证明曲线上挠层上同调的基本有限性与消失结果，
见定理 \ref{etale-cohomology-theorem-vanishing-affine-curves}。
第 \ref{etale-cohomology-section-vanishing-torsion-coefficients} 节将讨论
Noether 环上挠模的可构造层。

\begin{situation}
\label{etale-cohomology-situation-what-to-prove}
这里 $k$ 是代数闭域，$X$ 是维数 $\leq 1$ 的分离有限型概形（在 $k$ 上），
而 $\mathcal{F}$ 是 $X_\etale$ 上的挠阿贝尔层。
\end{situation}

\noindent
在情形 \ref{etale-cohomology-situation-what-to-prove} 中，要证明下列断言：
\begin{enumerate}
\item
\label{etale-cohomology-item-vanishing}
$H^q_\etale(X, \mathcal{F}) = 0$ 对 $q > 2$ 成立；
\item
\label{etale-cohomology-item-vanishing-affine}
$H^q_\etale(X, \mathcal{F}) = 0$ 对 $q > 1$ 成立，若 $X$ 仿射；
\item
\label{etale-cohomology-item-vanishing-p-p}
$H^q_\etale(X, \mathcal{F}) = 0$ 对 $q > 1$ 成立，若
$p = \text{char}(k) > 0$ 且 $\mathcal{F}$ 是 $p$-幂挠的；
\item
\label{etale-cohomology-item-finite-prime-to-p}
$H^q_\etale(X, \mathcal{F})$ 有限，若 $\mathcal{F}$
可构造且是与 $\text{char}(k)$ 互素的挠层；
\item
\label{etale-cohomology-item-finite-proper}
$H^q_\etale(X, \mathcal{F})$ 有限，若 $X$ 真且
$\mathcal{F}$ 可构造；
\item
\label{etale-cohomology-item-base-change-prime-to-p}
$H^q_\etale(X, \mathcal{F}) \to
H^q_\etale(X_{k'}, \mathcal{F}|_{X_{k'}})$ 对代数闭域的任意扩张
$k'/k$ 都是同构，前提是 $\mathcal{F}$ 是与
$\text{char}(k)$ 互素的挠层；
\item
\label{etale-cohomology-item-base-change-proper}
$H^q_\etale(X, \mathcal{F}) \to
H^q_\etale(X_{k'}, \mathcal{F}|_{X_{k'}})$ 对代数闭域的任意扩张
$k'/k$ 都是同构，前提是 $X$ 真；
\item
\label{etale-cohomology-item-surjective}
$H^2_\etale(X, \mathcal{F}) \to H^2_\etale(U, \mathcal{F})$
对所有开子集 $U \subset X$ 都是满射。
\end{enumerate}
给定任意情形 \ref{etale-cohomology-situation-what-to-prove}，若适用于该情形的断言都为真，
就说“断言 (\ref{etale-cohomology-item-vanishing})--(\ref{etale-cohomology-item-surjective}) 成立”。
证明从常系数上同调计算的如下推论开始。

\begin{lemma}
\label{etale-cohomology-lemma-constant-smooth-statements}
在情形 \ref{etale-cohomology-situation-what-to-prove} 中，假设 $X$ 光滑，且
$\mathcal{F} = \underline{\mathbf{Z}/\ell\mathbf{Z}}$，其中 $\ell$ 是素数。
则断言 (\ref{etale-cohomology-item-vanishing})--(\ref{etale-cohomology-item-surjective})
对 $\mathcal{F}$ 成立。
\end{lemma}

\begin{proof}
由于 $X$ 光滑，$X$ 是光滑曲线与点的有限不交并。
点（$k$ 的谱）上任意层的高阶 \'etale 上同调平凡，例如可由
引理 \ref{etale-cohomology-lemma-compare-cohomology-point}
或 \ref{etale-cohomology-lemma-affine-only-closed-points} 得到。
因此可以假设 $X$ 是光滑曲线。

\medskip\noindent
情形 I：$\ell$ 不同于 $k$ 的特征。
这一情形由引理 \ref{etale-cohomology-lemma-cohomology-smooth-projective-curve}
（射影情形）和引理 \ref{etale-cohomology-lemma-vanishing-cohomology-mu-smooth-curve}
（仿射情形）得到。关于上同调和代数闭基域扩张的断言
(\ref{etale-cohomology-item-base-change-prime-to-p}) 来自如下事实：亏格 $g$
与“穿孔”数 $r$ 在从 $k$ 到 $k'$ 时不变。
断言 (\ref{etale-cohomology-item-surjective}) 也成立，因为
$H^2_\etale(U, \mathcal{F})$ 只要 $U \not = X$ 就为零；此时 $U$ 仿射
（簇，引理 \ref{varieties-lemma-proper-minus-point} 和
\ref{varieties-lemma-curve-affine-projective}）。

\medskip\noindent
情形 II：$\ell$ 等于 $k$ 的特征。
消失性由引理 \ref{etale-cohomology-lemma-vanishing-variety-char-p-p} 得到。
断言 (\ref{etale-cohomology-item-finite-proper}) 和 (\ref{etale-cohomology-item-base-change-proper})
由引理 \ref{etale-cohomology-lemma-finiteness-proper-variety-char-p-p} 得到。
\end{proof}

\begin{remark}[代数闭基域扩张下的不变性]
\label{etale-cohomology-remark-invariance}
设 $k$ 是特征 $p > 0$ 的代数闭域。在第
\ref{etale-cohomology-section-artin-schreier} 节中已经看到正合列
$$
k[x] \to k[x] \to
H^1_\etale(\mathbf{A}^1_k, \mathbf{Z}/p\mathbf{Z}) \to 0
$$
，其中第一个箭头把 $f(x)$ 送到 $f^p - f$。余核的一组代表元由多项式
$$
\sum\nolimits_{p \not | n} \lambda_n x^n
$$
组成，其中 $\lambda_n \in k$。（若 $k$ 不是代数闭的，
还必须加入某些常数。）特别地，当 $k'/k$ 是代数闭扩张时，映射
$$
H^1_\etale(\mathbf{A}^1_k, \mathbf{Z}/p\mathbf{Z})
\to
H^1_\etale(\mathbf{A}^1_{k'}, \mathbf{Z}/p\mathbf{Z})
$$
一般不是同构。特别地，\'etale 基本群之间的映射
$\pi_1(\mathbf{A}^1_{k'}) \to \pi_1(\mathbf{A}^1_k)$
（在此插入未来的参考文献）也不是同构。因此，仿射直线的
\'etale 同伦型依赖于代数闭基域。由上面的引理
\ref{etale-cohomology-lemma-constant-smooth-statements} 可知，这一现象只在特征
$p$ 且系数为 $p$-幂挠时发生。
\end{remark}

\begin{lemma}
\label{etale-cohomology-lemma-ses-statements}
设 $k$ 是代数闭域。设 $X$ 是 $k$ 上维数 $\leq 1$ 的分离有限型概形。
设
$0 \to \mathcal{F}_1 \to \mathcal{F} \to \mathcal{F}_2 \to 0$
是 $X$ 上挠阿贝尔层的短正合列。若断言
(\ref{etale-cohomology-item-vanishing})--(\ref{etale-cohomology-item-surjective}) 对
$\mathcal{F}_1$ 和 $\mathcal{F}_2$ 成立，则它们也对 $\mathcal{F}$ 成立。
\end{lemma}

\begin{proof}
这大都直接来自定义和上同调长正合列。还要注意，$\mathcal{F}$
可构造（相应地，是与 $k$ 的特征互素的挠层）当且仅当
$\mathcal{F}_1$ 与 $\mathcal{F}_2$ 都可构造
（相应地，是与 $k$ 的特征互素的挠层）。见命题
\ref{etale-cohomology-proposition-constructible-over-noetherian}。略去若干细节。
\end{proof}

\begin{lemma}
\label{etale-cohomology-lemma-finite-pushforward-statements}
设 $k$ 是代数闭域。设 $f : X \to Y$ 是 $k$ 上维数 $\leq 1$
的分离有限型概形之间的有限态射。设 $\mathcal{F}$ 是 $X$ 上的挠阿贝尔层。
若断言 (\ref{etale-cohomology-item-vanishing})--(\ref{etale-cohomology-item-surjective})
对 $\mathcal{F}$ 成立，则它们也对 $f_*\mathcal{F}$ 成立。
\end{lemma}

\begin{proof}
事实上，有 $H^q_\etale(X, \mathcal{F}) = H^q_\etale(Y, f_*\mathcal{F})$，
这是由 $R^qf_*$ 对 $q > 0$ 的消失
（命题 \ref{etale-cohomology-proposition-finite-higher-direct-image-zero}）和 Leray 谱序列
（位点上的上同调，引理 \ref{sites-cohomology-lemma-apply-Leray}）得到的。
对 (\ref{etale-cohomology-item-surjective})，使用 $f_*$ 的形成与任意基变换交换这一事实
（引理 \ref{etale-cohomology-lemma-finite-pushforward-commutes-with-base-change}）。
\end{proof}

\begin{lemma}
\label{etale-cohomology-lemma-restrict-to-open}
在情形 \ref{etale-cohomology-situation-what-to-prove} 中，假设 $\mathcal{F}$ 可构造。
设 $j : X' \to X$ 是稠密开子概形的包含。则断言
(\ref{etale-cohomology-item-vanishing})--(\ref{etale-cohomology-item-surjective}) 对 $\mathcal{F}$ 成立
当且仅当它们对 $j_!j^{-1}\mathcal{F}$ 成立。
\end{lemma}

\begin{proof}
由于 $X'$ 稠密，$Z = X \setminus X'$ 的维数为 $0$，
故它是有限集 $Z = \{x_1, \ldots, x_n\}$，其元素都是 $k$-有理点。
考虑短正合列
$$
0 \to j_!j^{-1}\mathcal{F} \to \mathcal{F} \to i_*i^{-1}\mathcal{F} \to 0
$$
，见引理 \ref{etale-cohomology-lemma-ses-associated-to-open}。注意
$H^q_\etale(X, i_*i^{-1}\mathcal{F}) = H^q_\etale(Z, i^*\mathcal{F})$.
事实上，$i : Z \to X$ 是闭浸入，因而有限；于是由命题
\ref{etale-cohomology-proposition-finite-higher-direct-image-zero}，
$R^qi_*$ 对 $q > 0$ 消失，故该等式由 Leray 谱序列得到
（位点上的上同调，引理 \ref{sites-cohomology-lemma-apply-Leray}）。
由于 $Z$ 是代数闭域之谱的不交并，可知
$H^q_\etale(Z, i^*\mathcal{F}) = 0$ 对 $q > 0$ 成立，并且
$$
H^0_\etale(Z, i^{-1}\mathcal{F}) =
\bigoplus\nolimits_{i = 1, \ldots, n} \mathcal{F}_{x_i}
$$
是有限的，因为 $\mathcal{F}_{x_i}$ 由 $\mathcal{F}$
可构造这一假设而有限。上同调长正合列给出正合列
$$
0 \to
H^0_\etale(X, j_!j^{-1}\mathcal{F}) \to
H^0_\etale(X, \mathcal{F}) \to
H^0_\etale(Z, i^{-1}\mathcal{F}) \to
H^1_\etale(X, j_!j^{-1}\mathcal{F}) \to
H^1_\etale(X, \mathcal{F}) \to 0
$$
以及同构 $H^q_\etale(X, j_!j^{-1}\mathcal{F}) \to
H^q_\etale(X, \mathcal{F})$，对 $q > 1$ 成立。

\medskip\noindent
此时容易推出，(\ref{etale-cohomology-item-vanishing})--(\ref{etale-cohomology-item-surjective})
中的每个断言对 $\mathcal{F}$ 成立当且仅当它对
$j_!j^{-1}\mathcal{F}$ 成立。补充几点以帮助读者：
(a) 若 $\mathcal{F}$ 是与 $k$ 的特征互素的挠层，则
$j_!j^{-1}\mathcal{F}$ 也是；
(b) 层 $j_!j^{-1}\mathcal{F}$ 可构造；
(c) 有 $H^0_\etale(Z, i^{-1}\mathcal{F}) =
H^0_\etale(Z_{k'}, i^{-1}\mathcal{F}|_{Z_{k'}})$；以及 (d)
若 $U \subset X$ 是开集，则 $U' = U \cap X'$ 在 $U$ 中稠密。
\end{proof}

\begin{lemma}
\label{etale-cohomology-lemma-even-easier}
在情形 \ref{etale-cohomology-situation-what-to-prove} 中，假设 $X$ 光滑。
设 $j : U \to X$ 是开浸入。设 $\ell$ 是素数，并令
$\mathcal{F} = j_!\underline{\mathbf{Z}/\ell\mathbf{Z}}$。
则断言 (\ref{etale-cohomology-item-vanishing})--(\ref{etale-cohomology-item-surjective})
对 $\mathcal{F}$ 成立。
\end{lemma}

\begin{proof}
由于 $X$ 光滑，它是光滑曲线的不交并，故可以假设 $X$
是一条曲线（即不可约）。此时要么 $U = \emptyset$，无需证明；
要么 $U \subset X$ 稠密。在后一情形，引理由
引理 \ref{etale-cohomology-lemma-constant-smooth-statements} 和
\ref{etale-cohomology-lemma-restrict-to-open} 得到。
\end{proof}

\begin{lemma}
\label{etale-cohomology-lemma-somewhat-easier}
在情形 \ref{etale-cohomology-situation-what-to-prove} 中，假设 $X$ 约化。
设 $j : U \to X$ 是开浸入。设 $\ell$ 是素数，并令
$\mathcal{F} = j_!\underline{\mathbf{Z}/\ell\mathbf{Z}}$。
则断言 (\ref{etale-cohomology-item-vanishing})--(\ref{etale-cohomology-item-surjective})
对 $\mathcal{F}$ 成立。
\end{lemma}

\begin{proof}
与引理 \ref{etale-cohomology-lemma-even-easier} 的区别是，这里不假设 $X$ 光滑。
设 $\nu : X^\nu \to X$ 是正规化态射。则 $\nu$ 有限
（簇，引理 \ref{varieties-lemma-normalization-locally-algebraic}），
且 $X^\nu$ 光滑
（簇，引理 \ref{varieties-lemma-regular-point-on-curve}）。
设 $j^\nu : U^\nu \to X^\nu$ 是 $U$ 的逆像。
由引理 \ref{etale-cohomology-lemma-even-easier}，结论对
$j^\nu_!\underline{\mathbf{Z}/\ell\mathbf{Z}}$ 成立。
由引理 \ref{etale-cohomology-lemma-finite-pushforward-statements}，结论对
$\nu_*j^\nu_!\underline{\mathbf{Z}/\ell\mathbf{Z}}$ 成立。
一般而言，
$\nu_*j^\nu_!\underline{\mathbf{Z}/\ell\mathbf{Z}}$ 不等于
$j_!\underline{\mathbf{Z}/\ell\mathbf{Z}}$，
但可以如下绕过这一点。由于 $X$ 约化，态射 $\nu : X^\nu \to X$
在某个稠密开集 $j' : X' \to X$ 上是同构
（簇，引理 \ref{varieties-lemma-normalization-locally-algebraic}）。
在这个开集上有等式
$$
(j')^{-1}(\nu_*j^\nu_!\underline{\mathbf{Z}/\ell\mathbf{Z}})
=
(j')^{-1}(j_!\underline{\mathbf{Z}/\ell\mathbf{Z}})
$$
。对 $j' : X' \to X$ 和上述两个层使用两次
引理 \ref{etale-cohomology-lemma-restrict-to-open} 即得结论。
\end{proof}

\begin{lemma}
\label{etale-cohomology-lemma-vanishing-easier}
在情形 \ref{etale-cohomology-situation-what-to-prove} 中，假设 $X$ 约化。
设 $j : U \to X$ 是开浸入，且 $U$ 连通。设 $\ell$ 是素数。
设 $\mathcal{G}$ 是 $\mathbf{F}_\ell$-向量空间的有限局部常值层
（在 $U$ 上）。
令 $\mathcal{F} = j_!\mathcal{G}$。则断言
(\ref{etale-cohomology-item-vanishing})--(\ref{etale-cohomology-item-surjective}) 对 $\mathcal{F}$ 成立。
\end{lemma}

\begin{proof}
设 $f : V \to U$ 是引理 \ref{etale-cohomology-lemma-pullback-filtered} 中那样的有限
\'etale 态射，其次数与 $\ell$ 互素。第 \ref{etale-cohomology-section-trace-method} 节
的讨论给出映射
$$
\mathcal{G} \to f_*f^{-1}\mathcal{G} \to \mathcal{G}
$$
，其复合是同构。因此只需对
$\mathcal{F} = j_!f_*f^{-1}\mathcal{G}$ 证明引理。
由 Zariski 主定理（态射进阶，引理
\ref{more-morphisms-lemma-quasi-finite-separated-pass-through-finite})
，可以选取图
$$
\xymatrix{
V \ar[r]_{j'} \ar[d]_f & Y \ar[d]^{\overline{f}} \\
U \ar[r]^j & X
}
$$
，其中 $\overline{f} : Y \to X$ 有限，而 $j'$ 是像稠密的开浸入。
可以把 $Y$ 替换为其约化（这不改变 $V$，因为 $V$
作为 $U$ 上的 \'etale 概形是约化的）。
由于 $f$ 有限且 $V$ 在 $Y$ 中稠密，有 $V = U \times_X Y$。
由引理 \ref{etale-cohomology-lemma-compatible-shriek-push-finite}，有
$$
j_!f_*f^{-1}\mathcal{G} = \overline{f}_*j'_!f^{-1}\mathcal{G}
$$
由引理 \ref{etale-cohomology-lemma-finite-pushforward-statements}，只需考虑
$j'_!f^{-1}\mathcal{G}$。利用引理 \ref{etale-cohomology-lemma-pullback-filtered}
给出的滤过、$j'_!$ 的正合性以及引理 \ref{etale-cohomology-lemma-ses-statements}，
归结到
$\mathcal{F} = j'_!\underline{\mathbf{Z}/\ell\mathbf{Z}}$
的情形，这正是引理 \ref{etale-cohomology-lemma-somewhat-easier}。
\end{proof}


%10.20.09

\begin{theorem}
\label{etale-cohomology-theorem-vanishing-affine-curves}
若 $k$ 是代数闭域，$X$ 是维数 $\leq 1$ 的分离有限型概形（在 $k$ 上），
而 $\mathcal{F}$ 是 $X_\etale$ 上的挠阿贝尔层，则
\begin{enumerate}
\item
$H^q_\etale(X, \mathcal{F}) = 0$ 对 $q > 2$ 成立；
\item
$H^q_\etale(X, \mathcal{F}) = 0$ 对 $q > 1$ 成立，若 $X$ 仿射；
\item
$H^q_\etale(X, \mathcal{F}) = 0$ 对 $q > 1$ 成立，若
$p = \text{char}(k) > 0$ 且 $\mathcal{F}$ 是 $p$-幂挠的；
\item
$H^q_\etale(X, \mathcal{F})$ 有限，若 $\mathcal{F}$
可构造且是与 $\text{char}(k)$ 互素的挠层；
\item
$H^q_\etale(X, \mathcal{F})$ 有限，若 $X$ 真且
$\mathcal{F}$ 可构造；
\item
$H^q_\etale(X, \mathcal{F}) \to
H^q_\etale(X_{k'}, \mathcal{F}|_{X_{k'}})$ 对代数闭域的任意扩张
$k'/k$ 都是同构，前提是 $\mathcal{F}$ 是与
$\text{char}(k)$ 互素的挠层；
\item
$H^q_\etale(X, \mathcal{F}) \to
H^q_\etale(X_{k'}, \mathcal{F}|_{X_{k'}})$ 对代数闭域的任意扩张
$k'/k$ 都是同构，前提是 $X$ 真；
\item
$H^2_\etale(X, \mathcal{F}) \to H^2_\etale(U, \mathcal{F})$
对所有开子集 $U \subset X$ 都是满射。
\end{enumerate}
\end{theorem}

\begin{proof}
定理所说的正是：在情形 \ref{etale-cohomology-situation-what-to-prove} 中，
断言 (\ref{etale-cohomology-item-vanishing})--(\ref{etale-cohomology-item-surjective}) 成立。
第一步是把 $X$ 替换为其约化；命题
\ref{etale-cohomology-proposition-topological-invariance} 说明这是允许的。
由引理 \ref{etale-cohomology-lemma-torsion-colimit-constructible}，可以把
$\mathcal{F}$ 写成可构造阿贝尔层的滤过余极限。
取上同调与余极限交换，见引理 \ref{etale-cohomology-lemma-colimit}。
此外，沿 $X_{k'} \to X$ 的拉回作为左伴随与余极限交换。
因此只需对可构造层证明这些断言。

\medskip\noindent
本段中不再另行说明地使用引理 \ref{etale-cohomology-lemma-ses-statements}。写作
$\mathcal{F} = \mathcal{F}_1 \oplus \ldots \oplus \mathcal{F}_r$
，其中 $\mathcal{F}_i$ 对某个素数 $\ell_i$ 是 $\ell_i$-准素的，
便可假设 $\ell^n$ 零化 $\mathcal{F}$，其中 $\ell$ 为某个素数。
现在考虑正合列
$$
0 \to \mathcal{F}[\ell] \to \mathcal{F} \to \mathcal{F}/\mathcal{F}[\ell] \to 0.
$$
由此只需假设 $\mathcal{F}$ 是 $\ell$-挠的。这意味着
$\mathcal{F}$ 是 $\mathbf{F}_\ell$-向量空间的可构造层，
其中 $\ell$ 为某个素数。

\medskip\noindent
按定义，这意味着存在稠密开集 $U \subset X$，使得
$\mathcal{F}|_U$ 是 $\mathbf{F}_\ell$-向量空间的有限局部常值层。
由于 $\dim(X) \leq 1$，缩小 $U$ 后，可以假设
$U = U_1 \amalg \ldots \amalg U_n$ 是不可约概形的不交并
（只需去掉位于 $\geq 2$ 个分支之交中的 $U$ 的闭点）。
由引理 \ref{etale-cohomology-lemma-restrict-to-open}，归结到
$\mathcal{F} = j_!\mathcal{G}$ 的情形，其中 $\mathcal{G}$
是 $\mathbf{F}_\ell$-向量空间的有限局部常值层，定义在 $U$ 上。

\medskip\noindent
由于所选 $U = U_1 \amalg \ldots \amalg U_n$ 中每个 $U_i$ 都不可约，有
$$
j_!\mathcal{G} =
j_{1!}(\mathcal{G}|_{U_1}) \oplus \ldots \oplus
j_{n!}(\mathcal{G}|_{U_n})
$$
，其中 $j_i : U_i \to X$ 是包含态射。
$j_{i!}(\mathcal{G}|_{U_i})$ 的情形由引理
\ref{etale-cohomology-lemma-vanishing-easier} 处理。
\end{proof}

\begin{theorem}
\label{etale-cohomology-theorem-vanishing-curves}
设 $X$ 是维数为 $1$ 的有限型概形（在代数闭域 $k$ 上）。
设 $\mathcal{F}$ 是 $X_\etale$ 上的挠层。则
$$
H_\etale^q(X, \mathcal{F}) = 0, \quad \forall q \geq 3.
$$
若 $X$ 仿射，则还有 $H_\etale^2(X, \mathcal{F}) = 0$。
\end{theorem}

\begin{proof}
若 $X$ 分离，结论立即来自更精确的定理
\ref{etale-cohomology-theorem-vanishing-affine-curves}。
若 $X$ 不分离，选取仿射开覆盖
$X = X_1 \cup \ldots \cup X_n$。对 $n$ 归纳，可以假设消失性在
$U = X_1 \cup \ldots \cup X_{n - 1}$ 上成立。
于是 Mayer--Vietoris（引理 \ref{etale-cohomology-lemma-mayer-vietoris}）给出
$$
H^2_\etale(U, \mathcal{F}) \oplus H^2_\etale(X_n, \mathcal{F}) \to
H^2_\etale(U \cap X_n, \mathcal{F}) \to
H^3_\etale(X, \mathcal{F}) \to 0
$$
然而，$U \cap X_n$ 是仿射概形的开子概形，并且由维数假设它也仿射，
故群 $H^2_\etale(U \cap X_n, \mathcal{F})$ 由定理
\ref{etale-cohomology-theorem-vanishing-affine-curves} 消失。
\end{proof}

\begin{lemma}
\label{etale-cohomology-lemma-base-change-dim-1-separably-closed}
设 $k'/k$ 是可分闭域的扩张。设 $X$ 是 $k$ 上维数
$\leq 1$ 的真概形。设 $\mathcal{F}$ 是 $X$ 上的挠阿贝尔层。
则映射 $H^q_\etale(X, \mathcal{F}) \to
H^q_\etale(X_{k'}, \mathcal{F}|_{X_{k'}})$ 对 $q \geq 0$ 是同构。
\end{lemma}

\begin{proof}
定理 \ref{etale-cohomology-theorem-vanishing-affine-curves} 已证明代数闭域的情形。
给定引理中的 $k \subset k'$，可以选取图
$$
\xymatrix{
k' \ar[r] & \overline{k}' \\
k \ar[u] \ar[r] & \overline{k} \ar[u]
}
$$
，其中 $k \subset \overline{k}$ 和 $k' \subset \overline{k}'$
分别为代数闭包。由于 $k$ 与 $k'$ 可分闭，域扩张
$\overline{k}/k$ 与 $\overline{k}'/k'$
是代数纯不可分的。此时态射
$X_{\overline{k}} \to X$ 与 $X_{\overline{k}'} \to X_{k'}$
是普遍同胚。因此，$\mathcal{F}$ 的上同调可在
$X_{\overline{k}}$ 上计算，而 $\mathcal{F}|_{X_{k'}}$ 的上同调可在
$X_{\overline{k}'}$ 上计算，见命题
\ref{etale-cohomology-proposition-topological-invariance}。
故一般情形由代数闭域的情形推出。
\end{proof}






\section{曲线上挠模的上同调}
\label{etale-cohomology-section-vanishing-torsion-coefficients}

\noindent
本节对 Noether 环上的挠模可构造层重复第
\ref{etale-cohomology-section-vanishing-torsion} 节的论证。先从最有意思的一步开始。

\begin{lemma}
\label{etale-cohomology-lemma-constant-statements-coefficients}
设 $\Lambda$ 是 Noether 环，$M$ 是有限 $\Lambda$-模且被整数
$n > 0$ 零化，$k$ 是代数闭域，而 $X$ 是维数
$\leq 1$ 的分离有限型概形（在 $k$ 上）。则
\begin{enumerate}
\item $H^q_\etale(X, \underline{M})$ 是有限 $\Lambda$-模，
若 $n$ 与 $\text{char}(k)$ 互素；
\item $H^q_\etale(X, \underline{M})$ 是有限 $\Lambda$-模，
若 $X$ 真。
\end{enumerate}
\end{lemma}

\begin{proof}
若 $n = \ell n'$，其中 $\ell$ 是某个素数，则得到短正合列
$0 \to M[\ell] \to M \to M' \to 0$（有限 $\Lambda$-模），
且 $M'$ 被 $n'$ 零化。这给出相应的常值层短正合列，
进而给出上同调模的正合列
$$
H^q_\etale(X, \underline{M[n]}) \to
H^q_\etale(X, \underline{M}) \to
H^q_\etale(X, \underline{M'})
$$
因此，只要能证明 $M$ 被素数零化的情形，就可对 $n$ 归纳而得结论。

\medskip\noindent
设 $\ell$ 是素数，且 $\ell$ 零化 $M$。于是可把 $\Lambda$ 替换为
$\mathbf{F}_\ell$-代数 $\Lambda/\ell \Lambda$。事实上，
$\mathcal{F}$ 作为 $\Lambda$-模层的上同调与
$\mathcal{F}$ 作为 $\Lambda/\ell \Lambda$-模层的上同调相同；
例如可由位点上的上同调，引理
\ref{sites-cohomology-lemma-cohomology-modules-abelian-agree} 得到。

\medskip\noindent
假设 $\ell$ 是素数，并且 $\ell$ 同时零化 $M$ 和 $\Lambda$。
把问题归结到 $M$ 是有限自由 $\Lambda$-模的情形。
事实上，选取短正合列
$$
0 \to N \to \Lambda^{\oplus m} \to M \to 0
$$
这确定了正合列
$$
H^q_\etale(X, \underline{\Lambda^{\oplus m}}) \to
H^q_\etale(X, \underline{M}) \to
H^{q + 1}_\etale(X, \underline{N})
$$
对 $q$ 下降归纳便得到 $M$ 的结论，只要已知
$\Lambda^{\oplus m}$ 的结论。这里使用定理
\ref{etale-cohomology-theorem-vanishing-affine-curves}
所给出的上同调群在次数 $> 2$ 时消失。

\medskip\noindent
设 $\ell$ 是素数，并假设 $\ell$ 零化 $\Lambda$。尚需证明上同调群
$H^q_\etale(X, \underline{\Lambda})$ 是有限 $\Lambda$-模。
这里用一个技巧证明；“正确”的论证使用稍后才会证明的系数定理。
选取基 $\Lambda = \bigoplus_{i \in I} \mathbf{F}_\ell e_i$，
使得 $e_0 = 1$，其中 $0 \in I$。这一基的选取确定层同构
$$
\underline{\Lambda} = \bigoplus \underline{\mathbf{F}_\ell} e_i
$$
（在 $X_\etale$ 上）。因此
$$
H^q_\etale(X, \underline{\Lambda}) =
H^q_\etale(X, \bigoplus \underline{\mathbf{F}_\ell} e_i) =
\bigoplus H^q_\etale(X, \underline{\mathbf{F}_\ell})e_i
$$
，因为由定理 \ref{etale-cohomology-theorem-colimit}
（或引理 \ref{etale-cohomology-lemma-colimit}，或
引理 \ref{etale-cohomology-lemma-direct-sum-bounded-below-cohomology}），
在 $X$ 上取上同调与直和交换。已经知道
$H^q_\etale(X, \underline{\mathbf{F}_\ell})$ 是有限维
$\mathbf{F}_\ell$-向量空间
（由定理 \ref{etale-cohomology-theorem-vanishing-affine-curves}），
故 $H^q_\etale(X, \underline{\Lambda})$ 是同秩的自由
$\Lambda$-模。具体而言，给定一组基
$\xi_1, \ldots, \xi_m$（属于
$H^q_\etale(X, \underline{\mathbf{F}_\ell})$），则
$\xi_1 e_0, \ldots, \xi_m e_0$ 构成一组 $\Lambda$-基
（属于 $H^q_\etale(X, \underline{\Lambda})$）。
\end{proof}

\begin{lemma}
\label{etale-cohomology-lemma-finite-pushforward-coefficients}
设 $\Lambda$ 是 Noether 环，$k$ 是代数闭域，
$f : X \to Y$ 是 $k$ 上维数 $\leq 1$ 的分离有限型概形之间的有限态射，
而 $\mathcal{F}$ 是 $\Lambda$-模层（在 $X_\etale$ 上）。若
$H^q_\etale(X, \mathcal{F})$ 是有限 $\Lambda$-模，则
$H^q_\etale(Y, f_*\mathcal{F})$ 也如此。
\end{lemma}

\begin{proof}
事实上，有 $H^q_\etale(X, \mathcal{F}) = H^q_\etale(Y, f_*\mathcal{F})$，
这是由 $R^qf_*$ 对 $q > 0$ 的消失
（命题 \ref{etale-cohomology-proposition-finite-higher-direct-image-zero}）和 Leray 谱序列
（位点上的上同调，引理 \ref{sites-cohomology-lemma-apply-Leray}）得到的。
\end{proof}

\begin{lemma}
\label{etale-cohomology-lemma-restrict-to-open-coefficients}
设 $\Lambda$ 是 Noether 环，$k$ 是代数闭域，
$X$ 是 $k$ 上维数 $\leq 1$ 的分离有限型概形，
$\mathcal{F}$ 是可构造 $\Lambda$-模层（在 $X_\etale$ 上），
而 $j : X' \to X$ 是稠密开子概形的包含。则
$H^q_\etale(X, \mathcal{F})$ 是有限 $\Lambda$-模，当且仅当
$H^q_\etale(X, j_!j^{-1}\mathcal{F})$ 是有限 $\Lambda$-模。
\end{lemma}

\begin{proof}
由于 $X'$ 稠密，$Z = X \setminus X'$ 的维数为 $0$，
故它是有限集 $Z = \{x_1, \ldots, x_n\}$，其元素都是 $k$-有理点。
考虑短正合列
$$
0 \to j_!j^{-1}\mathcal{F} \to \mathcal{F} \to i_*i^{-1}\mathcal{F} \to 0
$$
，见引理 \ref{etale-cohomology-lemma-ses-associated-to-open}。注意
$H^q_\etale(X, i_*i^{-1}\mathcal{F}) = H^q_\etale(Z, i^*\mathcal{F})$.
事实上，$i : Z \to X$ 是闭浸入，因而有限；于是由命题
\ref{etale-cohomology-proposition-finite-higher-direct-image-zero}，
$R^qi_*$ 对 $q > 0$ 消失，故该等式由 Leray 谱序列得到
（位点上的上同调，引理 \ref{sites-cohomology-lemma-apply-Leray}）。
由于 $Z$ 是代数闭域之谱的不交并，可知
$H^q_\etale(Z, i^*\mathcal{F}) = 0$ 对 $q > 0$ 成立，并且
$$
H^0_\etale(Z, i^{-1}\mathcal{F}) =
\bigoplus\nolimits_{i = 1, \ldots, n} \mathcal{F}_{x_i}
$$
是有限 $\Lambda$-模；$\mathcal{F}_{x_i}$ 的有限性来自
$\mathcal{F}$ 是可构造 $\Lambda$-模层这一假设。
上同调长正合列给出正合列
$$
0 \to
H^0_\etale(X, j_!j^{-1}\mathcal{F}) \to
H^0_\etale(X, \mathcal{F}) \to
H^0_\etale(Z, i^{-1}\mathcal{F}) \to
H^1_\etale(X, j_!j^{-1}\mathcal{F}) \to
H^1_\etale(X, \mathcal{F}) \to 0
$$
以及同构 $H^0_\etale(X, j_!j^{-1}\mathcal{F}) \to
H^0_\etale(X, \mathcal{F})$，对 $q > 1$ 成立。引理由此容易推出。
\end{proof}

\begin{lemma}
\label{etale-cohomology-lemma-somewhat-easier-coefficients}
设 $\Lambda$ 是 Noether 环，$M$ 是有限 $\Lambda$-模且被整数
$n > 0$ 零化，$k$ 是代数闭域，$X$ 是维数
$\leq 1$ 的分离有限型概形（在 $k$ 上），而
$j : U \to X$ 是开浸入。则
\begin{enumerate}
\item $H^q_\etale(X, j_!\underline{M})$ 是有限 $\Lambda$-模，
若 $n$ 与 $\text{char}(k)$ 互素；
\item $H^q_\etale(X, j_!\underline{M})$ 是有限 $\Lambda$-模，
若 $X$ 真。
\end{enumerate}
\end{lemma}

\begin{proof}
由于 $\dim(X) \leq 1$，存在开集 $V \subset X$，它与 $U$ 不交，
使得 $X' = U \cup V$ 是 $X$ 中的稠密开集（略去细节）。
若 $j' : X' \to X$ 表示包含态射，则
$j_!\underline{M}$ 是 $j'_!\underline{M}$ 的直和项。
因此只需证明 $U$ 在 $X$ 中开且稠密的情形。这一情形由引理
\ref{etale-cohomology-lemma-restrict-to-open-coefficients} 和
\ref{etale-cohomology-lemma-constant-statements-coefficients} 得到。
\end{proof}

\begin{lemma}
\label{etale-cohomology-lemma-ses-statements-coefficients}
设 $\Lambda$ 是 Noether 环，$k$ 是代数闭域，
$X$ 是 $k$ 上维数 $\leq 1$ 的分离有限型概形，而
$0 \to \mathcal{F}_1 \to \mathcal{F} \to \mathcal{F}_2 \to 0$
是 $\Lambda$-模层的短正合列（在 $X_\etale$ 上）。若
$H^q_\etale(X, \mathcal{F}_i)$（$i = 1, 2$）是有限
$\Lambda$-模，则 $H^q_\etale(X, \mathcal{F})$ 是有限
$\Lambda$-模。
\end{lemma}

\begin{proof}
这直接来自上同调长正合列。
\end{proof}

\begin{lemma}
\label{etale-cohomology-lemma-vanishing-easier-coefficients}
设 $\Lambda$ 是 Noether 环，$k$ 是代数闭域，
$X$ 是维数 $\leq 1$ 的分离有限型概形（在 $k$ 上），
$j : U \to X$ 是开浸入且 $U$ 连通，$\ell$ 是素数，
$n > 0$，而 $\mathcal{G}$ 是有限型局部常值
$\Lambda$-模层（在 $U_\etale$ 上），并被 $\ell^n$ 零化。则
\begin{enumerate}
\item $H^q_\etale(X, j_!\mathcal{G})$ 是有限 $\Lambda$-模，
若 $\ell$ 与 $\text{char}(k)$ 互素；
\item $H^q_\etale(X, j_!\mathcal{G})$ 是有限 $\Lambda$-模，
若 $X$ 真。
\end{enumerate}
\end{lemma}

\begin{proof}
设 $f : V \to U$ 是引理 \ref{etale-cohomology-lemma-pullback-filtered-modules}
中那样的有限 \'etale 态射，其次数与 $\ell$ 互素。
第 \ref{etale-cohomology-section-trace-method} 节的讨论给出映射
$$
\mathcal{G} \to f_*f^{-1}\mathcal{G} \to \mathcal{G}
$$
，其复合是同构。因此只需证明
$H^q_\etale(X, j_!f_*f^{-1}\mathcal{G})$ 的有限性。
由 Zariski 主定理（态射进阶，引理
\ref{more-morphisms-lemma-quasi-finite-separated-pass-through-finite}），
可以选取图
$$
\xymatrix{
V \ar[r]_{j'} \ar[d]_f & Y \ar[d]^{\overline{f}} \\
U \ar[r]^j & X
}
$$
，其中 $\overline{f} : Y \to X$ 有限，而 $j'$ 是像稠密的开浸入。
由于 $f$ 有限且 $V$ 在 $Y$ 中稠密，有 $V = U \times_X Y$。
由引理 \ref{etale-cohomology-lemma-compatible-shriek-push-finite}，有
$$
j_!f_*f^{-1}\mathcal{G} = \overline{f}_*j'_!f^{-1}\mathcal{G}
$$
由引理 \ref{etale-cohomology-lemma-finite-pushforward-coefficients}，只需考虑
$j'_!f^{-1}\mathcal{G}$。利用引理
\ref{etale-cohomology-lemma-pullback-filtered-modules} 给出的滤过、
$j'_!$ 的正合性以及引理 \ref{etale-cohomology-lemma-ses-statements-coefficients}，
归结到 $\mathcal{F} = j'_!\underline{M}$ 的情形，其中所用的是有限
$\Lambda$-模 $M$；这正是引理
\ref{etale-cohomology-lemma-somewhat-easier-coefficients}。
\end{proof}

\begin{theorem}
\label{etale-cohomology-theorem-vanishing-affine-curves-coefficients}
设 $\Lambda$ 是 Noether 环，$k$ 是代数闭域，
$X$ 是维数 $\leq 1$ 的分离有限型概形（在 $k$ 上），而
$\mathcal{F}$ 是可构造且为挠的 $\Lambda$-模层
（在 $X_\etale$ 上）。则
\begin{enumerate}
\item
\label{etale-cohomology-item-finite-prime-to-p-coefficients}
$H^q_\etale(X, \mathcal{F})$ 是有限 $\Lambda$-模，
若 $\mathcal{F}$ 是与 $\text{char}(k)$ 互素的挠层；
\item
\label{etale-cohomology-item-finite-proper-coefficients}
$H^q_\etale(X, \mathcal{F})$ 是有限 $\Lambda$-模，若 $X$ 真。
\end{enumerate}
\end{theorem}

\begin{proof}
不再另行说明。写作
$\mathcal{F} = \mathcal{F}_1 \oplus \ldots \oplus \mathcal{F}_r$，
其中 $\mathcal{F}_i$ 被 $\ell_i^{n_i}$ 零化，这里的
$\ell_i$ 是素数，而 $n_i > 0$ 是整数。由引理
\ref{etale-cohomology-lemma-ses-statements-coefficients}，只需对
$\mathcal{F}_i$ 证明定理。因此可以并且确实假设
$\ell^n$ 零化 $\mathcal{F}$，其中 $\ell$ 是素数而 $n > 0$ 是整数。

\medskip\noindent
由于 $\mathcal{F}$ 作为 $\Lambda$-模层是可构造的，存在稠密开集
$U \subset X$，使得 $\mathcal{F}|_U$ 是有限型局部常值
$\Lambda$-模层。由于 $\dim(X) \leq 1$，缩小 $U$ 后，可以假设
$U = U_1 \amalg \ldots \amalg U_n$ 是不可约概形的不交并
（只需去掉位于 $\geq 2$ 个分支之交中的 $U$ 的闭点）。
由引理 \ref{etale-cohomology-lemma-restrict-to-open-coefficients}，归结到
$\mathcal{F} = j_!\mathcal{G}$ 的情形，其中 $\mathcal{G}$
是有限型局部常值 $\Lambda$-模层，定义在 $U$ 上
（且被 $\ell^n$ 零化）。

\medskip\noindent
由于所选 $U = U_1 \amalg \ldots \amalg U_n$ 中每个 $U_i$
都不可约，有
$$
j_!\mathcal{G} =
j_{1!}(\mathcal{G}|_{U_1}) \oplus \ldots \oplus
j_{n!}(\mathcal{G}|_{U_n})
$$
，其中 $j_i : U_i \to X$ 是包含态射。
$j_{i!}(\mathcal{G}|_{U_i})$ 的情形由引理
\ref{etale-cohomology-lemma-vanishing-easier-coefficients} 处理。
\end{proof}








\section{真概形的第一上同调}
\label{etale-cohomology-section-finite-etale-over-proper}

\noindent
在基本群，第 \ref{pione-section-finite-etale-over-proper} 节中，
我们已经在某种意义下看到：取 $R^1f_*\underline{G}$ 与基变换交换，
若 $f : X \to Y$ 是真态射，而 $G$ 是有限群（不必交换）。
本节推出这些结果的一个有用推论。

\begin{lemma}
\label{etale-cohomology-lemma-proper-over-henselian-and-h1}
设 $A$ 是 Hensel 局部环。设 $X$ 是 $A$ 上的真概形，
其闭纤维为 $X_0$。设 $M$ 是有限阿贝尔群。则
$H^1_\etale(X, \underline{M}) = H^1_\etale(X_0, \underline{M})$。
\end{lemma}

\begin{proof}
由位点上的上同调，引理 \ref{sites-cohomology-lemma-torsors-h1}，
$H^1_\etale(X, \underline{M})$ 的一个元素对应于一个
$\underline{M}$-挠子 $\mathcal{F}$（在 $X_\etale$ 上）。
这种挠子显然是有限局部常值层。因此，由引理
\ref{etale-cohomology-lemma-characterize-finite-locally-constant}，
$\mathcal{F}$ 可由概形 $V$ 表示，而该概形在 $X$ 上有限
\'etale。反过来，若概形 $V$ 在 $X$ 上有限 \'etale，并带有
$M$-作用，使它成为 $M$-挠子（在 $X$ 上），则它给出一个上同调类。
$X_0$ 上的上同调类与 $X_0$ 上有限 \'etale 挠子之间也有同样的对应。
因此，引理是基本群，引理
\ref{pione-lemma-finite-etale-on-proper-over-henselian}
中范畴等价的推论。
\end{proof}

\noindent
下列技术引理是真基变换定理证明中的关键材料。
这一论证逐字适用于 $A$ 上任意特殊纤维维数 $\leq 1$ 的真概形；
但事实上，该结论是真基变换定理的推论，而在其证明中我们只需要
这里这个特殊版本。

\begin{lemma}
\label{etale-cohomology-lemma-efface-cohomology-on-fibre-by-finite-cover}
设 $A$ 是 Hensel 局部环。令 $X = \mathbf{P}^1_A$。
设 $X_0 \subset X$ 是闭纤维。设 $\ell$ 是素数。设
$\mathcal{I}$ 是 $\mathbf{Z}/\ell\mathbf{Z}$-模的内射层
（在 $X_\etale$ 上）。则
$H^q_\etale(X_0, \mathcal{I}|_{X_0}) = 0$ 对 $q > 0$ 成立。
\end{lemma}

\begin{proof}
注意 $X$ 是可由 $2$ 个仿射开集覆盖的分离概形。因此对 $q > 1$，
结论来自 Gabber 的真基变换定理的仿射版本，见引理
\ref{etale-cohomology-lemma-vanishing-restriction-injective}。故可假设 $q = 1$。令
$\xi \in H^1_\etale(X_0, \mathcal{I}|_{X_0})$。
目标：证明 $\xi$ 是 $0$。
由引理 \ref{etale-cohomology-lemma-torsion-colimit-constructible} 和
\ref{etale-cohomology-lemma-colimit}，可以找到映射 $\mathcal{F} \to \mathcal{I}$，
其中 $\mathcal{F}$ 是可构造的
$\mathbf{Z}/\ell\mathbf{Z}$-模层，而 $\xi$ 来自某个元素 $\zeta$
（属于 $H^1_\etale(X_0, \mathcal{F}|_{X_0})$）。
假设有层的单射 $\mathcal{F} \to \mathcal{F}'$，这些层是
$\mathbf{Z}/\ell\mathbf{Z}$-模层（在 $X_\etale$ 上）。
由于 $\mathcal{I}$ 内射，可以把给定映射
$\mathcal{F} \to \mathcal{I}$ 延拓为映射
$\mathcal{F}' \to \mathcal{I}$。在这种情况下，可以把
$\mathcal{F}$ 替换为 $\mathcal{F}'$，并把 $\zeta$ 替换为
$\zeta$ 在 $H^1_\etale(X_0, \mathcal{F}'|_{X_0})$ 中的像。
此外，若 $\mathcal{F} = \mathcal{F}_1 \oplus \mathcal{F}_2$ 是直和，
则可以把 $\mathcal{F}$ 替换为 $\mathcal{F}_i$，并把 $\zeta$
替换为 $\zeta$ 在 $H^1_\etale(X_0, \mathcal{F}_i|_{X_0})$ 中的像。

\medskip\noindent
由引理 \ref{etale-cohomology-lemma-constructible-maps-into-constant-general}
及以上说明，可以假设 $\mathcal{F}$ 形如 $f_*\underline{M}$，
其中 $M$ 是有限 $\mathbf{Z}/\ell\mathbf{Z}$-模，而
$f : Y \to X$ 是有限表现的有限态射
（由引理 \ref{etale-cohomology-lemma-finite-pushforward-constructible}，
这种层仍然可构造，但我们不会用到这一点）。
由于 $f_*$ 的形成与任意基变换交换
（引理 \ref{etale-cohomology-lemma-finite-pushforward-commutes-with-base-change}），
$f_*\underline{M}$ 在 $X_0$ 上的限制等于
$\underline{M}$ 沿特殊纤维的诱导态射 $Y_0 \to X_0$ 的推前。
由 Leray 谱序列（命题 \ref{etale-cohomology-proposition-leray}）和高阶直像的消失
（命题 \ref{etale-cohomology-proposition-finite-higher-direct-image-zero}），得到
$$
H^1_\etale(X_0, f_*\underline{M}|_{X_0}) = H^1_\etale(Y_0, \underline{M}).
$$
由于 $Y \to \Spec(A)$ 真，可用引理
\ref{etale-cohomology-lemma-proper-over-henselian-and-h1} 看出
$H^1_\etale(Y_0, \underline{M})$ 等于
$H^1_\etale(Y, \underline{M})$。因此，上同调类 $\zeta$
可提升为上同调类
$$
\tilde\zeta \in H^1_\etale(Y, \underline{M}) = H^1_\etale(X, f_*\underline{M})
$$
然而，$\tilde \zeta$ 在 $H^1_\etale(X, \mathcal{I})$ 中的像为零，
因为 $\mathcal{I}$ 内射；再由下图的交换性
$$
\xymatrix{
H^1_\etale(X, f_*\underline{M}) \ar[r] \ar[d] &
H^1_\etale(X, \mathcal{I}) \ar[d] \\
H^1_\etale(X_0, (f_*\underline{M})|_{X_0}) \ar[r] &
H^1_\etale(X_0, \mathcal{I}|_{X_0})
}
$$
可知，$\xi$（$\zeta$ 的像）也为零。
\end{proof}











\section{基变换的预备知识}
\label{etale-cohomology-section-base-change-preliminaries}

\noindent
若你只关心光滑基变换定理或真基变换定理，可以直接跳到相应各节。
本节及随后几节考虑概形的交换图
$$
\xymatrix{
X \ar[d]_f & Y \ar[l]^h \ar[d]^e \\
S & T \ar[l]_g
}
$$
；通常假设此图笛卡尔，即 $Y = X \times_S T$。
上述交换图诱导小 \'etale 位点的交换图
$$
\xymatrix{
X_\etale \ar[d]_{f_{small}} & Y_\etale \ar[d]^{e_{small}} \ar[l]^{h_{small}} \\
S_\etale & T_\etale \ar[l]_{g_{small}}
}
$$
。采用记号
$$
f^{-1} = f_{small}^{-1}, \quad
g_* = g_{small, *}, \quad
e^{-1} = e_{small}^{-1}, \text{ 且}\quad
h_* = h_{small, *}.
$$
由位点，第 \ref{sites-section-pullback} 节，得到基变换（或拉回）映射
$$
f^{-1}g_*\mathcal{F}
\longrightarrow
h_*e^{-1}\mathcal{F}
$$
，其中 $\mathcal{F}$ 是 $T_\etale$ 上的层。若 $\mathcal{F}$
是 $T_\etale$ 上的阿贝尔层，则得到导出基变换映射
$$
f^{-1}Rg_*\mathcal{F}
\longrightarrow
Rh_*e^{-1}\mathcal{F}
$$
，见位点上的上同调，引理
\ref{sites-cohomology-lemma-base-change-map-flat-case}。
最后，若 $K$ 是 $D(T_\etale)$ 的任意对象，则有基变换映射
$$
f^{-1}Rg_*K
\longrightarrow
Rh_*e^{-1}K
$$
，见位点上的上同调，评注
\ref{sites-cohomology-remark-base-change}。

\begin{lemma}
\label{etale-cohomology-lemma-base-change-local}
考虑概形的笛卡尔图
$$
\xymatrix{
X \ar[d]_f & Y \ar[l]^h \ar[d]^e \\
S & T \ar[l]_g
}
$$
设 $\{U_i \to X\}$ 是 \'etale 覆盖，使得 $U_i \to S$
可分解为 $U_i \to V_i \to S$，其中 $V_i \to S$ 为 \'etale，
并考虑笛卡尔图
$$
\xymatrix{
U_i \ar[d]_{f_i} & U_i \times_X Y \ar[l]^{h_i} \ar[d]^{e_i} \\
V_i & V_i \times_S T \ar[l]_{g_i}
}
$$
设 $\mathcal{F}$ 是 $T_\etale$ 上的层。设 $K$ 属于 $D(T_\etale)$。
令 $K_i = K|_{V_i \times_S T}$，并令
$\mathcal{F}_i = \mathcal{F}|_{V_i \times_S T}$。
\begin{enumerate}
\item 若 $f_i^{-1}g_{i, *}\mathcal{F}_i = h_{i, *}e_i^{-1}\mathcal{F}_i$
对所有 $i$ 成立，则 $f^{-1}g_*\mathcal{F} = h_*e^{-1}\mathcal{F}$。
\item 若 $f_i^{-1}Rg_{i, *}K_i = Rh_{i, *}e_i^{-1}K_i$
对所有 $i$ 成立，则 $f^{-1}Rg_*K = Rh_*e^{-1}K$。
\item 若 $\mathcal{F}$ 是阿贝尔层，并且
$f_i^{-1}R^qg_{i, *}\mathcal{F}_i = R^qh_{i, *}e_i^{-1}\mathcal{F}_i$
对所有 $i$ 成立，则
$f^{-1}R^qg_*\mathcal{F} = R^qh_*e^{-1}\mathcal{F}$。
\end{enumerate}
\end{lemma}

\begin{proof}
证明 (1)。首先注意
$$
(f^{-1}g_*\mathcal{F})|_{U_i} =
f_i^{-1}(g_*\mathcal{F}|_{V_i}) =
f_i^{-1}g_{i, *}\mathcal{F}_i
$$
第一个等式成立是因为 $U_i \to X \to S$ 等于
$U_i \to V_i \to S$；第二个等式成立是因为由位点，引理
\ref{sites-lemma-localize-morphism-strong}，有
$g_*\mathcal{F}|_{V_i} = g_{i, *}\mathcal{F}_i$。类似地有
$$
(h_*e^{-1}\mathcal{F})|_{U_i} =
h_{i, *}(e^{-1}\mathcal{F}|_{U_i \times_X Y}) =
h_{i, *}e_i^{-1}\mathcal{F}_i
$$
因此，若基变换映射
$f_i^{-1}g_{i, *}\mathcal{F}_i \to h_{i, *}e_i^{-1}\mathcal{F}_i$
对所有 $i$ 都是同构，则基变换映射
$f^{-1}g_*\mathcal{F} \to h_*e^{-1}\mathcal{F}$
在 $U_i$ 上的限制对所有 $i$ 都是同构。由于 $\{U_i \to X\}$
是 \'etale 覆盖，可知它本身是同构。

\medskip\noindent
对其余两个断言，只需把对位点，引理
\ref{sites-lemma-localize-morphism-strong} 的引用替换为对
位点上的上同调，引理
\ref{sites-cohomology-lemma-restrict-direct-image-open} 的引用。
\end{proof}

\begin{lemma}
\label{etale-cohomology-lemma-base-change-compose}
考虑概形笛卡尔图组成的塔
$$
\xymatrix{
W \ar[d]_i & Z \ar[l]^j \ar[d]^k \\
X \ar[d]_f & Y \ar[l]^h \ar[d]^e \\
S & T \ar[l]_g
}
$$
设 $K$ 属于 $D(T_\etale)$。若
$$
f^{-1}Rg_*K \to Rh_*e^{-1}K
\quad\text{且}\quad
i^{-1}Rh_*e^{-1}K \to Rj_*k^{-1}e^{-1}K
$$
是同构，则
$(f \circ i)^{-1}Rg_*K \to Rj_*(e \circ k)^{-1}K$
是同构。类似地，若 $\mathcal{F}$ 是 $T_\etale$ 上的阿贝尔层，并且
$$
f^{-1}R^qg_*\mathcal{F} \to R^qh_*e^{-1}\mathcal{F}
\quad\text{且}\quad
i^{-1}R^qh_*e^{-1}\mathcal{F} \to R^qj_*k^{-1}e^{-1}\mathcal{F}
$$
是同构，则
$(f \circ i)^{-1}R^qg_*\mathcal{F} \to R^qj_*(e \circ k)^{-1}\mathcal{F}$
是同构。
\end{lemma}

\begin{proof}
这是形式的，只要验证这些基变换映射的复合就是外矩形的基变换映射；
见位点上的上同调，评注
\ref{sites-cohomology-remark-compose-base-change-horizontal}。
\end{proof}

\begin{lemma}
\label{etale-cohomology-lemma-base-change-Rf-star-colim}
设 $I$ 是有向集。考虑概形笛卡尔图的逆系统
$$
\xymatrix{
X_i \ar[d]_{f_i} & Y_i \ar[l]^{h_i} \ar[d]^{e_i} \\
S_i & T_i \ar[l]_{g_i}
}
$$
，其转移态射仿射，且 $g_i$ 拟紧拟分离。令 $X = \lim X_i$、
$S = \lim S_i$、$T = \lim T_i$ 和 $Y = \lim Y_i$，得到笛卡尔图
$$
\xymatrix{
X \ar[d]_f & Y \ar[l]^h \ar[d]^e \\
S & T \ar[l]_g
}
$$
设 $(\mathcal{F}_i, \varphi_{i'i})$ 是定义
\ref{etale-cohomology-definition-inverse-system-sheaves} 中那样的层系统（在
$(T_i)$ 上）。令 $\mathcal{F} = \colim p_i^{-1}\mathcal{F}_i$
（在 $T$ 上），其中 $p_i : T \to T_i$ 是投影。则有
\begin{enumerate}
\item 若 $f_i^{-1}g_{i, *}\mathcal{F}_i = h_{i, *}e_i^{-1}\mathcal{F}_i$
对所有 $i$ 成立，则
$f^{-1}g_*\mathcal{F} = h_*e^{-1}\mathcal{F}$。
\item 若 $\mathcal{F}_i$ 对所有 $i$ 都是阿贝尔层，并且
$f_i^{-1}R^qg_{i, *}\mathcal{F}_i = R^qh_{i, *}e_i^{-1}\mathcal{F}_i$
对所有 $i$ 成立，则
$f^{-1}R^qg_*\mathcal{F} = R^qh_*e^{-1}\mathcal{F}$。
\end{enumerate}
\end{lemma}

\begin{proof}
证明 (2)，略去 (1) 的证明。以下不再另行说明地使用：
层的拉回作为左伴随与余极限交换。注意 $h_i$ 作为 $g_i$ 的基变换
是拟紧拟分离的。以 $q_i : Y \to Y_i$ 表示投影，注意
$e^{-1}\mathcal{F} = \colim e^{-1}p_i^{-1}\mathcal{F}_i =
\colim q_i^{-1}e_i^{-1}\mathcal{F}_i$.
由引理 \ref{etale-cohomology-lemma-relative-colimit-general}，这给出
$$
R^qh_*e^{-1}\mathcal{F} = \colim r_i^{-1}R^qh_{i, *}e_i^{-1}\mathcal{F}_i
$$
，其中 $r_i : X \to X_i$ 是投影。类似地，有
$$
f^{-1}Rg_*\mathcal{F} =
f^{-1}\colim s_i^{-1}R^qg_{i, *}\mathcal{F}_i =
\colim r_i^{-1}f_i^{-1}R^qg_{i, *}\mathcal{F}_i
$$
，其中 $s_i : S \to S_i$ 是投影。引理得证。
\end{proof}

\begin{lemma}
\label{etale-cohomology-lemma-base-change-Rf-star-colim-complexes}
设 $I$、$X_i$、$Y_i$、$S_i$、$T_i$、$f_i$、$h_i$、$e_i$、$g_i$、
$X$、$Y$、$S$、$T$、$f$、$h$、$e$、$g$ 如引理
\ref{etale-cohomology-lemma-base-change-Rf-star-colim} 的陈述中所述。
设 $0 \in I$，并设 $K_0 \in D^+(T_{0, \etale})$。
对 $i \in I$、$i \geq 0$，以 $K_i$ 表示 $K_0$ 到 $T_i$ 的拉回。
以 $K$ 表示 $K_0$ 到 $T$ 的拉回。若
$f_i^{-1}Rg_{i, *}K_i = Rh_{i, *}e_i^{-1}K_i$
对所有 $i \geq 0$ 成立，则 $f^{-1}Rg_*K = Rh_*e^{-1}K$。
\end{lemma}

\begin{proof}
只需证明基变换映射
$f^{-1}Rg_*K \to Rh_*e^{-1}K$ 在上同调层上诱导同构。
换言之，必须证明 $f^{-1}R^pg_*K \to R^ph_*e^{-1}K$
对所有 $p \in \mathbf{Z}$ 都是同构，前提是已知
$f_i^{-1}R^pg_{i, *}K_i \to R^ph_{i, *}e_i^{-1}K_i$
对所有 $i \geq 0$ 和 $p \in \mathbf{Z}$ 都是同构。
此时可以完全照搬引理
\ref{etale-cohomology-lemma-base-change-Rf-star-colim} 的证明，只需把对引理
\ref{etale-cohomology-lemma-relative-colimit-general} 的引用替换为对引理
\ref{etale-cohomology-lemma-relative-colimit-general-complexes} 的引用。
\end{proof}

\begin{lemma}
\label{etale-cohomology-lemma-base-change-f-star-general-stalks}
考虑概形的笛卡尔图
$$
\xymatrix{
X \ar[d]_f & Y \ar[l]^h \ar[d]^e \\
S & T \ar[l]_g
}
$$
，其中 $g : T \to S$ 拟紧拟分离。设 $\mathcal{F}$
是 $T_\etale$ 上的阿贝尔层。设 $q \geq 0$。下列条件等价：
\begin{enumerate}
\item 对每个几何点 $\overline{x}$（属于 $X$），若其像为
$\overline{s} = f(\overline{x})$，则
$$
H^q(\Spec(\mathcal{O}^{sh}_{X, \overline{x}}) \times_S T, \mathcal{F})
=
H^q(\Spec(\mathcal{O}^{sh}_{S, \overline{s}}) \times_S T, \mathcal{F})
$$
\item $f^{-1}R^qg_*\mathcal{F} \to R^qh_*e^{-1}\mathcal{F}$
是同构。
\end{enumerate}
\end{lemma}

\begin{proof}
由于 $Y = X \times_S T$，有
$\Spec(\mathcal{O}^{sh}_{X, \overline{x}}) \times_X Y =
\Spec(\mathcal{O}^{sh}_{X, \overline{x}}) \times_S T$。
因此，由定理 \ref{etale-cohomology-theorem-higher-direct-images}
（以及引理 \ref{etale-cohomology-lemma-stalk-pullback}），(1) 中的映射就是
(2) 中映射在 $\overline{x}$ 处的茎映射。于是结论由定理
\ref{etale-cohomology-theorem-exactness-stalks} 得到。
\end{proof}

\begin{lemma}
\label{etale-cohomology-lemma-check-stalks-better}
设 $f : X \to S$ 是概形态射。设 $\overline{x}$ 是 $X$ 的几何点，
其像为 $\overline{s}$（在 $S$ 中）。设
$\Spec(K) \to \Spec(\mathcal{O}^{sh}_{S, \overline{s}})$ 是态射，
其中 $K$ 是可分闭域。设 $\mathcal{F}$ 是
$\Spec(K)_\etale$ 上的阿贝尔层。设 $q \geq 0$。下列条件等价：
\begin{enumerate}
\item
$H^q(\Spec(\mathcal{O}^{sh}_{X, \overline{x}}) \times_S \Spec(K), \mathcal{F}) =
H^q(\Spec(\mathcal{O}^{sh}_{S, \overline{s}}) \times_S \Spec(K), \mathcal{F})$
\item
$H^q(\Spec(\mathcal{O}^{sh}_{X, \overline{x}})
\times_{\Spec(\mathcal{O}^{sh}_{S, \overline{s}})} \Spec(K), \mathcal{F}) =
H^q(\Spec(K), \mathcal{F})$
\end{enumerate}
\end{lemma}

\begin{proof}
注意 $\Spec(K) \times_S \Spec(\mathcal{O}^{sh}_{S, \overline{s}})$
是 $K$ 上 \'etale 代数的滤过余极限之谱。由于 $K$ 可分闭，
每个 \'etale $K$-代数都是若干个 $K$ 的有限直积。因此可写成
$$
\Spec(K) \times_S \Spec(\mathcal{O}^{sh}_{S, \overline{s}}) =
\lim_{i \in I} \coprod\nolimits_{a \in A_i} \Spec(K)
$$
这个余滤过极限的每一项都是 $\Spec(K)$ 之不交并，
由有限集 $A_i$ 标号。注意 $A_i$ 非空，因为已给定
$\Spec(K) \to \Spec(\mathcal{O}^{sh}_{S, \overline{s}})$。于是
\begin{align*}
\Spec(\mathcal{O}^{sh}_{X, \overline{x}}) \times_S \Spec(K)
& =
\Spec(\mathcal{O}^{sh}_{X, \overline{x}})
\times_{\Spec(\mathcal{O}^{sh}_{S, \overline{s}})}
\left(
\Spec(\mathcal{O}^{sh}_{S, \overline{s}})
\times_S \Spec(K)\right) \\
& =
\lim_{i \in I} \coprod\nolimits_{a \in A_i}
\Spec(\mathcal{O}^{sh}_{X, \overline{x}})
\times_{\Spec(\mathcal{O}^{sh}_{S, \overline{s}})} \Spec(K)
\end{align*}
由于在当前情形取上同调与概形的极限交换
（定理 \ref{etale-cohomology-theorem-colimit}），结论随即成立。
\end{proof}













\section{正像的基变换}
\label{etale-cohomology-section-base-change-f-star}

\noindent
本节属于预备内容，初读时可略过。
本节讨论哪些态射 $f : X \to S$ 具有如下性质：对所有（集合）层都有
$f^{-1}g_* = h_*e^{-1}$，而且这对每个笛卡尔图
$$
\xymatrix{
X \ar[d]_f & Y \ar[l]^h \ar[d]^e \\
S & T \ar[l]_g
}
$$
都成立，只要 $g$ 拟紧拟分离。

\begin{lemma}
\label{etale-cohomology-lemma-base-change-f-star-general}
考虑概形的笛卡尔图
$$
\xymatrix{
X \ar[d]_f & Y \ar[l]^h \ar[d]^e \\
S & T \ar[l]_g
}
$$
。假设 $f$ 平坦，并且每个对象 $U$（属于 $X_\etale$）都有覆盖
$\{U_i \to U\}$，使得 $U_i \to S$ 可分解为
$U_i \to V_i \to S$，其中 $V_i \to S$ 为 \'etale，
而 $U_i \to V_i$ 拟紧且有几何连通纤维。
则对任意集合层 $\mathcal{F}$（在 $T_\etale$ 上），有
$f^{-1}g_*\mathcal{F} = h_*e^{-1}\mathcal{F}$。
\end{lemma}

\begin{proof}
设 $U \to X$ 是 \'etale 态射，且 $U \to S$ 可分解为
$U \to V \to S$，其中 $V \to S$ 为 \'etale，而
$U \to V$ 拟紧且有几何连通纤维。注意 $U \to V$ 平坦
（见平坦性的进一步结果，引理 \ref{flat-lemma-etale-flat-up-down}）。
我们断言
\begin{align*}
f^{-1}g_*\mathcal{F}(U)
& = g_*\mathcal{F}(V) \\
& = \mathcal{F}(V \times_S T) \\
& = e^{-1}\mathcal{F}(U \times_X Y) \\
& = h_*e^{-1}\mathcal{F}(U)
\end{align*}
。事实上，把 $U$ 看作 $X_\etale$ 的对象，把 $V$ 看作
$S_\etale$ 的对象，由引理 \ref{etale-cohomology-lemma-sections-upstairs}
得到第一个等式\footnote{严格地说，还用到了
$f^{-1}g_*\mathcal{F}$ 在 $U_\etale$ 上的限制，是经
$U \to V$ 拉回 $g_*\mathcal{F}$ 在 $V_\etale$ 上的限制所得；
见位点，引理 \ref{sites-lemma-localize-morphism-strong}。}。
把 $V \times_S T$ 看作 $T_\etale$ 的对象，由 $g_*$ 的定义
得到第二个等式。注意
$U \times_X Y = U \times_S T$（因为 $Y = X \times_S T$），
故 $U \times_X Y \to V \times_S T$ 是 $U \to V$ 的基变换，
从而有几何连通纤维。把 $U \times_X Y$ 看作 $Y_\etale$ 的对象，
与上述相同，由引理 \ref{etale-cohomology-lemma-sections-upstairs} 得到第三个等式。
最后，第四个等式来自 $h_*$ 的定义。

\medskip\noindent
由假设，$X_\etale$ 的每个对象都有一个 \'etale 覆盖，
上一段的论证适用于该覆盖，故引理成立。
\end{proof}

\begin{lemma}
\label{etale-cohomology-lemma-fppf-reduced-fibres-base-change-f-star}
考虑概形的笛卡尔图
$$
\xymatrix{
X \ar[d]_f & Y \ar[l]^h \ar[d]^e \\
S & T \ar[l]_g
}
$$
，其中 $f$ 平坦、局部有限表现，且有几何既约纤维。
则有 $f^{-1}g_*\mathcal{F} = h_*e^{-1}\mathcal{F}$，
其中 $\mathcal{F}$ 是 $T_\etale$ 上的任意层。
\end{lemma}

\begin{proof}
结合引理 \ref{etale-cohomology-lemma-base-change-f-star-general} 与
态射的进一步结果，引理
\ref{more-morphisms-lemma-cover-by-geometrically-connected} 即得。
\end{proof}

\begin{lemma}
\label{etale-cohomology-lemma-base-change-f-star-field}
考虑概形的笛卡尔图
$$
\xymatrix{
X \ar[d]_f & Y \ar[l]^h \ar[d]^e \\
S & T \ar[l]_g
}
$$
。假设 $S$ 是可分闭域之谱。
则有 $f^{-1}g_*\mathcal{F} = h_*e^{-1}\mathcal{F}$，
其中 $\mathcal{F}$ 是 $T_\etale$ 上的任意层。
\end{lemma}

\begin{proof}
可以在 $X$ 上局部论证，故可假设 $X$ 仿射。
于是可将 $X$ 写成 $S$ 上有限型仿射概形的余滤过极限。
由引理 \ref{etale-cohomology-lemma-base-change-Rf-star-colim}，可假设 $X$
在 $S$ 上有限型。此时引理 \ref{etale-cohomology-lemma-base-change-f-star-general}
适用，因为可分闭域上的任意有限型概形都是有限多个连通且
几何连通概形之不交并（见簇，引理
\ref{varieties-lemma-separably-closed-field-connected-components}）。
\end{proof}

\begin{lemma}
\label{etale-cohomology-lemma-base-change-f-star-valuation}
考虑概形的笛卡尔图
$$
\xymatrix{
X \ar[d]_f & Y \ar[l]^h \ar[d]^e \\
S & T \ar[l]_g
}
$$
。假设
\begin{enumerate}
\item $f$ 平坦且开；
\item $S$ 的剩余域都是可分代数闭的；
\item 给定 \'etale 态射 $U \to X$，其中 $U$ 为仿射概形，
可将 $U$ 写成 $X$ 的有限多个开子概形之不交并
（例如，当 $X$ 是函数域可分闭的正规整概形时）；
\item 任一纤维 $X_s$（属于 $f$）的每个非空开子集都连通
（例如，当 $X_s$ 不可约或为空时）。
\end{enumerate}
则对任意集合层 $\mathcal{F}$（在 $T_\etale$ 上），有
$f^{-1}g_*\mathcal{F} = h_*e^{-1}\mathcal{F}$。
\end{lemma}

\begin{proof}
略。提示：这些假设几乎直接蕴含引理
\ref{etale-cohomology-lemma-base-change-f-star-general} 的条件。
第 (3) 项括号内的例子来自引理
\ref{etale-cohomology-lemma-normal-scheme-with-alg-closed-function-field}。
\end{proof}

\noindent
下面的引理其实不属于这里，但似乎也没有更合适的位置。

\begin{lemma}
\label{etale-cohomology-lemma-fppf-reduced-fibres-pullback-products}
设 $f : X \to S$ 是平坦、局部有限表现且有几何既约纤维的概形态射。
则 $f^{-1} : \Sh(S_\etale) \to \Sh(X_\etale)$ 与积交换。
\end{lemma}

\begin{proof}
设 $I$ 是集合，并设 $\mathcal{G}_i$ 是 $S_\etale$ 上的层
（$i \in I$）。设 $U \to X$ 是 \'etale 态射，使得
$U \to S$ 可分解为 $U \to V \to S$，其中 $V \to S$ 为
\'etale，而 $U \to V$ 平坦、有限表现且有几何连通纤维。于是
\begin{align*}
f^{-1}(\prod \mathcal{G}_i)(U)
& =
(\prod \mathcal{G}_i)(V) \\
& =
\prod \mathcal{G}_i(V) \\
& =
\prod f^{-1}\mathcal{G}_i(U) \\
& =
(\prod f^{-1}\mathcal{G}_i)(U)
\end{align*}
，其中第一和第三个等式使用了引理
\ref{etale-cohomology-lemma-sections-upstairs}（还使用了如下事实：
$f^{-1}\mathcal{G}$ 在 $U_\etale$ 上的限制，是经
$U \to V$ 拉回 $\mathcal{G}$ 在 $V_\etale$ 上的限制所得；
见位点，引理 \ref{sites-lemma-localize-morphism-strong}）。
由态射的进一步结果，引理
\ref{more-morphisms-lemma-cover-by-geometrically-connected}，
每个对象 $U$（属于 $X_\etale$）都有 \'etale 覆盖
$\{U_i \to U\}$，使上一段的讨论适用于 $U_i$。引理得证。
\end{proof}

\begin{lemma}
\label{etale-cohomology-lemma-base-change-f-star}
设 $f : X \to S$ 是平坦的概形态射，并且对每个几何点
$\overline{x}$（属于 $X$），映射
$$
\mathcal{O}_{S, f(\overline{x})}^{sh}
\longrightarrow
\mathcal{O}_{X, \overline{x}}^{sh}
$$
都有几何连通纤维。则对每个概形的笛卡尔图
$$
\xymatrix{
X \ar[d]_f & Y \ar[l]^h \ar[d]^e \\
S & T \ar[l]_g
}
$$
，若 $g$ 拟紧拟分离，则有
$f^{-1}g_*\mathcal{F} = h_*e^{-1}\mathcal{F}$
，其中 $\mathcal{F}$ 是 $T_\etale$ 上的任意集合层。
\end{lemma}

\begin{proof}
只需在茎上检验等式；见定理 \ref{etale-cohomology-theorem-exactness-stalks}。
由定理 \ref{etale-cohomology-theorem-higher-direct-images}，有
$$
(h_*e^{-1}\mathcal{F})_{\overline{x}} =
\Gamma(\Spec(\mathcal{O}_{X, \overline{x}}^{sh}) \times_X Y, e^{-1}\mathcal{F})
$$
；类似地有
$$
(f^{-1}g_*^{-1}\mathcal{F})_{\overline{x}} =
(g_*^{-1}\mathcal{F})_{f(\overline{x})} =
\Gamma(\Spec(\mathcal{O}_{S, f(\overline{x})}^{sh}) \times_S T, \mathcal{F})
$$
。把引理 \ref{etale-cohomology-lemma-sections-upstairs} 应用于态射
$$
\Spec(\mathcal{O}_{X, \overline{x}}^{sh}) \times_X Y
\longrightarrow
\Spec(\mathcal{O}_{S, f(\overline{x})}^{sh}) \times_S T
$$
，可知这些集合相等；该态射是
$\Spec(\mathcal{O}_{X, \overline{x}}^{sh}) \to
\Spec(\mathcal{O}_{S, f(\overline{x})}^{sh})$
的基变换，因为 $Y = X \times_S T$。
\end{proof}







\section{高阶正像的基变换}
\label{etale-cohomology-section-base-change}

\noindent
本节是第 \ref{etale-cohomology-section-base-change-f-star} 节关于高阶正像的对应版本。
本节属于预备内容，初次阅读时应当略过。

\begin{remark}
\label{etale-cohomology-remark-base-change-holds}
设 $f : X \to S$ 是概形态射。设 $n$ 是整数。
我们称 $BC(f, n, q_0)$ 成立，如果对于每个交换图
$$
\xymatrix{
X \ar[d]_f & X' \ar[l] \ar[d]_{f'} & Y \ar[l]^h \ar[d]^e \\
S & S' \ar[l] & T \ar[l]_g
}
$$
其中 $X' = X \times_S S'$、$Y = X' \times_{S'} T$，且
$g$ 拟紧且拟分离，并且对于每个 Abel 层 $\mathcal{F}$，若它在
$T_\etale$ 上且被 $n$ 零化，则基变换映射
$$
(f')^{-1}R^qg_*\mathcal{F}
\longrightarrow
R^qh_*e^{-1}\mathcal{F}
$$
对 $q \leq q_0$ 都是同构。
\end{remark}

\begin{lemma}
\label{etale-cohomology-lemma-base-change-q-injective}
设 $f : X \to S$ 与 $n$ 如评注 \ref{etale-cohomology-remark-base-change-holds} 中所述，
并假设对某个 $q \geq 1$ 有 $BC(f, n, q - 1)$。则对于每个交换图
$$
\xymatrix{
X \ar[d]_f & X' \ar[l] \ar[d]_{f'} & Y \ar[l]^h \ar[d]^e \\
S & S' \ar[l] & T \ar[l]_g
}
$$
其中 $X' = X \times_S S'$、$Y = X' \times_{S'} T$，且
$g$ 拟紧且拟分离，并且对于每个 Abel 层 $\mathcal{F}$，若它在
$T_\etale$ 上且被 $n$ 零化，则有
\begin{enumerate}
\item 基变换映射
$(f')^{-1}R^qg_*\mathcal{F}\to R^qh_*e^{-1}\mathcal{F}$
是单射；
\item 如果 $\mathcal{F} \subset \mathcal{G}$，其中 $\mathcal{G}$
是 $T_\etale$ 上被 $n$ 零化的层，则
$$
\Coker\left(
(f')^{-1}R^qg_*\mathcal{F}\to R^qh_*e^{-1}\mathcal{F}
\right)
\subset
\Coker\left(
(f')^{-1}R^qg_*\mathcal{G}\to R^qh_*e^{-1}\mathcal{G}
\right)
$$
\item 如果 (2) 中的层 $\mathcal{G}$ 是 $\mathbf{Z}/n\mathbf{Z}$-模的内射层，则
$$
\Coker\left((f')^{-1}R^qg_*\mathcal{F}\to R^qh_*e^{-1}\mathcal{F} \right)
\subset R^qh_*e^{-1}\mathcal{G}
$$
\end{enumerate}
\end{lemma}

\begin{proof}
取短正合列
$0 \to \mathcal{F} \to \mathcal{I} \to \mathcal{Q} \to 0$
，其中 $\mathcal{I}$ 是 $\mathbf{Z}/n\mathbf{Z}$-模的内射层。
考虑由此诱导的图
$$
\xymatrix{
(f')^{-1}R^{q - 1}g_*\mathcal{I} \ar[d]_{\cong} \ar[r] &
(f')^{-1}R^{q - 1}g_*\mathcal{Q} \ar[d]_{\cong} \ar[r] &
(f')^{-1}R^qg_*\mathcal{F} \ar[d] \ar[r] &
0 \ar[d] \\
R^{q - 1}h_*e^{-1}\mathcal{I} \ar[r] &
R^{q - 1}h_*e^{-1}\mathcal{Q} \ar[r] &
R^qh_*e^{-1}\mathcal{F} \ar[r] &
R^qh_*e^{-1}\mathcal{I}
}
$$
，其两行正合。右上角为零，这是因为 $\mathcal{I}$ 是内射的。由
$BC(f, n, q - 1)$，左侧两个竖直箭头都是同构。因此 (1) 成立。
上图还表明
$$
\Coker\left(
(f')^{-1}R^qg_*\mathcal{F}\to R^qh_*e^{-1}\mathcal{F}
\right)
\subset
R^qh_*e^{-1}\mathcal{I}
$$
，故 (3) 成立。为证明 (2)，取
$\mathcal{F} \subset \mathcal{G} \subset \mathcal{I}$ 即可。
\end{proof}

\begin{lemma}
\label{etale-cohomology-lemma-base-change-q-integral-top}
设 $f : X \to S$ 与 $n$ 如评注 \ref{etale-cohomology-remark-base-change-holds} 中所述，
并假设对某个 $q \geq 1$ 有 $BC(f, n, q - 1)$。考虑交换图
$$
\vcenter{
\xymatrix{
X \ar[d]_f &
X' \ar[d]_{f'} \ar[l] &
Y \ar[l]^h \ar[d]^e &
Y' \ar[l]^{\pi'} \ar[d]^{e'} \\
S &
S' \ar[l] &
T \ar[l]_g &
T' \ar[l]_\pi
}
}
\quad\text{且}\quad
\vcenter{
\xymatrix{
X' \ar[d]_{f'} & & Y' \ar[ll]^{h' = h \circ \pi'} \ar[d]^{e'} \\
S' & & T' \ar[ll]_{g' = g \circ \pi}
}
}
$$
其中所有方块都是笛卡尔的，$g$ 拟紧且拟分离，而 $\pi$ 是整满射。
设 $\mathcal{F}$ 是 $T_\etale$ 上被 $n$ 零化的 Abel 层，并令
$\mathcal{F}' = \pi^{-1}\mathcal{F}$。如果基变换映射
$$
(f')^{-1}R^qg'_*\mathcal{F}' \longrightarrow R^qh'_*(e')^{-1}\mathcal{F}'
$$
是同构，则基变换映射
$(f')^{-1}R^qg_*\mathcal{F} \to R^qh_*e^{-1}\mathcal{F}$
也是同构。
\end{lemma}

\begin{proof}
注意到 $\mathcal{F} \to \pi_*\pi^{-1}\mathcal{F}'$ 是单射，
因为 $\pi$ 是满射（可在茎上检验）。因此，由引理
\ref{etale-cohomology-lemma-base-change-q-injective}，只需证明基变换映射
$$
(f')^{-1}R^qg_*\pi_*\mathcal{F}'
\longrightarrow
R^qh_*e^{-1}\pi_*\mathcal{F}'
$$
是同构。由假设，这一点成立，因为
$R^qg_*\pi_*\mathcal{F}' = R^qg'_*\mathcal{F}'$，
$e^{-1}\pi_*\mathcal{F}' =\pi'_*(e')^{-1}\mathcal{F}'$，并且
$R^qh_*\pi'_*(e')^{-1}\mathcal{F}' = R^qh'_*(e')^{-1}\mathcal{F}'$。
这些等式来自引理
\ref{etale-cohomology-lemma-integral-pushforward-commutes-with-base-change}、
\ref{etale-cohomology-lemma-what-integral} 以及相对 Leray 谱序列
（Cohomology on Sites，引理 \ref{sites-cohomology-lemma-relative-Leray}）。
\end{proof}

\begin{lemma}
\label{etale-cohomology-lemma-base-change-q-integral-bottom}
设 $f : X \to S$ 与 $n$ 如评注 \ref{etale-cohomology-remark-base-change-holds} 中所述，
并假设对某个 $q \geq 1$ 有 $BC(f, n, q - 1)$。考虑交换图
$$
\vcenter{
\xymatrix{
X \ar[d]_f &
X' \ar[d]_{f'} \ar[l] &
X'' \ar[l]^{\pi'} \ar[d]_{f''} &
Y \ar[l]^{h'} \ar[d]^e \\
S &
S' \ar[l] &
S'' \ar[l]_\pi &
T \ar[l]_{g'}
}
}
\quad\text{且}\quad
\vcenter{
\xymatrix{
X' \ar[d]_{f'} & & Y \ar[ll]^{h = h' \circ \pi'} \ar[d]^e \\
S' & & T \ar[ll]_{g = g' \circ \pi}
}
}
$$
其中所有方块都是笛卡尔的，$g'$ 拟紧且拟分离，而 $\pi$ 是整态射。
设 $\mathcal{F}$ 是 $T_\etale$ 上被 $n$ 零化的 Abel 层。如果基变换映射
$$
(f')^{-1}R^qg_*\mathcal{F} \longrightarrow R^qh_*e^{-1}\mathcal{F}
$$
是同构，则基变换映射
$(f'')^{-1}R^qg'_*\mathcal{F} \to R^qh'_*e^{-1}\mathcal{F}$
也是同构。
\end{lemma}

\begin{proof}
由于 $\pi$ 与 $\pi'$ 都是整态射，有 $R\pi_* = \pi_*$ 和
$R\pi'_* = \pi'_*$，见引理 \ref{etale-cohomology-lemma-what-integral}。
还有 $(f')^{-1}\pi_* = \pi'_*(f'')^{-1}$。因此
$\pi'_*(f'')^{-1}R^qg'_*\mathcal{F} = (f')^{-1}R^qg_*\mathcal{F}$
以及
$\pi'_*R^qh'_*e^{-1}\mathcal{F} = R^qh_*e^{-1}\mathcal{F}$.
所以，假设意味着对所考察的映射应用函子 $\pi'_*$ 后得到同构。
于是由引理 \ref{etale-cohomology-lemma-what-integral}，该映射本身是同构。
\end{proof}

\begin{lemma}
\label{etale-cohomology-lemma-formal-argument}
设 $T$ 是拟紧且拟分离的概形。设 $P$ 是 $T$ 上拟紧拟分离概形的一条性质。
假设
\begin{enumerate}
\item 如果 $T'' \to T'$ 是 $T$ 上拟紧拟分离概形之间的加厚，则
$P(T'')$ 成立当且仅当 $P(T')$ 成立。
\item 如果 $T' = \lim T_i$ 是 $T$ 上拟紧拟分离概形构成的逆系统之极限，
其转移态射为仿射态射，并且 $P(T_i)$ 对所有 $i$ 都成立，则 $P(T')$ 成立。
\item 如果 $Z \subset T'$ 是闭子概形，其补集 $V \subset T'$ 拟紧，
并且 $P(T')$ 成立，则 $P(V)$ 或 $P(Z)$ 至少有一个成立。
\end{enumerate}
那么，由 $P(T)$ 可推出：$P(\Spec(K))$ 成立，其中存在态射
$\Spec(K) \to T$，且 $K$ 是域。
\end{lemma}

\begin{proof}
考虑集合 $\mathfrak T$，其元素是闭子概形 $T' \subset T$，且满足 $P(T')$。
由假设 (2)，这个集合有极小元，记为 $T'$。由假设 (1)，$T'$ 是既约的。
设 $\eta \in T'$ 是 $T'$ 的某个不可约分支的泛点。则
$\eta = \Spec(K)$，其中 $K$ 是某个域，并且 $\eta = \lim V$，其中极限取遍
仿射开子概形 $V \subset T'$；这些开子概形都包含 $\eta$。由假设 (3) 以及 $T'$ 的极小性，
$P(V)$ 对所有这些 $V$ 都成立。因此由 (2)，$P(\eta)$ 成立，证明完毕。
\end{proof}

\begin{lemma}
\label{etale-cohomology-lemma-base-change-does-not-hold-pre}
设 $f : X \to S$ 与 $n$ 如评注 \ref{etale-cohomology-remark-base-change-holds} 中所述，
并假设对某个 $q \geq 1$，$BC(f, n, q - 1)$ 成立而
$BC(f, n, q)$ 不成立。那么存在交换图
$$
\xymatrix{
X \ar[d]_f & X' \ar[d]_{f'} \ar[l] & Y \ar[l]^h \ar[d]^e \\
S & S' \ar[l] & \Spec(K) \ar[l]_g
}
$$
，其中 $X' = X \times_S S'$、$Y = X' \times_{S'} \Spec(K)$，
$K$ 是域，并且存在 Abel 层 $\mathcal{F}$，它在 $\Spec(K)$ 上且被 $n$ 零化，使得
$(f')^{-1}R^qg_*\mathcal{F} \to R^qh_*e^{-1}\mathcal{F}$
不是同构。
\end{lemma}

\begin{proof}
取交换图
$$
\xymatrix{
X \ar[d]_f & X' \ar[l] \ar[d]_{f'} & Y \ar[l]^h \ar[d]^e \\
S & S' \ar[l] & T \ar[l]_g
}
$$
，其中 $X' = X \times_S S'$、$Y = X' \times_{S'} T$，
$g$ 拟紧且拟分离；再取 Abel 层 $\mathcal{F}$，它在 $T_\etale$ 上且被
$n$ 零化，使基变换映射
$(f')^{-1}R^qg_*\mathcal{F} \to R^qh_*e^{-1}\mathcal{F}$
不是同构。当然，可以并且确实用 $S'$ 的某个仿射开集替换 $S'$；
这使得 $T$ 拟紧且拟分离。由引理 \ref{etale-cohomology-lemma-base-change-q-injective}，映射
$(f')^{-1}R^qg_*\mathcal{F} \to R^qh_*e^{-1}\mathcal{F}$
是单射。取一个几何点 $\overline{x}$，它是 $X'$ 的几何点；再取元素
$\xi$，它属于 $(R^qh_*q^{-1}\mathcal{F})_{\overline{x}}$，
使它不属于下列映射的像：
$((f')^{-1}R^qg_*\mathcal{F})_{\overline{x}} \to
(R^qh_*e^{-1}\mathcal{F})_{\overline{x}}$.

\medskip\noindent
考虑态射 $\pi : T' \to T$，其中 $T'$ 拟紧且拟分离，并记
$\mathcal{F}' = \pi^{-1}\mathcal{F}$。记
$\pi' : Y' = Y \times_T T' \to Y$ 为 $\pi$ 的基变换，
$e' : Y' \to T'$ 为 $e$ 的基变换。图示为
$$
\vcenter{
\xymatrix{
X' \ar[d]_{f'} & Y \ar[l]^h \ar[d]^e & Y' \ar[l]^{\pi'} \ar[d]^{e'} \\
S' & T \ar[l]_g & T' \ar[l]_\pi
}
}
\quad\text{且}\quad
\vcenter{
\xymatrix{
X' \ar[d]_{f'} & & Y' \ar[ll]^{h' = h \circ \pi'} \ar[d]^{e'} \\
S' & & T' \ar[ll]_{g' = g \circ \pi}
}
}
$$
利用拉回映射，得到典范交换图
$$
\xymatrix{
(f')^{-1}R^qg_*\mathcal{F} \ar[r] \ar[d] &
(f')^{-1}R^qg'_*\mathcal{F}' \ar[d] \\
R^qh_*e^{-1}\mathcal{F} \ar[r] &
R^qh'_*(e')^{-1}\mathcal{F}'
}
$$
，这是 $X'$ 上 Abel 层的交换图。令 $P(T')$ 表示下列性质：
\begin{itemize}
\item 像 $\xi'$ 是 $\xi$ 在
$(Rh'_*(e')^{-1}\mathcal{F}')_{\overline{x}}$ 中的像，并且它
不属于下列映射的像：
$(f^{-1}R^qg'_*\mathcal{F}')_{\overline{x}} \to
(R^qh'_*(e')^{-1}\mathcal{F}')_{\overline{x}}$.
\end{itemize}
我们断言，性质 $P$ 满足引理 \ref{etale-cohomology-lemma-formal-argument} 的假设
(1)、(2)、(3)，从而即可证明本引理。

\medskip\noindent
性质 $P$ 满足引理 \ref{etale-cohomology-lemma-formal-argument} 的条件 (1)，
因为概形与该概形的加厚具有相同的 \'etale 拓扑。见命题
\ref{etale-cohomology-proposition-topological-invariance}。

\medskip\noindent
设 $I$ 是有向集，而 $T_i$ 是以 $I$ 为指标的逆系统，
其对象是 $T$ 上拟紧拟分离的概形，转移态射为仿射态射。
令 $T' = \lim T_i$。分别用 $\mathcal{F}'$ 与 $\mathcal{F}_i$
表示 $\mathcal{F}$ 到 $T'$ 与 $T_i$ 的拉回。考虑图
$$
\vcenter{
\xymatrix{
X \ar[d]_{f'} & Y \ar[l]^h \ar[d]^e & Y_i \ar[l]^{\pi_i'} \ar[d]^{e_i} \\
S & T \ar[l]_g & T_i \ar[l]_{\pi_i}
}
}
\quad\text{且}\quad
\vcenter{
\xymatrix{
X \ar[d]_{f'} & & Y_i \ar[ll]^{h_i = h \circ \pi_i'} \ar[d]^{e_i} \\
S & & T_i \ar[ll]_{g_i = g \circ \pi_i}
}
}
$$
，记号与上一段相同。显然，$\mathcal{F}'$（定义在 $T'$ 上）是各
$\mathcal{F}_i$ 到 $T'$ 的拉回之余极限，而 $(e')^{-1}\mathcal{F}'$
是各 $e_i^{-1}\mathcal{F}_i$ 到 $Y'$ 的拉回之余极限。
由引理 \ref{etale-cohomology-lemma-relative-colimit-general}，有
$$
R^qh'_*(e')^{-1}\mathcal{F}' = \colim R^qh_{i, *}e_i^{-1}\mathcal{F}_i
\quad\text{且}\quad
(f')^{-1}R^qg'_*\mathcal{F}' = \colim (f')^{-1}R^qg_{i, *}\mathcal{F}_i
$$
由此可见，如果 $P(T_i)$ 对所有 $i$ 都成立，则 $P(T')$ 成立。
所以性质 $P$ 满足引理 \ref{etale-cohomology-lemma-formal-argument} 的条件 (2)。

\medskip\noindent
最有意思的是引理 \ref{etale-cohomology-lemma-formal-argument} 的条件 (3)。
设 $T'$ 是 $T$ 上拟紧拟分离的概形，并且 $P(T')$ 成立。
设 $Z \subset T'$ 是闭子概形，其补集 $V \subset T'$ 拟紧。考虑图
$$
\xymatrix{
Y' \times_{T'} Z \ar[d]_{e_Z} \ar[r]_{i'} &
Y' \ar[d]_{e'} &
Y' \times_{T'} V \ar[l]^{j'} \ar[d]^{e_V} \\
Z \ar[r]^i &
T' &
V \ar[l]_j
}
$$
取单射 $j^{-1}\mathcal{F}' \to \mathcal{J}$，其中 $\mathcal{J}$
是 $\mathbf{Z}/n\mathbf{Z}$-模的内射层，定义在 $V$ 上。在茎上考察可知，映射
$$
\mathcal{F}' \to \mathcal{G} = j_*\mathcal{J} \oplus i_*i^{-1}\mathcal{F}'
$$
是单射。因此，$\xi'$ 映到下列群中的一个非零元素：
\begin{align*}
& \Coker\left(
((f')^{-1}R^qg'_*\mathcal{G})_{\overline{x}}
\to
(R^qh'_*(e')^{-1}\mathcal{G})_{\overline{x}}
\right) = \\
&
\Coker\left(
((f')^{-1}R^qg'_*j_*\mathcal{J})_{\overline{x}}
\to
(R^qh'_*(e')^{-1}j_*\mathcal{J})_{\overline{x}}
\right) \oplus \\
& \Coker\left(
((f')^{-1}R^qg'_*i_*i^{-1}\mathcal{F}')_{\overline{x}}
\to
(R^qh'_*(e')^{-1}i_*i^{-1}\mathcal{F}')_{\overline{x}}
\right)
\end{align*}
这里使用了引理 \ref{etale-cohomology-lemma-base-change-q-injective} 的 (2)。
如果 $\xi'$ 在第二个直和项中的像非零，则利用
$$
(f')^{-1}R^qg'_*i_*i^{-1}\mathcal{F}' =
(f')^{-1}R^q(g' \circ i)_*i^{-1}\mathcal{F}'
$$
（因为由命题 \ref{etale-cohomology-proposition-finite-higher-direct-image-zero}，
$Ri_* = i_*$）以及
$$
R^qh'_*(e')^{-1}i_*i^{-1}\mathcal{F} =
R^qh'_*i'_*e_Z^{-1}i^{-1}\mathcal{F} =
R^q(h' \circ i')_*e_Z^{-1}i^{-1}\mathcal{F}'
$$
（第一个等式由引理 \ref{etale-cohomology-lemma-finite-pushforward-commutes-with-base-change}，
第二个等式则来自命题 \ref{etale-cohomology-proposition-finite-higher-direct-image-zero}
所给出的 $Ri'_* = i'_*$），可知 $P(Z)$ 成立。
最后，假设 $\xi'$ 在第一个直和项中的像非零。有
$$
(e')^{-1}j_*\mathcal{J} = j'_*e_V^{-1}\mathcal{J}
\quad\text{且}\quad
R^aj'_*e_V^{-1}\mathcal{J} = 0, \quad a = 1, \ldots, q - 1
$$
；这是把 $BC(f, n, q - 1)$ 应用于图
$$
\xymatrix{
X \ar[d]_f & Y' \ar[l] \ar[d]_{e'} & Y \ar[l]^{j'} \ar[d]^{e_V} \\
S & T' \ar[l] & V \ar[l]_j
}
$$
并使用 $\mathcal{J}$ 的内射性得到的。
由 $h' \circ j'$ 的相对 Leray 谱序列
（Cohomology on Sites，引理 \ref{sites-cohomology-lemma-relative-Leray}），可知
$$
R^qh'_*(e')^{-1}j_*\mathcal{J} =
R^qh'_*j'_*e_V^{-1}\mathcal{J}
\longrightarrow
R^q(h' \circ j')_* e_V^{-1}\mathcal{J}
$$
是单射。因此，$\xi$ 映到
$(R^q(h' \circ j')_* e_V^{-1}\mathcal{J})_{\overline{x}}$
中的一个非零元素。把引理 \ref{etale-cohomology-lemma-base-change-q-injective} 的 (3)
应用于单射 $j^{-1}\mathcal{F}' \to \mathcal{J}$，即得 $P(V)$ 成立。
\end{proof}

\begin{lemma}
\label{etale-cohomology-lemma-base-change-does-not-hold}
设 $f : X \to S$ 与 $n$ 如评注 \ref{etale-cohomology-remark-base-change-holds} 中所述，
并假设对某个 $q \geq 1$，$BC(f, n, q - 1)$ 成立而
$BC(f, n, q)$ 不成立。那么存在交换图
$$
\xymatrix{
X \ar[d]_f & X' \ar[d] \ar[l] & Y \ar[l]^h \ar[d] \\
S & S' \ar[l] & \Spec(K) \ar[l]
}
$$
，其中两个方块都是笛卡尔的，并且
\begin{enumerate}
\item $S'$ 是仿射、整且正规的概形，其函数域代数闭；
\item $K$ 代数闭，并且 $\Spec(K) \to S'$ 是支配态射
（换言之，$K$ 是 $S'$ 的函数域的一个扩域）。
\end{enumerate}
此外，存在整数 $d | n$，使得 $R^qh_*(\mathbf{Z}/d\mathbf{Z})$ 非零。
\end{lemma}

\noindent
反过来，引理中 $R^qh_*(\mathbf{Z}/d\mathbf{Z})$ 的非消失蕴含
$BC(f, n, q)$ 不成立，因为引理
\ref{etale-cohomology-lemma-Rf-star-zero-normal-with-alg-closed-function-field} 表明
$R^q(\Spec(K) \to S')_*\mathbf{Z}/d\mathbf{Z} = 0$。

\begin{proof}
首先取引理 \ref{etale-cohomology-lemma-base-change-does-not-hold-pre} 中的图与
$\mathcal{F}$。可以并且确实假设 $S'$ 是仿射的
（这是显然的；若有疑问，可参见该引理的证明）。
由引理 \ref{etale-cohomology-lemma-base-change-q-integral-top}，可以假设 $K$ 代数闭。
此时 $\mathcal{F}$ 对应于一个 $\mathbf{Z}/n\mathbf{Z}$-模。
这样的模是若干 $\mathbf{Z}/d\mathbf{Z}$ 的直和，其中 $d | n$ 可变，
因此可以假设 $\mathcal{F}$ 是常值层，其值为 $\mathbf{Z}/d\mathbf{Z}$。
由引理 \ref{etale-cohomology-lemma-base-change-q-integral-bottom}，可以把 $S'$ 替换为
$S'$ 在 $\Spec(K)$ 中的正规化，证明完毕。
\end{proof}










\section{光滑基变换}
\label{etale-cohomology-section-smooth-base-change}

\noindent
本节证明光滑基变换定理。

\begin{lemma}
\label{etale-cohomology-lemma-smooth-base-change-fields}
设 $K/k$ 是域扩张。设 $X$ 是 $k$ 上的光滑仿射曲线，
并有有理点 $x \in X(k)$。设 $\mathcal{F}$ 是 $\Spec(K)$ 上的 Abel 层，
它被整数 $n$ 零化，而该整数在 $k$ 中可逆。设 $q > 0$，且
$$
\xi \in H^q(X_K, (X_K \to \Spec(K))^{-1}\mathcal{F})
$$
则存在
\begin{enumerate}
\item 有限扩张 $K'/K$ 与 $k'/k$，其中 $k' \subset K'$；
\item 有限 \'etale Galois 覆盖 $Z \to X_{k'}$，其 Galois 群为 $G$。
\end{enumerate}
使得 $G$ 的阶整除 $n$ 的某个幂，$Z \to X_{k'}$ 在 $x_{k'}$ 上分裂，
并且 $\xi$ 在 $H^q(Z_{K'}, (Z_{K'} \to \Spec(K))^{-1}\mathcal{F})$ 中变为零。
\end{lemma}

\begin{proof}
当 $q > 1$ 时，已知 $\xi$ 在
$H^q(X_{\overline{K}}, (X_{\overline{K}} \to \Spec(K))^{-1}\mathcal{F})$
中变为零（定理 \ref{etale-cohomology-theorem-vanishing-affine-curves}）。
由引理 \ref{etale-cohomology-lemma-directed-colimit-cohomology}，这意味着存在有限扩张
$K'/K$，使得 $\xi$ 在
$H^q(X_{K'}, (X_{K'} \to \Spec(K))^{-1}\mathcal{F})$ 中变为零。
因此在这种情形下可以取 $k' = k$ 以及 $Z = X$。

\medskip\noindent
假设 $q = 1$。回顾，$\mathcal{F}$ 对应于离散模 $M$，它带有连续
$\text{Gal}_K$-作用，见引理
\ref{etale-cohomology-lemma-equivalence-abelian-sheaves-point}。由于 $M$ 是 $n$-挠的，
它是有限 $\text{Gal}_K$-稳定子群的并。因此可归约到 $M$ 是被 $n$
零化的有限 Abel 群之情形，见引理 \ref{etale-cohomology-lemma-colimit}。以一个有限扩张替换
$K$ 后，可以假设 $\text{Gal}_K$ 在 $M$ 上的作用平凡。
所以可以假设 $\mathcal{F} = \underline{M}$ 是常值层，其值为有限 Abel 群
$M$，且该群被 $n$ 零化。

\medskip\noindent
可以把 $M$ 写成循环群的直和。任意两个有限 \'etale Galois 覆盖，
若其 Galois 群的阶在 $k$ 中可逆，就都可被第三个 Galois 群阶也在
$k$ 中可逆的覆盖控制（Fundamental Groups，第
\ref{pione-section-finite-etale-under-galois} 节）。因此只需在
$M = \mathbf{Z}/d\mathbf{Z}$ 且 $d | n$ 时证明本引理。

\medskip\noindent
假设 $M = \mathbf{Z}/d\mathbf{Z}$，其中 $d | n$。
在这种情形下，$\overline{\xi} = \xi|_{X_{\overline{K}}}$ 是下列群的元素：
$$
H^1(X_{\overline{k}}, \mathbf{Z}/d\mathbf{Z}) =
H^1(X_{\overline{K}}, \mathbf{Z}/d\mathbf{Z})
$$
见定理 \ref{etale-cohomology-theorem-vanishing-affine-curves}。这个群分类
$\mathbf{Z}/d\mathbf{Z}$-挠子，见 Cohomology on Sites，引理
\ref{sites-cohomology-lemma-torsors-h1}。与 $\overline{\xi}$ 对应的挠子
（看作 $X_{\overline{k}, \etale}$ 上的层）又给出有限 \'etale 态射
$T \to X_{\overline{k}}$，它带有 $\mathbf{Z}/d\mathbf{Z}$ 的作用，且该作用
在 $T$ 位于 $x_{\overline{k}}$ 上方的纤维上可迁，见引理
\ref{etale-cohomology-lemma-characterize-finite-locally-constant}。取一个连通分支
$T' \subset T$（如果 $\overline{\xi}$ 的阶为 $d$，则 $T$ 已经连通）。
于是 $T' \to X_{\overline{k}}$ 是有限 \'etale Galois 覆盖，其 Galois 群是子群
$G \subset \mathbf{Z}/d\mathbf{Z}$（略去一个小细节）。此外，元素
$\overline{\xi}$ 在映射 $H^1(X_{\overline{k}}, \mathbf{Z}/d\mathbf{Z}) \to
H^1(T', \mathbf{Z}/d\mathbf{Z})$ 下映为零，因为这是 $T$ 的定义性质之一。

\medskip\noindent
接下来使用极限论证，选取有限扩张 $k'/k$，它包含于 $\overline{k}$，使得
$T' \to X_{\overline{k}}$ 下降为有限 \'etale Galois 覆盖
$Z \to X_{k'}$，其 Galois 群为 $G$。见 Limits，引理
\ref{limits-lemma-descend-finite-presentation}、
\ref{limits-lemma-descend-finite-finite-presentation} 和
\ref{limits-lemma-descend-etale}。增大 $k'$ 后，可以假设 $Z$ 在 $x_{k'}$
上分裂。按照构造，$\xi$ 在
$H^1(Z_{\overline{K}}, \mathbf{Z}/d\mathbf{Z})$ 中的像为零。
因此由引理 \ref{etale-cohomology-lemma-directed-colimit-cohomology}，可找到有限子扩张
$\overline{K}/K'/K$，它包含 $k'$，使得 $\xi$ 在 $H^1(Z_{K'}, \mathbf{Z}/d\mathbf{Z})$
中变为零，证明完毕。
\end{proof}

\begin{theorem}[光滑基变换]
\label{etale-cohomology-theorem-smooth-base-change}
考虑概形的笛卡尔图
$$
\xymatrix{
X \ar[d]_f & Y \ar[l]^h \ar[d]^e \\
S & T \ar[l]_g
}
$$
，其中 $f$ 光滑，而 $g$ 拟紧且拟分离。那么
$$
f^{-1}R^qg_*\mathcal{F} = R^qh_*e^{-1}\mathcal{F}
$$
对任意 $q$ 以及任意 Abel 层 $\mathcal{F}$ 都成立；这里该层定义在
$T_\etale$ 上，并且它在几何点处的所有茎都是挠群，其阶在 $S$ 上可逆。
\end{theorem}

\begin{proof}[光滑基变换的第一种证明]
这个证明很长，但比下面给出的第二种证明更直接（使用的一般理论较少）。

\medskip\noindent
本定理在 $X_\etale$ 上是局部的。更准确地，假设有 \'etale 态射
$U \to X$，使得 $U \to S$ 分解为 $U \to V \to S$，其中
$V \to S$ 是 \'etale 的。于是可以考虑笛卡尔方块
$$
\xymatrix{
U \ar[d]_{f'} & U \times_X Y \ar[l]^{h'} \ar[d]^{e'} \\
V & V \times_S T \ar[l]_{g'}
}
$$
并令 $\mathcal{F}' = \mathcal{F}|_{V \times_S T}$。于是
$f^{-1}R^qg_*\mathcal{F}|_U = (f')^{-1}R^qg'_*\mathcal{F}'$，且
$R^qh_*e^{-1}\mathcal{F}|_U = R^qh'_*(e')^{-1}\mathcal{F}'$；
这是网站态射与局部化相容的结果，见 Sites，引理
\ref{sites-lemma-localize-morphism-strong} 以及 Cohomology on Sites，引理
\ref{sites-cohomology-lemma-restrict-direct-image-open}。因此，只需构造
$X$ 的一个 \'etale 覆盖，其成员为 $U \to X$，并构造如上的分解
$U \to V \to S$，使定理对带有 $f'$、$h'$、$g'$、$e'$ 的图成立。

\medskip\noindent
由光滑态射的局部结构，见 Morphisms，引理
\ref{morphisms-lemma-smooth-etale-over-affine-space}，可以假设 $X$ 与
$S$ 都是仿射的，并且 $X \to S$ 经过 \'etale 态射
$X \to \mathbf{A}^d_S$ 分解。如果有一叠笛卡尔图
$$
\xymatrix{
W \ar[d]_i & Z \ar[l]^j \ar[d]^k \\
X \ar[d]_f & Y \ar[l]^h \ar[d]^e \\
S & T \ar[l]_g
}
$$
，并且定理对下面和上面的方块成立，那么它也对外侧矩形成立；这是形式的。
把 $X \to S$ 写成复合
$$
X \to \mathbf{A}^{d - 1}_S \to \mathbf{A}^{d - 2}_S \to \ldots \to
\mathbf{A}^1_S \to S
$$
，可知只需在 $X$ 与 $S$ 都仿射且 $X \to S$ 的相对维数为 $1$ 时证明定理。

\medskip\noindent
对每个 $n \geq 1$，若它在 $S$ 上可逆，则令 $\mathcal{F}[n]$ 为
$\mathcal{F}$ 的被 $n$ 零化的截面所成的子层。由假设，
$\mathcal{F} = \colim \mathcal{F}[n]$，这是因为 $\mathcal{F}$ 的茎具有所述性质。函子 $e^{-1}$ 与 $f^{-1}$
作为左伴随与余极限交换。由引理 \ref{etale-cohomology-lemma-relative-colimit}，函子
$R^qh_*$ 与 $R^qg_*$ 和滤过余极限交换。因此只需对 $\mathcal{F}[n]$
证明定理。从现在起固定整数 $n$，使用 $\mathbf{Z}/n\mathbf{Z}$-模层，
并假设 $S$ 是 $\Spec(\mathbf{Z}[1/n])$ 上的概形。

\medskip\noindent
接下来归约到 $T$ 仿射的情形。由于 $g$ 拟紧且拟分离而 $S$ 仿射，
概形 $T$ 拟紧且拟分离。因此可以使用 Cohomology of Schemes，引理
\ref{coherent-lemma-induction-principle} 的归纳原理。故只需证明：如果
$T = W \cup W'$ 是开覆盖，并且定理对下列方块成立
$$
\xymatrix{
X \ar[d] & e^{-1}(W) \ar[l]^i \ar[d] \\
S & W \ar[l]_a
}
\quad
\xymatrix{
X \ar[d] & e^{-1}(W') \ar[l]^j \ar[d] \\
S & W' \ar[l]_b
}
\quad
\xymatrix{
X \ar[d] & e^{-1}(W \cap W') \ar[l]^-k \ar[d] \\
S & W \cap W' \ar[l]_c
}
$$
，则定理对原图成立。为此考虑图
$$
\xymatrix{
f^{-1}R^{q - 1}c_*\mathcal{F}|_{W \cap W'} \ar[d]^{\cong} \ar[r] &
f^{-1}R^qg_*\mathcal{F} \ar[d] \ar[r] &
f^{-1}R^qa_*\mathcal{F}|_W \oplus
f^{-1}R^qb_*\mathcal{F}|_{W'} \ar[d]_{\cong} \\
R^qk_*e^{-1}\mathcal{F}|_{e^{-1}(W \cap W')} \ar[r] &
R^qh_*e^{-1}\mathcal{F} \ar[r] &
R^qi_*e^{-1}\mathcal{F}|_{e^{-1}(W)} \oplus
R^qj_*e^{-1}\mathcal{F}|_{e^{-1}(W')}
}
$$
，其两行是引理 \ref{etale-cohomology-lemma-relative-mayer-vietoris} 的长正合列。
于是 $5$-引理给出所需结论。

\medskip\noindent
综上，可以假设 $S$、$X$、$T$ 与 $Y$ 都仿射，$\mathcal{F}$ 是
$n$-挠的，$X \to S$ 光滑且相对维数为 $1$，并且 $S$ 是
$\mathbf{Z}[1/n]$ 上的概形。对 $q$ 作归纳来证明定理。初始情形
$q = 0$ 由引理 \ref{etale-cohomology-lemma-fppf-reduced-fibres-base-change-f-star} 处理。
假设 $q > 0$，并且定理对所有更小次数都成立。取短正合列
$0 \to \mathcal{F} \to \mathcal{I} \to \mathcal{Q} \to 0$
，其中 $\mathcal{I}$ 是 $\mathbf{Z}/n\mathbf{Z}$-模的内射层。考虑诱导图
$$
\xymatrix{
f^{-1}R^{q - 1}g_*\mathcal{I} \ar[d]_{\cong} \ar[r] &
f^{-1}R^{q - 1}g_*\mathcal{Q} \ar[d]_{\cong} \ar[r] &
f^{-1}R^qg_*\mathcal{F} \ar[d] \ar[r] &
0 \ar[d] \\
R^{q - 1}h_*e^{-1}\mathcal{I} \ar[r] &
R^{q - 1}h_*e^{-1}\mathcal{Q} \ar[r] &
R^qh_*e^{-1}\mathcal{F} \ar[r] &
R^qh_*e^{-1}\mathcal{I}
}
$$
，其两行正合。右上角为零，因为 $\mathcal{I}$ 是内射的。
由归纳假设，左侧两个竖直箭头都是同构。因此只需证明
$R^qh_*e^{-1}\mathcal{I} = 0$。

\medskip\noindent
写成 $S = \Spec(A)$ 与 $T = \Spec(B)$，并设态射 $T \to S$
由环映射 $A \to B$ 给出。可以把
$A \to B = \colim_{i \in I} (A_i \to B_i)$ 写成
$\mathbf{Z}[1/n]$ 上有限型环映射的滤过余极限（见 Algebra，引理
\ref{algebra-lemma-limit-no-condition}）。对 $i \in I$，令
$S_i = \Spec(A_i)$ 和 $T_i = \Spec(B_i)$。当 $i$ 足够大时，可以找到
光滑态射 $X_i \to S_i$，其相对维数为 $1$，并使得
$X = X_i \times_{S_i} S$，见 Limits，引理
\ref{limits-lemma-descend-finite-presentation}、
\ref{limits-lemma-descend-smooth} 和
\ref{limits-lemma-descend-dimension-d}。令 $Y_i = X_i \times_{S_i} T_i$，得到方块
$$
\xymatrix{
X_i \ar[d]_{f_i} & Y_i \ar[l]^{h_i} \ar[d]^{e_i} \\
S_i & T_i \ar[l]_{g_i}
}
$$
注意，$\mathcal{I}_i = (T \to T_i)_*\mathcal{I}$ 是
$\mathbf{Z}/n\mathbf{Z}$-模的内射层，定义在 $T_i$ 上，见 Cohomology on Sites，引理
\ref{sites-cohomology-lemma-pushforward-injective-flat}。由引理
\ref{etale-cohomology-lemma-linus-hamann}，有
$\mathcal{I} = \colim (T \to T_i)^{-1}\mathcal{I}_i$。沿 $e$ 拉回得到
$e^{-1}\mathcal{I} = \colim (Y \to Y_i)^{-1}e_i^{-1}\mathcal{I}_i$。
把引理 \ref{etale-cohomology-lemma-relative-colimit-general} 应用于态射系统
$Y_i \to X_i$（其极限为 $Y \to X$），得到
$$
R^qh_*e^{-1}\mathcal{I} =
\colim (X \to X_i)^{-1} R^qh_{i, *} e_i^{-1}\mathcal{I}_i
$$
这把问题归约到 $T$ 与 $S$ 都是 $\mathbf{Z}[1/n]$ 上有限型仿射概形的情形。

\medskip\noindent
综上，有整数 $q \geq 1$，使定理在次数 $< q$ 时成立；概形 $S$ 与 $T$
都是 $\mathbf{Z}[1/n]$ 上有限型的仿射概形；$X \to S$ 光滑且相对维数为
$1$，其中 $X$ 仿射；$\mathcal{I}$ 是 $\mathbf{Z}/n\mathbf{Z}$-模的内射层；
需要证明 $R^qh_*e^{-1}\mathcal{I} = 0$。这将通过对 $\dim(T)$ 作归纳来完成。

\medskip\noindent
初始情形是 $T = \emptyset$，即 $\dim(T) < 0$。若不喜欢这种约定，
也可以把 $\dim(T) = 0$ 作为初始情形。此时 $T \to S$ 是有限的
（事实上 $T \to \Spec(\mathbf{Z}[1/n])$ 也是有限的，因为目标是 Jacobson 的；
细节略去），所以 $h$ 也有限，从而其高阶正像消失（见下文的参考文献）。

\medskip\noindent
假设 $\dim(T) = d \geq 0$，并且对所有 $T$ 维数更低的情形已知结论。
取 $U$，它在 $X$ 上仿射且 \'etale；再取截面 $\xi$，它是
$R^qh_*q^{-1}\mathcal{I}$ 在 $U$ 上的截面。需要证明 $\xi$ 为零。当然，
可以替换 $X$，所用概形为 $U$（并相应地替换 $Y$，所用概形为
$U \times_X Y$），并假设
$\xi \in H^0(X, R^qh_*e^{-1}\mathcal{I})$。此外，由于
$R^qh_*e^{-1}\mathcal{I}$ 是层，只需证明 $\xi$ 局部为零，这里的局部性取在 $X$ 上。
因此，可以用一个 \'etale 覆盖的成员替换 $X$。特别地，利用引理
\ref{etale-cohomology-lemma-higher-direct-images}，可以假设 $\xi$ 是某个元素
$\tilde \xi \in H^q(Y, e^{-1}\mathcal{I})$ 的像。用 $\tilde \xi$ 表述，
任务是证明 $\tilde \xi$ 在 $H^q(U_i \times_X Y, e^{-1}\mathcal{I})$
中变为零，其中 $\{U_i \to X\}$ 是某个 \'etale 覆盖。

\medskip\noindent
由 More on Morphisms，引理 \ref{more-morphisms-lemma-cover-smooth-by-special}，
可以假设 $X \to S$ 分解为 $X \to V \to S$，其中 $V \to S$ 是
\'etale 的，而 $X \to V$ 是仿射概形之间相对维数为 $1$ 的光滑态射，
它有截面且纤维几何连通。注意
$\dim(V \times_S T) \leq \dim(T) = d$，例如可用 More on Algebra，引理
\ref{more-algebra-lemma-dimension-etale-extension}。因此，可以替换 $S$，
所用概形为 $V$；并替换 $T$，所用概形为 $V \times_S T$
（与证明第一段中的讨论完全相同）。
所以可以假设 $X \to S$ 光滑且相对维数为 $1$，纤维几何连通，并有截面
$\sigma : S \to X$。

\medskip\noindent
设 $\pi : T' \to T$ 是有限满射。下面将用到
$\dim(T') \leq \dim(T) = d$，见 Algebra，引理
\ref{algebra-lemma-integral-dim-up}。选取单射 $\pi^{-1}\mathcal{I} \to \mathcal{I}'$，
其中目标是 $\mathbf{Z}/n\mathbf{Z}$-模的内射层。于是
$\mathcal{I} \to \pi_*\mathcal{I}'$ 是单射，因而有分裂（因为
$\mathcal{I}$ 是 $\mathbf{Z}/n\mathbf{Z}$-模的内射层）。记
$\pi' : Y' = Y \times_T T' \to Y$ 为 $\pi$ 的基变换，
$e' : Y' \to T'$ 为 $e$ 的基变换。图示为
$$
\xymatrix{
X \ar[d]_f & Y \ar[l]^h \ar[d]^e & Y' \ar[l]^{\pi'} \ar[d]^{e'} \\
S & T \ar[l]_g & T' \ar[l]_\pi
}
$$
由命题 \ref{etale-cohomology-proposition-finite-higher-direct-image-zero} 和引理
\ref{etale-cohomology-lemma-finite-pushforward-commutes-with-base-change}，有
$R\pi'_*(e')^{-1}\mathcal{I}' = e^{-1}\pi_*\mathcal{I}'$.
因此由 Leray 谱序列（Cohomology on Sites，引理
\ref{sites-cohomology-lemma-Leray}），有
$$
H^q(Y', (e')^{-1}\mathcal{I}') =
H^q(Y, e^{-1}\pi_*\mathcal{I}') \supset H^q(Y, e^{-1}\mathcal{I})
$$
，而且在沿任意 \'etale 态射 $U \to X$ 作基变换后这仍然成立。
因此可以替换 $T$，所用概形为 $T'$；替换 $\mathcal{I}$，所用层为
$\mathcal{I}'$；并用 $\tilde \xi$ 在
$H^q(Y', (e')^{-1}\mathcal{I}')$ 中的像替换它。

\medskip\noindent
假设有分解 $T \to S' \to S$，其中 $\pi : S' \to S$ 有限。
令 $X' = S' \times_S X$，可以考虑诱导图
$$
\xymatrix{
X \ar[d]_f & X' \ar[l]^{\pi'} \ar[d]^{f'} & Y \ar[l]^{h'} \ar[d]^e \\
S & S' \ar[l]_\pi & T \ar[l]_g
}
$$
由于 $\pi'$ 的高阶正像消失，由 Leray 谱序列可知
$R^qh_*e^{-1}\mathcal{I} = \pi'_*R^qh'_*e^{-1}\mathcal{I}$。因此
$H^0(X, R^qh_*e^{-1}\mathcal{I}) = H^0(X', R^qh'_*e^{-1}\mathcal{I})$。
所以 $\xi$ 为零当且仅当 $R^qh'_*e^{-1}\mathcal{I}$ 的相应截面为零\footnote{也可用另一种方式理解这一步。也就是说，需要证明存在 \'etale 覆盖
$\{U_i \to X\}$，使得 $\tilde \xi$ 在
$H^q(U_i \times_X Y, e^{-1}\mathcal{I})$ 中变为零。然而，如果能证明存在
\'etale 覆盖 $\{U'_j \to X'\}$，使得 $\tilde \xi$ 在
$H^q(U'_i \times_{X'} Y, e^{-1}\mathcal{I})$ 中变为零，那么由
$X' \to X$ 的性质 (B)（引理 \ref{etale-cohomology-lemma-finite-B}），存在 \'etale 覆盖
$\{U_i \to X\}$，使 $U_i \times_X X'$ 是若干 $X'$ 上概形的不交并，
其中每个概形都经过 $U'_j$ 分解（对某个 $j$）。于是如所需，$\tilde \xi$ 在
$H^q(U_i \times_X Y, e^{-1}\mathcal{I})$ 中变为零。}。
因此可以替换 $S$，所用概形为 $S'$；并替换 $X$，所用概形为 $X'$。注意
$\sigma : S \to X$ 基变换为 $\sigma' : S' \to X'$；所以替换后，
$X \to S$ 仍然有截面 $\sigma$，且纤维几何连通。

\medskip\noindent
我们将使用 $S$ 与 $T$ 都是 Nagata 概形这一事实，见 Algebra，命题
\ref{algebra-proposition-ubiquity-nagata}；这保证各种正规化都是有限的，
见 Morphisms，引理 \ref{morphisms-lemma-nagata-normalization-finite} 和
\ref{morphisms-lemma-nagata-normalization}。特别地，可以先用正规化替换
$T$，再替换 $S$，所用概形为 $S$ 在 $T$ 中的正规化。于是 $T \to S$ 是整正规概形
之间支配态射的不交并，见 Morphisms，引理
\ref{morphisms-lemma-normal-normalization}。显然可以逐个连通分支论证，
故可假设 $T \to S$ 是整正规概形之间的支配态射。

\medskip\noindent
设 $s \in S$ 与 $t \in T$ 为泛点。由引理
\ref{etale-cohomology-lemma-smooth-base-change-fields}，存在有限域扩张 $K/\kappa(t)$ 和
$k/\kappa(s)$，其中 $k$ 包含于 $K$；还存在有限 \'etale Galois 覆盖
$Z \to X_k$，其 Galois 群为 $G$，群阶整除 $n$ 的某个幂，且该覆盖在
$\sigma(\Spec(k))$ 上分裂，使得 $\tilde \xi$ 在
$H^q(Z_K, e^{-1}\mathcal{I}|_{Z_K})$ 中的像为零。令 $T' \to T$ 为
$T$ 在 $\Spec(K)$ 中的正规化，并令 $S' \to S$ 为 $S$ 在
$\Spec(k)$ 中的正规化。于是得到交换图
$$
\xymatrix{
S' \ar[d] & T' \ar[l] \ar[d] \\
S & T \ar[l]
}
$$
，其竖直箭头有限。由上述论证，可以并且确实替换 $S$ 与 $T$，
所用概形分别为 $S'$ 与 $T'$（并相应地替换 $X$，所用概形为
$X \times_S S'$；替换 $Y$，所用概形为 $Y \times_T T'$）。替换后，
得到有限 \'etale Galois 覆盖 $Z \to X_s$；它是 $X \to S$ 的泛纤维上的覆盖，
其 Galois 群 $G$ 的阶整除 $n$ 的某个幂，
该覆盖在 $\sigma(s)$ 上分裂，并且 $\tilde \xi$ 在
$H^q(Z_t, (Z_t \to Y)^{-1}e^{-1}\mathcal{I})$ 中映为零。这里
$Z_t = Z \times_S t = Z \times_s t = Z \times_{X_s} Y_t$。由于 $n$ 在
$S$ 上可逆，由 Fundamental Groups，引理 \ref{pione-lemma-extend-covering}，
可以找到有限 \'etale 态射 $U \to X$，其到 $X_s$ 的限制为 $Z$。

\medskip\noindent
此时替换 $X$，所用概形为 $U$；并替换 $Y$，所用概形为 $U \times_X Y$。替换后，
$X \to S$ 的纤维可能不再几何连通（仍有截面，但不再使用它）；
所得好处是，替换后 $\tilde \xi$ 在 $H^q(Y_t, e^{-1}\mathcal{I})$ 中映为零，
即 $\tilde \xi$ 在 $Y \to T$ 的泛纤维上的限制为零。

\medskip\noindent
回顾，$t$ 是 $T$ 的函数域之谱；也就是说，作为概形，$t$ 是 $T$ 的
非空仿射开子概形之极限。由引理 \ref{etale-cohomology-lemma-directed-colimit-cohomology}，
存在非空开子概形 $V \subset T$，使得 $\tilde \xi$ 在
$H^q(Y \times_T V, e^{-1}\mathcal{I}|_{Y \times_T V})$ 中映为零。

\medskip\noindent
记 $Z = T \setminus V$。考虑图
$$
\xymatrix{
Y \times_T Z \ar[d]_{e_Z} \ar[r]_{i'} &
Y \ar[d]_e &
Y \times_T V \ar[l]^{j'} \ar[d]^{e_V} \\
Z \ar[r]^i &
T &
V \ar[l]_j
}
$$
取单射 $i^{-1}\mathcal{I} \to \mathcal{I}'$，其目标为
$\mathbf{Z}/n\mathbf{Z}$-模的内射层，定义在 $Z$ 上。在茎上考察可知，映射
$$
\mathcal{I} \to j_*\mathcal{I}|_V \oplus i_*\mathcal{I}'
$$
是单射，并且由于 $\mathcal{I}$ 是 $\mathbf{Z}/n\mathbf{Z}$-模的内射层，
该单射分裂。因此只需证明 $\tilde \xi$ 在
$$
H^q(Y, e^{-1}j_*\mathcal{I}|_V) \oplus
H^q(Y, e^{-1}i_*\mathcal{I}')
$$
中映为零，至少在用一个 \'etale 覆盖的成员替换 $X$ 后如此。注意
$$
e^{-1}j_*\mathcal{I}|_V = j'_*e_V^{-1}\mathcal{I}|_V,\quad
e^{-1}i_*\mathcal{I}' = i'_*e_Z^{-1}\mathcal{I}'
$$
由关于 $q$ 的归纳假设，有
$$
R^aj'_*e_V^{-1}\mathcal{I}|_V = 0, \quad a = 1, \ldots, q - 1
$$
由 $j'$ 的 Leray 谱序列及上述消失，可知
$$
H^q(Y, j'_*(e_V^{-1}\mathcal{I}|_V))
\longrightarrow
H^q(Y \times_T V, e_V^{-1}\mathcal{I}_V) =
H^q(Y \times_T V, e^{-1}\mathcal{I}|_{Y \times_T V})
$$
是单射。因此，$\tilde \xi$ 在上述第一个直和项中的像为零，因为已知
$\tilde \xi$ 在 $H^q(Y \times_T V, e^{-1}\mathcal{I}|_{Y \times_T V})$
中消失。由于 $\dim(Z) < \dim(T) = d$，由归纳，$\tilde \xi$
在第二个直和项
$$
H^q(Y, e^{-1}i_*\mathcal{I}') =
H^q(Y, i'_*e_Z^{-1}\mathcal{I}') =
H^q(Y \times_T Z, e_Z^{-1}\mathcal{I}')
$$
中的像在用一个合适的 \'etale 覆盖的成员替换 $X$ 后消失。
这就完成了光滑基变换定理的证明。
\end{proof}

\begin{proof}[光滑基变换的第二种证明]
这个证明与较长的第一种证明相同；它较短，仅仅因为我们已把所用论证
分别整理成了若干引理。

\medskip\noindent
$q = 0$ 的情形是引理
\ref{etale-cohomology-lemma-fppf-reduced-fibres-base-change-f-star}。因此可以假设
$q > 0$，并且结论对所有更小次数都成立。

\medskip\noindent
对每个 $n \geq 1$，若它在 $S$ 上可逆，则令 $\mathcal{F}[n]$ 为
$\mathcal{F}$ 的被 $n$ 零化的截面所成的子层。由假设，
$\mathcal{F} = \colim \mathcal{F}[n]$，这是因为 $\mathcal{F}$ 的茎具有所述性质。函子 $e^{-1}$ 与 $f^{-1}$
作为左伴随与余极限交换。由引理 \ref{etale-cohomology-lemma-relative-colimit}，函子
$R^qh_*$ 与 $R^qg_*$ 和滤过余极限交换。因此只需对 $\mathcal{F}[n]$
证明定理。从现在起固定整数 $n$，并假设它在 $S$ 上可逆；使用
$\mathbf{Z}/n\mathbf{Z}$-模层。

\medskip\noindent
由引理 \ref{etale-cohomology-lemma-base-change-local}，问题在 $X$ 与 $S$ 上都是 \'etale
局部的。由光滑态射的局部结构，见 Morphisms，引理
\ref{morphisms-lemma-smooth-etale-over-affine-space}，可以假设 $X$ 与 $S$
都仿射，并且 $X \to S$ 经过 \'etale 态射
$X \to \mathbf{A}^d_S$ 分解。把 $X \to S$ 写成复合
$$
X \to \mathbf{A}^{d - 1}_S \to \mathbf{A}^{d - 2}_S \to \ldots \to
\mathbf{A}^1_S \to S
$$
，由引理 \ref{etale-cohomology-lemma-base-change-compose} 可知，只需在 $X$ 与 $S$
都仿射且 $X \to S$ 的相对维数为 $1$ 时证明定理。

\medskip\noindent
由引理 \ref{etale-cohomology-lemma-base-change-does-not-hold}，只需证明
$R^qh_*\mathbf{Z}/d\mathbf{Z} = 0$ 对 $d | n$ 成立，只要有笛卡尔图
$$
\xymatrix{
X \ar[d] & Y \ar[d] \ar[l]^h \\
S & \Spec(K) \ar[l]
}
$$
，其中 $X \to S$ 仿射、光滑且相对维数为 $1$，$S$ 是正规整环
$A$ 的谱，其分式域 $L$ 代数闭，而 $K/L$ 是代数闭域的扩张。

\medskip\noindent
回顾，$R^qh_*\mathbf{Z}/d\mathbf{Z}$ 是下列预层的相伴层：
$$
U \longmapsto
H^q(U \times_X Y, \mathbf{Z}/d\mathbf{Z}) =
H^q(U \times_S \Spec(K), \mathbf{Z}/d\mathbf{Z})
$$
；它定义在 $X_\etale$ 上（引理 \ref{etale-cohomology-lemma-higher-direct-images}）。
因此只需证明：给定 $U$ 以及
$\xi \in H^q(U \times_S \Spec(K), \mathbf{Z}/d\mathbf{Z})$，存在 \'etale 覆盖
$\{U_i \to U\}$，使得 $\xi$ 在
$H^q(U_i \times_S \Spec(K), \mathbf{Z}/d\mathbf{Z})$ 中变为零。

\medskip\noindent
当然，可以取 $U$ 仿射。此时 $U \times_S \Spec(K)$ 是 $K$ 上的
（光滑）仿射曲线，因而由定理 \ref{etale-cohomology-theorem-vanishing-affine-curves}，
当 $q > 1$ 时有消失。

\medskip\noindent
最后一种情形：$q = 1$。可以像 More on Morphisms，引理
\ref{more-morphisms-lemma-cover-smooth-by-special} 中那样，用一个 \'etale
覆盖的成员替换 $U$。于是 $U \to S$ 分解为 $U \to V \to S$，其中
$U \to V$ 的纤维几何连通，$U$、$V$ 仿射，$V \to S$ 是 \'etale 的，
并且有截面 $\sigma : V \to U$。由引理
\ref{etale-cohomology-lemma-normal-scheme-with-alg-closed-function-field}，$V$ 同构于 $S$
的（有限）多个（仿射）开子概形之不交并。显然，可以用其中一个替换
$S$，并替换 $X$，所用的是 $U$ 的相应分支。因此可以假设 $X \to S$ 的纤维
几何连通，有截面 $\sigma$，并且
$\xi \in H^1(Y, \mathbf{Z}/d\mathbf{Z})$。由于 $K$ 与 $L$ 都代数闭，有
$$
H^1(X_L, \mathbf{Z}/d\mathbf{Z}) =
H^1(Y, \mathbf{Z}/d\mathbf{Z})
$$
，见引理 \ref{etale-cohomology-lemma-base-change-dim-1-separably-closed}。因此存在有限
\'etale Galois 覆盖 $Z \to X_L$，其 Galois 群为
$G \subset \mathbf{Z}/d\mathbf{Z}$，并且该覆盖零化 $\xi$。
这可以从引理 \ref{etale-cohomology-lemma-smooth-base-change-fields} 的陈述或证明看出，
也可直接使用 $\xi$ 对应于一个 $\mathbf{Z}/d\mathbf{Z}$-挠子，
它定义在 $X_L$ 上。最后，由 Fundamental Groups，引理
\ref{pione-lemma-extend-covering-general}，得到（必然满的）有限 \'etale 态射
$X' \to X$，其到 $X_L$ 的限制是 $Z \to X_L$。由于 $\xi$ 在
$X'_K$ 中变为零，证明完毕。
\end{proof}

\noindent
下面这个光滑基变换定理的直接推论，才是实践中经常使用的形式。

\begin{lemma}
\label{etale-cohomology-lemma-smooth-base-change-general}
设 $S$ 是概形。设 $S' = \lim S_i$ 是概形 $S_i$ 构成的有向逆极限，
其中各概形在 $S$ 上光滑，转移态射仿射。设 $f : X \to S$
拟紧且拟分离，并作纤维方块
$$
\xymatrix{
X' \ar[d]_{f'} \ar[r]_{g'} & X \ar[d]^f \\
S' \ar[r]^g & S
}
$$
那么
$$
g^{-1}Rf_*E = R(f')_*(g')^{-1}E
$$
对任意 $E \in D^+(X_\etale)$ 都成立，只要其上同调层 $H^q(E)$ 的茎
都是挠群，且其阶在 $S$ 上可逆。
\end{lemma}

\begin{proof}
考虑谱序列
$$
E_2^{p, q} = R^pf_*H^q(E)
\quad\text{且}\quad
{E'}_2^{p, q} = R^pf'_*H^q((g')^{-1}E) = R^pf'_*(g')^{-1}H^q(E)
$$
，它们收敛到 $R^nf_*E$ 与 $R^nf'_*(g')^{-1}E$。
这些谱序列在 Derived Categories，引理
\ref{derived-lemma-two-ss-complex-functor} 中构造。把光滑基变换定理
（定理 \ref{etale-cohomology-theorem-smooth-base-change}）与引理
\ref{etale-cohomology-lemma-base-change-Rf-star-colim} 结合，可知
$$
g^{-1}R^pf_*H^q(E) = R^p(f')_*(g')^{-1}H^q(E)
$$
综合以上各点即得本引理。
\end{proof}












\section{光滑基变换的应用}
\label{etale-cohomology-section-applications-smooth-base-change}

\noindent
本节讨论光滑基变换定理的一些颇为直接的推论。

\begin{lemma}
\label{etale-cohomology-lemma-base-change-field-extension}
设 $L/K$ 是域扩张。设 $g : T \to S$ 是 $K$ 上概形之间的拟紧拟分离态射。
记 $g_L : T_L \to S_L$ 为 $g$ 到 $\Spec(L)$ 的基变。
设 $E \in D^+(T_\etale)$ 的上同调层的茎都是挠群，且其阶在 $K$ 中可逆。
设 $E_L$ 为 $E$ 到 $(T_L)_\etale$ 的拉回。则
$Rg_{L, *}E_L$ 是 $Rg_*E$ 到 $S_L$ 的拉回。
\end{lemma}

\begin{proof}
若 $L/K$ 可分，则 $L$ 是光滑 $K$-代数的滤过余极限，见
Algebra, Lemma \ref{algebra-lemma-colimit-syntomic}。
因此在这种情形下，本引理立即由
Lemma \ref{etale-cohomology-lemma-smooth-base-change-general} 得出。
一般情形下，令 $K'$ 和 $L'$ 分别为 $K$ 和 $L$ 的完美闭包
（Algebra, Definition \ref{algebra-definition-perfection}）。
态射 $\Spec(K') \to \Spec(K)$ 和 $\Spec(L') \to \Spec(L)$
都是泛同胚，因为 $K'/K$ 和 $L'/L$ 都是纯不可分扩张
（见 Algebra, Lemma \ref{algebra-lemma-p-ring-map}）。
由 \'etale 上同调的拓扑不变性，我们于是有
$(T_{K'})_\etale = T_\etale$、
$(S_{K'})_\etale = S_\etale$、
$(T_{L'})_\etale = (T_L)\etale$ 以及
$(S_{L'})_\etale = (S_L)_\etale$，见
Proposition \ref{etale-cohomology-proposition-topological-invariance}。
这将本引理归结为可分域扩张 $L'/K'$ 的情形
（由完美域的定义可知，见 Algebra, Definition \ref{algebra-definition-perfect}）。
\end{proof}

\begin{lemma}
\label{etale-cohomology-lemma-smooth-base-change-separably-closed}
设 $K/k$ 是可分闭域的扩张。设 $X$ 是 $k$ 上的拟紧拟分离概形。
设 $E \in D^+(X_\etale)$ 的上同调层的茎都是挠群，且其阶在 $k$ 中可逆。则
\begin{enumerate}
\item 映射 $H^q_\etale(X, E) \to H^q_\etale(X_K, E|_{X_K})$ 是同构；并且
\item $E \to R(X_K \to X)_*E|_{X_K}$ 是同构。
\end{enumerate}
\end{lemma}

\begin{proof}
证明 (1)。首先，令 $\overline{k}$ 和 $\overline{K}$
分别为 $k$ 和 $K$ 的代数闭包。态射
$\Spec(\overline{k}) \to \Spec(k)$ 和
$\Spec(\overline{K}) \to \Spec(K)$ 都是泛同胚，因为
$\overline{k}/k$ 和 $\overline{K}/K$ 都是纯不可分扩张
（见 Algebra, Lemma \ref{algebra-lemma-p-ring-map}）。
因此，由 \'etale 上同调的拓扑不变性，
$H^q_\etale(X, \mathcal{F}) =
H^q_\etale(X_{\overline{k}}, \mathcal{F}_{X_{\overline{k}}})$，见
Proposition \ref{etale-cohomology-proposition-topological-invariance}。
对 $X_K$ 和 $X_{\overline{K}}$ 也同样成立。
因此可以假设 $k$ 和 $K$ 都是代数闭域。
在这种情形下，$K$ 是光滑 $k$-代数的极限，见
Algebra, Lemma \ref{algebra-lemma-colimit-syntomic}。
由此，本引理是 Theorem \ref{etale-cohomology-theorem-smooth-base-change}
按 Lemma \ref{etale-cohomology-lemma-smooth-base-change-general} 中的形式所述的一个特例。

\medskip\noindent
证明 (2)。取 $U$ 为 $X_\etale$ 中任意拟紧拟分离对象。上述结论表明，
映射 $E \to R(X_K \to X)_*E|_{X_K}$ 的限制诱导上同调上的同构。
由于 $X_\etale$ 的每个对象都有一个由此类 $U$ 构成的 \'etale 覆盖，
这就证明了所需结论。
\end{proof}

\begin{lemma}
\label{etale-cohomology-lemma-base-change-does-not-hold-post}
设 $f : X \to S$ 和 $n$ 如 Remark \ref{etale-cohomology-remark-base-change-holds} 中所述。
假设 $n$ 在 $S$ 上可逆，并且对某个 $q \geq 1$，
$BC(f, n, q - 1)$ 成立，而 $BC(f, n, q)$ 不成立。
则存在交换图
$$
\xymatrix{
X \ar[d]_f & X' \ar[d] \ar[l] & Y \ar[l]^h \ar[d] \\
S & S' \ar[l] & \Spec(K) \ar[l]
}
$$
其中两个方块都是笛卡尔的，$S'$ 是仿射、整且正规的，
其函数域 $K$ 为代数闭域；并且存在整数 $d | n$，使得
$R^qh_*(\mathbf{Z}/d\mathbf{Z})$ 非零。
\end{lemma}

\begin{proof}
首先按照 Lemma \ref{etale-cohomology-lemma-base-change-does-not-hold}
选取一个图和 $\mathcal{F}$。
可以并且确实假设 $S'$ 是仿射的（这是显然的；如有疑问，
可参见该引理的证明）。令 $K'$ 为 $S'$ 的函数域，并令
$Y' = X' \times_{S'} \Spec(K')$，得到交换图
$$
\xymatrix{
X \ar[d]_f &
X' \ar[d] \ar[l] &
Y' \ar[l]^{h'} \ar[d] &
Y \ar[l] \ar[d] \\
S &
S' \ar[l] &
\Spec(K') \ar[l] &
\Spec(K) \ar[l]
}
$$
由 Lemma \ref{etale-cohomology-lemma-smooth-base-change-separably-closed}，
总正像 $R(Y \to Y')_*\mathbf{Z}/d\mathbf{Z}$
同构于 $\mathbf{Z}/d\mathbf{Z}$，此同构位于 $D(Y'_\etale)$ 中；这里用到了
$n$ 在 $S$ 上可逆这一条件。因此，由相对 Leray 谱序列可得
$Rh'_*\mathbf{Z}/d\mathbf{Z} = Rh_*\mathbf{Z}/d\mathbf{Z}$。
证毕。
\end{proof}










\section{真基变换定理}
\label{etale-cohomology-section-proper-base-change}

\noindent
本节陈述并证明真基变换定理。
我们的方法大致遵循 \cite[XII, Theorem 5.1]{SGA4} 中的证明，
并使用 Gabber（源自仿射情形）的思想稍稍简化论证。

\begin{lemma}
\label{etale-cohomology-lemma-zariski-h0-proper-over-henselian-pair}
设 $(A, I)$ 是 Hensel 对。设 $f : X \to \Spec(A)$ 是概形之间的真态射。
令 $Z = X \times_{\Spec(A)} \Spec(A/I)$。对任意层 $\mathcal{F}$，
若它定义在与 $X$ 相伴的拓扑空间上，则
$\Gamma(X, \mathcal{F}) = \Gamma(Z, \mathcal{F}|_Z)$。
\end{lemma}

\begin{proof}
我们用 Lemma \ref{etale-cohomology-lemma-h0-topological} 来证明这一点。首先，由 Properties, Lemma
\ref{properties-lemma-quasi-compact-quasi-separated-spectral}，
$X$ 的底拓扑空间是谱空间。设 $Y \subset X$ 是不可约闭子概形。
为完成证明，我们说明 $Y \cap Z = Y \times_{\Spec(A)} \Spec(A/I)$ 连通。
把 $X$ 替换为 $Y$ 后，可以假设 $X$ 不可约，而需要证明 $Z$ 连通。
设 $X \to \Spec(B) \to \Spec(A)$ 为 $f$ 的 Stein 分解
（More on Morphisms, Theorem
\ref{more-morphisms-theorem-stein-factorization-general}）。
此时 $A \to B$ 是整的，并且 $(B, IB)$ 是 Hensel 对
（More on Algebra, Lemma \ref{more-algebra-lemma-integral-over-henselian-pair}）。
因此可以假设 $X \to \Spec(A)$ 的纤维几何连通。另一方面，
$T \subset \Spec(A)$ 是 $f$ 的像，因而不可约；又因 $X$ 在 $A$ 上真，该像是闭的。
故由 More on Algebra, Lemma
\ref{more-algebra-lemma-irreducible-henselian-pair-connected}，
$T \cap V(I)$ 连通。现在
$Y \times_{\Spec(A)} \Spec(A/I) \to T \cap V(I)$
是具有连通纤维的满闭映射。结论于是由 Topology, Lemma
\ref{topology-lemma-connected-fibres-quotient-topology-connected-components} 得出。
\end{proof}

\begin{lemma}
\label{etale-cohomology-lemma-h0-proper-over-henselian-pair}
设 $(A, I)$ 是 Hensel 对。设 $f : X \to \Spec(A)$ 是概形之间的真态射。
令 $i : Z \to X$ 为 $X \times_{\Spec(A)} \Spec(A/I)$ 到 $X$ 的闭浸入。
对任意层 $\mathcal{F}$，若它定义在 $X_\etale$ 上，则
$\Gamma(X, \mathcal{F}) = \Gamma(Z, i_{small}^{-1}\mathcal{F})$。
\end{lemma}

\begin{proof}
这由 Lemma \ref{etale-cohomology-lemma-gabber-h0}、
\ref{etale-cohomology-lemma-zariski-h0-proper-over-henselian-pair}，
以及在 $X$ 上有限的任意概形在 $\Spec(A)$ 上都是真这一事实得出。
\end{proof}

\begin{lemma}
\label{etale-cohomology-lemma-h0-proper-over-henselian-local}
设 $A$ 是 Hensel 局部环。设 $f : X \to \Spec(A)$ 是概形之间的真态射。
令 $X_0 \subset X$ 为 $f$ 在闭点上方的纤维。对任意层 $\mathcal{F}$，
若它定义在 $X_\etale$ 上，则
$\Gamma(X, \mathcal{F}) = \Gamma(X_0, \mathcal{F}|_{X_0})$。
\end{lemma}

\begin{proof}
这是 Lemma \ref{etale-cohomology-lemma-h0-proper-over-henselian-pair} 的特例。
\end{proof}

\noindent
设 $f : X \to S$ 是概形之间的态射，并设
$\overline{s} : \Spec(k) \to S$ 是几何点。$f$ 在 $\overline{s}$ 处的纤维是概形
$X_{\overline{s}} = \Spec(k) \times_{\overline{s}, S} X$，
视为 $\Spec(k)$ 上的概形。若 $\mathcal{F}$ 是 $X_\etale$ 上的层，则记
$\mathcal{F}_{\overline{s}} = p_{small}^{-1}\mathcal{F}$
为 $\mathcal{F}$ 到 $(X_{\overline{s}})_\etale$ 的拉回。
下文将考虑集合
$$
\Gamma(X_{\overline{s}}, \mathcal{F}_{\overline{s}})
$$
令 $s \in S$ 为 $\overline{s}$ 的像点。令 $\kappa(s)^{sep}$
为 Definition \ref{etale-cohomology-definition-algebraic-geometric-point} 中所述的
$\kappa(s)$ 在 $k$ 中的可分代数闭包。
由 Lemma \ref{etale-cohomology-lemma-sections-base-field-extension}，拉回给出双射
$$
\Gamma(X_{\kappa(s)^{sep}},
p_{sep}^{-1} \mathcal{F})
\longrightarrow
\Gamma(X_{\overline{s}}, \mathcal{F}_{\overline{s}})
$$
其中 $p_{sep} : X_{\kappa(s)^{sep}} = \Spec(\kappa(s)^{sep}) \times_S X \to X$
是投影。

\begin{lemma}
\label{etale-cohomology-lemma-proper-pushforward-stalk}
设 $f : X \to S$ 是概形之间的真态射，并设
$\overline{s} \to S$ 是几何点。对任意层 $\mathcal{F}$，
若它定义在 $X_\etale$ 上，则典范映射
$$
(f_*\mathcal{F})_{\overline{s}} \longrightarrow
\Gamma(X_{\overline{s}}, \mathcal{F}_{\overline{s}})
$$
是双射。
\end{lemma}

\begin{proof}
由 Theorem \ref{etale-cohomology-theorem-higher-direct-images}（用于集合层）有
$$
(f_*\mathcal{F})_{\overline{s}} =
\Gamma(X \times_S \Spec(\mathcal{O}_{S, \overline{s}}^{sh}),
p_{small}^{-1}\mathcal{F})
$$
其中 $p : X \times_S \Spec(\mathcal{O}_{S, \overline{s}}^{sh}) \to X$
是投影。严格 Hensel 局部环 $\mathcal{O}_{S, \overline{s}}^{sh}$
的剩余域是 $\kappa(s)^{sep}$，故本引理由上述讨论和
Lemma \ref{etale-cohomology-lemma-h0-proper-over-henselian-local} 得出。
\end{proof}

\begin{lemma}
\label{etale-cohomology-lemma-proper-base-change-f-star}
设 $f : X \to Y$ 是概形之间的真态射，$g : Y' \to Y$
是概形之间的态射。置 $X' = Y' \times_Y X$，其投影为
$f' : X' \to Y'$ 和 $g' : X' \to X$。设 $\mathcal{F}$ 是
$X_\etale$ 上的任意层。则 $g^{-1}f_*\mathcal{F} = f'_*(g')^{-1}\mathcal{F}$。
\end{lemma}

\begin{proof}
存在典范映射 $g^{-1}f_*\mathcal{F} \to f'_*(g')^{-1}\mathcal{F}$。
具体地说，它是下述映射的伴随：
$$
f_*\mathcal{F} \longrightarrow
g_*f'_*(g')^{-1}\mathcal{F} = f_*g'_*(g')^{-1}\mathcal{F}
$$
它由 $f_*$ 作用于典范映射
$\mathcal{F} \to g'_*(g')^{-1}\mathcal{F}$ 得到。要检验该映射为同构，
可以计算其在茎上的作用。设 $y' : \Spec(k) \to Y'$ 是几何点，
其像 $y$ 位于 $Y$ 中。由 Lemma \ref{etale-cohomology-lemma-proper-pushforward-stalk}，
两个茎分别为 $\Gamma(X'_{y'}, \mathcal{F}_{y'})$ 和
$\Gamma(X_y, \mathcal{F}_y)$。这里，层 $\mathcal{F}_y$ 和
$\mathcal{F}_{y'}$ 分别是 $\mathcal{F}$ 沿投影 $X_y \to X$
和 $X'_{y'} \to X$ 的拉回。由于作为概形有 $X_y = X'_{y'}$，
可见这些茎相同。我们略去此同构与上述映射相容的验证。
\end{proof}

\noindent
现在开始讨论真基变换定理。为此先引入一些记号。考虑交换图
\begin{equation}
\label{etale-cohomology-equation-base-change-diagram}
\vcenter{
\xymatrix{
X' \ar[r]_{g'} \ar[d]_{f'} & X \ar[d]^f \\
Y' \ar[r]^g & Y
}
}
\end{equation}
其中各箭头都是概形态射。于是得到位点的交换图
$$
\xymatrix{
X'_\etale \ar[r]_{g'_{small}} \ar[d]_{f'_{small}} &
X_\etale \ar[d]^{f_{small}} \\
Y'_\etale \ar[r]^{g_{small}} &
Y_\etale
}
$$
对任意对象 $E$，若它属于 $D(X_\etale)$，则得到典范基变换映射
\begin{equation}
\label{etale-cohomology-equation-base-change}
g_{small}^{-1}Rf_{small, *}E \longrightarrow Rf'_{small, *}(g'_{small})^{-1}E
\end{equation}
它位于 $D(Y'_\etale)$ 中。见 Cohomology on Sites, Remark
\ref{sites-cohomology-remark-base-change}，其中以常值层
$\mathbf{Z}$ 作为环层。此公式中通常省略下标 ${}_{small}$。
例如，若 $E = \mathcal{F}[0]$，其中 $\mathcal{F}$ 是
$X_\etale$ 上的阿贝尔层，则基变换映射为
\begin{equation}
\label{etale-cohomology-equation-base-change-sheaf}
g^{-1}Rf_*\mathcal{F} \longrightarrow Rf'_*(g')^{-1}\mathcal{F}
\end{equation}
它位于 $D(Y'_\etale)$ 中。

\medskip\noindent
在上述一般性下，映射 (\ref{etale-cohomology-equation-base-change}) 不可能总是同构。
我们的目标是证明：如果图 (\ref{etale-cohomology-equation-base-change-diagram}) 是笛卡尔的，
$f : X \to Y$ 为真态射，$E$ 的上同调层都是挠层，且 $E$ 下有界，
那么该映射是同构。为研究这个问题，引入如下术语。
称{\it 对 $f : X \to Y$，上同调与基变换可交换}，如果
(\ref{etale-cohomology-equation-base-change-sheaf}) 对每个图
(\ref{etale-cohomology-equation-base-change-diagram})（其中 $X' = Y' \times_Y X$）
以及每个挠阿贝尔层 $\mathcal{F}$ 都是同构。

\begin{lemma}
\label{etale-cohomology-lemma-proper-base-change-in-terms-of-injectives}
设 $f : X \to Y$ 是概形之间的真态射。下列条件等价：
\begin{enumerate}
\item 对 $f$，上同调与基变换可交换（见上文）；
\item 对每个素数 $\ell$、每个内射
$\mathbf{Z}/\ell\mathbf{Z}$-模层 $\mathcal{I}$（定义在 $X_\etale$ 上），以及每个满足
$X' = Y' \times_Y X$ 的图 (\ref{etale-cohomology-equation-base-change-diagram})，
层 $R^qf'_*(g')^{-1}\mathcal{I}$ 在 $q > 0$ 时为零。
\end{enumerate}
\end{lemma}

\begin{proof}
显然 (1) 蕴含 (2)。反过来，假设 (2)，并设 $\mathcal{F}$ 是
$X_\etale$ 上的挠阿贝尔层。设 $Y' \to Y$ 是概形之间的态射，并令
$X' = Y' \times_Y X$，其投影 $g' : X' \to X$ 和
$f' : X' \to Y'$ 如图 (\ref{etale-cohomology-equation-base-change-diagram}) 所示。
我们要证明层的映射
$$
g^{-1}R^qf_*\mathcal{F} \longrightarrow R^qf'_*(g')^{-1}\mathcal{F}
$$
对所有 $q \geq 0$ 都是同构。

\medskip\noindent
对每个 $n \geq 1$，令 $\mathcal{F}[n]$ 为 $\mathcal{F}$ 中由
$n$ 零化的截面构成的子层。于是
$\mathcal{F} = \colim \mathcal{F}[n]$.
函子 $g^{-1}$ 和 $(g')^{-1}$ 与任意余极限可交换（因为它们是左伴随）。
由 Lemma \ref{etale-cohomology-lemma-relative-colimit}，沿 $f$ 或 $f'$ 取高阶正像与滤过余极限可交换。
因此有
$$
g^{-1}R^qf_*\mathcal{F} = \colim g^{-1}R^qf_*\mathcal{F}[n]
\quad\text{且}\quad
R^qf'_*(g')^{-1}\mathcal{F} =
\colim R^qf'_*(g')^{-1}\mathcal{F}[n]
$$
故只需对 $\mathcal{F}$ 被正整数 $n$ 零化的情形证明结论。

\medskip\noindent
若 $n = \ell n'$，其中 $\ell$ 是某个素数，则得到短正合列
$$
0 \to \mathcal{F}[\ell] \to \mathcal{F} \to
\mathcal{F}/\mathcal{F}[\ell] \to 0
$$
注意 $\mathcal{F}/\mathcal{F}[\ell]$ 被 $n'$ 零化。此外，若结论对
$\mathcal{F}[\ell]$ 和 $\mathcal{F}/\mathcal{F}[\ell]$ 都成立，
则由高阶正像的长正合列（以及 $5$ 引理），结论成立。
这样便归结为 $\mathcal{F}$ 被某个素数 $\ell$ 零化的情形。

\medskip\noindent
假设 $\mathcal{F}$ 被素数 $\ell$ 零化。选取内射分解
$\mathcal{F} \to \mathcal{I}^\bullet$，它位于
$D(X_\etale, \mathbf{Z}/\ell\mathbf{Z})$ 中。应用条件 (2) 和 Leray 非循环引理
（Derived Categories, Lemma \ref{derived-lemma-leray-acyclicity}），可知
$$
f'_*(g')^{-1}\mathcal{I}^\bullet
$$
计算 $Rf'_*(g')^{-1}\mathcal{F}$。再应用
Lemma \ref{etale-cohomology-lemma-proper-base-change-f-star} 即得结论。
\end{proof}

\begin{lemma}
\label{etale-cohomology-lemma-sandwich}
设 $f : X \to Y$ 和 $g : Y \to Z$ 是概形之间的真态射。假设
\begin{enumerate}
\item 对 $f$，上同调与基变换可交换；
\item 对 $g \circ f$，上同调与基变换可交换；并且
\item $f$ 是满射。
\end{enumerate}
则对 $g$，上同调与基变换可交换。
\end{lemma}

\begin{proof}
我们将使用 Lemma \ref{etale-cohomology-lemma-proper-base-change-in-terms-of-injectives}
中的等价性而不再特别说明。设 $\ell$ 是素数。
设 $\mathcal{I}$ 是一个内射层，它是 $\mathbf{Z}/\ell\mathbf{Z}$-模层，
并定义在 $Y_\etale$ 上。选取层的单射
$f^{-1}\mathcal{I} \to \mathcal{J}$，其中 $\mathcal{J}$ 是一个内射层，
它是 $\mathbf{Z}/\ell\mathbf{Z}$-模层，并定义在 $X_\etale$ 上。
由于 $f$ 是满射，映射 $\mathcal{I} \to f_*\mathcal{J}$ 是单射
（在几何点处看茎即可）。由于 $\mathcal{I}$ 内射，可知
$\mathcal{I}$ 是 $f_*\mathcal{J}$ 的直和项。因此只需对
$f_*\mathcal{J}$ 证明所需的消失。

\medskip\noindent
设 $Z' \to Z$ 是概形之间的态射，并置
$Y' = Z' \times_Z Y$ 和 $X' = Z' \times_Z X = Y' \times_ Y X$。
记 $a : X' \to X$、$b : Y' \to Y$ 和 $c : Z' \to Z$ 为投影；
$f' : X' \to Y'$ 和 $g' : Y' \to Z'$ 也同样记。由
Lemma \ref{etale-cohomology-lemma-proper-base-change-f-star} 有
$b^{-1}f_*\mathcal{J} = f'_*a^{-1}\mathcal{J}$.
另一方面，已知 $R^qf'_*a^{-1}\mathcal{J}$ 和
$R^q(g' \circ f')_*a^{-1}\mathcal{J}$ 在 $q > 0$ 时为零。
利用谱序列
(Cohomology on Sites，引理 \ref{sites-cohomology-lemma-relative-Leray})
$$
R^pg'_* R^qf'_* a^{-1}\mathcal{J} \Rightarrow
R^{p + q}(g' \circ f')_* a^{-1}\mathcal{J}
$$
可得
$ R^pg'_*(b^{-1}f_*\mathcal{J}) = R^pg'_*(f'_*a^{-1}\mathcal{J}) = 0$
在 $p > 0$ 时成立，这正是所需结论。
\end{proof}

\begin{lemma}
\label{etale-cohomology-lemma-composition}
设 $f : X \to Y$ 和 $g : Y \to Z$ 是概形之间的真态射。假设
\begin{enumerate}
\item 对 $f$，上同调与基变换可交换；并且
\item 对 $g$，上同调与基变换可交换。
\end{enumerate}
则对 $g \circ f$，上同调与基变换可交换。
\end{lemma}

\begin{proof}
我们将使用 Lemma \ref{etale-cohomology-lemma-proper-base-change-in-terms-of-injectives}
中的等价性而不再特别说明。设 $\ell$ 是素数。
设 $\mathcal{I}$ 是一个内射层，它是 $\mathbf{Z}/\ell\mathbf{Z}$-模层，
并定义在 $X_\etale$ 上。于是 $f_*\mathcal{I}$ 是一个内射层，
它是 $\mathbf{Z}/\ell\mathbf{Z}$-模层，并定义在 $Y_\etale$ 上
（Cohomology on Sites, Lemma
\ref{sites-cohomology-lemma-pushforward-injective-flat}）。
结论可由此形式地得出，但我们也把论证写明。

\medskip\noindent
设 $Z' \to Z$ 是概形之间的态射，并置
$Y' = Z' \times_Z Y$ 和 $X' = Z' \times_Z X = Y' \times_ Y X$。
记 $a : X' \to X$、$b : Y' \to Y$ 和 $c : Z' \to Z$ 为投影；
$f' : X' \to Y'$ 和 $g' : Y' \to Z'$ 也同样记。由
Lemma \ref{etale-cohomology-lemma-proper-base-change-f-star} 有
$b^{-1}f_*\mathcal{I} = f'_*a^{-1}\mathcal{I}$.
另一方面，已知 $R^qf'_*a^{-1}\mathcal{I}$ 和
$R^q(g')_*b^{-1}f_*\mathcal{I}$ 在 $q > 0$ 时为零。
利用谱序列
(Cohomology on Sites，引理 \ref{sites-cohomology-lemma-relative-Leray})
$$
R^pg'_* R^qf'_* a^{-1}\mathcal{I} \Rightarrow
R^{p + q}(g' \circ f')_* a^{-1}\mathcal{I}
$$
可得 $R^p(g' \circ f')_*a^{-1}\mathcal{I} = 0$ 在
$p > 0$ 时成立，这正是所需结论。
\end{proof}

\begin{lemma}
\label{etale-cohomology-lemma-finite}
\begin{slogan}
\'Etale 上同调的真基变换对有限态射成立。
\end{slogan}
设 $f : X \to Y$ 是概形之间的有限态射。
则对 $f$，上同调与基变换可交换。
\end{lemma}

\begin{proof}
注意有限态射是真态射，见
Morphisms, Lemma \ref{morphisms-lemma-finite-proper}。
此外，有限态射的基变换仍然有限，见
Morphisms, Lemma \ref{morphisms-lemma-base-change-finite}。
因此，将 Lemma \ref{etale-cohomology-lemma-proper-base-change-in-terms-of-injectives}
与 Proposition \ref{etale-cohomology-proposition-finite-higher-direct-image-zero}
结合便得结论。
\end{proof}

\begin{lemma}
\label{etale-cohomology-lemma-reduce-to-P1}
要证明每个概形真态射的上同调都与基变换可交换，
只需证明该结论对态射 $\mathbf{P}^1_S \to S$ 成立，其中 $S$ 是任意概形。
\end{lemma}

\begin{proof}
设 $f : X \to Y$ 是概形之间的真态射。设 $Y = \bigcup Y_i$
是仿射开覆盖，并置 $X_i = f^{-1}(Y_i)$。如果能够证明对
$X_i \to Y_i$ 上同调与基变换可交换，那么对 $f$ 也可交换。
这是因为高阶正像的形成与基底上的 Zariski（甚至 \'etale）局部化可交换，见
Lemma \ref{etale-cohomology-lemma-higher-direct-images}。因此可以假设 $Y$ 仿射。

\medskip\noindent
设 $Y$ 是仿射概形，并设 $X \to Y$ 是真态射。由 Chow 引理，存在交换图
$$
\xymatrix{
X \ar[rd] & X' \ar[d] \ar[l]^\pi \ar[r] & \mathbf{P}^n_Y \ar[dl] \\
& Y &
}
$$
其中 $X' \to \mathbf{P}^n_Y$ 是浸入，而 $\pi : X' \to X$
是真满射，见 Limits, Lemma \ref{limits-lemma-chow-finite-type}。
由于 $X \to Y$ 真，可知 $X' \to Y$ 真
（Morphisms, Lemma \ref{morphisms-lemma-composition-proper}）。
故 $X' \to \mathbf{P}^n_Y$ 是闭浸入
（Morphisms, Lemma \ref{morphisms-lemma-image-proper-scheme-closed}）。
从而 $X' \to X \times_Y \mathbf{P}^n_Y = \mathbf{P}^n_X$
是闭浸入（因为它是像闭的浸入）。

\medskip\noindent
由 Lemma \ref{etale-cohomology-lemma-sandwich}，只需证明对 $\pi$ 和 $X' \to Y$，
上同调与基变换可交换。这两个态射都可分解为一个闭浸入，
再接投影 $\mathbf{P}^n_S \to S$（其中 $S$ 为某个概形）。
由 Lemma \ref{etale-cohomology-lemma-finite}，结论对闭浸入成立（因为闭浸入有限）。
由 Lemma \ref{etale-cohomology-lemma-composition}，只需对投影
$\mathbf{P}^n_S \to S$ 证明结论。

\medskip\noindent
对每个 $n \geq 1$，存在有限满态射
$$
\mathbf{P}^1_S \times_S \ldots \times_S \mathbf{P}^1_S
\longrightarrow
\mathbf{P}^n_S
$$
它在坐标上由下式给出：
$$
((x_1 : y_1), (x_2 : y_2), \ldots, (x_n : y_n))
\longmapsto
(F_0 : \ldots : F_n)
$$
其中 $F_0, \ldots, F_n$ 是变量 $x_1, \ldots, y_n$
的整系数多项式，满足
$$
\prod (x_i t + y_i) = F_0 t^n + F_1 t^{n - 1} + \ldots + F_n
$$
再次应用 Lemmas \ref{etale-cohomology-lemma-sandwich}, \ref{etale-cohomology-lemma-finite}, and
\ref{etale-cohomology-lemma-composition}，即知本引理成立。
\end{proof}

\begin{theorem}[真基变换]
\label{etale-cohomology-theorem-proper-base-change}
设 $f : X \to Y$ 是概形之间的真态射，$g : Y' \to Y$
是概形之间的态射。置 $X' = Y' \times_Y X$，并考虑笛卡尔图
$$
\xymatrix{
X' \ar[r]_{g'} \ar[d]_{f'} & X \ar[d]^f \\
Y' \ar[r]^g & Y
}
$$
设 $\mathcal{F}$ 是 $X_\etale$ 上的挠阿贝尔层。则基变换映射
$$
g^{-1}Rf_*\mathcal{F} \longrightarrow Rf'_*(g')^{-1}\mathcal{F}
$$
是同构。
\end{theorem}

\begin{proof}
按照上面引入的术语，这就是说，对概形之间的每个真态射，
上同调与基变换可交换。由 Lemma \ref{etale-cohomology-lemma-reduce-to-P1}，
只需证明上同调对态射 $\mathbf{P}^1_S \to S$ 与基变换可交换，
其中 $S$ 是任意概形。

\medskip\noindent
设 $S$ 是严格 Hensel 局部环的谱，其闭点为 $s$。置
$X = \mathbf{P}^1_S$ 和 $X_0 = X_s = \mathbf{P}^1_s$。
设 $\mathcal{F}$ 是一个 $\mathbf{Z}/\ell\mathbf{Z}$-模层，
定义在 $X_\etale$ 上。证明的关键是
$$
H^q_\etale(X, \mathcal{F}) = H^q_\etale(X_0, \mathcal{F}|_{X_0}).
$$
具体地，选取分解 $\mathcal{F} \to \mathcal{I}^\bullet$，
其各项是内射 $\mathbf{Z}/\ell\mathbf{Z}$-模层。于是
$\mathcal{I}^\bullet|_{X_0}$ 是 $\mathcal{F}|_{X_0}$ 的一个分解，
其各项对右函子 $H^0_\etale(X_0, -)$ 非循环，见
Lemma \ref{etale-cohomology-lemma-efface-cohomology-on-fibre-by-finite-cover}。
Leray 非循环引理说明，右端由复形
$H^0_\etale(X_0, \mathcal{I}^\bullet|_{X_0})$ 计算；由
Lemma \ref{etale-cohomology-lemma-h0-proper-over-henselian-local}，它等于
$H^0_\etale(X, \mathcal{I}^\bullet)$。该复形计算左端。

\medskip\noindent
现在设 $S$ 是一般概形，$\mathcal{F}$ 是一个
$\mathbf{Z}/\ell\mathbf{Z}$-模层，定义在 $X_\etale$ 上。
设 $\overline{s} : \Spec(k) \to S$ 是 $S$ 的几何点，
它位于 $s \in S$ 上方。于是
$$
(R^qf_*\mathcal{F})_{\overline{s}} =
H^q_\etale(\mathbf{P}^1_{\mathcal{O}_{S, \overline{s}}^{sh}},
\mathcal{F}|_{\mathbf{P}^1_{\mathcal{O}_{S, \overline{s}}^{sh}}}) =
H^q_\etale(\mathbf{P}^1_{\kappa(s)^{sep}},
\mathcal{F}|_{\mathbf{P}^1_{\kappa(s)^{sep}}})
$$
其中 $\kappa(s)^{sep}$ 是 $\mathcal{O}_{S, \overline{s}}^{sh}$ 的剩余域，
即 $\kappa(s)$ 在 $k$ 中的可分代数闭包。
第一个等号由 Theorem \ref{etale-cohomology-theorem-higher-direct-images} 得出，
第二个等号由上一段的显示公式得出。

\medskip\noindent
最后，考虑概形之间的任意态射 $g : T \to S$，其中
$S$ 和 $\mathcal{F}$ 如上。令 $f' : \mathbf{P}^1_T \to T$ 为投影，
并令 $g' : \mathbf{P}^1_T \to \mathbf{P}^1_S$ 为 $g$ 诱导的态射。
考虑基变换映射
$$
g^{-1}R^qf_*\mathcal{F}
\longrightarrow
R^qf'_*(g')^{-1}\mathcal{F}
$$
设 $\overline{t}$ 是 $T$ 的几何点，其像为
$\overline{s} = g(\overline{t})$。由上述讨论，该映射在
$\overline{t}$ 处的茎上的映射为
$$
H^q_\etale(\mathbf{P}^1_{\kappa(s)^{sep}},
\mathcal{F}|_{\mathbf{P}^1_{\kappa(s)^{sep}}})
\longrightarrow
H^q_\etale(\mathbf{P}^1_{\kappa(t)^{sep}},
\mathcal{F}|_{\mathbf{P}^1_{\kappa(t)^{sep}}})
$$
由于 $\kappa(s)^{sep} \subset \kappa(t)^{sep}$，由
Lemma \ref{etale-cohomology-lemma-base-change-dim-1-separably-closed}，该映射是同构。

\medskip\noindent
这证明了对于 $\mathbf{P}^1_S \to S$ 和
$\mathbf{Z}/\ell\mathbf{Z}$-模层，上同调与基变换可交换。
特别地，对内射 $\mathbf{Z}/\ell\mathbf{Z}$-模层，
任意基变换后的高阶正像都为零。换言之，
引理 \ref{etale-cohomology-lemma-proper-base-change-in-terms-of-injectives}
的条件 (2) 成立，证明完成。
\end{proof}

\begin{lemma}
\label{etale-cohomology-lemma-proper-base-change}
设 $f : X \to Y$ 是概形之间的真态射，$g : Y' \to Y$
是概形之间的态射。置 $X' = Y' \times_Y X$，并记
$f' : X' \to Y'$ 和 $g' : X' \to X$ 为投影。
设 $E \in D^+(X_\etale)$ 的上同调层都是挠层。
则基变换映射 (\ref{etale-cohomology-equation-base-change})
$g^{-1}Rf_*E \to Rf'_*(g')^{-1}E$ 是同构。
\end{lemma}

\begin{proof}
这是利用下列谱序列从真基变换定理
（Theorem \ref{etale-cohomology-theorem-proper-base-change}）直接得到的推论：
$$
E_2^{p, q} = R^pf_*H^q(E)
\quad\text{且}\quad
{E'}_2^{p, q} = R^pf'_*(g')^{-1}H^q(E)
$$
它们分别收敛于 $R^nf_*E$ 和 $R^nf'_*(g')^{-1}E$。
这些谱序列在 Derived Categories, Lemma
\ref{derived-lemma-two-ss-complex-functor} 中构造。略去若干细节。
\end{proof}

\begin{lemma}
\label{etale-cohomology-lemma-proper-base-change-stalk}
设 $f : X \to Y$ 是概形之间的真态射，并设 $\overline{y} \to Y$
是几何点。
\begin{enumerate}
\item 对挠阿贝尔层 $\mathcal{F}$，若它定义在 $X_\etale$ 上，则有
$(R^nf_*\mathcal{F})_{\overline{y}} =
H^n_\etale(X_{\overline{y}}, \mathcal{F}_{\overline{y}})$.
\item 对上同调层均为挠层的 $E \in D^+(X_\etale)$，有
$(R^nf_*E)_{\overline{y}} =
H^n_\etale(X_{\overline{y}}, E|_{X_{\overline{y}}})$.
\end{enumerate}
\end{lemma}

\begin{proof}
在陈述中，$\mathcal{F}_{\overline{y}}$ 表示 $\mathcal{F}$
到概形论纤维 $X_{\overline{y}} = \overline{y} \times_Y X$ 的拉回。
由于沿 $\overline{y} \to Y$ 拉回给出 $\mathcal{F}$ 的茎，
本引理的第一条是 Theorem \ref{etale-cohomology-theorem-proper-base-change} 的特例；
第二条是 Lemma \ref{etale-cohomology-lemma-proper-base-change} 的特例。
\end{proof}






\section{真基变换的应用}
\label{etale-cohomology-section-applications-proper-base-change}

\noindent
本节讨论真基变换定理的一些颇为直接的推论。

\begin{lemma}
\label{etale-cohomology-lemma-base-change-separably-closed}
设 $K/k$ 是可分闭域的扩张。设 $X$ 是 $k$ 上的真概形。
设 $\mathcal{F}$ 是 $X_\etale$ 上的挠阿贝尔层。则映射
$H^q_\etale(X, \mathcal{F}) \to
H^q_\etale(X_K, \mathcal{F}|_{X_K})$ 对 $q \geq 0$ 是同构。
\end{lemma}

\begin{proof}
考察茎可知，这是 Theorem \ref{etale-cohomology-theorem-proper-base-change} 的特例。
\end{proof}

\begin{lemma}
\label{etale-cohomology-lemma-cohomological-dimension-proper}
设 $f : X \to Y$ 是概形之间的真态射，其所有纤维的维数均为 $\leq n$。
则对任意挠阿贝尔层 $\mathcal{F}$，若它定义在 $X_\etale$ 上，
都有 $R^qf_*\mathcal{F} = 0$，只要 $q > 2n$。
\end{lemma}

\begin{proof}
我们对 $n$ 归纳，同时为所有真态射证明此结论。

\medskip\noindent
若 $n = 0$，则 $f$ 是有限态射
(More on Morphisms，引理 \ref{more-morphisms-lemma-characterize-finite})
，而结论由 Proposition \ref{etale-cohomology-proposition-finite-higher-direct-image-zero} 成立。

\medskip\noindent
若 $n > 0$，则由 Lemma \ref{etale-cohomology-lemma-proper-base-change-stalk}，
只需证明 $H^i_\etale(X, \mathcal{F}) = 0$，其中 $i > 2n$，
$X$ 是真概形，满足 $\dim(X) \leq n$，且定义在代数闭域 $k$ 上，
而 $\mathcal{F}$ 是 $X$ 上的挠阿贝尔层。

\medskip\noindent
若 $n = 1$，结论由 Theorem \ref{etale-cohomology-theorem-vanishing-curves} 得出。
假设 $n > 1$。由 Proposition \ref{etale-cohomology-proposition-topological-invariance}，
可以用约化概形代替 $X$。令 $\nu : X^\nu \to X$ 为正规化。
这是满的双有理有限态射
（见 Varieties, Lemma \ref{varieties-lemma-normalization-locally-algebraic}），
因而在稠密开子集 $U \subset X$ 上是同构
（Morphisms, Lemma \ref{morphisms-lemma-birational-birational}）。
于是 $c : \mathcal{F} \to \nu_*\nu^{-1}\mathcal{F}$ 是单射
（因为 $\nu$ 满），并且在 $U$ 上是同构。记 $i : Z \to X$
为 $U$ 的补集的包含映射。由于 $U$ 在 $X$ 中稠密，有
$\dim(Z) < \dim(X) = n$。由
Proposition \ref{etale-cohomology-proposition-closed-immersion-pushforward}，有
$\Coker(c) = i_*\mathcal{G}$，其中 $\mathcal{G}$ 是
$Z_\etale$ 上的某个挠阿贝尔层。于是
$H^q_\etale(X, \Coker(c)) = H^q_\etale(Z, \mathcal{G})$
（由 Proposition \ref{etale-cohomology-proposition-finite-higher-direct-image-zero}
和 Leray 谱序列）；归纳假设说明 $c$ 的余核只在次数
$\leq 2(n - 1)$ 上有非零上同调。因此只需对
$\nu_*\nu^{-1}\mathcal{F}$ 证明结论。由于 $\nu$ 有限，
这归结为证明 $H^i_\etale(X^\nu, \nu^{-1}\mathcal{F})$
在 $i > 2n$ 时为零。下一段处理这种情形。

\medskip\noindent
假设 $X$ 是 $k$ 上维数为 $n$ 的整正规真概形。
选取非常值有理函数 $f$；它定义在 $X$ 上。子概形
$X' \subset X \times \mathbf{P}^1_k$ 是 $f$ 的图，并位于下图中：
$$
X \xleftarrow{b} X' \xrightarrow{f} \mathbf{P}^1_k
$$
注意 $b$ 在开子概形 $U \subset X$ 上是同构；其补集是闭子概形
$Z \subset X$，其余维为 $\geq 2$。具体地，$U$ 是 $f$ 的定义域，
它包含所有余维为 $1$ 且位于 $X$ 中的点，见
Morphisms，引理 \ref{morphisms-lemma-rational-map-from-reduced-to-separated}
和 \ref{morphisms-lemma-extend-across}
（并结合 Serre 正规性判据，见
Properties, Lemma \ref{properties-lemma-criterion-normal}）。
此外，$b$ 的纤维维数为 $\leq 1$（因为它们是 $\mathbf{P}^1$ 的闭子概形）。
因此，由归纳假设，$R^ib_*b^{-1}\mathcal{F}$ 仅当
$i \in \{0, 1, 2\}$ 时才可能非零。选取可区别三角
$$
\mathcal{F} \to Rb_*b^{-1}\mathcal{F} \to Q \to \mathcal{F}[1]
$$
像前面一样，利用 $\mathcal{F} \to b_*b^{-1}\mathcal{F}$ 是单射这一事实，
并结合刚才的结论，可知 $Q$ 仅在次数 $0, 1, 2$ 上有非零上同调层，
且这些层支撑在 $Z$ 上。此外，由 Lemma \ref{etale-cohomology-lemma-torsion-direct-image}，
这些上同调层都是挠层。利用归纳假设和
Derived Categories，引理 \ref{derived-lemma-two-ss-complex-functor}
的谱序列，可知 $H^i(X, Q)$ 在
$i > 2 + 2\dim(Z) \leq 2 + 2(n - 2) = 2n - 2$ 时为零。
因此只需证明 $H^i(X', b^{-1}\mathcal{F}) = 0$ 在
$i > 2n$ 时成立。此时使用态射
$$
f : X' \to \mathbf{P}^1_k
$$
其纤维维数为 $< n$。因此，由归纳假设，
$R^if_*b^{-1}\mathcal{F} = 0$ 在 $i > 2(n - 1)$ 时成立。
最后应用 Leray 谱序列
$$
H^i(\mathbf{P}^1_k, R^jf_*b^{-1}\mathcal{F})
\Rightarrow
H^{i + j}(X', b^{-1}\mathcal{F})
$$
以及事实 $\dim(\mathbf{P}^1_k) = 1$ 即得结论。
\end{proof}

\noindent
使用模 $n$ 系数时，可以对无界复形作真基变换。

\begin{lemma}
\label{etale-cohomology-lemma-proper-base-change-mod-n}
设 $f : X \to Y$ 是概形之间的真态射，$g : Y' \to Y$
是概形之间的态射。置 $X' = Y' \times_Y X$，并记
$f' : X' \to Y'$ 和 $g' : X' \to X$ 为投影。
设 $n \geq 1$ 是整数。设 $E \in D(X_\etale, \mathbf{Z}/n\mathbf{Z})$。
则基变换映射 (\ref{etale-cohomology-equation-base-change})
$g^{-1}Rf_*E \to Rf'_*(g')^{-1}E$ 是同构。
\end{lemma}

\begin{proof}
只需在 $Y$ 和 $Y'$ 都拟紧时证明。由 Morphisms, Lemma
\ref{morphisms-lemma-morphism-finite-type-bounded-dimension}，
可知 $f : X \to Y$ 和 $f' : X' \to Y'$ 的纤维维数有界。
因此 Lemma \ref{etale-cohomology-lemma-cohomological-dimension-proper} 说明
$$
f_* :
\textit{Mod}(X_\etale, \mathbf{Z}/n\mathbf{Z})
\longrightarrow
\textit{Mod}(Y_\etale, \mathbf{Z}/n\mathbf{Z})
$$
以及
$$
f'_* :
\textit{Mod}(X'_\etale, \mathbf{Z}/n\mathbf{Z})
\longrightarrow
\textit{Mod}(Y'_\etale, \mathbf{Z}/n\mathbf{Z})
$$
在 Derived Categories, Lemma
\ref{derived-lemma-unbounded-right-derived} 的意义下具有有限上同调维数。
选取 K-内射复形 $\mathcal{I}^\bullet$，它由
$\mathbf{Z}/n\mathbf{Z}$-模层组成；其每一项 $\mathcal{I}^n$ 都是内射
$\mathbf{Z}/n\mathbf{Z}$-模层，并且它表示 $E$。见 Injectives, Theorem
\ref{injectives-theorem-K-injective-embedding-grothendieck}。
由通常的真基变换定理，有
$R^qf'_*(g')^{-1}\mathcal{I}^n = 0$，只要 $q > 0$，见
Theorem \ref{etale-cohomology-theorem-proper-base-change}。因此，由
Derived Categories, Lemma \ref{derived-lemma-unbounded-right-derived}，
可以计算 $Rf'_*(g')^{-1}E$，所用复形为
$f'_*(g')^{-1}\mathcal{I}^\bullet$。再次应用通常的真基变换定理，
可知它等于 $g^{-1}f_*\mathcal{I}^\bullet$，这正是所需结论。
\end{proof}

\begin{lemma}
\label{etale-cohomology-lemma-pull-out-constant}
设 $X$ 是拟紧拟分离概形。设 $E \in D^+(X_\etale)$ 且
$K \in D^+(\mathbf{Z})$。则
$$
R\Gamma(X, E \otimes_\mathbf{Z}^\mathbf{L} \underline{K}) =
R\Gamma(X, E) \otimes_\mathbf{Z}^\mathbf{L} K
$$
\end{lemma}

\begin{proof}
假设 $H^i(E) = 0$ 对 $i \geq a$ 成立，且 $H^j(K) = 0$
对 $j \geq b$ 成立。可以表示 $K$，所用的是下有界复形
$K^\bullet$，其各项为无挠 $\mathbf{Z}$-模。（选取 K-平坦复形
$L^\bullet$ 表示 $K$，
再取 $K^\bullet = \tau_{\geq b - 1}L^\bullet$。这是可行的，因为
$\mathbf{Z}$ 的整体维数为 $1$。见 More on Algebra, Lemma
\ref{more-algebra-lemma-last-one-flat}。）可以表示 $E$，所用的是下有界复形
$\mathcal{E}^\bullet$。于是
$E \otimes_\mathbf{Z}^\mathbf{L} \underline{K}$ 由下列复形表示：
$$
\text{Tot}(\mathcal{E}^\bullet \otimes_\mathbf{Z} \underline{K}^\bullet)
$$
利用可区别三角
$$
\sigma_{\geq -b + n + 1}K^\bullet
\to K^\bullet \to
\sigma_{\leq -b + n}K^\bullet
$$
以及显然的消失
$$
H^n(X,
\text{Tot}(\mathcal{E}^\bullet \otimes_\mathbf{Z}
\sigma_{\geq -a + n + 1}\underline{K}^\bullet) = 0
$$
和
$$
H^n(R\Gamma(X, E)
\otimes_\mathbf{Z}^\mathbf{L} \sigma_{\geq -a + n + 1}K^\bullet) = 0
$$
可归结为 $K^\bullet$ 是由平坦 $\mathbf{Z}$-模组成的有界复形的情形。
重复此论证，又归结为 $K^\bullet$ 等于位于某个次数上的单个平坦
$\mathbf{Z}$-模的情形。接着，对 $\mathcal{E}^\bullet$ 使用笨截断，
以完全相同的方式归结为 $\mathcal{E}^\bullet$ 是位于某个次数上的
单个阿贝尔层的情形。因此只需证明
$$
H^n(X, \mathcal{E} \otimes_\mathbf{Z} \underline{M}) =
H^n(X, \mathcal{E}) \otimes_\mathbf{Z} M
$$
其中 $M$ 是平坦 $\mathbf{Z}$-模，而 $\mathcal{E}$ 是 $X$ 上的阿贝尔层。
在这种情形下，把 $M$ 写成有限自由 $\mathbf{Z}$-模的滤过余极限
（Lazard 定理，见 Algebra, Theorem \ref{algebra-theorem-lazard}）。
由 Theorem \ref{etale-cohomology-theorem-colimit}，这归结为有限自由
$\mathbf{Z}$-模 $M$ 的情形，此时结论显然成立。
\end{proof}

\begin{lemma}
\label{etale-cohomology-lemma-projection-formula-proper}
设 $f : X \to Y$ 是概形之间的真态射。
设 $E \in D^+(X_\etale)$ 的上同调层都是挠层。
设 $K \in D^+(Y_\etale)$。则
$$
Rf_*E \otimes_\mathbf{Z}^\mathbf{L} K =
Rf_*(E \otimes_\mathbf{Z}^\mathbf{L} f^{-1}K)
$$
该等式成立于 $D^+(Y_\etale)$ 中。
\end{lemma}

\begin{proof}
由 Cohomology on Sites, Section
\ref{sites-cohomology-section-projection-formula}，存在从左端到右端的典范映射。
我们在茎上检验该等式。回顾，导出张量积的计算与拉回可交换。见
Cohomology on Sites，引理
\ref{sites-cohomology-lemma-pullback-tensor-product}。因此有
$$
(E \otimes_\mathbf{Z}^\mathbf{L} f^{-1}K)_{\overline{x}}
=
E_{\overline{x}} \otimes_\mathbf{Z}^\mathbf{L} K_{\overline{y}}
$$
其中 $\overline{y}$ 是 $\overline{x}$ 在 $Y$ 中的像。
由于 $\mathbf{Z}$ 的整体维数为 $1$，可知该复形在次数
$< - 1 + a + b$ 上的上同调为零，只要
$H^i(E) = 0$ 对 $i \geq a$ 成立且 $H^j(K) = 0$
对 $j \geq b$ 成立。此外，由于 $H^i(E)$ 对每个 $i$
都是挠阿贝尔层，复形
$E \otimes_\mathbf{Z}^\mathbf{L} K$ 的上同调层也是如此。事实上有
$$
(E \otimes_\mathbf{Z}^\mathbf{L} f^{-1}K)
\otimes_{\mathbf{Z}}^\mathbf{L} \mathbf{Q} =
(E \otimes_\mathbf{Z}^\mathbf{L} \mathbf{Q})
\otimes_{\mathbf{Q}}^\mathbf{L}
(f^{-1}K \otimes_{\mathbf{Z}}^\mathbf{L} \mathbf{Q})
$$
它在导出范畴中为零。这样，Lemma \ref{etale-cohomology-lemma-proper-base-change-stalk}
可应用于两端，从而只需证明
$$
R\Gamma(X_{\overline{y}},
E|_{X_{\overline{y}}} \otimes_\mathbf{Z}^\mathbf{L}
(X_{\overline{y}} \to \overline{y})^{-1}K_{\overline{y}}) =
R\Gamma(X_{\overline{y}},
E|_{X_{\overline{y}}}) \otimes_\mathbf{Z}^\mathbf{L} K_{\overline{y}}
$$
这已在 Lemma \ref{etale-cohomology-lemma-pull-out-constant} 中证明。
\end{proof}








\section{局部无环性}
\label{etale-cohomology-section-local-acyclicity}

\noindent
本节从光滑基变换定理推出光滑态射的局部无环性。在 SGA 4
或 SGA 4.5 中，作者先证明光滑态射的一个局部无环性版本，
再由此推出光滑基变换定理。

\medskip\noindent
我们将采用 Deligne \cite[Definition 2.12, page 242]{SGA4.5}
给出的局部无环性表述。设 $f : X \to S$ 是概形态射。
设 $\overline{x}$ 是 $X$ 的几何点，其像
$\overline{s} = f(\overline{x})$ 是 $S$ 的点。设 $\overline{t}$ 是
$\Spec(\mathcal{O}^{sh}_{S, \overline{s}})$ 的几何点。
我们得到交换图
$$
\xymatrix{
F_{\overline{x}, \overline{t}} =
\overline{t} \times_{\Spec(\mathcal{O}^{sh}_{S, \overline{s}})}
\Spec(\mathcal{O}^{sh}_{X, \overline{x}}) \ar[r] \ar[d] &
\Spec(\mathcal{O}^{sh}_{X, \overline{x}}) \ar[r] \ar[d] &
X \ar[d] \\
\overline{t} \ar[r] &
\Spec(\mathcal{O}^{sh}_{S, \overline{s}}) \ar[r] &
S
}
$$
概形 $F_{\overline{x}, \overline{t}}$ 称为
{\it $f$ 在 $\overline{x}$ 处的一个消失圈簇}。
设 $K$ 是 $D(X_\etale)$ 的对象。对任意概形态射
$g : Y\to X$，我们以 $R\Gamma(Y, K)$ 记
$R\Gamma(Y_\etale, g_{small}^{-1}K)$。由于
$\mathcal{O}^{sh}_{X, \overline{x}}$ 是严格 Hensel 的，故有
$K_{\overline{x}} = R\Gamma(\Spec(\mathcal{O}^{sh}_{X, \overline{x}}), K)$.
因此得到典范映射
\begin{equation}
\label{etale-cohomology-equation-alpha-K}
\alpha_{K, \overline{x}, \overline{t}} :
K_{\overline{x}}
\longrightarrow
R\Gamma(F_{\overline{x}, \overline{t}}, K)
\end{equation}
它由沿
$F_{\overline{x}, \overline{t}} \to \Spec(\mathcal{O}^{sh}_{X, \overline{x}})$
拉回上同调得到。

\begin{definition}
\label{etale-cohomology-definition-locally-acyclic}
\begin{reference}
\cite[Definition 2.12, page 242]{SGA4.5} 及
\cite[Definition (1.3), page 54]{SGA4.5}
\end{reference}
设 $f : X \to S$ 是概形态射。
设 $K$ 是 $D(X_\etale)$ 的对象。
\begin{enumerate}
\item
设 $\overline{x}$ 是 $X$ 的几何点，其像为
$\overline{s} = f(\overline{x})$。
称 $f$ {\it 在 $\overline{x}$ 处相对于 $K$ 局部无环}，如果对每个几何点
$\overline{t}$（它是
$\Spec(\mathcal{O}^{sh}_{S, \overline{s}})$ 的几何点），映射 (\ref{etale-cohomology-equation-alpha-K}) 都是同构\footnote{我们不要求
$\overline{t}$ 是
$\Spec(\mathcal{O}^{sh}_{S, \overline{s}})$ 的代数几何点。通常利用
引理 \ref{etale-cohomology-lemma-smooth-base-change-separably-closed}
可约化到这种情形。}。
\item 称 $f$ {\it 相对于 $K$ 局部无环}，如果 $f$ 在 $\overline{x}$ 处
相对于 $K$ 局部无环，其中 $\overline{x}$ 取遍 $X$ 的所有几何点。
\item 称 $f$ {\it 相对于 $K$ 普遍局部无环}，如果对任意概形态射
$S' \to S$，基变换 $f' : X' \to S'$ 都相对于 $K$ 在 $X'$ 上的拉回局部无环。
\item 称 $f$ {\it 局部无环}，如果令 $\overline{x}$ 取遍 $X$ 的所有几何点，
并令 $n$ 取遍与 $\kappa(\overline{x})$ 的特征互素的所有整数，
态射 $f$ 都在 $\overline{x}$ 处相对于取值为
$\mathbf{Z}/n\mathbf{Z}$ 的常值层局部无环。
\item 称 $f$ {\it 普遍局部无环}，如果对任意概形态射
$S' \to S$，基变换 $f' : X' \to S'$ 都局部无环。
\end{enumerate}
\end{definition}

\noindent
设 $M$ 是阿贝尔群。那么 $f : X \to S$ 相对于常值层
$\underline{M}$ 的局部无环性归结为要求
$$
H^q(F_{\overline{x}, \overline{t}}, \underline{M}) =
\left\{
\begin{matrix}
M & \text{若} & q = 0 \\
0 & \text{若} & q \not = 0
\end{matrix}
\right.
$$
对任意几何点 $\overline{x}$（它是 $X$ 的几何点）以及任意几何点
$\overline{t}$（它是 $\Spec(\mathcal{O}^{sh}_{S, f(\overline{x})})$ 的几何点）
都成立。由此可见，局部无环等价于几何纤维
$F_{\overline{x}, \overline{t}}$ 的高阶上同调群消失；这些纤维来自
$\overline{x}$ 和 $\overline{s}$ 处严格 Hensel 化之间的映射。

\begin{proposition}
\label{etale-cohomology-proposition-smooth-locally-acyclic}
设 $f : X \to S$ 是光滑概形态射。则 $f$ 普遍局部无环。
\end{proposition}

\begin{proof}
由于光滑态射的基变换仍光滑，只须证明光滑态射局部无环。
设 $\overline{x}$ 是 $X$ 的几何点，其像为
$\overline{s} = f(\overline{x})$。设
$\overline{t}$ 是
$\Spec(\mathcal{O}^{sh}_{S, f(\overline{x})})$ 的几何点。
由于我们要证明的是环同态
$\mathcal{O}^{sh}_{S, \overline{s}} \to \mathcal{O}^{sh}_{X, \overline{x}}$
的一个性质（见 Definition \ref{etale-cohomology-definition-locally-acyclic} 后的讨论），
所以可且确实把 $f : X \to S$ 替换为基变换
$X \times_S \Spec(\mathcal{O}^{sh}_{S, \overline{s}}) \to
\Spec(\mathcal{O}^{sh}_{S, \overline{s}})$.
因此可且确实假设 $S$ 是严格 Hensel 局部环的谱，并且
$\overline{s}$ 位于 $S$ 的闭点之上。

\medskip\noindent
我们将把 Lemma \ref{etale-cohomology-lemma-base-change-f-star-general-stalks}
应用于图
$$
\xymatrix{
X \ar[d]_f & X_{\overline{t}} \ar[l]^h \ar[d]^e \\
S & \overline{t} \ar[l]_g
}
$$
以及层 $\mathcal{F} = \underline{M}$，其中 $M = \mathbf{Z}/n\mathbf{Z}$，
$n$ 是与 $\overline{x}$ 的剩余域特征互素的某个整数。由光滑基变换可知，映射
$f^{-1}R^qg_*\mathcal{F} \to R^qh_*e^{-1}\mathcal{F}$
是同构，见 Theorem \ref{etale-cohomology-theorem-smooth-base-change}（由 $n$ 的选择，
关于挠性的假设成立）。因此 Lemma \ref{etale-cohomology-lemma-base-change-f-star-general-stalks}
给出下式中间的等号：
$$
H^q(F_{\overline{x}, \overline{t}}, \underline{M}) =
H^q(\Spec(\mathcal{O}^{sh}_{X, \overline{x}}) \times_S \overline{t},
\underline{M})
=
H^q(\Spec(\mathcal{O}^{sh}_{S, \overline{s}}) \times_S \overline{t},
\underline{M}) =
H^q(\overline{t}, \underline{M})
$$
外侧的两个等号用到了
$S = \Spec(\mathcal{O}^{sh}_{S, \overline{s}})$。
由于 $\overline{t}$ 是可分闭域的谱，故有
$$
H^q(F_{\overline{x}, \overline{t}}, \underline{M}) =
\left\{
\begin{matrix}
M & \text{若} & q = 0 \\
0 & \text{若} & q \not = 0
\end{matrix}
\right.
$$
这正是所要证明的结论（见 Definition \ref{etale-cohomology-definition-locally-acyclic}
后的讨论）。
\end{proof}

\begin{lemma}
\label{etale-cohomology-lemma-locally-acyclic-locally-constant}
设 $f : X \to S$ 是概形态射。设 $\mathcal{F}$ 是 $X_\etale$ 上的
局部常值阿贝尔层，使得每当 $\overline{x}$ 是 $X$ 的几何点时，阿贝尔群
$\mathcal{F}_{\overline{x}}$ 都是挠群，且每个元素的阶均与
$\overline{x}$ 的剩余域特征互素。如果 $f$ 局部无环，则 $f$
相对于 $\mathcal{F}$ 局部无环。
\end{lemma}

\begin{proof}
事实上，设 $\overline{x}$ 是 $X$ 的几何点。由于 $\mathcal{F}$ 局部常值，
故 $\mathcal{F}$ 在 $\Spec(\mathcal{O}^{sh}_{X, \overline{x}})$ 上的限制
同构于常值层 $\underline{M}$，其中 $M = \mathcal{F}_{\overline{x}}$。
由假设，可把 $M = \colim M_i$ 写成有限阿贝尔群 $M_i$ 的滤过余极限，
并且它们的阶均与 $\overline{x}$ 的剩余域特征互素。取一个几何点
$\overline{t}$（它是 $\Spec(\mathcal{O}^{sh}_{S, f(\overline{x})})$ 的几何点）。
由于 $F_{\overline{x}, \overline{t}}$ 是仿射的，
由 Lemma \ref{etale-cohomology-lemma-colimit} 有
$$
H^q(F_{\overline{x}, \overline{t}}, \underline{M}) =
\colim H^q(F_{\overline{x}, \overline{t}}, \underline{M_i})
$$
对每个 $i$，可把 $M_i = \bigoplus \mathbf{Z}/n_{i, j}\mathbf{Z}$
写成有限直和，其中整数 $n_{i, j}$ 均与 $\overline{x}$ 的剩余域特征互素。
由于 $f$ 局部无环，故
$$
H^q(F_{\overline{x}, \overline{t}},
\underline{\mathbf{Z}/n_{i, j}\mathbf{Z}}) =
\left\{
\begin{matrix}
\mathbf{Z}/n_{i, j}\mathbf{Z} & \text{若} & q = 0 \\
0 & \text{若} & q \not = 0
\end{matrix}
\right.
$$
见 Definition \ref{etale-cohomology-definition-locally-acyclic} 后的讨论。
取直和与余极限可得
$$
H^q(F_{\overline{x}, \overline{t}}, \underline{M}) =
\left\{
\begin{matrix}
M & \text{若} & q = 0 \\
0 & \text{若} & q \not = 0
\end{matrix}
\right.
$$
结论得证。
\end{proof}

\begin{lemma}
\label{etale-cohomology-lemma-locally-acyclic-quasi-finite-base-change}
设
$$
\xymatrix{
X' \ar[r]_{g'} \ar[d]_{f'} & X \ar[d]^f \\
S' \ar[r]^g & S
}
$$
是概形的笛卡儿图。设 $K$ 是 $D(X_\etale)$ 的对象。
设 $\overline{x}'$ 是 $X'$ 的几何点，其像 $\overline{x}$ 是 $X$ 的几何点。如果
\begin{enumerate}
\item $f$ 在 $\overline{x}$ 处相对于 $K$ 局部无环；且
\item $g$ 局部拟有限，或者 $S' = \lim S_i$ 是在 $S$ 上局部拟有限、
过渡态射为仿射态射的概形的有向逆极限，或者 $g : S' \to S$ 是整态射，
\end{enumerate}
则 $f'$ 在 $\overline{x}'$ 处相对于 $(g')^{-1}K$ 局部无环。
\end{lemma}

\begin{proof}
以 $\overline{s}'$ 和 $\overline{s}$ 分别表示 $\overline{x}'$ 和
$\overline{x}$ 在 $S'$ 和 $S$ 中的像。设 $\overline{t}'$ 是
$\Spec(\mathcal{O}^{sh}_{S', \overline{s}'})$ 的几何点，并以
$\overline{t}$ 表示它在 $\Spec(\mathcal{O}^{sh}_{S, \overline{s}})$ 中的像。
由 Algebra, Lemma
\ref{algebra-lemma-base-change-strict-henselization-quasi-finite}
以及对 $g$ 的假设，下列同态是同构：
$$
\mathcal{O}^{sh}_{X, \overline{x}}
\otimes_{\mathcal{O}^{sh}_{S, \overline{s}}}
\mathcal{O}^{sh}_{S', \overline{s}'}
\longrightarrow
\mathcal{O}^{sh}_{X', \overline{x}'}
$$
由于按照约定
$\kappa(\overline{t}) = \kappa(\overline{t}')$，可得
$$
F_{\overline{x}', \overline{t}'} =
\Spec\left(
\mathcal{O}^{sh}_{X', \overline{x}'}
\otimes_{\mathcal{O}^{sh}_{S', \overline{s}'}}
\kappa(\overline{t}')\right) =
\Spec\left(
\mathcal{O}^{sh}_{X, \overline{x}}
\otimes_{\mathcal{O}^{sh}_{S, \overline{s}}}
\kappa(\overline{t})\right) =
F_{\overline{x}, \overline{t}}
$$
换言之，$f'$ 在 $\overline{x}'$ 处的消失圈簇是
$f$ 在 $\overline{x}$ 处的消失圈簇的实例。本引理由此和定义立即得到。
\end{proof}










\section{反特殊化映射}
\label{etale-cohomology-section-cospecialization}

\noindent
设 $f : X \to S$ 是概形态射。设 $\overline{x}$ 是 $X$ 的几何点，
其像 $\overline{s} = f(\overline{x})$ 是 $S$ 的点。设 $\overline{t}$ 是
$\Spec(\mathcal{O}^{sh}_{S, \overline{s}})$ 的几何点。
设 $K \in D(X_\etale)$。对任意概形态射 $g : Y \to X$，
我们以 $K|_Y$ 记 $g_{small}^{-1}K$，并以 $R\Gamma(Y, K)$
记 $R\Gamma(Y_\etale, g_{small}^{-1}K)$。我们断言，如果
\begin{enumerate}
\item $K$ 下有界，即 $K \in D^+(X_\etale)$；
\item $f$ 相对于 $K$ 局部无环，
\end{enumerate}
那么存在一个{\it 反特殊化映射}
$$
cosp :
R\Gamma(X_{\overline{t}}, K)
\longrightarrow
R\Gamma(X_{\overline{s}}, K)
$$
它与 Section \ref{etale-cohomology-section-specialization}，尤其是
Remark \ref{etale-cohomology-remark-specialization-map-and-fibres} 中考察的特殊化映射密切相关。

\medskip\noindent
为构造该映射，考察态射
$$
X_{\overline{t}}
\xrightarrow{h}
X \times_S \Spec(\mathcal{O}^{sh}_{S, \overline{s}})
\xleftarrow{i}
X_{\overline{s}}
$$
$h^{-1}$ 与 $Rh_*$ 之间伴随的单位给出映射
$$
\beta_{K, \overline{s}, \overline{t}} :
K|_{X \times_S \Spec(\mathcal{O}^{sh}_{S, \overline{s}})}
\longrightarrow
Rh_*(K|_{X_{\overline{t}}})
$$
它位于 $D((X \times_S \Spec(\mathcal{O}^{sh}_{S, \overline{s}}))_\etale)$ 中。
下面的 Lemma \ref{etale-cohomology-lemma-beta-is-isomorphism} 表明，在上述假设下，拉回
$i^{-1}\beta_{K, \overline{s}, \overline{t}}$ 是同构。
因此可把反特殊化映射定义为复合
\begin{align*}
R\Gamma(X_{\overline{t}}, K)
& =
R\Gamma(X \times_S \Spec(\mathcal{O}^{sh}_{S, \overline{s}}),
Rh_*(K|_{X_{\overline{t}}})) \\
&
\xrightarrow{i^{-1}}
R\Gamma(X_{\overline{s}}, i^{-1}Rh_*(K|_{X_{\overline{t}}})) \\
&
\xrightarrow{(i^{-1}\beta_{K, \overline{s}, \overline{t}})^{-1}}
R\Gamma(X_{\overline{s}},
i^{-1}(K|_{X \times_S \Spec(\mathcal{O}^{sh}_{S, \overline{s}})})) \\
& =
R\Gamma(X_{\overline{s}}, K)
\end{align*}

\begin{lemma}
\label{etale-cohomology-lemma-beta-is-isomorphism}
映射 $i^{-1}\beta_{K, \overline{s}, \overline{t}}$ 是同构。
\end{lemma}

\begin{proof}
映射 $h$、$i$、$\beta_{K, \overline{s}, \overline{t}}$ 的构造只依赖于
$X$ 和 $K$ 到 $\Spec(\mathcal{O}^{sh}_{S, \overline{s}})$ 的基变换。
因此可且确实假设 $S$ 是闭点为 $\overline{s}$ 的严格 Hensel 概形。
注意，$f$ 相对于 $K$ 的局部无环性在此基变换下保持（例如可用
Lemma \ref{etale-cohomology-lemma-locally-acyclic-quasi-finite-base-change}，也可在这个非常特殊的情形中
直接比较严格 Hensel 环）。

\medskip\noindent
设 $\overline{x}$ 是 $X_{\overline{s}}$ 的几何点。等价地，设
$\overline{x}$ 是在 $f$ 下的像为 $\overline{s}$ 的几何点。下面计算
$i^{-1}\beta_{K, \overline{s}, \overline{t}}$ 在 $\overline{x}$ 处的茎。
首先有
$$
(i^{-1}\beta_{K, \overline{s}, \overline{t}})_{\overline{x}} =
(\beta_{K, \overline{s}, \overline{t}})_{\overline{x}}
$$
因为拉回保持茎，见 Lemma \ref{etale-cohomology-lemma-stalk-pullback}。
由于处在 $S = \Spec(\mathcal{O}^{sh}_{S, \overline{s}})$ 的情形，
故 $h : X_{\overline{t}} \to X$ 满足
$X_{\overline{t}} \times_X \Spec(\mathcal{O}^{sh}_{X, \overline{x}})
= F_{\overline{x}, \overline{t}}$。因此
$$
(\beta_{K, \overline{s}, \overline{t}})_{\overline{x}} :
K_{\overline{x}}
\longrightarrow
Rh_*(K|_{X_{\overline{t}}})_{\overline{x}} =
R\Gamma(F_{\overline{x}, \overline{t}}, K)
$$
其中等号来自 Theorem \ref{etale-cohomology-theorem-higher-direct-images}。
于是映射 $(\beta_{K, \overline{s}, \overline{t}})_{\overline{x}}$
正是 Definition \ref{etale-cohomology-definition-locally-acyclic} 中使用的映射
$\alpha_{K, \overline{x}, \overline{t}}$。由 Theorem \ref{etale-cohomology-theorem-exactness-stalks}，
可在茎上检验一个映射是否为同构，故结论成立。
\end{proof}

\noindent
反特殊化映射一旦存在，便意在成为特殊化映射的逆。

\begin{lemma}
\label{etale-cohomology-lemma-specialization-cospecialization}
在上述情形中，如果还假设 $f$ 拟紧且拟分离，则图
$$
\xymatrix{
(Rf_*K)_{\overline{s}} \ar[r] \ar[d]_{sp} &
R\Gamma(X_{\overline{s}}, K) \\
(Rf_*K)_{\overline{t}} \ar[r] &
R\Gamma(X_{\overline{t}}, K) \ar[u]_{cosp}
}
$$
交换。
\end{lemma}

\begin{proof}
与 Lemma \ref{etale-cohomology-lemma-beta-is-isomorphism} 的证明相同，
可把 $S$ 替换为 $\Spec(\mathcal{O}^{sh}_{S, \overline{s}})$。
这时映射简化为 $h : X_{\overline{t}} \to X$、
$i : X_{\overline{s}} \to X$ 以及
$\beta_{K, \overline{s}, \overline{t}} : K \to Rh_*(K|_{X_{\overline{t}}})$。
由 Theorem \ref{etale-cohomology-theorem-higher-direct-images}，
$(Rf_*K)_{\overline{s}} = R\Gamma(X, K)$，故 $sp$ 与基变换映射
$(Rf_*K)_{\overline{t}} \to R\Gamma(X_{\overline{t}}, K)$
的复合就是沿 $h$ 拉回上同调。这等于映射
$$
R\Gamma(X, K) \xrightarrow{\beta_{K, \overline{s}, \overline{t}}}
R\Gamma(X, Rh_*(K|_{X_{\overline{t}}})) =
R\Gamma(X_{\overline{t}}, K)
$$
而映射 $cosp$ 先把此陈列公式中的等号 $=$ 反向使用，
再沿 $i$ 拉回，最后应用 $i^{-1}\beta_{K, \overline{s}, \overline{t}}$ 的逆。
因此得到所需的交换性。
\end{proof}

\begin{lemma}
\label{etale-cohomology-lemma-sp-isom-proper-torsion-loc-ac}
设 $f : X \to S$ 是概形态射。设 $K \in D(X_\etale)$。假设
\begin{enumerate}
\item $K$ 下有界，即 $K \in D^+(X_\etale)$；
\item $f$ 相对于 $K$ 局部无环；
\item $f$ 是固有态射；且
\item $K$ 的上同调层是挠层。
\end{enumerate}
则对每个几何点 $\overline{s}$（它是 $S$ 的几何点）以及每个几何点
$\overline{t}$（它是 $\Spec(\mathcal{O}^{sh}_{S, \overline{s}})$ 的几何点），
特殊化映射
$sp : (Rf_*K)_{\overline{s}} \to (Rf_*K)_{\overline{t}}$
与反特殊化映射
$cosp : R\Gamma(X_{\overline{t}}, K) \to R\Gamma(X_{\overline{s}}, K)$
都是同构。
\end{lemma}

\begin{proof}
由固有基变换定理（采用 Lemma \ref{etale-cohomology-lemma-proper-base-change-stalk} 的形式），有
$(Rf_*K)_{\overline{s}} = R\Gamma(X_{\overline{s}}, K)$，
对 $\overline{t}$ 也类似。“正确”的证明应说明
Lemma \ref{etale-cohomology-lemma-specialization-cospecialization} 中的论证表明，
此时 $sp$ 与 $cosp$ 是互逆同构。这里改为直接证明 $cosp$ 是同构。
由上面的讨论，$cosp$ 是同构当且仅当沿 $i$ 的拉回
$$
R\Gamma(X \times_S \Spec(\mathcal{O}^{sh}_{S, \overline{s}}),
Rh_*(K|_{X_{\overline{t}}}))
\longrightarrow
R\Gamma(X_{\overline{s}}, i^{-1}Rh_*(K|_{X_{\overline{t}}}))
$$
在 $D^+(\textit{Ab})$ 中是同构。这由固有态射
$f' : X \times_S \Spec(\mathcal{O}^{sh}_{S, \overline{s}}) \to
\Spec(\mathcal{O}^{sh}_{S, \overline{s}})$
沿态射 $\overline{s} \to \Spec(\mathcal{O}^{sh}_{S, \overline{s}})$
对复形 $K' = Rh_*(K|_{X_{\overline{t}}})$ 应用固有基变换定理得到。
由 Lemma \ref{etale-cohomology-lemma-torsion-direct-image}，复形 $K'$ 下有界且有挠上同调层。
由于 $\Spec(\mathcal{O}^{sh}_{S, \overline{s}})$ 严格 Hensel，
且 $\overline{s}$ 位于闭点之上，故上述箭头的源等于
$(Rf'_*K')_{\overline{s}}$，靶等于
$R\Gamma(X_{\overline{s}}, K')$，并且由已经使用过的
Lemma \ref{etale-cohomology-lemma-proper-base-change-stalk}，上述映射是同构。
于是 Lemma \ref{etale-cohomology-lemma-specialization-cospecialization} 的图中四个箭头有三个是同构，
结论随之成立。
\end{proof}

\begin{lemma}
\label{etale-cohomology-lemma-sp-isom-proper-loc-cst-torsion}
设 $f : X \to S$ 是概形态射。设 $\mathcal{F}$ 是 $X_\etale$ 上的阿贝尔层。假设
\begin{enumerate}
\item $f$ 光滑且固有；
\item $\mathcal{F}$ 局部常值；且
\item $\mathcal{F}_{\overline{x}}$ 是挠群，且每个元素的阶均与
$\overline{x}$ 的剩余特征互素，其中 $\overline{x}$ 取遍 $X$ 的所有几何点。
\end{enumerate}
则对每个几何点 $\overline{s}$（它是 $S$ 的几何点）以及每个几何点
$\overline{t}$（它是 $\Spec(\mathcal{O}^{sh}_{S, \overline{s}})$ 的几何点），
特殊化映射
$sp : (Rf_*\mathcal{F})_{\overline{s}} \to (Rf_*\mathcal{F})_{\overline{t}}$
是同构。
\end{lemma}

\begin{proof}
这由 Lemmas \ref{etale-cohomology-lemma-sp-isom-proper-torsion-loc-ac}、
\ref{etale-cohomology-lemma-locally-acyclic-locally-constant} 以及
Proposition \ref{etale-cohomology-proposition-smooth-locally-acyclic} 得到。
\end{proof}

















\section{上同调维数}
\label{etale-cohomology-section-cd}

\noindent
利用 Section \ref{etale-cohomology-section-vanishing-torsion} 中的结果，
再加上对固有基变换定理的一次看似无害的应用（见
Proposition \ref{etale-cohomology-proposition-cd-affine} 的证明），可以推出概形和域的
上同调维数的一些界。

\begin{definition}
\label{etale-cohomology-definition-cd}
设 $X$ 是拟紧拟分离概形。{\it $X$ 的上同调维数}是集合中最小的元素
$$
\text{cd}(X) \in \{0, 1, 2, \ldots\} \cup \{\infty\}
$$
，它满足：对任意阿贝尔挠层 $\mathcal{F}$（定义在 $X_\etale$ 上），
都有 $H^i_\etale(X, \mathcal{F}) = 0$，只要 $i > \text{cd}(X)$。
如果 $X = \Spec(A)$，有时也称它为 $A$ 的上同调维数。
\end{definition}

\noindent
如果概形具有特征 $p$，则对 $p$-幂挠层的上同调消失通常可以得到更尖锐的界。
我们将在别处讨论这一点（此处插入将来的参考文献）。

\begin{lemma}
\label{etale-cohomology-lemma-cd-limit}
设 $X = \lim X_i$ 是一族拟紧拟分离概形的有向逆极限，且过渡态射均为仿射态射。
则 $\text{cd}(X) \leq \max \text{cd}(X_i)$。
\end{lemma}

\begin{proof}
以 $f_i : X \to X_i$ 表示投影。设 $\mathcal{F}$ 是 $X_\etale$ 上的阿贝尔挠层。
由 Lemma \ref{etale-cohomology-lemma-linus-hamann}，有
$\mathcal{F} = \lim f_i^{-1}f_{i, *}\mathcal{F}$。
因此由 Theorem \ref{etale-cohomology-theorem-colimit}，
$H^q_\etale(X, \mathcal{F}) =
\colim H^q_\etale(X_i, f_{i, *}\mathcal{F})$。引理得证。
\end{proof}

\begin{lemma}
\label{etale-cohomology-lemma-cd-curve-over-field}
设 $K$ 是域。设 $X$ 是 $1$ 维仿射概形，且在 $K$ 上有限型。则
$\text{cd}(X) \leq 1 + \text{cd}(K)$。
\end{lemma}

\begin{proof}
设 $\mathcal{F}$ 是 $X_\etale$ 上的阿贝尔挠层。考察态射
$f : X \to \Spec(K)$ 的 Leray 谱序列，得到
$$
E_2^{p, q} = H^p(\Spec(K), R^qf_*\mathcal{F})
$$
，它收敛到 $H^{p + q}_\etale(X, \mathcal{F})$。
$R^qf_*\mathcal{F}$ 在几何点
$\Spec(\overline{K}) \to \Spec(K)$ 处的茎，是 $\mathcal{F}$ 到
$X_{\overline{K}}$ 的拉回之上同调。因此由
Theorem \ref{etale-cohomology-theorem-vanishing-affine-curves}，它在次数 $\geq 2$ 中消失。
\end{proof}

\begin{lemma}
\label{etale-cohomology-lemma-cd-field-extension}
设 $L/K$ 是域扩张。则
$\text{cd}(L) \leq \text{cd}(K) + \text{trdeg}_K(L)$。
\end{lemma}

\begin{proof}
如果 $\text{trdeg}_K(L) = \infty$，结论显然。否则可找到扩张序列
$L= L_r/L_{r - 1}/ \ldots /L_1/L_0 = K$，使得
$\text{trdeg}_{L_i}(L_{i + 1}) = 1$ 且 $r = \text{trdeg}_K(L)$。
因此只须证明 $r = 1$ 的情形。此时可把 $L = \colim A_i$
写成其有限型 $K$-子代数的滤过余极限。由 Lemma \ref{etale-cohomology-lemma-cd-limit}，
只须证明 $\text{cd}(A_i) \leq 1 + \text{cd}(K)$，这由
Lemma \ref{etale-cohomology-lemma-cd-curve-over-field} 得到。
\end{proof}

\begin{lemma}
\label{etale-cohomology-lemma-strictly-henselian}
设 $K$ 是域。设 $X$ 是 $K$ 上的有限型概形。设 $x \in X$。
置 $a = \text{trdeg}_K(\kappa(x))$ 且 $d = \dim_x(X)$。则存在映射
$$
K(t_1, \ldots, t_a)^{sep} \longrightarrow \mathcal{O}_{X, x}^{sh}
$$
使得
\begin{enumerate}
\item $\mathcal{O}_{X, x}^{sh}$ 的剩余域是
$K(t_1, \ldots, t_a)^{sep}$ 的纯不可分扩张；
\item $\mathcal{O}_{X, x}^{sh}$ 是一族有限型
$K(t_1, \ldots, t_a)^{sep}$-代数的滤过余极限，且这些代数的维数均为
$\leq d - a$。
\end{enumerate}
\end{lemma}

\begin{proof}
可假设 $X$ 是仿射的。由 Noether 正规化，在必要时再次缩小 $X$ 后，
可选取有限态射 $\pi : X \to \mathbf{A}^d_K$，见
Algebra, Lemma \ref{algebra-lemma-Noether-normalization-at-point}。
由于 $\kappa(x)$ 是 $\pi(x)$ 的剩余域的有限扩张，该剩余域的超越次数也为
$a$（相对于 $K$）。因此可找到有限态射
$\pi' : \mathbf{A}^d_K \to \mathbf{A}^d_K$，使得 $\pi'(\pi(x))$
对应于线性子空间 $\mathbf{A}^a_K \subset \mathbf{A}^d_K$ 的一般点；
该子空间由令最后 $d - a$ 个坐标为零给出。因此复合
$$
X \xrightarrow{\pi' \circ \pi} \mathbf{A}^d_K \xrightarrow{p} \mathbf{A}^a_K
$$
由 $\pi' \circ \pi$ 和投影 $p$ 组成，后者投到前 $a$ 个坐标；它把 $x$
映到一般点 $\eta \in \mathbf{A}^a_K$。所诱导的 \'etale 局部环映射
$$
K(t_1, \ldots, t_a)^{sep} =
\mathcal{O}_{\mathbf{A}^a_k, \eta}^{sh}
\longrightarrow \mathcal{O}_{X, x}^{sh}
$$
满足 (1)，因为 $\mathcal{O}_{X, x}^{sh}$ 的剩余域显然是可分闭域
$K(t_1, \ldots, t_a)^{sep}$ 的代数扩张。另一方面，如果 $X = \Spec(B)$，
则 $\mathcal{O}_{X, x}^{sh} = \colim B_j$ 是 \'etale $B$-代数 $B_j$ 的滤过余极限。
注意，$B_j$ 在 $K[t_1, \ldots, t_d]$ 上拟有限，因为 $B$ 在
$K[t_1, \ldots, t_d]$ 上有限。类似地，可把
$K(t_1, \ldots, t_a)^{sep} = \colim A_i$ 写成 \'etale
$K[t_1, \ldots, t_a]$-代数的滤过余极限。对每个 $i$，可找到 $j$，使得
$A_i \to K(t_1, \ldots, t_a)^{sep} \to \mathcal{O}_{X, x}^{sh}$
分解经过映射 $\psi_{i, j} : A_i \to B_j$。于是 $B_j$ 在
$A_i[t_{a + 1}, \ldots, t_d]$ 上拟有限。因此
$$
B_{i, j} = B_j \otimes_{\psi_{i, j}, A_i} K(t_1, \ldots, t_a)^{sep}
$$
的维数为 $\leq d - a$，因为它在
$K(t_1, \ldots, t_a)^{sep}[t_{a + 1}, \ldots, t_d]$ 上拟有限。
现在 (2) 得证，因为 $\mathcal{O}_{X, x}^{sh}$ 是滤过余极限\footnote{设
$R$ 是环。设 $A = \colim_{i \in I} A_i$ 是有限呈现
$R$-代数的滤过余极限。设 $B = \colim_{j \in J} B_j$ 是
$R$-代数的滤过余极限。设 $A \to B$ 是 $R$-代数映射。假设对所有
$i \in I$，存在 $j \in J$ 以及 $R$-代数映射
$\psi_{i, j} : A_i \to B_j$。若
$(i', j', \psi_{i', j'}) \geq (i, j, \psi_{i, j})$，意思是
$i' \geq i$、$j' \geq j$，并且 $\psi_{i, j}$ 与
$\psi_{i', j'}$ 相容，则这些三元组的集合构成有向集，并且
$B = \colim B_j \otimes_{\psi_{i, j} A_i} A$。}，其组成代数为
$B_{i, j}$。略去若干细节。
\end{proof}

\begin{proposition}
\label{etale-cohomology-proposition-cd-affine}
设 $K$ 是域。设 $X$ 是 $K$ 上有限型的仿射概形。则有
$\text{cd}(X) \leq \dim(X) + \text{cd}(K)$。
\end{proposition}

\begin{proof}
对 $\dim(X)$ 作归纳。设 $\mathcal{F}$ 是 $X_\etale$ 上的阿贝尔挠层。

\medskip\noindent
情形 $\dim(X) = 0$。此时结构态射 $f : X \to \Spec(K)$ 是有限的。
因此 $R^if_*\mathcal{F} = 0$，只要 $i > 0$，见
Proposition \ref{etale-cohomology-proposition-finite-higher-direct-image-zero}。
因此 $H^i_\etale(X, \mathcal{F}) = H^i_\etale(\Spec(K), f_*\mathcal{F})$；
这是由 $f$ 的 Leray 谱序列（Cohomology on Sites, Lemma
\ref{sites-cohomology-lemma-Leray}）得到的，
结论显然。

\medskip\noindent
情形 $\dim(X) = 1$。这就是 Lemma \ref{etale-cohomology-lemma-cd-curve-over-field}。

\medskip\noindent
假设 $d = \dim(X) > 1$，并且命题对域上维数 $< d$ 的有限型仿射概形成立。
由 Noether 正规化（例如见 Varieties, Lemma
\ref{varieties-lemma-noether-normalization-affine}），存在有限态射
$f : X \to \mathbf{A}^d_K$。回顾由
Proposition \ref{etale-cohomology-proposition-finite-higher-direct-image-zero}，
$R^if_*\mathcal{F} = 0$，只要 $i > 0$。由 $f$ 的 Leray 谱序列（Cohomology on Sites,
Lemma \ref{sites-cohomology-lemma-Leray}），只须对 $\pi_*\mathcal{F}$
在 $\mathbf{A}^d_K$ 上证明结论。

\medskip\noindent
插曲 I。设 $j : X \to Y$ 是光滑 $d$ 维簇的开浸入，这些簇定义在
$K$ 上（不要求仿射），其补集是某个有效 Cartier 除子 $D$ 的支集。
层 $R^qj_*\mathcal{F}$（其中 $q > 0$）支撑在 $D$ 上。我们断言
$(R^qj_*\mathcal{F})_{\overline{y}} = 0$，只要
$a = \text{trdeg}_K(\kappa(y)) > d - q$。事实上，由
Theorem \ref{etale-cohomology-theorem-higher-direct-images} 有
$$
(R^qj_*\mathcal{F})_{\overline{y}} =
H^q(\Spec(\mathcal{O}_{Y, y}^{sh}) \times_Y X, \mathcal{F})
$$
选取 $f \in \mathfrak m_y \subset \mathcal{O}_{Y, y}$ 作为 $D$ 的局部方程。
于是
$$
\Spec(\mathcal{O}_{Y, y}^{sh}) \times_Y X =
\Spec(\mathcal{O}_{Y, y}^{sh}[1/f])
$$
利用 Lemma \ref{etale-cohomology-lemma-strictly-henselian} 得到嵌入
$$
K(t_1, \ldots, t_a)^{sep}(x) =
K(t_1, \ldots, t_a)^{sep}[x]_{(x)}[1/x]
\longrightarrow
\mathcal{O}_{Y, y}^{sh}[1/f]
$$
由于在 $K$ 上，$\mathcal{O}_{Y, y}^{sh}$ 的分式域的超越次数为 $d$，
可见 $\mathcal{O}_{Y, y}^{sh}[1/f]$ 是一族有限型
$(d - a - 1)$ 维代数的滤过余极限；这些代数定义在域
$K(t_1, \ldots, t_a)^{sep}(x)$ 上，而该域由
Lemma \ref{etale-cohomology-lemma-cd-field-extension} 具有上同调维数 $1$。
因此由归纳假设和 Lemma \ref{etale-cohomology-lemma-cd-limit} 得到所需消失。

\medskip\noindent
插曲 II。设 $Z$ 是定义在 $K$ 上的光滑簇，其维数为 $d - 1$。设
$E_a \subset Z$ 是由满足条件的点 $z \in Z$ 构成的集合，该条件为
$\text{trdeg}_K(\kappa(z)) \leq a$。注意 $E_a$ 对特殊化封闭，见
Varieties, Lemma \ref{varieties-lemma-dimension-locally-algebraic}。
设 $\mathcal{G}$ 是 $Z$ 上的阿贝尔挠层，其支集包含在 $E_a$ 中。
我们断言 $H^b_\etale(Z, \mathcal{G}) = 0$，只要
$b > a + \text{cd}(K)$。事实上，可写
$\mathcal{G} = \colim \mathcal{G}_i$，其中 $\mathcal{G}_i$ 是支撑在闭子概形
$Z_i$ 上的阿贝尔挠层，该闭子概形包含在 $E_a$ 中，见
Lemma \ref{etale-cohomology-lemma-support-in-subset}。于是归纳假设给出
$\mathcal{G}_i$ 的所需消失\footnote{这里先用 Proposition
\ref{etale-cohomology-proposition-closed-immersion-pushforward}，把 $\mathcal{G}_i$ 写成
$Z_i$ 上某层的正像；归纳假设给出该层在 $Z_i$ 上的消失，而
$Z_i \to Z$ 的 Leray 谱序列给出 $\mathcal{G}_i$ 的消失。}。
最后由 Theorem \ref{etale-cohomology-theorem-colimit} 得到结论。

\medskip\noindent
考察交换图
$$
\xymatrix{
\mathbf{A}^d_K \ar[rd]_f \ar[rr]_-j & &
\mathbf{P}^1_K \times_K \mathbf{A}^{d - 1}_K \ar[ld]^g \\
& \mathbf{A}^{d - 1}_K
}
$$
注意 $j$ 是光滑 $d$ 维簇的开浸入，其补集为有效 Cartier 除子 $D$。
因此可以使用插曲 I 中得到的结果。下面考察相对 Leray 谱序列
$$
E_2^{p, q} = R^pg_*R^qj_*\mathcal{F} \Rightarrow R^{p + q}f_*\mathcal{F}
$$
层 $R^qj_*\mathcal{F}$ 在 $q > 0$ 时支撑在 $D$ 上；又因
$g|_D : D \to \mathbf{A}^{d - 1}_K$ 是同构，故
$R^pg_*R^qj_*\mathcal{F} = 0$，只要 $p > 0$ 且 $q > 0$。此外，
$R^qj_*\mathcal{F} = 0$，只要
$q > d$。另一方面，$g$ 是相对维数为 $1$ 的固有态射。
因此由 Lemma \ref{etale-cohomology-lemma-cohomological-dimension-proper}，
$R^pg_*j_*\mathcal{F} = 0$，只要 $p > 2$。于是谱序列的 $E_2$-页如下：
$$
\begin{matrix}
g_*R^dj_*\mathcal{F} & 0 & 0 \\
\ldots & \ldots & \ldots \\
g_*R^2j_*\mathcal{F} & 0 & 0 \\
g_*R^1j_*\mathcal{F} & 0 & 0 \\
g_*j_*\mathcal{F} & R^1g_*j_*\mathcal{F} & R^2g_*j_*\mathcal{F}
\end{matrix}
$$
因此 $R^qf_*\mathcal{F} = g_*R^qj_*\mathcal{F}$，只要 $q > 2$。
由插曲 I，$R^qf_*\mathcal{F}$ 在 $q > 2$ 时的支集包含在
$\mathbf{A}^{d - 1}_K$ 的如下点集之中：其剩余域的超越次数为
$\leq d - q$。由插曲 II，
$$
H^p(\mathbf{A}^{d - 1}_K, R^qf_*\mathcal{F}) = 0
\text{ 当 }p > d - q + \text{cd}(K)\text{ 且 }q > 2
$$
另一方面，由 Theorem \ref{etale-cohomology-theorem-higher-direct-images} 有
$R^2f_*\mathcal{F}_{\overline{\eta}} =
H^2(\mathbf{A}^1_{\overline{\eta}}, \mathcal{F}) = 0$
（由维数 $1$ 的情形消失），其中 $\eta$ 是 $\mathbf{A}^{d - 1}_K$ 的一般点。
因此再次由插曲 II，
$$
H^p(\mathbf{A}^{d - 1}_K, R^2f_*\mathcal{F}) = 0
\text{ 对 }p > d - 2 + \text{cd}(K)
$$
最后，由归纳假设有
$$
H^p(\mathbf{A}^{d - 1}_K, R^qf_*\mathcal{F}) = 0
\text{ 当 }p > d - 1 + \text{cd}(K)\text{ 且 }q = 0, 1
$$
。把以上结论与 Leray 谱序列
$H^p(\mathbf{A}^{d - 1}_K, R^qf_*\mathcal{F}) \Rightarrow
H^{p + q}(\mathbf{A}^d_K, \mathcal{F})$ 结合起来即得结论。
\end{proof}

\begin{lemma}
\label{etale-cohomology-lemma-interlude-II}
设 $K$ 是域。设 $X$ 是在 $K$ 上有限型的仿射概形。设
$E_a \subset X$ 是由点 $x \in X$ 构成的集合，这些点满足
$\text{trdeg}_K(\kappa(x)) \leq a$。设 $\mathcal{F}$ 是定义在
$X_\etale$ 上的阿贝尔挠层，其支集包含在 $E_a$ 中。则
$H^b_\etale(X, \mathcal{F}) = 0$，只要 $b > a + \text{cd}(K)$。
\end{lemma}

\begin{proof}
可写 $\mathcal{F} = \colim \mathcal{F}_i$，其中 $\mathcal{F}_i$ 是支撑在
闭子概形 $Z_i$ 上的阿贝尔挠层，该闭子概形包含在 $E_a$ 中，见
Lemma \ref{etale-cohomology-lemma-support-in-subset}。于是 Proposition
\ref{etale-cohomology-proposition-cd-affine} 给出 $\mathcal{F}_i$ 的所需消失。
略去细节；提示：先用 Proposition \ref{etale-cohomology-proposition-closed-immersion-pushforward}
把 $\mathcal{F}_i$ 写成 $Z_i$ 上某层的正像，再使用该层在 $Z_i$ 上的消失，
并用 $Z_i \to Z$ 的 Leray 谱序列得到 $\mathcal{F}_i$ 的消失。
最后由 Theorem \ref{etale-cohomology-theorem-colimit} 得到结论。
\end{proof}

\begin{lemma}
\label{etale-cohomology-lemma-interlude-I}
设 $f : X \to Y$ 是定义在域 $K$ 上的有限型概形之间的仿射态射。设
$E_a(X)$ 是由点 $x \in X$ 构成的集合，这些点满足
$\text{trdeg}_K(\kappa(x)) \leq a$。
设 $\mathcal{F}$ 是 $X_\etale$ 上的阿贝尔挠层，其支集包含在 $E_a$ 中。
则 $R^qf_*\mathcal{F}$ 的支集包含在 $E_{a - q}(Y)$ 中。
\end{lemma}

\begin{proof}
问题在 $Y$ 上是局部的，故可假设 $Y$ 仿射。于是 $X$ 也仿射，并可选取图
$$
\xymatrix{
X \ar[d]_f \ar[r]_i & \mathbf{A}^{n + m}_K \ar[d]^{\text{pr}} \\
Y \ar[r]^j & \mathbf{A}^n_K
}
$$
，其中水平箭头是闭浸入，右侧竖直箭头是投影（略去细节）。
有 $j_*R^qf_*\mathcal{F} = R^q\text{pr}_*i_*\mathcal{F}$；
这是由 $i$ 和 $j$ 的高阶正像消失得到的，见
Proposition \ref{etale-cohomology-proposition-finite-higher-direct-image-zero}。
此外，命题中对 $j_*$ 的茎的描述表明，只须证明
$j_*R^qf_*\mathcal{F}$ 的消失。因此可假设 $f$ 是投影态射
$\text{pr} : \mathbf{A}^{n + m}_K \to \mathbf{A}^n_K$，并设
$\mathcal{F}$ 是 $\mathbf{A}^{n + m}_K$ 上满足本引理假设的阿贝尔挠层。

\medskip\noindent
设 $y$ 是 $\mathbf{A}^n_K$ 的点。由 Theorem
\ref{etale-cohomology-theorem-higher-direct-images} 有
$$
(R^q\text{pr}_*\mathcal{F})_{\overline{y}} =
H^q(\mathbf{A}^{n + m}_K \times_{A^n_K} \Spec(\mathcal{O}_{Y, y}^{sh}),
\mathcal{F}) =
H^q(\mathbf{A}^m_{\mathcal{O}_{Y, y}^{sh}}, \mathcal{F})
$$
记 $b = \text{trdeg}_K(\kappa(y))$。由 Lemma \ref{etale-cohomology-lemma-strictly-henselian}
得到嵌入
$$
L = K(t_1, \ldots, t_b)^{sep} \longrightarrow \mathcal{O}_{Y, y}^{sh}
$$
把 $\mathcal{O}_{Y, y}^{sh} = \colim B_i$ 写成有限型 $L$-子代数
$B_i \subset \mathcal{O}_{Y, y}^{sh}$ 的滤过余极限；这些子代数包含正则函数环
$K[T_1, \ldots, T_n]$，即 $\mathbf{A}^n_K$ 上的正则函数环。于是
$$
\mathbf{A}^m_{\mathcal{O}_{Y, y}^{sh}} =
\lim \mathbf{A}^m_{B_i}
$$
如果 $z \in \mathbf{A}^m_{B_i}$ 是 $\mathcal{F}$ 的支集中的点，
并且 $x$ 是 $z$ 在 $\mathbf{A}^{m + n}_K$ 中的像，则有
$\text{trdeg}_K(\kappa(x)) \leq a$；这是由引理中对 $\mathcal{F}$ 的假设得到的。
由于 $\mathcal{O}_{Y, y}^{sh}$ 是 $K[T_1, \ldots, T_n]$ 上 \'etale 代数的
滤过余极限，并且 $B_i \subset \mathcal{O}_{Y, y}^{sh}$，可见
$\kappa(z)/\kappa(x)$ 是代数扩张（略去若干细节）。于是
$\text{trdeg}_K(\kappa(z)) \leq a$，从而
$\text{trdeg}_L(\kappa(z)) \leq a - b$。由 Lemma \ref{etale-cohomology-lemma-interlude-II}，
$$
H^q(\mathbf{A}^m_{B_i}, \mathcal{F}) = 0\text{ 对 }q > a - b
$$
因此由 Theorem \ref{etale-cohomology-theorem-colimit}，
$(Rf_*\mathcal{F})_{\overline{y}} = 0$，只要 $q > a - b$，这正是所需结论。
\end{proof}








\section{有限上同调维数}
\label{etale-cohomology-section-finite-cd}

\noindent
我们继续 Section \ref{etale-cohomology-section-cd} 中开始的讨论。

\begin{definition}
\label{etale-cohomology-definition-cd-f}
设 $f : X \to Y$ 是概形之间的拟紧拟分离态射。
{\it $f$ 的上同调维数}是集合中最小的元素
$$
\text{cd}(f) \in \{0, 1, 2, \ldots\} \cup \{\infty\}
$$
，它满足：对任意阿贝尔挠层 $\mathcal{F}$（定义在 $X_\etale$ 上），
都有 $R^if_*\mathcal{F} = 0$，只要 $i > \text{cd}(f)$。
\end{definition}

\begin{lemma}
\label{etale-cohomology-lemma-finite-cd}
设 $K$ 是域。
\begin{enumerate}
\item 如果 $f : X \to Y$ 是 $K$ 上有限型概形之间的态射，
则 $\text{cd}(f) < \infty$。
\item 如果 $\text{cd}(K) < \infty$，则 $\text{cd}(X) < \infty$；
这里 $X$ 是 $K$ 上任意有限型概形。
\end{enumerate}
\end{lemma}

\begin{proof}
证明 (1)。可以假设 $Y$ 是仿射的。我们将使用
Cohomology of Schemes，引理 \ref{coherent-lemma-induction-principle}
的归纳原理。如果 $X$ 也是仿射的，则结论由
Lemma \ref{etale-cohomology-lemma-interlude-I} 成立。因此只须证明：若
$X = U \cup V$，且结论对 $U \to Y$、$V \to Y$ 和
$U \cap V \to Y$ 成立，则它对 $f$ 也成立。这由相对
Mayer--Vietoris 序列得到，见 Lemma \ref{etale-cohomology-lemma-relative-mayer-vietoris}。

\medskip\noindent
证明 (2)。我们将使用
Cohomology of Schemes，引理 \ref{coherent-lemma-induction-principle}
的归纳原理。如果 $X$ 是仿射的，则结论由
Proposition \ref{etale-cohomology-proposition-cd-affine} 成立。因此只须证明：若
$X = U \cup V$，且结论对 $U$、$V$ 和 $U \cap V $ 成立，
则它对 $X$ 也成立。这由 Mayer--Vietoris 序列得到，见
Lemma \ref{etale-cohomology-lemma-mayer-vietoris}。
\end{proof}

\begin{lemma}
\label{etale-cohomology-lemma-finite-cd-mod-n-direct-sums}
上同调与直和。设 $n \geq 1$ 是整数。
\begin{enumerate}
\item 设 $f : X \to Y$ 是拟紧拟分离的概形态射，且
$\text{cd}(f) < \infty$。则函子
$$
Rf_* :
D(X_\etale, \mathbf{Z}/n\mathbf{Z})
\longrightarrow
D(Y_\etale, \mathbf{Z}/n\mathbf{Z})
$$
与直和可交换。
\item 设 $X$ 是拟紧拟分离概形，且
$\text{cd}(X) < \infty$。则函子
$$
R\Gamma(X, -) :
D(X_\etale, \mathbf{Z}/n\mathbf{Z})
\longrightarrow
D(\mathbf{Z}/n\mathbf{Z})
$$
与直和可交换。
\end{enumerate}
\end{lemma}

\begin{proof}
证明 (1)。由于 $\text{cd}(f) < \infty$，可知
$$
f_* :
\textit{Mod}(X_\etale, \mathbf{Z}/n\mathbf{Z})
\longrightarrow
\textit{Mod}(Y_\etale, \mathbf{Z}/n\mathbf{Z})
$$
在 Derived Categories, Lemma
\ref{derived-lemma-unbounded-right-derived} 的意义下具有有限上同调维数。
设 $I$ 是集合，并对 $i \in I$ 令 $E_i$ 为
$D(X_\etale, \mathbf{Z}/n\mathbf{Z})$ 的对象。
选取 K-内射复形 $\mathcal{I}_i^\bullet$；它由
$\mathbf{Z}/n\mathbf{Z}$-模组成，每一项
$\mathcal{I}_i^n$ 都是 $\mathbf{Z}/n\mathbf{Z}$-模的内射层，
而且该复形表示 $E_i$。见 Injectives, Theorem
\ref{injectives-theorem-K-injective-embedding-grothendieck}。
于是 $\bigoplus E_i$ 由复形
$\bigoplus \mathcal{I}_i^\bullet$（逐项直和）表示，见
Injectives, Lemma \ref{injectives-lemma-derived-products}。
由 Lemma \ref{etale-cohomology-lemma-relative-colimit} 有
$$
R^qf_*(\bigoplus \mathcal{I}_i^n) =
\bigoplus R^qf_*(\mathcal{I}_i^n) = 0
$$
，只要 $q > 0$，且对任意 $n$ 都成立。因此由
Derived Categories, Lemma \ref{derived-lemma-unbounded-right-derived}，
可知，可以用复形计算 $Rf_*(\bigoplus E_i)$：
$$
f_*(\bigoplus \mathcal{I}_i^\bullet) =
\bigoplus f_*(\mathcal{I}_i^\bullet)
$$
（等式仍由 Lemma \ref{etale-cohomology-lemma-relative-colimit} 得到）。根据已经使用过的
Injectives, Lemma \ref{injectives-lemma-derived-products}，该复形表示
$\bigoplus Rf_*E_i$。

\medskip\noindent
证明 (2)。它与 (1) 的证明完全相同，故略去。
\end{proof}

\begin{lemma}
\label{etale-cohomology-lemma-proper-mod-n-direct-sums}
设 $f : X \to Y$ 是概形之间的固有态射。设 $n \geq 1$ 是整数。则函子
$$
Rf_* :
D(X_\etale, \mathbf{Z}/n\mathbf{Z})
\longrightarrow
D(Y_\etale, \mathbf{Z}/n\mathbf{Z})
$$
与直和可交换。
\end{lemma}

\begin{proof}
只须证明 $Y$ 拟紧的情形。由
Morphisms，引理 \ref{morphisms-lemma-morphism-finite-type-bounded-dimension}
可知 $f : X \to Y$ 的纤维维数有界。因此
Lemma \ref{etale-cohomology-lemma-cohomological-dimension-proper} 蕴含
$\text{cd}(f) < \infty$。于是结论由
Lemma \ref{etale-cohomology-lemma-finite-cd-mod-n-direct-sums} 得到。
\end{proof}

\begin{lemma}
\label{etale-cohomology-lemma-pull-out-constant-mod-n}
设 $X$ 是拟紧拟分离概形，且 $\text{cd}(X) < \infty$。
设 $\Lambda$ 是挠环。设 $E \in D(X_\etale, \Lambda)$ 且
$K \in D(\Lambda)$。则
$$
R\Gamma(X, E \otimes_\Lambda^\mathbf{L} \underline{K}) =
R\Gamma(X, E) \otimes_\Lambda^\mathbf{L} K
$$
\end{lemma}

\begin{proof}
由 Cohomology on Sites, Section
\ref{sites-cohomology-section-projection-formula}，存在从右向左的典范映射。
令 $T(K)$ 表示引理的陈述对 $K \in D(\Lambda)$ 成立这一性质。
我们将验证 More on Algebra, Remark
\ref{more-algebra-remark-P-resolution} 的条件 (1)、(2) 和 (3) 对
$T$ 成立，从而得到结论。
性质 (1) 成立，因为等式两边都与直和可交换，见
Lemma \ref{etale-cohomology-lemma-finite-cd-mod-n-direct-sums}。
性质 (2) 成立，因为我们在比较三角范畴之间的正合函子，并且可以使用
Derived Categories, Lemma
\ref{derived-lemma-third-isomorphism-triangle}。
性质 (3) 是说，引理对 $K = \Lambda[k]$ 成立，其中移位
$k \in \mathbf{Z}$ 任意；这是显然的。
\end{proof}

\begin{lemma}
\label{etale-cohomology-lemma-projection-formula-proper-mod-n}
设 $f : X \to Y$ 是概形之间的固有态射。设 $\Lambda$ 是挠环。
设 $E \in D(X_\etale, \Lambda)$ 且
$K \in D(Y_\etale, \Lambda)$。则
$$
Rf_*E \otimes_\Lambda^\mathbf{L} K =
Rf_*(E \otimes_\Lambda^\mathbf{L} f^{-1}K)
$$
在 $D(Y_\etale, \Lambda)$ 中成立。
\end{lemma}

\begin{proof}
由 Cohomology on Sites, Section
\ref{sites-cohomology-section-projection-formula}，存在从左向右的典范映射。
我们在 $\overline{y}$ 处的茎上检验这个等式。
由固有基变换（采用 Lemma \ref{etale-cohomology-lemma-proper-base-change-mod-n} 中
$Y' = \overline{y}$ 的形式），问题约化到 $Y$ 是代数闭域之谱的情形。
这已在 Lemma \ref{etale-cohomology-lemma-pull-out-constant-mod-n} 中证明；其中由
Lemma \ref{etale-cohomology-lemma-cohomological-dimension-proper} 使用了
$\text{cd}(X) < \infty$。
\end{proof}






\section{\'etale 上同调中的 K\"unneth 公式}
\label{etale-cohomology-section-kunneth}

\noindent
我们先在其中一个因子固有的情形证明 K\"unneth 公式，再用该公式证明
开浸入的一项基变换性质。由此又得到“沿指向域之谱的态射作基变换”
（类似于光滑基变换）。最后再用它得到更一般的 K\"unneth 公式。

\begin{remark}
\label{etale-cohomology-remark-define-kunneth-map}
考察概形范畴中的 Cartesian 图：
$$
\xymatrix{
X \times_S Y \ar[d]_p \ar[r]_q \ar[rd]_c & Y \ar[d]^g \\
X \ar[r]^f & S
}
$$
设 $\Lambda$ 是环，并设 $E \in D(X_\etale, \Lambda)$ 且
$K \in D(Y_\etale, \Lambda)$。则存在典范映射
$$
Rf_*E \otimes_\Lambda^\mathbf{L} Rg_*K
\longrightarrow
Rc_*(p^{-1}E \otimes_\Lambda^\mathbf{L} q^{-1}K)
$$
例如，可以用典范映射 $Rf_*E \to Rc_*p^{-1}E$、
$Rg_*K \to Rc_*q^{-1}K$，以及 Cohomology on Sites,
Remark \ref{sites-cohomology-remark-cup-product} 中定义的相对杯积来定义它。
也可以取下列映射的伴随：
$$
c^{-1}(Rf_*E \otimes_\Lambda^\mathbf{L} Rg_*K)
=
p^{-1}f^{-1}Rf_*E \otimes_\Lambda^\mathbf{L} q^{-1} g^{-1}Rg_*K
\to
p^{-1}E \otimes_\Lambda^\mathbf{L} q^{-1}K
$$
这里使用伴随映射 $f^{-1}Rf_*E \to E$ 和
$g^{-1}Rg_*K \to K$。
\end{remark}


\begin{lemma}
\label{etale-cohomology-lemma-kunneth-one-proper}
设 $k$ 是可分闭域。设 $X$ 是 $k$ 上的固有概形。
设 $Y$ 是 $k$ 上的拟紧拟分离概形。
\begin{enumerate}
\item 如果 $E \in D^+(X_\etale)$ 的上同调层为挠层，且
$K \in D^+(Y_\etale)$，则
$$
R\Gamma(X \times_{\Spec(k)} Y,
\text{pr}_1^{-1}E
\otimes_\mathbf{Z}^\mathbf{L}
\text{pr}_2^{-1}K
)
=
R\Gamma(X, E)
\otimes_\mathbf{Z}^\mathbf{L}
R\Gamma(Y, K)
$$
\item 如果 $n \geq 1$ 是整数，$Y$ 在 $k$ 上有限型，
$E \in D(X_\etale, \mathbf{Z}/n\mathbf{Z})$，且
$K \in D(Y_\etale, \mathbf{Z}/n\mathbf{Z})$，则
$$
R\Gamma(X \times_{\Spec(k)} Y,
\text{pr}_1^{-1}E
\otimes_{\mathbf{Z}/n\mathbf{Z}}^\mathbf{L}
\text{pr}_2^{-1}K
)
=
R\Gamma(X, E)
\otimes_{\mathbf{Z}/n\mathbf{Z}}^\mathbf{L}
R\Gamma(Y, K)
$$
\end{enumerate}
\end{lemma}

\begin{proof}
证明 (1)。由 Lemma \ref{etale-cohomology-lemma-projection-formula-proper} 有
$$
R\text{pr}_{2, *}(
\text{pr}_1^{-1}E
\otimes_\mathbf{Z}^\mathbf{L}
\text{pr}_2^{-1}K) =
R\text{pr}_{2, *}(\text{pr}_1^{-1}E)
\otimes_\mathbf{Z}^\mathbf{L}
K
$$
由固有基变换（采用 Lemma \ref{etale-cohomology-lemma-proper-base-change} 的形式），
它等于对象
$$
\underline{R\Gamma(X, E)}
\otimes_\mathbf{Z}^\mathbf{L}
K
$$
，该对象属于 $D(Y_\etale)$。对此对象取 $R\Gamma(Y, -)$，
再由 $\text{pr}_2$ 的 Leray 谱序列，便得到 (1) 中等式的左边。
于是结论由 Lemma \ref{etale-cohomology-lemma-pull-out-constant} 得到。

\medskip\noindent
证明 (2)。它与 (1) 的证明完全相同，只是要使用
Lemmas \ref{etale-cohomology-lemma-projection-formula-proper-mod-n}、
\ref{etale-cohomology-lemma-proper-base-change-mod-n} 和
\ref{etale-cohomology-lemma-pull-out-constant-mod-n}，以及由
Lemma \ref{etale-cohomology-lemma-finite-cd} 得到的 $\text{cd}(Y) < \infty$。
\end{proof}

\begin{lemma}
\label{etale-cohomology-lemma-supported-in-closed-points}
设 $K$ 是可分闭域。设 $X$ 是 $K$ 上的有限型概形。设
$\mathcal{F}$ 是 $X_\etale$ 上的阿贝尔层，其支集包含在 $X$ 的闭点集内。
则 $H^q(X, \mathcal{F}) = 0$，只要 $q > 0$；并且
$\mathcal{F}$ 由全局截面生成。
\end{lemma}

\begin{proof}
（如果 $\mathcal{F}$ 是挠层，则消失性立即由
Lemma \ref{etale-cohomology-lemma-interlude-II} 得到。）
由 Lemma \ref{etale-cohomology-lemma-support-in-subset}，可把 $\mathcal{F}$
写成可构层 $\mathcal{F}_i$ 的滤过余极限；这些是
$\mathbf{Z}$-模层，其支集 $Z_i \subset X$ 是闭点的有限集。
由 Proposition \ref{etale-cohomology-proposition-closed-immersion-pushforward}，
这种层具有形式 $(Z_i \to X)_*\mathcal{G}_i$，其中
$\mathcal{G}_i$ 是 $Z_i$ 上的层。由于 $K$ 可分闭，概形
$Z_i$ 是若干可分闭域之谱的有限不交并。回顾
$H^q(Z_i, \mathcal{G}_i) = H^q(X, \mathcal{F}_i)$；这是由
$Z_i \to X$ 的 Leray 谱序列以及该态射的高阶正像消失得到的
（Proposition \ref{etale-cohomology-proposition-finite-higher-direct-image-zero}）。
由 Lemmas \ref{etale-cohomology-lemma-equivalence-abelian-sheaves-point} 和
\ref{etale-cohomology-lemma-compare-cohomology-point}，可知
$H^q(Z_i, \mathcal{G}_i)$ 在 $q > 0$ 时为零，而且
$H^0(Z_i, \mathcal{G}_i)$ 生成 $\mathcal{G}_i$。
因此 $H^q(X, \mathcal{F}_i)$ 在 $q > 0$ 时消失，并且
$\mathcal{F}_i$ 由全局截面生成。由 Theorem \ref{etale-cohomology-theorem-colimit} 可知
$H^q(X, \mathcal{F}) = 0$，只要 $q > 0$。
现在证明完成，因为阿贝尔群层中由全局截面生成的层的滤过余极限仍由
全局截面生成（略去细节）。
\end{proof}

\begin{lemma}
\label{etale-cohomology-lemma-vanishing-closed-points}
设 $K$ 是可分闭域。设 $X$ 是 $K$ 上的有限型概形。设
$Q \in D(X_\etale)$。假设 $Q_{\overline{x}}$ 只有在 $x$ 是
$X$ 的闭点时才可能非零。则
$$
Q = 0 \Leftrightarrow H^i(X, Q) = 0 \text{ 对所有 }i
$$
\end{lemma}

\begin{proof}
从左向右的蕴含是平凡的。因此只须证明反向蕴含。

\medskip\noindent
先假设 $Q$ 下有界；这种情形足以处理几乎所有应用。如果 $Q$ 非零，
就考察最小的 $i$，使得上同调层 $H^i(Q)$ 非零。
由 Lemma \ref{etale-cohomology-lemma-supported-in-closed-points} 有
$H^i(X, Q) = H^0(X, H^i(Q)) \not = 0$，结论随即得到。

\medskip\noindent
一般情形。令 $\mathcal{B} \subset \Ob(X_\etale)$ 为拟紧对象的集合。
由 Lemma \ref{etale-cohomology-lemma-supported-in-closed-points}，Cohomology on Sites,
Lemma \ref{sites-cohomology-lemma-cohomology-over-U-trivial} 的假设成立。
因此有 $H^q(U, Q) = H^0(U, H^q(Q))$，对所有
$U \in \mathcal{B}$ 都成立。特别地，这对 $U = X$ 成立。
于是结论由 Lemma \ref{etale-cohomology-lemma-supported-in-closed-points} 得到，因为
$Q$ 在 $D(X_\etale)$ 中为零，当且仅当 $H^q(Q)$ 对所有 $q$ 都为零。
\end{proof}

\begin{lemma}
\label{etale-cohomology-lemma-kunneth-localize-on-X}
设 $K$ 是域。设 $j : U \to X$ 是 $K$ 上有限型概形之间的开浸入。
设 $Y$ 是 $K$ 上的有限型概形。考察图
$$
\xymatrix{
Y \times_{\Spec(K)} X \ar[d]_q  &
Y \times_{\Spec(K)} U \ar[l]^h \ar[d]^p \\
X & U \ar[l]_j
}
$$
则基变换映射 $q^{-1}Rj_*\mathcal{F} \to Rh_*p^{-1}\mathcal{F}$
是同构，其中 $\mathcal{F}$ 是 $U_\etale$ 上的阿贝尔层，
其各茎为阶在 $K$ 中可逆的挠群。
\end{lemma}

\begin{proof}
写成 $\mathcal{F} = \colim \mathcal{F}[n]$，其中余极限取遍
$K$ 中可逆整数构成的乘法系。由于在当前情形上同调与滤过余极限可交换
（精确参考见 Lemma \ref{etale-cohomology-lemma-base-change-Rf-star-colim}），
只须对 $\mathcal{F}[n]$ 证明引理。因此可以假设 $\mathcal{F}$ 是
$\mathbf{Z}/n\mathbf{Z}$-模层，其中某个 $n$ 在 $K$ 中可逆
（证明最后才会用到这一点）。在证明中，用 $X \times_K Y$ 简记
$\Spec(K)$ 上的纤维积。
对 $\dim(X) + \dim(Y)$ 作归纳来证明引理。如果 $\dim(X) \leq 0$，
引理是平凡的，因为此时 $U$ 是 $X$ 的既开又闭子概形。
选取一点 $z \in X \times_K Y$。我们将证明在 $\overline{z}$ 处的茎映射是同构。

\medskip\noindent
假设 $z \mapsto x \in X$，并假设 $\text{trdeg}_K(\kappa(x)) > 0$。
置 $X' = \Spec(\mathcal{O}_{X, x}^{sh})$，并以
$U' \subset X'$ 表示 $U$ 的逆像。考察基变换
$$
\xymatrix{
Y \times_K X' \ar[d]_{q'}  &
Y \times_K U' \ar[l]^{h'} \ar[d]^{p'} \\
X' & U' \ar[l]_{j'}
}
$$
，它由 $X' \to X$ 对原图作基变换得到。注意 $X' \to X$ 是
\'etale 态射的滤过余极限。由采用 Lemma
\ref{etale-cohomology-lemma-smooth-base-change-general} 形式的光滑基变换，
$q^{-1}Rj_*\mathcal{F} \to Rh_*p^{-1}\mathcal{F}$ 在 $X'$ 上、
再到 $Y \times_K X'$ 的拉回是映射
$(q')^{-1}Rj'_*\mathcal{F}' \to Rj'_*(p')^{-1}\mathcal{F}'$，其中
$\mathcal{F}'$ 是 $\mathcal{F}$ 到 $U'$ 的拉回。
（在这一步使用本质上平凡的 \'etale 基变换已经足够。）
因此只须证明 $(q')^{-1}Rj'_*\mathcal{F}' \to Rj'_*(p')^{-1}\mathcal{F}'$
是同构，便可证明原映射在 $\overline{z}$ 处的茎上是同构。
由 Lemma \ref{etale-cohomology-lemma-strictly-henselian}，存在可分闭域 $L/K$，使得
$X' = \lim X_i$，其中 $X_i$ 是 $L$ 上有限型的仿射概形，且
$\dim(X_i) < \dim(X)$。当 $i$ 充分大时，存在开子概形
$U_i \subset X_i$，其限制为 $U'$（在 $X'$ 中）。
可以对图应用归纳假设：
$$
\vcenter{
\xymatrix{
Y \times_K X_i \ar[d]_{q_i}  &
Y \times_K U_i \ar[l]^{h_i} \ar[d]^{p_i} \\
X_i & U_i \ar[l]_{j_i}
}
}
\quad\text{equal to}\quad
\vcenter{
\xymatrix{
Y_L \times_L X_i \ar[d]_{q_i}  &
Y_L \times_L U_i \ar[l]^{h_i} \ar[d]^{p_i} \\
X_i & U_i \ar[l]_{j_i}
}
}
$$
这里底域为 $L$，并取 $\mathcal{F}$ 到这些图的拉回。
由 Lemma \ref{etale-cohomology-lemma-base-change-Rf-star-colim} 可知，映射
$(q')^{-1}Rj'_*\mathcal{F}' \to Rj'_*(p')^{-1}\mathcal{F}$ 是同构。

\medskip\noindent
假设 $z \mapsto y \in Y$，并假设 $\text{trdeg}_K(\kappa(y)) > 0$。
令 $Y' = \Spec(\mathcal{O}_{X, x}^{sh})$。
由 Lemma \ref{etale-cohomology-lemma-strictly-henselian}，存在可分闭域 $L/K$，使得
$Y' = \lim Y_i$，其中 $Y_i$ 是 $L$ 上有限型的仿射概形，且
$\dim(Y_i) < \dim(Y)$。特别地，$Y'$ 是 $L$ 上的概形。
用下标 $L$ 表示从 $K$ 上概形到 $L$ 上概形的基变换。考察交换图
$$
\vcenter{
\xymatrix{
Y' \times_K X \ar[d]_f  &
Y' \times_K U \ar[l]^{h'} \ar[d]^{f'} \\
Y \times_K X \ar[d]_q  &
Y \times_K U \ar[l]^h \ar[d]^p \\
X & U \ar[l]_j
}
}
\quad\text{且}\quad
\vcenter{
\xymatrix{
Y' \times_L X_L \ar[d]_{q'}  &
Y' \times_L U_L \ar[l]^{h'} \ar[d]^{p'} \\
X_L \ar[d] &
U_L \ar[l]^{j_L} \ar[d] \\
X & U \ar[l]_j
}
}
$$
并注意，左图与右图的顶行相同，底行也相同。
由光滑基变换可知
$f^{-1}Rh_*p^{-1}\mathcal{F} = Rh'_*(f')^{-1}p^{-1}\mathcal{F}$
（与上一段类似）。由 $\Spec(L) \to \Spec(K)$ 的光滑基变换
（Lemma \ref{etale-cohomology-lemma-base-change-field-extension}），可知
$Rj_{L, *}\mathcal{F}_L$ 是 $Rj_*\mathcal{F}$ 到 $X_L$ 的拉回。
合并这两个观察可知，只须证明右图上方方形的基变换映射是同构，
便可证明原映射在 $\overline{z}$ 处的茎上是同构\footnote{这里使用：
基变换映射的“竖直复合”仍是基变换映射，见 Cohomology on Sites,
Remark \ref{sites-cohomology-remark-compose-base-change}。}。
然后利用 $Y' = \lim Y_i$，完全照上一段论证，可见归纳假设迫使
$Y' \times_K X$ 上的映射成为同构。

\medskip\noindent
因此，任何使 $\dim(X) + \dim(Y)$ 最小的反例，都只会使映射
$q^{-1}Rj_*\mathcal{F} \to Rh_*p^{-1}\mathcal{F}$
在 $X \times_K Y$ 的闭点之茎上不是同构（因为一点
$z$（属于 $X \times_K Y$）是闭点，当且仅当 $z$ 在 $X$ 与 $Y$
中的像都是闭点）。由于只须局部证明同构，可以假设 $X$ 和 $Y$ 仿射。
但这时可选取稠密开浸入 $Y \to Y'$，使 $Y'$ 射影。
（选取闭浸入 $Y \to \mathbf{A}^n_K$，并令 $Y'$ 为 $Y$ 在
$\mathbf{P}^n_K$ 中的概形论闭包。）于是 $\dim(Y') = \dim(Y)$，
从而得到一个“最小”反例，其中 $Y$ 在 $K$ 上射影。
下一段将说明这是不可能的。

\medskip\noindent
考察引理陈述中的图，并假设 $q^{-1}Rj_*\mathcal{F} \to Rh_*p^{-1}\mathcal{F}$
在 $X \times_K Y$ 的所有非闭点处都是同构，且 $Y$ 固有。
该映射到 $(X \times_K Y)_{K^{sep}}$ 的限制，就是把引理之图基变换到
$K^{sep}$ 后所得的相应映射。因此可以并且确实假设 $K$ 可分闭。
选取正规三角
$$
q^{-1}Rj_*\mathcal{F} \to Rh_*p^{-1}\mathcal{F} \to Q \to
(q^{-1}Rj_*\mathcal{F})[1]
$$
，它位于 $D((X \times_K Y)_\etale)$ 中。由于 $Q$ 支撑在闭点上，
由 Lemma \ref{etale-cohomology-lemma-vanishing-closed-points}，只须证明
$H^i(X \times_K Y, Q) = 0$ 对所有 $i$ 成立。
因此只须证明
$q^{-1}Rj_*\mathcal{F} \to Rh_*p^{-1}\mathcal{F}$
在上同调上诱导同构。回顾 $\mathcal{F}$ 被 $n$ 零化，且该整数
在 $K$ 中可逆。由 Lemma \ref{etale-cohomology-lemma-kunneth-one-proper} 的 K\"unneth 公式有
\begin{align*}
R\Gamma(X \times_K Y, q^{-1}Rj_*\mathcal{F})
&=
R\Gamma(X, Rj_*\mathcal{F})
\otimes_{\mathbf{Z}/n\mathbf{Z}}^\mathbf{L}
R\Gamma(Y, \mathbf{Z}/n\mathbf{Z}) \\
& =
R\Gamma(U, \mathcal{F})
\otimes_{\mathbf{Z}/n\mathbf{Z}}^\mathbf{L}
R\Gamma(Y, \mathbf{Z}/n\mathbf{Z})
\end{align*}
并且
$$
R\Gamma(X \times_K Y, Rh_*p^{-1}\mathcal{F}) =
R\Gamma(U \times_K Y, p^{-1}\mathcal{F}) =
R\Gamma(U, \mathcal{F})
\otimes_{\mathbf{Z}/n\mathbf{Z}}^\mathbf{L}
R\Gamma(Y, \mathbf{Z}/n\mathbf{Z})
$$
证明至此完成。
\end{proof}

\begin{lemma}
\label{etale-cohomology-lemma-punctual-base-change}
设 $K$ 是域。对任意交换图
$$
\xymatrix{
X \ar[d] & X' \ar[l] \ar[d]_{f'} & Y \ar[l]^h \ar[d]^e \\
\Spec(K) & S' \ar[l] & T \ar[l]_g
}
$$
，其中各概形定义在 $K$ 上，
$X' = X \times_{\Spec(K)} S'$ 且 $Y = X' \times_{S'} T$，
$g$ 拟紧拟分离；再对任意阿贝尔层 $\mathcal{F}$（定义在
$T_\etale$ 上），若其各茎是阶在 $K$ 中可逆的挠群，则基变换映射
$$
(f')^{-1}Rg_*\mathcal{F}
\longrightarrow
Rh_*e^{-1}\mathcal{F}
$$
是同构。
\end{lemma}

\begin{proof}
问题在 $X$ 上是局部的，故可假设 $X$ 仿射。由 Limits, Lemma
\ref{limits-lemma-relative-approximation}，可把
$X = \lim X_i$ 写成余滤过极限，过渡态射仿射，而其中的概形
$X_i$ 在 $K$ 上有限型。记
$X'_i =  X_i \times_{\Spec(K)} S'$ 且 $Y_i = X'_i \times_{S'} T$。
由 Lemma \ref{etale-cohomology-lemma-base-change-Rf-star-colim}，只须对顶点为
$X_i, Y_i, S_i, T_i$ 的方形证明断言。因此可以假设 $X$ 在 $K$ 上有限型。
类似地，可写 $\mathcal{F} = \colim \mathcal{F}[n]$，其中余极限取遍
$K$ 中可逆整数构成的乘法系。上面使用的同一引理把问题约化到：
$\mathcal{F}$ 是 $\mathbf{Z}/n\mathbf{Z}$-模层，其中某个
$n$ 在 $K$ 中可逆。

\medskip\noindent
可以把 $K$ 替换为其代数闭包 $\overline{K}$。事实上，依
Lemma \ref{etale-cohomology-lemma-base-change-field-extension}，正像的形成与到
$\overline{K}$ 的基变换可交换（对 $g$ 与 $h$ 都成立）；而且只须在限制到
$X'_{\overline{K}}$ 后证明两者相同。接着，由 \'etale 上同调的拓扑不变性，
可以用既约化代替 $X$，见
Proposition \ref{etale-cohomology-proposition-topological-invariance}。
作此替换后，态射 $X \to \Spec(K)$ 平坦、有限呈现，并具有几何既约纤维；
任意基变换也有同样性质，特别是 $X' \to S'$。因此由
Lemma \ref{etale-cohomology-lemma-fppf-reduced-fibres-base-change-f-star}，
$(f')^{-1}g_*\mathcal{F} \to Rh_*e^{-1}\mathcal{F}$ 是同构。

\medskip\noindent
此时可以应用 Lemma \ref{etale-cohomology-lemma-base-change-does-not-hold-post}，
可见只须证明如下断言：给定交换图
$$
\xymatrix{
X \ar[d]_f & X' \ar[d] \ar[l] & Y \ar[l]^h \ar[d] \\
\Spec(K) & S' \ar[l] & \Spec(L) \ar[l]
}
$$
，两方形均 Cartesian；其中 $S'$ 仿射、整且正规，其函数域
$K$ 代数闭。则 $R^qh_*(\mathbf{Z}/d\mathbf{Z})$ 为零，只要
$q > 0$ 且 $d | n$。注意，该消失性等价于
$$
(f')^{-1}R^q(\Spec(L) \to S')_*\mathbf{Z}/d\mathbf{Z}
\longrightarrow
R^qh_*\mathbf{Z}/d\mathbf{Z}
$$
是同构，因为例如由
Lemma \ref{etale-cohomology-lemma-Rf-star-zero-normal-with-alg-closed-function-field}，
左边为零。

\medskip\noindent
写 $S' = \Spec(B)$，于是 $L$ 是 $B$ 的分式域。把
$B = \bigcup_{i \in I} B_i$ 写成其有限型 $K$-子代数 $B_i$ 的并。
令 $J$ 为二元组 $(i, g)$ 的集合，其中 $i \in I$ 且
$g \in B_i$ 非零；其序为：$(i', g') \geq (i, g)$ 当且仅当
$i' \geq i$，并且 $g$ 在 $(B_{i'})_{g'}$ 中的像可逆。于是
$L = \colim_{(i, g) \in J} (B_i)_g$。
对 $j = (i, g) \in J$，置 $S_j = \Spec(B_i)$ 且
$U_j = \Spec((B_i)_g)$。于是
$$
\vcenter{
\xymatrix{
X' \ar[d] & Y \ar[l]^h \ar[d] \\
S' & \Spec(L) \ar[l]
}
}
\quad\text{是下列对象的余极限}\quad
\vcenter{
\xymatrix{
X \times_K S_j \ar[d] & X \times_K U_j \ar[l]^{h_j} \ar[d] \\
S_j & U_j \ar[l]
}
}
$$
因此可以应用 Lemma \ref{etale-cohomology-lemma-base-change-Rf-star-colim}，
可见只须证明右侧诸图中的基变换成立；这正是
Lemma \ref{etale-cohomology-lemma-kunneth-localize-on-X} 中已经证明的内容。
\end{proof}

\begin{lemma}
\label{etale-cohomology-lemma-punctual-base-change-upgrade}
设 $K$ 是域。设 $n \geq 1$ 在 $K$ 中可逆。考察交换图
$$
\xymatrix{
X \ar[d] & X' \ar[l]^p \ar[d]_{f'} & Y \ar[l]^h \ar[d]^e \\
\Spec(K) & S' \ar[l] & T \ar[l]_g
}
$$
，其中
$X' = X \times_{\Spec(K)} S'$ 且 $Y = X' \times_{S'} T$，
$g$ 拟紧拟分离。典范映射
$$
p^{-1}E \otimes_{\mathbf{Z}/n\mathbf{Z}}^\mathbf{L} (f')^{-1}Rg_*F
\longrightarrow
Rh_*(h^{-1}p^{-1}E \otimes_{\mathbf{Z}/n\mathbf{Z}}^\mathbf{L} e^{-1}F)
$$
是同构，只要 $E$ 属于 $D^+(X_\etale, \mathbf{Z}/n\mathbf{Z})$，
其 Tor 振幅包含在 $[a, \infty]$ 内，其中 $a \in \mathbf{Z}$ 为某个整数；
并且 $F$ 属于 $D^+(T_\etale, \mathbf{Z}/n\mathbf{Z})$。
\end{lemma}

\begin{proof}
本引理把 Lemma \ref{etale-cohomology-lemma-punctual-base-change} 推广到导出范畴中的对象；
断言成立，是因为在 Lemma \ref{etale-cohomology-lemma-punctual-base-change} 中，
概形 $X$ 定义在 $K$ 上且是任意的。我们强烈建议读者跳过这个繁琐的证明
（另一选择是只读最后一段）。

\medskip\noindent
可以把 $E$ 表示为下有界 K-平坦复形 $\mathcal{E}^\bullet$；
该复形由平坦 $\mathbf{Z}/n\mathbf{Z}$-模组成。见 Cohomology on Sites,
Lemma \ref{sites-cohomology-lemma-bounded-below-tor-amplitude}。
选取整数 $b$，使得 $H^i(F) = 0$ 对 $i < b$ 成立。
选取充分大的整数 $N$，并考察短正合列
$$
0 \to \sigma_{\geq N + 1}\mathcal{E}^\bullet \to
\mathcal{E}^\bullet \to
\sigma_{\leq N}\mathcal{E}^\bullet \to 0
$$
，其中使用朴素截断。它产生正规三角
$E'' \to E \to E' \to E''[1]$，该三角位于
$D(X_\etale, \mathbf{Z}/n\mathbf{Z})$ 中。固定 $F$ 后，引理陈述中箭头的两边
都是关于 $E$ 的正合函子。注意
$$
p^{-1}E'' \otimes_{\mathbf{Z}/n\mathbf{Z}}^\mathbf{L} (f')^{-1}Rg_*F
\quad\text{且}\quad
Rh_*(h^{-1}p^{-1}E'' \otimes_{\mathbf{Z}/n\mathbf{Z}}^\mathbf{L} e^{-1}F)
$$
位于次数 $\geq N + b$。因此，如果能对对象 $E'$ 证明引理，
便可知引理在次数 $\leq N + b$ 中成立，从而得到结论。略去若干细节。
于是可以假设 $E$ 由平坦 $\mathbf{Z}/n\mathbf{Z}$-模组成的有界复形表示。
再作一次同类论证，可以假设 $E$ 由单个平坦
$\mathbf{Z}/n\mathbf{Z}$-模 $\mathcal{E}$ 给出。

\medskip\noindent
接着，对变量 $F$ 使用同一论证，把问题约化到 $F$ 由单个
$\mathbf{Z}/n\mathbf{Z}$-模层 $\mathcal{F}$ 给出的情形。
设 $\mathcal{F}$ 被整数 $m | n$ 零化。如果 $\ell$ 是整除
$m$ 的素数，且 $m > \ell$，则可以考察短正合列
$0 \to \mathcal{F}[\ell] \to \mathcal{F} \to
\mathcal{F}/\mathcal{F}[\ell] \to 0$
，从而约化到更小的 $m$。最终，问题约化到 $\mathcal{F}$ 被某个素数
$\ell$ 零化的情形，该素数整除 $n$。这时注意
$$
p^{-1}\mathcal{E}
\otimes_{\mathbf{Z}/n\mathbf{Z}}^\mathbf{L}
(f')^{-1}Rg_*\mathcal{F}
=
p^{-1}(\mathcal{E}/\ell \mathcal{E})
\otimes_{\mathbf{F}_\ell}^\mathbf{L}
(f')^{-1}Rg_*\mathcal{F}
$$
，这是由 $\mathcal{E}$ 的平坦性得到的。另一项同理。
因此问题约化到处理 $\mathbf{F}_\ell$-向量空间层的情形，下面讨论它。

\medskip\noindent
假设 $\ell$ 是在 $K$ 中可逆的素数。假设 $\mathcal{E}$、
$\mathcal{F}$ 是 $\mathbf{F}_\ell$-向量空间层，分别定义在
$X_\etale$ 与 $T_\etale$ 上。我们要证明
$$
p^{-1}\mathcal{E} \otimes_{\mathbf{F}_\ell} (f')^{-1}R^qg_*\mathcal{F}
\longrightarrow
R^qh_*(h^{-1}p^{-1}\mathcal{E} \otimes_{\mathbf{F}_\ell} e^{-1}\mathcal{F})
$$
对每个 $q \geq 0$ 都是同构。这个问题在 $X$ 上是局部的，故可假设
$X$ 仿射。可把 $\mathcal{E}$ 写成可构
$\mathbf{F}_\ell$-向量空间层的滤过余极限，这些层定义在 $X_\etale$ 上，见
Lemma \ref{etale-cohomology-lemma-torsion-colimit-constructible}。
由于张量积与滤过余极限可交换，高阶正像也如此
（Lemma \ref{etale-cohomology-lemma-relative-colimit}），可以假设 $\mathcal{E}$ 是
可构 $\mathbf{F}_\ell$-向量空间层，定义在 $X_\etale$ 上。
于是可以选取整数 $m$ 和有限且有限呈现的态射
$\pi_i : X_i \to X$，其中 $i = 1, \ldots, m$，使得存在单射
$$
\mathcal{E} \to
\bigoplus\nolimits_{i = 1, \ldots, m}
\pi_{i, *}\mathbf{F}_\ell
$$
，见 Lemma \ref{etale-cohomology-lemma-constructible-maps-into-constant-general}。
注意该直和也是可构层
（Lemma \ref{etale-cohomology-lemma-finite-pushforward-constructible}），
故其余核也可构（Lemma \ref{etale-cohomology-lemma-constructible-abelian}）。
由维数移位，即对 $q$ 作归纳，在可构
$\mathbf{F}_\ell$-向量空间层（定义在 $X_\etale$ 上）的范畴中，只须对层
$\pi_{i, *}\mathbf{F}_\ell$ 证明结论
（略去细节；提示：先证明 $q = 0$ 时的单射性，并让
$\mathcal{E}$ 遍历所有可构层）。
为了证明这种情形，把引理中的图扩充为
$$
\xymatrix{
X_i \ar[d]^{\pi_i} &
X'_i \ar[l]^{p_i} \ar[d]^{\pi'_i} &
Y_i \ar[l]^{h_i} \ar[d]^{\rho_i} \\
X \ar[d] & X' \ar[l]^p \ar[d]_{f'} & Y \ar[l]^h \ar[d]^e \\
\Spec(K) & S' \ar[l] & T \ar[l]_g
}
$$
，其中所有方形都 Cartesian。在以下等式中，将使用
$R\pi_{i, *} = \pi_{i, *}$，以及 $\pi'_i$、$\rho_i$ 的类似等式；
还将使用 Leray 谱序列、Lemma
\ref{etale-cohomology-lemma-finite-pushforward-commutes-with-base-change}，并对
$\pi_i$、$\pi'_i$、$\rho_i$ 使用
引理 \ref{etale-cohomology-lemma-projection-formula-proper-mod-n}
（对有限态射而言，该引理几乎是平凡的）。由此可见
\begin{align*}
p^{-1}\pi_{i, *}\mathbf{F}_\ell
\otimes_{\mathbf{F}_\ell} (f')^{-1}R^qg_*\mathcal{F}
& =
\pi'_{i, *}\mathbf{F}_\ell
\otimes_{\mathbf{F}_\ell} (f')^{-1}R^qg_*\mathcal{F} \\
& =
\pi'_{i, *}((\pi'_i)^{-1} (f')^{-1}R^qg_*\mathcal{F})
\end{align*}
类似地，有
\begin{align*}
R^qh_*(h^{-1}p^{-1} \pi_{i, *}\mathbf{F}_\ell
\otimes_{\mathbf{F}_\ell} e^{-1}\mathcal{F})
& =
R^qh_*(\rho_{i, *}\mathbf{F}_\ell
\otimes_{\mathbf{F}_\ell} e^{-1}\mathcal{F}) \\
& =
R^qh_*(\rho_i^{-1}e^{-1}\mathcal{F}) \\
& =
\pi'_{i, *}R^qh_{i, *} \rho_i^{-1}e^{-1}\mathcal{F})
\end{align*}
由于
$R^qh_{i, *} \rho_i^{-1}e^{-1}\mathcal{F} = 
(\pi'_i)^{-1} (f')^{-1}R^qg_*\mathcal{F}$，这由
Lemma \ref{etale-cohomology-lemma-punctual-base-change} 得到，故结论得证。
\end{proof}

\begin{lemma}
\label{etale-cohomology-lemma-punctual-base-change-upgrade-unbounded}
设 $K$ 是域。设 $n \geq 1$ 在 $K$ 中可逆。考察交换图
$$
\xymatrix{
X \ar[d] & X' \ar[l]^p \ar[d]_{f'} & Y \ar[l]^h \ar[d]^e \\
\Spec(K) & S' \ar[l] & T \ar[l]_g
}
$$
，其中各概形在 $K$ 上有限型，
$X' = X \times_{\Spec(K)} S'$ 且 $Y = X' \times_{S'} T$。
典范映射
$$
p^{-1}E \otimes_{\mathbf{Z}/n\mathbf{Z}}^\mathbf{L} (f')^{-1}Rg_*F
\longrightarrow
Rh_*(h^{-1}p^{-1}E \otimes_{\mathbf{Z}/n\mathbf{Z}}^\mathbf{L} e^{-1}F)
$$
是同构，其中 $E$ 属于 $D(X_\etale, \mathbf{Z}/n\mathbf{Z})$，
$F$ 属于 $D(T_\etale, \mathbf{Z}/n\mathbf{Z})$。
\end{lemma}

\begin{proof}
我们将利用这些函子与直和可交换这一事实，把问题约化到
Lemma \ref{etale-cohomology-lemma-punctual-base-change-upgrade}。建议读者跳过证明。
回顾：导出张量积与直和可交换；（导出）拉回也与直和可交换；而
$Rh_*$ 与 $Rg_*$ 同样与直和可交换，见 Lemmas \ref{etale-cohomology-lemma-finite-cd} 和
\ref{etale-cohomology-lemma-finite-cd-mod-n-direct-sums}
（这里用到了各概形在 $K$ 上有限型）。

\medskip\noindent
为完成证明，可作如下论证。先写
$E = \text{hocolim} \tau_{\leq N} E$。由于这些函子与直和可交换，
它们也与同伦余极限可交换。因此只须对上有界的 $E$ 证明引理。
类似地，对 $F$ 可以假设 $F$ 上有界。于是可以把 $E$ 表示为上有界复形
$\mathcal{E}^\bullet$，其项为 $\mathbf{Z}/n\mathbf{Z}$-模层。于是
$$
\mathcal{E}^\bullet = \colim \sigma_{\geq -N}\mathcal{E}^\bullet
$$
（朴素截断）。故可假设 $\mathcal{E}^\bullet$ 是
$\mathbf{Z}/n\mathbf{Z}$-模层组成的有界复形。对 $F$，选取由平坦（！）
$\mathbf{Z}/n\mathbf{Z}$-模层组成的上有界复形。
这样便约化到 $F$ 由平坦 $\mathbf{Z}/n\mathbf{Z}$-模层组成的有界复形表示
这一情形。此时应用 Lemma \ref{etale-cohomology-lemma-punctual-base-change-upgrade} 即得结论。
\end{proof}

\begin{lemma}
\label{etale-cohomology-lemma-kunneth}
设 $k$ 是可分闭域。设 $X$ 与 $Y$ 是 $k$ 上的有限型概形。
设 $n \geq 1$ 是在 $k$ 中可逆的整数。则对
$E \in D(X_\etale, \mathbf{Z}/n\mathbf{Z})$ 与
$K \in D(Y_\etale, \mathbf{Z}/n\mathbf{Z})$ 有
$$
R\Gamma(X \times_{\Spec(k)} Y,
\text{pr}_1^{-1}E
\otimes_{\mathbf{Z}/n\mathbf{Z}}^\mathbf{L}
\text{pr}_2^{-1}K
)
=
R\Gamma(X, E)
\otimes_{\mathbf{Z}/n\mathbf{Z}}^\mathbf{L}
R\Gamma(Y, K)
$$
\end{lemma}

\begin{proof}
由 Lemma \ref{etale-cohomology-lemma-punctual-base-change-upgrade-unbounded} 有
$$
R\text{pr}_{1, *}(
\text{pr}_1^{-1}E
\otimes_{\mathbf{Z}/n\mathbf{Z}}^\mathbf{L}
\text{pr}_2^{-1}K) =
E \otimes_{\mathbf{Z}/n\mathbf{Z}}^\mathbf{L}
\underline{R\Gamma(Y, K)}
$$
结论由 Lemma \ref{etale-cohomology-lemma-pull-out-constant-mod-n} 得到；
可以使用该引理，是因为由 Lemma \ref{etale-cohomology-lemma-finite-cd} 有
$\text{cd}(X) < \infty$。
\end{proof}














\section{混沌拓扑与 Zariski 拓扑的比较}
\label{etale-cohomology-section-compare-chaotic-Zariski}

\noindent
在构造仿射概形的结构层时，我们先构造它在仿射开子集上的取值，
然后再延拓到所有开子集。对于构造概形 $X$ 上的阿贝尔群复形，
类似的构造往往也很有用。回顾 $X_{affine, Zar}$ 表示 $X$ 的
仿射开子集所成的范畴，其拓扑由标准 Zariski 覆盖给出，见
Topologies, Definition \ref{topologies-definition-big-small-Zariski}。
提醒读者，$X_{affine, Zar}$ 的拓扑斯就是 $X$ 的小 Zariski 拓扑斯，
见 Topologies, Lemma \ref{topologies-lemma-alternative-zariski}。
本节用 $X_{affine}$ 表示同一底范畴，但赋予混沌拓扑；也就是说，
层与预层相同。于是得到位点态射
$$
\epsilon : X_{affine, Zar} \longrightarrow X_{affine}
$$
，如 Cohomology on Sites, Section \ref{sites-cohomology-section-compare}
中所述。

\begin{lemma}
\label{etale-cohomology-lemma-check-zar}
在上述情形中，设 $K$ 是 $D^+(X_{affine})$ 的对象。则 $K$ 属于
（全忠实）函子
$R\epsilon_* ; D(X_{affine, Zar}) \to D(X_{affine})$ 的本质像，
当且仅当下列两个条件成立：
\begin{enumerate}
\item $R\Gamma(\emptyset, K)$ 在 $D(\textit{Ab})$ 中为零；
\item 如果 $U = V \cup W$，其中 $U, V, W \subset X$ 是仿射开子集，
而 $V, W \subset U$ 是标准开子集
（Algebra, Definition \ref{algebra-definition-Zariski-topology}），
则 Cohomology on Sites, Lemma
\ref{sites-cohomology-lemma-c-square} 中的映射
$c^K_{U, V, W, V \cap W}$ 是拟同构。
\end{enumerate}
\end{lemma}

\begin{proof}
（由 Cohomology on Sites, Section
\ref{sites-cohomology-section-compare} 中的讨论，函子
$R\epsilon_*$ 是全忠实的。）
除去一个与空集有关的障碍，结论来自非常一般的 Cohomology on Sites,
Lemma \ref{sites-cohomology-lemma-descent-squares}；该引理的假设由
Schemes, Lemma \ref{schemes-lemma-sheaf-on-affines} 和
Cohomology on Sites，引理
\ref{sites-cohomology-lemma-descent-squares-helper} 保证。

\medskip\noindent
为绕开这一障碍，记 $X_{affine, almost-chaotic}$ 为如下位点：
其中空对象 $\emptyset$ 有两个覆盖，即
$\{\emptyset \to \emptyset\}$ 与空覆盖
（有关讨论见 Sites, Example \ref{sites-example-site-topological}）。
于是有位点态射
$$
X_{affine, Zar} \to X_{affine, almost-chaotic} \to X_{affine}
$$
上述论证适用于第一个箭头。然后留给读者验证：对象 $K$ 若属于
$D^+(X_{affine})$，则它属于（全忠实）函子
$D(X_{affine, almost-chaotic}) \to D(X_{affine})$ 的本质像，
当且仅当 $R\Gamma(\emptyset, K)$ 在 $D(\textit{Ab})$ 中为零。
\end{proof}









\section{大拓扑斯与小拓扑斯的比较}
\label{etale-cohomology-section-compare}

\noindent
设 $S$ 是概形。在 Topologies, Lemma
\ref{topologies-lemma-at-the-bottom-etale} 中，我们引入了比较态射
$\pi_S : (\Sch/S)_\etale \to S_\etale$ 和
$i_S : \Sh(S_\etale) \to \Sh((\Sch/S)_\etale)$
，满足 $\pi_S \circ i_S = \text{id}$ 且
$\pi_{S, *} = i_S^{-1}$。更一般地，如果 $f : T \to S$ 是
$(\Sch/S)_\etale$ 的对象，则存在态射
$i_f : \Sh(T_\etale) \to \Sh((\Sch/S)_\etale)$，使得
$f_{small} = \pi_S \circ i_f$，见 Topologies, Lemmas
\ref{topologies-lemma-put-in-T-etale} 和
\ref{topologies-lemma-morphism-big-small-etale}。在 Descent, Remark
\ref{descent-remark-change-topologies-ringed} 中，我们把这些态射延拓为
环化网站的态射
$$
\pi_S : ((\Sch/S)_\etale, \mathcal{O}) \to (S_\etale, \mathcal{O}_S)
$$
以及环化拓扑斯的态射
$$
i_S : (\Sh(S_\etale), \mathcal{O}_S) \to (\Sh((\Sch/S)_\etale), \mathcal{O})
$$
和
$$
i_f : (\Sh(T_\etale), \mathcal{O}_T) \to (\Sh((\Sch/S)_\etale, \mathcal{O}))
$$
。注意，限制 $i_S^{-1} = \pi_{S, *}$（见 Topologies, Definition
\ref{topologies-definition-restriction-small-etale}）把
$\mathcal{O}$ 变为 $\mathcal{O}_S$。类似地，$i_f^{-1}$ 把
$\mathcal{O}$ 变为 $\mathcal{O}_T$。见 Descent, Remark
\ref{descent-remark-change-topologies-ringed}。因此有
$i_S^*\mathcal{F} = i_S^{-1}\mathcal{F}$ 和
$i_f^*\mathcal{F} = i_f^{-1}\mathcal{F}$，这对任意
$\mathcal{O}$-模 $\mathcal{F}$（定义在 $(\Sch/S)_\etale$ 上）成立。
特别地，$i_S^*$ 与
$i_f^*$ 都是正合函子。函子 $i_S^*$ 常记作
$\mathcal{F} \mapsto \mathcal{F}|_{S_\etale}$（这与 Topologies,
Definition \ref{topologies-definition-restriction-small-etale}
中的记号并不冲突）。

\begin{lemma}
\label{etale-cohomology-lemma-compare-injectives}
设 $S$ 是概形。设 $T$ 是 $(\Sch/S)_\etale$ 的对象。
\begin{enumerate}
\item 如果 $\mathcal{I}$ 在 $\textit{Ab}((\Sch/S)_\etale)$ 中是内射对象，则
\begin{enumerate}
\item $i_f^{-1}\mathcal{I}$ 在 $\textit{Ab}(T_\etale)$ 中是内射对象；
\item $\mathcal{I}|_{S_\etale}$ 在 $\textit{Ab}(S_\etale)$ 中是内射对象。
\end{enumerate}
\item 如果 $\mathcal{I}^\bullet$ 是
$\textit{Ab}((\Sch/S)_\etale)$ 中的 K-内射复形，则
\begin{enumerate}
\item $i_f^{-1}\mathcal{I}^\bullet$ 是
$\textit{Ab}(T_\etale)$ 中的 K-内射复形；
\item $\mathcal{I}^\bullet|_{S_\etale}$ 是
$\textit{Ab}(S_\etale)$ 中的 K-内射复形。
\end{enumerate}
\end{enumerate}
对模而言，相应断言并不成立。
\end{lemma}

\begin{proof}
(1)(b) 与 (2)(b) 形式地来自如下事实：限制函子
$\pi_{S, *} = i_S^{-1}$ 是正合函子 $\pi_S^{-1}$ 的右伴随。
见 Homology, Lemma \ref{homology-lemma-adjoint-preserve-injectives} 和
Derived Categories, Lemma
\ref{derived-lemma-adjoint-preserve-K-injectives}。

\medskip\noindent
(1)(a) 与 (2)(a) 有两种证法。第一种证法：可以利用
$i_f^{-1}$ 是正合函子 $i_{f, !}$ 的右伴随。对集合层，该函子构造于
Topologies, Lemma \ref{topologies-lemma-put-in-T-etale}；对阿贝尔层，
构造于 Modules on Sites, Lemma
\ref{sites-modules-lemma-g-shriek-adjoint}。Modules on Sites, Lemma
\ref{sites-modules-lemma-exactness-lower-shriek} 证明了它是正合的。
第二种证法：由 Topologies, Lemma
\ref{topologies-lemma-morphism-big-small-etale} 可知
$i_f = i_T \circ f_{big}$。由于 $f_{big}$ 是局部化，沿它的拉回保持
内射对象与 K-内射对象，见 Cohomology on Sites, Lemmas
\ref{sites-cohomology-lemma-cohomology-of-open} 和
\ref{sites-cohomology-lemma-restrict-K-injective-to-open}。
然后对函子 $i_T^{-1}$ 应用已经证明的 (1)(b) 与 (2)(b)，即可得出结论。

\medskip\noindent
令 $S = \Spec(\mathbf{Z})$，并考察映射
$2 : \mathcal{O}_S \to \mathcal{O}_S$。这是 $\mathcal{O}_S$-模的单射，
定义在 $S_\etale$ 上。然而，拉回
$\pi_S^*(2) : \mathcal{O} \to \mathcal{O}$ 不是单射；在
$\Spec(\mathbf{F}_2)$ 上取值即可看出这一点。现在选取单射
$\alpha : \mathcal{O} \to \mathcal{I}$，其中 $\mathcal{O}$-模
$\mathcal{I}$ 是内射的，并定义在 $(\Sch/S)_\etale$ 上。
考察图
$$
\xymatrix{
\mathcal{O}_S \ar[d]_2 \ar[rr]_{\alpha|_{S_\etale}} & &
\mathcal{I}|_{S_\etale} \\
\mathcal{O}_S \ar@{..>}[rru]
}
$$
虚线箭头不可能存在于 $\mathcal{O}_S$-模范畴中，因为由伴随性，
它将意味着单射 $\alpha$ 通过非单射 $\pi_S^*(2)$ 分解，这是不可能的。
因此 $\mathcal{I}|_{S_\etale}$ 不是内射 $\mathcal{O}_S$-模。
\end{proof}

\noindent
设 $f : T \to S$ 是概形态射。Topologies, Lemma
\ref{topologies-lemma-morphism-big-small-etale} (3) 中的交换图诱导出
下列环化网站交换图：
$$
\xymatrix{
(T_\etale, \mathcal{O}_T) \ar[d]_{f_{small}} &
((\Sch/T)_\etale, \mathcal{O}) \ar[d]^{f_{big}} \ar[l]^{\pi_T} \\
(S_\etale, \mathcal{O}_S) &
((\Sch/S)_\etale, \mathcal{O}) \ar[l]_{\pi_S}
}
$$
把 $f_{small}^\sharp$、$f_{big}^\sharp$、$\pi_S^\sharp$ 和
$\pi_T^\sharp$ 的定义写出来，就很容易看出这一点。特别地，这意味着
\begin{equation}
\label{etale-cohomology-equation-compare-big-small}
(f_{big, *}\mathcal{F})|_{S_\etale} =
f_{small, *}(\mathcal{F}|_{T_\etale})
\end{equation}
对任意层 $\mathcal{F}$（定义在 $(\Sch/T)_\etale$ 上）都成立；如果
$\mathcal{F}$ 是 $\mathcal{O}$-模层，则
(\ref{etale-cohomology-equation-compare-big-small}) 是 $\mathcal{O}_S$-模的同构，
定义在 $S_\etale$ 上。

\begin{lemma}
\label{etale-cohomology-lemma-compare-higher-direct-image}
设 $f : T \to S$ 是概形态射。
\begin{enumerate}
\item 对 $K$（属于 $D((\Sch/T)_\etale)$），有
$
(Rf_{big, *}K)|_{S_\etale} = Rf_{small, *}(K|_{T_\etale})
$
，该等式位于 $D(S_\etale)$ 中。
\item 对 $K$（属于 $D((\Sch/T)_\etale, \mathcal{O})$），有
$
(Rf_{big, *}K)|_{S_\etale} = Rf_{small, *}(K|_{T_\etale})
$
，该等式位于 $D(\textit{Mod}(S_\etale, \mathcal{O}_S))$ 中。
\end{enumerate}
更一般地，设 $g : S' \to S$ 是 $(\Sch/S)_\etale$ 的对象。考察纤维积
$$
\xymatrix{
T' \ar[r]_{g'} \ar[d]_{f'} & T \ar[d]^f \\
S' \ar[r]^g & S
}
$$
则
\begin{enumerate}
\item[(3)] 对 $K$（属于 $D((\Sch/T)_\etale)$），有
$i_g^{-1}(Rf_{big, *}K) = Rf'_{small, *}(i_{g'}^{-1}K)$
，该等式位于 $D(S'_\etale)$ 中。
\item[(4)] 对 $K$（属于 $D((\Sch/T)_\etale, \mathcal{O})$），有
$i_g^*(Rf_{big, *}K) = Rf'_{small, *}(i_{g'}^*K)$
，该等式位于 $D(\textit{Mod}(S'_\etale, \mathcal{O}_{S'}))$ 中。
\item[(5)] 对 $K$（属于 $D((\Sch/T)_\etale)$），有
$g_{big}^{-1}(Rf_{big, *}K) = Rf'_{big, *}((g'_{big})^{-1}K)$
，该等式位于 $D((\Sch/S')_\etale)$ 中。
\item[(6)] 对 $K$（属于 $D((\Sch/T)_\etale, \mathcal{O})$），有
$g_{big}^*(Rf_{big, *}K) = Rf'_{big, *}((g'_{big})^*K)$
，该等式位于 $D(\textit{Mod}(S'_\etale, \mathcal{O}_{S'}))$ 中。
\end{enumerate}
\end{lemma}

\begin{proof}
选取表示 $K$ 的阿贝尔层 K-内射复形后，(1) 由
Lemma \ref{etale-cohomology-lemma-compare-injectives} 和
(\ref{etale-cohomology-equation-compare-big-small}) 得到。

\medskip\noindent
选取表示 $K$ 的阿贝尔层 K-内射复形后，(3) 由
Lemma \ref{etale-cohomology-lemma-compare-injectives} 和 Topologies, Lemma
\ref{topologies-lemma-morphism-big-small-cartesian-diagram-etale} 得到。

\medskip\noindent
(5) 就是 Cohomology on Sites, Lemma
\ref{sites-cohomology-lemma-localize-cartesian-square}.

\medskip\noindent
(6) 就是 Cohomology on Sites, Lemma
\ref{sites-cohomology-lemma-localize-cartesian-square-modules}.

\medskip\noindent
(2) 可证如下。上面已经看到，作为环化网站的态射，有
$\pi_S \circ f_{big} = f_{small} \circ \pi_T$。因此由
Cohomology on Sites，引理
\ref{sites-cohomology-lemma-derived-pushforward-composition} 得到
$R\pi_{S, *} \circ Rf_{big, *} = Rf_{small, *} \circ R\pi_{T, *}$
。由于限制函子 $\pi_{S, *}$ 与 $\pi_{T, *}$ 都正合，结论随即得到。

\medskip\noindent
(4) 由 (6) 以及把 (2) 应用于 $f' : T' \to S'$ 得到。
\end{proof}

\noindent
设 $S$ 是概形，并设 $\mathcal{H}$ 是 $(\Sch/S)_\etale$ 上的阿贝尔层。
回顾，$H^n_\etale(U, \mathcal{H})$ 表示 $\mathcal{H}$ 在对象
$U$ 上的上同调，其中该对象属于 $(\Sch/S)_\etale$。

\begin{lemma}
\label{etale-cohomology-lemma-compare-cohomology}
设 $f : T \to S$ 是概形态射。则
\begin{enumerate}
\item 对 $K$（属于 $D(S_\etale)$），有
$H^n_\etale(S, \pi_S^{-1}K) = H^n(S_\etale, K)$。
\item 对 $K$（属于 $D(S_\etale, \mathcal{O}_S)$），有
$H^n_\etale(S, L\pi_S^*K) = H^n(S_\etale, K)$。
\item 对 $K$（属于 $D(S_\etale)$），有
$H^n_\etale(T, \pi_S^{-1}K) = H^n(T_\etale, f_{small}^{-1}K)$。
\item 对 $K$（属于 $D(S_\etale, \mathcal{O}_S)$），有
$H^n_\etale(T, L\pi_S^*K) = H^n(T_\etale, Lf_{small}^*K)$。
\item 对 $M$（属于 $D((\Sch/S)_\etale)$），有
$H^n_\etale(T, M) = H^n(T_\etale, i_f^{-1}M)$。
\item 对 $M$（属于 $D((\Sch/S)_\etale, \mathcal{O})$），有
$H^n_\etale(T, M) = H^n(T_\etale, i_f^*M)$。
\end{enumerate}
\end{lemma}

\begin{proof}
为证明 (5)，用阿贝尔层 K-内射复形表示 $M$，应用
Lemma \ref{etale-cohomology-lemma-compare-injectives}，再展开定义。由于
$i_f^{-1}\pi_S^{-1} = f_{small}^{-1}$，(3) 由此得到。(1) 是 (3) 的特例。

\medskip\noindent
(6) 由非常一般的 Cohomology on Sites, Lemma
\ref{sites-cohomology-lemma-pullback-same-cohomology} 得到。又因为
$Lf_{small}^* = i_f^* \circ L\pi_S^*$，所以 (4) 也随之得到。
(2) 是 (4) 的特例。
\end{proof}

\begin{lemma}
\label{etale-cohomology-lemma-cohomological-descent-etale}
设 $S$ 是概形。对 $K \in D(S_\etale)$，映射
$$
K \longrightarrow R\pi_{S, *}\pi_S^{-1}K
$$
是同构。
\end{lemma}

\begin{proof}
这是因为 $\pi_S^{-1}$ 与 $\pi_{S, *} = i_S^{-1}$ 都是正合函子，
而复合 $\pi_{S, *} \circ \pi_S^{-1}$ 是恒等函子。
\end{proof}

\begin{lemma}
\label{etale-cohomology-lemma-compare-higher-direct-image-proper}
设 $f : T \to S$ 是概形的固有态射。则有
\begin{enumerate}
\item $\pi_S^{-1} \circ f_{small, *} = f_{big, *} \circ \pi_T^{-1}$
，这是函子 $\Sh(T_\etale) \to \Sh((\Sch/S)_\etale)$ 的等式；
\item $\pi_S^{-1}Rf_{small, *}K = Rf_{big, *}\pi_T^{-1}K$
，其中 $K$ 属于 $D^+(T_\etale)$，且其上同调层为挠层；
\item $\pi_S^{-1}Rf_{small, *}K = Rf_{big, *}\pi_T^{-1}K$
，其中 $K$ 属于 $D(T_\etale, \mathbf{Z}/n\mathbf{Z})$；
\item $\pi_S^{-1}Rf_{small, *}K = Rf_{big, *}\pi_T^{-1}K$
，这对所有 $K$（属于 $D(T_\etale)$）都成立，只要 $f$ 有限。
\end{enumerate}
\end{lemma}

\begin{proof}
证明 (1)。设 $\mathcal{F}$ 是 $T_\etale$ 上的层。设
$g : S' \to S$ 是 $(\Sch/S)_\etale$ 的对象。考察纤维积
$$
\xymatrix{
T' \ar[r]_{f'} \ar[d]_{g'} & S' \ar[d]^g \\
T \ar[r]^f & S
}
$$
则有
$$
(f_{big, *}\pi_T^{-1}\mathcal{F})(S') =
(\pi_T^{-1}\mathcal{F})(T') =
((g'_{small})^{-1}\mathcal{F})(T')  =
(f'_{small, *}(g'_{small})^{-1}\mathcal{F})(S')
$$
，其中第二个等式来自 Lemma \ref{etale-cohomology-lemma-describe-pullback}。另一方面，
$$
(\pi_S^{-1}f_{small, *}\mathcal{F})(S') =
(g_{small}^{-1}f_{small, *}\mathcal{F})(S')
$$
，这里仍使用了 Lemma \ref{etale-cohomology-lemma-describe-pullback}。因此，由集合层的
固有基变换（Lemma \ref{etale-cohomology-lemma-proper-base-change-f-star}），
这两个集合典范同构。该同构与限制映射相容，并定义同构
$\pi_S^{-1}f_{small, *}\mathcal{F} = f_{big, *}\pi_T^{-1}\mathcal{F}$。
于是得到函子的同构
$\pi_S^{-1} \circ f_{small, *} = f_{big, *} \circ \pi_T^{-1}$。

\medskip\noindent
证明 (2)。存在典范基变换映射
$\pi_S^{-1}Rf_{small, *}K \to Rf_{big, *}\pi_T^{-1}K$，它对任意
$K$（属于 $D(T_\etale)$）都有定义，见 Cohomology on Sites, Remark
\ref{sites-cohomology-remark-base-change}。为证明它是同构，只须证明：
基变换映射沿
$i_g : \Sh(S'_\etale) \to \Sh((\Sch/S)_\etale)$ 的拉回是同构，
这对任意对象 $g : S' \to S$（属于 $(\Sch/S)_\etale$）都成立。
令 $T', g', f'$ 如上一段所述。基变换映射的拉回为
\begin{align*}
g_{small}^{-1}Rf_{small, *}K
& =
i_g^{-1}\pi_S^{-1}Rf_{small, *}K \\
& \to
i_g^{-1}Rf_{big, *}\pi_T^{-1}K \\
& =
Rf'_{small, *}(i_{g'}^{-1}\pi_T^{-1}K) \\
& =
Rf'_{small, *}((g'_{small})^{-1}K)
\end{align*}
；这里使用了 $\pi_S \circ i_g = g_{small}$、
$\pi_T \circ i_{g'} = g'_{small}$，以及
Lemma \ref{etale-cohomology-lemma-compare-higher-direct-image}。如果 $K$ 下有界且
$K$ 的上同调层为挠层，则由固有基变换定理
（Lemma \ref{etale-cohomology-lemma-proper-base-change}），该映射是同构。

\medskip\noindent
(3) 的证明与 (2) 完全相同，只须用 Lemma
\ref{etale-cohomology-lemma-proper-base-change-mod-n} 代替 Lemma
\ref{etale-cohomology-lemma-proper-base-change}。

\medskip\noindent
证明 (4)。如果 $f$ 有限，则函子 $f_{small, *}$ 与 $f_{big, *}$
都是正合的。对 $f_{small}$，这由 Proposition
\ref{etale-cohomology-proposition-finite-higher-direct-image-zero} 得到。由于任意基变换
$f'$（来自 $f$）也有限，Lemma
\ref{etale-cohomology-lemma-compare-higher-direct-image} (3) 说明 $f_{big, *}$
也正合（因为高阶导出函子为零）。所以这种情形由 (1) 得到。
\end{proof}




\section{fppf 拓扑与 \'etale 拓扑的比较}
\label{etale-cohomology-section-fppf-etale}

\noindent
本节的一个范本，是 Cohomology on Sites, Section
\ref{sites-cohomology-section-cohomology-LC} 中讨论的局部紧拓扑空间上通常拓扑与
qc 拓扑的比较。我们先回顾 Topologies, Sections
\ref{topologies-section-change-topologies} 和
\ref{topologies-section-etale} 中的一些材料。

\medskip\noindent
设 $S$ 是概形，并设 $(\Sch/S)_{fppf}$ 是一个 fppf 网站。
在同一个底层范畴上还有第二个拓扑，即 \'etale 拓扑，因而还有第二个网站
$(\Sch/S)_\etale$。恒等函子
$(\Sch/S)_\etale \to (\Sch/S)_{fppf}$ 是连续的，并定义网站态射
$$
\epsilon_S : (\Sch/S)_{fppf} \longrightarrow (\Sch/S)_\etale
$$
。见 Cohomology on Sites, Section \ref{sites-cohomology-section-compare}。
请注意，$\epsilon_{S, *}$ 在底层预层上是恒等函子，而
$\epsilon_S^{-1}$ 把一个 \'etale 层变为它的 fppf 层化。令
$S_\etale$ 为小 \'etale 网站。存在网站态射
$$
\pi_S : (\Sch/S)_\etale \longrightarrow S_\etale
$$
，它由连续函子
$S_\etale \to (\Sch/S)_\etale$, $U \mapsto U$ 给出。
事实上，$S_\etale$ 有纤维积和终对象，上述函子与它们可交换，因而可以应用
Sites, Proposition \ref{sites-proposition-get-morphism}。

\begin{lemma}
\label{etale-cohomology-lemma-describe-pullback-pi-fppf}
沿用上述记号。设 $\mathcal{F}$ 是 $S_\etale$ 上的层。规则
$$
(\Sch/S)_{fppf} \longrightarrow \textit{Sets},\quad
(f : X \to S) \longmapsto \Gamma(X, f_{small}^{-1}\mathcal{F})
$$
给出一个层，因而当然也给出 $(\Sch/S)_\etale$ 上的层。
事实上，这个层等于 $\pi_S^{-1}\mathcal{F}$（在
$(\Sch/S)_\etale$ 上），而等于
$\epsilon_S^{-1}\pi_S^{-1}\mathcal{F}$（在 $(\Sch/S)_{fppf}$ 上）。
\end{lemma}

\begin{proof}
关于 \'etale 拓扑的断言就是 Lemma \ref{etale-cohomology-lemma-describe-pullback} 的内容。
为完成证明，只须说明 $\pi_S^{-1}\mathcal{F}$ 对 fppf 拓扑是层。
这也已在 Lemma \ref{etale-cohomology-lemma-describe-pullback} 中证明。
\end{proof}

\noindent
在 Lemma \ref{etale-cohomology-lemma-describe-pullback-pi-fppf} 的情形中，
$\epsilon_S$ 与 $\pi_S$ 的复合以及上述等式确定网站态射
$$
a_S : (\Sch/S)_{fppf} \longrightarrow S_\etale
$$

\begin{lemma}
\label{etale-cohomology-lemma-push-pull-fppf-etale}
沿用上述记号。设 $f : X \to Y$ 是 $(\Sch/S)_{fppf}$ 中的态射。
则有拓扑斯的交换图
$$
\xymatrix{
\Sh((\Sch/X)_{fppf}) \ar[rr]_{f_{big, fppf}} \ar[d]_{\epsilon_X} & &
\Sh((\Sch/Y)_{fppf}) \ar[d]^{\epsilon_Y} \\
\Sh((\Sch/X)_\etale) \ar[rr]^{f_{big, \etale}} & &
\Sh((\Sch/Y)_\etale)
}
$$
以及
$$
\xymatrix{
\Sh((\Sch/X)_{fppf}) \ar[rr]_{f_{big, fppf}} \ar[d]_{a_X} & &
\Sh((\Sch/Y)_{fppf}) \ar[d]^{a_Y} \\
\Sh(X_\etale) \ar[rr]^{f_{small}} & &
\Sh(Y_\etale)
}
$$
，其中 $a_X = \pi_X \circ \epsilon_X$ 且 $a_Y = \pi_X \circ \epsilon_X$。
\end{lemma}

\begin{proof}
这些图的交换性来自 Topologies, Section
\ref{topologies-section-change-topologies} 中的讨论。
\end{proof}

\begin{lemma}
\label{etale-cohomology-lemma-proper-push-pull-fppf-etale}
在 Lemma \ref{etale-cohomology-lemma-push-pull-fppf-etale} 中，如果 $f$ 是固有的，则有
$a_Y^{-1} \circ f_{small, *} = f_{big, fppf, *} \circ a_X^{-1}$.
\end{lemma}

\begin{proof}
可以重复 Lemma \ref{etale-cohomology-lemma-compare-higher-direct-image-proper} (1) 的证明；
这里我们改为由该结论推出本结果。由于 $\epsilon_{Y, *}$ 在底层预层上是
恒等函子，它反映同构。Lemma \ref{etale-cohomology-lemma-describe-pullback-pi-fppf}
中的描述说明 $\epsilon_{Y, *} \circ a_Y^{-1} = \pi_Y^{-1}$，
对 $X$ 也类似。为证明典范映射
$a_Y^{-1}f_{small, *}\mathcal{F} \to f_{big, fppf, *}a_X^{-1}\mathcal{F}$
是同构，只须证明
\begin{align*}
\pi_Y^{-1}f_{small, *}\mathcal{F}
& =
\epsilon_{Y, *}a_Y^{-1}f_{small, *}\mathcal{F} \\
& \to 
\epsilon_{Y, *}f_{big, fppf, *}a_X^{-1}\mathcal{F} \\
& =
f_{big, \etale, *} \epsilon_{X, *}a_X^{-1}\mathcal{F} \\
& =
f_{big, \etale, *}\pi_X^{-1}\mathcal{F}
\end{align*}
是同构。这就是 Lemma \ref{etale-cohomology-lemma-compare-higher-direct-image-proper} (1)。
\end{proof}

\begin{lemma}
\label{etale-cohomology-lemma-descent-sheaf-fppf-etale}
在 Lemma \ref{etale-cohomology-lemma-push-pull-fppf-etale} 中，假设 $f$ 平坦、局部有限表现且满射。
则函子
$$
\Sh(Y_\etale) \longrightarrow
\left\{
(\mathcal{G}, \mathcal{H}, \alpha)
\middle|
\begin{matrix}
\mathcal{G} \in \Sh(X_\etale),\ \mathcal{H} \in \Sh((\Sch/Y)_{fppf}), \\
\alpha : a_X^{-1}\mathcal{G} \to f_{big, fppf}^{-1}\mathcal{H}
\text{ 为同构}
\end{matrix}
\right\}
$$
把 $\mathcal{F}$ 送到
$(f_{small}^{-1}\mathcal{F}, a_Y^{-1}\mathcal{F}, can)$，它是一个等价。
\end{lemma}

\begin{proof}
函子 $a_X^{-1}$ 完全忠实（由 Lemma
\ref{etale-cohomology-lemma-describe-pullback-pi-fppf}，有
$a_{X, *}a_X^{-1} = \text{id}$）。因此遗忘函子
$(\mathcal{G}, \mathcal{H}, \alpha) \mapsto \mathcal{H}$ 把三元组范畴
等同于 $\Sh((\Sch/Y)_{fppf})$ 的一个全子范畴。此外，函子
$a_Y^{-1}$ 完全忠实，故引理中的函子也完全忠实。

\medskip\noindent
假设有一个 \'etale 覆盖 $\{Y_i \to Y\}$。令
$f_i : X_i \to Y_i$ 为 $f$ 的基变换，并记
$f_{ij} = f_i \times f_j : X_i \times_X X_j  \to Y_i \times_Y Y_j$。
我们断言：如果引理对 $f_i$ 和 $f_{ij}$ 都成立（对所有 $i, j$），
那么它对 $f$ 也成立。为说明这一点，注意给定的 \'etale 覆盖在四个网站
$Y_\etale, X_\etale, (\Sch/Y)_{fppf}, (\Sch/X)_{fppf}$ 中都确定了
终对象的一个 \'etale 覆盖。因此在四种情形的每一种中，层范畴都等价于
这个覆盖的黏合数据范畴（Sites, Lemma
\ref{sites-lemma-mapping-property-glue}）。随后，一个庞大的范畴交换图完成
该断言的证明；我们略去细节。该断言表明，可以在 $Y$ 上 \'etale 局部地工作。

\medskip\noindent
注意，$\{X \to Y\}$ 是一个 fppf 覆盖。在 $Y$ 上 \'etale 局部地工作，
可以假设存在态射 $s : X' \to X$，使得复合
$f' = f \circ s : X' \to Y$ 满射、有限且局部自由，见
More on Morphisms，引理
\ref{more-morphisms-lemma-dominate-fppf-etale-locally}。
我们断言：如果引理对 $f'$ 成立，那么它对 $f$ 也成立。
事实上，给定一个三元组 $(\mathcal{G}, \mathcal{H}, \alpha)$（针对 $f$），
沿 $s$ 拉回便得到三元组
$(s_{small}^{-1}\mathcal{G}, \mathcal{H}, s_{big, fppf}^{-1}\alpha)$（针对 $f'$）。
这个三元组的一个解给出层 $\mathcal{F}$（定义在 $Y_\etale$ 上），满足
$a_Y^{-1}\mathcal{F} = \mathcal{H}$。由证明的第一段，这意味着该三元组
属于本质像。于是归约到下一段所述的情形。

\medskip\noindent
假设 $f$ 满射、有限且局部自由。设
$(\mathcal{G}, \mathcal{H}, \alpha)$ 是一个三元组。在此情形考虑三元组
$$
(\mathcal{G}_1, \mathcal{H}_1, \alpha_1) =
(f_{small}^{-1}f_{small, *}\mathcal{G},
f_{big, fppf, *}f_{big, fppf}^{-1}\mathcal{H}, \alpha_1)
$$
，其中 $\alpha_1$ 来自如下等同：
\begin{align*}
a_X^{-1}f_{small}^{-1}f_{small, *}\mathcal{G}
& =
f_{big, fppf}^{-1}a_Y^{-1}f_{small, *}\mathcal{G} \\
& =
f_{big, fppf}^{-1}f_{big, fppf, *}a_X^{-1}\mathcal{G} \\
& \to
f_{big, fppf}^{-1}f_{big, fppf, *}f_{big, fppf}^{-1}\mathcal{H}
\end{align*}
；其中第三个等式是 Lemma \ref{etale-cohomology-lemma-proper-push-pull-fppf-etale}，
箭头由 $\alpha$ 给出。这个三元组属于我们函子的像，因为
$\mathcal{F}_1 = f_{small, *}\mathcal{F}$ 是一个解（为此再次使用
Lemma \ref{etale-cohomology-lemma-proper-push-pull-fppf-etale}；细节从略）。存在三元组的典范映射
$$
(\mathcal{G}, \mathcal{H}, \alpha)
\to
(\mathcal{G}_1, \mathcal{H}_1, \alpha_1)
$$
，它在第二个分量上使用单位
$\text{id} \to f_{big, fppf, *}f_{big, fppf}^{-1}$（由证明的第一段，
只须在第二个分量上指定态射）。由于 $\{f : X \to Y\}$ 是 fppf 覆盖，
映射 $\mathcal{H} \to \mathcal{H}_1$ 是单射（细节从略）。置
$$
\mathcal{G}_2 = \mathcal{G}_1 \amalg_\mathcal{G} \mathcal{G}_1\quad
\mathcal{H}_2 = \mathcal{H}_1 \amalg_\mathcal{H} \mathcal{H}_1
$$
，并令 $\alpha_2$ 为诱导的同构（拉回函子是正合的，因此这是有意义的）。
于是 $\mathcal{H}$ 是两个映射
$\mathcal{H}_1 \to \mathcal{H}_2$ 的等化子。对三元组
$(\mathcal{G}_2, \mathcal{H}_2, \alpha_2)$ 重复上述论证，可得三元组的单态射
$$
(\mathcal{G}_2, \mathcal{H}_2, \alpha_2)
\to
(\mathcal{G}_3, \mathcal{H}_3, \alpha_3)
$$
，使得最后这个三元组属于我们函子的像。设它对应于
$\mathcal{F}_3$（属于 $\Sh(Y_\etale)$）。由完全忠实性，得到两个映射
$\mathcal{F}_1 \to \mathcal{F}_3$；令 $\mathcal{F}$ 为这两个映射的等化子。
由所涉及拉回函子的正合性，便有
$a_Y^{-1}\mathcal{F} = \mathcal{H}$，如所要求。
\end{proof}

\begin{lemma}
\label{etale-cohomology-lemma-compare-fppf-etale}
考虑比较态射
$\epsilon : (\Sch/S)_{fppf} \to (\Sch/S)_\etale$。
令 $\mathcal{P}$ 表示概形的有限态射类。对
$X$（属于 $(\Sch/S)_\etale$），记
$\mathcal{A}'_X \subset \textit{Ab}((\Sch/X)_\etale)$
为由形如 $\pi_X^{-1}\mathcal{F}$ 的层构成的全子范畴，其中
$\mathcal{F}$ 属于 $\textit{Ab}(X_\etale)$。则 Cohomology on Sites 中的性质
(\ref{sites-cohomology-item-base-change-P}),
(\ref{sites-cohomology-item-restriction-A}),
(\ref{sites-cohomology-item-A-sheaf}),
(\ref{sites-cohomology-item-A-and-P})，以及
(\ref{sites-cohomology-item-refine-tau-by-P})
（见 Cohomology on Sites, Situation
\ref{sites-cohomology-situation-compare}）都成立。
\end{lemma}

\begin{proof}
我们先验证 Homology, Lemma
\ref{homology-lemma-characterize-weak-serre-subcategory} 的条件
(1)、(2)、(3)、(4)，从而说明
$\mathcal{A}'_X \subset \textit{Ab}((\Sch/X)_\etale)$ 是弱 Serre 子范畴。
(1)、(2)、(3) 是立即的，因为例如由 Lemma
\ref{etale-cohomology-lemma-cohomological-descent-etale}，$\pi_X^{-1}$ 正合且完全忠实。
如果
$0 \to \pi_X^{-1}\mathcal{F} \to \mathcal{G} \to \pi_X^{-1}\mathcal{F}' \to 0$
是 $\textit{Ab}((\Sch/X)_\etale)$ 中的短正合列，则由 Lemma
\ref{etale-cohomology-lemma-cohomological-descent-etale}，
$0 \to \mathcal{F} \to \pi_{X, *}\mathcal{G} \to \mathcal{F}' \to 0$
正合。因此 $\mathcal{G} = \pi_X^{-1}\pi_{X, *}\mathcal{G}$ 属于
$\mathcal{A}'_X$，这就验证了最后一个条件。

\medskip\noindent
Cohomology on Sites, Property (\ref{sites-cohomology-item-base-change-P})
成立，因为概形有纤维积，而且概形的有限态射经基变换后仍是有限态射；
见 Morphisms, Lemma \ref{morphisms-lemma-base-change-finite}。

\medskip\noindent
Cohomology on Sites, Property (\ref{sites-cohomology-item-restriction-A})
来自 Topologies, Lemma
\ref{topologies-lemma-morphism-big-small-etale} 中的交换图 (3)。

\medskip\noindent
Cohomology on Sites, Property (\ref{sites-cohomology-item-A-sheaf})
就是 Lemma \ref{etale-cohomology-lemma-describe-pullback-pi-fppf}。

\medskip\noindent
Cohomology on Sites, Property (\ref{sites-cohomology-item-A-and-P})
由 Lemma \ref{etale-cohomology-lemma-compare-higher-direct-image-proper} (4) 成立。

\medskip\noindent
Cohomology on Sites, Property (\ref{sites-cohomology-item-refine-tau-by-P})
由 More on Morphisms, Lemma
\ref{more-morphisms-lemma-dominate-fppf-etale-locally} 推出。
\end{proof}

\begin{lemma}
\label{etale-cohomology-lemma-V-C-all-n-etale-fppf}
沿用上述记号。
\begin{enumerate}
\item 对 $X \in \Ob((\Sch/S)_{fppf})$ 和阿贝尔层
$\mathcal{F}$（定义在 $X_\etale$ 上），有
$\epsilon_{X, *}a_X^{-1}\mathcal{F} = \pi_X^{-1}\mathcal{F}$
，并且 $R^i\epsilon_{X, *}(a_X^{-1}\mathcal{F}) = 0$ 对 $i > 0$ 成立。
\item 对有限态射 $f : X \to Y$（属于 $(\Sch/S)_{fppf}$）和
阿贝尔层 $\mathcal{F}$（定义在 $X$ 上），有
$a_Y^{-1}(R^if_{small, *}\mathcal{F}) =
R^if_{big, fppf, *}(a_X^{-1}\mathcal{F})$
，对所有 $i$ 成立。
\item 对概形 $X$ 和 $K$（属于 $D^+(X_\etale)$），映射
$\pi_X^{-1}K \to R\epsilon_{X, *}(a_X^{-1}K)$ 是同构。
\item 对概形的有限态射 $f : X \to Y$ 和 $K$（属于
$D^+(X_\etale)$），有
$a_Y^{-1}(Rf_{small, *}K) = Rf_{big, fppf, *}(a_X^{-1}K)$.
\item 对概形的固有态射 $f : X \to Y$ 和 $K$（属于
$D^+(X_\etale)$，且上同调层为挠层），有
$a_Y^{-1}(Rf_{small, *}K) = Rf_{big, fppf, *}(a_X^{-1}K)$.
\end{enumerate}
\end{lemma}

\begin{proof}
由 Lemma \ref{etale-cohomology-lemma-compare-fppf-etale}，Cohomology on Sites, Section
\ref{sites-cohomology-section-compare-general} 中的各引理都适用于当前情形。
为翻译这些结果，注意 Cohomology on Sites, Lemma
\ref{sites-cohomology-lemma-A} 中的范畴 $\mathcal{A}_X$，正是
$a_X^{-1} : \textit{Ab}(X_\etale) \to \textit{Ab}((\Sch/X)_{fppf})$
的本质像。

\medskip\noindent
(1) 等价于 $(V_n)$ 对所有 $n$ 都成立；这由 Cohomology on Sites, Lemma
\ref{sites-cohomology-lemma-V-C-all-n-general} 成立。

\medskip\noindent
(2) 来自对 Cohomology on Sites, Lemma
\ref{sites-cohomology-lemma-V-implies-C-general} 的结论应用
$\epsilon_Y^{-1}$。

\medskip\noindent
(3) 来自 Cohomology on Sites, Lemma
\ref{sites-cohomology-lemma-V-C-all-n-general} (1)，因为
$\pi_X^{-1}K$ 属于 $D^+_{\mathcal{A}'_X}((\Sch/X)_\etale)$，并且
$a_X^{-1} = \epsilon_X^{-1} \circ a_X^{-1}$。

\medskip\noindent
出于同样的理由，(4) 来自 Cohomology on Sites, Lemma
\ref{sites-cohomology-lemma-V-C-all-n-general} (2)。

\medskip\noindent
证明 (5)。我们使用
\begin{align*}
R\epsilon_{Y, *}Rf_{big, fppf, *}a_X^{-1}K
& =
Rf_{big, \etale, *}R\epsilon_{X, *}a_X^{-1}K \\
& =
Rf_{big, \etale, *}\pi_X^{-1}K \\
& =
\pi_Y^{-1}Rf_{small, *}K \\
& =
R\epsilon_{Y, *} a_Y^{-1}Rf_{small, *}K
\end{align*}
。第一个等式来自 Lemma \ref{etale-cohomology-lemma-push-pull-fppf-etale} 中的交换图以及
Cohomology on Sites，引理
\ref{sites-cohomology-lemma-derived-pushforward-composition}。
第二个等式是 (3)。第三个等式是 Lemma
\ref{etale-cohomology-lemma-compare-higher-direct-image-proper} (2)。第四个等式再次使用 (3)。
因此基变换映射
$a_Y^{-1}(Rf_{small, *}K) \to Rf_{big, fppf, *}(a_X^{-1}K)$
诱导同构
$$
R\epsilon_{Y, *}a_Y^{-1}Rf_{small, *}K \to
R\epsilon_{Y, *}Rf_{big, fppf, *}a_X^{-1}K
$$
。以下观察完成证明：设映射 $\alpha : a_Y^{-1}L \to M$ 中
$L$ 属于 $D^+(Y_\etale)$，$M$ 属于 $D^+((\Sch/Y)_{fppf})$，并且
$R\epsilon_{Y, *}\alpha$ 是同构，则该映射是同构。事实上，
我们对 $i$ 归纳证明 $H^i(\alpha)$ 是同构。对所有充分小的 $i$，
这成立。若它对 $i \leq i_0$ 成立，那么可知
$R^j\epsilon_{Y, *}H^i(M) = 0$（其中 $j > 0$ 且 $i \leq i_0$）；
这是由 (1) 得出的，因为在这个范围内
$H^i(M) = a_Y^{-1}H^i(L)$。于是由谱序列论证，
$\epsilon_{Y, *}H^{i_0 + 1}(M) = H^{i_0 + 1}(R\epsilon_{Y, *}M)$。
因此
$\epsilon_{Y, *}H^{i_0 + 1}(M) = \pi_Y^{-1}H^{i_0 + 1}(L) =
\epsilon_{Y, *}a_Y^{-1}H^{i_0 + 1}(L)$。
这说明 $H^{i_0 + 1}(\alpha)$ 是同构（因为 $\epsilon_{Y, *}$
在底层预层上是恒等函子，所以它反映同构），如所要求。
\end{proof}

\begin{lemma}
\label{etale-cohomology-lemma-cohomological-descent-etale-fppf}
设 $X$ 是概形。对 $K \in D^+(X_\etale)$，映射
$$
K \longrightarrow Ra_{X, *}a_X^{-1}K
$$
是同构，其中
$a_X : \Sh((\Sch/X)_{fppf}) \to \Sh(X_\etale)$ 如上所述。
\end{lemma}

\begin{proof}
我们先把断言归约到 $K$ 由单个阿贝尔层给出的情形。也就是说，取
$K$ 的一个下有界复形表示 $\mathcal{F}^\bullet$。由单个层的情形可知
$\mathcal{F}^n = a_{X, *} a_X^{-1} \mathcal{F}^n$
，而且层 $R^qa_{X, *}a_X^{-1}\mathcal{F}^n$ 对 $q > 0$ 为零。
对 $a_X^{-1}\mathcal{F}^\bullet$ 和函子 $a_{X, *}$ 应用 Leray 无环性引理
（Derived Categories, Lemma \ref{derived-lemma-leray-acyclicity}），
便得到结论。以下假设 $K = \mathcal{F}$。

\medskip\noindent
由 Lemma \ref{etale-cohomology-lemma-describe-pullback-pi-fppf}，有
$a_{X, *}a_X^{-1}\mathcal{F} = \mathcal{F}$。因此只须证明
$R^qa_{X, *}a_X^{-1}\mathcal{F} = 0$ 对 $q > 0$ 成立。为此可以使用
$a_X = \epsilon_X \circ \pi_X$ 和 Leray 谱序列
（Cohomology on Sites, Lemma \ref{sites-cohomology-lemma-relative-Leray}）。
由 Lemma \ref{etale-cohomology-lemma-V-C-all-n-etale-fppf}，有
$R^i\epsilon_{X, *}(a_X^{-1}\mathcal{F}) = 0$（当 $i > 0$），并且
$\epsilon_{X, *}a_X^{-1}\mathcal{F} = \pi_X^{-1}\mathcal{F}$。
由 Lemma \ref{etale-cohomology-lemma-cohomological-descent-etale}，有
$R^j\pi_{X, *}(\pi_X^{-1}\mathcal{F}) = 0$（当 $j > 0$）。证明完毕。
\end{proof}

\begin{lemma}
\label{etale-cohomology-lemma-compare-cohomology-etale-fppf}
设 $X$ 是概形，并设
$a_X : \Sh((\Sch/X)_{fppf}) \to \Sh(X_\etale)$ 如上。则：
\begin{enumerate}
\item $H^q(X_\etale, \mathcal{F}) = H^q_{fppf}(X, a_X^{-1}\mathcal{F})$
，其中 $\mathcal{F}$ 是 $X_\etale$ 上的阿贝尔层；
\item $H^q(X_\etale, K) = H^q_{fppf}(X, a_X^{-1}K)$，其中
$K \in D^+(X_\etale)$。
\end{enumerate}
例如，如果 $A$ 是阿贝尔群，则
$H^q_\etale(X, \underline{A}) = H^q_{fppf}(X, \underline{A})$.
\end{lemma}

\begin{proof}
由 Lemma \ref{etale-cohomology-lemma-cohomological-descent-etale-fppf} 和
Cohomology on Sites, Remark \ref{sites-cohomology-remark-before-Leray}
即得此结论。
\end{proof}







\section{fppf 拓扑与 \'etale 拓扑的比较：模}
\label{etale-cohomology-section-fppf-etale-modules}

\noindent
我们继续 Section \ref{etale-cohomology-section-fppf-etale} 中的讨论，但在本节简要讨论模层的情形。

\medskip\noindent
设 $S$ 是概形。Section \ref{etale-cohomology-section-fppf-etale} 中引入的网站态射
$\epsilon_S$、$\pi_S$ 及其复合 $a_S$，都有向环化网站态射的自然提升。
第一个记作
$$
\epsilon_S :
((\Sch/S)_{fppf}, \mathcal{O})
\longrightarrow
((\Sch/S)_\etale, \mathcal{O})
$$
。注意，我们可以对结构层使用相同的符号，因为这些层确有相同的底层预层。
第二个是
$$
\pi_S :
((\Sch/S)_\etale, \mathcal{O})
\longrightarrow
(S_\etale, \mathcal{O}_S)
$$
。第三个是态射
$$
a_S :
((\Sch/S)_{fppf}, \mathcal{O})
\longrightarrow
(S_\etale, \mathcal{O}_S)
$$
。我们已经知道，概形 $S$ 上的拟凝聚模范畴与
$(S_\etale, \mathcal{O}_S)$ 上的拟凝聚模范畴相同，见
Descent, Proposition \ref{descent-proposition-equivalence-quasi-coherent}。
由于我们关心的是陈述 \'etale 上同调与 fppf 上同调之间的比较，
在本节余下部分，我们将通过小 \'etale 网站来理解拟凝聚层。
下面回顾关于这些网站上拟凝聚模的已有结果。

\begin{lemma}
\label{etale-cohomology-lemma-review-quasi-coherent}
设 $S$ 是概形。设 $\mathcal{F}$ 是拟凝聚 $\mathcal{O}_S$-模
（定义在 $S_\etale$ 上）。
\begin{enumerate}
\item 规则
$$
\mathcal{F}^a : (\Sch/S)_\etale \longrightarrow \textit{Ab},\quad
(f : T \to S) \longmapsto \Gamma(T, f_{small}^*\mathcal{F})
$$
对 fppf 覆盖满足层条件，因而当然也对 \'etale 覆盖满足层条件；
\item $\mathcal{F}^a = \pi_S^*\mathcal{F}$ 在 $(\Sch/S)_\etale$ 上成立；
\item $\mathcal{F}^a = a_S^*\mathcal{F}$ 在 $(\Sch/S)_{fppf}$ 上成立；
\item 规则 $\mathcal{F} \mapsto \mathcal{F}^a$ 定义拟凝聚
$\mathcal{O}_S$-模与
$((\Sch/S)_\etale, \mathcal{O})$ 上拟凝聚模之间的等价；
\item 规则 $\mathcal{F} \mapsto \mathcal{F}^a$ 定义拟凝聚
$\mathcal{O}_S$-模与
$((\Sch/S)_{fppf}, \mathcal{O})$ 上拟凝聚模之间的等价；
\item 有 $\epsilon_{S, *}a_S^*\mathcal{F} = \pi_S^*\mathcal{F}$
和 $a_{S, *}a_S^*\mathcal{F} = \mathcal{F}$；
\item 有 $R^i\epsilon_{S, *}(a_S^*\mathcal{F}) = 0$
和 $R^ia_{S, *}(a_S^*\mathcal{F}) = 0$，其中 $i > 0$。
\end{enumerate}
\end{lemma}

\begin{proof}
我们强烈建议读者根据 Descent, Sections
\ref{descent-section-quasi-coherent-sheaves}、
\ref{descent-section-quasi-coherent-cohomology} 和
\ref{descent-section-quasi-coherent-sheaves-bis} 中的材料，自己找出这些结果的证明。

\medskip\noindent
我们先说明，本引理中的记号 $\mathcal{F}^a$ 为何与我们先前在 Sections
\ref{etale-cohomology-section-quasi-coherent}、\ref{etale-cohomology-section-cohomology-quasi-coherent}
以及 Descent, Section \ref{descent-section-quasi-coherent-sheaves}
中使用的记号相容。事实上，由 Descent, Proposition
\ref{descent-proposition-equivalence-quasi-coherent} 可知，存在一个拟凝聚模
$\mathcal{F}_0$，定义在概形 $S$ 上（换言之，小 Zariski 网站上），
使得 $\mathcal{F}$ 是下述规则的限制：
$$
\mathcal{F}_0^a : (\Sch/S)_\etale \longrightarrow \textit{Ab},\quad
(f : T \to S) \longmapsto \Gamma(T, f^*\mathcal{F})
$$
到子范畴 $S_\etale \subset (\Sch/S)_\etale$；这里 $f^*$ 表示概形上
模层的通常拉回。由于 $\mathcal{F}_0^a$ 是沿环化网站态射
$$
((\Sch/S)_\etale, \mathcal{O}) \longrightarrow (S_{Zar}, \mathcal{O}_{S_{Zar}})
$$
的拉回（Descent, Remark
\ref{descent-remark-change-topologies-ringed-sites}），由拉回的复合立即得到
$\mathcal{F}^a = \mathcal{F}_0^a$。Descent, Lemma
\ref{descent-lemma-sheaf-condition-holds} 于是证明，即使对 fpqc 覆盖也满足
层条件（另见 Proposition \ref{etale-cohomology-proposition-quasi-coherent-sheaf-fpqc}）。
再由 Descent, Remark \ref{descent-remark-change-topologies-ringed-sites}
得到 (2) 和 (3)，由 Descent, Proposition
\ref{descent-proposition-equivalence-quasi-coherent} 得到 (4) 和 (5)
（另见元结果 Theorem \ref{etale-cohomology-theorem-quasi-coherent}）。

\medskip\noindent
(6) 立即来自层 $\mathcal{F}^a = \pi_S^*\mathcal{F} = a_S^*\mathcal{F}$
的描述。

\medskip\noindent
对任意阿贝尔层 $\mathcal{H}$（定义在 $(\Sch/S)_{fppf}$ 上），高阶直接像
$R^p\epsilon_{S, *}\mathcal{H}$ 是预层
$U \mapsto H^p_{fppf}(U, \mathcal{H})$ 所关联的层（定义在
$(\Sch/S)_\etale$ 上）。见 Cohomology on Sites,
Lemma \ref{sites-cohomology-lemma-higher-direct-images}。因此，为证明
$R^p\epsilon_{S, *}a_S^*\mathcal{F} = R^p\epsilon_{S, *}\mathcal{F}^a = 0$
对 $p > 0$ 成立，只须证明任意概形 $U$（在 $S$ 上）都有 \'etale 覆盖
$\{U_i \to U\}_{i \in I}$，使得
$H^p_{fppf}(U_i, \mathcal{F}^a) = 0$ 对 $p > 0$ 成立。
如果取仿射开覆盖，则所需消灭性来自与通常上同调的比较
（Descent, Proposition \ref{descent-proposition-same-cohomology-quasi-coherent}
或 Theorem \ref{etale-cohomology-theorem-zariski-fpqc-quasi-coherent}），以及 Cohomology of
Schemes, Lemma \ref{coherent-lemma-quasi-coherent-affine-cohomology-zero}
所给出的仿射概形上拟凝聚层上同调的消灭性。

\medskip\noindent
为证明 $R^pa_{S, *}a_S^{-1}\mathcal{F} = R^pa_{S, *}\mathcal{F}^a = 0$
对 $p > 0$ 成立，完全同样地论证即可。证明完毕。
\end{proof}

\begin{lemma}
\label{etale-cohomology-lemma-cohomological-descent-etale-fppf-modules}
设 $S$ 是概形。对 $\mathcal{F}$（一个拟凝聚 $\mathcal{O}_S$-模，定义在
$S_\etale$ 上），映射
$$
\pi_S^*\mathcal{F} \longrightarrow R\epsilon_{S, *}(a_S^*\mathcal{F})
\quad\text{且}\quad
\mathcal{F} \longrightarrow Ra_{S, *}(a_S^*\mathcal{F})
$$
都是同构，其中
$a_S : \Sh((\Sch/S)_{fppf}) \to \Sh(S_\etale)$ 如上所述。
\end{lemma}

\begin{proof}
这立即来自 Lemma \ref{etale-cohomology-lemma-review-quasi-coherent} 的 (6) 和 (7)。
\end{proof}

\begin{lemma}
\label{etale-cohomology-lemma-cohomological-descent-complex-modules}
设 $S = \Spec(A)$ 是仿射概形。设 $M^\bullet$ 是 $A$-模复形。
考虑预层复形 $\mathcal{F}^\bullet$，其为 $\mathcal{O}$-模，定义在
$(\textit{Aff}/S)_{fppf}$ 上，并由规则
$$
(U/S) = (\Spec(B)/\Spec(A)) \longmapsto M^\bullet \otimes_A B
$$
给出。这是一个模复形，并且典范映射
$$
M^\bullet \longrightarrow
R\Gamma((\textit{Aff}/S)_{fppf}, \mathcal{F}^\bullet)
$$
是拟同构。
\end{lemma}

\begin{proof}
每个 $\mathcal{F}^n$ 都是模层，因为它等于 Lemma
\ref{etale-cohomology-lemma-review-quasi-coherent} 中的模
$\mathcal{G}^n = (\widetilde{M}^n)^a$ 在
$(\textit{Aff}/S)_{fppf} \subset (\Sch/S)_{fppf}$ 上的限制。
由于这个包含定义环化拓扑斯的等价（Topologies, Lemma
\ref{topologies-lemma-affine-big-site-fppf}），有
$$
R\Gamma((\textit{Aff}/S)_{fppf}, \mathcal{F}^\bullet) =
R\Gamma((\Sch/S)_{fppf}, \mathcal{G}^\bullet)
$$
。注意，例如由 Derived Categories of Schemes, Lemma
\ref{perfect-lemma-affine-compare-bounded}，有
$M^\bullet = R\Gamma(S, \widetilde{M}^\bullet)$。因此要证明比较映射
$$
R\Gamma(S, \widetilde{M}^\bullet)
\longrightarrow
R\Gamma((\Sch/S)_{fppf}, (\widetilde{M}^\bullet)^a)
$$
是同构。如果 $M^\bullet$ 下有界，则这由 Descent, Proposition
\ref{descent-proposition-same-cohomology-quasi-coherent} 和 Derived
Categories, Lemma \ref{derived-lemma-two-ss-complex-functor} 的第一个谱序列得到。
对一般情形，写成 $M^\bullet = \lim M_n^\bullet$，其中
$M_n^\bullet = \tau_{\geq -n}M^\bullet$。于是系统 $M_n^p$ 最终恒等于
$M^p$。我们断言
$$
(\widetilde{M}^\bullet)^a = R\lim (\widetilde{M}_n^\bullet)^a
$$
。也就是说，只须证明从左到右的自然映射在任意仿射对象
$U = \Spec(B)$（属于 $(\Sch/S)_{fppf}$）上诱导上同调同构。
对 $i \in \mathbf{Z}$ 和 $n > |i|$，有
$$
H^i(U, (\widetilde{M}_n^\bullet)^a) =
H^i(\tau_{\geq -n}M^\bullet \otimes_A B) =
H^i(M^\bullet \otimes_A B)
$$
。第一个等式由上面处理过的下有界情形成立。因此由 Cohomology on Sites,
Lemma \ref{sites-cohomology-lemma-RGamma-commutes-with-Rlim} 可知该断言成立。
最后得到
\begin{align*}
R\Gamma((\Sch/S)_{fppf}, (\widetilde{M}^\bullet)^a)
& =
R\Gamma((\Sch/S)_{fppf}, R\lim (\widetilde{M}_n^\bullet)^a) \\
& =
R\lim R\Gamma((\Sch/S)_{fppf}, (\widetilde{M}_n^\bullet)^a) \\
& =
R\lim M_n^\bullet \\
& =
M^\bullet
\end{align*}
，如所要求。第二个等式成立是因为 $R\lim$ 与 $R\Gamma$ 可交换，见
Cohomology on Sites，引理
\ref{sites-cohomology-lemma-RGamma-commutes-with-Rlim}.
\end{proof}










\section{ph 拓扑与 \'etale 拓扑的比较}
\label{etale-cohomology-section-ph-etale}

\noindent
本节的一个范本，是 Cohomology on Sites, Section
\ref{sites-cohomology-section-cohomology-LC} 中讨论的局部紧拓扑空间上通常拓扑与
qc 拓扑的比较。我们先回顾 Topologies, Sections
\ref{topologies-section-change-topologies} 和
\ref{topologies-section-etale} 中的一些材料。

\medskip\noindent
设 $S$ 是概形，并设 $(\Sch/S)_{ph}$ 是一个 ph 网站。
在同一个底层范畴上还有第二个拓扑，即 \'etale 拓扑，因而还有第二个网站
$(\Sch/S)_\etale$。恒等函子
$(\Sch/S)_\etale \to (\Sch/S)_{ph}$ 是连续的
（由 More on Morphisms, Lemma \ref{more-morphisms-lemma-fppf-ph} 和
Topologies，引理
\ref{topologies-lemma-zariski-etale-smooth-syntomic-fppf}），并定义网站态射
$$
\epsilon_S : (\Sch/S)_{ph} \longrightarrow (\Sch/S)_\etale
$$
。见 Cohomology on Sites, Section \ref{sites-cohomology-section-compare}。
请注意，$\epsilon_{S, *}$ 在底层预层上是恒等函子，而
$\epsilon_S^{-1}$ 把一个 \'etale 层变为它的 ph 层化。
令 $S_\etale$ 为小 \'etale 网站。存在网站态射
$$
\pi_S : (\Sch/S)_\etale \longrightarrow S_\etale
$$
，它由连续函子
$S_\etale \to (\Sch/S)_\etale$, $U \mapsto U$ 给出。
事实上，$S_\etale$ 有纤维积和终对象，上述函子与它们可交换，因而可以应用
Sites, Proposition \ref{sites-proposition-get-morphism}。

\begin{lemma}
\label{etale-cohomology-lemma-describe-pullback-pi-ph}
沿用上述记号。设 $\mathcal{F}$ 是 $S_\etale$ 上的层。规则
$$
(\Sch/S)_{ph} \longrightarrow \textit{Sets},\quad
(f : X \to S) \longmapsto \Gamma(X, f_{small}^{-1}\mathcal{F})
$$
给出一个层，因而当然也给出 $(\Sch/S)_\etale$ 上的层。
事实上，这个层等于 $\pi_S^{-1}\mathcal{F}$（在
$(\Sch/S)_\etale$ 上），而等于
$\epsilon_S^{-1}\pi_S^{-1}\mathcal{F}$（在 $(\Sch/S)_{ph}$ 上）。
\end{lemma}

\begin{proof}
关于 \'etale 拓扑的断言就是 Lemma \ref{etale-cohomology-lemma-describe-pullback} 的内容。
为完成证明，只须说明 $\pi_S^{-1}\mathcal{F}$ 对 ph 拓扑是层。
由 Topologies, Lemma \ref{topologies-lemma-characterize-sheaf}，
只须说明：给定固有满射态射 $V \to U$（其为 $S$ 上概形之间的态射），
有等化子图
$$
\xymatrix{
(\pi_S^{-1}\mathcal{F})(U) \ar[r] &
(\pi_S^{-1}\mathcal{F})(V) \ar@<1ex>[r] \ar@<-1ex>[r] &
(\pi_S^{-1}\mathcal{F})(V \times_U V)
}
$$
。置 $\mathcal{G} = \pi_S^{-1}\mathcal{F}|_{U_\etale}$。考虑交换图
$$
\xymatrix{
V \times_U V \ar[r] \ar[rd]_g \ar[d] & V \ar[d]^f \\
V \ar[r]^f & U
}
$$
。我们有
$$
(\pi_S^{-1}\mathcal{F})(V) = \Gamma(V, f^{-1}\mathcal{G}) =
\Gamma(U, f_*f^{-1}\mathcal{G})
$$
；这里用 $f_*$ 和 $f^{-1}$ 表示小 \'etale 网站之间的函子性。其次，有
$$
(\pi_S^{-1}\mathcal{F})(V \times_U V) =
\Gamma(V \times_U V, g^{-1}\mathcal{G}) =
\Gamma(U, g_*g^{-1}\mathcal{G})
$$
。等化子图中的两个映射来自两个映射
$$
f_*f^{-1}\mathcal{G} \longrightarrow g_*g^{-1}\mathcal{G}
$$
。因此只须证明 $\mathcal{G}$ 是这两个层映射的等化子。
设 $\overline{u}$ 是 $U$ 的几何点。置
$\Omega = \mathcal{G}_{\overline{u}}$。由 Lemma
\ref{etale-cohomology-lemma-proper-pushforward-stalk} 在 $\overline{u}$ 处取茎，得到两个映射
$$
H^0(V_{\overline{u}}, \underline{\Omega}) \longrightarrow
H^0((V \times_U V)_{\overline{u}}, \underline{\Omega}) =
H^0(V_{\overline{u}} \times_{\overline{u}} V_{\overline{u}},
\underline{\Omega})
$$
；其中 $\underline{\Omega}$ 表示取值为 $\Omega$ 的常值层。当然，
这些映射是沿投影映射的拉回。显然，由沿到第一因子的投影拉回得到的截面，
在第一投影的纤维上为常值；而由沿到第一因子的投影拉回得到的截面，
在第一投影的纤维上为常值。这些拉回映射之像的交中的截面，在整个
$V_{\overline{u}} \times_{\overline{u}} V_{\overline{u}}$ 上为常值；
也就是说，它们来自 $\Omega$ 中的元素，如所要求。
\end{proof}

\noindent
在 Lemma \ref{etale-cohomology-lemma-describe-pullback-pi-ph} 的情形中，
$\epsilon_S$ 与 $\pi_S$ 的复合以及上述等式确定网站态射
$$
a_S : (\Sch/S)_{ph} \longrightarrow S_\etale
$$

\begin{lemma}
\label{etale-cohomology-lemma-push-pull-ph-etale}
沿用上述记号。设 $f : X \to Y$ 是 $(\Sch/S)_{ph}$ 中的态射。
则有拓扑斯的交换图
$$
\xymatrix{
\Sh((\Sch/X)_{ph}) \ar[rr]_{f_{big, ph}} \ar[d]_{\epsilon_X} & &
\Sh((\Sch/Y)_{ph}) \ar[d]^{\epsilon_Y} \\
\Sh((\Sch/X)_\etale) \ar[rr]^{f_{big, \etale}} & &
\Sh((\Sch/Y)_\etale)
}
$$
以及
$$
\xymatrix{
\Sh((\Sch/X)_{ph}) \ar[rr]_{f_{big, ph}} \ar[d]_{a_X} & &
\Sh((\Sch/Y)_{ph}) \ar[d]^{a_Y} \\
\Sh(X_\etale) \ar[rr]^{f_{small}} & &
\Sh(Y_\etale)
}
$$
，其中 $a_X = \pi_X \circ \epsilon_X$ 且 $a_Y = \pi_X \circ \epsilon_X$。
\end{lemma}

\begin{proof}
这些图的交换性来自 Topologies, Section
\ref{topologies-section-change-topologies} 中的讨论。
\end{proof}

\begin{lemma}
\label{etale-cohomology-lemma-proper-push-pull-ph-etale}
在 Lemma \ref{etale-cohomology-lemma-push-pull-ph-etale} 中，如果 $f$ 是固有的，则有
$a_Y^{-1} \circ f_{small, *} = f_{big, ph, *} \circ a_X^{-1}$.
\end{lemma}

\begin{proof}
可以重复 Lemma \ref{etale-cohomology-lemma-compare-higher-direct-image-proper} (1) 的证明；
这里我们改为由该结论推出本结果。由于 $\epsilon_{Y, *}$ 在底层预层上是
恒等函子，它反映同构。Lemma \ref{etale-cohomology-lemma-describe-pullback-pi-ph}
中的描述说明 $\epsilon_{Y, *} \circ a_Y^{-1} = \pi_Y^{-1}$，
对 $X$ 也类似。为证明典范映射
$a_Y^{-1}f_{small, *}\mathcal{F} \to f_{big, ph, *}a_X^{-1}\mathcal{F}$
是同构，只须证明
\begin{align*}
\pi_Y^{-1}f_{small, *}\mathcal{F}
& =
\epsilon_{Y, *}a_Y^{-1}f_{small, *}\mathcal{F} \\
& \to 
\epsilon_{Y, *}f_{big, ph, *}a_X^{-1}\mathcal{F} \\
& =
f_{big, \etale, *} \epsilon_{X, *}a_X^{-1}\mathcal{F} \\
& =
f_{big, \etale, *}\pi_X^{-1}\mathcal{F}
\end{align*}
是同构。这就是 Lemma \ref{etale-cohomology-lemma-compare-higher-direct-image-proper} (1)。
\end{proof}

\begin{lemma}
\label{etale-cohomology-lemma-compare-ph-etale}
考虑比较态射
$\epsilon : (\Sch/S)_{ph} \to (\Sch/S)_\etale$。
令 $\mathcal{P}$ 表示概形的固有态射类。对
$X$（属于 $(\Sch/S)_\etale$），记
$\mathcal{A}'_X \subset \textit{Ab}((\Sch/X)_\etale)$
为由形如 $\pi_X^{-1}\mathcal{F}$ 的层构成的全子范畴，其中
$\mathcal{F}$ 是 $X_\etale$ 上的挠阿贝尔层。则 Cohomology on Sites 中的性质
(\ref{sites-cohomology-item-base-change-P}),
(\ref{sites-cohomology-item-restriction-A}),
(\ref{sites-cohomology-item-A-sheaf}),
(\ref{sites-cohomology-item-A-and-P})，以及
(\ref{sites-cohomology-item-refine-tau-by-P})
（见 Cohomology on Sites, Situation
\ref{sites-cohomology-situation-compare}）都成立。
\end{lemma}

\begin{proof}
我们先验证 Homology, Lemma
\ref{homology-lemma-characterize-weak-serre-subcategory} 的条件
(1)、(2)、(3)、(4)，从而说明
$\mathcal{A}'_X \subset \textit{Ab}((\Sch/X)_\etale)$ 是弱 Serre 子范畴。
(1)、(2)、(3) 是立即的，因为例如由 Lemma
\ref{etale-cohomology-lemma-cohomological-descent-etale}，$\pi_X^{-1}$ 正合且完全忠实。
如果
$0 \to \pi_X^{-1}\mathcal{F} \to \mathcal{G} \to \pi_X^{-1}\mathcal{F}' \to 0$
是 $\textit{Ab}((\Sch/X)_\etale)$ 中的短正合列，则由 Lemma
\ref{etale-cohomology-lemma-cohomological-descent-etale}，
$0 \to \mathcal{F} \to \pi_{X, *}\mathcal{G} \to \mathcal{F}' \to 0$
正合。特别地，$\pi_{X, *}\mathcal{G}$ 是 $X_\etale$ 上的挠阿贝尔层。
因此 $\mathcal{G} = \pi_X^{-1}\pi_{X, *}\mathcal{G}$ 属于
$\mathcal{A}'_X$，这就验证了最后一个条件。

\medskip\noindent
Cohomology on Sites, Property (\ref{sites-cohomology-item-base-change-P})
成立，因为概形有纤维积，而且概形的固有态射经基变换后仍是固有态射；
见 Morphisms, Lemma \ref{morphisms-lemma-base-change-proper}。

\medskip\noindent
Cohomology on Sites, Property (\ref{sites-cohomology-item-restriction-A})
来自 Topologies, Lemma
\ref{topologies-lemma-morphism-big-small-etale} 中的交换图 (3)。

\medskip\noindent
Cohomology on Sites, Property (\ref{sites-cohomology-item-A-sheaf})
就是 Lemma \ref{etale-cohomology-lemma-describe-pullback-pi-ph}。

\medskip\noindent
Cohomology on Sites, Property (\ref{sites-cohomology-item-A-and-P})
由 Lemma \ref{etale-cohomology-lemma-compare-higher-direct-image-proper} (2) 以及如下事实成立：
$R^if_{small}\mathcal{F}$ 是挠层，只要 $\mathcal{F}$ 是挠阿贝尔层
（定义在 $X_\etale$ 上）；见 Lemma
\ref{etale-cohomology-lemma-torsion-direct-image}。

\medskip\noindent
Cohomology on Sites, Property (\ref{sites-cohomology-item-refine-tau-by-P})
来自 More on Morphisms, Lemma
\ref{more-morphisms-lemma-dominate-fppf-etale-locally}，并结合以下事实：
有限态射是固有的，而满射固有态射定义 ph 覆盖；见
Topologies, Lemma \ref{topologies-lemma-surjective-proper-ph}。
\end{proof}

\begin{lemma}
\label{etale-cohomology-lemma-V-C-all-n-etale-ph}
沿用上述记号。
\begin{enumerate}
\item 对 $X \in \Ob((\Sch/S)_{ph})$ 和挠阿贝尔层 $\mathcal{F}$
（定义在 $X_\etale$ 上），有
$\epsilon_{X, *}a_X^{-1}\mathcal{F} = \pi_X^{-1}\mathcal{F}$
，并且 $R^i\epsilon_{X, *}(a_X^{-1}\mathcal{F}) = 0$ 对 $i > 0$ 成立。
\item 对固有态射 $f : X \to Y$（属于 $(\Sch/S)_{ph}$）和挠阿贝尔层
$\mathcal{F}$（定义在 $X$ 上），有
$a_Y^{-1}(R^if_{small, *}\mathcal{F}) =
R^if_{big, ph, *}(a_X^{-1}\mathcal{F})$
，对所有 $i$ 成立。
\item 对概形 $X$ 和 $K$（属于 $D^+(X_\etale)$，且上同调层为挠层），映射
$\pi_X^{-1}K \to R\epsilon_{X, *}(a_X^{-1}K)$ 是同构。
\item 对概形的固有态射 $f : X \to Y$ 和 $K$（属于
$D^+(X_\etale)$，且上同调层为挠层），有
$a_Y^{-1}(Rf_{small, *}K) = Rf_{big, ph, *}(a_X^{-1}K)$.
\end{enumerate}
\end{lemma}

\begin{proof}
由 Lemma \ref{etale-cohomology-lemma-compare-ph-etale}，Cohomology on Sites, Section
\ref{sites-cohomology-section-compare-general} 中的各引理都适用于当前情形。
为翻译这些结果，注意 Cohomology on Sites, Lemma
\ref{sites-cohomology-lemma-A} 中的范畴 $\mathcal{A}_X$，是
$\textit{Ab}((\Sch/X)_{ph})$ 中由形如 $a_X^{-1}\mathcal{F}$ 的层构成的
全子范畴，其中 $\mathcal{F}$ 是 $X_\etale$ 上的挠阿贝尔层。

\medskip\noindent
(1) 等价于 $(V_n)$ 对所有 $n$ 都成立；这由 Cohomology on Sites, Lemma
\ref{sites-cohomology-lemma-V-C-all-n-general} 成立。

\medskip\noindent
(2) 来自对 Cohomology on Sites, Lemma
\ref{sites-cohomology-lemma-V-implies-C-general} 的结论应用
$\epsilon_Y^{-1}$。

\medskip\noindent
(3) 来自 Cohomology on Sites, Lemma
\ref{sites-cohomology-lemma-V-C-all-n-general} (1)，因为
$\pi_X^{-1}K$ 属于 $D^+_{\mathcal{A}'_X}((\Sch/X)_\etale)$，并且
$a_X^{-1} = \epsilon_X^{-1} \circ a_X^{-1}$。

\medskip\noindent
出于同样的理由，(4) 来自 Cohomology on Sites, Lemma
\ref{sites-cohomology-lemma-V-C-all-n-general} (2)。
\end{proof}

\begin{lemma}
\label{etale-cohomology-lemma-cohomological-descent-etale-ph}
设 $X$ 是概形。对 $K \in D^+(X_\etale)$（其上同调层为挠层），映射
$$
K \longrightarrow Ra_{X, *}a_X^{-1}K
$$
是同构，其中 $a_X : \Sh((\Sch/X)_{ph}) \to \Sh(X_\etale)$ 如上所述。
\end{lemma}

\begin{proof}
我们先把断言归约到 $K$ 由单个阿贝尔层给出的情形。也就是说，把 $K$
表示为挠阿贝尔层的下有界复形 $\mathcal{F}^\bullet$。由 Cohomology on Sites, Lemma
\ref{sites-cohomology-lemma-torsion}，这是可能的。由单个层的情形可知
$\mathcal{F}^n = a_{X, *} a_X^{-1} \mathcal{F}^n$
，而且层 $R^qa_{X, *}a_X^{-1}\mathcal{F}^n$ 对 $q > 0$ 为零。
对 $a_X^{-1}\mathcal{F}^\bullet$ 和函子 $a_{X, *}$ 应用 Leray 无环性引理
（Derived Categories, Lemma \ref{derived-lemma-leray-acyclicity}），
便得到结论。以下假设 $K = \mathcal{F}$，其中 $\mathcal{F}$ 是挠阿贝尔层。

\medskip\noindent
由 Lemma \ref{etale-cohomology-lemma-describe-pullback-pi-ph}，有
$a_{X, *}a_X^{-1}\mathcal{F} = \mathcal{F}$。因此只须证明
$R^qa_{X, *}a_X^{-1}\mathcal{F} = 0$ 对 $q > 0$ 成立。
为此可以使用 $a_X = \epsilon_X \circ \pi_X$ 和 Leray 谱序列
（Cohomology on Sites, Lemma \ref{sites-cohomology-lemma-relative-Leray}）。
由 Lemma \ref{etale-cohomology-lemma-V-C-all-n-etale-ph}，有
$R^i\epsilon_{X, *}(a_X^{-1}\mathcal{F}) = 0$（当 $i > 0$），并且
$\epsilon_{X, *}a_X^{-1}\mathcal{F} = \pi_X^{-1}\mathcal{F}$。
由 Lemma \ref{etale-cohomology-lemma-cohomological-descent-etale}，有
$R^j\pi_{X, *}(\pi_X^{-1}\mathcal{F}) = 0$（当 $j > 0$）。证明完毕。
\end{proof}

\begin{lemma}
\label{etale-cohomology-lemma-compare-cohomology-etale-ph}
设 $X$ 是概形，并设
$a_X : \Sh((\Sch/X)_{ph}) \to \Sh(X_\etale)$ 如上。则：
\begin{enumerate}
\item $H^q(X_\etale, \mathcal{F}) = H^q_{ph}(X, a_X^{-1}\mathcal{F})$
，其中 $\mathcal{F}$ 是 $X_\etale$ 上的挠阿贝尔层；
\item $H^q(X_\etale, K) = H^q_{ph}(X, a_X^{-1}K)$
，其中 $K \in D^+(X_\etale)$ 且其上同调层为挠层。
\end{enumerate}
例如，如果 $A$ 是挠阿贝尔群，则
$H^q_\etale(X, \underline{A}) = H^q_{ph}(X, \underline{A})$.
\end{lemma}

\begin{proof}
由 Lemma \ref{etale-cohomology-lemma-cohomological-descent-etale-ph} 和
Cohomology on Sites, Remark \ref{sites-cohomology-remark-before-Leray}
即得此结论。
\end{proof}












\section{h 拓扑与 \'etale 拓扑的比较}
\label{etale-cohomology-section-h-etale}

\noindent
本节的一个范本，是 Cohomology on Sites, Section
\ref{sites-cohomology-section-cohomology-LC} 中讨论的局部紧拓扑空间上通常拓扑与
qc 拓扑的比较。此外，本节与比较 ph 拓扑和 \'etale 拓扑的那一节几乎逐字相同。
我们先回顾 Topologies, Sections
\ref{topologies-section-change-topologies}、
\ref{topologies-section-etale} 以及 More on Flatness, Section
\ref{flat-section-h} 中的一些材料。

\medskip\noindent
设 $S$ 是概形，并设 $(\Sch/S)_h$ 是一个 h 网站。
在同一个底层范畴上还有第二个拓扑，即 \'etale 拓扑，因而还有第二个网站
$(\Sch/S)_\etale$。恒等函子
$(\Sch/S)_\etale \to (\Sch/S)_h$ 是连续的
（由 More on Flatness, Lemma \ref{flat-lemma-zariski-h} 和
Topologies，引理
\ref{topologies-lemma-zariski-etale-smooth-syntomic-fppf}），并定义网站态射
$$
\epsilon_S : (\Sch/S)_h \longrightarrow (\Sch/S)_\etale
$$
。见 Cohomology on Sites, Section \ref{sites-cohomology-section-compare}。
请注意，$\epsilon_{S, *}$ 在底层预层上是恒等函子，而
$\epsilon_S^{-1}$ 把一个 \'etale 层变为它的 h 层化。
令 $S_\etale$ 为小 \'etale 网站。存在网站态射
$$
\pi_S : (\Sch/S)_\etale \longrightarrow S_\etale
$$
，它由连续函子
$S_\etale \to (\Sch/S)_\etale$, $U \mapsto U$ 给出。
事实上，$S_\etale$ 有纤维积和终对象，上述函子与它们可交换，因而可以应用
Sites, Proposition \ref{sites-proposition-get-morphism}。

\begin{lemma}
\label{etale-cohomology-lemma-describe-pullback-pi-h}
沿用上述记号。设 $\mathcal{F}$ 是 $S_\etale$ 上的层。规则
$$
(\Sch/S)_h \longrightarrow \textit{Sets},\quad
(f : X \to S) \longmapsto \Gamma(X, f_{small}^{-1}\mathcal{F})
$$
给出一个层，因而当然也给出 $(\Sch/S)_\etale$ 上的层。
事实上，这个层等于 $\pi_S^{-1}\mathcal{F}$（在
$(\Sch/S)_\etale$ 上），而等于
$\epsilon_S^{-1}\pi_S^{-1}\mathcal{F}$（在 $(\Sch/S)_h$ 上）。
\end{lemma}

\begin{proof}
关于 \'etale 拓扑的断言就是 Lemma \ref{etale-cohomology-lemma-describe-pullback} 的内容。
为完成证明，只须说明 $\pi_S^{-1}\mathcal{F}$ 对 h 拓扑是层。
然而，我们已在 Lemma \ref{etale-cohomology-lemma-describe-pullback-pi-ph} 中证明，
$\pi_S^{-1}\mathcal{F}$ 即使对更强的 ph 拓扑也是层。
\end{proof}

\noindent
在 Lemma \ref{etale-cohomology-lemma-describe-pullback-pi-h} 的情形中，
$\epsilon_S$ 与 $\pi_S$ 的复合以及上述等式确定网站态射
$$
a_S : (\Sch/S)_h \longrightarrow S_\etale
$$

\begin{lemma}
\label{etale-cohomology-lemma-push-pull-h-etale}
沿用上述记号。设 $f : X \to Y$ 是 $(\Sch/S)_h$ 中的态射。
则有拓扑斯的交换图
$$
\xymatrix{
\Sh((\Sch/X)_h) \ar[rr]_{f_{big, h}} \ar[d]_{\epsilon_X} & &
\Sh((\Sch/Y)_h) \ar[d]^{\epsilon_Y} \\
\Sh((\Sch/X)_\etale) \ar[rr]^{f_{big, \etale}} & &
\Sh((\Sch/Y)_\etale)
}
$$
以及
$$
\xymatrix{
\Sh((\Sch/X)_h) \ar[rr]_{f_{big, h}} \ar[d]_{a_X} & &
\Sh((\Sch/Y)_h) \ar[d]^{a_Y} \\
\Sh(X_\etale) \ar[rr]^{f_{small}} & &
\Sh(Y_\etale)
}
$$
，其中 $a_X = \pi_X \circ \epsilon_X$ 且 $a_Y = \pi_X \circ \epsilon_X$。
\end{lemma}

\begin{proof}
这些图的交换性与 Topologies, Section
\ref{topologies-section-change-topologies} 中所述类似。
\end{proof}

\begin{lemma}
\label{etale-cohomology-lemma-proper-push-pull-h-etale}
在 Lemma \ref{etale-cohomology-lemma-push-pull-h-etale} 中，如果 $f$ 是固有的，则有
$a_Y^{-1} \circ f_{small, *} = f_{big, h, *} \circ a_X^{-1}$.
\end{lemma}

\begin{proof}
可以重复 Lemma \ref{etale-cohomology-lemma-compare-higher-direct-image-proper} (1) 的证明；
这里我们改为由该结论推出本结果。由于 $\epsilon_{Y, *}$ 在底层预层上是
恒等函子，它反映同构。Lemma \ref{etale-cohomology-lemma-describe-pullback-pi-h}
中的描述说明 $\epsilon_{Y, *} \circ a_Y^{-1} = \pi_Y^{-1}$，
对 $X$ 也类似。为证明典范映射
$a_Y^{-1}f_{small, *}\mathcal{F} \to f_{big, h, *}a_X^{-1}\mathcal{F}$
是同构，只须证明
\begin{align*}
\pi_Y^{-1}f_{small, *}\mathcal{F}
& =
\epsilon_{Y, *}a_Y^{-1}f_{small, *}\mathcal{F} \\
& \to 
\epsilon_{Y, *}f_{big, h, *}a_X^{-1}\mathcal{F} \\
& =
f_{big, \etale, *} \epsilon_{X, *}a_X^{-1}\mathcal{F} \\
& =
f_{big, \etale, *}\pi_X^{-1}\mathcal{F}
\end{align*}
是同构。这就是 Lemma \ref{etale-cohomology-lemma-compare-higher-direct-image-proper} (1)。
\end{proof}

\begin{lemma}
\label{etale-cohomology-lemma-compare-h-etale}
考虑比较态射 $\epsilon : (\Sch/S)_h \to (\Sch/S)_\etale$。
令 $\mathcal{P}$ 表示固有态射类。对 $X$（属于 $(\Sch/S)_\etale$），记
$\mathcal{A}'_X \subset \textit{Ab}((\Sch/X)_\etale)$
为由形如 $\pi_X^{-1}\mathcal{F}$ 的层构成的全子范畴，其中
$\mathcal{F}$ 是 $X_\etale$ 上的挠阿贝尔层。则 Cohomology on Sites 中的性质
(\ref{sites-cohomology-item-base-change-P}),
(\ref{sites-cohomology-item-restriction-A}),
(\ref{sites-cohomology-item-A-sheaf}),
(\ref{sites-cohomology-item-A-and-P})，以及
(\ref{sites-cohomology-item-refine-tau-by-P})
（见 Cohomology on Sites, Situation
\ref{sites-cohomology-situation-compare}）都成立。
\end{lemma}

\begin{proof}
我们先验证 Homology, Lemma
\ref{homology-lemma-characterize-weak-serre-subcategory} 的条件
(1)、(2)、(3)、(4)，从而说明
$\mathcal{A}'_X \subset \textit{Ab}((\Sch/X)_\etale)$ 是弱 Serre 子范畴。
(1)、(2)、(3) 是立即的，因为例如由 Lemma
\ref{etale-cohomology-lemma-cohomological-descent-etale}，$\pi_X^{-1}$ 正合且完全忠实。
如果
$0 \to \pi_X^{-1}\mathcal{F} \to \mathcal{G} \to \pi_X^{-1}\mathcal{F}' \to 0$
是 $\textit{Ab}((\Sch/X)_\etale)$ 中的短正合列，则由 Lemma
\ref{etale-cohomology-lemma-cohomological-descent-etale}，
$0 \to \mathcal{F} \to \pi_{X, *}\mathcal{G} \to \mathcal{F}' \to 0$
正合。特别地，$\pi_{X, *}\mathcal{G}$ 是 $X_\etale$ 上的挠阿贝尔层。
因此 $\mathcal{G} = \pi_X^{-1}\pi_{X, *}\mathcal{G}$ 属于
$\mathcal{A}'_X$，这就验证了最后一个条件。

\medskip\noindent
Cohomology on Sites, Property (\ref{sites-cohomology-item-base-change-P})
成立，因为概形有纤维积，概形的固有态射经基变换后仍是固有态射（见
Morphisms, Lemma \ref{morphisms-lemma-base-change-proper}），而有限表现态射经
基变换后仍是有限表现态射（见 Morphisms, Lemma
\ref{morphisms-lemma-base-change-finite-presentation}）。

\medskip\noindent
Cohomology on Sites, Property (\ref{sites-cohomology-item-restriction-A})
来自 Topologies, Lemma
\ref{topologies-lemma-morphism-big-small-etale} 中的交换图 (3)。

\medskip\noindent
Cohomology on Sites, Property (\ref{sites-cohomology-item-A-sheaf})
就是 Lemma \ref{etale-cohomology-lemma-describe-pullback-pi-h}。

\medskip\noindent
Cohomology on Sites, Property (\ref{sites-cohomology-item-A-and-P})
由 Lemma \ref{etale-cohomology-lemma-compare-higher-direct-image-proper} (2) 以及如下事实成立：
$R^if_{small}\mathcal{F}$ 是挠层，只要 $\mathcal{F}$ 是挠阿贝尔层
（定义在 $X_\etale$ 上）；见 Lemma \ref{etale-cohomology-lemma-torsion-direct-image}。

\medskip\noindent
Cohomology on Sites, Property (\ref{sites-cohomology-item-refine-tau-by-P})
由 More on Morphisms, Lemma
\ref{more-morphisms-lemma-dominate-fppf-etale-locally} 推出，并结合以下事实：
满射有限局部自由态射是满射、固有且有限表现的，因而由 More on Flatness,
Lemma \ref{flat-lemma-surjective-proper-finite-presentation-h}，它定义 h 覆盖。
\end{proof}

\begin{lemma}
\label{etale-cohomology-lemma-V-C-all-n-etale-h}
沿用上述记号。
\begin{enumerate}
\item 对 $X \in \Ob((\Sch/S)_{h})$ 和挠阿贝尔层 $\mathcal{F}$
（定义在 $X_\etale$ 上），有
$\epsilon_{X, *}a_X^{-1}\mathcal{F} = \pi_X^{-1}\mathcal{F}$
，并且 $R^i\epsilon_{X, *}(a_X^{-1}\mathcal{F}) = 0$ 对 $i > 0$ 成立。
\item 对固有态射 $f : X \to Y$（属于 $(\Sch/S)_h$）和挠阿贝尔层
$\mathcal{F}$（定义在 $X$ 上），有
$a_Y^{-1}(R^if_{small, *}\mathcal{F}) =
R^if_{big, h, *}(a_X^{-1}\mathcal{F})$
，对所有 $i$ 成立。
\item 对概形 $X$ 和 $K$（属于 $D^+(X_\etale)$，且上同调层为挠层），映射
$\pi_X^{-1}K \to R\epsilon_{X, *}(a_X^{-1}K)$ 是同构。
\item 对概形的固有态射 $f : X \to Y$ 和 $K$（属于
$D^+(X_\etale)$，且上同调层为挠层），有
$a_Y^{-1}(Rf_{small, *}K) = Rf_{big, h, *}(a_X^{-1}K)$.
\end{enumerate}
\end{lemma}

\begin{proof}
由 Lemma \ref{etale-cohomology-lemma-compare-h-etale}，Cohomology on Sites, Section
\ref{sites-cohomology-section-compare-general} 中的各引理都适用于当前情形。
为翻译这些结果，注意 Cohomology on Sites, Lemma
\ref{sites-cohomology-lemma-A} 中的范畴 $\mathcal{A}_X$，是
$\textit{Ab}((\Sch/X)_h)$ 中由形如 $a_X^{-1}\mathcal{F}$ 的层构成的
全子范畴，其中 $\mathcal{F}$ 是 $X_\etale$ 上的挠阿贝尔层。

\medskip\noindent
(1) 等价于 $(V_n)$ 对所有 $n$ 都成立；这由 Cohomology on Sites, Lemma
\ref{sites-cohomology-lemma-V-C-all-n-general} 成立。

\medskip\noindent
(2) 来自对 Cohomology on Sites, Lemma
\ref{sites-cohomology-lemma-V-implies-C-general} 的结论应用
$\epsilon_Y^{-1}$。

\medskip\noindent
(3) 来自 Cohomology on Sites, Lemma
\ref{sites-cohomology-lemma-V-C-all-n-general} (1)，因为
$\pi_X^{-1}K$ 属于 $D^+_{\mathcal{A}'_X}((\Sch/X)_\etale)$，并且
$a_X^{-1} = \epsilon_X^{-1} \circ a_X^{-1}$。

\medskip\noindent
出于同样的理由，(4) 来自 Cohomology on Sites, Lemma
\ref{sites-cohomology-lemma-V-C-all-n-general} (2)。
\end{proof}

\begin{lemma}
\label{etale-cohomology-lemma-cohomological-descent-etale-h}
设 $X$ 是概形。对 $K \in D^+(X_\etale)$（其上同调层为挠层），映射
$$
K \longrightarrow Ra_{X, *}a_X^{-1}K
$$
是同构，其中 $a_X : \Sh((\Sch/X)_h) \to \Sh(X_\etale)$ 如上所述。
\end{lemma}

\begin{proof}
我们先把断言归约到 $K$ 由单个阿贝尔层给出的情形。也就是说，把 $K$
表示为挠阿贝尔层的下有界复形 $\mathcal{F}^\bullet$。由 Cohomology on Sites,
Lemma \ref{sites-cohomology-lemma-torsion}，这是可能的。由单个层的情形可知
$\mathcal{F}^n = a_{X, *} a_X^{-1} \mathcal{F}^n$
，而且层 $R^qa_{X, *}a_X^{-1}\mathcal{F}^n$ 对 $q > 0$ 为零。
对 $a_X^{-1}\mathcal{F}^\bullet$ 和函子 $a_{X, *}$ 应用 Leray 无环性引理
（Derived Categories, Lemma \ref{derived-lemma-leray-acyclicity}），
便得到结论。以下假设 $K = \mathcal{F}$，其中 $\mathcal{F}$ 是挠阿贝尔层。

\medskip\noindent
由 Lemma \ref{etale-cohomology-lemma-describe-pullback-pi-h}，有
$a_{X, *}a_X^{-1}\mathcal{F} = \mathcal{F}$。因此只须证明
$R^qa_{X, *}a_X^{-1}\mathcal{F} = 0$ 对 $q > 0$ 成立。
为此可以使用 $a_X = \epsilon_X \circ \pi_X$ 和 Leray 谱序列
（Cohomology on Sites, Lemma \ref{sites-cohomology-lemma-relative-Leray}）。
由 Lemma \ref{etale-cohomology-lemma-V-C-all-n-etale-h}，有
$R^i\epsilon_{X, *}(a_X^{-1}\mathcal{F}) = 0$（当 $i > 0$），并且
$\epsilon_{X, *}a_X^{-1}\mathcal{F} = \pi_X^{-1}\mathcal{F}$。
由 Lemma \ref{etale-cohomology-lemma-cohomological-descent-etale}，有
$R^j\pi_{X, *}(\pi_X^{-1}\mathcal{F}) = 0$（当 $j > 0$）。证明完毕。
\end{proof}

\begin{lemma}
\label{etale-cohomology-lemma-compare-cohomology-etale-h}
设 $X$ 是概形，并设
$a_X : \Sh((\Sch/X)_h) \to \Sh(X_\etale)$ 如上。则：
\begin{enumerate}
\item $H^q(X_\etale, \mathcal{F}) = H^q_h(X, a_X^{-1}\mathcal{F})$
，其中 $\mathcal{F}$ 是 $X_\etale$ 上的挠阿贝尔层；
\item $H^q(X_\etale, K) = H^q_h(X, a_X^{-1}K)$
，其中 $K \in D^+(X_\etale)$ 且其上同调层为挠层。
\end{enumerate}
例如，如果 $A$ 是挠阿贝尔群，则
$H^q_\etale(X, \underline{A}) = H^q_h(X, \underline{A})$.
\end{lemma}

\begin{proof}
由 Lemma \ref{etale-cohomology-lemma-cohomological-descent-etale-h} 和
Cohomology on Sites, Remark \ref{sites-cohomology-remark-before-Leray}
即得此结论。
\end{proof}











\section{\'etale 层的下降}
\label{etale-cohomology-section-glueing-etale}

\noindent
我们证明 \'etale 层在 fppf 拓扑和 h 拓扑中能够``黏合''，
并证明若干相关结果. 我们已经证明了下列相关结论
\begin{enumerate}
\item 引理 \ref{etale-cohomology-lemma-describe-pullback} 告诉我们，概形 $S$ 的小 \'etale
网站上的一个层确定了 $S$ 的大 \'etale 网站上的一个层，后者对 fpqc 覆盖满足
层条件（因而当然也对 Zariski、\'etale、光滑、合成以及 fppf 覆盖满足层条件），
\item 引理 \ref{etale-cohomology-lemma-describe-pullback-pi-fppf} 是上一点对 fppf 拓扑的重述，
\item 引理 \ref{etale-cohomology-lemma-describe-pullback-pi-ph} 对 ph 拓扑证明了同样的结论，
\item 引理 \ref{etale-cohomology-lemma-describe-pullback-pi-h} 对 h 拓扑证明了同样的结论，
\item 引理 \ref{etale-cohomology-lemma-descent-sheaf-fppf-etale} 是 \'etale 层的 fppf 下降的一个版本，以及
\item 注 \ref{etale-cohomology-remark-cohomological-descent-finite} 告诉我们，对有限满射态射有
\'etale 层的下降（我们将在下文澄清并加强这一点）.
\end{enumerate}
在关于单纯空间的章节中，我们还会证明这方面的一些附加结果；例如见
Simplicial Spaces，第 \ref{spaces-simplicial-section-fppf-hypercovering}
节和第 \ref{spaces-simplicial-section-proper-hypercovering-spaces} 节.

\medskip\noindent
为了方便地表述结果，我们需要一些记号. 设
$\mathcal{U} = \{f_i : X_i \to X\}$
是具有固定靶的概形态射族. 相对于 $\mathcal{U}$ 的 \'etale 层的一个
{\it 下降数据}是一个族
$((\mathcal{F}_i)_{i \in I}, (\varphi_{ij})_{i, j \in I})$，其中
\begin{enumerate}
\item $\mathcal{F}_i$ 属于 $\Sh(X_{i, \etale})$，并且
\item $\varphi_{ij} :
\text{pr}_{0, small}^{-1} \mathcal{F}_i
\longrightarrow
\text{pr}_{1, small}^{-1} \mathcal{F}_j$
是 $\Sh((X_i \times_X X_j)_\etale)$ 中的一个同构.
\end{enumerate}
并且满足{\it 余循环条件}：图表
$$
\xymatrix{
\text{pr}_{0, small}^{-1}\mathcal{F}_i
\ar[dr]_{\text{pr}_{02, small}^{-1}\varphi_{ik}}
\ar[rr]^{\text{pr}_{01, small}^{-1}\varphi_{ij}} & &
\text{pr}_{1, small}^{-1}\mathcal{F}_j
\ar[dl]^{\text{pr}_{12, small}^{-1}\varphi_{jk}} \\
& \text{pr}_{2, small}^{-1}\mathcal{F}_k
}
$$
在 $\Sh((X_i \times_X X_j \times_X X_k)_\etale)$ 中交换.
自然有{\it 下降数据的态射}这一概念，由此得到下降数据范畴.
称下降数据
$((\mathcal{F}_i)_{i \in I}, (\varphi_{ij})_{i, j \in I})$
为{\it 有效的}，如果存在一个 $\mathcal{F}$ 属于 $\Sh(X_\etale)$，并且存在同构
$\varphi_i : f_{i, small}^{-1} \mathcal{F} \to \mathcal{F}_i$
（位于 $\Sh(X_{i, \etale})$ 中），它们与 $\varphi_{ij}$ 相容，也就是说
$$
\varphi_{ij} =
\text{pr}_{1, small}^{-1} (\varphi_j) \circ
\text{pr}_{0, small}^{-1} (\varphi_i^{-1})
$$
也可以如下表述. 给定一个对象 $\mathcal{F}$ 属于 $\Sh(X_\etale)$，我们得到
{\it 典范下降数据}
$(f_{i, small}^{-1}\mathcal{F}_i, c_{ij})$，其中 $c_{ij}$ 是典范同构
$$
c_{ij} : \text{pr}_{0, small}^{-1} f_{i, small}^{-1}\mathcal{F}
\longrightarrow
\text{pr}_{1, small}^{-1} f_{j, small}^{-1}\mathcal{F}
$$
下降数据
$((\mathcal{F}_i)_{i \in I}, (\varphi_{ij})_{i, j \in I})$
有效，当且仅当它同构于某个 $\mathcal{F}$ 所对应的典范下降数据，且该对象属于
$\Sh(X_\etale)$.

\medskip\noindent
如果该族由单个态射 $\{X \to Y\}$ 组成，那么我们把下降数据视为一对
$(\mathcal{F}, \varphi)$，其中 $\mathcal{F}$ 是 $\Sh(X_\etale)$ 的一个对象，
而 $\varphi$ 是同构
$$
\text{pr}_{0, small}^{-1} \mathcal{F}
\longrightarrow
\text{pr}_{1, small}^{-1} \mathcal{F}
$$
，它位于 $\Sh((X \times_Y X)_\etale)$ 中并满足余循环条件：
$$
\xymatrix{
\text{pr}_{0, small}^{-1}\mathcal{F}
\ar[dr]_{\text{pr}_{02, small}^{-1}\varphi}
\ar[rr]^{\text{pr}_{01, small}^{-1}\varphi} & &
\text{pr}_{1, small}^{-1}\mathcal{F}
\ar[dl]^{\text{pr}_{12, small}^{-1}\varphi} \\
& \text{pr}_{2, small}^{-1}\mathcal{F}
}
$$
该图在 $\Sh((X \times_Y X \times_Y X)_\etale)$ 中交换.
与前面完全一样，也有下降数据的态射和有效性的概念.

\medskip\noindent
我们先证明满的整态射的有效下降.

\begin{lemma}
\label{etale-cohomology-lemma-glue-etale-sheaf-section}
设 $f : X \to Y$ 是有截面的概形态射. 则函子
$$
\Sh(Y_\etale)
\longrightarrow
\text{\'etale 层关于 }\{X \to Y\}\text{ 的下降数据}
$$
把 $\mathcal{G}$（它属于 $\Sh(Y_\etale)$）送到典范下降数据；它是范畴等价.
\end{lemma}

\begin{proof}
这是形式的，并且只依赖于拉回函子的函子性. 我们略去细节. 提示：如果
$s : Y \to X$ 是一个截面，那么一个拟逆是把 $(\mathcal{F}, \varphi)$
送到 $s_{small}^{-1}\mathcal{F}$ 的函子.
\end{proof}

\begin{lemma}
\label{etale-cohomology-lemma-glue-etale-sheaf-integral-surjective}
设 $f : X \to Y$ 是概形的满整态射. 函子
$$
\Sh(Y_\etale)
\longrightarrow
\text{\'etale 层关于 }\{X \to Y\}\text{ 的下降数据}
$$
是范畴等价.
\end{lemma}

\begin{proof}
在本证明中，我们省略拉回和正像函子的下标 ${}_{small}$.
记 $X_1 = X \times_Y X$，并记 $f_1 : X_1 \to Y$ 为态射
$f \circ \text{pr}_0 = f \circ \text{pr}_1$.
设 $(\mathcal{F}, \varphi)$ 是 $\{X \to Y\}$ 的一个下降数据.
令 $\mathcal{F}_1 = \text{pr}_0^{-1}\mathcal{F}$.
可把 $\varphi$ 视为定义了一个同构
$\mathcal{F}_1 \to \text{pr}_1^{-1}\mathcal{F}$.
我们断言，把下降数据
$(\mathcal{F}, \varphi)$
送到层
$$
\mathcal{G} =
\text{等化子}\left(
\xymatrix{
f_*\mathcal{F}
\ar@<1ex>[r] \ar@<-1ex>[r] &
f_{1, *}\mathcal{F}_1
}
\right)
$$
的规则是引理陈述中函子的拟逆. 两个箭头中的第一个来自映射
$$
f_*\mathcal{F} \to
f_*\text{pr}_{0, *}\text{pr}_0^{-1}\mathcal{F} =
f_{1, *}\mathcal{F}_1
$$
，而第二个箭头来自映射
$$
f_*\mathcal{F} \to
f_* \text{pr}_{1, *}\text{pr}_1^{-1}\mathcal{F}
\xleftarrow{\varphi}
f_* \text{pr}_{0, *} \text{pr}_0^{-1}\mathcal{F} =
f_{1, *}\mathcal{F}_1
$$
，其中向左的箭头可逆. 为了证明这确实可行，我们必须证明典范映射
$f^{-1}\mathcal{G} \to \mathcal{F}$ 是同构；细节略去. 为此，只需在任意一族
态射 $\Spec(k) \to Y$ 拉回之后核验，其中 $k$ 是代数闭域. 的确，相应的基变换
$X_k \to X$ 是联合满的，并且可以通过考察几何点处的茎来核验
$X_\etale$ 上层之间的一个映射是否为同构；见定理 \ref{etale-cohomology-theorem-exactness-stalks}.
由引理 \ref{etale-cohomology-lemma-integral-pushforward-commutes-with-base-change}，
$\mathcal{G}$ 从下降数据 $(\mathcal{F}, \varphi)$ 的上述构造与任意基变换交换.
因此可以假设 $Y$ 是一个代数闭域的谱（注意基变换保持态射 $f$ 的这些性质；见
Morphisms，引理 \ref{morphisms-lemma-base-change-surjective}
和 \ref{morphisms-lemma-base-change-finite}）.
在这种情形下，态射 $X \to Y$ 有一个截面，故由引理
\ref{etale-cohomology-lemma-glue-etale-sheaf-section} 可知该函子是等价.
另一方面，读者可以证明，该函子是等价当且仅当上述构造是拟逆；细节略去.
这就完成了证明.
\end{proof}

\begin{lemma}
\label{etale-cohomology-lemma-glue-etale-sheaf-proper-surjective}
设 $f : X \to Y$ 是概形的满真态射. 函子
$$
\Sh(Y_\etale)
\longrightarrow
\text{\'etale 层关于 }\{X \to Y\}\text{ 的下降数据}
$$
是范畴等价.
\end{lemma}

\begin{proof}
引理 \ref{etale-cohomology-lemma-glue-etale-sheaf-integral-surjective} 中给出的证明原封不动地适用，
但应当把对引理 \ref{etale-cohomology-lemma-integral-pushforward-commutes-with-base-change}
的引用替换为对引理 \ref{etale-cohomology-lemma-proper-base-change-f-star} 的引用.
\end{proof}

\begin{lemma}
\label{etale-cohomology-lemma-glue-etale-sheaf-check-after-base-change}
设 $f : X \to Y$ 是概形态射. 设 $Z \to Y$ 是概形的满整态射或满真态射.
如果函子
$$
\Sh(Z_\etale)
\longrightarrow
\text{\'etale 层关于 }\{X \times_Y Z \to Z\}\text{ 的下降数据}
$$
和
$$
\Sh((Z \times_Y Z)_\etale)
\longrightarrow
\text{\'etale 层关于 }
\{X \times_Y (Z \times_Y Z) \to Z \times_Y Z\}\text{ 的下降数据}
$$
是范畴等价，那么
$$
\Sh(Y_\etale)
\longrightarrow
\text{\'etale 层关于 }\{X \to Y\}\text{ 的下降数据}
$$
是等价.
\end{lemma}

\begin{proof}
这是定义以及引理 \ref{etale-cohomology-lemma-glue-etale-sheaf-integral-surjective}
和 \ref{etale-cohomology-lemma-glue-etale-sheaf-proper-surjective} 的形式推论. 细节略去.
\end{proof}

\begin{lemma}
\label{etale-cohomology-lemma-glue-etale-sheaf-fppf-cover}
设 $f : X \to Y$ 是满、平且局部有限表现的概形态射. 函子
$$
\Sh(Y_\etale)
\longrightarrow
\text{\'etale 层关于 }\{X \to Y\}\text{ 的下降数据}
$$
是范畴等价.
\end{lemma}

\begin{proof}
与引理 \ref{etale-cohomology-lemma-glue-etale-sheaf-integral-surjective} 的证明完全一样，
我们断言，一个拟逆由把 $(\mathcal{F}, \varphi)$ 送到
$$
\mathcal{G} =
\text{等化子}\left(
\xymatrix{
f_*\mathcal{F}
\ar@<1ex>[r] \ar@<-1ex>[r] &
f_{1, *}\mathcal{F}_1
}
\right)
$$
的函子给出；为了证明这一点，只需证明 $f^{-1}\mathcal{G} \to \mathcal{F}$
是同构. 这一点可局部核验，因此可以并且确实假设 $Y$ 是仿射的. 随后可找到一个
有限满态射 $Z \to Y$，使得存在开覆盖 $Z = \bigcup W_i$，其中
$W_i \to Y$ 经由 $X$ 分解. 见 More on Morphisms，引理
\ref{more-morphisms-lemma-dominate-fppf-finite}.
应用引理 \ref{etale-cohomology-lemma-glue-etale-sheaf-check-after-base-change}，可知只需在把
$Y$ 分别替换为 $Z$ 和 $Z \times_Y Z$、并把 $f$ 替换为其基变换之后证明本引理.
因此可以假设 $f$ 在 Zariski 局部有截面. 当然，利用该问题在 $Y$ 上是局部的，
我们便归结到有一个截面的情形，而这正是引理
\ref{etale-cohomology-lemma-glue-etale-sheaf-section}.
\end{proof}

\begin{lemma}
\label{etale-cohomology-lemma-glue-etale-sheaf-fppf}
设 $\{f_i : X_i \to X\}$ 是概形的一个 fppf 覆盖. 函子
$$
\Sh(X_\etale)
\longrightarrow
\text{\'etale 层关于 }\{f_i : X_i \to X\}\text{ 的下降数据}
$$
是范畴等价.
\end{lemma}

\begin{proof}
对态射 $f : \coprod X_i \to X$，有引理
\ref{etale-cohomology-lemma-glue-etale-sheaf-fppf-cover}. 随后一个形式论证表明，$f$ 的下降数据
与该覆盖的下降数据是一回事；比较 Descent，引理
\ref{descent-lemma-family-is-one}. 细节略去.
\end{proof}

\begin{lemma}
\label{etale-cohomology-lemma-glue-etale-sheaf-modification}
设 $f : X' \to X$ 是概形的真态射. 设 $i : Z \to X$ 是闭浸入.
令 $E = Z \times_X X'$. 图表为
$$
\xymatrix{
E \ar[d]_g \ar[r]_j & X' \ar[d]^f \\
Z \ar[r]^i & X
}
$$
如果 $f$ 在 $X \setminus Z$ 上是同构，那么函子
$$
\Sh(X_\etale)
\longrightarrow
\Sh(X'_\etale) \times_{\Sh(E_\etale)} \Sh(Z_\etale)
$$
是范畴等价.
\end{lemma}

\begin{proof}
我们将使用 Categories，例 \ref{categories-example-2-fibre-product-categories}
中构造的 $2$-纤维积范畴. 该函子把 $\mathcal{F}$ 送到三元组
$(f^{-1}\mathcal{F}, i^{-1}\mathcal{F}, c)$，其中
$c : j^{-1}f^{-1}\mathcal{F} \to g^{-1}i^{-1}\mathcal{F}$ 是典范同构.
我们将构造一个拟逆函子. 设 $(\mathcal{F}', \mathcal{G}, \alpha)$
是箭头右端的一个对象. 我们得到同构
$$
i^{-1}f_*\mathcal{F}' = g_*j^{-1}\mathcal{F}'
\xrightarrow{g_*\alpha}
g_*g^{-1}\mathcal{G}
$$
第一个等式是引理 \ref{etale-cohomology-lemma-proper-base-change-f-star}.
利用它，我们得到映射
$i_*\mathcal{G} \to i_*g_*g^{-1}\mathcal{G}$
和 $f'_*\mathcal{F}' \to i_*g_*g^{-1}\mathcal{G}$. 令
$$
\mathcal{F} = f_*\mathcal{F}' \times_{i_*g_*g^{-1}\mathcal{G}} i_*\mathcal{G}
$$
我们断言，$\mathcal{F}$ 是箭头左端的一个对象，并且它在右端的像同构于我们起初的
三元组. 计算 $\mathcal{F}$ 在几何点 $\overline{x}$ 处的茎，其中该点属于 $X$.

\medskip\noindent
如果 $\overline{x}$ 不在 $Z$ 中，那么一方面，$\overline{x}$ 来自唯一的几何点
$\overline{x}'$，其中该点属于 $X'$，并且
$\mathcal{F}'_{\overline{x}'} = (f_*\mathcal{F}')_{\overline{x}}$
；另一方面，$(i_*\mathcal{G})_{\overline{x}}$ 和
$(i_*g_*g^{-1}\mathcal{G})_{\overline{x}}$ 都是单点集.
所以可见 $\mathcal{F}_{\overline{x}}$ 等于 $\mathcal{F}'_{\overline{x}'}$.

\medskip\noindent
如果 $\overline{x}$ 在 $Z$ 中，也就是说，$\overline{x}$ 是一个几何点
$\overline{z}$ 的像，其中该点属于 $Z$，那么得到
$(i_*\mathcal{G})_{\overline{x}} = \mathcal{G}_{\overline{z}}$
以及
$$
(i_*g_*g^{-1}\mathcal{G})_{\overline{x}} =
(g_*g^{-1}\mathcal{G})_{\overline{z}} =
\Gamma(E_{\overline{z}}, g^{-1}\mathcal{G}|_{E_{\overline{z}}})
$$
（由上面使用的正像真基变换），并且类似地
$$
(f_*\mathcal{F}')_{\overline{x}} =
\Gamma(X'_{\overline{x}}, \mathcal{F}'|_{X'_{\overline{x}}})
$$
因为有等同 $E_{\overline{z}} = X'_{\overline{x}}$，并且 $\alpha$ 定义了层
$\mathcal{F}'|_{X'_{\overline{x}}}$ 与
$g^{-1}\mathcal{G}|_{E_{\overline{z}}}$ 之间的同构，所以得到
$$
\mathcal{F}_{\overline{x}} = \mathcal{G}_{\overline{z}}
$$
.

\medskip\noindent
为了完成证明，注意有典范映射
$i^{-1}\mathcal{F} \to \mathcal{G}$ 和 $f^{-1}\mathcal{F} \to \mathcal{F}'$，
它们与 $\alpha$ 相容，并且在茎上给出我们上面看到的同构.
我们略去这些映射的细致构造.
\end{proof}

\begin{lemma}
\label{etale-cohomology-lemma-h-descent-etale-sheaves}
设 $S$ 是概形. 则群胚化范畴
$$
p : \mathcal{S} \longrightarrow (\Sch/S)_h
$$
在 $U$ 上的纤维范畴是层范畴 $\Sh(U_\etale)$，即 $U$ 的小 \'etale
网站上的层所成的范畴；它是一个群胚叠.
\end{lemma}

\begin{proof}
为证明本引理，我们将核验 More on Flatness，引理
\ref{flat-lemma-refine-check-h-stack} 的条件 (1)、(2) 和 (3).

\medskip\noindent
条件 (1) 成立，因为层可以黏合（并且 Zariski 覆盖是 \'etale 覆盖）.
见 Sites，引理 \ref{sites-lemma-glue-sheaves}.

\medskip\noindent
为看出条件 (2)，设 $f : X \to Y$ 是 $S$ 上的满、平、真且有限表现的态射，
其中 $Y$ 是仿射的. 由引理 \ref{etale-cohomology-lemma-glue-etale-sheaf-fppf-cover}
或引理 \ref{etale-cohomology-lemma-glue-etale-sheaf-proper-surjective}，
我们有 $\{X \to Y\}$ 的下降.

\medskip\noindent
条件 (3) 立即从更一般的引理 \ref{etale-cohomology-lemma-glue-etale-sheaf-modification} 推出.
\end{proof}



















\section{爆破方形与 \'etale 上同调}
\label{etale-cohomology-section-blow-up-square}

\noindent
爆破方形在 More on Flatness，第 \ref{flat-section-blow-up-ph} 节中引入.
利用真基变换定理，可以看出爆破方形有一个 Mayer--Vietoris 型结论.

\begin{lemma}
\label{etale-cohomology-lemma-blow-up-square-cohomological-descent}
设 $X$ 是概形，并设 $Z \subset X$ 是由有限型拟凝聚理想截出的闭子概形.
考虑相应的爆破方形
$$
\xymatrix{
E \ar[d]_\pi \ar[r]_j & X' \ar[d]^b \\
Z \ar[r]^i & X
}
$$
对具有挠上同调层的 $K \in D^+(X_\etale)$，有可区别三角
$$
K \to Ri_*(K|_Z) \oplus Rb_*(K|_{X'}) \to Rc_*(K|_E) \to K[1]
$$
，它位于 $D(X_\etale)$ 中，其中 $c = i \circ \pi = b \circ j$.
\end{lemma}

\begin{proof}
记号 $K|_{X'}$ 表示 $b_{small}^{-1}K$.
选取阿贝尔层的一个下有界复形 $\mathcal{F}^\bullet$，它表示 $K$.
注意，$i_*(\mathcal{F}^\bullet|_Z)$ 表示 $Ri_*(K|_Z)$，因为 $i_*$ 正合
（命题 \ref{etale-cohomology-proposition-finite-higher-direct-image-zero}）.
选取拟同构
$b_{small}^{-1}\mathcal{F}^\bullet \to \mathcal{I}^\bullet$
，其中 $\mathcal{I}^\bullet$ 是 $X'_\etale$ 上内射阿贝尔层的下有界复形.
此映射伴随于一个映射 $\mathcal{F}^\bullet \to b_*(\mathcal{I}^\bullet)$，
并且 $b_*(\mathcal{I}^\bullet)$ 表示 $Rb_*(K|_{X'})$.
由引理 \ref{etale-cohomology-lemma-proper-base-change-f-star} 有
$\pi_*(\mathcal{I}^\bullet|_E) = (b_*\mathcal{I}^\bullet)|_Z$，而由引理
\ref{etale-cohomology-lemma-proper-base-change}，此复形表示 $R\pi_*(K|_E)$. 因此映射
$$
Ri_*(K|_Z) \oplus Rb_*(K|_{X'}) \to Rc_*(K|_E)
$$
由下有界复形之间的满射映射
$$
i_*(\mathcal{F}^\bullet|_Z) \oplus
b_*(\mathcal{I}^\bullet)
\to
i_*\left(b_*(\mathcal{I}^\bullet)|_Z\right)
$$
表示. 为得到所需的可区别三角，只需证明典范映射
$\mathcal{F}^\bullet \to i_*(\mathcal{F}^\bullet|_Z) \oplus
b_*(\mathcal{I}^\bullet)$
拟同构地映到上面所示复形映射的核（也就是说，复形的一个短正合列在导出范畴中
确定一个可区别三角；见 Derived Categories，第
\ref{derived-section-canonical-delta-functor} 节）.
可以在一个几何点 $\overline{x}$ 处的茎上核验这一点，其中该点属于 $X$.
如果 $\overline{x}$ 不在 $Z$ 中，那么 $X' \to X$ 在 $\overline{x}$ 的一个开邻域上
是同构. 因而，在此情形下，如果用 $\overline{x}'$ 表示 $X'$ 的相应几何点，
那么必须证明
$$
\mathcal{F}^\bullet_{\overline{x}} \to \mathcal{I}^\bullet_{\overline{x}'}
$$
是拟同构. 由我们对 $\mathcal{I}^\bullet$ 的选取，这一点成立.
如果 $\overline{x}$ 在 $Z$ 中，那么
$b_(\mathcal{I}^\bullet)_{\overline{x}} \to 
i_*\left(b_*(\mathcal{I}^\bullet)|_Z\right)_{\overline{x}}$
是阿贝尔群复形之间的同构. 因而该核等于
$i_*(\mathcal{F}^\bullet|_Z)_{\overline{x}} =
\mathcal{F}^\bullet_{\overline{x}}$，正如所需.
\end{proof}

\begin{lemma}
\label{etale-cohomology-lemma-blow-up-square-etale-cohomology}
设 $X$ 是概形，并设 $K \in D^+(X_\etale)$ 具有挠上同调层.
设 $Z \subset X$ 是由有限型拟凝聚理想截出的闭子概形. 考虑相应的爆破方形
$$
\xymatrix{
E \ar[d] \ar[r] & X' \ar[d]^b \\
Z \ar[r] & X
}
$$
则有典范长正合列
$$
H^p_\etale(X, K) \to
H^p_\etale(X', K|_{X'}) \oplus
H^p_\etale(Z, K|_Z) \to
H^p_\etale(E, K|_E) \to
H^{p + 1}_\etale(X, K)
$$
\end{lemma}

\begin{proof}[第一个证明]
这立即从引理 \ref{etale-cohomology-lemma-blow-up-square-cohomological-descent} 以及等式
$$
R\Gamma(X, Rb_*(K|_{X'})) = R\Gamma(X', K|_{X'})
$$
推出（见 Cohomology on Sites，第 \ref{sites-cohomology-section-leray} 节），
其他项也类似.
\end{proof}

\begin{proof}[第二个证明]
由引理 \ref{etale-cohomology-lemma-compare-cohomology-etale-ph}，这些上同调群是
$X, X', E, Z$ 关于网站 $(\Sch/X)_{ph}$ 上某个阿贝尔层复形的上同调.
（具体来说，是导出范畴的对象 $a_X^{-1}K$；见上面的引理
\ref{etale-cohomology-lemma-describe-pullback-pi-ph}，并回忆 $K|_{X'} = b_{small}^{-1}K$.）
由 More on Flatness，引理 \ref{flat-lemma-blow-up-square-ph}，可表预层图表的
ph 层化是余笛卡儿的. 因而，把非常一般的 Cohomology on Sites，引理
\ref{sites-cohomology-lemma-square-triangle} 应用于网站 $(\Sch/X)_{ph}$
和本引理的交换图表，即得结论.
\end{proof}

\begin{lemma}
\label{etale-cohomology-lemma-blow-up-square-equivalence}
设 $X$ 是概形，并设 $Z \subset X$ 是由有限型拟凝聚理想截出的闭子概形.
考虑相应的爆破方形
$$
\xymatrix{
E \ar[d]_\pi \ar[r]_j & X' \ar[d]^b \\
Z \ar[r]^i & X
}
$$
假设给定
\begin{enumerate}
\item 一个对象 $K'$，它属于 $D^+(X'_\etale)$ 并具有挠上同调层，
\item 一个对象 $L$，它属于 $D^+(Z_\etale)$ 并具有挠上同调层，以及
\item 一个同构 $\gamma : K'|_E \to L|_E$.
\end{enumerate}
那么存在一个对象 $K$ 属于 $D^+(X_\etale)$，并存在同构
$f : K|_{X'} \to K'$、$g : K|_Z \to L$，使得
$\gamma = g|_E \circ f^{-1}|_E$. 此外，给定
\begin{enumerate}
\item 一个对象 $M$，它属于 $D^+(X_\etale)$ 并具有挠上同调层，
\item 一个态射 $\alpha : K' \to M|_{X'}$，它属于 $D(X'_\etale)$，
\item 一个态射 $\beta : L \to M|_Z$，它属于 $D(Z_\etale)$，
\end{enumerate}
使得
$$
\alpha|_E  = \beta|_E \circ \gamma.
$$
那么存在一个态射 $M \to K$ 属于 $D(X_\etale)$，它在 $X'$ 上的限制是
$a \circ f$，而在 $Z$ 上的限制是 $b \circ g$.
\end{lemma}

\begin{proof}
如果 $K$ 存在，那么引理 \ref{etale-cohomology-lemma-blow-up-square-cohomological-descent}
给出了它嵌入其中的一个可区别三角. 因此我们只需选取一个可区别三角
$$
K \to Ri_*(L) \oplus Rb_*(K') \to Rc_*(L|_E) \to K[1]
$$
，其中 $c = i \circ \pi = b \circ j$. 这里，映射
$Ri_*(L) \to Rc_*(L|_E)$ 是把 $Ri_*$ 应用于伴随映射
$E \to R\pi_*(L|_E)$ 所得的映射. 映射 $Rb_*(K') \to Rc_*(L|_E)$
是典范映射 $Rb_*(K') \to Rc_*(K'|_E)) = R$ 与 $Rc_*(\gamma)$ 的复合.
引理陈述中的映射 $g$ 和 $f$ 是这些映射的伴随. 如果把这个可区别三角限制到
$Z$，那么由基变换定理（引理 \ref{etale-cohomology-lemma-proper-base-change}），映射
$Rb_*(K) \to Rc_*(L|_E)$ 变成同构，因而映射 $g : K|_Z \to L$ 是同构.
观察该可区别三角，可见 $f : K|_{X'} \to K'$ 在
$X' \setminus E = X \setminus Z$ 上是同构. 此外，由构造有
$\gamma \circ f|_E = g|_E$. 由于 $\gamma$ 和 $g$ 是同构，遂推出 $f$
在 $E$ 的几何点处的茎上也诱导同构. 因此 $f$ 是同构.

\medskip\noindent
对于最后一个陈述，可以把 $K'$ 替换为 $K|_{X'}$，把 $L$ 替换为 $K|_Z$，
并把 $\gamma$ 替换为典范等同. 注意，$\alpha$ 和 $\beta$ 诱导交换方形
$$
\xymatrix{
K \ar[r] \ar@{..>}[d] &
Ri_*(K|_Z) \oplus Rb_*(K|_{X'}) \ar[r] \ar[d]_{\beta \oplus \alpha} &
Rc_*(K|_E) \ar[r] \ar[d]_{\alpha|_E} &
K[1] \ar@{..>}[d] \\
M \ar[r] &
Ri_*(M|_Z) \oplus Rb_*(M|_{X'}) \ar[r] &
Rc_*(M|_E) \ar[r] &
M[1]
}
$$
因此，由导出范畴的公理，得到一条虚线箭头，从而产生可区别三角之间的一个态射.
\end{proof}














\section{近似爆破方形与 h 拓扑}
\label{etale-cohomology-section-blow-up-h}

\noindent
本节继续 More on Flatness，第 \ref{flat-section-blow-up-h} 节中的讨论.
为方便读者，回忆一个近似爆破方形是交换图表
\begin{equation}
\label{etale-cohomology-equation-almost-blow-up-square}
\vcenter{
\xymatrix{
E \ar[d] \ar[r] & X' \ar[d]^b \\
Z \ar[r] & X
}
}
\end{equation}
；它由概形组成，并满足下列条件：
\begin{enumerate}
\item $Z \to X$ 是有限表现闭浸入，
\item $E = b^{-1}(Z)$ 是 $X'$ 的局部主闭子概形，
\item $b$ 是真且有限表现的，
\item 闭子概形 $X'' \subset X'$ 由 $\mathcal{O}_{X'}$ 的支集位于 $E$ 上的截面
所成的拟凝聚理想截出（Properties，引理
\ref{properties-lemma-sections-supported-on-closed-subset}）；它是 $X$ 沿 $Z$ 的爆破.
\end{enumerate}
由此可知，态射 $b$ 诱导同构 $X' \setminus E \to X \setminus Z$.

\medskip\noindent
我们将给出一个判据，判断大 fppf 网站 $(\Sch/S)_{fppf}$ 的导出范畴中的对象
是否具有``h 层性''，其中 $S$ 是概形. 在同一个底范畴上还有
第二个拓扑，即 h 拓扑（More on Flatness，第 \ref{flat-section-h} 节）.
回忆 fppf 覆盖是 h 覆盖（More on Flatness，引理
\ref{flat-lemma-zariski-h}）. 因而可以考虑态射
$$
\epsilon : (\Sch/S)_h \longrightarrow (\Sch/S)_{fppf}
$$
见 Cohomology on Sites，第 \ref{sites-cohomology-section-compare} 节.
特别地，有完全忠实函子
$$
R\epsilon_* : D((\Sch/S)_h) \longrightarrow D((\Sch/S)_{fppf})
$$
我们可以问：这个函子的本质像是什么？

\begin{lemma}
\label{etale-cohomology-lemma-blow-up-square-h}
沿用上述记号，如果 $K$ 属于 $R\epsilon_*$ 的本质像，那么 Cohomology on Sites，
引理 \ref{sites-cohomology-lemma-c-square} 中的映射 $c^K_{X, Z, X', E}$ 是拟同构.
\end{lemma}

\begin{proof}
用 ${}^\#$ 表示 h 拓扑中的层化. 我们已经在 More on Flatness，引理
\ref{flat-lemma-blow-up-square-h} 中看到
$h_X^\# = h_Z^\# \amalg_{h_E^\#} h_{X'}^\#$. 另一方面，映射
$h_E^\# \to h_{X'}^\#$ 是单射，因为 $E \to X'$ 是单态射.
所以本引理是 Cohomology on Sites，引理
\ref{sites-cohomology-lemma-descent-squares-helper} 的特殊情形（该引理本身是
Cohomology on Sites，引理 \ref{sites-cohomology-lemma-square-triangle} 的形式推论）.
\end{proof}

\begin{proposition}
\label{etale-cohomology-proposition-check-h}
设 $K$ 是 $D^+((\Sch/S)_{fppf})$ 的一个对象. 则 $K$ 属于
$R\epsilon_* : D((\Sch/S)_h) \to D((\Sch/S)_{fppf})$
的本质像，当且仅当 $c^K_{X, X', Z, E}$ 对每个近似爆破方形
(\ref{etale-cohomology-equation-almost-blow-up-square}) 都是拟同构；这里该方形位于
$(\Sch/S)_h$ 中，并且 $X$ 是仿射的.
\end{proposition}

\begin{proof}
应用 Cohomology on Sites，引理
\ref{sites-cohomology-lemma-descent-squares} 即可证明；其假设由引理
\ref{etale-cohomology-lemma-blow-up-square-h} 和 More on Flatness，命题
\ref{flat-proposition-check-h} 成立.
\end{proof}

\begin{lemma}
\label{etale-cohomology-lemma-refine-check-h}
设 $K$ 是 $D^+((\Sch/S)_{fppf})$ 的一个对象. 则 $K$ 属于
$R\epsilon_* : D((\Sch/S)_h) \to D((\Sch/S)_{fppf})$ 的本质像，当且仅当
对 More on Flatness，例 \ref{flat-example-one-generator} 和
\ref{flat-example-two-generators} 中的每个近似爆破方形，
$c^K_{X, X', Z, E}$ 都是拟同构.
\end{lemma}

\begin{proof}
应用 Cohomology on Sites，引理
\ref{sites-cohomology-lemma-descent-squares} 即可证明；其假设由引理
\ref{etale-cohomology-lemma-blow-up-square-h} 和 More on Flatness，引理
\ref{flat-lemma-refine-check-h} 成立.
\end{proof}






\section{h 拓扑中结构层的上同调}
\label{etale-cohomology-section-cohomology-O-h}

\noindent
设 $p$ 是素数. 设 $(\mathcal{C}, \mathcal{O})$ 是满足
$p\mathcal{O} = 0$ 的环化网站. 我们把 $\colim_F \mathcal{O}$ 定义为下列系统在
环层范畴中的余极限
$$
\mathcal{O} \xrightarrow{F} \mathcal{O} \xrightarrow{F}
\mathcal{O} \xrightarrow{F} \ldots
$$
，其中 $F : \mathcal{O} \to \mathcal{O}$，$f \mapsto f^p$ 是 Frobenius 端同态.

\begin{lemma}
\label{etale-cohomology-lemma-h-sheaf-colim-F}
设 $p$ 是素数. 设 $S$ 是 $\mathbf{F}_p$ 上的概形. 考虑层
$\mathcal{O}^{perf} = \colim_F \mathcal{O}$，它位于 $(\Sch/S)_{fppf}$ 上.
那么 $\mathcal{O}^{perf}$ 属于
$R\epsilon_* : D((\Sch/S)_h) \to D((\Sch/S)_{fppf})$ 的本质像.
\end{lemma}

\begin{proof}
使用引理 \ref{etale-cohomology-lemma-refine-check-h} 的判据来证明. 在核验这些条件之前，注意对
一个拟紧拟分离对象 $X$ 属于 $(\Sch/S)_{fppf}$，有
$$
H^i_{fppf}(X, \mathcal{O}^{perf}) = \colim_F H^i_{fppf}(X, \mathcal{O})
$$
见 Cohomology on Sites，引理
\ref{sites-cohomology-lemma-colim-works-over-collection}.
还将使用 $H^i_{fppf}(X, \mathcal{O}) = H^i(X, \mathcal{O})$；见 Descent，
命题 \ref{descent-proposition-same-cohomology-quasi-coherent}.

\medskip\noindent
令 $A, f, J$ 如 More on Flatness，例 \ref{flat-example-one-generator}
中所述，并考虑相伴的近似爆破方形. 因为 $X$, $X'$, $Z$, $E$ 都是仿射的，
$\mathcal{O}$ 没有高阶上同调. 因而只需核验
$$
0 \to
\mathcal{O}^{perf}(X) \to
\mathcal{O}^{perf}(X') \oplus \mathcal{O}^{perf}(Z) \to
\mathcal{O}^{perf}(E) \to 0
$$
是短正合列. 这已经在 More on Flatness，引理
\ref{flat-lemma-h-sheaf-colim-F}（的证明）中证明.

\medskip\noindent
令 $X, X', Z, E$ 如 More on Flatness，例
\ref{flat-example-two-generators} 中所述. 因为 $X$ 和 $Z$ 是仿射的，所以有
$H^p(X, \mathcal{O}_X) = H^p(Z, \mathcal{O}_X) = 0$，这里 $p > 0$.
由 More on Flatness，引理 \ref{flat-lemma-funny-blow-up}，有
$H^p(X', \mathcal{O}_{X'}) = 0$，这里 $p > 0$.
因为 $E = \mathbf{P}^1_Z$ 并且 $Z$ 是仿射的，也有
$H^p(E, \mathcal{O}_E) = 0$，这里 $p > 0$. 与上一段一样，归结为核验
$$
0 \to
\mathcal{O}^{perf}(X) \to
\mathcal{O}^{perf}(X') \oplus \mathcal{O}^{perf}(Z) \to
\mathcal{O}^{perf}(E) \to 0
$$
是短正合列. 这已经在 More on Flatness，引理
\ref{flat-lemma-h-sheaf-colim-F}（的证明）中证明.
\end{proof}

\begin{proposition}
\label{etale-cohomology-proposition-h-cohomology-structure-sheaf}
设 $p$ 是素数. 设 $S$ 是 $\mathbf{F}_p$ 上拟紧拟分离的概形. 则
$$
H^i((\Sch/S)_h, \mathcal{O}^h) =
\colim_F H^i(S, \mathcal{O})
$$
这里左端的 $\mathcal{O}^h$ 是指结构层的 h 层化.
\end{proposition}

\begin{proof}
这只是引理 \ref{etale-cohomology-lemma-h-sheaf-colim-F} 的重述. 回忆
$\mathcal{O}^h = \mathcal{O}^{perf} = \colim_F \mathcal{O}$, see
；见 More on Flatness，引理 \ref{flat-lemma-char-p}.
由引理 \ref{etale-cohomology-lemma-h-sheaf-colim-F} 可知，$\mathcal{O}^{perf}$ 作为
$D((\Sch/S)_{fppf})$ 的一个对象，具有形式 $R\epsilon_*K$，其中某个
$K \in D((\Sch/S)_h)$. 随后 $K = \epsilon^{-1}\mathcal{O}^{perf}$，
而它实际上等于 $\mathcal{O}^{perf}$，因为 $\mathcal{O}^{perf}$ 是 h 层.
见 Cohomology on Sites，第 \ref{sites-cohomology-section-compare} 节.
因而 $R\epsilon_*\mathcal{O}^{perf} = \mathcal{O}^{perf}$
（请原谅这个令人困惑的记号）. 于是本命题由 Leray
$$
R\Gamma_h(S, \mathcal{O}^{perf}) =
R\Gamma_{fppf}(S, R\epsilon_*\mathcal{O}^{perf}) =
R\Gamma_{fppf}(S, \mathcal{O}^{perf})
$$
以及等式
$$
H^i_{fppf}(S, \mathcal{O}^{perf}) =
H^i_{fppf}(S, \colim_F \mathcal{O}) =
\colim_F H^i_{fppf}(S, \mathcal{O})
$$
推出，因为 $S$ 拟紧且拟分离（见引理 \ref{etale-cohomology-lemma-h-sheaf-colim-F} 的证明）.
\end{proof}
% Generated mechanically from the frozen crystalline.tex target.
% Do not edit this generated file; edit the bound locale source instead.

% 好，从这里开始。
%

\setcounter{chapter}{59}
\chapter{晶体上同调}



\phantomsection
\label{crystalline-section-phantom}




\section{引言}
\label{crystalline-section-introduction}

\noindent
本章以 Johan de Jong 于 2012 年在哥伦比亚大学所作的一系列讲座为基础。
本章旨在快速介绍晶体上同调。参考文献为专著 \cite{Berthelot}。

\medskip\noindent
关于除幂环的较初等纯代数讨论已经移至一个预备章节，因为它在讨论交换代数中的
Tate 分解时也很有用。参见“除幂代数”第
\ref{dpa-section-introduction} 节。













\section{除幂包络}
\label{crystalline-section-divided-power-envelope}

\noindent
以下引理中的构造称为除幂包络。它将在后文发挥重要作用。

\begin{lemma}
\label{crystalline-lemma-divided-power-envelope}
设 $(A, I, \gamma)$ 为除幂环，$A \to B$ 为环同态。设
$J \subset B$ 为满足 $IB \subset J$ 的理想。则存在除幂环同态
$$
(A, I, \gamma) \longrightarrow (D, \bar J, \bar \gamma)
$$
使得
$$
\Hom_{(A, I, \gamma)}((D, \bar J, \bar \gamma), (C, K, \delta)) =
\Hom_{(A, I)}((B, J), (C, K))
$$
并且该等式关于 $(A, I, \gamma)$ 上的除幂代数 $(C, K, \delta)$
具有函子性。这里，左端表示 $(A, I, \gamma)$ 上除幂环的态射，
右端表示 $(A, I)$ 上的（环，理想）对的态射。
\end{lemma}

\begin{proof}
以 $\mathcal{C}$ 表示由除幂环 $(C, K, \delta)$ 组成的范畴。考虑如下定义的函子
$F : \mathcal{C} \longrightarrow \textit{Sets}$：
$$
F(C, K, \delta) =
\left\{
(\varphi, \psi)
\middle|
\begin{matrix}
\varphi : (A, I, \gamma) \to (C, K, \delta)
\text{ 为除幂环同态} \\
\psi : (B, J) \to (C, K)\text{ 为满足 }
\psi(J) \subset K\text{ 的 }A\text{-代数同态}
\end{matrix}
\right\}
$$
我们将证明“除幂代数”引理 \ref{dpa-lemma-a-version-of-brown}
适用于此函子，从而得到本引理。设
$(\varphi, \psi) \in F(C, K, \delta)$。令 $C' \subset C$ 为由
$\varphi(A)$、$\psi(B)$ 以及所有 $f \in J$ 所对应的
$\delta_n(\psi(f))$ 生成的子环。令 $K' \subset K \cap C'$ 为
$C'$ 中由 $\varphi(I)$ 以及 $f \in J$ 所对应的
$\delta_n(\psi(f))$ 生成的理想。于是
$(C', K', \delta|_{K'})$ 是除幂环，并且 $C'$ 的基数由基数
$\kappa = |A| \otimes |B|^{\aleph_0}$ 控制。此外，$\varphi$ 分解为
$A \to C' \to C$，而 $\psi$ 分解为 $B \to C' \to C$。
这就证明“除幂代数”引理 \ref{dpa-lemma-a-version-of-brown}
的假设 (1) 成立。假设 (2) 是显然的，因为除幂环范畴中的极限与遗忘函子
$(C, K, \delta) \mapsto (C, K)$ 可交换；参见“除幂代数”引理
\ref{dpa-lemma-limits} 及其证明。
\end{proof}

\begin{definition}
\label{crystalline-definition-divided-power-envelope}
设 $(A, I, \gamma)$ 为除幂环，$A \to B$ 为环同态，并设
$J \subset B$ 为满足 $IB \subset J$ 的理想。引理
\ref{crystalline-lemma-divided-power-envelope} 中构造的除幂代数
$(D, \bar J, \bar\gamma)$ 称为{\it $B$ 中 $J$ 相对于
$(A, I, \gamma)$ 的除幂包络}，记作 $D_B(J)$ 或
$D_{B, \gamma}(J)$。
\end{definition}

\noindent
设 $(A, I, \gamma) \to (C, K, \delta)$ 为除幂环同态。
$D_{B, \gamma}(J) = (D, \bar J, \bar \gamma)$ 的泛性质为
$$
\begin{matrix}
\text{把 }J\text{ 映入 }K\text{ 的环同态 }B \to C
\end{matrix}
\longleftrightarrow
\begin{matrix}
\text{除幂同态} \\
(D, \bar J, \bar \gamma) \to (C, K, \delta)
\end{matrix}
$$
该对应由预复合与 $\text{id}_D$ 对应的映射 $B \to D$ 给出。
下面列出由泛性质直接推出的 $(D, \bar J, \bar \gamma)$ 的若干性质。
存在 $A$-代数同态
\begin{equation}
\label{crystalline-equation-divided-power-envelope}
B \longrightarrow D \longrightarrow B/J
\end{equation}
第一个箭头把 $J$ 映入 $\bar J$，而 $\bar J$ 是第二个箭头的核。
当 $n > 0$ 且 $x$ 属于 $J \to D$ 的像时，元素
$\bar\gamma_n(x)$ 作为 $D$ 中的理想生成 $\bar J$，并作为
$B$-代数生成 $D$。

\begin{lemma}
\label{crystalline-lemma-divided-power-envelop-quotient}
设 $(A, I, \gamma)$ 为除幂环。设
$\varphi : B' \to B$ 为核等于 $K$ 的 $A$-代数满射，
$IB \subset J \subset B$ 为理想，并令 $J' \subset B'$ 为 $J$ 的逆像。
记 $D_{B', \gamma}(J') = (D', \bar J', \bar\gamma)$。则
$D_{B, \gamma}(J) = (D'/K', \bar J'/K', \bar\gamma)$，其中
$K'$ 是由 $n \geq 1$、$k \in K$ 时的元素 $\bar\gamma_n(k)$
生成的理想。
\end{lemma}

\begin{proof}
记 $D_{B, \gamma}(J) = (D, \bar J, \bar \gamma)$。$D'$ 的泛性质给出
除幂代数同态 $D' \to D$。由于 $B' \to B$ 和 $J' \to J$ 均为满射，
可知 $D' \to D$ 也是满射（见上文的说明）。显然，当 $n \geq 1$、
$k \in K$ 时，$\bar\gamma_n(k)$ 属于其核；因而得到同态
$D'/K' \to D$。反过来，$\bar J'/K' \subset D'/K'$ 上存在除幂结构；
参见“除幂代数”引理 \ref{dpa-lemma-kernel}。于是 $D$ 的泛性质给出逆映射
$D \to D'/K'$，从而得证。
\end{proof}

\noindent
在定义 \ref{crystalline-definition-divided-power-envelope} 的情形中，可以选取满射
$P \to B$，其中 $P$ 是 $A$ 上的多项式代数，并令 $J' \subset P$
为 $J$ 的逆像。前一引理用 $D_{P, \gamma}(J')$ 描述了
$D_{B, \gamma}(J)$。注意，由“除幂代数”引理
\ref{dpa-lemma-gamma-extends}，$\gamma$ 可延拓为 $IP$ 上的除幂结构
$\gamma'$。因此 $D_{P, \gamma}(J') = D_{P, \gamma'}(J')$ 是下一引理
所述除幂包络特殊情形的一个例子。

\begin{lemma}
\label{crystalline-lemma-describe-divided-power-envelope}
设 $(B, I, \gamma)$ 为除幂代数，$I \subset J \subset B$ 为理想。
设 $(D, \bar J, \bar \gamma)$ 为 $J$ 相对于 $\gamma$ 的除幂包络。
选取元素 $f_t \in J$（$t \in T$），使得 $J = I + (f_t)$。则存在满射
$$
\Psi :
(B\langle x_t \rangle, IB\langle x_t \rangle + B\langle x_t \rangle_+, \delta)
\longrightarrow
(D, \bar J, \bar \gamma)
$$
它是除幂环同态，并把 $x_t$ 映为 $f_t$ 在 $D$ 中的像。$\Psi$ 的核由
元素 $x_t-f_t$ 以及所有如下元素生成：
$$
\delta_n\left(\sum r_t x_t - r_0\right)
$$
其中 $r_t \in B$、$r_0 \in I$，并且在 $B$ 中
$\sum r_t f_t = r_0$。
\end{lemma}

\begin{proof}
在本引理的陈述中，把 $B\langle x_t \rangle$ 看成以如下理想为指定理想的除幂环：
$J' = IB\langle x_t \rangle + B\langle x_t \rangle_{+}$；参见“除幂代数”注
\ref{dpa-remark-divided-power-polynomial-algebra}。
$\Psi$ 的存在性来自除幂多项式环的泛性质。其像是 $D$ 的除幂子环，
故由 $D$ 的泛性质等于 $D$；这就说明 $\Psi$ 是满射。显然，
$x_t-f_t$ 属于其核。令
$$
\mathcal{R} = \{(r_0, r_t) \in I \oplus \bigoplus\nolimits_{t \in T} B
\mid \sum r_t f_t = r_0 \text{ 于 }B\text{ 中}\}
$$
若 $(r_0,r_t)\in\mathcal R$，则显然
$\sum r_t x_t-r_0$ 属于该核。由于 $\Psi$ 是除幂环同态且
$\sum r_tx_t-r_0\in J'$，故
$\delta_n(\sum r_t x_t-r_0)$ 也属于该核。令
$K\subset B\langle x_t\rangle$ 为由 $x_t-f_t$ 以及
$(r_0,r_t)\in\mathcal R$ 时的元素
$\delta_n(\sum r_t x_t-r_0)$ 生成的理想。
为了证明 $K=\Ker(\Psi)$，只需证明 $\delta$ 可延拓至
$B\langle x_t\rangle/K$。事实上，若能如此，$D$ 的泛性质便给出
$\Psi$ 的逆映射 $D\to B\langle x_t\rangle/K$。因此，需要证明
$K\cap J'$ 在 $\delta_n$ 下保持不变；参见“除幂代数”引理
\ref{dpa-lemma-kernel}。令 $K'\subset B\langle x_t\rangle$
为由下列元素生成的理想：
\begin{enumerate}
\item 当 $m>0$ 且 $(r_0,r_t)\in\mathcal R$ 时的
$\delta_m(\sum r_t x_t-r_0)$；
\item 当 $m>0$ 且 $t',t\in I$ 时的 $x_{t'}^{[m]}(x_t-f_t)$。
\end{enumerate}
我们断言 $K'=K\cap J'$。根据“除幂代数”引理
\ref{dpa-lemma-kernel} (2)(c) 的判据，以及对上列元素的
$\delta_n$ 的计算（留给读者），该断言说明当 $n>0$ 时
$K\cap J'$ 在 $\delta_n$ 下保持不变。为证明该断言，先注意到
$K'\subset K\cap J'$。反过来，若 $h\in K\cap J'$，则模 $K'$ 可写成
$$
h = \sum r_t (x_t - f_t)
$$
其中 $r_t\in B$。由于 $h\in K\cap J'\subset J'$，可知
$r_0=\sum r_t f_t\in I$，从而 $(r_0,r_t)\in\mathcal R$，并有
$$
h = \sum r_t x_t - r_0
$$
属于 $K'$，正合所需。
\end{proof}

\begin{lemma}
\label{crystalline-lemma-divided-power-envelope-add-variables}
设 $(A,I,\gamma)$ 为除幂环，$B$ 为 $A$-代数，
$IB\subset J\subset B$ 为理想，并设 $x_i$ 为一族变量。则
$$
D_{B[x_i], \gamma}(JB[x_i] + (x_i)) = D_{B, \gamma}(J) \langle x_i \rangle
$$
\end{lemma}

\begin{proof}
一种证明方法是由引理 \ref{crystalline-lemma-describe-divided-power-envelope} 推出，
因为 $B[x_i]$ 中变量 $x_i$ 之间的任何关系都是平凡的。另一方面，
本引理也可由除幂多项式代数的泛性质与除幂包络的泛性质推出。
\end{proof}

\noindent
若 $B\to B'$ 在 $V(IB')\subset\Spec(B')$ 的所有素点处均平坦，
则下一引理的条件 (1) 与 (2) 成立；而这两个条件也与该平坦性条件
密切相关，参见“代数”引理 \ref{algebra-lemma-what-does-it-mean}。
特别地，该引理说明取除幂包络与局部化可交换。

\begin{lemma}
\label{crystalline-lemma-flat-base-change-divided-power-envelope}
设 $(A,I,\gamma)$ 为除幂环，$B\to B'$ 为 $A$-代数同态。假设
\begin{enumerate}
\item $B/IB \to B'/IB'$ 是平坦的；
\item $\text{Tor}_1^B(B', B/IB) = 0$.
\end{enumerate}
则对任意理想 $IB\subset J\subset B$，典范映射
$$
D_B(J) \otimes_B B' \longrightarrow D_{B'}(JB')
$$
是同构。
\end{lemma}

\begin{proof}
令 $D=D_B(J)$，并以 $\bar J\subset D$ 表示其除幂理想，
相应的除幂结构记作 $\bar\gamma$。$D$ 的泛性质给出 $B$-代数同态
$D\to D_{B'}(JB')$，从而给出引理中的映射。只需证明
$\bar\gamma$ 可延拓至 $D\otimes_B B'$；因为一旦如此，
$D_{B'}(JB')$ 的泛性质便给出引理中映射的逆映射
$D_{B'}(JB')\to D\otimes_B B'$。

\medskip\noindent
选取满射 $P\to B'$，其中 $P$ 是 $B$ 上的多项式代数。特别地，
$B\to P$ 平坦，故由“代数”引理 \ref{algebra-lemma-flat-base-change}，
$D\to D\otimes_B P$ 也平坦。于是由“除幂代数”引理
\ref{dpa-lemma-gamma-extends}，$\bar\gamma$ 可延拓至
$D\otimes_B P$；仍以 $\bar\gamma$ 表示该延拓。令
$\mathfrak a=\Ker(P\to B')$，于是有短正合列
$$
0 \to \mathfrak a \to P \to B' \to 0
$$
因此，$\text{Tor}_1^B(B',B/IB)=0$ 蕴含
$\mathfrak a\cap IP=I\mathfrak a$。现有如下交换图：
$$
\xymatrix{
B/J \otimes_B \mathfrak a \ar[r]_\beta &
B/J \otimes_B P \ar[r] &
B/J \otimes_B B' \\
D \otimes_B \mathfrak a \ar[r]^\alpha \ar[u] &
D \otimes_B P \ar[r] \ar[u] &
D \otimes_B B' \ar[u] \\
\bar J \otimes_B \mathfrak a \ar[r] \ar[u] &
\bar J \otimes_B P \ar[r] \ar[u] &
\bar J \otimes_B B' \ar[u]
}
$$
即使在上方和右方补上 $0$，该图仍然正合。我们须证明理想
$\bar J \otimes_B P$ 上的除幂保持理想
$\Im(\alpha) \cap \bar J \otimes_B P$；参见“除幂代数”引理
\ref{dpa-lemma-kernel}。考虑正合列
$$
0 \to \mathfrak a/I\mathfrak a \to P/IP \to B'/IB' \to 0
$$
（这里使用了上面得到的 $\mathfrak a\cap IP=I\mathfrak a$）。
由于 $B'/IB'$ 在 $B/IB$ 上平坦，对该序列施加
$B/J\otimes_{B/IB}-$ 后仍为正合；参见“代数”引理
\ref{algebra-lemma-flat-tor-zero}。因此
$$
\Ker(B/J \otimes_{B/IB} \mathfrak a/I\mathfrak a \to
B/J \otimes_{B/IB} P/IP) =
\Ker(\mathfrak a/J\mathfrak a \to P/JP)
$$
为零。因此 $\beta$ 是单射，从而
$\Im(\alpha) \cap \bar J \otimes_B P$ 是
$\bar J\otimes\mathfrak a$ 的像。若 $f\in\bar J$ 且
$a\in\mathfrak a$，则
$\bar\gamma_n(f \otimes a) = \bar\gamma_n(f) \otimes a^n$
故结论显然成立。
\end{proof}

\noindent
下一引理是 \cite[命题 2.1.7]{dJ-crystalline} 的一个特殊情形，
而后者又推广了 \cite[命题 2.8.2]{Berthelot}。

\begin{lemma}
\label{crystalline-lemma-flat-extension-divided-power-envelope}
设 $(B, I, \gamma) \to (B', I', \gamma')$ 为除幂环同态，
$I \subset J \subset B$ 与 $I' \subset J' \subset B'$ 为理想。假设
\begin{enumerate}
\item $B/I \to B'/I'$ 是平坦的；
\item $J' = JB' + I'$.
\end{enumerate}
则典范映射
$$
D_{B, \gamma}(J) \otimes_B B' \longrightarrow D_{B', \gamma'}(J')
$$
是同构。
\end{lemma}

\begin{proof}
令 $D = D_{B, \gamma}(J)$。选取生成 $J/I$ 的元素 $f_t \in J$。
如引理 \ref{crystalline-lemma-describe-divided-power-envelope} 的证明中那样，令
$\mathcal{R} = \{(r_0, r_t) \in I \oplus \bigoplus\nolimits_{t \in T} B
\mid \sum r_t f_t = r_0 \text{ 于 }B\text{ 中}\}$。该引理说明
$$
D = B\langle x_t \rangle/ K
$$
其中 $K$ 由元素 $x_t - f_t$ 以及 $(r_0, r_t) \in \mathcal{R}$ 时的
$\delta_n(\sum r_t x_t - r_0)$ 生成。因此
\begin{equation}
\label{crystalline-equation-base-change}
D \otimes_B B' = B'\langle x_t \rangle/K'
\end{equation}
其中 $K'$ 由上述 $K$ 的生成元在 $B'\langle x_t \rangle$ 中的像生成。
令 $f'_t \in B'$ 为 $f_t$ 的像。由假设 (1)，元素 $f'_t \in J'$
生成 $J'/I'$，并且 $x_t - f'_t \in K'$。令
$$
\mathcal{R}' =
\{(r'_0, r'_t) \in I' \oplus \bigoplus\nolimits_{t \in T} B'
\mid \sum r'_t f'_t = r'_0 \text{ 于 }B'\text{ 中}\}
$$
为完成证明，须对 $(r'_0, r'_t) \in \mathcal{R}'$ 证明
$\delta'_n(\sum r'_t x_t - r'_0) \in K'$。因为一旦如此，
$D \otimes_B B'$ 的表示式 (\ref{crystalline-equation-base-change}) 就与引理
\ref{crystalline-lemma-describe-divided-power-envelope} 由生成元 $f'_t$ 得到的
$D_{B', \gamma'}(J')$ 的表示完全相同。设
$(r'_0, r'_t) \in \mathcal{R}'$。则
$\sum r'_t f'_t = 0$ 在 $B'/I'$ 中成立。由假设 (1)，
$B/I \to B'/I'$ 平坦，故可应用平坦性的方程判据（“代数”引理
\ref{algebra-lemma-flat-eq}），得到某个 $m > 0$，以及
$r_{jt} \in B$ 和 $c_j \in B'$（$j = 1, \ldots, m$），使得
$$
r_{j0} = \sum\nolimits_t r_{jt} f_t  \in I \text{，其中 } j = 1, \ldots, m
$$
以及
$$
i'_t  = r'_t - \sum\nolimits_j c_j r_{jt} \in I' \text{，对所有 }t
$$
注意，这还蕴含
$r'_0 = \sum_t i'_t f_t + \sum_j c_j r_{j0}$。于是
\begin{align*}
\delta'_n(\sum\nolimits_t r'_t x_t  - r'_0)
& =
\delta'_n(
\sum\nolimits_t i'_t x_t +
\sum\nolimits_{t, j} c_j r_{jt} x_t -
\sum\nolimits_t i'_t f_t -
\sum\nolimits_j c_j r_{j0}) \\
& =
\delta'_n(
\sum\nolimits_t i'_t(x_t - f_t) +
\sum\nolimits_j c_j (\sum\nolimits_t r_{jt} x_t  - r_{j0}))
\end{align*}
由于 $\delta_n(a + b) = \sum_{m = 0, \ldots, n} \delta_m(a) \delta_{n - m}(b)$，
并且当 $m > 0$ 时，$\delta_m(\sum i'_t(x_t - f_t))$ 属于由
$x_t - f_t \in K'$ 生成的理想，只需证明
$\delta_n(\sum c_j (\sum r_{jt} x_t  - r_{j0}))$ 属于 $K'$。为此使用
$$
\delta_n(\sum\nolimits_j c_j (\sum\nolimits_t r_{jt} x_t  - r_{j0}))
=
\sum c_1^{n_1} \ldots c_m^{n_m}
\delta_{n_1}(\sum r_{1t} x_t  - r_{10}) \ldots
\delta_{n_m}(\sum r_{mt} x_t  - r_{m0})
$$
其中求和遍及所有满足 $n_1 + \ldots + n_m = n$ 的指标。这正是所需结论。
\end{proof}





\section{若干显式除幂加厚}
\label{crystalline-section-explicit-thickenings}

\noindent
本节的构造将帮助我们定义晶体位点上模晶体的联络。

\begin{lemma}
\label{crystalline-lemma-divided-power-first-order-thickening}
设 $(A, I, \gamma)$ 为除幂环，$M$ 为 $A$-模。令
$B = A \oplus M$ 为 $A$-代数，其中 $M$ 是平方为零的理想；并令
$J = I \oplus M$。对 $x \in I$ 与 $z \in M$，令
$$
\delta_n(x + z) = \gamma_n(x) + \gamma_{n - 1}(x)z
$$
则 $\delta$ 是除幂结构，并且 $A \to B$ 是从
$(A, I, \gamma)$ 到 $(B, J, \delta)$ 的除幂环同态。
\end{lemma}

\begin{proof}
须检验“除幂代数”定义 \ref{dpa-definition-divided-powers}
中的条件 (1)--(5)。这里直接证明；下一引理的证明给出了一种避免计算的方法。
条件 (1) 与 (3) 显然。条件 (2) 由下式推出：
\begin{align*}
\delta_n(x + z)\delta_m(x + z)
& =
(\gamma_n(x) + \gamma_{n - 1}(x)z)(\gamma_m(x) + \gamma_{m - 1}(x)z) \\
& = \gamma_n(x)\gamma_m(x) + \gamma_n(x)\gamma_{m - 1}(x)z +
\gamma_{n - 1}(x)\gamma_m(x)z \\
& =
\frac{(n + m)!}{n!m!} \gamma_{n + m}(x) +
\left(\frac{(n + m - 1)!}{n!(m - 1)!} +
\frac{(n + m - 1)!}{(n - 1)!m!}\right)
\gamma_{n + m - 1}(x) z \\
& =
\frac{(n + m)!}{n!m!}\delta_{n + m}(x + z)
\end{align*}
条件 (5) 由下式推出：
\begin{align*}
\delta_n(\delta_m(x + z))
& =
\delta_n(\gamma_m(x) + \gamma_{m - 1}(x)z) \\
& =
\gamma_n(\gamma_m(x)) + \gamma_{n - 1}(\gamma_m(x))\gamma_{m - 1}(x)z \\
& =
\frac{(nm)!}{n! (m!)^n} \gamma_{nm}(x) +
\frac{((n - 1)m)!}{(n - 1)! (m!)^{n - 1}}
\gamma_{(n - 1)m}(x) \gamma_{m - 1}(x) z \\
& = \frac{(nm)!}{n! (m!)^n}(\gamma_{nm}(x) + \gamma_{nm - 1}(x) z)
\end{align*}
这里使用了初等数论。为证明 (4)，须验证
$$
\delta_n(x + x' + z + z')
=
\gamma_n(x + x') + \gamma_{n - 1}(x + x')(z + z')
$$
等于
$$
\sum\nolimits_{i = 0}^n
(\gamma_i(x) + \gamma_{i - 1}(x)z)
(\gamma_{n - i}(x') + \gamma_{n - i - 1}(x')z')
$$
分别比较 $1$、$z$ 与 $z'$ 的系数，再使用 $\gamma$ 的条件 (4)，
便容易得到该等式。
\end{proof}

\begin{lemma}
\label{crystalline-lemma-divided-power-second-order-thickening}
设 $(A, I, \gamma)$ 为除幂环，$M,N$ 为 $A$-模，并设
$q : M \times M \to N$ 为 $A$-双线性映射。令
$B = A \oplus M \oplus N$ 为具有如下乘法的 $A$-代数：
$$
(x, z, w)\cdot (x', z', w') = (xx', xz' + x'z, xw' + x'w + q(z, z') + q(z', z))
$$
并令 $J = I \oplus M \oplus N$。对 $(x, z, w) \in J$，令
$$
\delta_n(x, z, w) = (\gamma_n(x), \gamma_{n - 1}(x)z,
\gamma_{n - 1}(x)w + \gamma_{n - 2}(x)q(z, z))
$$
则 $\delta$ 是除幂结构，并且 $A \to B$ 是从
$(A, I, \gamma)$ 到 $(B, J, \delta)$ 的除幂环同态。
\end{lemma}

\begin{proof}
先证明“除幂代数”定义 \ref{dpa-definition-divided-powers}
中的性质 (4)。在 $J$ 中选取 $(x,z,w)$ 与 $(x',z',w')$，并选取映射
$$
A_0 = \mathbf{Z}\langle s, s'\rangle \longrightarrow A,\quad
s \longmapsto x,
s' \longmapsto x'
$$
除幂多项式环的泛性质保证可以如此选取。令
$M_0 = A_0 \oplus A_0$ 以及
$N_0 = A_0 \oplus A_0 \oplus M_0 \otimes_{A_0} M_0$，并令
$q_0 : M_0 \times M_0 \to N_0$ 为自然映射。把 $M_0 \to M$
定义为将 $M_0$ 的两个基向量分别映至 $z$ 与 $z'$ 的 $A_0$-线性映射。
把 $N_0 \to N$ 定义为如下 $A_0$-线性映射：将 $N_0$ 的前两个
基向量分别映至 $w$ 与 $w'$，并在最后一个直和项上使用
$M_0 \otimes_{A_0} M_0 \to M \otimes_A M \xrightarrow{q} N$
于是，只需对情形 $(A_0,M_0,N_0,q_0)$ 证明恒等式 (4)。
其他恒等式也同理。这就把问题归约到 $A$ 为无
$\mathbf{Z}$-挠环且各模均无 $A$-挠的情形。在此情形中，只须证明
$$
n! \delta_n(x, z, w) = (x, z, w)^n
$$
在环 $A$ 中成立；参见“除幂代数”引理 \ref{dpa-lemma-silly}。
为此，注意到
$$
(x, z, w)^2 = (x^2, 2xz, 2xw + 2q(z, z))
$$
并且由归纳法
$$
(x, z, w)^n = (x^n, nx^{n - 1}z, nx^{n - 1}w + n(n - 1)x^{n - 2}q(z, z))
$$
另一方面，
$$
n! \delta_n(x, z, w) = (n!\gamma_n(x), n!\gamma_{n - 1}(x)z,
n!\gamma_{n - 1}(x)w + n!\gamma_{n - 2}(x) q(z, z))
$$
两式相符，证明完成。
\end{proof}




\section{兼容性}
\label{crystalline-section-compatibility}

\noindent
本节并非必读内容；它说明这里的讨论如何与 \cite{Berthelot} 中的讨论衔接。
考虑以下技术性概念。

\begin{definition}
\label{crystalline-definition-compatible}
设 $(A, I, \gamma)$ 与 $(B, J, \delta)$ 为除幂环，
$A \to B$ 为环同态。若 $J + IB$ 上存在除幂结构 $\bar\gamma$，使得
$$
(A, I, \gamma) \to (B, J + IB, \bar \gamma)\quad\text{且}\quad
(B, J, \delta) \to (B, J + IB, \bar \gamma)
$$
均为除幂环同态，则称{\it $\delta$ 与 $\gamma$ 兼容}。
\end{definition}

\noindent
设 $p$ 为素数，$(A,I,\gamma)$ 为除幂环，$A\to C$ 为环同态，
并且 $p$ 在 $C$ 中幂零。假设 $\gamma$ 可延拓至 $IC$（参见“除幂代数”
引理 \ref{dpa-lemma-gamma-extends}）。在此情形中，\cite{Berthelot}
所定义的 $\Spec(C)$ 在 $\Spec(A)$ 上的（大仿射）晶体位点，
是由以下系统组成的范畴的反范畴：
$$
(B, J, \delta, A \to B, C \to B/J)
$$
其中
\begin{enumerate}
\item $(B,J,\delta)$ 是除幂环，且 $p$ 在 $B$ 中幂零；
\item $\delta$ 与 $\gamma$ 兼容；
\item 下图
$$
\xymatrix{
B \ar[r] & B/J \\
A \ar[u] \ar[r] & C \ar[u]
}
$$
交换。
\end{enumerate}
在 \cite{Berthelot} 中，条件“$\gamma$ 可延拓至 $C$ 且 $\delta$ 与
$\gamma$ 兼容”用于保证 $\Spec(C)$ 的晶体上同调与
$\Spec(C/IC)$ 的晶体上同调相同。我们将只考虑满足 $IC=0$ 的 $C$，
从而避开这一问题\footnote{当然，这需要付出代价。}。在此情形中，
对上述系统 $(B,J,\delta,A\to B,C\to B/J)$，上图的交换性蕴含
$IB\subset J$；而兼容性等价于
$(A,I,\gamma)\to(B,J,\delta)$ 是除幂环同态。




\section{仿射晶体位点}
\label{crystalline-section-affine-site}

\noindent
本节讨论晶体位点的代数版本。以下是我们讨论这些内容时采用的基本情形。

\begin{situation}
\label{crystalline-situation-affine}
这里，$p$ 是素数，$(A,I,\gamma)$ 是除幂环，$A$ 是
$\mathbf{Z}_{(p)}$-代数；并且 $A\to C$ 是满足 $IC=0$ 且
$p$ 在 $C$ 中幂零的环同态。
\end{situation}

\noindent
通常，素数 $p$ 包含在除幂理想 $I$ 中。

\begin{definition}
\label{crystalline-definition-affine-thickening}
在情形 \ref{crystalline-situation-affine} 中。
\begin{enumerate}
\item $C$ 在 $(A,I,\gamma)$ 上的一个{\it 除幂加厚}，是一个除幂代数同态
$(A,I,\gamma)\to(B,J,\delta)$（其中 $p$ 在 $B$ 中幂零），
以及一个使下图交换的环同态 $C\to B/J$：
$$
\xymatrix{
B \ar[r] & B/J \\
& C \ar[u] \\
A \ar[uu] \ar[r] & A/I \ar[u]
}
$$
\item 一个{\it 除幂加厚同态}
$$
(B, J, \delta, C \to B/J) \longrightarrow (B', J', \delta', C \to B'/J')
$$
是除幂 $A$-代数同态 $\varphi:B\to B'$，使得复合映射
$C\to B/J\to B'/J'$ 等于给定映射 $C\to B'/J'$。
\item 以 $\text{CRIS}(C/A,I,\gamma)$，或简记为 $\text{CRIS}(C/A)$，
表示 $C$ 在 $(A,I,\gamma)$ 上的除幂加厚所组成的范畴。
\item 以 $\text{Cris}(C/A,I,\gamma)$，或简记为 $\text{Cris}(C/A)$，
表示由满足 $C\to B/J$ 为同构的对象 $(B,J,\delta,C\to B/J)$
组成的全子范畴。通常将这种对象记作 $(B\to C,\delta)$，其中约定
$J=\Ker(B\to C)$。
\end{enumerate}
\end{definition}

\noindent
注意，对上述除幂加厚 $(B,J,\delta)$，理想 $J$ 是局部幂零的；
参见“除幂代数”引理 \ref{dpa-lemma-nil}。有典范函子
\begin{equation}
\label{crystalline-equation-forget-affine}
\text{CRIS}(C/A) \longrightarrow C\text{-代数},\quad
(B, J, \delta) \longmapsto B/J
\end{equation}
一般而言，该范畴没有等化子或纤维积，也没有始对象（即空余极限）。

\begin{lemma}
\label{crystalline-lemma-affine-thickenings-colimits}
在情形 \ref{crystalline-situation-affine} 中。
\begin{enumerate}
\item $\text{CRIS}(C/A)$ 有有限积（但不一定有无限积）；
\item $\text{CRIS}(C/A)$ 有所有有限非空余极限，且
(\ref{crystalline-equation-forget-affine}) 与这些余极限可交换；
\item $\text{Cris}(C/A)$ 有所有有限非空余极限，且
$\text{Cris}(C/A) \to \text{CRIS}(C/A)$ 与这些余极限可交换。
\end{enumerate}
\end{lemma}

\begin{proof}
空积，即 $C$ 在 $(A,I,\gamma)$ 上的除幂加厚范畴中的终对象，
是如下零环：将其视为 $A$-代数，赋予零理想以及零理想上唯一的除幂，
再赋予从 $C$ 到零环的唯一同态。若
$(B_t,J_t,\delta_t)_{t\in T}$ 是 $\text{CRIS}(C/A)$ 中的一族对象，
则可按“除幂代数”引理 \ref{dpa-lemma-limits} 构造其积
$(\prod_t B_t,\prod_t J_t,\prod_t\delta_t)$。映射
$C\to\prod B_t/\prod J_t=\prod B_t/J_t$ 是显然的。然而，只有当
$T$ 有限时，才能保证 $p$ 在 $\prod_tB_t$ 中幂零。

\medskip\noindent
给定 $\text{CRIS}(C/A)$ 中的两个对象 $(B,J,\gamma)$ 与
$(B',J',\gamma')$，可以在除幂环范畴中构造余笛卡尔图
$$
\xymatrix{
(B, J, \delta) \ar[r] & (B'', J'', \delta'') \\
(A, I, \gamma) \ar[r] \ar[u] & (B', J', \delta') \ar[u]
}
$$
于是有
$$
B''/J'' = B/J \otimes_{A/I} B'/J' \longleftarrow C \otimes_{A/I} C
$$
参见“除幂代数”注 \ref{dpa-remark-pushout}。以
$J''\subset K\subset B''$ 表示使下图为推出的理想：
$$
\xymatrix{
B''/J'' \ar[r] & B''/K \\
C \otimes_{A/I} C \ar[r] \ar[u] & C \ar[u]
}
$$
也就是说，$B''/K\cong B/J\otimes_C B'/J'$。令
$D_{B''}(K)=(D,\bar K,\bar\delta)$ 为 $B''$ 中 $K$ 相对于
$(B'',J'',\delta'')$ 的除幂包络。容易验证，
$(D,\bar K,\bar\delta)$ 是 $(B,J,\delta)$ 与
$(B',J',\delta')$ 在 $\text{CRIS}(C/A)$ 中的余积。

\medskip\noindent
下面讨论余等化子。设
$\alpha,\beta:(B,J,\delta)\to(B',J',\delta')$ 为
$\text{CRIS}(C/A)$ 中的态射。考虑
$B''=B'/(\alpha(b)-\beta(b))$，并令 $J''\subset B''$ 为 $J'$ 的像。
令 $D_{B''}(J'')=(D,\bar J,\bar\delta)$ 为 $B''$ 中 $J''$ 相对于
$(B',J',\delta')$ 的除幂包络。容易验证，
$(D,\bar J,\bar\delta)$ 是 $(B,J,\delta)$ 与
$(B',J',\delta')$ 在 $\text{CRIS}(C/A)$ 中的余等化子。

\medskip\noindent
由“范畴”引理 \ref{categories-lemma-almost-finite-colimits-exist}，
$\text{CRIS}(C/A)$ 有所有有限非空余极限。上述构造表明
(\ref{crystalline-equation-forget-affine}) 与它们可交换。由于
$\text{Cris}(C/A)$ 是 (\ref{crystalline-equation-forget-affine}) 在 $C$ 上的纤维范畴，
这形式地蕴含 (3)。
\end{proof}

\begin{remark}
\label{crystalline-remark-completed-affine-site}
在情形 \ref{crystalline-situation-affine} 中，以 $\text{Cris}^\wedge(C/A)$
表示如下范畴：其对象是满足下列条件的对 $(B\to C,\delta)$：
\begin{enumerate}
\item $B$ 是 $p$-进完备的 $A$-代数；
\item $B \to C$ 是 $A$-代数满射；
\item $\delta$ 是 $\Ker(B \to C)$ 上的除幂结构；
\item $A \to B$ 是除幂环同态。
\end{enumerate}
态射按定义 \ref{crystalline-definition-affine-thickening} 定义。于是
$\text{Cris}(C/A)\subset\text{Cris}^\wedge(C/A)$ 是由满足
$p$ 在 $B$ 中幂零的对象组成的全子范畴。反过来，
$\text{Cris}^\wedge(C/A)$ 的任意对象 $(B\to C,\delta)$ 等于极限
$$
(B \to C, \delta) = \lim_e (B/p^eB \to C, \delta)
$$
其中当 $e\gg0$ 时，对象 $(B/p^eB\to C,\delta)$ 属于
$\text{Cris}(C/A)$；参见“除幂代数”引理
\ref{dpa-lemma-extend-to-completion}。特别地，
$\text{Cris}^\wedge(C/A)$ 是 $\text{Cris}(C/A)$ 的 pro-对象范畴的
全子范畴；参见“范畴”注 \ref{categories-remark-pro-category}。
\end{remark}

\begin{lemma}
\label{crystalline-lemma-list-properties}
在情形 \ref{crystalline-situation-affine} 中，设 $P\to C$ 为核等于 $J$ 的
$A$-代数满射，并记 $D_{P,\gamma}(J)=(D,\bar J,\bar\gamma)$。
令 $(D^\wedge,J^\wedge,\bar\gamma^\wedge)$ 为 $D$ 的 $p$-进完备化；
参见“除幂代数”引理 \ref{dpa-lemma-extend-to-completion}。
对每个 $e\geq1$，令 $P_e=P/p^eP$，令 $J_e\subset P_e$ 为 $J$ 的像，
并记 $D_{P_e,\gamma}(J_e)=(D_e,\bar J_e,\bar\gamma)$。
则对所有充分大的 $e$，有：
\begin{enumerate}
\item $p^eD\subset\bar J$ 与 $p^eD^\wedge\subset\bar J^\wedge$
均在除幂下保持不变；
\item 作为除幂环，$D^\wedge/p^eD^\wedge=D/p^eD=D_e$；
\item $(D_e,\bar J_e,\bar\gamma)$ 是 $\text{Cris}(C/A)$ 的对象；
\item $(D^\wedge,\bar J^\wedge,\bar\gamma^\wedge)$ 等于
$\lim_e(D_e,\bar J_e,\bar\gamma)$；
\item $(D^\wedge,\bar J^\wedge,\bar\gamma^\wedge)$ 是
$\text{Cris}^\wedge(C/A)$ 的对象。
\end{enumerate}
\end{lemma}

\begin{proof}
(1) 来自“除幂代数”引理 \ref{dpa-lemma-extend-to-completion}。
$p$-进完备化的一般性质给出
$D/p^eD=D^\wedge/p^eD^\wedge$。由于 $D/p^eD$ 是除幂环，且
$P\to D/p^eD$ 经过 $P_e$ 分解，$D_e$ 的泛性质给出映射
$D_e\to D/p^eD$。反过来，$D$ 的泛性质给出映射 $D\to D_e$，
它经过 $D/p^eD$ 分解。略去验证这两个映射互逆的步骤，从而得到 (2)。
当 $e$ 充分大时，$p^eC=0$，故 (3) 成立。(4) 来自“除幂代数”引理
\ref{dpa-lemma-extend-to-completion}，而 (5) 由定义显然成立。
\end{proof}

\begin{lemma}
\label{crystalline-lemma-set-generators}
在情形 \ref{crystalline-situation-affine} 中，设 $P$ 为 $A$ 上的多项式代数，
$P\to C$ 为核等于 $J$ 的 $A$-代数满射。取引理
\ref{crystalline-lemma-list-properties} 中的 $(D_e,\bar J_e,\bar\gamma)$。
对 $\text{CRIS}(C/A)$ 的每个对象 $(B,J_B,\delta)$，存在某个 $e$
以及 $\text{CRIS}(C/A)$ 中的态射 $D_e\to B$。
\end{lemma}

\begin{proof}
可以找到提升映射 $C\to B/J_B$ 的 $A$-代数同态 $P\to B$。
由 $\text{CRIS}(C/A)$ 的定义，对某个 $e$ 有 $p^eB=0$，故
$P\to B$ 分解为 $P\to P_e\to B$。由除幂包络的泛性质，
$P_e\to B$ 经过 $D_e$ 分解。
\end{proof}

\begin{lemma}
\label{crystalline-lemma-generator-completion}
在情形 \ref{crystalline-situation-affine} 中，设 $P$ 为 $A$ 上的多项式代数，
$P\to C$ 为核等于 $J$ 的 $A$-代数满射。令
$(D,\bar J,\bar\gamma)$ 为 $D_{P,\gamma}(J)$ 的 $p$-进完备化。
对 $\text{Cris}^\wedge(C/A)$ 的每个对象 $(B\to C,\delta)$，
存在 $\text{Cris}^\wedge(C/A)$ 中的态射 $D\to B$。
\end{lemma}

\begin{proof}
可以找到与到 $C$ 的映射兼容的 $A$-代数同态 $P\to B$。
由 $\text{Cris}^\wedge(C/A)$ 的定义，$P\to B$ 分解为
$P\to D_{P,\gamma}(J)\to B$。由于 $B$ 是 $p$-进完备的，
该映射可进一步经过 $D$ 分解。
\end{proof}



\section{微分模}
\label{crystalline-section-differentials}

\noindent
本节建立除幂环的微分模理论。

\begin{definition}
\label{crystalline-definition-derivation}
设 $A$ 为环，$(B,J,\delta)$ 为除幂环，$A\to B$ 为环同态，
$M$ 为 $B$-模。取值于 $M$ 的一个{\it 除幂 $A$-导子}，是满足如下条件的
映射 $\theta:B\to M$：它可加，在 $A$ 的元素上为零，满足莱布尼茨法则
$\theta(bb')=b\theta(b')+b'\theta(b)$，并且
$$
\theta(\delta_n(x)) = \delta_{n - 1}(x)\theta(x)
$$
对所有 $n\geq1$ 与 $x\in J$ 成立。
\end{definition}

\noindent
在该定义的情形中，与通常的导子一样，存在{\it 泛除幂 $A$-导子}
$$
\text{d}_{B/A, \delta} : B \to \Omega_{B/A, \delta}
$$
使得任意除幂 $A$-导子 $\theta:B\to M$ 都可唯一地写成
$\theta=\xi\circ d_{B/A,\delta}$，其中 $\xi:\Omega_{B/A,\delta}\to M$
为 $B$-线性映射。若
$(A,I,\gamma)\to(B,J,\delta)$ 是除幂环同态，则可忘掉 $A$ 上的除幂，
并考虑 $B$ 在 $A$ 上的除幂导子。下面给出（除幂）微分泛模的若干基本性质。

\begin{lemma}
\label{crystalline-lemma-omega}
设 $A$ 为环，$(B,J,\delta)$ 为除幂环，$A\to B$ 为环同态。
\begin{enumerate}
\item 考虑带有除幂理想 $(JB[x],\delta')$ 的 $B[x]$，其中
$\delta'$ 是 $\delta$ 到 $B[x]$ 的延拓。则
$$
\Omega_{B[x]/A, \delta'} =
\Omega_{B/A, \delta} \otimes_B B[x] \oplus B[x]\text{d}x.
$$
\item 考虑带有除幂理想
$(JB\langle x \rangle+B\langle x \rangle_{+},\delta')$ 的
$B\langle x \rangle$。则
$$
\Omega_{B\langle x\rangle/A, \delta'} =
\Omega_{B/A, \delta} \otimes_B B\langle x \rangle \oplus
B\langle x\rangle \text{d}x.
$$
\item 设 $K\subset J$ 为对所有 $n>0$ 均在 $\delta_n$ 下保持不变的理想。
令 $B'=B/K$，并以 $\delta'$ 表示 $J/K$ 上诱导的除幂。则
$\Omega_{B'/A,\delta'}$ 是 $\Omega_{B/A,\delta}\otimes_BB'$ 对如下
$B'$-子模的商：该子模由 $k\in K$ 时的 $\text{d}k$ 生成。
\end{enumerate}
\end{lemma}

\begin{proof}
这些结论直接来自 $\Omega_{B/A,\delta}$ 的构造：它是由元素
$\text{d}b$ 生成的自由 $B$-模对下列关系的商：
\begin{enumerate}
\item $\text{d}(b + b') = \text{d}b + \text{d}b'$, $b, b' \in B$,
\item $\text{d}a = 0$, $a \in A$,
\item $\text{d}(bb') = b \text{d}b' + b' \text{d}b$, $b, b' \in B$,
\item $\text{d}\delta_n(f) = \delta_{n - 1}(f)\text{d}f$, $f \in J$, $n > 1$.
\end{enumerate}
注意，最后一个关系解释了为何除幂多项式代数与通常的多项式代数给出
“相同”的答案：在前一种情形中，$x$ 是除幂理想的元素，因此
$\text{d}x^{[n]}=x^{[n-1]}\text{d}x$。
\end{proof}

\noindent
设 $(A,I,\gamma)$ 为除幂环。在这一背景下，$I$ 的幂应由如下除幂给出：
$$
I^{[n]} = \text{由下列元素生成的理想：}
\gamma_{e_1}(x_1) \ldots \gamma_{e_t}(x_t)
\text{ 其中 }\sum e_j \geq n\text{ 且 }x_j \in I.
$$
当然，$I^n\subset I^{[n]}$。注意 $I^{[1]}=I$；有时也令 $I^{[0]}=A$。

\begin{lemma}
\label{crystalline-lemma-diagonal-and-differentials}
设 $(A,I,\gamma)\to(B,J,\delta)$ 为除幂环同态。令
$(B(1),J(1),\delta(1))$ 为 $(B,J,\delta)$ 与自身在
$(A,I,\gamma)$ 上的余积，即使得下图
$$
\xymatrix{
(B, J, \delta) \ar[r] & (B(1), J(1), \delta(1)) \\
(A, I, \gamma) \ar[r] \ar[u] & (B, J, \delta) \ar[u]
}
$$
为余笛卡尔图。记 $K=\Ker(B(1)\to B)$。则
$K\cap J(1)\subset J(1)$ 在除幂结构下保持不变，并且典范地有
$$
\Omega_{B/A, \delta} = K/ \left(K^2 + (K \cap J(1))^{[2]}\right)
$$
\end{lemma}

\begin{proof}
$K\cap J(1)\subset J(1)$ 在除幂结构下保持不变，是因为
$B(1)\to B$ 是除幂环同态。

\medskip\noindent
回忆 $K/K^2$ 具有典范的 $B$-模结构。以
$s_0,s_1:B\to B(1)$ 表示两个余投影，并考虑由
$b\mapsto s_1(b)-s_0(b)$ 给出的映射
$\text{d}:B\to K/(K^2+(K\cap J(1))^{[2]})$。显然，$\text{d}$ 可加，
在 $A$ 上为零，并满足莱布尼茨法则。我们断言 $\text{d}$ 是除幂
$A$-导子。取 $x\in J$，令 $y=s_1(x)$、$z=s_0(x)$。下文以
$\delta$ 代替 $\delta(1)$ 表示 $J(1)$ 上的除幂结构。须证明对
$n\geq1$，模 $K^2+(K\cap J(1))^{[2]}$ 有
$\delta_n(y)-\delta_n(z)=\delta_{n-1}(y)(y-z)$。当 $n=1$ 时等式成立。
设 $n>1$。注意，当 $i>1$ 时，$\delta_i(y-z)$ 属于
$(K\cap J(1))^{[2]}$。模 $K^2+(K\cap J(1))^{[2]}$ 计算，得到
$$
\delta_n(z) = \delta_n(z - y + y) =
\sum\nolimits_{i = 0}^n \delta_i(z - y)\delta_{n - i}(y) =
\delta_{n - 1}(y) \delta_1(z - y) + \delta_n(y)
$$
这就证明了所需等式。

\medskip\noindent
设 $M$ 为 $B$-模，$\theta:B\to M$ 为除幂 $A$-导子。令
$D=B\oplus M$，其中 $M$ 是平方为零的理想。对 $n>1$，令
$\delta_n(x+m)=\delta_n(x)+\delta_{n-1}(x)m$，从而在
$J\oplus M\subset D$ 上定义除幂结构；参见引理
\ref{crystalline-lemma-divided-power-first-order-thickening}。有两个除幂代数同态
$B\to D$：第一个是包含映射，第二个由 $b\mapsto b+\theta(b)$ 给出。
因此，得到 $(A,I,\gamma)$ 上除幂代数的典范同态 $B(1)\to D$。
它诱导映射 $K\to M$；该映射在 $K^2$ 上为零（因为 $M$ 是平方为零的理想），
并且在 $(K\cap J(1))^{[2]}$ 上也为零，因为 $M^{[2]}=0$。
由构造，复合映射 $B\to K/K^2+(K\cap J(1))^{[2]}\to M$ 等于
$\theta$。因此 $\text{d}$ 是泛除幂 $A$-导子，从而得证。
\end{proof}

\begin{remark}
\label{crystalline-remark-filtration-differentials}
设 $A\to B$ 为环同态，$(J,\delta)$ 为 $B$ 上的除幂结构。
泛模 $\Omega_{B/A,\delta}$ 还带有一点额外结构：即由
$\text{d}_{B/A,\delta}(J)$ 生成的 $B$-子模
$N\subset\Omega_{B/A,\delta}$。在引理
\ref{crystalline-lemma-diagonal-and-differentials} 给出的同构下，它对应于
$K\cap J(1)$ 在 $\Omega_{B/A,\delta}$ 中的像。考虑 $A$-代数
$D=B\oplus\Omega^1_{B/A,\delta}$，其理想为 $\bar J=J\oplus N$，
除幂 $\bar\delta$ 如该引理证明中所定义。则
$(D,\bar J,\bar\delta)$ 是除幂环，并且由 $b\mapsto b$ 与
$b\mapsto b+\text{d}_{B/A,\delta}(b)$ 给出的两个映射 $B\to D$
都是 $A$ 上的除幂环同态。此外，$N$ 是使这一结论成立的
$\Omega_{B/A,\delta}$ 的最小子模。
\end{remark}

\begin{lemma}
\label{crystalline-lemma-diagonal-and-differentials-affine-site}
在情形 \ref{crystalline-situation-affine} 中，设 $(B,J,\delta)$ 为
$\text{CRIS}(C/A)$ 的对象。令 $(B(1),J(1),\delta(1))$ 为
$(B,J,\delta)$ 与自身在 $\text{CRIS}(C/A)$ 中的余积，并记
$K=\Ker(B(1)\to B)$。则 $K\cap J(1)\subset J(1)$ 在除幂结构下
保持不变，并且典范地有
$$
\Omega_{B/A, \delta} = K/ \left(K^2 + (K \cap J(1))^{[2]}\right)
$$
\end{lemma}

\begin{proof}
逐字照搬引理 \ref{crystalline-lemma-diagonal-and-differentials} 的证明即可。
唯一需要检验的是，除幂环 $D=B\oplus M$ 是
$\text{CRIS}(C/A)$ 的对象，且两个映射 $B\to D$ 是
$\text{CRIS}(C/A)$ 中的态射。由于 $D/(J\oplus M)=B/J$，
可以利用 $C\to B/J$ 将 $D$ 视为 $\text{CRIS}(C/A)$ 的对象；
关于态射的断言由构造显然成立。
\end{proof}

\begin{lemma}
\label{crystalline-lemma-module-differentials-divided-power-envelope}
设 $(A,I,\gamma)$ 为除幂环，$A\to B$ 为环同态，
$IB\subset J\subset B$ 为理想。令
$D_{B,\gamma}(J)=(D,\bar J,\bar\gamma)$ 为除幂包络。则
$$
\Omega_{D/A, \bar\gamma} = \Omega_{B/A} \otimes_B D
$$
\end{lemma}

\begin{proof}[第一种证明]
设 $M$ 为 $D$-模。我们断言，$A$-导子 $\vartheta:B\to M$
与除幂 $A$-导子 $\theta:D\to M$ 是同一回事。由米田引理，
该断言蕴含结论。

\medskip\noindent
考虑 $D$ 的平方零加厚 $D\oplus M$。若令 $M$ 上的高阶除幂运算为零，
则 $\bar J\oplus M$ 上有除幂结构 $\delta$。换言之，对任意
$x\in\bar J$ 与 $m\in M$，令
$\delta_n(x+m)=\bar\gamma_n(x)+\bar\gamma_{n-1}(x)m$；参见引理
\ref{crystalline-lemma-divided-power-first-order-thickening}。考虑 $A$-代数映射
$B\to D\oplus M$，其第一分量由 $B\to D$ 给出，第二分量为
$\vartheta$。由泛性质，得到对应的除幂代数同态
$D\to D\oplus M$；其第二分量正是与 $\vartheta$ 对应的除幂
$A$-导子 $\theta$。
\end{proof}

\begin{proof}[第二种证明]
先设 $B$ 在 $A$ 上平坦。此时 $\gamma$ 可延拓为 $IB$ 上的除幂结构
$\gamma'$；参见“除幂代数”引理 \ref{dpa-lemma-gamma-extends}。
因此 $D=D_{B,\gamma'}(J)$ 等于除幂环 $(D',J',\delta)$ 对元素
$x_t-f_t$ 与 $\delta_n(\sum r_t x_t-r_0)$ 所生成理想的商，其中
$D'=B\langle x_t\rangle$ 且
$J'=IB\langle x_t\rangle+B\langle x_t\rangle_+$；记号与说明参见引理
\ref{crystalline-lemma-describe-divided-power-envelope}。以
$\text{d}:D'\to\Omega_{D'/A,\delta}$ 表示泛导子。注意
$$
\Omega_{D'/A, \delta} =
\Omega_{B/A} \otimes_B D' \oplus \bigoplus D' \text{d}x_t,
$$
参见引理 \ref{crystalline-lemma-omega}。由此可知，$\Omega_{D/A,\bar\gamma}$
是 $\Omega_{D'/A,\delta}\otimes_{D'}D$ 对如下子模的商：该子模由
$\text{d}$ 作用于上列 $D'\to D$ 的核的生成元所得元素生成；
参见引理 \ref{crystalline-lemma-omega}。由于
$\text{d}(x_t-f_t)=-\text{d}f_t+\text{d}x_t$，在该商中有
$\text{d}x_t=\text{d}f_t$。特别地，映射
$\Omega_{B/A}\otimes_BD\to\Omega_{D/A,\gamma}$ 是满射；其核由
$\text{d}$ 作用于元素 $\delta_n(\sum r_t x_t-r_0)$ 所得的像给出。
然而，给定 $B$ 中的关系 $\sum r_tf_t-r_0=0$，其中
$r_t\in B$ 且 $r_0\in IB$，可知
\begin{align*}
\text{d}\delta_n(\sum r_t x_t - r_0)
& =
\delta_{n - 1}(\sum r_t x_t - r_0)\text{d}(\sum r_t x_t - r_0)
\\
& =
\delta_{n - 1}(\sum r_t x_t - r_0)
\left(
\sum r_t\text{d}(x_t - f_t) + \sum (x_t - f_t)\text{d}r_t
\right)
\end{align*}
因为 $\sum r_tf_t-r_0=0$ 在 $B$ 中成立。因此，它在
$\Omega_{B/A}\otimes_A D$ 中已经为零，从而当 $B$ 在 $A$ 上平坦时得证。

\medskip\noindent
在一般情形中，把 $B$ 写成多项式环的商 $P\to B$，并令
$J'\subset P$ 为 $J$ 的逆像。采用引理
\ref{crystalline-lemma-divided-power-envelop-quotient} 的记号，有 $D=D'/K'$。
由证明第一段处理的情形，
$\Omega_{D'/A,\bar\gamma'}=\Omega_{P/A}\otimes_PD'$。于是
$\Omega_{D/A,\bar\gamma}$ 是 $\Omega_{P/A}\otimes_PD$ 对如下子模的商：
该子模由 $k\in\Ker(P\to B)$ 时的 $\text{d}\bar\gamma_n'(k)$ 生成；
参见引理 \ref{crystalline-lemma-omega} 以及引理
\ref{crystalline-lemma-divided-power-envelop-quotient} 对 $K'$ 的描述。由于
$\text{d}\bar\gamma_n'(k)=\bar\gamma'_{n-1}(k)\text{d}k$，仍只需再对
$k\in\Ker(P\to B)$ 时的 $\text{d}k$ 所生成的子模取商。又因为
$\Omega_{B/A}$ 正是 $\Omega_{P/A}\otimes_AB$ 对这些元素所生成子模的商
（“代数”引理 \ref{algebra-lemma-differential-seq}），故得证。
\end{proof}

\begin{remark}
\label{crystalline-remark-divided-powers-de-rham-complex}
设 $A\to B$ 为环同态，$(J,\delta)$ 为 $B$ 上的除幂结构。令
$\Omega_{B/A,\delta}^i=\wedge_B^i\Omega_{B/A,\delta}$，其中
$\Omega_{B/A,\delta}$ 是泛除幂 $A$-导子
$\text{d}=\text{d}_{B/A}:B\to\Omega_{B/A,\delta}$ 的值域。注意，
$\Omega_{B/A,\delta}$ 是 $\Omega_{B/A}$ 对如下 $B$-子模的商：
该子模由 $x\in J$ 时的元素
$\text{d}\delta_n(x)-\delta_{n-1}(x)\text{d}x$ 生成。我们断言“代数”
引理 \ref{algebra-lemma-de-rham-complex} 适用。为此，只需验证
$\Omega_B$ 的元素
$\text{d}\delta_n(x)-\delta_{n-1}(x)\text{d}x$ 在
$\Omega^2_{B/A,\delta}$ 中的像为零。事实上，
$$
\text{d}(\delta_{n - 1}(x)) \wedge \text{d}x
= \delta_{n - 2}(x) \text{d}x \wedge \text{d}x = 0
$$
在 $\Omega^2_{B/A,\delta}$ 中成立，正合所需。因此得到{\it 除幂德拉姆复形}
$$
\Omega^0_{B/A, \delta} \to \Omega^1_{B/A, \delta} \to
\Omega^2_{B/A, \delta} \to \ldots
$$
它将在后文发挥重要作用。
\end{remark}

\begin{remark}
\label{crystalline-remark-connection}
设 $A\to B$ 为环同态，$\Omega_{B/A}\to\Omega$ 为满足“代数”引理
\ref{algebra-lemma-de-rham-complex} 之假设的商。设 $M$ 为 $B$-模。
一个{\it 联络}是可加映射
$$
\nabla : M \longrightarrow M \otimes_B \Omega
$$
使得对 $b\in B$ 与 $m\in M$ 有
$\nabla(bm)=b\nabla(m)+m\otimes\text{d}b$。在此情形中，可按规则
$$
\nabla : M \otimes_B \Omega^i \longrightarrow M \otimes_B \Omega^{i + 1}
$$
定义这些映射：
$\nabla(m\otimes\omega)=\nabla(m)\wedge\omega+m\otimes\text{d}\omega$。
这是良定义的，因为若 $b\in B$，则
\begin{align*}
\nabla(bm \otimes \omega) - \nabla(m \otimes b\omega)
& =
\nabla(bm) \wedge \omega + bm \otimes \text{d}\omega
- \nabla(m) \wedge b\omega - m \otimes \text{d}(b\omega) \\
& =
b\nabla(m) \wedge \omega + m \otimes \text{d}b \wedge \omega
+ bm \otimes \text{d}\omega \\
& \ \ \ \ \ \ - b\nabla(m) \wedge \omega - bm \otimes \text{d}(\omega)
- m \otimes \text{d}b \wedge \omega = 0
\end{align*}
依照惯例，当且仅当复合映射
$$
M \xrightarrow{\nabla} M \otimes_B \Omega^1
\xrightarrow{\nabla} M \otimes_B \Omega^2
$$
为零时，称该联络{\it 可积}。在这种情形下，得到复形
$$
M \xrightarrow{\nabla} M \otimes_B \Omega^1
\xrightarrow{\nabla} M \otimes_B \Omega^2
\xrightarrow{\nabla} M \otimes_B \Omega^3
\xrightarrow{\nabla} M \otimes_B \Omega^4 \to \ldots
$$
称为该联络的德拉姆复形。
\end{remark}

\begin{remark}
\label{crystalline-remark-base-change-connection}
考虑如下环的交换图：
$$
\xymatrix{
B \ar[r]_\varphi & B' \\
A \ar[u] \ar[r] & A' \ar[u]
}
$$
设 $\Omega_{B/A}\to\Omega$ 与 $\Omega_{B'/A'}\to\Omega'$ 为满足
“代数”引理 \ref{algebra-lemma-de-rham-complex} 之假设的商。
假设存在映射 $\varphi:\Omega\to\Omega'$，使下图交换：
$$
\xymatrix{
\Omega_{B/A} \ar[r] \ar[d] &
\Omega_{B'/A'} \ar[d] \\
\Omega \ar[r]^{\varphi} &
\Omega'
}
$$
其中上方水平箭头是由 $\varphi:B\to B'$ 诱导的典范映射
$\Omega_{B/A}\to\Omega_{B'/A'}$。在此情形中，给定任意一对
$(M,\nabla)$，其中 $M$ 为 $B$-模且
$\nabla:M\to M\otimes_B\Omega$ 为联络，可得到其{\it 基变换}
$(M\otimes_BB',\nabla')$，其中
$$
\nabla' :
M \otimes_B B'
\longrightarrow
(M \otimes_B B') \otimes_{B'} \Omega' = M \otimes_B \Omega'
$$
按下式定义：
$$
\nabla'(m \otimes b') =
\sum m_i \otimes b'\text{d}\varphi(b_i) + m \otimes \text{d}b' 
$$
这里 $\nabla(m)=\sum m_i\otimes\text{d}b_i$。若 $\nabla$ 可积，
则 $\nabla'$ 也可积；此时有德拉姆复形之间的典范映射（注
\ref{crystalline-remark-connection}）
\begin{equation}
\label{crystalline-equation-base-change-map-complexes}
M \otimes_B \Omega^\bullet
\longrightarrow
(M \otimes_B B') \otimes_{B'} (\Omega')^\bullet =
M \otimes_B (\Omega')^\bullet
\end{equation}
它把 $m\otimes\eta$ 映为 $m\otimes\varphi(\eta)$。
\end{remark}

\begin{lemma}
\label{crystalline-lemma-differentials-completion}
设 $A\to B$ 为环同态，$(J,\delta)$ 为 $B$ 上的除幂结构，$p$ 为素数。
假设 $A$ 是 $\mathbf{Z}_{(p)}$-代数，且 $p$ 在 $B/J$ 中幂零。则
$$
\lim_e \Omega_{B_e/A, \bar\delta} =
\lim_e \Omega_{B/A, \delta}/p^e\Omega_{B/A, \delta} =
\lim_e \Omega_{B^\wedge/A, \delta^\wedge}/p^e \Omega_{B^\wedge/A, \delta^\wedge}
$$
记号与说明参见证明。
\end{lemma}

\begin{proof}
由“除幂代数”引理 \ref{dpa-lemma-extend-to-completion}，对所有充分大的
$e$，$\delta$ 可延拓至 $B_e=B/p^eB$，故第一个极限有意义。
该引理还在完备化 $B^\wedge=\lim_eB_e$ 上给出除幂结构
$\delta^\wedge$，故最后一个极限也有意义。由引理 \ref{crystalline-lemma-omega}
以及恒有 $\text{d}p^e=0$，满射
$\Omega_{B/A,\delta}\to\Omega_{B_e/A,\bar\delta}$ 的核为
$p^e\Omega_{B/A,\delta}$。映射
$\Omega_{B^\wedge/A,\delta^\wedge}\to\Omega_{B_e/A,\bar\delta}$
的核也同理。因此结论显然。
\end{proof}



\section{除幂概形}
\label{crystalline-section-divided-power-schemes}

\noindent
下面说明如何将前述概念整体化。

\begin{definition}
\label{crystalline-definition-divided-power-structure}
设 $\mathcal C$ 为位点，$\mathcal O$ 为 $\mathcal C$ 上的环层，
$\mathcal I\subset\mathcal O$ 为理想层。$\mathcal I$ 上的一个
{\it 除幂结构 $\gamma$}，是映射序列
$\gamma_n:\mathcal I\to\mathcal I$（$n\geq1$），使得对
$\mathcal C$ 的任意对象 $U$，三元组
$$
(\mathcal{O}(U), \mathcal{I}(U), \gamma)
$$
为除幂环。
\end{definition}

\noindent
当然，这尤其适用于拓扑空间上的环层。不过，稍作一般化是有益的，
因为晶体位点的结构层正是定义在一个……位点上的！本章有时把上述定义中的
三元组 $(\mathcal C,\mathcal I,\gamma)$ 称为{\it 除幂拓扑斯}。
给定另一个 $(\mathcal C',\mathcal I',\gamma')$，以及环化拓扑斯的态射
$(f, f^\sharp) : (\Sh(\mathcal{C}), \mathcal{O}) \to
(\Sh(\mathcal{C}'), \mathcal{O}')$
若 $f^\sharp(f^{-1}\mathcal I')\subset\mathcal I$，且对所有 $n\geq1$
下图均交换，便称 $(f,f^\sharp)$ 诱导一个{\it 除幂拓扑斯态射}：
$$
\xymatrix{
f^{-1}\mathcal{I}' \ar[d]_{f^{-1}\gamma'_n} \ar[r]_{f^\sharp} &
\mathcal{I} \ar[d]^{\gamma_n} \\
f^{-1}\mathcal{I}' \ar[r]^{f^\sharp} & \mathcal{I}
}
$$
若 $f$ 来自由函子 $u:\mathcal C'\to\mathcal C$ 诱导的位点态射，
这恰好意味着对所有 $U'\in\Ob(\mathcal C')$，
$$
(\mathcal{O}'(U'), \mathcal{I}'(U'), \gamma')
\longrightarrow
(\mathcal{O}(u(U')), \mathcal{I}(u(U')), \gamma)
$$
均为除幂环同态。

\medskip\noindent
在概形的情形中，还要求除幂理想为{\bf 拟凝聚}的。除此以外，
定义与拓扑斯情形完全相同。具体如下。

\begin{definition}
\label{crystalline-definition-divided-power-scheme}
一个{\it 除幂概形}是三元组 $(S,\mathcal I,\gamma)$，其中 $S$ 是概形，
$\mathcal I$ 是拟凝聚理想层，而 $\gamma$ 是 $\mathcal I$ 上的除幂结构。
一个{\it 除幂概形态射}
$(S, \mathcal{I}, \gamma) \to (S', \mathcal{I}', \gamma')$，
是满足 $f^{-1}\mathcal I'\mathcal O_S\subset\mathcal I$ 的概形态射
$f:S\to S'$，并且
$$
(\mathcal{O}_{S'}(U'), \mathcal{I}'(U'), \gamma')
\longrightarrow
(\mathcal{O}_S(f^{-1}U'), \mathcal{I}(f^{-1}U'), \gamma)
$$
对每个开集 $U'\subset S'$ 都是除幂环同态。
\end{definition}

\noindent
回忆拟凝聚理想层与闭浸入之间存在一一对应，见
“态射”第 \ref{morphisms-section-closed-immersions} 节。
因此，给定除幂概形 $(T,\mathcal J,\gamma)$，由 $\mathcal J$ 定义出
一个典范闭浸入 $U\to T$。反过来，给定闭浸入 $U\to T$，以及与
$U\to T$ 相伴的理想层 $\mathcal J$ 上的除幂结构 $\gamma$，便得到
除幂概形 $(T,\mathcal J,\gamma)$。在许多情形中，我们只想考虑
$U\to T$ 为加厚时的这类三元组 $(U,T,\gamma)$；见
“态射进阶”定义 \ref{more-morphisms-definition-thickening}。

\begin{definition}
\label{crystalline-definition-divided-power-thickening}
若 $U\to T$ 是加厚，则称上述三元组 $(U,T,\gamma)$ 为一个
{\it 除幂加厚}。
\end{definition}

\noindent
当三个除幂概形中有一个是除幂加厚时，其纤维积存在。正式陈述如下。

\begin{lemma}
\label{crystalline-lemma-fibre-product}
设 $f:(T,\mathcal J,\delta)\to(S,\mathcal I,\gamma)$ 与
$f':(T',\mathcal J',\delta')\to(S,\mathcal I,\gamma)$ 是除幂概形态射。
则存在一个除幂概形 $(T'',\mathcal J'',\delta'')$ 以及笛卡尔图
$$
\xymatrix{
T \ar[d]_f & T'' \ar[d] \ar[l] \\
S & T' \ar[l]_{f'}
}
$$
于除幂概形范畴中成立。态射 $T''\to T\times_S T'$ 是闭浸入，
而态射 $T''_0\to T_0\times_{S_0}T'_0$ 是同构。
\end{lemma}

\begin{proof}
证明梗概。注意，借助 $\Spec(-)$，最后两个断言与
“除幂代数”注 \ref{dpa-remark-pushout}
中关于除幂代数推出的结论相容。因此我们如下把 $T''$ 构造成
$T\times_S T'$ 的闭子概形：对任意仿射开集 $U\subset S$、$V\subset T$、
$V'\subset T'$，若 $f(V),f'(V')\subset U$，便考虑由
“除幂代数”注 \ref{dpa-remark-pushout}
中的构造所确定的 $V\times_U V'$ 的闭子概形。由于各
$V\times_U V'$ 构成 $T\times_S T'$ 的一个开覆盖，可以依次进行：
(1) 证明这些闭子概形能够粘合；(2) 证明所得除幂结构能够粘合；
(3) 证明粘合所得对象是除幂概形范畴中的纤维积。
断言 (1) 等价于证明除幂环范畴中的推出构造与局部化交换
（适当表述后）；这由“除幂代数”引理
\ref{dpa-lemma-gamma-extends} 得出，因为该结果说明除幂可由平坦性
延拓到局部化上。
\end{proof}

\noindent
我们作如下观察。设 $(U,T,\gamma)$ 是除幂概形，$T$ 是
$\mathbf Z_{(p)}$ 上的概形，并且 $p$ 在 $U$ 上局部幂零。则
\begin{enumerate}
\item $p$ 在 $T$ 上局部幂零 $\Leftrightarrow U\to T$ 是加厚
（见“除幂代数”引理 \ref{dpa-lemma-nil}）；以及
\item 当 $e\gg0$ 时，$p^e\mathcal O_T$ 在 $T$ 上局部地被 $\gamma$ 保持
（见“除幂代数”引理
\ref{dpa-lemma-extend-to-completion}）。
\end{enumerate}
这表明，在以下假设下可以得到关于除幂加厚的良好结果。

\begin{situation}
\label{crystalline-situation-global}
这里 $p$ 是素数，$(S,\mathcal I,\gamma)$ 是 $\mathbf Z_{(p)}$ 上的
除幂概形。置 $S_0=V(\mathcal I)\subset S$。最后，$X\to S_0$
是概形态射，并且 $p$ 在 $X$ 上局部幂零。
\end{situation}

\noindent
我们将在此情形下定义大、小晶体位点。









\section{大晶体位点}
\label{crystalline-section-big-site}

\noindent
我们先定义大位点。给定除幂概形 $(S,\mathcal I,\gamma)$，若 $T$ 配备
一个除幂概形态射 $T\to S$，便称 $(T,\mathcal J,\delta)$ 是
$(S,\mathcal I,\gamma)$ 上的除幂概形。类似地，若 $T$ 配备一个
除幂概形态射 $T\to S$，便称除幂加厚 $(U,T,\delta)$ 是
$(S,\mathcal I,\gamma)$ 上的除幂加厚。

\begin{definition}
\label{crystalline-definition-divided-power-thickening-X}
在 情形 \ref{crystalline-situation-global} 中：
\begin{enumerate}
\item 一个{$X$ 相对于 $(S,\mathcal I,\gamma)$ 的\it 除幂加厚}，
由 $(S,\mathcal I,\gamma)$ 上的除幂加厚 $(U,T,\delta)$ 与一个
$S$-态射 $U\to X$ 给出。
\item {$X$ 相对于 $(S,\mathcal I,\gamma)$ 的\it 除幂加厚态射}
按显然方式定义。
\end{enumerate}
$X$ 相对于 $(S,\mathcal I,\gamma)$ 的除幂加厚所成的范畴记为
$\text{CRIS}(X/S,\mathcal I,\gamma)$，或简记为 $\text{CRIS}(X/S)$。
\end{definition}

\noindent
对 $\text{CRIS}(X/S)$ 中任意 $(U,T,\delta)$，$p$ 在 $T$ 上局部幂零；
见情形 \ref{crystalline-situation-global} 之前的讨论。与 $(U,T,\delta)$
相伴的全部数据可由下列交换图清楚表示：
$$
\xymatrix{
T \ar[dd] & U \ar[l] \ar[d] \\
& X \ar[d] \\
S & S_0 \ar[l]
}
$$
其中 $S_0=V(\mathcal I)\subset S$。$\text{CRIS}(X/S)$ 的态射也可类似地
表示成较大的交换图。特别地，有典范遗忘函子
\begin{equation}
\label{crystalline-equation-forget}
\text{CRIS}(X/S) \longrightarrow \Sch/X,\quad
(U, T, \delta) \longmapsto U
\end{equation}
以及它的单侧逆（同时也是左伴随）
\begin{equation}
\label{crystalline-equation-endow-trivial}
\Sch/X \longrightarrow \text{CRIS}(X/S),\quad
U \longmapsto (U, U, \emptyset)
\end{equation}
后者有时很有用。

\begin{lemma}
\label{crystalline-lemma-divided-power-thickening-fibre-products}
在情形 \ref{crystalline-situation-global} 中，范畴 $\text{CRIS}(X/S)$
具有一切非空有限极限，特别是二元积与纤维积。
函子 (\ref{crystalline-equation-forget}) 与这些极限交换。
\end{lemma}

\begin{proof}
略。提示：仿射情形见引理 \ref{crystalline-lemma-affine-thickenings-colimits}；
另见“除幂代数”注 \ref{dpa-remark-pushout}。
\end{proof}

\begin{lemma}
\label{crystalline-lemma-divided-power-thickening-base-change-flat}
在情形 \ref{crystalline-situation-global} 中，设
$$
\xymatrix{
(U_3, T_3, \delta_3) \ar[d] \ar[r] & (U_2, T_2, \delta_2) \ar[d] \\
(U_1, T_1, \delta_1) \ar[r] & (U, T, \delta)
}
$$
是 $X$ 相对于 $(S,\mathcal I,\gamma)$ 的除幂加厚范畴中的纤维方块。
若 $T_2\to T$ 平坦且 $U_2=T_2\times_TU$，则作为概形有
$T_3=T_1\times_TT_2$。
\end{lemma}

\begin{proof}
这是因为除幂结构沿平坦环同态可唯一延拓。见
“除幂代数”引理 \ref{dpa-lemma-gamma-extends}。
\end{proof}

\noindent
上述引理表明，除幂加厚的平坦态射经过基变换后仍是平坦态射，
并且事实上就是该态射的“通常”基变换。因此下述定义是合理的。

\begin{definition}
\label{crystalline-definition-big-crystalline-site}
在 情形 \ref{crystalline-situation-global} 中：
\begin{enumerate}
\item $X/S$ 的除幂加厚态射族
$\{(U_i,T_i,\delta_i)\to(U,T,\delta)\}$ 称为一个
{\it Zariski、\'etale、光滑、syntomic 或 fppf 覆盖}，当且仅当
\begin{enumerate}
\item 对所有 $i$ 都有 $U_i=U\times_TT_i$；并且
\item $\{T_i\to T\}$ 是 Zariski、\'etale、光滑、syntomic 或 fppf 覆盖。
\end{enumerate}
\item $X$ 在 $(S,\mathcal I,\gamma)$ 上的{\it 大晶体位点}，
是赋予 Zariski 拓扑的范畴 $\text{CRIS}(X/S)$。
\item $\text{CRIS}(X/S)$ 上的层拓扑斯记为
$(X/S)_{\text{CRIS}}$，有时也记为
$(X/S, \mathcal{I}, \gamma)_{\text{CRIS}}$\footnote{这与我们用
$\Sh(\mathcal C)$ 表示位点 $\mathcal C$ 所对应拓扑斯的约定冲突。}。
\end{enumerate}
\end{definition}

\noindent
这些拓扑斯具有一些显然的函子性。

\begin{remark}[函子性]
\label{crystalline-remark-functoriality-big-cris}
设 $p$ 是素数，并设
$(S,\mathcal I,\gamma)\to(S',\mathcal I',\gamma')$
是 $\mathbf Z_{(p)}$ 上的除幂概形态射。置
$S_0=V(\mathcal I)$ 与 $S'_0=V(\mathcal I')$。设
$$
\xymatrix{
X \ar[r]_f \ar[d] & Y \ar[d] \\
S_0 \ar[r] & S'_0
}
$$
是概形态射的交换图，并假设 $p$ 在 $X$ 与 $Y$ 上局部幂零。
则得到连续且余连续的函子
$$
\text{CRIS}(X/S) \longrightarrow \text{CRIS}(Y/S')
$$
它把 $(U,T,\delta)$ 对应到同一个 $(U,T,\delta)$，并以
$U\to X\to Y$ 作为从 $U$ 到 $Y$ 的 $S'$-态射。因此得到拓扑斯态射
$$
f_{\text{CRIS}} : (X/S)_{\text{CRIS}} \longrightarrow (Y/S')_{\text{CRIS}}
$$
见“位点论”第 \ref{sites-section-cocontinuous-morphism-topoi} 节。
\end{remark}

\begin{remark}[与 Zariski 位点的比较]
\label{crystalline-remark-compare-big-zariski}
在 情形 \ref{crystalline-situation-global} 中，函子 (\ref{crystalline-equation-forget})
是余连续的（细节略），并且与积和纤维积交换
（引理 \ref{crystalline-lemma-divided-power-thickening-fibre-products}）。
因此得到拓扑斯态射
$$
U_{X/S} : (X/S)_{\text{CRIS}} \longrightarrow \Sh((\Sch/X)_{Zar})
$$
从 $X/S$ 的大晶体拓扑斯指向 $X$ 的大 Zariski 拓扑斯。
见“位点论”第 \ref{sites-section-cocontinuous-morphism-topoi} 节。
\end{remark}

\begin{remark}[结构态射]
\label{crystalline-remark-big-structure-morphism}
在 情形 \ref{crystalline-situation-global} 中，考虑闭子概形
$S_0=V(\mathcal I)\subset S$。若假设 $p$ 在 $S_0$ 上局部幂零
（实际中总是如此），则以 $S_0$ 代替 $X$，便得到 定义
\ref{crystalline-definition-divided-power-thickening-X} 中的情形，从而得到位点
$\text{CRIS}(S_0/S)$。若 $f:X\to S_0$ 是 $X$ 在 $S$ 上的结构态射，
则得到环化拓扑斯态射的交换图
$$
\xymatrix{
(X/S)_{\text{CRIS}}
\ar[r]_{f_{\text{CRIS}}} \ar[d]_{U_{X/S}} &
(S_0/S)_{\text{CRIS}} \ar[d]^{U_{S_0/S}} \\
\Sh((\Sch/X)_{Zar}) \ar[r]^{f_{big}} & \Sh((\Sch/S_0)_{Zar}) \ar[rd] \\
& & \Sh((\Sch/S)_{Zar})
}
$$
这来自 注 \ref{crystalline-remark-functoriality-big-cris}。我们把复合
$(X/S)_{\text{CRIS}}\to\Sh((\Sch/S)_{Zar})$ 看作大晶体位点的结构态射。
即使 $p$ 在 $S_0$ 上并非局部幂零，结构态射
$$
(X/S)_{\text{CRIS}} \longrightarrow \Sh((\Sch/S)_{Zar})
$$
仍有定义，因为可以沿上图的下方路径得到它。因此，它是与下述余连续函子
$\text{CRIS}(X/S)\to(\Sch/S)_{Zar}$ 对应的拓扑斯态射：
$(U,T,\delta)/S\mapsto U/S$；见
“位点论”第 \ref{sites-section-cocontinuous-morphism-topoi} 节。
\end{remark}

\begin{remark}[相容性]
\label{crystalline-remark-compatibilities-big-cris}
上述态射满足许多相容性。例如，在 注
\ref{crystalline-remark-functoriality-big-cris} 的情形中，得到环化拓扑斯的交换图
$$
\xymatrix{
(X/S)_{\text{CRIS}} \ar[d] \ar[r] & (Y/S')_{\text{CRIS}} \ar[d] \\
\Sh((\Sch/S)_{Zar}) \ar[r] & \Sh((\Sch/S')_{Zar})
}
$$
其中竖直箭头为结构态射。
\end{remark}




\section{晶体位点}
\label{crystalline-section-site}

\noindent
由于 (\ref{crystalline-equation-forget}) 与积和纤维积交换，考察满足 $U\to X$
为开浸入的那些 $(U,T,\delta)$，便定义出一个在纤维积（更一般地，
非空有限极限）下保持的全子范畴。因此下述定义是合理的。

\begin{definition}
\label{crystalline-definition-crystalline-site}
在 情形 \ref{crystalline-situation-global} 中：
\begin{enumerate}
\item $X$ 在 $(S,\mathcal I,\gamma)$ 上的（小）{\it 晶体位点}，
记为 $\text{Cris}(X/S,\mathcal I,\gamma)$，或简记为
$\text{Cris}(X/S)$，是 $\text{CRIS}(X/S)$ 中由所有满足 $U\to X$
为开浸入的 $(U,T,\delta)$ 所成的全子范畴，并赋予 Zariski 拓扑。
\item $\text{Cris}(X/S)$ 上的层拓扑斯记为 $(X/S)_{\text{cris}}$，
有时也记为 $(X/S,\mathcal I,\gamma)_{\text{cris}}$\footnote{这与我们用
$\Sh(\mathcal C)$ 表示位点 $\mathcal C$ 所对应拓扑斯的约定冲突。}。
\end{enumerate}
\end{definition}

\noindent
对 $\text{Cris}(X/S)$ 中任意 $(U,T,\delta)$，态射 $U\to X$
定义了 $X$ 的小 Zariski 位点 $X_{Zar}$ 的一个对象。因此有典范遗忘函子
\begin{equation}
\label{crystalline-equation-forget-small}
\text{Cris}(X/S) \longrightarrow X_{Zar},\quad
(U, T, \delta) \longmapsto U
\end{equation}
以及左伴随
\begin{equation}
\label{crystalline-equation-endow-trivial-small}
X_{Zar} \longrightarrow \text{Cris}(X/S),\quad
U \longmapsto (U, U, \emptyset)
\end{equation}
后者有时很有用。

\medskip\noindent
正如可以比较一个概形的小、大 Zariski 位点一样，我们也可以比较
小、大晶体位点；见 “拓扑” 引理
\ref{topologies-lemma-at-the-bottom}。

\begin{lemma}
\label{crystalline-lemma-compare-big-small}
假设同定义 \ref{crystalline-definition-divided-power-thickening-X}。包含函子
$$
\text{Cris}(X/S) \to \text{CRIS}(X/S)
$$
与非空有限极限交换，并且是全忠实、连续且余连续的。存在拓扑斯态射
$$
(X/S)_{\text{cris}} \xrightarrow{i} (X/S)_{\text{CRIS}}
\xrightarrow{\pi} (X/S)_{\text{cris}}
$$
其复合为恒等态射，其中第一个由包含函子诱导。此外，$\pi_*=i^{-1}$。
\end{lemma}

\begin{proof}
第一个断言见引理 \ref{crystalline-lemma-divided-power-thickening-fibre-products}。
由此得到拓扑斯态射
$i:(X/S)_{\text{cris}}\to(X/S)_{\text{CRIS}}$，以及左伴随 $i_!$，
满足 $i^{-1}i_!=i^{-1}i_*=\text{id}$；见 “位点论” 引理
\ref{sites-lemma-when-shriek}、\ref{sites-lemma-preserve-equalizers}
及 \ref{sites-lemma-back-and-forth}。
我们断言 $i_!$ 是正合的。若此断言成立，便可由
$\pi^{-1}=i_!$ 与 $\pi_*=i^{-1}$ 定义 $\pi$，其余皆明。
为证明该断言，注意我们已知 $i_!$ 右正合且保持纤维积（见上述引文）。
因此只需证明 $i_!*=*$，其中 $*$ 表示集合层范畴的终对象。
为此，只需给出 $\text{Cris}(X/S)$ 的一族对象
$(U_i,T_i,\delta_i)$，$i\in I$，使得
$$
\coprod\nolimits_{i \in I} h_{(U_i, T_i, \delta_i)} \to *
$$
在 $(X/S)_{\text{CRIS}}$ 中满射（细节略；提示：利用
$\text{Cris}(X/S)$ 具有积，且函子
$\text{Cris}(X/S)\to\text{CRIS}(X/S)$ 与积交换）。
在仿射情形，这由引理 \ref{crystalline-lemma-set-generators} 得出。
一般情形的证明从略。
\end{proof}

\begin{remark}[函子性]
\label{crystalline-remark-functoriality-cris}
设 $p$ 是素数，并设
$(S,\mathcal I,\gamma)\to(S',\mathcal I',\gamma')$
是 $\mathbf Z_{(p)}$ 上的除幂概形态射。设
$$
\xymatrix{
X \ar[r]_f \ar[d] & Y \ar[d] \\
S_0 \ar[r] & S'_0
}
$$
是概形态射的交换图，并假设 $p$ 在 $X$ 与 $Y$ 上局部幂零。
仿照 “拓扑” 引理 \ref{topologies-lemma-morphism-big-small}，定义
$$
f_{\text{cris}} : (X/S)_{\text{cris}} \longrightarrow (Y/S')_{\text{cris}}
$$
公式为 $f_{\text{cris}}=\pi_Y\circ f_{\text{CRIS}}\circ i_X$；
其中 $i_X$ 与 $\pi_Y$ 分别是 引理 \ref{crystalline-lemma-compare-big-small}
中关于 $X$ 与 $Y$ 的态射，而 $f_{\text{CRIS}}$ 如
注 \ref{crystalline-remark-functoriality-big-cris} 所定义。
\end{remark}

\begin{remark}[与 Zariski 位点的比较]
\label{crystalline-remark-compare-zariski}
在 情形 \ref{crystalline-situation-global} 中，函子
(\ref{crystalline-equation-forget-small}) 连续、余连续，并且与积和纤维积交换。
因此得到拓扑斯态射
$$
u_{X/S} : (X/S)_{\text{cris}} \longrightarrow \Sh(X_{Zar})
$$
它联系 $X/S$ 的小晶体拓扑斯与 $X$ 的小 Zariski 拓扑斯。
见“位点论”第 \ref{sites-section-cocontinuous-morphism-topoi} 节。
\end{remark}

\begin{lemma}
\label{crystalline-lemma-localize}
在 情形 \ref{crystalline-situation-global} 中，设 $X'\subset X$ 与
$S'\subset S$ 是开子概形，并且 $X'$ 映入 $S'$。则存在全忠实函子
$\text{Cris}(X'/S')\to\text{Cris}(X/S)$，它诱导一个拓扑斯态射，
使下图交换：
$$
\xymatrix{
(X'/S')_{\text{cris}} \ar[r] \ar[d]_{u_{X'/S'}} &
(X/S)_{\text{cris}} \ar[d]^{u_{X/S}} \\
\Sh(X'_{Zar}) \ar[r] & \Sh(X_{Zar})
}
$$
此外，此图是 “位点论” 引理 \ref{sites-lemma-localize-morphism-topoi}
所述拓扑斯态射局部化的一个例子。
\end{lemma}

\begin{proof}
把 $\text{Cris}(X'/S')$ 的对象看作 $X$ 的除幂加厚
$(U,T,\delta)$，其中 $U\to X$ 通过 $X'\subset X$ 分解，便得到上述
全忠实函子（此时 $T\to S$ 自动通过 $S'$ 分解）。该函子显然余连续，
因而得到所示拓扑斯态射。把 $X'$ 看成 $X_{Zar}$ 的对象，并记其相伴的
可表层为 $h_{X'}\in\Sh(X_{Zar})$。显然 $\Sh(X'_{Zar})$ 是局部化
$\Sh(X_{Zar})/h_{X'}$。另一方面，借助上述函子，范畴
$\text{Cris}(X/S)/u_{X/S}^{-1}h_{X'}$
（见 “位点论” 引理 \ref{sites-lemma-localize-topos-site}）
与 $\text{Cris}(X'/S')$ 典范等同。证毕。
\end{proof}

\begin{remark}[结构态射]
\label{crystalline-remark-structure-morphism}
在 情形 \ref{crystalline-situation-global} 中，考虑闭子概形
$S_0=V(\mathcal I)\subset S$。若假设 $p$ 在 $S_0$ 上局部幂零
（实际中总是如此），则以 $S_0$ 代替 $X$，便得到 定义
\ref{crystalline-definition-divided-power-thickening-X} 中的情形，从而得到位点
$\text{Cris}(S_0/S)$。若 $f:X\to S_0$ 是 $X$ 在 $S$ 上的结构态射，
则得到拓扑斯的交换图
$$
\xymatrix{
(X/S)_{\text{cris}}
\ar[r]_{f_{\text{cris}}} \ar[d]_{u_{X/S}} &
(S_0/S)_{\text{cris}} \ar[d]^{u_{S_0/S}} \\
\Sh(X_{Zar}) \ar[r]^{f_{small}} & \Sh(S_{0, Zar}) \ar[rd] \\
& & \Sh(S_{Zar})
}
$$
见注 \ref{crystalline-remark-functoriality-cris}。我们把复合
$(X/S)_{\text{cris}}\to\Sh(S_{Zar})$ 看作晶体位点的结构态射。
即使 $p$ 在 $S_0$ 上并非局部幂零，结构态射
$$
\tau_{X/S} : (X/S)_{\text{cris}} \longrightarrow \Sh(S_{Zar})
$$
仍有定义，因为可以沿上图的下方路径得到它。
\end{remark}

\begin{remark}[相容性]
\label{crystalline-remark-compatibilities}
上述态射满足许多相容性。例如，在 注
\ref{crystalline-remark-functoriality-cris} 的情形中，得到环化拓扑斯的交换图
$$
\xymatrix{
(X/S)_{\text{cris}} \ar[d] \ar[r] & (Y/S')_{\text{cris}} \ar[d] \\
\Sh((\Sch/S)_{Zar}) \ar[r] & \Sh((\Sch/S')_{Zar})
}
$$
其中竖直箭头为结构态射。
\end{remark}



\section{晶体位点上的层}
\label{crystalline-section-sheaves}

\noindent
记号与假设同情形 \ref{crystalline-situation-global}。为了在本节同时讨论
$X/S$ 的小、大晶体位点，令
$$
\mathcal{C} = \text{CRIS}(X/S)
\quad\text{或}\quad
\mathcal{C} = \text{Cris}(X/S).
$$
$\mathcal C$ 上的层 $\mathcal F$ 对 $\mathcal C$ 的每个对象
$(U,T,\delta)$ 给出一个{\it 限制} $\mathcal F_T$。具体地，
$\mathcal F_T$ 是概形 $T$ 上由下式定义的 Zariski 层：
$$
\mathcal{F}_T(W) = \mathcal{F}(U \cap W, W, \delta|_W)
$$
其中 $W\subset T$ 为开集。此外，若 $f:T\to T'$ 是 $\mathcal C$
的对象 $(U,T,\delta)$ 与 $(U',T',\delta')$ 之间的态射，则有典范
{\it 比较}映射
\begin{equation}
\label{crystalline-equation-comparison}
c_f : f^{-1}\mathcal{F}_{T'} \longrightarrow \mathcal{F}_T.
\end{equation}
事实上，若 $W'\subset T'$ 为开集，则 $f$ 诱导 $\mathcal C$ 中的态射
$$
f|_{f^{-1}W'} :
(U \cap f^{-1}(W'), f^{-1}W', \delta|_{f^{-1}W'})
\longrightarrow
(U' \cap W', W', \delta|_{W'})
$$
因而可用 $\mathcal F$ 的限制映射 $(f|_{f^{-1}W'})^*$ 定义映射
$\mathcal F_{T'}(W')\to\mathcal F_T(f^{-1}W')$。这些映射显然与进一步
限制相容，故定义了从 $\mathcal F_{T'}$ 到 $\mathcal F_T$ 的一个
$f$-映射（见“层”第 \ref{sheaves-section-presheaves-functorial} 节，
尤其是 “层” 定义 \ref{sheaves-definition-f-map}），
从而得到 (\ref{crystalline-equation-comparison}) 中的映射 $c_f$。
注意，若 $f$ 是开浸入，则 $c_f$ 是同构，因为此时 $\mathcal F_T$
正是 $\mathcal F_{T'}$ 在 $T$ 上的限制。

\medskip\noindent
反过来，若对 $\mathcal C$ 的每个对象 $(U,T,\delta)$ 给定 Zariski 层
$\mathcal F_T$，并给定上述比较映射 $c_f$，且它们 (a) 对开浸入为同构，
(b) 满足适当的上闭链条件，则得到 $\mathcal C$ 上的一个层。其证明与
“拓扑” 引理 \ref{topologies-lemma-characterize-sheaf-big} 完全相同。

\medskip\noindent
$\mathcal C$ 上的{\it 结构层}是由下式定义的层 $\mathcal O_{X/S}$：
$$
\mathcal{O}_{X/S} :
(U, T, \delta)
\longmapsto
\Gamma(T, \mathcal{O}_T)
$$
由 $\mathcal C$ 中覆盖的定义可知这是一个层。设 $\mathcal F$ 是
$\mathcal O_{X/S}$-模层。此时比较映射 (\ref{crystalline-equation-comparison})
定义一个 $\mathcal O_T$-模的比较映射
\begin{equation}
\label{crystalline-equation-comparison-modules}
c_f : f^*\mathcal{F}_{T'} \longrightarrow \mathcal{F}_T
\end{equation}

\medskip\noindent
另一类例子从 $(\Sch/X)_{Zar}$ 或 $X_{Zar}$ 上的层 $\mathcal G$ 出发
（取决于 $\mathcal C=\text{CRIS}(X/S)$ 还是
$\mathcal C=\text{Cris}(X/S)$）。则由下式定义的
$\underline{\mathcal G}$ 是 $\mathcal C$ 上的层：
$$
\underline{\mathcal{G}} :
(U, T, \delta)
\longmapsto
\mathcal{G}(U)
$$
特别地，取 $\mathcal G=\mathbf G_a=\mathcal O_X$，得到
$$
\underline{\mathbf{G}_a} :
(U, T, \delta)
\longmapsto
\Gamma(U, \mathcal{O}_U)
$$
对每个对象 $(U,T,\delta)$，典范映射
$\Gamma(T,\mathcal O_T)\to\Gamma(U,\mathcal O_U)$ 定义层的满射
$\mathcal O_{X/S}\to\underline{\mathbf G_a}$。记其核为
$\mathcal J_{X/S}$，从而得到短正合列
$$
0 \to
\mathcal{J}_{X/S} \to
\mathcal{O}_{X/S} \to
\underline{\mathbf{G}_a} \to 0
$$
注意，$\mathcal J_{X/S}$ 配备典范除幂结构。事实上，对每个对象
$(U,T,\delta)$，第三个分量 $\delta$ 本来就是
$\mathcal O_T\to\mathcal O_U$ 之核上的除幂结构。因此（大）晶体拓扑斯
是除幂拓扑斯。





\section{模中的晶体}
\label{crystalline-section-crystals}

\noindent
晶体原来是一种非常一般的对象。不过其定义可能稍难理解，因此我们先在
晶体位点上的模这一情形中给出定义。

\begin{definition}
\label{crystalline-definition-modules}
在情形 \ref{crystalline-situation-global} 中，令
$\mathcal C=\text{CRIS}(X/S)$ 或 $\mathcal C=\text{Cris}(X/S)$，
并令 $\mathcal F$ 是 $\mathcal C$ 上的 $\mathcal O_{X/S}$-模层。
\begin{enumerate}
\item 若对 $\mathcal C$ 的每个对象 $(U,T,\delta)$，限制 $\mathcal F_T$
都是拟凝聚 $\mathcal O_T$-模，则称 $\mathcal F$ {\it 局部拟凝聚}。
\item 若 $\mathcal F$ 在“位点上的模”定义
\ref{sites-modules-definition-site-local} 的意义下拟凝聚，则称
$\mathcal F$ {\it 拟凝聚}。
\item 若所有比较映射 (\ref{crystalline-equation-comparison-modules}) 都是同构，
则称 $\mathcal F$ 是一个{$\mathcal O_{X/S}$-模中的\it 晶体}。
\end{enumerate}
\end{definition}

\noindent
这些概念之间有如下关系。

\begin{lemma}
\label{crystalline-lemma-crystal-quasi-coherent-modules}
采用 定义 \ref{crystalline-definition-modules} 中的记号
$X/S,\mathcal I,\gamma,\mathcal C,\mathcal F$。以下条件等价：
\begin{enumerate}
\item $\mathcal F$ 拟凝聚；
\item $\mathcal F$ 局部拟凝聚，并且是 $\mathcal O_{X/S}$-模中的晶体。
\end{enumerate}
\end{lemma}

\begin{proof}
假设 (1)。设 $f:(U',T',\delta')\to(U,T,\delta)$ 是 $\mathcal C$ 中的态射。
我们要证明：(a) $\mathcal F_T$ 是拟凝聚 $\mathcal O_T$-模；
(b) $c_f:f^*\mathcal F_T\to\mathcal F_{T'}$ 是同构。该假设意味着可以
找到一个覆盖 $\{(T_i,U_i,\delta_i)\to(T,U,\delta)\}$，使得对每个 $i$，
$\mathcal F$ 在 $\mathcal C/(T_i,U_i,\delta_i)$ 上的限制具有全局表现。
由于只需在 Zariski 局部证明 (a) 与 (b)，可以把
$f:(T',U',\delta')\to(T,U,\delta)$ 换成到 $(T_i,U_i,\delta_i)$ 的基变换，
并假设 $\mathcal F$ 在 $\mathcal C/(T,U,\delta)$ 上的限制具有全局表现
$$
\bigoplus\nolimits_{j \in J}
\mathcal{O}_{X/S}|_{\mathcal{C}/(U, T, \delta)} \longrightarrow
\bigoplus\nolimits_{i \in I}
\mathcal{O}_{X/S}|_{\mathcal{C}/(U, T, \delta)} \longrightarrow
\mathcal{F}|_{\mathcal{C}/(U, T, \delta)}
\longrightarrow 0
$$
显然，这给出表现
$$
\bigoplus\nolimits_{j \in J} \mathcal{O}_T \longrightarrow
\bigoplus\nolimits_{i \in I} \mathcal{O}_T \longrightarrow
\mathcal{F}_T
\longrightarrow 0
$$
故 (a) 成立。此外，该表现限制到 $T'$ 后给出 $\mathcal F_{T'}$
的同类表现，故 (b) 成立。

\medskip\noindent
假设 (2)。设 $(U,T,\delta)$ 是 $\mathcal C$ 的对象。我们要找
$(U,T,\delta)$ 的一个覆盖，使 $\mathcal F$ 限制到 $\mathcal C$
在覆盖各成员处的局部化后具有全局表现。因此可以假设 $T$ 仿射。
此时可选取表现
$$
\bigoplus\nolimits_{j \in J} \mathcal{O}_T \longrightarrow
\bigoplus\nolimits_{i \in I} \mathcal{O}_T \longrightarrow
\mathcal{F}_T
\longrightarrow 0
$$
因为已假设 $\mathcal F_T$ 是拟凝聚 $\mathcal O_T$-模。由
$\mathcal F$ 的晶体性质，对 $\mathcal C$ 的任意态射
$f:(U',T',\delta')\to(U,T,\delta)$，此表现拉回为 $\mathcal F_{T'}$
的表现。于是得到 $\mathcal F|_{\mathcal C/(U,T,\delta)}$ 所需的表现。
\end{proof}

\begin{definition}
\label{crystalline-definition-crystal-quasi-coherent-modules}
若 $\mathcal F$ 满足 引理 \ref{crystalline-lemma-crystal-quasi-coherent-modules}
中的等价条件，则称 $\mathcal F$ 是一个{\it 拟凝聚模中的晶体}。
若此外 $\mathcal F$ 还是有限局部自由的，则称它是一个
{\it 有限局部自由模中的晶体}。
\end{definition}

\noindent
当然，引理 \ref{crystalline-lemma-crystal-quasi-coherent-modules} 表明这种说法略显
累赘，因为拟凝聚模总是晶体。不过这是文献中的标准术语。

\begin{remark}
\label{crystalline-remark-crystal}
为表述晶体的一般概念，我们采用栈与强笛卡尔态射的语言；见
“栈” 定义 \ref{stacks-definition-stack} 及 “范畴”
定义 \ref{categories-definition-cartesian-over-C}。
在 情形 \ref{crystalline-situation-global} 中，设
$p:\mathcal C\to\text{Cris}(X/S)$ 是一个栈。一个
{$X$ 上相对于 $S$ 的、取值于 $\mathcal C$ 对象的\it 晶体}，
是一个{\it 笛卡尔截面} $\sigma:\text{Cris}(X/S)\to\mathcal C$；
亦即函子 $\sigma$ 满足 $p\circ\sigma=\text{id}$，且对
$\text{Cris}(X/S)$ 的每个态射 $f$，$\sigma(f)$ 都是强笛卡尔的。
大晶体位点的情形类似。
\end{remark}





\section{微分层}
\label{crystalline-section-differentials-sheaf}

\noindent
本节只使用（小）晶体位点，因为这似乎更自然。我们如下整体化
定义 \ref{crystalline-definition-derivation}。

\begin{definition}
\label{crystalline-definition-global-derivation}
在 情形 \ref{crystalline-situation-global} 中，设 $\mathcal F$ 是
$\text{Cris}(X/S)$ 上的 $\mathcal O_{X/S}$-模层。一个
{$S$-{\it 导子} $D:\mathcal O_{X/S}\to\mathcal F$} 是层的映射，
使得对 $\text{Cris}(X/S)$ 的每个对象 $(U,T,\delta)$，映射
$$
D : \Gamma(T, \mathcal{O}_T) \longrightarrow \Gamma(T, \mathcal{F})
$$
是除幂 $\Gamma(V,\mathcal O_V)$-导子，其中 $V\subset S$ 是任意开集，
使得 $T\to S$ 通过 $V$ 分解。
\end{definition}

\noindent
这意味着 $D$ 可加，满足 Leibniz 法则，消去来自 $S$ 的函数，并且对
除幂理想 $\mathcal J_{X/S}$ 的局部截面 $f$ 满足
$D(f^{[n]})=f^{[n-1]}D(f)$。这是下面将描述的一个非常一般概念的特例。

\medskip\noindent
请把以下讨论与“位点上的模”第
\ref{sites-modules-section-differentials} 节比较。设 $\mathcal C$ 是位点，
$\mathcal A\to\mathcal B$ 是 $\mathcal C$ 上环层的映射，
$\mathcal J\subset\mathcal B$ 是理想层，$\delta$ 是 $\mathcal J$
上的除幂结构，而 $\mathcal F$ 是 $\mathcal B$-模层。此时可以定义
{\it 除幂 $\mathcal A$-导子} $D:\mathcal B\to\mathcal F$：这意味着
$D$ 是 $\mathcal A$-线性的，满足 Leibniz 法则，并且对 $\mathcal J$
的局部截面 $x$ 满足 $D(\delta_n(x))=\delta_{n-1}(x)D(x)$。
在这种情形中，存在一个{\it 泛除幂 $\mathcal A$-导子}
$$
\text{d}_{\mathcal{B}/\mathcal{A}, \delta} :
\mathcal{B}
\longrightarrow
\Omega_{\mathcal{B}/\mathcal{A}, \delta}
$$
此外，$\text{d}_{\mathcal B/\mathcal A,\delta}$ 是复合
$$
\mathcal{B}
\longrightarrow
\Omega_{\mathcal{B}/\mathcal{A}}
\longrightarrow
\Omega_{\mathcal{B}/\mathcal{A}, \delta}
$$
其中第一个映射是在“位点上的模”引理
\ref{sites-modules-lemma-universal-module} 的证明中构造的泛导子，
第二个箭头则是对由下列局部截面生成的子模取商：
$\text{d}_{\mathcal{B}/\mathcal{A}}(\delta_n(x)) -
\delta_{n - 1}(x)\text{d}_{\mathcal{B}/\mathcal{A}}(x)$.

\medskip\noindent
我们把它转化为如下相对概念。设
$(f,f^\sharp):(\Sh(\mathcal C),\mathcal O)\to
(\Sh(\mathcal C'),\mathcal O')$ 是环化拓扑斯态射，
$\mathcal J\subset\mathcal O$ 是理想层，$\delta$ 是 $\mathcal J$
上的除幂结构，而 $\mathcal F$ 是 $\mathcal O$-模层。若 $D$ 是上述意义下
的除幂 $f^{-1}\mathcal O'$-导子，便称 $D:\mathcal O\to\mathcal F$
为除幂 $\mathcal O'$-导子。此外，记
$$
\Omega_{\mathcal{O}/\mathcal{O}', \delta} =
\Omega_{\mathcal{O}/f^{-1}\mathcal{O}', \delta}
$$
它是泛除幂 $\mathcal O'$-导子的取值模。

\medskip\noindent
把这应用于结构态射
$$
(X/S)_{\text{Cris}} \longrightarrow \Sh(S_{Zar})
$$
（见注 \ref{crystalline-remark-structure-morphism}），便恢复上述 定义
\ref{crystalline-definition-global-derivation} 的概念。特别地，存在泛除幂导子
$$
d_{X/S} : \mathcal{O}_{X/S} \to \Omega_{X/S}
$$
注意，我们在记号中省略了表示微分模与除幂相容的修饰符（似乎不太可能
有人会考虑晶体位点结构层的通常微分模）。

\begin{lemma}
\label{crystalline-lemma-module-differentials-divided-power-scheme}
设 $(T,\mathcal J,\delta)$ 是除幂概形，并设 $T\to S$ 是概形态射。
上述商 $\Omega_{T/S}\to\Omega_{T/S,\delta}$ 是拟凝聚
$\mathcal O_T$-模。对映入仿射开集 $V\subset S$ 的仿射开集
$W\subset T$，有
$$
\Gamma(W, \Omega_{T/S, \delta}) =
\Omega_{\Gamma(W, \mathcal{O}_W)/\Gamma(V, \mathcal{O}_V), \delta}
$$
其中右端如第 \ref{crystalline-section-differentials} 节中所构造。
\end{lemma}

\begin{proof}
略。
\end{proof}

\begin{lemma}
\label{crystalline-lemma-module-of-differentials}
在 情形 \ref{crystalline-situation-global} 中，对
$\text{Cris}(X/S)$ 中的 $(U,T,\delta)$，$\Omega_{X/S}$ 在 $T$ 上的限制
$(\Omega_{X/S})_T$ 等于 $\Omega_{T/S,\delta}$，而限制
$\text{d}_{X/S}|_T$ 等于 $\text{d}_{T/S,\delta}$。
\end{lemma}

\begin{proof}
略。
\end{proof}

\begin{lemma}
\label{crystalline-lemma-module-of-differentials-on-affine}
在 情形 \ref{crystalline-situation-global} 中，对 $\text{Cris}(X/S)$
中映入仿射开集 $V\subset S$ 的任意仿射对象 $(U,T,\delta)$，有
$$
\Gamma((U, T, \delta), \Omega_{X/S}) =
\Omega_{\Gamma(T, \mathcal{O}_T)/\Gamma(V, \mathcal{O}_V), \delta}
$$
其中右端如第 \ref{crystalline-section-differentials} 节中所构造。
\end{lemma}

\begin{proof}
合用 引理 \ref{crystalline-lemma-module-differentials-divided-power-scheme}
与 \ref{crystalline-lemma-module-of-differentials}。
\end{proof}

\begin{lemma}
\label{crystalline-lemma-describe-omega-small}
在 情形 \ref{crystalline-situation-global} 中，设 $(U,T,\delta)$ 是
$\text{Cris}(X/S)$ 的对象。令
$$
(U(1), T(1), \delta(1)) = (U, T, \delta) \times (U, T, \delta)
$$
为 $\text{Cris}(X/S)$ 中的积。令
$\mathcal K\subset\mathcal O_{T(1)}$ 是与闭浸入
$\Delta:T\to T(1)$ 对应的拟凝聚理想层。则
$\mathcal K\subset\mathcal J_{T(1)}$ 被 $\mathcal J_{T(1)}$
上的除幂结构保持，并且
$$
(\Omega_{X/S})_T = \mathcal{K}/\mathcal{K}^{[2]}
$$
\end{lemma}

\begin{proof}
由于 $U\to X$ 是开浸入，且 (\ref{crystalline-equation-forget-small}) 与积交换，
故 $U=U(1)$，从而 $\mathcal K\subset\mathcal J_{T(1)}$。有了这一事实，
在 $T$ 上仿射局部地工作，并使用 引理
\ref{crystalline-lemma-module-of-differentials-on-affine} 与
\ref{crystalline-lemma-diagonal-and-differentials-affine-site}，即得结论。
\end{proof}

\noindent
$\Omega_{X/S}$ 并不是拟凝聚 $\mathcal O_{X/S}$-模中的晶体，
但它满足两个紧密相关的性质（与 引理
\ref{crystalline-lemma-crystal-quasi-coherent-modules} 比较）。

\begin{lemma}
\label{crystalline-lemma-omega-locally-quasi-coherent}
在 情形 \ref{crystalline-situation-global} 中，微分层 $\Omega_{X/S}$
具有以下两个性质：
\begin{enumerate}
\item $\Omega_{X/S}$ 局部拟凝聚；
\item 对 $\text{Cris}(X/S)$ 中任意态射
$(U,T,\delta)\to(U',T',\delta')$，若 $f:T\to T'$ 是闭浸入，
则映射 $c_f:f^*(\Omega_{X/S})_{T'}\to(\Omega_{X/S})_T$ 满射。
\end{enumerate}
\end{lemma}

\begin{proof}
(1) 由引理 \ref{crystalline-lemma-module-differentials-divided-power-scheme}
与 \ref{crystalline-lemma-module-of-differentials} 合用得出。
(2) 由以下事实得出：$(\Omega_{X/S})_T=\Omega_{T/S,\delta}$
是 $\Omega_{T/S}$ 的商，且
$f^*\Omega_{T'/S}\to\Omega_{T/S}$ 满射。
\end{proof}







\section{两个泛加厚}
\label{crystalline-section-universal-thickenings}

\noindent
本节的构造将帮助我们在晶体位点上的模晶体上定义联络。在某种意义下，
这里的构造是在第 \ref{crystalline-section-explicit-thickenings} 节中给出的
“层化、泛化”版本。

\begin{remark}
\label{crystalline-remark-first-order-thickening}
在 情形 \ref{crystalline-situation-global} 中，设 $(U,T,\delta)$ 是
$\text{Cris}(X/S)$ 的对象。记
$\Omega_{T/S,\delta}=(\Omega_{X/S})_T$；见
引理 \ref{crystalline-lemma-module-of-differentials}。我们显式描述 $T$
的一个一阶加厚 $T'$。具体地，置
$$
\mathcal{O}_{T'} = \mathcal{O}_T \oplus \Omega_{T/S, \delta}
$$
并赋予代数结构，使 $\Omega_{T/S,\delta}$ 是平方为零的理想。
令 $\mathcal J\subset\mathcal O_T$ 是闭浸入 $U\to T$ 的理想层，置
$\mathcal J'=\mathcal J\oplus\Omega_{T/S,\delta}$。由下式在
$\mathcal J'$ 上定义除幂结构：
$$
\delta_n'(f, \omega) = (\delta_n(f), \delta_{n - 1}(f)\omega),
$$
见引理 \ref{crystalline-lemma-divided-power-first-order-thickening}。有两个环映射
$$
p_0, p_1 : \mathcal{O}_T \to \mathcal{O}_{T'}
$$
第一个由 $f\mapsto(f,0)$ 给出，第二个由
$f\mapsto(f,\text{d}_{T/S,\delta}f)$ 给出。注意二者都与
$\mathcal J$ 和 $\mathcal J'$ 上的除幂结构相容，商映射
$\mathcal O_{T'}\to\mathcal O_T$ 亦然。因此得到
$\text{Cris}(X/S)$ 的对象 $(U,T',\delta')$ 及交换图
$$
\xymatrix{
& T \ar[ld]_{\text{id}} \ar[d]^i \ar[rd]^{\text{id}} \\
T & T' \ar[l]_{p_0} \ar[r]^{p_1} & T
}
$$
于 $\text{Cris}(X/S)$ 中成立；其中 $i$ 是一阶加厚，其理想层与
$\Omega_{T/S,\delta}$ 等同，并且
$p_1-p_0:\mathcal O_T\to\mathcal O_{T'}$ 与泛导子
$\text{d}_{T/S,\delta}$ 复合包含
$\Omega_{T/S,\delta}\to\mathcal O_{T'}$ 后的映射等同。
\end{remark}

\begin{remark}
\label{crystalline-remark-second-order-thickening}
在 情形 \ref{crystalline-situation-global} 中，设 $(U,T,\delta)$ 是
$\text{Cris}(X/S)$ 的对象。记
$\Omega_{T/S,\delta}=(\Omega_{X/S})_T$；见
引理 \ref{crystalline-lemma-module-of-differentials}。并以
$\Omega^2_{T/S,\delta}$ 表示其二次外幂。我们显式描述 $T$
的一个二阶加厚 $T''$。具体地，置
$$
\mathcal{O}_{T''} =
\mathcal{O}_T \oplus \Omega_{T/S, \delta} \oplus \Omega_{T/S, \delta}
\oplus \Omega^2_{T/S, \delta}
$$
其代数结构定义如下：
$$
(f, \omega_1, \omega_2, \eta) \cdot
(f', \omega_1', \omega_2', \eta') =
(ff', f\omega_1' + f'\omega_1, f\omega_2' + f'\omega_2,
f\eta' + f'\eta + \omega_1 \wedge \omega_2' + \omega_1' \wedge \omega_2).
$$
令 $\mathcal J\subset\mathcal O_T$ 是闭浸入 $U\to T$ 的理想层，
并令 $\mathcal J''$ 是 $\mathcal J$ 在投影
$\mathcal O_{T''}\to\mathcal O_T$ 下的逆像。由下式在
$\mathcal J''$ 上定义除幂结构：
$$
\delta_n''(f, \omega_1, \omega_2, \eta) =
(\delta_n(f), \delta_{n - 1}(f)\omega_1, \delta_{n - 1}(f)\omega_2,
\delta_{n - 1}(f)\eta + \delta_{n - 2}(f)\omega_1 \wedge \omega_2)
$$
见引理 \ref{crystalline-lemma-divided-power-second-order-thickening}。有三个环映射
$q_0,q_1,q_2:\mathcal O_T\to\mathcal O_{T''}$，定义为
\begin{align*}
q_0(f) & = (f, 0, 0, 0), \\
q_1(f) & = (f, \text{d}f, 0, 0), \\
q_2(f) & = (f, \text{d}f, \text{d}f, 0)
\end{align*}
其中 $\text d=\text d_{T/S,\delta}$。注意三者都与 $\mathcal J$
及 $\mathcal J''$ 上的除幂结构相容。另有三个环映射
$q_{01},q_{12},q_{02}:\mathcal O_{T'}\to\mathcal O_{T''}$，
其中 $\mathcal O_{T'}$ 如 注 \ref{crystalline-remark-first-order-thickening}
所定义。具体地，置
\begin{align*}
q_{01}(f, \omega) & = (f, \omega, 0, 0), \\
q_{12}(f, \omega) & =
(f, \text{d}f, \omega, \text{d}\omega), \\
q_{02}(f, \omega) & = (f, \omega, \omega, 0)
\end{align*}
这些映射也与给定的除幂结构相容。我们验证 $q_{12}$：注意
$q_{12}$ 是环同态，因为
\begin{align*}
q_{12}(f, \omega)q_{12}(g, \eta) & =
(f, \text{d}f, \omega, \text{d}\omega)(g, \text{d}g, \eta, \text{d}\eta) \\
& =
(fg, f\text{d}g + g \text{d}f, f\eta + g\omega,
f\text{d}\eta + g\text{d}\omega + \text{d}f \wedge \eta +
\text{d}g \wedge \omega) \\
& = q_{12}(fg, f\eta + g\omega) = q_{12}((f, \omega)(g, \eta))
\end{align*}
并且 $q_{12}$ 与除幂相容，因为
\begin{align*}
\delta_n''(q_{12}(f, \omega)) & =
\delta_n''((f, \text{d}f, \omega, \text{d}\omega)) \\
& =
(\delta_n(f), \delta_{n - 1}(f)\text{d}f, \delta_{n - 1}(f)\omega,
\delta_{n - 1}(f)\text{d}\omega + \delta_{n - 2}(f)\text{d}(f) \wedge \omega)
\\
& = q_{12}((\delta_n(f), \delta_{n - 1}(f)\omega)) =
q_{12}(\delta'_n(f, \omega))
\end{align*}
对 $q_{01}$ 与 $q_{02}$ 的验证更为容易。注意
$q_0=q_{01}\circ p_0$、$q_1=q_{01}\circ p_1$、
$q_1=q_{12}\circ p_0$、$q_2=q_{12}\circ p_1$、
$q_0=q_{02}\circ p_0$，以及 $q_2=q_{02}\circ p_1$。
因此 $(U,T'',\delta'')$ 是 $\text{Cris}(X/S)$ 的对象，并得到态射
$$
\xymatrix{
T''
\ar@<2ex>[r]
\ar@<0ex>[r]
\ar@<-2ex>[r]
&
T'
\ar@<1ex>[r]
\ar@<-1ex>[r]
&
T
}
$$
于 $\text{Cris}(X/S)$ 中满足上述关系。在应用中，我们用
$q_i:T''\to T$ 与 $q_{ij}:T''\to T'$ 表示与上述环映射相伴的态射。
\end{remark}






\section{德拉姆复形}
\label{crystalline-section-de-Rham}

\noindent
在 情形 \ref{crystalline-situation-global} 中，我们在（小）晶体位点上工作，
对 $i\geq0$ 定义
$\Omega^i_{X/S}=\wedge^i_{\mathcal O_{X/S}}\Omega_{X/S}$。
泛 $S$-导子 $\text d_{X/S}$ 给出 $\text{Cris}(X/S)$ 上的
{\it 德拉姆复形}
$$
\mathcal{O}_{X/S} \to \Omega^1_{X/S} \to \Omega^2_{X/S} \to \ldots
$$
见引理 \ref{crystalline-lemma-module-of-differentials-on-affine} 与
注 \ref{crystalline-remark-divided-powers-de-rham-complex}。


\section{联络}
\label{crystalline-section-connections}

\noindent
在 情形 \ref{crystalline-situation-global} 中，给定 $\text{Cris}(X/S)$
上的 $\mathcal O_{X/S}$-模 $\mathcal F$，一个{\it 联络}是交换群层的映射
$$
\nabla :
\mathcal{F}
\longrightarrow
\mathcal{F} \otimes_{\mathcal{O}_{X/S}} \Omega_{X/S}
$$
使得对 $\mathcal F$ 与 $\mathcal O_{X/S}$ 的局部截面 $s,f$，有
$\nabla(fs)=f\nabla(s)+s\otimes\text df$。给定联络，存在典范映射
$
\nabla :
\mathcal{F} \otimes_{\mathcal{O}_{X/S}} \Omega^i_{X/S}
\longrightarrow
\mathcal{F} \otimes_{\mathcal{O}_{X/S}} \Omega^{i + 1}_{X/S}
$
其定义为 $\nabla(s\otimes\omega)=\nabla(s)\wedge\omega+
s\otimes\text d\omega$，同注 \ref{crystalline-remark-connection}。
若 $\nabla\circ\nabla=0$，则称该联络{\it 可积}。若 $\nabla$ 可积，
便得到 $\text{Cris}(X/S)$ 上的{\it 德拉姆复形}
$$
\mathcal{F} \to
\mathcal{F} \otimes_{\mathcal{O}_{X/S}} \Omega^1_{X/S} \to
\mathcal{F} \otimes_{\mathcal{O}_{X/S}} \Omega^2_{X/S} \to \ldots
$$
任意 $\mathcal O_{X/S}$-模中的晶体都典范地配备一个可积联络。

\begin{lemma}
\label{crystalline-lemma-automatic-connection}
在 情形 \ref{crystalline-situation-global} 中，设 $\mathcal F$ 是
$\text{Cris}(X/S)$ 上的 $\mathcal O_{X/S}$-模中的晶体。
则 $\mathcal F$ 典范地配备一个可积联络。
\end{lemma}

\begin{proof}
设 $(U,T,\delta)$ 是 $\text{Cris}(X/S)$ 的对象。令
$(U,T',\delta')$ 是 注 \ref{crystalline-remark-first-order-thickening}
中用 $(\Omega_{X/S})_T=\Omega_{T/S,\delta}$ 构造的 $T$ 的无穷小加厚。
它配备投影 $p_0,p_1:T'\to T$ 与对角映射 $i:T\to T'$。
由假设得到 $\mathcal O_{T'}$-模的同构
$$
p_0^*\mathcal{F}_T \xrightarrow{c_0}
\mathcal{F}_{T'} \xleftarrow{c_1}
p_1^*\mathcal{F}_T
$$
把 $c=c_1^{-1}\circ c_0$ 沿 $i$ 拉回到 $T$，得到 $\mathcal F_T$
的恒等映射。因此，若 $s\in\Gamma(T,\mathcal F_T)$，则
$\nabla(s)=p_1^*s-c(p_0^*s)$ 是 $p_1^*\mathcal F_T$ 的截面，沿 $i$
拉回后为零。所以 $\nabla(s)$ 是下列模的截面：
$$
\mathcal{F}_T
\otimes_{\mathcal{O}_T}
\Omega_{T/S, \delta}
$$
因为按构造 $\mathcal O_{T'}=\mathcal O_T\oplus\Omega_{T/S,\delta}$，
这正是 $p_1^*\mathcal F_T\to\mathcal F_T$ 的核。利用注
\ref{crystalline-remark-first-order-thickening} 中对 $\text d$ 的描述，容易验证
$\nabla(fs)=f\nabla(s)+s\otimes\text d(f)$。

\medskip\noindent
由此得到的映射族
$$
\nabla : \Gamma(T, \mathcal{F}_T) \to
\Gamma(T, \mathcal{F}_T \otimes_{\mathcal{O}_T} \Omega_{T/S, \delta})
$$
关于 $T$ 具有函子性，因为 $T'$ 的构造关于 $T$ 具有函子性。
因此得到一个联络。

\medskip\noindent
为证明该联络可积，考虑 注 \ref{crystalline-remark-second-order-thickening}
中构造的对象 $(U,T'',\delta'')$。由于 $\mathcal F$ 是层，可知
$$
\xymatrix{
q_0^*\mathcal{F}_T \ar[rr]_{q_{01}^*c} \ar[rd]_{q_{02}^*c} & &
q_1^*\mathcal{F}_T \ar[ld]^{q_{12}^*c} \\
& q_2^*\mathcal{F}_T
}
$$
是 $\mathcal O_{T''}$-模的交换图。对
$s\in\Gamma(T,\mathcal F_T)$，有
$c(p_0^*s)=p_1^*s-\nabla(s)$。写成
$\nabla(s)=\sum p_1^*s_i\cdot\omega_i$，其中 $s_i$ 是
$\mathcal F_T$ 的局部截面，$\omega_i$ 是 $\Omega_{T/S,\delta}$
的局部截面。把 $\omega_i$ 看作 $\mathcal O_{T'}$ 的结构层的局部截面，
所以这里写乘积而非张量积。一方面
\begin{align*}
q_{12}^*c \circ q_{01}^*c(q_0^*s) & = 
q_{12}^*c(q_1^*s - \sum q_1^*s_i \cdot q_{01}^*\omega_i) \\
& =
q_2^*s - \sum q_2^*s_i \cdot q_{12}^*\omega_i -
\sum q_2^*s_i \cdot q_{01}^*\omega_i +
\sum q_{12}^*\nabla(s_i) \cdot q_{01}^*\omega_i
\end{align*}
另一方面
$$
q_{02}^*c(q_0^*s) = q_2^*s - \sum q_2^*s_i \cdot q_{02}^*\omega_i.
$$
由注 \ref{crystalline-remark-second-order-thickening} 中的公式可知
$q_{01}^*\omega_i + q_{12}^*\omega_i - q_{02}^*\omega_i = \text{d}\omega_i$.
因此上述两个表达式之差为
$$
\sum q_2^*s_i \cdot \text{d}\omega_i -
\sum q_{12}^*\nabla(s_i) \cdot q_{01}^*\omega_i
$$
由 $\mathcal O_{T''}$ 上乘法的定义，注意
$q_{12}^*\omega\cdot q_{01}^*\omega'=\omega'\wedge\omega
=-\omega\wedge\omega'$。因此，上式正是 $\nabla^2(s)$，看作
$q_2^*\mathcal F$ 的子层
$\mathcal F_T\otimes\Omega^2_{T/S,\delta}$ 的截面。
交换性遂给出可积条件。
\end{proof}




\section{余单纯代数}
\label{crystalline-section-cosimplicial}

\noindent
本节应移至别处。一个{\it 余单纯环}是环范畴中的余单纯对象。
给定环 $R$，一个{\it 余单纯 $R$-代数}是 $R$-代数范畴中的余单纯对象。
余单纯环 $A_*$ 中的一个{\it 余单纯理想}，是对每个 $n$ 给定理想
$I_n\subset A_n$，使得对 $\Delta$ 中所有 $f:[n]\to[m]$，都有
$A(f)(I_n)\subset I_m$。

\medskip\noindent
设 $A_*$ 是余单纯环。令 $\mathcal C$ 为如下偶对 $(A,M)$ 所成的范畴：
$A$ 是环，$M$ 是 $A$-模。态射 $(A,M)\to(A',M')$ 由环映射
$A\to A'$ 与 $A$-模映射 $M\to M'$ 组成；这里借助 $A\to A'$ 及
$M'$ 的 $A'$-模结构，把 $M'$ 看作 $A$-模。于是可以把
{$A_*$ 上的\it 余单纯模 $M_*$} 定义为 $\mathcal C$ 的余单纯对象
$(A_*,M_*)$，其第一分量等于 $A_*$。{$A_*$ 上余单纯模的一个\it 同态
$\varphi_*:M_*\to N_*$}，是 $\mathcal C$ 中余单纯对象的态射
$(A_*,M_*)\to(A_*,N_*)$，其第一分量为 $1_{A_*}$。

\medskip\noindent
在 $A_*$ 上余单纯模同态 $\varphi_*,\psi_*:M_*\to N_*$ 之间的
一个{\it 同伦}，是相伴映射 $(A_*,M_*)\to(A_*,N_*)$ 之间的同伦，
其第一分量为平凡同伦（与 “单纯方法” 例
\ref{simplicial-example-trivial-homotopy} 对偶）。具体而言，这样的同伦是
$$
h : M_* \longrightarrow \Hom(\Delta[1], N_*)
$$
作为余单纯交换群同态的 $\varphi_*$ 与 $\psi_*$ 之间的同伦，
并且对每个 $n$，映射
$h_n:M_n\to\prod_{\alpha\in\Delta[1]_n}N_n$ 是 $A_n$-线性的。
下述引理是 “单纯方法” 引理
\ref{simplicial-lemma-functorial-homotopy} 对余单纯模的版本。

\begin{lemma}
\label{crystalline-lemma-homotopy-tensor}
设 $A_*$ 是余单纯环，并设 $\varphi_*,\psi_*:K_*\to M_*$
是余单纯 $A_*$-模的同态。
\begin{enumerate}
\item
\label{crystalline-item-tensor}
若 $\varphi_*$ 与 $\psi_*$ 同伦，则对任意余单纯 $A_*$-模 $L_*$，
$$
\varphi_* \otimes 1, \psi_* \otimes 1 :
K_* \otimes_{A_*} L_* \longrightarrow M_* \otimes_{A_*} L_*
$$
同伦。
\item
\label{crystalline-item-wedge}
若 $\varphi_*$ 与 $\psi_*$ 同伦，则
$$
\wedge^i(\varphi_*), \wedge^i(\psi_*) :
\wedge^i(K_*) \longrightarrow \wedge^i(M_*)
$$
同伦。
\item
\label{crystalline-item-base-change}
若 $\varphi_*$ 与 $\psi_*$ 同伦，且 $A_*\to B_*$ 是余单纯环同态，则
$$
\varphi_* \otimes 1, \psi_* \otimes 1 :
K_* \otimes_{A_*} B_* \longrightarrow M_* \otimes_{A_*} B_*
$$
作为余单纯 $B_*$-模同态是同伦的。
\item
\label{crystalline-item-completion}
若 $I_*\subset A_*$ 是余单纯理想，则完备化之间的诱导映射
$$
\varphi^\wedge_*, \psi^\wedge_* :
K_*^\wedge \longrightarrow M_*^\wedge
$$
同伦。
\item 可按需在此添加更多情形，例如对称幂。
\end{enumerate}
\end{lemma}

\begin{proof}
设 $h:M_*\longrightarrow\Hom(\Delta[1],N_*)$ 是给定同伦。在次数 $n$，有
$$
h_n = (h_{n, \alpha}) :
K_n \longrightarrow
\prod\nolimits_{\alpha \in \Delta[1]_n} K_n
$$
见“单纯方法”第 \ref{simplicial-section-homotopy-cosimplicial} 节。
一族 $h_{n,\alpha}$ 构成同伦的充要条件是：对每个
$f:[n]\to[m]$，都有
$$
h_{m, \alpha} \circ M_*(f) = N_*(f) \circ h_{n, \alpha \circ f}
$$
见 “单纯方法” Equation
(\ref{simplicial-equation-property-homotopy-cosimplicial})。
此外还应有 $\psi_n=h_{n,0:[n]\to[1]}$ 与
$\varphi_n=h_{n,1:[n]\to[1]}$。

\medskip\noindent
在引理的各情形中，我们都可给出相应映射。
情形 (\ref{crystalline-item-tensor})：使用同伦 $h\otimes1$，其次数 $n$ 分量定义为
$$
(h \otimes 1)_{n, \alpha} = h_{n, \alpha} \otimes 1_{L_n} :
K_n \otimes_{A_n} L_n
\longrightarrow
M_n \otimes_{A_n} L_n.
$$
情形 (\ref{crystalline-item-wedge})：使用同伦 $\wedge^ih$，其次数 $n$ 分量定义为
$$
\wedge^i(h)_{n, \alpha} = \wedge^i(h_{n, \alpha}) :
\wedge_{A_n}(K_n)
\longrightarrow
\wedge^i_{A_n}(M_n).
$$
情形 (\ref{crystalline-item-base-change})：使用同伦 $h\otimes1$，其次数 $n$ 分量定义为
$$
(h \otimes 1)_{n, \alpha} = h_{n, \alpha} \otimes 1 :
K_n \otimes_{A_n} B_n
\longrightarrow
M_n \otimes_{A_n} B_n.
$$
情形 (\ref{crystalline-item-completion})：使用同伦 $h^\wedge$，其次数 $n$ 分量定义为
$$
(h^\wedge)_{n, \alpha} = h_{n, \alpha}^\wedge :
K_n^\wedge
\longrightarrow
M_n^\wedge.
$$
这是成立的，因为每个 $h_{n,\alpha}$ 都是 $A_n$-线性的。
\end{proof}






\section{拟凝聚模中的晶体}
\label{crystalline-section-quasi-coherent-crystals}

\noindent
在 情形 \ref{crystalline-situation-affine} 中，置 $X=\Spec(C)$ 与
$S=\Spec(A)$。我们将分类 $\text{Cris}(X/S)$ 上拟凝聚模中的晶体。
在此之前先固定一些记号。

\medskip\noindent
选取 $A$ 上多项式环 $P=A[x_i]$ 及 $A$-代数满射 $P\to C$，其核为
$J=\Ker(P\to C)$。置
\begin{equation}
\label{crystalline-equation-D}
D = \lim_e D_{P, \gamma}(J) / p^eD_{P, \gamma}(J)
\end{equation}
这是 $p$-进完备化的除幂包络。该环配备除幂理想 $\bar J$ 与除幂结构
$\bar\gamma$；见引理 \ref{crystalline-lemma-list-properties}。置
$D_e=D/p^eD$，并以 $\bar J_e$ 表示 $\bar J$ 在 $D_e$ 中的像。
我们用简记
\begin{equation}
\label{crystalline-equation-omega-D}
\Omega_D = \lim_e \Omega_{D_e/A, \bar\gamma} =
\lim_e \Omega_{D/A, \bar\gamma}/p^e\Omega_{D/A, \bar\gamma}
\end{equation}
表示除幂微分模的 $p$-进完备化；见引理
\ref{crystalline-lemma-differentials-completion}。它也是
$\Omega_{D_{P,\gamma}(J)/A,\bar\gamma}$ 的 $p$-进完备化，
而后者以 $\text d x_i$ 为自由基；见引理
\ref{crystalline-lemma-module-differentials-divided-power-envelope}。
因此，$\Omega_D$ 的任意元素可唯一写成和
$\sum f_i\text d x_i$，并且对每个 $e$，只有有限多个 $f_i$
不属于 $p^eD$。此外，各映射
$\text d_{D_e/A,\bar\gamma}:D_e\to\Omega_{D_e/A,\bar\gamma}$
拼合成 $p$-进完备化上的一个除幂 $A$-导子
\begin{equation}
\label{crystalline-equation-derivation-D}
\text{d} : D \longrightarrow \Omega_D
\end{equation}

\medskip\noindent
我们还需要“$\Spec(D)$ 的积 $\Spec(D(n))$”；解释见命题
\ref{crystalline-proposition-compute-cohomology} 及其证明。其形式定义如下。
对 $n\geq0$，令
$J(n)=\Ker(P\otimes_A\ldots\otimes_AP\to C)$，其中张量积有
$n+1$ 个因子。置
\begin{equation}
\label{crystalline-equation-Dn}
D(n) = \lim_e
D_{P \otimes_A \ldots \otimes_A P, \gamma}(J(n))/
p^eD_{P \otimes_A \ldots \otimes_A P, \gamma}(J(n))
\end{equation}
即除幂包络的 $p$-进完备化。以 $\bar J(n)$ 表示其除幂理想，
以 $\bar\gamma(n)$ 表示其除幂。另引入
$D(n)_e=D(n)/p^eD(n)$、$p$-进完备化的微分模
\begin{equation}
\label{crystalline-equation-omega-Dn}
\Omega_{D(n)} = \lim_e \Omega_{D(n)_e/A, \bar\gamma} =
\lim_e \Omega_{D(n)/A, \bar\gamma}/p^e\Omega_{D(n)/A, \bar\gamma}
\end{equation}
以及导子
\begin{equation}
\label{crystalline-equation-derivation-Dn}
\text{d} : D(n) \longrightarrow \Omega_{D(n)}
\end{equation}
当然 $D=D(0)$。注意，各环 $D(0),D(1),D(2),\ldots$
构成除幂环范畴中的一个余单纯对象。

\begin{lemma}
\label{crystalline-lemma-structure-Dn}
设 $D$ 与 $D(n)$ 如 (\ref{crystalline-equation-D}) 和 (\ref{crystalline-equation-Dn})。
余积典范映射 $P\to P\otimes_A\ldots\otimes_AP$，
$f\mapsto f\otimes1\otimes\ldots\otimes1$，诱导 $D$-代数同构
\begin{equation}
\label{crystalline-equation-structure-Dn}
D(n) = \lim_e D\langle \xi_i(j) \rangle/p^eD\langle \xi_i(j) \rangle
\end{equation}
其中
$$
\xi_i(j) = x_i \otimes 1 \otimes \ldots \otimes 1 -
1 \otimes \ldots \otimes 1 \otimes x_i \otimes 1 \otimes \ldots \otimes 1
$$
对 $j=1,\ldots,n$，第二个 $x_i$ 位于第 $j+1$ 个位置；回忆 $D(n)$
从 $P$ 在 $A$ 上的 $n+1$ 重张量积出发构造。
\end{lemma}

\begin{proof}
我们有
$$
P \otimes_A \ldots \otimes_A P = P[\xi_i(j)]
$$
而 $J(n)$ 由 $J$ 与各元素 $\xi_i(j)$ 生成。因此引理由
引理 \ref{crystalline-lemma-divided-power-envelope-add-variables} 得出。
\end{proof}

\begin{lemma}
\label{crystalline-lemma-property-Dn}
设 $D$ 与 $D(n)$ 如 (\ref{crystalline-equation-D}) 和 (\ref{crystalline-equation-Dn})。
则 $(D,\bar J,\bar\gamma)$ 与
$(D(n),\bar J(n),\bar\gamma(n))$ 是 $\text{Cris}^\wedge(C/A)$
的对象（见注 \ref{crystalline-remark-completed-affine-site}），并且
$$
D(n) = \coprod\nolimits_{j = 0, \ldots, n} D
$$
为 $\text{Cris}^\wedge(C/A)$ 中的余积。
\end{lemma}

\begin{proof}
第一个断言显然。对于第二个断言，若 $(B\to C,\delta)$ 是
$\text{Cris}^\wedge(C/A)$ 的对象，则
$$
\Mor_{\text{Cris}^\wedge(C/A)}(D, B) = 
\Hom_A((P, J), (B, \Ker(B \to C)))
$$
对 $D(n)$ 也有类似等式，其中以
$(P\otimes_A\ldots\otimes_AP,J(n))$ 代替 $(P,J)$。
由于 $P\otimes_A\ldots\otimes_AP$ 是余积，所述余积性质随即得出。
\end{proof}

\noindent
在下述引理中，我们考虑满足以下条件的偶对 $(M,\nabla)$：
\begin{enumerate}
\item
\label{crystalline-item-complete}
$M$ 是 $p$-进完备的 $D$-模；
\item
\label{crystalline-item-connection}
$\nabla:M\to M\otimes^\wedge_D\Omega_D$ 是联络，亦即
$\nabla(fm)=m\otimes\text df+f\nabla(m)$；
\item
\label{crystalline-item-integrable}
$\nabla$ 可积（见注 \ref{crystalline-remark-connection}）；
\item
\label{crystalline-item-topologically-quasi-nilpotent}
$\nabla$ {\it 拓扑拟幂零}：若对某些算子 $\theta_i:M\to M$ 写成
$\nabla(m)=\sum\theta_i(m)\text d x_i$，则对任意 $m\in M$，
只有有限多个偶对 $(i,k)$ 满足 $\theta_i^k(m)\notin pM$。
\end{enumerate}
文献中有时把算子 $\theta_i$ 记为 $\nabla_{\partial/\partial x_i}$。
下述引理构造从 $\text{Cris}(X/S)$ 上拟凝聚模中的晶体到这类偶对所成
范畴的函子。我们将在 命题 \ref{crystalline-proposition-crystals-on-affine}
中证明此函子是等价。

\begin{lemma}
\label{crystalline-lemma-crystals-on-affine}
在上述情形中，存在函子
$$
\begin{matrix}
\text{Cris}(X/S)\text{ 上拟凝聚} \\
\mathcal{O}_{X/S}\text{-模中的晶体}
\end{matrix}
\longrightarrow
\begin{matrix}
\text{满足 (\ref{crystalline-item-complete})、(\ref{crystalline-item-connection})、} \\
\text{(\ref{crystalline-item-integrable}) 及
(\ref{crystalline-item-topologically-quasi-nilpotent}) 的偶对 }(M,\nabla)
\end{matrix}
$$
\end{lemma}

\begin{proof}
设 $\mathcal F$ 是 $X/S$ 上拟凝聚模中的晶体。置
$T_e=\Spec(D_e)$，则当 $e\gg0$ 时，$(X,T_e,\bar\gamma)$ 是
$\text{Cris}(X/S)$ 的对象。我们有闭浸入
$$
(X, T_e, \bar\gamma) \to (X, T_{e + 1}, \bar\gamma) \to \ldots
$$
置
$$
M =
\lim_e \Gamma((X, T_e, \bar\gamma), \mathcal{F}) =
\lim_e \Gamma(T_e, \mathcal{F}_{T_e}) = \lim_e M_e
$$
由于 $\mathcal F$ 局部拟凝聚，有
$\mathcal F_{T_e}=\widetilde{M_e}$。由于 $\mathcal F$ 是晶体，有
$M_e=M_{e+1}/p^eM_{e+1}$。故 $M_e=M/p^eM$，且 $M$ 是 $p$-进完备的；
见 “代数” 引理 \ref{algebra-lemma-limit-complete}。

\medskip\noindent
由引理 \ref{crystalline-lemma-automatic-connection}，$\mathcal F$ 配备典范可积联络
$\nabla:\mathcal F\to\mathcal F\otimes\Omega_{X/S}$。
在上述对象 $T_e$ 上取此联络，得到典范可积联络
$$
\nabla : M \longrightarrow M \otimes^\wedge_D \Omega_D
$$
为看出它拓扑幂零，我们具体解释此条件。

\medskip\noindent
现在对各环 $D(n)$ 作同样处理，得到 $p$-进完备的 $D(n)$-模 $M(n)$。
再次利用 $\mathcal F$ 的晶体性质，得到同构
$$
M \otimes^\wedge_{D, p_0} D(1) \rightarrow M(1)
\leftarrow M \otimes^\wedge_{D, p_1} D(1)
$$
与 引理 \ref{crystalline-lemma-automatic-connection} 的证明比较。以 $c$
表示从左到右的复合。取 $m\in M$，并记
$\xi_i=x_i\otimes1-1\otimes x_i$。利用
(\ref{crystalline-equation-structure-Dn})，可唯一写成
$$
c(m \otimes 1) = \sum\nolimits_K \theta_K(m) \otimes \prod \xi_i^{[k_i]}
$$
其中 $\theta_K(m)\in M$，求和遍历满足 $k_i\geq0$ 且
$\sum k_i<\infty$ 的多重指标 $K=(k_i)$。当 $K$ 在第 $i$ 个位置为 $1$、
其余位置为零时，置 $\theta_i=\theta_K$。我们有
$$
\nabla(m) = \sum \theta_i(m) \text{d}x_i.
$$
这可由与 $\nabla$ 的定义比较看出。具体地，在 引理
\ref{crystalline-lemma-automatic-connection} 中定义等式为
$p_1^*m=\nabla(m)-c(p_0^*m)$；而符号仍相符，因为 Stacks Project
始终在模去对角理想平方后采用
$\text df=p_1(f)-p_0(f)$，所以
$\xi_i=x_i\otimes1-1\otimes x_i$ 模去对角理想平方后映到
$-\text d x_i$。

\medskip\noindent
以 $q_i:D\to D(2)$ 与 $q_{ij}:D(1)\to D(2)$ 表示对应于指标 $i,j$
的余积典范映射。与 引理 \ref{crystalline-lemma-automatic-connection}
证明的最后一段相同，可知
$$
q_{02}^*c = q_{12}^*c \circ q_{01}^*c.
$$
这意味着
$$
\sum\nolimits_{K''} \theta_{K''}(m) \otimes \prod {\zeta''_i}^{[k''_i]}
=
\sum\nolimits_{K', K} \theta_{K'}(\theta_K(m))
\otimes \prod {\zeta'_i}^{[k'_i]} \prod \zeta_i^{[k_i]}
$$
在 $M\otimes^\wedge_{D,q_2}D(2)$ 中成立，其中
\begin{align*}
\zeta_i & = x_i \otimes 1 \otimes 1 - 1 \otimes x_i \otimes 1,\\
\zeta'_i & = 1 \otimes x_i \otimes 1 - 1 \otimes 1 \otimes x_i,\\
\zeta''_i & = x_i \otimes 1 \otimes 1 - 1 \otimes 1 \otimes x_i.
\end{align*}
特别地，$\zeta''_i=\zeta_i+\zeta'_i$，并且由引理
\ref{crystalline-lemma-structure-Dn}，$D(2)$ 是 $q_2(D)$ 上以
$\zeta_i,\zeta'_i$ 为变量的除幂多项式环的 $p$-进完备化。
比较上式系数，立即得到
$\theta_i\circ\theta_j=\theta_j\circ\theta_i$
（这也给出 $\nabla$ 可积性的另一证明），以及
$$
\theta_K(m) = (\prod \theta_i^{k_i})(m).
$$
特别地，由于上述表示 $c(m\otimes1)$ 的和必须 $p$-进收敛，
可知对每个 $i$ 与每个 $m\in M$，只有有限多个 $\theta_i^k(m)$
模 $p$ 非零。
\end{proof}

\begin{proposition}
\label{crystalline-proposition-crystals-on-affine}
函子
$$
\begin{matrix}
\text{Cris}(X/S)\text{ 上拟凝聚} \\
\mathcal{O}_{X/S}\text{-模中的晶体}
\end{matrix}
\longrightarrow
\begin{matrix}
\text{满足 (\ref{crystalline-item-complete})、(\ref{crystalline-item-connection})、} \\
\text{(\ref{crystalline-item-integrable}) 及
(\ref{crystalline-item-topologically-quasi-nilpotent}) 的偶对 }(M,\nabla)
\end{matrix}
$$
是 引理 \ref{crystalline-lemma-crystals-on-affine} 中的范畴等价。
\end{proposition}

\begin{proof}
给定 $(M,\nabla)$。我们将构造拟凝聚模中的晶体 $\mathcal F$。
写成 $\nabla(m)=\sum\theta_i(m)\text d x_i$。则
$\theta_i\circ\theta_j=\theta_j\circ\theta_i$；对任意满足
$k_i\geq0$ 且 $\sum k_i<\infty$ 的多重指标 $K=(k_i)$，可置
$\theta_K(m)=(\prod\theta_i^{k_i})(m)$。

\medskip\noindent
设 $(U,T,\delta)$ 是 $\text{Cris}(X/S)$ 中 $T$ 仿射的任意对象。
写 $T=\Spec(B)$，并设 $U\to T$ 的理想为 $J_B\subset B$。
由引理 \ref{crystalline-lemma-set-generators}，存在整数 $e$ 及态射
$$
f : (U, T, \delta) \longrightarrow (X, T_e, \bar\gamma)
$$
其中 $T_e=\Spec(D_e)$，同引理 \ref{crystalline-lemma-crystals-on-affine} 的证明。
选取这样的 $e$ 与 $f$；相应的除幂 $A$-代数映射也记为 $f:D\to B$。
我们把 $\mathcal F_T$ 定义为与下列 $B$-模相伴的拟凝聚
$\mathcal O_T$-模层：
$$
M \otimes_{D, f} B.
$$
但必须证明它与 $f$ 的选取无关。设 $g:D\to B$ 是另一个这样的态射。
由于 $f$ 与 $g$ 是 $\text{Cris}(X/S)$ 中的态射，映射
$f-g:D\to B$ 的像包含在除幂理想 $J_B$ 中。记
$\xi_i=f(x_i)-g(x_i)\in J_B$。仿照引理
\ref{crystalline-lemma-crystals-on-affine} 的证明，由下式定义同构
$$
c_{f, g} : M \otimes_{D, f} B \longrightarrow M \otimes_{D, g} B
$$
$$
m \otimes 1 \longmapsto 
\sum\nolimits_K \theta_K(m) \otimes \prod \xi_i^{[k_i]}
$$
由上述讨论及 $\nabla$ 拓扑拟幂零（故该和有限！），此定义有意义。
计算表明
$$
c_{g, h} \circ c_{f, g} = c_{f, h}
$$
对任意第三个态射
$h:(U,T,\delta)\longrightarrow(X,T_e,\bar\gamma)$ 成立。
并且 $c_{f,f}=1$。因此这些映射都是同构，从而模 $\mathcal F_T$
与 $f$ 的选取无关。

\medskip\noindent
若 $a:(U',T',\delta')\to(U,T,\delta)$ 是 $\text{Cris}(X/S)$
中仿射对象之间的态射，则取 $f'=f\circ a$，显然存在典范同构
$a^*\mathcal F_T\to\mathcal F_{T'}$。略去此映射与 $f$ 的选取无关的验证。
以这些映射为限制映射，显然得到 $\text{Cris}(X/S)$ 中仿射对象所成
全子范畴上的一个拟凝聚模晶体。略去它延拓为整个
$\text{Cris}(X/S)$ 上晶体的证明。我们也略去以下证明：此过程为函子，
并且与 引理 \ref{crystalline-lemma-crystals-on-affine} 中构造的函子互为拟逆。
\end{proof}

\begin{lemma}
\label{crystalline-lemma-crystals-on-affine-smooth}
在 情形 \ref{crystalline-situation-affine} 中，设 $A\to P'\to C$ 是环映射，
$A\to P'$ 光滑，而 $P'\to C$ 满射且核为 $J'$。令 $D'$ 为
$D_{P',\gamma}(J')$ 的 $p$-进完备化。则存在如上的数据
$A\to P\to C$ 的一种选取，以及除幂 $A$-代数同态
$$
a : D \longrightarrow D',\quad b : D' \longrightarrow D
$$
它们与映射 $D\to C$ 及 $D'\to C$ 相容，并且
$a\circ b=\text{id}_{D'}$。这些映射诱导以下两类偶对所成范畴之间的等价：
在 $D$ 上满足
(\ref{crystalline-item-complete})、(\ref{crystalline-item-connection})、
(\ref{crystalline-item-integrable}) 与 (\ref{crystalline-item-topologically-quasi-nilpotent})
的偶对 $(M,\nabla)$，以及在 $D'$ 上满足
(\ref{crystalline-item-complete})、(\ref{crystalline-item-connection})、
(\ref{crystalline-item-integrable}) 与
(\ref{crystalline-item-topologically-quasi-nilpotent})\footnote{对 $D'$ 上的
$(M',\nabla')$ 表述此条件颇为棘手，见证明。} 的偶对 $(M',\nabla')$。
特别地，命题 \ref{crystalline-proposition-crystals-on-affine} 的范畴等价
对取值于 $D'$ 上偶对的相应函子也成立。
\end{lemma}

\begin{proof}
选取 $A$ 上多项式代数 $P=A[y_1,\ldots,y_m]$ 及满射 $P\to P'$，
并以复合 $P\to P'\to C$ 定义 $P\to C$。由除幂包络与完备化的函子性，
得到满射 $a:D\to D'$。取足够大的 $e$，使 $D_e$ 是 $C$ 在 $A$
上的除幂加厚。于是 $D_e\to C$ 是核局部幂零的满射；见
“除幂代数”引理 \ref{dpa-lemma-nil}。置
$D'_e=D'/p^eD'$，则 $D_e\to D'_e$ 的核局部幂零。因此由
“代数” 引理 \ref{algebra-lemma-smooth-strong-lift}，可找到映射
$P'\to D'_e$ 的提升 $\beta_e:P'\to D_e$。注意，对任意 $i\geq0$，
$D_{e+i+1}\to D_{e+i}\times_{D'_{e+i}}D'_{e+i+1}$
是核平方为零的满射，因为 $p^{e+i}D\to p^{e+i}D'$ 满射。
依次把通常的提升性质（“代数” 命题
\ref{algebra-proposition-smooth-formally-smooth}）应用于各图
$$
\xymatrix{
P' \ar[r] & D_{e + i} \times_{D'_{e + i}} D'_{e + i + 1} \\
A \ar[u] \ar[r] & D_{e + i + 1} \ar[u]
}
$$
可知存在 $A$-代数映射 $\beta:P'\to D$，它与 $a$ 的复合为给定映射
$P'\to D'$。由除幂包络的泛性质得到映射
$D_{P',\gamma}(J')\to D$。由于 $D$ 是 $p$-进完备的，得到
$b:D'\to D$，满足 $a\circ b=\text{id}_{D'}$。

\medskip\noindent
考虑基变换函子
$$
F : (M, \nabla) \longmapsto
(M \otimes^\wedge_{D, a} D', \nabla')
\quad\text{且}\quad
G : (M', \nabla') \longmapsto
(M' \otimes^\wedge_{D', b} D, \nabla)
$$
它们作用于满足 (\ref{crystalline-item-complete})、(\ref{crystalline-item-connection}) 与
(\ref{crystalline-item-integrable}) 的带联络模；见注
\ref{crystalline-remark-base-change-connection}。由于
$a\circ b=\text{id}_{D'}$，$F\circ G$ 是恒等函子。若
$G(M',\nabla')$ 满足性质
(\ref{crystalline-item-topologically-quasi-nilpotent})，我们便称
$(M',\nabla')$ 也满足该性质。形式论证表明，为完成证明，只需在
$(M,\nabla)$ 满足四个条件 (\ref{crystalline-item-complete})、
(\ref{crystalline-item-connection})、(\ref{crystalline-item-integrable}) 与
(\ref{crystalline-item-topologically-quasi-nilpotent}) 时，证明
$G(F(M,\nabla))$ 同构于 $(M,\nabla)$。为此使用 命题
\ref{crystalline-proposition-crystals-on-affine} 证明中的函子性同构
$$
c_{\text{id}_D, b \circ a} :
M \otimes_{D, \text{id}_D} D
\longrightarrow
M \otimes_{D, b \circ a} D
$$
（它要求 $\nabla$ 拓扑拟幂零，而这已在假设中）。
最后还需证明此映射是水平的，即与联络相容；此验证从略。

\medskip\noindent
证明中的最后一个断言随即得出。
\end{proof}

\begin{remark}
\label{crystalline-remark-equivalence-more-general}
命题 \ref{crystalline-proposition-crystals-on-affine} 的等价在下述情形仍成立：
从满射 $P\to C$ 出发，其中 $P/A$ 满足 “代数” 引理
\ref{algebra-lemma-smooth-strong-lift} 的强提升性质。其证明可仿照
引理 \ref{crystalline-lemma-crystals-on-affine-smooth} 的证明。
（若以后需要，将在此补充细节。）此结果大概也有直接证明；不过采用
多项式环的优点在于，各环 $D(n)$ 是除幂多项式环的 $p$-进完备化，
从而简化代数运算。
\end{remark}



\section{上同调的一般评注}
\label{crystalline-section-cohomology-lqc}

\noindent
本节作一些准备，把仿射概形晶体位点上模的上同调转化为代数问题。

\begin{lemma}
\label{crystalline-lemma-vanishing-lqc}
在 情形 \ref{crystalline-situation-global} 中，设 $\mathcal F$ 是
$\text{Cris}(X/S)$ 上局部拟凝聚的 $\mathcal O_{X/S}$-模。则
$$
H^p((U, T, \delta), \mathcal{F}) = 0
$$
对所有 $p>0$ 以及所有 $T$ 或 $U$ 仿射的 $(U,T,\delta)$ 成立。
\end{lemma}

\begin{proof}
由于 $U\to T$ 是加厚，$U$ 仿射当且仅当 $T$ 仿射；见“极限”
引理 \ref{limits-lemma-affine}。把“位点上同调”引理
\ref{sites-cohomology-lemma-cech-vanish-collection} 应用于仿射对象
$(U,T,\delta)$ 所成的族 $\mathcal B$，以及仿射开覆盖
$\mathcal U=\{(U_i,T_i,\delta_i)\to(U,T,\delta)\}$ 所成的族
$\text{Cov}$。这种覆盖的 {\v C}ech 复形
${\check C}^*(\mathcal U,\mathcal F)$，恰是拟凝聚
$\mathcal O_T$-模 $\mathcal F_T$ 关于仿射概形 $T$ 的仿射开覆盖
$\{T_i\to T\}$ 的 {\v C}ech 复形（这里使用了 $\mathcal F$
局部拟凝聚的假设）。因此，由“概形上同调”引理
\ref{coherent-lemma-cech-cohomology-quasi-coherent} 与
\ref{coherent-lemma-quasi-coherent-affine-cohomology-zero}，
其 {\v C}ech 上同调为零。于是“位点上同调”引理
\ref{sites-cohomology-lemma-cech-vanish-collection} 的假设均成立，结论得证。
\end{proof}

\begin{lemma}
\label{crystalline-lemma-compare}
在情形 \ref{crystalline-situation-global} 中，进一步假设 $X$ 与 $S$
都是仿射概形。考虑 $\text{Cris}(X/S)$ 中由除幂加厚
$(X,T,\delta)$ 所成的全子范畴 $\mathcal C$，并赋予混沌拓扑
（见 “位点论” 例 \ref{sites-example-indiscrete}）。对任意局部拟凝聚的
$\mathcal O_{X/S}$-模 $\mathcal F$，有
$$
R\Gamma(\mathcal{C}, \mathcal{F}|_\mathcal{C}) =
R\Gamma(\text{Cris}(X/S), \mathcal{F})
$$
\end{lemma}

\begin{proof}
以 $\text{AffineCris}(X/S)$ 表示 $\text{Cris}(X/S)$ 中由 $U$ 与 $T$
均仿射的对象 $(U,T,\delta)$ 所成的全子范畴。规定
$\text{AffineCris}(X/S)$ 的态射族
$\{(U_i,T_i,\delta_i)\to(U,T,\delta)\}_{i\in I}$ 是覆盖，当且仅当它在
$\text{Cris}(X/S)$ 中是覆盖，由此把该范畴变成位点。在此定义下，包含函子
$$
\text{AffineCris}(X/S) \longrightarrow \text{Cris}(X/S)
$$
是 “位点论” 定义 \ref{sites-definition-special-cocontinuous-functor}
意义下的特殊余连续函子。证明与 “拓扑” 引理
\ref{topologies-lemma-affine-big-site-Zariski} 的证明完全相同。
因此，沿上述包含函子限制后，$\text{Cris}(X/S)$ 上的层拓扑斯与
$\text{AffineCris}(X/S)$ 上的层拓扑斯相同。于是只需对包含
$\mathcal C\subset\text{AffineCris}(X/S)$ 证明相应断言。

\medskip\noindent
下文不再说明地使用 $\mathcal C$ 与 $\text{AffineCris}(X/S)$
具有积和纤维积这一事实（细节略；见引理
\ref{crystalline-lemma-divided-power-thickening-fibre-products}）。包含函子
$u:\mathcal C\to\text{AffineCris}(X/S)$ 全忠实、连续，并与积和纤维积交换。
我们断言它定义环化位点态射
$$
f :
(\text{AffineCris}(X/S), \mathcal{O}_{X/S})
\longrightarrow
(\Sh(\mathcal{C}), \mathcal{O}_{X/S}|_\mathcal{C})
$$
为此使用 “位点论” 引理 \ref{sites-lemma-directed-morphism}。由于
$\mathcal C$ 具有纤维积且 $u$ 与之交换，各范畴
$\mathcal I^u_{(U,T,\delta)}$ 是有向范畴的不交并（由 “位点论” 引理
\ref{sites-lemma-almost-directed} 与 “范畴” 引理
\ref{categories-lemma-split-into-directed}）。因此只需证明
$\mathcal I^u_{(U,T,\delta)}$ 连通。非空性由引理
\ref{crystalline-lemma-set-generators} 得出：由于 $U$ 与 $T$ 仿射，该引理说明
$\mathcal C$ 中至少有一个对象 $(X,T',\delta')$ 及一个除幂加厚态射
$(U,T,\delta)\to(X,T',\delta')$。连通性由 $\mathcal C$ 具有积且
$u$ 与积交换得出（与 “位点论” 引理 \ref{sites-lemma-directed}
的证明比较）。

\medskip\noindent
注意 $f_*\mathcal F=\mathcal F|_\mathcal C$。因此，若对 $p>0$ 有
$R^pf_*\mathcal F=0$，引理即得证；见“位点上同调”引理
\ref{sites-cohomology-lemma-apply-Leray}。由“位点上同调”引理
\ref{sites-cohomology-lemma-higher-direct-images}，只需对所有
$(X,T,\delta)$ 证明
$H^p(\text{AffineCris}(X/S)/(X,T,\delta),\mathcal F)=0$。
这由引理 \ref{crystalline-lemma-vanishing-lqc} 得出，因为位点
$\text{AffineCris}(X/S)/(X,T,\delta)$ 的拓扑斯等价于该引理中使用的
位点 $\text{Cris}(X/S)/(X,T,\delta)$ 的拓扑斯。
\end{proof}

\begin{lemma}
\label{crystalline-lemma-complete}
在 情形 \ref{crystalline-situation-affine} 中，置
$\mathcal C=(\text{Cris}(C/A))^{opp}$ 与
$\mathcal C^\wedge=(\text{Cris}^\wedge(C/A))^{opp}$，并赋予混沌拓扑；
记号见注 \ref{crystalline-remark-completed-affine-site}。存在拓扑斯态射
$$
g : \Sh(\mathcal{C}) \longrightarrow \Sh(\mathcal{C}^\wedge)
$$
使得若 $\mathcal F$ 是 $\mathcal C$ 上的交换群层，则
$$
R^pg_*\mathcal{F}(B \to C, \delta) =
\left\{
\begin{matrix}
\lim_e \mathcal{F}(B_e \to C, \delta) & \text{若 }p = 0 \\
R^1\lim_e \mathcal{F}(B_e \to C, \delta) & \text{若 }p = 1 \\
0 & \text{其他情形}
\end{matrix}
\right.
$$
其中当 $e\gg0$ 时，$B_e=B/p^eB$。
\end{lemma}

\begin{proof}
范畴之间的任意函子都在同方向定义混沌拓扑斯之间的态射；例如，可把这种
函子看作位点之间的余连续函子，见 “位点论” 节
\ref{sites-section-cocontinuous-morphism-topoi}。略去对 $g_*\mathcal F$
之描述的证明。注意，命题中只有当 $e$ 足够大时，
$(B_e\to C,\delta)$ 才是 $\text{Cris}(C/A)$ 的对象。
设 $\mathcal I$ 是 $\mathcal C$ 上的内射交换群层。则转移映射
$$
\mathcal{I}(B_e \to C, \delta) \leftarrow
\mathcal{I}(B_{e + 1} \to C, \delta)
$$
满射，因为态射
$$
(B_e \to C, \delta)
\longrightarrow
(B_{e + 1} \to C, \delta)
$$
在范畴 $\mathcal C$ 中是单态。因此，对内射交换群层，引理所示公式两边
相同。取 $\mathcal F$ 的一个内射消解即可得到结论（层是预层，故正合性
在各对象上截面群的层面检验）。
\end{proof}

\begin{lemma}
\label{crystalline-lemma-category-with-covering}
设 $\mathcal C$ 是赋予混沌拓扑的范畴。设 $X$ 是 $\mathcal C$ 的对象，
并且 $\mathcal C$ 的每个对象都有一个指向 $X$ 的态射。假设
$\mathcal C$ 具有二元积。则对 $\mathcal C$ 上每个交换群层
$\mathcal F$，总上同调 $R\Gamma(\mathcal C,\mathcal F)$ 由复形
$$
\mathcal{F}(X) \to \mathcal{F}(X \times X) \to
\mathcal{F}(X \times X \times X) \to \ldots
$$
表示；该复形与余单纯交换群 $[n]\mapsto\mathcal F(X^n)$ 相伴。
\end{lemma}

\begin{proof}
由于任意预层在 $\mathcal C$ 上都是层，对所有 $q>0$ 有
$H^q(X^p,\mathcal F)=0$。关于 $X$ 的假设是 $h_X\to*$ 满射。
利用 $H^q(X,\mathcal F)=H^q(h_X,\mathcal F)$ 与
$H^q(\mathcal C,\mathcal F)=H^q(*,\mathcal F)$，可见该断言是
“位点上同调”引理
\ref{sites-cohomology-lemma-cech-to-cohomology-sheaf-sets} 的特例。
\end{proof}









\section{余单纯准备}
\label{crystalline-section-cohomology}

\noindent
本节比较晶体上同调与德拉姆上同调。我们遵循 \cite{Bhatt}。

\begin{example}
\label{crystalline-example-cosimplicial-module}
设 $A_*$ 是任意余单纯环。考虑由下式定义的余单纯模 $M_*$：
$$
M_n = \bigoplus\nolimits_{i = 0, ..., n} A_n e_i
$$
对映射 $f:[n]\to[m]$，把 $M_*(f):M_n\to M_m$ 定义为将 $e_i$
映到 $e_{f(i)}$ 的唯一 $A_*(f)$-线性映射。我们断言 $M_*$
上的恒等映射同伦于 $0$。事实上，一个同伦由余单纯模映射
$$
h : M_* \longrightarrow \Hom(\Delta[1], M_*)
$$
给出；见第 \ref{crystalline-section-cosimplicial} 节。对
$j\in\{0,\ldots,n+1\}$，令 $\alpha^n_j:[n]\to[1]$ 为由
$\alpha^n_j(i)=0\Leftrightarrow i<j$ 定义的映射。于是
$\Delta[1]_n=\{\alpha^n_0,\ldots,\alpha^n_{n+1}\}$，相应地
$\Hom(\Delta[1],M_*)_n=\prod_{j=0,\ldots,n+1}M_n$；见“单纯方法”第
\ref{simplicial-section-homotopy} 节与第
\ref{simplicial-section-homotopy-cosimplicial} 节。我们不用这个积表示，
而把 $\Hom(\Delta[1],M_*)_n$ 的元素看作函数
$\Delta[1]_n\to M_n$。用此记号，把 $h$ 的次数 $n$ 分量定义为
$$
h_n(e_i)(\alpha^n_j) =
\left\{
\begin{matrix}
e_{i} & \text{若} & i < j \\
0 & \text{其他情形} 
\end{matrix}
\right.
$$
先验证 $h$ 是余单纯模态射。即对 $f:[n]\to[m]$，证明
\begin{equation}
\label{crystalline-equation-cosimplicial-morphism}
h_m \circ M_*(f) = \Hom(\Delta[1], M_*)(f) \circ h_n
\end{equation}
(\ref{crystalline-equation-cosimplicial-morphism}) 左端先在 $e_i$ 处取值，
再在 $\alpha^m_j$ 处取值，所得为
$$
h_m(e_{f(i)})(\alpha^m_j) =
\left\{
\begin{matrix}
e_{f(i)} & \text{若} & f(i) < j \\
0 & \text{其他情形}
\end{matrix}
\right.
$$
注意 $\alpha^m_j\circ f=\alpha^n_{j'}$，其中
$0\leq j'\leq n+1$ 是使 $f(i)<j$ 当且仅当 $i<j'$ 的唯一指标。
因此，(\ref{crystalline-equation-cosimplicial-morphism}) 右端先在 $e_i$ 处取值，
再在 $\alpha^m_j$ 处取值，所得为
$$
M_*(f)(h_n(e_i)(\alpha^m_j \circ f) =
M_*(f)(h_n(e_i)(\alpha^n_{j'})) =
\left\{
\begin{matrix}
e_{f(i)} & \text{若} & i < j' \\
0 & \text{其他情形} 
\end{matrix}
\right.
$$
由对 $j'$ 的描述可知两结果相等。因此 $h$ 是余单纯模映射。
令 $0:\Delta[0]\to\Delta[1]$ 与 $1:\Delta[0]\to\Delta[1]$
为显然映射，并以
$ev_0,ev_1:\Hom(\Delta[1],M_*)\to M_*$ 表示相应取值映射。
读者容易验证，复合
$$
ev_0 \circ h, ev_1 \circ h : M_* \longrightarrow M_*
$$
分别为 $1$ 与 $0$，故 $h$ 是所需的 $1$ 与 $0$ 之间的同伦。
\end{example}

\begin{lemma}
\label{crystalline-lemma-vanishing-omega-1}
采用 (\ref{crystalline-equation-omega-Dn}) 中的记号，复形
$$
\Omega_{D(0)} \to \Omega_{D(1)} \to \Omega_{D(2)} \to \ldots
$$
作为余单纯 $D(*)$-模同伦于零。
\end{lemma}

\begin{proof}
我们将使用 “单纯方法” 引理
\ref{simplicial-lemma-functorial-homotopy} 的原则，更具体地说，
使用 引理 \ref{crystalline-lemma-homotopy-tensor}：任意函子都把（余）单纯对象之间
的同伦映射变为同伦映射。引理中的复形等于余单纯模
$$
M_* = \left(
\Omega_{P/A} \to
\Omega_{P \otimes_A P/A} \to
\Omega_{P \otimes_A P \otimes_A P/A} \to \ldots
\right)
$$
沿余单纯环映射 $P\otimes_A\ldots\otimes_AP\to D(n)$ 基变换后再作
$p$-进完备化；这由引理
\ref{crystalline-lemma-module-differentials-divided-power-envelope} 得出，
见 (\ref{crystalline-equation-omega-D}) 后的评注。因此，只需证明余单纯模 $M_*$
同伦于零（这里使用基变换与 $p$-进完备化）。利用
$\mathbf Z\to A$ 的基变换，甚至可假设 $A=\mathbf Z$ 且
$P=\mathbf Z[\{x_i\}_{i\in I}]$。此时 $P^{\otimes n+1}$ 是以各元素
$$
x_i(e) = 1 \otimes \ldots \otimes x_i \otimes \ldots \otimes 1
$$
为变量的多项式代数，其中 $x_i$ 位于第 $e$ 个位置。复形的各模以
$\text d x_i(e)$ 为自由生成元。注意，若 $f:[n]\to[m]$，则
$$
M_*(f)(\text{d}x_i(e)) = \text{d}x_i(f(e))
$$
因此，$M_*$ 是按 $I$ 求和的、例
\ref{crystalline-example-cosimplicial-module} 中所研究模的副本之直和，结论得证。
\end{proof}

\begin{lemma}
\label{crystalline-lemma-vanishing}
采用 (\ref{crystalline-equation-Dn}) 与 (\ref{crystalline-equation-omega-Dn}) 中的记号，
给定任意余单纯 $D(*)$-模 $M_*$ 及 $i>0$，余单纯模
$$
M_0 \otimes^\wedge_{D(0)} \Omega^i_{D(0)} \to
M_1 \otimes^\wedge_{D(1)} \Omega^i_{D(1)} \to
M_2 \otimes^\wedge_{D(2)} \Omega^i_{D(2)} \to \ldots
$$
同伦于零，其中 $\Omega^i_{D(n)}$ 是 $\Omega_{D(n)}$
的第 $i$ 个外幂的 $p$-进完备化。
\end{lemma}

\begin{proof}
由引理 \ref{crystalline-lemma-vanishing-omega-1}，$\Omega_{D(*)}$ 的自同态
$0$ 与 $1$ 同伦。应用函子 $\wedge^i$，可知余单纯模
$\wedge^i\Omega_{D(*)}$ 也如此；见引理
\ref{crystalline-lemma-homotopy-tensor}。再次应用同一引理，得到其 $p$-进完备化
$\Omega^i_{D(*)}$ 同伦等价于零。与 $M_*$ 作张量积，得到
$M_*\otimes_{D(*)}\Omega^i_{D(*)}$ 同伦于零；再次见引理
\ref{crystalline-lemma-homotopy-tensor}。最后应用 $p$-进完备化函子即完成证明。
\end{proof}



\section{除幂 Poincar\'e 引理}
\label{crystalline-section-poincare}

\noindent
这里只给出最简单的版本。

\begin{lemma}
\label{crystalline-lemma-poincare}
设 $A$ 是环，$P=A\langle x_i\rangle$ 是 $A$ 上的除幂多项式环。
对任意 $A$-模 $M$，复形
$$
0 \to M \to
M \otimes_A P \to
M \otimes_A \Omega^1_{P/A, \delta} \to
M \otimes_A \Omega^2_{P/A, \delta} \to \ldots
$$
正合。令 $D$ 为 $P$ 的 $p$-进完备化，并令 $\Omega^i_D$ 为
$\Omega_{D/A,\delta}$ 第 $i$ 个外幂的 $p$-进完备化。对任意
$p$-进完备的 $A$-模 $M$，复形
$$
0 \to M \to
M \otimes^\wedge_A D \to
M \otimes^\wedge_A \Omega^1_D \to
M \otimes^\wedge_A \Omega^2_D \to \ldots
$$
正合。
\end{lemma}

\begin{proof}
只需证明 $A$-模复形
$$
E :
(0 \to A \to P \to \Omega^1_{P/A, \delta} \to
\Omega^2_{P/A, \delta} \to \ldots)
$$
同伦等价于零。对每个多重指标 $K=(k_i)$，考虑子复形 $E(K)$，
其次数 $j$ 分量为
$$
\bigoplus\nolimits_{I = \{i_1, \ldots, i_j\} \subset \text{Supp}(K)}
A
\prod\nolimits_{i \not \in I} x_i^{[k_i]}
\prod\nolimits_{i \in I} x_i^{[k_i - 1]}
\text{d}x_{i_1} \wedge \ldots \wedge \text{d}x_{i_j}
$$
由于 $E=\bigoplus E(K)$，只需证明每个复形 $E(K)$ 同伦于零。
若 $K=0$，则 $E(K):(A\to A)$ 同伦于零。若 $K$ 的支集
$S$ 非空且有限，则复形 $E(K)$ 同构于
$$
0 \to A \to
\bigoplus\nolimits_{s \in S} A \to
\wedge^2(\bigoplus\nolimits_{s \in S} A) \to
\ldots \to \wedge^{\# S}(\bigoplus\nolimits_{s \in S} A) \to 0
$$
后者同伦于零，例如由“代数进阶”引理
\ref{more-algebra-lemma-homotopy-koszul-abstract} 可知。
\end{proof}

\noindent
例 \ref{crystalline-example-integrate} 解释了下述引理的另一种（更直接的）方法。

\begin{lemma}
\label{crystalline-lemma-relative-poincare}
设 $A$ 是环，$(B,I,\delta)$ 是除幂环，其中 $B$ 是 $A$-代数。
令 $P=B\langle x_i\rangle$ 为 $B$ 上除幂多项式环，照例配备除幂理想
$J=IP+B\langle x_i\rangle_+$。设 $M$ 是配备可积联络
$\nabla:M\to M\otimes_B\Omega^1_{B/A,\delta}$ 的 $B$-模。
则德拉姆复形的映射
$$
M \otimes_B \Omega^*_{B/A, \delta}
\longrightarrow
M \otimes_P \Omega^*_{P/A, \delta}
$$
是拟同构。令 $D$、相应地 $D'$ 分别为 $B$、相应地 $P$ 的
$p$-进完备化，并令 $\Omega^i_D$、相应地 $\Omega^i_{D'}$ 分别为
$\Omega^i_{B/A,\delta}$、相应地 $\Omega^i_{P/A,\delta}$ 的
$p$-进完备化。设 $M$ 是配备整联络
$\nabla:M\to M\otimes^\wedge_D\Omega^1_D$ 的 $p$-进完备 $D$-模。
则德拉姆复形的映射
$$
M \otimes^\wedge_D \Omega^*_D
\longrightarrow
M \otimes^\wedge_D \Omega^*_{D'}
$$
是拟同构。
\end{lemma}

\begin{proof}
考虑 $\Omega^*_{B/A,\delta}$ 上由子复形
$F^i(\Omega^*_{B/A,\delta})=\sigma_{\geq i}\Omega^*_{B/A,\delta}$
给出的递减滤过 $F^*$；见 “同调” 节
\ref{homology-section-truncations}。由下式，它在
$\Omega^*_{P/A,\delta}$ 上诱导递减滤过 $F^*$：
$$
F^i(\Omega^*_{P/A, \delta}) =
F^i(\Omega^*_{B/A, \delta}) \wedge \Omega^*_{P/A, \delta}.
$$
我们有可分裂短正合列
$$
0 \to \Omega^1_{B/A, \delta} \otimes_B P \to
\Omega^1_{P/A, \delta} \to
\Omega^1_{P/B, \delta} \to 0
$$
其中最后一个模以 $\text d x_i$ 为自由基。由此可知
$F^i(\Omega^*_{P/A,\delta})\to\Omega^*_{P/A,\delta}$
逐项为可分裂单射，并且
$$
\text{gr}^i_F(\Omega^*_{P/A, \delta}) =
\Omega^i_{B/A, \delta} \otimes_B \Omega^*_{P/B, \delta}
$$
为复形的等式。因此可由下式在
$M\otimes_B\Omega^*_{B/A,\delta}$ 上定义滤过 $F^*$：
$$
F^i(M \otimes_B \Omega^*_{P/A, \delta}) =
M \otimes_B F^i(\Omega^*_{P/A, \delta})
$$
并且复形满足
$$
\text{gr}^i_F(M \otimes_B \Omega^*_{P/A, \delta}) =
M \otimes_B \Omega^i_{B/A, \delta} \otimes_B \Omega^*_{P/B, \delta}
$$
由引理 \ref{crystalline-lemma-poincare}，这些复形中的每一个都拟同构于置于次数
$0$ 的 $M\otimes_B\Omega^i_{B/A,\delta}$。因此，引理中第一个所示映射
是滤过复形的态射，并在分次片上诱导拟同构。这蕴含它自身是拟同构；
例如可使用与滤过复形相伴的谱序列，见 “同调” 节
\ref{homology-section-filtered-complex}。

\medskip\noindent
第二个拟同构的证明完全相同。
\end{proof}



\section{仿射情形的上同调}
\label{crystalline-section-cohomology-affine}

\noindent
回到第 \ref{crystalline-section-quasi-coherent-crystals} 节中研究的情形。
从 $(A,I,\gamma)$ 与 $A/I\to C$ 出发，置 $X=\Spec(C)$ 与
$S=\Spec(A)$。然后选取 $A$ 上多项式环 $P$ 及核为 $J$ 的满射
$P\to C$。由 (\ref{crystalline-equation-D}) 与 (\ref{crystalline-equation-Dn}) 得到
$D$ 和 $D(n)$。置 $T(n)_e=\Spec(D(n)/p^eD(n))$，使
$(X,T(n)_e,\delta(n))$ 是 $\text{Cris}(X/S)$ 的对象。
设 $\mathcal F$ 是 $\mathcal O_{X/S}$-模层，并对
$n=0,1,2,3,\ldots$ 置
$$
M(n) = \lim_e \Gamma((X, T(n)_e, \delta(n)), \mathcal{F})
$$
这构成余单纯环 $D(0),D(1),D(2),\ldots$ 上的余单纯模。

\begin{proposition}
\label{crystalline-proposition-compute-cohomology}
采用上述记号，假设：
\begin{enumerate}
\item $\mathcal F$ 局部拟凝聚；
\item 对 $\text{Cris}(X/S)$ 中任意态射
$(U,T,\delta)\to(U',T',\delta')$，若 $f:T\to T'$ 是闭浸入，
则映射 $c_f:f^*\mathcal F_{T'}\to\mathcal F_T$ 满射。
\end{enumerate}
则复形
$$
M(0) \to M(1) \to M(2) \to \ldots
$$
计算 $R\Gamma(\text{Cris}(X/S),\mathcal F)$。
\end{proposition}

\begin{proof}
利用假设 (1) 与 引理 \ref{crystalline-lemma-compare}，可知
$R\Gamma(\text{Cris}(X/S),\mathcal F)$ 同构于
$R\Gamma(\mathcal C,\mathcal F)$。注意，引理 \ref{crystalline-lemma-compare}
与 \ref{crystalline-lemma-complete} 中使用的范畴 $\mathcal C$ 相同。
设 $f:T\to T'$ 是 (2) 中的闭浸入。映射
$c_f:f^*\mathcal F_{T'}\to\mathcal F_T$ 满射，当且仅当
$\mathcal F_{T'}\to f_*\mathcal F_T$ 满射。因此，若 $\mathcal F$
满足 (1) 与 (2)，便得到 $T'$ 上拟凝聚 $\mathcal O_{T'}$-模的短正合列
$$
0 \to \mathcal{K} \to \mathcal{F}_{T'} \to f_*\mathcal{F}_T \to 0
$$
见“概形”第 \ref{schemes-section-quasi-coherent} 节，尤其是
引理 \ref{schemes-lemma-push-forward-quasi-coherent}。因此，若 $T'$ 仿射，
则由 $H^1(T',\mathcal K)$ 消失可知限制映射
$\mathcal F(U',T',\delta')\to\mathcal F(U,T,\delta)$ 满射；见
“概形上同调”引理
\ref{coherent-lemma-quasi-coherent-affine-cohomology-zero}。
于是 引理 \ref{crystalline-lemma-complete} 中各逆系统的转移映射满射。
由此得到对所有 $p\geq1$，
$R^pg_*(\mathcal F|_\mathcal C)=0$，其中 $g$ 如 引理
\ref{crystalline-lemma-complete}。范畴 $\mathcal C^\wedge$ 的对象 $D$ 满足 引理
\ref{crystalline-lemma-category-with-covering} 的假设。这由引理
\ref{crystalline-lemma-generator-completion} 及下式可知：
$$
D \times \ldots \times D = D(n)
$$
在 $\mathcal C$ 中成立，因为 $D(n)$ 是
$\text{Cris}^\wedge(C/A)$ 中 $D$ 的 $n+1$ 重余积；见
引理 \ref{crystalline-lemma-property-Dn}。结论得证。
\end{proof}

\begin{lemma}
\label{crystalline-lemma-cohomology-is-zero}
假设与记号同命题 \ref{crystalline-proposition-compute-cohomology}。则
$$
H^j(\text{Cris}(X/S), \mathcal{F} \otimes_{\mathcal{O}_{X/S}} \Omega^i_{X/S})
= 0
$$
对所有 $i>0$ 与所有 $j\geq0$ 成立。
\end{lemma}

\begin{proof}
由引理 \ref{crystalline-lemma-omega-locally-quasi-coherent}，
$\mathcal H=\mathcal F\otimes_{\mathcal O_{X/S}}\Omega^i_{X/S}$
也满足命题 \ref{crystalline-proposition-compute-cohomology} 的假设 (1) 与 (2)。记
$$
M(n)_e=\Gamma((X,T(n)_e,\delta(n)),\mathcal F)
$$
于是 $M(n)=\lim_eM(n)_e$。则
\begin{align*}
\lim_e \Gamma((X, T(n)_e, \delta(n)), \mathcal{H}) & =
\lim_e M(n)_e \otimes_{D(n)_e} \Omega_{D(n)}/p^e\Omega_{D(n)} \\
& = \lim_e M(n)_e \otimes_{D(n)} \Omega_{D(n)}
\end{align*}
由引理 \ref{crystalline-lemma-vanishing}，余单纯模
$$
M(0)_e \otimes_{D(0)} \Omega^i_{D(0)} \to
M(1)_e \otimes_{D(1)} \Omega^i_{D(1)} \to
M(2)_e \otimes_{D(2)} \Omega^i_{D(2)} \to \ldots
$$
同伦于零。由于转移映射 $M(n)_{e+1}\to M(n)_e$ 满射，
相伴复形的逆极限非周期\footnote{事实上，因为各同伦彼此拼合，
它们甚至同伦于零，不过这里不需要这一点。之所以采用这种迂回论证，
是因为在引理的假设下，我们不知道
$M(n)_e=M(n)_{e+1}/p^eM(n)_{e+1}$，所以极限
$\lim_eM(n)_e\otimes_{D(n)}\Omega^i_{D(n)}$ 并不是
$M(n)\otimes_{D(n)}\Omega^i_{D(n)}$ 的 $p$-进完备化。
若 $\mathcal F$ 是晶体，该等式则成立。}。
于是由命题 \ref{crystalline-proposition-compute-cohomology}，
$\mathcal H$ 的上同调消失。
\end{proof}

\begin{proposition}
\label{crystalline-proposition-compute-cohomology-crystal}
假设同命题 \ref{crystalline-proposition-compute-cohomology}，但现在假设
$\mathcal F$ 是拟凝聚模中的晶体。令 $(M,\nabla)$ 为相应的 $D$
上带联络模；见命题 \ref{crystalline-proposition-crystals-on-affine}。
则复形
$$
M \otimes^\wedge_D \Omega^*_D
$$
计算 $R\Gamma(\text{Cris}(X/S),\mathcal F)$。
\end{proposition}

\begin{proof}
我们使用与双复形 $K^{*,*}$ 相伴的两个谱序列来证明，其中
$$
K^{a, b} = M \otimes_D^\wedge \Omega^a_{D(b)}
$$
已有结果如下。引理 \ref{crystalline-lemma-vanishing} 表明每一列
$K^{a,*}$（$a>0$）非周期。命题
\ref{crystalline-proposition-compute-cohomology} 表明第一列 $K^{0,*}$ 拟同构于
$R\Gamma(\text{Cris}(X/S),\mathcal F)$。因此，与双复形相伴的第一个
谱序列表明，$R\Gamma(\text{Cris}(X/S),\mathcal F)$ 与
$\text{Tot}(K^{*,*})$ 之间存在典范拟同构。

\medskip\noindent
接着考虑各行 $K^{*,b}$。由引理 \ref{crystalline-lemma-structure-Dn}，
$b+1$ 个映射 $D\to D(b)$ 中每一个都把 $D(b)$ 表示为 $D$
上某个除幂多项式代数的 $p$-进完备化。因此引理
\ref{crystalline-lemma-relative-poincare} 表明映射
$$
M \otimes^\wedge_D\Omega^*_D
\longrightarrow
M \otimes^\wedge_{D(b)} \Omega^*_{D(b)} = K^{*, b}
$$
是拟同构。注意，这些映射在上同调上定义{\it 同一个}映射（甚至在导出范畴
中定义同一个映射），因为其逆由余对角映射 $D(b)\to D$ 给出，
它对应于乘法映射 $P\otimes_A\ldots\otimes_AP\to P$。
因此，第二个谱序列的 $E_1$ 页为
$$
E_1^{a, b} = H^a(M \otimes^\wedge_D\Omega^*_D)
$$
其微分为
$$
E_1^{a, 0} \xrightarrow{0}
E_1^{a, 1} \xrightarrow{1}
E_1^{a, 2} \xrightarrow{0}
E_1^{a, 3} \xrightarrow{1} \ldots
$$
因为每个微分都是给定等同
$H^a(M\otimes^\wedge_D\Omega^*_D)=E_1^{a,0}=E_1^{a,1}=\ldots$
的交错和。因此，$E_2$ 页在第一行等于
$H^a(M\otimes^\wedge_D\Omega^*_D)$，在其他各处为零。
由此，$M\otimes^\wedge_D\Omega^*_D$ 与第一行的等同诱导它与
$\text{Tot}(K^{*,*})$ 之间的拟同构。
\end{proof}

\begin{lemma}
\label{crystalline-lemma-compute-cohomology-crystal-smooth}
假设同命题 \ref{crystalline-proposition-compute-cohomology-crystal}。
设 $A \to P' \to C$ 为环映射，其中 $A \to P'$ 光滑，而 $P' \to C$
为核等于 $J'$ 的满射。令 $D'$ 为 $D_{P', \gamma}(J')$ 的 $p$-进完备化。
设 $(M', \nabla')$ 为与 $\mathcal{F}$ 对应的 $D'$ 上偶对，见
引理 \ref{crystalline-lemma-crystals-on-affine-smooth}。则复形
$$
M' \otimes^\wedge_{D'} \Omega^*_{D'}
$$
计算 $R\Gamma(\text{Cris}(X/S), \mathcal{F})$。
\end{lemma}

\begin{proof}
依引理 \ref{crystalline-lemma-crystals-on-affine-smooth} 选取
$a : D \to D'$ 与 $b : D' \to D$。
注意，基变换 $M = M' \otimes_{D', b} D$ 连同其联络 $\nabla$
对应于 $\mathcal{F}$。因此由
命题 \ref{crystalline-proposition-compute-cohomology-crystal}，
$M \otimes^\wedge_D \Omega_D^*$ 计算 $\mathcal{F}$ 的晶体上同调。
故只需证明由 $a$ 与 $b$ 诱导的基变换映射
$$
M' \otimes^\wedge_{D'} \Omega^*_{D'}
\longrightarrow
M \otimes^\wedge_D \Omega^*_D
\quad\text{且}\quad
M \otimes^\wedge_D \Omega^*_D
\longrightarrow
M' \otimes^\wedge_{D'} \Omega^*_{D'}
$$
均为拟同构。由于 $a \circ b = \text{id}_{D'}$，其中一个方向的复合
是复形 $M' \otimes^\wedge_{D'} \Omega^*_{D'}$ 上的恒等映射。
因此只需证明映射
$$
M \otimes^\wedge_D \Omega^*_D
\longrightarrow
M \otimes^\wedge_D \Omega^*_D
$$
由 $b \circ a : D \to D$ 诱导时为拟同构。（注意，两端确为同一复形，
因为 $M = M' \otimes^\wedge_{D', b} D$，从而
$M \otimes^\wedge_{D, b \circ a} D =
M' \otimes^\wedge_{D', b \circ a \circ b} D =
M' \otimes^\wedge_{D', b} D = M$。）事实上我们断言：对任意与到 $C$
的增广相容的除幂 $A$-代数同态 $\rho : D \to D$，其诱导映射
$M \otimes^\wedge_D \Omega^*_D \to M \otimes^\wedge_{D, \rho} \Omega^*_D$
均为拟同构。

\medskip\noindent
写成 $\rho(x_i) = x_i + z_i$。由于 $\rho$ 与到 $C$ 的增广相容，
各元素 $z_i$ 均属于 $D$ 的除幂理想。因此可把映射 $\rho$ 分解为复合
$$
D \xrightarrow{\sigma} D\langle \xi_i \rangle^\wedge \xrightarrow{\tau} D
$$
其中第一条映射由 $x_i \mapsto x_i + \xi_i$ 给出，第二条映射是把
$\xi_i$ 映到 $z_i$ 的除幂 $D$-代数映射。（这里用到了多项式代数、
除幂多项式代数、除幂包络以及 $p$-进完备化的泛性质。）注意，
$D\langle \xi_i \rangle^\wedge$ 存在一个{\it 自同构} $\alpha$，满足
$\alpha(x_i) = x_i - \xi_i$ 且 $\alpha(\xi_i) = \xi_i$。
把 引理 \ref{crystalline-lemma-relative-poincare} 应用于 $\alpha \circ \sigma$
（它把 $x_i$ 映到 $x_i$），并利用 $\alpha$ 为同构，可知 $\sigma$
诱导 $M \otimes^\wedge_D \Omega^*_D$ 与下述复形之间的拟同构：
$M \otimes^\wedge_{D, \sigma} \Omega^*_{D\langle x_i \rangle^\wedge}$.
另一方面，映射 $\tau$ 以 $D \to D\langle x_i \rangle^\wedge$、
$x_i \mapsto x_i$ 为左逆。再次使用 引理 \ref{crystalline-lemma-relative-poincare}，
可知 $\tau$ 诱导下述两复形之间的拟同构：
$M \otimes^\wedge_{D, \sigma} \Omega^*_{D\langle x_i \rangle^\wedge}$
与 $M \otimes^\wedge_{D, \tau \circ \sigma} \Omega^*_D$。复合这两个
拟同构，即得 $\rho$ 所诱导的拟同构
$M \otimes^\wedge_D \Omega^*_D \to M \otimes^\wedge_{D, \rho} \Omega^*_D$，
如所要求。
\end{proof}










\section{两个反例}
\label{crystalline-section-examples}

\noindent
在讨论晶体上同调的一些成功应用之前，先给出两个例子，以说明当所论
概形在基上不固有或为奇异概形时，晶体上同调为何表现不佳。第一个例子
见 \cite{BO}。

\begin{example}
\label{crystalline-example-torsion}
令 $A = \mathbf{Z}_p$，并在除幂理想 $(p)$ 上赋予其唯一的除幂结构
$\gamma$。令 $C = \mathbf{F}_p[x, y]/(x^2, xy, y^2)$。选取表示
$$
C = P/J = \mathbf{Z}_p[x, y]/(x^2, xy, y^2, p)
$$
令 $D = D_{P, \gamma}(J)^\wedge$，其除幂理想
$(\bar J, \bar \gamma)$ 如第 \ref{crystalline-section-quasi-coherent-crystals} 节。
仍以 $x,y$ 表示它们在 $D$ 中的像。考虑元素
$$
\tau = \bar\gamma_p(x^2)\bar\gamma_p(y^2) - \bar\gamma_p(xy)^2 \in D
$$
注意 $p\tau=0$，因为在 $D$ 中
$$
p! \bar\gamma_p(x^2) \bar\gamma_p(y^2) =
x^{2p} \bar\gamma_p(y^2) = \bar\gamma_p(x^2y^2) =
x^py^p \bar\gamma_p(xy) = p! \bar\gamma_p(xy)^2
$$
。又注意到 $\Omega_D$ 中 $\text{d}\tau=0$，因为
\begin{align*}
\text{d}(\bar\gamma_p(x^2) \bar\gamma_p(y^2))
& =
\bar\gamma_{p - 1}(x^2)\bar\gamma_p(y^2)\text{d}x^2 +
\bar\gamma_p(x^2)\bar\gamma_{p - 1}(y^2)\text{d}y^2 \\
& =
2 x \bar\gamma_{p - 1}(x^2)\bar\gamma_p(y^2)\text{d}x +
2 y \bar\gamma_p(x^2)\bar\gamma_{p - 1}(y^2)\text{d}y \\
& =
2/(p - 1)!( x^{2p - 1} \bar\gamma_p(y^2)\text{d}x +
y^{2p - 1} \bar\gamma_p(x^2)\text{d}y ) \\
& =
2/(p - 1)!
(x^{p - 1} \bar\gamma_p(xy^2)\text{d}x +
y^{p - 1} \bar\gamma_p(x^2y)\text{d}y) \\
& =
2/(p - 1)!
(x^{p - 1}y^p \bar\gamma_p(xy)\text{d}x +
x^py^{p - 1} \bar\gamma_p(xy)\text{d}y) \\
& =
2 \bar\gamma_{p - 1}(xy) \bar\gamma_p(xy)(y\text{d}x + x \text{d}y) \\
& = 
\text{d}(\bar\gamma_p(xy)^2)
\end{align*}
最后断言 $D$ 中 $\tau\neq0$。为此，只需构造 $\text{Cris}(C/S)$ 的一个
对象 $(B \to \mathbf{F}_p[x, y]/(x^2, xy, y^2), \delta)$，使得 $\tau$
在 $B$ 中的像非零。取
$$
B = \mathbf{F}_p[x, y, u, v]/(x^3, x^2y, xy^2, y^3, xu, yu, xv, yv, u^2, v^2)
$$
并取其到 $C$ 的显然满射。令 $K=\Ker(B\to C)$，考虑映射
$$
\delta_p : K \longrightarrow K,\quad
ax^2 + bxy + cy^2 + du + ev + fuv \longmapsto a^pu + c^pv
$$
可检验它满足“除幂代数”引理
\ref{dpa-lemma-need-only-gamma-p} 的假设 (1)、(2)、(3)，故定义一个除幂结构。
此外，$\tau$ 映到 $uv$，而后者在 $B$ 中非零。置
$X=\Spec(C)$ 与 $S=\Spec(A)$。由此得到以下结论：
\begin{enumerate}
\item $H^0(\text{Cris}(X/S), \mathcal{O}_{X/S})$ 含有 $p$-挠；
\item Frobenius 拉回
$F^* : H^0(\text{Cris}(X/S), \mathcal{O}_{X/S})
\to H^0(\text{Cris}(X/S), \mathcal{O}_{X/S})$ 不是单射。
\end{enumerate}
具体地，由命题 \ref{crystalline-proposition-compute-cohomology-crystal}，
$\tau$ 定义 $H^0(\text{Cris}(X/S), \mathcal{O}_{X/S})$ 的一个非零挠元。
同样，$F^*(\tau)=\sigma(\tau)$，其中 $\sigma:D\to D$ 由 $P$ 上任一
Frobenius 提升诱导。若取 $\sigma(x)=x^p$ 与 $\sigma(y)=y^p$，则直接
计算可得 $F^*(\tau)=0$。
\end{example}

\noindent
下一个例子表明，即使对仿射 $n$-空间，晶体上同调也不给出所期望的结果。

\begin{example}
\label{crystalline-example-affine-n-space}
令 $A=\mathbf Z_p$，并在除幂理想 $(p)$ 上赋予其唯一的除幂结构
$\gamma$。令 $C=\mathbf F_p[x_1,\ldots,x_r]$。选取表示
$$
C = P/J = P/pP\quad\text{其中}\quad P = \mathbf{Z}_p[x_1, \ldots, x_r]
$$
由 “除幂代数” 引理 \ref{dpa-lemma-gamma-extends}，
$pP$ 带有除幂。故置 $D=P^\wedge$ 并取除幂理想 $(p)$，便得到
第 \ref{crystalline-section-quasi-coherent-crystals} 节中的情形。于是
$R\Gamma(\text{Cris}(X/S), \mathcal{O}_{X/S})$ 由复形
$$
D \to \Omega^1_D \to \Omega^2_D \to \ldots \to \Omega^r_D
$$
表示，见命题 \ref{crystalline-proposition-compute-cohomology-crystal}。
假设 $r>0$，可得：
\begin{enumerate}
\item $X=\mathbf A^r_{\mathbf F_p}$ 在 $S=\Spec(\mathbf Z_p)$ 上的晶体
结构层之晶体上同调仅在次数 $0,\ldots,r$ 可能非零。
\item $H^0(\text{Cris}(X/S), \mathcal{O}_{X/S})=\mathbf Z_p$。
\item 上同调群 $H^r(\text{Cris}(X/S), \mathcal{O}_{X/S})$ 是无限群，
且不是挠阿贝尔群。
\item 上同调群 $H^r(\text{Cris}(X/S), \mathcal{O}_{X/S})$
对 $p$-进拓扑不分离。
\end{enumerate}
前两条尚属合理，第 (3)、(4) 条却令人不安！只要具体写出上面展示的
复形，便立即可见这些断言。这里仅考察 $r=1$。此时所考察的是下述
$p$-进完备模的两项复形：
$$
\text{d} :
D = \left(
\bigoplus\nolimits_{n \geq 0} \mathbf{Z}_p x^n
\right)^\wedge
\longrightarrow
\Omega^1_D = \left(
\bigoplus\nolimits_{n \geq 1} \mathbf{Z}_p x^{n - 1}\text{d}x
\right)^\wedge
$$
除右端缺少第一个直和项外，该映射由
$\text{diag}(0,1,2,3,4,\ldots)$ 给出。显然，
$\bigoplus_{n>0}\mathbf Z_p/n\mathbf Z_p$ 是余核的一个子群，故余核
为无限群。事实上，元素
$$
\omega = \sum\nolimits_{e > 0} p^e x^{p^{2e} - 1}\text{d}x
$$
显然不是余核中的挠元。然而情况还更糟。考虑元素
$$
\eta = \sum\nolimits_{e > 0} p^e x^{p^e - 1}\text{d}x
$$
对每个 $t>0$，元素 $\eta$ 模 $\text d$ 的像同余于
$\sum_{e>t}p^ex^{p^e-1}\text d x$，而后者可被 $p^t$ 整除。因此
$\eta$ 在 $H^1(\text{Cris}(X/S),\mathcal O_{X/S})$ 中的上同调类属于
$\bigcap p^tH^1(\text{Cris}(X/S),\mathcal O_{X/S})$。但 $\eta$ 不在
$\text d$ 的像中，因为否则它必须是某个
$a+\sum_{e>0}x^{p^e}$（其中 $a\in\mathbf Z_p$）的像，而这个元素不属于
左端。因此 $\eta$ 的上同调类非零，第 (4) 条得证。（事实上，$\eta$
的上同调类不是挠元。）好的一面是，作为一个 $p$-进完备模复形的
上同调，各群 $H^i(\text{Cris}(X/S),\mathcal O_{X/S})$ 都是导出完备的
$\mathbf Z_p$-模；见“代数进阶”第
\ref{more-algebra-section-derived-completion} 节。
\end{example}








\section{应用}
\label{crystalline-section-applications}

\noindent
本节汇集前述各节内容的一些应用。

\begin{proposition}
\label{crystalline-proposition-compare-with-de-Rham}
处于 情形 \ref{crystalline-situation-global}。设 $\mathcal F$ 为
$\text{Cris}(X/S)$ 上拟凝聚模中的晶体。复形的截断映射
$$
(\mathcal{F} \to
\mathcal{F} \otimes_{\mathcal{O}_{X/S}} \Omega^1_{X/S} \to
\mathcal{F} \otimes_{\mathcal{O}_{X/S}} \Omega^2_{X/S} \to \ldots)
\longrightarrow \mathcal{F}[0],
$$
虽非拟同构，但施加 $Ru_{X/S,*}$ 后成为拟同构。事实上，对任意 $i>0$，
有
$$
Ru_{X/S, *}(\mathcal{F} \otimes_{\mathcal{O}_{X/S}} \Omega^i_{X/S}) = 0.
$$
\end{proposition}

\begin{proof}
由引理 \ref{crystalline-lemma-automatic-connection} 得到命题中所示的德拉姆复形。
简记 $\mathcal H=\mathcal F\otimes\Omega^i_{X/S}$。设 $X'\subset X$
为仿射开子概形，且映入某个仿射开子概形 $S'\subset S$。则
$$
(Ru_{X/S, *}\mathcal{H})|_{X'_{Zar}} =
Ru_{X'/S', *}(\mathcal{H}|_{\text{Cris}(X'/S')}),
$$
见引理 \ref{crystalline-lemma-localize}。因此引理 \ref{crystalline-lemma-cohomology-is-zero}
表明，$Ru_{X/S,*}\mathcal H$ 是 $X_{Zar}$ 上的层复形，且其在任意仿射
开集上的上同调均为零。由于 $X$ 的拓扑有一个由仿射开集组成的基，
这说明 $Ru_{X/S,*}\mathcal H$ 拟同构于零。
\end{proof}

\begin{remark}
\label{crystalline-remark-vanishing}
命题 \ref{crystalline-proposition-compare-with-de-Rham} 的证明表明，结论
$$
Ru_{X/S, *}(\mathcal{F} \otimes_{\mathcal{O}_{X/S}} \Omega^i_{X/S}) = 0
$$
对任意满足 命题 \ref{crystalline-proposition-compute-cohomology} 中条件 (1)、(2)
的 $\mathcal O_{X/S}$-模 $\mathcal F$，在 $i>0$ 时均成立。这适用于以下
非晶体：所有 $\Omega^i_{X/S}$，以及任意形如
$\underline{\mathcal F}$ 的层，其中 $\mathcal F$ 是拟凝聚
$\mathcal O_X$-模。特别地，它适用于层
$\underline{\mathcal O_X}=\underline{\mathbf G_a}$。但要构造德拉姆复形，
需要类似 引理 \ref{crystalline-lemma-automatic-connection} 的结果，而这要求
$\mathcal F$ 为晶体。因此，（目前）满足 命题
\ref{crystalline-proposition-compare-with-de-Rham} 完整断言的模层，恰组成拟凝聚模中
晶体的范畴。
\end{remark}

\noindent
处于 情形 \ref{crystalline-situation-global}。设 $\mathcal F$ 为
$\text{Cris}(X/S)$ 上拟凝聚模中的晶体，并设 $(U,T,\delta)$ 为
$\text{Cris}(X/S)$ 的对象。命题
\ref{crystalline-proposition-compare-with-de-Rham} 使我们能够构造典范映射
\begin{equation}
\label{crystalline-equation-restriction}
R\Gamma(\text{Cris}(X/S), \mathcal{F})
\longrightarrow
R\Gamma(T, \mathcal{F}_T \otimes_{\mathcal{O}_T} \Omega^*_{T/S, \delta})
\end{equation}
确切地说，有 $R\Gamma(\text{Cris}(X/S),\mathcal F)=
R\Gamma(\text{Cris}(X/S),\mathcal F\otimes\Omega^*_{X/S})$；可把整体
上同调类限制到 $T$，且由引理 \ref{crystalline-lemma-module-of-differentials}，
$\Omega_{X/S}$ 限制为 $\Omega_{T/S,\delta}$。

















\section{若干进一步结果}
\label{crystalline-section-missing}

\noindent
本节列出一些尚缺证明的结果。我们先把它们表述为一系列备注；只有在补入
详细证明后，才会将其改写为正式的引理和命题。

\begin{remark}[高阶正像]
\label{crystalline-remark-compute-direct-image}
设 $p$ 为素数，并设
$(S,\mathcal I,\gamma)\to(S',\mathcal I',\gamma')$ 为
$\mathbf Z_{(p)}$ 上除幂概形的态射。设
$$
\xymatrix{
X \ar[r]_f \ar[d] & X' \ar[d] \\
S_0 \ar[r] & S'_0
}
$$
为概形态射的交换图，并假设 $p$ 在 $X$ 与 $X'$ 上局部幂零。设
$\mathcal F$ 为 $\text{Cris}(X/S)$ 上的 $\mathcal O_{X/S}$-模。则
$Rf_{\text{cris},*}\mathcal F$ 可如下计算。

\medskip\noindent
给定 $\text{Cris}(X'/S')$ 的对象 $(U',T',\delta')$，置
$U=X\times_{X'}U'=f^{-1}(U')$（它是 $X$ 的开子概形）。以
$(T_0,T,\delta)$ 表示 $S$ 上的除幂概形，使得
$$
\xymatrix{
T \ar[r] \ar[d] & T' \ar[d] \\
S \ar[r] & S'
}
$$
在除幂概形范畴中为笛卡尔图，见引理 \ref{crystalline-lemma-fibre-product}。
存在诱导态射 $U\to T_0$，从而得到态射
$(U/T)_{\text{cris}}\to(X/S)_{\text{cris}}$，见注
\ref{crystalline-remark-functoriality-cris}。令 $\mathcal F_U$ 为 $\mathcal F$ 的拉回，
并令 $\tau_{U/T}:(U/T)_{\text{cris}}\to T_{Zar}$ 为结构态射。则有
\begin{equation}
\label{crystalline-equation-identify-pushforward}
\left(Rf_{\text{cris}, *}\mathcal{F}\right)_{T'} =
R(T \to T')_*\left(R\tau_{U/T, *} \mathcal{F}_U \right)
\end{equation}
其中左端是限制（见第 \ref{crystalline-section-sheaves} 节）。

\medskip\noindent
提示：首先证明 $\text{Cris}(U/T)$ 是 $\text{Cris}(X/S)$ 在集合层
$f_{\text{cris}}^{-1}h_{(U',T',\delta')}$ 处的局部化（取 “位点论” 引理
\ref{sites-lemma-localize-topos-site} 的意义）。其次，利用内射
$\mathcal O_{X/S}$-模的拉回是 $\text{Cris}(U/T)$ 上的内射
$\mathcal O_{U/T}$-模，把断言化归到 $\mathcal F$ 为内射模以及模的正像
这一情形。最后检验结果对普通正像成立。
\end{remark}

\begin{remark}[Mayer-Vietoris]
\label{crystalline-remark-mayer-vietoris}
在 注 \ref{crystalline-remark-compute-direct-image} 的情形中，设有开覆盖
$X=X'\cup X''$，并记 $X'''=X'\cap X''$。令 $f'$、$f''$ 与 $f''$
分别为 $f$ 在 $X'$、$X''$ 与 $X'''$ 上的限制。再令
$\mathcal F'$、$\mathcal F''$ 与 $\mathcal F'''$ 分别为 $\mathcal F$
在 $X'$、$X''$ 与 $X'''$ 的晶体位点上的限制。则存在可区别三角
$$
Rf_{\text{cris}, *}\mathcal{F}
\longrightarrow
Rf'_{\text{cris}, *}\mathcal{F}' \oplus Rf''_{\text{cris}, *}\mathcal{F}''
\longrightarrow
Rf'''_{\text{cris}, *}\mathcal{F}'''
\longrightarrow
Rf_{\text{cris}, *}\mathcal{F}[1]
$$
于 $D(\mathcal O_{X'/S'})$ 中。

\medskip\noindent
提示：这是下述事实的形式推论：子范畴 $\text{Cris}(X'/S)$、
$\text{Cris}(X''/S)$、$\text{Cris}(X'''/S)$ 对应于
$\text{Cris}(X/S)$ 的终层之开子对象，且第三个是前两个的交。
\end{remark}

\begin{remark}[{\v C}ech 复形]
\label{crystalline-remark-cech-complex}
设 $p$ 为素数，并设 $(A,I,\gamma)$ 为除幂环，其中 $A$ 是
$\mathbf Z_{(p)}$-代数。置 $S=\Spec(A)$ 与 $S_0=\Spec(A/I)$。设 $X$
为 $S_0$ 上的分离概形\footnote{此假设并非严格必要；利用超覆盖可将
本备注的构造推广到一般情形。}，且 $p$ 在 $X$ 上局部幂零。设
$\mathcal F$ 为拟凝聚 $\mathcal O_{X/S}$-模中的晶体。

\medskip\noindent
选取 $X$ 的仿射开覆盖
$X=\bigcup_{\lambda\in\Lambda}U_\lambda$。写成
$U_\lambda=\Spec(C_\lambda)$。选取 $A$ 上的多项式代数 $P_\lambda$
及满射 $P_\lambda\to C_\lambda$。固定这些选择后，可构造一个计算
$R\Gamma(\text{Cris}(X/S),\mathcal F)$ 的 {\v C}ech 复形。

\medskip\noindent
给定 $n\geq0$ 及 $\lambda_0,\ldots,\lambda_n\in\Lambda$，记
$U_{\lambda_0\ldots\lambda_n}=U_{\lambda_0}\cap\ldots\cap U_{\lambda_n}$。
由假设它是仿射概形。写成
$U_{\lambda_0\ldots\lambda_n}=\Spec(C_{\lambda_0\ldots\lambda_n})$，并置
$$
P_{\lambda_0 \ldots \lambda_n} =
P_{\lambda_0} \otimes_A \ldots \otimes_A P_{\lambda_n}
$$
它带有到 $C_{\lambda_0\ldots\lambda_n}$ 的典范满射。以
$J_{\lambda_0\ldots\lambda_n}$ 表示其核，并令
$D_{\lambda_0\ldots\lambda_n}$ 为
$P_{\lambda_0\ldots\lambda_n}$ 中 $J_{\lambda_0\ldots\lambda_n}$
相对于 $\gamma$ 的除幂包络之 $p$-进完备化。令
$M_{\lambda_0\ldots\lambda_n}$ 为通过 命题
\ref{crystalline-proposition-crystals-on-affine} 与 $\mathcal F$ 在
$\text{Cris}(U_{\lambda_0\ldots\lambda_n}/S)$ 上的限制相对应的
$P_{\lambda_0\ldots\lambda_n}$-模。由构造得到余单纯除幂环 $D(*)$，
其次数 $n$ 部分为环
$$
D(n) =
\prod\nolimits_{\lambda_0 \ldots \lambda_n}
D_{\lambda_0 \ldots \lambda_n}
$$
（这里使用除幂包络的函子性，以及在以类似方式定义的环 $P(*)$ 上的
平凡余单纯结构）。由于 $M_{\lambda_0\ldots\lambda_n}$ 是 $\mathcal F$
在对象 $\Spec(D_{\lambda_0\ldots\lambda_n})$ 上的“值”，故由规则
$$
M(n) = \prod\nolimits_{\lambda_0 \ldots \lambda_n}
M_{\lambda_0 \ldots \lambda_n}
$$
定义的 $M(*)$ 构成余单纯 $D(*)$-模。现断言
$$
R\Gamma(\text{Cris}(X/S), \mathcal{F}) = s(M(*))
$$
其中 $s(-)$ 表示与余单纯模相伴的上链复形（见 “单纯方法” 节
\ref{simplicial-section-dold-kan-cosimplicial}）。

\medskip\noindent
提示：其证明类似于 命题 \ref{crystalline-proposition-compute-cohomology} 的
证明（特别地，该结果对满足该命题假设的任意模成立）。
\end{remark}

\begin{remark}[交错 {\v C}ech 复形]
\label{crystalline-remark-alternating-cech-complex}
设 $p$ 为素数，并设 $(A,I,\gamma)$ 为除幂环，其中 $A$ 是
$\mathbf Z_{(p)}$-代数。置 $S=\Spec(A)$ 与 $S_0=\Spec(A/I)$。设 $X$
为 $S_0$ 上的分离拟紧概形，且 $p$ 在 $X$ 上局部幂零。设
$\mathcal F$ 为拟凝聚 $\mathcal O_{X/S}$-模中的晶体。

\medskip\noindent
选取 $X$ 的有限仿射开覆盖
$X=\bigcup_{\lambda\in\Lambda}U_\lambda$，并在 $\Lambda$ 上选取一个全序。
写成 $U_\lambda=\Spec(C_\lambda)$。选取 $A$ 上的多项式代数
$P_\lambda$ 及满射 $P_\lambda\to C_\lambda$。固定这些选择后，可构造
一个计算 $R\Gamma(\text{Cris}(X/S),\mathcal F)$ 的交错 {\v C}ech 复形。

\medskip\noindent
采用 注 \ref{crystalline-remark-cech-complex} 中引入的记号。以
$\Omega_{\lambda_0\ldots\lambda_n}$ 表示
$D_{\lambda_0\ldots\lambda_n}$ 在 $A$ 上与除幂结构相容的微分模之
$p$-进完备化。令 $\nabla$ 为 命题
\ref{crystalline-proposition-crystals-on-affine} 所给出的
$M_{\lambda_0\ldots\lambda_n}$ 上可积联络。考虑双复形
$M^{\bullet,\bullet}$，其各项为
$$
M^{n, m} =
\bigoplus\nolimits_{\lambda_0 < \ldots < \lambda_n}
M_{\lambda_0 \ldots \lambda_n}
\otimes^\wedge_{D_{\lambda_0 \ldots \lambda_n}}
\Omega^m_{D_{\lambda_0 \ldots \lambda_n}}.
$$
微分 $d_1$（使 $n$ 增加）取通常的 {\v C}ech 微分，而微分 $d_2$
取联络，即德拉姆复形的微分。我们断言
$$
R\Gamma(\text{Cris}(X/S), \mathcal{F}) = \text{Tot}(M^{\bullet, \bullet})
$$
其中 $\text{Tot}(-)$ 表示与双复形相伴的全复形，见 “同调”
定义 \ref{homology-definition-associated-simple-complex}。

\medskip\noindent
提示：由命题 \ref{crystalline-proposition-compare-with-de-Rham} 有
$$
R\Gamma(\text{Cris}(X/S), \mathcal{F}) = R\Gamma(\text{Cris}(X/S),
\mathcal{F} \otimes_{\mathcal{O}_{X/S}} \Omega_{X/S}^\bullet)
$$
。公式右端正是覆盖 $X=\bigcup_{\lambda\in\Lambda}U_\lambda$
（它诱导 $\text{Cris}(X/S)$ 的终层之一个开覆盖）与复形
$\mathcal F\otimes_{\mathcal O_{X/S}}\Omega_{X/S}^\bullet$ 的交错
{\v C}ech 复形；见命题
\ref{crystalline-proposition-compute-cohomology-crystal}。结论现由位点上同调的一个
一般结果推出：若交错 {\v C}ech 复形在每个部分上给出正确答案，它就
计算上同调（以后在此插入引用）。
\end{remark}

\begin{remark}[拟凝聚性]
\label{crystalline-remark-quasi-coherent}
在 注 \ref{crystalline-remark-compute-direct-image} 的情形中，假设 $S\to S'$
与 $X\to S_0$ 都拟紧且拟分离。则对拟凝聚 $\mathcal O_{X/S}$-模中的
晶体 $\mathcal F$，各层 $R^if_{\text{cris},*}\mathcal F$ 均局部拟凝聚。

\medskip\noindent
提示：须证明其在 $T'$ 上的限制是拟凝聚 $\mathcal O_{T'}$-模，其中
$(U',T',\delta')$ 为 $\text{Cris}(X'/S')$ 的任意对象。只需在 $T'$
仿射时证明。利用公式 (\ref{crystalline-equation-identify-pushforward})、$T\to T'$
拟紧且拟分离这一事实（因为 $T$ 在 $T'$ 沿 $S\to S'$ 的基变换上
仿射），以及“概形上同调”引理
\ref{coherent-lemma-quasi-coherence-higher-direct-images}，可知只需证明
各层 $R^i\tau_{U/T,*}\mathcal F_U$ 拟凝聚。注意 $U\to T_0$ 也拟紧且
拟分离；见“概形”引理 \ref{schemes-lemma-quasi-compact-permanence}
与 \ref{schemes-lemma-quasi-compact-permanence}。

\medskip\noindent
这把问题化为：当 $p$ 在 $S$ 上局部幂零时，证明
$R^i\tau_{X/S,*}\mathcal F$ 在 $S$ 上拟凝聚。这里 $\tau_{X/S}$ 是结构
态射，见注 \ref{crystalline-remark-structure-morphism}。可以在 $S$ 上局部地
工作，故可假设 $S$ 仿射（见引理 \ref{crystalline-lemma-localize}）。对覆盖 $X$
的仿射开集数目作归纳，并使用 Mayer--Vietoris（注
\ref{crystalline-remark-mayer-vietoris}），可把问题化归到 $X$ 也仿射的情形
（如“概形上同调”引理
\ref{coherent-lemma-quasi-coherence-higher-direct-images} 的证明）。设
$X=\Spec(C)$、$S=\Spec(A)$，使得 $(A,I,\gamma)$ 与 $A\to C$ 如
情形 \ref{crystalline-situation-affine}。依第
\ref{crystalline-section-quasi-coherent-crystals} 节选取 $A$ 上多项式代数 $P$ 及满射
$P\to C$。令 $(M,\nabla)$ 为与 $\mathcal F$ 对应的模，见命题
\ref{crystalline-proposition-crystals-on-affine}。应用命题
\ref{crystalline-proposition-compute-cohomology-crystal}，可知
$R\Gamma(\text{Cris}(X/S),\mathcal F)$ 由 $M\otimes_D\Omega_D^*$ 表示。
注意，由于 $p$ 在 $A$ 中幂零，无需完备化！须证明这与在
$S=\Spec(A)$ 中取主开集相容。设 $g\in A$。类似地，
$R\Gamma(\text{Cris}(X_g/S_g),\mathcal F)$ 由
$M_g\otimes_{D_g}\Omega_{D_g}^*$ 计算（这里再次使用无需 $p$-进完备化）。
由于局部化是 $A$-模上的正合函子，结论随即成立。
\end{remark}

\begin{remark}[有界性]
\label{crystalline-remark-bounded-cohomology}
在 注 \ref{crystalline-remark-compute-direct-image} 的情形中，假设 $S\to S'$
拟紧且拟分离，而 $X\to S_0$ 为有限型且拟分离。则存在整数 $i_0$，
使得对拟凝聚 $\mathcal O_{X/S}$-模中的任意晶体 $\mathcal F$，当
$i>i_0$ 时均有 $R^if_{\text{cris},*}\mathcal F=0$。

\medskip\noindent
提示：仿照注 \ref{crystalline-remark-quasi-coherent} 的论证（使用“概形上同调”引理
\ref{coherent-lemma-quasi-coherence-higher-direct-images}），问题化为在
命题 \ref{crystalline-proposition-compute-cohomology-crystal} 的情形中，当 $C$
是 $A$ 上有限型代数时，证明 $i\gg0$ 有
$H^i(\text{Cris}(X/S),\mathcal F)=0$。这是显然的，因为可以选取有限
多项式代数，而 $i\gg0$ 时 $\Omega_D^i=0$。
\end{remark}

\begin{remark}[具体的有界性]
\label{crystalline-remark-bounded-cohomology-over-point}
在 情形 \ref{crystalline-situation-global} 中，设 $\mathcal F$ 为拟凝聚
$\mathcal O_{X/S}$-模中的晶体。假设 $S_0$ 仅有一个点，且
$X\to S_0$ 为有限表示。
\begin{enumerate}
\item 若 $\dim X=d$ 且 $X/S_0$ 的嵌入维数为 $e$，则 $i>d+e$ 时
$H^i(\text{Cris}(X/S),\mathcal F)=0$。
\item 若 $X$ 分离且可由 $q$ 个仿射开集覆盖，并且 $X/S_0$ 的嵌入维数
为 $e$，则 $i>q+e$ 时 $H^i(\text{Cris}(X/S),\mathcal F)=0$。
\end{enumerate}
提示：在情形 (1) 中可使用
$$
H^i(\text{Cris}(X/S), \mathcal{F}) = H^i(X_{Zar}, Ru_{X/S, *}\mathcal{F})
$$
以及如下事实：$Ru_{X/S,*}\mathcal F$ 局部地由一个德拉姆复形计算，
该复形利用 $X$ 到 $S$ 上维数为 $e$ 的光滑概形中的嵌入构造
（见引理 \ref{crystalline-lemma-compute-cohomology-crystal-smooth}）。这些德拉姆
复形在所有次数 $>e$ 处为零。因此 (1) 由 “上同调” 命题
\ref{cohomology-proposition-vanishing-Noetherian} 推出。在情形 (2) 中，
使用交错 {\v C}ech 复形（见注
\ref{crystalline-remark-alternating-cech-complex}）化归到 $X$ 仿射的情形。仿射
情形下，利用与 $X$ 到 $S$ 上维数为 $e$ 的光滑概形中某个嵌入相伴的
德拉姆复形证明结果（构造这样的嵌入需要一些工作）。
\end{remark}

\begin{remark}[基变换映射]
\label{crystalline-remark-base-change}
在 注 \ref{crystalline-remark-compute-direct-image} 的情形中，假设
$S=\Spec(A)$ 与 $S'=\Spec(A')$ 仿射。设 $\mathcal F'$ 为
$\mathcal O_{X'/S'}$-模，并设 $\mathcal F$ 为 $\mathcal F'$ 的拉回。
则存在典范基变换映射
$$
L(S' \to S)^*R\tau_{X'/S', *}\mathcal{F}'
\longrightarrow
R\tau_{X/S, *}\mathcal{F}
$$
其中 $\tau_{X/S}$ 与 $\tau_{X'/S'}$ 是结构态射，见注
\ref{crystalline-remark-structure-morphism}。在整体截面上，这给出基变换映射
\begin{equation}
\label{crystalline-equation-base-change-map}
R\Gamma(\text{Cris}(X'/S'), \mathcal{F}') \otimes^\mathbf{L}_{A'} A
\longrightarrow
R\Gamma(\text{Cris}(X/S), \mathcal{F})
\end{equation}
于 $D(A)$ 中。

\medskip\noindent
提示：把“位点上同调”注
\ref{sites-cohomology-remark-base-change} 的极一般基变换映射与典范映射
$Lf_{\text{cris}}^*\mathcal{F}' \to
f_{\text{cris}}^*\mathcal{F}' = \mathcal{F}$.
\end{remark}

\begin{remark}[基变换同构]
\label{crystalline-remark-base-change-isomorphism}
若以下条件全部满足，则映射 (\ref{crystalline-equation-base-change-map}) 为同构：
\begin{enumerate}
\item $p$ 在 $A'$ 中幂零；
\item $\mathcal F'$ 是拟凝聚 $\mathcal O_{X'/S'}$-模中的晶体；
\item $X'\to S'_0$ 为拟紧拟分离态射；
\item $X = X' \times_{S'_0} S_0$,
\item $\mathcal F'$ 是平坦 $\mathcal O_{X'/S'}$-模；
\item $X'\to S'_0$ 为局部完全交态射（见“态射进阶”
定义 \ref{more-morphisms-definition-lci}；例如，当 $X'\to S'_0$
为 syntomic 或光滑态射时成立）；
\item $X'$ 与 $S_0$ 在 $S'_0$ 上 Tor 独立（见“代数进阶”
定义 \ref{more-algebra-definition-tor-independent}；例如，当
$S_0\to S'_0$ 或 $X'\to S'_0$ 中任一个平坦时成立）。
\end{enumerate}
提示：条件 (1) 意味着在以下论证中，$p$-进完备化不起作用，因而可以忽略。
利用条件 (3) 与 Mayer--Vietoris（见注
\ref{crystalline-remark-mayer-vietoris}），问题化归到 $X'$ 仿射的情形。事实上，
由条件 (6)，进一步缩小后可假设 $X'=\Spec(C')$，并给定表示
$C'=A'/I'[x_1,\ldots,x_n]/(\bar f'_1,\ldots,\bar f'_c)$，其中
$\bar f'_1,\ldots,\bar f'_c$ 是 $A'/I'$ 中的 Koszul-正则列。（这意味着
光滑局部地，$\bar f'_1,\ldots,\bar f'_c$ 构成正则列；见“代数进阶”引理
\ref{more-algebra-lemma-Koszul-regular-flat-locally-regular}。）选取
$\bar f'_i$ 到元素 $f'_i\in A'[x_1,\ldots,x_n]$ 的提升。由 (4)，
$X=\Spec(C)$，其中
$C=A/I[x_1,\ldots,x_n]/(\bar f_1,\ldots,\bar f_c)$，而
$f_i\in A[x_1,\ldots,x_n]$ 是 $f'_i$ 的像。由性质 (7)，
$\bar f_1,\ldots,\bar f_c$ 是 $A/I[x_1,\ldots,x_n]$ 中的
Koszul-正则列。$A'[x_1,\ldots,x_n]$ 中理想
$I'A'[x_1,\ldots,x_n]+(f'_1,\ldots,f'_c)$ 相对于 $\gamma'$ 的除幂包络为
$$
D' = A'[x_1, \ldots, x_n]\langle \xi_1, \ldots, \xi_c \rangle/(\xi_i - f'_i)
$$
，见引理 \ref{crystalline-lemma-describe-divided-power-envelope}。然后可检验
$\xi_1-f'_1,\ldots,\xi_n-f'_n$ 是环
$A'[x_1,\ldots,x_n]\langle\xi_1,\ldots,\xi_c\rangle$ 中的
Koszul-正则列。类似地，$A[x_1,\ldots,x_n]$ 中理想
$IA[x_1,\ldots,x_n]+(f_1,\ldots,f_c)$ 相对于 $\gamma$ 的除幂包络为
$$
D = A[x_1, \ldots, x_n]\langle \xi_1, \ldots, \xi_c\rangle/(\xi_i - f_i)
$$
，且 $\xi_1-f_1,\ldots,\xi_n-f_n$ 是环
$A[x_1,\ldots,x_n]\langle\xi_1,\ldots,\xi_c\rangle$ 中的
Koszul-正则列。于是 $D'\otimes_{A'}^\mathbf L A=D$。条件 (2) 蕴含
$\mathcal F'$ 对应于一个由带联络 $D'$-模组成的偶对 $(M',\nabla)$，
见命题 \ref{crystalline-proposition-crystals-on-affine}。则
$M=M'\otimes_{D'}D$ 对应于拉回 $\mathcal F$。由假设 (5)，$M'$ 是平坦
$D'$-模，故
$$
M = M' \otimes_{D'} D = M' \otimes_{D'} D' \otimes_{A'}^\mathbf{L} A
= M' \otimes_{A'}^\mathbf{L} A
$$
由于微分模 $\Omega_{D'}$ 与 $\Omega_D$（定义见节
\ref{crystalline-section-quasi-coherent-crystals}）是由相同生成元生成的自由
$D'$-模，可得
$$
M \otimes_D \Omega^\bullet_D =
M' \otimes_{D'} \Omega^\bullet_{D'} \otimes_{D'} D =
M' \otimes_{D'} \Omega^\bullet_{D'} \otimes_{A'}^\mathbf{L} A
$$
，再由命题 \ref{crystalline-proposition-compute-cohomology-crystal} 即得所需。
\end{remark}

\begin{remark}[Rlim]
\label{crystalline-remark-rlim}
设 $p$ 为素数，并设 $(A,I,\gamma)$ 为除幂环，其中 $A$ 是
$\mathbf Z_{(p)}$-代数，且 $p$ 在 $A/I$ 中幂零。置 $S=\Spec(A)$ 与
$S_0=\Spec(A/I)$。设 $X$ 为 $S_0$ 上概形，且 $p$ 在 $X$ 上局部幂零。
设 $\mathcal F$ 为任意 $\mathcal O_{X/S}$-模。当 $e\gg0$ 时，
$(p^e)\subset I$ 被 $\gamma$ 保持，见“除幂代数”引理
\ref{dpa-lemma-extend-to-completion}。对 $e\gg0$ 置
$S_e=\Spec(A/p^eA)$。则 $\text{Cris}(X/S_e)$ 是
$\text{Cris}(X/S)$ 的全子范畴，并以 $\mathcal F_e$ 表示 $\mathcal F$
在 $\text{Cris}(X/S_e)$ 上的限制。于是
$$
R\Gamma(\text{Cris}(X/S), \mathcal{F}) =
R\lim_e R\Gamma(\text{Cris}(X/S_e), \mathcal{F}_e)
$$

\medskip\noindent
提示：只需对内射的 $\mathcal F$ 证明。此时各层 $\mathcal F_e$ 也都是
内射模，转移映射 $\Gamma(\mathcal F_{e+1})\to\Gamma(\mathcal F_e)$
满射，并且 $\Gamma(\mathcal F)=\lim_e\Gamma(\mathcal F_e)$；这是因为按
$\text{Cris}(X/S)$ 的定义，其任意对象局部地都是某个范畴
$\text{Cris}(X/S_e)$ 的对象。
\end{remark}

\begin{remark}[比较]
\label{crystalline-remark-comparison}
设 $p$ 为素数，并设 $(A,I,\gamma)$ 为除幂环，其中 $p$ 在 $A$ 中幂零。
置 $S=\Spec(A)$ 与 $S_0=\Spec(A/I)$。设 $Y$ 为 $S$ 上光滑概形，并置
$X=Y\times_SS_0$。设 $\mathcal F$ 为拟凝聚 $\mathcal O_{X/S}$-模中的
晶体。则
\begin{enumerate}
\item $\gamma$ 延拓为 $X$ 在 $Y$ 中的理想上的除幂结构，使
$(X,Y,\gamma)$ 成为 $\text{Cris}(X/S)$ 的对象；
\item 限制 $\mathcal F_Y$（见第 \ref{crystalline-section-sheaves} 节）带有典范
可积联络
$\nabla : \mathcal{F}_Y \to
\mathcal{F}_Y \otimes_{\mathcal{O}_Y} \Omega_{Y/S}$；
\item 有
$$
R\Gamma(\text{Cris}(X/S), \mathcal{F}) =
R\Gamma(Y, \mathcal{F}_Y \otimes_{\mathcal{O}_Y} \Omega^\bullet_{Y/S})
$$
于 $D(A)$ 中。
\end{enumerate}
提示：(1) 见“除幂代数”引理
\ref{dpa-lemma-gamma-extends}；(2) 见引理
\ref{crystalline-lemma-automatic-connection}。对 (3)，注意存在映射
(\ref{crystalline-equation-restriction})。当 $X$ 仿射时，此映射为同构，见引理
\ref{crystalline-lemma-compute-cohomology-crystal-smooth}。这说明在
$Y_{Zar}=X_{Zar}$ 上，复形 $Ru_{X/S,*}\mathcal F$ 与
$\mathcal F_Y\otimes\Omega^\bullet_{Y/S}$ 拟同构。由于
$R\Gamma(\text{Cris}(X/S),\mathcal F)=
R\Gamma(X_{Zar},Ru_{X/S,*}\mathcal F)$，结论成立。
\end{remark}

\begin{remark}[完美性]
\label{crystalline-remark-perfect}
设 $p$ 为素数，并设 $(A,I,\gamma)$ 为除幂环，其中 $p$ 在 $A$ 中幂零。
置 $S=\Spec(A)$ 与 $S_0=\Spec(A/I)$。设 $X$ 为 $S_0$ 上固有光滑概形。
设 $\mathcal F$ 为有限局部自由拟凝聚 $\mathcal O_{X/S}$-模中的晶体。
则 $R\Gamma(\text{Cris}(X/S),\mathcal F)$ 是 $D(A)$ 中的完美对象。

\medskip\noindent
提示：由注 \ref{crystalline-remark-base-change-isomorphism} 有
$$
R\Gamma(\text{Cris}(X/S), \mathcal{F}) \otimes_A^\mathbf{L} A/I
\cong
R\Gamma(\text{Cris}(X/S_0), \mathcal{F}|_{\text{Cris}(X/S_0)})
$$
由注 \ref{crystalline-remark-comparison} 有
$$
R\Gamma(\text{Cris}(X/S_0), \mathcal{F}|_{\text{Cris}(X/S_0)}) =
R\Gamma(X, \mathcal{F}_X \otimes \Omega^\bullet_{X/S_0})
$$
在德拉姆复形上使用朴素滤过，可知只要各复形
$$
R\Gamma(X, \mathcal{F}_X \otimes \Omega^q_{X/S_0})
$$
在 $D(A/I)$ 中完美，前一展示复形就在 $D(A/I)$ 中完美；见“代数进阶”
引理 \ref{more-algebra-lemma-two-out-of-three-perfect}。利用
$\mathcal F_X\otimes\Omega^q_{X/S_0}$ 是有限局部自由层以及
$X\to S_0$ 固有且平坦，凝聚上同调中的标准论证可知此条件成立
（以后在此插入引用）。应用“代数进阶”引理
\ref{more-algebra-lemma-perfect-modulo-nilpotent-ideal}，可知
$$
R\Gamma(\text{Cris}(X/S), \mathcal{F}) \otimes_A^\mathbf{L} A/I^n
$$
对所有 $n$ 都是 $D(A/I^n)$ 中的完美对象。除非 $A$ 为 Noether 环，
这还不够。确切地说，尽管由 $p$ 幂零的假设，$I$ 局部幂零，见
“除幂代数”引理 \ref{dpa-lemma-nil}，仍不能推出存在 $n$
使 $I^n=0$。反例为 $\mathbf F_p\langle x\rangle$。当
$\mathcal F=\mathcal O_{X/S}$ 时，一般情形可用
\url{https://math.columbia.edu/~dejong/wordpress/?p=2227}
中的论证证明。当系数 $\mathcal F$ 非平凡时，\cite{Faltings-very} 的
论证似乎如下。由“代数进阶”引理
\ref{more-algebra-lemma-perfect-modulo-nilpotent-ideal}，化归到 $pA=0$
的情形。此时 Frobenius 映射 $A\to A$、$a\mapsto a^p$ 分解为
$A\to A/I\xrightarrow{\varphi}A$（因为 $x\in I$ 时 $x^p=0$）。置
$X^{(1)}=X\otimes_{A/I,\varphi}A$。$X$ 的绝对 Frobenius 态射经过态射
$F_X:X\to X^{(1)}$ 分解（某种相对 Frobenius）。仿射局部地，若
$X=\Spec(C)$，则 $X^{(1)}=\Spec(C\otimes_{A/I,\varphi}A)$，而 $F_X$
对应于 $C\otimes_{A/I,\varphi}A\to C$、$c\otimes a\mapsto c^pa$。
这定义环化拓扑斯的态射
$$
(X/S)_{\text{cris}}
\xrightarrow{(F_X)_{\text{cris}}}
(X^{(1)}/S)_{\text{cris}}
\xrightarrow{u_{X^{(1)}/S}}
\Sh(X^{(1)}_{Zar})
$$
，其复合记为 $\text{Frob}_X$。然后通过局部计算证明
$R\text{Frob}_{X,*}\mathcal F$ 可由 $\mathcal O_{X^{(1)}}$-模的一个
完美复形表示（！）。
\end{remark}

\begin{remark}[完备情形的完美性]
\label{crystalline-remark-complete-perfect}
设 $p$ 为素数，并设 $(A,I,\gamma)$ 为除幂环，其中 $A$ 是 $p$-进完备环，
且 $p$ 在 $A/I$ 中幂零。置 $S=\Spec(A)$ 与 $S_0=\Spec(A/I)$。设 $X$
为 $S_0$ 上固有光滑概形，并设 $\mathcal F$ 为有限局部自由拟凝聚
$\mathcal O_{X/S}$-模中的晶体。则
$R\Gamma(\text{Cris}(X/S),\mathcal F)$ 是 $D(A)$ 中的完美对象。

\medskip\noindent
提示：已知 $K=R\Gamma(\text{Cris}(X/S),\mathcal F)$ 是 $A/p^eA$ 上
各上同调的导出极限 $K=R\lim K_e$，见注 \ref{crystalline-remark-rlim}。
由注 \ref{crystalline-remark-perfect}，每个 $K_e$ 都是 $D(A/p^eA)$ 中的完美
复形。由于 $A$ 是 $p$-进完备的，结论由“代数进阶”引理
\ref{more-algebra-lemma-Rlim-perfect-gives-complete} 推出。
\end{remark}

\begin{remark}[完备情形的比较]
\label{crystalline-remark-complete-comparison}
设 $p$ 为素数，并设 $(A,I,\gamma)$ 为除幂环，其中 $A$ 为 Noether
$p$-进完备环，且 $p$ 在 $A/I$ 中幂零。置 $S=\Spec(A)$ 与
$S_0=\Spec(A/I)$。设 $Y$ 为 $S$ 上固有光滑概形，并置
$X=Y\times_SS_0$。设 $\mathcal F$ 为拟凝聚 $\mathcal O_{X/S}$-模中的
有限型晶体。则
\begin{enumerate}
\item 存在带可积联络的凝聚 $\mathcal O_Y$-模 $\mathcal F_Y$
$$
\nabla :
\mathcal{F}_Y
\longrightarrow
\mathcal{F}_Y \otimes_{\mathcal{O}_Y} \Omega_{Y/S}
$$
，使得 $\mathcal F_Y/p^e\mathcal F_Y$ 是 注
\ref{crystalline-remark-comparison} 中所得的 $A/p^eA$ 上带联络模；
\item 有
$$
R\Gamma(\text{Cris}(X/S), \mathcal{F}) =
R\Gamma(Y, \mathcal{F}_Y \otimes_{\mathcal{O}_Y} \Omega^\bullet_{Y/S})
$$
于 $D(A)$ 中。
\end{enumerate}
提示：$\mathcal F_Y$ 的存在性即 Grothendieck 存在定理（以后在此插入
引用）。两端都由模 $p^e$ 版本的 $R\lim$ 计算，故得到上同调的同构
（左端见注 \ref{crystalline-remark-rlim}；右端使用形式函数定理，见“概形上同调”
定理 \ref{coherent-theorem-formal-functions}）。由注
\ref{crystalline-remark-comparison}，各模 $p^e$ 版本彼此同构。
\end{remark}




\section{沿纯不可分映射拉回}
\label{crystalline-section-pull-back-along-pth-root}

\noindent
所谓 $\alpha_p$-覆盖，是指形如
$$
X' = \Spec(C[z]/(z^p - c)) \longrightarrow \Spec(C) = X
$$
的态射，其中 $C$ 为 $\mathbf F_p$-代数且 $c\in C$。等价地，$X'$ 是
$X$ 上的 $\alpha_p$-挠子。所谓{\it 迭代 $\alpha_p$-覆盖}
\footnote{这不是标准记号。}，是指特征 $p$ 的概形态射，它在目标上
局部地是有限多个 $\alpha_p$-覆盖的复合。本节证明，沿这类态射的拉回
在逆转素数 $p$ 后诱导晶体上同调的拟同构。事实上，我们将证明该结果
的一个精确版本。先给出一个需要若干记号才能表述的预备引理。

\medskip\noindent
假设给定环映射 $B\to B'$ 及满足 注
\ref{crystalline-remark-base-change-connection} 假设的商
$\Omega_B\to\Omega$ 与 $\Omega_{B'}\to\Omega'$。于是
(\ref{crystalline-equation-base-change-map-complexes}) 给出复形的典范映射
$$
c_M^\bullet :
M \otimes_B \Omega^\bullet
\longrightarrow
M \otimes_B (\Omega')^\bullet
$$
，其中 $M$ 为任意带可积联络
$\nabla:M\to M\otimes_B\Omega_B$ 的 $B$-模。

\medskip\noindent
再假设给定 $a\in B$、$z\in B'$ 及映射 $\theta:B'\to B'$，满足：
\begin{enumerate}
\item
\label{crystalline-item-d-a-zero}
$\text{d}(a) = 0$,
\item
\label{crystalline-item-direct-sum}
$\Omega'=B'\otimes_B\Omega\oplus B'\text d z$；对所有 $f\in B'$ 写成
$\text{d}(f) = \text{d}_1(f) + \partial_z(f) \text{d}z$
，其中 $\text d_1(f)\in B'\otimes\Omega$ 且 $\partial_z(f)\in B'$；
\item
\label{crystalline-item-theta-linear}
$\theta:B'\to B'$ 为 $B$-线性；
\item
\label{crystalline-item-integrate}
$\partial_z \circ \theta = a$,
\item
\label{crystalline-item-injective}
$B\to B'$ 泛单射（因而 $\Omega\to\Omega'$ 单射）；
\item
\label{crystalline-item-factor}
对所有 $f\in B'$，$af-\theta(\partial_z(f))\in B$；
\item
\label{crystalline-item-horizontal}
$(\theta \otimes 1)(\text{d}_1(f)) - \text{d}_1(\theta(f)) \in \Omega$
对所有 $f\in B'$，其中
$\theta \otimes 1 : B' \otimes \Omega \to B' \otimes \Omega$
\end{enumerate}
这些条件在逻辑上并非彼此独立。例如，假设 (\ref{crystalline-item-integrate}) 蕴含
$\partial_z(af-\theta(\partial_z(f)))=0$。因此，若 $B\to B'$ 的像恰为
被 $\partial_z$ 消去的元素全体，则 (\ref{crystalline-item-factor}) 随即成立。
对条件 (\ref{crystalline-item-horizontal}) 可作类似论证。

\begin{lemma}
\label{crystalline-lemma-find-homotopy}
在上述情形中，存在复形映射
$$
e_M^\bullet :
M \otimes_B (\Omega')^\bullet
\longrightarrow
M \otimes_B \Omega^\bullet
$$
，使 $c_M^\bullet\circ e_M^\bullet$ 与
$e_M^\bullet\circ c_M^\bullet$ 都同伦于乘以 $a$ 的映射。
\end{lemma}

\begin{proof}
本证明中的一切张量积均取于 $B$ 上。假设 (\ref{crystalline-item-direct-sum}) 蕴含
$$
M \otimes (\Omega')^i =
(B' \otimes M \otimes \Omega^i)
\oplus
(B' \text{d}z \otimes M \otimes \Omega^{i - 1})
$$
对所有 $i\geq0$ 成立。$M\otimes(\Omega')^i$ 的一组加法生成元由形如
$f\omega$ 与 $f\text d z\wedge\eta$ 的元素组成，其中 $f\in B'$、
$\omega\in M\otimes\Omega^i$ 且 $\eta\in M\otimes\Omega^{i-1}$。

\medskip\noindent
对 $f\in B'$，记
$$
\epsilon(f) = af - \theta(\partial_z(f))
\quad\text{且}\quad
\epsilon'(f) = (\theta \otimes 1)(\text{d}_1(f)) - \text{d}_1(\theta(f))
$$
由假设 (\ref{crystalline-item-factor}) 与 (\ref{crystalline-item-horizontal})，有
$\epsilon(f)\in B$ 且 $\epsilon'(f)\in\Omega$。按规则
$e^i_M(f\omega)=\epsilon(f)\omega$ 与
$e^i_M(f\text d z\wedge\eta)=\epsilon'(f)\wedge\eta$ 定义
$e_M^\bullet$。下面将看到，映射族 $e_M^i$ 构成复形映射。

\medskip\noindent
定义
$$
h^i : M \otimes_B (\Omega')^i \longrightarrow M \otimes_B (\Omega')^{i - 1}
$$
，对上述元素规定 $h^i(f\omega)=0$ 以及
$h^i(f\text d z\wedge\eta)=\theta(f)\eta$。我们断言
$$
\text{d} \circ h + h \circ \text{d} = a - c_M^\bullet \circ e_M^\bullet
$$
由 (\ref{crystalline-item-d-a-zero})，乘以 $a$ 是复形映射。又由假设
(\ref{crystalline-item-injective})，$c_M^\bullet$ 是单射复形映射，故可知
$e_M^\bullet$ 是复形映射。为证明断言，计算
\begin{align*}
(\text{d} \circ h + h \circ \text{d})(f \omega)
& =
h\left(\text{d}(f) \wedge \omega + f \nabla(\omega)\right)
\\
& =
\theta(\partial_z(f)) \omega
\\
& =
a f\omega - \epsilon(f)\omega 
\\
& =
a f \omega - c^i_M(e^i_M(f\omega))
\end{align*}
第二个等号成立是因为 $\nabla(\omega)$ 中不出现 $\text d z$，第三个
等号则由假设 (6) 得到。类似地，
\begin{align*}
(\text{d} \circ h + h \circ \text{d})(f \text{d}z \wedge \eta)
& =
\text{d}(\theta(f) \eta) +
h\left(\text{d}(f) \wedge \text{d}z \wedge \eta -
f \text{d}z \wedge \nabla(\eta)\right)
\\
& =
\text{d}(\theta(f)) \wedge \eta + \theta(f) \nabla(\eta)
- (\theta \otimes 1)(\text{d}_1(f)) \wedge \eta
- \theta(f) \nabla(\eta)
\\
& =
\text{d}_1(\theta(f)) \wedge \eta +
\partial_z(\theta(f)) \text{d}z \wedge \eta -
(\theta \otimes 1)(\text{d}_1(f)) \wedge \eta
\\
& =
a f \text{d}z \wedge \eta - \epsilon'(f) \wedge \eta \\
& = a f \text{d}z \wedge \eta - c^i_M(e^i_M(f \text{d}z \wedge \eta))
\end{align*}
第二个等号成立是因为
$\text{d}(f) \wedge \text{d}z \wedge \eta =
- \text{d}z \wedge \text{d}_1(f) \wedge \eta$.
第四个等号由假设 (\ref{crystalline-item-integrate}) 得到。另一方面，由定义立即有
$e^i_M(c^i_M(\omega))=\epsilon(1)\omega=a\omega$。引理得证。
\end{proof}

\begin{example}
\label{crystalline-example-integrate}
上述情形的一个标准例子如下：$B'=B\langle z\rangle$ 是除幂环
$(B,J,\delta)$ 上的除幂多项式环，而
$J'=B'_++JB'\subset B'$ 带除幂 $\delta'$。具体地，取
$\Omega=\Omega_{B,\delta}$ 与 $\Omega'=\Omega_{B',\delta'}$。
此时可取 $a=1$ 以及
$$
\theta( \sum b_m z^{[m]} ) = \sum b_m z^{[m + 1]}
$$
注意
$$
f - \theta(\partial_z(f)) = f(0)
$$
等于常数项。因此在此情形下，引理 \ref{crystalline-lemma-find-homotopy}
恢复晶体 Poincar\'e 引理（引理 \ref{crystalline-lemma-relative-poincare}）。
\end{example}

\begin{lemma}
\label{crystalline-lemma-computation}
处于 情形 \ref{crystalline-situation-affine}。假设 $D$ 与 $\Omega_D$ 如
(\ref{crystalline-equation-D}) 和 (\ref{crystalline-equation-omega-D})。设 $\lambda\in D$，并令
$D'$ 为下列环的 $p$-进完备化：
$$
D[z]\langle \xi \rangle/(\xi - (z^p - \lambda))
$$
；并令 $\Omega_{D'}$ 为 $D'$ 在 $A$ 上的除幂微分模之 $p$-进完备化。
对任意满足 (\ref{crystalline-item-complete})、(\ref{crystalline-item-connection})、
(\ref{crystalline-item-integrable}) 与 (\ref{crystalline-item-topologically-quasi-nilpotent})
的 $D$ 上偶对 $(M,\nabla)$，复形的典范映射
(\ref{crystalline-equation-base-change-map-complexes})
$$
c_M^\bullet : M \otimes_D^\wedge \Omega^\bullet_D
\longrightarrow
M \otimes_D^\wedge \Omega^\bullet_{D'}
$$
具有如下性质：存在反向映射 $e_M^\bullet$，使
$c_M^\bullet\circ e_M^\bullet$ 与 $e_M^\bullet\circ c_M^\bullet$
都同伦于乘以 $p$ 的映射。
\end{lemma}

\begin{proof}
取 $a=p$，使用 引理 \ref{crystalline-lemma-find-homotopy} 来证明。故须找出
$\theta:D'\to D'$ 并证明
(\ref{crystalline-item-d-a-zero}), (\ref{crystalline-item-direct-sum}), (\ref{crystalline-item-theta-linear}),
(\ref{crystalline-item-integrate}), (\ref{crystalline-item-injective}), (\ref{crystalline-item-factor}),
(\ref{crystalline-item-horizontal})。先汇集关于环 $D,D'$ 与模
$\Omega_D,\Omega_{D'}$ 的若干信息。

\medskip\noindent
写成
$$
D[z]\langle \xi \rangle/(\xi - (z^p - \lambda))
=
D\langle \xi \rangle[z]/(z^p - \xi - \lambda)
$$
可见 $D'$ 是下列自由 $D$-模的 $p$-进完备化：
$$
\bigoplus\nolimits_{i = 0, \ldots, p - 1}
\bigoplus\nolimits_{n \geq 0}
z^i \xi^{[n]} D
$$
，其中 $\xi^{[0]}=1$。因此 $D\to D'$ 有连续的 $D$-线性截面；特别地，
$D\to D'$ 泛单射，即 (\ref{crystalline-item-injective}) 成立。把 $D'$ 看作 $A$ 上
具有除幂理想 $\overline J'=\overline JD'+(\xi)$ 的除幂代数。于是 $D'$
也是 $D$ 中由 $z^p-\lambda$ 生成的理想之除幂包络的 $p$-进完备化，
见引理 \ref{crystalline-lemma-describe-divided-power-envelope}。故
$$
\Omega_{D'} = \Omega_D \otimes_D^\wedge D' \oplus D'\text{d}z
$$
，由引理 \ref{crystalline-lemma-module-differentials-divided-power-envelope} 得到。
这证明了 (\ref{crystalline-item-direct-sum})。注意 (\ref{crystalline-item-d-a-zero}) 显然成立。

\medskip\noindent
现在构造 $\theta$。（我们编写了一个 PARI/gp 脚本 theta.gp 来检验本证明
中的若干公式；该脚本位于 Stacks project 的 scripts 子目录中。）在此
之前先计算元素 $z^i\xi^{[n]}$ 的导数。有
$\text d z^i=iz^{i-1}\text d z$。当 $n\geq1$ 时，
$$
\text{d}\xi^{[n]} =
\xi^{[n - 1]} \text{d}\xi =
- \xi^{[n - 1]}\text{d}\lambda + p z^{p - 1} \xi^{[n - 1]}\text{d}z
$$
，因为 $\xi=z^p-\lambda$。当 $0<i<p$ 且 $n\geq1$ 时，有
\begin{align*}
\text{d}(z^i\xi^{[n]})
& =
iz^{i - 1}\xi^{[n]}\text{d}z + z^i\xi^{[n - 1]}\text{d}\xi \\
& =
iz^{i - 1}\xi^{[n]}\text{d}z + z^i\xi^{[n - 1]}\text{d}(z^p - \lambda) \\
& =
- z^i\xi^{[n - 1]}\text{d}\lambda +
(iz^{i - 1}\xi^{[n]} + pz^{i + p - 1}\xi^{[n - 1]})\text{d}z \\
& =
- z^i\xi^{[n - 1]}\text{d}\lambda +
(iz^{i - 1}\xi^{[n]} + pz^{i - 1}(\xi + \lambda)\xi^{[n - 1]})\text{d}z \\
& =
- z^i\xi^{[n - 1]}\text{d}\lambda +
((i + pn)z^{i - 1}\xi^{[n]} + p\lambda z^{i - 1}\xi^{[n - 1]})\text{d}z
\end{align*}
最后一个等号来自 $\xi\xi^{[n-1]}=n\xi^{[n]}$。因此
\begin{align*}
\partial_z(z^i) & = i z^{i - 1} \\
\partial_z(\xi^{[n]}) & = p z^{p - 1} \xi^{[n - 1]} \\
\partial_z(z^i\xi^{[n]}) & =
(i + pn) z^{i - 1} \xi^{[n]} + p \lambda z^{i - 1}\xi^{[n - 1]}
\end{align*}
受这些公式启发，按以下规则定义 $\theta$：
$$
\begin{matrix}
\theta(z^j)
& = & p\frac{z^{j + 1}}{j + 1}
& j = 0, \ldots p - 1, \\
\theta(z^{p - 1}\xi^{[m]})
& = & \xi^{[m + 1]}
& m \geq 1, \\
\theta(z^j \xi^{[m]})
& = &
\frac{p z^{j + 1} \xi^{[m]} - \theta(p\lambda z^j \xi^{[m - 1]})}{(j + 1 + pm)}
& 0 \leq j < p - 1, m \geq 1
\end{matrix}
$$
在最后一行中，对 $m$ 作归纳来定义所选的 $\theta$。展开可得（当
$0\leq j<p-1$ 且 $1\leq m$ 时）
$$
\theta(z^j \xi^{[m]}) =
\textstyle{\frac{p z^{j + 1} \xi^{[m]}}{(j + 1 + pm)} -
\frac{p^2 \lambda z^{j + 1} \xi^{[m - 1]}}{(j + 1 + pm)(j + 1 + p(m - 1))} +
\ldots +
\frac{(-1)^m p^{m + 1} \lambda^m z^{j + 1}}
{(j + 1 + pm) \ldots (j + 1)}}
$$
，不过下文不会使用此表达式。显然，$\theta$ 唯一延拓为 $D'$ 上
$p$-进连续的 $D$-线性映射。由构造，(\ref{crystalline-item-theta-linear}) 与
(\ref{crystalline-item-integrate}) 成立。尚须证明 (\ref{crystalline-item-factor}) 与
(\ref{crystalline-item-horizontal})。

\medskip\noindent
证明 (\ref{crystalline-item-factor}) 与 (\ref{crystalline-item-horizontal})。由于 $\theta$ 是
$D$-线性且连续的，只需证明
$p - \theta \circ \partial_z$ 与
$(\theta \otimes 1) \circ \text{d}_1 - \text{d}_1 \circ \theta$
分别作用于元素 $z^i\xi^{[n]}$ 时，所得分别属于 $D$ 与
$\Omega_D$\footnote{可直接计算：$p-\theta\circ\partial_z$ 作用于
$z^i\xi^{[n]}$ 时，除 $1$ 映到 $p$、$\xi$ 映到 $-p\lambda$ 外均为零。
而 $(\theta\otimes1)\circ\text d_1-\text d_1\circ\theta$ 作用于
$z^i\xi^{[n]}$ 时，除 $z^{p-1}\xi$ 映到 $-\lambda$ 外均为零。}。
置 $D_0=\mathbf Z_{(p)}[\lambda]$ 与
$D_0'=\mathbf Z_{(p)}[z,\lambda]\langle\xi\rangle/(\xi-z^p+\lambda)$。
注意上述每个表达式都是 $D_0'$ 或 $\Omega_{D_0'}$ 的元素。因此只需在
$D_0\to D_0'$ 的情形证明。注意 $D_0,D_0'$ 都是无挠环，且
$D_0\otimes\mathbf Q=\mathbf Q[\lambda]$、
$D'_0\otimes\mathbf Q=\mathbf Q[z,\lambda]$。故 $D_0\subset D'_0$ 是由
$\partial_z$ 消去的元素组成的子环，而 (\ref{crystalline-item-factor}) 由
(\ref{crystalline-item-integrate}) 推出；见紧邻 引理
\ref{crystalline-lemma-find-homotopy} 之前的讨论。类似地，
$\text{d}_1(f) = \partial_\lambda(f)\text{d}\lambda$，故
$$
\left((\theta \otimes 1) \circ \text{d}_1 - \text{d}_1 \circ \theta\right)(f)
=
\left(\theta(\partial_\lambda(f)) - \partial_\lambda(\theta(f))\right)
\text{d}\lambda
$$
对系数施加 $\partial_z$，得到
\begin{align*}
\partial_z\left(
\theta(\partial_\lambda(f)) - \partial_\lambda(\theta(f))
\right)
& =
p \partial_\lambda(f) - \partial_z(\partial_\lambda(\theta(f))) \\
& =
p \partial_\lambda(f) - \partial_\lambda(\partial_z(\theta(f))) \\
& =
p \partial_\lambda(f) - \partial_\lambda(p f) = 0
\end{align*}
故系数如所要求不依赖于 $z$。引理证毕。
\end{proof}

\noindent
注意，迭代 $\alpha_p$-覆盖 $X'\to X$（定义见本节引言）是有限局部自由的。
因此若 $X$ 连通，则 $X'\to X$ 的次数恒定且为 $p$ 的幂。

\begin{lemma}
\label{crystalline-lemma-pullback-along-p-power-cover}
设 $p$ 为素数，并设 $(S,\mathcal I,\gamma)$ 为 $\mathbf Z_{(p)}$ 上满足
$p\in\mathcal I$ 的除幂概形。置 $S_0=V(\mathcal I)\subset S$。设
$f:X'\to X$ 为 $S_0$ 上概形之间次数恒为 $q$ 的迭代
$\alpha_p$-覆盖。设 $\mathcal F$ 为 $X$ 上拟凝聚层中的任意晶体，并置
$\mathcal F'=f_{\text{cris}}^*\mathcal F$。在可区别三角
$$
Ru_{X/S, *}\mathcal{F}
\longrightarrow
f_*Ru_{X'/S, *}\mathcal{F}'
\longrightarrow
E
\longrightarrow
Ru_{X/S, *}\mathcal{F}[1]
$$
中，对象 $E$ 的各上同调层均被 $q$ 消去。
\end{lemma}

\begin{proof}
注意 $X'\to X$ 是同胚，故可等同 $X$ 与 $X'$ 的底拓扑空间。问题显然
在 $X$ 上是局部的，因此可假设 $X,X',S$ 均仿射，且 $X'\to X$ 写成复合
$$
X' = X_n \to X_{n - 1} \to X_{n - 2} \to \ldots \to X_0 = X
$$
，其中每个态射 $X_{i+1}\to X_i$ 都是 $\alpha_p$-覆盖。以
$\mathcal F_i$ 表示 $\mathcal F$ 到 $X_i$ 的拉回。只需证明每个映射
$$
R\Gamma(\text{Cris}(X_i/S), \mathcal{F}_i)
\longrightarrow
R\Gamma(\text{Cris}(X_{i + 1}/S), \mathcal{F}_{i + 1})
$$
都嵌入一个三角，其第三项的各上同调群被 $p$ 消去。（这里使用三角范畴
$D(X)$ 的公理 TR4；略去细节。）

\medskip\noindent
因此可假设 $S=\Spec(A)$、$X=\Spec(C)$、$X'=\Spec(C')$，且对某个
$c\in C$ 有 $C'=C[z]/(z^p-c)$。选取 $A$ 上多项式代数 $P$ 及满射
$P\to C$。令 $D$ 为 $P$ 中 $\Ker(P\to C)$ 的除幂包络之 $p$-进完备化，
如 (\ref{crystalline-equation-D})。置 $P'=P[z]$，并取把 $z$ 映到 $C'$ 中 $z$ 的类
的满射 $P'\to C'$。选取 $c\in C$ 的提升 $\lambda\in D$。于是，$P'$ 中
$\Ker(P'\to C')$ 的除幂包络之 $p$-进完备化 $D'$ 同构于下列环的
$p$-进完备化：
$D[z]\langle \xi \rangle/(\xi - (z^p - \lambda))$, see
，见引理 \ref{crystalline-lemma-computation} 及其证明。因此，利用命题
\ref{crystalline-proposition-compute-cohomology-crystal} 中拟凝聚模中晶体的上同调
计算，本引理立即推出结论。
\end{proof}

\noindent
下述引理中的界大概并非最优。

\begin{lemma}
\label{crystalline-lemma-pullback-along-p-power-cover-cohomology}
采用 引理 \ref{crystalline-lemma-pullback-along-p-power-cover} 的记号与假设，则映射
$$
f^* :
H^i(\text{Cris}(X/S), \mathcal{F})
\longrightarrow
H^i(\text{Cris}(X'/S), \mathcal{F}')
$$
的核与余核均被 $q^{i+1}$ 消去。
\end{lemma}

\begin{proof}
对象 $E$ 仅在次数 $-1$ 及以上可能有非零上同调层，故谱序列
$H^a(\mathcal H^b(E))\Rightarrow H^{a+b}(E)$ 收敛。把这一事实与可区别
三角所伴的上同调长正合列结合，即得所述界。
\end{proof}

\noindent
在 情形 \ref{crystalline-situation-global} 中，假设 $p\in\mathcal I$。置
$$
X^{(1)} = X \times_{S_0, F_{S_0}} S_0.
$$
以 $F_{X/S_0}:X\to X^{(1)}$ 表示相对 Frobenius 态射。

\begin{lemma}
\label{crystalline-lemma-pullback-relative-frobenius}
在上述情形中，假设 $X\to S_0$ 光滑且相对维数为 $d$。则 $F_{X/S_0}$
是次数 $p^d$ 的迭代 $\alpha_p$-覆盖。因此引理
\ref{crystalline-lemma-pullback-along-p-power-cover} 与
\ref{crystalline-lemma-pullback-along-p-power-cover-cohomology} 适用于此情形。
特别地，对 $\text{Cris}(X^{(1)}/S)$ 上拟凝聚模中的任意晶体
$\mathcal G$，映射
$$
F_{X/S_0}^* : H^i(\text{Cris}(X^{(1)}/S), \mathcal{G})
\longrightarrow
H^i(\text{Cris}(X/S), F_{X/S_0, \text{cris}}^*\mathcal{G})
$$
的核与余核均被 $p^{d(i+1)}$ 消去。
\end{lemma}

\begin{proof}
只需证明第一条断言。为此可假设 $X$ 在 $\mathbf A^d_{S_0}$ 上为
\'etale，见 “态射” 引理
\ref{morphisms-lemma-smooth-etale-over-affine-space}。以
$\varphi:X\to\mathbf A^d_{S_0}$ 表示这一 \'etale 态射。此时
$X/S_0$ 的相对 Frobenius 嵌入交换图
$$
\xymatrix{
X \ar[d] \ar[r] & X^{(1)} \ar[d] \\
\mathbf{A}^d_{S_0} \ar[r] & \mathbf{A}^d_{S_0}
}
$$
，其中下方横箭头是 $\mathbf A^d_{S_0}$ 在 $S_0$ 上的相对 Frobenius
态射。它把所有坐标提升到第 $p$ 次幂，故为迭代 $\alpha_p$-覆盖。
最后注意该图为纤维方块；见 \'Etale “态射” 引理
\ref{etale-lemma-relative-frobenius-etale}。
\end{proof}













\section{Frobenius 在晶体上同调上的作用}
\label{crystalline-section-frobenius}

\noindent
本节证明，在逆转素数 $p$ 后，Frobenius 拉回在晶体上同调上诱导拟同构。
但即使只是表述此结论，也需要处于一个特殊情形。

\begin{situation}
\label{crystalline-situation-F-crystal}
在 情形 \ref{crystalline-situation-global} 中，假设：
\begin{enumerate}
\item 对某个满足 $p\in I$ 的除幂环 $(A,I,\gamma)$，有
$S=\Spec(A)$；
\item 给定除幂环同态 $\sigma:A\to A$，使得对所有 $x\in A$，
$\sigma(x)=x^p\bmod pA$。
\end{enumerate}
\end{situation}

\noindent
在 情形 \ref{crystalline-situation-F-crystal} 中，态射
$\Spec(\sigma):S\to S$ 是绝对 Frobenius $F_{S_0}:S_0\to S_0$ 的提升，
并且图
$$
\xymatrix{
X \ar[d] \ar[r]_{F_X} & X \ar[d] \\
S_0 \ar[r]^{F_{S_0}} & S_0
}
$$
交换，其中 $F_X:X\to X$ 是 $X$ 的绝对 Frobenius 态射。因此得到晶体
拓扑斯的态射
$$
(F_X)_{\text{cris}} :
(X/S)_{\text{cris}}
\longrightarrow
(X/S)_{\text{cris}}
$$
，见注 \ref{crystalline-remark-functoriality-cris}。下面依照 Saavedra 的记号
给出关于 $F$-晶体的术语，见 \cite{Saavedra}。

\begin{definition}
\label{crystalline-definition-F-crystal}
在 情形 \ref{crystalline-situation-F-crystal} 中，所谓{\it $X/S$ 上（相对于
$\sigma$）的 $F$-晶体}，是指偶对 $(\mathcal E,F_\mathcal E)$，它由有限
局部自由 $\mathcal O_{X/S}$-模中的晶体 $\mathcal E$ 与映射
$$
F_\mathcal{E} : (F_X)_{\text{cris}}^*\mathcal{E} \longrightarrow \mathcal{E}
$$
组成。若存在整数 $i\geq0$ 及映射
$V:\mathcal E\to(F_X)_{\text{cris}}^*\mathcal E$，使
$V\circ F_\mathcal E=p^i\text{id}$，则称该 $F$-晶体{\it 非退化}。
\end{definition}

\begin{remark}
\label{crystalline-remark-F-crystal-variants}
设 $(\mathcal E,F)$ 为 定义 \ref{crystalline-definition-F-crystal} 中的
$F$-晶体。文献中常把非退化条件纳入 $F$-晶体的定义，并且还常假设
$F\circ V=p^n\text{id}$。下述结果所需的是：存在整数 $j\geq0$，使
$\Ker(F)$ 与 $\Coker(F)$ 均被 $p^j$ 消去。若 $\mathcal E$ 的秩有界
（例如当 $X$ 拟紧时），则这两个条件都由定义中所表述的非退化条件
推出。确切地说，设 $R$ 为环，$r\geq1$ 为整数，并设
$K,L\in\text{Mat}(r\times r,R)$ 满足 $KL=p^i1_{r\times r}$。则
$\det(K)\det(L)=p^{ri}$。令 $L'$ 为 $L$ 的伴随矩阵，即
$L'L=LL'=\det(L)$。置 $K'=p^{ri}K$ 与 $j=ri+i$。由于 $KL=p^i$，有
$K'L=p^j1_{r\times r}$，并且
$$
L K' = L K \det(L) \det(M) = L K L L' \det(M) = L p^i L' \det(M) =
p^j 1_{r \times r}
$$
因此，若 $V$ 如 定义 \ref{crystalline-definition-F-crystal}，则取
$V'=p^NV$，其中 $N>i\cdot\text{rank}(\mathcal E)$，便有
$V'\circ F=p^{N+i}$ 及 $F\circ V'=p^{N+i}$。
\end{remark}

\begin{theorem}
\label{crystalline-theorem-cohomology-F-crystal}
在 情形 \ref{crystalline-situation-F-crystal} 中，设
$(\mathcal E,F_\mathcal E)$ 为非退化 $F$-晶体。假设 $A$ 为 $p$-进完备
Noether 环，且 $X\to S_0$ 固有光滑。则典范映射
$$
F_\mathcal{E} \circ (F_X)_{\text{cris}}^* :
R\Gamma(\text{Cris}(X/S), \mathcal{E}) \otimes^\mathbf{L}_{A, \sigma} A
\longrightarrow
R\Gamma(\text{Cris}(X/S), \mathcal{E})
$$
在逆转 $p$ 后成为同构。
\end{theorem}

\begin{proof}
先把该箭头写成三个箭头的复合。置
$$
X^{(1)} = X \times_{S_0, F_{S_0}} S_0
$$
并以 $F_{X/S_0}:X\to X^{(1)}$ 表示相对 Frobenius 态射。以
$\mathcal E^{(1)}$ 表示 $\mathcal E$ 沿 $\Spec(\sigma)$ 的基变换；换言之，
它是经由与交换图
$$
\xymatrix{
X^{(1)} \ar[r] \ar[d] & X \ar[d] \\
S \ar[r]^{\Spec(\sigma)} & S
}
$$
相伴的晶体拓扑斯态射，把 $\mathcal E$ 拉回到
$\text{Cris}(X^{(1)}/S)$ 所得的晶体。于是有基变换映射
\begin{equation}
\label{crystalline-equation-base-change-sigma}
R\Gamma(\text{Cris}(X/S), \mathcal{E}) \otimes^\mathbf{L}_{A, \sigma} A
\longrightarrow
R\Gamma(\text{Cris}(X^{(1)}/S), \mathcal{E}^{(1)})
\end{equation}
，见注 \ref{crystalline-remark-base-change}。注意
$F_{X/S_0}:X\to X^{(1)}$ 与投影 $X^{(1)}\to X$ 的复合是绝对
Frobenius 态射 $F_X$。故
$F_{X/S_0}^*\mathcal E^{(1)}=(F_X)_{\text{cris}}^*\mathcal E$。因此沿
$F_{X/S_0}$ 的拉回给出映射
\begin{equation}
\label{crystalline-equation-to-prove}
F_{X/S_0}^* :
R\Gamma(\text{Cris}(X^{(1)}/S), \mathcal{E}^{(1)})
\longrightarrow
R\Gamma(\text{Cris}(X/S), (F_X)^*_{\text{cris}}\mathcal{E})
\end{equation}
最后可用 $F_\mathcal E$ 得到映射
\begin{equation}
\label{crystalline-equation-F-E}
R\Gamma(\text{Cris}(X/S), (F_X)^*_{\text{cris}}\mathcal{E})
\longrightarrow
R\Gamma(\text{Cris}(X/S), \mathcal{E})
\end{equation}
定理中的映射是上述三个映射 (\ref{crystalline-equation-base-change-sigma})、
(\ref{crystalline-equation-to-prove}) 与 (\ref{crystalline-equation-F-E}) 的复合。由注
\ref{crystalline-remark-base-change-isomorphism}，第一个映射模 $p$ 的任意幂都是
拟同构。由于所涉复形在 $D(A)$ 中完美，它本身即为拟同构，见注
\ref{crystalline-remark-complete-perfect}。第三个映射在逆转 $p$ 后为拟同构，因为
$F_\mathcal E$ 有一个相差 $p$ 的某次幂的逆，见注
\ref{crystalline-remark-F-crystal-variants}。最后，由引理
\ref{crystalline-lemma-pullback-relative-frobenius}，第二个映射在逆转 $p$ 后为同构。
\end{proof}
% Generated mechanically from the frozen proetale.tex target.
% Do not edit this generated file; edit the bound locale source instead.

% 好，从这里开始。
%

\setcounter{chapter}{60}
\chapter{Pro-\'etale 上同调}



\phantomsection
\label{proetale-section-phantom}


\section{引言}
\label{proetale-section-introduction}

\noindent
本章及更多相关内容可见预印本 \cite{BS}。

\medskip\noindent
本章的目标是介绍 pro-\'etale 拓扑，并发展该拓扑下阿贝尔层上同调的
基本理论。另一个目标是说明，使用 pro-\'etale 拓扑可以简化代数几何中
$\ell$-进上同调的引入。

\medskip\noindent
下面简要回顾我们所理解的 $\ell$-进 \'etale 上同调史。
在 \cite[报告 V 与 VI]{SGA5} 中，Grothendieck 等人发展出一套理论，
把 $\ell$-进层处理为 $\mathbf{Z}/\ell^n\mathbf{Z}$-模层的逆系统。
Deligne 在关于 Weil 猜想的第二篇论文（\cite{WeilII}）中，把
$\ell$-进层的导出范畴引入为若干范畴的一个 2-极限；这些范畴由概形
$X$ 的 \'etale 位点上 $\mathbf{Z}/\ell^n\mathbf{Z}$-模层的复形组成。
Beilinson、Bernstein 与 Deligne 的论文（\cite{BBD}）以这种方法为
基础，建立了优美的变态层理论。Uwe Jannsen 在题为“连续 \'Etale
上同调”的论文（\cite{Jannsen}）中，讨论了域上簇的一个 $\ell$-进层
上同调的重要变体。随后 Torsten Ekedahl 的论文（\cite{Ekedahl}）
讨论了便利地处理定义为极限的导出范畴所需的进形式体系。

\medskip\noindent
事实表明，在概形的 pro-\'etale 位点上工作，可以避开这些作者遇到的
一部分技术困难。代价是必须处理非 Noether 概形；即使我们只关心
代数闭域上簇的 $\ell$-进层及其上同调，也仍然如此。

\medskip\noindent
概形的（小）pro-\'etale 位点具有足够多的拟紧 w-可缩对象，这是一个
非常重要而显著的特征。这些对象的存在，对阿贝尔层的导出范畴以及层的
逆系统产生若干有用而（或许）不寻常的推论。正是这一特征使我们能够
处理 $\ell$-进层所涉及的种种微妙之处；不过，正如后文将看到的，
它还有不少其他益处。




\section{若干拓扑学}
\label{proetale-section-topology}

\noindent
先作若干准备。我们已在“拓扑学”第
\ref{topology-section-spectral} 节定义了{\it 谱空间}及谱空间之间的
{\it 谱映射}。环的谱是谱空间；见“代数”引理
\ref{algebra-lemma-spec-spectral}。

\begin{lemma}
\label{proetale-lemma-spectral-split}
设 $X$ 为谱空间，并设 $X_0\subset X$ 为闭点集。下列条件等价：
\begin{enumerate}
\item $X$ 的每个开覆盖都可由有限不交并分解
$X=\coprod U_i$ 加细，其中各 $U_i$ 在 $X$ 中既开又闭。
\item 复合映射 $X_0\to X\to\pi_0(X)$ 为双射。
\end{enumerate}
此外，若 $X_0$ 在 $X$ 中闭，且 $X$ 的每个点都特化到 $X_0$ 中唯一的
一点，则上述条件成立。
\end{lemma}

\begin{proof}
以下将直接使用：$X_0$ 拟紧（“拓扑学”引理
\ref{topology-lemma-closed-points-quasi-compact}），且 $\pi_0(X)$
是预有限空间（“拓扑学”引理 \ref{topology-lemma-spectral-pi0}）。
考虑图
$$
\xymatrix{
X_0 \ar[rd]_f \ar[r] & X \ar[d]^\pi \\
& \pi_0(X)
}
$$
若 (2) 成立，则由“拓扑学”引理 \ref{topology-lemma-bijective-map}，
连续双射 $f:X_0\to\pi_0(X)$ 是同胚。给定开覆盖
$X=\bigcup U_i$，便得到开覆盖
$\pi_0(X)=\bigcup f(X_0\cap U_i)$。由“拓扑学”引理
\ref{topology-lemma-profinite-refine-open-covering}，可找到形如
$\pi_0(X)=\coprod V_j$ 的有限开覆盖来加细它。由于
$X_0\to\pi_0(X)$ 为双射，$X$ 的每个连通分支都有唯一闭点，因而恰为
特化到该闭点的所有点组成的集合。因此，$\pi^{-1}(V_j)$ 是特化到
$f^{-1}(V_j)$ 中各点的点集。若
$f^{-1}(V_j)\subset X_0\cap U_i\subset U_i$，则
$\pi^{-1}(V_j)\subset U_i$（因为开集 $U_i$ 在泛化下封闭）。由此，
开覆盖 $X=\coprod\pi^{-1}(V_j)$ 加细了原来的覆盖。这证明 (2)
蕴涵 (1)。

\medskip\noindent
假设 (1)。设 $x,y\in X$ 为闭点。于是有开覆盖
$X=(X\setminus\{x\})\cup(X\setminus\{y\})$。由 (1)，存在不交并分解
$X=U\amalg V$，其中 $U$ 与 $V$ 均开（亦均闭），且 $x\in U$、$y\in V$。
特别地，$X$ 的每个连通分支至多有一个闭点。由“拓扑学”引理
\ref{topology-lemma-quasi-compact-closed-point}，每个连通分支
（它是闭的）也确有一个闭点。因此 $X_0\to\pi_0(X)$ 为双射，
这证明 (1) 蕴涵 (2)。

\medskip\noindent
假设 $X_0$ 在 $X$ 中闭，且每个点都特化到 $X_0$ 中唯一的一点。
则 $X_0$ 是由闭点组成的谱空间（“拓扑学”引理
\ref{topology-lemma-spectral-sub}），故为预有限空间（“拓扑学”引理
\ref{topology-lemma-characterize-profinite-spectral}）。设
$x,y\in X_0$ 互异。由“拓扑学”引理
\ref{topology-lemma-profinite-refine-open-covering}，可找到不交并分解
$X_0=U_0\amalg V_0$，其中 $U_0$ 与 $V_0$ 均开且闭，并且
$x\in U_0$、$y\in V_0$。令 $U\subset X$（相应地 $V\subset X$）
为特化到 $U_0$（相应地 $V_0$）的点集。注意
$X=U\amalg V$。由“拓扑学”引理
\ref{topology-lemma-make-spectral-space}，$U$ 是拟紧开子集的交，
故 $U$ 在可构造拓扑中闭。又因 $U$ 在特化下封闭，由“拓扑学”引理
\ref{topology-lemma-constructible-stable-specialization-closed}，
$U$ 闭。由对称性，$V$ 也闭，故 $U$ 与 $V$ 均开且闭。这证明 $x,y$
不在 $X$ 的同一连通分支中。换言之，$X_0\to\pi_0(X)$ 为单射。
由“拓扑学”引理 \ref{topology-lemma-quasi-compact-closed-point}
以及连通分支为闭集这一事实，该映射也满。于是最后一个条件蕴涵 (2)。
\end{proof}

\begin{example}
\label{proetale-example-not-w-local}
设 $T$ 为预有限空间。取一点 $t\in T$，并假设
$T\setminus\{t\}$ 不拟紧。令 $X=T\times\{0,1\}$。在 $X$ 上考虑
由下列集合给出子基的拓扑：当 $U\subset T$ 开时的
$U\times\{0,1\}$，集合 $X\setminus\{(t,0)\}$，以及当
$U\subset T$ 开且 $t\not \in U$ 时的 $U\times\{1\}$。
$X$ 的闭点集为 $X_0=T\times\{0\}$，而 $(t,1)$ 属于 $X_0$ 的闭包。
此外，$X_0\to\pi_0(X)$ 是双射。此例说明，引理
\ref{proetale-lemma-spectral-split} 的条件 (1) 与 (2) 并不蕴涵闭点集为闭集。
\end{example}

\noindent
事实表明，使用满足引理 \ref{proetale-lemma-spectral-split} 末句所述稍强性质的
谱空间更为方便。我们给这一性质命名。

\begin{definition}
\label{proetale-definition-w-local}
若谱空间 $X$ 的闭点集 $X_0$ 为闭集，且 $X$ 的每个点都特化到唯一的
闭点，则称该空间为{\it w-局部的}。w-局部空间之间的连续映射
$f:X\to Y$ 称为{\it w-局部的}，如果它是谱映射，并把 $X$ 的每个
闭点映到 $Y$ 的闭点。
\end{definition}

\noindent
在引理 \ref{proetale-lemma-spectral-split} 的证明中已经看到，此时
$X_0\to\pi_0(X)$ 是同胚，且 $X_0\cong\pi_0(X)$ 为预有限空间。
此外，$X$ 的一个连通分支恰是特化到某个给定 $x\in X_0$ 的所有点
组成的集合。

\begin{lemma}
\label{proetale-lemma-closed-subspace-w-local}
设 $X$ 为 w-局部谱空间。若 $Y\subset X$ 闭，则 $Y$ 为 w-局部的。
\end{lemma}

\begin{proof}
闭点子集 $Y_0\subset Y$ 为闭集，因为 $Y_0=X_0\cap Y$。由于 $X$
是 $w$-局部的，每个 $y\in Y$ 都特化到 $X_0$ 中唯一的一点。又因
$\overline{\{y\}}\subset Y$，该特化点属于 $Y$，从而也属于 $Y_0$。
故 $Y$ 为 $w$-局部的。
\end{proof}

\begin{lemma}
\label{proetale-lemma-silly}
设 $X$ 为谱空间，并设
$$
\xymatrix{
Y \ar[r] \ar[d] & T \ar[d] \\
X \ar[r] & \pi_0(X)
}
$$
为拓扑空间范畴中的笛卡尔图，其中 $T$ 为预有限空间。则 $Y$ 为
谱空间，且 $T=\pi_0(Y)$。若再假设 $X$ 为 w-局部的，则 $Y$ 为
w-局部的，$Y\to X$ 为 w-局部映射，并且 $Y$ 的闭点集是 $X$
的闭点集的逆像。
\end{lemma}

\begin{proof}
注意，$\pi_0(X)$ 是预有限空间，因而是 Hausdorff 的，所以 $Y$ 是
$X\times T$ 的闭子空间（使用“拓扑学”引理
\ref{topology-lemma-spectral-pi0} 与
\ref{topology-lemma-fibre-product-closed}）。又因 $X\times T$
为谱空间（“拓扑学”引理
\ref{topology-lemma-product-spectral-spaces}），故 $Y$ 为谱空间
（“拓扑学”引理 \ref{topology-lemma-spectral-sub}）。考虑典范分解
$Y\to\pi_0(Y)\to T$（“拓扑学”引理
\ref{topology-lemma-space-connected-components}）。显然
$\pi_0(Y)\to T$ 满射。$Y\to T$ 的各纤维同胚于
$X\to\pi_0(X)$ 的各纤维，因而连通。于是 $\pi_0(Y)\to T$ 单射。
由“拓扑学”引理 \ref{topology-lemma-bijective-map}，
$\pi_0(Y)\to T$ 是同胚。

\medskip\noindent
再假设 $X$ 为 w-局部的，并以 $X_0\subset X$ 表示闭点集。令
$Y_0\subset Y$ 为 $X_0$ 在 $Y$ 中的逆像。由引理
\ref{proetale-lemma-spectral-split}，$X_0\to\pi_0(X)$ 为双射，故
该逆像双射地映到 $T$。此外，$Y_0$ 是谱空间 $Y$ 的闭子集，因而拟紧。
于是由“拓扑学”引理 \ref{topology-lemma-bijective-map}，
$Y_0\to\pi_0(Y)=T$ 为同胚。因此 $Y_0$ 的所有点在 $Y$ 中均为闭点。
反之，若 $y\in Y$ 为闭点，则它在 $Y\to\pi_0(Y)=T$ 的相应纤维中闭，
从而其在 $X$ 中的像 $x$ 在 $X\to\pi_0(X)$ 的同胚纤维中闭。
这蕴涵 $x\in X_0$，故 $y\in Y_0$。所以 $Y_0$ 正是 $Y$ 的闭点集；
对每个 $y\in Y_0$，$y$ 的所有泛化点组成 $Y\to\pi_0(Y)$ 的相应
纤维。引理得证。
\end{proof}




\section{局部同构}
\label{proetale-section-local-isomorphism}

\noindent
先给出一个定义。

\begin{definition}
\label{proetale-definition-local-isomorphism}
设 $\varphi:A\to B$ 为环映射。
\begin{enumerate}
\item 若对每个素理想 $\mathfrak q\subset B$，都存在
$g\in B$、$g\not \in\mathfrak q$，使 $A\to B_g$ 诱导开浸入
$\Spec(B_g)\to\Spec(A)$，则称 $A\to B$ 为{\it 局部同构}。
\item 若对每个素理想 $\mathfrak q\subset B$，典范映射
$A_{\varphi^{-1}(\mathfrak q)}\to B_\mathfrak q$ 都是同构，则称
$A\to B$ {\it 诱导局部环同构}。
\end{enumerate}
\end{definition}

\noindent
列出若干初等性质。

\begin{lemma}
\label{proetale-lemma-base-change-local-isomorphism}
设 $A\to B$ 与 $A\to A'$ 为环映射，并令
$B'=B\otimes_AA'$ 为 $B$ 的基变换。
\begin{enumerate}
\item 若 $A\to B$ 为局部同构，则 $A'\to B'$ 为局部同构。
\item 若 $A\to B$ 诱导局部环同构，则 $A'\to B'$ 也诱导局部环同构。
\end{enumerate}
\end{lemma}

\begin{proof}
略。
\end{proof}

\begin{lemma}
\label{proetale-lemma-composition-local-isomorphism}
设 $A\to B$ 与 $B\to C$ 为环映射。
\begin{enumerate}
\item 若 $A\to B$ 与 $B\to C$ 均为局部同构，则 $A\to C$ 为局部同构。
\item 若 $A\to B$ 与 $B\to C$ 均诱导局部环同构，则 $A\to C$
也诱导局部环同构。
\end{enumerate}
\end{lemma}

\begin{proof}
略。
\end{proof}

\begin{lemma}
\label{proetale-lemma-local-isomorphism-permanence}
设 $A$ 为环，并设 $B\to C$ 为 $A$-代数同态。
\begin{enumerate}
\item 若 $A\to B$ 与 $A\to C$ 均为局部同构，则 $B\to C$ 为局部同构。
\item 若 $A\to B$ 与 $A\to C$ 均诱导局部环同构，则 $B\to C$
也诱导局部环同构。
\end{enumerate}
\end{lemma}

\begin{proof}
略。
\end{proof}

\begin{lemma}
\label{proetale-lemma-local-isomorphism-implies}
设 $A\to B$ 为局部同构。则
\begin{enumerate}
\item $A\to B$ 是 \'etale 的；
\item $A\to B$ 诱导局部环同构；
\item $A\to B$ 拟有限。
\end{enumerate}
\end{lemma}

\begin{proof}
略。
\end{proof}

\begin{lemma}
\label{proetale-lemma-structure-local-isomorphism}
设 $A\to B$ 为局部同构。则存在 $n\geq0$、
$g_1,\ldots,g_n\in B$ 与 $f_1,\ldots,f_n\in A$，使得
$(g_1,\ldots,g_n)=B$ 且 $A_{f_i}\cong B_{g_i}$。
\end{lemma}

\begin{proof}
略。
\end{proof}

\begin{lemma}
\label{proetale-lemma-fully-faithful-spaces-over-X}
设
$p:(Y,\mathcal{O}_Y)\to(X,\mathcal{O}_X)$ 与
$q:(Z,\mathcal{O}_Z)\to(X,\mathcal{O}_X)$ 为局部环空间的态射。
若 $\mathcal{O}_Y=p^{-1}\mathcal{O}_X$，则
$$
\Mor_{\text{LRS}/(X, \mathcal{O}_X)}((Z, \mathcal{O}_Z), (Y, \mathcal{O}_Y))
\longrightarrow
\Mor_{\textit{Top}/X}(Z, Y),\quad
(f, f^\sharp) \longmapsto f
$$
为双射。这里 $\text{LRS}/(X,\mathcal{O}_X)$ 是 $X$ 上局部环空间的
范畴，$\textit{Top}/X$ 是 $X$ 上拓扑空间的范畴。
\end{lemma}

\begin{proof}
由定义立即得到。
\end{proof}

\begin{lemma}
\label{proetale-lemma-local-isomorphism-fully-faithful}
设 $A$ 为环，并置 $X=\Spec(A)$。函子
$$
B \longmapsto \Spec(B)
$$
从满足 $A\to B$ 诱导局部环同构的 $A$-代数 $B$ 所成范畴，到
$X$ 上拓扑空间的范畴，是全忠实的。
\end{lemma}

\begin{proof}
这由引理 \ref{proetale-lemma-fully-faithful-spaces-over-X} 以及下述事实得到：
若 $A\to B$ 诱导局部环同构，则 $\Spec(A)$ 的结构层沿
$p:\Spec(B)\to\Spec(A)$ 的拉回等于 $\Spec(B)$ 的结构层。
\end{proof}




\section{Ind-Zariski 代数}
\label{proetale-section-ind-zariski}

\noindent
先给出一个定义；它与论文 \cite{BS} 中相应定义的比较见注
\ref{proetale-remark-slightly-stronger}。

\begin{definition}
\label{proetale-definition-ind-zariski}
若 $B$ 可写成滤过余极限 $B=\colim B_i$，且每个 $A\to B_i$ 都是
局部同构，则称环映射 $A\to B$ 为 {\it ind-Zariski} 的。
\end{definition}

\noindent
局部化 $A\to S^{-1}A$ 是 ind-Zariski 映射的一个例子；见“代数”引理
\ref{algebra-lemma-localization-colimit}。ind-Zariski 代数范畴对若干
自然运算封闭。

\begin{lemma}
\label{proetale-lemma-base-change-ind-zariski}
设 $A\to B$ 与 $A\to A'$ 为环映射，并令
$B'=B\otimes_AA'$ 为 $B$ 的基变换。若 $A\to B$ 为 ind-Zariski，
则 $A'\to B'$ 为 ind-Zariski。
\end{lemma}

\begin{proof}
略。
\end{proof}

\begin{lemma}
\label{proetale-lemma-composition-ind-zariski}
设 $A\to B$ 与 $B\to C$ 为环映射。若 $A\to B$ 与 $B\to C$
均为 ind-Zariski，则 $A\to C$ 为 ind-Zariski。
\end{lemma}

\begin{proof}
略。
\end{proof}

\begin{lemma}
\label{proetale-lemma-ind-zariski-permanence}
设 $A$ 为环，并设 $B\to C$ 为 $A$-代数同态。若 $A\to B$ 与
$A\to C$ 均为 ind-Zariski，则 $B\to C$ 为 ind-Zariski。
\end{lemma}

\begin{proof}
略。
\end{proof}

\begin{lemma}
\label{proetale-lemma-ind-ind-zariski}
ind-Zariski $A$-代数的滤过余极限在 $A$ 上仍为 ind-Zariski。
\end{lemma}

\begin{proof}
略。
\end{proof}

\begin{lemma}
\label{proetale-lemma-ind-zariski-implies}
设 $A\to B$ 为 ind-Zariski。则 $A\to B$ 诱导局部环同构。
\end{lemma}

\begin{proof}
略。
\end{proof}







\section{构造 w-局部仿射概形}
\label{proetale-section-construction}

\noindent
若仿射概形 $X$ 的底拓扑空间为 w-局部的（定义
\ref{proetale-definition-w-local}），则称该概形为 {\it w-局部的}。事实表明，
对任意环 $A$，都存在典范的忠实平坦 ind-Zariski 环映射
$A\to A_w$，使 $\Spec(A_w)$ 为 w-局部的。构造 $A_w$ 的关键是
下述简单引理。

\begin{lemma}
\label{proetale-lemma-localization}
设 $A$ 为环，置 $X=\Spec(A)$。设 $Z\subset X$ 为形如
$D(f)\cap V(I)$ 的局部闭子概形，其中 $f\in A$ 且 $I\subset A$ 为理想。则
\begin{enumerate}
\item 存在乘法子集 $S\subset A$，使得 $\Spec(S^{-1}A)$ 同胚地映到
$X$ 中特化到 $Z$ 的点集；
\item $A$-代数 $A_Z^\sim=S^{-1}A$ 仅依赖于底层局部闭子集
$Z\subset X$；
\item $Z$ 是 $\Spec(A_Z^\sim)$ 的闭子概形。
\end{enumerate}
若 $A\to A'$ 为环映射，且同样形状的局部闭子概形
$Z'\subset X'=\Spec(A')$ 映入 $Z$，则存在唯一的 $A$-代数映射
$A_Z^\sim\to(A')_{Z'}^\sim$。
\end{lemma}

\begin{proof}
令 $S\subset A$ 为所有在
$\Gamma(Z,\mathcal{O}_Z)=(A/I)_f$ 中映为可逆元的元素所成的乘法集。
若 $A$ 的素理想 $\mathfrak p$ 不特化到 $Z$，则 $\mathfrak p$ 在
$(A/I)_f$ 中生成单位理想。因此可对某个 $n\geq0$、
$g\in\mathfrak p$、$h\in I$ 写成 $f^n=g+h$。于是 $g\in S$，
故 $\mathfrak p$ 不属于 $S^{-1}A$ 的谱。反之，若 $\mathfrak p$
确实特化到 $Z$，即
$\mathfrak p\subset\mathfrak q\supset I$ 且
$f\not \in\mathfrak q$，则 $S^{-1}A$ 映到 $A_\mathfrak q$，所以
$\mathfrak p$ 属于 $S^{-1}A$ 的谱。这证明 (1)。

\medskip\noindent
局部化 $S^{-1}A$ 的同构类仅依赖于相应的子集
$\Spec(S^{-1}A)\subset\Spec(A)$，故 (2) 成立。由构造，
$S^{-1}A$ 满射到 $(A/I)_f$，故 (3) 成立。最后的断言则由下述事实
推出：与 $Z'$ 对应的乘法子集 $S'\subset A'$ 包含乘法子集 $S$ 的像。
\end{proof}

\noindent
设 $A$ 为环，并设 $E\subset A$ 为有限子集。考察 $E$ 中各元素的
消失情形，可得到 $X=\Spec(A)$ 到局部闭子概形的分层。更准确地，
对不交并分解 $E=E'\amalg E''$，置
\begin{equation}
\label{proetale-equation-stratum}
Z(E', E'') =
\bigcap\nolimits_{f \in E'} D(f) \cap \bigcap\nolimits_{f \in E''} V(f) =
D(\prod\nolimits_{f \in E'} f) \cap V( \sum\nolimits_{f \in E''} fA)
\end{equation}
$Z(E',E'')$ 的点恰为满足下述条件的 $x\in X$：每个 $f\in E'$ 在
$\kappa(x)$ 中的像非零，而每个 $f\in E''$ 在 $\kappa(x)$ 中的像
为零。因此，集合意义下显然有
\begin{equation}
\label{proetale-equation-stratify}
X = \coprod\nolimits_{E = E' \amalg E''} Z(E', E'')
\end{equation}
。注意，每个分层片都是可构造的。

\begin{lemma}
\label{proetale-lemma-refine}
设 $X=\Spec(A)$ 如上。给定由可构造子集组成的任意有限分层
$X=\coprod T_i$，存在有限子集 $E\subset A$，使分层
(\ref{proetale-equation-stratify}) 加细 $X=\coprod T_i$。
\end{lemma}

\begin{proof}
可写成有限并
$T_i=\bigcup_jU_{i,j}\cap V_{i,j}^c$，其中 $U_{i,j}$ 与 $V_{i,j}$ 是
$X$ 中的拟紧开集。再写成
$U_{i,j}=\bigcup D(f_{i,j,k})$ 与
$V_{i,j}=\bigcup D(g_{i,j,l})$。置
$E=\{f_{i,j,k}\}\cup\{g_{i,j,l}\}$ 即可，因为
(\ref{proetale-equation-stratify}) 的各分层片由 $E$ 中元素的消失型标记，
显然加细给定分层。
\end{proof}

\noindent
继续上述讨论。给定有限子集 $E\subset A$，置
\begin{equation}
\label{proetale-equation-ring}
A_E = \prod\nolimits_{E = E' \amalg E''} A_{Z(E', E'')}^\sim
\end{equation}
，其中沿用引理 \ref{proetale-lemma-localization} 的记号。此定义有意义，因为
(\ref{proetale-equation-stratum}) 表明每个 $Z(E',E'')$ 都具有所需的形状。
取此环的谱，记作
\begin{equation}
\label{proetale-equation-spectrum}
X_E = \Spec(A_E) = \coprod\nolimits_{E = E' \amalg E''} X_{E', E''}
\end{equation}
，其中 $X_{E',E''}=\Spec(A_{Z(E',E'')}^\sim)$。注意
\begin{equation}
\label{proetale-equation-closed}
Z_E = \coprod\nolimits_{E = E' \amalg E''} Z(E', E'')
\longrightarrow
X_E
\end{equation}
是闭子概形。由构造，闭子概形 $Z_E$ 包含仿射概形 $X_E$ 的所有闭点，
因为 $X_{E',E''}$ 的每个点都特化到 $Z(E',E'')$ 的某个点。

\medskip\noindent
令 $I(A)$ 为 $A$ 的所有有限子集按包含关系组成的偏序集。这是有向偏序集。
对 $E_1\subset E_2$，存在典范过渡映射
$A_{E_1}\to A_{E_2}$，它是 $A$-代数映射。事实上，给定分解
$E_2=E'_2\amalg E''_2$，置
$E'_1=E_1\cap E'_2$ 与 $E''_1=E_1\cap E''_2$。注意
$Z(E'_2,E''_2)\subset Z(E'_1,E''_1)$，故由引理
\ref{proetale-lemma-localization}，存在唯一的 $A$-代数映射
$A_{Z(E'_1,E''_1)}^\sim\to A_{Z(E'_2,E''_2)}^\sim$。把这些映射合在一起，
便得到所需的环映射 $A_{E_1}\to A_{E_2}$。注意，对应的仿射概形映射
\begin{equation}
\label{proetale-equation-transition}
X_{E_2} \longrightarrow X_{E_1}
\end{equation}
把 $Z_{E_2}$ 映入 $Z_{E_1}$。由唯一性，得到一个 $A$-代数系统，
其指标集为 $I(A)$，并置
\begin{equation}
\label{proetale-equation-colimit-ring}
A_w = \colim_{E \in I(A)} A_E
\end{equation}
这个 $A$-代数在 $A$ 上是 ind-Zariski 且忠实平坦的。最后，置
$X_w=\Spec(A_w)$，并令其带有闭子概形
$Z=\lim_{E\in I(A)}Z_E$。用公式表示为
\begin{equation}
\label{proetale-equation-final}
X_w = \lim_{E \in I(A)} X_E \supset Z = \lim_{E \in I(A)} Z_E
\end{equation}

\begin{lemma}
\label{proetale-lemma-make-w-local}
设 $X=\Spec(A)$ 为仿射概形，并采用如上的 $A\to A_w$、
$X_w=\Spec(A_w)$ 与 $Z\subset X_w$。则
\begin{enumerate}
\item $A\to A_w$ 是 ind-Zariski 且忠实平坦的；
\item $X_w\to X$ 诱导双射 $Z\to X$；
\item $Z$ 是 $X_w$ 的闭点集；
\item $Z$ 是约化概形；
\item $X_w$ 的每个点都特化到 $Z$ 的唯一一点。
\end{enumerate}
特别地，$X_w$ 是 w-局部的（定义 \ref{proetale-definition-w-local}）。
\end{lemma}

\begin{proof}
由构造，映射 $A\to A_w$ 是 ind-Zariski 的。对每个 $E$，态射
$Z_E\to X$ 都是双射，故 (2) 成立。由于 $Z\subset X_w$，可知
$X_w\to X$ 满射；再由“代数”引理 \ref{algebra-lemma-ff-rings}，
$A\to A_w$ 忠实平坦。这证明了 (1)。

\medskip\noindent
设 $y\in X_w$ 且 $y\not\in Z$。则存在 $E$，使 $y$ 在 $X_E$ 中的像
不属于 $Z_E$。于是对每个 $E\subset E'$，$y$ 在 $X_{E'}$ 中的像也不
属于 $Z_{E'}$。令 $T_{E'}\subset X_{E'}$ 为 $y$ 的像的闭包所给出的
约化闭子概形。显然，$T=\lim_{E\subset E'}T_{E'}$ 是 $y$ 在 $X_w$
中的闭包。对每个 $E\subset E'$，概形 $T_{E'}\cap Z_{E'}$ 由
$X_{E'}$ 的构造非空。因此
$\lim T_{E'}\cap Z_{E'}$ 非空，从而 $T\cap Z$ 非空。故 $y$ 不是闭点。
这表明 $X_w$ 的每个闭点都属于 $Z$。

\medskip\noindent
设 $y\in X_w$ 特化到 $z,z'\in Z$。我们来证明 $z=z'$；这将完成 (3)
的证明，并推出 (5)。令 $x,x'\in X$ 分别为 $z$ 与 $z'$ 的像。由于
$Z\to X$ 是双射，只需证明 $x=x'$。若 $x \not = x'$，则存在 $f\in A$，
使 $x\in D(f)$ 且 $x'\in V(f)$（或反之）。置 $E=\{f\}$，于是
$$
X_E = \Spec(A_f) \amalg \Spec(A_{V(f)}^\sim)
$$
可见 $z$ 与 $z'$ 分别映到 $x_E$ 与 $x'_E$，而后二者位于上述 $X_E$
分解的不同部分。因此，$x_E$ 与 $x'_E$ 不可能是同一点的特化，
这就是所需的矛盾。

\medskip\noindent
回忆，对有限子集 $E\subset A$，$Z_E$ 是局部闭子概形
$Z(E',E'')$ 的不交并；其中每一项都同构于 $(A/I)_f$ 的谱，
其中 $I$ 是由 $E''$ 生成的理想，$f$ 是 $E'$ 中各元素的乘积。
任意幂零元 $b$ 在 $(A/I)_f$ 中都是 $g/f^n$ 的类，其中 $g\in A$。
置 $E'=E\cup\{g\}$，读者可验证，$b$ 在此系统的过渡映射
$Z_{E'}\to Z_E$ 下拉回为零。这证明了 (4)。
\end{proof}

\begin{remark}
\label{proetale-remark-size-w}
设 $A$ 为环，$\kappa$ 为不小于 $A$ 的基数的无限基数。则
$A_w$（引理 \ref{proetale-lemma-make-w-local}）的基数至多为 $\kappa$。
事实上，每个 $A_E$ 的基数至多为 $\kappa$，而 $A$ 的有限子集所成的
集合的基数也至多为 $\kappa$。由于
$\kappa\otimes\kappa=\kappa$，结论成立；见“集合”第
\ref{sets-section-cardinals} 节。
\end{remark}

\begin{lemma}[此构造的泛性质]
\label{proetale-lemma-universal}
设 $A$ 为环，并设 $A\to A_w$ 为引理 \ref{proetale-lemma-make-w-local} 中构造的
环映射。对任意环映射 $A\to B$，若 $\Spec(B)$ 是 w-局部的，则存在唯一分解
$A\to A_w\to B$，使得 $\Spec(B)\to\Spec(A_w)$ 是 w-局部的。
\end{lemma}

\begin{proof}
记 $Y=\Spec(B)$，并记 $Y_0\subset Y$ 为闭点集。记给定态射为
$f:Y\to X$。回忆，$Y_0$ 是预有限的；特别地，$Y_0$ 的每个可构造子集
都是开闭的。设 $E\subset A$ 为有限子集。回忆，
$A_w=\colim A_E$，且 $\Spec(A_w)$ 的闭点集是闭子集
$Z_E\subset X_E=\Spec(A_E)$ 的逆极限。因此，只需证明存在唯一分解
$A\to A_E\to B$，使 $Y\to X_E$ 把 $Y_0$ 映入 $Z_E$。
由于 $Z_E\to X=\Spec(A)$ 是双射，且各分层片 $Z(E',E'')$ 可构造，故
$$
Y_0 = \coprod f^{-1}(Z(E', E'')) \cap Y_0
$$
是到开闭子集的不交并分解。由于 $Y_0=\pi_0(Y)$，得到 $Y$ 到开闭部分的
相应分解。因此，只需在 $f(Y_0)\subset Z(E',E'')$ 对某个分解
$E=E'\amalg E''$ 成立时构造该分解。在这种情况下，$f(Y)$ 包含于
$X$ 中特化到 $Z(E',E'')$ 的点集，而该点集同胚于 $X_{E',E''}$。
于是得到唯一的连续映射 $Y\to X_{E',E''}$，它位于 $X$ 上。由引理
\ref{proetale-lemma-fully-faithful-spaces-over-X}，它对应于唯一的概形态射
$Y\to X_{E',E''}$，该态射位于 $X$ 上。证明完毕。
\end{proof}

\noindent
回忆，环的谱为预有限空间，当且仅当其中每个点都是闭点。事实上，
还有一系列与此等价的条件；见“代数”引理
\ref{algebra-lemma-ring-with-only-minimal-primes} 或“拓扑学”引理
\ref{topology-lemma-characterize-profinite-spectral}。

\begin{lemma}
\label{proetale-lemma-profinite-goes-up}
设 $A$ 为环且 $\Spec(A)$ 为预有限空间，并设 $A\to B$ 为环映射。
在下列每种情形下，$\Spec(B)$ 都是预有限空间：
\begin{enumerate}
\item 若 $\mathfrak q,\mathfrak q'\subset B$ 位于 $A$ 的同一个素理想之上，
则既不成立 $\mathfrak q\subset\mathfrak q'$，也不成立
$\mathfrak q'\subset\mathfrak q$；
\item $A\to B$ 诱导代数的剩余域扩张；
\item $A\to B$ 是局部同构；
\item $A\to B$ 诱导局部环同构；
\item $A\to B$ 是弱 \'etale 的；
\item $A\to B$ 是拟有限的；
\item $A\to B$ 是非分歧的；
\item $A\to B$ 是 \'etale 的；
\item $B$ 是具有 (1)--(8) 中性质的 $A$-代数的滤过余极限；
\item 等等。
\end{enumerate}
\end{lemma}

\begin{proof}
由上述参考结果（“代数”引理
\ref{algebra-lemma-ring-with-only-minimal-primes} 或“拓扑学”引理
\ref{topology-lemma-characterize-profinite-spectral}），$\Spec(A)$ 的不同点
之间没有特化关系，而 $\Spec(B)$ 是预有限空间，当且仅当 $\Spec(B)$ 的
不同点之间没有特化关系。这些特化关系只能发生在
$\Spec(B)\to\Spec(A)$ 的纤维中。因此 (1) 成立。

\medskip\noindent
由“代数”引理 \ref{algebra-lemma-finite-residue-extension-closed}，
(2) 的假设蕴含 $B$ 的所有素理想都是极大理想，故 (2) 成立。
若 $A\to B$ 是局部同构或诱导局部环同构，则剩余域扩张均为平凡扩张，
所以 (3) 与 (4) 由 (2) 得出。若 $A\to B$ 是弱 \'etale 的，则
“代数续篇”引理
\ref{more-algebra-lemma-weakly-etale-residue-field-extensions} 表明它诱导
可分代数的剩余域扩张，故 (5) 由 (2) 得出。若 $A\to B$ 是拟有限的，
则各纤维都是有限离散拓扑空间，故 (6) 由 (1) 得出。因此 (3) 由 (1)
得出。情形 (7) 与 (8) 由此得出，因为非分歧及 \'etale 环映射是拟有限的
（“代数”引理 \ref{algebra-lemma-unramified-quasi-finite} 与
\ref{algebra-lemma-etale-quasi-finite}）。若
$B=\colim B_i$ 是 $A$-代数的滤过余极限，则由“极限”引理
\ref{limits-lemma-inverse-limit-top}，在拓扑空间范畴中有
$\Spec(B)=\lim\Spec(B_i)$。因此，若每个 $\Spec(B_i)$ 都是预有限空间，
则由“拓扑学”引理 \ref{topology-lemma-directed-inverse-limit-profinite}，
$\Spec(B)$ 亦然。这证明了 (9)。
\end{proof}

\begin{lemma}
\label{proetale-lemma-localize-along-closed-profinite}
设 $A$ 为环，并设 $V(I)\subset\Spec(A)$ 为一个作为拓扑空间是预有限的
闭子集。则存在 ind-Zariski 环映射 $A\to B$，使得 $\Spec(B)$ 是
w-局部的，其闭点集为 $V(IB)$，且 $A/I\cong B/IB$。
\end{lemma}

\begin{proof}
取引理 \ref{proetale-lemma-make-w-local} 中的 $A\to A_w$ 与
$Z\subset Y=\Spec(A_w)$。令 $T\subset Z$ 为 $V(I)$ 的逆像。
由“拓扑学”引理 \ref{topology-lemma-bijective-map}，
$T\to V(I)$ 是同胚。令 $B=(A_w)_T^\sim$；见引理
\ref{proetale-lemma-localization}。显然，$B$ 是 w-局部的，且闭点集为 $V(IB)$。
环映射 $A/I\to B/IB$ 是 ind-Zariski 的，并在底层拓扑空间上诱导同胚。
故由引理 \ref{proetale-lemma-local-isomorphism-fully-faithful}，它是同构。
\end{proof}

\begin{lemma}
\label{proetale-lemma-w-local-algebraic-residue-field-extensions}
设 $A$ 为环且 $X=\Spec(A)$ 是 w-局部的。设 $I\subset A$ 为定义闭点集
$X_0$ 的根理想；该闭点集取自 $X$。设环映射 $A\to B$ 在各素理想处诱导代数的
剩余域扩张。则
\begin{enumerate}
\item $Z=V(IB)$ 的每个点都是 $\Spec(B)$ 的闭点；
\item 存在 ind-Zariski 环映射 $B\to C$，使得
\begin{enumerate}
\item $B/IB\to C/IC$ 是同构；
\item 空间 $Y=\Spec(C)$ 是 w-局部的；
\item 诱导映射 $p:Y\to X$ 是 w-局部的；
\item $p^{-1}(X_0)$ 是 $Y$ 的闭点集。
\end{enumerate}
\end{enumerate}
\end{lemma}

\begin{proof}
把引理 \ref{proetale-lemma-profinite-goes-up} 应用于 $A/I\to B/IB$，可知
$Z=V(IB)=\Spec(B/IB)$ 的所有点都是闭点；事实上，$\Spec(B/IB)$ 是
预有限空间。最后把引理 \ref{proetale-lemma-localize-along-closed-profinite}
应用于 $IB\subset B$，即得结论。
\end{proof}





\section{诱导局部环同构与 ind-Zariski 的比较}
\label{proetale-section-connected-components}

\noindent
ind-Zariski 环映射 $A \to B$ 诱导局部环同构（引理
\ref{proetale-lemma-ind-zariski-implies}），但其逆命题不成立（“例”第
\ref{examples-section-not-ind-etale} 节）。不过，诱导局部环同构的环映射
可以用 ind-Zariski 环映射来描述；这给出一种结构定理，见命题
\ref{proetale-proposition-maps-wich-identify-local-rings}。

\medskip\noindent
设 $A$ 为环，令 $X = \Spec(A)$。由“拓扑学”引理
\ref{topology-lemma-spectral-pi0}（以及“代数”引理
\ref{algebra-lemma-spec-spectral}），连通分支空间 $\pi_0(X)$ 是
预有限空间。

\begin{lemma}
\label{proetale-lemma-construct}
设 $A$ 为环，令 $X = \Spec(A)$。设 $T \subset \pi_0(X)$ 为闭子集。
则存在满的 ind-Zariski 环映射 $A \to B$，使得
$\Spec(B) \to \Spec(A)$ 诱导 $\Spec(B)$ 与 $T$ 在 $X$ 中的逆像之间的
同胚。
\end{lemma}

\begin{proof}
令 $Z \subset X$ 为 $T$ 的逆像。则 $Z$ 是 $X$ 中所有包含 $Z$ 的开闭
子集的交 $Z = \bigcap Z_\alpha$；见“拓扑学”引理
\ref{topology-lemma-closed-union-connected-components}。对每个 $\alpha$，
有 $Z_\alpha = \Spec(A_\alpha)$，其中 $A \to A_\alpha$ 是局部同构
（在一个幂等元处的局部化）。置 $B = \colim A_\alpha$ 即得结论。
\end{proof}

\begin{lemma}
\label{proetale-lemma-construct-profinite}
设 $A$ 为环并令 $X = \Spec(A)$。设 $T$ 为预有限空间，且
$T \to \pi_0(X)$ 为连续映射。则存在 ind-Zariski 环映射
$A \to B$，使得在记 $Y = \Spec(B)$ 时，下图
$$
\xymatrix{
Y \ar[r] \ar[d] & \pi_0(Y) \ar[d] \\
X \ar[r] & \pi_0(X)
}
$$
在拓扑空间范畴中是笛卡尔图，并且作为 $\pi_0(X)$ 上的空间有
$\pi_0(Y) = T$。
\end{lemma}

\begin{proof}
具体地，把 $T = \lim T_i$ 写成以有向集为指标的有限离散空间逆系统的
极限（见“拓扑学”引理 \ref{topology-lemma-profinite}）。对每个 $i$，令
$Z_i = \Im(T \to \pi_0(X) \times T_i)$；这是闭子集。注意，
$X \times T_i$ 是 $A_i = \prod_{t \in T_i} A$ 的谱，且
$A \to A_i$ 是局部同构。由引理 \ref{proetale-lemma-construct}，
$Z_i \subset \pi_0(X \times T_i) = \pi_0(X) \times T_i$ 对应于一个
ind-Zariski 满射 $A_i \to B_i$，使得作为 $X \times T_i$ 的子集有
$\Spec(B_i) = X \times_{\pi_0(X)} Z_i$。过渡映射 $T_i \to T_{i'}$
诱导映射 $Z_i \to Z_{i'}$ 及
$X \times_{\pi_0(X)} Z_i \to X \times_{\pi_0(X)} Z_{i'}$，从而诱导环映射
$B_{i'} \to B_i$（引理 \ref{proetale-lemma-local-isomorphism-fully-faithful} 与
\ref{proetale-lemma-ind-zariski-implies}）。置 $B = \colim B_i$。由于
$T = \lim Z_i$，有
$X \times_{\pi_0(X)} T = \lim  X \times_{\pi_0(X)} Z_i$，因而
$Y = \Spec(B) = \lim \Spec(B_i)$ 位于拓扑空间的笛卡尔图
$$
\xymatrix{
Y \ar[r] \ar[d] & T \ar[d] \\
X \ar[r] & \pi_0(X)
}
$$
中。由引理 \ref{proetale-lemma-silly}，可得 $T = \pi_0(Y)$。
\end{proof}

\begin{example}
\label{proetale-example-construct-space}
设 $k$ 为域，$T$ 为预有限拓扑空间。存在 ind-Zariski 环映射
$k \to A$，使得 $\Spec(A)$ 同胚于 $T$。事实上，只需把引理
\ref{proetale-lemma-construct-profinite} 应用于
$T \to \pi_0(\Spec(k)) = \{*\}$。在这种情况下，若把 $T = \lim T_i$
写成每个 $T_i$ 都有限的滤过极限，则有
$$
A = \colim \text{Map}(T_i, k)
$$
\end{example}

\begin{lemma}
\label{proetale-lemma-w-local-morphism-equal-points-stalks-is-iso}
设 $A \to B$ 为满足下列条件的环映射：
\begin{enumerate}
\item $A \to B$ 诱导局部环同构；
\item 拓扑空间 $\Spec(B)$、$\Spec(A)$ 是 w-局部的；
\item $\Spec(B) \to \Spec(A)$ 是 w-局部的；
\item $\pi_0(\Spec(B)) \to \pi_0(\Spec(A))$ 是双射。
\end{enumerate}
则 $A \to B$ 是同构。
\end{lemma}

\begin{proof}
令 $X_0 \subset X = \Spec(A)$ 与 $Y_0 \subset Y = \Spec(B)$ 分别为
闭点集。由假设，$Y_0$ 映入 $X_0$，且诱导映射 $Y_0 \to X_0$ 是双射。
作为空间，$\Spec(A)$ 是 $A$ 在各闭点处的局部环之谱的不交并；$B$ 也
同样。因此 $X \to Y$ 是双射。由于 $A \to B$ 平坦，下行定理成立
（“代数”引理 \ref{algebra-lemma-flat-going-down}）。故“代数”引理
\ref{algebra-lemma-unique-prime-over-localize-below} 表明，对位于
$\mathfrak p \subset A$ 之上的任意素理想 $\mathfrak q \subset B$，有
$B_\mathfrak q = B_\mathfrak p$。由假设，$B_\mathfrak q = A_\mathfrak p$，
故对 $A$ 的每个素理想 $\mathfrak p$，都有
$A_\mathfrak p = B_\mathfrak p$。于是由“代数”引理
\ref{algebra-lemma-characterize-zero-local}，$A = B$。
\end{proof}

\begin{lemma}
\label{proetale-lemma-w-local-morphism-equal-stalks-is-ind-zariski}
设 $A \to B$ 为满足下列条件的环映射：
\begin{enumerate}
\item $A \to B$ 诱导局部环同构；
\item 拓扑空间 $\Spec(B)$、$\Spec(A)$ 是 w-局部的；
\item $\Spec(B) \to \Spec(A)$ 是 w-局部的。
\end{enumerate}
则 $A \to B$ 是 ind-Zariski 的。
\end{lemma}

\begin{proof}
置 $X = \Spec(A)$ 与 $Y = \Spec(B)$。令 $X_0 \subset X$ 与
$Y_0 \subset Y$ 为闭点集。取仿射概形的 ind-Zariski 态射
$A \to A'$，使得在记 $X' = \Spec(A')$ 时，下图
$$
\xymatrix{
X' \ar[r] \ar[d] & \pi_0(X') \ar[d] \\
X \ar[r] & \pi_0(X)
}
$$
在拓扑空间范畴中是笛卡尔图，并且作为 $\pi_0(X)$ 上的空间有
$\pi_0(X') = \pi_0(Y)$；见引理 \ref{proetale-lemma-construct-profinite}。
由引理 \ref{proetale-lemma-silly}，$X'$ 是 w-局部的，且闭点集
$X'_0 \subset X'$ 是 $X_0$ 的逆像。

\medskip\noindent
在 $X$ 上得到底层拓扑空间的连续映射 $Y \to X'$，它将 $\pi_0(Y)$
与 $\pi_0(X')$ 等同。由引理
\ref{proetale-lemma-local-isomorphism-fully-faithful}（以及引理
\ref{proetale-lemma-ind-zariski-implies}），它对应于 $X$ 上的仿射概形态射
$Y \to X'$。由于 $Y \to X$ 把 $Y_0$ 映入 $X_0$，可见
$Y \to X'$ 把 $Y_0$ 映入 $X'_0$，即 $Y \to X'$ 是 w-局部的。
由引理 \ref{proetale-lemma-w-local-morphism-equal-points-stalks-is-iso}，
$Y \cong X'$，得证。
\end{proof}

\noindent
下面的命题是我们稍后将证明的那类结果的一个预备。

\begin{proposition}
\label{proetale-proposition-maps-wich-identify-local-rings}
设 $A \to B$ 为诱导局部环同构的环映射。则存在忠实平坦的
ind-Zariski 环映射 $B \to B'$，使得 $A \to B'$ 是 ind-Zariski 的。
\end{proposition}

\begin{proof}
分别令 $A \to A_w$、$B \to B_w$ 为引理
\ref{proetale-lemma-make-w-local} 对 $A$、$B$ 构造的忠实平坦 ind-Zariski 环映射。
由于 $\Spec(B_w)$ 是 w-局部的，由引理 \ref{proetale-lemma-universal}，存在唯一分解
$A \to A_w \to B_w$，使得 $\Spec(B_w) \to \Spec(A_w)$ 是 w-局部的。
注意，$A_w \to B_w$ 诱导局部环同构；见引理
\ref{proetale-lemma-local-isomorphism-permanence}。由引理
\ref{proetale-lemma-w-local-morphism-equal-stalks-is-ind-zariski}，这意味着
$A_w \to B_w$ 是 ind-Zariski 的。由于 $B \to B_w$ 是忠实平坦的
ind-Zariski 映射（引理 \ref{proetale-lemma-make-w-local}），而复合
$A \to B \to B_w$ 是 ind-Zariski 的（引理
\ref{proetale-lemma-composition-ind-zariski}），命题得证。
\end{proof}

\noindent
上述命题使我们可以刻画仿射概形的 pro-Zariski 位点中的仿射弱可缩对象。

\begin{lemma}
\label{proetale-lemma-w-local-extremally-disconnected}
设 $A$ 为环。下列条件等价：
\begin{enumerate}
\item 每个诱导局部环同构的忠实平坦环映射 $A \to B$ 都有回缩；
\item 每个忠实平坦 ind-Zariski 环映射 $A \to B$ 都有回缩；
\item $A$ 满足
\begin{enumerate}
\item $\Spec(A)$ 是 w-局部的；
\item $\pi_0(\Spec(A))$ 是极端不连通的。
\end{enumerate}
\end{enumerate}
\end{lemma}

\begin{proof}
(1) 与 (2) 的等价性立即由命题
\ref{proetale-proposition-maps-wich-identify-local-rings} 得出。

\medskip\noindent
假设 (3)(a) 与 (3)(b)。设 $A \to B$ 忠实平坦且为 ind-Zariski。
下文将不再明说地使用如下事实：平坦映射 $A \to B$ 忠实平坦，当且仅当
$\Spec(A)$ 的每个闭点都属于 $\Spec(B) \to \Spec(A)$ 的像。
我们来证明 $A \to B$ 有回缩。

\medskip\noindent
设 $I \subset A$ 为理想，使得 $V(I) \subset \Spec(A)$ 是
$\Spec(A)$ 的闭点集。可把 $B$ 替换为引理
\ref{proetale-lemma-w-local-algebraic-residue-field-extensions} 对 $A \to B$ 与
$I \subset A$ 构造的环 $C$。因此可以假设 $\Spec(B)$ 是 w-局部的，
且 $\Spec(B)$ 的闭点集为 $V(IB)$。

\medskip\noindent
假设 $\Spec(B)$ 是 w-局部的，且 $\Spec(B)$ 的闭点集为 $V(IB)$。
对满连续映射 $V(IB) \to V(I)$ 选取一个连续截面。这是可能的，因为
$V(I) \cong \pi_0(\Spec(A))$ 极端不连通；见“拓扑学”命题
\ref{topology-proposition-projective-in-category-hausdorff-qc}。该截面的像是
闭子空间 $T \subset \pi_0(\Spec(B)) \cong V(IB)$，并同胚地映到
$\pi_0(A)$ 上。把 $B$ 替换为引理 \ref{proetale-lemma-construct} 中构造的
ind-Zariski 商环，可假设
$\pi_0(\Spec(B)) \to \pi_0(\Spec(A))$ 为双射。此时由引理
\ref{proetale-lemma-w-local-morphism-equal-points-stalks-is-iso}，$A \to B$ 是同构。

\medskip\noindent
假设 (1)，等价地假设 (2)。令 $A \to A_w$ 为引理
\ref{proetale-lemma-make-w-local} 中构造的环映射。由 (1)，存在回缩
$A_w \to A$。因此 $\Spec(A)$ 同胚于 $\Spec(A_w)$ 的一个闭子集。
由引理 \ref{proetale-lemma-closed-subspace-w-local}，(3)(a) 成立。最后，令
$T \to \pi_0(A)$ 为满映射，其中 $T$ 是极端不连通、拟紧的 Hausdorff
拓扑空间（“拓扑学”引理 \ref{topology-lemma-existence-projective-cover}）。
按适用于 $T \to \pi_0(\Spec(A))$ 的方式，选取引理
\ref{proetale-lemma-construct-profinite} 中的 $A \to B$。由 (1)，存在回缩
$B \to A$。于是 $T = \pi_0(\Spec(B)) \to \pi_0(\Spec(A))$ 有截面。
利用“拓扑学”命题
\ref{topology-proposition-projective-in-category-hausdorff-qc} 的形式范畴论
论证可知，$\pi_0(\Spec(A))$ 是极端不连通的。
\end{proof}

\begin{lemma}
\label{proetale-lemma-find-Zariski-w-contractible}
设 $A$ 为环。则存在忠实平坦的 ind-Zariski 环映射 $A \to B$，使得
$B$ 满足引理 \ref{proetale-lemma-w-local-extremally-disconnected} 中的等价条件。
\end{lemma}

\begin{proof}
首先应用引理 \ref{proetale-lemma-make-w-local}，可假设 $\Spec(A)$ 是 w-局部的。
选取极端不连通空间 $T$ 及满连续映射
$T \to \pi_0(\Spec(A))$；见“拓扑学”引理
\ref{topology-lemma-existence-projective-cover}。注意，$T$ 是预有限空间。
应用引理 \ref{proetale-lemma-construct-profinite}，得到 ind-Zariski 环映射
$A \to B$，使得 $\pi_0(\Spec(B)) \to \pi_0(\Spec(A))$ 实现映射
$T \to \pi_0(\Spec(A))$，并使下图
$$
\xymatrix{
\Spec(B) \ar[r] \ar[d] & \pi_0(\Spec(B)) \ar[d] \\
\Spec(A) \ar[r] & \pi_0(\Spec(A))
}
$$
在拓扑空间范畴中是笛卡尔图。注意，$\Spec(B)$ 是 w-局部的，
$\Spec(B) \to \Spec(A)$ 是 w-局部的，而且 $\Spec(B)$ 的闭点集是
$\Spec(A)$ 的闭点集的逆像；见引理 \ref{proetale-lemma-silly}。因此，
引理 \ref{proetale-lemma-w-local-extremally-disconnected} 的条件 (3) 对 $B$ 成立。
\end{proof}

\begin{remark}
\label{proetale-remark-slightly-stronger}
在引理 \ref{proetale-lemma-construct}、\ref{proetale-lemma-construct-profinite}、命题
\ref{proetale-proposition-maps-wich-identify-local-rings} 以及引理
\ref{proetale-lemma-find-Zariski-w-contractible} 中，我们都找到了具有某些性质的
ind-Zariski 环映射。论文 \cite{BS} 的作者使用 ind-(Zariski 局部化)
这一概念，它是主局部化的有限直积所成的滤过余极限。在上述每个结果中，
都可以把 ind-Zariski 替换为 ind-(Zariski 局部化)。但本文不需要这一点，
而且这里定义的环的 ind-Zariski 同态在形式上具有稍好的性质。此外，
ind-Zariski 环映射这一概念正是下一节所定义的 ind-\'etale 环映射的
自然类比。
\end{remark}




\section{Ind-\'etale 代数}
\label{proetale-section-ind-etale}

\noindent
先给出定义。

\begin{definition}
\label{proetale-definition-ind-etale}
若 $B$ 可写成 \'etale $A$-代数的滤过余极限，则称环映射 $A \to B$
为 {\it ind-\'etale} 的。
\end{definition}

\noindent
ind-\'etale 代数范畴在若干自然运算下封闭。

\begin{lemma}
\label{proetale-lemma-base-change-ind-etale}
设 $A \to B$ 与 $A \to A'$ 为环映射，并令 $B' = B \otimes_A A'$ 为
$B$ 的基变换。若 $A \to B$ 是 ind-\'etale 的，则 $A' \to B'$ 也是
ind-\'etale 的。
\end{lemma}

\begin{proof}
这就是“代数”引理 \ref{algebra-lemma-base-change-colimit-etale}。
\end{proof}

\begin{lemma}
\label{proetale-lemma-composition-ind-etale}
设 $A \to B$ 与 $B \to C$ 为环映射。若 $A \to B$ 与 $B \to C$ 都是
ind-\'etale 的，则 $A \to C$ 也是 ind-\'etale 的。
\end{lemma}

\begin{proof}
这就是“代数”引理 \ref{algebra-lemma-composition-colimit-etale}。
\end{proof}

\begin{lemma}
\label{proetale-lemma-ind-ind-etale}
ind-\'etale $A$-代数的滤过余极限在 $A$ 上仍为 ind-\'etale。
\end{lemma}

\begin{proof}
这就是“代数”引理 \ref{algebra-lemma-colimit-colimit-etale}。
\end{proof}

\begin{lemma}
\label{proetale-lemma-ind-etale-permanence}
设 $A$ 为环。设 $B \to C$ 为 ind-\'etale $A$-代数之间的 $A$-代数映射。
则 $C$ 是 ind-\'etale $B$-代数。
\end{lemma}

\begin{proof}
这就是“代数”引理 \ref{algebra-lemma-colimits-of-etale}。
\end{proof}

\begin{lemma}
\label{proetale-lemma-ind-etale-implies}
设 $A \to B$ 是 ind-\'etale 的。则 $A \to B$ 是弱 \'etale 的
（“代数续篇”定义 \ref{more-algebra-definition-weakly-etale}）。
\end{lemma}

\begin{proof}
这由“代数续篇”引理 \ref{more-algebra-lemma-when-weakly-etale} 得出。
\end{proof}

\begin{lemma}
\label{proetale-lemma-lift-ind-etale}
设 $A$ 为环，$I \subset A$ 为理想。基变换函子
$$
\text{ind-\'etale }A\text{-代数}
\longrightarrow
\text{ind-\'etale }A/I\text{-代数},\quad
C \longmapsto C/IC
$$
有全忠实的右伴随 $v$。特别地，给定 ind-\'etale $A/I$-代数
$\overline{C}$，存在 ind-\'etale $A$-代数 $C = v(\overline{C})$，使得
$\overline{C} = C/IC$。
\end{lemma}

\begin{proof}
设 $\overline{C}$ 为 ind-\'etale $A/I$-代数。考虑由分解
$A \to B \to \overline{C}$ 组成的范畴 $\mathcal{C}$，其中
$A \to B$ 是 \'etale 的。（本证明忽略若干集合论问题。）我们将证明
此范畴是有向的，并且 $C = \colim_\mathcal{C} B$ 是满足
$\overline{C} = C/IC$ 的 ind-\'etale $A$-代数。

\medskip\noindent
先证明 $\mathcal{C}$ 是有向的（“范畴”定义
\ref{categories-definition-directed}）。该范畴非空，因为
$A \to A \to \overline{C}$ 是一个对象。设
$A \to B \to \overline{C}$ 与 $A \to B' \to \overline{C}$ 是
$\mathcal{C}$ 的两个对象。则 $A \to B \otimes_A B' \to \overline{C}$
也是对象（使用“代数”引理 \ref{algebra-lemma-etale}）。设
$f, g : B \to B'$ 是 $\mathcal{C}$ 的对象
$A \to B \to \overline{C}$ 与 $A \to B' \to \overline{C}$ 之间的
两个映射。则其一个余等化子为
$A \to B' \otimes_{f, B, g} B' \to \overline{C}$。由“代数”引理
\ref{algebra-lemma-etale} 与 \ref{algebra-lemma-map-between-etale}，
它是 $\mathcal{C}$ 的对象。因此范畴 $\mathcal{C}$ 是有向的。

\medskip\noindent
把 $\overline{C} = \colim \overline{B_i}$ 写成滤过余极限，其中
$\overline{B_i}$ 在 $A/I$ 上是 \'etale 的。对每个 $i$，存在 \'etale 映射
$A \to B_i$，使得 $\overline{B_i} = B_i/IB_i$；见“代数”引理
\ref{algebra-lemma-lift-etale}。因此 $C \to \overline{C}$ 满射。由于
$C/IC \to \overline{C}$ 是 ind-\'etale 的（引理
\ref{proetale-lemma-ind-etale-permanence}），它是平坦的。因此 $\overline{C}$ 是
$C/IC$ 在某个乘法子集 $S \subset C/IC$ 处的局部化（“代数”引理
\ref{algebra-lemma-pure}）。取 $f \in C$，使其映到 $S \subset C/IC$ 的
一个元素。选取 $\mathcal{C}$ 中的 $A \to B \to \overline{C}$ 及
$g \in B$，使其在余极限中映到 $f$。于是
$A \to B_g \to \overline{C}$ 也是 $\mathcal{C}$ 的对象。因此 $f$ 是
$C$ 中的可逆元，从而 $C/IC = \overline{C}$。

\medskip\noindent
其次，断言对 $A$ 上的 ind-\'etale 代数 $D$ 有
$$
\Mor_A(D, C) = \Mor_{A/I}(D/ID, \overline{C})
$$
事实上，设 $D/ID \to \overline{C}$ 为 $A/I$-代数映射。把
$D = \colim_{i \in I} D_i$ 写成以有向集 $I$ 为指标的余极限，其中
$D_i$ 在 $A$ 上是 \'etale 的。由 $\mathcal{C}$ 的选取，得到函子
$I \to \mathcal{C}$，从而得到与到 $\overline{C}$ 的映射相容的映射
$D \to C$。断言得证。

\medskip\noindent
因此，由下述规则定义的函子 $v$
$$
\overline{C} 
\longmapsto
v(\overline{C}) = \colim_{A \to B \to \overline{C}} B
$$
是基变换函子 $u$ 的右伴随，正如引理所要求的。函子 $v$ 是全忠实的，
因为由构造有 $u \circ v = \text{id}$；见“范畴”引理
\ref{categories-lemma-adjoint-fully-faithful}。
\end{proof}









\section{构造 ind-\'etale 代数}
\label{proetale-section-construction-ind-etale}

\noindent
设 $A$ 为环。回顾任意 \'etale 环映射 $A \to B$ 都同构于相对维数为 $0$ 的
标准光滑环映射。这种环映射形如
$$
A \longrightarrow A[x_1, \ldots, x_n]/(f_1, \ldots, f_n)
$$
其中以 $\partial f_i/\partial x_j$ 为元素的 $n \times n$ 矩阵之行列式在商环中
可逆。见“代数”引理 \ref{algebra-lemma-etale-standard-smooth}。

\medskip\noindent
令 $S(A)$ 为所有相对维数为 $0$ 的{\it 忠实平坦}\footnote{在已有平坦性的
情况下，例如对光滑或 \'etale 环映射而言，这仅表示所诱导的谱映射是满射。
见“代数”引理 \ref{algebra-lemma-ff-rings}。}标准光滑 $A$-代数所成的集合。
令 $I(A)$ 为 $S(A)$ 的有限子集 $E$ 按包含关系组成的偏序集。注意 $I(A)$
是有向偏序集。对 $E = \{A \to B_1, \ldots, A \to B_n\}$，置
$$
B_E = B_1 \otimes_A \ldots \otimes_A B_n
$$
则 $B_E$ 是忠实平坦的 \'etale $A$-代数。若 $E \subset E'$，则有 \'etale
$A$-代数的典范过渡映射 $B_E \to B_{E'}$。具体地，设
$E = \{A \to B_1, \ldots, A \to B_n\}$ 且
$E' = \{A \to B_1, \ldots, A \to B_{n + m}\}$，则 $B_E \to B_{E'}$
把 $b_1 \otimes \ldots \otimes b_n$ 映到 $B_{E'}$ 的元素
$b_1 \otimes \ldots \otimes b_n \otimes 1 \otimes \ldots \otimes 1$。
此构造给出以 $I(A)$ 为指标的忠实平坦 \'etale $A$-代数系统，并置
$$
T(A) = \colim_{E \in I(A)} B_E
$$
注意，$T(A)$ 是忠实平坦的 ind-\'etale $A$-代数（“代数”引理
\ref{algebra-lemma-colimit-faithfully-flat}）。由构造，对任意忠实平坦的
\'etale $A$-代数 $B$，都有一个（不唯一的）$A$-代数映射 $B \to T(A)$。
事实上，选取某个 $(A \to B_0) \in S(A)$ 及同构 $B \cong B_0$。则典范映射
$$
B \to B_0 \to 
T(A) = \colim_{E \in I(A)} B_E
$$
就是所需映射。

\begin{lemma}
\label{proetale-lemma-first-construction}
给定环 $A$，存在忠实平坦的 ind-\'etale $A$-代数 $C$，使得每个忠实平坦的
\'etale 环映射 $C \to B$ 都有回缩。
\end{lemma}

\begin{proof}
置 $T^1(A) = T(A)$ 且 $T^{n + 1}(A) = T(T^n(A))$。令
$$
C = \colim T^n(A)
$$
此代数在每个 $T^n(A)$ 上忠实平坦，特别地在 $A$ 上忠实平坦；见“代数”引理
\ref{algebra-lemma-colimit-faithfully-flat}。此外，由引理
\ref{proetale-lemma-ind-ind-etale}，$C$ 在 $A$ 上是 ind-\'etale 的。若 $C \to B$
是 \'etale 的，则存在 $n$ 及 \'etale 环映射 $T^n(A) \to B'$，使得
$B = C \otimes_{T^n(A)} B'$；见“代数”引理 \ref{algebra-lemma-etale}。
若 $C \to B$ 忠实平坦，则 $\Spec(B) \to \Spec(C) \to \Spec(T^n(A))$
满射，因而 $\Spec(B') \to \Spec(T^n(A))$ 满射。换言之，
$T^n(A) \to B'$ 忠实平坦。由上述构造，存在 $T^n(A)$-代数映射
$B' \to T^{n + 1}(A)$。它诱导 $C$-代数映射 $B \to C$，证明完毕。
\end{proof}

\begin{remark}
\label{proetale-remark-size-T}
设 $A$ 为环，$\kappa$ 为不小于 $A$ 的基数的无限基数。则 $T(A)$ 的基数至多为
$\kappa$。事实上，每个 $B_E$ 的基数至多为 $\kappa$，指标集 $I(A)$ 的基数
也至多为 $\kappa$。由于 $\kappa \otimes \kappa = \kappa$，结论成立；见
“集合”第 \ref{sets-section-cardinals} 节。因此，引理
\ref{proetale-lemma-first-construction} 的证明中构造的环之基数也至多为 $\kappa$。
\end{remark}

\begin{remark}
\label{proetale-remark-first-construction-functorial}
构造 $A \mapsto T(A)$ 在如下意义下是函子性的：若 $A \to A'$ 为环映射，
则可以构造交换图
$$
\xymatrix{
A \ar[r] \ar[d] & T(A) \ar[d] \\
A' \ar[r] & T(A')
}
$$
。具体地，给定 $S(A)$ 中的
$(A \to A[x_1, \ldots, x_n]/(f_1, \ldots, f_n))$，利用环映射
$\varphi : A \to A'$ 得到 $S(A')$ 中的相应元素
$(A' \to A'[x_1, \ldots, x_n]/(f^\varphi_1, \ldots, f^\varphi_n))$；
这里 $f^\varphi$ 表示对多项式 $f$ 的系数施以 $\varphi$ 所得的多项式。此外，
环范畴中有交换图
$$
\xymatrix{
A \ar[r] \ar[d] & A[x_1, \ldots, x_n]/(f_1, \ldots, f_n) \ar[d] \\
A' \ar[r] & A'[x_1, \ldots, x_n]/(f^\varphi_1, \ldots, f^\varphi_n)
}
$$
。对有限集 $E \subset S(A)$，置 $E' = \varphi(E)$，并按显然方式定义
$B_E \to B_{E'}$。取余极限便得到所需映射 $T(A) \to T(A')$；见“范畴”引理
\ref{categories-lemma-functorial-colimit}。
\end{remark}

\begin{lemma}
\label{proetale-lemma-have-sections-quotient}
设环 $A$ 满足每个忠实平坦的 \'etale 环映射 $A \to B$ 都有回缩。则每个商环
$A/I$ 也具有同样的性质。
\end{lemma}

\begin{proof}
设 $A/I \to \overline{B}$ 忠实平坦且为 \'etale。由“代数”引理
\ref{algebra-lemma-lift-etale}，可对某个 \'etale 环映射 $A \to B$ 写成
$\overline{B} = B/IB$。映射 $\Spec(B) \to \Spec(A)$ 的像 $U$ 是开集且包含
$V(I)$。因此补集 $Z = \Spec(A) \setminus U$ 拟紧且与 $V(I)$ 不交。故对某个
$r \geq 0$ 及 $f_i \in I$，有
$Z \subset D(f_1) \cup \ldots \cup D(f_r)$。于是
$A \to B' = B \times \prod A_{f_i}$ 忠实平坦且为 \'etale，并且
$\overline{B} = B'/IB'$。因此，$A \to B'$ 的回缩 $B' \to A$ 诱导
$A/I \to \overline{B}$ 的回缩。
\end{proof}

\begin{lemma}
\label{proetale-lemma-have-sections-strictly-henselian}
设环 $A$ 满足每个忠实平坦的 \'etale 环映射 $A \to B$ 都有回缩。则 $A$ 在每个
极大理想处的局部环都是严格 Hensel 环。
\end{lemma}

\begin{proof}
设 $\mathfrak m$ 为 $A$ 的极大理想。设 $A \to B$ 为 \'etale 环映射，且
$\mathfrak q \subset B$ 为位于 $\mathfrak m$ 之上的素理想。根据“代数”引理
\ref{algebra-lemma-strict-henselization-different} 对严格 Hensel 化
$A_\mathfrak m^{sh}$ 的描述，只需证明 $A_\mathfrak m = B_\mathfrak q$。
注意，位于 $\mathfrak m$ 之上的素理想
$\mathfrak q = \mathfrak q_1, \mathfrak q_2, \ldots, \mathfrak q_n$ 只有有限个；
又因 \'etale 环映射是拟有限的，它们之间不存在特化关系，见“代数”引理
\ref{algebra-lemma-etale-quasi-finite}。因此 $\mathfrak q_i$ 是极大理想，并可找到
$g \in \mathfrak q_2 \cap \ldots \cap \mathfrak q_n$，
$g \not \in \mathfrak q$（“代数”引理 \ref{algebra-lemma-silly}）。把 $B$ 替换为
$B_g$ 后，$\mathfrak q$ 是 $B$ 中唯一位于 $\mathfrak m$ 之上的素理想。
映射 $\Spec(B) \to \Spec(A)$ 的像 $U \subset \Spec(A)$ 是开集（“代数”命题
\ref{algebra-proposition-fppf-open}）。因此补集 $\Spec(A) \setminus U$ 是闭集，
并可找到 $f \in A$、$f \not \in \mathfrak p$，使得
$\Spec(A) = U \cup D(f)$。环映射 $A \to B \times A_f$ 忠实平坦且为
\'etale，故由关于 $A$ 的假设，它有回缩 $\sigma : B \times A_f \to A$。
注意，$\sigma$ 是 \'etale 的，因而作为 \'etale $A$-代数之间的映射是平坦的
（“代数”引理 \ref{algebra-lemma-map-between-etale}）。由于 $\mathfrak q$ 是
$B \times A_f$ 中唯一位于 $A$ 之上的素理想，可知
$A_\mathfrak p \to B_\mathfrak q$ 有一个同样平坦的回缩。因此
$A_\mathfrak p \to B_\mathfrak q \to A_\mathfrak p$
是复合为恒等映射的平坦局部环映射。由于局部环的平坦局部同态是单射，得出
这些映射都是同构，正合所需。
\end{proof}

\begin{lemma}
\label{proetale-lemma-have-sections-localize}
设环 $A$ 满足每个忠实平坦的 \'etale 环映射 $A \to B$ 都有回缩。设
$Z \subset \Spec(A)$ 为闭子概形，并令 $A \to A_Z^\sim$ 如引理
\ref{proetale-lemma-localization} 中所构造。则每个忠实平坦的 \'etale 环映射
$A_Z^\sim \to C$ 都有回缩。
\end{lemma}

\begin{proof}
存在 \'etale 环映射 $A \to B'$，使得作为 $A_Z^\sim$-代数有
$C = B' \otimes_A A_Z^\sim$。映射 $\Spec(B') \to \Spec(A)$ 的像
$U' \subset \Spec(A)$ 是开集且包含 $V(I)$，因此可找到 $f \in I$，使得
$\Spec(A) = U' \cup D(f)$。于是 $A \to B' \times A_f$ 是 \'etale 且忠实平坦的。
由假设存在回缩 $B' \times A_f \to A$。局部化后得到所需回缩
$C \to A_Z^\sim$。
\end{proof}

\begin{lemma}
\label{proetale-lemma-get-w-local-algebraic-residue-field-extensions}
设 $A \to B$ 为在剩余域上诱导代数扩张的环映射。存在交换图
$$
\xymatrix{
B \ar[r] & D \\
A \ar[r] \ar[u] & C \ar[u]
}
$$
满足下列性质：
\begin{enumerate}
\item $A \to C$ 忠实平坦且为 ind-\'etale；
\item $B \to D$ 忠实平坦且为 ind-\'etale；
\item $\Spec(C)$ 是 w-局部的；
\item $\Spec(D)$ 是 w-局部的；
\item $\Spec(D) \to \Spec(C)$ 是 w-局部的；
\item $\Spec(D)$ 的闭点集是 $\Spec(C)$ 的闭点集的逆像；
\item $\Spec(C)$ 的闭点集满射到 $\Spec(A)$；
\item $\Spec(D)$ 的闭点集满射到 $\Spec(B)$；
\item 对极大理想 $\mathfrak m \subset C$，局部环 $C_\mathfrak m$ 是严格
Hensel 环。
\end{enumerate}
\end{lemma}

\begin{proof}
存在忠实平坦的 ind-Zariski 环映射 $A \to A'$，使得 $\Spec(A')$ 是
w-局部的，并且 $\Spec(A')$ 的闭点集映满 $\Spec(A)$；见引理
\ref{proetale-lemma-make-w-local}。设 $I \subset A'$ 为使 $V(I)$ 等于 $\Spec(A')$
的闭点集的理想。按引理 \ref{proetale-lemma-first-construction} 选取 $A' \to C'$。
注意，由引理 \ref{proetale-lemma-have-sections-strictly-henselian}，在极大理想
$\mathfrak m' \subset C'$ 处的局部环 $C'_{\mathfrak m'}$ 都是严格 Hensel 环。
把引理 \ref{proetale-lemma-w-local-algebraic-residue-field-extensions} 应用于
$A' \to C'$ 与 $I \subset A'$，得到满足 $C'/IC' \cong C/IC$ 的
$C' \to C$。注意，由于 $A' \to C'$ 忠实平坦，$\Spec(C'/IC')$ 满射到
$A'$ 的闭点集，特别地满射到 $\Spec(A)$。此外，由于
$V(IC) \subset \Spec(C)$ 是 $C$ 的闭点集，而 $C' \to C$ 是 ind-Zariski 的
（并诱导局部环同构），得到性质 (1)、(3)、(7) 与 (9)。

\medskip\noindent
用 $J \subset C$ 表示使 $V(J)$ 等于 $\Spec(C)$ 的闭点集的理想。置
$D' = B \otimes_A C$。环映射 $C \to D'$ 在剩余域上诱导代数扩张。注意，
由于 $V(J) \to \Spec(A)$ 满射，映射 $T = V(JD') \to \Spec(B)$ 也满射。
把引理 \ref{proetale-lemma-w-local-algebraic-residue-field-extensions} 应用于
$C \to D'$ 与 $J \subset C$，得到满足 $D'/JD' \cong D/JD$ 的
$D' \to D$。引理中其余所有性质都直接由引理
\ref{proetale-lemma-w-local-algebraic-residue-field-extensions} 的结论得出。
\end{proof}








\section{弱 \'etale 与 pro-\'etale 的比较}
\label{proetale-section-weakly-etale}

\noindent
回顾：若 $A \to B$ 平坦且 $B \otimes_A B \to B$ 平坦，则称环同态
$A \to B$ 为{\it 弱 \'etale} 的。我们已在“代数续篇”第
\ref{more-algebra-section-weakly-etale} 节证明了这种环映射的一些性质。特别地，
若 $A \to B$ 是局部同态，且 $A$ 是严格 Hensel 局部环，则 $A = B$；见
“代数续篇”定理 \ref{more-algebra-theorem-olivier}。利用该定理与上述工作，
得到弱 \'etale 环映射的如下结构定理。

\begin{proposition}
\label{proetale-proposition-weakly-etale}
设 $A \to B$ 为弱 \'etale 环映射。则存在忠实平坦的 ind-\'etale 环映射
$B \to B'$，使得 $A \to B'$ 是 ind-\'etale 的。
\end{proposition}

\begin{proof}
环映射 $A \to B$ 在剩余域上诱导（可分）代数扩张；见“代数续篇”引理
\ref{more-algebra-lemma-weakly-etale-residue-field-extensions}。因此可应用引理
\ref{proetale-lemma-get-w-local-algebraic-residue-field-extensions} 并选取图
$$
\xymatrix{
B \ar[r] & D \\
A \ar[r] \ar[u] & C \ar[u]
}
$$
并使其具有该引理所列的性质。由“代数续篇”引理
\ref{more-algebra-lemma-weakly-etale-permanence}，$C \to D$ 是弱 \'etale 的。
选取极大理想 $\mathfrak m \subset D$。由构造，它位于某个极大理想
$\mathfrak m' \subset C$ 之上。由“代数续篇”定理
\ref{more-algebra-theorem-olivier}，环映射
$C_{\mathfrak m'} \to D_\mathfrak m$ 是同构。由于 $\Spec(C)$ 的每个点都特化到
一个闭点，得出 $C \to D$ 诱导局部环同构。因此命题
\ref{proetale-proposition-maps-wich-identify-local-rings} 可用于环映射 $C \to D$。
选取忠实平坦且为 ind-Zariski 的 $D \to D'$，使得 $C \to D'$ 是
ind-Zariski 的。于是 $B \to D'$ 满足命题的要求。
\end{proof}





\section{V 拓扑与 pro-h 拓扑}
\label{proetale-section-V-versus-pro-h}

\noindent
V 拓扑在“拓扑”第 \ref{topologies-section-V} 节中引入。h 拓扑在
“平坦性续篇”第 \ref{flat-section-h} 节中引入。作为一种中间拓扑的 ph 拓扑
则在“拓扑”第 \ref{topologies-section-ph} 节中引入。

\medskip\noindent
给定合适概形范畴 $\mathcal{C}$ 上的拓扑 $\tau$，可按如下方式在
$\mathcal{C}$ 上引入“pro-$\tau$ 拓扑”。回顾，对 $\mathcal{C}$ 中的 $X$，用
$h_X$ 表示与 $X$ 相伴的可表预层。暂称 $\mathcal{C}$ 的态射 $X \to Y$ 为
$\tau$-盖\footnote{不要把它与覆盖的概念混淆。例如，若 $\tau = \etale$，
则任何有截面的态射 $X \to Y$ 都是 $\tau$-盖；但我们对 \'etale 覆盖
$\{V_i \to Y\}_{i \in I}$ 的定义要求每个 $V_i \to Y$ 都是 \'etale 的。}，
如果 $h_X \to h_Y$ 的 $\tau$-层化是满射。于是可把 pro-$\tau$ 拓扑定义为满足
下列条件的最粗拓扑：
\begin{enumerate}
\item pro-$\tau$ 拓扑细于 $\tau$ 拓扑；
\item 若 $Y$ 为仿射概形，且 $X = \lim X_\lambda$ 是 $Y$ 上仿射概形
$X_\lambda$ 的有向极限，并且对所有 $\lambda$，$h_{X_\lambda} \to h_Y$
都是 $\tau$-盖，则 $X \to Y$ 是 pro-$\tau$-盖。
\end{enumerate}
采用这种刻意严谨的表述，是因为不想指定 pro-$\tau$ 覆盖的一种选取：对不同的
$\tau$，适合采用不同的覆盖族集合。例如，在第 \ref{proetale-section-proetale} 节中将会
看到，为定义 pro-\'etale 拓扑，考察具有某种有限性条件的弱 \'etale 态射族很
有效。更一般地，本段提出的构造主要用于说明本节结果的动机；下文绝不会以此
方法隐式定义 pro-$\tau$ 拓扑。

\medskip\noindent
下一个引理说明 pro-V 拓扑等于 V 拓扑。

\begin{lemma}
\label{proetale-lemma-pro-V-V}
设 $Y$ 为仿射概形，$X = \lim X_i$ 为 $Y$ 上仿射概形的有向极限。下列条件等价：
\begin{enumerate}
\item $\{X \to Y\}$ 是标准 V 覆盖（“拓扑”定义
\ref{topologies-definition-standard-V-covering}）；
\item 对所有 $i$，$\{X_i \to Y\}$ 都是标准 V 覆盖。
\end{enumerate}
\end{lemma}

\begin{proof}
单元族 $\{X \to Y\}$ 是标准 V 覆盖，当且仅当对任意态射
$g : \Spec(V) \to Y$，存在赋值环扩张 $V \subset W$ 及交换图
$$
\xymatrix{
\Spec(W) \ar[r] \ar[d] & X \ar[d] \\
\Spec(V) \ar[r]^g & Y
}
$$
。因此 (1) $\Rightarrow$ (2) 由定义立即成立。反过来，假设 (2)，并给定如上的
$g : \Spec(V) \to Y$。写成 $\Spec(V) \times_Y X_i = \Spec(A_i)$。由于
$\{X_i \to Y\}$ 是标准 V 覆盖，可以选取赋值环 $W_i$ 及环映射
$A_i \to W_i$，使得复合 $V \to A_i \to W_i$ 是赋值环扩张。特别地，
$A_i$ 对其 $V$-挠的商 $A'_i$ 是忠实平坦的 $V$-代数。其平坦性由“代数续篇”
引理 \ref{more-algebra-lemma-valuation-ring-torsion-free-flat} 得出；谱上的满射性
则是因为 $A_i \to W_i$ 经由 $A'_i$ 分解。因此
$$
A = \colim A'_i
$$
是忠实平坦的 $V$-代数（“代数”引理
\ref{algebra-lemma-colimit-faithfully-flat}）。由于
$\{\Spec(A) \to \Spec(V)\}$ 是标准 fpqc 覆盖，它也是标准 V 覆盖（“拓扑”引理
\ref{topologies-lemma-standard-fpqc-standard-V}），故可选取
$\Spec(W) \to \Spec(A)$，使得 $V \to W$ 是赋值环扩张。再与态射
$\Spec(A) \to X = \Spec(\colim A_i)$ 复合，证明即告完成。
\end{proof}

\noindent
下一个引理说明 pro-h 拓扑、pro-ph 拓扑与 V 拓扑彼此相等。

\begin{lemma}
\label{proetale-lemma-pro-h-V}
设 $X \to Y$ 为仿射概形的态射。下列条件等价：
\begin{enumerate}
\item $\{X \to Y\}$ 是标准 V 覆盖（“拓扑”定义
\ref{topologies-definition-standard-V-covering}）；
\item $X = \lim X_i$ 是 $Y$ 上仿射概形的有向极限，并且对每个 $i$，
$\{X_i \to Y\}$ 都是 ph 覆盖；
\item $X = \lim X_i$ 是 $Y$ 上仿射概形的有向极限，并且对每个 $i$，
$\{X_i \to Y\}$ 都是 h 覆盖。
\end{enumerate}
\end{lemma}

\begin{proof}
证明 (2) $\Rightarrow$ (1)。回顾，由仿射概形间单个箭头给出的 V 覆盖是标准
V 覆盖；见“拓扑”定义 \ref{topologies-definition-V-covering} 与引理
\ref{topologies-lemma-refine-standard-V}。还要回顾，任意 ph 覆盖都是 V 覆盖；
见“拓扑”引理
\ref{topologies-lemma-zariski-etale-smooth-syntomic-fppf-fpqc-ph-V}。因此，若
$X = \lim X_i$ 如 (2) 所述，则对每个 $i$，$\{X_i \to Y\}$ 都是标准 V 覆盖。
由引理 \ref{proetale-lemma-pro-V-V}，(1) 成立。

\medskip\noindent
证明 (3) $\Rightarrow$ (2)。这是清楚的，因为 h 覆盖总是 ph 覆盖；见“平坦性
续篇”定义 \ref{flat-definition-h-covering}。

\medskip\noindent
证明 (1) $\Rightarrow$ (3)。这是有实质内容的方向，不过证明中的实质内容蕴含在
“平坦性续篇”引理 \ref{flat-lemma-equivalence-h-v-locally-finite-presentation}
中。写 $X = \Spec(A)$ 与 $Y = \Spec(R)$。可把 $A = \colim A_i$ 写成余极限，
其中 $A_i$ 在 $R$ 上有限表示；见“代数”引理
\ref{algebra-lemma-ring-colimit-fp}。置 $X_i = \Spec(A_i)$。由 (1) 与“拓扑”引理
\ref{topologies-lemma-refine-standard-V}，对所有 $i$，$\{X_i \to Y\}$ 都是标准
V 覆盖。因此，由“平坦性续篇”定义 \ref{flat-definition-h-covering}，
$\{X_i \to Y\}$ 是 h 覆盖。证明完毕。
\end{proof}

\noindent
粗略地说，下一个引理表明，保极限的 h 层对 V 覆盖满足层条件。另请与注
\ref{proetale-remark-h-limit-preserving} 比较。

\begin{lemma}
\label{proetale-lemma-h-limit-preserving}
设 $S$ 为概形，$F$ 为定义在 $S$ 上所有概形的范畴上的反变函子。若
\begin{enumerate}
\item $F$ 对 h 拓扑满足层性质；
\item $F$ 保极限（“极限”注 \ref{limits-remark-limit-preserving}），
\end{enumerate}
则 $F$ 对 V 拓扑满足层性质。
\end{lemma}

\begin{proof}
通过验证“拓扑”引理 \ref{topologies-lemma-sheaf-property-V} 的 (1) 与 (2') 来
证明。对 Zariski 覆盖的层性质来自 $F$ 对所有 h 覆盖都具有层性质这一事实。
最后，设 $X \to Y$ 为 $S$ 上仿射概形的态射，并且 $\{X \to Y\}$ 是 V 覆盖。
由引理 \ref{proetale-lemma-pro-h-V}，可把 $X = \lim X_i$ 写成 $Y$ 上仿射概形的有向
极限，使得对每个 $i$，$\{X_i \to Y\}$ 都是 h 覆盖。于是得到
\begin{align*}
&
\text{等化子}(
\xymatrix{
F(X) \ar@<1ex>[r] \ar@<-1ex>[r] &
F(X \times_Y X)
}
)
\\
& =
\text{等化子}(
\xymatrix{
\colim F(X_i) \ar@<1ex>[r] \ar@<-1ex>[r] &
\colim F(X_i \times_Y X_i)
}
)
\\
& =
\colim
\text{等化子}(
\xymatrix{
F(X_i) \ar@<1ex>[r] \ar@<-1ex>[r] &
F(X_i \times_Y X_i)
}
)
\\
& =
\colim F(Y) = F(Y)
\end{align*}
这正是所要证明的。第一个等式成立，是因为 $F$ 保极限，且
$X = \lim X_i$、$X \times_Y X = \lim X_i \times_Y X_i$。第二个等式成立，
是因为滤过余极限是正合的。第三个等式成立，是因为 $F$ 对 h 覆盖满足层性质。
\end{proof}

\begin{remark}
\label{proetale-remark-h-limit-preserving}
设 $S$ 为包含在大位点 $\Sch_h$ 中的概形。设 $F$ 为 $(\Sch/S)_h$ 上的集合层，
并且只要 $T = \lim T_i$ 是 $(\Sch/S)_h$ 中仿射概形的有向极限，就有
$F(T) = \colim F(T_i)$。在这种情况下，$F$ 唯一延拓为 $S$ 上所有概形的范畴
上的反变函子 $F'$，使得 (a) $F'$ 对 h 拓扑满足层性质，并且 (b) $F'$ 保极限。
见“平坦性续篇”引理 \ref{flat-lemma-extend-sheaf-h}。此时，引理
\ref{proetale-lemma-h-limit-preserving} 表明 $F'$ 对 V 拓扑满足层性质。
\end{remark}







\section{构造 w-可缩覆盖}
\label{proetale-section-w-contractible}

\noindent
本节构造仿射概形的 w-可缩覆盖。

\begin{definition}
\label{proetale-definition-w-contractible}
设 $A$ 为环。若每个忠实平坦的弱 \'etale 环映射 $A \to B$ 都有回缩，则称
$A$ 是{\it w-可缩的}。
\end{definition}

\noindent
注意，由命题 \ref{proetale-proposition-weakly-etale}，一个等价定义是要求每个忠实平坦的
ind-\'etale 环映射 $A \to B$ 都有回缩。下面这个关键观察将使我们能够构造
w-可缩环。

\begin{lemma}
\label{proetale-lemma-w-local-strictly-henselian-extremally-disconnected}
设 $A$ 为环。下列条件等价：
\begin{enumerate}
\item $A$ 是 w-可缩的；
\item 每个忠实平坦的 ind-\'etale 环映射 $A \to B$ 都有回缩；
\item $A$ 满足
\begin{enumerate}
\item $\Spec(A)$ 是 w-局部的；
\item $\pi_0(\Spec(A))$ 是极端不连通的；
\item 对每个极大理想 $\mathfrak m \subset A$，局部环 $A_\mathfrak m$
都是严格 Hensel 环。
\end{enumerate}
\end{enumerate}
\end{lemma}

\begin{proof}
(1) 与 (2) 的等价性立即由命题 \ref{proetale-proposition-weakly-etale} 得出。

\medskip\noindent
假设 (3)(a)、(3)(b) 与 (3)(c)。设 $A \to B$ 忠实平坦且为 ind-\'etale。
下文将不再明说地使用如下事实：平坦映射 $A \to B$ 忠实平坦，当且仅当
$\Spec(A)$ 的每个闭点都属于 $\Spec(B) \to \Spec(A)$ 的像。我们来证明
$A \to B$ 有回缩。

\medskip\noindent
设 $I \subset A$ 为理想，使得 $V(I) \subset \Spec(A)$ 是 $\Spec(A)$ 的闭点集。
可把 $B$ 替换为引理 \ref{proetale-lemma-w-local-algebraic-residue-field-extensions}
对 $A \to B$ 与 $I \subset A$ 构造的环 $C$。因此可以假设 $\Spec(B)$ 是
w-局部的，且 $\Spec(B)$ 的闭点集为 $V(IB)$。此时，由条件 (3)(c)，
$A \to B$ 诱导局部环同构，因为只需在 $B$ 中位于 $A$ 的极大理想之上的极大
理想处检验。于是由引理 \ref{proetale-lemma-w-local-extremally-disconnected}，
$A \to B$ 有回缩。

\medskip\noindent
假设 (1)，等价地假设 (2)。由引理
\ref{proetale-lemma-have-sections-strictly-henselian} 得到 (3)(c)。性质 (3)(a) 与
(3)(b) 由引理 \ref{proetale-lemma-w-local-extremally-disconnected} 得出。
\end{proof}

\begin{proposition}
\label{proetale-proposition-find-w-contractible}
对每个环 $A$，都存在忠实平坦的 ind-\'etale 环映射 $A \to D$，使得 $D$ 是
w-可缩的。
\end{proposition}

\begin{proof}
把引理 \ref{proetale-lemma-get-w-local-algebraic-residue-field-extensions} 应用于
$\text{id}_A : A \to A$，得到忠实平坦的 ind-\'etale 环映射 $A \to C$，使得
$C$ 是 w-局部的，且 $C$ 在每个极大理想处的局部环都是严格 Hensel 环。选取
极端不连通空间 $T$ 及满连续映射 $T \to \pi_0(\Spec(C))$；见“拓扑学”引理
\ref{topology-lemma-existence-projective-cover}。注意，$T$ 是预有限的。应用引理
\ref{proetale-lemma-construct-profinite}，得到 ind-Zariski 环映射 $C \to D$，使得
$\pi_0(\Spec(D)) \to \pi_0(\Spec(C))$ 实现 $T \to \pi_0(\Spec(C))$，并且
$$
\xymatrix{
\Spec(D) \ar[r] \ar[d] & \pi_0(\Spec(D)) \ar[d] \\
\Spec(C) \ar[r] & \pi_0(\Spec(C))
}
$$
在拓扑空间范畴中是笛卡尔图。注意，$\Spec(D)$ 是 w-局部的，
$\Spec(D) \to \Spec(C)$ 是 w-局部的，而且 $\Spec(D)$ 的闭点集是
$\Spec(C)$ 的闭点集的逆像；见引理 \ref{proetale-lemma-silly}。因此，$D$ 在其极大理想
处的局部环仍然都是严格 Hensel 环（因为它们同构于 $C$ 在相应极大理想处的
局部环）。由引理 \ref{proetale-lemma-w-local-strictly-henselian-extremally-disconnected}
可知，$D$ 是 w-可缩的。
\end{proof}

\begin{remark}
\label{proetale-remark-size-w-contractible}
设 $A$ 为环，$\kappa$ 为不小于 $A$ 的基数的无限基数。则命题
\ref{proetale-proposition-find-w-contractible} 中构造的环 $D$ 的基数至多为
$$
\kappa^{2^{2^{2^\kappa}}}.
$$
事实上，环映射 $A \to D$ 构造为复合
$$
A \to A_w = A' \to C' \to C \to D.
$$
其中，构造的前三步在引理
\ref{proetale-lemma-get-w-local-algebraic-residue-field-extensions} 证明的第一段中完成。
对第一步，由注 \ref{proetale-remark-size-w} 有 $|A_w| \leq \kappa$。由注
\ref{proetale-remark-size-T} 有 $|C'| \leq \kappa$。然后 $|C| \leq \kappa$，因为
$C$ 是 $(C')_w$ 的局部化（它由对 $C'$ 应用引理
\ref{proetale-lemma-localize-along-closed-profinite} 构造；该应用位于引理
\ref{proetale-lemma-w-local-algebraic-residue-field-extensions} 的证明中）。因此 $C$ 至多有
$2^\kappa$ 个极大理想。最后，环映射 $C \to D$ 诱导局部环同构，并且由
“拓扑学”注 \ref{topology-remark-size-projective-cover}，$D$ 的极大理想集的基数
至多为 $2^{2^{2^\kappa}}$。由于
$D \subset \prod_{\mathfrak m \subset D} D_\mathfrak m$，可见 $D$ 的基数至多为
上述基数。
\end{remark}

\begin{lemma}
\label{proetale-lemma-finite-finitely-presented-over-extremally-disconnected}
设 $A \to B$ 为拟有限且有限表示的环映射。若 $A$ 的剩余域都是可分代数闭域，
且 $\Spec(A)$ 是 Hausdorff 且极端不连通的，则 $\Spec(B)$ 是极端不连通的。
\end{lemma}

\begin{proof}
置 $X = \Spec(A)$ 与 $Y = \Spec(B)$。按“\'Etale 上同调”引理
\ref{etale-cohomology-lemma-decompose-quasi-finite-morphism} 选取有限分割
$X = \coprod X_i$ 及 $X'_i \to X_i$。由“拓扑学”命题
\ref{topology-proposition-projective-in-category-hausdorff-qc}，拓扑空间映射
$\coprod X_i \to X$（其源是拓扑空间范畴中的不交并）有截面。因此在拓扑意义下，
$X$ 是各层 $X_i$ 的不交并。故可把 $X$ 替换为各 $X_i$，并假设存在满的有限局部
自由态射 $X' \to X$，使 $(X' \times_X Y)_{red}$ 同构于有限多个
$X'_{red}$ 副本的不交并。考虑图
$$
\xymatrix{
\coprod_{i = 1, \ldots, r} X' \ar[r] \ar[d] & Y \ar[d] \\
X' \ar[r] & X
}
$$
关于 $A$ 的剩余域的假设蕴含：在底层点集上，此图是纤维积图（略去细节）。
由于 $X$ 极端不连通，而 $X'$ 是 Hausdorff 的（引理
\ref{proetale-lemma-profinite-goes-up}），连续映射 $X' \to X$ 有连续截面 $\sigma$。于是
$\coprod_{i = 1, \ldots, r} \sigma(X) \to Y$ 是双射连续映射。由“拓扑学”引理
\ref{topology-lemma-bijective-map}，它是同胚，证明完毕。
\end{proof}

\begin{lemma}
\label{proetale-lemma-finite-finitely-presented-over-w-contractible}
设 $A \to B$ 为有限且有限表示的环映射。若 $A$ 是 w-可缩的，则 $B$ 也是。
\end{lemma}

\begin{proof}
使用引理 \ref{proetale-lemma-w-local-strictly-henselian-extremally-disconnected} 的判据。
置 $X = \Spec(A)$、$Y = \Spec(B)$，并用 $f : Y \to X$ 表示所诱导的态射。
由于 $f : Y \to X$ 是有限态射，$Y$ 的闭点集 $Y_0$ 是 $X$ 的闭点集 $X_0$
的逆像。设 $y \in Y$ 的像为 $x \in X$，则 $x$ 特化到唯一闭点
$x_0 \in X$。设 $f^{-1}(\{x_0\}) = \{y_1, \ldots, y_n\}$，其中 $y_i$ 在
$Y$ 中闭。由于 $R = \mathcal{O}_{X, x_0}$ 是严格 Hensel 环，且 $f$ 有限，可知
$Y \times_{f, X} \Spec(R)$ 等于
$\coprod_{i = 1, \ldots, n} \Spec(R_i)$，其中每个 $R_i$ 都是在 $R$ 上有限的
局部环，其极大理想对应于 $y_i$；见“代数”引理
\ref{algebra-lemma-characterize-henselian} 的 (10)。于是 $y$ 恰好属于这些
$\Spec(R_i)$ 中的一个，从而 $y$ 恰好特化到一个 $y_i$。换言之，$Y$ 的每个点
都特化到 $Y_0$ 的唯一一点。因此 $Y$ 是 w-局部的。对每个像为 $x \in X_0$
的 $y \in Y_0$，把“代数”引理 \ref{algebra-lemma-finite-over-henselian} 应用于
$\mathcal{O}_{X, x} \to B \otimes_A \mathcal{O}_{X, x}$，可知
$\mathcal{O}_{Y, y}$ 是严格 Hensel 环。还需证明 $Y_0$ 极端不连通。为此考察
$X_0 \times_X Y \to X_0$，其中 $X_0 \subset X$ 带有约化诱导概形结构。注意，
$X_0 \times_X Y$ 的底层拓扑空间就是 $Y_0$。现在所需结论由引理
\ref{proetale-lemma-finite-finitely-presented-over-extremally-disconnected} 得出。
\end{proof}

\begin{lemma}
\label{proetale-lemma-localization-w-contractible}
设 $A$ 为环，$Z \subset \Spec(A)$ 为形如 $Z = V(f_1, \ldots, f_r)$ 的闭子集。
置 $B = A_Z^\sim$，见引理 \ref{proetale-lemma-localization}。若 $A$ 是 w-可缩的，
则 $B$ 也是。
\end{lemma}

\begin{proof}
设 $A_Z^\sim \to B$ 为忠实平坦的弱 \'etale 环映射。考虑环映射
$$
A \longrightarrow A_{f_1} \times \ldots \times A_{f_r} \times B
$$
，它忠实平坦且为弱 \'etale。若 $A$ 是 w-可缩的，则存在回缩 $\sigma$。考虑态射
$$
\Spec(A_Z^\sim) \to \Spec(A) \xrightarrow{\Spec(\sigma)}
\coprod \Spec(A_{f_i}) \amalg \Spec(B)
$$
$Z \subset \Spec(A_Z^\sim)$ 的每个点都映入分支 $\Spec(B)$。由于
$\Spec(A_Z^\sim)$ 的每个点都特化到 $Z$ 的一点，便得到所需态射
$\Spec(A_Z^\sim) \to \Spec(B)$。
\end{proof}






\section{pro-\'etale 位点}
\label{proetale-section-proetale}

\noindent
本节只讨论各种 pro-\'etale 位点及其间态射的实际定义与构造。弱可缩对象的存在性
将在第 \ref{proetale-section-weakly-contractible} 节中处理。

\medskip\noindent
pro-\'etale 拓扑在一点上类似 fpqc 拓扑（见“拓扑”第
\ref{topologies-section-fpqc} 节）：概形的小 pro-\'etale 位点上的层拓扑斯依赖于
底层概形范畴的选取。因此不能谈论概形的{\it 那个} pro-\'etale 拓扑斯。不过，
扩大底层概形范畴并不改变层的上同调群；见第 \ref{proetale-section-change-universe} 节。

\medskip\noindent
我们将用概形的弱 \'etale 态射定义 pro-\'etale 覆盖；见“态射续篇”第
\ref{more-morphisms-section-weakly-etale} 节。原因是，一方面，定义概形的
pro-\'etale 态射这一概念颇为棘手；另一方面，命题
\ref{proetale-proposition-weakly-etale} 保证下述定义给出相同的层\footnote{准确地说，
用这里选取的覆盖得到的 pro-\'etale 拓扑，与第 \ref{proetale-section-V-versus-pro-h} 节
所述一般步骤从 $\tau = \etale$ 出发得到的拓扑相同。}。

\begin{definition}
\label{proetale-definition-fpqc-covering}
设 $T$ 为概形。$T$ 的一个{\it pro-\'etale 覆盖}是概形态射族
$\{f_i : T_i \to T\}_{i \in I}$，使得每个 $f_i$ 都是弱 \'etale 的，并且对
每个仿射开集 $U \subset T$，都存在 $n \geq 0$、映射
$a : \{1, \ldots, n\} \to I$ 及仿射开集 $V_j \subset T_{a(j)}$，
$j = 1, \ldots, n$，满足 $\bigcup_{j = 1}^n f_{a(j)}(V_j) = U$。
\end{definition}

\noindent
当然，此条件蕴含 $T = \bigcup f_i(T_i)$。下面的引理使我们能够识别
pro-\'etale 覆盖，也使我们能把许多关于 pro-\'etale 覆盖的引理归约为 fpqc
覆盖的相应结果。

\begin{lemma}
\label{proetale-lemma-recognize-proetale-covering}
设 $T$ 为概形，$\{f_i : T_i \to T\}_{i \in I}$ 为以 $T$ 为靶的概形态射族。
下列条件等价：
\begin{enumerate}
\item $\{f_i : T_i \to T\}_{i \in I}$ 是 pro-\'etale 覆盖；
\item 每个 $f_i$ 都是弱 \'etale 的，且 $\{f_i : T_i \to T\}_{i \in I}$
是 fpqc 覆盖；
\item 每个 $f_i$ 都是弱 \'etale 的，并且对每个仿射开集 $U \subset T$，
存在除有限多个外均为空的拟紧开集 $U_i \subset T_i$，使得
$U = \bigcup f_i(U_i)$；
\item 每个 $f_i$ 都是弱 \'etale 的，并且存在仿射开覆盖
$T = \bigcup_{\alpha \in A} U_\alpha$，使得对每个 $\alpha \in A$，都存在
$i_{\alpha, 1}, \ldots, i_{\alpha, n(\alpha)} \in I$ 及拟紧开集
$U_{\alpha, j} \subset T_{i_{\alpha, j}}$，满足
$U_\alpha = \bigcup_{j = 1, \ldots, n(\alpha)}
f_{i_{\alpha, j}}(U_{\alpha, j})$。
\end{enumerate}
若 $T$ 拟分离，则上述条件还等价于：
\begin{enumerate}
\item[(5)] 每个 $f_i$ 都是弱 \'etale 的，并且对每个 $t \in T$，都存在
$i_1, \ldots, i_n \in I$ 及拟紧开集 $U_j \subset T_{i_j}$，使得
$\bigcup_{j = 1, \ldots, n} f_{i_j}(U_j)$ 是 $t$ 在 $T$ 中的一个邻域
（不一定是开集）。
\end{enumerate}
\end{lemma}

\begin{proof}
(1) 与 (2) 的等价性由定义立即成立。因此本引理由“拓扑”引理
\ref{topologies-lemma-recognize-fpqc-covering} 得出。
\end{proof}

\begin{lemma}
\label{proetale-lemma-etale-proetale}
任意 \'etale 覆盖及任意 Zariski 覆盖都是 pro-\'etale 覆盖。
\end{lemma}

\begin{proof}
这由 fpqc 覆盖的相应结果（“拓扑”引理
\ref{topologies-lemma-zariski-etale-smooth-syntomic-fppf-fpqc}）、引理
\ref{proetale-lemma-recognize-proetale-covering}，以及 \'etale 态射是弱 \'etale 态射这一
事实得出；见“态射续篇”引理 \ref{more-morphisms-lemma-when-weakly-etale}。
\end{proof}

\begin{lemma}
\label{proetale-lemma-proetale}
设 $T$ 为概形。
\begin{enumerate}
\item 若 $T' \to T$ 是同构，则 $\{T' \to T\}$ 是 $T$ 的 pro-\'etale 覆盖。
\item 若 $\{T_i \to T\}_{i\in I}$ 是 pro-\'etale 覆盖，且对每个 $i$ 都有
pro-\'etale 覆盖 $\{T_{ij} \to T_i\}_{j\in J_i}$，则
$\{T_{ij} \to T\}_{i \in I, j\in J_i}$ 是 pro-\'etale 覆盖。
\item 若 $\{T_i \to T\}_{i\in I}$ 是 pro-\'etale 覆盖，而 $T' \to T$
是概形态射，则 $\{T' \times_T T_i \to T'\}_{i\in I}$ 是 pro-\'etale 覆盖。
\end{enumerate}
\end{lemma}

\begin{proof}
这由以下事实得出：弱 \'etale 态射的复合与基变换仍为弱 \'etale（“态射续篇”引理
\ref{more-morphisms-lemma-composition-weakly-etale} 与
\ref{more-morphisms-lemma-base-change-weakly-etale}）、引理
\ref{proetale-lemma-recognize-proetale-covering}，以及 fpqc 覆盖的相应结果（见“拓扑”引理
\ref{topologies-lemma-fpqc}）。
\end{proof}

\begin{lemma}
\label{proetale-lemma-proetale-affine}
设 $T$ 为仿射概形，$\{T_i \to T\}_{i \in I}$ 为 $T$ 的 pro-\'etale 覆盖。
则存在作为 $\{T_i \to T\}_{i \in I}$ 之加细的 pro-\'etale 覆盖
$\{U_j \to T\}_{j = 1, \ldots, n}$，使得每个 $U_j$ 都是仿射概形。此外，
可以选取每个 $U_j$ 为某个 $T_i$ 中的仿射开集。
\end{lemma}

\begin{proof}
这直接由定义得出。
\end{proof}

\noindent
因此如下定义仿射概形的相应标准覆盖。

\begin{definition}
\label{proetale-definition-standard-proetale}
设 $T$ 为仿射概形。$T$ 的一个{\it 标准 pro-\'etale 覆盖}是态射族
$\{f_i : T_i \to T\}_{i = 1, \ldots, n}$，其中每个 $T_j$ 都是仿射的，
每个 $f_i$ 都是弱 \'etale 的，并且 $T = \bigcup f_i(T_i)$。
\end{definition}

\noindent
我们依照“拓扑”第 \ref{topologies-section-procedure} 节给出的一般方案，构造将要
使用的大 pro-\'etale 位点。不过，由于需要稍大的环来容纳某些构造的规模，
我们对这些构造略作修改。

\begin{definition}
\label{proetale-definition-big-proetale-site}
一个{\it 大 pro-\'etale 位点}是按如下方式构造的、如“位点”定义
\ref{sites-definition-site} 所述的任意位点 $\Sch_\proetale$：
\begin{enumerate}
\item 选取任意概形集 $S_0$，以及这些概形间任意的 pro-\'etale 覆盖集
$\text{Cov}_0$。
\item 把“集合”式 (\ref{sets-equation-bound}) 中的函数 $Bound$ 改为
$$
Bound(\kappa) = \max\{\kappa^{2^{2^{2^\kappa}}}, \kappa^{\aleph_0}, \kappa^+\}.
$$
\item 取从集合 $S_0$ 与函数 $Bound$ 出发、按“集合”引理
\ref{sets-lemma-construct-category} 构造的任意范畴 $\Sch_\alpha$ 为底层范畴。
\item 从范畴 $\Sch_\alpha$、pro-\'etale 覆盖类及上面选取的集合
$\text{Cov}_0$ 出发，按“集合”引理 \ref{sets-lemma-coverings-site} 选取任意覆盖集。
\end{enumerate}
\end{definition}

\noindent
关于大位点定义的动机与说明，见“拓扑”定义
\ref{topologies-definition-big-zariski-site} 后的各条注。

\medskip\noindent
由引理 \ref{proetale-lemma-proetale-induced} 将可看出，大 pro-\'etale 位点
$\Sch_\proetale$ 上的拓扑在某种意义下由所有概形的范畴上的 pro-\'etale 拓扑
诱导而来。

\begin{definition}
\label{proetale-definition-big-small-proetale}
设 $S$ 为概形，$\Sch_\proetale$ 为包含 $S$ 的大 pro-\'etale 位点。
\begin{enumerate}
\item $S$ 的{\it 大 pro-\'etale 位点}记为 $(\Sch/S)_\proetale$，它是“位点”
第 \ref{sites-section-localize} 节中引入的位点 $\Sch_\proetale/S$。
\item $S$ 的{\it 小 pro-\'etale 位点}记为 $S_\proetale$，它是
$(\Sch/S)_\proetale$ 的全子范畴，其对象为使 $U \to S$ 弱 \'etale 的
$U/S$。$S_\proetale$ 的覆盖是 $(\Sch/S)_\proetale$ 的任意覆盖
$\{U_i \to U\}$，其中 $U \in \Ob(S_\proetale)$。
\item $S$ 的{\it 大仿射 pro-\'etale 位点}记为
$(\textit{Aff}/S)_\proetale$，它是 $(\Sch/S)_\proetale$ 的全子范畴，其对象为
仿射的 $U/S$。$(\textit{Aff}/S)_\proetale$ 的覆盖是
$(\Sch/S)_\proetale$ 的任意标准 pro-\'etale 覆盖 $\{U_i \to U\}$。
\end{enumerate}
\end{definition}

\noindent
小 pro-\'etale 位点和大仿射 pro-\'etale 位点是否确为位点，并非完全显然。我们现在
验证这一点。

\begin{lemma}
\label{proetale-lemma-verify-site-proetale}
设 $S$ 为概形，$\Sch_\proetale$ 为包含 $S$ 的大 pro-\'etale 位点。则
$S_\proetale$ 与 $(\textit{Aff}/S)_\proetale$ 都是位点。
\end{lemma}

\begin{proof}
先证明 $S_\proetale$ 是位点。它是一个范畴，并配有一组具有固定目标的态射族。
因此须证明“位点”定义 \ref{sites-definition-site} 的性质 (1)、(2) 与 (3)。
由于 $(\Sch/S)_\proetale$ 是位点，只须证明：若 $\{U_i \to U\}$ 是
$(\Sch/S)_\proetale$ 的覆盖且 $U \in \Ob(S_\proetale)$，则也有
$U_i \in \Ob(S_\proetale)$。这由定义以及弱 \'etale 态射的复合仍为弱 \'etale
态射这一事实得出。

\medskip\noindent
同理，为证明 $(\textit{Aff}/S)_\proetale$ 是位点，只须证明仿射概形的标准
pro-\'etale 覆盖全体满足“位点”定义 \ref{sites-definition-site} 的性质 (1)、(2)
与 (3)。这由引理 \ref{proetale-lemma-recognize-proetale-covering} 以及标准 fpqc 覆盖的
相应结果（“拓扑”引理 \ref{topologies-lemma-fpqc-affine-axioms}）得出。
\end{proof}

\begin{lemma}
\label{proetale-lemma-fibre-products-proetale}
设 $S$ 为概形，$\Sch_\proetale$ 为包含 $S$ 的大 pro-\'etale 位点，并设
$\Sch$ 为所有概形组成的范畴。
\begin{enumerate}
\item 范畴 $\Sch_\proetale$、$(\Sch/S)_\proetale$、$S_\proetale$ 与
$(\textit{Aff}/S)_\proetale$ 都有纤维积，且与 $\Sch$ 中的纤维积一致。
\item 范畴 $\Sch_\proetale$、$(\Sch/S)_\proetale$ 与 $S_\proetale$ 都有
等化子，且与 $\Sch$ 中的等化子一致。
\item 范畴 $(\Sch/S)_\proetale$ 与 $S_\proetale$ 都有终对象，即 $S/S$。
\item 范畴 $\Sch_\proetale$ 有与 $\Sch$ 的终对象一致的终对象，即
$\Spec(\mathbf{Z})$。
\end{enumerate}
\end{lemma}

\begin{proof}
按构造，范畴 $\Sch_\proetale$ 包含 $\Spec(\mathbf{Z})$，并且对积与纤维积封闭；
见“集合”引理 \ref{sets-lemma-what-is-in-it}。设有概形态射 $U \to S$、
$V \to U$、$W \to U$，其中 $U, V, W \in \Ob(\Sch_\proetale)$。
$\Sch_\proetale$ 中的纤维积 $V \times_U W$ 也是 $\Sch$ 中的纤维积；在所有
$S$-概形组成的范畴中，它又是 $V/S$ 与 $W/S$ 在 $U/S$ 上的纤维积，因而也是
$(\Sch/S)_\proetale$ 中的纤维积。这证明了关于 $(\Sch/S)_\proetale$ 的结论。
若 $U \to S$、$V \to U$ 与 $W \to U$ 都是弱 \'etale 的，则
$V \times_U W \to S$ 也是弱 \'etale 的（见“态射续篇”第
\ref{more-morphisms-section-weakly-etale} 节），故 $S_\proetale$ 中也有纤维积。
若 $U, V, W$ 都是仿射的，则 $V \times_U W$ 也是仿射的，故
$(\textit{Aff}/S)_\proetale$ 中也有纤维积。

\medskip\noindent
设 $a, b : U \to V$ 是 $\Sch_\proetale$ 中的两个态射。此时 $a$ 与 $b$ 的等化子
（在概形范畴中）为
$$
V
\times_{\Delta_{V/\Spec(\mathbf{Z})}, V \times_{\Spec(\mathbf{Z})} V, (a, b)}
(U \times_{\Spec(\mathbf{Z})} U)
$$
由上可知它是 $\Sch_\proetale$ 的对象。因此 $\Sch_\proetale$ 有等化子。若 $a$ 与
$b$ 是 $S$ 上的态射，则其等化子（在概形范畴中）还可写为
$$
V \times_{\Delta_{V/S}, V \times_S V, (a, b)} (U \times_S U)
$$
故 $(\Sch/S)_\proetale$ 有等化子。此外，若 $U$ 与 $V$ 在 $S$ 上弱 \'etale，
则上述等化子作为在 $S$ 上弱 \'etale 的概形之纤维积，也弱 \'etale。
因此 $S_\proetale$ 有等化子。关于终对象的断言是显然的。
\end{proof}

\noindent
下面验证大仿射 pro-\'etale 位点与大 pro-\'etale 位点定义同一拓扑斯。

\begin{lemma}
\label{proetale-lemma-affine-big-site-proetale}
设 $S$ 为概形，$\Sch_\proetale$ 为包含 $S$ 的大 pro-\'etale 位点。函子
$(\textit{Aff}/S)_\proetale \to (\Sch/S)_\proetale$ 是特殊余连续函子。因此它诱导
从 $\Sh((\textit{Aff}/S)_\proetale)$ 到 $\Sh((\Sch/S)_\proetale)$ 的拓扑斯等价。
\end{lemma}

\begin{proof}
特殊余连续函子的概念见“位点”定义
\ref{sites-definition-special-cocontinuous-functor}。因此须验证“位点”引理
\ref{sites-lemma-equivalence} 的假设 (1)--(5)。将包含函子记为
$u : (\textit{Aff}/S)_\proetale \to (\Sch/S)_\proetale$。余连续性恰好是说：当
$T$ 仿射时，$T/S$ 的任意 pro-\'etale 覆盖均可由 $T$ 的一个标准 pro-\'etale
覆盖加细。这正是引理 \ref{proetale-lemma-proetale-affine} 的内容，故 (1) 成立。
由于标准 pro-\'etale 覆盖就是 pro-\'etale 覆盖，$u$ 显然连续，故 (2) 成立。
(3) 与 (4) 立即由 $u$ 全忠实得出。最后，(5) 由每个概形都有仿射开覆盖得出。
\end{proof}

\begin{lemma}
\label{proetale-lemma-put-in-T}
设 $\Sch_\proetale$ 为大 pro-\'etale 位点，$f : T \to S$ 为
$\Sch_\proetale$ 中的态射。函子 $T_\proetale \to (\Sch/S)_\proetale$ 余连续，
并诱导拓扑斯态射
$$
i_f :
\Sh(T_\proetale)
\longrightarrow
\Sh((\Sch/S)_\proetale)
$$
对于 $(\Sch/S)_\proetale$ 上的层 $\mathcal{G}$，有公式
$(i_f^{-1}\mathcal{G})(U/T) = \mathcal{G}(U/S)$。函子 $i_f^{-1}$ 还有一个左伴随
$i_{f, !}$，后者保持纤维积与等化子。
\end{lemma}

\begin{proof}
将该函子记为 $u : T_\proetale \to (\Sch/S)_\proetale$。换言之，给定与
$T_\proetale$ 的一个对象对应的弱 \'etale 态射 $j : U \to T$，令
$u(U \to T) = (f \circ j : U \to S)$。由引理
\ref{proetale-lemma-fibre-products-proetale}，此函子保持纤维积。此外，引理
\ref{proetale-lemma-fibre-products-proetale} 说明 $T_\proetale$ 有等化子且 $u$ 保持等化子。
它显然余连续。又因 $u$ 将覆盖变为
覆盖并保持纤维积，它也是连续的。因此结论由“位点”引理
\ref{sites-lemma-when-shriek} 与 \ref{sites-lemma-preserve-equalizers} 得出。
\end{proof}

\begin{lemma}
\label{proetale-lemma-at-the-bottom}
设 $S$ 为概形，$\Sch_\proetale$ 为包含 $S$ 的大 pro-\'etale 位点。包含函子
$S_\proetale \to (\Sch/S)_\proetale$ 满足“位点”引理
\ref{sites-lemma-bigger-site} 的假设，因而诱导位点态射
$$
\pi_S : (\Sch/S)_\proetale \longrightarrow S_\proetale
$$
及拓扑斯态射
$$
i_S : \Sh(S_\proetale) \longrightarrow \Sh((\Sch/S)_\proetale)
$$
使得 $\pi_S \circ i_S = \text{id}$。此外，$i_S = i_{\text{id}_S}$，其中
$i_{\text{id}_S}$ 如引理 \ref{proetale-lemma-put-in-T} 所述。特别地，函子
$i_S^{-1} = \pi_{S, *}$ 由规则
$i_S^{-1}(\mathcal{G})(U/S) = \mathcal{G}(U/S)$ 描述。
\end{lemma}

\begin{proof}
在此情形，函子 $u : S_\proetale \to (\Sch/S)_\proetale$ 除具有上面引理
\ref{proetale-lemma-put-in-T} 的证明中所见的性质外，还全忠实并将终对象变为终对象。
本引理由“位点”引理 \ref{sites-lemma-bigger-site} 得出。
\end{proof}

\begin{definition}
\label{proetale-definition-restriction-small-proetale}
在引理 \ref{proetale-lemma-at-the-bottom} 的情形中，函子 $i_S^{-1} = \pi_{S, *}$ 常称为
{\it 到小 pro-\'etale 位点的限制}。对于大 pro-\'etale 位点上的层
$\mathcal{F}$，将此限制记为 $\mathcal{F}|_{S_\proetale}$。
\end{definition}

\noindent
采用此记号，对于大位点上的层 $\mathcal{F}$ 和大位点上的层 $\mathcal{G}$，有
\begin{align*}
\Mor_{\Sh(S_\proetale)}(\mathcal{F}|_{S_\proetale}, \mathcal{G})
& =
\Mor_{\Sh((\Sch/S)_\proetale)}(\mathcal{F},
i_{S, *}\mathcal{G}) \\
\Mor_{\Sh(S_\proetale)}(\mathcal{G}, \mathcal{F}|_{S_\proetale})
& =
\Mor_{\Sh((\Sch/S)_\proetale)}(\pi_S^{-1}\mathcal{G}, \mathcal{F})
\end{align*}
此外，有 $(i_{S, *}\mathcal{G})|_{S_\proetale} = \mathcal{G}$，并且有
$(\pi_S^{-1}\mathcal{G})|_{S_\proetale} = \mathcal{G}$。

\begin{lemma}
\label{proetale-lemma-morphism-big}
设 $\Sch_\proetale$ 为大 pro-\'etale 位点，$f : T \to S$ 为
$\Sch_\proetale$ 中的态射。函子
$$
u : (\Sch/T)_\proetale \longrightarrow (\Sch/S)_\proetale, \quad
V/T \longmapsto V/S
$$
余连续，并有连续的右伴随
$$
v : (\Sch/S)_\proetale \longrightarrow (\Sch/T)_\proetale, \quad
(U \to S) \longmapsto (U \times_S T \to T).
$$
它们诱导同一个拓扑斯态射
$$
f_{big} :
\Sh((\Sch/T)_\proetale)
\longrightarrow
\Sh((\Sch/S)_\proetale)
$$
我们有 $f_{big}^{-1}(\mathcal{G})(U/T) = \mathcal{G}(U/S)$。
我们有 $f_{big, *}(\mathcal{F})(U/S) = \mathcal{F}(U \times_S T/T)$。
此外，$f_{big}^{-1}$ 有一个左伴随 $f_{big!}$，后者保持纤维积与等化子。
\end{lemma}

\begin{proof}
函子 $u$ 余连续、连续，并保持纤维积与等化子（细节从略；比较引理
\ref{proetale-lemma-put-in-T} 的证明）。因此可应用“位点”引理
\ref{sites-lemma-when-shriek} 与 \ref{sites-lemma-preserve-equalizers}，从而得到
$f_{big}^{-1}$ 的公式及 $f_{big!}$ 的存在性。此外，函子 $v$ 是右伴随，因为给定
$U/T$ 与 $V/S$，有 $\Mor_S(u(U), V) = \Mor_T(U, V \times_S T)$，正如所需。
于是可应用“位点”引理 \ref{sites-lemma-have-functor-other-way} 与
\ref{sites-lemma-have-functor-other-way-morphism}，得到 $f_{big, *}$ 的公式。
\end{proof}

\begin{lemma}
\label{proetale-lemma-morphism-big-small}
设 $\Sch_\proetale$ 为大 pro-\'etale 位点，$f : T \to S$ 为
$\Sch_\proetale$ 中的态射。
\begin{enumerate}
\item 有 $i_f = f_{big} \circ i_T$，其中 $i_f$ 如引理
\ref{proetale-lemma-put-in-T} 所述，而 $i_T$ 如引理 \ref{proetale-lemma-at-the-bottom} 所述。
\item 函子 $S_\proetale \to T_\proetale$，
$(U \to S) \mapsto (U \times_S T \to T)$ 连续，并诱导拓扑斯态射
$$
f_{small} : \Sh(T_\proetale) \longrightarrow \Sh(S_\proetale).
$$
有 $f_{small, *}(\mathcal{F})(U/S) = \mathcal{F}(U \times_S T/T)$。
\item 有如下位点态射的交换图
$$
\xymatrix{
T_\proetale \ar[d]_{f_{small}} &
(\Sch/T)_\proetale \ar[d]^{f_{big}} \ar[l]^{\pi_T}\\
S_\proetale &
(\Sch/S)_\proetale \ar[l]_{\pi_S}
}
$$
故作为拓扑斯态射，有 $f_{small} \circ \pi_T = \pi_S \circ f_{big}$。
\item 有 $f_{small} = \pi_S \circ f_{big} \circ i_T = \pi_S \circ i_f$。
\end{enumerate}
\end{lemma}

\begin{proof}
等式 $i_f = f_{big} \circ i_T$ 由等式
$i_f^{-1} = i_T^{-1} \circ f_{big}^{-1}$ 得出；后者由上面对这些函子的描述显然成立。
这证明了 (1)。

\medskip\noindent
函子 $u : S_\proetale \to T_\proetale$，
$u(U \to S) = (U \times_S T \to T)$ 将覆盖变为覆盖并保持纤维积；见引理
\ref{proetale-lemma-proetale} 与 \ref{proetale-lemma-fibre-products-proetale}。此外，
$S_\proetale$、$T_\proetale$ 都有终对象，即 $S/S$ 与 $T/T$，并且
$u(S/S) = T/T$。因此由“位点”命题 \ref{sites-proposition-get-morphism}，函子
$u$ 对应一个位点态射 $T_\proetale \to S_\proetale$。它进而给出该拓扑斯态射；
见“位点”引理 \ref{sites-lemma-morphism-sites-topoi}。直像的描述由这些引文显然得出。

\medskip\noindent
(3) 由下述事实得出：$\pi_S$ 与 $\pi_T$ 由包含函子给出，而 $f_{small}$ 与
$f_{big}$ 由基变换函子 $U \mapsto U \times_S T$ 给出。

\medskip\noindent
(4) 由 (3) 前复合 $i_T$ 得出。
\end{proof}

\noindent
在该引理的情形中，采用定义 \ref{proetale-definition-restriction-small-proetale} 的术语，
对于 $T$ 的大 pro-\'etale 位点上的层 $\mathcal{F}$，有
\begin{equation}
\label{proetale-equation-compare-big-small}
(f_{big, *}\mathcal{F})|_{S_\proetale} =
f_{small, *}(\mathcal{F}|_{T_\proetale}),
\end{equation}
此等式由引理中位点图的交换性显然得出，因为到 $T$（分别地，$S$）的小
pro-\'etale 位点的限制由 $\pi_{T, *}$（分别地，$\pi_{S, *}$）给出。涉及拉回与
限制的类似公式并不成立。

\begin{lemma}
\label{proetale-lemma-composition-proetale}
给定 $\Sch_\proetale$ 中的概形 $X$、$Y$、$Y$ 以及态射 $f : X \to Y$、
$g : Y \to Z$，有 $g_{big} \circ f_{big} = (g \circ f)_{big}$ 与
$g_{small} \circ f_{small} = (g \circ f)_{small}$。
\end{lemma}

\begin{proof}
对于大位点上的函子，这由引理 \ref{proetale-lemma-morphism-big} 中对直像与拉回的简单描述
得出。对于小位点上的函子，这由引理 \ref{proetale-lemma-morphism-big-small} 中对直像函子的
描述得出。
\end{proof}

\begin{lemma}
\label{proetale-lemma-morphism-big-small-cartesian-diagram}
设 $\Sch_\proetale$ 为大 pro-\'etale 位点。考虑 $\Sch_\proetale$ 中的笛卡尔图
$$
\xymatrix{
T' \ar[r]_{g'} \ar[d]_{f'} & T \ar[d]^f \\
S' \ar[r]^g & S
}
$$
则 $i_g^{-1} \circ f_{big, *} = f'_{small, *} \circ (i_{g'})^{-1}$，并且
$g_{big}^{-1} \circ f_{big, *} = f'_{big, *} \circ (g'_{big})^{-1}$。
\end{lemma}

\begin{proof}
由于该图是笛卡尔的，对 $U'/S'$ 有 $U' \times_{S'} T' = U' \times_S T$。
因此，$i_g^{-1} \circ f_{big, *}$ 与
$f'_{small, *} \circ (i_{g'})^{-1}$ 都将 $(\Sch/T)_\proetale$ 上的层
$\mathcal{F}$ 送到 $S'_\proetale$ 上的层
$U' \mapsto \mathcal{F}(U' \times_{S'} T')$（使用引理
\ref{proetale-lemma-put-in-T} 与 \ref{proetale-lemma-morphism-big}）。第二个等式可同样证明，也可由
非常一般的“位点”引理 \ref{sites-lemma-localize-morphism} 推出。
\end{proof}

\noindent
可以把 $S$ 的大 pro-\'etale 位点上的层看作一族层：对每个 $S$ 上的概形，取其小
pro-\'etale 位点上的一个层。

\begin{lemma}
\label{proetale-lemma-characterize-sheaf-big}
设 $S$ 为大 pro-\'etale 位点 $\Sch_\proetale$ 所包含的概形。大 pro-\'etale 位点
$(\Sch/S)_\proetale$ 上的层 $\mathcal{F}$ 由下列数据给出：
\begin{enumerate}
\item 对每个 $T/S \in \Ob((\Sch/S)_\proetale)$，取 $T_\proetale$ 上的层
$\mathcal{F}_T$；
\item 对 $(\Sch/S)_\proetale$ 中的每个 $f : T' \to T$，取映射
$c_f : f_{small}^{-1}\mathcal{F}_T \to \mathcal{F}_{T'}$。
\end{enumerate}
这些数据须满足下列条件：
\begin{enumerate}
\item[(a)] 对 $(\Sch/S)_\proetale$ 中任意 $f : T' \to T$ 与
$g : T'' \to T'$，复合 $c_g \circ g_{small}^{-1}c_f$ 等于 $c_{f \circ g}$；
\item[(b)] 若 $(\Sch/S)_\proetale$ 中的 $f : T' \to T$ 是弱 \'etale 的，
则 $c_f$ 是同构。
\end{enumerate}
\end{lemma}

\begin{proof}
与“拓扑”引理 \ref{topologies-lemma-characterize-sheaf-big-etale} 的证明相同。
\end{proof}

\begin{lemma}
\label{proetale-lemma-alternative}
设 $S$ 为概形。以 $S_{affine, \proetale}$ 表示 $S_\proetale$ 中由仿射对象组成的
全子范畴。$S_{affine, \proetale}$ 的覆盖取为标准 pro-\'etale 覆盖；见定义
\ref{proetale-definition-standard-proetale}。则限制
$$
\mathcal{F} \longmapsto \mathcal{F}|_{S_{affine, \etale}}
$$
定义拓扑斯等价 $\Sh(S_\proetale) \cong \Sh(S_{affine, \proetale})$。
\end{lemma}

\begin{proof}
这可直接由定义证明，是一道很好的练习。不过，它也立即由“位点”引理
\ref{sites-lemma-equivalence} 得出：只须验证包含函子
$S_{affine, \proetale} \to S_\proetale$ 是特殊余连续函子（见“位点”定义
\ref{sites-definition-special-cocontinuous-functor}）。
\end{proof}

\begin{lemma}
\label{proetale-lemma-affine-alternative}
设 $S$ 为仿射概形。以 $S_{app}$ 表示 $S_\proetale$ 中由满足
$\mathcal{O}(S) \to \mathcal{O}(U)$ 为 ind-\'etale 的仿射对象 $U$ 组成的全子范畴。
$S_{app}$ 的覆盖取为标准 pro-\'etale 覆盖；见定义
\ref{proetale-definition-standard-proetale}。则限制
$$
\mathcal{F} \longmapsto \mathcal{F}|_{S_{app}}
$$
定义拓扑斯等价 $\Sh(S_\proetale) \cong \Sh(S_{app})$。
\end{lemma}

\begin{proof}
由引理 \ref{proetale-lemma-alternative}，可用 $S_{affine, \proetale}$ 代替
$S_\proetale$。只须验证包含函子 $S_{app} \to S_{affine, \proetale}$ 是特殊
余连续函子，本引理便由“位点”引理 \ref{sites-lemma-equivalence} 得出；见“位点”
定义 \ref{sites-definition-special-cocontinuous-functor}。“位点”引理
\ref{sites-lemma-equivalence} 的各条件立即由定义及下列事实得出：(a)
$S_{affine, \proetale}$ 的任意对象 $U$ 都有一个覆盖 $\{V \to U\}$，其中 $V$
在 $X$ 上 ind-\'etale（命题 \ref{proetale-proposition-weakly-etale}）；以及 (b) 函子 $u$
全忠实。
\end{proof}

\begin{lemma}
\label{proetale-lemma-proetale-subcanonical}
设 $S$ 为概形。pro-\'etale 位点 $\Sch_\proetale$、$S_\proetale$、
$(\Sch/S)_\proetale$、$S_{affine, \proetale}$ 与
$(\textit{Aff}/S)_\proetale$ 上的拓扑均为次典范的。
\end{lemma}

\begin{proof}
合用引理 \ref{proetale-lemma-recognize-proetale-covering} 与“下降”引理
\ref{descent-lemma-fpqc-universal-effective-epimorphisms}。
\end{proof}





\section{弱可缩对象}
\label{proetale-section-weakly-contractible}

\noindent
本节证明一个关键事实：我们的 pro-\'etale 位点含有许多弱可缩对象。事实上，引理
\ref{proetale-lemma-get-many-weakly-contractible} 的证明正是定义
\ref{proetale-definition-big-proetale-site} 中函数 $Bound$ 采用该形式的原因（不过，对忽略
集合论问题的读者而言，这条信息没有实际内容）。

\noindent
先用 pro-\'etale 覆盖表述 w-可缩环的概念。

\begin{lemma}
\label{proetale-lemma-w-contractible-proetale-cover}
设 $T = \Spec(A)$ 为仿射概形。下列条件等价：
\begin{enumerate}
\item $A$ 是 w-可缩的；
\item $T$ 的每个 pro-\'etale 覆盖都可由形如
$T = \coprod_{i = 1, \ldots, n} U_i$ 的 Zariski 覆盖加细。
\end{enumerate}
\end{lemma}

\begin{proof}
假设 $A$ 是 w-可缩的。由引理 \ref{proetale-lemma-proetale-affine}，只须证明 $T$ 的某个
Zariski 覆盖能够加细每个标准 pro-\'etale 覆盖
$\{f_i : T_i \to T\}_{i = 1, \ldots, n}$。态射 $\coprod T_i \to T$ 是仿射概形间
满的弱 \'etale 态射。因此由定义 \ref{proetale-definition-w-contractible}，存在 $T$ 上的态射
$\sigma : T \to \coprod T_i$。于是 Zariski 覆盖
$T = \coprod \sigma^{-1}(T_i)$ 加细 $\{f_i : T_i \to T\}$。

\medskip\noindent
反之，假设 (2)。若 $A \to B$ 忠实平坦且弱 \'etale，则
$\{\Spec(B) \to T\}$ 是 pro-\'etale 覆盖。因此存在 Zariski 覆盖
$T = \coprod U_i$ 以及 $T$ 上的态射 $U_i \to \Spec(B)$。由于
$T = \coprod U_i$，得到 $T \to \Spec(B)$，也即一个 $A$-代数映射 $B \to A$。
这说明 $A$ 是 w-可缩的。
\end{proof}

\begin{lemma}
\label{proetale-lemma-w-contractible-is-weakly-contractible}
设 $\Sch_\proetale$ 为定义 \ref{proetale-definition-big-proetale-site} 所述的大
pro-\'etale 位点，并设 $T = \Spec(A)$ 为 $\Sch_\proetale$ 的仿射对象。下列条件等价：
\begin{enumerate}
\item $A$ 是 w-可缩的；
\item $T$ 是 $\Sch_\proetale$ 的弱可缩对象（“位点”定义
\ref{sites-definition-w-contractible}）；
\item $T$ 的每个 pro-\'etale 覆盖都可由形如
$T = \coprod_{i = 1, \ldots, n} U_i$ 的 Zariski 覆盖加细。
\end{enumerate}
\end{lemma}

\begin{proof}
引理 \ref{proetale-lemma-w-contractible-proetale-cover} 已证明 (1) 与 (3) 等价。

\medskip\noindent
假设 (3)，并设 $\mathcal{F} \to \mathcal{G}$ 是 $\Sch_\proetale$ 上层的满射。
取 $s \in \mathcal{G}(T)$。为证明 (2)，我们将证明 $s$ 属于
$\mathcal{F}(T) \to \mathcal{G}(T)$ 的像。可找到 $\Sch_\proetale$ 的覆盖
$\{T_i \to T\}$，使 $s$ 在 $T_i$ 上提升为 $\mathcal{F}$ 的截面（“位点”定义
\ref{sites-definition-sheaves-injective-surjective}）。由 (3)，可以假设有由既开又闭的
子集组成的有限覆盖 $T = \coprod_{j = 1, \ldots, m} U_j$，并有映到
$s|_{U_j}$ 的 $t_j \in \mathcal{F}(U_j)$。由于 Zariski 覆盖是
$\Sch_\proetale$ 中的覆盖（引理 \ref{proetale-lemma-etale-proetale}），可知
$\mathcal{F}(T) = \prod \mathcal{F}(U_j)$。因此
$t = (t_1, \ldots, t_m) \in \mathcal{F}(T)$ 是映到 $s$ 的截面。

\medskip\noindent
假设 (2)。取命题 \ref{proetale-proposition-find-w-contractible} 所述的 $A \to D$。
则 $\{V \to T\}$ 是 $\Sch_\proetale$ 的覆盖。（注意，由注
\ref{proetale-remark-size-w-contractible}、定义 \ref{proetale-definition-big-proetale-site} 中对函数
$Bound$ 的选取以及“集合”引理 \ref{sets-lemma-bound-size} 对仿射概形大小的计算，
$V = \Spec(D)$ 是 $\Sch_\proetale$ 的对象。）由于 $\Sch_\proetale$ 上的拓扑是
次典范的（引理 \ref{proetale-lemma-proetale-subcanonical}），$h_V \to h_T$ 是层的满射
（“位点”引理 \ref{sites-lemma-covering-surjective-after-sheafification}）。由于假设
$T$ 弱可缩，存在元素 $f \in h_V(T) = \Mor(T, V)$，其在 $h_T(T)$ 中的像为
$\text{id}_T$。因此 $A \to D$ 有回缩 $\sigma : D \to A$。现若 $A \to B$
忠实平坦且弱 \'etale，则 $D \to D \otimes_A B$ 具有相同性质，故存在回缩
$D \otimes_A B \to D$；将其与 $\sigma$ 复合，得到 $A \to B$ 的回缩
$B \to D \otimes_A B \to D \to A$。因此 $A$ 是 w-可缩的，(1) 成立。
\end{proof}

\begin{lemma}
\label{proetale-lemma-get-many-weakly-contractible}
设 $\Sch_\proetale$ 为定义 \ref{proetale-definition-big-proetale-site} 所述的大
pro-\'etale 位点。对 $\Sch_\proetale$ 的每个对象 $T$，存在 $\Sch_\proetale$
中的覆盖 $\{T_i \to T\}$，使每个 $T_i$ 都是仿射的，并且是某个 w-可缩环的谱。
特别地，$T_i$ 在 $\Sch_\proetale$ 中弱可缩。
\end{lemma}

\begin{proof}
对于不关心集合论问题的读者，本引理是引理
\ref{proetale-lemma-w-contractible-is-weakly-contractible} 与命题
\ref{proetale-proposition-find-w-contractible} 的平凡推论。下面给出细节。选取仿射开覆盖
$T = \bigcup U_i$，并写成 $U_i = \Spec(A_i)$。依照命题
\ref{proetale-proposition-find-w-contractible}，选取忠实平坦的 ind-\'etale 环映射
$A_i \to D_i$，使 $D_i$ 为 w-可缩的。态射族 $\{\Spec(D_i) \to T\}$ 是
pro-\'etale 覆盖。若能证明 $\Spec(D_i)$ 同构于 $\Sch_\proetale$ 的某个对象，记为
$T_i$，则由定义 \ref{proetale-definition-big-proetale-site} 中 $\Sch_\proetale$ 的构造，
更确切地说，由最后一步对“集合”引理 \ref{sets-lemma-coverings-site} 的应用，
$\{T_i \to T\}$ 将与 $\Sch_\proetale$ 的一个覆盖组合等价。由定义
\ref{proetale-definition-big-proetale-site} 中 $\Sch_\proetale$ 的构造，更确切地说，由第
(3) 步对“集合”引理 \ref{sets-lemma-construct-category} 的应用，为证明
$\Spec(D_i)$ 同构于 $\Sch_\proetale$ 的一个对象，只须证明
$|D_i| \leq Bound(\text{size}(T))$。由“集合”引理
\ref{sets-lemma-bound-affine} 与 \ref{sets-lemma-bound-finite-type}，有
$|A_i| \leq \text{size}(U_i) \leq \text{size}(T)$；故由注
\ref{proetale-remark-size-w-contractible}，有
$|D_i| \leq \kappa^{2^{2^{2^\kappa}}}$，其中 $\kappa = \text{size}(T)$。因此，
由定义 \ref{proetale-definition-big-proetale-site} 中对函数 $Bound$ 的选取，结论成立。
\end{proof}

\begin{lemma}
\label{proetale-lemma-proetale-enough-w-contractible}
设 $S$ 为概形。pro-\'etale 位点 $S_\proetale$、$(\Sch/S)_\proetale$、
$S_{affine, \proetale}$ 与 $(\textit{Aff}/S)_\proetale$，以及在 $S$ 仿射时的
$S_{app}$，都有足够多（仿射的）拟紧弱可缩对象；见“位点”定义
\ref{sites-definition-w-contractible}。
\end{lemma}

\begin{proof}
立即由引理 \ref{proetale-lemma-get-many-weakly-contractible} 得出。
\end{proof}

\begin{lemma}
\label{proetale-lemma-weakly-contractible-cover}
设 $S$ 为概形。pro-\'etale 位点 $\Sch_\proetale$、$S_\proetale$、
$(\Sch/S)_\proetale$ 具有下列性质：对任意对象 $U$，存在覆盖
$\{V \to U\}$，其中 $V$ 为弱可缩对象。若 $U$ 拟紧，则可选取 $V$ 为仿射且
弱可缩的。
\end{lemma}

\begin{proof}
设 $V = \coprod_{j \in J} V_j$ 是 $(\Sch/S)_\proetale$ 的对象，并且是弱可缩对象
$V_j$ 的不交并。由于不交并分解是 pro-\'etale 覆盖，对任意 pro-\'etale 层
$\mathcal{F}$，有 $\mathcal{F}(V) = \prod_{j \in J} \mathcal{F}(V_j)$。设
$\mathcal{F} \to \mathcal{G}$ 是集合层的满射。由于 $V_j$ 弱可缩，映射
$\mathcal{F}(V_j) \to \mathcal{G}(V_j)$ 是满射；见“位点”定义
\ref{sites-definition-w-contractible}。因此，作为集合满射之积，
$\mathcal{F}(V) \to \mathcal{G}(V)$ 是满射，从而 $V$ 弱可缩。

\medskip\noindent
依照引理 \ref{proetale-lemma-get-many-weakly-contractible}，选取覆盖
$\{U_i \to U\}_{i \in I}$，其中 $U_i$ 仿射且弱可缩。取
$V = \coprod_{i \in I} U_i$（这里有一个集合论问题，将在下文处理）。则
$\{V \to U\}$ 是所需的、由弱可缩对象给出的 pro-\'etale 覆盖（要验证它是覆盖，
使用引理 \ref{proetale-lemma-recognize-proetale-covering}）。若 $U$ 拟紧，则由引理
\ref{proetale-lemma-recognize-proetale-covering} 立即可知，可选取有限子集 $I' \subset I$，
使 $\{U_i \to U\}_{i \in I'}$ 仍为覆盖；于是
$\{\coprod_{i \in I'} U_i \to U\}$ 是所需的、由仿射弱可缩对象给出的覆盖。

\medskip\noindent
本段处理集合论问题，我们强烈建议读者跳过。为确定该不交并位于我们的部分宇宙中，
需要控制指标集 $I$ 的基数。由引理 \ref{proetale-lemma-get-many-weakly-contractible} 的
证明中覆盖 $\{U_i \to U\}_{i \in I}$ 的构造立即可见
$|I| \leq \text{size}(U)$，其中概形的大小如“集合”第
\ref{sets-section-categories-schemes} 节所定义。此外，对每个 $i$，有
$\text{size}(U_i) \leq Bound(\text{size}(U))$；这由引理
\ref{proetale-lemma-get-many-weakly-contractible} 的证明中对
$\Gamma(U_i, \mathcal{O}_{U_i})$ 基数的估计以及“集合”引理
\ref{sets-lemma-bound-affine} 得出。因而由“集合”引理 \ref{sets-lemma-bound-size}，
有 $\text{size}(\coprod_{i \in I} U_i)) \leq Bound(\text{size}(U))$。所以，由通过
“集合”引理 \ref{sets-lemma-construct-category} 对大 pro-\'etale 位点的构造，
$\coprod_{i \in I} U_i$ 同构于我们位点的一个对象，证明完毕。
\end{proof}






\section{弱可缩超覆盖}
\label{proetale-section-hypercovering}

\noindent
第 \ref{proetale-section-weakly-contractible} 节的结果蕴含由弱可缩对象组成的超覆盖之存在性。

\begin{lemma}
\label{proetale-lemma-w-contractible-hypercovering}
设 $X$ 为概形。
\begin{enumerate}
\item 对 $X_\proetale$ 的每个对象 $U$，存在 $X_\proetale$ 中 $U$ 的超覆盖
$K$，使每一项 $K_n$ 都由 $X_\proetale$ 中覆盖 $U$ 的一个弱可缩对象组成。
\item 对 $X_\proetale$ 的每个拟紧拟分离对象 $U$，存在 $X_\proetale$ 中 $U$
的超覆盖 $K$，使每一项 $K_n$ 都由 $X_\proetale$ 中覆盖 $U$ 的一个仿射弱可缩
对象组成。
\end{enumerate}
\end{lemma}

\begin{proof}
令 $\mathcal{B} \subset \Ob(X_\proetale)$ 为 $X_\proetale$ 的弱可缩对象集。由引理
\ref{proetale-lemma-weakly-contractible-cover}，$X_\proetale$ 的每个对象 $T$ 都有覆盖
$\{T_i \to T\}_{i \in I}$，其中 $I$ 有限且 $T_i \in \mathcal{B}$。由“超覆盖”引理
\ref{hypercovering-lemma-w-contractible}，得到 $U$ 的超覆盖 $K$，使得
$K_n = \{U_{n, i}\}_{i \in I_n}$，其中 $I_n$ 有限且 $U_{n, i}$ 弱可缩。随后可将
$K$ 替换为 $U$ 的如下超覆盖：其第 $n$ 次项由 $\{U_n\}$ 给出，其中
$U_n = \coprod_{i \in I_n} U_{n, i}$。“超覆盖”注
\ref{hypercovering-remark-take-unions-hypercovering-X} 保证此替换可行。

\medskip\noindent
令 $X_{qcqs, \proetale} \subset X_\proetale$ 为由拟紧拟分离对象组成的全子范畴，
并将 $X_{qcqs, \proetale}$ 的覆盖取为有限 pro-\'etale 覆盖。于是
$X_{qcqs, \proetale}$ 是有纤维积的位点，且包含函子
$X_{qcqs, \proetale} \to X_\proetale$ 连续并保持纤维积。特别地，若 $K$ 是
$X_{qcqs, \proetale}$ 中对象 $U$ 的超覆盖，则由“超覆盖”引理
\ref{hypercovering-lemma-hypercovering-continuous-functor}，$K$ 也是
$X_\proetale$ 中 $U$ 的超覆盖。令
$\mathcal{B} \subset \Ob(X_{qcqs, \proetale})$ 为仿射弱可缩对象集。由引理
\ref{proetale-lemma-get-many-weakly-contractible} 以及仿射概形的有限并仍为仿射这一事实，
对 $X_{qcqs, \proetale}$ 的每个对象 $U$，存在 $X_{qcqs, \proetale}$ 的覆盖
$\{V \to U\}$，其中 $V \in \mathcal{B}$。由“超覆盖”引理
\ref{hypercovering-lemma-w-contractible}，得到 $U$ 的超覆盖 $K$，使得
$K_n = \{U_{n, i}\}_{i \in I_n}$，其中 $I_n$ 有限且 $U_{n, i}$ 仿射而弱可缩。
随后可将 $K$ 替换为 $U$ 的如下超覆盖：其第 $n$ 次项由 $\{U_n\}$ 给出，其中
$U_n = \coprod_{i \in I_n} U_{n, i}$。“超覆盖”注
\ref{hypercovering-remark-take-unions-hypercovering-X} 保证此替换可行。
\end{proof}

\noindent
在下一个引理中，我们使用与位点中超覆盖 $K$ 相伴的 {\v C}ech 复形
$s(\mathcal{F}(K))$；见“超覆盖”第 \ref{hypercovering-section-hyper-cech} 节。
若 $K$ 是 $U$ 的超覆盖且 $K_n = \{U_n \to U\}$，则 {\v C}ech 复形形如
$$
s(\mathcal{F}(K)) =
\left(
\mathcal{F}(U_0) \to \mathcal{F}(U_1) \to \mathcal{F}(U_2) \to \ldots
\right)
$$
其中 $s(\mathcal{F}(U_n))$ 置于上同调次数 $n$。

\begin{lemma}
\label{proetale-lemma-compute-cohomology}
设 $X$ 为概形。设 $E \in D^+(X_\proetale)$ 由阿贝尔层的下有界复形
$\mathcal{E}^\bullet$ 表示。设 $K$ 是 $U \in \Ob(X_\proetale)$ 的超覆盖，且
$K_n = \{U_n \to U\}$，其中 $U_n$ 是 $X_\proetale$ 的弱可缩对象。则
$$
R\Gamma(U, E) = \text{Tot}(s(\mathcal{E}^\bullet(K)))
$$
在 $D(\textit{Ab})$ 中成立。
\end{lemma}

\begin{proof}
若 $\mathcal{E}$ 是 $X_\proetale$ 上的阿贝尔层，则“超覆盖”引理
\ref{hypercovering-lemma-cech-spectral-sequence} 的谱序列蕴含
$$
R\Gamma(X_\proetale, \mathcal{E}) = s(\mathcal{E}(K))
$$
因为任意层在 $U_n$ 上的高阶上同调群均消失；见“位点上同调”引理
\ref{sites-cohomology-lemma-w-contractible}。

\medskip\noindent
若 $\mathcal{E}^\bullet$ 下有界，则可选取内射分解
$\mathcal{E}^\bullet \to \mathcal{I}^\bullet$，并考虑复形映射
$$
\text{Tot}(s(\mathcal{E}^\bullet(K)))
\longrightarrow
\text{Tot}(s(\mathcal{I}^\bullet(K)))
$$
对每个 $n$，映射 $\mathcal{E}^\bullet(U_n) \to \mathcal{I}^\bullet(U_n)$ 是拟同构，
因为在 $U_n$ 上取截面是正合的。因此由“同调”引理
\ref{homology-lemma-first-quadrant-ss} 的一个谱序列，上述映射是拟同构。利用第一段的
结果可知，对每个 $p$，复形 $s(\mathcal{I}^p(K))$ 在次数 $n > 0$ 无同调，
并在次数 $0$ 计算 $\mathcal{I}^p(U)$。于是“同调”引理
\ref{homology-lemma-first-quadrant-ss} 的另一个谱序列说明
$\text{Tot}(s(\mathcal{I}^\bullet(K)))$ 计算
$R\Gamma(U, E) = \mathcal{I}^\bullet(U)$。
\end{proof}

\begin{lemma}
\label{proetale-lemma-quasi-compact-quasi-separated-commutes-direct-sums}
设 $X$ 为拟紧拟分离概形。函子
$R\Gamma(X, -) : D^+(X_\proetale) \to D(\textit{Ab})$ 与直和及同伦余极限交换。
\end{lemma}

\begin{proof}
该断言含义如下：设有 $D^+(X_\proetale)$ 的一族对象 $E_i$，使得
$\bigoplus E_i$ 是 $D^+(X_\proetale)$ 的对象。则
$R\Gamma(X, \bigoplus E_i) = \bigoplus R\Gamma(X, E_i)$。为验证这一点，选取
$X$ 的超覆盖 $K$，满足 $K_n = \{U_n \to X\}$，其中 $U_n$ 是仿射弱可缩概形；
见引理 \ref{proetale-lemma-w-contractible-hypercovering}。取整数 $N$，使得当 $p < N$ 时
$H^p(E_i) = 0$。选取表示 $E_i$ 的阿贝尔层复形 $\mathcal{E}_i^\bullet$，使得
当 $p < N$ 时 $\mathcal{E}_i^p = 0$。逐项直和
$\bigoplus \mathcal{E}_i^\bullet$ 在 $D(X_\proetale)$ 中表示 $\bigoplus E_i$；
见“内射对象”引理 \ref{injectives-lemma-derived-products}。由引理
\ref{proetale-lemma-compute-cohomology}，有
$$
R\Gamma(X, \bigoplus E_i) =
\text{Tot}(s((\bigoplus \mathcal{E}^\bullet_i)(K)))
$$
以及
$$
R\Gamma(X, E_i) = \text{Tot}(s(\mathcal{E}^\bullet_i(K)))
$$
由于每个 $U_n$ 拟紧，有
$$
\text{Tot}(s((\bigoplus \mathcal{E}^\bullet_i)(K))) =
\bigoplus \text{Tot}(s(\mathcal{E}^\bullet_i(K)))
$$
这由“位点上的模”引理 \ref{sites-modules-lemma-sections-over-quasi-compact} 得出。
关于同伦余极限的断言形式地源于如下事实：$R\Gamma$ 是三角范畴间的正合函子，
且它与直和交换，这一点刚刚已经证明。
\end{proof}

\begin{remark}
\label{proetale-remark-extend-to-all}
设 $X$ 为概形。由于 $X_\proetale$ 有足够多的弱可缩对象，由“位点上同调”命题
\ref{sites-cohomology-proposition-enough-weakly-contractibles}，对
$D(X_\proetale)$ 中所有 $K$，有 $K = R\lim \tau_{\geq -n}K$。又由“内射对象”引理
\ref{injectives-lemma-RF-commutes-with-Rlim}，$R\Gamma$ 与 $R\lim$ 交换，故
$$
R\Gamma(X, K) = R\lim R\Gamma(X, \tau_{\geq -n}K)
$$
在 $D(\textit{Ab})$ 中成立。这有时使我们能够把关于下有界复形的结果推广到所有复形。
\end{remark}







\section{紧生成}
\label{proetale-section-compact-generation}

\noindent
本节证明，与我们的 pro-\'etale 位点相伴的若干导出范畴是“导出范畴”定义
\ref{derived-definition-compactly-generated} 所述的紧生成范畴。

\begin{lemma}
\label{proetale-lemma-enough-compact-proetale}
设 $S$ 为概形，$\Lambda$ 为环。
\begin{enumerate}
\item $D(S_\proetale)$ 是紧生成的；
\item $D(S_\proetale, \Lambda)$ 是紧生成的；
\item 对 $S_\proetale$ 上任意环层 $\mathcal{A}$，
$D(S_\proetale, \mathcal{A})$ 是紧生成的；
\item $D((\Sch/S)_\proetale)$ 是紧生成的；
\item $D((\Sch/S)_\proetale, \Lambda)$ 是紧生成的；
\item 对 $(\Sch/S)_\proetale$ 上任意环层 $\mathcal{A}$，
$D((\Sch/S)_\proetale, \mathcal{A})$ 是紧生成的。
\end{enumerate}
\end{lemma}

\begin{proof}
证明 (3)。设 $U$ 是 $S_\proetale$ 的仿射弱可缩对象。则
$j_{U!}\mathcal{A}_U$ 是导出范畴 $D(S_\proetale, \mathcal{A})$ 的紧对象；见
“位点上同调”引理
\ref{sites-cohomology-lemma-quasi-compact-weakly-contractible-compact}。选取集合 $I$，
并对每个 $i \in I$ 选取 $S_\proetale$ 的仿射弱可缩对象 $U_i$，使
$S_\proetale$ 的每个仿射弱可缩对象都同构于某个 $U_i$。这是可行的，因为
$\Ob(S_\proetale)$ 是集合。为完成 (3) 的证明，只须证明
$\bigoplus j_{U_i, !}\mathcal{A}_{U_i}$ 是 $D(S_\proetale, \mathcal{A})$ 的生成元；
见“导出范畴”定义 \ref{derived-definition-generators}。为此，设 $K$ 是
$D(S_\proetale, \mathcal{A})$ 的非零对象。则在我们的位点 $S_\proetale$ 中存在
对象 $T$，且 $H^n(K)(T)$ 中存在非零元素 $\xi$。换言之，$\xi$ 是 $K$ 的第
$n$ 个上同调层的非零截面。可以假设 $K$ 由 $\mathcal{A}$-模层的复形
$\mathcal{K}^\bullet$ 表示，而 $\xi$ 是满足 $\text{d}(s) = 0$ 的截面
$s \in \mathcal{K}^n(T)$ 的类。确切地说，$\xi$ 局部表示为一个截面的类（因此，
用 $T$ 的某个覆盖成员替换 $T$ 后便得到此结论）。接着，依照引理
\ref{proetale-lemma-get-many-weakly-contractible} 选取覆盖
$\{T_j \to T\}_{j \in J}$。由于 $H^n(K)$ 是层，对某个 $j$，限制
$\xi|_{T_j}$ 仍非零。因此 $s|_{T_j}$ 在 $D(S_\proetale, \mathcal{A})$ 中定义非零映射
$j_{T_j, !}\mathcal{A}_{T_j} \to K$。由于对某个 $i \in I$ 有
$T_j \cong U_i$，结论得证。

\medskip\noindent
完全相同的论证适用于 $S$ 的大 pro-\'etale 位点。
\end{proof}







\section{拓扑的比较}
\label{proetale-section-compare-topologies}

\noindent
本节是“\'Etale 上同调”第
\ref{etale-cohomology-section-compare-topologies} 节的类似版本。

\begin{lemma}
\label{proetale-lemma-presheaf-value-weakly-contractible}
设 $X$ 为概形，$\mathcal{F}$ 为 $X_\proetale$ 上的集合预层，并且它将有限不交并
送到积。若 $W$ 是 $X_\proetale$ 的仿射弱可缩对象，则
$\mathcal{F}^\#(W) = \mathcal{F}(W)$。
\end{lemma}

\begin{proof}
回忆 $\mathcal{F}^\#$ 等于 $(\mathcal{F}^+)^+$；见“位点”定理
\ref{sites-theorem-plus}。这里 $\mathcal{F}^+$ 是将 $X_\proetale$ 的对象 $U$ 送到
$\colim H^0(\mathcal{U}, \mathcal{F})$ 的预层，其中余极限取遍 $U$ 的所有
pro-\'etale 覆盖 $\mathcal{U}$。因此只须证明：(a) $\mathcal{F}^+$ 将有限不交并
送到积；(b) 它将 $W$ 送到 $\mathcal{F}(W)$。若 $U = U_1 \amalg U_2$，则给定
$U$ 的 pro-\'etale 覆盖 $\mathcal{U} = \{f_j : V_j \to U\}$，得到 pro-\'etale 覆盖
$\mathcal{U}_i = \{f_j^{-1}(U_i) \to U_i\}$，并且显然有
$$
H^0(\mathcal{U}, \mathcal{F}) =
H^0(\mathcal{U}_1, \mathcal{F}) \times
H^0(\mathcal{U}_2, \mathcal{F})
$$
因为 $\mathcal{F}$ 将有限不交并送到积（这包括 $\mathcal{F}$ 将空概形送到单点集
这一条件）。这证明了 (a)。最后，由引理
\ref{proetale-lemma-w-contractible-is-weakly-contractible}，$W$ 的任意 pro-\'etale 覆盖都可由
有限不交并分解 $W = W_1 \amalg \ldots W_n$ 加细。因此恰因 $\mathcal{F}$ 在
$W$ 上的值是 $\mathcal{F}$ 在各 $W_j$ 上之值的积，有
$\mathcal{F}^+(W) = \mathcal{F}(W)$。这证明了 (b)。
\end{proof}

\begin{lemma}
\label{proetale-lemma-small-pullback-weakly-contractible}
设 $f : X \to Y$ 为概形态射，$\mathcal{F}$ 为 $Y_\proetale$ 上的集合层。若
$W$ 是 $X_\proetale$ 的仿射弱可缩对象，则
$$
f_{small}^{-1}\mathcal{F}(W) = \colim_{W \to V} \mathcal{F}(V)
$$
其中余极限取遍 $Y$ 上满足 $V \in Y_\proetale$ 的态射 $W \to V$。
\end{lemma}

\begin{proof}
回忆 $f_{small}^{-1}\mathcal{F}$ 是与下列预层相伴的层：
$$
u_p\mathcal{F} : U \mapsto \colim_{U \to V} \mathcal{F}(V)
$$
它在 $X_\etale$ 上；见“位点”第 \ref{sites-section-morphism-sites} 节与第
\ref{sites-section-continuous-functors} 节。记号中省略了如下事实：余极限取遍范畴
$\{U \to V, V \in Y_\proetale\}$ 的对偶范畴。由引理
\ref{proetale-lemma-presheaf-value-weakly-contractible}，只须证明 $u_p\mathcal{F}$ 将有限不交并
送到积。设 $U = U_1 \amalg U_2$ 是既开又闭的子概形之不交并。有函子
$$
\{U_1 \to V_1\} \times \{U_2 \to V_2\} \longrightarrow
\{U \to V\},\quad
(U_1 \to V_1, U_2 \to V_2) \longmapsto (U \to V_1 \amalg V_2)
$$
它是初始的（“范畴”定义 \ref{categories-definition-initial}）。因此，对偶范畴上的
相应函子是共尾的；由“范畴”引理 \ref{categories-lemma-cofinal}，可知
$u_p\mathcal{F}$ 在 $U$ 上的值，是在该积范畴上取各
$\mathcal{F}(V_1 \amalg V_2)$ 的余极限。由于 $\mathcal{F}$ 是层，它将不交并送到积，
故 $u_p\mathcal{F}$ 也如此。
\end{proof}

\begin{lemma}
\label{proetale-lemma-describe-pullback}
设 $S$ 为概形。考虑引理 \ref{proetale-lemma-at-the-bottom} 的态射
$$
\pi_S : (\Sch/S)_\proetale \longrightarrow S_\proetale
$$
。设 $\mathcal{F}$ 为 $S_\proetale$ 上的层。则 $\pi_S^{-1}\mathcal{F}$ 由规则
$$
(\pi_S^{-1}\mathcal{F})(T) = \Gamma(T_\proetale, f_{small}^{-1}\mathcal{F})
$$
给出，其中 $f : T \to S$。此外，$\pi_S^{-1}\mathcal{F}$ 对 fpqc 覆盖满足层条件。
\end{lemma}

\begin{proof}
注意，我们有态射 $i_f : \Sh(T_\proetale) \to \Sh(\Sch/S)_\proetale)$，使得作为
态射 $T_\proetale \to S_\proetale$，有 $\pi_S \circ i_f = f_{small}$；见引理
\ref{proetale-lemma-put-in-T}。由于拉回具有传递性，可得
$i_f^{-1} \pi_S^{-1}\mathcal{F} = f_{small}^{-1}\mathcal{F}$，正如所需。

\medskip\noindent
设 $\{g_i : T_i \to T\}_{i \in I}$ 为 fpqc 覆盖。最后一个断言的含义如下：给定
$T_\proetale$ 上的层 $\mathcal{G}$ 及截面
$s_i \in \Gamma(T_i, g_{i, small}^{-1}\mathcal{G})$，若它们到
$T_i \times_T T_j$ 上的拉回一致，则 $T$ 上存在 $\mathcal{G}$ 的唯一截面 $s$，
其到 $T_i$ 上的拉回等于 $s_i$。我们只证明 $T$ 仿射且覆盖由仿射概形间的单个
满平坦态射 $T' \to T$ 给出的情形，并略去一般情形到此情形的归约。

\medskip\noindent
设 $g : T' \to T$ 是仿射概形间的满平坦态射，并设
$s' \in g_{small}^{-1}\mathcal{G}(T')$ 为满足
$\text{pr}_0^*s' = \text{pr}_1^*s'$ 于 $T' \times_T T'$ 的截面。选取满的弱
\'etale 态射 $W \to T'$，其中 $W$ 仿射且弱可缩；见引理
\ref{proetale-lemma-weakly-contractible-cover}。由引理
\ref{proetale-lemma-small-pullback-weakly-contractible}，限制 $s'|_W$ 是
$\colim_{W \to U} \mathcal{G}(U)$ 的元素。选取与 $s'$ 对应的
$\phi : W \to U_0$ 与 $s_0 \in \mathcal{G}(U_0)$。选取满的弱 \'etale 态射
$V \to W \times_T W$，其中 $V$ 仿射且弱可缩。将诱导态射记为
$a, b : V \to W$。由于 $a^*(s'|_W) = b^*(s'|_W)$，且范畴
$\{V \to U, U \in T_\proetale\}$ 是余滤的（这是显然的；如有疑问，见“位点”引理
\ref{sites-lemma-directed-morphism}），两个态射
$\phi \circ a , \phi \circ b : V \to U_0$ 必须相等。由“下降”第
\ref{descent-section-fpqc-universal-effective-epimorphisms} 节中的结果（特别是“下降”
引理 \ref{descent-lemma-fpqc-universal-effective-epimorphisms}），存在唯一态射
$T \to U_0$，使 $\phi$ 是此态射与结构态射 $W \to T$ 的复合（略去一个小细节）。
于是可令 $s$ 为 $s_0$ 沿此态射的拉回。我们略去验证 $s$ 到 $T'$ 上的拉回为
$s'$。
\end{proof}









\section{大、小拓扑斯的比较}
\label{proetale-section-compare}

\noindent
本节是“\'Etale 上同调”第 \ref{etale-cohomology-section-compare} 节的类似版本。
下文常以 $\mathcal{F} \mapsto \mathcal{F}|_{S_\proetale}$ 表示与引理
\ref{proetale-lemma-at-the-bottom} 的拓扑斯态射
$i_S : \Sh(S_\proetale) \to \Sh((\Sch/S)_\proetale)$ 相应的拉回函子 $i_S^{-1}$。

\begin{lemma}
\label{proetale-lemma-compare-injectives}
设 $S$ 为概形，$T$ 为 $(\Sch/S)_\proetale$ 的对象。
\begin{enumerate}
\item 若 $\mathcal{I}$ 在 $\textit{Ab}((\Sch/S)_\proetale)$ 中内射，则
\begin{enumerate}
\item $i_f^{-1}\mathcal{I}$ 在 $\textit{Ab}(T_\proetale)$ 中内射；
\item $\mathcal{I}|_{S_\proetale}$ 在 $\textit{Ab}(S_\proetale)$ 中内射。
\end{enumerate}
\item 若 $\mathcal{I}^\bullet$ 是 $\textit{Ab}((\Sch/S)_\proetale)$ 中的 K-内射复形，
则
\begin{enumerate}
\item $i_f^{-1}\mathcal{I}^\bullet$ 是 $\textit{Ab}(T_\proetale)$ 中的 K-内射复形；
\item $\mathcal{I}^\bullet|_{S_\proetale}$ 是 $\textit{Ab}(S_\proetale)$ 中的
K-内射复形。
\end{enumerate}
\end{enumerate}
\end{lemma}

\begin{proof}
证明 (1)(a) 与 (2)(a)：$i_f^{-1}$ 是正合函子 $i_{f, !}$ 的右伴随。事实上，回忆
$i_f$ 对应余连续函子 $u : T_\proetale \to (\Sch/S)_\proetale$，后者连续并保持
纤维积与等化子；见引理 \ref{proetale-lemma-put-in-T} 及其证明。因此，由“位点上的模”引理
\ref{sites-modules-lemma-g-shriek-adjoint} 得到 $i_{f, !}$。“位点上的模”引理
\ref{sites-modules-lemma-exactness-lower-shriek} 证明了它是正合的。于是由“同调”引理
\ref{homology-lemma-adjoint-preserve-injectives} 与“导出范畴”引理
\ref{derived-lemma-adjoint-preserve-K-injectives}，得到 (1)(a) 与 (2)(a)。

\medskip\noindent
由于 $i_S = i_{\text{id}_S}$，(1)(b) 与 (2)(b) 分别是 (1)(a) 与 (2)(a) 的特例。
\end{proof}

\begin{lemma}
\label{proetale-lemma-compare-higher-direct-image}
设 $f : T \to S$ 为概形态射。对 $D((\Sch/T)_\proetale)$ 中的 $K$，有
$$
(Rf_{big, *}K)|_{S_\proetale} = Rf_{small, *}(K|_{T_\proetale})
$$
在 $D(S_\proetale)$ 中成立。更一般地，设
$S' \in \Ob((\Sch/S)_\proetale)$，其结构态射为 $g : S' \to S$。考虑纤维积图
$$
\xymatrix{
T' \ar[r]_{g'} \ar[d]_{f'} & T \ar[d]^f \\
S' \ar[r]^g & S
}
$$
则对 $D((\Sch/T)_\proetale)$ 中的 $K$，有
$$
i_g^{-1}(Rf_{big, *}K) = Rf'_{small, *}(i_{g'}^{-1}K)
$$
在 $D(S'_\proetale)$ 中成立，并且
$$
g_{big}^{-1}(Rf_{big, *}K) = Rf'_{big, *}((g'_{big})^{-1}K)
$$
在 $D((\Sch/S')_\proetale)$ 中成立。
\end{lemma}

\begin{proof}
选取表示 $K$ 的阿贝尔层 K-内射复形后，第一个等式由引理
\ref{proetale-lemma-compare-injectives} 与 (\ref{proetale-equation-compare-big-small}) 得出。选取表示
$K$ 的阿贝尔层 K-内射复形后，第二个等式由引理 \ref{proetale-lemma-compare-injectives} 与
引理 \ref{proetale-lemma-morphism-big-small-cartesian-diagram} 得出。类似地，选取表示 $K$ 的
阿贝尔层 K-内射复形后，第三个等式由“位点上同调”引理
\ref{sites-cohomology-lemma-cohomology-of-open}、
\ref{sites-cohomology-lemma-restrict-K-injective-to-open} 以及引理
\ref{proetale-lemma-morphism-big-small-cartesian-diagram} 得出。
\end{proof}

\noindent
设 $S$ 为概形，$\mathcal{H}$ 为 $(\Sch/S)_\proetale$ 上的阿贝尔层。回忆
$H^n_\proetale(U, \mathcal{H})$ 表示 $\mathcal{H}$ 在
$(\Sch/S)_\proetale$ 的对象 $U$ 上的上同调。

\begin{lemma}
\label{proetale-lemma-compare-cohomology-big-small}
设 $f : T \to S$ 为概形态射。对 $D(S_\proetale)$ 中的 $K$，有
$$
H^n_\proetale(S, \pi_S^{-1}K) = H^n(S_\proetale, K)
$$
以及
$$
H^n_\proetale(T, \pi_S^{-1}K) = H^n(T_\proetale, f_{small}^{-1}K).
$$
对 $D((\Sch/S)_\proetale)$ 中的 $M$，有
$$
H^n_\proetale(T, M) = H^n(T_\proetale, i_f^{-1}M).
$$
\end{lemma}

\begin{proof}
为证明最后一个等式，用阿贝尔层的 K-内射复形表示 $M$，应用引理
\ref{proetale-lemma-compare-injectives} 并展开定义即可。由于
$i_f^{-1} \circ \pi_S^{-1} = f_{small}^{-1}$，第二个等式由此得出。第一个等式是
第二个等式的特例。
\end{proof}

\begin{lemma}
\label{proetale-lemma-cohomological-descent-proetale}
设 $S$ 为概形。对 $K \in D(S_\proetale)$，映射
$$
K \longrightarrow R\pi_{S, *}\pi_S^{-1}K
$$
是同构。
\end{lemma}

\begin{proof}
这是因为 $\pi_S^{-1}$ 与 $\pi_{S, *} = i_S^{-1}$ 都是正合函子，并且复合
$\pi_{S, *} \circ \pi_S^{-1}$ 是恒等函子。
\end{proof}






















\section{pro-\'etale 位点的点}
\label{proetale-section-points}

\noindent
先应用 Deligne 判据证明有足够多的点。

\begin{lemma}
\label{proetale-lemma-points-proetale}
设 $S$ 为概形。pro-\'etale 位点 $\Sch_\proetale$、$S_\proetale$、
$(\Sch/S)_\proetale$、$S_{affine, \proetale}$ 与
$(\textit{Aff}/S)_\proetale$ 都有足够多的点。
\end{lemma}

\begin{proof}
$S$ 的大 pro-\'etale 拓扑斯等价于 $(\textit{Aff}/S)_\proetale$ 所定义的拓扑斯；
见引理 \ref{proetale-lemma-affine-big-site-proetale}。$S_\proetale$ 上的层拓扑斯等价于与
$S_{affine, \proetale}$ 相伴的拓扑斯；见引理 \ref{proetale-lemma-alternative}。关于位点
$(\textit{Aff}/S)_\proetale$ 与 $S_{affine, \proetale}$ 的结论立即由 Deligne 的结果
“位点”引理 \ref{sites-lemma-criterion-points} 得出。$\Sch_\proetale$ 的情形则由
它等于 $(\Sch/\Spec(\mathbf{Z}))_\proetale$ 得到。
\end{proof}

\noindent
设 $S$ 为概形，$\overline{s} : \Spec(k) \to S$ 为几何点。定义
$\overline{s}$ 的一个{\it pro-\'etale 邻域}为交换图
$$
\xymatrix{
\Spec(k) \ar[r]_-{\overline{u}} \ar[rd]_{\overline{s}} & U \ar[d] \\
& S
}
$$
，其中 $U \to S$ 弱 \'etale。

\begin{lemma}
\label{proetale-lemma-cofinal-etale}
设 $S$ 为概形，$\overline{s} : \Spec(k) \to S$ 为几何点。$\overline{s}$ 的
pro-\'etale 邻域范畴是余滤的。
\end{lemma}

\begin{proof}
证明与“\'Etale 上同调”引理 \ref{etale-cohomology-lemma-cofinal-etale} 的证明相同，
但须使用“态射续篇”引理
\ref{more-morphisms-lemma-composition-weakly-etale}、
\ref{more-morphisms-lemma-base-change-weakly-etale} 与
\ref{more-morphisms-lemma-weakly-etale-permanence} 中证明的关于弱 \'etale 态射的
相应事实。
\end{proof}

\begin{lemma}
\label{proetale-lemma-geometric-lift-to-cover}
设 $S$ 为概形，$\overline{s}$ 为 $S$ 的几何点，并设
$\mathcal{U} = \{\varphi_i : S_i \to S\}_{i\in I}$ 为 pro-\'etale 覆盖。则存在
$i \in I$ 及 $S_i$ 的几何点 $\overline{s}_i$，后者映到 $\overline{s}$。
\end{lemma}

\begin{proof}
这立即由如下事实得出：$\coprod \varphi_i$ 是满的，而弱 \'etale 态射诱导的剩余域
扩张是可分代数扩张（例如见“态射续篇”引理
\ref{more-morphisms-lemma-weakly-etale-strictly-henselian-local-rings}。
\end{proof}

\noindent
设 $S$ 为概形，$\overline{s}$ 为 $S$ 的几何点。对 $\Sh(S_\proetale)$ 中的
$\mathcal{F}$，以公式
$$
\mathcal{F}_{\overline{s}} = \colim_{(U, \overline{u})} \mathcal{F}(U)
$$
定义{\it $\mathcal{F}$ 在 $\overline{s}$ 处的茎}，其中余极限取遍 $\overline{s}$ 的
所有满足 $U \in \Ob(S_\proetale)$ 的 pro-\'etale 邻域
$(U, \overline{u})$。由上面两个引理，函子
$$
S_\proetale \textit{Sets},\quad
U \longmapsto
\{\overline{u}\text{ 是 }U\text{ 的几何点，且映到 }\overline{s}\}
$$
定义位点 $S_\proetale$ 的一个点；见“位点”定义 \ref{sites-definition-point} 与引理
\ref{sites-lemma-neighbourhoods-cofiltered}。因此函子
$\mathcal{F} \mapsto \mathcal{F}_{\overline{s}}$ 定义拓扑斯
$\Sh(S_\proetale)$ 的一个点；见“位点”定义 \ref{sites-definition-point-topos} 与引理
\ref{sites-lemma-point-site-topos}。特别地，此函子正合并与任意余极限交换。事实上，
此函子还有另一种描述。

\begin{lemma}
\label{proetale-lemma-classical-point}
在上述情形中，概形 $\Spec(\mathcal{O}_{S, \overline{s}}^{sh})$ 是
$X_\proetale$ 的对象，并且有关于 $\mathcal{F}$ 函子性的典范同构
$$
\mathcal{F}(\Spec(\mathcal{O}_{S, \overline{s}}^{sh})) =
\mathcal{F}_{\overline{s}}
$$
\end{lemma}

\begin{proof}
第一个断言由严格 Hensel 化作为 $S$ 上 \'etale 代数之滤过余极限的构造显然得出；
也可使用“态射续篇”引理
\ref{more-morphisms-lemma-weakly-etale-strictly-henselian-local-rings} 对弱 \'etale
态射的刻画。第二个断言由 Olivier 定理（“代数续篇”定理
\ref{more-algebra-theorem-olivier}）得出：概形
$\Spec(\mathcal{O}_{S, \overline{s}}^{sh})$ 是 $\overline{s}$ 的 pro-\'etale 邻域范畴
的始对象。
\end{proof}

\noindent
与 $S$ 的 \'etale 拓扑斯的情形不同，$\Sh(S_\proetale)$ 的每个点并不都具有这种
形式，与几何点相伴的点族也不是保守的。事实上，设 $S = \Spec(k)$，其中 $k$ 是
代数闭域。令 $A$ 为非零阿贝尔群。考虑 $S_\proetale$ 上如下定义的层 $\mathcal{F}$：
$$
\mathcal{F}(U) =
\frac{\{\text{映射 }U \to A\}}{\{\text{局部常值映射}\}}
$$
当 $U$ 仿射时用此式定义，一般情形则取层化；见例
\ref{proetale-example-functions-mod-locally-constant-functions}。若 $U = S = \Spec(k)$，则
$\mathcal{F}(U) = 0$，但一般而言 $\mathcal{F}$ 并非零层。事实上，
$S_\proetale$ 包含具有无限多个点的仿射对象。例如，设 $E = \lim E_n$ 是具有满
转移映射的有限集之逆极限，例如
$E = \mathbf{Z}_p = \lim \mathbf{Z}/p^n\mathbf{Z}$。概形
$U = \Spec(\colim \text{Map}(E_n, k))$ 是 $S_\proetale$ 的对象，因为
$\colim \text{Map}(E_n, k)$ 在 $k$ 上弱 \'etale（甚至 ind-Zariski）。由于存在并非
局部常值的映射 $E \to A$，故 $\mathcal{F}(U)$ 非零。因此 $\mathcal{F}$ 是非零
阿贝尔层，但它在 $S$ 的唯一几何点处的茎为零。既然已知 $S_\proetale$ 有足够多的
点，便可断定 pro-\'etale 位点必有一个点并非来自上述构造。

\medskip\noindent
使用仿射弱可缩对象，可以替代使用点的论证。首先，由引理
\ref{proetale-lemma-proetale-enough-w-contractible}，有足够多的仿射弱可缩对象。其次，若
$W \in \Ob(S_\proetale)$ 仿射且弱可缩，则函子
$$
\Sh(S_\proetale) \longrightarrow \textit{Sets},\quad
\mathcal{F} \longmapsto \mathcal{F}(W)
$$
是与所有极限交换的正合函子 $\Sh(S_\proetale) \to \textit{Sets}$。函子
$$
\textit{Ab}(S_\proetale) \longrightarrow \textit{Ab},\quad
\mathcal{F} \longmapsto \mathcal{F}(W)
$$
正合并与直和交换（因为 $W$ 拟紧；见“位点”引理
\ref{sites-lemma-directed-colimits-sections}），故与所有极限和余极限交换。此外，可在
$S_\proetale$ 的这些仿射弱可缩对象上求值，以检验阿贝尔层复形的正合性；见
“位点上同调”命题 \ref{sites-cohomology-proposition-enough-weakly-contractibles}。

\medskip\noindent
最后要指出，对仿射弱可缩的 $W$，函子 $\mathcal{F} \mapsto \mathcal{F}(W)$ 一般并非
$S_\proetale$ 某个点的茎函子，因为当 $W$ 不连通时，它不保持集合层的余积。事实上，
只要 $W$ 有多于 $1$ 个闭点，$W$ 就不连通；也就是说，只要 $W$ 不是严格 Hensel
局部环的谱（这正是上面讨论的特殊情形），它便不连通。





\section{与 \'etale 位点的比较}
\label{proetale-section-comparison}

\noindent
设 $X$ 为概形。适当选择位点后\footnote{如定义
\ref{proetale-definition-big-proetale-site} 所述，选择一个含有 $X$ 的大 pro-\'etale 位点
$\Sch_\proetale$。然后令 $\Sch_\etale$ 为与 $\Sch_\proetale$ 具有相同底范畴的
位点，但其覆盖恰为同时也是 \'etale 覆盖的那些 pro-\'etale 覆盖。作此选择后，
令 $X_\etale$ 与 $X_\proetale$ 分别为定义
\ref{proetale-definition-big-small-proetale} 以及“概形上的拓扑”定义
\ref{topologies-definition-big-small-etale} 中所定义的子范畴。另参见
“概形上的拓扑”注记 \ref{topologies-remark-choice-sites}。}，将 $U/X$ 映到 $U/X$
的函子 $u : X_\etale \to X_\proetale$ 定义了一个位点态射
$$
\epsilon : X_\proetale \longrightarrow X_\etale
$$
这由“位点”命题 \ref{sites-proposition-get-morphism} 得出。

\begin{lemma}
\label{proetale-lemma-describe-pullback-from-etale}
沿用上述记号。设 $\mathcal{F}$ 为 $X_\etale$ 上的层。规则
$$
X_\proetale \longrightarrow \textit{Sets},\quad
(f : Y \to X) \longmapsto \Gamma(Y_\etale, f_\etale^{-1}\mathcal{F})
$$
是一个层，并且等于 $\epsilon^{-1}\mathcal{F}$。这里
$f_\etale : Y_\etale \to X_\etale$ 是在“Étale 上同调”第
\ref{etale-cohomology-section-functoriality} 节中构造的小 \'etale 位点态射。
\end{lemma}

\begin{proof}
由引理 \ref{proetale-lemma-recognize-proetale-covering}，每个 pro-\'etale 覆盖都是 fpqc
覆盖。因此根据“Étale 上同调”引理
\ref{etale-cohomology-lemma-describe-pullback}，该公式定义了 $X_\proetale$
上的一个层。令 $a : \Sh(X_\etale) \to \Sh(X_\proetale)$ 为把 $\mathcal{F}$
映到引理中公式所给之层的函子。为证明 $a = \epsilon^{-1}$，只需证明 $a$
是 $\epsilon_*$ 的左伴随。

\medskip\noindent
设 $\mathcal{G}$ 为 $\Sh(X_\proetale)$ 的对象。回忆
$\epsilon_*\mathcal{G}$ 就是把 $\mathcal{G}$ 限制到全子范畴 $X_\etale$
所得的层。设 $f : Y \to X$ 为 $X_\proetale$ 的一个对象。我们把
$Y_\etale$ 看作 $X_\proetale$ 的子范畴。层 $\mathcal{G}$ 的限制映射定义了映射
$$
\epsilon_*\mathcal{G} = \mathcal{G}|_{X_\etale}
\longrightarrow
f_{\etale, *}(\mathcal{G}|_{Y_\etale})
$$
具体而言，对 $X_\etale$ 中的 $U$，
$f_{\etale, *}(\mathcal{G}|_{Y_\etale})$ 在 $U$ 上的值为
$\mathcal{G}(Y \times_X U)$，并且存在限制映射
$\mathcal{G}(U) \to \mathcal{G}(Y \times_X U)$。由伴随性，这确定了映射
$$
f_\etale^{-1}(\epsilon_*\mathcal{G}) \to \mathcal{G}|_{Y_\etale}
$$
对 $X_\proetale$ 中所有 $f : Y \to X$ 汇集这些映射，我们得到典范映射
$a(\epsilon_*\mathcal{G}) \to \mathcal{G}$。

\medskip\noindent
设 $\mathcal{F}$ 为 $\Sh(X_\etale)$ 的对象。显然有
$\mathcal{F} = \epsilon_*a(\mathcal{F})$。

\medskip\noindent 
我们断言，映射 $\mathcal{F} \to \epsilon_*a(\mathcal{F})$ 与
$a(\epsilon_*\mathcal{G}) \to \mathcal{G}$ 分别是该伴随的单位与余单位（参见
“范畴”第 \ref{categories-section-adjoint} 节）。为此，只需证明相应的映射
$$
\Mor_{\Sh(X_\proetale)}(a(\mathcal{F}), \mathcal{G}) \to
\Mor_{\Sh(X_\etale)}(\mathcal{F}, \epsilon^{-1}\mathcal{G})
$$
以及
$$
\Mor_{\Sh(X_\etale)}(\mathcal{F}, \epsilon^{-1}\mathcal{G}) \to
\Mor_{\Sh(X_\proetale)}(a(\mathcal{F}), \mathcal{G})
$$
互为逆映射。细节验证从略。
\end{proof}

\begin{lemma}
\label{proetale-lemma-fully-faithful}
设 $X$ 为概形。对 $X_\etale$ 上的每个层 $\mathcal{F}$，伴随映射
$\mathcal{F} \to \epsilon_*\epsilon^{-1}\mathcal{F}$ 是同构；换言之，对
$X_\etale$ 中的 $U$，有
$\epsilon^{-1}\mathcal{F}(U) = \mathcal{F}(U)$。
\end{lemma}

\begin{proof}
立即由引理 \ref{proetale-lemma-describe-pullback-from-etale} 中对
$\epsilon^{-1}$ 的描述得出。
\end{proof}

\begin{lemma}
\label{proetale-lemma-limit-pullback}
设 $X$ 为概形。设 $Y = \lim Y_i$ 是 $X_\proetale$ 中拟紧拟分离对象所成有向逆
系统的极限，并且转移态射为仿射态射。对 $X_\etale$ 上任意层 $\mathcal{F}$，有
$$
\epsilon^{-1}\mathcal{F}(Y) = \colim \epsilon^{-1}\mathcal{F}(Y_i)
$$
此外，如果 $Y_i$ 属于 $X_\etale$，则有
$\epsilon^{-1}\mathcal{F}(Y) = \colim \mathcal{F}(Y_i)$。
\end{lemma}

\begin{proof}
根据引理 \ref{proetale-lemma-describe-pullback-from-etale} 中对
$\epsilon^{-1}\mathcal{F}$ 的描述，所示公式是“Étale 上同调”定理
\ref{etale-cohomology-theorem-colimit} 的特殊情形。（当 $X$、$Y$ 与各
$Y_i$ 均为仿射概形时，可参见更易读的“Étale 上同调”引理
\ref{etale-cohomology-lemma-directed-colimit-cohomology}。）最后的断言立即由此及
引理 \ref{proetale-lemma-fully-faithful} 得出。
\end{proof}

\begin{lemma}
\label{proetale-lemma-affine-vanishing}
设 $X$ 为仿射概形。对 $X_\etale$ 上的内射阿贝尔层 $\mathcal{I}$，当
$p > 0$ 时有 $H^p(X_\proetale, \epsilon^{-1}\mathcal{I}) = 0$。
\end{lemma}

\begin{proof}
我们将用“位点上的上同调”引理
\ref{sites-cohomology-lemma-cech-vanish-collection} 证明这一点。令
$\mathcal{B} \subset \Ob(X_\proetale)$ 为 $X_\proetale$ 中满足
$\mathcal{O}(X) \to \mathcal{O}(U)$ 为 ind-\'etale 的仿射对象 $U$ 所成的集合。
令 $\text{Cov}$ 为形如 $\{U_i \to U\}_{i = 1, \ldots, n}$ 的 pro-\'etale
覆盖所成的集合，其中 $U \in \mathcal{B}$，并且对
$i = 1, \ldots, n$，$\mathcal{O}(U) \to \mathcal{O}(U_i)$ 为 ind-\'etale。
 “位点上的上同调”引理
\ref{sites-cohomology-lemma-cech-vanish-collection} 的性质 (1) 与 (2)，由引理
\ref{proetale-lemma-composition-ind-etale}、\ref{proetale-lemma-base-change-ind-etale}、
\ref{proetale-lemma-proetale-affine} 以及命题 \ref{proetale-proposition-weakly-etale} 可知，对
$\mathcal{B}$ 和 $\text{Cov}$ 成立。

\medskip\noindent
为检验条件 (3)，设
$\mathcal{U} = \{U_i \to U\}_{i = 1, \ldots, n}$ 是 $\text{Cov}$ 的元素。
我们必须证明 $\epsilon^{-1}\mathcal{I}$ 关于 $\mathcal{U}$ 的高阶 Čech
上同调群为零。首先，将 $U_i = \lim_{a \in A_i} U_{i, a}$ 写成有向逆极限，
其中 $U_{i, a} \to U$ 为 \'etale 态射且 $U_{i, a}$ 为仿射概形。我们将
$A_1 \times \ldots \times A_n$ 视为有向集，其序定义为：
$(a_1, \ldots, a_n) \geq (a_1', \ldots, a_n')$ 当且仅当对
$i = 1, \ldots, n$ 有 $a_i \geq a_i'$。注意，对所有
$a_1, \ldots, a_n \in A_1 \times \ldots \times A_n$，
$\mathcal{U}_{(a_1, \ldots, a_n)} = \{U_{i, a_i} \to U\}_{i = 1, \ldots, n}$
都是 \'etale 覆盖。又注意到
$$
U_{i_0} \times_U U_{i_1} \times_U \ldots \times_U U_{i_p}
=
\lim_{(a_1, \ldots, a_n) \in A_1 \times \ldots \times A_n}
U_{i_0, a_{i_0}} \times_U
U_{i_1, a_{i_1}} \times_U \ldots \times_U U_{i_p, a_{i_p}}
$$
对所有 $i_0, \ldots, i_p \in \{1, \ldots, n\}$ 都成立，因为极限与纤维积
可交换。因此由引理 \ref{proetale-lemma-limit-pullback} 以及滤过余极限的正合性，有
$$
\check{H}^p(\mathcal{U}, \epsilon^{-1}\mathcal{I}) =
\colim \check{H}^p(\mathcal{U}_{(a_1, \ldots, a_n)}, \epsilon^{-1}\mathcal{I}) 
$$
所以只需对 $U$ 的 \'etale 覆盖证明该消没性即可！

\medskip\noindent
设 $\mathcal{U} = \{U_i \to U\}_{i = 1, \ldots, n}$ 为 \'etale 覆盖，其中
$U_i$ 为仿射概形。将 $U = \lim_{b \in B} U_b$ 写成有向逆极限，其中 $U_b$
为仿射概形且 $U_b \to X$ 为 \'etale 态射。由“极限”引理
\ref{limits-lemma-descend-finite-presentation}、
\ref{limits-lemma-limit-affine} 和 \ref{limits-lemma-descend-etale}，可以选择
$b_0 \in B$，使得对 $i = 1, \ldots, n$，存在仿射概形之间的 \'etale 态射
$U_{i, b_0} \to U_{b_0}$，并且
$U_i = U \times_{U_{b_0}} U_{i, b_0}$。对 $b \geq b_0$，置
$U_{i, b} = U_b \times_{U_{b_0}} U_{i, b_0}$。当 $b$ 足够大时，族
$\mathcal{U}_b = \{U_{i, b} \to U_b\}_{i = 1, \ldots, n}$ 是 \'etale 覆盖，
参见“极限”引理 \ref{limits-lemma-descend-surjective}。与前面完全相同，可得
$$
\check{H}^p(\mathcal{U}, \epsilon^{-1}\mathcal{I}) =
\colim \check{H}^p(\mathcal{U}_b, \epsilon^{-1}\mathcal{I}) =
\colim \check{H}^p(\mathcal{U}_b, \mathcal{I})
$$
其中最后一个等号由引理 \ref{proetale-lemma-fully-faithful} 得出。右端每个 {\v C}ech
复形在正次数均无上同调（“位点上的上同调”引理
\ref{sites-cohomology-lemma-injective-trivial-cech}），故左端复形亦然。这就证明了
“位点上的上同调”引理 \ref{sites-cohomology-lemma-cech-vanish-collection}
的条件 (3)。由于 $X \in \mathcal{B}$，引理得证。
\end{proof}

\begin{lemma}
\label{proetale-lemma-relative-comparison}
设 $X$ 为概形。
\begin{enumerate}
\item 对 $X_\etale$ 上的阿贝尔层 $\mathcal{F}$，有
$R\epsilon_*(\epsilon^{-1}\mathcal{F}) = \mathcal{F}$。
\item 对 $K \in D^+(X_\etale)$，映射 $K \to R\epsilon_*\epsilon^{-1}K$
是同构。
\end{enumerate}
\end{lemma}

\begin{proof}
设 $\mathcal{I}$ 为 $X_\etale$ 上的内射阿贝尔层。回忆
$R^q\epsilon_*(\epsilon^{-1}\mathcal{I})$ 是与
$U \mapsto H^q(U_\proetale, \epsilon^{-1}\mathcal{I})$ 相关联的层，参见
“位点上的上同调”引理 \ref{sites-cohomology-lemma-higher-direct-images}。
由引理 \ref{proetale-lemma-affine-vanishing}，当 $q > 0$ 且 $U$ 为 $X$ 上的仿射
\'etale 概形时，该值为零。由于 $X_\etale$ 的每个对象都可由仿射对象覆盖，故当
$q > 0$ 时，$R^q\epsilon_*(\epsilon^{-1}\mathcal{I}) = 0$。

\medskip\noindent
设 $K \in D^+(X_\etale)$。选择一个表示 $K$ 的、由 $X_\etale$ 上内射阿贝尔层
组成的下有界复形 $\mathcal{I}^\bullet$。于是 $\epsilon^{-1}K$ 由
$\epsilon^{-1}\mathcal{I}^\bullet$ 表示。由 Leray 无环性引理（“衍生范畴”引理
\ref{derived-lemma-leray-acyclicity}），
$R\epsilon_*\epsilon^{-1}K$ 由 $\epsilon_*\epsilon^{-1}\mathcal{I}^\bullet$
表示。由引理 \ref{proetale-lemma-fully-faithful}，可得
$R\epsilon_*\epsilon^{-1}\mathcal{I}^\bullet = \mathcal{I}^\bullet$，从而 (2)
得证。(1) 是 (2) 的特殊情形。
\end{proof}

\begin{lemma}
\label{proetale-lemma-compare-cohomology}
设 $X$ 为概形。
\begin{enumerate}
\item 对 $X_\etale$ 上的阿贝尔层 $\mathcal{F}$，有
$$
H^i(X_\etale, \mathcal{F}) =
H^i(X_\proetale, \epsilon^{-1}\mathcal{F})
$$
对所有 $i$ 均成立。
\item 对 $K \in D^+(X_\etale)$，有
$$
R\Gamma(X_\etale, K) = R\Gamma(X_\proetale, \epsilon^{-1}K)
$$
\end{enumerate}
\end{lemma}

\begin{proof}
这是引理 \ref{proetale-lemma-relative-comparison} 与 Leray 谱序列（“位点上的上同调”
引理 \ref{sites-cohomology-lemma-apply-Leray}）的直接推论。
\end{proof}

\begin{lemma}
\label{proetale-lemma-compare-cohomology-nonabelian}
设 $X$ 为概形。设 $\mathcal{G}$ 为 $X_\etale$ 上的（可能非交换）群层。则有
$$
H^1(X_\etale, \mathcal{G}) =
H^1(X_\proetale, \epsilon^{-1}\mathcal{G})
$$
其中 $H^1$ 定义为挠子的同构类所成的集合（参见“位点上的上同调”第
\ref{sites-cohomology-section-h1-torsors} 节）。
\end{lemma}

\begin{proof}
由引理 \ref{proetale-lemma-fully-faithful}，函子 $\epsilon^{-1}$ 是全忠实的，故映射
$H^1(X_\etale, \mathcal{G}) \to H^1(X_\proetale, \epsilon^{-1}\mathcal{G})$
显然是单射。为证明满射性，只需证明任意 $\epsilon^{-1}\mathcal{G}$-挠子
$\mathcal{F}$ 都是 \'etale 局部平凡的。为此可以假设 $X$ 为仿射概形。因此，
我们归约为对仿射概形 $X$ 证明满射性。

\medskip\noindent
选择一个覆盖 $\{U \to X\}$，使得 (a) $U$ 为仿射概形，(b)
$\mathcal{O}(X) \to \mathcal{O}(U)$ 为 ind-\'etale，并且 (c)
$\mathcal{F}(U)$ 非空。由命题 \ref{proetale-proposition-weakly-etale}，以及 $X$ 的标准
pro-\'etale 覆盖在 $X$ 的所有 pro-\'etale 覆盖中共尾这一事实（引理
\ref{proetale-lemma-proetale-affine}），可以作出这种选择。将 $U = \lim U_i$ 写成
$X$ 上仿射 \'etale 概形的极限。选取 $s \in \mathcal{F}(U)$。令
$g \in \epsilon^{-1}\mathcal{G}(U \times_X U)$ 为满足
$g \cdot \text{pr}_1^*s = \text{pr}_2^*s$ 在 $\mathcal{F}(U \times_X U)$
中成立的唯一截面。于是 $g$ 满足余循环条件
$$
\text{pr}_{12}^*g \cdot \text{pr}_{23}^*g = \text{pr}_{13}^*g
$$
于 $\epsilon^{-1}\mathcal{G}(U \times_X U \times_X U)$ 中。由引理
\ref{proetale-lemma-limit-pullback}，有
$$
\epsilon^{-1}\mathcal{G}(U \times_X U) =
\colim \mathcal{G}(U_i \times_X U_i)
$$
以及
$$
\epsilon^{-1}\mathcal{G}(U \times_X U \times_X U) =
\colim \mathcal{G}(U_i \times_X U_i \times_X U_i)
$$
因此可以找到一个 $i$ 以及映到 $g$ 且满足余循环条件的元素
$g_i \in \mathcal{G}(U_i \times_X U_i)$。于是余循环 $g_i$ 定义了
$X_\etale$ 上的一个 $\mathcal{G}$-挠子；按构造，其拉回与 $\mathcal{F}$
同构。这里略去了一些细节（即挠子与 1-余循环之间的关系；这一内容应当加入位点
上同调一章）。
\end{proof}

\begin{lemma}
\label{proetale-lemma-compare-derived}
设 $X$ 为概形，$\Lambda$ 为环。
\begin{enumerate}
\item 全忠实函子
$\epsilon^{-1} : \textit{Mod}(X_\etale, \Lambda) \to
\textit{Mod}(X_\proetale, \Lambda)$
的本质像是弱 Serre 子范畴 $\mathcal{C}$。
\item 函子 $\epsilon^{-1}$ 定义了 $D^+(X_\etale, \Lambda)$ 与
$D^+_\mathcal{C}(X_\proetale, \Lambda)$ 之间的范畴等价，其拟逆由
$R\epsilon_*$ 给出。
\end{enumerate}
\end{lemma}

\begin{proof}
为证明 (1)，我们将验证“同调”引理
\ref{homology-lemma-characterize-weak-serre-subcategory} 的条件 (1)--(4)。由于
$\epsilon^{-1}$ 全忠实（引理 \ref{proetale-lemma-fully-faithful}）且正合，除了条件 (4)
之外，其余条件均显然成立。然而，如果
$$
0 \to \epsilon^{-1}\mathcal{F}_1 \to \mathcal{G} \to
\epsilon^{-1}\mathcal{F}_2 \to 0
$$
是 $X_\proetale$ 上 $\Lambda$-模层的短正合列，则得到
$$
0 \to \epsilon_*\epsilon^{-1}\mathcal{F}_1 \to \epsilon_*\mathcal{G} \to
\epsilon_*\epsilon^{-1}\mathcal{F}_2 \to
R^1\epsilon_*\epsilon^{-1}\mathcal{F}_1
$$
由引理 \ref{proetale-lemma-relative-comparison}，这等同于短正合列
$$
0 \to \mathcal{F}_1 \to \epsilon_*\mathcal{G} \to \mathcal{F}_2 \to 0
$$
将其拉回，可得 $\mathcal{G} = \epsilon^{-1}\epsilon_*\mathcal{G}$。这证明了 (1)。

\medskip\noindent
(2) 由 (1) 以及“位点上的上同调”引理
\ref{sites-cohomology-lemma-equivalence-bounded} 得出。
\end{proof}

\noindent
设 $\Lambda$ 为环。在“位点上的模”第
\ref{sites-modules-section-locally-constant} 节中，我们已经定义了位点上局部常值
$\Lambda$-模层的概念。如果 $M$ 是 $\Lambda$-模，则 $\underline{M}$ 作为
$\underline{\Lambda}$-模层是有限表现的，当且仅当 $M$ 是有限表现
$\Lambda$-模；参见“位点上的模”引理
\ref{sites-modules-lemma-locally-constant-finite-type}。

\begin{lemma}
\label{proetale-lemma-compare-locally-constant}
设 $X$ 为概形，$\Lambda$ 为环。函子 $\epsilon^{-1}$ 定义了范畴等价
$$
\left\{
\begin{matrix}
\text{局部常值层}\\
\text{为 }\Lambda\text{-模层，定义于 }X_\etale\\
\text{且为有限表现}
\end{matrix}
\right\}
\longleftrightarrow
\left\{
\begin{matrix}
\text{局部常值层}\\
\text{为 }\Lambda\text{-模层，定义于 }X_\proetale\\
\text{且为有限表现}
\end{matrix}
\right\}
$$
\end{lemma}

\begin{proof}
设 $\mathcal{F}$ 为 $X_\proetale$ 上有限表现的局部常值 $\Lambda$-模层。选择
pro-\'etale 覆盖 $\{U_i \to X\}$，使得 $\mathcal{F}|_{U_i}$ 为常值层，即
$\mathcal{F}|_{U_i} \cong \underline{M_i}_{U_i}$。注意，如果 $M_i$ 与 $M_j$
不同构，则 $U_i \times_X U_j$ 为空。对每个 $\Lambda$-模 $M$，令
$I_M = \{i \in I \mid M_i \cong M\}$。由于 pro-\'etale 覆盖是 fpqc 覆盖，
并由“下降”引理 \ref{descent-lemma-open-fpqc-covering}，可知
$U_M = \bigcup_{i \in I_M} \Im(U_i \to X)$ 是 $X$ 的开子集。于是
$X = \coprod U_M$ 是 $X$ 的不交开覆盖。可以用某个 $M$ 所对应的 $U_M$
替换 $X$，并假设对所有 $i$ 都有 $M_i = M$。

\medskip\noindent
考虑层 $\mathcal{I} = \mathit{Isom}(\underline{M}, \mathcal{F})$。该层是
$\mathcal{G} = \mathit{Isom}(\underline{M}, \underline{M})$ 的挠子。由
“位点上的模”引理 \ref{sites-modules-lemma-locally-constant}，有
$\mathcal{G} = \underline{G}$，其中 $G = \mathit{Isom}_\Lambda(M, M)$。由引理
\ref{proetale-lemma-compare-cohomology-nonabelian}，\'etale 拓扑与 pro-\'etale 拓扑的挠子
相同，因此 $\mathcal{I}$ 在 $X$ 上 \'etale 局部地有截面。于是
$\mathcal{F}$ 在 \'etale 局部为常值层，这正是所需结论。
\end{proof}

\begin{lemma}
\label{proetale-lemma-compare-locally-constant-derived}
设 $X$ 为概形，$\Lambda$ 为 Noether 环。令 $D_{flc}(X_\etale, \Lambda)$（相应地，
$D_{flc}(X_\proetale, \Lambda)$）为 $D(X_\etale, \Lambda)$（相应地，
$D(X_\proetale, \Lambda)$）的全子范畴，其对象是上同调层均为有限型局部常值
$\Lambda$-模层的那些复形。则
$$
\epsilon^{-1} :
D_{flc}^+(X_\etale, \Lambda)
\longrightarrow
D_{flc}^+(X_\proetale, \Lambda)
$$
是范畴等价。
\end{lemma}

\begin{proof}
由“位点上的模”引理
\ref{sites-modules-lemma-kernel-finite-locally-constant} 与“衍生范畴”
第 \ref{derived-section-triangulated-sub} 节，范畴
$D_{flc}(X_\etale, \Lambda)$ 和 $D_{flc}(X_\proetale, \Lambda)$ 分别是
$D(X_\etale, \Lambda)$ 和 $D(X_\proetale, \Lambda)$ 的严格全、饱和三角子范畴。
结合引理 \ref{proetale-lemma-compare-derived} 与
\ref{proetale-lemma-compare-locally-constant} 即得本引理的断言。
\end{proof}

\begin{lemma}
\label{proetale-lemma-compare-truncations}
设 $X$ 为概形，$\Lambda$ 为 Noether 环。设 $K$ 为
$D(X_\proetale, \Lambda)$ 的对象。置
$K_n = K \otimes_\Lambda^\mathbf{L} \underline{\Lambda/I^n}$。如果 $K_1$
满足
\begin{enumerate}
\item 属于 $\epsilon^{-1} :D(X_\etale, \Lambda/I) \to
D(X_\proetale, \Lambda/I)$ 的本质像，并且
\item 对某个 $a \in \mathbf{Z}$，具有位于 $[a,\infty)$ 内的 Tor 振幅，
\end{enumerate}
则作为 $D(X_\proetale, \Lambda/I^n)$ 的对象，$K_n$ 也满足 (1) 与 (2)。
\end{lemma}

\begin{proof}
关于 $K_n$ 的断言 (2) 由更一般的“位点上的上同调”引理
\ref{sites-cohomology-lemma-bounded} 得出。关于 $K_n$ 的断言 (1) 则由下列
特异三角以及对 $n$ 的归纳得出：
$$
K \otimes_\Lambda^\mathbf{L} \underline{I^n/I^{n + 1}} \to
K_{n + 1} \to
K_n \to
K \otimes_\Lambda^\mathbf{L} \underline{I^n/I^{n + 1}}[1]
$$
以及同构
$$
K \otimes_\Lambda^\mathbf{L} \underline{I^n/I^{n + 1}} =
K_1 \otimes_{\Lambda/I}^\mathbf{L} \underline{I^n/I^{n + 1}}
$$
再结合引理 \ref{proetale-lemma-compare-derived} 所证明的事实：$\epsilon^{-1}$ 的本质像是
$D^+(X_\proetale, \Lambda/I^n)$ 的三角子范畴。
\end{proof}

\begin{example}
\label{proetale-example-functions-mod-locally-constant-functions}
设 $X$ 为概形，$A$ 为阿贝尔群。以 $fun(-, A)$ 表示 $X_\proetale$ 上将 $U$
映到所有映射 $U \to A$（点集之间的映射）所成集合的层。考虑
$X_\proetale$ 上的层列
$$
0 \to \underline{A} \to fun(-, A) \to \mathcal{F} \to 0
$$
由于常值层是从终拓扑斯拉回的，可知
$\underline{A} = \epsilon^{-1}\underline{A}$。然而，如果 $A$ 有多于一个元素，
则 $fun(-, A)$ 与 $\mathcal{F}$ 都不是从 $X$ 的 \'etale 位点拉回的。为了在某些
情形下算出 $\mathcal{F}$ 的值，假设 $X$ 的所有点都是闭点且剩余域可分闭，并设
$U$ 为仿射概形。于是 $U$ 的所有点也都是闭点且剩余域可分闭，并且有
$$
H^1_\proetale(U, \underline{A}) =
H^1_\etale(U, \underline{A}) = 0
$$
这是由引理 \ref{proetale-lemma-compare-cohomology} 与“Étale 上同调”引理
\ref{etale-cohomology-lemma-affine-only-closed-points} 得出的。因此在此情形下有
$$
\mathcal{F}(U) = fun(U, A)/\underline{A}(U)
$$
\end{example}





\section{常值 Noether 情形下的衍生完备化}
\label{proetale-section-derived-completion-noetherian}

\noindent
我们继续“代数化与形式几何”第
\ref{algebraization-section-derived-completion} 节中开始的讨论；假定读者至少已读过
该节的一部分。

\medskip\noindent
设 $\mathcal{C}$ 为位点，$\Lambda$ 为 Noether 环，并设 $I \subset \Lambda$
为理想。回忆“位点上的模”引理
\ref{sites-modules-lemma-completion-flat} 告诉我们
$$
\underline{\Lambda}^\wedge = \lim \underline{\Lambda/I^n}
$$
是平坦 $\underline{\Lambda}$-代数，并且映射
$\underline{\Lambda} \to \underline{\Lambda}^\wedge$ 识别模 $I$ 的商。因此
“代数化与形式几何”引理
\ref{algebraization-lemma-restriction-derived-complete-equivalence} 告诉我们
$$
D_{comp}(\mathcal{C}, \Lambda) =
D_{comp}(\mathcal{C}, \underline{\Lambda}^\wedge)
$$
特别地，$D_{comp}(\mathcal{C}, \Lambda)$ 的对象 $K$ 的上同调层 $H^i(K)$
是 $\underline{\Lambda}^\wedge$-模层。为使记号方便，我们通常使用
$D_{comp}(\mathcal{C}, \Lambda)$。

\begin{lemma}
\label{proetale-lemma-naive-completion}
设 $\mathcal{C}$ 为位点，$\Lambda$ 为 Noether 环，并设 $I \subset \Lambda$
为理想。“代数化与形式几何”命题
\ref{algebraization-proposition-derived-completion} 中包含函子
$D_{comp}(\mathcal{C}, \Lambda) \to D(\mathcal{C}, \Lambda)$
的左伴随将 $K$ 映到
$$
K^\wedge = R\lim(K \otimes_\Lambda^\mathbf{L} \underline{\Lambda/I^n})
$$
特别地，$K$ 衍生完备，当且仅当
$K = R\lim(K \otimes_\Lambda^\mathbf{L} \underline{\Lambda/I^n})$。
\end{lemma}

\begin{proof}
选择 $I$ 的生成元 $f_1, \ldots, f_r$。由“代数化与形式几何”
引理 \ref{algebraization-lemma-derived-completion-koszul}，有
$$
K^\wedge = 
R\lim (K \otimes_\Lambda^\mathbf{L} \underline{K_n})
$$
其中 $K_n = K(\Lambda, f_1^n, \ldots, f_r^n)$。在“代数补充”引理
\ref{more-algebra-lemma-sequence-Koszul-complexes} 中，我们已经看到
$D(\Lambda)$ 中的 pro-系统 $\{K_n\}$ 与 $\{\Lambda/I^n\}$ 同构。因此引理得证。
\end{proof}

\begin{lemma}
\label{proetale-lemma-pushforward-Noetherian-case}
设 $\Lambda$ 为 Noether 环，$I \subset \Lambda$ 为理想，并设
$f : \Sh(\mathcal{D}) \to \Sh(\mathcal{C})$ 为拓扑斯态射。则
\begin{enumerate}
\item $Rf_*$ 将 $D_{comp}(\mathcal{D}, \Lambda)$ 映入
$D_{comp}(\mathcal{C}, \Lambda)$；
\item 映射 $Rf_* : D_{comp}(\mathcal{D}, \Lambda) \to
D_{comp}(\mathcal{C}, \Lambda)$ 有左伴随
$Lf_{comp}^* : D_{comp}(\mathcal{C}, \Lambda) \to
D_{comp}(\mathcal{D}, \Lambda)$，该左伴随由 $Lf^*$ 后接衍生完备化给出；
\item $Rf_*$ 与衍生完备化可交换；
\item 对 $D_{comp}(\mathcal{D}, \Lambda)$ 中的 $K$，有
$Rf_*K = R\lim Rf_*(K \otimes^\mathbf{L}_\Lambda \underline{\Lambda/I^n})$.
\item 对 $D_{comp}(\mathcal{C}, \Lambda)$ 中的 $M$，有
$Lf^*_{comp}M =
R\lim Lf^*(M \otimes^\mathbf{L}_\Lambda \underline{\Lambda/I^n})$.
\end{enumerate}
\end{lemma}

\begin{proof}
我们已经在“代数化与形式几何”引理
\ref{algebraization-lemma-pushforward-derived-complete-adjoint} 中看到了 (1) 与 (2)。
(3) 由“代数化与形式几何”引理
\ref{algebraization-lemma-pushforward-commutes-with-derived-completion} 得出。为证明
(4)，设 $K$ 衍生完备。则
$$
Rf_*K = Rf_*( R\lim K \otimes^\mathbf{L}_\Lambda \underline{\Lambda/I^n}) =
R\lim Rf_*(K \otimes^\mathbf{L}_\Lambda \underline{\Lambda/I^n})
$$
第一个等号由引理 \ref{proetale-lemma-naive-completion} 得出，第二个等号则因为 $Rf_*$
与 $R\lim$ 可交换（“位点上的上同调”引理
\ref{sites-cohomology-lemma-Rf-commutes-with-Rlim}）。这证明了 (4)。为证明 (5)，
由引理 \ref{proetale-lemma-naive-completion}，有
$$
Lf_{comp}^*M =
R\lim ( Lf^*M \otimes_\Lambda^\mathbf{L} \underline{\Lambda/I^n})
$$
由“位点上的上同调”引理
\ref{sites-cohomology-lemma-pullback-tensor-product}，$Lf^*$ 与衍生张量积可交换；
又因为 $Lf^*\underline{\Lambda/I^n} = \underline{\Lambda/I^n}$，故得到 (5)。
\end{proof}







\section{衍生完备化与弱可缩对象}
\label{proetale-section-derived-completion-proetale}

\noindent
我们继续第 \ref{proetale-section-derived-completion-noetherian} 节的讨论。本节将说明弱可缩
对象的存在如何简化对衍生完备模的研究。

\medskip\noindent
设 $\mathcal{C}$ 为位点，$\Lambda$ 为 Noether 环，并设 $I \subset \Lambda$
为理想。尽管关于 $D_{comp}(\mathcal{C}, \Lambda)$ 的一般理论相当完善，显式给出
衍生完备复形的例子却很困难。我们知道
\begin{enumerate}
\item $D(\mathcal{C}, \Lambda/I^n)$ 的每个对象 $M$ 限制为
$D(\mathcal{C}, \Lambda)$ 的一个衍生完备对象；并且
\item 对每个 $K \in D(\mathcal{C}, \Lambda)$，衍生完备化
$K^\wedge = R\lim (K \otimes_\Lambda^\mathbf{L} \underline{\Lambda/I^n})$
是衍生完备的。
\end{enumerate}
第一类对象平凡地完备，或许并不有趣。(2) 的问题在于，一般的衍生完备化颇为
神秘，即使在 $K = \underline{\Lambda}$ 的情形也是如此。具体而言，由同伦极限的
定义，在 $D(\mathcal{C}, \Lambda)$ 中有特异三角
$$
R\lim(\underline{\Lambda/I^n}) \to
\prod \underline{\Lambda/I^n} \to
\prod \underline{\Lambda/I^n} \to
R\lim(\underline{\Lambda/I^n})[1]
$$
其中乘积取在 $D(\mathcal{C}, \Lambda)$ 中。这些乘积通过取内射分解的乘积来计算
（“内射对象”引理 \ref{injectives-lemma-derived-products}），所以可知层
$H^p(\prod \underline{\Lambda/I^n})$ 是下列预层的层化：
$$
U \longmapsto \prod H^p(U, \Lambda/I^n).
$$
作为一个显式例子，如果 $X = \Spec(\mathbf{C}[t, t^{-1}])$、
$\mathcal{C} = X_\etale$、$\Lambda = \mathbf{Z}$、$I = (2)$ 且 $p = 1$，
那么得到下列预层的层化：
$$
U \mapsto \prod H^1(U_\etale, \mathbf{Z}/2^n\mathbf{Z})
$$
这里 $U$ 为 $X$ 上的 \'etale 概形。注意，
$H^1(X_\etale, \mathbf{Z}/m\mathbf{Z})$ 是阶为 $m$ 的循环群，其生成元
$\alpha_m$ 由方程 $t = s^m$ 所给的有限 \'etale
$\mathbf{Z}/m\mathbf{Z}$-覆盖给出（参见“Étale 上同调”第
\ref{etale-cohomology-section-computation} 节）。于是上述预层的截面
$$
\alpha = (\alpha_{2^n}) \in \prod H^1(X_\etale, \mathbf{Z}/2^n\mathbf{Z})
$$
在 $X$ 上任何非空 \'etale 概形上的限制都不为零，因而与该预层相关联的层不为零。

\medskip\noindent
然而，在 pro-\'etale 位点上不会发生这种现象。原因是我们有足够多的（拟紧）弱可缩
对象。在下面的命题中，我们汇集关于常值 Noether 情形下衍生完备化的一些结果；
这些结果适用于具有足够多弱可缩对象的位点（参见“位点”定义
\ref{sites-definition-w-contractible}）。

\begin{proposition}
\label{proetale-proposition-enough-weakly-contractibles}
设 $\mathcal{C}$ 为位点，并假设 $\mathcal{C}$ 有足够多的弱可缩对象。设
$\Lambda$ 为 Noether 环，$I \subset \Lambda$ 为理想。
\begin{enumerate}
\item 衍生完备 $\Lambda$-模层所成的范畴是
$\textit{Mod}(\mathcal{C}, \Lambda)$ 的弱 Serre 子范畴。
\item $\Lambda$-模层 $\mathcal{F}$ 满足
$\mathcal{F} = \lim \mathcal{F}/I^n\mathcal{F}$，当且仅当 $\mathcal{F}$
衍生完备且 $\bigcap I^n\mathcal{F} = 0$。
\item 层 $\underline{\Lambda}^\wedge$ 是衍生完备的。
\item 如果 $\ldots \to \mathcal{F}_3 \to \mathcal{F}_2 \to \mathcal{F}_1$
是衍生完备 $\Lambda$-模层的逆系统，则 $\lim \mathcal{F}_n$ 衍生完备。
\item 对象 $K \in D(\mathcal{C}, \Lambda)$ 衍生完备，当且仅当其每个上同调层
$H^p(K)$ 都衍生完备。
\item 对象 $K \in D_{comp}(\mathcal{C}, \Lambda)$ 上有界，当且仅当
$K \otimes_\Lambda^\mathbf{L} \underline{\Lambda/I}$ 上有界。
\item 如果 $K \otimes_\Lambda^\mathbf{L} \underline{\Lambda/I}$ 具有有限 Tor
维数，则对象 $K \in D_{comp}(\mathcal{C}, \Lambda)$ 有界。
\end{enumerate}
\end{proposition}

\begin{proof}
取子集 $\mathcal{B} \subset \Ob(\mathcal{C})$，使得每个 $U \in \mathcal{B}$
均弱可缩，并且 $\mathcal{C}$ 的每个对象都有一个由 $\mathcal{B}$ 中元素组成的
覆盖。下文将不加说明地使用“位点上的上同调”引理
\ref{sites-cohomology-lemma-w-contractible} 与命题
\ref{sites-cohomology-proposition-enough-weakly-contractibles} 的结果。

\medskip\noindent
回忆 $R\lim$ 与 $R\Gamma(U, -)$ 可交换，参见“内射对象”引理
\ref{injectives-lemma-RF-commutes-with-Rlim}。设 $f \in I$。回忆
$T(K, f)$ 是系统
$$
\ldots \xrightarrow{f} K \xrightarrow{f} K \xrightarrow{f} K
$$
在 $D(\mathcal{C}, \Lambda)$ 中的同伦极限。因此
$$
R\Gamma(U, T(K, f)) = T(R\Gamma(U, K), f).
$$
由于可以通过在 $U \in \mathcal{B}$ 上取值来检验 $D(\mathcal{C}, \Lambda)$
中对象之间的映射是否为同构，故可断定：$D(\mathcal{C}, \Lambda)$ 的对象 $K$
衍生完备，当且仅当对每个 $U \in \mathcal{B}$，对象 $R\Gamma(U, K)$ 作为
$D(\Lambda)$ 的对象衍生完备。

\medskip\noindent
上述注记表明，(1) 与 (5) 由环上模的相应结果得出，参见“代数补充”
引理 \ref{more-algebra-lemma-hom-from-Af} 与
\ref{more-algebra-lemma-serre-subcategory}。同理，(2) 可由“代数补充”命题
\ref{more-algebra-proposition-derived-complete-modules} 得出，因为对
$U \in \mathcal{B}$ 有 $(I^n\mathcal{F})(U) = I^n \cdot \mathcal{F}(U)$
（这是因为在 $U$ 上取值是正合的）。

\medskip\noindent
证明 (4)。同伦极限 $R\lim \mathcal{F}_n$ 属于 $D_{comp}(X, \Lambda)$（参见
“代数化与形式几何”定义
\ref{algebraization-definition-derived-complete} 后的讨论）。由刚刚证明的 (5)，
可得
$\lim \mathcal{F}_n = H^0(R\lim \mathcal{F}_n)$
是衍生完备的。(3) 是 (4) 的特殊情形。

\medskip\noindent
证明 (6) 与 (7)。结论由引理 \ref{proetale-lemma-naive-completion}、“位点上的上同调”
引理 \ref{sites-cohomology-lemma-bounded}，以及“位点上的上同调”命题
\ref{sites-cohomology-proposition-enough-weakly-contractibles} 中对同伦极限的计算得出。
\end{proof}






\section{一点的上同调}
\label{proetale-section-cohomology-point}

\noindent
设 $\Lambda$ 为关于理想 $I \subset \Lambda$ 完备的 Noether 环，并设 $k$ 为域。
本节中我们“计算”
$$
H^i(\Spec(k)_\proetale, \underline{\Lambda}^\wedge)
$$
其中仍如前面一样，$\underline{\Lambda}^\wedge = \lim_m
\underline{\Lambda/I^m}$。设 $k^{sep}$ 为 $k$ 的一个可分代数闭包。则
$$
\mathcal{U} = \{\Spec(k^{sep}) \to \Spec(k)\}
$$
是 $\Spec(k)$ 的 pro-\'etale 覆盖。我们将对该覆盖使用从 {\v C}ech 到上同调的
谱序列。置 $U_0 = \Spec(k^{sep})$，并令
\begin{align*}
U_n & =
\Spec(k^{sep}) \times_{\Spec(k)}
\Spec(k^{sep}) \times_{\Spec(k)} \ldots
\times_{\Spec(k)} \Spec(k^{sep}) \\
& =
\Spec(k^{sep} \otimes_k k^{sep} \otimes_k \ldots \otimes_k k^{sep})
\end{align*}
（共有 $n + 1$ 个因子）。注意，$U_0$ 的底拓扑空间 $|U_0|$ 是单点集，并且当
$n \geq 1$ 时有
$$
|U_n| = G \times \ldots \times G\quad (n\text{ 个因子})
$$
作为预有限空间，其中 $G = \text{Gal}(k^{sep}/k)$。具体而言，$U_n$ 的每个点
都有剩余域 $k^{sep}$，并把 $(\sigma_1, \ldots, \sigma_n)$ 与下列满射所对应的点
等同：
$$
k^{sep} \otimes_k k^{sep} \otimes_k \ldots \otimes_k k^{sep}
\longrightarrow k^{sep}, \quad
\lambda_0 \otimes \lambda_1 \otimes \ldots \lambda_n
\longmapsto \lambda_0 \sigma_1(\lambda_1) \ldots \sigma_n(\lambda_n)
$$
于是可计算得
\begin{align*}
R\Gamma((U_n)_\proetale, \underline{\Lambda}^\wedge)
& =
R\lim_m R\Gamma((U_n)_\proetale, \underline{\Lambda/I^m}) \\
& =
R\lim_m R\Gamma((U_n)_\etale, \underline{\Lambda/I^m}) \\
& =
\lim_m H^0(U_n, \underline{\Lambda/I^m}) \\
& = 
\text{Maps}_{cont}(G \times \ldots \times G, \Lambda)
\end{align*}
第一个等号成立，是因为 $R\Gamma$ 与衍生极限可交换，并且由命题
\ref{proetale-proposition-enough-weakly-contractibles}，$\Lambda^\wedge$ 是各层
$\underline{\Lambda/I^m}$ 的衍生极限。第二个等号由引理
\ref{proetale-lemma-compare-cohomology} 得出。第三个等号由“Étale 上同调”引理
\ref{etale-cohomology-lemma-affine-only-closed-points} 得出。第四个等号利用
“Étale 上同调”注记
\ref{etale-cohomology-remark-constant-locally-constant-maps} 识别常值层
$\underline{\Lambda/I^m}$ 的截面。接着利用 $\Lambda$ 关于 $I$ 完备，因而作为
拓扑空间等于 $\lim_m \Lambda/I^m$ 这一事实，得到
$\lim_m \text{Map}_{cont}(G, \Lambda/I^m) = \text{Map}_{cont}(G, \Lambda)$；
对 $G$ 的更高次幂也同样成立。此时“位点上的上同调”引理
\ref{sites-cohomology-lemma-cech-cohomology} 与
\ref{sites-cohomology-lemma-cech-spectral-sequence-application} 告诉我们，
$$
\Lambda \to \text{Maps}_{cont}(G, \Lambda) \to
\text{Maps}_{cont}(G \times G, \Lambda) \to \ldots
$$
该复形计算 pro-\'etale 上同调。换言之，可知
$$
H^i(\Spec(k)_\proetale, \underline{\Lambda}^\wedge)
=
H^i_{cont}(G, \Lambda)
$$
其中右端是 Tate 连续上同调，参见“Étale 上同调”第
\ref{etale-cohomology-section-continuous-group-cohomology} 节。当然，结果理应如此。

\begin{lemma}
\label{proetale-lemma-more-general-point}
设 $k$ 为域，并设 $G = \text{Gal}(k^{sep}/k)$ 为其绝对 Galois 群。此外，
\begin{enumerate}
\item 设 $M$ 为带连续 $G$-作用的预有限阿贝尔群；或者
\item 设 $\Lambda$ 为 Noether 环、$I \subset \Lambda$ 为理想，并设 $M$ 为带连续
$G$-作用的 $I$-进完备 $\Lambda$-模。
\end{enumerate}
则在 $\Spec(k)_\proetale$ 上有一个与 $M$ 相关联的典范层
$\underline{M}^\wedge$，使得
$$
H^i(\Spec(k), \underline{M}^\wedge) = H^i_{cont}(G, M)
$$
等式作为阿贝尔群或 $\Lambda$-模成立。
\end{lemma}

\begin{proof}
证明情形 (2)。置 $M_n = M/I^nM$。由于假设 $M$ 是 $I$-进完备的，故
$M = \lim M_n$。由于 $G$ 的作用连续，我们在 $M_n$ 上得到 $G$ 的连续作用。
由“Étale 上同调”定理
\ref{etale-cohomology-theorem-equivalence-sheaves-point}，该作用对应于
$\Spec(k)_\etale$ 上的一个（局部常值）$\Lambda/I^n$-模层
$\underline{M_n}$。通过比较态射 $\epsilon$ 将其拉回到
$\Spec(k)_\proetale$，并取极限
$$
\underline{M}^\wedge = \lim \epsilon^{-1}\underline{M_n}
$$
便得到引理中所述的层。与本节开头完全相同的论证给出它与 Tate 连续 Galois
上同调的比较。
\end{proof}






\section{pro-\'etale 位点的函子性}
\label{proetale-section-morphism}

\noindent
设 $f : X \to Y$ 为概形态射。函子 $Y_\proetale \to X_\proetale$，
$V \mapsto X \times_Y V$，诱导位点态射
$f_\proetale : X_\proetale \to Y_\proetale$；参见“位点”命题
\ref{sites-proposition-get-morphism}。事实上，我们得到位点态射的交换图
$$
\xymatrix{
X_\proetale \ar[r]_\epsilon \ar[d]_{f_\proetale} &
X_\etale \ar[d]^{f_\etale} \\
Y_\proetale \ar[r]^\epsilon & Y_\etale
}
$$
其中 $\epsilon$ 如第 \ref{proetale-section-comparison} 节所定义。特别地，有
$\epsilon^{-1} f_\etale^{-1} = f_\proetale^{-1} \epsilon^{-1}$。关于直像有如下
相应结果。

\begin{lemma}
\label{proetale-lemma-morphism-comparison}
设 $f : X \to Y$ 为拟紧拟分离的概形态射。
\begin{enumerate}
\item 设 $\mathcal{F}$ 为 $X_\etale$ 上的集合层。则有
$f_{\proetale, *}\epsilon^{-1}\mathcal{F} =
\epsilon^{-1}f_{\etale, *}\mathcal{F}$.
\item 设 $\mathcal{F}$ 为 $X_\etale$ 上的阿贝尔层。则有
$Rf_{\proetale, *}\epsilon^{-1}\mathcal{F} =
\epsilon^{-1}Rf_{\etale, *}\mathcal{F}$.
\end{enumerate}
\end{lemma}

\begin{proof}
证明 (1)。设 $\mathcal{F}$ 为 $X_\etale$ 上的集合层。存在典范映射
$\epsilon^{-1}f_{\etale, *}\mathcal{F} \to
f_{\proetale, *}\epsilon^{-1}\mathcal{F}$，参见“位点”第
\ref{sites-section-pullback} 节。为证明它是同构，可以在 $Y$ 上（Zariski）
局部地工作，因而可假设 $Y$ 为仿射概形。此时，$Y_\proetale$ 的每个对象都有一个
由形如 $V = \lim V_i$ 的对象组成的覆盖，其中各 $V_i$ 是 $Y$ 上的仿射
\'etale 概形（例如由命题 \ref{proetale-proposition-weakly-etale} 可知）。将映射
$\epsilon^{-1}f_{\etale, *}\mathcal{F} \to
f_{\proetale, *}\epsilon^{-1}\mathcal{F}$
在 $V$ 上取值，得到映射
$$
\colim \Gamma(X \times_Y V_i, \mathcal{F})
\longrightarrow
\Gamma(X \times_Y V, \epsilon^*\mathcal{F})
$$
左端参见引理 \ref{proetale-lemma-limit-pullback}。由引理
\ref{proetale-lemma-describe-pullback-from-etale}，有
$$
\Gamma(X \times_Y V, \epsilon^*\mathcal{F}) =
\Gamma(X \times_Y V, g_\etale^{-1}\mathcal{F})
$$
其中 $g : X \times_Y V \to X$ 是投影。因此结论由“Étale 上同调”引理
\ref{etale-cohomology-lemma-directed-colimit-cohomology} 得出。

\medskip\noindent
证明 (2)。与上面完全相同的论证表明，只需证明映射
$$
\colim H^i_\etale(X \times_Y V_i, \mathcal{F})
\longrightarrow
H^i_\etale(X \times_Y V, \mathcal{F})
$$
是同构；这再次由“Étale 上同调”引理
\ref{etale-cohomology-lemma-directed-colimit-cohomology} 得出。
\end{proof}







\section{有限态射与 pro-\'etale 位点}
\label{proetale-section-finite}

\noindent
概形的有限态射是否在 pro-\'etale 阿贝尔层上确定正合直像并不明显。

\begin{lemma}
\label{proetale-lemma-finite}
设 $f : Z \to X$ 为局部有限表现的有限概形态射。则
$f_{\proetale, *} : \textit{Ab}(Z_\proetale) \to \textit{Ab}(X_\proetale)$
是正合的。
\end{lemma}

\begin{proof}
为证明这一点，可以在 $X$ 上（Zariski）局部地工作，并假设 $X$ 为仿射概形，
记作 $X = \Spec(A)$。于是对某个有限表现的有限 $A$-代数 $B$，有
$Z = \Spec(B)$。命题 \ref{proetale-proposition-find-w-contractible} 的证明中的构造给出
忠实平坦的 ind-\'etale 环映射 $A \to D$，其中 $D$ 是 w-可缩的。可以通过在这样
的对象 $U = \Spec(D)$ 上取值来检验层列的正合性。此时，
$f_{\proetale, *}\mathcal{F}$ 在 $U$ 上的值等于 $\mathcal{F}$ 在
$V = \Spec(D \otimes_A B)$ 上的值。由引理
\ref{proetale-lemma-finite-finitely-presented-over-w-contractible}，$D \otimes_A B$ 是
w-可缩的，故在 $V$ 上取值是正合的。
\end{proof}








\section{闭浸入与 pro-\'etale 位点}
\label{proetale-section-closed-immersion}

\noindent
概形的闭浸入是否在 pro-\'etale 阿贝尔层上确定正合直像并不明显（而且很可能不
成立）。

\begin{lemma}
\label{proetale-lemma-closed-immersion-affines}
设 $i : Z \to X$ 为仿射概形之间的闭浸入态射。以 $X_{app}$ 和 $Z_{app}$
表示引理 \ref{proetale-lemma-affine-alternative} 中引入的位点。基变换函子
$$
u : X_{app} \to Z_{app},\quad U \longmapsto u(U) = U \times_X Z
$$
是连续的，并有一个全忠实左伴随 $v$。对 $Z_{app}$ 中的 $V$，态射
$V \to v(V)$ 是闭浸入，它将 $V$ 与
$u(v(V)) = v(V) \times_X Z$ 等同，并且 $v(V)$ 的每个点都特殊化到 $V$
中的一个点。函子 $v$ 余连续，并将覆盖映到覆盖。
\end{lemma}

\begin{proof}
该伴随的存在性立即由引理 \ref{proetale-lemma-lift-ind-etale} 及各定义得出。由
$X_{app}$ 中覆盖的定义，$u$ 显然是连续的。

\medskip\noindent
写成 $X = \Spec(A)$ 与 $Z = \Spec(A/I)$。设 $V = \Spec(\overline{C})$
为 $Z_{app}$ 的对象，并令 $v(V) = \Spec(C)$。我们已在引理
\ref{proetale-lemma-lift-ind-etale} 的陈述中看到，$V$ 等于
$v(V) \times_X Z = \Spec(C/IC)$。任意 $g \in C$，若其在
$C/IC = \overline{C}$ 中的像可逆，则它在 $C$ 中可逆。事实上，我们有
$A$-代数映射 $C \to C_g \to C/IC$，并由伴随性得到 $C$-代数映射
$C_g \to C$。因此 $v(V)$ 的每个点都特殊化到 $V$ 中的一个点。

\medskip\noindent
设 $\{V_i \to V\}$ 是 $Z_{app}$ 中的覆盖。于是
$\{v(V_i) \to v(V)\}$ 是 $Z_{app}$ 中的有限态射族，并且
$V \subset v(V)$ 的每个点都属于某个映射 $v(V_i) \to v(V)$ 的像。由于各态射
$v(V_i) \to v(V)$ 平坦（因为它们是弱 \'etale 的），可断定
$\{v(V_i) \to v(V)\}$ 联合满射。这证明了 $v$ 将覆盖映到覆盖。

\medskip\noindent
设 $V$ 为 $Z_{app}$ 的对象，并设 $\{U_i \to v(V)\}$ 是 $X_{app}$ 中的覆盖。
于是 $\{u(U_i) \to u(v(V)) = V\}$ 是 $Z_{app}$ 的覆盖。由伴随性，得到态射
$v(u(U_i)) \to U_i$。因此族 $\{v(u(U_i)) \to v(V)\}$ 加细给定覆盖，从而可知
$v$ 是余连续的。
\end{proof}

\begin{lemma}
\label{proetale-lemma-closed-immersion-affines-apply}
设 $Z \to X$ 为仿射概形之间的闭浸入态射。相应的拓扑斯态射
$i = i_\proetale$ 等于与引理 \ref{proetale-lemma-closed-immersion-affines} 中全忠实余连续
函子 $v : Z_{app} \to X_{app}$ 相关联的拓扑斯态射。因此
\begin{enumerate}
\item $i^{-1}\mathcal{F}$ 是与预层 $V \mapsto \mathcal{F}(v(V))$ 相关联的层；
\item 对 $Z_{app}$ 的弱可缩对象 $V$，有
$i^{-1}\mathcal{F}(V) = \mathcal{F}(v(V))$；
\item $i^{-1} : \Sh(X_\proetale) \to \Sh(Z_\proetale)$
有左伴随 $i^{Sh}_!$；
\item $i^{-1} : \textit{Ab}(X_\proetale) \to \textit{Ab}(Z_\proetale)$
有左伴随 $i_!$；
\item $\text{id} \to i^{-1}i^{Sh}_!$、$\text{id} \to i^{-1}i_!$ 与
$i^{-1}i_* \to \text{id}$ 都是同构；并且
\item $i_*$、$i^{Sh}_!$ 与 $i_!$ 都是全忠实的。
\end{enumerate}
\end{lemma}

\begin{proof}
由引理 \ref{proetale-lemma-affine-alternative}，可以用位点态射
$u : X_{app} \to Z_{app}$，$V \mapsto V \times_X Z$ 来描述 $i_\proetale$。
引理的第一个陈述由“位点”引理 \ref{sites-lemma-have-functor-other-way} 与
\ref{sites-lemma-have-functor-other-way-morphism} 得出（但要交换 $u$ 与 $v$
的角色）。

\medskip\noindent
证明 (1)。由 $i$ 作为与 $v$ 相关联之拓扑斯态射的描述，该结论按构造成立；参见
“位点”引理 \ref{sites-lemma-cocontinuous-morphism-topoi}。

\medskip\noindent
证明 (2)。由引理 \ref{proetale-lemma-closed-immersion-affines}，函子 $v$ 将覆盖映到覆盖，
所以预层 $\mathcal{G} : V \mapsto \mathcal{F}(v(V))$ 是分离预层（“位点”定义
\ref{sites-definition-separated}）。因此 $\mathcal{G}$ 的层化是
$\mathcal{G}^+$，参见“位点”定理 \ref{sites-theorem-plus}。接着，设 $V$ 为
$Z_{app}$ 的弱可缩对象。设
$\mathcal{V} = \{V_i \to V\}_{i = 1, \ldots, n}$ 是 $Z_{app}$ 中任意覆盖，并置
$\mathcal{V}' = \{\coprod V_i \to V\}$。由于 $v$ 与有限不交并可交换（因为它是
左伴随，或直接由构造可知），而 $\mathcal{F}$ 将有限不交并映到乘积，所以
$$
H^0(\mathcal{V}, \mathcal{G}) = H^0(\mathcal{V}', \mathcal{G})
$$
（记号同“位点”第 \ref{sites-section-sheafification} 节；另与“Étale 上同调”
引理 \ref{etale-cohomology-lemma-cech-complex} 比较）。因此可假设
覆盖由单个态射给出，即 $\{V' \to V\}$。由于 $V$ 弱可缩，该覆盖可由平凡覆盖
$\{V \to V\}$ 加细。因此 $\mathcal{G}^+ = i^{-1}\mathcal{F}$ 在 $V$ 上的值
就是 $\mathcal{F}(v(V))$，从而 (2) 得证。

\medskip\noindent
证明 (3)。$Z_{app}$ 的每个对象都有弱可缩对象组成的覆盖（引理
\ref{proetale-lemma-proetale-enough-w-contractible}）。由上可知，如果 $i^{Sh}_!$ 存在，
那么对弱可缩的 $V$ 应有 $i^{Sh}_!h_V = h_{v(V)}$。于是 $i^{Sh}_!$ 的存在性由
“位点”引理 \ref{sites-lemma-existence-lower-shriek} 得出。

\medskip\noindent
证明 (4)。同样地，对 $Z_{app}$ 中弱可缩的 $V$，置
$i_!\mathbf{Z}_V = \mathbf{Z}_{v(V)}$，对直和作类似处理，再应用“同调”
引理 \ref{homology-lemma-partially-defined-adjoint}，即得 $i_!$ 的存在性。细节从略。

\medskip\noindent
证明 (5)。设 $V$ 为 $Z_{app}$ 的可缩对象。则
$i^{-1}i^{Sh}_!h_V = i^{-1}h_{v(V)} = h_{u(v(V))} = h_V$。（一般而言，
$i^{-1}h_U = h_{u(U)}$。）由于可缩的 $V$ 所对应的层 $h_V$ 生成
$\Sh(Z_{app})$（“位点”引理
\ref{sites-lemma-sheaf-coequalizer-representable}），可得
$\text{id} \to i^{-1}i^{Sh}_!$ 是同构。对映射
$\text{id} \to i^{-1}i_!$ 亦同。然后
$(i^{-1}i_*\mathcal{H})(V) = i_*\mathcal{H}(v(V)) =
\mathcal{H}(u(v(V))) = \mathcal{H}(V)$，由此可知
$i^{-1}i_* \to \text{id}$ 是同构。

\medskip\noindent
现在 (6) 中的全忠实性断言由“范畴”引理
\ref{categories-lemma-adjoint-fully-faithful} 得出。
\end{proof}

\begin{lemma}
\label{proetale-lemma-closed-immersion}
设 $i : Z \to X$ 为概形的闭浸入。则
\begin{enumerate}
\item $i_\proetale^{-1}$ 与极限可交换；
\item $i_{\proetale, *}$ 是全忠实的；并且
\item $i_\proetale^{-1}i_{\proetale, *} \cong \text{id}_{\Sh(Z_\proetale)}$.
\end{enumerate}
\end{lemma}

\begin{proof}
由“位点”引理 \ref{sites-lemma-exactness-properties}，断言 (2) 与 (3) 等价。
(1) 与 (3) 在 $X$ 上是（Zariski）局部的，因此可假设 $X$ 为仿射概形。在此
情形，结论由引理 \ref{proetale-lemma-closed-immersion-affines-apply} 得出。
\end{proof}

\begin{lemma}
\label{proetale-lemma-thickening}
设 $i : Z \to X$ 为整、泛单射且满射的概形态射。则 $i_{\proetale, *}$ 与
$i_\proetale^{-1}$ 是 pro-\'etale 拓扑斯范畴之间互为拟逆的等价。
\end{lemma}

\begin{proof}
立即归约到 $X$ 为仿射概形的情形。此时 $Z$ 也为仿射概形。置
$A = \mathcal{O}(X)$ 与 $B = \mathcal{O}(Z)$。于是 $A$ 上与 $B$ 上的 \'etale
代数范畴等价，参见“Étale 上同调”定理
\ref{etale-cohomology-theorem-topological-invariance} 及注记
\ref{etale-cohomology-remark-affine-inside-equivalence}。因此 $A$ 上与 $B$ 上的
ind-\'etale 代数范畴等价。换言之，引理 \ref{proetale-lemma-affine-alternative} 中的范畴
$X_{app}$ 与 $Z_{app}$ 等价。我们略去该等价双向地将覆盖映到覆盖的验证。现在由
引理 \ref{proetale-lemma-affine-alternative} 所述 pro-\'etale 拓扑斯是 $X_{app}$ 上层的
拓扑斯，即得结论。
\end{proof}

\begin{lemma}
\label{proetale-lemma-compute-i-star}
设 $i : Z \to X$ 为概形的闭浸入。设 $U \to X$ 为 $X_\proetale$ 的对象，并且
\begin{enumerate}
\item $U$ 是仿射且弱可缩的；并且
\item $U$ 的每个点都特殊化到 $U \times_X Z$ 中的一个点。
\end{enumerate}
则对 $X_\proetale$ 上所有阿贝尔层，均有
$i_\proetale^{-1}\mathcal{F}(U \times_X Z) = \mathcal{F}(U)$。
\end{lemma}

\begin{proof}
由于拉回与限制可交换，可以用 $U$ 替换 $X$。因此可假设 $X$ 仿射且弱可缩，
并且 $X$ 的每个点都特殊化到 $Z$ 中的一个点。由引理
\ref{proetale-lemma-closed-immersion-affines-apply} 的 (1)，只需在此情形证明
$v(Z) = X$。所以我们必须证明：若 $A$ 是 w-可缩环，$I \subset A$ 是含于 $A$
之 Jacobson 根中的理想，并且 $A \to B \to A/I$ 是一个分解，其中
$A \to B$ 为 ind-\'etale，则存在唯一的收缩 $B \to A$，与到 $A/I$ 的映射相容。
注意，作为 $A/I$-代数，有 $B/IB = A/I \times R$。对 $B$ 作局部化替换后，
可以假设 $B/IB = A/I$。注意 $\Spec(B) \to \Spec(A)$ 是满射，因为其像
含有 $V(I)$，从而含有所有闭点，并且在特殊化下封闭。由于 $A$ 是 w-可缩的，存在
收缩 $B \to A$。由于 $B/IB = A/I$，该收缩与到 $A/I$ 的映射相容。唯一性的证明
从略（提示：利用 $A$ 与 $B$ 在 $A$ 的极大理想处有同构的局部环）。
\end{proof}

\begin{lemma}
\label{proetale-lemma-closed-immersion-complement-retrocompact-exact}
设 $i : Z \to X$ 为概形的闭浸入。如果 $X \setminus i(Z)$ 是 $X$ 的逆紧开子集，
则 $i_{\proetale, *}$ 是正合的。
\end{lemma}

\begin{proof}
该问题在 $X$ 上是局部的，因此可假设 $X$ 为仿射概形。记
$X = \Spec(A)$ 与 $Z = \Spec(A/I)$。存在 $f_1, \ldots, f_r \in I$，使得在
集合论意义下 $Z = V(f_1, \ldots, f_r)$，参见“代数”引理
\ref{algebra-lemma-qc-open}。由引理 \ref{proetale-lemma-thickening}，可以假设
$Z = \Spec(A/(f_1, \ldots, f_r))$。在此情形，由引理 \ref{proetale-lemma-finite}，函子
$i_{\proetale, *}$ 是正合的。
\end{proof}





\section{零延拓}
\label{proetale-section-extension-by-zero}

\noindent
“位点上的模”第 \ref{sites-modules-section-localize} 节中的一般材料使我们可以
作如下定义。

\begin{definition}
\label{proetale-definition-extension-zero}
设 $j : U \to X$ 为概形的弱 \'etale 态射。
\begin{enumerate}
\item 限制函子
$j^{-1} : \Sh(X_\proetale) \to \Sh(U_\proetale)$
有左伴随
$j_!^{Sh} : \Sh(X_\proetale) \to \Sh(U_\proetale)$.
\item 限制函子
$j^{-1} : \textit{Ab}(X_\proetale) \to \textit{Ab}(U_\proetale)$
有一个记作
$j_! : \textit{Ab}(U_\proetale) \to \textit{Ab}(X_\proetale)$
的左伴随，称为{\it 零延拓}。
\item 设 $\Lambda$ 为环。函子
$j^{-1} : \textit{Mod}(X_\proetale, \Lambda) \to
\textit{Mod}(U_\proetale, \Lambda)$
有左伴随
$j_! : \textit{Mod}(U_\proetale, \Lambda) \to
\textit{Mod}(X_\proetale, \Lambda)$
，称为{\it 零延拓}。
\end{enumerate}
\end{definition}

\noindent
照例，将其与 \'etale 情形进行比较。

\begin{lemma}
\label{proetale-lemma-jshriek-comparison}
设 $j : U \to X$ 为概形的 \'etale 态射，$\mathcal{G}$ 为 $U_\etale$ 上的阿贝尔
层。则作为 $X_\proetale$ 上的层，有
$\epsilon^{-1} j_!\mathcal{G} = j_!\epsilon^{-1}\mathcal{G}$。
\end{lemma}

\begin{proof}
这是因为两个函子都是 $j_\proetale^{-1} \epsilon_* =
\epsilon_* j_\etale^{-1}$ 的左伴随。该等式由第 \ref{proetale-section-morphism} 节中的讨论
成立。
\end{proof}

\begin{lemma}
\label{proetale-lemma-jshriek-zero}
设 $j : U \to X$ 为概形的弱 \'etale 态射，并设 $i : Z \to X$ 为满足
$U \times_X Z = \emptyset$ 的闭浸入。设 $V \to X$ 为 $X_\proetale$ 的仿射对象，
使得 $V$ 的每个点都特殊化到 $V_Z = Z \times_X V$ 中的一个点。则对
$U_\proetale$ 上所有阿贝尔层，均有 $j_!\mathcal{F}(V) = 0$。
\end{lemma}

\begin{proof}
设 $\{V_i \to V\}$ 为 pro-\'etale 覆盖。若能将该覆盖加细为一个覆盖，使其中各
成员都没有到 $U$ 的 $X$-态射，则引理得证（参见“位点上的模”第
\ref{sites-modules-section-localize} 节中 $j_!$ 的构造）。先加细该覆盖，得到
$V_i$ 均为仿射概形的有限覆盖。对每个 $i$，令 $V_i = \Spec(A_i)$，并令
$Z_i \subset V_i$ 为 $Z$ 的逆像。沿用引理 \ref{proetale-lemma-localization} 中的记号，置
$W_i = \Spec(A_{i, Z_i}^\sim)$。于是 $\coprod W_i \to V$ 是弱 \'etale 态射，
且其像包含 $V_Z$ 的所有点。由我们关于特殊化的假设，该像因而包含 $V$ 的所有
点。所以 $\{W_i \to V\}$ 是给定覆盖的一个 pro-\'etale 加细。但是 $W_i$ 中的
每个点都特殊化到位于 $Z$ 上方的一个点，故不存在 $X$ 上的态射 $W_i \to U$。
\end{proof}

\begin{lemma}
\label{proetale-lemma-open-immersion}
设 $j : U \to X$ 为概形的开浸入。则
$\text{id} \cong j^{-1}j_!$ 与 $j^{-1}j_* \cong \text{id}$，并且函子 $j_!$
和 $j_*$ 都是全忠实的。
\end{lemma}

\begin{proof}
参见“位点上的模”引理 \ref{sites-modules-lemma-restrict-back}（集合层的情形
参见“位点”引理 \ref{sites-lemma-restrict-back}），以及“范畴”引理
\ref{categories-lemma-adjoint-fully-faithful}。
\end{proof}

\noindent
下面说明零延拓与到补闭子概形的限制之间的关系。

\begin{lemma}
\label{proetale-lemma-ses-associated-to-open}
设 $X$ 为概形，$Z \subset X$ 为闭子概形，并令 $U \subset X$ 为其补集。以
$i : Z \to X$ 与 $j : U \to X$ 表示包含态射。假设 $j$ 是拟紧态射。对
$X_\proetale$ 上每个阿贝尔层，都有典范短正合列
$$
0 \to j_!j^{-1}\mathcal{F} \to \mathcal{F} \to i_*i^{-1}\mathcal{F} \to 0
$$
于 $X_\proetale$ 上，其中所有函子均针对 pro-\'etale 拓扑。
\end{lemma}

\begin{proof}
由所涉及函子的伴随性质得到这些映射。只需证明 $X_\proetale$ 有足够多的对象
（“位点”定义 \ref{sites-definition-w-contractible}），使得该序列在这些对象上
取值得到短正合列。设 $V = \Spec(A)$ 为 $X_\proetale$ 的仿射对象，并且 $A$
是 w-可缩的（这类对象足够多）。于是 $V \times_X Z$ 由理想 $I \subset A$
截出。$j$ 拟紧这一假设蕴含存在 $f_1, \ldots, f_r \in I$，使得
$V(I) = V(f_1, \ldots, f_r)$。我们得到忠实平坦的 ind-Zariski 环映射
$$
A \longrightarrow A_{f_1} \times \ldots \times A_{f_r} \times
A_{V(I)}^\sim
$$
其中 $A_{V(I)}^\sim$ 如引理 \ref{proetale-lemma-localization} 所定义。由于
$V_i = \Spec(A_{f_i}) \to X$ 经由 $U$ 分解，有
$$
j_!j^{-1}\mathcal{F}(V_i) = \mathcal{F}(V_i)
\quad\text{且}\quad
i_*i^{-1}\mathcal{F}(V_i) = 0
$$
另一方面，对概形 $V^\sim = \Spec(A_{V(I)}^\sim)$，有
$$
j_!j^{-1}\mathcal{F}(V^\sim) = 0
\quad\text{且}\quad
\mathcal{F}(V^\sim) = i_*i^{-1}\mathcal{F}(V^\sim)
$$
第一个等号由引理 \ref{proetale-lemma-jshriek-zero} 得出，第二个等号由引理
\ref{proetale-lemma-compute-i-star} 与 \ref{proetale-lemma-localization-w-contractible} 得出。因此该
序列在 $\Spec(A_{f_1} \times \ldots \times A_{f_r} \times A_{V(I)}^\sim)$
上取值得到正合列，引理得证。
\end{proof}

\begin{lemma}
\label{proetale-lemma-j-shriek-limits}
设 $j : U \to X$ 为概形的拟紧开浸入态射。函子
$j_! : \textit{Ab}(U_\proetale) \to \textit{Ab}(X_\proetale)$
与极限可交换。
\end{lemma}

\begin{proof}
由于 $j_!$ 正合，只需证明 $j_!$ 与乘积可交换。该问题在 $X$ 上是局部的，因此
可假设 $X$ 为仿射概形。设 $\mathcal{G}$ 为 $U_\proetale$ 上的阿贝尔层。我们有
$j^{-1}j_*\mathcal{G} = \mathcal{G}$。因此应用引理
\ref{proetale-lemma-ses-associated-to-open} 的正合列，得到
$$
0 \to j_!\mathcal{G} \to j_*\mathcal{G} \to i_*i^{-1}j_*\mathcal{G} \to 0
$$
其中 $i : Z \to X$ 是补集 $Z = X \setminus U$ 所带诱导既约概形结构的包含态射。
函子 $j_*$ 与 $i_*$ 作为右伴随都与乘积可交换。由引理
\ref{proetale-lemma-closed-immersion}，函子 $i^{-1}$ 与乘积可交换。因此 $j_!$ 也与乘积
可交换，因为在 pro-\'etale 位点上乘积是正合的（“位点上的上同调”命题
\ref{sites-cohomology-proposition-enough-weakly-contractibles}）。
\end{proof}




\section{pro-\'etale 位点上的可构造层}
\label{proetale-section-constructible}

\noindent
我们只讨论 Noether 环上的可构造 $\Lambda$-模层。以后还打算讨论集合、群等的可构造
层。

\begin{definition}
\label{proetale-definition-constructible}
设 $X$ 为概形，$\Lambda$ 为 Noether 环。如果对每个仿射开子集 $U \subset X$，
都存在将 $U$ 分解为可构造局部闭子概形的有限分解 $U = \coprod_i U_i$，使得对
所有 $i$，$\mathcal{F}|_{U_i}$ 都是有限型且局部常值的，则称
$X_\proetale$ 上的 $\Lambda$-模层为{\it 可构造的}。
\end{definition}

\noindent
这仍未给出任何“新的”对象。

\begin{lemma}
\label{proetale-lemma-compare-constructible}
设 $X$ 为概形，$\Lambda$ 为 Noether 环。函子 $\epsilon^{-1}$ 定义了范畴等价
$$
\left\{
\begin{matrix}
\text{可构造层}\\
\Lambda\text{-模层，定义于 }X_\etale\\
\end{matrix}
\right\}
\longleftrightarrow
\left\{
\begin{matrix}
\text{可构造层}\\
\Lambda\text{-模层，定义于 }X_\proetale\\
\end{matrix}
\right\}
$$
即 $X_\etale$ 上的可构造 $\Lambda$-模层与 $X_\proetale$ 上的可构造
$\Lambda$-模层之间的等价。
\end{lemma}

\begin{proof}
由引理 \ref{proetale-lemma-fully-faithful}，函子 $\epsilon^{-1}$ 全忠实，并与到各分层的拉回
（限制）可交换。因此，可构造 \'etale 层在 $\epsilon^{-1}$ 下的像是可构造
pro-\'etale 层。为完成证明，设 $\mathcal{F}$ 是定义
\ref{proetale-definition-constructible} 所述 $X_\proetale$ 上的可构造 $\Lambda$-模层。
存在典范映射
$$
\epsilon^{-1}\epsilon_*\mathcal{F} \longrightarrow \mathcal{F}
$$
我们将证明该映射是同构。这将证明 $\mathcal{F}$ 属于 $\epsilon^{-1}$ 的本质像，
从而完成证明（细节从略）。

\medskip\noindent
由于只需在 $X$ 上局部地证明这一点，可以假设 $X$ 拟紧拟分离，并且有一个由
可构造局部闭分层给出的有限分割 $X = \coprod_{i = 1, \ldots, n} X_i$，使得
$\mathcal{F}|_{X_i}$ 是有限型局部常值层。我们对 $n$ 归纳。基本情形
$n = 1$ 由引理 \ref{proetale-lemma-compare-locally-constant} 得出。取一点 $x \in X$；
则对某个 $k$ 有 $x \in X_k$。只需证明所示映射在 $x$ 的某个开邻域中是同构。
因此可以假设 $X_k$ 在 $X$ 中闭。置 $Z = X_k$，并以 $U \subset X$ 表示 $Z$
的补集。注意，归纳假设适用于 $\mathcal{F}$ 到 $Z$ 及到 $U$ 的限制。由引理
\ref{proetale-lemma-ses-associated-to-open}，在 $X_\proetale$ 上有短正合列
$$
0 \to j_!j^{-1}\mathcal{F} \to \mathcal{F} \to i_*i^{-1}\mathcal{F} \to 0
$$
函子性给出交换图
$$
\xymatrix{
0 \ar[r] &
\epsilon^{-1}\epsilon_*j_!j^{-1}\mathcal{F} \ar[r] \ar[d] &
\epsilon^{-1}\epsilon_*\mathcal{F} \ar[r] \ar[d] &
\epsilon^{-1}\epsilon_*i_*i^{-1}\mathcal{F} \ar[r] \ar[d] & 0 \\
0 \ar[r] &
j_!j^{-1}\mathcal{F} \ar[r] &
\mathcal{F} \ar[r] &
i_*i^{-1}\mathcal{F} \ar[r] & 0
}
$$
由归纳，一方面我们知道
$\epsilon^{-1}\epsilon_*i^{-1}\mathcal{F} \to i^{-1}\mathcal{F}$
以及
$\epsilon^{-1}\epsilon_*j^{-1}\mathcal{F} \to j^{-1}\mathcal{F}$
是同构；另一方面，知道
$i^{-1}\mathcal{F} = \epsilon^{-1}\mathcal{A}$
以及
$j^{-1}\mathcal{F} = \epsilon^{-1}\mathcal{B}$
其中有某些可构造 $\Lambda$-模层 $\mathcal{A}$ 定义在 $Z_\etale$ 上，而
$\mathcal{B}$ 定义在 $U_\etale$ 上。于是
$$
\epsilon^{-1}\epsilon_*j_!j^{-1}\mathcal{F} =
\epsilon^{-1}\epsilon_*j_!\epsilon^{-1}\mathcal{B} =
\epsilon^{-1}\epsilon_*\epsilon^{-1}j_!\mathcal{B} =
\epsilon^{-1}j_!\mathcal{B} =
j_!\epsilon^{-1}\mathcal{B} =
j_!j^{-1}\mathcal{F}
$$
其中第二个等号由引理 \ref{proetale-lemma-jshriek-comparison} 得出，第三个等号由引理
\ref{proetale-lemma-fully-faithful} 得出，第四个等号再次由引理
\ref{proetale-lemma-jshriek-comparison} 得出。类似地，有
$$
\epsilon^{-1}\epsilon_*i_*i^{-1}\mathcal{F} =
\epsilon^{-1}\epsilon_*i_*\epsilon^{-1}\mathcal{A} =
\epsilon^{-1}\epsilon_*\epsilon^{-1}i_*\mathcal{A} =
\epsilon^{-1}i_*\mathcal{A} =
i_*\epsilon^{-1}\mathcal{A} =
i_*i^{-1}\mathcal{F}
$$
这次使用引理 \ref{proetale-lemma-morphism-comparison}。由五引理可知，上面大图中央的竖直
映射是同构。
\end{proof}

\begin{lemma}
\label{proetale-lemma-constructible-serre}
设 $X$ 为概形，$\Lambda$ 为 Noether 环。$X_\proetale$ 上可构造
$\Lambda$-模层所成的范畴是 $\textit{Mod}(X_\proetale, \Lambda)$ 的弱 Serre
子范畴。
\end{lemma}

\begin{proof}
这是引理 \ref{proetale-lemma-compare-constructible}、\ref{proetale-lemma-compare-derived}，以及
\'etale 位点上的相应结果（“Étale 上同调”引理
\ref{etale-cohomology-lemma-constructible-abelian}）的形式推论。
\end{proof}

\begin{lemma}
\label{proetale-lemma-compare-constructible-derived}
设 $X$ 为概形，$\Lambda$ 为 Noether 环。令 $D_c(X_\etale, \Lambda)$（相应地，
$D_c(X_\proetale, \Lambda)$）为 $D(X_\etale, \Lambda)$（相应地，
$D(X_\proetale, \Lambda)$）的全子范畴，其对象是上同调层均为可构造
$\Lambda$-模层的那些复形。则
$$
\epsilon^{-1} :
D_c^+(X_\etale, \Lambda)
\longrightarrow
D_c^+(X_\proetale, \Lambda)
$$
是范畴等价。
\end{lemma}

\begin{proof}
由“Étale 上同调”引理
\ref{etale-cohomology-lemma-constructible-abelian}、引理
\ref{proetale-lemma-constructible-serre} 与“衍生范畴”第
\ref{derived-section-triangulated-sub} 节，范畴 $D_c(X_\etale, \Lambda)$ 与
$D_c(X_\proetale, \Lambda)$ 分别是 $D(X_\etale, \Lambda)$ 与
$D(X_\proetale, \Lambda)$ 的严格全、饱和三角子范畴。结合引理
\ref{proetale-lemma-compare-derived} 与 \ref{proetale-lemma-compare-constructible} 即得本引理的断言。
\end{proof}

\begin{lemma}
\label{proetale-lemma-tensor-c}
设 $X$ 为概形，$\Lambda$ 为 Noether 环。设
$K, L \in D_c^-(X_\proetale, \Lambda)$。则
$K \otimes_\Lambda^\mathbf{L} L$ 属于 $D_c^-(X_\proetale, \Lambda)$。
\end{lemma}

\begin{proof}
注意，$H^i(K \otimes_\Lambda^\mathbf{L} L)$ 与
$H^i(\tau_{\geq i - 1}K \otimes_\Lambda^\mathbf{L} \tau_{\geq i - 1}L)$
相同。因此可以假设 $K$ 与 $L$ 有界。在此情形，可应用引理
\ref{proetale-lemma-compare-constructible-derived} 归约到 \'etale 位点的情形；参见
“Étale 上同调”引理 \ref{etale-cohomology-lemma-tensor-c}。
\end{proof}

\begin{lemma}
\label{proetale-lemma-compare-truncations-constructible}
设 $X$ 为概形，$\Lambda$ 为 Noether 环，并设 $I \subset \Lambda$ 为理想。设
$K$ 为 $D(X_\proetale, \Lambda)$ 的对象。置
$K_n = K \otimes_\Lambda^\mathbf{L} \underline{\Lambda/I^n}$。如果 $K_1$
属于 $D^-_c(X_\proetale, \Lambda/I)$，那么对所有 $n$，$K_n$ 都属于
$D^-_c(X_\proetale, \Lambda/I^n)$。
\end{lemma}

\begin{proof}
考虑特异三角
$$
K \otimes_\Lambda^\mathbf{L} \underline{I^n/I^{n + 1}} \to
K_{n + 1} \to
K_n \to
K \otimes_\Lambda^\mathbf{L} \underline{I^n/I^{n + 1}}[1]
$$
以及同构
$$
K \otimes_\Lambda^\mathbf{L} \underline{I^n/I^{n + 1}} =
K_1 \otimes_{\Lambda/I}^\mathbf{L} \underline{I^n/I^{n + 1}}
$$
由引理 \ref{proetale-lemma-tensor-c}，该张量积具有可构造上同调层（并且当 $K_1$ 的上同调
消没时也消没）。因此，利用引理 \ref{proetale-lemma-constructible-serre} 对 $n$ 归纳，
可知每个 $K_n$ 都具有可构造上同调层。
\end{proof}








\section{可构造进层}
\label{proetale-section-adic}

\noindent
本节定义可构造 $\Lambda$-层的概念以及若干变体。

\begin{definition}
\label{proetale-definition-adic}
设 $\Lambda$ 为 Noether 环，$I \subset \Lambda$ 为理想。设 $X$ 为概形，并设
$\mathcal{F}$ 为 $X_\proetale$ 上的 $\Lambda$-模层。
\begin{enumerate}
\item 如果 $\mathcal{F} = \lim \mathcal{F}/I^n\mathcal{F}$，并且每个
$\mathcal{F}/I^n\mathcal{F}$ 都是可构造 $\Lambda/I^n$-模层，则称
$\mathcal{F}$ 为{\it 可构造 $\Lambda$-层}。
\item 如果 $\mathcal{F}$ 是可构造 $\Lambda$-层，并且每个
$\mathcal{F}/I^n\mathcal{F}$ 都局部常值，则称 $\mathcal{F}$ 为 {\it lisse}。
\item 如果存在 $I$-进完备 $\Lambda$-模 $M$，其中 $M/IM$ 有限，并且
$\mathcal{F}$ 局部同构于
$$
\underline{M}^\wedge = \lim \underline{M/I^nM}.
$$
则称 $\mathcal{F}$ 为{\it 进 lisse}\footnote{这可能不是标准记号。}。
\item 如果对每个仿射开子集 $U \subset X$，都存在将 $U = \coprod U_i$ 分解为
可构造局部闭子概形的分解，使得 $\mathcal{F}|_{U_i}$ 是进 lisse 的，则称
$\mathcal{F}$ 为{\it 进可构造的}\footnote{这可能不是标准记号。}。
\end{enumerate}
\end{definition}

\noindent
当 $\Lambda = \mathbf{Z}_\ell$ 且 $I = (\ell)$ 时，可构造 $\Lambda$-层的定义
等价于 \cite[Expos\'e VI, Definition 1.1.1]{SGA5} 中的定义。显然有蕴含关系
$$
\xymatrix{
\text{进 lisse} \ar@{=>}[r] \ar@{=>}[d] &
\text{进可构造} \ar@{=>}[d] \\
\text{lisse 可构造 }\Lambda\text{-层} \ar@{=>}[r] &
\text{可构造 }\Lambda\text{-层}
}
$$
在某些情形，竖直箭头可以反向（参见引理 \ref{proetale-lemma-Noetherian-constructible} 与
\ref{proetale-lemma-Noetherian-adic-constructible}）。一般而言，进可构造层范畴和可构造
$\Lambda$-层范畴在核与余核下都不封闭。

\medskip\noindent
具体而言，设 $X$ 为仿射概形，其底拓扑空间 $|X|$ 与
$\Lambda = \mathbf{Z}_\ell$ 同胚，参见例 \ref{proetale-example-construct-space}。以
$f : |X| \to \mathbf{Z}_\ell = \Lambda$ 表示一个同胚。可以把 $f$ 看作
$X$ 上 $\underline{\Lambda}^\wedge$ 的一个截面；乘以 $f$ 便定义了二项复形
$$
\underline{\Lambda}^\wedge \xrightarrow{f} \underline{\Lambda}^\wedge
$$
于 $X_\proetale$ 上。层 $\underline{\Lambda}^\wedge$ 是进 lisse 的。然而，上述
映射的余核不是进可构造的，因为该余核之茎的同构型取无穷多个值：
$\mathbf{Z}/\ell^n\mathbf{Z}$ 与 $\mathbf{Z}_\ell$。该余核是可构造
$\mathbf{Z}_\ell$-层。然而，核甚至不是可构造 $\mathbf{Z}_\ell$-层，因为它在
一个非拟紧开集上为零，但自身不为零。

\begin{lemma}
\label{proetale-lemma-Noetherian-constructible}
设 $X$ 为 Noether 概形，$\Lambda$ 为 Noether 环，$I \subset \Lambda$ 为理想，
并设 $\mathcal{F}$ 为 $X_\proetale$ 上的可构造 $\Lambda$-层。则存在由局部闭子概形
给出的有限分割 $X = \coprod X_i$，使得限制 $\mathcal{F}|_{X_i}$ 为 lisse。
\end{lemma}

\begin{proof}
令 $R = \bigoplus I^n/I^{n + 1}$。注意 $R$ 是 Noether 环。由于每个层
$\mathcal{F}/I^n\mathcal{F}$ 都是可构造 $\Lambda/I^n\Lambda$-模层，
$I^n\mathcal{F}/I^{n + 1}\mathcal{F}$ 也为可构造 $\Lambda/I$-模层，因而由引理
\ref{proetale-lemma-compare-constructible}，它是 $X_\etale$ 上某个可构造层
$\mathcal{G}_n$ 的拉回。置 $\mathcal{G} = \bigoplus \mathcal{G}_n$。这是
$X_\etale$ 上的 $R$-模层，并且映射
$$
\mathcal{G}_0 \otimes_{\Lambda/I} \underline{R}
\longrightarrow
\mathcal{G}
$$
是满射，因为各映射
$$
\mathcal{F}/I\mathcal{F} \otimes \underline{I^n/I^{n + 1}} \to
I^n\mathcal{F}/I^{n + 1}\mathcal{F}
$$
都是满射。因此，由“Étale 上同调”命题
\ref{etale-cohomology-proposition-constructible-over-noetherian}，
$\mathcal{G}$ 是可构造 $R$-模层。选择分割 $X = \coprod X_i$，使得
$\mathcal{G}|_{X_i}$ 是有限型局部常值 $R$-模层（“Étale 上同调”引理
\ref{etale-cohomology-lemma-constructible-quasi-compact-quasi-separated}）。我们断言
这就是引理中所需的分割。事实上，用 $X_i$ 替换 $X$ 后，可以假设
$\mathcal{G}$ 局部常值。于是每个层
 $I^n\mathcal{F}/I^{n + 1}\mathcal{F}$ 都局部常值。使用短正合列
$$
0 \to I^n\mathcal{F}/I^{n + 1}\mathcal{F} \to
\mathcal{F}/I^{n + 1}\mathcal{F} \to \mathcal{F}/I^n\mathcal{F} \to 0
$$
以及归纳和“位点上的模”引理
\ref{sites-modules-lemma-kernel-finite-locally-constant}，引理得证。
\end{proof}

\begin{lemma}
\label{proetale-lemma-weird}
设 $X$ 为弱可缩仿射概形，$\Lambda$ 为 Noether 环，$I \subset \Lambda$ 为理想。
设 $\mathcal{F}$ 为 $X_\proetale$ 上满足下列条件的 $\Lambda$-模层：
\begin{enumerate}
\item $\mathcal{F} = \lim \mathcal{F}/I^n\mathcal{F}$,
\item $\mathcal{F}/I^n\mathcal{F}$ 是常值 $\Lambda/I^n$-模层；
\item $\mathcal{F}/I\mathcal{F}$ 是有限型的。
\end{enumerate}
则 $\mathcal{F} \cong \underline{M}^\wedge$，其中 $M$ 是有限
$\Lambda^\wedge$-模。
\end{lemma}

\begin{proof}
选择 $\Lambda/I^n$-模 $M_n$，使得
$\mathcal{F}/I^n\mathcal{F} \cong \underline{M_n}$。由于有满射
$\mathcal{F}/I^{n + 1}\mathcal{F} \to \mathcal{F}/I^n\mathcal{F}$，可知存在满射
$M_{n + 1} \to M_n$，它们诱导同构
$M_{n + 1}/I^nM_{n + 1} \to M_n$。固定这类满射的一种选择，并置
$M = \lim M_n$。于是 $M$ 是 $I$-进完备 $\Lambda$-模，且
$M/I^nM = M_n$；参见“代数”引理 \ref{algebra-lemma-limit-complete}。由于
$M_1$ 是有限型 $\Lambda$-模（“位点上的模”引理
\ref{sites-modules-lemma-locally-constant-finite-type}），可知 $M$ 是有限
$\Lambda^\wedge$-模。考虑 $X_\proetale$ 上的层
$$
\mathcal{I}_n = \mathit{Isom}(\underline{M_n}, \mathcal{F}/I^n\mathcal{F})
$$
模去 $I^n$ 定义了转移映射
$$
\mathcal{I}_{n + 1} \longrightarrow \mathcal{I}_n
$$
由我们对 $M_n$ 的选择，层 $\mathcal{I}_n$ 是下列群层的挠子：
$$
\mathit{Isom}(\underline{M_n}, \underline{M_n}) =
\underline{\text{Isom}_\Lambda(M_n, M_n)}
$$
（“位点上的模”引理 \ref{sites-modules-lemma-locally-constant}），因为
$\mathcal{F}/I^n\mathcal{F}$ 在（\'etale）局部同构于 $\underline{M_n}$。由
“代数补充”引理 \ref{more-algebra-lemma-hom-systems-ML}，层系统
$(\mathcal{I}_n)$ 满足 Mittag-Leffler 条件。对每个 $n$，令
$\mathcal{I}'_n \subset \mathcal{I}_n$ 为当 $N \gg n$ 时
$\mathcal{I}_N \to \mathcal{I}_n$ 的像。则
$$
\ldots \to \mathcal{I}'_3 \to \mathcal{I}'_2 \to \mathcal{I}'_1 \to *
$$
是 $X_\proetale$ 上集合层的序列，其转移映射均为满射。由于 $*(X)$ 是单点集
（非空），并且在 $X$ 上取值将集合层的满射映到集合的满射，可以选取
$s \in \lim \mathcal{I}'_n(X)$。这些截面定义同构
$\underline{M}^\wedge \to \lim \mathcal{F}/I^n\mathcal{F} = \mathcal{F}$，证明
完成。
\end{proof}

\begin{lemma}
\label{proetale-lemma-connected-lisse}
设 $X$ 为连通概形，$\Lambda$ 为 Noether 环，$I \subset \Lambda$ 为理想。如果
$\mathcal{F}$ 是 $X_\proetale$ 上 lisse 可构造 $\Lambda$-层，则 $\mathcal{F}$
是进 lisse 的。
\end{lemma}

\begin{proof}
由引理 \ref{proetale-lemma-compare-locally-constant}，对某个局部常值
$\Lambda/I^n$-模层 $\mathcal{G}_n$，有
$\mathcal{F}/I^n\mathcal{F} = \epsilon^{-1}\mathcal{G}_n$。由“Étale 上同调”
引理 \ref{etale-cohomology-lemma-connected-locally-constant}，存在有限
$\Lambda/I^n$-模 $M_n$，使得 $\mathcal{G}_n$ 局部同构于 $\underline{M_n}$。
选择覆盖 $\{W_t \to X\}_{t \in T}$，使得每个 $W_t$ 都仿射且弱可缩。于是
$\mathcal{F}|_{W_t}$ 满足引理 \ref{proetale-lemma-weird} 的假设，因此对某个有限
$\Lambda^\wedge$-模 $N_t$，有
$\mathcal{F}|_{W_t} \cong \underline{N_t}^\wedge$。注意，对所有 $t$ 与 $n$，
$N_t/I^nN_t \cong M_n$。因此对所有 $t, t' \in T$，有
$N_t \cong N_{t'}$；参见“代数补充”引理
\ref{more-algebra-lemma-isomorphic-completions}。这证明了 $\mathcal{F}$ 是进 lisse
的。
\end{proof}

\begin{lemma}
\label{proetale-lemma-Noetherian-adic-constructible}
设 $X$ 为 Noether 概形，$\Lambda$ 为 Noether 环，$I \subset \Lambda$ 为理想，
并设 $\mathcal{F}$ 为 $X_\proetale$ 上的可构造 $\Lambda$-层。则
$\mathcal{F}$ 是进可构造的。
\end{lemma}

\begin{proof}
这是引理 \ref{proetale-lemma-Noetherian-constructible}、\ref{proetale-lemma-connected-lisse}、
Noether 概形局部连通这一事实（“拓扑”引理
\ref{topology-lemma-locally-Noetherian-locally-connected}），以及各定义的推论。
\end{proof}

\noindent
在衍生完备 $\Lambda$-模层的范畴中识别可构造 $\Lambda$-层将很有用。事实证明，
“代数补充”引理 \ref{more-algebra-lemma-derived-complete-finite} 的朴素类比
在此情形不成立。不过，下面给出“代数补充”引理
\ref{more-algebra-lemma-derived-complete-zero} 的类比。

\begin{lemma}
\label{proetale-lemma-derived-complete-zero}
设 $X$ 为概形，$\Lambda$ 为环，并设 $I \subset \Lambda$ 为有限生成理想。设
$\mathcal{F}$ 为 $X_\proetale$ 上的 $\Lambda$-模层。如果 $\mathcal{F}$ 衍生完备
且 $\mathcal{F}/I\mathcal{F} = 0$，则 $\mathcal{F} = 0$。
\end{lemma}

\begin{proof}
假设 $\mathcal{F}/I\mathcal{F}$ 为零。令 $I = (f_1, \ldots, f_r)$。设 $i < r$
为使得 $\mathcal{G} = \mathcal{F}/(f_1, \ldots, f_i)\mathcal{F}$ 非零的最大整数。
如果这样的 $i$ 不存在，则 $\mathcal{F} = 0$，这正是所要证明的。于是
$\mathcal{G}$ 作为衍生完备模之间某个映射的余核是衍生完备的，参见命题
\ref{proetale-proposition-enough-weakly-contractibles}。由 $i$ 的选择，
$f_{i + 1} : \mathcal{G} \to \mathcal{G}$ 是满射。因此
$$
\lim (\ldots \to \mathcal{G} \xrightarrow{f_{i + 1}} \mathcal{G}
\xrightarrow{f_{i + 1}} \mathcal{G})
$$
非零，这与 $\mathcal{G}$ 的衍生完备性矛盾。
\end{proof}

\begin{lemma}
\label{proetale-lemma-derived-complete-limit}
设 $X$ 为弱可缩仿射概形，$\Lambda$ 为 Noether 环，$I \subset \Lambda$ 为理想。
设 $\mathcal{F}$ 为 $X_\proetale$ 上的衍生完备 $\Lambda$-模层，并且
$\mathcal{F}/I\mathcal{F}$ 是有限型局部常值 $\Lambda/I$-模层。则存在整数 $t$
以及满射
$$
(\underline{\Lambda}^\wedge)^{\oplus t} \to \mathcal{F}
$$
\end{lemma}

\begin{proof}
由于 $X$ 弱可缩，存在有限不交开覆盖 $X = \coprod U_i$，使得
$\mathcal{F}/I\mathcal{F}|_{U_i}$ 同构于与有限 $\Lambda/I$-模 $M_i$ 相关联的
常值层。选取 $M_i$ 的有限多个生成元 $m_{ij}$。可以找到截面
$s_{ij} \in \mathcal{F}(X)$，其限制是将 $m_{ij}$ 看作
$\mathcal{F}/I\mathcal{F}$ 在 $U_i$ 上的截面所得者。令 $t$ 为各 $s_{ij}$ 的
总数。于是得到映射
$$
\alpha : \underline{\Lambda}^{\oplus t} \longrightarrow \mathcal{F}
$$
按构造，该映射模去 $I$ 后为满射。由引理 \ref{proetale-lemma-naive-completion}，
$\underline{\Lambda}^{\oplus t}$ 的衍生完备化是层
$(\underline{\Lambda}^\wedge)^{\oplus t}$。由于 $\mathcal{F}$ 衍生完备，可知
$\alpha$ 经由映射
$$
\alpha^\wedge :
(\underline{\Lambda}^\wedge)^{\oplus t}
\longrightarrow
\mathcal{F}
$$
分解。由命题 \ref{proetale-proposition-enough-weakly-contractibles}，
$\mathcal{Q} = \Coker(\alpha^\wedge)$ 是衍生完备 $\Lambda$-模层。按构造，
$\mathcal{Q}/I\mathcal{Q} = 0$。由引理 \ref{proetale-lemma-derived-complete-zero}，有
$\mathcal{Q} = 0$，这正是所要证明的。
\end{proof}








\section{合适的衍生范畴}
\label{proetale-section-suitable-derived}

\noindent
设 $X$ 为概形。我们将看到，对于许多概形 $X$，可以通过考察
$D(X_\proetale, \Lambda)$ 中满足如下性质的衍生完备对象 $K$，得到一个合适的
$\ell$-进层衍生范畴：$K \otimes_\Lambda^\mathbf{L} \mathbf{F}_\ell$
有界且其上同调层可构造。一般定义如下。

\begin{definition}
\label{proetale-definition-Dbc}
设 $\Lambda$ 为 Noether 环，$I \subset \Lambda$ 为理想。设 $X$ 为概形。
若 $D(X_\proetale, \Lambda)$ 的对象 $K$ 满足
\begin{enumerate}
\item $K$ 关于 $I$ 衍生完备；
\item $K \otimes_\Lambda^\mathbf{L} \underline{\Lambda/I}$
的上同调层可构造，并且局部具有有限 Tor 维数，
\end{enumerate}
则称该对象{\it 可构造}。以 $D_{cons}(X, \Lambda)$ 表示
$D(X_\proetale, \Lambda)$ 中由可构造 $K$ 组成的全子范畴。
\end{definition}

\noindent
回忆在我们的约定下，具有有限 Tor 维数的复形是有界的
（位点上的上同调，定义 \ref{sites-cohomology-definition-tor-amplitude}）。
事实上，可将目前证明的一切汇总为下面的引理。

\begin{lemma}
\label{proetale-lemma-describe-constructible-complexes}
在上述情形中，设 $K$ 属于 $D_{cons}(X, \Lambda)$，且 $X$ 拟紧。令
$K_n = K \otimes_\Lambda^\mathbf{L} \underline{\Lambda/I^n}$.
则存在 $a, b$，使得
\begin{enumerate}
\item $K = R\lim K_n$，且 $H^i(K) = 0$ 对所有 $i \not \in [a, b]$ 成立；
\item 每个 $K_n$ 的 Tor 振幅均包含于 $[a, b]$；
\item 每个 $K_n$ 均具有可构造上同调层；
\item 对某个
$L_n \in D_{ctf}(X_\etale, \Lambda/I^n)$
（Étale 上同调，定义 \ref{etale-cohomology-definition-ctf}），有
$K_n = \epsilon^{-1}L_n$。
\end{enumerate}
\end{lemma}

\begin{proof}
由于 $K$ 衍生完备，由引理 \ref{proetale-lemma-naive-completion} 有
$K = R\lim K_n$。按照局部具有有限 Tor 维数的定义，可以找到
$a, b$，使得 $K_1$ 的 Tor 振幅包含于 $[a, b]$。第 (2) 点由
位点上的上同调，引理 \ref{sites-cohomology-lemma-bounded} 得出。
接着，由于 $K$ 衍生完备，第 (1) 点可由位点上的上同调，命题
\ref{sites-cohomology-proposition-enough-weakly-contractibles}
对极限的描述以及以下事实得出：
$H^b(K_{n + 1}) \to H^b(K_n)$ 是满射，因为
$K_n = K_{n + 1} \otimes^\mathbf{L}_\Lambda \underline{\Lambda/I^n}$。
第 (3) 点由引理 \ref{proetale-lemma-compare-truncations-constructible} 得出。
第 (4) 点由引理 \ref{proetale-lemma-compare-constructible-derived} 以及
$L_n$ 具有有限 Tor 维数这一事实得出，因为 $K_n$ 具有有限 Tor 维数
（略去一个简短论证）。
\end{proof}

\begin{lemma}
\label{proetale-lemma-local-structure-constructible-complex}
设 $X$ 为弱可缩仿射概形。设 $\Lambda$ 为 Noether 环，
$I \subset \Lambda$ 为理想。设 $K$ 为 $D_{cons}(X, \Lambda)$ 的对象，
并且
$K \otimes_\Lambda^\mathbf{L} \underline{\Lambda/I}$ 的上同调层局部常值。
则存在有限不交开覆盖 $X = \coprod U_i$，并且对每个
$i$，存在一族有限的有限投射 $\Lambda^\wedge$-模
$M^a, \ldots, M^b$，使得在某些层的 $\Lambda$-模映射
$(\underline{M^i})^\wedge \to (\underline{M^{i + 1}})^\wedge$ 下，
$K|_{U_i}$ 在 $D(U_{i, \proetale}, \Lambda)$ 中由复形
$$
(\underline{M^a})^\wedge \to \ldots \to (\underline{M^b})^\wedge
$$
表示。
\end{lemma}

\begin{proof}
我们自由使用引理 \ref{proetale-lemma-describe-constructible-complexes} 的结论。
取该引理中的 $a, b$。我们对 $b - a$ 作归纳来证明本引理。
令 $\mathcal{F} = H^b(K)$。注意，由命题
\ref{proetale-proposition-enough-weakly-contractibles}，$\mathcal{F}$ 是衍生完备
$\Lambda$-模层。此外，$\mathcal{F}/I\mathcal{F}$ 是有限型
$\Lambda/I$-模的局部常值层。应用引理
\ref{proetale-lemma-derived-complete-limit}，得到满射
$\rho : (\underline{\Lambda}^\wedge)^{\oplus t} \to \mathcal{F}$.

\medskip\noindent
若 $a = b$，则 $K = \mathcal{F}[-b]$。在此情形下，
$$
\mathcal{F} \otimes_\Lambda^\mathbf{L} \underline{\Lambda/I} =
\mathcal{F}/I\mathcal{F}
$$
由于 $X$ 弱可缩且 $\mathcal{F}/I\mathcal{F}$ 局部常值，可以找到由仿射开集
$U_i$ 构成的有限不交并分解 $X = \coprod U_i$，以及
$\Lambda/I$-模 $\overline{M}_i$，使得
$\mathcal{F}/I\mathcal{F}$ 在 $U_i$ 上限制为
$\underline{\overline{M}_i}$。细分此覆盖后，可以假设映射
$$
\rho|_{U_i} \bmod I :
\underline{\Lambda/I}^{\oplus t}
\longrightarrow
\underline{\overline{M}_i}
$$
等于某个满射模同态
$\alpha_i : \Lambda/I^{\oplus t} \to \overline{M}_i$ 所诱导的
$\underline{\alpha_i}$；见位点上的模，引理
\ref{sites-modules-lemma-morphism-locally-constant}。注意，每个
$\overline{M}_i$ 都是有限 $\Lambda/I$-模。由于
$\mathcal{F}/I\mathcal{F}$ 的 Tor 振幅包含于 $[0, 0]$，可知
$\overline{M}_i$ 是平坦 $\Lambda/I$-模。因此 $\overline{M}_i$ 有限投射
（代数，引理 \ref{algebra-lemma-finite-projective}）。于是可以找到投影算子
$\overline{p}_i : (\Lambda/I)^{\oplus t} \to (\Lambda/I)^{\oplus t}$
使其像在映射 $\alpha_i$ 下同构地映到 $\overline{M}_i$。可以将
$\overline{p}_i$ 提升为投影算子
$p_i : (\Lambda^\wedge)^{\oplus t} \to
(\Lambda^\wedge)^{\oplus t}$\footnote{证明：由代数，引理
\ref{algebra-lemma-lift-idempotents-noncommutative}，可以把
$\overline{p}_i$ 提升为相容的投影算子系
$p_{i, n} : (\Lambda/I^n)^{\oplus t} \to (\Lambda/I^n)^{\oplus t}$；
再令 $p_i = \lim p_{i, n}$ 即可，因为
$\Lambda^\wedge = \lim \Lambda/I^n$。}。在 $M_i/IM_i = \overline{M}_i$
下，$M_i = \Im(p_i)$ 是有限且 $I$-进完备的 $\Lambda^\wedge$-模。
在 $U_i$ 上考察映射
$$
\underline{M_i}^\wedge \to
(\underline{\Lambda}^\wedge)^{\oplus t} \to
\mathcal{F}|_{U_i}
$$
按构造，其复合模 $I$ 诱导同构。源与靶均衍生完备，因而余核
$\mathcal{Q}$ 和核 $\mathcal{K}$ 也衍生完备。按构造有
$\mathcal{Q}/I\mathcal{Q} = 0$，故由引理
\ref{proetale-lemma-derived-complete-zero}，$\mathcal{Q}$ 为零。于是
$$
0 \to \mathcal{K}/I\mathcal{K} \to
\underline{\overline{M}_i}
\to \mathcal{F}/I\mathcal{F} \to 0
$$
因 $\text{Tor}_1$ 消失而正合，见本段开头；还要使用位点上的模，引理
\ref{sites-modules-lemma-completion-flat} 中的
$\underline{\Lambda}^\wedge/I\overline{\Lambda}^\wedge$，从而看出
$\underline{M_i}^\wedge/I\underline{M_i}^\wedge = \underline{\overline{M}_i}$。
所以按构造 $\mathcal{K}/I\mathcal{K} = 0$，与前同理可得
$\mathcal{K} = 0$。这证明了 $a = b$ 时的结论。

\medskip\noindent
若 $b > a$，则把映射 $\rho$ 提升为映射
$$
\tilde \rho : (\underline{\Lambda}^\wedge)^{\oplus t}[-b] \longrightarrow K
$$
于 $D(X_\proetale, \Lambda)$ 中。由第
\ref{proetale-section-derived-completion-noetherian} 节开头的讨论，可以把
$K$ 看成 $\underline{\Lambda}^\wedge$-模复形；又因为
$X_\proetale$ 弱可缩，所以将元素
$\rho(e_s) \in H^0(X, \mathcal{F})$ 提升为 $H^b(X, K)$ 的元素不存在障碍。
把 $\tilde \rho$ 嵌入到区分三角
$$
(\underline{\Lambda}^\wedge)^{\oplus t}[-b] \to K \to L \to
(\underline{\Lambda}^\wedge)^{\oplus t}[-b + 1]
$$
中，可知 $L$ 是 $D_{cons}(X, \Lambda)$ 的对象，并且
$L \otimes_\Lambda^\mathbf{L} \underline{\Lambda/I}$ 的 Tor 振幅包含于
$[a, b - 1]$（略去细节）。由归纳假设，可以如引理所述在局部描述 $L$；
设 $L$ 同构于
$$
(\underline{M^a})^\wedge \to \ldots \to (\underline{M^{b - 1}})^\wedge
$$
映射
$L \to (\underline{\Lambda}^\wedge)^{\oplus t}[-b + 1]$
对应于映射
$(\underline{M^{b - 1}})^\wedge \to (\underline{\Lambda}^\wedge)^{\oplus t}$
，它使我们可以将此复形延长一项。由三角范畴的性质，相应复形在衍生范畴中
同构于 $K$。证明完毕。
\end{proof}

\noindent
受可构造 $\Lambda$-层情形的启发，我们引入如下概念。

\begin{definition}
\label{proetale-definition-adic-constructible}
设 $X$ 为概形。设 $\Lambda$ 为 Noether 环，$I \subset \Lambda$ 为理想。
设 $K \in D(X_\proetale, \Lambda)$。
\begin{enumerate}
\item 若存在由有限投射 $\Lambda^\wedge$-模组成的有限复形 $M^\bullet$，
使得 $K$ 局部同构于
$$
\underline{M^a}^\wedge \to \ldots \to \underline{M^b}^\wedge
$$
，则称 $K$ 为{\it 进 lisse}\footnote{这可能不是标准记法。}。
\item 若对每个仿射开集 $U \subset X$，都存在分解为可构造局部闭子概形的
分解 $U = \coprod U_i$，使得 $K|_{U_i}$ 为进 lisse，则称 $K$
为{\it 进可构造}\footnote{这可能不是标准记法。}。
\end{enumerate}
\end{definition}

\noindent
引理 \ref{proetale-lemma-local-structure-constructible-complex} 中得到的局部结构
与进 lisse 复形的结构之间的差别在于：引理
\ref{proetale-lemma-local-structure-constructible-complex} 中的映射
$\underline{M^i}^\wedge \to \underline{M^{i + 1}}^\wedge$ 未必是常值的，
而上述定义要求它们为常值映射。

\begin{lemma}
\label{proetale-lemma-weakly-contractible-locally-constant-ML}
设 $X$ 为弱可缩仿射概形。设 $\Lambda$ 为 Noether 环，
$I \subset \Lambda$ 为理想。设 $K$ 为 $D_{cons}(X, \Lambda)$ 的对象，且
$K \otimes_\Lambda^\mathbf{L} \underline{\Lambda/I^n}$
在 $D(X_\proetale, \Lambda/I^n)$ 中同构于某个常值
$\Lambda/I^n$-模层复形。则
$$
H^0(X, K \otimes_\Lambda^\mathbf{L} \Lambda/I^n)
$$
满足 Mittag-Leffler 条件。
\end{lemma}

\begin{proof}
设 $K \otimes_\Lambda^\mathbf{L} \underline{\Lambda/I^n}$ 同构于
$\underline{E_n}$，其中 $E_n$ 是 $D(\Lambda/I^n)$ 的某个对象。
由于 $K \otimes_\Lambda^\mathbf{L} \underline{\Lambda/I}$ 具有有限 Tor
维数及有限型上同调层，可知 $E_1$ 完美（见代数补充，引理
\ref{more-algebra-lemma-perfect}）。转移映射
$$
K \otimes_\Lambda^\mathbf{L} \underline{\Lambda/I^{n + 1}}
\to
K \otimes_\Lambda^\mathbf{L} \underline{\Lambda/I^n}
$$
局部来自 $D(\Lambda/I^{n + 1})$ 中的复形映射
$E_{n + 1} \to E_n$（这样的映射可能有许多不同选择）；见位点上的上同调，
引理 \ref{sites-cohomology-lemma-locally-constant-map}。对每个 $n$
选取一个这样的映射，并注意到它诱导 $D(\Lambda/I^n)$ 中的同构
$E_{n + 1} \otimes_{\Lambda/I^{n + 1}}^\mathbf{L} \Lambda/I^n \to E_n$
。由代数补充，引理
\ref{more-algebra-lemma-Rlim-perfect-gives-complete}，可以找到由有限投射
$\Lambda^\wedge$-模组成的有限复形 $M^\bullet$，以及
$D(\Lambda/I^n)$ 中与转移映射相容的同构
$M^\bullet/I^nM^\bullet \to E_n$。

\medskip\noindent
现在注意，对于每个有限指标组 $n > m > k$，三重映射
$$
H^0(X, K \otimes_\Lambda^\mathbf{L} \Lambda/I^n)
\to
H^0(X, K \otimes_\Lambda^\mathbf{L} \Lambda/I^m)
\to
H^0(X, K \otimes_\Lambda^\mathbf{L} \Lambda/I^k)
$$
同构于
$$
H^0(X, \underline{M^\bullet/I^nM^\bullet})
\to
H^0(X, \underline{M^\bullet/I^mM^\bullet})
\to
H^0(X, \underline{M^\bullet/I^kM^\bullet})
$$
事实上，任选同构
$$
\underline{M^\bullet/I^nM^\bullet} \to
K \otimes_\Lambda^\mathbf{L} \Lambda/I^n
$$
，它模 $I^m$ 和 $I^k$ 诱导类似同构，因而断言成立。因此，为证明本引理，
只需证明系统
$H^0(X, \underline{M^\bullet/I^nM^\bullet})$ 满足 Mittag-Leffler 条件。
由于在 $X$ 上取截面是正合的，只需证明 $\Lambda$-模系统
$$
H^0(M^\bullet/I^nM^\bullet)
$$
满足 Mittag-Leffler 条件。令 $A = \Lambda^\wedge$，考察谱序列
$$
\text{Tor}_{-p}^A(H^q(M^\bullet), A/I^nA) \Rightarrow
H^{p + q}(M^\bullet/I^nM^\bullet)
$$
由代数补充，引理 \ref{more-algebra-lemma-tor-strictly-pro-zero}，当
$p > 0$ 时，pro-系统
$\{\text{Tor}_{-p}^A(H^q(M^\bullet), A/I^nA)\}$ 为零。因此 pro-系统
$\{H^0(M^\bullet/I^nM^\bullet)\}$ 等于 pro-系统
$\{H^0(M^\bullet)/I^nH^0(M^\bullet)\}$，引理得证。
\end{proof}

\begin{lemma}
\label{proetale-lemma-connected-adic-lisse}
设 $X$ 为连通概形。设 $\Lambda$ 为 Noether 环，$I \subset \Lambda$
为理想。若 $K$ 属于 $D_{cons}(X, \Lambda)$，且
$K \otimes_\Lambda \underline{\Lambda/I}$ 具有局部常值上同调层，
则 $K$ 为进 lisse（定义 \ref{proetale-definition-adic-constructible}）。
\end{lemma}

\begin{proof}
记 $K_n = K \otimes_\Lambda^\mathbf{L} \underline{\Lambda/I^n}$。
以下将不再说明地使用引理 \ref{proetale-lemma-describe-constructible-complexes}
的结论。由位点上的上同调，引理
\ref{sites-cohomology-lemma-locally-constant-bounded}，可知对所有 $n$，
$K_n$ 均具有局部常值上同调层。有 $K_n = \epsilon^{-1}L_n$，其中某个
$L_n$ 属于 $D_{ctf}(X_\etale, \Lambda/I^n)$ 且具有局部常值上同调层。由 Étale 上同调，引理
\ref{etale-cohomology-lemma-connected-ctf-locally-constant}，存在完美对象
$M_n \in D(\Lambda/I^n)$，使得 $L_n$ 在 étale 局部同构于
$\underline{M_n}$。与 $K_{n + 1} \to K_n$ 对应的映射
$L_{n + 1} \to L_n$ 诱导同构
$L_{n + 1} \otimes_{\Lambda/I^{n + 1}}^\mathbf{L} \underline{\Lambda/I^n}
\to L_n$。在 $X$ 上局部考察，可知 $D(\Lambda/I^{n + 1})$ 中存在映射
$M_{n + 1} \to M_n$，其诱导同构
$M_{n + 1} \otimes_{\Lambda/I^{n + 1}} \Lambda/I^n \to M_n$；见
位点上的上同调，引理 \ref{sites-cohomology-lemma-locally-constant-map}。
固定一组这样的映射。由代数补充，引理
\ref{more-algebra-lemma-Rlim-perfect-gives-complete}，可以找到由有限投射
$\Lambda^\wedge$-模组成的有限复形 $M^\bullet$，以及
$D(\Lambda/I^n)$ 中与转移映射相容的同构
$M^\bullet/I^nM^\bullet \to M_n$。为完成证明，我们将证明 $K$ 局部同构于
$$
\underline{M^\bullet}^\wedge =
\lim \underline{M^\bullet/I^nM^\bullet} =
R\lim \underline{M^\bullet/I^nM^\bullet}
$$
令 $E^\bullet$ 为 $M^\bullet$ 的对偶复形，见代数补充，引理
\ref{more-algebra-lemma-dual-perfect-complex} 及其证明。考察
$D(X_\proetale, \Lambda/I^n)$ 的对象
$$
H_n =
R\SheafHom_{\Lambda/I^n}(\underline{M^\bullet/I^nM^\bullet}, K_n) =
\underline{E^\bullet/I^nE^\bullet} \otimes_{\Lambda/I^n}^\mathbf{L} K_n
$$
。模去 $I^n$ 定义转移映射 $H_{n + 1} \to H_n$。令
$H = R\lim H_n$。则 $H$ 是 $D_{cons}(X, \Lambda)$ 的对象（略去细节），
且 $H \otimes_\Lambda^\mathbf{L} \underline{\Lambda/I^n} = H_n$。
选取覆盖 $\{W_t \to X\}_{t \in T}$，其中每个 $W_t$ 都是仿射且弱可缩的。
由我们对 $M^\bullet$ 的选择，可知
\begin{align*}
H_n|_{W_t} & \cong 
R\SheafHom_{\Lambda/I^n}(\underline{M^\bullet/I^nM^\bullet},
\underline{M^\bullet/I^nM^\bullet}) \\
& =
\underline{
\text{Tot}(E^\bullet/I^nE^\bullet \otimes_{\Lambda/I^n} M^\bullet/I^nM^\bullet)
}
\end{align*}
因此可以将引理 \ref{proetale-lemma-weakly-contractible-locally-constant-ML}
应用于 $H = R\lim H_n$。于是系统 $H^0(W_t, H_n)$ 满足
Mittag-Leffler 条件。由于对所有 $n \gg 1$，$H^0(W_t, H_n)$ 中都存在
一个元素，它映到
$$
H^0(W_t, H_1) = \Hom(\underline{M^\bullet/IM^\bullet}, K_1)
$$
中的同构，故可以在逆极限中找到元素 $(\varphi_{t, n})$，它模 $I$
给出同构。于是
$$
R\lim \varphi_{t, n} :
\underline{M^\bullet}^\wedge|_{W_t} =
R\lim \underline{M^\bullet/I^nM^\bullet}|_{W_t}
\longrightarrow
R\lim K_n|_{W_t} = K|_{W_t}
$$
是同构。证明完毕。
\end{proof}

\begin{proposition}
\label{proetale-proposition-Noetherian-adic-constructible}
设 $X$ 为 Noether 概形。设 $\Lambda$ 为 Noether 环，$I \subset \Lambda$
为理想。设 $K$ 为 $D_{cons}(X, \Lambda)$ 的对象。则 $K$ 进可构造
（定义 \ref{proetale-definition-adic-constructible}）。
\end{proposition}

\begin{proof}
这由引理 \ref{proetale-lemma-connected-adic-lisse}、Noether 概形局部连通这一事实
（拓扑，引理 \ref{topology-lemma-locally-Noetherian-locally-connected}）
以及相关定义得出。
\end{proof}







\section{固有基变换}
\label{proetale-section-proper-base-change}

\noindent
本节说明如何利用 étale 上同调的固有基变换定理，证明 pro-étale 位点上
衍生完备对象的固有基变换定理。这遵循如下总体思路：只用 pro-étale
拓扑处理“极限问题”，而其余一切则用已对 étale 拓扑证明的结果处理。

\begin{theorem}
\label{proetale-theorem-proper-base-change}
设 $f : X \to Y$ 为概形的固有态射，$g : Y' \to Y$ 为概形态射，
它们给出基变换图
$$
\xymatrix{
X' \ar[r]_{g'} \ar[d]_{f'} & X \ar[d]^f \\
Y' \ar[r]^g & Y
}
$$
设 $\Lambda$ 为 Noether 环，$I \subset \Lambda$ 为理想，并且
$\Lambda/I$ 是挠的。设 $K$ 为 $D(X_\proetale)$ 的对象，满足
\begin{enumerate}
\item $K$ 衍生完备；
\item $K \otimes_\Lambda^\mathbf{L} \underline{\Lambda/I^n}$ 下有界，
且其上同调层来自 $X_\etale$；
\item $\Lambda/I^n$ 是完美 $\Lambda$-模\footnote{若 $K$ 为可构造复形，
则可去掉此假设；见 \cite{BS}。}。
\end{enumerate}
则基变换映射
$$
Lg_{comp}^*Rf_*K \longrightarrow Rf'_*L(g')^*_{comp}K
$$
是同构。
\end{theorem}

\begin{proof}
略去基变换映射的构造（它只使用衍生直像与完备化衍生拉回的形式性质；
参见位点上的上同调，评注 \ref{sites-cohomology-remark-base-change}）。
记 $K_n = K \otimes^\mathbf{L}_\Lambda \underline{\Lambda/I^n}$。
由于 $K$ 衍生完备，由引理 \ref{proetale-lemma-naive-completion} 有
$K = R\lim K_n$。由引理 \ref{proetale-lemma-pushforward-Noetherian-case} 和
\ref{proetale-lemma-naive-completion}，可以展开左边：
$$
Lg_{comp}^* Rf_* K =
R\lim Lg^*(Rf_*K)\otimes^\mathbf{L}_\Lambda \underline{\Lambda/I^n} =
R\lim Lg^* Rf_* K_n
$$
最后一个等式来自 $\Lambda/I^n$ 是完美模以及投影公式（位点上的上同调，
引理 \ref{sites-cohomology-lemma-projection-formula}）。利用引理
\ref{proetale-lemma-pushforward-Noetherian-case}，可以展开右边：
$$
Rf'_* L(g')^*_{comp} K =
Rf'_* R\lim L(g')^* K_n  =
R\lim Rf'_* L(g')^* K_n
$$
最后一个等式是因为 $Rf'_*$ 与 $R\lim$ 可交换（位点上的上同调，引理
\ref{sites-cohomology-lemma-Rf-commutes-with-Rlim}）。因此，只需证明映射
$$
Lg^* Rf_* K_n \longrightarrow Rf'_* L(g')^* K_n
$$
均为同构。由引理 \ref{proetale-lemma-compare-derived} 和第二个条件，可以对某个
$L_n \in D^+(X_\etale, \Lambda/I^n)$ 写成
$K_n = \epsilon^{-1}L_n$。由引理 \ref{proetale-lemma-morphism-comparison} 以及
$\epsilon^{-1}$ 与拉回可交换这一事实，得到
$$
Lg^* Rf_* K_n =
Lg^* Rf_* \epsilon^*L_n =
Lg^* \epsilon^{-1} Rf_* L_n =
\epsilon^{-1} Lg^* Rf_* L_n
$$
以及
$$
Rf'_* L(g')^* K_n =
Rf'_* L(g')^* \epsilon^{-1} L_n =
Rf'_* \epsilon^{-1} L(g')^* L_n =
\epsilon^{-1} Rf'_* L(g')^* L_n
$$
（这里也使用了 $L_n$ 下有界）。最后，由 étale 上同调的固有基变换定理
（Étale 上同调，定理
\ref{etale-cohomology-theorem-proper-base-change}），有
$$
Lg^* Rf_* L_n = Rf'_* L(g')^* L_n
$$
（再次使用了 $L_n$ 下有界），定理得证。
\end{proof}






\section{更换部分宇宙}
\label{proetale-section-change-universe}

\noindent
我们建议读者跳过本节：这里将证明 pro-étale 拓扑中层的上同调不依赖于
部分宇宙的选择。具体而言，下面引理
\ref{proetale-lemma-proetale-cohomology-independent-partial-universe} 的函子
$g_*$ 是小 pro-étale 拓扑斯的一个嵌入，并且不改变上同调。
对大 pro-étale 位点，引理 \ref{proetale-lemma-change-alpha} 和
\ref{proetale-lemma-cohomology-enlarge-partial-universe} 给出本质上相同的结论。

\medskip\noindent
但首先，按照第 \ref{proetale-section-proetale} 节中的承诺，我们证明大 pro-étale
位点 $\Sch_\proetale$ 上的拓扑在某种意义下由全体概形之范畴上的
pro-étale 拓扑诱导。

\begin{lemma}
\label{proetale-lemma-proetale-induced}
设 $\Sch_\proetale$ 为定义 \ref{proetale-definition-big-proetale-site} 所述的大
pro-étale 位点。设 $T \in \Ob(\Sch_\proetale)$。任取 $T$ 的 pro-étale
覆盖 $\{T_i \to T\}_{i \in I}$。则位点 $\Sch_\proetale$ 中存在 $T$
的覆盖 $\{U_j \to T\}_{j \in J}$，它细分 $\{T_i \to T\}_{i \in I}$。
\end{lemma}

\begin{proof}
首先，取引理 \ref{proetale-lemma-get-many-weakly-contractible} 所述的覆盖
$\{V_k \to T\}$。于是 pro-étale 覆盖
$\{T_i \times_T V_k \to V_k\}$ 可由有限不交开覆盖
$V_k = V_{k, 1} \amalg \ldots \amalg V_{k, n_k}$ 细分，见引理
\ref{proetale-lemma-w-contractible-proetale-cover}。于是
$\{V_{k, i} \to T\}$ 是 $\Sch_\proetale$ 中的覆盖，并细分
$\{T_i \to T\}_{i \in I}$。
\end{proof}

\noindent
我们先陈述并证明小 pro-étale 位点的比较。注意，我们并未断言概形的小
pro-étale 拓扑斯不依赖于部分宇宙的选择；这并不成立，与小 étale
拓扑斯的情形不同（Étale 上同调，引理
\ref{etale-cohomology-lemma-etale-topos-independent-partial-universe}）。

\begin{lemma}
\label{proetale-lemma-proetale-cohomology-independent-partial-universe}
设 $S$ 为概形。设 $S_\proetale \subset S_\proetale'$ 为按定义
\ref{proetale-definition-big-small-proetale} 构造的 $S$ 的两个小 pro-étale 位点。
则包含函子满足位点，引理 \ref{sites-lemma-bigger-site} 的假设。
因而存在拓扑斯态射
$$
\xymatrix{
\Sh(S_\proetale) \ar[r]^g &
\Sh(S_\proetale') \ar[r]^f &
\Sh(S_\proetale)
}
$$
其复合同构于恒等态射，且 $f_* = g^{-1}$。此外，
\begin{enumerate}
\item 对 $\mathcal{F}' \in \textit{Ab}(S_\proetale')$，有
$H^p(S_\proetale', \mathcal{F}') = H^p(S_\proetale, g^{-1}\mathcal{F}')$,
\item 对 $\mathcal{F} \in \textit{Ab}(S_\proetale)$，有
$$
H^p(S_\proetale, \mathcal{F}) =
H^p(S_\proetale', g_*\mathcal{F}) =
H^p(S_\proetale', f^{-1}\mathcal{F}).
$$
\end{enumerate}
\end{lemma}

\begin{proof}
包含函子全忠实且连续。我们已经看到，$S_\proetale$ 和
$S_\proetale'$ 有纤维积及终对象，并且此函子与它们可交换
（引理 \ref{proetale-lemma-fibre-products-proetale}）。由引理
\ref{proetale-lemma-proetale-induced}，包含函子余连续。因此 $f$ 与 $g$
的存在性由位点，引理 \ref{sites-lemma-bigger-site} 得出。
第 (1) 点中的等式即位点上的上同调，引理
\ref{sites-cohomology-lemma-cohomology-bigger-site}。第 (2) 点由第 (1) 点
以及 $\mathcal{F} = g^{-1}g_*\mathcal{F} = g^{-1}f^{-1}\mathcal{F}$ 得出。
\end{proof}

\noindent
下面证明大 pro-étale 拓扑斯的相应结果。

\begin{lemma}
\label{proetale-lemma-change-alpha}
设给定定义 \ref{proetale-definition-big-proetale-site} 所述的大位点
$\Sch_\proetale$ 和 $\Sch'_\proetale$。假设 $\Sch_\proetale$ 包含于
$\Sch'_\proetale$。包含函子 $\Sch_\proetale \to \Sch'_\proetale$
满足位点，引理 \ref{sites-lemma-bigger-site} 的假设。存在拓扑斯态射
\begin{eqnarray*}
g : \Sh(\Sch_\proetale) &
\longrightarrow &
\Sh(\Sch'_\proetale) \\
f : \Sh(\Sch'_\proetale) &
\longrightarrow &
\Sh(\Sch_\proetale)
\end{eqnarray*}
使得 $f \circ g \cong \text{id}$。对 $\Sch_\proetale$ 的任意对象 $S$，
包含函子 $(\Sch/S)_\proetale \to (\Sch'/S)_\proetale$ 也满足位点，引理
\ref{sites-lemma-bigger-site} 的假设。因而类似地得到态射
\begin{eqnarray*}
g : \Sh((\Sch/S)_\proetale) &
\longrightarrow &
\Sh((\Sch'/S)_\proetale) \\
f : \Sh((\Sch'/S)_\proetale) &
\longrightarrow &
\Sh((\Sch/S)_\proetale)
\end{eqnarray*}
且 $f \circ g \cong \text{id}$。
\end{lemma}

\begin{proof}
对函子
$\Sch_\proetale \to \Sch'_\proetale$ 和
$(\Sch/S)_\proetale \to (\Sch'/S)_\proetale$，位点，引理
\ref{sites-lemma-bigger-site} 的假设 (b)、(c) 和 (e) 是立即成立的。
性质 (a) 由引理 \ref{proetale-lemma-proetale-induced} 成立。性质 (d) 成立，是因为
范畴 $\Sch_\proetale$、$\Sch'_\proetale$ 中的纤维积存在，并且与概形范畴
中的纤维积相容。
\end{proof}

\begin{lemma}
\label{proetale-lemma-cohomology-enlarge-partial-universe}
设 $S$ 为概形。设 $(\Sch/S)_\proetale$ 和 $(\Sch'/S)_\proetale$ 为定义
\ref{proetale-definition-big-small-proetale} 所述的 $S$ 的两个大 pro-étale 位点。
假设前者包含于后者。在此情形下，
\begin{enumerate}
\item 对定义于 $(\Sch'/S)_\proetale$ 上的任意阿贝尔层 $\mathcal{F}'$，
以及 $(\Sch/S)_\proetale$ 的任意对象 $U$，有
$$
H^p(U, \mathcal{F}'|_{(\Sch/S)_\proetale}) =
H^p(U, \mathcal{F}')
$$
换言之，在较大位点中计算的 $\mathcal{F}'$ 在 $U$ 上的上同调，与把
$\mathcal{F}'$ 限制到较小位点后在 $U$ 上计算的上同调一致。
\item 对 $(\Sch/S)_\proetale$ 上的任意阿贝尔层 $\mathcal{F}$，存在
$(\Sch/S)_\proetale'$ 上的阿贝尔层 $\mathcal{F}'$，其到
$(\Sch/S)_\proetale$ 的限制同构于 $\mathcal{F}$。
\end{enumerate}
\end{lemma}

\begin{proof}
由引理 \ref{proetale-lemma-change-alpha}，包含函子
$(\Sch/S)_\proetale \to (\Sch'/S)_\proetale$ 满足位点，引理
\ref{sites-lemma-bigger-site} 的假设。这推出第 (2) 点，而第 (1) 点由
位点上的上同调，引理 \ref{sites-cohomology-lemma-cohomology-bigger-site}
得出。
\end{proof}
% Generated mechanically from the frozen relative-cycles.tex target.
% Do not edit this generated file; edit the bound locale source instead.

% OK, start here.
%

\setcounter{chapter}{61}
\chapter{相对循环}



\phantomsection
\label{relative-cycles-section-phantom}




\section{引言}
\label{relative-cycles-section-introduction}

\noindent
基础参考文献是 \cite{SV}。

\medskip\noindent
本章只定义 \cite{SV} 中所谓的普遍整相对循环。与原文相比，这一选择使理论的
展开略为简洁，但当然也会损失一些内容。

\medskip\noindent
固定 Noether 概形之间的有限型态射 $X \to S$。$X/S$ 各纤维上的
$r$-循环族 $\alpha$，就是一个集合 $\alpha = (\alpha_s)_{s \in S}$，
其中 $\alpha_s \in Z_r(X_s)$。沿任意态射 $g : S' \to S$ 对 $\alpha$
作基变换而得到 $g^*\alpha$ 的方式是显然的。若 $\alpha$ 与特化相容，
则称 $\alpha$ 为 {\it $X/S$ 上的相对 $r$-循环}。换言之，对任意态射
$g : S' \to S$，其中 $S'$ 是一个离散赋值环的谱，我们要求
$g^*\alpha$ 的一般纤维特化为 $g^*\alpha$ 的闭纤维。见第
\ref{relative-cycles-section-families-specialization} 节。






\section{约定与记号}
\label{relative-cycles-section-conventions}

\noindent
关于固定 Noether 基上局部有限型概形之循环的约定与记号，请参阅
“Chow 同调与陈类”一章；见 Chow 同调，第 \ref{chow-section-setup} 节及后文。

\medskip\noindent
特别地，若 $X$ 在域 $k$ 上局部有限型，则 $Z_r(X)$ 表示 $r$ 维循环群；
见 Chow 同调，例 \ref{chow-example-field} 与第
\ref{chow-section-cycles} 节。给定满足 $\dim(Z) = r$ 的整闭子概形
$Z \subset X$，有 $[Z] \in Z_r(X)$；若 $X$ 拟紧，则 $Z_r(X)$ 是以这些
类为基的自由阿贝尔群。




\section{相对于域的循环}
\label{relative-cycles-section-relative-fields}

\noindent
设 $k$ 为域，$X$ 为 $k$ 上的局部代数概形，并设 $r \geq 0$ 为整数。
在此情形下，我们有 $X$ 上的 $r$-循环群 $Z_r(X)$；见第
\ref{relative-cycles-section-conventions} 节。

\medskip\noindent
{\bf 基变换。} 对任意域扩张 $k'/k$，都有基变换映射
$Z_r(X) \to Z_r(X_{k'})$；见 Chow 同调，第
\ref{chow-section-change-base} 节。具体而言，给定维数为 $r$ 的整闭子概形
$Z \subset X$，我们把 $[Z] \in Z_r(X)$ 映到与闭子概形
$Z_{k'} \subset X_{k'}$ 相伴的 $r$-循环
$[Z_{k'}]_r \in Z_r(X_{k'})$（一般而言，$Z_{k'}$ 当然既不必不可约，
也不必既约）。基变换映射 $Z_r(X) \to Z_r(X_{k'})$ 总是单射。

\begin{lemma}
\label{relative-cycles-lemma-multiplicities-field-extension}
设 $K/k$ 为域扩张，$Z$ 为 $k$ 上的整局部代数概形。不可约分支
$Z' \subset Z_K$ 的重数 $m_{Z', Z_K}$ 等于 $1$ 或 $k$ 的特征的某个幂。
\end{lemma}

\begin{proof}
若 $k$ 的特征为零，则 $k$ 完美，并且由“代数簇”引理
\ref{varieties-lemma-geometrically-reduced}，$X_K$ 既约，故重数总是 $1$。
现假设 $k$ 的特征为 $p > 0$。设 $L$ 为 $Z$ 的函数域。由于 $Z$ 在 $k$
上局部代数，域扩张 $L/k$ 有限生成。环 $K \otimes_k L$ 是 Noether 环
（代数，引理 \ref{algebra-lemma-Noetherian-field-extension}）。转化为代数
语言，我们必须证明：对每个极小素理想 $\mathfrak q$，Artin 局部环
$(K \otimes_k L)_\mathfrak q$ 的长度都是 $p$ 的幂。

\medskip\noindent
设 $L'/L$ 为有限纯不可分扩张，例如次数为 $p^n$。则
$K \otimes_k L \subset K \otimes_k L'$ 是秩为 $p^n$ 的有限自由环映射；
它在谱上诱导同胚，并诱导纯不可分的剩余域扩张。因此，对上述每个极小素
理想 $\mathfrak q$，都有唯一的极小素理想
$\mathfrak q' \subset K \otimes_k L'$ 位于其上，并且
$$
p^n \text{length}((K \otimes_k L)_\mathfrak q) =
[\kappa(\mathfrak q') : \kappa(\mathfrak q)]
\text{length}((K \otimes_k L')_{\mathfrak q'})
$$
这是把代数，引理 \ref{algebra-lemma-pushdown-module} 应用于
$M = (K \otimes_k L')_{\mathfrak q'} \cong
(K \otimes_k L)_{\mathfrak q}^{\oplus p^n}$ 所得。由于
$[\kappa(\mathfrak q') : \kappa(\mathfrak q)]$ 是 $p$ 的幂，可知只需对
$L'$ 与 $\mathfrak q'$ 证明该断言。

\medskip\noindent
由前一段和代数，引理 \ref{algebra-lemma-make-separable}，可以假设存在
子域塔 $L/k'/k$，使得 $L/k'$ 可分且 $k'/k$ 有限纯不可分。于是
$K \otimes_k k'$ 是 Artin 局部环。前一段的论证（取 $L = k$、
$L' = k'$）表明 $\text{length}(K \otimes_k k')$ 是 $p$ 的幂。又由于
$L/k'$ 是一个光滑 $k'$-代数的局部化（代数，引理
\ref{algebra-lemma-localization-smooth-separable}），故
$S = (K \otimes_k L)_\mathfrak q$ 是某个光滑
$R = K \otimes_k k'$-代数在极小素理想处的局部化。因此
$R \to S$ 是 Artin 局部环之间的平坦局部同态，且
$\mathfrak m_R S = \mathfrak m_S$。由代数，引理
\ref{algebra-lemma-pullback-module} 得
$\text{length}(K \otimes_k k') = \text{length}(R) =
\text{length}(S) = \text{length}((K \otimes_k L)_\mathfrak q)$
从而证明完成。
\end{proof}

\begin{lemma}
\label{relative-cycles-lemma-how-different}
设 $k$ 为特征 $p > 0$ 的域，其完美闭包为 $k^{perf}$。设 $X$ 为 $k$ 上的
代数概形，$r \geq 0$ 为整数。单射
$Z_r(X) \to Z_r(X_{k^{perf}})$ 的余核是 $p$-幂挠模（代数续篇，定义
\ref{more-algebra-definition-f-power-torsion}）。
\end{lemma}

\begin{proof}
由于 $X$ 拟紧，阿贝尔群 $Z_r(X)$ 是自由的，其基由维数为 $r$ 的整闭子概形
给出；$Z_r(X_{k^{perf}})$ 亦然。由于
$X_{k^{perf}} \to X$ 是同胚，映射 $Z_r(X) \to Z_r(X_{k^{perf}})$
为单射且余核为挠群。由引理 \ref{relative-cycles-lemma-multiplicities-field-extension}，
余核中的每个元素都是 $p$-幂挠的。
\end{proof}





\section{循环的特化}
\label{relative-cycles-section-specialization}

\noindent
设 $R$ 为离散赋值环，其分式域为 $K$，剩余域为 $\kappa$。设 $X$ 为
$R$ 上局部有限型的概形，并设 $r \geq 0$。存在特化映射
$$
sp_{X/R} : Z_r(X_K) \longrightarrow Z_r(X_\kappa)
$$
其定义如下。对维数为 $r$ 的整闭子概形 $Z \subset X_K$，以
$\overline{Z}$ 表示 $Z \to X$ 的概形论像。令 $sp_{X/R}$ 为满足下式的
唯一 $\mathbf{Z}$-线性映射：
$$
sp_{X/R}([Z]) = [\overline{Z}_\kappa]_r
$$
简要说明该定义为何良定。首先，态射 $X_K \to X$ 拟紧，故态射
$Z \to X$ 拟紧。因此，由“态射”引理
\ref{morphisms-lemma-flat-base-change-scheme-theoretic-image}，取
$Z \to X$ 的概形论像与平坦基变换可交换。特别地，基变换回 $X_K$ 后，
有 $Z = \overline{Z}_K$。由于 $Z$ 是整的，$\overline{Z}$ 当然也是整的；
事实上，它等于以 $Z$ 的一般点（之像）为一般点的唯一整闭子概形。由
“代数簇”引理
\ref{varieties-lemma-dominate-valuation-ring-dimension-fibres}，
$Z_\kappa$ 是维数为 $r$ 的等维概形。

\begin{lemma}
\label{relative-cycles-lemma-specialization-module}
设 $R$ 为离散赋值环，其分式域为 $K$，剩余域为 $\kappa$。设 $X$ 为
$R$ 上局部有限型的概形，$r \geq 0$。设 $\mathcal{F}$ 为在 $R$ 上平坦的
凝聚 $\mathcal{O}_X$-模。假设
$\dim(\text{Supp}(\mathcal{F}_K)) \leq r$。则
$\dim(\text{Supp}(\mathcal{F}_\kappa)) \leq r$，并且
$$
sp_{X/R}([\mathcal{F}_K]_r) = [\mathcal{F}_\kappa]_r
$$
\end{lemma}

\begin{proof}
关于维数的断言由“态射续篇”引理
\ref{more-morphisms-lemma-relative-dimension-support-flat} 得出。
设 $x$ 为维数 $r$ 的整闭子概形 $Z \subset X_\kappa$ 的一般点。为完成
证明，我们将说明等式左边 (L) 与右边 (R) 中 $[Z]$ 的系数相同。

\medskip\noindent
令 $A = \mathcal{O}_{X, x}$、$M = \mathcal{F}_x$。注意，$M$ 是在
$R$ 上平坦的有限 $A$-模。取素元 $\pi \in R$，使得
$A/\pi A = \mathcal{O}_{X_\kappa, x}$。由 Chow 同调，引理
\ref{chow-lemma-additivity-divisors-restricted}，有
$$
\sum\nolimits_i \text{length}_A(A/(\pi, \mathfrak q_i))
\text{length}_{A_{\mathfrak q_i}}(M_{\mathfrak q_i}) =
\text{length}_A(M/\pi M)
$$
其中求和遍历 $M$ 的支撑中的极小素理想 $\mathfrak q_i$。由于 $\pi$ 是
$M$ 上的非零因子，有 $\pi \not \in \mathfrak q_i$；因此，这些素理想
对应于 $\mathcal{F}_K$ 的支撑中那些特化到所选
$x \in X_\kappa$ 的一般点 $y_i \in X_K$。所以左边就是 (L) 中 $[Z]$
的系数。当然，$\text{length}_A(M/\pi M)$ 是 (R) 中 $[Z]$ 的系数，
证明完成。
\end{proof}

\begin{lemma}
\label{relative-cycles-lemma-specialization-closed}
设 $R$ 为离散赋值环，其分式域为 $K$，剩余域为 $\kappa$。设 $X$ 为
$R$ 上局部有限型的概形，$r \geq 0$。设 $W \subset X$ 为在 $R$ 上平坦的
闭子概形。假设 $\dim(W_K) \leq r$。则 $\dim(W_\kappa) \leq r$，并且
$$
sp_{X/R}([W_K]_r) = [W_\kappa]_r
$$
\end{lemma}

\begin{proof}
取 $\mathcal{F} = \mathcal{O}_W$，这就是引理
\ref{relative-cycles-lemma-specialization-module} 的特殊情形。见 Chow 同调，引理
\ref{chow-lemma-cycle-closed-coherent}。
\end{proof}

\begin{lemma}
\label{relative-cycles-lemma-specialization-extension}
设 $R'/R$ 为离散赋值环的扩张，它诱导分式域扩张 $K'/K$ 与剩余域扩张
$\kappa'/\kappa$（代数续篇，定义
\ref{more-algebra-definition-extension-discrete-valuation-rings}）。
设 $X$ 在 $R$ 上局部有限型，并记 $X' = X_{R'}$。则图
$$
\xymatrix{
Z_r(X'_{K'}) \ar[rr]_{sp_{X'/R'}} & & Z_r(X'_{\kappa'}) \\
Z_r(X_K) \ar[rr]^{sp_{X/R}} \ar[u] & & Z_r(X_\kappa) \ar[u]
}
$$
可交换，其中 $r \geq 0$，竖直箭头是基变换映射。
\end{lemma}

\begin{proof}
注意，$X'_{K'} = X_{K'} = X_K \times_{\Spec(K)} \Spec(K')$；闭纤维亦然，
故竖直箭头确有意义（见第 \ref{relative-cycles-section-relative-fields} 节）。现设
$Z \subset X_K$ 为整闭子概形，其概形论像为 $\overline{Z} \subset X$。
则 $Z_{K'} \subset X_{K'}$ 是闭子概形，其概形论像为
$\overline{Z}_{R'} \subset X_{R'}$。按定义，$[Z]$ 的基变换是
$[Z_{K'}]_r = [\overline{Z}_{K'}]_r$。我们有
$$
sp_{X/R}([Z]) = [\overline{Z}_\kappa]_r
\quad\text{且}\quad
sp_{X'/R'}([\overline{Z}_{K'}]_r) = [(\overline{Z}_{R'})_{\kappa'}]_r
$$
这是由引理 \ref{relative-cycles-lemma-specialization-module} 得出的。又因为
$(\overline{Z}_{R'})_{\kappa'} = (\overline{Z}_\kappa)_{\kappa'}$，
故得结论。
\end{proof}

\begin{lemma}
\label{relative-cycles-lemma-specialization-flat-pullback}
设 $R$ 为离散赋值环，其分式域为 $K$，剩余域为 $\kappa$。设 $X$ 为
$R$ 上局部有限型的概形。设 $f : X' \to X$ 为局部有限型、平坦且相对
维数为 $e$ 的态射。则图
$$
\xymatrix{
Z_{r + e}(X'_K) \ar[rr]_{sp_{X'/R}} & & Z_{r + e}(X'_\kappa) \\
Z_r(X_K) \ar[rr]^{sp_{X/R}} \ar[u] & & Z_r(X_\kappa) \ar[u]
}
$$
可交换，其中 $r \geq 0$，竖直箭头由平坦拉回给出。
\end{lemma}

\begin{proof}
设 $Z \subset X$ 为支配 $R$ 的整闭子概形。由 $sp_{X/R}$ 的构造，有
$sp_{X/R}([Z_K]) = [Z_\kappa]_r$，而这刻画了特化映射。置
$Z' = f^{-1}(Z) = X' \times_X Z$。由于 $R$ 是赋值环，$Z$ 在 $R$ 上
平坦。因此 $Z'$ 在 $R$ 上平坦，并且由引理
\ref{relative-cycles-lemma-specialization-closed}，有
$sp_{X'/R}([Z'_K]_{r + e}) = [Z'_\kappa]_{r + e}$
又由 Chow 同调，引理 \ref{chow-lemma-pullback-coherent}，有
$f_K^*[Z_K] = [Z'_K]_{r + e}$ 及
$f_\kappa^*[Z_\kappa]_r = [Z'_\kappa]_{r + e}$，故得结论。
\end{proof}

\begin{lemma}
\label{relative-cycles-lemma-specialization-proper-pushforward}
设 $R$ 为离散赋值环，其分式域为 $K$，剩余域为 $\kappa$。设
$f : X \to Y$ 为 $R$ 上局部有限型概形之间的固有态射。则图
$$
\xymatrix{
Z_r(X_K) \ar[rr]_{sp_{X/R}} \ar[d] & & Z_r(X_\kappa) \ar[d] \\
Z_r(Y_K) \ar[rr]^{sp_{Y/R}} & & Z_r(Y_\kappa)
}
$$
可交换，其中 $r \geq 0$，竖直箭头由固有前推给出。
\end{lemma}

\begin{proof}
设 $Z \subset X$ 为支配 $R$ 的整闭子概形。由 $sp_{X/R}$ 的构造，有
$sp_{X/R}([Z_K]) = [Z_\kappa]_r$，而这刻画了特化映射。置
$Z' = f(Z) \subset Y$。则 $Z'$ 是 $Y$ 中支配 $R$ 的整闭子概形。因此
$sp_{Y/R}([Z'_K]) = [Z'_\kappa]_r$。

\medskip\noindent
可把 $[Z]$ 看成 $Z_{r + 1}(X)$ 的元素。按定义，若
$\dim(Z') < r + 1$，则 $f_*[Z] = 0$；若 $Z \to Z'$ 是次数为 $d$ 的
一般有限态射，则 $f_*[Z] = d[Z']$。由于固有前推与沿
$Y_K \to Y$ 的平坦拉回可交换（Chow 同调，引理
\ref{chow-lemma-flat-pullback-proper-pushforward}），相应地有
$f_{K, *}[Z_K] = 0$ 或 $f_{K, *}[Z_K] = d[Z'_K]$。对下列交换图应用
Chow 同调，引理 \ref{chow-lemma-closed-in-X-gysin}：
$$
\xymatrix{
X_\kappa \ar[d] \ar[r]_i & X \ar[d] \\
Y_\kappa \ar[r]^j & Y
}
$$
由此得到 $f_{\kappa, *}[Z_\kappa]_r = 0$，或
$f_{\kappa, *}[Z_\kappa] = d[Z'_\kappa]_r$，因为显然有
$i^*[Z] = [Z_k]_r$ 及 $j^*[Z'] = [Z'_\kappa]_r$。综合以上结果即得结论。
\end{proof}









\section{纤维上的循环族}
\label{relative-cycles-section-cycles-fibres}

\noindent
设 $f : X \to S$ 为概形的局部有限型态射，$r \geq 0$ 为整数。
{\it $X/S$ 各纤维上的 $r$-循环族 $\alpha$} 是一个族
$$
\alpha = (\alpha_s)_{s \in S}
$$
它以概形 $S$ 的点 $s$ 为指标，其中 $\alpha_s \in Z_r(X_s)$ 是 $f$ 在
$s$ 处的概形论纤维 $X_s$ 上的 $r$-循环。对纤维上的 $r$-循环族，可以
施行多种构造。

\medskip\noindent
{\bf 基变换。} 设
$$
\xymatrix{
X' \ar[r] \ar[d] & X \ar[d]^f \\
S' \ar[r]^g & S
}
$$
为概形态射的笛卡尔方块，且 $f$ 局部有限型。设 $r \geq 0$ 为整数。
给定 $X/S$ 各纤维上的 $r$-循环族 $\alpha$，定义 $\alpha$ 的
{\it 基变换} $g^*\alpha$ 为循环族
$$
g^*\alpha = (\alpha'_{s'})_{s' \in S'}
$$
其中 $\alpha'_{s'} \in Z_r(X'_{s'})$ 是循环 $\alpha_s$ 的基变换，
$s = g(s')$；这里按第 \ref{relative-cycles-section-relative-fields} 节，使用概形论纤维的
恒等式 $X'_{s'} = X_s \times_{\Spec(\kappa(s))} \Spec(\kappa(s'))$。

\medskip\noindent
{\bf 限制。} 设 $f : X \to S$ 为概形的局部有限型态射，$r \geq 0$ 为整数。
设 $U \subset X$ 与 $V \subset S$ 为开子概形，并且 $f(U) \subset V$。
给定 $X/S$ 各纤维上的 $r$-循环族 $\alpha$，可以把 $\alpha$ 的
{\it 限制} $\alpha|_U$ 定义为 $U/V$ 各纤维上的 $r$-循环族
$$
\alpha|_U = (\alpha_s|_{U_s})_{s \in V}
$$
即限制到各概形论纤维所得的循环族。

\medskip\noindent
{\bf 平坦拉回。} 设 $X \to S$ 为概形的局部有限型态射，$r, e \geq 0$
为整数。设 $f : X' \to X$ 为相对维数 $e$ 的局部有限型平坦态射。
给定 $X/S$ 各纤维上的 $r$-循环族 $\alpha$，定义 $\alpha$ 的
{\it 平坦拉回} $f^*\alpha$ 为各纤维上的 $(r + e)$-循环族
$$
f^*\alpha = (f_s^*\alpha_s)_{s \in S}
$$
其中 $f_s^*\alpha_s \in Z_{r + e}(X'_s)$ 是沿概形论纤维间相对维数为
$e$ 的平坦态射 $f_s : X'_s \to X_s$，对 $Z_r(X_s)$ 中的循环
$\alpha_s$ 作平坦拉回所得。

\medskip\noindent
{\bf 固有前推。} 设
$$
\xymatrix{
X \ar[rr]_f \ar[rd] & & Y \ar[ld] \\
& S
}
$$
为概形态射的交换图，其中 $X$ 与 $Y$ 在 $S$ 上局部有限型，且 $f$ 固有。
设 $r \geq 0$ 为整数。给定 $X/S$ 各纤维上的 $r$-循环族 $\alpha$，
定义 $\alpha$ 的{\it 固有前推} $f_*\alpha$ 为 $Y/S$ 各纤维上的
$r$-循环族
$$
f_*\alpha = (f_{s, *}\alpha_s)_{s \in S}
$$
其中 $f_{s, *}\alpha_s \in Z_r(Y_s)$ 是沿概形论纤维间的固有态射
$f_s : X_s \to Y_s$，对 $Z_r(X_s)$ 中的循环 $\alpha_s$ 作固有前推所得。

\begin{lemma}
\label{relative-cycles-lemma-compatibilities}
上述运算满足如下相容性：
(1) 基变换具有函子性；
(2) 限制是基变换与平坦拉回（的一个特殊情形）的复合；
(3) 平坦拉回与基变换可交换；
(4) 平坦拉回具有函子性；
(5) 固有前推与基变换可交换；
(6) 固有前推具有函子性；
(7) 固有前推与平坦拉回可交换。
\end{lemma}

\begin{proof}
把 Chow 同调一章中证明的相应结果应用于所涉概形在 $S$ 上的各纤维，
即可直接得到这些相容性。我们略去精确陈述与详细证明，只列出若干出处。
第 (1) 点：Chow 同调，引理 \ref{chow-lemma-compose-base-change}。
第 (2) 点：显然。
第 (3) 点：Chow 同调，引理
\ref{chow-lemma-pullback-base-change-pullback}。
第 (4) 点：Chow 同调，引理 \ref{chow-lemma-compose-flat-pullback}。
第 (5) 点：Chow 同调，引理
\ref{chow-lemma-pullback-base-change-pushforward}。
第 (6) 点：Chow 同调，引理 \ref{chow-lemma-compose-pushforward}。
第 (7) 点：Chow 同调，引理
\ref{chow-lemma-flat-pullback-proper-pushforward}。
\end{proof}

\begin{example}
\label{relative-cycles-example-family-associated-module}
设 $f : X \to S$ 为概形的局部有限型态射，$r \geq 0$ 为整数。设
$\mathcal{F}$ 为有限型拟凝聚 $\mathcal{O}_X$-模。对 $s \in S$，以
$\mathcal{F}_s$ 表示 $\mathcal{F}$ 到 $X_s$ 的拉回。假设对所有
$s \in S$ 都有 $\dim(\text{Supp}(\mathcal{F}_s)) \leq r$。于是可将
$X/S$ 各纤维上的 $r$-循环族 $[\mathcal{F}/X/S]_r$ 与
$\mathcal{F}$ 相伴；其定义公式为
$$
[\mathcal{F}/X/S]_r = ([\mathcal{F}_s]_r)_{s \in S}
$$
其中 $[\mathcal{F}_s]_r$ 由 Chow 同调，定义
\ref{chow-definition-cycle-associated-to-coherent-sheaf} 给出。
\end{example}

\begin{lemma}
\label{relative-cycles-lemma-family-associated-module}
例 \ref{relative-cycles-example-family-associated-module} 中的构造与基变换、限制和平坦拉回
相容。
\end{lemma}

\begin{proof}
见 Chow 同调，引理 \ref{chow-lemma-pullback-coherent-base-change} 与
\ref{chow-lemma-pullback-coherent}。
\end{proof}

\begin{example}
\label{relative-cycles-example-family-associated-closed}
设 $f : X \to S$ 为概形的局部有限型态射，$r \geq 0$ 为整数。设
$Z \subset X$ 为闭子概形。对 $s \in S$，以 $Z_s$ 表示 $Z$ 在 $X_s$
中的逆像；等价地，它是 $Z$ 在 $s$ 处的概形论纤维，并视为 $X_s$ 的闭
子概形。假设对所有 $s \in S$ 都有 $\dim(Z_s) \leq r$。于是可将
$X/S$ 各纤维上的 $r$-循环族 $[Z/X/S]_r$ 与 $Z$ 相伴；其定义公式为
$$
[Z/X/S]_r = ([Z_s]_r)_{s \in S}
$$
其中 $[Z_s]_r$ 由 Chow 同调，定义
\ref{chow-definition-cycle-associated-to-closed-subscheme} 给出。
\end{example}

\begin{lemma}
\label{relative-cycles-lemma-family-associated-closed}
例 \ref{relative-cycles-example-family-associated-closed} 中的构造与基变换、限制和平坦拉回
相容。
\end{lemma}

\begin{proof}
取 $\mathcal{F} = (Z \to X)_*\mathcal{O}_Z$，这就是引理
\ref{relative-cycles-lemma-family-associated-module} 的特殊情形。见 Chow 同调，引理
\ref{chow-lemma-cycle-closed-coherent}。
\end{proof}

\begin{remark}[支撑]
\label{relative-cycles-remark-supports-family}
设 $f : X \to S$ 为概形的局部有限型态射，$r \geq 0$ 为整数。设
$\alpha$ 为 $X/S$ 各纤维上的 $r$-循环族。定义 $\alpha$ 的{\it 支撑}为
$$
\text{Supp}(\alpha) =
\bigcup\nolimits_{s \in S} \text{Supp}(\alpha_s) \subset X
$$
这里 $\text{Supp}(\alpha_s) \subset X_s$ 是循环 $\alpha_s$ 的支撑；
见 Chow 同调，定义 \ref{chow-definition-support-cycle}。支撑
$\text{Supp}(\alpha)$ 很少是 $X$ 的闭子集。
\end{remark}

\begin{lemma}
\label{relative-cycles-lemma-support-family}
按评注 \ref{relative-cycles-remark-supports-family} 取支撑与基变换、限制和平坦拉回相容。
\end{lemma}

\begin{proof}
略。
\end{proof}

\begin{lemma}
\label{relative-cycles-lemma-coequalizer-dim-r}
设 $f : X \to S$ 为概形的局部有限型态射，$r \geq 0$ 为整数。设
$g : S' \to S$ 为概形的满射。令 $S'' = S' \times_S S'$，并以
$f' : X' \to S'$ 和 $f'' : X'' \to S''$ 表示 $f$ 的基变换。设
$x \in X$ 满足 $\text{trdeg}_{\kappa(f(x))}(\kappa(x)) = r$。
\begin{enumerate}
\item 存在映到 $x$ 的 $x' \in X'$，使得
$\text{trdeg}_{\kappa(f'(x'))}(\kappa(x')) = r$。
\item 若 $x'_1, x'_2 \in X'$ 均满足 (1)，则存在 $x'' \in X''$，使得
$\text{trdeg}_{\kappa(f''(x''))}(\kappa(x'')) = r$ 且
$\text{pr}_i(x'') = x'_i$。
\end{enumerate}
\end{lemma}

\begin{proof}
(1) 即“态射”引理
\ref{morphisms-lemma-dimension-fibre-after-base-change}。设
$x'_1, x'_2$ 如 (2) 所述。由于 $X'' = X' \times_X X'$，存在
$x'' \in X''$ 同时映到 $x'_1$ 和 $x'_2$（例如见“下降”引理
\ref{descent-lemma-equiv-fibre-product}）。分别以 $s'' \in S''$、
$s'_i \in S'$ 和 $s \in S$ 表示 $x''$、$x'_i$ 和 $x$ 的像。
记 $k = \kappa(s)$，并令 $Z \subset X_k$ 为以 $x$ 为一般点的整闭子概形。
于是 $x'_i$ 是 $Z_{\kappa(s'_i)}$ 某个不可约分支的一般点。取包含
$x''$ 的不可约分支 $Z'' \subset Z_{\kappa(s'')}$，并以
$\xi'' \in Z''$ 表示其一般点。由于 $\xi'' \leadsto x''$，在两个投影下，
$\xi''$ 也必映到 $x'_i$。另一方面，由于它是 $Z$ 的基变换的一个
不可约分支的一般点，故
$\text{trdeg}_{\kappa(s'')}(\kappa(\xi'')) = r$。
\end{proof}

\begin{lemma}
\label{relative-cycles-lemma-descend-family}
设 $f : X \to S$ 为概形的局部有限型态射，$r \geq 0$ 为整数。设
$g : S' \to S$ 为概形态射，并令 $X' = S' \times_S X$。假设对每个
$s \in S$，都存在点 $s' \in S'$，满足 $g(s') = s$，并且域扩张
$\kappa(s')/\kappa(s)$ 可分。则：
\begin{enumerate}
\item 对 $X/S$ 各纤维上的两个 $r$-循环族 $\alpha_1$ 和 $\alpha_2$，若
$g^*\alpha_1 = g^*\alpha_2$，则 $\alpha_1 = \alpha_2$。
\item 给定 $X'/S'$ 各纤维上的 $r$-循环族 $\alpha'$，若作为
$(S' \times_S S') \times_S X / (S' \times_S S')$ 各纤维上的
$r$-循环族，有 $\text{pr}_1^*\alpha' = \text{pr}_2^*\alpha'$，则存在唯一的
$X/S$ 各纤维上的 $r$-循环族 $\alpha$，使得 $g^*\alpha = \alpha'$。
\end{enumerate}
\end{lemma}

\begin{proof}
(1) 由第 \ref{relative-cycles-section-relative-fields} 节所述基变换映射的单射性得出。
（只要 $S' \to S$ 满射，这一论证就成立。）

\medskip\noindent
设 $\alpha'$ 如 (2) 所述，并以
$\alpha'' = \text{pr}_1^*\alpha' = \text{pr}_2^*\alpha'$
表示其公共值。

\medskip\noindent
令 $(X/S)^{(r)}$ 为所有满足
$\text{trdeg}_{\kappa(f(x))}(\kappa(x)) = r$ 的 $x \in X$ 所成之集，
并类似定义 $(X'/S')^{(r)}$ 和 $(X''/S'')^{(r)}$。取系数后，可以把
$\alpha'$ 和 $\alpha''$ 看成函数
$\alpha' : (X'/S')^{(r)} \to \mathbf{Z}$ 和
$\alpha'' : (X''/S'')^{(r)} \to \mathbf{Z}$。给定函数
$$
\varphi : (X/S)^{(r)} \to \mathbf{Z}
$$
仿照基变换运算，定义
$g^*\varphi : (X'/S')^{(r)} \to \mathbf{Z}$。具体地，设
$x' \in (X'/S')^{(r)}$ 映到 $x \in X$、$s' \in S'$ 和 $s \in Z$。
分别以 $Z' \subset X'_{s'}$ 和 $Z \subset X_s$ 表示以 $x'$ 和 $x$ 为
一般点的整闭子概形。注意 $\dim(Z') = r$。若 $\dim(Z) < r$，则令
$(g^*\varphi)(x') = 0$。若 $\dim(Z) = r$，则 $Z'$ 是 $Z_{s'}$ 的一个
不可约分支，因而具有重数 $m_{Z', Z_{s'}}$；将其记作 $m(x', g)$。
定义
$$
(g^*\varphi)(x') = m(x', g) \varphi(x)
$$
注意，系数 $m(x', g)$ 总是正整数（例如见引理
\ref{relative-cycles-lemma-multiplicities-field-extension}）。类似地，有基变换映射
$$
\text{pr}_1^*, \text{pr}_2^* :
\text{Map}((X'/S')^{(r)}, \mathbf{Z})
\longrightarrow
\text{Map}((X''/S'')^{(r)}, \mathbf{Z})
$$
由基变换的结合性，
$\text{pr}_1^* \circ g^* = \text{pr}_2^* \circ g^*$（略去一个小细节）。
明确地说，就集合映射而言，该等式意味着对
$x'' \in (X''/S'')^{(r)}$ 有
$$
m(x'', \text{pr}_1) m(\text{pr}_1(x''), g) =
m(x'', \text{pr}_2) m(\text{pr}_2(x''), g)
$$
只要 $\text{pr}_1(x'')$ 和 $\text{pr}_2(x'')$ 属于
$(X''/S'')^{(r)}$。由引理 \ref{relative-cycles-lemma-coequalizer-dim-r}、上述等式和一个
初等论证\footnote{给定 $x \in (X/S)^{(r)}$，选取映到 $x$ 的
$x' \in (X'/S')^{(r)}$，并令
$\alpha(x) = \alpha'(x')/m(x', g)$。由该公式和引理，这一取值良定。}，
存在唯一映射
$$
\alpha : (X/S)^{(r)} \to \mathbf{Q}
$$
使得 $g^*\alpha = \alpha'$。为完成证明，只需说明 $\alpha$ 取整数值
（略去一个小细节：还须说明 $\alpha$ 在每条纤维上确定局部有限和；这由
$\alpha'$ 的相应性质得出）。给定像为 $s \in S$ 的任意
$x \in (X/S)^{(r)}$，可以选取点 $s' \in S'$，使得
$\kappa(s')/\kappa(s)$ 可分。于是可选取 $x' \in (X'/S')^{(r)}$，
使其映到 $s$ 和 $x$；此时 $Z_{s'}$ 约化，故 $m(x', g) = 1$。
因此 $\alpha(x) = \alpha'(x')$ 是整数。
\end{proof}

\begin{lemma}
\label{relative-cycles-lemma-pullback-universally-bijective}
设 $g : S' \to S$ 为诱导剩余域同构的概形双射态射。设
$f : X \to S$ 局部有限型，并令 $X' = S' \times_S X$。设 $r \geq 0$。
则沿 $g$ 的基变换在 $X/S$ 各纤维上的 $r$-循环族之群与
$X'/S'$ 各纤维上的 $r$-循环族之群之间确定一个双射。
\end{lemma}

\begin{proof}
略。
\end{proof}








\section{相对循环}
\label{relative-cycles-section-families-specialization}

\noindent
我们将使用如下定义；它与 \cite{SV} 中定义的比较见第
\ref{relative-cycles-section-compare} 节。

\begin{definition}
\label{relative-cycles-definition-relative-cycles}
设 $S$ 为局部 Noether 概形，$f : X \to S$ 为概形的局部有限型态射，
$r \geq 0$ 为整数。$X/S$ 上的一个{\it 相对 $r$-循环}，是指
$X/S$ 各纤维上的一个 $r$-循环族 $\alpha$，使得对每个态射
$g : S' \to S$，只要 $S'$ 是离散赋值环的谱，就有
$$
sp_{X'/S'}(\alpha_\eta) = \alpha_0
$$
其中 $sp_{X'/S'}$ 如第 \ref{relative-cycles-section-specialization} 节所定义，而
$\alpha_\eta$（相应地，$\alpha_0$）是 $\alpha$ 的基变换
$g^*\alpha$ 在 $S'$ 的一般点（相应地，闭点）处的值。$X/S$ 上所有
相对 $r$-循环所成的群记作 $z(X/S, r)$。
\end{definition}

\begin{lemma}
\label{relative-cycles-lemma-relative-cycle-functoriality}
设 $\alpha$ 是定义 \ref{relative-cycles-definition-relative-cycles} 所述的 $X/S$ 上相对
$r$-循环。则 $\alpha$ 的任意限制、基变换、平坦拉回或固有前推仍为相对
$r$-循环。
\end{lemma}

\begin{proof}
平坦拉回的情形使用引理 \ref{relative-cycles-lemma-specialization-flat-pullback}。
限制是平坦拉回的特殊情形。基变换的情形使用基变换的传递性。
固有前推的情形使用引理 \ref{relative-cycles-lemma-specialization-proper-pushforward}。
\end{proof}

\begin{lemma}
\label{relative-cycles-lemma-relative-cycles-h-descent}
设 $f : X \to S$ 为概形态射。假设 $S$ 局部 Noether，且 $f$ 局部有限型。
设 $r \geq 0$ 为整数，$\alpha$ 为 $X/S$ 各纤维上的 $r$-循环族。设
$\{g_i : S_i \to S\}$ 为一个 h-覆盖（“平坦性续篇”定义
\ref{flat-definition-h-covering}）。则 $\alpha$ 是相对 $r$-循环，当且仅当
每个基变换 $g_i^*\alpha$ 都是相对 $r$-循环。
\end{lemma}

\begin{proof}
若 $\alpha$ 是相对 $r$-循环，则由引理
\ref{relative-cycles-lemma-relative-cycle-functoriality}，每个基变换 $g_i^*\alpha$ 都是相对
$r$-循环。反之，假设每个 $g_i^*\alpha$ 都是相对 $r$-循环。设
$g : S' \to S$ 为态射，其中 $S'$ 是离散赋值环的谱。把 $S$ 替换为
$S'$、把 $X$ 替换为 $X' = X \times_S S'$、把 $\alpha$ 替换为
$\alpha' = g^*\alpha$，并利用 h-覆盖的基变换仍为 h-覆盖（“平坦性续篇”
引理 \ref{flat-lemma-h}），便归约到下一段所研究的问题。

\medskip\noindent
假设 $S$ 是离散赋值环的谱，其闭点为 $0$、一般点为 $\eta$。需要证明
$sp_{X/S}(\alpha_\eta) = \alpha_0$。由于 h-覆盖按定义是 V-覆盖，存在
$i$ 以及 $S_i$ 中点的特化 $s' \leadsto s$，满足
$g_i(s') = \eta$ 和 $g_i(s) = 0$；见“拓扑”引理
\ref{topologies-lemma-refine-qcqs-V}。由“性质”引理
\ref{properties-lemma-locally-Noetherian-specialization-dvr}，存在从离散
赋值环的谱 $S'$ 出发的态射 $h : S' \to S_i$，它把一般点 $\eta'$ 映到
$s'$，把闭点 $0'$ 映到 $s$。记 $\alpha' = h^*g_i^*\alpha$。由假设，
$sp_{X'/S'}(\alpha'_{\eta'}) = \alpha'_{0'}$。由于
$g = g_i \circ h : S' \to S$ 是由离散赋值环扩张诱导的概形态射，
$sp_{X/S}$ 与 $sp_{X'/S'}$ 和纤维上的基变换映射相容；见引理
\ref{relative-cycles-lemma-specialization-extension}。第 \ref{relative-cycles-section-relative-fields} 节
已经说明基变换映射 $Z_r(X_0) \to Z_r(X'_{0'})$ 是单射，故
$sp_{X/S}(\alpha_\eta) = \alpha_0$。
\end{proof}

\begin{lemma}
\label{relative-cycles-lemma-families-specialization-fppf-descent}
设 $f : X \to S$ 为概形态射。假设 $S$ 局部 Noether，且 $f$ 局部有限型。
设 $r, e \geq 0$ 为整数，$\alpha$ 为 $X/S$ 各纤维上的 $r$-循环族。
设 $\{f_i : X_i \to X\}$ 为一族联合满射的局部有限型平坦态射，且其相对
维数为 $e$。则 $\alpha$ 是相对 $r$-循环，当且仅当每个平坦拉回
$f_i^*\alpha$ 都是相对 $r$-循环。
\end{lemma}

\begin{proof}
若 $\alpha$ 是相对 $r$-循环，则由引理
\ref{relative-cycles-lemma-relative-cycle-functoriality}，每个拉回 $f_i^*\alpha$ 都是相对
$r$-循环。反之，假设每个 $f_i^*\alpha$ 都是相对 $r$-循环。设
$g : S' \to S$ 为态射，其中 $S'$ 是离散赋值环的谱。把 $S$ 替换为
$S'$、把 $X$ 替换为 $X' = X \times_S S'$、把 $\alpha$ 替换为
$\alpha' = g^*\alpha$，便归约到下一段所研究的问题。

\medskip\noindent
假设 $S$ 是离散赋值环的谱，其闭点为 $0$、一般点为 $\eta$。需要证明
$sp_{X/S}(\alpha_\eta) = \alpha_0$。以
$f_{i, 0} : X_{i, 0} \to X_0$ 表示 $f_i$ 到 $S$ 的闭点上的基变换；
$f_{i, \eta}$ 的记号类似。注意
$$
f_{i, 0}^*sp_{X/S}(\alpha_\eta) =
sp_{X_i/S}(f_{i, \eta}^*\alpha_\eta) = f_{i, 0}^*\alpha_0
$$
其中第一个等式由引理 \ref{relative-cycles-lemma-specialization-flat-pullback} 得出，
第二个等式由假设得出。由于映射族
$f_{i, 0}^* : Z_r(X_0) \to Z_r(X_{i, 0})$ 联合单射（这是因为
$f_{i, 0}$ 联合满射），结论随即成立。
\end{proof}

\begin{lemma}
\label{relative-cycles-lemma-check-after-closed}
设 $S$ 为局部 Noether 概形，$i : X \to Y$ 为 $S$ 上局部有限型概形之间的
闭浸入。设 $r \geq 0$，并令 $\alpha$ 为 $X/S$ 各纤维上的 $r$-循环族。
则 $\alpha$ 是 $X/S$ 上的相对 $r$-循环，当且仅当 $i_*\alpha$ 是
$Y/S$ 上的相对 $r$-循环。
\end{lemma}

\begin{proof}
由于基变换与 $i_*$ 可交换（引理 \ref{relative-cycles-lemma-compatibilities}），只需证明：
若 $S$ 是离散赋值环的谱，其一般点为 $\eta$、闭点为 $0$，则
$sp_{X/S}(\alpha_\eta) = \alpha_0$ 当且仅当
$sp_{Y/S}(i_{\eta, *}\alpha_\eta) = i_{0, *}\alpha_0$。这是因为
$i_{0, *} : Z_r(X_0) \to Z_r(Y_0)$ 是单射，而且由引理
\ref{relative-cycles-lemma-specialization-proper-pushforward}，有
$i_{0, *}sp_{X/S}(\alpha_\eta) =
sp_{Y/S}(i_{\eta, *}\alpha_\eta)$。
\end{proof}

\noindent
下面的引理将在引理 \ref{relative-cycles-lemma-relative-cycles-equal} 中得到加强。

\begin{lemma}
\label{relative-cycles-lemma-uniqueness-extension}
设 $f : X \to S$ 为概形态射。假设 $S$ 局部 Noether，且 $f$ 局部有限型。
设 $r \geq 0$，并令 $\alpha$ 和 $\beta$ 为 $X/S$ 上的相对 $r$-循环。
下列条件等价：
\begin{enumerate}
\item $\alpha = \beta$；
\item 对 $S$ 的任一不可约分支的一般点 $\eta \in S$，都有
$\alpha_\eta = \beta_\eta$。
\end{enumerate}
\end{lemma}

\begin{proof}
(1) $\Rightarrow$ (2) 显然。假设 (2)。对每个 $s \in S$，都可以找到
(2) 所述的一个 $\eta$，使其特化到 $s$。由“性质”引理
\ref{properties-lemma-locally-Noetherian-specialization-dvr}，存在从离散
赋值环的谱 $S'$ 出发的态射 $g : S' \to S$，它把一般点 $\eta'$ 映到
$\eta$，把闭点 $0$ 映到 $s$。于是 $\alpha_s$ 和 $\beta_s$ 是
$Z_r(X_s)$ 中的元素，并基变换为 $Z_r(X_{0'})$ 中同一个元素，即
$sp_{X_{S'}/S'}(\alpha_{\eta'})$，其中 $\alpha_{\eta'}$ 是
$\alpha_\eta$ 的基变换。第 \ref{relative-cycles-section-relative-fields} 节已经说明
基变换映射 $Z_r(X_s) \to Z_r(X_{0'})$ 是单射，故
$\alpha_s = \beta_s$。
\end{proof}

\begin{lemma}
\label{relative-cycles-lemma-family-associated-module-specialization}
在例 \ref{relative-cycles-example-family-associated-module} 的情形中，假设 $S$ 局部
Noether，且 $\mathcal{F}$ 在维数 $\geq r$ 处在 $S$ 上平坦
（“平坦性续篇”定义 \ref{flat-definition-flat-dimension-n}）。则
$[\mathcal{F}/X/S]_r$ 是 $X/S$ 上的相对 $r$-循环。
\end{lemma}

\begin{proof}
由“平坦性续篇”引理 \ref{flat-lemma-pre-flat-dimension-n}，关于
$\mathcal{F}$ 的假设在任意基变换下保持成立。此外，由引理
\ref{relative-cycles-lemma-family-associated-module}，$[\mathcal{F}/X/S]_r$ 的构造与任意
基变换相容。由于与特化相容这一条件是在基变换到离散赋值环的谱之后
检验的，问题归约到 $S$ 是赋值环的谱的情形。此时集合
$U = \{x \in X \mid \mathcal{F}\text{ 在点 }x\text{ 上平坦，底为 }S\}$
由“平坦性续篇”引理 \ref{flat-lemma-finite-type-flat} 是 $X$ 中的开集。
根据假设，$X$ 中 $U$ 的补集在 $S$ 上各纤维的维数 $< r$，故沿浸入
$U \subset X$ 的限制在任意基变换后都诱导纤维上 $r$-循环群的同构，
并与特化映射及由 $\mathcal{F}$ 所伴随相对循环的构造相容。因此，只需证明
$[\mathcal{F}|_U / U /S]_r$ 与特化相容。由于 $\mathcal{F}|_U$ 在 $S$
上平坦，这由引理 \ref{relative-cycles-lemma-specialization-module} 和定义得出。
\end{proof}

\begin{lemma}
\label{relative-cycles-lemma-family-associated-closed-specialization}
在例 \ref{relative-cycles-example-family-associated-closed} 的情形中，假设 $S$ 局部
Noether，且 $Z$ 在维数 $\geq r$ 处在 $S$ 上平坦。则
$[Z/X/S]_r$ 是 $X/S$ 上的相对 $r$-循环。
\end{lemma}

\begin{proof}
该假设意味着 $\mathcal{O}_Z$ 在维数 $\geq r$ 处在 $S$ 上平坦。于是对
$\mathcal{F} = (Z \to X)_*\mathcal{O}_Z$ 应用引理
\ref{relative-cycles-lemma-family-associated-module-specialization} 即得结论。
\end{proof}

\noindent
设 $S$ 为局部 Noether 概形，$f : X \to S$ 为有限型态射，并设
$r \geq 0$。以 $Hilb(X/S, r)$ 表示所有闭子概形 $Z \subset X$ 所成的
集合，其中 $Z \to S$ 平坦且相对维数 $\leq r$。由引理
\ref{relative-cycles-lemma-family-associated-closed-specialization}，对每个
$Z \in Hilb(X/S, r)$，都有元素 $[Z/X/S]_r \in z(X/S, r)$。因此得到群同态
\begin{equation}
\label{relative-cycles-equation-cycle-classes}
\text{由下列集合生成的自由 Abel 群： }Hilb(X/S, r) \longrightarrow z(X/S, r)
\end{equation}
它把 $\sum n_i[Z_i]$ 映到 $\sum n_i[Z_i/X/S]_r$。相对 $r$-循环的一个
关键特征是：在适当拓扑下，它们局部地（在 $X$ 和 $S$ 上）属于该映射的像。

\begin{lemma}
\label{relative-cycles-lemma-get-cycles}
设 $f : X \to S$ 为概形的有限型态射，且 $S$ Noether。设 $r \geq 0$，
$\alpha$ 为 $X/S$ 上的相对 $r$-循环。则存在固有完全分解态射
（“态射续篇”定义 \ref{more-morphisms-definition-cd-morphism}）
$g : S' \to S$，使得 $g^*\alpha$ 属于 (\ref{relative-cycles-equation-cycle-classes}) 的像。
\end{lemma}

\begin{proof}
由 Noether 归纳法，可以假设：对任意不是同构的闭浸入
$g : S' \to S$，结论对 $\alpha$ 的拉回成立。

\medskip\noindent
令 $S_1 \subset S$ 为一个不可约分支（视作整闭子概形）。令
$S_2 \subset S$ 为 $S'$ 的补集之闭包（视作约化闭子概形）。若
$S_2 \not = \emptyset$，则结论对沿 $S_1 \to S$ 和 $S_2 \to S$ 的
$\alpha$ 拉回成立。若 $g_1 : S'_1 \to S_1$ 和 $g_2 : S'_2 \to S_2$
是相应的固有完全分解态射，则
$S' = S'_1 \amalg S'_2 \to S$ 是固有完全分解态射，且结论对 $S$
成立\footnote{确实，$S'_1 \times_S X$ 中任一在 $S'_1$ 上平坦且相对维数
$\leq r$ 的闭子概形，都可以视作 $S' \times_S X$ 中一个在 $S'$ 上平坦
且相对维数 $\leq r$ 的闭子概形。}。因此可以假设 $S' \to S$ 是双射，
从而归约到下一段所述情形。

\medskip\noindent
假设 $S$ 整。令 $\eta \in S$ 为一般点，并令 $K = \kappa(\eta)$ 为
$S$ 的函数域。于是 $\alpha_\eta$ 是 $X_K$ 上的 $r$-循环。写成
$\alpha_\eta = \sum n_i[Y_i]$。取 $Y_i$ 的闭包，得到整闭子概形
$Z_i \subset X$，其到 $\eta$ 的基变换为 $Y_i$。由一般平坦性（例如
“态射”命题 \ref{morphisms-proposition-generic-flatness}），对每个 $i$，
$Z_i$ 在 $S$ 的某个非空开集 $U$ 上平坦。应用“平坦性续篇”引理
\ref{flat-lemma-flat-after-blowing-up}，可以找到 $U$-容许的吹起
$g : S' \to S$，使得 $Z_i$ 的严格变换 $Z'_i \subset X_{S'}$ 在 $S'$
上平坦。于是 $\beta = \sum n_i[Z'_i/X_{S'}/S']_r$ 属于
(\ref{relative-cycles-equation-cycle-classes}) 的像；由引理
\ref{relative-cycles-lemma-uniqueness-extension}，$\beta = g^*\alpha$。

\medskip\noindent
但这尚未完成证明，因为 $S' \to S$ 未必完全分解。修正这一点很容易：
以 $T \subset S$ 表示 $U$ 的补集（视作闭子概形）。由 Noether 归纳法，
可以找到固有完全分解态射 $T' \to T$，使得
$(T' \to S)^*\alpha$ 属于 (\ref{relative-cycles-equation-cycle-classes}) 的像。于是
$S' \amalg T' \to S$ 满足要求。
\end{proof}

\begin{lemma}
\label{relative-cycles-lemma-get-cycles-dvr}
设 $f : X \to S$ 为概形的有限型态射，其中 $S$ 是离散赋值环的谱。
设 $r \geq 0$。则 (\ref{relative-cycles-equation-cycle-classes}) 是满射。
\end{lemma}

\begin{proof}
这当然由引理 \ref{relative-cycles-lemma-get-cycles} 得出，也可以如下直接看出。设
$\alpha$ 是 $X/S$ 上的相对 $r$-循环。写
$\alpha_\eta = \sum n_i[Z_i]$（该和有限）。如第
\ref{relative-cycles-section-specialization} 节，以 $\overline{Z}_i \subset X$ 表示
$Z_i$ 的闭包。于是 $\alpha = \sum n_i[\overline{Z}_i/X/S]$。
\end{proof}

\begin{lemma}
\label{relative-cycles-lemma-relative-cycle-smooth}
设 $f : X \to S$ 为概形态射，$r \geq 0$。假设 $S$ 局部 Noether，且
$f$ 光滑、相对维数为 $r$。设 $\alpha \in z(X/S, r)$。则 $\alpha$ 的支撑
在 $X$ 中既开又闭（更精确的结论见证明）。
\end{lemma}

\begin{proof}
设 $x \in X$ 的像为 $s \in S$。由于 $f$ 光滑，$X_s$ 中存在唯一一个包含
$x$ 的不可约分支 $Z(x)$。于是 $\dim(Z(x)) = r$。以 $n_x$ 表示循环
$\alpha_s$ 中 $Z(x)$ 的系数。我们将证明函数 $x \mapsto n_x$ 在 $X$ 上
局部常值。

\medskip\noindent
设 $g : S' \to S$ 为局部 Noether 概形之间的态射。以 $X'$ 表示 $X$ 的
基变换，以 $\alpha' = g^*\alpha$ 表示 $\alpha$ 的基变换。设
$x' \in X'$ 映到 $s' \in S'$、$x \in X$ 和 $s \in S$。我们断言
$n_{x'} = n_x$。事实上，由于 $Z(x)$ 在 $\kappa(s)$ 上光滑，概形
$Z(x) \times_{\Spec(\kappa(s))} \Spec(\kappa(s'))$
是约化的。由于 $Z(x')$ 是该概形的一个不可约分支，根据第
\ref{relative-cycles-section-relative-fields} 节的基变换定义，$\alpha'_{s'}$ 中
$Z(x')$ 的系数 $n_{x'}$ 等于 $\alpha_s$ 中 $Z(x)$ 的系数 $n_x$，
从而证明了该断言。

\medskip\noindent
由于 $X$ 局部 Noether，为证明 $x \mapsto n_x$ 局部常值，只需证明：
若 $x' \leadsto x$ 是 $X$ 中的一个特化，则 $n_{x'} = n_x$。选取态射
$S' \to X$，其中 $S'$ 是离散赋值环的谱，它把一般点 $\eta$ 映到 $x'$，
把闭点 $0$ 映到 $x$；见“性质”引理
\ref{properties-lemma-locally-Noetherian-specialization-dvr}。于是 $f$ 沿
$S' \to S$ 的基变换 $X' \to S'$ 有一个截面 $\sigma : S' \to X'$，使得
$\sigma(\eta) \leadsto \sigma(0)$ 是 $X'$ 中点的特化，并映到 $X$ 中的
$x' \leadsto x$。因此归约到下一段的断言。

\medskip\noindent
设 $S$ 是离散赋值环的谱，其一般点为 $\eta$、闭点为 $0$，并且给定截面
$\sigma : S  \to X$。断言：$n_{\sigma(\eta)} = n_{\sigma(0)}$。由
“态射续篇”第 \ref{more-morphisms-section-connected-components} 节中的
讨论，特别是引理
\ref{more-morphisms-lemma-connected-along-section-open}，把 $X$ 替换为一个
开子概形后，可以假设 $X \to S$ 的纤维连通。由于这些纤维光滑，它们
不可约。于是 $\alpha_\eta = n[X_\eta]$，其中 $n = n_{\sigma(\eta)}$；关系
$sp_{X/S}(\alpha_\eta) = \alpha_0$ 蕴含 $\alpha_0 = n[X_0]$，即所需的
$n_{\sigma(0)} = n$。
\end{proof}

\begin{lemma}
\label{relative-cycles-lemma-relative-cycles-equal}
设 $f : X \to S$ 为概形态射。假设 $S$ 局部 Noether，且 $f$ 局部有限型。
设 $r \geq 0$ 且 $\alpha, \beta \in z(X/S, r)$。则集合
$E = \{s \in S : \alpha_s = \beta_s\}$ 在 $S$ 中闭。
\end{lemma}

\begin{proof}
问题在 $S$ 上是局部的，故可假设 $S$ 仿射。取仿射开覆盖
$X = \bigcup U_i$，并令
$E_i = \{s \in S : \alpha_s|_{U_{i, s}} = \beta_s|_{U_{i, s}}\}$。
于是 $E = \bigcap E_i$。因此只需对 $U_i \to S$ 以及 $\alpha$、$\beta$
到 $U_i$ 的限制证明该引理。这就归约到下一段所述情形。

\medskip\noindent
假设 $X$ 和 $S$ 拟紧。令 $\gamma = \alpha - \beta$，则
$E = \{s \in S : \gamma_s = 0\}$。由引理
\ref{relative-cycles-lemma-family-associated-closed-specialization}，存在有限族联合满射的
固有态射 $\{g_i : S_i \to S\}$，使得 $g_i^*\gamma$ 属于
(\ref{relative-cycles-equation-cycle-classes}) 的像。注意，$E_i = g_i^{-1}(E)$ 是所有满足
$(g_i^*\gamma)_t = 0$ 的点 $t \in S_i$ 所成之集。若对所有 $i$，$E_i$ 都闭，
则 $E = \bigcup g_i(E_i)$ 也闭。这就归约到下一段所述情形。

\medskip\noindent
假设 $X$ 和 $S$ 拟紧，且
$\gamma = \sum n_i[Z_i/X/S]_r$，其中只有有限个闭子概形 $Z_i \subset X$，
它们在 $S$ 上平坦且相对维数 $\leq r$。令
$X' = \bigcup Z_i$（概形论并）。于是 $i : X' \to X$ 是闭浸入，且
$X'$ 在 $S$ 上的相对维数 $\leq r$。此外，$\gamma = i_*\gamma'$，其中
$\gamma' = \sum n_i[Z_i/X'/S]_r$。显然
$E = E' = \{s \in S : \gamma'_s = 0\}$，故归约到下一段所述情形。

\medskip\noindent
假设 $X$ 在 $S$ 上的相对维数 $\leq r$。取 $s \in S$ 且
$s \not \in E$。我们将证明存在 $s$ 的开邻域 $V \subset S$，使得
$E \cap V$ 为空。假设 $s \not \in E$ 意味着存在维数为 $r$ 的整闭子概形
$Z \subset X_s$，使得 $\gamma_s$ 中 $[Z]$ 的系数 $n$ 非零。令
$x \in Z$ 为一般点。由于 $\dim(Z) = r$，$x$ 是 $X_s$ 的一个不可约分支
（即 $Z$）的一般点。因此，把 $X$ 替换为 $x$ 的一个开邻域后，可以假设
$Z$ 是 $X_s$ 的唯一不可约分支。特别地，$\gamma_s = n[Z]$。

\medskip\noindent
现在应用“态射续篇”引理
\ref{more-morphisms-lemma-local-structure-finite-type}，得到图
$$
\xymatrix{
X \ar[dd] & X' \ar[l]^g \ar[d]^\pi & x \ar@{|->}[dd] &
x' \ar@{|->}[l]  \ar@{|->}[d] \\
& Y \ar[d]^h & & y \ar@{|->}[d] \\
S \ar@{=}[r] & S & s & s \ar@{=}[l]
}
$$
并具有该引理列出的全部性质。令 $\gamma' = g^*\gamma$ 为平坦拉回。
注意 $E \subset E' = \{s \in S: \gamma'_s = 0\}$，并且
$s \not \in E'$，因为 $\gamma'_s$ 中 $Z'$ 的系数非零；这里
$Z' \subset X'_s$ 是 $x'$ 的闭包。类似地，令
$\gamma'' = \pi_*\gamma'$。于是
$E' \subset E'' = \{s \in S: \gamma''_s = 0\}$，且
$s \not \in E''$，因为 $\gamma''_s$ 中 $Z''$ 的系数非零；这里
$Z'' \subset Y_s$ 是 $y$ 的闭包。由引理
\ref{relative-cycles-lemma-relative-cycle-smooth} 和 $Y \to S$ 的开性，$s$ 的某个开邻域
与 $E''$ 不交，证明完成。
\end{proof}

\begin{lemma}
\label{relative-cycles-lemma-descend-through-limit}
设 $S = \lim_{i \in I} S_i$ 是一个 Noether 概形有向逆系统的极限，且转移
态射均为仿射态射。设 $0 \in I$，并令 $X_0 \to S_0$ 为概形的有限型态射。
对 $i \geq 0$，令 $X_i = S_i \times_{S_0} X_0$，并令
$X = S \times_{S_0} X_0$。若 $S$ 也 Noether，则
$$
z(X/S, r) = \colim_{i \geq 0} z(X_i/S_i, r)
$$
其中转移映射由相对 $r$-循环的基变换给出。
\end{lemma}

\begin{proof}
假设 $i \geq 0$ 且 $\alpha_i, \beta_i \in z(X_i/S_i, r)$ 映到
$z(X/S, r)$ 中同一个元素。则 $S \to S_i$ 的像包含在引理
\ref{relative-cycles-lemma-relative-cycles-equal} 的闭子集 $E \subset S_i$ 中。因此，对某个
$j \geq i$，态射 $S_j \to S_i$ 的像包含在 $E$ 中；见“极限”引理
\ref{limits-lemma-limit-contained-in-constructible}。所以 $\alpha_i$ 和
$\beta_i$ 到 $S_j$ 的基变换相等。故该映射为单射。

\medskip\noindent
设 $\alpha \in z(X/S, r)$。应用引理 \ref{relative-cycles-lemma-get-cycles}，得到固有完全
分解态射 $g : S' \to S$，使得 $g^*\alpha$ 属于
(\ref{relative-cycles-equation-cycle-classes}) 的像。令 $X' = S' \times_S X$。可以写成
$g^*\alpha = \sum n_a [Z_a/X'/S']_r$，其中 $Z_a \subset X'$ 是在 $S'$
上平坦且相对维数 $\leq r$ 的闭子概形。

\medskip\noindent
现在运用“极限”第 \ref{limits-section-descending-relative} 节的工具。
可以找到 $i \geq 0$，使得存在：
\begin{enumerate}
\item 固有完全分解态射 $g_i : S'_i \to S_i$，它到 $S$ 的基变换为
$g : S' \to S$；
\item 令 $X'_i = S'_i \times_{S_i} X_i$，则有闭子概形
$Z_{ai} \subset X'_i$，它们在 $S'_i$ 上平坦、相对维数 $\leq r$，且到
$S'$ 的基变换为 $Z_a$。
\end{enumerate}
这里使用了“极限”引理
\ref{limits-lemma-descend-finite-presentation},
\ref{limits-lemma-descend-closed-immersion-finite-presentation},
\ref{limits-lemma-descend-flat-finite-presentation},
\ref{limits-lemma-eventually-proper} 和
\ref{limits-lemma-limit-dimension}
以及“态射续篇”引理 \ref{more-morphisms-lemma-descend-cd}。考虑
$\alpha'_i = \sum n_a [Z_{ai}/X'_i/S'_i]_r \in z(X'_i/S'_i, r)$.
依构造，$\alpha'_i$ 在 $z(X'/S', r)$ 中的像等于基变换 $g^*\alpha$。

\medskip\noindent
令 $S''_i = S'_i \times_{S_i} S'_i$、
$X''_i = S''_i \times_{S_i} X_i$，并令 $S'' = S' \times_S S'$、
$X'' = S'' \times_S X$。以 $\text{pr}_1, \text{pr}_2 : S'' \to S'$ 和
$\text{pr}_1, \text{pr}_2 : S''_i \to S'_i$ 表示投影。两个基变换
$\text{pr}_1^*\alpha'_i$ 和 $\text{pr}_1^*\alpha'_i$ 映到
$z(X''/S'', r)$ 中同一个元素，因为
$\text{pr}_1^*g^*\alpha = \text{pr}_1^*g^*\alpha$。因此增大 $i$ 后，
由证明的第一段，可以假设
$\text{pr}_1^*\alpha'_i = \text{pr}_1^*\alpha'_i$。由引理
\ref{relative-cycles-lemma-descend-family}，得到 $X_i/S_i$ 各纤维上唯一的 $r$-循环族
$\alpha_i$，使得 $g_i^*\alpha_i = \alpha'_i$（这里用到
$S'_i \to S_i$ 完全分解）。由引理
\ref{relative-cycles-lemma-relative-cycles-h-descent}，$\alpha_i \in z(X_i/S_i, r)$。
引理 \ref{relative-cycles-lemma-descend-family} 中的唯一性说明 $\alpha_i$ 在
$z(X/S, r)$ 中的像为 $\alpha$，证明完成。
\end{proof}

\begin{lemma}
\label{relative-cycles-lemma-thickening}
设 $S$ 为局部 Noether 概形，$i : X \to X'$ 为 $S$ 上局部有限型概形之间的
加厚。设 $r \geq 0$。则 $i_* : z(X/S, r) \to z(X'/S, r)$ 是双射。
\end{lemma}

\begin{proof}
由于 $i_s : X_s \to X'_s$ 是加厚，显然 $i_*$ 在 $X/S$ 各纤维上的
$r$-循环族与 $X'/S$ 各纤维上的 $r$-循环族之间诱导双射。此外，对
$X/S$ 各纤维上的 $r$-循环族 $\alpha$，由引理
\ref{relative-cycles-lemma-check-after-closed}，
$\alpha \in z(X/S, r) \Leftrightarrow i_*\alpha \in z(X'/S, r)$。
引理得证。
\end{proof}

\begin{lemma}
\label{relative-cycles-lemma-extend-to-larger}
设 $S$ 为局部 Noether 概形，$X$ 为 $S$ 上局部有限型概形。设
$r \geq 0$。设 $U \subset X$ 为开集，使得 $X \setminus U$ 在 $S$ 上的
相对维数 $< r$；换言之，对所有 $s \in S$，都有
$\dim(X_s \setminus U_s) < r$。则限制确定双射
$z(X/S, r) \to z(U/S, r)$。
\end{lemma}

\begin{proof}
由维数假设，$Z_r(X_s) \to Z_r(U_s)$ 是双射，所以限制在 $X/S$ 各纤维上的
$r$-循环族与 $U/S$ 各纤维上的 $r$-循环族之间诱导双射。这些限制映射
$Z_r(X_s) \to Z_r(U_s)$ 与基变换和特化相容；见引理
\ref{relative-cycles-lemma-compatibilities} 和 \ref{relative-cycles-lemma-specialization-flat-pullback}。
引理由此容易得出；细节略。
\end{proof}

\begin{lemma}
\label{relative-cycles-lemma-seminormalize-base}
设 $g : S' \to S$ 为局部 Noether 概形之间诱导剩余域同构的泛同胚。设
$f : X \to S$ 局部有限型，并令 $X' = S' \times_S X$。设 $r \geq 0$。
则沿 $g$ 的基变换确定双射 $z(X/S, r) \to z(X'/S', r)$。
\end{lemma}

\begin{proof}
由引理 \ref{relative-cycles-lemma-pullback-universally-bijective}，$X/S$ 各纤维上的
$r$-循环族之群与 $X'/S'$ 各纤维上的 $r$-循环族之群之间存在双射。设
$\alpha$ 为 $X/S$ 各纤维上的一个 $r$-循环族，$\alpha' = g^*\alpha$ 为其
基变换。若 $R$ 是离散赋值环，则任意态射 $h : \Spec(R) \to S$ 都唯一地
分解为 $g \circ h'$，其中 $h' : \Spec(R) \to S'$。事实上，态射
$S' \times_S \Spec(R) \to \Spec(R)$ 是在剩余域上诱导双射的泛同胚，
因而有截面（例如因为 $R$ 是半正规环；见“态射”第
\ref{morphisms-section-seminormalization} 节）。因此，$\alpha$ 与特化相容
（即为相对 $r$-循环）这一条件，等价于 $\alpha'$ 与特化相容。
\end{proof}






\section{等维相对循环}
\label{relative-cycles-section-equidimensional}

\noindent
定义如下。

\begin{definition}
\label{relative-cycles-definition-equidimensional}
设 $f : X \to S$ 为概形态射。假设 $S$ 局部 Noether，且 $f$ 局部有限型。
设 $r \geq 0$ 为整数。如果 $X/S$ 上相对 $r$-循环 $\alpha$ 的支撑
（评注 \ref{relative-cycles-remark-supports-family}）包含在闭子集 $W \subset X$ 中，且后者
在 $S$ 上的相对维数 $\leq r$，则称 $\alpha$ {\it 等维}。
$X/S$ 上所有等维相对 $r$-循环所成的群记作 $z_{equi}(X/S, r)$。
\end{definition}

\begin{example}
\label{relative-cycles-example-not-equidimensional}
\begin{reference}
\cite[Example 3.1.9]{SV}
\end{reference}
存在并非等维的相对 $r$-循环。具体地，设 $k$ 为域，并令
$X = \Spec(k[x, y, t])$ 位于 $S = \Spec(k[x, y])$ 之上。设 $s$ 为 $S$
的一点，并以 $a, b \in \kappa(s)$ 表示 $x$ 和 $y$ 的像。考虑如下定义的
$X/S$ 各纤维上的 $0$-循环族 $\alpha$：
\begin{enumerate}
\item 若 $b = 0$，则 $\alpha_s = 0$；否则
\item $\alpha_s = [p] - [q]$，其中 $p$（相应地，$q$）是
$\Spec(\kappa(s)[t])$ 的 $\kappa(s)$-有理点，满足 $t = a/b$
（相应地，$t = (a + b^2)/b$）。
\end{enumerate}
留给读者证明它与特化相容；思路是，对 $\kappa(s)$ 上任意满足
$v(b) > 0$ 的赋值 $v$，$a/b$ 与 $(a + b^2)/b = a/b + b$ 在其剩余域上的
$\mathbf{P}^1$ 中趋于同一点。另一方面，$\alpha$ 的支撑之闭包包含
$(0, 0)$ 上的整条纤维。
\end{example}

\begin{lemma}
\label{relative-cycles-lemma-equidimensional-functoriality}
设 $f : X \to S$ 为概形态射。假设 $S$ 局部 Noether，且 $f$ 局部有限型。
设 $r \geq 0$ 为整数，并令 $\alpha$ 为 $X/S$ 上的相对 $r$-循环。若
$\alpha$ 等维，则 $\alpha$ 的任意限制、基变换或平坦拉回都等维。
\end{lemma}

\begin{proof}
略。
\end{proof}

\begin{lemma}
\label{relative-cycles-lemma-check-equidimensional}
设 $f : X \to S$ 为概形态射。假设 $S$ 局部 Noether，且 $f$ 局部有限型。
设 $r \geq 0$ 为整数，并令 $\alpha$ 为 $X/S$ 上的相对 $r$-循环。则检验
$\alpha$ 等维时，可以在 $X$ 和 $S$ 上 Zariski 局部地进行。
\end{lemma}

\begin{proof}
确实，$\alpha$ 等维这一条件恰好意味着 $\alpha$ 的支撑之闭包在 $S$ 上的
相对维数 $\leq r$。由于取闭包与限制到开集可交换，引理得证
（略去一个小细节）。
\end{proof}

\begin{lemma}
\label{relative-cycles-lemma-equidimensional-fppf-descent}
设 $f : X \to S$ 为概形态射。假设 $S$ 局部 Noether，且 $f$ 局部有限型。
设 $r \geq 0$ 为整数，并令 $\alpha$ 为 $X/S$ 上的相对 $r$-循环。设
$\{g_i : S_i \to S\}$ 为 fppf-覆盖。则 $\alpha$ 等维，当且仅当每个
基变换 $g_i^*\alpha$ 都等维。
\end{lemma}

\begin{proof}
若 $\alpha$ 等维，则由引理 \ref{relative-cycles-lemma-equidimensional-functoriality}，
每个 $g_i^*\alpha$ 也等维。反之，假设每个 $g_i^*\alpha$ 都等维。以 $W$
表示 $\text{Supp}(\alpha)$ 在 $X$ 中的闭包。由于
$g_i : S_i \to S$ 泛开（因为它平坦且局部有限表示），态射
$f_i : X_i = S_i \times_S X \to X$ 也泛开。记 $\alpha_i = g_i^*\alpha$。
由引理 \ref{relative-cycles-lemma-support-family}，
$\text{Supp}(\alpha_i) = f_i^{-1}(\text{Supp}(\alpha))$。由于 $f_i$ 开，
$W_i = f_i^{-1}(W)$ 是 $\text{Supp}(\alpha_i)$ 的闭包。因此根据假设，
态射 $W_i \to S_i$ 的相对维数 $\leq r$。由“态射”引理
\ref{morphisms-lemma-dimension-fibre-after-base-change}（以及态射
$S_i \to S$ 联合满射这一事实），可知 $W \to S$ 的相对维数 $\leq r$。
\end{proof}

\begin{lemma}
\label{relative-cycles-lemma-equidimensional-descent-pullbacks}
设 $f : X \to S$ 为概形态射。假设 $S$ 局部 Noether，且 $f$ 局部有限型。
设 $r, e \geq 0$ 为整数，并令 $\alpha$ 为 $X/S$ 上的相对 $r$-循环。
设 $\{f_i : X_i \to X\}$ 为一族联合满射的局部有限型平坦态射，且其相对
维数为 $e$。则 $\alpha$ 等维，当且仅当每个平坦拉回 $f_i^*\alpha$ 都等维。
\end{lemma}

\begin{proof}
略。提示：仿照引理 \ref{relative-cycles-lemma-equidimensional-fppf-descent} 的证明，
可证 $\alpha$ 的支撑之闭包 $W$ 沿 $f_i$ 的逆像，是 $f_i^*\alpha$ 的支撑之
闭包 $W_i$。于是，若对所有 $i$，$W_i \to S$ 的相对维数
$\leq r + e$，则 $W \to S$ 的相对维数 $\leq r$。
\end{proof}





\section{权函数与相对零循环}
\label{relative-cycles-section-weightings}

\noindent
本节说明，对拟有限态射，相对零循环与“态射续篇”第
\ref{more-morphisms-section-weightings} 节所定义的权函数密切相关。

\medskip\noindent
设 $S$ 为局部 Noether 概形，$f : X \to S$ 为概形的局部拟有限态射。则
$z(X/S, 0) = z_{equi}(X/S, 0)$，并且当 $r > 0$ 时 $z(X/S, r) = 0$。
给定 $\alpha \in z(X/S, 0)$，定义映射
$$
w_\alpha : X \longrightarrow \mathbf{Z},\quad
x \mapsto \alpha(x) [\kappa(x) : \kappa(s)]_i \quad\text{其中 }s = f(x)
$$
这里 $\alpha(x)$ 表示纤维 $X_s$ 上的 $0$-循环 $\alpha_s$ 中 $x$ 的系数，
而 $[K : k]_i$ 表示有限域扩张的不可分次数。

\begin{lemma}
\label{relative-cycles-lemma-weightings-pre}
在上述情形中，若 $g : S' \to S$ 为局部 Noether 概形之间的态射，则
$w_{g^*\alpha} = w_\alpha \circ g'$，其中 $g' : X' \to X$ 是投影
$X' = S' \times_S X \to X$。
\end{lemma}

\begin{proof}
设 $x' \in X'$ 在 $S', S, X$ 中的像分别为 $s', s, x$。则沿
$\kappa(s')/\kappa(s)$ 对 $[x]$ 作基变换后，$[x']$ 的系数等于局部环
$(\kappa(s') \otimes_{\kappa(s)} \kappa(x))_\mathfrak q$ 的长度；这里
$\mathfrak q$ 是与 $x'$ 对应的素理想。因此，只需证明
$$
[\kappa(x) : \kappa(s)]_i =
\text{length}((\kappa(s') \otimes_{\kappa(s)} \kappa(x))_\mathfrak q)
[\kappa(x') : \kappa(s')]_i
$$
即可得到与基变换的相容性。设 $k/\kappa(s')$ 为一个代数闭包。选取位于
$\mathfrak q$ 之上的素理想
$\mathfrak p \subset k \otimes_{\kappa(s)} \kappa(x)$。假设能够证明
$$
[\kappa(x) : \kappa(s)]_i =
\text{length}((k \otimes_{\kappa(s)} \kappa(x))_\mathfrak p)
\quad\text{且}\quad
[\kappa(x') : \kappa(s')]_i =
\text{length}((k \otimes_{\kappa(s')} \kappa(x'))_\mathfrak p)
$$
则结论成立，因为
$$
\text{length}((\kappa(s') \otimes_{\kappa(s)} \kappa(x))_\mathfrak q)
\text{length}((k \otimes_{\kappa(s')} \kappa(x'))_\mathfrak p)
=
\text{length}((k \otimes_{\kappa(s)} \kappa(x))_\mathfrak p)
$$
由“代数”引理 \ref{algebra-lemma-pullback-module} 以及
$\kappa(s') \otimes_{\kappa(s)} \kappa(x) \to k \otimes_{\kappa(s)} \kappa(x)$
的平坦性得出。为证明上述两个等式，只需证明第一个。令
$\kappa(x)/\kappa/\kappa(s)$ 为“域”引理
\ref{fields-lemma-separable-first} 中构造的子域。于是
$$
k \otimes_{\kappa(s)} \kappa(x) =
\prod\nolimits_{\sigma : \kappa \to k}
k \otimes_{\sigma, \kappa} \kappa(x)
$$
且每个因子都是次数为
$[\kappa(x) : \kappa] = [\kappa(x) : \kappa(s)]_i$ 的局部环，正合所需。
\end{proof}

\begin{lemma}
\label{relative-cycles-lemma-weightings}
设 $S$ 为局部 Noether 概形，$f : X \to S$ 为概形的局部拟有限态射。
\begin{enumerate}
\item 对 $\alpha \in z(X/S, 0)$，上述构造的映射
$w_\alpha : X \to \mathbf{Z}$ 是权函数。
\item 若 $X$ 拟紧，则给定权函数 $w : X \to \mathbf{Z}$，存在整数
$n > 0$，使得对某个 $\alpha \in z(X/S, 0)$ 有 $nw = w_\alpha$。
\item 若 $S$ 是 $\mathbf{F}_p$ 上的概形，则 (2) 中的整数 $n$ 可以取为
素数 $p$ 的幂。
\end{enumerate}
\end{lemma}

\begin{proof}
设 $\alpha \in z(X/S, 0)$，并选取图
$$
\xymatrix{
X \ar[d]_f & U \ar[l]^h \ar[d]^\pi \\
Y & V \ar[l]_g
}
$$
如“态射续篇”定义 \ref{more-morphisms-definition-weighting} 所述。以
$\beta \in z(U/V, 0)$ 表示基变换 $g^*\alpha$ 的限制。由与基变换的相容性
（引理 \ref{relative-cycles-lemma-weightings-pre}），$w_\beta = w_\alpha \circ h$，所以只需
证明 $\int_\pi w_\beta$ 在 $V$ 上局部常值。注意
\begin{align*}
\left( \int_\pi w_\beta \right)(v)
& =
\sum\nolimits_{u \in U, \pi(u) = v} 
\beta(u) [\kappa(u) : \kappa(v)]_i [\kappa(u) : \kappa(v)]_s \\
& =
\sum\nolimits_{u \in U, \pi(u) = v} \beta(u)[\kappa(u) : \kappa(v)]
\end{align*}
最后一个表达式是 $\pi_*\beta \in z(V/V, 0)$ 中 $v$ 的系数。由引理
\ref{relative-cycles-lemma-relative-cycle-smooth}，该函数在 $V$ 上局部常值。

\medskip\noindent
反之，设 $w : X \to S$ 为权函数，且 $X$ 拟紧。选取充分可除的整数 $n$。
令 $\alpha$ 为 $X/S$ 各纤维上的 $0$-循环族，使得对 $s \in S$ 有
$$
\alpha_s =
\sum\nolimits_{f(x) = s} \frac{n w(x)}{[\kappa(x) : \kappa(s)]_i} [x]
$$
作为 $X_s$ 上的零循环。这一定义有意义，因为 $f$ 的纤维普遍有界
（“态射”引理
\ref{morphisms-lemma-locally-quasi-finite-qc-source-universally-bounded}），
故可以选取 $n$，使右边对所有 $s \in S$ 都取整数。只要证明 $\alpha$
是相对 $0$-循环，引理的最后一个断言也随之成立。为此，需要证明
$\alpha$ 与沿离散赋值环的特化相容。注意，$w_\alpha$ 的构造以及引理
\ref{relative-cycles-lemma-weightings-pre} 的证明都没有使用 $\alpha$ 是相对循环这一点。
因此 $w_\alpha = nw$，且该等式在基变换后仍成立。此外，权函数的基变换
仍是权函数；见“态射续篇”引理
\ref{more-morphisms-lemma-weighting-base-change}。从而归约到下一段所述问题。

\medskip\noindent
假设 $S$ 是离散赋值环的谱，其一般点为 $\eta$、闭点为 $0$。设
$w : X \to S$ 为权函数，且 $X$ 在 $S$ 上拟有限。令 $\alpha$ 为上一段
构造的 $X/S$ 各纤维上的 $0$-循环族（取适当的 $n$）。需要证明
$sp_{X/S}(\alpha_\eta) = \alpha_0$。令 $\beta \in z(X/S, 0)$ 为 $X/S$
上的相对 $0$-循环，满足 $\beta_\eta = \alpha_\eta$ 且
$\beta_0 = sp_{X/S}(\alpha_\eta)$。于是
$w' = w_\beta - nw : X \to \mathbf{Z}$ 是权函数（使用上面的结论），
并在 $X$ 中映到 $\eta$ 的点处为零。由“态射续篇”引理
\ref{more-morphisms-lemma-weighting-specialization}，$w' = 0$。这意味着
$w_\beta = nw$，故如所需，$\alpha = \beta$。
\end{proof}

\begin{example}
\label{relative-cycles-example-n-must-be-positive-sometimes}
设 $p$ 为素数，$k'/k$ 为次数 $p^f$ 的纯不可分域扩张。令
$X = \Spec(k')$、$S = \Spec(k)$。把 $X$ 的唯一一点 $x$ 映到 $1$ 的映射
$w : X \to \mathbf{Z}$ 是 $X \to S$ 的权函数。另一方面，群
$z(X/S, 0)$ 由循环 $\alpha = [x]$ 自由生成。因此在此情形中，引理
\ref{relative-cycles-lemma-weightings} 所需的最小整数 $n$ 是 $n = p^f$。
\end{example}

\begin{example}
\label{relative-cycles-example-Merkurjev}
上述讨论可以用来“解释” A.S.~Merkurjev 的一个例子
\cite[Example 3.5.10]{SV}。设 $p$ 为素数，并令
$k = \mathbf{F}_p(a, b)$ 为 $\mathbf{F}_p$ 关于 $a$ 和 $b$ 的纯超越扩张。
令
$$
A = k[x, y, z]/(a x^p + b y^p - z^p)
$$
这是 $k$ 上有限型的正规整环。谱 $S = \Spec(A)$ 除了唯一的奇点 $s$ 外
正则；该点对应于 $A$ 的极大理想 $(x, y, z)$，其剩余域为 $k$。令
$k' = k(a^{1/p}, b^{1/p})$、$B = k'[x, y]$。利用嵌入 $k \to k'$，考虑
映射 $A \to B$：它把 $x$ 映到 $x$、把 $y$ 映到 $y$、把 $z$ 映到
$a^{1/p}x + b^{1/p}y$。令 $X = \Spec(B)$，并考虑与环映射 $A \to B$
对应的态射 $f : X \to S$。由“态射续篇”引理
\ref{more-morphisms-lemma-weighting-quasi-finite-Noetherian} 为 $f$ 产生的
权函数 $w : X \to \mathbf{Z}$ 恒等于 $1$，因为 $B$ 的分式域在 $A$ 上
纯不可分！令 $V = f^{-1}(U) \subset X$。由奇迹平坦性，$V \to U$ 平坦，
次数为 $p$。因此循环 $\alpha = [V/V/U]_0 \in z(V/U, 0)$ 满足
$w_\alpha = pw|_V$。然而，位于 $s$ 上方的 $X$ 的唯一一点 $x$ 的剩余域
为 $k'$，它在 $k$ 上的次数为 $p^2$；故不存在 $z(X/S, 0)$ 中限制为
$\alpha$ 的元素 $\beta$。事实上，由“态射续篇”引理
\ref{more-morphisms-lemma-weighting-specialization}，权函数
$pw - w_\beta$ 必为零；这将由上面引理证明中的公式蕴含
$$
\beta_s = \frac{pw(x)}{[\kappa(x) : \kappa(s)]_i} [x] =
\frac{1}{p}[x]
$$
于 $Z_0(X_s)$ 中，而这是不可能的\footnote{在 \cite{SV} 中，确实存在
相对循环群 $Cycl(X/S, 0)$ 中与 $pw$ 对应的元素，而且该循环确实特化为
$X_s$ 上系数非整数的循环。}。当然，循环 $p\alpha$ 是 $z(X/S, 0)$ 中
唯一一个元素的限制。
\end{example}









\section{有效相对循环}
\label{relative-cycles-section-effective}


\noindent
定义如下。

\begin{definition}
\label{relative-cycles-definition-effective}
设 $f : X \to S$ 为概形态射。假设 $S$ 局部 Noether，且 $f$ 局部有限型。
设 $r \geq 0$ 为整数。若对所有 $s \in S$，$\alpha_s$ 都是有效循环
（Chow 同调，定义 \ref{chow-definition-effective-cycle}），则称 $X/S$ 上的
相对 $r$-循环 $\alpha$ {\it 有效}。$X/S$ 上所有有效相对 $r$-循环所成的
幺半群记作 $z^{eff}(X/S, r)$。
\end{definition}

\noindent
下面将证明有效相对循环等维；见引理
\ref{relative-cycles-lemma-effective-equidimensional}。

\begin{lemma}
\label{relative-cycles-lemma-effective-functoriality}
设 $f : X \to S$ 为概形态射。假设 $S$ 局部 Noether，且 $f$ 局部有限型。
设 $r \geq 0$ 为整数，并令 $\alpha$ 为 $X/S$ 上的相对 $r$-循环。若
$\alpha$ 有效，则 $\alpha$ 的任意限制、基变换、平坦拉回或固有前推都有效。
\end{lemma}

\begin{proof}
略。
\end{proof}

\begin{lemma}
\label{relative-cycles-lemma-check-effective}
设 $f : X \to S$ 为概形态射。假设 $S$ 局部 Noether，且 $f$ 局部有限型。
设 $r \geq 0$ 为整数，并令 $\alpha$ 为 $X/S$ 上的相对 $r$-循环。则检验
$\alpha$ 有效时，可以在 $X$ 和 $S$ 上 Zariski 局部地进行。
\end{lemma}

\begin{proof}
略。
\end{proof}

\begin{lemma}
\label{relative-cycles-lemma-effective-descent}
设 $f : X \to S$ 为概形态射。假设 $S$ 局部 Noether，且 $f$ 局部有限型。
设 $r \geq 0$ 为整数，并令 $\alpha$ 为 $X/S$ 上的相对 $r$-循环。设
$g : S' \to S$ 为满射。则 $\alpha$ 有效，当且仅当基变换 $g^*\alpha$
有效。
\end{lemma}

\begin{proof}
略。
\end{proof}

\begin{lemma}
\label{relative-cycles-lemma-effective-descent-pullbacks}
设 $f : X \to S$ 为概形态射。假设 $S$ 局部 Noether，且 $f$ 局部有限型。
设 $r, e \geq 0$ 为整数，并令 $\alpha$ 为 $X/S$ 上的相对 $r$-循环。
设 $\{f_i : X_i \to X\}$ 为一族联合满射的局部有限型平坦态射，且其相对
维数为 $e$。则 $\alpha$ 有效，当且仅当每个平坦拉回 $f_i^*\alpha$ 都有效。
\end{lemma}

\begin{proof}
略。
\end{proof}

\begin{lemma}
\label{relative-cycles-lemma-effective-support-closed}
设 $f : X \to S$ 为概形态射。假设 $S$ 局部 Noether，且 $f$ 局部有限型。
设 $r, e \geq 0$ 为整数，并令 $\alpha$ 为 $X/S$ 上的相对 $r$-循环。
若 $\alpha$ 有效，则 $\text{Supp}(\alpha)$ 在 $X$ 中闭。
\end{lemma}

\begin{proof}
设 $g : S' \to S$ 是一个不可约分支的浸入，并把该分支视作整闭子概形。
由引理 \ref{relative-cycles-lemma-effective-functoriality} 和
\ref{relative-cycles-lemma-support-family}，只需证明基变换 $g^*\alpha$ 的支撑在
$S' \times_S S$ 中闭。因此可以假设 $S$ 为一般点是 $\eta$ 的整概形。
我们将证明 $\text{Supp}(\alpha)$ 是 $\text{Supp}(\alpha_\eta)$ 的闭包。
为此，任取 $s \in S$。存在态射 $g : S' \to S$，其中 $S'$ 是离散赋值环
的谱，它把一般点 $\eta' \in S'$ 映到 $\eta$，把闭点 $0 \in S'$ 映到
$s$；见“性质”引理
\ref{properties-lemma-locally-Noetherian-specialization-dvr}。于是只需证明
$g^*\alpha$ 的支撑等于 $\text{Supp}((g^\alpha)_{\eta'})$ 的闭包。这就
归约到下一段所述情形。

\medskip\noindent
此处 $S$ 是离散赋值环的谱，其一般点为 $\eta$、闭点为 $0$。需要证明
$\text{Supp}(\alpha)$ 是 $\text{Supp}(\alpha_\eta)$ 的闭包。由于 $\alpha$
有效，可以写成 $\alpha_\eta = \sum n_i[Z_i]$，其中 $n_i > 0$，且
$Z_i \subset X_\eta$ 是维数为 $r$ 的整闭子概形。由于
$\alpha_0 = sp_{X/S}(\alpha_\eta)$，有
$\alpha_0 = \sum n_i [\overline{Z}_{i, 0}]_r$，其中 $\overline{Z}_i$ 是
$Z_i$ 的闭包。由“代数簇”引理
\ref{varieties-lemma-dominate-valuation-ring-dimension-fibres}，
$\overline{Z}_{i, 0}$ 是维数为 $r$ 的等维概形。由于 $n_i > 0$，
$\text{Supp}(\alpha_0)$ 等于所有 $\overline{Z}_{i, 0}$ 的并；它就是
$\bigcup \overline{Z}_i$ 在 $0$ 上的纤维，而后者又是 $\bigcup Z_i$ 的
闭包，正合所需。
\end{proof}

\begin{lemma}
\label{relative-cycles-lemma-effective-equidimensional}
设 $f : X \to S$ 为概形态射。假设 $S$ 局部 Noether，且 $f$ 局部有限型。
设 $r, e \geq 0$ 为整数，并令 $\alpha$ 为 $X/S$ 上的相对 $r$-循环。
若 $\alpha$ 有效，则 $\alpha$ 等维。
\end{lemma}

\begin{proof}
假设 $\alpha$ 有效。由引理 \ref{relative-cycles-lemma-effective-support-closed}，支撑
$\text{Supp}(\alpha)$ 在 $X$ 中闭。由于 $\text{Supp}(\alpha) \to S$ 的
纤维是循环 $\alpha_s$ 的支撑，因而维数为 $r$，所以 $\alpha$ 等维。
\end{proof}

\begin{remark}
\label{relative-cycles-remark-representable}
设 $f : X \to S$ 为概形态射，其中 $S$ 局部 Noether，且 $f$ 局部有限型。
我们可以问，下述反变函子
$$
\begin{matrix}
\text{局部有限型的概形 }S'\text{，} \\
\text{其底概形为 }S
\end{matrix}
\longrightarrow
z^{eff}(X'/S', r)\text{，其中 }X' = S' \times_S X
$$
是否可表。由于 $z(X'/S', r) = z(X'_{red}/S'_{red}, r)$，答案不可能为是
（我们把构造一个实际反例留给读者）。更合适的问题是：能否在左端找到一个
子范畴，使该函子限制在其上后可表。引理 \ref{relative-cycles-lemma-seminormalize-base}
提示我们至少应限制到 $S$ 上半正规概形的范畴。

\medskip\noindent
若 $S/\Spec(\mathbf{Q})$ 是 Nagata 的，且 $f$ 为投射态射，则事实证明，
$S' \mapsto z^{eff}(X'/S', r)$ 在半正规的 $S'$ 所成范畴上可表。
粗略地说，这正是 \cite[Theorem 3.21]{KRC} 的内容。

\medskip\noindent
若 $S$ 有正特征点，那么即使把半正规性换成弱正规性，上述结论也不再成立；
若对任意双有理泛同胚 $T' \to T$ 都存在截面，则称局部 Noether 概形 $T$
为弱正规的。一个例子是取 $S = \Spec(k)$，其中 $k$ 为特征 $2$ 的域，并考虑
$X = \mathbf{A}^2_k$ 上次数为 $2$ 的 $0$-循环。具体而言，在
$W = X \times_S X$ 上有一个典范相对 $0$-循环
$\alpha \in z^{eff}(X_W/W, 0)$：对 $w = (x_1, x_2) \in W = X^2$，有
$\alpha_w = [x_1] + [x_2]$。此循环在交换两因子的对合
$\sigma : W \to W$ 下不变。由于 $W$ 光滑（从而正规，进而弱正规），若
$z(-/-, r)$ 在 $k$ 上有限型弱正规概形的范畴上由 $M$ 表示，我们便得到从
$W$ 到 $M$ 的一个 $\sigma$-不变态射。它继而定义从商概形
$\text{Sym}^2_S(X) = W/\langle \sigma \rangle$ 到 $M$ 的态射。由于
$\text{Sym}^2_S(X)$ 正规，由 $M$ 的模性质，我们会在
$X \times_S \text{Sym}^2_S(X) / \text{Sym}^2_S(X)$ 上得到一个相对
$0$-循环 $\beta$，其到 $W$ 的拉回为 $\alpha$。然而，这样的循环 $\beta$
并不存在。事实上，写 $X = \Spec(k[u, v])$，则概形
$\text{Sym}^2_S(X)$ 是下述环的谱：
$$
k[u_1 + u_2, u_1u_2, v_1 + v_2, v_1v_2, u_1v_1 + u_2v_2]
\subset
k[u_1, u_2, v_1, v_2]
$$
对角线 $u_1 = u_2, v_1 = v_2$ 在 $\text{Sym}^2_S(X)$ 中的像是闭子概形
$V = \Spec(k[u_1^2, v_1^2])$；这里用到了 $k$ 的特征为 $2$。考察 $V$ 的
一般点 $\eta$，则 $\beta_\eta$ 会是
$\mathbf{A}^2_{k(u_1^2, v_1^2)}$ 上次数为 $2$ 的零维循环，而它到
$\mathbf{A}^2_{k(u_1, u_2)}$ 的拉回将是
$2[\text{坐标为 } (u_1, v_2)]$。这显然不可能。

\medskip\noindent
上述讨论并不与 \cite[Theorem 4.13]{KRC} 矛盾，因为该定理中的 Chow 簇只对
一个函子给出粗表示（事实上是两个不同的函子；经过一些工作可以看出，在
$X$ 投射时只有其中一个与我们的函子一致）。类似地，
\cite[Section 4.4]{SV} 证明：当 $X/S$ 投射时，预层
$S' \mapsto z^{eff}(S' \times_S X/S', r)$ 的 $h$-层化等于某个可表函子的
$h$-层化。
\end{remark}

\begin{remark}
\label{relative-cycles-remark-equidimensional-over-geometrically-unibranch}
设 $f : X \to S$ 为概形态射，$r \geq 0$，并设 $Z \subset X$ 为闭子概形。
假设
\begin{enumerate}
\item $S$ 是 Noether 且几何单枝的，
\item $f$ 有限型，并且
\item $Z \to S$ 的相对维数 $\leq r$。
\end{enumerate}
则对所有充分可除的整数 $n \geq 1$，存在唯一的 $X/S$ 上有效相对
$r$-循环 $\alpha$，使得对 $S$ 的每个一般点 $\eta$ 都有
$\alpha_\eta = n[Z_\eta]_r$。这是 \cite[Theorem 3.4.2]{SV} 的一种改述。
若以后需要这一结果，我们将在此准确陈述并证明它。
\end{remark}











\section{固有相对循环}
\label{relative-cycles-section-proper}

\noindent
在我们的设定中，下面大概是恰当的定义。

\begin{definition}
\label{relative-cycles-definition-proper}
设 $f : X \to S$ 为概形态射。假设 $S$ 局部 Noether，且 $f$ 局部有限型。
设 $r \geq 0$ 为整数。若 $\alpha$ 的支撑（注记
\ref{relative-cycles-remark-supports-family}）包含在某个在 $S$ 上固有的闭子集
$W \subset X$ 中（概形上同调，定义
\ref{coherent-definition-proper-over-base}），则称 $X/S$ 上的相对
$r$-循环 $\alpha$ 为{\it 固有相对循环}。$X/S$ 上所有固有相对
$r$-循环所成的群记为 $c(X/S, r)$。
\end{definition}

\noindent
由概形上同调，引理 \ref{coherent-lemma-closed-closed-proper-over-base}，
这恰好意味着支撑的闭包在底上固有。要看出这些循环构成一个群，可用
概形上同调，引理 \ref{coherent-lemma-union-closed-proper-over-base}。

\begin{lemma}
\label{relative-cycles-lemma-proper-functoriality}
设 $f : X \to S$ 为概形态射。假设 $S$ 局部 Noether，且 $f$ 局部有限型。
设 $r \geq 0$ 为整数，并令 $\alpha$ 为 $X/S$ 上的相对 $r$-循环。若
$\alpha$ 固有，则 $\alpha$ 的任意基变换都固有。
\end{lemma}

\begin{proof}
略。
\end{proof}

\begin{lemma}
\label{relative-cycles-lemma-proper-h-descent}
设 $f : X \to S$ 为概形态射。假设 $S$ 局部 Noether，且 $f$ 局部有限型。
设 $r \geq 0$ 为整数，令 $\alpha$ 为 $X/S$ 上的相对 $r$-循环，并设
$\{g_i : S_i \to S\}$ 为一个 h-覆盖。则 $\alpha$ 固有，当且仅当每个
基变换 $g_i^*\alpha$ 都固有。
\end{lemma}

\begin{proof}
若 $\alpha$ 固有，则由引理 \ref{relative-cycles-lemma-proper-functoriality}，每个
$g_i^*\alpha$ 也固有。反之，假设每个 $g_i^*\alpha$ 都固有。为证明
$\alpha$ 固有，显然只需在 $S$ 上仿射局部地工作，故可设 $S$ 仿射。
于是可用一族 $\{T_j \to S\}$ 加细覆盖 $\{S_i \to S\}$，其中
$g : T \to S$ 是固有满态射，且 $T = \bigcup T_j$ 为开覆盖。
因此，$\beta = g^*\alpha$ 在 $Y = T \times_S X$ 上作为 $T$ 上的循环是
固有的。由引理 \ref{relative-cycles-lemma-support-family}，$\beta$ 的支撑是 $\alpha$ 的
支撑在态射 $f : Y \to X$ 下的逆像。因此，$f^{-1}\text{Supp}(\alpha)$
在该概形中的闭包 $W \subset Y$ 在 $T$ 上固有。又因态射 $T \to S$
固有，所以 $W$ 在 $S$ 上固有。于是由概形上同调，引理
\ref{coherent-lemma-functoriality-closed-proper-over-base}，像
$f(W) \subset X$ 是在 $S$ 上固有的闭子集。由于 $f(W)$ 包含
$\text{Supp}(\alpha)$，故 $\alpha$ 固有。
\end{proof}



\section{固有且等维的相对循环}
\label{relative-cycles-section-proper-equidimensional}

\noindent
设 $f : X \to S$ 为概形态射。假设 $S$ 局部 Noether，且 $f$ 局部有限型。
设 $r \geq 0$ 为整数。若 $X/S$ 上的相对 $r$-循环 $\alpha$ 既等维
（定义 \ref{relative-cycles-definition-equidimensional}）又固有（定义
\ref{relative-cycles-definition-proper}），则称 $\alpha$ 为{\it 固有且等维的相对循环}。
$X/S$ 上所有固有且等维的相对 $r$-循环所成的群记为
$c_{equi}(X/S, r)$。

\medskip\noindent
类似地，若 $X/S$ 上的相对 $r$-循环 $\alpha$ 既有效（定义
\ref{relative-cycles-definition-effective}）又固有（定义 \ref{relative-cycles-definition-proper}），
则称 $\alpha$ 为{\it 固有且有效的相对循环}。$X/S$ 上所有固有且有效的相对
$r$-循环所成的幺半群记为 $c^{eff}(X/S, r)$。注意，由引理
\ref{relative-cycles-lemma-effective-equidimensional}，这些循环都是等维的。

\medskip\noindent
于是有如下包含映射图：
$$
\xymatrix{
c^{eff}(X/S, r) \ar[r] \ar[d] &
c_{equi}(X/S, r) \ar[r] \ar[d] &
c(X/S, r) \ar[d] \\
z^{eff}(X/S, r) \ar[r] &
z_{equi}(X/S, r) \ar[r] &
z(X/S, r)
}
$$



\section{循环上的作用}
\label{relative-cycles-section-action}

\noindent
设 $S$ 为局部 Noether、普遍链状的概形，并赋予维数函数 $\delta$，参见
Chow 同调，第 \ref{chow-section-setup} 节。设 $X \to Y$ 为 $S$ 上的
概形态射，且二者均在 $S$ 上局部有限型。设 $r \geq 0$。最后，令
$\alpha$ 为 $X/Y$ 的纤维上的 $r$-循环族。对 $e \in \mathbf{Z}$，
我们将构造运算
$$
\alpha \cap - : Z_e(Y) \longrightarrow Z_{r + e}(X)
$$
具体而言，给定 $\beta \in Z_e(Y)$，写成 $\beta = \sum n_i[Z_i]$，其中
$Z_i \subset Y$ 是 $\delta$-维数为 $e$ 的整闭子概形，且族 $Z_i$ 在概形
$Y$ 中局部有限。令 $y_i \in Z_i$ 为一般点，并写
$\alpha_{y_i} = \sum m_{ij} [V_{ij}]$。于是 $V_{ij} \subset X_{y_i}$
是维数为 $r$ 的整闭子概形，且族 $V_{ij}$ 在概形 $X_{y_i}$ 中局部有限。
定义
$$
\alpha \cap \beta = \sum n_i m_{ij} [\overline{V}_{ij}]
\quad\in\quad
Z_{r + e}(X)
$$
这里，$\overline{V}_{ij} \subset X$ 是态射 $V_{ij} \to X_{y_i} \to X$
的概形论像；等价地，$\overline{V}_{ij} \subset X$ 是一个支配地映到
$Z_i \subset Y$、且一般纤维为 $V_{ij}$ 的整闭子概形。容易推出
$\dim_\delta(\overline{V}_{ij}) = r + e$，并且闭子概形族
$\overline{V}_{ij} \subset X$ 局部有限（验证从略）。因此
$\alpha \cap \beta$ 的确是 $Z_{r + e}(X)$ 的元素。

\begin{lemma}
\label{relative-cycles-lemma-action-bilinear}
上述构造是双线性的，即
$(\alpha_1 + \alpha_2) \cap \beta = \alpha_1 \cap \beta +
\alpha_2 \cap \beta$，且 $\alpha \cap (\beta_1 + \beta_2) =
\alpha \cap \beta_1 + \alpha \cap \beta_2$.
\end{lemma}

\begin{proof}
略。
\end{proof}

\begin{lemma}
\label{relative-cycles-lemma-action-opens}
若 $U \subset X$ 与 $V \subset Y$ 为开集，且 $f(U) \subset V$，则
$(\alpha \cap \beta)|_U$ 等于 $\alpha|_U \cap \beta|_V$。
\end{lemma}

\begin{proof}
由上面对 $\alpha \cap \beta$ 的显式描述立即可得。
\end{proof}

\begin{lemma}
\label{relative-cycles-lemma-action-base-change}
构造 $\alpha \cap \beta$ 与平坦基变换和平坦拉回相容（详见证明）。
\end{lemma}

\begin{proof}
令 $(S, \delta)$、$(S', \delta')$、$g : S' \to S$ 与
$c \in \mathbf{Z}$ 如 Chow 同调，情形 \ref{chow-situation-setup-base-change}
所述。设 $X \to Y$ 为 $S$ 上局部有限型的概形态射，记 $X' \to Y'$
为 $X \to Y$ 沿 $g$ 的基变换。令 $\alpha$ 为 $X/Y$ 的纤维上的
$r$-循环族，并令 $\beta \in Z_e(Y)$。记 $\alpha'$ 为 $\alpha$ 沿
$Y' \to Y$ 的基变换，记 $\beta' = g^*\beta \in Z_{e + c}(Y')$ 为
$\beta$ 沿 $g$ 的拉回，参见 Chow 同调，第 \ref{chow-section-change-base}
节。与基变换相容，是指 $\alpha' \cap \beta'$ 是
$\alpha \cap \beta$ 的基变换。

\medskip\noindent
先证与基变换相容。由于要证明的是 $X'$ 上循环的等式，由引理
\ref{relative-cycles-lemma-action-opens}，可以在 $Y$ 上局部地工作。因此可设 $Y$ 仿射。
特别地，$\beta$ 是素循环的有限线性组合。由于 $- \cap -$ 关于第二个
变量线性（引理 \ref{relative-cycles-lemma-action-bilinear}），只需在
$\beta = [Z]$ 时证明该等式，其中 $Z \subset Y$ 是 $\delta$-维数为
$e$ 的整闭子概形。

\medskip\noindent
令 $y \in Z$ 为一般点。写 $\alpha_y = \sum m_j [V_j]$，并令
$\overline{V}_j$ 为 $V_j$ 在 $X$ 中的闭包。则
$$
\alpha \cap \beta = \sum m_j[\overline{V}_j]
$$
$\beta$ 的基变换是 $Y' = Y \times_S S'$ 上的循环
$\beta' = \sum [Z \times_S S']_{e + c}$。令
$Z'_a \subset Z \times_S S'$ 为不可约分支，记 $y'_a \in Z'_a$ 为其
一般点，并记 $n_a$ 为 $Z'_a$ 在 $Z \times_S S'$ 中的重数。我们有
$$
\beta' = \sum [Z \times_S S']_{e + c} = \sum n_a[Z'_a]
$$
由于 $\alpha'$ 是 $\alpha$ 沿 $Y' \to Y$ 的基变换，故
$\alpha'_{y'_a} = \sum m_j [V_{j, \kappa(y'_a)}]_r$。令
$V'_{jab} \subset V_{j, \kappa(y'_a)}$ 为不可约分支，并记 $m_{jab}$
为 $V'_{jab}$ 在 $V_{j, \kappa(y'_a)}$ 中的重数。我们有
$$
\alpha'_{y'_a} = \sum m_j [V_{j, \kappa(y'_a)}]_r =
\sum m_j m_{jab} [V'_{jab}]
$$
因此
$$
\alpha' \cap \beta' = \sum n_a m_j m_{jab} [\overline{V}'_{jab}]
$$
其中 $\overline{V}'_{jab}$ 是 $V'_{jab}$ 在 $X'$ 中的闭包。因此，为证明
所需等式，只需证明
\begin{enumerate}
\item $\overline{V}_j \times_S S'$ 的不可约分支是概形
$\overline{V}'_{jab}$，并且
\item $\overline{V}'_{jab}$ 在 $\overline{V}_j \times_S S'$ 中的重数等于
$n_a m_{jab}$。
\end{enumerate}
注意，$V_j \to \overline{V}_j$ 是整概形间的双有理态射。态射
$V_j \times_S S' \to V_j$ 与
$\overline{V}_j \times_S S' \to \overline{V}_j$ 都是平坦的，因而把
不可约分支的一般点分别映到 $V_j$ 与 $\overline{V}_j$ 的（唯一）一般点。
由此，$V_j \times_S S' \to \overline{V}_j \times_S S'$ 是双有理态射，
故在不可约分支之间诱导双射，并识别其重数。这意味着只需证明
$V_j \times_S S'$ 的不可约分支是概形 $V'_{jab}$，且 $V'_{jab}$ 在
$V_j \times_S S'$ 中的重数等于 $n_a m_{jab}$。然而，这正是在断言图
$$
\xymatrix{
Z_r(V_j) \ar[r] & Z_{r + c}(V_j \times_S S') \\
Z_0(\Spec(\kappa(y))) \ar[r] \ar[u] &
Z_c(\Spec(\kappa(y)) \times_S S') \ar[u]
}
$$
交换，其中水平箭头是沿
$\Spec(\kappa(y)) \times_S S' \to \Spec(\kappa(y))$ 的基变换，竖直箭头
是平坦拉回。这已在 Chow 同调，引理
\ref{chow-lemma-pullback-base-change-pullback} 中证明。

\medskip\noindent
引理中关于平坦拉回的陈述含义如下。令 $(S, \delta)$、$X \to Y$、
$\alpha$ 与 $\beta$ 如上面构造 $\alpha \cap \beta$ 时所述。设
$Y' \to Y$ 是局部有限型、相对维数为 $c$ 的平坦态射。令 $\alpha'$ 为
$\alpha$ 沿 $Y' \to Y$ 的基变换，令 $\beta'$ 为 $\beta$ 的平坦拉回。
与平坦拉回相容，是指 $\alpha' \cap \beta'$ 是
$\alpha \cap \beta$ 沿 $X \times_Y Y' \to Y$ 的平坦拉回。事实上，令
$S = Y$ 且 $S' = Y'$，这就是上面讨论的特例。
\end{proof}

\begin{lemma}
\label{relative-cycles-lemma-action-coherent}
令 $(S, \delta)$ 与 $f : X \to Y$ 如上。设 $\mathcal{F}$ 为凝聚
$\mathcal{O}_X$-模，且对所有 $y \in Y$ 都有
$\dim(\text{Supp}(\mathcal{F}_y)) \leq r$。设 $\mathcal{G}$ 为凝聚
$\mathcal{O}_Y$-模，且 $\dim_\delta(\text{Supp}(\mathcal{G})) \leq e$。
令 $\alpha = [\mathcal{F}/X/Y]_r$（例
\ref{relative-cycles-example-family-associated-module}），并令 $\beta = [\mathcal{G}]_e$
（Chow 同调，定义
\ref{chow-definition-cycle-associated-to-coherent-sheaf}）。若
$\mathcal{F}$ 在 $Y$ 上平坦，则 $\alpha \cap \beta =
[\mathcal{F} \otimes_{\mathcal{O}_X} f^*\mathcal{G}]_{r + e}$。
\end{lemma}

\begin{proof}
注意
$$
\text{Supp}(\mathcal{F} \otimes_{\mathcal{O}_X} f^*\mathcal{G}) =
\text{Supp}(\mathcal{F}) \cap f^{-1}\text{Supp}(\mathcal{G}) =
\bigcup\nolimits_{y \in \text{Supp}(\mathcal{G})} \text{Supp}(\mathcal{F}_y)
$$
由此，它是 $\delta$-维数 $\leq r + e$ 的闭子集。因此表达式
$[\mathcal{F} \otimes_{\mathcal{O}_X} f^*\mathcal{G}]_{r + e}$ 有意义。

\medskip\noindent
沿用构造 $\alpha \cap \beta$ 时引入的记号
$\beta = \sum n_i[Z_i]$、$y_i \in Z_i$、
$\alpha_{y_i} = \sum m_{ij} [V_{ij}]$ 与 $\overline{V}_{ij}$。由于
$\beta = [\mathcal{G}]_e$，可知 $Z_i$ 是
$\text{Supp}(\mathcal{G})$ 中 $\delta$-维数为 $e$ 的不可约分支。
类似地，$V_{ij}$ 是 $\text{Supp}(\mathcal{F}_{y_i})$ 中维数为 $r$ 的
不可约分支。结合这一点与第一段中的等式可知，$\overline{V}_{ij}$ 是
$\text{Supp}(\mathcal{F} \otimes_{\mathcal{O}_X} f^*\mathcal{G})$
中 $\delta$-维数为 $r + e$ 的不可约分支。因此，为证明本引理，只需证明
$$
\text{length}_{\mathcal{O}_{X, \xi_{ij}}}(
(\mathcal{F} \otimes_{\mathcal{O}_X} f^*\mathcal{G})_{\xi_{ij}})
=
\text{length}_{\mathcal{O}_{X_{y_i}, \xi_{ij}}}((\mathcal{F}_{y_i})_{\xi_{ij}})
\cdot
\text{length}_{\mathcal{O}_{Y, y_i}}(\mathcal{G}_{y_i})
$$
由证明第一段，左端等于下述 $B = \mathcal{O}_{X, \xi_{ij}}$-模的长度：
$$
\mathcal{G}_{y_i}
\otimes_{\mathcal{O}_{Y, y_i}}
\mathcal{F}_{\xi_{ij}} =
M \otimes_A N
$$
这里，$M = \mathcal{G}_{y_i}$ 是有限长度的
$A = \mathcal{O}_{Y, y_i}$-模，而 $N = \mathcal{F}_{\xi_{ij}}$ 是有限
$B$-模，并且 $N/\mathfrak m_AN$ 长度有限。由于 $\mathcal{F}$ 在 $Y$
上平坦，模 $N$ 是 $A$-平坦的。公式右端等于
$$
\text{length}_B(N/\mathfrak m_A N) \cdot \text{length}_A(M)
$$
因此，公式左端与右端都关于 $M$ 可加（使用 $N$ 在 $A$ 上的平坦性）。
故只需在 $M = \kappa_A$ 为剩余域时证明该公式，而此时结论立即成立。
\end{proof}

\begin{lemma}
\label{relative-cycles-lemma-action-closed}
令 $(S, \delta)$ 与 $f : X \to Y$ 如上。设 $Z \subset X$ 为在 $Y$
上相对维数 $\leq r$ 的闭子概形。令 $\alpha = [Z/X/Y]_r$（例
\ref{relative-cycles-example-family-associated-closed}）。设 $W \subset Y$ 为
$\delta$-维数 $\leq e$ 的闭子概形。令 $\beta = [W]_e$（Chow 同调，
定义 \ref{chow-definition-cycle-associated-to-closed-subscheme}）。若
$Z$ 在 $Y$ 上平坦，则 $\alpha \cap \beta = [Z \times_Y W]_{r + e}$。
\end{lemma}

\begin{proof}
在引理 \ref{relative-cycles-lemma-action-coherent} 中取
$\mathcal{F} = \mathcal{O}_Z$ 和 $\mathcal{F} = \mathcal{O}_W$ 即得。
\end{proof}

\begin{lemma}
\label{relative-cycles-lemma-flat-pullback-as-action}
令 $(S, \delta)$ 与 $f : X \to Y$ 如上。假设 $f$ 平坦，且相对维数为
$e$。对 $\beta \in Z_r(Y)$，在 $Z_{e + r}(X)$ 中有
$f^*\beta = [X/X/Y]_e \cap \beta$。
\end{lemma}

\begin{proof}
只需在 $X$ 上局部地证明该等式，故可设 $Y$ 仿射。此时 $\beta$ 是素循环的
有限整数线性组合。由于等式两端都关于 $\beta$ 可加，可设
$\beta = [W]$，其中 $W$ 是 $Y$ 中 $\delta$-维数为 $r$ 的整闭子概形。
令 $f^{-1}(W) = W \times_Y X$ 为 $W$ 在 $X$ 中的概形论逆像。由引理
\ref{relative-cycles-lemma-action-closed}，有
$[X/X/Y]_e \cap \beta = [W \times_Y X]_{r + e}$；而由 Chow 同调，定义
\ref{chow-definition-flat-pullback}，有
$f^*\beta = [f^{-1}(W)]_{r + e}$。故得到所需等式。
\end{proof}

\begin{lemma}
\label{relative-cycles-lemma-action-push-pull}
令 $(S, \delta)$ 如上。设
$$
\xymatrix{
X' \ar[r]_f \ar[d] & X \ar[d] \\
Y' \ar[r]^g & Y
}
$$
是 $S$ 上局部有限型概形的笛卡尔图，且 $g$ 固有。设 $r, e \geq 0$，
令 $\alpha$ 为 $X/Y$ 的纤维上的 $r$-循环族，并令
$\beta' \in Z_e(Y')$。则
$f_*(g^*\alpha \cap \beta') = \alpha \cap g_*\beta'$。
\end{lemma}

\begin{proof}
由于要证明的是 $X$ 上循环的等式，由引理 \ref{relative-cycles-lemma-action-opens}，可在
$Y$ 上局部地工作。因此可设 $Y$ 仿射，从而 $Y'$ 拟紧。特别地，
$\beta'$ 是素循环的有限线性组合。由于 $- \cap -$ 关于第二个变量线性
（引理 \ref{relative-cycles-lemma-action-bilinear}），只需在 $\beta' = [Z']$ 时证明该
等式，其中 $Z' \subset Y'$ 是 $\delta$-维数为 $e$ 的整闭子概形。令
$Z = g(Z')$。这是 $Y$ 中 $\delta$-维数 $\leq e$ 的整闭子概形。为简便
起见，我们假设 $Z$ 的 $\delta$-维数等于 $e$，并把另一种（更容易的）
情形留给读者。令 $y \in Z$ 与 $y' \in Z'$ 为一般点。写
$\alpha_y = \sum m_j[V_j]$，其中 $V_j \subset X_y$ 是维数为 $r$ 的
整闭子概形。

\medskip\noindent
先假设 $g$ 为闭浸入。于是 $g_*\beta' = [Z]$ 且
$(g^*\alpha)_{y'} = \sum n_j[V_j]$；这是有意义的，因为 $V_j$ 包含在
$X_y$ 的闭子概形 $X'_{y'}$ 中。因此，此时等式显然：两边都得到
$\sum m_j[\overline{V}_j]$，其中 $\overline{V}_j$ 是 $V_j$ 在闭子概形
$X' \subset X$ 中的闭包。

\medskip\noindent
回到上述 $\beta' = [Z']$ 的一般情形。令 $W = Z \times_Y X$ 且
$W' = Z' \times_{Y'} X'$。考虑笛卡尔方块
$$
\xymatrix{
W \ar[r] \ar[d] & X \ar[d] \\
Z \ar[r] & Y
}
\quad
\xymatrix{
W' \ar[r] \ar[d] & X' \ar[d] \\
Z' \ar[r] & Y'
}
\quad
\xymatrix{
W' \ar[r] \ar[d] & W \ar[d] \\
Z' \ar[r] & Z
}
$$
由上一段已知前两个方块的结论，一个形式论证表明，只需对最后一个方块与
元素 $\beta' = [Z'] \in Z_e(Z')$ 证明结论。于是归结为下一段讨论的情形。

\medskip\noindent
假设 $Y' \to Y$ 是 $\delta$-维数为 $e$ 的整概形之间的一般有限态射，且
$\beta' = [Y']$。此时，$f_*(g^*\alpha \cap \beta')$ 与
$\alpha \cap g_*\beta'$ 都可写成支配 $Y$ 的素循环之和。因此，为检验
等式，可以把 $Y$ 换成一个非空开子概形。作此替换后，可设 $g$ 有限且
平坦，譬如次数为 $d \geq 1$。当然，这意味着
$g_*\beta' = g_*[Y'] = d[Y]$。又有 $\beta' = [Y'] = g^*[Y]$。因此
$$
f_*(g^*\alpha \cap \beta') =
f_*(g^*\alpha \cap g^*[Y]) =
f_*f^*(\alpha \cap [Y]) =
d (\alpha \cap [Y]) =
\alpha \cap g_*\beta'
$$
正合所需。第二个等号来自引理 \ref{relative-cycles-lemma-action-base-change}，第三个等号
来自 Chow 同调，引理 \ref{chow-lemma-finite-flat}。
\end{proof}







\section{Chow 群上的作用}
\label{relative-cycles-section-action-chow}

\noindent
当 $\alpha$ 为相对 $r$-循环时，第 \ref{relative-cycles-section-action} 节的运算
$\alpha \cap -$ 可经由有理等价因子化，并定义一个双变类。

\begin{lemma}
\label{relative-cycles-lemma-closed-in-X-gysin}
令 $(S, \delta)$ 如第 \ref{relative-cycles-section-action} 节所述。设
$f : X' \to X$ 为 $S$ 上局部有限型概形之间的固有态射。令
$(\mathcal{L}, s, i : D \to X)$ 如 Chow 同调，定义
\ref{chow-definition-gysin-homomorphism} 所述。构造图
$$
\xymatrix{
D' \ar[d]_g \ar[r]_{i'} & X' \ar[d]^f \\
D \ar[r]^i & X
}
$$
如 Chow 同调，注记 \ref{chow-remark-pullback-pairs} 所述。若
$\mathcal{L}|_D \cong \mathcal{O}_D$，则对任意
$\alpha' \in Z_{k + 1}(X')$，在 $Z_k(D)$ 中有
$i^*f_*\alpha' = g_*(i')^*\alpha'$。
\end{lemma}

\begin{proof}
该陈述有意义，因为所有运算都在循环层面定义；关于 Gysin 映射，参见
Chow 同调，注记 \ref{chow-remark-gysin-on-cycles}。假设对某个整闭子概形
$W' \subset X'$ 有 $\alpha = [W']$。令 $W = f(W') \subset X$。若
$W' \not \subset D'$，则 $W \not \subset D$，并且
$$
[W' \cap D']_k = \text{div}_{\mathcal{L}'|_{W'}}({s'|_{W'}})
\quad\text{且}\quad
[W \cap D]_k = \text{div}_{\mathcal{L}|_W}(s|_W)
$$
因此，由 Chow 同调，引理 \ref{chow-lemma-equal-c1-as-cycles}，第一个循环
在 $f_*$ 下的像等于第二个循环。故等式在循环层面成立。若
$W' \subset D'$，则 $W \subset D$，且按构造两端都为零。
\end{proof}

\begin{lemma}
\label{relative-cycles-lemma-action-gysin}
令 $(S, \delta)$ 如第 \ref{relative-cycles-section-action} 节所述。设 $X \to Y$ 为
$S$ 上局部有限型概形之间的态射。设 $r \geq 0$，并令
$\alpha \in z(X/Y, r)$ 为 $X/Y$ 上的相对 $r$-循环。令
$(\mathcal{L}, s, i : D \to Y)$ 如 Chow 同调，定义
\ref{chow-definition-gysin-homomorphism} 所述。构造笛卡尔图
$$
\xymatrix{
E \ar[d] \ar[r]_j & X \ar[d] \\
D \ar[r]^i & Y
}
$$
参见 Chow 同调，注记 \ref{chow-remark-pullback-pairs}。若
$\mathcal{L}|_D \cong \mathcal{O}_D$，则对 $e \in \mathbf{Z}$，图
$$
\xymatrix{
Z_e(D) \ar[rr]_{i^*\alpha \cap -} & &
Z_{e + r}(E) \\
Z_{e + 1}(Y) \ar[u]^{i^*} \ar[rr]^{\alpha \cap -} & &
Z_{r + e + 1}(X) \ar[u]_{j^*}
}
$$
交换，其中竖直箭头 $i^*$ 与 $j^*$ 是 Chow 同调，注记
\ref{chow-remark-gysin-on-cycles} 中的循环 Gysin 映射。
\end{lemma}

\begin{proof}
先作预备说明。假设 $g : Y' \to Y$ 为包络（Chow 同调，定义
\ref{chow-definition-envelope}）。分别用 $D', i', E', j', X', \alpha'$
表示 $D, i, E, j, X, \alpha$ 沿 $g$ 的基变换，并记
$f : X' \to X$ 为投影。假设本引理对
$D', i', E', j', X', Y', \alpha'$ 成立。于是，若
$\beta' \in Z_{e + 1}(Y')$，则
\begin{align*}
i^*\alpha \cap i^*g_*\beta'
& =
i^*\alpha \cap f_*(i')^*\beta' \\
& =
f_*(f^*i^*\alpha \cap (i')^*\beta') \\
& =
f_*((i')^*\alpha' \cap (i')^*\beta') \\
& =
f_*((j')^*(\alpha' \cap \beta')) \\
& =
j^*(f_*(f^*\alpha \cap \beta')) \\
& =
j^*(\alpha \cap g_*\beta')
\end{align*}
这里，第一个等号来自引理 \ref{relative-cycles-lemma-closed-in-X-gysin}，第二个等号来自
引理 \ref{relative-cycles-lemma-action-push-pull}，第三个等号是 $\alpha'$ 的定义，第四个
等号来自本引理对 $D', i', E', j', X', \alpha'$ 成立的假设，第五个等号
来自引理 \ref{relative-cycles-lemma-closed-in-X-gysin}，第六个等号来自引理
\ref{relative-cycles-lemma-action-push-pull}。因此，本引理对
$g_* : Z_{e + 1}(Y') \to Z_e(Y)$ 的像成立。然而，由于 $g$ 完全分解，
此映射满射，故本引理对 $D, i, E, j, X, Y, \alpha$ 成立。

\medskip\noindent
令 $\beta \in Z_{e + 1}(Y)$。我们须证明，在 $E$ 上作为循环有
$(D \to Y)^*\alpha \cap i^*\beta = j^*(\alpha \cap \beta)$。这一问题在
$E$ 上是局部的，因而可以把 $X$ 与 $Y$ 换成开子概形。（这里用到运算
$i^*$、$j^*$、$\alpha \cap - $ 与 $(D \to Y)^*\alpha \cap -$ 的构造
均与局部化交换。对 Gysin 映射这是显然的；对其余运算则由引理
\ref{relative-cycles-lemma-action-opens} 得到。）因此可设 $X$ 与 $Y$ 仿射，并归结到
下一段讨论的情形。

\medskip\noindent
假设 $X$ 与 $Y$ 拟紧。由证明第一段与引理 \ref{relative-cycles-lemma-get-cycles}，还可设
$\alpha$ 属于 (\ref{relative-cycles-equation-cycle-classes}) 的像。由所涉运算的线性，
可设 $\alpha = [Z/X/Y]_r$，其中 $Z \subset X$ 是在 $Y$ 上平坦且相对维数
$\leq r$ 的闭子概形。又因 $Y$ 拟紧，循环 $\beta$ 是素循环的有限线性
组合。由于所涉运算是线性的，只需在 $\beta = [W]$ 时证明等式，其中
$W \subset Y$ 是 $\delta$-维数为 $e + 1$ 的整闭子概形。

\medskip\noindent
若 $W \subset D$，则一方面 $i^*[W] = 0$，另一方面
$\alpha \cap [W]$ 支撑在 $E$ 上，故也有
$j^*(\alpha \cap [W]) = 0$。因此此时等式成立。

\medskip\noindent
设 $W \not \subset D$。于是 $i^*[W] = [D \cap W]_e$。注意，
$\alpha = [Z/X/Y]_r$ 沿 $i$ 的拉回 $i^*\alpha$ 是
$[(E \cap Z)/E/D]_r$，而且
$(E \cap Z) = E \times_Y Z = D \times_Y Z$ 在 $D$ 上平坦。因此，两次
使用引理 \ref{relative-cycles-lemma-action-closed} 可得
$$
i^*\alpha \cap i^*[W] =
[(E \cap Z) \times_D (D \cap W)]_{r + e} =
[E \cap (Z \times_Y W)]_{r + e} =
j^*(\alpha \cap [W])
$$
正合所需。
\end{proof}

\begin{proposition}
\label{relative-cycles-proposition-get-bivariant-class}
令 $(S, \delta)$ 如第 \ref{relative-cycles-section-action} 节所述。设 $X \to Y$ 为
$S$ 上局部有限型概形之间的态射。设 $r \geq 0$，并令
$\alpha \in z(X/Y, r)$ 为 $X/Y$ 上的相对 $r$-循环。对每个局部有限型
态射 $g : Y' \to Y$ 与每个 $e \in \mathbf{Z}$，考虑运算
$$
g^*\alpha \cap - : Z_e(Y') \to Z_{r + e}(X')
$$
其中 $X' = Y' \times_Y X$。这族运算可经由有理等价因子化，从而定义双变类
$c(\alpha) \in A^{-r}(X \to Y)$。
\end{proposition}

\begin{proof}
由引理 \ref{relative-cycles-lemma-action-gysin} 与 Chow 同调，引理
\ref{chow-lemma-factors-through-rational-equivalence}，该运算可经由有理等价
因子化。由 Chow 同调，引理 \ref{chow-lemma-bivariant-weaker} 以及引理
\ref{relative-cycles-lemma-action-push-pull}、\ref{relative-cycles-lemma-action-base-change} 和
\ref{relative-cycles-lemma-action-gysin}，所得 Chow 群上的运算是双变类。
\end{proof}

\begin{remark}
\label{relative-cycles-remark-characterize-relative-cycles}
令 $(S, \delta)$ 如第 \ref{relative-cycles-section-action} 节所述。设 $X \to Y$ 为
$S$ 上局部有限型概形之间的态射，并设 $r \geq 0$。设 $c$ 是一条规则，
它给每个局部有限型态射 $g : Y' \to Y$ 与每个 $e \in \mathbf{Z}$ 配置运算
$$
c \cap - : Z_e(Y') \to Z_{r + e}(X')
$$
并且它按引理 \ref{relative-cycles-lemma-action-gysin} 的意义与固有前推、平坦拉回和
Gysin 映射相容。则我们断言，存在 $X/Y$ 上的相对 $r$-循环 $\alpha$，
使得对上述每个 $g$ 都有 $c \cap = g^*\alpha \cap -$。若以后需要这一
结论，我们将在此仔细陈述并证明它。
\end{remark}









\section{纤维上循环族的复合}
\label{relative-cycles-section-compose-families}

\noindent
设 $X \to Y \to S$ 为概形态射，且二者均局部有限型。设
$r, e \geq 0$。令 $\alpha$ 为 $X/Y$ 的纤维上的 $r$-循环族，令
$\beta$ 为 $Y/S$ 的纤维上的 $e$-循环族。定义
$$
(\alpha \circ \beta)_s = (Y_s \to Y)^*\alpha \cap \beta_s
$$
便得到 $X/S$ 的纤维上的 $(r + e)$-循环族 $\alpha \circ \beta$。更准确地，
表达式 $(Y_s \to Y)^*\alpha$ 表示 $\alpha$ 沿 $Y_s \to Y$ 的基变换，
所得为 $X_s/Y_s$ 的纤维上的 $r$-循环族；运算 $- \cap -$ 已在第
\ref{relative-cycles-section-action} 节中定义并研究\footnote{明确地说，我们把
$s = \Spec(\kappa(s))$ 用作底概形，并取 $\delta(s) = 0$。}。

\begin{lemma}
\label{relative-cycles-lemma-construction-bilinear}
上述构造是双线性的，即有
$(\alpha_1 + \alpha_2) \circ \beta \alpha_1 \circ \beta +
\alpha_1 \circ \beta$，且 $\alpha \circ (\beta_1 + \beta_2) =
\alpha \circ \beta_1 + \alpha \circ \beta_2$.
\end{lemma}

\begin{proof}
略。提示：由引理 \ref{relative-cycles-lemma-action-bilinear}，该构造在纤维上双线性。
\end{proof}

\begin{lemma}
\label{relative-cycles-lemma-construction-opens}
若 $U \subset X$ 与 $V \subset Y$ 为开集，且 $f(U) \subset V$，则
$(\alpha \circ \beta)|_U$ 等于 $\alpha|_U \circ \beta|_V$。
\end{lemma}

\begin{proof}
略。提示：在纤维上使用引理 \ref{relative-cycles-lemma-action-opens}。
\end{proof}

\begin{lemma}
\label{relative-cycles-lemma-construction-base-change}
$\alpha \circ \beta$ 的构造与基变换相容。
\end{lemma}

\begin{proof}
设 $g : S' \to S$ 为概形态射。记 $X' \to Y'$ 为 $X \to Y$ 沿 $g$ 的
基变换，记 $\alpha'$ 为 $\alpha$ 沿 $Y' \to Y$ 的基变换，并记
$\beta'$ 为 $\beta$ 沿 $S' \to S$ 的基变换。断言的含义是，
$\alpha' \circ \beta'$ 是 $\alpha \circ \beta$ 沿
$g : S' \to S$ 的基变换。

\medskip\noindent
设 $s' \in S'$ 为像是 $s \in S$ 的点。则
$$
(\alpha' \circ \beta')_{s'} = (Y'_{s'} \to Y')^*\alpha' \cap \beta'_{s'}
$$
注意
$$
(Y'_{s'} \to Y')^*\alpha' =
(Y'_{s'} \to Y')^*(Y' \to Y)^*\alpha =
(Y'_{s'} \to Y_s)^*(Y_s \to Y)^*\alpha
$$
而 $\beta'_{s'}$ 是 $\beta_s$ 沿
$s' = \Spec(\kappa(s')) \to \Spec(\kappa(s)) = s$ 的基变换。因此，把引理
\ref{relative-cycles-lemma-action-base-change} 应用于 $(Y_s \to Y)^*\alpha$、$\beta_s$、
$X_s \to Y_s \to s$ 以及沿 $s' \to s$ 的基变换，即得结论。
\end{proof}

\begin{lemma}
\label{relative-cycles-lemma-construction-coherent}
设 $f : X \to Y$ 与 $Y \to S$ 为概形态射，且二者均局部有限型。设
$r, e \geq 0$。设 $\mathcal{F}$ 为有限型拟凝聚 $\mathcal{O}_X$-模，且对
所有 $y \in Y$ 有 $\dim(\text{Supp}(\mathcal{F}_y)) \leq r$。设
$\mathcal{G}$ 为有限型拟凝聚 $\mathcal{O}_Y$-模，且对所有 $s \in S$ 有
$\dim(\text{Supp}(\mathcal{G}_s)) \leq e$。若
$\alpha = [\mathcal{F}/X/Y]_r$ 且 $\beta = [\mathcal{G}/Y/S]_e$（例
\ref{relative-cycles-example-family-associated-module}），并且 $\mathcal{F}$ 在 $Y$ 上
平坦，则 $\alpha \circ \beta =
[\mathcal{F} \otimes_{\mathcal{O}_X} f^*\mathcal{G}/X/S]_{r + e}$。
\end{lemma}

\begin{proof}
首先，$\mathcal{F} \otimes_{\mathcal{O}_X} f^*\mathcal{G}$ 是有限型拟凝聚
$\mathcal{O}_X$-模。设 $s \in S$。注意
$$
(\mathcal{F} \otimes_{\mathcal{O}_X} f^*\mathcal{G})_s =
\mathcal{F}_s \otimes_{\mathcal{O}_{X_s}} f_s^*\mathcal{G}_s
$$
这来自张量积的右正合性。此外，$\mathcal{F}_s$ 作为平坦模的基变换，在
$Y_s$ 上平坦。因此，由引理 \ref{relative-cycles-lemma-action-coherent} 得到等式
$(\alpha \circ \beta)_s =
[(\mathcal{F} \otimes_{\mathcal{O}_X} f^*\mathcal{G})_s]_{r + e}$。
\end{proof}

\begin{lemma}
\label{relative-cycles-lemma-construction-closed}
设 $f : X \to Y$ 与 $Y \to S$ 为概形态射，且二者均局部有限型。设
$r, e \geq 0$。设 $Z \subset X$ 为在 $Y$ 上相对维数 $\leq r$ 的闭子概形，
并设 $W \subset Y$ 为在 $S$ 上相对维数 $\leq e$ 的闭子概形。若
$\alpha = [Z/X/Y]_r$ 且 $\beta = [W/Y/S]_e$（例
\ref{relative-cycles-example-family-associated-closed}），并且 $Z$ 在 $Y$ 上平坦，则
$\alpha \circ \beta = [Z \times_Y W/X/S]_{r + e}$。
\end{lemma}

\begin{proof}
在引理 \ref{relative-cycles-lemma-construction-coherent} 中取
$\mathcal{F} = \mathcal{O}_Z$ 与 $\mathcal{F} = \mathcal{O}_W$ 即得。
\end{proof}

\begin{lemma}
\label{relative-cycles-lemma-flat-pullback-as-composition}
设 $f : X \to Y$ 与 $Y \to S$ 为概形态射，且二者均局部有限型。假设
$f$ 平坦，且相对维数为 $e$。令 $\beta$ 为 $Y/S$ 的纤维上的
$r$-循环族。则作为 $X/S$ 的纤维上的 $(e + r)$-循环族，有
$f^*\beta = [X/X/Y]_e \circ \beta$。
\end{lemma}

\begin{proof}
展开定义，并使用 $[X/X/Y]_e$ 的构造与基变换相容这一事实（引理
\ref{relative-cycles-lemma-family-associated-closed}），即可由引理
\ref{relative-cycles-lemma-flat-pullback-as-action} 得到结论。
\end{proof}

\begin{lemma}
\label{relative-cycles-lemma-construction-push-pull}
设 $S$ 为概形。设
$$
\xymatrix{
X' \ar[r]_f \ar[d] & X \ar[d] \\
Y' \ar[r]^g & Y
}
$$
是 $S$ 上局部有限型概形的笛卡尔图，且 $g$ 固有。设 $r, e \geq 0$。
令 $\alpha$ 为 $X/Y$ 的纤维上的 $r$-循环族，并令 $\beta'$ 为 $Y'/S$
的纤维上的 $e$-循环族。则
$f_*(g^*(\alpha) \circ \beta') = \alpha \circ g_*\beta'$。
\end{lemma}

\begin{proof}
展开定义后，结论来自引理 \ref{relative-cycles-lemma-action-push-pull}。
\end{proof}

\begin{lemma}
\label{relative-cycles-lemma-construction-composition}
令 $(S, \delta)$ 如 Chow 同调，情形 \ref{chow-situation-setup} 所述。设
$X \to Y \to Z$ 为 $S$ 上局部有限型概形之间的态射。设
$r, s, e \geq 0$。则
$$
(\alpha \circ \beta) \cap \gamma = \alpha \cap (\beta \cap \gamma)
\quad\text{在}\quad Z_{r + s + e}(X)
$$
其中 $\alpha$ 为 $X/Y$ 的纤维上的 $r$-循环族，$\beta$ 为 $Y/Z$ 的纤维
上的 $s$-循环族，且 $\gamma \in Z_e(Z)$。
\end{lemma}

\begin{proof}
由于要证明的是 $X$ 上循环的等式，由引理 \ref{relative-cycles-lemma-action-opens}，可在
$Z$ 上局部地工作。因此可设 $Z$ 仿射。特别地，$\gamma$ 是素循环的有限
线性组合。由于 $- \cap -$ 关于第二个变量线性（引理
\ref{relative-cycles-lemma-action-bilinear}），只需在 $\gamma = [W]$ 时证明等式，其中
$W \subset Z$ 是 $\delta$-维数为 $e$ 的整闭子概形。

\medskip\noindent
令 $z \in W$ 为一般点。在 $Z_s(Y_z)$ 中写
$\beta_z = \sum m_j[V_j]$。于是 $\beta \cap \gamma$ 等于
$\sum m_j[\overline{V}_j]$，其中 $\overline{V}_j \subset Y$ 是一个在
$Y \to Z$ 下映入 $W$、且一般纤维为 $V_j$ 的整闭子概形。令
$y_j \in V_j$ 为一般点。我们也把它看成 $\overline{V}_j$ 的一般点
（它映到 $W$ 中的 $z$）。在 $Z_r(X_{y_j})$ 中写
$\alpha_{y_j} = \sum n_{jk} [W_{jk}]$。则
$\alpha \cap (\beta \cap \gamma)$ 等于
$$
\sum m_j n_{jk} [\overline{W}_{jk}]
$$
其中 $\overline{W}_{jk} \subset X$ 是一个在 $X \to Y$ 下映入
$\overline{V}_j$、且一般纤维为 $W_{jk}$ 的整闭子概形。

\medskip\noindent
另一方面，考虑
$$
(\alpha \circ \beta)_z = (Y_z \to Y)^*\alpha \cap \beta_z =
(Y_z \to Y)^*\alpha \cap (\sum m_j [V_j])
$$
由 $- \cap -$ 的构造，它等于 $X_z$ 上的循环
$$
\sum m_j n_{jk} [(\overline{W}_{jk})_z]
$$
因此按定义得到
$$
(\alpha \circ \beta) \cap [W] =
\sum m_j n_{jk} [\widetilde{W}_{jk}]
$$
其中 $\widetilde{W}_{jk} \subset X$ 是一个在 $X \to Z$ 下映入 $W$、且
一般纤维为 $(\overline{W}_{jk})_z$ 的整闭子概形。显然必有
$\widetilde{W}_{jk} = \overline{W}_{jk}$，证明完成。
\end{proof}









\section{相对循环的复合}
\label{relative-cycles-section-compose}

\noindent
设 $S$ 为局部 Noether 概形。设 $X \to Y$ 为 $S$ 上局部有限型概形之间的
态射。我们将利用第 \ref{relative-cycles-section-compose-families} 节的构造定义映射
$$
z(X/Y, r) \otimes_\mathbf{Z} z(Y/S, e) \longrightarrow z(X/S, r + e),\quad
\alpha \otimes \beta \longmapsto \alpha \circ \beta
$$
我们已经知道该构造是双线性的（引理 \ref{relative-cycles-lemma-construction-bilinear}），
故只要证明下述结论，就得到所显示的箭头。

\begin{lemma}
\label{relative-cycles-lemma-well-defined}
若 $\alpha$ 与 $\beta$ 都是相对循环，则 $\alpha \circ \beta$ 也是相对循环。
\end{lemma}

\begin{proof}
由引理 \ref{relative-cycles-lemma-construction-base-change}，$\alpha \circ \beta$ 的构造与
基变换相容。因此可设 $S$ 是某个离散赋值环的谱，其一般点为 $\eta$、
闭点为 $0$；我们须证明
$sp_{X/S}((\alpha \circ \beta)_\eta) = (\alpha \circ \beta)_0$。由于要
证明的是循环的等式，可在 $Y$ 与 $X$ 上局部地工作（这里使用引理
\ref{relative-cycles-lemma-construction-opens} 与 \ref{relative-cycles-lemma-specialization-flat-pullback}
看出各构造与限制交换）。故可设 $X$ 与 $Y$ 仿射。由引理
\ref{relative-cycles-lemma-get-cycles}，可找到固有完全分解态射 $g : Y' \to Y$，使得
$g^*\alpha$ 属于 (\ref{relative-cycles-equation-cycle-classes}) 的像。

\medskip\noindent
由于态射族 $g_\eta : Y'_\eta \to Y_\eta$ 完全分解，可找到
$\beta'_\eta \in Z_e(Y'_\eta)$，使得
$\beta_\eta = \sum g_{\eta, *}\beta'_\eta$，参见 Chow 同调，引理
\ref{chow-lemma-envelope}。令 $\beta'_0 = sp_{Y'/S}(\beta'_\eta)$，从而
$\beta' = (\beta'_\eta, \beta'_0)$ 是 $Y'/S$ 上的相对 $e$-循环。于是
$g_*\beta'$ 与 $\beta$ 是 $Y/S$ 上的相对 $e$-循环（引理
\ref{relative-cycles-lemma-relative-cycle-functoriality}），并且它们在 $\eta$ 处取值相同，
故二者相等（引理 \ref{relative-cycles-lemma-uniqueness-extension}）。由线性（引理
\ref{relative-cycles-lemma-construction-bilinear}），只需证明
$\alpha \circ g_*\beta'$ 是相对 $(r + e)$-循环。

\medskip\noindent
令 $X' = X \times_Y Y'$，并记 $f : X' \to X$ 为投影。由引理
\ref{relative-cycles-lemma-construction-push-pull}，有
$\alpha \circ g_*\beta' = f_*(g^*\alpha \circ \beta')$。由引理
\ref{relative-cycles-lemma-relative-cycle-functoriality}，只需证明
$g^*\alpha \circ \beta'$ 是相对 $(r + e)$-循环。使用引理
\ref{relative-cycles-lemma-get-cycles-dvr} 与双线性，即归结到下一段讨论的情形。

\medskip\noindent
假设 $\alpha = [Z/X/Y]_r$ 且 $\beta = [W/Y/S]$，其中 $Z \subset X$
是在 $Y$ 上平坦且相对维数 $\leq r$ 的闭子概形，而 $W \subset Y$ 是在
$S$ 上平坦且相对维数 $\leq e$ 的闭子概形。由引理
\ref{relative-cycles-lemma-construction-closed} 可知
$$
\alpha \circ \beta = [Z \times_X W/X/S]_{r + e}
$$
而 $Z \times_X W \subset X$ 是在 $S$ 上平坦、相对维数
$\leq r + e$ 的闭子概形。由引理
\ref{relative-cycles-lemma-family-associated-closed-specialization}，这是相对
$(r + e)$-循环。
\end{proof}

\begin{lemma}
\label{relative-cycles-lemma-composition-fundamental-cycles}
设 $f : X \to Y$ 与 $g : Y \to S$ 为概形态射。假设 $S$ 局部 Noether，
$g$ 局部有限型且平坦、相对维数为 $e \ge 0$，并且 $f$ 局部有限型且平坦、
相对维数为 $r \geq 0$。则在 $z(X/S, r + e)$ 中有
$[X/X/Y]_r \circ [Y/Y/S]_e = [X/X/S]_{r + e}$。
\end{lemma}

\begin{proof}
这是引理 \ref{relative-cycles-lemma-construction-closed} 的特例。
\end{proof}






\section{与 Suslin--Voevodsky 的比较}
\label{relative-cycles-section-compare}

\noindent
除循环记号取自 Chow 同调，第 \ref{chow-section-cycles} 节及后续内容外，
我们尽量采用与 \cite{SV} 相同的记号。比较如下：
\begin{enumerate}
\item \cite[Section 3.1]{SV} 中有“相对循环”、“维数为 $r$ 的相对循环”
以及“维数为 $r$ 的等维相对循环”的概念。本章中没有相应概念。因此，各群
$Cycl(X/S, r)$, $Cycl_{equi}(X/S, r)$,
$PropCycl(X/S, r)$ 与 $PropCycl_{equi}(X/S, r)$
在本章中没有对应物。
\item 在 \cite[page 36]{SV} 页底定义了各群
$z(X/S, r)$, $c(X/S, r)$, $z_{equi}(X/S, r)$, $c_{equi}(X/S, r)$
。当 $S$ 是分离 Noether 概形且 $X \to S$ 分离有限型时，这些定义与我们的
概念一致。
\item 在 \cite{SV} 中，符号 $z(X/S, r)$ 有时用来表示 $S$ 上有限型概形
范畴上的预层 $S' \mapsto z(S' \times_S X/S', r)$。符号
$c(X/S, r)$、$z_{equi}(X/S, r)$ 与 $c_{equi}(X/S, r)$。
也作类似使用。
\item \cite{SV} 中定义的基变换、平坦拉回和固有前推，在两套定义都适用时
与我们的定义一致。
\item 对 $\alpha \in z(X/S, r)$，第 \ref{relative-cycles-section-action} 节定义的运算
$\alpha \cap - : Z_e(S) \to Z_{e + r}(X)$，在两者都有定义时，与
\cite[Section 3.7]{SV} 中的运算 $Cor(\alpha, -)$ 一致。
\item 对 $X \to Y \to S$，第 \ref{relative-cycles-section-compose} 节定义的复合律
$z(X/Y, r) \otimes_\mathbf{Z} z(Y/S, e) \longrightarrow z(X/S, r + e)$
与 \cite[Corollary 3.7.5]{SV} 中的运算 $Cor_{X/Y}(-, -)$ 一致。
\end{enumerate}











\section{非 Noether 情形下的相对循环}
\label{relative-cycles-section-non-noetherian}

\noindent
我们强烈建议读者跳过本节。

\medskip\noindent
设 $f : X \to S$ 为有限表示的概形态射。设 $r \geq 0$。用
$Hilb(X/S, r)$ 表示所有满足下述条件的闭子概形 $Z \subset X$ 所成的集合：
$Z \to S$ 平坦、有限表示且相对维数 $\leq r$。考虑群同态
\begin{equation}
\label{relative-cycles-equation-cycle-classes-general}
\begin{matrix}
\text{自由 Abel 群} \\
\text{以 }Hilb(X/S, r)
\end{matrix}
\longrightarrow
\begin{matrix}
\text{纤维上的 }r\text{-循环族}\\
\text{对应于 }X/S
\end{matrix}
\end{equation}
它把 $\sum n_i[Z_i]$ 映到 $\sum n_i[Z_i/X/S]_r$。

\begin{lemma}
\label{relative-cycles-lemma-relative-r-cycle-general}
设 $S$ 为拟紧拟分离概形，$f : X \to S$ 为有限表示态射。设
$r \geq 0$，并令 $\alpha$ 为 $X/S$ 的纤维上的 $r$-循环族。下列条件等价：
\begin{enumerate}
\item 存在笛卡尔图
$$
\xymatrix{
X \ar[r] \ar[d] & X_0 \ar[d] \\
S \ar[r] & S_0
}
$$
其中 $X_0 \to S_0$ 是 Noether 概形之间的有限型态射，且存在
$\alpha_0 \in z(X_0/S_0, r)$，使得 $\alpha$ 是 $\alpha_0$ 沿
$S \to S_0$ 的基变换；
\item 存在有限表示的固有完全分解态射 $g : S' \to S$，使得
$g^*\alpha$ 属于 (\ref{relative-cycles-equation-cycle-classes-general}) 的像。
\end{enumerate}
\end{lemma}

\begin{proof}
给定 (1) 中的图与 $\alpha_0 \in z(X_0/S_0, r)$。由引理
\ref{relative-cycles-lemma-get-cycles}，存在固有完全分解态射 $g_0 : S'_0 \to S_0$，
使得 $g_0^*\alpha_0$ 属于 (\ref{relative-cycles-equation-cycle-classes-general}) 的像。
事实上，由于 $S'_0$ 是 Noether 的，$S'_0 \times_{S_0} X_0$ 的每个闭子概形
在 $S'_0$ 上都有限表示。令 $S' = S \times_{S_0} S'_0$，并沿
$S' \to S'_0$ 作基变换，即知 (2) 成立。

\medskip\noindent
反之，假设 (2) 成立。选取有限表示的固有完全分解态射 $g : S' \to S$，
使得 $g^*\alpha$ 属于 (\ref{relative-cycles-equation-cycle-classes-general}) 的像。令
$X' = S' \times_S X$。对某些闭子概形 $Z_a \subset X'$，写
$g^*\alpha = \sum n_a [Z_a/X'/S']_r$，其中这些子概形在 $S'$ 上平坦、
有限表示且相对维数 $\leq r$。

\medskip\noindent
把 $S = \lim S_i$ 写成转移态射为仿射态射的有向极限，其中 $S_i$ 在
$\mathbf{Z}$ 上有限型，参见极限，命题 \ref{limits-proposition-approximate}。
可找到充分大的 $i$，使得存在
\begin{enumerate}
\item 固有完全分解态射 $g_i : S'_i \to S_i$，它到 $S$ 的基变换是
$g : S' \to S$；
\item 令 $X'_i = S'_i \times_{S_i} X_i$，存在闭子概形
$Z_{ai} \subset X'_i$，其在 $S'_i$ 上平坦且相对维数 $\leq r$，并且到
$S'$ 的基变换是 $Z_a$。
\end{enumerate}
为此要使用极限中的各引理
\ref{limits-lemma-descend-finite-presentation},
\ref{limits-lemma-descend-closed-immersion-finite-presentation},
\ref{limits-lemma-descend-flat-finite-presentation},
\ref{limits-lemma-descend-surjective},
\ref{limits-lemma-eventually-proper} 和
\ref{limits-lemma-limit-dimension}
以及更多态射，引理 \ref{more-morphisms-lemma-descend-cd}。考虑
$\alpha'_i = \sum n_a [Z_{ai}/X'_i/S_i]_r \in z(X'_i/S'_i, r)$。按构造，
$\alpha'_i$ 基变换为 $X'/S'$ 的纤维上的 $r$-循环族后，与基变换
$g^*\alpha$ 一致。

\medskip\noindent
令 $S''_i = S'_i \times_{S_i} S'_i$、$X''_i = S''_i \times_{S_i} X_i$，
并令 $S'' = S' \times_S S'$、$X'' = S'' \times_S X$。用
$\text{pr}_1, \text{pr}_2 : S'' \to S'$ 与
$\text{pr}_1, \text{pr}_2 : S''_i \to S'_i$ 表示投影。$X''_i/S''_i$
上的相对 $r$-循环 $\text{pr}_1^*\alpha'_i$ 与
$\text{pr}_1^*\alpha'_i$ 基变换后，成为 $X''/S''$ 的纤维上同一个
$r$-循环族，因为
$\text{pr}_1^*g^*\alpha = \text{pr}_1^*g^*\alpha$。因此态射
$S'' \to S''_i$ 映入
$E = \{s \in S''_i : (\text{pr}_1^*\alpha'_i)_s =
(\text{pr}_1^*\alpha'_i)_s\}$。由引理 \ref{relative-cycles-lemma-relative-cycles-equal}，
这是闭子集。由于 $S'' = \lim_{i' \geq i} S''_{i'}$，由极限，引理
\ref{limits-lemma-limit-contained-in-constructible} 可知，对某个
$i' \geq i$，态射 $S''_{i'} \to S''_i$ 映入 $E$。因此，把 $i$ 换成
$i'$ 后，可设 $\text{pr}_1^*\alpha'_i = \text{pr}_1^*\alpha'_i$。由引理
\ref{relative-cycles-lemma-descend-family}，得到 $X_i/S_i$ 的纤维上唯一的 $r$-循环族
$\alpha_i$，满足 $g_i^*\alpha_i = \alpha'_i$（这里使用了
$S'_i \to S_i$ 完全分解）。由引理 \ref{relative-cycles-lemma-relative-cycles-h-descent}，
有 $\alpha_i \in z(X_i/S_i, r)$。引理 \ref{relative-cycles-lemma-descend-family} 中的
唯一性蕴含 $\alpha_i$ 的基变换是 $\alpha$，故 (1) 成立。
\end{proof}

\noindent
{\bf 讨论。}
若 $f : X \to S$、$r$ 与 $\alpha$ 如引理
\ref{relative-cycles-lemma-relative-r-cycle-general} 所述，则当引理
\ref{relative-cycles-lemma-relative-r-cycle-general} 的等价条件 (1) 与
(2) 成立时，可以称 $\alpha$ 为{\it $X/S$ 上的相对 $r$-循环}。这个定义有
许多良好性质；例如，当 $S$ 为 Noether 概形时，它与先前的定义不冲突，
且第 \ref{relative-cycles-section-families-specialization} 节的大多数结果都推广到这一
设定。

\medskip\noindent
还可作如下进一步推广。设 $S$ 任意，且 $f : X \to S$ 局部有限表示。设
$r \geq 0$，并令 $\alpha$ 为 $X/S$ 的纤维上的 $r$-循环族 $\alpha$。
如果对任意满足 $f(U) \subset V$ 的仿射开集 $U \subset X$ 与
$V \subset S$，限制 $\alpha|_U$ 都是上一段意义下 $U/V$ 上的相对
$r$-循环，则称 $\alpha$ 为{\it $X/S$ 上的相对 $r$-循环}。同样，先前的
许多结果都推广到这一设定。

\medskip\noindent
若以后需要这些推广，我们将在此仔细陈述并证明它们。
% Generated mechanically from the frozen more-etale.tex target.
% Do not edit this generated file; edit the bound locale source instead.

% OK, start here.
%

\setcounter{chapter}{62}
\chapter{进一步的 \'Etale 上同调}



\phantomsection
\label{more-etale-section-phantom}



\section{引言}
\label{more-etale-section-introduction}

\noindent
本章是概形 \'etale 上同调系列的第二章。第一章见“\'Etale 上同调”
第 \ref{etale-cohomology-section-introduction} 节。

\medskip\noindent
粗略地说，本章与前一章的分界如下：凡涉及“叹号函子”（紧支撑上同调及其
右伴随）以及使用这些材料的内容，都归入本章。






\section{延拓截面}
\label{more-etale-section-growing}

\noindent
本节讨论如下类型的结果。

\begin{lemma}
\label{more-etale-lemma-section-support-in-locally-closed-pre}
设 $X$ 为概形，$\mathcal{F}$ 为 $X_\etale$ 上的阿贝尔层。
设 $\varphi : U' \to U$ 为 $X_\etale$ 中的态射。设 $Z' \subset U'$ 为
闭子概形，使得 $Z' \to U' \to U$ 是闭浸入，其像为 $Z \subset U$。
则存在典范双射
$$
\{s \in \mathcal{F}(U) \mid \text{Supp}(s) \subset Z\} =
\{s' \in \mathcal{F}(U') \mid \text{Supp}(s') \subset Z'\}
$$
若 $\varphi^{-1}(Z) = Z'$，则该双射由限制给出。
\end{lemma}

\begin{proof}
考虑 $U'$ 的闭子概形 $Z'' = \varphi^{-1}(Z)$。由于 $Z'$ 在 $U'$ 中闭，
故 $Z' \subset Z''$ 闭。另一方面，$Z' \to Z''$ 是 \'etale 态射
（这是 $Z$ 上两个 \'etale 概形之间的态射），因而是开映射。因此，
对某个闭子集 $T$ 有 $Z'' = Z' \amalg T$。开覆盖
$U' = (U' \setminus T) \cup (U' \setminus Z')$ 表明
$$
\{s' \in \mathcal{F}(U') \mid \text{Supp}(s') \subset Z'\} =
\{s' \in \mathcal{F}(U' \setminus T) \mid \text{Supp}(s') \subset Z'\}
$$
而 \'etale 覆盖
$\{U' \setminus T \to U, U \setminus Z \to U\}$ 表明
$$
\{s \in \mathcal{F}(U) \mid \text{Supp}(s) \subset Z\} =
\{s' \in \mathcal{F}(U' \setminus T) \mid \text{Supp}(s') \subset Z'\}
$$
证明完成。
\end{proof}

\begin{lemma}
\label{more-etale-lemma-section-support-in-locally-closed}
设 $X$ 为概形，$Z \subset X$ 为局部闭子概形，$\mathcal{F}$ 为
$X_\etale$ 上的阿贝尔层。给定开子概形 $U, U' \subset X$，其中
$Z$ 均为闭子概形，则存在典范双射
$$
\{s \in \mathcal{F}(U) \mid \text{Supp}(s) \subset Z\} =
\{s \in \mathcal{F}(U') \mid \text{Supp}(s) \subset Z\}
$$
若 $U' \subset U$，则该双射由限制给出。
\end{lemma}

\begin{proof}
由于 $Z$ 是 $U \cap U'$ 的闭子概形，只须证明 $U' \subset U$ 的情形；
此时结论是引理 \ref{more-etale-lemma-section-support-in-locally-closed-pre} 的特例。
\end{proof}

\noindent
下面引入一个稍非标准、但后文很有用的记号。具体而言，在上述引理
\ref{more-etale-lemma-section-support-in-locally-closed} 的情形下，记
$$
H_Z(\mathcal{F}) = \{s \in \mathcal{F}(U) \mid \text{Supp}(s) \subset Z\}
$$
其中 $U \subset X$ 可任取为一个以 $Z$ 为闭子概形的开子概形。
若读者对这一写法的不精确性感到不安，可以取
$U = X \setminus \partial Z$，其中
$\partial Z = \overline{Z}\setminus Z$ 是 $Z$ 在 $X$ 中的“边界”。
不过，在下文许多论证中，能够灵活选择不同开子概形会发挥作用。
这一构造具有如下性质：
\begin{enumerate}
\item
\label{more-etale-item-inclusion}
若 $Z \subset Z'$ 是 $X$ 的局部闭子概形，且 $Z$ 在 $Z'$ 中闭，
则存在自然单射
$$
H_Z(\mathcal{F}) \to H_{Z'}(\mathcal{F}).
$$
\item
\label{more-etale-item-pullback}
若 $f : Y \to X$ 是概形态射，且 $Z \subset X$ 是局部闭子概形，
则存在自然拉回映射
$f^* : H_Z(\mathcal{F}) \to H_{f^{-1}Z}(f^{-1}\mathcal{F})$。
\end{enumerate}
把这一记号推广到如下情形会很方便：设 $W \in X_\etale$，且
$Z \subset W$ 为局部闭子概形。此时记
$$
H_Z(\mathcal{F}) =
\{s \in \mathcal{F}(U) \mid \text{Supp}(s) \subset Z\} =
H_Z(\mathcal{F}|_{W_\etale})
$$
其中 $U \subset W$ 与上面一样，可任取为一个以 $Z$ 为闭子概形的
开子概形\footnote{事实上，引理
\ref{more-etale-lemma-section-support-in-locally-closed-pre}
表明：若 $X$ 上的 $Z$ 同构于 $X_\etale$ 中某个对象 $W$ 的局部闭子概形，
则 $W$ 的选择无关紧要。}。







\section{紧支撑截面}
\label{more-etale-section-compact-support}

\noindent
本节可参考 \cite[报告 XVII，第 6 节]{SGA4}。
设 $f : X \to Y$ 为分离且局部有限型的概形态射。本节定义函子
$f_! : \textit{Ab}(X_\etale) \to \textit{Ab}(Y_\etale)$
如下：令 $f_!\mathcal{F} \subset f_*\mathcal{F}$ 为由相对于 $Y$ 具有
固有支撑（按适当方式定义）的截面所成的子层。

\medskip\noindent
警告：函子 $f_!$ 是导出范畴上某个函子 $Rf_!$ 的第零上同调层
（此处待补未来引用），但 $Rf_!$ 不是 $f_!$ 的导出函子。

\begin{lemma}
\label{more-etale-lemma-f-shriek-separated}
设 $f : X \to Y$ 为局部有限型概形态射，$\mathcal{F}$ 为
$X_\etale$ 上的阿贝尔层。规则
$$
Y_\etale \longrightarrow \textit{Ab},\quad
V \longmapsto \{s \in f_*\mathcal{F}(V) = \mathcal{F}(X_V) \mid
\text{Supp}(s) \subset X_V \text{ 固有地映到 }V\}
$$
定义了 $f_*\mathcal{F}$ 的阿贝尔子层。
\end{lemma}

\noindent
警告：若 $f$ 不分离，这个层并非“正确的”层。

\begin{proof}
回忆截面的支撑是闭的（“\'Etale 上同调”引理
\ref{etale-cohomology-lemma-support-section-closed}），故“概形上同调”
第 \ref{coherent-section-proper-over-base} 节的材料适用。由上面的引理及
“概形上同调”引理 \ref{coherent-lemma-union-closed-proper-over-base}，
可知我们在 $f_*\mathcal{F}(V)$ 中取出的子集是子群。由“概形上同调”
引理 \ref{coherent-lemma-base-change-closed-proper-over-base}，
可知该规则定义了子预层。
最后，设 $s \in f_*\mathcal{F}(V)$，并设有 \'etale 覆盖
$\{V_i \to V\}$，使得 $s|_{V_i}$ 的支撑在 $V_i$ 上固有。
注意，$s|_{V_i}$ 的支撑是 $s|_V$ 的支撑的逆像（使用以茎刻画支撑的方式
以及“\'Etale 上同调”引理
\ref{etale-cohomology-lemma-stalk-pullback}）。于是由“下降”引理
\ref{descent-lemma-descending-property-proper-over-base}，$s$ 的支撑在
$V$ 上固有。这证明该规则满足层条件。
\end{proof}

\begin{lemma}
\label{more-etale-lemma-separated-etale-shriek}
设 $j : U \to X$ 为分离 \'etale 态射，$\mathcal{F}$ 为
$U_\etale$ 上的阿贝尔层。“\'Etale 上同调”引理
\ref{etale-cohomology-lemma-shriek-into-star-separated-etale}
中的单射 $j_!\mathcal{F} \to j_*\mathcal{F}$ 之像，正是引理
\ref{more-etale-lemma-f-shriek-separated} 的子层。
\end{lemma}

\noindent
也可以把这个引理移到后面，再用两个层的茎的描述来证明。

\begin{proof}
在“\'Etale 上同调”引理
\ref{etale-cohomology-lemma-shriek-into-star-separated-etale} 的证明中，
$j_!\mathcal{F} \to j_*\mathcal{F}$ 是通过构造预层间的映射
$j_{p!}\mathcal{F} \to j_*\mathcal{F}$ 得到的；后者的像显然包含在引理
\ref{more-etale-lemma-f-shriek-separated} 的子层中。由于 $j_!\mathcal{F}$ 是
$j_{p!}\mathcal{F}$ 的层化，故 $j_!\mathcal{F} \to j_*\mathcal{F}$ 的像
也包含在该子层中。
反之，设 $s \in j_*\mathcal{F}(V)$ 的支撑 $Z$ 在 $V$ 上固有。则
$Z \to V$ 有限，且具有闭像 $Z' \subset V$；见“更多态射”引理
\ref{more-morphisms-lemma-characterize-finite}。$s$ 在
$V \setminus Z'$ 上的限制为零，而零截面属于
$j_!\mathcal{F} \to j_*\mathcal{F}$ 的像。
另一方面，若 $v \in Z'$，则可找到 \'etale 邻域
$(V', v') \to (V, v)$，使得有分解
$U_{V'} = W \amalg U'_1 \amalg \ldots \amalg U'_n$，
其中各项都是既开又闭的子概形，$U'_i \to V'$ 为同构，且
$T_{V'} \subset U'_1 \amalg \ldots \amalg U'_n$；见“\'Etale 态射”
引理 \ref{etale-lemma-etale-etale-local-technical}。把同构
$U'_i \to V'$ 取逆，并拉回 $s$，便得到 $n$ 个态射
$\varphi'_i : V' \to U$ 以及 $V'$ 上的截面 $s'_i$。
于是 $V'$ 上 $j_{p!}\mathcal{F}$ 的截面
$\sum (\varphi'_i, s'_i)$（参见“\'Etale 上同调”引理
\ref{etale-cohomology-lemma-shriek-into-star-separated-etale} 的证明中
$j_{p!}\mathcal{F}(V')$ 的公式）按构造映到 $s$ 在 $V'$ 上的限制。
因此，$s$ 在 \'etale 局部属于 $j_!\mathcal{F} \to j_*\mathcal{F}$ 的像，
证明完成。
\end{proof}

\begin{definition}
\label{more-etale-definition-f-shriek-separated}
设 $f : X \to Y$ 为分离（！）且局部有限型的概形态射，$\mathcal{F}$ 为
$X_\etale$ 上的阿贝尔层。引理 \ref{more-etale-lemma-f-shriek-separated} 所构造的
子层 $f_!\mathcal{F} \subset f_*\mathcal{F}$ 称为{\it 紧支撑直像}。
\end{definition}

\noindent
由引理 \ref{more-etale-lemma-separated-etale-shriek}，这不与“\'Etale 上同调”定义
\ref{etale-cohomology-definition-extension-zero} 冲突，因为两种定义均适用时
彼此一致。下面作一个基本检验。

\begin{lemma}
\label{more-etale-lemma-proper-f-shriek}
设 $f : X \to Y$ 为概形的固有态射。则 $f_! = f_*$。
\end{lemma}

\begin{proof}
由 $f_!$ 的构造立即得到。
\end{proof}

\noindent
下面的观察非常有用。

\begin{remark}[关于开嵌入的协变性]
\label{more-etale-remark-covariance-f-shriek-separated}
设 $f : X \to Y$ 为分离且局部有限型的概形态射，$\mathcal{F}$ 为
$X_\etale$ 上的阿贝尔层。设 $X' \subset X$ 为开子概形，并以
$f' : X' \to Y$ 表示 $f$ 的限制。存在典范单射
$$
f'_!(\mathcal{F}|_{X'}) \longrightarrow f_!\mathcal{F}
$$
具体地，设 $V \in Y_\etale$，并考虑截面
$s' \in f'_*(\mathcal{F}|_{X'})(V) = \mathcal{F}(X' \times_Y V)$，
其支撑 $Z'$ 在 $V$ 上固有。则 $Z'$ 在 $X \times_Y V$ 中也闭；见
“概形上同调”引理
\ref{coherent-lemma-functoriality-closed-proper-over-base}。因此存在唯一截面
$s \in \mathcal{F}(X \times_Y V) = f_*\mathcal{F}(V)$，使得它在
$X' \times_Y V$ 上的限制为 $s'$，而在
$X \times_Y V \setminus Z'$ 上的限制为零；见引理
\ref{more-etale-lemma-section-support-in-locally-closed}。这一构造与限制映射相容，
因而诱导所需的层映射 $f'_!(\mathcal{F}|_{X'}) \to f_!\mathcal{F}$，
它显然是单射。按构造得到交换图
$$
\xymatrix{
f'_!(\mathcal{F}|_{X'}) \ar[r] \ar[d] &
f_!\mathcal{F} \ar[d] \\
f'_*(\mathcal{F}|_{X'}) &
f_*\mathcal{F} \ar[l]
}
$$
它关于 $\mathcal{F}$ 具有函子性。显然，若 $X'' \subset X'$ 为开子概形，
且 $f'' = f|_{X''} : X'' \to Y$，则刚构造的典范映射之复合
$f''_!\mathcal{F}|_{X''} \to f'_!\mathcal{F}|_{X'} \to f_!\mathcal{F}$
就是典范映射 $f''_!\mathcal{F}|_{X''} \to f_!\mathcal{F}$。
\end{remark}

\begin{lemma}
\label{more-etale-lemma-compactify-f-shriek-separated}
设 $Y$ 为概形。设 $j : X \to \overline{X}$ 为 $Y$ 上概形的开浸入，
且 $\overline{X}$ 在 $Y$ 上固有。以 $f : X \to Y$ 和
$\overline{f} : \overline{X} \to Y$ 表示结构态射。对
$\mathcal{F} \in \textit{Ab}(X_\etale)$，存在典范同构（见证明）
$$
f_!\mathcal{F} \longrightarrow \overline{f}_!j_!\mathcal{F}
$$
由引理 \ref{more-etale-lemma-proper-f-shriek} 有
$\overline{f}_! = \overline{f}_*$，故作为
$\textit{Ab}(X_\etale) \to \textit{Ab}(Y_\etale)$ 的函子，有
$\overline{f}_* \circ j_! = f_!$。
\end{lemma}

\begin{proof}
由“\'Etale 上同调”引理
\ref{etale-cohomology-lemma-jshriek-open}，有
$(j_!\mathcal{F})|_X = \mathcal{F}$。因此，所示箭头就是评注
\ref{more-etale-remark-covariance-f-shriek-separated} 中对
$\mathcal{G} = j_!\mathcal{F}$ 所得的单射
$f_!(\mathcal{G}|_X) \to \overline{f}_!\mathcal{G}$
。由该映射的显式构造，现在只须证明：若 $V \in Y_\etale$，且
$s \in \overline{f}_!\mathcal{G}(V) = \overline{f}_*\mathcal{G}(V) =
\mathcal{G}(\overline{X}_V)$ 为截面，则 $s$ 的支撑包含在开子概形
$X_V \subset \overline{X}_V$ 中。这立即来自如下事实：
在 $\overline{X} \setminus X$ 的几何点处，$\mathcal{G}$ 的茎为零。
\end{proof}

\noindent
我们希望把 $f_!\mathcal{F}$ 的茎与纤维上的紧支撑截面联系起来。
为叙述这一点，需要先作一个定义。

\begin{definition}
\label{more-etale-definition-compact-support}
设 $X$ 为域 $k$ 上分离且局部有限型的概形，$\mathcal{F}$ 为
$X_\etale$ 上的阿贝尔层。令
$H^0_c(X, \mathcal{F}) \subset H^0(X, \mathcal{F})$
为支撑在 $k$ 上固有的截面之集。$H^0_c(X, \mathcal{F})$ 的元素称为
{\it 紧支撑截面}。
\end{definition}

\noindent
警告：若 $X$ 在 $k$ 上不分离，这个定义并非“正确的”定义。

\begin{lemma}
\label{more-etale-lemma-proper-compact-support}
设 $X$ 为域 $k$ 上的固有概形。则
$H^0_c(X, \mathcal{F}) = H^0(X, \mathcal{F})$。
\end{lemma}

\begin{proof}
由 $H^0_c$ 的构造立即得到。
\end{proof}

\begin{remark}[开嵌入与紧支撑截面]
\label{more-etale-remark-covariance-compact-support}
设 $X$ 为域 $k$ 上分离且局部有限型的概形，$\mathcal{F}$ 为
$X_\etale$ 上的阿贝尔层。完全类似于评注
\ref{more-etale-remark-covariance-f-shriek-separated}，对开子概形
$X' \subset X$ 存在单射
$$
H^0_c(X', \mathcal{F}|_{X'}) \longrightarrow H^0_c(X, \mathcal{F})
$$
这些映射使 $H^0_c$ 成为 $X$ 的 Zariski 位点上的一个“余层”。
\end{remark}

\begin{lemma}
\label{more-etale-lemma-compactify-compact-support}
设 $k$ 为域。设 $j : X \to \overline{X}$ 为 $k$ 上概形的开浸入，且
$\overline{X}$ 在 $k$ 上固有。对 $\mathcal{F} \in \textit{Ab}(X_\etale)$，
存在典范同构（见证明）
$$
H^0_c(X, \mathcal{F}) \longrightarrow
H^0_c(\overline{X}, j_!\mathcal{F}) =
H^0(\overline{X}, j_!\mathcal{F})
$$
其中右端等号由引理 \ref{more-etale-lemma-proper-compact-support} 得到。
\end{lemma}

\begin{proof}
由“\'Etale 上同调”引理
\ref{etale-cohomology-lemma-jshriek-open}，有
$(j_!\mathcal{F})|_X = \mathcal{F}$。因此，所示箭头就是评注
\ref{more-etale-remark-covariance-compact-support} 中对
$\mathcal{G} = j_!\mathcal{F}$ 所得的单射
$H^0_c(X, \mathcal{G}|_X) \to H^0_c(\overline{X}, \mathcal{G})$
。由该映射的显式构造，现在只须证明：若
$s \in H^0(\overline{X}, \mathcal{G})$ 为截面，则 $s$ 的支撑包含在开子概形
$X$ 中。这立即来自如下事实：在 $\overline{X} \setminus X$ 的几何点处，
$\mathcal{G}$ 的茎为零。
\end{proof}

\begin{lemma}
\label{more-etale-lemma-stalk-f-shriek-separated}
设 $f : X \to Y$ 为分离且局部有限型的概形态射，$\mathcal{F}$ 为
$X_\etale$ 上的阿贝尔层。则对任意几何点
$\overline{y} : \Spec(k) \to Y$，存在典范同构
$$
(f_!\mathcal{F})_{\overline{y}}
\longrightarrow
H^0_c(X_{\overline{y}}, \mathcal{F}|_{X_{\overline{y}}})
$$
\end{lemma}

\begin{proof}
回忆
$(f_*\mathcal{F})_{\overline{y}} = \colim f_*\mathcal{F}(V)$，其中余极限取遍
$\overline{y}$ 的 \'etale 邻域 $(V, \overline{v})$。若
$s \in f_*\mathcal{F}(V) = \mathcal{F}(X_V)$，则可将 $s$ 拉回为
$(X_V)_{\overline{v}} = X_{\overline{y}}$ 上的 $\mathcal{F}$ 的截面。
于是得到典范映射
$$
c_{\overline{y}} :
(f_*\mathcal{F})_{\overline{y}}
\longrightarrow 
H^0(X_{\overline{y}}, \mathcal{F}|_{X_{\overline{y}}})
$$
我们断言，此映射在子群 $(f_!\mathcal{F})_{\overline{y}}$ 与
$H^0_c(X_{\overline{y}}, \mathcal{F}|_{X_{\overline{y}}})$ 之间诱导双射。
该断言蕴含本引理，而且稍为精确一些：它说明引理中的等同是将
$\mathcal{F}$ 的截面拉回到 $f$ 的几何纤维所得。

\medskip\noindent
注意，任意元素
$s \in (f_!\mathcal{F})_{\overline{y}} \subset (f_*\mathcal{F})_{\overline{y}}$
都由 $c_{\overline{y}}$ 映到
$H^0_c(X_{\overline{y}}, \mathcal{F}|_{X_{\overline{y}}}) \subset
H^0(X_{\overline{y}}, \mathcal{F}|_{X_{\overline{y}}})$ 中的元素。
这是因为取截面的支撑与拉回可交换，且固有性在基变换下保持。
这至少构造出了引理所述的映射。为证明它是同构，可在 $Y$ 上作
Zariski 局部检验，因此以下设 $Y$ 为仿射概形。

\medskip\noindent
下面要用到如下观察：给定开子概形 $X' \subset X$，并令
$f' = f|_{X'}$，则得到交换图
$$
\xymatrix{
(f'_!(\mathcal{F}|_{X'}))_{\overline{y}} \ar[r] \ar[d] &
H^0_c(X'_{\overline{y}}, \mathcal{F}|_{X'_{\overline{y}}}) \ar[d] \\
(f_!\mathcal{F})_{\overline{y}} \ar[r] &
H^0_c(X_{\overline{y}}, \mathcal{F}|_{X_{\overline{y}}})
}
$$
其中横箭头是上述构造的映射，竖箭头由评注
\ref{more-etale-remark-covariance-f-shriek-separated} 和
\ref{more-etale-remark-covariance-compact-support} 给出。其理由如下：给定
$\overline{y}$ 的 \'etale 邻域 $(V, \overline{v})$ 以及截面
$s \in f_*\mathcal{F}(V) = \mathcal{F}(X_V)$，若其支撑 $Z$ 恰好包含在
$X'_V$ 中且在 $V$ 上固有，则 $s$ 同时给出
$(f'_!(\mathcal{F}|_{X'}))_{\overline{y}}$ 与
$(f_!\mathcal{F})_{\overline{y}}$ 中由图中竖箭头对应的两个元素；这两个元素
经横箭头都映到 $s$ 在 $X_{\overline{y}}$ 上的拉回
$s|_{X_{\overline{y}}}$。后者的支撑 $Z_{\overline{y}}$ 包含在
$X'_{\overline{y}}$ 中，故这一限制给出
$H^0_c(X'_{\overline{y}}, \mathcal{F}|_{X'_{\overline{y}}})$
与 $H^0_c(X_{\overline{y}}, \mathcal{F}|_{X_{\overline{y}}})$
中相容的一对元素。

\medskip\noindent
设 $s \in (f_!\mathcal{F})_{\overline{y}}$ 在
$H^0_c(X_{\overline{y}}, \mathcal{F}|_{X_{\overline{y}}})$ 中映为零。
设 $s$ 由 $s \in f_*\mathcal{F}(V) = \mathcal{F}(X_V)$ 表示，其支撑
$Z$ 在 $V$ 上固有。可设 $V$ 为仿射概形，因而 $Z$ 拟紧。于是可取包含
$Z$ 之像的拟紧开子概形 $X' \subset X$。此时 $Z$ 包含在 $X'_V$ 中，
从而 $s$ 是某个元素 $s' \in f'_!(\mathcal{F}|_{X'})(V)$ 的像，其中
$f' = f|_{X'}$ 如上一段所述。于是 $s'$ 在
$H^0_c(X'_{\overline{y}}, \mathcal{F}|_{X'_{\overline{y}}})$ 中映为零。
因此，为证明单射性，可用 $X'$ 替换 $X$，即可以假设 $X$ 拟紧。
这个情形将在下面证明。

\medskip\noindent
设 $t \in H^0_c(X_{\overline{y}}, \mathcal{F}|_{X_{\overline{y}}})$。
则 $t$ 的支撑包含在某个拟紧开子概形
$W \subset X_{\overline{y}}$ 中。因此，可找到拟紧开子概形
$X' \subset X$，使 $X'_{\overline{y}}$ 包含 $W$。显然，$t$ 属于单射
$H^0_c(X'_{\overline{y}}, \mathcal{F}|_{X'_{\overline{y}}}) \to
H^0_c(X_{\overline{y}}, \mathcal{F}|_{X_{\overline{y}}})$ 的像。
因此，为证明满射性，可用 $X'$ 替换 $X$，即可以假设 $X$ 拟紧。
这个情形将在下面证明。

\medskip\noindent
在证明的最后一段，我们证明 $X$ 拟紧且 $Y$ 仿射的情形。由
“关于平坦性的进一步结果”定理 \ref{flat-theorem-nagata}，存在 $Y$ 上的紧化
$j : X \to \overline{X}$。令 $\mathcal{G} = j_!\mathcal{F}$；由
“\'Etale 上同调”引理 \ref{etale-cohomology-lemma-jshriek-open}，有
$\mathcal{F} = \mathcal{G}|_X$。由上述讨论得到交换图
$$
\xymatrix{
(f_!\mathcal{F})_{\overline{y}} \ar[r] \ar[d] &
H^0_c(X_{\overline{y}}, \mathcal{F}|_{X_{\overline{y}}}) \ar[d] \\
(\overline{f}_!\mathcal{G})_{\overline{y}} \ar[r] &
H^0_c(\overline{X}_{\overline{y}}, \mathcal{G}|_{\overline{X}_{\overline{y}}})
}
$$
由引理 \ref{more-etale-lemma-compactify-f-shriek-separated} 和
\ref{more-etale-lemma-compactify-compact-support}，竖箭头均为同构。这把问题化归到固有态射
$\overline{X} \to Y$ 的情形。对固有态射，由引理
\ref{more-etale-lemma-proper-f-shriek}、\ref{more-etale-lemma-proper-compact-support} 以及直像的
固有基变换，我们的映射是同构；见“\'Etale 上同调”引理
\ref{etale-cohomology-lemma-proper-pushforward-stalk}。
\end{proof}

\begin{lemma}
\label{more-etale-lemma-base-change-f-shriek-separated}
考虑概形的笛卡尔方块
$$
\xymatrix{
X' \ar[r]_{g'} \ar[d]_{f'} & X \ar[d]^f \\
Y' \ar[r]^g & Y
}
$$
，其中 $f$ 分离且局部有限型。则对 $X_\etale$ 上任意阿贝尔层
$\mathcal{F}$，有
$f'_!(g')^{-1}\mathcal{F} = g^{-1}f_!\mathcal{F}$.
\end{lemma}

\begin{proof}
在非常一般的情形下，存在拉回映射
$g^{-1}f_*\mathcal{F} \to f'_*(g')^{-1}\mathcal{F}$；见“位点”第
\ref{sites-section-pullback} 节。我们断言，此映射把
$g^{-1}f_!\mathcal{F}$ 映入子层 $f'_!(g')^{-1}\mathcal{F}$，并诱导引理中的同构。

\medskip\noindent
取几何点 $\overline{y}': \Spec(k) \to Y'$，并以
$\overline{y} = g \circ \overline{y}'$ 表示其在 $Y$ 中的像。存在交换图
$$
\xymatrix{
(f_*\mathcal{F})_{\overline{y}} \ar[r] \ar[d] &
H^0(X_{\overline{y}}, \mathcal{F}|_{X_{\overline{y}}}) \ar[d] \\
(f'_*(g')^{-1}\mathcal{F})_{\overline{y}'} \ar[r] &
H^0(X'_{\overline{y}'}, (g')^{-1}\mathcal{F}|_{X'_{\overline{y}'}})
}
$$
其中横映射用于引理 \ref{more-etale-lemma-stalk-f-shriek-separated} 的证明，
竖映射是上述拉回映射。此图交换，因为所涉及的四个映射都是沿概形态射拉回
局部截面而得，且相应的概形态射图交换。由于引理所述方块为笛卡尔方块，
有 $X'_{\overline{y}'} = X_{\overline{y}}$。因此，由引理
\ref{more-etale-lemma-stalk-f-shriek-separated} 及其证明得到交换图
$$
\xymatrix{
(f_*\mathcal{F})_{\overline{y}} \ar[rrr] \ar[ddd] & & &
H^0(X_{\overline{y}}, \mathcal{F}|_{X_{\overline{y}}}) \ar[ddd] \\
& (f_!\mathcal{F})_{\overline{y}} \ar[r] \ar@{..>}[d] \ar[lu] &
H^0_c(X_{\overline{y}}, \mathcal{F}|_{X_{\overline{y}}}) \ar[d] \ar[ru] \\
& (f'_!(g')^{-1}\mathcal{F})_{\overline{y}'} \ar[r] \ar[ld] &
H^0_c(X'_{\overline{y}'}, (g')^{-1}\mathcal{F}|_{X'_{\overline{y}'}}) \ar[rd]\\
(f'_*(g')^{-1}\mathcal{F})_{\overline{y}'} \ar[rrr] & & &
H^0(X'_{\overline{y}'}, (g')^{-1}\mathcal{F}|_{X'_{\overline{y}'}})
}
$$
其中内方块的横箭头均为同构，右侧两条竖箭头均为恒等。另有，
东南、西南、东北、西北方向的箭头均为单射。因此，存在唯一的双射虚线箭头
使图交换。我们得出
$g^{-1}f_!\mathcal{F} \subset g^{-1}f_*\mathcal{F} \to f'_*(g')^{-1}\mathcal{F}$
的像包含在子层
$f'_!(g')^{-1}\mathcal{F} \subset f'_*(g')^{-1}\mathcal{F}$
中，因为这在茎上成立；见“\'Etale 上同调”定理
\ref{etale-cohomology-theorem-exactness-stalks}。同一定理继而说明所诱导的映射
为同构，证明完毕。
\end{proof}

\begin{lemma}
\label{more-etale-lemma-f-shriek-composition}
设 $f : X \to Y$ 与 $g : Y \to Z$ 为可复合的概形态射，且二者均分离并
局部有限型。设 $\mathcal{F}$ 为 $X_\etale$ 上的阿贝尔层。则作为
$(g \circ f)_*\mathcal{F}$ 的子层，有
$g_!f_!\mathcal{F} = (g \circ f)_!\mathcal{F}$。
\end{lemma}

\begin{proof}
我们强烈建议读者自行证明此结论。设 $W \in Z_\etale$，且
$s \in (g \circ f)_*\mathcal{F}(W) = \mathcal{F}(X_W)$。以
$T \subset X_W$ 表示 $s$ 的支撑；这是闭子集。注意，$s$ 是
$(g \circ f)_!\mathcal{F}$ 的截面，当且仅当 $T$ 在 $W$ 上固有。我们有
$f_!\mathcal{F} \subset f_*\mathcal{F}$，因而
$g_!f_!\mathcal{F} \subset g_!f_*\mathcal{F} \subset g_*f_*\mathcal{F}$.
另一方面，$s$ 是 $g_!f_!\mathcal{F}$ 的截面，当且仅当：(a) $T$ 在
$Y_W$ 上固有；(b) 将 $s$ 看作 $f_!\mathcal{F}$ 的截面时，其支撑
$T'$ 在 $W$ 上固有。若 (a) 成立，则 $T$ 在 $Y_W$ 中的像为闭集；又因
$f_!\mathcal{F} \subset f_*\mathcal{F}$，可知 $T' \subset Y_W$ 正是
$T$ 的像（略去细节；可考察茎）。

\medskip\noindent
归结起来，我们须证明：闭子集 $T \subset X_W$ 在 $W$ 上固有，当且仅当
$T$ 在 $Y_W$ 上固有且 $T$ 在 $Y_W$ 中的像在 $W$ 上固有。赋予 $T$
诱导的既约闭子概形结构。若 $T$ 在 $W$ 上固有，则由“态射”引理
\ref{morphisms-lemma-image-proper-scheme-closed}，$T \to Y_W$ 固有；
并由“概形上同调”引理
\ref{coherent-lemma-functoriality-closed-proper-over-base}，$T$ 在 $Y_W$
中的像在 $W$ 上固有。反之，若 $T$ 在 $Y_W$ 上固有，且 $T$ 在 $Y_W$
中的像在 $W$ 上固有，则态射 $T \to W$ 是固有态射的复合，因而固有
（这里给 $T$ 在 $Y_W$ 中的闭像赋予诱导的既约概形结构，从而把问题化为
概形态射的问题）；见“态射”引理 \ref{morphisms-lemma-composition-proper}。
\end{proof}

\begin{remark}
\label{more-etale-remark-f-shriek-base-change-composition}
上述构造的函子间同构满足以下两个性质：
\begin{enumerate}
\item 设 $f : X \to Y$、$g : Y \to Z$ 和 $h : Z \to T$ 为可复合的概形态射，
且均分离并局部有限型。则图
$$
\xymatrix{
(h \circ g \circ f)_! \ar[r] \ar[d] &
(h \circ g)_! \circ f_! \ar[d] \\
h_! \circ (g \circ f)_! \ar[r] &
h_! \circ g_! \circ f_!
}
$$
交换，其中箭头来自引理 \ref{more-etale-lemma-f-shriek-composition}。
\item 假设有概形图
$$
\xymatrix{
X' \ar[d]_{f'} \ar[r]_c & X \ar[d]^f \\
Y' \ar[d]_{g'} \ar[r]_b & Y \ar[d]^g \\
Z' \ar[r]^a & Z
}
$$
，其中两个方块均为笛卡尔方块，且 $f$ 与 $g$ 分离并局部有限型。则图
$$
\xymatrix{
a^{-1} \circ (g \circ f)_! \ar[d] \ar[rr] & &
(g' \circ f')_! \circ c^{-1} \ar[d] \\
a^{-1} \circ g_! \circ f_! \ar[r] &
g'_! \circ b^{-1} \circ f_! \ar[r] &
g'_! \circ f'_! \circ c^{-1}
}
$$
交换，其中横箭头来自引理 \ref{more-etale-lemma-base-change-f-shriek-separated}，
箭头来自引理 \ref{more-etale-lemma-f-shriek-composition}。
\end{enumerate}
(1) 成立是因为直像有类似的交换图。(2) 成立是因为直像的基变换映射具有
非常一般的相容性（“位点”评注 \ref{sites-remark-compose-base-change}），
且引理 \ref{more-etale-lemma-base-change-f-shriek-separated} 与
\ref{more-etale-lemma-f-shriek-composition} 中的同构均由相应的直像映射构造。
\end{remark}

\begin{lemma}
\label{more-etale-lemma-colim-f-shriek-separated}
设 $f : X \to Y$ 为分离且局部有限型的概形态射。设
$X = \bigcup_{i \in I} X_i$ 为开覆盖，且对任意 $i, j \in I$，都存在
$k$ 使 $X_i \cup X_j \subset X_k$。以 $f_i : X_i \to Y$ 表示 $f$
的限制。则
$$
f_!\mathcal{F} = \colim_{i \in I} f_{i, !}(\mathcal{F}|_{X_i})
$$
关于 $\mathcal{F} \in \textit{Ab}(X_\etale)$ 函子性地成立，其中转移映射
由评注 \ref{more-etale-remark-covariance-f-shriek-separated} 构造。
\end{lemma}

\begin{proof}
只须证明，从右到左的典范映射在 $Y_\etale$ 的拟紧对象 $V$ 上取值时为双射。
注意，右端余极限是有向的，且转移映射均为单射。因此可用“位点”引理
\ref{sites-lemma-directed-colimits-sections} 计算该余极限。于是结论归结为如下观察：
在 $V$ 上固有的闭子集 $Z \subset X_V$ 是拟紧的，故对某个 $i$，
它包含在 $X_{i, V}$ 中。
\end{proof}

\begin{lemma}
\label{more-etale-lemma-f-shriek-separated-direct-sums}
设 $f : X \to Y$ 为分离且局部有限型的概形态射。则函子 $f_!$ 与直和可交换。
\end{lemma}

\begin{proof}
令 $\mathcal{F} = \bigoplus \mathcal{F}_i$。为证明映射
$\bigoplus f_!\mathcal{F}_i \to f_!\mathcal{F}$ 为同构，只须证明这些层在
$Y_\etale$ 的任意拟紧对象 $V$ 上具有相同的截面。用 $V$ 替换 $Y$ 后，
只须证明
$H^0(Y, f_!\mathcal{F}) \subset H^0(X, \mathcal{F})$
等于
$\bigoplus H^0(Y, f_!\mathcal{F}_i)
\subset \bigoplus H^0(X, \mathcal{F}_i)
\subset H^0(X, \bigoplus \mathcal{F}_i)$.
把 $X$ 写成其拟紧开子概形的并，并应用引理
\ref{more-etale-lemma-colim-f-shriek-separated}，便可化归到 $X$ 也拟紧的情形。
此时，由“\'Etale 上同调”定理 \ref{etale-cohomology-theorem-colimit}，有
$H^0(X, \mathcal{F}) = \bigoplus H^0(X, \mathcal{F}_i)$。
考察截面的支撑即可容易得出结论。
\end{proof}

\begin{lemma}
\label{more-etale-lemma-lqf-f-shriek-separated-colimits}
设 $f : X \to Y$ 为分离且局部拟有限的概形态射。则
\begin{enumerate}
\item 对 $\textit{Ab}(X_\etale)$ 中的 $\mathcal{F}$ 及几何点
$\overline{y} : \Spec(k) \to Y$，有
$$
(f_!\mathcal{F})_{\overline{y}} =
\bigoplus\nolimits_{f(\overline{x}) = \overline{y}} \mathcal{F}_{\overline{x}}
$$
关于 $\mathcal{F}$ 函子性地成立；
\item 函子 $f_!$ 是正合的。
\end{enumerate}
\end{lemma}

\begin{proof}
由构造，函子 $f_!$ 左正合。右正合性可在茎上检验（“\'Etale 上同调”定理
\ref{etale-cohomology-theorem-exactness-stalks}）。因此只须证明 (1)。

\medskip\noindent
设 $\overline{y} : \Spec(k) \to Y$ 为几何点。概形 $X_{\overline{y}}$
的底拓扑空间是离散的（“态射”引理
\ref{morphisms-lemma-locally-quasi-finite-fibres}），且各点处的剩余域都等于
$k$（它们是 $k$ 的有限扩张）。因此，集合
$\{\overline{x} : \Spec(k) \to X : f(\overline{x}) = \overline{y}\}$
等于 $X_{\overline{y}}$ 的点集。于是茎的计算由更一般的引理
\ref{more-etale-lemma-stalk-f-shriek-separated} 得出。
\end{proof}









\section{有限支撑截面}
\label{more-etale-section-finite-support}

\noindent
本节把第 \ref{more-etale-section-compact-support} 节的构造推广到未必分离的局部拟有限态射。

\medskip\noindent
设 $f : X \to Y$ 为局部拟有限的概形态射，$\mathcal{F}$ 为
$X_\etale$ 上的阿贝尔层。给定 $Y_\etale$ 中的 $V$，以
$X_V = X \times_Y V$ 表示基变换。我们考虑有限形式和所成的群
\begin{equation}
\label{more-etale-equation-formal-sum}
s = \sum\nolimits_{i = 1, \ldots, n} (Z_i, s_i)
\end{equation}
，其中 $Z_i \subset X_V$ 为局部闭子概形，且态射 $Z_i \to V$ 有限
\footnote{由于 $f$ 局部拟有限，态射 $Z_i \to V$ 有限当且仅当它固有。}，
并且 $s_i \in H_{Z_i}(\mathcal{F})$。这里与第 \ref{more-etale-section-growing} 节一样，令
$$
H_{Z_i}(\mathcal{F}) =
\{s_i \in \mathcal{F}(U_i) \mid \text{Supp}(s_i) \subset Z_i\}
$$
，其中 $U_i \subset X_V$ 为把 $Z_i$ 包含为闭子概形的开子概形。我们对这些
形式和模去以下关系：
\begin{enumerate}
\item
\label{more-etale-item-sum}
$(Z, s) + (Z, s') = (Z, s + s')$,
\item
\label{more-etale-item-sub}
$(Z, s) = (Z', s)$（当 $Z \subset Z'$ 时）。
\end{enumerate}
注意，第二个关系确有意义：由于 $Z \to V$ 有限而 $Z' \to V$ 分离，
包含映射 $Z \to Z'$ 为闭浸入，故可使用 (\ref{more-etale-item-inclusion}) 中讨论的映射。

\medskip\noindent
以 $f_{p!}\mathcal{F}(V)$ 表示形式和 (\ref{more-etale-equation-formal-sum}) 所成阿贝尔群
模去这些关系的商。第一个关系说明，$f_{p!}\mathcal{F}(V)$ 是所有在 $V$
上有限的局部闭子概形 $Z \subset X_V$ 所对应的阿贝尔群
$H_Z(\mathcal{F})$ 之直和的商。第二个关系说明，我们实际上取的是余极限
\begin{equation}
\label{more-etale-equation-colimit-definition}
f_{p!}\mathcal{F}(V) = \colim_Z H_Z(\mathcal{F})
\end{equation}
这个公式为理解该构造提供了方便的抽象方式。

\medskip\noindent
接着注意，该构造可自然地成为 $Y_\etale$ 上的阿贝尔群预层
$f_{p!}\mathcal{F}$。具体地，给定 $Y_\etale$ 中的 $V' \to V$，得到基变换态射
$X_{V'} \to X_V$。若 $Z \subset X_V$ 是在 $V$ 上有限的局部闭子概形，则其
概形论逆像 $Z' \subset X_{V'}$ 在 $V'$ 上有限。此外，若 $U \subset X_V$
为开子概形且 $Z$ 在 $U$ 中闭，则逆像 $U' \subset X_{V'}$ 为开子概形，
且 $Z'$ 在 $U'$ 中闭。因此，$\mathcal{F}$ 的限制映射
$\mathcal{F}(U) \to \mathcal{F}(U')$ 把 $H_Z(\mathcal{F})$ 映入
$H_{Z'}(\mathcal{F})$；这是上文 (\ref{more-etale-item-pullback}) 中所述函子性的特例。
显然，这些映射与 $X_V$ 中此类局部闭子概形的包含
$Z_1 \subset Z_2$ 相容，于是得到映射
$$
f_{p!}\mathcal{F}(V) = \colim_Z H_Z(\mathcal{F})
\longrightarrow
\colim_{Z'} H_{Z'}(\mathcal{F}) =
f_{p!}\mathcal{F}(V')
$$
这些映射确实使 $f_{p!}\mathcal{F}$ 成为 $Y_\etale$ 上的阿贝尔群预层。
略去细节。

\medskip\noindent
最后注意，$f_{p!}\mathcal{F}$ 的构造关于
$\textit{Ab}(X_\etale)$ 中的 $\mathcal{F}$ 具有函子性。因此，给定局部拟有限态射
$f : X \to Y$，我们构造了函子
$$
f_{p!} :
\textit{Ab}(X_\etale)
\longrightarrow
\textit{PAb}(Y_\etale)
$$
，从 $X_\etale$ 上阿贝尔层的范畴到 $Y_\etale$ 上阿贝尔预层的范畴。
在把 $f_!$ 定义为此函子的层化之前，先检验在两种构造均适用时，它与
第 \ref{more-etale-section-compact-support} 节的构造以及“\'Etale 上同调”第
\ref{etale-cohomology-section-extension-by-zero} 节的构造一致。

\begin{lemma}
\label{more-etale-lemma-finite-support-f-shriek-separated}
设 $f : X \to Y$ 为分离且局部拟有限的概形态射。关于
$\mathcal{F} \in \textit{Ab}(X_\etale)$ 函子性地存在阿贝尔预层的典范同构（！）
$$
f_{p!}\mathcal{F} \longrightarrow f_!\mathcal{F}
$$
；它把定义 \ref{more-etale-definition-f-shriek-separated} 中的层
$f_!\mathcal{F}$ 与上述构造的预层 $f_{p!}\mathcal{F}$ 等同起来。
\end{lemma}

\begin{proof}
设 $V$ 为 $Y_\etale$ 的对象。若 $Z \subset X_V$ 局部闭且在 $V$ 上有限，
则因 $f$ 分离，态射 $Z \to X_V$ 为闭浸入。此外，若 $Z_i$（其中
$i = 1, \ldots, n$）为 $X_V$ 中在 $V$ 上有限的闭子概形，则
$Z_1 \cup \ldots \cup Z_n$（概形论并）是一个在 $V$ 上有限的闭子概形。
因此在此情形，定义 $f_{p!}\mathcal{F}(V)$ 的余极限
(\ref{more-etale-equation-colimit-definition}) 是有向的，并且
$f_{!p}\mathcal{F}(V)$ 恰等于 $\mathcal{F}(X_V)$ 中支撑在 $V$ 上有限的截面之集。
由于 $X_V$ 中在 $V$ 上固有的任意闭子集事实上都在 $V$ 上有限
（因为 $f$ 局部拟有限），由定义即知该集合等于 $f_!\mathcal{F}(V)$。
\end{proof}

\begin{lemma}
\label{more-etale-lemma-finite-support-stalk}
设 $f : X \to Y$ 为局部拟有限的概形态射，
$\overline{y} : \Spec(k) \to Y$ 为几何点。关于
$\textit{Ab}(X_\etale)$ 中的 $\mathcal{F}$ 函子性地有
$$
(f_{p!}\mathcal{F})_{\overline{y}} =
\bigoplus\nolimits_{f(\overline{x}) = \overline{y}} \mathcal{F}_{\overline{x}}
$$
\end{lemma}

\begin{proof}
回忆，预层在 $\overline{y}$ 处的茎定义为取遍 $\overline{y}$ 的 \'etale 邻域
$(V, \overline{v})$ 的通常余极限；见“\'Etale 上同调”定义
\ref{etale-cohomology-definition-stalk}。相应地，设
$s = \sum_{i = 1, \ldots, n} (Z_i, s_i)$ 如
(\ref{more-etale-equation-formal-sum}) 所示，是 $f_{p!}\mathcal{F}(V)$ 的元素，其中
$(V, \overline{v})$ 为 $\overline{y}$ 的 \'etale 邻域。由于
$$
X_{\overline{y}} = (X_V)_{\overline{v}} \supset Z_{i, \overline{v}}
$$
，且 $s_i$ 是 $X_V$ 中 $Z_i$ 的某个开邻域上的 $\mathcal{F}$ 截面，
故可把 $s$ 映到
$$
\sum\nolimits_{i = 1, \ldots, n} 
\sum\nolimits_{\overline{x} \in Z_{i, \overline{v}}}
\left(\text{取 }s_i\text{ 的类，位于 }\mathcal{F}_{\overline{x}}\right)
\quad\in\quad
\bigoplus\nolimits_{f(\overline{x}) = \overline{y}} \mathcal{F}_{\overline{x}}
$$
略去如下验证：该构造与限制映射相容；关系 (\ref{more-etale-item-sum})
$(Z, s) + (Z, s') - (Z, s + s')$ 以及当 $Z \subset Z'$ 时的关系
(\ref{more-etale-item-sub}) $(Z, s) - (Z', s)$ 均映为零。于是得到映射
$$
(f_{p!}\mathcal{F})_{\overline{y}}
\longrightarrow
\bigoplus\nolimits_{f(\overline{x}) = \overline{y}} \mathcal{F}_{\overline{x}}
$$

\medskip\noindent
先证明此箭头满射。只须取满足 $f(\overline{x}) = \overline{y}$ 的
$\overline{x}$，并证明直和项 $\mathcal{F}_{\overline{x}}$ 中的任意元素 $s$
属于其像。设 $s$ 由元素 $s \in \mathcal{F}(U)$ 表示，其中
$(U, \overline{u})$ 是 $\overline{x}$ 的 \'etale 邻域。由于 $f$ 局部拟有限，
态射 $U \to Y$ 也局部拟有限。由“关于态射的进一步结果”引理
\ref{more-morphisms-lemma-etale-makes-quasi-finite-finite-multiple-points-var}
，可以找到 $\overline{y}$ 的 \'etale 邻域 $(V, \overline{v})$、开子概形
$$
W \subset U \times_Y V,
$$
以及映到 $\overline{u}$ 和 $\overline{v}$ 的几何点 $\overline{w}$，使得
$W \to V$ 有限，且 $\overline{w}$ 是 $W$ 中唯一映到 $\overline{v}$ 的几何点。
（这里略去在当前使用的几何点语言与该引理陈述所用的点及剩余域扩张语言之间
作转换。）注意，$W \to X_V = X \times_Y V$ 是 \'etale 的。取
$\overline{w}$ 的像 $\overline{w}'$ 的仿射开邻域
$W' \subset X_V$。由于 $\overline{w}$ 是 $W$ 中位于 $\overline{v}$ 上方的
唯一点，且 $W \to V$ 为闭映射，用 $\overline{v}$ 的一个开邻域替换 $V$ 后，
可设 $W \to X_V$ 的像包含在 $W'$ 中。于是 $W \to W'$ 有限且 \'etale，
并且 $W$ 中恰有一个位于 $\overline{w}'$ 上方的几何点 $\overline{w}$。
因此，在 $W'$ 中 $\overline{w}'$ 的某个开邻域上，$W \to W'$ 是开浸入；
见“\'Etale 态射”引理 \ref{etale-lemma-finite-etale-one-point}。缩小 $V$ 与
$W'$ 后，可设 $W \to W'$ 为同构。于是可把 $s$ 看作 $\mathcal{F}$ 在
$W' \subset X_V$ 上的截面 $s'$，且该开子概形在 $V$ 上有限。因此按定义，
$(W', s')$ 给出 $j_{p!}\mathcal{F}(V)$ 的一个元素，并按所需映到 $s$。

\medskip\noindent
再证明此箭头单射。设
$s = \sum_{i = 1, \ldots, n} (Z_i, s_i)$ 如
(\ref{more-etale-equation-formal-sum}) 所示，是 $f_{p!}\mathcal{F}(V)$ 的元素，其中
$(V, \overline{v})$ 为 $\overline{y}$ 的 \'etale 邻域。假设 $s$ 在上述构造的
映射下映为零。首先，用 $(V, \overline{v})$ 自身的一个 \'etale 邻域替换它后，
可设存在分解
$Z_i = Z_{i, 1} \amalg \ldots \amalg Z_{i, m_i}$，其中各项均为开闭
子概形，使每个 $Z_{i, j}$ 在 $\overline{v}$ 上方恰有一个几何点。在显然的直和分解
$$
H_{Z_i}(\mathcal{F}) = \bigoplus H_{Z_{i, j}}(\mathcal{F})
$$
下，设元素 $s_i$ 对应于 $\sum s_{i, j}$。利用关系 (\ref{more-etale-item-sum}) 与
(\ref{more-etale-item-sub})，可把 $s$ 替换为
$\sum_{i = 1, \ldots, n} \sum_{j = 1, \ldots, m_i} (Z_{i, j}, s_{i, j})$.
换言之，可设 $Z_i$ 在 $\overline{v}$ 上方恰有一个几何点。令
$\overline{x}_1, \ldots, \overline{x}_m$ 为 $X$ 中位于 $\overline{y}$ 上方、
与各 $Z_i$ 中位于 $\overline{v}$ 上方的几何点相对应的几何点；注意，对某个
$j \in \{1, \ldots, m\}$，可能有多个指标 $i$，使 $\overline{x}_j$
对应于 $Z_i$ 的一个点。对 $X_V \to V$ 应用“关于态射的进一步结果”引理
\ref{more-morphisms-lemma-etale-makes-quasi-finite-finite-multiple-points-var}
，并用 $(V, \overline{v})$ 自身的一个 \'etale 邻域替换它后，可设存在开子概形
$$
W_j \subset X \times_Y V,\quad j = 1, \ldots, m
$$
以及 $W_j$ 中映到 $\overline{x}_j$ 和 $\overline{v}$ 的几何点
$\overline{w}_j$，使得 $W_j \to V$ 有限，且 $\overline{w}_j$ 是 $W_j$ 中
唯一映到 $\overline{v}$ 的几何点。缩小 $V$ 后，可设对某个 $j$ 有
$Z_i \subset W_j$，并且有映射
$H_{Z_i}(\mathcal{F}) \to H_{W_j}(\mathcal{F})$。因此，由关系
(\ref{more-etale-item-sub})，我们的元素等价于形如
$$
\sum\nolimits_{j = 1, \ldots, m} (W_j, t_j)
$$
的元素，其中 $t_j \in H_{W_j}(\mathcal{F})$。显然，此元素恰映到直和项
$\mathcal{F}_{\overline{x}_j}$ 中 $t_j$ 的类。由于 $s$ 映为零，$t_j$ 在
$\mathcal{F}_{\overline{x}_j}$ 中映为零。这说明 $t_j$ 在 $W_j$ 中
$\overline{w}_j$ 的某个开邻域上的限制为零；见“\'Etale 上同调”引理
\ref{etale-cohomology-lemma-zero-over-image}。再次缩小 $V$，便按所需得到
对所有 $j$ 都有 $t_j = 0$。
\end{proof}

\begin{lemma}
\label{more-etale-lemma-finite-support-etale-shriek}
设 $f = j : U \to X$ 为概形的 \'etale 态射。以 $j_{p!}$ 表示
“\'Etale 上同调”公式 (\ref{etale-cohomology-equation-j-p-shriek})
中的构造，以 $f_{p!}$ 表示上述构造。关于
$\mathcal{F} \in \textit{Ab}(X_\etale)$ 函子性地存在阿贝尔预层的典范映射
$$
j_{p!}\mathcal{F} \longrightarrow f_{p!}\mathcal{F}
$$
；它把“\'Etale 上同调”定义
\ref{etale-cohomology-definition-extension-zero} 中的层
$j_!\mathcal{F} = (j_{p!}\mathcal{F})^\#$ 与 $(f_{p!}\mathcal{F})^\#$ 等同起来。
\end{lemma}

\begin{proof}
阅读本引理的证明前，请先阅读“\'Etale 上同调”引理
\ref{etale-cohomology-lemma-shriek-into-star-separated-etale} 的证明。
设 $V$ 为 $X_\etale$ 的对象。回忆
$$
j_{p!}\mathcal{F}(V) =
\bigoplus\nolimits_{\varphi : V \to U} \mathcal{F}(V \xrightarrow{\varphi} U)
$$
。给定 $\varphi$，得到开子概形 $Z_\varphi \subset U_V = U \times_X V$，
即 $\varphi$ 的图的像。通过 $\varphi$，得到 $U$ 上的同构
$V \to Z_\varphi$，并可把元素
$$
s_\varphi \in \mathcal{F}(V \xrightarrow{\varphi} U) =
\mathcal{F}(Z_\varphi) = H_{Z_\varphi}(\mathcal{F})
$$
看作 $\mathcal{F}$ 在 $Z_{\varphi}$ 上的截面。由于
$Z_\varphi \subset U_V$ 为开子概形，事实上有
$H_{Z_\varphi}(\mathcal{F}) = \mathcal{F}(Z_\varphi)$，故可把
$s_\varphi$ 看作 $H_{Z_\varphi}(\mathcal{F})$ 的元素。这样，映射
$j_{p!}\mathcal{F} \to f_{p!}\mathcal{F}$ 由规则
$$
\sum\nolimits_{i = 1, \ldots, n} s_{\varphi_i}
\longmapsto
\sum\nolimits_{i = 1, \ldots, n} (Z_{\varphi_i}, s_{\varphi_i})
$$
定义，其中右端是 (\ref{more-etale-equation-formal-sum}) 中的形式和。略去它与限制映射相容
且关于 $\mathcal{F}$ 具有函子性的验证。

\medskip\noindent
为完成证明，我们断言，对给定的几何点
$\overline{y} : \Spec(k) \to Y$，存在交换图
$$
\xymatrix{
(j_{p!}\mathcal{F})_{\overline{y}} \ar[r] \ar[d] &
\bigoplus_{j(\overline{x}) = \overline{y}} \mathcal{F}_{\overline{x}}
\ar@{=}[d] \\
(f_{p!}\mathcal{F})_{\overline{y}} \ar[r] &
\bigoplus_{f(\overline{x}) = \overline{y}} \mathcal{F}_{\overline{x}}
}
$$
其中上方横箭头构造于“\'Etale 上同调”命题
\ref{etale-cohomology-proposition-describe-jshriek} 的证明中，下方横箭头构造于
引理 \ref{more-etale-lemma-finite-support-stalk} 的证明中，右侧竖箭头是显然的等号，
左侧竖箭头是上一段所定义映射在茎上的映射。根据这里及所引文献中涉及的
全部箭头的显式描述，该断言直接成立。由于横箭头均为同构，左侧竖箭头也为同构。
因此，由“\'Etale 上同调”定理
\ref{etale-cohomology-theorem-exactness-stalks}，我们的映射在层化上诱导同构。
\end{proof}

\begin{definition}
\label{more-etale-definition-f-shriek-lqf}
设 $f : X \to Y$ 为局部拟有限的概形态射。定义{\it 紧支撑直像}为函子
$$
f_! : \textit{Ab}(X_\etale) \longrightarrow \textit{Ab}(Y_\etale)
$$
，其定义公式为 $f_!\mathcal{F} = (f_{p!}\mathcal{F})^\#$；换言之，
$f_!\mathcal{F}$ 是上述构造的预层 $f_{p!}\mathcal{F}$ 的层化。
\end{definition}

\noindent
由引理 \ref{more-etale-lemma-finite-support-f-shriek-separated}，它与定义
\ref{more-etale-definition-f-shriek-separated} 不冲突（当两个定义均适用时）；由引理
\ref{more-etale-lemma-finite-support-etale-shriek}，它与“\'Etale 上同调”定义
\ref{etale-cohomology-definition-extension-zero} 也不冲突（当两个定义均适用时）。

\begin{lemma}
\label{more-etale-lemma-lqf-f-shriek-stalk}
设 $f : X \to Y$ 为局部拟有限的概形态射。则
\begin{enumerate}
\item 对 $\textit{Ab}(X_\etale)$ 中的 $\mathcal{F}$ 及几何点
$\overline{y} : \Spec(k) \to Y$，有
$$
(f_!\mathcal{F})_{\overline{y}} =
\bigoplus\nolimits_{f(\overline{x}) = \overline{y}} \mathcal{F}_{\overline{x}}
$$
关于 $\mathcal{F}$ 函子性地成立；
\item 函子 $f_! : \textit{Ab}(X_\etale) \to \textit{Ab}(Y_\etale)$
正合且与直和可交换。
\end{enumerate}
\end{lemma}

\begin{proof}
茎的公式由引理 \ref{more-etale-lemma-finite-support-stalk} 立即得到（事实上二者等价）。
再结合正合性可在茎上检验这一事实，便立即得到函子的正合性；见
“\'Etale 上同调”定理 \ref{etale-cohomology-theorem-exactness-stalks}。
\end{proof}

\begin{remark}[关于开嵌入的协变性]
\label{more-etale-remark-covariance-lqf-f-shriek}
设 $f : X \to Y$ 为局部拟有限的概形态射，$\mathcal{F}$ 为
$X_\etale$ 上的阿贝尔层。设 $X' \subset X$ 为开子概形，并以
$f' : X' \to Y$ 表示 $f$ 的限制。我们断言存在典范映射
$$
f'_!(\mathcal{F}|_{X'}) \longrightarrow f_!\mathcal{F}
$$
。具体地，此映射是典范映射
$$
f'_{p!}(\mathcal{F}|_{X'}) \to f_{p!}\mathcal{F}
$$
的层化；后者构造如下。设 $V \in Y_\etale$，并考虑如
(\ref{more-etale-equation-formal-sum}) 所示的截面
$s' = \sum_{i = 1, \ldots, n} (Z'_i, s'_i)$，它定义了
$f'_{p!}(\mathcal{F}|_{X'})(V)$ 的一个元素。此时
$Z'_i \subset X'_V$ 也可看作 $X_V$ 的局部闭子概形，且有
$H_{Z'_i}(\mathcal{F}|_{X'}) = H_{Z'_i}(\mathcal{F})$。把 $s'$ 映到完全相同的和
$s = \sum_{i = 1, \ldots, n} (Z'_i, s'_i)$
，但现在将它看作 $f_{p!}\mathcal{F}(V)$ 的元素。略去该构造与限制映射相容且
关于 $\mathcal{F}$ 具有函子性的验证。此构造具有以下性质：
\begin{enumerate}
\item 映射 $f'_{p!}\mathcal{F}' \to f_{p!}\mathcal{F}$ 与
$f'_!\mathcal{F}' \to f_!\mathcal{F}$ 同引理
\ref{more-etale-lemma-finite-support-stalk} 和 \ref{more-etale-lemma-lqf-f-shriek-stalk}
给出的茎的描述相容。
\item 若 $f$ 分离，则映射
$f'_{p!}\mathcal{F}' \to f_{p!}\mathcal{F}$ 与评注
\ref{more-etale-remark-covariance-f-shriek-separated} 中构造的映射相同；这里使用引理
\ref{more-etale-lemma-finite-support-f-shriek-separated} 中的同构。
\item 若 $X'' \subset X'$ 是另一个开子概形，则复合
$f''_{p!}(\mathcal{F}|_{X''}) \to f'_{p!}(\mathcal{F}|_{X'}) \to
f_{p!}\mathcal{F}$ 是包含 $X'' \subset X$ 所对应的映射
$f''_{p!}(\mathcal{F}|_{X''}) \to f_{p!}\mathcal{F}$。层化后可知，
对于
$f''_!(\mathcal{F}|_{X''}) \to f'_!(\mathcal{F}|_{X'}) \to f_!\mathcal{F}$.
也有同样的结论。
\item 映射 $f'_!\mathcal{F}' \to f_!\mathcal{F}$ 为单射，因为这可在茎上检验。
\end{enumerate}
把元素表示为上述有限和并考察这些元素的像，即可容易证明所有这些陈述。
\end{remark}

\begin{lemma}
\label{more-etale-lemma-lqf-colimit-f-shriek}
设 $f : X \to Y$ 为局部拟有限的概形态射，
$X = \bigcup_{i \in I} X_i$ 为开覆盖。则存在正合复形
$$
\ldots \to
\bigoplus\nolimits_{i_0, i_1, i_2} f_{i_0i_1i_2, !}
\mathcal{F}|_{X_{i_0i_1i_2}} \to
\bigoplus\nolimits_{i_0, i_1} f_{i_0i_1, !} \mathcal{F}|_{X_{i_0i_1}} \to
\bigoplus\nolimits_{i_0} f_{i_0, !} \mathcal{F}|_{X_{i_0}}
\to f_!\mathcal{F} \to 0
$$
，它关于 $\mathcal{F} \in \textit{Ab}(X_\etale)$ 函子性地成立；细节见证明。
\end{lemma}

\begin{proof}
这里照例令 $X_{i_0 \ldots i_p} = X_{i_0} \cap \ldots \cap X_{i_p}$，
并以 $f_{i_0 \ldots i_p}$ 表示 $f$ 在 $X_{i_0 \ldots i_p}$ 上的限制。
复形中的映射是评注 \ref{more-etale-remark-covariance-lqf-f-shriek} 中构造的映射，
符号规则与 {\v C}ech 复形相同。由引理 \ref{more-etale-lemma-lqf-f-shriek-stalk}
对茎的描述，正合性容易得到。略去细节。
\end{proof}

\begin{remark}[另一种构造]
\label{more-etale-remark-alternative-lqf-f-shriek}
引理 \ref{more-etale-lemma-lqf-colimit-f-shriek} 给出了局部拟有限态射 $f$ 的函子
$f_!$ 的另一种构造。具体地，给定概形的局部拟有限态射
$f : X \to Y$，可选取开覆盖 $X = \bigcup_{i \in I} X_i$，使每个
$f_i : X_i \to Y$ 都分离。例如可取 $X$ 的仿射开覆盖。于是可把
$f_!\mathcal{F}$ 定义为该引理复形倒数第二个映射的余核，即
$$
f_!\mathcal{F} = \Coker\left(
\bigoplus\nolimits_{i_0, i_1} f_{i_0i_1, !} \mathcal{F}|_{X_{i_0i_1}} \to
\bigoplus\nolimits_{i_0} f_{i_0, !} \mathcal{F}|_{X_{i_0}}
\right)
$$
。由于态射 $f_{i_0}$ 与 $f_{i_0 i_1}$ 分离，可使用第
\ref{more-etale-section-compact-support} 节中 $f_{i_0, !}$ 和 $f_{i_0i_1, !}$ 的构造。
随后可计算 $f_!$ 的茎（使用分离情形，即引理
\ref{more-etale-lemma-lqf-f-shriek-separated-colimits}），并得到引理
\ref{more-etale-lemma-lqf-f-shriek-stalk} 的结论。完成这一步后，本节其他全部结果也可由此推出。
\end{remark}

\begin{remark}
\label{more-etale-remark-construct-map-presheaves-downstairs}
设 $g : Y' \to Y$ 为概形态射。对 $Y'_\etale$ 上的阿贝尔预层
$\mathcal{G}'$，以 $g_*\mathcal{G}'$ 表示预层
$V \mapsto \mathcal{G}'(Y' \times_Y V)$。若
$\alpha : \mathcal{G} \to g_*\mathcal{G}'$ 是 $Y_\etale$ 上阿贝尔预层的映射，
则存在唯一的 $Y_\etale$ 上阿贝尔层映射
$\alpha^\# : \mathcal{G}^\# \to g_*((\mathcal{G}')^\#)$
使图
$$
\xymatrix{
\mathcal{G} \ar[d] \ar[r]_\alpha & g_*\mathcal{G}' \ar[d] \\
\mathcal{G}^\# \ar[r]^-{\alpha^\#} & g_*((\mathcal{G}')^\#)
}
$$
交换，其中竖映射来自典范映射 $\mathcal{G} \to \mathcal{G}^\#$ 与
$\mathcal{G}' \to (\mathcal{G}')^\#$。若
$\alpha' : g^{-1}\mathcal{G}^\# \to (\mathcal{G}')^\#$
是 $\alpha^\#$ 的伴随映射，则对几何点
$\overline{y}' : \Spec(k) \to Y'$，其在 $Y$ 中的像为
$\overline{y} = g \circ \overline{y}'$，映射
$$
\alpha'_{\overline{y}'} :
\mathcal{G}_{\overline{y}} =
(\mathcal{G}^\#)_{\overline{y}} =
(g^{-1}\mathcal{G}^\#)_{\overline{y}'}
\longrightarrow
(\mathcal{G}')^\#_{\overline{y}'} =
\mathcal{G}'_{\overline{y}'}
$$
按如下方式给出：把 $\mathcal{G}$ 在 \'etale 邻域
$(V, \overline{v})$ 上的截面 $s$ 在茎中的类，映到截面
$\alpha(s)$ 在 $g_*\mathcal{G}'(V) = \mathcal{G}'(Y' \times_Y V)$ 中的类；
后者取于 \'etale 邻域 $(Y' \times_Y V, (\overline{y}', \overline{v}))$
上，并视为 $\mathcal{G}'$ 在 $\overline{y}'$ 处茎中的类。
\end{remark}

\begin{lemma}
\label{more-etale-lemma-lqf-base-change-f-shriek}
考虑概形的笛卡尔方块
$$
\xymatrix{
X' \ar[r]_{g'} \ar[d]_{f'} & X \ar[d]^f \\
Y' \ar[r]^g & Y
}
$$
，其中 $f$ 局部拟有限。存在关于 $\textit{Ab}(X_\etale)$ 中的
$\mathcal{F}$ 函子性的同构
$g^{-1}f_!\mathcal{F} \to f'_!(g')^{-1}\mathcal{F}$；它与引理
\ref{more-etale-lemma-lqf-f-shriek-stalk} 给出的茎的描述相容（精确陈述见证明）。
\end{lemma}

\begin{proof}
采用评注 \ref{more-etale-remark-construct-map-presheaves-downstairs} 中的约定，我们显式构造
$Y_\etale$ 上阿贝尔预层的映射
$$
c : f_{p!}\mathcal{F} \longrightarrow g_*f'_{p!}(g')^{-1}\mathcal{F}
$$
。由评注 \ref{more-etale-remark-construct-map-presheaves-downstairs} 中的讨论，这确定了典范映射
$g^{-1}f_!\mathcal{F} \to f'_!(g')^{-1}\mathcal{F}$。最后，我们将证明此映射
在茎上诱导同构，再由“\'Etale 上同调”定理
\ref{etale-cohomology-theorem-exactness-stalks} 得出结论。

\medskip\noindent
构造映射 $c$。设 $V \in Y_\etale$，并考虑如
(\ref{more-etale-equation-formal-sum}) 所示的截面
$s = \sum_{i = 1, \ldots, n} (Z_i, s_i)$，它定义了
$f_{p!}\mathcal{F}(V)$ 的一个元素。$g_*f'_{p!}(g')^{-1}\mathcal{F}$ 在 $V$
上的值为 $f'_{p!}(g')^{-1}\mathcal{F}(V')$，其中
$V' = V \times_Y Y'$。以 $Z'_i \subset X'_{V'}$ 表示 $Z_i$ 到 $V'$ 的基变换。
由 (\ref{more-etale-item-pullback})，存在拉回映射
$H_{Z_i}(\mathcal{F}) \to H_{Z'_i}((g')^{-1}\mathcal{F})$.
以 $s'_i \in H_{Z'_i}((g')^{-1}\mathcal{F})$ 表示 $s_i$ 在拉回下的像，并令
$c(s) = \sum_{i = 1, \ldots, n} (Z'_i, s'_i)$；它如
(\ref{more-etale-equation-formal-sum}) 所示，定义了
$f'_{p!}(g')^{-1}\mathcal{F}(V')$ 的一个元素。略去如下验证：此构造与关系
(\ref{more-etale-item-sum})、(\ref{more-etale-item-sub}) 以及限制映射相容。该构造显然关于
$\mathcal{F}$ 具有函子性。

\medskip\noindent
设 $\overline{y}' : \Spec(k) \to Y'$ 为几何点，其在 $Y$ 中的像为
$\overline{y} = g \circ \overline{y}'$。由纤维积的传递性，有
$X'_{\overline{y}'} = X_{\overline{y}}$。因此，$g'$ 给出双射
$\{f'(\overline{x}') = \overline{y}'\} \to \{f(\overline{x}) = \overline{y}\}$
；若 $\overline{x}'$ 映到 $\overline{x}$，则
$((g')^{-1}\mathcal{F})_{\overline{x}'} = \mathcal{F}_{\overline{x}}$
；这由“\'Etale 上同调”引理 \ref{etale-cohomology-lemma-stalk-pullback} 得到。
现在断言图
$$
\xymatrix{
(g^{-1}f_!\mathcal{F})_{\overline{y}'} \ar@{=}[r] \ar[d] &
(f_!\mathcal{F})_{\overline{y}} \ar[r] \ar[ld] &
\bigoplus\nolimits_{f(\overline{x}) = \overline{y}}
\mathcal{F}_{\overline{x}} \ar[d]
\\
(f'_!(g')^{-1}\mathcal{F})_{\overline{y}'} \ar[rr] & &
\bigoplus\nolimits_{f'(\overline{x}') = \overline{y}'}
(g')^{-1}\mathcal{F}_{\overline{x}'}
}
$$
交换，其中横箭头由引理 \ref{more-etale-lemma-finite-support-stalk} 的证明给出，
右侧竖箭头依上述讨论是等号。西南方向的箭头在评注
\ref{more-etale-remark-construct-map-presheaves-downstairs} 中描述为拉回映射，即恰由上述构造
$c$ 给出。引理 \ref{more-etale-lemma-finite-support-stalk} 的证明对和
$\sum (Z_i, z_i)$ 在 $\overline{x}$ 处茎中之像的简单描述，立即说明该图交换。
本引理得证。
\end{proof}

\begin{lemma}
\label{more-etale-lemma-lqf-separated-shriek-composition}
设 $f' : X \to Y'$ 与 $g : Y' \to Y$ 为可复合的概形态射，其中 $f'$ 和
$f = g \circ f'$ 局部拟有限，且 $g$ 分离并局部有限型。则存在函子的典范同构
$g_! \circ f'_! = f_!$。此同构与以下各项相容：
\begin{enumerate}
\item[(a)] 评注 \ref{more-etale-remark-covariance-f-shriek-separated} 与
\ref{more-etale-remark-covariance-lqf-f-shriek} 所述关于开嵌入的协变性；
\item[(b)] 引理 \ref{more-etale-lemma-lqf-base-change-f-shriek} 与
\ref{more-etale-lemma-base-change-f-shriek-separated} 的基变换同构；
\item[(c)] 当 $f'$ 分离时，它等于引理
\ref{more-etale-lemma-f-shriek-composition} 的同构；这里使用引理
\ref{more-etale-lemma-finite-support-f-shriek-separated} 的等同。
\end{enumerate}
\end{lemma}

\begin{proof}
设 $\mathcal{F}$ 为 $X_\etale$ 上的阿贝尔层。采用评注
\ref{more-etale-remark-construct-map-presheaves-downstairs} 中的约定，我们将显式构造
$Y_\etale$ 上阿贝尔预层的映射
$$
c : f_{p!}\mathcal{F} \longrightarrow g_*f'_{p!}\mathcal{F}
$$
。由评注 \ref{more-etale-remark-construct-map-presheaves-downstairs} 中的讨论，这确定了典范映射
$c^\# : f_!\mathcal{F} \to g_*f'_!\mathcal{F}$。我们将证明 $c^\#$ 的像包含在
子层 $g_!f'_!\mathcal{F}$ 中，从而得到映射
$c' : f_!\mathcal{F} \to g_!f'_!\mathcal{F}$。接着证明 (a)、(b) 和 (c)。
最后，(b) 将使我们能够证明 $c'$ 为同构。

\medskip\noindent
构造映射 $c$。设 $V \in Y_\etale$，并设
$s = \sum (Z_i, s_i)$ 为如 (\ref{more-etale-equation-formal-sum}) 所示的和，
它定义了 $f_{p!}\mathcal{F}(V)$ 的一个元素。回忆，
$Z_i \subset X_V = X \times_Y V$ 是在 $V$ 上有限的局部闭子概形。
令 $V' = Y' \times_Y V$，则 $X_{V'} = X \times_{Y'} V' = X_V$。
因此 $Z_i \subset X_{V'}$ 局部闭；又因 $g$ 分离，$Z_i$ 在 $V'$ 上有限
（“态射”引理 \ref{morphisms-lemma-finite-permanence}）。故可令
$c(s) = \sum (Z_i, s_i)$，但现在把它看作
$f'_{p!}\mathcal{F}(V') = (g_*f'_{p!}\mathcal{F})(V)$ 的元素。此构造显然与关系
(\ref{more-etale-item-sum})、(\ref{more-etale-item-sub}) 以及限制映射相容，因而得到映射 $c$。

\medskip\noindent
注意，在上述讨论中，$f'_!\mathcal{F}$ 在 $V'$ 上的截面
$c(s) = \sum (Z_i, s_i)$ 在
$V' \setminus \Im(\coprod Z_i \to V')$ 上限制为零。由于
$\Im(\coprod Z_i \to V')$ 在 $V$ 上固有（例如由“态射”引理
\ref{morphisms-lemma-scheme-theoretic-image-is-proper}），可知 $c(s)$ 定义了
$g_!f'_!\mathcal{F} \subset g_*f'_!\mathcal{F}$ 在 $V$ 上的截面。
由于 $f_!\mathcal{F}$ 的每个局部截面在局部都来自 $f_{p!}\mathcal{F}$
的局部截面，可知 $c^\#$ 的像包含在 $g_!f'_!\mathcal{F}$ 中。因此得到诱导映射
$c' : f_!\mathcal{F} \to g_!f'_!\mathcal{F}$，它按证明首段所述分解 $c^\#$。

\medskip\noindent
证明 (a)。设 $Y'_1 \subset Y'$ 为开子概形，并令
$X_1 = (f')^{-1}(W')$。得到图
$$
\xymatrix{
X_1 \ar[d]_{f'_1} \ar[r]_a \ar@/_2em/[dd]_{f_1} &
X \ar[d]^{f'} \ar@/^2em/[dd]^f \\
Y'_1 \ar[d]_{g_1} \ar[r]_{b'} &
Y' \ar[d]^g \\
Y \ar@{=}[r] &
Y
}
$$
，其中横箭头均为开浸入。我们断言图
$$
\xymatrix{
f_{1, !}\mathcal{F}|_{X_1} \ar[r]_{c'_1} \ar[dd] &
g_{1, !}f'_{1, !}\mathcal{F}|_{X_1} \ar@{=}[d] \\
& g_{1, !}(f'_!\mathcal{F})|_{Y'_1} \ar[d] \\
f_!\mathcal{F} \ar[r]^{c'} &
g_!f'_!\mathcal{F} \ar[r] & g_*f'_!\mathcal{F}
}
$$
交换，其中左侧竖箭头来自评注 \ref{more-etale-remark-covariance-lqf-f-shriek}，
右侧竖箭头来自评注 \ref{more-etale-remark-covariance-f-shriek-separated}。
图中的等号来自如下事实：$f'_1$ 是 $f'$ 在 $Y'_1$ 上的限制，且我们对
$f'_!$ 的构造在基底上是局部的。最后，为证明交换性，取 $Y_\etale$ 的对象
$V$ 以及如 (\ref{more-etale-equation-formal-sum}) 所示的形式和
$s_1 = \sum (Z_{1, i}, s_{1, i})$，它定义了
$f_{1, p!}\mathcal{F}|_{X_1}(V)$ 的一个元素。回忆，这意味着
$Z_{1, i} \subset X_1 \times_Y V$ 局部闭且在 $V$ 上有限，并且
$s_{1, i} \in H_{Z_{1, i}}(\mathcal{F})$。沿所涉及的映射追踪此截面，
只须证明沿图的两条路径最终得到
$g_*f'_!\mathcal{F}(V) = f'_!\mathcal{F}(Y' \times_Y V)$.
中的同一元素。读者立即可见，两条路径均得到元素
$\sum (Z_{1, i}, s_{1, i})$，其中现在把 $Z_{1, i}$ 看作
$X \times_{Y'} (Y' \times_Y V) = X \times_Y V$ 的局部闭子概形，
且它在 $Y' \times_Y V$ 上有限。

\medskip\noindent
证明 (b)。设 $b : Y_1 \to Y$ 为概形态射。构造交换图
$$
\xymatrix{
X_1 \ar[d]_{f'_1} \ar[r]_a \ar@/_2em/[dd]_{f_1} &
X \ar[d]^{f'} \ar@/^2em/[dd]^f \\
Y'_1 \ar[d]_{g_1} \ar[r]_{b'} &
Y' \ar[d]^g \\
Y_1 \ar[r]^b &
Y
}
$$
，其中各方块为笛卡尔方块。我们断言该构造与引理
\ref{more-etale-lemma-lqf-base-change-f-shriek} 和
\ref{more-etale-lemma-base-change-f-shriek-separated} 的基变换映射相容；即下图的上方矩形
$$
\xymatrix{
b^{-1}f_!\mathcal{F} \ar[rr] \ar[d]_{b^{-1}c'} & &
f_{1, !}a^{-1}\mathcal{F} \ar[d]^{c_1'} \\
b^{-1}g_!f'_!\mathcal{F} \ar[r] \ar[d] &
g_{1, !}(b')^{-1}f'_!\mathcal{F} \ar[r] \ar[d] &
g_{1, !}f'_{1, !}a^{-1}\mathcal{F} \ar[d] \\
b^{-1}g_*f'_!\mathcal{F} \ar[r] &
g_{1, *}(b')^{-1}f'_!\mathcal{F} \ar[r] &
g_{1, *}f'_{1, !}a^{-1}\mathcal{F}
}
$$
交换。验证完全例行，我们建议读者略过。由于从中间一行到最下一行的箭头均为
单射，只须证明外框交换。为此，只须取 $b^{-1}f_!\mathcal{F}$ 的局部截面，
并证明沿任一路径最终得到 $g_{1, *}f'_{1, !}a^{-1}\mathcal{F}$ 的同一局部截面。
事实上，只需对如下局部截面作检验：它是上述 $f_{p!}\mathcal{F}(V)$ 的截面
$s = \sum (Z_i, s_i)$ 沿 $b$ 的拉回（因为这些拉回生成阿贝尔层
$b^{-1}f_!\mathcal{F}$）。以 $V_1$、$V'_1$ 和 $Z_{1, i}$ 分别表示
$V$、$V' = Y' \times_Y V$ 和 $Z_i$ 沿 $Y_1 \to Y$ 的基变换。回忆，
$Z_i$ 是 $X_V = X_{V'}$ 的局部闭子概形，因而 $Z_{1, i}$ 是
$(X_1)_{V_1} = (X_1)_{V'_1}$ 的局部闭子概形。此时 $b^{-1}c'$ 把 $s$
的拉回映到局部截面 $c(s) \sum (Z_i, s_i)$ 的拉回，其中后者视为
$f'_{p!}\mathcal{F}(V') = (g_*f'_{p!}\mathcal{F})(V)$ 的元素。
下方两个基变换映射的复合，恰把它映到
$\sum (Z_{i, 1}, s_{1, i})$，后者视为
$f'_{1, p!}a^{-1}\mathcal{F}(V'_1) =
g_{1, *}f'_{1, p!}a^{-1}\mathcal{F}(V_1)$ 的元素。另一方面，图上方的基变换映射
把 $s$ 的拉回映到 $\sum (Z_{1, i}, s_{1, i})$，后者视为
$f_{1, !}a^{-1}\mathcal{F}(V_1)$ 的元素。最后，按 $c'_1$ 的构造，它确实把
该元素映到 $\sum (Z_{i, 1}, s_{1, i})$，后者视为
$f'_{1, p!}a^{-1}\mathcal{F}(V'_1) =
g_{1, *}f'_{1, p!}a^{-1}\mathcal{F}(V_1)$ 的元素。交换性得证。

\medskip\noindent
证明 (c)。比较两个映射的定义即可。略去细节。

\medskip\noindent
为完成证明，只须证明 $c'$ 沿任意几何点
$\overline{y} : \Spec(k) \to Y$ 的拉回为同构。沿 $\overline{y}$ 拉回等同于取
$\overline{y}$ 处的茎（“\'Etale 上同调”评注
\ref{etale-cohomology-remark-stalk-pullback}），因而可应用“\'Etale 上同调”定理
\ref{etale-cohomology-theorem-exactness-stalks}。由刚证明的相容性 (b)，
可设 $Y$ 是 $k$ 的谱，而我们须证明 $c'$ 为同构。只须证明诱导映射
$$
\bigoplus\nolimits_{x \in X} \mathcal{F}_x = H^0(Y, f_!\mathcal{F})
\longrightarrow
H^0(Y, g_!f'_!\mathcal{F}) = H^0_c(Y', f'_!\mathcal{F})
$$
为同构。等号由引理 \ref{more-etale-lemma-lqf-f-shriek-stalk} 与
\ref{more-etale-lemma-stalk-f-shriek-separated} 成立。回忆，$X$ 是剩余域为 $k$ 的
Artin 局部环之谱的不交并；见“簇”引理
\ref{varieties-lemma-algebraic-scheme-dim-0}。由于左右两端都与直和可交换
（略去细节），可设 $\mathcal{F}$ 是支撑在某个 $x \in X$ 处的摩天层
$x_*A$。由引理 \ref{more-etale-lemma-lqf-f-shriek-stalk}，$f'_!\mathcal{F}$ 是
$x$ 在 $Y$ 中的像 $y'$ 处的摩天层。在这种情形下，我们的构造显然给出恒等映射
$A \to H^0_c(Y', y'_*A) = A$，正合所需。
\end{proof}

\begin{lemma}
\label{more-etale-lemma-lqf-shriek-composition}
设 $f : X \to Y$ 与 $g : Y \to Z$ 为可复合的局部拟有限概形态射。
则存在函子的典范同构
$$
(g \circ f)_! \longrightarrow g_! \circ f_!
$$
。这些同构满足以下性质：
\begin{enumerate}
\item 若 $f$ 与 $g$ 分离，则此同构与引理
\ref{more-etale-lemma-f-shriek-composition} 的同构一致。
\item 若 $g$ 分离，则此同构与引理
\ref{more-etale-lemma-lqf-separated-shriek-composition} 的同构一致。
\item 对几何点 $\overline{z} : \Spec(k) \to Z$，图
$$
\xymatrix{
((g \circ f)_!\mathcal{F})_{\overline{z}} \ar[d] \ar[rr] & &
\bigoplus\nolimits_{g(f(\overline{x})) = \overline{z}}
\mathcal{F}_{\overline{x}} \ar@{=}[d] \\
(g_!f_!\mathcal{F})_{\overline{z}} \ar[r] &
\bigoplus\nolimits_{g(\overline{y}) = \overline{z}}
(f_!\mathcal{F})_{\overline{y}} \ar[r] &
\bigoplus\nolimits_{g(f(\overline{x})) = \overline{z}}
\mathcal{F}_{\overline{x}}
}
$$
交换，其中横箭头由引理 \ref{more-etale-lemma-lqf-f-shriek-stalk} 给出。
\item 设 $h : Z \to T$ 为第三个局部拟有限概形态射。则图
$$
\xymatrix{
(h \circ g \circ f)_! \ar[r] \ar[d] &
(h \circ g)_! \circ f_! \ar[d] \\
h_! \circ (g \circ f)_! \ar[r] &
h_! \circ g_! \circ f_!
}
$$
交换。
\item 假设有概形图
$$
\xymatrix{
X' \ar[d]_{f'} \ar[r]_c & X \ar[d]^f \\
Y' \ar[d]_{g'} \ar[r]_b & Y \ar[d]^g \\
Z' \ar[r]^a & Z
}
$$
，其中两个方块均为笛卡尔方块，且 $f$ 与 $g$ 局部拟有限。则图
$$
\xymatrix{
a^{-1} \circ (g \circ f)_! \ar[d] \ar[rr] & &
(g' \circ f')_! \circ c^{-1} \ar[d] \\
a^{-1} \circ g_! \circ f_! \ar[r] &
g'_! \circ b^{-1} \circ f_! \ar[r] &
g'_! \circ f'_! \circ c^{-1}
}
$$
交换，其中横箭头来自引理 \ref{more-etale-lemma-lqf-base-change-f-shriek}。
\end{enumerate}
\end{lemma}

\begin{proof}
若 $f$ 与 $g$ 分离，则这是引理 \ref{more-etale-lemma-f-shriek-composition} 的特例。
若 $g$ 分离，则这是引理 \ref{more-etale-lemma-lqf-separated-shriek-composition}
的特例；而且它与 $f$、$g$ 均分离的情形一致。

\medskip\noindent
一般情形下的构造。选取开覆盖 $Y = \bigcup Y_i$，使 $g$ 的限制
$g_i : Y_i \to Z$ 分离。令 $X_i = f^{-1}(Y_i)$，并以
$f_i : X_i \to Y_i$ 表示 $f$ 的限制。再令 $h = g \circ f$，并以
$h_i : X_i \to Z$ 表示 $h$ 的限制。考虑图
$$
\xymatrix{
\bigoplus\nolimits_{i_0, i_1}
h_{i_0i_1, !}\mathcal{F}|_{X_{i_0i_1}} \ar[r] \ar[d] &
\bigoplus\nolimits_{i_0} h_{i_0, !}\mathcal{F}|_{X_{i_0}} \ar[r] \ar[d] &
h_!\mathcal{F} \ar[r] \ar@{..>}[dd] &
0 \\
\bigoplus\nolimits_{i_0, i_1}
g_{i_0i_1, !} f_{i_0i_1, !}\mathcal{F}|_{X_{i_0i_1}} \ar[r] \ar[d] &
\bigoplus\nolimits_{i_0}
g_{i_0, !} f_{i_0, !}\mathcal{F}|_{X_{i_0}} \ar[d] \\
\bigoplus\nolimits_{i_0, i_1}
g_{i_0i_1, !} (f_!\mathcal{F})|_{Y_{i_0i_1}} \ar[r] &
\bigoplus\nolimits_{i_0}
g_{i_0, !} (f_!\mathcal{F})|_{Y_{i_0}} \ar[r] &
g_!f_!\mathcal{F} \ar[r] &
0
}
$$
由引理 \ref{more-etale-lemma-lqf-colimit-f-shriek}，图的最上行与最下行均正合。
由引理 \ref{more-etale-lemma-lqf-separated-shriek-composition}，左上方块交换。
左下方块中的竖箭头来自等式
$(f_!\mathcal{F})|_{Y_{i_0i_1}} = f_{i_0i_1, !}\mathcal{F}|_{X_{i_0i_1}}$ 与
$(f_!\mathcal{F})|_{Y_{i_0}} = f_{i_0, !}\mathcal{F}|_{X_{i_0}}$
，因为 $f_!$ 的构造在基底上是局部的。此外，这些等式当然与以下等同相容：
$((f_!\mathcal{F})|_{Y_{i_0}})|_{Y_{i_0i_1}} =
(f_!\mathcal{F})|_{Y_{i_0i_1}}$ 以及
$(f_{i_0, !}\mathcal{F}|_{X_{i_0}})|_{Y_{i_0i_1}} =
f_{i_0i_1, !}\mathcal{F}|_{X_{i_0i_1}}$
。这些等同连同关于开嵌入 $Y_{i_0i_1} \subset Y_{i_0}$ 的协变性一起，
用于定义左下方块的横映射。因此这个方块也交换。由此可知，存在图中所示的
唯一虚线箭头，而且该箭头是同构。

\medskip\noindent
证明性质 (1)--(5)。固定开覆盖 $Y = \bigcup Y_i$。注意，若
$Y \to Z$ 恰好分离，则使用引理
\ref{more-etale-lemma-lqf-separated-shriek-composition} 的映射，便得到适配于上述大图的
虚线箭头（由该引理自身的性质）。这证明了 (2)；再由引理
\ref{more-etale-lemma-lqf-separated-shriek-composition} 与引理
\ref{more-etale-lemma-f-shriek-composition} 的映射之相容性，也证明了 (1)。接着，对 $Z$
上的任意概形 $Z'$，采用开覆盖 $Y' = \bigcup b^{-1}(Y_i)$ 构造的映射
$(g' \circ f')_! \to g'_! \circ f'_!$ 具有 (5) 中的相容性。这由引理
\ref{more-etale-lemma-lqf-separated-shriek-composition} 中所构造映射的相应相容性立即得到。
特别地，可考虑几何点 $\overline{z} : \Spec(k) \to Z$。由于
$X_{\overline{z}} \to Y_{\overline{z}} \to \Spec(k)$ 均为分离态射，
映射
$(g \circ f)_!\mathcal{F} \to g_! f_! \mathcal{F}$
沿 $\overline{z}$ 的基变换等于引理 \ref{more-etale-lemma-f-shriek-composition} 的映射。
读者立即可见，由此得到性质 (3)。当然，性质 (3) 保证采用开覆盖
$Y = \bigcup Y_i$ 构造的函子变换
$(g \circ f)_! \to g_! \circ f_!$ 不依赖于该开覆盖的选择。最后，利用已证明的
性质 (3) 考察茎，即得性质 (4)。
\end{proof}








\section{局部拟有限态射的权函数与迹映射}
\label{more-etale-section-weightings}

\noindent
本节可参考 \cite[报告 XVII，命题 6.2.5]{SGA4}。

\medskip\noindent
设 $f : X \to Y$ 为局部拟有限的概形态射。设
$w : X \to \mathbf{Z}$ 为 $f$ 的权函数；见“关于态射的进一步结果”定义
\ref{more-morphisms-definition-weighting}。设 $\mathcal{F}$ 为
$Y_\etale$ 上的阿贝尔层。本节将证明存在 $Y_\etale$ 上阿贝尔层的映射
$$
\text{Tr}_{f, w, \mathcal{F}} :
f_!f^{-1}\mathcal{F}
\longrightarrow
\mathcal{F}
$$
，它由以下性质刻画：在 $Y$ 的几何点 $\overline{y}$ 处的茎上，得到映射
$$
\bigoplus\nolimits_{f(\overline{x}) = \overline{y}} w(\overline{x}) :
(f_!f^{-1}\mathcal{F})_{\overline{y}} =
\bigoplus\nolimits_{f(\overline{x}) = \overline{y}}
\mathcal{F}_{\overline{y}}
\longrightarrow
\mathcal{F}_{\overline{y}}
$$
如此处所示，在对应于 $\overline{x}$ 的直和项上，该箭头由乘以整数
$w(\overline{x})$ 给出。箭头左侧的等号由引理
\ref{more-etale-lemma-lqf-f-shriek-stalk} 与“\'Etale 上同调”引理
\ref{etale-cohomology-lemma-stalk-pullback} 得到。

\medskip\noindent
若态射 $f : X \to Y$ 平坦、局部拟有限且局部有限表示，则存在典范权函数，
并得到其构造与基变换相容的典范迹映射；见例
\ref{more-etale-example-trace-for-flat-quasi-finite}。若 $Y$ 为局部 Noether 单枝概形，
且 $f : X \to Y$ 局部拟有限，则也可为 $f$ 定义一个（自然的）权函数，
从而在此情形下也有迹映射；见例
\ref{more-etale-example-trace-for-quasi-finite-over-normal}。

\begin{lemma}
\label{more-etale-lemma-f-shriek-projection}
设 $f : X \to Y$ 为局部拟有限的概形态射，$\Lambda$ 为环。
设 $\mathcal{F}$ 为 $X_\etale$ 上的 $\Lambda$-模层，$\mathcal{G}$ 为
$Y_\etale$ 上的 $\Lambda$-模层。存在 $Y_\etale$ 上 $\Lambda$-模层的典范同构
$$
can :
f_!\mathcal{F} \otimes_\Lambda \mathcal{G}
\longrightarrow
f_!(\mathcal{F} \otimes_\Lambda f^{-1}\mathcal{G})
$$
\end{lemma}

\begin{proof}
回忆，由定义 \ref{more-etale-definition-f-shriek-lqf}，有
$f_!\mathcal{F} = (f_{p!}\mathcal{F})^\#$，其中 $f_{p!}\mathcal{F}$ 是第
\ref{more-etale-section-finite-support} 节构造的预层。因此，为构造该箭头，只须构造
$Y_\etale$ 上预层的映射
$$
f_{p!}\mathcal{F} \otimes_{p, \Lambda} \mathcal{G}
\longrightarrow
f_{p!}(\mathcal{F} \otimes_\Lambda f^{-1}\mathcal{G})
$$
。这里符号 $\otimes_{p, \Lambda}$ 表示预层张量积；见“位点上的模”第
\ref{sites-modules-section-tensor-product} 节。设 $V$ 为 $Y_\etale$ 的对象。
回忆
$$
f_{p!}\mathcal{F}(V) = \colim_Z H_Z(\mathcal{F})
\quad\text{且}\quad
f_{p!}(\mathcal{F} \otimes_\Lambda f^{-1}\mathcal{G})(V) =
\colim_Z H_Z(\mathcal{F} \otimes_\Lambda f^{-1}\mathcal{G})
$$
；见第 \ref{more-etale-section-finite-support} 节。我们的映射在纯张量上由规则
$$
(Z, s) \otimes t \longmapsto (Z, s \otimes f^{-1}t)
$$
定义（记号见下文），再线性延拓到整个
$(f_{p!}\mathcal{F} \otimes_{p, \Lambda} \mathcal{G})(V) =
f_{p!}\mathcal{F}(V) \otimes_\Lambda \mathcal{G}(V)$.
所用记号如下：
\begin{enumerate}
\item $Z \subset X_V$ 是在 $V$ 上有限的局部闭子概形；
\item $s \in H_Z(\mathcal{F})$，这意味着对某个开子概形
$U \subset X_V$，其中 $Z \subset U$ 闭，有 $s \in \mathcal{F}(U)$ 且
$\text{Supp}(s) \subset Z$；
\item $t \in \mathcal{G}(V)$，其像为
$f^{-1}t \in f^{-1}\mathcal{G}(U)$。
\end{enumerate}
由于 $s \in \mathcal{F}(U)$ 的支撑包含在 $Z$ 中，
$s \otimes f^{-1}t$ 的支撑显然也包含在 $Z$ 中。因此，有序对
$(Z, s \otimes f^{-1}t)$ 确有意义。立即可见，该构造与余极限
$\colim_Z H_Z(\mathcal{F})$ 的转移映射可交换，且与限制映射相容。
最后，同样显然的是，该构造与引理 \ref{more-etale-lemma-lqf-f-shriek-stalk}
中 $f_!$ 的茎的等同相容。换言之，我们构造的映射 $can$ 在几何点
$\overline{y}$ 处的茎上适配于交换图
$$
\xymatrix{
(f_!\mathcal{F} \otimes_\Lambda \mathcal{G})_{\overline{y}}
\ar[r]_-{can_{\overline{y}}} \ar[d] &
f_!(\mathcal{F} \otimes_\Lambda f^{-1}\mathcal{G})_{\overline{y}} \ar[d] \\
(\bigoplus \mathcal{F}_{\overline{x}})
\otimes_\Lambda \mathcal{G}_{\overline{y}} \ar[r] &
\bigoplus
(\mathcal{F}_{\overline{x}} \otimes_\Lambda \mathcal{G}_{\overline{y}})
}
$$
，其中直和取遍位于 $\overline{y}$ 上方的几何点 $\overline{x}$，竖箭头是引理
\ref{more-etale-lemma-lqf-f-shriek-stalk} 的等同，而下方横箭头是显然的同构。
因此 $can$ 是同构，正合所需。
\end{proof}

\begin{lemma}
\label{more-etale-lemma-trace-map-exists}
设 $f : X \to Y$ 为局部拟有限的概形态射，
$w : X \to \mathbf{Z}$ 为 $f$ 的权函数。对 $Y$ 上任意阿贝尔层
$\mathcal{F}$，存在唯一的迹映射
$\text{Tr}_{f, w, \mathcal{F}} : f_!f^{-1}\mathcal{F} \to \mathcal{F}$
，它在茎上具有预先指定的行为。
\end{lemma}

\begin{proof}
由引理 \ref{more-etale-lemma-f-shriek-projection}，有等同
$f_!f^{-1}\mathcal{F} = f_!\underline{\mathbf{Z}} \otimes \mathcal{F}$
，它与这些层在几何点处茎的描述相容。因此，只须构造映射
$$
\text{Tr}_{f, w, \underline{\mathbf{Z}}} :
f_!\underline{\mathbf{Z}}
\longrightarrow
\underline{\mathbf{Z}}
$$
，使其在茎上具有预先指定的行为。由定义 \ref{more-etale-definition-f-shriek-lqf}，有
$f_!\underline{\mathbf{Z}} = (f_{p!}\underline{\mathbf{Z}})^\#$
，其中 $f_{p!}\underline{\mathbf{Z}}$ 是第 \ref{more-etale-section-finite-support} 节
构造的预层。因此只须构造 $Y_\etale$ 上预层的映射
$$
f_{p!}\underline{\mathbf{Z}} \longrightarrow \underline{\mathbf{Z}}
$$
。设 $V$ 为 $Y_\etale$ 的对象。由第 \ref{more-etale-section-finite-support} 节回忆
$$
f_{p!}\underline{\mathbf{Z}}(V) = \colim_Z H_Z(\underline{\mathbf{Z}})
$$
。这里余极限取遍所有在 $V$ 上有限的局部闭子概形 $Z \subset X_V$
所成的（偏序）集合。对每个这样的 $Z$，定义映射
$$
H_Z(\underline{\mathbf{Z}}) \longrightarrow \underline{\mathbf{Z}}(V)
$$
，使其与定义余极限的映射相容。

\medskip\noindent
设 $Z \subset X_V$ 局部闭且在 $V$ 上有限。取开子集
$U \subset X_V$，使它把 $Z$ 包含为闭子集。$H_Z(\underline{\mathbf{Z}})$
的元素 $s$ 是截面 $s \in \underline{\mathbf{Z}}(U)$，其支撑包含在 $Z$ 中。
令 $U_n \subset U$ 为 $s$ 取值 $n \in \mathbf{Z}$ 的开闭子集。由支撑条件，
当 $n \not = 0$ 时有 $Z \cap U_n = U_n$。因此，对 $n \not = 0$，开集
$U_n$ 在 $Z$ 中也闭（作为其余各项之并的补集）；又因 $Z$ 在 $V$ 上有限，
可知 $U_n \to V$ 有限。由权函数的定义，这意味着函数
$\int_{U_n \to V} w|_{U_n}$ 在 $V$ 上局部常值，故可把它看作
$\underline{\mathbf{Z}}(V)$ 的元素。我们的构造把 $(Z, s)$ 映到元素
$$
\sum\nolimits_{n \in \mathbf{Z},\ n \not = 0}
n \left(\int_{U_n \to V} w|_{U_n}\right)
\quad \in \quad \underline{\mathbf{Z}}(V)
$$
。该和在 $V$ 上局部有限，因而确有意义；略去细节（在整个讨论中，读者可先取
仿射开集，并保证论证中出现的所有概形都拟紧，从而使该和为有限和）。略去如下
验证：此构造与余极限中的映射以及定义
$f_{p!}\underline{\mathbf{Z}}$ 的限制映射相容。

\medskip\noindent
设 $\overline{y}$ 为 $Y$ 的几何点，位于点 $y \in Y$ 上方。在
$\overline{y}$ 处取茎，上述构造确定映射
$$
(f_!\underline{\mathbf{Z}})_{\overline{y}}
= \bigoplus\nolimits_{f(\overline{x}) = \overline{y}} \mathbf{Z}
\longrightarrow
\mathbf{Z} = \underline{\mathbf{Z}}_{\overline{y}}
$$
。为完成证明，将说明在对应于 $\overline{x}$ 的直和项上，此映射由乘以
$w(\overline{x})$ 给出。具体地，取位于 $\overline{y}$ 上方的
$\overline{x}$。可找到 \'etale 邻域
$(V, \overline{v}) \to (Y, \overline{y})$，使 $X_V$ 含有在 $V$ 上有限的
开子概形 $U$，并且 $U$ 只包含几何点 $\overline{x}$，不包含 $X$ 中位于
$\overline{y}$ 上方的其他几何点。这由“关于态射的进一步结果”引理
\ref{more-morphisms-lemma-etale-makes-quasi-finite-finite-multiple-points-var};
得出；略去一些细节。于是 $(U, 1)$ 定义了 $f_!\underline{\mathbf{Z}}$ 在
$V$ 上的截面；它在对应于 $\overline{x}$ 的直和项中映到 $1$，在其他直和项中
映到零（见引理 \ref{more-etale-lemma-finite-support-stalk} 的证明）。上述构造把
$(U, 1)$ 映到 $\int_{U \to V} w|_U$，后者在 $\overline{v}$ 的某个邻域上
恒取值 $w(\overline{x})$，正合所需。
\end{proof}

\begin{lemma}
\label{more-etale-lemma-properties-trace-map}
设 $f : X \to Y$ 为局部拟有限的概形态射，
$w : X \to \mathbf{Z}$ 为 $f$ 的权函数。上述构造的迹映射具有以下性质：
\begin{enumerate}
\item $\text{Tr}_{f, w, \mathcal{F}}$ 关于 $\mathcal{F}$ 具有函子性；
\item $\text{Tr}_{f, w, \mathcal{F}}$ 与任意基变换相容；
\item 给定环 $\Lambda$ 及 $D(Y_\etale, \Lambda)$ 中的 $K$，得到
$\text{Tr}_{f, w, K} : f_!f^{-1}K \to K$；它关于 $K$ 具有函子性，
且与任意基变换相容。
\end{enumerate}
\end{lemma}

\begin{proof}
(1) 可由引理 \ref{more-etale-lemma-trace-map-exists} 的证明中迹映射的构造得到；
也可更简单地注意到，该映射的刻画迫使这一性质在所有茎上成立。设
$$
\xymatrix{
X' \ar[r]_{g'} \ar[d]_{f'} & X \ar[d]^f \\
Y' \ar[r]^g & Y
}
$$
为概形的笛卡尔图。由“关于态射的进一步结果”引理
\ref{more-morphisms-lemma-weighting-base-change}，函数
$w' = w \circ g' : X' \to \mathbf{Z}$ 是 $f'$ 的权函数。陈述 (2) 意味着图
$$
\xymatrix{
g^{-1}f_!f^{-1}\mathcal{F}
\ar[rr]_-{g^{-1}\text{Tr}_{f, w, \mathcal{F}}} \ar@{=}[d] & &
g^{-1}\mathcal{F} \ar@{=}[d] \\
f'_!(f')^{-1}g^{-1}\mathcal{F}
\ar[rr]^-{\text{Tr}_{f', w', g^{-1}\mathcal{F}}} & &
g^{-1}\mathcal{F}
}
$$
交换，其中左侧竖等号由
$$
g^{-1}f_!f^{-1}\mathcal{F} =
f'_!(g')^{-1}f^{-1}\mathcal{F} =
f'_!(f')^{-1}g^{-1}\mathcal{F}
$$
给出；其中第一个等号来自引理 \ref{more-etale-lemma-lqf-base-change-f-shriek}
（下叹号的基变换）。该图的交换性来自迹映射在茎上作用的刻画，以及引理
\ref{more-etale-lemma-lqf-base-change-f-shriek} 的基变换映射保持茎的描述这一事实。

\medskip\noindent
由 (1) 与 (2)，(3) 随即得到，因为函子
$f^{-1} : D(Y_\etale, \Lambda) \to D(X_\etale, \Lambda)$ 与
$f_! : D(X_\etale, \Lambda) \to D(Y_\etale, \Lambda)$
是通过对表示所论对象的任意模复形逐项应用 $f^{-1}$ 与 $f_!$ 得到的。
\end{proof}

\begin{lemma}
\label{more-etale-lemma-trace-composition}
设 $f : X \to Y$ 与 $g : Y \to Z$ 为局部拟有限态射。设
$w_f : X \to \mathbf{Z}$ 为 $f$ 的权函数，
$w_g : Y \to \mathbf{Z}$ 为 $g$ 的权函数。对
$K \in D(Z_\etale, \Lambda)$，复合
$$
(g \circ f)_!(g \circ f)^{-1}K =
g_! f_! f^{-1} g^{-1}K
\xrightarrow{g_! \text{Tr}_{f, w_f, g^{-1}K}}
g_!g^{-1}K
\xrightarrow{\text{Tr}_{g, w_g, K}}
K
$$
等于 $\text{Tr}_{g \circ f, w_{g \circ f}, K}$，其中
$w_{g \circ f}(x) = w_f(x) w_g(f(x))$。
\end{lemma}

\begin{proof}
由引理 \ref{more-etale-lemma-lqf-shriek-composition}，有
$(g \circ f)_! = g_! \circ f_!$。在“关于态射的进一步结果”引理
\ref{more-morphisms-lemma-weighting-composition} 中已知 $w_{g \circ f}$
是 $g \circ f$ 的权函数，故该陈述确有意义。在茎上计算即可检验等式。
略去细节。
\end{proof}

\begin{example}[平坦拟有限态射的迹]
\label{more-etale-example-trace-for-flat-quasi-finite}
设 $f : X \to Y$ 为平坦、局部拟有限且局部有限表示的概形态射。
令
$$
w(x) = \text{length}_{\mathcal{O}_{X, x}}
(\mathcal{O}_{X, x}/\mathfrak m_{f(x)} \mathcal{O}_{X, x})
[\kappa(x) : \kappa(f(x))]_i
$$
，便得到典范正权函数 $w : X \to \mathbf{Z}$；见“关于态射的进一步结果”引理
\ref{more-morphisms-lemma-weighting-flat-quasi-finite}。因此，由引理
\ref{more-etale-lemma-trace-map-exists} 与 \ref{more-etale-lemma-properties-trace-map}，对 $f$ 得到迹映射
$$
\text{Tr}_{f, K} : f_!f^{-1}K \longrightarrow K
$$
；它关于 $D(Y_\etale, \Lambda)$ 中的 $K$ 具有函子性，且与任意基变换相容。
注意，$f$ 的任意基变换 $f' : X' \to Y'$ 满足相同性质，而且 $w$
限制为 $f'$ 的典范权函数。
\end{example}

\begin{remark}
\label{more-etale-remark-trace-etale-counit}
设 $j : U \to X$ 为概形的 \'etale 态射。则例
\ref{more-etale-example-trace-for-flat-quasi-finite} 的迹映射
$\text{Tr} : j_!j^{-1}K \to K$ 等于 $j_!$ 与 $j^{-1}$ 之间伴随的余单位。
在“\'Etale 上同调”第 \ref{etale-cohomology-section-trace-method} 节中，
我们已对这个余单位使用术语“迹”。
\end{remark}

\begin{example}[正规基底上拟有限态射的迹]
\label{more-etale-example-trace-for-quasi-finite-over-normal}
设 $Y$ 为几何单枝且局部 Noether 的概形；例如 $Y$ 可以是正规簇。
设 $f : X \to Y$ 为局部拟有限的概形态射。则存在 $f$ 的正权函数
$w : X \to \mathbf{Z}$；粗略地说，它把 $x$ 映到
$\mathcal{O}_{X, x}^{sh}$ 在 $\mathcal{O}_{Y, f(x)}^{sh}$ 上的“泛可分次数”。
见“关于态射的进一步结果”引理
\ref{more-morphisms-lemma-weighting-quasi-finite-Noetherian}。因此，由引理
\ref{more-etale-lemma-trace-map-exists} 与 \ref{more-etale-lemma-properties-trace-map}，对 $f$ 和 $w$
得到迹映射
$$
\text{Tr}_{f, w, K} : f_!f^{-1}K \longrightarrow K
$$
；它关于 $D(Y_\etale, \Lambda)$ 中的 $K$ 具有函子性，且与任意基变换相容。
不过，在此情形下，给定 $f$ 的基变换 $f' : X' \to Y'$，$w$ 在 $X'$
上的限制一般没有关于态射 $f'$ 的“自然”解释。
\end{example}
















\section{局部拟有限态射的上叹号}
\label{more-etale-section-duality-locally-quasi-finite}

\noindent
对概形的局部拟有限态射 $f : X \to Y$，函子
$f_! : \textit{Ab}(X_\etale) \to \textit{Ab}(Y_\etale)$ 与直和可交换且正合；
见引理 \ref{more-etale-lemma-lqf-f-shriek-stalk}。这表明它应有一个右伴随，记为 $f^!$。

\medskip\noindent
警告：这个函子是非导出版本！

\begin{lemma}
\label{more-etale-lemma-lqf-f-upper-shriek}
设 $f : X \to Y$ 为局部拟有限的概形态射。
\begin{enumerate}
\item 函子 $f_! : \textit{Ab}(X_\etale) \to \textit{Ab}(Y_\etale)$
有右伴随 $f^! : \textit{Ab}(Y_\etale) \to \textit{Ab}(X_\etale)$。
\item 有
$f^!(\overline{y}_*A) = \prod_{f(\overline{x}) = \overline{y}} \overline{x}_*A$.
\item 若 $\Lambda$ 为环，则函子
$f_! : \textit{Mod}(X_\etale, \Lambda) \to \textit{Mod}(Y_\etale, \Lambda)$
有右伴随
$f^! : \textit{Mod}(Y_\etale, \Lambda) \to \textit{Mod}(X_\etale, \Lambda)$
，且在底阿贝尔层上它与 $f^!$ 一致。
\end{enumerate}
\end{lemma}

\begin{proof}
证明 (1)。令 $E \subset \Ob(\textit{Ab}(Y_\etale))$ 为摩天层之积所成的类。
我们断言：
\begin{enumerate}
\item[(a)] $\textit{Ab}(Y_\etale)$ 中的每个 $\mathcal{G}$ 都是 $E$
中某个元素的子层；
\item[(b)] 对每个 $\mathcal{G} \in E$，存在
$\textit{Ab}(X_\etale)$ 的对象 $\mathcal{H}$，使得
$\Hom(f_!\mathcal{F}, \mathcal{G}) = \Hom(\mathcal{F}, \mathcal{H})$
关于 $\mathcal{F}$ 函子性地成立。
\end{enumerate}
一旦验证此断言，“同调”引理 \ref{homology-lemma-partially-defined-adjoint}
的对偶便给出伴随函子 $f^!$。

\medskip\noindent
(a) 成立，因为可把 $\mathcal{G}$ 映到层
$\prod \overline{y}_*\mathcal{G}_{\overline{y}}$，其中积取遍 $Y$ 的所有几何点。
由“\'Etale 上同调”定理 \ref{etale-cohomology-theorem-exactness-stalks}，
这是单射。（在拓扑空间上的阿贝尔层范畴中，这正是 Godement 分解的第一步。）

\medskip\noindent
如下可见 (b) 与引理的 (2)。假设对某些阿贝尔群 $A_{\overline{y}}$，有
$\mathcal{G} = \prod \overline{y}_*A_{\overline{y}}$。则
$$
\Hom(f_!\mathcal{F}, \mathcal{G}) =
\prod \Hom(f_!\mathcal{F}, \overline{y}_*A_{\overline{y}})
$$
。因此只须找到 $X_\etale$ 上的阿贝尔层 $\mathcal{H}_{\overline{y}}$，
它们表示函子
$\mathcal{F} \mapsto \Hom(f_!\mathcal{F}, \overline{y}_*A_{\overline{y}})$
，再取 $\mathcal{H} = \prod \mathcal{H}_{\overline{y}}$。这把问题化归到
$\mathcal{H} = \overline{y}_*A$ 的情形，其中
$\overline{y} : \Spec(k) \to Y$ 为某个固定几何点，$A$ 为某个固定阿贝尔群。
我们断言，此时取
$\mathcal{H} = \prod_{f(\overline{x}) = \overline{y}} \overline{x}_*A$ 即可。
这将完成引理 (1) 与 (2) 的证明。事实上，由引理
\ref{more-etale-lemma-lqf-f-shriek-stalk} 对茎的描述，一方面有
$$
\Hom(f_!\mathcal{F}, \overline{y}_*A) =
\Hom_{\textit{Ab}}((f_!\mathcal{F})_{\overline{y}}, A) =
\Hom_{\textit{Ab}}(\bigoplus\nolimits_{f(\overline{x}) = \overline{y}}
\mathcal{F}_{\overline{x}}, A)
$$
，另一方面有
$$
\Hom(\mathcal{F}, \mathcal{H}) =
\prod\nolimits_{f(\overline{x}) = \overline{y}}
\Hom(\mathcal{F}, \overline{x}_*A) =
\prod\nolimits_{f(\overline{x}) = \overline{y}}
\Hom_{\textit{Ab}}(\mathcal{F}_{\overline{x}}, A)
$$
。留给读者验证它们作为 $\mathcal{F}$ 的函子彼此等同。

\medskip\noindent
证明 (3)。注意，$\textit{Mod}(X_\etale, \Lambda)$ 的对象等同于
$\textit{Ab}(X_\etale)$ 的对象 $\mathcal{F}$ 连同映射
$\Lambda \to \text{End}(\mathcal{F})$。因此，(1) 中的函子 $f_!$ 与 $f^!$
定义了 (3) 中的函子 $f_!$ 与 $f^!$。直接计算可知它们互为伴随。
\end{proof}

\begin{lemma}
\label{more-etale-lemma-etale-upper-shriek}
设 $j : U \to X$ 为 \'etale 态射。则 $j^! = j^{-1}$。
\end{lemma}

\begin{proof}
这是因为第 \ref{more-etale-section-finite-support} 节定义的 $j_!$ 与“\'Etale 上同调”
第 \ref{etale-cohomology-section-extension-by-zero} 节定义的 $j_!$ 一致；
见引理 \ref{more-etale-lemma-finite-support-etale-shriek}。最后，在“\'Etale 上同调”
第 \ref{etale-cohomology-section-extension-by-zero} 节中，函子 $j_!$
定义为 $j^{-1}$ 的左伴随，故由伴随函子的唯一性得出结论。
\end{proof}

\begin{lemma}
\label{more-etale-lemma-upper-shriek-restriction}
设 $f : X \to Y$ 与 $g : Y \to Z$ 为分离且局部拟有限的态射。
存在典范同构 $(g \circ f)^! \to f^! \circ g^!$。给定第三个局部拟有限态射
$h : Z \to T$，图
$$
\xymatrix{
(h \circ g \circ f)^! \ar[r] \ar[d] &
f^! \circ (h \circ g)^! \ar[d] \\
(g \circ f)^! \circ h^! \ar[r] & f^! \circ g^! \circ h^!
}
$$
交换。
\end{lemma}

\begin{proof}
由伴随函子的唯一性，这立即转化为函子 $f_!$ 的相应（对偶）陈述。
见引理 \ref{more-etale-lemma-lqf-shriek-composition}。
\end{proof}

\begin{lemma}
\label{more-etale-lemma-upper-shriek-restriction-etale}
设 $j : U \to X$ 与 $j' : V \to U$ 为 \'etale 态射。同构
$(j \circ j')^{-1} = (j')^{-1} \circ j^{-1}$ 与引理
\ref{more-etale-lemma-upper-shriek-restriction} 的同构
$(j \circ j')^! = (j')^! \circ j^!$，通过引理
\ref{more-etale-lemma-etale-upper-shriek} 的同构彼此一致。
\end{lemma}

\begin{proof}
略。
\end{proof}

\begin{lemma}
\label{more-etale-lemma-lqf-base-change-upper-shriek}
考虑概形的笛卡尔方块
$$
\xymatrix{
X' \ar[r]_{g'} \ar[d]_{f'} & X \ar[d]^f \\
Y' \ar[r]^g & Y
}
$$
，其中 $f$ 局部拟有限。对 $Y'_\etale$ 上任意阿贝尔层 $\mathcal{F}$，有
$(g')_*(f')^!\mathcal{F} = f^!g_*\mathcal{F}$。
\end{lemma}

\begin{proof}
由伴随函子的唯一性，这由函子 $f_!$ 的相应（对偶）陈述得出。
见引理 \ref{more-etale-lemma-lqf-base-change-f-shriek}。
\end{proof}

\begin{remark}
\label{more-etale-remark-pointed-sets}
本节内容可推广到带基点集合的层。具体地，对位点 $\mathcal{C}$，以
$\Sh^*(\mathcal{C})$ 表示带基点集合的层之范畴。本节及前一节的构造经相应修改后
适用于带基点集合的层。因此，给定概形的局部拟有限态射 $f : X \to Y$，
得到一对伴随函子
$$
f_! : \Sh^*(X_\etale) \longrightarrow \Sh^*(Y_\etale)
\quad\text{且}\quad
f^! : \Sh^*(Y_\etale) \longrightarrow \Sh^*(X_\etale)
$$
，使得对 $Y$ 的每个几何点 $\overline{y}$，都有同构
$$
(f_!\mathcal{F})_{\overline{y}} =
\coprod\nolimits_{f(\overline{x}) = \overline{y}}
\mathcal{F}_{\overline{x}}
$$
（余积取于带基点集合的范畴中），它关于
$\mathcal{F} \in \Sh^*(X_\etale)$ 具有函子性；还有同构
$$
f^!(\overline{y}_*S) =
\prod\nolimits_{f(\overline{x}) = \overline{y}}
\overline{x}_*S
$$
，它关于带基点集合 $S$ 具有函子性。若
$F : \textit{Ab}(X_\etale) \to \Sh^*(X_\etale)$ 与
$F : \textit{Ab}(Y_\etale) \to \Sh^*(Y_\etale)$
表示遗忘函子，则各构造之间的相容性保证存在典范映射
$$
f_!F(\mathcal{F}) \longrightarrow F(f_!\mathcal{F})
$$
，它关于 $\mathcal{F} \in \textit{Ab}(X_\etale)$ 具有函子性；以及
$$
F(f^!\mathcal{G}) \longrightarrow f^!F(\mathcal{G})
$$
，它关于 $\mathcal{G} \in \textit{Ab}(Y_\etale)$ 具有函子性；二者分别在茎上、
在摩天层上给出显然的映射。事实上，变换
$F \circ f^! \to f^! \circ F$ 是同构（因为 $f^!$ 与积可交换）。
\end{remark}







\section{局部拟有限态射的导出上叹号}
\label{more-etale-section-derived-duality-locally-quasi-finite}

\noindent
取第 \ref{more-etale-section-duality-locally-quasi-finite} 节中各函子的导出版本，
可得如下结论。

\begin{lemma}
\label{more-etale-lemma-lqf-shriek-derived}
设 $f : X \to Y$ 是概形的局部拟有限态射，$\Lambda$ 是环。
定义 \ref{more-etale-definition-f-shriek-lqf} 与引理
\ref{more-etale-lemma-lqf-f-upper-shriek} 中的函子 $f_!$ 和 $f^!$
在导出范畴上诱导出一对伴随函子
$f_! : D(X_\etale, \Lambda) \to D(Y_\etale, \Lambda)$ 与
$Rf^! : D(Y_\etale, \Lambda) \to D(X_\etale, \Lambda)$。
\end{lemma}

\noindent
在分离情形，函子 $f_!$ 定义于第 \ref{more-etale-section-compact-support} 节。

\begin{proof}
这立即由“导出范畴”引理
\ref{derived-lemma-derived-adjoint-functors}、$f_!$ 的正合性
（引理 \ref{more-etale-lemma-lqf-f-shriek-stalk}，因而 $Lf_! = f_!$）以及如下事实得出：
$Y_\etale$ 上有足够多的 $\Lambda$-模 K-内射复形，故 $Rf^!$ 有定义。
\end{proof}

\begin{remark}
\label{more-etale-remark-lqf-shriek-derived-torsion}
设 $f : X \to Y$ 是概形的局部拟有限态射，$\Lambda$ 是环。
引理 \ref{more-etale-lemma-lqf-shriek-derived} 中的函子
$f_! : D(X_\etale, \Lambda) \to D(Y_\etale, \Lambda)$
把上同调层为挠层的复形送到上同调层为挠层的复形。
由 $f_!$ 的茎的描述，此事立即成立；见引理
\ref{more-etale-lemma-lqf-f-shriek-stalk}。
\end{remark}

\begin{lemma}
\label{more-etale-lemma-mayer-vietoris-shriek}
设 $X$ 是概形，$X = U \cup V$，其中 $U$ 与 $V$ 均为开集。
设 $\Lambda$ 是环且 $K \in D(X_\etale, \Lambda)$。则有特异三角
$$
j_{U \cap V!}K|_{U \cap V} \to
j_{U!}K|_U \oplus j_{V!}K|_V \to K \to 
j_{U \cap V!}K|_{U \cap V}[1]
$$
成立于 $D(X_\etale, \Lambda)$ 中，其中记号含义显然。
\end{lemma}

\begin{proof}
由于这里所用的限制函子与下叹号函子均正合，只需证明：对
$X_\etale$ 上任意阿贝尔层 $\mathcal{F}$，序列
$$
0 \to j_{U \cap V!}\mathcal{F}|_{U \cap V} \to
j_{U!}\mathcal{F}|_U \oplus j_{V!}\mathcal{F}|_V \to \mathcal{F} \to 0
$$
正合。这可在茎上验证。
\end{proof}

\begin{lemma}
\label{more-etale-lemma-triangle-associated-to-open}
设 $X$ 是概形，$Z \subset X$ 是闭子概形，$U \subset X$ 是其补集。
以 $i : Z \to X$ 和 $j : U \to X$ 表示嵌入态射。设 $\Lambda$ 是环且
$K \in D(X_\etale, \Lambda)$。则有特异三角
$$
j_!j^{-1}K \to K \to i_*i^{-1}K \to j_!j^{-1}K[1]
$$
成立于 $D(X_\etale, \Lambda)$ 中。
\end{lemma}

\begin{proof}
这立即由“\'Etale 上同调”引理
\ref{etale-cohomology-lemma-ses-associated-to-open} 以及如下事实得出：
函子 $j_!$、$j^{-1}$、$i_*$、$i^{-1}$ 均正合，故其导出版本可通过将这些函子
作用于任意一个表示 $K$ 的层复形来计算。
\end{proof}





\section{通过紧化构造导出下叹号的预备结果}
\label{more-etale-section-prelim}

\noindent
本节证明若干关于函子之间某些自然同构之存在性的引理；这些同构
立即来自固有基变换。

\begin{lemma}
\label{more-etale-lemma-shriek-proper-and-open}
考虑概形的交换图
$$
\xymatrix{
X' \ar[r]_{g'} \ar[d]_{f'} & X \ar[d]^f \\
Y' \ar[r]^g & Y
}
$$
其中 $f$ 与 $f'$ 固有，而 $g$ 与 $g'$ 分离且局部拟有限。
设 $\Lambda$ 是环。关于 $K \in D(X'_\etale, \Lambda)$ 函子性地，
在 $D(Y_\etale, \Lambda)$ 中有典范映射
$$
g_!Rf'_*K \longrightarrow Rf_*(g'_!K)
$$
。在以下任一情形，该映射是同构：
(a) $K$ 下有界且其上同调层为挠层；或
(b) $\Lambda$ 是挠环。
\end{lemma}

\begin{proof}
以 $X'_\etale$ 上的 $\Lambda$-模层 K-内射复形
$\mathcal{J}^\bullet$ 表示 $K$。选取到 $X_\etale$ 上的
$\Lambda$-模层 K-内射复形 $\mathcal{I}^\bullet$ 的拟同构
$g'_!\mathcal{J}^\bullet \to \mathcal{I}^\bullet$。于是可考虑映射
$$
g_!f'_*\mathcal{J}^\bullet =
g_!f'_!\mathcal{J}^\bullet =
f_!g'_!\mathcal{J}^\bullet =
f_*g'_!\mathcal{J}^\bullet \to
f_*\mathcal{I}^\bullet
$$
其中第一与第三个等号来自引理 \ref{more-etale-lemma-proper-f-shriek}，
第二个等号来自引理 \ref{more-etale-lemma-f-shriek-composition}；该引理告诉我们，
$g_! \circ f'_!$ 与 $f_! \circ g'_!$ 作为
$(g \circ f')_* = (f \circ g')_*$ 的子层，均等于
$(g \circ f')_! = (f \circ g')_!$。

\medskip\noindent
设 $\Lambda$ 为挠环，即处于情形 (b)。沿用上述记号，只需证明
$f_*g'_!\mathcal{J}^\bullet \to f_*\mathcal{I}^\bullet$ 是同构。
该问题在 $Y$ 上是局部的，故可设 $f$ 的纤维维数有界；见“态射”引理
\ref{morphisms-lemma-morphism-finite-type-bounded-dimension}。
于是 $Rf_*$ 的上同调维数有限；见“\'Etale 上同调”引理
\ref{etale-cohomology-lemma-cohomological-dimension-proper}。
因此由“导出范畴”引理 \ref{derived-lemma-unbounded-right-derived}，
若能证明对 $q > 0$ 以及 $X'_\etale$ 上任意内射 $\Lambda$-模层
$\mathcal{J}$，均有 $R^qf_*(g'_!\mathcal{J}) = 0$，结论即得。

\medskip\noindent
由“\'Etale 上同调”引理
\ref{etale-cohomology-lemma-proper-base-change-stalk}，
$R^qf_*(g'_!\mathcal{J})$ 在几何点 $\overline{y}$ 处的茎等于
$H^q(X_{\overline{y}}, (g'_!\mathcal{J})|_{X_{\overline{y}}})$。
由于 $g'_!$ 的构造与基变换交换
（引理 \ref{more-etale-lemma-base-change-f-shriek-separated}），它等于
$$
H^q(X_{\overline{y}}, g'_{\overline{y}, !}(\mathcal{J}|_{X'_{\overline{y}}}))
$$
，其中 $g'_{\overline{y}} : X'_{\overline{y}} \to X_{\overline{y}}$
是几何纤维之间所诱导的态射。由于 $Y' \to Y$ 局部拟有限，
$X'_{\overline{y}}$ 是各纤维 $X'_{\overline{y}'}$ 的不交并；这里
$\overline{y}'$ 遍历 $Y'$ 中位于 $\overline{y}$ 上方的几何点。
记 $g'_{\overline{y}'} : X'_{\overline{y}'} \to X_{\overline{y}}$
为 $g'_{\overline{y}}$ 在 $X'_{\overline{y}'}$ 上的限制。于是前述上同调群等于
$$
H^q(X_{\overline{y}},
\bigoplus\nolimits_{\overline{y}'/\overline{y}}
g'_{\overline{y}', !}(\mathcal{J}|_{X'_{\overline{y}'}}))
$$
；这例如可由引理 \ref{more-etale-lemma-colim-f-shriek-separated} 得到（但从第
\ref{more-etale-section-compact-support} 节中 $g'_{\overline{y}, !}$ 的定义也显然可见）。
由于在 $X_{\overline{y}}$ 上取 \'etale 上同调与直和交换
（“\'Etale 上同调”定理 \ref{etale-cohomology-theorem-colimit}），
我们归结为证明
$$
H^q(X_{\overline{y}}, g'_{\overline{y}', !}(\mathcal{J}|_{X'_{\overline{y}'}}))
$$
为零。注意，$g_{\overline{y}'} : X'_{\overline{y}'} \to X_{\overline{y}}$
是 $\overline{y}$ 上固有概形之间的态射，故其本身固有。它又局部拟有限，
因而 $g_{\overline{y}'}$ 有限。于是
$g'_{\overline{y}', !} = g'_{\overline{y}', *} = Rg'_{\overline{y}', *}$。
由 Leray，所需证明化为
$$
H^q(X'_{\overline{y}'}, \mathcal{J}|_{X'_{\overline{y}'}})
$$
为零。由于 $\Lambda$ 为挠环，这由固有基变换（“\'Etale 上同调”引理
\ref{etale-cohomology-lemma-proper-base-change-stalk}）得出，因为
$\mathcal{J}$ 在 $f'$ 下的高阶直像为零。

\medskip\noindent
现证情形 (a)。我们用标准论证从情形 (b) 推出它。我们将证明，对所有
$p \in \mathbf{Z}$，所诱导的映射
$g_! R^pf'_* K \to R^pf_*(g'_!K)$ 均为同构。固定整数
$p_0 \in \mathbf{Z}$，并取整数 $a$，使得当 $j < a$ 时 $H^j(K) = 0$。
我们对 $a$ 作降序归纳，证明当 $p \leq p_0$ 时
$g_! R^pf'_* K \to R^pf_*(g'_!K)$ 为同构。若 $a > p_0$，则由平凡消没，
上述映射的左右两端在 $p \leq p_0$ 时均为零；见“导出范畴”引理
\ref{derived-lemma-negative-vanishing}（并使用 $g_!$ 与 $g'_!$ 均为正合函子）。
设 $a \leq p_0$。考虑特异三角
$$
H^a(K)[-a] \to K \to \tau_{\geq a + 1}K
$$
由归纳假设，结论对 $\tau_{\geq a + 1}K$ 成立。下一段将证明结论对
$H^a(K)[-a]$ 成立。随后，对如下特异三角映射所伴的上同调层长正合列之间的
映射应用五引理
$$
\xymatrix{
g_! Rf'_*(H^a(K)[-a]) \ar[d] \ar[r] &
g_! Rf'_* K \ar[r] \ar[d] &
g_! Rf'_* \tau_{\geq a + 1} K \ar[d] \\
Rf_*(g'_!(H^a(K)[-a])) \ar[r] &
Rf_*(g'_!K) \ar[r] &
Rf_*(g'_!\tau_{|geq a + 1}K)
}
$$
，即得 $K$ 的结论。略去一些细节。

\medskip\noindent
设 $\mathcal{F}$ 是 $X'_\etale$ 上的挠阿贝尔层。为完成证明，我们证明对所有
$p$，映射 $g_! Rf'_*\mathcal{F} \to R^pf_*(g'_!\mathcal{F})$ 均为同构。
可写成 $\mathcal{F} = \bigcup \mathcal{F}[n]$，其中
$\mathcal{F}[n] = \Ker(n : \mathcal{F} \to \mathcal{F})$。
由情形 (b)，同构对 $\mathcal{F}[n]$ 成立。函子 $g_!$、$g'_!$、
$R^pf_*$、$R^pf'_*$ 均与滤过余极限交换（这由引理
\ref{more-etale-lemma-lqf-f-shriek-separated-colimits} 和“\'Etale 上同调”引理
\ref{etale-cohomology-lemma-relative-colimit-general} 得出），故证明完成。
\end{proof}

\begin{lemma}
\label{more-etale-lemma-shriek-proper-and-open-compose}
考虑概形的交换图
$$
\xymatrix{
X' \ar[r]_k \ar[d]_{f'} & X \ar[d]^f \\
Y' \ar[r]_l \ar[d]_{g'} & Y \ar[d]^g \\
Z' \ar[r]^m & Z
}
$$
其中 $f$、$f'$、$g$、$g'$ 固有，而 $k$、$l$、$m$ 分离且局部拟有限。
则引理 \ref{more-etale-lemma-shriek-proper-and-open} 对两个方块给出的同构复合后，
给出外侧矩形的同构（精确陈述见证明）。
\end{lemma}

\begin{proof}
该陈述的含义如下。若写成
$R(g \circ f)_* = Rg_* \circ Rf_*$ 与
$R(g' \circ f')_* = Rg'_* \circ Rf'_*$,
，则外侧矩形的同构
$m_! \circ Rg'_* \circ Rf'_* \to Rg_* \circ Rf_* \circ k_!$
等于两个方块的映射之复合
$$
m_! \circ Rg'_* \circ Rf'_* \to
Rg_* \circ l_! \circ Rf'_* \to
Rg_* \circ Rf_* \circ k_!
$$
。为证明此事，选取 $X'_\etale$ 上的 $\Lambda$-模 K-内射复形
$\mathcal{J}^\bullet$，以及到 $X_\etale$ 上的 $\Lambda$-模 K-内射复形
$\mathcal{I}^\bullet$ 的拟同构
$k_!\mathcal{J}^\bullet \to \mathcal{I}^\bullet$。
引理 \ref{more-etale-lemma-shriek-proper-and-open} 的证明表明，典范映射
$$
a : l_!f'_*\mathcal{J}^\bullet \to f_*\mathcal{I}^\bullet
$$
是拟同构，对它应用 $Rg_*$ 即产生第二个箭头。由“位点上同调”引理
\ref{sites-cohomology-lemma-K-injective-flat}，复形
$f_*\mathcal{I}^\bullet$（相应地，$f'_*\mathcal{J}^\bullet$）是
$Y_\etale$（相应地，$Y'_\etale$）上的 $\Lambda$-模 K-内射复形。
（使用这一点算是取巧，本可避免。）特别地，同一论证给出典范映射
$$
b : m_!g'_*f'_*\mathcal{J}^\bullet \to g_*f_*\mathcal{I}^\bullet
$$
也是拟同构，且该拟同构表示第一个箭头。最后，由于
$f'_*\mathcal{J}^\bullet$ 是 K-内射的，引理
\ref{more-etale-lemma-shriek-proper-and-open} 的证明表明
$g_*l_!f'_!\mathcal{J}^\bullet$ 表示
$Rg_*(l_!f'_*\mathcal{J}^\bullet)$。因此 $Rg_*(a) = g_*(a)$，且复合
$g_*(a) \circ b$ 正是引理 \ref{more-etale-lemma-shriek-proper-and-open}
对该矩形给出的箭头。
\end{proof}

\begin{lemma}
\label{more-etale-lemma-shriek-proper-and-open-compose-horizontal}
考虑概形的交换图
$$
\xymatrix{
X'' \ar[r]_{g'} \ar[d]_{f''} &
X' \ar[r]_g \ar[d]_{f'} &
X \ar[d]^f \\
Y'' \ar[r]^{h'} &
Y' \ar[r]^h &
Y
}
$$
其中 $f$、$f'$、$f''$ 固有，而 $g$、$g'$、$h$、$h'$ 分离且局部拟有限。
则引理 \ref{more-etale-lemma-shriek-proper-and-open} 对两个方块给出的同构复合后，
给出外侧矩形的同构（精确陈述见证明）。
\end{lemma}

\begin{proof}
该陈述的含义如下。若用引理 \ref{more-etale-lemma-f-shriek-composition} 的等式写成
$(h \circ h')_! = h_! \circ h'_!$ 与
$(g \circ g')_! = g_! \circ g'_!$
，则外侧矩形的同构
$h_! \circ h'_! \circ Rf''_* \to Rf_* \circ g_! \circ g'_!$
等于两个方块的映射之复合
$$
h_! \circ h'_! \circ Rf''_* \to
h_! \circ Rf'_* \circ g'_! \to
Rf_* \circ g_! \circ g'_!
$$
。为证明此事，选取 $X''_\etale$ 上的 $\Lambda$-模 K-内射复形
$\mathcal{I}^\bullet$，以及到 $X'_\etale$ 上的 $\Lambda$-模 K-内射复形
$\mathcal{J}^\bullet$ 的拟同构
$g'_!\mathcal{I}^\bullet \to \mathcal{J}^\bullet$。继而选取到
$X_\etale$ 上的 $\Lambda$-模 K-内射复形 $\mathcal{K}^\bullet$ 的拟同构
$g_!\mathcal{J}^\bullet \to \mathcal{K}^\bullet$。
引理 \ref{more-etale-lemma-shriek-proper-and-open} 的证明表明，典范映射
$$
h'_!f''_*\mathcal{I}^\bullet \to f'_*\mathcal{J}^\bullet
\quad\text{且}\quad
h_!f'_*\mathcal{J}^\bullet \to f_*\mathcal{K}^\bullet
$$
均为拟同构，且分别定义上述第一与第二个箭头。由于 $g_!$ 是正合函子
（引理 \ref{more-etale-lemma-lqf-f-shriek-separated-colimits}），
$g_!g'_!\mathcal{I}^\bullet \to \mathcal{K}^\bullet$ 是拟同构，
因而典范映射
$$
h_!h'_!f''_*\mathcal{I}^\bullet \to f_*\mathcal{K}^\bullet
$$
也是拟同构，并在导出范畴中表示外侧矩形的映射。显然，该映射是另两个映射的
复合，证明完成。
\end{proof}

\begin{remark}
\label{more-etale-remark-going-around}
考虑交换图
$$
\xymatrix{
X'' \ar[r]_{k'} \ar[d]_{f''} & X' \ar[r]_k \ar[d]_{f'} & X \ar[d]^f \\
Y'' \ar[r]^{l'} \ar[d]_{g''} & Y' \ar[r]^l \ar[d]_{g'} & Y \ar[d]^g \\
Z'' \ar[r]^{m'} & Z' \ar[r]^m & Z
}
$$
，其中竖箭头均固有，横箭头均分离且局部拟有限。如下把图中的方块标记为
$A$、$B$、$C$、$D$：
$$
\begin{matrix}
A & B \\
C & D
\end{matrix}
$$
于是，引理 \ref{more-etale-lemma-shriek-proper-and-open} 对各方块给出的映射为
（这里使用 $Rf_* = f_*$ 等）
$$
\begin{matrix}
\gamma_A : l'_! \circ f''_* \to f'_* \circ k'_! &
\gamma_B : l_! \circ f'_* \to f_* \circ k_! \\
\gamma_C : m'_! \circ g''_* \to g'_* \circ l'_! &
\gamma_D : m_! \circ g'_* \to g_* \circ l_!
\end{matrix}
$$
对各 $2 \times 1$ 与 $1 \times 2$ 矩形，另有四个映射
$$
\begin{matrix}
\gamma_{A + B} :
(l \circ l')_! \circ f''_* \to f_* \circ (k \circ k')_*  \\
\gamma_{C + D} :
(m \circ m')_! \circ g''_* \to g_* \circ (l \circ l')_! \\
\gamma_{A + C} :
m'_! \circ (g'' \circ f'')_* \to (g' \circ f')_* \circ k'_! \\
\gamma_{B + D} :
m_! \circ (g' \circ f')_* \to (g \circ f)_* \circ k_!
\end{matrix}
$$
由引理 \ref{more-etale-lemma-shriek-proper-and-open-compose-horizontal}，有
$$
\gamma_{A + B} = \gamma_B \circ \gamma_A, \quad
\gamma_{C + D} = \gamma_D \circ \gamma_C
$$
；由引理 \ref{more-etale-lemma-shriek-proper-and-open-compose}，有
$$
\gamma_{A + C} = \gamma_A \circ \gamma_C, \quad
\gamma_{B + D} = \gamma_B \circ \gamma_D
$$
更准确地说，此处应写成
$\gamma_{A + B} = (\gamma_B \star \text{id}_{k'_!}) \circ
(\text{id}_{l_!} \star \gamma_A)$，其中采用
“范畴”第 \ref{categories-section-formal-cat-cat} 节的记号
；其余情形亦同。综上，我们（先验地）得到两个变换
$$
m_! \circ m'_! \circ g''_* \circ f''_*
\longrightarrow
g_* \circ f_* \circ k_! \circ k'_!
$$
，即
$$
\gamma_B \circ \gamma_D \circ \gamma_A \circ \gamma_C =
\gamma_{B + D} \circ \gamma_{A + C}
$$
以及
$$
\gamma_B \circ \gamma_A \circ \gamma_D \circ \gamma_C =
\gamma_{A + B} \circ \gamma_{C + D}
$$
本评注的要点是指出这两个变换相等。为此，只需证明
$$
\xymatrix{
m_! \circ g'_* \circ l'_! \circ f''_* \ar[r]_{\gamma_D} \ar[d]_{\gamma_A} &
g_* \circ l_! \circ l'_! \circ f''_* \ar[d]^{\gamma_A} \\
m_! \circ g'_* \circ f'_* \circ k'_! \ar[r]^{\gamma_D} &
g_* \circ l_! \circ f'_* \circ k'_!
}
$$
交换。其原因是方块 $A$ 与 $D$ 仅在一点相交；更准确地说，这由“范畴”引理
\ref{categories-lemma-properties-2-cat-cats} 得出，或者更直接地，由“范畴”定义
\ref{categories-definition-horizontal-composition} 之前的讨论得出。
\end{remark}

\begin{lemma}
\label{more-etale-lemma-shriek-proper-and-open-base-change}
设 $b : Y_1 \to Y$ 是概形的态射。考虑概形的交换图
$$
\vcenter{
\xymatrix{
X' \ar[r]_{g'} \ar[d]_{f'} & X \ar[d]^f \\
Y' \ar[r]^g & Y
}
}
\quad\text{且令}\quad
\vcenter{
\xymatrix{
X'_1 \ar[r]_{g'_1} \ar[d]_{f'_1} & X_1 \ar[d]^{f_1} \\
Y'_1 \ar[r]^{g_1} & Y_1
}
}
$$
为沿 $b$ 的基变换。设 $f$ 与 $f'$ 固有，而 $g$ 与 $g'$ 分离且局部拟有限。
对环 $\Lambda$ 以及 $D(X'_\etale, \Lambda)$ 中的 $K$，有交换图
$$
\xymatrix{
b^{-1}g_!Rf'_*K \ar[d] \ar[r] &
g_{1, !}(b')^{-1}Rf'_*K \ar[r] &
g_{1, !}Rf'_{1, *}(a')^{-1}K \ar[d] \\
b^{-1}Rf_*g'_!K \ar[r] &
Rf_{1, *}a^{-1}g'_!K \ar[r] &
Rf_{1, *}g'_{1, !}(a')^{-1}K
}
$$
成立于 $D(Y_{1, \etale}, \Lambda)$ 中；其中
$a : X_1 \to X$、$a' : X'_1 \to X'$、$b' : Y'_1 \to Y'$ 是投影，
竖映射是引理 \ref{more-etale-lemma-shriek-proper-and-open} 的箭头，横箭头是基变换映射
（来自“\'Etale 上同调”第
\ref{etale-cohomology-section-base-change-preliminaries} 节）与引理
\ref{more-etale-lemma-base-change-f-shriek-separated} 的基变换映射。
\end{lemma}

\begin{proof}
以 $X'_\etale$ 上的 $\Lambda$-模层 K-内射复形
$\mathcal{J}^\bullet$ 表示 $K$。选取到 $X_\etale$ 上的
$\Lambda$-模层 K-内射复形 $\mathcal{I}^\bullet$ 的拟同构
$g'_!\mathcal{J}^\bullet \to \mathcal{I}^\bullet$。
引理 \ref{more-etale-lemma-shriek-proper-and-open} 的证明把
$g_!Rf'_*K \to Rf_*g'_!K$ 构造为
$$
g_!f'_*\mathcal{J}^\bullet =
g_!f'_!\mathcal{J}^\bullet =
f_!g'_!\mathcal{J}^\bullet =
f_*g'_!\mathcal{J}^\bullet \to
f_*\mathcal{I}^\bullet
$$
选取到 $X'_{1, \etale}$ 上的 $\Lambda$-模层 K-内射复形
$\mathcal{J}_1^\bullet$ 的拟同构
$(a')^{-1}\mathcal{J}^\bullet \to \mathcal{J}_1^\bullet$。
于是可选取复形图
$$
\xymatrix{
g'_{1, !}\mathcal{J}_1^\bullet \ar[rr] & &
\mathcal{I}_1^\bullet \\
g'_{1, !}(a')^{-1}\mathcal{J}^\bullet \ar[u] \ar@{=}[r] &
a^{-1}g'_!\mathcal{J}^\bullet \ar[r] &
a^{-1}\mathcal{I}^\bullet \ar[u]
}
$$
，它在同伦意义下交换，且所有箭头均为拟同构；等号来自引理
\ref{more-etale-lemma-proper-f-shriek}，而 $\mathcal{I}_1^\bullet$ 是
$X_{1, \etale}$ 上的 $\Lambda$-模层 K-内射复形。映射
$g_{1, !}Rf'_{1, *}(a')^{-1}K \to Rf_{1, *}g'_{1, !}(a')^{-1}K$ 由下式给出：
$$
g_{1, !}f'_{1, *}\mathcal{J}_1^\bullet =
g_{1, !}f'_{1, !}\mathcal{J}_1^\bullet =
f_{1, !}g'_{1, !}\mathcal{J}_1^\bullet =
f_{1, *}g'_{1, !}\mathcal{J}_1^\bullet \to
f_{1, *}\mathcal{I}_1^\bullet
$$
两个箭头中跨越这 $3$ 个等号的等同均与拉回映射相容；换言之，
阿贝尔层复形的图
$$
\xymatrix{
b^{-1}g_!f'_*\mathcal{J}^\bullet \ar@{=}[d] \ar[r] &
g_{1, !}(b')^{-1}f'_*\mathcal{J}^\bullet \ar[r] &
g_{1, !}f'_{1, *}(a')^{-1}\mathcal{J}^\bullet \ar@{=}[d] \\
b^{-1}f_*g'_!\mathcal{J}^\bullet \ar[r] &
f_{1, *}a^{-1}g'_!\mathcal{J}^\bullet \ar[r] &
f_{1, *}g'_{1, !}(a')^{-1}\mathcal{J}^\bullet
}
$$
交换。要证明这一点，只需把图中的 $g_!, g_{1, !}, g'_!, g'_{1, !}$
分别替换为 $g_*, g_{1, *}, g'_*, g'_{1, *}$ 后证明其交换；这是因为下叹号
函子被定义为 $*$ 函子的子函子，而基变换映射的定义与此相容，见引理
\ref{more-etale-lemma-base-change-f-shriek-separated} 的证明。对这个新图，交换性来自
拉回映射与图的横向、纵向叠合之相容性；见“位点”评注
\ref{sites-remark-compose-base-change} 与
\ref{sites-remark-compose-base-change-horizontal}。因此沿图任一方向绕行，
得到的都是 $f \circ g' = g \circ f'$ 沿 $b$ 基变换的拉回映射。
当然，下图交换：
$$
\xymatrix{
g_{1, !}f'_{1, *}(a')^{-1}\mathcal{J}^\bullet \ar@{=}[d] \ar[r] &
g_{1, !}f'_{1, *}\mathcal{J}_1^\bullet \ar@{=}[d] \\
f_{1, *}g'_{1, !}(a')^{-1}\mathcal{J}^\bullet \ar[r] &
f_{1, *}g'_{1, !}\mathcal{J}_1^\bullet 
}
$$
故要完成证明，只需证明
$$
\xymatrix{
b^{-1}f_*g'_!\mathcal{J}^\bullet \ar[r] \ar[d] &
f_{1, *}a^{-1}g'_!\mathcal{J}^\bullet \ar[r] \ar[d] &
f_{1, *}g'_{1, !}(a')^{-1}\mathcal{J}^\bullet \ar[r] &
f_{1, *}g'_{1, !}\mathcal{J}_1^\bullet \ar[d] \\
b^{-1}f_*\mathcal{I}^\bullet \ar[r] &
f_{1, *}a^{-1}\mathcal{I}^\bullet \ar[rr] & &
f_{1, *}\mathcal{I}_1^\bullet
}
$$
在导出范畴中交换；这由我们先前对映射的选取成立。
\end{proof}

\begin{lemma}
\label{more-etale-lemma-shriek-lqf-and-proper}
考虑概形的交换图
$$
\xymatrix{
X \ar[r]_f \ar[rd]_g & Y \ar[d]^h \\
& Z
}
$$
，其中 $f$ 与 $g$ 局部拟有限，而 $h$ 固有。设 $\Lambda$ 是环。
关于 $K \in D(X_\etale, \Lambda)$ 函子性地，在
$D(Z_\etale, \Lambda)$ 中有典范映射
$$
g_!K \longrightarrow Rh_*(f_!K)
$$
。在以下任一情形，该映射是同构：
(a) $K$ 下有界且其上同调层为挠层；或
(b) $\Lambda$ 是挠环。
\end{lemma}

\begin{proof}
若 $f$ 与 $g$ 分离，这是引理 \ref{more-etale-lemma-shriek-proper-and-open} 的特例。
我们建议读者跳过一般情形的证明，因为我们主要使用 $f$ 与 $g$ 分离的情形。

\medskip\noindent
以 $X_\etale$ 上的 $\Lambda$-模层复形 $\mathcal{K}^\bullet$ 表示 $K$。
选取到 $Y_\etale$ 上的 $\Lambda$-模层 K-内射复形
$\mathcal{I}^\bullet$ 的拟同构
$f_!\mathcal{K}^\bullet \to \mathcal{I}^\bullet$。考虑映射
$$
g_!\mathcal{K}^\bullet =
h_!f_!\mathcal{K}^\bullet =
h_*f_!\mathcal{K}^\bullet
\longrightarrow
h_*\mathcal{I}^\bullet
$$
，其中各等号由引理 \ref{more-etale-lemma-lqf-separated-shriek-composition}
和 \ref{more-etale-lemma-proper-f-shriek} 给出。这个复形映射确定引理陈述中的映射
$g_!K \to Rh_*(f_!K)$。

\medskip\noindent
设 $\Lambda$ 为挠环，即处于情形 (b)。为检验该映射是同构，可在 $Z$ 上
局部地工作。因此可设 $h$ 的纤维维数有界；见“态射”引理
\ref{morphisms-lemma-morphism-finite-type-bounded-dimension}。
于是 $Rh_*$ 的上同调维数有限；见“\'Etale 上同调”引理
\ref{etale-cohomology-lemma-cohomological-dimension-proper}。
故由“导出范畴”引理 \ref{derived-lemma-unbounded-right-derived}，
若能证明对 $q > 0$ 以及 $X_\etale$ 上任意 $\Lambda$-模层
$\mathcal{F}$，均有 $R^qh_*(f_!\mathcal{F}) = 0$，则
$h_*f_!\mathcal{K}^\bullet \to h_*\mathcal{I}^\bullet$ 是拟同构。

\medskip\noindent
注意，$\mathcal{G} = f_!\mathcal{F}$ 是 $Y$ 上的 $\Lambda$-模层；其茎仅可能在
满足 $\kappa(y)/\kappa(h(y))$ 为有限扩张的点 $y \in Y$ 处非零。这由引理
\ref{more-etale-lemma-lqf-f-shriek-stalk} 中对 $f_!\mathcal{F}$ 的茎的描述，以及
$f$ 与 $g$ 均局部拟有限这一事实得出。因此，由固有基变换定理
（“\'Etale 上同调”引理
\ref{etale-cohomology-lemma-proper-base-change-stalk}），只需证明
$H^q(Y_{\overline{z}}, \mathcal{H}) = 0$；这里 $\mathcal{H}$ 是
$\kappa(\overline{z})$ 上固有概形 $Y_{\overline{z}}$ 的一个层，且其支撑包含于
闭点集。所需消没于是由“\'Etale 上同调”引理
\ref{etale-cohomology-lemma-supported-in-closed-points} 得出。

\medskip\noindent
情形 (a) 可用引理 \ref{more-etale-lemma-shriek-proper-and-open} 的证明中完全相同的论证
从情形 (b) 推出（以引理 \ref{more-etale-lemma-lqf-f-shriek-stalk} 代替引理
\ref{more-etale-lemma-lqf-f-shriek-separated-colimits}）。
\end{proof}







\section{通过紧化构造导出下叹号}
\label{more-etale-section-derived-lower-shriek-compactification}

\noindent
设 $f : X \to Y$ 是概形的有限型分离态射，其中 $Y$ 拟紧且拟分离。
选取在 $Y$ 上的紧化 $j : X \to \overline{X}$；见“平坦性的进一步结果”
定理 \ref{flat-theorem-nagata}。设 $\Lambda$ 是环。以
$D^+_{tors}(X_\etale, \Lambda)$ 表示 $D(X_\etale, \Lambda)$ 的严格全且
饱和的三角子范畴，其对象为下有界且上同调层为挠层的对象 $K$。
我们将考虑函子
$$
Rf_! = R\overline{f}_* \circ j_! :
D^+_{tors}(X_\etale, \Lambda)
\longrightarrow
D^+_{tors}(Y_\etale, \Lambda)
$$
，其中 $\overline{f} : \overline{X} \to Y$ 是结构态射。该定义是有意义的：
由评注 \ref{more-etale-remark-lqf-shriek-derived-torsion}，函子 $j_!$ 把
$D^+_{tors}(X_\etale, \Lambda)$ 送入
$D^+_{tors}(\overline{X}_\etale, \Lambda)$；由“\'Etale 上同调”引理
\ref{etale-cohomology-lemma-torsion-direct-image}，$R\overline{f}_*$ 把
$D^+_{tors}(\overline{X}_\etale, \Lambda)$ 送入
$D^+_{tors}(Y_\etale, \Lambda)$。若 $\Lambda$ 是挠环，则定义
$$
Rf_! = R\overline{f}_* \circ j_! :
D(X_\etale, \Lambda)
\longrightarrow
D(Y_\etale, \Lambda)
$$
下面是必不可少的引理。

\begin{lemma}
\label{more-etale-lemma-shriek-well-defined}
设 $f : X \to Y$ 是拟紧拟分离概形之间的有限型分离态射。上述构造的函子
$Rf_!$ 在典范同构意义下与紧化的选取无关。
\end{lemma}

\begin{proof}
我们将在 $\Lambda$ 为挠环时对函子
$Rf_! : D(X_\etale, \Lambda) \to D(Y_\etale, \Lambda)$ 证明此事；
函子 $Rf_! : D^+_{tors}(X_\etale, \Lambda) \to
D^+_{tors}(Y_\etale, \Lambda)$ 的情形完全同样地证明。

\medskip\noindent
考虑 $X$ 在 $Y$ 上的紧化所成的范畴。由“平坦性的进一步结果”定理
\ref{flat-theorem-nagata} 以及引理
\ref{flat-lemma-compactifications-cofiltered} 和
\ref{flat-lemma-compactifyable}，该范畴是余滤过的。对紧化的每个选取
$$
j : X \to \overline{X},\quad \overline{f} : \overline{X} \to Y
$$
，上述构造配以函子 $R\overline{f}_* \circ j_! :
D(X_\etale, \Lambda) \to D(Y_\etale, \Lambda)$。更具体一些：给定
$X_\etale$ 上的 $\Lambda$-模层复形 $\mathcal{K}^\bullet$，我们选取一个
到 $\overline{X}_\etale$ 上某个 $\Lambda$-模层 K-内射复形的拟同构
$j_!\mathcal{K}^\bullet \to \mathcal{I}^\bullet$。我们的函子遂把
$\mathcal{K}^\bullet$ 送到 $\overline{f}_*\mathcal{I}^\bullet$。

\medskip\noindent
设给定 $Y$ 上紧化 $j_i : X \to \overline{X}_i$ 之间的态射
$g : \overline{X}_1 \to \overline{X}_2$。则得到同构
$$
R\overline{f}_{2, *} \circ j_{2, !} =
R\overline{f}_{2, *} \circ Rg_* \circ j_{1, !} =
R\overline{f}_{1, *} \circ j_{1, !}
$$
；第一个等号使用引理 \ref{more-etale-lemma-shriek-lqf-and-proper}。

\medskip\noindent
为完成证明，因 $X$ 在 $Y$ 上的紧化范畴余滤过，只需证明：$X$ 在 $Y$ 上
各紧化态射的复合，被送到函子同构的复合\footnote{具体地，若
$\alpha, \beta : F \to G$ 是函子态射，而 $\gamma : G \to H$ 是函子同构，
且 $\gamma \circ \alpha = \gamma \circ \beta$，则可推出
$\alpha = \beta$。}。为此，设 $j_3 : X \to \overline{X}_3$ 是第三个紧化，
而 $h : \overline{X}_2 \to \overline{X}_3$ 是紧化态射。我们必须证明复合
$$
R\overline{f}_{3, *} \circ j_{3, !} =
R\overline{f}_{3, *} \circ Rh_* \circ j_{2, !} =
R\overline{f}_{2, *} \circ j_{2, !} =
R\overline{f}_{2, *} \circ Rg_* \circ j_{1, !} =
R\overline{f}_{1, *} \circ j_{1, !}
$$
等于仅使用 $j_3$、$g \circ h$ 与 $j_1$ 构造出的函子同构。计算表明，只需证明
引理 \ref{more-etale-lemma-shriek-lqf-and-proper} 的映射之复合
$$
j_{3, !} \to Rh_* \circ j_{2, !} \to Rh_* \circ Rg_* \circ j_{1, !}
$$
通过等同 $R(h \circ g)_* = Rh_* \circ Rg_*$，与相应映射
$j_{3, !} \to R(h \circ g)_* \circ j_{1, !}$ 一致。
由于引理 \ref{more-etale-lemma-shriek-lqf-and-proper} 的映射是引理
\ref{more-etale-lemma-shriek-proper-and-open} 的映射的特例（因为 $j_1$ 与 $j_2$ 分离），
这立即由引理 \ref{more-etale-lemma-shriek-proper-and-open-compose} 得出。
\end{proof}

\begin{lemma}
\label{more-etale-lemma-shriek-composition}
设 $f : X \to Y$ 与 $g : Y \to Z$ 是拟紧拟分离概形之间的有限型分离态射。
则有典范同构 $Rg_! \circ Rf_! \to R(g \circ f)_!$。
\end{lemma}

\begin{proof}
选取 $Y$ 在 $Z$ 上的紧化 $i : Y \to \overline{Y}$，再选取 $X$ 在
$\overline{Y}$ 上的紧化 $X \to \overline{X}$。这里两次使用“平坦性的进一步结果”
定理 \ref{flat-theorem-nagata} 和引理 \ref{flat-lemma-compactifyable}。
令 $U$ 为 $Y$ 在 $\overline{X}$ 中的逆像，于是得到交换图
$$
\xymatrix{
X \ar[r]_j \ar[d]_f &
U \ar[dl]^{f'} \ar[r]_{j'} &
\overline{X} \ar[dl]^{\overline{f}} \\
Y \ar[r]_i \ar[d]_g &
\overline{Y} \ar[dl]^{\overline{g}} \\
Z
}
$$
于是有
\begin{align*}
R(g \circ f)_!
& =
R(\overline{g} \circ \overline{f})_* \circ (j' \circ j)_! \\
& =
R\overline{g}_* \circ R\overline{f}_* \circ j'_! \circ j_! \\
& =
R\overline{g}_* \circ i_! \circ Rf'_* \circ j_! \\
& =
Rg_! \circ Rf_!
\end{align*}
第一个等号是 $R(g \circ f)_!$ 的定义。第二个等号使用等同
$R(\overline{g} \circ \overline{f})_* = R\overline{g}_* \circ R\overline{f}_*$
以及引理 \ref{more-etale-lemma-f-shriek-composition} 的等同
$(j' \circ j)_! = j'_! \circ j_!$。第三个等号所用的等同
$i_! \circ Rf'_* \to R\overline{f}_* \circ j_!$ 来自引理
\ref{more-etale-lemma-shriek-proper-and-open}。最后的第四个等号是 $Rg_!$ 与 $Rf_!$ 的定义。
为完成证明，我们要说明该同构与所作选取无关。

\medskip\noindent
设有两个图
$$
\vcenter{
\xymatrix{
X \ar[r]_{j_1} \ar[d] &
U_1 \ar[dl]^{f_1} \ar[r]_{j'_1} &
\overline{X}_1 \ar[dl]^{\overline{f}_1} \\
Y \ar[r]_{i_1} \ar[d] &
\overline{Y}_1 \ar[dl]^{\overline{g}_1} \\
Z
}
}
\quad\text{且}\quad
\vcenter{
\xymatrix{
X \ar[r]_{j_2} \ar[d] &
U_2 \ar[dl]^{f_2} \ar[r]_{j'_2} &
\overline{X}_2 \ar[dl]^{\overline{f}_2} \\
Y \ar[r]_{i_2} \ar[d] &
\overline{Y}_2 \ar[dl]^{\overline{g}_2} \\
Z
}
}
$$
首先可选取 $Y$ 在 $Z$ 上的紧化 $i : Y \to \overline{Y}$，使其同时支配
$\overline{Y}_1$ 与 $\overline{Y}_2$；见“平坦性的进一步结果”引理
\ref{flat-lemma-compactifications-cofiltered}。由“平坦性的进一步结果”引理
\ref{flat-lemma-right-multiplicative-system} 以及“范畴”引理
\ref{categories-lemma-morphisms-right-localization} 和
\ref{categories-lemma-equality-morphisms-right-localization}，可选取 $X$ 在
$\overline{Y}$ 上的紧化 $X \to \overline{X}$，并配有态射
$\overline{X} \to \overline{X}_1$ 与 $\overline{X} \to \overline{X}_2$，
使复合 $\overline{X} \to \overline{Y} \to \overline{Y}_1$ 等于复合
$\overline{X} \to \overline{X}_1 \to \overline{Y}_1$，且复合
$\overline{X} \to \overline{Y} \to \overline{Y}_2$ 等于复合
$\overline{X} \to \overline{X}_2 \to \overline{Y}_2$。因此，只需在有如下
交换图时，比较这两个图所确定的映射：
$$
\xymatrix{
X \ar[rr]_{j_1} \ar@{=}[d] & &
U_1 \ar[d]^{h'} \ar[ddll] \ar[rr]_{j'_1} & &
\overline{X}_1 \ar[d]^h \ar[ddll] \\
X \ar'[r][rr]^-{j_2} \ar[d] & &
U_2 \ar'[dl][ddll] \ar'[r][rr]^-{j'_2} & &
\overline{X}_2 \ar[ddll] \\
Y \ar[rr]^{i_1} \ar@{=}[d] & & \overline{Y}_1 \ar[d]^k \\
Y \ar[rr]^{i_2} \ar[d] & & \overline{Y}_2 \ar[dll] \\
Z
}
$$
下列每个方块
$$
\xymatrix{
X \ar[r]_{j_1} \ar[d]_{\text{id}} \ar@{}[dr]|A &
U_1 \ar[d]^{h'} \\
X \ar[r]^{j_2} &
U_2
}
\quad
\xymatrix{
U_2 \ar[r]_{j_2'} \ar[d]_{f_2} \ar@{}[dr]|B &
\overline{X}_2 \ar[d]^{\overline{f}_2} \\
Y \ar[r]^{i_2} &
\overline{Y}_2
}
\quad
\xymatrix{
U_1 \ar[r]_{j_1'} \ar[d]_{f_1} \ar@{}[dr]|C &
\overline{X}_1 \ar[d]^{\overline{f}_1} \\
Y \ar[r]^{i_1} &
\overline{Y}_1
}
\quad
\xymatrix{
Y \ar[r]_{i_1} \ar[d]_{\text{id}} \ar@{}[dr]|D &
\overline{Y}_1 \ar[d]^k \\
Y \ar[r]^{i_2} &
\overline{Y}_2
}
\quad
\xymatrix{
X \ar[r]_{j_1' \circ j_1} \ar[d]_{\text{id}} \ar@{}[dr]|E &
\overline{X}_1 \ar[d]^h \\
X \ar[r]^{j_2} &
\overline{X}_2
}
$$
均产生如下同构：
\begin{align*}
\gamma_A & :
j_{2, !} \to Rh'_* \circ j_{1, !}  \\
\gamma_B & :
i_{2, !} \circ Rf_{2, *} \to R\overline{f}_{2, *} \circ j'_{2, !} \\
\gamma_C & :
i_{1, !} \circ Rf_{1, *} \to R\overline{f}_{1, *} \circ j'_{1, !} \\
\gamma_D & :
i_{2, !} \to Rk_* \circ i_{1, !} \\
\gamma_E & :
j_{2, !} \to Rh_* \circ (j'_1 \circ j_1)_!
\end{align*}
；这是通过应用引理 \ref{more-etale-lemma-shriek-proper-and-open} 的映射得到的
（当左侧竖箭头为恒等时，它与引理 \ref{more-etale-lemma-shriek-lqf-and-proper}
的映射相同）。记
\begin{align*}
F_1 & = Rf_{1, *} \circ j_{1, !} \\
F_2 & = Rf_{2, *} \circ j_{2, !} \\
G_1 & = R\overline{g}_{1, *} \circ i_{1, !} \\
G_2 & = R\overline{g}_{2, *} \circ i_{2, !} \\
C_1 & = R(\overline{g}_1 \circ \overline{f}_1)_* \circ (j'_1 \circ j_1)_! \\
C_2 & = R(\overline{g}_2 \circ \overline{f}_2)_* \circ (j'_2 \circ j_2)_!
\end{align*}
本证明第一段以及引理 \ref{more-etale-lemma-shriek-well-defined} 中给出的构造使用
\begin{enumerate}
\item $\gamma_C$ 作为映射 $G_1 \circ F_1 \to C_1$；
\item $\gamma_B$ 作为映射 $G_2 \circ F_2 \to C_2 $；
\item $\gamma_A$ 作为映射 $F_2 \to F_1$；
\item $\gamma_D$ 作为映射 $G_2 \to G_1$；以及
\item $\gamma_E$ 作为映射 $C_2 \to C_1$。
\end{enumerate}
因此必须证明图
$$
\xymatrix{
C_2 \ar[rr]_{\gamma_E} & &
C_1 \\
G_2 \circ F_2 \ar[rr]^{\gamma_D \circ \gamma_A} \ar[u]^{\gamma_B} & &
G_1 \circ F_1 \ar[u]_{\gamma_C}
}
$$
交换。我们将使用引理 \ref{more-etale-lemma-shriek-proper-and-open-compose} 与
\ref{more-etale-lemma-shriek-proper-and-open-compose-horizontal}，并沿用评注
\ref{more-etale-remark-going-around} 中（略有滥用的）记号；特别地，从记号中略去与恒等变换的
$\star$ 乘积。可写成 $\gamma_E = \gamma_F \circ \gamma_A$，其中
$$
\xymatrix{
U_1 \ar[r]_{j'_1} \ar[d]_{h'} \ar@{}[rd]|F &
\overline{X}_1 \ar[d]^h \\
U_2 \ar[r]^{j'_2} &
\overline{X}_2
}
$$
于是
$$
\gamma_E \circ \gamma_B = \gamma_F \circ \gamma_A  \circ \gamma_B
= \gamma_F \circ \gamma_B \circ \gamma_A
$$
；最后一个等号成立，是因为方块 $A$ 与 $B$ 仅在一点相交（类似于评注
\ref{more-etale-remark-going-around} 末尾的论证）。因此只需证明
$\gamma_C \circ \gamma_D = \gamma_F \circ \gamma_B$。两者都等于方块
$$
\xymatrix{
U_1 \ar[r] \ar[d] & \overline{X}_1 \ar[d] \\
Y \ar[r] & \overline{Y}_2
}
$$
的映射，故结论成立。
\end{proof}

\begin{lemma}
\label{more-etale-lemma-pseudo-functor}
设 $f : X \to Y$、$g : Y \to Z$、$h : Z \to T$ 是拟紧拟分离概形之间的
有限型分离态射。则引理 \ref{more-etale-lemma-shriek-composition} 的各同构所成的图
$$
\xymatrix{
Rh_! \circ Rg_! \circ Rf_! \ar[r]_{\gamma_C} \ar[d]^{\gamma_A} &
R(h \circ g)_! \circ Rf_! \ar[d]_{\gamma_{A + B}} \\
Rh_! \circ R(g \circ f)_! \ar[r]^{\gamma_{B + C}} &
R(h \circ g \circ f)_!
}
$$
交换（各 $\gamma$ 的含义见证明）。
\end{lemma}

\begin{proof}
为此，先选取 $Z$ 在 $T$ 上的紧化 $\overline{Z}$，再选取 $Y$ 在
$\overline{Z}$ 上的紧化 $\overline{Y}$，继而选取 $X$ 在 $\overline{Y}$ 上的
紧化 $\overline{X}$。这里使用“平坦性的进一步结果”定理
\ref{flat-theorem-nagata} 和引理 \ref{flat-lemma-compactifyable}。
令 $W \subset \overline{Y}$ 为 $Z$ 在 $\overline{Y} \to \overline{Z}$ 下的逆像，
并令 $U \subset V \subset \overline{X}$ 为 $Y \subset W$ 在
$\overline{X} \to \overline{Y}$ 下的逆像。由此得到图
$$
\xymatrix{
X \ar[d]_f \ar[r] & U \ar[r] \ar[d] \ar@{}[dr]|A &
V \ar[d] \ar[r] \ar@{}[rd]|B & \overline{X} \ar[d] \\
Y \ar[d]_g \ar[r] & Y \ar[r] \ar[d] & W \ar[r] \ar[d] \ar@{}[rd]|C &
\overline{Y} \ar[d] \\
Z \ar[d]_h \ar[r] & Z \ar[d] \ar[r] & Z \ar[d] \ar[r] & \overline{Z} \ar[d] \\
T \ar[r] & T \ar[r] & T \ar[r] & T
}
$$
不引入大量记号而完全仿照引理 \ref{more-etale-lemma-shriek-composition} 的证明论证，
可见第一个显示图中的各映射，按引理陈述的图所示，分别使用引理
\ref{more-etale-lemma-shriek-proper-and-open} 对矩形 $A + B$、$B + C$、$A$ 与 $C$
给出的映射。由引理 \ref{more-etale-lemma-shriek-proper-and-open-compose} 和
\ref{more-etale-lemma-shriek-proper-and-open-compose-horizontal}，有
$\gamma_{A + B} = \gamma_B \circ \gamma_A$ 以及
$\gamma_{B + C} = \gamma_B \circ \gamma_C$。因此，只要
$\gamma_A \circ \gamma_C = \gamma_C \circ \gamma_A$，所需等式即成立。
后一个等式确实成立，因为方块 $A$ 与 $C$ 仅在一点相交（类似于评注
\ref{more-etale-remark-going-around} 末尾的论证）。
\end{proof}

\begin{lemma}
\label{more-etale-lemma-base-change-shriek}
考虑笛卡尔方块
$$
\xymatrix{
X' \ar[r]_{g'} \ar[d]_{f'} & X \ar[d]^f \\
Y' \ar[r]^g & Y
}
$$
，其中各概形拟紧且拟分离，而 $f$ 分离且有限型。则有典范同构
$$
g^{-1} \circ Rf_! \to Rf'_! \circ (g')^{-1}
$$
此外，这些同构与引理 \ref{more-etale-lemma-shriek-composition} 的同构相容。
\end{lemma}

\begin{proof}
选取在 $Y$ 上的紧化 $j : X \to \overline{X}$，并以
$\overline{f} : \overline{X} \to Y$ 表示结构态射。以
$j' : X' \to \overline{X}'$ 与 $\overline{f}' : \overline{X}' \to Y'$
表示 $j$ 与 $\overline{f}$ 的基变换。由于
$Rf_! = R\overline{f}_* \circ j_!$ 且
$Rf'_! = R\overline{f}'_* \circ j'_!$，可通过下式构造该同构：
$$
g^{-1} \circ R\overline{f}_* \circ j_! \to
R\overline{f}'_* \circ (\overline{g}')^{-1} \circ j_! \to
R\overline{f}'_* \circ j'_! \circ (g')^{-1} 
$$
其中第一个箭头是固有基变换定理给出的同构（在下有界挠情形，使用
“\'Etale 上同调”引理 \ref{etale-cohomology-lemma-proper-base-change}；
在 $\Lambda$ 为挠环的情形，使用“\'Etale 上同调”引理
\ref{etale-cohomology-lemma-proper-base-change-mod-n}），第二个箭头是引理
\ref{more-etale-lemma-base-change-f-shriek-separated} 的同构。

\medskip\noindent
为完成证明，必须说明两点：第一，如此得到的函子同构与紧化的选取无关；
第二，若把引理中的两个基变换图竖向叠合，则这些基变换同构与引理
\ref{more-etale-lemma-shriek-composition} 的同构相容。一个略去的直接论证表明，
若能证明以下各同构与基变换相容，则两点均成立：
\begin{enumerate}
\item 当 $f : X \to Y$ 与 $g : Y \to Z$ 固有时，
$Rg_* \circ Rf_* = R(g \circ f)_*$；
\item 当 $f : X \to Y$ 与 $g : Y \to Z$ 分离且拟有限时，
$g_! \circ f_! = (g \circ f)_!$；以及
\item 当 $f : X \to Y$ 与 $f' : X' \to Y'$ 固有，而
$g : Y' \to Y$ 与 $g' : X' \to X$ 分离且拟有限，并满足
$f \circ g' = g \circ f'$ 时，$g_! \circ Rf'_* = Rf_* \circ g'_!$。
\end{enumerate}
对 (1)，这由“位点上同调”评注
\ref{sites-cohomology-remark-compose-base-change} 成立；对 (2)，由评注
\ref{more-etale-remark-f-shriek-base-change-composition} 成立；对 (3)，由引理
\ref{more-etale-lemma-shriek-proper-and-open-base-change} 成立。
\end{proof}

\begin{remark}
\label{more-etale-remark-shriek-to-star}
设 $f : X \to Y$ 是概形的有限型分离态射，其中 $Y$ 拟紧且拟分离。
下面将构造映射
$$
Rf_!K \longrightarrow Rf_*K
$$
；它关于 $D^+_{tors}(X_\etale, \Lambda)$ 中的 $K$ 函子性地成立；若
$\Lambda$ 为挠环，则同样关于 $D(X_\etale, \Lambda)$ 中的对象函子性地成立。
两种情形下，该函子变换均与下列同构相容：
\begin{enumerate}
\item 引理 \ref{more-etale-lemma-shriek-composition} 的同构
$Rg_! \circ Rf_! \to R(g \circ f)_!$，以及“位点上同调”引理
\ref{sites-cohomology-lemma-derived-pushforward-composition} 的同构
$Rg_* \circ Rf_* \to R(g \circ f)_*$；以及
\item 引理 \ref{more-etale-lemma-base-change-shriek} 的同构
$g^{-1} \circ Rf_! \to Rf'_! \circ (g')^{-1}$，以及“位点上同调”评注
\ref{sites-cohomology-remark-base-change} 的基变换映射。
\end{enumerate}
具体地，选取在 $Y$ 上的紧化 $j : X \to \overline{X}$，并以
$\overline{f} : \overline{X} \to Y$ 表示结构态射。由于
$Rf_! = R\overline{f}_* \circ j_!$ 且
$Rf_* = R\overline{f}_* \circ Rj_*$，只需构造函子变换 $j_! \to Rj_*$。
为此，使用“\'Etale 上同调”引理
\ref{etale-cohomology-lemma-shriek-into-star-separated-etale} 的典范变换
$j_! \to j_*$。我们略去所得变换与紧化选取无关的证明，也略去相容性
(1) 与 (2) 的证明。
\end{remark}













\section{导出下叹号的性质}
\label{more-etale-section-derived-lower-shriek-properties}

\noindent
下面列出导出下叹号的一些性质。

\begin{lemma}
\label{more-etale-lemma-derived-lower-shriek-commute-direct-sums}
设 $f : X \to Y$ 是拟紧拟分离概形之间的有限型分离态射。设 $\Lambda$ 是环。
\begin{enumerate}
\item 设 $K_i \in D^+_{tors}(X_\etale, \Lambda)$，$i \in I$，是一族对象。
假设给定 $a \in \mathbf{Z}$，使得当 $n < a$ 且 $i \in I$ 时
$H^n(K_i) = 0$。则 $Rf_!(\bigoplus_i K_i) = \bigoplus_i Rf_!K_i$。
\item 若 $\Lambda$ 为挠环，则函子
$Rf_! : D(X_\etale, \Lambda) \to D(Y_\etale, \Lambda)$
与直和交换。
\end{enumerate}
\end{lemma}

\begin{proof}
由构造，只需在 $f$ 为开浸入以及 $f$ 为固有态射时证明。对概形的任意开浸入
$j : U \to X$，函子 $j_! : D(U_\etale) \to D(X_\etale)$ 是拉回
$j^{-1} : D(X_\etale) \to D(U_\etale)$ 的左伴随，因而与直和交换；
见“位点上同调”引理 \ref{sites-cohomology-lemma-adjoint-lower-shriek-restrict}。
在固有情形有 $Rf_! = Rf_*$；在下有界情形，结论由“\'Etale 上同调”引理
\ref{etale-cohomology-lemma-direct-sum-bounded-below-pushforward} 得出；
当系数环 $\Lambda$ 为挠环时，结论由“\'Etale 上同调”引理
\ref{etale-cohomology-lemma-proper-mod-n-direct-sums} 得出。
\end{proof}

\begin{lemma}
\label{more-etale-lemma-derived-lower-shriek-bounded}
设 $f : X \to Y$ 是拟紧拟分离概形之间的有限型分离态射，$\Lambda$ 是环。
第 \ref{more-etale-section-derived-lower-shriek-compactification} 节中构造的函子 $Rf_!$
在如下意义下有界：存在整数 $N$，使得对
$E \in D^+_{tors}(X_\etale, \Lambda)$，或者在 $\Lambda$ 为挠环时对
$E \in D(X_\etale, \Lambda)$，有
\begin{enumerate}
\item 当 $i \leq a$ 时，$H^i(Rf_!(\tau_{\leq a}E) \to H^i(Rf_!(E))$ 是同构；
\item 当 $i \geq b$ 时，$H^i(Rf_!(E)) \to H^i(Rf_!(\tau_{\geq b - N}E))$ 是同构；
\item 若对某个 $-\infty \leq a \leq b \leq \infty$，当
$i \not \in [a, b]$ 时 $H^i(E) = 0$，则当 $i \not \in [a, b + N]$ 时
$H^i(Rf_!(E)) = 0$。
\end{enumerate}
\end{lemma}

\begin{proof}
设 $\Lambda$ 为挠环，并考虑函子
$Rf_! : D(X_\etale, \Lambda) \to D(Y_\etale, \Lambda)$。
由构造，只需在 $f$ 为开浸入以及 $f$ 为固有态射时证明。对概形的任意开浸入
$j : U \to X$，函子 $j_! : D(U_\etale) \to D(X_\etale)$ 正合，
故在此情形陈述以 $N = 0$ 成立。若 $f$ 固有，则 $Rf_! = Rf_*$，
即它是右导出函子。因此左侧的界由“导出范畴”引理
\ref{derived-lemma-negative-vanishing} 得出。此外，在此情形，
由“态射”引理 \ref{morphisms-lemma-morphism-finite-type-bounded-dimension}
和“\'Etale 上同调”引理
\ref{etale-cohomology-lemma-cohomological-dimension-proper}，
$f_* : \textit{Mod}(X_\etale, \Lambda) \to
\textit{Mod}(Y_\etale, \Lambda)$ 的上同调维数有界。因此由“导出范畴”引理
\ref{derived-lemma-unbounded-right-derived} 得到结论。

\medskip\noindent
其次，设 $\Lambda$ 任意，并考虑函子
$Rf_! : D^+_{tors}(X_\etale, \Lambda) \to D^+_{tors}(Y_\etale, \Lambda)$。
仍然立即化归到 $f$ 固有且 $Rf_! = Rf_*$ 的情形。(1) 仍是立即的。
为证明 (3)，可对 $b - a$ 作归纳，并使用截断的特异三角以及“\'Etale 上同调”引理
\ref{etale-cohomology-lemma-cohomological-dimension-proper}。
(2) 由 (3) 得出。略去细节。
\end{proof}

\begin{lemma}
\label{more-etale-lemma-Rf-shriek-for-quasi-finite}
设 $f : X \to Y$ 是拟紧拟分离概形之间的拟有限分离态射。则第
\ref{more-etale-section-derived-lower-shriek-compactification} 节中构造的函子 $Rf_!$
与第 \ref{more-etale-section-derived-duality-locally-quasi-finite} 节中构造的函子
$f_! : D(X_\etale, \Lambda) \to D(Y_\etale, \Lambda)$ 在二者共同定义域上的
限制一致。
\end{lemma}

\begin{proof}
由 Zariski 主定理（“态射的进一步结果”引理
\ref{more-morphisms-lemma-quasi-finite-separated-pass-through-finite}），
可找到开浸入 $j : X \to \overline{X}$ 与有限态射
$\overline{f} : \overline{X} \to Y$，使 $f = \overline{f} \circ j$。
由构造有 $Rf_! = R\overline{f}_* \circ j_!$。由于 $\overline{f}$ 有限，
由“\'Etale 上同调”命题
\ref{etale-cohomology-proposition-finite-higher-direct-image-zero}，有
$R\overline{f}_* = \overline{f}_*$。又例如由引理
\ref{more-etale-lemma-compactify-f-shriek-separated}，有
$\overline{f}_* \circ j_! = f_!$，故引理成立。
\end{proof}

\begin{lemma}
\label{more-etale-lemma-relative-mayer-vietoris}
设 $f : X \to Y$ 是拟紧拟分离概形之间的有限型分离态射。设 $U$ 与 $V$
是 $X$ 的拟紧开集，且 $X = U \cup V$。以 $a : U \to Y$、
$b : V \to Y$ 与 $c : U \cap V \to Y$ 表示 $f$ 的限制。设 $\Lambda$ 是环。
对 $D^+_{tors}(X_\etale, \Lambda)$ 中的 $K$，或者当 $\Lambda$ 为挠环时，
对 $K \in D(X_\etale, \Lambda)$，有特异三角
$$
Rc_!(K|_{U \cap V}) \to
Ra_!(K|_U) \oplus Rb_!(K|_V) \to
Rf_!K \to 
Rc_!(K|_{U \cap V})[1]
$$
成立于 $D(Y_\etale, \Lambda)$ 中。
\end{lemma}

\begin{proof}
这由引理 \ref{more-etale-lemma-mayer-vietoris-shriek}、引理
\ref{more-etale-lemma-shriek-composition} 给出的 $Rf_! \circ Rj_{U!} = Ra_!$，以及引理
\ref{more-etale-lemma-Rf-shriek-for-quasi-finite} 给出的 $Rj_{U!} = j_{U!}$ 得出。
\end{proof}

\begin{lemma}
\label{more-etale-lemma-relative-triangle-associated-to-open}
设 $f : X \to Y$ 是拟紧拟分离概形之间的有限型分离态射。设 $U$ 是 $X$
的拟紧开集，其补集为 $Z \subset X$。以 $g : U \to Y$ 和 $h : Z \to Y$
表示 $f$ 的限制。设 $\Lambda$ 是环。对 $D^+_{tors}(X_\etale, \Lambda)$
中的 $K$，或者当 $\Lambda$ 为挠环时，对 $K \in D(X_\etale, \Lambda)$，
有特异三角
$$
Rg_!(K|_U) \to
Rf_!K \to
Rh_!(K|_Z) \to
Rg_!(K|_U)[1]
$$
成立于 $D(Y_\etale, \Lambda)$ 中。
\end{lemma}

\begin{proof}
这由引理 \ref{more-etale-lemma-triangle-associated-to-open}、引理
\ref{more-etale-lemma-shriek-composition} 给出的 $Rf_! \circ Rj_! = Rg_!$ 与
$Rf_! \circ Ri_!$，以及引理 \ref{more-etale-lemma-Rf-shriek-for-quasi-finite} 给出的
$Rj_! = j_!$ 和 $Ri_! = i_! = i_*$ 得出。
\end{proof}

\begin{lemma}
\label{more-etale-lemma-shriek-and-thickening}
设 $f' : X' \to Y$ 是拟紧拟分离概形之间的有限型分离态射。设
$i : X \to X'$ 是增厚，并记 $f = f' \circ i$。设 $\Lambda$ 是环。
对 $D^+_{tors}(X'_\etale, \Lambda)$ 中的 $K'$，或者当 $\Lambda$ 为挠环时，
对 $K' \in D(X'_\etale, \Lambda)$，有 $Rf_!i^{-1}K' = Rf'_!K'$。
\end{lemma}

\begin{proof}
这成立是因为，由小 \'etale 拓扑斯的拓扑不变性（“\'Etale 上同调”定理
\ref{etale-cohomology-theorem-topological-invariance}），$i^{-1}$ 与
$i_* = i_!$ 是范畴的互逆等价，并且可以应用引理
\ref{more-etale-lemma-shriek-composition}。
\end{proof}

\begin{lemma}
\label{more-etale-lemma-projection-formula-Rf-lower-shriek}
设 $f : X \to Y$ 是拟紧拟分离概形之间的有限型分离态射。设 $\Lambda$ 是挠环，
$E \in D(X_\etale, \Lambda)$ 且 $K \in D(Y_\etale, \Lambda)$。则
$$
Rf_!E \otimes_\Lambda^\mathbf{L} K =
Rf_!(E \otimes_\Lambda^\mathbf{L} f^{-1}K)
$$
成立于 $D(Y_\etale, \Lambda)$ 中。
\end{lemma}

\begin{proof}
按 $Rf_!$ 的构造选取 $j : X \to \overline{X}$ 与
$\overline{f} : \overline{X} \to Y$。由“位点上同调”引理
\ref{sites-cohomology-lemma-j-shriek-and-tensor}，有
$j_!E \otimes_\Lambda^\mathbf{L} \overline{f}^{-1}K =
j_!(E \otimes_\Lambda^\mathbf{L} f^{-1}K)$。再应用“\'Etale 上同调”引理
\ref{etale-cohomology-lemma-projection-formula-proper-mod-n}，并使用
$f^{-1} = j^{-1} \circ \overline{f}^{-1}$ 与
$Rf_! = R\overline{f}_*j_!$，即得结论。
\end{proof}

\begin{remark}
\label{more-etale-remark-Rf-lower-shriek-change-of-rings}
设 $\Lambda_1 \to \Lambda_2$ 是挠环同态，$f : X \to Y$ 是拟紧拟分离概形之间的
有限型分离态射。图
$$
\xymatrix{
D(X_\etale, \Lambda_2) \ar[r]_{res} \ar[d]_{Rf_!} &
D(X_\etale, \Lambda_1) \ar[d]^{Rf_!} \\
D(Y_\etale, \Lambda_2) \ar[r]^{res} &
D(Y_\etale, \Lambda_1)
}
$$
交换；其中 $res$ 是“限制”函子，它通过给定环映射把 $\Lambda_2$-模视为
$\Lambda_1$-模。对第 \ref{more-etale-section-derived-lower-shriek-compactification} 节中
那样的分解 $f = \overline{f} \circ j$，写成
$Rf_! = R\overline{f}_* \circ j_!$。直接检验可见结论对 $j_!$ 成立，
而由“位点上同调”引理
\ref{sites-cohomology-lemma-modules-abelian-unbounded}，结论对
$R\overline{f}_*$ 成立。另一方面，图
$$
\xymatrix{
D(X_\etale, \Lambda_1)
\ar[r]_{- \otimes_{\Lambda_1}^\mathbf{L} \Lambda_2} \ar[d]_{Rf_!} &
D(X_\etale, \Lambda_2) \ar[d]^{Rf_!} \\
D(Y_\etale, \Lambda_1)
\ar[r]^{- \otimes_{\Lambda_1}^\mathbf{L} \Lambda_2} &
D(Y_\etale, \Lambda_2)
}
$$
也交换；这由引理 \ref{more-etale-lemma-projection-formula-Rf-lower-shriek} 得出。
\end{remark}

\begin{remark}
\label{more-etale-remark-projection-formula-for-internal-hom}
设 $f : X \to Y$ 是拟紧拟分离概形之间的有限型分离态射，$\Lambda$ 是挠系数环，
而 $K$ 与 $L$ 是 $D(X_\etale, \Lambda)$ 的对象。我们断言有典范映射
$$
\alpha :
Rf_*R\SheafHom_\Lambda(K, L)
\longrightarrow
R\SheafHom_\Lambda(Rf_!K, Rf_!L)
$$
关于 $K$ 与 $L$ 函子性地成立。具体地，按 $Rf_!$ 的构造选取
$j : X \to \overline{X}$ 与 $\overline{f} : \overline{X} \to Y$。
先定义映射
$$
\beta :
Rj_*R\SheafHom_\Lambda(K, L)
\longrightarrow
R\SheafHom_\Lambda(j_!K, j_!L)
$$
由导出范畴中内部 Hom 的构造，这等价于定义映射
$$
\beta' :
Rj_*R\SheafHom_\Lambda(K, L) \otimes_\Lambda^\mathbf{L} j_!K
\longrightarrow
j_!L
$$
见“位点上同调”第 \ref{sites-cohomology-section-internal-hom} 节。
$\beta'$ 的源等于
$$
j_!\left(R\SheafHom_\Lambda(K, L) \otimes_\Lambda^\mathbf{L} K\right)
$$
；这由“位点上同调”引理
\ref{sites-cohomology-lemma-j-shriek-and-tensor} 得出。因此可令
$\beta' = j_!\beta''$，其中
$\beta'' : R\SheafHom_\Lambda(K, L) \otimes_\Lambda^\mathbf{L} K \to L$
通过上述内部 Hom 的泛性质对应于 $R\SheafHom_\Lambda(K, L)$ 上的恒等态射。
由“位点上同调”评注
\ref{sites-cohomology-remark-projection-formula-for-internal-hom}，有典范映射
$$
\gamma :
R\overline{f}_*R\SheafHom_\Lambda(j_!K, j_!L)
\longrightarrow
R\SheafHom_\Lambda(R\overline{f}_*j_!K, R\overline{f}_*j_!L)
$$
由于 $Rf_! = R\overline{f}_*j_!$，并且（由 Leray）
$Rf_* = R\overline{f}_* Rj_*$，得到所需映射
$\alpha = \gamma \circ R\overline{f}_*\beta$。
\end{remark}









\section{导出上叹号}
\label{more-etale-section-derived-upper-shriek}

\noindent
我们用一个 Brown 可表示性定理得到 $Rf^!$。

\begin{lemma}
\label{more-etale-lemma-upper-shriek-derived}
设 $f : X \to Y$ 是拟紧拟分离概形之间的有限型分离态射，$\Lambda$ 是挠系数环。
函子
$Rf_! : D(X_\etale, \Lambda) \to D(Y_\etale, \Lambda)$
有右伴随
$Rf^! : D(Y_\etale, \Lambda) \to D(X_\etale, \Lambda)$.
\end{lemma}

\begin{proof}
这由“内射对象”命题 \ref{injectives-proposition-brown} 以及上面的引理
\ref{more-etale-lemma-derived-lower-shriek-commute-direct-sums} 得出。
\end{proof}

\begin{lemma}
\label{more-etale-lemma-Rf-upper-shriek-for-quasi-finite}
设 $f : X \to Y$ 是拟紧拟分离概形之间的拟有限分离态射，$\Lambda$ 是挠系数环。
引理 \ref{more-etale-lemma-upper-shriek-derived} 的函子
$Rf^! : D(Y_\etale, \Lambda) \to D(X_\etale, \Lambda)$ 与引理
\ref{more-etale-lemma-lqf-shriek-derived} 的函子 $Rf^!$ 相同。
\end{lemma}

\begin{proof}
由引理 \ref{more-etale-lemma-Rf-shriek-for-quasi-finite}，有 $Rf_! = f_!$；
结论遂由伴随的唯一性得出。
\end{proof}

\begin{lemma}
\label{more-etale-lemma-Rf-upper-shriek-for-etale}
设 $j : U \to X$ 是拟紧拟分离概形之间的分离 \'etale 态射，$\Lambda$ 是挠系数环。
函子 $Rj^! : D(X_\etale, \Lambda) \to D(U_\etale, \Lambda)$ 等于
$j^{-1}$。
\end{lemma}

\begin{proof}
这是因为 $Rj^!$ 与 $j^{-1}$ 都是 $Rj_! = j_!$ 的右伴随。例如见引理
\ref{more-etale-lemma-Rf-upper-shriek-for-quasi-finite} 与
\ref{more-etale-lemma-etale-upper-shriek}。
\end{proof}

\begin{lemma}
\label{more-etale-lemma-twisted-inverse-image-bounded-below}
设 $f : X \to Y$ 是拟紧拟分离概形之间的有限型分离态射，$\Lambda$ 是挠环。
函子 $Rf^!$ 把 $D^+(Y_\etale, \Lambda)$ 送入
$D^+(X_\etale, \Lambda)$。更准确地，存在整数 $N \geq 0$，使得若
$K \in D(Y_\etale, \Lambda)$ 在 $i < a$ 时满足 $H^i(K) = 0$，则在
$i < a - N$ 时 $H^i(Rf^!K) = 0$。
\end{lemma}

\begin{proof}
令 $N$ 为引理 \ref{more-etale-lemma-derived-lower-shriek-bounded} 中得到的整数。
由构造，对 $K \in D(Y_\etale, \Lambda)$ 与
$L \in \in D(X_\etale, \Lambda)$，有
$\Hom_X(L, Rf^!K) = \Hom_Y(Rf_!L, K)$。设当 $i < a$ 时
$H^i(K) = 0$。取 $L = \tau_{\leq a - N - 1}Rf^!K$。由引理
\ref{more-etale-lemma-derived-lower-shriek-bounded}，复形 $Rf_!L$ 的上同调层在次数
$\leq a - 1$ 消没。故由“导出范畴”引理
\ref{derived-lemma-negative-exts}，有 $\Hom_Y(Rf_!L, K) = 0$。
因此典范映射 $\tau_{\leq a - N - 1}Rf^!K \to Rf^!K$ 为零，
这蕴含当 $i \leq a - N - 1$ 时 $H^i(Rf^!K) = 0$。
\end{proof}

\noindent
设 $f : X \to Y$ 是拟分离拟紧概形之间的有限型分离态射，$\Lambda$ 是挠系数环。
对每个 $K \in D(Y_\etale, \Lambda)$ 与 $L \in D(X_\etale, \Lambda)$，
得到典范映射
\begin{equation}
\label{more-etale-equation-sheafy-trace}
Rf_*R\SheafHom_\Lambda(L, Rf^!K) \longrightarrow R\SheafHom_\Lambda(Rf_!L, K)
\end{equation}
具体地，该映射构造为复合
$$
Rf_*R\SheafHom_\Lambda(L, Rf^!K) \to
R\SheafHom_\Lambda(Rf_!L, Rf_!Rf^!K) \to
R\SheafHom_\Lambda(Rf_!L, K)
$$
，其中第一个箭头来自评注 \ref{more-etale-remark-projection-formula-for-internal-hom}，
第二个箭头是该伴随的余单位 $Rf_!Rf^!K \to K$。

\begin{lemma}
\label{more-etale-lemma-iso-on-RSheafHom}
设 $f : X \to Y$ 是拟紧拟分离概形之间的有限型分离态射，$\Lambda$ 是挠环。
对每个 $K \in D(Y_\etale, \Lambda)$ 与 $L \in D(X_\etale, \Lambda)$，
映射 (\ref{more-etale-equation-sheafy-trace})
$$
Rf_*R\SheafHom_\Lambda(L, Rf^!K) \longrightarrow R\SheafHom_\Lambda(Rf_!L, K)
$$
是同构。
\end{lemma}

\begin{proof}
为证明该引理，必须说明：对任意 $M \in D(Y_\etale, \Lambda)$，映射
(\ref{more-etale-equation-sheafy-trace}) 诱导双射
$$
\Hom_Y(M, Rf_*R\SheafHom_\Lambda(L, Rf^!K))
\longrightarrow
\Hom_Y(M, R\SheafHom_\Lambda(Rf_!L, K))
$$
为此使用如下等式链：
\begin{align*}
\Hom_Y(M, Rf_*R\SheafHom_\Lambda(L, Rf^!K))
& =
\Hom_X(f^{-1}M, R\SheafHom_\Lambda(L, Rf^!K)) \\
& =
\Hom_X(f^{-1}M \otimes_\Lambda^\mathbf{L} L, Rf^!K) \\
& =
\Hom_Y(Rf_!(f^{-1}M \otimes_\Lambda^\mathbf{L} L), K) \\
& =
\Hom_Y(M \otimes_\Lambda^\mathbf{L} Rf_!L, K) \\
& =
\Hom_Y(M, R\SheafHom_\Lambda(Rf_!L, K))
\end{align*}
第一个等号由“位点上同调”引理 \ref{sites-cohomology-lemma-adjoint} 成立。
第二个等号由“位点上同调”引理
\ref{sites-cohomology-lemma-internal-hom} 成立。第三个等号由 $Rf^!$ 的构造成立。
第四个等号由引理 \ref{more-etale-lemma-projection-formula-Rf-lower-shriek} 成立
（这是关键一步）。第五个等号由“位点上同调”引理
\ref{sites-cohomology-lemma-internal-hom} 成立。
\end{proof}

\begin{lemma}
\label{more-etale-lemma-iso-global-hom}
设 $f : X \to Y$ 是拟分离拟紧概形之间的有限型分离态射，$\Lambda$ 是挠环。
对每个 $K \in D(Y_\etale, \Lambda)$ 与 $L \in D(X_\etale, \Lambda)$，
映射 (\ref{more-etale-equation-sheafy-trace}) 诱导同构
$$
R\Hom_X(L, Rf^!K) \longrightarrow R\Hom_Y(Rf_!L, K)
$$
，其中两端为全局导出 Hom。
\end{lemma}

\begin{proof}
由“位点上同调”第 \ref{sites-cohomology-section-global-RHom} 节中的构造，有
$$
R\Hom_X(L, Rf^!K) =
R\Gamma(X, R\SheafHom_\Lambda(L, Rf^!K)) =
R\Gamma(Y, Rf_*R\SheafHom_\Lambda(L, Rf^!K))
$$
（第二个等号由 Leray 成立），以及
$$
R\Hom_Y(Rf_!L, K) = R\Gamma(Y, R\SheafHom_\Lambda(Rf_!L, K))
$$
因此该引理是引理 \ref{more-etale-lemma-iso-on-RSheafHom} 的推论。
\end{proof}

\begin{lemma}
\label{more-etale-lemma-cartesian-square-Rf-upper-shriek}
考虑笛卡尔方块
$$
\xymatrix{
X' \ar[r]_{g'} \ar[d]_{f'} & X \ar[d]^f \\
Y' \ar[r]^g & Y
}
$$
，其中各概形拟紧且拟分离，而 $f$ 分离且有限型。则有
$Rf^! \circ Rg_* = Rg'_* \circ R(f')^!$。
\end{lemma}

\begin{proof}
由伴随函子的唯一性，这从导出下叹号的基变换得出：由引理
\ref{more-etale-lemma-base-change-shriek}，有
$g^{-1} \circ Rf_! = Rf'_! \circ (g')^{-1}$。
\end{proof}

\begin{remark}
\label{more-etale-remark-Rf-upper-shriek-change-of-rings}
设 $\Lambda_1 \to \Lambda_2$ 是挠环同态，$f : X \to Y$ 是拟紧拟分离概形之间的
有限型分离态射。图
$$
\xymatrix{
D(X_\etale, \Lambda_2) \ar[r]_{res} &
D(X_\etale, \Lambda_1) \\
D(Y_\etale, \Lambda_2) \ar[r]^{res} \ar[u]^{Rf^!} &
D(Y_\etale, \Lambda_1) \ar[u]_{Rf^!}
}
$$
交换；其中 $res$ 是“限制”函子，它通过给定环映射把 $\Lambda_2$-模视为
$\Lambda_1$-模。这由伴随的唯一性、评注
\ref{more-etale-remark-Rf-lower-shriek-change-of-rings} 中第二个交换图，以及等式
$$
\Hom_{\Lambda_2}(K_1 \otimes_{\Lambda_1}^\mathbf{L} \Lambda_2, K_2) =
\Hom_{\Lambda_1}(K_1, res(K_2))
$$
成立这一事实得出。无论对象位于 $X_\etale$ 上还是 $Y_\etale$ 上，该等式都是
“位点上同调”引理 \ref{sites-cohomology-lemma-adjoint} 的一个非常特殊的情形。
\end{remark}











\section{紧支撑上同调}
\label{more-etale-section-compactly-supported-cohomology}

\noindent
设 $k$ 是域，$\Lambda$ 是环，$X$ 是 $k$ 上的有限型分离概形，其结构态射为
$f : X \to \Spec(k)$。在第
\ref{more-etale-section-derived-lower-shriek-compactification} 节中，我们定义了函子
$Rf_! : D^+_{tors}(X_\etale, \Lambda) \to D^+_{tors}(\Spec(k), \Lambda)$
；若 $\Lambda$ 是挠环，还定义了函子
$Rf_! : D(X_\etale, \Lambda) \to D(\Spec(k), \Lambda)$。
再与 $\Spec(k)$ 上的全局截面函子复合，便得到我们所谓的紧支撑上同调。

\begin{definition}
\label{more-etale-definition-cohomology-compact-support}
设 $X$ 是域 $k$ 上的有限型分离概形，$\Lambda$ 是环。设 $K$ 是
$D^+_{tors}(X_\etale, \Lambda)$ 的对象；若 $\Lambda$ 为挠环，也可设它是
$D(X_\etale, \Lambda)$ 的对象。$K$ 的 {\it 带紧支撑上同调}，亦称 $K$ 的
{\it 紧支撑上同调}，是
$$
R\Gamma_c(X, K) = R\Gamma(\Spec(k), Rf_!K)
$$
，其中 $f : X \to \Spec(k)$ 是结构态射。我们记
$H^i_c(X, K) = H^i(R\Gamma_c(X, K))$。
\end{definition}

\noindent
该定义与定义 \ref{more-etale-definition-compact-support} 不冲突；这将在引理
\ref{more-etale-lemma-compact-support-h0} 中检验。该定义的效用在于如下结果。

\begin{lemma}
\label{more-etale-lemma-stalk-R-f-shriek}
设 $f : X \to Y$ 是概形的有限型分离态射，其中 $Y$ 拟紧且拟分离。
设 $K$ 是 $D^+_{tors}(X_\etale, \Lambda)$ 的对象；若 $\Lambda$ 为挠环，
也可设它是 $D(X_\etale, \Lambda)$ 的对象。则对任意几何点
$\overline{y} : \Spec(k) \to Y$，在 $D(\Lambda)$ 中有典范同构
$$
(Rf_!K)_{\overline{y}}
\longrightarrow
R\Gamma_c(X_{\overline{y}}, K|_{X_{\overline{y}}})
$$
\end{lemma}

\begin{proof}
这是引理 \ref{more-etale-lemma-base-change-shriek} 与各定义的直接推论。
\end{proof}

\begin{lemma}
\label{more-etale-lemma-compact-support-h0}
设 $X$ 是域 $k$ 上的有限型分离概形。若 $\mathcal{F}$ 是挠阿贝尔层，
则定义 \ref{more-etale-definition-compact-support} 中定义的阿贝尔群
$H^0_c(X, \mathcal{F})$ 与定义
\ref{more-etale-definition-cohomology-compact-support} 中定义的阿贝尔群
$H^0_c(X, \mathcal{F})$ 一致。
\end{lemma}

\begin{proof}
选取在 $k$ 上的紧化 $j : X \to \overline{X}$。两种情形中，该群都被定义为
$H^0(\overline{X}, j_!\mathcal{F})$。对第一个版本，这是引理
\ref{more-etale-lemma-compactify-compact-support} 的结论；对第二个版本，则由构造成立。
\end{proof}

\begin{lemma}
\label{more-etale-lemma-finiteness-curves}
设 $k$ 是代数闭域，$X$ 是 $k$ 上维数 $\leq 1$ 的有限型有限型分离概形，
$\Lambda$ 是 Noether 环。设 $\mathcal{F}$ 是 $X$ 上的可构造挠
$\Lambda$-模层。则 $H^q_c(X, \mathcal{F})$ 是有限 $\Lambda$-模。
\end{lemma}

\begin{proof}
这是“\'Etale 上同调”定理
\ref{etale-cohomology-theorem-vanishing-affine-curves-coefficients} 的推论。
具体地，选取紧化 $j : X \to \overline{X}$。以 $X$ 的概形论闭包替换
$\overline{X}$ 后，可设 $\dim(\overline{X}) \leq 1$。于是
$H^q_c(X, \mathcal{F}) = H^q(\overline{X}, j_!\mathcal{F})$，可应用该定理。
\end{proof}

\begin{remark}[紧支撑上同调的协变性]
\label{more-etale-remark-covariance-compactly-supported}
设 $k$ 是域，$f : X \to Y$ 是 $k$ 上有限型分离概形之间的态射。若
$X$、$Y$ 与 $f$ 满足下列条件之一：
\begin{enumerate}
\item $f$ 为 \'etale；或
\item $f$ 平坦且拟有限；或
\item $f$ 拟有限且 $Y$ 几何单枝；或
\item $f$ 拟有限，且存在 $f$ 的加权 $w : X \to \mathbf{Z}$，
\end{enumerate}
则紧支撑上同调关于 $f$ 协变。更准确地，设 $\Lambda$ 是环，$K$ 是
$D^+_{tors}(Y_\etale, \Lambda)$ 的对象；若 $\Lambda$ 为挠环，也可设它是
$D(Y_\etale, \Lambda)$ 的对象。在假设 (1)--(4) 之一成立时，有典范映射
$$
\text{Tr}_{f, w, K} : f_!f^{-1}K \longrightarrow K
$$
迹映射的存在性见第 \ref{more-etale-section-weightings} 节；情形 (2) 与 (3) 分别见例
\ref{more-etale-example-trace-for-flat-quasi-finite} 和
\ref{more-etale-example-trace-for-quasi-finite-over-normal}。若
$p : X \to \Spec(k)$ 与 $q : Y \to \Spec(k)$ 表示结构态射，则由引理
\ref{more-etale-lemma-shriek-composition} 以及引理
\ref{more-etale-lemma-Rf-shriek-for-quasi-finite} 对拟有限分离态射 $f$ 给出的
$Rf_! = f_!$，有 $Rq_! \circ f_! = Rp_!$。因此可考虑映射
\begin{align*}
R\Gamma_c(X, f^{-1}K)
& =
R\Gamma(\Spec(k), Rp_!f^{-1}K) \\
& =
R\Gamma(\Spec(k), Rq_!f_!f^{-1}K) \\
& \xrightarrow{Rq_!\text{Tr}_{f, w, K}}
R\Gamma(\Spec(k), Rq_!K) \\
& =
R\Gamma_c(Y, K)
\end{align*}
特别地，若 $\Lambda$ 是挠环，则得到箭头
$$
\text{Tr}_f : R\Gamma_c(X, \Lambda) \longrightarrow R\Gamma_c(Y, \Lambda)
$$
该映射还有许多附加性质；例如，它与基域扩张相容。
\end{remark}










\section{一个可构造性结果}
\label{more-etale-section-family-smooth-curves-cohomology}

\noindent
我们“计算”常系数光滑射影曲线族的上同调。

\begin{lemma}
\label{more-etale-lemma-frobenius-linear-on-vb}
设 $p$ 是素数，$S$ 是 $\mathbf{F}_p$ 上的概形。设 $\mathcal{E}$ 是有限局部自由
$\mathcal{O}_S$-模，并把它视为 $S_\etale$ 上的 $\mathcal{O}_S$-模。
设 $F : \mathcal{E} \to \mathcal{E}$ 是 $S_\etale$ 上阿贝尔层的同态，
且对 $S_\etale$ 上 $\mathcal{O}_S$、$\mathcal{E}$ 的局部截面 $a$、$e$，满足
$F(a e) = a^pF(e)$。则
$$
\Coker(F - 1 : \mathcal{E} \to \mathcal{E})
$$
为零，且
$$
\Ker(F - 1 : \mathcal{E} \to \mathcal{E})
$$
是 $S_\etale$ 上的可构造阿贝尔层。
\end{lemma}

\noindent
该引理是“\'Etale 上同调”引理
\ref{etale-cohomology-lemma-F-1} 的推广。

\begin{proof}
可设 $S = \Spec(A)$，其中 $A$ 是 $\mathbf{F}_p$-代数，并设 $\mathcal{E}$ 是
自由 $A$-模 $Ae_1 \oplus \ldots \oplus Ae_n$ 所伴的拟凝聚模。写成
$F(e_i) = \sum a_{ij} e_j$。

\medskip\noindent
$F - 1$ 的满射性。只需证明：以一个忠实平坦 \'etale 扩张替换 $A$ 后，
任意元素 $\sum a_i e_i$（$a_i \in A$）都属于 $F - 1$ 的像。注意
$$
F(\sum x_ie_i) - \sum x_i e_i = \sum x_i^p a_{ij} e_j - \sum x_i e_i
$$
考虑 $A$-代数
$$
A' = A[x_1, \ldots, x_n]/(a_i + x_i - \sum\nolimits_j a_{ji} x_j^p)
$$
计算表明，$\text{d}x_i$ 在 $\Omega_{A'/A}$ 中为零，因而
$\Omega_{A'/A} = 0$。由于 $A'$ 是 $A$ 上有限型的，这蕴含
$\Spec(A') \to \Spec(A)$ 非分歧，因而拟有限。又因 $A'$ 由 $n$ 个元素生成，
且由 $n$ 个方程截出，故 $A'$ 是 $A$ 上的全局相对完全交。因此 $A'$ 在 $A$
上平坦，遂知 $A \to A'$ 是 \'etale 的（作为平坦且非分歧的环映射）。最后，
读者可通过直接验证 $\Spec(A') \to \Spec(A)$ 的所有几何纤维非空来证明
$A \to A'$ 忠实平坦；不过这也由“\'Etale 上同调”引理
\ref{etale-cohomology-lemma-F-1} 得出。最后，元素
$\sum x_i e_i \in A'e_1 \oplus \ldots \oplus A'e_n$ 在 $F - 1$ 下映到
$\sum a_i e_i$。

\medskip\noindent
核的可构造性。上述计算表明，$\Ker(F - 1)$ 由概形
$$
\Spec(A[x_1, \ldots, x_n]/(x_i - \sum\nolimits_j a_{ji} x_j^p))
$$
在 $S = \Spec(A)$ 上表示。由于这是 $S$ 上仿射且 \'etale 的概形，结论由
“\'Etale 上同调”引理 \ref{etale-cohomology-lemma-jshriek-constructible} 得出。
\end{proof}

\begin{lemma}
\label{more-etale-lemma-proper-smooth-family-curves}
设 $f : X \to S$ 是概形的固有光滑态射，其纤维几何连通且维数为 $1$。
设 $\ell$ 是素数。则 $R^qf_*\underline{\mathbf{Z}/\ell\mathbf{Z}}$
是一个可构造对象。
\end{lemma}

\begin{proof}
可设 $S$ 仿射，写成 $S = \Spec(A)$。若把 $A = \bigcup A_i$ 写成其有限型
$\mathbf{Z}$-子代数之并，则可找到某个 $i$ 以及有限型态射
$f_i : X_i \to S_i = \Spec(A_i)$，使其到 $S$ 的基变换为
$f : X \to S$；见“极限”引理 \ref{limits-lemma-descend-finite-presentation}。
增大 $i$ 后，可设 $f_i : X_i \to S_i$ 光滑、固有且相对维数为 $1$；
见“极限”引理
\ref{limits-lemma-eventually-proper}
\ref{limits-lemma-descend-smooth} 与
\ref{limits-lemma-descend-dimension-d}。由“态射的进一步结果”引理
\ref{more-morphisms-lemma-proper-flat-geom-red}，得到开子概形
$U_i \subset S_i$，使 $f_i : X_i \to S_i$ 在 $U_i$ 上的纤维几何连通。
于是 $S \to S_i$ 映入 $U_i$。可以用
$f_i : f_i^{-1}(U_i) \to U_i$ 替换 $X \to S$，从而化归到下一段讨论的情形。

\medskip\noindent
设 $S$ 为 Noether 概形。可写成 $S = U \cup Z$，其中 $U$ 是由 $\ell$
不消没所定义的开子概形，而 $Z = V(\ell) \subset S$。由于
$R^qf_*\underline{\mathbf{Z}/\ell\mathbf{Z}}$ 的构造与任意基变换交换
（“\'Etale 上同调”定理 \ref{etale-cohomology-theorem-proper-base-change}），
只需分别在 $U$ 与 $Z$ 上证明。于是化归到两个情形：
(a) $\ell$ 在 $S$ 上可逆；(b) $\ell$ 在 $S$ 上为零。

\medskip\noindent
情形 (a)。我们断言，此时层 $R^qf_*\underline{\mathbf{Z}/\ell\mathbf{Z}}$
在 $S$ 上有限且局部常值。首先，由固有基变换（采用“\'Etale 上同调”引理
\ref{etale-cohomology-lemma-proper-base-change-stalk} 的形式）以及有限性
（“\'Etale 上同调”定理
\ref{etale-cohomology-theorem-vanishing-affine-curves}），可知
$R^qf_*\underline{\mathbf{Z}/\ell\mathbf{Z}}$ 的茎有限。由“\'Etale 上同调”引理
\ref{etale-cohomology-lemma-sp-isom-proper-loc-cst-torsion}，
所有特殊化映射均为同构。故由“\'Etale 上同调”引理
\ref{etale-cohomology-lemma-characterize-locally-constant-module}，断言成立。

\medskip\noindent
情形 (b)。这里 $\ell = p$ 是素数，且 $S$ 是 $\Spec(\mathbf{F}_p)$ 上的概形。
由上述相同文献，已知 $R^qf_*\underline{\mathbf{Z}/p\mathbf{Z}}$ 的茎有限，
且在 $q \geq 2$ 时为零。由“\'Etale 上同调”引理
\ref{etale-cohomology-lemma-sections-upstairs}，有
$f_*\underline{\mathbf{Z}/p\mathbf{Z}} =
\underline{\mathbf{Z}/p\mathbf{Z}}$。只剩证明
$R^1f_*\underline{\mathbf{Z}/p\mathbf{Z}}$ 可构造。考虑 Artin--Schreier 序列
$$
0 \to \underline{\mathbf{Z}/p\mathbf{Z}} \to \mathcal{O}_X
\xrightarrow{F - 1} \mathcal{O}_X
\to 0
$$
见“\'Etale 上同调”第 \ref{etale-cohomology-section-artin-schreier} 节。
回忆，$f_*\mathcal{O}_X = \mathcal{O}_S$，且
$R^1f_*\mathcal{O}_X$ 是有限局部自由 $\mathcal{O}_S$-模，其秩等于
$X \to S$ 各纤维的亏格；见“代数曲线”引理
\ref{curves-lemma-genus-in-nodal-family-of-curves}。故有短正合列
$$
0 \to
\Coker(F - 1 : \mathcal{O}_S \to \mathcal{O}_S)
\to R^1f_*\underline{\mathbf{Z}/p\mathbf{Z}} \to
\Ker(F - 1 : R^1f_*\mathcal{O}_X \to R^1f_*\mathcal{O}_X)
\to 0
$$
应用引理 \ref{more-etale-lemma-frobenius-linear-on-vb} 即得结论。
\end{proof}

\begin{lemma}
\label{more-etale-lemma-proper-smooth-family-curves-modules}
设 $f : X \to S$ 是概形的固有光滑态射，其纤维几何连通且维数为 $1$。
设 $\Lambda$ 是 Noether 环，$M$ 是被整数 $n > 0$ 消去的有限 $\Lambda$-模。
则 $R^qf_*\underline{M}$ 是 $S$ 上的可构造 $\Lambda$-模层。
\end{lemma}

\begin{proof}
若对某个素数 $\ell$ 有 $n = \ell n'$，则得到有限 $\Lambda$-模的短正合列
$0 \to M[\ell] \to M \to M' \to 0$，且 $M'$ 被 $n'$ 消去。
这产生相应的常值层短正合列，继而产生正合列
$$
R^{q - 1}f_*\underline{M'} \to
R^qf_*\underline{M[n]} \to
R^qf_*\underline{M} \to
R^qf_*\underline{M'} \to
R^{q + 1}f_*\underline{M[n]}
$$
因此，若能在 $M$ 被某个素数消去时证明该结果，则对 $n$ 作归纳，并应用
“\'Etale 上同调”引理 \ref{etale-cohomology-lemma-constructible-abelian}，
即得一般情形。

\medskip\noindent
设 $\ell$ 是素数，且 $\ell$ 消去 $M$。于是可用 $\mathbf{F}_\ell$-代数
$\Lambda/\ell \Lambda$ 替换 $\Lambda$。具体地，把 $\underline{M}$ 视为
$\Lambda$-模层时所得的层 $R^qf_*\underline{M}$，与把 $\underline{M}$ 视为
$\Lambda/\ell \Lambda$-模层时算得的层 $R^qf_*\underline{M}$ 相同；
见“位点上同调”引理 \ref{sites-cohomology-lemma-modules-abelian-unbounded}。

\medskip\noindent
设 $\ell$ 是素数，且 $\ell$ 消去 $M$ 与 $\Lambda$。我们化归到 $M$ 为有限自由
$\Lambda$-模的情形。具体地，选取消解
$$
\ldots \to
\Lambda^{\oplus m_2} \to
\Lambda^{\oplus m_1} \to
\Lambda^{\oplus m_0} \to
M \to 0
$$
回忆，$f_*$ 在 $\Lambda$-模层上的上同调维数有限；见“\'Etale 上同调”引理
\ref{etale-cohomology-lemma-cohomological-dimension-proper} 以及“导出范畴”引理
\ref{derived-lemma-unbounded-right-derived}。因此，对某个充分大的整数 $a$，
$R^qf_*\underline{M}$ 是对象
$$
Rf_*(\underline{\Lambda^{\oplus m_a}} \to \ldots \to
\underline{\Lambda^{\oplus m_0}})
$$
在 $D(S_\etale, \Lambda)$ 中的第 $q$ 个上同调层。使用“导出范畴”引理
\ref{derived-lemma-two-ss-complex-functor} 的第一个谱序列（或者使用截断论证），
可知只需证明 $R^qf_*\underline{\Lambda})$ 可构造。

\medskip\noindent
此时终于可以使用等式
$$
(Rf_*\underline{\mathbf{Z}/\ell\mathbf{Z}})
\otimes_{\mathbf{Z}/\ell\mathbf{Z}}^\mathbf{L} \underline{\Lambda} =
Rf_*\underline{\Lambda}
$$
；它由“\'Etale 上同调”引理
\ref{etale-cohomology-lemma-projection-formula-proper-mod-n} 成立。
由于域 $\mathbf{Z}/\ell\mathbf{Z}$ 上任意模都平坦，得到
$$
(R^qf_*\underline{\mathbf{Z}/\ell\mathbf{Z}})
\otimes_{\mathbf{Z}/\ell\mathbf{Z}} \underline{\Lambda} =
R^qf_*\underline{\Lambda}
$$
故由“\'Etale 上同调”引理
\ref{etale-cohomology-lemma-tensor-constructible}，只需对
$R^qf_*\underline{\mathbf{Z}/\ell\mathbf{Z}}$ 证明结论。这正是引理
\ref{more-etale-lemma-proper-smooth-family-curves}。
\end{proof}






\section{上同调可构造的复形}
\label{more-etale-section-Dc}

\noindent
我们继续“\'Etale 上同调”第 \ref{etale-cohomology-section-Dc} 节开始的讨论。
特别地，对概形 $X$ 与 Noether 环 $\Lambda$，以 $D_c(X_\etale, \Lambda)$
表示 $D(X_\etale, \Lambda)$ 的严格全且饱和的三角子范畴；其中对象的
上同调层均为可构造 $\Lambda$-模层。

\begin{lemma}
\label{more-etale-lemma-qf-f-shriek-constructible}
设 $f : X \to Y$ 是局部拟有限且有限表示的概形态射。引理
\ref{more-etale-lemma-lqf-shriek-derived} 的函子
$f_! : D(X_\etale, \Lambda) \to D(Y_\etale, \Lambda)$ 把
$D_c(X_\etale, \Lambda)$ 送入 $D_c(Y_\etale, \Lambda)$。
\end{lemma}

\begin{proof}
由于函子 $f_!$ 正合，只需证明：对 $X_\etale$ 上任意可构造
$\Lambda$-模层 $\mathcal{F}$，$f_!\mathcal{F}$ 可构造。该问题在 $Y$ 上局部，
故可且确实设 $Y$ 仿射。于是 $X$ 拟紧且拟分离；见“态射”定义
\ref{morphisms-definition-finite-presentation}。设
$X = \bigcup_{i = 1, \ldots, n} X_i$ 是有限仿射开覆盖。由引理
\ref{more-etale-lemma-lqf-colimit-f-shriek}，只需证明
$f_{i, !}\mathcal{F}|_{X_i}$ 与 $f_{ii', !}\mathcal{F}|_{X_i \cap X_{i'}}$
可构造；其中 $f_i : X_i \to Y$ 与
$f_{ii'} : X_i \cap X_{i'} \to Y$ 是 $f$ 的限制。由于 $X_i$ 和
$X_i \cap X_{i'}$ 拟紧且分离，这意味着可设 $f$ 分离。由 Zariski 主定理
（采用“态射的进一步结果”引理
\ref{more-morphisms-lemma-quasi-finite-separated-pass-through-finite-addendum}
的形式），可选取分解 $f = g \circ j$，其中 $j : X \to X'$ 是开浸入，
$g : X' \to Y$ 有限且有限表示。由引理
\ref{more-etale-lemma-f-shriek-composition}，有 $f_! = g_! \circ j_!$。
由“\'Etale 上同调”引理 \ref{etale-cohomology-lemma-jshriek-constructible}，
$j_!\mathcal{F}$ 在 $X'$ 上可构造。态射 $g$ 有限，故由引理
\ref{more-etale-lemma-proper-f-shriek}，有 $g_! = g_*$。因此由“\'Etale 上同调”引理
\ref{etale-cohomology-lemma-finite-pushforward-constructible}，
$f_!\mathcal{F} = g_!j_!\mathcal{F} = g_*j_!\mathcal{F}$ 可构造。
\end{proof}

\begin{lemma}
\label{more-etale-lemma-constant-shriek-rel-dim-1}
设 $S$ 是有限维 Noether 仿射概形，$f : X \to S$ 是相对维数为 $1$ 的
分离、仿射、光滑态射。设 $\Lambda$ 是挠 Noether 环，$M$ 是有限
$\Lambda$-模。则 $Rf_!\underline{M}$ 的上同调层可构造。
\end{lemma}

\begin{proof}
我们对 $d = \dim(S)$ 作归纳来证明该结果。

\medskip\noindent
起始情形。若 $d = 0$，则只需证明 $R^qf_!\underline{M}$ 的茎是有限
$\Lambda$-模。若 $\overline{s}$ 是 $S$ 的几何点，则由引理
\ref{more-etale-lemma-stalk-R-f-shriek}，有
$(R^qf_!\underline{M})_{\overline{s}} = H^q_c(X_{\overline{s}}, \underline{M})$。
由引理 \ref{more-etale-lemma-finiteness-curves}，这是有限 $\Lambda$-模。

\medskip\noindent
归纳步骤。只需找到稠密开集 $U \subset S$，使
$Rf_!\underline{M}|_U$ 的上同调层可构造。事实上，由归纳假设以及
$Rf_!\underline{M}$ 的构造与所有基变换交换这一事实（引理
\ref{more-etale-lemma-base-change-shriek}），$Rf_!\underline{M}$ 在补集
$S \setminus U$ 上的限制将有可构造上同调层。具体地，设
$\eta \in S$ 是 $S$ 某个不可约分支的泛点。只需找到 $\eta$ 的开邻域 $U$，
使 $Rf_!\underline{M}$ 在 $U$ 上的限制可构造。下一段将完成此事。

\medskip\noindent
给定泛点 $\eta \in S$，选取图
$$
\xymatrix{
\overline{Y}_1 \amalg \ldots \amalg \overline{Y}_n \ar[rd] &
Y_1 \amalg \ldots \amalg Y_n \ar[r]_-\nu \ar[d] \ar[l]^j &
X_V \ar[r] \ar[d] &
X_U \ar[r] \ar[d] &
X \ar[d]^f \\
& T_1 \amalg \ldots \amalg T_n \ar[r] &
V \ar[r] &
U \ar[r] &
S
}
$$
，如“态射的进一步结果”引理
\ref{more-morphisms-lemma-make-good-curves}。我们将证明
$Rf_!\underline{M}|_U$ 可构造。首先，由于 $V \to U$ 有限且满，只需证明
到 $V$ 的拉回可构造；见“\'Etale 上同调”引理
\ref{etale-cohomology-lemma-check-constructible}。由于 $Rf_!$ 的构造与基变换
交换，只需证明 $R(X_V \to V)_!\underline{M}$ 可构造。令 $W \subset X_V$
为“态射的进一步结果”引理
\ref{more-morphisms-lemma-make-good-curves} (4) 给出的开子概形。令
$Z \subset X_V$ 为 $W$ 在 $X_V$ 中的补集所带的既约诱导概形结构。
则 $Z \to V$ 的纤维维数为 $0$（因为 $W$ 在各纤维中稠密），故
$Z \to V$ 拟有限。由特异三角
$$
R(W \to V)_!\underline{M} \to
R(X_V \to V)_!\underline{M} \to
R(Z \to V)_!\underline{M} \to
\ldots
$$
以及引理 \ref{more-etale-lemma-relative-triangle-associated-to-open} 和
\ref{more-etale-lemma-qf-f-shriek-constructible}，可知只需证明
$R(W \to V)_!\underline{M}$ 的上同调层可构造。其次，有
$$
R(W \to V)_!\underline{M} = R(\nu^{-1}(W) \to V)_!\underline{M}
$$
，因为态射 $\nu : \nu^{-1}(W) \to W$ 是增厚，且可应用引理
\ref{more-etale-lemma-shriek-and-thickening}。接着，以
$Z' \subset \coprod \overline{Y}_i$ 表示开集 $j(\nu^{-1}(W))$ 的补集。
$Z' \to V$ 仍拟有限。再次使用特异三角
$$
R(\nu^{-1}(W) \to V)_!\underline{M} \to
R(\coprod \overline{Y}_i \to V)_!\underline{M} \to
R(Z' \to V)_!\underline{M} \to
\ldots
$$
，可知只需证明
$$
R(\coprod \overline{Y}_i \to V)_!\underline{M} =
\bigoplus\nolimits_i R(\overline{Y}_i \to V)_!\underline{M} =
\bigoplus\nolimits_i R(T_i \to V)_!R(\overline{Y}_i \to T_i)_!\underline{M}
$$
的上同调层可构造（第二个等号由引理
\ref{more-etale-lemma-shriek-composition} 成立）。对
$R(\overline{Y}_i \to T_i)_!\underline{M}$，结论是引理
\ref{more-etale-lemma-proper-smooth-family-curves-modules}。又因 $T_i \to V$ 有限 \'etale，
可应用引理 \ref{more-etale-lemma-qf-f-shriek-constructible}，故证明完成。
\end{proof}

\begin{lemma}
\label{more-etale-lemma-loc-constant-shriek-rel-dim-1}
设 $Y$ 是有限维 Noether 仿射概形，$\Lambda$ 是挠 Noether 环。
设 $\mathcal{F}$ 是开子概形 $U \subset \mathbf{A}^1_Y$ 上有限型局部常值的
$\Lambda$-模层。则 $Rf_!\mathcal{F}$ 的上同调层可构造，其中
$f : U \to Y$ 是结构态射。
\end{lemma}

\begin{proof}
可把 $\Lambda$ 分解为乘积
$\Lambda = \Lambda_1 \times \ldots \times \Lambda_r$
，其中对某个素数 $\ell_i$，$\Lambda_i$ 是 $\ell_i$-准素的。因此可设存在素数
$\ell$ 与整数 $n > 0$，使 $\ell^n$ 消去 $\Lambda$（因而也消去
$\mathcal{F}$）。

\medskip\noindent
由于 $U$ 为 Noether 概形，可见 $U$ 只有有限多个连通分支。因此可设 $U$ 连通。
令 $g : U' \to U$ 为“\'Etale 上同调”引理
\ref{etale-cohomology-lemma-pullback-filtered-modules} 中构造的有限 \'etale 覆盖。
“\'Etale 上同调”第 \ref{etale-cohomology-section-trace-method} 节的讨论给出映射
$$
\mathcal{F} \to g_*g^{-1}\mathcal{F} \to \mathcal{F}
$$
，其复合为同构。因此只需对 $g_*g^{-1}\mathcal{F}$ 证明结论。另一方面，
由引理 \ref{more-etale-lemma-shriek-composition}，有
$Rf_!g_*g^{-1}\mathcal{F} = R(f \circ g)_!g^{-1}\mathcal{F}$。
由于（根据我们对 $g$ 的选取）$g^{-1}\mathcal{F}$ 有有限滤过，其分次商是形如
$\underline{M}$ 的常值 $\Lambda$-模层，其中 $M$ 是有限 $\Lambda$-模，
这就化归到引理 \ref{more-etale-lemma-constant-shriek-rel-dim-1} 中已证明的情形。
\end{proof}

\begin{lemma}
\label{more-etale-lemma-constructible-shriek-rel-dim-1}
设 $Y$ 是仿射概形，$\Lambda$ 是 Noether 环。设 $\mathcal{F}$ 是
$\mathbf{A}^1_Y$ 上的可构造挠 $\Lambda$-模层。则 $Rf_!\mathcal{F}$ 的
上同调层可构造，其中 $f : \mathbf{A}^1_Y \to Y$ 是结构态射。
\end{lemma}

\begin{proof}
设 $\mathcal{F}$ 被 $n > 0$ 消去。于是可用 $\Lambda/n\Lambda$ 替换
$\Lambda$ 而不改变 $Rf_!\mathcal{F}$。因此可且确实设 $\Lambda$ 为挠环。

\medskip\noindent
写 $Y = \Spec(R)$。若把 $R = \bigcup R_i$ 写成其有限型
$\mathbf{Z}$-子代数之并，则可找到某个 $i$，使 $\mathcal{F}$ 是
$\mathbf{A}^1_{R_i}$ 上某个可构造 $\Lambda$-模层的拉回；见“\'Etale 上同调”引理
\ref{etale-cohomology-lemma-category-is-colimit}。因此可设 $Y$ 是有限维 Noether 概形。

\medskip\noindent
设 $Y$ 是维数 $d = \dim(Y)$ 的有限维 Noether 概形，且 $\Lambda$ 为挠环。
我们对 $d$ 作归纳证明该结果。

\medskip\noindent
起始情形。若 $d = 0$，则只需证明 $R^qf_!\mathcal{F}$ 的茎是有限
$\Lambda$-模。若 $\overline{y}$ 是 $Y$ 的几何点，则由引理
\ref{more-etale-lemma-stalk-R-f-shriek}，有
$(R^qf_!\mathcal{F})_{\overline{y}} = H^q_c(X_{\overline{y}}, \mathcal{F})$。
由引理 \ref{more-etale-lemma-finiteness-curves}，这是有限 $\Lambda$-模。

\medskip\noindent
归纳步骤。只需找到稠密开集 $V \subset Y$，使
$Rf_!\mathcal{F}|_V$ 的上同调层可构造。事实上，由归纳假设以及
$Rf_!\mathcal{F}$ 的构造与所有基变换交换这一事实（引理
\ref{more-etale-lemma-base-change-shriek}），$Rf_!\mathcal{F}$ 在补集
$Y \setminus V$ 上的限制将有可构造上同调层。由可构造 $\Lambda$-模层的定义，
存在稠密开子概形 $U \subset \mathbf{A}^1_Y$，使 $\mathcal{F}|_U$ 是有限型局部常值
$\Lambda$-模层。以 $Z \subset \mathbf{A}^1_Y$ 表示补集（视为既约闭子概形）。
注意，$U$ 包含 $\mathbf{A}^1_Y \to Y$ 在 $Y$ 各不可约分支的泛点
$\xi_1, \ldots, \xi_n$ 上方纤维的所有泛点。因此 $Z \to Y$ 在
$\xi_1, \ldots, \xi_n$ 上方的纤维有限。以一个稠密开集替换 $Y$（这是允许的）后，
可设 $Z \to Y$ 有限；见“态射”引理 \ref{morphisms-lemma-generically-finite}。
由引理 \ref{more-etale-lemma-relative-triangle-associated-to-open} 的特异三角以及
$Z \to Y$ 的结果（引理 \ref{more-etale-lemma-qf-f-shriek-constructible}），化归到证明
$R(U \to Y)_!\mathcal{F}$ 的上同调层可构造。这就是引理
\ref{more-etale-lemma-loc-constant-shriek-rel-dim-1}。
\end{proof}

\begin{theorem}
\label{more-etale-theorem-constructible-shriek}
设 $f : X \to Y$ 是拟紧拟分离概形之间有限表示的分离态射，$\Lambda$ 是
Noether 环。设 $K$ 是 $D^+_{tors, c}(X_\etale, \Lambda)$ 的对象；若
$\Lambda$ 为挠环，也可设它是 $D_c(X_\etale, \Lambda)$ 的对象。则
$Rf_!K$ 的上同调层可构造；即 $Rf_!K$ 属于
$D^+_{tors, c}(Y_\etale, \Lambda)$，或者在 $\Lambda$ 为挠环时属于
$D_c(Y_\etale, \Lambda)$。
\end{theorem}

\begin{proof}
该问题在 $Y$ 上局部，故可且确实设 $Y$ 仿射。由归纳原理和引理
\ref{more-etale-lemma-relative-mayer-vietoris}，化归到 $X$ 也仿射的情形。

\medskip\noindent
设 $X$ 与 $Y$ 仿射。由于 $X$ 有限表示，可选取有限表示的闭浸入
$i : X \to \mathbf{A}^n_Y$。若 $p : \mathbf{A}^n_Y \to Y$ 表示结构态射，
则由引理 \ref{more-etale-lemma-shriek-composition}，有 $Rf_! = Rp_! \circ Ri_!$。
由引理 \ref{more-etale-lemma-qf-f-shriek-constructible}，结论对 $Ri_! = i_!$ 成立。
因此可设 $f$ 是投影态射 $\mathbf{A}^n_Y \to Y$。由于可把 $f$ 视为复合
$$
X = \mathbf{A}^n_Y \to \mathbf{A}^{n - 1}_Y \to
\mathbf{A}^{n - 2}_S \to \ldots \to
\mathbf{A}^1_Y \to Y
$$
，可设 $n = 1$。

\medskip\noindent
设 $Y$ 仿射且 $X = \mathbf{A}^1_Y$。由于 $Rf_!$ 的上同调维数有限
（引理 \ref{more-etale-lemma-derived-lower-shriek-bounded}），可设 $K$ 下有界。
使用“导出范畴”引理 \ref{derived-lemma-two-ss-complex-functor} 的第一个谱序列
（或者使用截断论证），化归到引理
\ref{more-etale-lemma-constructible-shriek-rel-dim-1} 的结果。
\end{proof}








\section{应用}
\label{more-etale-section-applications}

\noindent
本节给出定理 \ref{more-etale-theorem-constructible-shriek} 的一些应用。

\begin{lemma}
\label{more-etale-lemma-finiteness-compactly-supported}
设 $k$ 是代数闭域，$X$ 是 $k$ 上的有限型分离概形，$\Lambda$ 是 Noether 环。
设 $K$ 是 $D^+_{tors, c}(X_\etale, \Lambda)$ 的对象；若 $\Lambda$ 为挠环，
也可设它是 $D_c(X_\etale, \Lambda)$ 的对象。则对所有
$i \in \mathbf{Z}$，$H^i_c(X, K)$ 都是有限 $\Lambda$-模。
\end{lemma}

\begin{proof}
这是定理 \ref{more-etale-theorem-constructible-shriek} 与第
\ref{more-etale-section-compactly-supported-cohomology} 节中紧支撑上同调定义的直接推论。
\end{proof}

\begin{proposition}
\label{more-etale-proposition-loc-cst-torsion}
设 $f : X \to S$ 是概形的光滑固有态射，$\Lambda$ 是 Noether 环。
设 $\mathcal{F}$ 是 $X_\etale$ 上有限型局部常值的 $\Lambda$-模层，并且对
$X$ 的每个几何点 $\overline{x}$，茎 $\mathcal{F}_{\overline{x}}$ 都被某个
与 $\overline{x}$ 的剩余特征互素的整数 $n > 0$ 消去。则对所有
$i \in \mathbf{Z}$，$R^if_*\mathcal{F}$ 都是 $S_\etale$ 上有限型局部常值的
$\Lambda$-模层。
\end{proposition}

\begin{proof}
该问题在 $S$ 上局部，故可设 $S$ 仿射。对 $X$ 的点 $x$，以
$n_x \geq 1$ 表示消去 $\mathcal{F}_{\overline{x}}$ 的最小整数；这里
$\overline{x}$ 是 $X$ 中位于 $x$ 上方的某个（等价地，任意）几何点。
由于 $X$ 拟紧（它在仿射概形上固有），存在有限 \'etale 覆盖
$\{U_j \to X\}_{j = 1, \ldots, m}$，使 $\mathcal{F}|_{U_j}$ 为常值层。
由于 $U_j \to X$ 是开映射，函数 $x \mapsto n_x$ 局部常值且仅取有限多个值。
因此得到如下开闭子概形的有限分解：
$X = X_1 \amalg \ldots \amalg X_N$，使得 $n_x = n$ 当且仅当
$x \in X_n$。于是只需对所诱导的态射 $X_n \to S$ 以及 $\mathcal{F}$ 在
$X_n$ 上的限制证明该引理。因此可且确实设存在整数 $n > 0$，使
$\mathcal{F}$ 被 $n$ 消去，且 $n$ 与 $X$ 所有剩余域的剩余特征互素。

\medskip\noindent
由于 $f$ 光滑且固有，像 $f(X) \subset S$ 开且闭。因此可用 $f(X)$ 替换
$S$ 并设 $f(X) = S$。特别地，可设 $n$ 在定义仿射概形 $S$ 的环中可逆。

\medskip\noindent
本段化归到 $S$ 为 Noether 概形的情形。对某个 $\mathbf{Z}[1/n]$-代数 $A$，
写 $S = \Spec(A)$。把 $A = \bigcup A_i$ 写成其有限型
$\mathbf{Z}[1/n]$-子代数之并。可找到某个 $i$ 以及有限型态射
$f_i : X_i \to S_i = \Spec(A_i)$，使其到 $S$ 的基变换为
$f : X \to S$；见“极限”引理 \ref{limits-lemma-descend-finite-presentation}。
增大 $i$ 后，可设 $f_i : X_i \to S_i$ 光滑且固有；见“极限”引理
\ref{limits-lemma-eventually-proper}
\ref{limits-lemma-descend-smooth} 与
\ref{limits-lemma-descend-dimension-d}.
由“\'Etale 上同调”引理
\ref{etale-cohomology-lemma-category-loc-cst-is-colimit}
，可知存在某个 $i$ 以及有限型局部常值的 $\Lambda$-模层
$\mathcal{F}_i$，其到 $X$ 的拉回与 $\mathcal{F}$ 同构。由于
$\mathcal{F}$ 被 $n$ 消去，可用 $\Ker(n : \mathcal{F}_i \to \mathcal{F}_i)$
替换 $\mathcal{F}_i$，并设同一性质对 $\mathcal{F}_i$ 成立。
这就化归到下一段讨论的情形。

\medskip\noindent
设有整数 $n \geq 1$，基概形 $S$ 为 Noether 概形且位于
$\mathbf{Z}[1/n]$ 上，而 $\mathcal{F}$ 为 $n$-挠层。由定理
\ref{more-etale-theorem-constructible-shriek}，各层 $R^if_*\mathcal{F}$ 都是可构造
$\Lambda$-模层。由“\'Etale 上同调”引理
\ref{etale-cohomology-lemma-sp-isom-proper-torsion-loc-ac}，
$R^if_*\mathcal{F}$ 的特殊化映射总是同构。结论遂由“\'Etale 上同调”引理
\ref{etale-cohomology-lemma-characterize-locally-constant-module} 得出。
\end{proof}







\section{导出上叹号的进一步结果}
\label{more-etale-section-more-derived-upper-shriek}

\noindent
设 $\Lambda$ 是挠环。考虑交换图
$$
\xymatrix{
U \ar[rr]_j \ar[rd]_g & & U' \ar[ld]^{g'} \\
& Y
}
$$
，其中各概形拟紧且拟分离，$g$ 与 $g'$ 分离且有限型，而 $j$ 为 \'etale。
这诱导典范映射
$$
Rg_!\Lambda \longrightarrow Rg'_!\Lambda
$$
成立于 $D(Y_\etale, \Lambda)$ 中。具体地，由引理
\ref{more-etale-lemma-shriek-composition} 与
\ref{more-etale-lemma-Rf-shriek-for-quasi-finite}，有 $Rg_! = Rg'_! \circ j_!$。
另一方面，由于 $j_!$ 是 $j^{-1}$ 的左伴随，有余单位
$\text{Tr}_j : j_!\Lambda = j_!j^{-1}\Lambda \to \Lambda$；
也称它为 $j$ 的迹映射，见评注 \ref{more-etale-remark-trace-etale-counit}。
上述映射构造为复合
$$
Rg_!\Lambda = Rg'_!j_!\Lambda
\xrightarrow{Rg'_! \text{Tr}_j}
Rg'_!\Lambda
$$
再给定一个 \'etale 态射 $j' : U' \to U''$，其中某个
$g'' : U'' \to Y$ 分离且有限型，则映射对 $j$ 与 $j'$ 的复合
$$
Rg_!\Lambda \longrightarrow Rg'_!\Lambda \longrightarrow Rg''_!\Lambda
$$
等于为 $j' \circ j$ 构造的映射
$Rg_!\Lambda \longrightarrow Rg''_!\Lambda$。这由迹映射的相应陈述得出；
更一般的情形见引理 \ref{more-etale-lemma-trace-composition}。

\medskip\noindent
设 $f : X \to Y$ 是拟紧拟分离概形之间的有限型分离态射。于是得到函子
$$
X_{affine, \etale}
\longrightarrow
\left\{
\begin{matrix}
\text{有限型分离概形，底为 }Y\\
\text{其间的态射为 \'etale 态射}
\end{matrix}
\right\}
$$
因此上述构造确定函子
$X_{affine, \etale}^{opp} \to D(Y_\etale, \Lambda)$，
它把 $U$ 送到 $R(U \to Y)_!\Lambda$。

\begin{lemma}
\label{more-etale-lemma-describe-Rf-upper-shriek}
设 $f : X \to Y$ 是拟紧拟分离概形之间的有限型分离态射，$\Lambda$ 是挠环。
设 $K \in D(Y_\etale, \Lambda)$。对 $n \in \mathbf{Z}$，上同调层
$H^n(Rf^!K)$ 在 $X_{affine, \etale}$ 上的限制，是预层
$$
U \longmapsto \Hom_Y(R(U \to Y)_!\Lambda, K[n])
$$
所伴的层。关于 $R(U \to Y)_!\Lambda$ 的函子性质，见上面的讨论。
\end{lemma}

\begin{proof}
设 $j : U \to X$ 是 $X_{affine, \etale}$ 的对象，并令 $g = f \circ j$。
回忆，对 $D(X_\etale, \Lambda)$ 中任意 $M$，有
$\Hom_X(j_!\Lambda, M[n]) = H^n(U, M)$。于是 $H^n(Rf^!K)$ 是预层
$$
U \mapsto H^n(U, Rf^!K) = \Hom_X(j_!\Lambda, Rf^!K[n]) =
\Hom_Y(Rf_!j_!\Lambda, K[n] = \Hom_Y(Rg_!\Lambda, K[n])
$$
所伴的层。略去如下验证：转移映射由我们上面构造的各对象
$Rg_!\Lambda = R(U \to Y)_!\Lambda$ 之间的转移映射给出。
\end{proof}
% Generated mechanically from the frozen trace.tex target.
% Do not edit this generated file; edit the bound locale source instead.

% OK, start here.
%

\setcounter{chapter}{63}
\chapter{迹公式}



\phantomsection
\label{trace-section-phantom}



\section{引言}
\label{trace-section-introduction}

\noindent
这是 Johan de Jong 于 2009 年秋季在哥伦比亚大学讲授的
\'etale 上同调课程第二部分的讲义。最初的笔记记录者是
Thibaut Pugin、Zachary Maddock 和 Min Lee。我们将逐步补充对
Stacks Project 其余部分中背景材料的引用，并为所有陈述给出严格证明。









\section{迹公式}
\label{trace-section-trace-formula}

\noindent
典型的 \'etale 上同调课程通常会陈述并证明固有基变换定理、
光滑基变换定理、纯性以及 Poincar\'e 对偶；这些内容都可见
\cite[Arcata]{SGA4.5}。我们在此改而按照 Deligne 在
\cite[Rapport]{SGA4.5} 中的论述，研究 Frobenius 的迹公式。
我们只考察维数 1 的情形，但借助固有基变换，这已足以处理一般情形。
由于所考虑的所有上同调群都是 \'etale 上同调，我们省略下标
$_\etale$。下面描述所要得到的公式。设 $X$ 是有限域 $k$ 上维数为
1 的有限型概形，$\ell$ 是素数，而 $\mathcal{F}$ 是可构造平坦
$\mathbf{Z}/\ell^n\mathbf{Z}$-层。则
\begin{equation}
\label{trace-equation-trace-formula-initial}
\sum\nolimits_{x \in X(k)}
\text{Tr}(\text{Frob} | \mathcal{F}_{\bar x}) =
\sum\nolimits_{i = 0}^2
(-1)^i \text{Tr}(\pi_X^* | H^i_c(X \otimes_k \bar k, \mathcal{F}))
\end{equation}
是 $\mathbf{Z}/\ell^n\mathbf{Z}$ 中的等式。正如后文将看到的，这一表述
按其当前写法略有不当。不过，我们仍先说明其中出现的各个符号。




\section{Frobenius 态射}
\label{trace-section-frobenii}

\noindent
本节将证明一个“令人费解的”定理。该定理的拓扑类比如下。

\begin{exercise}
\label{trace-exercise-baffling}
设 $X$ 是拓扑空间，$g : X \to X$ 是连续映射，并且对 $X$ 的每个开集
$U$ 都有 $g^{-1}(U) = U$。则 $g$ 在 $X$ 的上同调上诱导恒等映射
（对任意系数均成立）。
\end{exercise}

\noindent
现在转向 \'etale 位点上的陈述。

\begin{lemma}
\label{trace-lemma-baffling}
设 $X$ 是概形，$g : X \to X$ 是态射。假设对每个 \'etale 态射
$\varphi : U \to X$，都存在同构
$$
\xymatrix{
U \ar[rd]_\varphi \ar[rr]^-\sim & & {U
\times_{\varphi, X, g} X} \ar[ld]^{\text{pr}_2} \\
& X
}
$$
，且它关于 $U$ 具有函子性。则 $g$ 在上同调上诱导恒等映射
（对任意层均成立）。
\end{lemma}

\begin{proof}
证明是形式的，并无困难。
\end{proof}

\noindent
关于 Frobenius 态射的不同变体，见“代数簇”第
\ref{varieties-section-frobenius} 节。

\begin{theorem}[令人费解定理]
\label{trace-theorem-baffling}
设 $X$ 是特征 $p > 0$ 的概形。则绝对 Frobenius 通过拉回在上同调上
诱导平凡映射；换言之，对每个整数 $j\geq 0$，
$$
F_X^* : H^j (X, \underline{\mathbf{Z}/n\mathbf{Z}}) \longrightarrow H^j (X,
\underline{\mathbf{Z}/n\mathbf{Z}})
$$
是恒等映射。
\end{theorem}

\noindent
该定理完全是形式的。不过，回顾如何计算上同调类的拉回仍很有益。
这里只说明：当上同调与 {\v C}ech 上同调一致时，只需利用位点上的
纤维积拉回 {\v C}ech 上闭链即可。一般情形相当技术性，见“超覆叠”
定理 \ref{hypercovering-theorem-cohomology-hypercoverings}。为证明该定理，
只须验证 Frobenius 满足引理 \ref{trace-lemma-baffling} 的假设。

\begin{proof}[定理 \ref{trace-theorem-baffling} 的证明]
须验证上述函子性同构的存在性。对 \'etale 态射
$\varphi : U \to X$，考虑图
$$
\xymatrix{
U \ar@{-->}[rd] \ar@/^1pc/[rrd]^{F_U}
\ar@/_1pc/[rdd]_\varphi \\
& {U \times_{\varphi, X, F_X} X} \ar[r]_-{\text{pr}_1}
\ar[d]^{\text{pr}_2} & U \ar[d]^\varphi \\
& X \ar[r]^{F_X} & X.
}
$$
虚箭头既是 \'etale 态射又是泛同胚，因而是同构；见“\'Etale 态射”
引理 \ref{etale-lemma-relative-frobenius-etale}。
\end{proof}

%10.22.09

\begin{definition}
\label{trace-definition-geometric-frobenius}
设 $k$ 是含 $q = p^f$ 个元素的有限域，$X$ 是 $k$ 上概形。
$X$ 的{\it 几何 Frobenius} 是 $\Spec(k)$ 上的态射
$\pi_X : X \to X$，它等于 $F_X^f$。
\end{definition}

\noindent
由于 $\pi_X$ 是 $k$ 上态射，可将它基变换到任意 $k$-概形。特别地，
可将它基变换到代数闭包 $\bar k$，得到态射
$\pi_X : X_{\bar k} \to X_{\bar k}$。仍用 $\pi_X$ 表示这一基变换
不会造成混淆，因为 $X_{\bar k}$ 本身并没有自己的几何 Frobenius。

\begin{lemma}
\label{trace-lemma-sheaf-over-finite-field-has-frobenius-descent}
设 $\mathcal{F}$ 是 $X_\etale$ 上的层。则存在典范同构
$\pi_X^{-1} \mathcal{F} \cong \mathcal{F}$ 与
$\mathcal{F} \cong {\pi_X}_*\mathcal{F}$。
\end{lemma}

\noindent
这对 fppf 位点不成立。

\begin{proof}
设 $\varphi : U \to X$ 是 \'etale 态射。回忆
${\pi_X}_* \mathcal{F} (U) = \mathcal{F} (U \times_{\varphi, X, \pi_X} X)$。
由于 $\pi_X = F_X^f$，由定理 \ref{trace-theorem-baffling} 的证明可知，
存在函子性同构
$$
\xymatrix{
U \ar[rd]_{\varphi} \ar[rr]_-{\gamma_U}
& & U \times_{\varphi, X, \pi_X} X \ar[ld]^{\text{pr}_2} \\
& X
}
$$
，其中 $\gamma_U = (\varphi, F_U^f)$。现在定义同构
$$
\mathcal{F} (U) \longrightarrow {\pi_X}_* \mathcal{F} (U) =
\mathcal{F} (U \times_{\varphi, X, \pi_X} X)
$$
为 $\mathcal{F}$ 沿 $\gamma_U^{-1}$ 的限制映射。另一个同构类似。
\end{proof}

\begin{remark}
\label{trace-remark-may-be-confusing}
$F^f_U$ 可能等于 $\pi_U$，也可能不等于。
\end{remark}

\noindent
继续讨论含 $q = p^f$ 个元素的有限域 $k$ 上概形 $X$ 的层上同调。
固定 $k$ 的代数闭包 $\bar k$，并以
$G_k = \text{Gal}(\bar k/k)$ 表示 $k$ 的绝对 Galois 群。
设 $\mathcal{F}$ 是 $X_\etale$ 上的阿贝尔层。如下在上同调群
$H^j (X_{\bar k}, \mathcal{F}|_{X_{\bar k}})$ 上定义左
$G_k$-模结构：若 $\sigma \in G_k$，则图
$$
\xymatrix{
X_{\bar k} \ar[rd] \ar[rr]^{\Spec(\sigma) \times \text{id}_X} & &
X_{\bar k} \ar[ld] \\
& X
}
$$
交换。因此，对 $\xi \in H^j (X_{\bar k}, \mathcal{F}|_{X_{\bar k}})$，
可令
$$
\sigma \cdot \xi := (\Spec(\sigma) \times \text{id}_X)^*\xi \in
H^j(X_{\bar k}, (\Spec(\sigma) \times \text{id}_X)^{-1}
\mathcal{F}|_{X_{\bar k}})
= H^j (X_{\bar k}, \mathcal{F}|_{X_{\bar k}}),
$$
，其中最后一个等号来自前图的交换性。由此，后一个群成为
$G_k$-模。

\begin{lemma}
\label{trace-lemma-two-actions-agree}
在上述情形中，以 $\alpha : X \to \Spec(k)$ 表示结构态射。考虑茎
$(R^j\alpha_*\mathcal{F})_{\Spec(\bar k)}$，并赋予它“\'Etale 上同调”
第 \ref{etale-cohomology-section-galois-action-stalks} 节所述的自然
Galois 作用。则同一
$$
(R^j\alpha_*\mathcal{F})_{\Spec(\bar k)} \cong H^j (X_{\bar k},
\mathcal{F}|_{X_{\bar k}})
$$
由“\'Etale 上同调”定理
\ref{etale-cohomology-theorem-higher-direct-images} 给出，并且是
$G_k$-模同构。
\end{lemma}

\noindent
把
$(R^j\alpha_!\mathcal{F})_{\Spec(\bar k)}$
与 $H^j_c (X_{\bar k}, \mathcal{F}|_{X_{\bar k}})$ 相比较时，也有类似结论。

\begin{proof}
略。
\end{proof}

\begin{definition}
\label{trace-definition-arithmetic-frobenius}
{\it 算术 Frobenius} 是 $G_k$ 中的映射
$\text{frob}_k : \bar k \to \bar k$，$x \mapsto x^q$。
\end{definition}

\begin{theorem}
\label{trace-theorem-geometric-arithmetic-inverse}
设 $\mathcal{F}$ 是 $X_\etale$ 上的阿贝尔层。则对每个
$j\geq 0$，$\text{frob}_k$ 在上同调群 $H^j(X_{\bar k},
\mathcal{F}|_{X_{\bar k}})$ 上的作用等于映射 $\pi_X^*$ 的逆。
\end{theorem}

\noindent
映射 $\pi_X^*$ 定义为复合
$$
H^j(X_{\bar k}, \mathcal{F}|_{X_{\bar k}}) \xrightarrow{{\pi_X}_{\bar k}^*}
H^j(X_{\bar k}, (\pi_X^{-1} \mathcal{F})|_{X_{\bar k}}) \cong
H^j(X_{\bar k}, \mathcal{F}|_{X_{\bar k}}).
$$
，其中最后一个同构来自引理
\ref{trace-lemma-sheaf-over-finite-field-has-frobenius-descent} 的典范同构
$\pi_X^{-1} \mathcal{F} \cong \mathcal{F}$。

\begin{proof}
复合 $X_{\bar k} \xrightarrow{\Spec(\text{frob}_k)} X_{\bar k}
\xrightarrow{\pi_X} X_{\bar k}$ 等于 $F_{X_{\bar k}}^f$，所以结论来自
适当推广到非平凡系数的令人费解定理。注意，上述复合在如下意义下交换：
$F_{X_{\bar k}}^f = \pi_X \circ \Spec(\text{frob}_k) =
\Spec(\text{frob}_k) \circ \pi_X$.
\end{proof}

\begin{definition}
\label{trace-definition-geometric-frobenius-on-stalk}
若 $x \in X(k)$ 是有理点，$\bar x : \Spec(\bar k) \to X$ 是位于 $x$
上方的几何点，则以 $\pi_x : \mathcal{F}_{\bar x} \to
\mathcal{F}_{\bar x}$ 表示 $\text{frob}_k^{-1}$ 的作用，并称之为
{\it 几何 Frobenius}\footnote{该记号并非标准记号。在 \cite{SGA4.5}
中，此算子记为 $F_x$。我们将来很可能会更改这里的记号。}
\end{definition}

\noindent
现在可对迹公式 (\ref{trace-equation-trace-formula-initial}) 作出更精确的陈述
（尽管它仍是错误的）。设 $X$ 是有限域 $k$ 上维数为 1 的有限型概形，
$\ell$ 是素数，而 $\mathcal{F}$ 是可构造平坦
$\mathbf{Z}/\ell^n\mathbf{Z}$-层。则
\begin{equation}
\label{trace-equation-trace-formula-second}
\sum\nolimits_{x \in X(k)}
\text{Tr}(\pi_x | \mathcal{F}_{\bar x})
=
\sum\nolimits_{i = 0}^2
(-1)^i \text{Tr}(\pi_X^* | H^i_c(X_{\bar k}, \mathcal{F}))
\end{equation}
是 $\mathbf{Z}/\ell^n\mathbf{Z}$ 中的等式。该等式之所以错误，是因为
右端的迹对这里所考虑的这类层没有意义。在处理这一问题之前，我们先说明
为何会出现几何 Frobenius（除了它本身是一个自然态射这一事实之外）。

\medskip\noindent
先考虑 $X = \mathbf{P}^1_k$ 且
$\mathcal{F} = \underline{\mathbf{Z}/\ell\mathbf{Z}}$ 的情形。对任意点，
Galois 模 $\mathcal{F}_{\bar x}$ 都是平凡的。因此，对任意作用于
$\mathcal{F}_{\bar x}$ 的态射 $\varphi$，左端为
$$
\sum\nolimits_{x \in X(k)} \text{Tr}(\varphi | \mathcal{F}_{\bar x}) =
\#\mathbf{P}^1_k(k) = q+1.
$$
由于 $\mathbf{P}^1_k$ 是固有的，紧支撑上同调等于通常上同调，故对态射
$\pi : \mathbf{P}^1_k \to \mathbf{P}^1_k$，右端等于
$$
\text{Tr}(\pi^* | H^0 (\mathbf{P}^1_{\bar k},
\underline{\mathbf{Z}/\ell\mathbf{Z}})) + \text{Tr}(\pi^* | H^2
(\mathbf{P}^1_{\bar k}, \underline{\mathbf{Z}/\ell\mathbf{Z}})).
$$
Galois 模 $H^0 (\mathbf{P}^1_{\bar k},
\underline{\mathbf{Z}/\ell\mathbf{Z}}) = \mathbf{Z}/\ell\mathbf{Z}$ 是平凡的，
因为恒等的拉回仍是恒等。因此，无论 $\pi$ 如何，第一个迹都是 1。
对第二个迹，须计算映射 $\pi : \mathbf{P}^1_{\bar k} \to
\mathbf{P}^1_{\bar k}$ 所诱导的拉回
$\pi^* : H^2(\mathbf{P}^1_{\bar k}, \underline{\mathbf{Z}/\ell\mathbf{Z}}))$
。这是一个很好的练习，答案是乘以 $\pi$ 的次数（证明见“\'Etale 上同调”
引理 \ref{etale-cohomology-lemma-pullback-on-h2-curve}）。换言之，
其行为与熟悉的复上同调情形相同。特别地，若 $\pi$ 是几何 Frobenius，
则得到
$$
\text{Tr}(\pi_X^* | H^2 (\mathbf{P}^1_{\bar k},
\underline{\mathbf{Z}/\ell\mathbf{Z}})) = q
$$
；若 $\pi$ 是算术 Frobenius，则得到
$$
\text{Tr}(\text{frob}_k^* | H^2 (\mathbf{P}^1_{\bar k},
\underline{\mathbf{Z}/\ell\mathbf{Z}})) = q^{-1}.
$$
后一种选择显然不对。

\begin{remark}
\label{trace-remark-compute-degree-lifting}
次数可通过提升到特征 0（以某种显然的意义），再考察复系数情形来计算。
这种方法几乎从不适用，因为对并非射影空间的概形，提升通常是不可能的。
\end{remark}

\noindent
尚须解释为何必须考虑紧支撑上同调。事实上，由 Poincar\'e 对偶可知，
对光滑代数簇而言这并非严格必要，但此时须添入 $q$ 的某些幂。例如，
考虑 $X = \mathbf{A}^1_k$ 且
$\mathcal{F} = \underline{\mathbf{Z}/\ell\mathbf{Z}}$ 的情形。
茎上的作用仍是平凡的，故只须考察上同调上的作用。此时 $\pi_X^*$ 在
$H^0(\mathbf{A}^1_{\bar k}, \underline{\mathbf{Z}/\ell\mathbf{Z}})$
上作为恒等映射，而在
$H^2_c(\mathbf{A}^1_{\bar k}, \underline{\mathbf{Z}/\ell\mathbf{Z}})$
上作为乘以 $q$ 的映射。




\section{迹}
\label{trace-section-traces}

\noindent
现在说明如何取非交换环上模的自同态之迹。固定一个有限环 $\Lambda$，
其基数与 $p$ 互素。典型地，对某个有限群 $G$，$\Lambda$ 是群环
$(\mathbf{Z}/\ell^n\mathbf{Z})[G]$。约定所考虑的所有 $\Lambda$-模
都是左 $\Lambda$-模。

\medskip\noindent
引入如下记号：令 $\Lambda^\natural$ 为 $\Lambda$ 关于由交换子生成的
加法子群之商（即由形如 $ab-ba$、$a,b\in\Lambda$ 的元素生成的子群）。
注意，$\Lambda^\natural$ 并不是环。

\medskip\noindent
例如，模 $(\mathbf{Z}/\ell^n\mathbf{Z})[G]^\natural$ 是类函数模的对偶，
故
$$
(\mathbf{Z}/\ell^n\mathbf{Z})[G]^\natural
=
\bigoplus\nolimits_{\text{各共轭类，所取群为 }G}
\mathbf{Z}/\ell^n\mathbf{Z}.
$$
对自由 $\Lambda$-模，有 $\text{End}_\Lambda(\Lambda^{\oplus m}) =
\text{Mat}_m(\Lambda)$。注意，由于这些模是左模，以矩阵表示自同态给出的是
右作用：若 $a \in \text{End}(\Lambda^{\oplus m})$ 的矩阵为 $A$，且
$v \in \Lambda$，则 $a(v) = v A$。

\begin{definition}
\label{trace-definition-trace}
自同态 $a$ 的{\it 迹}是表示它的矩阵之对角元之和。这定义了加法映射
$\text{Tr} :
\text{End}_\Lambda(\Lambda^{\oplus m}) \to \Lambda^\natural$.
\end{definition}

\begin{exercise}
\label{trace-exercise-trace-is-trace}
给定映射
$$
\Lambda^{\oplus m} \xrightarrow{a} \Lambda^{\oplus n}
\quad\text{与}\quad
\Lambda^{\oplus n} \xrightarrow{b} \Lambda^{\oplus m}
$$
证明 $\text{Tr}(ab) = \text{Tr}(ba)$。
\end{exercise}

\noindent
如下把迹的定义推广到有限射影 $\Lambda$-模 $P$ 及 $P$ 的自同态 $\varphi$。
把 $P$ 写成自由 $\Lambda$-模的直和项；即考虑映射
$P \xrightarrow{a} \Lambda^{\oplus n} \xrightarrow{b} P$，满足
\begin{enumerate}
\item
$\Lambda^{\oplus n} = \Im(a) \oplus \Ker(b)$；并且
\item
$b\circ a = \text{id}_P$.
\end{enumerate}
令 $\text{Tr}(\varphi) = \text{Tr}(a\varphi b)$。利用上一练习，容易验证
该定义是良定的。








\section{为何需要导出范畴？}
\label{trace-section-derived-categories-why}

\noindent
有了上述迹的定义，现在讨论该公式按当前写法所面临的另一个问题。
设 $C$ 是 $k$ 上的光滑射影曲线。一方面，考虑 $C_\etale$ 上茎同构于
${(\mathbf{Z}/\ell^n\mathbf{Z})}^{\oplus m}$ 的有限局部常值层
$\mathcal{F}$；另一方面，在固定选择 $\bar c$ 后，考虑连续表示
$\rho : \pi_1 (C, \bar c) \to
\text{GL}_m(\mathbf{Z}/\ell^n\mathbf{Z}))$。两者之间存在对应。
以 $\mathcal{F}_\rho$ 表示对应于 $\rho$ 的层。此时
$H^2 (C_{\bar k}, \mathcal{F}_\rho)$ 是
$\rho(\pi_1 (C, \bar c))$ 在
${(\mathbf{Z}/\ell^n\mathbf{Z})}^{\oplus m}$ 上作用的余不变量群，并且有
短正合列
$$
0 \longrightarrow \pi_1 (C_{\bar k}, \bar c) \longrightarrow \pi_1 (C, \bar c)
\longrightarrow G_k \longrightarrow 0.
$$
例如，令 $\mathbf{Z} = \mathbf{Z} \sigma$ 通过
$\sigma(x) = (1+\ell) x$ 作用于 $\mathbf{Z}/\ell^2\mathbf{Z}$。
其余不变量为
$(\mathbf{Z}/\ell^2\mathbf{Z})_{\sigma} = \mathbf{Z}/\ell\mathbf{Z}$，
这不是平坦 $\mathbf{Z}/\ell^2\mathbf{Z}$-模。因此，至少在上一节的
意义下，我们无法取 $H^2(C_{\bar k}, \mathcal{F}_\rho)$ 上某个作用的迹。

\medskip\noindent
事实上，我们的目标是考察 $\ell$-进系数的迹公式。但
$\mathbf{Q}_\ell = \mathbf{Z}_\ell[1/\ell]$ 且
$\mathbf{Z}_\ell = \lim \mathbf{Z}/\ell^n\mathbf{Z}$；即使层是平坦
$\mathbf{Z}/\ell^n\mathbf{Z}$-层，各个上同调群也未必平坦，因而无法
计算迹。一种可能的补救办法，是在导出范畴
$D(\mathbf{Z}/\ell^n\mathbf{Z})$ 中考察全导出复形
$R\Gamma(C_{\bar k}, \mathcal{F}_\rho)$，并证明它是完美对象；这意味着
它拟同构于有限自由模组成的有限复形。对这类复形可以定义迹，但为此须先
介绍导出范畴。






\section{导出范畴}
\label{trace-section-derived-categories}

\noindent
先规定记号。设 $\mathcal{A}$ 是阿贝尔范畴，以
$\text{Comp}(\mathcal{A})$ 表示 $\mathcal{A}$ 中复形组成的阿贝尔范畴。
以 $K(\mathcal{A})$ 表示复形关于同伦所得的范畴：其对象是
$\mathcal{A}$ 中的复形，而态射是复形态射的同伦类。它不是阿贝尔范畴。
粗略地说，$D(A)$ 定义为把 $\text{Comp}(\mathcal{A})$ 中的全部拟同构
求逆所得的范畴；等价地，也可从 $K(\mathcal{A})$ 出发。还可只使用
下有界复形，类似地定义
$\text{Comp}^+(\mathcal{A}), K^+(\mathcal{A}), D^+(\mathcal{A})$。
同样地，使用上有界复形定义
$\text{Comp}^-(\mathcal{A}), K^-(\mathcal{A}), D^-(\mathcal{A})$，
使用有界复形定义
$\text{Comp}^b(\mathcal{A}), K^b(\mathcal{A}), D^b(\mathcal{A})$。

\begin{remark}
\label{trace-remarks-derived-categories}
关于导出范畴的几点说明。
\begin{enumerate}
\item
当 $\mathcal{A}$ 相当任意时，会出现一些集合论问题；我们欣然将其略去。
\item
范畴 $K(A)$ 与 $D(A)$ 都带有三角范畴结构。
\item
当 $\mathcal{A}$ 是加法范畴时，也可定义范畴
$\text{Comp}(\mathcal{A})$ 与 $K(\mathcal{A})$。
\end{enumerate}
\end{remark}

\noindent
把复形 $K^\bullet \mapsto H^i(K^\bullet)$ 的上同调函子
$H^i : \text{Comp}(\mathcal{A}) \to \mathcal{A}$ 可延拓为函子
$H^i : K(\mathcal{A}) \to \mathcal{A}$ 与
$H^i : D(\mathcal{A}) \to \mathcal{A}$。

\begin{lemma}
\label{trace-lemma-when-in-bounded}
$D(\mathcal{A})$ 的对象 $E$ 属于 $D^+(\mathcal{A})$，当且仅当对所有
$i \ll 0$ 都有 $H^i(E) =0 $。对 $D^-$ 与 $D^b$ 也有类似陈述。
\end{lemma}

\begin{proof}
提示：使用截断函子。见“导出范畴”引理
\ref{derived-lemma-complex-cohomology-bounded}。
\end{proof}

\begin{lemma}
\label{trace-lemma-derived-categories}
导出范畴中对象之间的态射。
\begin{enumerate}
\item
设 $I^\bullet \in \text{Comp}^+(\mathcal{A})$，并且对每个
$n \in \mathbf{Z}$，$I^n$ 都是内射的。则
$$
\Hom_{D(\mathcal{A})}(K^\bullet, I^\bullet)
=
\Hom_{K(\mathcal{A})}(K^\bullet, I^\bullet).
$$
\item
设 $P^\bullet \in \text{Comp}^-(\mathcal{A})$，并且对每个
$n \in \mathbf{Z}$，$P^n$ 都是射影的。则
$$
\Hom_{D(\mathcal{A})}(P^\bullet, K^\bullet)
=
\Hom_{K(\mathcal{A})}(P^\bullet, K^\bullet).
$$
\item
若 $\mathcal{A}$ 有足够多内射对象，且
$\mathcal{I} \subset \mathcal{A}$ 是内射对象组成的加法子范畴，则
$
D^+(\mathcal{A})\cong K^+(\mathcal{I})
$
（作为三角范畴）。
\item
若 $\mathcal{A}$ 有足够多射影对象，且
$\mathcal{P} \subset \mathcal{A}$ 是射影对象组成的加法子范畴，则
$
D^-(\mathcal{A}) \cong K^-(\mathcal{P}).
$
\end{enumerate}
\end{lemma}

\begin{proof}
略。
\end{proof}

\begin{definition}
\label{trace-definition-derived-functor}
设 $F: \mathcal{A} \to \mathcal{B}$ 是左正合函子，并假设
$\mathcal{A}$ 有足够多内射对象。把{\it $F$ 的全右导出函子}定义为使下图
交换的函子 $RF: D^+(\mathcal{A}) \to D^+(\mathcal{B})$：
$$
\xymatrix{
D^+(\mathcal{A}) \ar[r]^{RF} & D^+(\mathcal{B}) \\
K^+(\mathcal I) \ar[u] \ar[r]^F & K^+(\mathcal{B}). \ar[u]
}
$$
这是可行的，因为由上一引理，左边的竖直箭头可逆。类似地，设
$G: \mathcal{A} \to \mathcal{B}$ 是右正合函子，并假设 $\mathcal{A}$
有足够多射影对象。把{\it $G$ 的全左导出函子}定义为使下图交换的函子
$LG: D^-(\mathcal{A}) \to D^-(\mathcal{B})$：
$$
\xymatrix{
D^-(\mathcal{A}) \ar[r]^{LG} & D^-(\mathcal{B}) \\
K^-(\mathcal{P}) \ar[u] \ar[r]^G & K^-(\mathcal{B}). \ar[u]
}
$$
这也是可行的，因为由上一引理，左边的竖直箭头可逆。
\end{definition}

\begin{remark}
\label{trace-remark-cohomology-of-derived-functor}
在这些情形下，有 $R^iF(K^\bullet) = H^i(RF(K^\bullet))$；其中左端定义为
复形 $F(K^\bullet)$ 的第 $i$ 个同调。
\end{remark}




\section{滤过导出范畴}
\label{trace-section-filtered-derived-category}

\noindent
原来还得把所有内容再做一遍，并构造滤过导出范畴。

\begin{definition}
\label{trace-definition-filtered}
设 $\mathcal{A}$ 是阿贝尔范畴。
\begin{enumerate}
\item 以 $\text{Fil}(\mathcal{A})$ 表示 $\mathcal{A}$ 的滤过对象
$(A, F)$ 组成的范畴，其中 $F$ 是如下形式的滤过：
$$
A \supset \ldots \supset F^n A \supset F^{n+1}A \supset \ldots
\supset 0.
$$
这是一个加法范畴。
\item 以 $\text{Fil}^f(\mathcal{A})$ 表示 $\text{Fil}(\mathcal{A})$ 的
全子范畴，其对象 $(A, F)$ 具有有限滤过。它也是加法范畴。
\item 若对每个 $p$，
$\text{gr}^p(I) = \text{gr}_F^p(I) = F^pI/F^{p+1}I$ 都是
$\mathcal{A}$ 中的内射（相应地，射影）对象，则称对象
$I \in \text{Fil}^f(\mathcal{A})$ 是{\it 滤过内射的}（相应地，
{\it 滤过射影的}）。
\item 复形范畴
$\text{Comp}(\text{Fil}^f(\mathcal{A})) \supset
\text{Comp}^+(\text{Fil}^f(\mathcal{A}))$
及其同伦范畴
$K(\text{Fil}^f(\mathcal{A})) \supset K^+(\text{Fil}^f(\mathcal A))$
均如前定义。
\item 若 $\text{Comp}(\text{Fil}^f(\mathcal{A}))$ 中复形的态射
$\alpha : K^\bullet \to L^\bullet$ 满足
$$
\text{gr}^p(\alpha): \text{gr}^p(K^\bullet) \to \text{gr}^p(L^\bullet)
$$
对每个 $p \in \mathbf{Z}$ 都是拟同构，则称该态射为
{\it 滤过拟同构}。
\item 在 $K(\text{Fil}^f(\mathcal{A}))$（相应地，
$K^+(\text{Fil}^f(\mathcal{A}))$）中把滤过拟同构求逆，由此定义
$DF(\mathcal{A})$（相应地，$DF^+(\mathcal{A})$）。
\end{enumerate}
\end{definition}

\begin{lemma}
\label{trace-lemma-filtered-derived-category}
若 $\mathcal{A}$ 有足够多内射对象，则
$DF^+(\mathcal{A}) \cong K^+(\mathcal{I})$，其中 $\mathcal{I}$ 是
$\text{Fil}^f(\mathcal{A})$ 中由滤过内射对象组成的全加法子范畴。
类似地，若 $\mathcal{A}$ 有足够多射影对象，则
$DF^-(\mathcal{A}) \cong K^-(\mathcal{P})$，其中 $\mathcal P$ 是
$\text{Fil}^f(\mathcal{A})$ 中由滤过射影对象组成的全加法子范畴。
\end{lemma}

\begin{proof}
略。
\end{proof}





\section{滤过导出函子}
\label{trace-section-filtered-derived-functors}

\noindent
接下来还有滤过导出函子。

\begin{definition}
\label{trace-definition-filtered-derived-functors}
设 $T: \mathcal{A} \to \mathcal{B}$ 是左正合函子，并假设
$\mathcal{A}$ 有足够多内射对象。把 $RT: DF^+(\mathcal{A}) \to D
F^+(\mathcal{B})$ 定义为使下图交换的函子：
$$
\xymatrix{
DF^+(\mathcal{A}) \ar[r]^{RT} & DF^+(\mathcal{B}) \\
K^+(\mathcal{I}) \ar[u] \ar[r]^{T \quad} & K^+(\text{Fil}^f(\mathcal{B})).
\ar[u]}
$$
由上一引理，这是良定义的。设 $G: \mathcal{A} \to \mathcal{B}$ 是右正合
函子，并假设 $\mathcal{A}$ 有足够多射影对象。把
$LG: DF^-(\mathcal{A}) \to DF^-(\mathcal{B})$ 定义为使下图交换的函子：
$$
\xymatrix{
DF^-(\mathcal{A}) \ar[r]^{LG} & DF^-(\mathcal{B}) \\
K^-(\mathcal{P}) \ar[u] \ar[r]^{G \quad} & K^-(\text{Fil}^f(\mathcal{B})).
\ar[u]}
$$
同样，由上一引理，这是良定义的。函子 $RT$（相应地，$LG$）称为
$T$（相应地，$G$）的{\it 滤过导出函子}。
\end{definition}

\begin{proposition}
\label{trace-proposition-compare-filtered-graded}
在上述情形下，有
$$
\text{gr}^p \circ RT = RT \circ \text{gr}^p
$$
其中左边的 $RT$ 是滤过导出函子，而右边的同名函子是全导出函子。换言之，
有交换图
$$
\xymatrix{
DF^+(\mathcal{A}) \ar[r]^{RT} \ar[d]_{\text{gr}^p} & DF^+(\mathcal{B})
\ar[d]^{\text{gr}^p}\\
D^+(\mathcal{A}) \ar[r]^{RT} & D^+(\mathcal{B}).}
$$
\end{proposition}

\begin{proof}
略。
\end{proof}

\noindent
给定 $K^\bullet \in DF^+(\mathcal{B})$，得到谱序列
$$
 E_1^{p, q} = H^{p+q}(\text{gr}^p K^\bullet) \Rightarrow H^{p+q}(\text{忘却
滤过}(K^\bullet)).
$$






\section{滤过复形的应用}
\label{trace-section-applications-filtered}

\noindent
设 $\mathcal{A}$ 是有足够多内射对象的阿贝尔范畴，并设
$0 \to L \to M \to N \to 0$ 是 $\mathcal{A}$ 中的短正合列。令
$\widetilde M \in \text{Fil}^f(\mathcal{A})$ 为带有下述滤过的 $M$：
$$
F^1M = L, \ F^nM = M
\text{ 当 }n \leq 0\text{，且 }F^nM = 0\text{ 当 }n \geq 2.
$$
由定义，有
$$
\text{忘却滤过}(\widetilde M) = M, \quad
\text{gr}^0(\widetilde M) = N, \quad
\text{gr}^1(\widetilde M) = L
$$
并且对其余所有 $n \neq 0, 1$，都有 $\text{gr}^n(\widetilde M) = 0$。
设 $T: \mathcal{A} \to \mathcal{B}$ 是左正合函子，并假设
$\mathcal{A}$ 有足够多内射对象。则
$RT(\widetilde M) \in DF^+(\mathcal{B})$ 是满足下式的滤过复形：
$$
\text{gr}^p(RT(\widetilde M))
\stackrel{\text{拟同构}}{=}
\left\{
\begin{matrix}
0 & \text{若} & p \neq 0, 1, \\
RT(N) & \text{若} & p = 0, \\
RT(L) & \text{若} & p = 1.
\end{matrix}
\right.
$$
并且
$\text{忘却滤过}(RT(\widetilde M))\stackrel{\text{拟同构}}{ = } RT(M)$。
对 $RT(\widetilde M)$ 应用谱序列，得到
$$
E_1^{p, q} = R^{p+q}T(\text{gr}^p(\widetilde M)) \Rightarrow
R^{p+q}T(\text{忘却滤过}(\widetilde M)).
$$
展开该谱序列，得到长正合列
$$
\xymatrix{
0 \ar[r] & T(L) \ar[r] & T(M) \ar[r] & T(N) \ar@(rd, ul)[rdllllr] \\
& R^1T(L) \ar[r] & R^1T(M) \ar[r] & \ldots
}
$$
下面将如下使用这一点。设 $X/k$ 是有限型概形，并设 $\mathcal{F}$ 是平坦
可构造 $\mathbf{Z}/\ell^n \mathbf{Z}$-模。我们要证明迹
$$
\text{Tr}( \pi_X^\ast | R\Gamma_c(X_{\bar k}, \mathcal{F})) \in
\mathbf{Z}/\ell^n \mathbf{Z}
$$
对短正合列具有可加性。为此，仅使用
$R\Gamma_c(X_{\bar k}, -) \in D^+(\mathbf{Z}/\ell^n \mathbf{Z})$ 还不够；
我们必须使用滤过导出范畴。







%10.29.09
\section{完美性}
\label{trace-section-perfect}

\noindent
设 $\Lambda$ 是一个（可能非交换的）环，$\text{Mod}_{\Lambda}$ 是左
$\Lambda$-模范畴，$K(\Lambda) = K(\text{Mod}_\Lambda)$ 是其同伦范畴，
而 $D(\Lambda)= D(\text{Mod}_\Lambda)$ 是其导出范畴。

\begin{definition}
\label{trace-definition-perfect}
以 $K_{perf}(\Lambda)$ 表示如下范畴：对象是有限射影 $\Lambda$-模的
有界复形，态射是复形态射关于同伦所得的类。函子
$K_{perf}(\Lambda)\to D(\Lambda)$ 是全忠实的（见“导出范畴”引理
\ref{derived-lemma-morphisms-from-projective-complex}）。以
$D_{perf}(\Lambda)$ 表示它的本质像。若 $D(\Lambda)$ 的对象属于
$D_{perf}(\Lambda)$，则称其为{\it 完美的}。
\end{definition}

\begin{proposition}
\label{trace-proposition-trace-well-defined}
设 $K\in D_{perf}(\Lambda)$ 且 $f\in \text{End}_{D(\Lambda)}(K)$。则迹
$\text{Tr}(f)\in \Lambda^\natural$ 是良定义的。
\end{proposition}

\begin{proof}
本证明中将反复使用“导出范畴”引理
\ref{derived-lemma-morphisms-from-projective-complex}
而不再说明。设 $P^\bullet$ 是有限射影 $\Lambda$-模的有界复形，并设
$\alpha : P^\bullet \to K$ 是 $D(\Lambda)$ 中的同构。则
$\alpha^{-1}\circ f\circ \alpha$ 对应于复形态射
$f^\bullet : P^\bullet \to P^\bullet$，且后者在同伦意义下良定义。令
$$
\text{Tr}(f) = \sum_i (-1)^i \text{Tr}(f^i : P^i \to P^i) \in \Lambda^\natural.
$$
给定 $P^\bullet$ 与 $\alpha$，该值与 $f^\bullet$ 的选择无关。事实上，
任一其他选择均有形式
$\tilde{f}^\bullet = f^\bullet + dh +hd$，其中某个同伦由
$h^i : P^i \to P^{i-1}(i\in \mathbf{Z})$ 给出。但
\begin{eqnarray*}
\text{Tr}(dh) & = & \sum_i (-1)^i \text{Tr}(P^i\xrightarrow{dh} P^i) \\
& = & \sum_i (-1)^i \text{Tr}(P^{i-1}\xrightarrow{hd} P^{i-1}) \\
& = & -\sum_i (-1)^{i-1}\text{Tr}(P^{i-1}\xrightarrow{hd} P^{i-1}) \\
& = & - \text{Tr}(hd)
\end{eqnarray*}
因此 $\sum_i (-1)^i \text{Tr} ((dh+hd)|_{P^i})=0$。此外，该值也与
$(P^\bullet , \alpha)$ 的选择无关。设 $(Q^\bullet, \beta)$ 是另一选择。
复合
$$
Q^\bullet \xrightarrow{\beta} K \xrightarrow{\alpha^{-1}} P^\bullet
\quad\text{与}\quad
P^\bullet \xrightarrow{\alpha} K \xrightarrow{\beta^{-1}} Q^\bullet
$$
分别可由复形态射 $\gamma_1^\bullet$ 与 $\gamma_2^\bullet$ 表示，并且
$\gamma_1^\bullet \circ \gamma_2^\bullet$ 同伦于恒等。因此，复形态射
$\gamma_2^\bullet\circ f^\bullet\circ \gamma_1^\bullet : Q^\bullet\to Q^\bullet$
表示 $D(\Lambda)$ 中的态射 $\beta^{-1}\circ f\circ\beta$。现有
\begin{eqnarray*}
\text{Tr}(\gamma_2^\bullet\circ f^\bullet\circ\gamma_1^\bullet|_{Q^\bullet}) &
= & \text{Tr}(\gamma_1^\bullet \circ\gamma_2^\bullet \circ
f^\bullet|_{P^\bullet})\\
& = & \text{Tr}(f^\bullet|_{P^\bullet})
\end{eqnarray*}
这是因为 $\gamma_1^\bullet \circ \gamma_2^\bullet$ 同伦于恒等，且我们已证明
该值与 $f^\bullet$ 的选择无关。
\end{proof}




\section{滤过与完美复形}
\label{trace-section-filtrations-perfect}

\noindent
现在给出完美复形范畴的滤过版本。若对每个 $p$，
$\text{gr}^p_F (M)$ 都是有限射影的，则称
$\text{Fil}^f(\text{Mod}_\Lambda)$ 的对象 $(M, F)$ 是
{\it 滤过有限射影的}。接着考虑由
$\text{Fil}^f(\text{Mod}_\Lambda)$ 中滤过有限射影对象的有界复形组成的
同伦范畴 $KF_{\text{perf}}(\Lambda)$。有范畴图
$$
\begin{matrix}
KF(\Lambda) & \supset & KF_{\text{perf}}(\Lambda)\\
\downarrow & & \downarrow\\
DF(\Lambda) & \supset & DF_{\text{perf}}(\Lambda)
\end{matrix}
$$
其中右边的竖直函子是全忠实的，而与前面一样，范畴
$DF_{\text{perf}}(\Lambda)$ 是它的本质像。

\begin{lemma}[可加性]
\label{trace-lemma-additivity}
设 $K\in DF_{\text{perf}}(\Lambda)$ 且 $f\in
\text{End}_{DF}(K)$。则
$$
\text{Tr}(f|_K) =
\sum\nolimits_{p\in \mathbf{Z}} \text{Tr}(f|_{\text{gr}^p K}).
$$
\end{lemma}

\begin{proof}
由命题 \ref{trace-proposition-trace-well-defined}，可以假设我们有
$\text{Fil}^f(\text{Mod}_\Lambda)$ 中滤过有限射影对象的一个有界复形
$P^\bullet$，以及 $\text{Comp}(\text{Fil}^f(\text{Mod}_\Lambda))$ 中的映射
$f^\bullet : P^\bullet\to P^\bullet$。于是本引理由下述结果推出；其证明
留给读者。
\end{proof}

\begin{lemma}
\label{trace-lemma-additive-filtered-finite-projective}
设 $P \in \text{Fil}^f(\text{Mod}_\Lambda)$ 是滤过有限射影对象，并设
$f : P \to P$ 是 $\text{Fil}^f(\text{Mod}_\Lambda)$ 中的自同态。则
$$
\text{Tr}(f|_P) =
\sum\nolimits_p \text{Tr}(f|_{\text{gr}^p(P)}).
$$
\end{lemma}

\begin{proof}
略。
\end{proof}







\section{完美对象的刻画}
\label{trace-section-characterizing-perfect}

\noindent
交换情形见“代数续篇”第
\ref{more-algebra-section-pseudo-coherent},
\ref{more-algebra-section-tor} 以及
\ref{more-algebra-section-perfect} 节。

\begin{definition}
\label{trace-definition-finite-tor-dimension}
设 $\Lambda$ 是一个（可能非交换的）环。若存在 $a, b \in \mathbf{Z}$，
使得对任意右 $\Lambda$-模 $N$ 以及所有 $i \not \in [a, b]$，均有
$H^i(N \otimes_{\Lambda}^\mathbf{L} K) = 0$，则称对象
$K\in D(\Lambda)$ 具有{\it 有限 $\text{Tor}$-维数}。
\end{definition}

\noindent
特别地，取 $N = \Lambda$ 即知这蕴含 $K \in D^b(\Lambda)$。

\begin{lemma}
\label{trace-lemma-characterize-perfect}
设 $\Lambda$ 是左 Noether 环，且 $K\in D(\Lambda)$。则 $K$ 完美，当且仅当
下列两个条件成立：
\begin{enumerate}
\item
$K$ 具有有限 $\text{Tor}$-维数；并且
\item
对所有 $i \in \mathbf{Z}$，$H^i(K)$ 都是有限 $\Lambda$-模。
\end{enumerate}
\end{lemma}

\begin{proof}
交换情形的证明见“代数续篇”引理 \ref{more-algebra-lemma-perfect}。
\end{proof}

\noindent
强烈建议读者尝试证明这一点。该证明依赖如下事实：有限左表现环上的有限模
平坦，当且仅当它是射影模。

\begin{remark}
\label{trace-remark-variant}
该引理的一个变式是考虑 Noether 概形 $X$，以及如下复形组成的范畴
$D_{perf}(\mathcal{O}_X)$：这些复形局部拟同构于有限局部自由
$\mathcal{O}_X$-模的有限复形。$D_{perf}(\mathcal{O}_X)$ 的对象 $K$
可由具有凝聚上同调层和有界 tor 维数来刻画。
\end{remark}








\section{良好复形的上同调}
\label{trace-section-cohomology-ctf}

\noindent
下述命题是关于 $D_{ctf}(X, \Lambda)$ 中对象之紧支撑上同调的一个更一般结果
的特殊情形。

\begin{proposition}
\label{trace-proposition-projective-curve-constructible-cohomology}
设 $X$ 是域 $k$ 上的射影曲线，$\Lambda$ 是有限环，且
$K\in D_{ctf}(X, \Lambda)$。则
$R\Gamma(X_{\bar k}, K)\in D_{perf}(\Lambda)$。
\end{proposition}

\begin{proof}[证明概要。]
第一步是证明：
\begin{enumerate}
\item[(1)]
{\it $R\Gamma(X_{\bar k}, K)$ 的上同调是有界的。}
\end{enumerate}
考虑谱序列
$$
H^i(X_{\bar k}, \underline H^j(K))
\Rightarrow
H^{i+j} (R\Gamma(X_{\bar k}, K)).
$$
由于 $K$ 有界且 $\Lambda$ 有限，各层 $\underline H^j(K)$ 都是挠层。此外，
$X_{\bar k}$ 具有有限上同调维数，所以左端只有对有限多个 $i$ 与 $j$ 才
非零。因此右端也只有有限多个非零项。
\begin{enumerate}
\item[(2)]
{\it 上同调群 $H^{i+j} (R\Gamma(X_{\bar k}, K))$ 是有限的。}
\end{enumerate}
由于各层 $\underline H^j(K)$ 可构造，各群
$H^i(X_{\bar k}, \underline H^j(K))$ 均有限（见“\'Etale 上同调”第
\ref{etale-cohomology-section-vanishing-torsion} 节），故再次由该谱序列即得。
\begin{enumerate}
\item[(3)]
{\it $R\Gamma(X_{\bar k}, K)$ 具有有限 $\text{Tor}$-维数。}
\end{enumerate}
设 $N$ 是右 $\Lambda$-模（事实上，由于 $\Lambda$ 有限，只需假设 $N$
有限）。由投影公式（模的变换），
$$
N \otimes^\mathbf{L}_\Lambda R \Gamma(X_{\bar k}, K) = R\Gamma(X_{\bar k},
N \otimes^\mathbf{L}_\Lambda K).
$$
因此
$$
H^i (N \otimes^\mathbf{L}_\Lambda R\Gamma(X_{\bar k}, K)) = H^i(R\Gamma(X_{\bar
k}, N \otimes_{\Lambda}^\mathbf{L} K)).
$$
现在考虑谱序列
$$
H^i (X_{\bar k}, \underline H^j (N \otimes_{\Lambda}^\mathbf{L} K))
\Rightarrow
H^{i+j}(R\Gamma(X_{\bar k}, N \otimes_{\Lambda}^\mathbf{L} K)).
$$
由于 $K$ 具有有限 $\text{Tor}$-维数，当 $j$ 充分小时，
$\underline H^j(N \otimes_{\Lambda}^\mathbf{L} K)$ 一律为零；而当 $i < 0$
时，左端为零。因此正如所断言的，$R\Gamma(X_{\bar k}, K)$ 具有有限
$\text{Tor}$-维数。故由引理 \ref{trace-lemma-characterize-perfect}，它是完美复形。
\end{proof}





\section{Lefschetz 数}
\label{trace-section-lefschetz-numbers}

\noindent
有限 tor 维数可构造复形的总上同调是完美复形；这一事实是上同调具有良好
性质的关键技术原因，也使我们能够严格定义迹公式中出现的迹。

\begin{definition}
\label{trace-definition-global-lefschetz-number}
设 $\Lambda$ 是有限环，$X$ 是有限域 $k$ 上的射影曲线，且
$K \in D_{ctf}(X, \Lambda)$（例如 $K = \underline\Lambda$）。有典范映射
$c_K : \pi_X^{-1}K \to K$，其基变换 $c_K|_{X_{\bar k}}$ 在完美复形
$R\Gamma(X_{\bar k}, K|_{X_{\bar k}})$ 上诱导一个记为 $\pi_X^*$ 的作用。
$K$ 的{\it 全局 Lefschetz 数}是该作用的迹
$\text{Tr}(\pi_X^* |_{R\Gamma(X_{\bar k}, K)})$。它是
$\Lambda^\natural$ 的元素。
\end{definition}

\begin{definition}
\label{trace-definition-local-lefschetz-number}
设 $\Lambda, X, k, K$ 如定义 \ref{trace-definition-global-lefschetz-number}。
由于 $K\in D_{ctf}(X, \Lambda)$，对 $X$ 的任意几何点 $\bar x$，复形
$K_{\bar x}$ 都是完美复形（属于 $D_{perf}(\Lambda)$）。如第
\ref{trace-section-frobenii} 节所见，Frobenius $\pi_X$ 作用于 $K_{\bar x}$。
$K$ 的{\it 局部 Lefschetz 数}是总和
$$
\sum\nolimits_{x\in X(k)} \text{Tr}(\pi_x |_{K_{\overline{x}}})
$$
它同样是 $\Lambda^\natural$ 的元素。
\end{definition}

\noindent
现在终于可以精确表述迹公式。

\begin{theorem}[Lefschetz 迹公式]
\label{trace-theorem-trace}
设 $X$ 是有限域 $k$ 上的射影曲线，$\Lambda$ 是有限环，且
$K \in D_{ctf}(X, \Lambda)$。则 $K$ 的全局与局部 Lefschetz 数相等，即
\begin{equation}
\label{trace-equation-trace-formula}
\text{Tr}(\pi^*_X |_{R\Gamma(X_{\bar k}, K)})
=
\sum\nolimits_{x\in X(k)} \text{Tr}(\pi_X |_{K_{\bar x}})
\end{equation}
是 $\Lambda^\natural$ 中的等式。
\end{theorem}

\begin{proof}
见下文讨论。
\end{proof}

%11.5.09
\noindent
我们将使用迹公式而不加证明。不过，我们会相当详细地介绍
\cite{SGA4.5} 给出的定理证明（其中一些未得到充分说明的事项列在第
\ref{trace-section-list-skipped} 节）。

\medskip\noindent
我们只对曲线陈述了该公式；在某种较弱的意义下，它是下述结果的推论。

\begin{theorem}[Weil]
\label{trace-theorem-weil-trace-formula}
设 $C$ 是代数闭域 $k$ 上的非奇异射影曲线，并设
$\varphi : C \to C$ 是不同于恒等的 $C$ 的 $k$-自同态。令
$V(\varphi) = \Delta_C \cdot \Gamma_\varphi$，其中 $\Delta_C$ 是对角，
$\Gamma_\varphi$ 是 $\varphi$ 的图，而交数在 $C \times C$ 上取得。令
$J = \underline{\Picardfunctor}^0_{C/k}$ 为 $C$ 的雅可比簇，并以
$\varphi^* : J \to J$ 表示 $\varphi$ 通过拉回诱导的作用。则
$$
V(\varphi) = 1 - \text{Tr}_J(\varphi^*) + \deg \varphi.
$$
\end{theorem}

\begin{proof}
数 $V(\varphi)$ 是 $\varphi$ 的不动点数；它等于
$$
V(\varphi) =
\sum\nolimits_{c \in |C| : \varphi(c) = c} m_{\text{Fix}(\varphi)} (c)
$$
其中 $m_{\text{Fix}(\varphi)} (c)$ 是 $c$ 作为 $\varphi$ 不动点的重数，
即局部一致化参数在 $\varphi - \text{id}_C$ 下之像的阶或消没。该定理的
证明可见 \cite{Lang} 与 \cite{Weil}。
\end{proof}

\begin{example}
\label{trace-example-elliptic-curve}
设 $C = E$ 是椭圆曲线，且 $\varphi = [n]$ 是乘以 $n$。则在雅可比簇上，
$\varphi^* = \varphi^t$ 是乘以 $n$，故其迹为 $2n$、次数为 $n^2$。
另一方面，$\varphi$ 的不动点是满足 $n p = p$ 的点 $p \in E$，即
$(n-1)$-挠点；其基数为 $(n-1)^2$。因此该定理在此写成
$$
(n-1)^2 = 1 - 2n + n^2.
$$
\end{example}

\noindent
{\bf 雅可比簇。}
现在不加证明地讨论 Weil 证明中所用的曲线与其雅可比簇之间的对应。设
$C$ 是代数闭域 $k$ 上的非奇异射影曲线，并选取基点 $c_0 \in C(k)$。
以 $A^1(C \times C)$（或 $\Pic(C \times C)$，或
$\text{CaCl}(C \times C)$）表示 $C \times C$ 上余维 1 除子组成的阿贝尔群。
则
$$
A^1(C \times C) = \text{pr}_1^* (A^1(C)) \oplus \text{pr}_2^* (A^1(C)) \oplus R
$$
其中
$$
R = \{ Z \in A^1(C \times C) \ |
\ Z|_{C \times \{c_0\}} \sim_\text{rat} 0
\text{ 且 }
Z|_{\{c_0\} \times C} \sim_\text{rat} 0 \}.
$$
换言之，$R$ 是如下线丛组成的子群：沿任一投影拉回后都成为平凡线丛。
于是有阿贝尔群的典范同构 $R \cong \text{End}(J)$，它把 $R$ 中的除子
$Z$ 送到自同态
$$
\begin{matrix}
J & \to & J \\
\left[ \mathcal{O}_C(D) \right] & \mapsto & (\text{pr}_1 |_Z)_* (\text{pr}_2
|_Z)^* (D).
\end{matrix}
$$
上述对应如下。以 $\sigma$ 表示交换 $C \times C$ 两个因子的自同构。
$$
\begin{matrix}
\hline & \\
\text{End}(J) & R \\
& \\
\hline & \\
\text{复合 }\alpha, \beta &
{\text{pr}_{13}}_* ({\text{pr}_{12}}^*(\alpha) \circ {\text{pr}_{23}}^*(\beta))
\\
& \\
\text{id}_J &
\Delta_C - \{c_0\} \times C - C \times \{c_0\} \\
& \\
\varphi^* &
\Gamma_\varphi - C \times \{\varphi(c_0)\}
- \sum_{\varphi(c) = c_0} \{c\} \times C \\
& \\
{
\begin{matrix}
\text{迹形式} \\
\alpha, \beta \mapsto \text{Tr}(\alpha \beta)
\end{matrix}
}
&
\alpha, \beta \mapsto - \int_{C \times C} \alpha . \sigma^*\beta
\\
& \\
{
\begin{matrix}
\text{Rosati 对合} \\
\alpha \mapsto \alpha^\dagger
\end{matrix}
}
&
\alpha \mapsto \sigma^*\alpha
\\
& \\
{
\begin{matrix}
\text{Rosati 正性} \\
\text{Tr}(\alpha\alpha^\dagger) > 0
\end{matrix}
}
&
{
\begin{matrix}
\text{Hodge 指数定理，作用于 }C \times C \\
- \int_{C \times C} \alpha \sigma^*\alpha > 0.
\end{matrix}
}
\\
& \\
\hline 
\end{matrix}
$$
事实上，由 Kunneth 公式，子群 $R$ 对应于
$H^1(C)\otimes H^1(C)$ 中的 $1, 1$ hodge 类。

\medskip\noindent
{\bf Weil 的证明。}利用上述对应，可以证明迹公式。我们有
\begin{eqnarray*}
V(\varphi) & = & \int_{C \times C} \Gamma_\varphi.\Delta \\
& = & \int_{C \times C} \Gamma_\varphi. \left(\Delta_C - \{c_0\} \times C - C
\times \{c_0\}\right) + \int_{C \times C} \Gamma_\varphi. \left(\{c_0\} \times C
+ C \times \{c_0\}\right).
\end{eqnarray*}
一方面，
$$
\int_{C \times C} \Gamma_\varphi. \left(\{c_0\} \times C + C \times
\{c_0\}\right)
=
1 + \deg \varphi
$$
另一方面，由于 $R$ 是丰沛除子
$\{c_0\} \times C + C \times \{c_0\}$ 的正交，
\begin{eqnarray*}
& &
\int_{C \times C} \Gamma_\varphi. \left(\Delta_C - \{c_0\} \times C - C \times
\{c_0\}\right) \\
& = &
\int_{C \times C} \left(\Gamma_\varphi - C \times \{\varphi(c_0)\} -
\sum_{\varphi(c) = c_0} \{c\} \times C \right). \left(\Delta_C - \{c_0\} \times
C - C \times \{c_0\}\right) \\
& = & - \text{Tr}_J (\varphi^* \circ \text{id}_J).
\end{eqnarray*}
综上，
$$
V(\varphi) = 1 - \text{Tr}_J (\varphi^*) + \deg \varphi
$$
这就是迹公式。

\begin{lemma}
\label{trace-lemma-weil-mod}
考虑定理 \ref{trace-theorem-weil-trace-formula} 的情形，并设 $\ell$ 是在 $k$ 中
可逆的素数。则
$$
\sum\nolimits_{i = 0}^2
(-1)^i
\text{Tr}(\varphi^* |_{H^i (C, \underline{\mathbf{Z}/\ell^n \mathbf{Z}})})
=
V(\varphi) \mod \ell^n.
$$
\end{lemma}

\begin{proof}[证明概要]
先注意，该假设是有意义的，因为对所有 $i$，
$H^i(C, \underline{\mathbf{Z}/\ell^n \mathbf{Z}})$ 都是自由
$\mathbf{Z}/\ell^n \mathbf{Z}$-模。$\varphi^*$ 在第 0 次上同调上的迹为
1。在 $k$ 中选择一个本原 $\ell^n$ 次单位根，便给出同构
$$
H^i(C, \underline{\mathbf{Z}/\ell^n \mathbf{Z}}) \cong H^i(C, \mu_{\ell^n})
$$
且它与几何 Frobenius 的作用相容。另一方面，
$H^1(C, \mu_{\ell^n}) = J[\ell^n]$。因此
\begin{eqnarray*}
\text{Tr}(\varphi^* |_{H^1 (C, \underline{\mathbf{Z}/\ell^n \mathbf{Z}})})) & =
& \text{Tr}_J (\varphi^*) \mod \ell^n \\
& = & \text{Tr}_{\mathbf{Z}/\ell^n \mathbf{Z}} (\varphi^* : J[\ell^n] \to
J[\ell^n]).
\end{eqnarray*}
此外，$H^2(C, \mu_{\ell^n}) = \Pic(C)/\ell^n\Pic(C) \cong
\mathbf{Z}/\ell^n \mathbf{Z}$，而在这里 $\varphi^*$ 是乘以 $\deg \varphi$。
故
$$
\text{Tr} (\varphi^*|_{H^2 (C, \underline{\mathbf{Z}/\ell^n \mathbf{Z}})}) =
\deg \varphi.
$$
于是
$$
\sum_{i = 0}^2 (-1)^i
\text{Tr}(\varphi^* |_{H^i (C, \underline{\mathbf{Z}/\ell^n \mathbf{Z}})}) =
1 - \text{Tr}_J(\varphi^*) + \deg \varphi \mod \ell^n
$$
而本推论由定理 \ref{trace-theorem-weil-trace-formula} 得出。
\end{proof}

\noindent
证明该推论的另一种方法是证明
$$
X \mapsto H^* (X, \mathbf{Q}_\ell) =
\mathbf{Q}_\ell \otimes
\lim_n H^*(X, \mathbf{Z}/\ell^n\mathbf{Z})
$$
在 $k$ 上光滑射影簇上定义了 Weil 上同调理论。于是迹公式
$$
V(\varphi) = \sum_{i = 0}^2 (-1)^i
\text{Tr}(\varphi^* |_{H^i(C, \mathbf{Q}_\ell)})
$$
是这些公理的形式推论（这是一个线性代数练习，其证明与拓扑情形相同）。




%11.10.09
\section{预备知识与命题链}
\label{trace-section-preliminaries}

\noindent
记号：
固定本节所用记号。以 $A$ 表示交换环，以 $\Lambda$ 表示一个（可能
非交换的）环，并给定环映射 $A\to \Lambda$，其像包含于 $\Lambda$ 的中心。
令 $G$ 为有限群，令 $\Gamma$ 为一个{\it $G$ 被 $\mathbf{N}$ 的幺半群
扩张}，意思是存在正合列
$$
1\to G\to \tilde\Gamma\to \mathbf{Z}\to 1
$$
且 $\Gamma$ 由 $\tilde\Gamma$ 中像非负的那些元素组成。最后，令 $P$ 是
$A[\Gamma]$-模，并且作为 $A[G]$-模是有限射影的；令 $M$ 是
$\Lambda[\Gamma]$-模，并且作为 $\Lambda$-模是有限射影的。

\medskip\noindent
我们的目标是计算 $1 \in \mathbf{N}$ 在 $P \otimes_A M$ 关于 $G$ 的
余不变量上作为 $\Lambda$-线性作用的迹，即数
$$
\text{Tr}_{\Lambda}\left(1; \left(P \otimes_A M\right)_G\right) \in
\Lambda^\natural.
$$
元素 $1\in \mathbf{N}$ 将对应于 Frobenius。

\begin{lemma}
\label{trace-lemma-epsilon}
以 $e\in G$ 表示单位元。映射
$$
\begin{matrix}
\Lambda[G] & \longrightarrow & \Lambda^{\natural}\\
\sum \lambda_g\cdot g & \longmapsto & \lambda_e
\end{matrix}
$$
经由 $\Lambda[G]^\natural$ 分解。以
$\varepsilon : \Lambda[G]^\natural\to \Lambda^\natural$ 表示所诱导的映射。
\end{lemma}

\begin{proof}
只需证明该映射消去换位子。有
$$
\left(\sum\lambda_g g\right)\left(\sum\mu_g g\right)-\left(\sum \mu_g
g\right)\left(\sum\lambda_g g\right)
= \sum_g\left(\sum_{g_1g_2=g}
\lambda_{g_1}\mu_{g_2}-\mu_{g_1}\lambda_{g_2}\right)g
$$
$e$ 的系数为
$$
\sum_g\left(\lambda_g\mu_{g^{-1}}-\mu_g\lambda_{g^{-1}}\right) =
\sum_g\left(\lambda_g\mu_{g^{-1}}-\mu_{g^{-1}}\lambda_g\right)
$$
它是换位子之和，故在 $\Lambda^\natural$ 中为零。
\end{proof}

\begin{definition}
\label{trace-definition-trace-G}
设 $f : P\to P$ 是有限射影 $\Lambda[G]$-模 $P$ 的自同态。定义
$$
\text{Tr}_{\Lambda}^G(f; P) := \varepsilon\left(\text{Tr}_{\Lambda[G]}(f;
P)\right)
$$
并称其为{\it $f$ 在 $P$ 上的 $G$-迹}。
\end{definition}

\begin{lemma}
\label{trace-lemma-lambda-trace}
设 $f : P\to P$ 是有限射影 $\Lambda[G]$-模 $P$ 的自同态。则
$$
\text{Tr}_{\Lambda}(f; P) = \# G \cdot \text{Tr}_\Lambda^G(f; P).
$$
\end{lemma}

\begin{proof}
由可加性，归约到 $P = \Lambda[G]$ 的情形。此时，$f$ 由右乘
$\Lambda[G]$ 的某个元素 $\sum\lambda_g\cdot g$ 给出。在基
$(g)_{g \in G}$ 下，$f$ 的矩阵在位置 $(g_1, g_2)$ 的系数是
$\lambda_{g_2^{-1}g_1}$。特别地，所有对角系数都是 $\lambda_e$，这样的
系数共有 $\# G$ 个。
\end{proof}

\begin{lemma}
\label{trace-lemma-A-module-structure}
映射 $A\to \Lambda$ 在 $\Lambda^\natural$ 上定义了 $A$-模结构。
\end{lemma}

\begin{proof}
这是显然的。
\end{proof}

\begin{lemma}
\label{trace-lemma-diagonal-action-projective-module}
设 $P$ 是有限射影 $A[G]$-模，$M$ 是作为 $\Lambda$-模有限射影的
$\Lambda[G]$-模。则对角 $G$-作用所诱导的结构使 $P \otimes_A M$ 成为
有限射影 $\Lambda[G]$-模。
\end{lemma}

\noindent
注意，由于 $M$ 是 $\Lambda$-模，$P \otimes_A M$ 自然也具有这种模结构。
把此结构与对角作用合在一起，明确写成
$$
\left(\sum\lambda_g g\right)\left(p \otimes m\right)
=
\sum g p \otimes \lambda_g g m.
$$

\begin{proof}
对任意 $\Lambda[G]$-模 $N$，有
$$
\Hom_{\Lambda[G]}\left(P \otimes_A M, N\right)= \Hom_{A[G]}\left(P,
\Hom_{\Lambda}(M, N)\right)
$$
其中 $G$ 在 $\Hom_{\Lambda}(M, N)$ 上的作用由 $(g\cdot
\varphi)(m) = g \varphi (g^{-1} m) $ 给出。现在只需注意，由于 $P$ 与 $M$
的射影性，右端是正合函子的复合。
\end{proof}

\begin{lemma}
\label{trace-lemma-multiplicative-trace}
在引理 \ref{trace-lemma-diagonal-action-projective-module} 的假设下，设
$u\in \text{End}_{A[G]}(P)$ 且 $v\in \text{End}_{\Lambda[G]}(M)$。则
$$
\text{Tr}_\Lambda^G \left(u \otimes v; P \otimes_A M\right) = \text{Tr}_A^G(u;
P)\cdot \text{Tr}_\Lambda(v;M).
$$
\end{lemma}

\begin{proof}[证明概要]
归约到 $P=A[G]$ 的情形。此时，$u$ 是右乘 $A[G]$ 的某个元素
$a = \sum a_gg$；把它记作 $u = R_a$。有 $\Lambda[G]$-模的同构
$$
\begin{matrix}
\varphi : & A[G]\otimes_A M & \cong & \left(A[G]\otimes_A M\right)'\\
& g \otimes m & \longmapsto & g \otimes g^{-1}m
\end{matrix}
$$
其中 $\left(A[G]\otimes_A M\right)'$ 的模结构由左 $G$-作用与 $M$ 上的
$\Lambda$-线性结构共同给出。这一结构转移把 $u \otimes v$ 变成
$\sum_ga_gR_g \otimes g^{-1}v$。换言之，
$$
\varphi \circ (u \otimes v) \circ \varphi^{-1}
=
\sum_ga_gR_g \otimes g^{-1}v.
$$
明确展开等式两端后，只须证明
$$
\text{Tr}_\Lambda^G\left(\sum_g a_gR_g \otimes g^{-1}v\right) = a_e\cdot
\text{Tr}_\Lambda(v; M).
$$
为此，只须证明
$$
\text{Tr}_\Lambda^G\left(a_gR_g \otimes g^{-1}v\right) =
\left\{
\begin{matrix}
0 & \text{ 若 } g\neq e\\
a_e\text{Tr}_\Lambda\left(v; M\right) & \text{ 若 }g = e
\end{matrix}
\right.
$$
这又可归约到 $M=\Lambda$。
\end{proof}

\noindent
记号：
考虑幺半群扩张 $1 \to G\to \Gamma\to \mathbf{N} \to 1$，并设
$\gamma\in \Gamma$。记 $Z_\gamma = \{g\in G | g\gamma = \gamma g\}$。

\begin{lemma}
\label{trace-lemma-gamma-z-gamma-trace}
设 $P$ 是 $\Lambda[\Gamma]$-模，并且作为 $\Lambda[G]$-模有限射影；另设
$\gamma \in \Gamma$。则
$$
\text{Tr}_{\Lambda}(\gamma, P) =
\# Z_\gamma \cdot \text{Tr}_\Lambda^{Z_\gamma}\left(\gamma, P\right).
$$
\end{lemma}

\begin{proof}
这直接由引理 \ref{trace-lemma-lambda-trace} 得出。
\end{proof}

\begin{lemma}
\label{trace-lemma-weak-trace}
设 $P$ 是 $A[\Gamma]$-模，并且作为 $A[G]$-模有限射影。设 $M$ 是
$\Lambda[\Gamma]$-模，并且作为 $\Lambda$-模有限射影。则
$$
\text{Tr}_{\Lambda}^{Z_\gamma}(\gamma, P \otimes_A M) =
\text{Tr}_A^{Z_\gamma}(\gamma, P)\cdot \text{Tr}_\Lambda(\gamma, M).
$$
\end{lemma}

\begin{proof}
这直接由引理 \ref{trace-lemma-multiplicative-trace} 得出。
\end{proof}

\begin{lemma}
\label{trace-lemma-trivial-trace}
设 $P$ 是 $\Lambda[\Gamma]$-模，并且作为 $\Lambda[G]$-模有限射影。则余不变量
$P_G = \Lambda \otimes_{\Lambda[G]} P$
构成有限射影 $\Lambda$-模，并带有 $\Gamma/G = \mathbf{N}$ 的作用。此外，
$$
\text{Tr}_\Lambda(1; P_G) =
\sum\nolimits'_{\gamma \mapsto 1} \text{Tr}_\Lambda^{Z_\gamma}(\gamma, P)
$$
其中 $\sum_{\gamma\mapsto 1}'$ 表示对 $\Gamma$ 中的 $G$-共轭类求和。
\end{lemma}

\begin{proof}[证明概要]
先在两端乘以 $\# G$ 后证明此式。
$$
\# G\cdot \text{Tr}_\Lambda(1; P_G)
= \text{Tr}_\Lambda(\sum\nolimits_{\gamma\mapsto 1} \gamma, P_G)
= \text{Tr}_\Lambda(\sum\nolimits_{\gamma\mapsto 1} \gamma, P)
$$
第二个等式由考察交换三角形
$$
\xymatrix{
P^G \ar[rd]_a & & P_G \ar[ll]^c \\
& P \ar[ur]_b
}
$$
得出，其中 $a$ 是典范包含，$b$ 是典范满射，且
$c = \sum_{\gamma \mapsto 1} \gamma$。于是
$$
(\sum\nolimits_{\gamma \mapsto 1} \gamma) |_P = a \circ c \circ b
\quad\text{且}\quad
(\sum\nolimits_{\gamma \mapsto 1} \gamma) |_{P_G} = b \circ a \circ c
$$
故二者具有相同的迹。于是
$$
\# G\cdot \text{Tr}_\Lambda(1; P_G)
=
{\sum_{\gamma\mapsto 1}}'
\frac{\# G}{\# Z_\gamma}\text{Tr}_\Lambda(\gamma, P)
= \# G{\sum_{\gamma\mapsto 1}}' \text{Tr}_\Lambda^{Z_\gamma}(\gamma, P).
$$
最后，通过某个泛性论证归约到 $\Lambda$ 无挠的情形。细节见
\cite{SGA4.5}。
\end{proof}

\begin{remark}
\label{trace-remark-content-trivial-trace}
我们试着说明引理 \ref{trace-lemma-weak-trace} 中公式的内容。假设把 $\Lambda$
看作平凡 $\Gamma$-模时，它具有有限消解
$
0\to P_r\to \ldots \to P_1 \to P_0\to \Lambda\to 0
$
其中各 $\Lambda[\Gamma]$-模 $P_i$ 作为 $\Lambda[G]$-模都是有限射影的。
在这种情形下，
$$
H_*\left(\left(P_\bullet\right)_G\right) =
\text{Tor}_*^{\Lambda[G]}\left(\Lambda, \Lambda\right) = H_*(G, \Lambda)
$$
且
$$
\text{Tr}_\Lambda^{Z_\gamma}\left(\gamma, P_\bullet\right) =\frac{1}{\#
Z_\gamma}\text{Tr}_\Lambda(\gamma, P_\bullet)=\frac{1}{\#
Z_\gamma}\text{Tr}(\gamma, \Lambda) = \frac{1}{\# Z_\gamma}.
$$
因此，引理 \ref{trace-lemma-weak-trace} 给出
$$
\text{Tr}_\Lambda (1 , P_G)
= \text{Tr}\left(1 |_{H_*(G, \Lambda)}\right)
= {\sum_{\gamma\mapsto 1}}'\frac{1}{\# Z_\gamma}.
$$
这可解释为栈 $BG$ 上的点计数。若 $\Lambda = \mathbf{F}_\ell$，其中
$\ell$ 与 $\# G$ 互素，则 $H_*(G, \Lambda)$ 在次数 0 为
$\mathbf{F}_\ell$（在其他次数为 $0$），而公式写成
$$
1 =
\sum\nolimits_{
\frac{\sigma\text{-共轭}}{\text{类}\langle\gamma\rangle}
}
\frac{1}{\# Z_\gamma} \mod \ell.
$$
在某种意义下，这是 $G$ 的一个“平凡”迹公式。稍后将看到，在某些情形下，
(\ref{trace-equation-trace-formula}) 可视为某类群的一个高度非平凡的迹公式，见第
\ref{trace-section-abstract-trace-formula} 节。
\end{remark}





%11.12.09
\section{迹公式的证明}
\label{trace-section-proof-trace-formula}

\begin{theorem}
\label{trace-theorem-trace-formula-again}
设 $k$ 是有限域，$X$ 是 $k$ 上维数至多为 1 的分离有限型概形。设
$\Lambda$ 是基数与 $k$ 的基数互素的有限环，且
$K\in D_{ctf}(X, \Lambda)$。则
\begin{equation}
\label{trace-equation-trace-formula-again}
\text{Tr}(\pi_X^* |_{R\Gamma_c(X_{\bar k}, K)})
=
\sum\nolimits_{x\in X(k)}
\text{Tr}(\pi_x |_{K_{\bar x}})
\end{equation}
是 $\Lambda^{\natural}$ 中的等式。
\end{theorem}

\noindent
关于该陈述的一些说明见评注 \ref{trace-remark-on-trace-formula-again}。记号：简记
$$
T'(X, K) =
\sum\nolimits_{x\in X(k)}
\text{Tr}(\pi_x |_{K_{\bar x}})
$$
为 (\ref{trace-equation-trace-formula-again}) 的右端，并简记
$$
T''(X, K)
=\text{Tr}(\pi_x^* |_{R\Gamma_c(X_{\bar k}, K)})
$$
为其左端。

\begin{proof}[定理 \ref{trace-theorem-trace-formula-again} 的证明]
证明分若干步进行。

\medskip\noindent
第 1 步。{\it 设 $j : \mathcal{U}\hookrightarrow X$ 是开浸入，其补集为
$Y = X - \mathcal{U}$，并设 $i : Y \hookrightarrow X$。则
$T''(X, K) = T''(\mathcal{U}, j^{-1} K)+ T''(Y, i^{-1}K)$，且
$T'(X, K) = T'(\mathcal{U}, j^{-1} K)+ T'(Y, i^{-1}K)$.}

\medskip\noindent
对 $T'$，这是显然的。对 $T''$，使用正合列
$$
0\to j_!j^{-1} K \to K \to i_* i^{-1} K \to 0
$$
在 $K$ 上得到一个滤过。由此产生对象 $\widetilde K \in DF(X, \Lambda)$，
其分次片为 $j_!j^{-1}K$ 与 $i_*i^{-1}K$，二者均属于
$D_{ctf}(X, \Lambda)$。于是，由滤过导出的抽象形式论证（待插入引用），
$R\Gamma_c(X_{\bar k}, K)\in DF_{perf}(\Lambda)$，并且它在
$DF_{perf}(\Lambda)$ 中带有 $\pi_x^*$。由关于滤过复形之迹的讨论
（待插入引用），得到
\begin{eqnarray*}
\text{Tr}(\pi_X^* |_{R\Gamma_c(X_{\bar k}, K)})
& = & \text{Tr}(\pi_X^* |_{R\Gamma_c(X_{\bar k}, j_!j^{-1}K)}) +
\text{Tr}(\pi_X^* |_{R\Gamma_c(X_{\bar k}, i_*i^{-1}K)}) \\
& = & T''(U, i^{-1}K) + T''(Y, i^{-1}K).
\end{eqnarray*}

\noindent
第 2 步。{\it 若 $\dim X\leq 0$，则定理成立。}

\medskip\noindent
事实上，在这种情形下，
$$
R\Gamma_c(X_{\bar k}, K) = R\Gamma(X_{\bar k}, K) = \Gamma(X_{\bar k}, K) =
\bigoplus\nolimits_{\bar x \in X_{\bar k}} K_{\bar x}
$$
且它带有自同态 $\pi_X^*$。由于 $\pi_X : X_{\bar k}\to X_{\bar k}$ 的
不动点恰为位于某个 $k$-有理点 $x \in X(k)$ 上方的点
$\bar x \in X_{\bar k}$，故
$$
\text{Tr}\big(\pi_X^*|_{R\Gamma_c(X_{\bar k}, K)}\big) =
\sum\nolimits_{x \in X(k)} \text{Tr}(\pi_x|_{K_{\bar x}}).
$$
正如所需。

\medskip\noindent
第 3 步。{\it 只须在下列情形证明等式
$T'(\mathcal{U}, \mathcal{F}) = T''(\mathcal{U}, \mathcal{F})$
：
\begin{itemize}
\item $\mathcal{U}$ 是 $k$ 上光滑不可约仿射曲线；
\item $\mathcal{U}(k) = \emptyset$；
\item $K=\mathcal{F}$ 是 $\mathcal{U}$ 上有限局部常值的 $\Lambda$-模层，
其茎为有限射影 $\Lambda$-模；并且
\item $\Lambda$ 被素数 $\ell$ 的某个幂消去，且 $\ell \in k^*$。
\end{itemize}
}

\medskip\noindent
事实上，由第 2 步，可以丢掉任意有限点集。但有理点只有有限多个，故可假设
一个也没有\footnote{此处背景中应当响起邪恶的笑声。}。通过转到不可约分支，
并在必要时丢掉坏点，可以假设 $\mathcal{U}$ 光滑、不可约且仿射。关于
$\mathcal{F}$ 的假设来自展开 $D_{ctf}(X, \Lambda)$ 的定义，关于
$\Lambda$ 的假设则来自考察其准素分解。

\medskip\noindent
在证明的其余部分，考虑情形
$$
\xymatrix{
\mathcal{V} \ar[d]_f \ar[r] & Y \ar[d]^{\bar f} \\
\mathcal{U} \ar[r] & X
}
$$
其中 $\mathcal{U}$ 如上，$f$ 是有限 \'etale Galois 覆盖，$\mathcal{V}$
连通，而水平箭头是射影完备化。记
$G=\text{Aut}(\mathcal{V}|\mathcal{U})$；还可假设
$f^{-1}\mathcal{F} =\underline M$ 是常值的，其中模
$M = \Gamma(\mathcal{V}, f^{-1}\mathcal{F})$ 是在 $\Lambda$ 上有限射影的
$\Lambda[G]$-模。这对应于平凡幺半群扩张
$$
1\to G\to \Gamma = G \times \mathbf{N}\to \mathbf{N}\to 1.
$$
在此背景下，利用上述归约，只须证明
$T''(\mathcal{U}, \mathcal{F}) = 0$。

\medskip\noindent
第 4 步。{\it $G$ 在 $f_*f^{-1}\mathcal{F}$ 上有自然作用，而迹映射
$f_*f^{-1}\mathcal{F}\to \mathcal{F}$ 定义同构}
$$
(f_*f^{-1}\mathcal{F})\otimes_{\Lambda[G]} \Lambda =
(f_*f^{-1}\mathcal{F})_G \cong \mathcal{F}.
$$

\medskip\noindent
为证明这一点，只须在一个几何点处展开所有定义。

\medskip\noindent
第 5 步。{\it 令 $A = \mathbf{Z}/\ell^n \mathbf{Z}$，其中 $n\gg 0$。则
$f_*f^{-1}\mathcal{F} \cong (f_*\underline A)
\otimes_{\underline A} \underline M$，并取对角 $G$-作用。}

\medskip\noindent
第 6 步。{\it 有典范同构
$(f_*\underline A \otimes_{\underline A} \underline M)
\otimes_{\Lambda[G]} \underline \Lambda \cong \mathcal{F}$.
}

\medskip\noindent
事实上，由于关于 $\mathcal{F}$ 的射影性假设，这里是导出张量积。

\medskip\noindent
第 7 步。{\it 有典范同构
$$
R\Gamma_c(\mathcal{U}_{\bar k}, \mathcal{F})
= (R\Gamma_c(\mathcal{U}_{\bar k}, f_*A)\otimes_A^\mathbf{L}
M)\otimes_{\Lambda[G]}^\mathbf{L} \Lambda,
$$
且与 $\pi^*_\mathcal{U}$ 的作用相容。
}

\medskip\noindent
这来自普遍系数定理，即 $R\Gamma_c$ 与 $\otimes^\mathbf{L}$ 交换这一事实，
以及 $\mathcal{F}$ 作为 $\Lambda$-模的平坦性。

\medskip\noindent
有
\begin{eqnarray*}
\text{Tr}(
\pi_\mathcal{U}^* |_{R\Gamma_c(\mathcal{U}_{\bar k}, \mathcal{F})})
& = &
{\sum_{g \in G}}'
\text{Tr}_{\Lambda}^{Z_g}
\left(
(g, \pi_\mathcal{U}^*)
|_{R\Gamma_c(\mathcal{U}_{\bar k}, f_*A)\otimes_A^\mathbf{L} M}
\right) \\
& = &
{\sum_{g\in G}}'
\text{Tr}_A^{Z_g}
(
(g, \pi_\mathcal{U}^*) |_{R\Gamma_c(\mathcal{U}_{\bar k}, f_*A)}
)
\cdot
\text{Tr}_\Lambda(g|_M)
\end{eqnarray*}
其中 $\Gamma$ 通过 $G$ 作用于
$R\Gamma_c(\mathcal{U}_{\bar k}, \mathcal{F})$，而 $(e, 1)$ 通过
$\pi_\mathcal{U}^*$ 作用。因此幺半群扩张由
$\Gamma = G \times \mathbf{N} \to \mathbf{N}$，$\gamma \mapsto 1$ 给出。
第一个等式来自引理 \ref{trace-lemma-trivial-trace}，第二个来自引理
\ref{trace-lemma-weak-trace}。

\medskip\noindent
第 8 步。{\it 只须证明
$\text{Tr}_A^{Z_g}((g, \pi_\mathcal{U}^*)
|_{R\Gamma_c(\mathcal{U}_{\bar k}, f_*A)}) \in A$
在 $\Lambda$ 中的像为零。
}

\medskip\noindent
回忆
\begin{eqnarray*}
\# Z_g \cdot \text{Tr}_A^{Z_g}((g, \pi_\mathcal{U}^*)
|_{R\Gamma_c(\mathcal{U}_{\bar k}, f_*A)})
& = & \text{Tr}_A((g, \pi_\mathcal{U}^*)
|_{R\Gamma_c(\mathcal{U}_{\bar k}, f_*A)})\\
& = &
\text{Tr}_A((g^{-1}\pi_\mathcal{V})^* |_{R\Gamma_c(\mathcal{V}_{\bar k}, A)}).
\end{eqnarray*}
第一个等式是引理 \ref{trace-lemma-gamma-z-gamma-trace}，第二个来自 Leray
谱序列，其中使用了 $f$ 的有限性以及我们只在 $A$ 上取迹这一事实。现在，
由于 $A=\mathbf{Z}/\ell^n\mathbf{Z}$ 且 $n \gg 0$，并且对某个（固定的）
$a$ 有 $\# Z_g = \ell^a$，只须证明下述结果。

\medskip\noindent
第 9 步。{\it 有
$\text{Tr}_A((g^{-1}\pi_\mathcal{V})^* |_{R\Gamma_c(\mathcal{V}, A)}) = 0$
于 $A$。}

\medskip\noindent
再次由可加性，有
\begin{eqnarray*}
& \text{Tr}_A((g^{-1}\pi_\mathcal{V})^* |_{R\Gamma_c(\mathcal{V}_{\bar k} A)})
+
\text{Tr}_A((g^{-1}\pi_\mathcal{V})^*
|_{R\Gamma_c(Y-\mathcal {V})_{\bar k}, A)}) \\
& =
\text{Tr}_A((g^{-1}\pi_Y)^* |_{R\Gamma(Y_{\bar k}, A)})
\end{eqnarray*}
由 Weil 迹公式定理 \ref{trace-theorem-weil-trace-formula}，后一个迹是
$g^{-1}\pi_Y$ 在 $Y$ 上的不动点数。此外，由第 2 步已经证明的 0 维情形，
$$
\text{Tr}_A((g^{-1}\pi_\mathcal{V})^*|_{R\Gamma_c(Y-\mathcal{V})_{\bar k}, A)})
$$
是 $g^{-1}\pi_Y$ 在 $(Y-\mathcal{V})_{\bar k}$ 上的不动点数。因此，
$$
\text{Tr}_A((g^{-1}\pi_\mathcal{V})^* |_{R\Gamma_c(\mathcal{V}_{\bar k}, A)})
$$
是 $g^{-1}\pi_Y$ 在 $\mathcal{V}_{\bar k}$ 上的不动点数。但不存在这样的
点：若 $\bar y\in Y_{\bar k}$ 在 $g^{-1}\pi_Y$ 下不动，则
$\bar f(\bar y) \in X_{\bar k}$ 在 $\pi_X$ 下不动。然而 $\mathcal{U}$
没有 $k$-有理点，所以必有
$\bar f(\bar y)\in (X-\mathcal{U})_{\bar k}$，从而
$\bar y\notin \mathcal{V}_{\bar k}$，矛盾。
证明完毕。
\end{proof}

\begin{remark}
\label{trace-remark-on-trace-formula-again}
关于定理 \ref{trace-theorem-trace-formula-again} 的几点说明。
\begin{enumerate}
\item
该公式在任意维数都成立。由一个 d\'evissage 引理（它使用固有基变换等），
在那样的一般性下可将其归约到当前陈述。
\item
复形 $R\Gamma_c(X_{\bar k}, K)$ 的定义如下：选择开浸入
$j : X \hookrightarrow \bar X$，其中 $\bar X$ 是 $k$ 上维数至多为 1 的
射影概形，并令
$$
R\Gamma_c(X_{\bar k}, K) := R\Gamma(\bar X_{\bar k}, j_!K).
$$
它与 $\bar X$ 的选择无关，这由（在此插入引用）得出。把
$H^i_c(X_{\bar k}, K)$ 定义为 $R\Gamma_c(X_{\bar k}, K)$ 的第 $i$ 个
上同调群。
\end{enumerate}
\end{remark}

\begin{remark}
\label{trace-remark-stronger}
尽管我们所做的都是归约，而且主要是代数，但迹公式定理
\ref{trace-theorem-trace-formula-again} 比 Weil 几何迹公式（定理
\ref{trace-theorem-weil-trace-formula}）强得多，因为它适用于系数系统（层），
而不仅仅适用于常系数。
\end{remark}


%11.17.09
\section{应用}
\label{trace-section-applications}

\noindent
好了，既然已经说明了迹公式的证明，就来尝试用它做些事情。





\section{关于 l-进层}
\label{trace-section-l-adic-sheaves}

\begin{definition}
\label{trace-definition-l-adic-sheaf}
设 $X$ 是 Noether 概形。$X$ 上的一个{\it $\mathbf{Z}_\ell$-层}，或简称
{\it $\ell$-进层} $\mathcal{F}$，是逆系统
$\left\{\mathcal{F}_n\right\}_{n\geq 1}$，其中
\begin{enumerate}
\item
$\mathcal{F}_n$ 是 $X_\etale$ 上的可构造
$\mathbf{Z}/\ell^n\mathbf{Z}$-模；并且
\item
转移映射 $\mathcal{F}_{n+1}\to \mathcal{F}_n$ 诱导同构
$\mathcal{F}_{n+1} \otimes_{\mathbf{Z}/\ell^{n+1}\mathbf{Z}}
\mathbf{Z}/\ell^n\mathbf{Z} \cong \mathcal{F}_n$.
\end{enumerate}
若每个 $\mathcal{F}_n$ 都是局部常值的，则称 $\mathcal{F}$ 是
{\it lisse 的}。这类对象之间的{\it 态射}就是逆系统之间的态射。
\end{definition}

\begin{lemma}
\label{trace-lemma-eventually-constant}
设 $\{\mathcal{G}_n\}_{n\geq 1}$ 是可构造
$\mathbf{Z}/\ell^n\mathbf{Z}$-模组成的逆系统。假设对所有 $k\geq 1$，映射
$$
\mathcal{G}_{n+1}/\ell^k \mathcal{G}_{n+1}\to \mathcal{G}_n /\ell^k
\mathcal{G}_n
$$
对所有 $n\gg 0$ 都是同构（界可能依赖于 $k$）。换言之，假设系统
$\{\mathcal{G}_n/\ell^k\mathcal{G}_n\}_{n\geq 1}$
最终常值，并把相应的层记作 $\mathcal{F}_k$。则系统
$\left\{\mathcal{F}_k\right\}_{k\geq 1}$ 构成 $X$ 上的
$\mathbf{Z}_\ell$-层。
\end{lemma}

\begin{proof}
证明是显然的。
\end{proof}

\begin{lemma}
\label{trace-lemma-l-adic-abelian}
$X$ 上的 $\mathbf{Z}_\ell$-层所成的范畴是阿贝尔范畴。
\end{lemma}

\begin{proof}
设
$\Phi = \left\{\varphi_n\right\}_{n\geq 1} :
\left\{\mathcal{F}_n\right\}
\to
\left\{\mathcal{G}_n\right\}$
是 $\mathbf{Z}_\ell$-层的态射。令
$$
\Coker(\Phi) =
\left\{
\Coker\left(\mathcal{F}_n \xrightarrow{\varphi_n} \mathcal{G}_n\right)
\right\}_{n\geq 1}
$$
而 $\Ker(\Phi)$ 是把引理 \ref{trace-lemma-eventually-constant} 应用于逆系统
$$
\left\{
\bigcap_{m\geq n}
\Im\left(\Ker(\varphi_m) \to \Ker(\varphi_n)\right)
\right\}_{n \geq 1}.
$$
由此定义出阿贝尔范畴这一点留给读者。
\end{proof}

\begin{example}
\label{trace-example-kernel}
设 $X=\Spec(\mathbf{C})$，并令
$\Phi : \mathbf{Z}_\ell\to \mathbf{Z}_\ell$ 为乘以 $\ell$。更精确地，
$$
\Phi = \left\{ \mathbf{Z}/\ell^n\mathbf{Z} \xrightarrow{\ell}
\mathbf{Z}/\ell^n\mathbf{Z}\right\}_{n \geq 1}.
$$
为计算核，考虑逆系统
$$
\ldots\to \mathbf{Z}/\ell\mathbf{Z}\xrightarrow{0}
\mathbf{Z}/\ell\mathbf{Z}\xrightarrow{0}\mathbf{Z}/\ell\mathbf{Z}.
$$
由于这些像总为零，$\Ker(\Phi)$ 作为系统为零。
\end{example}

\begin{remark}
\label{trace-remark-stalk-l-adic-sheaf}
若 $\mathcal{F} = \left\{\mathcal{F}_n\right\}_{n\geq 1}$ 是 $X$ 上的
$\mathbf{Z}_\ell$-层，且 $\bar x$ 是几何点，则
$M_n = \left\{\mathcal{F}_{n, \bar x}\right\}$ 是有限
$\mathbf{Z}/\ell^n\mathbf{Z}$-模组成的逆系统，使得
$M_{n+1}\to M_n$ 是满射，且 $M_n = M_{n+1}/\ell^n M_{n+1}$。由此
$$
M = \lim_n M_n = \lim \mathcal{F}_{n, \bar x}
$$
是有限 $\mathbf{Z}_\ell$-模。这由“代数”引理
\ref{algebra-lemma-limit-complete}、
\ref{algebra-lemma-finite-over-complete-ring} 以及
$M/\ell M = M_1$ 是有限 $\mathbf{F}_\ell$-模这一事实得出。因此，对某些
$r, t\geq 0$、$e_i\geq 1$，有
$M\cong \mathbf{Z}_\ell^{\oplus r} \oplus \oplus_{i = 1}^t
\mathbf{Z}_\ell/\ell^{e_i}\mathbf{Z}_\ell$。模
$M = \mathcal{F}_{\bar x}$ 称为 $\mathcal{F}$ 在 $\bar x$ 处的{\it 茎}。
\end{remark}

\begin{definition}
\label{trace-definition-torsion-l-adic-sheaf}
若对某个 $n$，$\ell^n : \mathcal{F} \to \mathcal{F}$ 是零映射，则称
$\mathbf{Z}_\ell$-层 $\mathcal{F}$ 是{\it 挠的}。$X$ 上
$\mathbf{Q}_\ell$-层所成的阿贝尔范畴，是 $\mathbf{Z}_\ell$-层的阿贝尔
范畴关于挠层所成 Serre 子范畴的商。换言之，其对象是 $X$ 上的
$\mathbf{Z}_\ell$-层；若 $\mathcal{F}, \mathcal{G}$ 是两个这样的对象，则
$$
\Hom_{\mathbf{Q}_\ell} \left(\mathcal{F}, \mathcal{G} \right) =
\Hom_{\mathbf{Z}_\ell} \left(\mathcal{F}, \mathcal{G}\right)
\otimes_{\mathbf{Z}_\ell} \mathbf{Q}_\ell.
$$
以 $\mathcal{F} \mapsto \mathcal{F} \otimes \mathbf{Q}_\ell$ 表示商函子
（它是包含函子的右伴随）。若 $\mathcal{F} =
\mathcal{F}' \otimes \mathbf{Q}_\ell$，其中 $\mathcal{F}'$ 是
$\mathbf{Z}_\ell$-层，且 $\bar x$ 是几何点，则 $\mathcal{F}$ 在
$\bar x$ 处的{\it 茎}为 $\mathcal{F}_{\bar x} =
\mathcal{F}'_{\bar x} \otimes \mathbf{Q}_\ell$。
\end{definition}

\begin{remark}
\label{trace-remark-torsion-stalks}
由于 $\mathbf{Z}_\ell$-层只在 Noether 概形上定义，它是挠的，当且仅当其
各茎都是挠的。
\end{remark}

\begin{definition}
\label{trace-definition-cohomology-l-adic}
若 $X$ 是代数闭域 $k$ 上的分离有限型概形，且
$\mathcal{F} = \left\{\mathcal{F}_n\right\}_{n\geq 1}$ 是 $X$ 上的
$\mathbf{Z}_\ell$-层，则定义
$$
H^i(X, \mathcal{F}) := \lim_n H^i(X, \mathcal{F}_n)
\quad\text{且}\quad
H_c^i(X, \mathcal{F}) := \lim_n H_c^i(X, \mathcal{F}_n).
$$
若对某个 $\mathbf{Z}_\ell$-层 $\mathcal{F}'$，有
$\mathcal{F} = \mathcal{F}'\otimes \mathbf{Q}_\ell$，则令
$$
H_c^i(X , \mathcal{F}) := H_c^i(X,
\mathcal{F}')\otimes_{\mathbf{Z}_\ell}\mathbf{Q}_\ell.
$$
称它们为 $X$ 取系数于 $\mathcal{F}$ 的{\it $\ell$-进上同调}。
\end{definition}





\section{L-函数}
\label{trace-section-L-function}

\begin{definition}
\label{trace-definition-L-function-finite-ring}
设 $X$ 是有限域 $k$ 上的有限型概形。设 $\Lambda$ 是阶与 $k$ 的特征互素的
有限环，且 $\mathcal{F}$ 是 $X_\etale$ 上可构造平坦的 $\Lambda$-模。令
$$
L(X, \mathcal{F}) :=
\prod\nolimits_{x\in |X|}
\det(1 - \pi_x^*T^{\deg x} |_{\mathcal{F}_{\bar x}})^{-1} \in \Lambda [[ T ]]
$$
其中 $|X|$ 是 $X$ 的闭点集，$\deg x = [\kappa(x): k]$，而 $\bar x$ 是位于
$x$ 上方的几何点。显然，这一定义可推广到以
$K \in D_{ctf}(X, \Lambda)$ 代替 $\mathcal{F}$ 的情形。称它为
{\it $\mathcal{F}$ 的 $L$-函数}。
\end{definition}

\begin{remark}
\label{trace-remark-T}
直观地，应把 $T$ 看成 $T = t^f$，其中 $p^f = \# k$。于是这些定义与基域
的大小无关。
\end{remark}

\begin{definition}
\label{trace-definition-L-function-l-adic}
现在假设 $\mathcal{F}$ 是 $X$ 上的 $\mathbf{Q}_\ell$-层。在此情形定义
$$
L(X, \mathcal{F}) :=
\prod\nolimits_{x \in |X|}
\det(1 - \pi_x^*T^{\deg x} |_{\mathcal{F}_{\bar x}})^{-1}
\in \mathbf{Q}_\ell[[T]].
$$
注意，该乘积收敛，因为每个给定次数的点只有有限多个。称它为
{\it $\mathcal{F}$ 的 $L$-函数}。
\end{definition}




\section{上同调解释}
\label{trace-section-L-cohomological}

\noindent
Grothendieck 如下解释 $L$-函数。

\begin{theorem}[有限系数]
\label{trace-theorem-A}
设 $X$ 是有限域 $k$ 上的有限型概形。设 $\Lambda$ 是阶与 $k$ 的特征互素的
有限环，且 $\mathcal{F}$ 是 $X_\etale$ 上可构造平坦的 $\Lambda$-模。则
$$
L(X, \mathcal{F}) =
\det(1 - \pi_X^*\ T |_{R\Gamma_c(X_{\bar k}, \mathcal{F})})^{-1}
\in \Lambda[[T]].
$$
\end{theorem}

\begin{proof}
略。
\end{proof}

\noindent
到目前为止，我们甚至还不知道每个上同调群
$H^i_c(X_{\bar k}, \mathcal{F})$ 是否自由。

\begin{theorem}[进层]
\label{trace-theorem-B}
设 $X$ 是有限域 $k$ 上的有限型概形，且 $\mathcal{F}$ 是 $X$ 上的
$\mathbf{Q}_\ell$-层。则
$$
L(X, \mathcal{F}) =
\prod\nolimits_i
\det(1 - \pi_X^*T |_{H_c^i(X_{\bar k} , \mathcal{F})})^{(-1)^{i + 1}}
\in \mathbf{Q}_\ell[[T]].
$$
\end{theorem}

\begin{proof}
下文将概述其证明。
\end{proof}

\begin{remark}
\label{trace-remark-which-is-harder}
由于我们只发展了一些迹理论，而没有发展行列式理论，定理
\ref{trace-theorem-A} 比定理 \ref{trace-theorem-B} 更难证明。我们只证明后者；前者见
\cite{SGA4.5}。还要注意，对 $\mathbf{Z}_\ell$ 系数不存在该定理的更一般
版本，因为没有 $\ell$-挠。
\end{remark}

\noindent
把定理 \ref{trace-theorem-B} 的证明归约为迹公式。由于 $\mathbf{Q}_\ell$ 的特征
为 0，只需证明取对数导数后的等式。更精确地，对两端应用
$T\frac{d}{dT} \log$。一方面，有
\begin{align*}
T \frac{d}{dT} \log L(X, \mathcal{F})
& =
T\frac{d}{dT} \log
\prod_{x \in |X|} \det(1 - \pi_x^* T^{\deg x} |_{\mathcal{F}_{\bar x}})^{-1} \\
& =
\sum_{x \in |X|} T \frac{d}{dT} \log
( \det(1 - \pi_x^* T^{\deg x} |_{\mathcal{F}_{\bar x}})^{-1}) \\
& =
\sum_{x \in |X|} \deg x
\sum_{n \geq 1} \text{Tr}((\pi_x^n)^* |_{\mathcal{F}_{\bar x}}) T^{n\deg x}
\end{align*}
最后一个等式来自公式
$$
T\frac{d}{dT}\log\left(\det\left(1-fT|_M\right)^{-1}\right) = \sum_{n\geq 1}
\text{Tr}(f^n|_M)T^n
$$
它对任意交换环 $\Lambda$ 以及有限射影 $\Lambda$-模 $M$ 的任意自同态 $f$
都成立。另一方面，有
\begin{align*}
& T\frac{d}{dT} \log\left(
\prod\nolimits_i
\det(1-\pi_X^*T |_{H_c^i\left(X_{\bar k} , \mathcal{F}\right)})^{(-1)^{i+1}}
\right) \\
& =
\sum\nolimits_i (-1)^i \sum\nolimits_{n \geq 1}
\text{Tr}\left((\pi_X^n)^* |_{H_c^i(X_{\bar k},\mathcal{F})}\right) T^n
\end{align*}
这里再次使用了同一公式。现在比较 $T$ 的各次幂并使用 M\"obius 反演公式，
可见定理 \ref{trace-theorem-B} 是下述等式的推论：
$$
\sum_{d | n} d \sum_{x \in |X| \atop \deg x = d}
\text{Tr}((\pi_X^{n/d})^* |_{\mathcal{F}_{\bar x}}) =
\sum_i (-1)^i \text{Tr}((\pi^n_X)^* |_{H^i_c(X_{\bar k}, \mathcal{F})}).
$$
以 $k_n$ 表示 $k$ 的 $n$ 次扩张，并令
$X_n = X \times_{\Spec k} \Spec(k_n)$、
$_n\mathcal{F} = \mathcal{F}|_{X_n}$，则这归结为
$$
\sum_{x \in X_n(k_n)} \text{Tr}(\pi_X^* |_{_n\mathcal{F}_{\bar x}})
=
\sum_i (-1)^i \text{Tr}((\pi^n_X)^* |_{H^i_c({(X_n)}_{\bar k}, _n\mathcal{F})})
$$
而这由定理 \ref{trace-theorem-C} 得出。

%11.19.09


\begin{theorem}
\label{trace-theorem-D}
设 $X/k$ 如上，$\Lambda$ 是满足 $\#\Lambda \in k^*$ 的有限环，且
$K\in D_{ctf}(X, \Lambda)$。则
$R\Gamma_c(X_{\bar k}, K)\in D_{perf}(\Lambda)$，并且
$$
\sum_{x\in X(k)}\text{Tr}\left(\pi_x |_{K_{\bar x}}\right) =
\text{Tr}\left(\pi_X^* |_{R\Gamma_c(X_{\bar k}, K )}\right).
$$
\end{theorem}

\begin{proof}
注意，当 $\dim X \leq 1$ 时，我们已经证明了这一点（引用）。一般情形容易
由该情形与固有基变换定理推出。
\end{proof}

\begin{theorem}
\label{trace-theorem-C}
设 $X$ 是有限域 $k$ 上的分离有限型概形，且 $\mathcal{F}$ 是 $X$ 上的
$\mathbf{Q}_\ell$-层。则对所有 $i$，
$\dim_{\mathbf{Q}_\ell}H_c^i(X_{\bar k}, \mathcal{F})$ 都是有限的，而且
只有在 $0\leq i \leq 2 \dim X$ 时才可能非零。此外，有
$$
\sum_{x\in X(k)} \text{Tr}\left(\pi_x |_{\mathcal{F}_{\bar x}}\right) =
\sum_i (-1)^i\text{Tr}\left(\pi_X^* |_{H_c^i(X_{\bar k}, \mathcal{F})}\right).
$$
\end{theorem}

\begin{proof}
说明如何从定理 \ref{trace-theorem-D} 推出此结论。先用若干 \'etale 上同调论证，
把证明归约为随后要证明的一个代数陈述。

\medskip\noindent
设 $\mathcal{F}$ 如定理所述。可把 $\mathcal{F}$ 写成
$\mathcal{F}'\otimes \mathbf{Q}_\ell$，其中
$\mathcal{F}' = \left\{\mathcal{F}'_n\right\}$ 是无挠
$\mathbf{Z}_\ell$-层；也就是说，在 $\mathbf{Z}_\ell$-层范畴中，
$\ell : \mathcal{F}'\to \mathcal{F}'$ 的核为零。于是每个
$\mathcal{F}_n'$ 都是 $X_\etale$ 上平坦可构造的
$\mathbf{Z}/\ell^n\mathbf{Z}$-模，故
$\mathcal{F}'_n \in D_{ctf}(X, \mathbf{Z}/\ell^n\mathbf{Z})$，并且
$\mathcal{F}_{n+1}'
\otimes^{\mathbf{L}}_{\mathbf{Z}/\ell^{n+1}\mathbf{Z}}
\mathbf{Z}/\ell^n\mathbf{Z} = \mathcal{F}_n'$.
注意，由于 $\mathcal{F}'_n$ 平坦，最后一个等式对通常的（非导出）张量积
也成立（这是同一个等式）。因此，
\begin{enumerate}
\item
复形 $K_n = R\Gamma_c\left(X_{\bar k}, \mathcal{F}_n'\right)$ 是完美的，
且在 $D(\mathbf{Z}/\ell^n\mathbf{Z})$ 中带有自同态
$\pi_n : K_n\to K_n$；
\item
有识别
$$
K_{n+1}
\otimes^{\mathbf{L}}_{\mathbf{Z}/\ell^{n+1}\mathbf{Z}}
\mathbf{Z}/\ell^n\mathbf{Z}
=
K_n
$$
位于 $D_{perf}(\mathbf{Z}/\ell^n\mathbf{Z})$ 中，并与自同态
$\pi_{n+1}$ 和 $\pi_n$ 相容（见 \cite[Rapport 4.12]{SGA4.5}）；
\item
等式 $\text{Tr}\left(\pi_X^* |_{K_n}\right) =
\sum_{x\in X(k)} \text{Tr}\left(\pi_x |_{(\mathcal{F}'_n)_{\bar x}}\right)$
成立；并且
\item
对每个 $x\in X(k)$，各元素
$\text{Tr}(\pi_x |_{\mathcal{F}'_{n, \bar x}}) \in \mathbf{Z}/\ell^n\mathbf{Z}$
构成 $\mathbf{Z}_\ell$ 的一个元素，它等于
$\text{Tr}(\pi_x |_{\mathcal{F}_{\bar x}}) \in \mathbf{Q}_\ell$。
\end{enumerate}
所以只需证明下述代数引理。
\end{proof}

\begin{lemma}
\label{trace-lemma-piece-together}
假设给定
$K_n\in D_{perf}(\mathbf{Z}/\ell^n\mathbf{Z})$, $\pi_n : K_n\to K_n$
以及同构
$\varphi_n :
K_{n+1} \otimes^\mathbf{L}_{\mathbf{Z}/\ell^{n+1}\mathbf{Z}}
\mathbf{Z}/\ell^n\mathbf{Z}
\to K_n$
与 $\pi_{n+1}$ 和 $\pi_n$ 相容。则
\begin{enumerate}
\item
各元素 $t_n = \text{Tr}(\pi_n |_{K_n})\in \mathbf{Z}/\ell^n\mathbf{Z}$
构成 $\mathbf{Z}_\ell$ 的一个元素 $t_\infty = \{t_n\}$；
\item
$\mathbf{Z}_\ell$-模 $H_\infty^i = \lim_n H^i(k_n)$ 是有限的，并且只有
有限多个 $i$ 对应非零模；并且
\item
各算子 $H^i(\pi_n): H^i(K_n)\to H^i(K_n)$ 彼此相容，并定义满足下式的
$\pi_\infty^i : H_\infty^i\to H_\infty^i$：
$$
\sum (-1)^i \text{Tr}(
\pi_\infty^i |_{H_\infty^i \otimes_{\mathbf{Z}_\ell}\mathbf{Q}_\ell}) =
t_\infty.
$$
\end{enumerate}
\end{lemma}

\begin{proof}
由于 $\mathbf{Z}/\ell^n\mathbf{Z}$ 是局部环且 $K_n$ 完美，每个 $K_n$ 都可由
有限自由 $\mathbf{Z}/\ell^n \mathbf{Z}$-模的有限复形 $K_n^\bullet$ 表示，
并且映射 $K_n^p \to K_n^{p+1}$ 的像包含于 $\ell K_n^{p+1}$。这样的复形在
同构意义下是唯一的。此外，$\pi_n$ 可由复形态射
$\pi_n^\bullet : K_n^\bullet\to K_n^\bullet$ 表示（它在同伦意义下唯一）。
同理，同构
$\varphi_n : K_{n+1} \otimes_{\mathbf{Z}/\ell^{n+1}\mathbf{Z}}^{\mathbf{L}}
\mathbf{Z}/\ell^n\mathbf{Z}\to K_n$ 由复形态射表示：
$$
\varphi_n^\bullet :
K_{n+1}^\bullet
\otimes_{\mathbf{Z}/\ell^{n+1}\mathbf{Z}}
\mathbf{Z}/\ell^n\mathbf{Z} \to K_n^\bullet.
$$
事实上，$\varphi_n^\bullet$ 是复形同构，因而可见
\begin{itemize}
\item
存在与 $n$ 无关的 $a, b\in \mathbf{Z}$，使得对所有
$i\notin[a, b]$，都有 $K_n^i = 0$；并且
\item
$K_n^i$ 的秩与 $n$ 无关。
\end{itemize}
因此，模 $K_\infty^i = \lim_n \{K_n^i, \varphi_n^i\}$ 是有限自由
$\mathbf{Z}_\ell$-模，而 $K_\infty^\bullet$ 是有限自由
$\mathbf{Z}_\ell$-模的有限复形。对非零项数归纳，可以证明
$H^i\left(K_\infty^\bullet\right) = \lim_n H^i\left(K_n^\bullet\right)$
（对无界复形这并不成立）。由此得出
$H_\infty^i = H^i\left(K_\infty^\bullet\right)$ 是有限
$\mathbf{Z}_\ell$-模。这证明了 {\it ii}。为证明引理的其余部分，必须处理
下列诸图可能不交换的问题：
$$
\xymatrix{
{K_{n+1}^\bullet} \ar[d]_{\pi_{n+1}^\bullet} \ar[r]^{\varphi_n^\bullet} &
{K_n^\bullet} \ar[d]^{\pi_n^\bullet} \\
{K_{n+1}^\bullet} \ar[r]_{\varphi_n^\bullet} & {K_n^\bullet.}
}
$$
不过，该图在导出范畴中确实交换，因而在同伦意义下交换。对 $n\geq 2$
归纳地把 $\pi_n^\bullet$ 换成与之同伦的复形态射，使这些图交换。具体地，
若 $h^i : K_{n+1}^i \to K_n^{i-1}$ 是一个同伦，即
$$
\pi_n^\bullet \circ \varphi_n^\bullet -
\varphi_n^\bullet \circ \pi_{n + 1}^\bullet = dh + hd,
$$
则选择提升 $h^i$ 的 $\tilde h^i : K_{n+1}^i\to K_{n+1}^{i-1}$。这是可行的，
因为 $K_{n+1}^i$ 自由且 $K_{n+1}^{i-1}\to K_n^{i-1}$ 为满射。然后把
$\pi_n^\bullet$ 换成由下式定义的 $\tilde\pi_n^\bullet$：
$$
\tilde\pi_{n+1}^\bullet = \pi_{n+1}^\bullet + d\tilde h+\tilde hd.
$$
作此 $\{\pi_n^\bullet\}$ 选择后，上述诸图交换，而这些映射拼合成
$K_\infty^\bullet$ 的自同态
$\pi_\infty^\bullet = \lim_n\pi_n^\bullet$。于是 {\it i} 是显然的：
各元素 $t_n = \sum(-1)^i \text{Tr}\left(\pi_n^i |_{K_n^i}\right)$ 拼合为
$\mathbf{Z}_\ell$ 的一个元素 $t_\infty$。此外，
\begin{align*}
t_\infty
& =
\sum (-1)^i \text{Tr}_{\mathbf{Z}_\ell}(\pi_\infty^i |_{K_\infty^i}) \\
& =
\sum (-1)^i \text{Tr}_{\mathbf{Q}_\ell}(
\pi_\infty^i |_{K_\infty^i \otimes_{\mathbf{Z}_\ell}\mathbf{Q}_\ell}) \\
& =
\sum (-1)^i \text{Tr}(
\pi_\infty |_{H^i(K_\infty^\bullet \otimes \mathbf{Q}_\ell)})
\end{align*}
最后一个等式来自如下事实：$\mathbf{Q}_\ell$ 是域，所以复形
$K_\infty^\bullet \otimes \mathbf{Q}_\ell$ 拟同构于其上同调
$H^i(K_\infty^\bullet \otimes \mathbf{Q}_\ell)$。后者又等于
$H^i(K_\infty^\bullet)\otimes_{\mathbf{Z}}\mathbf{Q}_\ell = H_\infty^i \otimes
\mathbf{Q}_\ell$。这就完成了该引理以及定理 \ref{trace-theorem-C} 的证明。
\end{proof}




\section{上文应补充的事项清单}
\label{trace-section-list-skipped}

\noindent
到目前为止，讲课中略去了下列内容的证明：
\begin{enumerate}
\item 曲线及其雅可比簇；
\item 固有基变换定理；
\item 对 $R\Gamma_c$ 的讨论不充分；
\item 更一般地，给定有限型态射 $f : X \to S$，其中源是分离的
$S$-拟射影概形，讨论 \'etale 层上的 $Rf_!$；
\item 对 $\otimes^{\mathbf{L}}$ 的讨论；
\item 关于 $R\Gamma_c$ 为何与 $\otimes^{\mathbf{L}}$ 交换的讨论。
\end{enumerate}






%11.24.09
\section{L-函数的例子}
\label{trace-section-examples-L-functions}

\noindent
利用曲线情形的定理 \ref{trace-theorem-B} 给出若干 $L$-函数的例子。





\section{常值层}
\label{trace-section-L-function-constant-sheaf}

\noindent
设 $k$ 是有限域，$X$ 是 $k$ 上光滑且几何不可约的曲线，并设
$\mathcal{F} = \underline{\mathbf{Q}_\ell}$ 为常值层。若 $\bar x$ 是 $X$
的几何点，则 Galois 模 $\mathcal{F}_{\bar x} = \mathbf{Q}_\ell$ 是平凡的，
故
$$
\det(1-\pi_x^*\ T^{\deg x} |_{\mathcal{F}_{\bar x}})^{-1} =
\frac{1}{1-T^{\deg x}}.
$$
应用定理 \ref{trace-theorem-B}，得到
\begin{align*}
L(X, \mathcal{F})
& =
\prod_{i = 0}^2
\det(1 - \pi_X^*T |_{H_c^i(X_{\bar k}, \mathbf{Q}_\ell)})^{(-1)^{i+1}} \\
& =
\frac{\det(1 - \pi_X^*T |_{H_c^1(X_{\bar k}, \mathbf{Q}_\ell)})}{
\det(1 - \pi_X^*T |_{H_c^0(X_{\bar k}, \mathbf{Q}_\ell)})
\cdot \det(1 - \pi_X^*T |_{H_c^2(X_{\bar k}, \mathbf{Q}_\ell)})}.
\end{align*}
为计算后一个式子，分两种情形。


\medskip\noindent
{\bf 射影情形。}
假设 $X$ 是射影的，则 $H_c^i(X_{\bar k}, \mathbf{Q}_\ell) =
H^i(X_{\bar k}, \mathbf{Q}_\ell)$，并且
$$
H^i(X_{\bar k}, \mathbf{Q}_\ell) =
\left\{
\begin{matrix}
\mathbf{Q}_\ell & \pi_X^* = 1 & \text{若 }i = 0, \\
\mathbf{Q}_\ell^{2g} & \pi_X^* = ? & \text{若 }i = 1, \\
\mathbf{Q}_\ell & \pi_X^* = q & \text{若 }i = 2.
\end{matrix}
\right.
$$
$\pi_X^*$ 在 $H^2$ 上作用的上述识别来自“\'Etale 上同调”引理
\ref{etale-cohomology-lemma-pullback-on-h2-curve}，以及 $\pi_X$ 的次数为
$q = \#(k)$ 这一事实。关于 $\pi_X^*$ 在第 1 次上同调上的作用，我们所知
不多。把它在 $\bar{\mathbf{Q}}_\ell$ 中的特征值记为
$\alpha_1, \ldots, \alpha_{2g}$。综合起来，定理 \ref{trace-theorem-B} 给出等式
$$
\prod\nolimits_{x \in |X|} \frac{1}{1 - T^{\deg x}} =
\frac{\det(1- \pi_X^* T|_{H^1(X_{\bar k}, \mathbf{Q}_\ell)})}{(1-T)(1-qT)} =
\frac{(1 - \alpha_1 T) \ldots (1 - \alpha_{2g}T)}{(1-T)(1-qT)}
$$
由此推出下述结果。

\begin{lemma}
\label{trace-lemma-count-points-projective}
设 $X$ 是有限域 $k$ 上光滑、射影且几何不可约的曲线。则
\begin{enumerate}
\item $L$-函数 $L(X, \mathbf{Q}_\ell)$ 是有理函数；
\item $\pi_X^*$ 在 $H^1(X_{\bar k}, \mathbf{Q}_\ell)$ 上的特征值
$\alpha_1, \ldots, \alpha_{2g}$ 是与 $\ell$ 无关的代数整数；
\item $X$ 在 $k_n$ 上的有理点数（其中 $[k_n : k] = n$）为
$$
\# X(k_n) = 1 - \sum\nolimits_{i = 1}^{2g}\alpha_i^n + q^n,
$$
\item 对每个 $i$，$|\alpha_i| < q$。
\end{enumerate}
\end{lemma}

\begin{proof}
第 (3) 项是把定理 \ref{trace-theorem-C} 应用于 $X \otimes k_n$ 上的
$\mathcal{F} = \underline{\mathbf{Q}_\ell}$。对第 (4) 项，使用下述结果。
\end{proof}

\begin{exercise}
\label{trace-exercise-powers}
设 $\alpha_1, \ldots, \alpha_n \in \mathbf{C}$。对任意包含正实轴、形如
$C_\varepsilon = \{ z \in \mathbf{C} \ | \ |\arg z| < \varepsilon \}$ 的
圆锥扇形（其中 $\varepsilon > 0$），存在整数 $k \geq 1$，使得
$\alpha_1^k, \ldots, \alpha_n^k \in C_\varepsilon$。
\end{exercise}

\noindent
然后证明对所有 $i$，都有 $|\alpha_i| \leq q$。再用关于复数的初等论证
（如素数定理的证明中那样）证明 $|\alpha_i| < q$。事实上，Riemann 假设
断言对所有 $i$ 都有 $|\alpha_i| = \sqrt{q}$。稍后还会回到这一点。

\medskip\noindent
{\bf 仿射情形。}
现在假设 $X$ 是仿射的，比如
$X= \bar X-\left\{x_1, \ldots, x_n\right\}$，其中
$j : X \hookrightarrow \bar X$ 是非奇异射影完备化。于是
$H_c^0(X_{\bar k}, \mathbf{Q}_\ell) = 0$ 且
$H_c^2(X_{\bar k}, \mathbf{Q}_\ell) = H^2(\bar X_{\bar k}, \mathbf{Q}_\ell)$，
故定理 \ref{trace-theorem-B} 写成
$$
L(X, \mathbf{Q}_\ell) = \prod_{x \in |X|}\frac{1}{1 - T^{\deg x}} =
\frac{\det(1-\pi_X^*T |_{H_c^1(X_{\bar k}, \mathbf{Q}_\ell)})}{1 - qT}.
$$
另一方面，前一种情形给出
\begin{eqnarray*}
L(X, \mathbf{Q}_\ell) & = & L(\bar X,
\mathbf{Q}_\ell)\prod_{i = 1}^n\left(1-T^{\deg x_i}\right) \\
& = & \frac{\prod_{i = 1}^n(1-T^{\deg
x_i})\prod_{j = 1}^{2g}(1-\alpha_jT)}{(1-T)(1-qT)}.
\end{eqnarray*}
因此可见
$\dim H_c^1(X_{\bar k}, \mathbf{Q}_\ell) =
2g+\sum_{i = 1}^n \deg(x_i)-1$，而 $\pi_{\bar X}^*$ 在第 1 次上同调上
作用的特征值 $\alpha_1, \ldots, \alpha_{2g}$ 都是单位根。更精确地，每个
$x_i$ 给出全体 $\deg(x_i)$ 次单位根，并删去一次 1。为利用凝聚层直接看出
这一点，考虑 $\bar X$ 上的短正合列
$$
0\to j_!\mathbf{Q}_\ell\to \mathbf{Q}_\ell\to\bigoplus_{i = 1}^n
\mathbf{Q}_{\ell, x_i}\to 0.
$$
上同调长正合列写成
$$
0\to \mathbf{Q}_\ell \to \bigoplus_{i = 1}^n \mathbf{Q}_\ell^{\oplus \deg x_i}
\to H_c^1(X_{\bar k}, \mathbf{Q}_\ell) \to H_c^1(\bar X_{\bar k},
\mathbf{Q}_\ell)\to 0
$$
其中 Frobenius 在 $\bigoplus_{i = 1}^n \mathbf{Q}_\ell^{\oplus
\deg x_i}$ 上的作用是逐项循环置换；并且
$H_c^2(X_{\bar k}, \mathbf{Q}_\ell) = H_c^2(\bar X_{\bar k}, \mathbf{Q}_\ell)$。






\section{Legendre 族}
\label{trace-section-legendre-family}

\noindent
设 $k$ 是奇特征有限域，
$X = \Spec(k[\lambda, \frac{1}{\lambda(\lambda - 1)}])$，并考虑
$\mathbf{P}^2_X$ 上的椭圆曲线族 $f : E \to X$，其仿射方程为
$y^2 = x(x - 1)(x - \lambda)$。令
$\mathcal{F} = Rf_*^1\mathbf{Q}_\ell =
\left\{R^1f_*\mathbf{Z}/\ell^n\mathbf{Z}\right\}_{n\geq 1} \otimes
\mathbf{Q}_\ell$。在此情形下，下列陈述成立：
\begin{itemize}
\item 对每个 $n \geq 1$，层 $R^1f_*(\mathbf{Z}/\ell^n\mathbf{Z})$ 都是
有限局部常值的；事实上，它是在 $\mathbf{Z}/\ell^n\mathbf{Z}$ 上秩为 2
的自由模；
\item 系统 $\{R^1f_*\mathbf{Z}/\ell^n\mathbf{Z}\}_{n\geq 1}$ 是 lisse
$\ell$-进层；并且
\item 对所有 $x\in |X|$，
$\det(1 - \pi_x\ T^{\deg x} |_{\mathcal{F}_{\bar x}}) =
(1 - \alpha_x T^{\deg x})(1 - \beta_x T^{\deg x})$
，其中 $\alpha_x, \beta_x$ 是 $E_x$ 的几何 Frobenius 在
$H^1(E_{\bar x}, \mathbf{Q}_\ell)$ 上作用的特征值。
\end{itemize}
注意，$E_x$ 只在 $\kappa(x)$ 上定义，而不在 $k$ 上定义。这些事实的证明
使用固有基变换定理与光滑态射的局部无圈性。细节见 \cite{SGA4.5}。由此
$$
L(E/X) := L(X, \mathcal{F}) = \prod_{x\in |X|}
\frac{1}{(1-\alpha_xT^{\deg x})(1-\beta_xT^{\deg x })} .
$$
应用定理 \ref{trace-theorem-B}，得到
$$
L(E/X) =
\prod_{i = 0}^2
\det\left(1 - \pi_X^*T |_{H_c^i(X_{\bar k}, \mathcal{F})}\right)^{(-1)^{i+1}},
$$
特别可见这是一个有理函数。此外，不难证明
$H_c^0(X_{\bar k}, \mathcal{F}) = H_c^2(X_{\bar k}, \mathcal{F}) = 0$，故只有
$$
L(E/X) = \det(1 - \pi_X^*T |_{H_c^1(X, \mathcal{F})}).
$$
为明确计算该行列式，考虑固有态射 $f : E \to X$ 关于
$\mathbf{Q}_\ell$ 的 Leray 谱序列，即
$$
H_c^i(X_{\bar k}, R^jf_*\mathbf{Q}_\ell) \Rightarrow H_c^{i+j}(E_{\bar
k}, \mathbf{Q}_\ell)
$$
该谱序列退化。有 $f_*\mathbf{Q}_\ell = \mathbf{Q}_\ell$ 且
$R^1f_*\mathbf{Q}_\ell = \mathcal{F}$。层
$R^2f_*\mathbf{Q}_\ell = \mathbf{Q}_\ell(-1)$ 是 $\mathbf{Q}_\ell$ 的
{\it Tate 扭曲}；也就是说，它是如下 $\mathbf{Q}_\ell$-层：Galois 在
$\bar x$ 处茎上的作用由乘以 $\#\kappa(x)$ 给出。因此，对所有
$n\geq 1$，
\begin{align*}
\# E(k_n)
& =
\sum\nolimits_i (-1)^i
\text{Tr}({\pi_E^n}^* |_{H_c^i(E_{\bar k}, \mathbf{Q}_\ell)}) \\
& =
\sum\nolimits_{i, j} (-1)^{i+j}
\text{Tr}({\pi^n_X}^* |_{H_c^i(X_{\bar k}, R^jf_*\mathbf{Q}_\ell)}) \\
& =
(q^n - 2) +
\text{Tr}({\pi_X^n}^* |_{H_c^1(X_{\bar k}, \mathcal{F})}) + q^n(q^n - 2) \\
& =
q^{2n} - q^n - 2 +
\text{Tr}({\pi_X^n}^* |_{H_c^1(X_{\bar k}, \mathcal{F})})
\end{align*}
其中第一个等式来自定理 \ref{trace-theorem-C}，第二个来自 Leray 谱序列，
第三个则由写出 $\mathbf{Q}_\ell$ 在 $f$ 下的高阶直像而得。也可以写成
$$
\# E(k_n) = \sum_{x \in X(k_n)} \# E_x(k_n)
$$
并对每条曲线使用迹公式。也可直接计数求得 $k_n$-有理点数。零截面贡献
$q^n -2$ 个点（略去 $\lambda = 0, 1$ 的点），因此
$$
\# E(k_n) =
q^n - 2 + \# \{y^2 = x(x - 1)(x - \lambda), \lambda\neq 0, 1\}.
$$
现在有
$$
\begin{matrix}
\# \{y^2 = x(x - 1)(x - \lambda),\ \lambda\neq 0, 1\} \\
\\
\quad =
\# \{y^2 = x(x - 1)(x - \lambda)\text{ 位于 }\mathbf{A}^3\}
- \# \{y^2 = x^2(x - 1)\} - \# \{y^2 = x(x - 1)^2\}\\
\\
\quad = \# \{\lambda = \frac{-y^2}{x(x - 1)} + x,\ x\neq 0, 1\} +
\# \{y^2 = x(x - 1)(x - \lambda), x = 0, 1\} - 2(q^n - \varepsilon_n) \\
\\
\quad = q^n(q^n - 2)+2q^n - 2(q^n - \varepsilon_n)\\
\\
\quad = q^{2n}-2q^n+2\varepsilon_n
\end{matrix}
$$
其中当 $-1$ 在 $k_n$ 中是平方时 $\varepsilon_n = 1$，否则为 0；即
$$
\varepsilon_n = \frac{1}{2}\left(1+\left(\frac{-1}{k_n}\right)\right) =
\frac{1}{2}\left(1+(-1)^{\frac{q^n - 1}{2}}\right).
$$
于是 $ \# E(k_n) = q^{2n} - q^n - 2+ 2\varepsilon_n$。与前一公式比较，得到
$$
\text{Tr}({\pi_X^n}^* |_{H_c^1(X_{\bar k}, \mathcal{F})}) =
2 \varepsilon_n = 1 + (-1)^{\frac{q^n - 1}{2}},
$$
由复数的初等代数可知，若 $-1$ 在 $k_n^*$ 中是平方，则
$\dim H_c^1(X_{\bar k}, \mathcal{F}) = 2$，且特征值为 $1$ 和 $1$。因此，
在这种情形下，
$$
L(E/X) = (1 - T)^2.
$$




\section{指数和}
\label{trace-section-exponential-sums}

\noindent
数论中的一个标准问题是求如下形式的和：
$$
S_{a, b}(p) = \sum_{x\in \mathbf{F}_p - \left\{0, 1\right\}} e^{\frac{2\pi
ix^a(x - 1)^b}{p}}.
$$
在我们的语境中，这可如下解释为上同调和。考虑基概形
$S = \Spec(\mathbf{F}_p[x, \frac{1}{x(x - 1)}])$，以及 $S$ 上由方程
$y^{p - 1} = x^a(x - 1)^b$ 给出的仿射曲线
$f : X \to \mathbf{P}^1-\{0, 1, \infty\}$。这是群为 $\mathbf{F}_p^*$ 的
有限 \'etale Galois 覆盖，并有分解
$$
f_*(\bar{\mathbf{Q}}_\ell^*) =
\bigoplus_{\chi : \mathbf{F}_p^*\to \bar{\mathbf{Q}}_\ell^*} \mathcal{F}_\chi
$$
其中 $\chi$ 遍历 $\mathbf{F}_p^*$ 的特征标，而 $\mathcal{F}_\chi$ 是秩为 1
的 lisse $\mathbf{Q}_\ell$-层，$\mathbf{F}_p^*$ 在其各茎上通过 $\chi$
作用。相应地得到分解
$$
H_c^1(X_{\bar k}, \mathbf{Q}_\ell) = \bigoplus_\chi H^1(\mathbf{P}_{\bar
k}^1-\{0, 1, \infty\}, \mathcal{F}_\chi)
$$
而指数和的上同调解释，是对某个合适的 $\chi$，在
$\mathbf{P}^1 - \{0, 1, \infty\}$ 上把迹公式应用于 $\mathcal{F}_\chi$。
它写成
$$
S_{a, b}(p) =
-\text{Tr}(\pi_X^*
|_{H^1(\mathbf{P}_{\bar k}^1-\{0, 1, \infty\}, \mathcal{F}_\chi)}).
$$
Weil 的一般形式论证表明该和中应当有某种消去。粗略应用 Riemann--Hurwitz
公式，可见
$$
2g_X-2 \approx -2 (p-1) + 3(p-2) \approx p
$$
故 $g_X\approx p/2$，这也表明各 $\chi$-部分很小。



%12.01.09
\section{用基本群表述的迹公式}
\label{trace-section-trace-formula-fundamental-group}

\noindent
在以下各节中，我们将完全用曲线的基本群重新表述迹公式，曲线碰巧是
$\mathbf{P}^1$ 的情形除外。




\section{基本群}
\label{trace-section-fundamental-groups}

\noindent
这些材料在基本群一章中有更详细的讨论，见“基本群”第
\ref{pione-section-introduction} 节。设 $X$ 是连通概形，并设
$\overline{x}\to X$ 是几何点。考虑函子
$$
\begin{matrix}
F_{\overline{x}}: &
\text{有限 \'etale} \atop \text{概形，基为 } X &
\longrightarrow & \text{有限集合} \\
&
Y/X &
\longmapsto &
F_{\overline{x}}(Y) =
\left\{\text{几何点 }\overline y \atop \text{属于 } Y
\text{ 且上方是 }\overline{x}\right\} = Y_{\overline{x}}
\end{matrix}
$$
令
$$
\pi_1(X, \overline{x})
=
Aut(F_{\overline{x}})
=
\text{下述函子的自同构集合：}F_{\overline{x}}
$$
注意，对每个有限 \'etale 态射 $Y \to X$，都有作用
$$
\pi_1(X, \overline{x}) \times F_{\overline{x}}(Y) \to F_{\overline{x}}(Y)
$$

\begin{definition}
\label{trace-definition-open}
形如
$\text{Stab}(\overline y\in F_{\overline{x}}(Y))\subset \pi_1(X, \overline{x})$
的子群称为{\it 开的}。
\end{definition}

\begin{theorem}[Grothendieck]
\label{trace-theorem-fundamental-group}
设 $X$ 是连通概形。
\begin{enumerate}
\item $\pi_1(X, \overline{x})$ 上存在一个拓扑，使开子群组成
$e\in \pi_1(X, \overline x)$ 的开邻域基本系。
\item 对于第 (1) 项的拓扑，群 $\pi_1(X, \overline{x})$ 是预有限群。
\item 函子
$$
\begin{matrix}
\text{有限的概形} \atop \text{且在其上 \'etale：}X & \to &
\text{有限离散连续的} \atop \pi_1(X, \overline{x})\text{-集合}\\
Y / X& \mapsto & F_{\overline{x}}(Y) \text{，带自然作用}
\end{matrix}
$$
是范畴等价。
\end{enumerate}
\end{theorem}

\begin{proof}
见 \cite{SGA1}。
\end{proof}

\begin{proposition}
\label{trace-proposition-integral-normal-fundamental-group}
设 $X$ 是整的正规 Noether 概形。设 $\overline y\to X$ 是位于一般点
$\eta\in X$ 上方的代数几何点。则
$$
\pi_x(X, \overline \eta) = Gal(M/\kappa(\eta))
$$
（$\kappa(\eta)$ 是 $X$ 的函数域），其中
$$
\kappa(\overline \eta)\supset M\supset \kappa(\eta) = k(X)
$$
是满足下述条件的最大子扩张：对每个有限子扩张
$M\supset L\supset \kappa(\eta)$，$X$ 在 $L$ 中的正规化在 $X$ 上有限
\'etale。
\end{proposition}

\begin{proof}
略。
\end{proof}

\noindent
{\bf 基点变换。}对 $X$ 的任意几何点 $\overline{x}_1, \overline{x}_2$，
存在纤维函子的同构
$$
\mathcal{F}_{\overline{x}_1} \cong \mathcal{F}_{\overline{x}_2}
$$
（这是从 $\overline{x}_1$ 到 $\overline{x}_2$ 的一条道路。）由这条道路共轭
给出同构
$$
\pi_1(X, \overline{x}_1) \cong \pi_1(X, \overline{x}_2)
$$
它在内作用意义下良定义。

\medskip\noindent
{\bf 函子性。}对连通概形的任意态射 $X_1\to X_2$ 以及任意
$\overline{x}\in X_1$，有典范映射
$$
\pi_1(X_1, \overline{x}) \to \pi_1(X_2, \overline{x})
$$
（为什么？因为纤维函子……）

\medskip\noindent
{\bf 基域。}设 $X$ 是域 $k$ 上的簇。则得到
$$
\pi_1(X, \overline{x}) \to
\pi_1(Spec(k), \overline{x}) =^{\text{prop}} Gal(k^{\text{sep}}/k)
$$
该映射为满射，当且仅当 $X$ 在 $k$ 上几何连通。因此在几何连通情形，
得到预有限群的短正合列
$$
1 \to \pi_1(X_{\overline{k}}, \overline{x}) \to
\pi_1(X, \overline{x}) \to
Gal(k^{\text{sep}}/k) \to 1
$$
（$\pi_1(X_{\overline{k}}, \overline{x})$：$X$ 的几何基本群；
$\pi_1(X, \overline{x})$：$X$ 的算术基本群。）

\medskip\noindent
{\bf 比较。}若 $X$ 是 $\mathbf{C}$ 上的簇，则
$$
\pi_1(X, \overline{x}) =
\text{预有限完备化：}
\pi_1(X(\mathbf{C})(\text{通常拓扑}), x)
$$
（这里 $x\in X(\mathbf{C})$。）

\medskip\noindent
{\bf Frobenius 元。}设 $X$ 是 $k$ 上的簇，且 $\# k < \infty$。对任意闭点
$x \in X$，令
$$
F_x\in \pi_1(x, \overline{x}) =
\text{Gal}(\kappa(x)^{\text{sep}}/\kappa(x))
$$
为几何 Frobenius。设 $\overline\eta$ 是代数几何一般点。则
$$
\pi_1(X, \overline\eta) \leftarrow^{\cong}
\pi_1(X, \overline{x})
{\text{函子性} \atop \leftarrow} \pi_1(x, \overline{x})
$$

\noindent
易知：
$$
\begin{matrix}
\pi_1(X, \overline \eta) & \to^{\deg} \pi_1(\Spec(k), \overline \eta) * &
= Gal(k^{sep}/k) \\
& & || \\
& & \widehat{\mathbf{Z}}\cdot F_{\Spec(k)} \\
F_x & \mapsto & \deg(x)\cdot F_{\Spec(k)}
\end{matrix}
$$
回忆：$\deg(x) = [\kappa(x):k]$。

\medskip\noindent
{\bf 基本群与 lisse 层。}
设 $X$ 是连通概形，$\overline{x}$ 是几何点。有范畴等价
$$
\begin{matrix}
(\Lambda\text{ 为有限环}) &
\text{有限局部常值的层，类型为 }
\atop \Lambda\text{-模，定义于 }X_\etale & \leftrightarrow &
\text{有限离散的 }\Lambda\text{-模}
\atop \text{，带连续 }\pi_1(X, \overline{x})\text{-作用}\\
\\
(\ell\text{ 为素数}) & \text{lisse }\ell\text{-进} \atop \text{层} &
\leftrightarrow &
\text{有限生成 }\mathbf{Z}_\ell
\text{-模 }M\text{，带连续}
\atop \pi_1(X, \overline{x})\text{-作用，并在下述模上取 }
\ell\text{-进拓扑：}M
\end{matrix}
$$
特别地，lisse $\mathbf{Q}_l$-层对应于连续同态
$$
\pi_1(X, \overline{x}) \to \text{GL}_r(\mathbf{Q}_l), \quad r\geq 0
$$

\noindent
记号：带作用的模 $(M, \rho)$ 对应于层 $\mathcal{F}_\rho$。

\medskip\noindent
{\bf 迹公式。}设 $X$ 是 $k$ 上的簇，且 $\# k < \infty$。
\begin{enumerate}
\item $\Lambda$ 为有限环，$(\# \Lambda, \# k)=1$，
$$
\rho : \pi_1(X, \overline{x})\to \text{GL}_r(\Lambda)
$$
且该同态连续。对每个 $n\geq 1$，有
$$
\sum_{d|n}d\left(
\sum_{x\in |X|, \atop \deg(x)=d}
\text{Tr}(\rho(F_x^{n/d}))\right) =
\text{Tr}\left(
(\pi_x^n)^* |_{R\Gamma_c(X_{\overline{k}}, \mathcal{F}_\rho)}\right)
$$
\item $l\neq char(k)$ 为素数，且 $\rho : \pi_1(X, \overline{x})\to
\text{GL}_r(\mathbf{Q}_l)$。对任意 $n\geq 1$，
$$
\sum_{d|n} d\left(
\sum_{x\in |X| \atop \deg(x)=d}
\text{Tr}
\left(
\rho(F_x^{n/d})
\right)
\right) =
\sum_{i = 0}^{2\dim X}
(-1)^i
\text{Tr}\left(
\pi_X^* |_{H_c^i(X_{\overline{k}}, \mathcal{F}_\rho)}\right)
$$
\end{enumerate}

\noindent
{\bf Weil 猜想。}（Deligne--Weil I，1974）若 $X$ 在 $k$ 上光滑射影，且
$\# k = q$，则 $\pi_X^*$ 在 $H^i(X_{\overline{k}}, \mathbf{Q}_l)$ 上的
特征值是满足 $|\alpha|=q^{1/2}$ 的代数整数 $\alpha$。

\medskip\noindent
{\bf Deligne 猜想。}（几乎已由 Lafforgue 加上 $\ldots$ 完全证明。）设 $X$ 是有限域
$k$ 上的正规簇，
$$
\rho : \pi_1(X, \overline{x}) \longrightarrow \text{GL}_r(\mathbf{Q}_l)
$$
且该同态连续。假设 $\rho$ 不可约，$\det(\rho)$ 具有有限阶。则
\begin{enumerate}
\item 存在数域 $E$，使得对所有 $x\in |X|$（闭点），$\rho(F_x)$ 的特征
多项式的系数都在 $E$ 中。
\item 对任意 $x\in |X|$，$\rho(F_x)$ 的特征值
$\alpha_{x, i}$，$i = 1, \ldots, r$，其复绝对值为 $1$。
（这些是代数数，但不必是代数整数。）
\item 对 $E$ 的每个不整除 $p$ 的有限素位 $\lambda$（也许要先把 $E$
稍微扩大），存在
$$
\rho\lambda : \pi_1(X, \overline{x}) \to \text{GL}_r(E_\lambda)
$$
与 $\rho$ 相容。（即若干 $F_x$ 的特征多项式相同。）
\end{enumerate}

\begin{theorem}[Deligne, Weil II]
\label{trace-theorem-weil-II}
若层 $\mathcal{F}_\rho$ 所对应的 $\rho$ 满足上述猜想的结论，则
$\pi_X^*$ 在 $H_c^i(X_{\overline{k}}, \mathcal{F}_{\rho})$ 上的特征值是
代数数 $\alpha$，其绝对值满足
$$
|\alpha|=q^{w/2}, \text{ 其中 }w\in \mathbf{Z},\ w\leq i
$$
此外，若 $X$ 光滑且射影，则 $w = i$。
\end{theorem}

\begin{proof}
见 \cite{WeilII}。
\end{proof}




%12.03.09
\section{预有限群、上同调与同调}
\label{trace-section-profinite-cohomology}

\noindent
设 $G$ 是预有限群。

\medskip\noindent
{\bf 上同调。}
考虑具有连续 $G$-作用的离散模范畴。
这个范畴有足够多的内射对象，因而可以定义
$$
H^i(G, M) = R^iH^0(G, M) = R^i(M\mapsto M^G)
$$
此外还有导出版本 $RH^0(G, -)$。

\medskip\noindent
{\bf 同调。}
考虑具有连续 $G$-作用的紧阿贝尔群范畴。
这个范畴有足够多的投射对象，因而可以定义
$$
H_i(G, M) = L_iH_0(G, M)=L_i(M\mapsto M_G)
$$
并且也有导出版本。

\medskip\noindent
{\bf 平凡对偶。}
函子 $M\mapsto M^\wedge = \Hom_{cont}(M, S^1)$
互换上述两个范畴，并且
$$
H^i(G, M)^\wedge = H_i(G, M^\wedge)
$$
此外，该函子把挠离散 $G$-模映为具有连续 $G$-作用的预有限模，
反之亦然；若 $M$ 是具有连续 $G$-作用的离散模或预有限模，则
$M^\wedge = \Hom(M, \mathbf{Q}/\mathbf{Z})$。

\medskip\noindent
{\bf 同调补注。}
\begin{enumerate}
\item 对有限环 $\Lambda$ 考虑 $\Lambda$-模时，有等同
$$
H_i(G, M)=Tor_i^{\Lambda[[G]]}(M, \Lambda)
$$
其中 $\Lambda[[G]]$ 是 $G$ 的各有限商的群代数之极限。
\item 若 $G$ 是 $\Gamma$ 的正规子群，且 $\Gamma$ 也是预有限群，则
\begin{itemize}
\item $H^0(G, -)$：离散 $\Gamma$-模$\to$ 离散
$\Gamma/G$-模
\item $H_0(G, -)$：紧 $\Gamma$-模 $\to$ 紧
$\Gamma/G$-模
\end{itemize}
因此预有限群 $\Gamma/G$ 作用于以 $\Gamma$-模为系数的
$G$ 的上同调群。换言之，存在导出函子
$$
RH^0(G, -) :
D^{+}(\text{离散 }\Gamma\text{-模})
\longrightarrow
D^{+}(\text{离散 }\Gamma/G\text{-模})
$$
对 $LH_0(G, -)$ 亦然。
\end{enumerate}








\section{再论曲线的上同调}
\label{trace-section-cohomology-curves-revisited}

\noindent
设 $k$ 为域，$X$ 为 $k$ 上几何连通的光滑曲线。
我们有基本短正合列
$$
1 \to
\pi_1(X_{\overline{k}}, \overline \eta) \to
\pi_1(X, \overline\eta) \to
\text{Gal}(k^{^{sep}}/k) \to 1
$$
若 $\Lambda$ 是满足 $\#\Lambda\in k^*$ 的有限环，$M$ 是有限
$\Lambda$-模，并给定连续同态
$$
\rho : \pi_1(X, \overline\eta) \to \text{Aut}_{\Lambda}(M)
$$
则以 $\mathcal{F}_\rho$ 表示 $X_\etale$ 上相应的层。

\begin{lemma}
\label{trace-lemma-identify-h2c}
存在典范同构
$$
H_c^2(X_{\overline{k}}, \mathcal{F}_\rho)=(M)_{\pi_1(X_{\overline{k}},
\overline\eta)}(-1)
$$
作为 $\text{Gal}(k^{^{sep}}/k)$-模成立。
\end{lemma}

\noindent
这里，下标 ${}_{\pi_1(X_{\overline{k}}, \overline\eta)}$
表示余不变量，而 $(-1)$ 表示 Tate 扭曲，即
$\sigma\in \text{Gal}(k^{^{sep}}/k)$ 通过下式作用：
$$
\chi_{cycl}(\sigma)^{-1}.\sigma\text{ 于右端的作用}
$$
其中
$$
\chi_{cycl} :
\text{Gal}(k^{^{sep}}/k)
\to
\prod\nolimits_{l\neq char(k)}\mathbf{Z}_l^*
$$
是分圆特征标。

\medskip\noindent
改述（Deligne，Weil II，第 338 页）。对 $X$ 上任意有限局部常值层
$\mathcal{F}$，存在最大商 $\mathcal{F}\to
\mathcal{F}''$，使得 $\mathcal{F}''/X_{\overline{k}}$ 为常值层，因而
$$
\mathcal{F}'' = (X\to \Spec(k))^{-1}F''
$$
其中 $F''$ 是 $\Spec(k)$ 上的层，即一个
$\text{Gal}(k^{^{sep}}/k)$-模。于是
$$
H_c^2(X_{\overline{k}}, \mathcal{F})\to H_c^2(X_{\overline{k}},
\mathcal{F}'')\to F''(-1)
$$
是同构。

\begin{proof}[引理 \ref{trace-lemma-identify-h2c} 的证明]
令 $Y\to^{\varphi}X$ 为与
$\Ker(\rho) \subset \pi_1(X, \overline\eta)$ 对应的有限 \'etale Galois 覆叠。于是
$$
\text{Aut}(Y/X)=Ind(\rho)
$$
是 Galois 群。此时 $\varphi^*\mathcal{F}_\rho =\underline M_Y$，并且
$$
\varphi_*\varphi^*\mathcal{F}_\rho\to \mathcal{F}_\rho
$$
给出
\begin{align*}
& H_c^2(X_{\overline{k}}, \varphi_*\varphi^*\mathcal{F}_\rho) \to
H_c^2(X_{\overline{k}}, \mathcal{F}_\rho)\\
& =H_c^2(Y_{\overline{k}}, \varphi^*\mathcal{F}_\rho)\\
& =H_c^2(Y_{\overline{k}}, \underline M) = \oplus_{\text{不可约
分支 } \atop Y_{\overline{k}}}M
\end{align*}
$$
\Im(\rho) \to H_c^2(Y_{\overline{k}}, \underline M) =
\oplus_{\text{不可约分支 } \atop Y_{\overline{k}}}
M \to_{\Im(\rho) \text{等变}} H_c^2(X_{\overline{k}},
\mathcal{F}_{\rho}) \to^{\text{平凡的 }
\Im(\rho) \atop \text{作用}}
$$
对不可约曲线 $C/\overline{k}$，有 $H_c^2(C, \underline M)=M$。

\medskip\noindent
由于
$$
{\text{不可约分支的 } \atop \text{集合 }Y_k} =
\frac{Im(\rho)}{Im(\rho|_{\pi_1(X_{\overline{k}}, \overline \eta)})}
$$
可知 $H_c^2(X_{\overline{k}}, \mathcal{F}_\rho)$ 是
$M_{\pi_1(X_{\overline{k}}, \overline \eta)}$ 的一个商。另一方面，
存在满射
$$
\mathcal{F}_\rho\to \mathcal{F}'' = {\text{层，其在 }
X\text{ 上与下式相伴 } \atop (M)_{\pi_1(X_{\overline{k}}, \overline
\eta)}\leftarrow\pi_1(X, \overline \eta)}
$$
$$
H_c^2(X_{\overline{k}}, \mathcal{F}_\rho)\to
M_{\pi_1(X_{\overline{k}}, \overline\eta)}
$$
Galois 作用中的扭曲来自如下事实：
$H_c^2(X_{\overline{k}}, \mu_n)=^{\text{典范}} \mathbf{Z}/n\mathbf{Z}$。
\end{proof}

\begin{remark}
\label{trace-remark-projective}
因此可知，若 $X$ 还是投射的，则关于表示 $\rho$ 函子性地有等同
$$
H^0(X_{\overline{k}}, \mathcal{F}_\rho) =
M^{\pi_1(X_{\overline{k}}, \overline\eta)}
$$
以及
$$
H_c^2(X_{\overline{k}}, \mathcal{F}_\rho) =
M_{\pi_1(X_{\overline{k}}, \overline \eta)}(-1)
$$
当然，若 $X$ 非投射，则
$H^0_c(X_{\overline{k}}, \mathcal{F}_\rho) = 0$。
\end{remark}


\begin{proposition}
\label{trace-proposition-curve-kpi1}
设 $X/k$ 如前，但 $X_{\overline{k}}\neq \mathbf{P}^1_{\overline{k}}$。
诸函子
$
(M, \rho)\mapsto H_c^{2-i}(X_{\overline{k}}, \mathcal{F}_\rho)
$
是
$(M, \rho)\mapsto H_c^2(X_{\overline{k}}, \mathcal{F}_\rho)$
的左导出函子，故
$$
H_c^{2-i}(X_{\overline{k}}, \mathcal{F}_\rho) =
H_i(\pi_1(X_{\overline{k}}, \overline \eta), M)(-1)
$$
此外，还有如下导出版本：
$$
R\Gamma_c(X_{\overline{k}}, \mathcal{F}_\rho)
=
LH_0(\pi_1(X_{\overline{k}}, \overline \eta), M(-1))
=
M(-1)
\otimes_{\Lambda[[\pi_1(X_{\overline{k}}, \overline \eta)]]}^\mathbf{L}
\Lambda
$$
其成立于 $D(\Lambda[[\widehat{\mathbf{Z}}]])$ 中。
类似地，诸函子
$(M, \rho)\mapsto H^i(X_{\overline{k}}, \mathcal{F}_\rho)$
是
$(M, \rho)\mapsto M^{\pi_1(X_{\overline{k}}, \overline \eta)}$
的右导出函子，故
$$
H^i(X_{\overline{k}}, \mathcal{F}_\rho) =
H^i(\pi_1(X_{\overline{k}}, \overline \eta), M)
$$
在此情形下也有导出版本。
\end{proposition}

\begin{proof}
（思路）证明两端都是泛 $\delta$-函子。
\end{proof}

\begin{remark}
\label{trace-remark-poincare-groups}
由该命题及平凡对偶可得
$$
H^{2-i}_c(X_{\overline{k}}, \mathcal{F}_\rho)
\times
H^i(X_{\overline{k}}, \mathcal{F}_\rho^\wedge(1))
\to
\mathbf{Q}/\mathbf{Z}
$$
这一完美配对。若 $X$ 是投射的，这就是 Poincare 对偶。
\end{remark}





\section{抽象迹公式}
\label{trace-section-abstract-trace-formula}

\noindent
设给定预有限群的扩张
$$
1 \to G \to \Gamma \xrightarrow{\deg} \widehat{\mathbf{Z}} \to 1
$$
我们称 $\Gamma$ {\it 具有抽象迹公式}，当且仅当存在
\begin{enumerate}
\item 整数 $q\geq 1$；以及
\item 对每个 $d\geq 1$，有有限集 $S_d$，并且对每个 $x\in S_d$，
有满足 $\deg(F_x) = d$ 的共轭类 $F_x \in \Gamma$
\end{enumerate}
使得下列条件成立：
\begin{enumerate}
\item 对每个不整除 $q$ 的 $\ell$，有 $\text{cd}_\ell(G)<\infty$；并且
\item 对每个满足 $q\in \Lambda^*$ 的有限环 $\Lambda$、每个带连续
$\Gamma$-作用的有限投射 $\Lambda$-模 $M$ 以及每个 $n > 0$，有
$$
\sum\nolimits_{d|n}d \left(
\sum\nolimits_{x \in S_d}
\text{Tr}( F_x^{n/d} |_M)
\right)
=
q^n \text{Tr}(F^n|_{M \otimes_{\Lambda[[G]]}^{\mathbf{L}}\Lambda})
$$
其成立于 $\Lambda^\natural$ 中。
\end{enumerate}
这里，$M \otimes_{\Lambda[[G]]}^{\mathbf{L}}\Lambda = LH_0(G, M)$ 表示
导出同调，而 $F=1$ 于 $\Gamma/G = \widehat{\mathbf{Z}}$ 中。

\begin{remark}
\label{trace-remark-abstract-trace-formula}
关于这一概念，有如下几点观察。
\begin{enumerate}
\item 若所模拟的是投射曲线，则可以使用上同调，并且不需要因子 $q^n$。
\item 我所知道的例子只有 $\Gamma = \pi_1(X, \overline \eta)$，
其中 $X$ 是有限域上的光滑、几何不可约的 $K(\pi, 1)$。在此情形下，
$q = (\# k)^{\dim X}$。取上述命题为前提，我们已在本课程中对曲线证明了这一点。
\item 给定整数 $q$ 后，诸集合 $S_d$ 唯一确定。（可以把 $q$ 乘以整数
$m$，再以 $S_d$ 的 $m^d$ 份副本替换 $S_d$，公式不变。）
\end{enumerate}
\end{remark}

\begin{example}
\label{trace-example-commutative}
固定整数 $q\geq 1$，
$$
\begin{matrix}
1 &
\to &
G = \widehat{\mathbf{Z}}^{(q)} &
\to &
\Gamma &
\to &
\widehat{\mathbf{Z}} &
\to &
1 \\
&
&
= \prod_{l\not \mid q} \mathbf{Z}_l &
&
F &
\mapsto &
1
\end{matrix}
$$
并令 $FxF^{-1} = ux$，其中 $u \in (\widehat{\mathbf{Z}}^{(q)})^*$。
只需使用平凡模 $\mathbf{Z}/m\mathbf{Z}$，便可看出
$$
q^n - (qu)^n \equiv \sum\nolimits_{d|n} d\# S_d
$$
对所有 $(m, q)=1$ 均成立于 $\mathbf{Z}/m\mathbf{Z}$ 中（容许
$u \to u^{-1}$）；这蕴含 $qu = a\in \mathbf{Z}$ 且
$|a| < q$。特殊情形 $a = 1$ 确实出现于
$$
\Gamma = \pi_1^t(\mathbf{G}_{m, \mathbf{F}_p}, \overline \eta),
\quad
\# S_1 = q - 1,
\quad\text{且}\quad
\# S_2 = \frac{(q^2-1)-(q-1)}{2}
$$
\end{example}



%12.08.09
\section{自守形式与层}
\label{trace-section-automorphic}

\noindent
参考文献：尤其参见 Drinfeld 的杰出论文
\cite{D1}、\cite{D2} 和 \cite{D0}。

\medskip\noindent
{\bf 非分歧尖点形式。}
设 $k$ 为特征 $p$ 的有限域。
设 $X$ 为 $k$ 上几何不可约的投射光滑曲线。
令 $K = k(X)$ 为 $X$ 的函数域。
设 $v$ 为 $K$ 的一个位点，这等价于取闭点 $x\in X$。
令 $K_v$ 为 $K$ 在 $v$ 处的完备化；这也等价于取
$X$ 在 $x$ 处局部环之完备化的分式域。
以 $O_v\subset K_v$ 表示整数环。进一步令
$$
O = \prod\nolimits_v O_v \subset \mathbf{A} = \prod_v' K_v
$$
并令 $\Lambda$ 为任意满足 $p$ 在 $\Lambda$ 中可逆的环。

\begin{definition}
\label{trace-definition-unramified}
一个{\it 取值于 $\Lambda$ 的 $\text{GL}_2(\mathbf{A})$ 上非分歧尖点形式}
\footnote{这很可能不是标准记号。}
是一个函数
$$
f : \text{GL}_2(\mathbf{A}) \to \Lambda
$$
并满足
\begin{enumerate}
\item 对所有 $x\in \text{GL}_2(\mathbf{A})$ 及所有
$\gamma\in \text{GL}_2(K)$，有 $f(x\gamma) = f(x)$；
\item 对所有 $x\in \text{GL}_2(\mathbf{A})$ 及所有
$u\in \text{GL}_2(O)$，有 $f(ux) = f(x)$；
\item 对所有 $x\in \text{GL}_2(\mathbf{A})$，有
$$
\int_{\mathbf{A} \mod K} f
\left(x
\left(
\begin{matrix}
1 & z \\
0 & 1
\end{matrix}
\right)
\right) dz = 0
$$
关于当 $\Lambda$ 为使 $p$ 可逆的一般环时如何解释此式，
参见 \cite[Section 4.1]{dJ-conjecture}。
\end{enumerate}
\end{definition}

\noindent
{\bf Hecke 算子。}
对 $K$ 的一个位点 $v$ 及一个非分歧尖点形式 $f$，令
$$
T_v(f)(x) =
\int_{g\in M_v}f(g^{-1}x)dg,
$$
并令
$$
U_v(f)(x) =
f\left(
\left(
\begin{matrix}
\pi_v^{-1} & 0 \\
0 & \pi_v^{-1}
\end{matrix}
\right)x\right)
$$
这里使用如下记号：$\pi_v \in O_v$ 是一致化元，
$$
M_v =
\left\{
h\in Mat(2\times 2, O_v) | \det h = \pi_vO_v^*\right\}
$$
而 $dg = $ 是 $\text{GL}_2(K_v)$ 上满足
$\int_{\text{GL}_2(O_v)} dg = 1$ 的 Haar 测度。明确地，有
$$
T_v(f)(x) =
f\left(
\left(
\begin{matrix}
\pi_v^{-1}& 0 \\
0 & 1
\end{matrix}
\right)
x\right) +
\sum_{i = 1}^{q_v}
f\left(\left(
\begin{matrix}
1 & 0 \\
-\pi_v^{-1}\lambda_i
& \pi_v^{-1}
\end{matrix}
\right) x\right)
$$
其中 $\lambda_i\in O_v$ 构成
$O_v/(\pi_v)=\kappa_v$ 的一组代表元，$q_v = \#\kappa_v$。

\medskip\noindent
{\bf 特征形式。} 一个{\it 特征形式} $f$ 是满足如下条件的非分歧尖点形式：
$f$ 的某个值是单位元，并且对某些（唯一确定的）
$t_v, u_v \in \Lambda$，有 $T_vf = t_vf$ 及 $U_vf = u_vf$。

\begin{theorem}
\label{trace-theorem-drinfeld-make-rho}
给定取值于 $\overline{\mathbf{Q}}_l$ 且特征值满足
$u_v\in \overline{\mathbf{Z}}_l^*$ 的特征形式 $f$，则存在
$$
\rho : \pi_1(X)\to \text{GL}_2(E)
$$
连续绝对不可约表示，其中 $E$ 是包含于
$\overline{\mathbf{Q}}_l$ 中的 $\mathbf{Q}_\ell$ 的有限扩张，并且
对所有位点 $v$，均有 $t_v = \text{Tr}(\rho(F_v))$ 以及
$u_v = q_v^{-1}\det\left(\rho(F_v)\right)$。
\end{theorem}

\begin{proof}
见 \cite{D0}。
\end{proof}

\begin{theorem}
\label{trace-theorem-drinfeld-make-f}
设 $\mathbf{Q}_l \subset E$ 为有限扩张，并且
$$
\rho : \pi_1(X)\to \text{GL}_2(E)
$$
是绝对不可约连续表示。则存在取值于
$\overline{\mathbf{Q}}_l$ 的特征形式 $f$，其特征值 $t_v$、$u_v$
满足等式
$t_v = \text{Tr}(\rho(F_v))$ 及 $u_v = q_v^{-1}\det(\rho(F_v))$。
\end{theorem}

\begin{proof}
见 \cite{D1}。
\end{proof}

\begin{remark}
\label{trace-remark-lafforgue}
由于 Lafforgue 及许多其他数学家的工作，我们现在对
$\text{GL}_n$ 已有允许分歧的、类似上述两个定理的完整结果！
换言之，对有限域上曲线的函数域，$\text{GL}_n$ 的完整整体
Langlands 对应已经为人所知。同时，这并不意味着关于有限域上曲线的
基本群不再有大量有趣问题有待回答，正如下文将看到的那样。
\end{remark}

\noindent
{\bf 中心特征标。} 若 $f$ 是特征形式，则
$$
\begin{matrix}
\chi_f : &
O^*\backslash \mathbf{A}^*/K^* &
\to &
\Lambda^* \\
&
(1, \ldots, \pi_v, 1, \ldots, 1) &
\mapsto &
u_v^{-1}
\end{matrix}
$$
称为{\it 中心特征标}。按照上述规范化，它对应于 $\rho$ 的行列式。令
$$
C(\Lambda) =
\left\{
{\text{非分歧尖点形式 } f \text{，其系数在 }\Lambda}
\atop {\text{ 满足 } U_v f = \varphi_v^{-1}f\forall v}
\right\}
$$

\begin{proposition}
\label{trace-proposition-cusp-forms-finite}
若 $\Lambda$ 是 Noether 环，则 $C(\Lambda)$ 是有限生成
$\Lambda$-模。此外，若 $\Lambda$ 是以
$\mathbf{F} \subset \Lambda$ 为素子域的域，则
$$
C(\Lambda)=(C(\mathbf{F}))\otimes_{\mathbf{F}}\Lambda
$$
并与 $T_v$ 的作用相容。
\end{proposition}

\begin{proof}
见 \cite[Proposition 4.7]{dJ-conjecture}。
\end{proof}

\noindent
该命题立即蕴含下列引理。

\begin{lemma}
\label{trace-lemma-eigenvalues-algebraic}
特征值的代数性。
若 $\Lambda$ 是域，则对 $f\in C(\Lambda)$，其特征值 $t_v$
在素子域 $\mathbf{F} \subset \Lambda$ 上是代数的。
\end{lemma}

\begin{proof}
由命题 \ref{trace-proposition-cusp-forms-finite} 得出。
\end{proof}

\noindent
综合以上结果，可以施行下列十分有用的技巧。

\begin{lemma}
\label{trace-lemma-switch-l}
更换 $l$。设 $E$ 为数域。
从绝对不可约连续表示
$$
\rho : \pi_1(X)\to SL_2(E_\lambda)
$$
出发，其中 $\lambda$ 是 $E$ 的一个不位于 $p$ 之上的位点。
则对 $E$ 的任意另一个不位于 $p$ 之上的位点 $\lambda'$，
存在有限扩张 $E'_{\lambda'}$ 及绝对不可约连续表示
$$
\rho': \pi_1(X)\to SL_2(E'_{\lambda'})
$$
它与 $\rho$ 相容，即所有 Frobenius 元的特征多项式均相同。
\end{lemma}

\noindent
注意，这是 Deligne 猜想的一个实例！

\begin{proof}
为证明更换引理，应用定理 \ref{trace-theorem-drinfeld-make-f}，
得到与 $\rho$ 相伴的特征形式
$f\in C(\overline{\mathbf{Q}}_l)$。
再由命题 \ref{trace-proposition-cusp-forms-finite} 可知，可以选择
$f\in C(E')$，其中 $E \subset E'$ 为有限扩张。
接着完备化 $E'$，得到特征形式 $f\in C(E'_{\lambda'})$，
其中 $E'_{\lambda'}$ 是 $E_{\lambda'}$ 的有限扩张。
最后应用定理 \ref{trace-theorem-drinfeld-make-rho}，并在必要时再次稍微扩大
$E'_{\lambda'}$，便得到绝对不可约连续表示
$\rho': \pi_1(X) \to SL_2(E_{\lambda'}')$。
\end{proof}

\noindent
推测：若对某个（拓扑）环 $\Lambda$ 有
$$
\left(
{\rho : \pi_1(X)\to SL_2(\Lambda) \atop \text{ 绝对不可约}}
\right)
\leftrightarrow
\text{属于下式的特征形式 } C(\Lambda)
$$
则 $\rho(F_v)$ 的所有特征值都是代数的（当 $\Lambda$ 是有限环时，
这不会以一种简单方式成立）。基于 Langlands 对应不仅对域成立、
而且在更一般情形下也成立这一推测，可以得到下列猜想。

\medskip\noindent
{\bf 猜想。}
（见 \cite{dJ-conjecture}）
对任意连续表示
$$
\rho : \pi_1(X)\to \text{GL}_n(\mathbf{F}_l[[t]])
$$
均有 $\# \rho(\pi_1(X_{\overline{k}}))<\infty$。

\medskip\noindent
用层的语言改述如下：
“对任意 $\overline{\mathbf{F}_l((t))}$-模的光滑层，其几何
单值群都是有限的。”

\begin{theorem}
\label{trace-theorem-conjecture-n-2}
若 $n\leq 2$，该猜想成立。
\end{theorem}

\begin{proof}
见 \cite{dJ-conjecture}。
\end{proof}

\begin{theorem}
\label{trace-theorem-conjecture-l-bigger-2n}
在若干尚未证明的结果的条件下，若 $l > 2n$，该猜想成立。
\end{theorem}

\begin{proof}
见 \cite{Gaitsgory}。
\end{proof}

\noindent
事实证明，这一猜想确有用途。
参见 Drinfeld 关于 Kashiwara 猜想的工作。除此之外，还有如下
更为具体的应用。

\begin{theorem}
\label{trace-theorem-deformation-rings}
（见 \cite[Theorem 3.5]{dJ-conjecture}）
设
$$
\rho_0: \pi_1(X)\to \text{GL}_n(\mathbf{F}_l)
$$
连续，且 $l\neq p$。假设
\begin{enumerate}
\item 对 $X$，该猜想成立；
\item $\rho_0 |_{\pi_1(X_{\overline{k}})}$ 绝对不可约；并且
\item $l$ 不整除 $n$。
\end{enumerate}
则 $\rho_0$ 的泛形变环 $R_{\text{univ}}$ 在 $\mathbf{Z}_l$ 上有限平坦。
\end{theorem}

\noindent
说明：存在表示 $\rho_{\text{univ}}:
\pi_1(X)\to \text{GL}_n(R_{\text{univ}})$（泛形变环），其中
$R_{\text{univ}}$ 是完备局部环，剩余域为 $\mathbf{F}_l$，且
$(R_{\text{univ}}\to
\mathbf{F}_l)\circ\rho_{\text{univ}}\cong\rho_0$。
给定任意 $R\to \mathbf{F}_l$，其中 $R$ 是完备局部环，并给定
$\rho : \pi_1(X)\to \text{GL}_n(R)$，则存在
$\psi : R_{\text{univ}}\to R$，使得
$\psi\circ\rho_{\text{univ}}\cong \rho$。定理断言态射
$$
\Spec(R_{\text{univ}})
\longrightarrow
\Spec(\mathbf{Z}_l)
$$
是有限平坦的。特别地，这样的 $\rho_0$ 可提升为
$\rho : \pi_1(X) \to \text{GL}_n(\overline{\mathbf{Q}}_l)$。

\medskip\noindent
补注：
\begin{enumerate}
\item 关于形变的定理很容易。
\item 朝该猜想迈进的任何结果看来都很困难。
\item 若能有更多关于 $\pi_1(X)$ 的猜想，将会很有趣！
\end{enumerate}




%12.10.09
\section{计数点}
\label{trace-section-counting}

\noindent
设 $X$ 为 $k$ 上光滑、几何不可约的投射曲线，且
$q = \# k$。迹公式给出：存在代数整数
$w_1, \ldots, w_{2g}$，使得
$$
\# X(k_n) = q^n - \sum\nolimits_{i = 1}^{2g_X} w_i^n + 1.
$$
若 $\sigma\in \text{Aut}(X)$，则对每个 $i$，存在 $j$ 使得
$\sigma(w_i)=w_j$。

\medskip\noindent
{\bf Riemann 假设。} 对所有 $i$，有 $|\omega_i| = \sqrt{q}$。

\medskip\noindent
Emil Artin 于 1924 年对超椭圆曲线提出了这一命题。
Weil 于 1940 年证明了它。Weil 给出了两个证明：
\begin{itemize}
\item 使用 $X \times X$ 上的相交理论及 Hodge 指数定理；以及
\item 使用 $X$ 的 Jacobian。
\end{itemize}
另有一个由 Bombieri 给出的证明，其最初思想源自 Stephanov：
它使用函数域 $k(X)$ 及其 Frobenius 算子（1969）。出发点是，
给定 $f\in k(X)$，可观察到 $f^q - f$ 是一个在 $X$ 的所有
$\mathbf{F}_q$-有理点处消失的有理函数，并可尝试用这一思想给出
点数的上界。


\section{Chebotarev 的精确形式}
\label{trace-section-chebotarev}

\noindent
作为第一个应用，我们来证明曲线的有限 \'etale Galois 覆叠所对应的
Chebotarev 定理之精确形式。
设 $\varphi : Y \to X$ 为群为 $G$ 的有限 \'etale Galois 覆叠。
它对应于同态
$$
\pi_1(X) \longrightarrow G = \text{Aut}(Y/X)
$$
假设 $Y_{\overline{k}} = $ 不可约。若 $C\subset G$ 是共轭类，
则对所有 $n > 0$，有
$$
| \# \{x \in X(k_n) \mid F_x \in C\} - \frac{\# C}{\# G} \cdot \# X(k_n) |
\leq
(\# C)(2g - 2) \sqrt{q^n}
$$
（警告：使用前请仔细核对右端的系数 $\# C$。）

\begin{proof}[略证]
写成
$$
\varphi_*(\overline{\mathbf{Q}_l}) =
\oplus_{\pi \in \widehat{G}} \mathcal{F}_{\pi}
$$
其中 $\widehat{G}$ 是 $G$ 在 $\overline{\mathbf{Q}}_l$ 上
不可约表示的同构类之集合。对 $\pi \in \widehat{G}$，
令 $\chi_{\pi}: G \to \overline{\mathbf{Q}}_l$
为 $\pi$ 的特征标。于是
$$
H^*(Y_{\overline{k}}, \overline{\mathbf{Q}}_l) =
\oplus_{\pi\in \widehat{G}}
H^*(Y_{\overline{k}}, \overline{\mathbf{Q}}_l)_\pi
=_{(\varphi\text{ 有限 })}
\oplus_{\pi\in \widehat{G}}
H^*(X_{\overline{k}}, \mathcal{F}_\pi)
$$
若 $\pi\neq 1$，则
$$
H^0(X_{\overline{k}}, \mathcal{F}_\pi) =
H^2(X_{\overline{k}}, \mathcal{F}_\pi) = 0,\quad
\dim H^1(X_{\overline{k}}, \mathcal{F}_\pi) = (2g_X - 2)d_\pi^2
$$
（可由作用于……的迹公式得到），由此看出
$$
|\sum_{x \in X(k_n)} \chi_\pi(\mathcal{F}_x)| \leq
(2g_X - 2) d_\pi^2\sqrt{q^n}
$$
写 $1_C = \sum_\pi a_\pi \chi_\pi$，则
$a_\pi = \langle 1_C, \chi_\pi\rangle$，且
$a_1 = \langle 1_C, \chi_1\rangle = \frac{\# C}{\# G}$，其中
$$
\langle f, h\rangle = \frac{1}{\# G}\sum_{g \in G} f(g)\overline{h(g)}
$$
因此有关系式
$$
\frac{\# C}{\# G} = ||1_C||^2 = \sum|a_\pi|^2
$$
最后一步：
\begin{align*}
\#\left\{x \in X(k_n) \mid F_x \in C\right\}
& =
\sum_{x \in X(k_n)} 1_C(x) \\
& =
\sum_{x \in X(k_n)} \sum_\pi a_\pi \chi_\pi(F_x) \\
& =
\underbrace{\frac{\# C}{\# G} \# X(k_n)}_{
\text{对应于 }\pi = 1}
+
\underbrace{\sum_{\pi\neq 1}a_\pi\sum_{x\in X(k_n)}\chi_\pi(F_x)}_{
\text{误差项（界为 }E)}
\end{align*}
误差项可估计为
\begin{align*}
|E|
& \leq
\sum_{\pi \in \widehat{G}, \atop \pi \neq 1}
|a_\pi| (2g - 2) d_\pi^2 \sqrt{q^n} \\
& \leq
\sum_{\pi \neq 1} \frac{\# C}{\# G} (2g_X - 2) d_\pi^3 \sqrt{q^n}
\end{align*}
由 Weil 猜想，$\# X(k_n)\sim q^n$。
\end{proof}



\section{有多少素点完全分解？}
\label{trace-section-how-many}

\noindent
本节给出有限域上曲线的 Riemann 假设的第二个应用。对数论学家而言，
Ihara 的论文《整体域的无限非分歧 Galois 扩张中有多少素点完全分解？》
或许值得一读，见 \cite{Ihara}。
考虑基本正合列
$$
1 \to
\pi_1(X_{\overline{k}}) \to
\pi_1(X) \xrightarrow{\deg}
\widehat{\mathbf{Z}} \to 1
$$

\begin{proposition}
\label{trace-proposition-finite-set-frobenii-generate-topologically}
存在由 $X$ 的闭点组成的有限集 $x_1, \ldots, x_n$，
使得与这些点对应的{\bf 所有} Frobenius 元在拓扑意义下生成
$\pi_1(X)$。
\end{proposition}

\noindent
另一种表述是：存在 $x_1, \ldots, x_n\in |X|$，使得
$\pi_1(X)$ 中包含每个 $x_i$ 所对应的 $1$ 个 Frobenius 元的最小正规闭子群
$\Gamma$ 就是整个 $\pi_1(X)$，即 $\Gamma = \pi_1(X)$。

\begin{proof}
取 $N\gg 0$，并令
$$
\{x_1, \ldots, x_n\} =
{\text{所有闭点的集合，其位于}
\atop X \text{ 且次数} \leq N\text{（相对于 } k}
$$
令 $\Gamma\subset \pi_1(X)$ 为上述等价表述中与这些点对应的子群。
假设 $\Gamma \neq \pi_1(X)$。于是可以选取 $\pi_1(X)$ 的一个正规开子群
$U$，使其包含 $\Gamma$ 且 $U \neq \pi_1(X)$。由 $X$ 的 Riemann 假设，
我们的点集中将有某个次数为 $N$ 的 $x_{i_1}$，以及某个次数为
$N - 1$ 的 $x_{i_2}$。这表明
$\deg : \Gamma \to \widehat{\mathbf{Z}}$ 是满射，
故对 $U$ 亦然。这恰好意味着：若 $Y \to X$ 是与 $U$ 对应的有限
\'etale Galois 覆叠，则 $Y_{\overline{k}}$ 不可约。
令 $G = \text{Aut}(Y/X)$。图示为
$$
Y \to^G X,\quad G = \pi_1(X)/U
$$
由构造，$X$ 中所有次数 $\leq N$ 的点都在 $Y$ 中完全分解。因此特别有
$$
\# Y(k_N)\geq (\# G)\# X(k_N)
$$
对两端应用 Riemann 假设，得到
$$
q^N+1+2g_Yq^{N/2}\geq \# G\# X(k_N)\geq \#
G(q^N+1-2g_Xq^{N/2})
$$
由于 $2g_Y-2 = (\# G)(2g_X-2)$，这意味着
$$
q^N + 1 + (\# G)(2g_X - 1) + 1)q^{N/2}\geq
\# G (q^N + 1 - 2g_Xq^{N/2})
$$
由此可见，当 $N$ 足够大时，$G$ 必须是平凡群。
\end{proof}

\noindent
{\bf 奇怪的问题。}
令 $W_X = \deg^{-1}(\mathbf{Z})\subset \pi_1(X)$。
是否存在由 $X$ 的闭点组成的某个有限集 $x_1, \ldots, x_n$，
使得与这些点对应的所有 Frobenius 元在{\it 代数意义下}生成 $W_X$？

\medskip\noindent
由 Baire 纲论证，这等价于对所有 Frobenius 元提出同一问题。





\section{实际上究竟有多少个点？}
\label{trace-section-really}

\noindent
若曲线的亏格相对于 $q$ 很大，则公式
$\# X(k) = q - \sum \omega_i + 1$ 中的主项并非 $q$，
而是第二项 $\sum \omega_i$；后者先验地可以具有约
$2g_X\sqrt{q}$ 的大小。在论文 \cite{Drinfeld-number} 中，
Drinfeld 和 Vladut 证明，该最大值实际上至多约为
$g\sqrt{q}$，这正是 Ihara 此前所预言的。

\medskip\noindent
固定 $q$，并令 $k$ 为具有 $k$ 个元素的域。令
$$
A(q) = \limsup_{g_X \to \infty} \frac{\# X(k)}{g_X}
$$
其中 $X$ 遍历 $k$ 上所有几何不可约的光滑投射曲线。按此定义，
有下列结果：
\begin{itemize}
\item RH $\Rightarrow A(q)\leq 2\sqrt{q}$
\item Ihara $\Rightarrow A(q)\leq \sqrt{2q}$
\item DV $\Rightarrow A(q)\leq \sqrt{q}-1$（事实上，若 $q$
是平方数，则此界精确）
\end{itemize}

\begin{proof}
给定 $X$，令 $w_1, \ldots, w_{2g}$ 如前，并令 $g = g_X$。
置 $\alpha_i = \frac{w_i}{\sqrt{q}}$，于是 $|\alpha_i| = 1$。
若出现 $\alpha_i$，则 $\overline{\alpha}_i = \alpha_i^{-1}$ 也出现。于是
$$
N = \# X(k) \leq X(k_r) = q^r + 1 - (\sum_i \alpha_i^r) q^{r/2}
$$
改写后可见，对每个 $r \geq 1$，
$$
-\sum_i \alpha_i^r \geq Nq^{-r/2} - q^{r/2} - q^{-r/2}
$$
注意到
$$
0 \leq |\alpha_i^n +\alpha_i^{n-1} +\ldots +\alpha_i +1|^2
= (n + 1) + \sum_{j = 1}^n (n + 1 - j) (\alpha_i^j + \alpha_i^{-j})
$$
因此
\begin{align*}
2g(n+1) & \geq - \sum_i \left(\sum_{j = 1}^n (n+1-j)(\alpha_i^j
+\alpha_i^{-j})\right)\\
& =-\sum_{j = 1}^n (n+1-j)\left(\sum_i\alpha_i^j
+\sum_i\alpha_i^{-j}\right)
\end{align*}
取此式的一半，得到
\begin{align*}
g(n+1)& \geq - \sum_{j = 1}^n (n+1-j)(\sum_i\alpha_i^j)\\
& \geq N\sum_{j = 1}^n (n+1-j)q^{-j/2}-\sum_{j = 1}^n
(n+1-j)(q^{j/2}+q^{-j/2})
\end{align*}
这给出
$$
\frac{N}{g}\leq \left(\sum_{j = 1}^n \frac{n+1-j}{n+1}q^{-j/2} \right)^{-1}
\cdot
\left(
1 + \frac{1}{g} \sum_{j = 1}^n \frac{n + 1 - j}{n + 1}(q^{j/2} + q^{-j/2})
\right)
$$
固定 $n$，令 $g\to \infty$，
$$
A(q)\leq \left(\sum_{j = 1}^n \frac{n+1-j}{n+1}q^{-j/2}\right)^{-1}
$$
于是
$$
A(q)\leq \lim_{n\to\infty}(\ldots) = \left(\sum_{j = 1}^\infty
q^{-j/2}\right)^{-1}=\sqrt{q}-1
$$
\end{proof}
% Generated mechanically from the frozen spaces.tex target.
% Do not edit this generated file; edit the bound locale source instead.

% OK, start here.
%

\setcounter{chapter}{64}
\chapter{代数空间}



\phantomsection
\label{spaces-section-phantom}


\section{引言}
\label{spaces-section-introduction}

\noindent
代数空间最初由 Michael Artin 引入，
见 \cite{ArtinI}、\cite{ArtinII}、
\cite{Artin-Theorem-Representability},
\cite{Artin-Construction-Techniques},
\cite{Artin-Algebraic-Spaces},
\cite{Artin-Algebraic-Approximation},
\cite{Artin-Implicit-Function},
以及 \cite{ArtinVersal}。
其中一些基础材料由他与 Knutson 共同发展，
后者写成了专著 \cite{Kn}。
Artin 将代数空间定义为 \'etale 拓扑下的层，
并要求它在 \'etale 拓扑下局部可表示
（见 \cite[Definition 1.3]{Artin-Implicit-Function}）。
在 Artin 的大部分工作中，所考虑的概形范畴
由某个固定的优秀 Noether 基概形上的局部有限型概形组成。

\medskip\noindent
我们的定义与 Artin 的原始定义略有不同。
具体而言，我们所称的代数空间是 fppf 拓扑下的层，
其对角态射可表示，并且存在一个由概形给出的 \'etale “覆盖”。
用 fppf 拓扑代替 \'etale 拓扑只是技术上的处理，
几乎不产生实质差异；我们将在“自举”第
\ref{bootstrap-section-spaces-etale} 节证明，
若改用 \'etale 拓扑，所得代数空间范畴仍然相同。
在同一章中，我们还会证明对角态射的条件在某种意义下可以去掉，
见“自举”第 \ref{bootstrap-section-bootstrap} 节。

\medskip\noindent
定义代数空间之后，我们作若干基础性观察。
本章的主要结果是：按我们的定义，代数空间与 \'etale 等价关系是同一回事；
见第 \ref{spaces-section-presentations} 节的讨论以及定理
\ref{spaces-theorem-presentation}。在 Artin 的框架中，对应结果为
\cite[Theorem 1.5]{Artin-Implicit-Function} 或
\cite[Proposition II.1.7]{Kn}。换言之，由 \'etale 等价关系定义的层
具有可表示的对角态射。因此，我们的定义在广义上与 Artin 的原始定义一致。
这还意味着，只需写出一个 \'etale 等价关系，便可给出代数空间的例子。

\medskip\noindent
在第 \ref{spaces-section-separation} 节中，我们介绍文献中出现的
若干代数空间分离公理。最后，在第 \ref{spaces-section-examples} 节中，
我们给出一些怪异以及不那么怪异的代数空间例子。






\section{一般说明}
\label{spaces-section-general}

\noindent
我们在一个合适的大 fppf 位点 $\Sch_{fppf}$ 上工作，
如“拓扑”定义 \ref{topologies-definition-big-fppf-site} 所述。
因此，除非另有明确说明，所有概形均指 $\Sch_{fppf}$ 的对象。
第 \ref{spaces-section-change-big-site} 节讨论改变大 fppf 位点后会发生什么。

\medskip\noindent
我们总是相对于包含在 $\Sch_{fppf}$ 中的基概形 $S$ 工作，
并使用大 fppf 位点 $(\Sch/S)_{fppf}$；
见“拓扑”定义 \ref{topologies-definition-big-small-fppf}。
取 $S = \Spec(\mathbf{Z})$ 即可恢复绝对情形。

\medskip\noindent
若 $U,T$ 是 $S$ 上的概形，则以 $U(T)$ 表示相对于 $S$ 的
$T$-值点集。用公式表示即为 $U(T)=\Mor_S(T,U)$。

\medskip\noindent
注意，任意 fpqc 覆盖都是泛有效满态射；见“下降”引理
\ref{descent-lemma-fpqc-universal-effective-epimorphisms}。
因此，$\Sch_{fppf}$ 上的拓扑弱于典范拓扑，
并且所有可表示预层都是层。







\section{预层的可表示态射}
\label{spaces-section-representable}

\noindent
设 $S$ 是包含于 $\Sch_{fppf}$ 的概形，并设
$F,G:(\Sch/S)_{fppf}^{opp}\to\textit{Sets}$。
设 $a:F\to G$ 是可表示的函子变换；见“范畴”定义
\ref{categories-definition-representable-map-presheaves}。
这意味着，对每个 $U\in\Ob((\Sch/S)_{fppf})$ 以及任意
$\xi\in G(U)$，纤维积 $h_U\times_{\xi,G}F$ 都可表示。
选取一个表示对象 $V_\xi$ 和同构
$h_{V_\xi}\to h_U\times_G F$。
由 Yoneda 引理（见“范畴”引理 \ref{categories-lemma-yoneda}），
投影 $h_{V_\xi}\to h_U\times_G F\to h_U$ 来自唯一的概形态射
$a_\xi:V_\xi\to U$。这一情形可形象地表示为图表
$$
\xymatrix{
V_\xi \ar@{~>}[r] \ar[d]_{a_\xi} & h_{V_\xi} \ar[d] \ar[r] & F \ar[d]^a \\
U \ar@{~>}[r] & h_U \ar[r]^\xi & G
}
$$
其中波形箭头表示 Yoneda 嵌入。下面给出关于这一概念的若干引理，
它们在非常一般的情形下成立。

\begin{lemma}
\label{spaces-lemma-morphism-schemes-gives-representable-transformation}
设 $S$ 是包含于 $\Sch_{fppf}$ 的概形，$X$、$Y$ 是
$(\Sch/S)_{fppf}$ 的对象，并设 $f:X\to Y$ 是概形态射。则
$$
h_f : h_X \longrightarrow h_Y
$$
是可表示的函子变换。
\end{lemma}

\begin{proof}
这是形式的，只用到范畴 $(\Sch/S)_{fppf}$ 具有纤维积这一事实。
\end{proof}

\begin{lemma}
\label{spaces-lemma-composition-representable-transformations}
设 $S$ 是包含于 $\Sch_{fppf}$ 的概形，并设
$F,G,H:(\Sch/S)_{fppf}^{opp}\to\textit{Sets}$。
若 $a:F\to G$ 与 $b:G\to H$ 是可表示的函子变换，则
$$
b \circ a : F \longrightarrow H
$$
是可表示的函子变换。
\end{lemma}

\begin{proof}
这完全是形式的，并且在任意范畴中成立。
\end{proof}

\begin{lemma}
\label{spaces-lemma-base-change-representable-transformations}
设 $S$ 是包含于 $\Sch_{fppf}$ 的概形，并设
$F,G,H:(\Sch/S)_{fppf}^{opp}\to\textit{Sets}$。
设 $a:F\to G$ 是可表示的函子变换，
$b:H\to G$ 是任意函子变换。考虑纤维积图表
$$
\xymatrix{
H \times_{b, G, a} F \ar[r]_-{b'} \ar[d]_{a'} & F \ar[d]^a \\
H \ar[r]^b & G
}
$$
则基变换 $a'$ 是可表示的函子变换。
\end{lemma}

\begin{proof}
这完全是形式的，并且在任意范畴中成立。
\end{proof}

\begin{lemma}
\label{spaces-lemma-product-representable-transformations}
设 $S$ 是包含于 $\Sch_{fppf}$ 的概形。对 $i = 1, 2$，设
$F_i,G_i:(\Sch/S)_{fppf}^{opp}\to\textit{Sets}$；再对 $i = 1, 2$，设
$a_i:F_i\to G_i$ 是可表示的函子变换。则
$$
a_1 \times a_2 : F_1 \times F_2 \longrightarrow G_1 \times G_2
$$
是可表示的函子变换。
\end{lemma}

\begin{proof}
将 $a_1\times a_2$ 写成复合
$F_1\times F_2\to G_1\times F_2\to G_1\times G_2$。
第一支箭头是 $a_1$ 沿映射 $G_1\times F_2\to G_1$ 的基变换，
第二支箭头是 $a_2$ 沿映射 $G_1\times G_2\to G_2$ 的基变换。
因此，本引理是引理 \ref{spaces-lemma-composition-representable-transformations}
与 \ref{spaces-lemma-base-change-representable-transformations} 的形式推论。
\end{proof}

\begin{lemma}
\label{spaces-lemma-representable-transformation-to-sheaf}
设 $S$ 是包含于 $\Sch_{fppf}$ 的概形，并设
$F,G:(\Sch/S)_{fppf}^{opp}\to\textit{Sets}$。
设 $a:F\to G$ 是可表示的函子变换。若 $G$ 是层，则 $F$ 也是层。
\end{lemma}

\begin{proof}
设 $\{\varphi_i:T_i\to T\}$ 是位点 $(\Sch/S)_{fppf}$ 的一个覆盖，
并设 $s_i\in F(T_i)$ 满足层条件。则
$\sigma_i=a(s_i)\in G(T_i)$ 也满足层条件，故存在唯一的
$\sigma\in G(T)$ 使得 $\sigma_i=\sigma|_{T_i}$。由假设，
$F'=h_T\times_{\sigma,G,a}F$ 是可表示预层，因而是层
（见第 \ref{spaces-section-general} 节的说明）。又注意到
$(\varphi_i,s_i)\in F'(T_i)$ 也满足层条件，故它们来自唯一的
$(\text{id}_T,s)\in F'(T)$。显然，$s$ 正是所需的 $F$ 的截面。
\end{proof}

\begin{lemma}
\label{spaces-lemma-representable-transformation-diagonal}
设 $S$ 是包含于 $\Sch_{fppf}$ 的概形，并设
$F,G:(\Sch/S)_{fppf}^{opp}\to\textit{Sets}$。
设 $a:F\to G$ 是可表示的函子变换。则
$\Delta_{F/G}:F\to F\times_G F$ 可表示。
\end{lemma}

\begin{proof}
取 $U\in\Ob((\Sch/S)_{fppf})$，并设
$\xi=(\xi_1,\xi_2)\in(F\times_G F)(U)$。
令 $\xi'=a(\xi_1)=a(\xi_2)\in G(U)$。
由假设，存在概形 $V$ 和态射 $V\to U$ 表示纤维积
$h_U\times_{\xi',G}F$。特别地，元素 $\xi_1,\xi_2$ 给出 $U$ 上的态射
$f_1,f_2:U\to V$。由于 $V$ 表示纤维积
$h_U\times_{\xi',G}F$，并且
$\xi'=a\circ\xi_1=a\circ\xi_2$，可见对任意态射 $g:U'\to U$ 有
$$
g^*\xi_1 = g^*\xi_2
\Leftrightarrow
f_1 \circ g = f_2 \circ g.
$$
换言之，$h_U\times_{\xi,F\times_G F}F$ 由概形
$V\times_{\Delta,V\times V,(f_1,f_2)}U$ 表示。
\end{proof}










\section{概形态射的若干常用性质列表}
\label{spaces-section-lists}

\noindent
为便于引用，下面几则评注列出一些态射性质；
它们具有后续结果所需的若干稳定性。

\begin{remark}
\label{spaces-remark-list-properties-stable-base-change}
下列态射性质或类型在{\it 任意基变换下稳定}：
\begin{enumerate}
\item 闭浸入、开浸入和局部闭浸入，见“概形”引理
\ref{schemes-lemma-base-change-immersion}；
\item 拟紧，见“概形”引理
\ref{schemes-lemma-quasi-compact-preserved-base-change}；
\item 泛闭，见“概形”定义 \ref{schemes-definition-universally-closed}；
\item （拟）分离，见“概形”引理 \ref{schemes-lemma-separated-permanence}；
\item 单态射，见“概形”引理 \ref{schemes-lemma-base-change-monomorphism}；
\item 满射，见“态射”引理 \ref{morphisms-lemma-base-change-surjective}；
\item 泛单射，见“态射”引理 \ref{morphisms-lemma-universally-injective}；
\item 仿射，见“态射”引理 \ref{morphisms-lemma-base-change-affine}；
\item 拟仿射，见“态射”引理 \ref{morphisms-lemma-base-change-quasi-affine}；
\item （局部）有限型，见“态射”引理 \ref{morphisms-lemma-base-change-finite-type}；
\item （局部）拟有限，见“态射”引理 \ref{morphisms-lemma-base-change-quasi-finite}；
\item （局部）有限表示，见“态射”引理
\ref{morphisms-lemma-base-change-finite-presentation}；
\item 相对维数为 $d$ 的局部有限型，见“态射”引理
\ref{morphisms-lemma-base-change-relative-dimension-d}；
\item 泛开，见“态射”定义 \ref{morphisms-definition-open}；
\item 平坦，见“态射”引理 \ref{morphisms-lemma-base-change-flat}；
\item syntomic，见“态射”引理 \ref{morphisms-lemma-base-change-syntomic}；
\item 平滑，见“态射”引理 \ref{morphisms-lemma-base-change-smooth}；
\item 非分歧（相应地，G-非分歧），见“态射”引理
\ref{morphisms-lemma-base-change-unramified}；
\item \'etale，见“态射”引理 \ref{morphisms-lemma-base-change-etale}；
\item 固有，见“态射”引理 \ref{morphisms-lemma-base-change-proper}；
\item H-射影，见“态射”引理 \ref{morphisms-lemma-H-projective-base-change}；
\item （局部）射影，见“态射”引理 \ref{morphisms-lemma-base-change-projective}；
\item 有限或整，见“态射”引理 \ref{morphisms-lemma-base-change-finite}；
\item 有限局部自由，见“态射”引理
\ref{morphisms-lemma-base-change-finite-locally-free}；
\item 泛商，见“态射”引理
\ref{morphisms-lemma-base-change-universally-submersive}；
\item 泛同胚，见“态射”引理
\ref{morphisms-lemma-base-change-universal-homeomorphism}。
\end{enumerate}
可按需要补充。
\end{remark}

\begin{remark}
\label{spaces-remark-list-properties-stable-composition}
在基变换下稳定的态射性质中
（见评注 \ref{spaces-remark-list-properties-stable-base-change} 的列表），
下列性质还{\it 对复合稳定}：
\begin{enumerate}
\item 闭浸入、开浸入和局部闭浸入，见“概形”引理
\ref{schemes-lemma-composition-immersion}；
\item 拟紧，见“概形”引理 \ref{schemes-lemma-composition-quasi-compact}；
\item 泛闭，见“态射”引理 \ref{morphisms-lemma-composition-proper}；
\item （拟）分离，见“概形”引理 \ref{schemes-lemma-separated-permanence}；
\item 单态射，见“概形”引理 \ref{schemes-lemma-composition-monomorphism}；
\item 满射，见“态射”引理 \ref{morphisms-lemma-composition-surjective}；
\item 泛单射，见“态射”引理
\ref{morphisms-lemma-composition-universally-injective}；
\item 仿射，见“态射”引理 \ref{morphisms-lemma-composition-affine}；
\item 拟仿射，见“态射”引理 \ref{morphisms-lemma-composition-quasi-affine}；
\item （局部）有限型，见“态射”引理
\ref{morphisms-lemma-composition-finite-type}；
\item （局部）拟有限，见“态射”引理
\ref{morphisms-lemma-composition-quasi-finite}；
\item （局部）有限表示，见“态射”引理
\ref{morphisms-lemma-composition-finite-presentation}；
\item 泛开，见“态射”引理 \ref{morphisms-lemma-composition-open}；
\item 平坦，见“态射”引理 \ref{morphisms-lemma-composition-flat}；
\item syntomic，见“态射”引理 \ref{morphisms-lemma-composition-syntomic}；
\item 平滑，见“态射”引理 \ref{morphisms-lemma-composition-smooth}；
\item 非分歧（相应地，G-非分歧），见“态射”引理
\ref{morphisms-lemma-composition-unramified}；
\item \'etale，见“态射”引理 \ref{morphisms-lemma-composition-etale}；
\item 固有，见“态射”引理 \ref{morphisms-lemma-composition-proper}；
\item H-射影，见“态射”引理 \ref{morphisms-lemma-H-projective-composition}；
\item 有限或整，见“态射”引理 \ref{morphisms-lemma-composition-finite}；
\item 有限局部自由，见“态射”引理
\ref{morphisms-lemma-composition-finite-locally-free}；
\item 泛商，见“态射”引理
\ref{morphisms-lemma-composition-universally-submersive}；
\item 泛同胚，见“态射”引理
\ref{morphisms-lemma-composition-universal-homeomorphism}。
\end{enumerate}
可按需要补充。
\end{remark}

\begin{remark}
\label{spaces-remark-list-properties-fpqc-local-base}
在上述基变换下稳定的性质中
（见评注 \ref{spaces-remark-list-properties-stable-base-change}），
下列性质还是{\it 基上的 fpqc 局部性质}
（从而更是基上的 fppf 局部性质）：
\begin{enumerate}
\item 对浸入而言，下列情形成立：
\begin{enumerate}
\item 闭浸入，见“下降”引理
\ref{descent-lemma-descending-property-closed-immersion}；
\item 开浸入，见“下降”引理
\ref{descent-lemma-descending-property-open-immersion}；
\item 拟紧浸入，见“下降”引理
\ref{descent-lemma-descending-property-quasi-compact-immersion}；
\end{enumerate}
\item 拟紧，见“下降”引理 \ref{descent-lemma-descending-property-quasi-compact}；
\item 泛闭，见“下降”引理
\ref{descent-lemma-descending-property-universally-closed}；
\item （拟）分离，见“下降”引理
\ref{descent-lemma-descending-property-quasi-separated} 以及
\ref{descent-lemma-descending-property-separated}；
\item 单态射，见“下降”引理 \ref{descent-lemma-descending-property-monomorphism}；
\item 满射，见“下降”引理 \ref{descent-lemma-descending-property-surjective}；
\item 泛单射，见“下降”引理
\ref{descent-lemma-descending-property-universally-injective}；
\item 仿射，见“下降”引理 \ref{descent-lemma-descending-property-affine}；
\item 拟仿射，见“下降”引理 \ref{descent-lemma-descending-property-quasi-affine}；
\item （局部）有限型，见“下降”引理
\ref{descent-lemma-descending-property-locally-finite-type} 以及
\ref{descent-lemma-descending-property-finite-type}；
\item （局部）拟有限，见“下降”引理
\ref{descent-lemma-descending-property-quasi-finite}；
\item （局部）有限表示，见“下降”引理
\ref{descent-lemma-descending-property-locally-finite-presentation} 以及
\ref{descent-lemma-descending-property-finite-presentation}；
\item 相对维数为 $d$ 的局部有限型，见“下降”引理
\ref{descent-lemma-descending-property-relative-dimension-d}；
\item 泛开，见“下降”引理 \ref{descent-lemma-descending-property-universally-open}；
\item 平坦，见“下降”引理 \ref{descent-lemma-descending-property-flat}；
\item syntomic，见“下降”引理 \ref{descent-lemma-descending-property-syntomic}；
\item 平滑，见“下降”引理 \ref{descent-lemma-descending-property-smooth}；
\item 非分歧（相应地，G-非分歧），见“下降”引理
\ref{descent-lemma-descending-property-unramified}；
\item \'etale，见“下降”引理 \ref{descent-lemma-descending-property-etale}；
\item 固有，见“下降”引理 \ref{descent-lemma-descending-property-proper}；
\item 有限或整，见“下降”引理 \ref{descent-lemma-descending-property-finite}；
\item 有限局部自由，见“下降”引理
\ref{descent-lemma-descending-property-finite-locally-free}；
\item 泛商，见“下降”引理
\ref{descent-lemma-descending-property-universally-submersive}；
\item 泛同胚，见“下降”引理
\ref{descent-lemma-descending-property-universal-homeomorphism}。
\end{enumerate}
注意，“浸入”这一性质未必是基上的 fpqc 局部性质；
但我们已在“下降”引理
\ref{descent-lemma-descending-fppf-property-immersion} 中证明，
它是基上的 fppf 局部性质。
\end{remark}








\section{预层可表示态射的性质}
\label{spaces-section-representable-properties}

\noindent
下面的定义使这一方法得以成立。

\begin{definition}
\label{spaces-definition-relative-representable-property}
取如上所述的 $S$，并设 $a : F \to G$ 可表示。
设 $\mathcal{P}$ 是概形态射的一个性质，满足
\begin{enumerate}
\item 在任意基变换下保持，见“概形”定义
\ref{schemes-definition-preserved-by-base-change}；并且
\item 是基上的 fppf 局部性质，见“下降”定义
\ref{descent-definition-property-morphisms-local}。
\end{enumerate}
此时，若对每个 $U \in \Ob((\Sch/S)_{fppf})$ 和任意
$\xi \in G(U)$，所得概形态射 $V_\xi \to U$ 都具有性质
$\mathcal{P}$，则称 $a$ {\it 具有性质 $\mathcal{P}$}。
\end{definition}

\noindent
需要强调的是，我们只将这个定义用于在基变换下稳定、
且在基上关于 fppf 拓扑为局部的态射性质。
这并不是因为其他情形下该定义没有意义；而是因为对于所考虑的性质，
我们可能希望另给一个更合适的定义。

\begin{remark}
\label{spaces-remark-warning}
考虑性质 $\mathcal{P}=$“满射”。此时，若说
“设 $F \to G$ 是满映射”，可能会产生歧义。
我们可能指上面定义 \ref{spaces-definition-relative-representable-property}
给出的概念，也可能指预层的满映射，见“位点”定义
\ref{sites-definition-presheaves-injective-surjective}；
若 $F$ 与 $G$ 都是层，还可能指层的满映射，见“位点”定义
\ref{sites-definition-sheaves-injective-surjective}。
讨论代数空间的态射时，除非另有说明，我们总是取第一种含义。
引理 \ref{spaces-lemma-surjective-flat-locally-finite-presentation}
给出一个这里的满性蕴含作为层映射的满性的情形。
\end{remark}

\noindent
下面作一个基本检验。

\begin{lemma}
\label{spaces-lemma-morphism-schemes-gives-representable-transformation-property}
设 $S$、$X$、$Y$ 是 $\Sch_{fppf}$ 的对象，
$f : X \to Y$ 是概形态射，并设 $\mathcal{P}$ 如定义
\ref{spaces-definition-relative-representable-property} 所述。
则 $h_X \longrightarrow h_Y$ 具有性质 $\mathcal{P}$，
当且仅当 $f$ 具有性质 $\mathcal{P}$。
\end{lemma}

\begin{proof}
由引理 \ref{spaces-lemma-morphism-schemes-gives-representable-transformation}，
本引理的陈述是有意义的。证明从略。
\end{proof}

\begin{lemma}
\label{spaces-lemma-composition-representable-transformations-property}
设 $S$ 是包含于 $\Sch_{fppf}$ 的概形，并设
$F, G, H : (\Sch/S)_{fppf}^{opp} \to \textit{Sets}$。
设 $\mathcal{P}$ 是定义 \ref{spaces-definition-relative-representable-property}
所述且对复合稳定的性质。设 $a : F \to G$、$b : G \to H$
是可表示的函子变换。若 $a$ 与 $b$ 具有性质 $\mathcal{P}$，则
$b \circ a : F \longrightarrow H$ 也具有该性质。
\end{lemma}

\begin{proof}
由引理 \ref{spaces-lemma-composition-representable-transformations}，
本引理的陈述是有意义的。证明从略。
\end{proof}

\begin{lemma}
\label{spaces-lemma-base-change-representable-transformations-property}
设 $S$ 是包含于 $\Sch_{fppf}$ 的概形，并设
$F, G, H : (\Sch/S)_{fppf}^{opp} \to \textit{Sets}$。
设 $\mathcal{P}$ 是定义 \ref{spaces-definition-relative-representable-property}
所述的性质。设 $a : F \to G$ 是可表示的函子变换，
$b : H \to G$ 是任意函子变换。考虑纤维积图表
$$
\xymatrix{
H \times_{b, G, a} F \ar[r]_-{b'} \ar[d]_{a'} & F \ar[d]^a \\
H \ar[r]^b & G
}
$$
若 $a$ 具有性质 $\mathcal{P}$，则基变换 $a'$ 也具有性质
$\mathcal{P}$。
\end{lemma}

\begin{proof}
由引理 \ref{spaces-lemma-base-change-representable-transformations}，
本引理的陈述是有意义的。证明从略。
\end{proof}

\begin{lemma}
\label{spaces-lemma-descent-representable-transformations-property}
设 $S$ 是包含于 $\Sch_{fppf}$ 的概形，并设
$F, G, H : (\Sch/S)_{fppf}^{opp} \to \textit{Sets}$。
设 $\mathcal{P}$ 是定义 \ref{spaces-definition-relative-representable-property}
所述的性质。设 $a : F \to G$ 是可表示的函子变换，
$b : H \to G$ 是任意函子变换。考虑纤维积图表
$$
\xymatrix{
H \times_{b, G, a} F \ar[r]_-{b'} \ar[d]_{a'} & F \ar[d]^a \\
H \ar[r]^b & G
}
$$
假设 $b$ 诱导 fppf 层的满映射 $H^\# \to G^\#$。
此时，若 $a'$ 具有性质 $\mathcal{P}$，则 $a$ 也具有性质
$\mathcal{P}$。
\end{lemma}

\begin{proof}
首先，由引理 \ref{spaces-lemma-base-change-representable-transformations}，
变换 $a'$ 可表示。取 $U \in \Ob((\Sch/S)_{fppf})$ 以及
$\xi \in G(U)$。由假设，存在一个 fppf 覆盖
$\{U_i \to U\}_{i \in I}$ 和元素 $\xi_i \in H(U_i)$，
使其经 $b$ 映到 $\xi|_U$。由一般范畴论可知，对每个 $i$，
都有纤维积图表
$$
\xymatrix{
U_i \times_{\xi_i, H, a'} (H \times_{b, G, a} F) \ar[r] \ar[d] &
U \times_{\xi, G, a} F \ar[d] \\
U_i \ar[r] & U
}
$$
由假设，左侧竖直箭头是具有性质 $\mathcal{P}$ 的概形态射。
由于 $\mathcal{P}$ 关于 fppf 拓扑是局部的，右侧竖直箭头也具有
性质 $\mathcal{P}$，这正是所需结论。
\end{proof}

\begin{lemma}
\label{spaces-lemma-product-representable-transformations-property}
设 $S$ 是包含于 $\Sch_{fppf}$ 的概形。设
$F_i, G_i : (\Sch/S)_{fppf}^{opp} \to \textit{Sets}$，
$i = 1, 2$。设 $a_i : F_i \to G_i$，$i = 1, 2$，
是可表示的函子变换。设 $\mathcal{P}$ 是定义
\ref{spaces-definition-relative-representable-property} 所述且对复合稳定的性质。
若 $a_1$ 与 $a_2$ 具有性质 $\mathcal{P}$，则
$a_1 \times a_2 : F_1 \times F_2 \longrightarrow G_1 \times G_2$
也具有该性质。
\end{lemma}

\begin{proof}
由引理 \ref{spaces-lemma-product-representable-transformations}，
本引理的陈述是有意义的。证明从略。
\end{proof}

\begin{lemma}
\label{spaces-lemma-representable-transformations-property-implication}
设 $S$ 是包含于 $\Sch_{fppf}$ 的概形，并设
$F, G : (\Sch/S)_{fppf}^{opp} \to \textit{Sets}$。
设 $a : F \to G$ 是可表示的函子变换，$\mathcal{P}$、
$\mathcal{P}'$ 是定义 \ref{spaces-definition-relative-representable-property}
所述的性质。假设对任意概形态射 $f : X \to Y$ 都有
$\mathcal{P}(f) \Rightarrow \mathcal{P}'(f)$。
若 $a$ 具有性质 $\mathcal{P}$，则 $a$ 具有性质 $\mathcal{P}'$。
\end{lemma}

\begin{proof}
形式可得。
\end{proof}

\begin{lemma}
\label{spaces-lemma-surjective-flat-locally-finite-presentation}
设 $S$ 是概形，$F, G : (\Sch/S)_{fppf}^{opp} \to \textit{Sets}$
是层。设 $a : F \to G$ 可表示、平坦、局部有限表示且为满射。
则 $a : F \to G$ 作为层的映射是满射。
\end{lemma}

\begin{proof}
设 $T$ 是 $S$ 上的概形，$g : T \to G$ 是 $G$ 的一个 $T$-值点。
由假设，$T' = F \times_G T$（由）某个概形（表示），且态射
$T' \to T$ 平坦、局部有限表示并为满射。因此
$\{T' \to T\}$ 是一个 fppf 覆盖，并且 $g|_{T'} \in G(T')$
来自 $F(T')$ 的元素，即映射 $T' \to F$。
这证明该映射作为层的映射是满射；见“位点”定义
\ref{sites-definition-sheaves-injective-surjective}。
\end{proof}

\noindent
下面刻画对角态射可表示的那些函子。

\begin{lemma}
\label{spaces-lemma-representable-diagonal}
设 $S$ 是包含于 $\Sch_{fppf}$ 的概形，$F$ 是
$(\Sch/S)_{fppf}$ 上的集合预层。下列条件等价：
\begin{enumerate}
\item 对角态射 $F \to F \times F$ 可表示；
\item 对 $U \in \Ob((\Sch/S)_{fppf})$ 及任意 $a \in F(U)$，
映射 $a : h_U \to F$ 可表示；
\item 对每一对 $U, V \in \Ob((\Sch/S)_{fppf})$ 及任意
$a \in F(U)$、$b \in F(V)$，纤维积
$h_U \times_{a, F, b} h_V$ 可表示。
\end{enumerate}
\end{lemma}

\begin{proof}
这完全是形式的；见“范畴”引理
\ref{categories-lemma-representable-diagonal}。
它只依赖于范畴 $(\Sch/S)_{fppf}$ 具有对象对的积与纤维积这一事实；
见“拓扑”引理 \ref{topologies-lemma-fibre-products-fppf}。
\end{proof}

\noindent
在该引理的情形下，对于其中任意态射 $\xi : h_U \to F$，
以及定义 \ref{spaces-definition-relative-representable-property} 所述的任意性质，
说 $\xi$ 具有性质 $\mathcal{P}$ 是有意义的。
特别地，这适用于 $\mathcal{P}=$“满射”以及
$\mathcal{P}=$“\'etale”；见上面的评注
\ref{spaces-remark-list-properties-fpqc-local-base}。
下面定义代数空间时将使用这一评注。

\begin{lemma}
\label{spaces-lemma-transformation-diagonal-properties}
设 $S$ 是包含于 $\Sch_{fppf}$ 的概形，$F$ 是
$(\Sch/S)_{fppf}$ 上的集合预层，并设 $\mathcal{P}$ 是定义
\ref{spaces-definition-relative-representable-property} 所述的性质。
若对每个 $U, V \in \Ob((\Sch/S)_{fppf})$、$a \in F(U)$ 及
$b \in F(V)$，都有
\begin{enumerate}
\item $h_U \times_{a, F, b} h_V$ 可表示，设其由概形 $W$ 表示；并且
\item 与 $h_U \times_{a, F, b} h_V \to h_U \times h_V$
对应的态射 $W \to U \times_S V$ 具有性质 $\mathcal{P}$，
\end{enumerate}
则 $\Delta : F \to F \times F$ 可表示并具有性质 $\mathcal{P}$。
\end{lemma}

\begin{proof}
由引理 \ref{spaces-lemma-representable-diagonal}，$\Delta$ 可表示。
条件 (2) 可表述为变换
$h_U \times_{a, F, b} h_V \to h_{U \times_S V}$
具有性质 $\mathcal{P}$；见引理
\ref{spaces-lemma-morphism-schemes-gives-representable-transformation-property}。
取 $T \in \Ob((\Sch/S)_{fppf})$ 及 $(a, b) \in (F \times F)(T)$。
我们有交换图表
$$
\xymatrix{
F \times_{\Delta, F \times F, (a, b)} h_T \ar[d] \ar[r] &
h_T \ar[d]^{\Delta_{T/S}} \\
h_T \times_{a, F, b} h_T \ar[r] \ar[d] &
h_{T \times_S T} \ar[d]^{(a, b)} \\
F \ar[r]^\Delta & F \times F
}
$$
其中两个方块均为 Cartesian 方块。由此可见，态射
$F \times_{F \times F} h_T \to h_T$ 是一个具有性质
$\mathcal{P}$ 的态射沿 $\Delta_{T/S}$ 的基变换。
由于 $\mathcal{P}$ 在基变换下保持，证明完成。
\end{proof}


















\section{代数空间}
\label{spaces-section-algebraic-spaces}

\noindent
定义如下。

\begin{definition}
\label{spaces-definition-algebraic-space}
设 $S$ 是包含于 $\Sch_{fppf}$ 的概形。一个{\it $S$ 上的代数空间}
是预层
$$
F : (\Sch/S)^{opp}_{fppf} \longrightarrow \textit{Sets}
$$
并满足下列性质：
\begin{enumerate}
\item 预层 $F$ 是层。
\item 对角态射 $F  \to F \times F$ 可表示。
\item 存在概形 $U \in \Ob((\Sch/S)_{fppf})$ 以及满且 \'etale 的映射
$h_U \to F$\footnote{见引理 \ref{spaces-lemma-representable-diagonal}、
定义 \ref{spaces-definition-relative-representable-property} 以及
评注 \ref{spaces-remark-warning}。}。
\end{enumerate}
\end{definition}

\noindent
这与“通常”的定义（例如 Knutson 专著 \cite{Kn} 中的定义）
有两处不同。

\medskip\noindent
第一处是我们要求 $F$ 为 fppf 拓扑下的层。
这样做的一个理由是，代数空间的许多自然例子都满足 fppf 覆盖
（甚至 fpqc 覆盖）的层条件。此外，代数空间之所以如此有用，
部分原因在于 Michael Artin 关于代数空间的结果；
他的方法内置了一个条件，保证所得对象在 $S$ 上局部有限表示。
综合来看，fppf 拓扑似乎是自然的工作拓扑。
最终所得代数空间范畴仍然相同；见“自举”第
\ref{bootstrap-section-spaces-etale} 节。

\medskip\noindent
第二处是我们只要求 $F$ 的对角映射可表示，而 \cite{Kn}
还要求它拟紧。若对某个 $S$ 上的概形 $U$ 有 $F = h_U$，
这对应于 $U$ 拟分离这一条件。我们的观点是：先尝试仅在 $F$
的对角态射可表示的假设下证明后续若干结果，只在必要处增加额外假设。
无论如何，这带来一个令人满意的推论：下面的引理成立。

\begin{lemma}
\label{spaces-lemma-scheme-is-space}
概形是代数空间。更准确地说，给定概形
$T \in \Ob((\Sch/S)_{fppf})$，可表示函子 $h_T$ 是代数空间。
\end{lemma}

\begin{proof}
由第 \ref{spaces-section-general} 节的说明，函子 $h_T$ 是层。
由于 $(\Sch/S)_{fppf}$ 具有纤维积，对角态射
$h_T \to h_T \times h_T = h_{T \times T}$ 可表示。
恒等映射 $h_T \to h_T$ 满且 \'etale。
\end{proof}

\begin{definition}
\label{spaces-definition-morphism-algebraic-spaces}
设 $F$、$F'$ 是 $S$ 上的代数空间。一个{\it $S$ 上代数空间的态射
$f : F \to F'$}，是从 $F$ 到 $F'$ 的函子变换。
\end{definition}

\noindent
通过 Yoneda 嵌入 $T/S \mapsto h_T$，$S$ 上代数空间的范畴
以 $(\Sch/S)_{fppf}$ 为满子范畴。此后，我们不再区分概形 $T/S$
与它所表示的代数空间。因此，当我们说“设 $f : T \to F$
是从概形 $T$ 到代数空间 $F$ 的态射”时，意思是
$T \in \Ob((\Sch/S)_{fppf})$，$F$ 是 $S$ 上的代数空间，
而 $f : h_T \to F$ 是 $S$ 上代数空间的态射。






\section{代数空间的纤维积}
\label{spaces-section-fibre-products}

\noindent
$S$ 上代数空间的范畴既有积，也有纤维积。

\begin{lemma}
\label{spaces-lemma-product-spaces}
设 $S$ 是包含于 $\Sch_{fppf}$ 的概形，$F,G$ 是 $S$ 上的代数空间。
则 $F \times G$ 是代数空间，并且是 $S$ 上代数空间范畴中的积。
\end{lemma}

\begin{proof}
显然 $H = F \times G$ 是层。$H$ 的对角态射就是 $F$ 与 $G$
的对角态射之积，故由引理
\ref{spaces-lemma-product-representable-transformations} 可表示。
最后，若 $U \to F$ 与 $V \to G$ 是满 \'etale 态射，且
$U, V \in \Ob((\Sch/S)_{fppf})$，则由引理
\ref{spaces-lemma-product-representable-transformations-property}，
$U \times V \to F \times G$ 满且 \'etale。
\end{proof}

\begin{lemma}
\label{spaces-lemma-fibre-product-spaces-over-sheaf-with-representable-diagonal}
设 $S$ 是包含于 $\Sch_{fppf}$ 的概形。设 $H$ 是
$(\Sch/S)_{fppf}$ 上对角态射可表示的层，$F,G$ 是 $S$ 上的代数空间，
而 $F \to H$、$G \to H$ 是层的映射。则 $F \times_H G$ 是代数空间。
\end{lemma}

\begin{proof}
我们验证定义 \ref{spaces-definition-algebraic-space} 的三个条件。
层的纤维积是层，故 $F \times_H G$ 是层。
$F \times_H G$ 的对角态射是下图中的左侧竖直箭头：
$$
\xymatrix{
F \times_H G \ar[r] \ar[d]_\Delta &
F \times G \ar[d]^{\Delta_F \times \Delta_G} \\
(F \times F) \times_{(H \times H)} (G \times G) \ar[r] &
(F \times F) \times (G \times G)
}
$$
该图为 Cartesian 方块。右侧态射可表示，故其基变换 $\Delta$
也可表示；见引理 \ref{spaces-lemma-product-representable-transformations} 以及
\ref{spaces-lemma-base-change-representable-transformations}。
最后，取 $U, V \in \Ob((\Sch/S)_{fppf})$，并设
$a : U \to F$、$b : V \to G$ 满且 \'etale。
由于 $\Delta_H$ 可表示，$U \times_H V$ 是概形。态射
$$
U \times_H V \longrightarrow F \times_H G
$$
是两个基变换 $U \times_H V \to U \times_H G$ 与
$U \times_H G \to F \times_H G$ 的复合；它们分别来自满 \'etale 态射
$U \to F$ 与 $V \to G$，因而该态射满且 \'etale。见引理
\ref{spaces-lemma-composition-representable-transformations} 以及
\ref{spaces-lemma-base-change-representable-transformations}。
这证明定义 \ref{spaces-definition-algebraic-space} 的最后一个条件成立，
从而 $F \times_H G$ 是代数空间。
\end{proof}

\begin{lemma}
\label{spaces-lemma-fibre-product-spaces}
设 $S$ 是包含于 $\Sch_{fppf}$ 的概形，$F \to H$、$G \to H$
是 $S$ 上代数空间的态射。则 $F \times_H G$ 是代数空间，
并且是 $S$ 上代数空间范畴中的纤维积。
\end{lemma}

\begin{proof}
由更强的引理
\ref{spaces-lemma-fibre-product-spaces-over-sheaf-with-representable-diagonal}，
$F \times_H G$ 是代数空间。又因为 $S$ 上代数空间的范畴是
$(\Sch/S)_{fppf}$ 上集合（预）层范畴的满子范畴，显然
$F \times_H G$ 是该代数空间范畴中的纤维积。
\end{proof}






\section{代数空间的黏合}
\label{spaces-section-glueing-algebraic-spaces}

\noindent
从本节开始，我们将真正滥用记号，不再区分概形与其表示的空间。

\begin{lemma}
\label{spaces-lemma-coproduct-sheaves-open-and-closed}
设 $S \in \Ob(\Sch_{fppf})$，$F$ 与 $G$ 是
$(\Sch/S)_{fppf}^{opp}$ 上的层，并以 $F \amalg G$ 表示层范畴中的余积。
映射 $F \to F \amalg G$ 由开浸入兼闭浸入表示。
\end{lemma}

\begin{proof}
设 $U$ 是概形，$\xi \in (F \amalg G)(U)$。回忆，层范畴中的余积
是预层余积的层化（“位点”引理 \ref{sites-lemma-colimit-sheaves}）。
因此，存在 fppf 覆盖 $\{g_i : U_i \to U\}_{i \in I}$ 和不交并分解
$I = I' \amalg I''$，使得 $U_i \to U \to F \amalg G$ 经过 $F$
（相应地，$G$）分解，当且仅当 $i \in I'$（相应地，$i \in I''$）。
由于 $F$ 与 $G$ 在 $F \amalg G$ 中的交为空，若 $i \in I'$ 且
$j \in I''$，则 $U_i \times_U U_j$ 为空。因此
$U' = \bigcup_{i \in I'} g_i(U_i)$ 与
$U'' = \bigcup_{i \in I''} g_i(U_i)$ 是 $U$ 的不交开子概形
（“态射”引理 \ref{morphisms-lemma-fppf-open}），并且
$U = U' \amalg U''$。我们略去 $U' = U \times_{F \amalg G} F$
的验证。
\end{proof}

\begin{lemma}
\label{spaces-lemma-representable-sheaf-coproduct-sheaves}
设 $S \in \Ob(\Sch_{fppf})$、$U \in \Ob((\Sch/S)_{fppf})$。
给定集合 $I$ 以及 $\Ob((\Sch/S)_{fppf})$ 上的层 $F_i$，
若作为层有 $U \cong \coprod_{i\in I} F_i$，则每个 $F_i$
均由某个开闭子概形 $U_i$ 表示，并且作为概形有
$U \cong \coprod U_i$。
\end{lemma}

\begin{proof}
由引理 \ref{spaces-lemma-coproduct-sheaves-open-and-closed}，映射
$F_i \to U$ 由开浸入兼闭浸入表示。因此 $F_i$ 由 $U$ 的某个开闭子概形
$U_i$ 表示。由于作为层有 $U \cong \coprod F_i$，而该等式可在点上检验，
故 $U = \coprod U_i$。
\end{proof}

\begin{lemma}
\label{spaces-lemma-algebraic-space-coproduct-sheaves}
设 $S \in \Ob(\Sch_{fppf})$，$F$ 是 $S$ 上的代数空间。
给定集合 $I$ 以及 $\Ob((\Sch/S)_{fppf})$ 上的层 $F_i$，
若作为层有 $F \cong \coprod_{i\in I} F_i$，则每个 $F_i$
都是 $S$ 上的代数空间。
\end{lemma}

\begin{proof}
$F \to F \times F$ 的可表示性蕴含每个对角态射
$F_i \to F_i \times F_i$ 可表示（由定义以及等式
$F \times_{(F \times F)} (F_i \times F_i) = F_i$ 立即可得）。
在 $(\Sch/S)_{fppf}$ 中选取概形 $U$ 以及满 \'etale 态射
$U \to F$（由假设存在）。由引理
\ref{spaces-lemma-base-change-representable-transformations-property}，
基变换 $U \times_F F_i \to F_i$ 满且 \'etale。另一方面，由引理
\ref{spaces-lemma-coproduct-sheaves-open-and-closed}，$U \times_F F_i$ 是概形。
因此，我们已验证定义 \ref{spaces-definition-algebraic-space} 的全部条件，
$F_i$ 是代数空间。
\end{proof}

\noindent
不关心集合论问题的读者可以忽略下面引理中关于 $I$ 与 $F_i$ 大小的条件。

\begin{lemma}
\label{spaces-lemma-coproduct-algebraic-spaces}
设 $S \in \Ob(\Sch_{fppf})$。给定集合 $I$ 和代数空间 $F_i$，
$i \in I$。若 $I$ 与 $F_i$ 不太“大”，则
$F = \coprod_{i \in I} F_i$ 是代数空间。例如，若可选取满 \'etale 态射
$U_i \to F_i$，使 $\coprod_{i \in I} U_i$ 同构于
$(\Sch/S)_{fppf}$ 的一个对象，则 $F$ 是代数空间。
\end{lemma}

\begin{proof}
由构造，$F$ 是层。我们略去 $F$ 的对角态射可表示的验证。
最后，若 $U$ 是 $(\Sch/S)_{fppf}$ 中与
$\coprod_{i \in I} U_i$ 同构的对象，则容易验证所得映射
$U \to \coprod F_i$ 满且 \'etale。
\end{proof}

\noindent
下面是“概形”引理 \ref{schemes-lemma-glue-functors} 的对应结果。

\begin{lemma}
\label{spaces-lemma-glueing-algebraic-spaces}
设 $S \in \Ob(\Sch_{fppf})$，$F$ 是 $(\Sch/S)_{fppf}$ 上的集合预层。
假设
\begin{enumerate}
\item $F$ 是层；
\item 存在指标集 $I$ 以及子函子 $F_i \subset F$，使得
\begin{enumerate}
\item 每个 $F_i$ 都是代数空间；
\item 每个 $F_i \to F$ 都可表示；
\item 每个 $F_i \to F$ 都是开浸入（见定义
\ref{spaces-definition-relative-representable-property}）；
\item 映射 $\coprod F_i \to F$ 作为层映射是满射；并且
\item $\coprod F_i$ 是代数空间（集合论条件；见引理
\ref{spaces-lemma-coproduct-algebraic-spaces}）。
\end{enumerate}
\end{enumerate}
则 $F$ 是代数空间。
\end{lemma}

\begin{proof}
设 $T$ 是 $(\Sch/S)_{fppf}$ 的对象，$T \to F$ 是态射。
由假设 (2)(b) 与 (2)(c)，纤维积 $F_i \times_F T$
由开子概形 $V_i \subset T$ 表示。因此
$(\coprod F_i) \times_F T$ 由 $T$ 上的概形 $\coprod V_i$ 表示。
由假设 (2)(d)，存在 fppf 覆盖 $\{T_j \to T\}_{j \in J}$，
使 $T_j \to T \to F$ 经过 $F_i$ 分解，其中 $i = i(j)$。
故 $T_j \to T$ 经过开子概形 $V_{i(j)} \subset T$ 分解。
由于 $\{T_j \to T\}$ 联合满，$T = \bigcup V_i$ 是开覆盖。
特别地，函子变换 $\coprod F_i \to F$ 可表示，并且在定义
\ref{spaces-definition-relative-representable-property} 的意义下为满射
（相关讨论见评注 \ref{spaces-remark-warning}）。

\medskip\noindent
其次，设 $T' \to F$ 是来自 $(\Sch/S)_{fppf}$ 中另一对象的态射。
如上写成 $T' = \bigcup V'_i$，其中 $V'_i = T' \times_F F_i$。
为证明对角态射 $F \to F \times F$ 可表示，只需证明
$G = T \times_F T'$ 可表示；见引理 \ref{spaces-lemma-representable-diagonal}。
考虑子函子 $G_i = G \times_F F_i$。注意
$G_i = V_i \times_{F_i} V'_i$；由于 $F_i$ 是代数空间，故其可表示。
由上所述，$G_i$ 构成 $G$ 的 Zariski 覆盖。因此，由“概形”引理
\ref{schemes-lemma-glue-functors}，$G$ 可表示。

\medskip\noindent
选取概形 $U \in \Ob((\Sch/S)_{fppf})$ 和满 \'etale 态射
$U \to \coprod F_i$（由假设存在）。可写 $U = \coprod U_i$，
其中 $U_i$ 是 $F_i$ 的逆像；见引理
\ref{spaces-lemma-representable-sheaf-coproduct-sheaves}。
我们断言 $U \to F$ 满且 \'etale。由于 $\coprod F_i \to F$ 为满射
（见证明的第一段），应用引理
\ref{spaces-lemma-composition-representable-transformations-property} 即得满性。
考虑上述 $T \to F$ 对应的纤维积 $U \times_F T$。
我们必须证明 $U \times_F T \to T$ 为 \'etale。
由于 $U \times_F T = \coprod U_i \times_F T$，只需证明每个
$U_i \times_F T \to T$ 为 \'etale。又因
$U_i \times_F T = U_i \times_{F_i} V_i$，这由 $U_i \to F_i$
为 \'etale、$V_i \to T$ 为开浸入以及“态射”引理
\ref{morphisms-lemma-open-immersion-etale} 和
\ref{morphisms-lemma-composition-etale} 得出。
\end{proof}














\section{代数空间的表示}
\label{spaces-section-presentations}

\noindent
给定一个代数空间，我们可以找到它的一个“表示”。


\begin{lemma}
\label{spaces-lemma-space-presentation}
设 $F$ 是 $S$ 上的代数空间，$f : U \to F$ 是从概形到 $F$
的满 \'etale 态射。令 $R = U \times_F U$。则
\begin{enumerate}
\item $j : R \to U \times_S U$ 定义 $U$ 上相对于 $S$ 的等价关系
（见“群胚”定义 \ref{groupoids-definition-equivalence-relation}）；
\item 态射 $s, t : R \to U$ 为 \'etale；并且
\item 图表
$$
\xymatrix{
R \ar@<1ex>[r] \ar@<-1ex>[r] &
U \ar[r] &
F
}
$$
是 $\Sh((\Sch/S)_{fppf})$ 中的余等化子图表。
\end{enumerate}
\end{lemma}

\begin{proof}
设 $T/S$ 是 $(\Sch/S)_{fppf}$ 的对象。则
$R(T) = \{(a, b) \in U(T) \times U(T) \mid f \circ a = f \circ b\}$，
它在 $U(T)$ 上定义等价关系。由于态射 $U \to F$ 为 \'etale，
态射 $s, t : R \to U$ 也为 \'etale。

\medskip\noindent
为证明 (3)，先证明 $U \to F$ 是层的满射；见“位点”定义
\ref{sites-definition-sheaves-injective-surjective}。
对上述 $T$，取 $\xi \in F(T)$，并令 $V = T \times_{\xi, F, f}U$。
由假设，$V$ 是概形且 $V \to T$ 满而 \'etale。因此
$\{V \to T\}$ 是 fppf 拓扑的覆盖。由构造，$\xi|_V$ 经过 $U$ 分解，
故 $U \to F$ 为满射。由“位点”引理
\ref{sites-lemma-coequalizer-surjection}，满性蕴含 $F$ 是该图表的余等化子。
\end{proof}

\noindent
该引理提示如下定义。

\begin{definition}
\label{spaces-definition-etale-equivalence-relation}
设 $S$ 是概形，$U$ 是 $S$ 上的概形。$U$ 上相对于 $S$
的一个{\it \'etale 等价关系}，是等价关系
$j : R \to U \times_S U$，其中 $s, t : R \to U$ 是概形的
\'etale 态射。
\end{definition}

\begin{definition}
\label{spaces-definition-presentation}
设 $F$ 是 $S$ 上的代数空间。$F$ 的一个{\it 表示}由以下数据给出：
$S$ 上的概形 $U$、$U$ 上相对于 $S$ 的 \'etale 等价关系 $R$，
以及满足 $R = U \times_F U$ 的满 \'etale 态射 $U \to F$。
\end{definition}

\noindent
等价地，我们可以要求存在同构
$$
U/R \cong F
$$
其中商 $U/R$ 按“群胚”第 \ref{groupoids-section-quotient-sheaves} 节定义。
为构造代数空间，我们将研究逆向问题：对于哪些等价关系，商层 $U/R$
是代数空间。最终将证明，只要 $R$ 是 $U$ 上相对于 $S$ 的
\'etale 等价关系，情形总是如此；见定理 \ref{spaces-theorem-presentation}。





























\section{代数空间与等价关系}
\label{spaces-section-spaces-from-equivalence-relations}

\noindent
设给定 $S$ 上的概形 $U$，以及 $U$ 上相对于 $S$ 的
\'etale 等价关系 $R$。我们希望证明这定义了一个代数空间。
下面将给出一系列引理，证明商层 $U/R$
（见“群胚”定义 \ref{groupoids-definition-quotient-sheaf}）
具备定义 \ref{spaces-definition-algebraic-space} 所要求的全部性质。

\begin{lemma}
\label{spaces-lemma-pullback-etale-equivalence-relation}
设 $S$ 是概形，$U$ 是 $S$ 上的概形，并设
$j = (s, t) : R \to U \times_S U$ 是 $U$ 上相对于 $S$ 的
\'etale 等价关系。设 $U' \to U$ 为 \'etale 态射，$R'$ 为 $R$
在 $U'$ 上的限制；见“群胚”定义
\ref{groupoids-definition-restrict-relation}。
则 $j' : R' \to U' \times_S U'$ 也是 \'etale 等价关系。
\end{lemma}

\begin{proof}
由“群胚”引理 \ref{groupoids-lemma-restrict-groupoid}
中对 $s',t'$ 的描述可见，$s' , t' : R' \to U'$ 是 \'etale 态射，
因为它们是 \'etale 态射的基变换之复合
（见“态射”引理 \ref{morphisms-lemma-base-change-etale} 以及
\ref{morphisms-lemma-composition-etale}）。
\end{proof}

\noindent
我们将经常使用下面的引理寻找代数空间的开子空间。
该引理的一个略微加强版（假设更一般）是“自举”引理
\ref{bootstrap-lemma-better-finding-opens}。

\begin{lemma}
\label{spaces-lemma-finding-opens}
设 $S$ 是概形，$U$ 是 $S$ 上的概形，
$j = (s, t) : R \to U \times_S U$ 是预关系，
$g : U' \to U$ 是态射。假设
\begin{enumerate}
\item $j$ 是等价关系；
\item $s, t : R \to U$ 满、平坦且局部有限表示；
\item $g$ 平坦且局部有限表示。
\end{enumerate}
令 $R' = R|_{U'}$ 为 $R$ 在 $U'$ 上的限制。则
$U'/R' \to U/R$ 可表示，并且是开浸入。
\end{lemma}

\begin{proof}
由“群胚”引理 \ref{groupoids-lemma-restrict-relation}，态射
$j' = (s', t') : R' \to U' \times_S U'$ 定义等价关系。
由于 $g$ 平坦且局部有限表示，$g$ 还是泛开的
（“态射”引理 \ref{morphisms-lemma-fppf-open}）。同理，$s,t$
也泛开。令 $W^1 = g(U') \subset U$，并令
$W = t(s^{-1}(W^1))$。则 $W^1$ 与 $W$ 是 $U$ 中的开集。
此外，由于 $j$ 是等价关系，有 $t(s^{-1}(W)) = W$
（例如见“群胚”引理
\ref{groupoids-lemma-constructing-invariant-opens}）。

\medskip\noindent
由“群胚”引理
\ref{groupoids-lemma-quotient-pre-equivalence-relation-restrict}，
层映射 $F' = U'/R' \to F = U/R$ 是单射。
设 $a : T \to F$ 是从概形到 $U/R$ 的态射。
我们必须证明 $T \times_F F'$ 由 $T$ 的某个开子概形表示。

\medskip\noindent
态射 $a$ 由以下数据给出：$T$ 的 fppf 覆盖
$\{\varphi_j : T_j \to T\}_{j \in J}$ 以及态射
$a_j : T_j \to U$，使映射
$$
a_j \times a_{j'} :
T_j \times_T T_{j'}
\longrightarrow
U \times_S U
$$
经某些（唯一的）映射 $r_{jj'} : T_j \times_T T_{j'} \to R$
通过 $j : R \to U \times_S U$ 分解。系统 $(a_j)$ 在下述意义下对应于
$a$：图表
$$
\xymatrix{
T_j \ar[r]_{a_j} \ar[d] & U \ar[d] \\
T \ar[r]^a & F
}
$$
交换。

\medskip\noindent
考虑开子集 $W_j = a_j^{-1}(W) \subset T_j$。
由于 $t(s^{-1}(W)) = W$，可见
$$
W_j \times_T T_{j'} =
r_{jj'}^{-1}(t^{-1}(W)) = r_{jj'}^{-1}(s^{-1}(W)) =
T_j \times_T W_{j'}.
$$
由“下降”引理 \ref{descent-lemma-open-fpqc-covering}，
这意味着存在开集 $W_T \subset T$，使对所有 $j \in J$ 都有
$\varphi_j^{-1}(W_T) = W_j$。我们断言 $W_T \to T$ 表示
$T \times_F F' \to T$。

\medskip\noindent
首先证明 $W_T \to T \to F$ 是 $F'(W_T)$ 的元素。
由于 $\{W_j \to W_T\}_{j \in J}$ 是 $W_T$ 的 fppf 覆盖，
只需证明每个 $W_j \to U \to F$ 都是 $F'(W_j)$ 的元素
（因为 $F'$ 是 fppf 拓扑下的层）。考虑交换图表
$$
\xymatrix{
W'_j \ar[rr] \ar[dd] \ar[rd] & & U' \ar[d]^g \\
& s^{-1}(W^1) \ar[r]_s \ar[d]^t & W^1 \ar[d] \\
W_j \ar[r]^{a_j|_{W_j}} & W \ar[r] & F
}
$$
其中 $W'_j = W_j \times_W s^{-1}(W^1) \times_{W^1} U'$。
由于 $t$ 与 $g$ 满、平坦且局部有限表示，$W'_j \to W_j$
也具有这些性质。因此，元素 $W_j \to U \to F$ 在 $W'_j$ 上的限制
是 $F'$ 的元素，正如所需。

\medskip\noindent
设 $f : T' \to T$ 是满足 $a|_{T'} \in F'(T')$ 的概形态射。
我们必须证明 $f$ 经过开集 $W_T$ 分解。由于
$\{T' \times_T T_j \to T'\}$ 是 $T'$ 的 fppf 覆盖，
只需证明每个 $T' \times_T T_j \to T$ 都经过 $W_T$ 分解。
因此可假设对某个 $j$，$f$ 分解为
$\varphi_j \circ f_j : T' \to T_j \to T$。此时，条件
$a|_{T'} \in F'(T')$ 意味着存在某个 fppf 覆盖
$\{\psi_i : T'_i \to T'\}_{i \in I}$ 以及某些态射
$b_i : T'_i \to U'$，使得
$$
\xymatrix{
T'_i \ar[r]_{b_i} \ar[d]_{f_j \circ \psi_i} & U' \ar[r]_g & U \ar[d] \\
T_j \ar[r]^{a_j} & U \ar[r] & F
}
$$
交换。该交换性意味着存在态射 $r'_i : T'_i \to R$，使得
$t \circ r'_i = a_j \circ f_j \circ \psi_i$ 且
$s \circ r'_i = g \circ b_i$。这蕴含
$\Im(f_j \circ \psi_i) \subset W_j$，证明完成。
\end{proof}

\noindent
下面的引理看起来理应平凡，但其实并非完全平凡。

\begin{lemma}
\label{spaces-lemma-when-it-works-it-works}
设 $S$ 是概形，$U$ 是 $S$ 上的概形，并设
$j = (s, t) : R \to U \times_S U$ 是 $U$ 上相对于 $S$ 的
\'etale 等价关系。若商 $U/R$ 是代数空间，则
$U \to U/R$ 为 \'etale 满射。因此
$(U, R, U \to U/R)$ 是代数空间 $U/R$ 的一个表示。
\end{lemma}

\begin{proof}
以 $c : U \to U/R$ 表示所论态射。设 $T$ 是概形，
$a : T \to U/R$ 是态射。我们必须证明概形态射
$\pi : T \times_{a, U/R, c} U \to T$ 为 \'etale 满射。
态射 $a$ 对应于一个 fppf 覆盖 $\{\varphi_i : T_i \to T\}$
以及态射 $a_i : T_i \to U$，使
$a_i \times a_{i'} : T_i \times_T T_{i'} \to U \times_S U$
经过 $R$ 分解，并且 $c \circ a_i = a \circ \varphi_i$。因此
$$
T_i \times_{\varphi_i, T} T \times_{a, U/R, c} U =
T_i \times_{c \circ a_i, U/R, c} U =
T_i \times_{a_i, U} U \times_{c, U/R, c} U = T_i \times_{a_i, U, t} R.
$$
由于 $t$ 为 \'etale 满射，$\pi$ 到 $T_i$ 的基变换也满而 \'etale。
满且 \'etale 这一性质关于基上的 fpqc 拓扑是局部的
（见评注 \ref{spaces-remark-list-properties-fpqc-local-base}），故结论成立。
\end{proof}

\begin{lemma}
\label{spaces-lemma-presentation-quasi-compact}
设 $S$ 是概形，$U$ 是 $S$ 上的概形，并设
$j = (s, t) : R \to U \times_S U$ 是 $U$ 上相对于 $S$ 的
\'etale 等价关系。假设 $U$ 仿射。则商 $F = U/R$ 是代数空间，
且 $U \to F$ 为 \'etale 满射。
\end{lemma}

\begin{proof}
由于 $j : R \to U \times_S U$ 是单态射，$j$ 分离
（见“概形”引理 \ref{schemes-lemma-monomorphism-separated}）。
由于 $U$ 仿射，$U \times_S U$（它带有到仿射概形 $U \times U$
的单态射）是分离的。因此 $R$ 分离；特别地，态射 $s,t$
既分离又 \'etale。

\medskip\noindent
由于复合 $R \to U \times_S U \to U$ 局部有限型，
$j$ 局部有限型（见“态射”引理
\ref{morphisms-lemma-permanence-finite-type}）。又因 $j$ 是单态射，
其纤维有限；由“态射”引理 \ref{morphisms-lemma-finite-fibre}，
$j$ 局部拟有限。综上，$j$ 分离且局部拟有限。

\medskip\noindent
第一步证明商映射 $c : U \to F$ 可表示。
取概形 $T$ 和态射 $a : T \to F$。我们必须证明层
$G = T \times_{a, F, c} U$ 可表示。正如引理
\ref{spaces-lemma-finding-opens} 与 \ref{spaces-lemma-when-it-works-it-works}
的证明中所见，存在 fppf 覆盖
$\{\varphi_i : T_i \to T\}_{i \in I}$ 和态射 $a_i : T_i \to U$，
使 $a_i \times a_{i'} : T_i \times_T T_{i'} \to U \times_S U$
经过 $R$ 分解，并且 $c \circ a_i = a \circ \varphi_i$。
与引理 \ref{spaces-lemma-when-it-works-it-works} 的证明相同，有
\begin{eqnarray*}
T_i \times_{\varphi_i, T} G & = &
T_i \times_{\varphi_i, T} T \times_{a, U/R, c} U \\
& = & T_i \times_{c \circ a_i, U/R, c} U \\
& = & T_i \times_{a_i, U} U \times_{c, U/R, c} U \\
& = & T_i \times_{a_i, U, t} R
\end{eqnarray*}
由于 $t$ 分离且 \'etale，特别地，它分离且局部拟有限
（由“态射”引理 \ref{morphisms-lemma-unramified-quasi-finite} 以及
\ref{morphisms-lemma-flat-unramified-etale}），$G$ 在每个 $T_i$
上的限制都由一个分离且局部拟有限的概形态射
$X_i \to T_i$ 表示。由“下降”引理
\ref{descent-lemma-descent-data-sheaves}，我们得到相对于 fppf 覆盖
$\{T_i \to T\}$ 的下降数据 $(X_i, \varphi_{ii'})$。
由于每个 $X_i \to T_i$ 分离且局部拟有限，由“态射续论”引理
\ref{more-morphisms-lemma-separated-locally-quasi-finite-morphisms-fppf-descend}，
该下降数据有效。因此，由“下降”引理
\ref{descent-lemma-descent-data-sheaves} (2)，$G$ 如所需可表示。

\medskip\noindent
证明的第二步是说明 $U \to F$ 满且 \'etale。这由上文立即可得：
第一步中已见 $G = T \times_{a, F, c} U$ 是 $T$ 上的概形，
其基变换是满 \'etale 的概形态射 $X_i \to T_i$。因此
$G \to T$ 满且 \'etale（见评注
\ref{spaces-remark-list-properties-fpqc-local-base}）。
或者，也可在当前情形下重读引理
\ref{spaces-lemma-when-it-works-it-works} 的证明。

\medskip\noindent
第三步也是最后一步，是证明对角映射 $F \to F \times F$ 可表示。
首先注意图表
$$
\xymatrix{
R \ar[r] \ar[d]_j & F \ar[d]^\Delta \\
U \times_S U \ar[r] & F \times F
}
$$
是纤维积方块。由引理
\ref{spaces-lemma-product-representable-transformations}，态射
$U \times_S U \to F \times F$ 可表示（注意
$h_U \times h_U = h_{U \times_S U}$）。此外，由引理
\ref{spaces-lemma-product-representable-transformations-property}，态射
$U \times_S U \to F \times F$ 满且 \'etale（还要注意，\'etale 与满射
都出现在评注 \ref{spaces-remark-list-properties-fpqc-local-base} 以及
\ref{spaces-remark-list-properties-stable-composition} 的列表中）。
由引理 \ref{spaces-lemma-base-change-representable-transformations} 与上图，
或者将 $R \to F$ 写成 $R \to U \to F$ 并应用引理
\ref{spaces-lemma-morphism-schemes-gives-representable-transformation} 和
\ref{spaces-lemma-composition-representable-transformations}，可知
$R \to F$ 也可表示。设 $T$ 是概形，$a : T \to F \times F$
是态射。我们必须证明 $G = T \times_{a, F \times F, \Delta} F$
可表示。由上文，概形态射
$$
T' = (U \times_S U) \times_{F \times F, a} T \longrightarrow T
$$
满且 \'etale。因此 $\{T' \to T\}$ 是 $T$ 的 \'etale 覆盖。
还要注意
$$
T' \times_T G = T' \times_{U \times_S U, j} R
$$
这可由考察下列立方图看出：
$$
\xymatrix{
& R \ar[rr] \ar[dd] & & F \ar[dd] \\
T' \times_T G \ar[rr] \ar[dd] \ar[ru] & & G \ar[dd] \ar[ru] & \\
& U \times_S U \ar'[r][rr] & & F \times F \\
T' \ar[rr] \ar[ru] & & T \ar[ru]
}
$$
因此，$G$ 在 $T'$ 上的限制由某个概形 $X$ 表示；而且态射
$X \to T'$ 是态射 $j$ 的基变换。故 $X \to T'$ 分离且局部拟有限
（见证明的第二段）。由“下降”引理
\ref{descent-lemma-descent-data-sheaves}，得到相对于 fppf 覆盖
$\{T' \to T\}$ 的下降数据 $(X, \varphi)$。由于 $X \to T'$
分离且局部拟有限，由“态射续论”引理
\ref{more-morphisms-lemma-separated-locally-quasi-finite-morphisms-fppf-descend}，
该下降数据有效。因此，由“下降”引理
\ref{descent-lemma-descent-data-sheaves} (2)，$G$ 如所需可表示。
\end{proof}

\begin{theorem}
\label{spaces-theorem-presentation}
设 $S$ 是概形，$U$ 是 $S$ 上的概形，并设
$j = (s, t) : R \to U \times_S U$ 是 $U$ 上相对于 $S$ 的
\'etale 等价关系。则商 $U/R$ 是代数空间，且 $U \to U/R$
为 \'etale 满射；换言之，$(U, R, U \to U/R)$ 是 $U/R$ 的一个表示。
\end{theorem}

\begin{proof}
由引理 \ref{spaces-lemma-when-it-works-it-works}，只需证明 $U/R$
是代数空间。设 $U' \to U$ 为满 \'etale 态射，则
$\{U' \to U\}$ 特别是 fppf 覆盖。令 $R'$ 为 $R$ 在 $U'$ 上的限制；
见“群胚”定义 \ref{groupoids-definition-restrict-relation}。
由“群胚”引理 \ref{groupoids-lemma-quotient-groupoid-restrict}，
$U/R \cong U'/R'$。由引理
\ref{spaces-lemma-pullback-etale-equivalence-relation}，$R'$ 是 $U'$ 上的
\'etale 等价关系。因此，可用 $U'$ 代替 $U$。

\medskip\noindent
将前述说明应用于 $U' = \coprod U_i$，其中
$U = \bigcup U_i$ 是 $U$ 的仿射开覆盖。因此，不妨假设
$U = \coprod U_i$，其中每个 $U_i$ 都是仿射概形。

\medskip\noindent
考虑 $R$ 在 $U_i$ 上的限制 $R_i$。由引理
\ref{spaces-lemma-pullback-etale-equivalence-relation}，这是 \'etale 等价关系。
令 $F_i = U_i/R_i$ 且 $F = U/R$。显然 $\coprod F_i \to F$ 为满射。
由引理 \ref{spaces-lemma-finding-opens}，每个 $F_i \to F$ 都可表示且为开浸入。
将引理 \ref{spaces-lemma-presentation-quasi-compact} 应用于 $(U_i, R_i)$，
可见 $F_i$ 是代数空间。再由引理
\ref{spaces-lemma-when-it-works-it-works}，$U_i \to F_i$ 为 \'etale 满射。
由引理 \ref{spaces-lemma-coproduct-algebraic-spaces}，$\coprod F_i$ 是代数空间。
最后，我们已验证引理 \ref{spaces-lemma-glueing-algebraic-spaces} 的全部假设，
故 $F = U/R$ 是代数空间。
\end{proof}










\section{代数空间：补充整理}
\label{spaces-section-algebraic-spaces-retrofitted}

\noindent
我们开始建立处理代数空间的引理库。第一个结果说明，
定义 \ref{spaces-definition-algebraic-space} 中关于对角态射的条件可作如下弱化。

\begin{lemma}
\label{spaces-lemma-etale-locally-representable-gives-space}
设 $S$ 是包含于 $\Sch_{fppf}$ 的概形。设 $F$ 是
$(\Sch/S)_{fppf}$ 上的层，并且存在 $U \in \Ob((\Sch/S)_{fppf})$
以及可表示、满且 \'etale 的映射 $U \to F$。则 $F$ 是代数空间。
\end{lemma}

\begin{proof}
令 $R = U \times_F U$。由于假设 $U \to F$ 可表示，所得对象是概形；
由于假设 $U \to F$ 为 \'etale，投影 $s, t : R \to U$ 为 \'etale。
因 $R = U \times_F U$，映射 $j = (t, s) : R \to U \times_S U$
是单态射并定义等价关系。由定理 \ref{spaces-theorem-presentation}，
商层 $F' = U/R$ 是代数空间，且 $U \to F'$ 为 \'etale 满射。
仍因 $R = U \times_F U$，得到典范分解 $U \to F' \to F$，
并且 $F' \to F$ 是层的单射。另一方面，由引理
\ref{spaces-lemma-surjective-flat-locally-finite-presentation}，$U \to F$
作为层的映射是满射。因此 $F' \to F$ 也满，故
$F' = F$ 是代数空间。
\end{proof}

\begin{lemma}
\label{spaces-lemma-etale-locally-representable-by-space-gives-space}
设 $S$ 是包含于 $\Sch_{fppf}$ 的概形，$G$ 是 $S$ 上的代数空间，
$F$ 是 $(\Sch/S)_{fppf}$ 上的层，并设 $G \to F$ 是可表示、满且
\'etale 的函子变换。则 $F$ 是代数空间。
\end{lemma}

\begin{proof}
选取概形 $U$ 和满 \'etale 态射 $U \to G$。
由于 $G$ 是代数空间，$U \to G$ 可表示。因此复合
$U \to G \to F$ 可表示、满且 \'etale；见引理
\ref{spaces-lemma-composition-representable-transformations} 和
\ref{spaces-lemma-composition-representable-transformations-property}。
故由引理 \ref{spaces-lemma-etale-locally-representable-gives-space}，
$F$ 是代数空间。
\end{proof}

\begin{lemma}
\label{spaces-lemma-representable-over-space}
\begin{slogan}
若一个函子存在到代数空间的可表示态射，则它是代数空间。
\end{slogan}
设 $S$ 是包含于 $\Sch_{fppf}$ 的概形，$F$ 是 $S$ 上的代数空间，
$G \to F$ 是可表示的函子变换。则 $G$ 是代数空间。
\end{lemma}

\begin{proof}
由引理 \ref{spaces-lemma-representable-transformation-to-sheaf}，$G$ 是层。
图表
$$
\xymatrix{
G \times_F G \ar[r] \ar[d] & F \ar[d]^{\Delta_F} \\
G \times G \ar[r] & F \times F
}
$$
是 Cartesian 方块。因此，由引理
\ref{spaces-lemma-base-change-representable-transformations}，
$G \times_F G \to G \times G$ 可表示。由引理
\ref{spaces-lemma-representable-transformation-diagonal}，
$G \to G \times_F G$ 可表示。因此，作为可表示变换的复合，
$\Delta_G : G \to G \times G$ 可表示；见引理
\ref{spaces-lemma-composition-representable-transformations}。
最后，设 $U$ 是 $(\Sch/S)_{fppf}$ 的对象，且 $U \to F$ 满而
\'etale。由假设，$U \times_F G$ 由概形 $U'$ 表示。由引理
\ref{spaces-lemma-base-change-representable-transformations-property}，
态射 $U' \to G$ 满且 \'etale。这验证了定义
\ref{spaces-definition-algebraic-space} 的最后一个条件，证明完成。
\end{proof}

\begin{lemma}
\label{spaces-lemma-representable-morphisms-spaces-property}
设 $S$ 是包含于 $\Sch_{fppf}$ 的概形，$F$、$G$ 是 $S$ 上的代数空间，
而 $G \to F$ 是可表示态射。设 $U \in \Ob((\Sch/S)_{fppf})$，
$q : U \to F$ 满且 \'etale，并令 $V = G \times_F U$。
最后，设 $\mathcal{P}$ 是定义
\ref{spaces-definition-relative-representable-property} 所述的概形态射性质。
则 $G \to F$ 具有性质 $\mathcal{P}$，当且仅当
$V \to U$ 具有性质 $\mathcal{P}$。
\end{lemma}

\begin{proof}
（本引理可由引理
\ref{spaces-lemma-base-change-representable-transformations-property} 与
\ref{spaces-lemma-descent-representable-transformations-property} 推出，
但这里也给出直接证明。）由定义显然可见，若 $G \to F$ 具有性质
$\mathcal{P}$，则 $V \to U$ 具有性质 $\mathcal{P}$。
反之，假设 $V \to U$ 具有性质 $\mathcal{P}$。
设 $T \to F$ 是从概形到 $F$ 的态射。令
$T' = T \times_F G$；由于 $G \to F$ 可表示，该纤维积是概形。
我们必须证明 $T' \to T$ 具有性质 $\mathcal{P}$。考虑概形的交换图表
$$
\xymatrix{
V \ar[d] & T \times_F V \ar[d] \ar[l] \ar[r] &
T \times_F G \ar[d] \ar@{=}[r] & T' \\
U & T \times_F U \ar[l] \ar[r] & T
}
$$
其中两个方块均为纤维积方块。因此，中间的箭头作为 $V \to U$
的基变换，具有性质 $\mathcal{P}$。最后，
$\{T \times_F U \to T\}$ 是 fppf 覆盖，因为它满且 \'etale；
而该性质关于基上的 fppf 拓扑是局部的，故 $T' \to T$
具有性质 $\mathcal{P}$。
\end{proof}

\begin{lemma}
\label{spaces-lemma-morphism-sheaves-with-P-effective-descent-etale}
设 $S$ 是包含于 $\Sch_{fppf}$ 的概形，$G \to F$ 是
$(\Sch/S)_{fppf}$ 上预层的变换，$\mathcal{P}$ 是概形态射的性质。
假设
\begin{enumerate}
\item $\mathcal{P}$ 在任意基变换下保持，是基上的 fppf 局部性质，
并且 $\mathcal{P}$ 型态射关于 fppf 覆盖满足下降；见“下降”定义
\ref{descent-definition-descending-types-morphisms}；
\item $G$ 是层；
\item $F$ 是代数空间；
\item 存在 $U \in \Ob((\Sch/S)_{fppf})$ 以及满 \'etale 态射
$U \to F$，使 $V = G \times_F U$ 可表示；并且
\item $V \to U$ 具有 $\mathcal{P}$。
\end{enumerate}
则 $G$ 是代数空间，$G \to F$ 可表示并具有性质 $\mathcal{P}$。
\end{lemma}

\begin{proof}
设 $T$ 是概形，$T \to F$ 是态射。则 $U \times_F T \to T$
为 \'etale 满射，故 $\{U \times_F T \to T\}$ 是 \'etale 拓扑的覆盖。
考虑
$$
W = G \times_F (U \times_F T) = V \times_F T = V \times_U (U \times_F T).
$$
由于 $F$ 是代数空间，它是概形。态射 $W \to U \times_F T$
作为 $V \to U$ 的基变换，具有性质 $\mathcal{P}$。存在同构
\begin{align*}
W \times_T (U \times_F T) & =
(G \times_F (U \times_F T)) \times_T (U \times_F T) \\
& = (U \times_F T) \times_T (G \times_F (U \times_F T)) \\
& = (U \times_F T) \times_T W
\end{align*}
它位于 $(U \times_F T) \times_T (U \times_F T)$ 上。
中间的等式将 $((g, (u_1, t)), (u_2, t))$ 映到
$((u_1, t), (g, (u_2, t)))$。这定义了
$W/U \times_F T/T$ 的下降数据；见“下降”定义
\ref{descent-definition-descent-datum}。这也可由“下降”引理
\ref{descent-lemma-descent-data-sheaves} 看出。事实上，我们有层
$G \times_F T$，它到 $U \times_F T$ 的基变换由 $W$ 表示，
而上面的同构正是“下降”引理
\ref{descent-lemma-descent-data-sheaves} 的证明中出现的同构。
由关于 $\mathcal{P}$ 的假设，上述下降数据可表示。
因此，由“下降”引理 \ref{descent-lemma-descent-data-sheaves}
的最后一个陈述，$G \times_F T$ 可表示。这证明
$G \to F$ 是可表示的函子变换。

\medskip\noindent
由于 $G \to F$ 可表示，由引理 \ref{spaces-lemma-representable-over-space}，
$G$ 是代数空间。再由引理
\ref{spaces-lemma-representable-morphisms-spaces-property}，
$G \to F$ 具有性质 $\mathcal{P}$。
\end{proof}

\begin{lemma}
\label{spaces-lemma-lift-morphism-presentations}
设 $S$ 是包含于 $\Sch_{fppf}$ 的概形，$F,G$ 是 $S$ 上的代数空间，
$a : F \to G$ 是态射。给定任意 $V \in \Ob((\Sch/S)_{fppf})$
以及满 \'etale 态射 $q : V \to G$，存在
$U \in \Ob((\Sch/S)_{fppf})$ 以及交换图表
$$
\xymatrix{
U \ar[d]_p \ar[r]_\alpha &
V \ar[d]^q \\
F \ar[r]^a & G
}
$$
其中 $p$ 满且 \'etale。
\end{lemma}

\begin{proof}
首先选取 $W \in \Ob((\Sch/S)_{fppf})$ 以及满 \'etale 态射
$W \to F$。然后令 $U = W \times_G V$。由于 $G$ 是代数空间，
$U$ 同构于 $(\Sch/S)_{fppf}$ 的一个对象。由于 $q$ 为 \'etale 满射，
$U \to W$ 也为 \'etale 满射（见引理
\ref{spaces-lemma-base-change-representable-transformations-property}）。
因此，作为 \'etale 满射的复合，$U \to F$ 为 \'etale 满射
（见引理 \ref{spaces-lemma-composition-representable-transformations-property}）。
\end{proof}



























\section{代数空间的浸入与 Zariski 覆盖}
\label{spaces-section-Zariski}

\noindent
此处出现一个有趣现象。我们已经定义了代数空间开浸入的概念
（通过定义 \ref{spaces-definition-relative-representable-property}），
却尚未定义{\it 点}的概念\footnote{我们将在“空间的性质”第
\ref{spaces-properties-section-points} 节为代数空间关联一个拓扑空间，
其开集将恰好对应于下面定义的开子空间。}。
因此，代数空间的{\it Zariski 拓扑}已经定义，但空间本身尚未出现！

\medskip\noindent
或许有些多余，我们仍正式引入如下浸入概念。

\begin{definition}
\label{spaces-definition-immersion}
设 $S \in \Ob(\Sch_{fppf})$ 是概形，$F$ 是 $S$ 上的代数空间。
\begin{enumerate}
\item 若 $S$ 上代数空间的态射可表示，并且在定义
\ref{spaces-definition-relative-representable-property} 的意义下是开浸入，
则称其为{\it 开浸入}。
\item $F$ 的一个{\it 开子空间}是子函子 $F' \subset F$，
其中 $F'$ 是代数空间且 $F' \to F$ 是开浸入。
\item 若 $S$ 上代数空间的态射可表示，并且在定义
\ref{spaces-definition-relative-representable-property} 的意义下是闭浸入，
则称其为{\it 闭浸入}。
\item $F$ 的一个{\it 闭子空间}是子函子 $F' \subset F$，
其中 $F'$ 是代数空间且 $F' \to F$ 是闭浸入。
\item 若 $S$ 上代数空间的态射可表示，并且在定义
\ref{spaces-definition-relative-representable-property} 的意义下是浸入，
则称其为{\it 浸入}。
\item $F$ 的一个{\it 局部闭子空间}是子函子 $F' \subset F$，
其中 $F'$ 是代数空间且 $F' \to F$ 是浸入。
\end{enumerate}
\end{definition}

\noindent
这些定义是有意义的，因为浸入特别是单态射（见“概形”引理
\ref{schemes-lemma-immersions-monomorphisms} 以及引理
\ref{spaces-lemma-representable-transformations-property-implication}），
所以代数空间浸入 $G \to F$ 的像是子函子 $F' \subset F$，
且与 $G$（典范）同构。因此，“概形”第
\ref{schemes-section-immersions} 节的部分讨论可移用于代数空间。

\begin{lemma}
\label{spaces-lemma-composition-immersions}
设 $S \in \Ob(\Sch_{fppf})$ 是概形。$S$ 上代数空间的
（闭，相应地开）浸入之复合，仍是 $S$ 上代数空间的
（闭，相应地开）浸入。
\end{lemma}

\begin{proof}
见引理 \ref{spaces-lemma-composition-representable-transformations-property}，
以及评注 \ref{spaces-remark-list-properties-fpqc-local-base}
（见该评注最后一行）和 \ref{spaces-remark-list-properties-stable-composition}。
\end{proof}

\begin{lemma}
\label{spaces-lemma-base-change-immersions}
设 $S \in \Ob(\Sch_{fppf})$ 是概形。$S$ 上代数空间的
（闭，相应地开）浸入的基变换，仍是 $S$ 上代数空间的
（闭，相应地开）浸入。
\end{lemma}

\begin{proof}
见引理 \ref{spaces-lemma-base-change-representable-transformations-property}
以及评注 \ref{spaces-remark-list-properties-fpqc-local-base}
（见该评注最后一行）。
\end{proof}

\begin{lemma}
\label{spaces-lemma-sub-subspaces}
设 $S \in \Ob(\Sch_{fppf})$ 是概形，$F$ 是 $S$ 上的代数空间。
设 $F_1$、$F_2$ 是 $F$ 的局部闭子空间。若作为 $F$ 的子函子有
$F_1 \subset F_2$，则 $F_1$ 是 $F_2$ 的局部闭子空间。
闭子空间和开子空间的情形同理。
\end{lemma}

\begin{proof}
设 $T \to F_2$ 是态射，其中 $T$ 是概形。由于 $F_2 \to F$
是单态射，有 $T \times_{F_2} F_1 = T \times_F F_1$。
本引理由此形式地推出。
\end{proof}

\noindent
下面正式定义代数空间的 Zariski 开覆盖。注意，在引理
\ref{spaces-lemma-glueing-algebraic-spaces} 中，我们已经将这种开覆盖
作为构造代数空间的方法。

\begin{definition}
\label{spaces-definition-Zariski-open-covering}
设 $S \in \Ob(\Sch_{fppf})$ 是概形，$F$ 是 $S$ 上的代数空间。
$F$ 的一个{\it Zariski 覆盖} $\{F_i \subset F\}_{i \in I}$，
由集合 $I$ 以及一族开子空间 $F_i \subset F$ 给出，并要求
$\coprod F_i \to F$ 是层的满映射。
\end{definition}

\noindent
注意，若 $T$ 是概形且 $a : T \to F$ 是态射，则每个纤维积
$T \times_F F_i$ 都等同于某个开子概形 $T_i \subset T$。
定义的最后一个条件恰好表示 $T = \bigcup_{i \in I} T_i$。

\medskip\noindent
显然，$F$ 的开子空间之集 $F_{Zar}$ 是集合
（因为 $(\Sch/S)_{fppf}$ 是位点，因而是集合）。
此外，以子函子的包含为态射，可将 $F_{Zar}$ 组成范畴；
由引理 \ref{spaces-lemma-sub-subspaces}，这些包含自动为开浸入。
最后，定义 \ref{spaces-definition-Zariski-open-covering} 给出范畴
$F_{Zar}$ 中 Zariski 覆盖 $\{F_i \to F'\}_{i \in I}$ 的概念。
因此，正如拓扑空间的情形（见“位点”例
\ref{sites-example-site-topological}），通过适当选取一组覆盖，
可得到代数空间 $F$ 的一个 Zariski 位点。

\begin{definition}
\label{spaces-definition-small-Zariski-site}
设 $S \in \Ob(\Sch_{fppf})$ 是概形，$F$ 是 $S$ 上的代数空间。
代数空间 $F$ 的一个{\it 小 Zariski 位点 $F_{Zar}$}，
是上述位点之一。
\end{definition}

\noindent
这给出了某个陈述在代数空间上 Zariski 局部成立的含义，
我们将如此使用该概念。一般而言，Zariski 拓扑对我们的目的不够细。
例如，可以考虑代数空间上的 Zariski 层范畴；事实将表明，
即使对拟凝聚层，这也不是正确的对象。只有使用 $F$ 的 \'etale
或 fppf 位点定义拟凝聚层，才会得到所需结果。











\section{代数空间的分离条件}
\label{spaces-section-separation}

\noindent
代数空间 $F$ 的分离条件，是对角态射 $F \to F \times F$ 的条件。
先列出对角态射自动具有的性质。由于按定义对角态射可表示，
下面的引理（通过定义 \ref{spaces-definition-relative-representable-property}）
是有意义的。

\begin{lemma}
\label{spaces-lemma-properties-diagonal}
设 $S$ 是包含于 $\Sch_{fppf}$ 的概形，$F$ 是 $S$ 上的代数空间，
$\Delta : F \to F \times F$ 是对角态射。则
\begin{enumerate}
\item $\Delta$ 局部有限型；
\item $\Delta$ 是单态射；
\item $\Delta$ 分离；并且
\item $\Delta$ 局部拟有限。
\end{enumerate}
\end{lemma}

\begin{proof}
设 $F = U/R$ 是 $F$ 的一个表示。与引理
\ref{spaces-lemma-presentation-quasi-compact} 的证明相同，图表
$$
\xymatrix{
R \ar[r] \ar[d]_j & F \ar[d]^\Delta \\
U \times_S U \ar[r] & F \times F
}
$$
是 Cartesian 方块。因此，由引理
\ref{spaces-lemma-representable-morphisms-spaces-property}，只需证明 $j$
具有本引理所列性质。（注意，性质 (1)--(4) 都出现在评注
\ref{spaces-remark-list-properties-stable-base-change} 与
\ref{spaces-remark-list-properties-fpqc-local-base} 的列表中。）
由于 $j$ 是等价关系，它是单态射，故由“概形”引理
\ref{schemes-lemma-monomorphism-separated} 分离。
由于 $R$ 是 \'etale 等价关系，$s, t : R \to U$ 为 \'etale，
因而 $s, t$ 局部有限型。于是由“态射”引理
\ref{morphisms-lemma-permanence-finite-type}，$j$ 局部有限型。
最后，因其为单态射，纤维有限；由“态射”引理
\ref{morphisms-lemma-finite-fibre}，它局部拟有限。
\end{proof}

\noindent
下面给出相对于基概形 $S$ 的若干常见分离条件。
这些条件也有绝对版本，我们将在“空间的性质”第
\ref{spaces-properties-section-separation} 节讨论。此外，
“空间的态射”第 \ref{spaces-morphisms-section-separation-axioms} 节
将讨论代数空间态射的分离条件。

\begin{definition}
\label{spaces-definition-separated}
设 $S$ 是包含于 $\Sch_{fppf}$ 的概形，$F$ 是 $S$ 上的代数空间，
$\Delta : F \to F \times F$ 是对角态射。
\begin{enumerate}
\item 若 $\Delta$ 是闭浸入，则称 $F$ {\it 在 $S$ 上分离}。
\item 若 $\Delta$ 是浸入，则称 $F$
{\it 在 $S$ 上局部分离}\footnote{文献中，这通常指拟分离且局部分离的代数空间。}。
\item 若 $\Delta$ 拟紧，则称 $F$ {\it 在 $S$ 上拟分离}。
\item 若存在 Zariski 覆盖 $F = \bigcup_{i \in I} F_i$，使每个
$F_i$ 都拟分离，则称 $F$ {\it 在 $S$ 上 Zariski 局部拟分离}
\footnote{该定义由 B.\ Conrad 建议。}。
\end{enumerate}
\end{definition}

\noindent
注意，若对角态射拟紧（即 $F$ 分离或拟分离时），则对角态射实际上
拟有限且分离，因而拟仿射（由“态射续论”引理
\ref{more-morphisms-lemma-quasi-finite-separated-quasi-affine}）。








\section{代数空间的例子}
\label{spaces-section-examples}

\noindent
本节构造若干代数空间的例子，其中一些由 B.\ Conrad 建议。
由于我们尚无太多可用理论，部分讨论略显笨拙。

\begin{example}
\label{spaces-example-affine-line-involution}
设 $k$ 是特征 $\not = 2$ 的域，$U = \mathbf{A}^1_k$。令
$$
j : R = \Delta \amalg \Gamma \longrightarrow U \times_k U
$$
其中 $\Delta = \{(x, x) \mid x \in \mathbf{A}^1_k\}$，而
$\Gamma = \{(x, -x) \mid x \in \mathbf{A}^1_k, x \not = 0\}$。
显然 $s, t : R \to U$ 为 \'etale，故 $j$ 是 \'etale 等价关系。
由定理 \ref{spaces-theorem-presentation}，商 $X = U/R$ 是代数空间。
由于 $R$ 拟紧，$X$ 拟分离。另一方面，由于态射 $j$ 不是浸入，
$X$ 不局部分离。
\end{example}

\begin{example}
\label{spaces-example-non-representable-descent}
设 $k$ 是域，$k'/k$ 是 $2$ 次 Galois 扩张，且
$\text{Gal}(k'/k) = \{1, \sigma\}$。令 $S = \Spec(k[x])$
且 $U = \Spec(k'[x])$。注意
$$
U \times_S U =
\Spec((k' \otimes_k k')[x]) =
\Delta(U) \amalg \Delta'(U)
$$
其中 $\Delta' = (1, \sigma) : U \to U \times_S U$。取
$$
R = \Delta(U) \amalg \Delta'(U \setminus \{0_U\})
$$
其中 $0_U \in U$ 表示 $x$-坐标为零的 $k'$-有理点。
容易看出，$R$ 是 $U$ 上相对于 $S$ 的 \'etale 等价关系，
故由定理 \ref{spaces-theorem-presentation}，$X = U/R$ 是代数空间。
$X$ 具有下列性质（其中一些要到后文才有意义）：
\begin{enumerate}
\item $X \to S$ 在 $S \setminus \{0_S\}$ 上是同构；
\item 态射 $X \to S$ 为 \'etale（见“空间的性质”定义
\ref{spaces-properties-definition-etale}）；
\item $X \to S$ 在 $0_S$ 上的纤维 $0_X$ 同构于
$\Spec(k') = 0_U$；
\item $X$ 不是概形，因为若它是概形，则 $\mathcal{O}_{X, 0_X}$
将是一个局部整环 $(\mathcal{O}, \mathfrak m, \kappa)$，其分式域为
$k(x)$，满足 $x \in \mathfrak m$ 且剩余域 $\kappa = k'$，这是不可能的；
\item $X$ 不分离，但局部分离且拟分离；
\item 存在满、有限、\'etale 的态射 $S' \to S$，使基变换
$X' = S' \times_S X$ 为概形（具体而言，基变换到
$S' = \Spec(k'[x])$ 后，$U$ 分裂成 $S'$ 的两个副本，而 $X'$
同构于 $0$ 加倍的仿射直线；见“概形”例
\ref{schemes-example-affine-space-zero-doubled}）；并且
\item 若把 $X$ 看作 $\Spec(k)$ 上的有限型代数空间，则类似地，
基变换 $X_{k'}$ 是概形，但 $X$ 不是概形。
\end{enumerate}
特别地，这给出了一个相对于覆盖 $\{\Spec(k') \to \Spec(k)\}$
而无效的概形下降数据的例子。
\end{example}

\noindent
另见“例子”引理 \ref{examples-lemma-non-effective-descent-projective}；
它说明，即使对概形的射影态射，下降数据也未必有效。
该例给出 ${\mathbf C}$ 上一个不是概形的三维平滑分离代数空间。

\medskip\noindent
下面的引理提供一种方便的方法：将代数空间构造成概形对自由群作用的商。

\begin{lemma}
\label{spaces-lemma-quotient}
设 $U \to S$ 是 $\Sch_{fppf}$ 中的态射，$G$ 是抽象群，
$G \to \text{Aut}_S(U)$ 是群同态。假设
\begin{itemize}
\item[(*)] 若 $u \in U$ 是点，且对某个非恒等元素 $g \in G$
有 $g(u) = u$，则 $g$ 诱导 $\kappa(u)$ 的非平凡自同构。
\end{itemize}
则
$$
j :
R = \coprod\nolimits_{g \in G} U
\longrightarrow
U \times_S U,
\quad
(g, x) \longmapsto (g(x), x)
$$
是 \'etale 等价关系，因而
$$
F = U/R
$$
由定理 \ref{spaces-theorem-presentation} 是代数空间。
\end{lemma}

\begin{proof}
在引理陈述中，$\text{Aut}_S(U)$ 表示 $U$ 在 $S$ 上的自同构群。
假设 $(*)$ 成立。下面证明
$$
j :
R = \coprod\nolimits_{g \in G} U
\longrightarrow
U \times_S U,
\quad
(g, x) \longmapsto (g(x), x)
$$
是单态射。这意味着：若 $T$ 是非空概形，$h : T \to U$
是满足 $g \circ h = g' \circ h$ 的 $T$-值点，则 $g = g'$。
设 $T \not = \emptyset$，$h : T \to U$ 且
$g \circ h = g' \circ h$。取 $t \in T$，考虑复合
$\Spec(\kappa(t)) \to \Spec(\kappa(h(t))) \to U$。
于是 $g^{-1} \circ g'$ 固定 $u = h(t)$，并在其剩余域上恒等。
由 $(*)$，$g = g'$。

\medskip\noindent
因此，若 $(*)$ 成立，则 $j$ 是关系（见“群胚”定义
\ref{groupoids-definition-equivalence-relation}）。它还是等价关系，
因为对连通概形 $T$ 的 $T$-值点，有
$R(T) = G \times U(T) \to U(T) \times U(T)$
（回忆我们始终在 $S$ 上工作）。此外，由于 $R$ 是 $U$ 的若干副本的
不交积，态射 $s, t : R \to U$ 为 \'etale。这证明
$j : R \to U \times_S U$ 是 \'etale 等价关系。
\end{proof}

\noindent
给定概形 $U$ 以及群 $G$ 在 $U$ 上的作用，若引理
\ref{spaces-lemma-quotient} 的条件 $(*)$ 成立，则称 $G$ 在 $U$ 上的作用
为{\it 自由的}。这等价于“群胚”定义
\ref{groupoids-definition-free-action} 中常群概形 $G_S$ 在 $U$ 上的
自由作用概念。该引理可解释为：概形对群的自由作用之商在代数空间范畴中存在。

\begin{definition}
\label{spaces-definition-quotient}
记号 $U \to S$、$G$、$R$ 如引理 \ref{spaces-lemma-quotient}。
若 $G$ 在 $U$ 上的作用满足 $(*)$，则称 $G$ {\it 自由作用于}
概形 $U$。此时，代数空间 $U/R$ 记作 $U/G$，称为
{\it $U$ 对 $G$ 的商}。
\end{definition}

\noindent
该记号与“群胚”定义 \ref{groupoids-definition-quotient-sheaf}
中引入的 $U/G$ 一致。后文我们将在不对作用作任何假设的情况下，
把该商理解为代数栈；届时使用记号 $[U/G]$。
在讨论例子之前，再证明几个便于讨论的引理。下面的引理讨论 $G$
有限时该商的各种分离条件。

\begin{lemma}
\label{spaces-lemma-quotient-finite-separated}
记号与假设如引理 \ref{spaces-lemma-quotient}。假设 $G$ 有限。则
\begin{enumerate}
\item 若 $U \to S$ 拟分离，则 $U/G$ 在 $S$ 上拟分离；并且
\item 若 $U \to S$ 分离，则 $U/G$ 在 $S$ 上分离。
\end{enumerate}
\end{lemma}

\begin{proof}
在引理 \ref{spaces-lemma-properties-diagonal} 的证明中已见，只需对态射
$j : R \to U \times_S U$ 证明相应性质。若 $U \to S$ 拟分离，
则对每个映入 $S$ 中某个仿射开集的仿射开集 $V \subset U$，
开集 $g(V) \cap V$ 拟紧。因此 $j$ 拟紧。
若 $U \to S$ 分离，则对角态射 $\Delta_{U/S}$ 是闭浸入。
故 $j : R \to U \times_S U$ 是像彼此不交的有限个闭浸入的余积，
从而 $j$ 是闭浸入。
\end{proof}

\begin{lemma}
\label{spaces-lemma-quotient-field-map}
记号与假设如引理 \ref{spaces-lemma-quotient}。若
$\Spec(k) \to U/G$ 是态射，则存在
\begin{enumerate}
\item 有限 Galois 扩张 $k'/k$；
\item 有限子群 $H \subset G$；
\item 同构 $H \to \text{Gal}(k'/k)$；并且
\item $H$-等变态射 $\Spec(k') \to U$。
\end{enumerate}
反之，这样的数据确定一个态射 $\Spec(k) \to U/G$。
\end{lemma}

\begin{proof}
考虑纤维积 $V = \Spec(k) \times_{U/G} U$。有图表
$$
\xymatrix{
V \ar[r] \ar[d] & U \ar[d] \\
\Spec(k) \ar[r] & U/G
}
$$
于是 $V$ 是 $\Spec(k)$ 上的非空 \'etale 概形，因而是不交并
$V = \coprod_{i \in I} \Spec(k_i)$，其中各 $k_i$ 是 $k$ 的有限可分扩张
（“态射”引理 \ref{morphisms-lemma-etale-over-field}）。有
\begin{align*}
V \times_{\Spec(k)} V
& =
(\Spec(k) \times_{U/G} U) \times_{\Spec(k)}(\Spec(k) \times_{U/G} U) \\
& = 
\Spec(k) \times_{U/G} U \times_{U/G} U \\
& =
\Spec(k) \times_{U/G} U \times G \\
& =
V \times G
\end{align*}
$G$ 在 $U$ 上的作用诱导作用 $a : G \times V \to V$。
上述等式意味着映射
$G \times V \to V \times_{\Spec(k)} V$，
$(g, v) \mapsto (a(g, v), v)$ 是同构。特别地，对每个 $i$ 都有同构
$H_i \times \Spec(k_i) \to \Spec(k_i \otimes_k k_i)$，
其中 $H_i \subset G$ 是固定 $i \in I$ 的元素组成的子群。
因此 $H_i$ 有限，并且是 $k_i/k$ 的 Galois 群。
逆向构造从略。
\end{proof}

\noindent
例如，由该引理可知，若 $k'/k$ 是有限 Galois 扩张，则
$\Spec(k')/\text{Gal}(k'/k) \cong \Spec(k)$。
若扩张无限会怎样？下面给出一个例子。

\begin{example}
\label{spaces-example-Qbar}
令 $S = \Spec(\mathbf{Q})$、$U = \Spec(\overline{\mathbf{Q}})$，
并令 $G = \text{Gal}(\overline{\mathbf{Q}}/\mathbf{Q})$ 以显然方式作用于
$U$。由构造，引理 \ref{spaces-lemma-quotient} 的性质 $(*)$ 成立，
从而得到代数空间
$$
X = \Spec(\overline{\mathbf{Q}})/G
\longrightarrow
S = \Spec(\mathbf{Q}).
$$
当然，作为对 $S$ 的近似，这完全荒谬！事实上，由 Artin--Schreier 定理
（见 \cite[Theorem 17, page 316]{JacobsonIII}），
$\text{Gal}(\overline{\mathbf{Q}}/\mathbf{Q})$ 的有限子群只有 $\{1\}$，
以及二阶群
$\text{Gal}(\overline{\mathbf{Q}}/\overline{\mathbf{Q}} \cap \mathbf{R})$
的共轭。因此，若 $\Spec(k) \to X$ 是态射，且 $k$ 在
$\mathbf{Q}$ 上代数，则由引理 \ref{spaces-lemma-quotient-field-map}
与上述定理，$k$ 或为 $\overline{\mathbf{Q}}$，或同构于
$\overline{\mathbf{Q}} \cap \mathbf{R}$。
\end{example}

\noindent
上述例子的问题在于，Galois 群自带一个拓扑；构造
$\Spec(\overline{\mathbf{Q}})$ 的任何商时，都应以某种方式纳入该拓扑。
在我看来，下面的例子合理得多，而且可能确实会在“自然界”中出现。

\begin{example}
\label{spaces-example-affine-line-translation}
设 $k$ 是特征零的域，$U = \mathbf{A}^1_k$ 且 $G = \mathbf{Z}$。
取作用 $n(x) = x + n$，即 $\mathbf{Z}$ 通过平移作用于仿射直线。
唯一的不动点是泛点，并且显然 $\mathbf{Z}$ 嵌入域 $k(x)$ 的自同构群。
（这里用到了特征零假设。）考虑仿射直线的泛点到商的态射
$$
\gamma : \Spec(k(x)) \longrightarrow X = \mathbf{A}^1_k/\mathbf{Z}
$$
我们断言，该态射不经过任何从域的谱到 $X$ 的单态射
$\Spec(L) \to X$ 分解。（这与概形的情形相反；见“概形”第
\ref{schemes-section-points} 节。）事实上，由于 $\mathbf{Z}$
没有非平凡有限子群，由引理 \ref{spaces-lemma-quotient-field-map} 可知，
任何这样的分解都满足 $k(x) = L$。最后，$\gamma$ 不是单态射，因为
$$
\Spec(k(x)) \times_{\gamma, X, \gamma} \Spec(k(x))
\cong
\Spec(k(x)) \times \mathbf{Z}.
$$
\end{example}

\noindent
该例表明，为定义代数空间 $X$ 的点，我们应考虑从域的谱到 $X$
的态射之等价类，而不是从域的谱出发的单态射之集。

\medskip\noindent
最后给出一个着实糟糕的例子。

\begin{example}
\label{spaces-example-infinite-product}
设 $k$ 是域，令 $A = \prod_{n \in \mathbf{N}} k$ 为无限乘积。
令 $U = \Spec(A)$，视为 $S = \Spec(k)$ 上的概形。
注意，投影映射 $\text{pr}_n : A \to k$ 定义开浸入兼闭浸入
$f_n : S \to U$。令
$$
R = U \amalg \coprod\nolimits_{(n, m) \in \mathbf{N}^2, \ n \not = m} S
$$
其中态射 $j$ 在分支 $U$ 上等于 $\Delta_{U/S}$，
在与 $(n,m)$ 对应的分支 $S$ 上等于 $j = (f_n, f_m)$。
由上面的评注，显然 $s,t$ 为 \'etale；也显然 $j$ 是等价关系。
因此得到代数空间
$$
X = U/R.
$$
为理解其含义，专门考虑域 $k$ 为含 $q$ 个元素的有限域的情形。
先简要讨论与概形 $U$ 关联的拓扑空间 $|U|$。
$A$ 的所有元素都满足 $x^q = x$。因此 $A$ 的每个剩余域都同构于
$k$，且 $U$ 的所有点都是闭点。但 $U$ 上的拓扑不是离散拓扑。
设 $u_n \in |U|$ 是与 $f_n$ 对应的点。如上所述，点 $u_n$
都是开点（因而是孤立点）。但必定还存在其他点，因为 $U$ 拟紧，
见“代数”引理 \ref{algebra-lemma-quasi-compact}，
所以它不可能等于一个无限离散集。另一种看法是，真理想
$$
I =
\{x = (x_n) \in A \mid \text{除有限多个外的 }x_n\text{ 均为零}\}
$$
包含于某个极大理想。还要注意，$A$ 的每个元素 $x$ 都形如
$x = ue$，其中 $u$ 是单位，$e$ 是幂等元。因此 $A$ 的拓扑有一个
由开闭子集组成的基（见“代数”引理
\ref{algebra-lemma-idempotent-spec}）。所以 $|U|$ 上的拓扑完全不连通，
但并非平凡。最后，注意 $\{u_n\}$ 在 $|U|$ 中稠密。

\medskip\noindent
后文将定义与 $X$ 关联的拓扑空间 $|X|$；见“空间的性质”第
\ref{spaces-properties-section-points} 节。关于 $|X|$ 能说什么？
事实上，映射 $|U| \to |X|$ 满且连续。所有点 $u_n$ 都映到 $|X|$
中的同一点 $x_0$，而其他点都不被等同。由于 $\{u_n\}$ 在 $|U|$
中稠密，$x_0$ 在 $|X|$ 中的闭包就是 $|X|$。换言之，$|X|$ 不可约，
且 $x_0$ 是 $|X|$ 的泛点。这似乎颇为怪异，因为 $x_0$ 还是结构态射
$X \to S$ 的一个截面 $S \to X$ 的像（而在概形情形下，这将蕴含它是闭点；
见“态射”引理
\ref{morphisms-lemma-algebraic-residue-field-extension-closed-point-fibre}）。
\end{example}

\noindent
无论你认为该例中实际发生了什么，它都清楚表明：
定义代数空间的不可约分支、连通性等概念时必须谨慎。






\section{大站点的变更}
\label{spaces-section-change-big-site}

\noindent
本节简要讨论变更大站点时会发生什么。结论是，我们总可以任意扩大
大站点；因此，可以假定我们所要考虑的任意一组概形，都包含在我们
考察代数空间时所用的大 fppf 站点中。下面给出这一结论的精确表述。

\begin{lemma}
\label{spaces-lemma-change-big-site}
设给定大站点 $\Sch_{fppf}$ 和 $\Sch'_{fppf}$。
假定 $\Sch_{fppf}$ 包含于 $\Sch'_{fppf}$，参见
拓扑，章节 \ref{topologies-section-change-alpha}。
设 $S$ 是 $\Sch_{fppf}$ 的一个对象，并令
\begin{align*}
g : \Sh((\Sch/S)_{fppf})
\longrightarrow
\Sh((\Sch'/S)_{fppf}), \\
f : \Sh((\Sch'/S)_{fppf})
\longrightarrow
\Sh((\Sch/S)_{fppf})
\end{align*}
为拓扑，引理 \ref{topologies-lemma-change-alpha}
中的拓扑斯态射。设 $F$ 是 $(\Sch/S)_{fppf}$ 上的集合层。则
\begin{enumerate}
\item 若 $F$ 可由 $S$ 上的概形
$X \in \Ob((\Sch/S)_{fppf})$ 表示，则 $f^{-1}F$ 也可表示；
事实上，它由同一个概形 $X$ 表示，只是现在将该概形视为
$(\Sch'/S)_{fppf}$ 的对象；并且
\item 若 $F$ 是 $S$ 上的代数空间，则 $f^{-1}F$ 也是
$S$ 上的代数空间。
\end{enumerate}
\end{lemma}

\begin{proof}
设 $X \in \Ob((\Sch/S)_{fppf})$。以 $h_X$ 表示
$(\Sch/S)_{fppf}$ 上与 $X$ 相伴的可表示层，以 $h'_X$
表示 $(\Sch'/S)_{fppf}$ 上与 $X$ 相伴的可表示层。根据
拓扑，章节 \ref{topologies-section-change-alpha}
中对 $f^{-1}$ 的描述可知 $f^{-1}h_X = h'_X$。这证明了 (1)。

\medskip\noindent
再假定 $F$ 是 $S$ 上的代数空间。由引理
\ref{spaces-lemma-space-presentation}，这意味着对
$(\Sch/S)_{fppf}$ 中的某个 \'etale 等价关系
$R \to U \times_S U$，有 $F = h_U/h_R$。由于 $f^{-1}$
是正合函子，遂有 $f^{-1}F = h'_U/h'_R$。因此，由定理
\ref{spaces-theorem-presentation}，$f^{-1}F$ 是 $S$ 上的代数空间。
\end{proof}

\noindent
注意，此引理纯属集合论，几乎没有实质内容。此外，一般而言，大站点上
代数空间的限制并不一定是小站点上的代数空间（原因仅在于基数）。
因此，这类简单引理只能用来扩大基范畴，而不能用来缩小它。

\begin{lemma}
\label{spaces-lemma-fully-faithful}
假定 $\Sch_{fppf}$ 包含于 $\Sch'_{fppf}$。设 $S$ 是
$\Sch_{fppf}$ 的一个对象。以 $\textit{Spaces}/S$ 表示使用
$\Sch_{fppf}$ 定义的 $S$ 上代数空间范畴；类似地，以
$\textit{Spaces}'/S$ 表示使用 $\Sch'_{fppf}$ 定义的
$S$ 上代数空间范畴。引理 \ref{spaces-lemma-change-big-site}
的构造定义了全忠实函子
$$
\textit{Spaces}/S \longrightarrow \textit{Spaces}'/S
$$
其本质像由如下的 $X' \in \Ob(\textit{Spaces}'/S)$ 组成：存在
$U, R \in \Ob((\Sch/S)_{fppf})$\footnote{必须要求 $R$ 存在，
这是因为我们在集合，方程 (\ref{sets-equation-bound}) 中选择了函数
$Bound$。就 $U$ 的大小而言，纤维积 $U \times_{X'} U$ 的大小可能
增长得比 $Bound$ 更快。令 $S = \Spec(A)$、
$U = \Spec(A[x_i, i \in I])$ 以及
$R = \coprod_{(\lambda_i) \in A^I} \Spec(A[x_i, y_i]/(x_i - \lambda_i y_i))$
即可说明这一点。在此情形中，$R$ 的大小按 $\kappa^\kappa$ 增长，
其中 $\kappa$ 是 $U$ 的大小。}，以及态射
$$
U \longrightarrow X'
\quad\text{且}\quad
R \longrightarrow U \times_{X'} U
$$
它们位于 $\Sh((\Sch'/S)_{fppf})$ 中，并且作为层的映射是满射
（例如，若所示态射为满 \'etale 态射）。
\end{lemma}

\begin{proof}
由站点，引理 \ref{sites-lemma-bigger-site}，函子
$f^{-1} : \Sh((\Sch/S)_{fppf}) \to \Sh((\Sch'/S)_{fppf})$
是全忠实的（参见拓扑，章节
\ref{topologies-section-change-alpha} 中的讨论）。
因此，本引理所示的函子也是全忠实的。

\medskip\noindent
设 $X' \in \Ob(\textit{Spaces}'/S)$，并且存在
$U \in \Ob((\Sch/S)_{fppf})$ 以及
$\Sh((\Sch'/S)_{fppf})$ 中的映射 $U \to X'$，后者作为层的映射
是满射。取满 \'etale 态射 $U' \to X'$，其中
$U' \in \Ob((\Sch'/S)_{fppf})$。令
$\kappa = \text{size}(U)$，参见集合，章节
\ref{sets-section-categories-schemes}。
则 $U$ 有仿射开覆盖 $U = \bigcup_{i \in I} U_i$，其中
$|I| \leq \kappa$。注意，$U' \times_{X'} U \to U$ 是满
\'etale 态射。对每个 $i$，可取拟紧开子集
$U'_i \subset U'$，使得 $U'_i \times_{X'} U_i \to U_i$
为满射（这是因为概形 $U' \times_{X'} U_i$ 是 Zariski 开集
$W \times_{X'} U_i$ 的并，其中 $W \subset U'$ 为仿射开集；
又因为 $U' \times_{X'} U_i \to U_i$ 是 \'etale 态射，因而是开映射）。
于是 $\coprod_{i \in I} U'_i \to X$ 是满 \'etale 态射，
因为我们假定 $U \to X$ 是层的满射，从而 $\coprod U_i \to X$
也是层的满射（细节略）。
由于 $U'_i \times_{X'} U \to U'_i$ 是层的满射，且
$U'_i$ 拟紧，可找到拟紧开子集
$W_i \subset U'_i \times_{X'} U$，使得 $W_i \to U'_i$
作为层的映射是满射（细节略）。于是 $W_i \to U$ 是 \'etale
态射，故由集合，引理 \ref{sets-lemma-bound-finite-type}，
$\text{size}(W_i) \leq \text{size}(U)$。再由集合，引理
\ref{sets-lemma-bound-by-covering}，可得
$\text{size}(U'_i) \leq \text{size}(U)$。因此，由集合，引理
\ref{sets-lemma-bound-size}，$\coprod_{i \in I} U'_i$
同构于 $(\Sch/S)_{fppf}$ 的一个对象。

\medskip\noindent
现在令 $X'$、$U \to X'$ 和 $R \to U \times_{X'} U$ 如引理
所述。上一段表明，可以找到 $U' \in \Ob((\Sch/S)_{fppf})$
以及 $\Sh((\Sch'/S)_{fppf})$ 中的满 \'etale 态射
$U' \to X'$。于是 $U' \times_{X'} U \to U'$ 是层的满射；
换言之，可以找到 fppf 覆盖 $\{U'_i \to U'\}$，使
$U'_i \to U'$ 经由 $U' \times_{X'} U \to U'$ 分解。
由集合，引理 \ref{sets-lemma-bound-fppf-covering}，可以找到满、
平坦且局部有限表示的态射 $\tilde U \to U'$，满足
$\text{size}(\tilde U) \leq \text{size}(U')$，并且
$\tilde U \to U'$ 经由 $U' \times_{X'} U \to U'$ 分解。
于是考察
$$
\xymatrix{
U' \times_{X'} U' \ar[d] &
\tilde U \times_{X'} \tilde U  \ar[l] \ar[d] \ar[r] &
U \times_{X'} U \ar[d] \\
U' \times_S U' & \tilde U \times_S \tilde U \ar[l] \ar[r] & U \times_S U
}
$$
其中两个方块均为笛卡尔方块。我们知道，底行的对象可由
$(\Sch/S)_{fppf}$ 的对象表示。由上一段的论证，
$U \times_{X'} U$ 也具有这一性质（因为按假定，我们有层的满射
$R \to U \times_{X'} U$）。由于 $(\Sch/S)_{fppf}$
按构造对纤维积封闭，可知 $\tilde U \times_{X'} \tilde U$
可由 $(\Sch/S)_{fppf}$ 的一个对象表示。最后，因为
$\tilde U \to U'$ 是 fppf 层的满射，所以
$\tilde U \times_{X'} \tilde U \to U' \times_{X'} U'$
也是 fppf 层的满射。因此再次应用上一段的结论，可得
$R' = U' \times_{X'} U'$ 可由 $(\Sch/S)_{fppf}$ 的一个对象表示。
此时，引理 \ref{spaces-lemma-space-presentation} 和定理
\ref{spaces-theorem-presentation} 表明，
$X = h_{U'}/h_{R'}$ 是 $\textit{Spaces}/S$ 的一个对象，并满足
$f^{-1}X \cong X'$，正如所需。
\end{proof}



\section{基概形的变更}
\label{spaces-section-change-base-scheme}

\noindent
本节简要讨论变更基概形时会发生什么。结论是，给定基概形的态射
$S \to S'$，任何 $S$ 上的代数空间都可视为 $S'$ 上的代数空间。
反之，给定 $S'$ 上的代数空间 $F'$，存在它的基变换 $F'_S$，
这是 $S$ 上的代数空间。
这里只说明 $S \to S'$ 是所考察的大 fppf 站点中的态射这一情形。
若只有 $S$ 或 $S'$ 包含在大站点中，则先按章节
\ref{spaces-section-change-big-site} 扩大该大站点。

\begin{lemma}
\label{spaces-lemma-change-base-scheme}
设给定大站点 $\Sch_{fppf}$。令 $g : S \to S'$ 为
$\Sch_{fppf}$ 的态射，并令
$j : (\Sch/S)_{fppf} \to (\Sch/S')_{fppf}$
为相应的局部化函子。设 $F$ 是 $(\Sch/S)_{fppf}$ 上的集合层。则
\begin{enumerate}
\item 对 $S'$ 上的概形 $T'$，有
$j_!F(T'/S') =
\coprod\nolimits_{\varphi : T' \to S} F(T' \xrightarrow{\varphi} S),$
\item 若 $F$ 可由概形 $X \in \Ob((\Sch/S)_{fppf})$ 表示，
则 $j_!F$ 可由 $j(X)$ 表示，而后者就是把 $X$
视作 $S'$ 上的概形；并且
\item 若 $F$ 是 $S$ 上的代数空间，则 $j_!F$ 是 $S'$
上的代数空间；若 $F = U/R$ 是一个表示，则
$j_!F = j(U)/j(R)$ 也是一个表示。
\end{enumerate}
设 $F'$ 是 $(\Sch/S')_{fppf}$ 上的集合层。则
\begin{enumerate}
\item[(4)] 对 $S$ 上的概形 $T$，有
$j^{-1}F'(T/S) = F'(T/S')$；
\item[(5)] 若 $F'$ 可由概形
$X' \in \Ob((\Sch/S')_{fppf})$ 表示，则 $j^{-1}F'$ 可表示，
具体由 $X'_S = S \times_{S'} X'$ 表示；并且
\item[(6)] 若 $F'$ 是代数空间，则 $j^{-1}F'$ 是代数空间；
若 $F' = U'/R'$ 是一个表示，则
$j^{-1}F' = U'_S/R'_S$ 是一个表示。
\end{enumerate}
\end{lemma}

\begin{proof}
函子 $j_!$、$j_*$ 和 $j^{-1}$ 定义于站点，引理
\ref{sites-lemma-relocalize}；那里还证明了
$j = j_{S/S'}$ 是 $(\Sch/S')_{fppf}$ 在对象 $S/S'$
处的局部化。因此，关于局部化函子的全部结果都适用于 $j$。
(1) 中的公式即站点，引理
\ref{sites-lemma-describe-j-shriek-good-site}。
按定义，$j_!$ 是限制函子 $j^{-1}$ 的左伴随，故 $j_!$ 右正合。
由站点，引理
\ref{sites-lemma-j-shriek-commutes-equalizers-fibre-products}，
它还与纤维积和等化子交换。由站点，引理
\ref{sites-lemma-describe-j-shriek-representable} 可知
$j_!h_X = h_{j(X)}$，故 (2) 成立。
若 $F$ 是 $S$ 上的代数空间，则可写成 $F = U/R$
（引理 \ref{spaces-lemma-space-presentation}），并且有
$$
j_!F = j(U)/j(R)
$$
这是因为 $j_!$ 右正合，因而与余等化子交换；此外，由于
$j_!$ 与纤维积交换，有
$j(R) = j(U) \times_{j_!F} j(U)$。态射
$j(s), j(t) : j(R) \to j(U)$ 就是态射 $s, t : R \to U$
（只是视作 $S'$ 上概形的态射），所以仍为 \'etale 态射。
因此 $(j(U), j(R), s, t)$ 是一个 \'etale 等价关系。
由定理 \ref{spaces-theorem-presentation} 可知 $j_!F$ 是代数空间。

\medskip\noindent
证明 (4)、(5) 和 (6)。站点，章节 \ref{sites-section-localize}
给出了 $j^{-1}$ 的描述。由站点，引理
\ref{sites-lemma-localize-given-products}，与 $X'/S'$ 相伴的
可表示层的限制，是与
$X'_S = S \times_{S'} Y'$ 相伴的可表示层。
限制函子 $j^{-1}$ 正合，故 $j^{-1}F' = U'_S/R'_S$。
再次由正合性，层 $R'_S$ 仍是 $U'_S$ 上的等价关系。
最后，两个映射 $R'_S \to U'_S$ 是 \'etale 态射，因为它们是
\'etale 态射 $R' \to U'$ 的基变换。因此，由定理
\ref{spaces-theorem-presentation}，$j^{-1}F' = U'_S/R'_S$
是代数空间，即得所证。
\end{proof}

\noindent
注意，表示 $j_!F = j(U)/j(R)$ 正是 $F$ 的表示，只不过将其
视作由 $S'$ 上概形给出的表示。因此，下面的定义是合理的。

\begin{definition}
\label{spaces-definition-base-change}
设 $\Sch_{fppf}$ 为大 fppf 站点，并设 $S \to S'$ 为该站点的态射。
\begin{enumerate}
\item 若 $F'$ 是 $S'$ 上的代数空间，则
{\it $F'$ 到 $S$ 的基变换}是引理
\ref{spaces-lemma-change-base-scheme} 所述的代数空间 $j^{-1}F'$。
我们将它记作 $F'_S$。
\item 若 $F$ 是 $S$ 上的代数空间，则将 $F$
{\it 视作 $S'$ 上的代数空间}，所得对象是引理
\ref{spaces-lemma-change-base-scheme} 所述的 $S'$ 上代数空间
$j_!F$。我们通常仍将它简记为 $F$；若不这样记，则写作 $j_!F$。
\end{enumerate}
\end{definition}

\noindent
代数空间 $j_!F$ 带有 $S'$ 上代数空间的典范态射
$j_!F \to S$。这是因为层 $j_!F$ 映到 $h_S$（例如参见引理
\ref{spaces-lemma-change-base-scheme} 中的显式描述）。
事实上，站点，引理 \ref{sites-lemma-essential-image-j-shriek}
已经表明，$(\Sch/S)_{fppf}$ 上的层范畴等价于如下二元组
$(\mathcal{F}', \mathcal{F}' \to h_S)$ 的范畴：其中的层是
$(\Sch/S')_{fppf}$ 上的层，而
$\mathcal{F}' \to h_S$ 是层的映射。此等价把层 $\mathcal{F}$
送到二元组 $(j_!\mathcal{F}, j_!\mathcal{F} \to h_S)$。
结合上面的结果，得到下述关于代数空间范畴的结论。

\begin{lemma}
\label{spaces-lemma-category-of-spaces-over-smaller-base-scheme}
设 $\Sch_{fppf}$ 为大 fppf 站点，并设 $S \to S'$ 为该站点的态射。
上述构造给出范畴的等价
$$
\left\{
\begin{matrix}
\text{代数空间范畴}\\
\text{基概形为 }S
\end{matrix}
\right\}
\leftrightarrow
\left\{
\begin{matrix}
\text{二元组范畴 }(F', F' \to S)\text{，其中}\\
\text{含有代数空间 }F'\text{，基概形为 }S'\text{，以及}\\
\text{代数空间态射 }F' \to S\text{，基概形为 }S'
\end{matrix}
\right\}
$$
\end{lemma}

\begin{proof}
设 $F$ 是 $S$ 上的代数空间。从左向右的函子把 $F$ 送到二元组
$(j_!F, j_!F \to S)$；由引理 \ref{spaces-lemma-change-base-scheme}，
这是右侧范畴的对象。根据站点，引理
\ref{sites-lemma-essential-image-j-shriek}，该构造给出层范畴的等价，
所以为完成证明，只需说明：若 $F$ 是层且 $j_!F$ 是代数空间，则
$F$ 是代数空间。为此，按引理 \ref{spaces-lemma-space-presentation} 写成
$j_!F = U'/R'$，其中
$U', R' \in \Ob((\Sch/S')_{fppf})$。
于是复合 $U' \to j_!F \to S$ 和 $R' \to j_!F \to S$
是 $S'$ 上概形的态射。以 $U, R$ 表示 $(\Sch/S)_{fppf}$
中的相应对象。两个态射 $R' \to U'$ 都是 $S$ 上的态射，因而对应于
态射 $R \to U$。由于这些只是相同的态射（但视作 $S$ 上的态射），
我们得到 $S$ 上的一个 \'etale 等价关系。由于 $j_!$ 给出层范畴的
等价（见上述参考），可知 $F = U/R$；再由定理
\ref{spaces-theorem-presentation}，可知 $F$ 是代数空间。
\end{proof}

\noindent
下面的引理是对上述结论的略微改述。

\begin{lemma}
\label{spaces-lemma-rephrase}
设 $\Sch_{fppf}$ 为大 fppf 站点，设 $S \to S'$ 为该站点的态射，
并设 $F'$ 是 $(\Sch/S')_{fppf}$ 上的层。下列条件等价：
\begin{enumerate}
\item 限制 $F'|_{(\Sch/S)_{fppf}}$ 是 $S$ 上的代数空间；并且
\item 层 $h_S \times F'$ 是 $S'$ 上的代数空间。
\end{enumerate}
\end{lemma}

\begin{proof}
在站点，引理 \ref{sites-lemma-essential-image-j-shriek}
的范畴等价之下，该限制与该乘积彼此对应，因此上述引理
\ref{spaces-lemma-category-of-spaces-over-smaller-base-scheme} 给出结论。
\end{proof}

\noindent
最后用一个关于相容性的引理结束本节。

\begin{lemma}
\label{spaces-lemma-viewed-as-properties}
设 $\Sch_{fppf}$ 为大 fppf 站点，设 $S \to S'$ 为该站点的态射，
并设 $F$ 是 $S$ 上的代数空间。设 $T$ 是 $S$ 上的概形，
且 $f : T \to F$ 是 $S$ 上的态射。
对 $f$ 应用引理
\ref{spaces-lemma-category-of-spaces-over-smaller-base-scheme}
所述的范畴等价，得到 $S'$ 上的态射 $f' : T' \to F'$。
对定义 \ref{spaces-definition-relative-representable-property}
中的任意性质 $\mathcal{P}$，都有
$\mathcal{P}(f') \Leftrightarrow \mathcal{P}(f)$。
\end{lemma}

\begin{proof}
设 $U$ 是 $S$ 上的概形，且 $U \to F$ 是满 \'etale 态射。
以 $U'$ 表示把概形 $U$ 视作 $S'$ 上概形所得的对象。引理
\ref{spaces-lemma-change-base-scheme} 已表明 $U' \to F'$ 是满 \'etale
态射。由于
$$
j(T \times_{f, F} U) = T' \times_{f', F'} U'
$$
概形态射 $T \times_{f, F} U \to U$ 与概形态射
$T' \times_{f', F'} U' \to U'$ 相识别。它们是同一个态射，
只是视作位于不同的基概形上。因此，本引理由引理
\ref{spaces-lemma-representable-morphisms-spaces-property} 得出。
\end{proof}
% Generated mechanically from the frozen spaces-properties.tex target.
% Do not edit this generated file; edit the bound locale source instead.

% OK, start here.
%

\setcounter{chapter}{65}
\chapter{代数空间的性质}



\phantomsection
\label{spaces-properties-section-phantom}


\section{引言}
\label{spaces-properties-section-introduction}

\noindent
关于代数空间的简要介绍，请参见
“空间”第 \ref{spaces-section-introduction} 节；关于代数空间的基本定义与约定，
也请阅读该章的有关部分。本章开始介绍代数空间的一些基本概念和性质。
对于拟分离代数空间的情形，基础参考文献是 \cite{Kn}。

\medskip\noindent
由于我们在编排上决定先讨论代数空间本身的性质，稍后再讨论代数空间态射的性质，
有些地方的论述会略显不便。代数空间的{\it \'etale 态射}是这条规则的例外，
我们将在第 \ref{spaces-properties-section-etale-morphisms} 节中引入它。
但在到达该节之前，每当我们说某个态射具有某种性质时，均自动假定该态射的源是概形
（或者该态射是可表示的）。

\medskip\noindent
本章的部分内容（尤其是关于点的内容）将在“良好代数空间”一章中得到加强。


\section{约定}
\label{spaces-properties-section-conventions}

\noindent
贯穿本章的约定是：所有概形都包含在一个大 fppf 站点
$\Sch_{fppf}$ 中；所考虑的每个环 $A$ 都满足
$\Spec(A)$（同构于）该大站点的一个对象。

\medskip\noindent
设 $S$ 为概形，$X$ 为 $S$ 上的代数空间。在本章及后续各章中，
我们把 $X$ 与自身的积（在 $S$ 上代数空间的范畴中）记作
$X \times_S X$，而不记作 $X \times X$。这样做是为了避免在变更基概形时
产生混淆；参见“空间”第 \ref{spaces-section-change-base-scheme} 节。


\section{分离公理}
\label{spaces-properties-section-separation}

\noindent
本节汇集代数空间的所有“绝对”分离条件。在我们的表述中，任一代数空间都是某个
确定基概形上的代数空间。因此，$S$ 上代数空间 $X$ 的任一绝对性质，都对应于把
$X$ 视作 $\Spec(\mathbf{Z})$ 上的代数空间时所满足的一个条件。精确定义如下。

\begin{definition}
\label{spaces-properties-definition-separated}
（比较“空间”定义 \ref{spaces-definition-separated}。）
考虑大 fppf 站点
$\Sch_{fppf} = (\Sch/\Spec(\mathbf{Z}))_{fppf}$.
设 $X$ 是 $\Spec(\mathbf{Z})$ 上的代数空间，并设
$\Delta : X \to X \times X$ 为对角态射。
\begin{enumerate}
\item 若 $\Delta$ 是闭浸入，则称 $X$ {\it 分离}。
\item 若 $\Delta$ 是浸入，则称 $X$
{\it 局部分离}\footnote{文献中，这通常指同时拟分离且局部分离的代数空间。}。
\item 若 $\Delta$ 拟紧，则称 $X$ {\it 拟分离}。
\item 若存在 Zariski 覆盖
$X = \bigcup_{i \in I} X_i$（见“空间”定义
\ref{spaces-definition-Zariski-open-covering}），使得每个 $X_i$ 都拟分离，
则称 $X$ {\it Zariski 局部拟分离}\footnote{这一概念由 B.\ Conrad 建议。}。
\end{enumerate}
设 $S$ 是包含在 $\Sch_{fppf}$ 中的概形，$X$ 是 $S$ 上的代数空间。
若把 $X$ 视作 $\Spec(\mathbf{Z})$ 上的代数空间（见“空间”定义
\ref{spaces-definition-base-change}）时具有相应性质，则称 $X$
{\it 分离}、{\it 局部分离}、{\it 拟分离}或
{\it Zariski 局部拟分离}。
\end{definition}

\noindent
若 $S$ 上的代数空间 $X$ 按上述绝对意义分离，则它在 $S$ 上分离；
对上述其他绝对分离性质亦同。我们将在“空间的态射”第
\ref{spaces-morphisms-section-separation-axioms} 节中详细讨论这一点。
引理 \ref{spaces-properties-lemma-quasi-separated-quasi-compact-pieces} 将表明，
Zariski 局部分离与基概形无关（因而等价于绝对概念）。

\begin{lemma}
\label{spaces-properties-lemma-trivial-implications}
设 $S$ 为概形，$X$ 为 $S$ 上的代数空间。定义
\ref{spaces-properties-definition-separated} 的各分离公理之间有下列蕴含关系：
\begin{enumerate}
\item 分离蕴含其余所有条件；
\item 拟分离蕴含 Zariski 局部拟分离。
\end{enumerate}
\end{lemma}

\begin{proof}
略。
\end{proof}

\begin{lemma}
\label{spaces-properties-lemma-characterize-quasi-separated}
设 $S$ 为概形，$X$ 为 $S$ 上的代数空间。下列条件等价：
\begin{enumerate}
\item $X$ 是拟分离代数空间；
\item 对任意 $U \to X$、$V \to X$，若 $U$、$V$ 是拟紧概形，
则纤维积 $U \times_X V$ 拟紧；
\item 对任意 $U \to X$、$V \to X$，若 $U$、$V$ 是仿射概形，
则纤维积 $U \times_X V$ 拟紧。
\end{enumerate}
\end{lemma}

\begin{proof}
由“空间”引理
\ref{spaces-lemma-category-of-spaces-over-smaller-base-scheme}
可设 $S = \Spec(\mathbf{Z})$。由于
$U \times_X V = X \times_{X \times X} (U \times V)$，
且当 $U$、$V$ 拟紧时 $U \times V$ 也拟紧，所以 (1) 蕴含 (2)。
显然 (2) 蕴含 (3)。现假设 (3)。选取概形 $W$ 以及满 \'etale 态射
$W \to X$。于是 $W \times W \to X \times X$ 是满 \'etale 态射。
因此，只需证明
$$
j : W \times_X W = X \times_{(X \times X)} (W \times W) \to W \times W
$$
拟紧；见“空间”引理
\ref{spaces-lemma-descent-representable-transformations-property}。
若 $U \subset W$ 和 $V \subset W$ 是仿射开集，则由假设，
$j^{-1}(U \times V) = U \times_X V$ 拟紧。仿射开集 $U \times V$
构成 $W \times W$ 的仿射开覆盖（“概形”引理
\ref{schemes-lemma-affine-covering-fibre-product}），故结论由“概形”引理
\ref{schemes-lemma-quasi-compact-affine} 得出。
\end{proof}

\begin{lemma}
\label{spaces-properties-lemma-characterize-separated}
设 $S$ 为概形，$X$ 为 $S$ 上的代数空间。下列条件等价：
\begin{enumerate}
\item $X$ 是分离代数空间；
\item 对任意 $U \to X$、$V \to X$，若 $U$、$V$ 是仿射概形，
则纤维积 $U \times_X V$ 仿射，并且
$$
\mathcal{O}(U) \otimes_\mathbf{Z} \mathcal{O}(V)
\longrightarrow
\mathcal{O}(U \times_X V)
$$
是满射。
\end{enumerate}
\end{lemma}

\begin{proof}
由“空间”引理
\ref{spaces-lemma-category-of-spaces-over-smaller-base-scheme}
可设 $S = \Spec(\mathbf{Z})$。由于
$U \times_X V = X \times_{X \times X} (U \times V)$，
且当 $U$、$V$ 仿射时 $U \times V$ 也仿射，所以 (1) 蕴含 (2)。
现假设 (2)。选取概形 $W$ 以及满 \'etale 态射 $W \to X$。
于是 $W \times W \to X \times X$ 是满 \'etale 态射。因此，只需证明
$$
j : W \times_X W = X \times_{(X \times X)} (W \times W) \to W \times W
$$
是闭浸入；见“空间”引理
\ref{spaces-lemma-descent-representable-transformations-property}。
若 $U \subset W$ 和 $V \subset W$ 是仿射开集，则由假设，
$j^{-1}(U \times V) = U \times_X V$ 仿射；又因相应的环同态为满射，
映射 $U \times_X V \to U \times V$ 是闭浸入。仿射开集
$U \times V$ 构成 $W \times W$ 的仿射开覆盖
（“概形”引理 \ref{schemes-lemma-affine-covering-fibre-product}），
故结论由“态射”引理 \ref{morphisms-lemma-closed-immersion} 得出。
\end{proof}







\section{代数空间的点}
\label{spaces-properties-section-points}

\noindent
由“空间”例 \ref{spaces-example-affine-line-translation} 可清楚看出，
不应把代数空间的点定义为从域谱出发的单态射。
相反，我们完全按照“概形”第 \ref{schemes-section-points} 节的说明，
把它们定义为从域谱出发的态射之等价类。

\medskip\noindent
设 $S$ 为概形，$F$ 是 $(\Sch/S)_{fppf}$ 上的预层，$K$ 为域。
考虑态射
$$
\Spec(K) \longrightarrow F.
$$
由 Yoneda 引理，该态射由元素 $p \in F(\Spec(K))$ 给出。
若存在第三个域 $\Omega$ 以及交换图表
$$
\xymatrix{
\Spec(\Omega) \ar[r] \ar[d] &
\Spec(L) \ar[d]^q \\
\Spec(K) \ar[r]^p &
F.
}
$$
则称这样的两个二元组 $(\Spec(K), p)$ 和 $(\Spec(L), q)$
{\it 等价}。换言之，存在域扩张 $K \to \Omega$ 和 $L \to \Omega$，
使 $p$ 与 $q$ 映到 $F(\Spec(\Omega))$ 中的同一元素。
略去该关系确为等价关系的验证。

\begin{definition}
\label{spaces-properties-definition-points}
设 $S$ 为概形，$X$ 为 $S$ 上的代数空间。$X$ 的一个{\it 点}，
是从域谱到 $X$ 的态射之等价类。$X$ 的点集记作 $|X|$。
\end{definition}

\noindent
注意，若 $f : X \to Y$ 是 $S$ 上代数空间的态射，则它诱导映射
$|f| : |X| \to |Y|$，把代表元 $x : \Spec(K) \to X$
送到代表元 $f \circ x : \Spec(K) \to Y$。

\begin{lemma}
\label{spaces-properties-lemma-scheme-points}
设 $S$ 为概形，$X$ 为 $S$ 上的概形。把 $X$ 视作概形所得的点，
与把 $X$ 视作代数空间所得的点之间有典范一一对应。
\end{lemma}

\begin{proof}
这就是“概形”引理 \ref{schemes-lemma-characterize-points}。
\end{proof}

\begin{lemma}
\label{spaces-properties-lemma-points-cartesian}
设 $S$ 为概形，并设
$$
\xymatrix{
Z \times_Y X \ar[r] \ar[d] & X \ar[d] \\
Z \ar[r] & Y
}
$$
是 $S$ 上代数空间的 Cartesian 图表。则点集映射
$$
|Z \times_Y X|
\longrightarrow
|Z| \times_{|Y|} |X|
$$
是满射。
\end{lemma}

\begin{proof}
具体地，设给定域 $K$、$L$ 以及态射
$\Spec(K) \to X$、$\Spec(L) \to Z$。它们作为 $|Y|$ 中的元素相同，
意味着存在公共扩张 $M/K$ 和 $M/L$，使得
$\Spec(M) \to \Spec(K) \to X \to Y$ 与
$\Spec(M) \to \Spec(L) \to Z \to Y$ 相同。
这恰好是得到态射 $\Spec(M) \to Z \times_Y X$ 的条件。
\end{proof}

\begin{lemma}
\label{spaces-properties-lemma-characterize-surjective}
设 $S$ 为概形，$X$ 为 $S$ 上的代数空间，并设
$f : T \to X$ 是从概形到 $X$ 的态射。下列条件等价：
\begin{enumerate}
\item 按“空间”定义
\ref{spaces-definition-relative-representable-property}，
$f : T \to X$ 是满射；
\item $|f| : |T| \to |X|$ 是满射。
\end{enumerate}
\end{lemma}

\begin{proof}
假设 (1)。设 $x : \Spec(K) \to X$ 是从域谱到 $X$ 的态射。
由假设，概形态射 $\Spec(K) \times_X T \to \Spec(K)$ 是满射。
因此存在域扩张 $K'/K$ 以及态射
$\Spec(K') \to \Spec(K) \times_X T$，使下图左侧方块交换：
$$
\xymatrix{
\Spec(K') \ar[r] \ar[d] &
\Spec(K) \times_X T \ar[d] \ar[r] &
T \ar[d] \\
\Spec(K) \ar@{=}[r] &
\Spec(K) \ar[r]^-x & X
}
$$
这表明 $|f| : |T| \to |X|$ 是满射。

\medskip\noindent
假设 (2)。设 $Z \to X$ 为态射，其中 $Z$ 是概形。我们必须证明概形态射
$Z \times_X T \to T$ 是满射，即 $|Z \times_X T| \to |Z|$ 是满射。
这由 (2) 与引理 \ref{spaces-properties-lemma-points-cartesian} 得出。
\end{proof}

\begin{lemma}
\label{spaces-properties-lemma-points-presentation}
设 $S$ 为概形，$X$ 为 $S$ 上的代数空间，并设 $X = U/R$
为 $X$ 的一个表示；见“空间”定义 \ref{spaces-definition-presentation}。
则 $|R| \to |U| \times |U|$ 的像是等价关系，并且 $|X|$
是 $|U|$ 对该等价关系的商。
\end{lemma}

\begin{proof}
该假设意味着：$U$ 是概形，$p : U \to X$ 是满 \'etale 态射，
$R = U \times_X U$ 是概形，并在 $U$ 上定义一个 \'etale 等价关系，
使得作为层有 $X = U/R$。由引理
\ref{spaces-properties-lemma-characterize-surjective}，$|U| \to |X|$ 是满射。
由引理 \ref{spaces-properties-lemma-points-cartesian}，映射
$$
|R| \longrightarrow |U| \times_{|X|} |U|
$$
是满射。因此，$|R| \to |U| \times |U|$ 的像恰好是所有二元组
$(u_1, u_2) \in |U| \times |U|$ 所成的集合，其中 $u_1$ 与 $u_2$
在 $|X|$ 中的像相同。合并这两个陈述即得本引理。
\end{proof}

\begin{lemma}
\label{spaces-properties-lemma-topology-points}
设 $S$ 为概形。$S$ 上代数空间的点集上存在唯一的拓扑，满足下列性质：
\begin{enumerate}
\item 若 $X$ 是 $S$ 上的概形，则 $|X|$ 上的拓扑是通常拓扑
（经由引理 \ref{spaces-properties-lemma-scheme-points} 的识别）；
\item 对每个 $S$ 上代数空间的态射 $X \to Y$，映射
$|X| \to |Y|$ 连续；
\item 对每个 \'etale 态射 $U \to X$，若 $U$ 是概形，则拓扑空间映射
$|U| \to |X|$ 连续且开。
\end{enumerate}
\end{lemma}

\begin{proof}
设 $X$ 是 $S$ 上的代数空间，并设 $p : U \to X$ 为满 \'etale 态射，
其中 $U$ 是 $S$ 上的概形。定义 $W \subset |X|$ 为开集，当且仅当
$|p|^{-1}(W)$ 是 $|U|$ 的开子集。这在 $|X|$ 上定义了一个拓扑
（即 $|X|$ 上的商拓扑；见“拓扑”引理 \ref{topology-lemma-quotient}）。

\medskip\noindent
下面证明该拓扑与表示的选择无关。为此，只需说明：若 $U'$ 是概形，
且 $U' \to X$ 是 \'etale 态射，则映射 $|U'| \to |X|$
（其中 $|X|$ 的拓扑按上述 $U \to X$ 定义）开且连续。
这还将证明 (3) 成立。令 $U'' = U \times_X U'$，于是有交换图表
$$
\xymatrix{
U'' \ar[r] \ar[d] & U' \ar[d] \\
U \ar[r] & X
}
$$
由于 $U \to X$ 和 $U' \to X$ 都是 \'etale 态射，
$U'' \to U$ 与 $U'' \to U'$ 都是概形的 \'etale 态射。
此外，$U'' \to U'$ 是满射。因此得到点集映射的交换图表
$$
\xymatrix{
|U''| \ar[r] \ar[d] & |U'| \ar[d] \\
|U| \ar[r] & |X|
}
$$
下方横箭头是满射（见引理 \ref{spaces-properties-lemma-characterize-surjective}
或引理 \ref{spaces-properties-lemma-points-presentation}），并且按 $|X|$ 上拓扑的定义连续。
由“态射”引理 \ref{morphisms-lemma-etale-open}，上方横箭头满、连续且开；
左侧竖箭头也连续且开（再次使用“态射”引理
\ref{morphisms-lemma-etale-open}）。于是形式地可知，
右侧竖箭头连续且开。

\medskip\noindent
最后证明 (2)。设 $a : X \to Y$ 为代数空间的态射。根据“空间”引理
\ref{spaces-lemma-lift-morphism-presentations}，可以找到图表
$$
\xymatrix{
U \ar[d]_p \ar[r]_\alpha & V \ar[d]^q \\
X \ar[r]^a & Y
}
$$
其中 $U$ 与 $V$ 是概形，$p$ 与 $q$ 是满 \'etale 态射。它给出图表
$$
\xymatrix{
|U| \ar[d]_p \ar[r]_\alpha & |V| \ar[d]^q \\
|X| \ar[r]^a & |Y|
}
$$
其中除下方横箭头外的所有箭头均已知连续，并且两个竖箭头都满且开。
由此得到所需的下方横箭头连续。
\end{proof}

\begin{definition}
\label{spaces-properties-definition-topological-space}
设 $S$ 为概形，$X$ 为 $S$ 上的代数空间。$X$ 的底层{\it 拓扑空间}
是点集 $|X|$ 连同引理 \ref{spaces-properties-lemma-topology-points} 所构造的拓扑。
\end{definition}

\noindent
事实证明，该拓扑空间携带的信息与“空间”定义
\ref{spaces-definition-small-Zariski-site} 的小 Zariski 站点
$X_{Zar}$ 相同。

\begin{lemma}
\label{spaces-properties-lemma-open-subspaces}
设 $S$ 为概形，$X$ 为 $S$ 上的代数空间。
\begin{enumerate}
\item 规则 $X' \mapsto |X'|$ 在 $X$ 的开子空间 $X'$
（见“空间”定义 \ref{spaces-definition-immersion}）与拓扑空间 $|X|$
的开集之间定义一个保持包含关系的双射。
\item $X$ 的开子空间族 $\{X_i \subset X\}_{i \in I}$
是 Zariski 覆盖（见“空间”定义
\ref{spaces-definition-Zariski-open-covering}），当且仅当
$|X| = \bigcup |X_i|$。
\end{enumerate}
换言之，$X$ 的小 Zariski 站点 $X_{Zar}$ 典范地等同于与拓扑空间
$|X|$ 相关联的站点（见“站点”例 \ref{sites-example-site-topological}）。
\end{lemma}

\begin{proof}
为证明 (1)，构造该规则的逆。设 $W \subset |X|$ 为开集。
选取表示 $X = U/R$，它对应于满 \'etale 映射 $p : U \to X$
及 \'etale 映射 $s, t : R \to U$。由构造，$|p|^{-1}(W)$ 是 $U$ 的开集。
以 $W' \subset U$ 表示相应的开子概形。显然，
$R' = s^{-1}(W') = t^{-1}(W')$ 是 $R$ 的 Zariski 开集，
并在 $W'$ 上定义 \'etale 等价关系。由“空间”引理
\ref{spaces-lemma-finding-opens}，态射 $X' = W'/R' \to X$ 是开浸入。
再由“空间”引理 \ref{spaces-lemma-representable-over-space}，
$X'$ 是代数空间。按构造 $|X'| = W$，即 $X'$ 是与 $W$ 对应的
$X$ 的子空间。这证明了 (1)。

\medskip\noindent
为证明 (2)，注意：若 $\{X_i \subset X\}_{i \in I}$ 是开子空间族，
则它是 Zariski 覆盖，当且仅当
$U = \bigcup U \times_X X_i$ 是开覆盖。这由 Zariski 覆盖的定义以及
态射 $U \to X$ 作为 $(\Sch/S)_{fppf}$ 上预层的映射是满射这一事实得出。
另一方面，由引理 \ref{spaces-properties-lemma-points-presentation}（以及投影
$U \times_X X_i \to X_i$ 满且 \'etale 这一事实），
$|X| = \bigcup |X_i|$ 当且仅当
$U = \bigcup U \times_X X_i$。于是得到 (2) 的等价性。
\end{proof}

\begin{lemma}
\label{spaces-properties-lemma-factor-through-open-subspace}
设 $S$ 为概形，$X$、$Y$ 为 $S$ 上的代数空间，$X' \subset X$
为开子空间，并设 $f : Y \to X$ 为 $S$ 上代数空间的态射。
则 $f$ 经过 $X'$ 分解，当且仅当 $|f| : |Y| \to |X|$
经过 $|X'| \subset |X|$ 分解。
\end{lemma}

\begin{proof}
由“空间”引理 \ref{spaces-lemma-base-change-immersions}，
$Y' = Y \times_X X' \to Y$ 是开浸入。若 $|f|(|Y|) \subset |X'|$，
则显然 $|Y'| = |Y|$。故由引理 \ref{spaces-properties-lemma-open-subspaces}，
$Y' = Y$。
\end{proof}

\begin{lemma}
\label{spaces-properties-lemma-etale-image-open}
设 $S$ 为概形，$X$ 为 $S$ 上的代数空间。设 $U$ 为概形，
$f : U \to X$ 为 \'etale 态射。令 $X' \subset X$ 为通过引理
\ref{spaces-properties-lemma-open-subspaces} 与开集 $|f|(|U|) \subset |X|$ 对应的开子空间。
则 $f$ 经过满 \'etale 态射 $f' : U \to X'$ 分解。
此外，若 $R = U \times_X U$，则 $R = U \times_{X'} U$，
且 $X'$ 有表示 $X' = U/R$。
\end{lemma}

\begin{proof}
分解的存在性由引理 \ref{spaces-properties-lemma-factor-through-open-subspace} 得出。
根据引理 \ref{spaces-properties-lemma-characterize-surjective}，态射 $f'$ 是满射。
为说明 $f'$ 是 \'etale 态射，设 $T \to X'$ 为态射，其中 $T$ 是概形。
由于 $X' \to X$ 是层的单态射，有
$T \times_X U = T \times_{X'} U$。又因假设 $f$ 为 \'etale 态射，
投影 $T \times_{X'} U \to T$ 是 \'etale 态射。
由于 $X' \to X$ 是单态射，还有
$U \times_X U = U \times_{X'} U$。最后，
$X' = U/R$ 由“空间”引理 \ref{spaces-lemma-space-presentation} 得出。
\end{proof}

\begin{lemma}
\label{spaces-properties-lemma-equivalence-class-point-monomorphism}
设 $S$ 为概形，$X$ 为 $S$ 上的代数空间。设
$p : \Spec(K) \to X$ 与 $q : \Spec(L) \to X$ 为态射，
其中 $K$ 和 $L$ 是域。假设 $p$ 与 $q$ 确定 $|X|$ 中的同一点，
且 $p$ 是单态射。则 $q$ 唯一地经过 $p$ 分解。
\end{lemma}

\begin{proof}
由于 $p$ 与 $q$ 定义 $|X|$ 中的同一点，概形
$$
Y = \Spec(K) \times_{p, X, q} \Spec(L)
$$
非空。单态射的基变换仍是单态射，故投影态射 $Y \to \Spec(L)$
是单态射。因此 $Y = \Spec(L)$；见“概形”引理
\ref{schemes-lemma-mono-towards-spec-field}。由此可知 $q$ 经过 $p$ 分解。
唯一性来自 $p$ 是单态射这一事实。
\end{proof}

\begin{lemma}
\label{spaces-properties-lemma-points-monomorphism}
设 $S$ 为概形，$X$ 为 $S$ 上的代数空间。考虑映射
$$
\{\Spec(k) \to X \text{ 为单态射，其中 }k\text{ 是域}\}
\longrightarrow
|X|
$$
该映射是单射。
\end{lemma}

\begin{proof}
这由引理 \ref{spaces-properties-lemma-equivalence-class-point-monomorphism} 得出。
\end{proof}

\noindent
我们将在“良好空间”引理
\ref{decent-spaces-lemma-decent-points-monomorphism} 中看到：
当 $X$ 良好时，引理 \ref{spaces-properties-lemma-points-monomorphism} 的映射是双射。

















\section{拟紧代数空间}
\label{spaces-properties-section-quasi-compact}

\begin{definition}
\label{spaces-properties-definition-quasi-compact}
设 $S$ 为概形，$X$ 为 $S$ 上的代数空间。若存在满 \'etale 态射
$U \to X$，其中 $U$ 是拟紧概形，则称 $X$ {\it 拟紧}。
\end{definition}

\begin{lemma}
\label{spaces-properties-lemma-quasi-compact-space}
设 $S$ 为概形，$X$ 为 $S$ 上的代数空间。则 $X$ 拟紧，当且仅当
$|X|$ 拟紧。
\end{lemma}

\begin{proof}
选取概形 $U$ 以及满 \'etale 态射 $U \to X$。我们将使用引理
\ref{spaces-properties-lemma-characterize-surjective}。若 $U$ 拟紧，则因
$|U| \to |X|$ 是满射，可知 $|X|$ 拟紧。若 $|X|$ 拟紧，
则因 $|U| \to |X|$ 是开映射，存在拟紧开集 $U' \subset U$，
使 $|U'| \to |X|$ 是满射（并且仍为 \'etale 态射）。结论得证。
\end{proof}

\begin{lemma}
\label{spaces-properties-lemma-finite-disjoint-quasi-compact}
有限个拟紧代数空间的不交并是拟紧代数空间。
\end{lemma}

\begin{proof}
这由引理 \ref{spaces-properties-lemma-quasi-compact-space} 及相应的拓扑事实立即得出。
\end{proof}

\begin{example}
\label{spaces-properties-example-quasi-compact-not-very-reasonable}
空间 $\mathbf{A}^1_{\mathbf{Q}}/\mathbf{Z}$ 是拟紧代数空间。
\end{example}

\begin{lemma}
\label{spaces-properties-lemma-space-locally-quasi-compact}
设 $S$ 为概形，$X$ 为 $S$ 上的代数空间。$|X|$ 的每个点都有一个
由拟紧开邻域组成的基本邻域系。特别地，$|X|$ 按“拓扑”定义
\ref{topology-definition-locally-quasi-compact} 的意义局部拟紧。
\end{lemma}

\begin{proof}
这形式地来自下述事实：存在概形 $U$ 以及拓扑空间的满、开、连续映射
$U \to |X|$。更精确地说，若 $u \in U$ 映到 $x \in |X|$，
则 $u$ 的仿射邻域之像给出 $x$ 的一个拟紧开基本邻域系。
\end{proof}







\section{特殊覆盖}
\label{spaces-properties-section-special-coverings}

\noindent
本节汇集若干关于代数空间存在满 \'etale 覆盖的直接引理。

\begin{lemma}
\label{spaces-properties-lemma-cover-by-union-affines}
设 $S$ 为概形，$X$ 为 $S$ 上的代数空间。存在满 \'etale 态射
$U \to X$，其中 $U$ 是仿射概形的不交并。还可以假设其中每个仿射概形
都映入 $S$ 的某个仿射开集。
\end{lemma}

\begin{proof}
设 $V \to X$ 为满 \'etale 态射。取 Zariski 开覆盖
$V = \bigcup_{i \in I} V_i$，使每个 $V_i$ 都映入 $S$ 的某个仿射开集。
令 $U = \coprod_{i \in I} V_i$，并取诱导态射 $U \to V \to X$。
它是 \'etale 且满的，因为它是函子的 \'etale、满、可表示变换之复合；
这里使用一般原理“空间”引理
\ref{spaces-lemma-composition-representable-transformations-property}，
以及“态射”引理 \ref{morphisms-lemma-composition-surjective} 和
\ref{morphisms-lemma-composition-etale}。
\end{proof}

\begin{lemma}
\label{spaces-properties-lemma-union-of-quasi-compact}
设 $S$ 为概形，$X$ 为 $S$ 上的代数空间。存在 Zariski 覆盖
$X = \bigcup X_i$，使每个代数空间 $X_i$ 都有一个由仿射概形给出的
满 \'etale 覆盖。还可以假设每个 $X_i$ 都映入 $S$ 的某个仿射开集。
\end{lemma}

\begin{proof}
由引理 \ref{spaces-properties-lemma-cover-by-union-affines}，可找到满 \'etale 态射
$U = \coprod U_i \to X$，其中 $U_i$ 仿射且映入 $S$ 的某个仿射开集。
令 $X_i \subset X$ 为 $X$ 的开子空间，使 $U_i \to X$ 经过满
\'etale 态射 $U_i \to X_i$ 分解；见引理
\ref{spaces-properties-lemma-etale-image-open}。由于 $U = \bigcup U_i$，
可知 $X = \bigcup X_i$。又因 $U_i \to X_i$ 是满射，
$X_i \to S$ 的像位于 $S$ 的某个仿射开集中。
\end{proof}

\begin{lemma}
\label{spaces-properties-lemma-quasi-compact-affine-cover}
设 $S$ 为概形，$X$ 为 $S$ 上的代数空间。则 $X$ 拟紧，当且仅当
存在满 \'etale 态射 $U \to X$，其中 $U$ 是仿射概形。
\end{lemma}

\begin{proof}
若存在满 \'etale 态射 $U \to X$，且 $U$ 仿射，则由定义
\ref{spaces-properties-definition-quasi-compact}，$X$ 拟紧。反之，若 $X$ 拟紧，
则 $|X|$ 拟紧。令 $U = \coprod_{i \in I} U_i$ 为仿射概形的不交并，
并带有满 \'etale 映射 $\varphi : U \to X$
（引理 \ref{spaces-properties-lemma-cover-by-union-affines}）。于是
$|X| = \bigcup \varphi(|U_i|)$。由拟紧性，存在有限子集
$i_1, \ldots, i_n$，使 $|X| = \bigcup \varphi(|U_{i_j}|)$。
因此 $U_{i_1} \cup \ldots \cup U_{i_n}$ 是带有到 $X$
的有限满态射的仿射概形。
\end{proof}

\noindent
下面的引理将被“代数空间的态射”一章中关于分离态射的讨论取代。

\begin{lemma}
\label{spaces-properties-lemma-separated-cover}
设 $S$ 为概形，$X$ 为 $S$ 上的代数空间。设 $U$ 是分离概形，
且 $U \to X$ 为 \'etale 态射。则 $U \to X$ 分离，并且
$R = U \times_X U$ 是分离概形。
\end{lemma}

\begin{proof}
令 $X' \subset X$ 为开子概形，使 $U \to X$ 经过满 \'etale 态射
$U \to X'$ 分解；见引理 \ref{spaces-properties-lemma-etale-image-open}。
若 $U \to X'$ 分离，则 $U \to X$ 也分离；见“空间”引理
\ref{spaces-lemma-composition-representable-transformations-property}。
这里，开浸入 $X' \to X$ 分离，由“空间”引理
\ref{spaces-lemma-representable-transformations-property-implication}
及“概形”引理 \ref{schemes-lemma-immersions-monomorphisms} 可知。
此外，由于 $U \times_{X'} U = U \times_X U$，只需用 $X'$ 替换 $X$
后证明结论；换言之，可以假设 $U \to X$ 满。考虑交换图表
$$
\xymatrix{
R = U \times_X U \ar[r] \ar[d] & U \ar[d] \\
U \ar[r] & X
}
$$
在“空间”引理 \ref{spaces-lemma-properties-diagonal} 的证明中，
我们已看到 $j : R \to U \times_S U$ 分离。由于 $U$ 是分离概形，
概形态射 $U \to S$ 分离；见“概形”引理
\ref{schemes-lemma-compose-after-separated}。因此，作为基变换，
$U \times_S U \to U$ 分离；见“概形”引理
\ref{schemes-lemma-separated-permanence}。故概形 $U \times_S U$ 分离
（仍由该引理）。又因 $j$ 分离，以同样方式可知 $R$ 分离。
于是 $R \to U$ 是分离态射（再次使用“概形”引理
\ref{schemes-lemma-compose-after-separated}）。因此，由“空间”引理
\ref{spaces-lemma-representable-morphisms-spaces-property} 与上图，
可知 $U \to X$ 分离。
\end{proof}

\begin{lemma}
\label{spaces-properties-lemma-quasi-separated}
设 $S$ 为概形，$X$ 为 $S$ 上的代数空间。若存在拟分离概形 $U$
以及满 \'etale 态射 $U \to X$，使两个投影
$U \times_X U \to U$ 中的任一个拟紧，则 $X$ 拟分离。
\end{lemma}

\begin{proof}
可把 $X$ 视作 $\mathbf{Z}$ 上的代数空间。考虑 Cartesian 图表
$$
\xymatrix{
U \times_X U \ar[r] \ar[d]_j & X \ar[d]^\Delta \\
U \times U \ar[r] & X \times X
}
$$
由于 $U$ 拟分离，投影 $U \times U \to U$ 拟分离
（它是概形的拟分离态射的基变换；见“概形”引理
\ref{schemes-lemma-separated-permanence}）。因此，由本引理的假设与
“概形”引理 \ref{schemes-lemma-quasi-compact-permanence}，$j$ 拟紧。
再由“空间”引理 \ref{spaces-lemma-representable-morphisms-spaces-property}，
得到所需的 $\Delta$ 拟紧。
\end{proof}

\begin{lemma}
\label{spaces-properties-lemma-quasi-separated-quasi-compact-pieces}
设 $S$ 为概形，$X$ 为 $S$ 上的代数空间。下列条件等价：
\begin{enumerate}
\item $X$ 在 $S$ 上 Zariski 局部拟分离；
\item $X$ Zariski 局部拟分离；
\item 存在 Zariski 开覆盖 $X = \bigcup X_i$，使对每个 $i$，
都存在仿射概形 $U_i$ 以及拟紧的满 \'etale 态射 $U_i \to X_i$；
\item 存在 Zariski 开覆盖 $X = \bigcup X_i$，使对每个 $i$，
都存在映入 $S$ 的某个仿射开集的仿射概形 $U_i$，
以及拟紧的满 \'etale 态射 $U_i \to X_i$。
\end{enumerate}
\end{lemma}

\begin{proof}
假设 $U_i \to X_i \subset X$ 如 (3) 所述。为证明 (4)，
对每个 $i$ 选取有限仿射开覆盖
$U_i = U_{i1} \cup \ldots \cup U_{in_i}$，使每个 $U_{ij}$
都映入 $S$ 的某个仿射开集。复合
$U_{ij} \to U_i \to X_i$ 是 \'etale 且拟紧的
（见“空间”引理
\ref{spaces-lemma-composition-representable-transformations-property}）。
令 $X_{ij} \subset X_i$ 为与 $|U_{ij}| \to |X_i|$ 的像相对应的开子空间；
见引理 \ref{spaces-properties-lemma-etale-image-open}。注意，$X_{ij} \subset X_i$
是单态射，而 $U_{ij} \to X$ 拟紧，故 $U_{ij} \to X_{ij}$ 拟紧。
于是 $X = \bigcup X_{ij}$ 是 (4) 所述的覆盖。
(4) $\Rightarrow$ (3) 是立即的。

\medskip\noindent
假设 (4)。为证明 $X$ 在 $S$ 上 Zariski 局部拟分离，
只需证明 $X_i$ 在 $S$ 上拟分离。因此可假设存在映入 $S$
某个仿射开集的仿射概形 $U$，以及拟紧的满 \'etale 态射
$U \to X$。考虑纤维积方块
$$
\xymatrix{
U \times_X U \ar[r] \ar[d] & U \times_S U \ar[d] \\
X \ar[r]^-{\Delta_{X/S}} & X \times_S X
}
$$
右侧竖箭头是满 \'etale 态射（见“空间”引理
\ref{spaces-lemma-product-representable-transformations-property}）。
又因 $U$ 映入 $S$ 的某个仿射开集，$U \times_S U$ 仿射
（见“概形”第 \ref{schemes-section-fibre-products} 节）。
此外，投影 $U \times_X U \to U$ 是 $U \to X$ 的基变换，因而拟紧；
故 $U \times_X U$ 拟紧。由“空间”引理
\ref{spaces-lemma-representable-morphisms-spaces-property}，
得到所需的 $\Delta_{X/S}$ 拟紧。

\medskip\noindent
假设 (1)。为证明 (3)，可立即归约到 $X$ 在 $S$ 上拟分离的情形。
由引理 \ref{spaces-properties-lemma-union-of-quasi-compact}，可找到 Zariski 开覆盖
$X = \bigcup X_i$，使每个 $X_i$ 都映入 $S$ 的某个仿射开集，
并且存在仿射概形 $U_i$ 以及满 \'etale 态射 $U_i \to X_i$。
由于 $U_i \to S$ 映入 $S$ 的某个仿射开集，
$U_i \times_S U_i$ 仿射；见“概形”第
\ref{schemes-section-fibre-products} 节。由于 $X$ 在 $S$ 上拟分离，态射
$$
R_i = U_i \times_{X_i} U_i = U_i \times_X U_i
\longrightarrow
U_i \times_S U_i
$$
作为 $\Delta_{X/S}$ 的基变换均拟紧。因此 $R_i$ 是拟紧概形，
进而每个投影 $R_i \to U_i$ 拟紧。于是，对覆盖 $U_i \to X_i$
及态射 $U_i \to X_i$ 应用“空间”引理
\ref{spaces-lemma-representable-morphisms-spaces-property}，
可知态射 $U_i \to X_i$ 按所需拟紧。

\medskip\noindent
至此，(1)、(3) 与 (4) 等价。由于 (3) 不涉及基概形，
可知它们也与 (2) 等价。
\end{proof}

\noindent
下面的引理将十分有用。

\begin{lemma}
\label{spaces-properties-lemma-finite-fibres-presentation}
设 $S$ 为概形，$X$ 为 $S$ 上的代数空间，$U$ 为概形。
设 $\varphi : U \to X$ 为 \'etale 态射，并且投影
$R = U \times_X U \to U$ 拟紧；例如，$\varphi$ 拟紧时即如此。
则下列映射的纤维
$$
|U| \to |X|
\quad\text{且}\quad
|R| \to |X|
$$
都是有限集。
\end{lemma}

\begin{proof}
记 $R = U \times_X U$，并以 $s, t : R \to U$ 表示两个投影。
设 $u \in U$ 为点，$x \in |X|$ 为其像。由引理
\ref{spaces-properties-lemma-points-cartesian}，$|U| \to |X|$ 在 $x$ 上的纤维等于
$s(t^{-1}(\{u\}))$，而 $|R| \to |X|$ 在 $x$ 上的纤维是
$t^{-1}(s(t^{-1}(\{u\})))$。由于 $t : R \to U$ 是拟紧
\'etale 态射，其纤维有限；事实上，由“态射”引理
\ref{morphisms-lemma-etale-over-field}，其纤维是域谱的不交并，
并且拟紧。结论得证。
\end{proof}








\section{由概形性质定义的空间性质}
\label{spaces-properties-section-types-properties}

\noindent
概形的任何 \'etale 局部性质都通过下述引理给出代数空间的相应性质。

\begin{lemma}
\label{spaces-properties-lemma-type-property}
设 $S$ 为概形，$X$ 为 $S$ 上的代数空间。设 $\mathcal{P}$
是对 \'etale 拓扑为局部的概形性质；见“下降”定义
\ref{descent-definition-property-local}。下列条件等价：
\begin{enumerate}
\item 存在概形 $U$ 及满 \'etale 态射 $U \to X$，使概形 $U$
具有性质 $\mathcal{P}$；
\item 对每个概形 $U$ 及每个 \'etale 态射 $U \to X$，概形 $U$
都具有性质 $\mathcal{P}$。
\end{enumerate}
若 $X$ 可表示，则这些条件等价于 $\mathcal{P}(X)$。
\end{lemma}

\begin{proof}
(2) $\Rightarrow$ (1) 是立即的。反之，选取满 \'etale 态射
$U \to X$，其中 $U$ 是具有 $\mathcal{P}$ 的概形，并设 $V$
为 \'etale $X$-概形。于是 $U \times_X V \rightarrow V$
是概形的满 \'etale 态射，故 $V$ 从 $U \times_X V$ 继承
$\mathcal{P}$，后者又从 $U$ 继承 $\mathcal{P}$（见“下降”定义
\ref{descent-definition-property-local} 之后的讨论）。最后的断言由
(1) 及“下降”定义 \ref{descent-definition-property-local} 即得。
\end{proof}

\begin{definition}
\label{spaces-properties-definition-type-property}
设 $\mathcal{P}$ 为对 \'etale 拓扑局部的概形性质。设 $S$ 为概形，
$X$ 为 $S$ 上的代数空间。若引理 \ref{spaces-properties-lemma-type-property}
中的任一等价条件成立，则称 $X$ {\it 具有性质 $\mathcal{P}$}。
\end{definition}

\begin{remark}
\label{spaces-properties-remark-list-properties-local-etale-topology}
下列性质对 \'etale 拓扑是局部的（应注意，fpqc、fppf、syntomic
及光滑拓扑都强于 \'etale 拓扑）：
\begin{enumerate}
\item 局部 Noether；见“下降”引理
\ref{descent-lemma-Noetherian-local-fppf}；
\item Jacobson；见“下降”引理
\ref{descent-lemma-Jacobson-local-fppf}；
\item 局部 Noether 且满足 $(S_k)$；见“下降”引理
\ref{descent-lemma-Sk-local-syntomic}；
\item Cohen--Macaulay；见“下降”引理
\ref{descent-lemma-CM-local-syntomic}；
\item Gorenstein；见“概形上的对偶性”引理
\ref{duality-lemma-gorenstein-local-syntomic}；
\item 既约；见“下降”引理
\ref{descent-lemma-reduced-local-smooth}；
\item 正规；见“下降”引理
\ref{descent-lemma-normal-local-smooth}；
\item 局部 Noether 且满足 $(R_k)$；见“下降”引理
\ref{descent-lemma-Rk-local-smooth}；
\item 正则；见“下降”引理
\ref{descent-lemma-regular-local-smooth}；
\item Nagata；见“下降”引理
\ref{descent-lemma-Nagata-local-smooth}。
\end{enumerate}
\end{remark}

\noindent
概形芽的任何 \'etale 局部性质都给出代数空间的相应性质。下面是必需的引理。

\begin{lemma}
\label{spaces-properties-lemma-local-source-target-at-point}
设 $\mathcal{P}$ 是概形芽的 \'etale 局部性质；见“下降”定义
\ref{descent-definition-local-at-point}。设 $S$ 为概形，$X$ 为
$S$ 上的代数空间，$x \in |X|$ 为 $X$ 的一点。考虑以概形 $U$
为源的 \'etale 态射 $a : U \to X$。下列条件等价：
\begin{enumerate}
\item 对上述任意 $U \to X$ 及满足 $a(u) = x$ 的 $u \in U$，
都有 $\mathcal{P}(U, u)$；
\item 对上述某个 $U \to X$ 及某个满足 $a(u) = x$ 的 $u \in U$，
有 $\mathcal{P}(U, u)$。
\end{enumerate}
若 $X$ 可表示，则这些条件等价于 $\mathcal{P}(X, x)$。
\end{lemma}

\begin{proof}
从略。
\end{proof}

\begin{definition}
\label{spaces-properties-definition-property-at-point}
设 $S$ 为概形，$X$ 为 $S$ 上的代数空间，$x \in |X|$。设
$\mathcal{P}$ 为概形芽的 \'etale 局部性质。若引理
\ref{spaces-properties-lemma-local-source-target-at-point} 中任一等价条件成立，
则称 $X$ {\it 在 $x$ 处具有性质 $\mathcal{P}$}。
\end{definition}

\begin{remark}
\label{spaces-properties-remark-list-properties-local-ring-local-etale-topology}
设 $P$ 为局部环的性质。假设对任意 \'etale 环映射 $A \to B$，
以及 $B$ 中位于 $A$ 的素理想 $\mathfrak p$ 之上的素理想
$\mathfrak q$，都有
$P(A_\mathfrak p) \Leftrightarrow P(B_\mathfrak q)$.
规定 $\mathcal{P}(U, u) = P(\mathcal{O}_{U, u})$，便得到概形芽
$(U, u)$ 的一个 \'etale 局部性质。在这种情形下，我们用“$X$
在 $x$ 处的局部环具有 $P$”表示 $X$ 在 $x$ 处具有性质
$\mathcal{P}$。下列各项都是这样的性质 $P$：
\begin{enumerate}
\item Noether；见“代数进阶”引理
\ref{more-algebra-lemma-Noetherian-etale-extension}；
\item 维数为 $d$；见“代数进阶”引理
\ref{more-algebra-lemma-dimension-etale-extension}；
\item 正则；见“代数进阶”引理
\ref{more-algebra-lemma-regular-etale-extension}；
\item 离散赋值环；这由 (2)、(3) 及“代数”引理
\ref{algebra-lemma-characterize-dvr} 得出；
\item 既约；见“代数进阶”引理
\ref{more-algebra-lemma-henselization-reduced}；
\item 正规；见“代数进阶”引理
\ref{more-algebra-lemma-henselization-normal}；
\item Noether 且深度为 $k$；见“代数进阶”引理
\ref{more-algebra-lemma-henselization-depth}；
\item Noether 且 Cohen--Macaulay；见“代数进阶”引理
\ref{more-algebra-lemma-henselization-CM}；
\item Noether 且 Gorenstein；见“对偶复形”引理
\ref{dualizing-lemma-flat-under-gorenstein}。
\end{enumerate}
还有更多满足这一条件的性质，例如 G-环与 Nagata。若我们以后需要，
会将它们以及相应代数事实之详细证明的参考文献一并加入这里。
\end{remark}







\section{可构造集}
\label{spaces-properties-section-constructible}

\begin{lemma}
\label{spaces-properties-lemma-locally-constructible}
设 $S$ 为概形，$X$ 为 $S$ 上的代数空间，$E \subset |X|$
为子集。下列条件等价：
\begin{enumerate}
\item 对源 $U$ 为概形的每个 \'etale 态射 $U \to X$，$E$
在 $U$ 中的逆像都是 $U$ 的局部可构造子集；
\item 对源 $U$ 为仿射概形的每个 \'etale 态射 $U \to X$，$E$
在 $U$ 中的逆像都是 $U$ 的可构造子集；
\item 存在源 $U$ 为概形的满 \'etale 态射 $U \to X$，使 $E$
在 $U$ 中的逆像为 $U$ 的局部可构造子集。
\end{enumerate}
\end{lemma}

\begin{proof}
由“性质”引理 \ref{properties-lemma-locally-constructible}，
(1) 与 (2) 等价。(1) 蕴含 (3) 是立即的。因此，设我们有满
\'etale 态射 $\varphi : U \to X$，其中 $U$ 为概形且
$\varphi^{-1}(E)$ 局部可构造。再设
$\varphi' : U' \to X$ 为另一个以概形 $U'$ 为源的 \'etale 态射。
于是
$$
E'' = \text{pr}_1^{-1}(\varphi^{-1}(E)) = \text{pr}_2^{-1}((\varphi')^{-1}(E))
$$
其中 $\text{pr}_1 : U \times_X U' \to U$ 与
$\text{pr}_2 : U \times_X U' \to U'$ 为投影。由“态射”引理
\ref{morphisms-lemma-inverse-image-constructible}，$E''$
在 $U \times_X U'$ 中局部可构造。设 $W' \subset U'$ 为仿射开集。
由于 $\text{pr}_2$ 是 \'etale 的，因而是开映射，故可选取拟紧开集
$W'' \subset U \times_X U'$，使 $\text{pr}_2(W'') = W'$。
于是 $\text{pr}_2|_{W''} : W'' \to W'$ 拟紧。因为 $\varphi$
满，由引理 \ref{spaces-properties-lemma-points-cartesian}，
$W' \cap (\varphi')^{-1}(E) = \text{pr}_2(E'' \cap W'')$。
因此再由“态射”定理 \ref{morphisms-theorem-chevalley}，
$W' \cap (\varphi')^{-1}(E) = \text{pr}_2(E'' \cap W'')$
局部可构造，正合所需。
\end{proof}

\begin{definition}
\label{spaces-properties-definition-locally-constructible}
设 $S$ 为概形，$X$ 为 $S$ 上的代数空间，$E \subset |X|$
为子集。若引理 \ref{spaces-properties-lemma-locally-constructible} 中的等价条件成立，
则称 $E$ {\it \'etale 局部可构造}。
\end{definition}

\noindent
当然，若 $X$ 可表示，即 $X$ 是概形，则这恰好表示 $E$
是其底拓扑空间的局部可构造子集。







\section{一点处的维数}
\label{spaces-properties-section-dimension}

\noindent
可用“下降”引理 \ref{descent-lemma-dimension-at-point-local}
定义代数空间 $X$ 在一点 $x$ 处的维数。所得概念不同于拓扑意义下的概念
（即 $|X|$ 在 $x$ 处的维数）。

\begin{definition}
\label{spaces-properties-definition-dimension-at-point}
设 $S$ 为概形，$X$ 为 $S$ 上的代数空间，$x \in |X|$
为 $X$ 的一点。定义 {\it $X$ 在 $x$ 处的维数}为元素
$\dim_x(X) \in \{0, 1, 2, \ldots, \infty\}$，它满足：
对任意（等价地，某个）二元组 $(a : U \to X, u)$，都有
$\dim_x(X) = \dim_u(U)$；这里 $a : U \to X$ 是从概形到 $X$
的 \'etale 态射，$u \in U$ 且 $a(u) = x$。见定义
\ref{spaces-properties-definition-property-at-point}、引理
\ref{spaces-properties-lemma-local-source-target-at-point} 及“下降”引理
\ref{descent-lemma-dimension-at-point-local}。
\end{definition}

\noindent
警告：一般而言，{\bf 并非} $\dim_x(X) = \dim_x(|X|)$。
“空间”例 \ref{spaces-example-infinite-product} 中的代数空间 $X$
给出一个反例。具体地，设 $x \in |X|$ 是不等于 $|X|$ 的一般点
$x_0$ 的一点。此时 $\dim_x(X) = 0$，但 $\dim_x(|X|) = 1$。
特别地，$X$ 的维数（定义如下）不同于 $|X|$ 的维数。

\begin{definition}
\label{spaces-properties-definition-dimension}
设 $S$ 为概形，$X$ 为 $S$ 上的代数空间。$X$ 的{\it 维数}
$\dim(X)$ 定义为
$$
\dim(X) = \sup\nolimits_{x \in |X|} \dim_x(X)
$$
\end{definition}

\noindent
由“性质”引理 \ref{properties-lemma-dimension}，若 $X$ 为概形，
这就是通常的概念。还有另一个整数度量概形在一点处的维数，即局部环的维数。
这一不变量也与 \'etale 态射相容；见第
\ref{spaces-properties-section-dimension-local-ring} 节。






\section{局部环的维数}
\label{spaces-properties-section-dimension-local-ring}

\noindent
代数空间的局部环之维数是一个良定义的概念。

\begin{lemma}
\label{spaces-properties-lemma-pre-dimension-local-ring}
设 $S$ 为概形，$X$ 为 $S$ 上的代数空间，$x \in |X|$ 为一点，
$d \in \{0, 1, 2, \ldots, \infty\}$。下列条件等价：
\begin{enumerate}
\item 存在概形 $U$、\'etale 态射 $a : U \to X$ 及满足
$a(u) = x$ 的点 $u \in U$，使 $\dim(\mathcal{O}_{U, u}) = d$；
\item 对任意概形 $U$、任意 \'etale 态射 $a : U \to X$
及任意满足 $a(u) = x$ 的点 $u \in U$，都有
$\dim(\mathcal{O}_{U, u}) = d$。
\end{enumerate}
若 $X$ 为概形，则这些条件等价于 $\dim(\mathcal{O}_{X, x}) = d$。
\end{lemma}

\begin{proof}
结合引理 \ref{spaces-properties-lemma-local-source-target-at-point} 与“下降”引理
\ref{descent-lemma-dimension-local-ring-local}。
\end{proof}

\begin{definition}
\label{spaces-properties-definition-dimension-local-ring}
设 $S$ 为概形，$X$ 为 $S$ 上的代数空间，$x \in |X|$ 为一点。
{\it $X$ 在 $x$ 处之局部环的维数}是满足引理
\ref{spaces-properties-lemma-pre-dimension-local-ring} 中等价条件的元素
$d \in \{0, 1, 2, \ldots, \infty\}$。在这种情形下，我们也称
{\it $x$ 是 $X$ 上余维为 $d$ 的点}。
\end{definition}

\noindent
除下述引理外，还请读者参见引理 \ref{spaces-properties-lemma-dimension-local-ring}
与 \ref{spaces-properties-lemma-dimension-decent-invariant-under-etale}。

\begin{lemma}
\label{spaces-properties-lemma-dimension}
设 $S$ 为概形，$X$ 为 $S$ 上的代数空间。下列量相等：
\begin{enumerate}
\item $X$ 的维数；
\item $X$ 各局部环之维数的上确界；
\item 当 $x \in |X|$ 时，$\dim_x(X)$ 的上确界。
\end{enumerate}
\end{lemma}

\begin{proof}
由定义 \ref{spaces-properties-definition-dimension}，(1) 与 (3) 中的数相等。
设 $U \to X$ 为以概形 $U$ 为源的满 \'etale 态射。由定义，
当 $x \in |X|$ 时 $\dim_x(X)$ 的上确界，等于当 $u$ 遍历 $U$
各点时 $\dim_u(U)$ 的上确界。由“性质”引理
\ref{properties-lemma-dimension}，后者又等于
$\dim(\mathcal{O}_{U, u})$ 的上确界；而由定义，这就是 (2)。
\end{proof}




\section{一般点}
\label{spaces-properties-section-generic-points}

\noindent
设 $T$ 为拓扑空间。依照 EGA I 第二版，{\it $T$ 的极大点}是
$T$ 的某个不可约分支的一般点。若 $T = |X|$ 是与代数空间 $X$
相伴的拓扑空间，则至少有两种极大点概念：可以把 $T$ 看作拓扑空间
并取其极大点，也可以取 $U$ 的极大点在 \'etale 态射
$U \to X$ 下的像，其中 $U$ 是概形。第二种概念对应于余维 $0$ 的点集（引理
\ref{spaces-properties-lemma-codimension-0-points}）。对于一般的代数空间，
余维 $0$ 的点更便于使用；对拟分离代数空间，以及更一般的良好代数空间，
这两个概念一致（“良好空间”引理
\ref{decent-spaces-lemma-decent-generic-points}）。

\begin{lemma}
\label{spaces-properties-lemma-codimension-0-points}
设 $S$ 为概形，$X$ 为 $S$ 上的代数空间，$x \in |X|$。
考虑以概形 $U$ 为源的 \'etale 态射 $a : U \to X$。下列条件等价：
\begin{enumerate}
\item $x$ 是 $X$ 上余维为 $0$ 的点；
\item 对上述某个 $U \to X$ 及满足 $a(u) = x$ 的某个 $u \in U$，
点 $u$ 是 $U$ 的某个不可约分支的一般点；
\item 对上述任意 $U \to X$ 及映到 $x$ 的任意 $u \in U$，
点 $u$ 都是 $U$ 的某个不可约分支的一般点。
\end{enumerate}
若 $X$ 可表示，则这些条件等价于 $x$ 是 $|X|$
某个不可约分支的一般点。
\end{lemma}

\begin{proof}
注意，概形 $U$ 的一点 $u$ 是 $U$ 某个不可约分支的一般点，当且仅当
$\dim(\mathcal{O}_{U, u}) = 0$（“性质”引理
\ref{properties-lemma-generic-point}）。故结论由 $X$
上一点之余维的定义（定义 \ref{spaces-properties-definition-dimension-local-ring}）得出。
\end{proof}

\begin{lemma}
\label{spaces-properties-lemma-codimension-0-points-dense}
设 $S$ 为概形，$X$ 为 $S$ 上的代数空间。$X$ 的余维 $0$ 点集
在 $|X|$ 中稠密。
\end{lemma}

\begin{proof}
若 $U$ 为概形，则其各不可约分支之一般点的集合在 $U$ 中稠密
（这对任意拟清醒拓扑空间都成立）。因此，若 $U \to X$
为满 \'etale 态射，则 $X$ 的余维 $0$ 点集是 $|U|$ 中某个稠密子集的像
（引理 \ref{spaces-properties-lemma-codimension-0-points}）。由于 $|X|$
对 $|U| \to |X|$ 带商拓扑，结论得证。
\end{proof}





\section{既约空间}
\label{spaces-properties-section-reduced}

\noindent
我们已在第 \ref{spaces-properties-section-types-properties} 节定义了既约代数空间。
这里仅证明关于既约代数空间的几个简单引理。

\begin{lemma}
\label{spaces-properties-lemma-subspace-induced-topology}
设 $S$ 为概形，$Z \to X$ 为代数空间的浸入。则 $|Z| \to |X|$
是 $|Z|$ 到 $|X|$ 的某个局部闭子集上的同胚。
\end{lemma}

\begin{proof}
设 $U$ 为概形，$U \to X$ 为满 \'etale 态射。于是
$Z \times_X U \to U$ 是概形的浸入，因而给出 $|Z \times_X U|$
到 $|U|$ 的某个局部闭子集 $T'$ 上的同胚。由引理
\ref{spaces-properties-lemma-points-cartesian}，子集 $T'$ 是 $|Z| \to |X|$ 之像
$T$ 的逆像。映射 $|Z| \to |X|$ 为单射，因为函子变换
$Z \to X$ 是单射；见“空间”第 \ref{spaces-section-Zariski} 节。
由“拓扑”引理 \ref{topology-lemma-open-morphism-quotient-topology}，
$T$ 在 $|X|$ 中局部闭。此外，连续映射 $|Z| \to T$ 是同胚，
因为 $|Z \times_X U| \to T'$ 是同胚，且
$|Z \times_Y U| \to |Z|$ 是商映射。
\end{proof}

\noindent
下述引理将帮助我们构造（局部）闭子空间。

\begin{lemma}
\label{spaces-properties-lemma-subspaces-presentation}
设 $S$ 为概形，$j : R \to U \times_S U$ 为 \'etale 等价关系，
$X = U/R$ 为相伴的代数空间（“空间”定理
\ref{spaces-theorem-presentation}）。存在典范双射
$$
R\text{-不变的局部闭子概形 }Z'\text{，所在概形为 }U
\leftrightarrow
\text{局部闭子空间 }Z\text{，所在空间为 }X
$$
此外，$Z \to X$ 为闭（相应地，开）浸入，当且仅当
$Z' \to U$ 为闭（相应地，开）浸入。
\end{lemma}

\begin{proof}
以 $\varphi : U \to X$ 表示典范映射。该双射把 $Z \to X$
送到 $Z' = Z \times_X U \to U$。由定义立即可知，若 $Z \to X$
为浸入、闭浸入或开浸入，则 $Z' \to U$ 分别也是浸入、闭浸入或开浸入。
$Z'$ 为 $R$-不变的也是清楚的（见“群胚”定义
\ref{groupoids-definition-invariant-open}）。

\medskip\noindent
反之，设 $Z' \to U$ 为 $R$-不变的浸入。令 $R'$ 为 $R$
在 $Z'$ 上的限制；见“群胚”定义
\ref{groupoids-definition-restrict-groupoid}。在本情形中
$R' = R \times_{s, U} Z' = Z' \times_{U, t} R$，故 $R'$
是 $Z'$ 上的 \'etale 等价关系。由“空间”定理
\ref{spaces-theorem-presentation}，$Z = Z'/R'$ 是代数空间。
按构造有 $U \times_X Z = Z'$，故 $U \times_X Z \to Z$ 是浸入。
注意，性质“浸入”在基变换下保持，并且在基底上对 fppf 拓扑是局部的
（见“空间”第 \ref{spaces-section-lists} 节）。此外，浸入是分离且局部
拟有限的（见“概形”引理 \ref{schemes-lemma-immersions-monomorphisms}
与“态射”引理 \ref{morphisms-lemma-immersion-locally-quasi-finite}）。
故由“态射进阶”引理
\ref{more-morphisms-lemma-separated-locally-quasi-finite-morphisms-fppf-descend}，
浸入对 fppf 覆盖满足下降。这意味着对 $Z \to X$、
$\mathcal{P}=$“浸入”以及满 \'etale 态射 $U \to X$，“空间”引理
\ref{spaces-lemma-morphism-sheaves-with-P-effective-descent-etale}
的全部假设均成立。因此 $Z \to X$ 可表示且为浸入；这正是子空间的定义
（见“空间”定义 \ref{spaces-definition-immersion}）。

\medskip\noindent
显然，这两个构造互逆，结论得证。
\end{proof}

\begin{lemma}
\label{spaces-properties-lemma-reduced-closed-subspace}
设 $S$ 为概形，$X$ 为 $S$ 上的代数空间，$T \subset |X|$
为闭子集。存在唯一的闭子空间 $Z \subset X$，满足：
(a) $|Z| = T$；(b) $Z$ 既约。
\end{lemma}

\begin{proof}
设 $U \to X$ 为满 \'etale 态射，其中 $U$ 为概形。令
$R = U \times_X U$，从而 $X = U/R$；见“空间”引理
\ref{spaces-lemma-space-presentation}。照例以 $s, t : R \to U$
表示两个投影态射。由引理 \ref{spaces-properties-lemma-points-presentation}，
$T$ 对应于闭子集 $T' \subset |U|$，且
$s^{-1}(T') = t^{-1}(T')$。令 $Z' \subset U$ 为 $T'$
上诱导的既约概形结构。在此情形中，纤维积
$Z' \times_{U, t} R$ 与 $Z' \times_{U, s} R$ 是 $R$ 的闭子概形
（“概形”引理 \ref{schemes-lemma-base-change-immersion}），
并且在 $Z'$ 上 \'etale（“态射”引理
\ref{morphisms-lemma-base-change-etale}），故为既约的
（因为既约性对 \'etale 拓扑是局部的；见注
\ref{spaces-properties-remark-list-properties-local-etale-topology}）。由于二者具有相同的
底拓扑空间（见上），可知
$Z' \times_{U, t} R = Z' \times_{U, s} R$。于是可应用引理
\ref{spaces-properties-lemma-subspaces-presentation}，得到闭子空间 $Z \subset X$，
其到 $U$ 的拉回为 $Z'$。按构造，$|Z| = T$ 且 $Z$ 既约，
这证明了存在性。唯一性的证明从略。
\end{proof}

\begin{lemma}
\label{spaces-properties-lemma-map-into-reduction}
设 $S$ 为概形，$X$、$Y$ 为 $S$ 上的代数空间，$Z \subset X$
为闭子空间，并假设 $Y$ 既约。态射 $f : Y \to X$ 通过 $Z$
分解，当且仅当 $f(|Y|) \subset |Z|$。
\end{lemma}

\begin{proof}
假设 $f(|Y|) \subset |Z|$。选取图表
$$
\xymatrix{
V \ar[d]_b \ar[r]_h & U \ar[d]^a \\
Y \ar[r]^f & X
}
$$
其中 $U$、$V$ 为概形，两个竖直箭头均满且 \'etale。由引理
\ref{spaces-properties-lemma-type-property}，概形 $V$ 既约。因此，由“概形”引理
\ref{schemes-lemma-map-into-reduction}，$h$ 通过 $a^{-1}(Z)$ 分解，
故 $a \circ h$ 通过 $Z$ 分解。由于 $Z \subset X$ 是子层，而
$V \to Y$ 是 $(\Sch/S)_{fppf}$ 上层的满射，可知
$X \to Y$ 通过 $Z$ 分解。
\end{proof}

\begin{definition}
\label{spaces-properties-definition-reduced-induced-space}
设 $S$ 为概形，$X$ 为 $S$ 上的代数空间，$Z \subset |X|$
为闭子集。$Z$ 上的一个{\it 代数空间结构}由 $X$ 的闭子空间
$Z'$ 给出，其中 $|Z'|$ 等于 $Z$。$Z$ 上的{\it 诱导既约代数空间结构}
是引理 \ref{spaces-properties-lemma-reduced-closed-subspace} 中构造的结构。
{\it $X$ 的既约化 $X_{red}$}是 $|X|$ 上的诱导既约代数空间结构。
\end{definition}











\section{概形轨迹}
\label{spaces-properties-section-schematic}

\noindent
每个代数空间都有一个最大的、为概形的开子空间；这多少是清楚的，
但我们仍在下面写出证明。当然，这个子空间可能为空，例如
$X = \mathbf{A}^1_{\mathbf{Q}}/\mathbf{Z}$（普适反例）。
另一方面，例如当 $X$ 拟分离时，这个最大开子概形事实上在 $X$ 中稠密！

\begin{lemma}
\label{spaces-properties-lemma-subscheme}
设 $S$ 为概形，$X$ 为 $S$ 上的代数空间。存在一个最大的开子空间
$X' \subset X$，它是概形。
\end{lemma}

\begin{proof}
设 $U \to X$ 为满 \'etale 态射，其中 $U$ 为概形，并令
$R = U \times_X U$。$X$ 的开子空间与 $U$ 的 $R$-不变开子概形
按 $1 - 1$ 对应，因而它们组成一个集合。设 $X_i$（$i \in I$）为 $X$
中所有为概形（即可表示）的开子空间。考虑开子空间 $X' \subset X$，
其底点集为 $|X|$ 的开集 $\bigcup |X_i|$。由引理
\ref{spaces-properties-lemma-characterize-surjective}，
$$
\coprod X_i \longrightarrow X'
$$
是 $(\Sch/S)_{fppf}$ 上层的满射。但由于每个 $X_i \to X'$
均可由开浸入表示，该映射事实上在 Zariski 拓扑中满。具体地，
若 $T \to X'$ 是从概形到 $X'$ 的态射，则
$X_i \times_{X'} T$ 是 $T$ 的开子概形。因此可应用“概形”引理
\ref{schemes-lemma-glue-functors}，得出 $X'$ 是概形。
\end{proof}

\noindent
在本节余下部分中，若代数空间 $X$ 的开子空间 $X'$ 所对应的开子集
$|X'| \subset |X|$ 稠密，则称该开子空间{\it 稠密}。

\begin{lemma}
\label{spaces-properties-lemma-quasi-separated-finite-etale-cover-dense-open-scheme}
设 $S$ 为概形，$X$ 为 $S$ 上的代数空间。若存在有限、\'etale
且满的态射 $U \to X$，其中 $U$ 为拟分离概形，则 $X$
存在一个为概形的稠密开子空间 $X'$。更精确地，$X$
中每个余维为 $0$ 的点 $x \in |X|$ 都包含于 $X'$。
\end{lemma}

\begin{proof}
设 $X' \subset X$ 为最大的、是概形的开子空间（引理
\ref{spaces-properties-lemma-subscheme}）。设 $x \in |X|$ 为 $X$ 上余维为 $0$
的一点。由引理 \ref{spaces-properties-lemma-codimension-0-points-dense}，
只需证明 $x \in X'$。取引理陈述中的 $U \to X$。记
$R = U \times_X U$，并照例以 $s, t : R \to U$ 表示投影。
注意，$s,t$ 满、有限且 \'etale。由引理
\ref{spaces-properties-lemma-finite-fibres-presentation}，$|U| \to |X|$
在 $x$ 上的纤维有限，记为 $\{\eta_1, \ldots, \eta_n\}$。
由引理 \ref{spaces-properties-lemma-codimension-0-points}，每个 $\eta_i$
都是 $U$ 某个不可约分支的一般点。由“性质”引理
\ref{properties-lemma-maximal-points-affine}，可找到包含
$\{\eta_1, \ldots, \eta_n\}$
的仿射开集 $W \subset U$（这里用到了 $U$ 拟分离）。由“群胚”引理
\ref{groupoids-lemma-find-invariant-affine}，可假设 $W$ 为
$R$-不变。由于 $W \subset U$ 是 $R$-不变仿射开集，$R$
在 $W$ 上的限制 $R_W$ 等于
$R_W = s^{-1}(W) = t^{-1}(W)$（见“群胚”定义
\ref{groupoids-definition-invariant-open} 及其后的讨论）。
特别地，映射 $R_W \to W$ 也有限且 \'etale，故 $R_W$ 仿射。
于是由“群胚”命题 \ref{groupoids-proposition-finite-flat-equivalence}，
$W/R_W$ 是概形。另一方面，由“空间”引理
\ref{spaces-lemma-finding-opens}，$W/R_W$ 是 $X$ 的开子空间，
并且按构造包含 $x$。
\end{proof}

\noindent
我们将在“良好空间”定理
\ref{decent-spaces-theorem-decent-open-dense-scheme}
中把下述命题推广到良好代数空间。

\begin{proposition}
\label{spaces-properties-proposition-locally-quasi-separated-open-dense-scheme}
设 $S$ 为概形，$X$ 为 $S$ 上的代数空间。若 $X$ 是 Zariski
局部拟分离的（例如 $X$ 拟分离），则 $X$ 存在为概形的稠密开子空间
$X'$。更精确地，$X$ 上每个余维为 $0$ 的点 $x \in |X|$
都包含于 $X'$。
\end{proposition}

\begin{proof}
由引理 \ref{spaces-properties-lemma-subscheme}，该问题在 $X$ 上是局部的。
故由引理 \ref{spaces-properties-lemma-quasi-separated-quasi-compact-pieces}，
可假设存在仿射概形 $U$ 及满、拟紧、\'etale 态射 $U \to X$。
此外，$U \to X$ 分离（引理 \ref{spaces-properties-lemma-separated-cover}）。
令 $R = U \times_X U$，并照例以 $s, t : R \to U$ 表示投影。
则 $s,t$ 满、拟紧、分离且 \'etale。因此 $s,t$ 也拟有限且具有有限纤维
（“态射”引理 \ref{morphisms-lemma-etale-locally-quasi-finite}、
\ref{morphisms-lemma-quasi-finite-locally-quasi-compact} 及
\ref{morphisms-lemma-quasi-finite}）。
由“态射”引理 \ref{morphisms-lemma-generically-finite}，对作为 $U$
某个不可约分支一般点的每个 $\eta \in U$，存在 $\eta$ 的开邻域
$V \subset U$，使 $s^{-1}(V) \to V$ 有限。由“下降”引理
\ref{descent-lemma-descending-property-finite}，有限性在目标上是
fpqc 局部的（特别是 \'etale 局部的）。因此可应用“群胚进阶”引理
\ref{more-groupoids-lemma-property-invariant}：使 $s$ 有限的最大开集
$W \subset U$ 是 $R$-不变的。由上可知，$W$ 包含 $U$
每个不可约分支的一般点。$R$ 在 $W$ 上的限制 $R_W$ 等于
$R_W = s^{-1}(W) = t^{-1}(W)$（见“群胚”定义
\ref{groupoids-definition-invariant-open} 及其后的讨论）。
按构造，$s_W, t_W : R_W \to W$ 有限且 \'etale。考虑开子空间
$X' = W/R_W \subset X$（见“空间”引理
\ref{spaces-lemma-finding-opens}）。按构造，包含映射 $X' \to X$
在余维 $0$ 的点上诱导双射。于是问题归结为引理
\ref{spaces-properties-lemma-quasi-separated-finite-etale-cover-dense-open-scheme}。
\end{proof}




\section{得到概形}
\label{spaces-properties-section-getting-a-scheme}

\noindent
上一节中我们用到了如下事实：若态射 $s, t : R \to U$ 有限且
\'etale，则仿射概形 $U$ 对等价关系 $R$ 的商 $U/R$ 是概形。
这是下述结果的特殊情形。

\begin{proposition}
\label{spaces-properties-proposition-finite-flat-equivalence-global}
设 $S$ 为概形，$(U, R, s, t, c)$ 为 $S$ 上的群胚概形。假设：
\begin{enumerate}
\item $s, t : R \to U$ 有限局部自由；
\item $j = (t, s)$ 是等价关系；
\item $U$ 的每个非空闭子集 $Z$ 都包含一点 $u$，其 $R$-等价类
$t(s^{-1}(\{u\}))$ 包含于 $U$ 的某个仿射开集\footnote{令
$E \subset U$ 为满足如下条件的点 $u$ 之集：
$t(s^{-1}(\{u\}))$ 包含于 $U$ 的某个仿射开集。若 $E = U$，
或 $U$ 的每个有限型点都属于 $E$，或每个 $u \in U$
都特化到 $E$ 的一点，则条件 (3) 成立。}。
\end{enumerate}
则存在 $S$ 上概形之间的有限局部自由态射 $U \to M$，使得
$R = U \times_M U$，并且 $M$ 表示 fppf 拓扑中的商层 $U/R$。
\end{proposition}

\begin{proof}
由假设 (3) 与“群胚”引理 \ref{groupoids-lemma-find-invariant-affine}，
可找到开覆盖 $U = \bigcup U_i$，使每个 $U_i$ 都是 $U$
的 $R$-不变仿射开集。令 $R_i = R|_{U_i}$。考虑 fppf 层
$F = U/R$ 与 $F_i = U_i/R_i$。由“空间”引理
\ref{spaces-lemma-finding-opens}，态射 $F_i \to F$ 可表示且为开浸入。
由“群胚”命题 \ref{groupoids-proposition-finite-flat-equivalence}，
层 $F_i$ 可由仿射概形表示。若 $T$ 为概形且 $T \to F$ 为态射，
则 $V_i = F_i \times_F T$ 在 $T$ 中开；我们断言
$T = \bigcup V_i$。事实上，在 $T$ 上 fppf 局部地，可将
$T \to F$ 提升为态射 $f : T \to U$，此时
$f^{-1}(U_i) \subset V_i$。因此 $F$ 可由概形表示；见“概形”引理
\ref{schemes-lemma-glue-functors}。
\end{proof}

\noindent
例如，若 $U$ 同构于某个仿射概形的局部闭子概形，或同构于某个分次环
$A$ 的 $\text{Proj}(A)$ 的局部闭子概形，则由“性质”引理
\ref{properties-lemma-ample-finite-set-in-affine}，第三个假设成立。
特别地，这可应用于有限群及有限群概形在拟仿射或拟投射概形上的自由作用。
例如，有限群 $G$ 在拟投射簇 $X$ 上自由作用所得的商 $X/G$
是概形。下面给出详细陈述。

\begin{lemma}
\label{spaces-properties-lemma-quotient-scheme}
设 $S$ 为概形，$G \to S$ 为群概形，$X \to S$ 为概形的态射，
$a : G \times_S X \to X$ 为一个作用。假设：
\begin{enumerate}
\item $G \to S$ 有限局部自由；
\item 作用 $a$ 自由；
\item $X \to S$ 仿射、拟仿射、投射或拟投射，或者 $X$
同构于某个仿射概形的开子概形，或者 $X$ 同构于某个分次环 $A$
之 $\text{Proj}(A)$ 的开子概形，或者 $G \to S$ 是 radicial 的。
\end{enumerate}
则 fppf 商层 $X/G$ 是概形，且 $X \to X/G$ 是 fppf $G$-挠子。
\end{lemma}

\begin{proof}
先证明 $X/G$ 是概形。因为作用自由，态射
$j = (a, \text{pr}) : G \times_S X \to X \times_S X$
是单态射，因而是等价关系；见“群胚”引理
\ref{groupoids-lemma-free-action}。由于假设 $G \to S$ 有限局部自由，
映射 $s, t : G \times_S X \to X$ 有限局部自由。现在只需证明命题
\ref{spaces-properties-proposition-finite-flat-equivalence-global} 的最后一个假设成立。
由于 $G$ 的作用在 $S$ 上，只需证明 $X \to S$ 的任一纤维中的有限点集
都包含于 $X$ 的某个仿射开集。若 $X$ 同构于仿射概形的开子概形，
或同构于某个分次环 $A$ 之 $\text{Proj}(A)$ 的开子概形，
则这由“性质”引理 \ref{properties-lemma-ample-finite-set-in-affine}
得出。若 $X \to S$ 仿射、拟仿射、投射或拟投射，可将 $S$
换为仿射开集，从而回到刚处理的情形。若 $G \to S$ 是 radicial 的，
则 $G$ 作用在 $X$ 各点上的轨道均为单点集，条件显然成立。
若干细节从略。

\medskip\noindent
为证明 $X \to X/G$ 是 fppf $G$-挠子（“群胚”定义
\ref{groupoids-definition-principal-homogeneous-space}），须证明
$G \times_S X \to X \times_{X/G} X$ 是同构，且 $X \to X/G$
在 fppf 局部有截面。第二点由 $X \to X/G$ 作为 fppf 层的映射为满
（按构造）立即得出。第一点来自命题
\ref{spaces-properties-proposition-finite-flat-equivalence-global} 结论中的同构
$R = U \times_M U$（注意，在本情形中 $R = G \times_S X$）。
\end{proof}

\begin{lemma}
\label{spaces-properties-lemma-quotient-separated}
记号与假设同命题 \ref{spaces-properties-proposition-finite-flat-equivalence-global}。则：
\begin{enumerate}
\item 若 $U$ 在 $S$ 上拟分离，则 $U/R$ 在 $S$ 上拟分离；
\item 若 $U$ 拟分离，则 $U/R$ 拟分离；
\item 若 $U$ 在 $S$ 上分离，则 $U/R$ 在 $S$ 上分离；
\item 若 $U$ 分离，则 $U/R$ 分离；
\item 在此添加更多内容。
\end{enumerate}
在引理 \ref{spaces-properties-lemma-quotient-scheme} 的情形中也有类似结果。
\end{lemma}

\begin{proof}
由于 $M$ 表示商层，有概形的 Cartesian 图表
$$
\xymatrix{
R \ar[r]_-j \ar[d] & U \times_S U \ar[d] \\
M \ar[r] & M \times_S M
}
$$
由于 $U \times_S U \to M \times_S M$ 满且有限局部自由，为证明
$M \to M \times_S M$ 拟紧（相应地，为闭浸入），只需证明
$j : R \to U \times_S U$ 拟紧（相应地，为闭浸入）；见“下降”引理
\ref{descent-lemma-descending-property-quasi-compact} 与
\ref{descent-lemma-descending-property-closed-immersion}。
由于 $j : R \to U \times_S U$ 是 $U$ 上的态射，而 $R$ 在 $U$
上有限，只要投影 $U \times_S U \to U$ 拟分离，$j$ 便拟紧
（“概形”引理 \ref{schemes-lemma-quasi-compact-permanence}）。
由于 $j$ 为单态射且局部有限型，只要它固有，$j$ 便是闭浸入
（“\'Etale 态射”引理 \ref{etale-lemma-characterize-closed-immersion}）；
而只要投影 $U \times_S U \to U$ 分离，这一条件便成立
（“态射”引理 \ref{morphisms-lemma-image-proper-scheme-closed}）。
这证明了 (1) 与 (3)。为证明 (2) 与 (4)，将 $S$
换为 $\Spec(\mathbf{Z})$；见定义 \ref{spaces-properties-definition-separated}。
由于引理 \ref{spaces-properties-lemma-quotient-scheme} 是通过应用命题
\ref{spaces-properties-proposition-finite-flat-equivalence-global} 证明的，
最后的断言也清楚。
\end{proof}




\section{拟分离空间上的点}
\label{spaces-properties-section-points-quasi-separated}

\noindent
在 Stacks 项目所采用的一般性下，代数空间上的点可能表现得很差。
然而，对拟分离空间，其行为大多与概形上的点相似。本节证明几个相关结果；
“良好空间”一章还包含更多结果，例如见“良好空间”第
\ref{decent-spaces-section-points} 节。

\begin{lemma}
\label{spaces-properties-lemma-quasi-separated-sober}
设 $S$ 为概形，$X$ 为 $S$ 上 Zariski 局部拟分离的代数空间。
则拓扑空间 $|X|$ 是清醒的（见“拓扑”定义
\ref{topology-definition-generic-point}）。
\end{lemma}

\begin{proof}
结合“拓扑”引理 \ref{topology-lemma-sober-local} 与引理
\ref{spaces-properties-lemma-quasi-separated-quasi-compact-pieces}，可假设存在仿射概形
$U$ 及满、拟紧、\'etale 态射 $U \to X$。令
$R = U \times_X U$，投影映射为 $s, t : R \to U$。应用引理
\ref{spaces-properties-lemma-finite-fibres-presentation}，可知 $s,t$ 的纤维有限。
于是“拓扑”引理 \ref{topology-lemma-quotient-kolmogorov}
的全部假设均满足，从而 $|X|$ 是 Kolmogorov 空间\footnote{
事实上，这里还用到了“概形”引理
\ref{schemes-lemma-scheme-sober}（概形的清醒性）、
“态射”引理 \ref{morphisms-lemma-etale-flat} 与
\ref{morphisms-lemma-generalizations-lift-flat}
（沿 \'etale 态射可提升一般化）、引理
\ref{spaces-properties-lemma-points-presentation}（用表示描述代数空间上的点），以及引理
\ref{spaces-properties-lemma-topology-points}（商映射的开性）。}。

\medskip\noindent
还需证明每个不可约闭子集 $T \subset |X|$ 都有一般点。由引理
\ref{spaces-properties-lemma-reduced-closed-subspace}，存在闭子空间 $Z \subset X$，
使 $|Z| = |T|$。注意，$U \times_X Z \to Z$ 是从仿射概形到 $Z$
的拟紧、满、\'etale 态射，故由引理
\ref{spaces-properties-lemma-quasi-separated-quasi-compact-pieces}，$Z$ 是 Zariski
局部拟分离的。由命题
\ref{spaces-properties-proposition-locally-quasi-separated-open-dense-scheme}，
存在为概形的开稠密子空间 $Z' \subset Z$。这意味着
$|Z'| \subset T$ 开且稠密。因此拓扑空间 $|Z'|$ 不可约，
即 $Z'$ 是不可约概形。由“概形”引理
\ref{schemes-lemma-scheme-sober}，$|Z'|$ 是单个点
$\eta \in |Z'| \subset T$ 的闭包，因而也有
$T = \overline{\{\eta\}}$，结论得证。
\end{proof}

\begin{lemma}
\label{spaces-properties-lemma-quasi-compact-quasi-separated-spectral}
设 $S$ 为概形，$X$ 为 $S$ 上拟紧且拟分离的代数空间。
拓扑空间 $|X|$ 是谱空间。
\end{lemma}

\begin{proof}
由“拓扑”定义 \ref{topology-definition-spectral-space}，须验证
$|X|$ 清醒、拟紧、具有由拟紧开集组成的基，并且任意两个拟紧开集
之交拟紧。由引理 \ref{spaces-properties-lemma-quasi-separated-sober}，$|X|$ 清醒；
由引理 \ref{spaces-properties-lemma-quasi-compact-space}，$|X|$ 拟紧。
由引理 \ref{spaces-properties-lemma-quasi-compact-affine-cover}，存在仿射概形 $U$
及满 \'etale 态射 $f : U \to X$。因为
$|f| : |U| \to |X|$ 开且连续，而 $|U|$ 具有由拟紧开集组成的基，
故 $|X|$ 也有这样的基。最后，设 $A, B \subset |X|$ 为拟紧开集。
则存在开子空间 $X', X'' \subset X$，使
$A = |X'|$ 且 $B = |X''|$（引理 \ref{spaces-properties-lemma-open-subspaces}）；
并且可选取仿射概形 $V$、$W$ 及满 \'etale 态射
$V \to X'$、$W \to X''$（引理
\ref{spaces-properties-lemma-quasi-compact-affine-cover}）。于是 $A \cap B$
是 $|V \times_X W| \to |X|$ 的像（引理
\ref{spaces-properties-lemma-points-cartesian}）。由于 $X$ 拟分离，
$V \times_X W$ 拟紧（引理
\ref{spaces-properties-lemma-characterize-quasi-separated}），故 $A \cap B$ 拟紧，
证明完成。
\end{proof}

\noindent
下述引理可用于证明一个代数空间同构于某个域的谱。

\begin{lemma}
\label{spaces-properties-lemma-point-like-spaces}
设 $S$ 为概形，$k$ 为域，$X$ 为 $S$ 上的代数空间，并假设存在
满 \'etale 态射 $\Spec(k) \to X$。若 $X$ 拟分离，则
$X \cong \Spec(k')$，其中 $k/k'$ 为有限可分扩张。
\end{lemma}

\begin{proof}
令 $R = \Spec(k) \times_X \Spec(k)$，于是有纤维积图表
$$
\xymatrix{
R \ar[r]_-s \ar[d]_-t & \Spec(k) \ar[d] \\
\Spec(k) \ar[r] & X
}
$$
由“空间”引理 \ref{spaces-lemma-space-presentation}，
$X = \Spec(k)/R$ 是商层。由于 $\Spec(k) \to X$ 为 \'etale，
态射 $s$ 与 $t$ 均为 \'etale。因此
$R = \coprod_{i \in I} \Spec(k_i)$ 是域谱的不交并，并且 $s$ 与 $t$
都诱导有限可分域扩张 $s, t : k \subset k_i$；见“态射”引理
\ref{morphisms-lemma-etale-over-field}。因为
$$
R = \Spec(k) \times_X \Spec(k)
= (\Spec(k) \times_S \Spec(k)) \times_{X \times_S X, \Delta} X
$$
并且按假设 $\Delta$ 拟紧，可知
$R \to \Spec(k) \times_S \Spec(k)$ 拟紧。由于
$\Spec(k) \times_S \Spec(k)$ 仿射，$R$ 拟紧，故 $I$ 有限。
这蕴含 $s$ 与 $t$ 是有限局部自由态射。于是由“群胚”命题
\ref{groupoids-proposition-finite-flat-equivalence}，
$\Spec(k)/R$ 可由 $\Spec(k')$ 表示，其中 $k' \subset k$
有限局部自由，且
$$
k' = \{x \in k \mid s_i(x) = t_i(x)\text{ 对所有 }i \in I\}
$$
容易看出 $k'$ 是域。
\end{proof}

\begin{remark}
\label{spaces-properties-remark-cannot-decide-yet}
引理 \ref{spaces-properties-lemma-point-like-spaces} 对良好代数空间也成立；见“良好空间”
引理 \ref{decent-spaces-lemma-decent-point-like-spaces}。事实上，
只有一个点的良好代数空间是概形；见“良好空间”引理
\ref{decent-spaces-lemma-when-field}。当 $X$ 局部分离时，这也成立，
因为局部分离的代数空间是良好的；见“良好空间”引理
\ref{decent-spaces-lemma-locally-separated-decent}。
\end{remark}







\section{代数空间的 \'Etale 态射}
\label{spaces-properties-section-etale-morphisms}

\noindent
本节本应属于“代数空间的态射”一章，但为了定义代数空间的小 \'etale
站点，需要一个代数空间在另一个之上 \'etale 的概念。因此必须先对从概形
到代数空间的 \'etale 态射以及代数空间之间的 \'etale 态射做一些准备。
关于代数空间的 \'etale 态射的更多内容，见“空间的态射”第
\ref{spaces-morphisms-section-etale} 节。

\begin{lemma}
\label{spaces-properties-lemma-etale-over-space}
设 $S$ 为概形，$X$ 为 $S$ 上的代数空间，$U$、$U'$ 为 $S$ 上的概形。
\begin{enumerate}
\item 若 $U \to U'$ 是概形的 \'etale 态射，而 $U' \to X$
是从 $U'$ 到 $X$ 的 \'etale 态射，则复合 $U \to X$
是从 $U$ 到 $X$ 的 \'etale 态射。
\item 若 $\varphi : U \to X$ 与 $\varphi' : U' \to X$
是到 $X$ 的 \'etale 态射，而概形态射 $\chi : U \to U'$
满足 $\varphi = \varphi' \circ \chi$，则 $\chi$ 是概形的
\'etale 态射。
\item 若 $\chi : U \to U'$ 是概形的满 \'etale 态射，而
$\varphi' : U' \to X$ 是满足 $\varphi = \varphi' \circ \chi$
为 \'etale 的态射，则 $\varphi'$ 为 \'etale。
\end{enumerate}
\end{lemma}

\begin{proof}
回忆：从概形到代数空间的任意态射均可表示；据此，我们从“空间”定义
\ref{spaces-definition-relative-representable-property}
得到从概形到代数空间的 \'etale 态射之定义。

\medskip\noindent
引理的 (1) 由此、\'etale 态射在复合下保持这一事实
（“态射”引理 \ref{morphisms-lemma-composition-etale}）以及“空间”引理
\ref{spaces-lemma-composition-representable-transformations-property} 与
\ref{spaces-lemma-morphism-schemes-gives-representable-transformation-property}
（它们是形式的）得出。

\medskip\noindent
为证明 (2)，选取 $S$ 上的概形 $W$ 及满 \'etale 态射 $W \to X$。
考虑 $\chi$ 的基变换
$\chi_W : W \times_X U \to W \times_X U'$。由于
$W \times_X U$ 与 $W \times_X U'$ 在 $W$ 上 \'etale，
由“态射”引理 \ref{morphisms-lemma-etale-permanence}，
$\chi_W$ 为 \'etale。另一方面，在交换图表
$$
\xymatrix{
W \times_X U \ar[r] \ar[d] & W \times_X U' \ar[d] \\
U \ar[r] & U'
}
$$
中，两个竖直箭头都 \'etale 且满。故由“下降”引理
\ref{descent-lemma-syntomic-smooth-etale-permanence}，
$U \to U'$ 为 \'etale。

\medskip\noindent
为证明 (3)，选取 $S$ 上的概形 $W$ 及态射 $W \to X$。
如上考虑图表
$$
\xymatrix{
W \times_X U \ar[r] \ar[d] & W \times_X U' \ar[d] \ar[r] & W \ar[d] \\
U \ar[r] & U' \ar[r] & X
}
$$
现在知道 $W \times_X U \to W \times_X U'$ 满且 \'etale
（因为它是 $U \to U'$ 的基变换），而
$W \times_X U \to W$ 为 \'etale。因此，由“下降”引理
\ref{descent-lemma-syntomic-smooth-etale-permanence}，
$W \times_X U' \to W$ 为 \'etale。按定义，这意味着
$\varphi'$ 为 \'etale。
\end{proof}

\begin{definition}
\label{spaces-properties-definition-etale}
设 $S$ 为概形。称 $S$ 上代数空间之间的态射 $f : X \to Y$
为 {\it \'etale} 的，当且仅当对每个以概形 $U$ 为源的 \'etale
态射 $\varphi : U \to X$，复合 $f \circ \varphi$ 也为 \'etale。
\end{definition}

\noindent
若 $X$ 与 $Y$ 是概形，则这与概形的 \'etale 态射之通常概念一致。
事实上，只要 $X \to Y$ 是代数空间的可表示态射，这一概念便与通过
“空间”定义 \ref{spaces-definition-relative-representable-property}
给出的概念一致。这由下述引理 \ref{spaces-properties-lemma-etale-local} 与“空间”引理
\ref{spaces-lemma-representable-morphisms-spaces-property} 结合得出。

\begin{lemma}
\label{spaces-properties-lemma-etale-local}
设 $S$ 为概形，$f : X \to Y$ 为 $S$ 上代数空间的态射。
下列条件等价：
\begin{enumerate}
\item $f$ 为 \'etale；
\item 存在满 \'etale 态射 $\varphi : U \to X$，其中 $U$ 为概形，
使复合 $f \circ \varphi$ 为 \'etale（作为代数空间的态射）；
\item 存在满 \'etale 态射 $\psi : V \to Y$，其中 $V$ 为概形，
使基变换 $V \times_Y X \to V$ 为 \'etale（作为代数空间的态射）；
\item 存在交换图表
$$
\xymatrix{
U \ar[d] \ar[r] & V \ar[d] \\
X \ar[r] & Y
}
$$
其中 $U$、$V$ 为概形，两个竖直箭头为 \'etale，左竖直箭头满，
并且上方的水平箭头为 \'etale。
\end{enumerate}
\end{lemma}

\begin{proof}
先证明 (4) 蕴含 (1)。假设给定 (4) 中的图表。设 $W \to X$
为以概形 $W$ 为源的 \'etale 态射。则
$W \times_X U \to U$ 为 \'etale。由于
$W \times_X U \to U$ 与 $U \to V$ 均为概形的 \'etale 态射，
其复合 $W \times_X U \to V$ 为 \'etale。因此，由引理
\ref{spaces-properties-lemma-etale-over-space} (1)，$W \times_X U \to Y$ 为
\'etale。投影 $W \times_X U \to W$ 也满且 \'etale，故由引理
\ref{spaces-properties-lemma-etale-over-space} (3)，$W \to Y$ 为 \'etale。

\medskip\noindent
再证明 (1) 蕴含 (4)。假设 (1)，并选取交换图表
$$
\xymatrix{
U \ar[d] \ar[r] & V \ar[d] \\
X \ar[r] & Y
}
$$
其中 $U \to X$ 与 $V \to Y$ 满且 \'etale；见“空间”引理
\ref{spaces-lemma-lift-morphism-presentations}。按假设，
态射 $U \to Y$ 为 \'etale，故由引理
\ref{spaces-properties-lemma-etale-over-space} (2)，$U \to V$ 为 \'etale。

\medskip\noindent
(2) 与 (3) 也等价于 (1) 的证明从略。
\end{proof}

\begin{lemma}
\label{spaces-properties-lemma-composition-etale}
代数空间的两个 \'etale 态射之复合为 \'etale。
\end{lemma}

\begin{proof}
这由定义立即得出。
\end{proof}

\begin{lemma}
\label{spaces-properties-lemma-base-change-etale}
代数空间的 \'etale 态射沿代数空间的任意态射作基变换后仍为 \'etale。
\end{lemma}

\begin{proof}
设 $X \to Y$ 为 $S$ 上代数空间的 \'etale 态射，$Z \to Y$
为代数空间的态射。选取概形 $U$ 及满 \'etale 态射 $U \to X$；
再选取概形 $W$ 及满 \'etale 态射 $W \to Z$。于是 $U \to Y$
为 \'etale，故在图表
$$
\xymatrix{
W \times_Y U \ar[d] \ar[r] & W \ar[d] \\
Z \times_Y X \ar[r] & Z
}
$$
中，上方的水平箭头为 \'etale。此外，左侧竖直箭头满且 \'etale
（验证从略）。故由引理 \ref{spaces-properties-lemma-etale-local}，下方水平箭头为
\'etale。
\end{proof}

\begin{lemma}
\label{spaces-properties-lemma-etale-permanence}
设 $S$ 为概形，$X,Y,Z$ 为代数空间。设
$g : X \to Z$、$h : Y \to Z$ 为 \'etale 态射，而
$f : X \to Y$ 为满足 $h \circ f = g$ 的态射。则 $f$ 为 \'etale。
\end{lemma}

\begin{proof}
选取交换图表
$$
\xymatrix{
U \ar[d] \ar[r]_\chi & V \ar[d] \\
X \ar[r] & Y
}
$$
其中 $U \to X$ 与 $V \to Y$ 满且 \'etale；见“空间”引理
\ref{spaces-lemma-lift-morphism-presentations}。按假设，态射
$\varphi : U \to X \to Z$ 与 $\psi : V \to Y \to Z$ 为
\'etale。此外，由对 $f,g,h$ 的假设，
$\psi \circ \chi = \varphi$。故由引理
\ref{spaces-properties-lemma-etale-over-space} (2)，$U \to V$ 为 \'etale。
\end{proof}

\begin{lemma}
\label{spaces-properties-lemma-etale-open}
设 $S$ 为概形。若 $X \to Y$ 为 $S$ 上代数空间的 \'etale 态射，
则相伴的拓扑空间映射 $|X| \to |Y|$ 是开映射。
\end{lemma}

\begin{proof}
这由引理 \ref{spaces-properties-lemma-etale-local} 中的图表及引理
\ref{spaces-properties-lemma-topology-points} 即得。
\end{proof}

\noindent
最后给出一个有趣的引理。带有到概形之 \'etale 态射的代数空间不一定是概形；
见“空间”例 \ref{spaces-example-non-representable-descent}。
但若目标是某个域的谱，结论成立。

\begin{lemma}
\label{spaces-properties-lemma-etale-over-field-scheme}
设 $S$ 为概形，$X \to \Spec(k)$ 为 $S$ 上的 \'etale 态射，
其中 $k$ 为域。则 $X$ 是概形。
\end{lemma}

\begin{proof}
设 $U$ 为仿射概形，$U \to X$ 为 \'etale 态射。由定义
\ref{spaces-properties-definition-etale}，$U \to \Spec(k)$ 为 \'etale 态射。
故 $U = \coprod_{i = 1, \ldots, n} \Spec(k_i)$ 是有限个域谱的不交并，
其中 $k_i$ 为 $k$ 的有限可分扩张；见“态射”引理
\ref{morphisms-lemma-etale-over-field}。态射
$R = U \times_X U \to U \times_{\Spec(k)} U$ 为单态射，而
$U \times_{\Spec(k)} U$ 也为有限个 $k$ 的有限可分扩张之谱的不交并。
故由“概形”引理 \ref{schemes-lemma-mono-towards-spec-field}，$R$
同样是有限个 $k$ 的有限可分扩张之谱的不交并。这里 $U$ 与 $R$
均仿射，且两个投影 $R \to U$ 都有限局部自由。因此，由“群胚”命题
\ref{groupoids-proposition-finite-flat-equivalence}，$U/R$ 是概形。
由“空间”引理 \ref{spaces-lemma-finding-opens}，它也是 $X$
的开子空间。再由引理 \ref{spaces-properties-lemma-subscheme}，$X$ 是概形。
\end{proof}













\section{空间与 fpqc 覆盖}
\label{spaces-properties-section-fpqc}

\noindent
设 $S$ 为概形。$S$ 上的代数空间定义为 fppf 拓扑中满足若干附加性质的层。
因此，它是否满足 fpqc 拓扑的层性质并非立即清楚（见“拓扑”定义
\ref{topologies-definition-sheaf-property-fpqc}）。本节给出 Gabber
的论证，证明答案是肯定的。不过，当我们说代数空间 $X$ 满足 fpqc
拓扑的层性质时，实际上只考虑这样的 fpqc 覆盖
$\{f_i : T_i \to T\}_{i \in I}$：$T,T_i$ 是大站点
$(\Sch/S)_{fppf}$ 的对象（依照我们的约定，见第
\ref{spaces-properties-section-conventions} 节）。

\begin{proposition}[Gabber]
\label{spaces-properties-proposition-sheaf-fpqc}
设 $S$ 为概形，$X$ 为 $S$ 上的代数空间。则 $X$ 满足 fpqc
拓扑的层性质。
\end{proposition}

\begin{proof}
由于 $X$ 是 Zariski 拓扑的层，只需证明如下事实：给定仿射概形的
满平坦态射 $f : T' \to T$，$X(T)$ 是两个映射
$X(T') \to X(T' \times_T T')$ 的等化子。见“拓扑”引理
\ref{topologies-lemma-sheaf-property-fpqc}（这里略去了一个小论证，
因为所引引理是对全体概形范畴上定义的函子表述的）。

\medskip\noindent
设 $a, b : T \to X$ 为满足 $a \circ f = b \circ f$ 的两个态射。
须证明 $a = b$。考虑纤维积
$$
E = X \times_{\Delta_{X/S}, X \times_S X, (a, b)} T.
$$
由“空间”引理 \ref{spaces-lemma-properties-diagonal}，态射
$\Delta_{X/S}$ 是可表示单态射，故 $E \to T$ 是概形的单态射。
假设 $a \circ f = b \circ f$ 蕴含 $T' \to T$
（唯一地）通过 $E$ 分解。考虑交换图表
$$
\xymatrix{
T' \times_T E \ar[r] \ar[d] & E \ar[d] \\
T' \ar[r] \ar@/^5ex/[u] \ar[ru] & T
}
$$
投影 $T' \times_T E \to T'$ 是带截面的单态射，故为同构。
于是由“下降”引理 \ref{descent-lemma-descending-property-isomorphism}，
$E \to T$ 为同构。这意味着 $a = b$，正合所需。

\medskip\noindent
接着，设 $c : T' \to X$ 为一个态射，使得两个复合
$T' \times_T T' \to T' \to X$ 相同。须找到态射 $a : T \to X$，
使它与 $T' \to T$ 的复合为 $c$。选取仿射概形 $U$
及 \'etale 态射 $U \to X$，使 $|U| \to |X|$ 的像包含
$|c| : |T'| \to |X|$ 的像。这可由引理
\ref{spaces-properties-lemma-topology-points}、\ref{spaces-properties-lemma-cover-by-union-affines}、
有限个仿射概形的不交并仍仿射这一事实以及 $|T'|$ 拟紧这一事实做到
（略去一个小论证）。由于 $U \to X$ 分离（引理
\ref{spaces-properties-lemma-separated-cover}），
$$
V = U \times_{X, c} T' \longrightarrow T'
$$
是概形的满、\'etale、分离态射（要看出它满，使用引理
\ref{spaces-properties-lemma-points-cartesian} 及对 $U \to X$ 的选择）。
等式 $c \circ \text{pr}_0 = c \circ \text{pr}_1$ 意味着
$V/T'/T$ 上得到一个下降数据（“下降”定义
\ref{descent-definition-descent-datum}），因为
\begin{align*}
V \times_{T'} (T' \times_T T')
& =
U \times_{X, c \circ \text{pr}_0} (T' \times_T T') \\
& =
(T' \times_T T') \times_{c \circ \text{pr}_1, X} U \\
& =
(T' \times_T T') \times_{T'} V
\end{align*}
由“态射进阶”引理
\ref{more-morphisms-lemma-etale-separated-ind-quasi-affine}，
态射 $V \to T'$ 是 ind-拟仿射的（因为 \'etale 态射局部拟有限；
见“态射”引理 \ref{morphisms-lemma-etale-locally-quasi-finite}）。
由“群胚进阶”引理 \ref{more-groupoids-lemma-ind-quasi-affine}，
该下降数据有效。设 $W \to T$ 为态射，并有同构
$\alpha : T' \times_T W \to V$，它与 $V$ 上给定的下降数据以及
$T' \times_T W$ 上的典范下降数据相容。则 $W \to T$ 满且
\'etale（“下降”引理 \ref{descent-lemma-descending-property-surjective}
与 \ref{descent-lemma-descending-property-etale}）。考虑复合
$$
b' : T' \times_T W \longrightarrow V = U \times_{X, c} T' \longrightarrow U
$$
两个复合
$b' \circ (\text{pr}_0, 1), 
b' \circ (\text{pr}_1, 1) :
(T' \times_T T') \times_T W \to T' \times_T W \to U$
因对 $\alpha$ 的选择及 $c$ 的相应性质而相同（计算从略）。
故由“下降”引理
\ref{descent-lemma-fpqc-universal-effective-epimorphisms}，
$b'$ 下降为态射 $b : W \to U$。图表
$$
\xymatrix{
T' \times_T W \ar[r] \ar[d] & W \ar[r]_b & U \ar[d] \\
T' \ar[rr]^c &  & X
}
$$
交换。这意味着已在 $T$ 上 \'etale 局部地证明了 $a$ 的存在性，
即有 $a' : W \to X$。但由于第一段已证明唯一性，这一 \'etale
局部解满足粘合条件；也就是说，作为 $X(W \times_T W)$ 的元素，
有 $\text{pr}_0^*a' = \text{pr}_1^*a'$。由于 $X$ 是 \'etale 层，
存在唯一的 $a \in X(T)$，其到 $W$ 的限制为 $a'$。
\end{proof}









\section{代数空间的 \'etale 站点}
\label{spaces-properties-section-etale-site}

\noindent
本节定义代数空间的小 \'etale 站点。这是概形的小 \'etale 站点
$S_\etale$ 的类比。引理 \ref{spaces-properties-lemma-etale-over-space} 蕴含：
在下述定义中，$X$ 的 \'etale 站点之对象之间的任意态射都是
\'etale 的；在 $X_\etale$ 的某个对象上 \'etale 的任意概形也仍是
$X_\etale$ 的对象。

\begin{definition}
\label{spaces-properties-definition-etale-site}
设 $S$ 为概形，$\Sch_{fppf}$ 为包含 $S$ 的大 fppf 站点，
$\Sch_\etale$ 为相应的大 \'etale 站点（即有相同的底范畴）。
设 $X$ 为 $S$ 上的代数空间。$X$ 的{\it 小 \'etale 站点
$X_\etale$}定义如下：
\begin{enumerate}
\item $X_\etale$ 的对象是态射 $\varphi : U \to X$，其中
$U \in \Ob((\Sch/S)_\etale)$ 为概形，$\varphi$ 为 \'etale 态射；
\item 态射 $(\varphi : U \to X) \to (\varphi' : U' \to X)$
由满足 $\varphi = \varphi' \circ \chi$ 的概形态射
$\chi : U \to U'$ 给出；
\item $X_\etale$ 的态射族
$\{(U_i \to X) \to (U \to X)\}_{i \in I}$ 是覆盖，当且仅当
$\{U_i \to U\}_{i \in I}$ 是 $(\Sch/S)_\etale$ 的覆盖。
\end{enumerate}
\end{definition}

\noindent
这一选择的结果是：代数空间的 \'etale 站点一般没有终对象！另一方面，
若 $X$ 恰为概形，则上述定义与“拓扑”定义
\ref{topologies-definition-big-small-etale} 一致。

\medskip\noindent
上面是默认站点，但还会使用几个变体。具体地，可以考虑在 $X$ 上
\'etale 的所有{\it 代数空间} $U$，得到下述站点
$X_{spaces, \etale}$；也可以考虑在 $X$ 上 \'etale 的所有
{\it 仿射概形} $U$，得到下述站点 $X_{affine, \etale}$。
讨论小 \'etale 站点的函子性时会使用前一种概念；见引理
\ref{spaces-properties-lemma-functoriality-etale-site}。

\begin{definition}
\label{spaces-properties-definition-spaces-etale-site}
设 $S$ 为概形，$\Sch_{fppf}$ 为包含 $S$ 的大 fppf 站点，
$\Sch_\etale$ 为相应的大 \'etale 站点（即有相同的底范畴）。
设 $X$ 为 $S$ 上的代数空间。$X$ 的站点
{\it $X_{spaces, \etale}$} 定义如下：
\begin{enumerate}
\item $X_{spaces, \etale}$ 的对象是态射 $\varphi : U \to X$，
其中 $U$ 为 $S$ 上的代数空间，$\varphi$ 为 $S$
上代数空间的 \'etale 态射；
\item $X_{spaces, \etale}$ 的态射
$(\varphi : U \to X) \to (\varphi' : U' \to X)$
由满足 $\varphi = \varphi' \circ \chi$ 的代数空间态射
$\chi : U \to U'$ 给出；
\item $X_{spaces, \etale}$ 的态射族
$\{\varphi_i : (U_i \to X) \to (U \to X)\}_{i \in I}$ 是覆盖，
当且仅当 $|U| = \bigcup \varphi_i(|U_i|)$。
\end{enumerate}
照例，依“集合”引理 \ref{sets-lemma-coverings-site}，选取一组
此类覆盖，使其至少包含 $X_\etale$ 中的覆盖，从而把
$X_{spaces, \etale}$ 变为站点。
\end{definition}

\noindent
由于 $X$ 的恒等态射为 \'etale，显然 $X_{spaces, \etale}$
确有终对象。下面立即证明相应拓扑斯等于 $X$ 的小 \'etale 拓扑斯。

\begin{lemma}
\label{spaces-properties-lemma-compare-etale-sites}
函子
$$
X_\etale \longrightarrow X_{spaces, \etale}, \quad
U/X \longmapsto U/X
$$
是特殊余连续函子（“站点”定义
\ref{sites-definition-special-cocontinuous-functor}），因而诱导拓扑斯等价
$\Sh(X_\etale) \to \Sh(X_{spaces, \etale})$。
\end{lemma}

\begin{proof}
须证明该函子满足“站点”引理 \ref{sites-lemma-equivalence}
的假设 (1)--(5)。该函子连续且余连续是清楚的，这证明 (1) 与 (2)。
由于该函子全忠实，(3) 与 (4) 立即成立。按定义，代数空间有概形给出的覆盖，
故 (5) 成立。
\end{proof}

\begin{remark}
\label{spaces-properties-remark-explain-equivalence}
下面解释引理 \ref{spaces-properties-lemma-compare-etale-sites} 的含义。设 $S$ 为概形，
$X$ 为 $S$ 上的代数空间，$\mathcal{F}$ 为 $X$ 的小 \'etale
站点 $X_\etale$ 上的层。该引理说明：$X_{spaces, \etale}$
上存在唯一的层 $\mathcal{F}'$，其到子范畴 $X_\etale$
的限制为 $\mathcal{F}$。若 $U \to X$ 是代数空间的 \'etale 态射，
应如何计算 $\mathcal{F}'(U)$？由代数空间的定义，存在概形 $U'$
及满 \'etale 态射 $U' \to U$。于是 $\{U' \to U\}$ 是
$X_{spaces, \etale}$ 中的覆盖，故得到等化子图表
$$
\xymatrix{
\mathcal{F}'(U) \ar[r] &
\mathcal{F}(U') \ar@<1ex>[r] \ar@<-1ex>[r] &
\mathcal{F}(U' \times_U U').
}
$$
注意，$U' \times_U U'$ 是概形，因而这里可以写
$\mathcal{F}$ 而非 $\mathcal{F}'$。这便说明了给定层
$\mathcal{F}$ 时如何计算 $\mathcal{F}'$。
\end{remark}

\begin{definition}
\label{spaces-properties-definition-affine-etale-site}
设 $S$ 为概形，$\Sch_{fppf}$ 为包含 $S$ 的大 fppf 站点，
$\Sch_\etale$ 为相应的大 \'etale 站点（即有相同的底范畴）。
设 $X$ 为 $S$ 上的代数空间。$X$ 的站点
{\it $X_{affine, \etale}$} 定义如下：
\begin{enumerate}
\item $X_{affine, \etale}$ 的对象是态射 $\varphi : U \to X$，
其中 $U \in \Ob((\Sch/S)_\etale)$ 为仿射概形，$\varphi$
为 \'etale 态射；
\item $X_{affine, \etale}$ 的态射
$(\varphi : U \to X) \to (\varphi' : U' \to X)$
由满足 $\varphi = \varphi' \circ \chi$ 的概形态射
$\chi : U \to U'$ 给出；
\item $X_{affine, \etale}$ 的态射族
$\{\varphi_i : (U_i \to X) \to (U \to X)\}_{i \in I}$ 是覆盖，
当且仅当 $\{U_i \to U\}$ 是标准 \'etale 覆盖；见“拓扑”定义
\ref{topologies-definition-standard-etale}。
\end{enumerate}
照例，依“集合”引理 \ref{sets-lemma-coverings-site}，
选取一组此类覆盖，从而把 $X_{affine, \etale}$ 变为站点。
\end{definition}

\begin{lemma}
\label{spaces-properties-lemma-alternative}
设 $S$ 为概形，$X$ 为 $S$ 上的代数空间。函子
$X_{affine, \etale} \to X_\etale$ 特殊余连续，并诱导从
$\Sh(X_{affine, \etale})$ 到 $\Sh(X_\etale)$ 的拓扑斯等价。
\end{lemma}

\begin{proof}
从略。提示：与“拓扑”引理
\ref{topologies-lemma-affine-big-site-etale} 的证明比较。
\end{proof}

\begin{definition}
\label{spaces-properties-definition-etale-topos}
设 $S$ 为概形，$X$ 为 $S$ 上的代数空间。$X$ 的
{\it \'etale 拓扑斯}，更精确地说，$X$ 的{\it 小 \'etale 拓扑斯}，
是 $X_\etale$ 上集合层的范畴 $\Sh(X_\etale)$。
\end{definition}

\noindent
由引理 \ref{spaces-properties-lemma-compare-etale-sites}，
$\Sh(X_\etale) = \Sh(X_{spaces, \etale})$，故也可将其视为
$X_{spaces, \etale}$ 上集合层的范畴。类似地，由引理
\ref{spaces-properties-lemma-alternative}，
$\Sh(X_\etale) = \Sh(X_{affine, \etale})$。
事实证明，该拓扑斯关于代数空间的态射具有函子性。精确陈述如下。

\begin{lemma}
\label{spaces-properties-lemma-functoriality-etale-site}
设 $S$ 为概形，$f : X \to Y$ 为 $S$ 上代数空间的态射。
\begin{enumerate}
\item 连续函子
$$
Y_{spaces, \etale} \longrightarrow X_{spaces, \etale}, \quad
V \longmapsto X \times_Y V
$$
诱导站点态射
$$
f_{spaces, \etale} :
X_{spaces, \etale}
\to
Y_{spaces, \etale}.
$$
\item 规则 $f \mapsto f_{spaces, \etale}$ 与复合相容；换言之，
$(f \circ g)_{spaces, \etale}
= f_{spaces, \etale} \circ g_{spaces, \etale}$（见“站点”定义
\ref{sites-definition-composition-morphisms-sites}）。
\item 与 $f_{spaces, \etale}$ 相伴的拓扑斯态射通过引理
\ref{spaces-properties-lemma-compare-etale-sites} 诱导拓扑斯态射
$f_{small} : \Sh(X_\etale) \to \Sh(Y_\etale)$，其构造与复合相容。
\item 若 $f$ 是代数空间的可表示态射，则 $f_{small}$
来自站点态射 $X_\etale \to Y_\etale$，它对应于连续函子
$V \mapsto X \times_Y V$。
\end{enumerate}
\end{lemma}

\begin{proof}
先证明 (1) 中的函子满足“站点”命题
\ref{sites-proposition-get-morphism} 的假设。须证明
$Y_{spaces, \etale}$ 有终对象（即 $Y$），并且该函子把它变为
$X_{spaces, \etale}$ 中的终对象（即 $X$）。这很清楚，因为在任意范畴中
$X \times_Y Y = X$。接着须证明 $Y_{spaces, \etale}$ 有纤维积。
这是成立的，因为代数空间范畴有纤维积，并且若 $V$、$V'$
在 $Y$ 上 \'etale，则 $V \times_Y V'$ 也在 $Y$ 上 \'etale
（见上面的引理 \ref{spaces-properties-lemma-composition-etale} 与
\ref{spaces-properties-lemma-base-change-etale}）。因此该命题可用，并得到 (1)
所述站点态射。

\medskip\noindent
展开定义即得 (2)。利用上面引理
\ref{spaces-properties-lemma-compare-etale-sites} 对 $X$ 与 $Y$ 给出的等价，(3) 是清楚的。
(4) 也成立，因为若 $f$ 可表示，则上述函子组成交换图表
$$
\xymatrix{
X_\etale \ar[r] &
X_{spaces, \etale} \\
Y_\etale \ar[r] \ar[u] &
Y_{spaces, \etale} \ar[u]
}
$$
的范畴。
\end{proof}

\noindent
在描述 $X$ 上的层与 $Y$ 上的层之间的关系时，还可以比上述引理更进一步。
具体地，可用 $f$-映射来表述；比较“层”定义
\ref{sheaves-definition-f-map}。如下。

\begin{definition}
\label{spaces-properties-definition-f-map}
设 $S$ 为概形，$f : X \to Y$ 为 $S$ 上代数空间的态射，
$\mathcal{F}$ 为 $X_\etale$ 上的集合层，$\mathcal{G}$
为 $Y_\etale$ 上的集合层。一个
{\it $f$-映射 $\varphi : \mathcal{G} \to \mathcal{F}$}
是由交换图表
$$
\xymatrix{
U \ar[d]_g \ar[r] & X \ar[d]^f \\
V \ar[r] & Y
}
$$
指标化的一族映射
$\varphi_{(U, V, g)} : \mathcal{G}(V) \to \mathcal{F}(U)$，
其中 $U \in X_\etale$、$V \in Y_\etale$，并且只要给定扩充图表
$$
\xymatrix{
U' \ar[r] \ar[d]_{g'} & U \ar[d]_g \ar[r] & X \ar[d]^f \\
V' \ar[r] & V \ar[r] & Y
}
$$
其中 $V' \to V$ 与 $U' \to U$ 为概形的 \'etale 态射，图表
$$
\xymatrix{
\mathcal{G}(V)
\ar[rr]_{\varphi_{(U, V, g)}}
\ar[d]_{\text{限制 }\mathcal{G}} & &
\mathcal{F}(U)
\ar[d]^{\text{限制 }\mathcal{F}} \\
\mathcal{G}(V')
\ar[rr]^{\varphi_{(U', V', g')}} & &
\mathcal{F}(U')
}
$$
就交换。
\end{definition}

\begin{lemma}
\label{spaces-properties-lemma-f-map}
设 $S$ 为概形，$f : X \to Y$ 为 $S$ 上代数空间的态射，
$\mathcal{F}$ 为 $X_\etale$ 上的集合层，$\mathcal{G}$
为 $Y_\etale$ 上的集合层。下列三个集合之间存在典范双射：
\begin{enumerate}
\item 映射 $\mathcal{G} \to f_{small, *}\mathcal{F}$ 的集合；
\item 映射 $f_{small}^{-1}\mathcal{G} \to \mathcal{F}$ 的集合；
\item $f$-映射 $\varphi : \mathcal{G} \to \mathcal{F}$ 的集合。
\end{enumerate}
\end{lemma}

\begin{proof}
注意，(1) 与 (2) 相同，因为函子 $f_{small, *}$
与 $f_{small}^{-1}$ 互为伴随。设
$\alpha : f_{small}^{-1}\mathcal{G} \to \mathcal{F}$
为 $Y_\etale$ 上层的映射，并给定定义
\ref{spaces-properties-definition-f-map} 中的图表
$$
\xymatrix{
U \ar[d]_g \ar[r]_{j_U} & X \ar[d]^f \\
V \ar[r]^{j_V} & Y
}
$$
。由图表的交换性，还得到映射
$g_{small}^{-1}(j_V)^{-1}\mathcal{G} \to (j_U)^{-1}\mathcal{F}$
（关于局部化函子的描述，比较“站点”第
\ref{sites-section-localize} 节）。因此当然得到映射
$\varphi_{(V, U, g)} :
\mathcal{G}(V) = (j_V)^{-1}\mathcal{G}(V)
\to
(j_U)^{-1}\mathcal{F}(U) = \mathcal{F}(U)$.
这一规则与进一步的限制相容，并定义从 $\mathcal{G}$ 到
$\mathcal{F}$ 的 $f$-映射；其验证从略。

\medskip\noindent
反之，设给定 $f$-映射
$\varphi = (\varphi_{(U, V, g)})$。令 $\mathcal{G}'$
（相应地，$\mathcal{F}'$）表示 $\mathcal{G}$（相应地，
$\mathcal{F}$）到 $Y_{spaces, \etale}$（相应地，
$X_{spaces, \etale}$）的扩张；见引理
\ref{spaces-properties-lemma-compare-etale-sites}。于是须构造层的映射
$$
\mathcal{G}' \longrightarrow (f_{spaces, \etale})_*\mathcal{F}'
$$
为此，设 $V \to Y$ 为代数空间的 \'etale 态射。须构造集合的映射
$$
\mathcal{G}'(V) \to \mathcal{F}'(X \times_Y V)
$$
选取满 \'etale 态射 $V' \to V$，其中 $V'$ 为概形；随后选取
满 \'etale 态射 $U' \to X \times_U V'$，其中 $U'$ 为概形。
得到概形态射 $g' : U' \to V'$，以及概形态射
$$
g'' : U' \times_{X \times_Y V} U' \longrightarrow V' \times_V V'
$$
考虑图表
$$
\xymatrix{
\mathcal{F}'(X \times_Y V) \ar[r] &
\mathcal{F}(U') \ar@<1ex>[r] \ar@<-1ex>[r] &
\mathcal{F}(U' \times_{X \times_Y V} U') \\
\mathcal{G}'(X \times_Y V) \ar[r] \ar@{..>}[u] &
\mathcal{G}(V') \ar@<1ex>[r] \ar@<-1ex>[r] \ar[u]_{\varphi_{(U', V', g')}} &
\mathcal{G}(V' \times_V V') \ar[u]_{\varphi_{(U'', V'', g'')}}
}
$$
映射 $\varphi_{...}$ 与限制的相容性表明右侧两个方块交换。
$X_{spaces, \etale}$ 中覆盖的定义表明两条水平行都是等化子图表，
故得到虚线箭头。这些箭头与限制映射相容的证明留给读者。
\end{proof}

\noindent
若代数空间态射 $X \to Y$ 是 \'etale 的，则拓扑斯态射
$\Sh(X_\etale) \to \Sh(Y_\etale)$ 是一个局部化。精确地说：

\begin{lemma}
\label{spaces-properties-lemma-etale-morphism-topoi}
设 $S$ 为概形，$f : X \to Y$ 为 $S$ 上代数空间的态射。
假设 $f$ 是 \'etale 的。此时有函子
$$
j : X_\etale \to Y_\etale, \quad
(\varphi : U \to X) \mapsto (f \circ \varphi : U \to Y)
$$
它是余连续的。拓扑斯态射 $f_{small}$ 正是与 $j$ 相伴的拓扑斯态射，
参见“站点”引理 \ref{sites-lemma-cocontinuous-morphism-topoi}。
此外，$j$ 也是连续的，故可应用“站点”引理
\ref{sites-lemma-when-shriek}。特别地，对 $Y_\etale$ 上的每个层
$\mathcal{G}$，都有
$f_{small}^{-1}\mathcal{G}(U) = \mathcal{G}(jU)$。
\end{lemma}

\begin{proof}
由代数空间 \'etale 态射的定义本身（定义 \ref{spaces-properties-definition-etale}）可知，
所给规则确实定义了上述函子 $j$。$j$ 余连续且连续是清楚的，因为
$Y_\etale$ 中 $j(\varphi : U \to X)$ 的一个覆盖
$\{U_i \to U\}$，恰好就是 $X_\etale$ 中
$(\varphi : U \to X)$ 的一个覆盖。只须再证明 $j$ 诱导的拓扑斯态射
与 $f_{small}$ 相同。为此考虑范畴图
$$
\xymatrix{
X_\etale \ar[r] \ar[d]^j &
X_{spaces, \etale} \ar@/_/[d]_{j_{spaces}} \\
Y_\etale \ar[r] &
Y_{spaces, \etale} \ar@/_/[u]_{v : V \mapsto X \times_Y V}
}
$$
其中函子 $j_{spaces}$ 是 $j$ 到范畴 $X_{spaces, \etale}$ 的显然延拓。
因此内方块交换。事实上，$j_{spaces}$ 可与“站点”第
\ref{sites-section-localize} 节中讨论的局部化函子
$j_X : Y_{spaces, \etale}/X \to Y_{spaces, \etale}$
等同。于是由“站点”引理 \ref{sites-lemma-localize-given-products}，
余连续函子 $j_{spaces}$ 与图中的函子 $v$ 诱导同一个拓扑斯态射。
再由“站点”引理 \ref{sites-lemma-composition-cocontinuous}，
内方块（由站点之间的余连续函子组成）的交换性给出相伴拓扑斯态射的交换图。
最后根据引理 \ref{spaces-properties-lemma-functoriality-etale-site} 中 $f_{small}$ 的构造即得结论。
\end{proof}

\noindent
上面的引理说明，经由代数空间的 \'etale 态射 $f : X \to Y$
拉回 $\mathcal{G}$，不过是把 $\mathcal{G}$ 限制到范畴 $X_\etale$。
我们常用简记
\begin{equation}
\label{spaces-properties-equation-restrict}
\mathcal{G}|_{X_\etale} = f_{small}^{-1}\mathcal{G}
\end{equation}
来表示这一点。注意，在这种情形下，引理中的函子
$j : X_\etale \to Y_\etale$ 是忠实的，但一般并非全忠实。
第 \ref{spaces-properties-section-localize} 节将以更技术性的方式讨论此事。

\begin{lemma}
\label{spaces-properties-lemma-pushforward-etale-base-change}
设 $S$ 为概形，并设
$$
\xymatrix{
X' \ar[r] \ar[d]_{f'} & X \ar[d]^f \\
Y' \ar[r]^g & Y
}
$$
为 $S$ 上代数空间的笛卡儿方块。设 $\mathcal{F}$ 为 $X_\etale$ 上的层。
若 $g$ 是 \'etale 的，则
\begin{enumerate}
\item $f'_{small, *}(\mathcal{F}|_{X'}) = (f_{small, *}\mathcal{F})|_{Y'}$
在 $\Sh(Y'_\etale)$ 中成立\footnote{还可得
$(f')_{small}^{-1}(\mathcal{G}|_{Y'}) = (f_{small}^{-1}\mathcal{G})|_{X'}$
，这是因为该图交换且有式 (\ref{spaces-properties-equation-restrict})}；
\item 若 $\mathcal{F}$ 是阿贝尔层，则
$R^if'_{small, *}(\mathcal{F}|_{X'}) = (R^if_{small, *}\mathcal{F})|_{Y'}$.
\end{enumerate}
\end{lemma}

\begin{proof}
考虑如下函子图
$$
\xymatrix{
X'_{spaces, \etale} \ar[r]_j &
X_{spaces, \etale} \\
Y'_{spaces, \etale} \ar[r]^j \ar[u]^{V' \mapsto V' \times_{Y'} X'} &
Y_{spaces, \etale} \ar[u]_{V \mapsto V \times_Y X}
}
$$
横箭头均为局部化，而纵箭头诱导站点态射。因此“站点”引理
\ref{sites-lemma-localize-morphism} 的最后一个陈述给出 (1)。
为证明 (2)，把 (1) 应用于 $\mathcal{F}$ 的一个内射分解，并使用限制函子
正合且保持内射对象这一事实（见“站点上的上同调”引理
\ref{sites-cohomology-lemma-cohomology-of-open}）。
\end{proof}

\noindent
下面的引理表明，可把代数空间小 \'etale 站点上的一个层看作一族相容的层：
它们分别定义在该空间上 \'etale 的各概形的小 \'etale 站点上。
请注意，引理中的全部比较映射 $c_f$ 都是同构；这与“拓扑”引理
\ref{topologies-lemma-characterize-sheaf-big-etale} 以及
$X_\etale$ 的对象之间每个态射均为 \'etale 态射这一事实相符。

\begin{lemma}
\label{spaces-properties-lemma-characterize-sheaf-small-etale}
设 $S$ 为概形，$X$ 为 $S$ 上的代数空间。$X_\etale$ 上的一个层
$\mathcal{F}$ 由下列数据给出：
\begin{enumerate}
\item 对每个 $U \in \Ob(X_\etale)$，给定 $U_\etale$ 上的层
$\mathcal{F}_U$；
\item 对 $X_\etale$ 中每个 $f : U' \to U$，给定同构
$c_f : f_{small}^{-1}\mathcal{F}_U \to \mathcal{F}_{U'}$.
\end{enumerate}
这些数据须满足如下条件：对 $X_\etale$ 中任意
$f : U' \to U$ 与 $g : U'' \to U'$，复合
$c_g \circ g_{small}^{-1} c_f$ 等于 $c_{f \circ g}$。
\end{lemma}

\begin{proof}
可按引理 \ref{spaces-properties-lemma-etale-morphism-topoi} 解释 $g_{small}^{-1}$。
于是本引理由关于站点的一般事实推出，见“站点”引理
\ref{sites-lemma-glue-sheaves-absolute}。
\end{proof}

\noindent
设 $S$ 为概形，$X$ 为 $S$ 上的代数空间。取任意满的 \'etale 态射
$\varphi : U \to X$ 所给出的 $X$ 的表示 $X = U/R$，见“空间”定义
\ref{spaces-definition-presentation}。特别地，我们得到群胚
$(U, R, s, t, c, e, i)$，使得 $j = (t, s) : R \to U \times_S U$，
见“群胚”引理 \ref{groupoids-lemma-equivalence-groupoid}。

\begin{lemma}
\label{spaces-properties-lemma-descent-sheaf}
沿用上述 $S$、$\varphi : U \to X$ 及 $(U, R, s, t, c, e, i)$。
对 $X_\etale$ 上任意层 $\mathcal{F}$，层\footnote{在本引理及其证明中，
我们把 $\varphi_{small}^{-1}$ 简写为 $\varphi^{-1}$，其他拉回也同样简写。}
$\mathcal{G} = \varphi^{-1}\mathcal{F}$ 带有典范同构
$$
\alpha :
t^{-1}\mathcal{G}
\longrightarrow
s^{-1}\mathcal{G}
$$
使得图
$$
\xymatrix{
& \text{pr}_1^{-1}t^{-1}\mathcal{G} \ar[r]_-{\text{pr}_1^{-1}\alpha} &
\text{pr}_1^{-1}s^{-1}\mathcal{G} \ar@{=}[rd] & \\
\text{pr}_0^{-1}s^{-1}\mathcal{G} \ar@{=}[ru] & & &
c^{-1}s^{-1}\mathcal{G} \\
&
\text{pr}_0^{-1}t^{-1}\mathcal{G} \ar[lu]^{\text{pr}_0^{-1}\alpha} \ar@{=}[r] &
c^{-1}t^{-1}\mathcal{G} \ar[ru]_{c^{-1}\alpha}
}
$$
交换。函子 $\mathcal{F} \mapsto (\mathcal{G}, \alpha)$ 给出
$X_\etale$ 上的层与上述二元组 $(\mathcal{G}, \alpha)$ 所成范畴之间的等价。
\end{lemma}

\begin{proof}[引理 \ref{spaces-properties-lemma-descent-sheaf} 的第一种证明]
令 $\mathcal{C} = X_{spaces, \etale}$。由引理
\ref{spaces-properties-lemma-etale-morphism-topoi} 及其证明，
$U_{spaces, \etale} = \mathcal{C}/U$，且拉回函子 $\varphi^{-1}$
就是限制函子。此外，$\{U \to X\}$ 是站点 $\mathcal{C}$ 的覆盖，
且 $R = U \times_X U$。同构 $\alpha$ 就是典范等同
$$
\left(\mathcal{F}|_{\mathcal{C}/U}\right)|_{\mathcal{C}/U \times_X U}
=
\left(\mathcal{F}|_{\mathcal{C}/U}\right)|_{\mathcal{C}/U \times_X U}
$$
而该图的交换性就是胶合数据的上循环条件。因此本引理是层的胶合这一结果的
特殊情形，见“站点”第 \ref{sites-section-glueing-sheaves} 节。
\end{proof}

\begin{proof}[引理 \ref{spaces-properties-lemma-descent-sheaf} 的第二种证明]
$\alpha$ 的存在性来自 $\varphi \circ t = \varphi \circ s$，以及拉回关于态射
具有函子性这一事实，见引理 \ref{spaces-properties-lemma-functoriality-etale-site}。
完全同样地，即由拉回的函子性可见，同构 $\alpha$ 嵌入上述交换图。
构造 $\mathcal{F} \mapsto (\varphi^{-1}\mathcal{F}, \alpha)$
关于层 $\mathcal{F}$ 显然具有函子性。因此我们得到了所需函子。

\medskip\noindent
反过来，设 $(\mathcal{G}, \alpha)$ 为这样的二元组，并设 $V \to X$
为 $X_\etale$ 的对象。这时态射 $V' = U \times_X V \to V$
是概形间满的 \'etale 态射，故 $\{V' \to V\}$ 是 $V$ 的 \'etale 覆盖。
令 $\mathcal{G}' = (V' \to V)^{-1}\mathcal{G}$。由于
$R = U \times_X U$，其中 $t = \text{pr}_0$ 且 $s = \text{pr}_0$，
可见 $V' \times_V V' = R \times_X V$，而投影态射
$s', t' : V' \times_V V' \to V'$ 分别是 $t$ 与 $s$ 的拉回。
因此 $\alpha$ 拉回为同构
$\alpha' : (t')^{-1}\mathcal{G}' \to (s')^{-1}\mathcal{G}'$。
据此直接定义
$$
\xymatrix{
\mathcal{F}(V) \ar@{=}[r] &
\text{等化子}(\mathcal{G}(V') \ar@<1ex>[r] \ar@<-1ex>[r] &
\mathcal{G}(V' \times_V V').
}
$$
略去这确实定义一个层的验证。若存在态射 $V \to U$，则
$\mathcal{G}(V) = \mathcal{F}(V)$；这是因为此时上述等化子为
$H^0(\{V' \to V\}, \mathcal{G}) = \mathcal{G}(V)$。
\end{proof}







\section{小 \'etale 站点的点}
\label{spaces-properties-section-points-small-etale-site}

\noindent
本节是“\'Etale 上同调”第
\ref{etale-cohomology-section-stalks} 节的对应版本。

\begin{definition}
\label{spaces-properties-definition-geometric-point}
设 $S$ 为概形，$X$ 为 $S$ 上的代数空间。
\begin{enumerate}
\item $X$ 的一个{\it 几何点}是态射
$\overline{x} : \Spec(k) \to X$，其中 $k$ 为代数闭域。
我们常滥用记号写成 $\overline{x} = \Spec(k)$。
\item 每个几何点 $\overline{x}$ 都有相应的“像”点 $x \in |X|$。
我们称 $\overline{x}$ 为一个{\it 位于 $x$ 上方的几何点}。
\end{enumerate}
\end{definition}

\noindent
事实表明，可以完全仿照概形小 \'etale 站点的情形，在几何点处取
$X_\etale$ 上层的茎。为此定义如下 \'etale 邻域概念。

\begin{definition}
\label{spaces-properties-definition-etale-neighbourhood}
设 $S$ 为概形，$X$ 为 $S$ 上的代数空间，并设 $\overline{x}$
为 $X$ 的一个几何点。
\begin{enumerate}
\item $X$ 中 $\overline{x}$ 的一个{\it \'etale 邻域}是交换图
$$
\xymatrix{
& U \ar[d]^\varphi \\
{\bar x} \ar[r]^{\bar x} \ar[ur]^{\bar u} & X
}
$$
其中 $\varphi$ 是 $S$ 上代数空间的 \'etale 态射。我们用记号
$\varphi : (U, \overline{u}) \to (X, \overline{x})$ 表示这种情形。
\item 一个{\it \'etale 邻域态射}
$(U, \overline{u}) \to (U', \overline{u}')$
是满足 $\overline{u}' = h \circ \overline{u}$ 的 $X$-态射
$h : U \to U'$。
\end{enumerate}
\end{definition}

\noindent
注意，我们允许 $U$ 为代数空间。当取 $X_\etale$ 上层的茎时，必须限制到
属于 $X_\etale$ 的那些 $U$，故此时只考虑 $U$ 为概形的情形。
另一种做法是使用站点 $X_{space, \etale}$ 并考虑所有 \'etale 邻域。
由下述引理的最后一个断言，两种做法并无区别。

\begin{lemma}
\label{spaces-properties-lemma-cofinal-etale}
设 $S$ 为概形，$X$ 为 $S$ 上的代数空间，并设 $\overline{x}$
为 $X$ 的几何点。\'Etale 邻域范畴是余滤的。更精确地：
\begin{enumerate}
\item 设 $(U_i, \overline{u}_i)_{i = 1, 2}$ 为 $X$ 中
$\overline{x}$ 的两个 \'etale 邻域。则存在第三个 \'etale 邻域
$(U, \overline{u})$ 以及态射
$(U, \overline{u}) \to (U_i, \overline{u}_i)$, $i = 1, 2$.
\item 设 $h_1, h_2: (U, \overline{u}) \to (U', \overline{u}')$
为 $\overline{s}$ 的 \'etale 邻域之间的两个态射。则存在 \'etale 邻域
$(U'', \overline{u}'')$ 及态射
$h : (U'', \overline{u}'') \to (U, \overline{u})$
等化 $h_1$ 与 $h_2$，即满足 $h_1 \circ h = h_2 \circ h$。
\end{enumerate}
此外，给定任意 \'etale 邻域
$(U, \overline{u}) \to (X, \overline{x})$
，存在 \'etale 邻域态射
$(U', \overline{u}') \to (U, \overline{u})$
，其中 $U'$ 为概形。
\end{lemma}

\begin{proof}
对 (1)，考虑纤维积 $U = U_1 \times_X U_2$。由于 \'etale 态射在基变换及
复合下保持，它在 $U_1$ 与 $U_2$ 上都是 \'etale 的；见引理
\ref{spaces-properties-lemma-base-change-etale} 与 \ref{spaces-properties-lemma-composition-etale}。
由 $(\overline{u}_1, \overline{u}_2)$ 定义的映射 $\overline{u} \to U$
使其成为同时映到 $U_1$ 与 $U_2$ 的 \'etale 邻域。

\medskip\noindent
对 (2)，把 $U''$ 定义为纤维积
$$
\xymatrix{
U'' \ar[r] \ar[d] & U \ar[d]^{(h_1, h_2)} \\
U' \ar[r]^-\Delta & U' \times_X U'.
}
$$
由于 $\overline{u}$ 与 $\overline{u}'$ 在 $X$ 上都等于 $\overline{x}$，
可见 $\overline{u}'' = (\overline{u}, \overline{u}')$ 是 $U''$ 的几何点。
特别地，$U'' \not = \emptyset$。此外，由于 $U'$ 在 $X$ 上 \'etale，
纤维积 $U'\times_X U'$ 也同样 \'etale（如上面对
$U_1 \times_X U_2$ 的讨论）。故由引理 \ref{spaces-properties-lemma-etale-permanence}，
纵箭头 $(h_1, h_2)$ 是 \'etale 的。于是基变换表明 $U''$ 在 $U'$ 上
\'etale，进而也在 $X$ 上 \'etale（因为 \'etale 态射的复合仍为
\'etale 态射）。所以 $(U'', \overline{u}'')$ 解出了 (2) 所提问题。

\medskip\noindent
为证最后一个断言，任取满的 \'etale 态射 $U' \to U$，其中 $U'$ 为概形。
于是 $U' \times_U \overline{u}$ 是在
$\overline{u} = \Spec(k)$ 上满且 \'etale 的概形，其中 $k$ 代数闭。
由“态射”引理 \ref{morphisms-lemma-etale-over-field}，
$U' \times_U \overline{u} \to \overline{u}$ 有截面，后者给出所需的
$\overline{u}'$。
\end{proof}

\begin{lemma}
\label{spaces-properties-lemma-geometric-lift-to-usual}
设 $S$ 为概形，$X$ 为 $S$ 上的代数空间。设
$\overline{x} : \Spec(k) \to X$ 为 $X$ 的一个几何点，位于
$x \in |X|$ 上方。设 $\varphi : U \to X$ 为代数空间的 \'etale 态射，
并取满足 $\varphi(u) = x$ 的 $u \in |U|$。则存在位于 $u$ 上方的几何点
$\overline{u} : \Spec(k) \to U$，使得
$\overline{x} = \varphi \circ \overline{u}$。
\end{lemma}

\begin{proof}
选取仿射概形 $U'$、点 $u' \in U'$ 以及把 $u'$ 映到 $u$ 的 \'etale 态射
$U' \to U$。若能对 $(U', u') \to (X, x)$ 证明本引理，结论即随之成立。
故可设 $U$ 为概形；特别地，$U \to X$ 可表示。考察笛卡儿图
$$
\xymatrix{
\Spec(k) \times_{\overline{x}, X, \varphi} U
\ar[d]_{\text{pr}_1} \ar[r]_-{\text{pr}_2} & U
\ar[d]^\varphi \\
\Spec(k) \ar[r]^-{\overline{x}} & X
}
$$
投影 $\text{pr}_1$ 是 \'etale 态射的基变换，故为 \'etale；见引理
\ref{spaces-properties-lemma-base-change-etale}。因此概形
$\Spec(k) \times_{\overline{x}, X, \varphi} U$ 是若干 $k$ 的有限可分扩域
之谱的不交并，见“态射”引理 \ref{morphisms-lemma-etale-over-field}。
但 $k$ 代数闭，故所有这些扩张都平凡，于是
$\Spec(k) \times_{\overline{x}, X, \varphi} U$ 是若干
$\Spec(k)$ 的不交并；其中每一份都对应一个满足
$\varphi \circ \overline{u} = \overline{x}$ 的几何点 $\overline{u}$。
由引理 \ref{spaces-properties-lemma-points-cartesian}，映射
$$
|\Spec(k) \times_{\overline{x}, X, \varphi} U|
\longrightarrow
|\Spec(k)| \times_{|X|} |U|
$$
是满的，故可选取位于 $u$ 上方的 $\overline{u}$。
\end{proof}

\begin{lemma}
\label{spaces-properties-lemma-geometric-lift-to-cover}
设 $S$ 为概形，$X$ 为 $S$ 上的代数空间，$\overline{x}$ 为 $X$
的几何点，且 $(U, \overline{u})$ 为 $\overline{x}$ 的 \'etale 邻域。
设 $\{\varphi_i : U_i \to U\}_{i \in I}$ 为
$X_{spaces, \etale}$ 中的一个 \'etale 覆盖。则存在 $i \in I$ 及
$\overline{u}_i : \overline{x} \to U_i$，使得
$\varphi_i : (U_i, \overline{u}_i) \to (U, \overline{u})$
为 \'etale 邻域态射。
\end{lemma}

\begin{proof}
令 $u \in |U|$ 为 $\overline{u}$ 的像。由于
$|U| = \bigcup_{i \in I} \varphi_i(|U_i|)$，存在某个 $i$ 以及映到
$x$ 的点 $u_i \in U_i$。对 $(U_i, u_i) \to (U, u)$ 与
$\overline{u}$ 应用引理 \ref{spaces-properties-lemma-geometric-lift-to-usual}，即得所需几何点。
\end{proof}

\begin{definition}
\label{spaces-properties-definition-stalk}
设 $S$ 为概形，$X$ 为 $S$ 上的代数空间，$\mathcal{F}$ 为
$X_\etale$ 上的预层，并设 $\overline{x}$ 为 $X$ 的几何点。
$\mathcal{F}$ 在 $\overline{x}$ 处的{\it 茎}是
$$
\mathcal{F}_{\bar x}
=
\colim_{(U, \overline{u})} \mathcal{F}(U)
$$
其中 $(U, \overline{u})$ 遍历 $X$ 中 $\overline{x}$ 的所有 \'etale 邻域，
且 $U \in \Ob(X_\etale)$。
\end{definition}

\noindent
由引理 \ref{spaces-properties-lemma-cofinal-etale}，此余极限的指标范畴是滤过的，
它就是 $X_\etale$ 中 \'etale 邻域范畴的反范畴。更精确地，
引理 \ref{spaces-properties-lemma-cofinal-etale} 说明所有 \'etale 邻域所成范畴的反范畴
是滤过的，而由属于 $X_\etale$ 的那些邻域组成的全子范畴是共尾子范畴，
因而也滤过。

\medskip\noindent
这意味着，可把 $\mathcal{F}_{\overline{x}}$ 的一个元素看作三元组
$(U, \overline{u}, \sigma)$，其中 $U \in \Ob(X_\etale)$ 且
$\sigma \in \mathcal{F}(U)$。两个三元组
$(U, \overline{u}, \sigma)$、$(U', \overline{u}', \sigma')$
定义茎中的同一个元素，当且仅当存在第三个 \'etale 邻域
$(U'', \overline{u}'')$, $U'' \in \Ob(X_\etale)$
以及 \'etale 邻域态射
$h : (U'', \overline{u}'') \to (U, \overline{u})$,
$h' : (U'', \overline{u}'') \to (U', \overline{u}')$，使得
$h^*\sigma = (h')^*\sigma'$ 在 $\mathcal{F}(U'')$ 中成立。
见“范畴”第 \ref{categories-section-directed-colimits} 节。

\medskip\noindent
这还说明，若 $\mathcal{F}'$ 是 $X_{spaces, \etale}$ 上对应于
$X_\etale$ 上 $\mathcal{F}$ 的层，则
\begin{equation}
\label{spaces-properties-equation-stalk-spaces-etale}
\mathcal{F}_{\overline{x}} = \colim_{(U, \overline{u})} \mathcal{F}'(U)
\end{equation}
其中的余极限现在遍历 $\overline{x}$ 的所有 \'etale 邻域。
以后常在使用 $X_\etale$ 与使用 $X_{spaces, \etale}$ 的观点之间切换，
而不再特别说明。

\medskip\noindent
特别地，若 $\mathcal{F}$ 是阿贝尔群、环等的预层，则按通常在阿贝尔群、
环等的有向余极限上定义群结构的方法，$\mathcal{F}_{\overline{x}}$
自然也是阿贝尔群、环等。

\begin{lemma}
\label{spaces-properties-lemma-stalk-gives-point}
\begin{slogan}
代数空间的一个几何点给出其 \'etale 拓扑斯的一个点。
\end{slogan}
设 $S$ 为概形，$X$ 为 $S$ 上的代数空间，并设 $\overline{x}$
为 $X$ 的几何点。考虑函子
$$
u : X_\etale \longrightarrow \textit{Sets}, \quad
U \longmapsto |U_{\overline{x}}|
$$
则 $u$ 定义站点 $X_\etale$ 的一个点 $p$
（“站点”定义 \ref{sites-definition-point}），且与之相伴的茎函子
$\mathcal{F} \mapsto \mathcal{F}_p$
（“站点”式 \ref{sites-equation-stalk}）正是上面定义的函子
$\mathcal{F} \mapsto \mathcal{F}_{\overline{x}}$。
\end{lemma}

\begin{proof}
在引理 \ref{spaces-properties-lemma-geometric-lift-to-usual} 的证明中已经看到，概形
$U_{\overline{x}} = \overline{x} \times_X U$ 是若干与
$\overline{x}$ 同构的概形的不交并。因此也可把 $|U_{\overline{x}}|$
看作 $U$ 中位于 $\overline{x}$ 上方的几何点之集，即所有能嵌入定义
\ref{spaces-properties-definition-etale-neighbourhood} 之图的态射
$\overline{u} : \overline{x} \to U$ 所成的集合。由此，$u(X)$ 为单点集，且
$u(U \times_V W) = u(U) \times_{u(V)} u(W)$
对 $X_\etale$ 中任意态射 $U \to V$ 与 $W \to V$ 都成立。
再者，给定 $X_\etale$ 中的覆盖 $\{U_i \to U\}_{i \in I}$，
由引理 \ref{spaces-properties-lemma-geometric-lift-to-cover} 可知
$\coprod u(U_i) \to u(U)$ 是满的。因此可应用“站点”命题
\ref{sites-proposition-point-limits}，故 $p$ 是站点 $X_\etale$ 的点。
最后，我们的函子 $\mathcal{F} \mapsto \mathcal{F}_{\overline{s}}$
与“站点”式 \ref{sites-equation-stalk} 中相伴于 $p$ 的函子
$\mathcal{F} \mapsto \mathcal{F}_p$ 由完全相同的余极限给出，
这证明了最后一个断言。
\end{proof}

\begin{lemma}
\label{spaces-properties-lemma-stalk-exact}
设 $S$ 为概形，$X$ 为 $S$ 上的代数空间，并设 $\overline{x}$
为 $X$ 的几何点。
\begin{enumerate}
\item 茎函子
$\textit{PAb}(X_\etale) \to \textit{Ab}$,
$\mathcal{F}  \mapsto  \mathcal{F}_{\overline{x}}$
是正合的。
\item 对 $X_\etale$ 上任意集合预层 $\mathcal{F}$，有
$(\mathcal{F}^\#)_{\overline{x}} = \mathcal{F}_{\overline{x}}$。
\item 函子
$\textit{Ab}(X_\etale) \to \textit{Ab}$,
$\mathcal{F} \mapsto \mathcal{F}_{\overline{x}}$ 是正合的。
\item 类似地，由茎函子
$\mathcal{F} \mapsto \mathcal{F}_{\overline{x}}$ 给出的函子
$\textit{PSh}(X_\etale) \to \textit{Sets}$ and
$\Sh(X_\etale) \to \textit{Sets}$ 也是正合的（见“范畴”定义
\ref{categories-definition-exact}），并与任意余极限交换。
\end{enumerate}
\end{lemma}

\begin{proof}
该结果由“站点上的模”第 \ref{sites-modules-section-stalks} 节的一般理论推出。
这是因为 $\mathcal{F} \mapsto \mathcal{F}_{\overline{x}}$ 来自
$X$ 的小 \'etale 站点的一个点，见引理 \ref{spaces-properties-lemma-stalk-gives-point}。
在概形小 \'etale 站点的情形，这些陈述中一部分的直接证明可见
“\'Etale 上同调”引理 \ref{etale-cohomology-lemma-stalk-exact} 的证明。
\end{proof}

\noindent
下面将看到，茎函子
$\mathcal{F} \mapsto \mathcal{F}_{\overline{x}}$
其实就是沿态射 $\overline{x}$ 的拉回。在此意义下，下述引理是上面引理的推广。

\begin{lemma}
\label{spaces-properties-lemma-stalk-pullback}
设 $S$ 为概形，$f : X \to Y$ 为 $S$ 上代数空间的态射。
\begin{enumerate}
\item 函子
$f_{small}^{-1} :
\textit{Ab}(Y_\etale)
\to
\textit{Ab}(X_\etale)$
是正合的。
\item 函子
$f_{small}^{-1} :
\Sh(Y_\etale)
\to
\Sh(X_\etale)$
是正合的，即它与有限极限及有限余极限交换；见“范畴”定义
\ref{categories-definition-exact}。
\item 对代数空间的任意 \'etale 态射 $V \to Y$，有
$f_{small}^{-1}h_V = h_{X \times_Y V}$。
\item 设 $\overline{x} \to X$ 为几何点，$\mathcal{G}$ 为
$Y_\etale$ 上的层。则有典范等同
$$
(f_{small}^{-1}\mathcal{G})_{\overline{x}} = \mathcal{G}_{\overline{y}}.
$$
其中 $\overline{y} = f \circ \overline{x}$。
\end{enumerate}
\end{lemma}

\begin{proof}
回忆，在引理 \ref{spaces-properties-lemma-functoriality-etale-site} 中，$f_{small}$
是通过 $f_{spaces, small}$ 定义的。(1)、(2)、(3) 是下述事实的一般推论：
$f_{spaces, \etale} :
X_{spaces, \etale}
\to
Y_{spaces, \etale}$
是站点态射。关于 (2) 见“站点”定义
\ref{sites-definition-morphism-sites}，关于 (1) 见“站点上的模”引理
\ref{sites-modules-lemma-flat-pullback-exact}，关于 (3) 见“站点”引理
\ref{sites-lemma-pullback-representable-sheaf}。

\medskip\noindent
证明 (4)。该陈述是“站点”引理 \ref{sites-lemma-point-morphism-sites}
通过引理 \ref{spaces-properties-lemma-stalk-gives-point} 得到的特殊情形。这里也给出直接证明。
注意，由引理 \ref{spaces-properties-lemma-stalk-exact}，取茎与层化交换。设
$\mathcal{G}'$ 为 $Y_{spaces, \etale}$ 上限制到 $Y_\etale$ 后等于
$\mathcal{G}$ 的层。回忆，$f_{spaces, \etale}^{-1}\mathcal{G}'$
是与下列预层相伴的层：
$$
U \longrightarrow \colim_{U \to X \times_Y V} \mathcal{G}'(V),
$$
见“站点”第 \ref{sites-section-continuous-functors} 节与第
\ref{sites-section-functoriality-PSh} 节。因此
\begin{align*}
(f_{spaces, \etale}^{-1}\mathcal{G}')_{\overline{x}}
& =
\colim_{(U, \overline{u})} f_{spaces, \etale}^{-1}\mathcal{G}'(U)
\\
& = \colim_{(U, \overline{u})}
\colim_{a : U \to X \times_Y V} \mathcal{G}'(V) \\
& = \colim_{(V, \overline{v})} \mathcal{G}'(V) \\
& = \mathcal{G}'_{\overline{y}}
\end{align*}
在第三个等号中，二元组 $(U, \overline{u})$ 与映射
$a : U \to X \times_Y V$ 对应于二元组 $(V, a \circ \overline{u})$。
由于 $\mathcal{G}'$ 的茎（相应地，
$f_{spaces, \etale}^{-1}\mathcal{G}'$ 的茎）等于 $\mathcal{G}$ 的茎
（相应地，$f_{small}^{-1}\mathcal{G}$ 的茎），见式
(\ref{spaces-properties-equation-stalk-spaces-etale})，结论成立。
\end{proof}

\begin{remark}
\label{spaces-properties-remark-stalk-pullback}
本注是“\'Etale 上同调”注
\ref{etale-cohomology-remark-stalk-pullback} 的对应版本。
设 $S$ 为概形，$X$ 为 $S$ 上的代数空间，并设
$\overline{x} : \Spec(k) \to X$ 为 $X$ 的几何点。
由“\'Etale 上同调”定理
\ref{etale-cohomology-theorem-equivalence-sheaves-point}，
$\Spec(k)_\etale$ 上层的范畴等价于集合范畴（把层送到其全局截面）。
因此，对态射 $\overline{x}$ 应用引理 \ref{spaces-properties-lemma-stalk-pullback} 的 (4)，
可知函子
$$
\Sh(X_\etale) \longrightarrow \textit{Sets}, \quad
\mathcal{F} \longmapsto \mathcal{F}_{\overline{x}}
$$
同构于函子
$$
\Sh(X_\etale)
\longrightarrow
\Sh(\Spec(k)_\etale) = \textit{Sets},
\quad
\mathcal{F} \longmapsto \overline{x}^*\mathcal{F}
$$
因此可把茎函子看作沿几何态射的拉回函子，而不只是引理
\ref{spaces-properties-lemma-stalk-gives-point} 的结果中那种抽象拓扑斯态射。
\end{remark}

\begin{remark}
\label{spaces-properties-remark-map-stalks}
设 $S$ 为概形，$X$ 为 $S$ 上的代数空间，并取 $x \in |X|$。
我们断言，对位于 $x$ 上方的任意两个几何点 $\overline{x}$ 与
$\overline{x}'$，相应茎函子彼此同构。由 $|X|$ 的定义，可找到第三个几何点
$\overline{x}''$，使得存在交换图
$$
\xymatrix{
\overline{x}'' \ar[r] \ar[d] \ar[rd]^{\overline{x}''} &
\overline{x}' \ar[d]^{\overline{x}'} \\
\overline{x} \ar[r]^{\overline{x}} & X.
}
$$
由于茎函子 $\mathcal{F} \mapsto \mathcal{F}_{\overline{x}}$
由沿态射 $\overline{x}$ 的拉回给出（其他两个也同样），
结论由拉回的函子性得到。
\end{remark}




\noindent
下述定理说明，代数空间的小 \'etale 站点有足够多的点。

\begin{theorem}
\label{spaces-properties-theorem-exactness-stalks}
设 $S$ 为概形，$X$ 为 $S$ 上的代数空间。集合层的映射
$a : \mathcal{F} \to \mathcal{G}$ 是单射（相应地，满射），当且仅当茎映射
$a_{\overline{x}} : \mathcal{F}_{\overline{x}} \to \mathcal{G}_{\overline{x}}$
对 $X$ 的所有几何点都是单射（相应地，满射）。$X_\etale$ 上的一个阿贝尔层列
是正合的，当且仅当它在 $S$ 的每个几何点处的茎上正合。
\end{theorem}

\begin{proof}
若 $X$ 为概形，本定理已知成立；见“\'Etale 上同调”定理
\ref{etale-cohomology-theorem-exactness-stalks}。选取满的 \'etale 态射
$f : U \to X$，其中 $U$ 为概形。由于 $\{U \to X\}$ 是覆盖
（在 $X_{spaces, \etale}$ 中），可通过限制到 $U$ 来检验层映射是否为单射或满射。
现在，若 $\overline{u} : \Spec(k) \to U$ 是 $U$ 的几何点，则
$(\mathcal{F}|_U)_{\overline{u}} = \mathcal{F}_{\overline{x}}$
，其中 $\overline{x} = f \circ \overline{u}$。（这从定义
$\overline{u}$ 与 $\overline{x}$ 处茎的余极限显然可见，也可由引理
\ref{spaces-properties-lemma-stalk-pullback} 推出。）故 $U$ 的结果蕴含 $X$ 的结果，证毕。
\end{proof}

\noindent
初读时应跳过下述引理。

\begin{lemma}
\label{spaces-properties-lemma-points-small-etale-site}
设 $S$ 为概形，$X$ 为 $S$ 上的代数空间，并设
$p : \Sh(pt) \to \Sh(X_\etale)$ 为 $X$ 的小 \'etale 拓扑斯的一个点。
则存在 $X$ 的几何点 $\overline{x}$，使得茎函子
$\mathcal{F} \mapsto \mathcal{F}_p$ 同构于茎函子
$\mathcal{F} \mapsto \mathcal{F}_{\overline{x}}$。
\end{lemma}

\begin{proof}
由“站点”引理 \ref{sites-lemma-point-site-topos}，站点的点与相伴拓扑斯的点
一一对应。因此可设 $p$ 由定义站点 $X_\etale$ 的一个点的函子
$u : X_\etale \to \textit{Sets}$ 给出。取对象
$U \in \Ob(X_\etale)$，使其结构态射 $j : U \to X$ 为满射。
注意，$h_U$ 是满射到终层的层。由于取茎正合，
$(h_U)_p = u(U)$ 非空（使用“站点”引理
\ref{sites-lemma-points-recover}）。取 $x \in u(U)$。由“站点”引理
\ref{sites-lemma-point-localize}，得到点
$q : \Sh(pt) \to \Sh(U_\etale)$，使得 $p = j_{small} \circ q$，
因而函子性地有 $\mathcal{F}_p = (\mathcal{F}|_U)_q$。
由“\'Etale 上同调”引理
\ref{etale-cohomology-lemma-points-small-etale-site}，存在 $U$ 的几何点
$\overline{u}$，并且对 $\mathcal{G} \in \Sh(U_\etale)$ 有函子同构
$\mathcal{G}_q = \mathcal{G}_{\overline{u}}$。令
$\overline{x} = j \circ \overline{u}$。于是引理
\ref{spaces-properties-lemma-stalk-pullback} 表明，对 $X_\etale$ 上的 $\mathcal{F}$，函子性地有
$\mathcal{F}_{\overline{x}} \cong (\mathcal{F}|_U)_{\overline{u}}$，证毕。
\end{proof}





\section{阿贝尔层的支撑}
\label{spaces-properties-section-support}

\noindent
先讨论局部截面的支撑。

\begin{lemma}
\label{spaces-properties-lemma-support-subsheaf-final}
设 $S$ 为概形，$X$ 为 $S$ 上的代数空间。设 $\mathcal{F}$ 为
$X$ 的 \'etale 拓扑斯之终对象的子层（见“站点”例
\ref{sites-example-singleton-sheaf}）。则存在唯一的开子空间
$W \subset X$，使得 $\mathcal{F} = h_W$。
\end{lemma}

\begin{proof}
该条件意味着，对 $\Ob(X_{spaces, \etale})$ 中每个
$\varphi : U \to X$，$\mathcal{F}(U)$ 或为单点集，或为空集。
特别地，局部截面总能胶合。若 $\mathcal{F}(U) \not = \emptyset$，则
$\mathcal{F}(\varphi(U)) \not = \emptyset$，因为
$\varphi(U) \subset X$ 是开子空间（引理 \ref{spaces-properties-lemma-etale-open}），且
$\{\varphi : U \to \varphi(U)\}$ 是 $X_{spaces, \etale}$ 中的覆盖。
取
$W = \bigcup_{\varphi : U \to S, \mathcal{F}(U) \not = \emptyset} \varphi(U)$
即得结论。
\end{proof}

\begin{lemma}
\label{spaces-properties-lemma-zero-over-image}
设 $S$ 为概形，$X$ 为 $S$ 上的代数空间，$\mathcal{F}$ 为
$X_{spaces, \etale}$ 上的阿贝尔层，并设
$\sigma \in \mathcal{F}(U)$ 为局部截面。则存在开子空间
$W \subset U$，使得
\begin{enumerate}
\item $W \subset U$ 是使得 $\sigma|_W = 0$ 的 $U$ 的最大开子空间；
\item 对 $X_{spaces, \etale}$ 中每个 $\varphi : V \to U$，有
$$
\sigma|_V = 0 \Leftrightarrow \varphi(V) \subset W,
$$
\item 对 $U$ 的每个几何点 $\overline{u}$，有
$$
(U, \overline{u}, \sigma) = 0\text{ 于 }\mathcal{F}_{\overline{x}}
\Leftrightarrow
\overline{u} \in W
$$
其中 $\overline{x} = (U \to X) \circ \overline{u}$。
\end{enumerate}
\end{lemma}

\begin{proof}
由于 $\mathcal{F}$ 是 \'etale 拓扑中的层，$\mathcal{F}$ 到 $U_{Zar}$ 的限制是
$U$ 上 Zariski 拓扑中的层。因此存在满足 (1) 的 Zariski 开集 $W$，
见“模”引理 \ref{modules-lemma-support-section-closed}。设
$\varphi : V \to U$ 为 $X_{spaces, \etale}$ 中的箭头。注意，
$\varphi(V) \subset U$ 是开子空间（引理 \ref{spaces-properties-lemma-etale-open}），
且 $\{V \to \varphi(V)\}$ 是 \'etale 覆盖。因此若
$\sigma|_V = 0$，则由 $\mathcal{F}$ 的层条件，
$\sigma|_{\varphi(V)} = 0$。这证明了 (2)。为证明 (3)，须证若
$(U, \overline{u}, \sigma)$ 定义 $\mathcal{F}_{\overline{x}}$ 的零元，
则 $\overline{u} \in W$。该假设意味着存在 \'etale 邻域态射
$(V, \overline{v}) \to (U, \overline{u})$，使得
$\sigma|_V = 0$。故由 (2)，$V \to U$ 的像包含于 $W$，
从而 $\overline{u} \in W$。
\end{proof}

\noindent
设 $S$ 为概形，$X$ 为 $S$ 上的代数空间，取 $x \in |X|$，并设
$\mathcal{F}$ 为 $X_\etale$ 上的层。由注 \ref{spaces-properties-remark-map-stalks}，
$\mathcal{F}$ 在位于 $x$ 上方的几何点处之茎的同构类是良定义的。

\begin{definition}
\label{spaces-properties-definition-support}
设 $S$ 为概形，$X$ 为 $S$ 上的代数空间，并设 $\mathcal{F}$
为 $X_\etale$ 上的阿贝尔层。
\begin{enumerate}
\item $\mathcal{F}$ 的{\it 支撑}是这样的点 $x \in |X|$ 所成的集合：
对位于 $x$ 上方的任意（等价地，某个）几何点 $\overline{x}$，有
$\mathcal{F}_{\overline{x}} \not = 0$。
\item 设 $\sigma \in \mathcal{F}(U)$ 为截面。$\sigma$ 的{\it 支撑}
是闭子集 $U \setminus W$，其中 $W \subset U$ 是 $U$ 中使
$\sigma$ 限制为零的最大开子集（见引理 \ref{spaces-properties-lemma-zero-over-image}）。
\end{enumerate}
\end{definition}

\begin{lemma}
\label{spaces-properties-lemma-support-section-closed}
设 $S$ 为概形，$X$ 为 $S$ 上的代数空间，$\mathcal{F}$ 为
$X_\etale$ 上的阿贝尔层，并取 $U \in \Ob(X_\etale)$ 与
$\sigma \in \mathcal{F}(U)$。
\begin{enumerate}
\item $\sigma$ 的支撑在 $|X|$ 中闭。
\item $\sigma + \sigma'$ 的支撑包含于
$\sigma, \sigma' \in \mathcal{F}(X)$ 的支撑之并。
\item 若 $\varphi : \mathcal{F} \to \mathcal{G}$ 是 $X_\etale$
上阿贝尔层的映射，则 $\varphi(\sigma)$ 的支撑包含于
$\sigma \in \mathcal{F}(U)$ 的支撑。
\item $\mathcal{F}$ 的支撑是 $\mathcal{F}$ 所有局部截面的支撑之像的并。
\item 若 $\mathcal{F} \to \mathcal{G}$ 满，则 $\mathcal{G}$ 的支撑是
$\mathcal{F}$ 的支撑的子集。
\item 若 $\mathcal{F} \to \mathcal{G}$ 单，则 $\mathcal{F}$ 的支撑是
$\mathcal{G}$ 的支撑的子集。
\end{enumerate}
\end{lemma}

\begin{proof}
(1) 由定义成立。(2) 与 (3) 成立，因为它们对 $\mathcal{F}$ 与
$\mathcal{G}$ 到 $U_{Zar}$ 的限制成立；见“模”引理
\ref{modules-lemma-support-section-closed}。(4) 是引理
\ref{spaces-properties-lemma-zero-over-image} 的 (3) 的直接推论。(5) 与 (6) 由其余各项推出。
\end{proof}

\begin{lemma}
\label{spaces-properties-lemma-support-sheaf-rings-closed}
代数空间小 \'etale 站点上的环层之支撑是闭的。
\end{lemma}

\begin{proof}
这是因为，按我们的约定，一个环等于 $0$ 当且仅当 $1 = 0$；
因此环层的支撑就是单位截面的支撑。
\end{proof}






\section{代数空间的结构层}
\label{spaces-properties-section-structure-sheaf}

\noindent
代数空间的结构层就是下述引理中的环层。

\begin{lemma}
\label{spaces-properties-lemma-sheaf-condition-holds}
设 $S$ 为概形，$X$ 为 $S$ 上的代数空间。规则
$U \mapsto \Gamma(U, \mathcal{O}_U)$ 定义 $X_\etale$ 上的环层。
\end{lemma}

\begin{proof}
由覆盖的定义及“下降”引理 \ref{descent-lemma-sheaf-condition-holds} 立即可得。
\end{proof}

\begin{definition}
\label{spaces-properties-definition-structure-sheaf}
设 $S$ 为概形，$X$ 为 $S$ 上的代数空间。$X$ 的{\it 结构层}是引理
\ref{spaces-properties-lemma-sheaf-condition-holds} 所述小 \'etale 站点
$X_\etale$ 上的环层 $\mathcal{O}_X$。
\end{definition}

\noindent
根据引理 \ref{spaces-properties-lemma-characterize-sheaf-small-etale}，层 $\mathcal{O}_X$
对应于一族 \'etale 层 $(\mathcal{O}_X)_U$，其中 $U$ 遍历
$X_\etale$ 的对象。由该引理的证明及我们的定义，显然直接有
$(\mathcal{O}_X)_U = \mathcal{O}_U$；这里 $\mathcal{O}_U$ 是“下降”定义
\ref{descent-definition-structure-sheaf} 中引入的 $U_\etale$ 的结构层。
特别地，若 $X$ 为概形，我们就恢复了 $X$ 的小 \'etale 站点上的层
$\mathcal{O}_X$。

\medskip\noindent
通过引理 \ref{spaces-properties-lemma-compare-etale-sites} 的等价
$\Sh(X_\etale) = \Sh(X_{spaces, \etale})$
，也可把 $\mathcal{O}_X$ 看作 $X_{spaces, \etale}$ 上的环层。
注 \ref{spaces-properties-remark-explain-equivalence} 说明了当 $Y \to X$ 是
$X_{spaces, \etale}$ 的对象时，如何计算 $\mathcal{O}_X(Y)$，
特别是 $\mathcal{O}_X(X)$。

\begin{lemma}
\label{spaces-properties-lemma-morphism-ringed-topoi}
设 $S$ 为概形，$f : X \to Y$ 为 $S$ 上代数空间的态射。
则存在典范映射
$f^\sharp : f_{small}^{-1}\mathcal{O}_Y \to \mathcal{O}_X$，使得
$$
(f_{small}, f^\sharp) :
(\Sh(X_\etale), \mathcal{O}_X)
\longrightarrow
(\Sh(Y_\etale), \mathcal{O}_Y)
$$
为环化拓扑斯态射。此外，
\begin{enumerate}
\item 构造 $f \mapsto (f_{small}, f^\sharp)$ 与复合相容。
\item 若 $f$ 是概形态射，则 $f^\sharp$ 是“下降”注
\ref{descent-remark-change-topologies-ringed} 所述映射。
\end{enumerate}
\end{lemma}

\begin{proof}
由引理 \ref{spaces-properties-lemma-f-map}，只须给出从 $\mathcal{O}_Y$ 到
$\mathcal{O}_X$ 的一个 $f$-映射。换言之，对每个交换图
$$
\xymatrix{
U \ar[d]_g \ar[r] & X \ar[d]^f \\
V \ar[r] & Y
}
$$
其中 $U \in X_\etale$、$V \in Y_\etale$，须给出环同态
$
(f^\sharp)_{(U, V, g)} :
\Gamma(V, \mathcal{O}_V)
\to
\Gamma(U, \mathcal{O}_U).
$
当然，只需取 $(f^\sharp)_{(U, V, g)} = g^\sharp$。这与限制映射相容是清楚的，
故确实给出一个 $f$-映射。略去它与复合相容，以及它与“下降”注
\ref{descent-remark-change-topologies-ringed} 中构造相同的验证。
\end{proof}

\begin{lemma}
\label{spaces-properties-lemma-reduced-space}
设 $S$ 为概形，$X$ 为 $S$ 上的代数空间。下列条件等价：
\begin{enumerate}
\item $X$ 既约；
\item 对每个 $x \in |X|$，$X$ 在 $x$ 处的局部环既约
（注 \ref{spaces-properties-remark-list-properties-local-ring-local-etale-topology}）。
\end{enumerate}
在此情形，$\Gamma(X, \mathcal{O}_X)$ 是既约环；而且若
$f \in \Gamma(X, \mathcal{O}_X)$ 满足 $X = V(f)$，则 $f = 0$。
\end{lemma}

\begin{proof}
对 $X$ 上 \'etale 的仿射概形应用“性质”引理
\ref{properties-lemma-characterize-reduced}，即得 (1) 与 (2) 等价。
最后两个陈述由所引引理以及如下事实推出：对某个在 $X$ 上 \'etale 的既约概形
$U$，$\Gamma(X, \mathcal{O}_X)$ 是 $\Gamma(U, \mathcal{O}_U)$ 的子环。
\end{proof}







\section{结构层的茎}
\label{spaces-properties-section-stalks-structure-sheaf}

\noindent
本节是“\'Etale 上同调”第
\ref{etale-cohomology-section-stalks-structure-sheaf} 节的对应版本。

\begin{lemma}
\label{spaces-properties-lemma-describe-etale-local-ring}
设 $S$ 为概形，$X$ 为 $S$ 上的代数空间，$\overline{x}$ 为 $X$
的几何点，并设 $(U, \overline{u})$ 为 $\overline{x}$ 的 \'etale 邻域，
其中 $U$ 为概形。则
$$
\mathcal{O}_{X, \overline{x}} =
\mathcal{O}_{U, \overline{u}} =
\mathcal{O}_{U, u}^{sh}
$$
其中左端是 $X$ 的结构层之茎，而右端是 $U$ 在点 $u$ 处局部环的严格
Hensel 化；$\overline{u}$ 即以该点为中心。
\end{lemma}

\begin{proof}
已知 $U_\etale$ 上的结构层 $\mathcal{O}_U$ 是 $X$ 的结构层的限制。
故第一个等号由引理 \ref{spaces-properties-lemma-stalk-pullback} 的 (4) 得到。
第二个等号见“\'Etale 上同调”引理
\ref{etale-cohomology-lemma-describe-etale-local-ring}。
\end{proof}

\begin{definition}
\label{spaces-properties-definition-etale-local-rings}
设 $S$ 为概形，$X$ 为 $S$ 上的代数空间，并设 $\overline{x}$
为 $X$ 的几何点，位于点 $x \in |X|$ 上方。
\begin{enumerate}
\item $X$ 在 $\overline{x}$ 处的{\it \'etale 局部环}是
$X_\etale$ 上结构层 $\mathcal{O}_X$ 在 $\overline{x}$ 处的茎。
记作 $\mathcal{O}_{X, \overline{x}}$。
\item $X$ 在 $\overline{x}$ 处的{\it 严格 Hensel 化}是概形
$\Spec(\mathcal{O}_{X, \overline{x}})$。
\end{enumerate}
\end{definition}

\noindent
$X$ 在 $\overline{x}$ 处的严格 Hensel 化作为 $X$ 上概形的同构类型，
仅依赖于点 $x \in |X|$，而与位于 $x$ 上方的几何点之选取无关；
见注 \ref{spaces-properties-remark-map-stalks}。

\begin{lemma}
\label{spaces-properties-lemma-etale-site-locally-ringed}
设 $S$ 为概形，$X$ 为 $S$ 上的代数空间。赋予结构层
$\mathcal{O}_X$ 的小 \'etale 站点 $X_\etale$ 是局部环化站点；
见“站点上的模”定义 \ref{sites-modules-definition-locally-ringed}。
\end{lemma}

\begin{proof}
这是因为各茎 $\mathcal{O}_{X, \overline{x}}$ 都是局部环，并且
$S_\etale$ 有足够多的点；见引理 \ref{spaces-properties-lemma-describe-etale-local-ring}
与定理 \ref{spaces-properties-theorem-exactness-stalks}。这些事实蕴含小 \'etale 站点局部环化，
见“站点上的模”引理 \ref{sites-modules-lemma-locally-ringed-stalk} 与
\ref{sites-modules-lemma-ringed-stalk-not-zero}。
\end{proof}

\begin{lemma}
\label{spaces-properties-lemma-dimension-local-ring}
设 $S$ 为概形，$X$ 为 $S$ 上的代数空间。取点 $x \in |X|$ 及
$d \in \{0, 1, 2, \ldots, \infty\}$。下列条件等价：
\begin{enumerate}
\item $X$ 在 $x$ 处局部环的维数（定义
\ref{spaces-properties-definition-dimension-local-ring}）为 $d$；
\item 对位于 $x$ 上方的某个几何点 $\overline{x}$，有
$\dim(\mathcal{O}_{X, \overline{x}}) = d$；
\item 对位于 $x$ 上方的任意几何点 $\overline{x}$，有
$\dim(\mathcal{O}_{X, \overline{x}}) = d$。
\end{enumerate}
\end{lemma}

\begin{proof}
(2) 与 (3) 的等价性来自如下事实：$\mathcal{O}_{X, \overline{x}}$
的同构类型只依赖于 $x \in |X|$；见注 \ref{spaces-properties-remark-map-stalks}。
利用引理 \ref{spaces-properties-lemma-describe-etale-local-ring}，(1) 与 (2)$+$(3) 的等价性
归结为下述陈述：对任意局部环 $R$，有
$\dim(R) = \dim(R^{sh})$。这正是“更多代数”引理
\ref{more-algebra-lemma-henselization-dimension}。
\end{proof}

\begin{lemma}
\label{spaces-properties-lemma-dimension-decent-invariant-under-etale}
设 $S$ 为概形，$f : X \to Y$ 为 $S$ 上代数空间的 \'etale 态射，
并取 $x \in X$。则 (1) $\dim_x(X) = \dim_{f(x)}(Y)$；
(2) $X$ 在 $x$ 处局部环的维数等于 $Y$ 在 $f(x)$ 处局部环的维数。
若 $f$ 满，则 (3) $\dim(X) = \dim(Y)$。
\end{lemma}

\begin{proof}
选取概形 $U$、点 $u \in U$ 以及把 $u$ 映到 $x$ 的 \'etale 态射
$U \to X$。则复合 $U \to Y$ 也为 \'etale，并把 $u$ 映到 $f(x)$。
相关数值是依照概形 $U$ 在 $u$ 处的性态定义的，故 (1) 与 (2) 成立。
关于 (1) 见定义 \ref{spaces-properties-definition-dimension-at-point}。(3) 是 (1) 的直接推论，
见定义 \ref{spaces-properties-definition-dimension}。
\end{proof}

\begin{lemma}
\label{spaces-properties-lemma-reduced-local-ring}
设 $S$ 为概形，$X$ 为 $S$ 上的代数空间，并取点 $x \in |X|$。
下列条件等价：
\begin{enumerate}
\item $X$ 在 $x$ 处的局部环既约
（注 \ref{spaces-properties-remark-list-properties-local-ring-local-etale-topology}）；
\item 对位于 $x$ 上方的某个几何点 $\overline{x}$，
$\mathcal{O}_{X, \overline{x}}$ 既约；
\item 对位于 $x$ 上方的任意几何点 $\overline{x}$，
$\mathcal{O}_{X, \overline{x}}$ 既约。
\end{enumerate}
\end{lemma}

\begin{proof}
(2) 与 (3) 的等价性来自如下事实：$\mathcal{O}_{X, \overline{x}}$
的同构类型只依赖于 $x \in |X|$；见注 \ref{spaces-properties-remark-map-stalks}。
利用引理 \ref{spaces-properties-lemma-describe-etale-local-ring}，(1) 与 (2)$+$(3) 的等价性
归结为：一个局部环既约，当且仅当它的严格 Hensel 化既约。
这正是“更多代数”引理 \ref{more-algebra-lemma-henselization-reduced}。
\end{proof}











\section{局部不可约性}
\label{spaces-properties-section-irreducible-local-ring}

\noindent
代数空间上的一个点具有良定义的 \'etale 局部环；在概形的情形，
它对应于局部环的严格 Hensel 化。一般而言，仅由 \'etale 局部环无法看出
概形或代数空间有多少个不可约分支通过给定点；我们只能计算几何分支的数目。

\begin{lemma}
\label{spaces-properties-lemma-irreducible-local-ring}
设 $S$ 为概形，$X$ 为 $S$ 上的代数空间，并取点 $x \in |X|$。
下列条件等价：
\begin{enumerate}
\item 对任意概形 $U$、\'etale 态射 $a : U \to X$ 及满足
$a(u) = x$ 的 $u \in U$，局部环 $\mathcal{O}_{U, u}$ 有唯一极小素理想；
\item 对任意概形 $U$、\'etale 态射 $a : U \to X$ 及满足
$a(u) = x$ 的 $u \in U$，恰有一个 $U$ 的不可约分支通过 $u$；
\item 对任意概形 $U$、\'etale 态射 $a : U \to X$ 及满足
$a(u) = x$ 的 $u \in U$，局部环 $\mathcal{O}_{U, u}$ 单枝；
\item 对任意概形 $U$、\'etale 态射 $a : U \to X$ 及满足
$a(u) = x$ 的 $u \in U$，局部环 $\mathcal{O}_{U, u}$ 几何单枝；
\item 对位于 $x$ 上方的任意几何点 $\overline{x}$，
$\mathcal{O}_{X, \overline{x}}$ 有唯一极小素理想。
\end{enumerate}
\end{lemma}

\begin{proof}
(1) 与 (2) 等价，因为通过 $u$ 的 $U$ 的不可约分支与 $U$ 在 $u$ 处局部环的
极小素理想成 $1$-$1$ 对应。取 (1) 中的 $a : U \to X$ 与 $u \in U$。
由引理 \ref{spaces-properties-lemma-describe-etale-local-ring}，$\mathcal{O}_{X, \overline{x}}$
是 $\mathcal{O}_{U, u}$ 的严格 Hensel 化。特别地，由“更多代数”引理
\ref{more-algebra-lemma-geometrically-unibranch}，(4) 与 (5) 等价。
由“更多态射”引理 \ref{more-morphisms-lemma-nr-branches}，
(2)、(3)、(4) 等价。
\end{proof}

\begin{definition}
\label{spaces-properties-definition-unibranch}
设 $S$ 为概形，$X$ 为 $S$ 上的代数空间，并取 $x \in |X|$。
若引理 \ref{spaces-properties-lemma-irreducible-local-ring} 中的等价条件成立，
则称 $X$ 在 $x$ 处{\it 几何单枝}。若 $X$ 在每个
$x \in |X|$ 处都几何单枝，则称 $X${\it 几何单枝}。
\end{definition}

\noindent
这与概形的定义相容（“性质”定义 \ref{properties-definition-unibranch}）。

\begin{lemma}
\label{spaces-properties-lemma-nr-branches-local-ring}
设 $S$ 为概形，$X$ 为 $S$ 上的代数空间，取点 $x \in |X|$，
并取整数 $n \in \{1, 2, \ldots\}$。下列条件等价：
\begin{enumerate}
\item 对任意概形 $U$、\'etale 态射 $a : U \to X$ 及满足
$a(u) = x$ 的 $u \in U$，局部环 $\mathcal{O}_{U, u}$ 的极小素理想数
为 $\leq n$，且至少对一组选取 $U, a, u$，该数为 $n$；
\item 对任意概形 $U$、\'etale 态射 $a : U \to X$ 及满足
$a(u) = x$ 的 $u \in U$，通过 $u$ 的 $U$ 的不可约分支数为
$\leq n$，且至少对一组选取 $U, a, u$，该数为 $n$；
\item 对任意概形 $U$、\'etale 态射 $a : U \to X$ 及满足
$a(u) = x$ 的 $u \in U$，$U$ 在 $u$ 处的分支数为 $\leq n$，
且至少对一组选取 $U, a, u$，该数为 $n$；
\item 对任意概形 $U$、\'etale 态射 $a : U \to X$ 及满足
$a(u) = x$ 的 $u \in U$，$U$ 在 $u$ 处的几何分支数为 $n$；
\item $\mathcal{O}_{X, \overline{x}}$ 的极小素理想数为 $n$。
\end{enumerate}
\end{lemma}

\begin{proof}
(1) 与 (2) 等价，因为通过 $u$ 的 $U$ 的不可约分支与 $U$ 在 $u$ 处局部环的
极小素理想成 $1$-$1$ 对应。取 (1) 中的 $a : U \to X$ 与 $u \in U$。
由引理 \ref{spaces-properties-lemma-describe-etale-local-ring}，$\mathcal{O}_{X, \overline{x}}$
是 $\mathcal{O}_{U, u}$ 的严格 Hensel 化。回忆，$U$ 在 $u$ 处的
（几何）分支数，就是 $\mathcal{O}_{U, u}$ 的（严格）Hensel 化之极小素理想数。
特别地，(4) 与 (5) 等价。由“更多态射”引理
\ref{more-morphisms-lemma-nr-branches}，(2)、(3)、(4) 等价。
\end{proof}

\begin{definition}
\label{spaces-properties-definition-number-of-geometric-branches}
设 $S$ 为概形，$X$ 为 $S$ 上的代数空间，并取 $x \in |X|$。
若引理 \ref{spaces-properties-lemma-nr-branches-local-ring} 的等价条件成立，则
{\it $X$ 在 $x$ 处的几何分支数}为 $n \in \mathbf{N}$；否则为 $\infty$。
\end{definition}








\section{Noether 代数空间}
\label{spaces-properties-section-noetherian}

\noindent
我们已在第 \ref{spaces-properties-section-types-properties} 节定义了局部 Noether 代数空间。

\begin{definition}
\label{spaces-properties-definition-noetherian}
设 $S$ 为概形，$X$ 为 $S$ 上的代数空间。若 $X$ 拟紧、拟分离且局部
Noether，则称 $X$ 是{\it Noether}的。
\end{definition}

\noindent
注意，Noether 代数空间 $X$ 不仅拟紧且局部 Noether，还须拟分离。
这与 Noether 概形的定义并不冲突，因为局部 Noether 概形必拟分离；
见“性质”引理 \ref{properties-lemma-locally-Noetherian-quasi-separated}。
对代数空间则未必如此：$X = \mathbf{A}^1_k/\mathbf{Z}$（见“空间”例
\ref{spaces-example-affine-line-translation}）局部 Noether 且拟紧，
但不拟分离（因而按我们的定义不是 Noether 的）。

\medskip\noindent
上述选择的一个后果是：Noether 代数空间上的有限型代数空间未必自动 Noether；
也就是说，“态射”引理 \ref{morphisms-lemma-finite-type-noetherian}
的类似命题不成立。正确陈述是，Noether 代数空间上的有限表示代数空间是
Noether 的（见“空间的态射”引理
\ref{spaces-morphisms-lemma-finite-presentation-noetherian}）。

\medskip\noindent
Noether 代数空间 $X$ 非常接近概形。本节余下部分收集若干引理来说明这一点。

\begin{lemma}
\label{spaces-properties-lemma-Noetherian-topology}
设 $S$ 为概形，$X$ 为 $S$ 上的代数空间。
\begin{enumerate}
\item 若 $X$ 局部 Noether，则 $|X|$ 是局部 Noether 拓扑空间。
\item 若 $X$ 拟紧且局部 Noether，则 $|X|$ 是 Noether 拓扑空间。
\end{enumerate}
\end{lemma}

\begin{proof}
设 $X$ 局部 Noether。选取概形 $U$ 及满的 \'etale 态射 $U \to X$。
由 $X$ 局部 Noether 可知 $U$ 局部 Noether。根据“性质”引理
\ref{properties-lemma-Noetherian-topology}，这意味着 $|U|$ 是局部
Noether 拓扑空间。由于 $|U| \to |X|$ 是开满射，“拓扑”引理
\ref{topology-lemma-image-Noetherian} 表明 $|X|$ 局部 Noether。
这证明了 (1)。若 $X$ 拟紧且局部 Noether，则 $|X|$ 拟紧且局部
Noether，故由“拓扑”引理
\ref{topology-lemma-quasi-compact-locally-Noetherian-Noetherian}，
$|X|$ 是 Noether 的。
\end{proof}

\begin{lemma}
\label{spaces-properties-lemma-Noetherian-sober}
设 $S$ 为概形，$X$ 为 $S$ 上的代数空间。若 $X$ Noether，
则 $|X|$ 是清醒 Noether 拓扑空间。
\end{lemma}

\begin{proof}
拟分离代数空间的底拓扑空间是清醒的，见引理
\ref{spaces-properties-lemma-quasi-separated-sober}；它由引理
\ref{spaces-properties-lemma-Noetherian-topology} 又是 Noether 的。
\end{proof}

\begin{lemma}
\label{spaces-properties-lemma-Noetherian-local-ring-Noetherian}
设 $S$ 为概形，$X$ 为 $S$ 上的 Noether 代数空间，并设
$\overline{x}$ 为 $X$ 的几何点。则 $\mathcal{O}_{X, \overline{x}}$
是 Noether 局部环。
\end{lemma}

\begin{proof}
选取 $\overline{x}$ 的 \'etale 邻域 $(U, \overline{u})$，其中 $U$ 为概形。
由引理 \ref{spaces-properties-lemma-describe-etale-local-ring}，
$\mathcal{O}_{X, \overline{x}}$ 是 $U$ 在 $u$ 处局部环的严格 Hensel 化。
按我们对 Noether 空间的定义，概形 $U$ 局部 Noether。故结论由“更多代数”引理
\ref{more-algebra-lemma-henselization-noetherian} 得到。
\end{proof}








\section{正则代数空间}
\label{spaces-properties-section-regular}

\noindent
我们已在第 \ref{spaces-properties-section-types-properties} 节定义了正则代数空间。

\begin{lemma}
\label{spaces-properties-lemma-regular}
设 $S$ 为概形，$X$ 为 $S$ 上局部 Noether 的代数空间。下列条件等价：
\begin{enumerate}
\item $X$ 正则；
\item 每个 \'etale 局部环 $\mathcal{O}_{X, \overline{x}}$ 都正则。
\end{enumerate}
\end{lemma}

\begin{proof}
设 $U$ 为概形，且 $U \to X$ 为满的 \'etale 态射。依假设，$U$ 局部
Noether。此外，每个 \'etale 局部环 $\mathcal{O}_{X, \overline{x}}$
都是 $U$ 上某个局部环的严格 Hensel 化，反之亦然；见引理
\ref{spaces-properties-lemma-describe-etale-local-ring}。故由“更多代数”引理
\ref{more-algebra-lemma-henselization-regular}，(2) 等价于 $U$ 的每个局部环
都正则，即 $U$ 为正则概形（见“性质”引理
\ref{properties-lemma-characterize-regular}）。由定义
\ref{spaces-properties-definition-type-property}，这又等价于 (1)。
\end{proof}

\noindent
可利用“下降”引理 \ref{descent-lemma-regular-local-ring-local}
定义代数空间 $X$ 在点 $x$ 处正则的含义。

\begin{definition}
\label{spaces-properties-definition-regular-at-point}
设 $S$ 为概形，$X$ 为 $S$ 上的代数空间，并取点 $x \in |X|$。
若对任意（等价地，某个）二元组 $(a : U \to X, u)$，
$\mathcal{O}_{U, u}$ 都是正则局部环，则称{\it $X$ 在 $x$ 处正则}；
这里 $a : U \to X$ 是从概形到 $X$ 的 \'etale 态射，且点
$u \in U$ 满足 $a(u) = x$。
\end{definition}

\noindent
见定义 \ref{spaces-properties-definition-property-at-point}、引理
\ref{spaces-properties-lemma-local-source-target-at-point} 以及“下降”引理
\ref{descent-lemma-regular-local-ring-local}。

\begin{lemma}
\label{spaces-properties-lemma-regular-at-x}
设 $S$ 为概形，$X$ 为 $S$ 上的代数空间，并取点 $x \in |X|$。
下列条件等价：
\begin{enumerate}
\item $X$ 在 $x$ 处正则；
\item 对位于 $x$ 上方的任意（等价地，某个）几何点 $\overline{x}$，
\'etale 局部环 $\mathcal{O}_{X, \overline{x}}$ 正则。
\end{enumerate}
\end{lemma}

\begin{proof}
设 $U$ 为概形，$u \in U$ 为点，并设 $a : U \to X$ 为把 $u$ 映到
$x$ 的 \'etale 态射。对 $X$ 中位于 $x$ 上方的任意几何点
$\overline{x}$，\'etale 局部环 $\mathcal{O}_{X, \overline{x}}$
是 $U$ 在 $u$ 处某个局部环的严格 Hensel 化，见引理
\ref{spaces-properties-lemma-describe-etale-local-ring}。故结论由“更多代数”引理
\ref{more-algebra-lemma-henselization-regular} 得到。
\end{proof}

\begin{lemma}
\label{spaces-properties-lemma-regular-normal}
正则代数空间是正规代数空间。
\end{lemma}

\begin{proof}
这由定义及概形的情形推出；见“性质”引理
\ref{properties-lemma-regular-normal}。
\end{proof}

 








\section{代数空间上的模层}
\label{spaces-properties-section-modules}

\noindent
若 $X$ 为代数空间，则 $X$ 上的模层是 $X$ 的小 \'etale 站点上的
$\mathcal{O}_X$-模层，其中 $\mathcal{O}_X$ 为 $X$ 的结构层。
模层范畴记作 $\textit{Mod}(\mathcal{O}_X)$。

\medskip\noindent
给定代数空间的态射 $f : X \to Y$，由引理
\ref{spaces-properties-lemma-morphism-ringed-topoi} 得到环化拓扑斯态射，进而由“站点上的模”定义
\ref{sites-modules-definition-pushforward} 得到良定义的拉回函子与直像函子
\begin{equation}
\label{spaces-properties-equation-push-pull}
f^* :
\textit{Mod}(\mathcal{O}_Y)
\longrightarrow
\textit{Mod}(\mathcal{O}_X), \quad
f_* :
\textit{Mod}(\mathcal{O}_X)
\longrightarrow
\textit{Mod}(\mathcal{O}_Y)
\end{equation}
，它们按通常方式互为伴随。若 $g : Y \to Z$ 是 $S$ 上代数空间的另一态射，
则
$(g \circ f)^* = f^* \circ g^*$，且 $(g \circ f)_* = g_* \circ f_*$
，这只因环化拓扑斯态射按相应方式复合（由上述引理）。

\begin{lemma}
\label{spaces-properties-lemma-etale-exact-pullback}
设 $S$ 为概形，$f : X \to Y$ 为 $S$ 上代数空间的 \'etale 态射。
则 $f^{-1}\mathcal{O}_Y = \mathcal{O}_X$，并且对任意
$\mathcal{O}_Y$-模层 $\mathcal{G}$，有
$f^*\mathcal{G} = f_{small}^{-1}\mathcal{G}$。特别地，
$f^* : \textit{Mod}(\mathcal{O}_Y) \to \textit{Mod}(\mathcal{O}_X)$
正合。
\end{lemma}

\begin{proof}
由引理 \ref{spaces-properties-lemma-etale-morphism-topoi} 中对逆像的描述及结构层的定义，
$f_{small}^{-1}\mathcal{O}_Y = \mathcal{O}_X$ 是清楚的。由于拉回按定义为
$$
f^*\mathcal{G} =
f_{small}^{-1}\mathcal{G} \otimes_{f_{small}^{-1}\mathcal{O}_Y}
\mathcal{O}_X
$$
，可得 $f^*\mathcal{G} = f_{small}^{-1}\mathcal{G}$。正合性也是清楚的，
因为 $f_{small}$ 是拓扑斯态射，所以 $f_{small}^{-1}$ 正合。
\end{proof}

\noindent
继续沿用式 (\ref{spaces-properties-equation-restrict}) 中引入的记号滥用，在上述引理的情形写成
\begin{equation}
\label{spaces-properties-equation-restrict-modules}
\mathcal{G}|_{X_\etale}
= f^*\mathcal{G}
= f_{small}^{-1}\mathcal{G}
\end{equation}
。第 \ref{spaces-properties-section-localize} 节将以更技术性的方式讨论此事。

\begin{lemma}
\label{spaces-properties-lemma-pushforward-etale-base-change-modules}
设 $S$ 为概形，并设
$$
\xymatrix{
X' \ar[r] \ar[d]_{f'} & X \ar[d]^f \\
Y' \ar[r]^g & Y
}
$$
为 $S$ 上代数空间的笛卡儿方块。取
$\mathcal{F} \in \textit{Mod}(\mathcal{O}_X)$。若 $g$ 是 \'etale 的，则
$f'_*(\mathcal{F}|_{X'}) = (f_*\mathcal{F})|_{Y'}$\footnote{还有
$(f')^*(\mathcal{G}|_{Y'}) = (f^*\mathcal{G})|_{X'}$
，这是因为该图交换且有式 (\ref{spaces-properties-equation-restrict-modules})}，并且
$R^if'_*(\mathcal{F}|_{X'}) = (R^if_*\mathcal{F})|_{Y'}$ 在
$\textit{Mod}(\mathcal{O}_{Y'})$ 中成立。
\end{lemma}

\begin{proof}
这是引理 \ref{spaces-properties-lemma-pushforward-etale-base-change} 在模情形的重新表述。
\end{proof}

\begin{lemma}
\label{spaces-properties-lemma-characterize-module-small-etale}
设 $S$ 为概形，$X$ 为 $S$ 上的代数空间。一个 $\mathcal{O}_X$-模层
$\mathcal{F}$ 由下列数据给出：
\begin{enumerate}
\item 对每个 $U \in \Ob(X_\etale)$，给定 $U_\etale$ 上的
$\mathcal{O}_U$-模层 $\mathcal{F}_U$；
\item 对 $X_\etale$ 中每个 $f : U' \to U$，给定同构
$c_f : f_{small}^*\mathcal{F}_U \to \mathcal{F}_{U'}$.
\end{enumerate}
这些数据须满足如下条件：对 $X_\etale$ 中任意 $f : U' \to U$ 与
$g : U'' \to U'$，复合 $c_g \circ g_{small}^*c_f$ 等于
$c_{f \circ g}$。
\end{lemma}

\begin{proof}
合用引理 \ref{spaces-properties-lemma-etale-exact-pullback} 与
\ref{spaces-properties-lemma-characterize-sheaf-small-etale}，并使用如下事实：
$X_\etale$ 的对象之间任意态射都是概形的 \'etale 态射。
\end{proof}







\section{\'Etale 局部化}
\label{spaces-properties-section-localize}

\noindent
应不惜一切代价避免阅读本节。

\medskip\noindent
设 $X \to Y$ 为代数空间的 \'etale 态射。则 $X$ 是
$Y_{spaces, \etale}$ 的对象；由定义立即可得（也见引理
\ref{spaces-properties-lemma-etale-morphism-topoi} 的证明）
\begin{equation}
\label{spaces-properties-equation-localize}
X_{spaces, \etale} = Y_{spaces, \etale}/X
\end{equation}
其中右端是站点 $Y_{spaces, \etale}$ 在对象 $X$ 处的局部化，
见“站点”定义 \ref{sites-definition-localize}。此外，由引理
\ref{spaces-properties-lemma-etale-exact-pullback}，此等同与结构层相容。
因此环化站点 $(X_{spaces, \etale}, \mathcal{O}_X)$ 等同于环化站点
$(Y_{spaces, \etale}, \mathcal{O}_Y)$ 在对象 $X$ 处的局部化：
\begin{equation}
\label{spaces-properties-equation-localize-ringed}
(X_{spaces, \etale}, \mathcal{O}_X) =
(Y_{spaces, \etale}/X, \mathcal{O}_Y|_{Y_{spaces, \etale}/X})
\end{equation}
右端所用的环化站点局部化定义于“站点上的模”定义
\ref{sites-modules-definition-localize-ringed-site}。

\medskip\noindent
现在设 $X \to Y$ 是代数空间的 \'etale 态射，且 $X$ 为概形。
则 $X$ 是 $Y_\etale$ 的对象，从而
\begin{equation}
\label{spaces-properties-equation-localize-at-scheme}
X_\etale = Y_\etale/X
\end{equation}
以及
\begin{equation}
\label{spaces-properties-equation-localize-at-scheme-ringed}
(X_\etale, \mathcal{O}_X) =
(Y_\etale/X, \mathcal{O}_Y|_{Y_\etale/X})
\end{equation}
，与上面相同。

\medskip\noindent
最后，若 $X \to Y$ 是代数空间的 \'etale 态射，且 $X$ 为仿射概形，
则 $X$ 是 $Y_{affine, \etale}$ 的对象，并有
\begin{equation}
\label{spaces-properties-equation-localize-at-affine}
X_{affine, \etale} = Y_{affine, \etale}/X
\end{equation}
以及
\begin{equation}
\label{spaces-properties-equation-localize-at-affine-ringed}
(X_{affine, \etale}, \mathcal{O}_X) =
(Y_{affine, \etale}/X, \mathcal{O}_Y|_{Y_{affine, \etale}/X})
\end{equation}
，与上面相同。

\medskip\noindent
下面证明这些局部化与态射相容。

\begin{lemma}
\label{spaces-properties-lemma-relocalize-morphism}
设 $S$ 为概形，并设
$$
\xymatrix{
U \ar[d]_p \ar[r]_g & V \ar[d]^q \\
X \ar[r]^f & Y
}
$$
为 $S$ 上代数空间的交换图，其中 $p$ 与 $q$ 都是 \'etale 的。
通过式 (\ref{spaces-properties-equation-localize-ringed}) 分别对 $U \to X$ 与
$V \to Y$ 给出的等同，环化拓扑斯态射
$$
(g_{spaces, \etale}, g^\sharp) :
(\Sh(U_{spaces, \etale}), \mathcal{O}_U)
\longrightarrow
(\Sh(V_{spaces, \etale}), \mathcal{O}_V)
$$
与下述态射 $2$-同构：从环化站点态射
$(f_{spaces, \etale}, f^\sharp)$ 以及对应于 $g$ 的映射
$c : U \to V \times_Y X$ 出发，按“站点上的模”引理
\ref{sites-modules-lemma-relocalize-morphism-ringed-sites} 构造的态射
$(f_{spaces, \etale, c}, f_c^\sharp)$。
\end{lemma}

\begin{proof}
态射 $(f_{spaces, \etale, c}, f_c^\sharp)$ 定义为局部化与基变换映射的复合
$f' \circ j$。类似地，$g$ 是复合 $U \to V \times_Y X \to V$。
故只须在下列两种情形证明本引理：(1) $f = \text{id}$；
(2) $U = X \times_Y V$。在情形 (1)，态射 $g : U \to V$ 是 \'etale 的，
见引理 \ref{spaces-properties-lemma-etale-permanence}。于是由式
(\ref{spaces-properties-equation-localize}) 与 (\ref{spaces-properties-equation-localize-ringed}) 周围的讨论，
$(g_{spaces, \etale}, g^\sharp)$ 是局部化态射；这恰是本引理在该情形的内容。
在情形 (2)，态射 $g_{spaces, \etale}$ 来自函子
$V_{spaces, \etale} \to U_{spaces, \etale}$,
$V'/V \mapsto V' \times_V U/U$
所给的环化站点态射，而态射 $f'$ 也是由此定义的；见“站点”引理
\ref{sites-lemma-localize-morphism}。略去在此情形验证
$(f')^\sharp = g^\sharp$（二者都是 $f^\sharp$ 到
$U_{spaces, \etale}$ 的限制）。
\end{proof}

\begin{lemma}
\label{spaces-properties-lemma-relocalize-morphism-at-schemes}
采用与引理 \ref{spaces-properties-lemma-relocalize-morphism} 相同的记号及假设，
但另假设 $U$ 与 $V$ 都是概形。通过式
(\ref{spaces-properties-equation-localize-at-scheme-ringed}) 分别对 $U \to X$ 与
$V \to Y$ 给出的等同，环化拓扑斯态射
$$
(g_{small}, g^\sharp) :
(\Sh(U_\etale), \mathcal{O}_U)
\longrightarrow
(\Sh(V_\etale), \mathcal{O}_V)
$$
与下述态射 $2$-同构：从 $(f_{small}, f^\sharp)$ 及对应于 $g$ 的映射
$s : h_U \to f_{small}^{-1}h_V$ 出发，按“站点上的模”引理
\ref{sites-modules-lemma-relocalize-morphism-ringed-topoi} 构造的态射
$(f_{small, s}, f_s^\sharp)$。
\end{lemma}

\begin{proof}
注意，作为环化拓扑斯态射，$(g_{small}, g^\sharp)$ 与相伴于环化站点态射
$(g_{spaces, \etale}, g^\sharp)$ 的环化拓扑斯态射 $2$-同构。
故结论由引理 \ref{spaces-properties-lemma-relocalize-morphism} 及“站点上的模”引理
\ref{sites-modules-lemma-relocalize-morphism-compare} 得到。
\end{proof}

\noindent
最后讨论代数空间 $Y$ 的小 \'etale 站点 $Y_\etale$ 上的集合层，
与 $Y$ 上 \'etale 的代数空间之间的关系。设 $S$ 为概形，$Y$ 为
$S$ 上的代数空间，且 $\mathcal{F}$ 为 $\Sh(Y_\etale)$ 的对象。
考虑函子
$$
X : (\Sch/S)_{fppf}^{opp} \longrightarrow \textit{Sets}
$$
，其定义规则为
$$
X(T) = \{(y, s) \mid y : T \to Y\text{ 是 }S\text{ 上的态射且 }
s \in \Gamma(T, y_{small}^{-1}\mathcal{F})\}
$$
给定态射 $g : T' \to T$，限制映射把 $(y, s)$ 送到
$(y \circ g, g_{small}^{-1}s)$。这是有意义的，因为由引理
\ref{spaces-properties-lemma-functoriality-etale-site}，
$y_{small} \circ g_{small} = (y \circ g)_{small}$。
存在典范映射 $X \to Y$，它把二元组 $(y, s)$ 送到 $y$。

\begin{lemma}
\label{spaces-properties-lemma-sheaf-gives-space}
设 $S$ 为概形，$Y$ 为 $S$ 上的代数空间，并设 $\mathcal{F}$
为 $Y_\etale$ 上的集合层。只要满足一个集合论条件（见证明），就有：
\begin{enumerate}
\item 上述与 $\mathcal{F}$ 相伴的函子 $X$ 是代数空间；
\item 映射 $X \to Y$ 是代数空间的 \'etale 态射；
\item 通过等同 $\Sh(Y_\etale) = \Sh(Y_{spaces, \etale})$，有
$\mathcal{F} \cong h_X$；
\item 有 $\mathcal{F} \cong f_{small, !}*$。这里 $*$ 是范畴
$\Sh(X_\etale)$ 的终对象，而 $f_{small, !}$ 由引理
\ref{spaces-properties-lemma-etale-morphism-topoi} 存在。
\end{enumerate}
\end{lemma}

\begin{proof}
先证明 $X$ 是 fppf 拓扑的层。设 $\{g_i : T_i \to T\}$ 是
$(\Sch/S)_{fppf}$ 的覆盖，且 $(y_i, s_i) \in X(T_i)$ 满足胶合条件，
即 $(y_i, s_i)$ 与 $(y_j, s_j)$ 到 $T_i \times_T T_j$ 的限制相同。
由于 $Y$ 是 fppf 拓扑的层，各 $y_i$ 胶合为唯一态射
$y : T \to Y$，使得 $y_i = y \circ g_i$。于是
$y_{i, small}^{-1}\mathcal{F} = g_{i, small}^{-1}y_{small}^{-1}\mathcal{F}$.
故由“\'Etale 上同调”引理
\ref{etale-cohomology-lemma-describe-pullback}，各截面 $s_i$
唯一胶合为 $y_{small}^{-1}\mathcal{F}$ 的一个截面。

\medskip\noindent
把 $\mathcal{F} \in \Ob(\Sh(Y_\etale))$ 送到
$X \in \Ob((\Sch/S)_{fppf})$ 的构造保持有限极限及所有余极限，
因为每个函子 $y_{small}^{-1}$ 都有此性质。当然，若
$V \in \Ob(Y_\etale)$，则该构造把 $Y_\etale$ 上的可表层 $h_V$
送到由 $V$ 表示的可表函子。

\medskip\noindent
由“站点”引理 \ref{sites-lemma-sheaf-coequalizer-representable}，
可找到集合 $I$、对每个 $i \in I$ 找到 $Y_\etale$ 的对象 $V_i$，
以及层的满射
$$
\coprod h_{V_i} \longrightarrow \mathcal{F}
$$
于 $Y_\etale$ 上。我们需要的集合论条件是指标集 $I$ 不太大\footnote{只须
$\mathcal{F}$ 在 $Y$ 的各几何点处之茎的基数上确界，不超过
$(\Sch/S)_{fppf}$ 某个对象的大小。}。于是
$V = \coprod V_i$ 是 $(\Sch/S)_{fppf}$ 的对象，从而也是
$Y_\etale$ 的对象，并且有满射 $h_V \to \mathcal{F}$。

\medskip\noindent
注意，$h_V$ 在 $\Sh(Y_\etale)$ 中与自身的积是
$h_{V \times_Y V}$。考虑纤维积
$$
h_V \times_\mathcal{F} h_V \subset h_{V \times_Y V}
$$
存在 $V \times_Y V$ 的开子概形 $R$，使得
$h_V \times_\mathcal{F} h_V = h_R$；见引理
\ref{spaces-properties-lemma-support-subsheaf-final}（略去一个小细节）。由 Yoneda 引理，
得到 $Y_\etale$ 中两个态射 $s, t : R \to V$，并得到余等化子图
$$
\xymatrix{
h_R \ar@<1ex>[r] \ar@<-1ex>[r] &
h_V \ar[r] &
\mathcal{F}
}
$$
于 $\Sh(Y_\etale)$ 中。当然，态射 $s, t$ 都是 \'etale 的，
并定义一个 \'etale 等价关系 $(t, s) : R \to V \times_S V$。

\medskip\noindent
由前两段的讨论，得到余等化子图
$$
\xymatrix{
R \ar@<1ex>[r] \ar@<-1ex>[r] &
V \ar[r] &
X
}
$$
于 $(\Sch/S)_{fppf}$ 中。因此，由“空间”定理
\ref{spaces-theorem-presentation}，$X = V/R$ 是代数空间。这证明了 (1)。
(2) 来自 $V \to Y$ 是 \'etale 的。(3) 由 $X$ 与 $h_X$ 的定义立即可得。
略去 (4) 的证明；它可由下述比较得到：一边是与引理
\ref{spaces-properties-lemma-etale-morphism-topoi} 的余连续函子 $j$ 相伴的态射，
另一边是上面对 $X_{spaces, \etale}$ 作为 $Y_{spaces, \etale}$
在 $X$ 处之局部化的描述以及“站点”引理
\ref{sites-lemma-describe-j-shriek-representable}。
\end{proof}







\section{恢复态射}
\label{spaces-properties-section-morphisms}

\noindent
本节证明，把一个代数空间与其局部环化的小 \'etale 拓扑斯对应起来的规则，
在适当意义下是全忠实的，参见
定理 \ref{spaces-properties-theorem-fully-faithful}。

\begin{lemma}
\label{spaces-properties-lemma-morphism-locally-ringed}
设 $S$ 为概形。
设 $f : X \to Y$ 为 $S$ 上代数空间的态射。
与 $f$ 关联的环化拓扑斯态射 $(f_{small}, f^\sharp)$
是局部环化拓扑斯的态射，参见
《站点上的模》，
定义 \ref{sites-modules-definition-morphism-locally-ringed-topoi}。
\end{lemma}

\begin{proof}
注意该断言是有意义的，因为我们已经看到
$(X_\etale, \mathcal{O}_{X_\etale})$ 以及
$(Y_\etale, \mathcal{O}_{Y_\etale})$
均为局部环化站点，参见
引理 \ref{spaces-properties-lemma-etale-site-locally-ringed}。
此外，我们知道 $X_\etale$ 有足够多的点，参见
定理 \ref{spaces-properties-theorem-exactness-stalks}。
因此只需证明 $(f_{small}, f^\sharp)$
满足《站点上的模》
引理 \ref{sites-modules-lemma-locally-ringed-morphism}
的条件 (3)。为此，取 $X_\etale$ 的一个点 $p$。由
引理 \ref{spaces-properties-lemma-points-small-etale-site}，
$p$ 对应于 $X$ 的一个几何点 $\overline{x}$。
由
引理 \ref{spaces-properties-lemma-stalk-pullback}，
点 $q = f_{small} \circ p$ 对应于 $Y$ 的几何点
$\overline{y} = f \circ \overline{x}$。
故我们需要证明的断言是，所诱导的 \'etale 局部环同态
$$
\mathcal{O}_{Y, \overline{y}} \longrightarrow \mathcal{O}_{X, \overline{x}}
$$
是局部环同态。可以直接证明这一点；这里我们改由概形的相应结果推出。
为此选择交换图
$$
\xymatrix{
U \ar[d] \ar[r]_\psi & V \ar[d] \\
X \ar[r] & Y
}
$$
其中 $U$ 和 $V$ 为概形，竖直箭头为满的 \'etale 态射（参见
《空间》，引理 \ref{spaces-lemma-lift-morphism-presentations}）。
选择一个提升 $\overline{u} : \overline{x} \to U$（由
引理 \ref{spaces-properties-lemma-geometric-lift-to-cover} 可知这是可能的）。
令 $\overline{v} = \psi \circ \overline{u}$。于是得到 \'etale 局部环的交换图
$$
\xymatrix{
\mathcal{O}_{U, \overline{u}} &
\mathcal{O}_{V, \overline{v}} \ar[l] \\
\mathcal{O}_{X, \overline{x}} \ar[u] &
\mathcal{O}_{Y, \overline{y}}. \ar[l] \ar[u]
}
$$
由
《\'Etale 上同调》，引理 \ref{etale-cohomology-lemma-morphism-locally-ringed}，
上方的水平箭头是局部环同态。最后，由
引理 \ref{spaces-properties-lemma-describe-etale-local-ring}，
竖直箭头均为同构。结论随即成立。
\end{proof}

\begin{lemma}
\label{spaces-properties-lemma-2-morphism}
设 $S$ 为概形。
设 $X$、$Y$ 为 $S$ 上的代数空间。
设 $f : X \to Y$ 为 $S$ 上代数空间的态射。
设 $t$ 是从 $(f_{small}, f^\sharp)$ 到其自身的一个 $2$-态射，参见
《站点上的模》，
定义 \ref{sites-modules-definition-2-morphism-ringed-topoi}。
则 $t = \text{id}$。
\end{lemma}

\begin{proof}
令 $X'$（分别地，$Y'$）表示按同样方式把 $X$ 视为
$\Spec(\mathbf{Z})$ 上的代数空间，参见
《空间》，定义 \ref{spaces-definition-base-change}。
由构造可知 $(X_{small}, \mathcal{O})$
等于 $(X'_{small}, \mathcal{O})$，对 $Y$ 也同样成立。
因此可以处理 $X'$ 和 $Y'$。换言之，我们可以假设
$S = \Spec(\mathbf{Z})$。

\medskip\noindent
假设 $S = \Spec(\mathbf{Z})$，且 $f : X \to Y$ 和 $t$ 如引理中所述。
这意味着 $t : f^{-1}_{small} \to f^{-1}_{small}$
是函子的一个变换，并且下图交换：
$$
\xymatrix{
f_{small}^{-1}\mathcal{O}_Y
\ar[rd]_{f^\sharp}  & &
f_{small}^{-1}\mathcal{O}_Y \ar[ll]^t \ar[ld]^{f^\sharp} \\
& \mathcal{O}_X
}
$$
设 $V \to Y$ 为 \'etale 态射且 $V$ 为仿射概形。
写作 $V = \Spec(B)$。选取 $b_j \in B$（$j \in J$）作为
$B$ 这个 $\mathbf{Z}$-代数的一组生成元。令
$T = \Spec(\mathbf{Z}[\{x_j\}_{j \in J}])$.
下文中我们将使用
$\Mor_{\Sch}(U, T) = \prod_{j \in J} \Gamma(U, \mathcal{O}_U)$
对任意概形 $U$ 都成立，而不再另行说明。
满环同态 $\mathbf{Z}[x_j] \to B$（$x_j \mapsto b_j$）
对应于闭浸入 $V \to T$。
于是得到一个单态射
$$
i : V \longrightarrow T_Y = T \times Y
$$
这是 $Y$ 上代数空间的单态射。就 $Y_\etale$ 上的层而言，
态射 $i$ 诱导层的单射
$h_i : h_V \to \prod_{j \in J} \mathcal{O}_Y$。
$i$ 到 $X$ 的基变换 $i' : X \times_Y V \to T_X$
也为单态射（《空间》，
引理 \ref{spaces-lemma-base-change-representable-transformations-property}）。
因此 $i' : X \times_Y V \to T_X$ 为单态射，这又意味着
$h_{i'} : h_{X \times_Y V} \to \prod_{j \in J} \mathcal{O}_X$
是层的单射。
通过引理 \ref{spaces-properties-lemma-stalk-pullback} 给出的恒等
$f_{small}^{-1}h_V = h_{X \times_Y V}$，
映射 $h_{i'}$ 等于
$$
\xymatrix{
f_{small}^{-1}h_V \ar[r]^-{f^{-1}h_i} &
\prod_{j \in J} f_{small}^{-1}\mathcal{O}_Y
\ar[r]^{\prod f^\sharp} &
\prod_{j \in J} \mathcal{O}_X
}
$$
(验证略去）。这意味着映射
$t : f_{small}^{-1}h_V \to f_{small}^{-1}h_V$
适合于下列交换图：
$$
\xymatrix{
f_{small}^{-1}h_V \ar[r]^-{f^{-1}h_i} \ar[d]^t &
\prod_{j \in J} f_{small}^{-1}\mathcal{O}_Y
\ar[r]^-{\prod f^\sharp} \ar[d]^{\prod t} &
\prod_{j \in J} \mathcal{O}_X \ar[d]^{\text{id}}\\
f_{small}^{-1}h_V \ar[r]^-{f^{-1}h_i} &
\prod_{j \in J} f_{small}^{-1}\mathcal{O}_Y
\ar[r]^-{\prod f^\sharp} &
\prod_{j \in J} \mathcal{O}_X
}
$$
右方方块的交换性由上文解释的关于 $t$ 的假设保证。
由于上面讨论的水平箭头复合是单射，我们也得到左方竖直箭头
是恒等映射。$Y_\etale$ 上的任意集合层都可以由一族形如
$h_V$（其中 $V$ 为仿射概形）的层的一个（很大的）余积满射到它
（结合引理 \ref{spaces-properties-lemma-alternative} 与
《站点》，引理 \ref{sites-lemma-sheaf-coequalizer-representable}）。
因此得到 $t : f_{small}^{-1} \to f_{small}^{-1}$
是恒等变换，正如所需。
\end{proof}

\begin{lemma}
\label{spaces-properties-lemma-faithful}
设 $S$ 为概形。
设 $X$、$Y$ 为 $S$ 上的代数空间。
对于 $S$ 上代数空间的任意两个态射 $a, b : X \to Y$，
若存在一个 $2$-同构
$(a_{small}, a^\sharp) \cong (b_{small}, b^\sharp)$
使得它们在环化拓扑斯的 $2$-范畴中相同，则二者相等。
\end{lemma}

\begin{proof}
设 $t : a_{small}^{-1} \to b_{small}^{-1}$ 为该 $2$-同构。
等价地，可以把 $t$ 看作变换
$t : a_{spaces, \etale}^{-1} \to b_{spaces, \etale}^{-1}$
因为 $X_\etale$ 上的层与 $X_{spaces, \etale}$ 上的层没有区别。
选取交换图
$$
\xymatrix{
U \ar[d]_p \ar[r]_\alpha  & V \ar[d]^q \\
X \ar[r]^a & Y
}
$$
其中 $U$、$V$ 为概形，且 $p$、$q$ 为满的 \'etale 态射。
考虑图
$$
\xymatrix{
h_U \ar[r]_-\alpha \ar@{=}[d] & a_{spaces, \etale}^{-1}h_V \ar[d]^t \\
h_U \ar@{..>}[r] & b_{spaces, \etale}^{-1}h_V
}
$$
由于层 $b_{spaces, \etale}^{-1}h_V$ 同构于
$h_{V \times_{Y, b} X}$，可知虚线箭头来自一个概形态射
$\beta : U \to V$，并适合交换图
$$
\xymatrix{
U \ar[d]_p \ar[r]_\beta  & V \ar[d]^q \\
X \ar[r]^b & Y
}
$$
我们断言存在如下 $2$-同构序列：
\begin{align*}
(\alpha_{small}, \alpha^\sharp)
& \cong
(\alpha_{spaces, \etale}, \alpha^\sharp) \\
& \cong
(a_{spaces, \etale, c}, a_c^\sharp) \\
& \cong
(b_{spaces, \etale, d}, b_d^\sharp) \\
& \cong
(\beta_{spaces, \etale}, \beta^\sharp) \\
& \cong
(\beta_{small}, \beta^\sharp)
\end{align*}
第一个和最后一个 $2$-同构来自
$U_{spaces, \etale}$ 上的层与 $U_\etale$ 上的层之间的恒等，
对 $V$ 也一样。第二个和第四个
$2$-同构来自
引理 \ref{spaces-properties-lemma-relocalize-morphism}，
其中 $c : U \to X \times_{a, Y} V$ 由 $\alpha$ 诱导，
$d : U \to X \times_{b, Y} V$ 由 $\beta$ 诱导。
中间的 $2$-同构来自变换 $t$。具体地，函子
$a_{spaces, \etale, c}^{-1}$ 对应于函子
$$
(\mathcal{H} \to h_V) \longmapsto
(a_{spaces, \etale}^{-1}\mathcal{H}
\times_{a_{spaces, \etale}^{-1}h_V, \alpha}
h_U \to
h_U)
$$
对 $b_{spaces, \etale, d}^{-1}$ 也同样，参见
《站点》，引理 \ref{sites-lemma-relocalize-morphism}。
这里使用了如下恒等：$Y_{spaces, \etale}/V$ 上的层
对应于 $\Sh(Y_{spaces, \etale})$ 中的箭头
$(\mathcal{H} \to h_V)$，对 $U/X$ 也一样；参见
《站点》，引理 \ref{sites-lemma-essential-image-j-shriek}。
在此恒等下，结构层 $\mathcal{O}_V$ 对应于
对 $(\mathcal{O}_Y \times h_V \to h_V)$，而 $\mathcal{O}_U$ 也类似；参见
《站点上的模》，
引理 \ref{sites-modules-lemma-localize-compare}。
由于 $t$ 交换 $\alpha$ 与 $\beta$，可知 $t$ 诱导同构
$$
t :
a_{spaces, \etale}^{-1}\mathcal{H}
\times_{a_{spaces, \etale}^{-1}h_V, \alpha}
h_U
\longrightarrow
b_{spaces, \etale}^{-1}\mathcal{H}
\times_{b_{spaces, \etale}^{-1}h_V, \beta}
h_U
$$
它在 $h_U$ 上，并且关于 $(\mathcal{H} \to h_V)$ 具有函子性。
此外，$t$ 与 $a_c^\sharp$、$b_d^\sharp$ 相容，因为按照上面对
结构层 $\mathcal{O}_U$、$\mathcal{O}_V$ 的描述，$t$ 与
$a^\sharp$、$b^\sharp$ 相容。因此，环化拓扑斯态射
$(\alpha_{small}, \alpha^\sharp)$ 和 $(\beta_{small}, \beta^\sharp)$
是 $2$-同构的。由
《\'Etale 上同调》，引理 \ref{etale-cohomology-lemma-faithful}，
得到 $\alpha = \beta$！由于 $p : U \to X$ 是层的满射，遂有 $a=b$。
\end{proof}

\noindent
下面是本节的主要结果。

\begin{theorem}
\label{spaces-properties-theorem-fully-faithful}
设 $X$、$Y$ 为 $\Spec(\mathbf{Z})$ 上的代数空间。
设
$$
(g, g^\sharp) :
(\Sh(X_\etale), \mathcal{O}_X)
\longrightarrow
(\Sh(Y_\etale), \mathcal{O}_Y)
$$
为局部环化拓扑斯的态射。则存在唯一的代数空间态射
$f : X \to Y$，使得
$(g, g^\sharp)$ 与 $(f_{small}, f^\sharp)$ 同构。
换言之，构造
$$
\textit{Spaces}/\Spec(\mathbf{Z})
\longrightarrow \textit{Locally ringed topoi},
\quad
X \longrightarrow (X_\etale, \mathcal{O}_X)
$$
是全忠实的（右侧的态射取至 $2$-同构）。
\end{theorem}

\begin{proof}
唯一性已在引理 \ref{spaces-properties-lemma-faithful} 中证明。
因此只需证明存在性。
在本证明中，我们将自由使用等式
（\ref{spaces-properties-equation-localize-at-scheme-ringed}）中的恒等，
以及引理 \ref{spaces-properties-lemma-relocalize-morphism-at-schemes} 的结果。

\medskip\noindent
令 $U \in \Ob(X_\etale)$，令
$V \in \Ob(Y_\etale)$
，并令 $s \in g^{-1}h_V(U)$ 为一个截面。可以把
$s$ 看作层的映射 $s : h_U \to g^{-1}h_V$。由
《站点上的模》，
引理 \ref{sites-modules-lemma-relocalize-morphism-ringed-topoi}，
得到环化拓扑斯态射的交换图
$$
\xymatrix{
(\Sh(X_\etale/U), \mathcal{O}_U)
\ar[rr]_-{(j, j^\sharp)} \ar[d]_{(g_s, g_s^\sharp)} & &
(\Sh(X_\etale), \mathcal{O}_X) \ar[d]^{(g, g^\sharp)} \\
(\Sh(V_\etale), \mathcal{O}_V) \ar[rr] & &
(\Sh(Y_\etale), \mathcal{O}_Y).
}
$$
由
《\'Etale 上同调》，定理 \ref{etale-cohomology-theorem-fully-faithful}，
得到唯一的概形态射 $f_s : U \to V$，使得
$(g_s, g_s^\sharp)$ 与 $(f_{s, small}, f_s^\sharp)$ 为 $2$-同构。
刚才说明的构造 $(U, V, s) \leadsto f_s$ 满足如下函子性：
设有 $X_\etale$ 中的态射 $a : U' \to U$ 以及 $Y_\etale$ 中的态射
$b : V' \to V$，
并有映射 $s' : h_{U'} \to g^{-1}h_{V'}$，使下图交换：
$$
\xymatrix{
h_{U'} \ar[d]_a \ar[r]_{s'} & g^{-1}h_{V'} \ar[d]^{g^{-1}b} \\
h_U \ar[r]^s & g^{-1}h_V
}
$$
则下图（概形的图）交换：
$$
\xymatrix{
U' \ar[r]_-{f_{s'}} \ar[d]_a & u(V') \ar[d]^{u(b)} \\
U \ar[r]^-{f_s} & u(V)
}
$$
之所以如此，是因为对由《站点上的模》
引理 \ref{sites-modules-lemma-relocalize-morphism-ringed-topoi}
构造的态射 $(g_s, g_s^\sharp)$，同样的条件成立；
再利用《\'Etale 上同调》定理
\ref{etale-cohomology-theorem-fully-faithful} 的唯一性即可。

\medskip\noindent
问题在于将态射 $f_s$ 胶合为代数空间的态射。
为此，先选取概形 $V$ 以及满的 \'etale 态射 $V \to Y$。
这意味着 $h_V \to *$ 为满射，因而 $g^{-1}h_V \to *$ 也为满射。
于是存在概形 $U$、满的 \'etale 态射 $U \to X$ 以及态射
$s : h_U \to g^{-1}h_V$。接着令 $R = V \times_Y V$，
$R' = U \times_X U$。由于 $g^{-1}$ 正合，有
$g^{-1}h_R = g^{-1}h_V \times g^{-1}h_V$。
因此 $s$ 诱导态射 $s \times s : h_{R'} \to g^{-1}h_R$。
应用上述构造可知得到概形态射的交换图
$$
\xymatrix{
R' \ar@<1ex>[d] \ar@<-1ex>[d] \ar[rr]_{f_{s \times s}} & &
R \ar@<1ex>[d] \ar@<-1ex>[d] \\
U \ar[rr]^{f_s} & &
V
}
$$
由于 $X = U/R'$ 且 $Y = V/R$（参见
《空间》，引理 \ref{spaces-lemma-space-presentation}），
可知该图定义了代数空间态射 $f : X \to Y$，并且落在一个显然的交换图中。
现在还需证明 $(f_{small}, f^\sharp)$ 与 $(g, g^\sharp)$
是 $2$-同构的。
令 $t_V : f_{s, small}^{-1} \to g_s^{-1}$ 以及
$t_R : f_{s \times s, small}^{-1} \to g_{s \times s}^{-1}$
为上述构造给出的 $2$-同构。
令 $\mathcal{G}$ 为 $Y_\etale$ 上的层。则 $t_V$ 定义同构
$$
f_{small}^{-1}\mathcal{G}|_{U_\etale}
=
f_{s, small}^{-1}\mathcal{G}|_{V_\etale}
\xrightarrow{t_V}
g_s^{-1}\mathcal{G}|_{V_\etale}
=
g^{-1}\mathcal{G}|_{U_\etale}.
$$
此外，经任一投影 $R' \to U$ 拉回到 $R'$ 后，该同构就是
$$
f_{small}^{-1}\mathcal{G}|_{R'_\etale}
=
f_{s \times s, small}^{-1}\mathcal{G}|_{R_\etale}
\xrightarrow{t_R}
g_{s \times s}^{-1}\mathcal{G}|_{R_\etale}
=
g^{-1}\mathcal{G}|_{R'_\etale}.
$$
由于 $\{U \to X\}$ 是站点 $X_{spaces, \etale}$ 中的覆盖，
这意味着第一个显示的同构下降为层的同构
$t : f_{small}^{-1}\mathcal{G} \to g^{-1}\mathcal{G}$
（细节略去）。由于 $t_V$ 和 $t_R$ 是函子变换，该同构关于
$\mathcal{G}$ 具有函子性。最后，$t$ 与 $f^\sharp$、$g^\sharp$ 相容，
因为 $t_V$ 和 $t_R$ 具有这一性质（细节略去）。
定理得证。
\end{proof}

\begin{lemma}
\label{spaces-properties-lemma-isomorphism-ringed-topoi}
设 $X$、$Y$ 为 $\mathbf{Z}$ 上的代数空间。若
$$
(g, g^\sharp) :
(\Sh(X_\etale), \mathcal{O}_X)
\longrightarrow
(\Sh(Y_\etale), \mathcal{O}_Y)
$$
是环化拓扑斯的同构，则存在唯一的代数空间态射
$f : X \to Y$，使得 $(g, g^\sharp)$ 与 $(f_{small}, f^\sharp)$
同构，并且 $f$ 是代数空间的同构。
\end{lemma}

\begin{proof}
由定理 \ref{spaces-properties-theorem-fully-faithful}，只需证明 $(g, g^\sharp)$ 是
局部环化拓扑斯的态射。由《站点上的模》，
引理 \ref{sites-modules-lemma-locally-ringed-morphism}
（并且因为站点 $X_\etale$ 有足够多的点），只需检查
映射
$\mathcal{O}_{Y, q} \to \mathcal{O}_{X, p}$
由 $g^\sharp$ 诱导，在 $q = f \circ p$ 且 $p$ 为 $X_\etale$ 的任意点时，
是局部环同态。由于它是同构，结论显然。
\end{proof}














\section{代数空间上的拟凝聚层}
\label{spaces-properties-section-quasi-coherent}

\noindent
在《下降》的各节
\ref{descent-section-quasi-coherent-sheaves}、
\ref{descent-section-quasi-coherent-cohomology} 以及
\ref{descent-section-quasi-coherent-sheaves-bis} 中，我们已经看到，对于概形 $U$，
$U$ 上的拟凝聚 $\mathcal{O}_U$-模与 $U$ 的小 \'etale 站点上的拟凝聚
$\mathcal{O}$-模没有区别。因此，将以下定义应用于可表示代数空间时，
它与我们原先对概形上的拟凝聚层的概念相容
（《概形》，第 \ref{schemes-section-quasi-coherent} 节）。

\begin{definition}
\label{spaces-properties-definition-quasi-coherent}
设 $S$ 为概形，$X$ 为 $S$ 上的代数空间。
一个{\it 拟凝聚} $\mathcal{O}_X$-模
是环化站点 $(X_\etale, \mathcal{O}_X)$ 上的拟凝聚模，
其意义参见《站点上的模》，
定义 \ref{sites-modules-definition-site-local}。
$X$ 上拟凝聚层的范畴记作
$\QCoh(\mathcal{O}_X)$.
\end{definition}

\noindent
注意，拟凝聚性是内在的概念（参见《站点上的模》，
引理 \ref{sites-modules-lemma-special-locally-free}），
所以这等价于说，在 $X_{spaces, \etale}$ 上相应的 $\mathcal{O}_X$-模是拟凝聚的。

\medskip\noindent
如通常情形，拟凝聚层关于拉回具有良好性质。

\begin{lemma}
\label{spaces-properties-lemma-pullback-quasi-coherent}
设 $S$ 为概形。
设 $f : X \to Y$ 为 $S$ 上代数空间的态射。
拉回函子
$f^* : \textit{Mod}(\mathcal{O}_Y) \to \textit{Mod}(\mathcal{O}_X)$
保持拟凝聚层。
\end{lemma}

\begin{proof}
这是一般事实，参见《站点上的模》，引理
\ref{sites-modules-lemma-local-pullback}。
\end{proof}

\noindent
注意，如果 $X$、$Y$ 恰为概形，该拉回函子与模的拟凝聚层之间的通常拉回函子
一致，参见《下降》，命题
\ref{descent-proposition-equivalence-quasi-coherent-functorial}。
下面给出一个必要引理，将其与 $X$ 的小 \'etale 站点的对象上的拟凝聚层作比较。

\begin{lemma}
\label{spaces-properties-lemma-characterize-quasi-coherent-small-etale}
设 $S$ 为概形，$X$ 为 $S$ 上的代数空间。
一个拟凝聚 $\mathcal{O}_X$-模 $\mathcal{F}$
由以下数据给出：
\begin{enumerate}
\item 对每个 $U \in \Ob(X_\etale)$，在 $U_\etale$ 上的一个拟凝聚
$\mathcal{O}_U$-模 $\mathcal{F}_U$，
\item 对 $X_\etale$ 中的每个 $f : U' \to U$，一个同构
$c_f : f_{small}^*\mathcal{F}_U \to \mathcal{F}_{U'}$.
\end{enumerate}
这些数据须满足如下条件：给定 $X_\etale$ 中的任意
$f : U' \to U$ 和 $g : U'' \to U'$，复合
$c_g \circ g_{small}^*c_f$ 等于 $c_{f \circ g}$。
\end{lemma}

\begin{proof}
结合引理 \ref{spaces-properties-lemma-pullback-quasi-coherent} 与
\ref{spaces-properties-lemma-characterize-module-small-etale} 即得。
\end{proof}

\begin{lemma}
\label{spaces-properties-lemma-stalk-quasi-coherent}
设 $S$ 为概形。
设 $X$ 为 $S$ 上的代数空间。
设 $\mathcal{F}$ 为拟凝聚 $\mathcal{O}_X$-模。
设 $x \in |X|$ 为一点，$\overline{x}$ 为位于 $x$ 上方的几何点。
最后，设
$\varphi : (U, \overline{u}) \to (X, \overline{x})$
为 \'etale 邻域，其中 $U$ 为概形。则
$$
(\varphi^*\mathcal{F})_u \otimes_{\mathcal{O}_{U, u}}
\mathcal{O}_{X, \overline{x}} =
\mathcal{F}_{\overline{x}}
$$
其中 $u \in U$ 是 $\overline{u}$ 的像。
\end{lemma}

\begin{proof}
注意，由引理 \ref{spaces-properties-lemma-describe-etale-local-ring}，
$\mathcal{O}_{X, \overline{x}} = \mathcal{O}_{U, u}^{sh}$，
所以张量积是有意义的。此外，由定义 \ref{spaces-properties-definition-stalk} 可知
$$
\mathcal{F}_{\overline{u}} = \colim (\varphi^*\mathcal{F})_u
$$
其中余极限遍历引理中所述的
$\varphi : (U, \overline{u}) \to (X, \overline{x})$。
因此，引理陈述中的左侧到右侧有一个典范映射。
$\mathcal{O}_{X, \overline{x}}$ 也有类似的余极限描述，并且由
引理 \ref{spaces-properties-lemma-characterize-quasi-coherent-small-etale}，有
$$
((\varphi')^*\mathcal{F})_{u'} =
(\varphi^*\mathcal{F})_u \otimes_{\mathcal{O}_{U, u}} \mathcal{O}_{U', u'}
$$
只要 $(U', \overline{u}') \to (U, \overline{u})$ 是 \'etale 邻域的态射。
最后使用 $\otimes$ 与余极限交换即可完成证明。
\end{proof}

\begin{lemma}
\label{spaces-properties-lemma-stalk-pullback-quasi-coherent}
设 $S$ 为概形。设 $f : X \to Y$ 为 $S$ 上代数空间的态射。
设 $\mathcal{G}$ 为拟凝聚 $\mathcal{O}_Y$-模。
设 $\overline{x}$ 为 $X$ 的几何点，且
$\overline{y} = f \circ \overline{x}$ 为其在 $Y$ 中的像。
则存在典范同构
$$
(f^*\mathcal{G})_{\overline{x}} =
\mathcal{G}_{\overline{y}} \otimes_{\mathcal{O}_{Y, \overline{y}}}
\mathcal{O}_{X, \overline{x}}
$$
把拉回的茎同构于该茎与 $\overline{x}$ 处 $X$ 的局部环的张量积。
\end{lemma}

\begin{proof}
由于
$f^*\mathcal{G} = f_{small}^{-1}\mathcal{G}
\otimes_{f_{small}^{-1}\mathcal{O}_Y} \mathcal{O}_X$，
由引理 \ref{spaces-properties-lemma-stalk-pullback} 中对拉回茎的描述以及取茎与张量积交换的事实，
即得结论。也可以更直接地如下看出这一点。
选取交换图
$$
\xymatrix{
U \ar[d]_p \ar[r]_\alpha  & V \ar[d]^q \\
X \ar[r]^a & Y
}
$$
其中 $U$、$V$ 为概形，且 $p$、$q$ 为满的 \'etale 态射。
由引理 \ref{spaces-properties-lemma-geometric-lift-to-usual}，可以选取 $U$ 的几何点
$\overline{u}$，使得 $\overline{x} = p \circ \overline{u}$。
令 $\overline{v} = \alpha \circ \overline{u}$。于是
\begin{align*}
(f^*\mathcal{G})_{\overline{x}} & =
(p^*f^*\mathcal{G})_u \otimes_{\mathcal{O}_{U, u}}
\mathcal{O}_{X, \overline{x}} \\
& = (\alpha^*q^*\mathcal{G})_u \otimes_{\mathcal{O}_{U, u}}
\mathcal{O}_{X, \overline{x}} \\
& = (q^*\mathcal{G})_v \otimes_{\mathcal{O}_{V, v}}
\mathcal{O}_{U, u} \otimes_{\mathcal{O}_{U, u}}
\mathcal{O}_{X, \overline{x}} \\
& = (q^*\mathcal{G})_v \otimes_{\mathcal{O}_{V, v}}
\mathcal{O}_{X, \overline{x}} \\
& = (q^*\mathcal{G})_v \otimes_{\mathcal{O}_{V, v}}
\mathcal{O}_{Y, \overline{y}} \otimes_{\mathcal{O}_{Y, \overline{y}}}
\mathcal{O}_{X, \overline{x}} \\
& = \mathcal{G}_{\overline{y}} \otimes_{\mathcal{O}_{Y, \overline{y}}}
\mathcal{O}_{X, \overline{x}}
\end{align*}
这里使用了引理 \ref{spaces-properties-lemma-stalk-quasi-coherent}（两次）以及概形上拟凝聚层拉回的相应结果，
参见《层》，引理 \ref{sheaves-lemma-stalk-pullback-modules}。
\end{proof}

\begin{lemma}
\label{spaces-properties-lemma-characterize-quasi-coherent}
设 $S$ 为概形，$X$ 为 $S$ 上的代数空间。
设 $\mathcal{F}$ 为 $\mathcal{O}_X$-模层。以下条件等价：
\begin{enumerate}
\item $\mathcal{F}$ 是拟凝聚 $\mathcal{O}_X$-模；
\item 存在 $S$ 上代数空间的 \'etale 态射 $f : Y \to X$，其
$|f| : |Y| \to |X|$ 为满射，并且 $f^*\mathcal{F}$ 在 $Y$ 上拟凝聚；
\item 存在概形 $U$ 及满的 \'etale 态射
$\varphi : U \to X$，使得 $\varphi^*\mathcal{F}$ 是拟凝聚
$\mathcal{O}_U$-模；并且
\item 对每个仿射概形 $U$ 及 \'etale 态射 $\varphi : U \to X$，
限制 $\varphi^*\mathcal{F}$ 是拟凝聚 $\mathcal{O}_U$-模。
\end{enumerate}
\end{lemma}

\begin{proof}
考虑 $\text{id}_X$，显然 (1) 蕴含 (2)。
假设 $f : Y \to X$ 如 (2) 中所述，且 $V \to Y$ 是来自概形到 $Y$ 的满
\'etale 态射。则复合 $V \to X$ 也为满的 \'etale 态射，
并且由引理 \ref{spaces-properties-lemma-pullback-quasi-coherent}，$\mathcal{F}$
拉回到 $V$ 上仍拟凝聚。因此 (2) 蕴含 (3)。

\medskip\noindent
令 $U \to X$ 如 (3) 中所述。沿用等式
(\ref{spaces-properties-equation-restrict-modules}) 中引入的记号滥用。
由于 $\mathcal{F}|_{U_\etale}$ 拟凝聚，存在一个 \'etale 覆盖
$\{U_i \to U\}$，使得 $\mathcal{F}|_{U_{i, \etale}}$ 有全局表示，参见
《站点上的模》，定义 \ref{sites-modules-definition-global} 与
引理 \ref{sites-modules-lemma-local-final-object}。
令 $V \to X$ 为 $X_\etale$ 的一个对象。由于 $U \to X$ 满且 \'etale，
映射族 $\{U_i \times_X V \to V\}$ 是 $V$ 的 \'etale 覆盖。
通过态射 $U_i \times_X V \to U_i$，可以将
$\mathcal{F}|_{U_{i, \etale}}$ 的全局表示限制为
$\mathcal{F}|_{(U_i \times_X V)_\etale}$ 的全局表示。
因此 $X_\etale$ 上的层 $\mathcal{F}$ 满足《站点上的模》，定义
\ref{sites-modules-definition-site-local} 的条件，故为拟凝聚。

\medskip\noindent
 (3) 与 (4) 的等价性来自每个概形都有仿射开覆盖这一事实。
\end{proof}

\begin{lemma}
\label{spaces-properties-lemma-properties-quasi-coherent}
设 $S$ 为概形，$X$ 为 $S$ 上的代数空间。
$X$ 上拟凝聚层的范畴 $\QCoh(\mathcal{O}_X)$
具有以下性质：
\begin{enumerate}
\item 任意拟凝聚层的直和仍为拟凝聚。
\item 任意拟凝聚层的余极限仍为拟凝聚。
\item 拟凝聚层态射的核与余核是拟凝聚的。
\item 给定 $\mathcal{O}_X$-模的短正合序列
$0 \to \mathcal{F}_1 \to \mathcal{F}_2 \to \mathcal{F}_3 \to 0$
若其中任意两个是拟凝聚的，则第三个也是拟凝聚的。
\item 给定两个拟凝聚 $\mathcal{O}_X$-模，它们的张量积是拟凝聚的。
\item 给定两个拟凝聚 $\mathcal{O}_X$-模
$\mathcal{F}$、$\mathcal{G}$，其中 $\mathcal{F}$
是有限表示的（参见第 \ref{spaces-properties-section-properties-modules} 节），
则内部 Hom
$\SheafHom_{\mathcal{O}_X}(\mathcal{F}, \mathcal{G})$
是拟凝聚的。
\end{enumerate}
\end{lemma}

\begin{proof}
若 $X$ 为概形，这就是《下降》，引理
\ref{descent-lemma-properties-quasi-coherent} 的结论。
我们将通过 \'etale 局部化把本引理化归到这种情形。

\medskip\noindent
选取概形 $U$ 及满的 \'etale 态射 $\varphi : U \to X$。
我们约定记号 $\textit{Mod}(\mathcal{O}_U) =
\textit{Mod}(U_\etale, \mathcal{O}_U)$ 以及
$\QCoh(\mathcal{O}_U) = \QCoh(U_\etale, \mathcal{O}_U)$；换言之，
尽管 $U$ 是概形，我们仍将 $U$ 上的拟凝聚模视为 $U$ 的小 \'etale 站点上的模。
由引理 \ref{spaces-properties-lemma-pullback-quasi-coherent} 有交换图
$$
\xymatrix{
\QCoh(\mathcal{O}_X) \ar[r]_{\varphi^*} \ar[d] &
\QCoh(\mathcal{O}_U) \ar[d] \\
\textit{Mod}(\mathcal{O}_X) \ar[r]^{\varphi^*} &
\textit{Mod}(\mathcal{O}_U)
}
$$
下方水平箭头是限制函子
(\ref{spaces-properties-equation-restrict-modules})
$\mathcal{G} \mapsto \mathcal{G}|_{U_\etale}$。
该函子同时有左伴随和右伴随，参见《站点上的模》，第
\ref{sites-modules-section-localize} 节，故与所有极限和余极限交换。
此外，已知 $\textit{Mod}(\mathcal{O}_X)$ 的对象属于
$\QCoh(\mathcal{O}_X)$ 当且仅当其限制到 $U$ 后属于
$\QCoh(\mathcal{O}_U)$，参见引理 \ref{spaces-properties-lemma-characterize-quasi-coherent}。
完成这些准备后即可开始证明。

\medskip\noindent
性质 (1) 的证明。设 $\mathcal{F}_i$（$i \in I$）为一族拟凝聚
$\mathcal{O}_X$-模。由上面的讨论，有
$$
\Big(\bigoplus \mathcal{F}_i\Big)|_{U_\etale} =
\bigoplus \mathcal{F}_i|_{U_\etale}
$$
每个模 $\mathcal{F}_i|_{U_\etale}$ 都是拟凝聚的。
因此由概形情形，直和是拟凝聚的。
于是 $\bigoplus \mathcal{F}_i$ 作为限制到 $U$ 后为拟凝聚模的模，
也是拟凝聚的。

\medskip\noindent
性质 (2) 的证明。设 $\mathcal{I} \to \QCoh(\mathcal{O}_X)$、
$i \mapsto \mathcal{F}_i$ 为一个图表。则
$$
(\colim \mathcal{F}_i)|_{U_\etale} = \colim \mathcal{F}_i|_{U_\etale}
$$
由上面的讨论，这与上述相同方式给出结论。

\medskip\noindent
性质 (3) 的证明。设 $a : \mathcal{F} \to \mathcal{F}'$
为 $\QCoh(\mathcal{O}_X)$ 中的箭头。则
有 $\Ker(a)|_{U_\etale} = \Ker(a|_{U_\etale})$
以及 $\Coker(a)|_{U_\etale} = \Coker(a|_{U_\etale})$，
同样得到结论。

\medskip\noindent
性质 (4) 的证明。限制
$0 \to \mathcal{F}_1|_{U_\etale} \to \mathcal{F}_2|_{U_\etale}
\to \mathcal{F}_3|_{U_\etale} \to 0$
是短正合的。因此该序列具有二取三性质，结论如前。

\medskip\noindent
性质 (5) 的证明。设 $\mathcal{F}$、$\mathcal{G}$ 属于 $\QCoh(\mathcal{O}_X)$。
则有
$$
(\mathcal{F} \otimes_{\mathcal{O}_X} \mathcal{G})_{U_\etale} =
\mathcal{F}|_{U_\etale} \otimes_{\mathcal{O}_U} \mathcal{G}|_{U_\etale}
$$
结论如前。

\medskip\noindent
性质 (6) 的证明。设 $\mathcal{F}$、$\mathcal{G}$
属于 $\QCoh(\mathcal{O}_X)$，且 $\mathcal{F}$
有限表示。则有
$$
\SheafHom_{\mathcal{O}_X}(\mathcal{F}, \mathcal{G})|_{U_\etale} =
\SheafHom_{\mathcal{O}_U}(\mathcal{F}|_{U_\etale}, \mathcal{G}|_{U_\etale})
$$
事实上，限制是局部化，参见第 \ref{spaces-properties-section-localize} 节，尤其是公式
(\ref{spaces-properties-equation-localize-at-scheme-ringed}))；内部 Hom 的形成与局部化交换，参见
《站点上的模》，引理 \ref{sites-modules-lemma-internal-hom-restriction}。
因此结论如前。
\end{proof}

\noindent
一般而言，沿代数空间态射的拟凝聚层推前不一定是拟凝聚的。
我们将在《空间的态射》，第 \ref{spaces-morphisms-section-pushforward} 节中回到这一问题。




\section{模的性质}
\label{spaces-properties-section-properties-modules}

\noindent
在《站点上的模》的各节
\ref{sites-modules-section-global}、\ref{sites-modules-section-local} 以及
定义 \ref{sites-modules-definition-flat} 中，我们为任意环化拓扑斯上的
$\mathcal{O}$-模定义了一些内在性质。若 $X$ 为代数空间，我们将自由地把这些概念
应用于环化站点 $(X_\etale, \mathcal{O}_X)$ 上的模，或等价地应用于环化站点
$(X_{spaces, \etale}, \mathcal{O}_X)$ 上的模。

\medskip\noindent
全局性质 $\mathcal{P}$：
\begin{enumerate}
\item[(a)] {\it 自由}；
\item[(b)] {\it 有限自由}；
\item[(c)] {\it 由全局截面生成}；
\item[(d)] {\it 由有限多个全局截面生成}；
\item[(e)] 具有{\it 全局表示}；以及
\item[(f)] 具有{\it 全局有限表示}。
\end{enumerate}
局部性质 $\mathcal{P}$：
\begin{enumerate}
\item[(g)] {\it 局部自由}；
\item[(f)] {\it 有限局部自由}；
\item[(h)] {\it 局部由截面生成}；
\item[(i)] {\it 局部由 $r$ 个截面生成}；
\item[(j)] {\it 有限型}；
\item[(k)] {\it 拟凝聚}（参见第 \ref{spaces-properties-section-quasi-coherent} 节）；
\item[(l)] {\it 有限表示}；
\item[(m)] {\it 相干}；以及
\item[(n)] {\it 平坦}。
\end{enumerate}
下面是一些由定义立即得到的结果：
\begin{enumerate}
\item 除 $\mathcal{P}=$“相干”外，在每种情形下该性质都在拉回下保持，参见
《站点上的模》，引理 \ref{sites-modules-lemma-global-pullback}、
\ref{sites-modules-lemma-local-pullback} 以及
\ref{sites-modules-lemma-pullback-flat}。
\item 上述每个性质（包括相干性）都在代数空间的 \'etale 态射拉回下保持
（因为此时拉回由限制给出，参见引理 \ref{spaces-properties-lemma-etale-morphism-topoi}）。
\item 假设 $f : Y \to X$ 是代数空间的满 \'etale 态射。对于局部性质
(g)--(m)，若 $f^*\mathcal{F}$ 具有 $\mathcal{P}$，则 $\mathcal{F}$
也具有 $\mathcal{P}$。这是因为 $\{Y \to X\}$ 是
$X_{spaces, \etale}$ 中的覆盖，并应用《站点上的模》引理
\ref{sites-modules-lemma-local-final-object}。
\item 若 $X$ 为概形，$\mathcal{F}$ 为 $X_\etale$ 上的拟凝聚模，且
$\mathcal{P}$ 为除“相干”或“局部自由”外的任一性质，则 $X_\etale$ 上
$\mathcal{F}$ 的 $\mathcal{P}$ 等价于 $\mathcal{F}|_{X_{Zar}}$ 的相应性质；
也就是说，把它看作概形 $X$ 上的拟凝聚层时，$\mathcal{F}$ 具有 $\mathcal{P}$。
参见《下降》，引理 \ref{descent-lemma-equivalence-quasi-coherent-properties}。
\item 若 $X$ 为局部 Noether 概形，$\mathcal{F}$ 为 $X_\etale$ 上的
拟凝聚模，则 $\mathcal{F}$ 在 $X_\etale$ 上相干，当且仅当
$\mathcal{F}|_{X_{Zar}}$ 相干；也就是说，这正对应于通常意义下
概形 $X$ 上的相干层。参见《下降》，引理
\ref{descent-lemma-equivalence-quasi-coherent-properties}。
\end{enumerate}



\section{局部投射模}
\label{spaces-properties-section-locally-projective}

\noindent
回忆一下，在《性质》，第 \ref{properties-section-locally-projective} 节中，
我们定义了局部投射拟凝聚模的概念。

\begin{lemma}
\label{spaces-properties-lemma-locally-projective}
设 $S$ 为概形，$X$ 为 $S$ 上的代数空间。
设 $\mathcal{F}$ 为拟凝聚 $\mathcal{O}_X$-模。以下条件等价：
\begin{enumerate}
\item 存在概形 $U$ 及满的 \'etale 态射 $U \to X$，使得限制
$\mathcal{F}|_U$ 在 $U$ 上局部投射；并且
\item 对任意概形 $U$ 及任意 \'etale 态射 $U \to X$，限制
$\mathcal{F}|_U$ 在 $U$ 上局部投射。
\end{enumerate}
\end{lemma}

\begin{proof}
令 $U \to X$ 如 (1) 中所述，并令 $V \to X$ 为 \'etale 态射，其中
$V$ 为概形。则 $\{U \times_X V \to V\}$ 是概形的 fppf 覆盖。
因此，若 $\mathcal{F}|_U$ 局部投射，则 $\mathcal{F}|_{U \times_X V}$
局部投射（参见《性质》，引理
\ref{properties-lemma-locally-projective-pullback}），
从而 $\mathcal{F}|_V$ 局部投射，参见《下降》，引理
\ref{descent-lemma-locally-projective-descends}。
\end{proof}

\begin{definition}
\label{spaces-properties-definition-locally-projective}
设 $S$ 为概形，$X$ 为 $S$ 上的代数空间。
设 $\mathcal{F}$ 为拟凝聚 $\mathcal{O}_X$-模。
若引理 \ref{spaces-properties-lemma-locally-projective} 的等价条件成立，
则称 $\mathcal{F}$ 为{\it 局部投射的}。
\end{definition}

\begin{lemma}
\label{spaces-properties-lemma-locally-projective-pullback}
设 $S$ 为概形。
设 $f : X \to Y$ 为 $S$ 上代数空间的态射。
设 $\mathcal{G}$ 为拟凝聚 $\mathcal{O}_Y$-模。
若 $\mathcal{G}$ 在 $Y$ 上局部投射，则 $f^*\mathcal{G}$
在 $X$ 上局部投射。
\end{lemma}

\begin{proof}
选取满的 \'etale 态射 $V \to Y$，其中 $V$ 为概形。
选取满的 \'etale 态射 $U \to V \times_Y X$，其中
$U$ 为概形。记诱导的态射为 $\psi : U \to V$。则
$$
f^*\mathcal{G}|_U = \psi^*(\mathcal{G}|_V)
$$
因此由定义以及概形情形的结果可得引理，参见《性质》，引理
\ref{properties-lemma-locally-projective-pullback}。
\end{proof}





\section{拟凝聚层与表示}
\label{spaces-properties-section-quasi-coherent-presentation}

\noindent
设 $S$ 为概形，$X$ 为 $S$ 上的代数空间。
令 $X = U/R$ 为由任意满的 \'etale 态射 $\varphi : U \to X$
得到的 $X$ 的一个表示，参见《空间》，定义
\ref{spaces-definition-presentation}。
特别地，我们得到群胚 $(U, R, s, t, c)$，其中
$j = (t, s) : R \to U \times_S U$，参见《群胚》，引理
\ref{groupoids-lemma-equivalence-groupoid}。
在《群胚》，定义 \ref{groupoids-definition-groupoid-module} 中，
我们定义了任意群胚上的拟凝聚层概念。具备这些记号后，有如下观察。

\begin{proposition}
\label{spaces-properties-proposition-quasi-coherent}
设 $S$、$\varphi : U \to X$ 以及 $(U, R, s, t, c)$ 如上。
对任意拟凝聚 $\mathcal{O}_X$-模 $\mathcal{F}$，层
$\varphi^*\mathcal{F}$ 带有如下典范同构：
$$
\alpha : t^*\varphi^*\mathcal{F} \longrightarrow s^*\varphi^*\mathcal{F}
$$
它满足《群胚》，定义 \ref{groupoids-definition-groupoid-module} 的条件，
因而在 $(U, R, s, t, c)$ 上定义拟凝聚层。
函子 $\mathcal{F} \mapsto (\varphi^*\mathcal{F}, \alpha)$
定义范畴等价
$$
\begin{matrix}
\text{Quasi-coherent} \\
\mathcal{O}_X\text{-模}
\end{matrix}
\longleftrightarrow
\begin{matrix}
\text{Quasi-coherent modules}\\
\text{定义于 }(U, R, s, t, c)
\end{matrix}
$$
\end{proposition}

\begin{proof}
在命题陈述及本证明中，我们把概形上的拟凝聚层视为该概形小
\'etale 站点上的拟凝聚层。这由《下降》各节的结果是允许的：
\ref{descent-section-quasi-coherent-sheaves}、
\ref{descent-section-quasi-coherent-cohomology} 以及
\ref{descent-section-quasi-coherent-sheaves-bis}。

\medskip\noindent
由于 $\varphi \circ t = \varphi \circ s$ 且拉回关于态射具有函子性，
便得到 $\alpha$ 的存在性，参见等式 (\ref{spaces-properties-equation-push-pull}) 周围的讨论。
完全同样地，即利用拉回的函子性，可知同构 $\alpha$ 满足《群胚》，
定义 \ref{groupoids-definition-groupoid-module} 的条件 (1)。
要验证定义的条件 (2)，只需注意 $\alpha$ 是同构，这显然成立。
构造 $\mathcal{F} \mapsto (\varphi^*\mathcal{F}, \alpha)$
显然关于拟凝聚层 $\mathcal{F}$ 具有函子性。
因此得到引理显示公式中从左到右的函子。

\medskip\noindent
反之，设 $(\mathcal{F}, \alpha)$ 是 $(U, R, s, t, c)$ 上的拟凝聚层。
令 $V \to X$ 为 $X_\etale$ 的一个对象。此时态射
$V' = U \times_X V \to V$ 是概形的满 \'etale 态射，故
$\{V' \to V\}$ 是 $V$ 的 \'etale 覆盖。此外，拟凝聚层 $\mathcal{F}$
拉回为 $V'$ 上的拟凝聚层 $\mathcal{F}'$。
由于 $R = U \times_X U$ 且 $t = \text{pr}_0$、$s = \text{pr}_0$，
可知 $V' \times_V V' = R \times_X V$，且投影映射
$V' \times_V V' \to V'$ 等于 $t$、$s$ 的拉回。因此 $\alpha$ 拉回为同构
$\alpha' : \text{pr}_0^*\mathcal{F}' \to \text{pr}_1^*\mathcal{F}'$，
并且二元组 $(\mathcal{F}', \alpha')$ 是关于 $\{V' \to V\}$ 的
拟凝聚层下降数据。由《下降》，命题
\ref{descent-proposition-fpqc-descent-quasi-coherent}，
该下降数据有效，从而得到 $V_\etale$ 上的拟凝聚
$\mathcal{O}_V$-模 $\mathcal{F}_V$。
要看出这给出了 $X_\etale$ 上的拟凝聚层，需要证明（由
引理 \ref{spaces-properties-lemma-characterize-quasi-coherent-small-etale}）
对于 $X_\etale$ 中的任意态射 $f : V_1 \to V_2$，存在典范同构
$c_f : \mathcal{F}_{V_1} \to \mathcal{F}_{V_2}$
且与态射复合相容。验证略去。
我们也略去如下验证：这定义了一个从右侧范畴到左侧范畴的函子，
并且它是上述函子的逆。
\end{proof}

\begin{proposition}
\label{spaces-properties-proposition-coherator}
设 $S$ 为概形，$X$ 为 $S$ 上的代数空间。
\begin{enumerate}
\item 范畴 $\QCoh(\mathcal{O}_X)$ 是 Grothendieck 阿贝尔范畴。
因此 $\QCoh(\mathcal{O}_X)$ 有足够多的内射对象并且有所有极限。
\item 包含函子
$\QCoh(\mathcal{O}_X) \to \textit{Mod}(\mathcal{O}_X)$
有右伴随\footnote{该函子有时称为{\it 相干化函子}。}
$$
Q : \textit{Mod}(\mathcal{O}_X) \longrightarrow \QCoh(\mathcal{O}_X)
$$
使得对每个拟凝聚层 $\mathcal{F}$，伴随映射
$Q(\mathcal{F}) \to \mathcal{F}$ 是同构。
\end{enumerate}
\end{proposition}

\begin{proof}
本证明重复概形情形的证明，参见《性质》，命题
\ref{properties-proposition-coherator}。
建议读者先阅读该证明。

\medskip\noindent
性质 (1) 意味着 $\QCoh(\mathcal{O}_X)$：(a) 有所有余极限；
(b) 滤余极限正合；(c) 有生成元，参见《内射》，第
\ref{injectives-section-grothendieck-conditions} 节。
由引理 \ref{spaces-properties-lemma-properties-quasi-coherent}，
$\QCoh(\mathcal{O}_X)$ 中的余极限存在并且与
$\textit{Mod}(\mathcal{O}_X)$ 中的余极限一致。由《站点上的模》，引理
\ref{sites-modules-lemma-limits-colimits}，滤余极限正合。
因此 (a)、(b) 成立。

\medskip\noindent
为构造一个生成元，选取表示 $X = U/R$，使得
$(U, R, s, t, c)$ 是一个
\'etale 群胚概形，特别地 $s$ 和 $t$ 是概形的平坦态射。选取《群胚》引理
\ref{groupoids-lemma-colimit-kappa} 中的一个基数 $\kappa$。
选取一族 $(\mathcal{E}_t, \alpha_t)_{t \in T}$，其中的对象是
$(U, R, s, t, c)$ 上的 $\kappa$-生成拟凝聚模，且如《群胚》引理
\ref{groupoids-lemma-set-of-iso-classes} 所述。
令 $\mathcal{F}_t$ 为 $X$ 上的拟凝聚模，它通过范畴等价
命题 \ref{spaces-properties-proposition-quasi-coherent} 对应于拟凝聚模
$(\mathcal{E}_t, \alpha_t)$。
于是可见，每个拟凝聚模 $\mathcal{H}$ 都是其拟凝聚子模的
有向余极限，而这些子模同构于某个 $\mathcal{F}_t$。因此
$\bigoplus_t \mathcal{F}_t$ 是 $\QCoh(\mathcal{O}_X)$ 的一个生成元，
从而 (c) 成立。关于极限和内射对象的断言在任意
Grothendieck 阿贝尔范畴中成立，参见《内射》定理
\ref{injectives-theorem-injective-embedding-grothendieck} 以及
引理 \ref{injectives-lemma-grothendieck-products}。

\medskip\noindent
\medskip\noindent
(2) 的证明。为构造 $Q$，我们使用如下的一般过程。
给定 $\textit{Mod}(\mathcal{O}_X)$ 中的一个对象 $\mathcal{F}$，
考虑函子
$$
\QCoh(\mathcal{O}_X)^{opp} \longrightarrow \textit{Sets},\quad
\mathcal{G} \longmapsto \Hom_X(\mathcal{G}, \mathcal{F})
$$
该函子把余极限变为极限，因而是可表示的，参见《内射》引理
\ref{injectives-lemma-grothendieck-brown}。
因此存在一个拟凝聚层 $Q(\mathcal{F})$ 以及函子性同构
$\Hom_X(\mathcal{G}, \mathcal{F}) = \Hom_X(\mathcal{G}, Q(\mathcal{F}))$
对 $\QCoh(\mathcal{O}_X)$ 中的 $\mathcal{G}$ 成立。由 Yoneda 引理
（《范畴》，引理 \ref{categories-lemma-yoneda}），构造
$\mathcal{F} \leadsto Q(\mathcal{F})$ 关于 $\mathcal{F}$ 具有函子性。
按构造，$Q$ 是包含函子的右伴随。
当 $\mathcal{F}$ 为拟凝聚层时，$Q(\mathcal{F}) \to \mathcal{F}$
为同构这一事实，在形式上源于包含函子
$\QCoh(\mathcal{O}_X) \to \textit{Mod}(\mathcal{O}_X)$
是全忠实的。
\end{proof}






\section{指向概形的态射}
\label{spaces-properties-section-morphisms-to-schemes}

\noindent
这是《概形》引理 \ref{schemes-lemma-morphism-into-affine} 的类似结果。

\begin{lemma}
\label{spaces-properties-lemma-morphism-to-affine-scheme}
设 $X$ 是 $\mathbf{Z}$ 上的代数空间。
设 $T$ 是仿射概形。映射
$$
\Mor(X, T)
\longrightarrow
\Hom(\Gamma(T, \mathcal{O}_T), \Gamma(X, \mathcal{O}_X))
$$
将 $f$ 映到 $f^\sharp$（在全局截面上）的映射是双射。
\end{lemma}

\begin{proof}
我们构造该映射的逆。令
$\varphi : \Gamma(T, \mathcal{O}_T) \to \Gamma(X, \mathcal{O}_X)$
为环映射。选取表示 $X = U/R$，参见《空间》定义
\ref{spaces-definition-presentation}。由
《概形》引理 \ref{schemes-lemma-morphism-into-affine}，复合映射
$$
\Gamma(T, \mathcal{O}_T) \to \Gamma(X, \mathcal{O}_X) \to
\Gamma(U, \mathcal{O}_U)
$$
对应于唯一的概形态射 $g : U \to T$。由同一引理，两个复合
$R \to U \to T$ 相等。因此得到态射
$f : X = U/R \to T$，使得 $U \to X \to T$ 等于 $g$。按构造，图表
$$
\xymatrix{
\Gamma(U, \mathcal{O}_U) & \Gamma(X, \mathcal{O}_X) \ar[l] \\
& \Gamma(T, \mathcal{O}_T) \ar[lu]^{g^\sharp} \ar[u]^{\varphi}_{f^\sharp}
}
$$
交换。因此，由于 $U \to X$ 是 \'etale 覆盖且 $\mathcal{O}_X$ 是
$X_\etale$ 上的层，$f^\sharp$ 等于 $\varphi$。
$f$ 的唯一性由 $g$ 的唯一性推出。
\end{proof}





\section{自由作用的商}
\label{spaces-properties-section-quotient-by-free}

\noindent
设 $S$ 为概形。
设 $X$ 为 $S$ 上的代数空间。
设 $G$ 为抽象群。
令 $a : G \to \text{Aut}(X)$ 为同态，即 $a$ 是 $G$ 在 $X$ 上的
{\it 作用}。若对每个 $S$ 上的概形 $T$，映射
$$
G \times X(T) \longrightarrow X(T)
$$
是自由的，则称该作用为 {\it 自由}。（我们不能使用《空间》引理
\ref{spaces-lemma-quotient} 中的判据，因为点可能没有定义良好的剩余域。）
在作用自由的情形，我们将把商 $X/G$ 构造为代数空间。这是一般结果
《引导》引理 \ref{bootstrap-lemma-quotient-free-action} 的特殊情形；
该引理将在后文证明。

\begin{lemma}
\label{spaces-properties-lemma-quotient}
设 $S$ 为概形。
设 $X$ 为 $S$ 上的代数空间。
设 $G$ 为在 $X$ 上自由作用的抽象群。
则商层 $X/G$ 是代数空间。
\end{lemma}

\begin{proof}
该陈述意味着，与预层
$$
T \longmapsto X(T)/G
$$
相关联的层 $F$
是一个代数空间。为看出这一点，我们构造一个表示。选取概形 $U$ 及满的
\'etale 态射 $\varphi : U \to X$。令 $V = \coprod_{g \in G} U$，并令
$\psi : V \to X$ 在对应于 $g \in G$ 的分支上等于
$a(g) \circ \varphi$。令 $G$ 通过置换分支作用在 $V$ 上，即
$g_0 \in G$ 通过 $U$ 的恒等态射，将对应于 $g$ 的分支映到对应于
$g_0g$ 的分支。于是 $\psi$ 是 $G$-等变态射，也就是说我们归结为
下一段所处理的情形。

\medskip\noindent
假设 $U$ 上存在 $G$-作用，且 $U \to X$ 是满的、\'etale 且 $G$-等变的。
此时 $R = U \times_X U$ 上诱导出一个 $G$-作用，并且与投影映射
$t, s : R \to U$ 相容。现在我们断言
$$
X/G = U/\coprod\nolimits_{g \in G} R
$$
其中映射
$$
j : \coprod\nolimits_{g \in G} R
\longrightarrow
U \times_S U
$$
由 $(r, g) \mapsto (t(r), g(s(r)))$ 给出。注意 $j$ 是单态射：
若 $(t(r), g(s(r))) = (t(r'), g'(s(r')))$，则
$t(r) = t(r')$，所以 $r$ 与 $r'$ 在 $X$ 中分别经 $s$ 和 $t$ 的像相同；
于是 $g = g'$（因为 $G$ 在 $X$ 上自由作用），从而
$s(r) = s(r')$，最后 $r = r'$（因为 $R$ 是 $U$ 上的等价关系）。
此外，$j$ 是等价关系（细节略）。
两个投影 $\coprod\nolimits_{g \in G} R \to U$ 都是 \'etale 的，因为
$s$ 和 $t$ 是 \'etale 的。因此 $j$ 是 \'etale 等价关系，并且由
《空间》定理 \ref{spaces-theorem-presentation}，
$U/\coprod\nolimits_{g \in G} R$ 是代数空间。
存在映射
$$
U/\coprod\nolimits_{g \in G} R \longrightarrow X/G
$$
由映射 $U \to X$ 诱导。它是层的同构；证明略。
\end{proof}
% Generated mechanically from the frozen spaces-morphisms.tex target.
% Do not edit this generated file; edit the bound locale source instead.

% OK, start here.
%

\setcounter{chapter}{66}
\chapter{代数空间的态射}



\phantomsection
\label{spaces-morphisms-section-phantom}


\section{引言}
\label{spaces-morphisms-section-introduction}

\noindent
本章介绍代数空间的若干类态射。
参考文献为 \cite{Kn}。

\medskip\noindent
目标是把《概形》以及《概形的态射》各章中定义的概形态射类型
分别扩展到代数空间范畴。每种情形略有不同，最好分别处理。





\section{约定}
\label{spaces-morphisms-section-conventions}

\noindent
固定假设所有概形都包含于大 fppf 站点 $\Sch_{fppf}$ 中。
并且，所考虑的每个环 $A$ 都满足 $\Spec(A)$ 同构于该大站点的一个对象。

\medskip\noindent
设 $S$ 为概形，设 $X$ 为 $S$ 上的代数空间。
在本章及下一章中，我们将用 $X \times_S X$ 表示 $X$ 与自身的乘积
（在 $S$ 上代数空间范畴中），而不写作 $X \times X$。





\section{可表示态射的性质}
\label{spaces-morphisms-section-representable}

\noindent
设 $S$ 为概形。
设 $f : X \to Y$ 为代数空间的可表示态射。在《空间》第
\ref{spaces-section-representable-properties} 节中，我们定义了当
$f$ 具有性质 $\mathcal{P}$ 的含义，其中 $\mathcal{P}$ 是概形态射的性质，
且满足下列条件：
\begin{enumerate}
\item 在任意基变换下保持，参见《概形》定义
\ref{schemes-definition-preserved-by-base-change}；并且
\item 在底上是 fppf 局部的，参见《下降》定义
\ref{descent-definition-property-morphisms-local}。
\end{enumerate}
也就是说，在此情形下，我们称 $f$ 具有性质 $\mathcal{P}$，当且仅当对每个
概形 $U$ 及任意态射 $U \to Y$，概形态射 $X \times_Y U \to U$
具有性质 $\mathcal{P}$。

\medskip\noindent
根据《空间》第 \ref{spaces-section-lists} 节的列表，这适用于下列性质：
(1)(a) 闭浸入，
(1)(b) 开浸入，
(1)(c) 拟紧浸入，
(2) 拟紧，
(3) 万有闭，
(4) （拟）分离，
(5) 单态射，
(6) 满射，
(7) 万有单射，
(8) 仿射，
(9) 拟仿射，
(10) （局部）有限型，
(11) （局部）拟有限，
(12) （局部）有限表示，
(13) 相对维数 $d$ 的局部有限型，
(14) 万有开，
(15) 平坦，
(16) syntomic，
(17) 光滑，
(18) 非分歧（相应地，G-非分歧），
(19) \'etale，
(20) proper，
(21) 有限或整，
(22) 有限局部自由，
(23) 万有商映射，
(24) 万有同胚，以及
(25) 浸入。

\medskip\noindent
本章将为不一定可表示的代数空间态射重新定义这些概念。
每次这样做时，我们都会确保新定义与旧定义一致，以避免歧义。

\medskip\noindent
注意，只要 $X$ 是概形，上述定义就适用，因为从概形到代数空间的
态射是可表示的。特别地，当 $X$ 与 $Y$ 都是概形时也适用。
在《空间》引理
\ref{spaces-lemma-morphism-schemes-gives-representable-transformation-property}
中我们已经看到，此时各定义相符，不会产生歧义。

\medskip\noindent
此外，在《空间》引理
\ref{spaces-lemma-base-change-representable-transformations-property}
中我们已经看到，如此定义的代数空间可表示态射的性质在任意代数空间
态射的基变换下保持。最后，在《空间》引理
\ref{spaces-lemma-composition-representable-transformations-property} 及
\ref{spaces-lemma-product-representable-transformations-property}
中我们已经看到，若 $\mathcal{P}$ 在复合下保持（对于上述性质中的
(1)(a)、(1)(b)、(1)(c)、(2)--(25)，但不包括 (13)，均如此），
则可表示态射的乘积保持性质 $\mathcal{P}$，并且可表示态射的复合也保持
性质 $\mathcal{P}$。

\medskip\noindent
下文将使用这些事实；每次使用时，我们简称为“本节”的结果。















\section{分离公理}
\label{spaces-morphisms-section-separation-axioms}

\noindent
列出代数空间态射的对角的一些先验性质是有意义的。

\begin{lemma}
\label{spaces-morphisms-lemma-properties-diagonal}
设 $S$ 为包含于 $\Sch_{fppf}$ 中的概形。
设 $f : X \to Y$ 为 $S$ 上代数空间的态射。
设 $\Delta_{X/Y} : X \to X \times_Y X$ 为对角态射。则
\begin{enumerate}
\item $\Delta_{X/Y}$ 可表示；
\item $\Delta_{X/Y}$ 局部有限型；
\item $\Delta_{X/Y}$ 是单态射；
\item $\Delta_{X/Y}$ 是分离的；以及
\item $\Delta_{X/Y}$ 局部拟有限。
\end{enumerate}
\end{lemma}

\begin{proof}
我们将使用如下事实：$\Delta_{X/S}$ 可表示（按代数空间的定义），
且满足性质 (2)--(5)，参见《空间》引理
\ref{spaces-lemma-properties-diagonal}。
注意，我们有如下因子分解
$$
X
\longrightarrow
X \times_Y X
\longrightarrow
X \times_S X
$$
它是对角 $\Delta_{X/S} : X \to X \times_S X$ 的因子分解。由于
$X \times_Y X \to X \times_S X$ 是单态射，且 $\Delta_{X/S}$ 可表示，
形式上可得 $\Delta_{X/Y}$ 可表示。特别地，其余陈述现在都有意义，
参见第 \ref{spaces-morphisms-section-representable} 节。

\medskip\noindent
选取满的 \'etale 态射 $U \to X$，其中 $U$ 为概形。考虑图表
$$
\xymatrix{
R = U \times_X U \ar[r] \ar[d] &
U \times_Y U \ar[d] \ar[r] &
U \times_S U \ar[d] \\
X \ar[r] & X \times_Y X \ar[r] & X \times_S X
}
$$
两个方块都是笛卡尔的，因而外侧矩形也是笛卡尔的。
顶行由概形组成，竖直箭头是满的 \'etale 态射。由《空间》引理
\ref{spaces-lemma-representable-morphisms-spaces-property}，
$\Delta_{X/Y}$ 的性质 (2)--(5) 等价于 $R \to U \times_Y U$ 的相应性质。
在《空间》引理 \ref{spaces-lemma-properties-diagonal} 的证明中，
我们已经看到 $R \to U \times_S U$ 具有性质 (2)--(5)。
态射 $U \times_Y U \to U \times_S U$ 是概形的单态射。
这些事实推出 $R \to U \times_Y U$ 具有性质 (2)--(5)。

\medskip\noindent
具体地，对于 (3)，注意 $R \to U \times_Y U$ 是单态射，因为复合
$R \to U \times_S U$ 是单态射。对于 (2)，注意 $R \to U \times_Y U$
局部有限型，因为复合 $R \to U \times_S U$ 局部有限型（《态射》引理
\ref{morphisms-lemma-permanence-finite-type}）。局部有限型的单态射
局部拟有限，因为其纤维有限（《态射》引理
\ref{morphisms-lemma-finite-fibre}），故得 (5)。单态射是分离的
（《概形》引理 \ref{schemes-lemma-monomorphism-separated}），故得 (4)。
\end{proof}

\begin{definition}
\label{spaces-morphisms-definition-separated}
设 $S$ 为概形。
设 $f : X \to Y$ 为 $S$ 上代数空间的态射。
设 $\Delta_{X/Y} : X \to X \times_Y X$ 为对角态射。
\begin{enumerate}
\item 称 $f$ 为 {\it 分离}，若 $\Delta_{X/Y}$ 是闭浸入。
\item 称 $f$ 为 {\it 局部分离}\footnote{在文献中，该术语通常指拟分离且局部分离的态射。}，
若 $\Delta_{X/Y}$ 是浸入。
\item 称 $f$ 为 {\it 拟分离}，若 $\Delta_{X/Y}$ 拟紧。
\end{enumerate}
\end{definition}

\noindent
由于 $\Delta_{X/Y}$ 可表示，该定义是有意义的，因而我们知道它具有
定义中所述性质之一是什么意思。下文（引理
\ref{spaces-morphisms-lemma-match-separated}）将看到，该定义与已有的概形态射及可表示
态射的定义相符。

\begin{lemma}
\label{spaces-morphisms-lemma-trivial-implications}
设 $S$ 为概形，设 $f : X \to Y$ 为 $S$ 上代数空间的态射。
若 $f$ 分离，则 $f$ 局部分离且 $f$ 拟分离。
\end{lemma}

\begin{proof}
这由一般原理成立，即《空间》引理
\ref{spaces-lemma-representable-transformations-property-implication}；
因为概形的闭浸入是浸入且拟紧。
\end{proof}

\begin{lemma}
\label{spaces-morphisms-lemma-base-change-separated}
定义 \ref{spaces-morphisms-definition-separated} 中列出的所有分离公理在基变换下保持。
\end{lemma}

\begin{proof}
设 $f : X \to Y$ 与 $Y' \to Y$ 为代数空间的态射。
设 $f' : X' \to Y'$ 为 $f$ 沿 $Y' \to Y$ 的基变换。则
$\Delta_{X'/Y'}$ 是 $\Delta_{X/Y}$ 沿态射
$X' \times_{Y'} X' \to X \times_Y X$ 的基变换。由第
\ref{spaces-morphisms-section-representable} 节的结果，定义
\ref{spaces-morphisms-definition-separated} 中对角的每个相关性质在基变换下保持。
因此引理成立。
\end{proof}

\begin{lemma}
\label{spaces-morphisms-lemma-fibre-product-after-map}
\begin{slogan}
代数空间的 ``魔法图'' 的上方箭头具有良好的浸入类性质；
在分离性假设下，这些性质会加强。
\end{slogan}
设 $S$ 为概形。设 $f : X \to Z$、$g : Y \to Z$ 以及 $Z \to T$
为 $S$ 上代数空间的态射。考虑诱导的态射
$i : X \times_Z Y \to X \times_T Y$。则
\begin{enumerate}
\item $i$ 可表示、局有限型、局拟有限、分离且为单态射；
\item 若 $Z \to T$ 局部分离，则 $i$ 为浸入；
\item 若 $Z \to T$ 分离，则 $i$ 为闭浸入；
\item 若 $Z \to T$ 拟分离，则 $i$ 拟紧。
\end{enumerate}
\end{lemma}

\begin{proof}
由一般范畴论，以下图表
$$
\xymatrix{
X \times_Z Y \ar[r]_i \ar[d] & X \times_T Y \ar[d] \\
Z \ar[r]^-{\Delta_{Z/T}} \ar[r] & Z \times_T Z
}
$$
是纤维积图表。因此 $i$ 是对角态射 $\Delta_{Z/T}$ 的基变换。
于是由引理 \ref{spaces-morphisms-lemma-properties-diagonal} 以及
第 \ref{spaces-morphisms-section-representable} 节的内容可得本引理。
\end{proof}

\begin{lemma}
\label{spaces-morphisms-lemma-semi-diagonal}
\begin{slogan}
代数空间态射的图的性质，是目标的分离性质的结果。
\end{slogan}
设 $S$ 为概形，设 $T$ 为 $S$ 上的代数空间。
设 $g : X \to Y$ 为 $T$ 上代数空间的态射。
考虑 $g$ 的图 $i : X \to X \times_T Y$。则
\begin{enumerate}
\item $i$ 可表示、局有限型、局拟有限、分离且为单态射；
\item 若 $Y \to T$ 局部分离，则 $i$ 为浸入；
\item 若 $Y \to T$ 分离，则 $i$ 为闭浸入；
\item 若 $Y \to T$ 拟分离，则 $i$ 拟紧。
\end{enumerate}
\end{lemma}

\begin{proof}
这是引理 \ref{spaces-morphisms-lemma-fibre-product-after-map} 应用于态射
$X = X \times_Y Y \to X \times_T Y$ 的特殊情形。
\end{proof}

\begin{lemma}
\label{spaces-morphisms-lemma-section-immersion}
设 $S$ 为概形。
设 $f : X \to T$ 为 $S$ 上代数空间的态射。
设 $s : T \to X$ 为 $f$ 的截面（即满足公式
$f \circ s = \text{id}_T$）。则
\begin{enumerate}
\item $s$ 可表示、局有限型、局拟有限、分离且为单态射；
\item 若 $f$ 局部分离，则 $s$ 为浸入；
\item 若 $f$ 分离，则 $s$ 为闭浸入；
\item 若 $f$ 拟分离，则 $s$ 拟紧。
\end{enumerate}
\end{lemma}

\begin{proof}
这是引理 \ref{spaces-morphisms-lemma-semi-diagonal} 应用于 $g = s$ 的特殊情形，
此时态射为 $i = s : T \to T \times_T X$。
\end{proof}

\begin{lemma}
\label{spaces-morphisms-lemma-composition-separated}
定义 \ref{spaces-morphisms-definition-separated} 中列出的所有分离公理
在态射复合下保持。
\end{lemma}

\begin{proof}
设 $f : X \to Y$ 与 $g : Y \to Z$ 为满足所考察公理适用条件的
代数空间态射。对角 $\Delta_{X/Z}$ 是如下复合：
$$
X \longrightarrow X \times_Y X \longrightarrow X \times_Z X.
$$
我们的分离公理定义为要求对角具有某个性质 $\mathcal{P}$。
由上面的引理 \ref{spaces-morphisms-lemma-fibre-product-after-map} 可知，第二个箭头
也具有此性质。因此引理成立，因为具有性质 $\mathcal{P}$ 的（可表示）
态射的复合仍是具有性质 $\mathcal{P}$ 的态射，见第
\ref{spaces-morphisms-section-representable} 节。
\end{proof}

\begin{lemma}
\label{spaces-morphisms-lemma-separated-over-separated}
设 $S$ 为概形。
设 $f : X \to Y$ 为 $S$ 上代数空间的态射。
\begin{enumerate}
\item 若 $Y$ 分离且 $f$ 分离，则 $X$ 分离。
\item 若 $Y$ 拟分离且 $f$ 拟分离，则 $X$ 拟分离。
\item 若 $Y$ 局部分离且 $f$ 局部分离，则 $X$ 局部分离。
\item 若 $Y$ 在 $S$ 上分离且 $f$ 分离，则 $X$ 在 $S$ 上分离。
\item 若 $Y$ 在 $S$ 上拟分离且 $f$ 拟分离，则 $X$ 在 $S$ 上拟分离。
\item 若 $Y$ 在 $S$ 上局部分离且 $f$ 局部分离，则 $X$ 在 $S$ 上局部分离。
\end{enumerate}
\end{lemma}

\begin{proof}
第 (4)、(5)、(6) 部分立即由引理
\ref{spaces-morphisms-lemma-composition-separated} 以及《空间》定义
\ref{spaces-definition-separated} 得到。
将 $X$ 和 $Y$ 看作 $\Spec(\mathbf{Z})$ 上的代数空间后，
第 (1)、(2)、(3) 部分归约为第 (4)、(5)、(6) 部分；见
《空间的性质》定义 \ref{spaces-properties-definition-separated}。
\end{proof}

\begin{lemma}
\label{spaces-morphisms-lemma-compose-after-separated}
设 $S$ 为概形。
设 $f : X \to Y$ 与 $g : Y \to Z$ 为 $S$ 上代数空间的态射。
\begin{enumerate}
\item 若 $g \circ f$ 分离，则 $f$ 分离。
\item 若 $g \circ f$ 局部分离，则 $f$ 局部分离。
\item 若 $g \circ f$ 拟分离，则 $f$ 拟分离。
\end{enumerate}
\end{lemma}

\begin{proof}
考虑对角态射 $g \circ f$ 的分解
$$
X \to X \times_Y X \to X \times_Z X
$$
无论哪种情形，最后一个态射都是单态射。因此对于任意概形 $T$
以及态射 $T \to X \times_Y X$，都有等式
$$
X \times_{(X \times_Y X)} T = X \times_{(X \times_Z X)} T.
$$
所以结论显然。
\end{proof}

\begin{lemma}
\label{spaces-morphisms-lemma-separated-implies-morphism-separated}
设 $S$ 为概形，设 $X$ 为 $S$ 上的代数空间。
\begin{enumerate}
\item 若 $X$ 分离，则 $X$ 在 $S$ 上分离。
\item 若 $X$ 局部分离，则 $X$ 在 $S$ 上局部分离。
\item 若 $X$ 拟分离，则 $X$ 在 $S$ 上拟分离。
\end{enumerate}
设 $f : X \to Y$ 为 $S$ 上代数空间的态射。
\begin{enumerate}
\item[(4)] 若 $X$ 在 $S$ 上分离，则 $f$ 分离。
\item[(5)] 若 $X$ 在 $S$ 上局部分离，则 $f$ 局部分离。
\item[(6)] 若 $X$ 在 $S$ 上拟分离，则 $f$ 拟分离。
\end{enumerate}
\end{lemma}

\begin{proof}
第 (4)、(5)、(6) 部分立即由引理
\ref{spaces-morphisms-lemma-compose-after-separated} 以及《空间》定义
\ref{spaces-definition-separated} 得到。
将 $X$ 和 $Y$ 看作 $\Spec(\mathbf{Z})$ 上的代数空间后，
第 (1)、(2)、(3) 部分由第 (4)、(5)、(6) 部分得到；见
《空间的性质》定义 \ref{spaces-properties-definition-separated}。
\end{proof}

\begin{lemma}
\label{spaces-morphisms-lemma-separated-local}
设 $S$ 为概形。
设 $f : X \to Y$ 为 $S$ 上代数空间的态射。
设 $\mathcal{P}$ 为定义 \ref{spaces-morphisms-definition-separated} 中的任一分离公理。
以下命题等价：
\begin{enumerate}
\item $f$ 具有性质 $\mathcal{P}$；
\item 对每个概形 $Z$ 及态射 $Z \to Y$，$f$ 的基变换
$Z \times_Y X \to Z$ 具有性质 $\mathcal{P}$；
\item 对每个仿射概形 $Z$ 及每个态射 $Z \to Y$，$f$ 的基变换
$Z \times_Y X \to Z$ 具有性质 $\mathcal{P}$；
\item 对每个仿射概形 $Z$ 及每个态射 $Z \to Y$，代数空间
$Z \times_Y X$ 具有性质 $\mathcal{P}$（见《空间的性质》定义
\ref{spaces-properties-definition-separated}）；
\item 存在概形 $V$ 及满射 \'etale 态射 $V \to Y$，使得基变换
$V \times_Y X \to V$ 具有性质 $\mathcal{P}$；
\item 存在 Zariski 覆盖 $Y = \bigcup Y_i$，使得每个态射
$f^{-1}(Y_i) \to Y_i$ 都具有性质 $\mathcal{P}$。
\end{enumerate}
\end{lemma}

\begin{proof}
下文将反复使用引理 \ref{spaces-morphisms-lemma-base-change-separated}，不再另行说明。
特别地，显然 (1) 蕴含 (2)，而 (2) 蕴含 (3)。

\medskip\noindent
证明 (3) 与 (4) 等价。注意，若 $Z$ 为仿射概形，则态射
$Z \to \Spec(\mathbf{Z})$ 作为 $\Spec(\mathbf{Z})$ 上代数空间的态射是分离的。
若 $Z \times_Y X \to Z$ 具有性质 $\mathcal{P}$，则由复合（见引理
\ref{spaces-morphisms-lemma-composition-separated}），$Z \times_Y X \to \Spec(\mathbf{Z})$
具有性质 $\mathcal{P}$。因此代数空间 $Z \times_Y X$ 具有性质 $\mathcal{P}$。
反之，若代数空间 $Z \times_Y X$ 具有性质 $\mathcal{P}$，则
$Z \times_Y X \to \Spec(\mathbf{Z})$ 具有性质 $\mathcal{P}$，从而由引理
\ref{spaces-morphisms-lemma-compose-after-separated} 可知 $Z \times_Y X \to Z$ 具有性质
$\mathcal{P}$。

\medskip\noindent
证明 (3) 蕴含 (5)。假设 (3)。设 $V$ 为概形，且
$V \to Y$ 为满射 \'etale 态射。我们需证明
$V \times_Y X \to V$ 具有性质 $\mathcal{P}$。换言之，需证明态射
$$
V \times_Y X \longrightarrow
(V \times_Y X) \times_V (V \times_Y X) = V \times_Y X \times_Y X
$$
具有相应性质（即为闭浸入、浸入或拟紧）。设
$V = \bigcup V_j$ 为 $V$ 的仿射开覆盖。由假设，每个态射
$$
V_j \times_Y X \longrightarrow V_j \times_Y X \times_Y X
$$
都具有相应性质。由于“闭浸入”、“浸入”、“拟紧浸入”或“拟紧”
在目标上是 Zariski 局部的，且 $V_j$ 覆盖 $V$，故得所需结论。

\medskip\noindent
证明 (5) 蕴含 (1)。设 $V \to Y$ 如 (5) 中所述。
则有纤维积图表
$$
\xymatrix{
V \times_Y X \ar[r] \ar[d] &
X \ar[d] \\
V \times_Y X \times_Y X \ar[r] &
X \times_Y X
}
$$
由假设，左侧竖直箭头是闭浸入、浸入、拟紧浸入或拟紧。
由《空间》引理 \ref{spaces-lemma-descent-representable-transformations-property}
可知，右侧竖直箭头也分别是闭浸入、浸入、拟紧浸入或拟紧。

\medskip\noindent
(1) 蕴含 (6) 是显然的，只需取覆盖 $Y = Y$。
假设 $Y = \bigcup Y_i$ 如 (6) 中所述。取概形 $V_i$ 及满射 \'etale
态射 $V_i \to Y_i$。注意，态射 $V_i \times_Y X \to V_i$ 具有
$\mathcal{P}$，因为它们是态射 $f^{-1}(Y_i) \to Y_i$ 的基变换。
令 $V = \coprod V_i$。则 $V \to Y$ 是 (5) 中的态射（细节略）。
因此 (6) 蕴含 (5)，证明完毕。
\end{proof}

\begin{lemma}
\label{spaces-morphisms-lemma-match-separated}
设 $S$ 为概形。
设 $f : X \to Y$ 为 $S$ 上代数空间的可表示态射。
\begin{enumerate}
\item 态射 $f$ 局部分离。
\item 态射 $f$ 按照上述定义 \ref{spaces-morphisms-definition-separated} 的意义
（拟）分离，当且仅当 $f$ 按照第 \ref{spaces-morphisms-section-representable} 节的意义
（拟）分离。
\end{enumerate}
特别地，若 $f : X \to Y$ 为 $S$ 上概形的态射，则 $f$ 按照定义
\ref{spaces-morphisms-definition-separated} 的意义（拟）分离，当且仅当 $f$ 作为概形态射
（拟）分离。
\end{lemma}

\begin{proof}
这是引理 \ref{spaces-morphisms-lemma-separated-local} 中 (1) 与 (2) 的等价性，
结合任意概形态射都局部分离这一事实得到；见《概形》引理
\ref{schemes-lemma-diagonal-immersion}。
\end{proof}











\section{满射态射}
\label{spaces-morphisms-section-surjective}

\noindent
我们已在第 \ref{spaces-morphisms-section-representable} 节定义了代数空间的可表示态射
为满射的含义。

\begin{lemma}
\label{spaces-morphisms-lemma-surjective-representable}
设 $S$ 为概形。设 $f : X \to Y$ 为 $S$ 上代数空间的可表示态射。
则 $f$ 按照第 \ref{spaces-morphisms-section-representable} 节的意义为满射，当且仅当
$|f| : |X| \to |Y|$ 为满射。
\end{lemma}

\begin{proof}
事实上，若 $f : X \to Y$ 可表示，则它为满射，当且仅当对于每个概形
$T$ 及每个态射 $T \to Y$，$f$ 的基变换
$f_T : T \times_Y X \to T$ 是概形的满射态射，换言之，当且仅当
$|f_T|$ 为满射。由《空间的性质》引理
\ref{spaces-properties-lemma-points-cartesian}，映射
$|T \times_Y X| \to |T| \times_{|Y|} |X|$ 总是满射。
因此，若 $|f| : |X| \to |Y|$ 为满射，则
$|f_T| : |T \times_Y X| \to |T|$ 为满射。反之，若对上述每个
$T \to Y$，$|f_T|$ 都为满射，则取 $T$ 为某个域的谱，便得到
$|X| \to |Y|$ 为满射。
\end{proof}

\noindent
这为下述定义扫清了道路。

\begin{definition}
\label{spaces-morphisms-definition-surjective}
设 $S$ 为概形。设 $f : X \to Y$ 为 $S$ 上代数空间的态射。
若关联拓扑空间的映射 $|f| : |X| \to |Y|$ 为满射，
则称 $f$ 为 {\it 满射}。
\end{definition}

\begin{lemma}
\label{spaces-morphisms-lemma-surjective-local}
设 $S$ 为概形。
设 $f : X \to Y$ 为 $S$ 上代数空间的态射。
以下命题等价：
\begin{enumerate}
\item $f$ 为满射；
\item 对每个概形 $Z$ 及任意态射 $Z \to Y$，态射
$Z \times_Y X \to Z$ 为满射；
\item 对每个仿射概形 $Z$ 及任意态射 $Z \to Y$，态射
$Z \times_Y X \to Z$ 为满射；
\item 存在概形 $V$ 及满射 \'etale 态射 $V \to Y$，使得
$V \times_Y X \to V$ 为满射态射；
\item 存在概形 $U$ 及满射 \'etale 态射
$\varphi : U \to X$，使得复合 $f \circ \varphi$ 为满射；
\item 存在交换图表
$$
\xymatrix{
U \ar[d] \ar[r] & V \ar[d] \\
X \ar[r] & Y
}
$$
其中 $U$、$V$ 为概形，竖直箭头为满射 \'etale 态射，
且上方水平箭头为满射；以及
\item 存在 Zariski 覆盖 $Y = \bigcup Y_i$，使得每个态射
$f^{-1}(Y_i) \to Y_i$ 都为满射。
\end{enumerate}
\end{lemma}

\begin{proof}
略。
\end{proof}

\begin{lemma}
\label{spaces-morphisms-lemma-composition-surjective}
满射态射的复合是满射。
\end{lemma}

\begin{proof}
这直接由定义得到。
\end{proof}

\begin{lemma}
\label{spaces-morphisms-lemma-base-change-surjective}
满射态射的基变换是满射。
\end{lemma}

\begin{proof}
直接由《空间的性质》引理
\ref{spaces-properties-lemma-points-cartesian} 得到。
\end{proof}










\section{开态射}
\label{spaces-morphisms-section-open}

\noindent
对于代数空间的可表示态射，我们已在第 \ref{spaces-morphisms-section-representable} 节
定义了“普遍开”的含义。因此，在给出自然定义之前，先检查它在可表示
情形下与该定义一致。

\begin{lemma}
\label{spaces-morphisms-lemma-characterize-representable-universally-open}
设 $S$ 为概形。设 $f : X \to Y$ 为 $S$ 上代数空间的可表示态射。
以下命题等价：
\begin{enumerate}
\item $f$ 按照第 \ref{spaces-morphisms-section-representable} 节的意义是普遍开；
\item 对每个代数空间态射 $Z \to Y$，拓扑空间的映射
$|Z \times_Y X| \to |Z|$ 是开映射。
\end{enumerate}
\end{lemma}

\begin{proof}
假设 (1)，并设 $Z \to Y$ 如 (2) 中所述。取概形 $V$ 及满射 \'etale
态射 $V \to Y$。由假设，概形态射 $V \times_Y X \to V$ 普遍开。
由《空间的性质》第 \ref{spaces-properties-section-points} 节，
在交换图表
$$
\xymatrix{
|V \times_Y X| \ar[r] \ar[d] & |Z \times_Y X| \ar[d] \\
|V| \ar[r] & |Z|
}
$$
水平箭头是开且满射的，此外
$$
|V \times_Y X| \longrightarrow |V| \times_{|Z|} |Z \times_Y X|
$$
是满射。因此，由于左侧竖直箭头是开映射，右侧竖直箭头也是开映射。
这证明了 (2)。蕴含 (2) $\Rightarrow$ (1) 直接由定义得到。
\end{proof}

\noindent
因此可以采用下述自然定义。

\begin{definition}
\label{spaces-morphisms-definition-open}
设 $S$ 为概形。设 $f : X \to Y$ 为 $S$ 上代数空间的态射。
\begin{enumerate}
\item 若拓扑空间的映射 $|f| : |X| \to |Y|$ 为开映射，
则称 $f$ 为 {\it 开}。
\item 若对每个代数空间态射 $Z \to Y$，拓扑空间的态射
$$
|Z \times_Y X| \to |Z|
$$
为开映射，即基变换 $Z \times_Y X \to Z$ 为开态射，
则称 $f$ 为 {\it 普遍开}。
\end{enumerate}
\end{definition}

\noindent
注意，代数空间的 \'etale 态射是普遍开的；见《空间的性质》定义
\ref{spaces-properties-definition-etale} 以及引理
\ref{spaces-properties-lemma-etale-open} 和
\ref{spaces-properties-lemma-base-change-etale}。

\begin{lemma}
\label{spaces-morphisms-lemma-base-change-universally-open}
代数空间的普遍开态射沿任意代数空间态射的基变换仍是普遍开。
\end{lemma}

\begin{proof}
这直接由定义得到。
\end{proof}

\begin{lemma}
\label{spaces-morphisms-lemma-composition-open}
一对（普遍）开代数空间态射的复合是（普遍）开的。
\end{lemma}

\begin{proof}
略。
\end{proof}

\begin{lemma}
\label{spaces-morphisms-lemma-universally-open-local}
设 $S$ 为概形。设 $f : X \to Y$ 为 $S$ 上代数空间的态射。
以下命题等价：
\begin{enumerate}
\item $f$ 普遍开；
\item 对每个概形 $Z$ 及每个态射 $Z \to Y$，投影
$|Z \times_Y X| \to |Z|$ 为开映射；
\item 对每个仿射概形 $Z$ 及每个态射 $Z \to Y$，投影
$|Z \times_Y X| \to |Z|$ 为开映射；以及
\item 存在概形 $V$ 及满射 \'etale 态射 $V \to Y$，使得
$V \times_Y X \to V$ 为代数空间的普遍开态射；以及
\item 存在 Zariski 覆盖 $Y = \bigcup Y_i$，使得每个态射
$f^{-1}(Y_i) \to Y_i$ 都普遍开。
\end{enumerate}
\end{lemma}

\begin{proof}
我们略去 (1) 蕴含 (2) 及 (2) 蕴含 (3) 的证明。

\medskip\noindent
假设 (3)。取满射 \'etale 态射 $V \to Y$。我们将证明
$V \times_Y X \to V$ 是代数空间的普遍开态射。设 $Z \to V$ 为
从代数空间到 $V$ 的态射。取满射 \'etale 态射 $W \to Z$，其中
$W = \coprod W_i$ 是仿射概形的不交并，见《空间的性质》引理
\ref{spaces-properties-lemma-cover-by-union-affines}。于是有如下交换图表
$$
\xymatrix{
\coprod_i |W_i \times_Y X| \ar@{=}[r] \ar[d] &
|W \times_Y X| \ar[r] \ar[d] &
|Z \times_Y X| \ar[d] \ar@{=}[r] &
|Z \times_V (V \times_Y X)| \ar[ld] \\
\coprod |W_i| \ar@{=}[r] &
|W| \ar[r] &
|Z|
}
$$
我们需证明东南箭头为开映射。中间的水平箭头是满射且开映射
（《空间的性质》引理 \ref{spaces-properties-lemma-etale-open}）。
由假设 (3) 及 $W_i$ 为仿射概形这一事实，左侧竖直箭头是开映射。
因此右侧竖直箭头也是开映射。

\medskip\noindent
假设 $V \to Y$ 如 (4) 中所述。我们将证明 $f$ 普遍开。
设 $Z \to Y$ 为代数空间的态射。考虑图表
$$
\xymatrix{
|(V \times_Y Z) \times_V (V \times_Y X)| \ar@{=}[r] \ar[rd] &
|V \times_Y X| \ar[r] \ar[d] &
|Z \times_Y X| \ar[d] \\
 &
|V \times_Y Z| \ar[r] &
|Z|
}
$$
西南箭头由假设为开映射。水平箭头是满射且开映射，因为相应的
代数空间态射是 \'etale（见《空间的性质》引理
\ref{spaces-properties-lemma-etale-open}）。因此右侧竖直箭头为开映射。

\medskip\noindent
显然，取覆盖 $Y = Y$ 即得 (1) 蕴含 (5)。
假设 $Y = \bigcup Y_i$ 如 (5) 中所述。则对于任意 $Z \to Y$，
可得到相应的 Zariski 覆盖 $Z = \bigcup Z_i$，使得 $f$ 到 $Z_i$
的基变换为开映射。由简单的拓扑论证，$Z \times_Y X \to Z$
为开映射。因此 (1) 成立。
\end{proof}

\begin{lemma}
\label{spaces-morphisms-lemma-space-over-field-universally-open}
设 $S$ 为概形。设 $p : X \to \Spec(k)$ 为 $S$ 上代数空间的态射，
其中 $k$ 为域。则 $p : X \to \Spec(k)$ 普遍开。
\end{lemma}

\begin{proof}
取概形 $U$ 及满射 \'etale 态射 $U \to X$。由《态射》引理
\ref{morphisms-lemma-scheme-over-field-universally-open}，复合
$U \to \Spec(k)$（作为概形态射）普遍开。设 $Z \to \Spec(k)$
为概形态射。则
$U \times_{\Spec(k)} Z \to X \times_{\Spec(k)} Z$ 为满射，见引理
\ref{spaces-morphisms-lemma-base-change-surjective}。因此下列映射中的第一个
$$
|U \times_{\Spec(k)} Z| \to |X \times_{\Spec(k)} Z| \to |Z|
$$
是满射。由于复合由上述结果为开映射，故第二个映射也为开映射。
于是由引理 \ref{spaces-morphisms-lemma-universally-open-local}，$p$ 普遍开。
\end{proof}








\section{商映射态射}
\label{spaces-morphisms-section-submersive}

\noindent
对于代数空间的可表示态射，我们已在第 \ref{spaces-morphisms-section-representable} 节
定义了“普遍商映射”的含义。因此，在给出自然定义之前，先检查它在可表示
情形下与该定义一致。

\begin{lemma}
\label{spaces-morphisms-lemma-characterize-representable-universally-submersive}
设 $S$ 为概形。设 $f : X \to Y$ 为 $S$ 上代数空间的可表示态射。
以下命题等价：
\begin{enumerate}
\item $f$ 按照第 \ref{spaces-morphisms-section-representable} 节的意义是普遍商映射；
\item 对每个代数空间态射 $Z \to Y$，拓扑空间的映射
$|Z \times_Y X| \to |Z|$ 为商映射。
\end{enumerate}
\end{lemma}

\begin{proof}
假设 (1)，并设 $Z \to Y$ 如 (2) 中所述。取概形 $V$ 及满射 \'etale
态射 $V \to Y$。由假设，概形态射 $V \times_Y X \to V$ 是普遍商映射。
由《空间的性质》第 \ref{spaces-properties-section-points} 节，在交换图表
$$
\xymatrix{
|V \times_Y X| \ar[r] \ar[d] & |Z \times_Y X| \ar[d] \\
|V| \ar[r] & |Z|
}
$$
水平箭头是开且满射的，此外
$$
|V \times_Y X| \longrightarrow |V| \times_{|Z|} |Z \times_Y X|
$$
为满射。因此由于左侧竖直箭头为商映射，右侧竖直箭头也为商映射。
这证明了 (2)。蕴含 (2) $\Rightarrow$ (1) 直接由定义得到。
\end{proof}

\noindent
因此可以采用下述自然定义。

\begin{definition}
\label{spaces-morphisms-definition-submersive}
设 $S$ 为概形。
设 $f : X \to Y$ 为 $S$ 上代数空间的态射。
\begin{enumerate}
\item 若连续映射 $|X| \to |Y|$ 为商映射，见《拓扑》定义
\ref{topology-definition-submersive}，则称 $f$ 为 {\it 商映射}\footnote{这与
微分流形的浸没概念有很大不同。}。
\item 若对于每个代数空间态射 $Y' \to Y$，基变换
$Y' \times_Y X \to Y'$ 为商映射，则称 $f$ 为 {\it 普遍商映射}。
\end{enumerate}
\end{definition}

\noindent
注意，商映射态射特别是满射。

\begin{lemma}
\label{spaces-morphisms-lemma-base-change-universally-submersive}
代数空间的普遍商映射态射沿任意代数空间态射的基变换仍是普遍商映射。
\end{lemma}

\begin{proof}
这直接由定义得到。
\end{proof}

\begin{lemma}
\label{spaces-morphisms-lemma-composition-universally-submersive}
一对（普遍）商映射代数空间态射的复合是（普遍）商映射。
\end{lemma}

\begin{proof}
略。
\end{proof}












\section{拟紧态射}
\label{spaces-morphisms-section-quasi-compact}

\noindent
由第 \ref{spaces-morphisms-section-representable} 节，我们知道代数空间的可表示态射
为拟紧的含义。为给出一般代数空间态射的定义，先作如下观察。

\begin{lemma}
\label{spaces-morphisms-lemma-characterize-representable-quasi-compact}
设 $S$ 为概形。
设 $f : X \to Y$ 为 $S$ 上代数空间的可表示态射。
以下命题等价：
\begin{enumerate}
\item $f$ 按照第 \ref{spaces-morphisms-section-representable} 节的意义为拟紧；
\item 对每个拟紧代数空间 $Z$ 及任意态射 $Z \to Y$，代数空间
$Z \times_Y X$ 拟紧。
\end{enumerate}
\end{lemma}

\begin{proof}
假设 (1)，并设 $Z \to Y$ 为代数空间态射，其中 $Z$ 拟紧。由《空间的性质》
定义 \ref{spaces-properties-definition-quasi-compact}，存在拟紧概形 $U$
及满射 \'etale 态射 $U \to Z$。由于 $f$ 可表示且拟紧，按定义
$U \times_Y X$ 是概形，且 $U \times_Y X \to U$ 拟紧。因此
$U \times_Y X$ 是拟紧概形。态射 $U \times_Y X \to Z \times_Y X$
是 \'etale 且满射的（因为它是可表示的满射 \'etale 态射 $U \to Z$
的基变换，见第 \ref{spaces-morphisms-section-representable} 节）。所以按定义
$Z \times_Y X$ 拟紧。

\medskip\noindent
假设 (2)。设 $Z \to Y$ 为态射，其中 $Z$ 为概形。我们需证明
$p : Z \times_Y X \to Z$ 拟紧。设 $U \subset Z$ 为仿射开集。
则 $p^{-1}(U) = U \times_Y Z$，且由假设 (2)，概形
$U \times_Y Z$ 拟紧。因此 $p$ 拟紧，见《概形》第
\ref{schemes-section-quasi-compact} 节。
\end{proof}

\noindent
这促成下述定义。

\begin{definition}
\label{spaces-morphisms-definition-quasi-compact}
设 $S$ 为概形。
设 $f : X \to Y$ 为 $S$ 上代数空间的态射。
若对每个拟紧代数空间 $Z$ 及态射 $Z \to Y$，纤维积
$Z \times_Y X$ 拟紧，则称 $f$ 为 {\it 拟紧}。
\end{definition}

\noindent
由上面的引理 \ref{spaces-morphisms-lemma-characterize-representable-quasi-compact}，
该定义与代数空间可表示态射的既有概念一致。

\begin{lemma}
\label{spaces-morphisms-lemma-quasi-compact-is-quasi-compact}
设 $S$ 为概形。若 $f : X \to Y$ 为 $S$ 上代数空间的拟紧态射，则其
底层拓扑空间映射 $|f| : |X| \to |Y|$ 拟紧。
\end{lemma}

\begin{proof}
设 $V \subset |Y|$ 为拟紧开集。由《空间的性质》引理
\ref{spaces-properties-lemma-open-subspaces}，存在开子空间
$Y' \subset Y$，使 $V = |Y'|$。由《空间的性质》引理
\ref{spaces-properties-lemma-quasi-compact-space}，$Y'$ 是拟紧代数空间，
因而由定义 \ref{spaces-morphisms-definition-quasi-compact}，$X' = Y' \times_Y X$ 是拟紧
代数空间。另一方面，$X' \subset X$ 是开子空间（《空间》引理
\ref{spaces-lemma-base-change-immersions}），且由《空间的性质》引理
\ref{spaces-properties-lemma-points-cartesian}，
$|X'| = |f|^{-1}(|X'|) = |f|^{-1}(V)$。再次应用《空间的性质》引理
\ref{spaces-properties-lemma-quasi-compact-space}，即得 $|X'|$ 是
$|X|$ 的拟紧开集。
\end{proof}

\begin{lemma}
\label{spaces-morphisms-lemma-base-change-quasi-compact}
代数空间的拟紧态射沿任意代数空间态射的基变换仍是拟紧。
\end{lemma}

\begin{proof}
略。提示：纤维积的传递性。
\end{proof}

\begin{lemma}
\label{spaces-morphisms-lemma-composition-quasi-compact}
一对拟紧代数空间态射的复合是拟紧。
\end{lemma}

\begin{proof}
略。提示：纤维积的传递性。
\end{proof}

\begin{lemma}
\label{spaces-morphisms-lemma-surjection-from-quasi-compact}
\begin{slogan}
拟紧代数空间在满射态射下的像是拟紧的。
\end{slogan}
设 $S$ 为概形。
\begin{enumerate}
\item 若 $X \to Y$ 是 $S$ 上代数空间的满射态射，且 $X$ 拟紧，
则 $Y$ 拟紧。
\item 若
$$
\xymatrix{
X \ar[rr]_f \ar[rd]_p & &
Y \ar[dl]^q \\
& Z
}
$$
是 $S$ 上代数空间态射的交换图表，且 $f$ 满射、$p$ 拟紧，则 $q$ 拟紧。
\end{enumerate}
\end{lemma}

\begin{proof}
假设 $X$ 拟紧且 $X \to Y$ 满射。由定义 \ref{spaces-morphisms-definition-surjective}，
$|X| \to |Y|$ 满射。因此由《空间的性质》引理
\ref{spaces-properties-lemma-quasi-compact-space} 以及拟紧空间在连续
映射下的像仍拟紧这一拓扑事实，$Y$ 拟紧；见《拓扑》引理
\ref{topology-lemma-image-quasi-compact}。
设 $f,p,q$ 如 (2) 中所述。设 $T \to Z$ 为态射，且其源是拟紧代数空间。
由假设，$T \times_Z X$ 拟紧。由引理 \ref{spaces-morphisms-lemma-base-change-surjective}，
态射 $T \times_Z X \to T \times_Z Y$ 满射。因此由第 (1) 部分，
$T \times_Z Y$ 也拟紧。故 $q$ 拟紧。
\end{proof}

\begin{lemma}
\label{spaces-morphisms-lemma-descent-quasi-compact}
设 $S$ 为概形。
设 $f : X \to Y$ 为 $S$ 上代数空间的态射。
设 $g : Y' \to Y$ 为代数空间的普遍开满射，且其基变换
$f' : X' \to Y'$ 拟紧。则 $f$ 拟紧。
\end{lemma}

\begin{proof}
设 $Z \to Y$ 为代数空间态射，且 $Z$ 拟紧。由于 $g$ 普遍开且满射，
$Y' \times_Y Z \to Z$ 开且满射。由于 $|Y' \times_Y Z|$ 的每一点
都有拟紧开邻域的基本系（见《空间的性质》引理
\ref{spaces-properties-lemma-space-locally-quasi-compact}），可取
$W \subset |Y' \times_Y Z|$ 中的拟紧开集，使其映到 $Z$ 满射。记
$f'' : W \times_Y X \to W$ 为 $f'$ 沿 $W \to Y'$ 的基变换。
由假设，$W \times_Y X$ 拟紧。由于 $W \to Z$ 满射，
$W \times_Y X \to Z \times_Y X$ 满射。因此由引理
\ref{spaces-morphisms-lemma-surjection-from-quasi-compact}，$Z \times_Y X$ 拟紧。
故 $f$ 拟紧。
\end{proof}

\begin{lemma}
\label{spaces-morphisms-lemma-quasi-compact-local}
\begin{slogan}
代数空间的拟紧态射在拉回下保持，并且在目标上是局部性质。
\end{slogan}
设 $S$ 为概形。
设 $f : X \to Y$ 为 $S$ 上代数空间的态射。
以下命题等价：
\begin{enumerate}
\item $f$ 拟紧；
\item 对每个概形 $Z$ 及任意态射 $Z \to Y$，代数空间态射
$Z \times_Y X \to Z$ 拟紧；
\item 对每个仿射概形 $Z$ 及任意态射 $Z \to Y$，代数空间
$Z \times_Y X$ 拟紧；
\item 存在概形 $V$ 及满射 \'etale 态射 $V \to Y$，使得
$V \times_Y X \to V$ 为拟紧代数空间态射；以及
\item 存在代数空间的满射 \'etale 态射 $Y' \to Y$，使得
$Y' \times_Y X \to Y'$ 为拟紧代数空间态射；以及
\item 存在 Zariski 覆盖 $Y = \bigcup Y_i$，使得每个态射
$f^{-1}(Y_i) \to Y_i$ 拟紧。
\end{enumerate}
\end{lemma}

\begin{proof}
下文将使用引理 \ref{spaces-morphisms-lemma-base-change-quasi-compact}，不再另行说明。
显然 (1) 蕴含 (2)，(2) 蕴含 (3)。假设 (3)。设 $Z$ 为 $S$ 上的拟紧代数空间，
且 $Z \to Y$ 为态射。由《空间的性质》引理
\ref{spaces-properties-lemma-quasi-compact-affine-cover}，存在仿射概形 $U$
及满射 \'etale 态射 $U \to Z$。于是 $U \times_Y X \to Z \times_Y X$
是代数空间的满射态射，见引理 \ref{spaces-morphisms-lemma-base-change-surjective}。
由假设，$|U \times_Y X|$ 拟紧。它满射到 $|Z \times_Y X|$，故由《拓扑》引理
\ref{topology-lemma-image-quasi-compact}，$|Z \times_Y X|$ 拟紧。
这证明了 (3) 蕴含 (1)。

\medskip\noindent
蕴含 (1) $\Rightarrow$ (4) 及 (4) $\Rightarrow$ (5) 是显然的。
蕴含 (5) $\Rightarrow$ (1) 由引理 \ref{spaces-morphisms-lemma-descent-quasi-compact} 以及
代数空间的 \'etale 态射普遍开这一事实得到（见定义
\ref{spaces-morphisms-definition-open} 后的讨论）。

\medskip\noindent
显然，取覆盖 $Y = Y$ 即得 (1) 蕴含 (6)。假设 $Y = \bigcup Y_i$
如 (6) 中所述。设 $Z$ 为仿射概形，且 $Z \to Y$ 为态射。
则存在有限标准仿射覆盖 $Z = Z_1 \cup \ldots \cup Z_n$，使得每个
$Z_j \to Y$ 对某个 $i_j$ 经过 $Y_{i_j}$ 分解。因此代数空间
$$
Z_j \times_Y X = Z_j \times_{Y_{i_j}} f^{-1}(Y_{i_j})
$$
拟紧。由于
$Z \times_Y X = \bigcup_{j = 1, \ldots, n} Z_j \times_Y X$
是 Zariski 覆盖，故
$|Z \times_Y X| = \bigcup_{j = 1, \ldots, n} |Z_j \times_Y X|$
（见《空间的性质》引理 \ref{spaces-properties-lemma-open-subspaces}）
是拟紧空间的有限并，从而拟紧。因此 (6) 蕴含 (3)。
\end{proof}

\noindent
下述引理（以及下一个引理）特别保证：只要 $X$ 是拟紧代数空间，
态射 $X \to \Spec(A)$ 就是拟紧的。

\begin{lemma}
\label{spaces-morphisms-lemma-quasi-compact-permanence}
设 $S$ 为概形。
设 $f : X \to Y$ 与 $g : Y \to Z$ 为 $S$ 上代数空间的态射。
若 $g \circ f$ 拟紧且 $g$ 拟分离，则 $f$ 拟紧。
\end{lemma}

\begin{proof}
这是因为 $f$ 等于复合
$(1, f) : X \to X \times_Z Y \to Y$。第一个映射是拟紧的，
由引理 \ref{spaces-morphisms-lemma-section-immersion}，因为它是拟分离态射
$X \times_Z Y \to X$ 的截面（该态射是 $g$ 的基变换，见引理
\ref{spaces-morphisms-lemma-base-change-separated}）。第二个映射是拟紧的，因为它是
$g \circ f$ 的基变换，见引理 \ref{spaces-morphisms-lemma-base-change-quasi-compact}。
拟紧态射的复合是拟紧的，见引理 \ref{spaces-morphisms-lemma-composition-quasi-compact}。
\end{proof}

\begin{lemma}
\label{spaces-morphisms-lemma-quasi-compact-quasi-separated-permanence}
设 $f : X \to Y$ 为概形 $S$ 上代数空间的态射。
\begin{enumerate}
\item 若 $X$ 拟紧且 $Y$ 拟分离，则 $f$ 拟紧。
\item 若 $X$ 拟紧且拟分离，且 $Y$ 拟分离，则 $f$ 拟紧且拟分离。
\item 拟紧且拟分离代数空间的纤维积拟紧且拟分离。
\end{enumerate}
\end{lemma}

\begin{proof}
第 (1) 部分由引理 \ref{spaces-morphisms-lemma-quasi-compact-permanence}（取
$Z = S = \Spec(\mathbf{Z})$）得到。第 (2) 部分由 (1) 及引理
\ref{spaces-morphisms-lemma-compose-after-separated} 得到。对于 (3)，设 $X \to Y$ 与
$Z \to Y$ 为拟紧且拟分离代数空间的态射。由 (2) 及引理
\ref{spaces-morphisms-lemma-base-change-quasi-compact} 和
\ref{spaces-morphisms-lemma-base-change-separated}，作为 $X \to Y$ 的基变换，
$X \times_Y Z \to Z$ 拟紧且拟分离。因此 $X \times_Y Z$ 作为
$Z$ 上拟紧且拟分离的代数空间本身拟紧且拟分离；见引理
\ref{spaces-morphisms-lemma-separated-over-separated} 和
\ref{spaces-morphisms-lemma-composition-quasi-compact}。
\end{proof}


\section{普遍闭态射}
\label{spaces-morphisms-section-universally-closed}

\noindent
对于代数空间的可表示态射，我们已在第 \ref{spaces-morphisms-section-representable} 节
定义了“普遍闭”的含义。因此，在给出自然定义之前，先检查它在可表示
情形下与该定义一致。

\begin{lemma}
\label{spaces-morphisms-lemma-characterize-representable-universally-closed}
设 $S$ 为概形。设 $f : X \to Y$ 为 $S$ 上代数空间的可表示态射。
以下命题等价：
\begin{enumerate}
\item $f$ 按照第 \ref{spaces-morphisms-section-representable} 节的意义普遍闭；
\item 对每个代数空间态射 $Z \to Y$，拓扑空间的映射
$|Z \times_Y X| \to |Z|$ 为闭映射。
\end{enumerate}
\end{lemma}

\begin{proof}
假设 (1)，并设 $Z \to Y$ 如 (2) 中所述。取概形 $V$ 及满射 \'etale
态射 $V \to Y$。由假设，概形态射 $V \times_Y X \to V$ 普遍闭。
由《空间的性质》第 \ref{spaces-properties-section-points} 节，在交换图表
$$
\xymatrix{
|V \times_Y X| \ar[r] \ar[d] & |Z \times_Y X| \ar[d] \\
|V| \ar[r] & |Z|
}
$$
水平箭头是开且满射的，此外
$$
|V \times_Y X| \longrightarrow |V| \times_{|Z|} |Z \times_Y X|
$$
为满射。因此，由于左侧竖直箭头是闭映射，右侧竖直箭头也是闭映射。
这证明了 (2)。蕴含 (2) $\Rightarrow$ (1) 直接由定义得到。
\end{proof}

\noindent
因此可以采用下述自然定义。

\begin{definition}
\label{spaces-morphisms-definition-closed}
设 $S$ 为概形。设 $f : X \to Y$ 为 $S$ 上代数空间的态射。
\begin{enumerate}
\item 若拓扑空间映射 $|X| \to |Y|$ 为闭映射，则称 $f$ 为 {\it 闭}。
\item 若对每个代数空间态射 $Z \to Y$，拓扑空间的态射
$$
|Z \times_Y X| \to |Z|
$$
为闭映射，即基变换 $Z \times_Y X \to Z$ 为闭态射，则称 $f$ 为
{\it 普遍闭}。
\end{enumerate}
\end{definition}

\begin{lemma}
\label{spaces-morphisms-lemma-base-change-universally-closed}
代数空间的普遍闭态射沿任意代数空间态射的基变换仍是普遍闭。
\end{lemma}

\begin{proof}
这直接由定义得到。
\end{proof}

\begin{lemma}
\label{spaces-morphisms-lemma-composition-universally-closed}
一对（普遍）闭代数空间态射的复合是（普遍）闭的。
\end{lemma}

\begin{proof}
略。
\end{proof}

\begin{lemma}
\label{spaces-morphisms-lemma-universally-closed-local}
设 $S$ 为概形。设 $f : X \to Y$ 为 $S$ 上代数空间的态射。
以下命题等价：
\begin{enumerate}
\item $f$ 普遍闭；
\item 对每个概形 $Z$ 及每个态射 $Z \to Y$，投影
$|Z \times_Y X| \to |Z|$ 为闭映射；
\item 对每个仿射概形 $Z$ 及每个态射 $Z \to Y$，投影
$|Z \times_Y X| \to |Z|$ 为闭映射；
\item 存在概形 $V$ 及满射 \'etale 态射 $V \to Y$，使得
$V \times_Y X \to V$ 为代数空间的普遍闭态射；以及
\item 存在 Zariski 覆盖 $Y = \bigcup Y_i$，使得每个态射
$f^{-1}(Y_i) \to Y_i$ 都普遍闭。
\end{enumerate}
\end{lemma}

\begin{proof}
我们略去 (1) 蕴含 (2) 及 (2) 蕴含 (3) 的证明。

\medskip\noindent
假设 (3)。取满射 \'etale 态射 $V \to Y$。我们将证明
$V \times_Y X \to V$ 是代数空间的普遍闭态射。设 $Z \to V$ 为
从代数空间到 $V$ 的态射。取满射 \'etale 态射 $W \to Z$，其中
$W = \coprod W_i$ 是仿射概形的不交并，见《空间的性质》引理
\ref{spaces-properties-lemma-cover-by-union-affines}。于是有如下交换图表
$$
\xymatrix{
\coprod_i |W_i \times_Y X| \ar@{=}[r] \ar[d] &
|W \times_Y X| \ar[r] \ar[d] &
|Z \times_Y X| \ar[d] \ar@{=}[r] &
|Z \times_V (V \times_Y X)| \ar[ld] \\
\coprod |W_i| \ar@{=}[r] &
|W| \ar[r] &
|Z|
}
$$
我们需证明东南箭头为闭映射。中间的水平箭头是满射且开映射
（《空间的性质》引理 \ref{spaces-properties-lemma-etale-open}）。
由假设 (3) 及 $W_i$ 为仿射概形这一事实，左侧竖直箭头为闭映射。
因此右侧竖直箭头也为闭映射。

\medskip\noindent
假设 (4)。我们将证明 $f$ 普遍闭。设 $Z \to Y$ 为代数空间的态射。
考虑图表
$$
\xymatrix{
|(V \times_Y Z) \times_V (V \times_Y X)| \ar@{=}[r] \ar[rd] &
|V \times_Y X| \ar[r] \ar[d] &
|Z \times_Y X| \ar[d] \\
 &
|V \times_Y Z| \ar[r] &
|Z|
}
$$
西南箭头由假设为闭映射。水平箭头是满射且开映射，因为相应的代数空间
态射是 \'etale（见《空间的性质》引理
\ref{spaces-properties-lemma-etale-open}）。因此右侧竖直箭头为闭映射。

\medskip\noindent
显然，取覆盖 $Y = Y$ 即得 (1) 蕴含 (5)。假设 $Y = \bigcup Y_i$
如 (5) 中所述。则对于任意 $Z \to Y$，可得到相应的 Zariski 覆盖
$Z = \bigcup Z_i$，使得 $f$ 到 $Z_i$ 的基变换为闭映射。由简单的
拓扑论证，$Z \times_Y X \to Z$ 为闭映射。因此 (1) 成立。
\end{proof}

\begin{example}
\label{spaces-morphisms-example-strange-universally-closed}
普遍闭态射的奇异例子。
设 $\mathbf{Q} \subset k$ 为特征零的域。令
$X = \mathbf{A}^1_k/\mathbf{Z}$，如《空间》例
\ref{spaces-example-affine-line-translation}。我们断言结构态射
$p : X \to \Spec(k)$ 普遍闭。
事实上，若 $Z/k$ 为概形，且 $T \subset |X \times_k Z|$ 为闭集，
则 $T$ 对应于 $T' \subset |\mathbf{A}^1 \times Z|$ 中的
$\mathbf{Z}$-不变闭集。容易看出，这意味着 $T'$ 是 $Z$ 的某个子集
$T''$ 的逆像。由《态射》引理
\ref{morphisms-lemma-fpqc-quotient-topology}，$T'' \subset Z$ 为闭集。
显然 $T''$ 是 $T$ 的像。因此由引理 \ref{spaces-morphisms-lemma-universally-closed-local}，
$p$ 普遍闭。
\end{example}

\begin{lemma}
\label{spaces-morphisms-lemma-universally-closed-quasi-compact}
设 $S$ 为概形。
$S$ 上代数空间的普遍闭态射是拟紧的。
\end{lemma}

\begin{proof}
该证明重复概形情形的证明，见《态射》引理
\ref{morphisms-lemma-universally-closed-quasi-compact}。
设 $f : X \to Y$ 为 $S$ 上代数空间的态射，并假设 $f$ 不拟紧。
我们的目标是证明 $f$ 不普遍闭。由引理 \ref{spaces-morphisms-lemma-quasi-compact-local}，
存在仿射概形 $Z$ 及态射 $Z \to Y$，使得 $Z \times_Y X \to Z$
不拟紧。为达到目标，只需证明 $Z \times_Y X \to Z$ 不普遍闭，
因此可假设 $Y = \Spec(B)$，其中 $B$ 为某个环。

\medskip\noindent
写成 $X = \bigcup_{i \in I} X_i$，其中 $X_i$ 是 $X$ 的拟紧开子空间。
例如，取满射 \'etale 态射 $U \to X$，其中 $U$ 为概形；取仿射开覆盖
$U = \bigcup U_i$，并令 $X_i \subset X$ 为 $U_i$ 的像。稍后将使用
$X_i \to Y$ 拟紧这一事实，见引理 \ref{spaces-morphisms-lemma-quasi-compact-permanence}。
令 $T = \Spec(B[a_i ; i \in I])$，令 $T_i = D(a_i) \subset T$。
令 $Z \subset T \times_Y X$ 为约化闭子空间，其底层点集为
$|T \times_Y Z| \setminus \bigcup_{i \in I} |T_i \times_Y X_i|$，见
《空间的性质》引理 \ref{spaces-properties-lemma-reduced-closed-subspace}。
（注意，$T_i \times_Y X_i$ 是 $T \times_Y X$ 的开子空间，因为
$T_i \to T$ 与 $X_i \to X$ 是开浸入；见《空间》引理
\ref{spaces-lemma-base-change-immersions} 和
\ref{spaces-lemma-composition-immersions}。）图表如下：
$$
\xymatrix{
Z \ar[r] \ar[rd] &
T \times_Y X \ar[d]^{f_T} \ar[r]_q &
X \ar[d]^f \\
& T \ar[r]^p & Y
}
$$
只需证明像 $f_T(|Z|)$ 在 $|T|$ 中不是闭集。

\medskip\noindent
我们断言，存在点 $y \in Y$，使得 $Y$ 中不存在 $y$ 的仿射开邻域
$V$ 满足 $X_V$ 拟紧。否则，我们可以用有限多个这样的 $V$ 覆盖 $Y$；
对每个 $V$，由引理 \ref{spaces-morphisms-lemma-quasi-compact-permanence}，态射
$Y_V \to V$ 拟紧，继而由引理 \ref{spaces-morphisms-lemma-quasi-compact-local} 得到
$f$ 拟紧，矛盾。固定断言中的一个 $y \in Y$。

\medskip\noindent
令 $t \in T$ 为位于 $y$ 上方且满足 $\kappa(t) = \kappa(y)$ 的点，
并使得对所有 $i$，$a_i = 1$ 于 $\kappa(t)$ 中。若存在 $z \in |Z|$
满足 $f_T(z) = t$，则对某个 $i$ 有 $q(t) \in X_i$。于是按 $Z$ 的构造，
$f_T(z) \not \in T_i$，这与按构造 $t \in T_i$ 矛盾。因此
$t \in |T| \setminus f_T(|Z|)$。

\medskip\noindent
假设 $f_T(|Z|)$ 在 $|T|$ 中闭。则存在元素
$g \in B[a_i; i \in I]$，使 $f_T(|Z|) \subset V(g)$ 但 $t \not \in V(g)$。
于是 $g$ 在 $\kappa(t)$ 中的像非零。特别地，$g$ 的某个系数在
$\kappa(y)$ 中的像非零。因此，该系数在 $y$ 的某个仿射开邻域 $V$ 上可逆。
令 $J$ 为满足 $j \in I$ 且变量 $a_j$ 出现在 $g$ 中的有限集合。
由于 $X_V$ 不拟紧而每个 $X_{i, V}$ 拟紧，我们可以取点
$x \in |X_V| \setminus \bigcup_{j \in J} |X_{j, V}|$。
换言之，$x \in |X| \setminus \bigcup_{j \in J} |X_j|$，且 $x$ 位于某个
$v \in V$ 上方。由于 $g$ 有一个在 $V$ 上可逆的系数，我们可以找到位于
$v$ 上方的点 $t' \in T$，使 $t' \not \in V(g)$ 且对所有 $i \notin J$
都有 $t' \in V(a_i)$。这是因为
$V(a_i; i \in I \setminus J) = \Spec(B[a_j; j\in J])$，
且该概形中位于 $v$ 上方的点集与
$\Spec(\kappa(v)[a_j; j \in J])$ 双射，而按构造 $g$ 限制为该多项式环
中的非零元素。换言之，对于每个 $i \not \in J$，$t' \not \in T_i$。
由《空间的性质》引理 \ref{spaces-properties-lemma-points-cartesian}，
可找到 $X \times_Y T$ 中的点 $z$，其映射到 $x \in X$ 以及 $t' \in T$。
由于当 $j \in J$ 时 $x \not \in |X_j|$，且当
$i \in I \setminus J$ 时 $t' \not \in T_i$，可知 $z \in |Z|$。
另一方面，$f_T(z) = t' \not \in V(g)$，这与 $f_T(Z) \subset V(g)$ 矛盾。
因此假设“$f_T(|Z|)$ 闭”错误，故如所需，$f_T$ 非闭。
\end{proof}

\noindent
分离代数空间在满射普遍闭态射下的目标仍然分离。

\begin{lemma}
\label{spaces-morphisms-lemma-image-universally-closed-separated}
设 $S$ 为概形，设 $B$ 为 $S$ 上的代数空间。
设 $f : X \to Y$ 为 $B$ 上代数空间的满射普遍闭态射。
\begin{enumerate}
\item 若 $X$ 拟分离，则 $Y$ 拟分离。
\item 若 $X$ 分离，则 $Y$ 分离。
\item 若 $X$ 在 $B$ 上拟分离，则 $Y$ 在 $B$ 上拟分离。
\item 若 $X$ 在 $B$ 上分离，则 $Y$ 在 $B$ 上分离。
\end{enumerate}
\end{lemma}

\begin{proof}
第 (1)、(2) 部分是 $S = B = \Spec(\mathbf{Z})$ 时第 (3)、(4) 部分的
结果（见《空间的性质》定义 \ref{spaces-properties-definition-separated}）。
考虑交换图表
$$
\xymatrix{
X \ar[d] \ar[rr]_{\Delta_{X/B}} & & X \times_B X \ar[d] \\
Y \ar[rr]^{\Delta_{Y/B}} & & Y \times_B Y
}
$$
左侧竖直箭头为满射（即普遍满射）。右侧竖直箭头是普遍闭的，
因为它是普遍闭态射
$X \times_B X \to X \times_B Y \to Y \times_B Y$ 的复合。因此它也拟紧，
见引理 \ref{spaces-morphisms-lemma-universally-closed-quasi-compact}。

\medskip\noindent
假设 $X$ 在 $B$ 上拟分离，即 $\Delta_{X/B}$ 拟紧。若 $Z$ 拟紧且
$Z \to Y \times_B Y$ 为态射，则
$Z \times_{Y \times_B Y} X \to Z \times_{Y \times_B Y} Y$
满射，并且由上述说明 $Z \times_{Y \times_B Y} X$ 拟紧。
因此 $\Delta_{Y/B}$ 拟紧，即 $Y$ 在 $B$ 上拟分离。

\medskip\noindent
假设 $X$ 在 $B$ 上分离，即 $\Delta_{X/B}$ 为闭浸入。若 $Z$ 为仿射概形，
且 $Z \to Y \times_B Y$ 为态射，则
$Z \times_{Y \times_B Y} X \to Z \times_{Y \times_B Y} Y$
满射，并且由上述说明 $Z \times_{Y \times_B Y} X \to Z$ 普遍闭。
因此 $\Delta_{Y/B}$ 普遍闭。由此 $\Delta_{Y/B}$ 可表示、局有限型且为
单态射（见引理 \ref{spaces-morphisms-lemma-properties-diagonal}），并且普遍闭，
所以是闭浸入；见《\'Etale 态射》引理
\ref{etale-lemma-characterize-closed-immersion}（也见抽象原理
《空间》引理 \ref{spaces-lemma-representable-transformations-property-implication}）。
于是 $Y$ 在 $B$ 上分离。
\end{proof}













\section{单态射}
\label{spaces-morphisms-section-monomorphisms}

\noindent
按照第 \ref{spaces-morphisms-section-representable} 节，代数空间的可表示态射
$X \to Y$ 是单态射，是指对于每个概形 $Z$ 及态射 $Z \to Y$，
态射 $Z \times_Y X \to Z$ 可由概形的单态射表示。
这恰好意味着 $Z \times_Y X \to Z$ 是 $(\Sch/S)_{fppf}$ 上层的
单射。由于这应对所有 $Z$ 及所有 $Z \to Y$ 成立，因而等价于
$X \to Y$ 是 $(\Sch/S)_{fppf}$ 上层的单射。因此，对（可能不可表示的）
代数空间态射定义单态射如下：\footnote{我们不知道代数空间的任意单态射
是否可表示。讨论见《空间态射的更多内容》第
\ref{spaces-more-morphisms-section-monomorphisms} 节。}

\begin{definition}
\label{spaces-morphisms-definition-monomorphism}
设 $S$ 为概形。
若 $S$ 上的代数空间态射是层的单射，即 $(\Sch/S)_{fppf}$ 上层范畴中的
单态射，则称其为 {\it 单态射}。
\end{definition}

\noindent
下述引理表明，这也等价于它是 $S$ 上代数空间范畴中的单态射。

\begin{lemma}
\label{spaces-morphisms-lemma-monomorphism}
设 $S$ 为概形。
设 $j : X \to Y$ 为 $S$ 上代数空间的态射。
以下命题等价：
\begin{enumerate}
\item $j$ 是单态射（如定义 \ref{spaces-morphisms-definition-monomorphism}）；
\item $j$ 是 $S$ 上代数空间范畴中的单态射；
\item 对角态射 $\Delta_{X/Y} : X \to X \times_Y X$ 是同构。
\end{enumerate}
\end{lemma}

\begin{proof}
注意，$X \times_Y X$ 既是 $(\Sch/S)_{fppf}$ 上层范畴中的纤维积，
也是 $S$ 上代数空间范畴中的纤维积，见《空间》引理
\ref{spaces-lemma-fibre-product-spaces}。(1) 与 (3) 的等价性是任意
网站上层单射的一般刻画。(2) 与 (3) 的等价性是具有纤维积的任意范畴中
对单态射的刻画。
\end{proof}

\begin{lemma}
\label{spaces-morphisms-lemma-monomorphism-separated}
代数空间的单态射是分离的。
\end{lemma}

\begin{proof}
这是因为同构是闭浸入，且可应用上面的引理
\ref{spaces-morphisms-lemma-monomorphism}。
\end{proof}

\begin{lemma}
\label{spaces-morphisms-lemma-composition-monomorphism}
单态射的复合是单态射。
\end{lemma}

\begin{proof}
这是因为层单射的复合仍为单射。
\end{proof}

\begin{lemma}
\label{spaces-morphisms-lemma-base-change-monomorphism}
单态射的基变换是单态射。
\end{lemma}

\begin{proof}
这是层范畴中纤维积的一般事实。
\end{proof}

\begin{lemma}
\label{spaces-morphisms-lemma-monomorphism-local}
设 $S$ 为概形。
设 $f : X \to Y$ 为 $S$ 上代数空间的态射。
以下命题等价：
\begin{enumerate}
\item $f$ 是单态射；
\item 对每个概形 $Z$ 及态射 $Z \to Y$，$f$ 的基变换
$Z \times_Y X \to Z$ 是单态射；
\item 对每个仿射概形 $Z$ 及每个态射 $Z \to Y$，$f$ 的基变换
$Z \times_Y X \to Z$ 是单态射；
\item 存在概形 $V$ 及满射 \'etale 态射 $V \to Y$，使得基变换
$V \times_Y X \to V$ 是单态射；以及
\item 存在 Zariski 覆盖 $Y = \bigcup Y_i$，使得每个态射
$f^{-1}(Y_i) \to Y_i$ 是单态射。
\end{enumerate}
\end{lemma}

\begin{proof}
下文将不再另行说明单态射的基变换仍是单态射这一事实，见引理
\ref{spaces-morphisms-lemma-base-change-monomorphism}。特别地，显然
(1) $\Rightarrow$ (2) $\Rightarrow$ (3) $\Rightarrow$ (4)，其中 (4) 中
$V$ 取为与 $Y$ \'etale 的仿射概形的不交并（见《空间的性质》引理
\ref{spaces-properties-lemma-cover-by-union-affines}）。设 $V$ 为概形，
且 $V \to Y$ 为满射 \'etale 态射。若 $V \times_Y X \to V$ 是单态射，
则 $X \to Y$ 也是单态射。事实上，给定任意层的笛卡尔图表
$$
\vcenter{
\xymatrix{
\mathcal{F} \ar[r]_a \ar[d]_b & \mathcal{G} \ar[d]^c \\
\mathcal{H} \ar[r]^d & \mathcal{I}
}
}
\quad
\quad
\mathcal{F} = \mathcal{H} \times_\mathcal{I} \mathcal{G}
$$
若 $c$ 是层的满射且 $a$ 为单射，则 $d$ 也是单射。因此 (4) 蕴含 (1)。
(5) 与 (1) 的等价性证明略。
\end{proof}

\begin{lemma}
\label{spaces-morphisms-lemma-immersions-monomorphisms}
代数空间的浸入是单态射。
特别地，任意浸入都是分离的。
\end{lemma}

\begin{proof}
设 $f : X \to Y$ 为代数空间的浸入。对于任意态射 $Z \to Y$，其中 $Z$
可表示，基变换 $Z \times_Y X \to Z$ 是概形的浸入，因而为单态射，
见《概形》引理 \ref{schemes-lemma-immersions-monomorphisms}。
因此 $f$ 可表示且为单态射。
\end{proof}

\noindent
我们将在《良好空间》引理
\ref{decent-spaces-lemma-monomorphism-toward-disjoint-union-dim-0-rings}
中加强下述引理。

\begin{lemma}
\label{spaces-morphisms-lemma-monomorphism-toward-field}
设 $S$ 为概形。设 $k$ 为域，且 $Z \to \Spec(k)$ 为 $S$ 上代数空间的
单态射。则或者 $Z = \emptyset$，或者 $Z = \Spec(k)$。
\end{lemma}

\begin{proof}
由引理 \ref{spaces-morphisms-lemma-monomorphism-separated} 和
\ref{spaces-morphisms-lemma-separated-over-separated}，可知 $Z$ 是分离代数空间。因此存在
概形形式的开稠密子空间 $Z' \subset Z$，见《空间的性质》命题
\ref{spaces-properties-proposition-locally-quasi-separated-open-dense-scheme}。
由《概形》引理 \ref{schemes-lemma-mono-towards-spec-field}，或者
$Z' = \emptyset$，或者 $Z' \cong \Spec(k)$。在第一种情形，得到
$Z = \emptyset$；在第二种情形，由于 $Z \to \Spec(k)$ 是在 $Z'$ 上为
同构的单态射，得到 $Z' = Z = \Spec(k)$。
\end{proof}

\begin{lemma}
\label{spaces-morphisms-lemma-monomorphism-injective-points}
设 $S$ 为概形。若 $X \to Y$ 为 $S$ 上代数空间的单态射，则
$|X| \to |Y|$ 为单射。
\end{lemma}

\begin{proof}
直接由定义得到。
\end{proof}















\section{拟凝聚层的推前}
\label{spaces-morphisms-section-pushforward}

\noindent
我们先证明一个简单引理，将代数空间态射的模层推前与概形态射的模层
推前联系起来。

\begin{lemma}
\label{spaces-morphisms-lemma-compute-pushforward}
设 $S$ 为概形。
设 $f : X \to Y$ 为 $S$ 上代数空间的态射。
设 $U \to X$ 为从概形到 $X$ 的满射 \'etale 态射。
令 $R = U \times_X U$，并按通常记号令 $t, s : R \to U$ 为投影态射。
记 $a : U \to Y$ 与 $b : R \to Y$ 为诱导态射。对于
$\textit{Mod}(\mathcal{O}_X)$ 中任意对象 $\mathcal{F}$，存在正合序列
$$
0 \to f_*\mathcal{F} \to a_*(\mathcal{F}|_U) \to b_*(\mathcal{F}|_R)
$$
其中第二个箭头为差 $t^* - s^*$。
\end{lemma}

\begin{proof}
我们也把 $\mathcal{F}$ 记作其延拓到 $X_{spaces, \etale}$ 上的模层，见
《空间的性质》注记 \ref{spaces-properties-remark-explain-equivalence}。
设 $V \to Y$ 为 $Y_\etale$ 中的对象。则 $V \times_Y X$ 是
$X_{spaces, \etale}$ 中的对象，且按定义
$f_*\mathcal{F}(V) = \mathcal{F}(V \times_Y X)$。由于 $U \to X$
满射 \'etale，故 $\{V \times_Y U \to V \times_Y X\}$ 是覆盖。
此外有
$(V \times_Y U) \times_X (V \times_Y U) = V \times_Y R$.
因此由 $\mathcal{F}$ 在 $X_{spaces, \etale}$ 上的层条件，有短正合序列
$$
0 \to \mathcal{F}(V \times_Y X)
\to \mathcal{F}(V \times_Y U) \to \mathcal{F}(V \times_Y R)
$$
其中第二个箭头是沿 $t$ 或 $s$ 限制所得的差。该正合序列关于 $V$
函子性，故得到本引理。
\end{proof}

\noindent
设 $S$ 为概形。设 $f : X \to Y$ 为 $S$ 上可表示代数空间 $X$ 与 $Y$
之间的拟紧且拟分离态射。由《下降》命题
\ref{descent-proposition-equivalence-quasi-coherent-functorial}，函子
$f_* : \QCoh(\mathcal{O}_X) \to \QCoh(\mathcal{O}_Y)$
在将 $X$ 与 $Y$ 看作概形时与通常函子一致。

\medskip\noindent
更一般地，设 $f : X \to Y$ 为 $S$ 上代数空间的可表示、拟紧且拟分离态射。
设 $V$ 为概形，且 $V \to Y$ 为满射 \'etale 态射。令
$U = V \times_Y X$，令 $f' : U \to V$ 为 $f$ 的基变换。则对于任意
拟凝聚 $\mathcal{O}_X$-模 $\mathcal{F}$，有
\begin{equation}
\label{spaces-morphisms-equation-representable-pushforward}
f'_*(\mathcal{F}|_U) = (f_*\mathcal{F})|_V,
\end{equation}
见《空间的性质》引理
\ref{spaces-properties-lemma-pushforward-etale-base-change-modules}。
又由于 $f' : U \to V$ 是概形的拟紧且拟分离态射，由上一段的注记，
我们可以将 $\mathcal{F}|_U$ 看作概形 $U$ 上的拟凝聚层、将 $f'$ 看作
概形态射来计算 $f'_*(\mathcal{F}|_U)$。下文经常使用这一点，不再说明。

\medskip\noindent
下一层一般性是考虑任意拟紧且拟分离的代数空间态射。

\begin{lemma}
\label{spaces-morphisms-lemma-pushforward}
设 $S$ 为概形。
设 $f : X \to Y$ 为 $S$ 上代数空间的态射。
若 $f$ 拟紧且拟分离，则 $f_*$ 将拟凝聚 $\mathcal{O}_X$-模变为
拟凝聚 $\mathcal{O}_Y$-模。
\end{lemma}

\begin{proof}
设 $\mathcal{F}$ 为 $X$ 上的拟凝聚层。需证明 $f_*\mathcal{F}$ 是 $Y$ 上的
拟凝聚层。只需证明对于任意仿射概形 $V$ 及 \'etale 态射 $V \to Y$，
$f_*\mathcal{F}$ 到 $V$ 的限制是拟凝聚的；见《空间的性质》引理
\ref{spaces-properties-lemma-characterize-quasi-coherent}。
令 $f' : V \times_Y X \to V$ 为 $f$ 沿 $V \to Y$ 的基变换。注意由引理
\ref{spaces-morphisms-lemma-base-change-quasi-compact} 和
\ref{spaces-morphisms-lemma-base-change-separated}，$f'$ 也拟紧且拟分离。
由 (\ref{spaces-morphisms-equation-representable-pushforward})，$f_*\mathcal{F}$ 到 $V$ 的
限制是 $\mathcal{F}$ 到 $V \times_Y X$ 的限制的 $f'_*$。因此可将 $f$
替换为 $f'$，并假设 $Y$ 是仿射概形。

\medskip\noindent
假设 $Y$ 为仿射概形。由于 $f$ 拟紧，$X$ 拟紧。因此可取仿射概形 $U$
及满射 \'etale 态射 $U \to X$，见《空间的性质》引理
\ref{spaces-properties-lemma-quasi-compact-affine-cover}。
由引理 \ref{spaces-morphisms-lemma-compute-pushforward}，有正合序列
$$
0 \to f_*\mathcal{F} \to a_*(\mathcal{F}|_U) \to b_*(\mathcal{F}|_R).
$$
其中 $R = U \times_X U$。由于 $X \to Y$ 拟分离，
$R \to U \times_Y U$ 是拟紧单态射。这意味着 $R$ 是拟紧分离概形
（此时 $U$ 与 $Y$ 均为仿射）。因此 $a : U \to Y$ 与 $b : R \to Y$
是概形的拟紧且拟分离态射。于是由《下降》命题
\ref{descent-proposition-equivalence-quasi-coherent-functorial}，
层 $a_*(\mathcal{F}|_U)$ 与 $b_*(\mathcal{F}|_R)$ 拟凝聚（也见本引理前的
讨论）。这说明 $f_*\mathcal{F}$ 是拟凝聚模的核，因而自身拟凝聚；见
《空间的性质》引理 \ref{spaces-properties-lemma-properties-quasi-coherent}。
\end{proof}

\noindent
高阶推前见《空间的上同调》第
\ref{spaces-cohomology-section-higher-direct-image} 节。









\section{浸入}
\label{spaces-morphisms-section-immersions}

\noindent
代数空间的开浸入、闭浸入及局部闭浸入已在《空间》第
\ref{spaces-section-Zariski} 节定义。也就是说，若代数空间态射可表示，
且按照第 \ref{spaces-morphisms-section-representable} 节的意义为闭浸入（分别为开浸入、浸入），
则称其为 {\it 闭浸入}（分别为 {\it 开浸入}、{\it 浸入}）。

\medskip\noindent
特别地，这些类型的态射在 $S$ 上代数空间范畴中的基变换与态射复合下保持；
见《空间》引理 \ref{spaces-lemma-composition-immersions} 和
\ref{spaces-lemma-base-change-immersions}。

\begin{lemma}
\label{spaces-morphisms-lemma-closed-immersion-local}
设 $S$ 为概形。设 $f : X \to Y$ 为 $S$ 上代数空间的态射。
以下命题等价：
\begin{enumerate}
\item $f$ 是闭浸入（分别为开浸入、浸入）；
\item 对每个概形 $Z$ 及任意态射 $Z \to Y$，态射
$Z \times_Y X \to Z$ 是闭浸入（分别为开浸入、浸入）；
\item 对每个仿射概形 $Z$ 及任意态射 $Z \to Y$，态射
$Z \times_Y X \to Z$ 是闭浸入（分别为开浸入、浸入）；
\item 存在概形 $V$ 及满射 \'etale 态射 $V \to Y$，使得
$V \times_Y X \to V$ 是闭浸入（分别为开浸入、浸入）；以及
\item 存在 Zariski 覆盖 $Y = \bigcup Y_i$，使得每个态射
$f^{-1}(Y_i) \to Y_i$ 是闭浸入（分别为开浸入、浸入）。
\end{enumerate}
\end{lemma}

\begin{proof}
由于闭浸入（分别为开浸入、浸入）的基变换仍为同类态射，显然
(1) 蕴含 (2)，(2) 蕴含 (3)。又 (3) 蕴含 (4)，因为可取 $V$ 为仿射概形
的不交并；见《空间的性质》引理
\ref{spaces-properties-lemma-cover-by-union-affines}。

\medskip\noindent
假设 $V \to Y$ 如 (4) 中所述。令 $\mathcal{P}$ 为概形态射的性质
“闭浸入（分别为开浸入、浸入）”。注意，性质 $\mathcal{P}$ 在任意基变换下
保持，且在基上是 fppf 局部的（见第 \ref{spaces-morphisms-section-representable} 节）。
此外，$\mathcal{P}$ 类型的态射是分离且局拟有限的（在三种情形中均如此；见
《概形》引理 \ref{schemes-lemma-immersions-monomorphisms} 以及
《态射》引理 \ref{morphisms-lemma-immersion-locally-quasi-finite}）。
因此由《态射的更多内容》引理
\ref{more-morphisms-lemma-separated-locally-quasi-finite-morphisms-fppf-descend}，
$\mathcal{P}$ 类型态射满足 fppf 覆盖下的下降。于是可应用《空间》引理
\ref{spaces-lemma-morphism-sheaves-with-P-effective-descent-etale}，得到
$X \to Y$ 可表示且具有性质 $\mathcal{P}$，即 (1) 成立。

\medskip\noindent
由 $\mathcal{P}$ 在目标上是 Zariski 局部性质，(1) 与 (5) 等价（上面已知
$\mathcal{P}$ 实际上在目标上是 fppf 局部的）。
\end{proof}

\begin{lemma}
\label{spaces-morphisms-lemma-immersion-permanence}
设 $S$ 为概形。
设 $Z \to Y \to X$ 为 $S$ 上代数空间的态射。
\begin{enumerate}
\item 若 $Z \to X$ 可表示、局有限型、局拟有限、分离且为单态射，
则 $Z \to Y$ 也可表示、局有限型、局拟有限、分离且为单态射。
\item 若 $Z \to X$ 为浸入且 $Y \to X$ 局部分离，则 $Z \to Y$ 为浸入。
\item 若 $Z \to X$ 为闭浸入且 $Y \to X$ 分离，则 $Z \to Y$ 为闭浸入。
\end{enumerate}
\end{lemma}

\begin{proof}
每种情形都考虑交换图表
$$
\xymatrix{
Z \ar[r] \ar[rd] & Y \times_X Z \ar[r] \ar[d] & Z \ar[d] \\
& Y \ar[r] & X
}
$$
其中上方水平箭头的复合为恒等。证明 (1)。第一个水平箭头是
$Y \times_X Z \to Z$ 的截面，故由引理 \ref{spaces-morphisms-lemma-section-immersion}，
它可表示、局有限型、局拟有限、分离且为单态射。
态射 $Y \times_X Z \to Y$ 是 $Z \to X$ 的基变换，因而也可表示、
局有限型、局拟有限、分离且为单态射（这些概形态射性质在基变换下保持；见
《空间》注记 \ref{spaces-remark-list-properties-stable-base-change}）。
因此复合也具有同样性质（这些概形态射性质在复合下保持；见《空间》注记
\ref{spaces-remark-list-properties-stable-composition}）。这证明 (1)。
其余结论完全同理。
\end{proof}

\begin{lemma}
\label{spaces-morphisms-lemma-immersion-when-closed}
设 $S$ 为概形。设 $i : Z \to X$ 为 $S$ 上代数空间的浸入。则
$|i| : |Z| \to |X|$ 是到局部闭子集上的同胚，且 $i$ 为闭浸入，当且仅当
像 $|i|(|Z|) \subset |X|$ 是闭集。
\end{lemma}

\begin{proof}
第一条断言是《空间的性质》引理
\ref{spaces-properties-lemma-subspace-induced-topology}。设 $U$ 为概形，
且 $U \to X$ 为满射 \'etale 态射。由假设，$T = U \times_X Z$ 为概形，
且 $j : T \to U$ 为概形的浸入。由引理 \ref{spaces-morphisms-lemma-closed-immersion-local}，
$i$ 为闭浸入，当且仅当 $j$ 为闭浸入。由《概形》引理
\ref{schemes-lemma-immersion-when-closed}，这又当且仅当 $j(T)$ 在 $U$
中为闭集。然而子集 $j(T) \subset U$ 是 $|i|(|Z|) \subset |X|$ 的逆像，
见《空间的性质》引理 \ref{spaces-properties-lemma-points-cartesian}。
证明完毕。
\end{proof}

\begin{remark}
\label{spaces-morphisms-remark-immersion}
设 $S$ 为概形。设 $i : Z \to X$ 为 $S$ 上代数空间的浸入。由于 $i$ 是
单态射，可以将 $|Z|$ 看作 $|X|$ 的子集；本注记中均如此处理。
令 $\partial |Z|$ 为 $|Z|$ 在拓扑空间 $|X|$ 中的边界。公式为
$$
\partial |Z| = \overline{|Z|} \setminus |Z|.
$$
令 $\partial Z$ 为 $X$ 的约化闭子空间，满足
$|\partial Z| = \partial |Z|$，其由约化诱导闭子空间结构得到；见《空间的性质》
定义 \ref{spaces-properties-definition-reduced-induced-space}。由构造，
$|Z|$ 在 $|X| \setminus |\partial Z| = |X \setminus \partial Z|$ 中为闭集。
因此任何代数空间浸入都可分解为先闭浸入后开浸入（但一般不能反向分解；见
《态射》例 \ref{morphisms-example-thibaut}）。
\end{remark}

\begin{remark}
\label{spaces-morphisms-remark-space-structure-locally-closed-subset}
设 $S$ 为概形，设 $X$ 为 $S$ 上的代数空间。设 $T \subset |X|$ 为局部闭子集。
令 $\partial T$ 为 $T$ 在拓扑空间 $|X|$ 中的边界。公式为
$$
\partial T = \overline{T} \setminus T.
$$
令 $U \subset X$ 为 $X$ 的开子空间，满足 $|U| = |X| \setminus \partial T$，见《空间的性质》
引理 \ref{spaces-properties-lemma-open-subspaces}。令 $Z$ 为 $U$ 的约化闭子空间，
满足 $|Z| = T$，由约化诱导闭子空间结构得到，见《空间的性质》定义
\ref{spaces-properties-definition-reduced-induced-space}。由构造，
$Z \to U$ 是代数空间的闭浸入，而 $U \to X$ 是开浸入，故
$Z \to X$ 是 $S$ 上代数空间的浸入（见《空间》引理
\ref{spaces-lemma-composition-immersions}）。注意 $Z$ 是约化代数空间，且作为
$|X|$ 的子集有 $|Z| = T$。有时称 $Z$ 为 $T$ 上的 {\it 约化诱导子空间结构}。
\end{remark}

\begin{lemma}
\label{spaces-morphisms-lemma-factor-the-other-way}
设 $S$ 为概形。设 $Z \to X$ 为 $S$ 上代数空间的浸入，并假设 $Z \to X$
拟紧。存在分解 $Z \to \overline{Z} \to X$，其中
$Z \to \overline{Z}$ 为开浸入，$\overline{Z} \to X$ 为闭浸入。
\end{lemma}

\begin{proof}
设 $U$ 为概形，且 $U \to X$ 为满射 \'etale 态射。按通常记号，令
$R = U \times_X U$，投影为 $s, t : R \to U$。令 $T = Z \times_U X$。
令 $\overline{T} \subset U$ 为 $T \to U$ 的概形论像。注意，由于拟紧态射的
概形论像与平坦基变换交换，
$s^{-1}\overline{T} = t^{-1}\overline{T}$；见《态射》引理
\ref{morphisms-lemma-flat-base-change-scheme-theoretic-image}。
因此得到闭子空间 $\overline{Z} \subset X$，其到 $U$ 的拉回为 $\overline{T}$；
见《空间的性质》引理 \ref{spaces-properties-lemma-subspaces-presentation}。
由《态射》引理 \ref{morphisms-lemma-quasi-compact-immersion}，
$T \to \overline{T}$ 为开浸入。因此 $Z \to \overline{Z}$ 为开浸入，结论成立。
\end{proof}


















\section{闭浸入}
\label{spaces-morphisms-section-closed-immersions}

\noindent
本节阐明前面关于代数空间浸入所得的一些结果。见《空间》第
\ref{spaces-section-Zariski} 节及本章第 \ref{spaces-morphisms-section-immersions} 节。
本节是代数空间版本的《态射》第
\ref{morphisms-section-closed-immersions} 节。

\begin{lemma}
\label{spaces-morphisms-lemma-closed-immersion-ideals}
设 $S$ 为概形。
设 $X$ 为 $S$ 上的代数空间。对于每个闭浸入 $i : Z \to X$，层
$i_*\mathcal{O}_Z$ 是拟凝聚 $\mathcal{O}_X$-模，映射
$i^\sharp : \mathcal{O}_X \to i_*\mathcal{O}_Z$ 为满射，且其核是拟凝聚理想层。
对应规则
$Z \mapsto \Ker(\mathcal{O}_X \to i_*\mathcal{O}_Z)$
定义一个反包含双射
$$
\begin{matrix}
\text{闭子空间}\\
Z \subset X
\end{matrix}
\longrightarrow
\begin{matrix}
\text{拟凝聚层}\\
\text{理想 }\mathcal{I} \subset \mathcal{O}_X
\end{matrix}
$$
此外，给定对应于拟凝聚理想层 $\mathcal{I} \subset \mathcal{O}_X$ 的闭子概形
$Z$，代数空间态射 $h : Y \to X$ 可经由 $Z$ 分解，当且仅当映射
$h^*\mathcal{I} \to h^*\mathcal{O}_X = \mathcal{O}_Y$ 为零。
\end{lemma}

\begin{proof}
设 $U \to X$ 为满射 \'etale 态射，且源为概形。考虑图表
$$
\xymatrix{
U \times_X Z \ar[r] \ar[d]_{i'} & Z \ar[d]^i \\
U \ar[r] & X
}
$$
由《态射》引理 \ref{spaces-morphisms-lemma-closed-immersion-local}
可知 $i$ 为闭浸入
当且仅当 $i'$ 为闭浸入。由
《空间的性质》
引理 \ref{spaces-properties-lemma-pushforward-etale-base-change-modules}
表明 $i'_*\mathcal{O}_{U \times_X Z}$ 是 $i_*\mathcal{O}_Z$ 到 $U$ 的限制。
因此关于 $\mathcal{O}_X \to i_*\mathcal{O}_Z$ 的断言，等价于关于
$\mathcal{O}_U \to i'_*\mathcal{O}_{U \times_X Z}$ 的相应断言。
又由于 $i'$ 是概形的闭浸入，这些结果由《态射》引理
\ref{morphisms-lemma-closed-immersion} 得到。

\medskip\noindent
 现在证明：给定拟凝聚理想层 $\mathcal{I} \subset \mathcal{O}_X$，公式
$$
Z(T) = \{h : T \to X \mid h^*\mathcal{I} \to \mathcal{O}_T
\text{ 为零}\}
$$
定义 $X$ 的一个闭子空间。显然它是 $X$ 的子函子。
为证明 $Z \to X$ 可由闭浸入表示，设
$\varphi : U \to X$ 是从概形到 $X$ 的态射。那么
$Z \times_X U$ 由 $U$ 的相应子函子表示，该子函子对应于理想层
$\Im(\varphi^*\mathcal{I} \to \mathcal{O}_U)$。由《空间的性质》
引理 \ref{spaces-properties-lemma-pullback-quasi-coherent}，
$\mathcal{O}_U$-模 $\varphi^*\mathcal{I}$ 在 $U$ 上是拟凝聚的，
从而 $\Im(\varphi^*\mathcal{I} \to \mathcal{O}_U)$
是 $U$ 上的拟凝聚理想层。由《概形》引理
\ref{schemes-lemma-characterize-closed-subspace}，可知
$Z \times_X U$ 由与 $\Im(\varphi^*\mathcal{I} \to \mathcal{O}_U)$
关联的 $U$ 的闭子概形表示。因此 $Z$ 是 $X$ 的闭子空间。

\medskip\noindent
在上述 $Z$ 的公式中，输入 $T$ 是概形，这是因为代数空间是
$(\Sch/S)_{fppf}$ 上的层。我们省略验证当 $T$ 为代数空间时
同一公式仍然成立的过程。
\end{proof}

\begin{definition}
\label{spaces-morphisms-definition-inverse-image-closed-subspace}
设 $S$ 为概形。设 $f : Y \to X$ 是 $S$ 上代数空间的态射。
设 $Z \subset X$ 为闭子空间。闭子空间 $Z$ 的
{\it 逆像 $f^{-1}(Z)$} 定义为 $Y$ 的闭子空间 $Z \times_X Y$。
\end{definition}

\noindent
由引理 \ref{spaces-morphisms-lemma-closed-immersion-local}，此定义是有意义的。
若 $\mathcal{I} \subset \mathcal{O}_X$ 是通过引理
\ref{spaces-morphisms-lemma-closed-immersion-ideals} 与 $Z$ 对应的拟凝聚理想层，则
$f^{-1}\mathcal{I}\mathcal{O}_Y = \Im(f^*\mathcal{I} \to \mathcal{O}_Y)$
是与 $f^{-1}(Z)$ 对应的理想层。

\begin{lemma}
\label{spaces-morphisms-lemma-closed-immersion-quasi-compact}
代数空间的闭浸入是拟紧的。
\end{lemma}

\begin{proof}
这由一般原理及《概形》引理
\ref{schemes-lemma-closed-immersion-quasi-compact} 得到；见《空间》引理
\ref{spaces-lemma-representable-transformations-property-implication}。
\end{proof}

\begin{lemma}
\label{spaces-morphisms-lemma-closed-immersion-separated}
代数空间的闭浸入是分离的。
\end{lemma}

\begin{proof}
这由一般原理及《概形》引理
\ref{schemes-lemma-immersions-monomorphisms} 得到；见《空间》引理
\ref{spaces-lemma-representable-transformations-property-implication}。
\end{proof}

\begin{lemma}
\label{spaces-morphisms-lemma-closed-immersion-push-pull}
设 $S$ 为概形。设 $i : Z \to X$ 是 $S$ 上代数空间的闭浸入。
\begin{enumerate}
\item 函子
$$
i_{small, *} :
\Sh(Z_\etale)
\longrightarrow
\Sh(X_\etale)
$$
是全忠实的，其本质像由 $X_\etale$ 上满足下述条件的集合层
$\mathcal{F}$ 组成：其到 $X \setminus Z$ 的限制同构于 $*$；以及
\item 函子
$$
i_{small, *} :
\textit{Ab}(Z_\etale)
\longrightarrow
\textit{Ab}(X_\etale)
$$
是全忠实的，其本质像由 $X_\etale$ 上支集包含于 $|Z|$ 的阿贝尔层组成。
\end{enumerate}
在两种情形中，$i_{small}^{-1}$ 都是函子 $i_{small, *}$ 的左逆。
\end{lemma}

\begin{proof}
设 $U$ 为概形，且 $U \to X$ 为满射 \'etale 态射。令
$V = Z \times_X U$。则 $V$ 为概形，且 $i' : V \to U$ 是概形的闭浸入。由
《空间的性质》引理 \ref{spaces-properties-lemma-pushforward-etale-base-change}，
对 $Z$ 上的任意层 $\mathcal{G}$，有
$$
(i_{small}^{-1}i_{small, *}\mathcal{G})|_V =
(i')_{small}^{-1}i'_{small, *}(\mathcal{G}|_V)
$$
由《空间的上同调》命题
\ref{etale-cohomology-proposition-closed-immersion-pushforward}，映射
$(i')_{small}^{-1}i'_{small, *}(\mathcal{G}|_V) \to \mathcal{G}|_V$
是同构。由于 $V \to Z$ 是满射 \'etale 态射，这意味着
$i_{small}^{-1}i_{small, *}\mathcal{G} \to \mathcal{G}$ 是同构。
这显然推出 $i_{small, *}$ 是全忠实的；见《网站》引理
\ref{sites-lemma-exactness-properties}。
为证明关于本质像的断言，考虑 $X_\etale$ 上的集合层 $\mathcal{F}$，
其到 $X \setminus Z$ 的限制同构于 $*$。如《空间的上同调》命题
\ref{etale-cohomology-proposition-closed-immersion-pushforward} 的证明中那样，
我们考虑伴随映射
$$
\mathcal{F} \longrightarrow i_{small, *}i_{small}^{-1}\mathcal{F}.
$$
如第一部分，我们由概形情形的相应结果看到该映射到
$U$ 的限制是同构。由于 $U$ 是 $X$ 的 \'etale 覆盖，故它是同构。
\end{proof}

\begin{lemma}
\label{spaces-morphisms-lemma-stalk-push-closed}
设 $S$ 为概形。设 $i : Z \to X$ 是 $S$ 上代数空间的闭浸入。
设 $\overline{z}$ 为 $Z$ 的几何点，其在 $X$ 中的像为 $\overline{x}$。则对
$Z_\etale$ 上的任意层 $\mathcal{F}$，有
$(i_{small, *}\mathcal{F})_{\overline{z}} = \mathcal{F}_{\overline{x}}$。
\end{lemma}

\begin{proof}
取 $\overline{x}$ 的一个 \'etale 邻域 $(U, \overline{u})$。
那么茎 $(i_{small, *}\mathcal{F})_{\overline{z}}$ 是
$i_{small, *}\mathcal{F}|_U$ 在 $\overline{u}$ 处的茎。由《空间的性质》引理
\ref{spaces-properties-lemma-pushforward-etale-base-change}，可将 $X$ 替换为 $U$，
并将 $Z$ 替换为 $Z \times_X U$。此时 $Z \to X$ 是概形的闭浸入，结论就是
《空间的上同调》引理
\ref{etale-cohomology-lemma-stalk-pushforward-closed-immersion}。
\end{proof}

\noindent
下述引理在拓扑斯的闭浸入情形下更一般地成立（未来补充引用）。

\begin{lemma}
\label{spaces-morphisms-lemma-closed-immersion-rings}
设 $S$ 为概形。设 $i : Z \to X$ 是 $S$ 上代数空间的闭浸入。设
$\mathcal{A}$ 为 $X_\etale$ 上的环层，$\mathcal{B}$ 为 $Z_\etale$ 上的环层。
设 $\varphi : \mathcal{A} \to i_{small, *}\mathcal{B}$ 为环层同态，
从而得到环化拓扑斯的态射
$$
f : (\Sh(Z_\etale), \mathcal{B}) \longrightarrow (\Sh(X_\etale), \mathcal{A}).
$$
对于 $\mathcal{A}$-模层 $\mathcal{F}$ 和 $\mathcal{B}$-模层 $\mathcal{G}$，典范映射
$$
\mathcal{F} \otimes_\mathcal{A} f_*\mathcal{G}
\longrightarrow
f_*(f^*\mathcal{F} \otimes_\mathcal{B} \mathcal{G}).
$$
是同构。
\end{lemma}

\begin{proof}
该映射是下述映射的伴随映射：
$$
f^*\mathcal{F} \otimes_\mathcal{B}
f^* f_*\mathcal{G} =
f^*(\mathcal{F} \otimes_\mathcal{A} f_*\mathcal{G})
\longrightarrow
f^*\mathcal{F} \otimes_\mathcal{B} \mathcal{G}
$$
它来自 $\text{id} : f^*\mathcal{F} \to f^*\mathcal{F}$
以及伴随映射 $f^* f_*\mathcal{G} \to \mathcal{G}$。
为证明该映射是同构，可以在茎上检验（《空间的性质》定理
\ref{spaces-properties-theorem-exactness-stalks}）。
设 $\overline{z} : \Spec(k) \to Z$ 为几何点，其像为
$\overline{x} = i \circ \overline{z} : \Spec(k) \to X$。
逐点计算这些映射在茎上的作用可知，需要证明
$$
\mathcal{F}_{\overline{x}}
\otimes_{\mathcal{A}_{\overline{x}}}
\mathcal{G}_{\overline{z}} =
(\mathcal{F}_{\overline{x}}
\otimes_{\mathcal{A}_{\overline{x}}}
\mathcal{B}_{\overline{z}}) \otimes_{\mathcal{B}_{\overline{z}}}
\mathcal{G}_{\overline{z}}
$$
这成立。这里使用了张量积与取茎可交换、拉回下茎的行为（《空间的性质》引理
\ref{spaces-properties-lemma-stalk-pullback}），以及沿闭浸入推前时茎的行为
（引理 \ref{spaces-morphisms-lemma-stalk-push-closed}）。
\end{proof}






\section{闭浸入与拟凝聚层}
\label{spaces-morphisms-section-closed-immersions-quasi-coherent}

\noindent
本节是《态射》第 \ref{morphisms-section-closed-immersions-quasi-coherent} 节的代数空间版本。

\begin{lemma}
\label{spaces-morphisms-lemma-i-star-equivalence}
设 $S$ 为概形。设 $i : Z \to X$ 是 $S$ 上代数空间的闭浸入。
设 $\mathcal{I} \subset \mathcal{O}_X$ 为定义 $Z$ 的拟凝聚理想层。
\begin{enumerate}
\item 对任意 $\mathcal{O}_X$-模 $\mathcal{F}$，伴随映射
$\mathcal{F} \to i_*i^*\mathcal{F}$ 诱导同构
$\mathcal{F}/\mathcal{I}\mathcal{F} \cong i_*i^*\mathcal{F}$。
\item 函子 $i^*$ 是 $i_*$ 的左逆，即对任意
$\mathcal{O}_Z$-模 $\mathcal{G}$，伴随映射
$i^*i_*\mathcal{G} \to \mathcal{G}$ 是同构。
\item 函子
$$
i_* :
\QCoh(\mathcal{O}_Z)
\longrightarrow
\QCoh(\mathcal{O}_X)
$$
是正合且全忠实的，其本质像由满足 $\mathcal{I}\mathcal{F} = 0$ 的
拟凝聚 $\mathcal{O}_X$-模 $\mathcal{F}$ 组成。
\end{enumerate}
\end{lemma}

\begin{proof}
在本证明中我们只处理小 \'etale 位点上的层，并使用 $i_*, i^{-1}, \ldots$
表示阿贝尔群层的推前与拉回，而不使用 $i_{small, *}, i_{small}^{-1}$。

\medskip\noindent
设 $\mathcal{F}$ 为 $\mathcal{O}_X$-模。将引理
\ref{spaces-morphisms-lemma-closed-immersion-rings} 应用于
$\mathcal{A} = \mathcal{O}_X$ 及
$\mathcal{G} = \mathcal{B} = \mathcal{O}_Z$，可知
$i_*i^*\mathcal{F} = \mathcal{F} \otimes_{\mathcal{O}_X} \mathcal{O}_Z$.
由引理 \ref{spaces-morphisms-lemma-closed-immersion-ideals}，有短正合列
$$
0 \to \mathcal{I} \to \mathcal{O}_X \to i_*\mathcal{O}_Z \to 0
$$
由张量积的性质，
$\mathcal{F} \otimes_{\mathcal{O}_X} i_*\mathcal{O}_Z
= \mathcal{F}/\mathcal{I}\mathcal{F}$。这证明了 (1)（但我们省略验证该映射
由伴随映射诱导）。

\medskip\noindent
设 $\mathcal{G}$ 为任意 $\mathcal{O}_Z$-模。由引理
\ref{spaces-morphisms-lemma-closed-immersion-push-pull}，有 $i^{-1}i_*\mathcal{G} = \mathcal{G}$。
因此为证明 (2)，需证明典范映射
$\mathcal{G} \otimes_{i^{-1}\mathcal{O}_X} \mathcal{O}_Z \to \mathcal{G}$
是同构。若能证明 $i^{-1}\mathcal{O}_X \to \mathcal{O}_Z$ 为满射，这由张量积的一般性质得到。
由引理 \ref{spaces-morphisms-lemma-closed-immersion-push-pull}，只需证明
$i_*i^{-1}\mathcal{O}_X \to i_*\mathcal{O}_Z$
是满射。由于满射映射 $\mathcal{O}_X \to i_*\mathcal{O}_Z$
经由此映射因子化，故 (2) 成立。

\medskip\noindent
最后证明引理最重要的部分，即 (3)。闭浸入是拟紧且分离的，见引理
\ref{spaces-morphisms-lemma-closed-immersion-quasi-compact} 和
\ref{spaces-morphisms-lemma-closed-immersion-separated}。因此引理
\ref{spaces-morphisms-lemma-pushforward} 适用，$Z$ 上拟凝聚层的推前确实是 $X$ 上的拟凝聚层。
于是得到函子
$i^{QCoh}_* : \QCoh(\mathcal{O}_Z)
\to \QCoh(\mathcal{O}_X)$.
由 (2) 可知 $i^{QCoh}_*$ 是全忠实的，因为它具有左逆 $i^*$。

\medskip\noindent
现在转而描述函子 $i_*$ 的本质像。对任意 $\mathcal{O}_Z$-模，显然有
$\mathcal{I}(i_*\mathcal{G}) = 0$，因为 $\mathcal{I}$ 是映射
$\mathcal{O}_X \to i_*\mathcal{O}_Z$ 的核，而我们用该映射在 $i_*\mathcal{G}$ 上定义
$\mathcal{O}_X$-模结构。反之，设 $\mathcal{F}$ 是满足
$\mathcal{I}\mathcal{F} = 0$ 的任意拟凝聚 $\mathcal{O}_X$-模。则
$\mathcal{F}$ 是 $i_*\mathcal{O}_Z$-模，因为 $i_*\mathcal{O}_Z = \mathcal{O}_X/\mathcal{I}$。
特别地，其支集包含于 $|Z|$。应用引理
\ref{spaces-morphisms-lemma-closed-immersion-push-pull} 可知存在某个 $\mathcal{O}_Z$-模 $\mathcal{G}$ 使
$\mathcal{F} \cong i_*\mathcal{G}$。剩下的唯一细节是说明 $\mathcal{G}$ 为拟凝聚。
这是因为由 (2) 和《空间的性质》引理
\ref{spaces-properties-lemma-pullback-quasi-coherent}，$\mathcal{G} \cong i^*\mathcal{F}$。
\end{proof}

\noindent
设 $i : Z \to X$ 为代数空间的闭浸入。由上述引理，滥用记号时我们常把
$i_*\mathcal{F}$ 在 $X$ 上的层也记为 $\mathcal{F}$。

\begin{lemma}
\label{spaces-morphisms-lemma-largest-quasi-coherent-subsheaf}
设 $S$ 为概形，设 $X$ 为 $S$ 上的代数空间。
设 $\mathcal{F}$ 为拟凝聚 $\mathcal{O}_X$-模，设 $\mathcal{G} \subset \mathcal{F}$ 为
$\mathcal{O}_X$-子模。存在唯一的拟凝聚 $\mathcal{O}_X$-子模
$\mathcal{G}' \subset \mathcal{G}$，具有如下性质：对每个拟凝聚
$\mathcal{O}_X$-模 $\mathcal{H}$，映射
$$
\Hom_{\mathcal{O}_X}(\mathcal{H}, \mathcal{G}')
\longrightarrow
\Hom_{\mathcal{O}_X}(\mathcal{H}, \mathcal{G})
$$
是双射。特别地，$\mathcal{G}'$ 是包含于 $\mathcal{G}$ 的 $\mathcal{F}$ 的最大拟凝聚
$\mathcal{O}_X$-子模。
\end{lemma}

\begin{proof}
令 $\mathcal{G}_a$（$a \in A$）为包含于 $\mathcal{G}$ 的拟凝聚
$\mathcal{O}_X$-子模的集合。令映射
$$
\bigoplus\nolimits_{a \in A} \mathcal{G}_a \longrightarrow \mathcal{F}
$$
的像 $\mathcal{G}'$。它是拟凝聚的，因为 $X$ 上拟凝聚层之间的映射的像仍为拟凝聚层，
且拟凝聚层的直和是拟凝聚的；见《空间的性质》引理
\ref{spaces-properties-lemma-properties-quasi-coherent}。
模 $\mathcal{G}'$ 包含于 $\mathcal{G}$，故它是包含于 $\mathcal{G}$ 的最大拟凝聚
$\mathcal{O}_X$-模。

\medskip\noindent
为证明该公式，设 $\mathcal{H}$ 为拟凝聚 $\mathcal{O}_X$-模，且
$\alpha : \mathcal{H} \to \mathcal{G}$ 为 $\mathcal{O}_X$-模映射。复合映射
$\mathcal{H} \to \mathcal{G} \to \mathcal{F}$ 的像是拟凝聚的，因为它是拟凝聚层映射的像。
因此该像包含于 $\mathcal{G}'$，从而 $\alpha$ 经由 $\mathcal{G}'$ 分解，如所需。
\end{proof}

\begin{lemma}
\label{spaces-morphisms-lemma-i-upper-shriek}
设 $S$ 为概形。设 $i : Z \to X$ 是 $S$ 上代数空间的闭浸入。
存在函子\footnote{这很可能是非标准记号。}
$i^! : \QCoh(\mathcal{O}_X) \to \QCoh(\mathcal{O}_Z)$
它是 $i_*$ 的右伴随。（比较《模》引理
\ref{modules-lemma-i-star-right-adjoint}。）
\end{lemma}

\begin{proof}
给定拟凝聚 $\mathcal{O}_X$-模 $\mathcal{G}$，考虑 $\mathcal{G}$ 的子层
$\mathcal{H}_Z(\mathcal{G})$，它由被 $\mathcal{I}$ 湮灭的局部截面组成。由引理
\ref{spaces-morphisms-lemma-largest-quasi-coherent-subsheaf}，存在典范的最大拟凝聚
$\mathcal{O}_X$-子模 $\mathcal{H}_Z(\mathcal{G})'$。按构造有
$$
\Hom_{\mathcal{O}_X}(i_*\mathcal{F}, \mathcal{H}_Z(\mathcal{G})')
=
\Hom_{\mathcal{O}_X}(i_*\mathcal{F}, \mathcal{G})
$$
对任意拟凝聚 $\mathcal{O}_Z$-模 $\mathcal{F}$ 成立。
因此可令 $i^!\mathcal{G} = i^*(\mathcal{H}_Z(\mathcal{G})')$。细节略。
\end{proof}

\noindent
利用拟凝聚理想层与闭子空间之间的 $1$ 对 $1$ 对应（见引理
\ref{spaces-morphisms-lemma-closed-immersion-ideals}），可以定义闭子概形的概形论交与概形论并。

\begin{definition}
\label{spaces-morphisms-definition-scheme-theoretic-intersection-union}
设 $S$ 为概形，设 $X$ 为 $S$ 上的代数空间。设 $Z, Y \subset X$ 为闭子空间，
分别对应于拟凝聚理想层 $\mathcal{I}, \mathcal{J} \subset \mathcal{O}_X$。
$Z$ 与 $Y$ 的 {\it 概形论交} 是由 $\mathcal{I} + \mathcal{J}$
定义的 $X$ 的闭子空间。$Z$ 与 $Y$ 的 {\it 概形论并} 是由
$\mathcal{I} \cap \mathcal{J}$ 定义的 $X$ 的闭子空间。
\end{definition}

\noindent
显然，概形论交的构造与 'etale 局部化相容，概形论并亦然。

\begin{lemma}
\label{spaces-morphisms-lemma-scheme-theoretic-intersection}
设 $S$ 为概形，设 $X$ 为 $S$ 上的代数空间。设 $Z, Y \subset X$ 为闭子空间。
令 $Z \cap Y$ 为 $Z$ 与 $Y$ 的概形论交。则 $Z \cap Y \to Z$ 与
$Z \cap Y \to Y$ 是闭浸入，并且
$$
\xymatrix{
Z \cap Y \ar[r] \ar[d] & Z \ar[d] \\
Y \ar[r] & X
}
$$
是 $S$ 上代数空间的笛卡尔图，即 $Z \cap Y = Z \times_X Y$。
\end{lemma}

\begin{proof}
由引理 \ref{spaces-morphisms-lemma-closed-immersion-ideals}，态射 $Z \cap Y \to Z$ 与
$Z \cap Y \to Y$ 是闭浸入。由于概形论交的构造与 'etale 局部化相容，
由概形情形可知该图为笛卡尔图。见《态射》引理
\ref{morphisms-lemma-scheme-theoretic-intersection}。
\end{proof}

\begin{lemma}
\label{spaces-morphisms-lemma-scheme-theoretic-union}
设 $S$ 为概形，设 $X$ 为 $S$ 上的代数空间。设 $Y, Z \subset X$ 为闭子空间。
令 $Y \cup Z$ 为 $Y$ 与 $Z$ 的概形论并，令 $Y \cap Z$ 为 $Y$ 与 $Z$ 的概形论交。
则 $Y \to Y \cup Z$ 与 $Z \to Y \cup Z$ 是闭浸入，且存在短正合列
$$
0 \to \mathcal{O}_{Y \cup Z} \to \mathcal{O}_Y \times \mathcal{O}_Z
\to \mathcal{O}_{Y \cap Z} \to 0
$$
的 $\mathcal{O}_Z$-模，并且图
$$
\xymatrix{
Y \cap Z \ar[r] \ar[d] & Y \ar[d] \\
Z \ar[r] & Y \cup Z
}
$$
在 $S$ 上代数空间范畴中为余笛卡尔图，即
$Y \cup Z = Y \amalg_{Y \cap Z} Z$。
\end{lemma}

\begin{proof}
由引理 \ref{spaces-morphisms-lemma-closed-immersion-ideals}，态射 $Y \to Y \cup Z$ 与
$Z \to Y \cup Z$ 是闭浸入。在短正合列中，我们利用引理
\ref{spaces-morphisms-lemma-i-star-equivalence} 的等价性，把 $X$ 的闭子空间上的拟凝聚模视为 $X$ 上的拟凝聚模。
对列中的第一个映射，使用典范映射
$\mathcal{O}_{Y \cup Z} \to \mathcal{O}_Y$ 和
$\mathcal{O}_{Y \cup Z} \to \mathcal{O}_Z$；
对第二个映射，使用典范映射
$\mathcal{O}_Y \to \mathcal{O}_{Y \cap Z}$ 以及
典范映射的负映射
$\mathcal{O}_Z \to \mathcal{O}_{Y \cap Z}$。为检验正合性，可以在 \'etale 局部工作，并由概形情形推出正合性
（《态射》引理 \ref{morphisms-lemma-scheme-theoretic-union}）。

\medskip\noindent
为证明该图为余笛卡尔图，设给定 $S$ 上的代数空间 $T$ 以及态射
$f : Y \to T$、$g : Z \to T$，它们在 $Y \cap Z \to T$ 上相同。目标是证明存在唯一的
态射 $h : Y \cup Z \to T$，其与 $f$、$g$ 相容。
构造 $h$ 时，可以在 $Y \cup Z$ 上作 \'etale 局部工作（因为作为代数空间，$Y \cup Z$ 是 \'etale 层）。
因此可假设 $X$ 为概形。此时由《态射》引理
\ref{morphisms-lemma-scheme-theoretic-union}，$Y \cup Z$ 是概形范畴中沿 $Y \cap Z$
将 $Y$ 与 $Z$ 粘合所得的推出。取概形 $T'$ 及满射 \'etale 态射 $T' \to T$。
令 $Y' = T' \times_{T, f} Y$，$Z' = T' \times_{T, g} Z$。
则 $Y'$ 与 $Z'$ 为概形，并有概形的典范同构
$\varphi : Y' \times_Y (Y \cap Z) \to Z' \times_Z (Y \cap Z)$。由《态射补篇》引理
\ref{more-morphisms-lemma-pushout-along-closed-immersions}
推出 $W' = Y' \amalg_{Y' \times_Y (Y \cap Z), \varphi} Z'$ 在概形范畴中存在。
由《态射补篇》引理
\ref{more-morphisms-lemma-pushout-along-closed-immersions-properties-above}，
态射 $W' \to Y \cup Z$ 是 \'etale 的。该态射是满射，因为 $Y' \to Y$ 与 $Z' \to Z$ 是满射。
态射 $f' : Y' \to T'$ 与 $g' : Z' \to T'$ 粘合为唯一的概形态射
$h' : W' \to T'$。由唯一性，复合 $W' \to T' \to T$ 下降为所需的态射
$h : Y \cup Z \to T$。细节略。
\end{proof}






\section{模的支集}
\label{spaces-morphisms-section-support}

\noindent
本节汇集代数空间上拟凝聚模支集的一些基本结果。设 $X$ 为代数
空间。$X_\etale$ 上阿贝尔层的支集已在《空间的性质》第
\ref{spaces-properties-section-support} 节中定义。我们对模的支集使用同一定义。
下述引理说明，这与概形上拟凝聚模所定义的概念一致。

\begin{lemma}
\label{spaces-morphisms-lemma-support-covering}
设 $S$ 为概形，设 $X$ 为 $S$ 上的代数空间。设 $\mathcal{F}$ 为拟凝聚
$\mathcal{O}_X$-模。设 $U$ 为概形，且 $\varphi : U \to X$ 为 \'etale 态射。则
$$
\text{Supp}(\varphi^*\mathcal{F}) = |\varphi|^{-1}(\text{Supp}(\mathcal{F}))
$$
其中左边是 $\varphi^*\mathcal{F}$ 作为概形 $U$ 上拟凝聚模时的支集。
\end{lemma}

\begin{proof}
 设 $u\in U$ 为普通点，设 $\overline{x}$ 为位于 $u$ 上方的几何点。由《空间的性质》引理
\ref{spaces-properties-lemma-stalk-quasi-coherent}，有
$(\varphi^*\mathcal{F})_u \otimes_{\mathcal{O}_{U, u}}
\mathcal{O}_{X, \overline{x}} = \mathcal{F}_{\overline{x}}$.
由于由《空间的性质》引理
\ref{spaces-properties-lemma-describe-etale-local-ring}，$\mathcal{O}_{U, u} \to \mathcal{O}_{X, \overline{x}}$
是严格 Hensel 化，故它是忠实平坦的（见《代数补篇》引理
\ref{more-algebra-lemma-dumb-properties-henselization}）。因此
$(\varphi^*\mathcal{F})_u = 0$ 当且仅当 $\mathcal{F}_{\overline{x}} = 0$。引理得证。
\end{proof}

\noindent
对有限型拟凝聚模，支集是闭的，可在纤维上检验，并与基变换相容。

\begin{lemma}
\label{spaces-morphisms-lemma-support-finite-type}
设 $S$ 为概形，设 $X$ 为 $S$ 上的代数空间。设 $\mathcal{F}$ 为有限型拟凝聚
$\mathcal{O}_X$-模。则
\begin{enumerate}
\item $\mathcal{F}$ 的支集是闭集。
\item 对于位于 $x \in |X|$ 上方的几何点 $\overline{x}$，有
$$
x \in \text{Supp}(\mathcal{F})
\Leftrightarrow
\mathcal{F}_{\overline{x}} \not = 0
\Leftrightarrow
\mathcal{F}_{\overline{x}} \otimes_{\mathcal{O}_{X, \overline{x}}}
\kappa(\overline{x}) \not = 0.
$$
\item 对任意代数空间态射 $f : Y \to X$，拉回 $f^*\mathcal{F}$ 也是有限型的，且
$\text{Supp}(f^*\mathcal{F}) = f^{-1}(\text{Supp}(\mathcal{F}))$。
\end{enumerate}
\end{lemma}

\begin{proof}
取概形 $U$ 及满射 \'etale 态射 $\varphi : U \to X$。由引理
\ref{spaces-morphisms-lemma-support-covering}，$\mathcal{F}$ 的支集的逆像是 $\varphi^*\mathcal{F}$ 的支集，
后者由《态射》引理 \ref{morphisms-lemma-support-finite-type} 为闭集。
因此 (1) 由 $|X|$ 上拓扑的定义得到。

\medskip\noindent
(2) 中第一个等价关系就是支集的定义。第二个等价关系由 Nakayama 引理得到，见
《代数》引理 \ref{algebra-lemma-NAK}。

\medskip\noindent
设 $f : Y \to X$ 如 (3) 所述。由《空间的性质》第
\ref{spaces-properties-section-properties-modules} 节，$f^*\mathcal{F}$ 是有限型的。
为证明最后一个断言，设 $\overline{y}$ 为 $Y$ 的几何点，映到 $X$ 上的几何点 $\overline{x}$。回忆
$$
(f^*\mathcal{F})_{\overline{y}} =
\mathcal{F}_{\overline{x}} \otimes_{\mathcal{O}_{X, \overline{x}}}
\mathcal{O}_{Y, \overline{y}},
$$
见《空间的性质》引理
\ref{spaces-properties-lemma-stalk-pullback-quasi-coherent}。
因此 $(f^*\mathcal{F})_{\overline{y}} \otimes \kappa(\overline{y})$ 非零当且仅当
$\mathcal{F}_{\overline{x}} \otimes \kappa(\overline{x})$ 非零。
由 (2)，这意味着 $x \in \text{Supp}(\mathcal{F})$ 当且仅当
$y \in \text{Supp}(f^*\mathcal{F})$，这正是断言 (3)。
\end{proof}

\noindent
下一个任务是说明，有限型拟凝聚模的概形论支集（见《态射》定义
\ref{morphisms-definition-scheme-theoretic-support}）
也可以对代数空间上的有限型拟凝聚模定义。

\begin{lemma}
\label{spaces-morphisms-lemma-scheme-theoretic-support}
设 $S$ 为概形，设 $X$ 为 $S$ 上的代数空间。设 $\mathcal{F}$ 为有限型拟凝聚
$\mathcal{O}_X$-模。存在最小闭子空间 $i : Z \to X$，使得存在拟凝聚
$\mathcal{O}_Z$-模 $\mathcal{G}$ 满足 $i_*\mathcal{G} \cong \mathcal{F}$。此外：
\begin{enumerate}
\item 若 $U$ 为概形且 $\varphi : U \to X$ 为 \'etale 态射，则 $Z \times_X U$ 是
$\varphi^*\mathcal{F}$ 的概形论支集。
\item 拟凝聚层 $\mathcal{G}$ 在唯一同构意义下唯一。
\item 拟凝聚层 $\mathcal{G}$ 是有限型的。
\item $\mathcal{G}$ 与 $\mathcal{F}$ 的支集均为 $|Z|$。
\end{enumerate}
\end{lemma}

\begin{proof}
取概形 $U$ 及满射 \'etale 态射 $\varphi : U \to X$。令 $R = U \times_X U$，投影为
$s, t : R \to U$。令 $i' : Z' \to U$ 为 $\varphi^*\mathcal{F}$ 的概形论支集，
并令 $\mathcal{G}'$ 为（在唯一同构意义下）唯一的有限型拟凝聚
$\mathcal{O}_{Z'}$-模，满足 $i'_*\mathcal{G}' = \varphi^*\mathcal{F}$；见《态射》引理
\ref{morphisms-lemma-scheme-theoretic-support}。
由于 $s^*\varphi^*\mathcal{F} = t^*\varphi^*\mathcal{F}$，可知
由《态射》引理 \ref{morphisms-lemma-flat-pullback-support}，
$R' = s^{-1}Z' = t^{-1}Z'$ 是 $R$ 的闭子概形。于是可应用《空间的性质》引理
\ref{spaces-properties-lemma-subspaces-presentation}，得到闭子空间 $i : Z \to X$，
其到 $U$ 的拉回为 $Z'$。记投影为 $s', t' : R' \to Z'$，并记给定的闭浸入为
$j' : R' \to R$，则
$$
j'_* (s')^*\mathcal{G}' = s^* i'_*\mathcal{G}' =
s^*\varphi^*\mathcal{F} = t^*\varphi^*\mathcal{F} =
t^*i'_*\mathcal{G}' = j'_*(t')^*\mathcal{G}'
$$
（第一个和最后一个等式由《概形的上同调》引理
\ref{coherent-lemma-flat-base-change-cohomology} 给出）。因此将《态射》引理
\ref{morphisms-lemma-flat-pullback-support} 应用于 $R' \to R$ 的唯一性，给出同构
$\alpha : (t')^*\mathcal{G}' \to (s')^*\mathcal{G}'$
它经由 $j'_*$ 与典范同构
$t^*\varphi^*\mathcal{F} = s^*\varphi^*\mathcal{F}$ 相容。显然 $\alpha$ 满足上链条件，
因此可应用《空间的性质》命题
\ref{spaces-properties-proposition-quasi-coherent}
得到 $Z$ 上的拟凝聚模 $\mathcal{G}$，其到 $Z'$ 的限制为与 $\alpha$ 相容的 $\mathcal{G}'$。
再次使用上述命题的等价性（这次对 $X$ 使用），可知 $i_*\mathcal{G} \cong \mathcal{F}$。

\medskip\noindent
这就证明了存在性。将它与概形的结果比较，并使用引理
\ref{spaces-morphisms-lemma-support-covering}，即可得到引理的其他性质。详细证明略。
\end{proof}

\begin{definition}
\label{spaces-morphisms-definition-scheme-theoretic-support}
设 $S$ 为概形，设 $X$ 为 $S$ 上的代数空间。设 $\mathcal{F}$ 为有限型拟凝聚
$\mathcal{O}_X$-模。$\mathcal{F}$ 的 {\it 概形论支集} 是引理
\ref{spaces-morphisms-lemma-scheme-theoretic-support} 中构造的闭子空间 $Z \subset X$。
\end{definition}

\noindent
在此情形中，我们常通过引理 \ref{spaces-morphisms-lemma-i-star-equivalence} 的范畴等价，
把 $\mathcal{F}$ 视为 $Z$ 上的有限型拟凝聚层。






\section{概形论像}
\label{spaces-morphisms-section-scheme-theoretic-image}

\noindent
注意：本节的部分内容非常一般，其行为可能与预期不同。

\begin{lemma}
\label{spaces-morphisms-lemma-scheme-theoretic-image}
\begin{slogan}
代数空间态射的概形论像存在。
\end{slogan}
设 $S$ 为概形。设 $f : X \to Y$ 是 $S$ 上代数空间的态射。存在闭子空间
$Z \subset Y$，使得 $f$ 经由 $Z$ 分解，并且对任意其他闭子空间
$Z' \subset Y$，使 $f$ 经由 $Z'$ 分解，都有 $Z \subset Z'$。
\end{lemma}

\begin{proof}
令 $\mathcal{I} = \Ker(\mathcal{O}_Y \to f_*\mathcal{O}_X)$。若 $\mathcal{I}$ 是拟凝聚的，
则取由 $\mathcal{I}$ 定义的闭子概形 $Z$ 即可，见引理
\ref{spaces-morphisms-lemma-closed-immersion-ideals}。一般情形下，需要证明存在包含于
$\mathcal{I}$ 的最大拟凝聚理想层 $\mathcal{I}'$。这由引理
\ref{spaces-morphisms-lemma-largest-quasi-coherent-subsheaf} 得到。
\end{proof}

\noindent
假设上述引理 \ref{spaces-morphisms-lemma-scheme-theoretic-image} 的情形中 $X$ 和 $Y$ 可表示。
引理中得到的闭子空间 $Z \subset Y$ 与《态射》引理
\ref{morphisms-lemma-scheme-theoretic-image} 中得到的闭子概形 $Z \subset Y$ 一致。
原因是，闭子空间（或子概形）与 $X_\etale$ 和 $X$ 上的拟凝聚理想层成反包含对应。
由于 $X_\etale$ 与 $X$ 上拟凝聚模的范畴相同（《空间的性质》第
\ref{spaces-properties-section-quasi-coherent} 节），结论得证。因此下述定义与早先对概形态射的定义一致。

\begin{definition}
\label{spaces-morphisms-definition-scheme-theoretic-image}
设 $S$ 为概形。设 $f : X \to Y$ 是 $S$ 上代数空间的态射。$f$ 的
{\it 概形论像}是使 $f$ 经由其分解的最小闭子空间 $Z \subset Y$；见上面的引理
\ref{spaces-morphisms-lemma-scheme-theoretic-image}。
\end{definition}

\noindent
我们通常只用 $f : X \to Z$ 表示 $f$ 的分解。若态射 $f$ 不是拟紧的，
则一般而言，概形论像的构造与限制到 $Y$ 的开子空间不相容。

\begin{lemma}
\label{spaces-morphisms-lemma-quasi-compact-scheme-theoretic-image}
设 $S$ 为概形。设 $f : X \to Y$ 是 $S$ 上代数空间的态射。
设 $Z \subset Y$ 为 $f$ 的概形论像。若 $f$ 拟紧，则
\begin{enumerate}
\item 理想层
$\mathcal{I} = \Ker(\mathcal{O}_Y \to f_*\mathcal{O}_X)$
是拟凝聚的；
\item 概形论像 $Z$ 是与 $\mathcal{I}$ 对应的闭子空间；
\item 对任意 \'etale 态射 $V \to Y$，$X \times_Y V \to V$ 的概形论像等于
$Z \times_Y V$；以及
\item 像 $|f|(|X|) \subset |Z|$ 是 $|Z|$ 的稠密子集。
\end{enumerate}
\end{lemma}

\begin{proof}
由于 $\mathcal{I}$ 的构造与 \'etale 局部化相容，为证明 (3)，只需证明 (1) 和 (2)。
若 (1) 成立，则我们已在引理 \ref{spaces-morphisms-lemma-scheme-theoretic-image} 的证明中证明 (2)。
现在证明 $\mathcal{I}$ 为拟凝聚。拟凝聚性是 \'etale 局部性质，故可假设 $Y$ 为仿射概形。
由于 $f$ 拟紧，可取仿射概形 $U$ 及满射 \'etale 态射 $U \to X$。
记 $f'$ 为复合 $U \to X \to Y$。则 $f_*\mathcal{O}_X$ 是
$f'_*\mathcal{O}_U$ 的子层，从而 $\mathcal{I} = \Ker(\mathcal{O}_Y \to \mathcal{O}_{X'})$。
由引理 \ref{spaces-morphisms-lemma-pushforward}，层 $f'_*\mathcal{O}_U$ 在 $Y$ 上拟凝聚。
因此，作为凝聚模间映射的核，$\mathcal{I}$ 是拟凝聚的。
最后，(4) 由 (1)、(2) 和 (3) 得到，因为在 $|Y|$ 中任何不包含于
$|f|(|X|)$ 的闭包的点上，理想 $\mathcal{I}$ 都是单位理想。
\end{proof}

\begin{lemma}
\label{spaces-morphisms-lemma-scheme-theoretic-image-reduced}
设 $S$ 为概形。设 $f : X \to Y$ 是 $S$ 上代数空间的态射。假设 $X$ 约化。则
\begin{enumerate}
\item $f$ 的概形论像 $Z$ 是 $\overline{|f|(|X|)}$ 上的约化诱导代数空间结构；以及
\item 对任意 \'etale 态射 $V \to Y$，$X \times_Y V \to V$ 的概形论像等于
$Z \times_Y V$。
\end{enumerate}
\end{lemma}

\begin{proof}
(1) 成立，因为 $\overline{|f|(|X|)}$ 上的约化诱导代数空间结构是 $Y$ 中
使 $f$ 经由其分解的最小闭子空间；见《空间的性质》引理
\ref{spaces-properties-lemma-map-into-reduction}。
(2) 由 (1)、$|V| \to |Y|$ 为开映射，以及约化性在 \'etale 局部化下保持不变得到。
\end{proof}

\begin{lemma}
\label{spaces-morphisms-lemma-reach-points-scheme-theoretic-image}
设 $S$ 为概形。设 $f : X \to Y$ 是 $S$ 上代数空间的拟紧态射。
设 $Z$ 为 $f$ 的概形论像。设 $z \in |Z|$。存在分式域为 $K$ 的赋值环 $A$，
以及交换图
$$
\xymatrix{
\Spec(K) \ar[rr] \ar[d] & & X \ar[d] \ar[ld] \\
\Spec(A) \ar[r] & Z \ar[r] & Y
}
$$
使得 $\Spec(A)$ 的闭点映到 $z$。
\end{lemma}

\begin{proof}
取含点 $z' \in V$ 的仿射概形 $V$，以及把 $z'$ 映到 $z$ 的 \'etale 态射
$V \to Y$。令 $Z' \subset V$ 为 $X \times_Y V \to V$ 的概形论像。
由引理 \ref{spaces-morphisms-lemma-quasi-compact-scheme-theoretic-image}，有
$Z' = Z \times_Y V$，故 $z' \in Z'$。由于 $f$ 拟紧且 $V$ 仿射，
$X \times_Y V$ 拟紧。因此存在仿射概形 $W$ 及满射 \'etale 态射
$W \to X \times_Y V$。此时 $Z' \subset V$ 也是 $W \to V$ 的概形论像。
由《态射》引理 \ref{morphisms-lemma-reach-points-scheme-theoretic-image}，可取图
$$
\xymatrix{
\Spec(K) \ar[r] \ar[d] &
W \ar[r] \ar[d] &
X \times_Y V \ar[d] \ar[r] &
X \ar[d] \\
\Spec(A) \ar[r] &
Z' \ar[r] &
V \ar[r] &
Y
}
$$
使得 $\Spec(A)$ 的闭点映到 $z'$。再与 $Z' \to Z$ 及
$W \to X \times_Y V \to X$ 复合，即得所求。
\end{proof}

\begin{lemma}
\label{spaces-morphisms-lemma-factor-factor}
设 $S$ 为概形。设
$$
\xymatrix{
X_1 \ar[d] \ar[r]_{f_1} & Y_1 \ar[d] \\
X_2 \ar[r]^{f_2} & Y_2
}
$$
为 $S$ 上代数空间的交换图。令 $Z_i \subset Y_i$（$i = 1, 2$）为 $f_i$ 的概形论像。
则态射 $Y_1 \to Y_2$ 诱导态射 $Z_1 \to Z_2$，并得到交换图
$$
\xymatrix{
X_1 \ar[r] \ar[d] & Z_1 \ar[d] \ar[r] & Y_1 \ar[d] \\
X_2 \ar[r] & Z_2 \ar[r] & Y_2
}
$$
\end{lemma}

\begin{proof}
$Z_2$ 在 $Y_1$ 中的概形论逆像是 $Y_1$ 的闭子空间，且 $f_1$ 经由它分解。
因此 $Z_1$ 包含于该逆像中。引理得证。
\end{proof}

\begin{lemma}
\label{spaces-morphisms-lemma-scheme-theoretic-image-of-partial-section}
设 $S$ 为概形。设 $f : X \to Y$ 是 $S$ 上代数空间的分离态射。
设 $V \subset Y$ 为开子空间，且 $V \to Y$ 拟紧。设 $s : V \to X$ 为满足
$f \circ s = \text{id}_V$ 的态射。令 $Y'$ 为 $s$ 的概形论像。
则 $Y' \to Y$ 在 $V$ 上是同构。
\end{lemma}

\begin{proof}
由引理 \ref{spaces-morphisms-lemma-quasi-compact-permanence}，态射 $s : V \to X$ 拟紧。
因此由引理 \ref{spaces-morphisms-lemma-quasi-compact-scheme-theoretic-image}，$s$ 的概形论像 $Y'$ 的构造
与限制到开子空间相容。特别地，$Y' \cap f^{-1}(V)$ 是分离态射
$f^{-1}(V) \to V$ 的一个截面的概形论像。分离态射的截面是闭浸入
（引理 \ref{spaces-morphisms-lemma-section-immersion}），故
$Y' \cap f^{-1}(V) \to V$ 是同构，如所需。
\end{proof}








\section{概形论闭包与稠密性}
\label{spaces-morphisms-section-scheme-theoretic-closure}

\noindent
本节是《态射》第 \ref{morphisms-section-scheme-theoretic-closure} 节的代数空间版本。

\begin{lemma}
\label{spaces-morphisms-lemma-scheme-theoretically-dense-representable}
设 $S$ 为概形。设 $W \subset S$ 为概形论稠密的开子概形
（《态射》定义 \ref{morphisms-definition-scheme-theoretically-dense}）。
设 $f : X \to S$ 为概形态射，且平坦、局部有限表示并局部拟有限。
则 $f^{-1}(W)$ 在 $X$ 中概形论稠密。
\end{lemma}

\begin{proof}
我们将使用《态射》引理
\ref{morphisms-lemma-characterize-scheme-theoretically-dense} 的刻画。
设 $V \subset X$ 为开集，且 $g \in \Gamma(V, \mathcal{O}_V)$ 是在
$f^{-1}(W) \cap V$ 上限制为零的函数。需证明 $g = 0$。假设 $g \not = 0$ 以导出矛盾。
由《态射补篇》引理 \ref{more-morphisms-lemma-go-down-with-annihilators}，
可缩小 $V$，并找到开集 $U \subset S$ 使其落入交换图
$$
\xymatrix{
V \ar[r] \ar[d]_\pi & X \ar[d]^f \\
U \ar[r] & S,
}
$$
一个拟凝聚子层 $\mathcal{F} \subset \mathcal{O}_U$、整数 $r > 0$，以及单射
$\mathcal{O}_U$-模映射 $\mathcal{F}^{\oplus r} \to \pi_*\mathcal{O}_V$，
其像包含 $g|_V$。设 $(g_1, \ldots, g_r) \in \Gamma(U, \mathcal{F}^{\oplus r})$ 映到 $g$。
由于 $g|_{f^{-1}W \cap V} = 0$，可知 $g_i|_{W \cap U} = 0$。
又因 $\mathcal{F} \subset \mathcal{O}_U$ 且 $W$ 在 $S$ 中概形论稠密，故 $g_i = 0$。
这推出 $g = 0$，与假设矛盾。
\end{proof}

\begin{lemma}
\label{spaces-morphisms-lemma-scheme-theoretically-dense}
设 $S$ 为概形。设 $X$ 为 $S$ 上的代数空间。设 $U \subset X$ 为开子空间。
下列条件等价：
\begin{enumerate}
\item 对每个代数空间的 \'etale 态射 $\varphi : V \to X$，$\varphi^{-1}(U)$ 在 $V$ 中的
概形论闭包等于 $V$；
\item 存在概形 $V$ 及满射 \'etale 态射 $\varphi : V \to X$，使得 $\varphi^{-1}(U)$ 在 $V$ 中的
概形论闭包等于 $V$；
\end{enumerate}
\end{lemma}

\begin{proof}
注意，若 $V \to V'$ 是在 $X$ 上 \'etale 的代数空间态射，且 $Z \subset V$、
相应地 $Z' \subset V'$ 分别是 $U \times_X V$、$U \times_X V'$ 在 $V$、$V'$ 中的概形论闭包，
则 $Z$ 映到 $Z'$。因此若 $V \to V'$ 满射且 \'etale，则 $Z = V$ 推出 $Z' = V'$。
此外，注意 \'etale 态射是平坦、局部有限表示且局部拟有限的
（见《态射》第 \ref{morphisms-section-etale} 节）。因此由引理
\ref{spaces-morphisms-lemma-scheme-theoretically-dense-representable}，若 $V$ 与 $V'$ 为概形，
则 $Z' = V'$ 推出 $Z = V$。利用每个代数空间都有概形的 \'etale 覆盖作形式论证，
可知 (1) 与 (2) 等价。
\end{proof}

\noindent
由引理 \ref{spaces-morphisms-lemma-scheme-theoretically-dense}，下述定义与概形情形的定义相容。

\begin{definition}
\label{spaces-morphisms-definition-scheme-theoretically-dense}
设 $S$ 为概形。设 $X$ 为 $S$ 上的代数空间。设 $U \subset X$ 为开子空间。
\begin{enumerate}
\item 态射 $U \to X$ 的概形论像称为 $U$ 在 $X$ 中的 {\it 概形论闭包}。
\item 若引理 \ref{spaces-morphisms-lemma-scheme-theoretically-dense} 的等价条件成立，则称 $U$ 在 $X$ 中
{\it 概形论稠密}。
\end{enumerate}
\end{definition}

\noindent
按此定义，$U$ 在 $X$ 中概形论稠密并不等价于 $U$ 的概形论闭包为 $X$。
这略显不优雅，但在适当的有限性条件下，我们将看到它确实成立。

\begin{lemma}
\label{spaces-morphisms-lemma-scheme-theoretically-dense-quasi-compact}
设 $S$ 为概形，设 $X$ 为 $S$ 上的代数空间，设 $U \subset X$ 为开子空间。
若 $U \to X$ 拟紧，则 $U$ 在 $X$ 中概形论稠密，当且仅当 $U$ 在 $X$ 中的
概形论闭包为 $X$。
\end{lemma}

\begin{proof}
由引理 \ref{spaces-morphisms-lemma-quasi-compact-scheme-theoretic-image} 的 (3) 得到。
\end{proof}

\begin{lemma}
\label{spaces-morphisms-lemma-characterize-scheme-theoretically-dense}
设 $S$ 为概形。设 $j : U \to X$ 是 $S$ 上代数空间的开浸入。
则 $U$ 在 $X$ 中概形论稠密，当且仅当
$\mathcal{O}_X \to j_*\mathcal{O}_U$ 是单射。
\end{lemma}

\begin{proof}
若 $\mathcal{O}_X \to j_*\mathcal{O}_U$ 为单射，
则将其限制到任意在 $X$ 上 \'etale 的代数空间 $V$ 仍为单射。
因此 $U \times_X V$ 在 $V$ 中的概形论闭包等于 $V$，见引理
\ref{spaces-morphisms-lemma-scheme-theoretic-image} 的证明。
反之，假设 $U \times_X V$ 的概形论闭包等于 $V$，对所有在 $X$ 上 \'etale 的 $V$ 都成立。
若 $\mathcal{O}_X \to j_*\mathcal{O}_U$ 不是单射，则可找到一个仿射概形（记为 $V = \Spec(A)$），
它在 $X$ 上 \'etale，并有非零元素 $f \in A$，使得 $f$ 在
$\Gamma(V \times_X U, \mathcal{O})$ 中映为零。此时 $V \times_X U$ 在 $V$ 中的概形论闭包
显然包含于 $\Spec(A/(f))$，矛盾。
\end{proof}

\begin{lemma}
\label{spaces-morphisms-lemma-intersection-scheme-theoretically-dense}
设 $S$ 为概形，设 $X$ 为 $S$ 上的代数空间。若 $U$、$V$ 是 $X$ 中概形论稠密的
开子空间，则 $U \cap V$ 亦然。
\end{lemma}

\begin{proof}
设 $W \to X$ 为任意 \'etale 态射。考虑映射
$\mathcal{O}(W) \to \mathcal{O}(W \times_X V)
\to \mathcal{O}(W \times_X (V \cap U))$.
由引理 \ref{spaces-morphisms-lemma-characterize-scheme-theoretically-dense}，两个映射都是单射，
故复合映射为单射。再由引理 \ref{spaces-morphisms-lemma-characterize-scheme-theoretically-dense}，
$U \cap V$ 在 $X$ 中概形论稠密。
\end{proof}

\begin{lemma}
\label{spaces-morphisms-lemma-quasi-compact-immersion}
设 $S$ 为概形。设 $h : Z \to X$ 是 $S$ 上代数空间的浸入。假设 $Z \to X$ 拟紧或 $Z$ 约化。
设 $\overline{Z} \subset X$ 为 $h$ 的概形论像。则态射 $Z \to \overline{Z}$ 是开浸入，
并将 $Z$ 识别为 $\overline{Z}$ 中概形论稠密的开子空间。此外，$Z$ 在 $\overline{Z}$ 中拓扑稠密。
\end{lemma}

\begin{proof}
在两种情形中，$\overline{Z}$ 的构造都与 \'etale 局部化相容，见引理
\ref{spaces-morphisms-lemma-quasi-compact-scheme-theoretic-image} 和
\ref{spaces-morphisms-lemma-scheme-theoretic-image-reduced}。因此由概形情形可得本引理，见《态射》引理
\ref{morphisms-lemma-quasi-compact-immersion}。
\end{proof}

\begin{lemma}
\label{spaces-morphisms-lemma-equality-of-morphisms}
设 $S$ 为概形，设 $B$ 为 $S$ 上的代数空间。设 $f, g : X \to Y$ 为 $B$ 上代数空间的态射。
设 $U \subset X$ 为开子空间且 $f|_U = g|_U$。若 $U$ 在 $X$ 中的概形论闭包为 $X$，
且 $Y \to B$ 分离，则 $f = g$。
\end{lemma}

\begin{proof}
由于 $Y \to B$ 分离，纤维积
$Y \times_{\Delta, Y \times_B Y, (f, g)} X$ 是闭子空间 $Z \subset X$。
由于 $f|_U = g|_U$，有 $U \subset Z$。因假设 $U$ 在 $X$ 中概形论稠密，故 $Z = X$。
\end{proof}






\section{主导态射}
\label{spaces-morphisms-section-dominant}

\noindent
我们复制概形主导态射的定义，以得到代数空间主导态射的概念。注意，除非态射是拟紧的且代数空间满足某些分离性公理，该定义的性质不太好。

\begin{definition}
\label{spaces-morphisms-definition-dominant}
设 $S$ 为概形。若 $S$ 上代数空间的态射 $f : X \to Y$ 所诱导的
$|f| : |X| \to |Y|$ 的像在 $|Y|$ 中稠密，则称为{\it 主导}。
\end{definition}






\section{普遍单射态射}
\label{spaces-morphisms-section-universally-injective}

\noindent
在第 \ref{spaces-morphisms-section-representable} 节中，我们已经定义了代数空间的可表示态射普遍单射的含义。对 $S$ 上的域 $K$（回忆，这意味着给定结构态射 $\Spec(K) \to S$）以及 $S$ 上的代数空间 $X$，我们写作
$X(K) = \Mor_S(\Spec(K), X)$。我们先将可表示态射的条件转换为关于点函子的条件。

\begin{lemma}
\label{spaces-morphisms-lemma-universally-injective-representable}
设 $S$ 为概形。设 $f : X \to Y$ 为 $S$ 上代数空间的可表示态射。则 $f$ 普遍单射
（按第 \ref{spaces-morphisms-section-representable} 节的含义）当且仅当对所有域 $K$，映射
$X(K) \to Y(K)$ 是单射。
\end{lemma}

\begin{proof}
以下不再另行说明，我们将使用《态射》引理
\ref{morphisms-lemma-universally-injective}。
假设 $f$ 普遍单射。则对任意域 $K$ 及任意态射 $\Spec(K) \to Y$，概形态射
$\Spec(K) \times_Y X \to \Spec(K)$ 普遍单射。因此该态射至多有一个截面
$\Spec(K) \times_Y X \to \Spec(K)$，故 $X(K) \to Y(K)$ 为单射。反之，假设对每个域 $K$，映射 $X(K) \to Y(K)$ 为单射。
设 $T \to Y$ 为从概形到 $Y$ 的态射，考虑基变换 $f_T : T \times_Y X \to T$。
对任意域 $K$，有
$$
(T \times_Y X)(K) = T(K) \times_{Y(K)} X(K)
$$
按纤维积的定义，故 $X(K) \to Y(K)$ 的单射性保证
$(T \times_Y X)(K) \to T(K)$ 单射，这意味着 $f_T$ 普遍单射，如所需。
\end{proof}

\noindent
接下来，把域值点之间的变换为单射这一性质转化为更具几何性的条件。

\begin{lemma}
\label{spaces-morphisms-lemma-universally-injective}
设 $S$ 为概形。
设 $f : X \to Y$ 为 $S$ 上代数空间的态射。
以下各条件等价：
\begin{enumerate}
\item 对每个 $S$ 上的域 $K$，映射 $X(K) \to Y(K)$ 是单射；
\item 对每个 $S$ 上代数空间的态射 $Y' \to Y$，诱导映射
$|Y' \times_Y X| \to |Y'|$ 是单射；
\item 对角态射 $X \to X \times_Y X$ 是满射。
\end{enumerate}
\end{lemma}

\begin{proof}
假设 (1)。设 $g : Y' \to Y$ 为代数空间的态射，并记
$f' : Y' \times_Y X \to Y'$ 为 $f$ 的基变换。
设 $K_i$（$i = 1, 2$）为域，并设
$\varphi_i : \Spec(K_i) \to Y' \times_Y X$ 为态射，
使得 $f' \circ \varphi_1$ 与 $f' \circ \varphi_2$ 在 $|Y'|$ 中定义同一点。
按定义，这意味着存在域 $\Omega$ 以及嵌入
$\alpha_i : K_i \subset \Omega$，使得两个态射
$f' \circ \varphi_i \circ \alpha_i : \Spec(\Omega) \to Y'$ 相等。
相应的交换图如下：
$$
\xymatrix{
\Spec(\Omega) \ar@/_5ex/[ddrr] \ar[rd]^{\alpha_1} \ar[r]_{\alpha_2} &
\Spec(K_2) \ar[rd]^{\varphi_2} \\
& \Spec(K_1) \ar[r]^{\varphi_1} &
Y' \times_Y X  \ar[d]^{f'} \ar[r]^{g'} &
 X \ar[d]^f \\
& & Y' \ar[r]^g & Y.
}
$$
特别地，复合态射 $g \circ f' \circ \varphi_i \circ \alpha_i$
相等。由假设 (1)，这意味着态射
$g' \circ \varphi_i \circ \alpha_i$ 相等，其中 $g' : Y' \times_Y X \to X$
是投影。由纤维积的泛性质可知，态射
$\varphi_i \circ \alpha_i : \Spec(\Omega) \to Y' \times_Y X$ 相等。
换言之，$\varphi_1$ 与 $\varphi_2$ 在 $Y' \times_Y X$ 中定义同一点。
因此 (2) 成立。

\medskip\noindent
假设 (2)。设 $K$ 为 $S$ 上的域，设 $a, b : \Spec(K) \to X$
为满足 $f \circ a = f \circ b$ 的两个态射。记
$c : \Spec(K) \to Y$ 为公共值。由假设，
$|\Spec(K) \times_{c, Y} X| \to |\Spec(K)|$ 是单射。
这意味着存在域 $\Omega$ 以及嵌入
$\alpha_i : K \to \Omega$，使得
$$
\xymatrix{
\Spec(\Omega) \ar[r]_{\alpha_1} \ar[d]_{\alpha_2} &
\Spec(K) \ar[d]^a \\
\Spec(K) \ar[r]^-b &
\Spec(K) \times_{c, Y} X
}
$$
是交换的。与到 $\Spec(K)$ 的投影复合可得
$\alpha_1 = \alpha_2$。记公共值为 $\alpha$。
于是 $\{\alpha : \Spec(\Omega) \to \Spec(K)\}$
是 $\Spec(K)$ 的 fpqc 覆盖，并且两个态射
$a, b$ 在该覆盖的各成员上相等。由《代数空间的性质》命题
\ref{spaces-properties-proposition-sheaf-fpqc}
可知 $a = b$。因此 (1) 成立。

\medskip\noindent
假设 (3)。设 $x, x' \in |X|$ 为满足
$f(x) = f(x')$（在 $|Y|$ 中）的两个点。由《代数空间的性质》引理
\ref{spaces-properties-lemma-points-cartesian}
可知存在 $x'' \in |X \times_Y X|$，其两个投影
分别为 $x$ 与 $x'$。由假设以及《代数空间的性质》
引理 \ref{spaces-properties-lemma-characterize-surjective}
可知存在 $x''' \in |X|$，使得 $\Delta_{X/Y}(x''') = x''$。
因此 $x = x'$。换言之，$f$ 是单射。
由于条件 (3) 在基变换下保持不变，故 $f$ 满足 (2)。

\medskip\noindent
假设 (2)。特别地，$|X \times_Y X| \to |X|$ 是单射，
这立即推出
$|\Delta_{X/Y}| : |X| \to |X \times_Y X|$ 是满射；进而由《代数空间的性质》
引理 \ref{spaces-properties-lemma-characterize-surjective}
可知 $\Delta_{X/Y}$ 是满射。
\end{proof}

\noindent
根据上面的两个引理，以下定义与已经定义的代数空间可表示态射的普遍单射概念不冲突。

\begin{definition}
\label{spaces-morphisms-definition-universally-injective}
设 $S$ 为概形。设 $f : X \to Y$ 为 $S$ 上代数空间的态射。若对每个态射
$Y' \to Y$，诱导映射
$|Y' \times_Y X| \to |Y'|$ 都是单射，则称 $f$ 为{\it 普遍单射}。
\end{definition}

\noindent
为明确起见，这意味着引理 \ref{spaces-morphisms-lemma-universally-injective} 中的任一（或全部）等价条件均成立。

\begin{remark}
\label{spaces-morphisms-remark-universally-injective-not-separated}
概形的普遍单射态射是分离的，见《态射》引理
\ref{morphisms-lemma-universally-injective-separated}。
但代数空间的态射不一定如此。具体地，代数空间
$X = \mathbf{A}^1_k/\{x \sim -x \mid x \not = 0\}$
是在《代数空间》例 \ref{spaces-example-affine-line-involution} 中构造的，
它带有一个态射 $X \to \mathbf{A}^1_k$，将坐标为 $x$ 的点映到坐标为 $x^2$ 的点。
该态射在 $0$ 之外是同构，并且 $X$ 中位于 $0$ 上方的点唯一。由于 $X$ 不是分离的，
这是一个不分离的代数空间普遍单射态射。
\end{remark}

\begin{lemma}
\label{spaces-morphisms-lemma-base-change-universally-injective}
普遍单射态射的基变换仍是普遍单射。
\end{lemma}

\begin{proof}
省略。提示：这是形式性的。
\end{proof}

\begin{lemma}
\label{spaces-morphisms-lemma-universally-injective-local}
设 $S$ 为概形。
设 $f : X \to Y$ 为 $S$ 上代数空间的态射。
以下各条件等价：
\begin{enumerate}
\item $f$ 是普遍单射；
\item 对每个概形 $Z$ 及任意态射 $Z \to Y$，态射
$Z \times_Y X \to Z$ 是普遍单射；
\item 对每个仿射概形 $Z$ 及任意态射
$Z \to Y$，态射 $Z \times_Y X \to Z$ 是普遍单射；
\item 存在概形 $Z$ 及满射态射
$Z \to Y$，使得 $Z \times_Y X \to Z$ 是普遍单射；
\item 存在 Zariski 覆盖 $Y = \bigcup Y_i$，使得
每个态射 $f^{-1}(Y_i) \to Y_i$ 都是普遍单射。
\end{enumerate}
\end{lemma}

\begin{proof}
在本证明中，以下不再另行说明，我们使用普遍单射在基变换下保持这一事实
（引理 \ref{spaces-morphisms-lemma-base-change-universally-injective}）。
显然有 (1) $\Rightarrow$ (2) $\Rightarrow$ (3)
$\Rightarrow$ (4)。

\medskip\noindent
设 $g : Z \to Y$ 如 (4) 所述。设 $y : \Spec(K) \to Y$ 为
从域的谱到 $Y$ 的态射。由假设，可找到扩域
$\alpha : K \subset K'$ 以及态射
$z : \Spec(K') \to Z$，使得 $y \circ \alpha = g \circ z$
（这里对记号作了显然的滥用）。由假设，
态射 $Z \times_Y X \to Z$ 是普遍单射，因此
$g \circ z : \Spec(K') \to Y$ 到 $X$ 的态射提升至多有一个。
由于 $\{\alpha : \Spec(K') \to \Spec(K)\}$ 是
fpqc 覆盖，这意味着 $y : \Spec(K) \to Y$ 到 $X$ 的态射提升至多有一个，见
《代数空间的性质》命题
\ref{spaces-properties-proposition-sheaf-fpqc}。
因此 (1) 成立。

\medskip\noindent
我们省略 (5) 与 (1) 等价性的验证。
\end{proof}

\begin{lemma}
\label{spaces-morphisms-lemma-composition-universally-injective}
普遍单射态射的复合仍是普遍单射。
\end{lemma}

\begin{proof}
省略。
\end{proof}











\section{仿射态射}
\label{spaces-morphisms-section-affine}

\noindent
我们已经在第 \ref{spaces-morphisms-section-representable} 节定义了代数空间的可表示态射
为仿射的含义。

\begin{lemma}
\label{spaces-morphisms-lemma-affine-representable}
设 $S$ 为概形。设 $f : X \to Y$ 为 $S$ 上代数空间的可表示态射。则
$f$ 是仿射的（即第 \ref{spaces-morphisms-section-representable} 节的含义）
当且仅当对所有仿射概形 $Z$ 及态射 $Z \to Y$，概形 $X \times_Y Z$ 是仿射的。
\end{lemma}

\begin{proof}
这直接由概形仿射态射的定义得到（《态射》定义
\ref{morphisms-definition-affine}）。
\end{proof}

\noindent
这为下面的定义铺平了道路。

\begin{definition}
\label{spaces-morphisms-definition-affine}
设 $S$ 为概形。
设 $f : X \to Y$ 为 $S$ 上代数空间的态射。
若对每个仿射概形 $Z$ 及态射 $Z \to Y$，代数空间 $X \times_Y Z$
都可由仿射概形表示，则称 $f$ 为{\it 仿射}。
\end{definition}

\begin{lemma}
\label{spaces-morphisms-lemma-affine-local}
设 $S$ 为概形。
设 $f : X \to Y$ 为 $S$ 上代数空间的态射。
以下各条件等价：
\begin{enumerate}
\item $f$ 可表示且仿射；
\item $f$ 仿射；
\item 对每个仿射概形 $V$ 及 'etale 态射 $V \to Y$，
概形 $X \times_Y V$ 是仿射的；
\item 存在概形 $V$ 及满射 'etale 态射
$V \to Y$，使得 $V \times_Y X \to V$ 是仿射的；
\item 存在 Zariski 覆盖 $Y = \bigcup Y_i$，使得
每个态射 $f^{-1}(Y_i) \to Y_i$ 都是仿射的。
\end{enumerate}
\end{lemma}

\begin{proof}
显然 (1) 推出 (2)，(2) 推出 (3)；取 $V$ 为在 $Y$ 上 'etale 的仿射概形的不交并，
由此 (3) 推出 (4)，见《代数空间的性质》引理
\ref{spaces-properties-lemma-cover-by-union-affines}。
设 $V \to Y$ 如 (4) 所述。对 $V$ 的每个仿射开子概形 $W$，有
$W \times_Y X$ 是 $V \times_Y X$ 的仿射开子概形。因此由《代数空间的性质》引理
\ref{spaces-properties-lemma-subscheme}
可知 $V \times_Y X$ 是概形。此外，态射
$V \times_Y X \to V$ 是仿射的。于是可以应用《代数空间》引理
\ref{spaces-lemma-morphism-sheaves-with-P-effective-descent-etale}，
因为仿射态射类满足所有所需性质（见《态射》引理
\ref{morphisms-lemma-base-change-affine} 以及《下降》引理
\ref{descent-lemma-descending-property-affine} 与
\ref{descent-lemma-affine}）。应用该引理的结论是
$f$ 可表示且仿射，即 (1) 成立。

\medskip\noindent
由于仿射性在目标上是 Zariski 局部的，(1) 与 (5) 等价（上面的参考还表明，
仿射性事实上在目标上是 fpqc 局部的）。
\end{proof}

\begin{lemma}
\label{spaces-morphisms-lemma-composition-affine}
仿射态射的复合是仿射的。
\end{lemma}

\begin{proof}
省略。提示：纤维积的传递性。
\end{proof}

\begin{lemma}
\label{spaces-morphisms-lemma-base-change-affine}
仿射态射的基变换是仿射的。
\end{lemma}

\begin{proof}
省略。提示：纤维积的传递性。
\end{proof}

\begin{lemma}
\label{spaces-morphisms-lemma-closed-immersion-affine}
闭浸入是仿射的。
\end{lemma}

\begin{proof}
立即由概形态射的相应命题得到，见《态射》引理
\ref{morphisms-lemma-closed-immersion-affine}。
\end{proof}

\begin{lemma}
\label{spaces-morphisms-lemma-affine-equivalence-algebras}
设 $S$ 为概形。设 $X$ 为 $S$ 上的代数空间。
存在范畴的反等价
$$
\begin{matrix}
\text{代数空间} \\
\text{仿射于 }X
\end{matrix}
\longleftrightarrow
\begin{matrix}
\text{拟凝聚层} \\
\text{作为 }\mathcal{O}_X\text{-代数}
\end{matrix}
$$
它将 $f : Y \to X$ 对应到层 $f_*\mathcal{O}_Y$。
此外，该等价与任意基变换相容。
\end{lemma}

\begin{proof}
本引理是《态射》引理
\ref{morphisms-lemma-affine-equivalence-algebras}.
设 $\mathcal{A}$ 为拟凝聚的 $\mathcal{O}_X$-代数层。
我们将构造代数空间的仿射态射
 $\pi : Y = \underline{\Spec}_X(\mathcal{A}) \to X$，满足
$\pi_*\mathcal{O}_Y \cong \mathcal{A}$。为此，取一个概形
$U$ 以及满射 \'etale 态射 $\varphi : U \to X$。照例，记
$R = U \times_X U$，其投影为 $s, t : R \to U$。记
$\psi : R \to X$ 为复合 $\psi = \varphi \circ s = \varphi \circ t$。
由上述引理，存在概形的仿射态射
$\pi_0 : V \to U$ 与 $\pi_1 : W \to R$，满足
$\pi_{0, *}\mathcal{O}_V \cong \varphi^*\mathcal{A}$ 以及
$\pi_{1, *}\mathcal{O}_W \cong \psi^*\mathcal{A}$。
由于该构造与基变换相容，存在态射 $s', t' : W \to V$，使得下列图
$$
\vcenter{
\xymatrix{
W \ar[r]_{s'} \ar[d] & V \ar[d] \\
R \ar[r]^s & U
}
}
\quad\text{且}\quad
\vcenter{
\xymatrix{
W \ar[r]_{t'} \ar[d] & V \ar[d] \\
R \ar[r]^t & U
}
}
$$
这些图是笛卡尔的。因此 $s', t'$ 是 \'etale 态射。由上述事实形式地可得
$(t', s') : W \to V \times_S V$ 是单态射。我们省略验证
$W \to V \times_S V$ 是等价关系（提示：考虑将 $\mathcal{A}$
拉回到 $U \times_X U \times_X U = R \times_{s, U, t} R$）。
商层 $Y = V/W$ 是代数空间，见《代数空间》定理
\ref{spaces-theorem-presentation}。
由《群胚》引理 \ref{groupoids-lemma-criterion-fibre-product}
可知 $Y \times_X U \cong V$。因此由引理
\ref{spaces-morphisms-lemma-affine-local}，$Y \to X$ 是仿射的。最后，下面的同构
$$
(Y \times_X U \to U)_*\mathcal{O}_{Y \times_X U} =
\pi_{0, *}\mathcal{O}_V \cong \varphi^*\mathcal{A}
$$
与粘合同构相容，因此由《代数空间的性质》命题
\ref{spaces-properties-proposition-quasi-coherent}
有 $(Y \to X)_*\mathcal{O}_Y \cong \mathcal{A}$。
我们省略验证该构造与基变换相容。
\end{proof}

\begin{definition}
\label{spaces-morphisms-definition-relative-spec}
设 $S$ 为概形。设 $X$ 为 $S$ 上的代数空间。
设 $\mathcal{A}$ 为拟凝聚的
$\mathcal{O}_X$-代数层。$\mathcal{A}$ 在 $X$ 上的{\it 相对谱}，
或简称 $\mathcal{A}$ 在 $X$ 上的{\it 谱}，是仿射态射
$\underline{\Spec}(\mathcal{A}) \to X$，它在引理
\ref{spaces-morphisms-lemma-affine-equivalence-algebras} 的范畴等价下对应于 $\mathcal{A}$。
\end{definition}

\noindent
构造相对谱与任意基变换相容。

\begin{remark}
\label{spaces-morphisms-remark-factorization-quasi-compact-quasi-separated}
设 $S$ 为概形。设 $f : Y \to X$ 为 $S$ 上代数空间的拟紧且
拟分离态射。则 $f$ 有规范分解
$$
Y \longrightarrow \underline{\Spec}_X(f_*\mathcal{O}_Y) \longrightarrow X
$$
这是有意义的，因为由引理 \ref{spaces-morphisms-lemma-pushforward}，
$f_*\mathcal{O}_Y$ 是拟凝聚的。态射
$Y \to \underline{\Spec}_X(f_*\mathcal{O}_Y)$ 来自规范的
$\mathcal{O}_Y$-代数态射
$f^*f_*\mathcal{O}_Y \to \mathcal{O}_Y$，该态射又对应于一个在 $Y$ 上的规范态射
$Y \to Y \times_X \underline{\Spec}_X(f_*\mathcal{O}_Y)$（见
引理 \ref{spaces-morphisms-lemma-affine-equivalence-algebras}），从而得到上面的
$f$ 的分解。
\end{remark}

\begin{lemma}
\label{spaces-morphisms-lemma-affine-equivalence-modules}
设 $S$ 为概形。
设 $f : Y \to X$ 为 $S$ 上代数空间的仿射态射。
设 $\mathcal{A} = f_*\mathcal{O}_Y$。
函子 $\mathcal{F} \mapsto f_*\mathcal{F}$ 诱导出范畴等价
$$
\left\{
\begin{matrix}
\text{拟凝聚}\\
\mathcal{O}_Y\text{-模范畴}
\end{matrix}
\right\}
\longrightarrow
\left\{
\begin{matrix}
\text{拟凝聚}\\
\mathcal{A}\text{-模范畴}
\end{matrix}
\right\}
$$
此外，一个 $\mathcal{A}$-模作为 $\mathcal{O}_X$-模是拟凝聚的，当且仅当
它作为 $\mathcal{A}$-模是拟凝聚的。
\end{lemma}

\begin{proof}
省略。
\end{proof}

\begin{lemma}
\label{spaces-morphisms-lemma-affine-permanence}
设 $S$ 为概形。设 $B$ 为 $S$ 上的代数空间。
假设 $g : X \to Y$ 为 $B$ 上代数空间的态射。
\begin{enumerate}
\item 若 $X$ 在 $B$ 上仿射且 $\Delta : Y \to Y \times_B Y$ 仿射，
则 $g$ 仿射。
\item 若 $X$ 在 $B$ 上仿射且 $Y$ 在 $B$ 上分离，
则 $g$ 仿射。
\item 从仿射概形到对角态射在 $\mathbf{Z}$ 上仿射的代数空间的态射
（如《代数空间的性质》定义 \ref{spaces-properties-definition-separated}）是仿射的。
\item 从仿射概形到分离代数空间的态射是仿射的。
\end{enumerate}
\end{lemma}

\begin{proof}
证明 (1)。由引理 \ref{spaces-morphisms-lemma-base-change-affine}，基变换 $X \times_B Y \to Y$ 是仿射的。
态射 $(1, g) : X \to X \times_B Y$ 是 $Y \to Y \times_B Y$ 沿
$X \times_B Y \to Y \times_B Y$ 的基变换。因此由引理
\ref{spaces-morphisms-lemma-base-change-affine} 它是仿射的。
仿射态射的复合是仿射的（见引理 \ref{spaces-morphisms-lemma-composition-affine}），故 (1) 成立。
由于闭浸入是仿射的（见引理 \ref{spaces-morphisms-lemma-closed-immersion-affine}），且 $Y/B$ 分离
意味着 $\Delta$ 是闭浸入，(2) 由 (1) 得到。(3) 与 (4) 分别是 (1) 与 (2) 的特例。
\end{proof}

\begin{lemma}
\label{spaces-morphisms-lemma-Artinian-affine}
设 $S$ 为概形。设 $X$ 为 $S$ 上拟分离的代数空间。
设 $A$ 为 Artin 环。任意态射 $\Spec(A) \to X$ 都是仿射的。
\end{lemma}

\begin{proof}
设 $U \to X$ 为 'etale 态射且 $U$ 仿射。要证明本引理，需证明
$\Spec(A) \times_X U$ 仿射，见引理 \ref{spaces-morphisms-lemma-affine-local}。
由于 $X$ 拟分离，概形 $\Spec(A) \times_X U$ 拟紧。此外，
投影态射 $\Spec(A) \times_X U \to \Spec(A)$ 是 'etale 的。
因此该态射的纤维是有限离散的，而且 $\Spec(A)$ 上的拓扑是离散的。
于是 $\Spec(A) \times_X U$ 是一个概形，其底层拓扑空间是有限离散集。
由《概形》引理 \ref{schemes-lemma-scheme-finite-discrete-affine} 即得结论。
\end{proof}












\section{准仿射态射}
\label{spaces-morphisms-section-quasi-affine}

\noindent
我们已经在第 \ref{spaces-morphisms-section-representable} 节定义了代数空间的可表示态射
为准仿射的含义。

\begin{lemma}
\label{spaces-morphisms-lemma-quasi-affine-representable}
设 $S$ 为概形。设 $f : X \to Y$ 为 $S$ 上代数空间的可表示
态射。则 $f$ 是准仿射的（即第 \ref{spaces-morphisms-section-representable} 节的含义）
当且仅当对所有仿射概形 $Z$ 及态射 $Z \to Y$，概形 $X \times_Y Z$ 是准仿射的。
\end{lemma}

\begin{proof}
这直接由概形准仿射态射的定义得到
（《态射》定义 \ref{morphisms-definition-quasi-affine}）。
\end{proof}

\noindent
这为下面的定义铺平了道路。

\begin{definition}
\label{spaces-morphisms-definition-quasi-affine}
设 $S$ 为概形。
设 $f : X \to Y$ 为 $S$ 上代数空间的态射。
若对每个仿射概形 $Z$ 及态射 $Z \to Y$，代数空间 $X \times_Y Z$
都可由准仿射概形表示，则称 $f$ 为{\it 准仿射}。
\end{definition}

\begin{lemma}
\label{spaces-morphisms-lemma-quasi-affine-local}
设 $S$ 为概形。
设 $f : X \to Y$ 为 $S$ 上代数空间的态射。
以下各条件等价：
\begin{enumerate}
\item $f$ 可表示且准仿射；
\item $f$ 准仿射；
\item 存在概形 $V$ 及满射 \'etale 态射
$V \to Y$，使得 $V \times_Y X \to V$ 是准仿射的；
\item 存在 Zariski 覆盖 $Y = \bigcup Y_i$，使得
每个态射 $f^{-1}(Y_i) \to Y_i$ 都是准仿射的。
\end{enumerate}
\end{lemma}

\begin{proof}
显然 (1) 推出 (2)；取 $V$ 为在 $Y$ 上 'etale 的仿射概形的不交并，
由《代数空间的性质》引理
\ref{spaces-properties-lemma-cover-by-union-affines}，(2) 推出 (3)。
设 $V \to Y$ 如 (3) 所述。对 $V$ 的每个仿射开子概形 $W$，有
$W \times_Y X$ 是 $V \times_Y X$ 的准仿射开子概形。因此由《代数空间的性质》引理
\ref{spaces-properties-lemma-subscheme}
可知 $V \times_Y X$ 是概形。此外态射
$V \times_Y X \to V$ 是准仿射的。于是可以应用《代数空间》引理
\ref{spaces-lemma-morphism-sheaves-with-P-effective-descent-etale}，
因为准仿射态射类满足所有所需性质（见《态射》引理
\ref{morphisms-lemma-base-change-quasi-affine} 以及《下降》引理
\ref{descent-lemma-descending-property-quasi-affine} 与
\ref{descent-lemma-quasi-affine}）。应用该引理的结论是
$f$ 可表示且准仿射，即 (1) 成立。

\medskip\noindent
由于准仿射性在目标上是 Zariski 局部的，(1) 与 (4) 等价（上面的参考还表明，
准仿射性事实上在目标上是 fpqc 局部的）。
\end{proof}

\begin{lemma}
\label{spaces-morphisms-lemma-composition-quasi-affine}
准仿射态射的复合是准仿射的。
\end{lemma}

\begin{proof}
省略。
\end{proof}

\begin{lemma}
\label{spaces-morphisms-lemma-base-change-quasi-affine}
准仿射态射的基变换是准仿射的。
\end{lemma}

\begin{proof}
省略。
\end{proof}

\begin{lemma}
\label{spaces-morphisms-lemma-characterize-quasi-affine}
设 $S$ 为概形。
代数空间的拟紧且拟分离态射
$f : Y \to X$ 是准仿射的，当且仅当其规范分解
$Y \to \underline{\Spec}_X(f_*\mathcal{O}_Y)$
（注记 \ref{spaces-morphisms-remark-factorization-quasi-compact-quasi-separated}）
是开浸入。
\end{lemma}

\begin{proof}
设 $U \to X$ 为满射态射，其中 $U$ 是概形。由于可以在沿 $U$ 的基变换后检验
$f$ 是否准仿射（引理 \ref{spaces-morphisms-lemma-quasi-affine-local}），且
$f_*\mathcal{O}_Y|_V$ 等于 $(Y \times_X U \to U)_*\mathcal{O}_{Y \times_X U}$
（《代数空间的性质》引理
\ref{spaces-properties-lemma-pushforward-etale-base-change-modules}），并且
相对谱的构造与基变换相容（引理 \ref{spaces-morphisms-lemma-affine-equivalence-algebras}），
可知断言归结为 $X$ 为概形的情形。
若 $X$ 为概形且 $f$ 准仿射，或
$Y \to \underline{\Spec}_X(f_*\mathcal{O}_Y)$ 为开浸入，
则 $Y$ 也是概形。因此问题归结为《态射》引理
\ref{morphisms-lemma-characterize-quasi-affine}。
\end{proof}







\section{源与目标上 \'etale 局部的态射类型}
\label{spaces-morphisms-section-local-source-target}

\noindent
给定一种概形态射的性质，它在源与目标上是{\it \'etale 局部的}（见
《下降》定义 \ref{descent-definition-local-source-target}），
我们可以用它定义代数空间态射的相应性质，即施加下面引理中的任一等价条件。

\begin{lemma}
\label{spaces-morphisms-lemma-local-source-target}
设 $\mathcal{P}$ 为概形态射的一种在源与目标上 \'etale 局部的性质。
设 $S$ 为概形。
设 $f : X \to Y$ 为 $S$ 上代数空间的态射。
考虑交换图
$$
\xymatrix{
U \ar[d]_a \ar[r]_h & V \ar[d]^b \\
X \ar[r]^f & Y
}
$$
其中 $U$ 与 $V$ 为概形，竖直箭头为 \'etale 态射。
以下条件等价：
\begin{enumerate}
\item 对任意上述图，态射 $h$ 具有性质
$\mathcal{P}$；
\item 对某个上述图（其中 $a : U \to X$ 为满射），
态射 $h$ 具有性质 $\mathcal{P}$。
\end{enumerate}
若 $X$ 与 $Y$ 可表示，这还等价于 $f$（作为概形态射）具有性质 $\mathcal{P}$。
若 $\mathcal{P}$ 还在任意基变换下保持，且在基上 fppf 局部，
那么对于可表示态射 $f$，这也等价于 $f$ 按第 \ref{spaces-morphisms-section-representable} 节的含义
具有性质 $\mathcal{P}$。
\end{lemma}

\begin{proof}
我们证明 (1) 与 (2) 的等价性。
(1) $\Rightarrow$ (2) 立即成立（注意《代数空间》引理
\ref{spaces-lemma-lift-morphism-presentations}）。
设
$$
\xymatrix{
U \ar[d] \ar[r]_h & V \ar[d] \\
X \ar[r]^f & Y
}
\quad\quad
\xymatrix{
U' \ar[d] \ar[r]_{h'} & V' \ar[d] \\
X \ar[r]^f & Y
}
$$
为引理中的两个图。假设 $U \to X$ 满射且 $h$ 具有性质 $\mathcal{P}$。
为证明 (2) 推出 (1)，需证明 $h'$ 具有 $\mathcal{P}$。为此考虑图
$$
\xymatrix{
U \ar[d]_h &
U \times_X U' \ar[l] \ar[d]^{(h, h')} \ar[r] &
U' \ar[d]^{h'} \\
V &
V \times_Y V' \ar[l] \ar[r] &
V'
}
$$
由《下降》引理 \ref{descent-lemma-local-source-target-characterize}，
$h$ 具有 $\mathcal{P}$ 推出 $(h, h')$ 具有 $\mathcal{P}$；又由于
$U \times_X U' \to U'$ 是满射，由同一引理可知 $h'$ 具有 $\mathcal{P}$。

\medskip\noindent
若 $X$ 与 $Y$ 可表示，则可应用《下降》引理
\ref{descent-lemma-local-source-target-characterize}，从而 (1) 与 (2) 等价于 $f$
具有 $\mathcal{P}$。

\medskip\noindent
最后，设 $f$ 可表示，且 $U,V,a,b,h$ 如引理 (2) 所述，并假设
$\mathcal{P}$ 在任意基变换下保持。我们需证明对任意概形
$Z$ 及态射 $Z \to X$，基变换 $Z \times_Y X \to Z$
具有性质 $\mathcal{P}$。考虑图
$$
\xymatrix{
Z \times_Y U \ar[d] \ar[r] &
Z \times_Y V \ar[d] \\
Z \times_Y X \ar[r] &
Z
}
$$
注意上方水平箭头是 $h$ 的基变换，故具有性质 $\mathcal{P}$。
左侧竖直箭头是满射 \'etale，右侧竖直箭头是 \'etale。因此再次应用《下降》引理
\ref{descent-lemma-local-source-target-characterize}，可知 $Z \times_Y X \to Z$
具有性质 $\mathcal{P}$。
\end{proof}

\begin{definition}
\label{spaces-morphisms-definition-P}
设 $S$ 为概形。
设 $\mathcal{P}$ 为概形态射的一种在源与目标上 \'etale 局部的性质。
若引理 \ref{spaces-morphisms-lemma-local-source-target} 的等价条件成立，
则称 $S$ 上代数空间的态射 $f : X \to Y${\it 具有性质 $\mathcal{P}$}。
\end{definition}

\noindent
下面是几个显然的注记。

\begin{remark}
\label{spaces-morphisms-remark-composition-P}
设 $S$ 为概形。设 $\mathcal{P}$ 为概形态射的一种在源与目标上 \'etale 局部的性质，
并且假设 $\mathcal{P}$ 在复合下稳定。则具有性质 $\mathcal{P}$ 的代数空间态射类
在复合下稳定。
\end{remark}

\begin{remark}
\label{spaces-morphisms-remark-base-change-P}
设 $S$ 为概形。设 $\mathcal{P}$ 为概形态射的一种在源与目标上 \'etale 局部的性质，
并且假设 $\mathcal{P}$ 在基变换下稳定。则具有性质 $\mathcal{P}$ 的代数空间态射类
在基变换下稳定。
\end{remark}

\noindent
给定一种概形芽态射的性质，它在源与目标上是{\it \'etale 局部的}（见
《下降》定义 \ref{descent-definition-local-source-target-at-point}），
我们可以用它定义代数空间在一点处的态射性质，即施加下面引理中的任一等价条件。

\begin{lemma}
\label{spaces-morphisms-lemma-local-source-target-at-point}
设 $\mathcal{Q}$ 为概形芽态射的一种在源与目标上 \'etale 局部的性质。
设 $S$ 为概形。
设 $f : X \to Y$ 为 $S$ 上代数空间的态射。
设 $x \in |X|$ 为 $X$ 的一点。
考虑图
$$
\xymatrix{
U \ar[d]_a \ar[r]_h & V \ar[d]^b \\
X \ar[r]^f & Y
}
\quad\quad
\xymatrix{
u \ar[d] \ar[r] & v \ar[d] \\
x \ar[r] & y
}
$$
其中 $U$ 与 $V$ 为概形，$a,b$ 为 \'etale 态射，且 $u,v,x,y$
分别是相应空间的点。
以下条件等价：
\begin{enumerate}
\item 对任意上述图，有 $\mathcal{Q}((U, u) \to (V, v))$；
\item 对某个上述图，有 $\mathcal{Q}((U, u) \to (V, v))$。
\end{enumerate}
若 $X$ 与 $Y$ 可表示，这还等价于 $\mathcal{Q}((X, x) \to (Y, y))$。
\end{lemma}

\begin{proof}
省略。提示：与引理 \ref{spaces-morphisms-lemma-local-source-target} 的证明非常相似。
\end{proof}

\begin{definition}
\label{spaces-morphisms-definition-P-at-point}
设 $\mathcal{Q}$ 为概形芽态射的一种在源与目标上 \'etale 局部的性质。
设 $S$ 为概形。给定 $S$ 上代数空间的态射 $f : X \to Y$ 及一点
$x \in |X|$，若引理 \ref{spaces-morphisms-lemma-local-source-target-at-point} 的等价条件成立，
则称 $f${\it 在 $x$ 处具有性质 $\mathcal{Q}$}。
\end{definition}

\noindent
下面的引理不应被盲目地用来把态射的性质转化为一点处态射的性质。
例如，若 $\mathcal{P}$ 是平坦性，则下面引理中的性质
$Q$ 意味着“$f$ 在 $x$ 的某个开邻域内平坦”，这不同于“$f$ 在 $x$ 处平坦”。

\begin{lemma}
\label{spaces-morphisms-lemma-local-source-target-global-implies-local}
设 $\mathcal{P}$ 为概形态射的一种在源与目标上 \'etale 局部的性质。
考虑《下降》引理
\ref{descent-lemma-local-source-target-global-implies-local} 中与 $\mathcal{P}$
关联的芽态射性质 $\mathcal{Q}$。则
\begin{enumerate}
\item $\mathcal{Q}$ 在源与目标上 \'etale 局部。
\item 给定代数空间态射 $f : X \to Y$ 及 $x \in |X|$，以下条件等价：
\begin{enumerate}
\item $f$ 在 $x$ 处具有 $\mathcal{Q}$；
\item 存在 $x$ 的开邻域 $X' \subset X$，使得 $X' \to Y$ 具有 $\mathcal{P}$。
\end{enumerate}
\item 给定代数空间态射 $f : X \to Y$，以下条件等价：
\begin{enumerate}
\item $f$ 具有 $\mathcal{P}$；
\item 对每个 $x \in |X|$，态射 $f$ 在 $x$ 处具有 $\mathcal{Q}$。
\end{enumerate}
\end{enumerate}
\end{lemma}

\begin{proof}
关于 (1) 见《下降》引理
\ref{descent-lemma-local-source-target-global-implies-local}。(1)(a) $\Rightarrow$ (2)(b)
是因为对引理 \ref{spaces-morphisms-lemma-local-source-target-at-point} 中的图取
$X' = a(U) \subset X$。 (2)(b) $\Rightarrow$ (1)(a) 显然。
(3)(a) 与 (3)(b) 的等价性来自概形态射的相应结果，见《下降》引理
\ref{descent-lemma-local-source-target-local-implies-global}。
\end{proof}

\begin{remark}
\label{spaces-morphisms-remark-when-apply}
我们将把上面的引理 \ref{spaces-morphisms-lemma-local-source-target-global-implies-local}
应用于《下降》注记 \ref{descent-remark-list-local-source-target} 列出的所有情形，
但不包括“平坦”。在每种情形中，我们定义：若 $f$ 在 $x$ 的某个邻域内具有
$\mathcal{P}$，则称 $f$ 在 $x$ 处具有性质 $\mathcal{P}$。
\end{remark}




\section{有限型态射}
\label{spaces-morphisms-section-finite-type}

\noindent
概形态射的“局部有限型”性质在源与目标上 \'etale 局部，见
《下降》注记 \ref{descent-remark-list-local-source-target}。
它还在基变换下保持，并在目标上 fpqc 局部，见《态射》引理
\ref{morphisms-lemma-base-change-finite-type} 以及《下降》引理
\ref{descent-lemma-descending-property-locally-finite-type}。
因此由上面的引理 \ref{spaces-morphisms-lemma-local-source-target}，我们可以如下定义代数空间态射
为局部有限型；当态射可表示时，这与第 \ref{spaces-morphisms-section-representable} 节已有的定义一致。

\begin{definition}
\label{spaces-morphisms-definition-locally-finite-type}
设 $S$ 为概形。
设 $f : X \to Y$ 为 $S$ 上代数空间的态射。
\begin{enumerate}
\item 若引理 \ref{spaces-morphisms-lemma-local-source-target} 的等价条件在
$\mathcal{P} = \text{局部有限型}$ 下成立，则称 $f$
{\it 局部有限型}。
\item 设 $x \in |X|$。若存在 $x$ 的开邻域 $X' \subset X$，使得
$f|_{X'} : X' \to Y$ 是局部有限型，则称 $f${\it 在 $x$ 处有限型}。
\item 若 $f$ 局部有限型且拟紧，则称{\it 为有限型}。
\end{enumerate}
\end{definition}

\noindent
考虑《代数空间》例 \ref{spaces-example-affine-line-translation} 中的代数空间
$\mathbf{A}^1_k/\mathbf{Z}$。
态射 $\mathbf{A}^1_k/\mathbf{Z} \to \Spec(k)$ 是有限型的。

\begin{lemma}
\label{spaces-morphisms-lemma-composition-finite-type}
有限型态射的复合是有限型的。
局部有限型也有同样结论。
\end{lemma}

\begin{proof}
见注记 \ref{spaces-morphisms-remark-composition-P} 及《态射》引理
\ref{morphisms-lemma-composition-finite-type}。
\end{proof}

\begin{lemma}
\label{spaces-morphisms-lemma-base-change-finite-type}
有限型态射的基变换是有限型的。
局部有限型也有同样结论。
\end{lemma}

\begin{proof}
见注记 \ref{spaces-morphisms-remark-base-change-P} 及《态射》引理
\ref{morphisms-lemma-base-change-finite-type}。
\end{proof}

\begin{lemma}
\label{spaces-morphisms-lemma-finite-type-local}
设 $S$ 为概形。
设 $f : X \to Y$ 为 $S$ 上代数空间的态射。
以下各条件等价：
\begin{enumerate}
\item $f$ 局部有限型；
\item 对每个 $x \in |X|$，态射 $f$ 在 $x$ 处有限型；
\item 对每个概形 $Z$ 及任意态射 $Z \to Y$，态射
$Z \times_Y X \to Z$ 局部有限型；
\item 对每个仿射概形 $Z$ 及任意态射
$Z \to Y$，态射 $Z \times_Y X \to Z$ 局部有限型；
\item 存在概形 $V$ 及满射 \'etale 态射
$V \to Y$，使得 $V \times_Y X \to V$ 局部有限型；
\item 存在概形 $U$ 及满射 \'etale 态射
$\varphi : U \to X$，使得复合 $f \circ \varphi$
局部有限型；
\item 对每个交换图
$$
\xymatrix{
U \ar[d] \ar[r] & V \ar[d] \\
X \ar[r] & Y
}
$$
其中 $U$ 与 $V$ 为概形且竖直箭头为 \'etale，
上方水平箭头局部有限型；
\item 存在交换图
$$
\xymatrix{
U \ar[d] \ar[r] & V \ar[d] \\
X \ar[r] & Y
}
$$
其中 $U$ 与 $V$ 为概形，竖直箭头为 \'etale，$U \to X$
为满射，且上方水平箭头局部有限型；
\item 存在 Zariski 覆盖 $Y = \bigcup_{i \in I} Y_i$，
以及 $f^{-1}(Y_i) = \bigcup X_{ij}$，使得
每个态射 $X_{ij} \to Y_i$ 都局部有限型。
\end{enumerate}
\end{lemma}

\begin{proof}
条件 (2)、(3)、(4)、(5)、(6)、(7) 与 (9) 均以直接方式推出条件 (8)。例如，若 (5) 成立，
则可取概形 $V$ 为仿射概形的不交并，并取满射态射 $V \to Y$
（见《代数空间的性质》引理
\ref{spaces-properties-lemma-cover-by-union-affines}）。
由 (5)，$V \times_Y X \to V$ 局部有限型。
取概形 $U$ 及满射 \'etale 态射
$U \to V \times_Y X$。由引理 \ref{spaces-morphisms-lemma-composition-finite-type}，
$U \to V$ 局部有限型。因此 (8) 成立。

\medskip\noindent
条件 (1)、(7) 与 (8) 按定义等价。

\medskip\noindent
最后，我们证明 (1) 推出 (2)、(3)、(4)、(5)、(6) 与 (9)。(2) 立即成立。
(3)、(4)、(5) 与 (9) 来自局部有限型在基变换下保持这一事实，见
引理 \ref{spaces-morphisms-lemma-base-change-finite-type}。
对于 (6)，取 $U = X$ 即可。
\end{proof}

\begin{lemma}
\label{spaces-morphisms-lemma-locally-finite-type-locally-noetherian}
设 $S$ 为概形。
设 $f : X \to Y$ 为 $S$ 上代数空间的态射。
若 $f$ 局部有限型且 $Y$ 局部 Noether，则 $X$ 局部 Noether。
\end{lemma}

\begin{proof}
设有
$$
\xymatrix{
U \ar[d] \ar[r] & V \ar[d] \\
X \ar[r] & Y
}
$$
为交换图，其中 $U$ 与 $V$ 为概形且竖直箭头为满射 \'etale。若 $f$ 局部有限型，则
$U \to V$ 局部有限型。若 $Y$ 局部 Noether，则 $V$ 局部 Noether。由《态射》引理
\ref{morphisms-lemma-finite-type-noetherian}
可知 $U$ 局部 Noether，这意味着 $X$ 局部 Noether。
\end{proof}

\begin{lemma}
\label{spaces-morphisms-lemma-permanence-finite-type}
设 $S$ 为概形。
设 $f : X \to Y$、$g : Y \to Z$ 为 $S$ 上代数空间的态射。
若 $g \circ f : X \to Z$ 局部有限型，则 $f : X \to Y$
局部有限型。
\end{lemma}

\begin{proof}
可取交换图
$$
\xymatrix{
U \ar[r] \ar[d] & V \ar[r] \ar[d] & W \ar[d] \\
X \ar[r] & Y \ar[r] & Z
}
$$
其中 $U$、$V$、$W$ 为概形，竖直箭头为满射 \'etale，见
《代数空间》引理 \ref{spaces-lemma-lift-morphism-presentations}。
此时可应用引理 \ref{spaces-morphisms-lemma-finite-type-local}
以及《态射》引理 \ref{morphisms-lemma-permanence-finite-type} 得出结论。
\end{proof}

\begin{lemma}
\label{spaces-morphisms-lemma-immersion-locally-finite-type}
浸入是局部有限型的。
\end{lemma}

\begin{proof}
由一般原则《代数空间》引理
\ref{spaces-lemma-representable-transformations-property-implication}
以及《态射》引理 \ref{morphisms-lemma-immersion-locally-finite-type} 得到。
\end{proof}







\section{点与几何点}
\label{spaces-morphisms-section-points-fields}

\noindent
本节对点与几何点作一些注记（见《代数空间的性质》定义
\ref{spaces-properties-definition-geometric-point}）。
理解 $X$ 的几何点的一种方式是：
考虑 $S$ 的几何点 $\overline{s} : \Spec(k) \to S$
以及将 $\overline{s}$ 提升为到 $X$ 的态射 $\overline{x}$。
相应图为
$$
\xymatrix{
\Spec(k) \ar[r]_-{\overline{x}} \ar[rd]_{\overline{s}} & X \ar[d] \\
& S.
}
$$
我们常说“设 $k$ 为 $S$ 上的代数闭域”，表示 $\Spec(k)$ 带有态射
$\Spec(k) \to S$。此时记
$$
X(k) = \Mor_S(\Spec(k), X) =
\{\overline{x} \in X\text{ 位于 }\overline{s}\}
$$
这是 $X$ 的 $k$-值点集。在这种情况下，映射
$X(k) \to |X|$ 的像包含于子集 $|X_s| \subset |X|$ 中。
这里 $X_s = \Spec(\kappa(s)) \times_S X$，其中 $s \in S$
是与 $\overline{s}$ 对应的点。由于
$\Spec(\kappa(s)) \to S$ 是单态射，其基变换
$X_s \to X$ 是单态射，且 $|X_s|$ 确实是 $|X|$ 的子集。

\begin{lemma}
\label{spaces-morphisms-lemma-locally-finite-type-surjective-geometric-points}
设 $S$ 为概形。设 $f : X \to Y$ 为代数空间的态射。
假设 $f$ 在 $S$ 上局部有限型。以下条件等价：
\begin{enumerate}
\item $f$ 是满射；
\item 对每个 $S$ 上的代数闭域 $k$，诱导映射
$X(k) \to Y(k)$ 是满射。
\end{enumerate}
\end{lemma}

\begin{proof}
取交换图
$$
\xymatrix{
U \ar[d] \ar[r] & V \ar[d] \\
X \ar[r] & Y
}
$$
其中 $U$、$V$ 为 $S$ 上的概形，竖直箭头为满射且 \'etale，见
《代数空间》引理 \ref{spaces-lemma-lift-morphism-presentations}。
由于 $f$ 局部有限型，$U \to V$ 局部有限型。

\medskip\noindent
假设 (1)，并设 $\overline{y} \in Y(k)$。由引理
\ref{spaces-morphisms-lemma-composition-surjective} 与 \ref{spaces-morphisms-lemma-composition-finite-type}，
$U \to Y$ 是满射且局部有限型。
令 $Z = U \times_{Y, \overline{y}} \Spec(k)$。这是一个概形。
由引理 \ref{spaces-morphisms-lemma-base-change-surjective} 与
\ref{spaces-morphisms-lemma-base-change-finite-type}，投影 $Z \to \Spec(k)$
是满射且局部有限型。
由《簇》引理 \ref{varieties-lemma-locally-finite-type-Jacobson}，
$Z$ 有一个 $k$-值点 $\overline{z}$。点 $\overline{z}$ 的像
$\overline{x} \in X(k)$ 映到 $\overline{y}$，即得所需结论。

\medskip\noindent
假设 (2)。由《代数空间的性质》引理
\ref{spaces-properties-lemma-characterize-surjective}，只需证明 $|X| \to |Y|$ 满射。
设 $y \in |Y|$。取 $u \in U$，其映到 $y$。
取 $k \supset \kappa(u)$ 为代数闭包。记相应的点为
$\overline{u} \in U(k)$，其像为 $\overline{y} \in Y(k)$。
由假设，存在映到 $\overline{y}$ 的 $\overline{x} \in X(k)$。
显然，$\overline{x}$ 的像 $x \in |X|$ 映到 $y$。
\end{proof}

\noindent
为表述下一个引理，引入如下记号。对概形 $T$，记
$$
\lambda(T) = \sup\{\aleph_0, |\kappa(t)| ; t \in T\}.
$$
换言之，$\lambda(T)$ 是界住 $T$ 的所有剩余域基数的最小无限基数。
注意，若 $R$ 是环，则 $R$ 的任意剩余域 $\kappa(\mathfrak p)$
的基数都被 $R$ 的基数所界住（细节略）。
这意味着 $\lambda(T) \leq \text{size}(T)$，其中
$\text{size}(T)$ 是《集合》一节 \ref{sets-section-categories-schemes} 中引入的概形 $T$ 的大小。
若 $L/K$ 是有限生成域扩张，则
$|K| \leq |L| \leq \max\{\aleph_0, |K|\}$。因此，若 $T' \to T$
是局部有限型的概形态射，则 $\lambda(T') \leq \lambda(T)$；若 $T' \to T$ 还
是满射，则等号成立。接下来，设 $S$ 为概形，$X$ 为 $S$ 上的代数空间。此时定义
$$
\lambda(X) := \lambda(U)
$$
其中 $U$ 是任意一个在 $S$ 上、带有指向 $X$ 的满射 \'etale 态射的概形。
该定义与 $U$ 的选择无关，因为给定两个这样的概形 $U$ 与 $U'$，纤维积
$U \times_X U'$ 是概形，并且有分别指向 $U$ 与 $U'$ 的满射 \'etale 态射；
由上面的讨论，$\lambda(U) = \lambda(U \times_X U') =
\lambda(U')$。

\begin{lemma}
\label{spaces-morphisms-lemma-large-enough}
设 $S$ 为概形。设 $X$、$Y$ 为 $S$ 上的代数空间。
\begin{enumerate}
\item 当 $k$ 遍历所有 $S$ 上的代数闭域时，几何点
$\overline{y} \in Y(k)$ 的集合覆盖 $|Y|$。
\item 当 $k$ 遍历所有满足
$|k| \geq \lambda(Y)$ 且 $|k| > \lambda(X)$ 的 $S$ 上代数闭域时，几何点
$\overline{y} \in Y(k)$ 的集合覆盖 $|Y|$。
\item 对任意几何点
$\overline{s} : \Spec(k) \to S$，若
$k$ 的基数 $> \lambda(X)$，则映射
$$
X(k) \longrightarrow |X_s|
$$
是满射。
\item 设 $X \to Y$ 为 $S$ 上代数空间的态射。
对任意几何点 $\overline{s} : \Spec(k) \to S$，若
$k$ 的基数 $> \lambda(X)$，则映射
$$
X(k) \longrightarrow |X| \times_{|Y|} Y(k)
$$
是满射。
\item 设 $X \to Y$ 为 $S$ 上代数空间的态射。
以下条件等价：
\begin{enumerate}
\item 映射 $X \to Y$ 是满射；
\item 对所有满足 $|k| > \lambda(X)$ 且 $|k| \geq \lambda(Y)$ 的
$S$ 上代数闭域 $k$，映射 $X(k) \to Y(k)$ 是满射。
\end{enumerate}
\end{enumerate}
\end{lemma}

\begin{proof}
为证明 (1)，取满射 \'etale 态射 $V \to Y$，其中 $V$ 为概形。
对每个 $v \in V$，取代数闭包 $\kappa(v) \subset k_v$。考虑态射
$\overline{x} : \Spec(k_v) \to V \to Y$。
由 $|Y|$ 的构造，这些态射覆盖 $|Y|$。

\medskip\noindent
为证明 (2)，使用以下两个事实（证明略）：
(\romannumeral1) 若 $K$ 为域且 $\overline{K}$ 为其代数闭包，则
$|\overline{K}| \leq \max\{\aleph_0, |K|\}$。
(\romannumeral2) 对任意代数闭域 $k$ 及基数 $\aleph$，若
$\aleph \geq |k|$，则存在代数闭域扩张 $k'/k$，满足 $|k'| = \aleph$。
现在令 $\aleph = \max\{\lambda(X), \lambda(Y)\}^+$。
这里 $\lambda^+ > \lambda$ 表示下一个更大的基数，见《集合》一节
\ref{sets-section-cardinals}。
由 (\romannumeral1)，证明第一段中构造的域 $k_u$ 的基数都被
$\lambda(X)$ 所界住。因此由 (\romannumeral2)，可找到扩张
$k_u \subset k'_u$，使得 $|k'_u| = \aleph$。态射
$\overline{x}' : \Spec(k'_u) \to X$ 覆盖 $|X|$，如所需。
要真正完成 (2) 的证明，还需说明概形 $\Spec(k'_u)$
（同构于）$\Sch_{fppf}$ 的对象，因为我们的约定是所有概形都是
$\Sch_{fppf}$ 的对象；对集合论问题不感兴趣的读者可跳过本段其余部分。
由构造，存在 $\Sch_{fppf}$ 的对象 $T$，使得
$\lambda(X)$ 与 $\lambda(Y)$ 被 $\text{size}(T)$ 所界住。
按照我们在《拓扑》定义 \ref{topologies-definition-big-fppf-site} 中
将 $\Sch_{fppf}$ 构造为《集合》引理 \ref{sets-lemma-construct-category}
中构造的范畴 $\Sch_\alpha$ 的方式，大小满足 $\leq \text{size}(T)^+$ 的任意概形
都同构于 $\Sch_{fppf}$ 的对象。见《集合》方程
(\ref{sets-equation-bound}) 中函数 $Bound$ 的表达式。
由于 $\aleph \leq \text{size}(T)^+$，结论成立。

\medskip\noindent
第 (3) 部分中的记号 $X_s$ 表示纤维积
$\Spec(\kappa(s)) \times_S X$，其中 $s \in S$
是与 $\overline{s}$ 对应的点。因此 (2) 由 (4) 在
$Y = \Spec(\kappa(s))$ 时推出。

\medskip\noindent
证明 (4)。设 $X \to Y$ 为 $S$ 上代数空间的态射。
设 $k$ 为 $S$ 上基数满足
$> \lambda(X)$ 的代数闭域。设 $\overline{y} \in Y(k)$ 与 $x \in |X|$
映到 $|Y|$ 中的同一点 $y$。
我们要找一个 $\overline{x} \in X(k)$，使其映到 $x$ 与 $\overline{y}$。
取交换图
$$
\xymatrix{
U \ar[d] \ar[r] & V \ar[d] \\
X \ar[r] & Y
}
$$
其中 $U$、$V$ 为 $S$ 上的概形，竖直箭头为满射且 \'etale，见
《代数空间》引理 \ref{spaces-lemma-lift-morphism-presentations}。
取 $u \in |U|$ 映到 $x$，并记其像为 $v \in |V|$。
我们把 $u = \Spec(\kappa(u))$ 与
$v = \Spec(\kappa(v))$ 看作概形。
注意 $V \times_Y \Spec(k)$ 是在 $k$ 上 \'etale 的概形。
因此它是 $k$ 的有限可分扩张的谱的不交并，见
《态射》引理 \ref{morphisms-lemma-etale-over-field}。
由于 $v$ 映到 $y$，$v \times_Y \Spec(k)$ 是非空概形。
由于 $v \to V$ 是单态射，$v \times_Y \Spec(k) \to V \times_Y \Spec(k)$
也是单态射。因此由《概形》引理
\ref{schemes-lemma-mono-towards-spec-field}，$v \times_Y \Spec(k)$ 是
$k$ 的有限可分扩张的谱的不交并。
于是态射 $v \times_Y \Spec(k) \to \Spec(k)$ 有截面；也就是说，存在态射
$\overline{v} : \Spec(k) \to V$，它位于 $v$ 之上并位于
$\overline{y}$ 之上。最后考虑概形
$$
u \times_{V, \overline{v}} \Spec(k)
=
\Spec(\kappa(u) \otimes_{\kappa(v)} k)
$$
其中 $\kappa(v) \to k$ 是定义态射 $\overline{v}$ 的域映射。
由假设，$k$ 的基数大于 $\kappa(u)$ 的基数，因此可应用
《代数》引理 \ref{algebra-lemma-base-change-Jacobson}，
得到任意极大理想
$\mathfrak m \subset \kappa(u) \otimes_{\kappa(v)} k$
的剩余域都代数于 $k$，因而等于 $k$。
这样的极大理想给出态射
$\overline{u} : \Spec(k) \to U$，它位于 $u$ 之上并映到
$\overline{v}$。复合 $\Spec(k) \to U \to X$
就是所需的几何点 $\overline{x} \in X(k)$。
(4) 的证明完成。

\medskip\noindent
(5) 是 (2) 与 (4) 以及《代数空间的性质》引理
\ref{spaces-properties-lemma-characterize-surjective} 的形式推论。
\end{proof}








\section{有限型点}
\label{spaces-morphisms-section-points-finite-type}

\noindent
设 $S$ 为概形。设 $X$ 为 $S$ 上的代数空间。
有限型点 $x \in |X|$ 是可由局部有限型态射
$\Spec(k) \to X$ 表示的点。
对代数空间与代数栈而言，有限型点是闭点的适当替代。例如，总有“足够多”的有限型点。

\begin{lemma}
\label{spaces-morphisms-lemma-point-finite-type}
设 $S$ 为概形。设 $X$ 为 $S$ 上的代数空间。
设 $x \in |X|$。以下条件等价：
\begin{enumerate}
\item 存在局部有限型态射 $\Spec(k) \to X$，它表示 $x$。
\item 存在概形 $U$、闭点 $u \in U$ 以及 \'etale 态射
$\varphi : U \to X$，使得 $\varphi(u) = x$。
\end{enumerate}
\end{lemma}

\begin{proof}
设 $u \in U$ 与 $U \to X$ 如 (2) 所述。则
$\Spec(\kappa(u)) \to U$ 是有限型的，而 $U \to X$
可表示且局部有限型（由一般原则《代数空间》引理
\ref{spaces-lemma-representable-transformations-property-implication}
以及《态射》引理 \ref{morphisms-lemma-etale-locally-finite-presentation} 与
\ref{morphisms-lemma-finite-presentation-finite-type}）。因此由引理
\ref{spaces-morphisms-lemma-composition-finite-type}，(1) 成立。

\medskip\noindent
反之，假设 $\Spec(k) \to X$ 局部有限型并表示 $x$。
设 $U \to X$ 为满射 \'etale 态射，其中 $U$ 为概形。由假设，
$U \times_X \Spec(k) \to U$ 局部有限型。取
$U \times_X \Spec(k)$ 的有限型点 $v$（至少存在一个，见《态射》引理
\ref{morphisms-lemma-identify-finite-type-points}）。
由《态射》引理 \ref{morphisms-lemma-finite-type-points-morphism}，
$v$ 的像 $u \in U$ 是 $U$ 的有限型点。
因此由《态射》引理 \ref{morphisms-lemma-identify-finite-type-points}，
缩小 $U$ 后可假设 $u$ 是 $U$ 的闭点，即 (2) 成立。
\end{proof}

\begin{definition}
\label{spaces-morphisms-definition-finite-type-point}
设 $S$ 为概形。设 $X$ 为 $S$ 上的代数空间。
若引理 \ref{spaces-morphisms-lemma-point-finite-type} 的等价条件成立，
则称点 $x \in |X|$ 为{\it 有限型点}\footnote{这略微滥用了术语；更准确的说法或许是“局部有限型点”。}。
以 $X_{\text{ft-pts}}$ 表示 $X$ 的有限型点集。
\end{definition}

\noindent
有限型点集可描述如下。

\begin{lemma}
\label{spaces-morphisms-lemma-identify-finite-type-points}
设 $S$ 为概形。设 $X$ 为 $S$ 上的代数空间。则
$$
X_{\text{ft-pts}} =
\bigcup\nolimits_{\varphi : U \to X\text{ \'etale }} |\varphi|(U_0)
$$
其中 $U_0$ 是 $U$ 的闭点集。
这里可令 $U$ 遍历所有在 $X$ 上 \'etale 的概形，或所有在 $X$ 上 \'etale 的仿射概形。
\end{lemma}

\begin{proof}
立即由引理 \ref{spaces-morphisms-lemma-point-finite-type} 得到。
\end{proof}

\begin{lemma}
\label{spaces-morphisms-lemma-finite-type-points-morphism}
设 $S$ 为概形。
设 $f : X \to Y$ 为 $S$ 上代数空间的态射。
若 $f$ 局部有限型，则
$f(X_{\text{ft-pts}}) \subset Y_{\text{ft-pts}}$.
\end{lemma}

\begin{proof}
取 $x \in X_{\text{ft-pts}}$。以局部有限型态射
$x : \Spec(k) \to X$ 表示 $x$。由引理
\ref{spaces-morphisms-lemma-composition-finite-type}，$f \circ x$ 局部有限型。
因此 $f(x) \in Y_{\text{ft-pts}}$。
\end{proof}

\begin{lemma}
\label{spaces-morphisms-lemma-finite-type-points-surjective-morphism}
设 $S$ 为概形。
设 $f : X \to Y$ 为 $S$ 上代数空间的态射。
若 $f$ 局部有限型且满射，则
$f(X_{\text{ft-pts}}) = Y_{\text{ft-pts}}$.
\end{lemma}

\begin{proof}
由引理 \ref{spaces-morphisms-lemma-finite-type-points-morphism}，有
$f(X_{\text{ft-pts}}) \subset Y_{\text{ft-pts}}$。
设 $y \in |Y|$ 为有限型点。以局部有限型态射
$\Spec(k) \to Y$ 表示 $y$。
由于 $f$ 满射，代数空间
$X_k = \Spec(k) \times_Y X$ 非空，因而由引理
\ref{spaces-morphisms-lemma-identify-finite-type-points}，它有一个有限型点 $x \in |X_k|$。
态射 $X_k \to X$ 是 $\Spec(k) \to Y$ 的基变换，故局部有限型
（引理 \ref{spaces-morphisms-lemma-base-change-finite-type}）。
因此由引理 \ref{spaces-morphisms-lemma-finite-type-points-morphism}，$x$ 在 $X$ 中的像是有限型点，
而按构造它映到 $y$。
\end{proof}

\begin{lemma}
\label{spaces-morphisms-lemma-enough-finite-type-points}
设 $S$ 为概形。设 $X$ 为 $S$ 上的代数空间。
对任意局部闭子集 $T \subset |X|$，有
$$
T \not = \emptyset
\Rightarrow
T \cap X_{\text{ft-pts}} \not = \emptyset.
$$
特别地，对任意闭子集 $T \subset |X|$，
$T \cap X_{\text{ft-pts}}$ 在 $T$ 中稠密。
\end{lemma}

\begin{proof}
设 $i : Z \to X$ 为 $T$ 上诱导的既约子空间结构，见注记
\ref{spaces-morphisms-remark-space-structure-locally-closed-subset}。
任意浸入都局部有限型，见引理
\ref{spaces-morphisms-lemma-immersion-locally-finite-type}。
因此由引理 \ref{spaces-morphisms-lemma-finite-type-points-morphism}，有
$Z_{\text{ft-pts}} \subset X_{\text{ft-pts}} \cap T$。
最后，任意非空仿射概形 $U$ 若有指向 $Z$ 的 \'etale 态射，则至少有一个闭点。
因此由引理 \ref{spaces-morphisms-lemma-identify-finite-type-points}，$Z$ 至少有一个有限型点。
引理得证。
\end{proof}

\noindent
下面给出代数空间有限型点的另一个更具技术性的刻画。

\begin{lemma}
\label{spaces-morphisms-lemma-point-finite-type-monomorphism}
设 $S$ 为概形。设 $X$ 为 $S$ 上的代数空间。
设 $x \in |X|$。以下条件等价：
\begin{enumerate}
\item $x$ 是有限型点；
\item 存在代数空间 $Z$，其底层拓扑空间 $|Z|$ 为单点集，并存在局部有限型态射
$f : Z \to X$，使得 $\{x\} = |f|(|Z|)$；
\item 存在代数空间 $Z$ 与态射 $f : Z \to X$，满足以下性质：
\begin{enumerate}
\item 存在满射 \'etale 态射
$z : \Spec(k) \to Z$，其中 $k$ 为域；
\item $f$ 局部有限型；
\item $f$ 是单态射；
\item $x = f(z)$。
\end{enumerate}
\end{enumerate}
\end{lemma}

\begin{proof}
假设 $x$ 为有限型点。取仿射概形 $U$、闭点 $u \in U$ 以及 \'etale 态射
$\varphi : U \to X$，满足 $\varphi(u) = x$，见引理
\ref{spaces-morphisms-lemma-identify-finite-type-points}。
照例令 $u = \Spec(\kappa(u))$。投影态射
$u \times_X u \to u$ 是复合
$$
u \times_X u \to u \times_X U \to u \times_X X = u
$$
其中第一箭头是闭浸入（$u \to U$ 的基变换），第二箭头是 \'etale 态射
（\'etale 态射 $U \to X$ 的基变换）。因此 $u \times_X U$ 是
$k$ 的有限可分扩张的谱的不交并（见《态射》引理
\ref{morphisms-lemma-etale-over-field}），
从而闭子概形 $u \times_X u$ 是 $k$ 的有限可分扩张的不交并；也就是说，
$u \times_X u \to u$ 是 \'etale 态射。由《代数空间》定理
\ref{spaces-theorem-presentation}，$Z = u/u \times_X u$ 是代数空间。
按构造，图
$$
\xymatrix{
u \ar[d] \ar[r] & U \ar[d] \\
Z \ar[r] & X
}
$$
交换且竖直箭头为 \'etale。因此 $Z \to X$ 局部有限型（见引理
\ref{spaces-morphisms-lemma-finite-type-local}）。
按构造，态射 $Z \to X$ 是单态射，且 $z$ 的像为 $x$。因此 (3) 成立。

\medskip\noindent
显然 (3) 推出 (2)。
若 (2) 成立，则由引理 \ref{spaces-morphisms-lemma-finite-type-points-morphism}，$x$ 是 $X$ 的有限型点
（同时由引理 \ref{spaces-morphisms-lemma-enough-finite-type-points} 可知
$Z_{\text{ft-pts}}$ 非空，即 $Z$ 的唯一一点是 $Z$ 的有限型点）。
\end{proof}





\section{Nagata 空间}
\label{spaces-morphisms-section-nagata}

\noindent
Nagata 代数空间的定义见《空间的性质》第
\ref{spaces-properties-section-types-properties} 节。

\begin{lemma}
\label{spaces-morphisms-lemma-finite-type-nagata}
设 $S$ 为概形。
设 $f : X \to Y$ 为 $S$ 上代数空间的态射。
若 $Y$ 是 Nagata 的且 $f$ 局部有限型，则 $X$ 是 Nagata 的。
\end{lemma}

\begin{proof}
取概形 $V$ 及满 \'etale 态射 $V \to Y$。
取概形 $U$ 及满 \'etale 态射 $U \to X \times_Y V$。
若 $Y$ 是 Nagata 的，则 $V$ 是 Nagata 概形。
若 $X \to Y$ 局部有限型，则 $U \to V$ 局部有限型。故由
《态射》引理 \ref{morphisms-lemma-finite-type-nagata}，
$V$ 是 Nagata 概形。于是按定义 $X$ 是 Nagata 的。
\end{proof}

\begin{lemma}
\label{spaces-morphisms-lemma-ubiquity-nagata}
下列类型的代数空间是 Nagata 的。
\begin{enumerate}
\item Nagata 概形上任意局部有限型的代数空间。
\item 域上任意局部有限型的代数空间。
\item Noether 完备局部环上任意局部有限型的代数空间。
\item $\mathbf{Z}$ 上任意局部有限型的代数空间。
\item 特征零 Dedekind 环上任意局部有限型的代数空间。
\item 等等。
\end{enumerate}
\end{lemma}

\begin{proof}
第一条由引理 \ref{spaces-morphisms-lemma-finite-type-nagata} 成立。
因此其余各条也成立，见《态射》引理
\ref{morphisms-lemma-ubiquity-nagata}。
\end{proof}








\section{拟有限态射}
\label{spaces-morphisms-section-quasi-finite}

\noindent
概形态射的“局部拟有限”性质在源与靶上都是 \'etale 局部的，见
《下降》评注 \ref{descent-remark-list-local-source-target}。
此性质也在基变换下稳定，并且在靶上是 fpqc 局部的，见
《态射》引理 \ref{morphisms-lemma-base-change-quasi-finite} 以及
《下降》引理 \ref{descent-lemma-descending-property-quasi-finite}。
因此，由上面的引理 \ref{spaces-morphisms-lemma-local-source-target}，
我们可以如下定义代数空间的态射局部拟有限的含义；当该态射可表时，
这与第 \ref{spaces-morphisms-section-representable} 节中已有的概念一致。

\begin{definition}
\label{spaces-morphisms-definition-locally-quasi-finite}
设 $S$ 为概形。
设 $f : X \to Y$ 为 $S$ 上代数空间的态射。
\begin{enumerate}
\item 若引理 \ref{spaces-morphisms-lemma-local-source-target} 中取
$\mathcal{P} = \text{局部拟有限}$.
时其中的等价条件成立，则称 $f$ {\it 局部拟有限}。
\item 设 $x \in |X|$。若存在 $x$ 的开邻域 $X' \subset X$，使得
$f|_{X'} : X' \to Y$ 局部拟有限，则称 $f$ {\it 在 $x$ 处拟有限}。
\item 若代数空间的态射 $f : X \to Y$ 局部拟有限且拟紧，
则称其{\it 拟有限}。
\end{enumerate}
\end{definition}

\noindent
由《态射》引理
\ref{morphisms-lemma-quasi-finite-locally-quasi-compact}，
最后一条与概形态射的拟有限概念相容。

\begin{lemma}
\label{spaces-morphisms-lemma-base-change-quasi-finite-locus}
设 $S$ 为概形。设 $f : X \to Y$ 与 $g : Y' \to Y$ 为
$S$ 上代数空间的态射。将 $f$ 沿 $g$ 的基变换记为
$f' : X' \to Y'$，并将投影记为 $g' : X' \to X$。
假设 $f$ 局部有限型。分别以 $W \subset |X|$ 与 $W' \subset |X'|$
表示 $f$ 与 $f'$ 拟有限的点集。
\begin{enumerate}
\item $W \subset |X|$ 与 $W' \subset |X'|$ 都是开的；
\item $W' = (g')^{-1}(W)$，即形成 $f$ 拟有限的点集与基变换可交换；
\item 局部拟有限态射的基变换局部拟有限；
\item 拟有限态射的基变换拟有限。
\end{enumerate}
\end{lemma}

\begin{proof}
取概形 $V$ 及满 \'etale 态射 $V \to Y$。
取概形 $U$ 及满 \'etale 态射 $U \to V \times_Y X$。
取概形 $V'$ 及满 \'etale 态射 $V' \to Y' \times_Y V$。
令 $U' = V' \times_V U$，于是 $U' \to X'$ 也是满 \'etale 态射。图示为
$$
\vcenter{
\xymatrix{
U' \ar[d] \ar[r] & U \ar[d] \\
V' \ar[r] & V
}
}
\quad\text{位于下图之上}\quad
\vcenter{
\xymatrix{
X' \ar[d] \ar[r] & X \ar[d] \\
Y' \ar[r] & Y
}
}
$$
取 $u \in |U|$，其像为 $x \in |X|$。
“局部拟有限”性质在源与靶上都是 \'etale 局部的，见
《下降》评注 \ref{descent-remark-list-local-source-target}。
因此可应用引理 \ref{spaces-morphisms-lemma-local-source-target-at-point} 与
\ref{spaces-morphisms-lemma-local-source-target-global-implies-local}，从而
$f : X \to Y$ 在 $x$ 处拟有限，当且仅当 $U \to V$ 在 $u$ 处拟有限。
对 $f' : X' \to Y'$ 及态射 $U' \to V'$ 也同样成立。
故 (1)、(2)、(3) 归结为《态射》引理
\ref{morphisms-lemma-base-change-quasi-finite} 与
\ref{morphisms-lemma-quasi-finite-points-open}。
(4) 由 (3) 和引理 \ref{spaces-morphisms-lemma-base-change-quasi-compact} 得出。
\end{proof}

\begin{lemma}
\label{spaces-morphisms-lemma-composition-quasi-finite}
拟有限态射的复合拟有限。
对于局部拟有限态射也有同样结论。
\end{lemma}

\begin{proof}
见评注 \ref{spaces-morphisms-remark-composition-P} 和
《态射》引理 \ref{morphisms-lemma-composition-quasi-finite}。
\end{proof}

\begin{lemma}
\label{spaces-morphisms-lemma-base-change-quasi-finite}
拟有限态射的基变换拟有限。
对于局部拟有限态射也有同样结论。
\end{lemma}

\begin{proof}
这是引理 \ref{spaces-morphisms-lemma-base-change-quasi-finite-locus} 的直接推论。
\end{proof}

\noindent
下述引理把局部拟有限态射刻画为局部有限型且具有“离散纤维”的态射。
不过，这不同于要求 $|X| \to |Y|$ 具有离散纤维，参见
《例》第 \ref{examples-section-specializations-fibre-etale} 节中的讨论。

\begin{lemma}
\label{spaces-morphisms-lemma-locally-quasi-finite}
设 $S$ 为概形。设 $f : X \to Y$ 为代数空间的态射。
假设 $f$ 局部有限型。下列条件等价：
\begin{enumerate}
\item $f$ 局部拟有限；
\item 对每个态射 $\Spec(k) \to Y$，其中 $k$ 为域，空间 $|X_k|$
都是离散的。这里 $X_k = \Spec(k) \times_Y X$。
\end{enumerate}
\end{lemma}

\begin{proof}
假设 $f$ 局部拟有限。令 $\Spec(k) \to Y$ 如 (2) 所述。
取满 \'etale 态射 $U \to X$，其中 $U$ 为概形。由《空间的性质》引理
\ref{spaces-properties-lemma-base-change-etale}，
$U_k = \Spec(k) \times_Y U \to X_k$ 是代数空间的 \'etale 态射。
由引理 \ref{spaces-morphisms-lemma-base-change-quasi-finite}，
$X_k \to \Spec(k)$ 局部拟有限。按定义，这意味着
$U_k \to \Spec(k)$ 局部拟有限。故由《态射》引理
\ref{morphisms-lemma-locally-quasi-finite-fibres}，$|U_k|$ 是离散的。
由于 $|U_k| \to |X_k|$ 满且开，遂得 $|X_k|$ 离散。

\medskip\noindent
反之，假设 (2)。取满 \'etale 态射 $V \to Y$，其中 $V$ 为概形。
再取满 \'etale 态射 $U \to V \times_Y X$，其中 $U$ 为概形。
注意，由于 $f$ 局部有限型，$U \to V$ 也局部有限型。图示为
$$
\xymatrix{
U \ar[r] \ar[rd] & X \times_Y V \ar[d] \ar[r] & V \ar[d] \\
& X \ar[r] & Y
}
$$
若 $f$ 不局部拟有限，则 $U \to V$ 不局部拟有限。因此存在专门化
$u \leadsto u'$，其中某些 $u, u' \in U$ 位于同一点 $v \in V$ 上方，见
《态射》引理 \ref{morphisms-lemma-quasi-finite-at-point-characterize}。
我们断言 $u, u'$ 在 $X_v = \Spec(\kappa(v)) \times_Y X$ 中的像不同；
这将按所需方式与 $|X_v|$ 离散的假设矛盾。令
$d = \text{trdeg}_{\kappa(v)}(\kappa(u))$ 以及
$d' = \text{trdeg}_{\kappa(v)}(\kappa(u'))$。
由《态射》引理 \ref{morphisms-lemma-dimension-fibre-specialization}，
可见 $d > d'$。注意 $U_v$（即 $U \to V$ 在 $v$ 上的纤维）是
$U$ 与 $X_v$ 在 $X \times_Y V$ 上的纤维积，故 $U_v \to X_v$ 是
\'etale 的（它是 \'etale 态射 $U \to X \times_Y V$ 的基变换）。
若 $u, u' \in U_v$ 映到 $|X_v|$ 的同一元素，则由《空间的性质》引理
\ref{spaces-properties-lemma-points-cartesian}，存在一点
$r \in R_v = U_v \times_{X_v} U_v$，满足 $t(r) = u$ 且 $s(r) = u'$。
注意 $s, t : R_v \to U_v$ 是 $\kappa(v)$ 上概形的 \'etale 态射，故
$\kappa(u) \subset \kappa(r) \supset \kappa(u')$ 是 $\kappa(v)$ 上域的
有限可分扩张（见《态射》引理 \ref{morphisms-lemma-etale-over-field}）。
因此这些超越次数相等。这一矛盾完成证明。
\end{proof}

\begin{lemma}
\label{spaces-morphisms-lemma-quasi-finite-local}
设 $S$ 为概形。
设 $f : X \to Y$ 为 $S$ 上代数空间的态射。
下列条件等价：
\begin{enumerate}
\item $f$ 局部拟有限；
\item 对每个 $x \in |X|$，态射 $f$ 在 $x$ 处拟有限；
\item 对每个概形 $Z$ 及任意态射 $Z \to Y$，态射
$Z \times_Y X \to Z$ 局部拟有限；
\item 对每个仿射概形 $Z$ 及任意态射 $Z \to Y$，态射
$Z \times_Y X \to Z$ 局部拟有限；
\item 存在概形 $V$ 及满 \'etale 态射 $V \to Y$，使得
$V \times_Y X \to V$ 局部拟有限；
\item 存在概形 $U$ 及满 \'etale 态射 $\varphi : U \to X$，使得复合
$f \circ \varphi$ 局部拟有限；
\item 对每个交换图
$$
\xymatrix{
U \ar[d] \ar[r] & V \ar[d] \\
X \ar[r] & Y
}
$$
其中 $U$、$V$ 为概形，竖直箭头为 \'etale 态射，顶部水平箭头局部拟有限；
\item 存在交换图
$$
\xymatrix{
U \ar[d] \ar[r] & V \ar[d] \\
X \ar[r] & Y
}
$$
其中 $U$、$V$ 为概形，竖直箭头为 \'etale 态射，且 $U \to X$ 满，
使得顶部水平箭头局部拟有限；
\item 存在 Zariski 覆盖 $Y = \bigcup_{i \in I} Y_i$ 以及
$f^{-1}(Y_i) = \bigcup X_{ij}$，使得每个态射
$X_{ij} \to Y_i$ 局部拟有限。
\end{enumerate}
\end{lemma}

\begin{proof}
略。
\end{proof}

\begin{lemma}
\label{spaces-morphisms-lemma-immersion-quasi-finite}
浸入局部拟有限。
\end{lemma}

\begin{proof}
略。
\end{proof}

\begin{lemma}
\label{spaces-morphisms-lemma-permanence-quasi-finite}
设 $S$ 为概形。
设 $X \to Y \to Z$ 为 $S$ 上代数空间的态射。
若 $X \to Z$ 局部拟有限，则 $X \to Y$ 局部拟有限。
\end{lemma}

\begin{proof}
取交换图
$$
\xymatrix{
U \ar[d] \ar[r] & V \ar[d] \ar[r] & W \ar[d] \\
X \ar[r] & Y \ar[r] & Z
}
$$
使竖直箭头均为满 \'etale 态射。（见《空间》引理
\ref{spaces-lemma-lift-morphism-presentations}。）
对顶行应用《态射》引理
\ref{morphisms-lemma-permanence-quasi-finite}。
\end{proof}

\begin{lemma}
\label{spaces-morphisms-lemma-quasi-finite-at-a-finite-number-of-points}
设 $S$ 为概形。设 $f : X \to Y$ 为 $S$ 上代数空间的有限型态射。
设 $y \in |Y|$。在 $y$ 上方且使 $f$ 拟有限的 $|X|$ 中点至多有限个。
\end{lemma}

\begin{proof}
取域 $k$ 及由 $y$ 所确定的等价类中的态射 $\Spec(k) \to Y$。
纤维 $X_k = \Spec(k) \times_Y X$ 是域上的有限型代数空间，尤其拟紧。
映射 $|X_k| \to |X|$ 满射到 $|X| \to |Y|$ 在 $y$ 上的纤维
（《空间的性质》引理 \ref{spaces-properties-lemma-points-cartesian}）。
而且，由引理 \ref{spaces-morphisms-lemma-base-change-quasi-finite-locus}，
$X_k \to \Spec(k)$ 拟有限的点集映满到 $y$ 上方使 $f$ 拟有限的点集。
取仿射概形 $U$ 及满 \'etale 态射 $U \to X_k$
（《空间的性质》引理
\ref{spaces-properties-lemma-quasi-compact-affine-cover}）。
于是 $U \to \Spec(k)$ 是有限型态射，且使该态射拟有限的点至多有限个，
见《态射》引理
\ref{morphisms-lemma-quasi-finite-at-a-finite-number-of-points}。
由于 $X_k \to \Spec(k)$ 在一点 $x'$ 处拟有限，当且仅当该点是 $U$
中某个使 $U \to \Spec(k)$ 拟有限之点的像，结论随即成立。
\end{proof}

\begin{lemma}
\label{spaces-morphisms-lemma-monomorphism-loc-finite-type-loc-quasi-finite}
设 $S$ 为概形。
设 $f : X \to Y$ 为 $S$ 上代数空间的态射。
若 $f$ 局部有限型且为单态射，则 $f$ 分离并且局部拟有限。
\end{lemma}

\begin{proof}
单态射是分离的，见引理 \ref{spaces-morphisms-lemma-monomorphism-separated}。
由引理 \ref{spaces-morphisms-lemma-quasi-finite-local}，只需在沿 $Z \to Y$ 作基变换后
证明此引理，其中 $Z$ 仿射。因此可以假设 $Y$ 是仿射概形。
取仿射概形 $U$ 及 \'etale 态射 $U \to X$。由于 $X \to Y$ 局部有限型，
仿射概形的态射 $U \to Y$ 有限型。由于 $X \to Y$ 是单态射，有
$U \times_X U = U \times_Y U$。特别地，映射 $U \times_Y U \to U$
是 \'etale 的。设 $y \in Y$。则或者 $U_y$ 为空，或者对位于 $y$
上方的某个 $u \in U$，
$\Spec(\kappa(u)) \times_{\Spec(\kappa(y))} U_y$
同构于 $U \times_Y U \to U$ 在 $u$ 上的纤维。这说明 $U \to Y$ 的
纤维是有限离散集（因为 $U \times_Y U \to U$ 是仿射概形的 \'etale 态射，
见《态射》引理 \ref{morphisms-lemma-etale-over-field}）。
因此 $U \to Y$ 拟有限，见《态射》引理
\ref{morphisms-lemma-quasi-finite-at-point-characterize}。
由于 $U \to X$ 是任意一个以仿射概形 $U$ 为源的 \'etale 态射，
这说明 $X \to Y$ 局部拟有限。
\end{proof}













\section{有限表示态射}
\label{spaces-morphisms-section-finite-presentation}

\noindent
概形态射的“局部有限表示”性质在源与靶上都是 \'etale 局部的，见
《下降》评注 \ref{descent-remark-list-local-source-target}。
此性质也在基变换下稳定，并且在靶上是 fpqc 局部的，见
《态射》引理 \ref{morphisms-lemma-base-change-finite-presentation} 以及
《下降》引理
\ref{descent-lemma-descending-property-locally-finite-presentation}。
因此，由上面的引理 \ref{spaces-morphisms-lemma-local-source-target}，
我们可以如下定义代数空间的态射局部有限表示的含义；当该态射可表时，
这与第 \ref{spaces-morphisms-section-representable} 节中已有的概念一致。

\begin{definition}
\label{spaces-morphisms-definition-locally-finite-presentation}
设 $S$ 为概形。
设 $f : X \to Y$ 为 $S$ 上代数空间的态射。
\begin{enumerate}
\item 若引理 \ref{spaces-morphisms-lemma-local-source-target} 中取
$\mathcal{P} =$“局部有限表示”时其中的等价条件成立，则称 $f$
{\it 局部有限表示}。
\item 设 $x \in |X|$。若存在 $x$ 的开邻域 $X' \subset X$，使得
$f|_{X'} : X' \to Y$ 局部有限表示，则称 $f$ {\it 在 $x$ 处有限表示}
\footnote{使用“在 $x$ 处局部有限表示”似乎很别扭；不过，现用术语也可能
引起误解，因为“在 $x$ 处有限表示”{\bf 并不}意味着存在开邻域
$X' \subset X$，使得 $f|_{X'}$ 有限表示。}。
\item 若代数空间的态射 $f : X \to Y$ 局部有限表示、拟紧且拟分离，
则称其{\it 有限表示}。
\end{enumerate}
\end{definition}

\noindent
注意，有限表示态射{\bf 并不}只是局部有限表示的拟紧态射。

\begin{lemma}
\label{spaces-morphisms-lemma-composition-finite-presentation}
有限表示态射的复合有限表示。
对于局部有限表示态射也有同样结论。
\end{lemma}

\begin{proof}
见评注 \ref{spaces-morphisms-remark-composition-P} 和《态射》引理
\ref{morphisms-lemma-composition-finite-presentation}。
还要使用拟紧态射与拟分离态射的相应结论
（引理 \ref{spaces-morphisms-lemma-composition-quasi-compact} 与
\ref{spaces-morphisms-lemma-composition-separated}）。
\end{proof}

\begin{lemma}
\label{spaces-morphisms-lemma-base-change-finite-presentation}
有限表示态射的基变换有限表示。
对于局部有限表示态射也有同样结论。
\end{lemma}

\begin{proof}
见评注 \ref{spaces-morphisms-remark-base-change-P} 和《态射》引理
\ref{morphisms-lemma-base-change-finite-presentation}。
还要使用拟紧态射与拟分离态射的相应结论
（引理 \ref{spaces-morphisms-lemma-base-change-quasi-compact} 与
\ref{spaces-morphisms-lemma-base-change-separated}）。
\end{proof}

\begin{lemma}
\label{spaces-morphisms-lemma-finite-presentation-local}
设 $S$ 为概形。
设 $f : X \to Y$ 为 $S$ 上代数空间的态射。
下列条件等价：
\begin{enumerate}
\item $f$ 局部有限表示；
\item 对每个 $x \in |X|$，态射 $f$ 在 $x$ 处有限表示；
\item 对每个概形 $Z$ 及任意态射 $Z \to Y$，态射
$Z \times_Y X \to Z$ 局部有限表示；
\item 对每个仿射概形 $Z$ 及任意态射 $Z \to Y$，态射
$Z \times_Y X \to Z$ 局部有限表示；
\item 存在概形 $V$ 及满 \'etale 态射 $V \to Y$，使得
$V \times_Y X \to V$ 局部有限表示；
\item 存在概形 $U$ 及满 \'etale 态射 $\varphi : U \to X$，使得复合
$f \circ \varphi$ 局部有限表示；
\item 对每个交换图
$$
\xymatrix{
U \ar[d] \ar[r] & V \ar[d] \\
X \ar[r] & Y
}
$$
其中 $U$、$V$ 为概形，竖直箭头为 \'etale 态射，顶部水平箭头局部有限表示；
\item 存在交换图
$$
\xymatrix{
U \ar[d] \ar[r] & V \ar[d] \\
X \ar[r] & Y
}
$$
其中 $U$、$V$ 为概形，竖直箭头为 \'etale 态射，且 $U \to X$ 满，
使得顶部水平箭头局部有限表示；
\item 存在 Zariski 覆盖 $Y = \bigcup_{i \in I} Y_i$ 以及
$f^{-1}(Y_i) = \bigcup X_{ij}$，使得每个态射
$X_{ij} \to Y_i$ 局部有限表示。
\end{enumerate}
\end{lemma}

\begin{proof}
略。
\end{proof}

\begin{lemma}
\label{spaces-morphisms-lemma-finite-presentation-finite-type}
局部有限表示的态射局部有限型。
有限表示态射有限型。
\end{lemma}

\begin{proof}
设 $f : X \to Y$ 为局部有限表示的代数空间态射。这意味着存在
引理 \ref{spaces-morphisms-lemma-local-source-target} 中的图，其中 $h$ 局部有限表示，
且竖直箭头 $a$ 满。由《态射》引理
\ref{morphisms-lemma-finite-presentation-finite-type}，$h$ 局部有限型。
故按定义 $X \to Y$ 局部有限型。若 $f$ 有限表示，则它拟紧，
从而 $f$ 有限型。
\end{proof}

\begin{lemma}
\label{spaces-morphisms-lemma-finite-presentation-noetherian}
设 $S$ 为概形。
设 $f : X \to Y$ 为 $S$ 上代数空间的态射。
若 $f$ 有限表示且 $Y$ 是 Noether 的，则 $X$ 是 Noether 的。
\end{lemma}

\begin{proof}
假设 $f$ 有限表示且 $Y$ 是 Noether 的。由引理
\ref{spaces-morphisms-lemma-finite-presentation-finite-type} 与
\ref{spaces-morphisms-lemma-locally-finite-type-locally-noetherian}，可见 $X$ 局部 Noether。
由于 $f$ 拟紧且 $Y$ 拟紧，可见 $X$ 拟紧。
由于 $f$ 有限表示，它拟分离（见定义
\ref{spaces-morphisms-definition-locally-finite-presentation}）；又由于 $Y$ 是 Noether 的，
它拟分离（见《空间的性质》定义
\ref{spaces-properties-definition-noetherian}）。
故由引理 \ref{spaces-morphisms-lemma-separated-over-separated}，$X$ 拟分离。
于是我们已验证《空间的性质》定义
\ref{spaces-properties-definition-noetherian} 中的全部三个条件，证明完成。
\end{proof}

\begin{lemma}
\label{spaces-morphisms-lemma-noetherian-finite-type-finite-presentation}
设 $S$ 为概形。
设 $f : X \to Y$ 为 $S$ 上代数空间的态射。
\begin{enumerate}
\item 若 $Y$ 局部 Noether 且 $f$ 局部有限型，则 $f$ 局部有限表示。
\item 若 $Y$ 局部 Noether，且 $f$ 有限型并拟分离，则 $f$ 有限表示。
\end{enumerate}
\end{lemma}

\begin{proof}
假设 $f : X \to Y$ 局部有限型且 $Y$ 局部 Noether。
这意味着存在引理 \ref{spaces-morphisms-lemma-local-source-target} 中的图，
其中 $h$ 局部有限型，且竖直箭头 $a$ 满。由《态射》引理
\ref{morphisms-lemma-noetherian-finite-type-finite-presentation}，
$h$ 局部有限表示。故按定义 $X \to Y$ 局部有限表示。这证明 (1)。
若 $f$ 有限型且拟分离，则它也拟紧且拟分离，(2) 立即得出。
\end{proof}

\begin{lemma}
\label{spaces-morphisms-lemma-finite-presentation-quasi-compact-quasi-separated}
设 $S$ 为概形。设 $Y$ 为 $S$ 上拟紧且拟分离的代数空间。
若 $X$ 在 $Y$ 上有限表示，则 $X$ 拟紧且拟分离。
\end{lemma}

\begin{proof}
略。
\end{proof}

\begin{lemma}
\label{spaces-morphisms-lemma-finite-presentation-permanence}
设 $S$ 为概形。
设 $f : X \to Y$ 与 $Y \to Z$ 为 $S$ 上代数空间的态射。
若 $X$ 在 $Z$ 上局部有限表示，且 $Y$ 在 $Z$ 上局部有限型，
则 $f$ 局部有限表示。
\end{lemma}

\begin{proof}
取概形 $W$ 及满 \'etale 态射 $W \to Z$。再取概形 $V$ 及满 \'etale 态射
$V \to W \times_Z Y$。最后取概形 $U$ 及满 \'etale 态射
$U \to V \times_Y X$。按定义，$U$ 在 $W$ 上局部有限表示，而 $V$
在 $W$ 上局部有限型。由《态射》引理
\ref{morphisms-lemma-finite-presentation-permanence}，态射 $U \to V$
局部有限表示。因此 $f$ 局部有限表示。
\end{proof}

\begin{lemma}
\label{spaces-morphisms-lemma-diagonal-morphism-finite-type}
设 $S$ 为概形。设 $f : X \to Y$ 为 $S$ 上代数空间的态射，其对角态射为
$\Delta : X \to X \times_Y X$。若 $f$ 局部有限型，则 $\Delta$ 局部有限表示。
若 $f$ 拟分离且局部有限型，则 $\Delta$ 有限表示。
\end{lemma}

\begin{proof}
注意，$\Delta$ 是 $X$ 上的态射（通过第二投影 $X \times_Y X \to X$）。
假设 $f$ 局部有限型。注意，$X$ 在 $X$ 上有限表示，而 $X \times_Y X$
在 $X$ 上有限型（由引理 \ref{spaces-morphisms-lemma-base-change-finite-type}）。
因此第一条由引理 \ref{spaces-morphisms-lemma-finite-presentation-permanence} 成立。
第二条由第一条、定义以及对角态射分离这一事实
（引理 \ref{spaces-morphisms-lemma-properties-diagonal}）得出。
\end{proof}

\begin{lemma}
\label{spaces-morphisms-lemma-open-immersion-locally-finite-presentation}
代数空间的开浸入局部有限表示。
\end{lemma}

\begin{proof}
开浸入按定义可表，因此可使用《空间》引理
\ref{spaces-lemma-representable-transformations-property-implication}
的一般原理以及《态射》引理
\ref{morphisms-lemma-open-immersion-locally-finite-presentation}。
\end{proof}

\begin{lemma}
\label{spaces-morphisms-lemma-closed-immersion-finite-presentation}
闭浸入 $i : Z \to X$ 有限表示，当且仅当相伴的拟凝聚理想层
$\mathcal{I} = \Ker(\mathcal{O}_X \to i_*\mathcal{O}_Z)$
有限型（作为 $\mathcal{O}_X$-模）。
\end{lemma}

\begin{proof}
设 $U$ 为概形，并设 $U \to X$ 为满 \'etale 态射。由引理
\ref{spaces-morphisms-lemma-finite-presentation-local}，可见
$i' : Z \times_X U \to U$ 有限表示，当且仅当 $i$ 有限表示。
由《空间的性质》第
\ref{spaces-properties-section-properties-modules} 节，可见 $\mathcal{I}$
有限型，当且仅当
$\mathcal{I}|_U = \Ker(\mathcal{O}_U \to i'_*\mathcal{O}_{Z \times_X U})$
有限型。因此结论归结为概形情形，见《态射》引理
\ref{morphisms-lemma-closed-immersion-finite-presentation}。
\end{proof}







\section{可构造集}
\label{spaces-morphisms-section-constructible}

\noindent
本节接续《空间的性质》第
\ref{spaces-properties-section-constructible} 节。

\begin{lemma}
\label{spaces-morphisms-lemma-inverse-image-constructible}
设 $S$ 为概形。
设 $f : X \to Y$ 为 $S$ 上代数空间的态射。
设 $E \subset |Y|$ 为子集。若 $E$ 在 $Y$ 中 \'etale 局部可构造，
则 $f^{-1}(E)$ 在 $X$ 中 \'etale 局部可构造。
\end{lemma}

\begin{proof}
取概形 $V$ 及满 \'etale 态射 $\varphi : V \to Y$。
取概形 $U$ 及满 \'etale 态射 $U \to V \times_Y X$。
于是 $U \to X$ 满且 \'etale，而 $f^{-1}(E)$ 在 $U$ 中的原像是
$\varphi^{-1}(E)$ 沿 $U \to V$ 的原像。因此引理归结为态射 $U \to V$
的概形情形（《态射》引理
\ref{morphisms-lemma-inverse-image-constructible}）以及定义
（《空间的性质》定义
\ref{spaces-properties-definition-locally-constructible}）。
\end{proof}

\begin{theorem}[Chevalley 定理]
\label{spaces-morphisms-theorem-chevalley}
设 $S$ 为概形。
设 $f : X \to Y$ 为 $S$ 上代数空间的态射。
假设 $f$ 拟紧且局部有限表示。则 $|X|$ 的每个 \'etale 局部可构造子集的像
都是 $|Y|$ 的 \'etale 局部可构造子集。
\end{theorem}

\begin{proof}
设 $E \subset |X|$ 为 \'etale 局部可构造子集。
设 $V \to Y$ 为 \'etale 态射，且 $V$ 仿射。只需证明 $f(E)$ 在 $V$
中的原像可构造，见《空间的性质》定义
\ref{spaces-properties-definition-locally-constructible}。
由于 $f$ 拟紧，$V \times_Y X$ 是拟紧代数空间。取仿射概形 $U$ 及满
\'etale 态射 $U \to V \times_Y X$（《空间的性质》引理
\ref{spaces-properties-lemma-quasi-compact-affine-cover}）。
由《空间的性质》引理 \ref{spaces-properties-lemma-points-cartesian}，
$f(E)$ 在 $V$ 中的原像，是 $E$ 在 $U$ 中的原像沿 $U \to V$ 的像。
因此结论归结为概形情形，见《态射》引理
\ref{morphisms-lemma-chevalley}。
\end{proof}







\section{平坦态射}
\label{spaces-morphisms-section-flat}

\noindent
概形态射的“平坦”性质在源与靶上都是 \'etale 局部的，见
《下降》评注 \ref{descent-remark-list-local-source-target}。
此性质也在基变换下稳定，并且在靶上是 fpqc 局部的，见
《态射》引理 \ref{morphisms-lemma-base-change-flat} 与
《下降》引理 \ref{descent-lemma-descending-property-flat}。
因此，由上面的引理 \ref{spaces-morphisms-lemma-local-source-target}，
我们可以如下定义代数空间的平坦态射；当该态射可表时，
这与第 \ref{spaces-morphisms-section-representable} 节中已有的概念一致。

\begin{definition}
\label{spaces-morphisms-definition-flat}
设 $S$ 为概形。
设 $f : X \to Y$ 为 $S$ 上代数空间的态射。
\begin{enumerate}
\item 若引理 \ref{spaces-morphisms-lemma-local-source-target} 中取
$\mathcal{P} =$“平坦”时其中的等价条件成立，则称 $f$ {\it 平坦}。
\item 设 $x \in |X|$。若引理 \ref{spaces-morphisms-lemma-local-source-target-at-point}
中取 $\mathcal{Q} =$“所诱导的局部环映射平坦”时其中的等价条件成立，
则称 $f$ {\it 在 $x$ 处平坦}。
\end{enumerate}
注意，由《下降》引理 \ref{descent-lemma-flat-at-point}，第二条定义良好。
\end{definition}

\noindent
先作一个简单核对。

\begin{lemma}
\label{spaces-morphisms-lemma-flat-is-flat-at-all-points}
设 $S$ 为概形。设 $f : X \to Y$ 为 $S$ 上代数空间的态射。
则 $f$ 平坦，当且仅当 $f$ 在 $|X|$ 的每一点处平坦。
\end{lemma}

\begin{proof}
取交换图
$$
\xymatrix{
U \ar[d]_a \ar[r]_h & V \ar[d]^b \\
X \ar[r]^f & Y
}
$$
其中 $U$ 与 $V$ 为概形，竖直箭头为 \'etale 态射，且 $a$ 满。
按定义，$f$ 平坦，当且仅当 $h$ 平坦（定义 \ref{spaces-morphisms-definition-P}）。
按定义，$f$ 在 $x \in |X|$ 处平坦，当且仅当 $h$ 在映到 $x$ 的某个
（等价地，任意一个）$u \in U$ 处平坦（定义
\ref{spaces-morphisms-definition-P-at-point}）。因此本引理由如下事实得出：
概形态射平坦，当且仅当它在源的每一点处平坦
（《态射》定义 \ref{morphisms-definition-flat}）。
\end{proof}

\begin{lemma}
\label{spaces-morphisms-lemma-composition-flat}
平坦态射的复合平坦。
\end{lemma}

\begin{proof}
见评注 \ref{spaces-morphisms-remark-composition-P} 和
《态射》引理 \ref{morphisms-lemma-composition-flat}。
\end{proof}

\begin{lemma}
\label{spaces-morphisms-lemma-base-change-flat}
平坦态射的基变换平坦。
\end{lemma}

\begin{proof}
见评注 \ref{spaces-morphisms-remark-base-change-P} 和
《态射》引理 \ref{morphisms-lemma-base-change-flat}。
\end{proof}

\begin{lemma}
\label{spaces-morphisms-lemma-flat-local}
设 $S$ 为概形。
设 $f : X \to Y$ 为 $S$ 上代数空间的态射。
下列条件等价：
\begin{enumerate}
\item $f$ 平坦；
\item 对每个 $x \in |X|$，态射 $f$ 在 $x$ 处平坦；
\item 对每个概形 $Z$ 及任意态射 $Z \to Y$，态射
$Z \times_Y X \to Z$ 平坦；
\item 对每个仿射概形 $Z$ 及任意态射 $Z \to Y$，态射
$Z \times_Y X \to Z$ 平坦；
\item 存在概形 $V$ 及满 \'etale 态射 $V \to Y$，使得
$V \times_Y X \to V$ 平坦；
\item 存在概形 $U$ 及满 \'etale 态射 $\varphi : U \to X$，使得复合
$f \circ \varphi$ 平坦；
\item 对每个交换图
$$
\xymatrix{
U \ar[d] \ar[r] & V \ar[d] \\
X \ar[r] & Y
}
$$
其中 $U$、$V$ 为概形，竖直箭头为 \'etale 态射，顶部水平箭头平坦；
\item 存在交换图
$$
\xymatrix{
U \ar[d] \ar[r] & V \ar[d] \\
X \ar[r] & Y
}
$$
其中 $U$、$V$ 为概形，竖直箭头为 \'etale 态射，且 $U \to X$ 满，
使得顶部水平箭头平坦；
\item 存在 Zariski 覆盖 $Y = \bigcup Y_i$ 以及
$f^{-1}(Y_i) = \bigcup X_{ij}$，使得每个态射 $X_{ij} \to Y_i$ 平坦。
\end{enumerate}
\end{lemma}

\begin{proof}
略。
\end{proof}

\begin{lemma}
\label{spaces-morphisms-lemma-fppf-open}
局部有限表示的平坦态射是普遍开的。
\end{lemma}

\begin{proof}
设 $f : X \to Y$ 为 $S$ 上代数空间的局部有限表示平坦态射。取图
$$
\xymatrix{
U \ar[r]_\alpha \ar[d] & V \ar[d] \\
X \ar[r] & Y
}
$$
其中 $U$ 与 $V$ 为概形，竖直箭头满且 \'etale，见《空间》引理
\ref{spaces-lemma-lift-morphism-presentations}。
由引理 \ref{spaces-morphisms-lemma-flat-local} 与 \ref{spaces-morphisms-lemma-finite-presentation-local}，
态射 $\alpha$ 平坦且局部有限表示。因此由《态射》引理
\ref{morphisms-lemma-fppf-open}，可见 $\alpha$ 普遍开。
故由引理 \ref{spaces-morphisms-lemma-universally-open-local}，$X \to Y$ 普遍开。
\end{proof}

\begin{lemma}
\label{spaces-morphisms-lemma-fpqc-quotient-topology}
设 $S$ 为概形。
设 $f : X \to Y$ 为 $S$ 上代数空间的平坦、拟紧满态射。
子集 $T \subset |Y|$ 是开的（相应地，闭的），当且仅当
$f^{-1}(|T|)$ 在 $|X|$ 中是开的（相应地，闭的）。
换言之，$f$ 是商映射；事实上，它是普遍商映射。
\end{lemma}

\begin{proof}
取仿射概形 $V_i$ 及 \'etale 态射 $V_i \to Y$，使得
$V = \coprod V_i \to Y$ 满，见《空间的性质》引理
\ref{spaces-properties-lemma-cover-by-union-affines}。
对每个 $i$，代数空间 $V_i \times_Y X$ 拟紧。因此可取仿射概形 $U_i$
及满 \'etale 态射 $U_i \to V_i \times_Y X$，见《空间的性质》引理
\ref{spaces-properties-lemma-quasi-compact-affine-cover}。
于是复合 $U_i \to V_i \times_Y X \to V_i$ 是仿射概形之间的满平坦态射。
当然，$U = \coprod U_i \to X$ 满且 \'etale，并且
$U \to V \times_Y X$ 满。此外，态射 $U \to V$ 是各态射
$U_i \to V_i$ 的不交并。故 $U \to V$ 满、拟紧且平坦。考虑图
$$
\xymatrix{
U \ar[r] \ar[d] & X \ar[d] \\
V \ar[r] & Y
}
$$
按 $|Y|$ 上拓扑的定义，集合 $T$ 是闭的（相应地，开的），当且仅当
$g^{-1}(T) \subset |V|$ 是闭的（相应地，开的）。对 $f^{-1}(T)$ 及其在
$|U|$ 中的原像也同样成立。由于 $U \to V$ 拟紧、满且平坦，
结论由《态射》引理 \ref{morphisms-lemma-fpqc-quotient-topology} 得出。
\end{proof}

\begin{lemma}
\label{spaces-morphisms-lemma-flat-at-point-etale-local-rings}
设 $S$ 为概形。
设 $f : X \to Y$ 为 $S$ 上代数空间的态射。
设 $\overline{x}$ 为 $X$ 的一个几何点，位于点 $x \in |X|$ 上方。
令 $\overline{y} = f \circ \overline{x}$。下列条件等价：
\begin{enumerate}
\item $f$ 在 $x$ 处平坦；
\item \'etale 局部环之间的映射
$\mathcal{O}_{Y, \overline{y}} \to \mathcal{O}_{X, \overline{x}}$
平坦。
\end{enumerate}
\end{lemma}

\begin{proof}
取交换图
$$
\xymatrix{
U \ar[d]_a \ar[r]_h & V \ar[d]^b \\
X \ar[r]^f & Y
}
$$
其中 $U$ 与 $V$ 为概形，$a, b$ 为 \'etale 态射，且 $u \in U$ 映到 $x$。
可取位于 $u$ 上方的几何点 $\overline{u} : \Spec(k) \to U$，使得
$\overline{x} = a \circ \overline{u}$，见《空间的性质》引理
\ref{spaces-properties-lemma-geometric-lift-to-usual}。
令 $\overline{v} = h \circ \overline{u}$，其像为 $v \in V$。我们知道
$$
\mathcal{O}_{X, \overline{x}} = \mathcal{O}_{U, u}^{sh}
\quad\text{且}\quad
\mathcal{O}_{Y, \overline{y}} = \mathcal{O}_{V, v}^{sh}
$$
见《空间的性质》引理
\ref{spaces-properties-lemma-describe-etale-local-ring}。于是得到局部环的交换图
$$
\xymatrix{
\mathcal{O}_{U, u} \ar[r] &
\mathcal{O}_{X, \overline{x}} \\
\mathcal{O}_{V, v} \ar[u] \ar[r] &
\mathcal{O}_{Y, \overline{y}} \ar[u]
}
$$
其中水平箭头平坦。我们要证明左侧竖直箭头平坦，当且仅当右侧竖直箭头平坦。
《代数》引理 \ref{algebra-lemma-flatness-descends-more-general} 告诉我们，
$\mathcal{O}_{U, u}$ 在 $\mathcal{O}_{V, v}$ 上平坦，当且仅当
$\mathcal{O}_{X, \overline{x}}$ 在 $\mathcal{O}_{V, v}$ 上平坦。
因此结论由《平坦性的进一步内容》引理
\ref{flat-lemma-flat-up-down-henselization} 得出。
\end{proof}

\begin{lemma}
\label{spaces-morphisms-lemma-flat-morphism-sites}
设 $S$ 为概形。
设 $f : X \to Y$ 为 $S$ 上代数空间的态射。
则 $f$ 平坦，当且仅当与 $f$ 相伴的站点态射
$
(f_{small}, f^\sharp) :
(X_\etale, \mathcal{O}_X)
\to
(Y_\etale, \mathcal{O}_Y)
$
平坦。
\end{lemma}

\begin{proof}
$(f_{small}, f^\sharp)$ 的平坦性由 $\mathcal{O}_X$ 作为
$f_{small}^{-1}\mathcal{O}_Y$-模的平坦性定义。这可以在茎上检验，见
《站点上的模》引理 \ref{sites-modules-lemma-check-flat-stalks} 以及
《空间的性质》定理 \ref{spaces-properties-theorem-exactness-stalks}。
而我们已经看到，$f$ 的平坦性可以在茎上检验，见引理
\ref{spaces-morphisms-lemma-flat-at-point-etale-local-rings}。
\end{proof}

\begin{lemma}
\label{spaces-morphisms-lemma-flat-pullback-support}
设 $S$ 为概形。设 $f : Y \to X$ 为 $S$ 上代数空间的态射。
设 $\mathcal{F}$ 为有限型拟凝聚 $\mathcal{O}_X$-模，其概形论支集为
$Z \subset X$。若 $f$ 平坦，则 $f^{-1}(Z)$ 是 $f^*\mathcal{F}$ 的概形论支集。
\end{lemma}

\begin{proof}
利用引理 \ref{spaces-morphisms-lemma-scheme-theoretic-support} 给出的概形论支集刻画，
以及引理 \ref{spaces-morphisms-lemma-flat-local} 中用 \'etale 覆盖给出的平坦态射刻画，
可将问题归结为概形情形，即《态射》引理
\ref{morphisms-lemma-flat-pullback-support}。
\end{proof}

\begin{lemma}
\label{spaces-morphisms-lemma-flat-morphism-scheme-theoretically-dense-open}
设 $S$ 为概形。
设 $f : X \to Y$ 为 $S$ 上代数空间的平坦态射。
设 $V \to Y$ 为拟紧开浸入。若 $V$ 在 $Y$ 中概形论稠密，
则 $f^{-1}V$ 在 $X$ 中概形论稠密。
\end{lemma}

\begin{proof}
利用引理 \ref{spaces-morphisms-lemma-scheme-theoretically-dense} 中对概形论稠密开集的刻画，
以及引理 \ref{spaces-morphisms-lemma-flat-local} 中用 \'etale 覆盖给出的平坦态射刻画，
可将问题归结为概形情形，即《态射》引理
\ref{morphisms-lemma-flat-morphism-scheme-theoretically-dense-open}。
\end{proof}

\begin{lemma}
\label{spaces-morphisms-lemma-flat-base-change-scheme-theoretic-image}
设 $S$ 为概形。设 $f : X \to Y$ 为 $S$ 上代数空间的平坦态射。
设 $g : V \to Y$ 为代数空间的拟紧态射。设 $Z \subset Y$ 为 $g$
的概形论像，并设 $Z' \subset X$ 为基变换 $V \times_Y X \to X$
的概形论像。则 $Z' = f^{-1}Z$。
\end{lemma}

\begin{proof}
取满 \'etale 态射 $Y' \to Y$，使得 $Y'$ 是仿射概形的不交并
（《空间的性质》引理
\ref{spaces-properties-lemma-cover-by-union-affines}）。
取满 \'etale 态射 $X' \to X \times_Y Y'$，使得 $X'$ 是仿射概形的不交并。
由引理 \ref{spaces-morphisms-lemma-flat-local}，态射 $X' \to Y'$ 平坦。令
$V' = V \times_Y Y'$。由引理
\ref{spaces-morphisms-lemma-quasi-compact-scheme-theoretic-image}，$Z$ 在 $Y'$ 中的原像是
$V' \to Y'$ 的概形论像，而 $Z'$ 在 $X'$ 中的原像是
$V' \times_{Y'} X' \to X'$ 的概形论像。由于 $X' \to X$ 满且 \'etale，
只需对态射 $X' \to Y'$ 与 $V' \to Y'$ 证明结论。因此可以假设
$X$ 与 $Y$ 是仿射概形。此时 $V$ 是拟紧代数空间。取仿射概形 $W$
及满 \'etale 态射 $W \to V$（《空间的性质》引理
\ref{spaces-properties-lemma-quasi-compact-affine-cover}）。显然，
$V \to Y$ 的概形论像与 $W \to Y$ 的概形论像一致；对于
$V \times_Y X \to Y$ 与 $W \times_Y X \to X$ 也同样成立。
因此问题归结为概形情形，即《态射》引理
\ref{morphisms-lemma-flat-base-change-scheme-theoretic-image}。
\end{proof}





\section{平坦模}
\label{spaces-morphisms-section-flat-modules}

\noindent
本节定义模在一点处平坦的含义。为此要用到代数空间的小 \'etale 站点
$X_\etale$ 上层的茎这一概念，见《空间的性质》定义
\ref{spaces-properties-definition-stalk}。

\begin{lemma}
\label{spaces-morphisms-lemma-flat-at-point}
设 $S$ 为概形。设 $f : X \to Y$ 为 $S$ 上代数空间的态射。
设 $\mathcal{F}$ 为 $X$ 上的拟凝聚层。设 $x \in |X|$。下列条件等价：
\begin{enumerate}
\item 对某个交换图
$$
\xymatrix{
U \ar[d]_a \ar[r]_h & V \ar[d]^b \\
X \ar[r]^f & Y
}
$$
其中 $U$ 与 $V$ 为概形，$a, b$ 为 \'etale 态射，且 $u \in U$ 映到 $x$，
模 $a^*\mathcal{F}$ 在 $u$ 处相对于 $V$ 平坦；
\item 茎 $\mathcal{F}_{\overline{x}}$ 在 \'etale 局部环
$\mathcal{O}_{Y, \overline{y}}$ 上平坦，其中 $\overline{x}$ 是位于 $x$
上方的任意几何点，且 $\overline{y} = f \circ \overline{x}$。
\end{enumerate}
\end{lemma}

\begin{proof}
% P08-E239: frozen source says "geometric proof" rather than "geometric point".
在此证明中，我们固定 $x$ 上方的几何证明
$\overline{x} : \Spec(k) \to X$，并将其在 $Y$ 中的像记为
$\overline{y} = f \circ \overline{x}$。给定 (1) 中的图，可以找到位于 $u$
上方的几何点 $\overline{u} : \Spec(k) \to U$，使得
$\overline{x} = a \circ \overline{u}$，见《空间的性质》引理
\ref{spaces-properties-lemma-geometric-lift-to-usual}。
令 $\overline{v} = h \circ \overline{u}$，其像为 $v \in V$。我们知道
$$
\mathcal{O}_{X, \overline{x}} = \mathcal{O}_{U, u}^{sh}
\quad\text{且}\quad
\mathcal{O}_{Y, \overline{y}} = \mathcal{O}_{V, v}^{sh}
$$
见《空间的性质》引理
\ref{spaces-properties-lemma-describe-etale-local-ring}。于是得到交换图
$$
\xymatrix{
\mathcal{O}_{U, u} \ar[r] &
\mathcal{O}_{X, \overline{x}} \\
\mathcal{O}_{V, v} \ar[u] \ar[r] &
\mathcal{O}_{Y, \overline{y}} \ar[u]
}
$$
，其中各对象为局部环。最后，由《空间的性质》引理
\ref{spaces-properties-lemma-stalk-quasi-coherent}，有
$$
\mathcal{F}_{\overline{x}} =
(a^*\mathcal{F})_u \otimes_{\mathcal{O}_{U, u}}
\mathcal{O}_{X, \overline{x}}
$$
。因此《代数》引理 \ref{algebra-lemma-flatness-descends-more-general}
告诉我们，$(a^*\mathcal{F})_u$ 在 $\mathcal{O}_{V, v}$ 上平坦，当且仅当
$\mathcal{F}_{\overline{x}}$ 在 $\mathcal{O}_{V, v}$ 上平坦。
故结论由《平坦性的进一步内容》引理
\ref{flat-lemma-flat-up-down-henselization} 得出。
\end{proof}

\begin{definition}
\label{spaces-morphisms-definition-flat-module}
设 $S$ 为概形。
设 $f : X \to Y$ 为 $S$ 上代数空间的态射。
设 $\mathcal{F}$ 为 $X$ 上的拟凝聚层。
\begin{enumerate}
\item 设 $x \in |X|$。若引理 \ref{spaces-morphisms-lemma-flat-at-point} 中的等价条件成立，
则称 $\mathcal{F}$ {\it 在 $x$ 处相对于 $Y$ 平坦}。
\item 若 $\mathcal{F}$ 在每个 $x \in |X|$ 处相对于 $Y$ 平坦，
则称 $\mathcal{F}$ {\it 在 $Y$ 上平坦}。
\end{enumerate}
\end{definition}

\noindent
作出定义后，我们得到必需的基变换引理。该引理说明，
形成拟凝聚层的平坦点集与平坦基变换可交换。

\begin{lemma}
\label{spaces-morphisms-lemma-base-change-module-flat}
设 $S$ 为概形。设
$$
\xymatrix{
X' \ar[d]_{f'} \ar[r]_{g'} & X \ar[d]^f \\
Y' \ar[r]^g & Y
}
$$
为 $S$ 上代数空间的 Cartesian 图。设 $x' \in |X'|$ 的像为 $x \in |X|$。
设 $\mathcal{F}$ 为 $X$ 上的拟凝聚层，并记
$\mathcal{F}' = (g')^*\mathcal{F}$。
\begin{enumerate}
\item 若 $\mathcal{F}$ 在 $x$ 处相对于 $Y$ 平坦，则 $\mathcal{F}'$
在 $x'$ 处相对于 $Y'$ 平坦。
\item 若 $g$ 在 $f'(x')$ 处平坦，且 $\mathcal{F}'$ 在 $x'$ 处相对于
$Y'$ 平坦，则 $\mathcal{F}$ 在 $x$ 处相对于 $Y$ 平坦。
\end{enumerate}
特别地，若 $\mathcal{F}$ 在 $Y$ 上平坦，则 $\mathcal{F}'$ 在 $Y'$ 上平坦。
\end{lemma}

\begin{proof}
取概形 $V$ 及满 \'etale 态射 $V \to Y$。
取概形 $U$ 及满 \'etale 态射 $U \to V \times_Y X$。
取概形 $V'$ 及满 \'etale 态射 $V' \to V \times_Y Y'$。
于是 $U' = V' \times_V U$ 是概形，并带有满 \'etale 态射
$U' = V' \times_V U \to Y' \times_Y X = X'$。取 $u' \in U'$ 映到
$x' \in |X'|$。于是可以用 $\mathcal{F}'|_{U'}$ 在 $u'$ 处相对于 $V'$
的平坦性，检验 $\mathcal{F}'$ 在 $x'$ 处相对于 $Y'$ 的平坦性。
故本引理由《态射的进一步内容》引理
\ref{more-morphisms-lemma-flat-locus-base-change} 得出。
\end{proof}

\noindent
下述引理用模来讨论平坦态射的“复合”。它还说明平坦性满足一种自上而下的下降。

\begin{lemma}
\label{spaces-morphisms-lemma-composition-module-flat}
设 $S$ 为概形。设 $X \to Y \to Z$ 为 $S$ 上代数空间的态射。
设 $\mathcal{F}$ 为 $X$ 上的拟凝聚层。设 $x \in |X|$ 的像为 $y \in |Y|$。
\begin{enumerate}
\item 若 $\mathcal{F}$ 在 $x$ 处相对于 $Y$ 平坦，且 $Y$ 在 $y$ 处相对于
$Z$ 平坦，则 $\mathcal{F}$ 在 $x$ 处相对于 $Z$ 平坦。
\item 设 $x : \Spec(K) \to X$ 为 $x$ 的一个代表。若
\begin{enumerate}
\item $\mathcal{F}$ 在 $x$ 处相对于 $Y$ 平坦；
\item $x^*\mathcal{F} \not = 0$；
\item $\mathcal{F}$ 在 $x$ 处相对于 $Z$ 平坦；
\end{enumerate}
则 $Y$ 在 $y$ 处相对于 $Z$ 平坦。
\item 设 $\overline{x}$ 为 $X$ 的一个几何点，位于 $x$ 上方，
并在 $Y$ 中具有像 $\overline{y}$。若 $\mathcal{F}_{\overline{x}}$ 是
忠实平坦 $\mathcal{O}_{Y, \overline{y}}$-模，且 $\mathcal{F}$ 在 $x$
处相对于 $Z$ 平坦，则 $Y$ 在 $y$ 处相对于 $Z$ 平坦。
\end{enumerate}
\end{lemma}

\begin{proof}
取 (3) 中的 $\overline{x}$ 与 $\overline{y}$，并以 $\overline{z}$ 表示所诱导的
$Z$ 的几何点。利用引理 \ref{spaces-morphisms-lemma-flat-at-point} 与
\ref{spaces-morphisms-lemma-flat-at-point-etale-local-rings} 中的平坦性刻画，本引理归结为
关于局部环映射
$\mathcal{O}_{Z, \overline{z}} \to \mathcal{O}_{Y, \overline{y}}$
和模 $\mathcal{F}_{\overline{x}}$ 的纯代数问题。(1) 由《代数》引理
\ref{algebra-lemma-composition-flat} 得出。注意条件 (2)(b) 保证
$\mathcal{F}_{\overline{x}}/
\mathfrak m_{\overline{y}} \mathcal{F}_{\overline{x}}$
非零。因此 (2)(a) $+$ (2)(b) 蕴含 $\mathcal{F}_{\overline{x}}$ 是忠实平坦
$\mathcal{O}_{Y, \overline{y}}$-模，见《代数》引理
\ref{algebra-lemma-ff}。故 (2) 是 (3) 的特例。最后，(3) 由《代数》引理
\ref{algebra-lemma-flat-permanence} 得出。
\end{proof}

\noindent
有时基变换发生在“上层”。下面给出精确表述。

\begin{lemma}
\label{spaces-morphisms-lemma-flat-permanence}
设 $S$ 为概形。设 $f : X \to Y$、$g : Y \to Z$ 为 $S$ 上代数空间的态射。
设 $\mathcal{G}$ 为 $Y$ 上的拟凝聚层。设 $x \in |X|$ 的像为 $y \in |Y|$。
若 $f$ 在 $x$ 处平坦，则
$$
\mathcal{G}\text{ 相对于 }Z\text{ 在此点平坦：}y
\Leftrightarrow
f^*\mathcal{G}\text{ 相对于 }Z\text{ 在此点平坦：}x.
$$
特别地，若 $f$ 满且平坦，则 $\mathcal{G}$ 在 $Z$ 上平坦，
当且仅当 $f^*\mathcal{G}$ 在 $Z$ 上平坦。
\end{lemma}

\begin{proof}
取 $X$ 的几何点 $\overline{x}$，并分别以 $\overline{y}$、$\overline{z}$
表示它在 $Y$、$Z$ 中的像。利用引理 \ref{spaces-morphisms-lemma-flat-at-point} 与
\ref{spaces-morphisms-lemma-flat-at-point-etale-local-rings} 中的平坦性刻画，
以及《空间的性质》引理
\ref{spaces-properties-lemma-stalk-pullback-quasi-coherent} 对
$f^*\mathcal{G}$ 在 $\overline{x}$ 处之茎的描述，本引理归结为关于局部环映射
$\mathcal{O}_{Z, \overline{z}} \to \mathcal{O}_{Y, \overline{y}}
\to \mathcal{O}_{X, \overline{x}}$
和模 $\mathcal{G}_{\overline{y}}$ 的纯代数问题。该代数命题即《代数》引理
\ref{algebra-lemma-flatness-descends-more-general}。
\end{proof}

\begin{lemma}
\label{spaces-morphisms-lemma-pf-flat-module-open}
设 $S$ 为概形。
设 $f : X \to Y$ 为 $S$ 上代数空间的态射。
设 $\mathcal{F}$ 为拟凝聚 $\mathcal{O}_X$-模。假设 $f$ 局部有限表示，
$\mathcal{F}$ 有限型，$X = \text{Supp}(\mathcal{F})$，且
$\mathcal{F}$ 在 $Y$ 上平坦。则 $f$ 普遍开。
\end{lemma}

\begin{proof}
取满 \'etale 态射 $\varphi : V \to Y$，其中 $V$ 为概形。
取满 \'etale 态射 $U \to V \times_Y X$，其中 $U$ 为概形。
于是只需对 $U \to V$ 及拟凝聚 $\mathcal{O}_V$-模
$\varphi^*\mathcal{F}$ 证明引理。因此本引理归结为概形情形，见
《态射》引理 \ref{morphisms-lemma-pf-flat-module-open}。
\end{proof}






\section{泛平坦性}
\label{spaces-morphisms-section-generic-flatness}

\noindent
本节是《态射》第 \ref{morphisms-section-generic-flatness} 节的对应版本。

\begin{proposition}
\label{spaces-morphisms-proposition-generic-flatness-reduced}
设 $S$ 为概形。
设 $f : X \to Y$ 为 $S$ 上代数空间的态射。
设 $\mathcal{F}$ 为拟凝聚 $\mathcal{O}_X$-模层。假设
\begin{enumerate}
\item $Y$ 约化；
\item $f$ 有限型；
\item $\mathcal{F}$ 是有限型 $\mathcal{O}_X$-模。
\end{enumerate}
则存在稠密开子空间 $W \subset Y$，使得 $f$ 的基变换 $X_W \to W$
平坦、局部有限表示且拟紧，并使得 $\mathcal{F}|_{X_W}$ 在 $W$ 上平坦，
且在 $\mathcal{O}_{X_W}$ 上有限表示。
\end{proposition}

\begin{proof}
设 $V$ 为概形，并设 $V \to Y$ 为满 \'etale 态射。令
$X_V = V \times_Y X$，并以 $\mathcal{F}_V$ 表示 $\mathcal{F}$ 在 $X_V$
上的限制。假设结论对态射 $X_V \to V$ 与层 $\mathcal{F}_V$ 成立。
于是存在开子概形 $V' \subset V$，使得 $X_{V'} \to V'$ 平坦且有限表示，
并且 $\mathcal{F}_{V'}$ 是在 $V'$ 上平坦的有限表示
$\mathcal{O}_{X_{V'}}$-模。令 $W \subset Y$ 为 \'etale 态射 $V' \to Y$
的像，见《空间的性质》引理
\ref{spaces-properties-lemma-etale-image-open}。于是 $V' \to W$ 是满
\'etale 态射，故由引理 \ref{spaces-morphisms-lemma-finite-presentation-local}、
\ref{spaces-morphisms-lemma-flat-local} 与 \ref{spaces-morphisms-lemma-quasi-compact-local}，可见
$X_W \to W$ 平坦、局部有限表示且拟紧。由《空间的性质》第
\ref{spaces-properties-section-properties-modules} 节中的讨论，
可见 $\mathcal{F}_W$ 作为 $\mathcal{O}_{X_W}$-模有限表示；
由引理 \ref{spaces-morphisms-lemma-base-change-module-flat}，可见 $\mathcal{F}_W$
在 $W$ 上平坦。这一论证把命题归结为 $Y$ 是概形的情形。

\medskip\noindent
假设可以在 $Y$ 为仿射概形时证明命题。设 $f : X \to Y$ 为 $S$ 上
代数空间的有限型态射，其中 $Y$ 为概形；并设 $\mathcal{F}$ 为有限型拟凝聚
$\mathcal{O}_X$-模。取仿射开覆盖 $Y = \bigcup V_j$。按假设，可以找到
稠密开集 $W_j \subset V_j$，使得 $X_{W_j} \to W_j$ 平坦、局部有限表示且拟紧，
并使得 $\mathcal{F}|_{X_{W_j}}$ 在 $W_j$ 上平坦，且作为
$\mathcal{O}_{X_{W_j}}$-模有限表示。此时只需取 $W = \bigcup W_j$，
结论即得。因此命题归结为 $Y$ 是仿射概形的情形。

\medskip\noindent
设 $Y$ 为 $S$ 上的仿射概形，设 $f : X \to Y$ 为 $S$ 上代数空间的有限型态射，
并设 $\mathcal{F}$ 为有限型拟凝聚 $\mathcal{O}_X$-模。由于 $f$ 有限型，
它拟紧，故 $X$ 拟紧。因此可取仿射概形 $U$ 及满 \'etale 态射 $U \to X$，
见《空间的性质》引理
\ref{spaces-properties-lemma-quasi-compact-affine-cover}。注意，$U \to Y$
有限型（在此情形下，这正是 $f$ 有限型的含义）。因此可应用《态射》命题
\ref{morphisms-proposition-generic-flatness-reduced}，从而存在稠密开集
$W \subset Y$，使得 $U_W \to W$ 平坦且有限表示，并使得
$\mathcal{F}|_{U_W}$ 在 $W$ 上平坦，且作为 $\mathcal{O}_{U_W}$-模有限表示。
按我们的定义，这意味着 $f$ 的基变换 $X_W \to W$ 平坦、局部有限表示且拟紧，
并且 $\mathcal{F}|_{X_W}$ 在 $W$ 上平坦，且在 $\mathcal{O}_{X_W}$ 上有限表示。
\end{proof}

\noindent
不能把上面引理的结论加强为要求 $X_W \to W$ 有限表示，因为
$\mathbf{A}^1_{\mathbf{Q}}/\mathbf{Z} \to \Spec(\mathbf{Q})$ 给出了反例。
问题在于对角态射 $\Delta_{X/Y}$ 可能不拟紧，即 $f$ 可能不拟分离。
显然，这也是唯一的问题。

\begin{proposition}
\label{spaces-morphisms-proposition-generic-flatness-reduced-quasi-separated}
设 $S$ 为概形。
设 $f : X \to Y$ 为 $S$ 上代数空间的态射。
设 $\mathcal{F}$ 为拟凝聚 $\mathcal{O}_X$-模层。假设
\begin{enumerate}
\item $Y$ 约化；
\item $f$ 拟分离；
\item $f$ 有限型；
\item $\mathcal{F}$ 是有限型 $\mathcal{O}_X$-模。
\end{enumerate}
则存在稠密开子空间 $W \subset Y$，使得 $f$ 的基变换 $X_W \to W$
平坦且有限表示，并使得 $\mathcal{F}|_{X_W}$ 在 $W$ 上平坦，
且在 $\mathcal{O}_{X_W}$ 上有限表示。
\end{proposition}

\begin{proof}
这立即由命题 \ref{spaces-morphisms-proposition-generic-flatness-reduced} 以及下述事实得出：
“有限表示” $=$ “局部有限表示” $+$ “拟紧” $+$ “拟分离”。
\end{proof}















\section{相对维数}
\label{spaces-morphisms-section-relative-dimension}

\noindent
本节定义代数空间态射在一点处的相对维数，以及若干密切相关的性质。

\begin{definition}
\label{spaces-morphisms-definition-dimension-fibre}
设 $S$ 为概形。
设 $f : X \to Y$ 为 $S$ 上代数空间的态射。
设 $x \in |X|$。设 $d, r \in \{0, 1, 2, \ldots, \infty\}$。
\begin{enumerate}
\item 若引理 \ref{spaces-morphisms-lemma-local-source-target-at-point} 中的等价条件对
《下降》引理 \ref{descent-lemma-dimension-local-ring-fibre} 所描述的性质
$\mathcal{P}_d$ 成立，则称{\it $f$ 在 $x$ 处之纤维的局部环维数}为 $d$。
\item 若引理 \ref{spaces-morphisms-lemma-local-source-target-at-point} 中的等价条件对
《下降》引理 \ref{descent-lemma-transcendence-degree-at-point} 所描述的性质
$\mathcal{P}_r$ 成立，则称{\it $x/f(x)$ 的超越次数}为 $r$。
\item 若引理 \ref{spaces-morphisms-lemma-local-source-target-at-point} 中的等价条件对
《下降》引理 \ref{descent-lemma-dimension-at-point} 所描述的性质
$\mathcal{P}_d$ 成立，则称{\it $f$ 在 $x$ 处具有相对维数 $d$}。
\end{enumerate}
\end{definition}

\noindent
把其含义明确写出。具体地，取引理
\ref{spaces-morphisms-lemma-local-source-target-at-point} 中的图
$$
\xymatrix{
U \ar[d]_a \ar[r]_h & V \ar[d]^b \\
X \ar[r]^f & Y
}
\quad\quad
\xymatrix{
u \ar[d] \ar[r] & v \ar[d] \\
x \ar[r] & y
}
$$
。则有
$$
\begin{matrix}
\text{相对维数：态射 }f\text{ 于点 }x & = &
\dim_u (U_v) \\
\text{纤维局部环的维数：态射 }f\text{ 于点 }x & = &
\dim(\mathcal{O}_{U_v, u})\\
\text{超越次数：}x/f(x) & = &
\text{trdeg}_{\kappa(v)}(\kappa(u))
\end{matrix}
$$
注意，若 $Y = \Spec(k)$ 是域的谱，则 $X/Y$ 在 $x$ 处的相对维数等于
$\dim_x(X)$；$x/f(x)$ 的超越次数就是在 $k$ 上的超越次数；而 $f$ 在 $x$
处之纤维的局部环维数就是 $x$ 处局部环的维数。换言之，在这种情形下，
相对概念变成绝对概念。

\begin{definition}
\label{spaces-morphisms-definition-relative-dimension}
设 $S$ 为概形。
设 $f : X \to Y$ 为 $S$ 上代数空间的态射。
设 $d \in \{0, 1, 2, \ldots\}$。
\begin{enumerate}
\item 若 $f$ 在所有 $x \in |X|$ 处都具有相对维数 $\leq d$，
则称 $f$ {\it 具有相对维数 $\leq d$}。
\item 若 $f$ 在所有 $x \in |X|$ 处都具有相对维数 $d$，
则称 $f$ {\it 具有相对维数 $d$}。
\end{enumerate}
\end{definition}

\noindent
粗略地说，具有{\it 等于} $d$ 的相对维数意味着所有非空纤维都是维数为
$d$ 的等维空间。

\begin{lemma}
\label{spaces-morphisms-lemma-compare-tr-deg}
设 $S$ 为概形。设 $X \to Y \to Z$ 为 $S$ 上代数空间的态射。
设 $x \in |X|$ 的像分别为 $y \in |Y|$、$z \in |Z|$。
假设 $X \to Y$ 局部拟有限，且 $Y \to Z$ 局部有限型。
则 $x/z$ 的超越次数等于 $y/z$ 的超越次数。
\end{lemma}

\begin{proof}
可以选取交换图
$$
\xymatrix{
U \ar[d] \ar[r] & V \ar[d] \ar[r] & W \ar[d] \\
X \ar[r] & Y \ar[r] & Z
}
\quad\quad
\xymatrix{
u \ar[d] \ar[r] & v \ar[d] \ar[r] & w \ar[d] \\
x \ar[r] & y \ar[r] & z
}
$$
其中 $U, V, W$ 为概形，竖直箭头为 \'etale 态射。按定义，态射
$U \to V$ 局部拟有限，这蕴含 $\kappa(v) \subset \kappa(u)$ 有限，见
《态射》引理 \ref{morphisms-lemma-residue-field-quasi-finite}。
故结论显然。
\end{proof}

\begin{lemma}
\label{spaces-morphisms-lemma-jacobson-finite-type-points}
设 $S$ 为概形。设 $f : X \to Y$ 为 $S$ 上代数空间的态射。
若 $f$ 局部有限型，$Y$ 是 Jacobson 的（《空间的性质》评注
\ref{spaces-properties-remark-list-properties-local-etale-topology}），
且 $x \in |X|$ 是 $X$ 的有限型点，则 $x/f(x)$ 的超越次数为 $0$。
\end{lemma}

\begin{proof}
取概形 $V$ 及满 \'etale 态射 $V \to Y$。
取概形 $U$ 及满 \'etale 态射 $U \to X \times_Y V$。
由引理 \ref{spaces-morphisms-lemma-finite-type-points-surjective-morphism}，
可以找到映到 $x$ 的有限型点 $u \in U$。缩小 $U$ 后，可以假设
$u \in U$ 是闭点（《态射》引理
\ref{morphisms-lemma-identify-finite-type-points}）。
设 $v \in V$ 为 $u$ 的像。由《态射》引理
\ref{morphisms-lemma-jacobson-finite-type-points}，扩张
$\kappa(u)/\kappa(v)$ 有限。这就完成证明。
\end{proof}

\begin{lemma}
\label{spaces-morphisms-lemma-rel-dimension-dimension}
设 $S$ 为概形。设 $f : X \to Y$ 为 $S$ 上局部 Noether 代数空间的态射，
且它平坦、局部有限型并具有相对维数 $d$。对 $|X|$ 中每个点 $x$，
若其在 $|Y|$ 中的像为 $y$，则 $\dim_x(X) = \dim_y(Y) + d$。
\end{lemma}

\begin{proof}
由代数空间在一点处之维数的定义（《空间的性质》定义
\ref{spaces-properties-definition-dimension-at-point}）以及具有相对维数 $d$
的定义，问题归结为概形的相应命题（《态射》引理
\ref{morphisms-lemma-rel-dimension-dimension}）。
\end{proof}







\section{态射与纤维维数}
\label{spaces-morphisms-section-dimension-fibres}

\noindent
本节是《态射》第 \ref{morphisms-section-dimension-fibres} 节的对应版本。
由于我们没有代数空间在点处的局部环，本节的表述稍显别扭。

\begin{lemma}
\label{spaces-morphisms-lemma-dimension-fibre-at-a-point}
设 $S$ 为概形。
设 $f : X \to Y$ 为 $S$ 上代数空间的态射。
设 $x \in |X|$。假设 $f$ 局部有限型。则有
$$
\begin{matrix}
\text{相对维数：态射 }f\text{ 于点 }x \\
= \\
\text{纤维局部环的维数：态射 }f\text{ 于点 }x \\
+ \\
\text{超越次数：}x/f(x)
\end{matrix}
$$
，其中记号如定义 \ref{spaces-morphisms-definition-dimension-fibre} 所示。
\end{lemma}

\begin{proof}
这立即由《态射》引理
\ref{morphisms-lemma-dimension-fibre-at-a-point} 得出：将其应用于
$h : U \to V$ 与 $u \in U$，其中这些对象如引理
\ref{spaces-morphisms-lemma-local-source-target-at-point} 所示。
\end{proof}

\begin{lemma}
\label{spaces-morphisms-lemma-dimension-fibre-at-a-point-additive}
设 $S$ 为概形。
设 $f : X \to Y$ 与 $g : Y \to Z$ 为 $S$ 上代数空间的态射。
设 $x \in |X|$，并令 $y = f(x)$。假设 $f$ 与 $g$ 局部有限型。则
\begin{enumerate}
\item
$$
\begin{matrix}
\text{相对维数：态射 }g \circ f\text{ 于点 }x \\
\leq \\
\text{相对维数：态射 }f\text{ 于点 }x \\
+ \\
\text{相对维数：态射 }g\text{ 于点 }y
\end{matrix}
$$
\item 若存在来自某个域之谱的态射 $\Spec(k) \to Z$，它属于
$g(f(x)) = g(y)$ 的等价类，且态射 $X_k \to Y_k$ 在 $x$ 处平坦，
则 (1) 中取等号；例如，当 $f$ 在 $x$ 处平坦时即如此；
\item
$$
\begin{matrix}
\text{超越次数：}x/g(f(x)) \\
= \\
\text{超越次数：}x/f(x) \\
+ \\
\text{超越次数：}f(x)/g(f(x))
\end{matrix}
$$
\end{enumerate}
\end{lemma}

\begin{proof}
取图
$$
\xymatrix{
U \ar[d] \ar[r] & V \ar[d] \ar[r] & W \ar[d] \\
X \ar[r] & Y \ar[r] & Z
}
$$
其中 $U, V, W$ 为概形，竖直箭头满且 \'etale。（见《空间》引理
\ref{spaces-lemma-lift-morphism-presentations}。）取映到 $x$ 的 $u \in U$。
分别以 $v, w$ 表示 $u$ 在 $V, W$ 中的像。对顶行及点 $u, v, w$
应用《态射》引理
\ref{morphisms-lemma-dimension-fibre-at-a-point-additive}。细节从略。
\end{proof}

\begin{lemma}
\label{spaces-morphisms-lemma-dimension-fibre-after-base-change}
设 $S$ 为概形。设
$$
\xymatrix{
X' \ar[r]_{g'} \ar[d]_{f'} & X \ar[d]^f \\
Y' \ar[r]^g & Y
}
$$
为 $S$ 上代数空间的纤维积图。设 $x' \in |X'|$。令 $x = g'(x')$。
假设 $f$ 局部有限型。则
\begin{enumerate}
\item
$$
\begin{matrix}
\text{相对维数：态射 }f\text{ 于点 }x \\
= \\
\text{相对维数：态射 }f'\text{ 于点 }x'
\end{matrix}
$$
\item 有
$$
\begin{matrix}
\text{纤维局部环的维数：态射 }f'\text{ 于点 }x' \\
- \\
\text{纤维局部环的维数：态射 }f\text{ 于点 }x \\
= \\
\text{超越次数：}x/f(x) \\
- \\
\text{超越次数：}x'/f'(x')
\end{matrix}
$$
且共同的值为 $\geq 0$；
\item 给定 $x$ 以及映到同一 $y \in |Y|$ 的 $y' \in |Y'|$，
存在 $x'$ 的一个选取，使得 (2) 中的整数为 $0$。
\end{enumerate}
\end{lemma}

\begin{proof}
取满 \'etale 态射 $V \to Y$，其中 $V$ 为概形。
取满 \'etale 态射 $U \to V \times_Y X$，其中 $U$ 为概形。
取满 \'etale 态射 $V' \to V \times_Y Y'$，其中 $V'$ 为概形。
令 $U' = V' \times_V U$。于是所诱导的态射 $U' \to X'$ 也满且
\'etale（论证从略）。取映到 $x'$ 的 $u' \in U'$。此时，对涉及
$U', U, V', V$ 的概形图及点 $u'$ 应用《态射》引理
\ref{morphisms-lemma-dimension-fibre-after-base-change}，即可得到 (1) 与 (2)。
为证明 (3)，先取映到 $y$ 的 $v \in V$。再利用《空间的性质》引理
\ref{spaces-properties-lemma-points-cartesian}，可以选取映到 $y'$ 与 $v$
的 $v' \in V'$，以及映到 $x$ 与 $v$ 的 $u \in U$。最后，由《态射》引理
\ref{morphisms-lemma-dimension-fibre-after-base-change}，可以选取映到
$v'$ 与 $u$ 的 $u' \in U'$，使该整数为零。于是取 $u'$ 的像
$x' \in |X'|$ 即可。
\end{proof}

\begin{lemma}
\label{spaces-morphisms-lemma-openness-bounded-dimension-fibres}
设 $S$ 为概形。
设 $f : X \to Y$ 为 $S$ 上代数空间的态射。
设 $n \geq 0$。假设 $f$ 局部有限型。集合
$$
W_n = \{x \in |X|
\text{ 满足相对维数：态射 }f\text{ 于点 } x \leq n\}
$$
在 $|X|$ 中是开的。
\end{lemma}

\begin{proof}
取图
$$
\xymatrix{
U \ar[r]_h \ar[d]_a & V \ar[d] \\
X \ar[r] & Y
}
$$
其中 $U$ 与 $V$ 为概形，竖直箭头满且 \'etale，见《空间》引理
\ref{spaces-lemma-lift-morphism-presentations}。由《态射》引理
\ref{morphisms-lemma-openness-bounded-dimension-fibres}，$h$ 具有相对维数
$\leq n$ 的点集 $U_n$ 在 $U$ 中是开的。由代数空间态射在点处之相对维数的
定义，可见 $U_n = a^{-1}(W_n)$。本引理由 $|X|$ 上拓扑的定义得出。
\end{proof}

\begin{lemma}
\label{spaces-morphisms-lemma-openness-bounded-dimension-fibres-finite-presentation}
设 $S$ 为概形。
设 $f : X \to Y$ 为 $S$ 上代数空间的态射。
设 $n \geq 0$。假设 $f$ 局部有限表示。开集
$$
W_n = \{x \in |X|
\text{ 满足相对维数：态射 }f\text{ 于点 } x \leq n\}
$$
（引理 \ref{spaces-morphisms-lemma-openness-bounded-dimension-fibres}）在 $|X|$ 中逆紧。
（见《拓扑》定义 \ref{topology-definition-quasi-compact}。）
\end{lemma}

\begin{proof}
取图
$$
\xymatrix{
U \ar[r]_h \ar[d]_a & V \ar[d] \\
X \ar[r] & Y
}
$$
其中 $U$ 与 $V$ 为概形，竖直箭头满且 \'etale，见《空间》引理
\ref{spaces-lemma-lift-morphism-presentations}。在引理
\ref{spaces-morphisms-lemma-openness-bounded-dimension-fibres} 的证明中，我们已看到
$a^{-1}(W_n) = U_n$ 是态射 $h$ 的相应点集。由《态射》引理
\ref{morphisms-lemma-openness-bounded-dimension-fibres-finite-presentation}，
可见 $U_n$ 在 $U$ 中逆紧。本引理由 $|X|$ 上拓扑的定义得出；
可与《空间的性质》引理
\ref{spaces-properties-lemma-space-locally-quasi-compact} 及其证明比较。
\end{proof}

\begin{lemma}
\label{spaces-morphisms-lemma-locally-quasi-finite-rel-dimension-0}
设 $S$ 为概形。
设 $f : X \to Y$ 为 $S$ 上代数空间的态射。
假设 $f$ 局部有限型。则 $f$ 局部拟有限，当且仅当 $f$ 在每个
$x \in |X|$ 处具有相对维数 $0$。
\end{lemma}

\begin{proof}
取图
$$
\xymatrix{
U \ar[r]_h \ar[d]_a & V \ar[d] \\
X \ar[r] & Y
}
$$
其中 $U$ 与 $V$ 为概形，竖直箭头满且 \'etale，见《空间》引理
\ref{spaces-lemma-lift-morphism-presentations}。定义蕴含：$h$ 局部拟有限，
当且仅当 $f$ 局部拟有限；而 $f$ 在所有 $x \in |X|$ 处具有相对维数 $0$，
当且仅当 $h$ 在所有 $u \in U$ 处具有相对维数 $0$。
因此结论由 $h$ 的相应结果得出，即《态射》引理
\ref{morphisms-lemma-locally-quasi-finite-rel-dimension-0}。
\end{proof}

\begin{lemma}
\label{spaces-morphisms-lemma-locally-finite-type-quasi-finite-part}
设 $S$ 为概形。
设 $f : X \to Y$ 为 $S$ 上代数空间的态射。
假设 $f$ 局部有限型。则存在典范开子空间 $X' \subset X$，使得
$f|_{X'} : X' \to Y$ 局部拟有限，并且对任意 $x \in |X|$，
若 $x \not \in |X'|$，则 $f$ 在该点的相对维数为 $\geq 1$。
形成 $X'$ 与任意基变换可交换。
\end{lemma}

\begin{proof}
合用引理 \ref{spaces-morphisms-lemma-openness-bounded-dimension-fibres}、
\ref{spaces-morphisms-lemma-locally-quasi-finite-rel-dimension-0} 与
\ref{spaces-morphisms-lemma-dimension-fibre-after-base-change}。
\end{proof}

\begin{lemma}
\label{spaces-morphisms-lemma-quasi-finite-at-point}
设 $S$ 为概形。考虑 Cartesian 图
$$
\xymatrix{
X \ar[d] & F \ar[l]^p \ar[d] \\
Y & \Spec(k) \ar[l]
}
$$
其中 $X \to Y$ 是 $S$ 上代数空间的局部有限型态射，而 $k$ 是 $S$ 上的域。
设 $z \in |F|$ 满足 $\dim_z(F) = 0$。则以包含 $p(z)$ 的开子空间替换 $X$ 后，
态射
$$
X \longrightarrow Y
$$
局部拟有限。
\end{lemma}

\begin{proof}
令 $X' \subset X$ 为引理 \ref{spaces-morphisms-lemma-locally-finite-type-quasi-finite-part}
给出的、使 $f$ 局部拟有限的开子空间。由于形成 $X'$ 与任意基变换可交换，
可见 $z \in X' \times_Y \Spec(k)$。故引理显然。
\end{proof}





\section{维数公式}
\label{spaces-morphisms-section-dimension-formula}

\noindent
维数公式的对应版本（《态射》引理
\ref{morphisms-lemma-dimension-formula}）表述起来有些棘手，
因为我们必须定义整代数空间（稍后会这样做），还要定义普遍链状代数空间。
不过，下面的版本很直接。

\begin{lemma}
\label{spaces-morphisms-lemma-dimension-formula-general}
设 $S$ 为概形。设 $f : X \to Y$ 为 $S$ 上代数空间的态射。
假设 $Y$ 局部 Noether，且 $f$ 局部有限型。设 $x \in |X|$ 的像为
$y \in |Y|$。则有
\begin{align*}
& \text{局部环维数：空间 }X\text{ 于点 }x \leq \\
& \text{局部环维数：空间 }Y\text{ 于点 }y + E - \\
& \text{超越次数：}x/y
\end{align*}
这里 $E$ 是 $\xi/f(\xi)$ 的超越次数之最大值，其中 $\xi \in |X|$
遍历专门化到 $x$ 且使 $X$ 的局部环维数为 $0$ 的点。
\end{lemma}

\begin{proof}
取仿射概形 $V$、\'etale 态射 $V \to Y$，以及映到 $y$ 的点 $v \in V$。
取仿射概形 $U$、\'etale 态射 $U \to X \times_Y V$，以及点 $u \in U$，
它在 $V$ 中映到 $v$，在 $X$ 中映到 $x$。展开定义
\ref{spaces-morphisms-definition-dimension-fibre} 与《空间的性质》定义
\ref{spaces-properties-definition-dimension-local-ring}，我们要证明
$$
\dim(\mathcal{O}_{U, u}) \leq
\dim(\mathcal{O}_{V, v}) + E - \text{trdeg}_{\kappa(v)}(\kappa(u))
$$
。设 $\xi_U \in U$ 为 $U$ 的某个包含 $u$ 的不可约分支之泛点。
则 $\xi_U$ 映到一点 $\xi \in |X|$，它位于定义量 $E$ 时所用的列表中；
事实上，$E$ 的定义中使用的每个 $\xi$ 都以此方式出现（略去一个小细节）。
特别地，这样的 $\xi$ 只有有限个，因而可以取最大值（即它确实是最大值，
而非上确界）。$\xi$ 在 $f(\xi)$ 上的超越次数是
$\text{trdeg}_{\kappa(\xi_V)}(\kappa(\xi_U))$，其中 $\xi_V \in V$
是 $\xi_U$ 的像。因此本引理由《态射》引理
\ref{morphisms-lemma-dimension-formula-general} 得出。
\end{proof}

\begin{lemma}
\label{spaces-morphisms-lemma-alteration-dimension-general}
设 $S$ 为概形。
设 $f : X \to Y$ 为 $S$ 上代数空间的态射。假设 $Y$ 局部 Noether，
且 $f$ 局部有限型。则
$$
\dim(X) \leq \dim(Y) + E
$$
，其中 $E$ 是 $\xi/f(\xi)$ 的超越次数之上确界，而 $\xi$ 遍历使 $X$
的局部环维数为 $0$ 的点。
\end{lemma}

\begin{proof}
这是引理 \ref{spaces-morphisms-lemma-dimension-formula-general} 与《空间的性质》引理
\ref{spaces-properties-lemma-dimension} 的直接推论。
\end{proof}







\section{Syntomic 态射}
\label{spaces-morphisms-section-syntomic}

\noindent
概形态射的“syntomic”性质在源与靶上都是 \'etale 局部的，见
《下降》评注 \ref{descent-remark-list-local-source-target}。
此性质也在基变换下稳定，并且在靶上是 fpqc 局部的，见《态射》引理
\ref{morphisms-lemma-base-change-syntomic} 与《下降》引理
\ref{descent-lemma-descending-property-syntomic}。因此，由上面的引理
\ref{spaces-morphisms-lemma-local-source-target}，我们可以如下定义代数空间的 syntomic 态射；
当该态射可表时，这与第 \ref{spaces-morphisms-section-representable} 节中已有的概念一致。

\begin{definition}
\label{spaces-morphisms-definition-syntomic}
设 $S$ 为概形。
设 $f : X \to Y$ 为 $S$ 上代数空间的态射。
\begin{enumerate}
\item 若引理 \ref{spaces-morphisms-lemma-local-source-target} 中取
$\mathcal{P} =$“syntomic”时其中的等价条件成立，则称 $f$ {\it syntomic}。
\item 设 $x \in |X|$。若存在 $x$ 的开邻域 $X' \subset X$，使得
$f|_{X'} : X' \to Y$ syntomic，则称 $f$ {\it 在 $x$ 处 syntomic}。
\end{enumerate}
\end{definition}

\begin{lemma}
\label{spaces-morphisms-lemma-composition-syntomic}
Syntomic 态射的复合 syntomic。
\end{lemma}

\begin{proof}
见评注 \ref{spaces-morphisms-remark-composition-P} 和《态射》引理
\ref{morphisms-lemma-composition-syntomic}。
\end{proof}

\begin{lemma}
\label{spaces-morphisms-lemma-base-change-syntomic}
Syntomic 态射的基变换 syntomic。
\end{lemma}

\begin{proof}
见评注 \ref{spaces-morphisms-remark-base-change-P} 和《态射》引理
\ref{morphisms-lemma-base-change-syntomic}。
\end{proof}

\begin{lemma}
\label{spaces-morphisms-lemma-syntomic-local}
设 $S$ 为概形。
设 $f : X \to Y$ 为 $S$ 上代数空间的态射。
下列条件等价：
\begin{enumerate}
\item $f$ syntomic；
\item 对每个 $x \in |X|$，态射 $f$ 在 $x$ 处 syntomic；
\item 对每个概形 $Z$ 及任意态射 $Z \to Y$，态射
$Z \times_Y X \to Z$ syntomic；
\item 对每个仿射概形 $Z$ 及任意态射 $Z \to Y$，态射
$Z \times_Y X \to Z$ syntomic；
\item 存在概形 $V$ 及满 \'etale 态射 $V \to Y$，使得
$V \times_Y X \to V$ 为 syntomic 态射；
\item 存在概形 $U$ 及满 \'etale 态射 $\varphi : U \to X$，使得复合
$f \circ \varphi$ syntomic；
\item 对每个交换图
$$
\xymatrix{
U \ar[d] \ar[r] & V \ar[d] \\
X \ar[r] & Y
}
$$
其中 $U$、$V$ 为概形，竖直箭头为 \'etale 态射，顶部水平箭头 syntomic；
\item 存在交换图
$$
\xymatrix{
U \ar[d] \ar[r] & V \ar[d] \\
X \ar[r] & Y
}
$$
其中 $U$、$V$ 为概形，竖直箭头为 \'etale 态射，且 $U \to X$ 满，
使得顶部水平箭头 syntomic；
\item 存在 Zariski 覆盖 $Y = \bigcup_{i \in I} Y_i$ 以及
$f^{-1}(Y_i) = \bigcup X_{ij}$，使得每个态射
$X_{ij} \to Y_i$ syntomic。
\end{enumerate}
\end{lemma}

\begin{proof}
略。
\end{proof}

\begin{lemma}
\label{spaces-morphisms-lemma-syntomic-locally-finite-presentation}
Syntomic 态射局部有限表示。
\end{lemma}

\begin{proof}
这立即由概形情形得出（《态射》引理
\ref{morphisms-lemma-syntomic-locally-finite-presentation}）。
\end{proof}

\begin{lemma}
\label{spaces-morphisms-lemma-syntomic-flat}
Syntomic 态射平坦。
\end{lemma}

\begin{proof}
这立即由概形情形得出（《态射》引理
\ref{morphisms-lemma-syntomic-flat}）。
\end{proof}

\begin{lemma}
\label{spaces-morphisms-lemma-syntomic-open}
Syntomic 态射普遍开。
\end{lemma}

\begin{proof}
合用引理 \ref{spaces-morphisms-lemma-syntomic-locally-finite-presentation}、
\ref{spaces-morphisms-lemma-syntomic-flat} 与 \ref{spaces-morphisms-lemma-fppf-open}。
\end{proof}





\section{光滑态射}
\label{spaces-morphisms-section-smooth}

\noindent
概形态射的“光滑”性质在源与靶上都是 \'etale 局部的，见
《下降》评注 \ref{descent-remark-list-local-source-target}。
此性质也在基变换下稳定，并且在靶上是 fpqc 局部的，见《态射》引理
\ref{morphisms-lemma-base-change-smooth} 与《下降》引理
\ref{descent-lemma-descending-property-smooth}。因此，由上面的引理
\ref{spaces-morphisms-lemma-local-source-target}，我们可以如下定义代数空间的光滑态射；
当该态射可表时，这与第 \ref{spaces-morphisms-section-representable} 节中已有的概念一致。

\begin{definition}
\label{spaces-morphisms-definition-smooth}
设 $S$ 为概形。
设 $f : X \to Y$ 为 $S$ 上代数空间的态射。
\begin{enumerate}
\item 若引理 \ref{spaces-morphisms-lemma-local-source-target} 中取
$\mathcal{P} =$“光滑”时其中的等价条件成立，则称 $f$ {\it 光滑}。
\item 设 $x \in |X|$。若存在 $x$ 的开邻域 $X' \subset X$，使得
$f|_{X'} : X' \to Y$ 光滑，则称 $f$ {\it 在 $x$ 处光滑}。
\end{enumerate}
\end{definition}

\begin{lemma}
\label{spaces-morphisms-lemma-composition-smooth}
光滑态射的复合光滑。
\end{lemma}

\begin{proof}
见评注 \ref{spaces-morphisms-remark-composition-P} 和《态射》引理
\ref{morphisms-lemma-composition-smooth}。
\end{proof}

\begin{lemma}
\label{spaces-morphisms-lemma-base-change-smooth}
光滑态射的基变换光滑。
\end{lemma}

\begin{proof}
见评注 \ref{spaces-morphisms-remark-base-change-P} 和《态射》引理
\ref{morphisms-lemma-base-change-smooth}。
\end{proof}

\begin{lemma}
\label{spaces-morphisms-lemma-smooth-local}
设 $S$ 为概形。
设 $f : X \to Y$ 为 $S$ 上代数空间的态射。
下列条件等价：
\begin{enumerate}
\item $f$ 光滑；
\item 对每个 $x \in |X|$，态射 $f$ 在 $x$ 处光滑；
\item 对每个概形 $Z$ 及任意态射 $Z \to Y$，态射
$Z \times_Y X \to Z$ 光滑；
\item 对每个仿射概形 $Z$ 及任意态射 $Z \to Y$，态射
$Z \times_Y X \to Z$ 光滑；
\item 存在概形 $V$ 及满 \'etale 态射 $V \to Y$，使得
$V \times_Y X \to V$ 为光滑态射；
\item 存在概形 $U$ 及满 \'etale 态射 $\varphi : U \to X$，使得复合
$f \circ \varphi$ 光滑；
\item 对每个交换图
$$
\xymatrix{
U \ar[d] \ar[r] & V \ar[d] \\
X \ar[r] & Y
}
$$
其中 $U$、$V$ 为概形，竖直箭头为 \'etale 态射，顶部水平箭头光滑；
\item 存在交换图
$$
\xymatrix{
U \ar[d] \ar[r] & V \ar[d] \\
X \ar[r] & Y
}
$$
其中 $U$、$V$ 为概形，竖直箭头为 \'etale 态射，且 $U \to X$ 满，
使得顶部水平箭头光滑；
\item 存在 Zariski 覆盖 $Y = \bigcup_{i \in I} Y_i$ 以及
$f^{-1}(Y_i) = \bigcup X_{ij}$，使得每个态射 $X_{ij} \to Y_i$ 光滑。
\end{enumerate}
\end{lemma}

\begin{proof}
略。
\end{proof}

\begin{lemma}
\label{spaces-morphisms-lemma-smooth-locally-finite-presentation}
代数空间的光滑态射局部有限表示。
\end{lemma}

\begin{proof}
设 $X \to Y$ 为代数空间的光滑态射。按定义，这意味着存在引理
\ref{spaces-morphisms-lemma-local-source-target} 中的图，其中 $h$ 光滑，且竖直箭头 $a$ 满。
由《态射》引理 \ref{morphisms-lemma-smooth-locally-finite-presentation}，
$h$ 局部有限表示。因此按定义 $X \to Y$ 局部有限表示。
\end{proof}

\begin{lemma}
\label{spaces-morphisms-lemma-smooth-locally-finite-type}
代数空间的光滑态射局部有限型。
\end{lemma}

\begin{proof}
合用引理 \ref{spaces-morphisms-lemma-smooth-locally-finite-presentation} 与
\ref{spaces-morphisms-lemma-finite-presentation-finite-type}。
\end{proof}

\begin{lemma}
\label{spaces-morphisms-lemma-smooth-flat}
代数空间的光滑态射平坦。
\end{lemma}

\begin{proof}
设 $X \to Y$ 为代数空间的光滑态射。按定义，这意味着存在引理
\ref{spaces-morphisms-lemma-local-source-target} 中的图，其中 $h$ 光滑，且竖直箭头 $a$ 满。
由《态射》引理 \ref{morphisms-lemma-smooth-locally-finite-presentation}，
$h$ 平坦。因此按定义 $X \to Y$ 平坦。
\end{proof}

\begin{lemma}
\label{spaces-morphisms-lemma-smooth-syntomic}
代数空间的光滑态射 syntomic。
\end{lemma}

\begin{proof}
设 $X \to Y$ 为代数空间的光滑态射。按定义，这意味着存在引理
\ref{spaces-morphisms-lemma-local-source-target} 中的图，其中 $h$ 光滑，且竖直箭头 $a$ 满。
由《态射》引理 \ref{morphisms-lemma-smooth-syntomic}，$h$ syntomic。
因此按定义 $X \to Y$ syntomic。
\end{proof}

\begin{lemma}
\label{spaces-morphisms-lemma-where-smooth}
设 $S$ 为概形。设 $f : X \to Y$ 为 $S$ 上代数空间的态射。
存在最大的开子空间 $U \subset X$，使得 $f|_U : U \to Y$ 光滑。
此外，形成此开子空间与下列基变换可交换：
\begin{enumerate}
\item 平坦且局部有限表示的态射；
\item 当 $f$ 局部有限表示时的平坦态射。
\end{enumerate}
\end{lemma}

\begin{proof}
$U$ 的存在性由如下事实得出：光滑性质在源上是 Zariski 局部的
（甚至是 \'etale 局部的），见引理 \ref{spaces-morphisms-lemma-smooth-local}。
而且，该引理使我们可以把性质 (1)、(2) 转化为概形态射的情形。
概形情形即《态射》引理
\ref{morphisms-lemma-set-points-where-fibres-smooth}。略去若干细节。
\end{proof}

\begin{lemma}
\label{spaces-morphisms-lemma-smoothness-dimension-spaces}
设 $X$ 与 $Y$ 为概形 $S$ 上局部 Noether 的代数空间，并设
$f : X \to Y$ 为光滑态射。对每个像为 $y \in |Y|$ 的点 $x \in |X|$，
$$
\dim_x(X) = \dim_y(Y) + \dim_x(X_y)
$$
，其中 $\dim_x(X_y)$ 是定义 \ref{spaces-morphisms-definition-dimension-fibre} 中
$f$ 在 $x$ 处的相对维数。
\end{lemma}

\begin{proof}
由代数空间在一点处之维数的定义（《空间的性质》定义
\ref{spaces-properties-definition-dimension-at-point}），问题归结为概形的相应命题
（《态射》引理 \ref{morphisms-lemma-smoothness-dimension}）。
\end{proof}



\section{非分歧态射}
\label{spaces-morphisms-section-unramified}

\noindent
概形态射的“非分歧”性质（相应地，“G-非分歧”性质）在源与靶上都是
\'etale 局部的，见《下降》评注
\ref{descent-remark-list-local-source-target}。此性质也在基变换下稳定，
并且在靶上是 fpqc 局部的，见《态射》引理
\ref{morphisms-lemma-base-change-unramified} 与《下降》引理
\ref{descent-lemma-descending-property-unramified}。因此，由上面的引理
\ref{spaces-morphisms-lemma-local-source-target}，我们可以如下定义代数空间的非分歧态射
（相应地，G-非分歧态射）；当该态射可表时，这与第
\ref{spaces-morphisms-section-representable} 节中已有的概念一致。

\begin{definition}
\label{spaces-morphisms-definition-unramified}
设 $S$ 为概形。
设 $f : X \to Y$ 为 $S$ 上代数空间的态射。
\begin{enumerate}
\item 若引理 \ref{spaces-morphisms-lemma-local-source-target} 中取
$\mathcal{P} = \text{非分歧}$ 时其中的等价条件成立，则称 $f$ {\it 非分歧}。
\item 设 $x \in |X|$。若存在 $x$ 的开邻域 $X' \subset X$，使得
$f|_{X'} : X' \to Y$ 非分歧，则称 $f$ {\it 在 $x$ 处非分歧}。
\item 若引理 \ref{spaces-morphisms-lemma-local-source-target} 中取
$\mathcal{P} = \text{G-非分歧}$ 时其中的等价条件成立，则称 $f$
{\it G-非分歧}。
\item 设 $x \in |X|$。若存在 $x$ 的开邻域 $X' \subset X$，使得
$f|_{X'} : X' \to Y$ G-非分歧，则称 $f$ {\it 在 $x$ 处 G-非分歧}。
\end{enumerate}
\end{definition}

\noindent
由于下一个引理，从此以后我们只发展非分歧态射的理论；需要使用
G-非分歧态射时，我们将直接称其为“局部有限表示的非分歧态射”。

\begin{lemma}
\label{spaces-morphisms-lemma-unramified-G-unramified}
设 $S$ 为概形，$f : X \to Y$ 为 $S$ 上代数空间的态射。
则 $f$ G-非分歧，当且仅当 $f$ 非分歧且局部有限表示。
\end{lemma}

\begin{proof}
考虑引理 \ref{spaces-morphisms-lemma-local-source-target} 中的任意图。所要说的正是：态射
$h$ G-非分歧，当且仅当它非分歧且局部有限表示。这由《态射》定义
\ref{morphisms-definition-unramified} 立即可见。
\end{proof}

\begin{lemma}
\label{spaces-morphisms-lemma-composition-unramified}
非分歧态射的复合非分歧。
\end{lemma}

\begin{proof}
见评注 \ref{spaces-morphisms-remark-composition-P} 和《态射》引理
\ref{morphisms-lemma-composition-unramified}。
\end{proof}

\begin{lemma}
\label{spaces-morphisms-lemma-base-change-unramified}
非分歧态射的基变换非分歧。
\end{lemma}

\begin{proof}
见评注 \ref{spaces-morphisms-remark-base-change-P} 和《态射》引理
\ref{morphisms-lemma-base-change-unramified}。
\end{proof}

\begin{lemma}
\label{spaces-morphisms-lemma-unramified-local}
设 $S$ 为概形。
设 $f : X \to Y$ 为 $S$ 上代数空间的态射。
下列条件等价：
\begin{enumerate}
\item $f$ 非分歧；
\item 对每个 $x \in |X|$，态射 $f$ 在 $x$ 处非分歧；
\item 对每个概形 $Z$ 及任意态射 $Z \to Y$，态射
$Z \times_Y X \to Z$ 非分歧；
\item 对每个仿射概形 $Z$ 及任意态射 $Z \to Y$，态射
$Z \times_Y X \to Z$ 非分歧；
\item 存在概形 $V$ 及满 \'etale 态射 $V \to Y$，使得
$V \times_Y X \to V$ 为非分歧态射；
\item 存在概形 $U$ 及满 \'etale 态射 $\varphi : U \to X$，使得复合
$f \circ \varphi$ 非分歧；
\item 对每个交换图
$$
\xymatrix{
U \ar[d] \ar[r] & V \ar[d] \\
X \ar[r] & Y
}
$$
其中 $U$、$V$ 为概形，竖直箭头为 \'etale 态射，顶部水平箭头非分歧；
\item 存在交换图
$$
\xymatrix{
U \ar[d] \ar[r] & V \ar[d] \\
X \ar[r] & Y
}
$$
其中 $U$、$V$ 为概形，竖直箭头为 \'etale 态射，且 $U \to X$ 满，
使得顶部水平箭头非分歧；
\item 存在 Zariski 覆盖 $Y = \bigcup_{i \in I} Y_i$ 以及
$f^{-1}(Y_i) = \bigcup X_{ij}$，使得每个态射
$X_{ij} \to Y_i$ 非分歧。
\end{enumerate}
\end{lemma}

\begin{proof}
略。
\end{proof}


\begin{lemma}
\label{spaces-morphisms-lemma-unramified-locally-finite-type}
代数空间的非分歧态射局部有限型。
\end{lemma}

\begin{proof}
借助引理 \ref{spaces-morphisms-lemma-local-source-target} 中的图，这转化为《态射》引理
\ref{morphisms-lemma-unramified-locally-finite-type}。
\end{proof}

\begin{lemma}
\label{spaces-morphisms-lemma-unramified-quasi-finite}
若 $f$ 在 $x$ 处非分歧，则 $f$ 在 $x$ 处拟有限。
特别地，非分歧态射局部拟有限。
\end{lemma}

\begin{proof}
借助引理 \ref{spaces-morphisms-lemma-local-source-target} 中的图，这转化为《态射》引理
\ref{morphisms-lemma-unramified-quasi-finite}。
\end{proof}

\begin{lemma}
\label{spaces-morphisms-lemma-immersion-unramified}
代数空间的浸入非分歧。
\end{lemma}

\begin{proof}
设 $i : X \to Y$ 为代数空间的浸入。取概形 $V$ 及满 \'etale 态射
$V \to Y$。则 $V \times_Y X \to V$ 为概形的浸入，因而非分歧
（见《态射》引理 \ref{morphisms-lemma-open-immersion-unramified} 与
\ref{morphisms-lemma-closed-immersion-unramified}）。故按定义，$i$ 非分歧。
\end{proof}

\begin{lemma}
\label{spaces-morphisms-lemma-diagonal-unramified-morphism}
设 $S$ 为概形。
设 $f : X \to Y$ 为 $S$ 上代数空间的态射。
\begin{enumerate}
\item 若 $f$ 非分歧，则对角态射
$\Delta_{X/Y} : X \to X \times_Y X$ 为开浸入。
\item 若 $f$ 局部有限型且 $\Delta_{X/Y}$ 为开浸入，则 $f$ 非分歧。
\end{enumerate}
\end{lemma}

\begin{proof}
无论如何，$\Delta_{X/Y}$ 都是可表单态射，见引理
\ref{spaces-morphisms-lemma-properties-diagonal}。取概形 $V$ 及满 \'etale 态射 $V \to Y$。
再取概形 $U$ 及满 \'etale 态射 $U \to X \times_Y V$。考虑交换图
$$
\xymatrix{
U \ar[d] \ar[rr]_-{\Delta_{U/V}} & &
U \times_V U \ar[d] \ar[r] &
V \ar[d]^{\Delta_{V/Y}} \\
X \ar[rr]^-{\Delta_{X/Y}} & &
X \times_Y X \ar[r] &
V \times_Y V
}
$$
其右方形为笛卡尔方块。左侧竖直箭头满且 \'etale。右侧竖直箭头
作为 $Y$ 上 \'etale 概形之间的态射是 \'etale 的，见《空间的性质》引理
\ref{spaces-properties-lemma-etale-permanence}。因此中间的竖直箭头也
\'etale（但未必满）。

\medskip\noindent
设 $f$ 非分歧。则 $U \to V$ 非分歧，故由《态射》引理
\ref{morphisms-lemma-diagonal-unramified-morphism}，$\Delta_{U/V}$ 为开浸入。
考察上图左方形，由《空间的性质》引理
\ref{spaces-properties-lemma-etale-local} 可知 $\Delta_{X/Y}$ 是 \'etale 态射。
因此 $\Delta_{X/Y}$ 是可表的 \'etale 单态射；由《\'Etale 态射》定理
\ref{etale-theorem-etale-radicial-open}，它是开浸入。（关于从概形语言到
函子语言的转译，另见《空间》引理
\ref{spaces-lemma-representable-transformations-property-implication}。）

\medskip\noindent
设 $f$ 局部有限型，且 $\Delta_{X/Y}$ 为开浸入。由代数空间的态射局部
有限型之定义，$U \to V$ 也局部有限型。考察上面的图，由《空间的性质》
引理 \ref{spaces-properties-lemma-etale-permanence}，可知 $\Delta_{U/V}$
作为 $X \times_Y X$ 上 \'etale 概形之间的态射是 \'etale 的。另一方面，
$\Delta_{U/V}$ 是概形之间态射的对角态射，因而无论如何都是浸入，见
《概形》引理 \ref{schemes-lemma-diagonal-immersion}。故它是开浸入，进而由
《态射》引理 \ref{morphisms-lemma-diagonal-unramified-morphism}，
$U \to V$ 非分歧。这又按定义说明 $f$ 非分歧。
\end{proof}

\begin{lemma}
\label{spaces-morphisms-lemma-where-unramified}
设 $S$ 为概形。考虑交换图
$$
\xymatrix{
X \ar[rr]_f \ar[rd]_p & & Y \ar[ld]^q \\
& Z
}
$$
，其中各对象为 $S$ 上的代数空间。设 $X \to Z$ 局部有限型。则存在开子空间
$U(f) \subset X$，使得 $|U(f)| \subset |X|$ 正是 $f$ 非分歧处的点集。
此外，对代数空间的任意态射 $Z' \to Z$，若 $f' : X' \to Y'$ 是 $f$
沿 $Z' \to Z$ 的基变换，则 $U(f')$ 是 $U(f)$ 在投影 $X' \to X$ 下的逆像。
\end{lemma}

\begin{proof}
本引理是《态射》引理
\ref{morphisms-lemma-set-points-where-fibres-unramified} 的类似结果；事实上，
我们将由它推出本引理。由定义 \ref{spaces-morphisms-definition-unramified}，集合
$\{x \in |X| : f \text{ 非分歧于 }x\}$ 在 $X$ 中为开集。因此只需证明最后的
断言。由引理 \ref{spaces-morphisms-lemma-permanence-finite-type}，态射 $X \to Y$ 局部有限型；
由引理 \ref{spaces-morphisms-lemma-base-change-finite-type}，态射 $X' \to Y'$ 局部有限型。

\medskip\noindent
取概形 $W$ 及满 \'etale 态射 $W \to Z$。取概形 $V$ 及满 \'etale 态射
$V \to W \times_Z Y$。取概形 $U$ 及满 \'etale 态射
$U \to V \times_Y X$。最后，取概形 $W'$ 及满 \'etale 态射
$W' \to W \times_Z Z'$。令 $V' = W' \times_W V$、
$U' = W' \times_W U$；于是得到满 \'etale 态射 $V' \to Y'$ 与
$U' \to X'$。以下将不再另行说明地使用如下事实：代数空间的 \'etale 态射在
伴随拓扑空间上诱导开映射（见《空间的性质》引理
\ref{spaces-properties-lemma-etale-open}）。结合引理
\ref{spaces-morphisms-lemma-unramified-local}，可知 $U(f)$ 是 $U$ 中使态射 $U \to V$
非分歧的点集 $T$ 在 $|X|$ 中的像。类似地，$U(f')$ 是 $U'$ 中使态射
$U' \to V'$ 非分歧的点集 $T'$ 在 $|X'|$ 中的像。按构造，图
$$
\xymatrix{
U' \ar[r] \ar[d] & U \ar[d] \\
V' \ar[r] & V
}
$$
在概形范畴中为笛卡尔图。因此，上述《态射》引理
\ref{morphisms-lemma-set-points-where-fibres-unramified} 表明 $T'$ 是 $T$
的逆像。由于 $|U'| \to |X'|$ 满，这就证明了本引理。
\end{proof}

\begin{lemma}
\label{spaces-morphisms-lemma-permanence-unramified}
设 $S$ 为概形，$X \to Y \to Z$ 为 $S$ 上代数空间的态射。
若 $X \to Z$ 非分歧，则 $X \to Y$ 非分歧。
\end{lemma}

\begin{proof}
取交换图
$$
\xymatrix{
U \ar[d] \ar[r] & V \ar[d] \ar[r] & W \ar[d] \\
X \ar[r] & Y \ar[r] & Z
}
$$
，其中竖直箭头均 \'etale 且满（见《空间》引理
\ref{spaces-lemma-lift-morphism-presentations}）。对顶行应用《态射》引理
\ref{morphisms-lemma-unramified-permanence}。
\end{proof}














\section{\'Etale 态射}
\label{spaces-morphisms-section-etale}

\noindent
代数空间的 \'etale 态射这一概念已在《空间的性质》定义
\ref{spaces-properties-definition-etale} 中定义。下面说明态射在一点处
\'etale 的含义。

\begin{definition}
\label{spaces-morphisms-definition-etale}
设 $S$ 为概形，$f : X \to Y$ 为 $S$ 上代数空间的态射，并设
$x \in |X|$。若存在 $x$ 的开邻域 $X' \subset X$，使得
$f|_{X'} : X' \to Y$ 为 \'etale 态射，则称 $f$ {\it 在 $x$ 处 \'etale}。
\end{definition}

\begin{lemma}
\label{spaces-morphisms-lemma-etale-local}
设 $S$ 为概形。
设 $f : X \to Y$ 为 $S$ 上代数空间的态射。
下列条件等价：
\begin{enumerate}
\item $f$ 为 \'etale 态射；
\item 对每个 $x \in |X|$，态射 $f$ 在 $x$ 处 \'etale；
\item 对每个概形 $Z$ 及任意态射 $Z \to Y$，态射
$Z \times_Y X \to Z$ 为 \'etale 态射；
\item 对每个仿射概形 $Z$ 及任意态射 $Z \to Y$，态射
$Z \times_Y X \to Z$ 为 \'etale 态射；
\item 存在概形 $V$ 及满 \'etale 态射 $V \to Y$，使得
$V \times_Y X \to V$ 为 \'etale 态射；
\item 存在概形 $U$ 及满 \'etale 态射 $\varphi : U \to X$，使得复合
$f \circ \varphi$ 为 \'etale 态射；
\item 对每个交换图
$$
\xymatrix{
U \ar[d] \ar[r] & V \ar[d] \\
X \ar[r] & Y
}
$$
其中 $U$、$V$ 为概形，竖直箭头为 \'etale 态射，顶部水平箭头也为
\'etale 态射；
\item 存在交换图
$$
\xymatrix{
U \ar[d] \ar[r] & V \ar[d] \\
X \ar[r] & Y
}
$$
其中 $U$、$V$ 为概形，竖直箭头为 \'etale 态射，且 $U \to X$ 满，
使得顶部水平箭头为 \'etale 态射；
\item 存在 Zariski 覆盖 $Y = \bigcup Y_i$ 以及
$f^{-1}(Y_i) = \bigcup X_{ij}$，使得每个态射
$X_{ij} \to Y_i$ 为 \'etale 态射。
\end{enumerate}
\end{lemma}

\begin{proof}
合用《空间的性质》引理
\ref{spaces-properties-lemma-etale-local}、
\ref{spaces-properties-lemma-base-change-etale} 与
\ref{spaces-properties-lemma-composition-etale}。略去若干细节。
\end{proof}

\begin{lemma}
\label{spaces-morphisms-lemma-composition-etale}
代数空间的两个 \'etale 态射之复合仍为 \'etale 态射。
\end{lemma}

\begin{proof}
这是《空间的性质》引理
\ref{spaces-properties-lemma-composition-etale} 的复述。
\end{proof}

\begin{lemma}
\label{spaces-morphisms-lemma-base-change-etale}
代数空间的 \'etale 态射沿代数空间的任意态射作基变换后仍为 \'etale 态射。
\end{lemma}

\begin{proof}
这是《空间的性质》引理
\ref{spaces-properties-lemma-base-change-etale} 的复述。
\end{proof}

\begin{lemma}
\label{spaces-morphisms-lemma-etale-locally-quasi-finite}
代数空间的 \'etale 态射局部拟有限。
\end{lemma}

\begin{proof}
设 $X \to Y$ 为代数空间的 \'etale 态射，见《空间的性质》定义
\ref{spaces-properties-definition-etale}。由《空间的性质》引理
\ref{spaces-properties-lemma-etale-local}，这意味着存在引理
\ref{spaces-morphisms-lemma-local-source-target} 中的图，其中 $h$ 为 \'etale 态射，且竖直箭头
$a$ 满。由《态射》引理 \ref{morphisms-lemma-etale-locally-quasi-finite}，
$h$ 局部拟有限。因此按定义 $X \to Y$ 局部拟有限。
\end{proof}

\begin{lemma}
\label{spaces-morphisms-lemma-etale-smooth}
代数空间的 \'etale 态射光滑。
\end{lemma}

\begin{proof}
证明与引理 \ref{spaces-morphisms-lemma-etale-locally-quasi-finite} 的证明相同。这里使用如下事实：
概形的 \'etale 态射光滑（由概形的 \'etale 态射之定义）。
\end{proof}

\begin{lemma}
\label{spaces-morphisms-lemma-etale-flat}
代数空间的 \'etale 态射平坦。
\end{lemma}

\begin{proof}
证明与引理 \ref{spaces-morphisms-lemma-etale-locally-quasi-finite} 的证明相同。这里使用
《态射》引理 \ref{morphisms-lemma-etale-flat}。
\end{proof}

\begin{lemma}
\label{spaces-morphisms-lemma-etale-locally-finite-presentation}
\begin{slogan}
\'Etale 蕴含局部有限表示。
\end{slogan}
代数空间的 \'etale 态射局部有限表示。
\end{lemma}

\begin{proof}
证明与引理 \ref{spaces-morphisms-lemma-etale-locally-quasi-finite} 的证明相同。这里使用
《态射》引理 \ref{morphisms-lemma-etale-locally-finite-presentation}。
\end{proof}

\begin{lemma}
\label{spaces-morphisms-lemma-etale-locally-finite-type}
代数空间的 \'etale 态射局部有限型。
\end{lemma}

\begin{proof}
\'Etale 态射局部有限表示，而局部有限表示的态射局部有限型；见引理
\ref{spaces-morphisms-lemma-etale-locally-finite-presentation} 与
\ref{spaces-morphisms-lemma-finite-presentation-finite-type}。
\end{proof}

\begin{lemma}
\label{spaces-morphisms-lemma-etale-unramified}
代数空间的 \'etale 态射非分歧。
\end{lemma}

\begin{proof}
证明与引理 \ref{spaces-morphisms-lemma-etale-locally-quasi-finite} 的证明相同。这里使用
《态射》引理 \ref{morphisms-lemma-etale-smooth-unramified}。
\end{proof}

\begin{lemma}
\label{spaces-morphisms-lemma-etale-permanence}
设 $S$ 为概形。设 $X,Y$ 为在某代数空间 $Z$ 上 \'etale 的代数空间。
则 $Z$ 上的任意态射 $X \to Y$ 都是 \'etale 态射。
\end{lemma}

\begin{proof}
这是《空间的性质》引理
\ref{spaces-properties-lemma-etale-permanence} 的复述。
\end{proof}

\begin{lemma}
\label{spaces-morphisms-lemma-unramified-flat-lfp-etale}
代数空间之间局部有限表示、平坦且非分歧的态射是 \'etale 态射。
\end{lemma}

\begin{proof}
设 $X \to Y$ 为代数空间之间局部有限表示、平坦且非分歧的态射。由
《空间的性质》引理 \ref{spaces-properties-lemma-etale-local}，这意味着存在
引理 \ref{spaces-morphisms-lemma-local-source-target} 中的图，其中 $h$ 局部有限表示、平坦且
非分歧，且竖直箭头 $a$ 满。由《态射》引理
\ref{morphisms-lemma-flat-unramified-etale}，$h$ 为 \'etale 态射。因此按定义，
$X \to Y$ 为 \'etale 态射。
\end{proof}






\section{固有态射}
\label{spaces-morphisms-section-proper}

\noindent
固有态射的概念在代数几何中起着重要作用。下面给出代数空间的固有态射之定义。

\begin{definition}
\label{spaces-morphisms-definition-proper}
设 $S$ 为概形，$f : X \to Y$ 为 $S$ 上代数空间的态射。
若 $f$ 分离、有限型且普遍闭，则称 $f$ {\it 固有}。
\end{definition}

\begin{lemma}
\label{spaces-morphisms-lemma-proper-local}
设 $S$ 为概形，$f : X \to Y$ 为 $S$ 上代数空间的态射。下列条件等价：
\begin{enumerate}
\item $f$ 固有；
\item 对每个概形 $Z$ 及每个态射 $Z \to Y$，投影
$Z \times_Y X \to Z$ 固有；
\item 对每个仿射概形 $Z$ 及每个态射 $Z \to Y$，投影
$Z \times_Y X \to Z$ 固有；
\item 存在概形 $V$ 及满 \'etale 态射 $V \to Y$，使得
$V \times_Y X \to V$ 固有；
\item 存在 Zariski 覆盖 $Y = \bigcup Y_i$，使得每个态射
$f^{-1}(Y_i) \to Y_i$ 固有。
\end{enumerate}
\end{lemma}

\begin{proof}
合用引理 \ref{spaces-morphisms-lemma-separated-local}、
\ref{spaces-morphisms-lemma-finite-type-local},
\ref{spaces-morphisms-lemma-quasi-compact-local} 与
\ref{spaces-morphisms-lemma-universally-closed-local}。
\end{proof}

\begin{lemma}
\label{spaces-morphisms-lemma-base-change-proper}
固有态射的基变换固有。
\end{lemma}

\begin{proof}
见引理 \ref{spaces-morphisms-lemma-base-change-separated}、
\ref{spaces-morphisms-lemma-base-change-finite-type} 与
\ref{spaces-morphisms-lemma-base-change-universally-closed}。
\end{proof}

\begin{lemma}
\label{spaces-morphisms-lemma-composition-proper}
固有态射的复合固有。
\end{lemma}

\begin{proof}
见引理 \ref{spaces-morphisms-lemma-composition-separated}、
\ref{spaces-morphisms-lemma-composition-finite-type} 与
\ref{spaces-morphisms-lemma-composition-universally-closed}。
\end{proof}

\begin{lemma}
\label{spaces-morphisms-lemma-closed-immersion-proper}
代数空间的闭浸入是代数空间的固有态射。
\end{lemma}

\begin{proof}
闭浸入按定义可表，故这由《空间》引理
\ref{spaces-lemma-representable-transformations-property-implication}
及概形态射的相应结果得出；见《态射》引理
\ref{morphisms-lemma-closed-immersion-proper}。
\end{proof}

\begin{lemma}
\label{spaces-morphisms-lemma-universally-closed-permanence}
设 $S$ 为概形。考虑代数空间的交换图
$$
\xymatrix{
X \ar[rr] \ar[rd] & &
Y \ar[ld] \\
& B &
}
$$
，其中各对象在 $S$ 上。
\begin{enumerate}
\item 若 $X \to B$ 普遍闭且 $Y \to B$ 分离，则态射 $X \to Y$ 普遍闭。
特别地，$|X|$ 在 $|Y|$ 中的像为闭集。
\item 若 $X \to B$ 固有且 $Y \to B$ 分离，则态射 $X \to Y$ 固有。
\end{enumerate}
\end{lemma}

\begin{proof}
设 $X \to B$ 普遍闭且 $Y \to B$ 分离。把该态射分解为
$X \to X \times_B Y \to Y$。第一个态射是闭浸入，见引理
\ref{spaces-morphisms-lemma-semi-diagonal}，故普遍闭。投影 $X \times_B Y \to Y$ 是普遍闭态射的
基变换，因而普遍闭，见引理 \ref{spaces-morphisms-lemma-base-change-universally-closed}。
因此，$X \to Y$ 作为普遍闭态射的复合而普遍闭，见引理
\ref{spaces-morphisms-lemma-composition-universally-closed}。这证明了 (1)。为推出 (2)，合用 (1)
与引理 \ref{spaces-morphisms-lemma-compose-after-separated}、
\ref{spaces-morphisms-lemma-quasi-compact-permanence} 及
\ref{spaces-morphisms-lemma-permanence-finite-type}。
\end{proof}

\begin{lemma}
\label{spaces-morphisms-lemma-image-proper-is-proper}
设 $S$ 为概形，$B$ 为 $S$ 上的代数空间，并设
$f : X \to Y$ 为 $B$ 上代数空间的态射。若 $X$ 在 $B$ 上普遍闭且 $f$ 满，
则 $Y$ 在 $B$ 上普遍闭。特别地，若 $Y$ 还在 $B$ 上分离且有限型，
则 $Y$ 在 $B$ 上固有。
\end{lemma}

\begin{proof}
设 $X$ 普遍闭且 $f$ 满。以 $p : X \to B$、$q : Y \to B$ 表示结构态射。
设 $B' \to B$ 为 $S$ 上代数空间的态射。基变换
$f' : X_{B'} \to Y_{B'}$ 满（引理 \ref{spaces-morphisms-lemma-base-change-surjective}），
而基变换 $p' : X_{B'} \to B'$ 闭。若 $T \subset Y_{B'}$ 为闭集，
则 $(f')^{-1}(T) \subset X_{B'}$ 为闭集，故
$p'((f')^{-1}(T)) = q'(T)$ 为闭集。因此 $q'$ 闭。
\end{proof}

\begin{lemma}
\label{spaces-morphisms-lemma-scheme-theoretic-image-is-proper}
设 $S$ 为概形。设
$$
\xymatrix{
X \ar[rr]_h \ar[rd]_f & & Y \ar[ld]^g \\
& B
}
$$
为 $S$ 上代数空间态射的交换图。设
\begin{enumerate}
\item $X \to B$ 为固有态射；
\item $Y \to B$ 分离且局部有限型。
\end{enumerate}
则 $h$ 的概形论像 $Z \subset Y$ 在 $B$ 上固有，且 $X \to Z$ 满。
\end{lemma}

\begin{proof}
$h$ 的概形论像已在第 \ref{spaces-morphisms-section-scheme-theoretic-image} 节中构造。
注意，$h$ 拟紧（引理 \ref{spaces-morphisms-lemma-quasi-compact-quasi-separated-permanence}），
故 $|h|(|X|) \subset |Z|$ 稠密（引理
\ref{spaces-morphisms-lemma-quasi-compact-scheme-theoretic-image}）。另一方面，
$|h|(|X|)$ 在 $|Y|$ 中为闭集（引理
\ref{spaces-morphisms-lemma-universally-closed-permanence}），故 $X \to Z$ 满。
于是 $Z \to B$ 固有（引理 \ref{spaces-morphisms-lemma-image-proper-is-proper}）。
\end{proof}

\begin{lemma}
\label{spaces-morphisms-lemma-separated-diagonal-proper}
设 $S$ 为概形。
设 $f : X \to Y$ 为 $S$ 上代数空间的态射。
下列条件等价：
\begin{enumerate}
\item $f$ 分离；
\item $\Delta_{X/Y} : X \to X \times_Y X$ 普遍闭；
\item $\Delta_{X/Y} : X \to X \times_Y X$ 固有。
\end{enumerate}
\end{lemma}

\begin{proof}
(1) $\Rightarrow$ (3) 由引理 \ref{spaces-morphisms-lemma-closed-immersion-proper} 得出。
在余下的证明中，我们将不再另行说明地使用《空间》引理
\ref{spaces-lemma-representable-transformations-property-implication}。
回顾 $\Delta_{X/Y}$ 是局部有限型的可表单态射，见引理
\ref{spaces-morphisms-lemma-properties-diagonal}。由于概形态射中固有 $\Rightarrow$ 普遍闭，
故 (3) 蕴含 (2)。若 $\Delta_{X/Y}$ 普遍闭，则《\'Etale 态射》引理
\ref{etale-lemma-characterize-closed-immersion} 表明它是闭浸入。
因此 (2) $\Rightarrow$ (1)，证明完毕。
\end{proof}





\section{赋值判据}
\label{spaces-morphisms-section-valuative}

\noindent
本节介绍代数空间态射之赋值判据的基础。进一步结果可参见：
\begin{enumerate}
\item 普遍闭性的赋值判据见第
\ref{spaces-morphisms-section-valuative-criterion-universally-closed} 节；
\item 分离性的赋值判据见第 \ref{spaces-morphisms-section-valuative-separatedness} 节；
\item 固有性的赋值判据见第 \ref{spaces-morphisms-section-valuative-criterion-properness} 节；
\item 其他逆命题见《良好空间》第
\ref{decent-spaces-section-valuative-criterion-universally-closed}
节及《良好空间》引理
\ref{decent-spaces-lemma-re-characterize-universally-closed}；
\item 在 Noether 情形，只需对离散赋值环检查该判据；见《空间的上同调》第
\ref{spaces-cohomology-section-Noetherian-valuative-criterion} 节及
《空间的极限》第 \ref{spaces-limits-section-Noetherian-valuative-criterion} 节；
\item Noether 情形下赋值判据的精细版本见《空间的极限》第
\ref{spaces-limits-section-refined-valuative-criteria}.
\end{enumerate}
我们先正式陈述定义，再讨论它与概形态射情形的差别。

\begin{definition}
\label{spaces-morphisms-definition-valuative-criterion}
设 $S$ 为概形，$f : X \to Y$ 为 $S$ 上代数空间的态射。若对任意由实线箭头
构成的交换图
$$
\xymatrix{
\Spec(K) \ar[r] \ar[d] & X \ar[d] \\
\Spec(A) \ar[r] \ar@{-->}[ru] & Y
}
$$
，其中 $A$ 为分式域是 $K$ 的赋值环，虚线箭头至多存在一个（不要求存在），
则称 $f$ {\it 满足赋值判据的唯一性部分}。若对上述任意实线图，存在域扩张
$K'/K$、支配 $A$ 的赋值环 $A' \subset K'$ 以及态射
$\Spec(A') \to X$，使得下图交换
$$
\xymatrix{
\Spec(K') \ar[r] \ar[d] & \Spec(K) \ar[r] & X \ar[d] \\
\Spec(A') \ar[r] \ar[rru] & \Spec(A) \ar[r] & Y
}
$$
，则称 $f$ {\it 满足赋值判据的存在性部分}。若 $f$ 同时满足存在性部分与
唯一性部分，则称 $f$ {\it 满足赋值判据}。
\end{definition}

\noindent
对代数空间的态射，赋值判据之存在性部分的表述略有不同，因为可能需要扩张
赋值环的分式域。实践中，这一区别几乎从不起作用。
\begin{enumerate}
\item 检查赋值判据的唯一性部分从不涉及分式域扩张，故与概形情形完全相同。
\item 一般而言必须允许域扩张，见例 \ref{spaces-morphisms-example-finite-separable-needed}。
\item 对代数空间的态射，在赋值判据的存在性部分中总可以取有限可分扩张
$K'/K$，见引理 \ref{spaces-morphisms-lemma-finite-separable-enough}。
\item 若 $f : X \to Y$ 是代数空间的分离态射，则检查赋值判据的存在性部分时
总可取 $K = K'$，见引理 \ref{spaces-morphisms-lemma-usual-enough}。
\item 对拟紧拟分离态射 $f : X \to Y$，“$f$ 分离且普遍闭”等价于“$f$
满足通常的赋值判据”，见引理
\ref{spaces-morphisms-lemma-characterize-separated-and-universally-closed}。固有性的赋值判据就是
通常的判据，见引理 \ref{spaces-morphisms-lemma-characterize-proper}。
\end{enumerate}
作为理论的第一步，我们证明：当代数空间的态射可表时，该判据与概形态射所用
的表述相同。

\begin{lemma}
\label{spaces-morphisms-lemma-valuative-criterion-representable}
设 $S$ 为概形，$f : X \to Y$ 为 $S$ 上代数空间的态射，并设 $f$ 可表。
下列条件等价：
\begin{enumerate}
\item $f$ 满足定义 \ref{spaces-morphisms-definition-valuative-criterion} 中赋值判据的存在性部分；
\item 对任意由实线箭头构成的交换图
$$
\xymatrix{
\Spec(K) \ar[r] \ar[d] & X \ar[d] \\
\Spec(A) \ar[r] \ar@{-->}[ru] & Y
}
$$
，其中 $A$ 为分式域是 $K$ 的赋值环，都存在虚线箭头；换言之，$f$ 满足
《概形》定义 \ref{schemes-definition-valuative-criterion} 中赋值判据的存在性部分。
\end{enumerate}
\end{lemma}

\begin{proof}
只需证明：给定定义 \ref{spaces-morphisms-definition-valuative-criterion} 中形如
$$
\xymatrix{
\Spec(K') \ar[r] \ar[d] & \Spec(K) \ar[r] & X \ar[d] \\
\Spec(A') \ar[r] \ar[rru]^\varphi & \Spec(A) \ar[r] & Y
}
$$
的交换图后，也能找到可嵌入此图的态射 $\Spec(A) \to X$。
% P08-E243: frozen source fibre-product formula preserved; see ERRATA_CANDIDATES.jsonl.
令 $X_A = \Spec(A) \times_Y Y$。由于 $f$ 可表，$X_A$ 是概形。态射 $\varphi$
给出态射 $\varphi' : \Spec(A') \to X_A$。设 $x \in X_A$ 为
$\varphi' : \Spec(A') \to X_A$ 的闭点之像。于是有环的交换图
$$
\xymatrix{
K' & K \ar[l] & \mathcal{O}_{X_A, x} \ar[l] \ar[lld] \\
A' \ar[u] & A \ar[l] & A \ar[l] \ar[u]
}
$$
由于 $A$ 为赋值环且 $A'$ 支配 $A$，有 $K \cap A' = A$。因此环映射
$\mathcal{O}_{X_A, x} \to K$ 的像包含在 $A$ 中。于是得到所需态射
$\Spec(A) \to X_A$（见《概形》第 \ref{schemes-section-points} 节）。
\end{proof}

\begin{lemma}
\label{spaces-morphisms-lemma-finite-separable-enough}
设 $S$ 为概形，$f : X \to Y$ 为 $S$ 上代数空间的态射。下列条件等价：
\begin{enumerate}
\item $f$ 满足定义 \ref{spaces-morphisms-definition-valuative-criterion} 中赋值判据的存在性部分；
\item $f$ 满足定义 \ref{spaces-morphisms-definition-valuative-criterion} 中赋值判据的存在性部分，
但将其修改为要求扩张 $K'/K$ 有限可分。
\end{enumerate}
\end{lemma}

\begin{proof}
只需证明 (1) 蕴含 (2)。设给定定义 \ref{spaces-morphisms-definition-valuative-criterion} 中的图
$$
\xymatrix{
\Spec(K') \ar[r] \ar[d] & \Spec(K) \ar[r] & X \ar[d] \\
\Spec(A') \ar[r] \ar[rru] & \Spec(A) \ar[r] & Y
}
$$
，其中 $K \subset K'$ 任意。取概形 $U$ 及满 \'etale 态射 $U \to X$。于是
$$
\Spec(A') \times_X U \longrightarrow \Spec(A')
$$
满且 \'etale。令 $p$ 为 $\Spec(A') \times_X U$ 中映至 $\Spec(A')$ 闭点的点。
令 $p' \leadsto p$ 为 $p$ 的一个泛化，并使其映至 $\Spec(A')$ 的一般点。
这样的泛化存在，因为泛化可沿概形的平坦态射提升；见《态射》引理
\ref{morphisms-lemma-generalizations-lift-flat}。于是 $p'$ 对应于概形
$\Spec(K') \times_X U$ 的一个点。注意
$$
\Spec(K') \times_X U
=
\Spec(K') \times_{\Spec(K)} (\Spec(K) \times_X U)
$$
，故 $p'$ 映至一点 $q' \in \Spec(K) \times_X U$，其剩余域是 $K$ 的有限可分
扩张。最后，$p' \leadsto p$ 映至概形 $U$ 上的特化 $u' \leadsto u$。
用上述记号，得到环的图
$$
\xymatrix{
\kappa(p') & & \kappa(q') \ar[ll] & \kappa(u') \ar[l] \\
& \mathcal{O}_{\Spec(A') \times_X U, p} \ar[lu] & &
\mathcal{O}_{U, u} \ar[ll] \ar[u] \\
K' \ar[uu] & A' \ar[l] \ar[u] & A \ar[l] \ar'[u][uu]
}
$$
这意味着，由 $A$ 与 $\mathcal{O}_{U,u}$ 的像生成的环
$B \subset \kappa(q')$ 映至 $\kappa(p')$ 的一个子环，而该子环包含在
$\mathcal{O}_{\Spec(A') \times_X U,p} \to \kappa(p')$ 的像 $B'$ 中。
注意 $B'$ 是局部环。
% P08-E244: frozen source ideal containment preserved; see ERRATA_CANDIDATES.jsonl.
令 $\mathfrak m \subset B$ 为极大理想。按构造，$A \cap \mathfrak m$
（相应地，$\mathcal{O}_{U,u} \cap \mathfrak m$、$A' \cap \mathfrak m$）
是 $A$（相应地，$\mathcal{O}_{U,u}$、$A'$）的极大理想。令
$\mathfrak q = B \cap \mathfrak m$。这是素理想，且 $A \cap \mathfrak q$
是 $A$ 的极大理想。因此 $B_{\mathfrak q} \subset \kappa(q')$ 是支配 $A$
的局部环。由《代数》引理 \ref{algebra-lemma-dominate}，可找到分式域为
$\kappa(q')$ 且支配 $B_{\mathfrak q}$ 的赋值环
$A_1 \subset \kappa(q')$。局部环映射 $\mathcal{O}_{U,u} \to A_1$ 给出态射
$\Spec(A_1) \to U \to X$
，使得图
$$
\xymatrix{
\Spec(\kappa(q')) \ar[r] \ar[d] & \Spec(K) \ar[r] & X \ar[d] \\
\Spec(A_1) \ar[r] \ar[rru] & \Spec(A) \ar[r] & Y
}
$$
交换。由于 $A_1$ 的分式域为 $\kappa(q')$，且按构造
$\kappa(q')/K$ 有限可分，本引理得证。
\end{proof}

\begin{lemma}
\label{spaces-morphisms-lemma-push-down-solution}
设 $S$ 为概形，$f : X \to Y$ 为 $S$ 上代数空间的分离态射。设给定定义
\ref{spaces-morphisms-definition-valuative-criterion} 中的图
$$
\xymatrix{
\Spec(K') \ar[r] \ar[d] & \Spec(K) \ar[r] & X \ar[d] \\
\Spec(A') \ar[r] \ar[rru] & \Spec(A) \ar[r] \ar@{-->}[ru] & Y
}
$$
，其中 $K \subset K'$ 任意。则存在使该图交换的虚线箭头。
\end{lemma}

\begin{proof}
我们必须证明能找到嵌入此图的态射 $\Spec(A) \to X$。

\medskip\noindent
考虑 $X$ 的基变换 $X_A = \Spec(A) \times_Y X$。则
$X_A \to \Spec(A)$ 是代数空间的分离态射（引理
\ref{spaces-morphisms-lemma-base-change-separated}）。对上图的所有态射作基变换，得到
$$
\xymatrix{
\Spec(K') \ar[r] \ar[d] & \Spec(K) \ar[r] & X_A \ar[d] \\
\Spec(A') \ar[r] \ar[rru] & \Spec(A) \ar@{=}[r] & \Spec(A)
}
$$
因此可以用 $X_A$ 替换 $X$，设 $Y = \Spec(A)$，并设有上图。我们还可以且
确实用一个包含 $|\Spec(A')| \to |X|$ 之像的拟紧开子空间替换 $X$。

\medskip\noindent
由引理 \ref{spaces-morphisms-lemma-quasi-compact-permanence}，态射 $\Spec(A') \to X$ 拟紧。
令 $Z \subset X$ 为 $\Spec(A') \to X$ 的概形论像。则 $Z$ 是 $A$ 上的既约
（引理 \ref{spaces-morphisms-lemma-scheme-theoretic-image-reduced}）、拟紧（因为是 $X$ 的闭
子空间）且分离（同样因为是 $X$ 的闭子空间）的代数空间。考虑态射
$\Spec(A') \to X$ 沿概形的平坦态射 $\Spec(K) \to \Spec(A)$ 的基变换
$$
\Spec(K') = \Spec(A') \times_{\Spec(A)} \Spec(K) \to
X \times_{\Spec(A)} \Spec(K) = X_K
$$
。由引理 \ref{spaces-morphisms-lemma-flat-base-change-scheme-theoretic-image}，此态射的概形论像
是 $Z$ 的基变换 $Z_K$。另一方面，由假设（即上面的交换图），此态射分解经过
态射 $\Spec(K) \to Z_K$，而后者是结构态射 $Z_K \to \Spec(K)$ 的截面。
由于 $Z_K$ 分离，该截面为闭浸入（引理 \ref{spaces-morphisms-lemma-section-immersion}）。
所以 $Z_K = \Spec(K)$。

\medskip\noindent
取满 \'etale 态射 $V \to Z$，其中 $V$ 为仿射概形（《空间的性质》引理
\ref{spaces-properties-lemma-quasi-compact-affine-cover}）。记
$V = \Spec(B)$。由于 $Z$ 分离，$V \times_Z \Spec(A') = \Spec(C)$ 为仿射。
注意 $B \to C$ 单射，因为由引理
\ref{spaces-morphisms-lemma-quasi-compact-scheme-theoretic-image}，$V$ 是
$V \times_Z \Spec(A') \to V$ 的概形论像。另一方面，$A' \to C$ 是
\'etale 的，因为它对应于 $V \to Z$ 的基变换。由于 $A'$ 是无挠
$A$-模，$A' \to C$ 的平坦性蕴含 $C$ 是无挠 $A$-模，进而 $B$ 是无挠
$A$-模。注意，作为 $A$-模无挠等价于平坦（《关于代数的进一步结果》引理
\ref{more-algebra-lemma-valuation-ring-torsion-free-flat}）。接下来写成
$$
V \times_Z V = \Spec(B')
$$
。由于 $V \to Z$ 为 \'etale 态射，两个环映射 $B \to B'$ 都是 \'etale 的。
因 $Z_K = \Spec(K)$，典范满射 $B \otimes_A B \to B'$ 在 $A$ 上与 $K$
张量后成为同构。但由上述讨论，$B \otimes_A B$ 作为 $A$-模无挠。
故 $B' = B \otimes_A B$。由此可知，环映射 $A \to B$ 沿忠实平坦环映射
% P08-E245: frozen source surjectivity justification preserved; see ERRATA_CANDIDATES.jsonl.
$A \to B$ 的基变换是 \'etale 的（注意，因 $X \to \Spec(A)$ 满，
$\Spec(B) \to \Spec(A)$ 满）。因此 $A \to B$ 为 \'etale 映射
（《下降》引理 \ref{descent-lemma-descending-property-etale}）；换言之，
% P08-E246: frozen source codomain preserved; see ERRATA_CANDIDATES.jsonl.
$V \to X$ 为 \'etale 态射。由于
$V \times_Z V = V \times_{\Spec(A)} V$，可知代数空间
$Z = \Spec(A)$（例如由《空间》引理 \ref{spaces-lemma-space-presentation}），
证明完毕。
\end{proof}

\begin{lemma}
\label{spaces-morphisms-lemma-usual-enough}
设 $S$ 为概形，$f : X \to Y$ 为 $S$ 上代数空间的分离态射。下列条件等价：
\begin{enumerate}
\item $f$ 满足定义 \ref{spaces-morphisms-definition-valuative-criterion} 中赋值判据的存在性部分；
\item 对任意由实线箭头构成的交换图
$$
\xymatrix{
\Spec(K) \ar[r] \ar[d] & X \ar[d] \\
\Spec(A) \ar[r] \ar@{-->}[ru] & Y
}
$$
，其中 $A$ 为分式域是 $K$ 的赋值环，都存在虚线箭头；换言之，$f$ 满足
《概形》定义 \ref{schemes-definition-valuative-criterion} 中赋值判据的存在性部分。
\end{enumerate}
\end{lemma}

\begin{proof}
只需证明 (1) 蕴含 (2)。设给定 (2) 中的交换图
$$
\xymatrix{
\Spec(K) \ar[r] \ar[d] & X \ar[d] \\
\Spec(A) \ar[r] & Y
}
$$
。由 (1)，存在定义 \ref{spaces-morphisms-definition-valuative-criterion} 中的交换图
$$
\xymatrix{
\Spec(K') \ar[r] \ar[d] & \Spec(K) \ar[r] & X \ar[d] \\
\Spec(A') \ar[r] \ar[rru] & \Spec(A) \ar[r] & Y
}
$$
，其中 $K \subset K'$ 任意。由引理 \ref{spaces-morphisms-lemma-push-down-solution}，能找到嵌入
此图的态射 $\Spec(A) \to X$，即 (2) 成立。
\end{proof}

\begin{example}
\label{spaces-morphisms-example-finite-separable-needed}
考虑《空间》例 \ref{spaces-example-non-representable-descent} 中构造的代数空间
$X$。回顾它是把原点加倍之仿射直线的 Galois 扭曲；该扭曲对应于二次 Galois
域扩张 $k'/k$。因而它带有态射
$$
\pi : X \longrightarrow S = \mathbf{A}^1_k
$$
，且该态射拟紧。我们断言 $\pi$ 普遍闭。事实上，沿
$\Spec(k') \to \Spec(k)$ 作基变换后，态射 $\pi$ 等同于态射
$$
\text{原点加倍的仿射直线}
\longrightarrow
\text{仿射直线}
$$
，而它普遍闭（略去若干细节）。由于态射 $\Spec(k') \to \Spec(k)$ 普遍闭且满，
追图可见 $\pi$ 普遍闭。另一方面，考虑图
$$
\xymatrix{
\Spec(k((x))) \ar[r] \ar[d] & X \ar[d]^\pi \\
\Spec(k[[x]]) \ar[r] \ar@{..>}[ru] & \mathbf{A}^1_k
}
$$
由于 $X$ 中位于 $0 \in \mathbf{A}^1_k$ 上方的唯一一点对应于单态射
$\Spec(k') \to X$，显然不可能存在虚线箭头！这说明一般而言需要有限可分域扩张。
\end{example}

\begin{lemma}
\label{spaces-morphisms-lemma-base-change-valuative-criteria}
满足赋值判据之存在性部分（相应地，唯一性部分）的代数空间态射，沿代数空间的
任意态射作基变换后仍满足赋值判据的存在性部分（相应地，唯一性部分）。
\end{lemma}

\begin{proof}
设 $f : X \to Y$ 为概形 $S$ 上代数空间的态射，并设 $Z \to Y$ 为 $S$ 上
代数空间的任意态射。考虑下列由实线箭头构成的交换图
$$
\xymatrix{
\Spec(K) \ar[r] \ar[d] & Z \times_Y X \ar[r] \ar[d] & X \ar[d] \\
\Spec(A) \ar[r] \ar@{-->}[ru] \ar@{-->}[rru] & Z \ar[r] & Y
}
$$
。使该图交换的西北方向虚线箭头之集合，与使该图交换的西偏北方向虚线箭头之
集合一一对应。这证明了“唯一性”情形。对存在性部分，设 $f$ 满足赋值判据的
存在性部分。若给定上述实线交换图，则由假设，存在域扩张 $K'/K$、支配 $A$
的赋值环 $A' \subset K'$，以及嵌入下列交换图的态射 $\Spec(A') \to X$
$$
\xymatrix{
\Spec(K') \ar[r] \ar[d] &
\Spec(K) \ar[r] & Z \times_Y X \ar[r] & X \ar[d] \\
\Spec(A') \ar[r] \ar[rrru] & \Spec(A) \ar[r] & Z \ar[r] & Y
}
$$
。由上述说明，斜箭头对应于所需箭头
$\Spec(A') \to Z \times_Y X$。
\end{proof}

\begin{lemma}
\label{spaces-morphisms-lemma-composition-valuative-criteria}
两个满足赋值判据（存在性部分，相应地，唯一性部分）的代数空间态射之复合，
仍满足赋值判据（存在性部分，相应地，唯一性部分）。
\end{lemma}

\begin{proof}
设 $f : X \to Y$、$g : Y \to Z$ 为概形 $S$ 上代数空间的态射。考虑下列
由实线箭头构成的交换图
$$
\xymatrix{
\Spec(K) \ar[dd] \ar[r] & X \ar[d]^f \\
& Y \ar[d]^g \\
\Spec(A) \ar[r] \ar@{-->}[ru] \ar@{-->}[ruu] & Z
}
$$
若 $g$ 满足唯一性部分，则至多存在一个使图交换的西北方向虚线箭头。
若 $f$ 也满足唯一性部分，则至多存在一个使图交换的北偏西方向虚线箭头。
存在性情形的证明来自考察下图
$$
\xymatrix{
\Spec(K'') \ar[r] \ar[dd] &
\Spec(K') \ar[r] &
\Spec(K) \ar[r] &
X \ar[d]^f \\
& & & Y \ar[d]^g \\
\Spec(A'') \ar[r] \ar[rrruu] &
\Spec(A') \ar[r] \ar[rru] &
\Spec(A) \ar[r] &
Z
}
$$
具体而言，$g$ 的存在性部分给出扩张 $K'$、赋值环 $A'$ 及箭头
$\Spec(A') \to Y$；随后，$f$ 的存在性部分给出扩张 $K''$、赋值环 $A''$
及箭头 $\Spec(A'') \to X$。
\end{proof}






\section{普遍闭性的赋值判据}
\label{spaces-morphisms-section-valuative-criterion-universally-closed}

\noindent
对拟紧态射，赋值判据的存在性部分蕴含普遍闭性，见引理
\ref{spaces-morphisms-lemma-quasi-compact-existence-universally-closed}。在概形情形，这是一个
“当且仅当”的命题，但对代数空间的态射则不成立。例
\ref{spaces-morphisms-example-strange-universally-closed} 表明
$\mathbf{A}^1_k/\mathbf{Z} \to \Spec(k)$ 普遍闭，但容易看出赋值判据的存在性
部分不成立。我们将在《良好空间》第
\ref{decent-spaces-section-valuative-criterion-universally-closed} 节中重访这一主题，
并证明当态射的源为良好空间时逆命题成立（相对版本另见《良好空间》引理
\ref{decent-spaces-lemma-re-characterize-universally-closed}）。

\begin{lemma}
\label{spaces-morphisms-lemma-quasi-compact-existence-universally-closed}
设 $S$ 为概形，$f : X \to Y$ 为 $S$ 上代数空间的态射。设
\begin{enumerate}
\item $f$ 拟紧；
\item $f$ 满足赋值判据的存在性部分。
\end{enumerate}
则 $f$ 普遍闭。
\end{lemma}

\begin{proof}
由引理 \ref{spaces-morphisms-lemma-base-change-quasi-compact} 与
\ref{spaces-morphisms-lemma-base-change-valuative-criteria}，性质 (1)、(2) 在任意基变换下保持。
由引理 \ref{spaces-morphisms-lemma-universally-closed-local}，只需证明：每当 $T$ 是映入 $Y$ 的
$S$ 上仿射概形时，$|T \times_Y X| \to |T|$ 闭。因此，只需证明：若 $Y$
是仿射概形，$f : X \to Y$ 拟紧且满足赋值判据的存在性部分，则
$f : |X| \to |Y|$ 闭。此时 $X$ 是拟紧代数空间。由《空间的性质》引理
\ref{spaces-properties-lemma-quasi-compact-affine-cover}，存在仿射概形 $U$
及满 \'etale 态射 $\varphi : U \to X$。设 $T \subset |X|$ 闭。
逆像 $\varphi^{-1}(T) \subset U$ 闭，因而是某仿射闭子概形 $Z \subset U$
的点集。于是由《代数》引理
\ref{algebra-lemma-image-stable-specialization-closed}，只要
$f(T) = f(\varphi(|Z|)) \subset |Y|$ 在特化下封闭，它就是闭集。

\medskip\noindent
设 $y' \leadsto y$ 为 $Y$ 中的特化，且 $y' \in f(T)$。取一点
$x' \in T \subset |X|$，使其在 $f$ 下映至 $y'$。可以用某个域 $K$ 上的态射
$\Spec(K) \to X$ 表示 $x'$。于是有图
$$
\xymatrix{
\Spec(K) \ar[r]_-{x'} \ar[d] & X \ar[d]^f \\
\Spec(\mathcal{O}_{Y, y}) \ar[r] & Y,
}
$$
左侧竖直映射的存在性见《概形》第 \ref{schemes-section-points} 节。取支配环映射
$\mathcal{O}_{Y,y} \to K$ 之像的赋值环 $A \subset K$（这是可能的，因为当
$y' \not = y$ 时，该像是局部环而非域；见《代数》引理
\ref{algebra-lemma-dominate}）。由假设，存在域扩张 $K'/K$、支配 $A$ 的赋值环
$A' \subset K'$，以及嵌入该交换图的态射 $\Spec(A') \to X$。由于 $A'$ 支配
$A$，且 $A$ 支配 $\mathcal{O}_{Y,y}$，可知 $\Spec(A')$ 的闭点映至一点
$x \in X$，满足 $f(x) = y$，且该点是 $x'$ 的特化。由于 $T$ 闭，故
$x \in T$，从而如所需 $y \in f(T)$。
\end{proof}

\noindent
下一个引理将在《良好空间》引理
\ref{decent-spaces-lemma-re-characterize-universally-closed} 中得到推广。

\begin{lemma}
\label{spaces-morphisms-lemma-characterize-universally-closed-quasi-separated}
设 $S$ 为概形，$f : X \to Y$ 为 $S$ 上代数空间的态射。
\begin{enumerate}
\item 若 $f$ 拟分离且普遍闭，则 $f$ 满足赋值判据的存在性部分。
\item 若 $f$ 拟紧且拟分离，则 $f$ 普遍闭，当且仅当赋值判据的存在性部分成立。
\end{enumerate}
\end{lemma}

\begin{proof}
若 (1) 成立，则将其与引理
\ref{spaces-morphisms-lemma-quasi-compact-existence-universally-closed} 合用即得 (2)。
设 $f$ 拟分离且普遍闭。设给定定义
\ref{spaces-morphisms-definition-valuative-criterion} 中的图
$$
\xymatrix{
\Spec(K) \ar[r] \ar[d] & X \ar[d] \\
\Spec(A) \ar[r] & Y
}
$$
。形式论证表明，所需图
$$
\xymatrix{
\Spec(K') \ar[r] \ar[d] & \Spec(K) \ar[r] & X \ar[d] \\
\Spec(A') \ar[r] \ar[rru] & \Spec(A) \ar[r] & Y
}
$$
的存在性由态射 $X_A \to \Spec(A)$ 情形下的存在性得出。由于拟分离和普遍闭
都在基变换下保持，本引理由下一段中的结果得出。

\medskip\noindent
考虑由实线箭头构成的图
$$
\xymatrix{
\Spec(K) \ar[r]_-x \ar[d] & X \ar[d]^f \\
\Spec(A) \ar@{=}[r] \ar@{..>}[ru] & \Spec(A)
}
$$
，其中 $A$ 为分式域是 $K$ 的赋值环。由引理
\ref{spaces-morphisms-lemma-quasi-compact-permanence} 及 $f$ 拟分离这一事实，态射 $x$ 拟紧。
由于 $f$ 普遍闭，特别有 $|f|(\overline{\{x\}})$ 在 $\Spec(A)$ 中为闭集。
由于该像包含 $\Spec(A)$ 的一般点，存在位于 $x$ 的闭包中的点
$x' \in |X|$，映至 $\Spec(A)$ 的闭点。由引理
\ref{spaces-morphisms-lemma-reach-points-scheme-theoretic-image}，可找到交换图
$$
\xymatrix{
\Spec(K') \ar[r] \ar[d] & \Spec(K) \ar[d] \\
\Spec(A') \ar[r] & X
}
$$
，使得 $\Spec(A')$ 的闭点映至 $x' \in |X|$。由此，
$\Spec(A') \to \Spec(A)$ 把闭点映至闭点；换言之，$A'$ 支配 $A$，
证明完毕。
\end{proof}

\begin{lemma}
\label{spaces-morphisms-lemma-characterize-universally-closed-separated}
设 $S$ 为概形，$f : X \to Y$ 为 $S$ 上代数空间的态射。设 $f$ 拟紧且分离。
下列条件等价：
\begin{enumerate}
\item $f$ 普遍闭；
\item 定义 \ref{spaces-morphisms-definition-valuative-criterion} 中赋值判据的存在性部分成立；
\item 对任意由实线箭头构成的交换图
$$
\xymatrix{
\Spec(K) \ar[r] \ar[d] & X \ar[d] \\
\Spec(A) \ar[r] \ar@{-->}[ru] & Y
}
$$
，其中 $A$ 为分式域是 $K$ 的赋值环，都存在虚线箭头；换言之，$f$ 满足
《概形》定义 \ref{schemes-definition-valuative-criterion} 中赋值判据的存在性部分。
\end{enumerate}
\end{lemma}

\begin{proof}
由于 $f$ 分离，由引理 \ref{spaces-morphisms-lemma-usual-enough}，(2) 与 (3) 等价。
(3) 与 (1) 的等价性由引理
\ref{spaces-morphisms-lemma-characterize-universally-closed-quasi-separated} 得出。
\end{proof}

\begin{lemma}
\label{spaces-morphisms-lemma-lift-valuation-ring-through-flat-morphism}
设 $S$ 为概形，$f : X \to Y$ 为 $S$ 上代数空间的平坦态射。设
$\Spec(A) \to Y$ 为态射，其中 $A$ 为赋值环。若 $\Spec(A)$ 的闭点映至
$|Y|$ 中属于 $|X| \to |Y|$ 之像的一点，则存在交换图
$$
\xymatrix{
\Spec(A') \ar[r] \ar[d] & X \ar[d] \\
\Spec(A) \ar[r] & Y
}
$$
，其中 $A \to A'$ 为赋值环的扩张（《关于代数的进一步结果》定义
\ref{more-algebra-definition-extension-valuation-rings}）。
\end{lemma}

\begin{proof}
基变换 $X_A \to \Spec(A)$ 平坦（引理 \ref{spaces-morphisms-lemma-base-change-flat}），且
$\Spec(A)$ 的闭点属于 $|X_A| \to |\Spec(A)|$ 的像（《空间的性质》引理
\ref{spaces-properties-lemma-points-cartesian}）。故可设 $Y = \Spec(A)$。
设 $U \to X$ 为满 \'etale 态射，其中 $U$ 为概形。取 $u \in U$ 映至
$\Spec(A)$ 的闭点。考虑平坦局部环映射
$A \to B = \mathcal{O}_{U,u}$。由《代数》引理
\ref{algebra-lemma-ff-rings}，存在素理想 $\mathfrak q \subset B$，使得
$\mathfrak q$ 位于 $(0) \subset A$ 上方。由《代数》引理 \ref{algebra-lemma-dominate}，可找到
支配 $B/\mathfrak q$ 的赋值环 $A' \subset \kappa(\mathfrak q)$。
所诱导的态射 $\Spec(A') \to U \to X$ 即为本引理问题的解。
\end{proof}

\begin{lemma}
\label{spaces-morphisms-lemma-refined-valuative-criterion-universally-closed}
设 $S$ 为概形，并设 $f : X \to Y$ 与 $h : U \to X$ 为 $S$ 上代数空间的
态射。若
\begin{enumerate}
\item $f$ 与 $h$ 拟紧；
\item $|h|(|U|)$ 在 $|X|$ 中稠密；
\end{enumerate}
并且对任意由实线箭头构成的交换图
$$
\xymatrix{
\Spec(K) \ar[r] \ar[d] & U \ar[r] & X \ar[d] \\
\Spec(A) \ar[rr] \ar@{-->}[rru] & & Y
}
$$
，其中 $A$ 为分式域是 $K$ 的赋值环，都有
\begin{enumerate}
\item[(3)] 至多存在一个使该图交换的虚线箭头；
\item[(4)] 存在域扩张 $K'/K$、支配 $A$ 的赋值环
$A' \subset K'$ 及态射 $\Spec(A') \to X$，使得下图交换
$$
\xymatrix{
\Spec(K') \ar[r] \ar[d] & \Spec(K) \ar[r] & U \ar[r] & X \ar[d] \\
\Spec(A') \ar[r] \ar[rrru] & \Spec(A) \ar[rr] & & Y
}
$$
\end{enumerate}
，则 $f$ 普遍闭。若此外
\begin{enumerate}
\item[(5)] $f$ 拟分离，
\end{enumerate}
则 $f$ 分离且普遍闭。
\end{lemma}

\begin{proof}
设 (1)、(2)、(3)、(4) 成立。我们将验证 $f$ 满足赋值判据的存在性部分；
由引理 \ref{spaces-morphisms-lemma-quasi-compact-existence-universally-closed}，这将蕴含 $f$
普遍闭。为此，考虑交换图
\begin{equation}
\label{spaces-morphisms-equation-start-with}
\vcenter{
\xymatrix{
\Spec(K) \ar[r] \ar[d] & X \ar[d] \\
\Spec(A) \ar[r] & Y
}
}
\end{equation}
，其中 $A$ 为赋值环，$K$ 为 $A$ 的分式域。注意，赋值环与域均既约，故由
《空间的性质》引理 \ref{spaces-properties-lemma-map-into-reduction}，可以分别用
$U$、$X$、$S$ 的既约化替换它们。在此情形，$h(U)$ 稠密这一假设意味着
$h : U \to X$ 的概形论像是 $X$，见引理
\ref{spaces-morphisms-lemma-scheme-theoretic-image-reduced}。

\medskip\noindent
化归至 $Y$ 仿射的情形。取 \'etale 态射 $\Spec(R) \to Y$，使得
$\Spec(A)$ 的闭点映至 $\Im(|\Spec(R)| \to |Y|)$ 中的一个元素。由引理
\ref{spaces-morphisms-lemma-lift-valuation-ring-through-flat-morphism}，可找到赋值环的局部环映射
$A \to A'$ 以及嵌入下列交换图的态射 $\Spec(A') \to \Spec(R)$
$$
\xymatrix{
\Spec(A') \ar[r] \ar[d] & \Spec(R) \ar[d] \\
\Spec(A) \ar[r] & Y
}
$$
。由于定义 \ref{spaces-morphisms-definition-valuative-criterion} 允许赋值环扩张，显然可以用
$A'$ 替换 $A$，用 $\Spec(R)$ 替换 $Y$，用 $X \times_Y \Spec(R)$ 替换 $X$，
并用 $U \times_Y \Spec(R)$ 替换 $U$。

\medskip\noindent
从现在起，设 $Y = \Spec(R)$ 为仿射概形。设 $\Spec(B) \to X$ 为来自仿射
概形的 \'etale 态射，并使态射 $\Spec(K) \to X$ 属于
$|\Spec(B)| \to |X|$ 的像。由于可以用扩张 $K' \supset K$ 替换 $K$，并用
支配 $A$ 的赋值环 $A' \subset K'$ 替换 $A$（它由《代数》引理
\ref{algebra-lemma-dominate} 存在），由 $|X|$ 的定义，可以设态射
$\Spec(K) \to X$ 分解经过 $\Spec(B)$。换言之，可以把 $K$ 视为
$B$-代数。取 $B$ 上的多项式代数 $P$ 及 $B$-代数满射 $P \to K$。
于是 $\Spec(P) \to X$ 作为复合 $\Spec(P) \to \Spec(B) \to X$ 而平坦。
因此，由引理 \ref{spaces-morphisms-lemma-flat-base-change-scheme-theoretic-image}，态射
$U \times_X \Spec(P) \to \Spec(P)$ 的概形论像是 $\Spec(P)$。由引理
\ref{spaces-morphisms-lemma-reach-points-scheme-theoretic-image}，可找到交换图
$$
\xymatrix{
\Spec(K') \ar[r] \ar[d] & U \times_X \Spec(P) \ar[d] \\
\Spec(A') \ar[r] & \Spec(P)
}
$$
，其中 $A'$ 为赋值环，$K'$ 为 $A'$ 的分式域，并且 $\Spec(A')$ 的闭点映至
$\Spec(K) \subset \Spec(P)$。换言之，存在 $B$-代数映射
$\varphi : K \to A'/\mathfrak m_{A'}$。取支配 $\varphi(A)$ 且分式域为
$K'' = A'/\mathfrak m_{A'}$ 的赋值环
$A'' \subset A'/\mathfrak m_{A'}$（《代数》引理
\ref{algebra-lemma-dominate}）。令
$$
C = \{\lambda \in A' \mid \lambda \bmod \mathfrak m_{A'} \in A''\}.
$$
由《代数》引理 \ref{algebra-lemma-stack-valuation-rings}，这是赋值环。
由于 $C$ 是分式域为 $K'$ 的 $R$-代数，得到由实线箭头构成的交换图
$$
\xymatrix{
\Spec(K'_1) \ar@{-->}[r] \ar@{-->}[d] &
\Spec(K') \ar[r] \ar[d] & U \ar[r] & X \ar[d] \\
\Spec(C_1) \ar@{-->}[r] \ar@{-->}[rrru] & \Spec(C) \ar[rr] & & Y
}
$$
，如引理陈述所示。因此假设 (4) 产生 $C \to C_1$ 以及使该图交换的虚线箭头。
令 $A_1' = (C_1)_\mathfrak p$ 为 $C_1$ 在位于
$\mathfrak m_{A'} \subset C$ 上方的素理想 $\mathfrak p \subset C_1$ 处的局部化。
由《关于代数的进一步结果》引理
\ref{more-algebra-lemma-valuation-ring-torsion-free-flat}，$C \to C_1$ 平坦，
故这样的素理想 $\mathfrak p$ 由《代数》引理
\ref{algebra-lemma-local-flat-ff} 与 \ref{algebra-lemma-ff-rings} 存在。
注意，$A'$ 是 $C$ 在 $\mathfrak m_{A'}$ 处的局部化，而 $A'_1$ 是赋值环
（《代数》引理 \ref{algebra-lemma-make-valuation-rings}）。换言之，
$A' \to A'_1$ 是赋值环的局部环映射。假设 (3) 蕴含图
$$
\xymatrix{
\Spec(A'_1) \ar[r] \ar[d] & \Spec(C_1) \ar[r] & X \\
\Spec(A') \ar[r] & \Spec(P) \ar[r] & \Spec(B) \ar[u]
}
$$
交换。因此，态射 $\Spec(C_1) \to X$ 在 $\Spec(C_1/\mathfrak p)$ 上的限制，
在 $\Spec(C_1/\mathfrak p)$ 的一般点处限制为复合
$$
\Spec(\kappa(\mathfrak p)) \to
\Spec(A'/\mathfrak m_{A'}) = \Spec(K'') \to
\Spec(K) \to X
$$
。此外，$C_1/\mathfrak p$ 是赋值环（《代数》引理
\ref{algebra-lemma-make-valuation-rings}），它支配 $A''$，而后者支配 $A$。
因此，态射 $\Spec(C_1/\mathfrak p) \to X$ 如所需见证了图
(\ref{spaces-morphisms-equation-start-with}) 的赋值判据之存在性部分。

\medskip\noindent
接下来，设 (5) 也成立；换言之，态射
% P08-E247: frozen source diagonal base preserved; see ERRATA_CANDIDATES.jsonl.
$\Delta : X \to X \times_S X$ 拟紧。此时，$h$ 与 $\Delta$ 满足假设
(1)--(4)。故证明的第一部分表明 $\Delta$ 普遍闭。由引理
\ref{spaces-morphisms-lemma-separated-diagonal-proper}，可知 $f$ 分离。
\end{proof}

















\section{分离性的赋值判据}
\label{spaces-morphisms-section-valuative-separatedness}

\noindent
我们先证明一个逆命题，然后陈述该判据。

\begin{lemma}
\label{spaces-morphisms-lemma-separated-implies-valuative}
设 $S$ 为概形。设 $f : X \to Y$ 为 $S$ 上代数空间的态射。
若 $f$ 分离，则 $f$ 满足赋值判据的唯一性部分。
\end{lemma}

\begin{proof}
设给定定义 \ref{spaces-morphisms-definition-valuative-criterion} 中的图。假设有两个不同的态射
$a, b : \Spec(A) \to X$ 可嵌入该图。令 $Z \subset \Spec(A)$ 为 $a$ 与 $b$
的等化子。于是
$Z = \Spec(A) \times_{(a, b), X \times_Y X, \Delta} X$。
若 $f$ 分离，则 $\Delta$ 为闭浸入，故这是 $\Spec(A)$ 的闭子概形。
依假设，它包含 $\Spec(A)$ 的一般点。由于 $A$ 是整环，这蕴含
$Z = \Spec(A)$。因此如所需 $a = b$。
\end{proof}

% P08-E248: frozen source lemma-title grammar noted; see ERRATA_CANDIDATES.jsonl.
\begin{lemma}[分离性的赋值判据]
\label{spaces-morphisms-lemma-valuative-criterion-separatedness}
设 $S$ 为概形。设 $f : X \to Y$ 为 $S$ 上代数空间的态射。假设
\begin{enumerate}
\item 态射 $f$ 拟分离；
\item 态射 $f$ 满足赋值判据的唯一性部分。
\end{enumerate}
则 $f$ 分离。
\end{lemma}

\begin{proof}
假设 (1) 意味着 $\Delta_{X/Y}$ 拟紧。我们断言态射
$\Delta_{X/Y} : X \to X \times_Y X$ 满足赋值判据的存在性部分。
设给定由实线箭头构成的交换图
$$
\xymatrix{
\Spec(K) \ar[r] \ar[d] & X \ar[d] \\
\Spec(A) \ar[r] \ar@{-->}[ru] & X \times_Y X
}
$$
。右下方的箭头对应于 $Y$ 上的一对态射
$a, b : \Spec(A) \to X$。由假设 (2)，可知 $a = b$，故以 $a$ 为虚线箭头即可。
于是引理 \ref{spaces-morphisms-lemma-quasi-compact-existence-universally-closed} 适用，从而
$\Delta_{X/Y}$ 普遍闭。由于 $\Delta_{X/Y}$ 总是局部有限型且分离，由
《关于态射的进一步结果》引理 \ref{more-morphisms-lemma-characterize-finite}，
可知 $\Delta_{X/Y}$ 是有限态射（也可使用《空间》引理
\ref{spaces-lemma-representable-transformations-property-implication}
的一般原理）。此时 $\Delta_{X/Y}$ 是可表的有限单态射，因而由《态射》引理
\ref{morphisms-lemma-finite-monomorphism-closed}，它是闭浸入。
\end{proof}

\begin{lemma}
\label{spaces-morphisms-lemma-characterize-separated-and-universally-closed}
设 $S$ 为概形。设 $f : X \to Y$ 为 $S$ 上代数空间的态射。假设 $f$ 拟紧且
拟分离。则下列条件等价：
\begin{enumerate}
\item $f$ 分离且普遍闭；
\item 定义 \ref{spaces-morphisms-definition-valuative-criterion} 中的赋值判据成立；
\item 对任意由实线箭头构成的交换图
$$
\xymatrix{
\Spec(K) \ar[r] \ar[d] & X \ar[d] \\
\Spec(A) \ar[r] \ar@{-->}[ru] & Y
}
$$
，其中 $A$ 为分式域是 $K$ 的赋值环，都存在唯一的虚线箭头；换言之，$f$
满足《概形》定义 \ref{schemes-definition-valuative-criterion} 中的赋值判据。
\end{enumerate}
\end{lemma}

\begin{proof}
由于 $f$ 拟分离，赋值判据的唯一性部分蕴含 $f$ 分离
% P08-E249: frozen source spelling error noted; see ERRATA_CANDIDATES.jsonl.
（引理 \ref{spaces-morphisms-lemma-valuative-criterion-separatedness}）。反之，若 $f$ 分离，
则它满足赋值判据的唯一性部分（引理 \ref{spaces-morphisms-lemma-separated-implies-valuative}）。
综上，在三个情形中，态射 $f$ 均分离且满足赋值判据的唯一性部分。此时本引理是
引理 \ref{spaces-morphisms-lemma-characterize-universally-closed-separated} 的形式推论。
\end{proof}



\section{固有性的赋值判据}
\label{spaces-morphisms-section-valuative-criterion-properness}

\noindent
现给出如下命题。

\begin{lemma}[固有性的赋值判据]
\label{spaces-morphisms-lemma-characterize-proper}
设 $S$ 为概形。设 $f : X \to Y$ 为 $S$ 上代数空间的态射。假设 $f$ 有限型且
拟分离。则下列条件等价：
\begin{enumerate}
\item $f$ 固有；
\item 定义 \ref{spaces-morphisms-definition-valuative-criterion} 中的赋值判据成立；
\item 对任意由实线箭头构成的交换图
$$
\xymatrix{
\Spec(K) \ar[r] \ar[d] & X \ar[d] \\
\Spec(A) \ar[r] \ar@{-->}[ru] & Y
}
$$
，其中 $A$ 为分式域是 $K$ 的赋值环，都存在唯一的虚线箭头；换言之，$f$
满足《概形》定义 \ref{schemes-definition-valuative-criterion} 中的赋值判据。
\end{enumerate}
\end{lemma}

\begin{proof}
这是引理 \ref{spaces-morphisms-lemma-characterize-separated-and-universally-closed} 与各定义的
形式推论。
\end{proof}





\section{整态射与有限态射}
\label{spaces-morphisms-section-integral}

\noindent
我们已经在第 \ref{spaces-morphisms-section-representable} 节中定义了代数空间的可表示态射何谓
整态射（相应地，有限态射）。

\begin{lemma}
\label{spaces-morphisms-lemma-integral-representable}
设 $S$ 为概形。设 $f : X \to Y$ 为 $S$ 上代数空间的可表示态射。则
$f$ 为整态射，相应地为有限态射（依第 \ref{spaces-morphisms-section-representable} 节之意义），
当且仅当对所有仿射概形 $Z$ 及态射 $Z \to Y$，概形 $X \times_Y Z$ 为仿射的，
且在 $Z$ 上为整的，相应地为有限的。
\end{lemma}

\begin{proof}
这直接由概形的整（相应地，有限）态射之定义得出
（《态射》定义 \ref{morphisms-definition-integral}）。
\end{proof}

\noindent
由此可以作如下定义。

\begin{definition}
\label{spaces-morphisms-definition-integral}
设 $S$ 为概形。设 $f : X \to Y$ 为 $S$ 上代数空间的态射。
\begin{enumerate}
\item 若对每个仿射概形 $Z$ 及态射 $Z \to Y$，代数空间 $X \times_Y Z$
都可由一个在 $Z$ 上为整的仿射概形表示，则称 $f$ 为{\it 整态射}。
\item 若对每个仿射概形 $Z$ 及态射 $Z \to Y$，代数空间 $X \times_Y Z$
都可由一个在 $Z$ 上为有限的仿射概形表示，则称 $f$ 为{\it 有限态射}。
\end{enumerate}
\end{definition}

\begin{lemma}
\label{spaces-morphisms-lemma-integral-local}
设 $S$ 为概形。设 $f : X \to Y$ 为 $S$ 上代数空间的态射。下列条件等价：
\begin{enumerate}
\item $f$ 可表且为整态射（相应地，有限态射）；
\item $f$ 为整态射（相应地，有限态射）；
\item 存在概形 $V$ 及满 \'etale 态射 $V \to Y$，使得
$V \times_Y X \to V$ 为整态射（相应地，有限态射）；
\item 存在扎里斯基覆盖 $Y = \bigcup Y_i$，使每个态射
$f^{-1}(Y_i) \to Y_i$ 都为整态射（相应地，有限态射）。
\end{enumerate}
\end{lemma}

\begin{proof}
显然，(1) 蕴含 (2)；取 $V$ 为 $Y$ 上 \'etale 仿射概形的不交并，即知
(2) 蕴含 (3)，见《空间的性质》引理
\ref{spaces-properties-lemma-cover-by-union-affines}。
假设 $V \to Y$ 如 (3) 所示。则对 $V$ 的每个仿射开集 $W$，
$W \times_Y X$ 都是 $V \times_Y X$ 的仿射开集。因此，由
《空间的性质》引理 \ref{spaces-properties-lemma-subscheme}，可知
$V \times_Y X$ 是概形。此外，态射 $V \times_Y X \to V$ 是仿射的。
这意味着可以应用《空间》引理
\ref{spaces-lemma-morphism-sheaves-with-P-effective-descent-etale}，
因为整（相应地，有限）态射类满足所有所需性质（见《态射》引理
\ref{morphisms-lemma-base-change-finite}，以及《下降》引理
\ref{descent-lemma-descending-property-integral}、
\ref{descent-lemma-descending-property-finite} 和 \ref{descent-lemma-affine}）。
应用该引理所得的结论是 $f$ 可表且为整态射（相应地，有限态射），即 (1) 成立。

\medskip\noindent
(1) 与 (4) 的等价性由整态射（相应地，有限态射）这一性质在目标上是扎里斯基
局部的得出（上面的文献事实上表明，整态射或有限态射在目标上是 fpqc 局部的）。
\end{proof}

\begin{lemma}
\label{spaces-morphisms-lemma-composition-integral}
整（相应地，有限）态射的复合仍为整（相应地，有限）态射。
\end{lemma}

\begin{proof}
略。
\end{proof}

\begin{lemma}
\label{spaces-morphisms-lemma-base-change-integral}
整（相应地，有限）态射的基变换仍为整（相应地，有限）态射。
\end{lemma}

\begin{proof}
略。
\end{proof}

\begin{lemma}
\label{spaces-morphisms-lemma-finite-integral}
代数空间的有限态射是整态射。代数空间的局部有限型整态射是有限态射。
\end{lemma}

\begin{proof}
在两种情形中，态射都是可表的，并且可以在用映入 $Y$ 的仿射概形作基变换后
检验该条件，见引理 \ref{spaces-morphisms-lemma-integral-local}。因此，本引理由概形情形的
同一引理得出，见《态射》引理 \ref{morphisms-lemma-finite-integral}。
\end{proof}

\begin{lemma}
\label{spaces-morphisms-lemma-integral-universally-closed}
设 $S$ 为概形。设 $f : X \to Y$ 为 $S$ 上代数空间的态射。下列条件等价：
\begin{enumerate}
\item $f$ 为整态射；
\item $f$ 仿射且普遍闭。
\end{enumerate}
\end{lemma}

\begin{proof}
在两种情形中，态射都是可表的，并且可以在用映入 $Y$ 的仿射概形作基变换后
检验该条件，见引理 \ref{spaces-morphisms-lemma-integral-local}、\ref{spaces-morphisms-lemma-affine-local} 和
\ref{spaces-morphisms-lemma-universally-closed-local}。因此，结论由《态射》引理
\ref{morphisms-lemma-integral-universally-closed} 得出。
\end{proof}

\begin{lemma}
\label{spaces-morphisms-lemma-finite-quasi-finite}
代数空间的有限态射是拟有限态射。
\end{lemma}

\begin{proof}
设 $f : X \to Y$ 为代数空间的态射。由定义 \ref{spaces-morphisms-definition-integral} 及引理
\ref{spaces-morphisms-lemma-quasi-compact-local}、\ref{spaces-morphisms-lemma-quasi-finite-local}，这两个性质
都可在基变换至 $Y$ 上的仿射概形后检验；换言之，可设 $Y$ 仿射。
若 $f$ 有限，则 $X$ 是概形。因此，结论由概形的相应结果得出，见《态射》引理
\ref{morphisms-lemma-finite-quasi-finite}。
\end{proof}

\begin{lemma}
\label{spaces-morphisms-lemma-finite-proper}
设 $S$ 为概形。设 $f : X \to Y$ 为 $S$ 上代数空间的态射。下列条件等价：
\begin{enumerate}
\item $f$ 有限；
\item $f$ 仿射且固有。
\end{enumerate}
\end{lemma}

\begin{proof}
在两种情形中，态射都是可表的，并且可以在用映入 $Y$ 的仿射概形作基变换后
检验该条件，见引理 \ref{spaces-morphisms-lemma-integral-local}、\ref{spaces-morphisms-lemma-affine-local} 和
\ref{spaces-morphisms-lemma-proper-local}。因此，结论由《态射》引理
\ref{morphisms-lemma-finite-proper} 得出。
\end{proof}

\begin{lemma}
\label{spaces-morphisms-lemma-closed-immersion-finite}
闭浸入是有限的（因而更是整的）。
\end{lemma}

\begin{proof}
略。
\end{proof}

\begin{lemma}
\label{spaces-morphisms-lemma-finite-union-finite}
设 $S$ 为概形。设 $X_i \to Y$，$i = 1, \ldots, n$，为 $S$ 上代数空间的
有限态射。则 $X_1 \amalg \ldots \amalg X_n \to Y$ 也有限。
\end{lemma}

\begin{proof}
通过 \'etale 局部化，这由概形情形得出
（《态射》引理 \ref{morphisms-lemma-finite-union-finite}）。
\end{proof}

\begin{lemma}
\label{spaces-morphisms-lemma-finite-permanence}
设 $S$ 为概形。设 $f : X \to Y$ 与 $g : Y \to Z$ 为 $S$ 上代数空间的态射。
\begin{enumerate}
\item 若 $g \circ f$ 有限且 $g$ 分离，则 $f$ 有限。
\item 若 $g \circ f$ 为整态射且 $g$ 分离，则 $f$ 为整态射。
\end{enumerate}
\end{lemma}

\begin{proof}
假设 $g \circ f$ 有限（相应地，为整态射）且 $g$ 分离。由引理
\ref{spaces-morphisms-lemma-base-change-integral}，基变换 $X \times_Z Y \to Y$ 有限
（相应地，为整态射）。由于 $Y \to Z$ 分离，态射
$X \to X \times_Z Y$ 是闭浸入，见引理 \ref{spaces-morphisms-lemma-section-immersion}。
闭浸入是有限态射（相应地，整态射），见引理
\ref{spaces-morphisms-lemma-closed-immersion-finite}。有限态射（相应地，整态射）的复合仍为
有限态射（相应地，整态射），见引理
\ref{spaces-morphisms-lemma-composition-integral}。故结论成立。
\end{proof}





\section{有限局部自由态射}
\label{spaces-morphisms-section-finite-locally-free}

\noindent
我们已经在第 \ref{spaces-morphisms-section-representable} 节中定义了代数空间的可表示态射何谓
有限局部自由态射。

\begin{lemma}
\label{spaces-morphisms-lemma-finite-locally-free-representable}
设 $S$ 为概形。设 $f : X \to Y$ 为 $S$ 上代数空间的可表示态射。则 $f$
为有限局部自由态射（依第 \ref{spaces-morphisms-section-representable} 节之意义），当且仅当
$f$ 仿射且层 $f_*\mathcal{O}_X$ 是有限局部自由 $\mathcal{O}_Y$-模。
\end{lemma}

\begin{proof}
假设 $f$ 为有限局部自由态射（依第 \ref{spaces-morphisms-section-representable} 节之定义）。
这意味着，对每个源为概形的态射 $V \to Y$，基变换
$f' : V \times_Y X \to V$ 都是概形的有限局部自由态射。这又依概形的
有限局部自由态射之定义，这意味着 $f'_*\mathcal{O}_{V \times_Y X}$
是有限局部自由 $\mathcal{O}_V$-模。可以选取满 \'etale 态射 $V \to Y$。
由《空间的性质》引理
\ref{spaces-properties-lemma-pushforward-etale-base-change-modules}，
可知 $f_*\mathcal{O}_X$ 在 $V$ 上的限制有限局部自由。因此，把《站点上的模》
引理 \ref{sites-modules-lemma-local-final-object} 应用于
$Y_{spaces, \etale}$ 上的层 $f_*\mathcal{O}_X$，可知
$f_*\mathcal{O}_X$ 有限局部自由。

\medskip\noindent
反之，假设 $f$ 仿射，且 $f_*\mathcal{O}_X$ 是有限局部自由
$\mathcal{O}_Y$-模。设 $V$ 为概形，并设 $V \to Y$ 为满 \'etale 态射。
再次由《空间的性质》引理
\ref{spaces-properties-lemma-pushforward-etale-base-change-modules}，
可知 $f'_*\mathcal{O}_{V \times_Y X}$ 有限局部自由。因此，
$f' : V \times_Y X \to V$ 为有限局部自由态射（因为它也仿射）。
由《空间》引理
\ref{spaces-lemma-morphism-sheaves-with-P-effective-descent-etale}，
可知 $f$ 为有限局部自由态射（使用《态射》引理
\ref{morphisms-lemma-base-change-finite-locally-free}，
% P08-E251: frozen source citation separators noted; see ERRATA_CANDIDATES.jsonl.
以及《下降》引理 \ref{descent-lemma-descending-property-finite-locally-free}
和 \ref{descent-lemma-affine}）。故结论成立。
\end{proof}

\noindent
由此可以作如下定义。

\begin{definition}
\label{spaces-morphisms-definition-finite-locally-free}
设 $S$ 为概形。设 $f : X \to Y$ 为 $S$ 上代数空间的态射。若 $f$ 仿射且
$f_*\mathcal{O}_X$ 是有限局部自由 $\mathcal{O}_Y$-模，则称 $f$
为{\it 有限局部自由态射}。
% P08-E250: frozen source grammar noted; see ERRATA_CANDIDATES.jsonl.
在此情形，若层 $f_*\mathcal{O}_X$ 有限局部自由且秩为 $d$，则称 $f$
的{\it 秩}或{\it 次数}为 $d$。
\end{definition}

\begin{lemma}
\label{spaces-morphisms-lemma-finite-locally-free-local}
设 $S$ 为概形。设 $f : X \to Y$ 为 $S$ 上代数空间的态射。下列条件等价：
\begin{enumerate}
\item $f$ 可表且为有限局部自由态射；
\item $f$ 为有限局部自由态射；
\item 存在概形 $V$ 及满 \'etale 态射 $V \to Y$，使得
$V \times_Y X \to V$ 为有限局部自由态射；
\item 存在扎里斯基覆盖 $Y = \bigcup Y_i$，使得每个态射
$f^{-1}(Y_i) \to Y_i$ 都为有限局部自由态射。
\end{enumerate}
\end{lemma}

\begin{proof}
显然，(1) 蕴含 (2)；取 $V$ 为 $Y$ 上 \'etale 仿射概形的不交并，即知
(2) 蕴含 (3)，见《空间的性质》引理
\ref{spaces-properties-lemma-cover-by-union-affines}。
假设 $V \to Y$ 如 (3) 所示。则对 $V$ 的每个仿射开集 $W$，
$W \times_Y X$ 都是 $V \times_Y X$ 的仿射开集。因此，由《空间的性质》
引理 \ref{spaces-properties-lemma-subscheme}，可知 $V \times_Y X$ 是概形。
此外，态射 $V \times_Y X \to V$ 是仿射的。这意味着可以应用《空间》引理
\ref{spaces-lemma-morphism-sheaves-with-P-effective-descent-etale}，
因为有限局部自由态射类满足所有所需性质（见《态射》引理
\ref{morphisms-lemma-base-change-finite-locally-free}，
以及《下降》引理 \ref{descent-lemma-descending-property-finite-locally-free}
和 \ref{descent-lemma-affine}）。应用该引理所得的结论是 $f$ 可表且为
有限局部自由态射，即 (1) 成立。

\medskip\noindent
(1) 与 (4) 的等价性由有限局部自由这一性质在目标上是扎里斯基局部的得出
（上面的文献事实上表明，有限局部自由在目标上是 fpqc 局部的）。
\end{proof}

\begin{lemma}
\label{spaces-morphisms-lemma-composition-finite-locally-free}
有限局部自由态射的复合仍为有限局部自由态射。
\end{lemma}

\begin{proof}
略。
\end{proof}

\begin{lemma}
\label{spaces-morphisms-lemma-base-change-finite-locally-free}
有限局部自由态射的基变换仍为有限局部自由态射。
\end{lemma}

\begin{proof}
略。
\end{proof}

\begin{lemma}
\label{spaces-morphisms-lemma-finite-flat}
设 $S$ 为概形。设 $f : X \to Y$ 为 $S$ 上代数空间的态射。下列条件等价：
\begin{enumerate}
\item $f$ 为有限局部自由态射；
\item $f$ 有限、平坦且局部有限表示。
\end{enumerate}
若 $Y$ 局部诺特，则这些条件还等价于
\begin{enumerate}
\item[(3)] $f$ 有限且平坦。
\end{enumerate}
\end{lemma}

\begin{proof}
在三个情形中，态射都是可表的，并且可以在用满 \'etale 态射
$V \to Y$ 作基变换后检验该性质，见引理 \ref{spaces-morphisms-lemma-integral-local}、
\ref{spaces-morphisms-lemma-finite-locally-free-local}、\ref{spaces-morphisms-lemma-flat-local} 和
\ref{spaces-morphisms-lemma-finite-presentation-local}。若 $Y$ 局部诺特，则 $V$ 局部诺特。
因此，结论由概形情形的相应结果得出，见《态射》引理
\ref{morphisms-lemma-finite-flat}。
\end{proof}





\section{有理映射}
\label{spaces-morphisms-section-rational-maps}

\noindent
本节是《态射》第 \ref{morphisms-section-rational-maps} 节的对应版本。
我们将不再特别说明地使用如下事实：拓扑空间中稠密开集的交仍是稠密开集。

\begin{definition}
\label{spaces-morphisms-definition-rational-map}
设 $S$ 为概形。设 $X$、$Y$ 为 $S$ 上的代数空间。
\begin{enumerate}
\item 设 $f : U \to Y$、$g : V \to Y$ 为 $S$ 上代数空间的态射，分别定义在
$X$ 的稠密开子空间 $U$、$V$ 上。若对某个稠密开子空间
$W \subset U \cap V$ 有 $f|_W = g|_W$，则称 $f$ 与 $g$ {\it 等价}。
\item 从 $X$ 到 $Y$ 的{\it 有理映射}是 (1) 中所定义等价关系的一个等价类。
\item 给定 $S$ 上代数空间的态射 $X \to B$ 与 $Y \to B$。若该等价类有一个
代表元 $f : U \to Y$ 是 $B$ 上的态射，则称从 $X$ 到 $Y$ 的有理映射为
{\it 从 $X$ 到 $Y$ 的 $B$-有理映射}。
\end{enumerate}
\end{definition}

\noindent
对定义 (1) 中的两个态射 $f$、$g$，我们说它们定义同一个有理映射，而不说
它们等价。在以后将考虑的许多情形中，代数空间 $X$ 与 $Y$ 分别含有作为概形的
% P08-E252: frozen source agreement error noted; see ERRATA_CANDIDATES.jsonl.
稠密开子空间 $X'$ 与 $Y'$。此时，从 $X$ 到 $Y$ 的有理映射与《态射》定义
% P08-E253: frozen source cross-reference target preserved; see ERRATA_CANDIDATES.jsonl.
\ref{spaces-morphisms-definition-rational-map} 意义下从 $X'$ 到 $Y'$ 的 $S$-有理映射相同。
于是可以运用为概形建立的全部理论。

\begin{definition}
\label{spaces-morphisms-definition-rational-function}
设 $S$ 为概形。设 $X$ 为 $S$ 上的代数空间。$X$ 上的{\it 有理函数}是从
$X$ 到 $\mathbf{A}^1_S$ 的有理映射。
\end{definition}

\noindent
由《态射》定义 \ref{morphisms-definition-rational-function} 后的讨论可知，
当 $X$ 恰为概形时，这与该处定义的概念相同。

\medskip\noindent
回忆，对 $S$ 上任意概形 $T$，都有典范恒等式
$$
\Mor_S(T, \mathbf{A}^1_S) = \Mor(T, \mathbf{A}^1_\mathbf{Z}) =
\Gamma(T, \mathcal{O}_T)
$$
，见《概形》例 \ref{schemes-example-global-sections}。因此，
$\mathbf{A}^1_S$ 是 $S$ 上概形范畴中的环对象。换言之，加法与乘法定义态射
$$
+ : \mathbf{A}^1_S \times_S \mathbf{A}^1_S \to \mathbf{A}^1_S
\quad\text{且}\quad
* : \mathbf{A}^1_S \times_S \mathbf{A}^1_S \to \mathbf{A}^1_S
$$
，并满足环中加法与乘法的公理（照例是带 $1$ 的交换环）。因此，取值于
$\mathbf{A}^1_S$ 的有理映射集也有自然的环结构。

\begin{definition}
\label{spaces-morphisms-definition-ring-of-rational-functions}
设 $S$ 为概形。设 $X$ 为 $S$ 上的代数空间。$X$ 上的{\it 有理函数环}是环
$R(X)$；其元素为有理函数，加法与乘法如上所述。
\end{definition}

\noindent
稍后我们将定义整代数空间的函数域，见《域上的空间》第
\ref{spaces-over-fields-section-integral-spaces} 节。

\begin{definition}
\label{spaces-morphisms-definition-domain-of-definition}
设 $S$ 为概形。设 $\varphi$ 为 $S$ 上两个代数空间 $X$ 与 $Y$ 之间的有理映射。
若 $\varphi$ 有一个代表元 $(U, f)$ 满足 $x \in |U|$，则称 $\varphi$
{\it 在点 $x \in |X|$ 处有定义}。$\varphi$ 的{\it 定义域}是 $\varphi$
有定义的所有点之集。
\end{definition}

\noindent
通过《空间的性质》引理 \ref{spaces-properties-lemma-open-subspaces}，
把定义域视为 $X$ 的开子空间。依此定义，一般而言并非必有 $\varphi$ 的
某个代表元定义在整个定义域上。

\begin{lemma}
\label{spaces-morphisms-lemma-rational-map-from-reduced-to-separated}
设 $S$ 为概形。设 $X$ 与 $Y$ 为 $S$ 上的代数空间。假设 $X$ 既约且 $Y$
在 $S$ 上分离。设 $\varphi$ 为从 $X$ 到 $Y$ 的有理映射，其定义域为
$U \subset X$。则存在唯一的代数空间态射 $f : U \to Y$ 表示 $\varphi$。
\end{lemma}

\begin{proof}
设 $(V, g)$ 与 $(V', g')$ 为 $\varphi$ 的代表元。则 $g, g'$ 在某个稠密开子空间
$W \subset V \cap V'$ 上相等。另一方面，$g|_{V \cap V'}$ 与
$g'|_{V \cap V'}$ 的等化子 $E$ 是 $V \cap V'$ 的闭子空间，因为它是将
$\Delta : Y \to Y \times_S Y$ 沿由 $g|_{V \cap V'}$ 与 $g'|_{V \cap V'}$
给出的态射 $V \cap V' \to Y \times_S Y$ 作基变换所得。现由
$W \subset E$ 可知 $|E| = |V \cap V'|$。由于 $V \cap V'$ 既约，在概形论意义下
$E = V \cap V'$；换言之，$g|_{V \cap V'} = g'|_{V \cap V'}$，见
《空间的性质》引理 \ref{spaces-properties-lemma-map-into-reduction}。
因此，可以把 $\varphi$ 的各代表元 $g : V \to Y$ 粘合为态射
$f : U \to Y$，因为 $\coprod V \to U$ 是 fppf 层的满射，且
$\coprod_{V, V'} V \cap V' = (\coprod V) \times_U (\coprod V)$。
\end{proof}

\noindent
一般而言，复合有理映射没有意义。原因是第一个有理映射的代表元之像可能与
第二个有理映射的定义域没有交。不过，若假设所论空间不可约，并只考察
支配有理映射，则可以复合有理映射。

\begin{definition}
\label{spaces-morphisms-definition-dominant-rational}
设 $S$ 为概形。设 $X$ 与 $Y$ 为 $S$ 上的代数空间。假设 $|X|$ 与 $|Y|$
不可约。若从 $X$ 到 $Y$ 的有理映射的任意代表元 $f : U \to Y$ 都是定义
\ref{spaces-morphisms-definition-dominant} 意义下的支配态射，则称该有理映射是{\it 支配的}。
\end{definition}

\noindent
可以把不可约代数空间 $X$ 与 $Y$ 之间的支配有理映射 $\varphi$，同从 $Y$
到 $Z$ 的任意有理映射 $\psi$ 复合。具体地，选取代表元
$f : U \to Y$，其中 $|U| \subset |X|$ 开且稠密；以及代表元
$g : V \to Z$，其中 $|V| \subset |Y|$ 开且稠密。则
$W = |f|^{-1}(V) \subset |X|$ 开且非空（因为 $|f|$ 的像稠密，故必与非空开集
$V$ 相交）；又因 $|X|$ 不可约，该开集稠密。定义 $\psi \circ \varphi$ 为
$g \circ f|_W : W \to Z$ 的等价类。略去其良定义性的验证。

\medskip\noindent
由此得到一个范畴，其对象为 $S$ 上不可约代数空间，态射为支配有理映射。

\begin{definition}
\label{spaces-morphisms-definition-birational-spaces}
设 $S$ 为概形。设 $X$ 与 $Y$ 为 $S$ 上的代数空间，且 $|X|$ 与 $|Y|$
不可约。若 $X$ 与 $Y$ 在 $S$ 上不可约代数空间及支配有理映射所成的范畴中
同构，则称 $X$ 与 $Y$ {\it 双有理}。
\end{definition}

\noindent
若 $X$ 与 $Y$ 是双有理的不可约代数空间，则对所有代数空间 $Z$，从 $X$
到 $Z$ 的有理映射集与从 $Y$ 到 $Z$ 的有理映射集双射（关于 $Z$ 函子性地）。
% P08-E254: frozen source number agreement noted; see ERRATA_CANDIDATES.jsonl.
对“一般的”不可约代数空间，这只是一个可能的定义。另一种定义是要求
% P08-E255: frozen source missing complementizer noted; see ERRATA_CANDIDATES.jsonl.
$X$ 与 $Y$ 具有同构的有理函数环；有时这两个概念等价
% P08-E256: frozen source editorial placeholder preserved; see ERRATA_CANDIDATES.jsonl.
（此处插入未来文献）。

\begin{lemma}
\label{spaces-morphisms-lemma-birational}
设 $S$ 为概形。设 $X$ 与 $Y$ 为 $S$ 上的代数空间，且 $|X|$ 与 $|Y|$
不可约。
% P08-E257: frozen source number agreement noted; see ERRATA_CANDIDATES.jsonl.
则 $X$ 与 $Y$ 双有理，当且仅当存在非空开子空间 $U \subset X$ 与
$V \subset Y$，它们作为 $S$ 上代数空间同构。
\end{lemma}

\begin{proof}
假设 $X$ 与 $Y$ 双有理。设 $f : U \to Y$ 与 $g : V \to X$ 分别定义从 $X$
到 $Y$ 以及从 $Y$ 到 $X$ 的互逆支配有理映射。缩小 $U$ 后，可以假设
$f : U \to Y$ 分解经过 $V$。由于 $g \circ f$ 作为支配有理映射是恒等映射，
可知复合 $U \to V \to X$ 在 $U$ 的某个稠密开集上为恒等映射。因此，用更小的
开集替换 $U$ 后，可以假设 $U \to V \to X$ 是 $U$ 到 $X$ 的包含映射。
由对称性，存在开子空间 $V' \subset V$，使得
$g|_{V'} : V' \to X$ 分解经过 $U \subset X$，且 $V' \to U \to Y$
为恒等映射。$|U| \to |V|$ 下 $|V'|$ 的逆像是 $|U|$ 的开集，因而等于某个
开子空间 $U' \subset U$ 的 $|U'|$，见《空间的性质》引理
\ref{spaces-properties-lemma-open-subspaces}。于是 $U' \subset U \to V$
分解为 $U' \to V'$。类似地，$V' \to U$ 分解为 $V' \to U'$。
容易验证，$U' \to V'$ 与 $V' \to U'$ 是 $S$ 上代数空间的互逆态射，
证明完毕。
\end{proof}









\section{代数空间的相对正规化}
\label{spaces-morphisms-section-normalization-X-in-Y}

\noindent
本节是《态射》第 \ref{morphisms-section-normalization-X-in-Y} 节的对应版本。

\begin{lemma}
\label{spaces-morphisms-lemma-integral-closure}
设 $S$ 为概形。设 $X$ 为 $S$ 上的代数空间。设 $\mathcal{A}$ 为
$\mathcal{O}_X$-代数的拟凝聚层。存在 $\mathcal{O}_X$-代数的拟凝聚层
$\mathcal{A}' \subset \mathcal{A}$，使得对 $X_\etale$ 的任意仿射对象 $U$，
环 $\mathcal{A}'(U) \subset \mathcal{A}(U)$ 是 $\mathcal{O}_X(U)$
在 $\mathcal{A}(U)$ 中的整闭包。
\end{lemma}

\begin{proof}
设 $U$ 为 $X_\etale$ 的对象。则 $U$ 是概形。用 $\mathcal{A}|_U$ 表示其在
扎里斯基站点上的限制。于是 $\mathcal{A}|_U$ 是 $\mathcal{O}_U$-代数的
拟凝聚层，故可以应用《态射》引理 \ref{morphisms-lemma-integral-closure}，
找到拟凝聚子代数 $\mathcal{A}'_U \subset \mathcal{A}|_U$，使得
$\mathcal{A}'_U$ 在任意仿射开集 $W \subset U$ 上的值如引理陈述所示。
若 $f : U' \to U$ 是 $X_\etale$ 中的态射，则
$\mathcal{A}|_{U'} = f^*(\mathcal{A}|_U)$，其中 $f^*$ 指在扎里斯基拓扑中
沿态射 $f$ 的拉回；这是因为 $\mathcal{A}$ 拟凝聚（见《空间的性质》第
\ref{spaces-properties-section-quasi-coherent} 节的引言，以及该处指向
概形下降一章相关讨论的文献）。由于 $f$ 为 \'etale，《关于态射的进一步结果》
引理 \ref{more-morphisms-lemma-integral-closure-smooth-pullback} 给出典范同构
$f^*(\mathcal{A}'_U) = \mathcal{A}'_{U'}$。这立即说明，我们在 $X_\etale$
上得到 $\mathcal{O}_X$-代数的子预层
$\mathcal{A}' \subset \mathcal{A}$；它对扎里斯基拓扑成层，并在仿射对象上
具有所需的值。又因为每个 $\mathcal{A}'_U$ 在概形 $U$ 上拟凝聚，且对
\'etale 态射 $f : U' \to U$ 有
$\mathcal{A}'_{U'} = f^*(\mathcal{A}'_U)$，所以 $\mathcal{A}'$ 在
$X_\etale$ 上也拟凝聚（这是局部性质，而上面的文献正是以这种方式描述
$U_\etale$ 上的拟凝聚模）。
\end{proof}

\begin{definition}
\label{spaces-morphisms-definition-integral-closure}
设 $S$ 为概形。设 $X$ 为 $S$ 上的代数空间。设 $\mathcal{A}$ 为
$\mathcal{O}_X$-代数的拟凝聚层。$\mathcal{O}_X$
{\it 在 $\mathcal{A}$ 中的整闭包}是上面引理
\ref{spaces-morphisms-lemma-integral-closure} 中构造的拟凝聚 $\mathcal{O}_X$-子代数
$\mathcal{A}' \subset \mathcal{A}$。
\end{definition}

\noindent
特别地，我们将把它应用于
$\mathcal{A} = f_*\mathcal{O}_Y$ 的情形，其中 $f : Y \to X$
为代数空间的拟紧拟分离态射（见引理 \ref{spaces-morphisms-lemma-pushforward}）。
然后可以取该拟凝聚 $\mathcal{O}_X$-代数的相对谱（引理
\ref{spaces-morphisms-lemma-affine-equivalence-algebras}），得到 $X$ 在 $Y$ 中的正规化。

\begin{definition}
\label{spaces-morphisms-definition-normalization-X-in-Y}
设 $S$ 为概形。设 $f : Y \to X$ 为 $S$ 上代数空间的拟紧拟分离态射。设
$\mathcal{O}'$ 为 $\mathcal{O}_X$ 在 $f_*\mathcal{O}_Y$ 中的整闭包。
{\it $X$ 在 $Y$ 中的正规化}是 $S$ 上代数空间的态射
$$
\nu : X' = \underline{\Spec}_X(\mathcal{O}') \to X
$$
。它配有初始态射 $f$ 的自然分解
$$
Y \xrightarrow{f'} X' \xrightarrow{\nu} X
$$
。
\end{definition}

\noindent
为得到这一分解，使用注 \ref{spaces-morphisms-remark-factorization-quasi-compact-quasi-separated}
以及 $\underline{\Spec}$ 构造的函子性。

\begin{lemma}
\label{spaces-morphisms-lemma-properties-normalization}
设 $S$ 为概形。设 $f : Y \to X$ 为 $S$ 上代数空间的拟紧拟分离态射。设
$Y \to X' \to X$ 为 $X$ 在 $Y$ 中的正规化。
\begin{enumerate}
\item 若 $W \to X$ 为 $S$ 上代数空间的 \'etale 态射，则
$W \times_X X'$ 是 $W$ 在 $W \times_X Y$ 中的正规化。
\item 若 $Y$ 与 $X$ 可表，则 $X'$ 可表，并典范同构于《态射》第
\ref{morphisms-section-normalization} 节所构造的概形 $X$ 在概形 $Y$
中的正规化。
\end{enumerate}
\end{lemma}

\begin{proof}
由构造立即可知，形成 $X$ 在 $Y$ 中的正规化与 \'etale 基变换交换，即 (1)
成立。另一方面，若 $X$ 与 $Y$ 为概形，则对仿射开集 $U \subset X$，
$f_*\mathcal{O}_Y(U) = \mathcal{O}_Y(f^{-1}(U))$；因此，
$\nu^{-1}(U)$ 恰为概形相应构造中所得同一环的谱。
\end{proof}

\noindent
下面给出这一构造的刻画。

\begin{lemma}
\label{spaces-morphisms-lemma-characterize-normalization}
设 $S$ 为概形。设 $f : Y \to X$ 为 $S$ 上代数空间的拟紧拟分离态射。
分解 $f = \nu \circ f'$，其中 $\nu : X' \to X$ 是 $X$ 在 $Y$ 中的正规化，
由下列两个性质刻画：
\begin{enumerate}
\item 态射 $\nu$ 为整态射；
\item 对任意分解 $f = \pi \circ g$，其中 $\pi : Z \to X$ 为整态射，都存在交换图
$$
\xymatrix{
Y \ar[d]_{f'} \ar[r]_g & Z \ar[d]^\pi \\
X' \ar[ru]^h \ar[r]^\nu & X
}
$$
，其中态射 $h : X' \to Z$ 唯一。
\end{enumerate}
此外，在 (2) 中，态射 $h : X' \to Z$ 是 $Z$ 在 $Y$ 中的正规化。
\end{lemma}

\begin{proof}
设 $\mathcal{O}' \subset f_*\mathcal{O}_Y$ 为定义
\ref{spaces-morphisms-definition-normalization-X-in-Y} 中 $\mathcal{O}_X$ 的整闭包。
由构造，态射 $\nu$ 为整态射，这证明了 (1)。假设给定 (2) 中的分解
$f = \pi \circ g$，其中 $\pi : Z \to X$ 为整态射。由定义
\ref{spaces-morphisms-definition-integral}，$\pi$ 仿射；因而 $Z$ 是某个拟凝聚
$\mathcal{O}_X$-代数层 $\mathcal{B}$ 的相对谱。
% P08-E258: frozen source domain of g preserved; see ERRATA_CANDIDATES.jsonl.
态射 $g : X \to Z$ 对应于 $\mathcal{O}_X$-代数映射
$\chi : \mathcal{B} \to f_*\mathcal{O}_Y$。对每个 \'etale 映至 $X$ 的
仿射 $U$，$\mathcal{B}(U)$ 在 $\mathcal{O}_X(U)$ 上整
（由定义 \ref{spaces-morphisms-definition-integral}），故由引理
\ref{spaces-morphisms-lemma-integral-closure} 可知
$\chi(\mathcal{B}) \subset \mathcal{O}'$。由相对谱构造的函子性及引理
\ref{spaces-morphisms-lemma-affine-equivalence-algebras}，这给出唯一态射
$h : X' \to Z$。略去该图交换的验证。

\medskip\noindent
显然，(1)、(2) 刻画分解 $f = \nu \circ f'$，因为它们把该分解刻画为某范畴的
始对象。由引理 \ref{spaces-morphisms-lemma-finite-permanence}，(2) 中的态射 $h$ 为整态射。
给定分解 $g = \pi' \circ g'$，其中 $\pi' : Z' \to Z$ 为整态射，可得分解
$f = (\pi \circ \pi') \circ g'$ 以及态射 $h' : X' \to Z'$。唯一性蕴含
$\pi' \circ h' = h$。因此，刻画 (1)、(2) 可应用于态射
$h : X' \to Z$，从而得到本引理的最后一个陈述。
\end{proof}

\begin{lemma}
\label{spaces-morphisms-lemma-normalization-in-reduced}
设 $S$ 为概形。设 $f : Y \to X$ 为 $S$ 上代数空间的拟紧拟分离态射。
设 $X' \to X$ 为 $X$ 在 $Y$ 中的正规化。若 $Y$ 既约，则 $X'$ 也既约。
\end{lemma}

\begin{proof}
这由既约环的子环仍既约这一事实得出。略去一些细节。
\end{proof}

\begin{lemma}
\label{spaces-morphisms-lemma-normalization-generic}
设 $S$ 为概形。设 $f : Y \to X$ 为概形的拟紧拟分离态射。设 $X' \to X$
为 $X$ 在 $Y$ 中的正规化。若 $x' \in |X'|$ 是余维 $0$ 的点
（《空间的性质》定义
\ref{spaces-properties-definition-dimension-local-ring}），则 $x'$
是某个余维 $0$ 的点 $y \in |Y|$ 之像。
\end{lemma}

\begin{proof}
由引理 \ref{spaces-morphisms-lemma-properties-normalization} 及各定义，可设
$X = \Spec(A)$ 仿射。此时 $X' = \Spec(A')$，其中 $A'$ 是 $A$ 在
$\Gamma(Y, \mathcal{O}_Y)$ 中的整闭包，而 $x'$ 对应于 $A'$ 的极小素理想。
选取满 \'etale 态射 $V \to Y$，其中 $V = \Spec(B)$ 仿射。则
$A' \to B$ 为单射，故 $A'$ 的每个极小素理想都是 $B$ 的某个极小素理想之像，
见《代数》引理 \ref{algebra-lemma-injective-minimal-primes-in-image}。引理得证。
\end{proof}

\begin{lemma}
\label{spaces-morphisms-lemma-normalization-in-disjoint-union}
设 $S$ 为概形。设 $f : Y \to X$ 为 $S$ 上代数空间的拟紧拟分离态射。
假设 $Y = Y_1 \amalg Y_2$ 是两个代数空间的不交并。记
$f_i = f|_{Y_i}$。设 $X_i'$ 为 $X$ 在 $Y_i$ 中的正规化。则
$X_1' \amalg X_2'$ 是 $X$ 在 $Y$ 中的正规化。
\end{lemma}

\begin{proof}
略。
\end{proof}

\begin{lemma}
\label{spaces-morphisms-lemma-normalization-in-universally-closed}
设 $S$ 为概形。
% P08-E259: frozen source number agreement noted; see ERRATA_CANDIDATES.jsonl.
设 $f : X \to Y$ 为 $S$ 上代数空间的拟紧、拟分离且普遍闭的态射。
则 $f_*\mathcal{O}_X$ 在 $\mathcal{O}_Y$ 上整。换言之，$Y$ 在 $X$
中的正规化等于注 \ref{spaces-morphisms-remark-factorization-quasi-compact-quasi-separated}
中的分解
$$
X \longrightarrow \underline{\Spec}_Y(f_*\mathcal{O}_X)
\longrightarrow Y
$$
。
\end{lemma}

\begin{proof}
该问题在 $Y$ 上是 \'etale 局部的，故可化归至 $Y = \Spec(R)$ 仿射的情形。
设 $h \in \Gamma(X, \mathcal{O}_X)$。需要证明 $h$ 满足 $R$ 上的首一方程。
把 $h$ 视为下列交换图中的态射
$$
\xymatrix{
X \ar[rr]_h \ar[rd]_f & & \mathbf{A}^1_Y \ar[ld] \\
& Y &
}
$$
。设 $Z \subset \mathbf{A}^1_Y$ 为 $h$ 的概形论像，见定义
\ref{spaces-morphisms-definition-scheme-theoretic-image}。由于 $f$ 拟紧且
$\mathbf{A}^1_Y \to Y$ 分离，态射 $h$ 拟紧，见引理
\ref{spaces-morphisms-lemma-quasi-compact-permanence}。由引理
\ref{spaces-morphisms-lemma-quasi-compact-scheme-theoretic-image}，态射 $X \to Z$
在底层拓扑空间上有稠密像。由引理
\ref{spaces-morphisms-lemma-universally-closed-permanence}，态射 $X \to Z$ 闭。
因此，在集合论意义下 $h(X) = Z$。于是可以应用引理
\ref{spaces-morphisms-lemma-image-proper-is-proper}，得出 $Z \to Y$ 普遍闭（甚至固有）。
由于 $Z \subset \mathbf{A}^1_Y$，可知 $Z \to Y$ 仿射且固有，故由引理
\ref{spaces-morphisms-lemma-integral-universally-closed} 它为整态射。写成
$\mathbf{A}^1_Y = \Spec(R[T])$，可知 $Z$ 的理想 $I \subset R[T]$
含有首一多项式 $P(T) \in R[T]$。因此 $P(h) = 0$，结论成立。
\end{proof}

\begin{lemma}
\label{spaces-morphisms-lemma-normalization-in-integral}
设 $S$ 为概形。设 $f : Y \to X$ 为 $S$ 上代数空间的整态射。
% P08-E260: frozen source space-level terminology preserved; see ERRATA_CANDIDATES.jsonl.
则 $X$ 在 $Y$ 中的整闭包等于 $Y$。
\end{lemma}

\begin{proof}
由引理 \ref{spaces-morphisms-lemma-integral-universally-closed}，这是引理
\ref{spaces-morphisms-lemma-normalization-in-universally-closed} 的特例。
\end{proof}

\begin{lemma}
\label{spaces-morphisms-lemma-nagata-normalization-finite}
设 $S$ 为概形。设 $f : X \to Y$ 为 $S$ 上代数空间的态射。假设：
\begin{enumerate}
\item $Y$ 是纳伽塔的；
\item $f$ 拟分离且有限型；
\item $X$ 既约。
\end{enumerate}
则 $Y$ 在 $X$ 中的正规化 $\nu : Y' \to Y$ 是有限态射。
\end{lemma}

\begin{proof}
该问题在 $Y$ 上是 \'etale 局部的，见引理
\ref{spaces-morphisms-lemma-properties-normalization}。因此可设 $Y = \Spec(R)$ 仿射。
此时 $R$ 是诺特纳伽塔环，而我们需要证明 $R$ 在
$\Gamma(X, \mathcal{O}_X)$ 中的整闭包是有限 $R$-模。由于 $f$ 拟紧，
可知 $X$ 拟紧。选取仿射概形 $U$ 及满 \'etale 态射 $U \to X$
（《空间的性质》引理
\ref{spaces-properties-lemma-quasi-compact-affine-cover}）。
% P08-E261: frozen source sheaf subscript preserved; see ERRATA_CANDIDATES.jsonl.
则 $\Gamma(X, \mathcal{O}_X) \subset \Gamma(U, \mathcal{O}_X)$。
由于 $R$ 诺特，只需证明 $R$ 在 $\Gamma(U, \mathcal{O}_U)$ 中的整闭包是
有限 $R$-模。由于 $U \to Y$ 有限型，这由《态射》引理
\ref{morphisms-lemma-nagata-normalization-finite} 得出。
\end{proof}






\section{正规化}
\label{spaces-morphisms-section-normalization}

\noindent
本节是《态射》第 \ref{morphisms-section-normalization} 节的对应版本。

\begin{lemma}
\label{spaces-morphisms-lemma-prepare-normalization}
设 $S$ 为概形。设 $X$ 为 $S$ 上的代数空间。下列条件等价：
\begin{enumerate}
\item 存在满 \'etale 态射 $U \to X$，其中 $U$ 为概形，且 $U$ 的每个拟紧开集
只有有限多个不可约分支；
\item 对每个概形 $U$ 及每个 \'etale 态射 $U \to X$，$U$ 的每个拟紧开集
只有有限多个不可约分支；
\item 对每个在 $X$ 上 \'etale 的拟紧代数空间 $Y$，$Y$ 的余维 $0$ 点集
（《空间的性质》定义
\ref{spaces-properties-definition-dimension-local-ring}）有限；
\item 对每个在 $X$ 上 \'etale 的拟紧代数空间 $Y$，空间 $|Y|$
只有有限多个不可约分支。
\end{enumerate}
若 $X$ 可表，这意味着 $X$ 的每个拟紧开集只有有限多个不可约分支。
\end{lemma}

\begin{proof}
(1) 与 (2) 的等价性以及最后一个陈述，由《下降》引理
\ref{descent-lemma-locally-finite-nr-irred-local-fppf} 和《空间的性质》引理
\ref{spaces-properties-lemma-type-property} 得出。只考虑作为概形的 $Y$，
即可明显看出 (4) 蕴含 (1)、(2)。类似地，(3) 蕴含 (1)、(2)，因为对概形而言，
余维 $0$ 的点就是其各不可约分支的一般点；例如见《空间的性质》引理
\ref{spaces-properties-lemma-codimension-0-points}。

\medskip\noindent
反之，假设 (2)，并设 $Y \to X$ 为代数空间的 \'etale 态射，其中 $Y$ 拟紧。
可以选取仿射概形 $V$ 及满 \'etale 态射 $V \to Y$
（《空间的性质》引理
\ref{spaces-properties-lemma-quasi-compact-affine-cover}）。
由 (2)，$V$ 只有有限多个不可约分支；又因 $|V| \to |Y|$ 连续且满，由
《拓扑》引理 \ref{topology-lemma-surjective-continuous-irreducible-components}，
可知 $|Y|$ 只有有限多个不可约分支。故 (4) 成立。类似地，由《空间的性质》引理
\ref{spaces-properties-lemma-codimension-0-points}，$V$ 各不可约分支一般点的
像就是 $Y$ 的余维 $0$ 点，从而这些点只有有限多个，即 (3) 成立。
\end{proof}

\begin{lemma}
\label{spaces-morphisms-lemma-locally-noetherian-can-be-normalized}
设 $S$ 为概形。设 $X$ 为 $S$ 上局部诺特的代数空间。则 $X$ 满足引理
\ref{spaces-morphisms-lemma-prepare-normalization} 中的等价条件。
\end{lemma}

\begin{proof}
若 $U \to X$ 为 \'etale 态射且 $U$ 为概形，则 $U$ 是局部诺特概形，见
《空间的性质》第 \ref{spaces-properties-section-types-properties} 节。
局部诺特概形的不可约分支集局部有限（《除子》引理
\ref{divisors-lemma-components-locally-finite}）。因此，$X$ 满足该引理的
条件 (2)。
\end{proof}

\begin{lemma}
\label{spaces-morphisms-lemma-flat-image-point-codimension-0}
设 $S$ 为概形。设 $f : X \to Y$ 为 $S$ 上代数空间的平坦态射。则对
$x \in |X|$ 有：$x$ 余维 $0$ 于
$X \Rightarrow f(x)$ 余维 $0$ 于 $Y$。
\end{lemma}

\begin{proof}
通过《空间的性质》引理
\ref{spaces-properties-lemma-codimension-0-points} 及 \'etale 局部化，
这化为概形态射及不可约分支一般点的情形。此时，结论由一般化可沿概形的
平坦态射提升得出，见《态射》引理
\ref{morphisms-lemma-generalizations-lift-flat}。
\end{proof}

\begin{lemma}
\label{spaces-morphisms-lemma-flat-locally-finite-type-points-codimension-0}
设 $S$ 为概形。设 $f : X \to Y$ 为 $S$ 上代数空间的态射。假设 $f$
平坦且局部有限型，并假设 $Y$ 满足引理
\ref{spaces-morphisms-lemma-prepare-normalization} 中的等价条件。则 $X$ 也满足引理
\ref{spaces-morphisms-lemma-prepare-normalization} 中的等价条件，且对 $x \in |X|$ 有：
$x$ 余维 $0$ 于 $X \Rightarrow f(x)$ 余维 $0$ 于 $Y$。
\end{lemma}

\begin{proof}
最后一个陈述由引理 \ref{spaces-morphisms-lemma-flat-image-point-codimension-0} 得出。
选取满 \'etale 态射 $V \to Y$，其中 $V$ 为概形；再选取满 \'etale 态射
$U \to X \times_Y V$，其中 $U$ 为概形。只需证明 $U$ 的每个拟紧开集只有
有限多个不可约分支。以下不再特别说明地使用《空间的性质》引理
\ref{spaces-properties-lemma-codimension-0-points} 的结果。由已经证明的内容，
$U$ 的余维 $0$ 点位于
% P08-E262: frozen source repeated U preserved; see ERRATA_CANDIDATES.jsonl.
$U$ 中余维 $0$ 的点上方，而后者依假设局部有限。因此，只需证明：对
余维为 $0$ 的 $v \in V$，概形论纤维
$U_v = U \times_V v$ 的余维 $0$ 点局部有限。这成立是因为 $U_v$ 是
$\kappa(v)$ 上局部有限型的概形，因而局部诺特；例如可应用引理
\ref{spaces-morphisms-lemma-locally-noetherian-can-be-normalized}。
\end{proof}

\begin{lemma}
\label{spaces-morphisms-lemma-normalization}
设 $S$ 为概形。对每个满足引理 \ref{spaces-morphisms-lemma-prepare-normalization}
中等价条件的 $S$ 上代数空间 $X$，都存在代数空间的态射
$$
\nu_X : X^\nu \longrightarrow X
$$
，具有下列性质：
\begin{enumerate}
\item 若 $X$ 满足引理 \ref{spaces-morphisms-lemma-prepare-normalization} 中的等价条件，则
$X^\nu$ 正规且 $\nu_X$ 为整态射；
\item 若 $X$ 为概形，且其每个拟紧开集只有有限多个不可约分支，则
$\nu_X : X^\nu \to X$ 是《态射》第
\ref{morphisms-section-normalization} 节所构造的 $X$ 的正规化；
\item 若 $f : X \to Y$ 为 $S$ 上代数空间的态射，且二者都满足引理
\ref{spaces-morphisms-lemma-prepare-normalization} 中的等价条件，且 $X$ 的每个余维 $0$ 点
都被 $f$ 映至 $Y$ 的余维 $0$ 点，则存在唯一的 $S$ 上代数空间态射
$f^\nu : X^\nu \to Y^\nu$，使得
$\nu_Y \circ f^\nu = f \circ \nu_X$；
\item 若 $f : X \to Y$ 为代数空间的 \'etale 或光滑态射，且 $Y$ 满足引理
\ref{spaces-morphisms-lemma-prepare-normalization} 中的等价条件，则 (3) 的假设成立，且态射
$f^\nu$ 诱导同构 $X^\nu  \to X \times_Y Y^\nu$。
\end{enumerate}
\end{lemma}

\begin{proof}
考虑范畴 $\mathcal{C}$：其对象为 $S$ 上的概形 $U$，使得 $U$ 的每个拟紧开集
只有有限多个不可约分支；其态射为 $S$ 上概形的那些态射
$g : U \to V$，它们把 $U$ 每个不可约分支的一般点映至 $V$
某个不可约分支的一般点。我们已经证明：
\begin{enumerate}
\item[(a)] 对 $U \in \Ob(\mathcal{C})$，有《态射》定义
\ref{morphisms-definition-normalization} 中的正规化态射
$\nu_U : U^\nu \to U$；
\item[(b)] 对 $U \in \Ob(\mathcal{C})$，态射 $\nu_U$ 为整态射，且
$U^\nu$ 为正规概形，见《态射》引理
\ref{morphisms-lemma-normalization-normal}；
\item[(c)] 对每个 $g : U \to V \in \text{Arrows}(\mathcal{C})$，存在唯一态射
$g^\nu : U^\nu \to V^\nu$，使得
$\nu_V \circ g^\nu = g \circ \nu_U$；见《态射》引理
\ref{morphisms-lemma-normalization-normal} 的 (4)，把它应用于复合
% P08-E263: frozen source variables in the composition preserved; see ERRATA_CANDIDATES.jsonl.
$X^\nu \to X \to Y$；
\item[(d)] 若 $V \in \Ob(\mathcal{C})$ 且 $g : U \to V$ 为 \'etale 或光滑态射，
则 $U \in \Ob(\mathcal{C})$，$g \in \text{Arrows}(\mathcal{C})$，且态射
$g^\nu$ 诱导同构 $U^\nu \to U \times_V V^\nu$；见引理
\ref{spaces-morphisms-lemma-flat-locally-finite-type-points-codimension-0} 与
《关于态射的进一步结果》引理
\ref{more-morphisms-lemma-normalization-and-smooth}。
\end{enumerate}
我们的任务是把这一构造推广到 $S$ 上代数空间 $X$ 的相应范畴。

\medskip\noindent
设 $X$ 为满足引理 \ref{spaces-morphisms-lemma-prepare-normalization} 中等价条件的 $S$
上代数空间。设 $U \to X$ 为满 \'etale 态射，其中 $U$ 为概形。置
$R = U \times_X U$，其投影为 $s, t : R \to U$，并置
$j = (t, s) : R \to U \times_S U$，于是 $X = U/R$；见《空间》引理
\ref{spaces-lemma-space-presentation}。由对 $X$ 的假设，$U$ 与 $R$
是 $\mathcal{C}$ 的对象，而态射 $s$、$t$ 是 $S$ 上概形的 \'etale 态射。
由 (a)，有正规化态射 $\nu_U : U^\nu \to U$ 与
$\nu_R : R^\nu \to R$；由 (d)，有态射
$s^\nu : R^\nu \to U^\nu$、$t^\nu : R^\nu \to U^\nu$，它们分别定义同构
$R^\nu \to R \times_{s, U} U^\nu$ 与
$R^\nu \to U^\nu \times_{U, t} R$。因此 $s^\nu$ 与 $t^\nu$ 为 \'etale
态射（因为它们同构于 \'etale 态射的基变换）。所诱导的态射
$j^\nu = (t^\nu, s^\nu) : R^\nu \to U^\nu \times_S U^\nu$
是单态射，因为它等于复合
\begin{align*}
R^\nu
& \to
(U^\nu \times_{U, t} R) \times_R (R \times_{s, U} U^\nu) \\
& =
U^\nu \times_{U, t} R \times_{s, U} U^\nu \\
& \xrightarrow{j}
U^\nu \times_U (U \times_S U) \times_U U^\nu \\
& =
U^\nu \times_S U^\nu
\end{align*}
。第一个箭头是 $\nu_R$ 的对角态射。（这说明 $R^\nu$ 是 $R$ 在 $U^\nu$
上限制的子概形。）利用性质 (d) 对纤维积作形式计算，可知
$R^\nu \times_{s^\nu, U^\nu, t^\nu} R^\nu$ 是
$R \times_{s, U, t} R$ 的正规化。因此，由 (d)，\'etale 态射
$c : R \times_{s, U, t} R \to R$ 唯一延拓为 $c^\nu$。态射 $c^\nu$
与投影
$\text{pr}_{13} : U^\nu \times_S U^\nu \times_S U^\nu \to U^\nu \times_S U^\nu$
相容。类似地，存在态射 $i^\nu : R^\nu \to R^\nu$，它与交换因子的态射
$U^\nu \times_S U^\nu \to U^\nu \times_S U^\nu$ 相容；还存在态射
$e^\nu : U^\nu \to R^\nu$，它与对角态射
$U^\nu \to U^\nu \times_S U^\nu$ 相容。综上可知，
$j^\nu : R^\nu \to U^\nu \times_S U^\nu$ 是 \'etale 等价关系。
此时置 $X^\nu = U^\nu/R^\nu$
（《空间》定理 \ref{spaces-theorem-presentation}）。于是由《群胚》引理
\ref{groupoids-lemma-criterion-fibre-product}，有
$U^\nu = X^\nu \times_X U$。

\medskip\noindent
% P08-E264: frozen source grammar preserved semantically; see ERRATA_CANDIDATES.jsonl.
上一段已经证明：对每个满足引理 \ref{spaces-morphisms-lemma-prepare-normalization}
中等价条件的 $S$ 上代数空间 $X$，若选取满 \'etale 态射
$g : U \to X$，其中 $U$ 为概形，则得到代数空间的笛卡尔图
$$
\xymatrix{
X^\nu \ar[d]_{\nu_X} & U^\nu \ar[l]^{g^\nu} \ar[d]^{\nu_U} \\
X & U \ar[l]_g
}
$$
。这立即推出 $X^\nu$ 是正规代数空间，且 $\nu_X$ 是整态射。这就得到引理的
(1)。

\medskip\noindent
下面将证明态射 $\nu_X : X^\nu \to X$ 在唯一同构的意义下与 $g$ 的选择无关。
眼下，若 $X$ 是概形，就选择 $\text{id} : X \to X$，于是引理的 (2) 显然成立。

\medskip\noindent
还须证明 (3) 与 (4)。令 $g : U \to X$、$\nu_X : X^\nu \to X$ 以及
$g^\nu : U^\nu \to X^\nu$ 如上。设 $Z$ 为正规概形，并设
$h : Z \to U$ 与 $a : Z \to X^\nu$ 为 $S$ 上的态射，使得
$g \circ h = \nu_X \circ a$，且 $Z$ 的每个不可约分支都经由 $h$ 支配 $U$
的某个不可约分支。由《态射》引理
\ref{morphisms-lemma-normalization-normal} 的 (4)，存在唯一态射
$h^\nu : Z \to U^\nu$，使得 $h = \nu_U \circ h^\nu$。图示如下：
$$
\xymatrix{
X^\nu \ar[d]_{\nu_X} &
U^\nu \ar[l]^{g^\nu} \ar[d]^{\nu_U} &
Z \ar[l]^{h^\nu} \ar@/_1em/[ll]_a \ar[dl]^h
\\
X & U \ar[l]_g
}
$$
注意 $a = g^\nu \circ h^\nu$。事实上，四个顶点为 $X^\nu$、$X$、$U^\nu$、
$U$ 的方块是笛卡尔的，因此这立即由给定 $h$ 时 $h^\nu$ 的唯一性得出。
换言之，给定上述 $h : Z \to U$（而不指定 $a$），存在唯一态射
$a : Z \to X^\nu$，使得 $\nu_X \circ a = g \circ h$。

\medskip\noindent
令 $f : X \to Y$ 如引理陈述的 (3) 所示。假设已经选取满 \'etale 态射
$U \to X$ 与 $V \to Y$，其中 $U$ 和 $V$ 为概形，并且 $f$ 可提升为态射
$g : U \to V$。于是 $g \in \text{Arrows}(\mathcal{C})$，并得到与
$\nu_U$ 和 $\nu_V$ 相容的唯一态射 $g^\nu : U^\nu \to V^\nu$。然而，此时两个态射
$$
R^\nu = U^\nu \times_{X^\nu} U^\nu
\to U^\nu \to V^\nu \to Y^\nu
$$
% P08-E265: frozen source spelling "in stead" noted; see ERRATA_CANDIDATES.jsonl.
由上一段的论述必然相同（其中用 $Y$ 代替 $X$）。
由于 $X^\nu$ 是取 $U^\nu$ 关于 $R^\nu$ 的商而构造的，遂得到 (3) 所述的
（唯一）态射 $f^\nu : X^\nu \to Y^\nu$。

\medskip\noindent
为说明 $X^\nu$ 的构造与满 \'etale 态射 $g : U \to X$ 的选择无关，把上一段
的构造应用于 $\text{id} : X \to X$ 以及 $X$ 的两个 \'etale 覆盖之间的态射
$U' \to U$。这已经足够，因为给定 $X$ 的任意两个 \'etale 覆盖，总有第三个覆盖
同时支配二者。读者可验证：分别用 $U' \to X$ 与 $U \to X$ 构造的两个正规化之间
的态射，在基变换到 $U'$ 后成为同构，因而原来就是同构。细节从略。

\medskip\noindent
类似地可证 (4)，这里略去证明；提示：使用上面的 (d)。
\end{proof}

\noindent
由此得到下列定义。

\begin{definition}
\label{spaces-morphisms-definition-normalization}
设 $S$ 为概形。设 $X$ 为满足引理 \ref{spaces-morphisms-lemma-prepare-normalization}
中等价条件的 $S$ 上代数空间。定义 $X$ 的{\it 正规化}为态射
$$
\nu_X : X^\nu \longrightarrow X
$$
，它是在引理 \ref{spaces-morphisms-lemma-normalization} 中构造的。
\end{definition}

\noindent
该定义适用于局部诺特代数空间，见引理
\ref{spaces-morphisms-lemma-locally-noetherian-can-be-normalized}。通常只对约化代数空间定义正规化。
按照上述定义，$X$ 的正规化与 $X$ 的约化 $X_{red}$ 的正规化相同。

\begin{lemma}
\label{spaces-morphisms-lemma-normalization-reduced}
设 $S$ 为概形。设 $X$ 为满足引理 \ref{spaces-morphisms-lemma-prepare-normalization}
中等价条件的 $S$ 上代数空间。正规化态射 $\nu$ 经过约化 $X_{red}$ 分解，且
$X^\nu \to X_{red}$ 是 $X_{red}$ 的正规化。
\end{lemma}

\begin{proof}
这可在 $X$ 上 \'etale 局部检验，因而归结为概形的情形，即《态射》引理
\ref{morphisms-lemma-normalization-reduced}。略去若干细节。
\end{proof}

\begin{lemma}
\label{spaces-morphisms-lemma-normalization-normal}
设 $S$ 为概形。设 $X$ 为满足引理 \ref{spaces-morphisms-lemma-prepare-normalization}
中等价条件的 $S$ 上代数空间。
\begin{enumerate}
\item 正规化 $X^\nu$ 是正规的。
\item 态射 $\nu : X^\nu \to X$ 是满整态射。
\item 映射 $|\nu| : |X^\nu| \to |X|$ 在余维 $0$ 点集之间诱导双射
（《空间的性质》定义
\ref{spaces-properties-definition-dimension-local-ring}）。
\item 设 $Z \to X$ 为态射。假设 $Z$ 是正规代数空间，并且对 $z \in |Z|$ 有：
% P08-E266: frozen source uses undefined f(z); preserved; see ERRATA_CANDIDATES.jsonl.
$z$ 余维 $0$ 于 $Z \Rightarrow f(z)$ 余维 $0$ 于 $X$。则存在唯一分解
$Z \to X^\nu \to X$。
\end{enumerate}
\end{lemma}

\begin{proof}
(1)、(2)、(3) 由概形的相应结果（《态射》引理
\ref{morphisms-lemma-normalization-normal}）以及下述事实得出：概形的一点是某个
不可约分支的一般点，当且仅当其局部环的维数为零（《性质》引理
\ref{properties-lemma-generic-point}）。

\medskip\noindent
设 $Z \to X$ 为 (4) 中的态射。取概形 $U$ 以及满 \'etale 态射 $U \to X$。
选取概形 $V$ 及满 \'etale 态射 $V \to U \times_X Z$。关于余维 $0$ 点的条件保证
$V \to U$ 把 $V$ 各不可约分支的一般点映至 $U$ 各不可约分支的一般点。因此由
《态射》引理 \ref{morphisms-lemma-normalization-normal}，得到唯一分解
$V \to U^\nu \to U$。唯一性保证两个映射
$$
V \times_{U \times_X Z} V \to V \to U^\nu
$$
相同，因为它们都是映射 $V \times_{U \times_X Z} V \to V \to U$ 的唯一分解。
由于代数空间 $U \times_X Z$ 等于商 $V/V \times_{U \times_X Z} V$
（见《空间》第 \ref{spaces-section-presentations} 节），便得到典范态射
$U \times_X Z \to U^\nu$。图示如下：
$$
\xymatrix{
U \times_X Z \ar[r] \ar[d] & U^\nu \ar[r] \ar[d] & U \ar[d] \\
Z \ar@/_/[rr] \ar@{..>}[r] & X^\nu \ar[r] & X
}
$$
为得到虚线箭头，注意态射 $U \times_X Z \to U^\nu$ 的构造关于 \'etale 态射
$U \to X$ 是函子性的（略去精确表述与证明）。因此若置
$R = U \times_X U$，其投影为 $s, t : R \to U$，则得到态射
$R \times_X Z \to R^\nu$，它与 $s, t : R \to U$ 及
$s^\nu, t^\nu : R^\nu \to U^\nu$ 交换。回顾
$X^\nu = U^\nu/R^\nu$，见引理 \ref{spaces-morphisms-lemma-normalization} 的证明。由于
$X = U/R$，一个简单的层论论证表明
$Z = (U \times_X Z)/(R \times_X Z)$。于是态射
$U \times_X Z \to U^\nu$ 与 $R \times_X Z \to R^\nu$ 定义所需的态射
$Z \to X^\nu$。
\end{proof}

\begin{lemma}
\label{spaces-morphisms-lemma-nagata-normalization}
设 $S$ 为概形。设 $X$ 为 $S$ 上的 Nagata 代数空间。正规化
$\nu : X^\nu \to X$ 是有限态射。
\end{lemma}

\begin{proof}
由于 Nagata 代数空间 $X$ 是局部诺特的，定义
\ref{spaces-morphisms-definition-normalization} 适用。根据引理 \ref{spaces-morphisms-lemma-normalization}
中 $X^\nu$ 的构造，立即归结为概形的情形，即《态射》引理
\ref{morphisms-lemma-nagata-normalization}。
\end{proof}












\section{分离且局部拟有限的态射}
\label{spaces-morphisms-section-schemehood}

\noindent
本节证明：在某概形上局部拟有限且分离的代数空间是可表的。这推出分离且局部拟有限
的态射可表（见引理 \ref{spaces-morphisms-lemma-locally-quasi-finite-separated-representable}）。
不过先来一个引理……（它将被命题
\ref{spaces-morphisms-proposition-locally-quasi-finite-separated-over-scheme} 取代）。

\begin{lemma}
\label{spaces-morphisms-lemma-neighbourhood-scheme}
设 $S$ 为概形。考虑 $S$ 上代数空间的交换图
$$
\xymatrix{
V' \ar[r] \ar[rd] & T' \times_T X \ar[r] \ar[d] & X \ar[d] \\
& T' \ar[r] & T
}
$$
。假设
\begin{enumerate}
\item $T' \to T$ 是仿射概形的 \'etale 态射；
\item $X \to T$ 是分离且局部拟有限的态射；
\item $V'$ 是 $T' \times_T X$ 的开子空间；
\item $V' \to T'$ 是拟仿射的。
\end{enumerate}
在此情形下，$V'$ 在 $X$ 中的像 $U$ 是 $X$ 的可表拟紧开子空间。
\end{lemma}

\begin{proof}
先作几个简单观察。由引理 \ref{spaces-morphisms-lemma-quasi-affine-local}，$V'$ 可表。
它也是拟紧的（因为它是仿射概形上的拟仿射概形；见《态射》引理
\ref{morphisms-lemma-quasi-affine-separated}）。由于 $T' \times_T X \to X$
是 \'etale 的（《空间的性质》引理
\ref{spaces-properties-lemma-base-change-etale}），映射
$|T' \times_T X| \to |X|$ 是开的；见《空间的性质》引理
\ref{spaces-properties-lemma-etale-open}。令 $U \subset X$ 为与 $|V'|$ 的像
对应的开子空间；见《空间的性质》引理
\ref{spaces-properties-lemma-open-subspaces}。因为 $|V'|$ 拟紧，$|U|$ 也拟紧，
故由《空间的性质》引理 \ref{spaces-properties-lemma-quasi-compact-space}，
$U$ 是拟紧代数空间。

\medskip\noindent
由《态射》引理
\ref{morphisms-lemma-locally-quasi-finite-qc-source-universally-bounded}，态射
$T' \to T$ 是普遍有界的。因此可对限制 $T' \to T$ 各纤维次数的整数 $n$ 作归纳；
关于 \'etale 态射情形下该整数的描述，见《态射》引理
\ref{morphisms-lemma-etale-universally-bounded}。若 $n = 1$，则 $T' \to T$
是开浸入（见《\'Etale 态射》定理
\ref{etale-theorem-etale-radicial-open}），结论显然。假设 $n > 1$。

\medskip\noindent
考虑仿射概形 $T'' = T' \times_T T'$。由于 $T' \to T$ 是 \'etale 的，有分解
（分解成既开又闭的仿射子概形）$T'' = \Delta(T') \amalg T^*$。事实上，由
《态射》引理 \ref{morphisms-lemma-diagonal-unramified-morphism}，
$\Delta = \Delta_{T'/T}$ 是开的；又因为 $T' \to T$ 作为仿射概形间的态射是
分离的，它也是闭的。第二投影
$\text{pr}_1 : T' \times_T T' \to T'$ 是基变换，其各纤维的次数由 $n$ 限制；
见《态射》引理 \ref{morphisms-lemma-base-change-universally-bounded}。另一方面，
$\text{pr}_1|_{\Delta(T')} : \Delta(T') \to T'$ 是同构，每个纤维恰有一点。
因此应用《态射》引理 \ref{morphisms-lemma-etale-universally-bounded} 可知，限制
$\text{pr}_1|_{T^*} : T^* \to T'$ 各纤维的次数由 $n - 1$ 限制。于是把归纳假设
应用于图
$$
\xymatrix{
p_0^{-1}(V') \cap X^* \ar[r] \ar[rd] &
X^* \ar[r]_{p_1|_{X^*}} \ar[d] &
X' \ar[d] \\
& T^* \ar[r]^{\text{pr}_1|_{T^*}} & T'
}
$$
，得到 $p_1(p_0^{-1}(V') \cap X^*)$ 是拟紧概形。这里置
$X'' = T'' \times_T X$、$X^* = T^* \times_T X$、$X' = T' \times_T X$，而
$p_0, p_1 : X'' \to X'$ 是 $\text{pr}_0, \text{pr}_1$ 的基变换。引理的大多数
假设都通过基变换给出上图的相应假设。例如，
$p_0^{-1}(V') = T'' \times_{T'} V'$ 作为基变换，是 $T''$ 上的拟仿射概形。
略去一些验证。

\medskip\noindent
由《空间的性质》引理 \ref{spaces-properties-lemma-subscheme} 可知
$$
p_1(p_0^{-1}(V')) =
V' \cup p_1(p_0^{-1}(V') \cap X^*)
$$
是拟紧概形。此外，$p_1(p_0^{-1}(V'))$ 显然是证明第一段所讨论的拟紧开子空间
$U \subset X$ 的逆像。换言之，$T' \times_T U$ 是概形！注意
$T' \times_T U$ 拟紧，且在 $T'$ 上分离、局部拟有限，因为
$T' \times_T X \to T'$ 是原态射 $X \to T$ 的基变换，故局部拟有限且分离
（见引理 \ref{spaces-morphisms-lemma-base-change-separated} 与
\ref{spaces-morphisms-lemma-base-change-quasi-finite}）。于是由《关于态射的进一步结果》引理
\ref{more-morphisms-lemma-quasi-finite-separated-quasi-affine}，
$T' \times_T U \to T'$ 是拟仿射的。

\medskip\noindent
由《下降》引理 \ref{descent-lemma-descent-data-sheaves}，这给出
$T' \times_T U / T'$ 上相对于 \'etale 覆盖 $\{T' \to W\}$ 的下降资料，其中
$W \subset T$ 是态射 $T' \to T$ 的像。
% P08-E267: frozen source U' preserved; see ERRATA_CANDIDATES.jsonl.
因为 $U'$ 在 $T'$ 上拟仿射，由《下降》引理
\ref{descent-lemma-quasi-affine} 可知该资料有效；再由《下降》引理
\ref{descent-lemma-descent-data-sheaves} 的最后一部分，推出 $U$ 如所需是概形。
略去一些小细节。
\end{proof}

\begin{proposition}
\label{spaces-morphisms-proposition-locally-quasi-finite-separated-over-scheme}
设 $S$ 为概形。设 $f : X \to T$ 为 $S$ 上代数空间的态射。假设
\begin{enumerate}
\item $T$ 可表；
\item $f$ 局部拟有限；
\item $f$ 分离。
\end{enumerate}
则 $X$ 可表。
\end{proposition}

\begin{proof}
令 $T = \bigcup T_i$ 为概形 $T$ 的仿射开覆盖。若能证明开子空间
$X_i = f^{-1}(T_i)$ 可表，则由《空间的性质》引理
\ref{spaces-properties-lemma-subscheme}，$X$ 可表。注意
$X_i = T_i \times_T X$，且局部拟有限与分离在基变换下都稳定；见引理
\ref{spaces-morphisms-lemma-base-change-separated} 与 \ref{spaces-morphisms-lemma-base-change-quasi-finite}。
因此可假设 $T$ 是仿射概形。

\medskip\noindent
由《空间的性质》引理
\ref{spaces-properties-lemma-union-of-quasi-compact}，存在 Zariski 覆盖
$X = \bigcup X_i$，使每个 $X_i$ 都有来自某仿射概形的满 \'etale 覆盖。
再次由《空间的性质》引理 \ref{spaces-properties-lemma-subscheme}，只需对每个
$X_i$ 证明命题。因此可假设存在仿射概形 $U$ 及满 \'etale 态射 $U \to X$。
这样便归结为下一段的情形。

\medskip\noindent
假设有
$$
U \longrightarrow X \longrightarrow T
$$
，其中 $U$ 与 $T$ 为仿射概形，$U \to X$ 为满 \'etale 态射，而
$X \to T$ 分离且局部拟有限。由引理 \ref{spaces-morphisms-lemma-etale-locally-quasi-finite}
与 \ref{spaces-morphisms-lemma-composition-quasi-finite}，态射 $U \to T$ 局部拟有限。由于
$U$ 与 $T$ 都仿射，它是拟有限的。置 $R = U \times_X U$。则
$X = U/R$；见《空间》引理 \ref{spaces-lemma-space-presentation}。因为
$X \to T$ 分离，态射 $R \to U \times_T U$ 为闭浸入；见引理
\ref{spaces-morphisms-lemma-fibre-product-after-map}。特别地，$R$ 也是仿射概形。由于
$U \to X$ 是 \'etale 的，投影态射 $t, s : R \to U$ 也都是 \'etale 的。
特别地，$s$ 与 $t$ 拟有限、平坦且有限表示（见《态射》引理
\ref{morphisms-lemma-etale-locally-quasi-finite}、
\ref{morphisms-lemma-etale-flat} 与
\ref{morphisms-lemma-etale-locally-finite-presentation}）。

\medskip\noindent
令 $(U, R, s, t, c)$ 为 $U$ 上 \'etale 等价关系 $R$ 所对应的群胚。取一点
$u \in U$，并以 $p \in T$ 记其像。把《关于群胚的进一步结果》引理
\ref{more-groupoids-lemma-quasi-finite-over-base-j-proper} 应用于概形 $T$ 上的群胚
$(U, R, s, t, c)$ 以及上述点 $p$、$u$。由上一段的讨论，该引理的假设
(1)--(7) 全部满足。因此得到 \'etale 邻域 $(T', p') \to (T, p)$ 及不交并分解
$$
U_{T'} = U' \amalg W, \quad
R_{T'} = R' \amalg W'
$$
以及满足前述《关于群胚的进一步结果》引理
\ref{more-groupoids-lemma-quasi-finite-over-base-j-proper} 之结论
(a)、(b)、(c)、(d)、(e)、(f)、(g)、(h) 的 $u' \in U'$。必要时缩小 $T'$ 后，
可以并确实假设 $T'$ 仿射。结论 (h) 蕴含
$R' = U' \times_{X_{T'}} U'$，且投影映射与 $s'$、$t'$ 的限制相同。因此，结论
(g) 中的 $(U', R', s'|_{R'}, t'|_{R'}, c'|_{R' \times_{t', U', s'} R'})$
是 \'etale 等价关系。由《空间》引理 \ref{spaces-lemma-finding-opens} 可知，
$U'/R'$ 是 $X_{T'}$ 的开子空间。由结论 (d)，概形 $U'$、$R'$ 仿射，而态射
$s'|_{R'}, t'|_{R'}$ 为有限 \'etale 态射。因此《群胚》命题
\ref{groupoids-proposition-finite-flat-equivalence} 适用，从而 $U'/R'$ 是仿射概形。

\medskip\noindent
由此得出：对上述每一对点 $(u, p)$，都能找到满足
$\kappa(p) = \kappa(p')$ 的 \'etale 邻域 $(T', p') \to (T, p)$，以及映至 $u$
的一点 $u' \in U_{T'}$，使得 $u'$ 在 $|X_{T'}|$ 中的像 $x'$ 有一个开邻域
$V'$，它是 $X_{T'}$ 中的仿射概形。应用引理 \ref{spaces-morphisms-lemma-neighbourhood-scheme}，得到
作为概形且包含 $x$（即 $u$ 在 $|X|$ 中的像）的开子空间 $W \subset X$。
这对每个 $x$ 都成立，故由《空间的性质》引理
\ref{spaces-properties-lemma-subscheme}，$X$ 是概形。证明完毕。
\end{proof}



\section{应用}
\label{spaces-morphisms-section-applications}

\noindent
下列引理的另一证明，是把它看作代数空间（不可表）态射的 Zariski 主定理之推论；
见《关于空间态射的进一步结果》第
\ref{spaces-more-morphisms-section-structure-quasi-finite} 节。事实上，
《关于空间态射的进一步结果》引理
\ref{spaces-more-morphisms-lemma-quasi-finite-separated-quasi-affine} 蕴含：
代数空间的拟有限分离态射是拟仿射的，因而可表。

\begin{lemma}
\label{spaces-morphisms-lemma-locally-quasi-finite-separated-representable}
设 $S$ 为概形。设 $f : X \to Y$ 为 $S$ 上代数空间的态射。若 $f$ 局部拟有限
且分离，则 $f$ 可表。
\end{lemma}

\begin{proof}
这立即由命题 \ref{spaces-morphisms-proposition-locally-quasi-finite-separated-over-scheme}
以及局部拟有限性和分离性在任意基变换下保持这一事实得出；见引理
\ref{spaces-morphisms-lemma-base-change-quasi-finite} 与 \ref{spaces-morphisms-lemma-base-change-separated}。
\end{proof}

\begin{lemma}
\label{spaces-morphisms-lemma-etale-universally-injective-open}
\begin{slogan}
普遍单射的 \'etale 映射是开浸入。
\end{slogan}
设 $S$ 为概形。设 $f : X \to Y$ 为 $S$ 上代数空间的 \'etale 普遍单射态射。
则 $f$ 是开浸入。
\end{lemma}

\begin{proof}
设 $T \to Y$ 为从概形到 $Y$ 的态射。若能证明 $X \times_Y T \to T$ 是开浸入，
便得结论。由于 \'etale 与普遍单射都是在基变换下稳定的态射性质（见引理
\ref{spaces-morphisms-lemma-base-change-etale} 与
\ref{spaces-morphisms-lemma-base-change-universally-injective}），可假设 $Y$ 是概形。注意由引理
\ref{spaces-morphisms-lemma-universally-injective}，对角态射
$\Delta_{X/Y} : X \to X \times_Y X$ 是 \'etale 的满单态射。因此
$\Delta_{X/Y}$ 是同构（见《空间》引理
\ref{spaces-lemma-surjective-flat-locally-finite-presentation}）；特别地，$X$
在 $Y$ 上分离。再由命题
\ref{spaces-morphisms-proposition-locally-quasi-finite-separated-over-scheme}，$X$ 也是概形。
最后，由概形的 \'etale 态射基本定理，$X \to Y$ 是开浸入；见《\'Etale 态射》
定理 \ref{etale-theorem-etale-radicial-open}。
\end{proof}



\section{Zariski 主定理（可表情形）}
\label{spaces-morphisms-section-Zariski}

\noindent
这是可以利用正规化与 \'etale 局部化交换来证明的版本。要证明更强的版本
（关于不可表示态射），还需发展更多工具。见《关于空间态射的进一步结果》第
\ref{spaces-more-morphisms-section-structure-quasi-finite} 节。

\begin{lemma}
\label{spaces-morphisms-lemma-finite-type-separated}
设 $S$ 为概形。设 $f : X \to Y$ 为 $S$ 上代数空间的可表、有限型、分离态射。
令 $Y'$ 为 $Y$ 在 $X$ 中的正规化。图示如下：
$$
\xymatrix{
X \ar[rd]_f \ar[rr]_{f'} & & Y' \ar[ld]^\nu \\
& Y &
}
$$
则存在开子空间 $U' \subset Y'$，使得
\begin{enumerate}
\item $(f')^{-1}(U') \to U'$ 是同构；
\item $(f')^{-1}(U') \subset X$ 是 $f$ 为拟有限的点集。
\end{enumerate}
\end{lemma}

\begin{proof}
设 $W \to Y$ 为满 \'etale 态射，其中 $W$ 为概形。则 $W \times_Y X$ 也是概形。
由引理 \ref{spaces-morphisms-lemma-properties-normalization}，代数空间 $W \times_Y Y'$ 可表，
并且是概形 $W$ 在概形 $W \times_Y X$ 中的正规化。图示如下：
$$
\xymatrix{
W \times_Y X \ar[rd]_{(1, f)} \ar[rr]_{(1, f')} & &
W \times_Y Y' \ar[ld]^{(1, \nu)} \\
& W &
}
$$
由《关于态射的进一步结果》引理
\ref{more-morphisms-lemma-finite-type-separated}，本引理的结论在 $W$ 上成立。
令 $V' \subset W \times_Y Y'$ 为满足下列条件的开子概形：
\begin{enumerate}
\item $(1, f')^{-1}(V') \to V'$ 是同构；
\item $(1, f')^{-1}(V') \subset W \times_Y X$ 是 $(1, f)$ 为拟有限的点集。
\end{enumerate}
由引理 \ref{spaces-morphisms-lemma-locally-finite-type-quasi-finite-part}，存在使 $f$ 拟有限的最大
开点集 $U \subset X$，且 $W \times_Y U = (1, f')^{-1}(V')$。由引理
\ref{spaces-morphisms-lemma-closed-immersion-local}，态射 $f'|_U : U \to Y'$ 是开浸入，因为它
在 $W$ 上的基变换，是同构 $(1, f')^{-1}(V') \to V'$ 后接开浸入
$V' \to W \times_Y Y'$。置 $U' = \Im(U \to Y')$ 即完成证明
（略去对 $(f')^{-1}(U') = U$ 的验证）。
\end{proof}

\noindent
在下列引理中可以去掉可表的假设，因为已经证明局部拟有限的分离态射可表。

\begin{lemma}
\label{spaces-morphisms-lemma-quasi-finite-separated-quasi-affine}
设 $S$ 为概形。设 $f : X \to Y$ 为 $S$ 上代数空间的态射。假设 $f$ 拟有限
且分离。令 $Y'$ 为 $Y$ 在 $X$ 中的正规化。图示如下：
$$
\xymatrix{
X \ar[rd]_f \ar[rr]_{f'} & & Y' \ar[ld]^\nu \\
& Y &
}
$$
则 $f'$ 是拟紧开浸入，而 $\nu$ 是整态射。特别地，$f$ 拟仿射。
\end{lemma}

\begin{proof}
由引理 \ref{spaces-morphisms-lemma-locally-quasi-finite-separated-representable}，态射 $f$ 可表。
故可应用引理 \ref{spaces-morphisms-lemma-finite-type-separated}。于是存在开子空间
$U' \subset Y'$，使得 $(f')^{-1}(U') = X$（！），且 $X \to U'$ 是同构！
换言之，$f'$ 是开浸入。注意 $f'$ 拟紧，因为 $f$ 拟紧，而
$\nu : Y' \to Y$ 分离（引理 \ref{spaces-morphisms-lemma-quasi-compact-permanence}）。因此，对每个
仿射概形 $Z$ 及态射 $Z \to Y$，纤维积 $Z \times_Y X$ 都是仿射概形
$Z \times_Y Y'$ 的拟紧开子概形。故按定义，$f$ 拟仿射。
\end{proof}








\section{泛同胚}
\label{spaces-morphisms-section-universal-homeomorphisms}

\noindent
概形的泛同胚类在复合与任意基变换下封闭，并且在基底上是 fppf 局部的。见
《态射》引理 \ref{morphisms-lemma-composition-universal-homeomorphism}、
\ref{morphisms-lemma-base-change-universal-homeomorphism}，以及《下降》引理
\ref{descent-lemma-descending-property-universal-homeomorphism}。因此，把第
\ref{spaces-morphisms-section-representable} 节的讨论应用于这一概念，就知道代数空间的可表示态射
是泛同胚意味着什么。

\begin{lemma}
\label{spaces-morphisms-lemma-universal-homeomorphism-representable}
设 $S$ 为概形。设 $f : X \to Y$ 为 $S$ 上代数空间的可表示态射。则 $f$ 是泛同胚
（依第 \ref{spaces-morphisms-section-representable} 节之意义），当且仅当对代数空间的每个态射
$Z \to Y$，基变换映射 $Z \times_Y X \to Z$ 都诱导同胚
$|Z \times_Y X| \to |Z|$。
\end{lemma}

\begin{proof}
若对代数空间的每个态射 $Z \to Y$，基变换映射 $Z \times_Y X \to Z$ 都诱导
同胚 $|Z \times_Y X| \to |Z|$，那么当 $Z$ 是概形时也如此；这形式地蕴含
$f$ 是第 \ref{spaces-morphisms-section-representable} 节意义下的泛同胚。反之，若 $f$ 是第
\ref{spaces-morphisms-section-representable} 节意义下的泛同胚，则由《空间》引理
\ref{spaces-lemma-representable-transformations-property-implication} 与《态射》
引理 \ref{morphisms-lemma-universal-homeomorphism}，$X \to Y$ 是整的、普遍单射
且满。故由引理 \ref{spaces-morphisms-lemma-integral-universally-closed}，$f$ 普遍闭；它又普遍
单射且（普遍）满。也就是说，$f$ 是泛同胚。
\end{proof}

\begin{definition}
\label{spaces-morphisms-definition-universal-homeomorphism}
设 $S$ 为概形。称 $S$ 上代数空间的态射 $f : X \to Y$ 为{\it 泛同胚}，
当且仅当对代数空间的每个态射 $Z \to Y$，基变换 $Z \times_Y X \to Z$
都诱导同胚 $|Z \times_Y X| \to |Z|$。
\end{definition}

\noindent
由引理 \ref{spaces-morphisms-lemma-universal-homeomorphism-representable}，该定义与此前对代数空间
可表示态射所作的定义并不冲突。对代数空间的态射而言，泛同胚未必总是整的。

\begin{example}
\label{spaces-morphisms-example-universal-homeomorphism}
这是备注 \ref{spaces-morphisms-remark-universally-injective-not-separated} 的延续。考虑代数空间
$X = \mathbf{A}^1_k/\{x \sim -x \mid x \not = 0\}$.
存在态射
$$
\mathbf{A}^1_k \longrightarrow X \longrightarrow \mathbf{A}^1_k
$$
，使得第一个箭头是满 \'etale 态射，第二个箭头普遍单射，而复合是映射
$x \mapsto x^2$。因此该复合普遍闭。由此可知，映射
$X \to \mathbf{A}^1_k$ 是泛同胚，但 $X \to \mathbf{A}^1_k$ 不分离。
\end{example}

\noindent
设 $S$ 为概形。设 $f : X \to Y$ 为 $S$ 上代数空间的泛同胚。则 $f$ 普遍闭，
因而拟紧；见引理 \ref{spaces-morphisms-lemma-universally-closed-quasi-compact}。但 $f$ 不必分离
（见上例），甚至不必拟分离：一个例子是取无限维仿射空间
$\mathbf{A}^\infty = \Spec(k[x_1, x_2, \ldots])$ 关于下述等价关系的商：翻转
有限多个非零坐标的符号（细节从略）。

\medskip\noindent
先陈述几个必备引理。

\begin{lemma}
\label{spaces-morphisms-lemma-base-change-universal-homeomorphism}
代数空间的泛同胚沿代数空间的任意态射作基变换后仍是泛同胚。
\end{lemma}

\begin{proof}
这立即由定义得出。
\end{proof}

\begin{lemma}
\label{spaces-morphisms-lemma-composition-universal-homeomorphism}
代数空间的两个泛同胚之复合是泛同胚。
\end{lemma}

\begin{proof}
从略。
\end{proof}

\begin{lemma}
\label{spaces-morphisms-lemma-reduction-universal-homeomorphism}
设 $S$ 为概形。设 $X$ 为 $S$ 上的代数空间。典范闭浸入 $X_{red} \to X$
（见《空间的性质》定义
\ref{spaces-properties-definition-reduced-induced-space}）是泛同胚。
\end{lemma}

\begin{proof}
从略。
\end{proof}

\noindent
由于目前没有更合适的位置，暂把下列结果放在这里。

\begin{lemma}
\label{spaces-morphisms-lemma-integral-universally-injective-push-pull}
设 $S$ 为概形。设 $f : Y \to X$ 为 $S$ 上代数空间的普遍单射整态射。
\begin{enumerate}
\item 函子
$$
f_{small, *} : \Sh(Y_\etale) \longrightarrow \Sh(X_\etale)
$$
是全忠实的；其本质像由 $X_\etale$ 上的那些集合层 $\mathcal{F}$ 构成，它们在
$|X| \setminus f(|Y|)$ 上的限制同构于 $*$；
\item 函子
$$
f_{small, *} : \textit{Ab}(Y_\etale) \longrightarrow \textit{Ab}(X_\etale)
$$
% P08-E268: frozen source target-site mismatch preserved; see ERRATA_CANDIDATES.jsonl.
是全忠实的；其本质像由 $Y_\etale$ 上支撑包含于 $f(|Y|)$ 的那些阿贝尔层构成。
\end{enumerate}
在两种情形下，$f_{small}^{-1}$ 都是函子 $f_{small, *}$ 的左逆。
\end{lemma}

\begin{proof}
由于 $f$ 是整的，它普遍闭（引理 \ref{spaces-morphisms-lemma-integral-universally-closed}）。
特别地，$f(|Y|)$ 是 $|X|$ 的闭子集，所以上述陈述有意义。证明的其余部分与引理
\ref{spaces-morphisms-lemma-closed-immersion-push-pull} 的证明相同，只是使用《\'Etale 上同调》
命题 \ref{etale-cohomology-proposition-integral-universally-injective-pushforward}
来代替《\'Etale 上同调》命题
\ref{etale-cohomology-proposition-closed-immersion-pushforward}。
\end{proof}
% Generated mechanically from the frozen decent-spaces.tex target.
% Do not edit this generated file; edit the bound locale source instead.

% OK, start here.
%

\setcounter{chapter}{67}
\chapter{良好代数空间}



\phantomsection
\label{decent-spaces-section-phantom}


\section{引言}
\label{decent-spaces-section-introduction}

\noindent
本章研究一般代数空间的“局部”性质，也就是研究不一定拟分离的代数空间。
拟分离代数空间见 \cite{Kn}。对于更一般的代数空间，确实会出现本质上
新的现象，尤其涉及点及点的特化。另一方面，代数空间的大多数基本结果
不必顾及这些现象，因此我们把这部分材料单列一章，置于理论的标准展开
之后。



\section{约定}
\label{decent-spaces-section-conventions}

\noindent
固定假设所有概形都包含在一个大 fppf 位点 $\Sch_{fppf}$ 中，并且所考虑的
每个环 $A$ 都满足 $\Spec(A)$（同构于）该大位点的一个对象。

\medskip\noindent
设 $S$ 为概形，$X$ 为 $S$ 上的代数空间。在本章及下一章中，$X$ 与自身
的乘积（在 $S$ 上代数空间的范畴中）记作 $X \times_S X$，而不写作
$X \times X$。



\section{普遍有界纤维}
\label{decent-spaces-section-universally-bounded}

\noindent
我们简要讨论从概形到代数空间的态射具有普遍有界纤维是什么意思。关于
概形态射的类似定义与结果，参见《态射》第
\ref{morphisms-section-universally-bounded} 节。

\begin{definition}
\label{decent-spaces-definition-universally-bounded}
设 $S$ 为概形。设 $X$ 为 $S$ 上的代数空间，$U$ 为 $S$ 上的概形，
$f : U \to X$ 为 $S$ 上的态射。若存在整数 $n$，使得对每个域 $k$ 及
每个态射 $\Spec(k) \to X$，纤维积 $\Spec(k) \times_X U$ 都是 $k$
上的有限概形，且在 $k$ 上的次数 $\leq n$，则称 {\it $f$ 的纤维普遍
有界}\footnote{这大概不是标准术语。}。
\end{definition}

\noindent
% P08-E269: frozen source's undefined Y and mismatched factors preserved; see ERRATA_CANDIDATES.jsonl.
这个定义是有意义的，因为纤维积 $\Spec(k) \times_Y X$ 是概形。此外，
若 $Y$ 是概形，则这里的概念恢复为《态射》定义
\ref{morphisms-definition-universally-bounded} 中的概念；这是《态射》引理
\ref{morphisms-lemma-characterize-universally-bounded} 的结果。

\begin{lemma}
\label{decent-spaces-lemma-composition-universally-bounded}
设 $S$ 为概形，$X$ 为 $S$ 上的代数空间。设 $V \to U$ 为 $S$ 上概形的
态射，$U \to X$ 为从 $U$ 到 $X$ 的态射。若 $V \to U$ 与 $U \to X$
的纤维普遍有界，则 $V \to X$ 的纤维也普遍有界。
\end{lemma}

\begin{proof}
取适用于 $V \to U$ 的整数 $n$，并按定义
\ref{decent-spaces-definition-universally-bounded} 取适用于 $U \to X$ 的整数 $m$。
设 $\Spec(k) \to X$ 为态射，其中 $k$ 为域。考虑态射
$$
\Spec(k) \times_X V
\longrightarrow
\Spec(k) \times_X U
\longrightarrow
\Spec(k).
$$
由假设，概形 $\Spec(k) \times_X U$ 在 $k$ 上有限且次数至多为 $m$，
而整数 $n$ 限制第一个态射各纤维的次数。因此由《态射》引理
\ref{morphisms-lemma-composition-universally-bounded}，
$\Spec(k) \times_X V$ 在 $k$ 上有限且次数至多为 $nm$。
\end{proof}

\begin{lemma}
\label{decent-spaces-lemma-base-change-universally-bounded}
设 $S$ 为概形。设 $Y \to X$ 为 $S$ 上代数空间的可表示态射，
$U \to X$ 为从概形到 $X$ 的态射。若 $U \to X$ 的纤维普遍有界，
则 $U \times_X Y \to Y$ 的纤维普遍有界。
\end{lemma}

\begin{proof}
由定义及纤维积的性质即得。（注意，我们假设 $Y \to X$ 可表示，故
$U \times_X Y$ 是概形，定义确实适用。）
\end{proof}

\begin{lemma}
\label{decent-spaces-lemma-descent-universally-bounded}
设 $S$ 为概形。设 $g : Y \to X$ 为 $S$ 上代数空间的可表示态射，
$f : U \to X$ 为从概形到 $X$ 的态射，并设
$f' : U \times_X Y \to Y$ 为 $f$ 的基变换。若
$$
\Im(|f| : |U| \to |X|) \subset \Im(|g| : |Y| \to |X|)
$$
且 $f'$ 的纤维普遍有界，则 $f$ 的纤维普遍有界。
\end{lemma}

\begin{proof}
按定义 \ref{decent-spaces-definition-universally-bounded} 对态射 $f'$ 取整数
$n \geq 0$，它限制各纤维积 $\Spec(k) \times_Y (U \times_X Y)$ 的
次数。我们断言 $n$ 也适用于 $f$。事实上，设
$x : \Spec(k) \to X$ 是域谱给出的态射。则
$\Spec(k) \times_X U$ 或为空（此时无须证明），或 $x$ 属于 $|f|$ 的像。
由《空间的性质》引理
\ref{spaces-properties-lemma-points-cartesian} 及本引理的假设，这意味着
存在域扩张 $k'/k$ 与交换图表
$$
\xymatrix{
\Spec(k') \ar[r] \ar[d] & Y \ar[d] \\
\Spec(k) \ar[r] & X
}
$$
于是可见
$$
\Spec(k') \times_Y (U \times_X Y) =
\Spec(k') \times_{\Spec(k)} (\Spec(k) \times_X U)
$$
由于按假设概形 $\Spec(k') \times_Y (U \times_X Y)$ 在 $k'$ 上有限且
次数 $\leq n$，所以 $\Spec(k) \times_X U$ 在 $k$ 上也有限且次数
$\leq n$，如所需。（略去若干细节。）
\end{proof}

\begin{lemma}
\label{decent-spaces-lemma-universally-bounded-permanence}
设 $S$ 为概形，$X$ 为 $S$ 上的代数空间。考虑交换图表
$$
\xymatrix{
U \ar[rd]_g \ar[rr]_f & & V \ar[ld]^h \\
& X &
}
$$
其中 $U$ 与 $V$ 为概形。若 $g$ 的纤维普遍有界，且 $f$ 满射且平坦，
则 $h$ 的纤维也普遍有界。
\end{lemma}

\begin{proof}
假设 $g$ 的纤维普遍有界，且 $f$ 满射且平坦。按定义
\ref{decent-spaces-definition-universally-bounded}，设整数 $n \geq 0$ 限制各概形
$\Spec(k) \times_X U$ 的次数。我们断言 $n$ 也适用于 $h$。
设 $\Spec(k) \to X$ 为域谱到 $X$ 的态射。考虑概形态射
% P08-E270: frozen source's reversed base-change arrow preserved; see ERRATA_CANDIDATES.jsonl.
$$
\Spec(k) \times_X V \longrightarrow \Spec(k) \times_X U
$$
它平坦且满射。由假设，左边的概形在 $\Spec(k)$ 上有限且次数
$\leq n$。由《态射》引理
\ref{morphisms-lemma-universally-bounded-permanence}，右边概形的次数也由
$n$ 限制，如所需。
\end{proof}

\begin{lemma}
\label{decent-spaces-lemma-universally-bounded-finite-fibres}
设 $S$ 为概形，$X$ 为 $S$ 上的代数空间，$U$ 为 $S$ 上的概形。
设 $\varphi : U \to X$ 为 $S$ 上的态射。若 $\varphi$ 的纤维普遍有界，
则存在整数 $n$，使得 $|U| \to |X|$ 的每个纤维至多含 $n$ 个元素。
\end{lemma}

\begin{proof}
定义 \ref{decent-spaces-definition-universally-bounded} 中的整数 $n$ 即可。事实上，
取 $x \in |X|$，并用态射 $x : \Spec(k) \to X$ 表示 $x$。于是得到交换图表
$$
\xymatrix{
\Spec(k) \times_X U \ar[r] \ar[d] & U \ar[d] \\
\Spec(k) \ar[r]^x & X
}
$$
由《空间的性质》引理
\ref{spaces-properties-lemma-points-cartesian}，该图表表明 $x$ 在 $|U|$
中的逆像就是上方水平箭头的像。由于 $\Spec(k) \times_X U$ 在 $k$ 上
有限且次数 $\leq n$，它至多含 $n$ 个点。
\end{proof}








\section{有限性条件与点}
\label{decent-spaces-section-points-monomorphisms}

\noindent
本节进一步讨论何时能用从域谱到该空间的单态射表示点。

\begin{remark}
\label{decent-spaces-remark-recall}
在证明下一个引理之前，先回顾概形的 \'etale 态射的一些事实：
\begin{enumerate}
\item \'Etale 态射平坦，故一般化可沿 \'etale 态射提升
（《态射》引理 \ref{morphisms-lemma-etale-flat} 与
\ref{morphisms-lemma-generalizations-lift-flat}）。
\item \'Etale 态射非分歧，而非分歧态射局部拟有限，故其纤维离散
（《态射》引理 \ref{morphisms-lemma-flat-unramified-etale}、
\ref{morphisms-lemma-unramified-quasi-finite} 与
\ref{morphisms-lemma-quasi-finite-at-point-characterize}）。
\item 拟紧 \'etale 态射拟有限，因而特别具有有限纤维
（《态射》引理 \ref{morphisms-lemma-quasi-finite-locally-quasi-compact} 与
\ref{morphisms-lemma-quasi-finite}）。
\item 域 $k$ 上的 \'etale 概形是 $k$ 的有限可分域扩张之谱的不交并
（《态射》引理 \ref{morphisms-lemma-etale-over-field}）。
\end{enumerate}
关于 \'etale 态射的一般讨论，参见《\'Etale 态射》第
\ref{etale-section-etale-morphisms} 节。
\end{remark}

\begin{lemma}
\label{decent-spaces-lemma-U-finite-above-x}
设 $S$ 为概形，$X$ 为 $S$ 上的代数空间，并设 $x \in |X|$。下列条件等价：
\begin{enumerate}
\item 存在概形族 $U_i$ 及 \'etale 态射 $\varphi_i : U_i \to X$，使得
$\coprod \varphi_i : \coprod U_i \to X$ 满射，并且对每个 $i$，
$|U_i| \to |X|$ 在 $x$ 上的纤维有限；
\item 对每个仿射概形 $U$ 及 \'etale 态射 $\varphi : U \to X$，
$|U| \to |X|$ 在 $x$ 上的纤维有限。
\end{enumerate}
\end{lemma}

\begin{proof}
(2) $\Rightarrow$ (1) 显然。设 $\varphi_i : U_i \to X$ 是 (1) 中的
\'etale 态射族，并设 $\varphi : U \to X$ 是从仿射概形到 $X$ 的
\'etale 态射。考虑纤维积图表
$$
\xymatrix{
U \times_X U_i \ar[r]_-{p_i} \ar[d]_{q_i} & U_i \ar[d]^{\varphi_i} \\
U \ar[r]^\varphi & X
}
\quad \quad
\xymatrix{
\coprod U \times_X U_i \ar[r]_-{\coprod p_i} \ar[d]_{\coprod q_i} &
\coprod U_i \ar[d]^{\coprod \varphi_i} \\
U \ar[r]^\varphi & X
}
$$
因为 $q_i$ 是 \'etale 态射，所以它是开态射（见注
\ref{decent-spaces-remark-recall}）。而且态射 $\coprod q_i$ 满射。因此存在有限多个指标
$i_1, \ldots, i_n$ 以及拟紧开子集
$W_{i_j} \subset U \times_X U_{i_j}$，它们的像覆盖 $U$。
态射 $p_i$ 是 \'etale 的，因而局部拟有限（见上面对 \'etale 态射的注）。
故可应用《态射》引理
\ref{morphisms-lemma-locally-quasi-finite-qc-source-universally-bounded}，
得知下列态射的纤维有限：
% P08-E271: frozen source's codomain U_i is preserved; see ERRATA_CANDIDATES.jsonl.
$p_{i_j}|_{W_{i_j}} : W_{i_j} \to U_i$。
于是由《空间的性质》引理
\ref{spaces-properties-lemma-points-cartesian} 及关于 $\varphi_i$ 的假设，
可知 $\varphi$ 在 $x$ 上的纤维有限。换言之，(2) 成立。
\end{proof}

\begin{lemma}
\label{decent-spaces-lemma-R-finite-above-x}
设 $S$ 为概形，$X$ 为 $S$ 上的代数空间，并设 $x \in |X|$。下列条件等价：
\begin{enumerate}
\item 存在概形 $U$、\'etale 态射 $\varphi : U \to X$ 以及映到 $x$ 的点
$u, u' \in U$，使得令 $R = U \times_X U$ 后，态射
$$
|R| \to |U| \times_{|X|} |U|
$$
在 $(u, u')$ 上的纤维有限；
\item 对每个概形 $U$、\'etale 态射 $\varphi : U \to X$ 以及任意映到
$x$ 的点 $u, u' \in U$，令 $R = U \times_X U$ 后，态射
$$
|R| \to |U| \times_{|X|} |U|
$$
在 $(u, u')$ 上的纤维有限；
\item 存在 $x$ 的等价类中的态射 $\Spec(k) \to X$，其中 $k$ 为域，
使得投影 $\Spec(k) \times_X \Spec(k) \to \Spec(k)$ 是
\'etale 且拟紧的；
\item 存在 $x$ 的等价类中的单态射 $\Spec(k) \to X$，其中 $k$ 为域。
\end{enumerate}
\end{lemma}

\begin{proof}
假设 (1) 成立，即设 $\varphi : U \to X$ 是从概形到 $X$ 的 \'etale 态射，
$u, u'$ 是 $U$ 中位于 $x$ 上方的点，且
$|R| \to |U| \times_{|X|} |U|$ 在 $(u, u')$ 上的纤维是有限集。
在本证明中，把点 $u = \Spec(\kappa(u))$ 视为概形。注意，$u \to U$ 与
$u' \to U$ 都是单态射（见《概形》引理
\ref{schemes-lemma-injective-points-surjective-stalks}），故
$u \times_X u' \to R = U \times_X U$ 是单态射。采用这种语言，假设实际
意味着 $u \times_X u'$ 是一个底拓扑空间只有有限多个点的概形。
设 $\psi : W \to X$ 是从概形到 $X$ 的 \'etale 态射，并设
$w, w' \in W$ 是 $W$ 中映到 $x$ 的点。我们须证明 $w \times_X w'$ 是一个底拓扑
空间只有有限多个点的概形。考虑纤维积图表
$$
\xymatrix{
W \times_X U \ar[r]_p \ar[d]_q & U \ar[d]^\varphi \\
W \ar[r]^\psi & X
}
$$
因为 $x$ 是 $u$ 与 $u'$ 的像，所以可在 $W \times_X U$ 中取点
$\tilde w, \tilde w'$，使得 $q(\tilde w) = w$、
$q(\tilde w') = w'$、$u = p(\tilde w)$ 且 $u' = p(\tilde w')$；见
《空间的性质》引理 \ref{spaces-properties-lemma-points-cartesian}。
因为 $p$、$q$ 是 \'etale 态射，域扩张
$\kappa(w) \subset \kappa(\tilde w) \supset \kappa(u)$ 与
$\kappa(w') \subset \kappa(\tilde w') \supset \kappa(u')$
中的各个扩张均有限可分；见注 \ref{decent-spaces-remark-recall}。于是得到交换图表
% P08-E272: frozen source's lower-left X-product is preserved; see ERRATA_CANDIDATES.jsonl.
$$
\xymatrix{
w \times_X w' \ar[d] &
\tilde w \times_X \tilde w' \ar[l] \ar[d] \ar[r] &
u \times_X u' \ar[d] \\
w \times_X w' &
\tilde w \times_S \tilde w' \ar[l] \ar[r] &
u \times_S u'
}
$$
其中各方形都是纤维积方形。下方水平态射是 \'etale 且拟紧的：这是因为
形如 $\Spec(k) \times_S \Spec(k')$ 的概形都是仿射概形，并结合上面对域扩张
的观察可得。因此，上方水平箭头也是 \'etale 且拟紧的，故具有有限纤维。
上面已说明 $|u \times_X u'|$ 有限，因而 $|w \times_X w'|$ 有限。
换言之，(2) 成立。

\medskip\noindent
假设 (2) 成立。设 $U \to X$ 是从概形 $U$ 出发的 \'etale 态射，且 $x$
属于 $|U| \to |X|$ 的像。取映到 $x$ 的点 $u \in U$。上一段表明，
$u = \Spec(\kappa(u)) \to X$ 具有如下性质：$u \times_X u$ 的底拓扑空间
有限。另一方面，投影 $u \times_X u \to u$ 是下列复合：
$$
u \times_X u \longrightarrow
u \times_X U \longrightarrow
u \times_X X = u,
$$
也就是一个单态射（单态射 $u \to U$ 的基变换）与一个 \'etale 态射
（\'etale 态射 $U \to X$ 的基变换）的复合。因此，$u \times_X U$ 是
$\kappa(u)$ 的有限可分域扩张之谱的不交并（见注
\ref{decent-spaces-remark-recall}）。因为 $u \times_X u$ 有限，它在 $u \times_X U$
中的像是 $\kappa(u)$ 的有限可分域扩张之谱的有限不交并。由《概形》引理
\ref{schemes-lemma-mono-towards-spec-field}，可知 $u \times_X u$ 本身也是
$\kappa(u)$ 的有限可分域扩张之谱的有限不交并。换言之，
$u \times_X u \to u$ 拟紧且 \'etale。因此 (3) 成立。

\medskip\noindent
下面证明 (3) 蕴含 (4)。设 $\Spec(k) \to X$ 是 $x$ 的等价类中从域谱到
$X$ 的态射，使得两个投影
$t, s : R = \Spec(k) \times_X \Spec(k)  \to \Spec(k)$
拟紧且 \'etale。特别地，这意味着 $R$ 是 $\Spec(k)$ 上的 \'etale 等价关系。
由《空间》定理 \ref{spaces-theorem-presentation}，商层
$X' = \Spec(k)/R$ 是代数空间。由《群胚》引理
\ref{groupoids-lemma-quotient-groupoid-restrict}，映射 $X' \to X$ 是单态射。
由于 $s,t$ 拟紧，$R$ 拟紧，因而《空间的性质》引理
\ref{spaces-properties-lemma-point-like-spaces} 适用于 $X'$；于是对于某个域
$k'$，有 $X' = \Spec(k')$。故得到分解
$$
\Spec(k) \longrightarrow
\Spec(k') \longrightarrow X
$$
它表明 $\Spec(k') \to X$ 是映到 $x \in |X|$ 的单态射。换言之，(4) 成立。

\medskip\noindent
最后证明 (4) 蕴含 (1)。设 $\Spec(k) \to X$ 是 $x$ 的等价类中的单态射，
其中 $k$ 为域。设 $U \to X$ 是从概形 $U$ 到 $X$ 的满 \'etale 态射，
并设 $u \in U$ 是 $x$ 上方的点。由于 $\Spec(k) \times_X u$ 非空，且
$\Spec(k) \times_X u \to u$ 是单态射，由《概形》引理
\ref{schemes-lemma-mono-towards-spec-field} 可知
$\Spec(k) \times_X u = u$。因此 $u \to U \to X$ 经过
$\Spec(k) \to X$ 分解，如下图所示：
$$
\xymatrix{
u \ar[r] \ar[d] & U \ar[d] \\
\Spec(k) \ar[r] & X
}
$$
因为右边的竖直箭头是 \'etale 态射，这蕴含 $\kappa(u)/k$ 是有限可分扩张。
于是
$$
u \times_X u = u \times_{\Spec(k)} u
$$
是有限概形，再由本证明第一段对条件 (1) 含义的讨论即得结论。
\end{proof}

\begin{lemma}
\label{decent-spaces-lemma-weak-UR-finite-above-x}
设 $S$ 为概形，$X$ 为 $S$ 上的代数空间，并设 $x \in |X|$。
设 $U$ 为概形，$\varphi : U \to X$ 为 \'etale 态射。下列条件等价：
\begin{enumerate}
\item $x$ 属于 $|U| \to |X|$ 的像，并且令 $R = U \times_X U$ 后，下列
两个映射在 $x$ 上的纤维都有限：
$$
|U| \longrightarrow |X|
\quad\text{且}\quad
|R| \longrightarrow |X|
$$
\item 存在 $x$ 的等价类中的单态射 $\Spec(k) \to X$，其中 $k$ 为域，
并且纤维积 $\Spec(k) \times_X U$ 是 $k$ 上的非空有限概形。
\end{enumerate}
\end{lemma}

\begin{proof}
假设 (1) 成立。它显然蕴含引理 \ref{decent-spaces-lemma-R-finite-above-x} 的第一个条件，
因而得到 $x$ 的等价类中的单态射 $\Spec(k) \to X$。取纤维积可见，
$\Spec(k) \times_X U \to \Spec(k)$ 是 $\Spec(k)$ 上只有有限多个点的
\'etale 概形，因此是 $k$ 上的非空有限概形，即 (2) 成立。

\medskip\noindent
假设 (2) 成立。由假设，$x$ 属于 $|U| \to |X|$ 的像。由《空间的性质》引理
\ref{spaces-properties-lemma-points-cartesian}，$|U| \to |X|$ 在 $x$ 上的
纤维等于 $|\Spec(k) \times_X U|$，故它显然有限。
$|R| \to |X|$ 在 $x$ 上的纤维也显然有限，因为它等于下列概形的底集合：
$$
(\Spec(k) \times_X U) \times_{\Spec(k)} (\Spec(k) \times_X U)
$$
而该概形在 $k$ 上有限。因此 (1) 成立。
\end{proof}

\begin{lemma}
\label{decent-spaces-lemma-UR-finite-above-x}
设 $S$ 为概形，$X$ 为 $S$ 上的代数空间，并设 $x \in |X|$。下列条件等价：
\begin{enumerate}
\item 对每个仿射概形 $U$ 及任意 \'etale 态射 $\varphi : U \to X$，令
$R = U \times_X U$ 后，下列两个映射在 $x$ 上的纤维都有限：
$$
|U| \longrightarrow |X|
\quad\text{且}\quad
|R| \longrightarrow |X|
$$
\item 存在概形 $U_i$ 及 \'etale 态射 $U_i \to X$，使得
$\coprod U_i \to X$ 满射，并且对每个 $i$，令
$R_i = U_i \times_X U_i$ 后，下列两个映射在 $x$ 上的纤维都有限：
$$
|U_i| \longrightarrow |X|
\quad\text{且}\quad
|R_i| \longrightarrow |X|
$$
\item 存在 $x$ 的等价类中的单态射 $\Spec(k) \to X$，其中 $k$ 为域，
并且对任意仿射概形 $U$ 及 \'etale 态射 $U \to X$，纤维积
$\Spec(k) \times_X U$ 是 $k$ 上的有限概形；
\item 存在 $x$ 的等价类中的拟紧单态射 $\Spec(k) \to X$，其中 $k$ 为域；
\item 存在 $x$ 的等价类中的拟紧态射 $\Spec(k) \to X$，其中 $k$ 为域；
\item 每个 $x$ 的等价类中的态射 $\Spec(k) \to X$ 都拟紧，其中 $k$ 为域。
\end{enumerate}
\end{lemma}

\begin{proof}
对每个从仿射概形 $U$ 出发的 \'etale 态射 $U \to X$ 应用引理
\ref{decent-spaces-lemma-weak-UR-finite-above-x}，即得 (1) 与 (3) 等价。
(3) 蕴含 (2) 是显然的。假设 $U_i \to X$ 与 $R_i$ 如 (2) 所述。
由引理 \ref{decent-spaces-lemma-U-finite-above-x}，对任意仿射概形 $U$ 及 \'etale 态射
$U \to X$，$|U| \to |X|$ 在 $x$ 上的纤维有限。将此纤维记为
$\{u_1, \ldots, u_n\}$。取某个使 $x$ 属于 $|U_i| \to |X|$ 之像的 $i$，
则引理 \ref{decent-spaces-lemma-R-finite-above-x} 的 (1) 适用于 $U_i \to X$，故可知
$|R = U \times_X U| \to |U| \times_{|X|} |U|$ 在每个
$(u_a, u_b)$ 上的纤维有限，其中 $a, b \in \{1, \ldots, n\}$。
因此 $|R| \to |X|$ 在 $x$ 上的纤维有限。由此 (1) 成立。至此已知
(1)、(2)、(3) 等价。

\medskip\noindent
若 (4) 成立，则对任意仿射概形 $U$ 及 \'etale 态射 $U \to X$，概形
$\Spec(k) \times_X U$ 一方面在 $k$ 上是 \'etale 的（故由注
\ref{decent-spaces-remark-recall}，它是 $k$ 的有限可分扩张之谱的不交并），另一方面在
$U$ 上拟紧（因而拟紧）。由此 (3) 成立。反过来，若 $U_i \to X$ 如 (2)
所述，且 $\Spec(k) \to X$ 是 (3) 中的单态射，则
$$
\coprod \Spec(k) \times_X U_i
\longrightarrow
\coprod U_i
$$
拟紧（因为在每个 $U_i$ 上，$\Spec(k) \times_X U_i$ 是域谱的有限不交并）。
故由《空间的态射》引理 \ref{spaces-morphisms-lemma-quasi-compact-local}，
$\Spec(k) \to X$ 拟紧。

\medskip\noindent
(4) 蕴含 (5) 是立即的。反过来，设 $\Spec(k) \to X$ 是 $x$ 的等价类中的
拟紧态射。设 $U \to X$ 是 \'etale 态射且 $U$ 仿射。考虑纤维积
$$
\xymatrix{
F \ar[r] \ar[d] & U \ar[d] \\
\Spec(k) \ar[r] & X
}
$$
于是 $F \to U$ 拟紧，故 $F$ 拟紧。另一方面，$F \to \Spec(k)$ 是
\'etale 态射，所以 $F$ 是 $k$ 的有限可分扩张之谱的有限不交并（注
\ref{decent-spaces-remark-recall}）。因为 $|F| \to |U|$ 的像就是 $|U| \to |X|$ 在
$x$ 上的纤维（《空间的性质》引理
\ref{spaces-properties-lemma-points-cartesian}），所以 $|U| \to |X|$ 在
$x$ 上的纤维有限。概形
$F \times_{\Spec(k)} F$ 也拟紧且在 $\Spec(k)$ 上 \'etale，故同样是域谱
的有限并。存在单态射
$F \times_X F \to F \times_{\Spec(k)} F$，因此由《概形》引理
\ref{schemes-lemma-mono-towards-spec-field}，$F \times_X F$ 是域谱的
有限不交并。于是 $F \times_X F \to U \times_X U = R$ 的像有限。
由《空间的性质》引理 \ref{spaces-properties-lemma-points-cartesian}，
该像就是 $|R| \to |X|$ 在 $x$ 上的纤维，故 (1) 成立。至此已知
(1)--(5) 等价。

\medskip\noindent
(6) 蕴含 (5) 是显然的。反过来，设 $\Spec(k) \to X$ 如 (4) 所述，
并设 $\Spec(k') \to X$ 是另一个 $x$ 的等价类中的态射，其中 $k'$ 为域。
由《空间的性质》引理
\ref{spaces-properties-lemma-equivalence-class-point-monomorphism}，给定态射
具有分解 $\Spec(k') \to \Spec(k) \to X$。这是拟紧态射的复合，因而由
《空间的态射》引理 \ref{spaces-morphisms-lemma-composition-quasi-compact}
仍拟紧，如所需。
\end{proof}

\begin{lemma}
\label{decent-spaces-lemma-U-universally-bounded}
设 $S$ 为概形，$X$ 为 $S$ 上的代数空间。下列条件等价：
\begin{enumerate}
\item 存在概形 $U_i$ 及 \'etale 态射 $U_i \to X$，使得
$\coprod U_i \to X$ 满射，并且每个 $U_i \to X$ 的纤维普遍有界；
\item 对每个仿射概形 $U$ 及 \'etale 态射 $\varphi : U \to X$，
$U \to X$ 的纤维普遍有界。
\end{enumerate}
\end{lemma}

\begin{proof}
(2) $\Rightarrow$ (1) 显然。假设 (1) 成立。设
$(\varphi_i : U_i \to X)_{i \in I}$ 是从概形到 $X$ 的 \'etale 态射族，
它覆盖 $X$，并且每个 $\varphi_i$ 的纤维普遍有界。设
$\psi : U \to X$ 是从仿射概形到 $X$ 的 \'etale 态射。对每个 $i$，
考虑纤维积图表
$$
\xymatrix{
U \times_X U_i \ar[r]_{p_i} \ar[d]_{q_i} & U_i \ar[d]^{\varphi_i} \\
U \ar[r]^\psi & X
}
$$
因为 $q_i$ 是 \'etale 态射，所以它是开态射（见注
\ref{decent-spaces-remark-recall}）。而且由于态射族 $(\varphi_i)_{i \in I}$ 满射，
有 $U = \bigcup \Im(q_i)$。由于 $U$ 仿射，因而拟紧，所以可取有限多个
% P08-E273: frozen source says "finite finitely"; intended "find finitely" is rendered; see ERRATA_CANDIDATES.jsonl.
$i_1, \ldots, i_n \in I$ 以及拟紧开子集
$W_j \subset U \times_X U_{i_j}$，使得
% P08-E274: frozen source uses p although the covering map to U is q; formula preserved; see ERRATA_CANDIDATES.jsonl.
$U = \bigcup p_{i_j}(W_j)$.
态射 $p_{i_j}$ 是 \'etale 的，因而局部拟有限（见上面对 \'etale 态射的注）。
故可应用《态射》引理
\ref{morphisms-lemma-locally-quasi-finite-qc-source-universally-bounded}
得知 $p_{i_j}|_{W_j} : W_j \to U_{i_j}$ 的纤维普遍有界。因此由引理
\ref{decent-spaces-lemma-composition-universally-bounded}，$W_j \to X$ 的纤维普遍有界，
从而 $\coprod_{j = 1, \ldots, n} W_j \to X$ 的纤维也普遍有界。
由于 $\coprod_{j = 1, \ldots, n} W_j \to X$ 经过下列满 \'etale 态射分解：
$\coprod q_{i_j}|_{W_j} : \coprod_{j = 1, \ldots, n} W_j \to U$，
引理 \ref{decent-spaces-lemma-universally-bounded-permanence} 表明 $U \to X$ 的纤维
普遍有界。换言之，(2) 成立。
\end{proof}

\begin{lemma}
\label{decent-spaces-lemma-characterize-very-reasonable}
设 $S$ 为概形，$X$ 为 $S$ 上的代数空间。下列条件等价：
\begin{enumerate}
\item 存在 Zariski 覆盖 $X = \bigcup X_i$，并且对每个 $i$，存在概形
$U_i$ 及拟紧满 \'etale 态射 $U_i \to X_i$；
\item 存在概形 $U_i$ 及 \'etale 态射 $U_i \to X$，使得投影
$U_i \times_X U_i \to U_i$ 拟紧，且 $\coprod U_i \to X$ 满射。
\end{enumerate}
\end{lemma}

\begin{proof}
若 (1) 成立，则态射 $U_i \to X_i \to X$ 是 \'etale 的；这是综合《态射》引理
\ref{morphisms-lemma-composition-etale} 与《空间》引理
\ref{spaces-lemma-composition-representable-transformations-property} 与
\ref{spaces-lemma-morphism-schemes-gives-representable-transformation-property}
所得。而且因为 $U_i \times_X U_i = U_i \times_{X_i} U_i$，两个投影
$U_i \times_X U_i \to U_i$ 都拟紧。

\medskip\noindent
若 (2) 成立，令 $X_i \subset X$ 为与开映射 $|U_i| \to |X|$ 的像对应的
开子空间；见《空间的性质》引理
\ref{spaces-properties-lemma-etale-image-open}。态射 $U_i \to X_i$ 满射，
故 $U_i \to X_i$ 是满 \'etale 态射。投影 $U_i \times_{X_i} U_i \to U_i$ 拟紧，
因为 $U_i \times_{X_i} U_i = U_i \times_X U_i$。于是由《空间》引理
\ref{spaces-lemma-representable-morphisms-spaces-property}，态射
$U_i \to X_i$ 拟紧。
\end{proof}









\section{代数空间上的条件}
\label{decent-spaces-section-conditions}

\noindent
本节讨论上文所见代数空间的各种自然条件之间的关系。关于这些条件的含义，
可先阅读第 \ref{decent-spaces-section-reasonable-decent} 节以形成直观认识。

\begin{lemma}
\label{decent-spaces-lemma-bounded-fibres}
设 $S$ 为概形，$X$ 为 $S$ 上的代数空间。考虑 $X$ 的下列条件：
\begin{itemize}
\item[] $(\alpha)$ 对每个 $x \in |X|$，引理
\ref{decent-spaces-lemma-U-finite-above-x} 的等价条件成立。
\item[] $(\beta)$ 对每个 $x \in |X|$，引理
\ref{decent-spaces-lemma-R-finite-above-x} 的等价条件成立。
\item[] $(\gamma)$ 对每个 $x \in |X|$，引理
\ref{decent-spaces-lemma-UR-finite-above-x} 的等价条件成立。
\item[] $(\delta)$ 引理 \ref{decent-spaces-lemma-U-universally-bounded} 的等价条件成立。
\item[] $(\epsilon)$ 引理
\ref{decent-spaces-lemma-characterize-very-reasonable} 的等价条件成立。
\item[] $(\zeta)$ 空间 $X$ 在 Zariski 局部拟分离。
\item[] $(\eta)$ 空间 $X$ 拟分离。
\item[] $(\theta)$ 空间 $X$ 可表示，即 $X$ 是概形。
\item[] $(\iota)$ 空间 $X$ 是拟分离概形。
\end{itemize}
有
$$
\xymatrix{
& (\theta) \ar@{=>}[rd] & & & &  \\
(\iota) \ar@{=>}[ru] \ar@{=>}[rd] & &
(\zeta) \ar@{=>}[r] &
(\epsilon) \ar@{=>}[r] &
(\delta) \ar@{=>}[r] &
(\gamma) \ar@{<=>}[r] & (\alpha) + (\beta) \\
& (\eta) \ar@{=>}[ru] & & & &
}
$$
\end{lemma}

\begin{proof}
蕴含关系 $(\gamma) \Leftrightarrow (\alpha) + (\beta)$ 是立即的。
左侧菱形中的蕴含关系由定义显然成立。

\medskip\noindent
假设 $(\zeta)$ 成立，即 $X$ 在 Zariski 局部拟分离。则由《空间的性质》引理
\ref{spaces-properties-lemma-quasi-separated-quasi-compact-pieces}.
可知 $(\epsilon)$ 成立。

\medskip\noindent
假设 $(\epsilon)$ 成立。由引理
\ref{decent-spaces-lemma-characterize-very-reasonable}，存在 Zariski 开覆盖
$X = \bigcup X_i$，使得对每个 $i$，存在概形 $U_i$ 及拟紧满
\'etale 态射 $U_i \to X_i$。选取一个 $i$ 及仿射开子概形
$W \subset U_i$。只须证明 $W \to X$ 的纤维普遍有界，因为这时所有这些
态射 $W \to X$ 构成的族覆盖 $X$。为此考虑图表
$$
\xymatrix{
W \times_X U_i \ar[r]_-p \ar[d]_q & U_i \ar[d] \\
W \ar[r] & X
}
$$
由于 $W \to X$ 经过 $X_i$ 分解，有
$W \times_X U_i = W \times_{X_i} U_i$，故 $q$ 拟紧。又因 $W$ 仿射，
这蕴含概形 $W \times_X U_i$ 拟紧。于是可应用《态射》引理
\ref{morphisms-lemma-locally-quasi-finite-qc-source-universally-bounded}
得知 $p$ 的纤维普遍有界。再由引理
\ref{decent-spaces-lemma-descent-universally-bounded}，$W \to X$ 的纤维也普遍有界。

\medskip\noindent
假设 $(\delta)$ 成立。设 $U$ 为仿射概形，$U \to X$ 为 \'etale 态射。
由假设，态射 $U \to X$ 的纤维普遍有界。于是两个投影
$R = U \times_X U \to U$ 的纤维也普遍有界；见引理
\ref{decent-spaces-lemma-base-change-universally-bounded}。再由引理
\ref{decent-spaces-lemma-composition-universally-bounded}，$R \to X$ 的纤维也普遍有界。
% P08-E275: frozen source uses x in X instead of x in |X|; formula preserved; see ERRATA_CANDIDATES.jsonl.
因此，对任意 $x \in X$，$|U| \to |X|$ 与 $|R| \to |X|$ 在 $x$ 上的纤维
有限；见引理 \ref{decent-spaces-lemma-universally-bounded-finite-fibres}。换言之，引理
\ref{decent-spaces-lemma-UR-finite-above-x} 的等价条件成立。这证明了
$(\delta) \Rightarrow (\gamma)$。
\end{proof}

\begin{lemma}
\label{decent-spaces-lemma-properties-local}
设 $S$ 为概形。设 $\mathcal{P}$ 是引理 \ref{decent-spaces-lemma-bounded-fibres} 中所列
代数空间的性质 $(\alpha)$、$(\beta)$、$(\gamma)$、$(\delta)$、
$(\epsilon)$、$(\zeta)$、$(\theta)$ 之一。若 $X$ 是 $S$ 上的代数空间，
$X = \bigcup X_i$ 是 Zariski 开覆盖，且每个 $X_i$ 都具有
$\mathcal{P}$，则 $X$ 具有 $\mathcal{P}$。
\end{lemma}

\begin{proof}
% P08-E276: frozen source says “let ... is a”; intended “let ... be a” is rendered; see ERRATA_CANDIDATES.jsonl.
设 $X$ 为 $S$ 上的代数空间，并设 $X = \bigcup X_i$ 为 Zariski 开覆盖，
其中每个 $X_i$ 都具有 $\mathcal{P}$。

\medskip\noindent
考虑 $\mathcal{P} = (\alpha)$ 的情形。$X_i$ 的条件 $(\alpha)$ 意味着：
对每个 $x \in |X_i|$、每个仿射概形 $U$ 及 \'etale 态射
$\varphi : U \to X_i$，$\varphi : |U| \to |X_i|$ 在 $x$ 上的纤维有限。
% P08-E277: frozen source uses x in X instead of x in |X|; formula preserved; see ERRATA_CANDIDATES.jsonl.
考虑 $x \in X$、仿射概形 $U$ 以及 \'etale 态射 $U \to X$。由于
$X = \bigcup X_i$ 是 Zariski 开覆盖，存在有限仿射开覆盖
% P08-E278: frozen source says “there exits”; intended “there exists” is rendered; see ERRATA_CANDIDATES.jsonl.
$U = U_1 \cup \ldots \cup U_n$，使得每个 $U_j \to X$ 都经过某个
$X_{i_j}$ 分解。由假设，对 $j = 1, \ldots, n$，
$|U_j | \to |X_{i_j}|$ 在 $x$ 上的纤维有限。这显然意味着
$|U| \to |X|$ 在 $x$ 上的纤维有限。这证明了 $(\alpha)$ 的情形。

\medskip\noindent
考虑 $\mathcal{P} = (\beta)$ 的情形。$X_i$ 的条件 $(\beta)$ 意味着每个
$x \in |X_i|$ 都可由从域谱到 $X_i$ 的单态射表示。由于
$X_i \to X$ 是单态射，且 $X = \bigcup X_i$，故 $X$ 也有相同性质。

\medskip\noindent
考虑 $\mathcal{P} = (\gamma)$ 的情形。注意，由引理
\ref{decent-spaces-lemma-bounded-fibres}，$(\gamma) = (\alpha) + (\beta)$，故
$(\gamma)$ 的结论来自上面已经处理的情形。

\medskip\noindent
考虑 $\mathcal{P} = (\delta)$ 的情形。$X_i$ 的条件 $(\delta)$ 意味着存在
概形 $U_{ij}$ 及纤维普遍有界的 \'etale 态射 $U_{ij} \to X_i$，它们覆盖
$X_i$。这些概形还给出满 \'etale 态射 $\coprod U_{ij} \to X$，并且
$U_{ij} \to X$ 的纤维仍普遍有界。

\medskip\noindent
考虑 $\mathcal{P} = (\epsilon)$ 的情形。$X_i$ 的条件 $(\epsilon)$ 意味着
可找到集合 $J_i$ 及态射 $\varphi_{ij} : U_{ij} \to X_i$，使得每个
$\varphi_{ij}$ 都是 \'etale 态射，两个投影
$U_{ij} \times_{X_i} U_{ij} \to U_{ij}$ 都拟紧，并且
$\coprod_{j \in J_i} U_{ij} \to X_i$ 满射。此时复合
$U_{ij} \to X_i \to X$ 是 \'etale 态射；这是综合《态射》引理
\ref{morphisms-lemma-composition-etale} 与
\ref{morphisms-lemma-open-immersion-etale} 以及《空间》引理
\ref{spaces-lemma-composition-representable-transformations-property} 与
\ref{spaces-lemma-morphism-schemes-gives-representable-transformation-property}
所得。由于 $X_i \subset X$ 是子空间，有
$U_{ij} \times_{X_i} U_{ij} = U_{ij} \times_X U_{ij}$，故纤维积条件得以
保持。又显然 $\coprod_{i, j} U_{ij} \to X$ 满射。因此 $X$ 满足
$(\epsilon)$。

\medskip\noindent
考虑 $\mathcal{P} = (\zeta)$ 的情形。$X_i$ 的条件 $(\zeta)$ 意味着
$X_i$ 在 Zariski 局部拟分离。这立即表明 $X$ 在 Zariski 局部拟分离。

\medskip\noindent
关于 $(\theta)$，见《空间的性质》引理
\ref{spaces-properties-lemma-subscheme}。
\end{proof}

\begin{lemma}
\label{decent-spaces-lemma-representable-properties}
设 $S$ 为概形。设 $\mathcal{P}$ 是引理 \ref{decent-spaces-lemma-bounded-fibres} 中所列
代数空间的性质 $(\beta)$、$(\gamma)$、$(\delta)$、$(\epsilon)$、
$(\theta)$ 之一。设 $X$、$Y$ 为 $S$ 上的代数空间，并设
$X \to Y$ 为可表示态射。若 $Y$ 具有性质 $\mathcal{P}$，则 $X$ 也具有该性质。
\end{lemma}

\begin{proof}
设 $f : X \to Y$ 是代数空间的可表示态射，并假设 $Y$ 具有
$\mathcal{P}$。取 $x \in |X|$，并令 $y = f(x) \in |Y|$。

\medskip\noindent
考虑 $\mathcal{P} = (\beta)$ 的情形。$Y$ 的条件 $(\beta)$ 意味着存在表示
$y$ 的单态射 $\Spec(k) \to Y$。纤维积
$X_y = \Spec(k) \times_Y X$ 是概形，而 $x$ 对应于 $X_y$ 的一个点，
即对应于单态射 $\Spec(k') \to X_y$。由于 $X_y \to X$ 也是单态射，
可知 $x$ 由单态射 $\Spec(k') \to X_y \to X$ 表示。换言之，
$X$ 满足 $(\beta)$。

\medskip\noindent
考虑 $\mathcal{P} = (\gamma)$ 的情形。由于
$(\gamma) \Rightarrow (\beta)$，上一段表明 $y$ 与 $x$ 可由下图中的
单态射表示：
$$
\xymatrix{
\Spec(k') \ar[r]_-x \ar[d] & X \ar[d] \\
\Spec(k) \ar[r]^-y & Y
}
$$
此外，由引理 \ref{decent-spaces-lemma-UR-finite-above-x} 的 (2) 所给出的性质
$(\gamma)$ 的定义，存在概形 $V_i$ 及 \'etale 态射 $V_i \to Y$，使得
$\coprod V_i \to Y$ 满射；并且对每个 $i$，令
$R_i = V_i \times_Y V_i$ 后，下列两个映射在 $y$ 上的纤维都有限：
$$
|V_i| \longrightarrow |Y|
\quad\text{且}\quad
|R_i| \longrightarrow |Y|
$$
这意味着概形 $(V_i)_y$ 与 $(R_i)_y$ 在 $y = \Spec(k)$ 上有限。
由于 $X \to Y$ 可表示，纤维积 $U_i = V_i \times_Y X$ 是概形。
态射 $U_i \to X$ 是 \'etale 的，且 $\coprod U_i \to X$ 满射。
最后，对每个 $i$，有
$$
(U_i)_x =
(V_i \times_Y X)_x =
(V_i)_y \times_{\Spec(k)} \Spec(k')
$$
以及
$$
(U_i \times_X U_i)_x =
\left((V_i \times_Y X) \times_X (V_i \times_Y X)\right)_x =
(R_i)_y \times_{\Spec(k)} \Spec(k')
$$
因此它们作为有限概形 $(V_i)_y$ 与 $(R_i)_y$ 的基变换，在 $k'$ 上有限。
再次应用引理 \ref{decent-spaces-lemma-UR-finite-above-x} 的第二个条件，可知
$X$ 满足 $(\gamma)$。

\medskip\noindent
考虑 $\mathcal{P} = (\delta)$ 的情形。设 $V \to Y$ 是 \'etale 态射，其中
$V$ 为仿射概形。由于 $Y$ 具有性质 $(\delta)$，该态射的纤维普遍有界。
由引理 \ref{decent-spaces-lemma-base-change-universally-bounded}，基变换
$V \times_Y X \to X$ 的纤维也普遍有界。因此引理
\ref{decent-spaces-lemma-U-universally-bounded} 的第一部分适用，并且我们看到
% P08-E279: frozen source concludes Y instead of X; wording preserved; see ERRATA_CANDIDATES.jsonl.
$Y$ 也具有性质 $(\delta)$。

\medskip\noindent
考虑 $\mathcal{P} = (\epsilon)$ 的情形。我们将反复使用《空间》引理
\ref{spaces-lemma-base-change-representable-transformations-property}。
设 $V_i \to Y$ 如引理
\ref{decent-spaces-lemma-characterize-very-reasonable} 的 (2) 所述。令
$U_i = X \times_Y V_i$。态射 $U_i \to X$ 是 \'etale 的，且
$\coprod U_i \to X$ 满射。由于
$U_i \times_X U_i = X \times_Y (V_i \times_Y V_i)$，可知
% P08-E280: frozen source changes the base from X to Y in the projection; formula preserved; see ERRATA_CANDIDATES.jsonl.
投影 $U_i \times_Y U_i \to U_i$ 是投影
$V_i \times_Y V_i \to V_i$ 的基变换，因而也拟紧。因此 $X$ 满足引理
\ref{decent-spaces-lemma-characterize-very-reasonable} 的 (2)。

\medskip\noindent
考虑 $\mathcal{P} = (\theta)$ 的情形。此时结论就是《范畴》引理
\ref{categories-lemma-representable-over-representable}。
\end{proof}











\section{合理代数空间与良好代数空间}
\label{decent-spaces-section-reasonable-decent}

\noindent
在引理 \ref{decent-spaces-lemma-bounded-fibres} 中，我们已经看到代数空间的若干条件；
它们涉及从仿射概形到 $X$ 的 \'etale 态射的行为，以及由概形对 $X$ 作
某些特殊 \'etale 覆盖的存在性。现将这些不同类型的条件列表如下：
$$
\boxed{
\begin{matrix}
(\alpha) & \text{从仿射概形出发的 \'etale 态射具有有限纤维} \\
(\beta) & \text{点来自域谱的单态射} \\
(\gamma) & \text{点来自域谱的拟紧单态射} \\
(\delta) & \text{从仿射概形出发的 \'etale 态射具有普遍有界纤维} \\
(\epsilon) & \text{可由概形作 \'etale 覆盖，且各态射到其像拟紧}
\end{matrix}
}
$$

\medskip\noindent
下述定义中的条件并不完全是 $X$ 的对角线所满足的条件，但在某种意义上，
它们是 $X$ 的分离性条件。

\begin{definition}
\label{decent-spaces-definition-very-reasonable}
设 $S$ 为概形，$X$ 为 $S$ 上的代数空间。
\begin{enumerate}
\item 若对每个点 $x \in X$，引理
\ref{decent-spaces-lemma-UR-finite-above-x} 的等价条件成立，换言之，引理
\ref{decent-spaces-lemma-bounded-fibres} 的性质 $(\gamma)$ 成立，则称 $X$ 是
{\it 良好的}。
% P08-E281: frozen source uses x in X instead of x in |X|; formula preserved; see ERRATA_CANDIDATES.jsonl.
\item 若引理 \ref{decent-spaces-lemma-U-universally-bounded} 的等价条件成立，换言之，
引理 \ref{decent-spaces-lemma-bounded-fibres} 的性质 $(\delta)$ 成立，则称 $X$ 是
{\it 合理的}。
\item 若引理 \ref{decent-spaces-lemma-characterize-very-reasonable} 的等价条件成立，
即引理 \ref{decent-spaces-lemma-bounded-fibres} 的性质 $(\epsilon)$ 成立，则称 $X$ 是
{\it 非常合理的}。
\end{enumerate}
\end{definition}

\noindent
代数空间的这些条件之间有下列蕴含关系：
$$
\xymatrix{
\text{可表示} \ar@{=>}[rd] & & & \\
 & \text{非常合理} \ar@{=>}[r] &
\text{合理} \ar@{=>}[r] &
\text{良好} \\
\text{拟分离} \ar@{=>}[ru] & & &
}
$$
“非常合理代数空间”这一概念现已废弃。引入它，是因为过去证明某些结果需要
这一假设，而这些结果现在已经对良好空间这一类得到证明。良好空间构成这样
的空间 $X$ 的最大类别：其中 $|X|$ 的拓扑与 $X$ 自身的性质之间具有良好
联系。

\begin{example}
\label{decent-spaces-example-not-decent}
《空间》例 \ref{spaces-example-affine-line-translation} 中构造的代数空间
$\mathbf{A}^1_{\mathbf{Q}}/\mathbf{Z}$ 不是良好的，因为它的“一般点”
不能由从域谱出发的单态射表示。
\end{example}

\begin{remark}
\label{decent-spaces-remark-reasonable}
从技术上说，合理代数空间比非常合理代数空间更易处理。例如，若
$X \to Y$ 是代数空间的拟紧满 \'etale 态射，且 $X$ 合理，则 $Y$ 也合理；
见引理 \ref{decent-spaces-lemma-descent-conditions}。但我们不知道性质“非常合理”是否
也有这一结论。下面还将给出合理代数空间所具有的另一个技术性质。
\end{remark}

\begin{lemma}
\label{decent-spaces-lemma-fun-property-reasonable}
设 $S$ 为概形，$X$ 为拟紧合理代数空间。则存在由拟紧拟分离代数空间
$X_i$ 构成的有向系统，使得 $X = \colim_i X_i$（这是层范畴中的余极限）。
而且可将其安排为满足：
\begin{enumerate}
\item 对每个 $S$ 上的拟紧概形 $T$，有 $\colim X_i(T) = X(T)$；
\item 系统的转移态射 $X_i \to X_{i'}$ 与余投影 $X_i \to X$ 都满且
\'etale；
\item 若 $X$ 是概形，则代数空间 $X_i$ 都是概形，而转移态射
$X_i \to X_{i'}$ 与余投影 $X_i \to X$ 都是局部同构。
\end{enumerate}
\end{lemma}

\begin{proof}
略述证明。由《空间的性质》引理
\ref{spaces-properties-lemma-quasi-compact-affine-cover}，有
$X = U/R$，其中 $U$ 仿射。在此情形，合理性意味着 $U \to X$ 的纤维
普遍有界。因此存在整数 $N$，使 $U \to X$ 的“纤维”次数至多为 $N$；
见定义 \ref{decent-spaces-definition-universally-bounded}。将群胚的结构映射记为
$s, t : R \to U$ 与 $c : R \times_{s, U, t} R \to R$。

\medskip\noindent
断言：对每个拟紧开子集 $A \subset R$，存在开子集 $R' \subset R$，使得
\begin{enumerate}
\item $A \subset R'$,
\item $R'$ 拟紧；
\item $(U, R', s|_{R'}, t|_{R'}, c|_{R' \times_{s, U, t} R'})$ 是群胚概形。
\end{enumerate}
注意，$e : U \to R$ 是 \'etale 态射 $s : R \to U$ 的截面，因而是开态射；
见《\'Etale 态射》命题 \ref{etale-proposition-properties-sections}。
又因 $U$ 仿射，它拟紧。因此可用 $A \cup e(U) \subset R$ 代替 $A$，
并假设 $A$ 包含 $e(U)$。
% P08-E283: frozen construction does not close A under the groupoid inverse; preserved; see ERRATA_CANDIDATES.jsonl.
接着归纳定义 $A^1 = A$ 以及
$$
A^n = c(A^{n - 1} \times_{s, U, t} A) \subset R
$$
（$n \geq 2$）。归纳可见，对所有 $n \geq 2$，$A^n$ 都拟紧，因为它是拟紧
纤维积 $A^{n - 1} \times_{s, U, t} A$ 的像。若 $k$ 是 $S$ 上的代数闭域，
并考察 $k$-点，则
% P08-E282: frozen path formula omits the endpoint u_n = u'; formula preserved; see ERRATA_CANDIDATES.jsonl.
$$
A^n(k) = \left\{(u, u') \in U(k)
:
\begin{matrix}
\text{存在 } u = u_1, u_2, \ldots, u_n \in U(k)\text{，使得} \\
(u_i , u_{i + 1}) \in A \text{ 对所有 }i = 1, \ldots, n - 1.
\end{matrix}
\right\}
$$
由于 $U(k) \to X(k)$ 的纤维大小至多为 $N$，若 $n > N$，上列序列中便有
重复项，因而可将其缩短；这证明对所有 $n \geq N$，都有
$A^N = A^n$。于是 $R' = A^N$ 给出群胚概形
$(U, R', s|_{R'}, t|_{R'}, c|_{R' \times_{s, U, t} R'})$，从而证明上述断言。

\medskip\noindent
考虑 $(\Sch/S)_{fppf}$ 上的层映射
$$
\colim_{R' \subset R} U/R' \longrightarrow U/R
$$
其中 $R' \subset R$ 遍历 $R$ 的所有拟紧开子概形，它们如上给出 \'etale
等价关系。每个商 $U/R'$ 都是代数空间（见《空间》定理
\ref{spaces-theorem-presentation}）。由于 $R'$ 拟紧且 $U$ 仿射，态射
$R' \to U \times_{\Spec(\mathbf{Z})} U$ 拟紧，故 $U/R'$ 拟分离。最后，
若 $T$ 是拟紧概形，则
$$
\colim_{R' \subset R} U(T)/R'(T) \longrightarrow U(T)/R(T)
$$
是双射，因为由上述断言，从 $T$ 到 $R$ 的每个态射最终都落在某个开子关系
$R'$ 中。这显然蕴含层 $U/R'$ 的余极限为 $U/R$。换言之，代数空间
$X = U/R$ 是拟分离代数空间 $U/R'$ 的余极限。

\medskip\noindent
性质 (1)、(2) 来自上述讨论。若 $X$ 是概形，则选择 $U$ 为 $X$ 的有限多个
仿射开集之不交并即可得到 (3)。略去细节。
\end{proof}

\begin{lemma}
\label{decent-spaces-lemma-representable-named-properties}
设 $S$ 为概形，$X$、$Y$ 为 $S$ 上的代数空间，并设 $X \to Y$ 为可表示
态射。若 $Y$ 良好（相应地，合理），则 $X$ 也如此。
\end{lemma}

\begin{proof}
这是引理 \ref{decent-spaces-lemma-representable-properties} 的改述。
\end{proof}

\begin{lemma}
\label{decent-spaces-lemma-etale-named-properties}
设 $S$ 为概形，$X \to Y$ 为 $S$ 上代数空间的 \'etale 态射。若 $Y$ 良好，
相应地，合理，则 $X$ 也如此。
\end{lemma}

\begin{proof}
设 $U$ 为仿射概形，$U \to X$ 为 \'etale 态射。令
$R = U \times_X U$ 与 $R' = U \times_Y U$。注意，$R \to R'$ 是单态射。

\medskip\noindent
设 $x \in |X|$。为证明 $X$ 良好，须证明 $|U| \to |X|$ 与
$|R| \to |X|$ 在 $x$ 上的纤维有限。但若 $Y$ 良好，则
$|U| \to |Y|$ 与 $|R'| \to |Y|$ 的纤维有限。因此得到“良好”的结论。

\medskip\noindent
为证明 $X$ 合理，须证明 $U \to X$ 的纤维普遍有界。然而，若 $Y$ 合理，
则 $U \to Y$ 的纤维普遍有界，而这立即蕴含 $U \to X$ 的纤维也普遍有界。
因此得到“合理”的结论。
\end{proof}







\section{点与特化}
\label{decent-spaces-section-specializations}

\noindent
存在代数空间的 \'etale 态射 $f : X \to Y$，使得
$|f| : |X| \to |Y|$ 的某个纤维中的点之间存在非平凡特化；见《例》引理
\ref{examples-lemma-specializations-fibre-etale}。若该态射的源是概形，
则可通过要求 $Y$ 满足条件 ($\alpha$) 来避免这种情形。

\begin{lemma}
\label{decent-spaces-lemma-no-specializations-map-to-same-point}
设 $S$ 为概形，$X$ 为 $S$ 上的代数空间，并设 $U \to X$ 是从概形到
$X$ 的 \'etale 态射。假设 $u, u' \in |U|$ 映到 $|X|$ 的同一点 $x$，
且 $u' \leadsto u$。若序对 $(X, x)$ 满足引理
\ref{decent-spaces-lemma-U-finite-above-x} 的等价条件，则 $u = u'$。
\end{lemma}

\begin{proof}
假设序对 $(X, x)$ 满足引理 \ref{decent-spaces-lemma-U-finite-above-x} 的等价条件。
设 $U$ 为概形，$U \to X$ 为 \'etale 态射，并设 $u, u' \in |U|$ 映到
$|X|$ 中的 $x$，且 $u' \leadsto u$。我们可以且确实用 $u$ 的一个仿射
邻域替换 $U$。设 $t, s : R = U \times_X U \to U$ 为两个 \'etale 投影。

\medskip\noindent
取点 $r \in R$，使得 $t(r) = u$ 且 $s(r) = u'$。由《空间的性质》引理
\ref{spaces-properties-lemma-points-presentation}，这是可能的。
由于一般化可沿 \'etale 态射 $t$ 提升（注 \ref{decent-spaces-remark-recall}），可找到
特化 $r' \leadsto r$，使得 $t(r') = u'$。令 $u'' = s(r')$，则
$u'' \leadsto u'$。于是可重复此过程，找到 $r'' \leadsto r'$，使得
$t(r'') = u''$。令 $u''' = s(r'')$，如此继续。图示如下：
$$
\xymatrix{
& r'' \ar[rd]^s \ar[ld]_t \ar@{~>}[d] & \\
u'' \ar@{~>}[d] & r' \ar[rd]^s \ar[ld]_t \ar@{~>}[d] & u''' \ar@{~>}[d] \\
u' \ar@{~>}[d] & r \ar[rd]^s \ar[ld]_t & u'' \ar@{~>}[d] \\
u & & u'
}
$$
注 \ref{decent-spaces-remark-recall} 已说明，\'etale 态射 $s$ 的同一纤维中的点之间没有
特化。因此，若对某个 $n$ 有 $u^{(n + 1)} = u^{(n)}$，则也有
$r^{(n)} = r^{(n - 1)}$，进而取 $t$ 得
$u^{(n)} = u^{(n - 1)}$。这迫使整座塔坍缩，特别地 $u = u'$。所以，
若 $u \not = u'$，则所有这些特化都严格，且
$\{u, u', u'', \ldots\}$ 是 $U$ 中映到 $|X|$ 的点 $x$ 的无限点集。
由于已选 $U$ 为仿射概形，这与引理
\ref{decent-spaces-lemma-U-finite-above-x} 的第二部分矛盾，如所需。
\end{proof}

\begin{lemma}
\label{decent-spaces-lemma-no-specializations-map-to-same-point-Noetherian}
设 $S$ 为概形，$X$ 为 $S$ 上的代数空间。设 $U \to X$ 是从概形到 $X$
的 \'etale 态射。假设 $u, u' \in |U|$ 映到 $|X|$ 的同一点 $x$，且
$u' \leadsto u$。若 $X$ 局部 Noether，则 $u = u'$。
\end{lemma}

\begin{proof}
《概形》第 \ref{schemes-section-points} 节的讨论表明，
$\mathcal{O}_{U, u'}$ 是 Noether 局部环 $\mathcal{O}_{U, u}$ 的局部化。
由《空间的性质》引理
\ref{spaces-properties-lemma-pre-dimension-local-ring}，有
$\dim(\mathcal{O}_{U, u}) = \dim(\mathcal{O}_{U, u'})$。
Noether 局部环的维数理论于是给出 $u = u'$。
\end{proof}

\begin{lemma}
\label{decent-spaces-lemma-specialization}
设 $S$ 为概形，$X$ 为 $S$ 上的代数空间。设 $x, x' \in |X|$，并假设
$x' \leadsto x$，也就是说，$x$ 是 $x'$ 的一个特化。假设序对 $(X, x')$
满足引理 \ref{decent-spaces-lemma-UR-finite-above-x} 的等价条件。则对任意从概形 $U$
出发的 \'etale 态射 $\varphi : U \to X$ 以及任意满足 $\varphi(u) = x$ 的
$u \in U$，都存在点 $u'\in U$，使得 $u' \leadsto u$ 且
$\varphi(u') = x'$。
% P08-E284: frozen source omits “there” before “exists”; intended proposition translated; see ERRATA_CANDIDATES.jsonl.
\end{lemma}

\begin{proof}
我们可以用 $u$ 的一个仿射开邻域替换 $U$，因而可以假设 $U$ 仿射。
由于 $x$ 属于开映射 $|U| \to |X|$ 的像，$x'$ 也属于该像。因此，由
《空间的性质》引理
\ref{spaces-properties-lemma-etale-image-open}，可以用与
$|U| \to |X|$ 的像对应的 Zariski 开子空间替换 $X$。换言之，可以假设
$U \to X$ 满且 \'etale。设 $s, t : R = U \times_X U \to U$ 为两个投影。
根据序对 $(X, x')$ 满足引理 \ref{decent-spaces-lemma-UR-finite-above-x} 的等价条件这一
假设，$|U| \to |X|$ 和 $|R| \to |X|$ 在 $x'$ 上的纤维都是有限的。设
% P08-E285: frozen source types point sets as subsets of U and R rather than |U| and |R|; formulas preserved; see ERRATA_CANDIDATES.jsonl.
$\{u'_1, \ldots, u'_n\} \subset U$ 与
$\{r'_1, \ldots, r'_m\} \subset R$ 构成 $\{x'\}$ 的全部原像。考虑闭集
$$
T = \overline{\{u'_1\}} \cup \ldots \cup \overline{\{u'_n\}} \subset |U|,
\quad
T' = \overline{\{r'_1\}} \cup \ldots \cup \overline{\{r'_m\}} \subset |R|.
$$
显然 $s(T') \subset T$。由于 $R$ 是等价关系，且点集 $\{r_j'\}$ 按构造在
$R$ 的取逆运算下不变，故还有 $t(T') = s(T')$。任取点 $w \in T$，则对
某个 $i$ 有 $u'_i \leadsto w$。取 $r \in R$ 使得 $s(r) = w$。由注
\ref{decent-spaces-remark-recall}，一般化可沿 $s : R \to U$ 提升，故可找到
$r' \leadsto r$，使得 $s(r') = u_i'$。于是对某个 $j$ 有 $r' = r'_j$，
从而 $w \in s(T')$。因此 $T = s(T') = t(T')$ 是 $|U|$ 中的一个
$|R|$-不变闭集。这意味着 $T$ 是 $|X|$ 的闭子集（！）
$T'' = \varphi(T)$ 的原像；参见《空间的性质》引理
\ref{spaces-properties-lemma-points-presentation} 和
\ref{spaces-properties-lemma-topology-points}。所以
$T'' = \overline{\{x'\}}$。由于 $x \in T''$，$T$ 含有某个映到 $x$ 的点
$u_1$。也就是说，对某个 $i$，存在映到给定特化 $x' \leadsto x$ 的特化
$u'_i \leadsto u_1$。

\medskip\noindent
为完成证明，由《空间的性质》引理
\ref{spaces-properties-lemma-points-cartesian}，取点 $r \in R$，使得
$s(r) = u$ 且 $t(r) = u_1$。由于一般化可沿 $t$ 提升，且
$u'_i \leadsto u_1$，可找到特化 $r' \leadsto r$，使得
$t(r') = u'_i$。令 $u' = s(r')$，则 $u' \leadsto u$ 且
$\varphi(u') = x'$，如所需。
\end{proof}

\begin{lemma}
\label{decent-spaces-lemma-generalizations-lift-flat}
设 $S$ 为概形，$f : Y \to X$ 为 $S$ 上代数空间的平坦态射。设
$x, x' \in |X|$，并假设 $x' \leadsto x$，也就是说，$x$ 是 $x'$ 的一个
特化。假设序对 $(X, x')$ 满足引理 \ref{decent-spaces-lemma-UR-finite-above-x} 的等价
条件（例如，$X$ 良好、$X$ 拟分离或 $X$ 可表时即如此）。则对每个满足
$f(y) = x$ 的 $y \in |Y|$，都存在点 $y' \in |Y|$，使得
$y' \leadsto y$ 且 $f(y') = x'$。
\end{lemma}

\begin{proof}
（括号中的断言由良好空间的定义以及第
\ref{decent-spaces-section-reasonable-decent} 节所述各分离条件之间的蕴含关系得到。）
取概形 $V$ 及满 \'etale 态射 $V \to Y$，再取映到 $y$ 的 $v \in V$。
于是只需对 $V \to X$ 证明本引理，故可假设 $Y$ 为概形。取概形 $U$ 及
满 \'etale 态射 $U \to X$，并取映到 $x$ 的 $u \in U$。由引理
\ref{decent-spaces-lemma-specialization}，可选取映到 $x'$ 的 $u' \leadsto u$。由
《空间的性质》引理 \ref{spaces-properties-lemma-points-cartesian}，可选取
$z \in U \times_X Y$，使其映到 $y$ 和 $u$。于是归结为概形之间的平坦
态射 $U \times_X Y \to U$ 的情形，而这正是《态射》引理
\ref{morphisms-lemma-generalizations-lift-flat}。
\end{proof}






\section{用概形对代数空间分层}
\label{decent-spaces-section-stratifications}

\noindent
本节将证明：拟紧拟分离代数空间存在由局部闭子空间组成的有限分层，其中每一层
都是概形，而且各部分通过初等特异方块（elementary distinguished square）
粘合。我们先对合理代数空间证明一个
稍弱的结果。

\begin{lemma}
\label{decent-spaces-lemma-quasi-compact-reasonable-stratification}
设 $S$ 为概形。设 $W \to X$ 为从概形 $W$ 到代数空间 $X$ 的态射，并且
该态射平坦、局部有限表示、分离、局部拟有限且纤维普遍有界。存在既约闭子空间
$$
\emptyset = Z_{-1} \subset Z_0 \subset Z_1 \subset Z_2 \subset
\ldots \subset Z_n = X
$$
使得，若令 $X_r = Z_r \setminus Z_{r - 1}$，则分层
$X = \coprod_{r = 0, \ldots, n} X_r$ 由下列泛性质刻画：给定
$g : T \to X$，投影 $W \times_X T \to T$ 是次数为 $r$ 的有限局部自由
态射，当且仅当 $g(|T|) \subset |X_r|$。
\end{lemma}

\begin{proof}
取整数 $n$，使其界住 $W \to X$ 各纤维的次数。取概形 $U$ 及满
\'etale 态射 $U \to X$。将《态射进阶》引理
\ref{more-morphisms-lemma-stratify-flat-fp-lqf-universally-bounded}
用于 $W \times_X U \to U$，得到闭子集
$$
\emptyset = Y_{-1} \subset Y_0 \subset Y_1 \subset Y_2 \subset
\ldots \subset Y_n = U
$$
它们由本引理中针对态射 $W \times_X U \to U$ 所述的性质刻画。显然，这些
闭子集的形成与基变换交换。令 $R = U \times_X U$，其投影为
$s, t : R \to U$，可得
$$
s^{-1}(Y_r) = t^{-1}(Y_r)
$$
作为 $R$ 的闭子集相等。换言之，闭子集 $Y_r \subset U$ 是 $R$-不变的。
这意味着 $|Y_r|$ 是某个闭子集 $Z_r \subset |X|$ 的原像。仍用
$Z_r \subset X$ 表示赋予既约诱导代数空间结构的该子集；参见
《空间的性质》定义
\ref{spaces-properties-definition-reduced-induced-space}。

\medskip\noindent
设 $g : T \to X$ 为代数空间的态射。取概形 $V$ 及满 \'etale 态射
$V \to T$。为证明引理最后的断言，只须对复合态射 $V \to X$ 证明它
（由有限局部自由态射的定义可知；参见《空间的态射》第
\ref{spaces-morphisms-section-finite-locally-free} 节）。类似地，概形态射
$W \times_X V \to V$ 是次数为 $r$ 的有限局部自由态射，当且仅当概形态射
$$
W \times_X (U \times_X V)
\longrightarrow
U \times_X V
$$
是次数为 $r$ 的有限局部自由态射（参见《下降》引理
\ref{descent-lemma-descending-property-finite-locally-free}）。按构造，这又
等价于 $|U \times_X V| \to |U|$ 的像包含在 $|Y_r|$ 中，而后者等价于
$|V| \to |X|$ 的像包含在 $|Z_r|$ 中。
\end{proof}

\begin{lemma}
\label{decent-spaces-lemma-stratify-flat-fp-lqf}
设 $S$ 为概形。设 $W \to X$ 为从概形 $W$ 到代数空间 $X$ 的态射，并且
该态射平坦、局部有限表示、分离且局部拟有限。则存在开子空间
$$
X = X_0 \supset X_1 \supset X_2 \supset \ldots
$$
使得，对任意域 $k$，态射 $\Spec(k) \to X$ 经由 $X_d$ 分解，当且仅当
$W \times_X \Spec(k)$ 在 $k$ 上的次数为 $\geq d$。
\end{lemma}

\begin{proof}
取概形 $U$ 及满 \'etale 态射 $U \to X$。将《态射进阶》引理
\ref{more-morphisms-lemma-stratify-flat-fp-lqf} 用于
$W \times_X U \to U$，得到开子概形
$$
U = U_0 \supset U_1 \supset U_2 \supset \ldots
$$
它们由本引理中针对态射 $W \times_X U \to U$ 所述的性质刻画。
% P08-E286: frozen source calls the open subschemes U_d “closed subsets”; intended word translated; see ERRATA_CANDIDATES.jsonl.
显然，这些开子概形的形成与基变换交换。令 $R = U \times_X U$，其投影为
$s, t : R \to U$，可得
$$
s^{-1}(U_d) = t^{-1}(U_d)
$$
作为 $R$ 的开子概形相等。换言之，开子概形 $U_d \subset U$ 是
$R$-不变的。这意味着 $U_d$ 是某个开子空间 $X_d \subset X$ 的原像
（《空间的性质》引理
\ref{spaces-properties-lemma-subspaces-presentation}）。
\end{proof}

\begin{lemma}
\label{decent-spaces-lemma-filter-quasi-compact}
设 $S$ 为概形，$X$ 为 $S$ 上的拟紧代数空间。存在开子空间
$$
\ldots \subset U_4 \subset U_3 \subset U_2 \subset U_1 = X
$$
具有下列性质：
\begin{enumerate}
\item 令 $T_p = U_p \setminus U_{p + 1}$（赋予既约诱导子空间结构），则
存在分离概形 $V_p$ 及满 \'etale 态射 $f_p : V_p \to U_p$，使得
$f_p^{-1}(T_p) \to T_p$ 是同构；
\item 若 $x \in |X|$ 可由某个域给出的拟紧态射 $\Spec(k) \to X$ 表示，
则对某个 $p$ 有 $x \in T_p$。
\end{enumerate}
\end{lemma}

\begin{proof}
由《空间的性质》引理
\ref{spaces-properties-lemma-quasi-compact-affine-cover}，可取仿射概形 $U$
及满 \'etale 态射 $U \to X$。
% P08-E287: frozen source starts at p >= 0 although the displayed p-fold scheme fibre product and later filtration start at p = 1; preserved; see ERRATA_CANDIDATES.jsonl.
对 $p \geq 0$，令
$$
W_p = U \times_X \ldots \times_X U \setminus \text{所有对角线}
$$
其中纤维积有 $p$ 个因子。由于 $U$ 分离，态射 $U \to X$ 分离，且所有
纤维积 $U \times_X \ldots \times_X U$ 都是分离概形。由于 $U \to X$
分离，对角态射 $U \to U \times_X U$ 是闭浸入。由于 $U \to X$ 为
\'etale，对角态射 $U \to U \times_X U$ 又是开浸入；参见《空间的态射》
引理
\ref{spaces-morphisms-lemma-etale-unramified} 和
\ref{spaces-morphisms-lemma-diagonal-unramified-morphism}。
类似地，所有对角态射都是开闭浸入，且 $W_p$ 是
$U \times_X \ldots \times_X U$ 的开闭子概形。此外，态射
$$
U \times_X \ldots \times_X U \longrightarrow
U \times_{\Spec(\mathbf{Z})} \ldots \times_{\Spec(\mathbf{Z})} U
$$
局部拟有限且分离（《空间的态射》引理
\ref{spaces-morphisms-lemma-fibre-product-after-map}），而其目标是仿射概形。
因此，$U \times_X \ldots \times_X U$ 的任意有限点集都包含在某个仿射开中；
参见《态射进阶》引理
\ref{more-morphisms-lemma-separated-locally-quasi-finite-over-affine}。
所以 $W_p$ 也具有同样性质。对称群 $S_p$ 在 $X$ 上自由作用于 $W_p$
% P08-E288: frozen source says “fix point locus”; intended fixed-point locus translated; see ERRATA_CANDIDATES.jsonl.
（因为我们已从 $U \times_X \ldots \times_X U$ 中删去不动点轨迹）。由上述
事实及《空间的性质》命题
\ref{spaces-properties-proposition-finite-flat-equivalence-global}
可知，商 $V_p = W_p/S_p$ 是概形。由于 $S_p$ 在 $W_p$ 上的作用位于 $X$
上方，存在态射 $V_p \to X$。又因 $W_p \to X$ 为 \'etale，且
$W_p \to V_p$ 为满 \'etale 态射，故 $V_p \to X$ 也是 \'etale；参见
《空间的性质》引理 \ref{spaces-properties-lemma-etale-local}。再由
《空间的性质》引理 \ref{spaces-properties-lemma-quotient-separated} 可知
$V_p$ 是分离概形。

\medskip\noindent
令 $U_p \subset X$ 为 $V_p \to X$ 的像所给出的开子空间。按构造，当 $k$
代数闭时，态射 $\Spec(k) \to X$ 经由 $U_p$ 分解，当且仅当
$U \times_X \Spec(k)$ 有 $\geq p$ 个点；照例注意，概形论意义下
$U \times_X \Spec(k)$ 是若干（可能无限多个）$\Spec(k)$ 副本的不交并，
参见注 \ref{decent-spaces-remark-recall}。由此可知，$U_p$ 给出引理所述的 $X$ 的滤过。
此外，态射 $\Spec(k) \to X$ 经由 $T_p$ 分解，当且仅当
$U \times_X \Spec(k)$ 恰有 $p$ 个点。在这种情况下，
$V_p \times_X \Spec(k)$ 恰有一个点。令
$Z_p = f_p^{-1}(T_p) \subset V_p$，这是 $V_p$ 的闭子概形。于是
$Z_p \to T_p$ 是代数空间之间的 \'etale 态射，并且对任意代数闭域 $k$
在 $k$-值点上诱导双射。确切地说，这意味着 $Z_p \to T_p$ 普遍单射，
因而由《空间的态射》引理
\ref{spaces-morphisms-lemma-etale-universally-injective-open}
它是开浸入，进而是同构，这就证明了 (1)。

\medskip\noindent
设 $x : \Spec(k) \to X$ 为拟紧态射，其中 $k$ 为域。则复合态射
$\Spec(\overline{k}) \to \Spec(k) \to X$ 也拟紧（《空间的态射》引理
\ref{spaces-morphisms-lemma-composition-quasi-compact}）。在此情形下，概形
$U \times_X \Spec(\overline{k})$ 拟紧。鉴于上面已经看到它是若干
$\Spec(\overline{k})$ 副本的不交并，可知它只有有限多个点。若点数为
$p$，则确有 $x \in T_p$，证明完成。
\end{proof}

\begin{lemma}
\label{decent-spaces-lemma-filter-reasonable}
设 $S$ 为概形，$X$ 为 $S$ 上拟紧的合理代数空间。存在整数 $n$ 及开子空间
$$
\emptyset = U_{n + 1} \subset
U_n \subset U_{n - 1} \subset \ldots \subset U_1 = X
$$
具有下列性质：令 $T_p = U_p \setminus U_{p + 1}$（赋予既约诱导子空间
结构），则存在分离概形 $V_p$ 及满 \'etale 态射
$f_p : V_p \to U_p$，使得 $f_p^{-1}(T_p) \to T_p$ 是同构。
\end{lemma}

\begin{proof}
本引理的证明与引理 \ref{decent-spaces-lemma-filter-quasi-compact} 的证明完全相同。由
$X$ 合理，存在整数 $n$ 界住 $U \to X$ 各纤维的次数；参见定义
\ref{decent-spaces-definition-very-reasonable}。于是 $U_{n + 1} = \emptyset$，证明完成。
\end{proof}

\begin{lemma}
\label{decent-spaces-lemma-stratify-reasonable}
设 $S$ 为概形，$X$ 为 $S$ 上拟紧的合理代数空间。存在整数 $n$ 及开子空间
$$
\emptyset = U_{n + 1} \subset
U_n \subset U_{n - 1} \subset \ldots \subset U_1 = X
$$
使得每个 $T_p = U_p \setminus U_{p + 1}$（赋予既约诱导子空间结构）都是
概形。
\end{lemma}

\begin{proof}
这是引理 \ref{decent-spaces-lemma-filter-reasonable} 的直接推论。
\end{proof}

\noindent
下述结果几乎与 \cite[Proposition 5.7.8]{GruRay} 完全相同。

\begin{lemma}
\label{decent-spaces-lemma-filter-quasi-compact-quasi-separated}
\begin{reference}
本结果几乎与 \cite[Proposition 5.7.8]{GruRay} 完全相同。
\end{reference}
设 $X$ 为 $\Spec(\mathbf{Z})$ 上的拟紧拟分离代数空间。存在整数 $n$ 及
开子空间
$$
\emptyset = U_{n + 1} \subset
U_n \subset U_{n - 1} \subset \ldots \subset U_1 = X
$$
具有下列性质：令 $T_p = U_p \setminus U_{p + 1}$（赋予既约诱导子空间
结构），则存在拟紧分离概形 $V_p$ 及满 \'etale 态射
$f_p : V_p \to U_p$，使得 $f_p^{-1}(T_p) \to T_p$ 是同构。
\end{lemma}

\begin{proof}
本引理的证明与引理 \ref{decent-spaces-lemma-filter-quasi-compact} 的证明完全相同。注意，
拟分离空间是合理的；参见引理 \ref{decent-spaces-lemma-bounded-fibres} 和定义
\ref{decent-spaces-definition-very-reasonable}。所以，与引理
\ref{decent-spaces-lemma-filter-reasonable} 中一样，有 $U_{n + 1} = \emptyset$。在论证
末尾还应补充：由于 $X$ 拟分离，各概形
$U \times_X \ldots \times_X U$ 都拟紧。因此各概形 $W_p$ 拟紧，进而它们
对对称群 $S_p$ 的商 $V_p = W_p/S_p$ 也是拟紧概形。
\end{proof}

\noindent
下述引理或许应当放在别处。

\begin{lemma}
\label{decent-spaces-lemma-locally-constructible}
设 $S$ 为概形，$X$ 为 $S$ 上的拟分离代数空间，$E \subset |X|$ 为子集。
则 $E$ 是 \'etale 局部可构造的（《空间的性质》定义
\ref{spaces-properties-definition-locally-constructible}），当且仅当 $E$ 是
拓扑空间 $|X|$ 的局部可构造子集（《拓扑》定义
\ref{topology-definition-constructible}）。
\end{lemma}

\begin{proof}
先假设 $E \subset |X|$ 是拓扑空间 $|X|$ 的局部可构造子集。设
$f : U \to X$ 为 \'etale 态射，其中 $U$ 为概形。需要证明
$f^{-1}(E)$ 在 $U$ 中局部可构造。此问题关于 $U$ 和 $X$ 都是局部的，
故可假设 $X$ 拟紧、$E \subset |X|$ 可构造且 $U$ 仿射。在这种情况下，
$U \to X$ 拟紧，因而 $f : |U| \to |X|$ 拟紧。注意，$|X|$（相应地
$U$）的逆紧开集与 $|X|$（相应地 $U$）的拟紧开集是同一回事；参见
《拓扑》引理 \ref{topology-lemma-topology-quasi-separated-scheme}。于是由
《拓扑》引理 \ref{topology-lemma-inverse-images-constructibles}，
$f^{-1}(E)$ 可构造。

\medskip\noindent
反过来，假设 $E$ 是 \'etale 局部可构造的。我们要证明 $E$ 在拓扑空间
$|X|$ 中局部可构造。此问题关于 $X$ 是局部的，故可假设 $X$ 既拟紧又
拟分离。我们将证明此时 $E$ 在 $|X|$ 中可构造。选取开子空间
$$
\emptyset = U_{n + 1} \subset
U_n \subset U_{n - 1} \subset \ldots \subset U_1 = X
$$
及满 \'etale 态射 $f_p : V_p \to U_p$，使其诱导同构
$f_p^{-1}(T_p) \to T_p = U_p \setminus U_{p + 1}$，其中 $V_p$ 是引理
\ref{decent-spaces-lemma-filter-quasi-compact-quasi-separated} 所给出的拟紧分离概形。
按定义，$E$ 的原像 $E_p \subset V_p$ 在 $V_p$ 中局部可构造。由
《性质》引理
\ref{properties-lemma-constructible-quasi-compact-quasi-separated}，
$E_p$ 在 $V_p$ 中可构造。因此，由《拓扑》引理
\ref{topology-lemma-intersect-constructible-with-closed}，
$E_p \cap |f_p^{-1}(T_p)| = E \cap |T_p|$ 在 $|T_p|$ 中可构造（注意，
$V_p \setminus f_p^{-1}(T_p)$ 是拟紧空间 $U_{p + 1}$ 在拟紧态射 $f_p$
下的原像，因而拟紧）。所以
$$
E = (|T_n| \cap E) \cup (|T_{n - 1}| \cap E) \cup \ldots \cup
(|T_1| \cap E)
$$
由《拓扑》引理
\ref{topology-lemma-collate-constructible-from-constructible} 可构造。这里用到
$|T_p|$ 在 $|X|$ 中可构造，这由上述讨论显然可见。
\end{proof}





\section{由概形给出的整覆盖}
\label{decent-spaces-section-integral-cover}

\noindent
这里证明：给定任意拟紧拟分离代数空间 $X$，存在概形 $Y$ 及满整态射
$Y \to X$。在建立代数空间极限的若干理论后，我们将证明还可将其取为有限
态射；参见《空间的极限》第 \ref{spaces-limits-section-finite-cover} 节。

\begin{lemma}
\label{decent-spaces-lemma-extend-integral-morphism}
设 $S$ 为概形。设 $j : V \to Y$ 为 $S$ 上代数空间的拟紧开浸入，并设
$\pi : Z \to V$ 为整态射。则存在整态射 $\nu : Y' \to Y$，使得 $Z$ 在
$V$ 上同构于 $V$ 在 $Y'$ 中的原像。
\end{lemma}

\begin{proof}
由于 $j$ 和 $\pi$ 都拟紧且分离，$j \circ \pi$ 也如此。令
$\nu : Y' \to Y$ 为 $Y$ 在 $Z$ 中的正规化；参见《空间的态射》第
\ref{spaces-morphisms-section-normalization-X-in-Y} 节。当然，$\nu$ 是
整态射；参见《空间的态射》引理
\ref{spaces-morphisms-lemma-characterize-normalization}。最后的断言形式地
得自《空间的态射》引理
\ref{spaces-morphisms-lemma-properties-normalization} 和
\ref{spaces-morphisms-lemma-normalization-in-integral}。
\end{proof}

\begin{lemma}
\label{decent-spaces-lemma-there-is-a-scheme-integral-over}
设 $S$ 为概形，$X$ 为 $S$ 上的拟紧拟分离代数空间。
\begin{enumerate}
\item 存在满整态射 $Y \to X$，其中 $Y$ 为概形；
\item 给定满 \'etale 态射 $U \to X$，可以选取 $Y \to X$，使得对每个
$y \in Y$，都存在开邻域 $V \subset Y$，使 $V \to X$ 经由 $U$ 分解。
\end{enumerate}
\end{lemma}

\begin{proof}
(1) 是 (2) 中 $U = X$ 的特殊情形。取满 \'etale 态射 $U' \to U$，其中
$U'$ 为概形。显然可用 $U'$ 替换 $U$，故可假设 $U$ 为概形。由于 $X$
拟紧，存在有限多个仿射开集 $U_i \subset U$，使得
$U' = \coprod U_i \to X$ 满射。再次用 $U'$ 替换 $U$，可假设 $U$ 仿射。
由于 $X$ 拟分离，因而合理，存在整数 $d$ 界住 $U \to X$ 的几何纤维次数
（参见引理 \ref{decent-spaces-lemma-bounded-fibres}）。我们将对 $d$ 归纳，对所有满且
\'etale 地映到 $X$ 的拟紧分离概形 $U$ 证明本引理。若 $d = 1$，则
$U = X$，取 $Y = U$ 即得结果。假设 $d > 1$。

\medskip\noindent
应用《空间的态射》引理
\ref{spaces-morphisms-lemma-quasi-finite-separated-quasi-affine}，得到分解
$$
\xymatrix{
U \ar[rr]_j \ar[rd] & & Y \ar[ld]^\pi \\
& X
}
$$
其中 $\pi$ 为整态射，$j$ 为拟紧开浸入。我们可以且确实假设 $j(U)$ 在
$Y$ 中概形论稠密。于是 $U \times_X Y$ 是拟紧分离概形（因为它在 $U$
上为整），并且
$$
U \times_X Y = U \amalg W
$$
这里，第一个直和项是 $U \to U \times_X Y$ 的像（由《空间的态射》引理
\ref{spaces-morphisms-lemma-semi-diagonal}，它是闭的；又因为它是 $Y$ 上
\'etale 代数空间之间的 \'etale 态射，所以也是开的），第二个直和项是其
开闭补集。$W$ 的像 $V \subset Y$ 是包含 $Y \setminus U$ 的开子空间。

\medskip\noindent
\'etale 态射 $W \to Y$ 的几何纤维基数为 $< d$。事实上，直接考察即可
看出这对 $U \subset Y$ 的几何点成立。由于 $|U| \subset |Y|$ 稠密，由
引理 \ref{decent-spaces-lemma-quasi-compact-reasonable-stratification}，它对 $Y$ 的所有
几何点都成立（拟紧 \'etale 态射的纤维次数在特化下不会增加）。因此可将
归纳假设用于 $W \to V$，得到满整态射 $Z \to V$，其中 $Z$ 为概形，并且
该态射在 Zariski 局部经由 $W$ 分解。选取分解 $Z \to Z' \to Y$，其中
$Z' \to Y$ 为整态射，$Z \to Z'$ 为开浸入（引理
\ref{decent-spaces-lemma-extend-integral-morphism}）。用 $Z$ 在 $Z'$ 中的概形论闭包替换
$Z'$ 后，可假设 $Z$ 在 $Z'$ 中概形论稠密。此时有
$Z' \times_Y V = Z$。最后，在 $Y \setminus V$ 上赋予诱导闭子空间结构，
记为 $T \subset Y$。考虑态射
$$
Z' \amalg T \longrightarrow X
$$
按构造，这是满整态射。由于 $T \subset U$，显然态射 $T \to X$ 经由
$U$ 分解。另一方面，设 $z \in Z'$ 为一点。若 $z \not \in Z$，则 $z$
映到 $Y \setminus V \subset U$ 的一点，因而可找到 $z$ 的一个邻域，使得
该态射在此邻域上经由 $U$ 分解。若 $z \in Z$，则 $z$ 在 $Z$ 中有一个
开邻域（它也是 $z$ 在 $Z'$ 中的开邻域），该邻域经由
$W \subset U \times_X Y$，从而经由 $U$ 分解。
\end{proof}

\begin{lemma}
\label{decent-spaces-lemma-there-is-a-scheme-integral-over-refined}
设 $S$ 为概形，$X$ 为 $S$ 上的拟紧拟分离代数空间，并且 $|X|$ 只有有限
多个不可约分支。
\begin{enumerate}
\item 存在满整态射 $Y \to X$，其中 $Y$ 为概形，并且
% P08-E289: frozen statement names an undefined f instead of the displayed morphism Y -> X; preserved; see ERRATA_CANDIDATES.jsonl.
$f$ 在某个拟紧稠密开集 $U \subset X$ 上有限且 \'etale；
\item 给定满 \'etale 态射 $V \to X$，可以选取 $Y \to X$，使得对每个
$y \in Y$，都存在开邻域 $W \subset Y$，使 $W \to X$ 经由 $V$ 分解。
\end{enumerate}
\end{lemma}

\begin{proof}
本证明与引理 \ref{decent-spaces-lemma-there-is-a-scheme-integral-over} 的证明大致相同，
但为得到稠密拟紧开集 $U$，需要补充一些技术性说明（并且很遗憾，为了追踪
$U$ 还需改换记号）。

\medskip\noindent
(1) 是 (2) 中 $V = X$ 的特殊情形。

\medskip\noindent
证明 (2)。取满 \'etale 态射 $V' \to V$，其中 $V'$ 为概形。显然可以用
$V'$ 替换 $V$，故可假设 $V$ 为概形。由于 $X$ 拟紧，存在有限多个仿射
开集 $V_i \subset V$，使得 $V' = \coprod V_i \to X$ 满射。再次用 $V'$
替换 $V$ 后，可假设 $V$ 仿射。由于 $X$ 拟分离，因而合理，存在整数 $d$
界住 $V \to X$ 的几何纤维次数（参见引理
\ref{decent-spaces-lemma-bounded-fibres}）。

\medskip\noindent
我们将对 $d \geq 1$ 归纳，证明下列归纳假设 $(H_d)$：
\begin{itemize}
\item 对任意只有有限多个不可约分支的拟紧拟分离代数空间 $X$、任意
$m \geq 0$、任意拟紧分离概形 $V_j$（$j = 1, \ldots, m$），以及任意
\'etale 态射 $\varphi_j : V_j \to X$（$j = 1, \ldots, m$），若 $d$ 界住
$\varphi_j : V_j\to X$ 的几何纤维次数，且
$\varphi = \coprod \varphi_j : V = \coprod V_j \to X$ 满射，则本引理的
断言对 $\varphi : V \to X$ 成立。
\end{itemize}
若 $d = 1$，则每个 $\varphi_j$ 都是开浸入。因此 $X$ 是概形，并且
% P08-E290: frozen base case takes Y = V although the surjective union of open immersions V -> X need not be integral; formula preserved; see ERRATA_CANDIDATES.jsonl.
取 $Y = V$ 即得结果。假设 $d > 1$，假设 $(H_{d - 1})$，并令
$m$、$\varphi : V_j \to X$（$j = 1, \ldots, m$）如 $(H_d)$ 中所述。

\medskip\noindent
设 $\eta_1, \ldots, \eta_n \in |X|$ 为 $|X|$ 各不可约分支的一般点。由
《空间的性质》命题
\ref{spaces-properties-proposition-locally-quasi-separated-open-dense-scheme}，
存在开子概形 $U \subset X$，使得 $\eta_1, \ldots, \eta_n \in U$。缩小
$U$ 后，可假设 $U$ 仿射；再由《态射》引理
\ref{morphisms-lemma-generically-finite}，可假设每个
$\varphi_j : V_j \to X$ 在 $U$ 上有限且 \'etale。当然，$U$ 拟紧且在
$X$ 中稠密，并且 $\varphi_j^{-1}(U)$ 在 $V_j$ 中稠密。特别地，每个
$V_j$ 只有有限多个不可约分支。

\medskip\noindent
固定 $j \in \{1, \ldots, m\}$。如《空间的态射》引理
\ref{spaces-morphisms-lemma-quasi-finite-separated-quasi-affine} 中那样，
令 $Y_j$ 为 $X$ 在 $V_j$ 中的正规化。得到分解
$$
\xymatrix{
V_j \ar[rr] \ar[rd]_{\varphi_j} & & Y_j \ar[ld]^{\pi_j} \\
& X
}
$$
其中 $\pi_j$ 为整态射，$V_j \to Y_j$ 为拟紧开浸入。由于 $Y_j$ 是 $X$
在 $V_j$ 中的正规化，由《空间的态射》引理
\ref{spaces-morphisms-lemma-properties-normalization} 和
\ref{spaces-morphisms-lemma-normalization-in-integral}
可知 $\varphi_j^{-1}(U) \to \pi_j^{-1}(U)$ 是同构。因此 $\pi_j$ 在 $U$
上有限且 \'etale。注意，$V_j$ 在 $Y_j$ 中概形论稠密，因为 $Y_j$ 是
$X$ 在 $V_j$ 中的正规化（这由《空间的态射》引理
\ref{spaces-morphisms-lemma-characterize-normalization} 对相对正规化的刻画
得出）。由于 $V_j$ 拟紧，可知 $|V_j| \subset |Y_j|$ 稠密；参见
《空间的态射》第 \ref{spaces-morphisms-section-scheme-theoretic-closure} 节
（尤其是《空间的态射》引理
\ref{spaces-morphisms-lemma-quasi-compact-immersion}）。从而 $|Y_j|$ 只有
有限多个不可约分支。于是 $V_j \times_X Y_j$ 是拟紧分离概形
% P08-E291: frozen source says finite over V_j although only integrality of the normalization map is available here; preserved; see ERRATA_CANDIDATES.jsonl.
（因为它在 $V_j$ 上有限），并且
$$
V_j \times_X Y_j = V_j \amalg W_j
$$
这里，第一个直和项是 $V_j \to V_j \times_X Y_j$ 的像（由《空间的态射》
引理 \ref{spaces-morphisms-lemma-semi-diagonal}，它是闭的；又因为它是
% P08-E292: frozen source says the spaces are étale over Y instead of Y_j; formula preserved; see ERRATA_CANDIDATES.jsonl.
$Y$ 上 \'etale 代数空间之间的 \'etale 态射，所以也是开的），第二个
直和项是其开闭补集。

\medskip\noindent
\'etale 态射 $W_j \to Y_j$ 的几何纤维基数为 $< d$。事实上，直接考察
即可看出这对 $V_j \subset Y_j$ 的几何点成立。由于
$|V_j| \subset |Y_j|$ 稠密，由引理
\ref{decent-spaces-lemma-quasi-compact-reasonable-stratification}，它对 $Y_j$ 的所有几何点
都成立（拟紧 \'etale 态射的纤维次数在特化下不会增加）。
% P08-E294: frozen proof applies H_{d-1} without establishing the required surjectivity of V_j \amalg W_j -> Y_j; preserved; see ERRATA_CANDIDATES.jsonl.
将
$(H_{d - 1})$ 用于 $V_j \amalg W_j \to Y_j$，得到满整态射
$Y_j' \to Y_j$，其中 $Y_j'$ 为概形；该态射在 Zariski 局部经由
$V_j \amalg W_j$ 分解，并且在某个拟紧稠密开集 $U_j \subset Y_j$ 上有限
且 \'etale。缩小 $U$ 后，我们可以且确实假设
$\pi_j^{-1}(U) \subset U_j$（对所有 $j$ 可以且确实选取同一个 $U$；略去
若干细节）。

\medskip\noindent
我们断言
$$
Y = \coprod\nolimits_{j = 1, \ldots, m} Y'_j \longrightarrow X
$$
解决了问题。首先，该态射为整态射，因为在每个直和项上都有整态射的复合
% P08-E293: frozen source routes the summand through the coproduct Y rather than its component Y_j; formula preserved; see ERRATA_CANDIDATES.jsonl.
$Y'_j \to Y \to X$（《空间的态射》引理
\ref{spaces-morphisms-lemma-composition-integral}）。其次，该态射在 Zariski
局部经由 $V = \coprod V_j$ 分解，因为上面已经看到，每个
$Y'_j \to Y_j$ 在 Zariski 局部经由
$V_j \amalg W_j = V_j \times_X Y_j$ 分解。最后，由于
$Y'_j \to Y_j$ 和 $Y_j \to X$ 在 $U$ 上都有限且 \'etale，其复合也如此。
证明完成。
\end{proof}



















\section{概形轨迹}
\label{decent-spaces-section-schematic}

\noindent
本节证明良好代数空间有一个作为概形的稠密开子空间。我们先对合理代数空间
证明这一点。


\begin{proposition}
\label{decent-spaces-proposition-reasonable-open-dense-scheme}
设 $S$ 为概形，$X$ 为 $S$ 上的代数空间。若 $X$ 合理，则 $X$ 存在一个
作为概形的稠密开子空间。
\end{proposition}

\begin{proof}
由《空间的性质》引理 \ref{spaces-properties-lemma-subscheme}，此问题关于
$X$ 是局部的。因此可假设存在仿射概形 $U$ 及满 \'etale 态射
$U \to X$（《空间的性质》引理
\ref{spaces-properties-lemma-cover-by-union-affines}）。由于 $X$ 合理，存在
整数 $n$ 界住 $U \to X$ 各纤维的次数；参见定义
\ref{decent-spaces-definition-very-reasonable}。我们对 $n$ 归纳证明，只要
\begin{enumerate}
\item $U \to X$ 是满 \'etale 态射，且各纤维次数为 $\leq n$；
\item $U$ 同构于某个仿射概形的局部闭子概形，
\end{enumerate}
则概形轨迹在 $X$ 中稠密。

\medskip\noindent
令 $X_n \subset X$ 为闭子空间 $Z_{n - 1} \subset X$ 的补开子空间，其中
该闭子空间是引理 \ref{decent-spaces-lemma-quasi-compact-reasonable-stratification} 对态射
$U \to X$ 构造的。令 $U_n \subset U$ 为 $X_n$ 的原像。于是
$U_n \to X_n$ 是次数为 $n$ 的有限局部自由态射。因此，由《空间的性质》
命题
\ref{spaces-properties-proposition-finite-flat-equivalence-global}
以及 $U_n$ 的任意有限点集都包含在 $U_n$ 的某个仿射开集内这一事实（参见
《性质》引理 \ref{properties-lemma-ample-finite-set-in-affine}），$X_n$ 是
概形。

\medskip\noindent
令 $X' \subset X$ 为开子空间，使得 $|X'|$ 是 $|Z_{n - 1}|$ 在 $|X|$ 中的
内部（参见《拓扑》定义 \ref{topology-definition-nowhere-dense}）。令
$U' \subset U$ 为其原像。于是 $U' \to X'$ 满且 \'etale，纤维次数由
$n - 1$ 界住。由归纳假设，$X'$ 的概形轨迹是某个稠密开集
$X'' \subset X'$。由初等拓扑可知，$X'' \cup X_n \subset X$ 开且稠密，
从而得证。
\end{proof}

\begin{theorem}[David Rydh]
\label{decent-spaces-theorem-decent-open-dense-scheme}
设 $S$ 为概形，$X$ 为 $S$ 上的代数空间。若 $X$ 良好，则 $X$ 存在一个
作为概形的稠密开子空间。
\end{theorem}

\begin{proof}
假设定理对某个良好代数空间 $X$ 不成立。由《空间的性质》引理
\ref{spaces-properties-lemma-subscheme}，存在最大的作为概形的开子空间
$X' \subset X$。由于 $X'$ 在 $X$ 中不稠密，存在开子空间
$X'' \subset X$，使得 $|X''| \cap |X'| = \emptyset$。用 $X''$ 替换 $X$
后，得到一个非空良好代数空间 $X$，其中不含{\it 任何}作为概形的开子空间。

\medskip\noindent
取非空仿射概形 $U$ 及 \'etale 态射 $U \to X$。
% P08-E295: frozen source calls the image replacement an open subscheme although it is only known to be an open subspace; intended wording translated; see ERRATA_CANDIDATES.jsonl.
我们可以且确实用与 $|U| \to |X|$ 的像对应的开子空间替换 $X$。考虑引理
\ref{decent-spaces-lemma-stratify-flat-fp-lqf} 对态射 $U \to X$ 构造的开子空间序列
% P08-E296: frozen display omits the inclusion symbol after X_2; formula preserved; see ERRATA_CANDIDATES.jsonl.
$$
X = X_0 \supset X_1 \supset X_2 \ldots
$$
。注意，由于 $U \to X$ 满射，有 $X_0 = X_1$。令
$U = U_0 = U_1 \supset U_2 \ldots$ 为 $U$ 中相应的开子概形序列。

\medskip\noindent
取非空仿射开集 $V_1 \subset U_1$（例如 $V_1 = U_1$）。我们将归纳构造
非空仿射开集序列 $V_1 \supset V_2 \supset \ldots$，其中
$V_n \subset U_n$。事实上，在已构造 $V_1, \ldots, V_{n - 1}$ 后，除非
$V_{n - 1} \cap U_n = \emptyset$，总能选取 $V_n$。但若
$V_{n - 1} \cap U_n = \emptyset$，则满足
$|X'| = \Im(|V_{n - 1}| \to |X|)$ 的开子空间 $X' \subset X$ 包含在
$|X| \setminus |X_n|$ 中。因此，$V_{n - 1} \to X'$ 是纤维次数由
$n - 1$ 界住的 \'etale 态射。换言之，$X'$ 按定义合理，故由命题
\ref{decent-spaces-proposition-reasonable-open-dense-scheme}，$X'$ 含有非空开子概形。这是
矛盾，表明可以选取 $V_n$。

\medskip\noindent
由《极限》引理 \ref{limits-lemma-limit-nonempty}，极限
$V_\infty = \lim V_n$ 是非空概形。取态射 $\Spec(k) \to V_\infty$。按构造，
复合态射 $\Spec(k) \to V_\infty \to U \to X$ 的像包含在每个 $X_d$ 中。
换言之，
% P08-E297: frozen source says “the fibred” where the fibre U x_X Spec(k) is meant; intended noun translated; see ERRATA_CANDIDATES.jsonl.
纤维 $U \times_X \Spec(k)$ 的次数无限，这与良好空间的定义矛盾。此矛盾
完成定理的证明。
\end{proof}

\begin{lemma}
\label{decent-spaces-lemma-when-quotient-scheme-at-point}
设 $S$ 为概形，$X \to Y$ 为 $S$ 上代数空间的满有限局部自由态射。对
$y \in |Y|$，下列条件等价：
\begin{enumerate}
\item $y$ 属于 $Y$ 的概形轨迹；
\item 存在仿射开集 $U \subset X$，包含 $y$ 的原像。
\end{enumerate}
\end{lemma}

\begin{proof}
% P08-E298: frozen proof types the point as y in Y rather than y in |Y|; formula preserved; see ERRATA_CANDIDATES.jsonl.
若 $y \in Y$ 属于概形轨迹，则它有仿射开邻域 $V \subset Y$；$V$ 在 $X$
中的原像 $U$ 是在 $V$ 上有限的开集，因而仿射。所以 (1) 蕴含 (2)。

\medskip\noindent
反过来，假设给定 (2) 中的 $U \subset X$。令 $R = X \times_Y X$，并将投影
记为 $s, t : R \to X$。考虑
% P08-E299: frozen set difference lacks parentheses around the intersection; formula preserved; see ERRATA_CANDIDATES.jsonl.
$Z = R \setminus s^{-1}(U) \cap t^{-1}(U)$。这是 $R$ 的闭子集。像 $t(Z)$
是 $X$ 的闭子集，可以粗略地描述为 $X$ 中与 $X \setminus U$ 的某点
$R$-等价的点集。因此，$U' = X \setminus t(Z)$ 是 $X$ 中包含在 $U$ 内的
$R$-不变开子空间，并且包含 $X \to Y$ 在 $y$ 上的纤维。由于 $X \to Y$
是开态射（《空间的态射》引理
\ref{spaces-morphisms-lemma-fppf-open}），$U'$ 的像是开子空间
$V' \subset Y$。由于 $U'$ 是 $R$-不变的，且 $R = X \times_Y X$，可知
$U'$ 是 $V'$ 的原像（使用《空间的性质》引理
\ref{spaces-properties-lemma-points-cartesian}）。用 $V'$ 替换 $Y$、用 $U'$
替换 $X$ 后，可假设 $X$ 为同构于某个仿射概形的开子概形的概形。

\medskip\noindent
假设 $X$ 为同构于某个仿射概形的开子概形的概形。此时，fppf 商层 $X/R$
是概形；参见《空间的性质》命题
\ref{spaces-properties-proposition-finite-flat-equivalence-global}。由于 $Y$
是 fppf 拓扑中的层，得到一个使 $X \to Y$ 分解的典范映射 $X/R \to Y$。
由于 $X \to Y$ 满且有限局部自由，它作为层的映射也满（《空间》引理
\ref{spaces-lemma-surjective-flat-locally-finite-presentation}）。因此
$X/R \to Y$ 作为层的映射满。另一方面，由于作为层有
$R = X \times_Y X$，可知 $X/R \to Y$ 作为层的映射单。因此
$X/R \to Y$ 是同构，进而 $Y$ 可表。
\end{proof}

\noindent
至此，我们已有数种不同的方法证明下述引理。

\begin{lemma}
\label{decent-spaces-lemma-finite-etale-cover-dense-open-scheme}
设 $S$ 为概形，$X$ 为 $S$ 上的代数空间。若存在有限、\'etale 且满的态射
$U \to X$，其中 $U$ 为概形，则 $X$ 存在一个作为概形的稠密开子空间。
\end{lemma}

\begin{proof}[第一种证明]
态射 $U \to X$ 有限局部自由。因此，$X$ 可分解为开闭子空间
$X_d \subset X$，使得 $U \times_X X_d \to X_d$ 是次数为 $d$ 的有限局部
自由态射。故可假设 $U \to X$ 是次数为 $d$ 的有限局部自由态射。此时，令
$U_i \subset U$（$i \in I$）遍历仿射开集。对每个 $i$，态射
$U_i \to X$ 都是 \'etale 的，并有普遍有界纤维（具体地，由 $d$ 界住）。
换言之，$X$ 合理，结论由命题
\ref{decent-spaces-proposition-reasonable-open-dense-scheme} 得出。
\end{proof}

\begin{proof}[第二种证明]
此问题关于 $X$ 是局部的（《空间的性质》引理
\ref{spaces-properties-lemma-subscheme}），故可假设 $X$ 拟紧。于是 $U$ 拟紧，
从而存在稠密开子概形 $W \subset U$，且它分离（《性质》引理
\ref{properties-lemma-quasi-compact-dense-open-separated}）。令
$Z = U \setminus W$。令 $R = U \times_X U$，并将投影记为
$s, t : R \to U$。则 $t^{-1}(Z)$ 在 $R$ 中无处稠密（《拓扑》引理
\ref{topology-lemma-open-inverse-image-closed-nowhere-dense}），因而
$\Delta = s(t^{-1}(Z))$ 是 $U$ 的 $R$-不变无处稠密闭子集（《态射》引理
\ref{morphisms-lemma-image-nowhere-dense-finite}）。取
$u \in U \setminus \Delta$ 为某个不可约分支的一般点。由于这些点在
$U \setminus \Delta$ 中稠密，且 $\Delta$ 无处稠密，只需证明 $u$ 的像
$x \in X$ 属于 $X$ 的概形轨迹。

注意，$t(s^{-1}(\{u\})) \subset W$ 是 $W$ 的不可约分支的一般点所成的
有限集（比较《空间的性质》引理
\ref{spaces-properties-lemma-codimension-0-points}）。由《性质》引理
\ref{properties-lemma-maximal-points-affine}，可找到仿射开集 $V \subset W$，
使得 $t(s^{-1}(\{u\})) \subset V$。由于 $t(s^{-1}(\{u\}))$ 是
$|U| \to |X|$ 在 $x$ 上的纤维，结论由引理
\ref{decent-spaces-lemma-when-quotient-scheme-at-point} 得出。
\end{proof}

\begin{proof}[第三种证明]
（此证明本质上与第二种证明相同，但使用的引用较少。）
假设 $X$ 是代数空间，$U$ 是概形，且 $U \to X$ 是有限、\'etale 且满的态射。
记 $R = U \times_X U$，并照例将投影记为 $s, t : R \to U$。注意 $s,t$ 满、
有限且 \'etale。断言：$U$ 的所有 $R$-不变仿射开集之并在 $U$ 中拓扑稠密。

\medskip\noindent
证明该断言。设 $W \subset U$ 为仿射开集。令
$W' = t(s^{-1}(W)) \subset U$。由于 $s^{-1}(W)$ 仿射（因而拟紧），可知
$W' \subset U$ 是拟紧开集。由《性质》引理
\ref{properties-lemma-quasi-compact-dense-open-separated}，存在稠密开集
$W'' \subset W'$，且它是分离概形。令 $\Delta' = W' \setminus W''$。
% P08-E300: frozen source says Delta' is closed nowhere dense in W'', but it is the complement of W'' in W'; frozen formula preserved; see ERRATA_CANDIDATES.jsonl.
这是 $W''$ 的无处稠密闭子集。由于
$t|_{s^{-1}(W)} : s^{-1}(W) \to W'$ 是开态射（因为它是 \'etale 的），可知
原像 $(t|_{s^{-1}(W)})^{-1}(\Delta') \subset s^{-1}(W)$ 是无处稠密闭子集
（参见《拓扑》引理
\ref{topology-lemma-open-inverse-image-closed-nowhere-dense}）。因此，由《态射》
引理 \ref{morphisms-lemma-image-nowhere-dense-finite} 可知
$$
\Delta = s\left((t|_{s^{-1}(W)})^{-1}(\Delta')\right)
$$
是 $W$ 的无处稠密闭子集。取点 $\eta \in W$，使得
$\eta \not \in \Delta$，并且它是 $W$（因而也是 $U$）某个不可约分支的
一般点。按上述选择，有限集
$t(s^{-1}(\{\eta\})) = \{\eta_1, \ldots, \eta_n\}$ 包含在分离概形 $W''$ 中。
% P08-E301: frozen source has the malformed phrase “the fibres of s is are”; intended assertion translated; see ERRATA_CANDIDATES.jsonl.
注意，$s$ 的纤维是有限离散空间，并且一般化可沿 \'etale 态射 $t$ 提升；参见
《态射》引理 \ref{morphisms-lemma-etale-flat} 和
\ref{morphisms-lemma-generalizations-lift-flat}。由此可知，每个 $\eta_i$ 都是
$W''$ 某个不可约分支的一般点。因此，由《性质》引理
\ref{properties-lemma-maximal-points-affine}，可找到仿射开集 $V \subset W''$，
使得 $\{\eta_1, \ldots, \eta_n\} \subset V$。由《群胚》引理
\ref{groupoids-lemma-find-invariant-affine}，这说明 $\eta$ 包含在 $U$ 的某个
$R$-不变仿射开子概形中。由于 $W$ 是 $U$ 的任意仿射开集，且
$W \setminus \Delta$ 的不可约分支的一般点集在 $W$ 中稠密，断言得证。

\medskip\noindent
用该断言即可完成证明。事实上，若 $W \subset U$ 是 $R$-不变仿射开集，则
$R$ 在 $W$ 上的限制 $R_W$ 满足
$R_W = s^{-1}(W) = t^{-1}(W)$（参见《群胚》定义
\ref{groupoids-definition-invariant-open} 及其后的讨论）。特别地，映射
$R_W \to W$ 也有限且 \'etale；因此 $R_W$ 仿射。由《群胚》命题
\ref{groupoids-proposition-finite-flat-equivalence}，$W/R_W$ 是概形。另一方面，
由《空间》引理 \ref{spaces-lemma-finding-opens}，$W/R_W$ 是 $X$ 的开子空间。
% P08-E302: frozen source omits an article/plural marker in “contained in R-invariant affine open”; intended grammar translated; see ERRATA_CANDIDATES.jsonl.
% P08-E303: frozen source says “is open dense” rather than “is open and dense”; intended grammar translated; see ERRATA_CANDIDATES.jsonl.
故若稠密点集中的每一点都包含在 $U$ 的某个 $R$-不变仿射开集中，则 $X$ 的
概形轨迹（参见《空间的性质》引理
\ref{spaces-properties-lemma-subscheme}）必在 $X$ 中开且稠密。
\end{proof}













\section{剩余域与 Hensel 局部环}
\label{decent-spaces-section-residue-fields-henselian-local-rings}

\noindent
对良好代数空间，可以定义一点处的剩余域与 Hensel 局部环。例如，下述引理
说明良好空间上一点的剩余域确有定义。

\begin{lemma}
\label{decent-spaces-lemma-decent-points-monomorphism}
设 $S$ 为概形，$X$ 为 $S$ 上的代数空间。考虑映射
$$
\{\Spec(k) \to X \text{ 为单态射，其中 }k\text{ 为域}\}
\longrightarrow
|X|
$$
。此映射总为单射。若 $X$ 良好，则它是双射。
\end{lemma}

\begin{proof}
《空间的性质》引理
\ref{spaces-properties-lemma-points-monomorphism} 已表明该映射一般总为单射。
由引理 \ref{decent-spaces-lemma-bounded-fibres}，当 $X$ 良好时它为满射（其实也可以说，
这是良好性定义的一部分）。
\end{proof}

\noindent
设 $S$ 为概形，$X$ 为 $S$ 上的代数空间。若点 $x \in |X|$ 可由单态射
$\Spec(k) \to X$ 表示，则域 $k$ 在唯一同构意义下唯一。对良好代数空间，
由引理 \ref{decent-spaces-lemma-decent-points-monomorphism}，每一点都存在这种单态射，
所以下述定义有意义。

\begin{definition}
\label{decent-spaces-definition-residue-field}
设 $S$ 为概形，$X$ 为 $S$ 上的良好代数空间。设 $x \in |X|$。{\it $X$ 在
$x$ 处的剩余域}是唯一的域 $\kappa(x)$，它带有表示 $x$ 的单态射
$\Spec(\kappa(x)) \to X$。
\end{definition}

\noindent
设 $S$ 为概形，$f : X \to Y$ 为 $S$ 上良好代数空间的态射。设
$x \in |X|$ 为一点，并令 $y = f(x) \in |Y|$。则复合态射
$\Spec(\kappa(x)) \to Y$ 属于定义 $y$ 的等价类，因而经由
$\Spec(\kappa(y)) \to Y$ 分解。换言之，得到交换图
$$
\xymatrix{
\Spec(\kappa(x)) \ar[r]_-x \ar[d] & X \ar[d]^f \\
\Spec(\kappa(y)) \ar[r]^-y & Y
}
$$
左侧竖直态射对应于域同态 $\kappa(y) \to \kappa(x)$。我们常简称它为
$f$ 诱导的同态。

\begin{lemma}
\label{decent-spaces-lemma-identifies-residue-fields}
设 $S$ 为概形，$f : X \to Y$ 为 $S$ 上良好代数空间的态射。设
$x \in |X|$ 为一点，其像为 $y = f(x) \in |Y|$。下列条件等价：
\begin{enumerate}
\item $f$ 诱导同构 $\kappa(y) \to \kappa(x)$；
\item 诱导态射 $\Spec(\kappa(x)) \to Y$ 是单态射。
\end{enumerate}
\end{lemma}

\begin{proof}
由上述讨论立即得出。
\end{proof}

\noindent
下述引理说明良好代数空间上一点的 Hensel 局部环确有定义。

\begin{lemma}
\label{decent-spaces-lemma-decent-space-elementary-etale-neighbourhood}
设 $S$ 为概形，$X$ 为 $S$ 上的良好代数空间。对每一点 $x \in |X|$，
存在 \'etale 态射
$$
(U, u) \longrightarrow (X, x)
$$
，其中 $U$ 为仿射概形，$u$ 是 $U$ 中位于 $x$ 上方的唯一一点，并且诱导同态
$\kappa(x) \to \kappa(u)$ 是同构。
\end{lemma}

\begin{proof}
可用包含 $x$ 的拟紧开集替换 $X$，从而假设 $X$ 拟紧。回忆，由良好空间的
定义，$x$ 可由从某个域的谱出发的拟紧（单）态射表示。
% P08-E305: frozen source omits “of” in “from the spectrum a field”; intended grammar translated; see ERRATA_CANDIDATES.jsonl.
故引理由引理 \ref{decent-spaces-lemma-filter-quasi-compact} 得出。
\end{proof}

\begin{definition}
\label{decent-spaces-definition-elemenary-etale-neighbourhood}
设 $S$ 为概形，$X$ 为 $S$ 上的代数空间。设 $x \in X$ 为一点。
% P08-E304: frozen definition types the point as x in X rather than x in |X|; formula preserved; see ERRATA_CANDIDATES.jsonl.
{\it 初等 \'etale 邻域}是 \'etale 态射 $(U, u) \to (X, x)$，其中 $U$
为概形，$u \in U$ 为映到 $x$ 的一点，且态射
$u = \Spec(\kappa(u)) \to X$ 是单态射。{\it 初等 \'etale 邻域的态射}
$(U, u) \to (U', u')$ 定义为 $X$ 上将 $u$ 映到 $u'$ 的态射 $U \to U'$。
\end{definition}

\noindent
若 $X$ 不良好，则初等 \'etale 邻域范畴可能为空。

\begin{lemma}
\label{decent-spaces-lemma-elementary-etale-neighbourhoods}
设 $S$ 为概形，$X$ 为 $S$ 上的良好代数空间，并设 $x$ 为 $X$ 的一点。
$(X, x)$ 的初等 \'etale 邻域范畴是余滤的（参见《范畴》定义
\ref{categories-definition-codirected}）。
\end{lemma}

\begin{proof}
由引理 \ref{decent-spaces-lemma-decent-space-elementary-etale-neighbourhood}，该范畴非空。
假设有两个初等 \'etale 邻域 $(U_i, u_i) \to (X, x)$。考虑
$U = U_1 \times_X U_2$。由于 $\Spec(\kappa(u_i)) \to X$（$i = 1, 2$）
都是 $x$ 所在等价类中的单态射（引理
\ref{decent-spaces-lemma-identifies-residue-fields}），可知
% P08-E306: frozen source places a detached comma after the closing parenthesis; intended punctuation translated; see ERRATA_CANDIDATES.jsonl.
$$
u = \Spec(\kappa(u_1)) \times_X \Spec(\kappa(u_2))
$$
是某个域 $\kappa(u)$ 的谱，使得诱导映射 $\kappa(u_i) \to \kappa(u)$ 是
同构。于是 $u \to U$ 是 $U$ 的一点，并且
$(U, u) \to (X, x)$ 是支配 $(U_i, u_i)$ 的初等 \'etale 邻域。
若 $a, b : (U_1, u_1) \to (U_2, u_2)$ 是这两个初等 \'etale 邻域之间的
两个态射，则考虑概形
$$
U = U_1 \times_{(a, b), (U_2 \times_X U_2), \Delta} U_2
$$
。由《空间的性质》引理
\ref{spaces-properties-lemma-etale-permanence}，$U \to X$ 是 \'etale 的。
此外，完全同前可知 $U$ 有一点 $u$，使得
$(U, u) \to (X, x)$ 是初等 \'etale 邻域。最后，$U \to U_1$ 等化 $a$ 与
$b$，证明完毕。
\end{proof}

\begin{definition}
\label{decent-spaces-definition-henselian-local-ring}
设 $S$ 为概形，$X$ 为 $S$ 上的良好代数空间，并设 $x \in |X|$。{\it $X$
在 $x$ 处的 Hensel 局部环}为
% P08-E308: frozen source has a stray comma between the definition's subject and “is”; intended punctuation translated; see ERRATA_CANDIDATES.jsonl.
$$
\mathcal{O}_{X, x}^h = \colim \Gamma(U, \mathcal{O}_U)
$$
，其中余极限取遍初等 \'etale 邻域 $(U, u) \to (X, x)$。
\end{definition}

\noindent
下面是《空间的性质》引理
\ref{spaces-properties-lemma-describe-etale-local-ring} 的对应结果。

\begin{lemma}
\label{decent-spaces-lemma-describe-henselian-local-ring}
设 $S$ 为概形，$X$ 为 $S$ 上的良好代数空间，并设 $x \in |X|$。设
$(U, u) \to (X, x)$ 为初等 \'etale 邻域。则
$$
\mathcal{O}_{X, x}^h = \mathcal{O}_{U, u}^h
$$
。换言之，$X$ 在 $x$ 处的 Hensel 局部环等于 $U$ 在 $u$ 处的局部环
$\mathcal{O}_{U, u}$ 的 Hensel 化 $\mathcal{O}_{U, u}^h$。
\end{lemma}

\begin{proof}
由于 $(X, x)$ 的初等 \'etale 邻域范畴余滤（引理
\ref{decent-spaces-lemma-elementary-etale-neighbourhoods}），可知 $(U, u)$ 的初等
\'etale 邻域范畴在 $(X, x)$ 的初等 \'etale 邻域范畴中为初始的。
% P08-E307: frozen proof twice uses singular “neighbourhood” for a category of neighbourhoods; intended plurals translated; see ERRATA_CANDIDATES.jsonl.
于是该等式由《态射补篇》引理
\ref{more-morphisms-lemma-describe-henselization} 和《范畴》引理
\ref{categories-lemma-cofinal} 得出（由于定义 Hensel 局部环的余极限取在
初等 \'etale 邻域范畴的反范畴上，初始性转化为共尾性）。
\end{proof}

\begin{lemma}
\label{decent-spaces-lemma-henselian-local-ring-strict}
设 $S$ 为概形，$X$ 为 $S$ 上的良好代数空间。设 $\overline{x}$ 为 $X$
位于 $x \in |X|$ 上方的几何点。$X$ 在 $\overline{x}$ 处的 \'etale 局部环
$\mathcal{O}_{X, \overline{x}}$（《空间的性质》定义
\ref{spaces-properties-definition-etale-local-rings}）是 $X$ 在 $x$ 处的
Hensel 局部环 $\mathcal{O}_{X, x}^h$ 的严格 Hensel 化。
\end{lemma}

\begin{proof}
这由引理 \ref{decent-spaces-lemma-describe-henselian-local-ring}、《空间的性质》引理
\ref{spaces-properties-lemma-describe-etale-local-ring} 以及下述事实得出：
对局部环 $(R, \mathfrak m, \kappa)$ 和给定的 $\kappa$ 的可分代数闭包
$\kappa^{sep}$，有 $(R^h)^{sh} = R^{sh}$。该等式由《代数》引理
\ref{algebra-lemma-uniqueness-henselian} 得出。
\end{proof}

\begin{lemma}
\label{decent-spaces-lemma-residue-field-henselian-local-ring}
设 $S$ 为概形，$X$ 为 $S$ 上的良好代数空间，并设 $x \in |X|$。
$X$ 在 $x$ 处的 Hensel 局部环（定义
\ref{decent-spaces-definition-henselian-local-ring}）的剩余域，就是 $X$ 在 $x$ 处的
剩余域（定义 \ref{decent-spaces-definition-residue-field}）。
\end{lemma}

\begin{proof}
取初等 \'etale 邻域 $(U, u) \to (X, x)$。则 $\kappa(u) = \kappa(x)$，且
$\mathcal{O}_{X, x}^h = \mathcal{O}_{U, u}^h$
（引理 \ref{decent-spaces-lemma-describe-henselian-local-ring}）。由《代数》引理
\ref{algebra-lemma-henselization}，$\mathcal{O}_{U, u}^h$ 的剩余域为
$\kappa(u)$（该引理的结论就是局部环 Hensel 化的构造/定义；参见《代数》定义
\ref{algebra-definition-henselization}）。
\end{proof}

\begin{remark}
\label{decent-spaces-remark-functoriality-henselian-local-ring}
设 $S$ 为概形，$f : X \to Y$ 为 $S$ 上良好代数空间的态射。设
$x \in |X|$，其像为 $y \in |Y|$。取初等 \'etale 邻域
$(V, v) \to (Y, y)$（由引理
\ref{decent-spaces-lemma-decent-space-elementary-etale-neighbourhood} 可以如此选取）。
则 $V \times_Y X$ 是 $X$ 上的 \'etale 代数空间，它有唯一一点 $x'$，在
$X$ 中映到 $x$，在 $V$ 中映到 $v$。（细节从略；使用所有点都可由从域的谱
出发的单态射表示这一事实。）取初等 \'etale 邻域
$(U, u) \to (V \times_Y X, x')$。于是得到交换图
% P08-E310: frozen remark uses compatible geometric points \bar{x}, \bar{y} without introducing them; formulas preserved; see ERRATA_CANDIDATES.jsonl.
$$
\xymatrix{
\Spec(\mathcal{O}_{X, \overline{x}}) \ar[r] \ar[d] &
\Spec(\mathcal{O}_{X, x}^h) \ar[r] \ar[d] &
\Spec(\mathcal{O}_{U, u}) \ar[r] \ar[d] &
U \ar[r] \ar[d] &
X \ar[d] \\
\Spec(\mathcal{O}_{Y, \overline{y}}) \ar[r] &
\Spec(\mathcal{O}_{Y, y}^h) \ar[r] &
\Spec(\mathcal{O}_{V, v}) \ar[r] &
V \ar[r] &
Y
}
$$
它来自下列等同：
$\mathcal{O}_{X, \overline{x}} = \mathcal{O}_{U, u}^{sh}$,
$\mathcal{O}_{X, x}^h = \mathcal{O}_{U, u}^h$,
$\mathcal{O}_{Y, \overline{y}} = \mathcal{O}_{V, v}^{sh}$,
$\mathcal{O}_{Y, y}^h = \mathcal{O}_{V, v}^h$
；参见引理 \ref{decent-spaces-lemma-describe-henselian-local-ring}、《空间的性质》
引理 \ref{spaces-properties-lemma-describe-etale-local-ring}，以及《代数》
第 \ref{algebra-section-ind-etale} 节和第
\ref{algebra-section-henselization} 节所讨论的（严格）Hensel 化函子性。
% P08-E309: frozen source has the malformed phrase “see in Lemma”; intended grammar translated; see ERRATA_CANDIDATES.jsonl.
\end{remark}









\section{良好空间上的点}
\label{decent-spaces-section-points}

\noindent
本节证明良好代数空间上的点的若干性质。下述引理表明，点的特化在良好代数
空间上的表现良好。《空间》例
\ref{spaces-example-infinite-product} 表明，一般而言这{\bf 并不}成立。

\begin{lemma}
\label{decent-spaces-lemma-decent-no-specializations-map-to-same-point}
设 $S$ 为概形，$X$ 为 $S$ 上的良好代数空间，并设 $U \to X$ 为从概形到
$X$ 的 \'etale 态射。若 $u, u' \in |U|$ 映到 $|X|$ 的同一点，且
$u' \leadsto u$，则 $u = u'$。
\end{lemma}

\begin{proof}
结合引理 \ref{decent-spaces-lemma-bounded-fibres} 与
\ref{decent-spaces-lemma-no-specializations-map-to-same-point}。
\end{proof}

\begin{lemma}
\label{decent-spaces-lemma-decent-specialization}
设 $S$ 为概形，$X$ 为 $S$ 上的良好代数空间。设 $x, x' \in |X|$，并假设
$x' \leadsto x$，即 $x$ 是 $x'$ 的特化。则对从概形 $U$ 出发的任意
\'etale 态射 $\varphi : U \to X$ 及满足 $\varphi(u) = x$ 的任意
$u \in U$，存在点 $u'\in U$，使得 $u' \leadsto u$ 且
$\varphi(u') = x'$。
% P08-E311: frozen source omits “there” before “exists a point” and omits spacing in u'\in U; intended prose translated while formula bytes are preserved; see ERRATA_CANDIDATES.jsonl.
\end{lemma}

\begin{proof}
结合引理 \ref{decent-spaces-lemma-bounded-fibres} 与
\ref{decent-spaces-lemma-specialization}。
\end{proof}

\begin{lemma}
\label{decent-spaces-lemma-kolmogorov}
设 $S$ 为概形，$X$ 为 $S$ 上的良好代数空间。则 $|X|$ 是 Kolmogorov
空间（参见《拓扑》定义 \ref{topology-definition-generic-point}）。
\end{lemma}

\begin{proof}
设 $x_1, x_2 \in |X|$，且 $x_1 \leadsto x_2$ 与
$x_2 \leadsto x_1$。须证明 $x_1 = x_2$。取概形 $U$ 及 \'etale 态射
$U \to X$，使得 $x_1,x_2$ 都属于 $|U| \to |X|$ 的像。由引理
\ref{decent-spaces-lemma-decent-specialization}，可在 $U$ 中找到特化
$u_1 \leadsto u_2$，映到 $x_1 \leadsto x_2$。再由引理
\ref{decent-spaces-lemma-decent-specialization}，可找到映到 $x_2 \leadsto x_1$ 的
$u_2' \leadsto u_1$。这意味着 $u_2' \leadsto u_2$ 是 $U$ 中映到 $X$ 同一点
$x_2$ 的两个点之间的特化。除非 $u_2' = u_2$，这是不可能的；参见引理
\ref{decent-spaces-lemma-decent-no-specializations-map-to-same-point}。因此也有
$u_1 = u_2$，正是所需结论。
\end{proof}

\begin{proposition}
\label{decent-spaces-proposition-reasonable-sober}
设 $S$ 为概形，$X$ 为 $S$ 上的良好代数空间。则拓扑空间 $|X|$ 是清醒的
（参见《拓扑》定义 \ref{topology-definition-generic-point}）。
\end{proposition}

\begin{proof}
引理 \ref{decent-spaces-lemma-kolmogorov} 已表明 $|X|$ 是 Kolmogorov 空间。因此只需证明
每个不可约闭子集 $T \subset |X|$ 都有一般点。由《空间的性质》引理
\ref{spaces-properties-lemma-reduced-closed-subspace}，存在闭子空间
$Z \subset X$，满足
% P08-E312: frozen source writes |Z| = |T| although T is already a subset of |X|; formula preserved; see ERRATA_CANDIDATES.jsonl.
$|Z| = |T|$。按定义，这意味着 $Z \to X$ 是代数空间的可表态射。因此由引理
\ref{decent-spaces-lemma-representable-properties}，$Z$ 是良好代数空间。由定理
\ref{decent-spaces-theorem-decent-open-dense-scheme}，存在作为概形的稠密开子空间
$Z' \subset Z$。这意味着 $|Z'| \subset T$ 开且稠密。因此拓扑空间
$|Z'|$ 不可约，也就是说 $Z'$ 是不可约概形。由《概形》引理
\ref{schemes-lemma-scheme-sober}，$|Z'|$ 是某个点 $\eta \in T$ 的闭包，
因而也有 $T = \overline{\{\eta\}}$，证毕。
\end{proof}

\noindent
对良好代数空间，维数具有预期的性质。

\begin{lemma}
\label{decent-spaces-lemma-dimension-decent-space}
设 $S$ 为概形。按《空间的性质》第
\ref{spaces-properties-section-dimension} 节定义的维数在 $S$ 上的良好
代数空间 $X$ 上表现良好：
\begin{enumerate}
\item 若 $x \in |X|$，则 $\dim_x(|X|) = \dim_x(X)$；
\item $\dim(|X|) = \dim(X)$.
\end{enumerate}
\end{lemma}

\begin{proof}
证明 (1)。取带有点 $u \in U$ 的概形 $U$ 及将 $u$ 映到 $x$ 的 \'etale
态射 $h : U \to X$。按定义，$X$ 在 $x$ 处的维数是 $\dim_u(|U|)$。
因此可以选取 $U$，使得 $\dim_x(X) = \dim(|U|)$。设 $d$ 为整数。若
$\dim(U) \geq d$，则 $U$ 中存在非平凡特化序列
$u_d \leadsto \ldots \leadsto u_0$。取像后得到相应序列
$h(u_d) \leadsto \ldots \leadsto h(u_0)$
；由引理 \ref{decent-spaces-lemma-decent-no-specializations-map-to-same-point}，其中每个
特化都非平凡。因此，$|U|$ 在 $|X|$ 中的像的维数至少为 $d$。反过来，
假设 $x_d \leadsto \ldots \leadsto x_0$ 是 $|X|$ 中的特化序列，且 $x_0$
属于 $|U| \to |X|$ 的像。由引理 \ref{decent-spaces-lemma-decent-specialization}，可将其
提升为 $U$ 中的特化序列。

\medskip\noindent
(2) 立即由 (1)、《拓扑》引理
\ref{topology-lemma-dimension-supremum-local-dimensions} 及《空间的性质》
第 \ref{spaces-properties-section-dimension} 节得出。
\end{proof}

\begin{lemma}
\label{decent-spaces-lemma-dimension-local-ring-quasi-finite}
设 $S$ 为概形，$X \to Y$ 为 $S$ 上代数空间的局部拟有限态射。设
$x \in |X|$，其像为 $y \in |Y|$。则 $Y$ 在 $y$ 处的局部环之维数
$\geq$ $X$ 在 $x$ 处的局部环之维数。
% P08-E313: frozen source says “is >= to the dimension”; intended grammar translated while the comparison symbol is preserved; see ERRATA_CANDIDATES.jsonl.
\end{lemma}

\begin{proof}
代数空间上一点处局部环维数的定义见《空间的性质》定义
\ref{spaces-properties-definition-dimension-local-ring}。取 \'etale 态射
$(V, v) \to (Y, y)$，其中 $V$ 为概形。再取 \'etale 态射
$U \to V \times_Y X$ 及点 $u \in U$，它映到 $x \in |X|$ 与 $v \in V$。
于是 $U \to V$ 局部拟有限，而须证明
$$
\dim(\mathcal{O}_{V, v}) \geq \dim(\mathcal{O}_{U, u})
$$
。这正是《代数》引理 \ref{algebra-lemma-dimension-inequality-quasi-finite}。
\end{proof}

\begin{lemma}
\label{decent-spaces-lemma-dimension-quasi-finite}
设 $S$ 为概形，$X \to Y$ 为 $S$ 上代数空间的局部拟有限态射。则
$\dim(X) \leq \dim(Y)$。
\end{lemma}

\begin{proof}
这由引理 \ref{decent-spaces-lemma-dimension-local-ring-quasi-finite} 与《空间的性质》引理
\ref{spaces-properties-lemma-dimension} 得出。
\end{proof}

\noindent
下述引理比《空间的性质》引理
\ref{spaces-properties-lemma-point-like-spaces} 略强一点。我们将在引理
\ref{decent-spaces-lemma-when-field} 中加强它。

\begin{lemma}
\label{decent-spaces-lemma-decent-point-like-spaces}
设 $S$ 为概形，$k$ 为域，$X$ 为 $S$ 上的代数空间，并假设存在满
\'etale 态射 $\Spec(k) \to X$。若 $X$ 良好，则
$X \cong \Spec(k')$，其中 $k/k'$ 为有限可分扩张。
\end{lemma}

\begin{proof}
该假设蕴含 $|X| = \{x\}$ 是单点集。由于 $X$ 良好，可以找到像为 $x$ 的
拟紧单态射 $\Spec(k') \to X$。于是投影
$U = \Spec(k') \times_X \Spec(k) \to \Spec(k)$
是单态射，故 $U = \Spec(k)$；参见《概形》引理
\ref{schemes-lemma-mono-towards-spec-field}。因此投影
$\Spec(k) = U \to \Spec(k')$ 是 \'etale 的，证毕。
\end{proof}






\section{约化单点空间}
\label{decent-spaces-section-singleton}

\noindent
{\it 单点}空间是满足 $|X|$ 为单点集的代数空间 $X$。事实表明，它们可能比
仅仅作为域的谱更有趣；参见《空间》例 \ref{spaces-example-Qbar}。我们建立
少量工具，以便讨论这类空间。

\begin{lemma}
\label{decent-spaces-lemma-flat-cover-by-field}
设 $S$ 为概形，$Z$ 为 $S$ 上的代数空间。设 $k$ 为域，且
$\Spec(k) \to Z$ 满且平坦。则对任意域 $k'$，任意态射
$\Spec(k') \to Z$ 均满且平坦。
\end{lemma}

\begin{proof}
考虑纤维方块
$$
\xymatrix{
T \ar[d] \ar[r] & \Spec(k) \ar[d] \\
\Spec(k') \ar[r] & Z
}
$$
注意，$T \to \Spec(k')$ 平坦且满，故 $T$ 非空。另一方面，由于 $k$ 为域，
$T \to \Spec(k)$ 平坦。因此 $T \to Z$ 平坦且满。由《空间的态射》引理
\ref{spaces-morphisms-lemma-flat-permanence}，$\Spec(k') \to Z$ 平坦。
它也满，因为按假设 $|Z|$ 是单点集。
\end{proof}

\begin{lemma}
\label{decent-spaces-lemma-unique-point}
设 $S$ 为概形，$Z$ 为 $S$ 上的代数空间。下列条件等价：
\begin{enumerate}
\item $Z$ 约化，且 $|Z|$ 是单点集；
\item 存在满平坦态射 $\Spec(k) \to Z$，其中 $k$ 为域；
\item 存在局部有限型、满且平坦的态射 $\Spec(k) \to Z$，其中 $k$ 为域。
\end{enumerate}
\end{lemma}

\begin{proof}
假设 (1)。设 $W$ 为概形，且 $W \to Z$ 为满 \'etale 态射。则 $W$ 是
约化概形。设 $\eta \in W$ 为 $W$ 某个不可约分支的一般点。由于 $W$ 约化，
有 $\mathcal{O}_{W, \eta} = \kappa(\eta)$。因此典范态射
$\eta = \Spec(\kappa(\eta)) \to W$ 平坦。可知复合态射 $\eta \to Z$ 平坦
（参见《空间的态射》引理
\ref{spaces-morphisms-lemma-composition-flat}）。由于 $|Z|$ 是单点集，它也
满。换言之，(2) 成立。

\medskip\noindent
假设 (2)。设 $W$ 为概形，且 $W \to Z$ 为满 \'etale 态射。取域 $k$ 及
满平坦态射 $\Spec(k) \to Z$。则 $W \times_Z \Spec(k)$ 是 $k$ 上的
\'etale 概形。因此 $W \times_Z \Spec(k)$ 是若干域的谱之不交并（参见注
\ref{decent-spaces-remark-recall}），特别地它约化。由于
$W \times_Z \Spec(k) \to W$ 满且平坦，由《下降》引理
\ref{descent-lemma-descend-reduced} 可知 $W$ 约化。换言之，(1) 成立。

\medskip\noindent
显然 (3) 蕴含 (2)。最后假设 (2)。取非空仿射概形 $W$ 及 \'etale 态射
$W \to Z$。取闭点 $w \in W$，并令 $k = \kappa(w)$。复合态射
$$
\Spec(k) \xrightarrow{w} W \longrightarrow Z
$$
由《空间的态射》引理
\ref{spaces-morphisms-lemma-composition-finite-type} 和
\ref{spaces-morphisms-lemma-etale-locally-finite-type} 局部有限型；由引理
\ref{decent-spaces-lemma-flat-cover-by-field}，它也平坦且满。因此 (3) 成立。
\end{proof}

\noindent
下述引理挑出一类比前一引理稍好的单点代数空间。

\begin{lemma}
\label{decent-spaces-lemma-unique-point-better}
设 $S$ 为概形，$Z$ 为 $S$ 上的代数空间。下列条件等价：
\begin{enumerate}
\item $Z$ 约化且局部 Noether，并且 $|Z|$ 是单点集；
\item 存在局部有限表示、满且平坦的态射 $\Spec(k) \to Z$，其中 $k$ 为域。
\end{enumerate}
\end{lemma}

\begin{proof}
假设 (2) 成立。由引理 \ref{decent-spaces-lemma-unique-point}，$Z$ 约化且 $|Z|$ 是单点集。
设 $W$ 为概形，且 $W \to Z$ 为满 \'etale 态射。取域 $k$ 及局部有限表示、
满且平坦的态射 $\Spec(k) \to Z$。则 $W \times_Z \Spec(k)$ 是 $k$ 上的
\'etale 概形，因而是若干域的谱之不交并（参见注
\ref{decent-spaces-remark-recall}），所以局部 Noether。由于
$W \times_Z \Spec(k) \to W$ 平坦、满且局部有限表示，可知
$\{W \times_Z \Spec(k) \to W\}$ 是 fppf 覆盖，进而由《下降》引理
\ref{descent-lemma-Noetherian-local-fppf} 得知 $W$ 局部 Noether。换言之，
(1) 成立。

\medskip\noindent
假设 (1)。取非空仿射概形 $W$ 及 \'etale 态射 $W \to Z$。取闭点
$w \in W$，并令 $k = \kappa(w)$。由于 $W$ 局部 Noether，态射
$w : \Spec(k) \to W$ 有限表示；参见《态射》引理
\ref{morphisms-lemma-closed-immersion-finite-presentation}。因此复合态射
$$
\Spec(k) \xrightarrow{w} W \longrightarrow Z
$$
由《空间的态射》引理
\ref{spaces-morphisms-lemma-composition-finite-presentation} 和
\ref{spaces-morphisms-lemma-etale-locally-finite-presentation} 局部有限
表示；由引理 \ref{decent-spaces-lemma-flat-cover-by-field}，它也平坦且满。因此 (2)
成立。
\end{proof}

\begin{lemma}
\label{decent-spaces-lemma-monomorphism-into-point}
设 $S$ 为概形，$Z' \to Z$ 为 $S$ 上代数空间的单态射。假设存在域 $k$ 及
局部有限表示、满且平坦的态射 $\Spec(k) \to Z$。则 $Z'$ 为空，或
$Z' = Z$。
\end{lemma}

\begin{proof}
可假设 $Z'$ 非空。此时纤维积 $T = Z' \times_Z \Spec(k)$ 非空；参见
《空间的性质》引理 \ref{spaces-properties-lemma-points-cartesian}。现在
$T$ 是代数空间，且投影 $T \to \Spec(k)$ 是单态射。因此
$T = \Spec(k)$；参见《空间的态射》引理
\ref{spaces-morphisms-lemma-monomorphism-toward-field}。于是
$\Spec(k) \to Z$ 经由 $Z'$ 分解。但由于 $\Spec(k) \to Z$ 满、平坦且局部
有限表示，$\Spec(k) \to Z$ 作为 $(\Sch/S)_{fppf}$ 上层的映射也满（参见《空间》注
\ref{spaces-remark-warning}），故 $Z' = Z$。
\end{proof}

\noindent
下述引理表明，可以为代数空间的每一点配上一个典范的、约化且局部 Noether
的单点代数空间。

\begin{lemma}
\label{decent-spaces-lemma-find-singleton-from-point}
设 $S$ 为概形，$X$ 为 $S$ 上的代数空间，并设 $x \in |X|$。则存在唯一的
$S$ 上代数空间的单态射 $Z \to X$，使得代数空间 $Z$ 满足引理
\ref{decent-spaces-lemma-unique-point-better} 中的等价条件，并且 $|Z| \to |X|$ 的像为
$\{x\}$。
\end{lemma}

\begin{proof}
取概形 $U$ 及满 \'etale 态射 $U \to X$。令 $R = U \times_X U$，从而
$X = U/R$ 是一个表现（参见《空间》第
\ref{spaces-section-presentations} 节）。令
$$
U' = \coprod\nolimits_{u \in U\text{ 且位于 }x} \Spec(\kappa(u)).
$$
典范态射 $U' \to U$ 是单态射。令
$$
R' = U' \times_X U' = R \times_{(U \times_S U)} (U' \times_S U').
$$
由于 $U' \to U$ 是单态射，可知投影 $s', t' : R' \to U'$ 都分解为一个
单态射再接一个 \'etale 态射。又因 $U'$ 是若干域的谱之不交并，使用注
\ref{decent-spaces-remark-recall} 和《概形》引理
\ref{schemes-lemma-mono-towards-spec-field} 可知 $R'$ 是若干域的谱之不交并，
并且态射 $s', t' : R' \to U'$ 是 \'etale 的。因此，由《空间》定理
\ref{spaces-theorem-presentation}，$Z = U'/R'$ 是代数空间。由于 $R'$ 是
$R$ 沿 $U' \to U$ 的限制，由《群胚》引理
\ref{groupoids-lemma-quotient-groupoid-restrict} 可知 $Z \to X$ 是单态射。
由于 $Z \to X$ 是单态射，由《空间的态射》引理
\ref{spaces-morphisms-lemma-monomorphism-injective-points} 可知
$|Z| \to |X|$ 单。由《空间的性质》引理
\ref{spaces-properties-lemma-points-cartesian} 可知
$$
|U'| = |Z \times_X U'| \to |Z| \times_{|X|} |U'|
$$
满。这蕴含（由 $U'$ 的选取）$|Z| \to |X|$ 的像为 $\{x\}$。因此 $|Z|$
是单点集。最后，按构造 $U'$ 局部 Noether 且约化，也就是说，$Z$ 满足引理
\ref{decent-spaces-lemma-unique-point-better} 中的等价条件。

\medskip\noindent
下面证明 $Z \to X$ 的唯一性。假设 $Z' \to X$ 是另一个这样的代数空间单态射。
则投影
$$
Z' \longleftarrow Z' \times_X Z \longrightarrow Z
$$
都是单态射。由《空间的性质》引理
\ref{spaces-properties-lemma-points-cartesian}，中间的代数空间非空。因此由
引理 \ref{decent-spaces-lemma-monomorphism-into-point}，两个投影都是同构，证毕。
\end{proof}

\noindent
我们引入下述术语，它预示了稍后将引入的剩余胚；参见《叠的性质》定义
\ref{stacks-properties-definition-residual-gerbe}。

\begin{definition}
\label{decent-spaces-definition-residual-space}
设 $S$ 为概形，$X$ 为 $S$ 上的代数空间，并设 $x \in |X|$。{\it $X$ 在
$x$ 处的剩余空间}\footnote{这是非标准记号。}是引理
\ref{decent-spaces-lemma-find-singleton-from-point} 中构造的单态射 $Z_x \to X$。
\end{definition}

\noindent
特别地，我们知道 $Z_x$ 是局部 Noether、约化的单点代数空间，并且存在域及
满、平坦、局部有限表示的态射
$$
\Spec(k) \longrightarrow Z_x.
$$
剩余空间常由从域的谱出发的单态射给出。

\begin{lemma}
\label{decent-spaces-lemma-residual-space-monomorphism}
设 $S$ 为概形，$X$ 为 $S$ 上的代数空间，并设 $x \in |X|$。$X$ 在 $x$
处的剩余空间 $Z_x$ 同构于某个域的谱，当且仅当 $x$ 可由单态射
$\Spec(k) \to X$ 表示，其中 $k$ 为域。若 $X$ 良好，则这对所有
$x \in |X|$ 都成立。
\end{lemma}

\begin{proof}
由于 $Z_x \to X$ 是单态射，若对某个域 $k$ 有 $Z_x = \Spec(k)$，则 $x$
由单态射 $\Spec(k) = Z_x \to X$ 表示。反过来，若表示 $x$ 的
$\Spec(k) \to X$ 是单态射，则 $Z_x \times_X \Spec(k) \to \Spec(k)$ 是
源非空的单态射；非空性由《空间的性质》引理
\ref{spaces-properties-lemma-points-cartesian} 得出。因此，由《空间的态射》
引理 \ref{spaces-morphisms-lemma-monomorphism-toward-field}，
$Z_x \times_X \Spec(k) = \Spec(k)$。于是得到单态射
$\Spec(k) \to Z_x$。由引理 \ref{decent-spaces-lemma-monomorphism-into-point}，它是同构。
最后一个断言由引理 \ref{decent-spaces-lemma-decent-points-monomorphism} 得出。
\end{proof}

\noindent
由下述引理，剩余空间是正则代数空间。

\begin{lemma}
\label{decent-spaces-lemma-residual-space-regular}
约化且局部 Noether 的单点代数空间 $Z$ 正则。
\end{lemma}

\begin{proof}
设 $Z$ 为概形 $S$ 上约化且局部 Noether 的单点代数空间。设 $W \to Z$ 为满
\'etale 态射，其中 $W$ 为概形。设 $k$ 为域，并设 $\Spec(k) \to Z$ 满、
平坦且局部有限表示（参见引理 \ref{decent-spaces-lemma-unique-point-better}）。概形
$T = W \times_Z \Spec(k)$ 在 $k$ 上 \'etale，特别地正则；参见注
\ref{decent-spaces-remark-recall}。由于 $T \to W$ 局部有限表示、平坦且满，由《下降》引理
\ref{descent-lemma-descend-regular} 可知 $W$ 正则。按定义，这意味着 $Z$
正则。
\end{proof}

\begin{lemma}
\label{decent-spaces-lemma-factor-through-residual-space-Noetherian}
设 $S$ 为概形，$f : Y \to X$ 为 $S$ 上代数空间的态射，并设
$x \in |X|$ 为一点。假设
\begin{enumerate}
\item $|f|(|Y|)$ 包含在 $\{x\} \subset |X|$ 中；
\item $Y$ 约化；
\item $X$ 局部 Noether。
\end{enumerate}
则 $f$ 经由 $X$ 在 $x$ 处的剩余空间 $Z_x$ 分解。
\end{lemma}

\begin{proof}
先作一则说明：由于 $Z_x \to X$ 是单态射，只需找到满 \'etale 态射
$Y' \to Y$，使得 $Y' \to X$ 经由 $Z_x$ 分解。还要注意，$Y'$ 也约化。

\medskip\noindent
设 $U$ 为仿射概形，$U \to X$ 为 \'etale 态射，使得 $x$ 属于
$|U| \to |X|$ 的像。由于 $X$ 局部 Noether，$U$ 是 Noether 仿射概形。
由假设 (1)，$Y' = U \times_X Y \to Y$ 满且 \'etale。令
$E \subset |U|$ 为映到 $x$ 的点集。
% P08-E314: frozen source omits “by” in “Denote E ... the set”; intended grammar translated; see ERRATA_CANDIDATES.jsonl.
$E$ 的元素之间不存在非平凡特化；参见引理
\ref{decent-spaces-lemma-no-specializations-map-to-same-point-Noetherian}。态射
$Y' \to U$ 将 $|Y'|$ 映入 $E$。由引理
\ref{decent-spaces-lemma-find-singleton-from-point} 的证明中对 $Z_x$ 的构造，知道
$\coprod_{u \in E} u \to X$ 经由 $Z_x$ 分解。因此只需证明 $Y' \to U$ 经由
$\coprod_{u \in E} u \to X$ 分解。用概形给出的一个 \'etale 覆盖替换 $Y'$
后（先前的说明允许这样做），这由《态射》引理
\ref{morphisms-lemma-factor-through-set-of-points} 得出。
% P08-E315: frozen source says “which we are allowed” without the infinitive “to do”; intended grammar translated; see ERRATA_CANDIDATES.jsonl.
\end{proof}

\begin{lemma}
\label{decent-spaces-lemma-factor-through-residual-space-decent}
设 $S$ 为概形，$f : Y \to X$ 为 $S$ 上代数空间的态射，并设
$x \in |X|$ 为一点。假设
\begin{enumerate}
\item $|f|(|Y|)$ 包含在 $\{x\} \subset |X|$ 中；
\item $Y$ 约化；
\item $x$ 可由拟紧单态射 $x : \Spec(k) \to X$ 表示，其中 $k$ 为域
（例如，当 $X$ 良好时）。
\end{enumerate}
则 $f$ 经由 $X$ 在 $x$ 处的剩余空间 $Z_x = \Spec(k)$ 分解。
\end{lemma}

\begin{proof}
由引理 \ref{decent-spaces-lemma-residual-space-monomorphism}，有 $Z_x = \Spec(k)$。

\medskip\noindent
先作一则说明：由于 $\Spec(k) \to X$ 是单态射，只需找到满 \'etale 态射
$Y' \to Y$，使得 $Y' \to X$ 经由 $Z_x$ 分解。还要注意，$Y'$ 也约化。

\medskip\noindent
用 $x$ 的拟紧开邻域替换 $X$ 后，可假设 $X$ 拟紧。
% P08-E316: frozen source omits “is” in “assume X quasi-compact”; intended grammar translated; see ERRATA_CANDIDATES.jsonl.
由引理 \ref{decent-spaces-lemma-filter-quasi-compact}，$x$ 是 $T \subset U \subset X$ 的点，
其中 $T \to U$（相应地，$U \to X$）是闭浸入（相应地，开浸入），且 $T$ 是
概形。由《空间的性质》引理
\ref{spaces-properties-lemma-factor-through-open-subspace}，$f$ 经由 $U$
分解，故可假设 $U = X$。又因为 $Y$ 约化，由《空间的性质》引理
\ref{spaces-properties-lemma-map-into-reduction}，$f$ 经由 $T$ 分解。因此可
假设 $X = T$ 是概形。由先前的说明，也可假设 $Y$ 是概形。这便归结为
《态射》引理 \ref{morphisms-lemma-factor-through-point}。
\end{proof}

\begin{example}
\label{decent-spaces-example-does-not-factor-through-residual-gerbe}
下面给出引理 \ref{decent-spaces-lemma-factor-through-residual-space-Noetherian} 与
\ref{decent-spaces-lemma-factor-through-residual-space-decent} 在 $X$ 既非局部 Noether
又非良好时的反例。设 $k$ 为域，$G$ 为无限预有限群。令 $Y$ 为把 $G$ 看作
零维仿射 $k$-群概形所得的概形，即
$Y = \Spec(\text{局部常值映射 } G \to k)$。令 $\Gamma$ 为把 $G$ 看作离散
$k$-群概形所得的概形，并令其以平移作用于 $Y$。置 $X = Y/\Gamma$。这是
单点代数空间，其投影为 $q : Y \to X$。设 $e \in G$ 为单位元（任意元素
均可），并将其视为 $Y$ 的 $k$-点。得到 $k$-点
$x : \Spec(k) \to X$；它是单态射，因为它是 $X \to \Spec(k)$ 的截面。
我们断言，尽管 $Y$ 仿射且约化，并且 $|X| = \{x\}$，态射 $q$ 不经由任何
态射 $\Spec(K) \to X$ 分解，其中 $K$ 为域。否则，由《空间的性质》引理
\ref{spaces-properties-lemma-equivalence-class-point-monomorphism}，它将经由
$x$ 分解。现在，$q$ 沿 $x$ 的拉回为 $\Gamma \to \Spec(k)$，其中投影
$\Gamma \to Y$ 是轨道映射 $g \mapsto g \cdot e$。后者没有截面，故断言
成立。
\end{example}

\begin{lemma}
\label{decent-spaces-lemma-residual-space-closed}
设 $S$ 为概形，$X$ 为 $S$ 上的代数空间。设 $x \in |X|$，其剩余空间为
$Z_x \subset X$。假设 $X$ 局部 Noether。则 $x$ 是 $|X|$ 的闭点，当且
仅当态射 $Z_x \to X$ 是闭浸入。
\end{lemma}

\begin{proof}
若 $Z_x \to X$ 是闭浸入，则 $x$ 是 $|X|$ 的闭点；参见《空间的态射》引理
\ref{spaces-morphisms-lemma-immersion-when-closed}。反过来，假设 $x$ 是
$|X|$ 的闭点。令 $Z \subset X$ 为满足 $|Z| = \{x\}$ 的约化闭子空间
（《空间的性质》引理
\ref{spaces-properties-lemma-reduced-closed-subspace}）。由《空间的态射》
引理 \ref{spaces-morphisms-lemma-immersion-locally-finite-type} 和
\ref{spaces-morphisms-lemma-locally-finite-type-locally-noetherian}，$Z$
局部 Noether。又因 $Z$ 约化且 $|Z| = \{x\}$，按定义，$Z = Z_x$ 是剩余
空间。
% P08-E317: frozen source has the malformed sequence “it Z = Z_x”; intended grammar translated; see ERRATA_CANDIDATES.jsonl.
\end{proof}











\section{良好空间}
\label{decent-spaces-section-decent}

\noindent
本节汇集良好空间的若干有用事实。

\begin{lemma}
\label{decent-spaces-lemma-locally-Noetherian-decent-quasi-separated}
任意局部 Noether 的良好代数空间都是拟分离的。
\end{lemma}

\begin{proof}
具体而言，设 $X$ 为良好且局部 Noether 的代数空间（在某个基概形上，例如
在 $\mathbf{Z}$ 上）。设 $U \to X$ 与 $V \to X$ 为 \'etale 态射，其中
$U$ 与 $V$ 为仿射概形。须证明 $W = U \times_X V$ 拟紧（《空间的性质》
引理 \ref{spaces-properties-lemma-characterize-quasi-separated}）。由于 $X$
局部 Noether，概形 $U$ 与 $V$ 都 Noether，且 $W$ 局部 Noether。由于 $X$
良好，态射 $W \to U$ 的纤维有限。事实上，可以用拟紧单态射
$\Spec(k) \to X$ 表示任意 $x \in |X|$。于是 $U_k$ 与 $V_k$ 都是 $k$
的有限可分扩张之谱的有限不交并（注 \ref{decent-spaces-remark-recall}），从而
$W_k = U_k \times_{\Spec(k)} V_k$ 有限。令 $n$ 为 $W \to U$ 在 $U$
各不可约分支一般点处纤维次数的最大值。考虑《态射补篇》引理
\ref{more-morphisms-lemma-stratify-flat-fp-lqf} 对 $W \to U$ 给出的分层
$$
U = U_0 \supset U_1 \supset U_2 \supset \ldots
$$
。由上述 $n$ 的选取，可知 $U_{n + 1}$ 为空。因此 $W \to U$ 的纤维普遍
有界。于是可应用《态射补篇》引理
\ref{more-morphisms-lemma-stratify-flat-fp-lqf-universally-bounded}，得到由
闭子集组成的分层
$$
\emptyset = Z_{-1} \subset Z_0 \subset Z_1 \subset Z_2 \subset
\ldots \subset Z_n = U
$$
，使得令 $S_r = Z_r \setminus Z_{r - 1}$ 后，态射
$W \times_U S_r \to S_r$ 有限局部自由。由于 $U$ Noether，概形 $S_r$
都 Noether，进而概形 $W \times_U S_r$ 都 Noether，于是
% P08-E318: frozen source identifies W with a scheme coproduct over locally closed strata; the displayed formula is preserved, while the argument only needs their finite union to be quasi-compact; see ERRATA_CANDIDATES.jsonl.
$W = \coprod W \times_U S_r$ 如所需为拟紧的。
\end{proof}

\begin{lemma}
\label{decent-spaces-lemma-when-field}
设 $S$ 为概形，$X$ 为 $S$ 上的良好代数空间。
\begin{enumerate}
\item 若 $|X|$ 是单点集，则 $X$ 是概形。
\item 若 $|X|$ 是单点集且 $X$ 约化，则对某个域 $k$ 有
$X \cong \Spec(k)$。
\end{enumerate}
\end{lemma}

\begin{proof}
假设 $|X|$ 是单点集。由定理
\ref{decent-spaces-theorem-decent-open-dense-scheme} 立即可知 $X$ 是概形，但也可如下
直接论证。取仿射概形 $U$ 及满 \'etale 态射 $U \to X$。令
$R = U \times_X U$。由引理 \ref{decent-spaces-lemma-UR-finite-above-x}（以及良好空间的
定义），$U$ 与 $R$ 都只有有限多个点。由引理
\ref{decent-spaces-lemma-decent-no-specializations-map-to-same-point}，这些点在 $U$ 与 $R$
中全都闭。因此 $U$ 与 $R$ 是仿射概形。可将 $U$ 缩小为单点空间。于是 $U$
是某个 Hensel 局部环的谱；参见《代数》引理
\ref{algebra-lemma-local-dimension-zero-henselian}。投影 $R \to U$ 是
\'etale 的；又因为 $U$ 是 $0$ 维 Hensel 局部环的谱，所以它们有限
\'etale；参见《代数》引理 \ref{algebra-lemma-characterize-henselian}。
由《群胚》命题 \ref{groupoids-proposition-finite-flat-equivalence}，$X$
是概形。

\medskip\noindent
(2) 由 (1) 以及约化单点概形是域的谱这一事实得出。
\end{proof}

\begin{remark}
\label{decent-spaces-remark-one-point-decent-scheme}
我们将在《空间的极限》引理
\ref{spaces-limits-lemma-reduction-scheme} 中看到：若一个代数空间的约化
是概形，则该代数空间本身也是概形。
\end{remark}

\begin{lemma}
\label{decent-spaces-lemma-algebraic-residue-field-extension-closed-point}
设 $S$ 为概形，$X$ 为 $S$ 上的良好代数空间。考虑交换图
$$
\xymatrix{
\Spec(k) \ar[rr] \ar[rd] & & X \ar[ld] \\
& S
}
$$
。假设 $\Spec(k) \to S$ 的像点 $s \in S$ 是闭点，且扩张
$\kappa(s) \subset k$ 是代数的。则 $\Spec(k) \to X$ 的像 $x$ 是
$|X|$ 的闭点。
\end{lemma}

\begin{proof}
假设对某个 $x' \in |X|$ 有 $x \leadsto x'$。取一个 \'etale 态射
$U \to X$，其中 $U$ 是概形，并取映到 $x'$ 的一点 $u' \in U'$。
% P08-E319: frozen source places u' in undefined U' although the chosen scheme is U; formula preserved; see ERRATA_CANDIDATES.jsonl.
由引理 \ref{decent-spaces-lemma-decent-specialization}，可在 $U$ 中取特化
$u \leadsto u'$，使 $u$ 在 $X$ 中映到 $x$。于是 $u$ 是概形
$W = \Spec(k) \times_X U$ 的某一点 $w$ 的像。由于投影
$W \to \Spec(k)$ 是 \'etale 的，$\kappa(w) \supset k$ 是有限扩张，因而
$\kappa(w) \supset \kappa(s)$ 是代数扩张，继而
$\kappa(u) \supset \kappa(s)$ 是代数扩张。由《态射》引理
\ref{morphisms-lemma-algebraic-residue-field-extension-closed-point-fibre}，
$u$ 是 $U$ 的闭点。因此 $u = u'$，从而 $x = x'$。
\end{proof}

\begin{lemma}
\label{decent-spaces-lemma-finite-residue-field-extension-finite}
设 $S$ 为概形，$X$ 为 $S$ 上的良好代数空间。考虑交换图
$$
\xymatrix{
\Spec(k) \ar[rr] \ar[rd] & & X \ar[ld] \\
& S
}
$$
。假设 $\Spec(k) \to S$ 的像点 $s \in S$ 是闭点，且域扩张
$k/\kappa(s)$ 有限。则 $\Spec(k) \to X$ 是有限态射。若
$\kappa(s) = k$，则 $\Spec(k) \to X$ 是闭浸入。
\end{lemma}

\begin{proof}
由引理 \ref{decent-spaces-lemma-algebraic-residue-field-extension-closed-point}，像点
$x \in |X|$ 是闭点。令 $Z \subset X$ 为满足 $|Z| = \{x\}$ 的约化闭子空间
（《空间的性质》引理
\ref{spaces-properties-lemma-reduced-closed-subspace}）。由引理
\ref{decent-spaces-lemma-representable-named-properties}，$Z$ 是良好代数空间。再由引理
\ref{decent-spaces-lemma-when-field}，对某个域 $k'$ 有 $Z = \Spec(k')$。当然
$k \supset k' \supset \kappa(s)$。于是 $\Spec(k) \to Z$ 是概形的有限态射；
$Z \to X$ 是闭浸入，因而也是有限态射。因此 $\Spec(k) \to X$ 有限
（《空间的态射》引理
\ref{spaces-morphisms-lemma-composition-integral}）。若 $k = \kappa(s)$，则
$\Spec(k) = Z$，且 $\Spec(k) \to X$ 是闭浸入。
\end{proof}

\begin{lemma}
\label{decent-spaces-lemma-decent-space-closed-point}
设 $S$ 为概形，$X$ 为 $S$ 上的良好代数空间。设 $x \in |X|$ 为闭点。
则可由某个域的谱给出的闭浸入 $i : \Spec(k) \to X$ 表示 $x$。
\end{lemma}

\begin{proof}
由定义 \ref{decent-spaces-definition-very-reasonable}，可用拟紧单态射
$i : \Spec(k) \to X$ 表示 $x$，其中 $k$ 是域。取 \'etale 态射
$U \to X$，其中 $U$ 是仿射概形。由于 $x$ 闭且 $X$ 良好，$|U| \to |X|$
在 $x$ 上的纤维 $F$ 由闭点组成（引理
\ref{decent-spaces-lemma-decent-no-specializations-map-to-same-point}）。由于 $i$ 是单态射，
$U_k = U \times_X \Spec(k) \to U$ 也是单态射。特别地，映射
$|U_k| \to F$ 是单射。又因 $U_k$ 拟紧且在域上 \'etale，注
\ref{decent-spaces-remark-recall} 表明 $U_k$ 是若干域的谱的有限不交并。记
$U_k = \Spec(k_1) \amalg \ldots \amalg \Spec(k_r)$。由于
$\Spec(k_i) \to U$ 是单态射，其像 $u_i$ 的剩余域为
$\kappa(u_i) = k_i$。又因 $u_i \in F$ 是闭点，态射
$\Spec(k_i) \to U$ 是闭浸入。各 $u_i$ 两两不同，故 $U_k \to U$ 是闭浸入。
因此由《空间的态射》引理
\ref{spaces-morphisms-lemma-closed-immersion-local}，$i$ 是闭浸入。证毕。
\end{proof}





\section{局部分离空间}
\label{decent-spaces-section-locally-separated}

\noindent
事实证明，局部分离的代数空间是良好空间。

\begin{lemma}
\label{decent-spaces-lemma-infinite-number}
设 $A$ 为环，$k$ 为域。设 $\mathfrak p_n$（$n \geq 1$）为 $A$ 的一列
两两不同的素理想，并且对每个 $n$ 给定嵌入
$k \to \kappa(\mathfrak p_n)$。则下列态射之像的闭包
$$
\coprod\nolimits_{n \not = m}
\Spec(\kappa(\mathfrak p_n) \otimes_k \kappa(\mathfrak p_m))
\longrightarrow
\Spec(A \otimes A)
$$
与对角相交。
\end{lemma}

\begin{proof}
令 $k_n = \kappa(\mathfrak p_n)$。可设 $A = \prod k_n$。记
$x_n = \Spec(k_n)$ 为对应于 $A \to k_n$ 的开闭点。
% P08-E320: frozen source omits “by” in “Denote x_n ...”; intended naming construction translated; see ERRATA_CANDIDATES.jsonl.
于是 $\Spec(A) = Z \amalg \{x_n\}$，其中 $Z$ 是非空闭子集。具体地，
$Z = V(e_n; n \geq 1)$，这里 $e_n$ 是 $A$ 中对应于因子 $k_n$ 的幂等元；
由于各 $e_n$ 生成的理想不等于 $A$，$Z$ 非空。我们将证明像的闭包包含
$\Delta(Z)$。映射
$$
(\prod k_n) \otimes_k (\prod k_m)
\longrightarrow
\prod\nolimits_{n \not = m} k_n \otimes_k k_m
$$
的核是由 $e_n \otimes e_n$（$n \geq 1$）生成的理想。因此，谱上诱导映射
之像的闭包为 $V(e_n \otimes e_n; n \geq 1)$，它与
$\Delta(\Spec(A))$ 的交为 $\Delta(Z)$。故只须证明
$$
\coprod\nolimits_{n \not = m} \Spec(k_n \otimes_k k_m)
\longrightarrow
\Spec(\prod\nolimits_{n \not = m} k_n \otimes_k k_m)
$$
有稠密像。这是因为环同态族
$\prod_{n \not = m} k_n \otimes_k k_m \to k_n \otimes_k k_m$ 联合单射。
\end{proof}

\begin{lemma}[David Rydh]
\label{decent-spaces-lemma-locally-separated-decent}
局部分离的代数空间是良好的。
\end{lemma}

\begin{proof}
设 $S$ 为概形，$X$ 为 $S$ 上局部分离的代数空间。由《空间的性质》定义
\ref{spaces-properties-definition-separated}，可设
$S = \Spec(\mathbf{Z})$。未注明基底的纤维积均取在 $\mathbf{Z}$ 上。设
$x \in |X|$。取概形 $U$、\'etale 态射 $U \to X$ 以及映到 $|X|$ 中
$x$ 的一点 $u \in U$。一如通常，等同 $u = \Spec(\kappa(u))$。由于
$X$ 局部分离，态射
$$
u \times_X u \to u \times u
$$
是浸入（《空间的态射》引理
\ref{spaces-morphisms-lemma-fibre-product-after-map}）。因此《群胚补篇》引理
\ref{more-groupoids-lemma-locally-closed-image-is-closed} 表明它是闭浸入
（使用《概形》引理 \ref{schemes-lemma-immersion-when-closed}）。态射
$u \times_X u \to u \times_X U$ 是单态射（它是 $u \to U$ 的基变换），而
$u \times_X U \to u$ 是 \'etale 态射，故 $u \times_X u$ 是若干域的谱的
不交并（见注 \ref{decent-spaces-remark-recall} 及《概形》引理
\ref{schemes-lemma-mono-towards-spec-field}）。它在仿射概形 $u \times u$ 中
还是闭的，故 $u \times_X u$ 是若干域的谱的有限不交并。因此可用单态射
$\Spec(k) \to X$ 表示 $x$，其中 $k$ 是域；参见引理
\ref{decent-spaces-lemma-R-finite-above-x}。

\medskip\noindent
接下来，设 $U = \Spec(A)$ 为仿射概形，且 $U \to X$ 是 \'etale 态射。
为完成证明，只须证明 $F = U \times_X \Spec(k)$ 有限。将其写成 $k$
的有限可分扩张之谱的不交并
$F = \coprod_{i \in I} \Spec(k_i)$。须证明 $I$ 有限。令
$R = U \times_X U$。由于 $X$ 局部分离，态射
$j : R \to U \times U$ 是浸入。取开子概形 $U' \subset U \times U$，使
$j$ 分解经过闭浸入 $j' : R \to U'$。令 $e : U \to R$ 为对角映射。注意
$e$ 是两个在 $U$ 上 \'etale 的概形之间的态射，且
$\Delta = j \circ e$ 是闭浸入；由此可知，对某个开闭子概形
$W \subset R$ 有 $R = e(U) \amalg W$。由于 $j'$ 是闭浸入，
$j'(W) \subset U'$ 闭且与 $j'(e(U))$ 不交。因此在 $U \times U$ 中有
$\overline{j(W)} \cap \Delta(U) = \emptyset$。注意，对所有
$i \not = i'$、$i, i' \in I$，$W$ 都包含
$\Spec(k_i \otimes_k k_{i'})$。由引理 \ref{decent-spaces-lemma-infinite-number}，可知
$I$ 如所需为有限集。
\end{proof}








\section{赋值判据}
\label{decent-spaces-section-valuative-criterion-universally-closed}

\noindent
对于以良好空间为源的拟紧态射，赋值判据是该态射普遍闭的必要条件。


\begin{proposition}
\label{decent-spaces-proposition-characterize-universally-closed}
设 $S$ 为概形，$f : X \to Y$ 为 $S$ 上代数空间的态射。假设 $f$ 拟紧且
$X$ 良好。则 $f$ 普遍闭，当且仅当赋值判据的存在性部分成立。
\end{proposition}

\begin{proof}
《空间的态射》引理
\ref{spaces-morphisms-lemma-quasi-compact-existence-universally-closed}
已经给出一个方向。为证明另一个方向，假设 $f$ 普遍闭。设
$$
\xymatrix{
\Spec(K) \ar[r] \ar[d] & X \ar[d] \\
\Spec(A) \ar[r] & Y
}
$$
为《空间的态射》定义
\ref{spaces-morphisms-definition-valuative-criterion} 中的图。令
$X_A = \Spec(A) \times_Y X$，于是有
$$
\xymatrix{
\Spec(K) \ar[r] \ar[rd] & X_A \ar[d] \\
 & \Spec(A)
}
$$
。由《空间的态射》引理
\ref{spaces-morphisms-lemma-base-change-quasi-compact}，
$X_A \to \Spec(A)$ 拟紧。由于 $X_A \to X$ 可表，引理
\ref{decent-spaces-lemma-representable-properties} 表明 $X_A$ 也良好。此外，由于 $f$
普遍闭，$X_A \to \Spec(A)$ 普遍闭。因此可以而且确实将 $X$ 换成 $X_A$，
将 $Y$ 换成 $\Spec(A)$。

\medskip\noindent
设 $x' \in |X|$ 为 $\Spec(K) \to X$ 的等价类，设
$y \in |Y| = |\Spec(A)|$ 为闭点。令 $y' = f(x')$；它是
$\Spec(A)$ 的一般点。由于 $f$ 普遍闭，$f(\overline{\{x'\}})$ 包含
$\overline{\{y'\}}$，因而包含 $y$。取满足 $f(x) = y$ 的一点
$x \in \overline{\{x'\}}$。取概形 $U$ 以及 \'etale 态射
$\varphi : U \to X$，使得存在 $u \in U$ 满足 $\varphi(u) = x$。由引理
\ref{decent-spaces-lemma-specialization} 及 $X$ 良好的假设，在 $U$ 上存在特化
$u' \leadsto u$，满足 $\varphi(u') = x'$。这意味着存在公共域扩张
$K \subset K' \supset \kappa(u')$，使得
$$
\xymatrix{
\Spec(K') \ar[r] \ar[d] & U \ar[d] \\
\Spec(K) \ar[r] \ar[rd] & X \ar[d] \\
 & \Spec(A)
}
$$
交换。这给出如下交换的环图
$$
\xymatrix{
K' & \mathcal{O}_{U, u} \ar[l] \\
K \ar[u] & \\
 & A \ar[lu] \ar[uu]
}
$$
。由《代数》引理 \ref{algebra-lemma-dominate}，可找到赋值环
$A' \subset K'$，它支配 $\mathcal{O}_{U, u}$ 在 $K'$ 中的像。按构造，
$\mathcal{O}_{U, u}$ 支配 $A$，故 $A'$ 也支配 $A$。于是得到与
《空间的态射》定义
\ref{spaces-morphisms-definition-valuative-criterion} 中第二个图同型的图，
命题得证。
\end{proof}







\section{相对条件}
\label{decent-spaces-section-relative-conditions}

\noindent
这是（又）一个技术性章节，讨论代数空间上与点有关的条件。跳过本节或许是个
好主意。

\begin{definition}
\label{decent-spaces-definition-relative-conditions}
设 $S$ 为概形。若 $S$ 上的代数空间 $X$ 具有引理
\ref{decent-spaces-lemma-bounded-fibres} 中相应的性质，则称 $X$ {\it 具有性质 $(\beta)$}。
设 $f : X \to Y$ 为 $S$ 上代数空间的态射。
\begin{enumerate}
\item 若对任意概形 $T$ 及态射 $T \to Y$，纤维积 $T \times_Y X$ 都具有
性质 $(\beta)$，则称 $f$ {\it 具有性质 $(\beta)$}。
\item 若对任意概形 $T$ 及态射 $T \to Y$，纤维积 $T \times_Y X$ 都是
良好代数空间，则称 $f$ 是{\it 良好的}。
\item 若对任意概形 $T$ 及态射 $T \to Y$，纤维积 $T \times_Y X$ 都是
合理代数空间，则称 $f$ 是{\it 合理的}。
\item 若对任意概形 $T$ 及态射 $T \to Y$，纤维积 $T \times_Y X$ 都是
非常合理的代数空间，则称 $f$ 是{\it 非常合理的}。
\end{enumerate}
\end{definition}

\noindent
非正式讨论见注 \ref{decent-spaces-remark-very-reasonable}。事实将表明，非常合理态射这一
类并不那么有用，而良好态射与合理态射这两类是有用的。

\begin{lemma}
\label{decent-spaces-lemma-properties-trivial-implications}
设 $S$ 为概形，$f : X \to Y$ 为 $S$ 上代数空间的态射。关于 $f$ 的各条件
之间有下列蕴含关系：
$$
\xymatrix{
\text{可表} \ar@{=>}[rd] & & & & \\
& \text{非常合理} \ar@{=>}[r] & \text{合理} \ar@{=>}[r] &
\text{良好} \ar@{=>}[r] & (\beta) \\
\text{拟分离} \ar@{=>}[ru] & & & &
}
$$
\end{lemma}

\begin{proof}
这由各定义、引理 \ref{decent-spaces-lemma-bounded-fibres} 及《空间的态射》引理
\ref{spaces-morphisms-lemma-separated-local} 立即得出。
\end{proof}

\noindent
下面是另一个合理性检验。

\begin{lemma}
\label{decent-spaces-lemma-property-for-morphism-out-of-property}
设 $S$ 为概形，$f : X \to Y$ 为 $S$ 上代数空间的态射。若 $X$ 良好
（相应地，合理；具有引理 \ref{decent-spaces-lemma-bounded-fibres} 的性质 $(\beta)$），则
$f$ 良好（相应地，合理；具有性质 $(\beta)$）。
\end{lemma}

\begin{proof}
设 $T$ 为概形并给定态射 $T \to Y$。则 $T \to Y$ 可表，故其基变换
$T \times_Y X \to X$ 可表。因此若 $X$ 良好（或合理），则
$T \times_Y X$ 也如此；见引理
\ref{decent-spaces-lemma-representable-named-properties}。性质 $(\beta)$ 的情形类似；见引理
\ref{decent-spaces-lemma-representable-properties}。
\end{proof}

\begin{lemma}
\label{decent-spaces-lemma-base-change-relative-conditions}
具有性质 $(\beta)$、良好或合理这三种条件都在任意基变换下保持。
\end{lemma}

\begin{proof}
这由定义立即得出。
\end{proof}

\begin{lemma}
\label{decent-spaces-lemma-property-over-property}
设 $S$ 为概形，$f : X \to Y$ 为 $S$ 上代数空间的态射。设
$\omega \in \{\beta, decent, reasonable\}$。假设 $Y$ 具有性质 $(\omega)$，
且 $f : X \to Y$ 具有 $(\omega)$。则 $X$ 具有 $(\omega)$。
\end{lemma}

\begin{proof}
先证明 $\omega = \beta$ 的情形。此时须证明，任意 $x \in |X|$ 都可由从某个
域的谱到 $X$ 的单态射表示。令 $y = f(x) \in |Y|$。按假设，存在域 $k$ 及
表示 $y$ 的单态射 $\Spec(k) \to Y$。于是 $x$ 对应于
$\Spec(k) \times_Y X$ 的一点 $x'$。按假设，$x'$ 可由单态射
$\Spec(k') \to \Spec(k) \times_Y X$ 表示。显然，复合
$\Spec(k') \to X$ 是表示 $x$ 的单态射。

\medskip\noindent
再证明 $\omega = decent$ 的情形。设 $x \in |X|$ 且
$y = f(x) \in |Y|$。由前一段的结果，可取图
$$
\xymatrix{
\Spec(k') \ar[r]_x \ar[d] & X \ar[d]^f \\
\Spec(k) \ar[r]^y & Y
}
$$
，其中两个水平箭头都是单态射。
% P08-E321: frozen source omits “are” in “whose horizontal arrows monomorphisms”; intended grammar translated; see ERRATA_CANDIDATES.jsonl.
由于 $Y$ 良好，态射 $y$ 拟紧。由于 $f$ 良好，代数空间
$\Spec(k) \times_Y X$ 良好，故单态射
$\Spec(k') \to \Spec(k) \times_Y X$ 拟紧。于是单态射
$x : \Spec(k') \to X$ 是拟紧态射，因为它是拟紧态射的复合（使用
《空间的态射》引理
\ref{spaces-morphisms-lemma-base-change-quasi-compact} 和
\ref{spaces-morphisms-lemma-composition-quasi-compact}）。由于点 $x$ 任意，
这表明 $X$ 良好。

\medskip\noindent
最后证明 $\omega = reasonable$ 的情形。取 \'etale 态射 $V \to Y$，其中
$V$ 是仿射概形；再取 \'etale 态射 $U \to V \times_Y X$，其中 $U$ 是仿射
概形。按假设，$V \to Y$ 具有普遍有界纤维。由引理
\ref{decent-spaces-lemma-base-change-universally-bounded}，态射 $V \times_Y X \to X$
具有普遍有界纤维。由关于 $f$ 的假设，$U \to V \times_Y X$ 具有普遍
有界纤维。再由引理 \ref{decent-spaces-lemma-composition-universally-bounded}，复合
$U \to X$ 具有普遍有界纤维。因此存在足够多从概形出发、具有普遍有界纤维的
\'etale 态射 $U \to X$，从而 $X$ 合理。
% P08-E322: frozen source uses singular “there exists sufficiently many morphisms”; intended plural existence translated; see ERRATA_CANDIDATES.jsonl.
\end{proof}

\begin{lemma}
\label{decent-spaces-lemma-composition-relative-conditions}
具有性质 $(\beta)$、良好或合理这三种条件都在复合下保持。
\end{lemma}

\begin{proof}
设 $\omega \in \{\beta, decent, reasonable\}$。设 $f : X \to Y$ 与
$g : Y \to Z$ 为概形 $S$ 上代数空间的态射，并假设 $f$ 与 $g$ 都具有性质
$(\omega)$。须证明，对任意概形 $T$ 及态射 $T \to Z$，空间
$T \times_Z X$ 都具有 $(\omega)$。由引理
\ref{decent-spaces-lemma-base-change-relative-conditions}，这归结为以下断言：若 $Y$ 是
具有性质 $(\omega)$ 的代数空间，且 $f : X \to Y$ 是具有 $(\omega)$ 的态射，
则 $X$ 具有 $(\omega)$。这正是引理
\ref{decent-spaces-lemma-property-over-property} 的内容。
\end{proof}

\begin{lemma}
\label{decent-spaces-lemma-fibre-product-conditions}
设 $S$ 为概形，$f : X \to Y$、$g : Z \to Y$ 为 $S$ 上代数空间的态射。
若 $X$ 与 $Z$ 良好（相应地，合理；具有引理
\ref{decent-spaces-lemma-bounded-fibres} 的性质 $(\beta)$），则 $X \times_Y Z$ 也如此。
\end{lemma}

\begin{proof}
事实上，由引理 \ref{decent-spaces-lemma-property-for-morphism-out-of-property}，态射
$X \to Y$ 具有该性质。再由引理
\ref{decent-spaces-lemma-base-change-relative-conditions}，其基变换
$X \times_Y Z \to Z$ 具有该性质。最后，引理
\ref{decent-spaces-lemma-property-over-property} 表明 $X \times_Y Z$ 具有该性质。
\end{proof}

\begin{lemma}
\label{decent-spaces-lemma-descent-conditions}
设 $S$ 为概形，$f : X \to Y$ 为 $S$ 上代数空间的态射，并设
$\mathcal{P} \in \{(\beta), decent, reasonable\}$。假设
\begin{enumerate}
\item $f$ 拟紧；
\item $f$ 是 \'etale 态射；
\item $|f| : |X| \to |Y|$ 满；并且
\item 代数空间 $X$ 具有性质 $\mathcal{P}$。
\end{enumerate}
则 $Y$ 具有性质 $\mathcal{P}$。
\end{lemma}

\begin{proof}
先证明 $\mathcal{P} = (\beta)$ 的情形。设 $y \in |Y|$ 为一点。须证明 $y$
可由从域的谱出发的单态射表示。取一点 $x \in |X|$，使 $f(x) = y$。按假设，
可由单态射 $\Spec(k) \to X$ 表示 $x$，其中 $k$ 是域。由引理
\ref{decent-spaces-lemma-R-finite-above-x}，只须证明两个投影
$\Spec(k) \times_Y \Spec(k) \to \Spec(k)$ 都是拟紧 \'etale 态射。第一个
投影可分解为
$$
\Spec(k) \times_Y \Spec(k)
\longrightarrow
\Spec(k) \times_Y X
\longrightarrow
\Spec(k)
$$
。第一个态射是单态射，第二个态射拟紧且 \'etale。由《空间的性质》引理
\ref{spaces-properties-lemma-etale-over-field-scheme}，
$\Spec(k) \times_Y X$ 是概形。因此它是 $k$ 的有限可分域扩张之谱的有限
不交并。由《概形》引理 \ref{schemes-lemma-mono-towards-spec-field}，第一个
箭头将 $\Spec(k) \times_Y \Spec(k)$ 等同于 $k$ 的有限可分域扩张之谱的
有限不交并。因此该投影态射拟紧且 \'etale。

\medskip\noindent
再证明 $\mathcal{P} = decent$ 的情形。证明的第一段已经表明，每个
$y \in |Y|$ 都可由单态射 $y : \Spec(k) \to Y$ 表示。选取这样的 $y$。
取仿射概形 $U$ 及 \'etale 态射 $U \to X$，使 $|U| \to |Y|$ 的像包含
$y$。由引理 \ref{decent-spaces-lemma-UR-finite-above-x}，只须证明 $U_y$ 是 $k$ 上的
有限概形。纤维积 $X_y = \Spec(k) \times_Y X$ 是 $k$ 上拟紧的 \'etale
代数空间。因此由《空间的性质》引理
\ref{spaces-properties-lemma-etale-over-field-scheme}，它是概形，从而是
$k$ 的有限可分扩张之谱的有限不交并。记
$X_y = \{x_1, \ldots, x_n\}$，其中 $x_i$ 由
$x_i : \Spec(k_i) \to X$ 给出，且 $[k_i : k] < \infty$。按假设 $X$ 良好，
故概形 $U_{x_i} = \Spec(k_i) \times_X U$ 在 $k_i$ 上有限。最后注意，作为
概形有 $U_y = \coprod U_{x_i}$，故 $U_y$ 如所需在 $k$ 上有限。

\medskip\noindent
最后证明 $\mathcal{P} = reasonable$ 的情形。取仿射概形 $V$ 及 \'etale
态射 $V \to Y$。须证明 $V \to Y$ 的纤维普遍有界。
% P08-E323: frozen source has “We have the show”; intended infinitive translated; see ERRATA_CANDIDATES.jsonl.
代数空间 $V \times_Y X$ 拟紧。因此由《空间的性质》引理
\ref{spaces-properties-lemma-quasi-compact-affine-cover}，可找到仿射概形
$W$ 及满 \'etale 态射 $W \to V \times_Y X$。下图中实线部分为已给图：
$$
\xymatrix{
W \ar[r]  \ar[rd] &
V \times_Y X \ar[r] \ar[d] &
X \ar[d]_f & \Spec(k) \ar@{..>}[l]^x \ar@{..>}[ld]^y \\
 & V \ar[r] & Y
}
$$
由空间 $X$ 合理的假设，态射 $W \to X$ 具有普遍有界纤维。令 $n$ 为
$W \to X$ 的纤维次数的一个界。我们断言，同一整数也给出 $V \to Y$ 的
纤维之界。事实上，设 $y \in |Y|$ 为一点。则存在 $x \in |X|$，使
$f(x) = y$（见上文）。
% P08-E324: frozen source uses “a x” although the variable is conventionally read with an initial vowel sound; intended existential grammar translated; see ERRATA_CANDIDATES.jsonl.
这意味着可找到域 $k$ 以及由上图虚线箭头给出的态射 $x, y$。特别地，得到
满 \'etale 态射
$$
\Spec(k) \times_{x, X} W
\to
\Spec(k) \times_{x, X} (V \times_Y X) = \Spec(k) \times_{y, Y} V
$$
。这表明 $\Spec(k) \times_{y, Y} V$ 在 $k$ 上的次数不超过
$\Spec(k) \times_{x, X} W$ 在 $k$ 上的次数，即 $\leq n$，结论成立。
（论证的最后一部分与引理
\ref{decent-spaces-lemma-descent-universally-bounded} 证明中的论证相同。遗憾的是，该引理
不够一般，因为它只适用于可表态射。）
\end{proof}

\begin{lemma}
\label{decent-spaces-lemma-relative-conditions-local}
设 $S$ 为概形，$f : X \to Y$ 为 $S$ 上代数空间的态射，并设\linebreak
$\mathcal{P} \in \{(\beta), decent, reasonable, very\ reasonable\}$。下列条件
等价：
\begin{enumerate}
\item $f$ 具有 $\mathcal{P}$；
\item 对每个仿射概形 $Z$ 及每个态射 $Z \to Y$，$f$ 的基变换
$Z \times_Y X \to Z$ 具有 $\mathcal{P}$；
\item 对每个仿射概形 $Z$ 及每个态射 $Z \to Y$，代数空间
$Z \times_Y X$ 具有 $\mathcal{P}$；以及
\item 存在 Zariski 覆盖 $Y = \bigcup Y_i$，使每个态射
$f^{-1}(Y_i) \to Y_i$ 都具有 $\mathcal{P}$。
\end{enumerate}
若 $\mathcal{P} \in \{(\beta), decent, reasonable\}$，则上述条件还等价于
\begin{enumerate}
\item[(5)] 存在概形 $V$ 及满 \'etale 态射 $V \to Y$，使基变换
$V \times_Y X \to V$ 具有 $\mathcal{P}$。
\end{enumerate}
\end{lemma}

\begin{proof}
蕴含 (1) $\Rightarrow$ (2) $\Rightarrow$ (3) $\Rightarrow$ (4) 是平凡的。
蕴含 (3) $\Rightarrow$ (1) 可如下得到。设 $Z \to Y$ 为态射，其源是
$S$ 上的概形。考虑代数空间 $Z \times_Y X$。若假设 (3)，则对任意仿射开
子集 $W \subset Z$，$Z \times_Y X$ 的开子空间 $W \times_Y X$ 具有性质
$\mathcal{P}$。故由引理 \ref{decent-spaces-lemma-properties-local}，空间
$Z \times_Y X$ 具有性质 $\mathcal{P}$，即 (1) 成立。类似的论证（略去）
表明 (4) 蕴含 (1)。

\medskip\noindent
蕴含 (1) $\Rightarrow$ (5) 是平凡的。设 $V \to Y$ 为 (5) 中从概形出发的
\'etale 态射。设 $Z$ 为仿射概形，并给定态射 $Z \to Y$。考虑图
$$
\xymatrix{
Z \times_Y V \ar[r]_q \ar[d]_p & V \ar[d] \\
Z \ar[r] & Y
}
$$
。由于 $p$ 是 \'etale 态射，因而是开映射，可选取有限多个仿射开子概形
$W_i \subset Z \times_Y V$，使 $Z = \bigcup p(W_i)$。考虑交换图
$$
\xymatrix{
V \times_Y X \ar[d] &
(\coprod W_i) \times_Y X \ar[l] \ar[d] \ar[r] &
Z \times_Y X \ar[d] \\
V &
\coprod W_i \ar[l] \ar[r] &
Z
}
$$
已知 $V \times_Y X$ 具有性质 $\mathcal{P}$。由引理
\ref{decent-spaces-lemma-representable-properties}，
$(\coprod W_i) \times_Y X$ 具有性质 $\mathcal{P}$。注意，态射
$(\coprod W_i) \times_Y X \to Z \times_Y X$ 是
$\coprod W_i \to Z$ 的基变换，故拟紧且 \'etale。因此由引理
\ref{decent-spaces-lemma-descent-conditions}，$Z \times_Y X$ 具有性质 $\mathcal{P}$。
\end{proof}

\begin{remark}
\label{decent-spaces-remark-very-reasonable}
性质 $(\beta)$、良好、合理及非常合理的非正式描述已在第
\ref{decent-spaces-section-reasonable-decent} 节给出。非常粗略地说，一个态射具有其中
某种性质，是指该态射的纤维具有相应性质。良好性适用于证明关于 $|X|$ 上各点
之特化的结论。合理性稍强一些，并且在技术上很容易使用。
\end{remark}

\noindent
下面给出先前承诺的一个使用良好态射的引理。

\begin{lemma}
\label{decent-spaces-lemma-re-characterize-universally-closed}
设 $S$ 为概形，$f : X \to Y$ 为 $S$ 上代数空间的态射。假设 $f$ 拟紧且
良好。（例如，当 $f$ 可表或拟分离时即如此；见引理
\ref{decent-spaces-lemma-properties-trivial-implications}。）则 $f$ 普遍闭，当且仅当
赋值判据的存在性部分成立。
\end{lemma}

\begin{proof}
《空间的态射》引理
\ref{spaces-morphisms-lemma-quasi-compact-existence-universally-closed}
已经证明，任意满足赋值判据存在性部分的拟紧态射都是普遍闭的。为证明另一
方向，假设 $f$ 普遍闭。命题
\ref{decent-spaces-proposition-characterize-universally-closed} 的证明已经表明，只须对
任意赋值环 $A$ 及任意态射 $\Spec(A) \to Y$ 证明，基变换
$f_A : X_A \to \Spec(A)$ 满足赋值判据的存在性部分。按定义，代数空间
$X_A$ 具有性质 $(\gamma)$，故可将命题
\ref{decent-spaces-proposition-characterize-universally-closed} 应用于态射 $f_A$，结论成立。
\end{proof}





\section{纤维中的点}
\label{decent-spaces-section-points-fibres}

\noindent
设 $S$ 为概形。考虑 $S$ 上代数空间的 Cartesian 图
\begin{equation}
\label{decent-spaces-equation-points-fibres}
\xymatrix{
W \ar[r]_q \ar[d]_p & Z \ar[d]^g \\
X \ar[r]^f & Y
}
\end{equation}
。设 $x \in |X|$ 与 $z \in |Z|$ 是映到同一点 $y \in |Y|$ 的点。可以问：
集合
\begin{equation}
\label{decent-spaces-equation-fibre}
F_{x, z} = \{ w \in |W| \text{ 满足 }p(w) = x\text{ 且 }q(w) = z\}
\end{equation}
何时有限？

\begin{example}
\label{decent-spaces-example-schemes}
若 $X, Y, Z$ 都是概形，则集合 $F_{x, z}$ 等于
$\kappa(x) \otimes_{\kappa(y)} \kappa(z)$ 的谱（《概形》引理
\ref{schemes-lemma-points-fibre-product}）。
% P08-E325: frozen source equates the set F_{x,z} with a spectrum rather than its underlying point set; intended point-set reading translated without changing formulas; see ERRATA_CANDIDATES.jsonl.
因此，若 $\kappa(y) \subset \kappa(x)$ 有限，或若
$\kappa(y) \subset \kappa(z)$ 有限，就得到有限集。特别地，若 $g$ 在 $z$
处拟有限，则总是如此（《态射》引理
\ref{morphisms-lemma-residue-field-quasi-finite}）。
\end{example}

\begin{example}
\label{decent-spaces-example-not-finite}
设 $K$ 为特征 $0$ 的域，并赋予无限阶自同构 $\sigma$。令
$Y = \Spec(K)/\mathbf{Z}$ 及 $X = \mathbf{A}^1_K/\mathbf{Z}$，其中
$\mathbf{Z}$ 通过 $\sigma$ 作用于 $K$，并通过 $t \mapsto t + 1$ 作用于
$\mathbf{A}^1_K = \Spec(K[t])$。令 $Z = \Spec(K)$。则
$W = \mathbf{A}^1_K$。图为
$$
\xymatrix{
\mathbf{A}^1_K \ar[r]_q \ar[d]_p & \Spec(K) \ar[d]^g \\
\mathbf{A}^1_K/\mathbf{Z} \ar[r]^f & \Spec(K)/\mathbf{Z}
}
$$
取对应于 $t = 0$ 的 $x$，并取 $\Spec(K)$ 的唯一一点 $z$。于是作为集合有
$F_{x, z} = \mathbf{Z}$。
\end{example}

\begin{lemma}
\label{decent-spaces-lemma-surjective-on-fibres}
在 (\ref{decent-spaces-equation-points-fibres}) 的情形下，若 $Z' \to Z$ 为态射，且
$z' \in |Z'|$ 映到 $z$，则诱导映射 $F_{x, z'} \to F_{x, z}$ 是满射。
\end{lemma}

\begin{proof}
令 $W' = X \times_Y Z' = W \times_Z Z'$。由《空间的性质》引理
\ref{spaces-properties-lemma-points-cartesian}，
$|W'| \to |W| \times_{|Z|} |Z'|$ 是满射。这就给出
$F_{x, z'} \to F_{x, z}$ 的满性。
% P08-E326: frozen source concludes with the sentence fragment “Hence the surjectivity ...”; intended conclusion translated; see ERRATA_CANDIDATES.jsonl.
\end{proof}

\begin{lemma}
\label{decent-spaces-lemma-qf-and-qc-finite-fibre}
在图 (\ref{decent-spaces-equation-points-fibres}) 中，若 $f$ 是有限型态射且 $f$ 在 $x$
处拟有限，则集合 (\ref{decent-spaces-equation-fibre}) 有限。
\end{lemma}

\begin{proof}
态射 $q$ 在每个 $w \in F_{x, z}$ 处拟有限；见《空间的态射》引理
\ref{spaces-morphisms-lemma-base-change-quasi-finite-locus}。因此本引理由
《空间的态射》引理
\ref{spaces-morphisms-lemma-quasi-finite-at-a-finite-number-of-points} 得出。
\end{proof}

\begin{lemma}
\label{decent-spaces-lemma-decent-finite-fibre}
在图 (\ref{decent-spaces-equation-points-fibres}) 中，若 $y$ 可由单态射
$\Spec(k) \to Y$ 表示，其中 $k$ 是域，且 $g$ 在 $z$ 处拟有限，则集合
(\ref{decent-spaces-equation-fibre}) 有限。（特例：$Y$ 良好且 $g$ 是 \'etale 态射。）
\end{lemma}

\begin{proof}
两次应用引理 \ref{decent-spaces-lemma-surjective-on-fibres}，可将 $Z$ 换成
$Z_k = \Spec(k) \times_Y Z$，将 $X$ 换成
$X_k = \Spec(k) \times_Y X$。也可以而且确实将 $Y$ 换成 $\Spec(k)$。
注意，由《空间的态射》引理
\ref{spaces-morphisms-lemma-base-change-quasi-finite-locus}，
$Z_k \to \Spec(k)$ 在 $z$ 处拟有限。取概形 $V$、一点 $v \in V$ 以及
将 $v$ 映到 $z$ 的 \'etale 态射 $V \to Z_k$。再取概形 $U$、一点
$u \in U$ 以及将 $u$ 映到 $x$ 的 \'etale 态射 $U \to X_k$。再次由引理
\ref{decent-spaces-lemma-surjective-on-fibres}，只须对图
$$
\xymatrix{
U \times_{\Spec(k)} V \ar[r] \ar[d] & V \ar[d] \\
U \ar[r] & \Spec(k)
}
$$
证明 $F_{u, v}$ 有限。态射 $V \to \Spec(k)$ 在 $v$ 处拟有限（这由
《空间的态射》第 \ref{spaces-morphisms-section-local-source-target} 节的
一般讨论及在一点处拟有限的定义得出）。此时，有限性由例
\ref{decent-spaces-example-schemes} 得出。引理陈述中的括注则来自以下两个事实：良好
空间上的点可由从域的谱出发的单态射表示；代数空间的 \'etale 态射局部
拟有限。
\end{proof}

\begin{lemma}
\label{decent-spaces-lemma-topology-fibre}
\begin{slogan}
代数空间的域值点之纤维具有预期的 Zariski 拓扑。
\end{slogan}
设 $S$ 为概形，$f : X \to Y$ 为 $S$ 上代数空间的态射。设
$y \in |Y|$，并假设 $y$ 由拟紧单态射 $\Spec(k) \to Y$ 表示。则
$|X_k| \to |X|$ 是到赋予诱导拓扑的 $f^{-1}(\{y\}) \subset |X|$ 上的
同胚。
\end{lemma}

\begin{proof}
以下将不再另行说明地使用《空间的性质》引理
\ref{spaces-properties-lemma-etale-open} 及《空间的态射》引理
\ref{spaces-morphisms-lemma-monomorphism-injective-points}。设
$V \to Y$ 为从仿射概形 $V$ 出发的 \'etale 态射，使得存在映到 $y$ 的
$v \in V$。由于 $\Spec(k) \to Y$ 拟紧，$V$ 中只有有限多个点映到 $y$
（引理 \ref{decent-spaces-lemma-UR-finite-above-x}）。缩小 $V$ 后，可设 $v$ 是唯一
这样的点。取概形 $U$ 及满 \'etale 态射 $U \to X$。考虑交换图
$$
\xymatrix{
U \ar[d] & U_V \ar[l] \ar[d] & U_v \ar[l] \ar[d] \\
X \ar[d] & X_V \ar[l] \ar[d] & X_v \ar[l] \ar[d] \\
Y & V \ar[l] & v \ar[l]
}
$$
由于 $U_v \to U_V$ 将 $U_v$ 等同于 $U_V$ 中赋予诱导拓扑的一个子集
（《概形》引理 \ref{schemes-lemma-fibre-topological}），又由于
$|U_V| \to |X_V|$ 与 $|U_v| \to |X_v|$ 都是满开映射，可知
$|X_v| \to |X_V|$ 是到其像（赋予诱导拓扑）上的同胚。另一方面，在开映射
$|X_V| \to |X|$ 下，$f^{-1}(\{y\})$ 的逆像等于 $|X_v|$。故
$|X_v| \to f^{-1}(\{y\})$ 是开映射。态射 $X_v \to X$ 分解经过 $X_k$；
由《空间的性质》引理
\ref{spaces-properties-lemma-points-cartesian}，$|X_k| \to |X|$ 单射，且像为
$f^{-1}(\{y\})$。由于 $X_v \to X_k$ 满，使用
$|X_v| \to |X_k| \to f^{-1}(\{y\})$ 即得本引理。
\end{proof}

\begin{lemma}
\label{decent-spaces-lemma-conditions-on-point-on-space-over-field}
设 $X$ 为域 $k$ 上局部有限型的代数空间，$x \in |X|$。考虑下列条件：
\begin{enumerate}
\item $\dim_x(|X|) = 0$；
\item $x$ 在 $|X|$ 中闭，且若在 $|X|$ 中有 $x' \leadsto x$，则 $x' = x$；
\item $x$ 是 $|X|$ 的孤立点；
\item $\dim_x(X) = 0$；
\item $X \to \Spec(k)$ 在 $x$ 处拟有限。
\end{enumerate}
则 (2)、(3)、(4) 与 (5) 等价。若 $X$ 良好，则 (1) 也与其余条件等价。
\end{lemma}

\begin{proof}
例如由《空间的态射》引理
\ref{spaces-morphisms-lemma-locally-finite-type-quasi-finite-part} 与
\ref{spaces-morphisms-lemma-quasi-finite-at-point}，(4) 与 (5) 等价。

\medskip\noindent
设 $U \to X$ 为 \'etale 态射，其中 $U$ 是仿射概形，并设 $u \in U$ 为
映到 $x$ 的一点。此外，若 $x$ 是闭点，例如在 (2) 或 (3) 的情形下，则
可以而且确实假设 $u$ 是闭点。按定义
$\dim_u(U) = \dim_x(X)$；若 $u$ 是闭点，则这又等于
$\dim(\mathcal{O}_{U, u})$，见《代数》引理
\ref{algebra-lemma-dimension-closed-point-finite-type-field}。

\medskip\noindent
若 $\dim_x(X) > 0$ 且 $u$ 闭，则由上述论证，可在 $U$ 中选取非平凡特化
$u' \leadsto u$。于是 $\kappa(u')$ 在 $k$ 上的超越次数大于
$\kappa(u)$ 在 $k$ 上的超越次数。因为 $x/k$ 与 $x'/k$ 的超越次数有定义，
所以它们在 $X$ 中的像 $x$ 与 $x'$ 不同；见《空间的态射》定义
\ref{spaces-morphisms-definition-dimension-fibre}。这尤其适用于 (2) 与 (3)
的情形，故 (2) 与 (3) 蕴含 (4)。

\medskip\noindent
反之，若 $X \to \Spec(k)$ 在 $x$ 处局部拟有限，则
$U \to \Spec(k)$ 在 $u$ 处局部拟有限，因而 $u$ 是 $U$ 的孤立点
（《态射》引理 \ref{morphisms-lemma-quasi-finite-at-point-characterize}）。
由于 $|U| \to |X|$ 连续且开，由此可知 (5) 蕴含 (2) 与 (3)。

\medskip\noindent
假设 $X$ 良好且 (1) 成立。由引理
\ref{decent-spaces-lemma-dimension-decent-space}，$\dim_x(X) = \dim_x(|X|)$，证明完成。
\end{proof}

\begin{lemma}
\label{decent-spaces-lemma-conditions-on-space-over-field}
设 $X$ 为域 $k$ 上局部有限型的代数空间。考虑下列条件：
\begin{enumerate}
\item $|X|$ 是有限集；
\item $|X|$ 是离散空间；
\item $\dim(|X|) = 0$；
\item $\dim(X) = 0$；
\item $X \to \Spec(k)$ 局部拟有限；
\end{enumerate}
则 (2)、(3)、(4) 与 (5) 等价。若 $X$ 良好，则 (1) 蕴含其余条件。
\end{lemma}

\begin{proof}
例如由《空间的态射》引理
\ref{spaces-morphisms-lemma-locally-finite-type-quasi-finite-part}，(4)
与 (5) 等价。

\medskip\noindent
设 $U \to X$ 为满 \'etale 态射，其中 $U$ 是概形。

\medskip\noindent
若 $\dim(U) > 0$，则在 $U$ 中选取非平凡特化 $u \leadsto u'$；此时
$\kappa(u)$ 在 $k$ 上的超越次数大于 $\kappa(u')$ 在 $k$ 上的超越次数。
因为 $x/k$ 与 $x'/k$ 的超越次数有定义，所以它们在 $X$ 中的像 $x$ 与
$x'$ 不同；见《空间的态射》定义
\ref{spaces-morphisms-definition-dimension-fibre}。故 (2) 与 (3) 蕴含 (4)。

\medskip\noindent
反之，若 $X \to \Spec(k)$ 局部拟有限，则 $U$ 局部 Noether
（《态射》引理 \ref{morphisms-lemma-finite-type-noetherian}），且维数为
$0$（《态射》引理
\ref{morphisms-lemma-locally-quasi-finite-rel-dimension-0}），因而是 Artin
局部环之谱的不交并（《性质》引理
\ref{properties-lemma-locally-Noetherian-dimension-0}）。所以 $U$ 是离散
拓扑空间；又因 $|U| \to |X|$ 连续且开，$|X|$ 也如此。换言之，(4) 蕴含
(2) 与 (3)。

\medskip\noindent
假设 $X$ 良好且 (1) 成立。此时可选取上述 $U$ 为仿射概形。
$|U| \to |X|$ 的纤维有限（这是良好空间定义性质的一部分）。故 $U$ 是
$k$ 上只有有限多个点的有限型概形。因此 $U$ 在 $k$ 上拟有限
（《态射》引理 \ref{morphisms-lemma-finite-fibre}）；按定义，这意味着
$X \to \Spec(k)$ 局部拟有限。
\end{proof}

\begin{lemma}
\label{decent-spaces-lemma-conditions-on-point-in-fibre-and-qf}
设 $S$ 为概形，$f : X \to Y$ 为 $S$ 上代数空间的局部有限型态射。设
$x \in |X|$ 的像为 $y \in |Y|$。令 $F = f^{-1}(\{y\})$，并赋予从 $|X|$
诱导的拓扑。设 $k$ 为域，且 $\Spec(k) \to Y$ 属于定义 $y$ 的等价类。令
$X_k = \Spec(k) \times_Y X$。设 $\tilde x \in |X_k|$ 映到
$x \in |X|$。考虑下列条件：
\begin{enumerate}
\item
\label{decent-spaces-item-fibre-at-x-dim-0}
$\dim_x(F) = 0$；
\item
\label{decent-spaces-item-isolated-in-fibre}
$x$ 在 $F$ 中孤立；
\item
\label{decent-spaces-item-no-specializations-in-fibre}
$x$ 在 $F$ 中闭，且若在 $F$ 中有 $x' \leadsto x$，则 $x = x'$；
\item
\label{decent-spaces-item-dimension-top-k-fibre}
$\dim_{\tilde x}(|X_k|) = 0$；
\item
\label{decent-spaces-item-isolated-in-k-fibre}
$\tilde x$ 在 $|X_k|$ 中孤立；
\item
\label{decent-spaces-item-no-specializations-in-k-fibre}
$\tilde x$ 在 $|X_k|$ 中闭，且若在 $|X_k|$ 中有
$\tilde x' \leadsto \tilde x$，则 $\tilde x = \tilde x'$；
\item
\label{decent-spaces-item-k-fibre-at-x-dim-0}
$\dim_{\tilde x}(X_k) = 0$；
\item
\label{decent-spaces-item-quasi-finite-at-x}
$f$ 在 $x$ 处拟有限。
\end{enumerate}
则有
$$
\xymatrix{
(\ref{decent-spaces-item-dimension-top-k-fibre}) \ar@{=>}[r]_{f\text{ 良好}} &
(\ref{decent-spaces-item-isolated-in-k-fibre}) \ar@{<=>}[r] &
(\ref{decent-spaces-item-no-specializations-in-k-fibre}) \ar@{<=>}[r] &
(\ref{decent-spaces-item-k-fibre-at-x-dim-0}) \ar@{<=>}[r] &
(\ref{decent-spaces-item-quasi-finite-at-x})
}
$$
若 $Y$ 良好，则条件 (\ref{decent-spaces-item-isolated-in-fibre}) 与
(\ref{decent-spaces-item-no-specializations-in-fibre}) 彼此等价，并且等价于条件
(\ref{decent-spaces-item-isolated-in-k-fibre})、
(\ref{decent-spaces-item-no-specializations-in-k-fibre})、
(\ref{decent-spaces-item-k-fibre-at-x-dim-0}) 以及
(\ref{decent-spaces-item-quasi-finite-at-x})。
若 $Y$ 与 $X$ 都良好，则所有条件等价。
\end{lemma}

\begin{proof}
由引理 \ref{decent-spaces-lemma-conditions-on-point-on-space-over-field}，条件
(\ref{decent-spaces-item-isolated-in-k-fibre}),
(\ref{decent-spaces-item-no-specializations-in-k-fibre}) 与
(\ref{decent-spaces-item-k-fibre-at-x-dim-0}) 彼此等价，并且等价于
$X_k \to \Spec(k)$ 在 $\tilde x$ 处拟有限这一条件。因此由《空间的态射》引理
\ref{spaces-morphisms-lemma-base-change-quasi-finite-locus}
，它们也等价于 (\ref{decent-spaces-item-quasi-finite-at-x})。若 $f$ 良好，则 $X_k$ 是
良好代数空间；引理
\ref{decent-spaces-lemma-conditions-on-point-on-space-over-field} 表明
(\ref{decent-spaces-item-dimension-top-k-fibre}) 蕴含
(\ref{decent-spaces-item-isolated-in-k-fibre})。

\medskip\noindent
若 $Y$ 良好，则可在 $y$ 的等价类中选取拟紧单态射
$\Spec(k') \to Y$。此时，引理 \ref{decent-spaces-lemma-topology-fibre} 表明
$|X_{k'}| \to F$ 是同胚。结合上述论证，即得引理中的其余断言；细节略去。
\end{proof}

\begin{lemma}
\label{decent-spaces-lemma-conditions-on-fibre-and-qf}
设 $S$ 为概形，$f : X \to Y$ 为 $S$ 上代数空间的局部有限型态射。设
$y \in |Y|$。设 $k$ 为域，且 $\Spec(k) \to Y$ 属于定义 $y$ 的等价类。
令 $X_k = \Spec(k) \times_Y X$，并令 $F = f^{-1}(\{y\})$ 赋予从 $|X|$
诱导的拓扑。考虑下列条件：
\begin{enumerate}
\item
\label{decent-spaces-item-fibre-finite}
$F$ 有限；
\item
\label{decent-spaces-item-fibre-discrete}
$F$ 是离散拓扑空间；
\item
\label{decent-spaces-item-fibre-no-specializations}
$\dim(F) = 0$；
\item
\label{decent-spaces-item-k-fibre-finite}
$|X_k|$ 是有限集；
\item
\label{decent-spaces-item-k-fibre-discrete}
$|X_k|$ 是离散空间；
\item
\label{decent-spaces-item-k-fibre-no-specializations}
$\dim(|X_k|) = 0$；
\item
\label{decent-spaces-item-k-fibre-dim-0}
$\dim(X_k) = 0$；
\item
\label{decent-spaces-item-quasi-finite-at-points-fibre}
$f$ 在 $|X|$ 中位于 $y$ 上方的所有点处拟有限。
\end{enumerate}
则有
$$
\xymatrix{
(\ref{decent-spaces-item-fibre-finite}) &
(\ref{decent-spaces-item-k-fibre-finite}) \ar@{=>}[l] \ar@{=>}[r]_{f\text{ 良好}} &
(\ref{decent-spaces-item-k-fibre-discrete}) \ar@{<=>}[r] &
(\ref{decent-spaces-item-k-fibre-no-specializations}) \ar@{<=>}[r] &
(\ref{decent-spaces-item-k-fibre-dim-0}) \ar@{<=>}[r] &
(\ref{decent-spaces-item-quasi-finite-at-points-fibre})
}
$$
若 $Y$ 良好，则条件 (\ref{decent-spaces-item-fibre-discrete}) 与
(\ref{decent-spaces-item-fibre-no-specializations})
彼此等价，并且等价于条件 (\ref{decent-spaces-item-k-fibre-discrete})、
(\ref{decent-spaces-item-k-fibre-no-specializations})、(\ref{decent-spaces-item-k-fibre-dim-0}) 及
(\ref{decent-spaces-item-quasi-finite-at-points-fibre}).
若 $Y$ 与 $X$ 都良好，则 (\ref{decent-spaces-item-fibre-finite}) 蕴含所有其他条件。
\end{lemma}

\begin{proof}
由引理 \ref{decent-spaces-lemma-conditions-on-space-over-field}，条件
(\ref{decent-spaces-item-k-fibre-discrete})、(\ref{decent-spaces-item-k-fibre-no-specializations}) 与
(\ref{decent-spaces-item-k-fibre-dim-0}) 彼此等价，并且等价于
$X_k \to \Spec(k)$ 局部拟有限这一条件。因此由《空间的态射》引理
\ref{spaces-morphisms-lemma-base-change-quasi-finite-locus}
，它们也等价于 (\ref{decent-spaces-item-quasi-finite-at-points-fibre})。若 $f$ 良好，则
$X_k$ 是良好代数空间；引理
\ref{decent-spaces-lemma-conditions-on-space-over-field} 表明
(\ref{decent-spaces-item-k-fibre-finite}) 蕴含 (\ref{decent-spaces-item-k-fibre-discrete})。

\medskip\noindent
由《空间的性质》引理 \ref{spaces-properties-lemma-points-cartesian}，映射
$|X_k| \to F$ 是满射，故有
(\ref{decent-spaces-item-k-fibre-finite}) $\Rightarrow$ (\ref{decent-spaces-item-fibre-finite}).

\medskip\noindent
若 $Y$ 良好，则可在 $y$ 的等价类中选取拟紧单态射
$\Spec(k') \to Y$。此时，引理 \ref{decent-spaces-lemma-topology-fibre} 表明
$|X_{k'}| \to F$ 是同胚。结合上述论证，即得引理中的其余断言；细节略去。
\end{proof}







\section{单态射}
\label{decent-spaces-section-monomorphisms}

\noindent
下面给出单态射可表示的另一种情形。更多信息见《空间态射进阶》第
\ref{spaces-more-morphisms-section-monomorphisms}
节。

\begin{lemma}
\label{decent-spaces-lemma-monomorphism-toward-disjoint-union-dim-0-rings}
设 $S$ 为概形。设 $Y$ 是 $S$ 上若干零维局部环之谱的无交并。
设 $f : X \to Y$ 是 $S$ 上代数空间的单态射。则 $f$ 可表示，即
$X$ 是概形。
\end{lemma}

\begin{proof}
立即可归约到 $Y = \Spec(A)$ 的情形，其中 $A$ 是零维局部环，即
$\Spec(A) = \{\mathfrak m_A\}$
是单点集。若 $X = \emptyset$，则无需证明。否则，选取非空仿射概形
$U = \Spec(B)$ 以及一个 \'etale 态射 $U \to X$。由于 $|X|$ 是单点集
（它是 $|Y|$ 的子集；见《空间态射》引理
\ref{spaces-morphisms-lemma-monomorphism-injective-points}），可知
$U \to X$ 满射。注意
$U \times_X U = U \times_Y U = \Spec(B \otimes_A B)$.
因此两个环同态 $B \to B \otimes_A B$ 都是 \'etale 的。由于
$$
(B \otimes_A B)/\mathfrak m_A(B \otimes_A B)
=
(B/\mathfrak m_AB) \otimes_{A/\mathfrak m_A} (B/\mathfrak m_AB)
$$
可知
$B/\mathfrak m_AB \to (B \otimes_A B)/\mathfrak m_A(B \otimes_A B)$
平坦；事实上，它是自由的，其秩等于 $B/\mathfrak m_AB$ 作为
$A/\mathfrak m_A$-向量空间的维数。由于
$B \to B \otimes_A B$ 是 \'etale 的，只有当该维数有限时才可能如此
（例如见《态射》引理 \ref{morphisms-lemma-etale-universally-bounded} 与
\ref{morphisms-lemma-locally-quasi-finite-qc-source-universally-bounded}）。
$B$ 的每个素理想都位于 $\mathfrak m_A$（即 $A$ 的唯一素理想）上方。
因此作为拓扑空间有 $\Spec(B) = \Spec(B/\mathfrak m_A)$；又因
$B/\mathfrak m_A B$ 是 Artin 环，此空间是有限离散集，见《代数》引理
\ref{algebra-lemma-finite-dimensional-algebra} 与
\ref{algebra-lemma-artinian-finite-length}。所以 $B$ 的所有素理想都是
极大理想，且 $B = B_1 \times \ldots \times B_n$ 是有限多个零维局部环
的乘积，见《代数》引理 \ref{algebra-lemma-product-local}。由《代数》
引理 \ref{algebra-lemma-local-dimension-zero-henselian}，所有局部环
$B_i$ 都是 Hensel 的，故 $B \to B \otimes_A B$ 是有限 \'etale 的。
于是由《群胚》命题 \ref{groupoids-proposition-finite-flat-equivalence}，
$X$ 是仿射概形。
\end{proof}








\section{一般点}
\label{decent-spaces-section-generic-points}

\noindent
本节接续《空间性质》第
\ref{spaces-properties-section-generic-points} 节。

\begin{lemma}
\label{decent-spaces-lemma-decent-generic-points}
设 $S$ 为概形，$X$ 为 $S$ 上的良好代数空间，并设 $x \in |X|$。
下列条件等价：
\begin{enumerate}
\item $x$ 是 $|X|$ 某个不可约分支的一般点；
\item 对任意带基点代数空间的 \'etale 态射
$(Y, y) \to (X, x)$，$y$ 是 $|Y|$ 某个不可约分支的一般点；
\item 对某个带基点代数空间的 \'etale 态射
$(Y, y) \to (X, x)$，$y$ 是 $|Y|$ 某个不可约分支的一般点；
\item $X$ 在 $x$ 处的局部环维数为零；
\item $x$ 是 $X$ 上余维为 $0$ 的点。
\end{enumerate}
\end{lemma}

\begin{proof}
由定义，条件 (4) 与 (5) 对任意代数空间等价；见《空间性质》定义
\ref{spaces-properties-definition-dimension-local-ring}。由引理
\ref{decent-spaces-lemma-etale-named-properties}，(2) 和 (3) 中的任何 $Y$ 都是良好的。
因此只需证明 (1) 与 (4) 等价；因为 $Y$ 在 $y$ 处的局部环与 $X$ 在 $x$
处的局部环维数相同，届时它们与 (2)、(3) 的等价性也随之成立。
取从仿射概形出发的 \'etale 态射 $f : U \to X$，并取映到 $x$ 的点
$u \in U$。

\medskip\noindent
假设 (1)。设 $u' \leadsto u$ 是 $U$ 中的特化。则
$f(u') = f(u) = x$。由引理
\ref{decent-spaces-lemma-decent-no-specializations-map-to-same-point}，有 $u' = u$。
因此 $u$ 是 $U$ 某个不可约分支的一般点。于是
$\dim(\mathcal{O}_{U, u}) = 0$，故 (4) 成立。

\medskip\noindent
假设 (4)。点 $x$ 包含在某个不可约分支 $T \subset |X|$ 中。由于
$|X|$ 是清醒的（命题 \ref{decent-spaces-proposition-reasonable-sober}），
% P08-E327: frozen English source omits the verb in “we T has”.
$T$ 有一般点 $x'$。当然 $x' \leadsto x$。由引理
\ref{decent-spaces-lemma-decent-specialization}，可将此特化提升为 $U$ 中的
$u' \leadsto u$。除非 $u' = u$，即 $x' = x$，否则这与
$\dim(\mathcal{O}_{U, u}) = 0$ 的假设矛盾。
\end{proof}

\begin{lemma}
\label{decent-spaces-lemma-codimension-local-ring}
设 $S$ 为概形，$X$ 为 $S$ 上的良好代数空间。设
$T \subset |X|$ 为不可约闭子集，并设 $\xi \in T$ 为其一般点
（命题 \ref{decent-spaces-proposition-reasonable-sober}）。则
$\text{codim}(T, |X|)$（《拓扑》定义
\ref{topology-definition-codimension}）等于 $X$ 在 $\xi$ 处的局部环维数
（《空间性质》定义
\ref{spaces-properties-definition-dimension-local-ring}）。
\end{lemma}

\begin{proof}
选取概形 $U$、点 $u \in U$ 以及把 $u$ 映到 $\xi$ 的 \'etale 态射
$U \to X$。由引理 \ref{decent-spaces-lemma-decent-specialization}，任意非平凡特化链
$\xi_e \leadsto \ldots \leadsto \xi_0 = \xi$
都可提升为 $U$ 中的链
$u_e \leadsto \ldots \leadsto u_0 = u$。反过来，由引理
\ref{decent-spaces-lemma-decent-no-specializations-map-to-same-point}，$U$ 中任意非平凡
特化链
$u_e \leadsto \ldots \leadsto u_0 = u$
都映成非平凡特化链
$\xi_e \leadsto \ldots \leadsto \xi_0 = \xi$。
因为 $|X|$ 与 $U$ 都是清醒拓扑空间，所以 $T$ 在 $|X|$ 中的余维等于
$\overline{\{u\}}$ 在 $U$ 中的余维。这样，引理归约到概形情形，即
《性质》引理 \ref{properties-lemma-codimension-local-ring}。
\end{proof}

\begin{lemma}
\label{decent-spaces-lemma-get-reasonable}
设 $S$ 为概形，$X$ 为 $S$ 上的代数空间。假设
\begin{enumerate}
\item $X$ 上每个拟紧 \'etale 概形都只有有限多个不可约分支；
\item $X$ 上每个余维为 $0$ 的点 $x \in |X|$ 都可由单态射
$\Spec(k) \to X$ 表示。
\end{enumerate}
则 $X$ 是合理代数空间。
\end{lemma}

\begin{proof}
设 $U$ 为仿射概形，$a : U \to X$ 为 \'etale 态射。需要证明 $a$ 的纤维
普遍有界。由假设 (1)，概形 $U$ 只有有限多个不可约分支。令
$u_1, \ldots, u_n \in U$ 为这些不可约分支的一般点，并令
$\{x_1, \ldots, x_m\} \subset |X|$ 为
$\{u_1, \ldots, u_n\}$ 的像。每个 $x_j$ 都是余维为 $0$ 的点。
由假设 (2)，可选取表示 $x_j$ 的单态射 $\Spec(k_j) \to X$。
由《空间性质》引理
\ref{spaces-properties-lemma-codimension-0-points}，有
$$
U \times_X \Spec(k_j) = \coprod\nolimits_{a(u_i) = x_j} \Spec(\kappa(u_i))
$$
这是 $\Spec(k_j)$ 上次数为
$d_j = \sum_{a(u_i) = x_j} [\kappa(u_i) : k_j]$ 的有限概形。令
$n = \max d_j$。

\medskip\noindent
注意，$a$ 是分离的（《空间性质》引理
\ref{spaces-properties-lemma-separated-cover}）。考虑引理
\ref{decent-spaces-lemma-stratify-flat-fp-lqf} 中与 $U \to X$ 相伴的分层
$$
X = X_0 \supset X_1 \supset X_2 \supset \ldots
$$
由上述 $n$ 的选取，可知 $X_{n + 1}$ 为空。事实上，否则
% P08-E328: frozen source says the open subset of U contains an x_i, but x_i lies in |X|.
$a^{-1}(X_{n + 1})$ 是 $U$ 的非空开集，因而会包含某个 $x_i$。
这将意味着 $X_{n + 1}$ 包含 $x_j = a(u_i)$，矛盾。因此
$U \to X$ 的纤维普遍有界（实际上由整数 $n$ 所界）。
\end{proof}

\begin{lemma}
\label{decent-spaces-lemma-finitely-many-irreducible-components}
设 $S$ 为概形，$X$ 为 $S$ 上的代数空间。下列条件等价：
\begin{enumerate}
\item $X$ 是良好的，且 $|X|$ 只有有限多个不可约分支；
\item $X$ 上每个拟紧 \'etale 概形都只有有限多个不可约分支，
$X$ 上只有有限多个余维为 $0$ 的点 $x \in |X|$，且每个这样的点都可
由单态射 $\Spec(k) \to X$ 表示；
\item 存在一个稠密开子空间 $X' \subset X$，它是概形；$X'$ 只有有限多个
不可约分支，其一般点为 $\{x'_1, \ldots, x'_m\}$，并且对
$j = 1, \ldots, m$，态射 $x'_j \to X$ 拟紧。
\end{enumerate}
此外，若这些条件成立，则 $X$ 是合理的，而各点 $x'_j \in |X|$ 正是
$|X|$ 的不可约分支的一般点。
\end{lemma}

\begin{proof}
在证明中，我们将不再特别说明而使用《空间性质》引理
\ref{spaces-properties-lemma-codimension-0-points}
。假设 (1)。由定理 \ref{decent-spaces-theorem-decent-open-dense-scheme}，$X$ 有一个
稠密开子概形 $X'$。由于 $|X'|$ 的每个不可约分支的闭包都是 $|X|$
的不可约分支，可知 $|X'|$ 只有有限多个不可约分支。因此 (3) 成立。

\medskip\noindent
假设 $X' \subset X$ 如 (3) 所述。令
$\{x'_1, \ldots, x'_m\}$ 为 $X'$ 的不可约分支的一般点。设
$a : U \to X$ 为 \'etale 态射，其中 $U$ 是拟紧概形。为证明 (2)，只需
证明 $U$ 只有有限多个不可约分支，而且它们的一般点位于
$\{x'_1, \ldots, x'_m\}$ 上方。只需对 $U$ 的某个有限仿射开覆盖中的
各成员证明此事，故可且确实假设 $U$ 仿射。注意
$U' = a^{-1}(X') \subset U$ 是稠密开集。由于 $U' \to X'$ 是概形的
\'etale 态射，$U'$ 的不可约分支的一般点正是位于
$\{x'_1, \ldots, x'_m\}$ 上方的点。又因 $x'_j \to X$ 拟紧，$U$ 中位于
$x'_j$ 上方的点只有有限多个（引理 \ref{decent-spaces-lemma-UR-finite-above-x}）。
因此 $U'$ 只有有限多个不可约分支，而这些不可约分支的闭包就是 $U$ 的
不可约分支。故 (2) 成立。

\medskip\noindent
假设 (2)。由引理 \ref{decent-spaces-lemma-get-reasonable} 可得 (1) 以及最后的断言。
（这里还用到合理代数空间是良好的；见定义
\ref{decent-spaces-definition-very-reasonable} 后的讨论。）
\end{proof}



\section{一般有限态射}
\label{decent-spaces-section-generically-finite}

\noindent
本节针对代数空间的态射，讨论《态射》第
\ref{morphisms-section-generically-finite} 节和《簇》第
\ref{varieties-section-generically-finite} 节中针对概形态射所讨论的内容。

\begin{lemma}
\label{decent-spaces-lemma-generically-finite}
设 $S$ 为概形，$f : X \to Y$ 为 $S$ 上代数空间的态射。假设 $f$
拟分离且为有限型。设 $y \in |Y|$ 是 $Y$ 上余维为 $0$ 的点。
下列条件等价：
\begin{enumerate}
\item 若 $\Spec(k) \to Y$ 表示 $y$，则空间 $|X_k|$ 有限；
\item $X \to Y$ 在 $|X|$ 中位于 $y$ 上方的所有点处拟有限；
\item 存在开子空间 $Y' \subset Y$，其中 $y \in |Y'|$，使得
$Y' \times_Y X \to Y'$ 有限。
\end{enumerate}
若 $Y$ 良好，则它们还等价于
\begin{enumerate}
\item[(4)] 集合 $f^{-1}(\{y\})$ 有限。
\end{enumerate}
\end{lemma}

\begin{proof}
(1) 与 (2) 的等价性来自引理 \ref{decent-spaces-lemma-conditions-on-fibre-and-qf}
（以及由引理 \ref{decent-spaces-lemma-properties-trivial-implications} 得到的事实：
拟分离态射是良好的）。

\medskip\noindent
假设等价条件 (1) 和 (2) 成立。选取仿射概形 $V$ 以及把某个点
$v \in V$ 映到 $y$ 的 \'etale 态射 $V \to Y$。由《空间性质》引理
\ref{spaces-properties-lemma-codimension-0-points}，$v$ 是 $V$ 某个
不可约分支的一般点。选取仿射概形 $U$ 以及满射 \'etale 态射
$U \to V \times_Y X$。于是 $U \to V$ 为有限型。由 (2)，态射
$U \to V$ 在位于 $v$ 上方的每个点处拟有限。因此 $U \to V$ 在 $v$
上方的纤维有限（《态射》引理
\ref{morphisms-lemma-quasi-finite-at-a-finite-number-of-points}）。
由《态射》引理 \ref{morphisms-lemma-generically-finite}，缩小 $V$ 后
可假设 $U \to V$ 有限。令
$$
R = U \times_{V \times_Y X} U
$$
由于 $f$ 拟分离，$V \times_Y X$ 拟分离，因而 $R$ 是拟紧概形。
此外，
% P08-E329: frozen source has plural “morphisms” with singular “is”.
态射 $R \to V$ 是 \'etale 态射 $R \to U$ 与有限态射
$U \to V$ 的复合，故拟有限。因此可再次应用《态射》引理
\ref{morphisms-lemma-generically-finite}；缩小 $V$ 后，可假设
$R \to V$ 也有限。当然，这意味着两个投影
% P08-E330: frozen source writes R \to V for the two projections; their targets are U.
$R \to V$ 都是有限 \'etale 的。于是
$V/R = V \times_Y X$ 是仿射概形；见《群胚》命题
\ref{groupoids-proposition-finite-flat-equivalence}。由《态射》引理
\ref{morphisms-lemma-image-proper-is-proper}，可知
$V \times_Y X \to V$ 是固有的；再由《态射》引理
\ref{morphisms-lemma-finite-proper}，可知 $V \times_Y X \to V$ 有限。
最后，令 $Y' \subset Y$ 为与 $|V| \to |Y|$ 的像相对应的 $Y$ 的开子空间。
由于到 $V$ 的基变换有限，且 $V \to Y'$ 是满射 \'etale 态射，由
《空间态射》引理 \ref{spaces-morphisms-lemma-integral-local} 可知
$Y' \times_Y X \to Y'$ 有限。

\medskip\noindent
若 $Y$ 良好且 $f$ 拟分离，则 $X$ 也良好；这里使用引理
\ref{decent-spaces-lemma-properties-trivial-implications} 与
\ref{decent-spaces-lemma-property-over-property}。因此应用引理
\ref{decent-spaces-lemma-conditions-on-fibre-and-qf} 可知 (4) 蕴含 (1) 和 (2)。
另一方面，由《空间态射》引理
\ref{spaces-morphisms-lemma-quasi-finite-at-a-finite-number-of-points}
可知 (2) 蕴含 (4)。
\end{proof}

\begin{lemma}
\label{decent-spaces-lemma-generically-finite-reprise}
设 $S$ 为概形，$f : X \to Y$ 为 $S$ 上代数空间的态射。假设 $f$
拟分离且局部有限型，并且 $Y$ 拟分离。设 $y \in |Y|$ 为 $Y$ 上余维
为 $0$ 的点。下列条件等价：
\begin{enumerate}
\item 集合 $f^{-1}(\{y\})$ 有限；
\item 若 $\Spec(k) \to Y$ 表示 $y$，则空间 $|X_k|$ 有限；
\item 存在开子空间 $X' \subset X$ 与 $Y' \subset Y$，满足
$f(X') \subset Y'$、$y \in |Y'|$ 及 $f^{-1}(\{y\}) \subset |X'|$，
使得 $f|_{X'} : X' \to Y'$ 有限。
\end{enumerate}
\end{lemma}

\begin{proof}
由于拟分离代数空间是良好的，(1) 与 (2) 的等价性来自引理
\ref{decent-spaces-lemma-conditions-on-fibre-and-qf}。为证明 (1) 和 (2) 蕴含 (3)，
可且确实用包含 $y$ 的拟紧开集替换 $Y$。由于 $f^{-1}(\{y\})$ 有限，
可以找到包含该纤维的拟紧开子空间
% P08-E331: frozen source has the malformed phrase “open subspace of X' subset X”.
$X' \subset X$。由《空间态射》引理
\ref{spaces-morphisms-lemma-quasi-compact-quasi-separated-permanence}
（这里用到 $Y$ 拟分离），限制 $f|_{X'} : X' \to Y$ 拟紧且拟分离。
将引理 \ref{decent-spaces-lemma-generically-finite} 应用于
$f|_{X'} : X' \to Y$，即得 (3)。略去 (3) 蕴含 (2) 的证明。
\end{proof}

\begin{lemma}
\label{decent-spaces-lemma-quasi-finiteness-over-generic-point}
设 $S$ 为概形，$f : X \to Y$ 为 $S$ 上代数空间的态射。假设 $f$
局部有限型。分别以 $X^0 \subset |X|$ 和 $Y^0 \subset |Y|$ 表示
$X$ 和 $Y$ 的余维为 $0$ 的点集。设 $y \in Y^0$。下列条件等价：
\begin{enumerate}
\item $f^{-1}(\{y\}) \subset X^0$；
\item $f$ 在位于 $y$ 上方的所有点处拟有限；
\item $f$ 在位于 $y$ 上方的所有 $x \in X^0$ 处拟有限。
\end{enumerate}
\end{lemma}

\begin{proof}
设 $V$ 为概形，且 $V \to Y$ 为满射 \'etale 态射。设 $U$ 为概形，
且 $U \to V \times_Y X$ 为满射 \'etale 态射。则 $f$ 在点
$u \in U$ 的像 $x$ 处拟有限，当且仅当 $U \to V$ 在 $u$ 处拟有限。
此外，$x \in X^0$ 当且仅当 $u$ 是 $U$ 某个不可约分支的一般点
（《空间性质》引理
\ref{spaces-properties-lemma-codimension-0-points}）。因此本引理归约到
态射 $U \to V$ 的情形，即《态射》引理
\ref{morphisms-lemma-quasi-finiteness-over-generic-point}。
\end{proof}

\begin{lemma}
\label{decent-spaces-lemma-finite-over-dense-open}
设 $S$ 为概形，$f : X \to Y$ 为 $S$ 上代数空间的态射。假设 $f$
局部有限型。分别以 $X^0 \subset |X|$ 和 $Y^0 \subset |Y|$ 表示
$X$ 和 $Y$ 的余维为 $0$ 的点集。假设
\begin{enumerate}
\item $Y$ 良好；
\item $X^0$ 与 $Y^0$ 有限，且 $f^{-1}(Y^0) = X^0$；
\item $f$ 拟紧或 $f$ 分离。
\end{enumerate}
则存在稠密开子空间 $V \subset Y$，使得 $f^{-1}(V) \to V$ 有限。
\end{lemma}

\begin{proof}
由引理 \ref{decent-spaces-lemma-finitely-many-irreducible-components} 和
\ref{decent-spaces-lemma-decent-generic-points}，可假设 $Y$ 是只有有限多个不可约分支
的概形。进一步缩小后，可假设 $Y$ 是一般点为 $y$ 的不可约仿射概形。
此时 $f$ 在 $y$ 上方的纤维有限。

\medskip\noindent
假设 $f$ 拟紧，且 $Y$ 仿射不可约。此时 $X$ 拟紧，可以选取仿射概形
$U$ 以及满射 \'etale 态射 $U \to X$。于是 $U \to Y$ 为有限型，而
$U \to Y$ 在 $y$ 上方的纤维就是 $U$ 的不可约分支的一般点集 $U^0$
（《空间性质》引理
\ref{spaces-properties-lemma-codimension-0-points}）。因此 $U^0$ 有限
（《态射》引理
\ref{morphisms-lemma-quasi-finite-at-a-finite-number-of-points}），
并且缩小 $Y$ 后，可假设 $U \to Y$ 有限（《态射》引理
\ref{morphisms-lemma-generically-finite}）。接着考虑
$R = U \times_X U$。由于投影 $s : R \to U$ 是 \'etale 的，可知
$R^0 = s^{-1}(U^0)$ 位于 $y$ 上方。由于
$R \to U \times_Y U$ 是单态射，且 $U \times_Y U \to Y$ 有限，
故 $R^0$ 有限。并且 $R$ 分离（《空间性质》引理
\ref{spaces-properties-lemma-separated-cover}）。因此可以再次缩小 $Y$，
达到 $R$ 在 $Y$ 上有限的情形（《态射》引理
\ref{morphisms-lemma-finite-over-dense-open}）。在此情形下，完全沿用引理
\ref{decent-spaces-lemma-generically-finite} 证明中的论证，即得 $X = U/R$ 在 $Y$
上有限（也可以直接应用该引理，因为此时立即可知 $X$ 也拟分离）。

\medskip\noindent
假设 $f$ 分离，且 $Y$ 仿射不可约。按引理
\ref{decent-spaces-lemma-generically-finite-reprise} 选取 $V \subset Y$ 和
$U \subset X$。由于 $f|_U : U \to V$ 有限，可知
$U \subset f^{-1}(V)$ 既开又闭（《空间态射》引理
\ref{spaces-morphisms-lemma-universally-closed-permanence} 与
\ref{spaces-morphisms-lemma-finite-proper}）。因此对 $X$ 的某个开子空间
$W$，有 $f^{-1}(V) = U \amalg W$。然而，由于 $U$ 包含 $X$ 中所有余维为
$0$ 的点，由《空间性质》引理
\ref{spaces-properties-lemma-codimension-0-points-dense} 可知
$W = \emptyset$，这正是所需结论。
\end{proof}






\section{双有理态射}
\label{decent-spaces-section-birational}

\noindent
下面对代数空间的双有理态射之定义，似乎最接近我们对概形的双有理态射
所作的定义（《态射》定义 \ref{morphisms-definition-birational}）。

\begin{definition}
\label{decent-spaces-definition-birational}
设 $S$ 为概形。设 $X$ 和 $Y$ 为 $S$ 上的代数空间。
% P08-E332: frozen English source omits “be” after X and Y.
假设 $X$ 和 $Y$ 良好，并且 $|X|$ 与 $|Y|$ 都只有有限多个不可约分支。
若态射 $f : X \to Y$ 满足
\begin{enumerate}
\item $|f|$ 在 $|X|$ 的不可约分支的一般点集与 $|Y|$ 的不可约分支的
一般点集之间诱导双射；
\item 对每个不可约分支的一般点 $x \in |X|$，局部环同态
$\mathcal{O}_{Y, f(x)} \to \mathcal{O}_{X, x}$ 是同构
（见下文说明），
\end{enumerate}
则称该态射为{\it 双有理}态射。
\end{definition}

\noindent
说明：由于 $X$ 和 $Y$ 良好，拓扑空间 $|X|$ 与 $|Y|$ 是清醒的
（命题 \ref{decent-spaces-proposition-reasonable-sober}），所以条件 (1) 有意义。
此外，因为我们假设 $|X|$ 和 $|Y|$ 只有有限多个不可约分支，所以一般点
$x_1, \ldots, x_n \in |X|$ 以及相应的
$y_1, \ldots, y_n \in |Y|$ 分别包含在 $|X|$ 与 $|Y|$ 的任意稠密开集
中。特别地，由定理 \ref{decent-spaces-theorem-decent-open-dense-scheme}，它们分别
包含在 $X$ 与 $Y$ 的概形轨迹中。因此可将
$\mathcal{O}_{X, x_i}$ 以及相应的 $\mathcal{O}_{Y, y_i}$ 定义为该概形
分别在 $x_i$ 与 $y_i$ 处的局部环。

\medskip\noindent
由此可知，若态射 $f : X \to Y$ 双有理，则存在稠密开子空间
$X' \subset X$ 与 $Y' \subset Y$，使得
\begin{enumerate}
\item $f(X') \subset Y'$；
\item $X'$ 和 $Y'$ 可表示；
\item 按《态射》定义 \ref{morphisms-definition-birational} 的意义，
$f|_{X'} : X' \to Y'$ 双有理。
\end{enumerate}
不过，我们确实要求 $X$ 与 $Y$ 良好且只有有限多个不可约分支。
引理 \ref{decent-spaces-lemma-finitely-many-irreducible-components} 给出了对具有有限多个
不可约分支的良好代数空间的其他刻画。在多数情形下，双有理态射在稠密
开集上是同构。

\begin{lemma}
\label{decent-spaces-lemma-birational-dominant}
设 $S$ 为概形。设 $f : X \to Y$ 为 $S$ 上代数空间的态射，其中二者
良好且只有有限多个不可约分支。若 $f$ 双有理，则 $f$ 支配。
\end{lemma}

\begin{proof}
由定义立即得到。见《空间态射》定义
\ref{spaces-morphisms-definition-dominant}。
\end{proof}

\begin{lemma}
\label{decent-spaces-lemma-birational-generic-fibres}
设 $S$ 为概形。设 $f : X \to Y$ 为 $S$ 上代数空间的双有理态射，其中
二者良好且只有有限多个不可约分支。若
% P08-E333: frozen source places the topological point y in Y rather than |Y|.
$y \in Y$ 是某个不可约分支的一般点，则基变换
$X \times_Y \Spec(\mathcal{O}_{Y, y}) \to \Spec(\mathcal{O}_{Y, y})$
是同构。
\end{lemma}

\begin{proof}
令 $X' \subset X$ 和 $Y' \subset Y$ 为最大的可表示开子空间；见引理
\ref{decent-spaces-lemma-finitely-many-irreducible-components}。由引理
\ref{decent-spaces-lemma-quasi-finiteness-over-generic-point}，$f$ 在 $y$ 上方的纤维
% P08-E334: frozen source has the malformed “is consists”.
由 $X$ 的余维为 $0$ 的点组成，因而包含在 $X'$ 中。因此
$X \times_Y \Spec(\mathcal{O}_{Y, y}) =
X' \times_{Y'} \Spec(\mathcal{O}_{Y', y})$，结论遂由《态射》引理
\ref{morphisms-lemma-birational-generic-fibres} 得出。
\end{proof}

\begin{lemma}
\label{decent-spaces-lemma-birational-birational}
设 $S$ 为概形。设 $f : X \to Y$ 为 $S$ 上代数空间的双有理态射，其中
二者良好且只有有限多个不可约分支。假设下列条件之一成立：
\begin{enumerate}
\item $f$ 局部有限型，且 $Y$ 既约（即整）；
% P08-E335: frozen source incorrectly equates reduced with integral.
\item $f$ 局部有限表示。
\end{enumerate}
则存在稠密开子空间 $U \subset X$ 和 $V \subset Y$，使得
$f(U) \subset V$，且 $f|_U : U \to V$ 是同构。
\end{lemma}

\begin{proof}
由引理 \ref{decent-spaces-lemma-finitely-many-irreducible-components}，可假设 $X$ 和
$Y$ 都是概形。在此情形下，结论就是《态射》引理
\ref{morphisms-lemma-birational-birational}。
\end{proof}

\begin{lemma}
\label{decent-spaces-lemma-birational-isomorphism-over-dense-open}
设 $S$ 为概形。设 $f : X \to Y$ 为 $S$ 上代数空间的双有理态射，其中
二者良好且只有有限多个不可约分支。假设
\begin{enumerate}
\item $f$ 拟紧或 $f$ 分离；
\item $f$ 局部有限型且 $Y$ 既约，或者 $f$ 局部有限表示。
\end{enumerate}
则存在稠密开子空间 $V \subset Y$，使得 $f^{-1}(V) \to V$ 是同构。
\end{lemma}

\begin{proof}
由引理 \ref{decent-spaces-lemma-finitely-many-irreducible-components}，可假设 $Y$ 是概形。
由引理 \ref{decent-spaces-lemma-finite-over-dense-open}，可假设 $f$ 有限。于是 $X$
也是概形，结论由《态射》引理
\ref{morphisms-lemma-birational-isomorphism-over-dense-open} 得出。
\end{proof}

\begin{lemma}
\label{decent-spaces-lemma-birational-etale-localization}
设 $S$ 为概形。设 $f : X \to Y$ 为 $S$ 上代数空间的态射，其中二者
良好且只有有限多个不可约分支。若 $f$ 双有理，而 $V \to Y$
是以 $V$ 为仿射概形的 \'etale 态射，则 $X \times_Y V$ 良好且只有
有限多个不可约分支，并且 $X \times_Y V \to V$ 双有理。
\end{lemma}

\begin{proof}
代数空间 $U = X \times_Y V$ 是良好的（引理
\ref{decent-spaces-lemma-etale-named-properties}）。$V$ 和 $U$ 的一般点分别是
$|V|$ 和 $|U|$ 中位于 $|Y|$ 与 $|X|$ 的一般点上方的元素
（引理 \ref{decent-spaces-lemma-decent-generic-points}）。由于 $Y$ 良好，可知 $V$
上只有有限多个一般点。令 $\xi \in |X|$ 为某个不可约分支的一般点。
由定义 \ref{decent-spaces-definition-birational} 后的讨论，有 Cartesian 方块
$$
\xymatrix{
\Spec(\mathcal{O}_{X, \xi}) \ar[d] \ar[r] & X \ar[d] \\
\Spec(\mathcal{O}_{Y, f(\xi)}) \ar[r] & Y
}
$$
其中水平态射是识别局部环的单态射，而左侧竖直箭头是同构。因此在图
$$
\xymatrix{
\Spec(\mathcal{O}_{X, \xi}) \times_X U \ar[d] \ar[r] & U \ar[d] \\
\Spec(\mathcal{O}_{Y, f(\xi)}) \times_Y V \ar[r] & V
}
$$
左侧竖直箭头是同构。水平箭头的像分别包含在 $U$ 与 $V$ 的概形轨迹中，
并识别局部环（略去一些细节）。由于水平箭头的像分别是 $|U|$ 与 $|V|$
中位于 $\xi$ 与 $f(\xi)$ 上方的点，结论得证。
% P08-E336: frozen source uses singular “image” with plural “are”.
\end{proof}

\begin{lemma}
\label{decent-spaces-lemma-birational-induced-morphism-normalizations}
设 $S$ 为概形。设 $f : X \to Y$ 为 $S$ 上两个代数空间之间的双有理态射，
其中二者良好且只有有限多个不可约分支。则正规化
$X^\nu \to X$ 与 $Y^\nu \to Y$ 存在，并且有 $S$ 上代数空间的交换图
$$
\xymatrix{
X^\nu \ar[r] \ar[d] &  Y^\nu \ar[d] \\
X \ar[r] & Y
}
$$
。态射 $X^\nu \to Y^\nu$ 双有理。
\end{lemma}

\begin{proof}
由引理 \ref{decent-spaces-lemma-finitely-many-irreducible-components}，$X$ 和 $Y$ 满足
《空间态射》引理 \ref{spaces-morphisms-lemma-prepare-normalization} 中的
等价条件，因此正规化有定义。由《空间态射》引理
\ref{spaces-morphisms-lemma-normalization-normal}，代数空间 $X^\nu$
正规，并把余维为 $0$ 的点映到余维为 $0$ 的点。由于 $f$ 把余维为 $0$
的点映到余维为 $0$ 的点（由引理 \ref{decent-spaces-lemma-decent-generic-points}，
在良好空间上这些点就是一般点），再由《空间态射》引理
\ref{spaces-morphisms-lemma-normalization-normal}，复合
$X^\nu \to X \to Y$ 经过 $Y^\nu$ 分解。

\medskip\noindent
例如由引理 \ref{decent-spaces-lemma-representable-named-properties}，$X^\nu$ 与 $Y^\nu$
是良好的。此外，态射 $X^\nu \to X$ 和 $Y^\nu \to Y$ 在不可约分支上
诱导双射（见上述引用），故 $X^\nu$ 与 $Y^\nu$ 都只有有限多个不可约
分支，而态射 $X^\nu \to Y^\nu$ 在它们的一般点之间诱导双射。因此为
证明 $X^\nu \to Y^\nu$ 双有理，只需证明它在这些点处的局部环上诱导
同构。为此，可用 $X$ 和 $Y$ 的一般点的开邻域分别替换它们，故可假设
$X$ 和 $Y$ 是一般点分别为 $x$ 和 $y$ 的不可约仿射概形。由于 $f$
双有理，同态
% P08-E337: frozen source reverses the natural local-ring map from Definition birational.
$\mathcal{O}_{X, x} \to \mathcal{O}_{Y, y}$ 是同构。令
$x^\nu \in X^\nu$ 和 $y^\nu \in Y^\nu$ 分别为位于 $x$ 和 $y$ 上方的点。
由正规化的构造，有
$\mathcal{O}_{X^\nu, x^\nu} = \mathcal{O}_{X, x}/\mathfrak m_x$；
$Y$ 上也类似。因此同态
$\mathcal{O}_{X^\nu, x^\nu} \to \mathcal{O}_{Y^\nu, y^\nu}$
也是同构。
\end{proof}

\begin{lemma}
\label{decent-spaces-lemma-finite-birational-over-normal}
设 $S$ 为概形，$f : X \to Y$ 为 $S$ 上代数空间的态射。假设
\begin{enumerate}
\item $X$ 与 $Y$ 良好且只有有限多个不可约分支；
\item $f$ 整且双有理；
\item $Y$ 正规；
\item $X$ 既约。
\end{enumerate}
则 $f$ 是同构。
\end{lemma}

\begin{proof}
设 $V \to Y$ 为以 $V$ 为仿射概形的 \'etale 态射。只需证明
$U = X \times_Y V \to V$ 是同构。由引理
\ref{decent-spaces-lemma-birational-etale-localization} 及其证明，$U$ 与 $V$ 良好且
只有有限多个不可约分支，并且 $U \to V$ 双有理。由《性质》引理
\ref{properties-lemma-normal-locally-finite-nr-irreducibles}
，$V$ 是有限多个整概形的无交并。因此可假设 $V$ 整。由于 $f$ 双有理，
$U$ 不可约且既约，即为整的（注意，由于 $f$ 整、因而可表示，$U$ 是
概形）。这样，可假设 $X$ 与 $Y$ 是整概形；结论遂由概形情形得出，见
《态射》引理 \ref{morphisms-lemma-finite-birational-over-normal}。
\end{proof}

\begin{lemma}
\label{decent-spaces-lemma-normalization-normal}
设 $S$ 为概形。设 $f : X \to Y$ 为 $S$ 上良好代数空间的整双有理态射，
且这些空间只有有限多个不可约分支。则存在分解
$Y^\nu \to X \to Y$，并且 $Y^\nu \to X$ 是 $X$ 的正规化。
\end{lemma}

\begin{proof}
考虑引理 \ref{decent-spaces-lemma-birational-induced-morphism-normalizations} 中的态射
$X^\nu \to Y^\nu$。由《空间态射》引理
\ref{spaces-morphisms-lemma-finite-permanence}，此态射是整的。因此由引理
\ref{decent-spaces-lemma-finite-birational-over-normal}，它是同构。
\end{proof}




\section{Jacobson 空间}
\label{decent-spaces-section-jacobson}

\noindent
我们已在《空间性质》注
\ref{spaces-properties-remark-list-properties-local-etale-topology}
中定义了代数空间的 Jacobson 性质。对于可表示代数空间，它与《性质》第
\ref{properties-section-jacobson} 节中讨论的性质一致。对于一般代数空间
% P08-E338: frozen source appends a stray |X| after “general algebraic spaces”.
$|X|$，Jacobson 性质与拓扑空间 $|X|$ 的行为之间的关系并不显然。
不过，良好（例如拟分离或局部分离）代数空间 $X$ 是 Jacobson 的，当且
仅当 $|X|$ 是 Jacobson 的（见引理 \ref{decent-spaces-lemma-decent-Jacobson}）。

\begin{lemma}
\label{decent-spaces-lemma-Jacobson-universally-Jacobson}
设 $S$ 为概形，$X$ 为 $S$ 上的 Jacobson 代数空间。任何在 $X$
上局部有限型的代数空间都是 Jacobson 的。
\end{lemma}

\begin{proof}
设 $U \to X$ 为满射 \'etale 态射，其中 $U$ 是概形。于是按定义
$U$ 是 Jacobson 的；而对局部有限型的概形态射 $V \to U$，概形情形
的相应结果（《态射》引理
\ref{morphisms-lemma-Jacobson-universally-Jacobson}）表明 $V$
是 Jacobson 的。因此，若 $Y \to X$ 是代数空间的局部有限型态射，
令 $V = U \times_X Y$，则按定义可知 $Y$ 是 Jacobson 的。
\end{proof}

\begin{lemma}
\label{decent-spaces-lemma-Jacobson-ft-points-lift-to-closed}
设 $S$ 为概形，$X$ 为 $S$ 上的 Jacobson 代数空间。设
$x \in X_{\text{ft-pts}}$，且 $g : W \to X$ 是局部有限型态射，其中
$W$ 为概形。若 $x \in \Im(|g|)$，则 $W$ 中存在映到 $x$ 的闭点。
\end{lemma}

\begin{proof}
取 \'etale 态射 $U \to X$，其中 $U$ 为概形，且有映到 $x$ 的闭点
$u \in U$；见《空间态射》引理
\ref{spaces-morphisms-lemma-identify-finite-type-points}。由引理
\ref{decent-spaces-lemma-Jacobson-universally-Jacobson}，$W$、$W \times_X U$ 和
$U$ 都是 Jacobson 概形。因此由《态射》引理
\ref{morphisms-lemma-jacobson-finite-type-points}，这些概形上的有限型点
恰好就是闭点。$u$ 的逆像 $T \subset W \times_X U$ 是非空
% P08-E339: frozen English source omits “is” in “as x in the image”.
（因为 $x$ 位于 $W \to X$ 的像中）闭子集。由《态射》引理
\ref{morphisms-lemma-enough-finite-type-points}，$W \times_X U$
中存在映到 $u$ 的闭点 $t$。由于 $W \times_X U \to W$ 局部有限型，
由《态射》引理 \ref{morphisms-lemma-jacobson-finite-type-points}，
$t$ 在 $W$ 中的像是闭的。
\end{proof}

\begin{lemma}
\label{decent-spaces-lemma-decent-Jacobson-ft-pts}
设 $S$ 为概形，$X$ 为 $S$ 上的良好 Jacobson 代数空间。则
$X_{\text{ft-pts}} \subset |X|$ 是闭点集。
\end{lemma}

\begin{proof}
若 $x \in |X|$ 闭，则可用闭浸入 $\Spec(k) \to X$ 表示 $x$；见引理
\ref{decent-spaces-lemma-decent-space-closed-point}。因此 $x$ 当然是有限型点。

\medskip\noindent
反过来，设 $x \in |X|$ 为有限型点。我们知道，$x$ 可由拟紧单态射
$\Spec(k) \to X$ 表示，其中 $k$ 为域（定义
\ref{decent-spaces-definition-very-reasonable}）。另一方面，按定义，存在表示 $x$
的局部有限型态射 $\Spec(k') \to X$（《态射》定义
\ref{morphisms-definition-finite-type-point}）。由此得到分解
$\Spec(k') \to \Spec(k) \to X$。设 $U \to X$ 为任意 \'etale 态射，
其中 $U$ 仿射，并考虑态射
$$
\Spec(k') \times_X U \to \Spec(k) \times_X U \to U
$$
。拟紧概形 $\Spec(k) \times_X U$ 在 $\Spec(k)$ 上 \'etale，因而是
有限多个域之谱的无交并（注 \ref{decent-spaces-remark-recall}）。此外，第一个态射
满射且局部有限型（《态射》引理
\ref{morphisms-lemma-permanence-finite-type}），故在有限型点上满射
（《态射》引理
\ref{morphisms-lemma-finite-type-points-surjective-morphism}）；而由于
$U$ 是 Jacobson 的，复合态射（它局部有限型）把有限型点映到闭点
（《态射》引理 \ref{morphisms-lemma-jacobson-finite-type-points}）。
因此 $\Spec(k) \times_X U \to U$ 的像是有限闭点集，从而是闭的。
这对每个仿射 $U$ 和 \'etale 态射 $U \to X$ 都成立，故
$x \in |X|$ 是闭的。
\end{proof}

\begin{lemma}
\label{decent-spaces-lemma-decent-Jacobson}
设 $S$ 为概形，$X$ 为 $S$ 上的良好代数空间。则 $X$ 是 Jacobson 的，
当且仅当 $|X|$ 是 Jacobson 的。
\end{lemma}

\begin{proof}
假设 $X$ 是 Jacobson 的，且 $T \subset |X|$ 为闭子集。由《空间态射》
引理
\ref{spaces-morphisms-lemma-enough-finite-type-points}
，$T \cap X_{\text{ft-pts}}$ 在 $T$ 中稠密。由引理
\ref{decent-spaces-lemma-decent-Jacobson-ft-pts}，$X_{\text{ft-pts}}$ 是 $|X|$ 的
闭点集。因此 $|X|$ 的确是 Jacobson 的。

\medskip\noindent
假设 $|X|$ 是 Jacobson 的。设 $f : U \to X$ 为 \'etale 态射，其中
$U$ 是仿射概形。需要证明 $U$ 是 Jacobson 的。若 $x \in |X|$ 闭，
则纤维 $F = f^{-1}(\{x\})$ 是 $U$ 的有限（由良好的定义）闭子集
（由 $|X|$ 上拓扑的构造）。由于 $F$ 的点之间不存在特化
（引理 \ref{decent-spaces-lemma-decent-no-specializations-map-to-same-point}），
$F$ 的每个点在 $U$ 中都闭。若 $U$ 不是 Jacobson 的，则存在非闭点
$u \in U$，使得 $\{u\}$ 局部闭（《拓扑》引理
\ref{topology-lemma-non-jacobson-Noetherian-characterize}）。我们将证明
$f(u) \in |X|$ 闭；由上述结论，$u$ 随即在 $U$ 中闭，得到矛盾并完成证明。
为证明这一点，可用 $u$ 的一个仿射开邻域替换 $U$。因此可假设
$\{u\}$ 在 $U$ 中闭。令 $R = U \times_X U$，其投影为
$s, t : R \to U$。于是 $s^{-1}(\{u\}) = \{r_1, \ldots, r_m\}$ 有限
（由良好空间的定义）。用 $u$ 的一个更小仿射开邻域替换 $U$ 后，
可假设对 $j = 1, \ldots, m$ 都有 $t(r_j) = u$。因此 $\{u\}$
是 $U$ 的 $R$-不变闭子集。于是 $\{f(u)\}$ 是 $X$ 的局部闭子集，
因为它在 $|X|$ 的开集 $|f|(|U|)$ 中闭。由于 $|X|$ 是 Jacobson 的，
可知 $f(u)$ 在 $|X|$ 中闭，这正是所需。
\end{proof}

\begin{lemma}
\label{decent-spaces-lemma-punctured-spec}
设 $S$ 为概形，$X$ 为 $S$ 上良好的局部 Noether 代数空间。设
$x \in |X|$。则
$$
W = \{x' \in |X| : x' \leadsto x,\ x' \not = x\}
$$
是 Noether、谱、清醒且 Jacobson 的拓扑空间。
\end{lemma}

\begin{proof}
可以用任意包含 $x$ 的开子空间替换。因此可假设 $X$ 拟紧。于是
$|X|$ 是 Noether 拓扑空间（《空间性质》引理
\ref{spaces-properties-lemma-Noetherian-topology}）。因此 $W$ 是
Noether 拓扑空间（《拓扑》引理 \ref{topology-lemma-Noetherian}）。

\medskip\noindent
结合引理 \ref{decent-spaces-lemma-locally-Noetherian-decent-quasi-separated} 与
《空间性质》引理
\ref{spaces-properties-lemma-quasi-compact-quasi-separated-spectral}，
可知 $|X|$ 是谱拓扑空间。由《拓扑》引理
\ref{topology-lemma-make-spectral-space}，$W \cup \{x\}$ 是谱拓扑空间。
现在 $W$ 是 $W \cup \{x\}$ 的拟紧开集，故由《拓扑》引理
\ref{topology-lemma-spectral-sub}，$W$ 是谱空间。

\medskip\noindent
设 $E \subset W$ 为不可约闭子集。若 $Z \subset |X|$ 是 $E$ 的闭包，
则 $x \in Z$。由命题 \ref{decent-spaces-proposition-reasonable-sober}，存在唯一
一般点 $\eta \in Z$。当然 $\eta \in W$，因而 $\eta \in E$。所以
$E$ 有唯一一般点，即 $W$ 是清醒的。

\medskip\noindent
设 $x' \in W$，且 $\{x'\}$ 在 $W$ 中局部闭。为完成证明，需要证明
$x'$ 是 $W$ 的闭点。否则，$W$ 中存在非平凡特化
$x' \leadsto x'_1$。设 $U$ 为仿射概形，$u \in U$ 为一点，并设
$U \to X$ 为把 $u$ 映到 $x$ 的 \'etale 态射。由引理
\ref{decent-spaces-lemma-decent-specialization}，可选取映到
$x' \leadsto x'_1 \leadsto x$ 的特化
$u' \leadsto u'_1 \leadsto u$。令
$\mathfrak p' \subset \mathcal{O}_{U, u}$ 为与 $u'$ 对应的素理想。
这些特化的存在表明
$\dim(\mathcal{O}_{U, u}/\mathfrak p') \geq 2$。因此由《代数》引理
\ref{algebra-lemma-Noetherian-local-domain-dim-2-infinite-opens}，
$\Spec(\mathcal{O}_{U, u}/\mathfrak p')$ 的每个非空开集都无限。
由引理 \ref{decent-spaces-lemma-decent-no-specializations-map-to-same-point}，得到连续映射
$$
\Spec(\mathcal{O}_{U, u}/\mathfrak p')
\setminus \{\mathfrak m_u/\mathfrak p'\}
\longrightarrow
W
$$
由于左端的一般点映到 $x'$，该映射的像包含在
$\overline{\{x'\}}$ 中。因此，所示箭头下 $\{x'\}$ 的逆像是非空开集，
从而无限。然而，由于 $X$ 良好，$U \to X$ 的纤维有限，由此得出
% P08-E340: frozen source says the singleton {x'} itself is infinite.
$\{x'\}$ 是无限的。此矛盾完成证明。
\end{proof}






\section{局部不可约性}
\label{decent-spaces-section-irreducible-local-ring}

\noindent
我们已在《空间性质》第
\ref{spaces-properties-section-irreducible-local-ring} 节中定义了代数空间
在一点处的几何分支数。代数空间在一点处的分支数只能对良好代数空间定义。

\begin{lemma}
\label{decent-spaces-lemma-irreducible-local-ring}
设 $S$ 为概形，$X$ 为 $S$ 上的良好代数空间，并设 $x \in |X|$ 为一点。
下列条件等价：
\begin{enumerate}
\item 对任意初等 \'etale 邻域 $(U, u) \to (X, x)$，局部环
$\mathcal{O}_{U, u}$ 有唯一极小素理想；
\item 对任意初等 \'etale 邻域 $(U, u) \to (X, x)$，通过 $u$ 的
$U$ 的不可约分支唯一；
\item 对任意初等 \'etale 邻域 $(U, u) \to (X, x)$，局部环
$\mathcal{O}_{U, u}$ 单枝；
\item Hensel 局部环 $\mathcal{O}_{X, x}^h$ 有唯一极小素理想。
\end{enumerate}
\end{lemma}

\begin{proof}
(1) 与 (2) 等价，是因为通过 $u$ 的 $U$ 的不可约分支与 $U$ 在 $u$
处的局部环之极小素理想成 $1$-$1$ 对应。环 $\mathcal{O}_{X, x}^h$ 是
$\mathcal{O}_{U, u}$ 的 Hensel 化；见定义
\ref{decent-spaces-definition-henselian-local-ring} 后的讨论。特别地，由《代数进阶》
引理 \ref{more-algebra-lemma-unibranch}，(3) 与 (4) 等价。(2) 与 (3)
的等价性来自《态射进阶》引理
\ref{more-morphisms-lemma-nr-branches}。
\end{proof}

\begin{definition}
\label{decent-spaces-definition-unibranch}
设 $S$ 为概形，$X$ 为 $S$ 上的良好代数空间，并设 $x \in |X|$。
若引理 \ref{decent-spaces-lemma-irreducible-local-ring} 的等价条件成立，
则称 $X$ {\it 在 $x$ 处单枝}。若 $X$ 在每个 $x \in |X|$ 处都单枝，
则称 $X$ {\it 单枝}。
\end{definition}

\noindent
这与概形的定义相符（《性质》定义
\ref{properties-definition-unibranch}）。

\begin{lemma}
\label{decent-spaces-lemma-nr-branches-local-ring}
设 $S$ 为概形，$X$ 为 $S$ 上的良好代数空间，并设 $x \in |X|$ 为一点。
设 $n \in \{1, 2, \ldots\}$ 为整数。下列条件等价：
\begin{enumerate}
\item 对任意初等 \'etale 邻域 $(U, u) \to (X, x)$，局部环
$\mathcal{O}_{U, u}$ 的极小素理想数为 $\leq n$，且至少对一个
$(U, u)$ 的选取，此数为 $n$；
\item 对任意初等 \'etale 邻域 $(U, u) \to (X, x)$，通过 $u$ 的 $U$
之不可约分支数为 $\leq n$，且至少对一个 $(U, u)$ 的选取，此数为 $n$；
% P08-E341: frozen source omits “of” after “the number”.
\item 对任意初等 \'etale 邻域 $(U, u) \to (X, x)$，$U$ 在 $u$ 处的
分支数为 $\leq n$，且至少对一个 $(U, u)$ 的选取，此数为 $n$；
\item $\mathcal{O}_{X, x}^h$ 的极小素理想数为 $n$。
\end{enumerate}
\end{lemma}

\begin{proof}
(1) 与 (2) 等价，是因为通过 $u$ 的 $U$ 的不可约分支与 $U$ 在 $u$
处的局部环之极小素理想成 $1$-$1$ 对应。环
% P08-E342: frozen source omits the superscript h on the henselization.
$\mathcal{O}_{X, x}$ 是 $\mathcal{O}_{U, u}$ 的 Hensel 化；见定义
\ref{decent-spaces-definition-henselian-local-ring} 后的讨论。特别地，由《代数进阶》
引理 \ref{more-algebra-lemma-unibranch}，(3) 与 (4) 等价。(2) 与 (3)
的等价性来自《态射进阶》引理
\ref{more-morphisms-lemma-nr-branches}。
\end{proof}

\begin{definition}
\label{decent-spaces-definition-number-of-branches}
设 $S$ 为概形，$X$ 为 $S$ 上的良好代数空间，并设 $x \in |X|$。
若引理 \ref{decent-spaces-lemma-nr-branches-local-ring} 的等价条件成立，则
{\it $X$ 在 $x$ 处的分支数}为 $n \in \mathbf{N}$；否则为 $\infty$。
\end{definition}






\section{链状代数空间}
\label{decent-spaces-section-catenary}

\noindent
本节把《性质》第 \ref{properties-section-catenary} 节和《态射》第
\ref{morphisms-section-universally-catenary} 节中的内容推广到代数空间。

\begin{definition}
\label{decent-spaces-definition-catenary}
设 $S$ 为概形，$X$ 为 $S$ 上的良好代数空间。若 $|X|$ 链状
（《拓扑》定义 \ref{topology-definition-catenary}），则称 $X$
为{\it 链状}代数空间。
\end{definition}

\noindent
若 $X$ 可表示，则这等价于表示 $X$ 的概形之相应概念。

\begin{lemma}
\label{decent-spaces-lemma-scheme-with-dimension-function}
设 $S$ 为局部 Noether 且普遍链状的概形。设
$\delta : S \to \mathbf{Z}$ 为维数函数。设 $X$ 为 $S$ 上的良好代数空间，
并且结构态射 $X \to S$ 局部有限型。设
$\delta_X : |X| \to \mathbf{Z}$ 为把 $x$ 映到 $\delta(f(x))$ 加上
$x/f(x)$ 的超越次数所得的映射。则 $\delta_X$ 是 $|X|$ 上的维数函数。
% P08-E343: frozen source uses f(x) without naming the structure morphism f.
\end{lemma}

\begin{proof}
设 $\varphi : U \to X$ 为满射 \'etale 态射，其中 $U$ 为概形。
由《态射》引理 \ref{morphisms-lemma-dimension-function-propagates}，
以类似方式定义的函数 $\delta_U$ 是 $U$ 上的维数函数。另一方面，
由《空间态射》定义
\ref{spaces-morphisms-definition-dimension-fibre} 中相对超越次数的定义，
有 $\delta_U(u) = \delta_X(\varphi(u))$。

\medskip\noindent
设 $x \leadsto x'$ 为 $|X|$ 中点的特化。由引理
\ref{decent-spaces-lemma-decent-specialization}，可以找到 $U$ 中点的特化
$u \leadsto u'$，满足 $\varphi(u) = x$ 且 $\varphi(u') = x'$。
% P08-E344: frozen source begins this sentence with lowercase “by”.
此外，$x = x'$ 当且仅当 $u = u'$；见引理
\ref{decent-spaces-lemma-decent-no-specializations-map-to-same-point}。因此
$\delta_U$ 是维数函数这一事实表明 $\delta_X$ 也是维数函数；见《拓扑》
定义 \ref{topology-definition-dimension-function}。
\end{proof}

\begin{lemma}
\label{decent-spaces-lemma-universally-catenary-scheme}
设 $S$ 为局部 Noether 且普遍链状的概形。设 $X$ 为 $S$ 上的代数空间，
且 $X$ 良好，结构态射 $X \to S$ 局部有限型。则 $X$ 链状。
\end{lemma}

\begin{proof}
此问题在 $S$ 上是局部的（使用《拓扑》引理
\ref{topology-lemma-catenary}）。因此可假设 $S$ 有维数函数；见《拓扑》
引理 \ref{topology-lemma-locally-dimension-function}。于是由引理
\ref{decent-spaces-lemma-scheme-with-dimension-function}，$|X|$ 有维数函数。
由于 $|X|$ 是清醒的（命题 \ref{decent-spaces-proposition-reasonable-sober}），
由《拓扑》引理 \ref{topology-lemma-dimension-function-catenary}，
$|X|$ 链状。
\end{proof}

\noindent
由引理 \ref{decent-spaces-lemma-universally-catenary-scheme}，下面的定义与可表示代数空间
已有的概念相容。

\begin{definition}
\label{decent-spaces-definition-universally-catenary}
设 $S$ 为概形，$X$ 为 $S$ 上良好且局部 Noether 的代数空间。
若对每个局部有限型且 $Y$ 良好的代数空间态射 $Y \to X$，代数空间
$Y$ 都链状，则称 $X$ 是{\it 普遍链状}的。
\end{definition}

\noindent
若 $X$ 是代数空间，则条件
``$X$ 良好且局部 Noether'' 等价于
``$X$ 拟分离且局部 Noether''。这就是引理
\ref{decent-spaces-lemma-locally-Noetherian-decent-quasi-separated}。因此也可把上述
定义理解为：$X$ 普遍链状，当且仅当对于所有拟分离且局部有限型的态射
$Y \to X$，$Y$ 都链状。

\begin{lemma}
\label{decent-spaces-lemma-universally-catenary}
设 $S$ 为概形，$X$ 为 $S$ 上良好、局部 Noether 且普遍链状的代数空间。
则任何在 $X$ 上局部有限型的良好代数空间都是普遍链状的。
\end{lemma}

\begin{proof}
这由定义以及下述事实形式地得出：局部有限型态射的复合仍局部有限型
（《空间态射》引理
\ref{spaces-morphisms-lemma-composition-finite-type}）。
\end{proof}

\begin{lemma}
\label{decent-spaces-lemma-check-dimension-function-finite-cover}
设 $S$ 为概形。设 $f : Y \to X$ 为良好且局部 Noether 的代数空间之间
的满射有限态射。设 $\delta : |X| \to \mathbf{Z}$ 为函数。若
$\delta \circ |f|$ 是维数函数，则 $\delta$ 是维数函数。
\end{lemma}

\begin{proof}
设
% P08-E345: frozen source uses the mapsto arrow rather than the specialization arrow.
$x \mapsto x'$，$x \not = x'$，为 $|X|$ 中的特化。选取
$y \in |Y|$ 使得 $|f|(y) = x$。由于 $|f|$ 是闭映射（《空间态射》引理
\ref{spaces-morphisms-lemma-finite-proper}），可找到特化
$y \leadsto y'$，使得 $|f|(y') = x'$。因此
$\delta(x) = \delta(|f|(y)) > \delta(|f|(y')) = \delta(x')$
（见《拓扑》定义 \ref{topology-definition-dimension-function}）。
若 $x \leadsto x'$ 是直接特化，则 $y \leadsto y'$ 也是直接特化：
事实上，若 $y \leadsto y'' \leadsto y'$，则 $|f|(y'')$ 必为 $x$ 或
$x'$；而由引理 \ref{decent-spaces-lemma-conditions-on-fibre-and-qf}，$|f|$ 的纤维
中的点之间不存在非平凡特化。
\end{proof}

\noindent
有关讨论将在《空间态射进阶》第
\ref{spaces-more-morphisms-section-catenary} 节中继续。
% Generated mechanically from the frozen spaces-cohomology.tex target.
% Do not edit this generated file; edit the bound locale source instead.

% OK, start here.
%

\setcounter{chapter}{68}
\chapter{代数空间的上同调}


\phantomsection
\label{spaces-cohomology-section-phantom}





\section{引言}
\label{spaces-cohomology-section-introduction}

\noindent
本章讨论代数空间的上同调。虽然我们会证明一些关于阿贝尔层上同调的结果，
但重点是拟凝聚层的上同调；换言之，我们将证明“概形的上同调”一章中
若干结果的类似命题。本章的一些结果可见 \cite{Kn}。

\medskip\noindent
本章缺少的一项重要工具是{\it 归纳原理}，也就是“概形的上同调”中
引理 \ref{coherent-lemma-induction-principle} 对拟紧拟分离代数空间的类似形式。
“空间的导出范畴”第 \ref{spaces-perfect-section-induction} 节对它作了精确表述并给出了详细证明。
本章不用归纳原理，而采用交错 {\v C}ech 复形，见
第 \ref{spaces-cohomology-section-alternating-cech} 节。这个复形旨在证明
命题 \ref{spaces-cohomology-proposition-vanishing} 一类的消失性命题；不过在某些情形下，
归纳原理更为有力，而且或许也更为“标准”。我们建议读者在读完本节的
部分内容之后，再去了解这一归纳原理。



\section{约定}
\label{spaces-cohomology-section-conventions}

\noindent
我们始终假设所有概形都包含在一个大 fppf 位点 $\Sch_{fppf}$ 中，
并且所考虑的每个环 $A$ 都满足：$\Spec(A)$（同构于）这个大位点的一个对象。

\medskip\noindent
设 $S$ 为概形，$X$ 为 $S$ 上的代数空间。在本章及后续各章中，
我们把 $X$ 与自身的积（取于 $S$ 上代数空间的范畴中）记作
$X \times_S X$，而不记作 $X \times X$。








\section{高阶正像}
\label{spaces-cohomology-section-higher-direct-image}

\noindent
设 $S$ 为概形，$X$ 为 $S$ 上的可表代数空间，并设 $\mathcal{F}$ 为 $X$ 上的
拟凝聚模（见“代数空间的性质”第
\ref{spaces-properties-section-quasi-coherent} 节）。由“下降”命题
\ref{descent-proposition-same-cohomology-quasi-coherent}，上同调群
$H^i(X, \mathcal{F})$ 与通常的上同调群一致；后者是在表示 $X$ 的概形上，
对相应拟凝聚模按 Zariski 拓扑计算得到的。

\medskip\noindent
更一般地，设 $f : X \to Y$ 是可表代数空间 $X$ 与 $Y$ 之间的拟紧拟分离态射，
$\mathcal{F}$ 是 $X$ 上的拟凝聚模。由“下降”引理
\ref{descent-lemma-higher-direct-images-small-etale}，层
$R^if_*\mathcal{F}$ 与通常的高阶正像一致；后者是在表示 $X$ 的概形到表示
$Y$ 的概形的态射上，对拟凝聚模按 Zariski 拓扑计算得到的。

\medskip\noindent
再一般些，设 $f : X \to Y$ 是 $S$ 上代数空间之间的可表、拟紧且拟分离的态射。
设 $V$ 为概形，$V \to Y$ 为平展满态射。令 $U = V \times_Y X$，并以
$f' : U \to V$ 表示 $f$ 的基变换。于是对任意拟凝聚
$\mathcal{O}_X$-模 $\mathcal{F}$，有
\begin{equation}
\label{spaces-cohomology-equation-representable-higher-direct-image}
R^if'_*(\mathcal{F}|_U) = (R^if_*\mathcal{F})|_V,
\end{equation}
见“代数空间的性质”引理
\ref{spaces-properties-lemma-pushforward-etale-base-change-modules}。
又因为 $f' : U \to V$ 是概形之间的拟紧拟分离态射，根据上一段的说明，
我们可以把 $\mathcal{F}|_U$ 看成概形 $U$ 上的拟凝聚层，把 $f'$ 看成概形之间
的态射，从而计算 $R^if'_*(\mathcal{F}|_U)$。下文将反复使用这一事实而不再说明。

\medskip\noindent
下面证明：对于代数空间的任意拟紧拟分离态射，拟凝聚层的高阶正像仍为拟凝聚层。
证明中会用到一个技巧；另一种“更好”的证明会使用相对 {\v C}ech 复形，
参见“代数栈上的层”第 \ref{stacks-sheaves-section-cech} 节、
第 \ref{stacks-sheaves-section-sheaf-cech-complex} 节及其后各节。

\begin{lemma}
\label{spaces-cohomology-lemma-higher-direct-image}
设 $S$ 为概形，$f : X \to Y$ 为 $S$ 上代数空间的态射。若 $f$ 拟紧且拟分离，
则 $R^if_*$ 把拟凝聚 $\mathcal{O}_X$-模变为拟凝聚 $\mathcal{O}_Y$-模。
\end{lemma}

\begin{proof}
取平展态射 $V \to Y$，其中 $V$ 为仿射概形。令 $U = V \times_Y X$，并以
$f' : U \to V$ 表示诱导态射。设 $\mathcal{F}$ 为拟凝聚
$\mathcal{O}_X$-模。由“代数空间的性质”引理
\ref{spaces-properties-lemma-pushforward-etale-base-change-modules}，有
$R^if'_*(\mathcal{F}|_U) = (R^if_*\mathcal{F})|_V$.
由于“是拟凝聚模”这一性质对于 $Y$ 上的平展拓扑是局部的（见“代数空间的性质”
引理 \ref{spaces-properties-lemma-characterize-quasi-coherent}），
我们可以用 $V$ 代替 $Y$；也就是说，可以假设 $Y$ 是仿射概形。

\medskip\noindent
假设 $Y$ 仿射。由于 $f$ 拟紧，$X$ 也是拟紧的。因此可以选取一个仿射概形
$U$ 以及一个平展满态射 $g : U \to X$，见“代数空间的性质”引理
\ref{spaces-properties-lemma-quasi-compact-affine-cover}。图示如下：
$$
\xymatrix{
U \ar[r]_g \ar[rd]_{f \circ g} & X \ar[d]^f \\
& Y
}
$$
因为 $X$ 拟分离，态射 $g : U \to X$ 可表、分离且拟紧。因此本引理对 $g$
成立（由引理前的讨论）。它对 $f \circ g : U \to Y$ 也成立，因为这是仿射概形
之间的态射。

\medskip\noindent
在上一段所述情形中，我们将对 $n$ 作归纳，证明命题 $IH_n$：对 $X$ 上任意
拟凝聚层 $\mathcal{F}$，当 $i \leq n$ 时，层 $R^if\mathcal{F}$ 拟凝聚。
$n = 0$ 的情形由“代数空间的态射”引理
\ref{spaces-morphisms-lemma-pushforward} 得到。假设 $IH_n$ 成立；以下证明
$IH_{n + 1}$。

\medskip\noindent
设 $\mathcal{H}$ 为拟凝聚 $\mathcal{O}_U$-模。考虑 Leray 谱序列
$$
E_2^{p, q} = R^pf_* R^qg_* \mathcal{H}
\Rightarrow
R^{p + q}(f \circ g)_*\mathcal{H}
$$
见“位点上的上同调”引理 \ref{sites-cohomology-lemma-relative-Leray}。
由 $IH_n$，$R^qg_*\mathcal{H}$ 拟凝聚，因而当 $p \leq n$ 时，所有层
$R^pf_*R^qg_*\mathcal{H}$ 都拟凝聚。各层
$R^{p + q}(f \circ g)_*\mathcal{H}$ 也都拟凝聚（事实上，当
$p + q > 0$ 时它们为零，但这里不需要这一点）。在总次数不超过
$n + 1$ 的项中，唯一尚不知是否拟凝聚的模是
$E_2^{n + 1, 0} = R^{n + 1}f_*g_*\mathcal{H}$。此外，微分
$d_r^{n + 1, 0} : E_r^{n + 1, 0} \to E_r^{n + 1 + r, 1 - r}$
为零，因为其目标为零。再利用 $\QCoh(\mathcal{O}_X)$ 是
$\textit{Mod}(\mathcal{O}_X)$ 的弱 Serre 子范畴（“代数空间的性质”引理
\ref{spaces-properties-lemma-properties-quasi-coherent}），可知
$R^{n + 1}f_*g_*\mathcal{H}$ 拟凝聚（细节从略）。

\medskip\noindent
设 $\mathcal{F}$ 为拟凝聚 $\mathcal{O}_X$-模，并令
$\mathcal{H} = g^*\mathcal{F}$。由于 $U \to X$ 是平展满态射，伴随映射
$\mathcal{F} \to g_*g^*\mathcal{F} = g_*\mathcal{H}$ 为单射。考虑正合列
$$
0 \to \mathcal{F} \to g_*\mathcal{H} \to \mathcal{G} \to 0
$$
其中 $\mathcal{G}$ 是第一个映射的余核，因而是拟凝聚的。应用上同调长正合列，得到
$$
R^nf_*g_*\mathcal{H} \to
R^nf_*\mathcal{G} \to
R^{n + 1}f_*\mathcal{F} \to
R^{n + 1}f_*g_*\mathcal{H} \to
R^{n + 1}f_*\mathcal{G}
$$
第一个箭头的余核拟凝聚，而上面已经证明
$R^{n + 1}f_*g_*\mathcal{H}$ 拟凝聚。因此 $R^{n + 1}f_*\mathcal{F}$
有一个两步滤过，其中第一步拟凝聚，第二步是某个拟凝聚层的子模。
由于 $\mathcal{F}$ 是任意拟凝聚 $\mathcal{O}_X$-模，这个结论也适用于
$\mathcal{G}$。所以可以选取正合列
$0 \to \mathcal{A} \to R^{n + 1}f_*\mathcal{G} \to \mathcal{B}$
其中 $\mathcal{A}$、$\mathcal{B}$ 都是拟凝聚 $\mathcal{O}_Y$-模。
于是复合映射
$R^{n + 1}f_*g_*\mathcal{H} \to R^{n + 1}f_*\mathcal{G}
\to \mathcal{B}$ 的核 $\mathcal{K}$ 拟凝聚，并得到映射
$\mathcal{K} \to \mathcal{A}$；其核 $\mathcal{K}'$ 也拟凝聚。因此
$R^{n + 1}f_*\mathcal{F}$ 位于正合列
$$
R^nf_*g_*\mathcal{H} \to
R^nf_*\mathcal{G} \to
R^{n + 1}f_*\mathcal{F} \to \mathcal{K}' \to 0
$$
中；其中除 $R^{n + 1}f_*\mathcal{F}$ 外，其余各模均拟凝聚。由此得出
$R^{n + 1}f_*\mathcal{F}$ 拟凝聚，即 $IH_{n + 1}$ 如期成立。
\end{proof}

\begin{lemma}
\label{spaces-cohomology-lemma-quasi-coherence-higher-direct-images-application}
设 $S$ 为概形，$f : X \to Y$ 为 $S$ 上代数空间的拟分离拟紧态射。
对任意拟凝聚 $\mathcal{O}_X$-模 $\mathcal{F}$ 以及 $Y_\etale$ 的任意仿射对象
$V$，有
$$
H^q(V \times_Y X, \mathcal{F}) = H^0(V, R^qf_*\mathcal{F})
$$
其中 $q \in \mathbf{Z}$ 任意。
\end{lemma}

\begin{proof}
由于 $Rf_*$ 的形成与平展局部化可交换（“代数空间的性质”引理
\ref{spaces-properties-lemma-pushforward-etale-base-change-modules}），
我们可以用 $V$ 代替 $Y$，并假设 $Y = V$ 仿射。考虑 Leray 谱序列
$E_2^{p, q} = H^p(Y, R^qf_*\mathcal{F})$
它收敛到 $H^{p + q}(X, \mathcal{F})$，见“位点上的上同调”引理
\ref{sites-cohomology-lemma-Leray}。由引理
\ref{spaces-cohomology-lemma-higher-direct-image}，各层 $R^qf_*\mathcal{F}$ 均拟凝聚。
再由“概形的上同调”引理
\ref{coherent-lemma-quasi-coherent-affine-cohomology-zero}，当 $p > 0$ 时
$E_2^{p, q} = 0$。所以该谱序列在 $E_2$ 页退化，结论成立。
\end{proof}




\section{有限态射}
\label{spaces-cohomology-section-finite-morphisms}

\noindent
下面给出若干对所有阿贝尔层都成立的结果（特别地，也适用于拟凝聚模）。
我们要{\bf 提醒}读者：对于概形的有限态射和 Zariski 拓扑，以下引理并不成立。

\begin{lemma}
\label{spaces-cohomology-lemma-finite-higher-direct-image-zero}
设 $S$ 为概形，$f : X \to Y$ 为代数空间的整态射（例如有限态射）。则
$f_* : \textit{Ab}(X_\etale) \to \textit{Ab}(Y_\etale)$
是正合函子，并且当 $p > 0$ 时 $R^pf_* = 0$。
\end{lemma}

\begin{proof}
由“代数空间的性质”引理
\ref{spaces-properties-lemma-pushforward-etale-base-change}，
可以在 $Y$ 的一个平展覆盖上计算高阶正像。因此可以假设 $Y$ 是概形，
进而 $X$ 也是概形（见“代数空间的态射”引理
\ref{spaces-morphisms-lemma-integral-local}）。此时可应用“平展上同调”引理
\ref{etale-cohomology-lemma-what-integral}。对于有限情形，读者也可以参阅
技术性较弱的“平展上同调”命题
\ref{etale-cohomology-proposition-finite-higher-direct-image-zero}。
\end{proof}

\begin{lemma}
\label{spaces-cohomology-lemma-stalk-push-finite}
设 $S$ 为概形，$f : X \to Y$ 为 $S$ 上代数空间的有限态射。设
$\overline{y}$ 是 $Y$ 的一个几何点，并且它在 $X$ 中的提升为
$\overline{x}_1, \ldots, \overline{x}_n$。则
$$
(f_*\mathcal{F})_{\overline{y}} =
\prod\nolimits_{i = 1, \ldots, n}
\mathcal{F}_{\overline{x}_i}
$$
对 $X_\etale$ 上任意层 $\mathcal{F}$ 都成立。
\end{lemma}

\begin{proof}
选取 $\overline{y}$ 的一个平展邻域 $(V, \overline{v})$。于是茎
$(f_*\mathcal{F})_{\overline{y}}$ 就是 $f_*\mathcal{F}|_V$ 在
$\overline{v}$ 处的茎。由“代数空间的性质”引理
\ref{spaces-properties-lemma-pushforward-etale-base-change}，
可以用 $V$ 代替 $Y$，并用 $X \times_Y V$ 代替 $X$。
% P09-ERR-0001: frozen source literally has Z -> X here; see corrections.md.
于是 $Z \to X$ 是概形之间的有限态射，而所需结论就是“平展上同调”命题
\ref{etale-cohomology-proposition-finite-higher-direct-image-zero}。
\end{proof}

\begin{lemma}
\label{spaces-cohomology-lemma-finite-rings}
设 $S$ 为概形，$\pi : X \to Y$ 为 $S$ 上代数空间的有限态射。设
$\mathcal{A}$ 是 $X_\etale$ 上的环层，$\mathcal{B}$ 是 $Y_\etale$ 上的环层，
并设 $\varphi : \mathcal{B} \to \pi_*\mathcal{A}$ 为环层同态，从而得到
带环拓扑斯的态射
$$
f = (\pi, \varphi) :
(\Sh(X_\etale), \mathcal{A})
\longrightarrow
(\Sh(Y_\etale), \mathcal{B}).
$$
若 $\mathcal{F}$ 是 $\mathcal{A}$-模层，$\mathcal{G}$ 是 $\mathcal{B}$-模层，
则典范映射
$$
\mathcal{G} \otimes_\mathcal{B} f_*\mathcal{F}
\longrightarrow
f_*(f^*\mathcal{G} \otimes_\mathcal{A} \mathcal{F}).
$$
是同构。
\end{lemma}

\begin{proof}
这个映射是下列映射的伴随映射：
$$
f^*\mathcal{G} \otimes_\mathcal{A}
f^* f_*\mathcal{F} =
f^*(\mathcal{G} \otimes_\mathcal{B} f_*\mathcal{F})
\longrightarrow
f^*\mathcal{G} \otimes_\mathcal{A} \mathcal{F}
$$
它由 $\text{id} : f^*\mathcal{G} \to f^*\mathcal{G}$ 和伴随映射
$f^* f_*\mathcal{F} \to \mathcal{F}$ 给出。要证明它是同构，只需在茎上检验
（见“代数空间的性质”定理
\ref{spaces-properties-theorem-exactness-stalks}）。设 $\overline{y}$ 为 $Y$
的几何点，$\overline{x}_1, \ldots, \overline{x}_n$ 为位于
$\overline{y}$ 上方的 $X$ 的几何点。把这些映射在茎上的作用写出后，只需证明
$$
\mathcal{G}_{\overline{y}}
\otimes_{\mathcal{B}_{\overline{y}}}
\left(
\bigoplus\nolimits_{i = 1, \ldots, n} \mathcal{F}_{\overline{x}_i}
\right) =
\bigoplus\nolimits_{i = 1, \ldots, n}
% P09-ERR-0002: frozen source literally has B_{xbar}; see corrections.md.
(\mathcal{G}_{\overline{y}}
\otimes_{\mathcal{B}_{\overline{x}}}
\mathcal{A}_{\overline{x}_i}) \otimes_{\mathcal{A}_{\overline{x}_i}}
\mathcal{F}_{\overline{x}_i}
$$
确实成立。这里使用了张量积与取茎可交换，以及“代数空间的性质”引理
\ref{spaces-properties-lemma-stalk-pullback} 所述拉回下茎的行为；还使用了
% P09-ERR-0003: frozen source says closed immersion; see corrections.md.
引理 \ref{spaces-cohomology-lemma-stalk-push-finite} 所述沿闭浸入作正像时茎的行为。
\end{proof}

\noindent
最后给出一个对有限态射而言极其一般的投影公式。

\begin{lemma}
\label{spaces-cohomology-lemma-projection-formula-finite}
沿用引理 \ref{spaces-cohomology-lemma-finite-rings} 中的
$S$、$X$、$Y$、$\pi$、$\mathcal{A}$、$\mathcal{B}$、$\varphi$ 和 $f$，则
$$
K \otimes_\mathcal{B}^\mathbf{L} Rf_*M =
Rf_*(Lf^*K \otimes_\mathcal{A}^\mathbf{L} M)
$$
对任意 $K \in D(\mathcal{B})$ 和 $M \in D(\mathcal{A})$，上述等式在
$D(\mathcal{B})$ 中成立。
\end{lemma}

\begin{proof}
由于 $f_*$ 正合（引理 \ref{spaces-cohomology-lemma-finite-higher-direct-image-zero}），
对任意代表复形逐项应用 $f_*$ 即可计算函子 $Rf_*$。选取一个代表 $K$ 的
$\mathcal{B}$-模复形 $\mathcal{K}^\bullet$，使其 K-平坦且各项平坦；见
“位点上的上同调”引理 \ref{sites-cohomology-lemma-K-flat-resolution}。
于是 $f^*\mathcal{K}^\bullet$ 也 K-平坦且各项平坦；见“位点上的上同调”引理
\ref{sites-cohomology-lemma-pullback-K-flat}。任取一个代表 $M$ 的
$\mathcal{A}$-模复形 $\mathcal{M}^\bullet$。现在只需证明
$$
\text{Tot}(\mathcal{K}^\bullet \otimes_\mathcal{B} f_*\mathcal{M}^\bullet)
=
f_*\text{Tot}(f^*\mathcal{K}^\bullet \otimes_\mathcal{A} \mathcal{M}^\bullet)
$$
因为按照上述选择，这两个复形分别代表引理中公式的右端和左端。
又因为 $f_*$ 与直和可交换（例如由引理
\ref{spaces-cohomology-lemma-stalk-push-finite} 中对茎的描述可知），问题归结为等式
$$
\mathcal{K}^n \otimes_\mathcal{B} f_*\mathcal{M}^m
=
f_*(f^*\mathcal{K}^n \otimes_\mathcal{A} \mathcal{M}^m)
$$
而这些等式由引理 \ref{spaces-cohomology-lemma-finite-rings} 成立。
\end{proof}









\section{余极限与上同调}
\label{spaces-cohomology-section-colimits}

\noindent
下面的引理特别适用于拟凝聚层的图式。

\begin{lemma}
\label{spaces-cohomology-lemma-colimits}
设 $S$ 为概形，$X$ 为 $S$ 上的代数空间。若 $X$ 拟紧且拟分离，则
$$
\colim_i H^p(X, \mathcal{F}_i)
\longrightarrow
H^p(X, \colim_i \mathcal{F}_i)
$$
对 $X_\etale$ 上阿贝尔层的每个滤过图式都是同构。
\end{lemma}

\begin{proof}
这由“位点上的上同调”引理
\ref{sites-cohomology-lemma-colim-works-over-collection} 得到。具体而言，
令 $\mathcal{B} \subset \Ob(X_{spaces, \etale})$ 为所有在 $X$ 上平展、
拟紧且拟分离的空间组成的集合。注意，若 $U \in \mathcal{B}$，则由于 $U$
拟紧，由 $U_i \in \mathcal{B}$ 构成的有限覆盖 $\{U_i \to U\}$ 的全体，
在 $X_{spaces, \etale}$ 中 $U$ 的所有覆盖组成的集合里是共尾的。由
“代数空间的态射”引理
\ref{spaces-morphisms-lemma-quasi-compact-quasi-separated-permanence}，
集合 $\mathcal{B}$ 满足“位点上的上同调”引理
\ref{sites-cohomology-lemma-colim-works-over-collection} 的全部假设。
因为 $X \in \mathcal{B}$，结论成立。
\end{proof}

\begin{lemma}
\label{spaces-cohomology-lemma-colimit-cohomology}
\begin{slogan}
拟紧拟分离态射的高阶正像与层的滤过余极限可交换。
\end{slogan}
设 $S$ 为概形，$f : X \to Y$ 为 $S$ 上代数空间的拟紧拟分离态射。设
$\mathcal{F} = \colim \mathcal{F}_i$ 是 $X_\etale$ 上阿贝尔层的滤过余极限。
则对任意 $p \geq 0$，有
$$
R^pf_*\mathcal{F} = \colim R^pf_*\mathcal{F}_i.
$$
\end{lemma}

\begin{proof}
我们要用到如下事实：拓扑斯态射 $f_{small} : X_{small} \to Y_{small}$
来自位点态射
$f_{spaces, \etale} : X_{spaces, \etale} \to Y_{spaces, \etale}$
它对应于连续函子 $V \longmapsto X \times_Y V$；见“代数空间的性质”引理
\ref{spaces-properties-lemma-functoriality-etale-site}。我们将对这个位点态射
应用“位点上的上同调”引理
\ref{sites-cohomology-lemma-higher-direct-image-colimit}。由于
$Y_{spaces, \etale}$ 的每个对象都有一个由仿射对象组成的覆盖，只需证明：
当 $V$ 仿射且在 $Y$ 上平展时，有
$H^p(X \times_Y V, \mathcal{F}) = \colim H^p(X \times_Y V, \mathcal{F}_i)$.
由于 $V$ 仿射，代数空间 $X \times_Y V$ 拟紧且拟分离。因此应用引理
\ref{spaces-cohomology-lemma-colimits} 即得结论。
\end{proof}

\noindent
下面的引理说明，有限表示模在拟紧拟分离代数空间上具有预期的性质。

\begin{lemma}
\label{spaces-cohomology-lemma-finite-presentation-quasi-compact-colimit}
设 $S$ 为概形，$X$ 为 $S$ 上拟紧拟分离的代数空间。设 $I$ 为有向集，
$(\mathcal{F}_i, \varphi_{ii'})$ 为由 $I$ 索引的 $\mathcal{O}_X$-模系统，
并设 $\mathcal{G}$ 为有限表示的 $\mathcal{O}_X$-模。则
$$
\colim_i \Hom_X(\mathcal{G}, \mathcal{F}_i)
=
\Hom_X(\mathcal{G}, \colim_i \mathcal{F}_i).
$$
特别地，$\Hom_X(\mathcal{G}, -)$ 在 $\QCoh(\mathcal{O}_X)$ 中与滤过余极限可交换。
\end{lemma}

\begin{proof}
所示等式是“位点上的模”引理
\ref{sites-modules-lemma-finite-presentation-quasi-compact-colimit}
的一个特例。为应用该引理，需要对位点 $X_\etale$ 验证“位点”引理
\ref{sites-lemma-directed-colimits-global-sections} 第 (4) 部分的假设。
为此，我们验证“位点”注记 \ref{sites-remark-stronger-conditions}
中的假设 (2)(a)、(2)(b)、(2)(c)。令
$\mathcal{B} \subset \Ob(X_\etale)$ 为仿射对象的集合，则：
\begin{enumerate}
\item 由于 $X$ 拟紧，存在 $U \in \mathcal{B}$，使 $U \to X$ 为满射
（见“代数空间的性质”引理
\ref{spaces-properties-lemma-quasi-compact-affine-cover}），
所以 $h_U^\# \to *$ 为满射。
\item 对 $U \in \mathcal{B}$，$U$ 的每个平展覆盖
$\{U_i \to U\}_{i \in I}$ 都可由一个有限平展覆盖
$\{U_j \to U\}_{j = 1, \ldots, m}$ 加细，其中 $U_j \in \mathcal{B}$
（见“拓扑”引理 \ref{topologies-lemma-etale-affine}）。
\item 对 $U, U' \in \Ob(X_\etale)$，有
$h_U^\# \times h_{U'}^\# = h_{U \times_X U'}^\#$.
若 $U, U' \in \mathcal{B}$，则因为 $X$ 拟分离，$U \times_X U'$ 拟紧；
例如见“代数空间的态射”引理
\ref{spaces-morphisms-lemma-quasi-compact-quasi-separated-permanence}。
因此可找到平展满态射 $U'' \to U \times_X U'$，其中
$U'' \in \mathcal{B}$（见“代数空间的性质”引理
\ref{spaces-properties-lemma-quasi-compact-affine-cover}）。换言之，存在态射
$U'' \to U$ 和 $U'' \to U'$，使得
% P09-ERR-0004: frozen source literally has lower-case u'; see corrections.md.
$h_{U''}^\# \to h_U^\# \times h_{u'}^\#$ 为满射。
\end{enumerate}
关于最后一个断言，注意包含函子
$\QCoh(\mathcal{O}_X) \to \textit{Mod}(\mathcal{O}_X)$ 与余极限可交换，
且有限表示模拟凝聚。参见“代数空间的性质”引理
\ref{spaces-properties-lemma-properties-quasi-coherent}。
\end{proof}





\section{交错 {\v C}ech 复形}
\label{spaces-cohomology-section-alternating-cech}

\noindent
设 $S$ 为概形，$f : U \to X$ 为 $S$ 上代数空间的平展态射。函子
$$
j : U_{spaces, \etale} \longrightarrow X_{spaces, \etale},\quad
V/U \longmapsto V/X
$$
诱导出 $U_{spaces, \etale}$ 与局部化 $X_{spaces, \etale}/U$ 之间的等价；
见“代数空间的性质”第 \ref{spaces-properties-section-localize} 节。因此存在函子
$$
f_! : \textit{Ab}(U_\etale) \longrightarrow
\textit{Ab}(X_\etale),\quad
f_! : \textit{Mod}(\mathcal{O}_U) \longrightarrow \textit{Mod}(\mathcal{O}_X),
$$
它们分别左伴随于
$$
f^{-1} : \textit{Ab}(X_\etale) \longrightarrow
\textit{Ab}(U_\etale),\quad
f^* : \textit{Mod}(\mathcal{O}_X) \longrightarrow \textit{Mod}(\mathcal{O}_U)
$$
见“位点上的模”第 \ref{sites-modules-section-localize} 节。
注意：先验地说，这个函子与紧支撑上同调毫无关系！在上述文献中，
我们把它称为“零延拓”函子。还要注意，对 $\mathcal{O}_X$-模层有
$f^* = f^{-1}$，所以 $f_!$ 的两个版本彼此一致。

\medskip\noindent
下面会用到这个构造，先回顾它的若干性质。给定 $U_\etale$ 上的阿贝尔层
$\mathcal{G}$，
% P09-ERR-0005: frozen source literally omits G after f_!; see ERRATA_CANDIDATES.jsonl.
层 $f_!$ 是下列预层的层化：
$$
V/X \longmapsto
f_!\mathcal{G}(V) =
\bigoplus\nolimits_{\varphi \in \Mor_X(V, U)}
\mathcal{G}(V \xrightarrow{\varphi} U),
$$
见“位点上的模”引理 \ref{sites-modules-lemma-extension-by-zero}。
此外，若 $\mathcal{G}$ 是 $\mathcal{O}_U$-模，则 $f_!\mathcal{G}$ 是同一个
阿贝尔群预层的层化，而该预层以显然的方式带有 $\mathcal{O}_X$-模结构
（见同处）。设 $\overline{x} : \Spec(k) \to X$ 为几何点，则有典范等同
$$
(f_!\mathcal{G})_{\overline{x}} =
\bigoplus\nolimits_{\overline{u}} \mathcal{G}_{\overline{u}}
$$
其中直和遍历所有满足 $f \circ \overline{u} = \overline{x}$ 的
$\overline{u} : \Spec(k) \to U$；见“位点上的模”引理
\ref{sites-modules-lemma-stalk-j-shriek} 和“代数空间的性质”引理
\ref{spaces-properties-lemma-points-small-etale-site}。下文将研究层
$f_!\underline{\mathbf{Z}}$。这里 $\underline{\mathbf{Z}}$ 表示
$X_\etale$ 或 $U_\etale$ 上的常值层。

\begin{lemma}
\label{spaces-cohomology-lemma-product-is-tensor-product}
设 $S$ 为概形，$f_i : U_i \to X$ 为 $S$ 上代数空间的平展态射。则有同构
$$
f_{1, !}\underline{\mathbf{Z}} \otimes_{\mathbf{Z}}
f_{2, !}\underline{\mathbf{Z}}
\longrightarrow
f_{12, !}\underline{\mathbf{Z}}
$$
其中 $f_{12} : U_1 \times_X U_2 \to X$ 是结构态射；并且有
$$
(f_1 \amalg f_2)_! \underline{\mathbf{Z}}
\longrightarrow
f_{1, !}\underline{\mathbf{Z}} \oplus
f_{2, !}\underline{\mathbf{Z}}
$$
\end{lemma}

\begin{proof}
一旦定义了这个映射，由上面对茎的描述可知它是同构。定义该映射时，
只需在预层层次上工作。因此需要定义映射
$$
\left(\bigoplus\nolimits_{\varphi_1 \in \Mor_X(V, U_1)} \mathbf{Z}\right)
\otimes_{\mathbf{Z}}
\left(\bigoplus\nolimits_{\varphi_2 \in \Mor_X(V, U_2)} \mathbf{Z}\right)
\longrightarrow
\bigoplus\nolimits_{\varphi \in \Mor_X(V, U_1 \times_X U_2)}
\mathbf{Z}
$$
使用显然的记号，把元素 $1_{\varphi_1} \otimes 1_{\varphi_2}$ 映到元素
$1_{\varphi_1 \times \varphi_2}$。第二个等式的证明从略。
\end{proof}

\noindent
另一个重要性质是迹映射
$$
\text{Tr}_f : f_!\underline{\mathbf{Z}} \longrightarrow \underline{\mathbf{Z}}.
$$
迹映射伴随于映射 $\mathbf{Z} \to f^{-1}\underline{\mathbf{Z}}$
（后者是同构）。若 $\overline{x}$ 如上，则 $\text{Tr}_f$ 在
$\overline{x}$ 处的茎上是映射
$$
(\text{Tr}_f)_{\overline{x}} :
(f_!\underline{\mathbf{Z}})_{\overline{x}} =
\bigoplus\nolimits_{\overline{u}} \mathbf{Z}
\longrightarrow
\mathbf{Z} = \underline{\mathbf{Z}}_{\overline{x}}
$$
它把给定的整数相加。这是因为它伴随于映射
$1 : \mathbf{Z} \to f^{-1}\underline{\mathbf{Z}}$。特别地，若 $f$ 既平展又
为满射，则 $\text{Tr}_f$ 也是满射。

\medskip\noindent
假设 $f : U \to X$ 是代数空间的平展满态射。考虑与上述迹映射相伴的
{\it Koszul 复形}
$$
\ldots \to \wedge^3f_!\underline{\mathbf{Z}} \to
\wedge^2f_!\underline{\mathbf{Z}} \to f_!\underline{\mathbf{Z}} \to
\underline{\mathbf{Z}} \to 0
$$
这里的外幂取在环层 $\underline{\mathbf{Z}}$ 上。映射由规则
$$
e_1 \wedge \ldots \wedge e_n \longmapsto
\sum\nolimits_{i = 1, \ldots, n} (-1)^{i + 1}
\text{Tr}_f(e_i)
e_1 \wedge \ldots \wedge \widehat{e_i} \wedge \ldots \wedge e_n
$$
定义，其中 $e_1, \ldots, e_n$ 是 $f_!\underline{\mathbf{Z}}$ 的局部截面。
设 $\overline{x}$ 为 $X$ 的几何点，并令
$M_{\overline{x}} = (f_!\underline{\mathbf{Z}})_{\overline{x}} =
\bigoplus_{\overline{u}} \mathbf{Z}$。则上述复形在 $\overline{x}$ 处的茎为复形
$$
\ldots \to \wedge^3 M_{\overline{x}} \to \wedge^2 M_{\overline{x}}
\to M_{\overline{x}} \to \mathbf{Z} \to 0
$$
它是正合的，因为 $M_{\overline{x}} \to \mathbf{Z}$ 为满射；见
“关于代数的更多内容”引理
\ref{more-algebra-lemma-homotopy-koszul-abstract}。因此，若以
$K^\bullet = K^\bullet(f)$ 表示满足
$K^i = \wedge^{i + 1}f_!\underline{\mathbf{Z}}$ 的复形，就得到拟同构
\begin{equation}
\label{spaces-cohomology-equation-quasi-isomorphism}
K^\bullet \longrightarrow \underline{\mathbf{Z}}[0]
\end{equation}
我们用复形 $K^\bullet$ 定义与 $f : U \to X$ 相伴的所谓交错 {\v C}ech 复形。

\begin{definition}
\label{spaces-cohomology-definition-alternating-cech-complex}
设 $S$ 为概形，$f : U \to X$ 为 $S$ 上代数空间的平展满态射，并设
$\mathcal{F}$ 为 $\textit{Ab}(X_\etale)$ 的对象。与 $\mathcal{F}$ 和 $f$
相伴的{\it 交错 {\v C}ech 复形}\footnote{这可能不是标准记号}
$\check{\mathcal{C}}^\bullet_{alt}(f, \mathcal{F})$ 是复形
$$
\Hom(K^0, \mathcal{F}) \to \Hom(K^1, \mathcal{F}) \to
\Hom(K^2, \mathcal{F}) \to \ldots
$$
其中 Hom 群在 $\textit{Ab}(X_\etale)$ 中计算。
\end{definition}

\noindent
读者可以验证：若 $U = \coprod U_i$，且 $f|_{U_i} : U_i \to X$ 是某个子空间
的开浸入，则 $\check{\mathcal{C}}_{alt}^\bullet(f, \mathcal{F})$ 与“上同调”
第 \ref{cohomology-section-alternating-cech} 节针对 Zariski 覆盖
$X = \bigcup U_i$ 及 $\mathcal{F}$ 在 $X$ 的 Zariski 位点上的限制所引入的
复形一致。不过更重要的是，把交错 {\v C}ech 复形的上同调与层上同调联系起来。

\begin{lemma}
\label{spaces-cohomology-lemma-alternating-cech-to-cohomology}
设 $S$ 为概形，$f : U \to X$ 为 $S$ 上代数空间的平展满态射，并设
$\mathcal{F}$ 为 $\textit{Ab}(X_\etale)$ 的对象。则在 $D(\textit{Ab})$ 中
存在典范映射
$$
\check{\mathcal{C}}^\bullet_{alt}(f, \mathcal{F})
\longrightarrow
R\Gamma(X, \mathcal{F})
$$
此外，存在一个 $E_1$ 页为
$$
E_1^{p, q} =
\Ext_{\textit{Ab}(X_\etale)}^q(K^p, \mathcal{F})
$$
的谱序列，它收敛到 $H^{p + q}(X, \mathcal{F})$，其中
$K^p = \wedge^{p + 1}f_!\underline{\mathbf{Z}}$。
\end{lemma}

\begin{proof}
回顾拟同构 $K^\bullet \to \underline{\mathbf{Z}}[0]$，见
(\ref{spaces-cohomology-equation-quasi-isomorphism})。在 $\textit{Ab}(X_\etale)$ 中选取单射
消解 $\mathcal{F} \to \mathcal{I}^\bullet$。考虑二重复形
$\Hom(K^\bullet, \mathcal{I}^\bullet)$，其项为
$\Hom(K^p, \mathcal{I}^q)$。微分
$d_1^{p, q} : A^{p, q} \to A^{p + 1, q}$
来自微分 $K^{p + 1} \to K^p$，而微分
$d_2^{p, q} : A^{p, q} \to A^{p, q + 1}$ 来自微分
$\mathcal{I}^q \to \mathcal{I}^{q + 1}$。以
$\text{Tot}(\Hom(K^\bullet, \mathcal{I}^\bullet))$ 表示相应的全复形；见
“同调代数”第 \ref{homology-section-double-complexes} 节。我们将使用与这个
二重复形相伴的两个谱序列 $({}'E_r, {}'d_r)$ 和 $({}''E_r, {}''d_r)$；见
“同调代数”第 \ref{homology-section-double-complex} 节。

\medskip\noindent
因为 $K^\bullet$ 是 $\underline{\mathbf{Z}}$ 的消解，复形
$$
\Hom(K^\bullet, \mathcal{I}^q) :
\Hom(K^0, \mathcal{I}^q) \to
\Hom(K^1, \mathcal{I}^q) \to
\Hom(K^2, \mathcal{I}^q) \to \ldots
$$
在正次数上无上同调，并且其 $H^0$ 等于 $\Gamma(X, \mathcal{I}^q)$。
因此由“同调代数”引理
\ref{homology-lemma-double-complex-gives-resolution}，自然映射
$$
\mathcal{I}^\bullet(X) \longrightarrow
\text{Tot}(\Hom(K^\bullet, \mathcal{I}^\bullet))
$$
是阿贝尔群复形的拟同构。特别地，有
$H^n(\text{Tot}(\Hom(K^\bullet, \mathcal{I}^\bullet))) = H^n(X, \mathcal{F})$.

\medskip\noindent
本引理中的映射 $\check{\mathcal{C}}^\bullet_{alt}(f, \mathcal{F}) \to
R\Gamma(X, \mathcal{F})$，是映射
$\check{\mathcal{C}}^\bullet_{alt}(f, \mathcal{F}) \to
\text{Tot}(\Hom(K^\bullet, \mathcal{I}^\bullet))$ 与上述拟同构之逆的复合。

\medskip\noindent
最后考虑谱序列 $({}'E_r, {}'d_r)$。有
$$
E_1^{p, q} = q\text{次上同调，属于 }
\Hom(K^p, \mathcal{I}^0) \to
\Hom(K^p, \mathcal{I}^1) \to
\Hom(K^p, \mathcal{I}^2) \to \ldots
$$
这就证明了引理。
\end{proof}

\noindent
由该引理可见，理解 Ext 群
$\Ext_{\textit{Ab}(X_\etale)}(K^p, \mathcal{F})$ 很重要；也就是说，
要理解函子 $\mathcal{F} \mapsto \Hom(K^p, \mathcal{F})$ 的右导出函子。

\begin{lemma}
\label{spaces-cohomology-lemma-compute}
设 $S$ 为概形，$f : U \to X$ 为 $S$ 上代数空间的满、平展且分离的态射。
对 $p \geq 0$，令
$$
W_p = U \times_X \ldots \times_X U \setminus \text{所有对角线}
$$
其中纤维积有 $p + 1$ 个因子。在 $X$ 上，$S_{p + 1}$ 自由作用于 $W_p$，并且
$$
\Hom(K^p, \mathcal{F}) = S_{p + 1}\text{-反不变元，属于 }
\mathcal{F}(W_p)
$$
这一等式关于 $\mathcal{F}$ 函子性地成立，其中
$K^p = \wedge^{p + 1}f_!\underline{\mathbf{Z}}$。
\end{lemma}

\begin{proof}
因为 $U \to X$ 分离，对角态射 $U \to U \times_X U$ 是闭浸入。又因为
$U \to X$ 平展，该对角态射也是开浸入；见“代数空间的态射”引理
\ref{spaces-morphisms-lemma-etale-unramified} 和
\ref{spaces-morphisms-lemma-diagonal-unramified-morphism}。因此 $W_p$ 是
$U^{p + 1} = U \times_X \ldots \times_X U$ 的既开又闭子空间。我们已经删去
这个作用的所有不动点，所以 $S_{p + 1}$ 在 $W_p$ 上的作用是自由的。
由引理 \ref{spaces-cohomology-lemma-product-is-tensor-product}，有
$$
(f_!\underline{\mathbf{Z}})^{\otimes p + 1} =
f^{p + 1}_!\underline{\mathbf{Z}} = (W_p \to X)_!\underline{\mathbf{Z}}
\oplus Rest
$$
其中 $f^{p + 1} : U^{p + 1} \to X$ 是结构态射。考察 $X$ 的几何点
$\overline{x}$ 上方的茎，可见
$$
\left(
\bigoplus\nolimits_{\overline{u} \mapsto \overline{x}} \mathbf{Z}
\right)^{\otimes p + 1}
\longrightarrow
(W_p \to X)_!\underline{\mathbf{Z}}_{\overline{x}}
$$
是一个商映射，其核由所有满足某个 $i \not = j$ 时
$\overline{u}_i = \overline{u}_j$ 的张量
$1_{\overline{u}_0} \otimes \ldots \otimes 1_{\overline{u}_p}$ 生成。
因此商映射
$$
(f_!\underline{\mathbf{Z}})^{\otimes p + 1}
\longrightarrow
\wedge^{p + 1}f_!\underline{\mathbf{Z}}
$$
经由 $(W_p \to X)_!\underline{\mathbf{Z}}$ 分解；也就是说，得到
$$
(f_!\underline{\mathbf{Z}})^{\otimes p + 1}
\longrightarrow
(W_p \to X)_!\underline{\mathbf{Z}}
\longrightarrow
\wedge^{p + 1}f_!\underline{\mathbf{Z}}
$$
这已经证明 $\Hom(K^p, \mathcal{F})$ 函子性地是下列群的子群：
$$
\Hom((W_p \to X)_!\underline{\mathbf{Z}}, \mathcal{F}) = \mathcal{F}(W_p)
$$
要把它认作 $S_{p + 1}$-反不变元，需要证明满射
$(W_p \to X)_!\underline{\mathbf{Z}}
\to \wedge^{p + 1}f_!\underline{\mathbf{Z}}$ 给出最大的
$S_{p + 1}$-反不变商。换言之，需要证明
$\wedge^{p + 1}f_!\underline{\mathbf{Z}}$ 是
$(W_p \to X)_!\underline{\mathbf{Z}}$ 对下述子层的商：该子层由局部截面
$s - \text{符号}(\sigma)\sigma(s)$ 生成，其中 $s$ 是
$(W_p \to X)_!\underline{\mathbf{Z}}$ 的局部截面。这可以在茎上检验，
而在茎上结论显然。
\end{proof}

\begin{lemma}
\label{spaces-cohomology-lemma-twist}
设 $S$ 为概形，$W$ 为 $S$ 上的代数空间，并设有限群 $G$ 自由作用于 $W$。
令 $U = W/G$，见“代数空间的性质”引理
\ref{spaces-properties-lemma-quotient}。设
$\chi : G \to \{+1, -1\}$ 为特征标。则 $U_\etale$ 上存在秩为 $1$ 的局部自由
$\mathbf{Z}$-模层 $\underline{\mathbf{Z}}(\chi)$，使得对 $U_\etale$ 上任意
阿贝尔层 $\mathcal{F}$ 都有
$$
H^0(W, \mathcal{F}|_W)^\chi =
H^0(U, \mathcal{F} \otimes_{\mathbf{Z}} \underline{\mathbf{Z}}(\chi))
$$
\end{lemma}

\begin{proof}
商态射 $q : W \to U$ 是 $G$-挠子；也就是说，存在平展满态射 $U' \to U$，
使得作为 $U'$ 上带 $G$-作用的空间，有
$W \times_U U' = \coprod_{g \in G} U'$（取 $U' = W$ 即可）。因此
$q_*\underline{\mathbf{Z}}$ 是带 $G$-作用的有限局部自由 $\mathbf{Z}$-模。
对 $U$ 的任意几何点 $\overline{u}$，得到 $G$-等变同构
$$
(q_*\underline{\mathbf{Z}})_{\overline{u}}
= \bigoplus\nolimits_{\overline{w} \mapsto \overline{u}} \mathbf{Z}
= \bigoplus\nolimits_{g \in G} \mathbf{Z} = \mathbf{Z}[G]
$$
第二个等号用到位于 $\overline{u}$ 上方的一个几何点 $\overline{w}_0$，
它把对应于 $g \in G$ 的直和项映到对应于 $g(\overline{w}_0)$ 的直和项。有
$$
H^0(W, \mathcal{F}|_W) =
H^0(U, \mathcal{F} \otimes_\mathbf{Z} q_*\underline{\mathbf{Z}})
$$
这是因为
$q_*\mathcal{F}|_W = \mathcal{F} \otimes_\mathbf{Z} q_*\underline{\mathbf{Z}}$；
限制到 $U'$ 后即可检验。令
$$
\underline{\mathbf{Z}}(\chi) =
(q_*\underline{\mathbf{Z}})^\chi \subset
q_*\underline{\mathbf{Z}}
$$
为按 $\chi$ 变换的截面所成的子层。对 $U$ 的任意几何点 $\overline{u}$，有
$$
\underline{\mathbf{Z}}(\chi)_{\overline{u}} =
\mathbf{Z} \cdot \sum\nolimits_g \chi(g) g
\subset
\mathbf{Z}[G] = (q_*\underline{\mathbf{Z}})_{\overline{u}}
$$
由此可知 $\underline{\mathbf{Z}}(\chi)$ 是秩为 $1$ 的局部自由层
（更确切地说，应限制到 $U'$ 后检验）。注意，对任意 $\mathbf{Z}$-模 $M$，
$M[G]$ 的 $\chi$-半不变元恰为形如
$m \cdot \sum\nolimits_g \chi(g) g$ 的元素。因此，对 $U$ 上任意阿贝尔层
$\mathcal{F}$，有
$$
\left(\mathcal{F} \otimes_\mathbf{Z} q_*\underline{\mathbf{Z}}\right)^\chi
=
\mathcal{F} \otimes_\mathbf{Z} \underline{\mathbf{Z}}(\chi)
$$
因为该等式在所有茎上成立。取整体截面即得引理结论。
\end{proof}

\noindent
现在把上述结果汇集起来，得到下面这个优美的结论。

\begin{lemma}
\label{spaces-cohomology-lemma-alternating-spectral-sequence}
设 $S$ 为概形，$f : U \to X$ 为 $S$ 上代数空间的满、平展且分离的态射。
对 $p \geq 0$，令
$$
W_p = U \times_X \ldots \times_X U \setminus \text{所有对角线}
$$
如引理 \ref{spaces-cohomology-lemma-compute}，其中有 $p + 1$ 个因子。设
$\chi_p : S_{p + 1} \to \{+1, -1\}$ 为符号特征标，并令
$U_p = W_p/S_{p + 1}$；$\underline{\mathbf{Z}}(\chi_p)$ 的含义同引理
\ref{spaces-cohomology-lemma-twist}。则引理 \ref{spaces-cohomology-lemma-alternating-cech-to-cohomology}
中的谱序列，其 $E_1$ 页为
$$
E_1^{p, q} =
H^q(U_p, \mathcal{F}|_{U_p} \otimes_\mathbf{Z} \underline{\mathbf{Z}}(\chi_p))
$$
并收敛到 $H^{p + q}(X, \mathcal{F})$。
\end{lemma}

\begin{proof}
注意，$S_{p + 1}$ 在 $W_p$ 上的作用是 $X$ 上的作用，所以确实得到态射
$U_p \to X$。由于 $W_p \to X$ 平展，且 $W_p \to U_p$ 是平展满态射，
$U_p \to X$ 也平展；见“代数空间的态射”引理
\ref{spaces-morphisms-lemma-etale-local}。因此
$\textit{Ab}(X_\etale)$ 的单射对象限制为
$\textit{Ab}(U_{p, \etale})$ 的单射对象；见“位点上的上同调”引理
\ref{sites-cohomology-lemma-cohomology-of-open}。此外，函子
$\mathcal{G} \mapsto
% P09-ERR-0006: frozen source has an extra closing parenthesis; see ERRATA_CANDIDATES.jsonl.
\mathcal{G} \otimes_\mathbf{Z} \underline{\mathbf{Z}}(\chi_p))$
% P09-ERR-0007: frozen source writes Ab(U_p), unlike the surrounding etale categories.
是 $\textit{Ab}(U_p)$ 的自等价，因而把单射对象变为单射对象，并且是正合的
（因为 $\underline{\mathbf{Z}}(\chi_p)$ 是可逆
$\underline{\mathbf{Z}}$-模）。因此，给定 $\textit{Ab}(X_\etale)$ 中的单射
消解 $\mathcal{F} \to \mathcal{I}^\bullet$，复形
$$
\Gamma(U_p,
\mathcal{I}^0|_{U_p} \otimes_\mathbf{Z} \underline{\mathbf{Z}}(\chi_p))
\to
\Gamma(U_p,
\mathcal{I}^1|_{U_p} \otimes_\mathbf{Z} \underline{\mathbf{Z}}(\chi_p))
\to
\Gamma(U_p,
\mathcal{I}^2|_{U_p} \otimes_\mathbf{Z} \underline{\mathbf{Z}}(\chi_p))
\to \ldots
$$
计算
$H^*(U_p,
\mathcal{F}|_{U_p} \otimes_\mathbf{Z} \underline{\mathbf{Z}}(\chi_p))$.
另一方面，由引理 \ref{spaces-cohomology-lemma-twist}，它等于下列复形中
$S_{p + 1}$-反不变元所成的复形：
$$
\Gamma(W_p, \mathcal{I}^0) \to
\Gamma(W_p, \mathcal{I}^1) \to
\Gamma(W_p, \mathcal{I}^2) \to \ldots
$$
由引理 \ref{spaces-cohomology-lemma-compute}，后者又等于复形
$$
\Hom(K^p, \mathcal{I}^0) \to
\Hom(K^p, \mathcal{I}^1) \to
\Hom(K^p, \mathcal{I}^2) \to \ldots
$$
它计算
$\Ext^*_{\textit{Ab}(X_\etale)}(K^p, \mathcal{F})$.
汇总以上各点即得结论。
\end{proof}





\section{拟凝聚层的高阶消失性}
\label{spaces-cohomology-section-higher-vanishing}

\noindent
本节证明：给定拟紧拟分离代数空间 $X$，存在整数 $n = n(X)$，使得 $X$ 上
任意拟凝聚层的上同调在次数超过 $n$ 时消失。

\begin{lemma}
\label{spaces-cohomology-lemma-quasi-coherent-twist}
沿用引理 \ref{spaces-cohomology-lemma-twist} 中的 $S$、$W$、$G$、$U$、$\chi$。
若 $\mathcal{F}$ 是拟凝聚 $\mathcal{O}_U$-模，则
$\mathcal{F} \otimes_{\mathbf{Z}} \underline{\mathbf{Z}}(\chi)$ 也是。
\end{lemma}

\begin{proof}
$\mathcal{O}_U$-模结构是显然的。要检验
$\mathcal{F} \otimes_{\mathbf{Z}} \underline{\mathbf{Z}}(\chi)$ 拟凝聚，
只需平展局部地检验。由于 $\underline{\mathbf{Z}}(\chi)$ 作为
$\underline{\mathbf{Z}}$-模是有限局部自由的，引理随即成立。
\end{proof}

\noindent
即使 $X$ 是概形，下面的命题也很有意义。它是“概形的上同调”引理
\ref{coherent-lemma-vanishing-nr-affines} 的自然推广。在陈述之前先注意：
给定从仿射概形到拟分离代数空间 $X$ 的平展态射 $f : U \to X$，$f$ 的纤维
普遍有界。特别地，存在整数 $d$，使得 $|U| \to |X|$ 的每个纤维至多有
$d$ 个点；这就是“合宜代数空间”引理
\ref{decent-spaces-lemma-bounded-fibres} 中的蕴含
$(\eta) \Rightarrow (\delta)$。

\begin{proposition}
\label{spaces-cohomology-proposition-vanishing}
设 $S$ 为概形，$X$ 为 $S$ 上拟紧且分离的代数空间。设 $U$ 为仿射概形，
$f : U \to X$ 为平展满态射，并设 $d$ 是 $|U| \to |X|$ 各纤维大小的上界。
则对任意拟凝聚 $\mathcal{O}_X$-模 $\mathcal{F}$，当 $q \geq d$ 时有
$H^q(X, \mathcal{F}) = 0$。
\end{proposition}

\begin{proof}
使用引理 \ref{spaces-cohomology-lemma-alternating-spectral-sequence} 的谱序列。由于 $U$ 分离，
$f$ 也分离，所以该引理适用；见“代数空间的态射”引理
\ref{spaces-morphisms-lemma-compose-after-separated}。因为 $X$ 分离，概形
$U \times_X \ldots \times_X U$ 是
$U \times_{\Spec(\mathbf{Z})} \ldots \times_{\Spec(\mathbf{Z})} U$ 的闭子概形，
因而仿射。所以 $W_p$ 仿射。进而由“群胚”命题
\ref{groupoids-proposition-finite-flat-equivalence}，
$U_p = W_p/S_{p + 1}$ 是仿射概形。
% P09-ERR-0008: frozen source switches from U_p to W_p in the next three loci.
第 \ref{spaces-cohomology-section-higher-direct-image} 节的讨论说明，$W_p$ 上拟凝聚层的上同调
（把 $W_p$ 看作代数空间）与底层仿射概形上相应拟凝聚层的上同调一致；
所以由“概形的上同调”引理
\ref{coherent-lemma-quasi-coherent-affine-cohomology-zero}，它在正次数上消失。
由引理 \ref{spaces-cohomology-lemma-quasi-coherent-twist}，层
$\mathcal{F}|_{U_p} \otimes_\mathbf{Z} \underline{\mathbf{Z}}(\chi_p)$
拟凝聚。因此
$H^q(W_p,
\mathcal{F}|_{U_p} \otimes_\mathbf{Z} \underline{\mathbf{Z}}(\chi_p))$
在 $q > 0$ 时为零。由整数 $d$ 的定义，当 $p \geq d$ 时
$W_p = \emptyset$。所以
$H^0(W_p,
\mathcal{F}|_{U_p} \otimes_\mathbf{Z} \underline{\mathbf{Z}}(\chi_p))$
在 $p \geq d$ 时也为零。这就证明了命题。
\end{proof}

\noindent
下面的引理证明：拟紧拟分离代数空间对于拟凝聚模具有有限上同调维数。
我们明确写出这个界，只是因为稍后要用它证明高阶正像的类似结果。

\begin{lemma}
\label{spaces-cohomology-lemma-vanishing-quasi-separated}
设 $S$ 为概形，$X$ 为 $S$ 上拟紧拟分离的代数空间。则可以选取：
\begin{enumerate}
\item 一个仿射概形 $U$；
\item 一个平展满态射 $f : U \to X$；
\item 一个整数 $d$，它界定 $U \to X$ 各纤维的次数；
\item 对每个 $p = 0, 1, \ldots, d$，从仿射概形 $V_p$ 出发的平展满态射
$V_p \to U_p$，其中 $U_p$ 的含义同引理
\ref{spaces-cohomology-lemma-alternating-spectral-sequence}；以及
\item 一个整数 $d_p$，它界定 $V_p \to U_p$ 各纤维的次数。
\end{enumerate}
此外，只要有 (1)--(5)，则对任意拟凝聚 $\mathcal{O}_X$-模
$\mathcal{F}$，当 $q \geq \max(d_p + p)$ 时有
$H^q(X, \mathcal{F}) = 0$。
\end{lemma}

\begin{proof}
由于 $X$ 拟紧，可以找到平展满态射 $U \to X$，其中 $U$ 仿射；见
“代数空间的性质”引理
\ref{spaces-properties-lemma-quasi-compact-affine-cover}。由“合宜代数空间”引理
\ref{decent-spaces-lemma-bounded-fibres}，$f$ 的纤维普遍有界，故可找到 $d$。
有 $U_p = W_p/S_{p + 1}$，且
$W_p \subset U \times_X \ldots \times_X U$ 既开又闭。由于 $X$ 拟分离，
概形 $W_p$ 拟紧，因此 $U_p$ 拟紧。由于 $U$ 分离，概形 $W_p$ 分离，
所以由“代数空间”引理 \ref{spaces-lemma-quotient-finite-separated}
的绝对形式，$U_p$ 分离。由“代数空间的性质”引理
\ref{spaces-properties-lemma-quasi-compact-affine-cover}，
% P09-ERR-0009: frozen source says V_p -> W_p, contrary to the lemma statement.
可以找到态射 $V_p \to W_p$。再由“合宜代数空间”引理
\ref{decent-spaces-lemma-bounded-fibres}，可以找到各整数 $d_p$。

\medskip\noindent
此时，证明使用谱序列
$$
E_1^{p, q} =
H^q(U_p, \mathcal{F}|_{U_p} \otimes_\mathbf{Z} \underline{\mathbf{Z}}(\chi_p))
\Rightarrow
H^{p + q}(X, \mathcal{F})
$$
见引理 \ref{spaces-cohomology-lemma-alternating-spectral-sequence}。由整数 $d$ 的定义，
% P09-ERR-0010: frozen source writes U_p = 0 for a space.
当 $p \geq d$ 时 $U_p = 0$。由命题 \ref{spaces-cohomology-proposition-vanishing} 和引理
\ref{spaces-cohomology-lemma-quasi-coherent-twist}，可知
$H^q(U_p,
\mathcal{F}|_{U_p} \otimes_\mathbf{Z} \underline{\mathbf{Z}}(\chi_p))$
当 $p = 0, \ldots, d$ 且 $q \geq d_p$ 时为零。引理随之成立。
\end{proof}









\section{高阶正像的消失性}
\label{spaces-cohomology-section-vanishing-higher-direct-images}

\noindent
应用第 \ref{spaces-cohomology-section-higher-vanishing} 节的结果，可得拟紧拟分离态射下
拟凝聚层高阶正像的消失性。这很有用，因为在某些情形下，它允许我们对上同调
次数作降归纳。

\begin{lemma}
\label{spaces-cohomology-lemma-vanishing-higher-direct-images}
设 $S$ 为概形，$f : X \to Y$ 为 $S$ 上代数空间的态射。假设：
\begin{enumerate}
\item $f$ 拟紧且拟分离；
\item $Y$ 拟紧。
\end{enumerate}
则存在整数 $n(X \to Y)$，使得对任意代数空间 $Y'$、任意态射
$Y' \to Y$，以及 $X' = Y' \times_Y X$ 上任意拟凝聚层
$\mathcal{F}'$，当 $i \geq n(X \to Y)$ 时，高阶正像
$R^if'_*\mathcal{F}'$ 为零。
\end{lemma}

\begin{proof}
取平展满态射 $V \to Y$，其中 $V$ 是仿射概形；见“代数空间的性质”引理
\ref{spaces-properties-lemma-quasi-compact-affine-cover}。假设已经对基变换
$f_V : V \times_Y X \to V$ 证明了结论，那么结论对 $f$ 也成立，并可取
$n(X \to Y) = n(X_V \to V)$。事实上，若 $Y' \to Y$ 与
$\mathcal{F}'$ 如引理所述，则 $R^if'_*\mathcal{F}'|_{V \times_Y Y'}$
等于 $R^if'_{V, *}\mathcal{F}'|_{X'_V}$，其中
$f'_V : X'_V = V \times_Y Y' \times_Y X \to V \times_Y Y' = Y'_V$；
见“代数空间的性质”引理
\ref{spaces-properties-lemma-pushforward-etale-base-change-modules}。
因此可以假设 $Y$ 是仿射概形。

\medskip\noindent
此外，要对所有 $Y' \to Y$ 与 $\mathcal{F}'$ 证明消失性，只需处理 $Y'$
为仿射概形的情形。此时由引理 \ref{spaces-cohomology-lemma-higher-direct-image}，
$R^if'_*\mathcal{F}'$ 拟凝聚。因此只需证明
$H^i(X', \mathcal{F}') = 0$；因为由“位点上的上同调”引理
\ref{sites-cohomology-lemma-apply-Leray} 以及仿射代数空间上拟凝聚层高阶
上同调的消失性（命题 \ref{spaces-cohomology-proposition-vanishing}），有
$H^i(X', \mathcal{F}') = H^0(Y', R^if'_*\mathcal{F}')$。

\medskip\noindent
按引理 \ref{spaces-cohomology-lemma-vanishing-quasi-separated} 选取
$U \to X$、$d$、$V_p \to U_p$ 与 $d_p$。对任意仿射概形 $Y'$ 及态射
$Y' \to Y$，记
$X' = Y' \times_Y X$、$U' = Y' \times_Y U$、
$V'_p = Y' \times_Y V_p$。则
% P09-ERR-0011: frozen source sets d'_p = d instead of d_p.
$U' \to X'$、$d' = d$、$V'_p \to U'_p$ 以及 $d'_p = d$，构成代数空间
$X'$ 的一组如引理 \ref{spaces-cohomology-lemma-vanishing-quasi-separated} 所述的选择
（细节从略）。所以当 $i \geq \max(p + d_p)$ 时，
$H^i(X', \mathcal{F}') = 0$，结论成立。
\end{proof}

\begin{lemma}
\label{spaces-cohomology-lemma-affine-vanishing-higher-direct-images}
设 $S$ 为概形，$f : X \to Y$ 为 $S$ 上代数空间的仿射态射。则对任意拟凝聚
$\mathcal{O}_X$-模 $\mathcal{F}$，当 $i > 0$ 时
$R^if_*\mathcal{F} = 0$。
\end{lemma}

\begin{proof}
回顾代数空间的仿射态射是可表的。因此结论由
(\ref{spaces-cohomology-equation-representable-higher-direct-image}) 和“概形的上同调”引理
\ref{coherent-lemma-relative-affine-vanishing} 得到。
\end{proof}

\begin{lemma}
\label{spaces-cohomology-lemma-relative-affine-cohomology}
设 $S$ 为概形，$f : X \to Y$ 为 $S$ 上代数空间的仿射态射，并设
$\mathcal{F}$ 为拟凝聚 $\mathcal{O}_X$-模。则对所有 $i \geq 0$，有
$H^i(X, \mathcal{F}) = H^i(Y, f_*\mathcal{F})$。
\end{lemma}

\begin{proof}
这由引理 \ref{spaces-cohomology-lemma-affine-vanishing-higher-direct-images} 和 Leray 谱序列得到。
参见“位点上的上同调”引理 \ref{sites-cohomology-lemma-apply-Leray}。
\end{proof}






\section{支撑在闭子空间上的上同调}
\label{spaces-cohomology-section-cohomology-support}

\noindent
本节是“上同调”第 \ref{cohomology-section-cohomology-support}、
\ref{cohomology-section-cohomology-support-bis} 节以及“平展上同调”第
\ref{etale-cohomology-section-cohomology-support} 节针对代数空间上阿贝尔层的
类似形式。

\medskip\noindent
设 $S$ 为概形，$X$ 为 $S$ 上的代数空间，$Z \subset X$ 为闭子空间，并设
$\mathcal{F}$ 为 $X_\etale$ 上的阿贝尔层。令
$$
\Gamma_Z(X, \mathcal{F}) =
\{s \in \mathcal{F}(X) \mid \text{Supp}(s) \subset Z\}
$$
为支撑在 $Z$ 中的截面（见“代数空间的性质”定义
\ref{spaces-properties-definition-support}）。这是一个左正合函子，但一般并不
正合。因此得到导出函子
$$
R\Gamma_Z(X, -) : D(X_\etale) \longrightarrow D(\textit{Ab})
$$
并把支撑在 $Z$ 中的上同调群定义为
$H^q_Z(X, \mathcal{F}) = R^q\Gamma_Z(X, \mathcal{F})$。

\medskip\noindent
设 $\mathcal{I}$ 为 $X_\etale$ 上的单射阿贝尔层，$U \subset X$ 为 $Z$ 的补
开子空间。则限制映射 $\mathcal{I}(X) \to \mathcal{I}(U)$ 为满射
（见“位点上的上同调”引理
\ref{sites-cohomology-lemma-restriction-along-monomorphism-surjective}），
其核为 $\Gamma_Z(X, \mathcal{I})$。立即可得：对 $K \in D(X_\etale)$，
在 $D(\textit{Ab})$ 中有区分三角
$$
R\Gamma_Z(X, K) \to R\Gamma(X, K) \to R\Gamma(U, K) \to R\Gamma_Z(X, K)[1]
$$
从而得到上同调长正合列
$$
\ldots \to H^i_Z(X, K) \to H^i(X, K) \to H^i(U, K) \to
H^{i + 1}_Z(X, K) \to \ldots
$$
其中 $K$ 是 $D(X_\etale)$ 的任意对象。

\medskip\noindent
对 $X_\etale$ 上的阿贝尔层 $\mathcal{F}$，可以考虑{\it 支撑在 $Z$ 中的截面
所成的子层}，记为 $\mathcal{H}_Z(\mathcal{F})$，其定义规则为
$$
\mathcal{H}_Z(\mathcal{F})(U) =
\{s \in \mathcal{F}(U) \mid \text{Supp}(s) \subset U \times_X Z\}
$$
这里使用“代数空间的性质”定义
\ref{spaces-properties-definition-support} 中的截面支撑。利用“代数空间的态射”
引理 \ref{spaces-morphisms-lemma-closed-immersion-push-pull} 的等价，
可把 $\mathcal{H}_Z(\mathcal{F})$ 看成 $Z_\etale$ 上的阿贝尔层。因此得到函子
$$
\textit{Ab}(X_\etale) \longrightarrow \textit{Ab}(Z_\etale),\quad
\mathcal{F} \longmapsto \mathcal{H}_Z(\mathcal{F})
$$
它左正合，但一般并不正合。

\begin{lemma}
\label{spaces-cohomology-lemma-sections-with-support-acyclic}
设 $S$ 为概形，$i : Z \to X$ 为 $S$ 上代数空间的闭浸入，并设
$\mathcal{I}$ 为 $X_\etale$ 上的单射阿贝尔层。则
$\mathcal{H}_Z(\mathcal{I})$ 是 $Z_\etale$ 上的单射阿贝尔层。
\end{lemma}

\begin{proof}
注意，对 $Z_\etale$ 上任意阿贝尔层 $\mathcal{G}$，有
$$
\Hom_Z(\mathcal{G}, \mathcal{H}_Z(\mathcal{F})) =
\Hom_X(i_*\mathcal{G}, \mathcal{F})
$$
因为 $i_*\mathcal{G}$ 的任意截面终究都支撑在 $Z$ 中。由于 $i_*$ 正合
（引理 \ref{spaces-cohomology-lemma-finite-higher-direct-image-zero}），且 $\mathcal{I}$ 在
$X_\etale$ 上单射，可知 $\mathcal{H}_Z(\mathcal{I})$ 在 $Z_\etale$ 上单射。
\end{proof}

\noindent
以
$$
R\mathcal{H}_Z : D(X_\etale) \longrightarrow D(Z_\etale)
$$
表示导出函子。令
$\mathcal{H}^q_Z(\mathcal{F}) = R^q\mathcal{H}_Z(\mathcal{F})$，从而
$\mathcal{H}^0_Z(\mathcal{F}) = \mathcal{H}_Z(\mathcal{F})$。
由上面的引理，有 Grothendieck 谱序列
$$
E_2^{p, q} = H^p(Z, \mathcal{H}^q_Z(\mathcal{F}))
\Rightarrow H^{p + q}_Z(X, \mathcal{F})
$$

\begin{lemma}
\label{spaces-cohomology-lemma-cohomology-with-support-sheaf-on-support}
设 $S$ 为概形，$i : Z \to X$ 为 $S$ 上代数空间的闭浸入，并设
$\mathcal{G}$ 为 $Z_\etale$ 上的单射阿贝尔层。则当 $p > 0$ 时，
$\mathcal{H}^p_Z(i_*\mathcal{G}) = 0$。
\end{lemma}

\begin{proof}
这是因为函子 $i_*$ 正合（引理
\ref{spaces-cohomology-lemma-finite-higher-direct-image-zero}），并把单射阿贝尔层变为单射
阿贝尔层（见“位点上的上同调”引理
\ref{sites-cohomology-lemma-pushforward-injective-flat}）。
\end{proof}

\begin{lemma}
\label{spaces-cohomology-lemma-etale-localization-sheaf-with-support}
设 $S$ 为概形，$f : X \to Y$ 为 $S$ 上代数空间的平展态射。设
$Z \subset Y$ 为闭子空间，且 $f^{-1}(Z) \to Z$ 是代数空间的同构。
% P09-ERR-0012: frozen source places F on X although it is then pulled back from Y.
设 $\mathcal{F}$ 为 $X$ 上的阿贝尔层。则
$$
\mathcal{H}^q_Z(\mathcal{F}) = \mathcal{H}^q_{f^{-1}(Z)}(f^{-1}\mathcal{F})
$$
作为 $Z = f^{-1}(Z)$ 上的阿贝尔层有上述等式，并且
$H^q_Z(Y, \mathcal{F}) = H^q_{f^{-1}(Z)}(X, f^{-1}\mathcal{F})$。
\end{lemma}

\begin{proof}
因为 $f$ 平展，$\mathcal{F}$ 的单射消解拉回为 $f^{-1}\mathcal{F}$ 的单射
消解。因此只需对 $\mathcal{H}_Z(-)$ 检验等式，而这由定义直接得到。
支撑上同调的证明相同。若干细节从略。
\end{proof}

\noindent
设 $S$ 为概形，$X$ 为 $S$ 上的代数空间，$T \subset |X|$ 为闭子集。
以 $D_T(X_\etale)$ 表示 $D(X_\etale)$ 的严格满且饱和的三角子范畴，
其对象的上同调层都支撑在 $T$ 上。

\begin{lemma}
\label{spaces-cohomology-lemma-complexes-with-support-on-closed}
设 $S$ 为概形，$i : Z \to X$ 为 $S$ 上代数空间的闭浸入。映射
$Ri_* = i_* : D(Z_\etale) \to D(X_\etale)$ 诱导等价
$D(Z_\etale) \to D_{|Z|}(X_\etale)$，其拟逆为
$$
% P09-ERR-0013: frozen source alternates D_Z and the defined D_|Z| notation.
i^{-1}|_{D_Z(X_\etale)} = R\mathcal{H}_Z|_{D_{|Z|}(X_\etale)}
$$
\end{lemma}

\begin{proof}
回顾 $i^{-1}$ 与 $i_*$ 是一对正合伴随函子，并且
% P09-ERR-0014: frozen source says identify functor; see ERRATA_CANDIDATES.jsonl.
$i^{-1}i_*$ 同构于阿贝尔层上的 identify 函子。见“代数空间的性质”引理
\ref{spaces-properties-lemma-stalk-pullback} 和“代数空间的态射”引理
\ref{spaces-morphisms-lemma-closed-immersion-push-pull}。因此
$i_* : D(Z_\etale) \to D_Z(X_\etale)$ 全忠实，而 $i^{-1}$ 给出左逆。
另一方面，设 $K$ 为 $D_Z(X_\etale)$ 的对象，并考虑伴随映射
$K \to i_*i^{-1}K$.
由 $i_*$ 与 $i^{-1}$ 的正合性，它在上同调层上诱导伴随映射
$H^n(K) \to i_*i^{-1}H^n(K)$。由于这些上同调层支撑在 $Z$ 上，
这些伴随映射是同构，从而 $D(Z_\etale) \to D_Z(X_\etale)$ 是等价。

\medskip\noindent
最后还要证明：若 $K$ 是 $D_Z(X_\etale)$ 的对象，则
$R\mathcal{H}_Z(K) = i^{-1}K$。可以使用刚刚证明的
$K = i_*i^{-1}K$。为 $i^{-1}K$ 选取一个 K-单射代表
$\mathcal{I}^\bullet$。由于 $i_*$ 是正合函子 $i^{-1}$ 的右伴随，复形
$i_*\mathcal{I}^\bullet$ 是 K-单射的（见“导出范畴”引理
\ref{derived-lemma-adjoint-preserve-K-injectives}）。于是
$R\mathcal{H}_Z(K)$ 由
$\mathcal{H}_Z(i_*\mathcal{I}^\bullet) = \mathcal{I}^\bullet$ 计算，
正合所需。
\end{proof}






\section{维数以上的消失性}
\label{spaces-cohomology-section-vanishing-above-dimension}

\noindent
设 $S$ 为概形，$X$ 为 $S$ 上拟紧拟分离的代数空间。此时 $|X|$ 是谱空间；
见“代数空间的性质”引理
\ref{spaces-properties-lemma-quasi-compact-quasi-separated-spectral}。
此外，$X$ 的维数（定义见“代数空间的性质”定义
\ref{spaces-properties-definition-dimension}）等于 $|X|$ 的 Krull 维数；
见“合宜代数空间”引理
\ref{decent-spaces-lemma-dimension-decent-space}。我们将证明：$X$ 上拟凝聚层
的上同调在超过维数的次数上消失。即便只考虑域上有限型的拟分离代数空间，
这个结果也已经很有意义。

\begin{lemma}
\label{spaces-cohomology-lemma-vanishing-above-dimension}
设 $S$ 为概形，$X$ 为 $S$ 上拟紧拟分离的代数空间，并假设对某个整数 $d$ 有
$\dim(X) \leq d$。
% P09-ERR-0015: frozen source repeats F in the noun phrase; see ERRATA_CANDIDATES.jsonl.
设 $\mathcal{F}$ 为 $X$ 上的拟凝聚层 $\mathcal{F}$。
\begin{enumerate}
\item $H^q(X, \mathcal{F}) = 0$ for $q > d$,
\item 对任意拟紧开子空间 $U \subset X$，
$H^d(X, \mathcal{F}) \to H^d(U, \mathcal{F})$ 为满射；
\item 对任意补集拟紧的闭子空间 $Z \subset X$，当 $q > d$ 时
$H^q_Z(X, \mathcal{F}) = 0$。
\end{enumerate}
\end{lemma}

\begin{proof}
由“代数空间的性质”引理
\ref{spaces-properties-lemma-dimension-decent-invariant-under-etale}，
每个在 $X$ 上平展的代数空间 $Y$ 都满足维数 $\leq d$。若 $Y$ 拟分离，
则由“合宜代数空间”引理
\ref{decent-spaces-lemma-dimension-decent-space}，$Y$ 的维数等于 $|Y|$ 的
Krull 维数。此外，若 $Y$ 是概形，则 $\mathcal{F}$ 在 $Y$ 上的平展上同调
（相应地，支撑在闭子概形上的平展上同调）与 $\mathcal{F}$ 的通常上同调
（相应地，支撑在该闭子概形上的通常上同调）一致。见“下降”命题
\ref{descent-proposition-same-cohomology-quasi-coherent} 和“平展上同调”引理
\ref{etale-cohomology-lemma-cohomology-with-support-quasi-coherent}。
下文将不加说明地使用这些事实。

\medskip\noindent
由“合宜代数空间”引理
\ref{decent-spaces-lemma-filter-quasi-compact-quasi-separated}，
存在整数 $n$ 以及开子空间
$$
\emptyset = U_{n + 1} \subset
U_n \subset U_{n - 1} \subset \ldots \subset U_1 = X
$$
它们具有如下性质：令 $T_p = U_p \setminus U_{p + 1}$，赋予其诱导的既约
子空间结构；则存在拟紧分离概形 $V_p$ 以及平展满态射
$f_p : V_p \to U_p$，使 $f_p^{-1}(T_p) \to T_p$ 为同构。

\medskip\noindent
由于 $U_n = V_n$ 是概形，由开头的说明和“上同调”命题
\ref{cohomology-proposition-cohomological-dimension-spectral}，
$\mathcal{F}$ 在 $U_n$ 上的上同调在次数 $> d$ 时消失。归纳地假设已经证明：
当 $q > d$ 时，$H^q(U_{p + 1}, \mathcal{F}|_{U_{p + 1}}) = 0$。
要推出 $\mathcal{F}$ 在 $U_p$ 上的上同调在次数 $> d$ 时消失，只需证明
当 $q > d$ 时 $H_{T_p}^q(U_p, \mathcal{F})$ 为零。而由引理
\ref{spaces-cohomology-lemma-etale-localization-sheaf-with-support}，有
$$
H^q_{T_p}(U_p, \mathcal{F}) = H^q_{f_p^{-1}(T_p)}(V_p, \mathcal{F})
$$
又因为 $V_p$ 是概形，所需消失性由“上同调”命题
\ref{cohomology-proposition-cohomological-dimension-spectral} 得到。
由此证明了 (1)。

\medskip\noindent
为证明 (2)，设 $U \subset X$ 为拟紧开子空间。考虑开子空间
$U' = U \cup U_n$，并令 $Z = U' \setminus U$。则
$g : U_n \to U'$ 是平展态射，且 $g^{-1}(Z) \to Z$ 为同构。因此由引理
\ref{spaces-cohomology-lemma-etale-localization-sheaf-with-support}，有
$H^q_Z(U', \mathcal{F}) = H^q_Z(U_n, \mathcal{F})$。由于 $U_n$ 是概形，
可应用“上同调”命题
\ref{cohomology-proposition-cohomological-dimension-spectral}，所以它在次数
$> d$ 时消失。于是 $H^d(U', \mathcal{F}) \to H^d(U, \mathcal{F})$
为满射。归纳地假设问题已经归结到 $U$ 包含 $U_{p + 1}$ 的情形。
令 $U' = U \cup U_p$、$Z = U' \setminus U$，并使用平展态射
$f_p : V_p \to U'$；它满足 $f_p^{-1}(Z) \to Z$ 为同构。换言之，再次有
$$
H^q_Z(U', \mathcal{F}) = H^q_{f_p^{-1}(Z)}(V_p, \mathcal{F})
$$
它在次数 $> d$ 时消失。因此
$H^d(U', \mathcal{F}) \to H^d(U, \mathcal{F})$ 为满射。最终到达
$U_1 = X \subset U$ 的阶段，证明完成。

\medskip\noindent
形式论证说明 (2) 蕴含 (3)。
\end{proof}








\section{上同调与基变换（一）}
\label{spaces-cohomology-section-cohomology-and-base-change}

\noindent
设 $S$ 为一个概形。设 $f : X \to Y$ 为 $S$ 上代数空间之间的态射，
并设 $\mathcal{F}$ 为 $X$ 上的拟凝聚层。再设 $g : Y' \to Y$ 为
$S$ 上代数空间之间的态射。记 $X' = X_{Y'} = Y' \times_Y X$ 为
$X$ 的基变换，并记 $f' : X' \to Y'$ 为 $f$ 的基变换。
还把投影记为 $g' : X' \to X$，并令
$\mathcal{F}' = (g')^*\mathcal{F}$。下图表示这一情形：
\begin{equation}
\label{spaces-cohomology-equation-base-change-diagram}
\vcenter{
\xymatrix{
\mathcal{F}' = (g')^*\mathcal{F} &
X' \ar[r]_{g'} \ar[d]_{f'} &
X \ar[d]^f &
\mathcal{F} \\
Rf'_*\mathcal{F}' &
Y' \ar[r]^g &
Y &
Rf_*\mathcal{F}
}
}
\end{equation}
下面是我们所考虑的基变换性质最简单的情形。

\begin{lemma}
\label{spaces-cohomology-lemma-affine-base-change}
设 $S$ 为一个概形。设 $f : X \to Y$ 为 $S$ 上代数空间之间的仿射态射，
并设 $\mathcal{F}$ 为拟凝聚 $\mathcal{O}_X$-模。此时
$f_*\mathcal{F} \cong Rf_*\mathcal{F}$ 是拟凝聚层，而且对于每个图
(\ref{spaces-cohomology-equation-base-change-diagram})，都有
$$
g^*f_*\mathcal{F} = f'_*(g')^*\mathcal{F}.
$$
\end{lemma}

\begin{proof}
由 (\ref{spaces-cohomology-equation-representable-higher-direct-image}) 附近的讨论，
问题归结为概形之间仿射态射的情形；该情形已在《概形的上同调》引理
\ref{coherent-lemma-affine-base-change} 中处理。
\end{proof}

\begin{lemma}[平坦基变换]
\label{spaces-cohomology-lemma-flat-base-change-cohomology}
设 $S$ 为一个概形。考虑 $S$ 上代数空间的笛卡尔图
$$
\xymatrix{
X' \ar[d]_{f'} \ar[r]_{g'} & X \ar[d]^f \\
Y' \ar[r]^g & Y
}
$$
设 $\mathcal{F}$ 为拟凝聚 $\mathcal{O}_X$-模，其拉回为
$\mathcal{F}' = (g')^*\mathcal{F}$。假设 $g$ 平坦，且 $f$ 拟紧、拟分离。
对于任意 $i \geq 0$，
\begin{enumerate}
\item 《站点上的上同调》引理
\ref{sites-cohomology-lemma-base-change-map-flat-case}
中的基变换映射
$$
g^*R^if_*\mathcal{F} \longrightarrow R^if'_*\mathcal{F}',
$$
是同构；
\item 如果 $Y = \Spec(A)$ 且 $Y' = \Spec(B)$，则
$H^i(X, \mathcal{F}) \otimes_A B = H^i(X', \mathcal{F}')$.
\end{enumerate}
\end{lemma}

\begin{proof}
由《空间的态射》引理 \ref{spaces-morphisms-lemma-base-change-flat}，
态射 $g'$ 平坦。注意，$g$ 和 $g'$ 的平坦性等价于相应小 \'etale 带环站点
之间态射的平坦性，见《空间的态射》引理
\ref{spaces-morphisms-lemma-flat-morphism-sites}。因此可以应用
《站点上的上同调》引理
\ref{sites-cohomology-lemma-base-change-map-flat-case}
得到基变换映射
$$
g^*R^pf_*\mathcal{F} \longrightarrow R^pf'_*\mathcal{F}'
$$
为了证明此映射是同构，可以在 $Y'$ 的 \'etale 拓扑下局部地验证。
因此可以假设 $Y$ 和 $Y'$ 都是仿射概形。不妨写成
$Y = \Spec(A)$ 和 $Y' = \Spec(B)$。此时真正要证明的是映射
$$
H^p(X, \mathcal{F}) \otimes_A B \longrightarrow H^p(X_B, \mathcal{F}_B)
$$
是同构，其中 $X_B = \Spec(B) \times_{\Spec(A)} X$，而
$\mathcal{F}_B$ 是 $\mathcal{F}$ 在 $X_B$ 上的拉回。换言之，只需证明 (2)。

\medskip\noindent
固定一个平坦环同态 $A \to B$，并设 $X$ 为 $A$ 上拟紧且拟分离的代数空间。
注意，$g' : X_B \to X$ 是 $\Spec(B) \to \Spec(A)$ 的基变换，因而是仿射的。
故由引理 \ref{spaces-cohomology-lemma-affine-vanishing-higher-direct-images}，
高阶正像 $R^i(g')_*\mathcal{F}_B$ 均为零。因此
$H^p(X_B, \mathcal{F}_B) = H^p(X, g'_*\mathcal{F}_B)$，见
《站点上的上同调》引理 \ref{sites-cohomology-lemma-apply-Leray}。
此外，有
$$
g'_*\mathcal{F}_B = \mathcal{F} \otimes_{\underline{A}} \underline{B}
$$
其中 $\underline{A}$、$\underline{B}$ 表示取值分别为 $A$、$B$ 的常值环层。
事实上，显然存在一个从右端到左端的映射。对于任意 \'etale 于 $X$ 的
仿射概形 $U$，有
\begin{align*}
g'_*\mathcal{F}_B(U) & = \mathcal{F}_B(\Spec(B) \times_{\Spec(A)} U) \\
& =
\Gamma(\Spec(B) \times_{\Spec(A)} U,
(\Spec(B) \times_{\Spec(A)} U \to U)^*\mathcal{F}|_U) \\
& =
B \otimes_A \mathcal{F}(U)
\end{align*}
所以该映射是同构。利用 Lazard 定理，把 $B = \colim M_i$ 写成有限自由
$A$-模 $M_i$ 的滤过余极限，见《代数》定理
\ref{algebra-theorem-lazard}。由此得到
\begin{align*}
H^p(X, g'_*\mathcal{F}_B) &
= H^p(X, \mathcal{F} \otimes_{\underline{A}} \underline{B}) \\
& = H^p(X, \colim_i \mathcal{F} \otimes_{\underline{A}} \underline{M_i}) \\
& = \colim_i H^p(X, \mathcal{F} \otimes_{\underline{A}} \underline{M_i}) \\
& = \colim_i H^p(X, \mathcal{F}) \otimes_A M_i \\
& = H^p(X, \mathcal{F}) \otimes_A \colim_i M_i \\
& = H^p(X, \mathcal{F}) \otimes_A B
\end{align*}
第一个等式来自上面已经证明的
$g'_*\mathcal{F}_B = \mathcal{F} \otimes_{\underline{A}} \underline{B}$。
第二个等式是因为 $\otimes$ 与余极限交换。第三个等式是因为 $X$ 上的
上同调与余极限交换（见引理 \ref{spaces-cohomology-lemma-colimits}）。第四个等式是因为
$M_i$ 有限自由（也就是说，上同调与有限直和交换）。第五个等式是因为
$\otimes$ 与余极限交换。第六个等式来自我们对该系统的选取。
\end{proof}

\begin{lemma}
\label{spaces-cohomology-lemma-etale-pull-push-split}
设 $f : X \to Y$ 为代数空间之间拟紧、分离的 \'etale 态射。则对于任意
拟凝聚 $\mathcal{O}_X$-模 $\mathcal{F}$，映射
$f^*f_*\mathcal{F} \to \mathcal{F}$ 是分裂的。
\end{lemma}

\begin{proof}
考虑笛卡尔图
$$
\xymatrix{
X \times_Y X \ar[r]_-p \ar[d]_q & X \ar[d]^f \\
X \ar[r]^f & Y
}
$$
由引理 \ref{spaces-cohomology-lemma-affine-base-change}，有
$f^*f_*\mathcal{F} = q_*p^*\mathcal{F}$。态射
$\Delta : X \to X \times_Y X$ 是开的（因为它是两个 \'etale 于 $Y$ 的
代数空间之间的态射），又因 $f$ 分离而是闭的。因此可见，
$p^*\mathcal{F}$ 是 $\Delta_*\mathcal{F}$ 与一个支撑在
$\Delta(X)$ 的（开闭）补集上的拟凝聚模的直和。追踪这些映射，读者可验证
$$
\mathcal{F} = q_*\Delta_*\mathcal{F} \to
q_*p^*\mathcal{F} = f^*f_*\mathcal{F} \to
\mathcal{F}
$$
是恒等映射。细节从略。
\end{proof}



\section{局部 Noether 代数空间上的凝聚模}
\label{spaces-cohomology-section-coherent}

\noindent
本节对应于《概形的上同调》\ref{coherent-section-coherent-sheaves} 节。
在《站点上的模》定义 \ref{sites-modules-definition-site-local} 中，
我们已经定义了任意带环拓扑斯上的凝聚模。这里用这一概念定义局部
Noether 代数空间上的凝聚模。虽然也可以在更一般的情形下研究凝聚模，
但我们不打算这样做。

\begin{definition}
\label{spaces-cohomology-definition-coherent}
设 $S$ 为一个概形，$X$ 为 $S$ 上的局部 Noether 代数空间。如果 $X$ 上的
拟凝聚模 $\mathcal{F}$ 按照《站点上的模》定义
\ref{sites-modules-definition-site-local} 的意义，是站点 $X_\etale$ 上的
凝聚 $\mathcal{O}_X$-模，则称 $\mathcal{F}$ 是{\it 凝聚的}。
\end{definition}

\noindent
这一定义与局部 Noether 概形上已有的凝聚模概念相容；见《空间的性质》
\ref{spaces-properties-section-properties-modules} 节中的断言 (5)（或更直接地，
见《下降》引理 \ref{descent-lemma-equivalence-quasi-coherent-properties}）。
因此，从现在起，如果 $X$ 是 $S$ 上的局部 Noether 概形，我们将不再区分
把 $X$ 看作概形时的凝聚模与把 $X$ 看作代数空间时的凝聚模；这与
《空间的性质》\ref{spaces-properties-section-quasi-coherent} 节所讨论的
拟凝聚模范畴的相应等同相容。

\medskip\noindent
在上述约定下，下面的引理给出了局部 Noether 代数空间上凝聚模的一种
直观刻画。

\begin{lemma}
\label{spaces-cohomology-lemma-coherent-Noetherian}
设 $S$ 为一个概形，$X$ 为 $S$ 上的局部 Noether 代数空间，并设
$\mathcal{F}$ 为一个 $\mathcal{O}_X$-模。下列条件等价：
\begin{enumerate}
\item $\mathcal{F}$ 是凝聚的；
\item $\mathcal{F}$ 是有限型拟凝聚 $\mathcal{O}_X$-模；
\item $\mathcal{F}$ 是有限表示的 $\mathcal{O}_X$-模；
\item 对于任意 \'etale 态射 $\varphi : U \to X$，其中 $U$ 是概形，
拉回 $\varphi^*\mathcal{F}$ 是 $U$ 上的凝聚模；
\item 存在满 \'etale 态射 $\varphi : U \to X$，其中 $U$ 是概形，
使得拉回 $\varphi^*\mathcal{F}$ 是 $U$ 上的凝聚模。
\end{enumerate}
特别地，$\mathcal{O}_X$ 是凝聚的，任意可逆 $\mathcal{O}_X$-模是凝聚的；
更一般地，任意有限局部自由 $\mathcal{O}_X$-模都是凝聚的。
\end{lemma}

\begin{proof}
首先，如果 $X$ 是局部 Noether 代数空间，而 $U \to X$ 是 \'etale 态射，
则 $U$ 是局部 Noether 的，见《空间的性质》
\ref{spaces-properties-section-types-properties} 节。于是本引理由
《空间的性质》\ref{spaces-properties-section-properties-modules} 节中的
(1)--(5) 以及局部 Noether 概形上凝聚模的相应结论推出；后者见
《概形的上同调》引理 \ref{coherent-lemma-coherent-Noetherian}。
\end{proof}

\begin{lemma}
\label{spaces-cohomology-lemma-coherent-abelian-Noetherian}
设 $S$ 为一个概形，$X$ 为 $S$ 上的局部 Noether 代数空间。
凝聚 $\mathcal{O}_X$-模所成的范畴是 Abel 范畴。更确切地说，凝聚
$\mathcal{O}_X$-模之间态射的核与余核都是凝聚的。凝聚层的任意扩张也是
凝聚的。
\end{lemma}

\begin{proof}
选取一个概形 $U$ 和满 \'etale 态射 $f : U \to X$。拉回 $f^*$ 等于限制函子，
因而是正合函子，见《空间的性质》公式
(\ref{spaces-properties-equation-restrict-modules})。由引理
\ref{spaces-cohomology-lemma-coherent-Noetherian}，可以通过检验 $f^*\mathcal{F}$ 是否凝聚来
检验 $\mathcal{O}_X$-模 $\mathcal{F}$ 是否凝聚。因此本引理由概形的情形
推出；该情形即《概形的上同调》引理
\ref{coherent-lemma-coherent-abelian-Noetherian}。
\end{proof}

\noindent
凝聚模构成拟凝聚 $\mathcal{O}_X$-模范畴的一个 Serre 子范畴。对于一般
带环拓扑斯上的模，这一结论并不成立。

\begin{lemma}
\label{spaces-cohomology-lemma-coherent-Noetherian-quasi-coherent-sub-quotient}
设 $S$ 为一个概形，$X$ 为 $S$ 上的局部 Noether 代数空间，并设
$\mathcal{F}$ 为凝聚 $\mathcal{O}_X$-模。$\mathcal{F}$ 的任意拟凝聚子模
都是凝聚的，$\mathcal{F}$ 的任意拟凝聚商模也都是凝聚的。
\end{lemma}

\begin{proof}
选取一个概形 $U$ 和满 \'etale 态射 $f : U \to X$。拉回 $f^*$ 等于限制函子，
因而是正合函子，见《空间的性质》公式
(\ref{spaces-properties-equation-restrict-modules})。由引理
\ref{spaces-cohomology-lemma-coherent-Noetherian}，可以通过检验 $f^*\mathcal{H}$ 是否凝聚来
检验 $\mathcal{O}_X$-模 $\mathcal{G}$ 是否凝聚。% P09-ERR-0016
因此本引理由概形的情形推出；该情形即《概形的上同调》引理
\ref{coherent-lemma-coherent-Noetherian-quasi-coherent-sub-quotient}。
\end{proof}

\begin{lemma}
\label{spaces-cohomology-lemma-tensor-hom-coherent}
设 $S$ 为一个概形。设 $X$ 为 $S$ 上的局部 Noether 代数空间，。% P09-ERR-0017
设 $\mathcal{F}$、$\mathcal{G}$ 为凝聚 $\mathcal{O}_X$-模。则
$\mathcal{O}_X$-模 $\mathcal{F} \otimes_{\mathcal{O}_X} \mathcal{G}$ 和
$\SheafHom_{\mathcal{O}_X}(\mathcal{F}, \mathcal{G})$ 都是凝聚的。
\end{lemma}

\begin{proof}
由引理 \ref{spaces-cohomology-lemma-coherent-Noetherian}，这可由概形的相应结论推出，见
《概形的上同调》引理 \ref{coherent-lemma-tensor-hom-coherent}。
\end{proof}

\begin{lemma}
\label{spaces-cohomology-lemma-local-isomorphism}
设 $S$ 为一个概形，$X$ 为 $S$ 上的局部 Noether 代数空间。设
$\mathcal{F}$、$\mathcal{G}$ 为凝聚 $\mathcal{O}_X$-模，且
$\varphi : \mathcal{G} \to \mathcal{F}$ 为 $\mathcal{O}_X$-模同态。
设 $\overline{x}$ 为 $X$ 上位于 $x \in |X|$ 之上的几何点。
\begin{enumerate}
\item 如果 $\mathcal{F}_{\overline{x}} = 0$，则存在 $x$ 的开邻域
$X' \subset X$，使得 $\mathcal{F}|_{X'} = 0$。
\item 如果 $\varphi_{\overline{x}} : \mathcal{G}_{\overline{x}} \to
\mathcal{F}_{\overline{x}}$ 是单射，则存在 $x$ 的开邻域 $X' \subset X$，
使得 $\varphi|_{X'}$ 是单射。
\item 如果 $\varphi_{\overline{x}} : \mathcal{G}_{\overline{x}} \to
\mathcal{F}_{\overline{x}}$ 是满射，则存在 $x$ 的开邻域 $X' \subset X$，
使得 $\varphi|_{X'}$ 是满射。
\item 如果 $\varphi_{\overline{x}} : \mathcal{G}_{\overline{x}} \to
\mathcal{F}_{\overline{x}}$ 是双射，则存在 $x$ 的开邻域 $X' \subset X$，
使得 $\varphi|_{X'}$ 是同构。
\end{enumerate}
\end{lemma}

\begin{proof}
设 $\varphi : U \to X$ 为 \'etale 态射，其中 $U$ 是概形，并设 $u \in U$
为映到 $x$ 的点。由《空间的性质》引理
\ref{spaces-properties-lemma-stalk-quasi-coherent} 及
\ref{spaces-properties-lemma-describe-etale-local-ring}
以及《代数进阶》引理
\ref{more-algebra-lemma-dumb-properties-henselization}，可见
$\varphi_{\overline{x}}$ 是单射、满射或双射，当且仅当
$\varphi_u : \varphi^*\mathcal{F}_u \to \varphi^*\mathcal{G}_u$
具有相应性质。% P09-ERR-0018
因此可以应用本引理的概形版本，得到（必要时缩小 $U$ 后）映射
$\varphi^*\mathcal{F} \to \varphi^*\mathcal{G}$ 是单射、满射或同构。
设 $X' \subset X$ 为与 $|\varphi|(|U|) \subset |X|$ 对应的开子空间，见
《空间的性质》引理 \ref{spaces-properties-lemma-open-subspaces}。由于
$\{U \to X'\}$ 是 \'etale 拓扑的覆盖，便得到 $\varphi|_{X'}$ 如所要求地
为单射、满射或同构。最后，考察映射 $\mathcal{F} \to 0$ 即可看出 (1)
由 (2) 推出。
\end{proof}

\begin{lemma}
\label{spaces-cohomology-lemma-coherent-support-closed}
设 $S$ 为一个概形，$X$ 为 $S$ 上的局部 Noether 代数空间。设
$\mathcal{F}$ 为凝聚 $\mathcal{O}_X$-模，$i : Z \to X$ 为
$\mathcal{F}$ 的概形论支撑，并设 $\mathcal{G}$ 为满足
$i_*\mathcal{G} = \mathcal{F}$ 的拟凝聚 $\mathcal{O}_Z$-模，见
《空间的态射》定义
\ref{spaces-morphisms-definition-scheme-theoretic-support}。则
$\mathcal{G}$ 是凝聚 $\mathcal{O}_Z$-模。
\end{lemma}

\begin{proof}
本引理的陈述是有意义的，因为凝聚模尤其是有限型的。此外，$Z \to X$
是闭浸入，因而是局部有限型的，所以 $Z$ 局部 Noether；见《空间的态射》引理
\ref{spaces-morphisms-lemma-immersion-locally-finite-type} 及
\ref{spaces-morphisms-lemma-locally-finite-type-locally-noetherian}。
最后，由于 $\mathcal{G}$ 是有限型的，引理
\ref{spaces-cohomology-lemma-coherent-Noetherian} 表明它是凝聚 $\mathcal{O}_Z$-模。
\end{proof}

\begin{lemma}
\label{spaces-cohomology-lemma-i-star-equivalence}
设 $S$ 为一个概形，$i : Z \to X$ 为 $S$ 上局部 Noether 代数空间之间的
闭浸入。设 $\mathcal{I} \subset \mathcal{O}_X$ 为定义 $Z$ 的拟凝聚理想层。
函子 $i_*$ 在被 $\mathcal{I}$ 零化的凝聚 $\mathcal{O}_X$-模范畴与
凝聚 $\mathcal{O}_Z$-模范畴之间诱导一个等价。
\end{lemma}

\begin{proof}
由《空间的态射》引理 \ref{spaces-morphisms-lemma-i-star-equivalence}，
此函子全忠实。设 $\mathcal{F}$ 为被 $\mathcal{I}$ 零化的凝聚
$\mathcal{O}_X$-模。由《空间的态射》引理
\ref{spaces-morphisms-lemma-i-star-equivalence}，可将
$\mathcal{F}$ 写成 $i_*\mathcal{G}$，其中 $\mathcal{G}$ 是 $Z$ 上的某个
拟凝聚层。为了检验 $\mathcal{G}$ 凝聚，可以在 \'etale 拓扑下局部地验证
（引理 \ref{spaces-cohomology-lemma-coherent-Noetherian}）。选取一个由概形给出的 \'etale
覆盖后，概形的情形（《概形的上同调》引理
\ref{coherent-lemma-i-star-equivalence}）表明 $\mathcal{G}$ 凝聚。
因此该函子全忠实，证明完成。% P09-ERR-0019
\end{proof}

\begin{lemma}
\label{spaces-cohomology-lemma-finite-pushforward-coherent}
设 $S$ 为一个概形，$f : X \to Y$ 为 $S$ 上代数空间之间的有限态射，且
$Y$ 局部 Noether。设 $\mathcal{F}$ 为凝聚 $\mathcal{O}_X$-模。假设 $f$
有限且 $Y$ 局部 Noether。% P09-ERR-0020
则当 $p > 0$ 时 $R^pf_*\mathcal{F} = 0$，并且 $f_*\mathcal{F}$ 凝聚。
\end{lemma}

\begin{proof}
选取一个概形 $V$ 和满 \'etale 态射 $V \to Y$。则
$V \times_Y X \to V$ 是局部 Noether 概形之间的有限态射。由
(\ref{spaces-cohomology-equation-representable-higher-direct-image})，问题归结为概形的情形；
该情形即《概形的上同调》引理
\ref{coherent-lemma-finite-pushforward-coherent}。
\end{proof}






\section{Noether 空间上的凝聚层}
\label{spaces-cohomology-section-coherent-quasi-compact}

\noindent
本节介绍 Noether 代数空间上凝聚层的一些性质。

\begin{lemma}
\label{spaces-cohomology-lemma-acc-coherent}
设 $S$ 为一个概形，$X$ 为 $S$ 上的 Noether 代数空间，并设
$\mathcal{F}$ 为凝聚 $\mathcal{O}_X$-模。$\mathcal{F}$ 的拟凝聚子模满足
升链条件。换言之，给定任意拟凝聚子模序列
$$
\mathcal{F}_1 \subset \mathcal{F}_2 \subset \ldots \subset \mathcal{F}
$$
则存在某个 $n \geq 0$，使得
$\mathcal{F}_n = \mathcal{F}_{n + 1} = \ldots $。
\end{lemma}

\begin{proof}
选取一个仿射概形 $U$ 和满 \'etale 态射 $U \to X$（见《空间的性质》引理
\ref{spaces-properties-lemma-quasi-compact-affine-cover}）。由
《空间的态射》引理
\ref{spaces-morphisms-lemma-locally-finite-type-locally-noetherian}，
$U$ 是 Noether 概形。如果
$\mathcal{F}_n|_U = \mathcal{F}_{n + 1}|_U = \ldots$，则
$\mathcal{F}_n = \mathcal{F}_{n + 1} = \ldots$。因此结论可由概形的情形
推出，见《概形的上同调》引理 \ref{coherent-lemma-acc-coherent}。
\end{proof}

\begin{lemma}
\label{spaces-cohomology-lemma-power-ideal-kills-sheaf}
设 $S$ 为一个概形，$X$ 为 $S$ 上的 Noether 代数空间。设 $\mathcal{F}$ 为
$X$ 上的凝聚层，并设 $\mathcal{I} \subset \mathcal{O}_X$ 为与闭子空间
$Z \subset X$ 对应的拟凝聚理想层。存在某个 $n \geq 0$ 使
$\mathcal{I}^n\mathcal{F} = 0$，当且仅当（在集合论意义下）
$\text{Supp}(\mathcal{F}) \subset Z$。
\end{lemma}

\begin{proof}
选取一个仿射概形 $U$ 和满 \'etale 态射 $U \to X$（见《空间的性质》引理
\ref{spaces-properties-lemma-quasi-compact-affine-cover}）。由
《空间的态射》引理
\ref{spaces-morphisms-lemma-locally-finite-type-locally-noetherian}，
$U$ 是 Noether 概形。注意，$\mathcal{I}^n\mathcal{F}|_U = 0$ 当且仅当
$\mathcal{I}^n\mathcal{F} = 0$，支撑的条件也同样如此。因此结论可由概形
的情形推出，见《概形的上同调》引理
\ref{coherent-lemma-power-ideal-kills-sheaf}。
\end{proof}

\begin{lemma}[Artin-Rees]
\label{spaces-cohomology-lemma-Artin-Rees}
设 $S$ 为一个概形，$X$ 为 $S$ 上的 Noether 代数空间。设 $\mathcal{F}$ 为
$X$ 上的凝聚层，$\mathcal{G} \subset \mathcal{F}$ 为拟凝聚子层，并设
$\mathcal{I} \subset \mathcal{O}_X$ 为拟凝聚理想层。则存在 $c \geq 0$，
使得对于所有 $n \geq c$，都有
$$
\mathcal{I}^{n - c}(\mathcal{I}^c\mathcal{F} \cap \mathcal{G})
=
\mathcal{I}^n\mathcal{F} \cap \mathcal{G}
$$
\end{lemma}

\begin{proof}
选取一个仿射概形 $U$ 和满 \'etale 态射 $U \to X$（见《空间的性质》引理
\ref{spaces-properties-lemma-quasi-compact-affine-cover}）。由
《空间的态射》引理
\ref{spaces-morphisms-lemma-locally-finite-type-locally-noetherian}，
$U$ 是 Noether 概形。本引理中的等式成立，当且仅当把它限制到 $U$ 后
成立。因此结论可由概形的情形推出，见《概形的上同调》引理
\ref{coherent-lemma-Artin-Rees}。
\end{proof}

\begin{lemma}
\label{spaces-cohomology-lemma-homs-over-open}
设 $S$ 为一个概形，$X$ 为 $S$ 上的 Noether 代数空间。设 $\mathcal{F}$ 为
拟凝聚 $\mathcal{O}_X$-模，$\mathcal{G}$ 为凝聚 $\mathcal{O}_X$-模，并设
$\mathcal{I} \subset \mathcal{O}_X$ 为拟凝聚理想层。记 $Z \subset X$ 为
相应的闭子空间，并令 $U = X \setminus Z$。存在典范同构
$$
\colim_n \Hom_{\mathcal{O}_X}(\mathcal{I}^n\mathcal{G}, \mathcal{F})
\longrightarrow
\Hom_{\mathcal{O}_U}(\mathcal{G}|_U, \mathcal{F}|_U).
$$
特别地，有同构
$$
\colim_n \Hom_{\mathcal{O}_X}(\mathcal{I}^n, \mathcal{F})
\longrightarrow
\Gamma(U, \mathcal{F}).
$$
\end{lemma}

\begin{proof}
设 $W$ 为仿射概形，并设 $W \to X$ 为满 \'etale 态射（见《空间的性质》
引理 \ref{spaces-properties-lemma-quasi-compact-affine-cover}）。令
$R = W \times_X W$。则 $W$ 和 $R$ 都是 Noether 概形，见《空间的态射》
引理 \ref{spaces-morphisms-lemma-locally-finite-type-locally-noetherian}。
因此，由《概形的上同调》引理 \ref{coherent-lemma-homs-over-open}，
该结论对于 $\mathcal{F}$、$\mathcal{G}$、$\mathcal{I}$、$U$、$Z$ 在
$W$ 和 $R$ 上的限制成立。% P09-ERR-0021
由此形式地推出该结论在 $X$ 上成立。
\end{proof}






\section{凝聚层的逐层剥离}
\label{spaces-cohomology-section-devissage}

\noindent
本节对应于《概形的上同调》\ref{coherent-section-devissage} 节。

\begin{lemma}
\label{spaces-cohomology-lemma-prepare-filter-support}
设 $S$ 为一个概形，$X$ 为 $S$ 上的 Noether 代数空间，并设 $\mathcal{F}$
为 $X$ 上的凝聚层。假设
$\text{Supp}(\mathcal{F}) = Z \cup Z'$，其中 $Z$、$Z'$ 都是闭集。
则存在凝聚层的短正合列
$$
0 \to \mathcal{G}' \to \mathcal{F} \to \mathcal{G} \to 0
$$
满足 $\text{Supp}(\mathcal{G}') \subset Z'$ 且
$\text{Supp}(\mathcal{G}) \subset Z$。
\end{lemma}

\begin{proof}
设 $\mathcal{I} \subset \mathcal{O}_X$ 为定义 $Z$ 上约化诱导闭子空间结构的
理想层，见《空间的性质》引理
\ref{spaces-properties-lemma-reduced-closed-subspace}。考虑子层
$\mathcal{G}'_n = \mathcal{I}^n\mathcal{F}$ 以及商层
$\mathcal{G}_n = \mathcal{F}/\mathcal{I}^n\mathcal{F}$。对于每个 $n$，都有
短正合列
$$
0 \to \mathcal{G}'_n \to \mathcal{F} \to \mathcal{G}_n \to 0
$$
对于 $Z' \setminus Z$ 的每个几何点 $\overline{x}$，有
$\mathcal{I}_{\overline{x}} = \mathcal{O}_{X, \overline{x}}$，因而
$\mathcal{G}_{n, \overline{x}} = 0$。所以
$\text{Supp}(\mathcal{G}_n) \subset Z$。注意，$X \setminus Z'$ 是 Noether
代数空间。因此由引理 \ref{spaces-cohomology-lemma-power-ideal-kills-sheaf}，存在 $n$ 使得
$\mathcal{G}'_n|_{X \setminus Z'} =
\mathcal{I}^n\mathcal{F}|_{X \setminus Z'} = 0$。对于这样的 $n$，有
$\text{Supp}(\mathcal{G}'_n) \subset Z'$。取
$\mathcal{G}' = \mathcal{G}'_n$ 和 $\mathcal{G} = \mathcal{G}_n$ 即可。
\end{proof}

\noindent
下文将自由使用《空间的态射》定义
\ref{spaces-morphisms-definition-scheme-theoretic-support} 中所定义的
有限型模的概形论支撑。

\begin{lemma}
\label{spaces-cohomology-lemma-prepare-filter-irreducible}
设 $S$ 为一个概形，$X$ 为 $S$ 上的 Noether 代数空间，并设 $\mathcal{F}$
为 $X$ 上的凝聚层。假设 $\mathcal{F}$ 的概形论支撑是约化的
$Z \subset X$，且 $|Z|$ 不可约。则存在整数 $r > 0$、非零理想层
$\mathcal{I} \subset \mathcal{O}_Z$，以及凝聚层之间的单射
$$
i_*\left(\mathcal{I}^{\oplus r}\right) \to \mathcal{F}
$$
其余核支撑在 $Z$ 的一个真闭子空间上。% P09-ERR-0022
\end{lemma}

\begin{proof}
由假设，存在支撑为 $Z$ 的凝聚 $\mathcal{O}_Z$-模 $\mathcal{G}$，使得
$\mathcal{F} \cong i_*\mathcal{G}$，见引理
\ref{spaces-cohomology-lemma-coherent-support-closed}。因此只需在 $Z = X$ 且
$i = \text{id}$ 的情形下证明本引理。

\medskip\noindent
由《空间的性质》命题
\ref{spaces-properties-proposition-locally-quasi-separated-open-dense-scheme}
存在一个稠密开子空间 $U \subset X$，且 $U$ 是概形。注意，$U$ 是 Noether
整概形。缩小 $U$ 后，可以假设
$\mathcal{F}|_U \cong \mathcal{O}_U^{\oplus r}$（例如由《概形的上同调》
引理 \ref{coherent-lemma-prepare-filter-irreducible}，或由直接的代数论证）。
设 $\mathcal{I} \subset \mathcal{O}_X$ 为拟凝聚理想层，其相应闭子空间是
$U$ 在 $X$ 中的补集（例如见《空间的性质》
\ref{spaces-properties-section-reduced} 节）。由引理
\ref{spaces-cohomology-lemma-homs-over-open}，存在 $n \geq 0$ 以及态射
$\mathcal{I}^n(\mathcal{O}_X^{\oplus r}) \to \mathcal{F}$，它在 $U$ 上恢复
上述同构。由于
$\mathcal{I}^n(\mathcal{O}_X^{\oplus r}) = (\mathcal{I}^n)^{\oplus r}$，
便得到引理所述的映射。该映射是单射：事实上，若 $\sigma$ 是某个概形
$W$ 上 $\mathcal{I}^{\oplus r}$ 的非零截面，其中 $W$ \'etale 于 $X$，
则由于 $X$ 约化，从而 $W$ 约化，$\sigma$ 的支撑包含 $W$ 的一个非空开集。
然而，$(\mathcal{I}^n)^{\oplus r} \to \mathcal{F}$ 的核在一个稠密开集上
为零，所以 $\sigma$ 不可能是该核的截面。
\end{proof}

\begin{lemma}
\label{spaces-cohomology-lemma-coherent-filter}
设 $S$ 为一个概形，$X$ 为 $S$ 上的 Noether 代数空间，并设 $\mathcal{F}$
为 $X$ 上的凝聚层。存在由凝聚子层组成的滤过
$$
0 = \mathcal{F}_0 \subset \mathcal{F}_1 \subset
\ldots \subset \mathcal{F}_m = \mathcal{F}
$$
使得对于每个 $j = 1, \ldots, m$，存在约化闭子空间 $Z_j \subset X$，
其中 $|Z_j|$ 不可约，并存在理想层
$\mathcal{I}_j \subset \mathcal{O}_{Z_j}$，使得
$$
\mathcal{F}_j/\mathcal{F}_{j - 1}
\cong (Z_j \to X)_* \mathcal{I}_j
$$
\end{lemma}

\begin{proof}
考虑集合
$$
\mathcal{T} =
\left\{
\begin{matrix}
T \subset |X|
\text{ 为闭集，且存在凝聚层 }
\mathcal{F} \\
\text{ 满足 }
\text{Supp}(\mathcal{F}) = T
\text{，而本引理对它不成立}
\end{matrix}
\right\}
$$
我们要证明 $\mathcal{T}$ 为空。否则，由于 $|X|$ 是 Noether 的（《空间的性质》
引理 \ref{spaces-properties-lemma-Noetherian-topology}），可以选取
$\mathcal{T}$ 的极小元 $T$。这意味着存在 $X$ 上支撑为 $T$ 的凝聚层
$\mathcal{F}$，而本引理对它不成立。显然 $T \not = \emptyset$，因为支撑
为空的层只能是零层，而本引理对零层成立（取 $m = 0$）。

\medskip\noindent
如果 $T$ 可约，则可写成 $T = Z_1 \cup Z_2$，其中 $Z_1$、$Z_2$ 都是
严格小于 $T$ 的闭集。于是可以应用引理
\ref{spaces-cohomology-lemma-prepare-filter-support}，得到凝聚层的短正合列
$$
0 \to
\mathcal{G}_1 \to
\mathcal{F} \to
\mathcal{G}_2 \to 0
$$
满足 $\text{Supp}(\mathcal{G}_i) \subset Z_i$。由 $T$ 的极小性，每个
$\mathcal{G}_i$ 都有引理陈述中的滤过。考察 $\mathcal{F}$ 上诱导的滤过，
便得到矛盾。因此 $T$ 不可约。

\medskip\noindent
假设 $T$ 不可约。设 $\mathcal{J}$ 为定义 $T$ 上约化诱导闭子空间结构的
理想层，见《空间的性质》引理
\ref{spaces-properties-lemma-reduced-closed-subspace}。由引理
\ref{spaces-cohomology-lemma-power-ideal-kills-sheaf}，存在 $n \geq 0$ 使得
$\mathcal{J}^n\mathcal{F} = 0$。因而得到滤过
$$
0 = \mathcal{I}^n\mathcal{F} \subset \mathcal{I}^{n - 1}\mathcal{F}
\subset \ldots \subset \mathcal{I}\mathcal{F} \subset \mathcal{F}
$$
其每个相继次商都被 $\mathcal{J}$ 零化。% P09-ERR-0023
因此，如果这些次商中的每一个都有引理陈述中的滤过，那么
$\mathcal{F}$ 也有。换言之，可以假设 $\mathcal{J}$ 零化 $\mathcal{F}$。

\medskip\noindent
假设 $T$ 不可约，且 $\mathcal{J}\mathcal{F} = 0$，其中 $\mathcal{J}$ 如上。
则 $\mathcal{F}$ 的概形论支撑为 $T$，见《空间的态射》引理
\ref{spaces-morphisms-lemma-i-star-equivalence}。因此可以应用引理
\ref{spaces-cohomology-lemma-prepare-filter-irreducible}，得到短正合列
$$
0 \to
i_*(\mathcal{I}^{\oplus r}) \to
\mathcal{F} \to
\mathcal{Q} \to 0
$$
其中 $\mathcal{Q}$ 的支撑是 $T$ 的真闭子集。由 $T$ 的极小性，
$\mathcal{Q}$ 有所需类型的滤过。于是显然 $\mathcal{F}$ 也有这样的滤过，
这给出最终的矛盾。
\end{proof}

\begin{lemma}
\label{spaces-cohomology-lemma-property-initial}
设 $S$ 为一个概形，$X$ 为 $S$ 上的 Noether 代数空间，并设
$\mathcal{P}$ 为 $X$ 上凝聚层的一条性质。假设：
\begin{enumerate}
\item 对于凝聚层的任意短正合列
$$
0 \to \mathcal{F}_1 \to \mathcal{F} \to \mathcal{F}_2 \to 0
$$
如果 $\mathcal{F}_i$（$i = 1, 2$）具有性质 $\mathcal{P}$，则
$\mathcal{F}$ 也具有性质 $\mathcal{P}$；
\item 对于每个满足 $|Z|$ 不可约的约化闭子空间 $Z \subset X$，以及每个
拟凝聚理想层 $\mathcal{I} \subset \mathcal{O}_Z$，$i_*\mathcal{I}$
都具有性质 $\mathcal{P}$。
\end{enumerate}
则 $X$ 上的每个凝聚层都具有性质 $\mathcal{P}$。
\end{lemma}

\begin{proof}
首先注意，如果凝聚层 $\mathcal{F}$ 有由凝聚子层组成的滤过
$$
0 = \mathcal{F}_0 \subset \mathcal{F}_1 \subset
\ldots \subset \mathcal{F}_m = \mathcal{F}
$$
并且每个 $\mathcal{F}_i/\mathcal{F}_{i - 1}$ 都具有性质 $\mathcal{P}$，
则 $\mathcal{F}$ 也具有性质 $\mathcal{P}$。这由 $\mathcal{P}$ 的性质 (1)
推出。另一方面，由引理 \ref{spaces-cohomology-lemma-coherent-filter}，可以为任意
$\mathcal{F}$ 构造滤过，使其相继次商如 (2) 所述。因此本引理成立。
\end{proof}

\noindent
下面给出上述引理的一个更实用的变体。

\begin{lemma}
\label{spaces-cohomology-lemma-property-higher-rank-cohomological}
设 $S$ 为一个概形，$X$ 为 $S$ 上的 Noether 代数空间，并设
$\mathcal{P}$ 为 $X$ 上凝聚层的一条性质。假设：
\begin{enumerate}
\item 对于凝聚层的任意短正合列
$$
0 \to \mathcal{F}_1 \to \mathcal{F} \to \mathcal{F}_2 \to 0
$$
如果 $\mathcal{F}_i$（$i = 1, 2$）具有性质 $\mathcal{P}$，则
$\mathcal{F}$ 也具有性质 $\mathcal{P}$；
\item 如果对于某个 $r \geq 1$，$\mathcal{F}^{\oplus r}$ 具有性质
$\mathcal{P}$，则 $\mathcal{F}$ 也具有性质 $\mathcal{P}$；
\item 对于每个满足 $|Z|$ 不可约的约化闭子空间 $i : Z \to X$，
存在 $Z$ 上的凝聚层 $\mathcal{G}$，使得
\begin{enumerate}
\item $\text{Supp}(\mathcal{G}) = Z$；
\item 对于每个非零拟凝聚理想层
$\mathcal{I} \subset \mathcal{O}_Z$，存在拟凝聚子层
$\mathcal{G}' \subset \mathcal{I}\mathcal{G}$，使得
$\text{Supp}(\mathcal{G}/\mathcal{G}')$ 是 $|Z|$ 中的真闭集，并且
$i_*\mathcal{G}'$ 具有性质 $\mathcal{P}$。
\end{enumerate}
\end{enumerate}
则 $X$ 上的每个凝聚层都具有性质 $\mathcal{P}$。
\end{lemma}

\begin{proof}
考虑集合
$$
\mathcal{T} =
\left\{
\begin{matrix}
T \subset |X|
\text{ 为非空闭集，且存在凝聚层 } \\
\mathcal{F}
\text{ 满足 }
\text{Supp}(\mathcal{F}) = T
\text{，而本引理对它不成立}
\end{matrix}
\right\}
$$
我们要证明 $\mathcal{T}$ 为空。否则，由于 $|X|$ 是 Noether 的（《空间的性质》
引理 \ref{spaces-properties-lemma-Noetherian-topology}），可以选取
$\mathcal{T}$ 的极小元 $T$。这意味着存在 $X$ 上支撑为 $T$ 的凝聚层
$\mathcal{F}$，而本引理对它不成立。

\medskip\noindent
如果 $T$ 可约，则可写成 $T = Z_1 \cup Z_2$，其中 $Z_1$、$Z_2$ 都是
严格小于 $T$ 的闭集。于是可以应用引理
\ref{spaces-cohomology-lemma-prepare-filter-support}，得到凝聚层的短正合列
$$
0 \to
\mathcal{G}_1 \to
\mathcal{F} \to
\mathcal{G}_2 \to 0
$$
满足 $\text{Supp}(\mathcal{G}_i) \subset Z_i$。由 $T$ 的极小性，每个
$\mathcal{G}_i$ 都具有性质 $\mathcal{P}$。于是由 (1)，$\mathcal{F}$
具有性质 $\mathcal{P}$，这给出矛盾。

\medskip\noindent
假设 $T$ 不可约。设 $\mathcal{J}$ 为定义 $T$ 上约化诱导闭子空间结构的
理想层，见《空间的性质》引理
\ref{spaces-properties-lemma-reduced-closed-subspace}。由引理
\ref{spaces-cohomology-lemma-power-ideal-kills-sheaf}，存在 $n \geq 0$ 使得
$\mathcal{J}^n\mathcal{F} = 0$。因而得到滤过
$$
0 = \mathcal{J}^n\mathcal{F} \subset \mathcal{J}^{n - 1}\mathcal{F}
\subset \ldots \subset \mathcal{J}\mathcal{F} \subset \mathcal{F}
$$
其每个相继次商都被 $\mathcal{J}$ 零化。因此，如果这些次商中的每一个
都有本引理陈述中的滤过，那么由 (1)，$\mathcal{F}$ 也有。% P09-ERR-0025
换言之，可以假设 $\mathcal{J}$ 零化 $\mathcal{F}$。

\medskip\noindent
假设 $T$ 不可约，且 $\mathcal{J}\mathcal{F} = 0$，其中 $\mathcal{J}$ 如上。
记 $i : Z \to X$ 为与 $\mathcal{J}$ 对应的闭子空间。则对某个凝聚
$\mathcal{O}_Z$-模 $\mathcal{H}$，有 $\mathcal{F} = i_*\mathcal{H}$；见
《空间的态射》引理 \ref{spaces-morphisms-lemma-i-star-equivalence} 以及
引理 \ref{spaces-cohomology-lemma-coherent-support-closed}。设 $\mathcal{G}$ 为满足
(3)(a) 和 (3)(b) 的 $Z$ 上凝聚层。应用引理
\ref{spaces-cohomology-lemma-prepare-filter-irreducible}，得到单射
$$
\mathcal{I}_1^{\oplus r_1} \to \mathcal{H}
\quad\text{以及}\quad
\mathcal{I}_2^{\oplus r_2} \to \mathcal{G}
$$
其余核的支撑是 $Z$ 中的真闭集。因此存在非空开集 $V \subset Z$，使得
$$
\mathcal{H}^{\oplus r_2}_V \cong \mathcal{G}^{\oplus r_1}_V
$$
设 $\mathcal{I} \subset \mathcal{O}_Z$ 为定义 $Z \setminus V$ 的拟凝聚
理想层，我们得到（引理 \ref{spaces-cohomology-lemma-homs-over-open}）映射
$$
\mathcal{I}^n\mathcal{G}^{\oplus r_1} \longrightarrow \mathcal{H}^{\oplus r_2}
$$
它在 $V$ 上是同构。其核支撑在 $Z \setminus V$ 上，因而被
$\mathcal{I}$ 的某个幂零化，见引理
\ref{spaces-cohomology-lemma-power-ideal-kills-sheaf}。因此增大 $n$ 后，可以假设所示映射
是单射，见引理 \ref{spaces-cohomology-lemma-Artin-Rees}。应用 (3)(b)，可找到
$\mathcal{G}' \subset \mathcal{I}^n\mathcal{G}$，使得
$$
(i_*\mathcal{G}')^{\oplus r_1} \longrightarrow
i_*\mathcal{H}^{\oplus r_2} = \mathcal{F}^{\oplus r_2}
$$
是单射，其余核支撑在 $Z$ 的真闭子集上，并且 $i_*\mathcal{G}'$ 具有性质
$\mathcal{P}$。由 (1)，$(i_*\mathcal{G}')^{\oplus r_1}$ 具有性质
$\mathcal{P}$。再由 (1) 以及 $T = |Z|$ 的极小性，
$\mathcal{F}^{\oplus r_2}$ 具有性质 $\mathcal{P}$。最后由 (2)，
$\mathcal{F}$ 具有性质 $\mathcal{P}$，这正是所需的矛盾。
\end{proof}

\begin{lemma}
\label{spaces-cohomology-lemma-property-higher-rank-cohomological-variant}
设 $S$ 为一个概形，$X$ 为 $S$ 上的 Noether 代数空间，并设
$\mathcal{P}$ 为 $X$ 上凝聚层的一条性质。假设：
\begin{enumerate}
\item 对于 $X$ 上凝聚层的任意短正合列，如果其中三个层中的两个具有
性质 $\mathcal{P}$，则第三个也具有性质 $\mathcal{P}$；
\item 如果对于某个 $r \geq 1$，$\mathcal{F}^{\oplus r}$ 具有性质
$\mathcal{P}$，则 $\mathcal{F}$ 也具有性质 $\mathcal{P}$；
\item 对于每个满足 $|Z|$ 不可约的约化闭子空间 $i : Z \to X$，
存在 $X$ 上概形论支撑为 $Z$ 的凝聚层 $\mathcal{G}$，并且
$\mathcal{G}$ 具有性质 $\mathcal{P}$。
\end{enumerate}
则 $X$ 上的每个凝聚层都具有性质 $\mathcal{P}$。
\end{lemma}

\begin{proof}
我们将证明引理 \ref{spaces-cohomology-lemma-property-initial} 的条件 (1) 和 (2) 成立。
条件 (1) 显然成立。为证明 (2)，令
$$
\mathcal{T} =
\left\{
\begin{matrix}
i : Z \to X \text{ 为满足 }|Z|\text{ 不可约的约化闭子空间，并且}\\
\text{ 对某个拟凝聚 }\mathcal{I} \subset \mathcal{O}_Z\text{，}
i_*\mathcal{I}\text{ 不具有 }\mathcal{P}
\end{matrix}
\right\}
$$
如果 $\mathcal{T}$ 非空，则由于 $X$ 是 Noether 的，可以在
$\mathcal{T}$ 中找到极小的 $i : Z \to X$。下面证明这会导致矛盾。

\medskip\noindent
设 $\mathcal{G}$ 为假设 (3) 中假定存在的、概形论支撑为 $Z$ 的层。% P09-ERR-0026
设 $\varphi : i_*\mathcal{I}^{\oplus r} \to \mathcal{G}$ 如引理
\ref{spaces-cohomology-lemma-prepare-filter-irreducible} 所述。设
$$
0 = \mathcal{F}_0 \subset \mathcal{F}_1 \subset
\ldots \subset \mathcal{F}_m = \Coker(\varphi)
$$
为引理 \ref{spaces-cohomology-lemma-coherent-filter} 所述的滤过。由 $Z$ 的极小性和假设 (1)，
$\Coker(\varphi)$ 具有性质 $\mathcal{P}$。由于 $\varphi$ 是单射，再次使用
假设 (1) 可知 $i_*\mathcal{I}^{\oplus r}$ 具有性质 $\mathcal{P}$。
使用假设 (2)，便得到 $i_*\mathcal{I}$ 具有性质 $\mathcal{P}$。

\medskip\noindent
最后，若 $\mathcal{J} \subset \mathcal{O}_Z$ 是另一个拟凝聚理想层，令
$\mathcal{K} = \mathcal{I} \cap \mathcal{J}$，并考虑短正合列
$$
0 \to \mathcal{K} \to \mathcal{I} \to \mathcal{I}/\mathcal{K} \to 0
\quad
\text{以及}
\quad
0 \to \mathcal{K} \to \mathcal{J} \to \mathcal{J}/\mathcal{K} \to 0
$$
沿用上述论证并使用 $Z$ 的极小性，可见
$i_*\mathcal{I}/\mathcal{K}$ 和 $i_*\mathcal{J}/\mathcal{K}$
具有性质 $\mathcal{P}$。% P09-ERR-0027
因此由假设 (1)，先得到 $i_*\mathcal{K}$ 具有性质 $\mathcal{P}$，再得到
$i_*\mathcal{J}$ 具有性质 $\mathcal{P}$。换言之，$Z$ 不是
$\mathcal{T}$ 的元素，这正是所需的矛盾。
\end{proof}






\section{凝聚模的极限}
\label{spaces-cohomology-section-limits}

\noindent
（局部 Noether 代数空间上的）凝聚模的余极限通常并不凝聚。但它是拟凝聚的，
因为代数空间上拟凝聚模的任意余极限都是拟凝聚的，见《空间的性质》引理
\ref{spaces-properties-lemma-properties-quasi-coherent}。反过来，如果该
代数空间是 Noether 的，那么每个拟凝聚模都是凝聚模的滤过余极限。

\begin{lemma}
\label{spaces-cohomology-lemma-directed-colimit-coherent}
设 $S$ 为一个概形，$X$ 为 $S$ 上的 Noether 代数空间。每个拟凝聚
$\mathcal{O}_X$-模都是其凝聚子模的滤过余极限。
\end{lemma}

\begin{proof}
设 $\mathcal{F}$ 为拟凝聚 $\mathcal{O}_X$-模。如果
$\mathcal{G}, \mathcal{H} \subset \mathcal{F}$ 是凝聚
$\mathcal{O}_X$-子模，则 $\mathcal{G} \oplus \mathcal{H} \to \mathcal{F}$
的像是另一个同时包含二者的凝聚 $\mathcal{O}_X$-子模（见引理
\ref{spaces-cohomology-lemma-coherent-abelian-Noetherian} 和
\ref{spaces-cohomology-lemma-coherent-Noetherian-quasi-coherent-sub-quotient}）。由此可见，
该系统是有向的。因此现在只需证明 $\mathcal{F}$ 可以写成凝聚模的滤过
余极限；一旦如此，取这些模在 $\mathcal{F}$ 中的像，便可断定这样的
凝聚子模足够多。

\medskip\noindent
设 $U$ 为仿射概形，$U \to X$ 为满 \'etale 态射。令
$R = U \times_X U$，从而照例有 $X = U/R$。由《空间的性质》命题
\ref{spaces-properties-proposition-quasi-coherent}
可见 $\QCoh(\mathcal{O}_X) = \QCoh(U, R, s, t, c)$。因此问题归结为证明
$\QCoh(U, R, s, t, c)$ 的相应结论。于是本结论由更一般的《群胚》引理
\ref{groupoids-lemma-colimit-coherent} 推出。
\end{proof}

\begin{lemma}
\label{spaces-cohomology-lemma-direct-colimit-finite-presentation}
设 $S$ 为一个概形，$f : X \to Y$ 为 $S$ 上代数空间之间的仿射态射，且
$Y$ 是 Noether 的。则每个拟凝聚 $\mathcal{O}_X$-模都是有限表示
$\mathcal{O}_X$-模的滤过余极限。
\end{lemma}

\begin{proof}
设 $\mathcal{F}$ 为拟凝聚 $\mathcal{O}_X$-模。写成
$f_*\mathcal{F} = \colim \mathcal{H}_i$，其中 $\mathcal{H}_i$ 是凝聚
$\mathcal{O}_Y$-模；见引理 \ref{spaces-cohomology-lemma-directed-colimit-coherent}。由引理
\ref{spaces-cohomology-lemma-coherent-Noetherian}，各 $\mathcal{H}_i$ 都是有限表示的
$\mathcal{O}_Y$-模。因此 $f^*\mathcal{H}_i$ 是有限表示的
$\mathcal{O}_X$-模，见《空间的性质》
\ref{spaces-properties-section-properties-modules} 节。我们断言映射
$$
\colim f^*\mathcal{H}_i = f^*f_*\mathcal{F} \to \mathcal{F}
$$
是满射，因为已假设 $f$ 仿射，即，选取一个概形 $V$ 和满 \'etale 态射
$V \to Y$。% P09-ERR-0030
令 $U = X \times_Y V$。则 $U$ 是概形，$f' : U \to V$ 仿射，且
$U \to X$ 是满 \'etale 态射。由《空间的性质》引理
\ref{spaces-properties-lemma-pushforward-etale-base-change-modules}
可见 $f'_*(\mathcal{F}|_U) = f_*\mathcal{F}|_V$，拉回也同样如此。因此
$f^*f_*\mathcal{F} \to \mathcal{F}$ 在 $U$ 上的限制是映射
$$
f^*f_*\mathcal{F}|_U = (f')^*(f_*\mathcal{F})|_V) =
(f')^*f'_*(\mathcal{F}|_U) \to \mathcal{F}|_U
$$
它是满射，因为 $f'$ 是概形之间的仿射态射。% P09-ERR-0031
所以该断言成立。

\medskip\noindent
由此得出，$X$ 上每个拟凝聚模都是有限表示模的某个滤过余极限的商。
特别地，$\mathcal{F}$ 是某个映射的余核
$$
\colim_{j \in J} \mathcal{G}_j \longrightarrow \colim_{i \in I} \mathcal{H}_i
$$
其中 $\mathcal{G}_j$ 和 $\mathcal{H}_i$ 都有限表示。注意，对于每个
$j \in I$，存在 $i \in I$ 以及态射
$\alpha : \mathcal{G}_j \to \mathcal{H}_i$，使得% P09-ERR-0032
$$
\xymatrix{
\mathcal{G}_j \ar[r]_\alpha \ar[d] & \mathcal{H}_i \ar[d] \\
\colim_{j \in J} \mathcal{G}_j \ar[r] &
\colim_{i \in I} \mathcal{H}_i
}
$$
交换，见引理 \ref{spaces-cohomology-lemma-finite-presentation-quasi-compact-colimit}。
在这种情形下，$\Coker(\alpha)$ 是有限表示的 $\mathcal{O}_X$-模，并带有
一个映射 $\Coker(\alpha) \to \mathcal{F}$。考虑由上述三元组
$(i, j, \alpha)$ 组成的集合 $K$。规定
$(i, j, \alpha) \leq (i', j', \alpha')$，当且仅当
$i \leq i'$、$j \leq j'$，且下图
$$
\xymatrix{
\mathcal{G}_j \ar[r]_\alpha \ar[d] & \mathcal{H}_i \ar[d] \\
\mathcal{G}_{j'} \ar[r]^{\alpha'} &
\mathcal{H}_{i'}
}
$$
交换。由上可知，$K$ 是有向偏序集，
$$
\mathcal{F} = \colim_{(i, j, \alpha) \in K} \Coker(\alpha),
$$
于是结论得证。
\end{proof}






\section{上同调的消失性}
\label{spaces-cohomology-section-vanishing}

\noindent
本节证明：如果一个拟紧且拟分离的代数空间上所有拟凝聚层的高阶上同调
都消失，则该代数空间仿射。我们将通过一系列引理来证明；一旦命题
\ref{spaces-cohomology-proposition-vanishing-affine} 得证，这些引理便都会失去独立作用。

\begin{situation}
\label{spaces-cohomology-situation-vanishing}
这里 $S$ 是一个概形，$X$ 是 $S$ 上拟紧且拟分离的代数空间，并满足下述
性质：对于每个拟凝聚 $\mathcal{O}_X$-模 $\mathcal{F}$，都有
$H^1(X, \mathcal{F}) = 0$。令 $A = \Gamma(X, \mathcal{O}_X)$。
\end{situation}

\noindent
我们希望证明典范态射
$$
p : X \longrightarrow \Spec(A)
$$
（见《空间的性质》引理
\ref{spaces-properties-lemma-morphism-to-affine-scheme}）是同构。
如果 $M$ 是 $A$-模，我们用 $M \otimes_A \mathcal{O}_X$ 表示拟凝聚模
$p^*\tilde M$。

\begin{lemma}
\label{spaces-cohomology-lemma-vanishing-compute}
在情形 \ref{spaces-cohomology-situation-vanishing} 中，对于 $A$-模 $M$，有
$p_*(M \otimes_A \mathcal{O}_X) = \widetilde{M}$ 且
$\Gamma(X, M \otimes_A \mathcal{O}_X) = M$。
\end{lemma}

\begin{proof}
由《空间的态射》引理 \ref{spaces-morphisms-lemma-pushforward}，
$p_*(M \otimes_A \mathcal{O}_X)$ 是 $\Spec(A)$ 上的拟凝聚模，故等式
$p_*(M \otimes_A \mathcal{O}_X) = \widetilde{M}$ 可由等式
$\Gamma(X, M \otimes_A \mathcal{O}_X) = M$ 推出。由引理
\ref{spaces-cohomology-lemma-colimits}，注意到
$\Gamma(X, \bigoplus_{i \in I} \mathcal{O}_X) =
\bigoplus_{i \in I} A$。因此本引理对自由模成立。选取短正合列
$F_1 \to F_0 \to M$，其中 $F_0,F_1$ 是自由 $A$-模。由于
$H^1(X, -)$ 为零，全局截面函子右正合。此外，拉回 $p^*$ 也右正合。
所以
$$
\Gamma(X, F_1 \otimes_A \mathcal{O}_X) \to
\Gamma(X, F_0 \otimes_A \mathcal{O}_X) \to
\Gamma(X, M \otimes_A \mathcal{O}_X) \to 0
$$
正合。结论随即成立。
\end{proof}

\noindent
下面的引理表明，对 $X \to \Spec(A)$ 沿
$\Spec(A') \to \Spec(A)$ 作基变换后，情形 \ref{spaces-cohomology-situation-vanishing}
仍然保持。

\begin{lemma}
\label{spaces-cohomology-lemma-vanishing-base-change}
在情形 \ref{spaces-cohomology-situation-vanishing} 中。
\begin{enumerate}
\item 给定代数空间之间的仿射态射 $X' \to X$，对于每个拟凝聚
$\mathcal{O}_{X'}$-模 $\mathcal{F}'$，都有
$H^1(X', \mathcal{F}') = 0$。
\item 给定 $A$-代数 $A'$，令 $X' = X \times_{\Spec(A)} \Spec(A')$，
则态射 $X' \to X$ 仿射，且 $\Gamma(X', \mathcal{O}_{X'}) = A'$。
\end{enumerate}
\end{lemma}

\begin{proof}
(1) 由引理 \ref{spaces-cohomology-lemma-affine-vanishing-higher-direct-images} 以及 Leray
谱序列（《站点上的上同调》引理 \ref{sites-cohomology-lemma-Leray}）推出。
设 $A \to A'$ 如 (2) 所述。由于仿射态射在基变换下保持（《空间的态射》
引理 \ref{spaces-morphisms-lemma-base-change-affine}），且仿射概形之间的
态射是仿射的，所以 $X' \to X$ 仿射。由引理
\ref{spaces-cohomology-lemma-affine-base-change}，有
$(X' \to X)_*\mathcal{O}_{X'} = A' \otimes_A \mathcal{O}_X$，因此
$$
\Gamma(X', \mathcal{O}_{X'}) =
\Gamma(X, (X' \to X)_*\mathcal{O}_{X'}) =
\Gamma(X, A' \otimes_A \mathcal{O}_X) = A'
$$
由引理 \ref{spaces-cohomology-lemma-vanishing-compute} 可得所需等式。
\end{proof}

\begin{lemma}
\label{spaces-cohomology-lemma-vanishing-separate-closed}
在情形 \ref{spaces-cohomology-situation-vanishing} 中，设 $Z_0, Z_1 \subset |X|$ 为不交的
闭子集。则存在 $a \in A$，使得 $Z_0 \subset V(a)$ 且
$Z_1 \subset V(a - 1)$。
\end{lemma}

\begin{proof}
赋予 $Z_0$、$Z_1$ 约化诱导子空间结构（《空间的性质》定义
\ref{spaces-properties-definition-reduced-induced-space}），并将相应闭浸入
记为 $i_0 : Z_0 \to X$ 和 $i_1 : Z_1 \to X$。由于
$Z_0 \cap Z_1 = \emptyset$，可见拟凝聚 $\mathcal{O}_X$-模的典范映射
$$
\mathcal{O}_X
\longrightarrow
i_{0, *}\mathcal{O}_{Z_0} \oplus i_{1, *}\mathcal{O}_{Z_1}
$$
是满射（考察几何点处的茎即可）。由于 $H^1(X, -)$ 在此映射的核上为零，
所诱导的全局截面映射是满射。因此可以找到 $a \in A$，它映到右端的
全局截面 $(0, 1)$。
\end{proof}

\begin{lemma}
\label{spaces-cohomology-lemma-vanishing-injective}
在情形 \ref{spaces-cohomology-situation-vanishing} 中，态射 $p : X \to \Spec(A)$ 是泛单的。
\end{lemma}

\begin{proof}
设 $A \to k$ 为环同态，其中 $k$ 是域。只需证明
$\Spec(k) \times_{\Spec(A)} X$ 至多有一个点（见《空间的态射》引理
\ref{spaces-morphisms-lemma-universally-injective-local}）。利用引理
\ref{spaces-cohomology-lemma-vanishing-base-change}，可以假设 $A$ 是域，而我们必须证明
$|X|$ 至多有一个点。

\medskip\noindent
把 $X$ 看作 $\Spec(k)$ 上的代数空间，并对任意扩张 $K/k$，用 $X(K)$
表示 $X$ 的 $K$-值点，见《空间的态射》
\ref{spaces-morphisms-section-points-fields} 节。如果 $K/k$ 是超越次数
足够大的代数闭域扩张，则 $X(K) \to |X|$ 是满射，见《空间的态射》引理
\ref{spaces-morphisms-lemma-large-enough}。因此用 $K$ 替换 $k$ 后，只需证明
（在 $A = k$ 的情形下）$X(k)$ 是单点集。

\medskip\noindent
设 $x, x' \in X(k)$。由《良态空间》引理
\ref{decent-spaces-lemma-algebraic-residue-field-extension-closed-point}，
$x$ 和 $x'$ 是 $|X|$ 的闭点。因此，如果 $x \not = x'$，由引理
\ref{spaces-cohomology-lemma-vanishing-separate-closed}，$x$ 和 $x'$ 映到 $\Spec(k)$ 的不同点。
% P09-ERR-0034
由此得到 $x = x'$，如所要求。
\end{proof}

\begin{lemma}
\label{spaces-cohomology-lemma-vanishing-separated}
在情形 \ref{spaces-cohomology-situation-vanishing} 中，态射 $p : X \to \Spec(A)$ 是分离的。
\end{lemma}

\begin{proof}
由《良态空间》引理
\ref{decent-spaces-lemma-there-is-a-scheme-integral-over}
可以找到一个概形 $Y$ 和满整态射 $Y \to X$。整态射是仿射的，所以可以
应用引理 \ref{spaces-cohomology-lemma-vanishing-base-change}，得到对于每个拟凝聚
$\mathcal{O}_Y$-模 $\mathcal{G}$，都有 $H^1(Y, \mathcal{G}) = 0$。
由于 $Y \to X$ 拟紧且 $X$ 拟紧，$Y$ 也拟紧。由于 $Y$ 是概形，可以应用
《概形的上同调》引理
\ref{coherent-lemma-quasi-compact-h1-zero-covering}，得到 $Y$ 仿射，因而
$Y$ 分离。注意，整态射是仿射且泛闭的，见《空间的态射》引理
\ref{spaces-morphisms-lemma-integral-universally-closed}。由《空间的态射》
引理 \ref{spaces-morphisms-lemma-image-universally-closed-separated}，
可见 $X$ 是分离代数空间。
\end{proof}

\begin{proposition}
\label{spaces-cohomology-proposition-vanishing-affine}
\begin{slogan}
代数空间情形下的 Serre 仿射性判别法。
\end{slogan}
一个拟紧且拟分离的代数空间仿射，当且仅当拟凝聚层的所有高阶上同调群
都消失。更确切地说，情形 \ref{spaces-cohomology-situation-vanishing} 中的任意代数空间
都是仿射概形。
\end{proposition}

\begin{proof}
选取一个仿射概形 $U = \Spec(B)$ 和满 \'etale 态射
$\varphi : U \to X$。令 $R = U \times_X U$。由于 $p$ 分离（引理
\ref{spaces-cohomology-lemma-vanishing-separated}），$R$ 是
$U \times_{\Spec(A)} U = \Spec(B \otimes_A B)$ 的闭子概形。因此
$R = \Spec(C)$ 也是仿射的，而且环同态
$$
B \otimes_A B \longrightarrow C
$$
是满射。照例把两个映射记为 $s, t : B \to C$。选取
$g_1, \ldots, g_m \in B$，使得 $s(g_1), \ldots, s(g_m)$ 在
$t : B \to C$ 上生成 $C$（这是可能的，因为 $t : B \to C$ 有限表示，
且所示映射是满射）。于是 $g_1, \ldots, g_m$ 给出
$\varphi_*\mathcal{O}_U$ 的全局截面，并且映射
$$
\mathcal{O}_X[z_1, \ldots, z_n] \longrightarrow \varphi_*\mathcal{O}_U,
\quad
z_j \longmapsto g_j
$$
是满射：限制到 $U$ 上即可检验。% P09-ERR-0033
事实上，由引理 \ref{spaces-cohomology-lemma-flat-base-change-cohomology}，
$\varphi^*\varphi_*\mathcal{O}_U = t_*\mathcal{O}_R$，因此恰好得到
$s(g_i)$ 在 $t : B \to C$ 上生成 $C$ 的条件。由于该映射的核的
$H^1$ 消失，可见
$$
\Gamma(X, \mathcal{O}_X[x_1, \ldots, x_n]) =
A[x_1, \ldots, x_n] \longrightarrow
\Gamma(X, \varphi_*\mathcal{O}_U) = \Gamma(U, \mathcal{O}_U) = B
$$
是满射。因此 $B$ 是有限型 $A$-代数，从而 $X \to \Spec(A)$ 有限型且分离。
由引理 \ref{spaces-cohomology-lemma-vanishing-injective} 以及《空间的态射》引理
\ref{spaces-morphisms-lemma-locally-quasi-finite}，它还是局部拟有限的。
于是由《空间的态射》引理
\ref{spaces-morphisms-lemma-locally-quasi-finite-separated-representable}，
$X \to \Spec(A)$ 可表，从而 $X$ 是概形。最后，应用《概形的上同调》引理
\ref{coherent-lemma-quasi-compact-h1-zero-covering}，可知 $X$ 仿射，
因而等于 $\Spec(A)$。
\end{proof}

\begin{lemma}
\label{spaces-cohomology-lemma-Noetherian-h1-zero}
设 $S$ 为一个概形，$X$ 为 $S$ 上的 Noether 代数空间。假设对于每个凝聚
$\mathcal{O}_X$-模 $\mathcal{F}$，都有 $H^1(X, \mathcal{F}) = 0$。
则 $X$ 是仿射概形。
\end{lemma}

\begin{proof}
由引理 \ref{spaces-cohomology-lemma-directed-colimit-coherent} 和 \ref{spaces-cohomology-lemma-colimits}，
该假设推出对于每个拟凝聚 $\mathcal{O}_X$-模 $\mathcal{F}$，都有
$H^1(X, \mathcal{F}) = 0$。于是由命题
\ref{spaces-cohomology-proposition-vanishing-affine}，$X$ 仿射。
\end{proof}

\begin{lemma}
\label{spaces-cohomology-lemma-Noetherian-h1-zero-invertible}
设 $S$ 为一个概形，$X$ 为 $S$ 上的 Noether 代数空间，并设
$\mathcal{L}$ 为可逆 $\mathcal{O}_X$-模。假设对于每个凝聚
$\mathcal{O}_X$-模 $\mathcal{F}$，存在 $n \geq 1$ 使得
$H^1(X, \mathcal{F} \otimes_{\mathcal{O}_X} \mathcal{L}^{\otimes n}) = 0$。
则 $X$ 是概形，而且 $\mathcal{L}$ 在 $X$ 上是丰沛的。
\end{lemma}

\begin{proof}
设 $s \in H^0(X, \mathcal{L}^{\otimes d})$ 为全局截面。设 $U \subset X$
为 $s$ 生成 $\mathcal{L}^{\otimes d}$ 的开子空间。特别地，有
$\mathcal{L}^{\otimes d}|_U \cong \mathcal{O}_U$。我们断言 $U$ 仿射。

\medskip\noindent
断言的证明。我们将证明，对于每个拟凝聚 $\mathcal{O}_U$-模
$\mathcal{F}$，都有 $H^1(U, \mathcal{F}) = 0$。由命题
\ref{spaces-cohomology-proposition-vanishing-affine}，这将证明该断言。记
$j : U \to X$ 为包含态射。由于 $j$ 在 \'etale 拓扑下局部仿射（由
《态射》引理 \ref{morphisms-lemma-affine-s-open}），《空间的态射》引理
\ref{spaces-morphisms-lemma-affine-local} 表明 $j$ 仿射。因此
$$
H^1(U, \mathcal{F}) = H^1(X, j_*\mathcal{F})
$$
由引理 \ref{spaces-cohomology-lemma-affine-vanishing-higher-direct-images}（以及《站点上的
上同调》引理 \ref{sites-cohomology-lemma-apply-Leray}）成立。把
$j_*\mathcal{F} = \colim \mathcal{F}_i$ 写成凝聚 $\mathcal{O}_X$-模的
滤过余极限，见引理 \ref{spaces-cohomology-lemma-directed-colimit-coherent}。则
$$
H^1(X, j_*\mathcal{F}) = \colim H^1(X, \mathcal{F}_i)
$$
由引理 \ref{spaces-cohomology-lemma-colimits} 成立。因此只需证明 $H^1(X, \mathcal{F}_i)$
在 $H^1(U, j^*\mathcal{F}_i)$ 中的像为零。由假设，存在 $n \geq 1$，使得
$$
H^1(X,
\mathcal{F}_i \otimes_{\mathcal{O}_X}
(\mathcal{O}_X \oplus \mathcal{L} \oplus \ldots
\oplus \mathcal{L}^{\otimes d - 1})
\otimes_{\mathcal{O}_X} \mathcal{L}^{\otimes n}) = 0
$$
因此存在 $a \geq 0$，使得
$H^1(X, \mathcal{F}_i \otimes_{\mathcal{O}_X} \mathcal{L}^{\otimes ad}) = 0$。
另一方面，映射
$$
s^a : \mathcal{F}_i \longrightarrow
\mathcal{F}_i \otimes_{\mathcal{O}_X} \mathcal{L}^{\otimes ad}
$$
限制到 $U$ 后是同构。考察交换图
$$
\xymatrix{
H^1(X, \mathcal{F}_i) \ar[r] \ar[d]_{s^a} & H^1(U, j^*\mathcal{F}_i)
\ar[d]^{\cong} \\
H^1(X, \mathcal{F}_i \otimes_{\mathcal{O}_X} \mathcal{L}^{\otimes ad}) \ar[r] &
H^1(U,
j^*(\mathcal{F}_i \otimes_{\mathcal{O}_X} \mathcal{L}^{\otimes ad}))
}
$$
可知映射 $H^1(X, \mathcal{F}_i) \to H^1(U, j^*\mathcal{F}_i)$ 为零，
断言得证。

\medskip\noindent
设 $x \in |X|$ 为闭点。由《良态空间》引理
\ref{decent-spaces-lemma-decent-space-closed-point}，可以用闭浸入
$i : \Spec(k) \to X$ 表示 $x$（这里还用到了拟分离代数空间是良态的，
见《良态空间》\ref{decent-spaces-section-reasonable-decent} 节）。因此
$\mathcal{O}_X \to i_*\mathcal{O}_{\Spec(k)}$ 是满射。设
$\mathcal{I} \subset \mathcal{O}_X$ 为其核，并选取 $d \geq 1$，使得
$H^1(X, \mathcal{I} \otimes_{\mathcal{O}_X} \mathcal{L}^{\otimes d}) = 0$。
则
$$
H^0(X, \mathcal{L}^{\otimes d}) \to
H^0(X,
i_*\mathcal{O}_{\Spec(k)} \otimes_{\mathcal{O}_X} \mathcal{L}^{\otimes d}) =
H^0(\Spec(k), i^*\mathcal{L}^{\otimes d}) \cong k
$$
由上同调长正合列，该映射是满射。因此存在
$s \in H^0(X, \mathcal{L}^{\otimes d})$，使得 $x \in U$，其中 $U$ 是上面
与 $s$ 对应的开子空间。由上述断言，$x$ 位于 $X$ 的概形轨迹中
（见《空间的性质》引理 \ref{spaces-properties-lemma-subscheme}）。

\medskip\noindent
为了得出 $X$ 是概形，只需证明 $|X|$ 中任何包含所有闭点的开子集都等于
$|X|$。这由 $|X|$ 是 Noether 拓扑空间这一事实推出，见《空间的性质》
引理 \ref{spaces-properties-lemma-Noetherian-sober}。最后，当 $X$ 是概形时，
可以应用《概形的上同调》引理
\ref{coherent-lemma-quasi-compact-h1-zero-invertible}，得到
$\mathcal{L}$ 是丰沛的。
\end{proof}









\section{有限态射与仿射性}
\label{spaces-cohomology-section-finite-affine}

\noindent
本节对应于《概形的上同调》\ref{coherent-section-finite-affine} 节。

\begin{lemma}
\label{spaces-cohomology-lemma-finite-morphism-Noetherian}
设 $S$ 为一个概形，$f : Y \to X$ 为 $S$ 上代数空间之间的态射。假设
$f$ 有限且为满射，并且 $X$ 局部 Noether。设 $i : Z \to X$ 为闭浸入。
记 $i' : Z' \to Y$ 为 $Z$ 的逆像（《空间的态射》
\ref{spaces-morphisms-section-closed-immersions} 节），并记
$f' : Z' \to Z$ 为诱导态射。则
$\mathcal{G} = f'_*\mathcal{O}_{Z'}$ 是支撑为 $Z$ 的凝聚
$\mathcal{O}_Z$-模。
\end{lemma}

\begin{proof}
注意，$f'$ 是 $f$ 的基变换，因而由《空间的态射》引理
\ref{spaces-morphisms-lemma-base-change-surjective} 和
\ref{spaces-morphisms-lemma-base-change-integral}，它有限且为满射。
由《空间的态射》引理
\ref{spaces-morphisms-lemma-locally-finite-type-locally-noetherian}（以及闭浸入
和有限态射都是有限型的这一事实），$Y$、$Z$、$Z'$ 都局部 Noether。
由引理 \ref{spaces-cohomology-lemma-finite-pushforward-coherent}，$\mathcal{G}$ 是凝聚
$\mathcal{O}_Z$-模。$\mathcal{G}$ 的支撑在 $|Z|$ 中闭，见《空间的态射》
引理 \ref{spaces-morphisms-lemma-support-finite-type}。因此，如果
$\mathcal{G}$ 的支撑不等于 $|Z|$，则用一个开子空间替换 $X$ 后，可以
假设 $\mathcal{G} = 0$ 而 $Z \not = \emptyset$。这将意味着
$f'_*\mathcal{O}_{Z'} = 0$。特别地，截面
$1 \in \Gamma(Z', \mathcal{O}_{Z'}) = \Gamma(Z, f'_*\mathcal{O}_{Z'})$
将为零，从而 $Z' = \emptyset$ 是空代数空间。这不可能，因为
$Z' \to Z$ 是满射。
\end{proof}

\begin{lemma}
\label{spaces-cohomology-lemma-affine-morphism-projection-ideal}
设 $S$ 为一个概形，$f : Y \to X$ 为 $S$ 上代数空间之间的态射。
设 $\mathcal{F}$ 为 $Y$ 上的拟凝聚层，$\mathcal{I}$ 为 $X$ 上的拟凝聚
理想层。如果 $f$ 仿射，则
$\mathcal{I}f_*\mathcal{F} = f_*(f^{-1}\mathcal{I}\mathcal{F})$
（记号的含义见证明）。
\end{lemma}

\begin{proof}
上述记号的含义如下。由于 $f^{-1}$ 是正合函子，
$f^{-1}\mathcal{I}$ 是 $f^{-1}\mathcal{O}_X$ 的理想层。通过
$Y_\etale$ 上的映射
$f^\sharp : f^{-1}\mathcal{O}_X \to \mathcal{O}_Y$，它作用在
$\mathcal{F}$ 上。于是 $f^{-1}\mathcal{I}\mathcal{F}$ 是由形如 $as$ 的
局部截面的和生成的子层，其中 $a$ 是 $f^{-1}\mathcal{I}$ 的局部截面，
$s$ 是 $\mathcal{F}$ 的局部截面。它是 $\mathcal{F}$ 的拟凝聚
$\mathcal{O}_Y$-子模，因为它也等于自然映射
$f^*\mathcal{I} \otimes_{\mathcal{O}_Y} \mathcal{F} \to \mathcal{F}$ 的像。

\medskip\noindent
说明这些记号后，证明便很直接。事实上，该问题在 $X$ 上是 \'etale 局部的，
因而可以假设 $X$ 是仿射概形。在这种情形下，结论由概形的相应结论推出，
见《概形的上同调》引理
\ref{coherent-lemma-affine-morphism-projection-ideal}。
\end{proof}

\begin{lemma}
\label{spaces-cohomology-lemma-image-affine-finite-morphism-affine-Noetherian}
设 $S$ 为一个概形，$f : Y \to X$ 为 $S$ 上代数空间之间的态射。假设：
\begin{enumerate}
\item $f$ 有限；
\item $f$ 为满射；
\item $Y$ 仿射；
\item $X$ 是 Noether 的。
\end{enumerate}
则 $X$ 仿射。
\end{lemma}

\begin{proof}
我们将证明，在本引理的假设下，对于任意凝聚 $\mathcal{O}_X$-模
$\mathcal{F}$，都有 $H^1(X, \mathcal{F}) = 0$。由引理
\ref{spaces-cohomology-lemma-directed-colimit-coherent} 和 \ref{spaces-cohomology-lemma-colimits}，这又推出
对于每个拟凝聚 $\mathcal{O}_X$-模 $\mathcal{F}$，都有
$H^1(X, \mathcal{F}) = 0$。于是由命题
\ref{spaces-cohomology-proposition-vanishing-affine}，$X$ 仿射。

\medskip\noindent
设 $\mathcal{P}$ 为 $X$ 上凝聚层 $\mathcal{F}$ 的下述性质：
$$
\mathcal{P}(\mathcal{F}) \Leftrightarrow H^1(X, \mathcal{F}) = 0.
$$
我们将应用引理 \ref{spaces-cohomology-lemma-property-higher-rank-cohomological}，因此必须为
$\mathcal{P}$ 验证该引理的 (1)、(2)、(3)。性质 (1) 由与层的短正合列
相伴的上同调长正合列推出。性质 (2) 来自 $H^1(X, -)$ 是可加函子。
为证明 (3)，设 $i : Z \to X$ 为满足 $|Z|$ 不可约的约化闭子空间。
设 $i' : Z' \to Y$ 和 $f' : Z' \to Z$ 如引理
\ref{spaces-cohomology-lemma-finite-morphism-Noetherian} 中所述，并令
$\mathcal{G} = f'_*\mathcal{O}_{Z'}$。我们断言 $\mathcal{G}$ 满足引理
\ref{spaces-cohomology-lemma-property-higher-rank-cohomological} 的性质 (3)(a) 和 (3)(b)，
这将完成证明。性质 (3)(a) 已在引理
\ref{spaces-cohomology-lemma-finite-morphism-Noetherian} 中证明。为证明 (3)(b)，设
$\mathcal{I}$ 为 $Z$ 上的非零拟凝聚理想层。记
$\mathcal{I}' \subset \mathcal{O}_{Z'}$ 为拟凝聚理想
$(f')^{-1}\mathcal{I} \mathcal{O}_{Z'}$，也就是
$(f')^*\mathcal{I} \to \mathcal{O}_{Z'}$ 的像。由引理
\ref{spaces-cohomology-lemma-affine-morphism-projection-ideal}，有
$f_*\mathcal{I}' = \mathcal{I} \mathcal{G}$。% P09-ERR-0035
我们断言其共同值
$\mathcal{G}' = \mathcal{I} \mathcal{G} = f'_*\mathcal{I}'$
满足 (3)(b) 中的条件。首先，显然 $\mathcal{G}/\mathcal{G}'$ 的支撑包含在
$\mathcal{O}_Z/\mathcal{I}$ 的支撑中；后者是 $|Z|$ 的真子空间，因为
$\mathcal{I}$ 是约化不可约代数空间 $Z$ 上的非零理想层。态射 $f'$ 仿射，
故由引理 \ref{spaces-cohomology-lemma-affine-vanishing-higher-direct-images}，
$R^1f'_*\mathcal{I}' = 0$。由于 $Z'$ 仿射（它是仿射概形的闭子概形），
有 $H^1(Z', \mathcal{I}') = 0$。因此 Leray 谱序列（取《站点上的上同调》
引理 \ref{sites-cohomology-lemma-apply-Leray} 的形式）推出
$H^1(Z, f'_*\mathcal{I}') = 0$。由于 $i : Z \to X$ 仿射，
$R^1i_*f'_*\mathcal{I}' = 0$；再次应用 Leray，得到
$H^1(X, i_*f'_*\mathcal{I}') = 0$。换言之，如所要求，有
$H^1(X, i_*\mathcal{G}') = 0$。
\end{proof}








\section{Chow 引理的弱形式}
\label{spaces-cohomology-section-weak-chow}

\noindent
本节快速证明下面的引理，以帮助我们证明恰当代数空间上凝聚模上同调的
基本结果。

\begin{lemma}
\label{spaces-cohomology-lemma-weak-chow}
设 $A$ 为一个环，$X$ 为 $\Spec(A)$ 上的代数空间，其结构态射
$X \to \Spec(A)$ 分离且为有限型。则存在恰当满态射 $X' \to X$，
其中 $X'$ 是 $\Spec(A)$ 上 H-拟射影的概形。
\end{lemma}

\begin{proof}
设 $W$ 为仿射概形，$f : W \to X$ 为满 \'etale 态射。存在整数 $d$，使得
$f$ 的所有几何纤维都至多有 $d$ 个点（因为 $X$ 是分离代数，因而合理，
见《良态空间》引理 \ref{decent-spaces-lemma-bounded-fibres}）。
% P09-ERR-0036
取最小的 $d$，便得到非空开集 $U \subset X$，使得
$f^{-1}(U) \to U$ 是次数为 $d$ 的有限 \'etale 态射，见《良态空间》引理
\ref{decent-spaces-lemma-quasi-compact-reasonable-stratification}。令
$$
V \subset W \times_X W \times_X \ldots \times_X W
$$
为所有对角线的补集（纤维积中有 $d$ 个因子）。由于 $W \to X$ 分离，
对角态射 $W \to W \times_X W$ 是闭浸入。由于 $W \to X$ 是 \'etale 的，
该对角态射也是开浸入，见《空间的态射》引理
\ref{spaces-morphisms-lemma-etale-unramified} 和
\ref{spaces-morphisms-lemma-diagonal-unramified-morphism}。因此这些对角线是
拟紧概形 $W \times_X \ldots \times_X W$ 的开闭子概形。特别地，$V$
是拟紧概形。选取开浸入 $W \subset Y$，其中 $Y$ 在 $A$ 上 H-射影
（这是可能的，因为 $W$ 仿射且在 $A$ 上为有限型；例如可以使用《态射》引理
\ref{morphisms-lemma-quasi-affine-finite-type-over-S} 和
\ref{morphisms-lemma-H-quasi-projective-open-H-projective}）。令
$$
Z \subset Y \times_A Y \times_A \ldots \times_A Y
$$
为复合态射
$V \to W \times_X \ldots \times_X W \to Y \times_A \ldots \times_A Y$
的概形论像。由于 $V$ 拟紧且 $Y \times_A \ldots \times_A Y$ 分离，
该态射拟紧。注意，$V \to Y \times_A \ldots \times_A Y$ 是浸入，故
$V \to Z$ 是开浸入，见《态射》引理
\ref{morphisms-lemma-quasi-compact-immersion}。投影态射给出 $d$ 个态射
$g_i : Z \to Y$。由于 $Y$ 在 $A$ 上射影，这些态射 $g_i$ 都是射影的，
见《态射》\ref{morphisms-section-projective} 节中的材料。令
$$
X' = \bigcup g_i^{-1}(W) \subset Z
$$
存在态射 $X' \to X$，它在 $g_i^{-1}(W)$ 上的限制是复合态射
$g_i^{-1}(W) \to W \to X$。事实上，这些态射在 $V$ 上相同，因而由
《空间的态射》引理 \ref{spaces-morphisms-lemma-equality-of-morphisms}，
它们在 $g_i^{-1}(W) \cap g_j^{-1}(W)$ 上相同。断言：态射
$X' \to X$ 恰当。

\medskip\noindent
如果断言成立，则对 $d$ 归纳即可证明本引理。事实上，按构造，$X'$ 在
$\Spec(A)$ 上 H-拟射影。由于 $V$ 满射到 $U$，$X' \to X$ 的像包含开集
$U$。记 $T$ 为 $X \setminus U$ 上的约化诱导代数空间结构。则
$T \times_X W$ 是 $W$ 的闭子概形，因而仿射。此外，态射
$T \times_X W \to T$ 是 \'etale 的，且每个几何纤维少于 $d$ 个点。
由归纳假设，存在恰当满态射 $T' \to T$，其中 $T'$ 是
$\Spec(A)$ 上 H-拟射影的概形。由于 $T$ 是 $X$ 的闭子空间，
$T' \to X$ 是恰当态射。因此取恰当满态射
$X' \amalg T' \to X$ 即得本引理。

\medskip\noindent
断言的证明。按构造，态射 $X' \to X$ 分离且为有限型。我们将为态射
$V \to X'$ 和 $X' \to X$ 验证《空间的态射》引理
\ref{spaces-morphisms-lemma-refined-valuative-criterion-universally-closed}
的条件 (1)--(4)。条件 (1) 和 (2) 已在上面验证。条件 (3) 成立，因为
$X' \to X$ 分离（其源是分离代数空间）。因此只需对图
$$
\xymatrix{
\Spec(K) \ar[r] \ar[d] & V \ar[d] \\
\Spec(R) \ar[r] & X
}
$$
检验到 $X'$ 的可提升性，其中 $R$ 是分式域为 $K$ 的赋值环。注意，上方
水平态射由 $W$ 的 $d$ 个两两不同的 $K$-值点
$w_1, \ldots, w_d$ 给出。事实上，它们恰好是该图所给点
$x \in X(K)$ 的全部逆像。由于 $W \to X$ 为满射，必要时用赋值环扩张
替换 $R$ 后，可以把态射 $\Spec(R) \to X$ 提升为态射
$w : \Spec(R) \to W$，见《空间的态射》引理
\ref{spaces-morphisms-lemma-lift-valuation-ring-through-flat-morphism}。
由于 $w_1, \ldots, w_d$ 是 $x$ 的全部逆像，我们看到
$w|_{\Spec(K)}$ 等于其中某一个，比如 $w_i$。% P09-ERR-0037
于是得到交换图
$$
\xymatrix{
\Spec(K) \ar[r] \ar[d] &  Z \ar[d]_{g_i}\\
\Spec(R) \ar[r]^w & Y
}
$$
由射影态射 $g_i$ 的恰当性赋值判别法，可以把 $w$ 提升为
$z : \Spec(R) \to Z$，见《态射》引理
\ref{morphisms-lemma-locally-projective-proper} 以及《概形》命题
\ref{schemes-proposition-characterize-universally-closed}。$z$ 的像位于
$g_i^{-1}(W) \subset X'$ 中，证明完成。
\end{proof}







\section{Noether 赋值判别法}
\label{spaces-cohomology-section-Noetherian-valuative-criterion}

\noindent
我们使用离散赋值环证明恰当性赋值判别法的一个版本。更精确（因而也更
技术化）的版本见《空间的极限》
\ref{spaces-limits-section-Noetherian-valuative-criterion} 节。

\begin{lemma}
\label{spaces-cohomology-lemma-check-separated-dvr}
设 $S$ 为一个概形，$f : X \to Y$ 为 $S$ 上代数空间之间的态射。假设：
\begin{enumerate}
\item $Y$ 局部 Noether；
\item $f$ 局部有限型且拟分离；
\item 对于每个交换图
$$
\xymatrix{
\Spec(K) \ar[r] \ar[d] & X \ar[d] \\
\Spec(A) \ar[r] \ar@{-->}[ru] & Y
}
$$
其中 $A$ 是离散赋值环，$K$ 是其分式域，使该图交换的虚线箭头至多有一个。
\end{enumerate}
则 $f$ 分离。
\end{lemma}

\begin{proof}
必须证明对角态射 $\Delta : X \to X \times_Y X$ 是闭浸入。我们已经知道
$\Delta$ 可表、分离，是单态射且局部有限型，见《空间的态射》引理
\ref{spaces-morphisms-lemma-properties-diagonal}。选取仿射概形 $U$ 和
\'etale 态射 $U \to X \times_Y X$。令
$V = X \times_{\Delta, X \times_Y X} U$。只需证明 $V \to U$ 是闭浸入
（《空间的态射》引理 \ref{spaces-morphisms-lemma-closed-immersion-local}）。
由于 $X \times_Y X$ 在 $Y$ 上局部有限型，$U$ 是 Noether 的（使用
《空间的态射》引理
\ref{spaces-morphisms-lemma-composition-finite-type},
\ref{spaces-morphisms-lemma-base-change-finite-type} 和
\ref{spaces-morphisms-lemma-locally-finite-type-locally-noetherian}）。
由于 $\Delta$ 可表，$V$ 是概形。此外，因为 $f$ 拟分离，$V$ 拟紧。
所以 $V \to U$ 为有限型。考虑概形态射的交换图
$$
\xymatrix{
\Spec(K) \ar[r] \ar[d] & V \ar[d] \\
\Spec(A) \ar[r] \ar@{-->}[ru] & U
}
$$
其中 $A$ 是分式域为 $K$ 的离散赋值环。可以把复合态射
$\Spec(A) \to U \to X \times_Y X$ 解释为一对态射
$a, b : \Spec(A) \to X$；它们复合到 $Y$ 后相同，且限制到
$\Spec(K)$ 后相等。因此假设 (3) 保证 $a = b$，从而得到图中的虚线箭头。
由《极限》引理 \ref{limits-lemma-Noetherian-dvr-valuative-proper}，
$V \to U$ 恰当。换言之，$\Delta$ 恰当。由于 $\Delta$ 是单态射，
《\'Etale 态射》引理 \ref{etale-lemma-characterize-closed-immersion}
表明 $\Delta$ 是闭浸入，如所要求。
\end{proof}

\begin{lemma}
\label{spaces-cohomology-lemma-check-proper-dvr}
设 $S$ 为一个概形，$f : X \to Y$ 为 $S$ 上代数空间之间的态射。假设：
\begin{enumerate}
\item $Y$ 局部 Noether；
\item $f$ 有限型且拟分离；
\item 对于每个交换图
$$
\xymatrix{
\Spec(K) \ar[r] \ar[d] & X \ar[d] \\
\Spec(A) \ar[r] \ar@{-->}[ru] & Y
}
$$
其中 $A$ 是离散赋值环，$K$ 是其分式域，存在唯一的虚线箭头使该图交换。
\end{enumerate}
则 $f$ 恰当。
\end{lemma}

\begin{proof}
由引理 \ref{spaces-cohomology-lemma-check-separated-dvr}，$f$ 分离，所以只需证明 $f$ 泛闭。
为此可以在 $Y$ 上 \'etale 局部地验证（《空间的态射》引理
\ref{spaces-morphisms-lemma-universally-closed-local}）。因此可以假设
$Y = \Spec(A)$ 是 Noether 仿射概形。按照 Chow 引理的弱形式（引理
\ref{spaces-cohomology-lemma-weak-chow}）选取 $X' \to X$。我们断言
$X' \to \Spec(A)$ 泛闭。由《空间的态射》引理
\ref{spaces-morphisms-lemma-image-proper-is-proper}，该断言蕴含本引理。
为证明该断言，根据《极限》引理
\ref{limits-lemma-check-universally-closed-Noetherian}，只需证明在每个实线
部分交换的图
$$
\xymatrix{
\Spec(K) \ar[r] \ar[d] & X' \ar[r] & X \ar[d] \\
\Spec(A) \ar[rr] \ar@{-->}[ru]^a \ar@{-->}[rru]_b & & Y
}
$$
中，可以找到虚线箭头 $a$，其中 $A$ 是分式域为 $K$ 的离散赋值环。
% P09-ERR-0038
由假设可以找到虚线箭头 $b$。于是态射
$X' \times_{X, b} \Spec(A) \to \Spec(A)$ 是概形之间的恰当态射；
由概形态射的赋值判别法，可以把 $b$ 提升为所需态射 $a$。
\end{proof}

\begin{remark}[完备离散赋值环的变体]
\label{spaces-cohomology-remark-variant}
在引理 \ref{spaces-cohomology-lemma-check-separated-dvr} 和 \ref{spaces-cohomology-lemma-check-proper-dvr} 中，
只需考虑完备离散赋值环。确切地说，在引理
\ref{spaces-cohomology-lemma-check-separated-dvr} 中，可以把条件 (3) 替换为下述条件：
给定任意交换图
$$
\xymatrix{
\Spec(K) \ar[r] \ar[d] & X \ar[d] \\
\Spec(A) \ar[r] \ar@{-->}[ru] & Y
}
$$
其中 $A$ 是分式域为 $K$ 的完备离散赋值环，使该图交换的虚线箭头至多
有一个。事实上，给定引理 \ref{spaces-cohomology-lemma-check-separated-dvr} (3) 中的任意图，
其完备化 $A^\wedge$ 是离散赋值环（《代数进阶》引理
\ref{more-algebra-lemma-completion-dvr}），而箭头
$\Spec(A^\wedge) \to X$ 的唯一性蕴含箭头 $\Spec(A) \to X$ 的唯一性；
例如这由《空间的性质》命题
\ref{spaces-properties-proposition-sheaf-fpqc} 推出。
类似地，在引理 \ref{spaces-cohomology-lemma-check-proper-dvr} 中，可以把条件 (3) 替换为
下述条件：给定任意交换图
$$
\xymatrix{
\Spec(K) \ar[r] \ar[d] & X \ar[d] \\
\Spec(A) \ar[r] & Y
}
$$
其中 $A$ 是分式域为 $K$ 的完备离散赋值环，则存在完备离散赋值环的扩张
$A \subset A'$，诱导分式域扩张 $K \subset K'$，并存在唯一箭头
$\Spec(A') \to X$ 使图
$$
\xymatrix{
\Spec(K') \ar[r] \ar[d] & \Spec(K) \ar[r] & X \ar[d] \\
\Spec(A') \ar[r] \ar[rru] & \Spec(A) \ar[r] & Y
}
$$
交换。事实上，给定引理 \ref{spaces-cohomology-lemma-check-proper-dvr} (3) 中的任意图，
只要对于离散赋值环的{\it 任意}扩张 $A \subset B$，存在交换图
$$
\xymatrix{
\Spec(L) \ar[r] \ar[d] & \Spec(K) \ar[r] & X \ar[d] \\
\Spec(B) \ar[r] \ar[rru] & \Spec(A) \ar[r] & Y
}
$$
就能推出存在箭头 $\Spec(A) \to X$ 嵌入该图。这已在《空间的态射》引理
\ref{spaces-morphisms-lemma-push-down-solution} 中证明。事实上，从这些考察
可以看出，只要某类离散赋值环满足：任给一个离散赋值环，都有它的一个
扩张属于该类，那么只在该类离散赋值环上的图中寻找虚线箭头就已足够。
例如，可以取剩余域代数闭的完备离散赋值环。
\end{remark}







\section{凝聚层的高阶正像}
\label{spaces-cohomology-section-proper-pushforward}

\noindent
本节证明一个基本事实：凝聚层在恰当态射下的高阶正像是凝聚的。首先证明
一个辅助引理。

\begin{lemma}
\label{spaces-cohomology-lemma-kill-by-twisting}
设 $S$ 为一个概形。考虑 $S$ 上代数空间的交换图
$$
\xymatrix{
X \ar[r]_i \ar[rd]_f & \mathbf{P}^n_Y \ar[d] \\
& Y
}
$$
假设 $i$ 是闭浸入，且 $Y$ 是 Noether 的。令
$\mathcal{L} = i^*\mathcal{O}_{\mathbf{P}^n_Y}(1)$。设 $\mathcal{F}$ 为
$X$ 上的凝聚模。则存在整数 $d_0$，使得对于所有 $d \geq d_0$ 和
所有 $p > 0$，都有
$R^pf_*(\mathcal{F} \otimes_{\mathcal{O}_X} \mathcal{L}^{\otimes d}) = 0$。
\end{lemma}

\begin{proof}
可以在 $Y$ 上 \'etale 局部地检验
$R^pf_*(\mathcal{F} \otimes \mathcal{L}^{\otimes d})$ 是否为零，见公式
(\ref{spaces-cohomology-equation-representable-higher-direct-image})。因此可以假设 $Y$ 是
Noether 环的谱。在这种情形下，$X$ 是概形，结论由《概形的上同调》引理
\ref{coherent-lemma-kill-by-twisting} 推出。
\end{proof}

\begin{lemma}
\label{spaces-cohomology-lemma-proper-pushforward-coherent}
设 $S$ 为一个概形，$f : X \to Y$ 为 $S$ 上代数空间之间的恰当态射，
且 $Y$ 局部 Noether。设 $\mathcal{F}$ 为凝聚 $\mathcal{O}_X$-模。
则对于所有 $i \geq 0$，$R^if_*\mathcal{F}$ 都是凝聚
$\mathcal{O}_Y$-模。
\end{lemma}

\begin{proof}
首先，由《空间的态射》引理
\ref{spaces-morphisms-lemma-locally-finite-type-locally-noetherian}，
$X$ 是局部 Noether 代数空间，故引理的陈述有意义。此外，计算
$R^if_*\mathcal{F}$ 与 $Y$ 上的 \'etale 局部化交换（《空间的性质》引理
\ref{spaces-properties-lemma-pushforward-etale-base-change-modules}），而且
可以在 $Y$ 上 \'etale 局部地检验 $R^if_*\mathcal{F}$ 是否凝聚（引理
\ref{spaces-cohomology-lemma-coherent-Noetherian}）。因此可以假设
$Y = \Spec(A)$ 是 Noether 仿射概形。

\medskip\noindent
假设 $Y = \Spec(A)$ 是仿射概形。注意，$f$ 局部有限表示（《空间的态射》
引理 \ref{spaces-morphisms-lemma-noetherian-finite-type-finite-presentation}）。
因此它有限表示，从而 $X$ 是 Noether 的（《空间的态射》引理
\ref{spaces-morphisms-lemma-finite-presentation-noetherian}）。于是引理
\ref{spaces-cohomology-lemma-property-higher-rank-cohomological-variant} 可应用于 $X$ 的
凝聚模范畴。对于 $X$ 上的凝聚层 $\mathcal{F}$，规定：$\mathcal{P}$
成立，当且仅当 $R^if_*\mathcal{F}$ 是 $\Spec(A)$ 上的凝聚模。
下面证明引理 \ref{spaces-cohomology-lemma-property-higher-rank-cohomological-variant} 的
条件 (1)、(2)、(3) 对此性质成立，从而完成本引理的证明。

\medskip\noindent
验证条件 (1)。设
$$
0 \to \mathcal{F}_1 \to \mathcal{F}_2 \to \mathcal{F}_3 \to 0
$$
为 $X$ 上凝聚层的短正合列。考虑高阶正像的长正合列
$$
R^{p - 1}f_*\mathcal{F}_3 \to
R^pf_*\mathcal{F}_1 \to
R^pf_*\mathcal{F}_2 \to
R^pf_*\mathcal{F}_3 \to
R^{p + 1}f_*\mathcal{F}_1
$$
显然，如果三个层 $\mathcal{F}_i$ 中的两个具有性质 $\mathcal{P}$，
则第三个的高阶正像在这个正合复形中夹在两个凝聚层之间。因此由引理
\ref{spaces-cohomology-lemma-coherent-abelian-Noetherian} 和
\ref{spaces-cohomology-lemma-coherent-Noetherian-quasi-coherent-sub-quotient}，这些高阶正像
也凝聚。所以第三个层也具有性质 $\mathcal{P}$。

\medskip\noindent
验证条件 (2)。这立即来自等式
$R^if_*(\mathcal{F}_1 \oplus \mathcal{F}_2) =
R^if_*\mathcal{F}_1 \oplus R^if_*\mathcal{F}_2$，以及凝聚模的直和项凝聚
这一事实（见上面引用的引理）。

\medskip\noindent
验证条件 (3)。设 $i : Z \to X$ 为闭浸入，其中 $Z$ 约化且 $|Z|$ 不可约。
令 $g = f \circ i : Z \to \Spec(A)$。设 $\mathcal{G}$ 为 $Z$ 上概形论
支撑等于 $Z$ 的凝聚模，并且对于所有 $p$，$R^pg_*\mathcal{G}$ 凝聚。
则 $\mathcal{F} = i_*\mathcal{G}$ 是 $X$ 上概形论支撑为 $Z$ 的凝聚模，
且 $R^pf_*\mathcal{F} = R^pg_*\mathcal{G}$。为看出这一点，使用 Leray
谱序列（《站点上的上同调》引理
\ref{sites-cohomology-lemma-relative-Leray}），并使用引理
\ref{spaces-cohomology-lemma-affine-vanishing-higher-direct-images} 以及闭浸入仿射这一事实
（《空间的态射》引理 \ref{spaces-morphisms-lemma-closed-immersion-affine}），
得到当 $q > 0$ 时 $R^qi_*\mathcal{G} = 0$。因此问题归结为寻找
$Z$ 上支撑等于 $Z$ 的凝聚层
$\mathcal{G}$，使得对于所有 $p$，$R^pg_*\mathcal{G}$ 凝聚。

\medskip\noindent
对态射 $Z \to \Spec(A)$ 应用引理 \ref{spaces-cohomology-lemma-weak-chow}，得到图
$$
\xymatrix{
Z \ar[rd]_g & Z' \ar[d]^-{g'} \ar[l]^\pi \ar[r]_i & \mathbf{P}^n_A \ar[dl] \\
& \Spec(A) &
}
$$
其中 $\pi : Z' \to Z$ 恰当且为满射，$i$ 是浸入。由于
$Z \to \Spec(A)$ 恰当，由《空间的态射》引理
\ref{spaces-morphisms-lemma-composition-proper}，$g'$ 恰当。因此 $i$ 是
闭浸入（《空间的态射》引理
\ref{spaces-morphisms-lemma-universally-closed-permanence} 和
\ref{spaces-morphisms-lemma-immersion-when-closed}）。由此可知态射
$i' = (i, \pi) : \mathbf{P}^n_A \times_{\Spec(A)} Z' = \mathbf{P}^n_Z$ 为
闭浸入（《空间的态射》引理
\ref{spaces-morphisms-lemma-semi-diagonal}）。% P09-ERR-0039
令
$$
\mathcal{L} =
i^*\mathcal{O}_{\mathbf{P}^n_A}(1) =
(i')^*\mathcal{O}_{\mathbf{P}^n_Z}(1)
$$
可以分别对 $\mathcal{L}$ 与 $\pi$、$\mathcal{L}$ 与 $g'$ 应用引理
\ref{spaces-cohomology-lemma-kill-by-twisting}。因此对于所有 $d \gg 0$ 和 $p > 0$，有
$R^p\pi_*\mathcal{L}^{\otimes d} = 0$ 且
$R^p(g')_*\mathcal{L}^{\otimes d} = 0$。令
$\mathcal{G} = \pi_*\mathcal{L}^{\otimes d}$。由 Leray 谱序列（《站点上的
上同调》引理 \ref{sites-cohomology-lemma-relative-Leray}），有
$$
E_2^{p, q} = R^pg_* R^q\pi_*\mathcal{L}^{\otimes d}
\Rightarrow
R^{p + q}(g')_*\mathcal{L}^{\otimes d}
$$
按 $d$ 的选取，$E_2^{p, q}$ 中仅 $q = 0$ 的项可能非零，而
$R^{p + q}(g')_*\mathcal{L}^{\otimes d}$ 中仅 $p = q = 0$ 的项可能非零。
这推出当 $p > 0$ 时 $R^pg_*\mathcal{G} = 0$，并且
$g_*\mathcal{G} = (g')_*\mathcal{L}^{\otimes d}$。应用《概形的上同调》
引理 \ref{coherent-lemma-locally-projective-pushforward}，可知
$g_*\mathcal{G} = (g')_*\mathcal{L}^{\otimes d}$ 凝聚。

\medskip\noindent
还必须检验 $\mathcal{G}$ 的支撑是 $Z$。这来自
$\mathcal{L}^{\otimes d}$ 有许多全局截面这一事实。下面详述。由于
$\mathcal{O}_{\mathbf{P}^n}(d)$ 对所有 $d \geq 0$ 都由全局截面生成，
$\mathcal{L}^{\otimes d}$ 也如此。选取 $z \in Z'$，使它映到 $Z$ 的一般点
$\xi$；由于 $\pi$ 为满射，这是可以做到的。（注意，$Z$ 确实有一般点：
$|Z|$ 不可约，且 $Z$ 是 Noether 的，因而拟分离；所以由《空间的性质》
引理 \ref{spaces-properties-lemma-quasi-separated-sober}，$|Z|$ 是 sober
拓扑空间。）选取在 $z$ 处不消失的
$s \in \Gamma(Z', \mathcal{L}^{\otimes d})$。由于
$\Gamma(Z, \mathcal{G}) = \Gamma(Z', \mathcal{L}^{\otimes d})$，可以把
$s$ 看作 $\mathcal{G}$ 的全局截面。选取 $Z'$ 上位于 $z$ 之上的几何点
$\overline{z}$，并把 $Z$ 上相应的几何点记为
$\overline{\xi} = g' \circ \overline{z}$。% P09-ERR-0040
伴随映射
$$
(g')^*\mathcal{G} =
(g')^*g'_*\mathcal{L}^{\otimes d} \longrightarrow \mathcal{L}^{\otimes d}
$$
诱导茎之间的映射
$\mathcal{G}_{\overline{\xi}} \to \mathcal{L}_{\overline{z}}$，见
《空间的性质》引理
\ref{spaces-properties-lemma-stalk-pullback-quasi-coherent}。此外，该伴随
映射把 $s$ 的拉回（把 $s$ 看作 $\mathcal{G}$ 的截面）送到 $s$
（把 $s$ 看作 $\mathcal{L}^{\otimes d}$ 的截面）。因此，$s$ 在下列箭头
源端向量空间中的像
$$
\mathcal{G}_{\overline{\xi}} \otimes \kappa(\overline{\xi})
\longrightarrow
\mathcal{L}^{\otimes d}_{\overline{z}} \otimes \kappa(\overline{z})
$$
非零，因为按 $s$ 的选取，它在箭头目标端的像非零。因此 $\xi$ 位于
$\mathcal{G}$ 的支撑中（《空间的态射》引理
\ref{spaces-morphisms-lemma-support-finite-type}）。由于 $|Z|$ 不可约且
$Z$ 约化，可知 $\mathcal{G}$ 的概形论支撑是整个 $Z$，如所要求。
\end{proof}

\begin{lemma}
\label{spaces-cohomology-lemma-proper-over-affine-cohomology-finite}
设 $A$ 为 Noether 环，$f : X \to \Spec(A)$ 为代数空间的恰当态射，
并设 $\mathcal{F}$ 为凝聚 $\mathcal{O}_X$-模。则对于所有 $i \geq 0$，
$H^i(X, \mathcal{F})$ 是有限 $A$-模。
\end{lemma}

\begin{proof}
这正是引理 \ref{spaces-cohomology-lemma-proper-pushforward-coherent} 的仿射情形。事实上，
由引理 \ref{spaces-cohomology-lemma-higher-direct-image}，$R^if_*\mathcal{F}$ 是拟凝聚层，
因而是与 $A$-模
$\Gamma(\Spec(A), R^if_*\mathcal{F}) = H^i(X, \mathcal{F})$ 相伴的拟凝聚层。
该等式由《站点上的上同调》引理
\ref{sites-cohomology-lemma-apply-Leray} 以及仿射概形上拟凝聚模的高阶
上同调群消失（《概形的上同调》引理
\ref{coherent-lemma-quasi-coherent-affine-cohomology-zero}）推出。由引理
\ref{spaces-cohomology-lemma-coherent-Noetherian}，$R^if_*\mathcal{F}$ 是凝聚层，当且仅当
$H^i(X, \mathcal{F})$ 是有限型 $A$-模。因此引理
\ref{spaces-cohomology-lemma-proper-pushforward-coherent} 给出结论。
\end{proof}

\begin{lemma}
\label{spaces-cohomology-lemma-graded-finiteness}
设 $A$ 为 Noether 环，$B$ 为有限生成的分次 $A$-代数，
$f : X \to \Spec(A)$ 为代数空间的恰当态射。令
$\mathcal{B} = f^*\widetilde B$。设 $\mathcal{F}$ 为有限型拟凝聚分次
$\mathcal{B}$-模。对于每个 $p \geq 0$，分次 $B$-模
$H^p(X, \mathcal{F})$ 都是有限 $B$-模。
\end{lemma}

\begin{proof}
为证明这一点，考虑纤维积图
$$
\xymatrix{
X' = \Spec(B) \times_{\Spec(A)} X
\ar[r]_-\pi \ar[d]_{f'} &
X \ar[d]^f \\
\Spec(B) \ar[r] &
\Spec(A)
}
$$
注意，$f'$ 是恰当态射，见《空间的态射》引理
\ref{spaces-morphisms-lemma-base-change-proper}。此外，$B$ 是有限生成
$A$-代数，因而是 Noether 的（《代数》引理
\ref{algebra-lemma-Noetherian-permanence}）。这推出 $X'$ 是 Noether
代数空间（《空间的态射》引理
\ref{spaces-morphisms-lemma-finite-presentation-noetherian}）。由《空间的态射》
引理 \ref{spaces-morphisms-lemma-affine-equivalence-algebras}，$X'$ 是拟凝聚
$\mathcal{O}_X$-代数 $\mathcal{B}$ 的相对谱。由于 $\mathcal{F}$ 是拟凝聚
$\mathcal{B}$-模，存在唯一的拟凝聚 $\mathcal{O}_{X'}$-模
$\mathcal{F}'$，使得 $\pi_*\mathcal{F}' = \mathcal{F}$，见《空间的态射》
引理 \ref{spaces-morphisms-lemma-affine-equivalence-modules}。由于
$\mathcal{F}$ 作为 $\mathcal{B}$-模是有限型的，$\mathcal{F}'$ 是有限型
$\mathcal{O}_{X'}$-模（细节从略）。换言之，$\mathcal{F}'$ 是凝聚
$\mathcal{O}_{X'}$-模（引理 \ref{spaces-cohomology-lemma-coherent-Noetherian}）。由于态射
$\pi : X' \to X$ 仿射，有
$$
H^p(X, \mathcal{F}) = H^p(X', \mathcal{F}')
$$
这由引理 \ref{spaces-cohomology-lemma-affine-vanishing-higher-direct-images} 和《站点上的上同调》
引理 \ref{sites-cohomology-lemma-apply-Leray} 推出。因此本引理由引理
\ref{spaces-cohomology-lemma-proper-over-affine-cohomology-finite} 推出。
\end{proof}









\section{丰沛可逆层与上同调}
\label{spaces-cohomology-section-ample-cohomology}

\noindent
下面用扭转后的上同调消失性，给出仿射基上恰当代数空间的丰沛性判别法。

\begin{lemma}
\label{spaces-cohomology-lemma-vanshing-gives-ample}
设 $R$ 为 Noether 环，$X$ 为 $R$ 上的恰当代数空间，并设 $\mathcal{L}$
为可逆 $\mathcal{O}_X$-模。下列条件等价：
\begin{enumerate}
\item $X$ 是概形，且 $\mathcal{L}$ 在 $X$ 上丰沛；
\item 对于每个凝聚 $\mathcal{O}_X$-模 $\mathcal{F}$，存在
$n_0 \geq 0$，使得对于所有 $n \geq n_0$ 和 $p > 0$，都有
$H^p(X, \mathcal{F} \otimes \mathcal{L}^{\otimes n}) = 0$；
\item 对于每个凝聚 $\mathcal{O}_X$-模 $\mathcal{F}$，存在
$n \geq 1$，使得
$H^1(X, \mathcal{F} \otimes \mathcal{L}^{\otimes n}) = 0$。
\end{enumerate}
\end{lemma}

\begin{proof}
(1) $\Rightarrow$ (2) 由《概形的上同调》引理
\ref{coherent-lemma-vanshing-gives-ample} 推出。(2) $\Rightarrow$ (3) 显然。
(3) $\Rightarrow$ (1) 即引理 \ref{spaces-cohomology-lemma-Noetherian-h1-zero-invertible}。
\end{proof}

\begin{lemma}
\label{spaces-cohomology-lemma-surjective-finite-morphism-ample}
设 $R$ 为 Noether 环，$f : Y \to X$ 为 $R$ 上恰当代数空间之间的态射。
设 $\mathcal{L}$ 为可逆 $\mathcal{O}_X$-模。假设 $f$ 有限且为满射。
下列条件等价：
\begin{enumerate}
\item $X$ 是概形且 $\mathcal{L}$ 丰沛；
\item $Y$ 是概形且 $f^*\mathcal{L}$ 丰沛。
\end{enumerate}
\end{lemma}

\begin{proof}
假设 (1) 成立。有限态射（由概形）可表，所以 $Y$ 是概形，见《空间的态射》
引理 \ref{spaces-morphisms-lemma-integral-local}。于是 (2) 由
《概形的上同调》引理
\ref{coherent-lemma-surjective-finite-morphism-ample} 推出。

\medskip\noindent
假设 (2) 成立。对凝聚 $\mathcal{O}_X$-模 $\mathcal{F}$，定义下述性质
$P$：存在 $n_0$，使得对于所有 $n \geq n_0$ 和 $p > 0$，都有
$H^p(X, \mathcal{F} \otimes \mathcal{L}^{\otimes n}) = 0$。我们将证明
任意凝聚 $\mathcal{O}_X$-模 $\mathcal{F}$ 都具有 $P$；由引理
\ref{spaces-cohomology-lemma-vanshing-gives-ample}，这蕴含 $\mathcal{L}$ 丰沛。我们将应用
引理 \ref{spaces-cohomology-lemma-property-higher-rank-cohomological}，因此必须为 $P$ 验证
该引理的 (1)、(2)、(3)。性质 (1) 来自与层的短正合列相伴的上同调
长正合列，以及与可逆层作张量积是正合函子这一事实。性质 (2) 来自
$H^p(X, -)$ 是可加函子。

\medskip\noindent
为证明 (3)，设 $i : Z \to X$ 为满足 $|Z|$ 不可约的约化闭子空间。
设 $i' : Z' \to Y$ 和 $f' : Z' \to Z$ 如引理
\ref{spaces-cohomology-lemma-finite-morphism-Noetherian} 中所述，并令
$\mathcal{G} = f'_*\mathcal{O}_{Z'}$。我们断言 $\mathcal{G}$ 满足引理
\ref{spaces-cohomology-lemma-property-higher-rank-cohomological} 的性质 (3)(a) 和 (3)(b)，
这将完成证明。性质 (3)(a) 已在引理
\ref{spaces-cohomology-lemma-finite-morphism-Noetherian} 中证明。为证明 (3)(b)，设
$\mathcal{I}$ 为 $Z$ 上的非零拟凝聚理想层。记
$\mathcal{I}' \subset \mathcal{O}_{Z'}$ 为拟凝聚理想
$(f')^{-1}\mathcal{I} \mathcal{O}_{Z'}$，也就是
$(f')^*\mathcal{I} \to \mathcal{O}_{Z'}$ 的像。由引理
\ref{spaces-cohomology-lemma-affine-morphism-projection-ideal}，有
$f_*\mathcal{I}' = \mathcal{I} \mathcal{G}$。% P09-ERR-0041
我们断言其共同值
$\mathcal{G}' = \mathcal{I} \mathcal{G} = f'_*\mathcal{I}'$ 满足 (3)(b)
中的条件。首先，显然 $\mathcal{G}/\mathcal{G}'$ 的支撑包含在
$\mathcal{O}_Z/\mathcal{I}$ 的支撑中；后者是 $|Z|$ 的真子空间，因为
$\mathcal{I}$ 是约化不可约代数空间 $Z$ 上的非零理想层。回忆 $f'_*$、
$i_*$、$i'_*$ 把凝聚模变为凝聚模，见引理
\ref{spaces-cohomology-lemma-finite-pushforward-coherent} 和 \ref{spaces-cohomology-lemma-i-star-equivalence}。
由于 $Y$ 是概形且 $\mathcal{L}$ 丰沛，由引理
\ref{spaces-cohomology-lemma-vanshing-gives-ample}，存在 $n_0$ 使得% P09-ERR-0042
$$
H^p(Y, i'_*\mathcal{I}' \otimes_{\mathcal{O}_Y} f^*\mathcal{L}^{\otimes n}) = 0
$$
对于 $n \geq n_0$ 和 $p > 0$ 成立。于是得到
\begin{align*}
H^p(X, i_*\mathcal{G}' \otimes_{\mathcal{O}_X} \mathcal{L}^{\otimes n})
& =
H^p(Z, \mathcal{G'} \otimes_{\mathcal{O}_Z} i^*\mathcal{L}^{\otimes n}) \\
& =
H^p(Z,
f'_*\mathcal{I}' \otimes_{\mathcal{O}_Z} i^*\mathcal{L}^{\otimes n})) \\
& =
H^p(Z, f'_*(\mathcal{I}' \otimes_{\mathcal{O}_{Z'}}
(f')^*i^*\mathcal{L}^{\otimes n})) \\
& =
H^p(Z, f'_*(\mathcal{I}' \otimes_{\mathcal{O}_{Z'}}
(i')^*f^*\mathcal{L}^{\otimes n})) \\
& =
H^p(Z',
\mathcal{I}' \otimes_{\mathcal{O}_{Z'}} (i')^*f^*\mathcal{L}^{\otimes n})) \\
& =
H^p(Y, i'_*\mathcal{I}' \otimes_{\mathcal{O}_Y} f^*\mathcal{L}^{\otimes n}) = 0
\end{align*}
% P09-ERR-0043
这里使用了投影公式和 Leray 谱序列（见《站点上的上同调》
\ref{sites-cohomology-section-projection-formula} 节和
\ref{sites-cohomology-section-leray} 节），以及引理
\ref{spaces-cohomology-lemma-finite-higher-direct-image-zero}。这验证了引理
\ref{spaces-cohomology-lemma-property-higher-rank-cohomological} 的性质 (3)(b)，如所要求。
\end{proof}










\section{形式函数定理}
\label{spaces-cohomology-section-theorem-formal-functions}

\noindent
本节对应于《概形的上同调》
\ref{coherent-section-theorem-formal-functions} 节。建议读者先阅读该节。

\begin{situation}
\label{spaces-cohomology-situation-formal-functions}
这里 $A$ 是 Noether 环，$I \subset A$ 是理想。此外，
$f : X \to \Spec(A)$ 是代数空间的恰当态射，$\mathcal{F}$ 是 $X$ 上的
凝聚层。
\end{situation}

\noindent
在这一情形下，用 $I^n\mathcal{F}$ 表示 $\mathcal{F}$ 的拟凝聚子模：
它作为 $\mathcal{O}_X$-模由 $\mathcal{F}$ 的局部截面与 $I^n$ 中元素的
乘积生成。换言之，它是映射
$f^*\widetilde{I} \otimes_{\mathcal{O}_X} \mathcal{F} \to \mathcal{F}$ 的像。
% P09-ERR-0044

\begin{lemma}
\label{spaces-cohomology-lemma-cohomology-powers-ideal-times-F}
在情形 \ref{spaces-cohomology-situation-formal-functions} 中。令
$B = \bigoplus_{n \geq 0} I^n$。则对于每个 $p \geq 0$，分次 $B$-模
$\bigoplus_{n \geq 0} H^p(X, I^n\mathcal{F})$ 是有限 $B$-模。
\end{lemma}

\begin{proof}
令 $\mathcal{B} = \bigoplus I^n\mathcal{O}_X = f^*\widetilde{B}$。
则 $\bigoplus I^n\mathcal{F}$ 是有限型分次 $\mathcal{B}$-模。因此结论由
引理 \ref{spaces-cohomology-lemma-graded-finiteness} 推出。
\end{proof}

\begin{lemma}
\label{spaces-cohomology-lemma-cohomology-powers-ideal-application}
在情形 \ref{spaces-cohomology-situation-formal-functions} 中。对于每个 $p \geq 0$，
存在整数 $c \geq 0$，使得：
\begin{enumerate}
\item 对于所有 $n \geq c$，乘法映射
$I^{n - c} \otimes H^p(X, I^c\mathcal{F}) \to H^p(X, I^n\mathcal{F})$
是满射；
\item 对于所有 $n \geq 0$、$m \geq c$，映射
$H^p(X, I^{n + m}\mathcal{F}) \to H^p(X, I^n\mathcal{F})$ 的像包含在
子模 $I^{m - c} H^p(X, I^n\mathcal{F})$ 中。
\end{enumerate}
\end{lemma}

\begin{proof}
由引理 \ref{spaces-cohomology-lemma-cohomology-powers-ideal-times-F}，可以找到
$d_1, \ldots, d_t \geq 0$ 以及
$x_i \in H^p(X, I^{d_i}\mathcal{F})$，使得
$\bigoplus_{n \geq 0} H^p(X, I^n\mathcal{F})$ 作为
$B = \bigoplus_{n \geq 0} I^n$-模由 $x_1, \ldots, x_t$ 生成。取
$c = \max\{d_i\}$。(1) 显然成立。对于 (2)，令
$b = \max(0, n - c)$。考虑 $A$-模的交换图
$$
\xymatrix{
I^{n + m - c - b} \otimes I^b \otimes
H^p(X, I^c\mathcal{F}) \ar[r] \ar[d] &
I^{n + m - c} \otimes H^p(X, I^c\mathcal{F}) \ar[r] &
H^p(X, I^{n + m}\mathcal{F}) \ar[d] \\
I^{n + m - c - b} \otimes H^p(X, I^n\mathcal{F}) \ar[rr] & &
H^p(X, I^n\mathcal{F})
}
$$
由本引理的 (1)，当 $n + m \geq c$ 时，水平箭头的复合是满射。
另一方面，显然 $n + m - c - b \geq m - c$。所以 (2) 成立。
\end{proof}

\begin{lemma}
\label{spaces-cohomology-lemma-ML-cohomology-powers-ideal}
在情形 \ref{spaces-cohomology-situation-formal-functions} 中。固定 $p \geq 0$。
\begin{enumerate}
\item 存在 $c_1 \geq 0$，使得对于所有 $n \geq c_1$，有
$$
\Ker(
H^p(X, \mathcal{F}) \to H^p(X, \mathcal{F}/I^n\mathcal{F})
)
\subset
I^{n - c_1}H^p(X, \mathcal{F}).
$$
\item 逆系统
$$
\left(H^p(X, \mathcal{F}/I^n\mathcal{F})\right)_{n \in \mathbf{N}}
$$
满足 Mittag-Leffler 条件（见《同调》定义
\ref{homology-definition-Mittag-Leffler}）。
\item 事实上，对于任意 $p$ 和 $n$，存在 $c_2(n) \geq n$，使得
$$
\Im(H^p(X, \mathcal{F}/I^k\mathcal{F})
\to H^p(X, \mathcal{F}/I^n\mathcal{F}))
=
\Im(H^p(X, \mathcal{F})
\to H^p(X, \mathcal{F}/I^n\mathcal{F}))
$$
对于所有 $k \geq c_2(n)$ 成立。
\end{enumerate}
\end{lemma}

\begin{proof}
令 $c_1 = \max\{c_p, c_{p + 1}\}$，其中 $c_p,c_{p +1}$ 是引理
\ref{spaces-cohomology-lemma-cohomology-powers-ideal-application} 分别对 $H^p$ 和
$H^{p + 1}$ 给出的整数。我们将在 (1)、(2)、(3) 的证明中使用这个常数。

\medskip\noindent
先证明 (1)。考虑短正合列
$$
0 \to I^n\mathcal{F} \to \mathcal{F} \to \mathcal{F}/I^n\mathcal{F} \to 0
$$
由上同调长正合列可见
$$
\Ker(
H^p(X, \mathcal{F}) \to H^p(X, \mathcal{F}/I^n\mathcal{F})
)
=
\Im(
H^p(X, I^n\mathcal{F}) \to H^p(X, \mathcal{F})
)
$$
因此按 $c_1$ 的选取，当 $n \geq c_1$ 时，它包含在
$I^{n - c_1}H^p(X, \mathcal{F})$ 中。

\medskip\noindent
注意，按 Mittag-Leffler 条件的定义，(3) 蕴含 (2)。

\medskip\noindent
下面证明 (3)。在余下的证明中固定 $n$。考虑交换图
$$
\xymatrix{
0 \ar[r] &
I^n\mathcal{F} \ar[r] &
\mathcal{F} \ar[r] &
\mathcal{F}/I^n\mathcal{F} \ar[r] &
0 \\
0 \ar[r] &
I^{n + m}\mathcal{F} \ar[r] \ar[u] &
\mathcal{F} \ar[r] \ar[u] &
\mathcal{F}/I^{n + m}\mathcal{F} \ar[r] \ar[u] &
0
}
$$
它诱导出下面的交换图
$$
\xymatrix{
H^p(X, I^n\mathcal{F}) \ar[r] &
H^p(X, \mathcal{F}) \ar[r] &
H^p(X, \mathcal{F}/I^n\mathcal{F}) \ar[r]_\delta &
H^{p + 1}(X, I^n\mathcal{F}) \\
H^p(X, I^{n + m}\mathcal{F}) \ar[r] \ar[u] &
H^p(X, \mathcal{F}) \ar[r] \ar[u]^1 &
H^p(X, \mathcal{F}/I^{n + m}\mathcal{F}) \ar[r] \ar[u] &
H^{p + 1}(X, I^{n + m}\mathcal{F}) \ar[u]^a
}
$$
如果 $m \geq c_1$，则 $a$ 的像包含在
$I^{m - c_1} H^{p + 1}(X, I^n\mathcal{F})$ 中。由 Artin--Rees 引理
（见《代数》引理 \ref{algebra-lemma-map-AR}），存在整数 $c_3(n)$，使得
$$
I^N H^{p + 1}(X, I^n\mathcal{F}) \cap \Im(\delta)
\subset
\delta\left(I^{N - c_3(n)}H^p(X, \mathcal{F}/I^n\mathcal{F})\right)
$$
对于所有 $N \geq c_3(n)$ 成立。由于
$H^p(X, \mathcal{F}/I^n\mathcal{F})$ 被 $I^n$ 零化，可见如果
$m \geq c_3(n) + c_1 + n$，则
$$
\Im(H^p(X, \mathcal{F}/I^{n + m}\mathcal{F})
\to H^p(X, \mathcal{F}/I^n\mathcal{F}))
=
\Im(H^p(X, \mathcal{F})
\to H^p(X, \mathcal{F}/I^n\mathcal{F}))
$$
换言之，取 $c_2(n) = c_3(n) + c_1 + n$，(3) 成立。
\end{proof}

\begin{theorem}[形式函数定理]
\label{spaces-cohomology-theorem-formal-functions}
在情形 \ref{spaces-cohomology-situation-formal-functions} 中，固定 $p \geq 0$。映射系统
$$
H^p(X, \mathcal{F})/I^nH^p(X, \mathcal{F})
\longrightarrow
H^p(X, \mathcal{F}/I^n\mathcal{F})
$$
诱导极限之间的同构
$$
H^p(X, \mathcal{F})^\wedge
\longrightarrow
\lim_n H^p(X, \mathcal{F}/I^n\mathcal{F})
$$
其中左端是 $A$-模 $H^p(X, \mathcal{F})$ 关于理想 $I$ 的完备化，见
《代数》\ref{algebra-section-completion} 节。此外，对于极限拓扑，
这实际上是同胚。
\end{theorem}

\begin{proof}
事实上，这立即由引理 \ref{spaces-cohomology-lemma-ML-cohomology-powers-ideal} 推出。
下面写出细节。令 $M = H^p(X, \mathcal{F})$ 且
$M_n = H^p(X, \mathcal{F}/I^n\mathcal{F})$，并记
$N_n = \Im(M \to M_n)$。由《同调》
\ref{homology-section-inverse-systems} 节对极限的描述，有
$$
\lim_n M_n
=
\{(x_n) \in \prod M_n \mid \varphi_i(x_n) = x_{n - 1}, \ n = 2, 3, \ldots\}
$$
选取元素 $x = (x_n) \in \lim_n M_n$。% P09-ERR-0045
由引理 \ref{spaces-cohomology-lemma-ML-cohomology-powers-ideal} (3)，对于所有 $n$ 都有
$x_n \in N_n$，因为按定义，对于所有 $m$，$x_n$ 是某个
$x_{n + m} \in M_{n + m}$ 的像。由引理
\ref{spaces-cohomology-lemma-ML-cohomology-powers-ideal} (1)，约化映射存在分解
$$
M \to N_n \to M/I^{n - c_1}M
$$
对于 $n \geq c_1$，记 $y_n \in M/I^{n - c_1}M$ 为 $x_n$ 的像。
当 $n' \geq n$ 时，复合映射 $M \to M_{n'} \to M_n$ 是给定的映射
$M \to M_n$，所以在典范映射
$M/I^{n' - c_1}M \to M/I^{n - c_1}M$ 下，$y_{n'}$ 映到 $y_n$。
因此 $y = (y_{n + c_1})$ 定义了 $\lim_n M/I^nM$ 的一个元素。
我们略去验证 $y$ 在本引理的映射
$$
M^\wedge = \lim_n M/I^nM \longrightarrow \lim_n M_n
$$
下映到 $x$，也略去关于拓扑的验证。
\end{proof}

\begin{lemma}
\label{spaces-cohomology-lemma-spell-out-theorem-formal-functions}
设 $A$ 为一个环，$I \subset A$ 为理想。假设 $A$ 是 Noether 的并关于
$I$ 完备。设 $f : X \to \Spec(A)$ 为代数空间的恰当态射，
$\mathcal{F}$ 为 $X$ 上的凝聚层。则
$$
H^p(X, \mathcal{F}) = \lim_n H^p(X, \mathcal{F}/I^n\mathcal{F})
$$
对于所有 $p \geq 0$ 成立。
\end{lemma}

\begin{proof}
这是形式函数定理（定理 \ref{spaces-cohomology-theorem-formal-functions}）在完备 Noether
基环情形下的改述。事实上，在这种情形下，$A$-模
$H^p(X, \mathcal{F})$ 是有限的（引理
\ref{spaces-cohomology-lemma-proper-over-affine-cohomology-finite}），因而是 $I$-进完备的
（《代数》引理 \ref{algebra-lemma-completion-tensor}）。所以左端不必再
作完备化。
\end{proof}

\begin{lemma}
\label{spaces-cohomology-lemma-formal-functions-stalk}
设 $S$ 为一个概形，$f : X \to Y$ 为 $S$ 上代数空间之间的态射，
$\mathcal{F}$ 为 $X$ 上的拟凝聚层。假设：
\begin{enumerate}
\item $Y$ 局部 Noether；
\item $f$ 恰当；
\item $\mathcal{F}$ 凝聚。
\end{enumerate}
设 $\overline{y}$ 为 $Y$ 的几何点。考虑纤维
$X_1 = X_{\overline{y}}$ 的“无穷小邻域”
$$
\xymatrix{
X_n =
\Spec(\mathcal{O}_{Y, \overline{y}}/\mathfrak m_{\overline{y}}^n) \times_Y X
\ar[r]_-{i_n} \ar[d]_{f_n} &
X \ar[d]^f \\
\Spec(\mathcal{O}_{Y, \overline{y}}/\mathfrak m_{\overline{y}}^n)
\ar[r]^-{c_n} & Y
}
$$
并令 $\mathcal{F}_n = i_n^*\mathcal{F}$。则有
$$
\left(R^pf_*\mathcal{F}\right)_{\overline{y}}^\wedge
\cong
\lim_n H^p(X_n, \mathcal{F}_n)
$$
作为 $\mathcal{O}_{Y, \overline{y}}^\wedge$-模的同构。
\end{lemma}

\begin{proof}
这只是形式函数定理（定理 \ref{spaces-cohomology-theorem-formal-functions}）一个特殊情形的
改述。下面详述。注意，$\mathcal{O}_{Y, \overline{y}}$ 是 Noether 局部环，
见《空间的性质》引理
\ref{spaces-properties-lemma-Noetherian-local-ring-Noetherian}。考虑典范态射
$c : \Spec(\mathcal{O}_{Y, \overline{y}}) \to Y$。它等同局部环，因而是
平坦态射。记 $f' : X' \to \Spec(\mathcal{O}_{Y, \overline{y}})$ 为 $f$
到这个局部环上的基变换。由引理 \ref{spaces-cohomology-lemma-flat-base-change-cohomology}，
有 $c^*R^pf_*\mathcal{F} = R^pf'_*\mathcal{F}'$。此外，对于所有
$n \geq 1$，有典范等同 $X_n = X'_n$。

\medskip\noindent
因此可以假设 $Y = \Spec(A)$，其中 $A$ 是极大理想为 $\mathfrak m$ 的
严格 Hensel Noether 局部环，并且 $\overline{y} \to Y$ 等于
$\Spec(A/\mathfrak m) \to Y$。于是
$$
\left(R^pf_*\mathcal{F}\right)_{\overline{y}} =
\Gamma(Y, R^pf_*\mathcal{F}) =
H^p(X, \mathcal{F})
$$
因为 $(Y, \overline{y})$ 是 $\overline{y}$ 的 \'etale 邻域范畴中的始对象。
各态射 $c_n$ 都是闭浸入，所以其基变换 $i_n$ 也都是闭浸入。注意
$i_{n, *}\mathcal{F}_n = i_{n, *}i_n^*\mathcal{F}
= \mathcal{F}/\mathfrak m^n\mathcal{F}$。由 $i_n$ 的 Leray 谱序列以及
引理 \ref{spaces-cohomology-lemma-finite-pushforward-coherent}，可见
$$
H^p(X_n, \mathcal{F}_n) =
H^p(X, i_{n, *}\mathcal{F}) =
H^p(X, \mathcal{F}/\mathfrak m^n\mathcal{F})
$$
% P09-ERR-0046
因此确实可以应用形式函数定理来计算引理陈述中的极限，结论得证。
\end{proof}

\noindent
下面的引理将在后一个引理中推广到维数 $>0$ 的纤维。

\begin{lemma}
\label{spaces-cohomology-lemma-higher-direct-images-zero-finite-fibre}
设 $S$ 为一个概形，$f : X \to Y$ 为 $S$ 上代数空间之间的态射，
$\overline{y}$ 为 $Y$ 的几何点。假设：
\begin{enumerate}
\item $Y$ 局部 Noether；
\item $f$ 恰当；
\item $X_{\overline{y}}$ 的底拓扑空间是离散的。
\end{enumerate}
则对于 $X$ 上的任意凝聚层 $\mathcal{F}$ 和所有 $p > 0$，都有
$(R^pf_*\mathcal{F})_{\overline{y}} = 0$。
\end{lemma}

\begin{proof}
设 $\kappa(\overline{y})$ 为局部环 $\mathcal{O}_{Y, \overline{y}}$ 的剩余域。
如引理 \ref{spaces-cohomology-lemma-formal-functions-stalk} 中那样，令
$X_{\overline{y}} = X_1 = \Spec(\kappa(\overline{y})) \times_Y X$。
由《空间的态射》引理 \ref{spaces-morphisms-lemma-quasi-finite-at-point}，
态射 $f : X \to Y$ 在 $X \to Y$ 位于 $\overline{y}$ 上的纤维的每个点处
拟有限。因此 $X_{\overline{y}} \to \overline{y}$ 分离且拟有限。于是由
《空间的态射》命题
\ref{spaces-morphisms-proposition-locally-quasi-finite-separated-over-scheme}，
$X_{\overline{y}}$ 是概形。由于它拟紧，其底拓扑空间是有限离散空间。
再由《概形》引理 \ref{schemes-lemma-scheme-finite-discrete-affine}，它是
仿射概形。由引理
\ref{spaces-cohomology-lemma-image-affine-finite-morphism-affine-Noetherian}，各代数空间
$X_n$ 也都是仿射概形。此外，每个 $X_n$ 的底拓扑与 $X_1$ 的相同。
% P09-ERR-0047
因此对于所有 $p > 0$，有 $H^p(X_n, \mathcal{F}_n) = 0$。由引理
\ref{spaces-cohomology-lemma-formal-functions-stalk}，可见
$(R^pf_*\mathcal{F})_{\overline{y}}^\wedge = 0$。由引理
\ref{spaces-cohomology-lemma-proper-pushforward-coherent}，$R^pf_*\mathcal{F}$ 凝聚，因而
$R^pf_*\mathcal{F}_{\overline{y}}$ 是有限
$\mathcal{O}_{Y, \overline{y}}$-模。% P09-ERR-0048
由《代数》引理 \ref{algebra-lemma-completion-tensor}，这推出
$(R^pf_*\mathcal{F})_{\overline{y}} = 0$。
\end{proof}

\begin{lemma}
\label{spaces-cohomology-lemma-higher-direct-images-zero-above-dimension-fibre}
\begin{slogan}
对于恰当映射，次数大于纤维维数的高阶正像之茎为零。
\end{slogan}
设 $S$ 为一个概形，$f : X \to Y$ 为 $S$ 上代数空间之间的态射，
$\overline{y}$ 为 $Y$ 的几何点。假设：
\begin{enumerate}
\item $Y$ 局部 Noether；
\item $f$ 恰当；
\item $\dim(X_{\overline{y}}) = d$.
\end{enumerate}
则对于 $X$ 上的任意凝聚层 $\mathcal{F}$ 和所有 $p > d$，都有
$(R^pf_*\mathcal{F})_{\overline{y}} = 0$。
\end{lemma}

\begin{proof}
设 $\kappa(\overline{y})$ 为局部环 $\mathcal{O}_{Y, \overline{y}}$ 的剩余域。
如引理 \ref{spaces-cohomology-lemma-formal-functions-stalk} 中那样，令
$X_{\overline{y}} = X_1 = \Spec(\kappa(\overline{y})) \times_Y X$。
此外，每个无穷小邻域 $X_n$ 的底拓扑空间都与 $X_{\overline{y}}$ 的相同。
因此由引理 \ref{spaces-cohomology-lemma-vanishing-above-dimension}，对于所有 $p > d$，有
$H^p(X_n, \mathcal{F}_n) = 0$。进而由引理
\ref{spaces-cohomology-lemma-formal-functions-stalk}，当 $p > d$ 时
$(R^pf_*\mathcal{F})_{\overline{y}}^\wedge = 0$。由引理
\ref{spaces-cohomology-lemma-proper-pushforward-coherent}，$R^pf_*\mathcal{F}$ 凝聚，因而
$R^pf_*\mathcal{F}_{\overline{y}}$ 是有限
$\mathcal{O}_{Y, \overline{y}}$-模。% P09-ERR-0048
由《代数》引理 \ref{algebra-lemma-completion-tensor}，这推出
$(R^pf_*\mathcal{F})_{\overline{y}} = 0$。
\end{proof}






\section{形式函数定理的应用}
\label{spaces-cohomology-section-applications-formal-functions}

\noindent
我们将在需要时在这里加入更多内容。

\begin{lemma}
\label{spaces-cohomology-lemma-characterize-finite}
（更一般的版本见《空间的态射进阶》引理
\ref{spaces-more-morphisms-lemma-characterize-finite}。）
设 $S$ 为一个概形，$f : X \to Y$ 为 $S$ 上代数空间之间的态射。
假设 $Y$ 局部 Noether。下列条件等价：
\begin{enumerate}
\item $f$ 有限；
\item $f$ 恰当，并且对于每个态射 $\Spec(k) \to Y$，其中 $k$ 是域，
$|X_k|$ 都是离散空间。
\end{enumerate}
\end{lemma}

\begin{proof}
由《空间的态射》引理 \ref{spaces-morphisms-lemma-finite-proper}，有限态射
恰当。由《空间的态射》引理
\ref{spaces-morphisms-lemma-finite-quasi-finite}，有限态射拟有限。
拟有限态射的纤维 $X_k$ 离散，见《空间的态射》引理
\ref{spaces-morphisms-lemma-locally-quasi-finite}。因此有限态射恰当且具有
离散纤维 $X_k$。

\medskip\noindent
假设 $f$ 恰当且具有离散纤维 $X_k$。我们要证明 $f$ 有限。事实上，只需
证明 $f$ 仿射。如果 $f$ 仿射，则由《空间的态射》引理
\ref{spaces-morphisms-lemma-integral-universally-closed}，$f$ 是整态射；
随后由《空间的态射》引理 \ref{spaces-morphisms-lemma-finite-integral}，
$f$ 有限。

\medskip\noindent
为证明 $f$ 仿射，可以假设 $Y$ 仿射；目标是证明 $X$ 也仿射。由于 $f$
恰当，$X$ 分离且拟紧。我们将证明，对于任意凝聚
$\mathcal{O}_X$-模 $\mathcal{F}$，都有 $H^1(X, \mathcal{F}) = 0$。
由引理 \ref{spaces-cohomology-lemma-directed-colimit-coherent} 和 \ref{spaces-cohomology-lemma-colimits}，
这推出对于每个拟凝聚 $\mathcal{O}_X$-模 $\mathcal{F}$，都有
$H^1(X, \mathcal{F}) = 0$。进而由命题
\ref{spaces-cohomology-proposition-vanishing-affine}，$X$ 仿射。由引理
\ref{spaces-cohomology-lemma-higher-direct-images-zero-finite-fibre}，对于 $Y$ 的所有几何点，
$R^1f_*\mathcal{F}$ 的茎都为零。换言之，$R^1f_*\mathcal{F} = 0$。
因此由 $f$ 的 Leray 谱序列，
$H^1(X , \mathcal{F}) = H^1(Y, f_*\mathcal{F})$。由于 $Y$ 仿射且
$f_*\mathcal{F}$ 拟凝聚（《空间的态射》引理
\ref{spaces-morphisms-lemma-pushforward}），《概形的上同调》引理
\ref{coherent-lemma-quasi-coherent-affine-cohomology-zero} 给出
$H^1(Y, f_*\mathcal{F}) = 0$。所以 $H^1(X, \mathcal{F}) = 0$，如所要求。
\end{proof}

\noindent
作为推论，得到下面的有用结果。

\begin{lemma}
\label{spaces-cohomology-lemma-proper-finite-fibre-finite-in-neighbourhood}
（更一般的版本见《空间的态射进阶》引理
\ref{spaces-more-morphisms-lemma-proper-finite-fibre-finite-in-neighbourhood}。）
设 $S$ 为一个概形，$f : X \to Y$ 为 $S$ 上代数空间之间的态射，
$\overline{y}$ 为 $Y$ 的几何点。假设：
\begin{enumerate}
\item $Y$ 局部 Noether；
\item $f$ 恰当；
\item $|X_{\overline{y}}|$ 有限。
\end{enumerate}
则存在 $\overline{y}$ 的开邻域 $V \subset Y$，使得
$f|_{f^{-1}(V)} : f^{-1}(V) \to V$ 有限。
\end{lemma}

\begin{proof}
由《空间的态射》引理 \ref{spaces-morphisms-lemma-quasi-finite-at-point}，
态射 $f$ 在 $X$ 中位于 $\overline{y}$ 上的所有几何点处拟有限。由
《空间的态射》引理
\ref{spaces-morphisms-lemma-locally-finite-type-quasi-finite-part}，$f$
拟有限的点所成集合是开子空间 $U \subset X$。令 $Z = X \setminus U$。
则 $\overline{y} \not \in f(Z)$。由于 $f$ 恰当，集合 $f(Z) \subset Y$ 闭。
选取 $\overline{y}$ 的任意开邻域 $V \subset Y$，使得
$Z \cap V = \emptyset$。% P09-ERR-0049
则 $f^{-1}(V) \to V$ 局部拟有限且恰当。因此
$f^{-1}(V) \to V$ 具有离散纤维 $X_k$（《空间的态射》引理
\ref{spaces-morphisms-lemma-locally-quasi-finite}），这些纤维拟紧，因而有限。
所以由引理 \ref{spaces-cohomology-lemma-characterize-finite}，$f^{-1}(V) \to V$ 有限。
\end{proof}
% Generated mechanically from the frozen spaces-limits.tex target.
% Do not edit this generated file; edit the bound locale source instead.

% OK, start here.
%

\setcounter{chapter}{69}
\chapter{代数空间的极限}



\phantomsection
\label{spaces-limits-section-phantom}


\section{引言}
\label{spaces-limits-section-introduction}

\noindent
本章汇集与代数空间的极限有关的材料。第一个主题是刻画基底 $S$ 上
局部有限表示的代数空间 $F$：它们正是保极限函子。随后我们研究以
有向集为指标（《范畴》定义 \ref{categories-definition-directed-set}）、
转移态射为仿射态射的逆系统的极限。依照 \cite{CLO}，我们讨论拟紧
拟分离代数空间的绝对 Noether 逼近。另一种方法由 David Rydh 给出
（见 \cite{rydh_approx}）；其结果还涵盖某些代数叠的绝对 Noether 逼近。


\section{约定}
\label{spaces-limits-section-conventions}

\noindent
我们始终假设所有概形都包含在一个 fppf 大位点 $\Sch_{fppf}$ 中，
并且所考虑的每个环 $A$ 都具有如下性质：$\Spec(A)$（同构于）
这个大位点中的一个对象。

\medskip\noindent
设 $S$ 为概形，$X$ 为 $S$ 上的代数空间。在本章及以下各章中，
我们用 $X \times_S X$ 表示 $X$ 与自身的积（在 $S$ 上代数空间的
范畴中），而不用 $X \times X$。











\section{有限表示态射}
\label{spaces-limits-section-finite-presentation}

\noindent
本节把《极限》命题
\ref{limits-proposition-characterize-locally-finite-presentation}
推广到代数空间的态射。下面定义的动机来自刚才引用的命题。

\begin{definition}
\label{spaces-limits-definition-locally-finite-presentation}
设 $S$ 为概形。
\begin{enumerate}
\item 称函子 $F : (\Sch/S)_{fppf}^{opp} \to \textit{Sets}$ {\it 保极限}或
{\it 局部有限表示}，如果对于 $S$ 上仿射概形的任意有向逆系统
$T_i$，只要其极限 $T = \lim T_i$ 是 $S$ 上的仿射概形，就有
$$
F(T) = \colim F(T_i).
$$
我们有时也说 $F$ {\it 在 $S$ 上局部有限表示}。
\item 设 $F, G : (\Sch/S)_{fppf}^{opp} \to \textit{Sets}$。称函子的自然变换
$a : F \to G$ {\it 保极限}或{\it 局部有限表示}，如果对于 $S$ 上
任意概形 $T$ 及任意 $y \in G(T)$，函子
$$
F_y : (\Sch/T)_{fppf}^{opp} \longrightarrow \textit{Sets}, \quad
T'/T \longmapsto \{x \in F(T') \mid a(x) = y|_{T'}\}
$$
在 $T$ 上局部有限表示\footnote{引理
\ref{spaces-limits-lemma-characterize-relative-limit-preserving} 中的刻画 (2)
或许更容易理解。}。我们有时也说 $F$ {\it 相对于 $G$ 保极限}。
\end{enumerate}
\end{definition}

\noindent
在某种意义下，函子 $F_y$ 是 $a : F \to G$ 在 $y$ 上的纤维，
但它是 $T$ 的 fppf 大位点上的预层。这个函子可写成
\begin{equation}
\label{spaces-limits-equation-fibre-map-functors}
F_y =
F|_{(\Sch/T)_{fppf}}
{\times}_{G|_{(\Sch/T)_{fppf}}}
*
\end{equation}
其中 $*$ 是 $(\Sch/T)_{fppf}$ 上（预）层范畴的终对象（见
《位点》例 \ref{sites-example-singleton-sheaf}），而映射
$* \to G|_{(\Sch/T)_{fppf}}$ 由 $y$ 给出。注意，如果
$j : (\Sch/T)_{fppf} \to (\Sch/S)_{fppf}$ 是局部化函子，
那么上述公式变为 $F_y = j^{-1}F \times_{j^{-1}G} *$，且
$j_!F_y$ 正是纤维积 $F \times_{G, y} T$。（关于局部化，见
《位点》第 \ref{sites-section-localize} 节；特别地，关于预层的
$j_!$，见《位点》注 \ref{sites-remark-localize-presheaves}。）

\medskip\noindent
此时，对于 $S$ 上代数空间的态射 $X \to Y$ 何谓局部有限表示，
我们暂时有两个定义：一个来自《空间的态射》定义
\ref{spaces-morphisms-definition-locally-finite-presentation}；另一个把
$X \to Y$ 看成函子的自然变换，从而应用定义
\ref{spaces-limits-definition-locally-finite-presentation}（对后一概念，我们将尽可能
使用“保极限”这一术语）。命题
\ref{spaces-limits-proposition-characterize-locally-finite-presentation} 将证明这两个
定义一致。

\begin{lemma}
\label{spaces-limits-lemma-characterize-relative-limit-preserving}
设 $S$ 为概形，$a : F \to G$ 为函子
$(\Sch/S)_{fppf}^{opp} \to \textit{Sets}$ 之间的自然变换。
下列条件等价：
\begin{enumerate}
\item $a : F \to G$ 保极限；
\item 对 $S$ 上仿射概形的任意有向逆系统 $T_i$，只要其极限
$T = \lim T_i$ 是 $S$ 上的仿射概形，下列集合图就是纤维积图：
$$
\xymatrix{
\colim_i F(T_i) \ar[r] \ar[d]_a & F(T) \ar[d]^a \\
\colim_i G(T_i) \ar[r] & G(T)
}
$$
\end{enumerate}
\end{lemma}

\begin{proof}
假设 (1)。考虑 (2) 中的 $T = \lim_{i \in I} T_i$。设
$(y, x_T)$ 为纤维积
$\colim_i G(T_i) \times_{G(T)} F(T)$
的一个元素。于是对某个 $i$，$y$ 来自 $y_i \in G(T_i)$。
考虑定义 \ref{spaces-limits-definition-locally-finite-presentation} 中
$(\Sch/T_i)_{fppf}$ 上的函子 $F_{y_i}$。可见
$x_T \in F_{y_i}(T)$。此外，$T = \lim_{i' \geq i} T_{i'}$ 是
$T_i$ 上仿射概形的有向系统。因此由 (1)，$x_T$ 是
$\colim_{i' \geq i} F_{y_i}(T_{i'})$ 中唯一元素 $x$ 的像。
所以 $x$ 是 $\colim F(T_i)$ 中映到数对 $(y,x_T)$ 的唯一元素。
这就证明了 (2)。

\medskip\noindent
假设 (2)。设 $T$ 为概形且 $y_T \in G(T)$。我们必须证明 $F_{y_T}$
保极限。设 $T' = \lim_{i \in I} T'_i$ 是 $T$ 上的仿射概形，且是
$T$ 上仿射概形 $T'_i$ 的有向极限。取 $x_{T'} \in F_{y_T}$。
% P09-ERR-0050
由于 $I$ 是有向集，可以选取 $i \in I$。把 $y_{T'}$ 的像记作
$y_i \in F(T'_i)$。% P09-ERR-0051
于是 $(y_i,x_{T'})$ 是纤维积
$\colim_i G(T'_i) \times_{G(T')} F(T')$
的一个元素。因此由 (2)，存在 $\colim_i F(T'_i)$ 中的唯一元素 $x$
映到 $(y_i,x_{T'})$。显然，$x$ 给出 $\colim_i F_y(T'_i)$ 中映到
$x_{T'}$ 的元素，结论成立。
\end{proof}

\begin{lemma}
\label{spaces-limits-lemma-composition-locally-finite-presentation}
设 $S$ 为 $\Sch_{fppf}$ 中的概形。设
$F, G, H : (\Sch/S)_{fppf}^{opp} \to \textit{Sets}$，并设
$a : F \to G$、$b : G \to H$ 为函子的自然变换。如果 $a$ 和 $b$
都保极限，则
$$
b \circ a : F \longrightarrow H
$$
也保极限。
\end{lemma}

\begin{proof}
令 $T = \lim_{i \in I}T_i$ 如引理
\ref{spaces-limits-lemma-characterize-relative-limit-preserving} 的刻画 (2) 所述。
考虑集合图
$$
\xymatrix{
\colim_i F(T_i) \ar[r] \ar[d]_a & F(T) \ar[d]^a \\
\colim_i G(T_i) \ar[r] \ar[d]_b & G(T) \ar[d]^b \\
\colim_i H(T_i) \ar[r] & H(T)
}
$$
依假设，两个方块都是纤维积方块。因此外侧矩形也是纤维积图，
引理得证。
\end{proof}

\begin{lemma}
\label{spaces-limits-lemma-locally-finite-presentation-permanence}
设 $S$ 为 $\Sch_{fppf}$ 中的概形。设
$F, G, H : (\Sch/S)_{fppf}^{opp} \to \textit{Sets}$，并设
$a : F \to G$、$b : G \to H$ 为函子的自然变换。如果 $b \circ a$
和 $b$ 保极限，则 $a$ 保极限。
\end{lemma}

\begin{proof}
令 $T = \lim_{i \in I}T_i$ 如引理
\ref{spaces-limits-lemma-characterize-relative-limit-preserving} 的刻画 (2) 所述。
考虑集合图
$$
\xymatrix{
\colim_i F(T_i) \ar[r] \ar[d]_a & F(T) \ar[d]^a \\
\colim_i G(T_i) \ar[r] \ar[d]_b & G(T) \ar[d]^b \\
\colim_i H(T_i) \ar[r] & H(T)
}
$$
依假设，下方方块和外侧矩形都是集合的纤维积。因此上方方块也是
纤维积方块，引理得证。
\end{proof}

\begin{lemma}
\label{spaces-limits-lemma-base-change-locally-finite-presentation}
设 $S$ 为 $\Sch_{fppf}$ 中的概形。设
$F, G, H : (\Sch/S)_{fppf}^{opp} \to \textit{Sets}$，并设
$a : F \to G$、$b : H \to G$ 为函子的自然变换。考虑纤维积图
$$
\xymatrix{
H \times_{b, G, a} F \ar[r]_-{b'} \ar[d]_{a'} & F \ar[d]^a \\
H \ar[r]^b & G
}
$$
如果 $a$ 保极限，则其基变换 $a'$ 也保极限。
\end{lemma}

\begin{proof}
略。提示：这是形式的。
\end{proof}

\begin{lemma}
\label{spaces-limits-lemma-fibre-product-locally-finite-presentation}
设 $S$ 为 $\Sch_{fppf}$ 中的概形。设
$E, F, G, H : (\Sch/S)_{fppf}^{opp} \to \textit{Sets}$，并设
$a : F \to G$、$b : H \to G$ 及 $c : G \to E$ 为函子的自然变换。
如果 $c$、$c \circ a$ 和 $c \circ b$ 保极限，则
$F \times_G H \to E$ 也保极限。
\end{lemma}

\begin{proof}
令 $T = \lim_{i \in I}T_i$ 如引理
\ref{spaces-limits-lemma-characterize-relative-limit-preserving} 的刻画 (2) 所述。
于是
$$
\colim (F \times_G H)(T_i) =
\colim F(T_i) \times_{\colim G(T_i)} \colim H(T_i)
$$
因为滤过余极限与有限积交换。因此我们的目标是证明
$$
\xymatrix{
\colim F(T_i) \times_{\colim G(T_i)} \colim H(T_i) \ar[r] \ar[d] &
F(T) \times_{G(T)} H(T) \ar[d] \\
\colim_i E(T_i) \ar[r] & E(T)
}
$$
是纤维积图。为此只需注意，对于集合映射 $E' \to E$、$F \to G$、
$H \to G$ 及 $G \to E$，有
$$
E' \times_E (F \times_G H) =
(E' \times_E F) \times_{(E' \times_E G)} (E' \times_E H)
$$
略去若干细节。
\end{proof}

\begin{lemma}
\label{spaces-limits-lemma-sheafify-finite-presentation}
设 $S$ 为 $\Sch_{fppf}$ 中的概形，
$F : (\Sch/S)_{fppf}^{opp} \to \textit{Sets}$ 为函子。
如果 $F$ 保极限，则其层化 $F^\#$ 也保极限。
\end{lemma}

\begin{proof}
假设 $F$ 保极限。由于 $F^\#=(F^+)^+$（见《位点》定理
\ref{sites-theorem-plus}），只需证明 $F^+$ 保极限。设 $T$ 为 $S$ 上
的仿射概形，并把 $T=\lim T_i$ 写成 $S$ 上仿射概形的一个逆系统的
有向极限。回忆 $F^+(T)$ 是 $\check H^0(\mathcal V,F)$ 的余极限，
其中该极限遍历 $(\Sch/S)_{fppf}$ 中 $T$ 的所有覆盖。% P09-ERR-0052
仿射概形的任意 fppf 覆盖都能由一个标准 fppf 覆盖加细，见《位点》
引理 \ref{topologies-lemma-fppf-affine}。因此可写成
$$
F^+(T)
=
\colim_{\mathcal{V}\text{ 标准覆盖 }T}
\check H^0(\mathcal{V}, F).
$$
余极限中的任意
$\mathcal V=\{T_k\to T\}_{k=1,\ldots,n}$，都可对某个 $i$ 写成
$\mathcal V_i\times_{T_i}T$，其中
$\mathcal V_i=\{T_{i,k}\to T_i\}_{k=1,\ldots,n}$ 是 $T_i$ 的一个
标准 fppf 覆盖。对 $i'\geq i$，记其基变换为
$\mathcal V_{i'}=\{T_{i',k}\to T_{i'}\}_{k=1,\ldots,n}$。于是
\begin{align*}
% P09-ERR-0053
\colim_{i' \geq i} \check H^0(\mathcal{V}_i, F)
& =
\colim_{i' \geq i}
\text{等化子}
\left(
\xymatrix{
\prod F(T_{i', k})
\ar@<1ex>[r] \ar@<-1ex>[r] &
\prod F(T_{i', k} \times_{T_{i'}} T_{i', l})
}
\right)
\\
& =
\text{等化子}
\left(
\xymatrix{
\colim_{i' \geq i}
\prod F(T_{i', k})
\ar@<1ex>[r] \ar@<-1ex>[r] &
% P09-ERR-0054
\colim_{k' \geq k}
\prod F(T_{i', k} \times_{T_{i'}} T_{i', l})
}
\right)
\\
& =
\text{等化子}
\left(
\xymatrix{
\prod F(T_k)
\ar@<1ex>[r] \ar@<-1ex>[r] &
\prod F(T_k \times_T T_l)
}
\right)
\\
& =
\check H^0(\mathcal{V}, F)
\end{align*}
这里第二个等号成立，是因为滤过余极限正合。第三个等号成立，是因为
$F$ 保极限，并且由《极限》引理 \ref{limits-lemma-scheme-over-limit}，
有 $\lim_{i'\geq i}T_{i',k}=T_k$ 及
$\lim_{i'\geq i}T_{i',k}\times_{T_{i'}}T_{i',l}=T_k\times_TT_l$。
同时对所有覆盖应用这一点，得到
\begin{align*}
F^+(T)
& =
\colim_{\mathcal{V}\text{ 标准覆盖 }T} \check H^0(\mathcal{V}, F) \\
& =
\colim_{i \in I}
\colim_{\mathcal{V}_i\text{ 标准覆盖 }T_i}
\check H^0(T \times_{T_i}\mathcal{V}_i, F) \\
& =
\colim_{i \in I} F^+(T_i)
\end{align*}
由《范畴》引理 \ref{categories-lemma-colimits-commute}，可以交换这些
余极限的次序。
\end{proof}

\begin{lemma}
\label{spaces-limits-lemma-sheaf-finite-presentation}
设 $S$ 为概形，$F : (\Sch/S)_{fppf}^{opp} \to \textit{Sets}$ 为函子。
假设
\begin{enumerate}
\item $F$ 是层；
\item 存在一个 fppf 覆盖 $\{U_j\to S\}_{j\in J}$，使得
$F|_{(\Sch/U_j)_{fppf}}$ 保极限。
\end{enumerate}
那么 $F$ 保极限。
\end{lemma}

\begin{proof}
设 $T$ 为 $S$ 上的仿射概形，$I$ 为有向集，并设 $T_i$ 为 $S$ 上
仿射概形的逆系统，且 $T=\lim T_i$。我们必须证明典范映射
$\colim F(T_i)\to F(T)$ 是双射。

\medskip\noindent
选取某个 $0\in I$，再选取一个标准 fppf 覆盖
$\{V_{0,k}\to T_0\}_{k=1,\ldots,m}$，它加细给定的 $S$ 的 fppf 覆盖
拉回得到的 $\{U_j\times_ST_0\to T_0\}$。对每个 $i\geq0$，令
$V_{i,k}=T_i\times_{T_0}V_{0,k}$，并令
$V_k=T\times_{T_0}V_{0,k}$。注意
$V_k=\lim_{i\geq0}V_{i,k}$，见《极限》引理
\ref{limits-lemma-scheme-over-limit}。

\medskip\noindent
假设 $x,x'\in\colim F(T_i)$ 映到 $F(T)$ 中的同一元素。设对某个
$i\in I$，$x,x'$ 分别由元素 $x_i,x'_i\in F(T_i)$ 给出（由于 $I$
有向，可以为二者选取同一个 $i$）。由假设 (2) 以及 $x_i,x'_i$
映到 $F(T)$ 中同一元素这一事实，可知
$$
x_i|_{V_{i', k}} = x'_i|_{V_{i', k}}
$$
对某个充分大的 $i'\in I$ 成立。由于 $k$ 遍历有限集
$\{1,\ldots,m\}$，可以对每个 $k$ 选取同一个 $i'$。因为
$\{V_{i',k}\to T_{i'}\}$ 是 fppf 覆盖且 $F$ 是层，所以如所需有
$x_i|_{T_{i'}}=x'_i|_{T_{i'}}$。这证明了映射
$\colim F(T_i)\to F(T)$ 是单射。

\medskip\noindent
为证明满射性，作类似论证。设 $x\in F(T)$。由假设 (2)，对每个 $k$
都可选取一个 $i$，使得 $x|_{V_k}$ 来自某个元素
$x_{i,k}\in F(V_{i,k})$。如前，可选取一个对所有 $k$ 都适用的 $i$。
由上面已经证明的单射性可知
$$
x_{i, k}|_{V_{i', k} \times_{T_{i'}} V_{i', l}}
=
x_{i, l}|_{V_{i', k} \times_{T_{i'}} V_{i', l}}
$$
对某个充分大的 $i'$ 成立。因此由 $F$ 的层条件，各元素
$x_{i,k}|_{V_{i',k}}$ 可以粘合为所需的元素 $x_{i'}\in F(T_{i'})$。
\end{proof}

\begin{lemma}
\label{spaces-limits-lemma-sheafify-finite-presentation-map}
设 $S$ 为 $\Sch_{fppf}$ 中的概形，并设
$F,G:(\Sch/S)_{fppf}^{opp}\to\textit{Sets}$ 为函子。如果自然变换
$a:F\to G$ 保极限，则诱导的层之间的自然变换
$F^\#\to G^\#$ 也保极限。
\end{lemma}

\begin{proof}
设 $T$ 为概形且 $y\in G^\#(T)$。我们必须证明函子
$F^\#_y : (\Sch/T)_{fppf}^{opp} \to \textit{Sets}$
——它由 $F^\#\to G^\#$ 和 $y$ 按照定义
\ref{spaces-limits-definition-locally-finite-presentation} 构造——保极限。由公式
(\ref{spaces-limits-equation-fibre-map-functors}) 可见 $F^\#_y$ 是层。选取一个
fppf 覆盖 $\{V_j\to T\}_{j\in J}$，使得 $y|_{V_j}$ 来自某个元素
$y_j\in F(V_j)$。% P09-ERR-0055
注意，$F^\#$ 在 $(\Sch/V_j)_{fppf}$ 上的限制正是
$F^\#_{y_j}$。% P09-ERR-0056
如果能证明 $F^\#_{y_j}$ 保极限，那么引理
\ref{spaces-limits-lemma-sheaf-finite-presentation} 保证 $F^\#_y$ 保极限，结论成立。
这把问题约化到 $y\in G(T)$ 的情形。

\medskip\noindent
设 $y\in G(T)$。此时我们断言 $F^\#_y=(F_y)^\#$。这由公式
(\ref{spaces-limits-equation-fibre-map-functors}) 得出。因此这种情形由引理
\ref{spaces-limits-lemma-sheafify-finite-presentation} 得出。
\end{proof}

\begin{proposition}
\label{spaces-limits-proposition-characterize-locally-finite-presentation}
设 $S$ 为概形，$f:X\to Y$ 为 $S$ 上代数空间的态射。下列条件等价：
\begin{enumerate}
\item $f$ 是局部有限表示的代数空间态射，见《空间的态射》定义
\ref{spaces-morphisms-definition-locally-finite-presentation}；
\item 作为函子的自然变换，态射 $f:X\to Y$ 保极限，见定义
\ref{spaces-limits-definition-locally-finite-presentation}。
\end{enumerate}
\end{proposition}

\begin{proof}
假设 (1)。设 $T$ 为概形且 $y\in Y(T)$。我们必须证明
$T\times_YX$ 在定义 \ref{spaces-limits-definition-locally-finite-presentation} 的意义下
在 $T$ 上保极限。因此问题约化为证明：如果代数空间 $X$ 作为代数空间
在 $S$ 上局部有限表示，那么作为函子
$X : (\Sch/S)_{fppf}^{opp} \to \textit{Sets}$.
它保极限。为此选取一个展示 $X=U/R$，见《空间》定义
\ref{spaces-definition-presentation}。由《空间的态射》定义
\ref{spaces-morphisms-definition-locally-finite-presentation}，$U$ 和 $R$
都是在 $S$ 上局部有限表示的概形。因此由《极限》命题
\ref{limits-proposition-characterize-locally-finite-presentation}，只要
$T=\lim_iT_i$ 位于 $(\Sch/S)_{fppf}$ 中，就有
$$
U(T) = \colim U(T_i), \quad
R(T) = \colim R(T_i)
$$
。由此可知预层
$$
(\Sch/S)_{fppf}^{opp} \longrightarrow \textit{Sets}, \quad
W \longmapsto U(W)/R(W)
$$
保极限。因此由引理 \ref{spaces-limits-lemma-sheafify-finite-presentation}，其层化
$X=U/R$ 也保极限。

\medskip\noindent
假设 (2)。选取概形 $V$ 和满平展态射 $V\to Y$。再选取概形 $U$
和满平展态射 $U\to V\times_YX$。由引理
\ref{spaces-limits-lemma-base-change-locally-finite-presentation}，函子间的自然变换
$V\times_YX\to V$ 保极限。由《空间的态射》引理
\ref{spaces-morphisms-lemma-etale-locally-finite-presentation}，代数空间的
态射 $U\to V\times_YX$ 局部有限表示，因而由证明的第一部分，
它作为函子的自然变换保极限。再由引理
\ref{spaces-limits-lemma-composition-locally-finite-presentation}，复合
$U\to V\times_YX\to V$ 作为函子的自然变换保极限。因此由《极限》
命题 \ref{limits-proposition-characterize-locally-finite-presentation}，
概形态射 $U\to V$ 局部有限表示（这里暂且不论集合论问题，见证明的
最后一段）。按照定义，这意味着 (1) 成立。

\medskip\noindent
集合论说明。设 $U\to V$ 为 $(\Sch/S)_{fppf}$ 中的态射。在《极限》
命题 \ref{limits-proposition-characterize-locally-finite-presentation} 的
陈述中，$U\to V$ 局部有限表示的刻画要求：对 $V$ 上仿射概形的
{\it 每个}有向逆系统 $(T_i,f_{ii'})$，都有
$U(T)=\colim V(T_i)$；% P09-ERR-0057
但在当前情形中，我们只能考虑 $V$ 上那些（同构于）
$(\Sch/S)_{fppf}$ 中对象的仿射概形 $T_i$。所以必须保证
$(\Sch/S)_{fppf}$ 中有足够多的仿射对象，使证明成立。考察《极限》
命题 \ref{limits-proposition-characterize-locally-finite-presentation} 中
(2) $\Rightarrow$ (1) 的证明，可见问题约化到 $U,V$ 都仿射的情形。
设 $U=\Spec(A)$、$V=\Spec(B)$。由 $(\Sch/S)_{fppf}$ 的构造，基数
$\leq|B|$ 的任意环的谱都同构于 $(\Sch/S)_{fppf}$ 中的一个对象。
因此只需注意，《代数》引理
\ref{algebra-lemma-characterize-finite-presentation} 的证明中“仅当”部分
只使用基数 $\leq|B|$ 的 $A$-代数。
\end{proof}

\begin{remark}
\label{spaces-limits-remark-limit-preserving}
下面是命题 \ref{spaces-limits-proposition-characterize-locally-finite-presentation} 的
一个重要特例。设 $S$ 为概形，$X$ 为 $S$ 上的代数空间。那么 $X$
在 $S$ 上局部有限表示，当且仅当 $X$ 作为函子
$(\Sch/S)^{opp}\to\textit{Sets}$ 保极限。与《极限》注
\ref{limits-remark-limit-preserving} 比较。事实上，下面的引理
\ref{spaces-limits-lemma-surjection-is-enough} 将说明，只要映射
$$
\colim X(T_i) \longrightarrow X(T)
$$
在 $T=\lim T_i$ 是 $S$ 上仿射概形的有向极限时恒为满射，就已经足够。
\end{remark}

\begin{lemma}
\label{spaces-limits-lemma-surjection-is-enough}
设 $S$ 为概形，$f:X\to Y$ 为 $S$ 上代数空间的态射。如果对于 $S$
上仿射概形的每个有向极限 $T=\lim_{i\in I}T_i$，映射
$$
\colim X(T_i) \longrightarrow X(T) \times_{Y(T)} \colim Y(T_i)
$$
都是满射，则 $f$ 局部有限表示。换言之，在命题
\ref{spaces-limits-proposition-characterize-locally-finite-presentation} 的 (2) 中，
只需检验引理 \ref{spaces-limits-lemma-characterize-relative-limit-preserving} 的
判据中的满射性。
\end{lemma}

\begin{proof}
选取概形 $V$ 和满平展态射 $g:V\to Y$。再选取概形 $U$ 和满平展态射
$h:U\to V\times_YX$。只需证明，对于引理中的 $T=\lim T_i$，映射
$$
\colim U(T_i) \longrightarrow U(T) \times_{V(T)} \colim V(T_i)
$$
是满射；因为这样一来，由《极限》引理
\ref{limits-lemma-surjection-is-enough}，$U\to V$ 就局部有限表示
（集合论说明与命题
\ref{spaces-limits-proposition-characterize-locally-finite-presentation} 的证明完全相同）。
因此取 $a:T\to U$ 和 $b_i:T_i\to V$，二者确定同一个态射
$T\to V$。图示为
$$
\xymatrix{
T \ar[d]_a \ar[rr]_{p_i} & & T_i \ar[d]^{b_i} \ar@{..>}[ld] \\
U \ar[r]^-h & X \times_Y V \ar[d] \ar[r] & V \ar[d]^g \\
& X \ar[r]^f & Y
}
$$
由引理的假设，增大 $i$ 后可以找到态射 $c_i:T_i\to X$，使得
$h\circ a=(b_i,c_i)\circ p_i:T_i\to V\times_YX$% P09-ERR-0058
且 $f\circ c_i=g\circ b_i$。因为 $h$ 是代数空间的平展态射（因而
局部有限表示），由命题
\ref{spaces-limits-proposition-characterize-locally-finite-presentation}，映射
$$
\colim U(T_i) \longrightarrow U(T) \times_{(X \times_Y V)(T)}
\colim (X \times_Y V)(T_i)
$$
是满射。因此再次增大 $i$ 后，可以找到所需的态射 $a_i:T_i\to U$，
使得 $a=a_i\circ p_i$ 且 $b_i=(U\to V)\circ a_i$。
\end{proof}















\section{代数空间的极限}
\label{spaces-limits-section-limits}

\noindent
下面的引理说明本章如何理解代数空间的极限。我们将不再另行说明地
使用如下事实：代数空间的仿射态射的基变换仍为仿射态射（见
《空间的态射》引理 \ref{spaces-morphisms-lemma-base-change-affine}）。

\begin{lemma}
\label{spaces-limits-lemma-directed-inverse-system-has-limit}
设 $S$ 为概形，$I$ 为有向集。设 $(X_i,f_{ii'})$ 是 $S$ 上代数空间
范畴中以 $I$ 为指标的逆系统。如果态射 $f_{ii'}:X_i\to X_{i'}$
都是仿射的，那么极限 $X=\lim_iX_i$（作为 fppf 层）是代数空间。
此外，
\begin{enumerate}
\item 每个态射 $f_i:X\to X_i$ 都是仿射的；
\item 对任意 $i\in I$ 及代数空间的任意态射 $T\to X_i$，都有
$$
X \times_{X_i} T = \lim_{i' \geq i} X_{i'} \times_{X_i} T.
$$
这里等式是 $S$ 上代数空间之间的等式。
\end{enumerate}
\end{lemma}

\begin{proof}
(2) 是极限 $X=\lim X_i$ 作为 $S$ 上代数空间存在这一事实的形式推论。
选取元素 $0\in I$（这是可能的，因为有向集非空）。选取概形 $U_0$
和满平展态射 $U_0\to X_0$。令 $R_0=U_0\times_{X_0}U_0$，从而
$X_0=U_0/R_0$。对 $i\geq0$，令 $U_i=X_i\times_{X_0}U_0$ 及
$R_i=X_i\times_{X_0}R_0=U_i\times_{X_i}U_i$。由《极限》引理
\ref{limits-lemma-directed-inverse-system-has-limit}，
$U=\lim_{i\geq0}U_i$ 和 $R=\lim_{i\geq0}R_i$ 都是概形。此外，
两个态射 $s,t:R\to U$ 是两个投影 $R_0\to U_0$ 沿态射 $U\to U_0$
的基变换，特别地它们平展。态射 $R\to U\times_SU$ 定义一个等价关系，
因为等价关系的有向极限仍为等价关系。% P09-ERR-0059
所以 $R\to U\times_SU$ 是平展等价关系。我们断言自然映射
\begin{equation}
\label{spaces-limits-equation-isomorphism-sheaves}
U/R \longrightarrow \lim X_i
\end{equation}
是 $S$ 上概形范畴中的 fppf 层的同构。由《空间》定理
\ref{spaces-theorem-presentation}，该断言推出 $X=\lim X_i$ 是代数空间。

\medskip\noindent
设 $Z$ 为概形且 $a:Z\to\lim X_i$ 为态射。于是 $a=(a_i)$，其中
$a_i:Z\to X_i$。令 $W_0=Z\times_{a_0,X_0}U_0$。由上面对
$U_i\to X_i$ 的选取，对所有 $i\geq0$ 都有
$W_0=Z\times_{a_i,X_i}U_i$。因此得到态射
$W_0\to\lim_{i\geq0}U_i=U$。由于 $W_0\to Z$ 满且平展，可知
(\ref{spaces-limits-equation-isomorphism-sheaves}) 是层的满射。最后，设 $Z$ 为概形，
且 $a,b:Z\to U/R$ 是两个被 (\ref{spaces-limits-equation-isomorphism-sheaves})
映为同一态射的态射。我们必须证明 $a=b$。用某个 fppf 覆盖的各成员
替换 $Z$ 后，可以假设存在给出 $a,b$ 的态射 $a',b':Z\to U$。
$a,b$ 被 (\ref{spaces-limits-equation-isomorphism-sheaves}) 映为同一态射这一条件
意味着：对每个 $i\geq0$，复合
$a_i',b_i':Z\to U\to U_i$ 作为到 $U_i/R_i=X_i$ 的态射相等。
因此 $(a_i',b_i'):Z\to U_i\times_SU_i$ 通过 $R_i$ 分解，设分解由
某个态射 $c_i:Z\to R_i$ 给出。由于 $R=\lim_{i\geq0}R_i$，可见
$c=\lim c_i:Z\to R$ 是一个态射，它说明 $a,b$ 作为从 $Z$ 到
$U/R$ 的态射相等。

\medskip\noindent
由上面已经看到 $U_i\times_{X_i}X=U$，而按构造 $U\to U_i$ 仿射，
所以 (1) 成立。
\end{proof}

\begin{lemma}
\label{spaces-limits-lemma-space-over-limit}
设 $S$ 为概形，$I$ 为有向集。设 $(X_i,f_{ii'})$ 是以 $I$ 为指标的
$S$ 上代数空间的逆系统，且转移态射均仿射。令 $X=\lim_iX_i$，
并取 $0\in I$。假设 $T\to X_0$ 是代数空间的态射。那么
$$
T \times_{X_0} X = \lim_{i \geq 0} T \times_{X_0} X_i
$$
这里等式是 $S$ 上代数空间之间的等式。
\end{lemma}

\begin{proof}
由引理 \ref{spaces-limits-lemma-directed-inverse-system-has-limit}，极限 $X$ 是代数空间。
该等式是形式的，见《范畴》引理
\ref{categories-lemma-colimits-commute}。
\end{proof}

\begin{lemma}
\label{spaces-limits-lemma-directed-inverse-system-closed-immersions}
设 $S$ 为概形，$I$ 为有向集。设
$(X_i,f_{i'i})\to(Y_i,g_{i'i})$ 是以 $I$ 为指标的 $S$ 上代数空间
逆系统之间的态射。假设
\begin{enumerate}
\item 态射 $f_{i'i}:X_{i'}\to X_i$ 仿射；
\item 态射 $g_{i'i}:Y_{i'}\to Y_i$ 仿射；
\item 态射 $X_i\to Y_i$ 是闭浸入。
\end{enumerate}
那么 $\lim X_i\to\lim Y_i$ 是闭浸入。
\end{lemma}

\begin{proof}
由引理 \ref{spaces-limits-lemma-directed-inverse-system-has-limit}，$\lim X_i$ 和
$\lim Y_i$ 都存在。取 $0\in I$，并选取仿射概形 $V_0$ 和平展态射
$V_0\to Y_0$。于是态射
$V_i=Y_i\times_{Y_0}V_0\to U_i=X_i\times_{Y_0}V_0$% P09-ERR-0060
是仿射概形的闭浸入。因此态射
$V=Y\times_{Y_0}V_0\to U=X\times_{Y_0}V_0$% P09-ERR-0060
是闭浸入；这是因为 $V=\lim V_i$、$U=\lim U_i$，并且仿射概形的
闭浸入的极限仍为闭浸入：满环同态的滤过余极限仍为满射。当我们改变
$V_0\to Y_0$ 的选取时，平展态射 $V\to Y$ 构成 $Y$ 的平展覆盖，
故引理成立。
\end{proof}

\begin{lemma}
\label{spaces-limits-lemma-directed-inverse-system-reduced}
设 $S$ 为概形，$I$ 为有向集。设 $(X_i,f_{i'i})$ 是以 $I$ 为指标的
$S$ 上代数空间的逆系统。如果每个 $X_i$ 都既约，那么 $X$ 既约。
% P09-ERR-0061
\end{lemma}

\begin{proof}
由引理 \ref{spaces-limits-lemma-directed-inverse-system-has-limit}，$\lim X_i$ 存在。
取 $0\in I$，并选取仿射概形 $V_0$ 和平展态射
$U_0\to X_0$。% P09-ERR-0062
于是仿射概形 $U_i=X_i\times_{X_0}U_0$ 既约。因此
$U=X\times_{X_0}U_0$ 作为既约仿射概形的极限也是既约仿射概形：
既约环的滤过余极限仍既约。当我们改变 $U_0\to X_0$ 的选取时，
平展态射 $U\to X$ 构成 $X$ 的平展覆盖，故引理成立。
\end{proof}

\begin{lemma}
\label{spaces-limits-lemma-better-characterize-relative-limit-preserving}
设 $S$ 为概形，$X\to Y$ 为 $S$ 上代数空间的态射。命题
\ref{spaces-limits-proposition-characterize-locally-finite-presentation} 中的等价条件
(1)、(2) 还等价于
\begin{enumerate}
\item[(3)] 对 $S$ 上拟紧拟分离代数空间 $T_i$ 的每个有向极限
$T=\lim T_i$，如果转移态射仿射，则集合图
$$
\xymatrix{
\colim_i \Mor(T_i, X) \ar[r] \ar[d] & \Mor(T, X) \ar[d] \\
\colim_i \Mor(T_i, Y) \ar[r] & \Mor(T, Y)
}
$$
是纤维积图。
\end{enumerate}
\end{lemma}

\begin{proof}
显然 (3) 蕴含 (2)。我们假设 (2) 并证明 (3)。这个证明相当形式化，
我们鼓励读者自行找出证明。

\medskip\noindent
先证明在另外假设每个 $T_i$ 都分离时 (3) 成立。选取 $i\in I$，
再选取满平展态射 $U_i\to T_i$，其中 $U_i$ 仿射。应用引理
\ref{spaces-limits-lemma-space-over-limit}，令 $U=U_i\times_{T_i}T$ 及
$U_{i'}=U_i\times_{T_i}T_{i'}$，可得
$U=\lim_{i'\geq i}U_{i'}$。当然，$U$ 和 $U_{i'}$ 都仿射（见引理
\ref{spaces-limits-lemma-directed-inverse-system-has-limit}）。由于 $T_i$ 分离，
纤维积 $V_i=U_i\times_{T_i}U_i$ 也是仿射概形；并且得到仿射概形
$V=V_i\times_{T_i}T$ 和 $V_{i'}=V_i\times_{T_i}T_{i'}$，满足
$V=\lim_{i'\geq i}V_{i'}$。注意，$U\to T$ 和 $U_i\to T_i$ 都满且
平展，并且 $V=U\times_TU$、
$V_{i'}=U_{i'}\times_{T_{i'}}U_{i'}$。

注意，$\Mor(T,X)$ 是两个映射 $\Mor(U,X)\to\Mor(V,X)$ 的等化子；
例如，这是因为 $X$ 作为 $(\Sch/S)_{fppf}$ 上的层，是两个映射
$h_V\to h_u$ 的余等化子。% P09-ERR-0063
类似地，$\Mor(T_{i'},X)$ 是两个映射
$\Mor(U_{i'},X)\to\Mor(V_{i'},X)$ 的等化子。当然，把 $X$ 换成 $Y$
也有同样结论。条件 (2) 说明，在 $U=\lim U_i$ 和 $V=\lim V_i$
的情形中，(3) 中的图是纤维积。% P09-ERR-0064
由形式推理可知，对 $T=\lim T_i$ 也有同样结论。

\medskip\noindent
在一般情形中，选取仿射概形 $U$、元素 $i\in I$ 以及满平展态射
$U\to T_i$。重复上一段的论证仍可完成证明：概形 $V_{i'}$、$V$
不再仿射，但仍拟紧且分离，所以可应用上一段的结果。
\end{proof}




\section{性质的下降}
\label{spaces-limits-section-descent}

\noindent
本节是《极限》第 \ref{limits-section-descent} 节的类似版本。

\begin{lemma}
\label{spaces-limits-lemma-inverse-limit-sets}
设 $S$ 为概形，$X=\lim_{i\in I}X_i$ 是 $S$ 上代数空间的有向逆系统
的极限，其转移态射仿射（引理
\ref{spaces-limits-lemma-directed-inverse-system-has-limit}）。如果每个 $X_i$ 都良态
（例如拟分离或局部分离），那么作为集合有 $|X|=\lim_i|X_i|$。
\end{lemma}

\begin{proof}
存在典范映射 $|X|\to\lim|X_i|$。选取 $0\in I$。如果
$W_0\subset X_0$ 是开子空间，那么
$f_0^{-1}W_0 = \lim_{i \geq 0} f_{i0}^{-1}W_0$，见
引理 \ref{spaces-limits-lemma-directed-inverse-system-has-limit}。因此，只要能对
$X_0$ 拟紧的逆系统证明本引理，一般情形便随之成立。所以可以并且确实
假设 $X_0$ 拟紧。

\medskip\noindent
选取仿射概形 $U_0$ 和满平展态射 $U_0\to X_0$。令
$U_i=X_i\times_{X_0}U_0$、$U=X\times_{X_0}U_0$，并令
$R_i=U_i\times_{X_i}U_i$、$R=U\times_XU$。回忆
$U=\lim U_i$ 且 $R=\lim R_i$，见引理
\ref{spaces-limits-lemma-directed-inverse-system-has-limit} 的证明。还回忆
$|X|=|U|/|R|$ 且 $|X_i|=|U_i|/|R_i|$。由《极限》引理
\ref{limits-lemma-topology-limit}，有 $|U|=\lim|U_i|$ 及
$|R|=\lim|R_i|$。

\medskip\noindent
证明 $|X|\to\lim|X_i|$ 的满射性。设 $(x_i)\in\lim|X_i|$。记
$S_i\subset|U_i|$ 为 $x_i$ 的逆像。由良态空间的定义（《良态空间》
定义 \ref{decent-spaces-definition-very-reasonable}），这是有限非空集。
所以 $\lim S_i$ 非空，见《范畴》引理
\ref{categories-lemma-nonempty-limit}。取
$(u_i)\in\lim S_i\subset\lim|U_i|$。由上所述，它确定一点
$u\in|U|$；后者映到某个 $x\in|X|$，而 $x$ 又映到 $\lim|X_i|$
中给定的元素 $(x_i)$。

\medskip\noindent
证明 $|X|\to\lim|X_i|$ 的单射性。假设 $x,x'\in|X|$ 映到
$\lim|X_i|$ 中的同一点。选取提升 $u,u'\in|U|$，并记其像为
$u_i,u'_i\in|U_i|$。对每个 $i$，令 $T_i\subset|R_i|$ 为映到
$(u_i,u'_i)\in|U_i|\times|U_i|$ 的点集。由良态空间的定义
（《良态空间》定义 \ref{decent-spaces-definition-very-reasonable}），
这是有限集。此外，由于假设 $x,x'$ 映到 $X_i$ 中同一点，$T_i$ 非空。
故 $\lim T_i$ 非空，见《范畴》引理
\ref{categories-lemma-nonempty-limit}。如前，令
$r\in|R|=\lim|R_i|$ 为对应于 $\lim T_i$ 中一个元素的点。按构造，
$r$ 在 $|U|\times|U|$ 中映到 $(u,u')$，从而如所需有
$x=x'$ 于 $|X|$ 中。

\medskip\noindent
关于括号中的断言：拟分离代数空间是良态的，见《良态空间》第
\ref{decent-spaces-section-reasonable-decent} 节（其中关键观察是
《空间的性质》引理
\ref{spaces-properties-lemma-finite-fibres-presentation}）。由《良态空间》
引理 \ref{decent-spaces-lemma-locally-separated-decent}，局部分离
代数空间也是良态的。
\end{proof}

\begin{lemma}
\label{spaces-limits-lemma-topology-limit}
在引理 \ref{spaces-limits-lemma-inverse-limit-sets} 的相同记号和假设下，作为拓扑空间
有 $|X|=\lim_i|X_i|$。
\end{lemma}

\begin{proof}
我们使用《拓扑》引理 \ref{topology-lemma-characterize-limit} 的判据。
引理 \ref{spaces-limits-lemma-inverse-limit-sets} 已说明作为集合有
$|X|=\lim_i|X_i|$。映射 $f_i:X\to X_i$ 是代数空间的态射，因而确定
连续映射 $|X|\to|X_i|$。所以对每个开集 $U_i\subset|X_i|$，
$f_i^{-1}(U_i)$ 都开。最后，设 $x\in|X|$，并设
$x\in V\subset|X|$ 是一个开邻域。我们必须找到某个 $i$ 及 $x$ 的像的
开邻域 $W_i\subset|X_i|$，使 $f_i^{-1}(W_i)\subset V$。
% P09-ERR-0065
选取 $0\in I$。选取概形 $U_0$ 和满平展态射 $U_0\to X_0$。令
$U=X\times_{X_0}U_0$，并对 $i\geq0$ 令
$U_i=X_i\times_{X_0}U_0$。由引理
\ref{spaces-limits-lemma-directed-inverse-system-has-limit}，在概形范畴中
$U=\lim_{i\geq0}U_i$。选取映到 $x$ 的 $u\in U$。由概形的相应结果
（《极限》引理 \ref{limits-lemma-inverse-limit-top}），可以找到
$i\geq0$ 和 $u$ 的像的开邻域 $E_i\subset U_i$，使其在 $U$ 中的逆像
包含于 $V$ 在 $U$ 中的逆像。此时可令 $W_i\subset|X_i|$ 等于
$E_i$ 的像。这是可行的，因为 $|U_i|\to|X_i|$ 是开映射。
\end{proof}

\begin{lemma}
\label{spaces-limits-lemma-limit-nonempty}
设 $S$ 为概形，$X=\lim_{i\in I}X_i$ 是 $S$ 上代数空间的有向逆系统
的极限，其转移态射仿射（引理
\ref{spaces-limits-lemma-directed-inverse-system-has-limit}）。如果每个 $X_i$ 都拟紧且
非空，那么 $|X|$ 非空。
\end{lemma}

\begin{proof}
选取 $0\in I$。选取仿射概形 $U_0$ 和满平展态射 $U_0\to X_0$。
令 $U_i=X_i\times_{X_0}U_0$ 及 $U=X\times_{X_0}U_0$。于是每个
$U_i$ 都是非空仿射概形。因此 $U=\lim U_i$ 非空（《极限》引理
\ref{limits-lemma-limit-nonempty}），从而 $X$ 非空。
\end{proof}

\begin{lemma}
\label{spaces-limits-lemma-inverse-limit-irreducibles}
设 $S$ 为概形，$X=\lim_{i\in I}X_i$ 是 $S$ 上代数空间的有向逆系统
的极限，其转移态射仿射（引理
\ref{spaces-limits-lemma-directed-inverse-system-has-limit}）。设 $x\in|X|$，其像为
$x_i\in|X_i|$。如果每个 $X_i$ 都良态，那么作为集合有
$\overline{\{x\}}=\lim_i\overline{\{x_i\}}$；赋予约化诱导概形结构后，
作为代数空间也有此等式。% P09-ERR-0066
\end{lemma}

\begin{proof}
令 $Z=\overline{\{x\}}\subset|X|$ 及
$Z_i=\overline{\{x_i\}}\subset|X_i|$。由于 $|X|\to|X_i|$ 连续，
对每个 $i$，$Z$ 都映入 $Z_i$。又因为作为集合有
$|X|=\lim|X_i|$（引理 \ref{spaces-limits-lemma-inverse-limit-sets}），所以得到单射
$Z\to\lim Z_i$。假设 $x'\in|X|$ 不属于 $Z$。则存在开子集
$U\subset|X|$，使 $x'\in U$ 而 $x\notin U$。由于作为拓扑空间有
$|X|=\lim|X_i|$（引理 \ref{spaces-limits-lemma-topology-limit}），可对某个子集
$J\subset I$ 和开集 $U_j\subset|X_j|$ 写成
$U=\bigcup_{j\in J}f_j^{-1}(U_j)$，见《拓扑》引理
\ref{topology-lemma-describe-limits}。于是对某个 $j\in J$，有
$f_j(x')\in U_j$ 且 $f_j(x)\notin U_j$。换言之，
$f_j(x')\notin Z_j$。所以作为集合有 $Z=\lim Z_i$。

\medskip\noindent
接着赋予 $Z$ 和 $Z_i$ 各自的约化诱导概形结构，见《空间的性质》定义
\ref{spaces-properties-definition-reduced-induced-space}。转移态射
$X_{i'}\to X_i$ 诱导仿射态射 $Z_{i'}\to Z_i$，投影 $X\to X_i$
诱导相容态射 $Z\to Z_i$。因此得到代数空间的态射
$Z\to\lim Z_i\to X$。由引理
\ref{spaces-limits-lemma-directed-inverse-system-closed-immersions}，
$\lim Z_i\to X$ 是闭浸入。由引理
\ref{spaces-limits-lemma-directed-inverse-system-reduced}，代数空间 $\lim Z_i$ 既约。
由上所述，$Z\to\lim Z_i$ 在点上是双射。由约化诱导闭子概形结构的
唯一性，可知这个态射是代数空间的同构。% P09-ERR-0067
\end{proof}

\begin{situation}
\label{spaces-limits-situation-descent}
设 $S$ 为概形，$X=\lim_{i\in I}X_i$ 是 $S$ 上代数空间的有向逆系统
的极限，其转移态射仿射（引理
\ref{spaces-limits-lemma-directed-inverse-system-has-limit}）。假设对每个 $i\in I$，
$X_i$ 都拟紧拟分离。再选取一个元素 $0\in I$。
\end{situation}

\begin{lemma}
\label{spaces-limits-lemma-descend-section}
记号和假设如情形 \ref{spaces-limits-situation-descent}。设 $\mathcal F_0$ 是 $X_0$
上的拟凝聚层。对 $i\geq0$，令
$\mathcal F_i=f_{0i}^*\mathcal F_0$，并令
$\mathcal F=f_0^*\mathcal F_0$。那么
$$
\Gamma(X, \mathcal{F}) = \colim_{i \geq 0} \Gamma(X_i, \mathcal{F}_i)
$$
\end{lemma}

\begin{proof}
选取满平展态射 $U_0\to X_0$，其中 $U_0$ 是仿射概形（《空间的性质》
引理 \ref{spaces-properties-lemma-quasi-compact-affine-cover}）。令
$U_i=X_i\times_{X_0}U_0$。令 $R_0=U_0\times_{X_0}U_0$ 及
$R_i=R_0\times_{X_0}X_i$。在引理
\ref{spaces-limits-lemma-directed-inverse-system-has-limit} 的证明中，我们已经看到
存在展示 $X=U/R$，满足 $U=\lim U_i$ 且 $R=\lim R_i$。注意 $U_i$
和 $U$ 仿射，而 $R_i$ 和 $R$ 拟紧且分离（因为 $X_i$ 拟分离）。
所以《极限》引理 \ref{limits-lemma-descend-section} 推出
$$
\mathcal{F}(U) = \colim \mathcal{F}_i(U_i)
\quad\text{且}\quad
\mathcal{F}(R) = \colim \mathcal{F}_i(R_i).
$$
由于 $\Gamma(X,\mathcal F)=\Ker(\mathcal F(U)\to\mathcal F(R))$，
并且类似地
$\Gamma(X_i, \mathcal{F}_i) =
\Ker(\mathcal{F}_i(U_i) \to \mathcal{F}_i(R_i))$
，引理随之成立。
\end{proof}

\begin{lemma}
\label{spaces-limits-lemma-descend-opens}
记号和假设如情形 \ref{spaces-limits-situation-descent}。对任意拟紧开子空间
$U\subset X$，存在某个 $i$ 及拟紧开集 $U_i\subset X_i$，使得
$U_i$ 在 $X$ 中的逆像是 $U$。
\end{lemma}

\begin{proof}
这由引理 \ref{spaces-limits-lemma-directed-inverse-system-has-limit} 中的极限构造
以及概形的相应结果——《极限》引理
\ref{limits-lemma-descend-opens}——形式地得出。
\end{proof}

\noindent
下面的引理将被更强的引理 \ref{spaces-limits-lemma-descend-isomorphism} 取代。

\begin{lemma}
\label{spaces-limits-lemma-descend-equality}
记号和假设如情形 \ref{spaces-limits-situation-descent}。设 $f_0:Y_0\to Z_0$ 为
$X_0$ 上代数空间的态射。假设 (a) $Y_0\to X_0$ 和 $Z_0\to X_0$
可表；(b) $Y_0,Z_0$ 拟紧拟分离；(c) $f_0$ 局部有限表示；并且
(d) $Y_0\times_{X_0}X\to Z_0\times_{X_0}X$ 是同构。那么存在
$i\geq0$，使 $Y_0\times_{X_0}X_i\to Z_0\times_{X_0}X_i$ 是同构。
\end{lemma}

\begin{proof}
选取仿射概形 $U_0$ 和满平展态射 $U_0\to X_0$。令
$U_i=U_0\times_{X_0}X_i$ 及 $U=U_0\times_{X_0}X$。应用《极限》
引理 \ref{limits-lemma-descend-isomorphism} 可知，对某个 $i\geq0$，
$Y_0\times_{X_0}U_i\to Z_0\times_{X_0}U_i$ 是概形的同构（细节略）。
由于 $U_i\to X_i$ 满且平展，因而
$Y_0\times_{X_0}X_i\to Z_0\times_{X_0}X_i$ 是同构（细节略）。
\end{proof}

\begin{lemma}
\label{spaces-limits-lemma-descend-separated}
记号和假设如情形 \ref{spaces-limits-situation-descent}。如果 $X$ 分离，则对某个
$i\in I$，$X_i$ 分离。
\end{lemma}

\begin{proof}
选取仿射概形 $U_0$ 和满平展态射 $U_0\to X_0$。对 $i\geq0$，令
$U_i=U_0\times_{X_0}X_i$，并令 $U=U_0\times_{X_0}X$。注意，
$U_i,U$ 是仿射概形，且配有满平展态射 $U_i\to X_i$ 和 $U\to X$。
令 $R_i=U_i\times_{X_i}U_i$、$R=U\times_XU$，相应投影为
$s_i,t_i:R_i\to U_i$ 及 $s,t:R\to U$。注意，$R_i,R$ 是拟紧分离
概形（因为代数空间 $X_i,X$ 拟分离）。映射 $s_i:R_i\to U_i$ 和
$s:R\to U$ 都有限型。

按照定义，$X_i$ 分离当且仅当
$(t_i,s_i):R_i\to U_i\times U_i$ 是闭浸入；而由于假设 $X$ 分离，
态射 $(t,s):R\to U\times U$ 是闭浸入。由于 $R\to U$ 有限型，
存在某个 $i$，使态射 $R\to U_i\times U$ 是闭浸入（《极限》引理
\ref{limits-lemma-finite-type-eventually-closed}）。固定这样的 $i\in I$。
对 $i'\geq i$ 的态射系统 $R_{i'}\to U_i\times U_{i'}$ 应用《极限》
引理 \ref{limits-lemma-descend-closed-immersion-finite-presentation}
（这是允许的，因为确有
$R_{i'}=R_i\times_{U_i\times U_i}U_i\times U_{i'}$），可知当 $i'$
充分大时，$R_{i'}\to U_i\times U_{i'}$ 是闭浸入。这立即推出
$R_{i'}\to U_{i'}\times U_{i'}$ 是闭浸入，从而完成引理的证明。
\end{proof}

\begin{lemma}
\label{spaces-limits-lemma-limit-is-affine}
记号和假设如情形 \ref{spaces-limits-situation-descent}。如果 $X$ 仿射，则存在某个
$i$，使 $X_i$ 仿射。
\end{lemma}

\begin{proof}
选取 $0\in I$。选取仿射概形 $U_0$ 和满平展态射 $U_0\to X_0$。
令 $U=U_0\times_{X_0}X$，并对 $i\geq0$ 令
$U_i=U_0\times_{X_0}X_i$。由于转移态射仿射，代数空间 $U_i,U$
都仿射。因此 $U\to X$ 是仿射概形之间的平展态射。于是可写
$X=\Spec(A)$、$U=\Spec(B)$ 及
$$
B = A[x_1, \ldots, x_n]/(g_1, \ldots, g_n)
$$
使得 $\Delta=\det(\partial g_\lambda/\partial x_\mu)$ 在 $B$ 中可逆，
见《代数》引理 \ref{algebra-lemma-etale-standard-smooth}。令
$A_i=\mathcal O_{X_i}(X_i)$。由引理 \ref{spaces-limits-lemma-descend-section}，
$A=\colim A_i$。增大 $0$ 后，可以假设存在
$g_{1,i},\ldots,g_{n,i}\in A_i[x_1,\ldots,x_n]$ 映到
$g_1,\ldots,g_n$。令
$$
B_i = A_i[x_1, \ldots, x_n]/(g_{1, i}, \ldots, g_{n, i})
$$
对所有 $i\geq0$ 如此定义。必要时增大 $0$，可以假设对所有 $i\geq0$，
$\Delta_i=\det(\partial g_{\lambda,i}/\partial x_\mu)$ 在 $B_i$ 中可逆。
所以 $A_i\to B_i$ 是平展环同态。再次增大 $0$，还可假设
$\Spec(B_i)\to\Spec(A_i)$ 满，见《极限》引理
\ref{limits-lemma-descend-surjective}。又一次增大 $0$ 后，可以选取元素
$h_{1,i},\ldots,h_{n,i}\in\mathcal O_{U_i}(U_i)$，它们映到
$B=\mathcal O_U(U)$ 中 $x_1,\ldots,x_n$ 的类，并使
$g_{\lambda,i}(h_{\nu,i})=0$ 于 $\mathcal O_{U_i}(U_i)$ 中。
% P09-ERR-0068
于是得到交换图
\begin{equation}
\label{spaces-limits-equation-to-show-cartesian}
\vcenter{
\xymatrix{
X_i \ar[d] & U_i \ar[l] \ar[d] \\
\Spec(A_i) & \Spec(B_i) \ar[l]
}
}
\end{equation}
按构造，$B_i=B_0\otimes_{A_0}A_i$ 且
$B=B_0\otimes_{A_0}A$。考虑态射
$$
f_0 : U_0 \longrightarrow X_0 \times_{\Spec(A_0)} \Spec(B_0)
$$
这是拟紧拟分离代数空间之间的态射，并且在 $X_0$ 上可表、分离且平展。
按我们的选取，$f_0$ 到 $X$ 的基变换是同构。因此由引理
\ref{spaces-limits-lemma-descend-equality}，存在某个 $i$，使 $f_0$ 到 $X_i$ 的
基变换是同构；换言之，图 (\ref{spaces-limits-equation-to-show-cartesian}) 是笛卡尔图。
于是，把《下降》引理 \ref{descent-lemma-descent-data-sheaves} 应用于
fppf 覆盖 $\{\Spec(B_i)\to\Spec(A_i)\}$，并结合《下降》引理
\ref{descent-lemma-affine}，可知 $X_i\to\Spec(A_i)$ 可由一个在
$\Spec(A_i)$ 上仿射的概形表示，正合所需。（当然，这还推出
$X_i=\Spec(A_i)$，但我们不需要这一点。）
\end{proof}

\begin{lemma}
\label{spaces-limits-lemma-limit-is-scheme}
记号和假设如情形 \ref{spaces-limits-situation-descent}。如果 $X$ 是概形，则存在某个
$i$，使 $X_i$ 是概形。
\end{lemma}

\begin{proof}
选取有限仿射开覆盖 $X=\bigcup W_j$。由引理
\ref{spaces-limits-lemma-descend-opens}，可以找到 $i\in I$ 及开子空间
$W_{j,i}\subset X_i$，其到 $X$ 的基变换为 $W_j\to X$。由引理
\ref{spaces-limits-lemma-limit-is-affine}，可假设每个 $W_{j,i}$ 都是仿射概形。
这意味着 $X_i$ 是概形（例如见《空间的性质》第
\ref{spaces-properties-section-schematic} 节）。
\end{proof}

\begin{lemma}
\label{spaces-limits-lemma-finite-type-eventually-closed}
设 $S$ 为概形，$B$ 为 $S$ 上的代数空间。设 $X=\lim X_i$ 是 $B$
上代数空间的有向极限，且转移态射仿射。设 $Y\to X$ 为 $B$ 上
代数空间的态射。
\begin{enumerate}
\item 如果 $Y\to X$ 是闭浸入，$X_i$ 拟紧，且 $Y\to B$ 局部有限型，
那么当 $i$ 充分大时，$Y\to X_i$ 是闭浸入。
\item 如果 $Y\to X$ 是浸入，$X_i$ 拟分离，$Y\to B$ 局部有限型，
且 $Y$ 拟紧，那么当 $i$ 充分大时，$Y\to X_i$ 是浸入。
\item 如果 $Y\to X$ 是同构，$X_i$ 拟紧，$X_i\to B$ 局部有限型，
转移态射 $X_{i'}\to X_i$ 是闭浸入，且 $Y\to B$ 局部有限表示，
那么当 $i$ 充分大时，$Y\to X_i$ 是同构。
\item 如果 $Y\to X$ 是单态射，$X_i$ 拟分离，$Y\to B$ 局部有限型，
且 $Y$ 拟紧，那么当 $i$ 充分大时，$Y\to X_i$ 是单态射。
\end{enumerate}
\end{lemma}

\begin{proof}
证明 (1)。选取 $0\in I$。由于 $X_0$ 拟紧，可以选取仿射概形 $W$
和平展态射 $W\to B$，使 $|X_0|\to|B|$ 的像包含在
$|W|\to|B|$ 的像中。选取仿射概形 $U_0$ 和平展态射
$U_0\to X_0\times_BW$，使 $U_0\to X_0$ 满。（由 $W$ 的选取和
$X_0$ 的拟紧性，这是可能的；细节略。）对 $i\geq0$，分别令
$V\to Y$、$U\to X$、$U_i\to X_i$ 为 $U_0\to X_0$ 的基变换。
只需证明当 $i$ 充分大时 $V\to U_i$ 是闭浸入。因此约化为在 $W$
上证明 $V\to U=\lim U_i$ 的结果。这由概形的情形——《极限》引理
\ref{limits-lemma-finite-type-eventually-closed}——得出。

\medskip\noindent
证明 (2)。选取 $0\in I$。选取拟紧开子空间 $X'_0\subset X_0$，使
$Y\to X_0$ 通过 $X'_0$ 分解。对 $i\geq0$，用 $X'_0$ 的逆像替换
$X_i$ 后，可以假设所有 $X_i'$ 都拟紧拟分离。% P09-ERR-0069
设 $U\subset X$ 是拟紧开集，使 $Y\to X$ 通过闭浸入 $Y\to U$ 分解
（这样的 $U$ 存在，因为 $Y$ 拟紧）。由引理
\ref{spaces-limits-lemma-descend-opens}，可假设 $U=\lim U_i$，其中
$U_i\subset X_i$ 是拟紧开集。由 (1)，对某个 $i$，$Y\to U_i$
是闭浸入。所以 (2) 成立。

\medskip\noindent
证明 (3)。选取 $0\in I$。选取仿射概形 $U_0$ 和满平展态射
$U_0\to X_0$。令 $U_i=X_i\times_{X_0}U_0$，
$U=X\times_{X_0}U_0=Y\times_{X_0}U_0$。于是 $U=\lim U_i$ 是仿射
概形的极限，该系统的转移映射是闭浸入，并且 $U\to U_0$ 有限表示
（因为 $U\to B$ 局部有限表示，$U_0\to B$ 局部有限型，并应用
《空间的态射》引理
\ref{spaces-morphisms-lemma-finite-presentation-permanence}）。
因此问题约化为下列代数事实：如果 $A=\lim A_i$ 是转移映射为满射的
$R$-代数的有向余极限，且 $A$ 在 $A_0$ 上有限表示，那么对某个 $i$，
$A=A_i$。事实上，写成 $A=A_0/(f_1,\ldots,f_n)$，再选取 $i$，使
$f_1,\ldots,f_n$ 在满射 $A_0\to A_i$ 下映为零。

\medskip\noindent
证明 (4)。令 $Z_i=Y\times_{X_i}Y$。由于转移态射
$X_{i'}\to X_i$ 仿射因而分离，转移态射 $Z_{i'}\to Z_i$ 是闭浸入，
见《空间的态射》引理
\ref{spaces-morphisms-lemma-fibre-product-after-map}。因为 $Y\to X$
是单态射，所以 $\lim Z_i=Y\times_XY=Y$。选取 $0\in I$。由于
$Y\to X_0$ 局部有限型（《空间的态射》引理
\ref{spaces-morphisms-lemma-permanence-finite-type}），态射 $Y\to Z_0$
局部有限表示（《空间的态射》引理
\ref{spaces-morphisms-lemma-diagonal-morphism-finite-type}）。态射
$Z_i\to Z_0$ 局部有限型（它们是闭浸入）。最后，因为 $X_i$ 拟分离
且 $Y$ 拟紧，$Z_i=Y\times_{X_i}Y$ 拟紧。因此可把 (3) 应用于
$Z_0$ 上的 $Y=\lim_{i\geq0}Z_i$，得到对某个 $i$ 有 $Y=Z_i$。
这就证明了 (4) 以及整个引理。
\end{proof}

\begin{lemma}
\label{spaces-limits-lemma-eventually-separated}
设 $S$ 为概形，$Y$ 为 $S$ 上的代数空间。设 $X=\lim X_i$ 是 $Y$
上代数空间的有向极限，且转移态射仿射。假设
\begin{enumerate}
\item $Y$ 拟分离；
\item $X_i$ 拟紧拟分离；
\item 态射 $X\to Y$ 分离。
\end{enumerate}
那么对所有充分大的 $i$，$X_i\to Y$ 分离。
\end{lemma}

\begin{proof}
取 $0\in I$。选取仿射概形 $W$ 和平展态射 $W\to Y$，使
$|W|\to|Y|$ 的像包含 $|X_0|\to|Y|$ 的像。由于 $X_0$ 拟紧，
这是可能的。只需检验对某个 $i\geq0$，
$W\times_YX_i\to W$ 分离；因为 $W\times_YX_i$ 在 $W$ 上的对角态射
是 $X_i\to X_i\times_YX_i$ 沿满平展态射
$(X_i\times_YX_i)\times_YW\to X_i\times_YX_i$ 的基变换。由于
$Y$ 拟分离，代数空间 $W\times_YX_i$ 拟紧（也拟分离）。所以可作
到 $W$ 的基变换，并假设 $Y$ 是仿射概形。当 $Y$ 是仿射概形时，
我们必须证明对充分大的 $i$，$X_i$ 是分离代数空间，而已知 $X$
是分离代数空间。因此这种情形由引理 \ref{spaces-limits-lemma-descend-separated}
得出。
\end{proof}

\begin{lemma}
\label{spaces-limits-lemma-eventually-affine}
设 $S$ 为概形，$Y$ 为 $S$ 上的代数空间。设 $X=\lim X_i$ 是 $Y$
上代数空间的有向极限，且转移态射仿射。假设
\begin{enumerate}
\item $Y$ 拟紧拟分离；
\item $X_i$ 拟紧拟分离；
\item $X\to Y$ 仿射。
\end{enumerate}
那么当 $i$ 充分大时，$X_i\to Y$ 仿射。
\end{lemma}

\begin{proof}
选取仿射概形 $W$ 和满平展态射 $W\to Y$。于是 $X\times_YW$ 仿射，
并且只需检验对某个 $i$，$X_i\times_YW$ 仿射（《空间的态射》引理
\ref{spaces-morphisms-lemma-affine-local}）。这由引理
\ref{spaces-limits-lemma-limit-is-affine} 得出。
\end{proof}

\begin{lemma}
\label{spaces-limits-lemma-eventually-finite}
设 $S$ 为概形，$Y$ 为 $S$ 上的代数空间。设 $X=\lim X_i$ 是 $Y$
上代数空间的有向极限，且转移态射仿射。假设
\begin{enumerate}
\item $Y$ 拟紧拟分离；
\item $X_i$ 拟紧拟分离；
\item 转移态射 $X_{i'}\to X_i$ 有限；
\item $X_i\to Y$ 局部有限型；
\item $X\to Y$ 整。
\end{enumerate}
那么当 $i$ 充分大时，$X_i\to Y$ 有限。
\end{lemma}

\begin{proof}
选取仿射概形 $W$ 和满平展态射 $W\to Y$。于是 $X\times_YW$ 在 $W$
上有限；只需检验对某个 $i$，$X_i\times_YW$ 在 $W$ 上有限
（《空间的态射》引理 \ref{spaces-morphisms-lemma-integral-local}）。
由引理 \ref{spaces-limits-lemma-limit-is-scheme}，这约化到概形的情形。对概形，
结论由《极限》引理 \ref{limits-lemma-eventually-finite} 得出。
\end{proof}

\begin{lemma}
\label{spaces-limits-lemma-eventually-closed-immersion}
设 $S$ 为概形，$Y$ 为 $S$ 上的代数空间。设 $X=\lim X_i$ 是 $Y$
上代数空间的有向极限，且转移态射仿射。假设
\begin{enumerate}
\item $Y$ 拟紧拟分离；
\item $X_i$ 拟紧拟分离；
\item 转移态射 $X_{i'}\to X_i$ 是闭浸入；
\item $X_i\to Y$ 局部有限型；
\item $X\to Y$ 是闭浸入。
\end{enumerate}
那么当 $i$ 充分大时，$X_i\to Y$ 是闭浸入。
\end{lemma}

\begin{proof}
选取仿射概形 $W$ 和满平展态射 $W\to Y$。于是 $X\times_YW$ 是 $W$
的闭子空间；只需检验对某个 $i$，$X_i\times_YW$ 是闭子空间 $W$。
% P09-ERR-0071
（《空间的态射》引理
\ref{spaces-morphisms-lemma-closed-immersion-local}）。由引理
\ref{spaces-limits-lemma-limit-is-scheme}，这约化到概形的情形。对概形，结论由
《极限》引理 \ref{limits-lemma-eventually-closed-immersion} 得出。
\end{proof}

















\section{态射性质的下降}
\label{spaces-limits-section-descent-of-properties}

\noindent
本节是第 \ref{spaces-limits-section-descent} 节针对态射性质的对应版本。我们将在如下情形中工作。

\begin{situation}
\label{spaces-limits-situation-descent-property}
设 $S$ 为概形。设 $B = \lim B_i$ 是 $S$ 上代数空间的一个有向逆系统的极限，
其转移态射均为仿射态射
（引理 \ref{spaces-limits-lemma-directed-inverse-system-has-limit}）。
取 $0 \in I$，并设 $f_0 : X_0 \to Y_0$ 为 $B_0$ 上代数空间的一个态射。
假设 $B_0$、$X_0$、$Y_0$ 均拟紧且拟分离。令
$f_i : X_i \to Y_i$ 为 $f_0$ 到 $B_i$ 的基变换，并令
$f : X \to Y$ 为 $f_0$ 到 $B$ 的基变换。
\end{situation}

\begin{lemma}
\label{spaces-limits-lemma-descend-etale}
沿用情形 \ref{spaces-limits-situation-descent-property} 的记号和假设。若
\begin{enumerate}
\item $f$ 平展，
\item $f_0$ 局部有限表示，
\end{enumerate}
则对某个 $i \geq 0$，$f_i$ 平展。
\end{lemma}

\begin{proof}
选取仿射概形 $V_0$ 以及满平展态射
$V_0 \to Y_0$。再选取仿射概形 $U_0$ 以及满平展态射
$U_0 \to V_0 \times_{Y_0} X_0$。图示为
$$
\xymatrix{
U_0 \ar[d] \ar[r] & V_0 \ar[d] \\
X_0 \ar[r] & Y_0
}
$$
按构造，竖直箭头均满且平展。将此图基变换到 $B_i$ 或 $B$，得到
$$
\vcenter{
\xymatrix{
U_i \ar[d] \ar[r] & V_i \ar[d] \\
X_i \ar[r] & Y_i
}
}
\quad\text{且}\quad
\vcenter{
\xymatrix{
U \ar[d] \ar[r] & V \ar[d] \\
X \ar[r] & Y
}
}
$$
注意 $U_i,V_i,U,V$ 均为仿射概形，竖直态射均为满平展态射，
而态射 $U_i \to V_i$ 的极限是 $U \to V$。回忆
$X_i \to Y_i$ 平展当且仅当 $U_i \to V_i$ 平展；类似地，
$X \to Y$ 平展当且仅当 $U \to V$ 平展
（《空间的态射》，引理 \ref{spaces-morphisms-lemma-etale-local}）。
由于 $f_0$ 局部有限表示，态射 $U_0 \to V_0$ 也是如此。
因此本引理由《极限》，引理 \ref{limits-lemma-descend-etale} 得出。
\end{proof}

\begin{lemma}
\label{spaces-limits-lemma-descend-smooth}
沿用情形 \ref{spaces-limits-situation-descent-property} 的记号和假设。若
\begin{enumerate}
\item $f$ 光滑，
\item $f_0$ 局部有限表示，
\end{enumerate}
则对某个 $i \geq 0$，$f_i$ 光滑。
\end{lemma}

\begin{proof}
选取仿射概形 $V_0$ 以及满平展态射
$V_0 \to Y_0$。再选取仿射概形 $U_0$ 以及满平展态射
$U_0 \to V_0 \times_{Y_0} X_0$。图示为
$$
\xymatrix{
U_0 \ar[d] \ar[r] & V_0 \ar[d] \\
X_0 \ar[r] & Y_0
}
$$
按构造，竖直箭头均满且平展。将此图基变换到 $B_i$ 或 $B$，得到
$$
\vcenter{
\xymatrix{
U_i \ar[d] \ar[r] & V_i \ar[d] \\
X_i \ar[r] & Y_i
}
}
\quad\text{且}\quad
\vcenter{
\xymatrix{
U \ar[d] \ar[r] & V \ar[d] \\
X \ar[r] & Y
}
}
$$
注意 $U_i,V_i,U,V$ 均为仿射概形，竖直态射均为满平展态射，
而态射 $U_i \to V_i$ 的极限是 $U \to V$。回忆
$X_i \to Y_i$ 光滑当且仅当 $U_i \to V_i$ 光滑；类似地，
$X \to Y$ 光滑当且仅当 $U \to V$ 光滑
（《空间的态射》，定义 \ref{spaces-morphisms-definition-smooth}）。
由于 $f_0$ 局部有限表示，态射 $U_0 \to V_0$ 也是如此。
因此本引理由《极限》，引理 \ref{limits-lemma-descend-smooth} 得出。
\end{proof}

\begin{lemma}
\label{spaces-limits-lemma-descend-surjective}
沿用情形 \ref{spaces-limits-situation-descent-property} 的记号和假设。若
\begin{enumerate}
\item $f$ 满，
\item $f_0$ 局部有限表示，
\end{enumerate}
则对某个 $i \geq 0$，$f_i$ 满。
\end{lemma}

\begin{proof}
选取仿射概形 $V_0$ 以及满平展态射
$V_0 \to Y_0$。再选取仿射概形 $U_0$ 以及满平展态射
$U_0 \to V_0 \times_{Y_0} X_0$。图示为
$$
\xymatrix{
U_0 \ar[d] \ar[r] & V_0 \ar[d] \\
X_0 \ar[r] & Y_0
}
$$
按构造，竖直箭头均满且平展。将此图基变换到 $B_i$ 或 $B$，得到
$$
\vcenter{
\xymatrix{
U_i \ar[d] \ar[r] & V_i \ar[d] \\
X_i \ar[r] & Y_i
}
}
\quad\text{且}\quad
\vcenter{
\xymatrix{
U \ar[d] \ar[r] & V \ar[d] \\
X \ar[r] & Y
}
}
$$
注意 $U_i,V_i,U,V$ 均为仿射概形，竖直态射均为满平展态射，
态射 $U_i \to V_i$ 的极限是 $U \to V$，而且态射
$U_i \to X_i \times_{Y_i} V_i$ 与
$U \to X \times_Y V$ 均为满态射（它们是
$U_0 \to X_0 \times_{Y_0} V_0$ 的基变换）。特别地，
$X_i \to Y_i$ 满当且仅当 $U_i \to V_i$ 满；类似地，
$X \to Y$ 满当且仅当 $U \to V$ 满。由于 $f_0$ 局部有限表示，
态射 $U_0 \to V_0$ 也是如此。因此本引理由概形情形
（《极限》，引理 \ref{limits-lemma-descend-surjective}）得出。
\end{proof}

\begin{lemma}
\label{spaces-limits-lemma-descend-universally-injective}
沿用情形 \ref{spaces-limits-situation-descent-property} 的记号和假设。若
\begin{enumerate}
\item $f$ 普遍单射，
\item $f_0$ 局部有限型，
\end{enumerate}
则对某个 $i \geq 0$，$f_i$ 普遍单射。
\end{lemma}

\begin{proof}
回忆，态射 $X \to Y$ 普遍单射，当且仅当对角态射
$X \to X \times_Y X$ 为满态射
（《空间的态射》，定义
\ref{spaces-morphisms-definition-universally-injective} 及
引理 \ref{spaces-morphisms-lemma-universally-injective}）。
% P09-ERR-0072
注意 $X_0 \to X_0 \times_{Y_0} X_0$ 局部有限表示
（《空间的态射》，引理
\ref{spaces-morphisms-lemma-diagonal-morphism-finite-type}）。
因此，将引理 \ref{spaces-limits-lemma-descend-surjective} 应用于态射
$X_0 \to X_0 \times_{Y_0} X_0$ 即得本引理。
\end{proof}

\begin{lemma}
\label{spaces-limits-lemma-descend-affine}
沿用情形 \ref{spaces-limits-situation-descent-property} 的记号和假设。若
$f$ 仿射，则对某个 $i \geq 0$，$f_i$ 仿射。
\end{lemma}

\begin{proof}
选取仿射概形 $V_0$ 以及满平展态射 $V_0 \to Y_0$。
令 $V_i = V_0 \times_{Y_0} Y_i$ 且 $V = V_0 \times_{Y_0} Y$。
由于 $f$ 仿射，$V \times_Y X = \lim V_i \times_{Y_i} X_i$ 仿射。
由引理 \ref{spaces-limits-lemma-limit-is-affine}，对某个 $i \geq 0$，
$V_i \times_{Y_i} X_i$ 仿射。对此 $i$，态射 $f_i$ 仿射
（《空间的态射》，引理 \ref{spaces-morphisms-lemma-affine-local}）。
\end{proof}

\begin{lemma}
\label{spaces-limits-lemma-descend-finite}
沿用情形 \ref{spaces-limits-situation-descent-property} 的记号和假设。若
\begin{enumerate}
\item $f$ 有限，
\item $f_0$ 局部有限型，
\end{enumerate}
则对某个 $i \geq 0$，$f_i$ 有限。
\end{lemma}

\begin{proof}
选取仿射概形 $V_0$ 以及满平展态射 $V_0 \to Y_0$。
令 $V_i = V_0 \times_{Y_0} Y_i$ 且 $V = V_0 \times_{Y_0} Y$。
由于 $f$ 有限，$V \times_Y X = \lim V_i \times_{Y_i} X_i$
是 $V$ 上的有限概形。由引理 \ref{spaces-limits-lemma-limit-is-affine}，
对某个 $i \geq 0$，$V_i \times_{Y_i} X_i$ 仿射。必要时增大 $i$，
由《极限》，引理 \ref{limits-lemma-descend-finite-finite-presentation}，
可使 $V_i \times_{Y_i} X_i \to V_i$ 有限。对此 $i$，态射 $f_i$ 有限
（《空间的态射》，引理 \ref{spaces-morphisms-lemma-integral-local}）。
\end{proof}

\begin{lemma}
\label{spaces-limits-lemma-descend-closed-immersion}
沿用情形 \ref{spaces-limits-situation-descent-property} 的记号和假设。若
\begin{enumerate}
\item $f$ 为闭浸入，
\item $f_0$ 局部有限型，
\end{enumerate}
则对某个 $i \geq 0$，$f_i$ 为闭浸入。
\end{lemma}

\begin{proof}
选取仿射概形 $V_0$ 以及满平展态射 $V_0 \to Y_0$。
令 $V_i = V_0 \times_{Y_0} Y_i$ 且 $V = V_0 \times_{Y_0} Y$。
由于 $f$ 为闭浸入，
$V \times_Y X = \lim V_i \times_{Y_i} X_i$
是仿射概形 $V$ 的闭子概形。由引理 \ref{spaces-limits-lemma-limit-is-affine}，
对某个 $i \geq 0$，$V_i \times_{Y_i} X_i$ 仿射。必要时增大 $i$，
由《极限》，引理
\ref{limits-lemma-descend-closed-immersion-finite-presentation}，
可使 $V_i \times_{Y_i} X_i \to V_i$ 为闭浸入。
% P09-ERR-0073
对此 $i$，态射 $f_i$ 为闭浸入
（《空间的态射》，引理 \ref{spaces-morphisms-lemma-integral-local}）。
\end{proof}

\begin{lemma}
\label{spaces-limits-lemma-descend-separated-morphism}
沿用情形 \ref{spaces-limits-situation-descent-property} 的记号和假设。
若 $f$ 分离，则对某个 $i \geq 0$，$f_i$ 分离。
\end{lemma}

\begin{proof}
将引理 \ref{spaces-limits-lemma-descend-closed-immersion} 应用于对角态射
$\Delta_{X_0/Y_0} : X_0 \to X_0 \times_{Y_0} X_0$。
（对角态射局部有限型，且纤维积
$X_0 \times_{Y_0} X_0$ 拟紧且拟分离。略去若干细节。）
\end{proof}

\begin{lemma}
\label{spaces-limits-lemma-descend-isomorphism}
沿用情形 \ref{spaces-limits-situation-descent-property} 的记号和假设。若
\begin{enumerate}
\item $f$ 为同构，
\item $f_0$ 局部有限表示，
\end{enumerate}
则对某个 $i \geq 0$，$f_i$ 为同构。
\end{lemma}

\begin{proof}
一个态射为同构，当且仅当它平展、普遍单射且满，参见
《空间的态射》，引理
\ref{spaces-morphisms-lemma-etale-universally-injective-open}。
因此，本引理由引理 \ref{spaces-limits-lemma-descend-etale}、
\ref{spaces-limits-lemma-descend-surjective} 及
\ref{spaces-limits-lemma-descend-universally-injective} 得出。
\end{proof}

\begin{lemma}
\label{spaces-limits-lemma-descend-monomorphism}
沿用情形 \ref{spaces-limits-situation-descent-property} 的记号和假设。若
\begin{enumerate}
\item $f$ 为单态射，
\item $f_0$ 局部有限型，
\end{enumerate}
则对某个 $i \geq 0$，$f_i$ 为单态射。
\end{lemma}

\begin{proof}
回忆，一个态射为单态射，当且仅当其对角态射为同构。由
《空间的态射》，引理
\ref{spaces-morphisms-lemma-diagonal-morphism-finite-type}，
态射 $X_0 \to X_0 \times_{Y_0} X_0$ 局部有限表示。
由于 $X_0 \times_{Y_0} X_0$ 拟紧且拟分离，引理
\ref{spaces-limits-lemma-descend-isomorphism} 表明，对某个 $i \geq 0$，
$\Delta_i : X_i \to X_i \times_{Y_i} X_i$ 为同构。对此 $i$，
态射 $f_i$ 为单态射。
\end{proof}

\begin{lemma}
\label{spaces-limits-lemma-descend-flat}
沿用情形 \ref{spaces-limits-situation-descent-property} 的记号和假设。
设 $\mathcal{F}_0$ 为拟凝聚 $\mathcal{O}_{X_0}$-模，
并以 $\mathcal{F}_i$ 表示其到 $X_i$ 的拉回，以 $\mathcal{F}$
表示其到 $X$ 的拉回。若
\begin{enumerate}
\item $\mathcal{F}$ 在 $Y$ 上平坦，
\item $\mathcal{F}_0$ 有限表示，并且
\item $f_0$ 局部有限表示，
\end{enumerate}
则对某个 $i \geq 0$，$\mathcal{F}_i$ 在 $Y_i$ 上平坦。
特别地，若 $f_0$ 局部有限表示且 $f$ 平坦，则对某个
$i \geq 0$，$f_i$ 平坦。
\end{lemma}

\begin{proof}
选取仿射概形 $V_0$ 以及满平展态射
$V_0 \to Y_0$。再选取仿射概形 $U_0$ 以及满平展态射
$U_0 \to V_0 \times_{Y_0} X_0$。图示为
$$
\xymatrix{
U_0 \ar[d] \ar[r] & V_0 \ar[d] \\
X_0 \ar[r] & Y_0
}
$$
按构造，竖直箭头均满且平展。将此图基变换到 $B_i$ 或 $B$，得到
$$
\vcenter{
\xymatrix{
U_i \ar[d] \ar[r] & V_i \ar[d] \\
X_i \ar[r] & Y_i
}
}
\quad\text{且}\quad
\vcenter{
\xymatrix{
U \ar[d] \ar[r] & V \ar[d] \\
X \ar[r] & Y
}
}
$$
注意 $U_i,V_i,U,V$ 均为仿射概形，竖直态射均为满平展态射，
而态射 $U_i \to V_i$ 的极限是 $U \to V$。回忆，
$\mathcal{F}_i$ 在 $Y_i$ 上平坦，当且仅当
$\mathcal{F}_i|_{U_i}$ 在 $V_i$ 上平坦；类似地，
$\mathcal{F}$ 在 $Y$ 上平坦，当且仅当
$\mathcal{F}|_U$ 在 $V$ 上平坦
（《空间的态射》，定义 \ref{spaces-morphisms-definition-flat}）。
由于 $f_0$ 局部有限表示，态射 $U_0 \to V_0$ 也是如此。
因此本引理由《极限》，引理
\ref{limits-lemma-descend-module-flat-finite-presentation} 得出。
\end{proof}

\begin{lemma}
\label{spaces-limits-lemma-eventually-proper}
沿用情形 \ref{spaces-limits-situation-descent-property} 的假设和记号。若
\begin{enumerate}
\item $f$ 固有，并且
\item $f_0$ 局部有限型，
\end{enumerate}
则存在 $i$，使得 $f_i$ 固有。
\end{lemma}

\begin{proof}
选取仿射概形 $V_0$ 以及满平展态射 $V_0 \to Y_0$。
令 $V_i = Y_i \times_{Y_0} V_0$ 且 $V = Y \times_{Y_0} V_0$。
只需证明 $f_i$ 到 $V_i$ 的基变换固有，参见《空间的态射》，
引理 \ref{spaces-morphisms-lemma-proper-local}。
因此可以假设 $Y_0$ 仿射。

\medskip\noindent
由引理 \ref{spaces-limits-lemma-descend-separated-morphism}，对某个 $i \geq 0$，
$f_i$ 分离。以 $i$ 代替 $0$，可以假设 $f_0$ 分离。
注意 $f_0$ 拟紧。因此 $f_0$ 分离且有限型。由
《空间的上同调》，引理 \ref{spaces-cohomology-lemma-weak-chow}，
可以选取图
$$
\xymatrix{
X_0 \ar[rd] & X_0' \ar[d] \ar[l]^\pi \ar[r] & \mathbf{P}^n_{Y_0} \ar[dl] \\
& Y_0 &
}
$$
其中 $X_0' \to \mathbf{P}^n_{Y_0}$ 为浸入，而
$\pi : X_0' \to X_0$ 固有且满。引入
$X' = X_0' \times_{Y_0} Y$ 及 $X_i' = X_0' \times_{Y_0} Y_i$。
由《空间的态射》，引理
\ref{spaces-morphisms-lemma-composition-proper} 及
\ref{spaces-morphisms-lemma-base-change-proper}，
$X' \to Y$ 固有。因此 $X' \to \mathbf{P}^n_Y$ 为闭浸入
（《空间的态射》，引理
\ref{spaces-morphisms-lemma-universally-closed-permanence}）。
由《空间的态射》，引理
\ref{spaces-morphisms-lemma-image-proper-is-proper}，
只需证明对某个 $i$，$X'_i \to Y_i$ 固有。由引理
\ref{spaces-limits-lemma-descend-closed-immersion}，当 $i$ 充分大时，
$X'_i \to \mathbf{P}^n_{Y_i}$ 为闭浸入。于是
$X'_i \to Y_i$ 固有，结论得证。
\end{proof}

\begin{lemma}
\label{spaces-limits-lemma-eventually-relative-dimension}
沿用情形 \ref{spaces-limits-situation-descent-property} 的假设和记号。
设 $d \geq 0$。若
\begin{enumerate}
\item $f$ 的相对维数 $\leq d$
（《空间的态射》，定义
\ref{spaces-morphisms-definition-relative-dimension}），
\item $f_0$ 局部有限型，
\end{enumerate}
则存在 $i$，使得 $f_i$ 的相对维数 $\leq d$。
\end{lemma}

\begin{proof}
选取仿射概形 $V_0$ 以及满平展态射
$V_0 \to Y_0$。再选取仿射概形 $U_0$ 以及满平展态射
$U_0 \to V_0 \times_{Y_0} X_0$。图示为
$$
\xymatrix{
U_0 \ar[d] \ar[r] & V_0 \ar[d] \\
X_0 \ar[r] & Y_0
}
$$
按构造，竖直箭头均满且平展。将此图基变换到 $B_i$ 或 $B$，得到
$$
\vcenter{
\xymatrix{
U_i \ar[d] \ar[r] & V_i \ar[d] \\
X_i \ar[r] & Y_i
}
}
\quad\text{且}\quad
\vcenter{
\xymatrix{
U \ar[d] \ar[r] & V \ar[d] \\
X \ar[r] & Y
}
}
$$
注意 $U_i,V_i,U,V$ 均为仿射概形，竖直态射均为满平展态射，
而态射 $U_i \to V_i$ 的极限是 $U \to V$。
在此情形中，$X_i \to Y_i$ 的相对维数 $\leq d$，当且仅当
$U_i \to V_i$ 的相对维数 $\leq d$
（定义见《态射》，定义
\ref{morphisms-definition-relative-dimension-d}）。
为说明这一等价性，注意代数空间态射的定义用到
《空间的态射》，定义
\ref{spaces-morphisms-definition-dimension-fibre}，
而后者使用平展局部化。$X \to Y$ 与 $U \to V$ 也是如此。
由于 $f_0$ 局部有限型，态射 $U_0 \to V_0$ 也是如此。
因此本引理由更一般的《极限》，引理
\ref{limits-lemma-limit-dimension} 得出。
\end{proof}














\section{相对对象的下降}
\label{spaces-limits-section-descending-relative}

\noindent
下述引理是本节这类结果的典型例子。

\begin{lemma}
\label{spaces-limits-lemma-descend-finite-presentation}
设 $S$ 为概形，$I$ 为有向集。设 $(X_i, f_{ii'})$ 是由 $S$ 上代数空间
组成的、以 $I$ 为指标的逆系统。假设
\begin{enumerate}
\item 态射 $f_{ii'} : X_i \to X_{i'}$ 均为仿射态射，
\item 空间 $X_i$ 均拟紧且拟分离。
\end{enumerate}
令 $X = \lim_i X_i$。则 $X$ 上有限表示代数空间的范畴，是
$X_i$ 上有限表示代数空间范畴在 $I$ 上的余极限。
\end{lemma}

\begin{proof}
取 $0 \in I$。选取满平展态射 $U_0 \to X_0$，其中 $U_0$ 为仿射概形
（《空间的性质》，引理
\ref{spaces-properties-lemma-quasi-compact-affine-cover}）。
令 $U_i = X_i \times_{X_0} U_0$。令
$R_0 = U_0 \times_{X_0} U_0$ 且
$R_i = R_0 \times_{X_0} X_i$。以
$s_i,t_i : R_i \to U_i$ 及 $s,t : R \to U$ 表示两个投影。
在引理 \ref{spaces-limits-lemma-directed-inverse-system-has-limit} 的证明中，我们已看到
存在展示 $X = U/R$，其中 $U = \lim U_i$ 且 $R = \lim R_i$。
注意 $U_i$ 与 $U$ 均仿射，而 $R_i$ 与 $R$ 均拟紧且分离
（因为 $X_i$ 拟分离）。设 $Y$ 为 $S$ 上的代数空间，并设
$Y \to X$ 为有限表示态射。令 $V = U \times_X Y$。
这是 $U$ 上的有限表示代数空间。选取仿射概形 $W$ 及满平展态射
$W \to V$。于是 $W \to Y$ 也满且平展。令
$R' = W \times_Y W$，从而 $Y = W/R'$
（参见《空间》，第 \ref{spaces-section-presentations} 节）。
注意 $W$ 是 $U$ 上有限表示的概形，而 $R'$ 是 $R$ 上有限表示的概形
（略去细节）。由《极限》，引理
\ref{limits-lemma-descend-finite-presentation}，可以找到指标 $i$ 及
有限表示的概形态射 $W_i \to U_i$，其到 $U$ 的基变换给出
$W \to U$。类似地，必要时增大 $i$ 后，可以找到 $R_i$ 上有限表示的概形
$R'_i$，其到 $R$ 的基变换为 $R'$。
投影态射 $s',t' : R' \to W$ 分别位于投影态射
$s,t : R \to U$ 之上。因此，可以把 $s'$（相应地 $t'$）
视为 $U$ 上有限表示概形之间的态射（其中 $R' \to U$ 的结构态射
由 $R' \to R$ 后接 $s$（相应地 $t$）给出）。于是再次应用
《极限》，引理 \ref{limits-lemma-descend-finite-presentation} 可见，
必要时增大 $i$ 后，存在态射 $s'_i,t'_i : R'_i \to W_i$，
% P09-ERR-0074
其到 $U$ 的基变换为 $S',t'$。由《极限》，引理
\ref{limits-lemma-descend-etale} 及
\ref{limits-lemma-descend-monomorphism}，可以假设
$s'_i,t'_i$ 均平展，而且
$j'_i : R'_i \to W_i \times_{X_i} W_i$ 为单态射（这里通过任一投影
把 $j'_i$ 看作 $U_i$ 上有限表示概形之间的态射——选哪个投影无关紧要）。
令 $Y_i = W_i/R'_i$（参见《空间》，定理
\ref{spaces-theorem-presentation}），便得到 $X_i$ 上的有限表示代数空间，
其到 $X$ 的基变换同构于 $Y$。

\medskip\noindent
这表明，$X$ 上的每个有限表示代数空间都来自某个 $X_i$ 上的有限表示
代数空间；换言之，本引理中的函子本质满。为证明它全忠实，取指标
$0 \in I$ 以及 $X_0$ 上两个有限表示代数空间 $Y_0,Z_0$。
令 $Y_i = X_i \times_{X_0} Y_0$、$Y = X \times_{X_0} Y_0$、
$Z_i = X_i \times_{X_0} Z_0$ 且 $Z = X \times_{X_0} Z_0$。
设 $\alpha : Y \to Z$ 为 $X$ 上代数空间的态射。选取满平展态射
$V_0 \to Y_0$，其中 $V_0$ 为仿射概形。令
$V_i = V_0 \times_{Y_0} Y_i$ 且 $V = V_0 \times_{Y_0} Y$；
它们是分别带有到 $Y_i$ 与 $Y$ 的满平展态射的仿射概形。
% P09-ERR-0075
由命题 \ref{spaces-limits-proposition-characterize-locally-finite-presentation}
（应用于依假设有限表示的 $Z_0 \to X_0$），复合
$V \to Y \to Z \to Z_0$ 来自某个 $i \geq 0$ 上一个（本质唯一的）态射
$V_i \to Z_0$。增大 $i$ 后，两个复合
$$
V_i \times_{Y_i} V_i \to V_i \to Z_0
$$
相等，因为它们在极限中相等。因而得到一个（本质唯一的）态射
$Y_i \to Z_0$。由于这是 $X_0$ 上的态射，它诱导出一个到
$Z_i = Z_0 \times_{X_0} X_i$ 的态射，正合所需。
\end{proof}

\begin{lemma}
\label{spaces-limits-lemma-descend-modules-finite-presentation}
沿用引理 \ref{spaces-limits-lemma-descend-finite-presentation} 的记号和假设。
有限表示 $\mathcal{O}_X$-模的范畴，是有限表示
$\mathcal{O}_{X_i}$-模范畴在 $I$ 上的余极限。
\end{lemma}

\begin{proof}
取 $0 \in I$。选取仿射概形 $U_0$ 及满平展态射
$U_0 \to X_0$。令 $U_i = X_i \times_{X_0} U_0$。
令 $R_0 = U_0 \times_{X_0} U_0$ 且
$R_i = R_0 \times_{X_0} X_i$。以
$s_i,t_i : R_i \to U_i$ 及 $s,t : R \to U$ 表示两个投影。
在引理 \ref{spaces-limits-lemma-directed-inverse-system-has-limit} 的证明中，我们已看到
存在展示 $X = U/R$，其中 $U = \lim U_i$ 且 $R = \lim R_i$。
注意 $U_i$ 与 $U$ 均仿射，而 $R_i$ 与 $R$ 均拟紧且分离
（因为 $X_i$ 拟分离）。此外，还有
% P09-ERR-0076
$R \times_{s, U, t} R = \colim R_i \times_{s_i, U_i, t_i} R_i$.
因此我们知道
$\QCoh(\mathcal{O}_U) = \colim \QCoh(\mathcal{O}_{U_i})$,
$\QCoh(\mathcal{O}_R) = \colim \QCoh(\mathcal{O}_{R_i})$,
以及
$\QCoh(\mathcal{O}_{R \times_{s, U, t} R}) =
\colim \QCoh(\mathcal{O}_{R_i \times_{s_i, U_i, t_i} R_i})$，由
《极限》，引理 \ref{limits-lemma-descend-modules-finite-presentation}。
我们有 $\QCoh(\mathcal{O}_X) = \QCoh(U, R, s, t, c)$ 及
$\QCoh(\mathcal{O}_{X_i}) = \QCoh(U_i, R_i, s_i, t_i, c_i)$,
参见《空间的性质》，命题
\ref{spaces-properties-proposition-quasi-coherent}。
因此形式地得到结论。
\end{proof}

\begin{lemma}
\label{spaces-limits-lemma-descend-invertible-modules}
沿用引理 \ref{spaces-limits-lemma-descend-finite-presentation} 的记号和假设。则
\begin{enumerate}
\item 任一有限局部自由 $\mathcal{O}_X$-模都是某个 $i$ 上一个有限局部自由
$\mathcal{O}_{X_i}$-模的拉回，
\item 任一可逆 $\mathcal{O}_X$-模都是某个 $i$ 上一个可逆
$\mathcal{O}_{X_i}$-模的拉回。
\end{enumerate}
\end{lemma}

\begin{proof}
证明 (2)。设 $\mathcal{L}$ 为可逆 $\mathcal{O}_X$-模。
由于可逆模有限表示，由引理
\ref{spaces-limits-lemma-descend-modules-finite-presentation}，可以找到 $i$ 以及
$X_i$ 上有限表示的模 $\mathcal{L}_i$ 与 $\mathcal{N}_i$，使得
% P09-ERR-0077
$f_i^*\mathcal{L}_i \cong \mathcal{L}$ 且
$f_i^*\mathcal{N}_i \cong \mathcal{L}^{\otimes -1}$。
由于拉回与张量积交换，
$f_i^*(\mathcal{L}_i \otimes_{\mathcal{O}_{X_i}} \mathcal{N}_i)$
同构于 $\mathcal{O}_X$。由于有限表示模的张量积仍有限表示，同一引理表明，
$f_{i'i}^*\mathcal{L}_i
\otimes_{\mathcal{O}_{X_{i'}}} f_{i'i}^*\mathcal{N}_i$
对某个 $i' \geq i$ 同构于 $\mathcal{O}_{X_{i'}}$。
从而 $f_{i'i}^*\mathcal{L}_i$ 可逆
（《位点上的模》，引理 \ref{sites-modules-lemma-invertible}），
证明完成。

\medskip\noindent
证明 (1)。略。提示：仿照 (2) 的证明，并使用如下事实：一个模
（在局部赋环位点上）有限局部自由，当且仅当它具有对偶；参见
《位点上的模》，第 \ref{sites-modules-section-duals} 节。
也可以仿照概形情形的证明，参见《极限》，引理
\ref{limits-lemma-descend-invertible-modules}。
\end{proof}














\section{绝对诺特逼近}
\label{spaces-limits-section-approximation}

\noindent
下述结果见 \cite[Theorem 1.2.2]{CLO}。证明中的一个关键要素是
《良态空间》，引理
\ref{decent-spaces-lemma-filter-quasi-compact-quasi-separated}。

\begin{proposition}
\label{spaces-limits-proposition-approximate}
\begin{reference}
我们的证明紧随 \cite[Theorem 1.2.2]{CLO} 中给出的证明。
\end{reference}
设 $X$ 为 $\Spec(\mathbf{Z})$ 上拟紧且拟分离的代数空间。
则存在有向集 $I$ 以及以 $I$ 为指标的代数空间逆系统
$(X_i, f_{ii'})$，使得
\begin{enumerate}
% P09-ERR-0078
\item 转移态射 $f_{ii'}$ 均为仿射态射
\item 每个 $X_i$ 均拟分离且在 $\mathbf{Z}$ 上有限型，并且
\item $X = \lim X_i$。
\end{enumerate}
\end{proposition}

\begin{proof}
应用《良态空间》，引理
\ref{decent-spaces-lemma-filter-quasi-compact-quasi-separated}
可得开子空间 $U_p \subset X$、概形 $V_p$ 以及具有所述性质的态射
$f_p : V_p \to U_p$。注意
$f_n : V_n \to U_n$ 是代数空间的平展态射，并且它在
$T_n = (V_n)_{red}$ 的逆像上的限制为同构。因而 $f_n$ 为同构，
例如可由《空间的态射》，引理
\ref{spaces-morphisms-lemma-etale-universally-injective-open} 得到。
特别地，$U_n$ 是拟紧分离概形。因此可将
$U_n = \lim U_{n,i}$ 写成 $\mathbf{Z}$ 上有限型概形的有向极限，
其转移态射均为仿射态射；参见《极限》，命题
\ref{limits-proposition-approximate}。于是对 $p$ 作降归纳，
便将问题归约为下一段所述问题。

\medskip\noindent
这里有 $U \subset X$、$U = \lim U_i$、$Z \subset X$ 以及
$f : V \to X$，它们具有下列性质：
\begin{enumerate}
\item $X$ 是拟紧且拟分离的代数空间，
\item $V$ 是拟紧分离概形，
\item $U \subset X$ 是拟紧开子空间，
\item $(U_i, g_{ii'})$ 是 $\mathbf{Z}$ 上有限型拟分离代数空间的
有向逆系统，其转移态射均为仿射态射，且极限为 $U$，
\item $Z \subset X$ 是闭子空间，满足 $|X| = |U| \amalg |Z|$，
\item $f : V \to X$ 是满平展态射，且
$f^{-1}(Z) \to Z$ 为同构。
\end{enumerate}
问题：证明本命题的结论对 $X$ 成立。

\medskip\noindent
注意 $W = f^{-1}(U) \subset V$ 是 $U$ 上平展的拟紧开子概形。
因此可应用引理 \ref{spaces-limits-lemma-descend-finite-presentation} 及
\ref{spaces-limits-lemma-descend-etale}，找到指标 $0 \in I$ 以及有限表示平展态射
$W_0 \to U_0$，其到 $U$ 的基变换给出 $W$。令
$W_i = W_0 \times_{U_0} U_i$，则 $W = \lim_{i \geq 0} W_i$。
增大 $0$ 后，可以假设 $W_i$ 均为概形；参见引理
\ref{spaces-limits-lemma-limit-is-scheme}。此外，$W_i$ 在 $\mathbf{Z}$ 上有限型。

\medskip\noindent
将《极限》，引理 \ref{limits-lemma-approximate} 应用于
$W = \lim_{i \geq 0} W_i$ 及包含 $W \subset V$。以该引理中得到的
有向集 $J$ 代替 $I$。这样便可将 $V$ 写成有向极限
$V = \lim V_i$，其中 $V_i$ 是 $\mathbf{Z}$ 上有限型概形，
转移映射均为仿射映射，并且每个 $V_i$ 都以 $W_i$ 为开子概形
（与转移态射相容）。对每个 $i$，可以在概形范畴中作推出
$$
\xymatrix{
W_i \ar[r] \ar[d]_\Delta & V_i \ar[d] \\
W_i \times_{U_i} W_i \ar[r] & R_i
}
$$
。事实上，左侧竖直箭头和上方水平箭头都是概形的开浸入。换言之，
可沿公共开子概形 $W_i$ 黏合 $V_i$ 与
$W_i \times_{U_i} W_i$ 来构造 $R_i$
（参见《概形》，第 \ref{schemes-section-glueing-schemes} 节）。
注意，平展投影 $W_i \times_{U_i} W_i \to W_i$ 延拓为平展态射
$s_i,t_i : R_i \to V_i$。显然，态射
$j_i = (t_i, s_i) : R_i \to V_i \times V_i$ 是 $V_i$ 上的
平展等价关系。注意 $W_i \times_{U_i} W_i$ 拟紧
（因为 $U_i$ 拟分离且 $W_i$ 拟紧），而 $V_i$ 拟紧，故 $R_i$ 拟紧。
当 $i \geq i'$ 时，图
\begin{equation}
\label{spaces-limits-equation-cartesian}
\vcenter{
\xymatrix{
R_i \ar[r] \ar[d]_{s_i} & R_{i'} \ar[d]^{s_{i'}} \\
V_i \ar[r] & V_{i'}
}
}
\end{equation}
是笛卡尔图，因为
$$
(W_{i'} \times_{U_{i'}} W_{i'}) \times_{U_{i'}} U_i =
W_{i'} \times_{U_{i'}} U_i \times_{U_i} U_i \times_{U_{i'}} W_{i'} =
W_i \times_{U_i} W_i.
$$
考虑代数空间 $X_i = V_i/R_i$（参见《空间》，定理
\ref{spaces-theorem-presentation}）。由于 $V_i$ 在 $\mathbf{Z}$ 上有限型
且 $R_i$ 拟紧，可见 $X_i$ 拟分离且在 $\mathbf{Z}$ 上有限型
（参见《空间的性质》，引理
\ref{spaces-properties-lemma-quasi-separated}，以及《空间的态射》，引理
\ref{spaces-morphisms-lemma-surjection-from-quasi-compact} 及
\ref{spaces-morphisms-lemma-finite-type-local}）。
由于上述 $R_i$ 的构造与转移态射相容，当 $i \geq i'$ 时，
得到代数空间的态射 $X_i \to X_{i'}$。交换图
$$
\xymatrix{
V_i \ar[r] \ar[d] & V_{i'} \ar[d] \\
X_i \ar[r] & X_{i'}
}
$$
是笛卡尔图，因为图 (\ref{spaces-limits-equation-cartesian}) 是笛卡尔图；参见
《群胚》，引理 \ref{groupoids-lemma-criterion-fibre-product}。
由于 $V_i \to V_{i'}$ 仿射，这蕴含 $X_i \to X_{i'}$ 仿射；参见
《空间的态射》，引理 \ref{spaces-morphisms-lemma-affine-local}。
因此，由引理 \ref{spaces-limits-lemma-directed-inverse-system-has-limit}，
可以形成极限 $X' = \lim X_i$。我们断言 $X \cong X'$；
这将完成命题的证明。

\medskip\noindent
证明该断言。令 $R = \lim R_i$。按构造，代数空间 $X'$ 带有满平展态射
$V \to X'$，使得
$$
V \times_{X'} V \cong R
$$
（使用引理 \ref{spaces-limits-lemma-directed-inverse-system-has-limit}）。
按构造，$\lim W_i \times_{U_i} W_i = W \times_U W$ 且
$V = \lim V_i$，所以 $R$ 是 $W \times_U W$ 与 $V$ 沿 $W$ 黏合所得的并。
性质 (6) 蕴含投影 $V \times_X V \to V$ 在
$f^{-1}(Z) \subset V$ 上为同构。因此，概形 $V \times_X V$
是开子概形 $\Delta_{V/X}(V)$ 与 $W \times_U W$ 的并，二者沿
$\Delta_{W/X}(W)$ 相交。由此得到唯一同构
$R \cong V \times_X V$，且与到 $V$ 的投影相容。
由于 $V \to X$ 与 $V \to X'$ 均满且平展，可见
% P09-ERR-0079
$$
X = V/ V \times_X V = V/R = V/V \times_{X'} V = X'
$$
；由《空间》，引理 \ref{spaces-lemma-space-presentation}，结论得证。
\end{proof}




\section{应用}
\label{spaces-limits-section-applications}

\noindent
下述引理也可以直接由《良态空间》，引理
\ref{decent-spaces-lemma-filter-quasi-compact-quasi-separated}
得出，而不必经过绝对诺特逼近。

\begin{lemma}
\label{spaces-limits-lemma-colimit-finitely-presented}
设 $S$ 为概形，$X$ 为 $S$ 上拟紧且拟分离的代数空间。
每个拟凝聚 $\mathcal{O}_X$-模都是有限表示
$\mathcal{O}_X$-模的滤过余极限。
\end{lemma}

\begin{proof}
可将 $X$ 视为 $\Spec(\mathbf{Z})$ 上的代数空间；参见《空间》，定义
\ref{spaces-definition-base-change} 及《空间的性质》，定义
\ref{spaces-properties-definition-separated}。因此可应用命题
\ref{spaces-limits-proposition-approximate}，写成 $X = \lim X_i$，其中 $X_i$
在 $\mathbf{Z}$ 上有限表示。于是 $X_i$ 是诺特代数空间；参见
《空间的态射》，引理
\ref{spaces-morphisms-lemma-finite-presentation-noetherian}。
态射 $X \to X_i$ 仿射；参见引理
\ref{spaces-limits-lemma-directed-inverse-system-has-limit}。结论由
《空间的上同调》，引理
\ref{spaces-cohomology-lemma-direct-colimit-finite-presentation} 得出。
\end{proof}

\noindent
本节余下部分由引理 \ref{spaces-limits-lemma-colimit-finitely-presented} 的直接应用组成。

\begin{lemma}
\label{spaces-limits-lemma-directed-colimit-finite-type}
设 $S$ 为概形，$X$ 为 $S$ 上拟紧且拟分离的代数空间。
设 $\mathcal{F}$ 为拟凝聚 $\mathcal{O}_X$-模。
则 $\mathcal{F}$ 是其有限型拟凝聚子模的有向余极限。
\end{lemma}

\begin{proof}
若 $\mathcal{G},\mathcal{H} \subset \mathcal{F}$ 是有限型拟凝聚
$\mathcal{O}_X$-子模，则 $\mathcal{G} \oplus \mathcal{H} \to \mathcal{F}$
的像仍是有限型拟凝聚 $\mathcal{O}_X$-子模，并包含二者。
由此可见该系统有向。为证明 $\mathcal{F}$ 是此系统的余极限，
按引理 \ref{spaces-limits-lemma-colimit-finitely-presented} 写成
$\mathcal{F} = \colim_i \mathcal{F}_i$，即有限表示拟凝聚层的有向余极限。
则像 $\mathcal{G}_i = \Im(\mathcal{F}_i \to \mathcal{F})$
是 $\mathcal{F}$ 的有限型拟凝聚子层。由于 $\mathcal{F}$
是这些子层的余极限，结论得证。
\end{proof}

\begin{lemma}
\label{spaces-limits-lemma-finite-directed-colimit-surjective-maps}
设 $S$ 为概形，$X$ 为 $S$ 上拟紧且拟分离的代数空间。
设 $\mathcal{F}$ 为有限型拟凝聚 $\mathcal{O}_X$-模。则可以写成
% P09-ERR-0080
$\mathcal{F} = \lim \mathcal{F}_i$，其中每个 $\mathcal{F}_i$
都是有限表示 $\mathcal{O}_X$-模，且所有转移映射
$\mathcal{F}_i \to \mathcal{F}_{i'}$ 均满。
\end{lemma}

\begin{proof}
将 $\mathcal{F} = \colim \mathcal{G}_i$ 写成有限表示
$\mathcal{O}_X$-模的滤过余极限
（引理 \ref{spaces-limits-lemma-colimit-finitely-presented}）。
我们断言，对某个 $i$，$\mathcal{G}_i \to \mathcal{F}$ 为满态射。
事实上，选取平展满射 $U \to X$，其中 $U$ 为仿射概形。
选取有限多个生成 $\mathcal{F}|_U$ 的截面
$s_k \in \mathcal{F}(U)$。由于 $U$ 仿射，当 $i$ 充分大时，
$s_k$ 属于 $\mathcal{G}_i \to \mathcal{F}$ 的像。因此，当 $i$
充分大时，$\mathcal{G}_i \to \mathcal{F}$ 为满态射。
选取这样的 $i$，并令 $\mathcal{K} \subset \mathcal{G}_i$
为映射 $\mathcal{G}_i \to \mathcal{F}$ 的核。写成
$\mathcal{K} = \colim \mathcal{K}_a$，即其有限型拟凝聚子模的滤过余极限
（引理 \ref{spaces-limits-lemma-directed-colimit-finite-type}）。于是
$\mathcal{F} = \colim \mathcal{G}_i/\mathcal{K}_a$ 即为本引理所求。
\end{proof}

\noindent
设 $X$ 为代数空间。在下述引理中，我们使用{\it 有限表示拟凝聚
$\mathcal{O}_X$-代数 $\mathcal{A}$}这一概念。它是指：对每个
$X$ 上平展的仿射概形 $U = \Spec(R)$，都有
$\mathcal{A}|_U = \widetilde{A}$，其中 $A$ 是作为 $R$-代数有限表示的
（交换）$R$-代数。

\begin{lemma}
\label{spaces-limits-lemma-algebra-directed-colimit-finite-presentation}
设 $S$ 为概形，$X$ 为 $S$ 上拟紧且拟分离的代数空间。
设 $\mathcal{A}$ 为拟凝聚 $\mathcal{O}_X$-代数。
则 $\mathcal{A}$ 是有限表示拟凝聚 $\mathcal{O}_X$-代数的有向余极限。
\end{lemma}

\begin{proof}
首先按引理 \ref{spaces-limits-lemma-colimit-finitely-presented}，将
$\mathcal{A} = \colim_i \mathcal{F}_i$ 写成有限表示拟凝聚层的有向余极限。
对每个 $i$，令 $\mathcal{B}_i = \text{Sym}(\mathcal{F}_i)$
为 $\mathcal{F}_i$ 在 $\mathcal{O}_X$ 上的对称代数。写成
$\mathcal{I}_i = \Ker(\mathcal{B}_i \to \mathcal{A})$。
再写成 $\mathcal{I}_i = \colim_j \mathcal{F}_{i,j}$，其中
$\mathcal{F}_{i,j}$ 是 $\mathcal{I}_i$ 的有限型拟凝聚子模；参见引理
\ref{spaces-limits-lemma-directed-colimit-finite-type}。令
$\mathcal{I}_{i,j} \subset \mathcal{I}_i$ 为由
$\mathcal{F}_{i,j}$ 生成的 $\mathcal{B}_i$-理想。令
$\mathcal{A}_{i,j} = \mathcal{B}_i/\mathcal{I}_{i,j}$。
则 $\mathcal{A}_{i,j}$ 是有限表示拟凝聚 $\mathcal{O}_X$-代数。
规定 $(i,j) \leq (i',j')$，当且仅当 $i \leq i'$ 且映射
$\mathcal{B}_i \to \mathcal{B}_{i'}$ 将理想
$\mathcal{I}_{i,j}$ 映入理想 $\mathcal{I}_{i',j'}$。
于是显然 $\mathcal{A} = \colim_{i,j} \mathcal{A}_{i,j}$。
\end{proof}

\noindent
设 $X$ 为代数空间。在下述引理中，我们使用{\it 有限型拟凝聚
$\mathcal{O}_X$-代数 $\mathcal{A}$}这一概念。它是指：对每个
$X$ 上平展的仿射概形 $U = \Spec(R)$，都有
$\mathcal{A}|_U = \widetilde{A}$，其中 $A$ 是作为 $R$-代数有限型的
（交换）$R$-代数。

\begin{lemma}
\label{spaces-limits-lemma-algebra-directed-colimit-finite-type}
设 $S$ 为概形，$X$ 为 $S$ 上拟紧且拟分离的代数空间。
设 $\mathcal{A}$ 为拟凝聚 $\mathcal{O}_X$-代数。
则 $\mathcal{A}$ 是其有限型拟凝聚 $\mathcal{O}_X$-子代数的有向余极限。
\end{lemma}

\begin{proof}
略。提示：与引理 \ref{spaces-limits-lemma-directed-colimit-finite-type} 的证明比较。
\end{proof}

\noindent
设 $X$ 为代数空间。在下述引理中，我们使用{\it 有限（相应地，整）
拟凝聚 $\mathcal{O}_X$-代数 $\mathcal{A}$}这一概念。它是指：
对每个 $X$ 上平展的仿射概形 $U = \Spec(R)$，都有
$\mathcal{A}|_U = \widetilde{A}$，其中 $A$ 是作为 $R$-代数有限
（相应地，整）的（交换）$R$-代数。

\begin{lemma}
\label{spaces-limits-lemma-finite-algebra-directed-colimit-finite-finitely-presented}
设 $S$ 为概形，$X$ 为 $S$ 上拟紧且拟分离的代数空间。
设 $\mathcal{A}$ 为有限拟凝聚 $\mathcal{O}_X$-代数。则
$\mathcal{A} = \colim \mathcal{A}_i$ 是有限且有限表示的拟凝聚
$\mathcal{O}_X$-代数的有向余极限，其转移映射均满。
\end{lemma}

\begin{proof}
由引理 \ref{spaces-limits-lemma-finite-directed-colimit-surjective-maps}，
存在有限表示 $\mathcal{O}_X$-模 $\mathcal{F}$ 及满射
$\mathcal{F} \to \mathcal{A}$。利用代数结构，得到满射
$$
\text{Sym}^*_{\mathcal{O}_X}(\mathcal{F}) \longrightarrow \mathcal{A}
$$
。以 $\mathcal{J}$ 表示其核。将
$\mathcal{J} = \colim \mathcal{E}_i$ 写成有限型
$\mathcal{O}_X$-子模 $\mathcal{E}_i$ 的滤过余极限
（引理 \ref{spaces-limits-lemma-directed-colimit-finite-type}）。令
$$
\mathcal{A}_i = \text{Sym}^*_{\mathcal{O}_X}(\mathcal{F})/(\mathcal{E}_i)
$$
，其中 $(\mathcal{E}_i)$ 表示由
$\mathcal{E}_i \to \text{Sym}^*_{\mathcal{O}_X}(\mathcal{F})$ 的像生成的
理想层。于是每个 $\mathcal{A}_i$ 都是有限表示
$\mathcal{O}_X$-代数，转移映射均满，且
$\mathcal{A} = \colim \mathcal{A}_i$。为完成证明，还须说明当 $i$
充分大时，$\mathcal{A}_i$ 是有限 $\mathcal{O}_X$-代数。
为此，选取平展满射 $U \to X$，其中 $U$ 为仿射概形。
取生成元 $f_1,\ldots,f_m \in \Gamma(U,\mathcal{F})$。
由于 $\mathcal{A}(U)$ 是有限 $\mathcal{O}_X(U)$-代数，对每个 $j$
都存在首一多项式 $P_j \in \mathcal{O}(U)[T]$，使得
$P_j(f_j)$ 在 $\mathcal{A}(U)$ 中为零。按构造
$\mathcal{A} = \colim \mathcal{A}_i$，故当 $i$ 充分大时，
$P_j(f_j) = 0$ 在 $\mathcal{A}_i(U)$ 中成立。对这样的 $i$，
代数 $\mathcal{A}_i$ 有限。
\end{proof}

\begin{lemma}
\label{spaces-limits-lemma-integral-algebra-directed-colimit-finite}
设 $S$ 为概形，$X$ 为 $S$ 上拟紧且拟分离的代数空间。
设 $\mathcal{A}$ 为整拟凝聚 $\mathcal{O}_X$-代数。则
\begin{enumerate}
\item $\mathcal{A}$ 是其有限拟凝聚 $\mathcal{O}_X$-子代数的有向余极限，
\item $\mathcal{A}$ 是有限且有限表示 $\mathcal{O}_X$-代数的有向余极限。
\end{enumerate}
\end{lemma}

\begin{proof}
由引理 \ref{spaces-limits-lemma-algebra-directed-colimit-finite-type}，有
$\mathcal{A} = \colim \mathcal{A}_i$，其中
% P09-ERR-0081
$\mathcal{A}_i \subset \mathcal{A}$ 遍历有限型拟凝聚
$\mathcal{O}_X$-子代数。$\mathcal{A}$ 的任一有限型拟凝聚
$\mathcal{O}_X$-子代数都是有限代数（在 $X$ 上平展的仿射概形上使用
《代数》，引理
\ref{algebra-lemma-characterize-finite-in-terms-of-integral}）。
这证明了 (1)。

\medskip\noindent
为证明 (2)，使用引理 \ref{spaces-limits-lemma-colimit-finitely-presented}，
将 $\mathcal{A} = \colim \mathcal{F}_i$ 写成有限表示
$\mathcal{O}_X$-模的余极限。对每个 $i$，令 $\mathcal{J}_i$
为映射
$$
\text{Sym}^*_{\mathcal{O}_X}(\mathcal{F}_i) \longrightarrow \mathcal{A}
$$
的核。当 $i' \geq i$ 时，有诱导映射
$\mathcal{J}_i \to \mathcal{J}_{i'}$，并且
$\mathcal{A} =
\colim \text{Sym}^*_{\mathcal{O}_X}(\mathcal{F}_i)/\mathcal{J}_i$.
此外，拟凝聚 $\mathcal{O}_X$-代数
$\text{Sym}^*_{\mathcal{O}_X}(\mathcal{F}_i)/\mathcal{J}_i$
都是有限代数（见上文）。将
$\mathcal{J}_i = \colim \mathcal{E}_{ik}$ 写成有限表示
$\mathcal{O}_X$-模的余极限。给定 $i' \geq i$ 与 $k$，存在 $k'$，
使得有映射 $\mathcal{E}_{ik} \to \mathcal{E}_{i'k'}$，从而图
$$
\xymatrix{
\mathcal{J}_i \ar[r] & \mathcal{J}_{i'} \\
\mathcal{E}_{ik} \ar[u] \ar[r] & \mathcal{E}_{i'k'} \ar[u]
}
$$
交换。这由《空间的上同调》，引理
\ref{spaces-cohomology-lemma-finite-presentation-quasi-compact-colimit}
得出。它诱导映射
$$
\mathcal{A}_{ik} =
\text{Sym}^*_{\mathcal{O}_X}(\mathcal{F}_i)/(\mathcal{E}_{ik})
\longrightarrow
\text{Sym}^*_{\mathcal{O}_X}(\mathcal{F}_{i'})/(\mathcal{E}_{i'k'}) =
\mathcal{A}_{i'k'}
$$
，其中 $(\mathcal{E}_{ik})$ 表示由 $\mathcal{E}_{ik}$ 生成的理想。
% P09-ERR-0082
拟凝聚 $\mathcal{O}_X$-代数 $\mathcal{A}_{ki}$
有限表示，并且当 $k$ 充分大时有限
（参见引理
\ref{spaces-limits-lemma-finite-algebra-directed-colimit-finite-finitely-presented}
的证明）。最后，有
$$
\colim \mathcal{A}_{ik} = \colim \mathcal{A}_i = \mathcal{A}
$$
。其中第一个等式已在引理
\ref{spaces-limits-lemma-finite-algebra-directed-colimit-finite-finitely-presented}
的证明中给出，第二个等式则源于 $\mathcal{A}$ 是模
$\mathcal{F}_i$ 的余极限。
\end{proof}

\begin{lemma}
\label{spaces-limits-lemma-extend}
设 $S$ 为概形，$X$ 为 $S$ 上拟紧且拟分离的代数空间。
设 $U \subset X$ 为拟紧开子空间。设 $\mathcal{F}$ 为拟凝聚
$\mathcal{O}_X$-模。设
$\mathcal{G} \subset \mathcal{F}|_U$ 为有限型拟凝聚
$\mathcal{O}_U$-子模。则存在有限型拟凝聚子模
$\mathcal{G}' \subset \mathcal{F}$，使得
$\mathcal{G}'|_U = \mathcal{G}$。
\end{lemma}

\begin{proof}
以 $j : U \to X$ 表示包含态射。由于 $X$ 拟分离且 $U$ 拟紧，
态射 $j$ 拟紧。因此
% P09-ERR-0083
$j_*\mathcal{G} \subset j_*\mathcal{F}|_U$ 的两端都是 $X$ 上的拟凝聚模
（《空间的态射》，引理
\ref{spaces-morphisms-lemma-pushforward}）。
令 $\mathcal{H} =
\Ker(j_*\mathcal{G} \oplus \mathcal{F} \to j_*\mathcal{F}|_U)$.
则 $\mathcal{H}|_U = \mathcal{G}$。由引理
\ref{spaces-limits-lemma-directed-colimit-finite-type}，可以找到有限型拟凝聚子模
$\mathcal{H}' \subset \mathcal{H}$，使得
$\mathcal{H}'|_U = \mathcal{H}|_U = \mathcal{G}$。
令 $\mathcal{G}' = \Im(\mathcal{H}' \to \mathcal{F})$ 即得结论。
\end{proof}









\section{相对逼近}
\label{spaces-limits-section-relative-approximation}

\noindent
我们讨论命题 \ref{spaces-limits-proposition-approximate} 在一个基底上的变体。

\begin{lemma}
\label{spaces-limits-lemma-approximate-morphism}
设 $f : X \to Y$ 为 $\mathbf{Z}$ 上拟紧且拟分离代数空间之间的态射。
则存在
% P09-ERR-0084
有向集 $I$ 以及以 $I$ 为指标的态射逆系统
% P09-ERR-0085
$(f_i : X_i \to Y_i)$，使得转移态射
$X_i \to X_{i'}$ 与 $Y_i \to Y_{i'}$ 均为仿射态射，
$X_i$ 与 $Y_i$ 均拟分离且在 $\mathbf{Z}$ 上有限型，并且
$(X \to Y) = \lim (X_i \to Y_i)$。
\end{lemma}

\begin{proof}
按命题 \ref{spaces-limits-proposition-approximate}，写成
$X = \lim_{a \in A} X_a$ 及 $Y = \lim_{b \in B} Y_b$；即
$X_a$ 与 $Y_b$ 均拟分离且在 $\mathbf{Z}$ 上有限型，
而转移态射均为仿射态射。

\medskip\noindent
固定 $b \in B$。将引理
\ref{spaces-limits-lemma-better-characterize-relative-limit-preserving}
应用于 $\mathbf{Z}$ 上的 $Y_b$ 与 $X = \lim X_a$，可知存在
$a \in A$ 及态射 $f_{a,b} : X_a \to Y_b$，使图
$$
\xymatrix{
X \ar[d] \ar[r] & Y \ar[d] \\
X_a \ar[r] & Y_b
}
$$
交换。令 $I$ 为以这种方式得到的三元组 $(a,b,f_{a,b})$ 的集合。

\medskip\noindent
设 $(a,b,f_{a,b})$ 与 $(a',b',f_{a',b'})$ 属于 $I$。令
% P09-ERR-0086
$b'' \leq \min(b,b')$。再次使用引理
\ref{spaces-limits-lemma-better-characterize-relative-limit-preserving}，存在
$a'' \geq \max(a,a')$，使两个复合
$X_{a''} \to X_a \to Y_b \to Y_{b''}$ 与
$X_{a''} \to X_{a'} \to Y_{b'} \to Y_{b''}$ 相等。
在 $I$ 上赋予预序
$$
(a, b, f_{a, b}) \geq (a', b', f_{a', b'})
\Leftrightarrow
a \geq a',\ b \geq b',\text{ 且 }
g_{b, b'} \circ f_{a, b} = f_{a', b'} \circ h_{a, a'}
$$
，其中 $h_{a,a'} : X_a \to X_{a'}$ 与
$g_{b,b'} : Y_b \to Y_{b'}$ 为转移态射。上述讨论表明 $I$ 有向，
而映射 $I \to A$、$(a,b,f_{a,b}) \mapsto a$ 与
% P09-ERR-0087
$I \to B$、$(a,b,f_{a,b})$ 都共尾。若对
$i = (a,b,f_{a,b})$ 令 $X_i = X_a$、$Y_i = Y_b$ 且
$f_i = f_{a,b}$，便得到以 $I$ 为指标的态射逆系统，并且
$$
\lim_{i \in I} X_i = \lim_{a \in A} X_a = X
\quad\text{且}\quad
% P09-ERR-0088
\lim_{i \in I} S_i = \lim_{b \in B} Y_b = Y
$$
；这由《范畴》，引理 \ref{categories-lemma-initial} 得出
（回忆，$I$ 上的极限其实是在与 $I$ 关联的对偶范畴上的极限，
因而共尾变为始）。证明完成。
\end{proof}

\begin{lemma}
\label{spaces-limits-lemma-relative-approximation}
设 $S$ 为概形，$f : X \to Y$ 为 $S$ 上代数空间的态射。假设
\begin{enumerate}
\item $X$ 拟紧且拟分离，并且
\item $Y$ 拟分离。
\end{enumerate}
则 $X = \lim X_i$ 是 $Y$ 上有限表示代数空间 $X_i$ 的一个有向逆系统
在 $Y$ 上的极限，其转移态射均为仿射态射。
\end{lemma}

\begin{proof}
由于 $|f|(|X|)$ 拟紧，可以将 $Y$ 替换为其点集包含
$|f|(|X|)$ 的拟紧开子空间。因此也可以假设 $Y$ 拟紧。
由引理 \ref{spaces-limits-lemma-approximate-morphism}，可将
$(X \to Y) = \lim (X_i \to Y_i)$ 写成某个态射有向逆系统的极限，
% P09-ERR-0089
其中的态射是 $\mathbf{Z}$ 上有限型概形的态射，转移态射均为仿射态射。
由于极限与极限交换
（《范畴》，引理 \ref{categories-lemma-colimits-commute}），
有 $X = \lim X_i \times_{Y_i} Y$。当 $i \geq i'$ 时，转移态射
$X_i \times_{Y_i} Y \to X_{i'} \times_{Y_{i'}} Y$ 仿射，因为它是复合
$$
X_i \times_{Y_i} Y \to X_i \times_{Y_{i'}} Y \to X_{i'} \times_{Y_{i'}} Y
$$
；其中第一个态射为闭浸入（由《空间的态射》，引理
\ref{spaces-morphisms-lemma-fibre-product-after-map}），第二个态射是
仿射态射的基变换（《空间的态射》，引理
\ref{spaces-morphisms-lemma-base-change-affine}），而仿射态射的复合仿射
（《空间的态射》，引理
\ref{spaces-morphisms-lemma-composition-affine}）。
态射 $f_i$ 有限表示（《空间的态射》，引理
\ref{spaces-morphisms-lemma-noetherian-finite-type-finite-presentation} 及
\ref{spaces-morphisms-lemma-finite-presentation-permanence}），
故基变换 $X_i \times_{f_i,Y_i} Y \to Y$ 有限表示
（《空间的态射》，引理
\ref{spaces-morphisms-lemma-base-change-finite-presentation}）。
\end{proof}






\section{有限型对象嵌入有限表示对象}
\label{spaces-limits-section-finite-type-closed-in-finite-presentation}

\noindent
本节是《极限》，第
\ref{limits-section-finite-type-closed-in-finite-presentation} 节的对应版本。

\begin{lemma}
\label{spaces-limits-lemma-affine-morphism-is-limit}
设 $S$ 为概形，$f : X \to Y$ 为 $S$ 上代数空间的仿射态射。
% P09-ERR-0090
若 $Y$ 拟紧且拟分离，则 $X$ 是一个有向极限
$X = \lim X_i$，其中每个 $X_i$ 在 $Y$ 上仿射且有限表示。
\end{lemma}

\begin{proof}
考虑拟凝聚 $\mathcal{O}_Y$-模
$\mathcal{A} = f_*\mathcal{O}_X$。由引理
\ref{spaces-limits-lemma-algebra-directed-colimit-finite-presentation}，
可将 $\mathcal{A} = \colim \mathcal{A}_i$ 写成有限表示
$\mathcal{O}_Y$-代数 $\mathcal{A}_i$ 的有向余极限。
令 $X_i = \underline{\Spec}_Y(\mathcal{A}_i)$；参见
《空间的态射》，定义
\ref{spaces-morphisms-definition-relative-spec}。
按构造，$X_i \to Y$ 仿射且有限表示，并且 $X = \lim X_i$。
\end{proof}

\begin{lemma}
\label{spaces-limits-lemma-integral-limit-finite-and-finite-presentation}
设 $S$ 为概形，$f : X \to Y$ 为 $S$ 上代数空间的整态射。
假设 $Y$ 拟紧且拟分离。则可以将 $X$ 写成有向极限
$X = \lim X_i$，其中 $X_i$ 在 $Y$ 上有限且有限表示。
\end{lemma}

\begin{proof}
考虑拟凝聚 $\mathcal{O}_Y$-模
$\mathcal{A} = f_*\mathcal{O}_X$。由引理
\ref{spaces-limits-lemma-integral-algebra-directed-colimit-finite}，
可将 $\mathcal{A} = \colim \mathcal{A}_i$ 写成有限且有限表示的
$\mathcal{O}_Y$-代数 $\mathcal{A}_i$ 的有向余极限。
令 $X_i = \underline{\Spec}_Y(\mathcal{A}_i)$；参见
《空间的态射》，定义
\ref{spaces-morphisms-definition-relative-spec}。
按构造，$X_i \to Y$ 有限且有限表示，并且 $X = \lim X_i$。
\end{proof}

\begin{lemma}
\label{spaces-limits-lemma-finite-in-finite-and-finite-presentation}
设 $S$ 为概形，$f : X \to Y$ 为 $S$ 上代数空间的有限态射。
假设 $Y$ 拟紧且拟分离。则可以将 $X$ 写成有向极限
$X = \lim X_i$，其中转移映射均为闭浸入，而对象
$X_i$ 在 $Y$ 上有限且有限表示。
\end{lemma}

\begin{proof}
考虑有限拟凝聚 $\mathcal{O}_Y$-模
$\mathcal{A} = f_*\mathcal{O}_X$。由引理
\ref{spaces-limits-lemma-finite-algebra-directed-colimit-finite-finitely-presented}，
可将 $\mathcal{A} = \colim \mathcal{A}_i$ 写成有限且有限表示的
$\mathcal{O}_Y$-代数 $\mathcal{A}_i$ 的有向余极限，
其转移映射均满。令
$X_i = \underline{\Spec}_Y(\mathcal{A}_i)$；参见《空间的态射》，定义
\ref{spaces-morphisms-definition-relative-spec}。
按构造，$X_i \to Y$ 有限且有限表示，转移映射均为闭浸入，
并且 $X = \lim X_i$。
\end{proof}

\begin{lemma}
\label{spaces-limits-lemma-closed-is-limit-closed-and-finite-presentation}
\begin{slogan}
拟紧拟分离代数空间的闭浸入可由有限表示闭浸入逼近。
\end{slogan}
设 $S$ 为概形，$f : X \to Y$ 为 $S$ 上代数空间的闭浸入。
假设 $Y$ 拟紧且拟分离。则可以将 $X$ 写成有向极限
$X = \lim X_i$，其中转移映射均为闭浸入，而态射
$X_i \to Y$ 均为有限表示闭浸入。
\end{lemma}

\begin{proof}
设 $\mathcal{I} \subset \mathcal{O}_Y$ 为将 $X$ 定义成 $Y$
的闭子空间的拟凝聚理想层。由引理
\ref{spaces-limits-lemma-directed-colimit-finite-type}，可将
$\mathcal{I} = \colim \mathcal{I}_i$ 写成其有限型拟凝聚子模的滤过余极限。
% P09-ERR-0091
令 $X_i$ 为由 $\mathcal{I}_i$ 截出的 $X$ 的闭子空间。
则 $X_i \to Y$ 为有限表示闭浸入，且 $X = \lim X_i$。略去若干细节。
\end{proof}

\begin{lemma}
\label{spaces-limits-lemma-quasi-affine-closed-in-finite-presentation}
设 $S$ 为概形，$f : X \to Y$ 为 $S$ 上代数空间的态射。假设
\begin{enumerate}
\item $f$ 局部有限型且拟仿射，并且
\item $Y$ 拟紧且拟分离。
\end{enumerate}
则存在有限表示态射 $f' : X' \to Y$，以及 $Y$ 上的闭浸入
$X \to X'$。
\end{lemma}

\begin{proof}
由《空间的态射》，引理
\ref{spaces-morphisms-lemma-characterize-quasi-affine}，可找到分解
$X \to Z \to Y$，其中 $X \to Z$ 是拟紧开浸入，而
$Z \to Y$ 仿射。写成 $Z = \lim Z_i$，其中 $Z_i$ 在 $Y$ 上仿射且
有限表示（引理 \ref{spaces-limits-lemma-affine-morphism-is-limit}）。
对某个 $0 \in I$，可找到拟紧开子空间 $U_0 \subset Z_0$，
使 $X$ 同构于 $U_0$ 在 $Z$ 中的逆像
（引理 \ref{spaces-limits-lemma-descend-opens}）。令 $U_i$ 为 $U_0$ 在 $Z_i$
中的逆像，则
% P09-ERR-0092
$U = \lim U_i$。由引理 \ref{spaces-limits-lemma-finite-type-eventually-closed}，
当某个 $i$ 充分大时，$X \to U_i$ 为闭浸入。令 $X'=U_i$ 即完成证明。
\end{proof}

\begin{lemma}
\label{spaces-limits-lemma-finite-type-closed-in-finite-presentation}
设 $S$ 为概形，$f : X \to Y$ 为 $S$ 上代数空间的态射。假设：
\begin{enumerate}
% P09-ERR-0093
\item $f$ 局部有限型。
\item $X$ 拟紧且拟分离，并且
\item $Y$ 拟紧且拟分离。
\end{enumerate}
则存在有限表示态射 $f' : X' \to Y$，以及 $Y$ 上代数空间的闭浸入
$X \to X'$。
\end{lemma}

\begin{proof}
由命题 \ref{spaces-limits-proposition-approximate}，可将 $X = \lim_i X_i$ 写成极限，
其中 $X_i$ 拟分离且在 $\mathbf{Z}$ 上有限型，而转移态射
$f_{ii'} : X_i \to X_{i'}$ 均为仿射态射。考虑交换图
$$
\xymatrix{
X \ar[r] \ar[rd] & X_{i, Y} \ar[r] \ar[d] & X_i \ar[d] \\
& Y \ar[r] & \Spec(\mathbf{Z})
}
$$
注意 $X_i$ 在 $\Spec(\mathbf{Z})$ 上有限表示；参见《空间的态射》，引理
\ref{spaces-morphisms-lemma-noetherian-finite-type-finite-presentation}。
因此，由《空间的态射》，引理
\ref{spaces-morphisms-lemma-base-change-finite-presentation}，
基变换 $X_{i,Y} \to Y$ 有限表示。注意
% P09-ERR-0094
$\lim X_{i,Y} = X \times Y$，并且 $X \to X \times Y$ 为单态射。
由引理 \ref{spaces-limits-lemma-finite-type-eventually-closed}，当 $i$ 充分大时，
$X \to X_{i,Y}$ 为单态射。固定这样的 $i$。注意
$X \to X_{i,Y}$ 局部有限型（《空间的态射》，引理
\ref{spaces-morphisms-lemma-permanence-finite-type}），且为单态射，
故分离且局部拟有限（《空间的态射》，引理
\ref{spaces-morphisms-lemma-monomorphism-loc-finite-type-loc-quasi-finite}）。
因此 $X \to X_{i,Y}$ 可表。进而 $X \to X_{i,Y}$ 拟仿射，因为可以使用
《空间》，引理
\ref{spaces-lemma-representable-transformations-property-implication}
中的原理，以及概形态射的相应结果《态射进阶》，引理
\ref{more-morphisms-lemma-quasi-finite-separated-quasi-affine}。
于是引理 \ref{spaces-limits-lemma-quasi-affine-closed-in-finite-presentation}
给出分解 $X \to X' \to X_{i,Y}$，其中 $X \to X'$ 为闭浸入，
$X' \to X_{i,Y}$ 有限表示。最后，$X' \to Y$ 作为有限表示态射的复合
而有限表示（《空间的态射》，引理
\ref{spaces-morphisms-lemma-composition-finite-presentation}）。
\end{proof}

\begin{proposition}
\label{spaces-limits-proposition-separated-closed-in-finite-presentation}
设 $S$ 为概形。
% P09-ERR-0095
$f : X \to Y$ 为 $S$ 上代数空间的态射。假设
\begin{enumerate}
\item $f$ 有限型且分离，并且
\item $Y$ 拟紧且拟分离。
\end{enumerate}
则存在有限表示分离态射 $f' : X' \to Y$，以及 $Y$ 上的闭浸入
$X \to X'$。
\end{proposition}

\begin{proof}
由引理 \ref{spaces-limits-lemma-finite-type-closed-in-finite-presentation}，
存在闭浸入 $X \to Z$，其中 $Z/Y$ 有限表示。设
$\mathcal{I} \subset \mathcal{O}_Z$ 为拟凝聚理想层，它将
% P09-ERR-0096
$X$ 定义成 $Y$ 的闭子概形。由引理
\ref{spaces-limits-lemma-directed-colimit-finite-type}，可将 $\mathcal{I}$ 写成其
有限型拟凝聚理想层的有向余极限
$\mathcal{I} = \colim_{a \in A} \mathcal{I}_a$。
令 $X_a \subset Z$ 为由 $\mathcal{I}_a$ 定义的闭子空间。
它们组成以 $A$ 为指标的逆系统。转移态射
$X_a \to X_{a'}$ 仿射，因为它们是闭浸入。每个 $X_a$ 都拟紧且
拟分离，因为它是 $Z$ 的闭子空间，而依我们的假设，$Z$ 拟紧且拟分离。
由 $\mathcal{I} = \colim_{a \in A} \mathcal{I}_a$ 立即得到
$X = \lim_a X_a$。每个态射 $X_a \to Z$ 都有限表示；参见
《态射》，引理 \ref{morphisms-lemma-closed-immersion-finite-presentation}。
因此态射 $X_a \to Y$ 均有限表示。于是只需证明对某个
$a \in A$，$X_a \to Y$ 分离。由于已经假设 $X \to Y$ 分离，
这由引理 \ref{spaces-limits-lemma-eventually-separated} 得出。
\end{proof}






\section{固有态射的逼近}
\label{spaces-limits-section-approximate-proper}

\begin{lemma}
\label{spaces-limits-lemma-proper-limit-of-proper-finite-presentation}
设 $S$ 为概形，$f : X \to Y$ 为 $S$ 上代数空间的固有态射，
且 $Y$ 拟紧且拟分离。则 $X = \lim X_i$ 是代数空间 $X_i$
的有向极限，其中 $X_i$ 在 $Y$ 上固有且有限表示，而转移态射及态射
$X \to X_i$ 均为闭浸入。
\end{lemma}

\begin{proof}
由命题 \ref{spaces-limits-proposition-separated-closed-in-finite-presentation}，
可找到闭浸入 $X \to X'$，其中 $X'$ 在 $Y$ 上分离且有限表示。
由引理 \ref{spaces-limits-lemma-closed-is-limit-closed-and-finite-presentation}，
可写成 $X = \lim X_i$，其中 $X_i \to X'$ 为有限表示闭浸入。
我们断言，当所有 $i$ 充分大时，态射 $X_i \to Y$ 固有；
这将完成证明。

\medskip\noindent
为证明该断言，可以假设 $Y$ 为仿射概形；参见《空间的态射》，引理
\ref{spaces-morphisms-lemma-proper-local}。接着，使用 Chow 引理的弱版本；
参见《空间的上同调》，引理
\ref{spaces-cohomology-lemma-weak-chow}，得到图
$$
\xymatrix{
X' \ar[rd] & X'' \ar[d] \ar[l]^\pi \ar[r] & \mathbf{P}^n_Y \ar[dl] \\
& Y &
}
$$
，其中 $X'' \to \mathbf{P}^n_Y$ 为浸入，而
$\pi : X'' \to X'$ 固有且满。以
$X'_i \subset X''$（相应地 $\pi^{-1}(X)$）表示
$X_i \subset X'$（相应地 $X \subset X'$）的概形论逆像。
则 $\lim X'_i = \pi^{-1}(X)$。由于 $\pi^{-1}(X) \to Y$ 固有
（《空间的态射》，引理
\ref{spaces-morphisms-lemma-composition-proper}），
$\pi^{-1}(X) \to \mathbf{P}^n_Y$ 为闭浸入
（《空间的态射》，引理
\ref{spaces-morphisms-lemma-universally-closed-permanence} 及
\ref{spaces-morphisms-lemma-immersion-when-closed}）。
因此，由引理 \ref{spaces-limits-lemma-eventually-closed-immersion}，当 $i$ 充分大时，
$X'_i \to \mathbf{P}^n_Y$ 为闭浸入。于是 $X'_i$ 在 $Y$ 上固有。
对这样的 $i$，由《空间的态射》，引理
\ref{spaces-morphisms-lemma-image-proper-is-proper}，
态射 $X_i \to Y$ 固有。
\end{proof}

\begin{lemma}
\label{spaces-limits-lemma-proper-limit-of-proper-finite-presentation-noetherian}
设 $f : X \to Y$ 为 $\mathbf{Z}$ 上代数空间的固有态射，且
$Y$ 拟紧且拟分离。则存在有向集 $I$ 以及以 $I$ 为指标的代数空间
态射逆系统 $(f_i : X_i \to Y_i)$，使得转移态射
$X_i \to X_{i'}$ 与 $Y_i \to Y_{i'}$ 均为仿射态射，
$f_i$ 固有且有限表示，$Y_i$ 在 $\mathbf{Z}$ 上有限表示，并且
$(X \to Y) = \lim (X_i \to Y_i)$。
\end{lemma}

\begin{proof}
由引理 \ref{spaces-limits-lemma-proper-limit-of-proper-finite-presentation}，
可写成 $X = \lim_{k \in K} X_k$，其中 $X_k \to Y$ 固有且有限表示。
接着，由绝对诺特逼近（命题 \ref{spaces-limits-proposition-approximate}），
可写成 $Y = \lim_{j \in J} Y_j$，其中 $Y_j$ 在 $\mathbf{Z}$
上有限表示。对每个 $k$，存在 $j$ 以及有限表示态射
$X_{k,j} \to Y_j$，使得作为 $Y$ 上的代数空间，
$X_k \cong Y \times_{Y_j} X_{k,j}$；参见引理
\ref{spaces-limits-lemma-descend-finite-presentation}。增大 $j$ 后，可以假设
$X_{k,j} \to Y_j$ 固有；参见引理 \ref{spaces-limits-lemma-eventually-proper}。
% P09-ERR-0097
集合 $I$ 将由这些序对 $(k,j)$ 组成，相应态射为
$X_{k,j} \to Y_j$。对每个 $k' \geq k$，可以找到 $j' \geq j$
及位于 $Y_{j'} \to Y_j$ 之上的态射
% P09-ERR-0098
$X_{j',k'} \to X_{j,k}$，其到 $Y$ 的基变换给出态射
$X_{k'} \to X_k$（再次由引理
\ref{spaces-limits-lemma-descend-finite-presentation} 得出）。
这些态射构成该系统的转移态射。略去若干细节。
\end{proof}

\noindent
回忆有限型拟凝聚模的概形论支撑；参见《空间的态射》，定义
\ref{spaces-morphisms-definition-scheme-theoretic-support}。

\begin{lemma}
\label{spaces-limits-lemma-eventually-proper-support}
沿用情形 \ref{spaces-limits-situation-descent-property} 的假设和记号。
设 $\mathcal{F}_0$ 为拟凝聚 $\mathcal{O}_{X_0}$-模。
以 $\mathcal{F}$ 与 $\mathcal{F}_i$ 表示 $\mathcal{F}_0$
到 $X$ 与 $X_i$ 的拉回。假设
\begin{enumerate}
\item $f_0$ 局部有限型，
\item $\mathcal{F}_0$ 有限型，
\item $\mathcal{F}$ 的概形论支撑在 $Y$ 上固有。
\end{enumerate}
则对某个 $i$，$\mathcal{F}_i$ 的概形论支撑在 $Y_i$ 上固有。
\end{lemma}

\begin{proof}
可以用 $\mathcal{F}_0$ 的概形论支撑替换 $X_0$。由《空间的态射》，引理
\ref{spaces-morphisms-lemma-support-finite-type}，这保证 $X_i$
是 $\mathcal{F}_i$ 的支撑，而 $X$ 是 $\mathcal{F}$ 的支撑。
于是，若 $Z \subset X$ 表示 $\mathcal{F}$ 的概形论支撑，则
$Z \to X$ 为普遍同胚。由于依假设 $Z \to Y$ 固有，故
$X \to Y$ 固有；参见《态射》，引理
\ref{morphisms-lemma-image-proper-is-proper}。由引理
\ref{spaces-limits-lemma-eventually-proper}，对某个 $i$，
% P09-ERR-0099
$X_i \to Y$ 固有。进而由《空间的态射》，引理
\ref{spaces-morphisms-lemma-closed-immersion-proper} 及
\ref{spaces-morphisms-lemma-composition-proper}，
$\mathcal{F}_i$ 的概形论支撑 $Z_i$ 在 $Y$ 上固有。
\end{proof}








\section{嵌入仿射空间}
\label{spaces-limits-section-embedding}

\noindent
下面给出若干技术性引理，稍后证明 Chow 引理时将用到。

\begin{lemma}
\label{spaces-limits-lemma-embedding-into-affine-over-ls-qs}
设 $S$ 为概形，$f : U \to X$ 为 $S$ 上代数空间的态射。假设 $U$
是仿射概形，$f$ 局部有限型，并且 $X$ 拟分离且局部分离。那么存在
$X$ 上的浸入 $U \to \mathbf{A}^n_X$。
\end{lemma}

\begin{proof}
记 $U = \Spec(A)$。把 $A = \colim A_i$ 写成有限型
$\mathbf{Z}$-子代数的滤过余极限。对每个 $i$，态射
$U \to U_i = \Spec(A_i)$ 诱导一个态射
$$
U \longrightarrow X \times U_i
$$
于 $X$ 之上。在极限中，态射 $U \to X \times U$ 是浸入，因为
$X$ 局部分离，见《空间的态射》引理
\ref{spaces-morphisms-lemma-semi-diagonal}。由引理
\ref{spaces-limits-lemma-finite-type-eventually-closed}，对某个 $i$，
$U \to X \times U_i$ 是浸入。由于 $U_i$ 同构于
$\mathbf{A}^n_{\mathbf{Z}}$ 的一个闭子概形，引理得证。
\end{proof}

\begin{remark}
\label{spaces-limits-remark-cannot-embed-in-general}
我们已在《例》第 \ref{examples-section-embedding-affines} 节看到，
若去掉 $X$ 局部分离的假设，引理
\ref{spaces-limits-lemma-embedding-into-affine-over-ls-qs} 不再成立。这引出如下问题：
若去掉 $X$ 拟分离的假设，引理
\ref{spaces-limits-lemma-embedding-into-affine-over-ls-qs} 是否仍成立？如果你知道答案，
请发送电子邮件至
\href{mailto:stacks.project@gmail.com}{stacks.project@gmail.com}。
\end{remark}

\begin{lemma}
\label{spaces-limits-lemma-embedding-into-affine-over-qs}
设 $S$ 为概形，$f : Y \to X$ 为 $S$ 上代数空间的态射。假设 $X$
是诺特的，并且 $f$ 有限表示。那么存在稠密开集 $V \subset Y$ 以及浸入
$V \to \mathbf{A}^n_X$。
\end{lemma}

\begin{proof}
这些假设蕴含 $Y$ 是诺特的（《空间的态射》引理
\ref{spaces-morphisms-lemma-finite-presentation-noetherian}）。因此 $Y$
拟分离，从而有一个稠密开子概形（《空间的性质》命题
\ref{spaces-properties-proposition-locally-quasi-separated-open-dense-scheme}）。
于是可以假设 $Y$ 是诺特概形。去掉 $Y$ 的不可约分支之间的交（使用
《拓扑》引理 \ref{topology-lemma-Noetherian} 和《性质》引理
\ref{properties-lemma-Noetherian-topology}），可以假设 $Y$ 是若干不可约
诺特概形的不交并。由于存在浸入
$$
\mathbf{A}^n_X \amalg \mathbf{A}^m_X
\longrightarrow
\mathbf{A}^{\max(n, m) + 1}_X
$$
（略去细节），只需在 $Y$ 不可约的情形证明结论。

\medskip\noindent
假设 $Y$ 是不可约概形。令 $T \subset |X|$ 为 $f : Y \to X$ 的像的闭包。
注意，由于 $|Y|$ 和 $|X|$ 都是良态拓扑空间（《空间的性质》引理
\ref{spaces-properties-lemma-quasi-separated-sober}），$T$ 不可约，并且有
唯一的泛点 $\xi$，它是 $Y$ 的泛点 $\eta$ 的像。
% P09-ERR-0100
设 $\mathcal{I} \subset X$ 为切出 $T$ 上约化诱导空间结构的拟凝聚理想层
（《空间的性质》定义
\ref{spaces-properties-definition-reduced-induced-space}）。由于
$\mathcal{O}_{Y, \eta}$ 是阿廷局部环，对某个 $n > 0$ 有
$f^{-1}\mathcal{I}^n \mathcal{O}_{Y, \eta} = 0$。由于
$f^{-1}\mathcal{I}\mathcal{O}_Y$ 是有限型拟凝聚理想，可知对某个非空开集
$V \subset Y$，有 $f^{-1}\mathcal{I}^n\mathcal{O}_V = 0$。令
$Z \subset X$ 为由 $\mathcal{I}^n$ 切出的闭子空间。按构造，
$V \to Y \to X$ 通过 $Z$ 分解。由于
$\mathbf{A}^n_Z \to \mathbf{A}^n_X$ 是浸入，可以用 $Z$ 代替 $X$、用
$V$ 代替 $Y$。这样便归结到如下情形：$Y$ 和 $X$ 不可约，并且
$Y \to X$ 把 $Y$ 的泛点映到 $X$ 的泛点。

\medskip\noindent
假设 $Y$ 和 $X$ 不可约，$Y$ 是概形，并且 $Y \to X$ 把 $Y$ 的泛点映到
$X$ 的泛点。由《空间的性质》命题
\ref{spaces-properties-proposition-locally-quasi-separated-open-dense-scheme}，
$X$ 有一个稠密开子概形 $U \subset X$。选取非空仿射开集 $V \subset Y$，
使其在 $X$ 中的像包含于 $U$。由《态射》引理
\ref{morphisms-lemma-quasi-affine-finite-type-over-S}，可以把 $V \to U$
分解为 $V \to \mathbf{A}^n_U \to U$。再与
$\mathbf{A}^n_U \to \mathbf{A}^n_X$ 复合，便得到所需的浸入。
\end{proof}








\section{支撑在闭子集上的截面}
\label{spaces-limits-section-sections-with-support-in-closed}

\noindent
本节是《性质》第
\ref{properties-section-sections-with-support-in-closed} 节的类似版本。

\begin{lemma}
\label{spaces-limits-lemma-quasi-coherent-finite-type-ideals}
设 $S$ 为概形，$X$ 为拟紧拟分离代数空间，$U \subset X$ 为开子空间。
下列条件等价：
\begin{enumerate}
\item $U \to X$ 拟紧；
\item $U$ 拟紧；
\item 存在有限型拟凝聚理想层
$\mathcal{I} \subset \mathcal{O}_X$，使得
$|X| \setminus |U| = |V(\mathcal{I})|$。
\end{enumerate}
\end{lemma}

\begin{proof}
设 $W$ 为仿射概形，$\varphi : W \to X$ 为满平展态射，见
《空间的性质》引理
\ref{spaces-properties-lemma-quasi-compact-affine-cover}。若 (1) 成立，则
$\varphi^{-1}(U) \to W$ 拟紧，故 $\varphi^{-1}(U)$ 拟紧，进而 $U$
拟紧（因为 $|\varphi^{-1}(U)| \to |U|$ 是满射）。若 (2) 成立，则
$\varphi^{-1}(U)$ 拟紧；这是因为 $X$ 拟分离，故 $\varphi$ 拟紧
（《空间的态射》引理
\ref{spaces-morphisms-lemma-quasi-compact-quasi-separated-permanence}）。
于是由《性质》引理
\ref{properties-lemma-quasi-coherent-finite-type-ideals}，
$\varphi^{-1}(U) \to W$ 是概形的拟紧态射。再由《空间的态射》引理
\ref{spaces-morphisms-lemma-quasi-compact-local}，$U \to X$ 拟紧。因此
(1) 与 (2) 等价。

\medskip\noindent
假设 (1) 和 (2) 成立。由《空间的性质》引理
\ref{spaces-properties-lemma-reduced-closed-subspace}，存在唯一的拟凝聚理想层
$\mathcal{J}$，它切出 $|X| \setminus |U|$ 上的约化诱导闭子空间结构。
% P09-ERR-0101
注意，$\mathcal{J}|_U = \mathcal{O}_U$，它是有限型
$\mathcal{O}_U$-模。由于 $U$ 拟紧，由引理
\ref{spaces-limits-lemma-directed-colimit-finite-type}，存在有限型拟凝聚子层
$\mathcal{I} \subset \mathcal{J}$，满足
$\mathcal{I}|_U = \mathcal{J}|_U$。于是
$|X| \setminus |U| = |V(\mathcal{I})|$，从而得到 (3)。反过来，若
$\mathcal{I}$ 如 (3) 中所述，则把关于概形的引理（《性质》引理
\ref{properties-lemma-quasi-coherent-finite-type-ideals}）应用于 $W$ 上的
$\varphi^{-1}\mathcal{I}$，可知 $\varphi^{-1}(U) \subset W$ 是拟紧开集。
因此 (2) 成立。
\end{proof}

\begin{lemma}
\label{spaces-limits-lemma-sections-annihilated-by-ideal}
设 $S$ 为概形，$X$ 为 $S$ 上的代数空间，
$\mathcal{I} \subset \mathcal{O}_X$ 为拟凝聚理想层，$\mathcal{F}$ 为拟凝聚
$\mathcal{O}_X$-模。考察如下 $\mathcal{O}_X$-模层 $\mathcal{F}'$：
它把 $X_\etale$ 的每个对象 $U$ 映为模
$$
\mathcal{F}'(U)
=
\{s \in \mathcal{F}(U) \mid
\mathcal{I}s = 0\}
$$
假设 $\mathcal{I}$ 有限型。那么
\begin{enumerate}
\item $\mathcal{F}'$ 是拟凝聚 $\mathcal{O}_X$-模层；
\item 对 $X_\etale$ 中的仿射对象 $U$，有
$\mathcal{F}'(U) = \{s \in \mathcal{F}(U) \mid \mathcal{I}(U)s = 0\}$；
\item $\mathcal{F}'_x = \{s \in \mathcal{F}_x \mid \mathcal{I}_x s = 0\}$。
\end{enumerate}
\end{lemma}

\begin{proof}
显然，定义 $\mathcal{F}'$ 的规则给出 $\mathcal{F}$ 的子层。因此可以在
$X$ 上平展局部地验证其余断言。这样便约化到关于概形的情形，即《性质》
引理 \ref{properties-lemma-sections-annihilated-by-ideal}。
\end{proof}

\begin{definition}
\label{spaces-limits-definition-subsheaf-sections-annihilated-by-ideal}
设 $S$ 为概形，$X$ 为 $S$ 上的代数空间，
$\mathcal{I} \subset \mathcal{O}_X$ 为有限型拟凝聚理想层，
$\mathcal{F}$ 为拟凝聚 $\mathcal{O}_X$-模。上面引理
\ref{spaces-limits-lemma-sections-annihilated-by-ideal} 中定义的子层
$\mathcal{F}' \subset \mathcal{F}$ 称为{\it 被 $\mathcal{I}$ 零化的截面子层}。
\end{definition}

\begin{lemma}
\label{spaces-limits-lemma-push-sections-annihilated-by-ideal}
设 $S$ 为概形，$f : X \to Y$ 为 $S$ 上代数空间的拟紧拟分离态射。
设 $\mathcal{I} \subset \mathcal{O}_Y$ 为有限型拟凝聚理想层，
$\mathcal{F}$ 为拟凝聚 $\mathcal{O}_X$-模，并设
$\mathcal{F}' \subset \mathcal{F}$ 为被
$f^{-1}\mathcal{I}\mathcal{O}_X$ 零化的截面子层。那么
$f_*\mathcal{F}' \subset f_*\mathcal{F}$ 是被 $\mathcal{I}$ 零化的截面子层。
\end{lemma}

\begin{proof}
略。提示：$f$ 拟紧拟分离的假设蕴含 $f_*\mathcal{F}$ 拟凝聚
（《空间的态射》引理 \ref{spaces-morphisms-lemma-pushforward}），
故引理 \ref{spaces-limits-lemma-sections-annihilated-by-ideal} 可应用于
$\mathcal{I}$ 和 $f_*\mathcal{F}$。
\end{proof}

\noindent
下面考察支撑在闭子集上的截面层。它同样未必总是拟凝聚层；但若该闭集的
补集在给定代数空间中是``逆紧的''，它便是拟凝聚层。

\begin{lemma}
\label{spaces-limits-lemma-sections-supported-on-closed-subset}
设 $S$ 为概形，$X$ 为 $S$ 上的代数空间。设 $T \subset |X|$ 为闭子集，
$U \subset X$ 为满足 $T \amalg |U| = |X|$ 的开子空间。设
$\mathcal{F}$ 为拟凝聚 $\mathcal{O}_X$-模。考察如下
$\mathcal{O}_X$-模层 $\mathcal{F}'$：它把 $X_\etale$ 的每个对象
$\varphi : W \to X$ 映为模
$$
\mathcal{F}'(W)
=
\{s \in \mathcal{F}(W) \mid
s\text{ 的支撑包含于 }|\varphi|^{-1}(T)\}
$$
若 $U \to X$ 拟紧，则
\begin{enumerate}
% P09-ERR-0102
\item 当 $W$ 仿射时，存在有限生成理想
$I \subset \mathcal{O}_X(W)$，使得 $|\varphi|^{-1}(T) = V(I)$；
\item 对 (1) 中的 $W$ 和 $I$，有
$\mathcal{F}'(W) = \{x \in \mathcal{F}(W) \mid
I^nx = 0 \text{ 对某个 } n\}$；
\item $\mathcal{F}'$ 是拟凝聚 $\mathcal{O}_X$-模层。
\end{enumerate}
\end{lemma}

\begin{proof}
显然，定义 $\mathcal{F}'$ 的规则给出 $\mathcal{F}$ 的子层。因此可以在
$X$ 上平展局部地验证其余断言。这样便约化到关于概形的情形，即《性质》
引理 \ref{properties-lemma-sections-supported-on-closed-subset}。
\end{proof}

\begin{definition}
\label{spaces-limits-definition-subsheaf-sections-supported-on-closed}
设 $S$ 为概形，$X$ 为 $S$ 上的代数空间。设 $T \subset |X|$ 为闭子集，
其补集对应于开子空间 $U \subset X$，并且包含态射 $U \to X$ 拟紧。
设 $\mathcal{F}$ 为拟凝聚 $\mathcal{O}_X$-模。上面引理
\ref{spaces-limits-lemma-sections-supported-on-closed-subset} 中定义的拟凝聚子层
$\mathcal{F}' \subset \mathcal{F}$ 称为{\it 支撑在 $T$ 上的截面子层}。
\end{definition}

\begin{lemma}
\label{spaces-limits-lemma-push-sections-supported-on-closed-subset}
设 $S$ 为概形，$f : X \to Y$ 为 $S$ 上代数空间的拟紧拟分离态射，
$T \subset |Y|$ 为闭子集。假设 $|Y| \setminus T$ 对应于开子空间
$V \subset Y$，并且 $V \to Y$ 拟紧。设 $\mathcal{F}$ 为拟凝聚
$\mathcal{O}_X$-模，$\mathcal{F}' \subset \mathcal{F}$ 为支撑在
$|f|^{-1}T$ 上的截面子层。那么
$f_*\mathcal{F}' \subset f_*\mathcal{F}$ 是支撑在 $T$ 上的截面子层。
\end{lemma}

\begin{proof}
略。提示：% P09-ERR-0103
$|X| \setminus |f|^{-1}T$ 是开子空间 $U = f^{-1}V \subset X$ 的支撑。
由于 $V \to Y$ 拟紧，其基变换 $U \to X$ 也拟紧。$f$ 拟紧拟分离的假设
蕴含 $f_*\mathcal{F}$ 拟凝聚。因此，引理
\ref{spaces-limits-lemma-sections-supported-on-closed-subset} 既可应用于 $T$ 和
$f_*\mathcal{F}$，也可应用于 $|f|^{-1}T$ 和 $\mathcal{F}$。所述拟凝聚模的
相等性由定义立即得到。
\end{proof}















\section{仿射空间的刻画}
\label{spaces-limits-section-affine}

\noindent
本节是《极限》第 \ref{limits-section-affine} 节的类似版本。

\begin{lemma}
\label{spaces-limits-lemma-affine}
设 $S$ 为概形，$f : X \to Y$ 为 $S$ 上代数空间的态射。假设 $f$ 满且有限，
并假设 $X$ 仿射。那么 $Y$ 仿射。
\end{lemma}

\begin{proof}
我们可以且确实把 $f : X \to Y$ 看作 $\Spec(\mathbf{Z})$ 上代数空间的态射
（见《空间》定义 \ref{spaces-definition-base-change}）。注意，有限态射是
仿射且普遍闭的，见《空间的态射》引理
\ref{spaces-morphisms-lemma-integral-universally-closed}。由《空间的态射》
引理 \ref{spaces-morphisms-lemma-image-universally-closed-separated}，$Y$
是分离代数空间。由于 $f$ 满且 $X$ 拟紧，$Y$ 也拟紧。

\medskip\noindent
由引理 \ref{spaces-limits-lemma-finite-in-finite-and-finite-presentation}，可以写成
$X = \lim X_a$，其中每个 $X_a \to Y$ 都有限且有限表示。由引理
\ref{spaces-limits-lemma-limit-is-affine}，当 $a$ 充分大时 $X_a$ 仿射。因此可以且确实
假设 $f : X \to Y$ 有限、满且有限表示。

\medskip\noindent
由命题 \ref{spaces-limits-proposition-approximate}，可以把 $Y = \lim Y_i$ 写成
$\mathbf{Z}$ 上有限表示代数空间的有向极限。由引理
\ref{spaces-limits-lemma-descend-finite-presentation}，可以找到 $0 \in I$ 以及有限表示态射
$X_0 \to Y_0$，使得当 $i \geq 0$ 时
$X_i = X_0 \times_{Y_0} Y_i$，并且 $X = \lim_i X_i$。由引理
\ref{spaces-limits-lemma-descend-finite}，当 $i$ 充分大时 $X_i \to Y_i$ 有限。由引理
\ref{spaces-limits-lemma-descend-surjective}，当 $i$ 充分大时 $X_i \to Y_i$ 满。由引理
\ref{spaces-limits-lemma-limit-is-affine}，当 $i$ 充分大时 $X_i$ 仿射。因此当 $i$
充分大时，可以应用《空间的上同调》引理
\ref{spaces-cohomology-lemma-image-affine-finite-morphism-affine-Noetherian}
得出 $Y_i$ 仿射。这蕴含 $Y$ 仿射，证明完成。
\end{proof}

\begin{proposition}
\label{spaces-limits-proposition-affine}
设 $S$ 为概形，$f : X \to Y$ 为 $S$ 上代数空间的态射。假设 $X$ 仿射，
并且 $f$ 满且普遍闭\footnote{整态射是普遍闭的，见《空间的态射》引理
\ref{spaces-morphisms-lemma-integral-universally-closed}。}。那么 $Y$ 仿射。
\end{proposition}

\begin{proof}
我们可以且确实把 $f : X \to Y$ 看作 $\Spec(\mathbf{Z})$ 上代数空间的态射
（见《空间》定义 \ref{spaces-definition-base-change}）。由《空间的态射》
引理 \ref{spaces-morphisms-lemma-image-universally-closed-separated}，$Y$
是分离代数空间。继而由《空间的态射》引理
\ref{spaces-morphisms-lemma-affine-permanence}，$f$ 仿射。再由
《空间的态射》引理
\ref{spaces-morphisms-lemma-integral-universally-closed}，$f$ 是整态射。

\medskip\noindent
由上一段，可以假设 $f : X \to Y$ 满且整，$X$ 仿射，而 $Y$ 分离。
由于 $f$ 满且 $X$ 拟紧，还可推出 $Y$ 拟紧。

\medskip\noindent
考察层 $\mathcal{A} = f_*\mathcal{O}_X$。这是拟凝聚
$\mathcal{O}_Y$-代数层，见《空间的态射》引理
\ref{spaces-morphisms-lemma-pushforward}。由引理
\ref{spaces-limits-lemma-colimit-finitely-presented}，可以把
$\mathcal{A} = \colim_i \mathcal{F}_i$ 写成有限型
$\mathcal{O}_Y$-模的滤过余极限。令
$\mathcal{A}_i \subset \mathcal{A}$ 为由 $\mathcal{F}_i$ 生成的
$\mathcal{O}_Y$-子代数。由于代数映射
$\mathcal{O}_Y \to \mathcal{A}$ 是整的，每个 $\mathcal{A}_i$ 都是有限拟凝聚
$\mathcal{O}_Y$-代数。因此
$$
X_i = \underline{\Spec}_Y(\mathcal{A}_i) \longrightarrow Y
$$
是代数空间的有限态射。这里 $\underline{\Spec}$ 是《空间的态射》引理
\ref{spaces-morphisms-lemma-affine-equivalence-algebras} 的构造。显然
$X = \lim_i X_i$。因此由引理 \ref{spaces-limits-lemma-limit-is-affine}，当 $i$
充分大时，概形 $X_i$ 仿射。此外，由于 $X \to Y$ 通过每个 $X_i$ 分解，
$X_i \to Y$ 满。于是由引理 \ref{spaces-limits-lemma-affine}，$Y$ 仿射。
\end{proof}

\noindent
下面这个上述结果的推论可见于 \cite{CLO}。

\begin{lemma}
\label{spaces-limits-lemma-reduction-scheme}
\begin{reference}
\cite[3.1.12]{CLO}
\end{reference}
设 $S$ 为概形，$X$ 为 $S$ 上的代数空间。若 $X_{red}$ 是概形，
则 $X$ 是概形。
\end{lemma}

\begin{proof}
设 $U' \subset X_{red}$ 为仿射开子概形。令 $U \subset X$ 为对应于开子集
$|U'| \subset |X_{red}| = |X|$ 的开子空间。那么 $U' \to U$ 满且整。
因此由命题 \ref{spaces-limits-proposition-affine}，$U$ 仿射。这样，$X$ 的每个点都包含在
$X$ 的一个开子概形中，即 $X$ 是概形。
\end{proof}

\begin{lemma}
\label{spaces-limits-lemma-integral-universally-bijective-scheme}
设 $S$ 为概形，$f : X \to Y$ 为 $S$ 上代数空间的态射。假设 $f$ 是整态射，
并诱导双射 $|X| \to |Y|$。那么 $X$ 是概形当且仅当 $Y$ 是概形。
\end{lemma}

\begin{proof}
按照定义，整态射可表。因此若 $Y$ 是概形，则 $X$ 也是概形。反过来，
假设 $X$ 是概形。设 $U \subset X$ 为仿射开集。整态射是闭的，并且 $|f|$
是双射，故 $|f|(|U|) \subset |Y|$ 是开集，因为它是
$|f|(|X| \setminus |U|)$ 的补集。令 $V \subset Y$ 为满足
$|V| = |f|(|U|)$ 的开子空间，见《空间的性质》引理
\ref{spaces-properties-lemma-open-subspaces}。那么 $U \to V$ 整且满，
所以由命题 \ref{spaces-limits-proposition-affine}，$V$ 是仿射概形。证明完成。
\end{proof}

\begin{lemma}
\label{spaces-limits-lemma-check-closed-infinitesimally}
设 $S$ 为概形，$f : X \to B$ 和 $B' \to B$ 为 $S$ 上代数空间的态射。
假设
\begin{enumerate}
\item $B' \to B$ 是闭浸入；
\item $|B'| \to |B|$ 是双射；
\item $X \times_B B' \to B'$ 是闭浸入；
\item $X \to B$ 有限型，或者 $B' \to B$ 有限表示。
\end{enumerate}
那么 $f : X \to B$ 是闭浸入。
\end{lemma}

\begin{proof}
假设 (1) 和 (2) 蕴含 $B_{red} = B'_{red}$。置
$X' = X \times_B B'$。那么 % P09-ERR-0104
$X' \to X$ 是闭浸入，并且 $X'_{red} = X_{red}$。设 $U \to B$ 为平展态射，
其中 $U$ 仿射。那么 $X' \times_B U \to X \times_B U$ 是代数空间的闭浸入，
并在底层约化空间上诱导同构。由于 $X' \times_B U$ 是概形（因为
$B' \to B$ 和 $X' \to B'$ 可表），由引理 \ref{spaces-limits-lemma-reduction-scheme}，
$X \times_B U$ 也是概形。因此 $X \to B$ 也可表。这样便约化到概形的情形，
见《态射》引理 \ref{morphisms-lemma-check-closed-infinitesimally}。
\end{proof}










\section{概形的有限覆盖}
\label{spaces-limits-section-finite-cover}

\noindent
作为本章极限结果的一个应用，我们证明：对任意拟紧拟分离代数空间 $X$，
存在概形 $Y$ 以及满有限态射 $Y \to X$。我们将依赖已经证明的如下结果：
% P09-ERR-0105
可以找到一个由概形给出的有限整覆盖；该结果在《良态空间》第
\ref{decent-spaces-section-integral-cover} 节中证明。

\begin{proposition}
\label{spaces-limits-proposition-there-is-a-scheme-finite-over}
设 $S$ 为概形，$X$ 为 $S$ 上的拟紧拟分离代数空间。
\begin{enumerate}
\item 存在有限表示的满有限态射 $Y \to X$，其中 $Y$ 是概形；
\item 给定满平展态射 $U \to X$，可以选取 $Y \to X$，使得对每个
$y \in Y$，存在开邻域 $V \subset Y$，且 $V \to X$ 通过 $U$ 分解。
\end{enumerate}
\end{proposition}

\begin{proof}
(1) 是 (2) 中 $U = X$ 的特殊情形。取《良态空间》引理
\ref{decent-spaces-lemma-there-is-a-scheme-integral-over} 中的
$Y \to X$。选取有限仿射开覆盖 $Y = \bigcup V_j$，使得 $V_j \to X$
通过 $U$ 分解。可以写成 $Y = \lim Y_i$，其中 $Y_i \to X$ 有限且
有限表示，见引理 \ref{spaces-limits-lemma-integral-limit-finite-and-finite-presentation}。
当 $i$ 充分大时，代数空间 $Y_i$ 是概形，见引理
\ref{spaces-limits-lemma-limit-is-scheme}。当 $i$ 充分大时，可以找到仿射开集
$V_{i, j} \subset Y_i$，其在 $Y$ 中的逆像恢复 $V_j$，见引理
\ref{spaces-limits-lemma-descend-opens}。进一步增大 $i$ 后，$X$ 上的态射
$V_j \to U$ 来自 $X$ 上的态射 $V_{i, j} \to U$，见命题
\ref{spaces-limits-proposition-characterize-locally-finite-presentation}。证明完成。
\end{proof}

\begin{lemma}
\label{spaces-limits-lemma-integral-limit-finite-and-finite-presentation-refined}
设 $S$ 为概形，$f : X \to Y$ 为 $S$ 上代数空间的整态射。假设 $Y$ 拟紧且
拟分离。设 $V \subset Y$ 为拟紧开子空间，使得 $f^{-1}(V) \to V$ 有限且
有限表示。那么可以把 $X$ 写成有向极限 $X = \lim X_i$，其中
$f_i : X_i \to Y$ 有限且有限表示，并且对所有 $i$，
$f^{-1}(V) \to f_i^{-1}(V)$ 都是同构。
\end{lemma}

\begin{proof}
本引理是命题 \ref{spaces-limits-proposition-there-is-a-scheme-finite-over} 的小幅加强。
考察整拟凝聚 $\mathcal{O}_Y$-代数
$\mathcal{A} = f_*\mathcal{O}_X$。下一段将把
$\mathcal{A} = \colim \mathcal{A}_i$ 写成有限且有限表示的
$\mathcal{O}_Y$-代数 $\mathcal{A}_i$ 的有向余极限，并使得
$\mathcal{A}_i|_V = \mathcal{A}|_V$。完成这一步后，置
$X_i = \underline{\Spec}_Y(\mathcal{A}_i)$，见《空间的态射》定义
\ref{spaces-morphisms-definition-relative-spec}。按构造，$X_i \to Y$
有限且有限表示，$X = \lim X_i$，并且
$f_i^{-1}(V) = f^{-1}(V)$。

\medskip\noindent
关于代数的断言，其证明类似于引理
\ref{spaces-limits-lemma-integral-algebra-directed-colimit-finite} 第 (2) 部分的证明。
首先使用引理 \ref{spaces-limits-lemma-colimit-finitely-presented}，把
$\mathcal{A} = \colim \mathcal{F}_i$ 写成有限表示
$\mathcal{O}_Y$-模的余极限。由于 $\mathcal{A}|_V$ 是有限型
$\mathcal{O}_V$-模，可以且确实假设对所有 $i$，
$\mathcal{F}_i|_V \to \mathcal{A}|_V$ 都是满射。对每个 $i$，令
$\mathcal{J}_i$ 为下列映射的核：
% P09-ERR-0106
$$
\text{Sym}^*_{\mathcal{O}_X}(\mathcal{F}_i) \longrightarrow \mathcal{A}
$$
当 $i' \geq i$ 时，有诱导映射
$\mathcal{J}_i \to \mathcal{J}_{i'}$。我们有
$\mathcal{A} =
\colim \text{Sym}^*_{\mathcal{O}_X}(\mathcal{F}_i)/\mathcal{J}_i$。
此外，拟凝聚 $\mathcal{O}_X$-代数
$\text{Sym}^*_{\mathcal{O}_X}(\mathcal{F}_i)/\mathcal{J}_i$ 是有限的
（因为它们是整拟凝聚 $\mathcal{O}_Y$-代数 $\mathcal{A}$ 在
$\mathcal{O}_X$ 上的有限型拟凝聚子代数）。由上述满性，
$\text{Sym}^*_{\mathcal{O}_X}(\mathcal{F}_i)/\mathcal{J}_i$ 在 $V$
上的限制是 $\mathcal{A}|_V$。由于 $\mathcal{A}|_V$ 作为
$\mathcal{O}_Y$-代数有限表示，$\mathcal{J}_i|_V$ 作为
$\text{Sym}^*_{\mathcal{O}_X}(\mathcal{F}_i)|_V$ 的理想层有限生成。
把 $\mathcal{J}_i = \colim \mathcal{E}_{ik}$ 写成有限表示
$\mathcal{O}_X$-模的余极限。由上面的有限生成断言，可以且确实假设
$\mathcal{E}_{ik}|_V$ 作为
% P09-ERR-0107
$\text{Sym}^*_{\mathcal{O}_X}(\mathcal{F}_i)|_V$ 的理想层生成
$\mathcal{J}_i|_V$。给定 $i' \geq i$ 和 $k$，存在 $k'$，使得有映射
$\mathcal{E}_{ik} \to \mathcal{E}_{i'k'}$，并使图
$$
\xymatrix{
\mathcal{J}_i \ar[r] & \mathcal{J}_{i'} \\
\mathcal{E}_{ik} \ar[u] \ar[r] & \mathcal{E}_{i'k'} \ar[u]
}
$$
交换。这由《空间的上同调》引理
\ref{spaces-cohomology-lemma-finite-presentation-quasi-compact-colimit}
推出。它诱导映射
$$
\mathcal{A}_{ik} =
\text{Sym}^*_{\mathcal{O}_X}(\mathcal{F}_i)/(\mathcal{E}_{ik})
\longrightarrow
\text{Sym}^*_{\mathcal{O}_X}(\mathcal{F}_{i'})/(\mathcal{E}_{i'k'}) =
\mathcal{A}_{i'k'}
$$
其中 $(\mathcal{E}_{ik})$ 表示由 $\mathcal{E}_{ik}$ 生成的理想。
% P09-ERR-0108
当 $k$ 充分大时，拟凝聚 $\mathcal{O}_X$-代数 $\mathcal{A}_{ki}$
有限表示且有限（见引理
\ref{spaces-limits-lemma-finite-algebra-directed-colimit-finite-finitely-presented} 的证明）。
此外，按构造有 $\mathcal{A}_{ik}|_V = \mathcal{A}|_V$。最后，
$$
\colim \mathcal{A}_{ik} = \colim \mathcal{A}_i = \mathcal{A}
$$
其中第一个等式已在引理
\ref{spaces-limits-lemma-finite-algebra-directed-colimit-finite-finitely-presented} 的证明中
给出；第二个等式成立是因为 $\mathcal{A}$ 是模 $\mathcal{F}_i$ 的余极限。
\end{proof}

\begin{lemma}
\label{spaces-limits-lemma-there-is-a-scheme-finite-over-refined}
设 $S$ 为概形，$X$ 为 $S$ 上的拟紧拟分离代数空间，并且 $|X|$ 只有有限多个
不可约分支。
\begin{enumerate}
\item 存在有限表示的满有限态射 $f : Y \to X$，其中 $Y$ 是概形，并且
$f$ 在某个拟紧稠密开集 $U \subset X$ 上有限平展；
\item 给定满平展态射 $V \to X$，可以选取 $Y \to X$，使得对每个
$y \in Y$，存在开邻域 $W \subset Y$，且 $W \to X$ 通过 $V$ 分解。
\end{enumerate}
\end{lemma}

\begin{proof}
(1) 是 (2) 中 $V = X$ 的特殊情形。

\medskip\noindent
证明 (2)。取《良态空间》引理
\ref{decent-spaces-lemma-there-is-a-scheme-integral-over-refined} 中的
$\pi : Y \to X$，并令 $U \subset X$ 为拟紧稠密开集，使得
$\pi^{-1}(U) \to U$ 有限平展。选取有限仿射开覆盖
$Y = \bigcup W_j$，使得 $W_j \to X$ 通过 $V$ 分解。可以写成
$Y = \lim Y_i$，其中 $\pi_i : Y_i \to X$ 有限且有限表示，并且
$\pi^{-1}(U) \to \pi_i^{-1}(U)$ 是同构，见引理
\ref{spaces-limits-lemma-integral-limit-finite-and-finite-presentation-refined}。当 $i$
充分大时，代数空间 $Y_i$ 是概形，见引理 \ref{spaces-limits-lemma-limit-is-scheme}。
当 $i$ 充分大时，可以找到仿射开集 $W_{i, j} \subset Y_i$，其在 $Y$
中的逆像恢复 $W_j$，见引理 \ref{spaces-limits-lemma-descend-opens}。进一步增大 $i$ 后，
$X$ 上的态射 $W_j \to V$ 来自 $X$ 上的态射
% P09-ERR-0109
$W_{i, j} \to U$，见命题
\ref{spaces-limits-proposition-characterize-locally-finite-presentation}。证明完成。
\end{proof}

\begin{lemma}
\label{spaces-limits-lemma-there-is-a-scheme-finite-over-filtered}
设 $S$ 为概形，$X$ 为 $S$ 上的拟紧拟分离代数空间。存在 $t \geq 0$ 以及
闭子空间
$$
X \supset Z_0 \supset Z_1 \supset \ldots \supset Z_t = \emptyset
$$
使得 $Z_i \to X$ 有限表示，$Z_0 \subset X$ 是加厚，并且对每个
$i = 0, \ldots t - 1$，存在概形 $Y_i$ 以及有限表示的满有限态射
$Y_i \to Z_i$，该态射在 $Z_i \setminus Z_{i + 1}$ 上有限平展。
\end{lemma}

\begin{proof}
可以把 $X$ 看作 $\Spec(\mathbf{Z})$ 上的代数空间，见《空间》定义
\ref{spaces-definition-base-change} 和《空间的性质》定义
\ref{spaces-properties-definition-separated}。因此可以应用命题
\ref{spaces-limits-proposition-approximate}。由此可找到仿射态射 $X \to X_0$，其中
$X_0$ 在 $\mathbf{Z}$ 上有限表示。若能对 $X_0$ 证明本引理，则把分层和
态射拉回到 $X$，便得到 $X$ 的结论；略去一些细节。这样便约化到下一段
讨论的情形。

\medskip\noindent
假设 $X$ 在 $\mathbf{Z}$ 上有限表示。那么 $X$ 是诺特的，而 $|X|$ 是
有限维诺特拓扑空间（有有限多个不可约分支）。因此可以对
$\dim(|X|)$ 作归纳。以 $X$ 为靶的任何有限态射都有限表示，所以在证明的
其余部分可以忽略这一要求。由引理
\ref{spaces-limits-lemma-there-is-a-scheme-finite-over-refined}，存在满有限态射
$Y \to X$，它在某个稠密开集 $U \subset X$ 上有限平展。置 $Z_0 = X$，
并令 $Z_1 \subset X$ 为满足 $|Z_1| = |X| \setminus |U|$ 的约化闭子空间。
由归纳假设，存在整数 $t \geq 0$ 以及滤过
$$
Z_1 \supset Z_{1, 0} \supset Z_{1, 1} \supset \ldots
\supset Z_{1, t} = \emptyset
$$
由闭子空间组成，其中 $Z_{1, 0} \to Z_1$ 是加厚，并且存在有限满态射
$Y_{1, i} \to Z_{1, i}$，它们在
$Z_{1, i} \setminus Z_{1, i + 1}$ 上有限平展。由于 $Z_1$ 约化，
$Z_1 = Z_{1, 0}$。因此对 $i \geq 1$，可以置
$Z_i = Z_{1, i - 1}$ 和 $Y_i = Y_{1, i - 1}$，引理得证。
\end{proof}











\section{得到概形}
\label{spaces-limits-section-representable}

\noindent
下面再给出几种证明代数空间是概形的技巧。第一种技巧基于如下事实：
可以证明存在一个不是概形的极小闭子空间。

\begin{lemma}
\label{spaces-limits-lemma-minimal-closed-subspace}
设 $S$ 为概形，$X$ 为 $S$ 上的拟紧拟分离代数空间。若 $X$ 不是概形，
则存在闭子空间 $Z \subset X$，使得 $Z$ 不是概形，但每个真闭子空间
$Z' \subset Z$ 都是概形。
\end{lemma}

\begin{proof}
使用佐恩引理证明。令 $\mathcal{Z}$ 为所有不是概形的闭子空间 $Z$ 所成的集，
并按包含关系排序。由假设，$\mathcal{Z}$ 包含 $X$，故非空。若
$Z_\alpha$ 是 $\mathcal{Z}$ 的一个全序子集，则
$Z = \bigcap Z_\alpha$ 属于 $\mathcal{Z}$。事实上，
$$
Z = \lim Z_\alpha
$$
且转移态射仿射。因此可应用引理 \ref{spaces-limits-lemma-limit-is-scheme} 得出：若 $Z$
是概形，则某个 $Z_\alpha$ 也是概形。（即使 $Z = \emptyset$，这一论证也
成立；但注意，由引理 \ref{spaces-limits-lemma-limit-nonempty}，这种情形不会发生。）
所以由佐恩引理，$\mathcal{Z}$ 有极小元。
\end{proof}

\noindent
现在可以证明这些极小非概形的一些性质。

\begin{lemma}
\label{spaces-limits-lemma-minimal-nonscheme}
设 $S$ 为概形，$X$ 为 $S$ 上的拟紧拟分离代数空间。假设每个真闭子空间
$Z \subset X$ 都是概形，但 $X$ 不是概形。那么 $X$ 约化且不可约。
\end{lemma}

\begin{proof}
由引理 \ref{spaces-limits-lemma-reduction-scheme}，$X$ 约化。选取闭子集
$T_1 \subset |X|$ 和 $T_2 \subset |X|$，使得
$|X| = T_1 \cup T_2$。若 $T_1$ 和 $T_2$ 都是真闭子集，则相应的约化诱导
闭子空间 $Z_1, Z_2 \subset X$（《空间的性质》定义
\ref{spaces-properties-definition-reduced-induced-space}）都是概形；而
$Z = Z_1 \times_X Z_2 = Z_1 \cap Z_2$ 作为 $Z_1$ 或 $Z_2$ 的闭子概形，
也是概形。注意，推出 $Z_1 \amalg_Z Z_2$ 在概形范畴中存在，见
《态射进阶》引理
\ref{more-morphisms-lemma-pushout-along-closed-immersions}。
% P09-ERR-0110
一种做法是证明 $Z_1 \amalg_Z Z_2$ 同构于 $X$，但这里不能使用这一点，
因为代数空间推出的材料在理论中稍后才出现。我们转而使用引理
\ref{spaces-limits-lemma-affine}，为每个点找到仿射邻域。具体地，设 $x \in |X|$。
若 $x \not \in Z_1$，则 $x$ 有一个是概形的邻域，即
$X \setminus Z_1$；当 $x \not \in Z_2$ 时亦同。若
$x \in Z = Z_1 \cap Z_2$，则选取包含 $z$ 的仿射开集
% P09-ERR-0111
$U \subset Z_1 \amalg_Z Z_2$。于是 $U_1 = Z_1 \cap U$ 和
$U_2 = Z_2 \cap U$ 是仿射开集，并且它们与 $Z$ 的交相同。由于
$|Z_1| = T_1$ 和 $|Z_2| = T_2$ 是 $|X|$ 的闭子集，且相交于 $|Z|$，
可找到开集 $W \subset |X|$，使得 $W \cap T_1 = |U_1|$ 且
$W \cap T_2 = |U_2|$。仍用 $W$ 表示 $X$ 中相应的开子空间。那么
$x \in |W|$，并且态射 $U_1 \amalg U_2 \to W$ 满且有限，其源是仿射概形。
因此由引理 \ref{spaces-limits-lemma-affine}，$W$ 是仿射概形。
\end{proof}

\noindent
下一个引理的关键在于，只需在 $X$ 的各点的像处检验条件。

\begin{lemma}
\label{spaces-limits-lemma-enough-local}
设 $f: X \to S$ 为从代数空间到概形 $S$ 的拟紧拟分离态射。若对每个
$x \in |X|$，记其像为 $s = f(x) \in S$，代数空间
$X \times_S \Spec(\mathcal{O}_{S,s})$ 都是概形，则 $X$ 是概形。
\end{lemma}

\begin{proof}
设 $x \in |X|$。只需找到 $s = f(x)$ 的开邻域 $U$，使得
$X \times_S U$ 是概形。由于 $X \times_S \Spec(\mathcal{O}_{S, s})$
是概形，并且
% P09-ERR-0112
$\mathcal{O}_{S, s} = \colim \mathcal{O}_S(U)$，其中余极限遍历 $S$ 中
$s$ 的仿射开邻域，所以
$$
X \times_S \Spec(\mathcal{O}_{S, s}) = \lim X \times_S U
$$
由引理 \ref{spaces-limits-lemma-limit-is-scheme}，对某个 $U$，$X \times_S U$ 是概形。
\end{proof}

\noindent
除了像引理 \ref{spaces-limits-lemma-enough-local} 那样限制到局部环之外，还可以限制到
基的闭子概形。

\begin{lemma}
\label{spaces-limits-lemma-maximal-ideal}
设 $\varphi : X \to \Spec(A)$ 为从代数空间到仿射概形的拟紧拟分离态射。
若 $X$ 不是概形，则存在理想 $I \subset A$，使得基变换 $X_{A/I}$
不是概形，但对每个满足 $I \subset I'$ 且 $I \not = I'$ 的 $I'$，
基变换 $X_{A/I'}$ 都是概形。
\end{lemma}

\begin{proof}
使用佐恩引理证明。令 $\mathcal{I}$ 为满足 $X_{A/I}$ 不是概形的理想 $I$
所成的集。由假设，$\mathcal{I}$ 包含 $(0)$。若 $I_\alpha$ 是
$\mathcal{I}$ 中的一条理想链，则 $I = \bigcup I_\alpha$ 属于
$\mathcal{I}$。事实上，$A/I = \colim A/I_\alpha$，从而
$$
X_{A/I} = \lim X_{A/I_\alpha}
$$
因此可应用引理 \ref{spaces-limits-lemma-limit-is-scheme} 得出：若 $X_{A/I}$ 是概形，
% P09-ERR-0113
则某个 $X_{A/I_\alpha}$ 也是概形。所以由佐恩引理，$\mathcal{I}$ 有极大元。
\end{proof}






\section{在闭纤维处黏合}
\label{spaces-limits-section-change-over-closed-points}

\noindent
把上面的理论应用于局部环的谱，可得到几个关于相对代数空间的良好黏合结果。
我们先证明一个辅助引理（它将在《自举》第
\ref{bootstrap-section-applications} 节中得到大幅推广）。

\begin{lemma}
\label{spaces-limits-lemma-relative-glueing}
设 $S = U \cup W$ 是概形的开覆盖。那么由基变换给出的函子
$$
FP_S \longrightarrow FP_U \times_{FP_{U \cap W}} FP_W
$$
是等价；其中 $FP_T$ 表示概形 $T$ 上有限表示代数空间的范畴。
\end{lemma}

\begin{proof}
首先，由于 $S = U \cup W$ 是 Zariski 覆盖，$(\Sch/S)_{fppf}$ 上的层范畴
等价于三元组 $(\mathcal{F}_U, \mathcal{F}_W, \varphi)$ 的范畴，其中
$\mathcal{F}_U$ 是 $(\Sch/U)_{fppf}$ 上的层，$\mathcal{F}_W$ 是
$(\Sch/W)_{fppf}$ 上的层，而
$$
\varphi :
\mathcal{F}_U|_{(\Sch/U \cap W)_{fppf}}
\longrightarrow
\mathcal{F}_W|_{(\Sch/U \cap W)_{fppf}}
$$
是同构。见《位点》引理 \ref{sites-lemma-mapping-property-glue}。（注意，
不需要其他黏合数据，因为 $U \times_S U = U$、$W \times_S W = W$；
同理，余圈条件自动成立。）现在，若 $(\Sch/S)_{fppf}$ 上的层
$\mathcal{F}$ 在此等价下映到
$(\mathcal{F}_U, \mathcal{F}_W, \varphi)$，则 $\mathcal{F}$ 是代数空间
当且仅当 $\mathcal{F}_U$ 和 $\mathcal{F}_W$ 都是代数空间。这由
《代数空间》引理 \ref{spaces-lemma-glueing-algebraic-spaces} 立即得到，
因为 $\mathcal{F}_U \to \mathcal{F}$ 和
$\mathcal{F}_W \to \mathcal{F}$ 可由开浸入表示，并覆盖 $\mathcal{F}$。
最后，在这种情形下，由《空间的态射》引理
\ref{spaces-morphisms-lemma-quasi-compact-local},
\ref{spaces-morphisms-lemma-separated-local}，以及
\ref{spaces-morphisms-lemma-finite-presentation-local}，代数空间
$\mathcal{F}$ 在 $S$ 上有限表示，当且仅当 $\mathcal{F}_U$ 在 $U$ 上
有限表示且 $\mathcal{F}_W$ 在 $W$ 上有限表示。
\end{proof}

\begin{lemma}
\label{spaces-limits-lemma-glueing-near-closed-point}
设 $S$ 为概形，$s \in S$ 为闭点，并且
$U = S \setminus \{s\} \to S$ 拟紧。置
$V = \Spec(\mathcal{O}_{S, s}) \setminus \{s\}$，则有范畴等价
$$
FP_S \longrightarrow FP_U \times_{FP_V} FP_{\Spec(\mathcal{O}_{S, s})}
$$
其中 $FP_T$ 表示 $T$ 上有限表示代数空间的范畴。
\end{lemma}

\begin{proof}
设 $W \subset S$ 为 $s$ 的开邻域。由引理 \ref{spaces-limits-lemma-relative-glueing}，函子
$$
FP_S \to FP_U \times_{FP_{W \setminus \{s\}}} FP_W
$$
是范畴等价。当 $W$ 遍历 $s$ 的仿射开邻域时，
$\mathcal{O}_{S, s} = \colim \mathcal{O}_W(W)$，从而
$\Spec(\mathcal{O}_{S, s}) = \lim W$。因此，
% P09-ERR-0114
$\Spec(\mathcal{O}_{S, s})$ 上有限表示代数空间的范畴，是 $W$ 上有限表示
代数空间的范畴的极限，其中 $W$ 遍历 $s$ 的仿射开邻域；见引理
\ref{spaces-limits-lemma-descend-finite-presentation}。对每个仿射开集 $s \in W$，
% P09-ERR-0115
由于 $U \to S$ 拟紧，$U \cap W$ 拟紧。因此
$V = \lim W \cap U = \lim W \setminus \{s\}$ 是拟紧拟分离概形的极限
（见《极限》引理
\ref{limits-lemma-directed-inverse-system-has-limit}）。因而 $V$ 上有限表示
代数空间的范畴也是 $W \cap U$ 上有限表示代数空间的各范畴的极限，
其中 $W$ 遍历 $s$ 的仿射开邻域。把这些结果结合起来即形式地得到本引理。
\end{proof}

\begin{lemma}
\label{spaces-limits-lemma-glueing-near-point}
设 $S$ 为概形，$U \subset S$ 为逆紧开集，$s \in S$ 为 $U$ 的补集中的点。
置 $V = \Spec(\mathcal{O}_{S, s}) \cap U$，则有范畴等价
$$
\colim_{s \in U' \supset U\text{ 为开集}} FP_{U'}
\longrightarrow
FP_U \times_{FP_V} FP_{\Spec(\mathcal{O}_{S, s})}
$$
其中 $FP_T$ 表示 $T$ 上有限表示代数空间的范畴。
\end{lemma}

\begin{proof}
设 $W \subset S$ 为 $s$ 的开邻域。由引理
\ref{spaces-limits-lemma-relative-glueing}，函子
$$
FP_{U \cup W}
\longrightarrow
FP_U \times_{FP_{U \cap W}} FP_W
$$
是范畴等价。当 $W$ 遍历 $s$ 的仿射开邻域时，有
$\mathcal{O}_{S, s} = \colim \mathcal{O}_W(W)$，从而
$\Spec(\mathcal{O}_{S, s}) = \lim W$。因此，
$\Spec(\mathcal{O}_{S, s})$ 上有限表示代数空间的范畴，是 $W$ 上有限表示
代数空间的范畴的极限，其中 $W$ 遍历 $s$ 的仿射开邻域；见引理
\ref{spaces-limits-lemma-descend-finite-presentation}。对每个仿射开集 $s \in W$，
由于 $U \to S$ 拟紧，$U \cap W$ 拟紧。因此
$V = \lim W \cap U$ 是拟紧拟分离概形的极限（见《极限》引理
\ref{limits-lemma-directed-inverse-system-has-limit}）。因而 $V$ 上有限表示
代数空间的范畴也是 $W \cap U$ 上有限表示代数空间的各范畴的极限，
其中 $W$ 遍历 $s$ 的仿射开邻域。把这些结果结合起来即形式地得到本引理。
\end{proof}

\begin{lemma}
\label{spaces-limits-lemma-glueing-near-multiple-closed-points}
设 $S$ 为概形，$s_1, \ldots, s_n \in S$ 为两两不同的闭点，并且
$U = S \setminus \{s_1, \ldots, s_n\} \to S$ 拟紧。置
$S_i = \Spec(\mathcal{O}_{S, s_i})$ 和 $U_i = S_i \setminus \{s_i\}$，
则有范畴等价
$$
FP_S \longrightarrow
FP_U \times_{(FP_{U_1} \times \ldots \times FP_{U_n})}
(FP_{S_1} \times \ldots \times FP_{S_n})
$$
其中 $FP_T$ 表示 $T$ 上有限表示代数空间的范畴。
\end{lemma}

\begin{proof}
当 $n = 1$ 时，这就是引理 \ref{spaces-limits-lemma-glueing-near-closed-point}。
当 $n > 1$ 时，可以用完全相同的方法证明本引理，也可以从该引理推出。
例如，假设 $f_i : X_i \to S_i$ 是 $FP_{S_i}$ 的对象，$f : X \to U$
是 $FP_U$ 的对象，并给定同构
$X_i \times_{S_i} U_i = X \times_U U_i$。由引理
\ref{spaces-limits-lemma-glueing-near-closed-point}，可以找到有限表示态射
$f' : X' \to U' = S \setminus \{s_1, \ldots, s_{n - 1}\}$，它在 $S_i$
上同构于 $X_i$，在 $U$ 上同构于 $X$，并且这些同构与给定同构
% P09-ERR-0116
$X_i \times_{S_n} U_n = X \times_U U_n$ 相容。然后可以对
$f_i : X_i \to S_i$（$i \leq n - 1$）、$f' : X' \to U'$ 以及诱导同构
$X_i \times_{S_i} U_i = X' \times_{U'} U_i$（$i \leq n - 1$）应用归纳。
这证明了本质满性。% P09-ERR-0117
全忠实性的证明从略。
\end{proof}







\section{对修改的应用}
\label{spaces-limits-section-modifications-at-a-point}

\noindent
利用极限，可以用亨泽尔局部环来描述良态代数空间在一个闭点之上的修改范畴。

\begin{lemma}
\label{spaces-limits-lemma-excision-modifications}
设 $S$ 为概形。考察 $S$ 上代数空间的分离平展态射
$f : V \to W$。假设存在闭子空间 $T \subset W$，使得
$f^{-1}T \to T$ 是同构。置 $W^0 = W \setminus T$ 和
$V^0 = f^{-1}W^0$，则基变换函子
$$
\left\{
\begin{matrix}
g : X \to W\text{ 为代数空间的态射} \\
g^{-1}(W^0) \to W^0\text{ 是同构}
\end{matrix}
\right\}
\longrightarrow
\left\{
\begin{matrix}
h : Y \to V\text{ 为代数空间的态射} \\
h^{-1}(V^0) \to V^0\text{ 是同构}
\end{matrix}
\right\}
$$
是范畴等价。
\end{lemma}

\begin{proof}
由于 $V \to W$ 分离，对 $V \times_W V$ 的某个开闭子空间 $U$，有
$V \times_W V = \Delta(V) \amalg U$。由 $f^{-1}T \to T$ 是同构这一
假设，$U \times_W T = \emptyset$；也就是说，% P09-ERR-0118
两个投影 $U \to V$ 的像都包含于 $V^0$。

\medskip\noindent
给定右侧范畴中的 $h : Y \to V$，考察 $(\Sch/S)_{fppf}$ 上由下式定义的
反变函子 $X$：
$$
X(T) = \{(w, y) \mid
w :  T \to W,\ y : T \times_{w, W} V \to Y\text{ 为 }V\text{ 上的态射}\}
$$
% P09-ERR-0119
用 $g : X \to W$ 表示把 $(w, y) \in X(T)$ 映到 $w \in W(T)$ 的映射。
由于 $h^{-1}V^0 \to V^0$ 是同构，若 $w : T \to W$ 的像包含于 $W^0$，
% P09-ERR-0120
则 $h$ 有唯一选择。换言之，$X \times_{g, W} W^0 = W^0$。另一方面，
考察 $X \times_{g, W, f} V$ 的一个 $T$-值点 $(w, y, v)$。那么
$w = f \circ v$，并且
% P09-ERR-0121
$$
y : T \times_{f \circ v, W} V \longrightarrow V
$$
是 $V$ 上的态射。考察态射
$$
T \times_{f \circ v, W} V \xrightarrow{(v, \text{id}_V)}
V \times_W V = V \amalg U
$$
$V$ 的逆像是通过
$(\text{id}_T, v) : T \to T \times_{f \circ v, W} V$ 嵌入的 $T$。
复合 $y' = y \circ (\text{id}_T, v) : T \to Y$ 是满足
$v = h \circ y'$ 的态射，并且它确定 $y$；这是因为第二个投影把 $U$
映入 $V^0$，所以 $y$ 在另一部分上的限制被唯一确定。由此
$X \times_{g, W, f} V \to Y$，$(w, y, v) \mapsto y'$ 是同构。

\medskip\noindent
所以只要证明 $X$ 是代数空间，就完成了证明。由于 $V \to W$ 分离且平展，
由《空间的态射》引理
\ref{spaces-morphisms-lemma-locally-quasi-finite-separated-representable}
（以及《空间的态射》引理
\ref{spaces-morphisms-lemma-etale-locally-quasi-finite}），它可表。当然，
$W^0 \to W$ 作为开浸入可表且平展。因此
$$
W^0 \amalg Y = X \times_{g, W} W^0 \amalg X \times_{g, W, f} V
= X \times_{g, W} (W^0 \amalg V) \longrightarrow X
$$
由《空间》引理 \ref{spaces-lemma-base-change-representable-transformations}
和 \ref{spaces-lemma-base-change-representable-transformations-property}
可表、满且平展。因此由《空间》引理
\ref{spaces-lemma-etale-locally-representable-by-space-gives-space}，$X$
是代数空间。
\end{proof}

\begin{lemma}
\label{spaces-limits-lemma-excision-modifications-properties}
记号和假设同引理 \ref{spaces-limits-lemma-excision-modifications}。设 $g : X \to W$
与 $h : Y \to V$ 在该等价下对应。那么 $g$ 拟紧、拟分离、分离、
局部有限表示、有限表示、局部有限型、有限型、固有、整、有限，
% P09-ERR-0122
以及在此添加更多性质，当且仅当 $h$ 具有相应性质。
\end{lemma}

\begin{proof}
若 $g$ 拟紧、拟分离、分离、局部有限表示、有限表示、局部有限型、有限型、
固有、有限，则作为 $g$ 的基变换，$h$ 也具有相应性质；见《空间的态射》
引理
\ref{spaces-morphisms-lemma-base-change-quasi-compact},
\ref{spaces-morphisms-lemma-base-change-separated},
\ref{spaces-morphisms-lemma-base-change-finite-presentation},
\ref{spaces-morphisms-lemma-base-change-finite-type},
\ref{spaces-morphisms-lemma-base-change-proper},
% P09-ERR-0123
\ref{spaces-morphisms-lemma-base-change-integral}。反过来，设 $P$ 是代数空间
态射的一项性质，它在基上平展局部，并且任意代数空间的恒等态射都具有
该性质。由于 $\{W^0 \to W, V \to W\}$ 是平展覆盖，要证明 $g$ 具有
$P$，只需证明 $h$ 具有 $P$。因此使用《空间的态射》引理
\ref{spaces-morphisms-lemma-quasi-compact-local},
\ref{spaces-morphisms-lemma-separated-local},
\ref{spaces-morphisms-lemma-finite-presentation-local},
\ref{spaces-morphisms-lemma-finite-type-local},
\ref{spaces-morphisms-lemma-proper-local},
\ref{spaces-morphisms-lemma-integral-local}.
即可得出结论。
\end{proof}

\begin{lemma}
\label{spaces-limits-lemma-modifications}
设 $S$ 为概形，$X$ 为 $S$ 上的良态代数空间。设 $x \in |X|$ 为闭点，
并且 $U = X \setminus \{x\} \to X$ 拟紧。置
$V = \Spec(\mathcal{O}_{X, x}^h) \setminus \{\mathfrak m_x^h\}$，
则基变换函子
$$
\left\{
\begin{matrix}
f : Y \to X\text{ 有限表示} \\
f^{-1}(U) \to U\text{ 是同构}
\end{matrix}
\right\}
\longrightarrow
\left\{
\begin{matrix}
g : Y \to \Spec(\mathcal{O}_{X, x}^h)\text{ 有限表示} \\
g^{-1}(V) \to V\text{ 是同构}
\end{matrix}
\right\}
$$
是范畴等价。
\end{lemma}

\begin{proof}
取《良态空间》引理
\ref{decent-spaces-lemma-decent-space-elementary-etale-neighbourhood} 中
$x$ 的初等平展邻域 $a : (W, w) \to (X, x)$，其中 $W$ 仿射。由于 $x$
是 $X$ 的闭点，而 $w$ 是 $W$ 中位于 $x$ 上方的唯一一点，$w$ 是 $W$
的闭点。由于 $a$ 平展，并把 $x$ 和 $w$ 处的剩余域等同起来，$a$
诱导同构 $a^{-1}x \to x$（视为 $X$ 和 $W$ 的闭子空间）。因此可以应用
引理 \ref{spaces-limits-lemma-excision-modifications} 和
\ref{spaces-limits-lemma-excision-modifications-properties}，把问题约化到 $X$ 是
仿射概形的情形。

\medskip\noindent
假设 $X$ 是仿射概形。回忆，$\mathcal{O}_{X, x}^h$ 是
$\Gamma(U, \mathcal{O}_U)$ 在仿射初等平展邻域
$(U, u) \to (X, x)$ 上的余极限。还回忆，这些邻域所成的范畴是余滤过的，
见《良态空间》引理
\ref{decent-spaces-lemma-elementary-etale-neighbourhoods} 或《态射进阶》引理
\ref{more-morphisms-lemma-elementary-etale-neighbourhoods}。于是
$\Spec(\mathcal{O}_{X, x}^h) = \lim U$，且
$V = \lim U \setminus \{u\}$（引理
\ref{spaces-limits-lemma-directed-inverse-system-has-limit}），其中两个极限取自同一范畴。
因此由引理 \ref{spaces-limits-lemma-descend-finite-presentation}，
% P09-ERR-0124
右侧范畴是各对 $(U, u)$ 所对应范畴的余极限。又由第一段的材料，
这些范畴中的每一个都等价于对 $(X, x)$ 所对应的范畴。证明完成。
\end{proof}






\section{普遍闭态射}
\label{spaces-limits-section-universally-closed}

\noindent
本节讨论拟紧（但未必分离）态射何时普遍闭。我们先证明一个引理，
它使我们能够在局部有限表示的基变换之后检验普遍闭性。

\begin{lemma}
\label{spaces-limits-lemma-separate}
设 $S$ 为概形，$f : X \to Y$ 和 $g : Z \to Y$ 为 $S$ 上代数空间的态射。
设 $z \in |Z|$，并设 $T \subset |X \times_Y Z|$ 为闭子集，且
$z \not \in \Im(T \to |Z|)$。若 $f$ 拟紧，则存在平展邻域
$(V, v) \to (Z, z)$、交换图
$$
\xymatrix{
V \ar[d] \ar[r]_a & Z' \ar[d]^b \\
Z \ar[r]^g & Y,
}
$$
以及闭子集 $T' \subset |X \times_Y Z'|$，使得
\begin{enumerate}
\item 态射 $b : Z' \to Y$ 局部有限表示；
\item 置 $z' = a(v)$，有 $z' \not \in \Im(T' \to |Z'|)$；
\item $T$ 在 $|X \times_Y V|$ 中的逆像经由
$|X \times_Y V| \to |X \times_Y Z'|$ 映入 $T'$。
\end{enumerate}
此外，可以假设 $V$ 和 $Z'$ 是仿射概形；若 $Z$ 是概形，还可以假设
$V$ 是 $z$ 的仿射开邻域。
\end{lemma}

\begin{proof}
我们将从概形态射的相应结果推出本引理。设 $y \in |Y|$ 为 $z$ 的像。
先选取仿射平展邻域 $(U, u) \to (Y, y)$，再选取仿射平展邻域
$(V, v) \to (Z, z)$，使态射 $V \to Y$ 通过 $U$ 分解。于是可以作如下替换：
\begin{enumerate}
\item 用 $X \times_Y U \to U$ 代替 $X \to Y$；
\item 用 $V \to U$ 代替 $Z \to Y$；
\item 用 $v$ 代替 $z$；
\item 用 $T$ 在
$|(X \times_Y U) \times_U V| = |X \times_Y V|$ 中的逆像代替 $T$。
\end{enumerate}
事实上，下面将证明：用 $v$ 的一个仿射开邻域代替 $V$ 后，对某个有限表示的
$Z' \to U$，存在态射 $a : V \to Z'$ 以及
$|(X \times_Y U) \times_U Z'| = |X \times_Y Z'|$ 的闭子集 $T'$，使得
$T$ 映入 $T'$ 且 $a(v) \not \in \Im(T' \to |Z'|)$。因此可以且确实
假设 $Z$ 和 $Y$ 是仿射概形，但要求所找到的解中 $V$ 是 $z$ 的开邻域。

\medskip\noindent
由于 $f$ 拟紧且 $Y$ 仿射，代数空间 $X$ 拟紧。选取仿射概形 $W$ 以及
满平展态射 $W \to X$。令 $T_W \subset |W \times_Y Z|$ 为 $T$ 的逆像。
那么 $z$ 不属于 $T_W$ 的像。由概形情形（《极限》引理
\ref{limits-lemma-separate}），可以找到 $z$ 的开邻域 $V \subset Z$、
% P09-ERR-0125
概形的交换图
$$
\xymatrix{
V \ar[d] \ar[r]_a & Z' \ar[d]^b \\
Z \ar[r]^g & Y,
}
$$
以及闭子集 $T' \subset |W \times_Y Z'|$，使得
\begin{enumerate}
\item 态射 $b : Z' \to Y$ 局部有限表示；
\item 置 $z' = a(z)$，有 % P09-ERR-0126
$z' \not \in \Im(T' \to Z')$；
\item $T_1 = T_W \cap |W \times_Y V|$ 经由
$|W \times_Y V| \to |W \times_Y Z'|$ 映入 $T'$。
\end{enumerate}
交换图
$$
\xymatrix{
W \times_Y Z \ar[d] &
W \times_Y V \ar[l] \ar[rr]_{a_1} \ar[d]_c & &
W \times_Y Z' \ar[d]^q \\
X \times_Y Z &
X \times_Y V \ar[l] \ar[rr]^{a_2} & &
X \times_Y Z'
}
$$
的方块是笛卡尔的，而竖直映射% P09-ERR-0127
是满、平展的，因而尤其是开的。考察左侧方块可知，
$T_1 = T_W \cap |W \times_Y V|$ 是
$T_2 = T \cap |X \times_Y V|$ 在 $c$ 下的逆像。由《空间的性质》引理
\ref{spaces-properties-lemma-points-cartesian}，有
$a_1(T_1) = q^{-1}(a_2(T_2))$。由《拓扑》引理
\ref{topology-lemma-open-morphism-quotient-topology}，有
$$
q^{-1}\left(\overline{a_2(T_2)}\right) =
\overline{q^{-1}(a_2(T_2))} =
\overline{a_1(T_1)} \subset T'
$$
由于 $q$ 满，而 $T'$ 的像不包含 $z'$，
$\overline{a_2(T_2)} \to |Z'|$ 的像也不包含 $z'$。因此可取上面含
$Z', V, a, b$ 的图以及闭子集
$\overline{a_2(T_2)} \subset |X \times_Y Z'|$ 作为引理所提问题的解。
\end{proof}

\begin{lemma}
\label{spaces-limits-lemma-test-universally-closed}
设 $S$ 为概形，$f : X \to Y$ 为 $S$ 上代数空间的拟紧态射。下列条件等价：
\begin{enumerate}
\item $f$ 普遍闭；
\item 对每个局部有限表示态射 $Z \to Y$，映射
$|X \times_Y Z| \to |Z|$ 都是闭的；
\item 存在概形 $V$ 以及满平展态射 $V \to Y$，使得对所有 $n \geq 0$，
$|\mathbf{A}^n \times (X \times_Y V)| \to |\mathbf{A}^n \times V|$
都是闭的。
\end{enumerate}
\end{lemma}

\begin{proof}
显然 (1) 蕴含 (2)。假设对 $S$ 上代数空间的某个态射 $Z \to Y$，映射
$|X \times_Y Z| \to |Z|$ 不是闭的。这意味着存在闭子集
$T \subset |X \times_Y Z|$，使得 $\Im(T \to |Z|)$ 不是闭集。选取
$z \in |Z|$，使其属于 $T$ 的像的闭包但不属于该像。应用引理
\ref{spaces-limits-lemma-separate}，可找到平展邻域 $(V, v) \to (Z, z)$、交换图
$$
\xymatrix{
V \ar[d] \ar[r]_a & Z' \ar[d]^b \\
Z \ar[r]^g & Y,
}
$$
以及闭子集 $T' \subset |X \times_Y Z'|$，使得
\begin{enumerate}
\item 态射 $b : Z' \to Y$ 局部有限表示；
\item 置 $z' = a(v)$，有 $z' \not \in \Im(T' \to |Z'|)$；
\item $T$ 在 $|X \times_Y V|$ 中的逆像经由
$|X \times_Y V| \to |X \times_Y Z'|$ 映入 $T'$。
\end{enumerate}
我们断言 $z'$ 属于 $\Im(T' \to |Z'|)$ 的闭包；这蕴含
$|X \times_Y Z'| \to |Z'|$ 不是闭映射。该断言说明 (2) 蕴含 (1)。
为验证该断言，考察% P09-ERR-0128
如下交换图：
$$
\xymatrix{
X \times_Y Z \ar[d] &
X \times_Y V \ar[l] \ar[d] \ar[r] &
X \times_Y Z' \ar[d] \\
Z & V \ar[l] \ar[r]^a & Z'
}
$$
令 $T_V \subset |X \times_Y V|$ 为 $T$ 的逆像。由《空间的性质》引理
\ref{spaces-properties-lemma-points-cartesian}，$T_V$ 在 $|V|$ 中的像是
$T$ 在 $|Z|$ 中的像的逆像。由于 $z$ 属于 $T \to |Z|$ 的像的闭包，
并且 $|V| \to |Z|$ 是开映射，$v$ 属于 $T_V \to |V|$ 的像的闭包。
由于 $T_V$ 在 $|X \times_Y Z'|$ 中的像包含于 % P09-ERR-0129
$|T'|$，立即可知 $z' = a(v)$ 属于 $T'$ 的像的闭包。

\medskip\noindent
显然 (1) 蕴含 (3)。取 (3) 中的 $V \to Y$。若能证明
$X \times_Y V \to V$ 普遍闭，则由《空间的态射》引理
\ref{spaces-morphisms-lemma-universally-closed-local}，$f$ 普遍闭。因此只需
证明如下断言：若 $f : X \to Y$ 是代数空间的拟紧态射，$Y$ 是概形，
并且对所有 $n$，$|\mathbf{A}^n \times X| \to
|\mathbf{A}^n \times Y|$ 都是闭的，则 $f$ 满足 (2)。设 $Z \to Y$
局部有限表示。我们要证明映射 $|X \times_Y Z| \to |Z|$ 是闭的。
这个问题在 $Z$ 上平展局部，故可以假设 $Z$ 仿射（略去一些细节）。
由于 $Y$ 是概形、$Z$ 仿射且 $Z \to Y$ 局部有限表示，可找到浸入
$Z \to \mathbf{A}^n \times Y$，见《态射》引理
\ref{morphisms-lemma-quasi-affine-finite-type-over-S}。考察笛卡尔图
$$
\vcenter{
\xymatrix{
X \times_Y Z \ar[d] \ar[r] & \mathbf{A}^n \times X \ar[d] \\
Z \ar[r] & \mathbf{A}^n \times Y
}
}
\quad
\begin{matrix}
\text{它诱导} \\
\text{笛卡尔方块}
\end{matrix}
\quad
\vcenter{
\xymatrix{
|X \times_Y Z| \ar[d] \ar[r] & |\mathbf{A}^n \times X| \ar[d] \\
|Z| \ar[r] & |\mathbf{A}^n \times Y|
}
}
$$
这是拓扑空间的笛卡尔方块，其水平箭头是到局部闭子集的同胚
（《空间的性质》引理
\ref{spaces-properties-lemma-subspace-induced-topology}）。因此，
% P09-ERR-0130
$|X \times_Y Z|$ 的每个闭子集 $T$ 都是
$|\mathbf{A}^n \times Y|$ 的某个闭子集 $T'$ 的拉回。由于假设断言
$T'$ 在 $|\mathbf{A}^n \times X|$ 中的像是闭集，可知 $T$ 在 $|Z|$
中的像是闭集，如所需。
\end{proof}

\begin{lemma}
\label{spaces-limits-lemma-limited-base-change}
设 $S$ 为概形，$f : X \to Y$ 为 $S$ 上代数空间的态射。
% P09-ERR-0131
假设 $f$ 分离且有限型。下列条件等价：
\begin{enumerate}
\item 态射 $f$ 固有；
\item 对任意局部有限表示态射 % P09-ERR-0132
$Y \to Z$，映射 $|X \times_Y Z| \to |Z|$ 是闭的；
\item 存在概形 $V$ 以及满平展态射 $V \to Y$，使得对所有 $n \geq 0$，
$|\mathbf{A}^n \times (X \times_Y V)| \to |\mathbf{A}^n \times V|$
都是闭的。
\end{enumerate}
\end{lemma}

\begin{proof}
由于固有态射恰好就是分离、有限型且普遍闭的态射，本引理是引理
\ref{spaces-limits-lemma-test-universally-closed} 的特殊情形。
\end{proof}


\section{诺特赋值判据}
\label{spaces-limits-section-Noetherian-valuative-criterion}

\noindent
我们已经在《空间的上同调》第
\ref{spaces-cohomology-section-Noetherian-valuative-criterion}.
节中证明了一些结果。概形的相应内容见《极限》第
\ref{limits-section-Noetherian-valuative-criterion} 节。

\medskip\noindent
本节的许多结果都可以（或许也应当）通过引用下述引理来证明，
尽管我们并非总是如此处理。

\begin{lemma}
\label{spaces-limits-lemma-reach-point-closure-Noetherian}
设 $S$ 为概形，$f : X \to Y$ 为 $S$ 上代数空间的态射。
% P09-ERR-0133
假设 $f$ 有限型且 $Y$ 局部诺特。设 $y \in |Y|$ 是 $|f|$ 的像的闭包中的点。
则存在交换图
$$
\xymatrix{
\Spec(K) \ar[r] \ar[d] & X \ar[d]^f \\
\Spec(A) \ar[r] & Y
}
$$
其中 $A$ 是离散赋值环，$K$ 是其分式域，并且 $\Spec(A)$ 的闭点映到 $y$。
此外，可以假设与 $\Spec(K) \to X$ 对应的点 $x \in |X|$ 是余维 $0$ 的点
\footnote{参见《空间的性质》第
\ref{spaces-properties-section-generic-points} 节中的讨论。}，
并且 $K$ 是某个在 $X$ 上平展的概形上的一点的剩余域。
\end{lemma}

\begin{proof}
选择仿射概形 $V$、点 $v \in V$ 以及把 $v$ 映到 $y$ 的平展态射
$V \to Y$。映射 $|V| \to |Y|$ 是开映射，而且由《空间的性质》引理
\ref{spaces-properties-lemma-points-cartesian}，
$|X \times_Y V| \to |V|$ 的像是 $|f|$ 的像的逆像。因此，点 $v$ 位于
$|X \times_Y V| \to |V|$ 的像的闭包中。若能对
$X \times_Y V \to V$ 和点 $v$ 证明本引理，那么关于 $f$ 和 $y$ 的结论
也随之成立。这样便归约到下一段所述的情形。

\medskip\noindent
设有引理中那样的 $f : X \to Y$ 和 $y \in |Y|$，其中 $Y$ 是仿射概形。
由于 $f$ 拟紧，可知 $X$ 拟紧。因此可以选择仿射概形 $W$ 和满平展态射
$W \to X$。于是 $|f|$ 的像与 $W \to Y$ 的像相同。这样便归约到概形的情形，
即《极限》引理 \ref{limits-lemma-reach-point-closure-Noetherian}。
\end{proof}

\noindent
首先陈述关于分离性的结果。我们将经常使用如下形状的交换实线图，
其中各箭头是基概形 $S$ 上代数空间的态射：
\begin{equation}
\label{spaces-limits-equation-valuative}
\vcenter{
\xymatrix{
\Spec(K) \ar[r] \ar[d] & X \ar[d] \\
\Spec(A) \ar[r] \ar@{-->}[ru] & Y
}
}
\end{equation}
这里 $A$ 是赋值环，$K$ 是其分式域。

\begin{lemma}
\label{spaces-limits-lemma-Noetherian-dvr-valuative-separation}
设 $S$ 为概形，$f : X \to Y$ 为 $S$ 上代数空间的态射。假设 $f$ 拟分离且
局部有限型，并且 $Y$ 局部诺特。下列条件等价：
\begin{enumerate}
\item 态射 $f$ 分离；
\item 对任意图 (\ref{spaces-limits-equation-valuative})，至多存在一条虚线箭头；
\item 对所有满足 $A$ 为离散赋值环的图 (\ref{spaces-limits-equation-valuative})，
至多存在一条虚线箭头；
\item 对所有满足 $A$ 为离散赋值环且 $\Spec(K) \to X$ 的像是 $X$ 上
余维 $0$ 的点的图 (\ref{spaces-limits-equation-valuative})，至多存在一条虚线箭头。
\end{enumerate}
\end{lemma}

\begin{proof}
由《空间的态射》引理
\ref{spaces-morphisms-lemma-separated-implies-valuative}，有
(1) $\Rightarrow$ (2)。(2) $\Rightarrow$ (3) 和
(3) $\Rightarrow$ (4) 是立即的。只需证明 (4) 蕴含 (1)。

\medskip\noindent
假设 (4)。必须证明对角态射 $\Delta : X \to X \times_Y X$ 是闭浸入。
我们已经知道 $\Delta$ 可表、分离、是单态射且局部有限型；参见
《空间的态射》引理
\ref{spaces-morphisms-lemma-properties-diagonal}。
选择仿射概形 $U$ 和平展态射 $U \to X \times_Y X$，并置
$V = X \times_{\Delta, X \times_Y X} U$。只需证明 $V \to U$ 是闭浸入
（《空间的态射》引理
\ref{spaces-morphisms-lemma-closed-immersion-local}）。
由于 $X \times_Y X$ 在 $Y$ 上局部有限型，可知 $U$ 诺特（使用
《空间的态射》诸引理
\ref{spaces-morphisms-lemma-composition-finite-type},
\ref{spaces-morphisms-lemma-base-change-finite-type} 和
\ref{spaces-morphisms-lemma-locally-finite-type-locally-noetherian}）。
注意，由于 $\Delta$ 可表，$V$ 是概形；又由于 $f$ 拟分离，$V$ 拟紧。
故 $V \to U$ 分离且有限型。考虑概形态射的交换图
$$
\xymatrix{
\Spec(K) \ar[r] \ar[d] & V \ar[d] \\
\Spec(A) \ar[r] \ar@{-->}[ru] & U
}
$$
其中 $A$ 是以 $K$ 为分式域的离散赋值环，而 $K$ 是诺特概形 $V$ 的某个
泛点的剩余域。由于 $V \to X$ 平展（它是平展态射
$U \to X \times_Y X$ 的基变换），可知 $\Spec(K) \to V \to X$ 的像是
余维 $0$ 的点；参见《空间的性质》第
\ref{spaces-properties-section-dimension-local-ring} 节。
可将复合 $\Spec(A) \to U \to X \times_Y X$ 解释为一对态射
$a, b : \Spec(A) \to X$：它们作为到 $Y$ 的态射相同，在限制到
$\Spec(K)$ 后也相同，并且该限制映到余维 $0$ 的点。因此假设 (4) 保证
$a = b$，于是得到图中的虚线箭头。由《极限》引理
\ref{limits-lemma-Noetherian-dvr-valuative-proper}，可知 $V \to U$ 固有。
换言之，$\Delta$ 固有。由于 $\Delta$ 是单态射，遂知 $\Delta$ 是闭浸入
（《平展态射》引理 \ref{etale-lemma-characterize-closed-immersion}），如所需。
\end{proof}

\begin{lemma}
\label{spaces-limits-lemma-Noetherian-dvr-valuative-proper}
设 $S$ 为概形，$f : X \to Y$ 为 $S$ 上代数空间的态射。假设 $f$ 拟分离且
有限型，并且 $Y$ 局部诺特。下列条件等价：
\begin{enumerate}
\item $f$ 固有；
\item $f$ 满足赋值判据，参见《空间的态射》定义
\ref{spaces-morphisms-definition-valuative-criterion}；
\item 对任意图 (\ref{spaces-limits-equation-valuative})，恰有一条虚线箭头；
\item 对所有满足 $A$ 为离散赋值环的图 (\ref{spaces-limits-equation-valuative})，
恰有一条虚线箭头；并且
\item 对所有满足 $A$ 为离散赋值环且 $\Spec(K) \to X$ 的像是 $X$ 上
余维 $0$ 的点的图 (\ref{spaces-limits-equation-valuative})，恰有一条虚线箭头
\footnote{还有一个更精细的表述：在存在性部分，只要求在扩张离散赋值环后
存在虚线箭头。}。 % P09-ERR-0134
\end{enumerate}
\end{lemma}

\begin{proof}
由《空间的态射》引理
\ref{spaces-morphisms-lemma-characterize-proper}，有
(1) $\Leftrightarrow$ (2) $\Leftrightarrow$ (3)。显然
(3) $\Rightarrow$ (4) $\Rightarrow$ (5)。最后证明 (5) 蕴含 (1)。

\medskip\noindent
假设 (5)。由引理 \ref{spaces-limits-lemma-Noetherian-dvr-valuative-separation}，可知
$f$ 分离。只需进一步证明 $f$ 普遍闭。设 $V \to Y$ 是平展态射，其中
$V$ 为仿射概形。只需证明基变换 $V \times_Y X \to V$ 普遍闭；参见
《空间的态射》引理
\ref{spaces-morphisms-lemma-universally-closed-local}。考虑交换图
$$
\xymatrix{
\Spec(K) \ar[r] \ar[d] & V \times_Y X \ar[d] \ar[r] & X \ar[d] \\
\Spec(A) \ar[r] \ar@{-->}[ru] \ar@{..>}[rru] & V \ar[r] & Y
}
$$
其箭头均为 $S$ 上代数空间的态射，其中 $A$ 是以 $K$ 为分式域的离散赋值环，
且 $\Spec(K) \to V \times_Y X$ 映到代数空间 $V \times_Y X$ 的余维 $0$
的点。由于 $V \times_Y X \to X$ 平展，$\Spec(K) \to X$ 的像是 $X$ 的
余维 $0$ 的点。因此由 (5)，得到图中两条虚线箭头中较长的一条，当然也就
得到较短的一条。这说明我们的假设对态射 $V \times_Y X \to V$ 成立，
从而归约到下一段所讨论的情形。

\medskip\noindent
% P09-ERR-0135
假设 $Y$ 为诺特仿射概形。在此情形下，$X$ 是 $Y$ 上有限型的分离诺特代数空间
（我们已经知道 $f$ 分离）。（特别地，由《空间的性质》命题
\ref{spaces-properties-proposition-locally-quasi-separated-open-dense-scheme}，
代数空间 $X$ 有一个为概形的稠密开子空间，不过严格说来我们不会用到这一点。）
按照 Chow 引理的弱形式（《空间的上同调》引理
\ref{spaces-cohomology-lemma-weak-chow}），选择 $Y$ 上拟射影概形 $X'$ 和固有
满态射 $X' \to X$。可以用支配 $X$ 的某个不可约分支的那些不可约分支的
不交并替换 $X'$；细节从略。特别地，可以假设概形 $X'$ 的泛点映到 $X$ 的
余维 $0$ 的点（在此情形下，这些恰是 $X$ 的泛点）。我们断言
$X' \to Y$ 固有。由《空间的态射》引理
\ref{spaces-morphisms-lemma-image-proper-is-proper}，该断言蕴含 $X$ 在 $Y$
上固有。为证明该断言，根据《极限》引理
\ref{limits-lemma-Noetherian-dvr-valuative-proper}，只需证明在每个交换实线图
$$
\xymatrix{
\Spec(K) \ar[r] \ar[d] & X' \ar[r] & X \ar[d] \\
\Spec(A) \ar[rr] \ar@{-->}[ru]^a \ar@{-->}[rru]_b & & Y
}
$$
中都能找到虚线箭头 $a$，其中 $A$ 是以 $K$ 为分式域的离散赋值环，且
$K$ 是 $X'$ 的某个泛点的剩余域（由于 $X'$ 分离，我们已经知道唯一性）。
由假设 (5)，可以找到虚线箭头 $b$。于是态射
$X' \times_{X, b} \Spec(A) \to \Spec(A)$ 是概形的固有态射；由概形态射的
赋值判据，可将 $b$ 提升为所需的态射 $a$。
\end{proof}

\begin{lemma}
\label{spaces-limits-lemma-check-universally-closed-Noetherian}
设 $S$ 为概形，$f : X \to Y$ 为 $S$ 上代数空间的态射。假设 $Y$ 局部诺特且
$f$ 有限型。则下列条件等价：
\begin{enumerate}
\item $f$ 普遍闭；
\item $f$ 满足赋值判据的存在性部分；
\item 存在概形 $V$ 和满平展态射 $V \to Y$，使得对所有 $n \geq 0$，
$|\mathbf{A}^n \times X \times_Y V| \to |\mathbf{A}^n \times V|$ 都是闭映射；
\item 对所有满足 $A$ 为离散赋值环的图 (\ref{spaces-limits-equation-valuative})，
% P09-ERR-0136
存在域的有限可分扩张 $K'/K$、支配 $A$ 的离散赋值环 $A' \subset K'$，
以及态射 $\Spec(A') \to X$，使下图交换：
$$
\xymatrix{
\Spec(K') \ar[r] \ar[d] & \Spec(K) \ar[r] & X \ar[d] \\
\Spec(A') \ar[r] \ar[rru] & \Spec(A) \ar[r] & Y
}
$$
\item 对所有满足 $A$ 为离散赋值环的图 (\ref{spaces-limits-equation-valuative})，
% P09-ERR-0137
存在域扩张 $K'/K$、支配 $A$ 的赋值环 $A' \subset K'$，以及态射
$\Spec(A') \to X$，使下图交换：
$$
\xymatrix{
\Spec(K') \ar[r] \ar[d] & \Spec(K) \ar[r] & X \ar[d] \\
\Spec(A') \ar[r] \ar[rru] & \Spec(A) \ar[r] & Y
}
$$
\end{enumerate}
\end{lemma}

\begin{proof}
由引理 \ref{spaces-limits-lemma-test-universally-closed} 和《空间的态射》引理
\ref{spaces-morphisms-lemma-quasi-compact-existence-universally-closed}，
(1)、(2) 与 (3) 等价。这些等价条件蕴含 (4)，因为《空间的态射》引理
\ref{spaces-morphisms-lemma-finite-separable-enough} 告诉我们，在赋值判据的
存在性部分，总可以选择有限可分的 $K'/K$；而由 Krull--Akizuki 定理
（《代数》引理 \ref{algebra-lemma-krull-akizuki}），这自动迫使 $A'$ 为
离散赋值环。(4) $\Rightarrow$ (5) 是立即的。在余下的证明中，
我们证明 (5) 蕴含 (1)。

\medskip\noindent
假设 (5)。% P09-ERR-0138
选择仿射概形 $V$ 和平展态射 $V \to Y$。只需证明 $f$ 到 $V$ 的基变换
普遍闭；参见《空间的态射》引理
\ref{spaces-morphisms-lemma-universally-closed-local}。完全如同引理
\ref{spaces-limits-lemma-Noetherian-dvr-valuative-proper} 的证明，可知假设 (5) 被此基变换
继承；细节从略。这样便归约到下一段所讨论的情形。

\medskip\noindent
假设 $Y$ 为诺特仿射概形且 (5) 成立。为证明 $f$ 普遍闭，只需证明对所有 $n$，
$|X \times \mathbf{A}^n| \to |Y \times \mathbf{A}^n|$ 都是闭映射
（由上述讨论）。由于假设 (5) 被乘积态射
$X \times \mathbf{A}^n \to Y \times \mathbf{A}^n$ 继承（细节从略），
问题归约为证明 $|X| \to |Y|$ 是闭映射。

\medskip\noindent
假设 $Y$ 为诺特仿射概形且 (5) 成立。设 $T \subset |X|$ 为闭子集。
必须证明 $T$ 在 $|Y|$ 中的像是闭集。可以用 $T$ 上的约化诱导闭子空间结构
替换 $X$；关于性质 (5) 被此替换保持的验证从略。因此归约为证明
$|X| \to |Y|$ 的像是闭集。

\medskip\noindent
设 $y \in |Y|$ 是 $|X| \to |Y|$ 的像的闭包中的点。由引理
\ref{spaces-limits-lemma-reach-point-closure-Noetherian}，可以选择交换图
$$
\xymatrix{
\Spec(K) \ar[r] \ar[d] & X \ar[d]^f \\
\Spec(A) \ar[r] & Y
}
$$
其中 $A$ 是离散赋值环，$K$ 是其分式域，并且 $\Spec(A)$ 的闭点映到 $y$。
由性质 (5) 立即可知，$y$ 位于 $|X| \to |Y|$ 的像中，证明完成。
\end{proof}


















\section{精细化的诺特赋值判据}
\label{spaces-limits-section-refined-valuative-criteria}

\noindent
本节是《极限》第 \ref{limits-section-refined-valuative-criteria} 节的对应版本。
在检验赋值判据时，通常不必考虑所有可能的赋值环图。

\begin{lemma}
\label{spaces-limits-lemma-refined-valuative-criterion-proper}
设 $S$ 为概形，$f : X \to Y$ 和 $h : U \to X$ 为 $S$ 上代数空间的态射。
假设 $Y$ 局部诺特，$f$ 与 $h$ 有限型，$f$ 分离，并且
$|h| : |U| \to |X|$ 的像在 $|X|$ 中稠密。若对任意交换实线图
$$
\xymatrix{
\Spec(K) \ar[r] \ar[d] & U \ar[r]^h & X \ar[d]^f \\
\Spec(A) \ar[rr] \ar@{-->}[rru] & & Y
}
$$
（其中 $A$ 是以 $K$ 为分式域的离散赋值环），都存在使图交换的虚线箭头，
则 $f$ 固有。
\end{lemma}

\begin{proof}
只需证明 $f$ 普遍闭。设 $V \to Y$ 是平展态射，其中 $V$ 是仿射概形。
由《空间的态射》引理
\ref{spaces-morphisms-lemma-universally-closed-local}，只需证明基变换
$X \times_Y V \to V$ 普遍闭。由《空间的性质》引理
\ref{spaces-properties-lemma-points-cartesian}，
$|U \times_Y V| \to |X \times_Y V|$ 的像 $I$ 是 $|h|$ 的像的逆像。
由于 $|X \times_Y V| \to |X|$ 是开映射（《空间的性质》引理
\ref{spaces-properties-lemma-etale-open}），可知 $I$ 在
$|X \times_Y V|$ 中稠密。因此，态射
$U \times_Y V \to X \times_Y V \to V$ 满足本引理的假设。故可以假设
$Y$ 是仿射概形。

\medskip\noindent
假设 $Y$ 是仿射概形。于是 $U$ 拟紧。% P09-ERR-0139
选择仿射概形以及满平展态射 $W \to U$。于是可以并且确实用 $W$ 替换 $U$，
从而假设 $U$ 仿射。由 Chow 引理的弱版本（《空间的上同调》引理
\ref{spaces-cohomology-lemma-weak-chow}），可以选择满固有态射 $X' \to X$，
其中 $X'$ 是概形。于是 $U' = X' \times_X U$ 是概形，且
$U' \to X'$ 有限型。可以用 $h' : U' \to X'$ 的概形论像替换 $X'$；
因而 $h'(U')$ 在 $X'$ 中稠密。我们断言，对于每个图
$$
\xymatrix{
% P09-ERR-0140
\Spec(K) \ar[r] \ar[d] & U' \ar[r]^h & X' \ar[d]^{f'} \\
\Spec(A) \ar[rr] \ar@{-->}[rru] & & Y
}
$$
（其中 $A$ 是以 $K$ 为分式域的离散赋值环），都存在使图交换的虚线箭头。
事实上，先由本引理的假设得到箭头 $\Spec(A) \to X$，再利用固有性的赋值判据
（《空间的态射》引理 \ref{spaces-morphisms-lemma-characterize-proper}）
将其提升为箭头 $\Spec(A) \to X'$。态射 $X' \to Y$ 是固有态射与分离态射
的复合，因而分离。因此由概形的情形，态射 $X' \to Y$ 固有（《极限》引理
\ref{limits-lemma-refined-valuative-criterion-proper}）。由《空间的态射》引理
\ref{spaces-morphisms-lemma-image-proper-is-proper}，可知 $X \to Y$ 固有。
\end{proof}

\begin{lemma}
\label{spaces-limits-lemma-refined-valuative-criterion-separated}
设 $S$ 为概形，$f : X \to Y$ 和 $h : U \to X$ 为 $S$ 上代数空间的态射。
假设 $Y$ 局部诺特，$f$ 局部有限型且拟分离，$h$ 有限型，并且
$|h| : |U| \to |X|$ 的像在 $|X|$ 中稠密。若对任意交换实线图
$$
\xymatrix{
\Spec(K) \ar[r] \ar[d] & U \ar[r]^h & X \ar[d]^f \\
\Spec(A) \ar[rr] \ar@{-->}[rru] & & Y
}
$$
（其中 $A$ 是以 $K$ 为分式域的离散赋值环），至多存在一条使图交换的
虚线箭头，则 $f$ 分离。
\end{lemma}

\begin{proof}
把引理 \ref{spaces-limits-lemma-refined-valuative-criterion-proper} 应用于态射
$U \to X$ 和 $\Delta : X \to X \times_Y X$。来检验其条件。注意，
由于 $f$ 拟分离，$\Delta$ 拟紧。当然，$\Delta$ 局部有限型且分离
（这对任意对角态射都成立）。最后，设给定交换实线图
$$
\xymatrix{
\Spec(K) \ar[r] \ar[d] & U \ar[r]^h & X \ar[d]^\Delta \\
\Spec(A) \ar[rr]^{(a, b)} \ar@{-->}[rru] & & X \times_Y X
}
$$
其中 $A$ 是以 $K$ 为分式域的离散赋值环。于是 $a$ 和 $b$ 给出本引理图中的
两条虚线箭头，因而必相等。因此可用 $a = b$ 作为虚线箭头，从而得到存在性。
这就完成了证明。
\end{proof}

\begin{lemma}
\label{spaces-limits-lemma-refined-valuative-criterion-universally-closed}
设 $S$ 为概形，$f : X \to Y$ 和 $h : U \to X$ 为 $S$ 上代数空间的态射。
假设 $Y$ 局部诺特，$f$ 与 $h$ 有限型，$f$ 拟分离，并且 $h(U)$ 在 $X$
中稠密。若对任意交换实线图
$$
\xymatrix{
\Spec(K) \ar[r] \ar[d] & U \ar[r]^h & X \ar[d]^f \\
\Spec(A) \ar[rr] \ar@{-->}[rru] & & Y
}
$$
（其中 $A$ 是以 $K$ 为分式域的离散赋值环），都存在唯一一条使图交换的
虚线箭头，则 $f$ 固有。
\end{lemma}

\begin{proof}
合并引理 \ref{spaces-limits-lemma-refined-valuative-criterion-separated} 与
\ref{spaces-limits-lemma-refined-valuative-criterion-proper} 即得。
\end{proof}







\section{下降有限型空间}
\label{spaces-limits-section-finite-type-quasi-separated}

\noindent
本节继续第 \ref{spaces-limits-section-finite-type-closed-in-finite-presentation} 节的主题，
并遵循第 \ref{spaces-limits-section-descending-relative} 节所讨论结果的思路。
它也是《极限》第 \ref{limits-section-finite-type-quasi-separated} 节关于
代数空间的对应版本。

\begin{situation}
\label{spaces-limits-situation-limit-noetherian}
设 $S$ 为概形，例如 $\Spec(\mathbf{Z})$。设
$B = \lim_{i \in I} B_i$ 是 $S$ 上诺特空间的有向逆系统的极限，
其转移态射 $B_{i'} \to B_i$（$i' \geq i$）均为仿射态射。
\end{situation}

\begin{lemma}
\label{spaces-limits-lemma-good-diagram}
在情形 \ref{spaces-limits-situation-limit-noetherian} 中，设 $X \to B$ 是代数空间的
拟分离有限型态射。则存在 $i \in I$ 及图
\begin{equation}
\label{spaces-limits-equation-good-diagram}
\vcenter{
\xymatrix{
X \ar[r] \ar[d] & W \ar[d] \\
B \ar[r] & B_i
}
}
\end{equation}
使得 $W \to B_i$ 有限型，并且诱导态射
$X \to B \times_{B_i} W$ 是闭浸入。
\end{lemma}

\begin{proof}
由引理 \ref{spaces-limits-lemma-finite-type-closed-in-finite-presentation}，可以找到
$B$ 上的闭浸入 $X \to X'$，其中 $X'$ 是 $B$ 上有限表示的代数空间。
由引理 \ref{spaces-limits-lemma-descend-finite-presentation}，可以找到 $i$ 和有限表示态射
$X'_i \to B_i$，其拉回为 $X'$。置 $W = X'_i$。
\end{proof}

\begin{lemma}
\label{spaces-limits-lemma-limit-from-good-diagram}
在情形 \ref{spaces-limits-situation-limit-noetherian} 中，设 $X \to B$ 是代数空间的
拟分离有限型态射。给定 $i \in I$ 及图
$$
\vcenter{
\xymatrix{
X \ar[r] \ar[d] & W \ar[d] \\
B \ar[r] & B_i
}
}
$$
如 (\ref{spaces-limits-equation-good-diagram})。对 $i' \geq i$，令 $X_{i'}$ 为
$X \to B_{i'} \times_{B_i} W$ 的概形论像。则
$X = \lim_{i' \geq i} X_{i'}$。
\end{lemma}

\begin{proof}
由于 $X$ 拟紧且拟分离，$X \to B_{i'} \times_{B_i} W$ 的概形论像的形成
与平展局部化交换（《空间的态射》引理
\ref{spaces-morphisms-lemma-quasi-compact-scheme-theoretic-image}）。
因此，可以并且确实假设 $W$ 仿射，并映入一个在 $B_i$ 上平展的仿射
$U_i$。于是
$$
B_{i'} \times_{B_i} W =
B_{i'} \times_{B_i} U_i \times_{U_i} W =
U_{i'} \times_{U_i} W
$$
其中 $U_{i'} = B_{i'} \times_{B_i} U_i$ 是仿射的，因为转移态射仿射。
故本引理归结为概形的情形，即《极限》引理
\ref{limits-lemma-limit-from-good-diagram}。
\end{proof}

\begin{lemma}
\label{spaces-limits-lemma-morphism-good-diagram}
在情形 \ref{spaces-limits-situation-limit-noetherian} 中，设 $f : X \to Y$ 是 $B$ 上
拟分离有限型代数空间的态射。设
$$
\vcenter{
\xymatrix{
X \ar[r] \ar[d] & W \ar[d] \\
B \ar[r] & B_{i_1}
}
}
\quad\text{且}\quad
\vcenter{
\xymatrix{
Y \ar[r] \ar[d] & V \ar[d] \\
B \ar[r] & B_{i_2}
}
}
$$
为形如 (\ref{spaces-limits-equation-good-diagram}) 的图。设
$X = \lim_{i \geq i_1} X_i$ 和 $Y = \lim_{i \geq i_2} Y_i$ 是引理
\ref{spaces-limits-lemma-limit-from-good-diagram} 中相应的极限描述。则存在
$i_0 \geq \max(i_1, i_2)$ 以及逆系统态射
$$
(f_i)_{i \geq i_0} : (X_i)_{i \geq i_0} \to (Y_i)_{i \geq i_0}
$$
它在 $(B_i)_{i \geq i_0}$ 上，并且 % P09-ERR-0141
$f = \lim_{i \geq i_0} f_i$。若
$(g_i)_{i \geq i_0} : (X_i)_{i \geq i_0} \to (Y_i)_{i \geq i_0}$
是在 $(B_i)_{i \geq i_0}$ 上的第二个逆系统态射，并且 % P09-ERR-0142
$f = \lim_{i \geq i_0} g_i$，则对所有 $i \gg i_0$ 都有 $f_i = g_i$。
\end{lemma}

\begin{proof}
由于 $V \to B_{i_2}$ 有限表示且 $X = \lim_{i \geq i_1} X_i$，可以应用命题
\ref{spaces-limits-proposition-characterize-locally-finite-presentation}（经引理
\ref{spaces-limits-lemma-better-characterize-relative-limit-preserving} 加强），找到
$i_0 \geq \max(i_1, i_2)$ 以及 $B_{i_2}$ 上的态射
$h : X_{i_0} \to V$，使得 $X \to X_{i_0} \to V$ 等于
$X \to Y \to V$。对 $i \geq i_0$，得到交换实线图
$$
\xymatrix{
X \ar[d] \ar[r] &
X_i \ar[r] \ar@{..>}[d] \ar@/_2pc/[dd] |!{[d];[ld]}\hole &
X_{i_0} \ar[d]^h \\
Y \ar[r] \ar[d] & Y_i \ar[r] \ar[d] & V \ar[d] \\
B \ar[r] & B_i \ar[r] & B_{i_0}
}
$$
由于 $X \to X_i$ 的像在概形论意义下稠密，而 $Y_i$ 是
$Y \to B_i \times_{B_{i_2}} V$ 的概形论像，可知由该图诱导的态射
$X_i \to B_i \times_{B_{i_2}} V$ 经由 $Y_i$ 分解（《空间的态射》引理
\ref{spaces-morphisms-lemma-factor-factor}）。这就证明了存在性。

\medskip\noindent
唯一性。对 $i \geq i_0$，令 $E_i \to X_i$ 为 $f_i$ 与 $g_i$ 的等化子。
有
$E_i = Y_i \times_{\Delta, Y_i \times_{B_i} Y_i, (f_i, g_i)} X_i$.
因此，作为 $Y_i$ 在 $B_i$ 上的对角态射的基变换，$E_i \to X_i$ 是有限表示
单态射；参见《空间的态射》诸引理
\ref{spaces-morphisms-lemma-properties-diagonal} 和
\ref{spaces-morphisms-lemma-diagonal-morphism-finite-type}。由于 $X_i$ 是
$B_i \times_{B_{i_0}} X_{i_0}$ 的闭子空间，而 $Y_i$ 也类似，可知
$$
E_i =
X_i \times_{(B_i \times_{B_{i_0}} X_{i_0})} (B_i \times_{B_{i_0}} E_{i_0}) =
X_i \times_{X_{i_0}} E_{i_0}
$$
类似地，有 $X = X \times_{X_{i_0}} E_{i_0}$。故由引理
\ref{spaces-limits-lemma-descend-isomorphism}，对充分大的 $i$ 有 $E_i = X_i$。
\end{proof}

\begin{remark}
\label{spaces-limits-remark-finite-type-gives-well-defined-system}
在情形 \ref{spaces-limits-situation-limit-noetherian} 中，引理 \ref{spaces-limits-lemma-good-diagram}、
\ref{spaces-limits-lemma-limit-from-good-diagram} 和 \ref{spaces-limits-lemma-morphism-good-diagram}
告诉我们：$B$ 上拟分离有限型代数空间的范畴等价于 $(B_i)_{i \in I}$ 上
某类代数空间逆系统，即把引理 \ref{spaces-limits-lemma-limit-from-good-diagram} 应用于
形如 (\ref{spaces-limits-equation-good-diagram}) 的图所产生的那些逆系统。例如，给定
有限型拟分离态射 $X \to B$，若选择两个形如
(\ref{spaces-limits-equation-good-diagram}) 的不同图 $X \to V_1 \to B_{i_1}$ 与
$X \to V_2 \to B_{i_2}$，则沿两个方向把引理
\ref{spaces-limits-lemma-morphism-good-diagram} 应用于 $\text{id}_X$，可见 $X$ 的相应极限
描述典范同构（至多需要缩小有向集 $I$）。其余依此类推。
\end{remark}

\begin{lemma}
\label{spaces-limits-lemma-morphism-good-diagram-flat}
记号与假设同引理 \ref{spaces-limits-lemma-morphism-good-diagram}。若 $f$ 平坦且有限表示，
则存在 $i_3 > i_0$，使得对 $i \geq i_3$，$f_i$ 平坦，
$X_i = Y_i \times_{Y_{i_3}} X_{i_3}$，且
$X = Y \times_{Y_{i_3}} X_{i_3}$。
\end{lemma}

\begin{proof}
由引理 \ref{spaces-limits-lemma-descend-finite-presentation}，可以选择 $i \geq i_2$ 以及
有限表示态射 $U \to Y_i$，使得 $X = Y \times_{Y_i} U$
（此处使用了 $f$ 有限表示）。增大 $i$ 后，可以假设 $U \to Y_i$ 平坦；
参见引理 \ref{spaces-limits-lemma-descend-flat}。如注
\ref{spaces-limits-remark-finite-type-gives-well-defined-system} 中所讨论，可以并且确实
用与 $X \to U \to B_i$ 对应的系统替换用来定义系统
$(X_i)_{i \geq i_1}$ 的初始图。因此，对 $i' \geq i$，$X_{i'}$ 被定义为
$X \to B_{i'} \times_{B_i} U$ 的概形论像。

\medskip\noindent
由于 $U \to Y_i$ 平坦（此处使用了 $f$ 平坦），又由于
$X = Y \times_{Y_i} U$，且 $Y \to Y_i$ 的概形论像为 $Y_i$，可知
$X \to U$ 的概形论像为 $U$（《空间的态射》引理
\ref{spaces-morphisms-lemma-flat-base-change-scheme-theoretic-image}）。
注意，由系统 $Y_j$ 的构造，对 $i' \geq i$，
$Y_{i'} \to B_{i'} \times_{B_i} Y_i$ 是闭浸入。于是与上述相同的论证表明，
$X \to B_{i'} \times_{B_i} U$ 的概形论像等于闭子空间
$Y_{i'} \times_{Y_i} U$。因此，对所有 $i' \geq i$，有
$X_{i'} = Y_{i'} \times_{Y_i} U$，故取 $i_3 = i$ 即得本引理。
\end{proof}

\begin{lemma}
\label{spaces-limits-lemma-morphism-good-diagram-smooth}
记号与假设同引理 \ref{spaces-limits-lemma-morphism-good-diagram}。若 $f$ 光滑，
则存在 $i_3 > i_0$，使得对 $i \geq i_3$，$f_i$ 光滑。
\end{lemma}

\begin{proof}
合并引理 \ref{spaces-limits-lemma-morphism-good-diagram-flat} 与
\ref{spaces-limits-lemma-descend-smooth} 即得。
\end{proof}

\begin{lemma}
\label{spaces-limits-lemma-morphism-good-diagram-proper}
记号与假设同引理 \ref{spaces-limits-lemma-morphism-good-diagram}。若 $f$ 固有，
则存在 $i_3 \geq i_0$，使得对 $i \geq i_3$，$f_i$ 固有。
\end{lemma}

\begin{proof}
由注 \ref{spaces-limits-remark-finite-type-gives-well-defined-system} 中的讨论，
选取适合形如 (\ref{spaces-limits-equation-good-diagram}) 的图的 $i_1$ 和 $W$，
不影响本引理的真假。因此如下选择 $W$。首先选择闭浸入 $X \to X'$，
使得 $X' \to Y$ 固有且有限表示；参见引理
\ref{spaces-limits-lemma-proper-limit-of-proper-finite-presentation}。然后选择
$i_3 \geq i_2$ 和固有态射 $W \to Y_{i_3}$，使得
$X' = Y \times_{Y_{i_3}} W$。这是可能的，因为
$Y = \lim_{i \geq i_2} Y_i$，以及 % P09-ERR-0143
引理 \ref{spaces-limits-lemma-relative-approximation} 和 \ref{spaces-limits-lemma-eventually-proper}。
按此选择 $W$ 后，由构造立即可知，对 $i \geq i_3$，代数空间 $X_i$ 是
$Y_i \times_{Y_{i_3}} W \subset B_i \times_{B_{i_3}} W$ 的闭子空间，
因而在 $Y_i$ 上固有。
\end{proof}

\begin{lemma}
\label{spaces-limits-lemma-good-diagram-fibre-product}
在情形 \ref{spaces-limits-situation-limit-noetherian} 中，设有笛卡尔图
$$
\xymatrix{
X^1 \ar[r]_p \ar[d]_q & X^3 \ar[d]^a \\
X^2 \ar[r]^b & X^4
}
$$
其中各代数空间在 $B$ 上拟分离且有限型。对每个 $j = 1, 2, 3, 4$，
选择 $i_j \in I$ 及图
$$
\xymatrix{
X^j \ar[r] \ar[d] & W^j \ar[d] \\
B \ar[r] & B_{i_j}
}
$$
形如 (\ref{spaces-limits-equation-good-diagram})。令
$X^j = \lim_{i \geq i_j} X^j_i$ 为引理 % P09-ERR-0144
\ref{spaces-limits-lemma-morphism-good-diagram} 中相应的极限描述。令
$(a_i)_{i \geq i_5}$、$(b_i)_{i \geq i_6}$、$(p_i)_{i \geq i_7}$ 和
$(q_i)_{i \geq i_8}$ 为引理 \ref{spaces-limits-lemma-morphism-good-diagram} 中构造的
相应逆系统态射。则存在 $i_9 \geq \max(i_5, i_6, i_7, i_8)$，使得对
$i \geq i_9$ 有 $a_i \circ p_i = b_i \circ q_i$，并且
$$
(q_i, p_i) : X^1_i \longrightarrow X^2_i \times_{b_i, X^4_i, a_i} X^3_i
$$
是闭浸入。若 $a$ 与 $b$ 平坦且有限表示，则存在
$i_{10} \geq \max(i_5, i_6, i_7, i_8, i_9)$，使得对 $i \geq i_{10}$，
最后显示的态射是同构。
\end{lemma}

\begin{proof}
由注 \ref{spaces-limits-remark-finite-type-gives-well-defined-system} 中的讨论，
选取适合形如 (\ref{spaces-limits-equation-good-diagram}) 的图的 $W^1$ 不影响本引理的真假。
因此可以选择 $W^1 = W^2 \times_{W^4} W^3$。于是由 $X^1_i$ 的构造立即可知，
$a_i \circ p_i = b_i \circ q_i$，并且
$$
(q_i, p_i) : X^1_i \longrightarrow X^2_i \times_{b_i, X^4_i, a_i} X^3_i
$$
是闭浸入。

\medskip\noindent
若 $a$ 与 $b$ 平坦且有限表示，则作为 $a$ 与 $b$ 的基变换，$p$ 与 $q$
也如此。因此，可将引理 \ref{spaces-limits-lemma-morphism-good-diagram-flat} 分别应用于
$a$、$b$、$p$、$q$ 以及 $a \circ p = b \circ q$。由此存在
$i_9 \in I$，使得 % P09-ERR-0145
$$
(q_i, p_i) : X^1_i \to X^2_i \times_{X^4_i} X^3_i
$$
对所有 $i \geq i_9$ 都是 $(q_{i_9}, p_{i_9})$ 沿态射 % P09-ERR-0146
$X^4_i \to X^4_{i_9}$ 的基变换。由引理 \ref{spaces-limits-lemma-descend-isomorphism}，
可知对所有充分大的 $i$，$(q_i, p_i)$ 是同构。
\end{proof}
% Generated mechanically from the frozen spaces-divisors.tex target.
% Do not edit this generated file; edit the bound locale source instead.

% OK, start here.
%

\setcounter{chapter}{70}
\chapter{代数空间上的除子}



\phantomsection
\label{spaces-divisors-section-phantom}


\section{引言}
\label{spaces-divisors-section-introduction}

\noindent
本章研究代数空间上的除子及相关主题。关于代数空间的基本参考文献是
\cite{Kn}。









\section{伴随点与弱伴随点}
\label{spaces-divisors-section-associated}

\noindent
在概形的情形，我们已经引入了两个相互竞争的伴随点概念，即通常的伴随点
（《除子》第 \ref{divisors-section-associated} 节）和弱伴随点
（《除子》第 \ref{divisors-section-weakly-associated} 节）。对于一般代数空间，
伴随点的概念基本无用，我们甚至不打算引入它。若代数空间局部诺特，则允许
用“伴随点”一词代替“弱伴随点”，因为对诺特概形两种概念相同
（《除子》引理 \ref{divisors-lemma-ass-weakly-ass}）。在给出定义之前，
先需要一个引理。

\begin{lemma}
\label{spaces-divisors-lemma-associated}
设 $S$ 为概形，$X$ 为 $S$ 上的代数空间，$\mathcal{F}$ 为拟凝聚
$\mathcal{O}_X$-模，且 $x \in |X|$。下列条件等价：
\begin{enumerate}
\item 对某个平展态射 $f : U \to X$（其中 $U$ 是概形）以及映到 $x$ 的
$u \in U$，点 $u$ 弱伴随于 $f^*\mathcal{F}$；
\item 对每个平展态射 $f : U \to X$（其中 $U$ 是概形）以及映到 $x$ 的
$u \in U$，点 $u$ 弱伴随于 $f^*\mathcal{F}$；
\item $\mathcal{O}_{X, \overline{x}}$ 的极大理想是茎
$\mathcal{F}_{\overline{x}}$ 的弱伴随素理想。
\end{enumerate}
若 $X$ 局部诺特，则这些条件还等价于：
\begin{enumerate}
\item[(4)] 对某个平展态射 $f : U \to X$（其中 $U$ 是概形）以及映到 $x$
的 $u \in U$，点 $u$ 伴随于 $f^*\mathcal{F}$；
\item[(5)] 对每个平展态射 $f : U \to X$（其中 $U$ 是概形）以及映到 $x$
的 $u \in U$，点 $u$ 伴随于 $f^*\mathcal{F}$；
\item[(6)] $\mathcal{O}_{X, \overline{x}}$ 的极大理想是茎
$\mathcal{F}_{\overline{x}}$ 的伴随素理想。
\end{enumerate}
\end{lemma}

\begin{proof}
选择带有点 $u$ 的概形 $U$ 以及把 $u$ 映到 $x$ 的平展态射
$f : U \to X$。把 $\overline{x}$ 提升为 $U$ 上位于 $u$ 上方的几何点。
回忆
$\mathcal{O}_{X, \overline{x}} = \mathcal{O}_{U, u}^{sh}$
，其中严格亨泽尔化是相对于所选的 $\overline{x}$ 的提升而言；参见
《空间的性质》引理
\ref{spaces-properties-lemma-describe-etale-local-ring}。最后，有
$$
\mathcal{F}_{\overline{x}} =
(f^*\mathcal{F})_u \otimes_{\mathcal{O}_{U, u}}
\mathcal{O}_{X, \overline{x}} =
(f^*\mathcal{F})_u \otimes_{\mathcal{O}_{U, u}}
\mathcal{O}_{U, u}^{sh}
$$
其中等式由《空间的性质》引理
\ref{spaces-properties-lemma-stalk-quasi-coherent} 得到。
因此，(1)、(2) 与 (3) 的等价性来自《平坦性的更多内容》引理
\ref{flat-lemma-weakly-associated-henselization}。若 $X$ 局部诺特，
则上述任意 $U$ 都局部诺特，故由《除子》引理
\ref{divisors-lemma-ass-weakly-ass}，(1)（相应地 (2)）等价于 (4)
（相应地 (5)）。另一方面，在局部诺特情形，局部环
$\mathcal{O}_{X, \overline{x}}$ 也诺特（《空间的性质》引理
\ref{spaces-properties-lemma-Noetherian-local-ring-Noetherian}）。
% P09-ERR-0147
因此，由同一引理（或由《代数》引理
\ref{algebra-lemma-ass-weakly-ass}）得到 (3) 与 (6) 的等价性。
\end{proof}

\begin{definition}
\label{spaces-divisors-definition-weakly-associated}
设 $S$ 为概形，$X$ 为 $S$ 上的代数空间，$\mathcal{F}$ 为 $X$ 上的
拟凝聚层，且 $x \in |X|$。
\begin{enumerate}
\item 若引理 \ref{spaces-divisors-lemma-associated} 的等价条件 (1)、(2)、(3) 成立，
则称 $x$ {\it 弱伴随于} $\mathcal{F}$。
\item % P09-ERR-0148
用 $\text{WeakAss}(\mathcal{F})$ 表示 $\mathcal{F}$ 的弱伴随点集合。
\item $X$ 的{\it 弱伴随点}是 $\mathcal{O}_X$ 的弱伴随点。
\end{enumerate}
若 $X$ 局部诺特，则当且仅当 $x$ 弱伴随于 $\mathcal{F}$ 时，称
{\it $x$ 伴随于 $\mathcal{F}$}，并置
$\text{Ass}(\mathcal{F}) = \text{WeakAss}(\mathcal{F})$。最后
（仍假设 $X$ 局部诺特），当且仅当 $x$ 是 $X$ 的弱伴随点时，称
{\it $x$ 是 $X$ 的伴随点}。
\end{definition}

\noindent
现在可以证明一些必备引理。

\begin{lemma}
\label{spaces-divisors-lemma-weakly-ass-support}
设 $S$ 为概形，$X$ 为 $S$ 上的代数空间，$\mathcal{F}$ 为拟凝聚
$\mathcal{O}_X$-模。则
$\text{WeakAss}(\mathcal{F}) \subset \text{Supp}(\mathcal{F})$。
\end{lemma}

\begin{proof}
这由定义立即得到。$X$ 上阿贝尔层的支撑定义见《空间的性质》定义
\ref{spaces-properties-definition-support}。
\end{proof}

\begin{lemma}
\label{spaces-divisors-lemma-ses-weakly-ass}
设 $S$ 为概形，$X$ 为 $S$ 上的代数空间。设
$0 \to \mathcal{F}_1 \to \mathcal{F}_2 \to \mathcal{F}_3 \to 0$
为 $X$ 上拟凝聚层的短正合列。则
$\text{WeakAss}(\mathcal{F}_2) \subset
\text{WeakAss}(\mathcal{F}_1) \cup \text{WeakAss}(\mathcal{F}_3)$
且
$\text{WeakAss}(\mathcal{F}_1) \subset \text{WeakAss}(\mathcal{F}_2)$。
\end{lemma}

\begin{proof}
对每个几何点 $\overline{x} \in X$，茎的序列
$0 \to \mathcal{F}_{1, \overline{x}} \to
\mathcal{F}_{2, \overline{x}} \to
\mathcal{F}_{3, \overline{x}} \to 0$
是 $\mathcal{O}_{X, \overline{x}}$-模的短正合列。因此本引理由
《代数》引理 \ref{algebra-lemma-weakly-ass} 得出。
\end{proof}

\begin{lemma}
\label{spaces-divisors-lemma-weakly-ass-zero}
设 $S$ 为概形，$X$ 为 $S$ 上的代数空间，$\mathcal{F}$ 为拟凝聚
$\mathcal{O}_X$-模。则
$$
\mathcal{F} = (0) \Leftrightarrow \text{WeakAss}(\mathcal{F}) = \emptyset
$$
\end{lemma}

\begin{proof}
取概形 $U$ 以及满平展态射 $f : U \to X$。则 $\mathcal{F}$ 为零当且仅当
$f^*\mathcal{F}$ 为零。因此，本引理由定义以及概形情形的对应引理得出；参见
《除子》引理 \ref{divisors-lemma-weakly-ass-zero}。
\end{proof}

\begin{lemma}
\label{spaces-divisors-lemma-minimal-support-in-weakly-ass}
设 $S$ 为概形，$X$ 为 $S$ 上的代数空间，$\mathcal{F}$ 为拟凝聚
$\mathcal{O}_X$-模，且 $x \in |X|$。若
\begin{enumerate}
\item $x \in \text{Supp}(\mathcal{F})$；
\item $x$ 是 $X$ 的余维 $0$ 点（《空间的性质》定义
\ref{spaces-properties-definition-dimension-local-ring}）。
\end{enumerate}
则 $x \in \text{WeakAss}(\mathcal{F})$。若 $\mathcal{F}$ 是有限型
$\mathcal{O}_X$-模，其概形论支撑为 $Z$（《空间的态射》定义
\ref{spaces-morphisms-definition-scheme-theoretic-support}），且 $x$ 是 $Z$
的余维 $0$ 点，则 $x \in \text{WeakAss}(\mathcal{F})$。
\end{lemma}

\begin{proof}
由于 $x \in \text{Supp}(\mathcal{F})$，茎
$\mathcal{F}_{\overline{x}}$ 非零。因此，由《代数》引理
\ref{algebra-lemma-weakly-ass-zero}，
$\text{WeakAss}(\mathcal{F}_{\overline{x}})$ 非空。另一方面，
$\mathcal{O}_{X, \overline{x}}$ 的谱仅含一个点。因此按定义，$x$ 是
$\mathcal{F}$ 的弱伴随点。最后一个断言成立是因为
% P09-ERR-0149
$\mathcal{O}_{X, \overline{x}} \to \mathcal{O}_{Z, \overline{z}}$
是满射，$\mathcal{O}_{Z, \overline{z}}$ 的谱仅含一个点，并且
$\mathcal{F}_{\overline{x}}$ 是非零的 $\mathcal{O}_{Z, \overline{z}}$-模。
\end{proof}

\begin{lemma}
\label{spaces-divisors-lemma-minimal-support-in-weakly-ass-decent}
设 $S$ 为概形，$X$ 为 $S$ 上的代数空间，$\mathcal{F}$ 为拟凝聚
$\mathcal{O}_X$-模，且 $x \in |X|$。若
\begin{enumerate}
\item $X$ 是良态的（例如拟分离或局部分离）；
\item $x \in \text{Supp}(\mathcal{F})$；
\item $x$ 不是 $\text{Supp}(\mathcal{F})$ 中另一个点的特化。
\end{enumerate}
则 $x \in \text{WeakAss}(\mathcal{F})$。
\end{lemma}

\begin{proof}
（拟分离代数空间是良态的，参见《良态空间》第
\ref{decent-spaces-section-reasonable-decent} 节；局部分离代数空间是良态的，
参见《良态空间》引理 \ref{decent-spaces-lemma-locally-separated-decent}。）
取概形 $U$、点 $u \in U$ 以及把 $u$ 映到 $x$ 的平展态射
$f : U \to X$。由《良态空间》引理
\ref{decent-spaces-lemma-decent-no-specializations-map-to-same-point}，
若 $u' \leadsto u$ 是非平凡特化，则 $f(u') \not = x$。由此可见，
$u \in \text{Supp}(f^*\mathcal{F})$ 不是
$\text{Supp}(f^*\mathcal{F})$ 中另一个点的特化。因此，由
% P09-ERR-0150
《除子》引理 \ref{spaces-divisors-lemma-minimal-support-in-weakly-ass}，
$u \in \text{WeakAss}(f^*\mathcal{F})$。
\end{proof}

\begin{lemma}
\label{spaces-divisors-lemma-finite-ass}
设 $S$ 为概形，$X$ 为 $S$ 上的局部诺特代数空间，$\mathcal{F}$ 为凝聚
$\mathcal{O}_X$-模。则对每个拟紧开子集 $W \subset |X|$，
$\text{Ass}(\mathcal{F}) \cap W$ 是有限集。
\end{lemma}

\begin{proof}
取拟紧概形 $U$ 以及平展态射 $U \to X$，使得 $W$ 是
$|U| \to |X|$ 的像。于是 $U$ 是诺特概形，应用《除子》引理
\ref{divisors-lemma-finite-ass} 即得结论。
\end{proof}

\begin{lemma}
\label{spaces-divisors-lemma-restriction-injective-open-contains-weakly-ass}
设 $S$ 为概形，$X$ 为 $S$ 上的代数空间，$\mathcal{F}$ 为拟凝聚
$\mathcal{O}_X$-模。若 $U \to X$ 是平展态射且
$\text{WeakAss}(\mathcal{F}) \subset \Im(|U| \to |X|)$，则
$\Gamma(X, \mathcal{F}) \to \Gamma(U, \mathcal{F})$ 是单射。
\end{lemma}

\begin{proof}
设 $s \in \Gamma(X, \mathcal{F})$ 是限制到 $U$ 上为零的截面。令
$\mathcal{F}' \subset \mathcal{F}$ 为 $s$ 所定义的映射
$\mathcal{O}_X \to \mathcal{F}$ 的像。则 $\mathcal{F}'|_U = 0$。这蕴含
$\text{WeakAss}(\mathcal{F}') \cap \Im(|U| \to |X|) = \emptyset$
（按弱伴随点的定义）。另一方面，由引理 \ref{spaces-divisors-lemma-ses-weakly-ass}，
$\text{WeakAss}(\mathcal{F}') \subset \text{WeakAss}(\mathcal{F})$
。故 $\text{WeakAss}(\mathcal{F}') = \emptyset$。于是由引理
\ref{spaces-divisors-lemma-weakly-ass-zero}，$\mathcal{F}' = 0$。
\end{proof}

\begin{lemma}
\label{spaces-divisors-lemma-weakass-pushforward}
设 $S$ 为概形，$f : X \to Y$ 为 $S$ 上代数空间间的拟紧拟分离态射，
$\mathcal{F}$ 为拟凝聚 $\mathcal{O}_X$-模。设 $y \in |Y|$ 不在
$|f|$ 的像中。则 $y$ 不弱伴随于 $f_*\mathcal{F}$。
\end{lemma}

\begin{proof}
由《空间的态射》引理 \ref{spaces-morphisms-lemma-pushforward}，
$\mathcal{O}_Y$-模 $f_*\mathcal{F}$ 拟凝聚，故本引理的陈述有意义。
取仿射概形 $V$、点 $v \in V$ 以及把 $v$ 映到 $y$ 的平展态射
$V \to Y$。可将 $f : X \to Y$、$\mathcal{F}$、$y$ 分别替换为
$X \times_Y V \to V$、$\mathcal{F}|_{X \times_Y V}$、$v$。
因此可设 $Y$ 为仿射概形。此时 $X$ 拟紧，故可取仿射概形 $U$ 以及满平展态射
$U \to X$。
% P09-ERR-0151
用 $g : U \to Y$ 表示合成态射。于是
$f_*\mathcal{F} \subset g_*(\mathcal{F}|_U)$。
由引理 \ref{spaces-divisors-lemma-ses-weakly-ass}，可归约到概形的情形，即《除子》引理
\ref{divisors-lemma-weakass-pushforward}。
\end{proof}

\begin{lemma}
\label{spaces-divisors-lemma-check-injective-on-weakass}
设 $S$ 为概形，$X$ 为 $S$ 上的代数空间。设
$\varphi : \mathcal{F} \to \mathcal{G}$ 为拟凝聚 $\mathcal{O}_X$-模间的
映射。假设对每个 $x \in |X|$，下列情形至少有一个发生：
\begin{enumerate}
\item $\mathcal{F}_{\overline{x}} \to \mathcal{G}_{\overline{x}}$
是单射；或
\item $x \not \in \text{WeakAss}(\mathcal{F})$。
\end{enumerate}
则 $\varphi$ 是单射。
\end{lemma}

\begin{proof}
这些假设蕴含 $\text{WeakAss}(\Ker(\varphi)) = \emptyset$，故由引理
\ref{spaces-divisors-lemma-weakly-ass-zero}，$\Ker(\varphi) = 0$。
\end{proof}

\begin{lemma}
\label{spaces-divisors-lemma-weakass-reduced}
设 $S$ 为概形，$X$ 为 $S$ 上的既约代数空间。则
% P09-ERR-0152
$X$ 的弱伴随点恰为 $X$ 的余维 $0$ 点。
\end{lemma}

\begin{proof}
在平展局部工作，本引理由《除子》引理
\ref{divisors-lemma-weakass-reduced} 及《空间的性质》引理
\ref{spaces-properties-lemma-codimension-0-points} 得出。
\end{proof}










\section{态射与弱伴随点}
\label{spaces-divisors-section-morphisms-weakly-associated}

\begin{lemma}
\label{spaces-divisors-lemma-weakly-ass-reverse-functorial}
设 $S$ 为概形，$f : X \to Y$ 为 $S$ 上代数空间间的仿射态射，
$\mathcal{F}$ 为拟凝聚 $\mathcal{O}_X$-模。则有
% P09-ERR-0153
$$
\text{WeakAss}_S(f_*\mathcal{F}) \subset f(\text{WeakAss}_X(\mathcal{F}))
$$
\end{lemma}

\begin{proof}
取概形 $V$ 以及满平展态射 $V \to Y$。置 $U = X \times_Y V$。则
$U \to V$ 是概形间的仿射态射。由弱伴随点的定义，问题归约为概形间的态射
$U \to V$；这种情形已由《除子》引理
\ref{divisors-lemma-weakly-ass-reverse-functorial} 处理。
\end{proof}

\begin{lemma}
\label{spaces-divisors-lemma-ass-functorial-equal}
设 $S$ 为概形，$f : X \to Y$ 为 $S$ 上代数空间间的仿射态射，
$\mathcal{F}$ 为拟凝聚 $\mathcal{O}_X$-模。若 $X$ 局部诺特，则有
$$
\text{WeakAss}_Y(f_*\mathcal{F}) =
f(\text{WeakAss}_X(\mathcal{F}))
$$
\end{lemma}

\begin{proof}
取概形 $V$ 以及满平展态射 $V \to Y$。置 $U = X \times_Y V$。则
$U \to V$ 是概形间的仿射态射，且 $U$ 局部诺特。由弱伴随点的定义，
问题归约为概形间的态射 $U \to V$；这种情形已由《除子》引理
\ref{divisors-lemma-ass-functorial-equal} 处理。
\end{proof}

\begin{lemma}
\label{spaces-divisors-lemma-weakly-associated-finite}
设 $S$ 为概形，$f : X \to Y$ 为 $S$ 上代数空间间的有限态射，
$\mathcal{F}$ 为拟凝聚 $\mathcal{O}_X$-模。则
$\text{WeakAss}(f_*\mathcal{F}) = f(\text{WeakAss}(\mathcal{F}))$。
\end{lemma}

\begin{proof}
取概形 $V$ 以及满平展态射 $V \to Y$。置 $U = X \times_Y V$。则
$U \to V$ 是概形间的有限态射。由弱伴随点的定义，问题归约为概形间的态射
$U \to V$；这种情形已由《除子》引理
\ref{divisors-lemma-weakly-associated-finite} 处理。
\end{proof}

\begin{lemma}
\label{spaces-divisors-lemma-weakly-ass-pullback}
设 $S$ 为概形，$f : X \to Y$ 为 $S$ 上代数空间间的态射，
$\mathcal{G}$ 为拟凝聚 $\mathcal{O}_Y$-模。设 $x \in |X|$ 且
$y = f(x) \in |Y|$。若
\begin{enumerate}
\item % P09-ERR-0154
$y \in \text{WeakAss}_S(\mathcal{G})$；
\item $f$ 在 $x$ 处平坦；且
\item $f$ 的纤维在 $x$ 处的局部环维数为零（《空间的态射》定义
\ref{spaces-morphisms-definition-dimension-fibre}）；
\end{enumerate}
则 $x \in \text{WeakAss}(f^*\mathcal{G})$。
\end{lemma}

\begin{proof}
取概形 $V$、点 $v \in V$ 以及把 $v$ 映到 $y$ 的平展态射 $V \to Y$。
取概形 $U$、点 $u \in U$ 以及平展态射
% P09-ERR-0155
$U \to V \times_Y X$，它把 $v$ 映到位于 $v$ 与 $x$ 上方的一个点。
这是可行的，因为由《空间的性质》引理
% P09-ERR-0156
\ref{spaces-properties-lemma-points-cartesian}，存在
$t \in |V \times_Y X|$ 映到 $(v, y)$。按定义可见，
$\mathcal{O}_{U_v, u}$ 的维数为零。因此 $u$ 是纤维 $U_v$ 的一个一般点。
由弱伴随点的定义，问题归约为概形间的态射 $U \to V$。这种情形已由
《除子》引理 \ref{divisors-lemma-weakly-ass-pullback} 处理。
\end{proof}

\begin{lemma}
\label{spaces-divisors-lemma-weakly-ass-change-fields}
设 $K/k$ 为域扩张，$X$ 为 $k$ 上的代数空间，$\mathcal{F}$ 为拟凝聚
$\mathcal{O}_X$-模。设 $y \in X_K$ 的像为 $x \in X$。若 $y$ 是拉回
$\mathcal{F}_K$ 的弱伴随点，则 $x$ 是 $\mathcal{F}$ 的弱伴随点。
\end{lemma}

\begin{proof}
这是将《除子》引理 \ref{divisors-lemma-weakly-ass-change-fields}
改述为代数空间的语言所得的陈述。略去这种改述的细节。
\end{proof}

\begin{lemma}
\label{spaces-divisors-lemma-finite-flat-weak-assassin-up-down}
设 $S$ 为概形，$f : X \to Y$ 为代数空间间的有限平坦态射，
$\mathcal{G}$ 为拟凝聚 $\mathcal{O}_Y$-模。设 $x \in |X|$ 的像为
$y \in |Y|$。则
% P09-ERR-0157
$$
x \in \text{WeakAss}(g^*\mathcal{G})
\Leftrightarrow
y \in \text{WeakAss}(\mathcal{G})
$$
\end{lemma}

\begin{proof}
由概形的情形（《平坦性的更多内容》引理
\ref{flat-lemma-finite-flat-weak-assassin-up-down}）经平展局部化立即得出。
\end{proof}

\begin{lemma}
\label{spaces-divisors-lemma-etale-weak-assassin-up-down}
设 $S$ 为概形，$f : X \to Y$ 为代数空间间的平展态射，
$\mathcal{G}$ 为拟凝聚 $\mathcal{O}_Y$-模。设 $x \in |X|$ 的像为
$y \in |Y|$。则
$$
x \in \text{WeakAss}(f^*\mathcal{G})
\Leftrightarrow
y \in \text{WeakAss}(\mathcal{G})
$$
\end{lemma}

\begin{proof}
这由弱伴随点的定义立即得到；事实上，概形情形的对应引理
（《平坦性的更多内容》引理
\ref{flat-lemma-etale-weak-assassin-up-down}）正是本定义的基础。
\end{proof}







\section{相对弱伴随点集}
\label{spaces-divisors-section-relative-weak-assassin}

\noindent
为了定义这个对象，需要先证明几个引理。

\begin{lemma}
\label{spaces-divisors-lemma-locally-noetherian-fibre}
设 $S$ 为概形，$f : X \to Y$ 为 $S$ 上代数空间间的态射，且
$y \in |Y|$。下列条件等价：
\begin{enumerate}
\item 对某个概形 $V$、点 $v \in V$ 以及把 $v$ 映到 $y$ 的平展态射
$V \to Y$，代数空间 $X_v$ 局部诺特；
\item 对每个概形 $V$、点 $v \in V$ 以及把 $v$ 映到 $y$ 的平展态射
$V \to Y$，代数空间 $X_v$ 局部诺特；且
\item 存在域 $k$ 以及表示 $y$ 的态射 $\Spec(k) \to Y$，使得 $X_k$
局部诺特。
\end{enumerate}
若存在域 $k_0$ 以及表示 $y$ 的单态射 $\Spec(k_0) \to Y$，则上述条件还
等价于：
\begin{enumerate}
\item[(4)] 代数空间 $X_{k_0}$ 局部诺特。
\end{enumerate}
\end{lemma}

\begin{proof}
注意 $X_v = v \times_Y X = \Spec(\kappa(v)) \times_Y X$。因此蕴含
(2) $\Rightarrow$ (1) $\Rightarrow$ (3) 是显然的。假设从某个域的谱出发的
态射 $\Spec(k) \to Y$ 使得 $X_k$ 局部诺特。设 $V \to Y$ 是从概形
$V$ 出发的平展态射，并且
% P09-ERR-0158
$v \in V$ 是映到 $y$ 的点。则概形 $v \times_Y \Spec(k)$ 非空。取一点
$w \in v \times_Y \Spec(k)$。考虑态射
$$
X_v \longleftarrow X_w \longrightarrow X_k
$$
由于 $V \to Y$ 平展，且可将 $w$ 视为 $V \times_Y \Spec(k)$ 的一个点，
故 $\kappa(w)/k$ 是有限可分域扩张（《态射》引理
\ref{morphisms-lemma-etale-over-field}）。因此，作为
$w \to \Spec(k)$ 的基变换，$X_w \to X_k$ 是有限平展态射。故
$X_w$ 局部诺特（《空间的态射》引理
\ref{spaces-morphisms-lemma-locally-finite-type-locally-noetherian}）。
态射 $X_w \to X_v$ 是满、仿射、平坦的，因为它是满、仿射、平坦态射
$w \to v$ 的基变换。于是 $X_w$ 局部诺特蕴含 $X_v$ 局部诺特。为看出
这一点，取满平展态射 $U \to X$，再利用 $U_w \to U_v$ 满、仿射且平坦。
在概形 $U_v$ 上仿射局部地工作，由《代数》引理
% P09-ERR-0159
\ref{algebra-lemma-descent-Noetherian} 可得 $U_w$ 局部诺特。

\medskip\noindent
最后，只需在有表示 $y$ 的单态射 $\Spec(k_0) \to Y$ 时证明 (3) 蕴含
(4)。此时 $\Spec(k) \to Y$ 分解为
$\Spec(k) \to \Spec(k_0) \to Y$。上述论证于是表明，$X_k$ 局部诺特
% P09-ERR-0160
蕴含 $X_{k_0}$ 局部诺特。
\end{proof}

\begin{definition}
\label{spaces-divisors-definition-locally-Noetherian-fibre}
设 $S$ 为概形，$f : X \to Y$ 为 $S$ 上代数空间间的态射，且
$y \in |Y|$。若引理 \ref{spaces-divisors-lemma-locally-noetherian-fibre} 的等价条件
(1)、(2)、(3) 成立，则称 {\it $f$ 在 $y$ 上的纤维局部诺特}。若这对每个
$y \in |Y|$ 都成立，则称 {\it $f$ 的各纤维局部诺特}。
\end{definition}

\noindent
当然，保证纤维局部诺特的通常方法是假设该态射局部有限型。

\begin{lemma}
\label{spaces-divisors-lemma-locally-finite-type-locally-Noetherian-fibres}
设 $S$ 为概形，$f : X \to Y$ 为 $S$ 上代数空间间的态射。若 $f$
局部有限型，则 $f$ 的各纤维局部诺特。
\end{lemma}

\begin{proof}
这由《空间的态射》引理
\ref{spaces-morphisms-lemma-locally-finite-type-locally-noetherian}
以及域的谱为诺特概形这一事实得出。
\end{proof}

\begin{lemma}
\label{spaces-divisors-lemma-relative-assassin}
设 $S$ 为概形，$f : X \to Y$ 为 $S$ 上代数空间间的态射。设
$x \in |X|$ 且 $y = f(x) \in |Y|$，并设 $\mathcal{F}$ 为拟凝聚
$\mathcal{O}_X$-模。考虑交换图
$$
\xymatrix{
X \ar[d] & X \times_Y V \ar[d] \ar[l] & X_v \ar[d] \ar[l] \\
Y & V \ar[l] & v \ar[l]
}
\quad
\xymatrix{
X \ar[d] & U \ar[d] \ar[l] & U_v \ar[d] \ar[l] \\
Y & V \ar[l] & v \ar[l]
}
\quad
\xymatrix{
x \ar@{|->}[d] &
x' \ar@{|->}[d] \ar@{|->}[l] &
u \ar@{|->}[ld] \ar@{|->}[l] \\
y &
v \ar@{|->}[l]
}
$$
其中 $V$ 和 $U$ 是概形，$V \to Y$ 与 $U \to X \times_Y V$ 平展，而
$v \in V$、$x' \in |X_v|$、$u \in U$ 是按最后一个图相关联的点。
% P09-ERR-0161
用 $\mathcal{F}|_{X_v}$ 和 $\mathcal{F}|_{U_v}$ 表示 $\mathcal{F}$
的拉回。下列条件等价：
\begin{enumerate}
\item 对某组如上的 $V, v, x'$，点 $x'$ 是
$\mathcal{F}|_{X_v}$ 的弱伴随点；
\item 对每组如上的 $V \to Y, v, x'$，点 $x'$ 是
$\mathcal{F}|_{X_v}$ 的弱伴随点；
\item 对某组如上的 $U, V, u, v$，点 $u$ 是
$\mathcal{F}|_{U_v}$ 的弱伴随点；
\item 对每组如上的 $U, V, u, v$，点 $u$ 是
$\mathcal{F}|_{U_v}$ 的弱伴随点；
\item 对某个域 $k$、某个表示 $y$ 的态射 $\Spec(k) \to Y$，以及某个映到
$x$ 的 $t \in |X_k|$，点 $t$ 是 $\mathcal{F}|_{X_k}$ 的弱伴随点。
\end{enumerate}
若存在域 $k_0$ 以及表示 $y$ 的单态射 $\Spec(k_0) \to Y$，则上述条件还
等价于：
\begin{enumerate}
\item[(6)] $x_0$ 是 $\mathcal{F}|_{X_{k_0}}$ 的弱伴随点，其中
$x_0 \in |X_{k_0}|$ 是映到 $x$ 的唯一点。
\end{enumerate}
若 $f$ 在 $y$ 上的纤维局部诺特，则在条件 (1)、(2)、(3)、(4)、(6) 中，
可以用“伴随”代替“弱伴随”。
\end{lemma}

\begin{proof}
注意，给定引理中如上的 $V, v, x'$，可以找到
$U \to X \times_Y V$ 和映到 $x'$ 的 $u \in U$，于是态射
$U_v \to X_v$ 平展。因此，(1) 与 (3) 等价以及 (2) 与 (4) 等价都是
显然的；其中每个条件都蕴含 (5)。我们将证明 (5) 蕴含 (2)。设给定
$V, v, x'$，以及
% P09-ERR-0162
$\Spec(k) \to X$ 和 $t \in |X_k|$，使得点 $t$ 是
$\mathcal{F}|_{X_k}$ 的弱伴随点。可以取一点
$w \in v \times_Y \Spec(k)$。于是得到态射
$$
X_v \longleftarrow X_w \longrightarrow X_k
$$
由于 $V \to Y$ 平展，且可将 $w$ 视为 $V \times_Y \Spec(k)$ 的一点，
故 $\kappa(w)/k$ 是有限可分域扩张（《态射》引理
\ref{morphisms-lemma-etale-over-field}）。因此，作为
$w \to \Spec(k)$ 的基变换，$X_w \to X_k$ 是有限平展态射。于是由引理
\ref{spaces-divisors-lemma-etale-weak-assassin-up-down}，$X_w$ 中位于 $t$ 上方的任一点
$x''$ 都是 $\mathcal{F}|_{X_w}$ 的弱伴随点。可以选取映到 $x'$ 的
$x''$（《空间的性质》引理
\ref{spaces-properties-lemma-points-cartesian}）。于是引理
\ref{spaces-divisors-lemma-weakly-ass-change-fields} 蕴含 $x'$ 是
$\mathcal{F}|_{X_v}$ 的弱伴随点。

\medskip\noindent
为完成证明，只需说明：若给定 (6) 中的 $\Spec(k_0) \to Y$，则等价条件
(1)--(5) 蕴含 (6)。此时，(5) 中的态射 $\Spec(k) \to Y$ 唯一分解为
$\Spec(k) \to \Spec(k_0) \to Y$。于是 $x_0$ 是 $t$ 在态射
$X_k \to X_{k_0}$ 下的像。因此，上述同一引理表明 (6) 成立。
\end{proof}

\begin{definition}
\label{spaces-divisors-definition-relative-weak-assassin}
设 $S$ 为概形，$f : X \to Y$ 为 $S$ 上代数空间间的态射，
$\mathcal{F}$ 为拟凝聚 $\mathcal{O}_X$-模。{\it $\mathcal{F}$ 在 $X$
中相对于 $Y$ 的相对弱伴随点集}是集合
$\text{WeakAss}_{X/Y}(\mathcal{F}) \subset |X|$，它由满足引理
\ref{spaces-divisors-lemma-relative-assassin} 中等价条件的那些 $x \in |X|$ 组成。
若 $f$ 的各纤维局部诺特（定义
\ref{spaces-divisors-definition-locally-Noetherian-fibre}），则使用记号
$\text{Ass}_{X/Y}(\mathcal{F})$。
\end{definition}

\noindent
用这个记号，可以表述若干已对概形证明的结果。

\begin{lemma}
\label{spaces-divisors-lemma-bourbaki}
设 $S$ 为概形，$f : X \to Y$ 为 $S$ 上代数空间间的态射，
$\mathcal{F}$ 为拟凝聚 $\mathcal{O}_X$-模，而 $\mathcal{G}$ 为拟凝聚
$\mathcal{O}_Y$-模。假设
\begin{enumerate}
\item $\mathcal{F}$ 在 $Y$ 上平坦；
\item $X$ 和 $Y$ 局部诺特；且
\item $f$ 的各纤维局部诺特。
\end{enumerate}
则
$$
\text{Ass}_X(\mathcal{F} \otimes_{\mathcal{O}_X} f^*\mathcal{G}) =
\{x \in \text{Ass}_{X/Y}(\mathcal{F})\text{ 使得 }
f(x) \in \text{Ass}_Y(\mathcal{G}) \}
$$
\end{lemma}

\begin{proof}
经平展局部化，这立即由概形情形的结果得出；参见《除子》引理
\ref{divisors-lemma-bourbaki}。概形情形的结果更一般，仅仅是因为我们没有
为非诺特代数空间定义伴随点（因此，为了使本结果甚至能够表述，必须假设
$X$ 以及 $X \to Y$ 的各纤维局部诺特）。
\end{proof}

\begin{lemma}
\label{spaces-divisors-lemma-base-change-relative-assassin}
设 $S$ 为概形。设
$$
\xymatrix{
X' \ar[d]_{f'} \ar[r]_{g'} & X \ar[d]^f \\
Y' \ar[r]^g & Y
}
$$
为 $S$ 上代数空间的笛卡尔图。设 $\mathcal{F}$ 为拟凝聚
$\mathcal{O}_X$-模，并置 $\mathcal{F}' = (g')^*\mathcal{F}$。若 $f$
局部有限型，则
\begin{enumerate}
\item $x' \in \text{Ass}_{X'/Y'}(\mathcal{F}')
\Rightarrow g'(x') \in \text{Ass}_{X/Y}(\mathcal{F})$；
\item 若 $x \in \text{Ass}_{X/Y}(\mathcal{F})$，则给定满足
$f(x) = g(y')$ 的 $y' \in |Y'|$，存在
$x' \in \text{Ass}_{X'/Y'}(\mathcal{F}')$，使得
$g'(x') = x$ 且 $f'(x') = y'$。
\end{enumerate}
\end{lemma}

\begin{proof}
经平展局部化，这由概形的情形得出。我们完整写出细节。取概形 $V$ 以及满平展
态射 $V \to Y$。取概形 $U$ 以及满平展态射
$U \to V \times_Y X$。取概形 $V'$ 以及满平展态射
$V' \to V \times_Y Y'$。则 $U' = V' \times_V U$ 是概形，并且态射
$U' \to X'$ 满且平展。

\medskip\noindent
(1) 的证明。取映到 $x'$ 的 $u' \in U'$。
% P09-ERR-0163
用 $v' \in V'$ 表示 $u'$ 的像。按定义，
$x' \in \text{Ass}_{X'/Y'}(\mathcal{F}')$ 等价于
$u' \in \text{Ass}(\mathcal{F}|_{U'_{v'}})$（因为 $U'_{v'}$ 局部诺特，
所以写 $\text{Ass}$ 而不是 $\text{WeakAss}$ 是有意义的）。应用《除子》
引理 \ref{divisors-lemma-base-change-relative-assassin} 可见，$u'$ 的像
$u \in U$ 属于 $\text{Ass}(\mathcal{F}|_{U_v})$，其中
$v \in V$ 是 $u$ 的像。这又意味着
$g'(x') \in \text{Ass}_{X/Y}(\mathcal{F})$。

\medskip\noindent
(2) 的证明。取映到 $x$ 的 $u \in U$。
% P09-ERR-0164
用 $v \in V$ 表示 $u$ 的像。按定义，
$x \in \text{Ass}_{X/Y}(\mathcal{F})$ 等价于
$u \in \text{Ass}(\mathcal{F}|_{U_v})$。取一点 $v' \in V'$，它映到
$y' \in |Y'|$ 及 $v \in V$（这由《空间的性质》引理
\ref{spaces-properties-lemma-points-cartesian} 可行）。令
$t \in \Spec(\kappa(v') \otimes_{\kappa(v)} \kappa(u))$
为某个不可约分支的一般点，并令 $u' \in U'$ 为 $t$ 的像。应用《除子》
引理 \ref{divisors-lemma-base-change-relative-assassin} 可见
$u' \in \text{Ass}(\mathcal{F}'|_{U'_{v'}})$。这又意味着
$x' \in \text{Ass}_{X'/Y'}(\mathcal{F}')$，其中 $x' \in |X'|$
是 $u'$ 的像。
\end{proof}

\begin{lemma}
\label{spaces-divisors-lemma-base-change-relative-assassin-quasi-finite}
采用引理 \ref{spaces-divisors-lemma-base-change-relative-assassin} 的记号与假设。假设 $g$
局部拟有限，或更一般地，假设对每个 $y' \in |Y'|$，$y'/g(y')$ 的超越次数
为 $0$。则 $\text{Ass}_{X'/Y'}(\mathcal{F}')$ 是
$\text{Ass}_{X/Y}(\mathcal{F})$ 的逆像。
\end{lemma}

\begin{proof}
点在其像上的超越次数定义见《空间的态射》定义
\ref{spaces-morphisms-definition-dimension-fibre}。设 $x' \in |X'|$
的像为 $x \in |X|$。取概形 $V$ 以及满平展态射 $V \to Y$。取概形
$U$ 以及满平展态射 $U \to V \times_Y X$。取概形 $V'$ 以及满平展态射
$V' \to V \times_Y Y'$。则 $U' = V' \times_V U$ 是概形，并且态射
$U' \to X'$ 满且平展。取映到 $x$ 的 $u \in U$。
% P09-ERR-0165
用 $v \in V$ 表示 $u$ 的像。按定义，
$x \in \text{Ass}_{X/Y}(\mathcal{F})$ 等价于
$u \in \text{Ass}(\mathcal{F}|_{U_v})$。取一点 $u' \in U'$，它映到
$x' \in |X'|$ 及 $u \in U$（这由《空间的性质》引理
\ref{spaces-properties-lemma-points-cartesian} 可行）。令 $v' \in V'$
为 $u'$ 的像。按定义，$x' \in \text{Ass}_{X'/Y'}(\mathcal{F}')$
等价于 $u' \in \text{Ass}(\mathcal{F}'|_{U'_{v'}})$。现在，将《除子》
注记 \ref{divisors-remark-base-change-relative-assassin} 中的讨论应用于
$u' \in \Spec(\kappa(v') \otimes_{\kappa(v)} \kappa(u))$，即得本引理。
\end{proof}

\begin{lemma}
\label{spaces-divisors-lemma-relative-weak-assassin-finite}
设 $S$ 为概形，$f : X \to Y$ 为 $S$ 上代数空间间的态射，
$i : Z \to X$ 为有限态射，而 $\mathcal{G}$ 为拟凝聚
$\mathcal{O}_Z$-模。则
$\text{WeakAss}_{X/Y}(i_*\mathcal{G}) =
i(\text{WeakAss}_{Z/Y}(\mathcal{G}))$。
\end{lemma}

\begin{proof}
经平展局部化，这由概形情形（《除子》引理
\ref{divisors-lemma-relative-weak-assassin-finite}）得出。略去细节。
\end{proof}

\begin{lemma}
\label{spaces-divisors-lemma-relative-assassin-constructible}
设 $Y$ 为概形，$X$ 为 $Y$ 上有限表示的代数空间，$\mathcal{F}$ 为有限表示的
拟凝聚 $\mathcal{O}_X$-模。设 $U \subset X$ 是开子空间，并且
$U \to Y$ 拟紧。则集合
$$
E = \{y \in Y \mid \text{Ass}_{X_y}(\mathcal{F}_y) \subset |U_y|\}
$$
在 $Y$ 中局部可构造。
\end{lemma}

\begin{proof}
注意，由于 $Y$ 是概形，取纤维
$X_y = \Spec(\kappa(y)) \times_Y X$ 是有意义的。（此外，按我们的定义，
集合 $\text{Ass}_{X_y}(\mathcal{F}_y)$ 恰为
$\text{Ass}_{X/Y}(\mathcal{F}) \to Y$ 在 $y$ 上的纤维，但这里不需要
这一点。）问题在 $Y$ 上是局部的；事实上，只需在 $Y$ 仿射时证明 $E$
可构造。在这种情形，$X$ 拟紧。取仿射概形 $W$ 以及满平展态射
$\varphi : W \to X$。则对所有 $y \in Y$，
$\text{Ass}_{X_y}(\mathcal{F}_y)$ 是
$\text{Ass}_{W_y}(\varphi^*\mathcal{F}_y)$ 的像。因此，本引理由开集
$\varphi^{-1}(U) \subset W$ 与态射 $W \to Y$ 的概形情形得出。概形情形
就是《态射的更多内容》引理
\ref{more-morphisms-lemma-relative-assassin-constructible}。
\end{proof}









\section{Fitting 理想}
\label{spaces-divisors-section-fitting-ideals}

\noindent
本节继续《除子》第 \ref{divisors-section-fitting-ideals} 节中的讨论。
设 $S$ 为概形，$X$ 为 $S$ 上的代数空间，而 $\mathcal{F}$ 为有限型拟凝聚
$\mathcal{O}_X$-模。在这种情形，可以构造 Fitting 理想
$$
0 = \text{Fit}_{-1}(\mathcal{F}) \subset \text{Fit}_0(\mathcal{F}) \subset
\text{Fit}_1(\mathcal{F}) \subset \ldots \subset \mathcal{O}_X
$$
，即由下述性质刻画的拟凝聚理想层序列：
% P09-ERR-0166
对每个平展于 $X$ 的仿射概形 $U = \Spec(A)$，若 $\mathcal{F}|_U$ 对应于
$A$-模 $M$，则 $\text{Fit}_i(\mathcal{F})|_U$ 对应于理想
$\text{Fit}_i(M) \subset A$。这是良定义的拟凝聚理想层，因为若
$A \to B$ 是平展环同态，则由《代数的更多内容》引理
\ref{more-algebra-lemma-fitting-ideal-basics} 的 (3)，$M \otimes_A B$
作为 $B$-模的第 $i$ 个 Fitting 理想等于 $\text{Fit}_i(M) B$。更准确地说
（或许），拟凝聚理想层
% P09-ERR-0167
$\text{Fit}_0(\mathcal{O}_X)$ 的存在性例如可由《空间的性质》引理
\ref{spaces-properties-lemma-characterize-quasi-coherent-small-etale}
中对拟凝聚层的描述，以及《除子》引理
\ref{divisors-lemma-base-change-fitting-ideal} 给出的拉回性质得到。

\medskip\noindent
以这种方式构造 Fitting 理想的优点是，可以立即看出 Fitting 理想的形成与
平展局部化交换，因而 Fitting 理想的许多性质立即归约为概形情形的对应性质。
我们将经常使用《空间的性质》第
\ref{spaces-properties-section-properties-modules} 节中的讨论，在概形上的
拟凝聚层与代数空间上的拟凝聚层的性质之间作转化。

\begin{lemma}
\label{spaces-divisors-lemma-base-change-fitting-ideal}
设 $S$ 为概形，$f : X \to Y$ 为 $S$ 上代数空间间的态射，
$\mathcal{F}$ 为有限型拟凝聚 $\mathcal{O}_Y$-模。则
$f^{-1}\text{Fit}_i(\mathcal{F}) \cdot \mathcal{O}_X =
\text{Fit}_i(f^*\mathcal{F})$。
\end{lemma}

\begin{proof}
经平展局部化，归约为《除子》引理
\ref{divisors-lemma-base-change-fitting-ideal}。
\end{proof}

\begin{lemma}
\label{spaces-divisors-lemma-fitting-ideal-of-finitely-presented}
设 $S$ 为概形，$X$ 为 $S$ 上的代数空间，$\mathcal{F}$ 为有限表示
$\mathcal{O}_X$-模。则 $\text{Fit}_r(\mathcal{F})$ 是有限型拟凝聚理想。
\end{lemma}

\begin{proof}
经平展局部化，归约为《除子》引理
\ref{divisors-lemma-fitting-ideal-of-finitely-presented}。
\end{proof}

\begin{lemma}
\label{spaces-divisors-lemma-on-subscheme-cut-out-by-Fit-0}
设 $S$ 为概形，$X$ 为 $S$ 上的代数空间，而 $\mathcal{F}$ 为有限型拟凝聚
$\mathcal{O}_X$-模。设 $Z_0 \subset X$ 为由
$\text{Fit}_0(\mathcal{F})$ 切出的闭子空间，$Z \subset X$ 为
$\mathcal{F}$ 的概形论支撑。则
\begin{enumerate}
\item 作为闭子空间，$Z \subset Z_0 \subset X$；
\item 作为 $|X|$ 的闭子集，
$|Z| = |Z_0| = \text{Supp}(\mathcal{F})$；
\item 存在有限型拟凝聚 $\mathcal{O}_{Z_0}$-模 $\mathcal{G}_0$，使得
$$
(Z_0 \to X)_*\mathcal{G}_0 = \mathcal{F}.
$$
\end{enumerate}
\end{lemma}

\begin{proof}
回忆 $Z$ 的形成与平展局部化交换；参见《空间的态射》定义
\ref{spaces-morphisms-definition-scheme-theoretic-support}（其中用
《空间的态射》引理
\ref{spaces-morphisms-lemma-scheme-theoretic-support} 定义 $Z$）。因此
(1) 和 (2) 由概形情形得出；参见《除子》引理
\ref{divisors-lemma-on-subscheme-cut-out-by-Fit-0}。为得到 (3) 中的
$\mathcal{G}_0$，可使用《空间的态射》引理
\ref{spaces-morphisms-lemma-scheme-theoretic-support} 中 $Z$ 上的
$\mathcal{G}$，并置 $\mathcal{G}_0 = (Z \to Z_0)_*\mathcal{G}$。
\end{proof}

\begin{lemma}
\label{spaces-divisors-lemma-fitting-ideal-generate-locally}
设 $S$ 为概形，$X$ 为 $S$ 上的代数空间，$\mathcal{F}$ 为有限型拟凝聚
$\mathcal{O}_X$-模，且 $x \in |X|$。则 $\mathcal{F}$ 在 $x$ 的某个平展
邻域中可由 $r$ 个元素生成，当且仅当
$\text{Fit}_r(\mathcal{F})_{\overline{x}} = \mathcal{O}_{X, \overline{x}}$。
\end{lemma}

\begin{proof}
经平展局部化，归约为《除子》引理
\ref{divisors-lemma-fitting-ideal-generate-locally}（还要使用《空间的性质》
第 \ref{spaces-properties-section-stalks-structure-sheaf} 节中对局部环的
描述，以及局部环的严格亨泽尔化是忠实平坦的这一事实，以看出严格亨泽尔化
上的等式等价于局部环上的等式）。
\end{proof}

\begin{lemma}
\label{spaces-divisors-lemma-fitting-ideal-finite-locally-free}
设 $S$ 为概形，$X$ 为 $S$ 上的代数空间，$\mathcal{F}$ 为有限型拟凝聚
$\mathcal{O}_X$-模。设 $r \geq 0$。下列条件等价：
\begin{enumerate}
\item $\mathcal{F}$ 是秩为 $r$ 的有限局部自由模；
\item $\text{Fit}_{r - 1}(\mathcal{F}) = 0$ 且
$\text{Fit}_r(\mathcal{F}) = \mathcal{O}_X$；且
\item 对 $k < r$ 有 $\text{Fit}_k(\mathcal{F}) = 0$，而对
$k \geq r$ 有 $\text{Fit}_k(\mathcal{F}) = \mathcal{O}_X$。
\end{enumerate}
\end{lemma}

\begin{proof}
经平展局部化，归约为《除子》引理
\ref{divisors-lemma-fitting-ideal-finite-locally-free}。
\end{proof}

\begin{lemma}
\label{spaces-divisors-lemma-locally-free-rank-r-pullback}
设 $S$ 为概形，$X$ 为 $S$ 上的代数空间，$\mathcal{F}$ 为有限型拟凝聚
$\mathcal{O}_X$-模。由 $\mathcal{F}$ 的 Fitting 理想定义的闭子空间
% P09-ERR-0168
$$
X = Z_{-1} \supset Z_0 \supset Z_1 \supset Z_2 \ldots
$$
具有下列性质：
\begin{enumerate}
\item 交集 $\bigcap Z_r$ 为空。
\item 由下述规则定义的函子
$(\Sch/X)^{opp} \to \textit{Sets}$
$$
T \longmapsto
\left\{
\begin{matrix}
\{*\} & \text{若 }\mathcal{F}_T\text{ 局部可由 }
\leq r\text{ 个截面生成} \\
\emptyset & \text{否则}
\end{matrix}
\right.
$$
可由开子空间 $X \setminus Z_r$ 表示。
\item 由下述规则定义的函子
$F_r : (\Sch/X)^{opp} \to \textit{Sets}$
% P09-ERR-0169
$$
T \longmapsto
\left\{
\begin{matrix}
\{*\} & \text{若 }\mathcal{F}_T\text{ 是局部自由且秩为 }r\\
\emptyset & \text{否则}
\end{matrix}
\right.
$$
可由 $X$ 的局部闭子空间 $Z_{r - 1} \setminus Z_r$ 表示。
\end{enumerate}
若 $\mathcal{F}$ 有限表示，则
$Z_r \to X$、$X \setminus Z_r \to X$ 及
$Z_{r - 1} \setminus Z_r \to X$ 都是有限表示态射。
\end{lemma}

\begin{proof}
经平展局部化，归约为《除子》引理
\ref{divisors-lemma-locally-free-rank-r-pullback}。
\end{proof}

\begin{lemma}
\label{spaces-divisors-lemma-finite-presentation-module}
设 $S$ 为概形，$X$ 为 $S$ 上的代数空间，$\mathcal{F}$ 为有限表示
$\mathcal{O}_X$-模。设
% P09-ERR-0170
$X = Z_{-1} \subset Z_0 \subset Z_1 \subset \ldots$
如引理 \ref{spaces-divisors-lemma-locally-free-rank-r-pullback} 中所述。置
$X_r = Z_{r - 1} \setminus Z_r$。则
$X' = \coprod_{r \geq 0} X_r$ 表示函子
% P09-ERR-0171
% P09-ERR-0172
$$
F_{flat} : \Sch/X \longrightarrow \textit{Sets},\quad\quad
T \longmapsto
\left\{
\begin{matrix}
\{*\} & \text{若 }\mathcal{F}_T\text{ 平坦于 }T\\
\emptyset & \text{否则}
\end{matrix}
\right.
$$
。此外，$\mathcal{F}|_{X_r}$ 是秩为 $r$ 的局部自由模，且态射
$X_r \to X$ 与 $X' \to X$ 都是有限表示态射。
\end{lemma}

\begin{proof}
经平展局部化，归约为《除子》引理
\ref{divisors-lemma-finite-presentation-module}。
\end{proof}














\section{有效 Cartier 除子}
\label{spaces-divisors-section-effective-Cartier-divisors}

\noindent
% P09-ERR-0173
出于某种原因，先于其他一切概念定义有效 Cartier 除子的概念似乎较为方便。注意，
在《空间的态射》第 \ref{spaces-morphisms-section-closed-immersions} 节中，
我们讨论了闭子空间与拟凝聚理想层之间的对应。此外，在《空间的性质》第
\ref{spaces-properties-section-properties-modules} 节中，我们讨论了拟凝聚模的
性质，特别是“局部由 $1$ 个元素生成”这一性质。这些参考结果表明，下面的定义
与概形情形的定义相容。

\begin{definition}
\label{spaces-divisors-definition-effective-Cartier-divisor}
设 $S$ 为概形，$X$ 为 $S$ 上的代数空间。
\begin{enumerate}
\item $X$ 的一个{\it 局部主闭子空间}是指其理想层局部由 $1$ 个元素生成的
闭子空间。
\item $X$ 上的一个{\it 有效 Cartier 除子}是指闭子空间 $D \subset X$，
使得理想层 $\mathcal{I}_D \subset \mathcal{O}_X$ 是可逆
$\mathcal{O}_X$-模。
\end{enumerate}
\end{definition}

\noindent
因此，有效 Cartier 除子是局部主闭子空间，但逆命题并非总成立。有效
Cartier 除子是在最强意义下具有纯余维 $1$ 的闭子空间；也就是说，它们局部由
一个非零因子切出。特别地，它们无处稠密。

\begin{lemma}
\label{spaces-divisors-lemma-characterize-effective-Cartier-divisor}
设 $S$ 为概形，$X$ 为 $S$ 上的代数空间，并设 $D \subset X$ 为闭子空间。
下列条件等价：
\begin{enumerate}
\item 子空间 $D$ 是 $X$ 上的有效 Cartier 除子。
\item 存在某个概形 $U$ 及满平展态射 $U \to X$，使得逆像
$D \times_X U$ 是 $U$ 上的有效 Cartier 除子。
\item 对任意概形 $U$ 及任意平展态射 $U \to X$，逆像
$D \times_X U$ 都是 $U$ 上的有效 Cartier 除子。
\item 对每个 $x \in |D|$，存在带点代数空间的平展态射
$(U, u) \to (X, x)$，使得 $U = \Spec(A)$ 且
$D \times_X U = \Spec(A/(f))$，其中 $f \in A$ 不是零因子。
\end{enumerate}
\end{lemma}

\begin{proof}
(1)--(3) 的等价性由定义 \ref{spaces-divisors-definition-effective-Cartier-divisor} 及其前面的
参考结果立即得到。假设 (1) 成立，并令 $x \in |D|$。选取概形 $W$ 及满平展
态射 $W \to X$。选取映到 $x$ 的点 $w \in D \times_X W$。由 (3)，
$D \times_X W$ 是 $W$ 上的有效 Cartier 除子。于是，按照《除子》引理
\ref{divisors-lemma-characterize-effective-Cartier-divisor}，在 $W$ 中选取 $w$
的一个仿射开邻域，即可得到仿射平展邻域 $U$。
% P09-ERR-0174

\medskip\noindent
假设 (4) 成立。由《除子》引理
\ref{divisors-lemma-characterize-effective-Cartier-divisor} 可知，
$\mathcal{I}_D|_U$ 可逆。所有这样的 $U$ 与 $X \setminus D$ 合在一起给出
$X$ 的一个平展覆盖，故 $\mathcal{I}_D$ 是可逆 $\mathcal{O}_X$-模。
\end{proof}

\begin{lemma}
\label{spaces-divisors-lemma-complement-locally-principal-closed-subscheme}
设 $S$ 为概形，$X$ 为 $S$ 上的代数空间。设 $Z \subset X$ 为局部主闭子空间，
并令 $U = X \setminus Z$。则 $U \to X$ 是仿射态射。
\end{lemma}

\begin{proof}
% P09-ERR-0175
该问题在 $X$ 上是平展局部的，见《空间的态射》引理
\ref{spaces-morphisms-lemma-affine-local} 和引理
\ref{spaces-divisors-lemma-characterize-effective-Cartier-divisor}。因此，本结论归约为概形情形，
即《除子》引理
\ref{divisors-lemma-complement-locally-principal-closed-subscheme}。
\end{proof}

\begin{lemma}
\label{spaces-divisors-lemma-complement-effective-Cartier-divisor}
设 $S$ 为概形，$X$ 为 $S$ 上的代数空间。设 $D \subset X$ 为有效 Cartier
除子，并令 $U = X \setminus D$。则 $U \to X$ 是仿射态射，且 $U$ 在 $X$
中是概形论稠密的。
\end{lemma}

\begin{proof}
仿射性由引理 \ref{spaces-divisors-lemma-complement-locally-principal-closed-subscheme} 给出。
由《空间的态射》定义
\ref{spaces-morphisms-definition-scheme-theoretically-dense}，稠密性问题在 $X$
上是平展局部的。因此，本结论归约为概形情形，即《除子》引理
\ref{divisors-lemma-complement-effective-Cartier-divisor}。
\end{proof}

\begin{lemma}
\label{spaces-divisors-lemma-effective-Cartier-makes-dimension-drop}
设 $S$ 为概形，$X$ 为 $S$ 上的代数空间。设 $D \subset X$ 为有效 Cartier
除子，并令 $x \in |D|$。若 $\dim_x(X) < \infty$，则
$\dim_x(D) < \dim_x(X)$。
\end{lemma}

\begin{proof}
% P09-ERR-0176
有效 Cartier 除子的定义以及代数空间在一点处的维数之定义（《空间的性质》
定义 \ref{spaces-properties-definition-dimension-at-point}）都是平展局部的。
因此，本引理归约为概形情形，即《除子》引理
\ref{divisors-lemma-effective-Cartier-makes-dimension-drop}。
\end{proof}

\begin{definition}
\label{spaces-divisors-definition-sum-effective-Cartier-divisors}
设 $S$ 为概形，$X$ 为 $S$ 上的代数空间。给定 $X$ 上的有效 Cartier 除子
$D_1$、$D_2$，我们令 $D = D_1 + D_2$ 为 $X$ 的闭子空间，它对应于拟凝聚
理想层
% P09-ERR-0177
$\mathcal{I}_{D_1}\mathcal{I}_{D_2} \subset \mathcal{O}_S$。
称它为{\it 有效 Cartier 除子 $D_1$ 与 $D_2$ 的和}。
\end{definition}

\noindent
显然，给定 $X$ 上有限多个有效 Cartier 除子 $D_i$ 及非负整数 $n_i$，我们可以
定义和 $\sum n_iD_i$。

\begin{lemma}
\label{spaces-divisors-lemma-sum-effective-Cartier-divisors}
两个有效 Cartier 除子的和仍是有效 Cartier 除子。
\end{lemma}

\begin{proof}
略。经平展局部化，这归约为如下简单的代数事实：若 $f_1, f_2 \in A$ 是环
$A$ 中的非零因子，则 $f_1f_2 \in A$ 也是非零因子。
\end{proof}

\begin{lemma}
\label{spaces-divisors-lemma-sum-closed-subschemes-effective-Cartier}
设 $S$ 为概形，$X$ 为 $S$ 上的代数空间。设 $Z, Y$ 是 $X$ 的两个闭子空间，
其理想层分别为 $\mathcal{I}$ 和 $\mathcal{J}$。若
$\mathcal{I}\mathcal{J}$ 定义了有效 Cartier 除子 $D \subset X$，则 $Z$ 与
$Y$ 都是有效 Cartier 除子，并且 $D = Z + Y$。
\end{lemma}

\begin{proof}
由引理 \ref{spaces-divisors-lemma-characterize-effective-Cartier-divisor}，这归约为概形情形，
即《除子》引理 \ref{divisors-lemma-sum-closed-subschemes-effective-Cartier}。
\end{proof}

\noindent
回忆一下，我们已在《空间的态射》定义
\ref{spaces-morphisms-definition-inverse-image-closed-subspace} 中定义了闭子空间
在任意代数空间态射下的逆像。

\begin{lemma}
\label{spaces-divisors-lemma-pullback-locally-principal}
设 $S$ 为概形，$f : X' \to X$ 为 $S$ 上代数空间的态射，并设
$Z \subset X$ 为局部主闭子空间。则逆像 $f^{-1}(Z)$ 是 $X'$ 的局部主闭
子空间。
\end{lemma}

\begin{proof}
略。
\end{proof}

\begin{definition}
\label{spaces-divisors-definition-pullback-effective-Cartier-divisor}
设 $S$ 为概形，$f : X' \to X$ 为 $S$ 上代数空间的态射，并设 $D \subset X$
为有效 Cartier 除子。若闭子空间 $f^{-1}(D) \subset X'$ 是有效 Cartier
除子，则称{\it $D$ 沿 $f$ 的拉回有定义}。在这种情形下，将它记作 $f^*D$
或 $f^{-1}(D)$，并称之为{\it 该有效 Cartier 除子的拉回}。
\end{definition}

\noindent
在实际应用中，$f^{-1}(D)$ 是有效 Cartier 除子这一条件经常成立。

\begin{lemma}
\label{spaces-divisors-lemma-pullback-effective-Cartier-defined}
设 $S$ 为概形，$f : X \to Y$ 为 $S$ 上代数空间的态射，并设 $D \subset Y$
为有效 Cartier 除子。在下列每一种情形中，$D$ 沿 $f$ 的拉回都有定义：
\begin{enumerate}
\item 对 $X$ 的任意弱伴随点 $x$，均有 $f(x) \not \in |D|$；
\item $f$ 是平坦的；以及
% P09-ERR-0178
\item 按需在此补充。
\end{enumerate}
\end{lemma}

\begin{proof}
经平展局部化，本引理归约为概形情形，见《除子》引理
\ref{divisors-lemma-pullback-effective-Cartier-defined}。
\end{proof}

\begin{lemma}
\label{spaces-divisors-lemma-pullback-effective-Cartier-divisors-additive}
设 $S$ 为概形，$f : X' \to X$ 为 $S$ 上代数空间的态射，并设 $D_1$、
$D_2$ 为 $X$ 上的有效 Cartier 除子。若 $D_1$ 与 $D_2$ 的拉回都有定义，
则 $D = D_1 + D_2$ 的拉回也有定义，并且
$f^*D = f^*D_1 + f^*D_2$。
\end{lemma}

\begin{proof}
略。
\end{proof}




\section{有效 Cartier 除子与可逆层}
\label{spaces-divisors-section-effective-Cartier-invertible}

\noindent
由于有效 Cartier 除子具有可逆理想层（定义
\ref{spaces-divisors-definition-effective-Cartier-divisor}），下面的定义是有意义的。

\begin{definition}
\label{spaces-divisors-definition-invertible-sheaf-effective-Cartier-divisor}
设 $S$ 为概形，$X$ 为 $S$ 上的代数空间，并设 $D \subset X$ 为有效 Cartier
除子，其理想层为 $\mathcal{I}_D$。
\begin{enumerate}
\item {\it 与 $D$ 相伴的可逆层 $\mathcal{O}_X(D)$} 定义为
$$
\mathcal{O}_X(D) =
\SheafHom_{\mathcal{O}_X}(\mathcal{I}_D, \mathcal{O}_X) =
\mathcal{I}_D^{\otimes -1}.
$$
\item 典范截面通常记作 $1$ 或 $1_D$，它是 $\mathcal{O}_X(D)$ 的整体截面，
对应于包含映射 $\mathcal{I}_D \to \mathcal{O}_X$。
\item 记
$\mathcal{O}_X(-D) = \mathcal{O}_X(D)^{\otimes -1} = \mathcal{I}_D$.
\item 给定另一个有效 Cartier 除子 $D' \subset X$，定义
$\mathcal{O}_X(D - D') =
\mathcal{O}_X(D) \otimes_{\mathcal{O}_X} \mathcal{O}_X(-D')$.
\end{enumerate}
\end{definition}

\noindent
作几点说明。下文将看到，赋值 $D \mapsto \mathcal{O}_X(D)$ 把有效 Cartier
除子的加法（定义 \ref{spaces-divisors-definition-sum-effective-Cartier-divisors}）变为 $X$ 的
Picard 群中的加法（引理
\ref{spaces-divisors-lemma-invertible-sheaf-sum-effective-Cartier-divisors}）。然而，上述定义中
的表达式 $D - D'$ 本身没有任何几何意义。更精确地，可以把 $X$ 上所有有效
Cartier 除子构成的集合看成交换幺半群 $\text{EffCart}(X)$，其零元为空的有效
Cartier 除子。于是，赋值 $(D, D') \mapsto \mathcal{O}_X(D - D')$ 定义了群同态
$$
\text{EffCart}(X)^{gp} \longrightarrow \Pic(X)
$$
，其中左端是 $\text{EffCart}(X)$ 的群化。换言之，在写
$\mathcal{O}_X(D - D')$ 时，可以把 $D - D'$ 看成
$\text{EffCart}(X)^{gp}$ 的一个元素。

\begin{lemma}
\label{spaces-divisors-lemma-conormal-effective-Cartier-divisor}
设 $S$ 为概形，$X$ 为 $S$ 上的代数空间，并设 $D \subset X$ 为有效 Cartier
除子。则余法层满足
% P09-ERR-0179
$\mathcal{C}_{D/X} = \mathcal{I}_D|D =
\mathcal{O}_X(D)^{\otimes -1}|_D$。
\end{lemma}

\begin{proof}
略。
\end{proof}

\begin{lemma}
\label{spaces-divisors-lemma-invertible-sheaf-sum-effective-Cartier-divisors}
设 $S$ 为概形，$X$ 为 $S$ 上的代数空间。设 $D_1$、$D_2$ 为 $X$ 上的有效
Cartier 除子，并令 $D = D_1 + D_2$。则存在唯一的同构
$$
\mathcal{O}_X(D_1) \otimes_{\mathcal{O}_X} \mathcal{O}_X(D_2)
\longrightarrow
\mathcal{O}_X(D)
$$
，它把 $1_{D_1} \otimes 1_{D_2}$ 映到 $1_D$。
\end{lemma}

\begin{proof}
略。
\end{proof}

\begin{definition}
\label{spaces-divisors-definition-regular-section}
设 $S$ 为概形，$X$ 为 $S$ 上的代数空间，并设 $\mathcal{L}$ 为 $X$ 上的可逆
层。若映射 $\mathcal{O}_X \to \mathcal{L}$，$f \mapsto fs$ 是单射，则称整体
截面 $s \in \Gamma(X, \mathcal{L})$ 为{\it 正则截面}。
\end{definition}

\begin{lemma}
\label{spaces-divisors-lemma-regular-section-structure-sheaf}
设 $S$ 为概形，$X$ 为 $S$ 上的代数空间，并设
$f \in \Gamma(X, \mathcal{O}_X)$。下列条件等价：
\begin{enumerate}
\item $f$ 是正则截面；
\item 对任意 $x \in X$，像 $f \in \mathcal{O}_{X, \overline{x}}$ 都不是
零因子；
\item 对任意平展于 $X$ 的仿射概形 $U = \Spec(A)$，限制 $f|_U$ 都是 $A$
中的非零因子；以及
\item 存在概形 $U$ 及满平展态射 $U \to X$，使得 $f|_U$ 是
$\mathcal{O}_U$ 的正则截面。
\end{enumerate}
\end{lemma}

\begin{proof}
略。
\end{proof}

\noindent
注意，可逆 $\mathcal{O}_X$-模 $\mathcal{L}$ 的整体截面 $s$ 可以看作
$\mathcal{O}_X$-模映射 $s : \mathcal{O}_X \to \mathcal{L}$。因此，其对偶是
映射 $s : \mathcal{L}^{\otimes -1} \to \mathcal{O}_X$。（关于对偶可逆层，
见《位点上的模》引理 \ref{sites-modules-lemma-constructions-invertible}。）

\begin{definition}
\label{spaces-divisors-definition-zero-scheme-s}
设 $S$ 为概形，$X$ 为 $S$ 上的代数空间。设 $\mathcal{L}$ 为可逆层，并设
$s \in \Gamma(X, \mathcal{L})$。$s$ 的{\it 零点概形}是闭子空间
$Z(s) \subset X$，它由拟凝聚理想层 $\mathcal{I} \subset \mathcal{O}_X$
定义；该理想层是映射
$s : \mathcal{L}^{\otimes -1} \to \mathcal{O}_X$ 的像。
\end{definition}

\begin{lemma}
\label{spaces-divisors-lemma-zero-scheme}
设 $S$ 为概形，$X$ 为 $S$ 上的代数空间。设 $\mathcal{L}$ 为可逆
$\mathcal{O}_X$-模，并设 $s \in \Gamma(X, \mathcal{L})$。
\begin{enumerate}
\item 考虑满足下述条件的闭浸入 $i : Z \to X$：
% P09-ERR-0180
$i^*s \in \Gamma(Z, i^*\mathcal{L}))$ 为零；并按包含关系对这些闭浸入排序。
零点概形 $Z(s)$ 是这个偏序集的最大元。
\item 对任意 $S$ 上代数空间的态射 $f : Y \to X$，$f^*s = 0$ 在
$\Gamma(Y, f^*\mathcal{L})$ 中成立，当且仅当 $f$ 经由 $Z(s)$ 分解。
\item 零点概形 $Z(s)$ 是 $X$ 的局部主闭子空间。
\item 零点概形 $Z(s)$ 是 $X$ 上的有效 Cartier 除子，当且仅当 $s$ 是
$\mathcal{L}$ 的正则截面。
\end{enumerate}
\end{lemma}

\begin{proof}
略。
\end{proof}

\begin{lemma}
\label{spaces-divisors-lemma-characterize-OD}
设 $S$ 为概形，$X$ 为 $S$ 上的代数空间。
\begin{enumerate}
\item 若 $D \subset X$ 是有效 Cartier 除子，则 $\mathcal{O}_X(D)$ 的典范
截面 $1_D$ 是正则的。
\item 反之，若 $s$ 是可逆层 $\mathcal{L}$ 的正则截面，则存在唯一的有效
Cartier 除子 $D = Z(s) \subset X$ 以及唯一的同构
$\mathcal{O}_X(D) \to \mathcal{L}$，该同构把 $1_D$ 映到 $s$。
\end{enumerate}
构造
$D \mapsto (\mathcal{O}_X(D), 1_D)$ 与 $(\mathcal{L}, s) \mapsto Z(s)$
给出下列集合之间的互逆映射：
$$
\left\{
\begin{matrix}
\text{有效 Cartier 除子，其所在空间为 }X
\end{matrix}
\right\}
\leftrightarrow
\left\{
\begin{matrix}
\text{二元组 }(\mathcal{L}, s)\text{，其中所含对象为一个可逆}\\
\mathcal{O}_X\text{-模和一个正则整体截面}
\end{matrix}
\right\}
$$
\end{lemma}

\begin{proof}
略。
\end{proof}





\section{诺特空间上的有效 Cartier 除子}
\label{spaces-divisors-section-Noetherian-effective-Cartier}

\noindent
在局部诺特情形下，关于有效 Cartier 除子和正则截面的大部分讨论都会有所简化。

\begin{lemma}
\label{spaces-divisors-lemma-effective-Cartier-divisor-Sk}
设 $S$ 为概形，$X$ 为 $S$ 上的局部诺特代数空间。设 $D \subset X$ 为有效
Cartier 除子。若 $X$ 满足 $(S_k)$，则 $D$ 满足 $(S_{k - 1})$。
\end{lemma}

\begin{proof}
由我们对代数空间的性质 $(S_k)$ 所作的定义（《空间的性质》第
\ref{spaces-properties-section-types-properties} 节）以及引理
\ref{spaces-divisors-lemma-characterize-effective-Cartier-divisor}，本结论归约为概形情形
（《除子》引理 \ref{divisors-lemma-effective-Cartier-divisor-Sk}）。
\end{proof}

\begin{lemma}
\label{spaces-divisors-lemma-normal-effective-Cartier-divisor-S1}
设 $S$ 为概形，$X$ 为 $S$ 上的局部诺特正规代数空间。设 $D \subset X$ 为
有效 Cartier 除子。则 $D$ 满足 $(S_1)$。
\end{lemma}

\begin{proof}
由我们对代数空间的正规性所作的定义（《空间的性质》第
\ref{spaces-properties-section-types-properties} 节）以及引理
\ref{spaces-divisors-lemma-characterize-effective-Cartier-divisor}，本结论归约为概形情形
（《除子》引理 \ref{divisors-lemma-normal-effective-Cartier-divisor-S1}）。
\end{proof}

\noindent
下面的引理有时可用于构造有效 Cartier 除子。

\begin{lemma}
\label{spaces-divisors-lemma-complement-open-affine-effective-cartier-divisor}
设 $S$ 为概形，$X$ 为 $S$ 上的正则诺特分离代数空间。设 $U \subset X$ 为
稠密仿射开子空间。则存在有效 Cartier 除子 $D \subset X$，使得
$U = X \setminus D$。
\end{lemma}

\begin{proof}
我们断言，$X \setminus U$ 上的约化诱导代数空间结构 $D$（《空间的性质》
定义 \ref{spaces-properties-definition-reduced-induced-space}）就是所求的有效
Cartier 除子。$D$ 的构造与平展局部化可交换，见《空间的性质》引理
\ref{spaces-properties-lemma-reduced-closed-subspace} 的证明。取满平展态射
$X' \to X$，其中 $X'$ 为仿射概形。由于 $X$ 分离，
$U' = X' \times_X U$ 是仿射的。由于 $|X'| \to |X|$ 是开映射，$U'$ 在 $X'$
中稠密。又由于 $D' = X' \times_X D$ 是 $X' \setminus U'$ 上诱导的约化概形
结构，由《除子》引理
\ref{divisors-lemma-complement-open-affine-effective-cartier-divisor} 及其证明，
$D'$ 是有效 Cartier 除子。这正是所要证明的。
\end{proof}

\begin{lemma}
\label{spaces-divisors-lemma-Noetherian-regular-separated-pic-effective-Cartier}
设 $S$ 为概形，$X$ 为 $S$ 上的正则诺特分离代数空间。则每个可逆
$\mathcal{O}_X$-模都同构于
$$
\mathcal{O}_X(D - D') =
\mathcal{O}_X(D) \otimes_{\mathcal{O}_X} \mathcal{O}_X(D')^{\otimes -1}
$$
，其中 $D, D'$ 是 $X$ 上的某些有效 Cartier 除子。
\end{lemma}

\begin{proof}
设 $\mathcal{L}$ 为可逆 $\mathcal{O}_X$-模。选取稠密仿射开子空间
$U \subset X$，使得 $\mathcal{L}|_U$ 平凡。这是可能的，因为 $X$ 有一个为
概形的稠密开子空间，见《空间的性质》命题
\ref{spaces-properties-proposition-locally-quasi-separated-open-dense-scheme}。
% P09-ERR-0181
将平凡化记作 $s : \mathcal{O}_U \to \mathcal{L}|_U$。$U$ 的补集是一个有效
Cartier 除子 $D$。我们断言，对某个 $n > 0$，映射 $s$ 唯一延拓为映射
$$
s : \mathcal{O}_X(-nD) \longrightarrow \mathcal{L}
$$
。该断言蕴含本引理，因为它表明
$\mathcal{L} \otimes_{\mathcal{O}_X} \mathcal{O}_X(nD)$ 有一个正则整体截面，
故由引理 \ref{spaces-divisors-lemma-characterize-OD}，它同构于某个有效 Cartier 除子 $D'$
所对应的 $\mathcal{O}_X(D')$。为证明该断言，可以在平展局部工作。因此可设
$X$ 为仿射诺特概形。由于
$\mathcal{O}_X(-nD) = \mathcal{I}^n$，其中
$\mathcal{I} = \mathcal{O}_X(-D)$ 是 $D$ 在 $X$ 中的理想层，此情形由
《概形的上同调》引理 \ref{coherent-lemma-homs-over-open} 得到。
\end{proof}

\noindent
下面的引理其实应当放在另一节中。

\begin{lemma}
\label{spaces-divisors-lemma-smooth-over-valuation-ring-effective-Cartier}
设 $R$ 为赋值环，其分式域为 $K$。设 $X$ 为 $R$ 上的代数空间，且
$X \to \Spec(R)$ 光滑。对每个有效 Cartier 除子 $D \subset X_K$，都存在
有效 Cartier 除子 $D' \subset X$，使得 $D'_K = D$。
\end{lemma}

\begin{proof}
令 $D' \subset X$ 为 $D \to X_K \to X$ 的概形论像。由于该态射拟紧，$D'$
的形成与平坦基变换可交换，见《空间的态射》引理
\ref{spaces-morphisms-lemma-flat-base-change-scheme-theoretic-image}。特别地，
$D'_K = D$。因此可以假设 $X$ 仿射，记 $X = \Spec(A)$。于是
$X_K = \Spec(A \otimes_R K)$，且 $D$ 对应于理想
$I \subset A \otimes_R K$。我们必须证明 $J = I \cap A$ 在 $X$ 中切出一个
有效 Cartier 除子。首先注意，$A/J$ 在 $R$ 上平坦（因为它是无挠 $R$-模，
见《代数进阶》引理
\ref{more-algebra-lemma-valuation-ring-torsion-free-flat}），因而由《代数进阶》
引理
\ref{more-algebra-lemma-flat-finite-type-valuation-ring-finite-presentation}
及《代数》引理 \ref{algebra-lemma-extension}，$J$ 是有限生成的。因此，只需证明
对每个素理想 $\mathfrak q \subset A$，$J_\mathfrak q \subset A_\mathfrak q$
均由单个元素生成。令 $\mathfrak p = R \cap \mathfrak q$。则
$R_\mathfrak p$ 是赋值环（《代数》引理
\ref{algebra-lemma-make-valuation-rings}）。另由《代数》引理
\ref{algebra-lemma-characterize-smooth-over-field}，
$A_\mathfrak q/\mathfrak p A_\mathfrak q$ 是正则环。因此，可以应用《代数
进阶》引理 \ref{more-algebra-lemma-picard-group-generic-fibre-regular} 得到
$I(A_\mathfrak q \otimes_R K)$ 由单个元素
$f \in A_\mathfrak p \otimes_R K$ 生成。清去分母后，可以假设
$f \in A_\mathfrak q$。令 $\mathfrak c \subset R_\mathfrak p$ 为 $f$ 的内容
理想（见《代数进阶》定义 \ref{more-algebra-definition-content-ideal} 及《平坦性
进阶》引理 \ref{flat-lemma-flat-finite-type-local-valuation-ring-has-content}）。
由于 $R_\mathfrak p$ 是赋值环，且 $\mathfrak c$ 有限生成（《代数进阶》引理
\ref{more-algebra-lemma-content-finitely-generated}），存在
$\pi \in R_\mathfrak p$ 使得 $\mathfrak c = (\pi)$（《代数》引理
\ref{algebra-lemma-characterize-valuation-ring}）。
% P09-ERR-0182
把 $f$ 替换为 $\pi^{-1}f$ 后，有 $f \in A_\mathfrak q$ 且
$f \not \in \mathfrak pA_\mathfrak q$。
% P09-ERR-0183
断言：$I_\mathfrak q = (f)$；这将完成证明。为证明该断言，注意
$f \in I_\mathfrak q$。因而有满射
$A_\mathfrak q/(f) \to A_\mathfrak q/I_\mathfrak q$，它在 $R$ 上与 $K$ 张量
后成为同构。因此，只要 $A_\mathfrak q/(f)$ 在 $R_\mathfrak p$ 上平坦，证明
就完成了。这由《代数》引理 \ref{algebra-lemma-grothendieck-general} 及我们对
$f$ 的选择得到。
\end{proof}









\section{相对有效 Cartier 除子}
\label{spaces-divisors-section-effective-Cartier-morphisms}

\noindent
下面的引理表明，在基上平坦的有效 Cartier 除子实际上是基上的一个“有效
Cartier 除子族”。例如，它在任意纤维上的限制都是有效 Cartier 除子。

\begin{lemma}
\label{spaces-divisors-lemma-relative-Cartier}
设 $S$ 为概形，$f : X \to Y$ 为 $S$ 上代数空间的态射，并设
$D \subset X$ 为闭子空间。假设
\begin{enumerate}
\item $D$ 是有效 Cartier 除子；并且
\item $D \to Y$ 是平坦态射。
\end{enumerate}
% P09-ERR-0184
则对每个概形态射 $g : Y' \to Y$，拉回 $(g')^{-1}D$ 都是
$X' = Y' \times_Y X$ 上的有效 Cartier 除子，其中
$g' : X' \to X$ 是投影。
\end{lemma}

\begin{proof}
由引理 \ref{spaces-divisors-lemma-characterize-effective-Cartier-divisor}，成为有效 Cartier
除子这一性质是平展局部的。因此，本引理立即归约为概形情形，即《除子》引理
\ref{divisors-lemma-relative-Cartier}。
% P09-ERR-0185
\end{proof}

\noindent
这个引理促成了下面的定义。

\begin{definition}
\label{spaces-divisors-definition-relative-effective-Cartier-divisor}
设 $S$ 为概形，$f : X \to Y$ 为 $S$ 上代数空间的态射。$X/Y$ 上的一个
{\it 相对有效 Cartier 除子}是指有效 Cartier 除子 $D \subset X$，使得
$D \to Y$ 是代数空间的平坦态射。
\end{definition}








\section{亚纯函数与截面}
\label{spaces-divisors-section-meromorphic-functions}

\noindent
本节对应于《除子》第 \ref{divisors-section-meromorphic-functions} 节。请注意：
% P09-ERR-0186
在代数空间情形下，处理这些材料甚至比在概形情形下更容易出错！

\medskip\noindent
设 $S$ 为概形，$X$ 为 $S$ 上的代数空间。对任意平展于 $X$ 的概形 $U$，
我们已经定义了集合
$\mathcal{S}(U) \subset \mathcal{O}_X(U)$
，它由 $U$ 上 $\mathcal{O}_X$ 的正则截面组成，见定义
\ref{spaces-divisors-definition-regular-section}。正则截面限制到平展的 $V/U$ 上仍是正则的。
因此，$\mathcal{S} : U \mapsto \mathcal{S}(U)$ 是 $\mathcal{O}_X$ 的一个
（集合）子层。若要标明它对 $X$ 的依赖性，有时记
$\mathcal{S} = \mathcal{S}_X$。此外，对每个 $U$，$\mathcal{S}(U)$ 都是环
$\mathcal{O}_X(U)$ 的乘法子集。因此，可以考虑 $X_\etale$ 上的环预层
$$
U \longmapsto \mathcal{S}(U)^{-1} \mathcal{O}_X(U),
$$
及其层化，见《位点上的模》第
\ref{sites-modules-section-localizing-sheaves-rings} 节。

\begin{definition}
\label{spaces-divisors-definition-sheaf-meromorphic-functions}
设 $S$ 为概形，$X$ 为 $S$ 上的代数空间。$X$ 上的{\it 亚纯函数层}是与上面
所示预层相伴的 $X_\etale$ 上的层 {\it $\mathcal{K}_X$}。$X$ 上的一个
{\it 亚纯函数}是 $\mathcal{K}_X$ 的整体截面。
\end{definition}

\noindent
由于每个 $\mathcal{S}(U)$ 的每个元素都是 $\mathcal{O}_X(U)$ 中的非零因子，
可见环层的自然映射 $\mathcal{O}_X \to \mathcal{K}_X$ 是单射。此外，由层化
与取茎的相容性，对 $X$ 的任意几何点 $\overline{x}$，有
$$
\mathcal{K}_{X, \overline{x}} =
\mathcal{S}_{\overline{x}}^{-1}\mathcal{O}_{X, \overline{x}}
$$
。集合 $\mathcal{S}_{\overline{x}}$ 是 $\mathcal{O}_{X, \overline{x}}$ 中
非零因子集合的子集，但一般并不等于后者。

\begin{lemma}
\label{spaces-divisors-lemma-describe-S}
设 $S$ 为概形，$X$ 为 $S$ 上的代数空间。若 $U$ 是平展于 $X$ 的仿射概形，
则集合 $\mathcal{S}_X(U)$ 正是 $\mathcal{O}_X(U)$ 中非零因子的集合。
\end{lemma}

\begin{proof}
由引理 \ref{spaces-divisors-lemma-regular-section-structure-sheaf} 得到。
\end{proof}

\noindent
接着，设 $\mathcal{F}$ 为 $X_\etale$ 上的 $\mathcal{O}_X$-模层。考虑预层
$U \mapsto \mathcal{S}(U)^{-1}\mathcal{F}(U)$。它的层化是
$\mathcal{F} \otimes_{\mathcal{O}_X} \mathcal{K}_X$，见《位点上的模》引理
\ref{sites-modules-lemma-simple-invert-module}。

\begin{definition}
\label{spaces-divisors-definition-meromorphic-section}
设 $S$ 为概形，$X$ 为 $S$ 上的代数空间，并设 $\mathcal{F}$ 为 $X_\etale$
上的 $\mathcal{O}_X$-模层。
\begin{enumerate}
% P09-ERR-0187
\item 用 $\mathcal{K}_X(\mathcal{F})$ 表示如下 $\mathcal{K}_X$-模层：它是预层
$U \mapsto \mathcal{S}(U)^{-1}\mathcal{F}(U)$ 的层化。等价地，
$\mathcal{K}_X(\mathcal{F}) =
\mathcal{F} \otimes_{\mathcal{O}_X} \mathcal{K}_X$（见上文）。
\item $\mathcal{F}$ 的一个{\it 亚纯截面}是
$\mathcal{K}_X(\mathcal{F})$ 的整体截面。
\end{enumerate}
\end{definition}

\noindent
特别地，对 $X$ 的任意几何点 $\overline{x}$，有
$$
\mathcal{K}_X(\mathcal{F})_{\overline{x}}
=
\mathcal{F}_{\overline{x}}
\otimes_{\mathcal{O}_{X, \overline{x}}} \mathcal{K}_{X, \overline{x}}
=
\mathcal{S}_{\overline{x}}^{-1}\mathcal{F}_{\overline{x}}
$$
。不过必须谨慎，因为如上所述，$\mathcal{S}_{\overline{x}}$ 未必是平展局部环
$\mathcal{O}_{X, \overline{x}}$ 中非零因子的集合。亚纯截面层一般不是拟凝聚
模，但确实与拟凝聚模共有一些性质。

\begin{lemma}
\label{spaces-divisors-lemma-meromorphic-quasi-coherent}
设 $S$ 为概形，$X$ 为 $S$ 上的代数空间。假设：
\begin{enumerate}
\item[(a)] $X$ 的每个弱伴随点都是余维 $0$ 的点；并且
\item[(b)] $X$ 满足《空间的态射》引理
\ref{spaces-morphisms-lemma-prepare-normalization} 中的等价条件。
\end{enumerate}
则：
\begin{enumerate}
\item $\mathcal{K}_X$ 是拟凝聚 $\mathcal{O}_X$-代数层；
\item 对仿射的 $U \in X_\etale$，$\mathcal{K}_X(U)$ 是
$\mathcal{O}_X(U)$ 的全分式环；
% P09-ERR-0188
\item 对几何点 $\overline{x}$，集合 $\mathcal{S}_{\overline{x}}$ 是
$\mathcal{O}_{X, \overline{x}}$ 中非零因子的集合；并且
\item 对几何点 $\overline{x}$，环 $\mathcal{K}_{X, \overline{x}}$ 是
$\mathcal{O}_{X, \overline{x}}$ 的全分式环。
\end{enumerate}
\end{lemma}

\begin{proof}
% P09-ERR-0189
由引理 \ref{spaces-divisors-lemma-regular-section-structure-sheaf} 可知，对仿射的
$U \in X_\etale$，$\mathcal{S}_X(U) \subset \mathcal{O}_X(U)$ 是
$\mathcal{O}_X(U)$ 中非零因子的集合。因此，预层
$\mathcal{S}^{-1}\mathcal{O}_X$ 在 $X_{affine, \etale}$ 上等于
$$
U \longmapsto Q(\mathcal{O}_X(U))
$$
，记号见《代数》例 \ref{algebra-example-localize-at-prime}。注意，$X$ 的余维
$0$ 点对应于 $U$ 的一般点，见《空间的性质》引理
\ref{spaces-properties-lemma-codimension-0-points}。因此，若
$U = \Spec(A)$，则环 $A$ 只有有限多个极小素理想，并且 $A$ 的每个弱伴随
素理想都是极小的。对 $A$ 的任意平展扩张，同样的结论成立（因为这种扩张的
谱是平展于 $X$ 的仿射概形，因而可在上一句中充当 $A$）。为证明我们的预层
是层且拟凝聚，只需证明，当 $A \to B$ 是平展环同态时，
$$
Q(A) \otimes_A B \longrightarrow Q(B)
$$
是同构，见《空间的性质》引理
\ref{spaces-properties-lemma-characterize-quasi-coherent-small-etale}。（为定义上述
箭头，注意 $A \to B$ 平坦，故它把非零因子映为非零因子。）由《代数》引理
\ref{algebra-lemma-total-ring-fractions-no-embedded-points} 及
\ref{algebra-lemma-weakly-ass-zero-divisors}
% P09-ERR-0190
，有
$$
Q(A) = \prod\nolimits_{\mathfrak p \subset A\text{ 为极小素理想}} A_\mathfrak p
\quad\text{且}\quad
Q(B) = \prod\nolimits_{\mathfrak q \subset B\text{ 为极小素理想}} B_\mathfrak q
$$
。由于 $A \to B$ 平展，$B$ 的极小素理想恰好是位于 $A$ 的极小素理想之上的
$B$ 的素理想（例如由《代数进阶》引理
\ref{more-algebra-lemma-dimension-etale-extension}）。由《代数》引理
\ref{algebra-lemma-local-dimension-zero-henselian}、
\ref{algebra-lemma-characterize-henselian} (13) 及
\ref{algebra-lemma-mop-up}，$A_\mathfrak p \otimes_A B$ 是有限多个局部环的
直积，这些局部环都是 $A_\mathfrak p$ 上的有限平展代数。这显然蕴含
$A_\mathfrak p \otimes_A B =
\prod_{\mathfrak q\text{ 位于 }\mathfrak p} B_\mathfrak q$，正如所需。

\medskip\noindent
至此，我们已经知道 (1) 和 (2) 成立。证明 (3)。设
$s \in \mathcal{O}_{X, \overline{x}}$ 为非零因子。可以找到平展邻域
$(U, \overline{u}) \to (X, \overline{x})$ 以及映到 $s$ 的
$f \in \mathcal{O}_X(U)$。令 $u \in U$ 为 $\overline{u}$ 所确定的点。由于
$\mathcal{O}_{U, u} \to \mathcal{O}_{X, \overline{x}}$ 忠实平坦（后者是前者
的严格亨泽尔化），可见 $f$ 在 $\mathcal{O}_{U, u}$ 中的像是非零因子。由
《除子》引理 \ref{divisors-lemma-meromorphic-weakass-finite}，缩小 $U$ 后可使
$f$ 成为非零因子，因而成为 $\mathcal{S}_X(U)$ 的截面。(4) 通过由 (3) 计算
茎得到。
\end{proof}

\begin{lemma}
\label{spaces-divisors-lemma-meromorphic-quasi-coherent-agree}
设 $S$ 为概形，$X$ 为 $S$ 上的代数空间。假设：
\begin{enumerate}
\item[(a)] $X$ 的每个弱伴随点都是余维 $0$ 的点；
\item[(b)] $X$ 满足《空间的态射》引理
\ref{spaces-morphisms-lemma-prepare-normalization} 中的等价条件；
\item[(c)] $X$ 可由概形 $X_0$ 表示（这是一个别扭但临时使用的记号）。
\end{enumerate}
则亚纯函数层 $\mathcal{K}_X$ 是与拟凝聚亚纯函数层
$\mathcal{K}_{X_0}$ 相伴的拟凝聚 $\mathcal{O}_X$-代数层。
\end{lemma}

\begin{proof}
关于 $\QCoh(\mathcal{O}_X)$ 与 $\QCoh(\mathcal{O}_{X_0})$ 之间的等价，见
《空间的性质》第 \ref{spaces-properties-section-quasi-coherent} 节。本引理成立，
因为由引理 \ref{spaces-divisors-lemma-meromorphic-quasi-coherent} 及《除子》引理
\ref{divisors-lemma-meromorphic-weakass-finite}，$\mathcal{K}_X$ 与
$\mathcal{K}_{X_0}$ 都拟凝聚，并且在 $X$ 与 $X_0$ 的相应仿射开集上取值相同。
\end{proof}

\begin{definition}
\label{spaces-divisors-definition-pullback-meromorphic-sections}
设 $S$ 为概形，$f : X \to Y$ 为 $S$ 上代数空间的态射。若对每个交换图
$$
\xymatrix{
U \ar[r] \ar[d] & X \ar[d] \\
V \ar[r] & Y
}
$$
，其中 $U \in X_\etale$ 且 $V \in Y_\etale$，以及任意截面
$s \in \mathcal{S}_Y(V)$，拉回 $f^\sharp(s) \in \mathcal{O}_X(U)$ 都属于
$\mathcal{S}_X(U)$，则称{\it 亚纯函数沿 $f$ 的拉回有定义}。
\end{definition}

\noindent
在这种情形下，存在诱导映射
$f^\sharp : f_{small}^{-1}\mathcal{K}_Y \to \mathcal{K}_X$；换言之，得到带环
拓扑斯态射的交换图
$$
\xymatrix{
(\Sh(X_\etale), \mathcal{K}_X) \ar[r] \ar[d]^{f_{small}} &
(\Sh(X_\etale), \mathcal{O}_X) \ar[d]^{f_{small}} \\
(\Sh(Y_\etale), \mathcal{K}_Y) \ar[r] &
(\Sh(Y_\etale), \mathcal{O}_Y)
}
$$
对截面 $s \in \Gamma(Y, \mathcal{K}_Y)$，有时记
$f^*(s) = f^\sharp(s)$。

\begin{lemma}
\label{spaces-divisors-lemma-pullback-meromorphic-sections-defined}
% P09-ERR-0191
设 $S$ 为概形，$f : X \to Y$ 为 $S$ 上代数空间的态射。在下列每一种情形中，
亚纯截面的拉回都有定义：
\begin{enumerate}
\item $X$ 的弱伴随点被映到 $Y$ 的余维 $0$ 点；
\item $f$ 是平坦的；
% P09-ERR-0192
\item 按需在此补充。
\end{enumerate}
\end{lemma}

\begin{proof}
经平展局部化，这转化为概形情形，见《除子》引理
\ref{divisors-lemma-pullback-meromorphic-sections-defined}。完成这一转化时，使用
引理 \ref{spaces-divisors-lemma-regular-section-structure-sheaf}（正则截面的描述）、定义
\ref{spaces-divisors-definition-weakly-associated}（弱伴随点的定义）以及《空间的性质》引理
\ref{spaces-properties-lemma-codimension-0-points}（余维 $0$ 点的描述）。
\end{proof}

\begin{lemma}
\label{spaces-divisors-lemma-compute-meromorphic}
设 $S$ 为概形，$X$ 为 $S$ 上的代数空间。假设：
\begin{enumerate}
\item[(a)] $X$ 的每个弱伴随点都是余维 $0$ 的点；
\item[(b)] $X$ 满足《空间的态射》引理
\ref{spaces-morphisms-lemma-prepare-normalization} 中的等价条件；
\item[(c)] $X$ 的每个余维 $0$ 点都可由单态射 $\Spec(k) \to X$ 表示。
\end{enumerate}
令 $X^0 \subset |X|$ 为 $X$ 的余维 $0$ 点所成的集合。则有
$$
\mathcal{K}_X =
\bigoplus\nolimits_{\eta \in X^0} j_{\eta, *}\mathcal{O}_{X, \eta} =
\prod\nolimits_{\eta \in X^0} j_{\eta, *}\mathcal{O}_{X, \eta}
$$
，其中 $j_\eta : \Spec(\mathcal{O}_{X, \eta}) \to X$ 是《概形》第
\ref{schemes-section-points} 节中的典范映射；这是有意义的，因为 $X^0$ 包含在
$X$ 的概形轨迹中。类似地，对每个拟凝聚 $\mathcal{O}_X$-模 $\mathcal{F}$，
得到公式
$$
\mathcal{K}_X(\mathcal{F}) =
\bigoplus\nolimits_{\eta \in X^0} j_{\eta, *}\mathcal{F}_\eta =
\prod\nolimits_{\eta \in X^0} j_{\eta, *}\mathcal{F}_\eta
$$
，它描述 $\mathcal{F}$ 的亚纯截面层。最后，$X$ 的有理函数环就是 $X$ 上的
亚纯函数环；写成公式即 $R(X) = \Gamma(X, \mathcal{K}_X)$。
\end{lemma}

\begin{proof}
由《良态空间》引理 \ref{decent-spaces-lemma-get-reasonable} 及第
\ref{decent-spaces-section-reasonable-decent} 节，$X$ 是良态的
\footnote{反之，若 $X$ 良态，则条件 (c) 自动成立。}。因此，
$X^0 \subset |X|$ 是不可约分支的一般点所成的集合（《良态空间》引理
\ref{decent-spaces-lemma-decent-generic-points}），且由 (b)，$X^0$ 在 $|X|$
中局部有限。于是 $X^0$ 包含在 $|X|$ 的每个稠密开子集中。特别地，$X^0$
包含在概形轨迹中（《良态空间》定理
\ref{decent-spaces-theorem-decent-open-dense-scheme}）。因此，局部环
$\mathcal{O}_{X, \eta}$ 与态射 $j_\eta$ 都有定义。

\medskip\noindent
注意，模层的局部有限直和等于其直积。这个事实连同 $X^0$ 在 $|X|$ 中局部
有限，解释了陈述中直和与直积之间的等式。又由于
$\mathcal{K}_X(\mathcal{F}) =
\mathcal{F} \otimes_{\mathcal{O}_X} \mathcal{K}_X$，第二个等式由第一个等式
得到。

\medskip\noindent
令
$j : Y = \coprod\nolimits_{\eta \in X^0} \Spec(\mathcal{O}_{X, \eta}) \to X$
% P09-ERR-0193
为诸态射 $j_\eta$ 的乘积。我们要证明
$\mathcal{K}_X = j_*\mathcal{O}_Y$。注意，$\mathcal{K}_Y = \mathcal{O}_Y$，
因为 $Y$ 是若干维数为 $0$ 的局部环之谱的不交并：在零维局部环中，任意非零
因子都是单位。接着，由引理
\ref{spaces-divisors-lemma-pullback-meromorphic-sections-defined}，亚纯函数沿 $j$ 的拉回有定义。
这给出映射
$$
\mathcal{K}_X \longrightarrow j_*\mathcal{O}_Y.
$$
设仿射的 $U \in X_\etale$。由引理 \ref{spaces-divisors-lemma-meromorphic-quasi-coherent}，
% P09-ERR-0194
左端的取值是 $\mathcal{O}_X(U)$ 的全分式环。另一方面，右端等于 $U$ 在各余维
$0$ 点（即 $U$ 的一般点）处的局部环之直积。这两个环相等（正如我们已在引理
\ref{spaces-divisors-lemma-meromorphic-quasi-coherent} 的证明中看到的），这是由《代数》引理
\ref{algebra-lemma-total-ring-fractions-no-embedded-points} 及
\ref{algebra-lemma-weakly-ass-zero-divisors} 得到的。因此，我们的映射是同构。

\medskip\noindent
最后，必须证明 $R(X) = \Gamma(X, \mathcal{K}_X)$。把概形情形（《除子》引理
\ref{divisors-lemma-meromorphic-weakass-finite}）应用于概形轨迹
$X' \subset X$ 即得。确切地说，由于它是 $X$ 的稠密开子空间（见上文），
按照定义，$X$ 的有理函数环与 $X'$ 的有理函数环相同。当把 $X'$ 看作概形时，
$R(X')$ 当然与其有理函数环一致。另一方面，由上面对 $\mathcal{K}_X$ 的描述，
以及上面所见的 $X^0 \subset |X'|$ 包含在每个稠密开集中的事实，可知
$\Gamma(X, \mathcal{K}_X) = \Gamma(X', \mathcal{K}_{X'})$。最后使用引理
\ref{spaces-divisors-lemma-meromorphic-quasi-coherent-agree} 中记录的相容性。
\end{proof}

\begin{definition}
\label{spaces-divisors-definition-regular-meromorphic-section}
设 $S$ 为概形，$X$ 为 $S$ 上的代数空间，并设 $\mathcal{L}$ 为可逆
$\mathcal{O}_X$-模。若诱导映射
$\mathcal{K}_X \to \mathcal{K}_X(\mathcal{L})$ 是单射，则称 $\mathcal{L}$ 的
亚纯截面 $s$ 是{\it 正则的}。
\end{definition}

\noindent
下面详细说明（正则）亚纯截面何时可以拉回。

\begin{lemma}
\label{spaces-divisors-lemma-meromorphic-sections-pullback}
设 $S$ 为概形，$f : X \to Y$ 为 $S$ 上代数空间的态射。假设亚纯函数沿 $f$
的拉回有定义（见定义 \ref{spaces-divisors-definition-pullback-meromorphic-sections}）。
\begin{enumerate}
\item 设 $\mathcal{F}$ 为 $\mathcal{O}_Y$-模层。存在亚纯截面的典范拉回映射
$f^* : \Gamma(Y, \mathcal{K}_Y(\mathcal{F})) \to
\Gamma(X, \mathcal{K}_X(f^*\mathcal{F}))$
，其定义在 $\mathcal{F}$ 的亚纯截面上。
% P09-ERR-0195
\item 设 $\mathcal{L}$ 为可逆 $\mathcal{O}_X$-模。$\mathcal{L}$ 的正则亚纯
截面 $s$ 拉回为 $f^*\mathcal{L}$ 的正则亚纯截面 $f^*s$。
\end{enumerate}
\end{lemma}

\begin{proof}
略。
\end{proof}

\begin{lemma}
\label{spaces-divisors-lemma-regular-meromorphic-section-exists}
设 $S$ 为概形，$X$ 为 $S$ 上的代数空间，且满足引理
\ref{spaces-divisors-lemma-compute-meromorphic} 的 (a)、(b)、(c)。则每个可逆
$\mathcal{O}_X$-模 $\mathcal{L}$ 都有正则亚纯截面。
\end{lemma}

\begin{proof}
使用引理 \ref{spaces-divisors-lemma-compute-meromorphic} 的记号。对所有
% P09-ERR-0196
$\eta \in X^0$，$\mathcal{L}$ 的茎 $\mathcal{L}_\eta$ 都有定义，并且它是秩
$1$ 的自由 $\mathcal{O}_{X, \eta}$-模。对每个 $\eta \in X^0$ 选取生成元
$s_\eta \in \mathcal{L}_\eta$。由引理 \ref{spaces-divisors-lemma-compute-meromorphic} 中对
$\mathcal{K}_X$ 与 $\mathcal{K}_X(\mathcal{L})$ 的描述，立即可知
$s = \prod s_\eta$ 是 $\mathcal{L}$ 的正则亚纯截面。
\end{proof}












\section{相对 Proj}
\label{spaces-divisors-section-relative-proj}

\noindent
本节在代数空间的框架下重新讨论相对 Proj 的构造。本节内容对应于概形情形下
《构造》第 \ref{constructions-section-relative-proj} 节及《除子》第
\ref{divisors-section-relative-proj} 节的内容。

\begin{situation}
\label{spaces-divisors-situation-relative-proj}
这里，$S$ 是概形，$X$ 是 $S$ 上的代数空间，而 $\mathcal{A}$ 是拟凝聚分次
$\mathcal{O}_X$-代数。
\end{situation}

\noindent
在情形 \ref{spaces-divisors-situation-relative-proj} 中，我们将定义函子
$F : (\Sch/S)_{fppf}^{opp} \to \textit{Sets}$，它最终将是一个代数空间。
我们作适当修改后沿用《构造》第
\ref{constructions-section-relative-proj} 节的步骤。首先，给定 $S$ 上的概形
$T$，定义一个{\it $T$ 上的四元组}为系统
$(d, f : T \to X, \mathcal{L}, \psi)$
\begin{enumerate}
\item $d \geq 1$ 是整数；
\item $f : T \to X$ 是 $S$ 上的态射；
\item $\mathcal{L}$ 是可逆 $\mathcal{O}_T$-模；并且
\item
$\psi : f^*\mathcal{A}^{(d)} \to \bigoplus_{n \geq 0}\mathcal{L}^{\otimes n}$
是分次 $\mathcal{O}_T$-代数同态，且
$f^*\mathcal{A}_d \to \mathcal{L}$ 是满的。
\end{enumerate}
称两个四元组 $(d, f, \mathcal{L}, \psi)$ 与
$(d', f', \mathcal{L}', \psi')$ {\it 等价}\footnote{此定义受到《构造》引理
\ref{constructions-lemma-equivalent-relative} 的启发。采用这一定义的优点是，
它显然定义了一个等价关系。}，当且仅当 $f = f'$，并且对某个正整数
$m = ad = a'd'$，存在同构
$\beta : \mathcal{L}^{\otimes a} \to (\mathcal{L}')^{\otimes a'}$
，使得 $\beta \circ \psi|_{f^*\mathcal{A}^{(m)}}$ 与
$\psi'|_{f^*\mathcal{A}^{(m)}}$ 作为如下分次环同态相同：
% P09-ERR-0197
$f^*\mathcal{A}^{(m)} \to \bigoplus_{n \geq 0} (\mathcal{L}')^{\otimes mn}$。
给定四元组 $(d, f, \mathcal{L}, \psi)$ 与态射 $h : T' \to T$，其拉回为
$(d, f \circ h, h^*\mathcal{L}, h^*\psi)$。拉回保持此等价关系。最后，对 $S$
上的{\it 拟紧}概形 $T$，令
$$
F(T) = \text{四元组等价类的集合，其基底为 }T
$$
；对 $S$ 上任意概形 $T$，令
$$
F(T)
=
\lim_{V \subset T\text{ 的拟紧开子集}} F(V).
$$
换言之，$F(T)$ 的元素 $\xi$ 对应于一族相容的元素选择
$\xi_V \in F(V)$，其中 $V$ 遍历 $T$ 的拟紧开子集。因此，我们定义了函子
\begin{equation}
\label{spaces-divisors-equation-proj}
% P09-ERR-0198
F : \Sch^{opp} \longrightarrow \textit{Sets}
\end{equation}
存在函子态射 $F \to X$，它把四元组 $(d, f, \mathcal{L}, \psi)$ 映到 $f$。

\begin{lemma}
\label{spaces-divisors-lemma-relative-proj}
% P09-ERR-0199
在情形 \ref{spaces-divisors-situation-relative-proj} 中，上述函子 $F$ 是代数空间。对任意态射
$g : Z \to X$，其中 $Z$ 是概形，存在与进一步基变换相容的典范同构
$\underline{\text{Proj}}_Z(g^*\mathcal{A}) = Z \times_X F$
。
\end{lemma}

\begin{proof}
只需证明第二个断言，见《空间》引理
\ref{spaces-lemma-representable-over-space}。设 $g : Z \to X$ 为态射，其中 $Z$
是概形。令 $F'$ 为与分次拟凝聚 $\mathcal{O}_Z$-代数
$g^*\mathcal{A}$ 相伴的四元组函子。于是存在典范同构
$F' = Z \times_X F$，它把 $F'$ 的四元组
$(d, f : T \to Z, \mathcal{L}, \psi)$ 映到
$(d, g \circ f, \mathcal{L}, \psi)$（略去细节，见《构造》引理
\ref{constructions-lemma-proj-base-change} 的证明）。由《构造》诸引理
\ref{constructions-lemma-equivalent-relative}、
\ref{constructions-lemma-relative-proj}、
% P09-ERR-0200
\ref{constructions-lemma-glueing-gives-functor-proj} 以及定义
\ref{constructions-definition-relative-proj}
可知，$F'$ 由 $\underline{\text{Proj}}_Z(g^*\mathcal{A})$ 表示。
\end{proof}

\noindent
上述引理说明以下定义是有意义的。

\begin{definition}
\label{spaces-divisors-definition-relative-proj}
设 $S$ 为概形，$X$ 为 $S$ 上的代数空间，并设 $\mathcal{A}$ 为拟凝聚分次
$\mathcal{O}_X$-代数层。$\mathcal{A}$ 在 $X$ 上的{\it 相对齐次谱}，或
$\mathcal{A}$ 在 $X$ 上的{\it 齐次谱}，或 $\mathcal{A}$ 在 $X$ 上的
{\it 相对 Proj}，是引理 \ref{spaces-divisors-lemma-relative-proj} 中 $X$ 上的代数空间 $F$。
将其记作 $\pi : \underline{\text{Proj}}_X(\mathcal{A}) \to X$。
\end{definition}

\noindent
特别地，由构造，相对 Proj 的结构态射是可表的。也可以通过黏合来理解相对
Proj。设 $\varphi : U \to X$ 为满平展态射，其中 $U$ 是概形。令
$R = U \times_X U$，其投影态射为 $s, t : R  \to U$。由引理
\ref{spaces-divisors-lemma-relative-proj}，存在 $U$ 上的典范同构
$$
\gamma : 
\underline{\text{Proj}}_U(\varphi^*\mathcal{A})
\longrightarrow
\underline{\text{Proj}}_X(\mathcal{A}) \times_X U
$$
。令 $\alpha : t^*\varphi^*\mathcal{A} \to s^*\varphi^*\mathcal{A}$ 为
《空间的性质》命题
\ref{spaces-properties-proposition-quasi-coherent} 中的典范同构。则图
$$
\xymatrix{
&
\underline{\text{Proj}}_U(\varphi^*\mathcal{A}) \times_{U, s} R
\ar@{=}[r] &
\underline{\text{Proj}}_R(s^*\varphi^*\mathcal{A})
\ar[dd]_{\text{诱导自 }\alpha} \\
\underline{\text{Proj}}_X(\mathcal{A}) \times_X R
\ar[ru]_{s^*\gamma} \ar[rd]^{t^*\gamma} \\
&
\underline{\text{Proj}}_U(\varphi^*\mathcal{A}) \times_{U, t} R
\ar@{=}[r] &
\underline{\text{Proj}}_R(t^*\varphi^*\mathcal{A})
}
$$
交换（等号来自《构造》引理
\ref{constructions-lemma-relative-proj-base-change}）。因此，若用
$\mathcal{A}_U$、$\mathcal{A}_R$ 表示 $\mathcal{A}$ 到 $U$、$R$ 的拉回，
% P09-ERR-0201
则 $P = \underline{\text{Proj}}_X(\mathcal{A})$ 有由概形
$P_U = \underline{\text{Proj}}_U(\mathcal{A}_U)$ 给出的平展覆盖，并且
$P_U \times_P P_U$ 等于
$P_R = \underline{\text{Proj}}_R(\mathcal{A}_R)$。利用这些说明，可以按通常
方式使用平展局部化，把概形情形中相对 Proj 的结果转移到代数空间情形。

\begin{lemma}
\label{spaces-divisors-lemma-twists-of-structure-sheaf}
在情形 \ref{spaces-divisors-situation-relative-proj} 中，相对 Proj 带有拟凝聚
$\mathbf{Z}$-分次代数层
$\bigoplus_{n \in \mathbf{Z}}
\mathcal{O}_{\underline{\text{Proj}}_X(\mathcal{A})}(n)$
及分次代数的典范同态
$$
\psi :
\pi^*\mathcal{A}
\longrightarrow
\bigoplus\nolimits_{n \geq 0}
\mathcal{O}_{\underline{\text{Proj}}_X(\mathcal{A})}(n)
$$
；它们基变换到 $X$ 上任意概形后，与《构造》引理
\ref{constructions-lemma-glue-relative-proj-twists} 相符。
\end{lemma}

\begin{proof}
如定义 \ref{spaces-divisors-definition-relative-proj} 后的讨论，选取概形 $U$ 及满平展态射
$U \to X$，令 $R = U \times_X U$，其投影为 $s, t : R \to U$，并令
$\mathcal{A}_U = \mathcal{A}|_U$、$\mathcal{A}_R = \mathcal{A}|_R$，以及
$\pi : P = \underline{\text{Proj}}_X(\mathcal{A}) \to X$、
$\pi_U : P_U = \underline{\text{Proj}}_U(\mathcal{A}_U)$ 以及
% P09-ERR-0202
$\pi_R : P_R = \underline{\text{Proj}}_U(\mathcal{A}_R)$。
% P09-ERR-0203
由《构造》引理 \ref{constructions-lemma-glue-relative-proj-twists}，有拟凝聚
$\mathbf{Z}$-分次 $\mathcal{O}_{P_U}$-代数层
$\bigoplus_{n \in \mathbf{Z}} \mathcal{O}_{P_U}(n)$
以及典范映射
$\psi_U : \pi_U^*\mathcal{A}_U \to \bigoplus_{n \geq 0} \mathcal{O}_{P_U}(n)$
；对 $P_R$ 亦然。由《构造》引理
\ref{constructions-lemma-relative-proj-base-change}，经任一投影
$P_R \to P_U$ 拉回 $\mathcal{O}_{P_U}(n)$ 与 $\psi_U$，分别得到
$\mathcal{O}_{P_R}(n)$ 与 $\psi_R$。由《空间的性质》命题
\ref{spaces-properties-proposition-quasi-coherent}
得到 $\mathcal{O}_{P}(n)$ 与 $\psi$。略去其与拉回到 $X$ 上任意概形相容的验证。
\end{proof}

\noindent
构造出相对 Proj 后，下面讨论它的一些基本性质。

\begin{lemma}
\label{spaces-divisors-lemma-relative-proj-base-change}
设 $S$ 为概形，$g : X' \to X$ 为 $S$ 上代数空间的态射，并设
$\mathcal{A}$ 为拟凝聚分次 $\mathcal{O}_X$-代数层。则存在典范同构
$$
r :
\underline{\text{Proj}}_{X'}(g^*\mathcal{A})
\longrightarrow
X' \times_X \underline{\text{Proj}}_X(\mathcal{A})
$$
以及相应的同构
$$
\theta :
r^*\text{pr}_2^*\left(\bigoplus\nolimits_{d \in \mathbf{Z}}
\mathcal{O}_{\underline{\text{Proj}}_X(\mathcal{A})}(d)\right)
\longrightarrow
\bigoplus\nolimits_{d \in \mathbf{Z}}
\mathcal{O}_{\underline{\text{Proj}}_{X'}(g^*\mathcal{A})}(d)
$$
，它是 $\mathbf{Z}$-分次
$\mathcal{O}_{\underline{\text{Proj}}_{X'}(g^*\mathcal{A})}$-代数的同构。
\end{lemma}

\begin{proof}
令 $F$ 为函子 (\ref{spaces-divisors-equation-proj})，并令 $F'$ 为在 $X'$ 上用
$g^*\mathcal{A}$ 定义的相应函子。我们断言存在函子的典范同构
$r : F' \to X' \times_X F$（当然，$r$ 正是本引理中的同构）。只需对 $S$ 上的
拟紧概形 $T$ 构造双射
$r : F'(T) \to X'(T) \times_{X(T)} F(T)$。首先，若
$\xi = (d', f', \mathcal{L}', \psi')$ 是 $F'$ 的一个 $T$ 上四元组，则可令
$r(\xi) = (f', (d', g \circ f', \mathcal{L}', \psi'))$。这是有意义的，因为
% P09-ERR-0204
$(g \circ f')^*\mathcal{A}^{(d)} = (f')^*(g^*\mathcal{A})^{(d)}$。
逆映射把配对 $(f', (d, f, \mathcal{L}, \psi))$ 映到四元组
$(d, f', \mathcal{L}, \psi)$。略去最后一个断言的证明（提示：通过平展局部化
约化到概形情形，并应用《构造》引理
\ref{constructions-lemma-relative-proj-base-change}）。
\end{proof}

\begin{lemma}
\label{spaces-divisors-lemma-relative-proj-separated}
在情形 \ref{spaces-divisors-situation-relative-proj} 中，态射
$\pi : \underline{\text{Proj}}_X(\mathcal{A}) \to X$
是分离的。
\end{lemma}

\begin{proof}
由《空间的态射》引理 \ref{spaces-morphisms-lemma-separated-local} 及相对 Proj
的构造，这归结为概形情形，即《构造》引理
\ref{constructions-lemma-relative-proj-separated}。
\end{proof}

\begin{lemma}
\label{spaces-divisors-lemma-relative-proj-quasi-compact}
在情形 \ref{spaces-divisors-situation-relative-proj} 中，若下列条件之一成立：
\begin{enumerate}
\item $\mathcal{A}$ 作为 $\mathcal{A}_0$-代数层是有限型的；
\item $\mathcal{A}$ 作为 $\mathcal{A}_0$-代数由 $\mathcal{A}_1$ 生成，且
$\mathcal{A}_1$ 是有限型 $\mathcal{A}_0$-模；
\item 存在有限型拟凝聚 $\mathcal{A}_0$-子模
$\mathcal{F} \subset \mathcal{A}_{+}$，使得
$\mathcal{A}_{+}/\mathcal{F}\mathcal{A}$ 是
$\mathcal{A}/\mathcal{F}\mathcal{A}$ 的局部幂零理想层；
\end{enumerate}
则 $\pi : \underline{\text{Proj}}_X(\mathcal{A}) \to X$ 是拟紧的。
\end{lemma}

\begin{proof}
由《空间的态射》引理
\ref{spaces-morphisms-lemma-quasi-compact-local} 及相对 Proj 的构造，这归结为
概形情形，即《除子》引理
\ref{divisors-lemma-relative-proj-quasi-compact}。
\end{proof}

\begin{lemma}
\label{spaces-divisors-lemma-relative-proj-finite-type}
在情形 \ref{spaces-divisors-situation-relative-proj} 中，若 $\mathcal{A}$ 作为
$\mathcal{O}_X$-代数层是有限型的，则
$\pi : \underline{\text{Proj}}_X(\mathcal{A}) \to X$ 是有限型的。
\end{lemma}

\begin{proof}
由《空间的态射》引理
\ref{spaces-morphisms-lemma-finite-type-local} 及相对 Proj 的构造，这归结为
概形情形，即《除子》引理
\ref{divisors-lemma-relative-proj-finite-type}。
\end{proof}

\begin{lemma}
\label{spaces-divisors-lemma-relative-proj-universally-closed}
在情形 \ref{spaces-divisors-situation-relative-proj} 中，若
$\mathcal{O}_X \to \mathcal{A}_0$
是整代数映射\footnote{换言之，$\mathcal{O}_X$ 在 $\mathcal{A}_0$ 中的整闭包
（见《空间的态射》定义
\ref{spaces-morphisms-definition-integral-closure}）等于
$\mathcal{A}_0$。}，且 $\mathcal{A}$ 作为 $\mathcal{A}_0$-代数是有限型的，则
$\pi : \underline{\text{Proj}}_X(\mathcal{A}) \to X$ 是普遍闭的。
\end{lemma}

\begin{proof}
由《空间的态射》引理
\ref{spaces-morphisms-lemma-universally-closed-local} 及相对 Proj 的构造，这归结为
概形情形，即《除子》引理
\ref{divisors-lemma-relative-proj-universally-closed}。
\end{proof}

\begin{lemma}
\label{spaces-divisors-lemma-relative-proj-proper}
在情形 \ref{spaces-divisors-situation-relative-proj} 中，下列条件等价：
\begin{enumerate}
\item $\mathcal{A}_0$ 是有限型 $\mathcal{O}_X$-模，且 $\mathcal{A}$ 作为
$\mathcal{A}_0$-代数是有限型的；
\item $\mathcal{A}_0$ 是有限型 $\mathcal{O}_X$-模，且 $\mathcal{A}$ 作为
$\mathcal{O}_X$-代数是有限型的。
\end{enumerate}
若这些条件成立，则
$\pi : \underline{\text{Proj}}_X(\mathcal{A}) \to X$
是固有的。
\end{lemma}

\begin{proof}
由《空间的态射》引理
\ref{spaces-morphisms-lemma-proper-local} 及相对 Proj 的构造，这归结为概形
情形，即《除子》引理
% P09-ERR-0205
\ref{divisors-lemma-relative-proj-universally-closed}。
\end{proof}

\begin{lemma}
\label{spaces-divisors-lemma-relative-proj-generated-in-degree-1}
设 $S$ 为概形，$X$ 为 $S$ 上的代数空间。设 $\mathcal{A}$ 为拟凝聚分次
$\mathcal{O}_X$-模层，
% P09-ERR-0206
并由 $\mathcal{A}_1$ 作为 $\mathcal{A}_0$-代数生成。令
$P = \underline{\text{Proj}}_X(\mathcal{A})$，则：
\begin{enumerate}
\item $P$ 表示函子 $F_1$；该函子把 $S$ 上的 $T$ 映到三元组
$(f, \mathcal{L}, \psi)$ 的同构类所成的集合，其中 $f : T \to X$ 是 $S$ 上
的态射，$\mathcal{L}$ 是可逆 $\mathcal{O}_T$-模，而
$\psi : f^*\mathcal{A} \to \bigoplus_{n \geq 0} \mathcal{L}^{\otimes n}$
是分次 $\mathcal{O}_T$-代数映射，并诱导满射
$f^*\mathcal{A}_1 \to \mathcal{L}$；
\item 典范映射 $\pi^*\mathcal{A}_1 \to \mathcal{O}_P(1)$ 是满的；并且
\item 每个 $\mathcal{O}_P(n)$ 都可逆，且乘法映射诱导同构
$\mathcal{O}_P(n) \otimes_{\mathcal{O}_P} \mathcal{O}_P(m) =
\mathcal{O}_P(n + m)$。
\end{enumerate}
\end{lemma}

\begin{proof}
略。概形情形见《构造》引理
\ref{constructions-lemma-apply-relative}。
\end{proof}







\section{相对 Proj 的函子性}
\label{spaces-divisors-section-functoriality-relative-proj}

\noindent
本节对应于《构造》第
\ref{constructions-section-functoriality-relative-proj} 节。

\begin{lemma}
\label{spaces-divisors-lemma-morphism-relative-proj}
设 $S$ 为概形，$X$ 为 $S$ 上的代数空间。设
$\psi : \mathcal{A} \to \mathcal{B}$ 为拟凝聚分次
$\mathcal{O}_X$-代数的映射。令
$P = \underline{\text{Proj}}_X(\mathcal{A}) \to X$ 以及
$Q = \underline{\text{Proj}}_X(\mathcal{B}) \to X$。
存在典范开子空间 $U(\psi) \subset Q$ 以及 $X$ 上代数空间的典范态射
$$
r_\psi :
U(\psi)
\longrightarrow
P
$$
和 $\mathbf{Z}$-分次 $\mathcal{O}_{U(\psi)}$-代数的映射
$$
\theta = \theta_\psi :
r_\psi^*\left(
\bigoplus\nolimits_{d \in \mathbf{Z}} \mathcal{O}_P(d)
\right)
\longrightarrow
\bigoplus\nolimits_{d \in \mathbf{Z}} \mathcal{O}_{U(\psi)}(d).
$$
三元组 $(U(\psi), r_\psi, \theta)$ 由下述性质刻画：对任意平展于 $X$ 的概形
$W$，三元组
$$
(U(\psi) \times_X W,\quad
r_\psi|_{U(\psi) \times_X W} : U(\psi) \times_X W \to  P \times_X W,\quad
\theta|_{U(\psi) \times_X W})
$$
等于《构造》引理 \ref{constructions-lemma-morphism-relative-proj} 中与
$\psi : \mathcal{A}|_W \to \mathcal{B}|_W$ 相伴的三元组。
\end{lemma}

\begin{proof}
本引理由平展局部化及概形情形得到，见定义
\ref{spaces-divisors-definition-relative-proj} 后的讨论。略去细节。
\end{proof}

\begin{lemma}
\label{spaces-divisors-lemma-morphism-relative-proj-transitive}
设 $S$ 为概形，$X$ 为 $S$ 上的代数空间。设 $\mathcal{A}$、$\mathcal{B}$、
$\mathcal{C}$ 为拟凝聚分次 $\mathcal{O}_X$-代数。令
$P = \underline{\text{Proj}}_X(\mathcal{A})$、
$Q = \underline{\text{Proj}}_X(\mathcal{B})$ 及
$R = \underline{\text{Proj}}_X(\mathcal{C})$。设
$\varphi : \mathcal{A} \to \mathcal{B}$、
$\psi : \mathcal{B} \to \mathcal{C}$ 为分次 $\mathcal{O}_X$-代数映射。则有
% P09-ERR-0207
$$
U(\psi \circ \varphi) = r_\varphi^{-1}(U(\psi))
\quad
\text{且}
\quad
r_{\psi \circ \varphi}
=
r_\varphi \circ r_\psi|_{U(\psi \circ \varphi)}.
$$
此外有
$$
\theta_\psi \circ r_\psi^*\theta_\varphi
=
\theta_{\psi \circ \varphi}
$$
，其中采用显然的记号。
\end{lemma}

\begin{proof}
略。
\end{proof}

\begin{lemma}
\label{spaces-divisors-lemma-surjective-graded-rings-map-relative-proj}
% P09-ERR-0209
沿用上述引理 \ref{spaces-divisors-lemma-morphism-relative-proj} 的假设与记号。假设当
$d \gg 0$ 时，$\mathcal{A}_d \to \mathcal{B}_d$ 是满的。则：
\begin{enumerate}
\item $U(\psi) = Q$；
% P09-ERR-0208
\item $r_\psi : Q \to R$ 是闭浸入；并且
\item 映射 $\theta : r_\psi^*\mathcal{O}_P(n) \to \mathcal{O}_Q(n)$
是满的，但一般不是同构（即使
$\mathcal{A} \to \mathcal{B}$ 是满的）。
\end{enumerate}
\end{lemma}

\begin{proof}
通过平展局部化，这由概形情形（《构造》引理
\ref{constructions-lemma-surjective-graded-rings-map-relative-proj}）得到。
\end{proof}

\begin{lemma}
\label{spaces-divisors-lemma-eventual-iso-graded-rings-map-relative-proj}
% P09-ERR-0209
沿用上述引理 \ref{spaces-divisors-lemma-morphism-relative-proj} 的假设与记号。假设对所有
$d \gg 0$，$\mathcal{A}_d \to \mathcal{B}_d$ 都是同构。则：
\begin{enumerate}
\item $U(\psi) = Q$；
\item $r_\psi : Q \to P$ 是同构；并且
\item 映射 $\theta : r_\psi^*\mathcal{O}_P(n) \to \mathcal{O}_Q(n)$
都是同构。
\end{enumerate}
\end{lemma}

\begin{proof}
通过平展局部化，这由概形情形（《构造》引理
\ref{constructions-lemma-eventual-iso-graded-rings-map-relative-proj}）得到。
\end{proof}

\begin{lemma}
\label{spaces-divisors-lemma-surjective-generated-degree-1-map-relative-proj}
% P09-ERR-0209
沿用上述引理 \ref{spaces-divisors-lemma-morphism-relative-proj} 的假设与记号。假设当
$d \gg 0$ 时，$\mathcal{A}_d \to \mathcal{B}_d$ 是满的，并且
$\mathcal{A}$ 在 $\mathcal{A}_0$ 上由 $\mathcal{A}_1$ 生成。则：
\begin{enumerate}
\item $U(\psi) = Q$；
\item $r_\psi : Q \to P$ 是闭浸入；并且
\item 映射 $\theta : r_\psi^*\mathcal{O}_P(n) \to \mathcal{O}_Q(n)$
都是同构。
\end{enumerate}
\end{lemma}

\begin{proof}
通过平展局部化，这由概形情形（《构造》引理
\ref{constructions-lemma-surjective-generated-degree-1-map-relative-proj}）得到。
\end{proof}










\section{可逆层与到相对 Proj 的态射}
\label{spaces-divisors-section-invertible-relative-proj}

\noindent
看来我们可能会在某处用到下面这个引理。
情形如下：
\begin{enumerate}
\item 设 $S$ 为概形，$Y$ 为 $S$ 上的代数空间。
\item 设 $\mathcal{A}$ 为拟凝聚分次 $\mathcal{O}_Y$-代数。
\item 以 $\pi : \underline{\text{Proj}}_Y(\mathcal{A}) \to Y$ 表示
$\mathcal{A}$ 在 $Y$ 上的相对 Proj。
\item 设 $f : X \to Y$ 为 $S$ 上代数空间的态射。
\item 设 $\mathcal{L}$ 为可逆 $\mathcal{O}_X$-模。
\item 设
$\psi : f^*\mathcal{A} \to \bigoplus_{d \geq 0} \mathcal{L}^{\otimes d}$
为分次 $\mathcal{O}_X$-代数的同态。
\end{enumerate}
给定这些数据，令 $U(\psi) \subset X$ 为满足下式的开子空间：
$$
|U(\psi)| = \bigcup\nolimits_{d \geq 1}
\{\text{满足下述条件的轨迹：}f^*\mathcal{A}_d \to \mathcal{L}^{\otimes d}
\text{ 为满射}\}
$$
$U(\psi) \subset X$ 的构造与沿任意态射 $X' \to X$ 的拉回可交换。

\begin{lemma}
\label{spaces-divisors-lemma-invertible-map-into-relative-proj}
% P09-ERR-0210
采用上述假设与记号。态射 $\psi$ 诱导一个 $Y$ 上代数空间的典范态射
$$
r_{\mathcal{L}, \psi} :
U(\psi) \longrightarrow \underline{\text{Proj}}_Y(\mathcal{A})
$$
，并伴随一个分次 $\mathcal{O}_{U(\psi)}$-代数映射
$$
\theta :
r_{\mathcal{L}, \psi}^*\left(
\bigoplus\nolimits_{d \geq 0}
\mathcal{O}_{\underline{\text{Proj}}_Y(\mathcal{A})}(d)
\right)
\longrightarrow
\bigoplus\nolimits_{d \geq 0} \mathcal{L}^{\otimes d}|_{U(\psi)}
$$
，它由下列性质刻画：
\begin{enumerate}
\item 对平展态射 $V \to Y$ 和 $d \geq 0$，下图
% P09-ERR-0211
$$
\xymatrix{
\mathcal{A}_d(V) \ar[d]_{\psi} \ar[r]_{\psi} &
\Gamma(V \times_Y X, \mathcal{L}^{\otimes d}) \ar[d]^{restrict} \\
\Gamma(V \times_Y \underline{\text{Proj}}_Y(\mathcal{A}),
\mathcal{O}_{\underline{\text{Proj}}_Y(\mathcal{A})}(d)) \ar[r]^-\theta &
\Gamma(V \times_Y U(\psi), \mathcal{L}^{\otimes d})
}
$$
可交换。
\item 对任意 $d \geq 1$ 以及任意态射 $W \to X$，其中 $W$ 为概形，若
$\psi|_W : f^*\mathcal{A}_d|_W \to \mathcal{L}^{\otimes d}|_W$ 为满射，
则 (a) $W \to X$ 分解经过 $U(\psi)$，且 (b) $W \to U(\psi)$ 与
$r_{\mathcal{L}, \psi}$ 的复合等同于态射
$W \to \underline{\text{Proj}}_Y(\mathcal{A})$；后者因
$\underline{\text{Proj}}_Y(\mathcal{A})$ 的构造而存在，见定义
\ref{spaces-divisors-definition-relative-proj}。
\item 考虑可交换图
$$
\xymatrix{
X' \ar[r]_{g'} \ar[d]_{f'} & X \ar[d]^f \\
Y' \ar[r]^g & Y
}
$$
，其中 $X'$ 与 $Y'$ 为概形。置 $\mathcal{A}' = g^*\mathcal{A}$ 及
$\mathcal{L}' = (g')^*\mathcal{L}$，并以
% P09-ERR-0212
$\psi' : (f')^*\mathcal{A} \to \bigoplus_{d \geq 0} (\mathcal{L}')^{\otimes d}$
表示 $\psi$ 的拉回。令 $U(\psi')$、
% P09-ERR-0213
$r_{\psi', \mathcal{L}'}$ 及 $\theta'$ 为
% P09-ERR-0214
《构造》引理 \ref{spaces-divisors-lemma-invertible-map-into-relative-proj} 中所构造的开子空间、
态射及同态。那么 $U(\psi') = (g')^{-1}(U(\psi))$，且在引理
\ref{spaces-divisors-lemma-relative-proj-base-change} 的同构
$\underline{\text{Proj}}_{Y'}(\mathcal{A}') =
Y' \times_Y \underline{\text{Proj}}_Y(\mathcal{A})$
下，$r_{\psi', \mathcal{L}'}$ 等同于 $r_{\psi, \mathcal{L}}$ 的基变换。
此外，$\theta'$ 是 $\theta$ 的拉回。
\end{enumerate}
\end{lemma}

\begin{proof}
略。提示：首先注意，对 $X$ 上的拟紧概形 $W$，下列条件等价：
\begin{enumerate}
\item $W \to X$ 分解经过 $U(\psi)$；以及
% P09-ERR-0215
\item 存在 $d$，使得
$\psi|_W : f^*\mathcal{A}_d|_W \to \mathcal{L}^{\otimes d}|_W$
为满射。
\end{enumerate}
这就把 $U(\psi)$ 描述为基范畴 $(\Sch/S)_{fppf}$ 上 $X$ 的子函子。
对这样的 $W$ 与 $d$，考虑四元组
$(d, W \to Y, \mathcal{L}|_W, \psi^{(d)}|_W)$。
由 $\underline{\text{Proj}}_Y(\mathcal{A})$ 的定义，得到态射
$W \to \underline{\text{Proj}}_Y(\mathcal{A})$。根据四元组的等价概念可见，
此态射与 $d$ 的选取无关。这显然定义了函子之间的变换
$r_{\psi, \mathcal{L}} : U(\psi) \to \underline{\text{Proj}}_Y(\mathcal{A})$，
即代数空间的态射。按构造，此态射满足 (2)。由于《构造》引理
\ref{constructions-lemma-invertible-map-into-relative-proj} 中构造的态射满足同一
性质，故 (3) 成立。

\medskip\noindent
为构造 $\theta$ 并验证引理中相容性 (1)，在 $Y$ 与 $X$ 上平展局部地工作，
论证同定义 \ref{spaces-divisors-definition-relative-proj} 后的讨论。
\end{proof}








\section{相对丰沛层}
\label{spaces-divisors-section-relatively-ample}

\noindent
本节是《态射》第 \ref{morphisms-section-relatively-ample} 节在代数空间情形下的
对应版本。我们如下定义相对丰沛可逆层。

\begin{definition}
\label{spaces-divisors-definition-relatively-ample}
设 $S$ 为概形，$f : X \to Y$ 为 $S$ 上代数空间的态射，
$\mathcal{L}$ 为可逆 $\mathcal{O}_X$-模。若 $f : X \to Y$ 可表，且对每个
从概形 $Z$ 出发的态射 $Z \to Y$，$\mathcal{L}$ 到
$X_Z = Z \times_Y X$ 上的拉回 $\mathcal{L}_Z$ 按照《态射》定义
\ref{morphisms-definition-relatively-ample} 在 $X_Z/Z$ 上丰沛，则称
$\mathcal{L}$ {\it 相对丰沛}，或称为 {\it $f$-相对丰沛}、
{\it 在 $X/Y$ 上丰沛}、或 {\it $f$-丰沛}。
\end{definition}

\noindent
我们几乎总会把有关相对丰沛可逆层的问题归约到概形情形。因此，本节主要包含
一些基本检验。

\begin{lemma}
\label{spaces-divisors-lemma-relatively-ample-sanity-check}
设 $S$ 为概形，$f : X \to Y$ 为 $S$ 上代数空间的态射，
$\mathcal{L}$ 为可逆 $\mathcal{O}_X$-模。假设 $Y$ 是概形。下列条件等价：
\begin{enumerate}
\item $\mathcal{L}$ 在定义 \ref{spaces-divisors-definition-relatively-ample} 的意义下在
$X/Y$ 上丰沛；以及
\item $X$ 是概形，且 $\mathcal{L}$ 在《态射》定义
\ref{morphisms-definition-relatively-ample} 的意义下在 $X/Y$ 上丰沛。
\end{enumerate}
\end{lemma}

\begin{proof}
这由定义及《态射》引理 \ref{morphisms-lemma-ample-base-change} 得出（该引理说明，
对概形而言，相对丰沛性在基变换下保持）。
\end{proof}

\begin{lemma}
\label{spaces-divisors-lemma-ample-base-change}
设 $S$ 为概形，$f : X \to Y$ 为 $S$ 上代数空间的态射，
$\mathcal{L}$ 为可逆 $\mathcal{O}_X$-模。设 $Y' \to Y$ 为 $S$ 上代数空间的
态射，$f' : X' \to Y'$ 为 $f$ 的基变换，并以 $\mathcal{L}'$ 表示
$\mathcal{L}$ 到 $X'$ 上的拉回。若 $\mathcal{L}$ 是 $f$-丰沛的，则
$\mathcal{L}'$ 是 $f'$-丰沛的。
\end{lemma}

\begin{proof}
这立即由定义得出！（提示：使用基变换的传递性。）
\end{proof}

\begin{lemma}
\label{spaces-divisors-lemma-relatively-ample-properties}
设 $S$ 为概形，$f : X \to Y$ 为 $S$ 上代数空间的态射。若存在一个
$f$-丰沛可逆层，则 $f$ 可表、拟紧且分离。
\end{lemma}

\begin{proof}
这由定义及《态射》引理 \ref{morphisms-lemma-relatively-ample-separated} 显然可得。
（若对此有疑问，可参看《代数空间》引理
\ref{spaces-lemma-representable-transformations-property-implication} 的原理。）
\end{proof}

\begin{lemma}
\label{spaces-divisors-lemma-descend-relatively-ample}
设 $V \to U$ 为仿射概形之间的满平展态射，$X$ 为 $U$ 上的代数空间，
$\mathcal{L}$ 为可逆 $\mathcal{O}_X$-模。令 $Y = V \times_U X$，并令
$\mathcal{N}$ 为 $\mathcal{L}$ 到 $Y$ 上的拉回。下列条件等价：
\begin{enumerate}
\item $\mathcal{L}$ 在 $X/U$ 上丰沛；以及
\item $\mathcal{N}$ 在 $Y/V$ 上丰沛。
\end{enumerate}
\end{lemma}

\begin{proof}
蕴含关系 (1) $\Rightarrow$ (2) 由引理 \ref{spaces-divisors-lemma-ample-base-change} 得出。
假设 (2)。这蕴含 $Y \to V$ 拟紧且分离（引理
\ref{spaces-divisors-lemma-relatively-ample-properties}），并且 $Y$ 是概形。因此，态射
$f : X \to U$ 拟紧且分离（《空间的态射》引理
\ref{spaces-morphisms-lemma-quasi-compact-local} 及
\ref{spaces-morphisms-lemma-separated-local}）。置
$\mathcal{A} = \bigoplus_{d \geq 0} f_*\mathcal{L}^{\otimes d}$。
这是拟凝聚分次 $\mathcal{O}_U$-代数层（《空间的态射》引理
\ref{spaces-morphisms-lemma-pushforward}）。由伴随性得到映射
$\psi : f^*\mathcal{A} \to \bigoplus_{d \geq 0} \mathcal{L}^{\otimes d}$。
应用引理 \ref{spaces-divisors-lemma-invertible-map-into-relative-proj}，得到开子空间
$U(\psi) \subset X$ 及态射
$$
r_{\mathcal{L}, \psi} : U(\psi) \to \underline{\text{Proj}}_U(\mathcal{A})
$$
% P09-ERR-0216
由于 $h : V \to U$ 平展，故
$\mathcal{A}|_V = (Y \to V)_*(\bigoplus_{d \geq 0} \mathcal{N}^{\otimes d})$，
见《空间的性质》引理
\ref{spaces-properties-lemma-pushforward-etale-base-change-modules}。因此，
$\psi$ 到 $Y$ 上的拉回 $\psi'$ 是《态射》引理
\ref{morphisms-lemma-characterize-relatively-ample} 第 (5) 部分中情形
$(Y \to V, \mathcal{N})$ 的伴随映射。由于 $\mathcal{N}$ 在 $Y/V$ 上丰沛，
由刚引用的引理可得 $U(\psi') = Y$，且 $r_{\mathcal{N}, \psi'}$ 是开浸入。
引理 \ref{spaces-divisors-lemma-invertible-map-into-relative-proj} 告诉我们，
$r_{\mathcal{L}, \psi}$ 的构造与基变换可交换，故 $U(\psi) = X$，并且有可交换图
$$
\xymatrix{
Y \ar[r]_-{r'} \ar[d] &
\underline{\text{Proj}}_V(\mathcal{A}|_V) \ar[d] \ar[r] &
V \ar[d] \\
X \ar[r]^-r &
\underline{\text{Proj}}_U(\mathcal{A}) \ar[r] &
U
}
$$
，其中两个方块都是纤维积。由《空间的态射》引理
\ref{spaces-morphisms-lemma-closed-immersion-local} 可知 $r$ 是开浸入。因此，
$X$ 是概形。于是可应用《态射》引理
\ref{morphisms-lemma-characterize-relatively-ample} 第 (5) 部分，得出
$\mathcal{L}$ 在 $X/U$ 上丰沛。
\end{proof}

\begin{lemma}
\label{spaces-divisors-lemma-relatively-ample-local}
设 $S$ 为概形，$f : X \to Y$ 为 $S$ 上代数空间的态射，
$\mathcal{L}$ 为可逆 $\mathcal{O}_X$-模。下列条件等价：
\begin{enumerate}
\item $\mathcal{L}$ 在 $X/Y$ 上丰沛；
\item 对每个概形 $Z$ 及每个态射 $Z \to Y$，代数空间
$X_Z = Z \times_Y X$ 都是概形，且拉回 $\mathcal{L}_Z$ 在 $X_Z/Z$ 上丰沛；
\item 对每个仿射概形 $Z$ 及每个态射 $Z \to Y$，代数空间
$X_Z = Z \times_Y X$ 都是概形，且拉回 $\mathcal{L}_Z$ 在 $X_Z/Z$ 上丰沛；
\item 存在概形 $V$ 及满平展态射 $V \to Y$，使得代数空间
$X_V = V \times_Y X$ 是概形，且拉回 $\mathcal{L}_V$ 在 $X_V/V$ 上丰沛。
\end{enumerate}
\end{lemma}

\begin{proof}
按定义，(1) 与 (2) 等价。蕴含关系 (2) $\Rightarrow$ (3) 是直接的。若 (3)
成立且 $Z \to Y$ 如 (2) 所述，则可见
% P09-ERR-0217
$X_Z \to Z$ 在 $Z$ 上仿射局部地可表。因此，例如由《空间的性质》引理
\ref{spaces-properties-lemma-subscheme}，$X_Z$ 是概形。于是
$\mathcal{L}_Z$ 在 $X_Z/Z$ 上丰沛，因为这在 $Z$ 上是局部的，并且可以使用
《态射》引理 \ref{morphisms-lemma-characterize-relatively-ample}。故 (1)、(2)、
(3) 等价。显然，这些条件蕴含 (4)。

\medskip\noindent
假设 (4)。设 $Z \to Y$ 为态射，其中 $Z$ 仿射。则
$U = V \times_Y Z \to Z$ 是满平展态射，且 $\mathcal{L}_Z$ 沿
$X_U \to X_Z$ 的拉回在 $X_U/U$ 上相对丰沛。
% P09-ERR-0218
当然，我们可以用一个仿射开集替换 $U$。于是由引理
\ref{spaces-divisors-lemma-descend-relatively-ample}，$\mathcal{L}_Z$ 在 $X_Z/Z$ 上丰沛。
故 (4) $\Rightarrow$ (3)，证明完毕。
\end{proof}

\begin{lemma}
\label{spaces-divisors-lemma-quasi-affine-O-ample}
设 $S$ 为概形，$f : X \to Y$ 为 $S$ 上代数空间的态射。则 $f$ 拟仿射，
当且仅当 $\mathcal{O}_X$ 是 $f$-相对丰沛的。
\end{lemma}

\begin{proof}
这由概形情形得出，见《态射》引理
\ref{morphisms-lemma-quasi-affine-O-ample}。
\end{proof}



\section{相对丰沛性与上同调}
\label{spaces-divisors-section-ample-and-proper}

\noindent
本节包含一些与《概形的上同调》第
\ref{coherent-section-applications-formal-functions} 节及第
\ref{coherent-section-ample-cohomology} 节中的结果相关的结论。

\medskip\noindent
下面的引理只是我们所能做的一例。

\begin{lemma}
\label{spaces-divisors-lemma-vanshing-gives-ample}
% P09-ERR-0219
设 $R$ 为诺特环，$X$ 为 $R$ 上的代数空间，并且结构态射
$f : X \to \Spec(R)$ 固有。设 $\mathcal{L}$ 为可逆 $\mathcal{O}_X$-模。
下列条件等价：
\begin{enumerate}
\item $\mathcal{L}$ 在 $X/R$ 上丰沛（定义
\ref{spaces-divisors-definition-relatively-ample}）；
\item 对每个凝聚 $\mathcal{O}_X$-模 $\mathcal{F}$，存在 $n_0 \geq 0$，使得
$H^p(X, \mathcal{F} \otimes \mathcal{L}^{\otimes n}) = 0$
对所有 $n \geq n_0$ 及 $p > 0$ 成立。
\end{enumerate}
\end{lemma}

\begin{proof}
蕴含关系 (1) $\Rightarrow$ (2) 由《概形的上同调》引理
\ref{coherent-lemma-coherent-proper-ample} 得出，因为假设 (1) 蕴含 $X$ 是概形。
蕴含关系 (2) $\Rightarrow$ (1) 即《空间的上同调》引理
\ref{spaces-cohomology-lemma-Noetherian-h1-zero-invertible}。
\end{proof}

\begin{lemma}
\label{spaces-divisors-lemma-ample-on-fibre}
设 $Y$ 为诺特概形，$X$ 为 $Y$ 上的代数空间，并且结构态射
$f : X \to Y$ 固有。设 $\mathcal{L}$ 为可逆 $\mathcal{O}_X$-模，
$\mathcal{F}$ 为凝聚 $\mathcal{O}_X$-模。设 $y \in Y$ 为一点，使得 $X_y$
是概形，且 $\mathcal{L}_y$ 在 $X_y$ 上丰沛。则存在 $d_0$，使得对所有
$d \geq d_0$ 都有
$$
R^pf_*(\mathcal{F} \otimes_{\mathcal{O}_X} \mathcal{L}^{\otimes d})_y = 0
\text{ 对 }p > 0
$$
，并且映射
$$
f_*(\mathcal{F} \otimes_{\mathcal{O}_X} \mathcal{L}^{\otimes d})_y
\longrightarrow
H^0(X_y, \mathcal{F}_y \otimes_{\mathcal{O}_{X_y}} \mathcal{L}_y^{\otimes d})
$$
是满射。
\end{lemma}

\begin{proof}
注意，$\mathcal{O}_{Y, y}$ 是诺特局部环。考虑典范态射
$c : \Spec(\mathcal{O}_{Y, y}) \to Y$，见《概形》公式
(\ref{schemes-equation-canonical-morphism})。由于它诱导局部环的等同，故为平坦
态射。暂以 $f' : X' \to \Spec(\mathcal{O}_{Y, y})$ 表示 $f$ 到该局部环上的
基变换。由《空间的上同调》引理
\ref{spaces-cohomology-lemma-flat-base-change-cohomology} 可见
$c^*R^pf_*\mathcal{F} = R^pf'_*\mathcal{F}'$。此外，纤维 $X_y$ 与 $X'_y$
被等同。因此，可以假设 $Y = \Spec(A)$ 是
诺特局部环 $(A, \mathfrak m, \kappa)$ 的谱，且 $y \in Y$ 对应于
$\mathfrak m$。在此情形下，对所有 $p \geq 0$ 都有
$R^pf_*(\mathcal{F} \otimes_{\mathcal{O}_X} \mathcal{L}^{\otimes d})_y =
H^p(X, \mathcal{F} \otimes_{\mathcal{O}_X} \mathcal{L}^{\otimes d})$
。以 $f_y : X_y \to \Spec(\kappa)$ 表示投影。

\medskip\noindent
令 $B = \text{Gr}_\mathfrak m(A) =
\bigoplus_{n \geq 0} \mathfrak m^n/\mathfrak m^{n + 1}$。考虑拟凝聚分次
$\mathcal{O}_{X_y}$-代数层 $\mathcal{B} = f_y^*\widetilde{B}$。我们将沿用
《空间的上同调》第 \ref{spaces-cohomology-section-theorem-formal-functions} 节的
记号，其中以 $\mathfrak m$ 代替 $I$。由于 $X_y$ 是由
$\mathfrak m\mathcal{O}_X$ 截出的 $X$ 的闭子空间，可以把
$\mathfrak m^n\mathcal{F}/\mathfrak m^{n + 1}\mathcal{F}$ 看作凝聚
$\mathcal{O}_{X_y}$-模，见《空间的上同调》引理
\ref{spaces-cohomology-lemma-i-star-equivalence}。于是
$\bigoplus_{n \geq 0} \mathfrak m^n\mathcal{F}/\mathfrak m^{n + 1}\mathcal{F}$
是有限型拟凝聚分次 $\mathcal{B}$-模，因为它在 $\mathcal{B}$ 上由 0 次部分生成，
% P09-ERR-0220
且其 0 次部分为
$\mathcal{F}_y = \mathcal{F}/\mathfrak m \mathcal{F}$
，这是凝聚 $\mathcal{O}_{X_y}$-模。因此，由《概形的上同调》引理
\ref{coherent-lemma-graded-finiteness} 第 (2) 部分，存在 $d_0$，使得
$$
H^p(X_y, \mathfrak m^n \mathcal{F}/ \mathfrak m^{n + 1}\mathcal{F}
\otimes_{\mathcal{O}_{X_y}} \mathcal{L}_y^{\otimes d}) = 0
$$
对所有 $p > 0$、$d \geq d_0$ 及 $n \geq 0$ 成立。由《空间的上同调》引理
\ref{spaces-cohomology-lemma-relative-affine-cohomology}，这等价于
$
H^p(X, \mathfrak m^n \mathcal{F}/ \mathfrak m^{n + 1}\mathcal{F}
\otimes_{\mathcal{O}_X} \mathcal{L}^{\otimes d}) = 0
$
对所有 $p > 0$、$d \geq d_0$ 及 $n \geq 0$ 成立。

\medskip\noindent
考虑凝聚 $\mathcal{O}_X$-模的短正合列
$$
0 \to \mathfrak m^n\mathcal{F}/\mathfrak m^{n + 1} \mathcal{F}
\to \mathcal{F}/\mathfrak m^{n + 1} \mathcal{F}
\to \mathcal{F}/\mathfrak m^n \mathcal{F} \to 0
$$
。与 $\mathcal{L}^{\otimes d}$ 作张量积是正合函子，因而得到短正合列
$$
0 \to
\mathfrak m^n\mathcal{F}/\mathfrak m^{n + 1} \mathcal{F}
\otimes_{\mathcal{O}_X} \mathcal{L}^{\otimes d}
\to \mathcal{F}/\mathfrak m^{n + 1} \mathcal{F}
\otimes_{\mathcal{O}_X} \mathcal{L}^{\otimes d}
\to \mathcal{F}/\mathfrak m^n \mathcal{F}
\otimes_{\mathcal{O}_X} \mathcal{L}^{\otimes d} \to 0
$$
。利用上同调长正合列及上述消失性，以归纳法得出：
\begin{enumerate}
\item $H^p(X, \mathcal{F}/\mathfrak m^n \mathcal{F}
\otimes_{\mathcal{O}_X} \mathcal{L}^{\otimes d}) = 0$
对所有 $p > 0$、$d \geq d_0$ 及 $n \geq 0$ 成立；以及
\item $H^0(X, \mathcal{F}/\mathfrak m^n \mathcal{F}
\otimes_{\mathcal{O}_X} \mathcal{L}^{\otimes d}) \to
H^0(X_y, \mathcal{F}_y \otimes_{\mathcal{O}_{X_y}} \mathcal{L}_y^{\otimes d})$
对所有 $d \geq d_0$ 及 $n \geq 1$ 都是满射。
\end{enumerate}
由形式函数定理（《空间的上同调》定理
\ref{spaces-cohomology-theorem-formal-functions}），可知
$H^p(X, \mathcal{F} \otimes_{\mathcal{O}_X} \mathcal{L}^{\otimes d})$
的 $\mathfrak m$-进完备化对所有 $d \geq d_0$ 及 $p > 0$ 都为零。
由《空间的上同调》引理
\ref{spaces-cohomology-lemma-proper-over-affine-cohomology-finite}，
$H^p(X, \mathcal{F} \otimes_{\mathcal{O}_X} \mathcal{L}^{\otimes d})$ 是有限
$A$-模，故由 Nakayama 引理（《代数》引理 \ref{algebra-lemma-NAK}），
$H^p(X, \mathcal{F} \otimes_{\mathcal{O}_X} \mathcal{L}^{\otimes d})$
对所有 $d \geq d_0$ 及 $p > 0$ 都为零。对 $p = 0$，由《空间的上同调》引理
\ref{spaces-cohomology-lemma-ML-cohomology-powers-ideal} 第 (3) 部分可知
$H^0(X, \mathcal{F} \otimes_{\mathcal{O}_X} \mathcal{L}^{\otimes d}) \to
H^0(X_y, \mathcal{F}_y \otimes_{\mathcal{O}_{X_y}} \mathcal{L}_y^{\otimes d})$
是满射，这就给出引理的最后一个断言。
\end{proof}

\begin{lemma}
\label{spaces-divisors-lemma-ample-in-neighbourhood}
（更一般的版本见《空间上的下降》引理
\ref{spaces-descent-lemma-ample-in-neighbourhood}。）设 $Y$ 为诺特概形，
$X$ 为 $Y$ 上的代数空间，并且结构态射 $f : X \to Y$ 固有。设
$\mathcal{L}$ 为可逆 $\mathcal{O}_X$-模。设 $y \in Y$ 为一点，使得 $X_y$
是概形，且 $\mathcal{L}_y$ 在 $X_y$ 上丰沛。则存在 $y$ 的开邻域
$V \subset Y$，使得 $\mathcal{L}|_{f^{-1}(V)}$ 在 $f^{-1}(V)/V$ 上丰沛
（如定义 \ref{spaces-divisors-definition-relatively-ample}）。
\end{lemma}

\begin{proof}
对 $\mathcal{F} = \mathcal{O}_X$，选取引理 \ref{spaces-divisors-lemma-ample-on-fibre} 中的
$d_0$。选取 $d \geq d_0$，使得存在 $r \geq 0$ 及截面
$s_{y, 0}, \ldots, s_{y, r} \in H^0(X_y, \mathcal{L}_y^{\otimes d})$
，它们定义闭浸入
$$
\varphi_y =
\varphi_{\mathcal{L}_y^{\otimes d}, (s_{y, 0}, \ldots, s_{y, r})} :
X_y \to \mathbf{P}^r_{\kappa(y)}.
$$
这由《态射》引理
\ref{morphisms-lemma-finite-type-over-affine-ample-very-ample} 保证；此外，为看出
$\varphi_y$ 是闭浸入，我们还使用《态射》引理
\ref{morphisms-lemma-image-proper-scheme-closed}，并使用《构造》第
\ref{constructions-section-projective-space} 节关于以可逆层和截面描述到射影空间
之态射的内容。由 $d_0$ 的选取，在用 $y$ 的一个开邻域替换 $Y$ 后，可以选取
$s_0, \ldots, s_r \in H^0(X, \mathcal{L}^{\otimes d})$
，它们映到 $s_{y, 0}, \ldots, s_{y, r}$。令 $X_{s_i} \subset X$ 为 $s_i$
生成 $\mathcal{L}^{\otimes d}$ 所在的开子空间。由于诸 $s_{y, i}$ 生成
$\mathcal{L}_y^{\otimes d}$，可见
$|X_y| \subset U = \bigcup |X_{s_i}|$。由于 $X \to Y$ 是闭的，存在开邻域
$y \in V \subset Y$，使得 $|f|^{-1}(V) \subset U$。用 $V$ 替换 $Y$ 后，
可以假设诸 $s_i$ 生成 $\mathcal{L}^{\otimes d}$。于是得到态射
$$
\varphi = \varphi_{\mathcal{L}^{\otimes d}, (s_0, \ldots, s_r)} :
X \longrightarrow \mathbf{P}^r_Y
$$
，满足 $\mathcal{L}^{\otimes d} \cong \varphi^*\mathcal{O}_{\mathbf{P}^r_Y}(1)$，
其到 $y$ 的基变换给出
$\varphi_y$（严格来说，我们需要写出一个证明，说明《构造》第
\ref{constructions-section-projective-space} 节给出的到射影空间之态射的构造也能
用于描述从代数空间到射影空间的态射；我们略去细节）。

\medskip\noindent
我们将用一个小技巧完成证明；“正确的”证明是直接说明，在基变换到 $y$ 的一个
开邻域后，$\varphi$ 是闭浸入。具体而言，由《空间的上同调》引理
\ref{spaces-cohomology-lemma-proper-finite-fibre-finite-in-neighbourhood} 可见，
% P09-ERR-0221
$\varphi$ 在 $\mathbf{P}^r_Y \to Y$ 位于 $y$ 上方的纤维
$\mathbf{P}^r_{\kappa(y)}$ 的一个开邻域上是有限的。利用
$\mathbf{P}^r_Y \to Y$ 是闭的，在缩小 $Y$ 后，可以假设 $\varphi$ 有限。
特别地，$X$ 是概形。于是，由非常一般的《态射》引理
\ref{morphisms-lemma-pullback-ample-tensor-relatively-ample}，
$\mathcal{L}^{\otimes d} \cong \varphi^*\mathcal{O}_{\mathbf{P}^r_Y}(1)$ 是丰沛的。
\end{proof}








\section{相对 Proj 的闭子空间}
\label{spaces-divisors-section-closed-in-relative-proj}

\noindent
这里给出一些关于相对 Proj 的闭子空间的辅助引理。本节是《除子》第
\ref{divisors-section-closed-in-relative-proj} 节的对应版本。

\begin{lemma}
\label{spaces-divisors-lemma-closed-subscheme-proj}
设 $S$ 为概形，$X$ 为 $S$ 上的代数空间，$\mathcal{A}$ 为拟凝聚分次
$\mathcal{O}_X$-代数。设
$\pi : P = \underline{\text{Proj}}_X(\mathcal{A}) \to X$ 为 $\mathcal{A}$ 的
相对 Proj，$i : Z \to P$ 为闭子空间。以 $\mathcal{I} \subset \mathcal{A}$
表示典范映射
$$
\mathcal{A}
\longrightarrow
\bigoplus\nolimits_{d \geq 0} \pi_*\left((i_*\mathcal{O}_Z)(d)\right)
$$
的核。若 $\pi$ 拟紧，则有同构
$Z = \underline{\text{Proj}}_X(\mathcal{A}/\mathcal{I})$。
\end{lemma}

\begin{proof}
由引理 \ref{spaces-divisors-lemma-relative-proj-separated}，态射 $\pi$ 分离。由于 $\pi$ 拟紧，
$\pi_*$ 把拟凝聚模变为拟凝聚模，见《空间的态射》引理
\ref{spaces-morphisms-lemma-pushforward}。因此，$\mathcal{I}$ 是拟凝聚
$\mathcal{O}_X$-模。特别地，$\mathcal{B} = \mathcal{A}/\mathcal{I}$ 是拟凝聚分次
$\mathcal{O}_X$-代数。函子性态射
$Z' = \underline{\text{Proj}}_X(\mathcal{B}) \to
\underline{\text{Proj}}_X(\mathcal{A})$ 处处有定义，并且是闭浸入，见引理
\ref{spaces-divisors-lemma-surjective-graded-rings-map-relative-proj}。因此，只需证明
$Z = Z'$（作为 $P$ 的闭子空间）。

\medskip\noindent
在此之后，该问题在基上是平展局部的，故通过平展局部化归约到概形情形
（《除子》引理 \ref{divisors-lemma-closed-subscheme-proj}）。
\end{proof}

\noindent
当闭子空间局部由有限多个方程截出时，可以用 $\mathcal{A}$ 的有限型理想层定义它。

\begin{lemma}
\label{spaces-divisors-lemma-closed-subscheme-proj-finite}
设 $S$ 为概形，$X$ 为 $S$ 上拟紧且拟分离的代数空间，$\mathcal{A}$ 为拟凝聚分次
$\mathcal{O}_X$-代数。设
$\pi : P = \underline{\text{Proj}}_X(\mathcal{A}) \to X$ 为 $\mathcal{A}$ 的
相对 Proj，$i : Z \to P$ 为闭子概形。若 $\pi$ 拟紧，且
% P09-ERR-0222
$i$ 有限表示，则存在 $d > 0$ 以及有限型拟凝聚 $\mathcal{O}_X$-子模
$\mathcal{F} \subset \mathcal{A}_d$，使得
$Z = \underline{\text{Proj}}_X(\mathcal{A}/\mathcal{F}\mathcal{A})$.
\end{lemma}

\begin{proof}
读者可以重做概形情形所用的论证。不过，我们将用一个技巧说明该引理由概形情形
得出。设 $\mathcal{I} \subset \mathcal{A}$ 为引理
\ref{spaces-divisors-lemma-closed-subscheme-proj} 中截出 $Z$ 的拟凝聚分次理想。选取仿射概形
$U$ 及满平展态射 $U \to X$，见《空间的性质》引理
\ref{spaces-properties-lemma-quasi-compact-affine-cover}。由概形情形（《除子》引理
\ref{divisors-lemma-closed-subscheme-proj-finite}），存在 $d > 0$ 及有限型拟凝聚
$\mathcal{O}_U$-子模
$\mathcal{F}' \subset \mathcal{I}_d|_U \subset \mathcal{A}_d|_U$
，使得 $Z \times_X U$ 等于
$\underline{\text{Proj}}_U(\mathcal{A}|_U/\mathcal{F}'\mathcal{A}|_U)$.
由《空间的极限》引理 \ref{spaces-limits-lemma-directed-colimit-finite-type}，可找到
有限型拟凝聚子模 $\mathcal{F} \subset \mathcal{I}_d$，使得
$\mathcal{F}' \subset \mathcal{F}|_U$。令
$Z' = \underline{\text{Proj}}_X(\mathcal{A}/\mathcal{F}\mathcal{A})$.
则 $Z' \to P$ 是闭浸入
% P09-ERR-0223
（引理 \ref{spaces-divisors-lemma-surjective-generated-degree-1-map-relative-proj}），且因为
$\mathcal{F}\mathcal{A} \subset \mathcal{I}$，有 $Z \subset Z'$。另一方面，由
$\mathcal{F}$ 的选取，$Z' \times_X U \subset Z \times_X U$。故如愿有
$Z = Z'$。
\end{proof}

\begin{lemma}
\label{spaces-divisors-lemma-closed-subscheme-proj-finite-type}
设 $S$ 为概形，$X$ 为 $S$ 上拟紧且拟分离的代数空间，$\mathcal{A}$ 为拟凝聚分次
$\mathcal{O}_X$-代数。设
$\pi : P = \underline{\text{Proj}}_X(\mathcal{A}) \to X$ 为 $\mathcal{A}$ 的
相对 Proj。设
% P09-ERR-0224
$i : Z \to X$ 为闭子空间，$U \subset X$ 为开集。假设：
\begin{enumerate}
\item $\pi$ 拟紧；
% P09-ERR-0222
\item $i$ 有限表示；
\item $|U| \cap |\pi|(|i|(|Z|)) = \emptyset$,
\item $U$ 拟紧；
\item 对所有 $n$，$\mathcal{A}_n$ 都是有限型 $\mathcal{O}_X$-模。
\end{enumerate}
则存在 $d > 0$ 及有限型拟凝聚 $\mathcal{O}_X$-子模
$\mathcal{F} \subset \mathcal{A}_d$，满足 (a)
$Z = \underline{\text{Proj}}_X(\mathcal{A}/\mathcal{F}\mathcal{A})$
，且 (b) $\mathcal{A}_d/\mathcal{F}$ 的支撑与 $U$ 不交。
\end{lemma}

\begin{proof}
我们采用与引理 \ref{spaces-divisors-lemma-closed-subscheme-proj-finite} 的证明相同的技巧，归约到
概形情形。设 $\mathcal{I} \subset \mathcal{A}$ 为引理
\ref{spaces-divisors-lemma-closed-subscheme-proj} 中截出 $Z$ 的拟凝聚分次理想。选取仿射概形
$W$ 及满平展态射 $W \to X$，见《空间的性质》引理
\ref{spaces-properties-lemma-quasi-compact-affine-cover}。由概形情形（《除子》引理
\ref{divisors-lemma-closed-subscheme-proj-finite-type}），存在 $d > 0$ 及有限型
拟凝聚 $\mathcal{O}_W$-子模
$\mathcal{F}' \subset \mathcal{I}_d|_W \subset \mathcal{A}_d|_W$
，使得 (a) $Z \times_X W$ 等于
$\underline{\text{Proj}}_W(\mathcal{A}|_W/\mathcal{F}'\mathcal{A}|_W)$
，且 (b) $\mathcal{A}_d|_W/\mathcal{F}'$ 的支撑与 $U \times_X W$ 不交。
由《空间的极限》引理 \ref{spaces-limits-lemma-directed-colimit-finite-type}，可找到
有限型拟凝聚子模 $\mathcal{F} \subset \mathcal{I}_d$，使得
$\mathcal{F}' \subset \mathcal{F}|_W$。令
$Z' = \underline{\text{Proj}}_X(\mathcal{A}/\mathcal{F}\mathcal{A})$.
则 $Z' \to P$ 是闭浸入（引理
% P09-ERR-0223
\ref{spaces-divisors-lemma-surjective-generated-degree-1-map-relative-proj}），且因为
$\mathcal{F}\mathcal{A} \subset \mathcal{I}$，有 $Z \subset Z'$。另一方面，由
$\mathcal{F}$ 的选取，$Z' \times_X W \subset Z \times_X W$。故 $Z = Z'$。
最后，$\mathcal{A}_d|_W/\mathcal{F}|_W$ 是
$\mathcal{A}_d|_W/\mathcal{F}'$ 的商，而后者的支撑包含于
% P09-ERR-0225
$W \setminus U \times_X W$，故可见 $\mathcal{A}_d/\mathcal{F}$ 的支撑包含于
$X \setminus U$。引理得证。
\end{proof}

\begin{lemma}
\label{spaces-divisors-lemma-conormal-sheaf-section-projective-bundle}
设 $S$ 为概形，$X$ 为 $S$ 上的代数空间，$\mathcal{E}$ 为拟凝聚
$\mathcal{O}_X$-模。存在双射
$$
\left\{
\begin{matrix}
\text{截面 }\sigma\text{，其所属的} \\
\text{态射为 } \mathbf{P}(\mathcal{E}) \to X
\end{matrix}
\right\}
\leftrightarrow
\left\{
\begin{matrix}
\text{满射 }\mathcal{E} \to \mathcal{L}\text{，其中} \\
\mathcal{L}\text{ 是可逆 }\mathcal{O}_X\text{-模}
\end{matrix}
\right\}
$$
在此情形下，$\sigma$ 是闭浸入，并且有典范同构
$$
\Ker(\mathcal{E} \to \mathcal{L})
\otimes_{\mathcal{O}_X} \mathcal{L}^{\otimes -1}
\longrightarrow
\mathcal{C}_{\sigma(X)/\mathbf{P}(\mathcal{E})}
$$
该双射和同构都与基变换相容。
\end{lemma}

\begin{proof}
由于这些构造与基变换相容，只需在 $X$ 上平展局部地验证该断言。因此，可以假设
$X$ 是概形，此时结果即《除子》引理
\ref{divisors-lemma-conormal-sheaf-section-projective-bundle}。
\end{proof}






\section{爆破}
\label{spaces-divisors-section-blowing-up}

\noindent
爆破是代数几何中的重要工具。

\begin{definition}
\label{spaces-divisors-definition-blow-up}
设 $S$ 为概形，$X$ 为 $S$ 上的代数空间。设
$\mathcal{I} \subset \mathcal{O}_X$ 为拟凝聚理想层，并设 $Z \subset X$ 为
$\mathcal{I}$ 所对应的闭子空间（《空间的态射》引理
\ref{spaces-morphisms-lemma-closed-immersion-ideals}）。沿 $Z$ 的
{\it $X$ 的爆破}，或在理想层 $\mathcal{I}$ 处的 {\it $X$ 的爆破}，是态射
$$
b :
\underline{\text{Proj}}_X
\left(\bigoplus\nolimits_{n \geq 0} \mathcal{I}^n\right)
\longrightarrow
X
$$
。爆破的 {\it 例外除子} 是逆像 $b^{-1}(Z)$。有时称 $Z$ 为爆破的
{\it 中心}。
\end{definition}

\noindent
稍后会看到，例外除子是有效 Cartier 除子。此外，爆破可以刻画为 $X$ 上使 $Z$
的逆像成为有效 Cartier 除子的“最小”代数空间。

\medskip\noindent
若 $b : X' \to X$ 是 $X$ 在 $Z$ 处的爆破，则常以 $\mathcal{O}_{X'}(n)$ 表示
结构层的扭转。注意，它们是可逆 $\mathcal{O}_{X'}$-模，并且
$\mathcal{O}_{X'}(n) = \mathcal{O}_{X'}(1)^{\otimes n}$
，因为 $X'$ 是一个由 $1$ 次部分生成的拟凝聚分次 $\mathcal{O}_X$-代数的相对
Proj，见引理 \ref{spaces-divisors-lemma-relative-proj-generated-in-degree-1}。

\begin{lemma}
\label{spaces-divisors-lemma-blowing-up-affine}
设 $S$ 为概形，$X$ 为 $S$ 上的代数空间，
$\mathcal{I} \subset \mathcal{O}_X$ 为拟凝聚理想层。设 $U = \Spec(A)$ 为
平展于 $X$ 的仿射概形，并设 $I \subset A$ 为 $\mathcal{I}|_U$ 所对应的理想。
若 $X' \to X$ 是 $X$ 在 $\mathcal{I}$ 处的爆破，则有典范同构
$$
U \times_X X' = \text{Proj}(\bigoplus\nolimits_{d \geq 0} I^d)
$$
，这是 $U$ 上概形的同构；右端是 $A$ 中 $I$ 的 Rees 代数的齐次谱。此外，
$U \times_X X'$ 有一个由仿射爆破代数 $A[\frac{I}{a}]$ 的谱组成的仿射开覆盖。
\end{lemma}

\begin{proof}
注意，限制 $\mathcal{I}|_U$ 等于 $\mathcal{I}$ 沿态射 $U \to X$ 的拉回，见
《空间的性质》第 \ref{spaces-properties-section-modules} 节。因此，将引理
\ref{spaces-divisors-lemma-relative-proj} 与《除子》引理
\ref{divisors-lemma-blowing-up-affine} 结合即得本引理。
\end{proof}

\begin{lemma}
\label{spaces-divisors-lemma-flat-base-change-blowing-up}
设 $S$ 为概形，$X_1 \to X_2$ 为 $S$ 上代数空间的平坦态射，
$Z_2 \subset X_2$ 为闭子空间。设 $Z_1$ 为 $Z_2$ 在 $X_1$ 中的逆像，并设
$X'_i$ 为 $X_i$ 在 $Z_i$ 处的爆破。则存在笛卡尔图
$$
\xymatrix{
X_1' \ar[r] \ar[d] & X_2' \ar[d] \\
X_1 \ar[r] & X_2
}
$$
，其中各项都是 $S$ 上的代数空间。
\end{lemma}

\begin{proof}
设 $\mathcal{I}_2$ 为 $Z_2$ 在 $X_2$ 中的理想层，以 $g : X_1 \to X_2$ 表示
给定态射。则 $Z_1$ 的理想层 $\mathcal{I}_1$ 是映射
$g^*\mathcal{I}_2 \to \mathcal{O}_{X_1}$
的像（见《空间的态射》定义
\ref{spaces-morphisms-definition-inverse-image-closed-subspace} 及定义后的讨论）。
由引理 \ref{spaces-divisors-lemma-relative-proj-base-change} 可见，
$X_1 \times_{X_2} X_2'$ 是 $\bigoplus_{n \geq 0} g^*\mathcal{I}_2^n$ 的
相对 Proj。由于 $g$ 平坦，映射
$g^*\mathcal{I}_2^n \to \mathcal{O}_{X_1}$ 是单射，其像为
$\mathcal{I}_1^n$。故 $X_1 \times_{X_2} X_2' = X_1'$。
\end{proof}

\begin{lemma}
\label{spaces-divisors-lemma-blowing-up-gives-effective-Cartier-divisor}
设 $S$ 为概形，$X$ 为 $S$ 上的代数空间，$Z \subset X$ 为闭子空间。
$X$ 在 $Z$ 处的爆破 $b : X' \to X$ 具有下列性质：
\begin{enumerate}
\item $b|_{b^{-1}(X \setminus Z)} : b^{-1}(X \setminus Z) \to X \setminus Z$
是同构；
\item 例外除子 $E = b^{-1}(Z)$ 是 $X'$ 上的有效 Cartier 除子；
\item 有典范同构
$\mathcal{O}_{X'}(-1) = \mathcal{O}_{X'}(E)$
\end{enumerate}
\end{lemma}

\begin{proof}
设 $U$ 为概形，$U \to X$ 为满平展态射。由于爆破与平坦基变换可交换（引理
\ref{spaces-divisors-lemma-flat-base-change-blowing-up}），可以在基变换到 $U$ 后证明各断言。
这把问题归约到概形情形，此时结果即《除子》引理
\ref{divisors-lemma-blowing-up-gives-effective-Cartier-divisor}。
\end{proof}

\begin{lemma}[爆破的泛性质]
\label{spaces-divisors-lemma-universal-property-blowing-up}
% P09-ERR-0226
\begin{slogan}
爆破一个闭子集，使之成为 Cartier 的。
\end{slogan}
设 $S$ 为概形，$X$ 为 $S$ 上的代数空间，$Z \subset X$ 为闭子空间。设
$\mathcal{C}$ 为 $(\textit{Spaces}/X)$ 的满子范畴，其对象是满足如下条件的
$Y \to X$：$Z$ 的逆像是 $Y$ 上的有效 Cartier 除子。则 $X$ 在 $Z$ 处的爆破
$b : X' \to X$ 是 $\mathcal{C}$ 的终对象。
\end{lemma}

\begin{proof}
由引理 \ref{spaces-divisors-lemma-blowing-up-gives-effective-Cartier-divisor} 可见，
$b : X' \to X$ 是 $\mathcal{C}$ 的对象。设 $f : Y \to X$ 为
$\mathcal{C}$ 的对象。我们须证明存在唯一的 $X$ 上态射 $Y \to X'$。令
$D = f^{-1}(Z)$。设 $\mathcal{I} \subset \mathcal{O}_X$ 为 $Z$ 的理想层，
$\mathcal{I}_D$ 为 $D$ 的理想层。则
$f^*\mathcal{I} \to \mathcal{I}_D$ 是到可逆 $\mathcal{O}_Y$-模的满射。
它延拓为分次 $\mathcal{O}_Y$-代数映射
$\psi : \bigoplus f^*\mathcal{I}^d \to \bigoplus \mathcal{I}_D^d$
。（注意，因为 $D$ 是有效 Cartier 除子，有
$\mathcal{I}_D^d = \mathcal{I}_D^{\otimes d}$。）由引理
% P09-ERR-0227
\ref{spaces-divisors-lemma-relative-proj-generated-in-degree-1}，三元组
$(f : Y \to X, \mathcal{I}_D, \psi)$ 定义了 $X$ 上态射 $Y \to X'$。其限制
$$
Y \setminus D \longrightarrow X' \setminus b^{-1}(Z) = X \setminus Z
$$
是唯一的。由引理 \ref{spaces-divisors-lemma-complement-effective-Cartier-divisor}，开子空间
$Y \setminus D$ 在 $Y$ 中概形论稠密。因此，由《空间的态射》引理
\ref{spaces-morphisms-lemma-equality-of-morphisms}，态射 $Y \to X'$ 唯一
（而且由引理 \ref{spaces-divisors-lemma-relative-proj-separated}，$b$ 分离）。
\end{proof}

\begin{lemma}
\label{spaces-divisors-lemma-blow-up-effective-Cartier-divisor}
设 $S$ 为概形，$X$ 为 $S$ 上的代数空间，$Z \subset X$ 为有效 Cartier
除子。$X$ 在 $Z$ 处的爆破是 $X$ 的恒等态射。
\end{lemma}

\begin{proof}
这立即由爆破的泛性质（引理 \ref{spaces-divisors-lemma-universal-property-blowing-up}）得出。
\end{proof}

\begin{lemma}
\label{spaces-divisors-lemma-blow-up-reduced-space}
设 $S$ 为概形，$X$ 为 $S$ 上的代数空间，
$\mathcal{I} \subset \mathcal{O}_X$ 为拟凝聚理想层。若 $X$ 约化，则 $X$
在 $\mathcal{I}$ 处的爆破 $X'$ 约化。
\end{lemma}

\begin{proof}
设 $U$ 为概形，$U \to X$ 为满平展态射。由于爆破与平坦基变换可交换（引理
\ref{spaces-divisors-lemma-flat-base-change-blowing-up}），可以在基变换到 $U$ 后证明该断言。
这把问题归约到概形情形，此时结果即《除子》引理
\ref{divisors-lemma-blow-up-reduced-scheme}。
\end{proof}

\begin{lemma}
\label{spaces-divisors-lemma-blowup-finite-nr-irreducibles}
设 $S$ 为概形，$X$ 为 $S$ 上的代数空间，$b : X' \to X$ 为 $X$ 在某闭子空间
处的爆破。若 $X$ 满足《空间的态射》引理
\ref{spaces-morphisms-lemma-prepare-normalization} 的等价条件，则 $X'$ 也满足。
\end{lemma}

\begin{proof}
这立即由陈述中所引用的引理、引理 \ref{spaces-divisors-lemma-blowing-up-affine} 中爆破的平展局部
描述以及《除子》引理
\ref{divisors-lemma-blow-up-and-irreducible-components} 得出。
\end{proof}

\begin{lemma}
\label{spaces-divisors-lemma-blow-up-pullback-effective-Cartier}
设 $S$ 为概形，$X$ 为 $S$ 上的代数空间，$b : X' \to X$ 为 $X$ 在某闭子空间
处的爆破。对 $X$ 上任意有效 Cartier 除子 $D$，拉回 $b^{-1}D$ 有定义
（见定义 \ref{spaces-divisors-definition-pullback-effective-Cartier-divisor}）。
\end{lemma}

\begin{proof}
由引理 \ref{spaces-divisors-lemma-blowing-up-affine} 及
\ref{spaces-divisors-lemma-characterize-effective-Cartier-divisor}，这归约为下列代数事实：
设 $A$ 为环，$I \subset A$ 为理想，$a \in I$，而 $x \in A$ 为非零因子。
则 $x$ 在 $A[\frac{I}{a}]$ 中的像是非零因子。事实上，假设
$x (y/a^n) = 0$ 在 $A[\frac{I}{a}]$ 中成立。则对某个 $m$，在 $A$ 中有
$a^mxy = 0$。由于 $x$ 是非零因子，故 $a^my = 0$。于是如所需，$y/a^n$
在 $A[\frac{I}{a}]$ 中为零。
\end{proof}

\begin{lemma}
\label{spaces-divisors-lemma-blowing-up-two-ideals}
设 $S$ 为概形，$X$ 为 $S$ 上的代数空间。设
% P09-ERR-0228
$\mathcal{I} \subset \mathcal{O}_X$ 及 $\mathcal{J}$ 为拟凝聚理想层。设
$b : X' \to X$ 为 $X$ 在 $\mathcal{I}$ 处的爆破，$b' : X'' \to X'$ 为
$X'$ 在 $b^{-1}\mathcal{J} \mathcal{O}_{X'}$ 处的爆破。则 $X'' \to X$ 典范
同构于 $X$ 在 $\mathcal{I}\mathcal{J}$ 处的爆破。
\end{lemma}

\begin{proof}
设 $E \subset X'$ 为 $b$ 的例外除子，由引理
\ref{spaces-divisors-lemma-blowing-up-gives-effective-Cartier-divisor}，它是有效 Cartier 除子。
由引理 \ref{spaces-divisors-lemma-blow-up-pullback-effective-Cartier}，$(b')^{-1}E$ 是
$X''$ 上的有效 Cartier 除子。设 $E' \subset X''$ 为 $b'$ 的例外除子（它也是
有效 Cartier 除子）。考虑有效 Cartier 除子
$E'' = E' + (b')^{-1}E$。按构造，$E''$ 的理想为
$(b \circ b')^{-1}\mathcal{I} (b \circ b')^{-1}\mathcal{J} \mathcal{O}_{X''}$.
因此，由引理 \ref{spaces-divisors-lemma-universal-property-blowing-up}，存在从 $X''$ 到
$X$ 在 $\mathcal{I}\mathcal{J}$ 处的爆破 $c : Y \to X$ 的典范态射。反过来，
由于 $\mathcal{I}\mathcal{J}$ 拉回为可逆理想，可见
$c^{-1}\mathcal{I}\mathcal{O}_Y$ 定义有效 Cartier 除子，见引理
\ref{spaces-divisors-lemma-sum-closed-subschemes-effective-Cartier}。因此，由引理
% P09-ERR-0229
\ref{spaces-divisors-lemma-universal-property-blowing-up}，存在 $X$ 上态射 $c' : Y \to X'$。
于是 $(c')^{-1}b^{-1}\mathcal{J}\mathcal{O}_Y = c^{-1}\mathcal{J}\mathcal{O}_Y$，
它也定义有效 Cartier 除子。因此，存在
$X'$ 上态射 $c'' : Y \to X''$。我们略去验证此态射与先前构造的态射
$X'' \to Y$ 互逆。
\end{proof}

\begin{lemma}
\label{spaces-divisors-lemma-blowing-up-projective}
设 $S$ 为概形，$X$ 为 $S$ 上的代数空间，
$\mathcal{I} \subset \mathcal{O}_X$ 为拟凝聚理想层。设 $b : X' \to X$
为 $X$ 在理想层 $\mathcal{I}$ 处的爆破。若 $\mathcal{I}$ 是有限型的，则
$b : X' \to X$ 是固有态射。
\end{lemma}

\begin{proof}
设 $U$ 为概形，$U \to X$ 为满平展态射。由于爆破与平坦基变换可交换（引理
\ref{spaces-divisors-lemma-flat-base-change-blowing-up}），可以在基变换到 $U$ 后证明该断言
（见《空间的态射》引理 \ref{spaces-morphisms-lemma-proper-local}）。这把问题归约到
概形情形。此时，由《除子》引理 \ref{divisors-lemma-blowing-up-projective}，
态射 $b$ 是射影的，故由《态射》引理
\ref{morphisms-lemma-locally-projective-proper}，它是固有的。
\end{proof}

\begin{lemma}
\label{spaces-divisors-lemma-composition-finite-type-blowups}
设 $S$ 为概形，$X$ 为 $S$ 上的代数空间。假设 $X$ 拟紧且拟分离。设
$Z \subset X$ 为有限表示闭子空间，$b : X' \to X$ 为中心为 $Z$ 的爆破。
设 $Z' \subset X'$ 为有限表示闭子空间，$X'' \to X'$ 为中心为 $Z'$ 的爆破。
则存在有限表示闭子空间 $Y \subset X$，使得：
\begin{enumerate}
\item $|Y| = |Z| \cup |b|(|Z'|)$；
\item 复合态射 $X'' \to X$ 同构于 $X$ 在 $Y$ 处的爆破。
\end{enumerate}
\end{lemma}

\begin{proof}
条件 $Z \to X$ 有限表示意味着 $Z$ 由有限型拟凝聚理想层
$\mathcal{I} \subset \mathcal{O}_X$ 截出，见《空间的态射》引理
\ref{spaces-morphisms-lemma-closed-immersion-finite-presentation}。写
$\mathcal{A} = \bigoplus_{n \geq 0} \mathcal{I}^n$，从而
$X' = \underline{\text{Proj}}(\mathcal{A})$。注意，由《空间的极限》引理
\ref{spaces-limits-lemma-quasi-coherent-finite-type-ideals}，$X \setminus Z$ 是
$X$ 的拟紧开子空间。由于 $b^{-1}(X \setminus Z) \to X \setminus Z$ 是同构
（引理 \ref{spaces-divisors-lemma-blowing-up-gives-effective-Cartier-divisor}），同一结果表明
$b^{-1}(X \setminus Z) \setminus Z'$ 是 $X'$ 的拟紧开子空间。因此，
$U = X \setminus (Z \cup b(Z'))$ 是 $X$ 的拟紧开子空间。由引理
\ref{spaces-divisors-lemma-closed-subscheme-proj-finite-type}，存在 $d > 0$ 及有限型
$\mathcal{O}_X$-子模 $\mathcal{F} \subset \mathcal{I}^d$，使得
$Z' = \underline{\text{Proj}}(\mathcal{A}/\mathcal{F}\mathcal{A})$，并且
$\mathcal{I}^d/\mathcal{F}$ 的支撑包含于 $X \setminus U$。

\medskip\noindent
由于 $\mathcal{F} \subset \mathcal{I}^d$ 是 $\mathcal{O}_X$-子模，可以把
$\mathcal{F} \subset \mathcal{I}^d \subset \mathcal{O}_X$ 看作 $X$ 上的有限型
拟凝聚理想层。为避免混淆，将其记为
$\mathcal{J} \subset \mathcal{O}_X$。由于 $\mathcal{I}^d / \mathcal{J}$ 及
$\mathcal{O}/\mathcal{I}^d$ 的支撑都包含于 $|X| \setminus |U|$，可见
$|V(\mathcal{J})|$ 包含于 $|X| \setminus |U|$。反过来，由于
$\mathcal{J} \subset \mathcal{I}^d$，可见
$|Z| \subset |V(\mathcal{J})|$。在
$X \setminus Z \cong X' \setminus b^{-1}(Z)$ 上，理想层 $\mathcal{J}$ 截出
$Z'$（见下方公式）。因此，$|V(\mathcal{J})|$ 等于
$|Z| \cup |b|(|Z'|)$，从而还有
$|V(\mathcal{I}\mathcal{J})| = |Z| \cup |b|(|Z'|)$。此外，
$\mathcal{I}\mathcal{J}$ 是两个有限型理想之积，故也是有限型理想。我们断言
$X'' \to X$ 同构于 $X$ 在 $\mathcal{I}\mathcal{J}$ 处的爆破；令
$Y = V(\mathcal{I}\mathcal{J})$，这便完成引理的证明。

\medskip\noindent
首先回顾，$X$ 在 $\mathcal{I}\mathcal{J}$ 处的爆破等同于 $X'$ 在
$b^{-1}\mathcal{J} \mathcal{O}_{X'}$ 处的爆破，见引理
\ref{spaces-divisors-lemma-blowing-up-two-ideals}。因此，只需证明 $X'$ 在
$b^{-1}\mathcal{J} \mathcal{O}_{X'}$ 处的爆破与 $X'$ 在 $Z'$ 处的爆破一致。
我们将证明
$$
b^{-1}\mathcal{J} \mathcal{O}_{X'} = \mathcal{I}_E^d \mathcal{I}_{Z'}
$$
% P09-ERR-0230
，这是 $X''$ 上理想层的等式。这将证明所需，因为 $\mathcal{I}_E^d$ 截出有效
Cartier 除子 $dE$，而我们可以使用引理
\ref{spaces-divisors-lemma-blow-up-effective-Cartier-divisor} 和
\ref{spaces-divisors-lemma-blowing-up-two-ideals}。

\medskip\noindent
为验证所示理想的等式，可以局部进行。采用引理
\ref{spaces-divisors-lemma-blowing-up-affine} 中的记号 $A$、$I$、$a \in I$，可见
% P09-ERR-0231
$\mathcal{F}$ 对应于一个 $R$-子模 $M \subset I^d$，后者同构映到理想
$J \subset R$。条件
$Z' = \underline{\text{Proj}}(\mathcal{A}/\mathcal{F}\mathcal{A})$
意味着 $Z' \cap \Spec(A[\frac{I}{a}])$ 由元素 $m/a^d$（$m \in M$）生成的
理想截出。设元素 $m \in M$ 对应于函数 $f \in J$。那么在仿射爆破代数
$A' = A[\frac{I}{a}]$ 中，有 $f = (a^dm)/a^d = a^d (m/a^d)$。
故等式成立。
\end{proof}









\section{严格变换}
\label{spaces-divisors-section-strict-transform}

\noindent
本节是《除子》第 \ref{divisors-section-strict-transform} 节的对应版本。设 $S$
为概形，$B$ 为 $S$ 上的代数空间，且 $Z \subset B$ 为闭子空间。设
$b : B' \to B$ 为 $B$ 在 $Z$ 处的爆破，并以 $E \subset B'$ 表示例外除子
% P09-ERR-0232
$E = b^{-1}Z$。下文常考虑 $B$ 上的代数空间 $X$，并作笛卡尔图
$$
\xymatrix{
\text{pr}_{B'}^{-1}E \ar[r] \ar[d] &
X \times_B B' \ar[r]_-{\text{pr}_X} \ar[d]_{\text{pr}_{B'}} &
X \ar[d]^f \\
E \ar[r] & B' \ar[r] & B
}
$$
。由于 $E$ 是有效 Cartier 除子（引理
\ref{spaces-divisors-lemma-blowing-up-gives-effective-Cartier-divisor}），可见
$\text{pr}_{B'}^{-1}E \subset X \times_B B'$ 是局部主的（引理
\ref{spaces-divisors-lemma-pullback-locally-principal}）。因此，
$\text{pr}_{B'}^{-1}E$ 在 $X \times_B B'$ 中之补集的包含态射是仿射的，特别地
是拟紧的（引理 \ref{spaces-divisors-lemma-complement-locally-principal-closed-subscheme}）。
因此，对拟凝聚 $\mathcal{O}_{X \times_B B'}$-模 $\mathcal{G}$，支撑于
$|\text{pr}_{B'}^{-1}E|$ 的截面子层是拟凝聚子模，见《空间的极限》定义
\ref{spaces-limits-definition-subsheaf-sections-supported-on-closed}。若
$\mathcal{G}$ 是拟凝聚代数层，例如
$\mathcal{G} = \mathcal{O}_{X \times_B B'}$，则此子层是 $\mathcal{G}$ 的理想。

\begin{definition}
\label{spaces-divisors-definition-strict-transform}
% P09-ERR-0233
沿用上述 $Z \subset B$ 与 $f : X \to B$。
\begin{enumerate}
\item 给定拟凝聚 $\mathcal{O}_X$-模 $\mathcal{F}$，$\mathcal{F}$ 关于 $B$
在 $Z$ 处爆破的{\it 严格变换}，是 $\text{pr}_X^*\mathcal{F}$ 模去支撑于
$|\text{pr}_{B'}^{-1}E|$ 的截面子模所得的商 $\mathcal{F}'$。
\item $X$ 的{\it 严格变换}是闭子空间
$X' \subset X \times_B B'$，它由 $\mathcal{O}_{X \times_B B'}$ 中支撑于
$|\text{pr}_{B'}^{-1}E|$ 的截面所成的拟凝聚理想截出。
\end{enumerate}
\end{definition}

\noindent
注意，沿一次爆破取严格变换依赖于该爆破所使用的闭子空间（而不仅依赖于态射
$B' \to B$）。

\begin{lemma}[平展局部化与严格变换]
\label{spaces-divisors-lemma-strict-transform-local}
处于定义 \ref{spaces-divisors-definition-strict-transform} 的情形。设
$$
\xymatrix{
U \ar[r] \ar[d] & X \ar[d] \\
V \ar[r] & B
}
$$
为态射的交换图，其中 $U$ 与 $V$ 为概形，水平箭头为平展态射。设
$V' \to V$ 为 $V$ 在 $Z \times_B V$ 处的爆破。则
\begin{enumerate}
\item $V' = V \times_B B'$，且态射
$V' \to B'$ 与 $U \times_V V' \to X \times_B B'$ 都是平展的；
\item $U$ 关于 $V' \to V$ 的严格变换 $U'$ 等于 $X' \times_X U$，其中
$X'$ 是 $X$ 关于 $B' \to B$ 的严格变换；并且
\item 对拟凝聚 $\mathcal{O}_X$-模 $\mathcal{F}$，严格变换
$\mathcal{F}'$ 在 $U \times_V V'$ 上的限制是 $\mathcal{F}|_U$ 关于
$V' \to V$ 的严格变换。
\end{enumerate}
\end{lemma}

\begin{proof}
(1) 由爆破与平坦基变换可交换（引理
\ref{spaces-divisors-lemma-flat-base-change-blowing-up}）、平展态射平坦以及平展态射的基变换仍平展
这些事实得出。(3) 随即由下述事实得出：取支撑于闭集的截面层与沿平展态射的拉回
可交换，见《空间的极限》引理
% P09-ERR-0234
\ref{spaces-limits-lemma-sections-supported-on-closed-subset}。(2) 由把 (3) 应用于
$\mathcal{F} = \mathcal{O}_X$ 得出。
\end{proof}

\begin{lemma}
\label{spaces-divisors-lemma-strict-transform}
处于定义 \ref{spaces-divisors-definition-strict-transform} 的情形。
\begin{enumerate}
\item $X$ 的严格变换 $X'$ 是 $X$ 在闭子空间 $f^{-1}Z$ 处的爆破；这里该闭子空间
位于 $X$ 中。
\item 对拟凝聚 $\mathcal{O}_X$-模 $\mathcal{F}$，严格变换
$\mathcal{F}'$ 典范同构于 $\mathcal{F}$ 关于爆破 $X' \to X$ 的严格变换沿
$X' \to X \times_B B'$ 的正像。
\end{enumerate}
\end{lemma}

\begin{proof}
设 $X'' \to X$ 为 $X$ 在 $f^{-1}Z$ 处的爆破。由爆破的泛性质（引理
\ref{spaces-divisors-lemma-universal-property-blowing-up}），存在交换图
$$
\xymatrix{
X'' \ar[r] \ar[d] & X \ar[d] \\
B' \ar[r] & B
}
$$
，从而得到态射 $i : X'' \to X \times_B B'$。引理的第一个断言是 $i$ 为像等于
$X'$ 的闭浸入。第二个断言是 $\mathcal{F}' = i_*\mathcal{F}''$，其中
$\mathcal{F}''$ 是 $\mathcal{F}$ 关于爆破 $X'' \to X$ 的严格变换。可以在 $X$
上平展局部地检验这些断言，因而归约到概形情形（《除子》引理
\ref{divisors-lemma-strict-transform}）。略去若干细节。
\end{proof}

\begin{lemma}
\label{spaces-divisors-lemma-strict-transform-flat}
处于定义 \ref{spaces-divisors-definition-strict-transform} 的情形。
\begin{enumerate}
\item 若 $X$ 在所有位于 $Z$ 上方的点处都平坦于 $B$，则 $X$ 的严格变换等于
基变换 $X \times_B B'$。
\item 设 $\mathcal{F}$ 为拟凝聚 $\mathcal{O}_X$-模。若 $\mathcal{F}$ 在所有
位于 $Z$ 上方的点处都平坦于 $B$，则 $\mathcal{F}$ 的严格变换
$\mathcal{F}'$ 等于拉回 $\text{pr}_X^*\mathcal{F}$。
\end{enumerate}
\end{lemma}

\begin{proof}
略。提示：由概形情形（《除子》引理
\ref{divisors-lemma-strict-transform-flat}）经平展局部化（引理
\ref{spaces-divisors-lemma-strict-transform-local}）得出。
\end{proof}

\begin{lemma}
\label{spaces-divisors-lemma-strict-transform-affine}
设 $S$ 为概形，$B$ 为 $S$ 上的代数空间，$Z \subset B$ 为闭子空间。设
$b : B' \to B$ 为 $B$ 在 $Z$ 处的爆破，$g : X \to Y$ 为 $B$ 上代数空间的
仿射态射。设 $\mathcal{F}$ 为 $X$ 上的拟凝聚层。设
$g' : X \times_B B' \to Y \times_B B'$ 为 $g$ 的基变换，$\mathcal{F}'$ 为
$\mathcal{F}$ 关于 $b$ 的严格变换。则 $g'_*\mathcal{F}'$ 是
$g_*\mathcal{F}$ 的严格变换。
\end{lemma}

\begin{proof}
略。提示：由概形情形（《除子》引理
\ref{divisors-lemma-strict-transform-affine}）经平展局部化（引理
\ref{spaces-divisors-lemma-strict-transform-local}）得出。
\end{proof}

\begin{lemma}
\label{spaces-divisors-lemma-strict-transform-different-centers}
设 $S$ 为概形，$B$ 为 $S$ 上的代数空间，$Z \subset B$ 为闭子空间。设
$D \subset B$ 为有效 Cartier 除子，$Z' \subset B$ 为由 $Z$ 与 $D$ 的理想层
之积截出的闭子空间。设 $B' \to B$ 为 $B$ 在 $Z$ 处的爆破。
\begin{enumerate}
\item $B$ 在 $Z'$ 处的爆破同构于 $B' \to B$。
\item 设 $f : X \to B$ 为代数空间的态射，$\mathcal{F}$ 为拟凝聚
$\mathcal{O}_X$-模。若 $\mathcal{F}$ 中支撑于 $|f^{-1}D|$ 的截面子层为零，
则 $\mathcal{F}$ 关于在 $Z$ 处爆破的严格变换，与 $\mathcal{F}$ 关于 $B$
在 $Z'$ 处爆破的严格变换一致。
\end{enumerate}
\end{lemma}

\begin{proof}
略。提示：由概形情形（《除子》引理
\ref{divisors-lemma-strict-transform-different-centers}）经平展局部化（引理
\ref{spaces-divisors-lemma-strict-transform-local}）得出。
\end{proof}

\begin{lemma}
\label{spaces-divisors-lemma-strict-transform-composition-blowups}
设 $S$ 为概形，$B$ 为 $S$ 上的代数空间，$Z \subset B$ 为闭子空间。设
$b : B' \to B$ 为中心为 $Z$ 的爆破。设 $Z' \subset B'$ 为闭子空间，
$B'' \to B'$ 为中心为 $Z'$ 的爆破。设
% P09-ERR-0235
$Y \subset B$ 为闭子概形，并且 $|Y| = |Z| \cup |b|(|Z'|)$，复合态射
$B'' \to B$ 同构于 $B$ 在 $Y$ 处的爆破。在此情形，对任意 $B$ 上概形 $X$
及 $\mathcal{F} \in \QCoh(\mathcal{O}_X)$，有
\begin{enumerate}
\item $\mathcal{F}$ 关于 $B$ 在 $Y$ 处爆破的严格变换，等于先取
$\mathcal{F}$ 关于 $B$ 在 $Z$ 处的爆破 $B' \to B$ 的严格变换，再取所得层关于
在 $Z'$ 处的爆破 $B'' \to B'$ 的严格变换；并且
\item $X$ 关于 $B$ 在 $Y$ 处爆破的严格变换，等于先取 $X$ 关于 $B$ 在 $Z$
处的爆破 $B' \to B$ 的严格变换，再取所得空间关于在 $Z'$ 处的爆破
$B'' \to B'$ 的严格变换。
\end{enumerate}
\end{lemma}

\begin{proof}
略。提示：由概形情形（《除子》引理
\ref{divisors-lemma-strict-transform-composition-blowups}）经平展局部化（引理
\ref{spaces-divisors-lemma-strict-transform-local}）得出。
\end{proof}

\begin{lemma}
\label{spaces-divisors-lemma-strict-transform-universally-injective}
处于定义 \ref{spaces-divisors-definition-strict-transform} 的情形。假设
$$
0 \to \mathcal{F}_1 \to \mathcal{F}_2 \to \mathcal{F}_3 \to 0
$$
是 $X$ 上拟凝聚层的正合列，并且在任意基变换 $T \to B$ 后仍正合。则关于任意
爆破 $B' \to B$ 的严格变换
% P09-ERR-0236
$\mathcal{F}_i'$ 也组成短正合列
$0 \to \mathcal{F}'_1 \to \mathcal{F}'_2 \to \mathcal{F}'_3 \to 0$。
\end{lemma}

\begin{proof}
略。提示：由概形情形（《除子》引理
\ref{divisors-lemma-strict-transform-universally-injective}）经平展局部化（引理
\ref{spaces-divisors-lemma-strict-transform-local}）得出。
\end{proof}

\begin{lemma}
\label{spaces-divisors-lemma-strict-transform-blowup-fitting-ideal}
设 $S$ 为概形，$B$ 为 $S$ 上的代数空间，$\mathcal{F}$ 为有限型拟凝聚
$\mathcal{O}_B$-模。设
% P09-ERR-0237
$Z_k \subset S$ 为由 $\text{Fit}_k(\mathcal{F})$ 截出的闭子概形，见第
\ref{spaces-divisors-section-fitting-ideals} 节。设 $B' \to B$ 为 $B$ 在 $Z_k$ 处的爆破，
$\mathcal{F}'$ 为 $\mathcal{F}$ 的严格变换。则 $\mathcal{F}'$ 局部可由
$\leq k$ 个截面生成。
\end{lemma}

\begin{proof}
略。由概形情形（《除子》引理
\ref{divisors-lemma-strict-transform-blowup-fitting-ideal}）经平展局部化（引理
\ref{spaces-divisors-lemma-strict-transform-local}）得出。
\end{proof}

\begin{lemma}
\label{spaces-divisors-lemma-strict-transform-blowup-fitting-ideal-locally-free}
设 $S$ 为概形，$B$ 为 $S$ 上的代数空间，$\mathcal{F}$ 为有限型拟凝聚
$\mathcal{O}_B$-模。设
% P09-ERR-0237
$Z_k \subset S$ 为由 $\text{Fit}_k(\mathcal{F})$ 截出的闭子概形，见第
\ref{spaces-divisors-section-fitting-ideals} 节。假设 $\mathcal{F}$ 在 $B \setminus Z_k$ 上
局部自由且秩为 $k$。设 $B' \to B$ 为 $B$ 在 $Z_k$ 处的爆破，$\mathcal{F}'$
为 $\mathcal{F}$ 的严格变换。则 $\mathcal{F}'$ 局部自由且秩为 $k$。
\end{lemma}

\begin{proof}
略。由概形情形（《除子》引理
\ref{divisors-lemma-strict-transform-blowup-fitting-ideal-locally-free}）
经平展局部化（引理 \ref{spaces-divisors-lemma-strict-transform-local}）得出。
\end{proof}












\section{容许爆破}
\label{spaces-divisors-section-admissible-blowups}

\noindent
为更好地控制所用的爆破，引入以下标准术语。

\begin{definition}
\label{spaces-divisors-definition-admissible-blowup}
设 $S$ 为概形，$X$ 为 $S$ 上的代数空间，$U \subset X$ 为开子空间。若存在
有限表示闭浸入 $Z \to X$，其中 $Z$ 与 $U$ 不交，并且 $X'$ 同构于 $X$ 在 $Z$ 处的爆破，
则称态射 $X' \to X$ 为{\it $U$-容许爆破}。
\end{definition}

\noindent
回顾：$Z \to X$ 有限表示，当且仅当理想层
$\mathcal{I}_Z \subset \mathcal{O}_X$ 为有限型，见《空间的态射》引理
\ref{spaces-morphisms-lemma-closed-immersion-finite-presentation}。特别地，
$U$-容许爆破是固有态射，见引理 \ref{spaces-divisors-lemma-blowing-up-projective}。注意，可能有多个
中心给出同一态射。因此，要求仅仅是存在某个与 $U$ 不交且产生 $X'$ 的中心。最后，
由于态射 $b : X' \to X$ 在 $U$ 上是同构（见引理
\ref{spaces-divisors-lemma-blowing-up-gives-effective-Cartier-divisor}），我们常滥用记号，也把 $U$
看作 $X'$ 的开子空间。

\begin{lemma}
\label{spaces-divisors-lemma-composition-admissible-blowups}
设 $S$ 为概形，$X$ 为 $S$ 上拟紧拟分离的代数空间，$U \subset X$ 为拟紧开子空间。
设 $b : X' \to X$ 为 $U$-容许爆破，$X'' \to X'$ 为 $U$-容许爆破。则复合态射
$X'' \to X$ 是 $U$-容许爆破。
\end{lemma}

\begin{proof}
立即由更精确的引理 \ref{spaces-divisors-lemma-composition-finite-type-blowups} 得出。
\end{proof}

\begin{lemma}
\label{spaces-divisors-lemma-extend-admissible-blowups}
设 $S$ 为概形。设
% P09-ERR-0238
$X$ 为拟紧拟分离代数空间，$U, V \subset X$ 为拟紧开子空间。设
$b : V' \to V$ 为 $U \cap V$-容许爆破。则存在 $U$-容许爆破 $X' \to X$，
其在 $V$ 上的限制为 $V'$。
\end{lemma}

\begin{proof}
设 $\mathcal{I} \subset \mathcal{O}_V$ 为有限型拟凝聚理想层，使得
$V(\mathcal{I})$ 与 $U \cap V$ 不交，且 $V'$ 同构于 $V$ 在 $\mathcal{I}$ 处的
爆破。设 $\mathcal{I}' \subset \mathcal{O}_{U \cup V}$ 为拟凝聚理想层，其在 $U$
上的限制为 $\mathcal{O}_U$，在 $V$ 上的限制为 $\mathcal{I}$。由《空间的极限》
引理 \ref{spaces-limits-lemma-extend}，存在有限型拟凝聚理想层
$\mathcal{J} \subset \mathcal{O}_X$，其在 $U \cup V$ 上的限制为
$\mathcal{I}'$。引理得证。
\end{proof}

\begin{lemma}
\label{spaces-divisors-lemma-dominate-admissible-blowups}
设 $S$ 为概形，$X$ 为 $S$ 上拟紧拟分离的代数空间，$U \subset X$ 为拟紧开子空间。
设
% P09-ERR-0239
$b_i : X_i \to X$（$i = 1, \ldots, n$）为 $U$-容许爆破。则存在
$U$-容许爆破 $b : X' \to X$，使得 (a) 对 $i = 1, \ldots, n$，$b$ 分解为
$X' \to X_i \to X$；且 (b) 每个态射 $X' \to X_i$ 都是 $U$-容许爆破。
\end{lemma}

\begin{proof}
设 $\mathcal{I}_i \subset \mathcal{O}_X$ 为有限型拟凝聚理想层，使得
$V(\mathcal{I}_i)$ 与 $U$ 不交，且 $X_i$ 同构于 $X$ 在
$\mathcal{I}_i$ 处的爆破。令
$\mathcal{I} = \mathcal{I}_1 \cdot \ldots \cdot \mathcal{I}_n$，并令 $X'$ 为
$X$ 在 $\mathcal{I}$ 处的爆破。则由引理 \ref{spaces-divisors-lemma-blowing-up-two-ideals}，
$X' \to X$ 经由 $b_i$ 分解。
\end{proof}

\begin{lemma}
\label{spaces-divisors-lemma-separate-disjoint-opens-by-blowing-up}
设 $S$ 为概形，$X$ 为 $S$ 上拟紧拟分离的代数空间，$U, V$ 为 $X$ 的两个不交
拟紧开子空间。则存在 $U \cup V$-容许爆破 $b : X' \to X$，使得 $X'$ 是开子空间
的不交并 $X' = X'_1 \amalg X'_2$，并且 $b^{-1}(U) \subset X'_1$、
$b^{-1}(V) \subset X'_2$。
\end{lemma}

\begin{proof}
选取有限型拟凝聚理想层 $\mathcal{I}$ 与 $\mathcal{J}$，分别使得
$X \setminus U = V(\mathcal{I})$ 与 $X \setminus V = V(\mathcal{J})$，见
《空间的极限》引理
\ref{spaces-limits-lemma-quasi-coherent-finite-type-ideals}。于是
$|V(\mathcal{I}\mathcal{J})| = |X|$，故 $\mathcal{I}\mathcal{J}$ 是局部幂零
理想层。由于 $\mathcal{I}$ 与 $\mathcal{J}$ 是有限型的，且 $X$ 拟紧，存在
$n > 0$ 使得 $\mathcal{I}^n \mathcal{J}^n = 0$。可以并且确实用
$\mathcal{I}^n$ 代替 $\mathcal{I}$，用 $\mathcal{J}^n$ 代替 $\mathcal{J}$。
于是 $\mathcal{I} \mathcal{J} = 0$。设 $b : X' \to X$ 为在
$\mathcal{I} + \mathcal{J}$ 处的爆破。由于
% P09-ERR-0240
$|V(\mathcal{I} + \mathcal{J})| = |X| \setminus |U| \cup |V|$，它是
$U \cup V$-容许的。我们将证明，如引理陈述中那样，$X'$ 是开子空间的不交并
$X' = X'_1 \amalg X'_2$。

\medskip\noindent
由于 $|V(\mathcal{I} + \mathcal{J})|$ 是 $|U \cup V|$ 的补集，可知
$V \cup U$ 在 $X'$ 中概形论稠密，见引理
\ref{spaces-divisors-lemma-blowing-up-gives-effective-Cartier-divisor} 及
\ref{spaces-divisors-lemma-complement-effective-Cartier-divisor}。因此，若存在这样的分解
$X' = X'_1 \amalg X'_2$ 为开闭子空间，则 $X'_1$ 是 $U$ 在 $X'$ 中的概形论闭包，
类似地，$X'_2$ 是 $V$ 在 $X'$ 中的概形论闭包。由于 $U \to X'$ 与
$V \to X'$ 拟紧，取概形论闭包与平展局部化可交换（《空间的态射》引理
\ref{spaces-morphisms-lemma-quasi-compact-scheme-theoretic-image}）。因此，为验证
$X'_1$ 与 $X'_2$ 的存在性，可以在 $X$ 上平展局部地进行。这把问题归约到概形
情形；该情形已在《除子》引理
\ref{divisors-lemma-separate-disjoint-opens-by-blowing-up} 的证明中处理。
\end{proof}
% Generated mechanically from the frozen spaces-over-fields.tex target.
% Do not edit this generated file; edit the bound locale source instead.

% OK, start here.
%

\setcounter{chapter}{71}
\chapter{域上的代数空间}



\phantomsection
\label{spaces-over-fields-section-phantom}


\section{引言}
\label{spaces-over-fields-section-introduction}

\noindent
本章是关于簇的一章在代数空间情形下的对应版本。代数空间的一份参考文献是
\cite{Kn}。


\section{约定}
\label{spaces-over-fields-section-conventions}

\noindent
始终假设所有概形都包含在一个大 fppf 位点 $\Sch_{fppf}$ 中，并且所考虑的每个环
$A$ 都满足：$\Spec(A)$（同构于）该大位点的一个对象。

\medskip\noindent
设 $S$ 为概形，$X$ 为 $S$ 上的代数空间。在本章及以下各章中，将用
$X \times_S X$ 表示 $X$ 与自身的积（在 $S$ 上代数空间的范畴中），而不用
$X \times X$。



\section{泛有限态射}
\label{spaces-over-fields-section-generically-finite}

\noindent
本节继续《良态空间》第 \ref{decent-spaces-section-generically-finite} 节中的讨论，
以及《簇》第 \ref{varieties-section-generically-finite} 节中关于代数空间态射的对应讨论。

\begin{lemma}
\label{spaces-over-fields-lemma-quasi-finite-in-codim-1}
设 $S$ 为概形。设 $f : X \to Y$ 为 $S$ 上代数空间的态射。假设 $f$ 局部有限型，
且 $Y$ 局部诺特。设 $y \in |Y|$ 为 $Y$ 上余维 $\leq 1$ 的点。设
$X^0 \subset |X|$ 为 $X$ 上余维 $0$ 的点所成的集合。再假设下列条件之一成立：
\begin{enumerate}
\item 对每个 $x \in X^0$，$x/f(x)$ 的超越次数为 $0$；
\item 对每个满足 $f(x) \leadsto y$ 的 $x \in X^0$，$x/f(x)$ 的超越次数为 $0$；
\item $f$ 在每个 $x \in X^0$ 处拟有限；
\item $f$ 在 $|X|$ 的一个稠密点集上拟有限；
% P09-ERR-0241
\item 此处需补充更多条件。
\end{enumerate}
则 $f$ 在 $X$ 中位于 $y$ 上方的每一点处拟有限。
\end{lemma}

\begin{proof}
我们要把证明归约到概形情形。为此，选取交换图
$$
\xymatrix{
U \ar[r] \ar[d]_g & X \ar[d]^f \\
V \ar[r] & Y
}
$$
其中 $U$, $V$ 为概形，水平箭头平展且满。选取映到 $y$ 的 $v \in V$。注意，
$V$ 局部诺特，并且 $\dim(\mathcal{O}_{V, v}) \leq 1$（见《空间的性质》定义
\ref{spaces-properties-definition-dimension-local-ring} 以及
注 \ref{spaces-properties-remark-list-properties-local-etale-topology}）。
$U \to V$ 在 $v$ 上的纤维 $U_v$ 满射到
$f^{-1}(\{y\}) \subset |X|$。$X^0$ 在 $U$ 中的逆像恰为 $U$ 的不可约分支的一般点集
（《空间的性质》引理 \ref{spaces-properties-lemma-codimension-0-points}）。
若 $\eta \in U$ 是这样的点，其像为 $x \in X^0$，则 $x / f(x)$ 的超越次数就是
$\kappa(\eta)$ 在 $\kappa(g(\eta))$ 上的超越次数（《空间的态射》定义
\ref{spaces-morphisms-definition-dimension-fibre}）。注意，$U \to V$ 在
$u \in U$ 处拟有限，当且仅当 $f$ 在 $u$ 于 $X$ 中的像处拟有限。

\medskip\noindent
情形 (1)。此时可应用《簇》引理
\ref{varieties-lemma-quasi-finite-in-codim-1} 的情形 (1)，从而 $U \to V$
在 $U_v$ 的所有点处拟有限。因此，$f$ 在位于 $y$ 上方的每一点处拟有限。

\medskip\noindent
情形 (2)。设 $u \in U$ 为某个不可约分支的一般点，且它在 $V$ 中的像特化到 $v$。
则 $u$ 的像 $x \in X^0$ 满足 $f(x) \leadsto y$。故可应用《簇》引理
\ref{varieties-lemma-quasi-finite-in-codim-1} 的情形 (2)，并如前得出结论。

\medskip\noindent
情形 (3) 由《簇》引理 \ref{varieties-lemma-quasi-finite-in-codim-1}
的情形 (3) 得出。

\medskip\noindent
在情形 (4) 中，由于 $|U| \to |X|$ 是开映射，$U \to V$ 拟有限的点集也稠密。
因此，可应用《簇》引理 \ref{varieties-lemma-quasi-finite-in-codim-1}
的情形 (4)。
\end{proof}

\begin{lemma}
\label{spaces-over-fields-lemma-finite-in-codim-1}
设 $S$ 为概形。设 $f : X \to Y$ 为 $S$ 上代数空间的态射。假设 $f$ 固有，且
$Y$ 局部诺特。设 $y \in Y$ 为 $Y$ 中余维 $\leq 1$ 的点。设
$X^0 \subset |X|$ 为 $X$ 上余维 $0$ 的点所成的集合。再假设下列条件之一成立：
\begin{enumerate}
\item 对每个 $x \in X^0$，$x/f(x)$ 的超越次数为 $0$；
\item 对每个满足 $f(x) \leadsto y$ 的 $x \in X^0$，$x/f(x)$ 的超越次数为 $0$；
\item $f$ 在每个 $x \in X^0$ 处拟有限；
\item $f$ 在 $|X|$ 的一个稠密点集上拟有限；
% P09-ERR-0241
\item 此处需补充更多条件。
\end{enumerate}
则存在包含 $y$ 的开子空间 $Y' \subset Y$，使得
$Y' \times_Y X \to Y'$ 有限。
\end{lemma}

\begin{proof}
由引理 \ref{spaces-over-fields-lemma-quasi-finite-in-codim-1}，态射 $f$ 在位于 $y$ 上方的每一点处
拟有限。设 $\overline{y} : \Spec(k) \to Y$ 为位于 $y$ 上方的几何点。则
$|X_{\overline{y}}|$ 是离散空间（《良态空间》引理
\ref{decent-spaces-lemma-conditions-on-fibre-and-qf}）。由于 $f$ 固有，
$X_{\overline{y}}$ 拟紧，从而 $|X_{\overline{y}}|$ 有限。因此，可应用
《空间的上同调》引理
\ref{spaces-cohomology-lemma-proper-finite-fibre-finite-in-neighbourhood}
得出结论。
\end{proof}

\begin{lemma}
\label{spaces-over-fields-lemma-modification-normal-iso-over-codimension-1}
设 $S$ 为概形，$X$ 为 $S$ 上的诺特代数空间。设 $f : Y \to X$ 为代数空间的
双有理固有态射，并且 $Y$ 约化。设 $U \subset X$ 为使 $f$ 在其上为同构的最大开子空间。
则 $U$ 包含：
\begin{enumerate}
\item $X$ 中每个余维 $0$ 的点；
\item 每个在 $X$ 上余维为 $1$ 且 $X$ 在 $x$ 处的局部环正规的
$x \in |X|$（《空间的性质》注
\ref{spaces-properties-remark-list-properties-local-ring-local-etale-topology}）；以及
\item 每个满足 $|Y| \to |X|$ 在 $x$ 上的纤维有限且 $X$ 在 $x$ 处的局部环正规的
$x \in |X|$。
\end{enumerate}
\end{lemma}

\begin{proof}
(1) 由《良态空间》引理
\ref{decent-spaces-lemma-birational-isomorphism-over-dense-open} 得出（还用到诺特代数空间
$X$ 与 $Y$ 拟分离，从而良态）。(2) 由 (3) 和引理
\ref{spaces-over-fields-lemma-finite-in-codim-1} 得出（还用到有限态射具有有限纤维）。设
$x \in |X|$ 如 (3) 所述。由《空间的上同调》引理
\ref{spaces-cohomology-lemma-proper-finite-fibre-finite-in-neighbourhood}
（由《良态空间》引理 \ref{decent-spaces-lemma-conditions-on-fibre-and-qf}
可知其适用），可假设 $f$ 有限。选取仿射概形 $X'$、平展态射 $X' \to X$，以及映到 $x$ 的点
% P09-ERR-0242
$x' \in X$。只需证明：存在 $x' \in X'$ 的开邻域 $U'$，使得
$Y \times_X X' \to X'$ 在 $U'$ 上为同构（此时 $U$ 包含 $U'$ 在 $X$ 中的像，见
《空间》引理 \ref{spaces-lemma-descent-representable-transformations-property}）。于是
% P09-ERR-0243
$Y \times_X X' \to X$ 是有限双有理态射（《良态空间》引理
\ref{decent-spaces-lemma-birational-etale-localization}）。由于有限态射是仿射的，问题归约为
诺特仿射概形的有限双有理态射 $Y \to X$，其中 $x \in X$ 且
$\mathcal{O}_{X, x}$ 是正规整环。该情形已在《簇》引理
\ref{varieties-lemma-modification-normal-iso-over-codimension-1} 中处理。
\end{proof}






\section{整代数空间}
\label{spaces-over-fields-section-integral-spaces}

\noindent
尚未定义整代数空间的概念。问题在于，整性不是概形的平展局部性质。可以用《性质》
引理 \ref{properties-lemma-characterize-integral} 中的性质，即 $X$ 约化且 $|X|$
不可约，来定义整代数空间。但这样一来，《空间》例
\ref{spaces-example-infinite-product} 中描述的代数空间将是整的，这似乎不妥。
为避免这种病态情形，还将假设 $X$ 是良态代数空间，尽管也许存在更弱的替代条件。

\begin{definition}
\label{spaces-over-fields-definition-integral-algebraic-space}
设 $S$ 为概形。若 $S$ 上的代数空间 $X$ 约化、良态且 $|X|$ 不可约，则称其为
{\it 整的}。
\end{definition}

\noindent
在此情形下，不可约拓扑空间 $|X|$ 是良态拓扑空间（《良态空间》命题
\ref{decent-spaces-proposition-reasonable-sober}），故它有唯一的一般点 $x$。
事实上，我们已在《良态空间》引理
\ref{decent-spaces-lemma-finitely-many-irreducible-components} 中刻画了具有有限多个
不可约分支的良态代数空间。应用该引理可知，若代数空间 $X$ 约化，具有不可约稠密
开子概形 $X'$，其一般点为 $x'$，且态射 $x' \to X$ 拟紧，则它是整的。

\begin{lemma}
\label{spaces-over-fields-lemma-integral-algebraic-space-rational-functions}
设 $S$ 为概形，$X$ 为 $S$ 上的整代数空间。设 $\eta \in |X|$ 为 $X$ 的一般点。
有典范等同
$$
R(X) = \mathcal{O}_{X, \eta}^h = \kappa(\eta)
$$
其中 $R(X)$ 是《空间的态射》定义
\ref{spaces-morphisms-definition-ring-of-rational-functions},
中定义的有理函数环，$\kappa(\eta)$ 是《良态空间》定义
\ref{decent-spaces-definition-residue-field} 中定义的剩余域，而
$\mathcal{O}_{X, \eta}^h$ 是《良态空间》定义
\ref{decent-spaces-definition-elemenary-etale-neighbourhood} 中定义的亨泽尔局部环。
特别地，这些环都是域。
\end{lemma}

\begin{proof}
由于 $X$ 在 $\eta$ 的某个开邻域中是概形（见上文讨论），这立即由概形的对应结果得出，
见《态射》引理 \ref{morphisms-lemma-integral-scheme-rational-functions}。还用到：
域的亨泽尔化是其自身，并且这些对象关于代数空间的定义与概形情形的定义相容。
细节从略。
\end{proof}

\noindent
这引出以下定义。

\begin{definition}
\label{spaces-over-fields-definition-function-field}
设 $S$ 为概形，$X$ 为 $S$ 上的整代数空间。$X$ 的{\it 函数域}，或
{\it 有理函数域}，是引理
\ref{spaces-over-fields-lemma-integral-algebraic-space-rational-functions} 中的域 $R(X)$。
\end{definition}

\noindent
有时也用 $k(X)$ 而非 $R(X)$ 表示该域。

\begin{lemma}
\label{spaces-over-fields-lemma-integral-sections}
设 $S$ 为概形，$X$ 为 $S$ 上的整代数空间。则
$\Gamma(X, \mathcal{O}_X)$ 是整环。
\end{lemma}

\begin{proof}
令 $R = \Gamma(X, \mathcal{O}_X)$。若 $f, g \in R$ 非零且 $fg = 0$，则
$X = V(f) \cup V(g)$，其中 $V(f)$ 表示由 $f$ 截出的 $X$ 的闭子空间。由于
$X$ 不可约，必有 $V(f) = X$ 或 $V(g) = X$。于是由《空间的性质》引理
\ref{spaces-properties-lemma-reduced-space}，必有 $f = 0$ 或 $g = 0$。
\end{proof}

\noindent
下面是关于正规整代数空间的一个引理。

\begin{lemma}
\label{spaces-over-fields-lemma-normal-integral-cover-by-affines}
设 $S$ 为概形，$X$ 为 $S$ 上的正规整代数空间。对每个 $x \in |X|$，存在正规整
仿射概形 $U$ 和平展态射 $U \to X$，使得 $x$ 属于其像。
\end{lemma}

\begin{proof}
选取仿射概形 $U$ 和平展态射 $U \to X$，使得 $x$ 属于其像。设
% P09-ERR-0244
$u_i$（$i \in I$）为 $U$ 的不可约分支的一般点。则每个 $u_i$ 都映到 $X$ 的一般点
（《良态空间》引理 \ref{decent-spaces-lemma-decent-generic-points}）。由良态空间的定义
（《良态空间》定义 \ref{decent-spaces-definition-very-reasonable}），$I$ 有限。因此
$U = \Spec(A)$，其中 $A$ 是具有有限多个极小素理想的正规环。由《代数》引理
\ref{algebra-lemma-characterize-reduced-ring-normal}，
$A = \prod_{i \in I} A_i$ 是正规整环的乘积。于是 $U = \coprod U_i$，其中
$U_i = \Spec(A_i)$；对某个 $i$，$x$ 属于 $U_i \to X$ 的像。引理得证。
\end{proof}

\begin{lemma}
\label{spaces-over-fields-lemma-normal-integral-sections}
设 $S$ 为概形，$X$ 为 $S$ 上的正规整代数空间。则
$\Gamma(X, \mathcal{O}_X)$ 是正规整环。
\end{lemma}

\begin{proof}
令 $R = \Gamma(X, \mathcal{O}_X)$。由引理 \ref{spaces-over-fields-lemma-integral-sections}，
$R$ 是整环。设 $f = a/b$ 为 $R$ 的分式域中在 $R$ 上整的元素。对任意平展态射
$U \to X$，其中 $U$ 是概形，至多存在一个
$f_U \in \Gamma(U, \mathcal{O}_U)$ 满足 $b|_U f_U = a|_U$。事实上，$U$
约化，且 $U$ 的一般点映到 $X$ 的一般点，这蕴含 $b|_U$ 是非零因子。对每个
$x \in |X|$，按引理 \ref{spaces-over-fields-lemma-normal-integral-cover-by-affines} 选取
$U \to X$。由《性质》引理 \ref{properties-lemma-normal-integral-sections}，
$\Gamma(U, \mathcal{O}_U)$ 是正规整环，故存在唯一的
$f_U \in \Gamma(U, \mathcal{O}_U)$ 满足 $b|_U f_U = a|_U$。由上述唯一性，
这些 $f_U$ 可黏合，并定义结构层的整体截面 $f$，即 $R$ 的元素。
\end{proof}

\begin{lemma}
\label{spaces-over-fields-lemma-decent-irreducible-closed}
设 $S$ 为概形，$X$ 为 $S$ 上的良态代数空间。下列集合之间存在典范双射：
\begin{enumerate}
\item $X$ 的点集，即 $|X|$；
\item $|X|$ 的不可约闭子集所成的集合；
\item $X$ 的整的闭子空间所成的集合。
\end{enumerate}
从 (1) 到 (2) 的双射把 $x$ 映到 $\overline{\{x\}}$；从 (3) 到 (2) 的双射把
$Z$ 映到 $|Z|$。
\end{lemma}

\begin{proof}
由《良态空间》命题 \ref{decent-spaces-proposition-reasonable-sober}，$|X|$
是良态拓扑空间，故上述映射给出 (1) 与 (2) 之间的双射。给定闭且不可约的
$T \subset |X|$，存在唯一的约化闭子空间 $Z \subset X$ 使得 $|Z| = T$；
即 $Z$ 是 $T$ 上的约化诱导子空间结构，见《空间的性质》定义
\ref{spaces-properties-definition-reduced-induced-space}。它是整代数空间，因为它良态、
约化且不可约。
\end{proof}






\section{整代数空间之间的态射}
\label{spaces-over-fields-section-morphisms-between-integral-spaces}

\noindent
以下引理刻画整代数空间之间有限次数的支配态射。

\begin{lemma}
\label{spaces-over-fields-lemma-finite-degree}
设 $S$ 为概形。设
% P09-ERR-0245
$X$, $Y$ 为 $S$ 上的整代数空间。设 $x \in |X|$ 与 $y \in |Y|$ 为一般点。
设 $f : X \to Y$ 局部有限型。假设 $f$ 支配（《空间的态射》定义
\ref{spaces-morphisms-definition-dominant}）。下列条件等价：
\begin{enumerate}
\item $x/y$ 的超越次数为 $0$；
\item 扩张 $\kappa(x)/\kappa(y)$（见证明）有限；
\item 存在非空仿射开集 $U \subset X$ 和 $V \subset Y$，使得
$f(U) \subset V$ 且 $f|_U : U \to V$ 有限；
\item $f$ 在 $x$ 处拟有限；
\item $x$ 是 $|X|$ 中映到 $y$ 的唯一点。
\end{enumerate}
若 $f$ 分离或 $f$ 拟紧，则这些条件还等价于：
\begin{enumerate}
\item[(6)] 存在非空仿射开集 $V \subset Y$，使得 $f^{-1}(V) \to V$ 有限。
\end{enumerate}
\end{lemma}

\begin{proof}
由初等拓扑，因 $f$ 支配，有 $f(x) = y$。设 $Y' \subset Y$ 为 $Y$ 的概形轨迹，
设 $X' \subset f^{-1}(Y')$ 为 $f^{-1}(Y')$ 的概形轨迹。由上文讨论，并用到
《良态空间》命题 \ref{decent-spaces-proposition-reasonable-sober} 及定理
\ref{decent-spaces-theorem-decent-open-dense-scheme}，可知 $x \in |X'|$ 且
$y \in |Y'|$。于是 $f|_{X'} : X' \to Y'$ 是局部有限型的整概形之间的态射。
因此，由《态射》引理 \ref{morphisms-lemma-finite-degree}，(1)、(2)、(3) 等价。

\medskip\noindent
对 $X \to Y \to Y$ 应用《空间的态射》引理
\ref{spaces-morphisms-lemma-compare-tr-deg}，条件 (4) 蕴含条件 (1)。另一方面，
有限态射拟有限，并且因为 $x$ 是一般点而有 $x \in U$，故条件 (3) 蕴含条件 (4)。
所以 (1) -- (4) 等价。

\medskip\noindent
假设等价条件 (1) -- (4) 成立。设 $x' \mapsto y$。则 $x \leadsto x'$ 是
$|X| \to |Y|$ 在 $y$ 上的纤维中的特化。若 $x' \not = x$，则由《良态空间》引理
\ref{decent-spaces-lemma-conditions-on-point-in-fibre-and-qf}，$f$ 在 $x$ 处不拟有限。
故 $x = x'$，条件 (5) 成立。反之，若 (5) 成立，则概形态射 $X' \to Y'$
（见上文）也满足 (5)，再由《态射》引理 \ref{morphisms-lemma-finite-degree}
可知 (1) 成立。

\medskip\noindent
注意，无需对 $f$ 作任何进一步假设，(6) 就蕴含等价条件 (1) -- (5)。为完成证明，
只需说明等价条件 (1) -- (5) 蕴含 (6)。这由《良态空间》引理
\ref{decent-spaces-lemma-finite-over-dense-open} 得出。
\end{proof}

\begin{definition}
\label{spaces-over-fields-definition-degree}
设 $S$ 为概形。设 $X$ 与 $Y$ 为 $S$ 上的整代数空间。设 $f : X \to Y$
局部有限型且支配。假设引理 \ref{spaces-over-fields-lemma-finite-degree} 的等价条件 (1) -- (5)
中任一条件成立。设 $x \in |X|$ 与 $y \in |Y|$ 为一般点。正整数
$$
\deg(X/Y) = [\kappa(x) : \kappa(y)]
$$
称为{\it $X$ 在 $Y$ 上的次数}。
\end{definition}

\begin{lemma}
\label{spaces-over-fields-lemma-degree-composition}
设 $S$ 为概形。设 $X$, $Y$, $Z$ 为 $S$ 上的整代数空间。设 $f : X \to Y$ 与
$g : Y \to Z$ 为局部有限型的支配态射。假设引理
\ref{spaces-over-fields-lemma-finite-degree} 的等价条件 (1) -- (5) 中任一条件对 $f$ 和 $g$ 成立。则
$$
\deg(X/Z) = \deg(X/Y) \deg(Y/Z).
$$
\end{lemma}

\begin{proof}
这由域的有限扩张塔中次数的乘法性得出，见《域》引理
\ref{fields-lemma-multiplicativity-degrees}。
\end{proof}










\section{Weil 除子}
\label{spaces-over-fields-section-Weil-divisors}

\noindent
本节是《除子》第 \ref{divisors-section-Weil-divisors} 节的对应版本。

\medskip\noindent
我们将对局部诺特整代数空间引入 Weil 除子以及 Weil 除子的有理等价。由于不假设
代数空间拟紧，定义 Weil 除子时必须稍加谨慎。例如，有理函数可能有无穷多个极点，
所以必须允许素除子的无穷和。在拟紧情形下，Weil 除子照常是有限和。下面的基本引理
将经常用于证明闭子空间族局部有限。

\begin{lemma}
\label{spaces-over-fields-lemma-components-locally-finite}
设 $S$ 为概形，$X$ 为 $S$ 上局部诺特的代数空间。若 $T \subset |X|$ 为闭子集，
则 $T$ 的不可约分支族局部有限。
\end{lemma}

\begin{proof}
拓扑空间 $|X|$ 局部诺特（《空间的性质》引理
\ref{spaces-properties-lemma-Noetherian-topology}）。诺特拓扑空间只有有限多个不可约分支，
且诺特空间的子空间仍为诺特空间（《拓扑》引理
\ref{topology-lemma-Noetherian}）。因此，本引理由局部有限的定义得出（《拓扑》定义
\ref{topology-definition-locally-finite}）。
\end{proof}

\noindent
设 $S$ 为概形，$X$ 为 $S$ 上的良态代数空间。设 $Z$ 为 $X$ 的整的闭子空间，
$\xi \in |Z|$ 为一般点。由《良态空间》引理
\ref{decent-spaces-lemma-codimension-local-ring}，$|Z|$ 在 $|X|$ 中的余维等于
$X$ 在 $\xi$ 处的局部环之维数。回顾：这也表述为{\it $\xi$ 是 $X$ 上余维为
$1$ 的点}，见《空间的性质》定义
\ref{spaces-properties-definition-dimension-local-ring}。

\begin{definition}
\label{spaces-over-fields-definition-Weil-divisor}
设 $S$ 为概形。设 $X$ 为 $S$ 上局部诺特的整代数空间。
\begin{enumerate}
\item {\it 素除子}是余维为 $1$ 的整的闭子空间 $Z \subset X$；即 $|Z|$ 的一般点
是 $X$ 上余维为 $1$ 的点。
\item {\it Weil 除子}是形式和 $D = \sum n_Z Z$，其中对 $X$ 的素除子求和，并且族
$\{|Z| : n_Z \not = 0\}$ 在 $|X|$ 中局部有限（《拓扑》定义
\ref{topology-definition-locally-finite}）。
\end{enumerate}
$X$ 上所有 Weil 除子所成的群记为 $\text{Div}(X)$。
\end{definition}

\noindent
下一步要定义与有理函数相伴的 Weil 除子。为此，需要定义局部诺特整代数空间 $X$
上的有理函数沿素除子 $Z$ 的消没阶。设 $\xi \in |Z|$ 为一般点。这里遇到的问题是：
局部环 $\mathcal{O}_{X, \xi}$ 不存在，而亨泽尔局部环
$\mathcal{O}_{X, \xi}^h$ 可能不是整环，见例 \ref{spaces-over-fields-example-not-a-domain}。
为绕开这一问题，使用以下引理。

\begin{lemma}
\label{spaces-over-fields-lemma-order-vanishing}
设 $S$ 为概形，$X$ 为 $S$ 上局部诺特的整代数空间。设 $Z \subset X$ 为素除子，
$\xi \in |Z|$ 为一般点。则亨泽尔局部环 $\mathcal{O}_{X, \xi}^h$
是约化的 $1$ 维诺特局部环，并且存在典范单射
$$
R(X) \longrightarrow Q(\mathcal{O}_{X, \xi}^h)
$$
从 $X$ 的函数域 $R(X)$ 到全分式环。
\end{lemma}

\begin{proof}
将使用《良态空间》第
\ref{decent-spaces-section-residue-fields-henselian-local-rings} 节的结果。设
$(U, u) \to (X, \xi)$ 为初等平展邻域。注意，$U$ 局部诺特且约化。因此，
$\mathcal{O}_{U, u}$ 是约化的 $1$ 维诺特环（其维数来自素除子的定义）。用 $u$ 的一个
仿射开邻域替换 $U$ 后，可假设 $U$ 诺特且仿射。再用更小的开集替换 $U$ 后，可假设
$U$ 的每个不可约分支都经过 $u$。由于 $U \to X$ 是开映射而 $X$ 不可约，
$U \to X$ 支配。因此，通过复合有理映射得到环同态 $R(X) \to R(U)$，见
《空间的态射》第 \ref{spaces-morphisms-section-rational-maps} 节。由于 $R(X)$ 是域，
该同态为单射。由对 $U$ 的选取，$R(U)$ 是全分式环
$Q(\mathcal{O}_{U, u})$，见《态射》引理
\ref{morphisms-lemma-integral-scheme-rational-functions} 及《代数》引理
\ref{algebra-lemma-total-ring-fractions-no-embedded-points}。

\medskip\noindent
至此，已证明把引理中的 $\mathcal{O}_{X, \xi}^h$ 换成
% P09-ERR-0246
$\mathcal{O}_{U, u}$ 后所得的全部陈述。然而，$\mathcal{O}_{X, \xi}^h$ 是
$\mathcal{O}_{U, u}$ 的亨泽尔化。因此，$\mathcal{O}_{X, \xi}^h$ 是约化的 $1$ 维诺特环，
见《更多代数》诸引理
\ref{more-algebra-lemma-henselization-reduced},
\ref{more-algebra-lemma-henselization-dimension} 及
\ref{more-algebra-lemma-henselization-noetherian}.
由《更多代数》引理 \ref{more-algebra-lemma-dumb-properties-henselization}，
$\mathcal{O}_{U, u} \to \mathcal{O}_{X, \xi}^h$ 忠实平坦，故它把非零因子映为非零因子。
因此得到典范同态
$Q(\mathcal{O}_{U, u}) \to Q(\mathcal{O}_{X, \xi}^h)$，进而得到所需同态。
略去该同态与下述邻域选取无关的验证：
% P09-ERR-0247
$(U, u) \to (X, x)$；稍好的做法是先注意到
$\colim Q(\mathcal{O}_{U, u}) = Q(\mathcal{O}_{X, \xi}^h)$.
\end{proof}

\begin{definition}
\label{spaces-over-fields-definition-order-vanishing}
设 $S$ 为概形，$X$ 为 $S$ 上局部诺特的整代数空间。设 $f \in R(X)^*$。对每个
素除子 $Z \subset X$，把{\it $f$ 沿 $Z$ 的消没阶}定义为整数
$$
\text{ord}_Z(f) =
\text{length}_{\mathcal{O}_{X, \xi}^h}
(\mathcal{O}_{X, \xi}^h/a \mathcal{O}_{X, \xi}^h) -
\text{length}_{\mathcal{O}_{X, \xi}^h}
(\mathcal{O}_{X, \xi}^h/b \mathcal{O}_{X, \xi}^h)
$$
其中 $a, b \in \mathcal{O}_{X, \xi}^h$ 是非零因子，并且 $f$ 在
$Q(\mathcal{O}_{X, \xi}^h)$ 中的像（引理 \ref{spaces-over-fields-lemma-order-vanishing}）等于 $a/b$。
由《代数》引理 \ref{algebra-lemma-ord-additive}，该定义是良定义的。
\end{definition}

\noindent
若 $\mathcal{O}_{X, \xi}^h$ 恰为整环，则有
$$
\text{ord}_Z(f) = \text{ord}_{\mathcal{O}_{X, \xi}^h}(f)
$$
其中右端是《代数》定义 \ref{algebra-definition-ord} 中的概念。注意，对
$f, g \in R(X)^*$，有
$$
\text{ord}_Z(fg) = \text{ord}_Z(f) + \text{ord}_Z(g).
$$
当然，可能有 $\text{ord}_Z(f) < 0$。在此情形下，称 $f$ 沿 $Z$ 有{\it 极点}，
并称 $-\text{ord}_Z(f) > 0$ 为{\it $f$ 沿 $Z$ 的极点阶数}。必须注意，条件
$\text{ord}_Z(f) \geq 0$ 与条件 $f \in \mathcal{O}_{X, \xi}^h$
{\bf 并不}等价，除非局部环
% P09-ERR-0248
$\mathcal{O}_{X, \xi}$ 是离散赋值环。

\begin{lemma}
\label{spaces-over-fields-lemma-order-vanishing-agrees}
设 $S$ 为概形，$X$ 为 $S$ 上局部诺特的整代数空间。设 $f \in R(X)^*$。若素除子
$Z \subset X$ 与 $X$ 的概形轨迹相交，则定义 \ref{spaces-over-fields-definition-order-vanishing}
中的消没阶 $\text{ord}_Z(f)$ 与《除子》定义
\ref{divisors-definition-order-vanishing} 中的消没阶一致。
\end{lemma}

\begin{proof}
缩小 $X$ 后，可假设 $X$ 是整诺特概形。若 $\xi \in Z$ 表示一般点，则
$\mathcal{O}_{X, \xi}^h$ 是 $\mathcal{O}_{X, \xi}$ 的亨泽尔化（《良态空间》引理
\ref{decent-spaces-lemma-describe-henselian-local-ring}）。为证明引理，以下等式既充分又必要：
$$
\text{length}_{\mathcal{O}_{X, \xi}}
(\mathcal{O}_{X, \xi}/a \mathcal{O}_{X, \xi}) =
\text{length}_{\mathcal{O}_{X, \xi}^h}
(\mathcal{O}_{X, \xi}^h/a \mathcal{O}_{X, \xi}^h)
$$
这立即由《代数》引理 \ref{algebra-lemma-pullback-module} 得出（并用到
$\mathcal{O}_{X, \xi} \to \mathcal{O}_{X, \xi}^h$ 是诺特局部环之间的平坦局部环同态）。
\end{proof}

\begin{lemma}
\label{spaces-over-fields-lemma-divisor-locally-finite}
设 $S$ 为概形，$X$ 为 $S$ 上局部诺特的整代数空间。设 $f \in R(X)^*$。则族
% P09-ERR-0249
$$
\{Z \subset X \mid Z\text{ 为素除子，其一般点为 }\xi
\text{ 且 }f\text{ 不属于 }\mathcal{O}_{X, \xi}\}
$$
以及
$$
\{Z \subset X \mid Z \text{ 为素除子且 }\text{ord}_Z(f) \not = 0\}
$$
在 $X$ 中局部有限。
\end{lemma}

\begin{proof}
存在非空开子空间 $U \subset X$，使 $f$ 对应于
$\Gamma(U, \mathcal{O}_X^*)$ 的一个截面。因此，引理中的集合里可能出现的素除子
全都对应于 $|X| \setminus |U|$ 的不可约分支。于是引理
\ref{spaces-over-fields-lemma-components-locally-finite} 给出所需结论。
\end{proof}

\noindent
该引理使我们可以作如下定义。

\begin{definition}
\label{spaces-over-fields-definition-principal-divisor}
设 $S$ 为概形。设 $X$ 为 $S$ 上局部诺特的整代数空间。设 $f \in R(X)^*$。
{\it 与 $f$ 相伴的主 Weil 除子}是 Weil 除子
$$
\text{div}(f) = \text{div}_X(f) = \sum \text{ord}_Z(f) [Z]
$$
其中对素除子求和，而 $\text{ord}_Z(f)$ 如定义
\ref{spaces-over-fields-definition-order-vanishing} 所述。由引理
\ref{spaces-over-fields-lemma-divisor-locally-finite}，这是有意义的。
\end{definition}

\begin{lemma}
\label{spaces-over-fields-lemma-div-additive}
设 $S$ 为概形。设 $X$ 为 $S$ 上局部诺特的整代数空间。设
$f, g \in R(X)^*$。则
$$
\text{div}_X(fg) = \text{div}_X(f) + \text{div}_X(g)
$$
作为 $X$ 上的 Weil 除子成立。
\end{lemma}

\begin{proof}
这由 $\text{ord}$ 函数的可加性显然得出。
\end{proof}

\noindent
% P09-ERR-0250
由上述引理可见，主 Weil 除子所成的集合构成所有 Weil 除子之群的一个子群。
这引出以下定义。

\begin{definition}
\label{spaces-over-fields-definition-class-group}
设 $S$ 为概形，$X$ 为 $S$ 上局部诺特的整代数空间。$X$ 的
{\it Weil 除子类群}是 Weil 除子群对主 Weil 除子子群的商，记作
$\text{Cl}(X)$。
\end{definition}

\noindent
由构造得到正合复形
\begin{equation}
\label{spaces-over-fields-equation-Weil-divisor-class}
R(X)^* \xrightarrow{\text{div}} \text{Div}(X) \to \text{Cl}(X) \to 0
\end{equation}
可把它看作 $\text{Cl}(X)$ 的一个表示。下一步要把 Weil 除子类群与 Picard 群联系起来。

\begin{example}
\label{spaces-over-fields-example-Z-mod-2}
这是《空间的态射》例
\ref{spaces-morphisms-example-universal-homeomorphism} 的继续。考虑代数空间
$X = \mathbf{A}^1_k/\{t \sim -t \mid t \not = 0\}$.
它是域 $k$ 上的光滑代数空间。存在普遍同胚
$$
X \longrightarrow \mathbf{A}^1_k = \Spec(k[t])
$$
它在 $\mathbf{A}^1_k \setminus \{0\}$ 上为同构。由此 $X$ 诺特且为整的。由于
$\dim(X) = 1$，$X$ 的素除子就是 $X$ 的闭点。考虑位于
$0 \in \mathbf{A}^1_k$ 上方的唯一闭点 $x \in |X|$。由于
% P09-ERR-0251
$X \setminus \{x\}$ 同构地映到 $\mathbf{A}^1 \setminus \{0\}$，
$\text{Cl}(X)$ 中不同于 $x$ 的闭点之类均为零。然而，$t$ 在 $X$ 上的除子为
$2[x]$。故 $\text{Cl}(X) = \mathbf{Z}/2\mathbf{Z}$。
\end{example}

\begin{example}
\label{spaces-over-fields-example-not-a-domain}
设 $k$ 为域。令
$$
U = \Spec(k[x, y]/(xy))
$$
为 $\mathbf{A}^2_k$ 中两条坐标轴的并。
% P09-ERR-0252
记 $\Delta : U \to U \times_k U$ 为对角态射，并记
$\Delta' : U \to U \times_k U$ 为映射 $u \mapsto (u, \sigma(u))$，其中
$\sigma : U \to U$，$(x, y) \mapsto (y, x)$，是交换两条坐标轴的自同构。令
$$
R = \Delta(U) \amalg \Delta'(U \setminus \{0_U\})
$$
其中 $0_U \in U$ 为原点。容易看出，$R$ 是 $U$ 上的平展等价关系。商
$X = U/R$ 是代数空间。态射 $U \to \mathbf{A}^1_k$，
$(x, y) \mapsto x + y$，是 $R$-不变的，故定义态射
$$
X \longrightarrow \mathbf{A}^1_k
$$
该态射是普遍同胚，并且在 $\mathbf{A}^1_k \setminus \{0\}$ 上为同构。因此，
$X$ 整且诺特。与例 \ref{spaces-over-fields-example-Z-mod-2} 完全相同，读者可证
$\text{Cl}(X) = \mathbf{Z}/2\mathbf{Z}$，其生成元对应于映到
$0 \in \mathbf{A}^1_k$ 的唯一闭点 $x \in |X|$。然而，在此情形下，$X$ 在 $x$
处的亨泽尔局部环不是整环，因为它是 $\mathcal{O}_{U, 0_U}$ 的亨泽尔化。
\end{example}














\section{与可逆模相伴的 Weil 除子类}
\label{spaces-over-fields-section-c1}

\noindent
本节完全按照第 \ref{spaces-over-fields-section-Weil-divisors} 节相同的步骤，在局部诺特整代数空间上
定义典范映射 $\Pic(X) \to \text{Cl}(X)$。

\medskip\noindent
设 $S$ 为概形，$X$ 为 $S$ 上局部诺特的整代数空间，$\mathcal{L}$ 为可逆
$\mathcal{O}_X$-模。由《空间上的除子》引理
\ref{spaces-divisors-lemma-regular-meromorphic-section-exists}，存在正则亚纯截面
$s \in \Gamma(X, \mathcal{K}_X(\mathcal{L}))$。事实上，由《空间上的除子》引理
\ref{spaces-divisors-lemma-compute-meromorphic}，这等价于
$\mathcal{L}_\eta$ 中的非零元素，其中 $\eta \in |X|$ 为一般点。同一引理还说明，
若 $\mathcal{L} = \mathcal{O}_X$，则 $s$ 等价于 $X$ 上的非零有理函数（故以下做法
与第 \ref{spaces-over-fields-section-Weil-divisors} 节的构造一致）。

\medskip\noindent
设 $Z \subset X$ 为素除子，$\xi \in |Z|$ 为一般点。将定义 $s$ 沿 $Z$ 的消没阶。
考虑典范态射
$$
c_\xi : \Spec(\mathcal{O}_{X, \xi}^h) \longrightarrow X
$$
其源是 $X$ 在 $\xi$ 处的亨泽尔局部环之谱（《良态空间》定义
% P09-ERR-0253
\ref{decent-spaces-definition-henselian-local-ring}）。拉回
$\mathcal{L}_\xi = c_\xi^*\mathcal{L}$ 是可逆模，因而平凡；选取
$\mathcal{L}_\xi$ 的一个生成元 $s_\xi$。由于 $c_\xi$ 平坦，亚纯函数及（正则）截面对
$c_\xi$ 的拉回有定义，见
《空间上的除子》定义
\ref{spaces-divisors-definition-pullback-meromorphic-sections} 以及
引理 \ref{spaces-divisors-lemma-pullback-meromorphic-sections-defined} 及
\ref{spaces-divisors-lemma-meromorphic-sections-pullback}.
因此得到
$$
c_\xi^*(s) = f s_\xi
$$
其中 $f \in Q(\mathcal{O}_{X, \xi}^h)$ 是某个非零因子。这里使用《除子》引理
\ref{divisors-lemma-locally-Noetherian-K}，借助 $\mathcal{O}_{X, \xi}^h$
的全分式环来表示
$\mathcal{L}_\xi \cong \mathcal{O}_{\Spec(\mathcal{O}_{X, \xi}^h)}$
的亚纯截面空间。约定把该元素记为
$$
s/s_\xi = f \in  Q(\mathcal{O}_{X, \xi}^h)
$$
注意，若改变 $s_\xi$ 的选取，则 $f = s/s_\xi$ 被 $uf$ 代替，其中
$u \in \mathcal{O}_{X, \xi}^h$ 是单位。

\begin{definition}
\label{spaces-over-fields-definition-order-vanishing-meromorphic}
设 $S$ 为概形。设
% P09-ERR-0254
$X$ 为 $S$ 上局部诺特的整代数空间，$\mathcal{L}$ 为可逆 $\mathcal{O}_X$-模。
设 $s \in \Gamma(X, \mathcal{K}_X(\mathcal{L}))$ 为 $\mathcal{L}$ 的正则亚纯截面。
对每个一般点为 $\xi \in |Z|$ 的素除子 $Z \subset X$，把
{\it $s$ 沿 $Z$ 的消没阶}定义为整数
$$
\text{ord}_{Z, \mathcal{L}}(s) =
\text{length}_{\mathcal{O}_{X, \xi}^h}
(\mathcal{O}_{X, \xi}^h/a \mathcal{O}_{X, \xi}^h) -
\text{length}_{\mathcal{O}_{X, \xi}^h}
(\mathcal{O}_{X, \xi}^h/b \mathcal{O}_{X, \xi}^h)
$$
其中 $a, b \in \mathcal{O}_{X, \xi}^h$ 为非零因子，并且上文构造的
$Q(\mathcal{O}_{X, \xi}^h)$ 中的元素 $s/s_\xi$ 等于 $a/b$。由上文及《代数》引理
\ref{algebra-lemma-ord-additive}，这是良定义的。
\end{definition}

\noindent
如上所述，$\mathcal{O}_X$ 的正则亚纯截面 $s$ 可写成 $s = f \cdot 1$，其中 $f$
是 $X$ 上的非零有理函数，并且有
$\text{ord}_Z(f) = \text{ord}_{Z, \mathcal{O}_X}(s)$。与主除子情形一样，有以下引理。

\begin{lemma}
\label{spaces-over-fields-lemma-divisor-meromorphic-locally-finite}
设 $S$ 为概形，$X$ 为 $S$ 上局部诺特的整代数空间，$\mathcal{L}$ 为可逆
$\mathcal{O}_X$-模。设
% P09-ERR-0255
$s \in \mathcal{K}_X(\mathcal{L})$ 为 $\mathcal{L}$ 的正则（即非零）亚纯截面。
则集合
$$
\{Z \subset X \mid Z \text{ 为素除子，其一般点为 }\xi
\text{ 且 }s\text{ 不属于 }\mathcal{L}_\xi\}
$$
以及
$$
\{Z \subset X \mid Z \text{ 为素除子且 }
\text{ord}_{Z, \mathcal{L}}(s) \not = 0\}
$$
在 $X$ 中局部有限。
\end{lemma}

\begin{proof}
存在非空开子空间 $U \subset X$，使 $s$ 对应于 $\Gamma(U, \mathcal{L})$
的一个截面，且该截面在 $U$ 上生成 $\mathcal{L}$。因此，引理中的集合里可能出现的
素除子全都对应于 $|X| \setminus |U|$ 的不可约分支。故引理
% P09-ERR-0256
\ref{spaces-over-fields-lemma-components-locally-finite} 给出所需结论。
\end{proof}

\begin{lemma}
\label{spaces-over-fields-lemma-divisor-meromorphic-well-defined}
设 $S$ 为概形。设
% P09-ERR-0257
$X$ 为 $S$ 上局部诺特的整代数空间，$\mathcal{L}$ 为可逆 $\mathcal{O}_X$-模。
设 $s, s' \in \mathcal{K}_X(\mathcal{L})$ 为 $\mathcal{L}$ 的非零亚纯截面。
则 $f = s/s'$ 是 $R(X)^*$ 的元素，并且有
$$
\sum \text{ord}_{Z, \mathcal{L}}(s)[Z]
=
\sum \text{ord}_{Z, \mathcal{L}}(s')[Z]
+
\text{div}(f)
$$
作为 Weil 除子成立。
\end{lemma}

\begin{proof}
这由定义显然。注意，引理 \ref{spaces-over-fields-lemma-divisor-meromorphic-locally-finite}
保证这些和确实是 Weil 除子。
\end{proof}

\begin{definition}
\label{spaces-over-fields-definition-divisor-invertible-sheaf}
设 $S$ 为概形，$X$ 为 $S$ 上局部诺特的整代数空间，$\mathcal{L}$ 为可逆
$\mathcal{O}_X$-模。
\begin{enumerate}
\item 对 $\mathcal{L}$ 的任意非零亚纯截面 $s$，把{\it 与 $s$ 相伴的 Weil 除子}
定义为
$$
\text{div}_\mathcal{L}(s) =
\sum \text{ord}_{Z, \mathcal{L}}(s) [Z] \in \text{Div}(X)
$$
其中对素除子求和。由引理 \ref{spaces-over-fields-lemma-divisor-meromorphic-locally-finite}，
这是良定义的。
\item 把{\it 与 $\mathcal{L}$ 相伴的 Weil 除子类}定义为
$\text{div}_\mathcal{L}(s)$ 在 $\text{Cl}(X)$ 中的像，其中 $s$ 是
$\mathcal{L}$ 在 $X$ 上的任意非零亚纯截面。由引理
\ref{spaces-over-fields-lemma-divisor-meromorphic-well-defined}，这是良定义的。
\end{enumerate}
\end{definition}

\noindent
正如所预期的，该构造关于可逆模可加。

\begin{lemma}
\label{spaces-over-fields-lemma-c1-additive}
设 $S$ 为概形，$X$ 为 $S$ 上局部诺特的整代数空间。设 $\mathcal{L}$、
$\mathcal{N}$ 为可逆 $\mathcal{O}_X$-模。设 $s$、$t$ 分别为 $\mathcal{L}$、
$\mathcal{N}$ 的非零亚纯截面。则 $st$ 是
$\mathcal{L} \otimes_{\mathcal{O}_X} \mathcal{N}$ 的非零亚纯截面，并且
$$
\text{div}_{\mathcal{L} \otimes \mathcal{N}}(st)
=
\text{div}_\mathcal{L}(s) + \text{div}_\mathcal{N}(t)
$$
在 $\text{Div}(X)$ 中成立。特别地，
$\mathcal{L} \otimes_{\mathcal{O}_X} \mathcal{N}$ 的 Weil 除子类是
$\mathcal{L}$ 与 $\mathcal{N}$ 的 Weil 除子类之和。
\end{lemma}

\begin{proof}
设 $s$、$t$ 分别为 $\mathcal{L}$、$\mathcal{N}$ 的非零亚纯截面。则 $st$ 是
$\mathcal{L} \otimes \mathcal{N}$ 的非零亚纯截面。设 $Z \subset X$ 为素除子，
$\xi \in |Z|$ 为其一般点。按本节前述记号，选取生成元
$s_\xi \in \mathcal{L}_\xi$ 和 $t_\xi \in \mathcal{N}_\xi$。则
$s_\xi \otimes t_\xi$ 是 $(\mathcal{L} \otimes \mathcal{N})_\xi$ 的生成元。
在 $Q(\mathcal{O}_{X, \xi}^h)$ 中，
$st/(s_\xi t_\xi) = (s/s_\xi)(t/t_\xi)$。应用《代数》引理
\ref{algebra-lemma-ord-additive} 的可加性，可得
$$
\text{div}_{\mathcal{L} \otimes \mathcal{N}, Z}(st)
=
\text{div}_{\mathcal{L}, Z}(s) + \text{div}_{\mathcal{N}, Z}(t)
$$
% P09-ERR-0258
略去一些细节。
\end{proof}

\noindent
设 $S$ 为概形，$X$ 为 $S$ 上局部诺特的整代数空间。由上述构造和引理，得到交换群同态
\begin{equation}
\label{spaces-over-fields-equation-c1}
\Pic(X) \longrightarrow \text{Cl}(X)
\end{equation}
它把可逆模映到其 Weil 除子类。

\begin{lemma}
\label{spaces-over-fields-lemma-normal-c1-injective}
设 $S$ 为概形，$X$ 为 $S$ 上局部诺特的整代数空间。若 $X$ 正规，则映射
(\ref{spaces-over-fields-equation-c1}) $\Pic(X) \to \text{Cl}(X)$ 为单射。
\end{lemma}

\begin{proof}
设 $\mathcal{L}$ 为可逆 $\mathcal{O}_X$-模，其相伴 Weil 除子类平凡。设 $s$ 为
$\mathcal{L}$ 的正则亚纯截面。该假设意味着，对某个 $f \in R(X)^*$，有
$\text{div}_\mathcal{L}(s) = \text{div}(f)$。于是 $t = f^{-1}s$ 是
$\mathcal{L}$ 的正则亚纯截面，且 $\text{div}_\mathcal{L}(t) = 0$，见引理
\ref{spaces-over-fields-lemma-divisor-meromorphic-well-defined}。断言 $t$ 定义 $\mathcal{L}$ 的一个平凡化。
该断言即完成引理的证明。由于目前尚无太多可用理论，对该断言的证明略显笨拙；建议读者
% P09-ERR-0259
略过此证明。

\medskip\noindent
可以平展局部地验证该断言。设 $U \in X_\etale$ 仿射，且
$\mathcal{L}|_U$ 平凡。设 $s_U \in \Gamma(U, \mathcal{L}|_U)$ 为一个平凡化。
由《性质》引理 \ref{properties-lemma-normal-locally-finite-nr-irreducibles}，
还可假设 $U$ 整。把 $U = \Spec(A)$ 写成正规诺特整环 $A$ 的谱，其分式域为 $K$。
可以对某个 $K$ 中的元素 $f$ 写成 $t|_U = f s_U$；例如见《空间上的除子》引理
\ref{spaces-divisors-lemma-meromorphic-quasi-coherent}。设高度为一的素理想
$\mathfrak p \subset A$ 对应于余维为 $1$ 的点 $u \in U$，后者映到余维为 $1$ 的点
$\xi \in |X|$。按本节开头所述选取 $\mathcal{L}_\xi$ 的平凡化 $s_\xi$。
选取 $U$ 位于 $u$ 上方的几何点 $\overline{u}$。则
$$
(\mathcal{O}_{X, \xi}^h)^{sh} =
\mathcal{O}_{X, \overline{u}} =
\mathcal{O}_{U, u}^{sh} = (A_\mathfrak p)^{sh}
$$
见《良态空间》诸引理
\ref{decent-spaces-lemma-henselian-local-ring-strict}
及《空间的性质》引理
\ref{spaces-properties-lemma-describe-etale-local-ring}.
$X$ 的正规性说明这些全都是离散赋值环。把 $\mathcal{L}$ 拉回到
$\Spec(\mathcal{O}_{X, \overline{u}})$ 后，平凡化 $s_U$ 与 $s_\xi$ 相差一个单位。
写 $t = f_\xi s_\xi$，其中 $f_\xi \in Q(\mathcal{O}_{X, \xi}^h)$。由此，
$f_\xi$ 与 $f$ 在 $Q(\mathcal{O}_{X, \overline{u}})$ 中相差一个单位。若
$Z \subset X$ 表示对应于 $\xi$ 的素除子（引理
\ref{spaces-over-fields-lemma-decent-irreducible-closed}），则
$0 = \text{ord}_{Z, \mathcal{L}}(t) =
\text{ord}_{\mathcal{O}_{X, \xi}^h}(f_\xi)$
而由于 $\mathcal{O}_{X, \xi}^h$ 是离散赋值环，$f_\xi$ 是单位。因此，$f$ 是
$\mathcal{O}_{X, \overline{u}}$ 中的单位，特别有 $f \in A_\mathfrak p$。
由《代数》引理
\ref{algebra-lemma-normal-domain-intersection-localizations-height-1}，这蕴含
$f \in A$。因此 $t \in \Gamma(X, \mathcal{L})$。把 $t^{-1}$ 视为
$\mathcal{L}^{\otimes -1}$ 的亚纯截面并重复这一论证，即完成证明。
\end{proof}

















\section{修改与 alteration}
\label{spaces-over-fields-section-modifications-alterations}

\noindent
利用整代数空间的概念，可以如下定义修改。

\begin{definition}
\label{spaces-over-fields-definition-modification}
设 $S$ 为概形，$X$ 为 $S$ 上的整代数空间。所谓{\it $X$ 的修改}，是指
$S$ 上代数空间的一个双有理固有态射 $f : X' \to X$，其中 $X'$ 为整空间。
\end{definition}

\noindent
关于代数空间的双有理态射，见《良态空间》定义
\ref{decent-spaces-definition-birational}。

\begin{lemma}
\label{spaces-over-fields-lemma-modification-iso-over-open}
设 $f : X' \to X$ 为定义 \ref{spaces-over-fields-definition-modification} 所述的修改。
则存在非空开子空间 $U \subset X$，使 $f^{-1}(U) \to U$ 为同构。
\end{lemma}

\begin{proof}
由引理 \ref{spaces-over-fields-lemma-finite-degree}，存在非空开子空间 $U \subset X$，使
$f^{-1}(U) \to U$ 有限。由泛平坦性（《空间的态射》命题
\ref{spaces-morphisms-proposition-generic-flatness-reduced})
可假设 $f^{-1}(U) \to U$ 平坦且为有限表示。因此 $f^{-1}(U) \to U$
有限局部自由（《空间的态射》引理 \ref{spaces-morphisms-lemma-finite-flat}）。
由于 $f$ 双有理，$X'$ 在 $X$ 上的次数为 $1$。故 $f^{-1}(U) \to U$
是次数为 $1$ 的有限局部自由态射，换言之，它是同构。
\end{proof}

\begin{definition}
\label{spaces-over-fields-definition-alteration}
设 $S$ 为概形，$X$ 为 $S$ 上的整代数空间。所谓{\it $X$ 的 alteration}，
是指 $S$ 上代数空间的一个固有支配态射 $f : Y \to X$，其中 $Y$ 为整空间，
并且存在某个非空开子空间 $U \subset X$，使 $f^{-1}(U) \to U$ 有限。
\end{definition}

\noindent
若 $f : Y \to X$ 是整代数空间之间的支配固有态射，并且一般点处诱导的
剩余域扩张有限，则它就是 alteration。准确陈述如下。
% P09-ERR-0260

\begin{lemma}
\label{spaces-over-fields-lemma-alteration-generically-finite}
设 $S$ 为概形，$f : X \to Y$ 为 $S$ 上整代数空间之间的固有支配态射。
则 $f$ 为 alteration，当且仅当引理 \ref{spaces-over-fields-lemma-finite-degree} 中相互等价的
条件 (1)--(6) 中任一条件成立。
\end{lemma}

\begin{proof}
这是陈述中所引引理的直接推论。
\end{proof}

\begin{lemma}
\label{spaces-over-fields-lemma-alteration-contained-in}
设 $S$ 为概形，$f : X \to Y$ 为 $S$ 上代数空间的固有满态射。假设 $Y$ 为整空间。
则存在整闭子空间 $X' \subset X$，使得 $f' = f|_{X'} : X' \to Y$ 为 alteration。
\end{lemma}

\begin{proof}
设 $V \subset Y$ 为非空仿射开子空间（《良态空间》定理
\ref{decent-spaces-theorem-decent-open-dense-scheme}）。设 $\eta \in V$ 为一般点。
则 $X_\eta$ 是 $\eta$ 上非空的固有代数空间。选取闭点 $x \in |X_\eta|$
（其存在性是因为 $|X_\eta|$ 是拟紧的良态拓扑空间，见《良态空间》命题
\ref{decent-spaces-proposition-reasonable-sober} 及《拓扑》引理
\ref{topology-lemma-quasi-compact-closed-point}）。
% P09-ERR-0261
令 $X'$ 为 $\overline{\{x\}} \subset |X|$ 上的约化诱导闭子空间结构
（《空间的性质》定义
\ref{spaces-properties-definition-reduced-induced-space}）。
% P09-ERR-0262
由于像包含 $\eta$，态射 $f' : X' \to Y$ 为满射。又因 $f'$ 是闭浸入与固有态射的复合，
故它是固有的。最后，纤维 $X'_\eta$ 只有一个点；对此可分别对 $X \to Y$ 和
$X' \to Y$ 以及点 $\eta$ 应用《良态空间》引理
\ref{decent-spaces-lemma-topology-fibre}。由于 $Y$ 良态且 $X' \to Y$ 分离，
可知 $X'$ 良态（《良态空间》诸引理
% P09-ERR-0263
\ref{decent-spaces-lemma-properties-trivial-implications} 及
\ref{decent-spaces-lemma-property-over-property}）。于是由引理
\ref{spaces-over-fields-lemma-alteration-generically-finite}，$f'$ 是 alteration。
\end{proof}








\section{概形轨迹}
\label{spaces-over-fields-section-schematic}

\noindent
关于代数空间的概形轨迹，我们已经证明了若干结果。以下列出相关参考：
\begin{enumerate}
\item 《空间的性质》第
\ref{spaces-properties-section-schematic} 节及第
\ref{spaces-properties-section-getting-a-scheme},
\item 《良态空间》第 \ref{decent-spaces-section-schematic} 节，
\item 《空间的性质》引理
\ref{spaces-properties-lemma-point-like-spaces}
$\Leftarrow$
《良态空间》引理 \ref{decent-spaces-lemma-decent-point-like-spaces}
$\Leftarrow$
《良态空间》引理 \ref{decent-spaces-lemma-when-field}，
\item 《空间的极限》第 \ref{spaces-limits-section-affine} 节，以及
\item 《空间的极限》第 \ref{spaces-limits-section-representable} 节。
\end{enumerate}
在某些情形下，特定类型的代数空间态射自动可表，例如分离的局部拟有限态射
（《空间的态射》引理
\ref{spaces-morphisms-lemma-locally-quasi-finite-separated-representable}）以及
平坦单态射（《空间态射进阶》引理
\ref{spaces-more-morphisms-lemma-flat-case}）。第
\ref{spaces-over-fields-section-schematic-and-field-extension} 节将研究基域扩张下概形轨迹的变化。

\begin{lemma}
\label{spaces-over-fields-lemma-locally-finite-type-dim-zero}
设 $S$ 为概形，$X$ 为 $S$ 上的代数空间。假设 $X$ 至少满足下列条件之一：
\begin{enumerate}
\item $X$ 拟分离且 $\dim(X) = 0$；
\item $X$ 在某个域 $k$ 上局部有限型，且 $\dim(X) = 0$；
\item $X$ 诺特且 $\dim(X) = 0$；或者
\item 在此补充更多内容。
% P09-ERR-0264
\end{enumerate}
则 $X$ 是分离概形，并且 $X$ 的任意拟紧开子空间均为仿射的。
\end{lemma}

\begin{proof}
若证明了 $X$ 的任意拟紧开子空间均为仿射的，则 $X$ 是分离概形。因此可假设
$X$ 拟紧，并以证明 $X$ 仿射为目标。情形 (2) 与 (3) 立即由情形 (1) 推出；不过我们将
分别证明 (2) 与 (3)，因为这些证明所用的理论少得多。
% P09-ERR-0265

\medskip\noindent
情形 (3) 的证明。设 $U$ 为仿射概形，$U \to X$ 为平展态射。令
$R = U \times_X U$。两个投影态射 $s, t : R \to U$ 都是概形的平展态射。
由《空间的性质》定义 \ref{spaces-properties-definition-dimension}，可知
$\dim(U) = 0$ 且 $\dim(R) = 0$。由于 $R$ 是维数为 $0$ 的局部诺特概形，
$R$ 是一些阿廷局部环之谱的不交并（《性质》引理
\ref{properties-lemma-locally-Noetherian-dimension-0}）。由于已假设 $X$ 诺特
（从而拟分离），可知 $R$ 拟紧。因此 $R$ 是仿射概形（使用《概形》引理
\ref{schemes-lemma-disjoint-union-affines}）。平展态射 $s, t : R \to U$ 诱导有限的
剩余域扩张。因此，由《代数》引理
\ref{algebra-lemma-essentially-of-finite-type-into-artinian-local}
（略去一处小细节），$s$ 与 $t$ 均为有限态射。于是《群胚》命题
\ref{groupoids-proposition-finite-flat-equivalence} 说明 $X = U/R$ 是仿射概形。

\medskip\noindent
情形 (2) 的证明——几乎与情形 (3) 的证明相同。设 $U$ 为仿射概形，
$U \to X$ 为平展满态射。令 $R = U \times_X U$。两个投影态射
$s, t : R \to U$ 都是概形的平展态射。由《空间的性质》定义
\ref{spaces-properties-definition-dimension}，可知 $\dim(U) = 0$，并且同样有
$\dim(R) = 0$。另一方面，态射 $U \to \Spec(k)$ 是平展态射 $U \to X$ 与
$X \to \Spec(k)$ 的复合，故为局部有限型，见《空间的态射》诸引理
\ref{spaces-morphisms-lemma-composition-finite-type} 及
\ref{spaces-morphisms-lemma-etale-locally-finite-type}.
同理，$R \to \Spec(k)$ 为局部有限型。因此由《簇》引理
\ref{varieties-lemma-algebraic-scheme-dim-0}，$U$ 与 $R$ 都是一些在 $k$ 上有限的
阿廷局部 $k$-代数之谱的不交并。故 $U \times_{\Spec(k)} U$ 也具有同样的性质。由于
$$
R = U \times_X U \longrightarrow U \times_{\Spec(k)} U
$$
是单态射，可知 $R$ 是有限（！）多个有限 $k$-代数之谱的不交并。因此 $R$ 仿射，
见《概形》引理 \ref{schemes-lemma-disjoint-union-affines}。再次应用《簇》引理
\ref{varieties-lemma-algebraic-scheme-dim-0}，可知 $R$ 在 $k$ 上有限。故 $s, t$
均为有限态射，见《态射》引理 \ref{morphisms-lemma-finite-permanence}。于是《群胚》命题
\ref{groupoids-proposition-finite-flat-equivalence} 说明 $X = U/R$ 是仿射概形。

\medskip\noindent
情形 (1) 的上同调证明。由《空间的上同调》引理
\ref{spaces-cohomology-lemma-vanishing-above-dimension}，$X$ 上任意拟凝聚层
$\mathcal{F}$ 的高阶上同调群均消没。因此由《空间的上同调》命题
\ref{spaces-cohomology-proposition-vanishing-affine}，$X$ 为仿射的（特别地，是概形）。

\medskip\noindent
情形 (1) 的几何证明。选取分层
$$
\emptyset = U_{n + 1} \subset
U_n \subset U_{n - 1} \subset \ldots \subset U_1 = X
$$
以及《良态空间》引理
\ref{decent-spaces-lemma-filter-quasi-compact-quasi-separated} 所述的平展态射
$f_p : V_p \to U_p$（下文将使用它们的全部性质）。由代数空间维数的定义，
$\dim(V_p) = 0$。因此《性质》引理 \ref{properties-lemma-dimension-zero}
可用于每个 $V_p$。于是 $f_p^{-1}(U_{p + 1}) \subset V_p$ 是拟紧开子集，故既仿射又闭。
由此 $|T_p| \subset |U_p|$（见前引处）既开又闭。因此 $X$ 是一些开闭子空间的不交并，
而这些子空间的约化结构均为概形。于是 $X$ 是概形（《空间的极限》引理
\ref{spaces-limits-lemma-reduction-scheme}）。最后应用上文已经引用过的概形情形，即完成证明。
\end{proof}

\noindent
下列引理说明，拟分离代数空间在余维 $1$ 以外是概形。

\begin{lemma}
\label{spaces-over-fields-lemma-generic-point-in-schematic-locus}
设 $S$ 为概形，$X$ 为 $S$ 上的拟分离代数空间。设 $x \in |X|$。下列条件等价：
\begin{enumerate}
\item $x$ 是 $X$ 上余维为 $0$ 的点；
\item $X$ 在 $x$ 处的局部环维数为 $0$；以及
\item $x$ 是 $|X|$ 的某个不可约分支的一般点。
\end{enumerate}
若这些条件成立，则存在包含 $x$ 的 $X$ 的开子空间，并且该开子空间为概形。
\end{lemma}

\begin{proof}
(1)、(2) 与 (3) 的等价性由《良态空间》引理
\ref{decent-spaces-lemma-decent-generic-points} 以及拟分离代数空间良态这一事实
（《良态空间》第 \ref{decent-spaces-section-reasonable-decent} 节）推出。
不过，下一段将给出这一等价性的更初等证明。

\medskip\noindent
注意，(1) 与 (2) 依定义等价（《空间的性质》定义
\ref{spaces-properties-definition-dimension-local-ring}）。为证明 (1) 与 (3) 等价，
可假设 $X$ 拟紧。选取
$$
\emptyset = U_{n + 1} \subset
U_n \subset U_{n - 1} \subset \ldots \subset U_1 = X
$$
以及《良态空间》引理
\ref{decent-spaces-lemma-filter-quasi-compact-quasi-separated} 所述的
$f_i : V_i \to U_i$。设 $x \in U_i$ 且 $x \not \in U_{i + 1}$。则存在唯一的
$y \in V_i$，使 $x = f_i(y)$。若 (1) 成立，则 $y$ 是 $V_i$ 的某个不可约分支的一般点
（《空间的性质》引理 \ref{spaces-properties-lemma-codimension-0-points}）。由于
$f_i^{-1}(U_{i + 1})$ 是 $V_i$ 中不包含 $y$ 的拟紧开子集，存在 $y$ 的开邻域
$W \subset V_i$，它与 $f_i^{-1}(V_i)$ 不交（见《性质》引理
% P09-ERR-0266
\ref{properties-lemma-generic-point-in-constructible}，或者更简单地见《代数》引理
\ref{algebra-lemma-standard-open-containing-maximal-point}）。于是
$f_i|_W : W \to X$ 是到其像的同构，故 $x = f_i(y)$ 是 $|X|$ 的一般点。
反之，假设 (3) 成立。则 $f_i$ 把 $\overline{\{y\}}$ 映到 $|U_i|$ 的不可约分支
$\overline{\{x\}}$ 上。由于 $|f_i|$ 在 $\overline{\{x\}}$ 上方为双射，可知
$\overline{\{y\}}$ 是 $U_i$ 的一个不可约分支。因此 $x$ 是余维为 $0$ 的点。
% P09-ERR-0267

\medskip\noindent
引理的最后一个陈述即《空间的性质》命题
\ref{spaces-properties-proposition-locally-quasi-separated-open-dense-scheme}.
\end{proof}

\noindent
下列引理说明，分离的局部诺特代数空间在余维 $1$ 处为概形，亦即在余维 $2$ 以外为概形。

\begin{lemma}
\label{spaces-over-fields-lemma-codim-1-point-in-schematic-locus}
\begin{slogan}
分离代数空间在余维 1 处为概形。
\end{slogan}
设 $S$ 为概形，$X$ 为 $S$ 上的代数空间。设 $x \in |X|$。若 $X$ 分离、局部诺特，
并且 $X$ 在 $x$ 处的局部环维数为 $\leq 1$（《空间的性质》定义
\ref{spaces-properties-definition-dimension-local-ring}),
则存在包含 $x$ 的 $X$ 的开子空间，并且该开子空间为概形。
\end{lemma}

\begin{proof}
（另有一种避开有限群胚材料的方法，见下述评注。）可以把 $X$ 换成 $x$ 的一个拟紧邻域，
因而可假设 $X$ 拟紧、分离且诺特。
% P09-ERR-0268
存在概形 $U$ 及有限满态射 $U \to X$，见《空间的极限》命题
\ref{spaces-limits-proposition-there-is-a-scheme-finite-over}.
令 $R = U \times_X U$。则 $j : R \to U \times_S U$ 是等价关系，并由此得到
$S$ 上的群胚概形 $(U, R, s, t, c)$，其中 $s, t$ 有限，且 $U$ 诺特而分离。
设 $\{u_1, \ldots, u_n\} \subset U$ 为映到 $x$ 的点集。由《良态空间》引理
\ref{decent-spaces-lemma-dimension-local-ring-quasi-finite}.
可得 $\dim(\mathcal{O}_{U, u_i}) \leq 1$。

\medskip\noindent
由《群胚进阶》引理
\ref{more-groupoids-lemma-find-affine-codimension-1}，存在包含轨道
$\{u_1, \ldots, u_n\}$ 的 $R$-不变仿射开子集 $W \subset U$。由于 $U \to X$
有限且满，连续映射 $|U| \to |X|$ 为闭满射，故由《拓扑》引理
\ref{topology-lemma-closed-morphism-quotient-topology}，它是商映射。因此 $f(W)$ 为开集，
并存在开子空间 $X' \subset X$，使 $f : W \to X'$ 为有限满态射。
% P09-ERR-0269
由《空间的上同调》引理
\ref{spaces-cohomology-lemma-image-affine-finite-morphism-affine-Noetherian}
可知 $X'$ 是仿射概形，证明完成。
\end{proof}

\begin{remark}
\label{spaces-over-fields-remark-alternate-proof-scheme-codim-1}
下面给出引理 \ref{spaces-over-fields-lemma-codim-1-point-in-schematic-locus} 的一个证明梗概，
其中不使用《群胚进阶》引理
\ref{more-groupoids-lemma-find-affine-codimension-1}.

\medskip\noindent
步骤 1。可假设 $X$ 是约化的诺特分离代数空间（例如使用《空间的上同调》引理
\ref{spaces-cohomology-lemma-image-affine-finite-morphism-affine-Noetherian}
或者《空间的极限》引理 \ref{spaces-limits-lemma-reduction-scheme}），
并可选取有限满态射 $Y \to X$，其中 $Y$ 为诺特概形（由《空间的极限》命题
\ref{spaces-limits-proposition-there-is-a-scheme-finite-over}).

\medskip\noindent
步骤 2。用 $x$ 的一个开邻域替换 $X$ 后，存在双有理有限态射 $X' \to X$ 以及闭子概形
$Y' \subset X' \times_X Y$，使 $Y' \to X'$ 为有限局部自由满态射。事实上，由于 $X$
约化，存在稠密开子空间 $U \subset X$，使 $Y$ 在其上方平坦（《空间的态射》命题
\ref{spaces-morphisms-proposition-generic-flatness-reduced}）。于是可以选取
$U$-容许爆破 $b : \tilde X \to X$，使 $Y$ 的严格变换 $\tilde Y$ 在 $\tilde X$ 上平坦，
见《空间态射进阶》引理
\ref{spaces-more-morphisms-lemma-flat-after-blowing-up}。
（若想避开关于爆破的结果，也可改用 Hilbert 概形。）然后令 $X' \subset \tilde X$ 为
$b^{-1}(U)$ 的概形论闭包，并令 $Y' = X' \times_{\tilde X} \tilde Y$。由于 $x$ 是余维为
$1$ 的点，可知 $X' \to X$ 在 $x$ 的某个邻域上有限（引理
\ref{spaces-over-fields-lemma-finite-in-codim-1}）。

\medskip\noindent
步骤 3。把 $X$ 缩小到 $x$ 的一个更小邻域后，$X'$ 是概形。这是因为 $Y'$ 是概形、
$Y' \to X'$ 有限局部自由，并且 $Y'$ 中每个由余维为 $1$ 的点组成的有限集都包含在某个
仿射开子集中。使用《空间的性质》命题
% P09-ERR-0270
\ref{spaces-properties-proposition-finite-flat-equivalence-global} 及《簇》命题
\ref{varieties-proposition-finite-set-of-points-of-codim-1-in-affine}.

\medskip\noindent
步骤 4。存在仿射开子集 $W' \subset X'$，它包含 $x$ 上方的所有点，并且是 $X$ 的某个
开子空间的原像。为证明这一点，设 $Z \subset X$ 为 $X' \to X$ 不是同构的点集的闭包。
可以假设 $x \in Z$，否则已经完成。此时 $x$ 是 $Z$ 的某个不可约分支的一般点；缩小
$X$ 后，可假设 $Z$ 是仿射概形（引理
\ref{spaces-over-fields-lemma-generic-point-in-schematic-locus}）。于是原像 $Z' \subset X'$ 也是仿射概形。
设 $x_1, \ldots, x_n \in Z'$ 为映到 $x$ 的各点。可以在 $X'$ 中找到仿射开子集 $W'$，
使它与 $Z'$ 的交集是 $Z$ 中某个包含 $x$ 的主开子集的原像。事实上，首先利用《簇》命题
\ref{varieties-proposition-finite-set-of-points-of-codim-1-in-affine}，选取包含
$x_1, \ldots, x_n$ 的仿射开子集 $W' \subset X'$。然后选取包含 $x$ 的主开子集
$D(f) \subset Z$，使其原像 $D(f|_{Z'})$ 包含在 $W' \cap Z'$ 中。再选取
$f' \in \Gamma(W', \mathcal{O}_{W'})$，其限制为 $f|_{Z'}$，并把 $W'$
替换为 $D(f') \subset W'$。由于 $X' \to X$ 在 $Z' \to Z$ 之外为同构，对 $W'$
的这一选取保证了 $W'$ 的像 $W \subset X$ 为开集，并且它在 $X'$ 中的原像为 $W'$。

\medskip\noindent
步骤 5。此时 $W' \to W$ 是有限满态射，而由《空间的上同调》引理
\ref{spaces-cohomology-lemma-image-affine-finite-morphism-affine-Noetherian}
可知 $W$ 是概形，证明完成。
\end{remark}





\section{概形轨迹与域扩张}
\label{spaces-over-fields-section-schematic-and-field-extension}

\noindent
域 $k$ 上不可表的代数空间，在基变换到 $k$ 的某个扩张后可能变得可表（即成为概形）。
见《空间》例 \ref{spaces-example-non-representable-descent}。本节研究这一问题。

\begin{lemma}
\label{spaces-over-fields-lemma-scheme-after-purely-inseparable-base-change}
设 $k$ 为域，$X$ 为 $k$ 上的代数空间。若存在纯不可分域扩张 $k'/k$，
使 $X_{k'}$ 为概形，则 $X$ 为概形。
\end{lemma}

\begin{proof}
态射 $X_{k'} \to X$ 是整的、满的且泛单。故本引理由《空间的极限》引理
\ref{spaces-limits-lemma-integral-universally-bijective-scheme}.
推出。
\end{proof}

\begin{lemma}
\label{spaces-over-fields-lemma-when-scheme-after-base-change}
设 $k$ 为域，其代数闭包为 $\overline{k}$。设 $X$ 为 $k$ 上的拟分离代数空间。
\begin{enumerate}
\item 若存在域扩张 $K/k$，使 $X_K$ 为概形，则 $X_{\overline{k}}$ 为概形。
\item 若 $X$ 拟紧，并存在域扩张 $K/k$，使 $X_K$ 为概形，则对于 $k$ 的某个
有限可分扩张 $k'$，$X_{k'}$ 为概形。
\end{enumerate}
\end{lemma}

\begin{proof}
由于每个代数空间都是其拟紧开子空间的并，引理的第一部分由第二部分推出
（略去一些细节）。因此假设 $X$ 拟紧，并给定扩张 $K/k$，使 $X_K$ 可表。
把 $K = \bigcup A$ 写成有限生成 $k$-子代数 $A$ 的余极限。由《空间的极限》引理
\ref{spaces-limits-lemma-limit-is-scheme}，对某个 $A$，$X_A$ 为概形。选取极大理想
$\mathfrak m \subset A$。由 Hilbert 零点定理（《代数》定理
\ref{algebra-theorem-nullstellensatz}），剩余域 $k' = A/\mathfrak m$ 是 $k$ 的有限扩张。
故 $X_{k'}$ 为概形。若 $k' \supset k$ 不可分，设 $k'/k''/k$ 为《域》引理
\ref{fields-lemma-separable-first} 中得到的子扩张。由于 $k'/k''$ 纯不可分，
由引理 \ref{spaces-over-fields-lemma-scheme-after-purely-inseparable-base-change}，代数空间
$X_{k''}$ 为概形。由于 $k''|k$ 可分，证明完成。
% P09-ERR-0271
\end{proof}

\begin{lemma}
\label{spaces-over-fields-lemma-base-change-by-Galois}
设 $k'/k$ 为有限 Galois 扩张，其 Galois 群为 $G$。设 $X$ 为 $k$ 上的代数空间。
则 $G$ 自由作用在代数空间 $X_{k'}$ 上，并且在《空间的性质》引理
\ref{spaces-properties-lemma-quotient} 的意义下有 $X = X_{k'}/G$。
\end{lemma}

\begin{proof}
略。提示：先证明 $\Spec(k) = \Spec(k')/G$，再使用取商与基变换的相容性。
\end{proof}

\begin{lemma}
\label{spaces-over-fields-lemma-when-quotient-scheme-at-point}
设 $S$ 为概形，$X$ 为 $S$ 上的代数空间，$G$ 为自由作用在 $X$ 上的有限群。
按《空间的性质》引理 \ref{spaces-properties-lemma-quotient} 令 $Y = X/G$。
对 $y \in |Y|$，下列条件等价：
\begin{enumerate}
\item $y$ 位于 $Y$ 的概形轨迹中；以及
\item 存在仿射开子集 $U \subset X$，它包含 $y$ 的原像。
\end{enumerate}
\end{lemma}

\begin{proof}
由《空间的性质》引理 \ref{spaces-properties-lemma-quotient} 中 $Y = X/G$ 的构造，
态射 $X \to Y$ 满且平展。显然有 $X \times_Y X = X \times G$，故态射
$X \to Y$ 甚至有限平展。它也是满的。因此，本引理由《良态空间》引理
\ref{decent-spaces-lemma-when-quotient-scheme-at-point} 推出。
\end{proof}

\begin{lemma}
\label{spaces-over-fields-lemma-scheme-after-purely-transcendental-base-change}
设 $k$ 为域，$X$ 为 $k$ 上的拟分离代数空间。若存在纯超越域扩张 $K/k$，
使 $X_K$ 为概形，则 $X$ 为概形。
\end{lemma}

\begin{proof}
由于每个代数空间都是其拟紧开子空间的并，可假设 $X$ 拟紧（略去一些细节）。
回忆《域》定义 \ref{fields-definition-transcendence}，关于扩张 $K/k$ 的假设表示
$K$ 是 $k$ 上某个多项式环（变量可能有无限多个）的分式域。因此 $K = \bigcup A$
是一些子代数的并，其中每个子代数都是 $k$ 上有限多项式代数的局部化。由《空间的极限》引理
\ref{spaces-limits-lemma-limit-is-scheme}，对某个 $A$，$X_A$ 是概形。写成
$$
A = k[x_1, \ldots, x_n][1/f]
$$
其中 $f \in k[x_1, \ldots, x_n]$ 非零。

\medskip\noindent
若 $k$ 无限，则可如下完成证明：选取 $a_1, \ldots, a_n \in k$，使
$f(a_1, \ldots, a_n) \not = 0$。于是 $(a_1, \ldots, a_n)$ 定义一个
$k$-代数映射 $A \to k$，它把 $x_i$ 映到 $a_i$，把 $1/f$ 映到
$1/f(a_1, \ldots, a_n)$。故基变换 $X_A \times_{\Spec(A)} \Spec(k) \cong X$
如所需地是概形。
% P09-ERR-0272

\medskip\noindent
本段在 $k$ 有限时完成证明。此时把 $X = \lim X_i$ 写成极限，其中 $X_i$ 在 $k$ 上为有限表示，
转移态射均为仿射态射（《空间的极限》引理
\ref{spaces-limits-lemma-relative-approximation}）。使用《空间的极限》引理
\ref{spaces-limits-lemma-limit-is-scheme}，可知对某个 $i$，$X_{i, A}$ 为概形。因此可假设
$X \to \Spec(k)$ 为有限表示。设 $x \in |X|$ 为闭点。可以用闭浸入
$\Spec(\kappa) \to X$ 表示 $x$（《良态空间》引理
\ref{decent-spaces-lemma-decent-space-closed-point}）。于是
$\Spec(\kappa) \to \Spec(k)$ 为有限型，故 $\kappa$ 是 $k$ 的有限扩张
（由 Hilbert 零点定理，见《代数》定理 \ref{algebra-theorem-nullstellensatz}；
略去一些细节）。设 $[\kappa : k] = d$。选取与 $d$ 互素的整数 $n \gg 0$，并令
$k'/k$ 为次数为 $n$ 的扩张。则 $k'/k$ 为 Galois 扩张，且
$G = \text{Aut}(k'/k)$ 是阶为 $n$ 的循环群。当 $n$ 足够大时，与上文相同的理由说明存在
$k$-代数同态 $A \to k'$。
% P09-ERR-0273
于是 $X_{k'}$ 是概形，且 $X = X_{k'}/G$（引理
\ref{spaces-over-fields-lemma-base-change-by-Galois}）。另一方面，由于 $n$ 与 $d$ 互素，有
$$
\Spec(\kappa) \times_{X} X_{k'} =
\Spec(\kappa) \times_{\Spec(k)} \Spec(k') =
\Spec(\kappa \otimes_k k')
$$
是某个域的谱。换言之，$X_{k'} \to X$ 在 $x$ 上方的纤维只有一个点。因此由引理
\ref{spaces-over-fields-lemma-when-quotient-scheme-at-point}，$x$ 如所需地位于 $X$ 的概形轨迹中。
\end{proof}

\begin{remark}
\label{spaces-over-fields-remark-when-does-the-argument-work}
设 $k$ 为有限域，$K/k$ 为几何不可约域扩张。则 $K$ 是几何不可约有限型 $k$-代数
$A$ 的极限。给定 $A$，Lang 与 Weil 的估计 \cite{LW} 说明，当 $n \gg 0$ 时，
存在 $k$-代数同态 $A \to k'$，其中 $k'/k$ 的次数为 $n$。
% P09-ERR-0274
分析引理 \ref{spaces-over-fields-lemma-scheme-after-purely-transcendental-base-change} 证明中的论证可知：
若 $X$ 是 $k$ 上的拟分离代数空间，且 $X_K$ 为概形，则 $X$ 为概形。若以后需要这一结果，
我们将在此作出准确陈述并给出证明。
% P09-ERR-0275
\end{remark}

\begin{lemma}
\label{spaces-over-fields-lemma-scheme-over-algebraic-closure-enough-affines}
设 $k$ 为域，其代数闭包为 $\overline{k}$。设 $X$ 为 $k$ 上满足下列条件的代数空间：
\begin{enumerate}
\item $X$ 良态且在 $k$ 上局部有限型；
\item $X_{\overline{k}}$ 为概形；以及
\item $X_{\overline{k}}$ 的任意有限个 $\overline{k}$-有理点都包含在某个仿射开子集中。
% P09-ERR-0276
\end{enumerate}
则 $X$ 为概形。
\end{lemma}

\begin{proof}
若 $K/k$ 为扩张，则基变换 $X_K$ 良态（《良态空间》引理
\ref{decent-spaces-lemma-representable-named-properties}），并且在 $K$ 上局部有限型
（《空间的态射》引理 \ref{spaces-morphisms-lemma-base-change-finite-type}）。
由引理 \ref{spaces-over-fields-lemma-scheme-after-purely-inseparable-base-change}，只需证明 $X$ 在基变换到
$k$ 的完美化后成为概形；因此可假设 $k$ 为完美域（这一步并非严格必要，但可使其余论证
更易理解）。用拟紧开子集覆盖 $X$ 可知，只需在 $X$ 拟紧时证明本引理（略去一处小细节）。
在此情形，$|X|$ 是良态拓扑空间（《良态空间》命题
\ref{decent-spaces-proposition-reasonable-sober}）。故只需证明 $|X|$ 中每个闭点都位于
$X$ 的概形轨迹中（使用《空间的性质》引理
\ref{spaces-properties-lemma-subscheme} 及《拓扑》引理
\ref{topology-lemma-quasi-compact-closed-point}）。

\medskip\noindent
设 $x \in |X|$ 为闭点。由《良态空间》引理
\ref{decent-spaces-lemma-decent-space-closed-point}，可找到表示 $x$ 的闭浸入
$\Spec(l) \to X$。态射 $\Spec(l) \to \Spec(k)$ 为有限型（《空间的态射》引理
\ref{spaces-morphisms-lemma-composition-finite-type}），因此由 Hilbert 零点定理
（《代数》定理 \ref{algebra-theorem-nullstellensatz}），$l$ 是 $k$ 的有限扩张。
由于 $k$ 完美，该扩张可分。因此概形
$$
\Spec(l) \times_X X_{\overline{k}} =
\Spec(l) \times_{\Spec(k)} \Spec(\overline{k}) =
\Spec(l \otimes_k \overline{k})
$$
是有限多个 $\overline{k}$-有理点的不交并。由假设 (3)，可以找到包含这些点的仿射开子集
$W \subset X_{\overline{k}}$。

\medskip\noindent
由引理 \ref{spaces-over-fields-lemma-when-scheme-after-base-change}，对某个有限扩张 $k'/k$，
$X_{k'}$ 为概形。扩大 $k'$ 后，可假设存在仿射开子集 $U' \subset X_{k'}$，其到
$\overline{k}$ 的基变换恢复 $W$（使用 $X_{\overline{k}}$ 是概形 $X_{k''}$ 的极限这一事实，
其中 $k' \subset k'' \subset \overline{k}$ 为有限扩张，并使用《极限》诸引理
\ref{limits-lemma-descend-opens} 及 \ref{limits-lemma-limit-affine}）。可以假设
$k'/k$ 为 Galois 扩张（取正规闭包，见《域》引理
\ref{fields-lemma-normal-closure}，并使用 $k$ 完美这一事实）。
% P09-ERR-0277
令 $G = \text{Gal}(k'/k)$。依构造，$G$-不变闭子概形
$\Spec(l) \times_X X_{k'}$ 包含在 $U'$ 中。因此由引理
\ref{spaces-over-fields-lemma-base-change-by-Galois} 及
\ref{spaces-over-fields-lemma-when-quotient-scheme-at-point}，$x$ 位于概形轨迹中。
\end{proof}

\noindent
下列两个引理应当移到别处。请将下一个引理与《良态空间》引理
\ref{decent-spaces-lemma-conditions-on-space-over-field} 比较。
% P09-ERR-0278

\begin{lemma}
\label{spaces-over-fields-lemma-locally-quasi-finite-over-field}
设 $k$ 为域，$X$ 为 $k$ 上的代数空间。下列条件等价：
\begin{enumerate}
\item $X$ 在 $k$ 上局部拟有限；
\item $X$ 在 $k$ 上局部有限型，且维数为 $0$；
\item $X$ 为概形，且在 $k$ 上局部拟有限；
\item $X$ 为概形、在 $k$ 上局部有限型且维数为 $0$；以及
\item $X$ 是一些阿廷局部 $k$-代数 $A$ 在 $k$ 上的谱的不交并，并且
$\dim_k(A) < \infty$。
\end{enumerate}
\end{lemma}

\begin{proof}
由于基底是域，$X/k$ 的相对维数与 $X$ 的维数相同。
% P09-ERR-0279
因此由《空间的态射》引理
\ref{spaces-morphisms-lemma-locally-quasi-finite-rel-dimension-0}，(1) 与 (2) 等价。
进而由引理 \ref{spaces-over-fields-lemma-locally-finite-type-dim-zero}（以及平凡的蕴含关系），
(1)--(4) 等价。最后，《簇》引理
\ref{varieties-lemma-algebraic-scheme-dim-0} 说明 (1)--(4) 与 (5) 等价。
% P09-ERR-0280
\end{proof}

\begin{lemma}
\label{spaces-over-fields-lemma-mono-towards-locally-quasi-finite-over-field}
设 $k$ 为域，$f : X \to Y$ 为 $k$ 上代数空间的单态射。若 $Y$ 在 $k$ 上局部拟有限，
则 $X$ 亦然。
\end{lemma}

\begin{proof}
假设 $Y$ 在 $k$ 上局部拟有限。由引理
\ref{spaces-over-fields-lemma-locally-quasi-finite-over-field}，有 $Y = \coprod \Spec(A_i)$，其中每个
$A_i$ 都是在 $k$ 上有限的阿廷局部环。由《良态空间》引理
\ref{decent-spaces-lemma-monomorphism-toward-disjoint-union-dim-0-rings}，
$X$ 为概形。考虑 $X_i = f^{-1}(\Spec(A_i))$。则 $X_i$ 或有一个点，或没有点。
若 $X_i$ 没有点，则无须证明。若 $X_i$ 有一个点，则 $X_i = \Spec(B_i)$，其中 $B_i$
为零维局部环，且 $A_i \to B_i$ 是环的满态射。特别地，
$A_i/\mathfrak m_{A_i} = B_i/\mathfrak m_{A_i}B_i$；由中山引理，即《代数》引理
\ref{algebra-lemma-NAK}（因为 $\mathfrak m_{A_i}$ 是幂零理想！），可知
$A_i \to B_i$ 为满射。因此 $B_i$ 是有限局部 $k$-代数，再由引理
\ref{spaces-over-fields-lemma-locally-quasi-finite-over-field} 可知 $X \to \Spec(k)$ 局部拟有限。
\end{proof}











\section{几何约化代数空间}
\label{spaces-over-fields-section-geometrically-reduced}

\noindent
若 $X$ 是域上的约化代数空间，则扩张基域后，$X$ 可能变得非约化。对于几何约化代数空间，
这种情形不会发生。

\begin{definition}
\label{spaces-over-fields-definition-geometrically-reduced}
设 $k$ 为域，$X$ 为 $k$ 上的代数空间。
\begin{enumerate}
\item 设 $x \in |X|$ 为点。若 $\mathcal{O}_{X, \overline{x}}$ 在 $k$ 上几何约化，
则称 $X${\it 在 $x$ 处几何约化}。
\item 若 $X$ 在 $X$ 的每个点处都几何约化，则称 $X$ 在 $k$ 上{\it 几何约化}。
\end{enumerate}
\end{definition}

\noindent
注意，若 $X$ 在 $x$ 处几何约化，则 $X$ 在 $x$ 处的局部环约化（《空间的性质》引理
\ref{spaces-properties-lemma-reduced-local-ring}）。同理，若 $X$ 在 $k$ 上几何约化，
则 $X$ 约化（由《空间的性质》引理
\ref{spaces-properties-lemma-reduced-space}）。特别地，下列引理说明该定义与域上概形的
相应性质并不冲突。

\begin{lemma}
\label{spaces-over-fields-lemma-geometrically-reduced-at-point}
设 $k$ 为域，$X$ 为 $k$ 上的代数空间。设 $x \in |X|$。下列条件等价：
\begin{enumerate}
\item $X$ 在 $x$ 处几何约化；
\item 对某个平展邻域 $(U, u) \to (X, x)$，其中 $U$ 为概形，$U$ 在 $u$ 处几何约化；
\item 对任意平展邻域 $(U, u) \to (X, x)$，其中 $U$ 为概形，$U$ 在 $u$ 处几何约化。
\end{enumerate}
\end{lemma}

\begin{proof}
回忆，局部环 $\mathcal{O}_{X, \overline{x}}$ 是 $\mathcal{O}_{U, u}$ 的严格亨泽尔化，
见《空间的性质》引理
\ref{spaces-properties-lemma-describe-etale-local-ring}。由《簇》引理
\ref{varieties-lemma-geometrically-reduced-at-point}，$U$ 在 $u$ 处几何约化，当且仅当
$\mathcal{O}_{U, u}$ 在 $k$ 上几何约化。因此需要证明：若 $A$ 为局部 $k$-代数，则
$A$ 在 $k$ 上几何约化，当且仅当 $A^{sh}$ 在 $k$ 上几何约化。使用几何约化代数的定义
（《代数》定义 \ref{algebra-definition-geometrically-reduced}）来验证这一点。设 $K/k$
为域扩张。由于 $A \to A^{sh}$ 忠实平坦（《代数进阶》引理
\ref{more-algebra-lemma-dumb-properties-henselization}），可知
$A \otimes_k K \to A^{sh} \otimes_k K$ 忠实平坦（《代数》引理
\ref{algebra-lemma-flat-base-change}）。故若 $A^{sh} \otimes_k K$ 约化，则由《代数》引理
\ref{algebra-lemma-descent-reduced}，$A \otimes_k K$ 也约化。反之，回忆 $A^{sh}$ 是平展
$A$-代数的余极限，见《代数》引理
% P09-ERR-0281
\ref{algebra-lemma-strict-henselization}。因此 $A^{sh} \otimes_k K$ 是平展
$A \otimes_k K$-代数的滤过余极限。由《代数》引理
\ref{algebra-lemma-reduced-goes-up} 即得结论。
\end{proof}

\begin{lemma}
\label{spaces-over-fields-lemma-geometrically-reduced}
设 $k$ 为域，$X$ 为 $k$ 上的代数空间。下列条件等价：
\begin{enumerate}
\item $X$ 几何约化；
\item 对某个平展满态射 $U \to X$，其中 $U$ 为概形，$U$ 几何约化；
\item 对任意平展态射 $U \to X$，其中 $U$ 为概形，$U$ 几何约化。
\end{enumerate}
\end{lemma}

\begin{proof}
这直接由定义与引理 \ref{spaces-over-fields-lemma-geometrically-reduced-at-point} 得出。
\end{proof}

\noindent
该概念在特征零的情形并无新意。

\begin{lemma}
\label{spaces-over-fields-lemma-perfect-reduced}
设 $X$ 为完美域 $k$ 上的代数空间（例如 $k$ 的特征为零）。
\begin{enumerate}
\item 对 $x \in |X|$，若 $\mathcal{O}_{X, \overline{x}}$ 约化，则 $X$ 在 $x$ 处几何约化。
\item 若 $X$ 约化，则 $X$ 在 $k$ 上几何约化。
\end{enumerate}
\end{lemma}

\begin{proof}
第一个陈述由《代数》引理
\ref{algebra-lemma-separable-extension-preserves-reducedness} 及完美域的定义
（《代数》定义 \ref{algebra-definition-perfect}）得出。第二个陈述由第一个陈述得出。
\end{proof}

\begin{lemma}
\label{spaces-over-fields-lemma-geometrically-reduced-positive-characteristic}
设 $k$ 为特征 $p > 0$ 的域，$X$ 为 $k$ 上的代数空间。下列条件等价：
\begin{enumerate}
\item $X$ 在 $k$ 上几何约化；
\item 对每个域扩张 $k'/k$，$X_{k'}$ 约化；
\item 对每个有限纯不可分域扩张 $k'/k$，$X_{k'}$ 约化；
\item $X_{k^{1/p}}$ 约化；
\item $X_{k^{perf}}$ 约化；以及
\item $X_{\bar k}$ 约化。
\end{enumerate}
\end{lemma}

\begin{proof}
选取平展满态射 $U \to X$，其中 $U$ 为概形。借助引理
\ref{spaces-over-fields-lemma-geometrically-reduced}，本引理由 $U$ 在 $k$ 上的相应结果推出。
见《簇》引理 \ref{varieties-lemma-geometrically-reduced}。
\end{proof}

\begin{lemma}
\label{spaces-over-fields-lemma-geometrically-reduced-upstairs}
设 $k$ 为域，$X$ 为 $k$ 上的代数空间。设 $k'/k$ 为域扩张，$x \in |X|$ 为点，
$x' \in |X_{k'}|$ 为 $x$ 上方的点。下列条件等价：
\begin{enumerate}
\item $X$ 在 $x$ 处几何约化；
\item $X_{k'}$ 在 $x'$ 处几何约化。
\end{enumerate}
特别地，$X$ 在 $k$ 上几何约化，当且仅当 $X_{k'}$ 在 $k'$ 上几何约化。
\end{lemma}

\begin{proof}
选取平展态射 $U \to X$，其中 $U$ 为概形；再选取映到 $x \in |X|$ 的点 $u \in U$。
由《空间的性质》引理 \ref{spaces-properties-lemma-points-cartesian}，可以选取同时映到
$u$ 与 $x'$ 的点 $u' \in U_{k'} = U \times_X X_{k'}$。借助引理
\ref{spaces-over-fields-lemma-geometrically-reduced-at-point}，本引理由 $U, u, u'$ 的相应引理推出，
后者即《簇》引理 \ref{varieties-lemma-geometrically-reduced-upstairs}。
\end{proof}

\begin{lemma}
\label{spaces-over-fields-lemma-geometrically-reduced-etale-local}
设 $k$ 为域，$f : X \to Y$ 为 $k$ 上代数空间的态射。设 $x \in |X|$ 为点，
其像为 $y \in |Y|$。
\begin{enumerate}
\item 若 $f$ 在 $x$ 处平展，则
$X$ 在 $x$ 处几何约化 $\Leftrightarrow$
$Y$ 在 $y$ 处几何约化；
\item 若 $f$ 平展且满，则
$X$ 几何约化 $\Leftrightarrow$
$Y$ 几何约化。
\end{enumerate}
\end{lemma}

\begin{proof}
若 $f$ 在 $x$ 处平展，则
$\mathcal{O}_{X, \overline{x}} = \mathcal{O}_{Y, \overline{y}}$，故 (1) 显然。
(2) 立即由 (1) 推出。
\end{proof}






\section{几何连通代数空间}
\label{spaces-over-fields-section-geometrically-connected}

\noindent
若 $X$ 是域上的连通代数空间，则扩张基域后，$X$ 可能变得不连通。对于几何连通
代数空间，这种情形不会发生。

\begin{definition}
\label{spaces-over-fields-definition-geometrically-connected}
设 $X$ 为域 $k$ 上的代数空间。若对 $k$ 的每个域扩张 $k'$，基变换 $X_{k'}$
都连通，则称 $X$ 在 $k$ 上{\it 几何连通}。
\end{definition}

\noindent
依约定，连通拓扑空间非空；因而几何连通代数空间当然非空。

\begin{lemma}
\label{spaces-over-fields-lemma-geometrically-connected-check-after-extension}
设 $X$ 为域 $k$ 上的代数空间，$k'/k$ 为域扩张。则 $X$ 在 $k$ 上几何连通，
当且仅当 $X_{k'}$ 在 $k'$ 上几何连通。
\end{lemma}

\begin{proof}
若 $X$ 在 $k$ 上几何连通，则 $X_{k'}$ 在 $k'$ 上几何连通是显然的。反之，对任意
域扩张 $k''/k$，存在 $k'/k$ 与 $k''/k$ 的共同扩域 $k'''/k$，见《域》引理
\ref{fields-lemma-common-extension-field}。态射 $X_{k'''} \to X_{k''}$ 是满的
（它是域的谱之间某个满态射的基变换），故 $X_{k'''}$ 连通蕴含 $X_{k''}$ 连通。
因此，若 $X_{k'}$ 在 $k'$ 上几何连通，则 $X$ 在 $k$ 上几何连通。
\end{proof}

\begin{lemma}
\label{spaces-over-fields-lemma-bijection-connected-components}
设 $k$ 为域，$X$, $Y$ 为 $k$ 上的代数空间。假设 $X$ 在 $k$ 上几何连通。则投影态射
$$
p : X \times_k Y \longrightarrow Y
$$
诱导连通分支之间的双射。
\end{lemma}

\begin{proof}
设 $y \in |Y|$ 由态射 $\Spec(K) \to Y$ 表示，其中 $K$ 为域。由《空间的性质》引理
\ref{spaces-properties-lemma-points-cartesian}，$|X \times_k Y| \to |Y|$
在 $y$ 上的纤维是 $|X_K| \to |X \times_k Y|$ 的像。因此，由 $X$ 几何连通的假设，
这些纤维连通。由《空间的态射》引理
\ref{spaces-morphisms-lemma-space-over-field-universally-open}，映射 $|p|$ 是开映射。
于是应用《拓扑》引理
\ref{topology-lemma-connected-fibres-connected-components} 即得结论。
\end{proof}

\begin{lemma}
\label{spaces-over-fields-lemma-separably-closed-field-connected-components}
设 $k'/k$ 为域扩张，$X$ 为 $k$ 上的代数空间。假设 $k$ 是可分代数闭域。则态射
$X_{k'} \to X$ 诱导连通分支之间的双射。特别地，$X$ 在 $k$ 上几何连通，
当且仅当 $X$ 连通。
\end{lemma}

\begin{proof}
由于 $k$ 可分代数闭，$k'$ 在 $k$ 上几何连通，见《代数》引理
\ref{algebra-lemma-separably-closed-connected-implies-geometric}。因此，由《簇》引理
\ref{varieties-lemma-affine-geometrically-connected}，$Z = \Spec(k')$ 在 $k$
上几何连通。由于 $X_{k'} = Z \times_k X$，该结论是引理
\ref{spaces-over-fields-lemma-bijection-connected-components} 的特殊情形。
\end{proof}

\begin{lemma}
\label{spaces-over-fields-lemma-characterize-geometrically-connected}
设 $k$ 为域，$X$ 为 $k$ 上的代数空间，$\overline{k}$ 为 $k$ 的可分代数闭包。
则 $X$ 几何连通，当且仅当基变换 $X_{\overline{k}}$ 连通。
\end{lemma}

\begin{proof}
假设 $X_{\overline{k}}$ 连通。设 $k'/k$ 为域扩张。存在域扩张
$\overline{k}'/\overline{k}$，使得 $k'$ 作为 $k$ 的扩张嵌入 $\overline{k}'$。
由引理 \ref{spaces-over-fields-lemma-separably-closed-field-connected-components}，
$X_{\overline{k}'}$ 连通。由于 $X_{\overline{k}'} \to X_{k'}$ 是满态射，
可得所需的 $X_{k'}$ 连通性。
\end{proof}

\noindent
设 $k$ 为域，$\overline{k}/k$ 为（可能无限的）Galois 扩张。例如，$\overline{k}$
可以是 $k$ 的可分代数闭包。对任意 $\sigma \in \text{Gal}(\overline{k}/k)$，
得到相应的自同构
$
\Spec(\sigma) :
\Spec(\overline{k})
\longrightarrow
\Spec(\overline{k})
$.
注意
$\Spec(\sigma) \circ \Spec(\tau) = \Spec(\tau \circ \sigma)$。因此得到相反群在概形
$\Spec(\overline{k})$ 上的作用
$$
\text{Gal}(\overline{k}/k)^{opp} \times \Spec(\overline{k})
\longrightarrow
\Spec(\overline{k})
$$
设 $X$ 为 $k$ 上的代数空间。由于按定义
$X_{\overline{k}} =
\Spec(\overline{k}) \times_{\Spec(k)} X$
，上述作用诱导典范作用
\begin{equation}
\label{spaces-over-fields-equation-galois-action-base-change-kbar}
\text{Gal}(\overline{k}/k)^{opp} \times X_{\overline{k}}
\longrightarrow
X_{\overline{k}}.
\end{equation}

\begin{lemma}
\label{spaces-over-fields-lemma-Galois-action-quasi-compact-open}
设 $k$ 为域，$X$ 为 $k$ 上的代数空间，$\overline{k}$ 为 $k$ 的（可能无限的）
Galois 扩张。设 $V \subset X_{\overline{k}}$ 为拟紧开子空间。则：
\begin{enumerate}
\item 存在有限子扩张 $\overline{k}/k'/k$ 及拟紧开子空间 $V' \subset X_{k'}$，
使得 $V = (V')_{\overline{k}}$；
\item 存在开子群 $H \subset \text{Gal}(\overline{k}/k)$，使得对所有
$\sigma \in H$ 都有 $\sigma(V) = V$。
\end{enumerate}
\end{lemma}

\begin{proof}
选取概形 $U$ 及平展满态射 $U \to X$。选取拟紧开子空间
$W \subset U_{\overline{k}}$，使其在 $X_{\overline{k}}$ 中的像为 $V$。
这是可能的，因为 $|U_{\overline{k}}| \to |X_{\overline{k}}|$ 连续，且
$|U_{\overline{k}}|$ 有由拟紧开集组成的基。对 $W \subset U_{\overline{k}}$
应用《簇》引理 \ref{varieties-lemma-Galois-action-quasi-compact-open}，即得本引理。
\end{proof}

\begin{lemma}
\label{spaces-over-fields-lemma-closed-fixed-by-Galois}
设 $k$ 为域，$\overline{k}/k$ 为（可能无限的）Galois 扩张，$X$ 为 $k$ 上的
代数空间。设 $\overline{T} \subset |X_{\overline{k}}|$ 具有下列性质：
\begin{enumerate}
\item $\overline{T}$ 是 $|X_{\overline{k}}|$ 的闭子集；
\item 对每个 $\sigma \in \text{Gal}(\overline{k}/k)$，都有
$\sigma(\overline{T}) = \overline{T}$。
\end{enumerate}
则存在闭子集 $T \subset |X|$，使其在 $|X_{k'}|$ 中的逆像为 $\overline{T}$。
% P09-ERR-0282
\end{lemma}

\begin{proof}
令 $T \subset |X|$ 为 $\overline{T}$ 的像。由于
$|X_{\overline{k}}| \to |X|$ 是满射，该陈述等价于 $T$ 闭且其逆像为
$\overline{T}$。选取概形 $U$ 及平展满态射 $U \to X$。由概形的情形
（见《簇》引理 \ref{varieties-lemma-closed-fixed-by-Galois}），存在闭子集
$T' \subset |U|$，其在 $|U_{\overline{k}}|$ 中的逆像就是 $\overline{T}$
的逆像。由于 $|U_{\overline{k}}| \to |X_{\overline{k}}|$ 是满射，
可知 $T'$ 是 $T$ 经 $|U| \to |X|$ 的逆像。由 $|X|$ 上拓扑的构造，
这说明 $T$ 闭。以同样方式可知，$\overline{T}$ 是 $T$ 的逆像。
\end{proof}

\begin{lemma}
\label{spaces-over-fields-lemma-characterize-geometrically-disconnected}
设 $k$ 为域，$X$ 为 $k$ 上的代数空间。下列条件等价：
\begin{enumerate}
\item $X$ 几何连通；
\item 对每个有限可分域扩张 $k'/k$，代数空间 $X_{k'}$ 连通。
\end{enumerate}
\end{lemma}

\begin{proof}
本证明与《簇》引理
\ref{varieties-lemma-characterize-geometrically-disconnected} 的证明相同，
区别仅在于：将《簇》引理
\ref{varieties-lemma-characterize-geometrically-connected} 替换为引理
\ref{spaces-over-fields-lemma-characterize-geometrically-connected}；将《簇》引理
\ref{varieties-lemma-Galois-action-quasi-compact-open} 替换为引理
\ref{spaces-over-fields-lemma-Galois-action-quasi-compact-open}；并将《簇》引理
\ref{varieties-lemma-closed-fixed-by-Galois} 替换为引理
\ref{spaces-over-fields-lemma-closed-fixed-by-Galois}。我们建议读者阅读那个证明，而不是下面这个证明。
% P09-ERR-0283

\medskip\noindent
由定义立即可知 (1) 蕴含 (2)。假设 $X$ 不是几何连通的。设
$k \subset \overline{k}$ 为 $k$ 的可分代数闭包。由引理
\ref{spaces-over-fields-lemma-characterize-geometrically-connected}，$X_{\overline{k}}$ 不连通。
设 $X_{\overline{k}} = \overline{U} \amalg \overline{V}$，其中
$\overline{U}$ 与 $\overline{V}$ 是 $X_{\overline{k}}$ 的非空开闭代数子空间。

\medskip\noindent
设 $W \subset X$ 为任意拟紧开子空间。则
$W_{\overline{k}} \cap \overline{U}$ 与 $W_{\overline{k}} \cap \overline{V}$
都是 $W_{\overline{k}}$ 的开闭子空间。特别地，
$W_{\overline{k}} \cap \overline{U}$ 与 $W_{\overline{k}} \cap \overline{V}$
拟紧；并且由引理 \ref{spaces-over-fields-lemma-Galois-action-quasi-compact-open}，
$W_{\overline{k}} \cap \overline{U}$ 与 $W_{\overline{k}} \cap \overline{V}$
都定义在某个有限子扩张上，
且在 $\text{Gal}(\overline{k}/k)$ 的某个开子群作用下不变。下文将直接使用这一点，
不再另行说明。

\medskip\noindent
选取拟紧开子空间 $W_0 \subset X$，使得
$W_{0, \overline{k}} \cap \overline{U}$ 与
$W_{0, \overline{k}} \cap \overline{V}$ 均非空。选取有限子扩张
$\overline{k}/k'/k$，以及分解 $W_{0, k'} = U_0' \amalg V_0'$，
其中两部分均开且闭，并且
$W_{0, \overline{k}} \cap \overline{U} = (U'_0)_{\overline{k}}$ 且
$W_{0, \overline{k}} \cap \overline{V} = (V'_0)_{\overline{k}}$。
令 $H = \text{Gal}(\overline{k}/k') \subset \text{Gal}(\overline{k}/k)$。
特别地，
$\sigma(W_{0, \overline{k}} \cap \overline{U}) =
W_{0, \overline{k}} \cap \overline{U}$，对 $\overline{V}$ 亦然。

\medskip\noindent
按上述方式选定 $W_0$ 与 $k'$ 后，对每个拟紧开子空间 $W \subset X$，令
$$
U_W =
\bigcap\nolimits_{\sigma \in H} \sigma(W_{\overline{k}} \cap \overline{U}),
\quad
V_W =
\bigcup\nolimits_{\sigma \in H} \sigma(W_{\overline{k}} \cap \overline{V}).
$$
由于 $W_{\overline{k}} \cap \overline{U}$ 与
$W_{\overline{k}} \cap \overline{V}$ 在 $\text{Gal}(\overline{k}/k)$ 的某个
开子群作用下不变，上述并与交实际上都是有限的。因此 $U_W$ 与 $V_W$ 均为开闭子空间。
此外，由构造有 $W_{\bar k} = U_W \amalg V_W$。

\medskip\noindent
我们断言：若 $W \subset W' \subset X$ 为拟紧开子空间，则
$W_{\overline{k}} \cap U_{W'} = U_W$ 且
$W_{\overline{k}} \cap V_{W'} = V_W$。验证从略。因此，定义
$U = \bigcup_{W \subset X} U_W$ 与 $V = \bigcup_{W \subset X} V_W$ 后，得到
$X_{\overline{k}} = U \amalg V$，它是开闭子集的不交并。
% P09-ERR-0284
显然，$V$ 非空，因为它（局部地）由取并构造。另一方面，$U$ 非空，因为由构造它包含
$W_0 \cap \overline{U}$。
% P09-ERR-0285
最后，由构造，$U, V \subset X_{\bar k}$ 闭且 $H$-不变。因此，由引理
\ref{spaces-over-fields-lemma-closed-fixed-by-Galois}，存在闭子集 $U', V' \subset X_{k'}$，使得
$U = (U')_{\bar k}$ 且 $V = (V')_{\bar k}$。显然
$X_{k'} = U' \amalg V'$，故 $X_{k'}$ 如所需是不连通的。
\end{proof}






\section{几何不可约代数空间}
\label{spaces-over-fields-section-geometrically-irreducible}

\noindent
《空间》例 \ref{spaces-example-infinite-product} 表明，最好不要在完全一般的情形下
考虑不可约代数空间\footnote{准确地说，若我们说“代数空间 $X$ 不可约”，大概是指
“拓扑空间 $|X|$ 不可约”。}。对于良态（例如拟分离）代数空间，这类灾难不会发生。
因此，仅在代数空间良态的假设下作出下列定义。

\begin{definition}
\label{spaces-over-fields-definition-geometrically-irreducible}
设 $k$ 为域，$X$ 为 $k$ 上的良态代数空间。若对 $k$ 的任意域扩张 $k'$，
拓扑空间 $|X_{k'}|$ 都不可约\footnote{不可约空间非空。}，则称 $X$
{\it 几何不可约}。
\end{definition}

\noindent
注意，$X_{k'}$ 是良态代数空间（《良态空间》引理
\ref{decent-spaces-lemma-representable-named-properties}）。因此，拓扑空间
$|X_{k'}|$ 是良态拓扑空间。《良态空间》命题
\ref{decent-spaces-proposition-reasonable-sober}。
% P09-ERR-0286





\section{几何整代数空间}
\label{spaces-over-fields-section-geometrically-integral}

\noindent
回忆，整代数空间按定义是良态的，见第 \ref{spaces-over-fields-section-integral-spaces} 节。

\begin{definition}
\label{spaces-over-fields-definition-geometrically-integral}
设 $X$ 为域 $k$ 上的代数空间。若对 $k$ 的每个域扩张 $k'$，代数空间
$X_{k'}$ 都整（定义 \ref{spaces-over-fields-definition-integral-algebraic-space}），
则称 $X$ 在 $k$ 上{\it 几何整}。
\end{definition}

\noindent
特别地，$X$ 是良态代数空间。它与几何约化及几何不可约之间有如下关系。

\begin{lemma}
\label{spaces-over-fields-lemma-geometrically-integral}
设 $k$ 为域，$X$ 为 $k$ 上的良态代数空间。则 $X$ 在 $k$ 上几何整，
当且仅当 $X$ 在 $k$ 上既几何约化又几何不可约。
\end{lemma}

\begin{proof}
这直接由定义得出，因为在良态性的条件下，我们的整性概念等价于约化且不可约。
\end{proof}

\begin{lemma}
\label{spaces-over-fields-lemma-proper-geometrically-reduced-global-sections}
设 $k$ 为域，$X$ 为 $k$ 上的固有代数空间。
\begin{enumerate}
\item $A = H^0(X, \mathcal{O}_X)$ 是有限维 $k$-代数；
\item $A = \prod_{i = 1, \ldots, n} A_i$ 是阿廷局部 $k$-代数的乘积，
$|X|$ 的每个连通分支对应一个因子；
\item 若 $X$ 约化，则 $A = \prod_{i = 1, \ldots, n} k_i$ 是域的乘积，
每个域都是 $k$ 的有限扩张；
\item 若 $X$ 几何约化，则 $k_i$ 是 $k$ 的有限可分扩张；
\item 若 $X$ 几何连通，则 $A$ 在 $k$ 上几何不可约；
\item 若 $X$ 几何不可约，则 $A$ 在 $k$ 上几何不可约；
\item 若 $X$ 几何约化且连通，则 $A = k$；以及
% P09-ERR-0288
\item 若 $X$ 几何整，则 $A = k$。
\end{enumerate}
\end{lemma}

\begin{proof}
由《空间的上同调》引理
\ref{spaces-cohomology-lemma-proper-over-affine-cohomology-finite}，
$A = H^0(X, \mathcal{O}_X)$ 是有限维 $k$-代数。这证明了 (1)。

\medskip\noindent
由《代数》引理 \ref{algebra-lemma-finite-dimensional-algebra} 与《代数》命题
\ref{algebra-proposition-dimension-zero-ring}，$A$ 是局部环的乘积。若
$X = Y \amalg Z$，其中 $Y$ 与 $Z$ 是 $X$ 的开子空间，则取 $\mathcal{O}_X$
在 $Y$ 上为 $1$、在 $Z$ 上为 $0$ 的截面，得到幂等元 $e \in A$。反之，若
$e \in A$ 为幂等元，则得到 $|X|$ 的相应分解。最后，由《空间的态射》引理
\ref{spaces-morphisms-lemma-finite-presentation-noetherian} 与《空间的性质》引理
\ref{spaces-properties-lemma-Noetherian-topology}，$|X|$ 是诺特拓扑空间，
故其连通分支都是开的。因此，$|X|$ 的连通分支与 $A$ 的本原幂等元之间存在 $1$ 对 $1$ 的对应。
这证明了 (2)。

\medskip\noindent
若 $X$ 约化，则 $A$ 约化（《空间的性质》引理
\ref{spaces-properties-lemma-reduced-space}）。因此，局部环 $A_i = k_i$
约化，从而是域（例如由《代数》引理
\ref{algebra-lemma-minimal-prime-reduced-ring}）。这证明了 (3)。

\medskip\noindent
若 $X$ 几何约化，则对
$A \otimes_k \overline{k} =
H^0(X_{\overline{k}}, \mathcal{O}_{X_{\overline{k}}})$
也有同样结论（该等式见《空间的上同调》引理
\ref{spaces-cohomology-lemma-flat-base-change-cohomology}）。
% P09-ERR-0287
这蕴含 $k_i \otimes_k \overline{k}$ 是域的乘积，因而例如由《代数》引理
\ref{algebra-lemma-characterize-separable-field-extensions} 与
\ref{algebra-lemma-geometrically-reduced-finite-purely-inseparable-extension}，
$k_i/k$ 可分。这证明了 (4)。

\medskip\noindent
若 $X$ 几何连通，则 $A \otimes_k \overline{k} =
H^0(X_{\overline{k}}, \mathcal{O}_{X_{\overline{k}}})$
由 (2) 是零维局部环，故其谱只有一个点，特别地不可约。因此 $A$ 几何不可约。
这证明了 (5)。当然，(5) 蕴含 (6)。

\medskip\noindent
若 $X$ 几何约化且连通，则 $A = k_1$ 是域，而扩张 $k_1/k$ 有限可分且几何不可约。
但此时 $k_1 \otimes_k \overline{k}$ 是 $[k_1 : k]$ 个 $\overline{k}$ 的乘积，
从而 $k_1 = k$。这证明了 (7)。当然，(7) 蕴含 (8)。
\end{proof}

\begin{lemma}
\label{spaces-over-fields-lemma-characterize-trivial-pic-integral}
设 $k$ 为域，$X$ 为 $k$ 上的固有整代数空间，$\mathcal{L}$ 为可逆
$\mathcal{O}_X$-模。若 $H^0(X, \mathcal{L})$ 与
$H^0(X, \mathcal{L}^{\otimes - 1})$ 均非零，则
$\mathcal{L} \cong \mathcal{O}_X$。
\end{lemma}

\begin{proof}
设 $s \in H^0(X, \mathcal{L})$ 与
$t \in H^0(X, \mathcal{L}^{\otimes - 1})$ 为非零截面。设 $x \in |X|$
为 $s$ 的支撑中的点。选取仿射平展邻域 $(U, u) \to (X, x)$，使得
$\mathcal{L}|_U \cong \mathcal{O}_U$。于是 $s|_U$ 对应于约化概形 $U$
（因为 $X$ 约化）上的非零正则函数，故它在 $U$ 的某个不可约分支的一般点处不消失。
% P09-ERR-0289
由《良态空间》引理 \ref{decent-spaces-lemma-decent-generic-points}，
可知 $|X|$ 的一般点 $\eta$ 属于 $s$ 的支撑。对 $t$ 亦然。于是 $st$ 必非零，
因为 $X$ 在 $\eta$ 处的局部环是域（由上述引理，该局部环维数为零；由于 $X$ 约化，
该局部环约化；再应用《代数》引理
\ref{algebra-lemma-minimal-prime-reduced-ring}）。然而，在引理
\ref{spaces-over-fields-lemma-proper-geometrically-reduced-global-sections} 中已经看到
$K = H^0(X, \mathcal{O}_X)$ 是域。因此 $st$ 处处非零，从而
$s : \mathcal{O}_X \to \mathcal{L}$ 是同构。
\end{proof}






\section{维数}
\label{spaces-over-fields-section-dimension}

\noindent
本节继续讨论维数。此前的材料列举如下：
\begin{enumerate}
\item 维数定义于《空间的性质》第
\ref{spaces-properties-section-dimension} 节；
\item 局部环的维数定义于《空间的性质》第
\ref{spaces-properties-section-dimension-local-ring} 节；
\item 《空间的性质》引理 \ref{spaces-properties-lemma-dimension-local-ring} 与
\ref{spaces-properties-lemma-dimension-decent-invariant-under-etale} 中的若干结果；
\item 相对维数定义于《空间的态射》第
\ref{spaces-morphisms-section-relative-dimension} 节；
\item 关于纤维维数的结果见《空间的态射》第
\ref{spaces-morphisms-section-dimension-fibres} 节；
\item 维数公式的一个弱形式见《空间的态射》第
\ref{spaces-morphisms-section-dimension-formula} 节；
\item 关于光滑性与维数的一个结果见《空间的态射》引理
\ref{spaces-morphisms-lemma-smoothness-dimension-spaces}；
\item 对良态空间，维数等于 $\dim(|X|)$，见《良态空间》引理
\ref{decent-spaces-lemma-dimension-decent-space}；
% P09-ERR-0290
\item 关于拟有限映射与维数的结果见《良态空间》引理
\ref{decent-spaces-lemma-dimension-local-ring-quasi-finite} 与
\ref{decent-spaces-lemma-dimension-quasi-finite}。
\end{enumerate}
在《空间的态射进阶》第 \ref{spaces-more-morphisms-section-dimension} 节中，
我们将讨论有限型态射的纤维维数跳跃。

\begin{lemma}
\label{spaces-over-fields-lemma-integral-dimension}
设 $S$ 为概形，$f : X \to Y$ 为代数空间的整态射。则
$\dim(X) \leq \dim(Y)$。若 $f$ 为满态射，则 $\dim(X) = \dim(Y)$。
\end{lemma}

\begin{proof}
选取平展满态射 $V \to Y$，其中 $V$ 为概形。则 $U = X \times_Y V$ 是概形，
而 $U \to V$ 是整态射（若 $f$ 为满态射，它也为满态射）。由《空间的性质》引理
\ref{spaces-properties-lemma-dimension-decent-invariant-under-etale}，
有 $\dim(X) = \dim(U)$ 与 $\dim(Y) = \dim(V)$。因此，该结论由概形的情形推出，
即《态射》引理 \ref{morphisms-lemma-integral-dimension}。
\end{proof}

\begin{lemma}
\label{spaces-over-fields-lemma-alteration-dimension}
设 $S$ 为概形，$f : X \to Y$ 为 $S$ 上代数空间的态射。假设：
\begin{enumerate}
\item $Y$ 局部诺特；
\item $X$ 与 $Y$ 是整代数空间；
\item $f$ 支配；以及
\item $f$ 局部有限型。
\end{enumerate}
若 $x \in |X|$ 与 $y \in |Y|$ 是一般点，则
$$
\dim(X) \leq \dim(Y) + \text{超越次数 }x/y.
$$
若 $f$ 固有，则等号成立。
\end{lemma}

\begin{proof}
回忆，$|X|$ 与 $|Y|$ 都是不可约良态拓扑空间，见定义
\ref{spaces-over-fields-definition-integral-algebraic-space} 后的讨论。因此，$f$ 支配意味着
$|f|$ 把 $x$ 映到 $y$。此外，$x \in |X|$ 是 $X$ 的局部环维数为 $0$
的唯一点，见《良态空间》引理
\ref{decent-spaces-lemma-decent-generic-points}。由《空间的态射》引理
\ref{spaces-morphisms-lemma-dimension-formula-general}，对任意点
$x' \in |X|$，$X$ 在该点处的局部环维数不超过 $Y$ 在
$y' = f(x')$ 处的局部环维数加上 $x/y$ 的超越次数。由于 $X$ 与 $Y$ 的维数分别是
各点 $x'$ 与 $y'$ 处局部环维数的上确界（《空间的性质》引理
\ref{spaces-properties-lemma-dimension}），故该不等式成立。

\medskip\noindent
假设 $f$ 固有。设 $V \subset Y$ 为非空拟紧开子空间。若能对态射
$f^{-1}(V) \to V$ 证明等式，则也得到 $X \to Y$ 的等式。因此可以假设
$X$ 与 $Y$ 拟紧。注意，$X$ 作为局部诺特良态代数空间是拟分离的，见《良态空间》引理
\ref{decent-spaces-lemma-locally-Noetherian-decent-quasi-separated}。因此，可以选取
有限满态射 $Y' \to Y$，其中 $Y'$ 为概形，见《空间的极限》命题
\ref{spaces-limits-proposition-there-is-a-scheme-finite-over}。用适当闭子概形替换
$Y'$ 后，可假设 $Y'$ 整，例如见更一般的引理
\ref{spaces-over-fields-lemma-alteration-contained-in}。由同一引理，可以选取闭子空间
$X' \subset X \times_Y Y'$，使得 $X'$ 整且 $X' \to X$ 为有限满态射。此时
$X'$ 也局部诺特（《空间的态射》引理
\ref{spaces-morphisms-lemma-locally-finite-type-locally-noetherian}），
再应用一次《空间的极限》命题
\ref{spaces-limits-proposition-there-is-a-scheme-finite-over}，可以选取有限满态射
$X'' \to X'$，其中 $X''$ 为概形。与前面相同，可假设 $X''$ 整。图示为
$$
\xymatrix{
X'' \ar[d] \ar[r] & X \ar[d]^f \\
Y' \ar[r] & Y
}
$$
由引理 \ref{spaces-over-fields-lemma-integral-dimension}，有 $\dim(X'') = \dim(X)$ 与
$\dim(Y') = \dim(Y)$。由于 $X$ 与 $Y$ 分别有 $x$ 与 $y$ 的概形开邻域，
容易看出一般点 $x'' \in X''$ 与 $y' \in Y'$ 分别是映到 $x$ 与 $y$ 的唯一点，
且剩余域扩张 $\kappa(x'')/\kappa(x)$ 与 $\kappa(y')/\kappa(y)$ 有限。
这蕴含 $x''/y'$ 与 $x/y$ 的超越次数相同。因此，等式由概形的情形推出，
即《态射》引理 \ref{morphisms-lemma-alteration-dimension}。
% P09-ERR-0291
\end{proof}





\section{域上的光滑空间}
\label{spaces-over-fields-section-smooth}

\noindent
本节是《簇》第 \ref{varieties-section-smooth} 节的对应版本。

\begin{lemma}
\label{spaces-over-fields-lemma-smooth-regular}
设 $k$ 为域，$X$ 为 $k$ 上光滑的代数空间。则 $X$ 是正则代数空间。
\end{lemma}

\begin{proof}
选取概形 $U$ 及平展满态射 $U \to X$。态射 $U \to \Spec(k)$ 是一个平展态射
（因而光滑）与一个光滑态射的复合，所以光滑（见《空间的态射》引理
\ref{spaces-morphisms-lemma-etale-smooth} 与
\ref{spaces-morphisms-lemma-composition-smooth}）。因此，由《簇》引理
\ref{varieties-lemma-smooth-regular}，$U$ 正则。按《空间的性质》定义
\ref{spaces-properties-definition-type-property}，这意味着 $X$ 正则。
\end{proof}

\begin{lemma}
\label{spaces-over-fields-lemma-smooth-separable-closed-points-dense}
设 $k$ 为域，$X$ 为 $\Spec(k)$ 上光滑的代数空间。所有可表示为态射
$\Spec(k') \to X$ 的像的点 $x \in |X|$ 构成 $|X|$ 中的稠密子集，
其中 $k' \supset k$ 有限可分。
% P09-ERR-0292
\end{lemma}

\begin{proof}
选取概形 $U$ 及平展满态射 $U \to X$。态射 $U \to \Spec(k)$ 是一个平展态射
（因而光滑）与一个光滑态射的复合，所以光滑（见《空间的态射》引理
\ref{spaces-morphisms-lemma-etale-smooth} 与
\ref{spaces-morphisms-lemma-composition-smooth}）。于是应用《簇》引理
\ref{varieties-lemma-smooth-separable-closed-points-dense} 可知，$U$ 中剩余域为
$k$ 的有限可分扩张的闭点稠密。由 $|X|$ 上拓扑的定义，这蕴含本引理。
\end{proof}







\section{欧拉示性数}
\label{spaces-over-fields-section-euler}

\noindent
本节证明域上固有代数空间的相干层之欧拉示性数的一些初等性质。

\begin{definition}
\label{spaces-over-fields-definition-euler-characteristic}
设 $k$ 为域，$X$ 为 $k$ 上的固有代数空间，$\mathcal{F}$ 为相干
$\mathcal{O}_X$-模。在此情形下，{\it $\mathcal{F}$ 的欧拉示性数}是整数
% P09-ERR-0293
$$
\chi(X, \mathcal{F}) = \sum\nolimits_i (-1)^i \dim_k H^i(X, \mathcal{F}).
$$
该公式的合理性见下文。
\end{definition}

\noindent
在该定义的情形下，向量空间 $H^i(X, \mathcal{F})$ 中只有有限多个非零
（《空间的上同调》引理
\ref{spaces-cohomology-lemma-vanishing-quasi-separated}），且这些空间都有限维
（《空间的上同调》引理
\ref{spaces-cohomology-lemma-proper-over-affine-cohomology-finite}）。
因此 $\chi(X, \mathcal{F}) \in \mathbf{Z}$ 定义良好。注意，该定义依赖于域 $k$，
而不只依赖于二元组 $(X, \mathcal{F})$。

\begin{lemma}
\label{spaces-over-fields-lemma-euler-characteristic-additive}
设 $k$ 为域，$X$ 为 $k$ 上的固有代数空间。设
$0 \to \mathcal{F}_1 \to \mathcal{F}_2 \to \mathcal{F}_3 \to 0$
为 $X$ 上相干模的短正合列。则
$$
\chi(X, \mathcal{F}_2) = \chi(X, \mathcal{F}_1) + \chi(X, \mathcal{F}_3)
$$
\end{lemma}

\begin{proof}
考虑与该引理中的短正合列相伴的上同调长正合列
$$
0 \to H^0(X, \mathcal{F}_1) \to H^0(X, \mathcal{F}_2) \to
H^0(X, \mathcal{F}_3) \to H^1(X, \mathcal{F}_1) \to \ldots
$$
。线性代数中的秩-零度定理表明
$$
0 = \dim H^0(X, \mathcal{F}_1) - \dim H^0(X, \mathcal{F}_2)
+ \dim H^0(X, \mathcal{F}_3) - \dim H^1(X, \mathcal{F}_1) + \ldots
$$
这立即蕴含本引理。
\end{proof}

\begin{lemma}
\label{spaces-over-fields-lemma-euler-characteristic-morphism}
设 $k$ 为域，$f : Y \to X$ 为 $k$ 上固有代数空间的态射，$\mathcal{G}$
为相干 $\mathcal{O}_Y$-模。则
$$
\chi(Y, \mathcal{G}) = \sum (-1)^i \chi(X, R^if_*\mathcal{G})
$$
\end{lemma}

\begin{proof}
该公式有意义：层 $R^if_*\mathcal{G}$ 相干，且只有有限多个非零，见《空间的上同调》引理
\ref{spaces-cohomology-lemma-proper-pushforward-coherent} 与
\ref{spaces-cohomology-lemma-vanishing-higher-direct-images}。由《位点上的上同调》引理
\ref{sites-cohomology-lemma-Leray}，存在谱序列，其
$$
E_2^{p, q} = H^p(X, R^qf_*\mathcal{G})
$$
收敛到 $H^{p + q}(Y, \mathcal{G})$。由 $X$ 上上同调的有限性，只有有限多个
$E_2^{p, q}$ 非零，且每个 $E_2^{p, q}$ 都是有限维向量空间。因此，对
$r \geq 2$，$E_r^{p, q}$ 亦然，并且
$$
\sum (-1)^{p + q} \dim_k E_r^{p, q}
$$
与 $r$ 无关。对充分大的 $r$，有 $E_r^{p, q} = E_\infty^{p, q}$；而收敛意味着
$H^n(Y, \mathcal{G})$ 上存在滤过，其分次片为满足 $p + 1 = n$ 的
$E_\infty^{p, q}$（这就是谱序列收敛的含义），故得结论。
% P09-ERR-0294
\end{proof}








\section{数值相交}
\label{spaces-over-fields-section-num}

\noindent
本节利用固有代数空间上相干层的欧拉示性数来得到可逆模的数值交数。
主要工具是下述引理。

\begin{lemma}
\label{spaces-over-fields-lemma-numerical-polynomial-from-euler}
设 $k$ 为域，$X$ 为 $k$ 上的固有代数空间，$\mathcal{F}$ 为相干
$\mathcal{O}_X$-模，而
$\mathcal{L}_1, \ldots, \mathcal{L}_r$ 为可逆 $\mathcal{O}_X$-模。映射
$$
(n_1, \ldots, n_r) \longmapsto
\chi(X, \mathcal{F} \otimes
\mathcal{L}_1^{\otimes n_1} \otimes \ldots \otimes
\mathcal{L}_r^{\otimes n_r})
$$
是关于 $n_1, \ldots, n_r$ 的数值多项式，其总次数不超过
$\mathcal{F}$ 的概形论支撑之维数。
\end{lemma}

\begin{proof}
设 $Z \subset X$ 为 $\mathcal{F}$ 的概形论支撑。则对某个相干
$\mathcal{O}_Z$-模 $\mathcal{G}$ 有
$\mathcal{F} = i_*\mathcal{G}$
（《空间的上同调》引理
\ref{spaces-cohomology-lemma-coherent-support-closed}），并且
$$
\chi(X, \mathcal{F} \otimes
\mathcal{L}_1^{\otimes n_1} \otimes \ldots \otimes
\mathcal{L}_r^{\otimes n_r}) =
\chi(Z, \mathcal{G} \otimes
i^*\mathcal{L}_1^{\otimes n_1} \otimes \ldots \otimes
i^*\mathcal{L}_r^{\otimes n_r})
$$
，这是投影公式
（《位点上的上同调》引理 \ref{sites-cohomology-lemma-projection-formula}）
与《空间的上同调》引理
\ref{spaces-cohomology-lemma-relative-affine-cohomology} 的结果。
由于 $|Z| = \text{Supp}(\mathcal{F})$，只需证明
$$
P_\mathcal{F}(n_1, \ldots, n_r) :
(n_1, \ldots, n_r)
\longmapsto
\chi(X, \mathcal{F} \otimes
\mathcal{L}_1^{\otimes n_1} \otimes \ldots \otimes
\mathcal{L}_r^{\otimes n_r})
$$
是关于 $n_1, \ldots, n_r$ 的总次数至多为 $\dim(X)$ 的数值多项式。
若上述结论成立，就称相干 $\mathcal{O}_X$-模 $\mathcal{F}$
具有性质 $\mathcal{P}$。

\medskip\noindent
我们将用逐层剥离法证明该断言；更确切地说，将核验《空间的上同调》引理
\ref{spaces-cohomology-lemma-property-higher-rank-cohomological-variant}
的条件 (1)、(2) 与 (3) 均成立。

\medskip\noindent
核验条件 (1)。设
$$
0 \to \mathcal{F}_1 \to \mathcal{F}_2 \to \mathcal{F}_3 \to 0
$$
为 $X$ 上相干层的短正合列。由引理
\ref{spaces-over-fields-lemma-euler-characteristic-additive}，有
$$
P_{\mathcal{F}_2}(n_1, \ldots, n_r) =
P_{\mathcal{F}_1}(n_1, \ldots, n_r) +
P_{\mathcal{F}_3}(n_1, \ldots, n_r)
$$
% P09-ERR-0295
于是显然，若三个层 $\mathcal{F}_i$ 中有两个具有性质 $\mathcal{P}$，
则第三个也具有该性质。

\medskip\noindent
条件 (2) 由下式立即得到：
$P_{\mathcal{F}^{\oplus m}}(n_1, \ldots, n_r) =
mP_\mathcal{F}(n_1, \ldots, n_r)$.

\medskip\noindent
证明 (3)。设 $i : Z \to X$ 为约化闭子空间，且 $|Z|$ 不可约。
我们须找出 $X$ 上支撑为 $Z$ 的相干模 $\mathcal{G}$，使
$\mathcal{P}$ 对 $\mathcal{G}$ 成立。下面给出两种构造：
一种使用 Chow 引理，另一种使用概形的有限覆盖。

\medskip\noindent
% P09-ERR-0296
用概形的有限覆盖证明 $\mathcal{G}$ 的存在性。
选取有限满态射 $\pi : Z' \to Z$，其中 $Z'$ 为概形，见《空间的极限》命题
\ref{spaces-limits-proposition-there-is-a-scheme-finite-over}。
置 $\mathcal{G} = i_*\pi_*\mathcal{O}_{Z'} = (i \circ \pi)_*\mathcal{O}_{Z'}$。
$Z'$ 在 $k$ 上固有，而且 $\mathcal{G}$ 的支撑为 $Y$（略去细节）。
% P09-ERR-0297
我们有
% P09-ERR-0298
$$
R(\pi \circ i)_*(\mathcal{O}_{Z'}) = \mathcal{G}
\quad\text{且}\quad
R(\pi \circ i)_*(\pi^*i^*(\mathcal{L}_1^{\otimes n_1} \otimes \ldots \otimes
\mathcal{L}_r^{\otimes n_r})
) = \mathcal{G} \otimes \mathcal{L}_1^{\otimes n_1} \otimes \ldots \otimes
\mathcal{L}_r^{\otimes n_r}
$$
。第一式成立是因为 $i \circ \pi$ 为仿射态射
（《空间的上同调》引理
\ref{spaces-cohomology-lemma-affine-vanishing-higher-direct-images}）；
第二式由第一式与投影公式
（《位点上的上同调》引理 \ref{sites-cohomology-lemma-projection-formula}）
得到。应用 Leray
（《位点上的上同调》引理 \ref{sites-cohomology-lemma-apply-Leray}），得到
$$
P_\mathcal{G}(n_1, \ldots, n_r) =
\chi(Z', \pi^*i^*(\mathcal{L}_1^{\otimes n_1} \otimes \ldots \otimes
\mathcal{L}_r^{\otimes n_r}))
$$
。由概形情形
（《簇》引理 \ref{varieties-lemma-numerical-polynomial-from-euler}），
这是关于 $n_1, \ldots, n_r$ 的次数至多为 $\dim(Z')$ 的数值多项式。
由于 $\dim(Z') \leq \dim(Z) \leq \dim(X)$，结论成立。
第一个不等式来自《良态空间》引理
\ref{decent-spaces-lemma-dimension-quasi-finite}。

\medskip\noindent
% P09-ERR-0299
用 Chow 引理证明 $\mathcal{G}$ 的存在性。对态射 $Z \to \Spec(k)$
应用《空间的上同调》引理 \ref{spaces-cohomology-lemma-weak-chow}。
于是得到 $\Spec(k)$ 上的满固有态射 $f : Y \to Z$，其中对某个 $m$，
$Y$ 是 $\mathbf{P}^m_k$ 的闭子概形。用某个闭子概形替换 $Y$ 后，
可设 $Y$ 为整概形且 $f : Y \to Z$ 为 alteration，见引理
\ref{spaces-over-fields-lemma-alteration-contained-in}。
% P09-ERR-0300
以 $\mathcal{O}_Y(n)$ 表示 $\mathcal{O}_{\mathbf{P}^m_k}(n)$ 的拉回。
选取 $n > 0$，使得当 $p > 0$ 时
$R^pf_*\mathcal{O}_Y(n) = 0$，见《空间的上同调》引理
\ref{spaces-cohomology-lemma-kill-by-twisting}。
我们断言 $\mathcal{G} = i_*f_*\mathcal{O}_Y(n)$ 具有性质 $\mathcal{P}$。
事实上，由概形情形
（《簇》引理 \ref{varieties-lemma-numerical-polynomial-from-euler}），
映射
$$
(n_1, \ldots, n_r)
\longmapsto
\chi(Y, \mathcal{O}_Y(n) \otimes
f^*i^*(\mathcal{L}_1^{\otimes n_1} \otimes \ldots \otimes
\mathcal{L}_r^{\otimes n_r}))
$$
是关于 $n_1, \ldots, n_r$ 的总次数至多为 $\dim(Y)$ 的数值多项式。
另一方面，由投影公式
（《位点上的上同调》引理 \ref{sites-cohomology-lemma-projection-formula}），
\begin{align*}
i_*Rf_*\left(
\mathcal{O}_Y(n) \otimes
f^*i^*(\mathcal{L}_1^{\otimes n_1} \otimes \ldots \otimes
\mathcal{L}_r^{\otimes n_r})\right)
& =
i_*Rf_*\mathcal{O}_Y(n) \otimes
\mathcal{L}_1^{\otimes n_1} \otimes \ldots \otimes
\mathcal{L}_r^{\otimes n_r} \\
& =
\mathcal{G} \otimes \mathcal{L}_1^{\otimes n_1} \otimes \ldots \otimes
\mathcal{L}_r^{\otimes n_r}
\end{align*}
，其中最后一个等式来自对 $n$ 的选取。由 Leray
（《位点上的上同调》引理 \ref{sites-cohomology-lemma-apply-Leray}），得到
$$
\chi(Y, \mathcal{O}_Y(n) \otimes
f^*i^*(\mathcal{L}_1^{\otimes n_1} \otimes \ldots \otimes
\mathcal{L}_r^{\otimes n_r})) =
P_\mathcal{G}(n_1, \ldots, n_r)
$$
。由于 $\dim(Y) \leq \dim(Z) \leq \dim(X)$，结论成立。
第一个不等式来自《空间的态射》引理
\ref{spaces-morphisms-lemma-alteration-dimension-general}
以及 $Y \to Z$ 为 alteration 这一事实（因而它在一般点处诱导的剩余域扩张有限）。
% P09-ERR-0301
\end{proof}

\noindent
下述引理大致表明，最高次系数只依赖于相干模在其支撑的一般点处的长度。
% P09-ERR-0302

\begin{lemma}
\label{spaces-over-fields-lemma-numerical-polynomial-leading-term}
设 $k$ 为域，$X$ 为 $k$ 上的固有代数空间，$\mathcal{F}$ 为相干
$\mathcal{O}_X$-模，而
$\mathcal{L}_1, \ldots, \mathcal{L}_r$ 为可逆 $\mathcal{O}_X$-模。
令 $d = \dim(\text{Supp}(\mathcal{F}))$。
设 $Z_i \subset X$ 是 $\text{Supp}(\mathcal{F})$ 的所有 $d$ 维不可约分支。
取 $Z_i$ 的几何一般点 $\overline{x}_i$，并置
$m_i = \text{length}_{\mathcal{O}_{X, \overline{x}_i}}
(\mathcal{F}_{\overline{x}_i})$.
则
$$
\chi(X, \mathcal{F} \otimes \mathcal{L}_1^{\otimes n_1} \otimes \ldots \otimes
\mathcal{L}_r^{\otimes n_r}) -
\sum\nolimits_i
m_i\ \chi(Z_i, \mathcal{L}_1^{\otimes n_1} \otimes \ldots \otimes
\mathcal{L}_r^{\otimes n_r}|_{Z_i})
$$
是关于 $n_1, \ldots, n_r$ 的总次数 $< d$ 的数值多项式。
\end{lemma}

\begin{proof}
先证明一个稍弱的断言。设 $\dim(X) = N$，并令 $X_i \subset X$
为所有 $N$ 维不可约分支。取 $X_i$ 的几何一般点 $\overline{x}_i$。
平展局部环 $\mathcal{O}_{X, \overline{x}_i}$ 是 $0$ 维 Noether 环，
故对每个相干 $\mathcal{O}_X$-模 $\mathcal{F}$，长度
$$
m_i(\mathcal{F}) = \text{length}_{\mathcal{O}_{X, \overline{x}_i}}
(\mathcal{F}_{\overline{x}_i})
$$
是整数且 $\geq 0$。我们断言
% P09-ERR-0303
$$
E(\mathcal{F}) =
\chi(X, \mathcal{F} \otimes \mathcal{L}_1^{\otimes n_1} \otimes \ldots \otimes
\mathcal{L}_r^{\otimes n_r}) -
\sum\nolimits_i
m_i(\mathcal{F})\ \chi(Z_i, \mathcal{L}_1^{\otimes n_1} \otimes \ldots \otimes
\mathcal{L}_r^{\otimes n_r}|_{Z_i})
$$
是关于 $n_1, \ldots, n_r$ 的总次数 $< N$ 的数值多项式。
我们用《空间的上同调》引理
\ref{spaces-cohomology-lemma-property-higher-rank-cohomological-variant}
证明这一点。对任意短正合列 $0 \to \mathcal{F}' \to \mathcal{F} \to
\mathcal{F}'' \to 0$，有
$E(\mathcal{F}) = E(\mathcal{F}') + E(\mathcal{F}'')$.
这来自欧拉示性数的可加性
（引理 \ref{spaces-over-fields-lemma-euler-characteristic-additive}）
与长度的可加性
（《代数》引理 \ref{algebra-lemma-length-additive}）。
这立即蕴含《空间的上同调》引理
\ref{spaces-cohomology-lemma-property-higher-rank-cohomological-variant}
的性质 (1) 与 (2)。最后，对任意不可约约化闭子空间
$Z \subset X$ 取 $\mathcal{G} = \mathcal{O}_Z$，即可验证性质 (3)。
% P09-ERR-0304
事实上，若对某个 $i_0$ 有 $Z = Z_{i_0}$，则
$m_i(\mathcal{G}) = \delta_{i_0i}$，从而 $E(\mathcal{G}) = 0$。
若对所有 $i$ 都有 $Z \not = Z_i$，则对所有 $i$ 都有
$m_i(\mathcal{G}) = 0$，且 $\dim(Z) < N$；此时结论来自引理
\ref{spaces-over-fields-lemma-numerical-polynomial-from-euler}。

\medskip\noindent
现在证明引理所述断言。
设 $Z \subset X$ 为 $\mathcal{F}$ 的概形论支撑。则对某个相干
$\mathcal{O}_Z$-模 $\mathcal{G}$ 有
$\mathcal{F} = i_*\mathcal{G}$
（《空间的上同调》引理
\ref{spaces-cohomology-lemma-coherent-support-closed}），并且
$$
\chi(X, \mathcal{F} \otimes
\mathcal{L}_1^{\otimes n_1} \otimes \ldots \otimes
\mathcal{L}_r^{\otimes n_r}) =
\chi(Z, \mathcal{G} \otimes
i^*\mathcal{L}_1^{\otimes n_1} \otimes \ldots \otimes
i^*\mathcal{L}_r^{\otimes n_r})
$$
，这是投影公式
（《位点上的上同调》引理 \ref{sites-cohomology-lemma-projection-formula}）
与《空间的上同调》引理
\ref{spaces-cohomology-lemma-relative-affine-cohomology} 的结果。
由于 $|Z| = \text{Supp}(\mathcal{F})$，对所有 $i$ 都有
$Z_i \subset Z$，且这些正是 $Z$ 的所有 $d$ 维不可约分支。
我们可以并且确实把 $\overline{x}_i$ 看作 $Z$ 的几何点。映射
$i^\sharp : \mathcal{O}_X \to i_*\mathcal{O}_Z$ 给出满射
$$
\mathcal{O}_{X, \overline{x}_i} \to \mathcal{O}_{Z, \overline{x}_i}
$$
。通过此映射，并由 $\mathcal{F} = i_*\mathcal{G}$，得到模同构
$\mathcal{G}_{\overline{x}_i} = \mathcal{F}_{\overline{x}_i}$。
因此
$$
m_i = \text{length}_{\mathcal{O}_{X, \overline{x}_i}}
(\mathcal{F}_{\overline{x}_i}) =
\text{length}_{\mathcal{O}_{Z, \overline{x}_i}}
(\mathcal{G}_{\overline{x}_i})
$$
。于是引理中的表达式等于
$$
\chi(Z, \mathcal{G} \otimes \mathcal{L}_1^{\otimes n_1} \otimes \ldots \otimes
\mathcal{L}_r^{\otimes n_r}) -
\sum\nolimits_i
m_i\ \chi(Z_i, \mathcal{L}_1^{\otimes n_1} \otimes \ldots \otimes
\mathcal{L}_r^{\otimes n_r}|_{Z_i})
$$
；结论来自第一段的讨论（在那里以 $Z$ 代替 $X$）。
% P09-ERR-0305
\end{proof}

\begin{definition}
\label{spaces-over-fields-definition-intersection-number}
设 $k$ 为域，$X$ 为 $k$ 上的固有代数空间，且
$i : Z \to X$ 为 $d$ 维闭子空间。设
$\mathcal{L}_1, \ldots, \mathcal{L}_d$ 为可逆
$\mathcal{O}_X$-模。定义{\it 交数}
$(\mathcal{L}_1 \cdots \mathcal{L}_d \cdot Z)$
为下述数值多项式中 $n_1 \ldots n_d$ 的系数：
$$
\chi(X, i_*\mathcal{O}_Z \otimes \mathcal{L}_1^{\otimes n_1} \otimes
\ldots \otimes \mathcal{L}_d^{\otimes n_d}) =
\chi(Z, \mathcal{L}_1^{\otimes n_1} \otimes
\ldots \otimes \mathcal{L}_d^{\otimes n_d}|_Z)
$$
。在特殊情形
$\mathcal{L}_1 = \ldots = \mathcal{L}_d = \mathcal{L}$ 下，
记作 $(\mathcal{L}^d \cdot Z)$。
\end{definition}

\noindent
定义中的等式来自投影公式
（《上同调》第 \ref{cohomology-section-projection-formula} 节）
与《概形的上同调》引理
\ref{coherent-lemma-relative-affine-cohomology}。
% P09-ERR-0306
下面证明关于这些交数的几个引理。

\begin{lemma}
\label{spaces-over-fields-lemma-intersection-number-integer}
在定义 \ref{spaces-over-fields-definition-intersection-number} 的情形下，交数
$(\mathcal{L}_1 \cdots \mathcal{L}_d \cdot Z)$
是整数。
\end{lemma}

\begin{proof}
任一关于 $n_1, \ldots, n_d$ 的 $e$ 次数值多项式，都可唯一地写成函数
${n_1 \choose k_1}{n_2 \choose k_2} \ldots {n_d \choose k_d}$ 的
$\mathbf{Z}$-线性组合，其中 $k_1 + \ldots + k_d \leq e$。
取 $e = d$ 即得。细节留作练习。
\end{proof}

\begin{lemma}
\label{spaces-over-fields-lemma-intersection-number-additive}
在定义 \ref{spaces-over-fields-definition-intersection-number} 的情形下，交数
$(\mathcal{L}_1 \cdots \mathcal{L}_d \cdot Z)$
是可加的：若 $\mathcal{L}_i = \mathcal{L}_i' \otimes \mathcal{L}_i''$，
则
$$
(\mathcal{L}_1 \cdots \mathcal{L}_i \cdots \mathcal{L}_d \cdot Z) =
(\mathcal{L}_1 \cdots \mathcal{L}_i' \cdots \mathcal{L}_d \cdot Z) +
(\mathcal{L}_1 \cdots \mathcal{L}_i'' \cdots \mathcal{L}_d \cdot Z)
$$
\end{lemma}

\begin{proof}
事实上，由引理 \ref{spaces-over-fields-lemma-numerical-polynomial-from-euler}，函数
$$
(n_1, \ldots, n_{i - 1}, n_i', n_i'', n_{i + 1}, \ldots, n_d)
\mapsto
\chi(Z, \mathcal{L}_1^{\otimes n_1} \otimes
\ldots \otimes (\mathcal{L}_i')^{\otimes n_i'} \otimes
(\mathcal{L}_i'')^{\otimes n_i''} \otimes \ldots \otimes
\mathcal{L}_d^{\otimes n_d}|_Z)
$$
是关于 $d + 1$ 个变量的总次数至多为 $d$ 的数值多项式。
\end{proof}

\begin{lemma}
\label{spaces-over-fields-lemma-intersection-number-in-terms-of-components}
在定义 \ref{spaces-over-fields-definition-intersection-number} 的情形下，设
$Z_i \subset Z$ 为所有 $d$ 维不可约分支。令
$m_i = \text{length}_{\mathcal{O}_{X, \overline{x}_i}}
(\mathcal{O}_{Z, \overline{x}_i})$
，其中 $\overline{x}_i$ 是 $Z_i$ 的几何一般点。则
$$
(\mathcal{L}_1 \cdots \mathcal{L}_d \cdot Z) =
\sum m_i(\mathcal{L}_1 \cdots \mathcal{L}_d \cdot Z_i)
$$
\end{lemma}

\begin{proof}
立即由引理 \ref{spaces-over-fields-lemma-numerical-polynomial-leading-term} 与定义得到。
\end{proof}

\begin{lemma}
\label{spaces-over-fields-lemma-intersection-number-and-pullback}
设 $k$ 为域，$f : Y \to X$ 为 $k$ 上固有代数空间的态射。
设 $Z \subset Y$ 为 $d$ 维整闭子空间，且
$\mathcal{L}_1, \ldots, \mathcal{L}_d$ 为可逆 $\mathcal{O}_X$-模。则
$$
(f^*\mathcal{L}_1 \cdots f^*\mathcal{L}_d \cdot Z) =
\deg(f|_Z : Z \to f(Z)) (\mathcal{L}_1 \cdots \mathcal{L}_d \cdot f(Z))
$$
，其中 $\deg(Z \to f(Z))$ 如定义 \ref{spaces-over-fields-definition-degree} 所述；
若 $\dim(f(Z)) < d$，则将其取为 $0$。
\end{lemma}

\begin{proof}
在该陈述中，$f(Z) \subset X$ 是 $f$ 的概形论像，也是在闭子集
$f(|Z|) \subset X$ 上赋予约化诱导代数空间结构所得之空间，见《空间的态射》引理
\ref{spaces-morphisms-lemma-scheme-theoretic-image-reduced}。
% P09-ERR-0307
于是 $Z$ 与 $f(Z)$ 都是 $k$ 上约化且固有（从而良态）的代数空间，
因而为整代数空间
（定义 \ref{spaces-over-fields-definition-integral-algebraic-space}）。
% P09-ERR-0308
左端由下述函数中 $n_1 \ldots n_d$ 的系数计算：
$$
\chi(Y, \mathcal{O}_Z \otimes f^*\mathcal{L}_1^{\otimes n_1} \otimes
\ldots \otimes f^*\mathcal{L}_d^{\otimes n_d}) =
\sum (-1)^i
\chi(X, R^if_*\mathcal{O}_Z \otimes
\mathcal{L}_1^{\otimes n_1} \otimes \ldots \otimes
\mathcal{L}_d^{\otimes n_d})
$$
该等式来自引理 \ref{spaces-over-fields-lemma-euler-characteristic-morphism} 与投影公式
（《上同调》引理 \ref{cohomology-lemma-projection-formula}）。
若 $f(Z)$ 的维数 $< d$，则由引理
\ref{spaces-over-fields-lemma-numerical-polynomial-from-euler}，右端是总次数 $< d$ 的多项式，
故结论成立。以下设 $\dim(f(Z)) = d$。由维数理论
（引理 \ref{spaces-over-fields-lemma-alteration-dimension}），引理
\ref{spaces-over-fields-lemma-finite-degree} 的等价条件 (1) -- (5) 成立。
因此 $\deg(Z \to f(Z))$ 定义良好。
再次应用引理 \ref{spaces-over-fields-lemma-finite-degree}，可知 $f : Z \to f(Z)$
在 $f(Z)$ 的某个非空开集 $V$ 上有限；必要时缩小 $V$ 后，可设 $V$
为概形。设 $\xi \in V$ 为一般点。于是
% P09-ERR-0309
$\deg(f : Z \to f(Z))$ 是 $f_*\mathcal{O}_Z$ 在 $\xi$ 处的茎作为
$\mathcal{O}_{X, \xi}$-模的长度，
% P09-ERR-0310
并且当 $i > 0$ 时，$R^if_*\mathcal{O}_X$ 在 $\xi$ 处的茎为零
% P09-ERR-0311
（例如由《空间的上同调》引理
\ref{spaces-cohomology-lemma-finite-higher-direct-image-zero}）。
因此，当 $i > 0$ 时，各项 $\chi(X, R^if_*\mathcal{O}_Z \otimes
\mathcal{L}_1^{\otimes n_1} \otimes \ldots \otimes
\mathcal{L}_d^{\otimes n_d})$ 的总次数均 $< d$，并且
$$
\chi(X, f_*\mathcal{O}_Z \otimes
\mathcal{L}_1^{\otimes n_1} \otimes \ldots \otimes
\mathcal{L}_d^{\otimes n_d})
=
\deg(f : Z \to f(Z)) \chi(f(Z),
\mathcal{L}_1^{\otimes n_1} \otimes \ldots \otimes
\mathcal{L}_d^{\otimes n_d}|_{f(Z)})
$$
由引理 \ref{spaces-over-fields-lemma-numerical-polynomial-leading-term}，上式在模去一个
总次数 $< d$ 的多项式后成立。所求结论随即得到。
\end{proof}

\begin{lemma}
\label{spaces-over-fields-lemma-numerical-intersection-effective-Cartier-divisor}
设 $k$ 为域，$X$ 为 $k$ 上的固有代数空间，而
$Z \subset X$ 为 $d$ 维闭子空间。设
$\mathcal{L}_1, \ldots, \mathcal{L}_d$
为可逆 $\mathcal{O}_X$-模。假设存在有效 Cartier 除子
$D \subset Z$，使得
$\mathcal{L}_1|_Z \cong \mathcal{O}_Z(D)$。则
$$
(\mathcal{L}_1 \cdots \mathcal{L}_d \cdot Z) =
(\mathcal{L}_2 \cdots \mathcal{L}_d \cdot D)
$$
\end{lemma}

\begin{proof}
可以用 $Z$ 替换 $X$，并用 $\mathcal{L}_i|_Z$ 替换 $\mathcal{L}_i$。
因此可设 $X = Z$ 且 $\mathcal{L}_1 = \mathcal{O}_X(D)$。
此时 $\mathcal{L}_1^{-1}$ 是 $D$ 的理想层，故可考虑短正合列
$$
0 \to \mathcal{L}_1^{\otimes -1} \to \mathcal{O}_X \to \mathcal{O}_D \to 0
$$
置
$P(n_1, \ldots, n_d) =
\chi(X, \mathcal{L}_1^{\otimes n_1} \otimes \ldots \otimes
\mathcal{L}_d^{\otimes n_d})$
并置
$Q(n_1, \ldots, n_d) =
\chi(D, \mathcal{L}_1^{\otimes n_1} \otimes \ldots \otimes
\mathcal{L}_d^{\otimes n_d}|_D)$.
由可加性
（引理 \ref{spaces-over-fields-lemma-euler-characteristic-additive}），得到
$$
P(n_1, \ldots, n_d) - P(n_1 - 1, n_2, \ldots, n_d) =
Q(n_1, \ldots, n_d)
$$
。由于 $P$ 的总次数至多为 $d$，可知 $P$ 中 $n_1 \ldots n_d$
的系数等于 $Q$ 中 $n_2 \ldots n_d$ 的系数。
\end{proof}
% Generated mechanically from the frozen spaces-topologies.tex target.
% Do not edit this generated file; edit the bound locale source instead.

% OK, start here.
%

\setcounter{chapter}{72}
\chapter{代数空间上的拓扑}


\phantomsection
\label{spaces-topologies-section-phantom}





\section{引言}
\label{spaces-topologies-section-introduction}

\noindent
本章介绍代数空间范畴上的若干拓扑。可与 \cite{SGA1}、
\cite{Ner}、\cite{LM-B} 及 \cite{Kn} 中的材料比较。
在此之前需要指出，可以选取许多不同的位点（定义见《位点》定义
\ref{sites-definition-site}），而它们在所依范畴上给出相同的层概念。
因此，我们的选择可能与参考文献中的选择略有不同，但最终得到相同的上同调群等。




\section{一般构造过程}
\label{spaces-topologies-section-procedure}

\noindent
本节说明构造下文所用位点的一般过程。读者若未读过《拓扑》第
\ref{topologies-section-procedure} 节，则这段讨论几乎无法理解。

\medskip\noindent
设 $S$ 为基概形。从 $S$ 以及任意一个希望包含的 $S$-概形集合出发，
按《集合》引理 \ref{sets-lemma-construct-category} 构造任一范畴
$\Sch_\alpha$。再从范畴 $\Sch_\alpha$ 与 fppf 覆盖类出发，
按《集合》引理 \ref{sets-lemma-coverings-site} 选取
$\Sch_\alpha$ 上任一覆盖集合 $\text{Cov}_{fppf}$。
以 $\Sch_{fppf}$ 表示所得大 fppf 位点，以
$(\Sch/S)_{fppf}$ 表示相应的 $S$ 的大 fppf 位点。
（以上完全依照《拓扑》第 \ref{topologies-section-fppf} 节的规定。）

\medskip\noindent
作出上述选择后，$S$ 上代数空间的范畴具有一个同构类集合。
这是因为任一 $S$ 上代数空间都可写成 $U/R$，其中
$j : R \to U \times_S U$ 是某个平展等价关系，且
$U, R \in \Ob((\Sch/S)_{fppf})$，见《空间》引理
\ref{spaces-lemma-space-presentation}。
因此，可在 $S$ 上代数空间的范畴中找出一个满子范畴
$\textit{Spaces}/S$；它的对象构成集合，并且每个代数空间都同构于
$\textit{Spaces}/S$ 的某个对象。固定这样一个范畴。

\medskip\noindent
在以下各节中，给定拓扑 $\tau$ 后，大位点
$(\textit{Spaces}/S)_\tau$（相应地，$S$ 上代数空间 $X$ 的大位点
$(\textit{Spaces}/X)_\tau$）的所依范畴是
$\textit{Spaces}/S$（相应地，是 $\textit{Spaces}/S$ 的子范畴
$\textit{Spaces}/X$，见《范畴》例
\ref{categories-example-category-over-X}）。
按通常方式把它变成位点：定义一类 $\tau$-覆盖，并应用《集合》引理
\ref{sets-lemma-coverings-site} 选取一个足够大的覆盖集合来定义该拓扑。

\medskip\noindent
需要指出，{\it 代数空间 $X$ 的小平展位点 $X_\etale$} 已在
《空间的性质》定义
\ref{spaces-properties-definition-etale-site} 中定义。
它的对象是在 $X$ 上平展的概形；由代数空间的定义，这样的对象很多。
% P09-ERR-0312
不过，从本章的视角看（比较《拓扑》定义
\ref{topologies-definition-big-small-etale}），更自然的是
《空间的性质》定义
\ref{spaces-properties-definition-spaces-etale-site} 中的位点
$X_{spaces, \etale}$。这两个位点定义同一拓扑斯，见《空间的性质》引理
\ref{spaces-properties-lemma-compare-etale-sites}。
本章不再重新定义它们，而直接使用。






\section{Zariski 拓扑}
\label{spaces-topologies-section-zariski}

\noindent
在《空间》第 \ref{spaces-section-Zariski} 节中，我们介绍了用开子空间
覆盖代数空间的 Zariski 覆盖概念。将开子空间替换为开浸入，得到如下相应概念。

\begin{definition}
\label{spaces-topologies-definition-zariski-covering}
设 $S$ 为概形，$X$ 为 $S$ 上的代数空间。
{\it $X$ 的 Zariski 覆盖}是 $S$ 上代数空间的一个态射族
$\{f_i : X_i \to X\}_{i \in I}$，其中每个 $f_i$ 都是开浸入，并且
$$
|X| = \bigcup\nolimits_{i \in I} |f_i|(|X_i|),
$$
即这些态射联合满。
\end{definition}

\noindent
尽管 Zariski 覆盖偶尔有用，但代数空间范畴上相应的拓扑确实太粗，
因而并不特别有用。尽管如此，它仍定义一个位点。

\begin{lemma}
\label{spaces-topologies-lemma-zariski}
设 $S$ 为概形，$X$ 为 $S$ 上的代数空间。
\begin{enumerate}
\item 若 $X' \to X$ 是同构，则 $\{X' \to X\}$ 是 $X$ 的 Zariski 覆盖。
\item 若 $\{X_i \to X\}_{i\in I}$ 是 Zariski 覆盖，并且对每个
$i$，$\{X_{ij} \to X_i\}_{j\in J_i}$ 是 Zariski 覆盖，则
$\{X_{ij} \to X\}_{i \in I, j\in J_i}$ 是 Zariski 覆盖。
\item 若 $\{X_i \to X\}_{i\in I}$ 是 Zariski 覆盖，而
$X' \to X$ 是代数空间的态射，则
$\{X' \times_X X_i \to X'\}_{i\in I}$ 是 Zariski 覆盖。
\end{enumerate}
\end{lemma}

\begin{proof}
略。
\end{proof}










\section{平展拓扑}
\label{spaces-topologies-section-etale}

\noindent
本节讨论代数空间的平展覆盖概念，并定义代数空间的大平展位点。
请与《拓扑》第 \ref{topologies-section-etale} 节比较。

\begin{definition}
\label{spaces-topologies-definition-etale-covering}
设 $S$ 为概形，$X$ 为 $S$ 上的代数空间。
{\it $X$ 的平展覆盖}是 $S$ 上代数空间的一个态射族
$\{f_i : X_i \to X\}_{i \in I}$，其中每个 $f_i$ 都平展，并且
$$
|X| = \bigcup\nolimits_{i \in I} |f_i|(|X_i|),
$$
即这些态射联合满。
\end{definition}

\noindent
这恰好就是《拓扑》定义 \ref{topologies-definition-etale-covering}。
特别地，若 $X$ 及所有 $X_i$ 都是概形，便得到通常的概形平展覆盖概念。

\begin{lemma}
\label{spaces-topologies-lemma-zariski-etale}
任一 Zariski 覆盖都是平展覆盖。
\end{lemma}

\begin{proof}
由定义以及开浸入是平展态射这一事实即可得出（由于浸入可表，
该事实由《态射》引理 \ref{morphisms-lemma-open-immersion-etale}
经《空间》引理
\ref{spaces-lemma-representable-transformations-property-implication} 得到）。
\end{proof}

\begin{lemma}
\label{spaces-topologies-lemma-etale}
设 $S$ 为概形，$X$ 为 $S$ 上的代数空间。
\begin{enumerate}
\item 若 $X' \to X$ 是同构，则 $\{X' \to X\}$ 是 $X$ 的平展覆盖。
\item 若 $\{X_i \to X\}_{i\in I}$ 是平展覆盖，并且对每个
$i$，$\{X_{ij} \to X_i\}_{j\in J_i}$ 是平展覆盖，则
$\{X_{ij} \to X\}_{i \in I, j\in J_i}$ 是平展覆盖。
\item 若 $\{X_i \to X\}_{i\in I}$ 是平展覆盖，而
$X' \to X$ 是代数空间的态射，则
$\{X' \times_X X_i \to X'\}_{i\in I}$ 是平展覆盖。
\end{enumerate}
\end{lemma}

\begin{proof}
略。
\end{proof}

\noindent
下一个引理说明位点
$(\textit{Spaces}/X)_\etale$ 与 $(\textit{Spaces}/X)_{smooth}$
具有相同的层范畴。

\begin{lemma}
\label{spaces-topologies-lemma-etale-dominates-smooth}
设 $S$ 为概形，$X$ 为 $S$ 上的代数空间。
设 $\{X_i \to X\}_{i \in I}$ 是 $X$ 的光滑覆盖。
则存在 $X$ 的平展覆盖 $\{U_j \to X\}_{j \in J}$，
它加细 $\{X_i \to X\}_{i \in I}$。
\end{lemma}

\begin{proof}
先取概形 $U$ 及满平展态射 $U \to X$。
对每个 $i$，取概形 $W_i$ 及满平展态射 $W_i \to X_i$。
于是 $\{W_i \to X\}_{i \in I}$ 是加细
$\{X_i \to X\}_{i \in I}$ 的光滑覆盖。因此
$\{W_i \times_X U \to U\}_{i \in I}$ 是概形的光滑覆盖。
由《态射续论》引理 \ref{more-morphisms-lemma-etale-dominates-smooth}，
可取加细 $\{W_i \times_X U \to U\}$ 的平展覆盖
$\{U_j \to U\}$。于是 $\{U_j \to X\}_{j \in J}$ 是加细
$\{X_i \to X\}_{i \in I}$ 的平展覆盖。
\end{proof}

\begin{definition}
\label{spaces-topologies-definition-big-etale-site}
设 $S$ 为概形。大平展位点 {\it $(\textit{Spaces}/S)_\etale$}
是按如下方式构造的任一位点：
\begin{enumerate}
\item 按《拓扑》第 \ref{topologies-section-etale} 节选取一个
大平展位点 $(\Sch/S)_\etale$。
\item 取 $S$ 上代数空间的范畴 $\textit{Spaces}/S$ 为所依范畴
（关于它为何是集合，见第 \ref{spaces-topologies-section-procedure} 节的讨论）。
\item 按《集合》引理 \ref{sets-lemma-coverings-site}，
从范畴 $\textit{Spaces}/S$ 及定义
\ref{spaces-topologies-definition-etale-covering} 的平展覆盖类出发，选取任一覆盖集合。
\end{enumerate}
\end{definition}

\noindent
定义完毕后，可通过局部化得到代数空间的平展位点。

\begin{definition}
\label{spaces-topologies-definition-big-small-etale}
设 $S$ 为概形，$(\textit{Spaces}/S)_\etale$ 如定义
\ref{spaces-topologies-definition-big-etale-site}。
设 $X$ 为 $S$ 上的代数空间，即 $(\textit{Spaces}/S)_\etale$
的对象。则 $X$ 的大平展位点
{\it $(\textit{Spaces}/X)_\etale$} 是《位点》第
\ref{sites-section-localize} 节所引入的位点
$(\textit{Spaces}/S)_\etale$ 在 $X$ 处的局部化。
\end{definition}

\noindent
回顾定义中的 $S$ 上代数空间 $X$，我们已经定义了小平展位点
$X_{spaces, \etale}$ 与 $X_\etale$，见《空间的性质》第
\ref{spaces-properties-section-etale-site} 节。
我们借助包含函子 $X_\etale \subset X_{spaces, \etale}$
（《空间的性质》引理
\ref{spaces-properties-lemma-compare-etale-sites}）
默认识别相应的拓扑斯，并称之为 $X$ 的小平展拓扑斯。
下面建立这些位点所关联的拓扑斯之间的一些关系。

\begin{lemma}
\label{spaces-topologies-lemma-put-in-T-etale}
设 $S$ 为概形，$f : Y \to X$ 为
$(\textit{Spaces}/S)_\etale$ 中的态射。包含函子
$Y_{spaces, \etale} \to (\textit{Spaces}/X)_\etale$
是余连续的，并诱导拓扑斯态射
$$
i_f :
\Sh(Y_\etale)
\longrightarrow
\Sh((\textit{Spaces}/X)_\etale)
$$
对 $(\textit{Spaces}/X)_\etale$ 上的层 $\mathcal{G}$，有公式
$(i_f^{-1}\mathcal{G})(U/Y) = \mathcal{G}(U/X)$。
函子 $i_f^{-1}$ 还具有左伴随 $i_{f, !}$，后者与纤维积和等化子可交换。
\end{lemma}

\begin{proof}
以 $u : Y_{spaces, \etale} \to (\textit{Spaces}/X)_\etale$ 表示该函子。
换言之，给定对应于 $Y_{spaces, \etale}$ 中一个对象的平展态射
$j : U \to Y$，置
% P09-ERR-0313
$u(U \to T) = (f \circ j : U \to S)$.
范畴 $Y_{spaces, \etale}$ 有纤维积和等化子，且 $u$ 与它们可交换。
% P09-ERR-0314
显然 $u$ 余连续。函子 $u$ 也连续，因为 $u$ 将覆盖变为覆盖并与
纤维积可交换。因此引理由《位点》引理
\ref{sites-lemma-when-shriek} 和
\ref{sites-lemma-preserve-equalizers} 得出。
\end{proof}

\begin{lemma}
\label{spaces-topologies-lemma-at-the-bottom-etale}
设 $S$ 为概形，$X$ 为 $(\textit{Spaces}/S)_\etale$ 的对象。
包含函子 $X_{spaces, \etale} \to (\textit{Spaces}/X)_\etale$
满足《位点》引理 \ref{sites-lemma-bigger-site} 的假设，因而诱导位点态射
$$
\pi_X :
(\textit{Spaces}/X)_\etale
\longrightarrow
X_{spaces, \etale}
$$
以及拓扑斯态射
$$
i_X :
\Sh(X_\etale)
\longrightarrow
\Sh((\textit{Spaces}/X)_\etale)
$$
使得 $\pi_X \circ i_X = \text{id}$。此外，
$i_X = i_{\text{id}_X}$，其中 $i_{\text{id}_X}$ 如引理
\ref{spaces-topologies-lemma-put-in-T-etale} 所述。特别地，函子
$i_X^{-1} = \pi_{X, *}$ 由规则
$i_X^{-1}(\mathcal{G})(U/X) = \mathcal{G}(U/X)$.
描述。
\end{lemma}

\begin{proof}
在此情形，函子
$u : X_{spaces, \etale} \to (\textit{Spaces}/X)_\etale$
除具有上面引理 \ref{spaces-topologies-lemma-put-in-T-etale} 的证明中所见性质外，
还是全忠实的，并将终对象变为终对象。
本引理由《位点》引理 \ref{sites-lemma-bigger-site} 得出。
\end{proof}

\begin{definition}
\label{spaces-topologies-definition-restriction-small-etale}
在引理 \ref{spaces-topologies-lemma-at-the-bottom-etale} 的情形中，函子
$i_X^{-1} = \pi_{X, *}$ 常称为{\it 限制到小平展位点}。
% P09-ERR-0315
对大平展位点上的层 $\mathcal{F}$，常以
$\mathcal{F}|_{X_\etale}$ 表示这一限制。
\end{definition}

\noindent
采用这一记号，对大位点上的层 $\mathcal{F}$ 及小位点上的层
$\mathcal{G}$，有
\begin{align*}
\Mor_{\Sh(X_\etale)}(
\mathcal{F}|_{X_\etale},
\mathcal{G})
& =
\Mor_{\Sh((\textit{Spaces}/X)_\etale)}(
\mathcal{F},
i_{X, *}\mathcal{G}) \\
\Mor_{\Sh(X_\etale)}(
\mathcal{G},
\mathcal{F}|_{X_\etale})
& =
\Mor_{\Sh((\textit{Spaces}/X)_\etale)}(
\pi_X^{-1}\mathcal{G},
\mathcal{F})
\end{align*}
此外，有 $(i_{X, *}\mathcal{G})|_{X_\etale} = \mathcal{G}$，
也有 $(\pi_X^{-1}\mathcal{G})|_{X_\etale} = \mathcal{G}$。

\begin{lemma}
\label{spaces-topologies-lemma-morphism-big-etale}
设 $S$ 为概形，$f : Y \to X$ 为
$(\textit{Spaces}/S)_\etale$ 中的态射。函子
$$
u :
(\textit{Spaces}/Y)_\etale
\longrightarrow
(\textit{Spaces}/X)_\etale,
\quad
V/Y \longmapsto V/X
$$
是余连续的，并有连续右伴随
$$
v :
(\textit{Spaces}/X)_\etale
\longrightarrow
(\textit{Spaces}/Y)_\etale,
\quad
(U \to X) \longmapsto (U \times_X Y \to Y).
$$
它们诱导同一拓扑斯态射
$$
f_{big} :
\Sh((\textit{Spaces}/Y)_\etale)
\longrightarrow
\Sh((\textit{Spaces}/X)_\etale)
$$
且有 $f_{big}^{-1}(\mathcal{G})(U/Y) = \mathcal{G}(U/X)$，
$f_{big, *}(\mathcal{F})(U/X) = \mathcal{F}(U \times_X Y/Y)$。
此外，$f_{big}^{-1}$ 有左伴随 $f_{big!}$，后者与纤维积和等化子可交换。
\end{lemma}

\begin{proof}
函子 $u$ 余连续、连续，并与纤维积及等化子可交换
（细节略去；比较引理 \ref{spaces-topologies-lemma-put-in-T-etale} 的证明）。
因此可应用《位点》引理 \ref{sites-lemma-when-shriek} 和
\ref{sites-lemma-preserve-equalizers}，得到 $f_{big}^{-1}$
的公式以及 $f_{big!}$ 的存在性。此外，函子 $v$ 是右伴随，
因为给定 $U/Y$ 与 $V/X$，有
$\Mor_X(u(U), V) = \Mor_Y(U, V \times_X Y)$，正合所需。
于是可应用《位点》引理
\ref{sites-lemma-have-functor-other-way} 和
\ref{sites-lemma-have-functor-other-way-morphism}，
得到 $f_{big, *}$ 的公式。
\end{proof}

\begin{lemma}
\label{spaces-topologies-lemma-morphism-big-small-etale}
设 $S$ 为概形，$f : Y \to X$ 为
$(\textit{Spaces}/S)_\etale$ 中的态射。
\begin{enumerate}
\item % P09-ERR-0316
有 $i_f = f_{big} \circ i_T$，其中 $i_f$ 如引理
\ref{spaces-topologies-lemma-put-in-T-etale}，$i_T$ 如引理
\ref{spaces-topologies-lemma-at-the-bottom-etale}。
\item % P09-ERR-0317
函子 $X_{spaces, \etale} \to T_{spaces, \etale}$，
$(U \to X) \mapsto (U \times_X Y \to Y)$ 连续，并诱导位点态射
$$
f_{spaces, \etale} : Y_{spaces, \etale} \longrightarrow X_{spaces, \etale}
$$
相应的小平展拓扑斯态射记作
$$
f_{small} : \Sh(Y_\etale) \to \Sh(X_\etale)
$$
有 $f_{small, *}(\mathcal{F})(U/X) = \mathcal{F}(U \times_X Y/Y)$。
\item 有如下位点态射的交换图
$$
\xymatrix{
Y_{spaces, \etale} \ar[d]_{f_{spaces, \etale}} &
(\textit{Spaces}/Y)_\etale \ar[d]^{f_{big}} \ar[l]^-{\pi_Y}\\
X_{spaces, \etale} &
(\textit{Spaces}/X)_\etale \ar[l]_-{\pi_X}
}
$$
因而作为拓扑斯态射，
$f_{small} \circ \pi_Y = \pi_X \circ f_{big}$。
\item 有 $f_{small} = \pi_X \circ f_{big} \circ i_Y = \pi_X \circ i_f$。
\end{enumerate}
\end{lemma}

\begin{proof}
% P09-ERR-0318
等式 $i_f = f_{big} \circ i_Y$ 由等式
$i_f^{-1} = i_T^{-1} \circ f_{big}^{-1}$ 得出；
后者由上述函子描述即知。因此得到 (1)。

\medskip\noindent
% P09-ERR-0319
在《空间的性质》引理
\ref{spaces-properties-lemma-functoriality-etale-site} 中已经说明，
函子 $u : X_{spaces, \etale} \to Y_{spaces, \etale}$，
$u(U \to X) = (U \times_X Y \to Y)$，给出一个位点态射及相应的
小平展拓扑斯态射。正像的描述是显然的。

\medskip\noindent
(3) 成立，因为 $\pi_X$ 和 $\pi_Y$ 由包含函子给出，而
$f_{spaces, \etale}$ 和 $f_{big}$ 由基变换函子
$U \mapsto U \times_X Y$ 给出。

\medskip\noindent
在 (3) 两端前复合 $i_Y$ 即得 (4)。
\end{proof}

\noindent
在该引理的情形中，采用定义
\ref{spaces-topologies-definition-restriction-small-etale} 的术语，
对 $Y$ 的大平展位点上的层 $\mathcal{F}$，有
$$
(f_{big, *}\mathcal{F})|_{X_\etale} =
f_{small, *}(\mathcal{F}|_{Y_\etale}),
$$
% P09-ERR-0320
该等式由引理中的位点图交换即知，因为限制到 $Y$（相应地 $X$）
的小平展位点由 $\pi_{Y, *}$（相应地 $\pi_{X, *}$）给出。
涉及拉回与限制的类似公式不成立。

\begin{lemma}
\label{spaces-topologies-lemma-composition-etale}
设 $S$ 为概形。给定 $(\textit{Spaces}/S)_\etale$ 中的态射
$f : X \to Y$、$g : Y \to Z$，有
$g_{big} \circ f_{big} = (g \circ f)_{big}$，且
$g_{small} \circ f_{small} = (g \circ f)_{small}$.
\end{lemma}

\begin{proof}
对大位点上的函子，此结论由引理 \ref{spaces-topologies-lemma-morphism-big-etale}
中正像与拉回的简单描述得出。对小位点上的函子，此结论由引理
\ref{spaces-topologies-lemma-morphism-big-small-etale} 中正像函子的描述得出。
\end{proof}

\begin{lemma}
\label{spaces-topologies-lemma-morphism-big-small-cartesian-diagram-etale}
设 $S$ 为概形。考虑 $(\textit{Spaces}/S)_\etale$ 中的笛卡尔图
$$
\xymatrix{
Y' \ar[r]_{g'} \ar[d]_{f'} & Y \ar[d]^f \\
X' \ar[r]^g & X
}
$$
则
$i_g^{-1} \circ f_{big, *} = f'_{small, *} \circ (i_{g'})^{-1}$
且 $g_{big}^{-1} \circ f_{big, *} = f'_{big, *} \circ (g'_{big})^{-1}$。
\end{lemma}

\begin{proof}
由于该图是笛卡儿的，对 $U'/X'$ 有
$U' \times_{X'} Y' = U' \times_X Y$。因此
$i_g^{-1} \circ f_{big, *}$ 与
$f'_{small, *} \circ (i_{g'})^{-1}$ 都把
$(\textit{Spaces}/Y)_\etale$ 上的层 $\mathcal{F}$ 送到
$X'_\etale$ 上的层
$U' \mapsto \mathcal{F}(U' \times_{X'} Y')$
（使用引理 \ref{spaces-topologies-lemma-put-in-T-etale} 和
\ref{spaces-topologies-lemma-morphism-big-etale}）。
第二个等式可用相同方法证明，也可由一般性的《位点》引理
\ref{sites-lemma-localize-morphism} 推出。
\end{proof}

\begin{remark}
\label{spaces-topologies-remark-change-topologies-ringed}
位点 $(\textit{Spaces}/X)_\etale$ 与 $X_{spaces, \etale}$
都带有结构层。对小平展位点，我们已在《空间的性质》第
\ref{spaces-properties-section-structure-sheaf} 节看到这一点。
大平展位点 $(\textit{Spaces}/X)_\etale$ 上的结构层
$\mathcal{O}$ 定义为：给对象 $U$ 指派 $U$ 的结构层的整体截面。
这是有意义的，因为 $U$ 本身毕竟是代数空间，因而具有结构层。
由于 $\mathcal{O}_U$ 是 $U$ 的平展位点上的层，如此定义的预层
$\mathcal{O}$ 满足 $U$ 的覆盖的层条件；也就是说，
$\mathcal{O}$ 是层。
可以把上面定义的态射 $i_f$、$\pi_X$、$i_X$、$f_{small}$ 及
$f_{big}$ 分别提升为带环位点、带环拓扑斯的态射。下面逐一处理。
\begin{enumerate}
\item 在引理 \ref{spaces-topologies-lemma-put-in-T-etale} 中，以 $\mathcal{O}$
表示 $(\textit{Spaces}/X)_\etale$ 上的结构层。
按构造有
$(i_f^{-1}\mathcal{O})(U/Y) = \mathcal{O}_U(U) = \mathcal{O}_Y(U)$。
% P09-ERR-0321
从而有同构 $i_f^\sharp : i_f^{-1}\mathcal{O} \to \mathcal{O}_Y$。
\item 在引理 \ref{spaces-topologies-lemma-at-the-bottom-etale} 中已经指出，
$i_X$ 是 $f = \text{id}_X$ 时 $i_f$ 的特例，故回到情形 (1)。
\item 在引理 \ref{spaces-topologies-lemma-at-the-bottom-etale} 中，态射
$\pi_X$ 满足 $(\pi_{X, *}\mathcal{O})(U) = \mathcal{O}(U) =
\mathcal{O}_X(U)$。因此可用它定义
$\pi_X^\sharp : \mathcal{O}_X \to \pi_{X, *}\mathcal{O}$。
\item 在引理 \ref{spaces-topologies-lemma-morphism-big-small-etale} 中，
《空间的性质》引理
\ref{spaces-properties-lemma-morphism-ringed-topoi}
讨论了把 $f_{small}$ 扩张为带环拓扑斯态射。
\item 在引理 \ref{spaces-topologies-lemma-morphism-big-small-etale} 中，
函子 $f_{big}^{-1}$ 只不过是通过包含函子
$(\textit{Spaces}/Y)_\etale \to (\textit{Spaces}/X)_\etale$
作限制。设 $\mathcal{O}_1$ 为 $(\textit{Spaces}/X)_\etale$
上的结构层，$\mathcal{O}_2$ 为 $(\textit{Spaces}/Y)_\etale$
上的结构层。得到典范同构
$f_{big}^\sharp : f_{big}^{-1}\mathcal{O}_1 \to \mathcal{O}_2$。
\end{enumerate}
此外，在这些定义下，复合也正确相容。略去详细陈述与证明。
\end{remark}









\section{光滑拓扑}
\label{spaces-topologies-section-smooth}

\noindent
本节讨论代数空间的光滑覆盖概念，并定义代数空间的大光滑位点。
请与《拓扑》第 \ref{topologies-section-smooth} 节比较。

\begin{definition}
\label{spaces-topologies-definition-smooth-covering}
设 $S$ 为概形，$X$ 为 $S$ 上的代数空间。
{\it $X$ 的光滑覆盖}是 $S$ 上代数空间的一个态射族
$\{f_i : X_i \to X\}_{i \in I}$，其中每个 $f_i$ 都光滑，并且
$$
|X| = \bigcup\nolimits_{i \in I} |f_i|(|X_i|),
$$
即这些态射联合满。
\end{definition}

\noindent
这恰好就是《拓扑》定义 \ref{topologies-definition-smooth-covering}。
特别地，若 $X$ 及所有 $X_i$ 都是概形，便得到通常的概形光滑覆盖概念。

\begin{lemma}
\label{spaces-topologies-lemma-zariski-etale-smooth}
任一平展覆盖都是光滑覆盖；因此更进一步，任一 Zariski 覆盖都是光滑覆盖。
\end{lemma}

\begin{proof}
这由定义、平展态射是光滑态射这一事实
（《空间的态射》引理 \ref{spaces-morphisms-lemma-etale-smooth}）
及引理 \ref{spaces-topologies-lemma-zariski-etale} 即得。
\end{proof}

\begin{lemma}
\label{spaces-topologies-lemma-smooth}
设 $S$ 为概形，$X$ 为 $S$ 上的代数空间。
\begin{enumerate}
\item 若 $X' \to X$ 是同构，则 $\{X' \to X\}$ 是 $X$ 的光滑覆盖。
\item 若 $\{X_i \to X\}_{i\in I}$ 是光滑覆盖，并且对每个
$i$，$\{X_{ij} \to X_i\}_{j\in J_i}$ 是光滑覆盖，则
$\{X_{ij} \to X\}_{i \in I, j\in J_i}$ 是光滑覆盖。
\item 若 $\{X_i \to X\}_{i\in I}$ 是光滑覆盖，而
$X' \to X$ 是代数空间的态射，则
$\{X' \times_X X_i \to X'\}_{i\in I}$ 是光滑覆盖。
\end{enumerate}
\end{lemma}

\begin{proof}
略。
\end{proof}

\noindent
待续……









\section{Syntomic 拓扑}
\label{spaces-topologies-section-syntomic}

\noindent
本节讨论代数空间的 syntomic 覆盖概念，并定义代数空间的大
syntomic 位点。请与《拓扑》第
\ref{topologies-section-syntomic} 节比较。

\begin{definition}
\label{spaces-topologies-definition-syntomic-covering}
设 $S$ 为概形，$X$ 为 $S$ 上的代数空间。
{\it $X$ 的 syntomic 覆盖}是 $S$ 上代数空间的一个态射族
$\{f_i : X_i \to X\}_{i \in I}$，其中每个 $f_i$ 都是 syntomic
态射，并且
$$
|X| = \bigcup\nolimits_{i \in I} |f_i|(|X_i|),
$$
即这些态射联合满。
\end{definition}

\noindent
这恰好就是《拓扑》定义 \ref{topologies-definition-syntomic-covering}。
特别地，若 $X$ 及所有 $X_i$ 都是概形，便得到通常的概形
syntomic 覆盖概念。

\begin{lemma}
\label{spaces-topologies-lemma-zariski-etale-smooth-syntomic}
任一光滑覆盖都是 syntomic 覆盖；因此更进一步，任一平展覆盖或
Zariski 覆盖都是 syntomic 覆盖。
\end{lemma}

\begin{proof}
这由定义、光滑态射是 syntomic 态射这一事实
（《空间的态射》引理
\ref{spaces-morphisms-lemma-smooth-syntomic}）
及引理 \ref{spaces-topologies-lemma-zariski-etale-smooth} 即得。
\end{proof}

\begin{lemma}
\label{spaces-topologies-lemma-syntomic}
设 $S$ 为概形，$X$ 为 $S$ 上的代数空间。
\begin{enumerate}
\item 若 $X' \to X$ 是同构，则 $\{X' \to X\}$ 是 $X$ 的
syntomic 覆盖。
\item 若 $\{X_i \to X\}_{i\in I}$ 是 syntomic 覆盖，并且对每个
$i$，$\{X_{ij} \to X_i\}_{j\in J_i}$ 是 syntomic 覆盖，则
$\{X_{ij} \to X\}_{i \in I, j\in J_i}$ 是 syntomic 覆盖。
\item 若 $\{X_i \to X\}_{i\in I}$ 是 syntomic 覆盖，而
$X' \to X$ 是代数空间的态射，则
$\{X' \times_X X_i \to X'\}_{i\in I}$ 是 syntomic 覆盖。
\end{enumerate}
\end{lemma}

\begin{proof}
略。
\end{proof}

\noindent
待续……










\section{Fppf 拓扑}
\label{spaces-topologies-section-fppf}

\noindent
本节讨论代数空间的 fppf 覆盖概念，并定义代数空间的大 fppf 位点。
请与《拓扑》第 \ref{topologies-section-fppf} 节比较。

\begin{definition}
\label{spaces-topologies-definition-fppf-covering}
设 $S$ 为概形，$X$ 为 $S$ 上的代数空间。
{\it $X$ 的 fppf 覆盖}是 $S$ 上代数空间的一个态射族
$\{f_i : X_i \to X\}_{i \in I}$，其中每个 $f_i$ 都平坦且局部有限表示，
并且
$$
|X| = \bigcup\nolimits_{i \in I} |f_i|(|X_i|),
$$
即这些态射联合满。
\end{definition}

\noindent
这恰好就是《拓扑》定义 \ref{topologies-definition-fppf-covering}。
特别地，若 $X$ 及所有 $X_i$ 都是概形，便得到通常的概形 fppf 覆盖概念。

\begin{lemma}
\label{spaces-topologies-lemma-zariski-etale-smooth-syntomic-fppf}
任一 syntomic 覆盖都是 fppf 覆盖；因此更进一步，任一光滑、平展或
Zariski 覆盖都是 fppf 覆盖。
\end{lemma}

\begin{proof}
这由定义、syntomic 态射平坦且局部有限表示这一事实
（《空间的态射》引理
\ref{spaces-morphisms-lemma-syntomic-locally-finite-presentation} 和
\ref{spaces-morphisms-lemma-syntomic-flat}）
及引理 \ref{spaces-topologies-lemma-zariski-etale-smooth-syntomic} 即得。
\end{proof}

\begin{lemma}
\label{spaces-topologies-lemma-fppf}
设 $S$ 为概形，$X$ 为 $S$ 上的代数空间。
\begin{enumerate}
\item 若 $X' \to X$ 是同构，则 $\{X' \to X\}$ 是 $X$ 的 fppf 覆盖。
\item 若 $\{X_i \to X\}_{i\in I}$ 是 fppf 覆盖，并且对每个
$i$，$\{X_{ij} \to X_i\}_{j\in J_i}$ 是 fppf 覆盖，则
$\{X_{ij} \to X\}_{i \in I, j\in J_i}$ 是 fppf 覆盖。
\item 若 $\{X_i \to X\}_{i\in I}$ 是 fppf 覆盖，而
$X' \to X$ 是代数空间的态射，则
$\{X' \times_X X_i \to X'\}_{i\in I}$ 是 fppf 覆盖。
\end{enumerate}
\end{lemma}

\begin{proof}
略。
\end{proof}

\begin{lemma}
\label{spaces-topologies-lemma-refine-fppf-schemes}
设 $S$ 为概形，$X$ 为 $S$ 上的代数空间。假设
$\mathcal{U} = \{f_i : X_i \to X\}_{i \in I}$ 是 $X$ 的 fppf 覆盖。
则存在 $\mathcal{U}$ 的加细
$\mathcal{V} = \{g_i : T_i \to X\}$，它是 fppf 覆盖且每个
$T_i$ 都是概形。
\end{lemma}

\begin{proof}
略。提示：对每个 $i$，取概形 $T_i$ 及满平展态射
$T_i \to X_i$，再验证 $\{T_i \to X\}$ 是 fppf 覆盖。
\end{proof}

\begin{lemma}
\label{spaces-topologies-lemma-fppf-covering-surjective}
设 $S$ 为概形。设 $\{f_i : X_i \to X\}_{i \in I}$
为 $S$ 上代数空间的 fppf 覆盖。则层态射
$$
\coprod X_i \longrightarrow X
$$
是满的。
\end{lemma}

\begin{proof}
这由《空间》引理
\ref{spaces-lemma-surjective-flat-locally-finite-presentation} 得出。
若对本引理的含义有疑问，另见《空间》注
\ref{spaces-remark-warning}。
\end{proof}

\begin{definition}
\label{spaces-topologies-definition-big-fppf-site}
设 $S$ 为概形。大 fppf 位点 {\it $(\textit{Spaces}/S)_{fppf}$}
是按如下方式构造的任一位点：
\begin{enumerate}
\item 按《拓扑》第 \ref{topologies-section-fppf} 节选取一个大
fppf 位点 $(\Sch/S)_{fppf}$。
\item 取 $S$ 上代数空间的范畴 $\textit{Spaces}/S$ 为所依范畴
（关于它为何是集合，见第 \ref{spaces-topologies-section-procedure} 节的讨论）。
\item 按《集合》引理 \ref{sets-lemma-coverings-site}，
从范畴 $\textit{Spaces}/S$ 及定义
\ref{spaces-topologies-definition-fppf-covering} 的 fppf 覆盖类出发，选取任一覆盖集合。
\end{enumerate}
\end{definition}

\noindent
定义完毕后，可通过局部化得到代数空间的 fppf 位点。

\begin{definition}
\label{spaces-topologies-definition-big-small-fppf}
设 $S$ 为概形，$(\textit{Spaces}/S)_{fppf}$ 如定义
\ref{spaces-topologies-definition-big-fppf-site}。
设 $X$ 为 $S$ 上的代数空间，即 $(\textit{Spaces}/S)_{fppf}$
的对象。则 $X$ 的大 fppf 位点
{\it $(\textit{Spaces}/X)_{fppf}$} 是《位点》第
\ref{sites-section-localize} 节所引入的位点
$(\textit{Spaces}/S)_{fppf}$ 在 $X$ 处的局部化。
\end{definition}

\noindent
下面建立这些位点所关联的拓扑斯之间的一些关系。

\begin{lemma}
\label{spaces-topologies-lemma-morphism-big-fppf}
设 $S$ 为概形，$f : Y \to X$ 为 $S$ 上代数空间的态射。函子
$$
u : (\textit{Spaces}/Y)_{fppf} \longrightarrow (\textit{Spaces}/X)_{fppf},
\quad
V/Y \longmapsto V/X
$$
是余连续的，并有连续右伴随
% P09-ERR-0322
$$
v : (\textit{Spaces}/X)_{fppf} \longrightarrow (\textit{Spaces}/Y)_{fppf},
\quad
(U \to Y) \longmapsto (U \times_X Y \to Y).
$$
它们诱导同一拓扑斯态射
$$
f_{big} :
\Sh((\textit{Spaces}/Y)_{fppf})
\longrightarrow
\Sh((\textit{Spaces}/X)_{fppf})
$$
且有 $f_{big}^{-1}(\mathcal{G})(U/Y) = \mathcal{G}(U/X)$，
$f_{big, *}(\mathcal{F})(U/X) = \mathcal{F}(U \times_X Y/Y)$。
此外，$f_{big}^{-1}$ 有左伴随 $f_{big!}$，后者与纤维积和等化子可交换。
\end{lemma}

\begin{proof}
函子 $u$ 余连续、连续，并与纤维积及等化子可交换。因此可应用
《位点》引理 \ref{sites-lemma-when-shriek} 和
\ref{sites-lemma-preserve-equalizers}，得到 $f_{big}^{-1}$
的公式以及 $f_{big!}$ 的存在性。此外，函子 $v$ 是右伴随，
因为给定
% P09-ERR-0323
$U/T$ 与 $V/X$，有
$\Mor_X(u(U), V) = \Mor_Y(U, V \times_X Y)$，正合所需。
于是可应用《位点》引理
\ref{sites-lemma-have-functor-other-way} 和
\ref{sites-lemma-have-functor-other-way-morphism}，
得到 $f_{big, *}$ 的公式。
\end{proof}

\begin{lemma}
\label{spaces-topologies-lemma-composition-fppf}
设 $S$ 为概形。给定 $S$ 上代数空间的态射
$f : X \to Y$、$g : Y \to Z$，有
$g_{big} \circ f_{big} = (g \circ f)_{big}$.
\end{lemma}

\begin{proof}
这由引理 \ref{spaces-topologies-lemma-morphism-big-fppf} 中大位点上函子的
正像与拉回的简单描述得出。
\end{proof}








\section{ph 拓扑}
\label{spaces-topologies-section-ph}

\noindent
本节定义 ph 拓扑。这是由平展覆盖和正常满态射生成的拓扑，见引理
\ref{spaces-topologies-lemma-characterize-sheaf}。

\begin{definition}
\label{spaces-topologies-definition-ph-covering}
设 $S$ 为概形，$X$ 为 $S$ 上的代数空间。
{\it $X$ 的 ph 覆盖}是 $S$ 上代数空间的一个态射族
$\{X_i \to X\}_{i \in I}$，
% P09-ERR-0324
使得 $f_i$ 局部有限型，并且对每个仿射 $U$ 及态射 $U \to X$，
存在加细族 $\{X_i \times_X U \to U\}_{i \in I}$ 的标准 ph 覆盖
$\{U_j \to U\}_{j = 1, \ldots, m}$。
\end{definition}

\noindent
% P09-ERR-0325
换言之，存在指标 $i_1, \ldots, i_m \in I$ 及态射
$h_j : U_j \to X_{i_j}$，使得
$f_{i_j} \circ h_j = h \circ g_j$。注意，若 $X$ 及所有 $X_i$ 均
可表，则由《拓扑》定义 \ref{topologies-definition-ph-covering}，
这与概形的 ph 覆盖相同。

\begin{lemma}
\label{spaces-topologies-lemma-zariski-etale-smooth-syntomic-fppf-ph}
任一 fppf 覆盖都是 ph 覆盖；因此更进一步，任一 syntomic、光滑、
平展或 Zariski 覆盖都是 ph 覆盖。
\end{lemma}

\begin{proof}
只需证明 fppf 覆盖是 ph 覆盖，其余结论便由引理
\ref{spaces-topologies-lemma-zariski-etale-smooth-syntomic-fppf} 得出。
设 $\{X_i \to X\}_{i \in I}$ 是基概形 $S$ 上代数空间的 fppf 覆盖。
设 $U$ 为仿射概形并给定态射 $U \to X$。可以将 fppf 覆盖
% P09-ERR-0326
$\{X_i \times_U U \to U\}_{i \in I}$ 加细为 fppf 覆盖
$\{T_i \to U\}_{i \in I}$，其中 $T_i$ 是概形
（引理 \ref{spaces-topologies-lemma-refine-fppf-schemes}）。
再由《态射续论》引理 \ref{more-morphisms-lemma-fppf-ph}
（以及概形的 ph 覆盖定义），可找出加细
$\{T_i \to U\}_{i \in I}$ 的标准 ph 覆盖
$\{U_j \to U\}_{j = 1, \ldots, m}$。
因此按定义，$\{X_i \to X\}_{i \in I}$ 是 ph 覆盖。
\end{proof}

\begin{lemma}
\label{spaces-topologies-lemma-surjective-proper-ph}
设 $S$ 为概形，$f : Y \to X$ 为 $S$ 上代数空间的固有满态射。
则 $\{Y \to X\}$ 是 ph 覆盖。
\end{lemma}

\begin{proof}
设 $U \to X$ 为态射，其中 $U$ 仿射。
由 Chow 引理（其弱形式见《空间的上同调》引理
\ref{spaces-cohomology-lemma-weak-chow}），存在概形的固有满态射
$V \to U$，它经 $Y \times_X U \to U$ 分解。
取 $V$ 的任一有限仿射开覆盖，便得到 $U$ 的一个标准 ph 覆盖，
它加细
% P09-ERR-0327
$\{X \times_Y U \to U\}$，正合所需。
\end{proof}

\begin{lemma}
\label{spaces-topologies-lemma-ph}
设 $S$ 为概形，$X$ 为 $S$ 上的代数空间。
\begin{enumerate}
\item 若 $X' \to X$ 是同构，则 $\{X' \to X\}$ 是 $X$ 的 ph 覆盖。
\item 若 $\{X_i \to X\}_{i\in I}$ 是 ph 覆盖，并且对每个
$i$，$\{X_{ij} \to X_i\}_{j\in J_i}$ 是 ph 覆盖，则
$\{X_{ij} \to X\}_{i \in I, j\in J_i}$ 是 ph 覆盖。
\item 若 $\{X_i \to X\}_{i\in I}$ 是 ph 覆盖，而
$X' \to X$ 是代数空间的态射，则
$\{X' \times_X X_i \to X'\}_{i\in I}$ 是 ph 覆盖。
\end{enumerate}
\end{lemma}

\begin{proof}
(1) 是显然的。考虑 (3) 中的 $g : X' \to X$ 及 ph 覆盖
$\{X_i \to X\}_{i\in I}$。由《空间的态射》引理
\ref{spaces-morphisms-lemma-base-change-finite-type}，态射
$X' \times_X X_i \to X'$ 局部有限型。
若 $h' : Z \to X'$ 是从仿射概形到 $X'$ 的态射，则置
$h = g \circ h' : Z \to X$。关于 $\{X_i \to X\}_{i\in I}$
的假设意味着，存在标准 ph 覆盖
$\{Z_j \to Z\}_{j = 1, \ldots, n}$ 以及盖住 $h$ 的态射
$Z_j \to X_{i(j)}$，其中某些 $i(j) \in I$。
由纤维积的泛性质，还得到 $h'$ 上的态射
$Z_j \to X' \times_X X_{i(j)}$。
因此 $\{X' \times_X X_i \to X'\}_{i\in I}$ 是 ph 覆盖，
这就证明了 (3)。

\medskip\noindent
设 $\{X_i \to X\}_{i\in I}$ 与
$\{X_{ij} \to X_i\}_{j\in J_i}$ 如 (2) 所述。
设 $h : Z \to X$ 为从仿射概形到 $X$ 的态射。
由假设，存在标准 ph 覆盖
$\{Z_j \to Z\}_{j = 1, \ldots, n}$ 及盖住 $h$ 的态射
$h_j : Z_j \to X_{i(j)}$，其中某些 $i(j) \in I$。
仍由假设，存在标准 ph 覆盖
$\{Z_{j, l} \to Z_j\}_{l = 1, \ldots, n(j)}$
以及盖住 $h_j$ 的态射 $Z_{j, l} \to X_{i(j)j(l)}$，
其中某些 $j(l) \in J_{i(j)}$。
由《拓扑》引理 \ref{topologies-lemma-refine-by-standard-ph}，
族 $\{Z_{j, l} \to Z\}$ 可被某个标准 ph 覆盖加细。
因此 $\{X_{ij} \to X\}_{i \in I, j\in J_i}$ 是 ph 覆盖。
\end{proof}

\begin{definition}
\label{spaces-topologies-definition-big-ph-site}
设 $S$ 为概形。大 ph 位点 {\it $(\textit{Spaces}/S)_{ph}$}
是按如下方式构造的任一位点：
\begin{enumerate}
\item 按《拓扑》第 \ref{topologies-section-ph} 节选取一个大 ph 位点
$(\Sch/S)_{ph}$。
\item 取 $S$ 上代数空间的范畴 $\textit{Spaces}/S$ 为所依范畴
（关于它为何是集合，见第 \ref{spaces-topologies-section-procedure} 节的讨论）。
\item 按《集合》引理 \ref{sets-lemma-coverings-site}，
从范畴 $\textit{Spaces}/S$ 及定义
\ref{spaces-topologies-definition-ph-covering} 的 ph 覆盖类出发，选取任一覆盖集合。
\end{enumerate}
\end{definition}

\noindent
定义完毕后，可通过局部化得到代数空间的 ph 位点。

\begin{definition}
\label{spaces-topologies-definition-big-small-ph}
设 $S$ 为概形，$(\textit{Spaces}/S)_{ph}$ 如定义
\ref{spaces-topologies-definition-big-ph-site}。
设 $X$ 为 $S$ 上的代数空间，即 $(\textit{Spaces}/S)_{ph}$ 的对象。
则 $X$ 的大 ph 位点 {\it $(\textit{Spaces}/X)_{ph}$} 是
《位点》第 \ref{sites-section-localize} 节所引入的位点
$(\textit{Spaces}/S)_{ph}$ 在 $X$ 处的局部化。
\end{definition}

\noindent
下面给出先前所承诺的 ph 层刻画。

\begin{lemma}
\label{spaces-topologies-lemma-characterize-sheaf}
设 $S$ 为概形，$X$ 为 $S$ 上的代数空间。
设 $\mathcal{F}$ 为 $(\textit{Spaces}/X)_{ph}$ 上的预层。
则 $\mathcal{F}$ 是层，当且仅当
\begin{enumerate}
\item $\mathcal{F}$ 对平展覆盖满足层条件；并且
\item 若 $f : V \to U$ 是 $(\textit{Spaces}/X)_{ph}$ 中的固有满态射，
则 $\mathcal{F}(U)$ 双射地映到两个映射
$\mathcal{F}(V) \to \mathcal{F}(V \times_U V)$ 的等化子。
\end{enumerate}
\end{lemma}

\begin{proof}
只需证明若 (1) 与 (2) 成立，则 $\mathcal{F}$ 是层。
设 $\{T_i \to T\}$ 为 ph 覆盖，即
$(\textit{Spaces}/X)_{ph}$ 中的覆盖。我们验证此覆盖的层条件。
设 $s_i \in \mathcal{F}(T_i)$ 为一些截面，它们限制到
$T_i \times_T T_{i'}$ 后相同。我们将证明存在唯一截面
% P09-ERR-0328
$s \in \mathcal{F}$，它在 $T_i$ 上限制为 $s_i$。
设 $\{U_j \to T\}$ 为平展覆盖，其中 $U_j$ 仿射。
由性质 (1)，为构造 $s$，只需构造截面
$s_j \in \mathcal{F}(U_j)$，使其在
% P09-ERR-0329
$U_j \cap U_{j'}$ 上相符。
考虑 ph 覆盖 $\{T_i \times_T U_j \to U_j\}$。
于是 $s_{ji} = s_i|_{T_i \times_T U_j}$ 是一些截面，它们在
$(T_i \times_T U_j) \times_{U_j} (T_{i'} \times_T U_j)$ 上相符。
取固有满态射 $V_j \to U_j$ 及有限仿射开覆盖
$V_j = \bigcup V_{jk}$，使标准 ph 覆盖
$\{V_{jk} \to U_j\}$ 加细
$\{T_i \times_T U_j \to U_j\}$。
若 $s_{jk} \in \mathcal{F}(V_{jk})$ 表示经隐含态射将
$s_{ji}$ 拉回到 $V_{jk}$ 所得的截面，则 $s_{jk}$ 黏合为截面
$s'_j \in \mathcal{F}(V_j)$。再次利用重叠处的一致性可知，
$s'_j$ 属于两个映射
$\mathcal{F}(V_j) \to \mathcal{F}(V_j \times_{U_j} V_j)$ 的等化子。
因此由 (2)，$s'_j$ 来自唯一截面
$s_j \in \mathcal{F}(U_j)$。略去验证这些截面
$s_j$ 具有全部所需性质。
\end{proof}

\noindent
下面建立这些位点所关联的拓扑斯之间的一些关系。

\begin{lemma}
\label{spaces-topologies-lemma-morphism-big-ph}
设 $S$ 为概形，$f : Y \to X$ 为 $S$ 上代数空间的态射。函子
$$
u : (\textit{Spaces}/Y)_{ph} \longrightarrow (\textit{Spaces}/X)_{ph},
\quad
V/Y \longmapsto V/X
$$
是余连续的，并有连续右伴随
% P09-ERR-0330
$$
v : (\textit{Spaces}/X)_{ph} \longrightarrow (\textit{Spaces}/Y)_{ph},
\quad
(U \to Y) \longmapsto (U \times_X Y \to Y).
$$
它们诱导同一拓扑斯态射
$$
f_{big} :
\Sh((\textit{Spaces}/Y)_{ph})
\longrightarrow
\Sh((\textit{Spaces}/X)_{ph})
$$
且有 $f_{big}^{-1}(\mathcal{G})(U/Y) = \mathcal{G}(U/X)$，
$f_{big, *}(\mathcal{F})(U/X) = \mathcal{F}(U \times_X Y/Y)$。
此外，$f_{big}^{-1}$ 有左伴随 $f_{big!}$，后者与纤维积和等化子可交换。
\end{lemma}

\begin{proof}
函子 $u$ 余连续、连续，并与纤维积及等化子可交换。因此可应用
《位点》引理 \ref{sites-lemma-when-shriek} 和
\ref{sites-lemma-preserve-equalizers}，得到 $f_{big}^{-1}$
的公式以及 $f_{big!}$ 的存在性。此外，函子 $v$ 是右伴随，
因为给定
% P09-ERR-0331
$U/T$ 与 $V/X$，有
$\Mor_X(u(U), V) = \Mor_Y(U, V \times_X Y)$，正合所需。
于是可应用《位点》引理
\ref{sites-lemma-have-functor-other-way} 和
\ref{sites-lemma-have-functor-other-way-morphism}，
得到 $f_{big, *}$ 的公式。
\end{proof}

\begin{lemma}
\label{spaces-topologies-lemma-composition-ph}
设 $S$ 为概形。给定 $S$ 上代数空间的态射
$f : X \to Y$、$g : Y \to Z$，有
$g_{big} \circ f_{big} = (g \circ f)_{big}$.
\end{lemma}

\begin{proof}
这由引理 \ref{spaces-topologies-lemma-morphism-big-ph} 中大位点上函子的
正像与拉回的简单描述得出。
\end{proof}

\begin{lemma}
\label{spaces-topologies-lemma-cech-enough}
设 $S$ 为概形，$X$ 为 $S$ 上的代数空间。
设 $P$ 为 $(\textit{Spaces}/X)_{fppf}$ 中对象的一种性质，
使得每当 $\{U_i \to U\}$ 是
$(\textit{Spaces}/X)_{fppf}$ 中的覆盖时，都有
$$
P(U_{i_0} \times_U \ldots \times_U U_{i_p})
\text{ 对所有 }
p \geq 0,\ i_0, \ldots, i_p \in I
\Rightarrow P(U)
$$
若对所有在 $X$ 上平坦、局部有限表示的仿射 $U$ 都有 $P(U)$，
则 $P(X)$。
\end{lemma}

\begin{proof}
设 $U$ 为 $X$ 上局部有限表示的分离代数空间。可选择平展覆盖
$\{U_i \to U\}_{i \in I}$，其中
% P09-ERR-0332
$V_i$ 仿射。由于 $U$ 分离，
$U_{i_0} \times_U \ldots \times_U U_{i_p}$ 总是仿射，
故总有 $P(U_{i_0} \times_U \ldots \times_U U_{i_p})$。
于是 $P(U)$ 成立。选择一个由仿射概形不交并而成的概形 $U$
以及满平展态射 $U \to X$。
则 $U \times_X \ldots \times_X U$（共有 $p + 1$ 个因子）
是 $X$ 上平展的分离代数空间。因此由上可知
$P(U \times_X \ldots \times_X U)$，从而 $P(X)$ 成立。
\end{proof}












\section{Fpqc 拓扑}
\label{spaces-topologies-section-fpqc}

\noindent
我们简要讨论代数空间的 fpqc 覆盖概念。请与《拓扑》第
\ref{topologies-section-fpqc} 节比较。
我们将在《空间的下降》命题
\ref{spaces-descent-proposition-fpqc-descent-quasi-coherent} 中证明，
% P09-ERR-0333
拟凝聚层可沿这些覆盖下降。

\begin{definition}
\label{spaces-topologies-definition-fpqc-covering}
设 $S$ 为概形，$X$ 为 $S$ 上的代数空间。
{\it $X$ 的 fpqc 覆盖}是代数空间的一个态射族
$\{f_i : X_i \to X\}_{i \in I}$，其中每个 $f_i$ 都平坦；
并且对每个仿射概形 $Z$ 及态射 $h : Z \to X$，都存在加细族
$\{X_i \times_X Z \to Z\}_{i \in I}$ 的标准 fpqc 覆盖
$\{g_j : Z_j \to Z\}_{j = 1, \ldots, m}$。
\end{definition}

\noindent
% P09-ERR-0334
换言之，存在指标 $i_1, \ldots, i_m \in I$ 及态射
$h_j : Z_j \to X_{i_j}$，使得
$f_{i_j} \circ h_j = h \circ g_j$。注意，若 $X$ 及所有 $X_i$ 均
可表，则由《拓扑》引理
\ref{topologies-lemma-fpqc-covering-affines-mapping-in}，
这与概形的 fpqc 覆盖相同。

\begin{lemma}
\label{spaces-topologies-lemma-zariski-etale-smooth-syntomic-fppf-fpqc}
任一 fppf 覆盖都是 fpqc 覆盖；因此更进一步，任一 syntomic、光滑、
平展或 Zariski 覆盖都是 fpqc 覆盖。
\end{lemma}

\begin{proof}
只需证明 fppf 覆盖是 fpqc 覆盖，其余结论便由引理
\ref{spaces-topologies-lemma-zariski-etale-smooth-syntomic-fppf} 得出。
设 $\{f_i : U_i \to U\}_{i \in I}$ 是 $S$ 上代数空间的 fppf 覆盖。
按定义，$f_i$ 平坦，这验证了定义
\ref{spaces-topologies-definition-fpqc-covering} 的第一个条件。
为验证第二个条件，设 $V \to U$ 为态射，其中 $V$ 仿射。
可选择平展覆盖 $\{V_{ij} \to V \times_U U_i\}$，其中 $V_{ij}$ 仿射。
于是复合
$f_{ij} : V_{ij} \to V \times_U U_i \to V$ 平坦且局部有限表示，
因为它们是具有这些性质的态射的复合
（《空间的态射》引理
\ref{spaces-morphisms-lemma-composition-finite-presentation},
\ref{spaces-morphisms-lemma-composition-flat},
\ref{spaces-morphisms-lemma-etale-flat} 及
\ref{spaces-morphisms-lemma-etale-locally-finite-presentation}).
因此这些态射是开的
（《空间的态射》引理 \ref{spaces-morphisms-lemma-fppf-open}），并且
$|V| = \bigcup_{i \in I} \bigcup_{j \in J_i} f_{ij}(|V_{ij}|)$
是 $|V|$ 的开覆盖。由于 $|V|$ 拟紧，该覆盖有有限加细。
设 $V_{i_1j_1}, \ldots, V_{i_Nj_N}$ 实现这一点。
则 $\{V_{i_kj_k} \to V\}_{k = 1, \ldots, N}$ 是 $V$ 的标准
fpqc 覆盖，并加细族
% P09-ERR-0335
$\{U_i \times_U V \to V\}$。证明完毕。
\end{proof}

\begin{lemma}
\label{spaces-topologies-lemma-fpqc}
设 $S$ 为概形，$X$ 为 $S$ 上的代数空间。
\begin{enumerate}
\item 若 $X' \to X$ 是同构，则 $\{X' \to X\}$ 是 $X$ 的 fpqc 覆盖。
\item 若 $\{X_i \to X\}_{i\in I}$ 是 fpqc 覆盖，并且对每个
$i$，$\{X_{ij} \to X_i\}_{j\in J_i}$ 是 fpqc 覆盖，则
$\{X_{ij} \to X\}_{i \in I, j\in J_i}$ 是 fpqc 覆盖。
\item 若 $\{X_i \to X\}_{i\in I}$ 是 fpqc 覆盖，而
$X' \to X$ 是代数空间的态射，则
$\{X' \times_X X_i \to X'\}_{i\in I}$ 是 fpqc 覆盖。
\end{enumerate}
\end{lemma}

\begin{proof}
(1) 是显然的。考虑 (3) 中的 $g : X' \to X$ 及 fpqc 覆盖
$\{X_i \to X\}_{i\in I}$。由《空间的态射》引理
\ref{spaces-morphisms-lemma-base-change-flat}，态射
$X' \times_X X_i \to X'$ 平坦。
若 $h' : Z \to X'$ 是从仿射概形到 $X'$ 的态射，则置
$h = g \circ h' : Z \to X$。关于 $\{X_i \to X\}_{i\in I}$
的假设意味着，存在标准 fpqc 覆盖
$\{Z_j \to Z\}_{j = 1, \ldots, n}$ 以及盖住 $h$ 的态射
$Z_j \to X_{i(j)}$，其中某些 $i(j) \in I$。
由纤维积的泛性质，还得到 $h'$ 上的态射
$Z_j \to X' \times_X X_{i(j)}$。
因此 $\{X' \times_X X_i \to X'\}_{i\in I}$ 是 fpqc 覆盖，
这就证明了 (3)。

\medskip\noindent
设 $\{X_i \to X\}_{i\in I}$ 与
$\{X_{ij} \to X_i\}_{j\in J_i}$ 如 (2) 所述。
设 $h : Z \to X$ 为从仿射概形到 $X$ 的态射。
由假设，存在标准 fpqc 覆盖
$\{Z_j \to Z\}_{j = 1, \ldots, n}$ 及盖住 $h$ 的态射
$h_j : Z_j \to X_{i(j)}$，其中某些 $i(j) \in I$。
仍由假设，存在标准 fpqc 覆盖
$\{Z_{j, l} \to Z_j\}_{l = 1, \ldots, n(j)}$
以及盖住 $h_j$ 的态射 $Z_{j, l} \to X_{i(j)j(l)}$，
其中某些 $j(l) \in J_{i(j)}$。
由《拓扑》引理 \ref{topologies-lemma-fpqc-affine-axioms}，
族 $\{Z_{j, l} \to Z\}$ 是标准 fpqc 覆盖。
因此 $\{X_{ij} \to X\}_{i \in I, j\in J_i}$ 是 fpqc 覆盖。
\end{proof}

\begin{lemma}
\label{spaces-topologies-lemma-recognize-fpqc-covering}
设 $S$ 为概形，$X$ 为 $S$ 上的代数空间。
假设 $\{f_i : X_i \to X\}_{i \in I}$ 是以 $X$ 为靶的代数空间态射族。
设 $U \to X$ 为从概形到 $X$ 的满平展态射。
则 $\{f_i : X_i \to X\}_{i \in I}$ 是 $X$ 的 fpqc 覆盖，
当且仅当 $\{U \times_X X_i \to U\}_{i \in I}$ 是 $U$ 的 fpqc 覆盖。
\end{lemma}

\begin{proof}
若 $\{X_i \to X\}_{i \in I}$ 是 fpqc 覆盖，则由引理
\ref{spaces-topologies-lemma-fpqc}，$\{U \times_X X_i \to U\}_{i \in I}$ 也是。
反之，假设 $\{U \times_X X_i \to U\}_{i \in I}$ 是 fpqc 覆盖。
设 $h : Z \to X$ 为从仿射概形到 $X$ 的态射。
则 $U \times_X Z \to Z$ 是概形的满平展态射，特别地是开态射。
因此可找出 $U \times_X Z$ 的有限多个仿射开集
$W_1, \ldots, W_t$，其像覆盖 $Z$。
对每个 $j$，把 $\{U \times_X X_i \to U\}_{i \in I}$
是 fpqc 覆盖这一条件应用于态射 $W_j \to U$，得到加细
$\{W_j \times_X X_i \to W_j\}_{i \in I}$ 的标准 fpqc 覆盖
$\{W_{jl} \to W_j\}$。
于是 $\{W_{jl} \to Z\}$ 是 $Z$ 的标准 fpqc 覆盖
（见《拓扑》引理 \ref{topologies-lemma-fpqc-affine-axioms}），
并加细
% P09-ERR-0336
$\{Z \times_X X_i \to X\}$，即得结论。
\end{proof}

\begin{lemma}
\label{spaces-topologies-lemma-refine-fpqc-schemes}
设 $S$ 为概形，$X$ 为 $S$ 上的代数空间。假设
$\mathcal{U} = \{f_i : X_i \to X\}_{i \in I}$ 是 $X$ 的 fpqc 覆盖。
则存在 $\mathcal{U}$ 的加细
$\mathcal{V} = \{g_i : T_i \to X\}$，它是 fpqc 覆盖且每个
$T_i$ 都是概形。
\end{lemma}

\begin{proof}
略。提示：对每个 $i$，取概形 $T_i$ 及满平展态射
$T_i \to X_i$，再验证 $\{T_i \to X\}$ 是 fpqc 覆盖。
\end{proof}

\noindent
待续……
% Generated mechanically from the frozen spaces-descent.tex target.
% Do not edit this generated file; edit the bound locale source instead.

% OK, start here.
%

\setcounter{chapter}{73}
\chapter{下降与代数空间}


\phantomsection
\label{spaces-descent-section-phantom}


\section{引言}
\label{spaces-descent-section-introduction}

\noindent
在关于代数空间上拓扑的章中（见《空间上的拓扑》第
\ref{spaces-topologies-section-introduction} 节），我们引入了代数空间的
平展、fppf、光滑、syntomic 和 fpqc 覆盖。
本章讨论代数空间上的哪些结构能够沿这类覆盖下降。
例如见 \cite{Gr-I}、\cite{Gr-II}、\cite{Gr-III}、
\cite{Gr-IV}、\cite{Gr-V} 和 \cite{Gr-VI}。



\section{约定}
\label{spaces-descent-section-conventions}

\noindent
始终假设所有概形都包含在某个大 fppf 位点 $\Sch_{fppf}$ 中。
并且所考虑的每个环 $A$ 都具有如下性质：
$\Spec(A)$（同构于）这个大位点的一个对象。

\medskip\noindent
设 $S$ 为概形，$X$ 为 $S$ 上的代数空间。
在本章及后续各章中，我们用 $X \times_S X$
表示 $X$ 与自身的积（在 $S$ 上代数空间的范畴中），
而不用 $X \times X$。




\section{拟凝聚层的下降资料}
\label{spaces-descent-section-equivalence}

\noindent
本节是《下降》第 \ref{descent-section-equivalence} 节
在代数空间情形下的类似版本。
先阅读该节是有益的。

\begin{definition}
\label{spaces-descent-definition-descent-datum-quasi-coherent}
设 $S$ 为概形。设 $\{f_i : X_i \to X\}_{i \in I}$
为 $S$ 上具有固定靶 $X$ 的代数空间态射族。
\begin{enumerate}
\item 关于给定族的一个{\it 拟凝聚层下降资料
$(\mathcal{F}_i, \varphi_{ij})$}由下列数据给出：
一个拟凝聚层 $\mathcal{F}_i$，它位于 $X_i$ 上，
对每个 $i \in I$；以及拟凝聚
$\mathcal{O}_{X_i \times_X X_j}$-模的同构
$\varphi_{ij} : \text{pr}_0^*\mathcal{F}_i \to \text{pr}_1^*\mathcal{F}_j$，
对每一对 $(i, j) \in I^2$，
使得对每个三元指标组 $(i, j, k) \in I^3$，图
$$
\xymatrix{
\text{pr}_0^*\mathcal{F}_i \ar[rd]_{\text{pr}_{01}^*\varphi_{ij}}
\ar[rr]_{\text{pr}_{02}^*\varphi_{ik}} & &
\text{pr}_2^*\mathcal{F}_k \\
& \text{pr}_1^*\mathcal{F}_j \ar[ru]_{\text{pr}_{12}^*\varphi_{jk}} &
}
$$
作为 $\mathcal{O}_{X_i \times_X X_j \times_X X_k}$-模的图交换。
这称为{\it 余圈条件}。
\item 一个{\it 下降资料态射 $\psi : (\mathcal{F}_i, \varphi_{ij}) \to
(\mathcal{F}'_i, \varphi'_{ij})$}由一个态射族
$\psi = (\psi_i)_{i\in I}$ 给出；这些态射是
$\mathcal{O}_{X_i}$-模态射
$\psi_i : \mathcal{F}_i \to \mathcal{F}'_i$，并且所有图
$$
\xymatrix{
\text{pr}_0^*\mathcal{F}_i \ar[r]_{\varphi_{ij}} \ar[d]_{\text{pr}_0^*\psi_i}
& \text{pr}_1^*\mathcal{F}_j \ar[d]^{\text{pr}_1^*\psi_j} \\
\text{pr}_0^*\mathcal{F}'_i \ar[r]^{\varphi'_{ij}} &
\text{pr}_1^*\mathcal{F}'_j \\
}
$$
交换。
\end{enumerate}
\end{definition}

\begin{lemma}
\label{spaces-descent-lemma-map-families}
设 $S$ 为概形。
设 $\mathcal{U} = \{U_i \to U\}_{i \in I}$ 和
$\mathcal{V} = \{V_j \to V\}_{j \in J}$
为 $S$ 上具有固定靶的代数空间态射族。
设 $(g, \alpha : I \to J, (g_i)) : \mathcal{U} \to \mathcal{V}$
为具有固定靶的映射族之间的一个态射，见
《位点》定义 \ref{sites-definition-morphism-coverings}。
设 $(\mathcal{F}_j, \varphi_{jj'})$ 为关于族
$\{V_j \to V\}_{j \in J}$ 的拟凝聚层下降资料。则
\begin{enumerate}
\item 系统
$$
\left(g_i^*\mathcal{F}_{\alpha(i)},
(g_i \times g_{i'})^*\varphi_{\alpha(i)\alpha(i')}\right)
$$
是关于族 $\{U_i \to U\}_{i \in I}$ 的下降资料。
\item 这一构造关于下降资料
$(\mathcal{F}_j, \varphi_{jj'})$ 是函子性的。
\item 给定具有固定靶的映射族之间的第二个态射
$(g', \alpha' : I \to J, (g'_i))$，
若 $g = g'$，则存在下降资料的函子性同构
$$
(g_i^*\mathcal{F}_{\alpha(i)},
(g_i \times g_{i'})^*\varphi_{\alpha(i)\alpha(i')})
\cong
((g'_i)^*\mathcal{F}_{\alpha'(i)},
(g'_i \times g'_{i'})^*\varphi_{\alpha'(i)\alpha'(i')}).
$$
\end{enumerate}
\end{lemma}

\begin{proof}
略。提示：给出第 (3) 部分下降资料同构的映射
$g_i^*\mathcal{F}_{\alpha(i)} \to (g'_i)^*\mathcal{F}_{\alpha'(i)}$
是把映射 $\varphi_{\alpha(i)\alpha'(i)}$ 沿态射
$(g_i, g'_i) : U_i \to V_{\alpha(i)} \times_V V_{\alpha'(i)}$
拉回所得的映射。
\end{proof}

\noindent
设 $g : U \to V$ 为代数空间的一个态射。
上述引理说明，只要具有靶 $V$ 与 $U$ 的映射族之间存在一个
“位于 $g$ 之上”的态射，相对于这两个族的下降资料范畴之间
便有一个良定义的拉回函子。

\begin{definition}
\label{spaces-descent-definition-descent-datum-effective-quasi-coherent}
设 $S$ 为概形。
设 $\{U_i \to U\}_{i \in I}$ 为 $S$ 上具有固定靶的
代数空间态射族。
\begin{enumerate}
\item 设 $\mathcal{F}$ 为拟凝聚 $\mathcal{O}_U$-模。
我们把 $\mathcal{F}$ 上关于覆盖 $\{U \to U\}$ 的
唯一下降资料称为{\it 平凡下降资料}。
% P09-ERR-0337
\item 平凡下降资料到 $\{U_i \to U\}$ 的拉回称为
{\it 典范下降资料}。记作 $(\mathcal{F}|_{U_i}, can)$。
\item 关于给定族的拟凝聚层下降资料
$(\mathcal{F}_i, \varphi_{ij})$ 称为{\it 有效的}，
如果存在一个拟凝聚层 $\mathcal{F}$ 位于 $U$ 上，使
$(\mathcal{F}_i, \varphi_{ij})$ 同构于
$(\mathcal{F}|_{U_i}, can)$。
\end{enumerate}
\end{definition}

\begin{lemma}
\label{spaces-descent-lemma-zariski-descent-effective}
设 $S$ 为概形，$U$ 为 $S$ 上的代数空间。
设 $\{U_i \to U\}$ 为 $U$ 的一个 Zariski 覆盖，见
《空间上的拓扑》定义
\ref{spaces-topologies-definition-zariski-covering}。
关于族 $\mathcal{U} = \{U_i \to U\}$ 的任意拟凝聚层下降资料
都是有效的。此外，从拟凝聚 $\mathcal{O}_U$-模范畴到
关于 $\{U_i \to U\}$ 的下降资料范畴的函子是全忠实的。
\end{lemma}

\begin{proof}
略。
\end{proof}









\section{拟凝聚层的 Fpqc 下降}
\label{spaces-descent-section-fpqc-descent-quasi-coherent}

\noindent
模的平坦下降的主要应用，是相应的关于 fpqc 覆盖的
拟凝聚层下降命题。

\begin{proposition}
\label{spaces-descent-proposition-fpqc-descent-quasi-coherent}
设 $S$ 为概形。
设 $\{X_i \to X\}$ 为 $S$ 上代数空间的一个 fpqc 覆盖，见
《空间上的拓扑》定义
\ref{spaces-topologies-definition-fpqc-covering}。
关于 $\{X_i \to X\}$ 的任意拟凝聚层下降资料都是有效的。
此外，从拟凝聚 $\mathcal{O}_X$-模范畴到关于
$\{X_i \to X\}$ 的下降资料范畴的函子是全忠实的。
\end{proposition}

\begin{proof}
这大致是概形的相应结果的一个形式推论，见
《下降》命题
\ref{descent-proposition-fpqc-descent-quasi-coherent}。
以下是一种证明策略：
\begin{enumerate}
\item $\{X_i \to X\}$ 是平凡覆盖 $\{X \to X\}$ 的一个加细，
这一事实经由引理 \ref{spaces-descent-lemma-map-families} 给出函子
$\QCoh(\mathcal{O}_X) \to DD(\{X_i \to X\})$，
它从拟凝聚 $\mathcal{O}_X$-模范畴映到给定族的下降资料范畴。
\item 为证明该命题，我们将构造一个拟逆函子
$back : DD(\{X_i \to X\}) \to \QCoh(\mathcal{O}_X)$。
\item 再次应用引理 \ref{spaces-descent-lemma-map-families} 可见，有函子
$DD(\{X_i \to X\}) \to DD(\{T_j \to X\})$，
如果 $\{T_j \to X\}$ 是给定族的一个加细。
因此，为构造函子 $back$，我们可以假设每个 $X_i$ 都是概形，
见《空间上的拓扑》引理
\ref{spaces-topologies-lemma-refine-fpqc-schemes}。
这把问题约化到所有 $X_i$ 都是概形的情形。
\item $X$ 上的拟凝聚层按定义就是拟凝聚
$\mathcal{O}_X$-模，位于 $X_\etale$ 上。现在，对任意
$U \in \Ob(X_\etale)$，我们得到由概形组成的一个 fppf 覆盖
$\{U_i \times_X X_i \to U\}$，以及覆盖态射
$g : \{U_i \times_X X_i \to U\} \to \{X_i \to X\}$，
它位于 $U \to X$ 之上。
% P09-ERR-0338
给定下降资料 $\xi = (\mathcal{F}_i, \varphi_{ij})$，
我们得到拟凝聚 $\mathcal{O}_U$-模
$\mathcal{F}_{\xi, U}$；它对应于引理
\ref{spaces-descent-lemma-map-families} 中的拉回 $g^*\xi$
到 $U$ 的覆盖，并使用概形的 fppf 覆盖的有效性，见
《下降》命题
\ref{descent-proposition-fpqc-descent-quasi-coherent}。
\item 验证 $\xi \mapsto \mathcal{F}_{\xi, U}$ 关于 $\xi$
是函子性的。略。
\item 验证 $\xi \mapsto \mathcal{F}_{\xi, U}$ 与态射
$U \to U'$ 相容，其中该态射属于位点 $X_\etale$，从而层系统
$\mathcal{F}_{\xi, U}$ 对应于一个拟凝聚层
$\mathcal{F}_\xi$，它位于 $X_\etale$ 上，见《空间的性质》引理
\ref{spaces-properties-lemma-characterize-quasi-coherent-small-etale}。
细节略。
\item 验证 $back : \xi \mapsto \mathcal{F}_\xi$
是第 (1) 项所构造函子的拟逆。略。
\end{enumerate}
证毕。
\end{proof}








\section{拟凝聚模与仿射对象}
\label{spaces-descent-section-alternative-quasi-coherent}

\noindent
设 $S$ 为概形，$X$ 为 $S$ 上的代数空间。
回忆，$X_{affine, \etale}$ 是 $X_\etale$ 的全子范畴，
其对象为仿射对象；通过宣告覆盖为标准平展覆盖，把它视为位点。
见《空间的性质》定义
\ref{spaces-properties-definition-affine-etale-site}。
由《空间的性质》引理 \ref{spaces-properties-lemma-alternative}，
我们有拓扑斯的等价
$g : \Sh(X_{affine, \etale}) \to \Sh(X_\etale)$，
其拉回函子由限制给出。
回忆，$\mathcal{O}_X$ 表示 $X_\etale$ 上的结构层。
于是得到带环拓扑斯的等价
\begin{equation}
\label{spaces-descent-equation-alternative-small-ringed}
(\Sh(X_{affine, \etale}), \mathcal{O}_X|_{X_{affine, \etale}})
\longrightarrow
(\Sh(X_\etale), \mathcal{O}_X).
\end{equation}
我们经常用 $\mathcal{O}_X$ 代替
$\mathcal{O}_X|_{X_{affine, \etale}}$。
% P09-ERR-0339
有了这些约定，我们也可以比较拟凝聚模。

\begin{lemma}
\label{spaces-descent-lemma-quasi-coherent-alternative-small}
设 $S$ 为概形，$X$ 为 $S$ 上的代数空间。
设 $\mathcal{F}$ 为 $\mathcal{O}_X$-模预层，
位于 $X_{affine, \etale}$ 上。下列条件等价：
\begin{enumerate}
\item 对每个态射 $U \to U'$，它属于 $X_{affine, \etale}$，映射
$\mathcal{F}(U') \otimes_{\mathcal{O}_X(U')} \mathcal{O}_X(U)
\to \mathcal{F}(U)$ 是同构；
\item $\mathcal{F}$ 是带环位点
$(X_{affine, \etale}, \mathcal{O}_X)$ 上的拟凝聚模，
其含义见《位点上的模》定义
\ref{sites-modules-definition-site-local}；
\item 经由等价 (\ref{spaces-descent-equation-alternative-small-ringed})，
$\mathcal{F}$ 对应于 $X$ 上的一个拟凝聚模。
\end{enumerate}
\end{lemma}

\begin{proof}
假设 (1) 成立。为证明 $\mathcal{F}$ 是层，设
$\mathcal{U} = \{U_i \to U\}_{i = 1, \ldots, n}$
为 $X_{affine, \etale}$ 的一个覆盖。
$\mathcal{F}$ 对 $\mathcal{U}$ 的层条件，
由我们关于 $\mathcal{F}$ 的假设，约化为证明
$$
0 \to \mathcal{F}(U) \to
\prod \mathcal{F}(U) \otimes_{\mathcal{O}_X(U)} \mathcal{O}_X(U_i) \to
\prod \mathcal{F}(U) \otimes_{\mathcal{O}_X(U)} \mathcal{O}_X(U_i \times_U U_j)
$$
正合。这是成立的，因为
$\mathcal{O}_X(U) \to \prod \mathcal{O}_X(U_i)$
是忠实平坦的（由《下降》引理
\ref{descent-lemma-standard-covering-Cech}，并且
$X_{affine, \etale}$ 中的覆盖是标准平展覆盖），
故可应用《下降》引理 \ref{descent-lemma-ff-exact}。
接着证明 $\mathcal{F}$ 在 $X_{affine, \etale}$ 上拟凝聚。
具体而言，对 $U$，它属于 $X_{affine, \etale}$，令
$R = \mathcal{O}_X(U)$，并选择一个展示
$$
\bigoplus\nolimits_{k \in K} R
\longrightarrow
\bigoplus\nolimits_{l \in L} R
\longrightarrow
\mathcal{F}(U)
\longrightarrow 0
$$
由自由 $R$-模给出。由性质 (1) 与张量积的右正合性可见，
对每个态射 $U' \to U$，它属于 $X_{affine, \etale}$，
我们得到展示
$$
\bigoplus\nolimits_{k \in K} \mathcal{O}_X(U')
\longrightarrow
\bigoplus\nolimits_{l \in L} \mathcal{O}_X(U')
\longrightarrow
\mathcal{F}(U')
\longrightarrow 0
$$
换言之，$\mathcal{F}$ 到局部化范畴
$X_{affine, etale}/U$ 的限制具有展示
% P09-ERR-0340
$$
\bigoplus\nolimits_{k \in K} \mathcal{O}_X|_{X_{affine, \etale}/U}
\longrightarrow
\bigoplus\nolimits_{l \in L} \mathcal{O}_X|_{X_{affine, \etale}/U}
\longrightarrow
\mathcal{F}|_{X_{affine, \etale}/U}
\longrightarrow 0
$$
这正是证明 $\mathcal{F}$ 拟凝聚所需要的。
谨为这套糟糕的记号致歉；由此完成 (1) 蕴含 (2) 的证明。

\medskip\noindent
由于拟凝聚模的概念是内蕴的（《位点上的模》引理
\ref{sites-modules-lemma-special-locally-free}），
等价 (\ref{spaces-descent-equation-alternative-small-ringed})
诱导拟凝聚模范畴之间的等价。
因此 (2) 与 (3) 等价。

\medskip\noindent
假设 (3) 并证明 (1)。具体而言，设
$\mathcal{G}$ 为 $X$ 上与 $\mathcal{F}$ 对应的拟凝聚模。
设 $h : U \to U' \to X$ 为 $X_{affine, \etale}$ 的一个态射。
把结构态射记作 $f : U \to X$ 和 $f' : U' \to X$，
% P09-ERR-0341
于是 $f = f' \circ h$。我们有
$\mathcal{F}(U') = \Gamma(U', (f')^*\mathcal{G})$ 以及
$\mathcal{F}(U) = \Gamma(U, f^*\mathcal{G}) = \Gamma(U, h^*(f')^*\mathcal{G})$。
因此 (1) 由《概形》引理
\ref{schemes-lemma-widetilde-pullback} 得出。
\end{proof}

\section{模的有限性性质的下降}
\label{spaces-descent-section-descent-finiteness}

\noindent
本节是《下降》第 \ref{descent-section-descent-finiteness} 节
在代数空间情形下的类似版本。
目标是说明，可以通过在一个覆盖的各成员上检验，来检验一个拟凝聚模
是否具有某个有限性条件。
% P09-ERR-0342
我们将反复使用如下证明模式。
假设 $X$ 是代数空间，且 $\{X_i \to X\}$
是一个 fppf（相应地，fpqc）覆盖。设 $U \to X$
为满平展态射，其中 $U$ 是概形。则存在一个
fppf（相应地，fpqc）覆盖 $\{Y_j \to X\}$，使得
\begin{enumerate}
\item $\{Y_j \to X\}$ 是 $\{X_i \to X\}$ 的一个加细；
\item 每个 $Y_j$ 都是概形；
\item 每个态射 $Y_j \to X$ 都经 $U$ 分解；并且
% P09-ERR-0343
\item $\{Y_j \to U\}$ 是 $U$ 的一个
fppf（相应地，fpqc）覆盖。
\end{enumerate}
具体而言，先用一个 fppf（相应地，fpqc）覆盖加细
$\{X_i \to X\}$，使每个 $X_i$ 都是概形，见
《空间上的拓扑》引理
\ref{spaces-topologies-lemma-refine-fppf-schemes}，
相应地，引理 \ref{spaces-topologies-lemma-refine-fpqc-schemes}。
然后令 $Y_i = U \times_X X_i$。拟凝聚
$\mathcal{O}_X$-模 $\mathcal{F}$ 是有限型、有限展示等等，
当且仅当拟凝聚 $\mathcal{O}_U$-模 $\mathcal{F}|_U$
是有限型、有限展示等等。
因此，我们可以利用加细 $\{Y_j \to X\}$ 的存在性，
把下列引理的证明约化到概形的情形。
我们将用“{\it 结果由概形情形经平展局部化得出}”来表示这一点。

\begin{lemma}
\label{spaces-descent-lemma-finite-type-descends}
设 $X$ 为概形 $S$ 上的代数空间。
设 $\mathcal{F}$ 为拟凝聚 $\mathcal{O}_X$-模。
设 $\{f_i : X_i \to X\}_{i \in I}$ 为一个 fpqc 覆盖，
使每个 $f_i^*\mathcal{F}$ 都是有限型
$\mathcal{O}_{X_i}$-模。
则 $\mathcal{F}$ 是有限型 $\mathcal{O}_X$-模。
\end{lemma}

\begin{proof}
这由概形情形经平展局部化得出，见
《下降》引理 \ref{descent-lemma-finite-type-descends}。
\end{proof}

\begin{lemma}
\label{spaces-descent-lemma-finite-presentation-descends}
设 $X$ 为概形 $S$ 上的代数空间。
设 $\mathcal{F}$ 为拟凝聚 $\mathcal{O}_X$-模。
设 $\{f_i : X_i \to X\}_{i \in I}$ 为一个 fpqc 覆盖，
使每个 $f_i^*\mathcal{F}$ 都是有限展示
$\mathcal{O}_{X_i}$-模。
则 $\mathcal{F}$ 是有限展示 $\mathcal{O}_X$-模。
\end{lemma}

\begin{proof}
这由概形情形经平展局部化得出，见
《下降》引理 \ref{descent-lemma-finite-presentation-descends}。
\end{proof}

\begin{lemma}
\label{spaces-descent-lemma-flat-descends}
设 $X$ 为概形 $S$ 上的代数空间。
设 $\mathcal{F}$ 为拟凝聚 $\mathcal{O}_X$-模。
设 $\{f_i : X_i \to X\}_{i \in I}$ 为一个 fpqc 覆盖，
使每个 $f_i^*\mathcal{F}$ 都是平坦
$\mathcal{O}_{X_i}$-模。
则 $\mathcal{F}$ 是平坦 $\mathcal{O}_X$-模。
\end{lemma}

\begin{proof}
这由概形情形经平展局部化得出，见
《下降》引理 \ref{descent-lemma-flat-descends}。
\end{proof}

\begin{lemma}
\label{spaces-descent-lemma-finite-locally-free-descends}
设 $X$ 为概形 $S$ 上的代数空间。
设 $\mathcal{F}$ 为拟凝聚 $\mathcal{O}_X$-模。
设 $\{f_i : X_i \to X\}_{i \in I}$ 为一个 fpqc 覆盖，
使每个 $f_i^*\mathcal{F}$ 都是有限局部自由
$\mathcal{O}_{X_i}$-模。
则 $\mathcal{F}$ 是有限局部自由 $\mathcal{O}_X$-模。
\end{lemma}

\begin{proof}
这由概形情形经平展局部化得出，见
《下降》引理 \ref{descent-lemma-finite-locally-free-descends}。
\end{proof}

\noindent
局部投射拟凝聚层的定义见《空间的性质》第
\ref{spaces-properties-section-locally-projective} 节。
那里也证明了这一概念在拉回下保持。

\begin{lemma}
\label{spaces-descent-lemma-locally-projective-descends}
设 $X$ 为概形 $S$ 上的代数空间。
设 $\mathcal{F}$ 为拟凝聚 $\mathcal{O}_X$-模。
设 $\{f_i : X_i \to X\}_{i \in I}$ 为一个 fpqc 覆盖，
使每个 $f_i^*\mathcal{F}$ 都是局部投射
$\mathcal{O}_{X_i}$-模。
则 $\mathcal{F}$ 是局部投射 $\mathcal{O}_X$-模。
\end{lemma}

\begin{proof}
这由概形情形经平展局部化得出，见
《下降》引理 \ref{descent-lemma-locally-projective-descends}。
\end{proof}

\noindent
我们还在这里加入两个与上述结果相关、但性质略有不同的结果。

\begin{lemma}
\label{spaces-descent-lemma-finite-over-finite-module}
设 $S$ 为概形。
设 $f : X \to Y$ 为 $S$ 上代数空间的一个态射。
设 $\mathcal{F}$ 为拟凝聚 $\mathcal{O}_X$-模。
假设 $f$ 是有限态射。
则 $\mathcal{F}$ 是有限型 $\mathcal{O}_X$-模，
当且仅当 $f_*\mathcal{F}$ 是有限型
$\mathcal{O}_Y$-模。
\end{lemma}

\begin{proof}
由于 $f$ 有限，它是可表的。选择概形 $V$ 以及满平展态射
$V \to Y$。于是 $U = V \times_Y X$ 是概形，
它有一个到 $X$ 的满平展态射以及一个有限态射
$\psi : U \to V$（即 $f$ 的基变换）。由于
$\psi_*(\mathcal{F}|_U) = f_*\mathcal{F}|_V$，
该引理的结论立即由概形版本
《下降》引理
\ref{descent-lemma-finite-over-finite-module}
得出。
% P09-ERR-0344
\end{proof}

\begin{lemma}
\label{spaces-descent-lemma-finite-finitely-presented-module}
设 $S$ 为概形。
设 $f : X \to Y$ 为 $S$ 上代数空间的一个态射。
设 $\mathcal{F}$ 为拟凝聚 $\mathcal{O}_X$-模。
假设 $f$ 有限且有限展示。
则 $\mathcal{F}$ 是有限展示 $\mathcal{O}_X$-模，
当且仅当 $f_*\mathcal{F}$ 是有限展示
$\mathcal{O}_Y$-模。
\end{lemma}

\begin{proof}
由于 $f$ 有限，它是可表的。选择概形 $V$ 以及满平展态射
$V \to Y$。于是 $U = V \times_Y X$ 是概形，
它有一个到 $X$ 的满平展态射以及一个有限态射
$\psi : U \to V$（即 $f$ 的基变换）。由于
$\psi_*(\mathcal{F}|_U) = f_*\mathcal{F}|_V$，
该引理的结论立即由概形版本
《下降》引理
\ref{descent-lemma-finite-finitely-presented-module}
得出。
% P09-ERR-0345
\end{proof}





\section{Fpqc 覆盖}
\label{spaces-descent-section-fpqc}

\noindent
本节是《下降》第
\ref{descent-section-fpqc-universal-effective-epimorphisms} 节
的类似版本。目前我们尚不知道，关于概形的 fpqc 覆盖的全部材料
是否也对代数空间成立。

\begin{lemma}
\label{spaces-descent-lemma-open-fpqc-covering}
设 $S$ 为概形。
设 $\{f_i : T_i \to T\}_{i \in I}$ 为 $S$ 上代数空间的
一个 fpqc 覆盖。
假设对每个 $i$ 都有开子空间 $W_i \subset T_i$，
使得对所有 $i, j \in I$，有
$\text{pr}_0^{-1}(W_i) = \text{pr}_1^{-1}(W_j)$，
它们作为 $T_i \times_T T_j$ 的开子空间相等。
则存在唯一的开子空间 $W \subset T$，使
$W_i = f_i^{-1}(W)$ 对每个 $i$ 都成立。
\end{lemma}

\begin{proof}
由《空间上的拓扑》引理
\ref{spaces-topologies-lemma-refine-fpqc-schemes}，
可以假设每个 $T_i$ 都是概形。
选择概形 $U$ 以及满平展态射 $U \to T$。
于是 $\{T_i \times_T U \to U\}$ 是 $U$ 的一个 fpqc 覆盖，
并且 $T_i \times_T U$ 对每个 $i$ 都是概形。
因此，由《下降》引理
\ref{descent-lemma-open-fpqc-covering}，
开集族 $W_i \times_T U$ 来自唯一的开子概形
$W' \subset U$。
由于 $U \to X$ 是开映射，我们可以把 $W \subset X$
定义为 $W'$ 的像这个 Zariski 开集，见
% P09-ERR-0346
《空间的性质》第 \ref{spaces-properties-section-points} 节。
我们略去验证这一构造确实有效，即 $W_i$
是 $W$ 的逆像，对每个 $i$。
\end{proof}

\begin{lemma}
\label{spaces-descent-lemma-fpqc-universal-effective-epimorphisms}
设 $S$ 为概形。设 $\{T_i \to T\}$ 为 $S$ 上代数空间的
一个 fpqc 覆盖，见《空间上的拓扑》定义
\ref{spaces-topologies-definition-fpqc-covering}。
则给定代数空间 $B$ 位于 $S$ 上，序列
$$
\xymatrix{
\Mor_S(T, B) \ar[r] &
\prod\nolimits_i \Mor_S(T_i, B) \ar@<1ex>[r] \ar@<-1ex>[r] &
\prod\nolimits_{i, j} \Mor_S(T_i \times_T T_j, B)
}
$$
是等化子图。
换言之，$S$ 上代数空间范畴中的每个可表函子
都对 fpqc 覆盖满足层条件。
\end{lemma}

\begin{proof}
我们知道，当 $\{T_i \to T\}$ 是概形的 fpqc 覆盖时，
结论成立，见《空间的性质》命题
\ref{spaces-properties-proposition-sheaf-fpqc}。
这是关键事实；我们鼓励读者略过证明中其余的形式部分。
选择概形 $U$ 以及满平展态射
$U \to T$。设 $U_i$ 为概形，并设
$U_i \to T_i \times_T U$ 为满平展态射。
则 $\{U_i \to U\}$ 是一个 fpqc 覆盖。
这由《空间上的拓扑》引理
\ref{spaces-topologies-lemma-fpqc} 和
\ref{spaces-topologies-lemma-recognize-fpqc-covering} 得出。
由上可知，对 $\{U_i \to U\}$，该结论成立。

\medskip\noindent
这意味着如下内容：假设 $b_i : T_i \to B$ 是一个态射族，
并且有 $b_i \circ \text{pr}_0 = b_j \circ \text{pr}_1$，
它们作为态射 $T_i \times_T T_j \to B$ 相等。
令 $a_i : U_i \to B$ 为 $U_i \to T_i$ 与 $b_i$ 的复合。
由上可得唯一的态射 $a : U \to B$，使得 $a_i$
是 $a$ 与 $U_i \to U$ 的复合。
唯一性保证 $a \circ \text{pr}_0 = a \circ \text{pr}_1$，
它们作为态射 $U \times_T U \to B$ 相等。
于是，由于作为层有 $T = U/(U \times_T U)$，
我们发现 $a$ 来自唯一的态射 $b : T \to B$。
追逐图表可见，$b$ 正是所求态射。
\end{proof}

\section{态射的有限性与光滑性性质的下降}
\label{spaces-descent-section-descent-finiteness-morphisms}

\noindent
下面这类引理有时很有用。

\begin{lemma}
\label{spaces-descent-lemma-curiosity}
设 $S$ 为概形。设 $X \to Y \to Z$ 为代数空间的态射。
% P09-ERR-0347
设 $P$ 为 $S$ 上代数空间态射的下列性质之一：
平坦、局部有限型、局部有限展示。
假设 $X \to Z$ 具有 $P$，并且
$X \to Y$ 是 $(\Sch/S)_{fppf}$ 上层的满射。
则 $Y \to Z$ 具有 $P$。
\end{lemma}

\begin{proof}
选择概形 $W$ 以及满平展态射 $W \to Z$。
选择概形 $V$ 以及满平展态射 $V \to W \times_Z Y$。
选择概形 $U$ 以及满平展态射 $U \to V \times_Y X$。
由假设，可以找到一个 fppf 覆盖 $\{V_i \to V\}$
以及提升 $V_i \to X$，它们提升态射 $V_i \to Y$。
由于 $U \to X$ 满平展，可见在 fppf 覆盖
$\{V_i \times_X U \to V\}$ 的各成员上存在到 $U$ 的提升。
因此 $U \to V$ 诱导 $(\Sch/S)_{fppf}$ 上层的一个满射。
按照代数空间态射具有性质 $P$ 的定义（见
《空间的态射》定义 \ref{spaces-morphisms-definition-flat}、
定义 \ref{spaces-morphisms-definition-locally-finite-type} 以及
定义 \ref{spaces-morphisms-definition-locally-finite-presentation}），
可见 $U \to W$ 具有 $P$，而我们必须证明 $V \to W$
具有 $P$。于是问题约化为概形态射的情形；该情形在
《下降》引理 \ref{descent-lemma-curiosity} 中处理。
\end{proof}

\noindent
上述引理更为标准的一个情形如下。
（其中“平坦”的版本由《空间的态射》引理
\ref{spaces-morphisms-lemma-flat-permanence} 得出。）

\begin{lemma}
\label{spaces-descent-lemma-flat-finitely-presented-permanence}
设 $S$ 为概形。设
$$
\xymatrix{
X \ar[rr]_f \ar[rd]_p & &
Y \ar[dl]^q \\
& B
}
$$
为 $S$ 上代数空间态射的一个交换图。
假设 $f$ 满、平坦且局部有限展示，并假设 $p$
局部有限展示（相应地，局部有限型）。
则 $q$ 局部有限展示（相应地，局部有限型）。
\end{lemma}

\begin{proof}
由于 $\{X \to Y\}$ 是一个 fppf 覆盖，它诱导 fppf 层的满射
（《空间上的拓扑》引理
\ref{spaces-topologies-lemma-fppf-covering-surjective}），
而本引理是引理 \ref{spaces-descent-lemma-curiosity} 的特例。
另一方面，一个更容易的论证是由概形的类似结果推出它。
具体而言，该问题在 $B$ 与 $Y$ 上都是平展局部的
（《空间的态射》引理
\ref{spaces-morphisms-lemma-finite-type-local} 和
\ref{spaces-morphisms-lemma-finite-presentation-local}）。
因此可以假设 $B$ 与 $Y$ 是仿射概形。
由于 $|X| \to |Y|$ 是开映射
（《空间的态射》引理
\ref{spaces-morphisms-lemma-fppf-open}），
可以选择仿射概形 $U$ 以及平展态射 $U \to X$，
使复合 $U \to Y$ 是满射。
在这种情况下，结论由《下降》引理
\ref{descent-lemma-flat-finitely-presented-permanence} 得出。
\end{proof}

\begin{lemma}
\label{spaces-descent-lemma-syntomic-smooth-etale-permanence}
设 $S$ 为概形。设
$$
\xymatrix{
X \ar[rr]_f \ar[rd]_p & &
Y \ar[dl]^q \\
& B
}
$$
为 $S$ 上代数空间态射的一个交换图。
假设
\begin{enumerate}
\item $f$ 满且为 syntomic（相应地，光滑；相应地，平展）；
\item $p$ 为 syntomic（相应地，光滑；相应地，平展）。
\end{enumerate}
则 $q$ 为 syntomic（相应地，光滑；相应地，平展）。
\end{lemma}

\begin{proof}
我们由概形的类似结果推出它。
具体而言，该问题在 $B$ 与 $Y$ 上都是平展局部的
（《空间的态射》引理
\ref{spaces-morphisms-lemma-syntomic-local}、
\ref{spaces-morphisms-lemma-smooth-local} 以及
\ref{spaces-morphisms-lemma-etale-local}）。
因此可以假设 $B$ 与 $Y$ 是仿射概形。
由于 $|X| \to |Y|$ 是开映射
（《空间的态射》引理
\ref{spaces-morphisms-lemma-fppf-open}），
可以选择仿射概形 $U$ 以及平展态射 $U \to X$，
使复合 $U \to Y$ 是满射。
在这种情况下，结论由《下降》引理
\ref{descent-lemma-syntomic-smooth-etale-permanence} 得出。
\end{proof}

\noindent
事实上，可以如下加强这一结果。

\begin{lemma}
\label{spaces-descent-lemma-smooth-permanence}
设 $S$ 为概形。设
$$
\xymatrix{
X \ar[rr]_f \ar[rd]_p & &
Y \ar[dl]^q \\
& B
}
$$
为 $S$ 上代数空间态射的一个交换图。假设
\begin{enumerate}
\item $f$ 满、平坦且局部有限展示；
\item $p$ 光滑（相应地，平展）。
\end{enumerate}
则 $q$ 光滑（相应地，平展）。
\end{lemma}

\begin{proof}
我们由概形的类似结果推出它。
具体而言，该问题在 $B$ 与 $Y$ 上都是平展局部的
（《空间的态射》引理
\ref{spaces-morphisms-lemma-smooth-local} 和
\ref{spaces-morphisms-lemma-etale-local}）。
因此可以假设 $B$ 与 $Y$ 是仿射概形。
由于 $|X| \to |Y|$ 是开映射
（《空间的态射》引理
\ref{spaces-morphisms-lemma-fppf-open}），
可以选择仿射概形 $U$ 以及平展态射 $U \to X$，
使复合 $U \to Y$ 是满射。
在这种情况下，结论由《下降》引理
\ref{descent-lemma-smooth-permanence} 得出。
\end{proof}

\begin{lemma}
\label{spaces-descent-lemma-syntomic-permanence}
设 $S$ 为概形。设
$$
\xymatrix{
X \ar[rr]_f \ar[rd]_p & &
Y \ar[dl]^q \\
& B
}
$$
为 $S$ 上代数空间态射的一个交换图。假设
\begin{enumerate}
\item $f$ 满、平坦且局部有限展示；
\item $p$ 为 syntomic。
\end{enumerate}
则 $q$ 与 $f$ 都是 syntomic。
\end{lemma}

\begin{proof}
我们由概形的类似结果推出它。
具体而言，该问题在 $B$ 与 $Y$ 上都是平展局部的
（《空间的态射》引理
\ref{spaces-morphisms-lemma-syntomic-local}）。
因此可以假设 $B$ 与 $Y$ 是仿射概形。
由于 $|X| \to |Y|$ 是开映射
（《空间的态射》引理
\ref{spaces-morphisms-lemma-fppf-open}），
可以选择仿射概形 $U$ 以及平展态射 $U \to X$，
使复合 $U \to Y$ 是满射。
在这种情况下，结论由《下降》引理
\ref{descent-lemma-syntomic-permanence} 得出。
\end{proof}

\section{空间性质的下降}
\label{spaces-descent-section-descending-properties-spaces}

\noindent
本节收录一些如下类型的结果。

\begin{lemma}
\label{spaces-descent-lemma-descend-unibranch}
设 $S$ 为概形。
设 $f : X \to Y$ 为 $S$ 上代数空间的一个态射。
设 $x \in |X|$。
若 $f$ 在 $x$ 处平坦，且 $X$ 在 $x$ 处几何单支，
则 $Y$ 在 $f(x)$ 处几何单支。
\end{lemma}

\begin{proof}
考虑平展局部环之间的映射
$\mathcal{O}_{Y, f(\overline{x})} \to \mathcal{O}_{X, \overline{x}}$。
由《空间的态射》引理
\ref{spaces-morphisms-lemma-flat-at-point-etale-local-rings}，
该映射平坦。因此，若 $\mathcal{O}_{X, \overline{x}}$
有唯一的极小素理想，则 $\mathcal{O}_{Y, f(\overline{x})}$
也有唯一的极小素理想（由降链性质，见
《代数》引理 \ref{algebra-lemma-flat-going-down}）。
\end{proof}

\begin{lemma}
\label{spaces-descent-lemma-descend-reduced}
\begin{slogan}
一个源约化的平坦满代数空间态射具有约化的靶。
\end{slogan}
设 $S$ 为概形。
设 $f : X \to Y$ 为 $S$ 上代数空间的一个态射。
若 $f$ 平坦且满，并且 $X$ 约化，则 $Y$ 约化。
\end{lemma}

\begin{proof}
选择概形 $V$ 以及满平展态射 $V \to Y$。
选择概形 $U$ 以及满平展态射
$U \to X \times_Y V$。由于 $f$ 满且平坦，概形态射
$U \to V$ 满且平坦。这样便把问题约化为概形的情形
（因为 $X$ 与 $Y$ 的约化性是用 $U$ 与 $V$ 的约化性定义的，
见《空间的性质》第
\ref{spaces-properties-section-types-properties} 节）。
概形的情形是《下降》引理
\ref{descent-lemma-descend-reduced}。
\end{proof}

\begin{lemma}
\label{spaces-descent-lemma-descend-locally-Noetherian}
设 $f : X \to Y$ 为代数空间的一个态射。
若 $f$ 局部有限展示、平坦且满，并且 $X$ 局部 Noether，
则 $Y$ 局部 Noether。
\end{lemma}

\begin{proof}
选择概形 $V$ 以及满平展态射 $V \to Y$。
选择概形 $U$ 以及满平展态射
$U \to X \times_Y V$。由于 $f$ 满、平坦且局部有限展示，
概形态射 $U \to V$ 满、平坦且局部有限展示。
这样便把问题约化为概形的情形（因为 $X$ 与 $Y$
局部 Noether 是用 $U$ 与 $V$ 局部 Noether 定义的，见
《空间的性质》第
\ref{spaces-properties-section-types-properties} 节）。
在概形情形下，结论由《下降》引理
\ref{descent-lemma-Noetherian-local-fppf} 得出。
\end{proof}

\begin{lemma}
\label{spaces-descent-lemma-descend-regular}
设 $f : X \to Y$ 为代数空间的一个态射。
若 $f$ 局部有限展示、平坦且满，并且 $X$ 正则，
则 $Y$ 正则。
\end{lemma}

\begin{proof}
由引理 \ref{spaces-descent-lemma-descend-locally-Noetherian}，
我们知道 $Y$ 局部 Noether。
选择概形 $V$ 以及满平展态射 $V \to Y$。
只须证明 $V$ 的所有局部环都是正则局部环，见
《性质》引理 \ref{properties-lemma-characterize-regular}。
选择概形 $U$ 以及满平展态射
$U \to X \times_Y V$。由于 $f$ 满且平坦，概形态射
$U \to V$ 满且平坦。由假设，$U$ 是正则概形，
尤其其所有局部环都正则（由上述引理）。
因此，本引理由《代数》引理
\ref{algebra-lemma-flat-under-regular} 得出。
\end{proof}

\begin{lemma}
\label{spaces-descent-lemma-reduced-local-smooth}
设 $f : X \to Y$ 为代数空间的一个光滑态射。
若 $Y$ 约化，则 $X$ 约化。若 $f$ 满且 $X$ 约化，
则 $Y$ 约化。
\end{lemma}

\begin{proof}
选择交换图
$$
\xymatrix{
U \ar[d] \ar[r] & V \ar[d] \\
X \ar[r] & Y
}
$$
其中 $U$ 与 $V$ 是概形，竖直箭头满且平展，并且
$U \to X \times_Y V$ 满平展。
按照《空间的性质》第
\ref{spaces-properties-section-types-properties} 节中
约化代数空间的定义，$X$ 是约化代数空间当且仅当
$U$ 是约化概形。对 $Y$ 与 $V$ 亦然。
态射 $U \to V$ 是概形的光滑态射，见
《空间的态射》引理
\ref{spaces-morphisms-lemma-smooth-local}。
由于对概形而言，约化性对光滑拓扑是局部的
（《下降》引理 \ref{descent-lemma-reduced-local-smooth}），
可见 $U$ 约化，如果 $V$ 约化。
另一方面，若 $X \to Y$ 满，则 $U \to V$ 满；
在这种情况下，若 $U$ 约化，则 $V$ 约化。
\end{proof}

\section{态射性质的下降}
\label{spaces-descent-section-descending-properties-morphisms}

\noindent
本节引入代数空间态射的一个性质何时在某个拓扑中关于靶是局部的
这一概念。请与《下降》第
\ref{descent-section-descending-properties-morphisms} 节比较。

\begin{definition}
\label{spaces-descent-definition-property-morphisms-local}
设 $S$ 为概形。
设 $\mathcal{P}$ 为 $S$ 上代数空间态射的一个性质。
设 $\tau \in \{fpqc, \allowbreak fppf, \allowbreak syntomic, \allowbreak smooth, \allowbreak \etale\}$。
我们称 $\mathcal{P}$ {\it 关于基是 $\tau$-局部的}，
或{\it 关于靶是 $\tau$-局部的}，或
{\it 对 $\tau$-拓扑关于基是局部的}，如果对代数空间的任意
$\tau$-覆盖
$\{Y_i \to Y\}_{i \in I}$ 以及代数空间的任意态射
$f : X \to Y$，都有
$$
f \text{ 具有 }\mathcal{P}
\Leftrightarrow
\text{每个 }Y_i \times_Y X \to Y_i\text{ 具有 }\mathcal{P}.
$$
\end{definition}

\noindent
为免疑义，由于同构总是覆盖，我们看到（或要求）
性质 $\mathcal{P}$ 对 $X \to Y$ 成立，当且仅当它对
任意箭头 $X' \to Y'$ 成立，只要该箭头与 $X \to Y$ 同构。
若一个性质关于靶是 $\tau$-局部的，则它在沿
$\tau$-覆盖中出现的态射作基变换时保持。
下面是形式化陈述。

\begin{lemma}
\label{spaces-descent-lemma-pullback-property-local-target}
设 $S$ 为概形。
设 $\tau \in \{fpqc, fppf, syntomic, smooth, \etale\}$。
设 $\mathcal{P}$ 为 $S$ 上代数空间态射的一个性质，
并且关于靶是 $\tau$-局部的。设 $f : X \to Y$
具有性质 $\mathcal{P}$。
对任意平坦、相应地平坦且局部有限展示、
相应地 syntomic、相应地平展的态射 $Y' \to Y$，
基变换 $f' : Y' \times_Y X \to Y'$，即 $f$ 的基变换，具有性质
$\mathcal{P}$。
% P09-ERR-0348
\end{lemma}

\begin{proof}
这是成立的，因为可以把 $Y' \to Y$ 放入一个构成
$\tau$-覆盖的态射族中。
\end{proof}

\noindent
上面的一个简单而常用的推论是：若
$f : X \to Y$ 具有性质 $\mathcal{P}$，该性质关于靶
是 $\tau$-局部的，并且 $f(X) \subset V$，
其中 $V \subset Y$ 是某个开子空间，则所诱导的态射 $X \to V$
也具有 $\mathcal{P}$。
% P09-ERR-0349
证明：对 $f$ 沿 $V \to Y$ 作基变换即得 $X \to V$。

\begin{lemma}
\label{spaces-descent-lemma-largest-open-of-the-base}
设 $S$ 为概形。
设 $\tau \in \{fppf, syntomic, smooth, \etale\}$。
设 $\mathcal{P}$ 为 $S$ 上代数空间态射的一个性质，
并且关于靶是 $\tau$-局部的。
对代数空间的任意态射 $f : X \to Y$ 位于 $S$ 上，
存在最大的开子空间 $W(f) \subset Y$，使限制
$X_{W(f)} \to W(f)$ 具有 $\mathcal{P}$。此外：
\begin{enumerate}
\item 若 $g : Y' \to Y$ 是平坦且局部有限展示、
syntomic、光滑或平展的代数空间态射，并且基变换
$f' : X_{Y'} \to Y'$ 具有 $\mathcal{P}$，则
$g$ 经 $W(f)$ 分解；
\item 若 $g : Y' \to Y$ 平坦且局部有限展示、
syntomic、光滑或平展，则 $W(f') = g^{-1}(W(f))$；并且
\item 若 $\{g_i : Y_i \to Y\}$ 是一个 $\tau$-覆盖，则
$g_i^{-1}(W(f)) = W(f_i)$，其中 $f_i$ 是 $f$
沿 $Y_i \to Y$ 的基变换。
\end{enumerate}
\end{lemma}

\begin{proof}
考虑 $W_{set} \subset |Y|$，它是像
$g(|Y'|) \subset |Y|$ 的并，其中这些像来自满足下列性质的态射
$g : Y' \to Y$：
\begin{enumerate}
\item $g$ 平坦且局部有限展示、syntomic、光滑或平展；并且
\item 基变换 $Y' \times_{g, Y} X \to Y'$ 具有性质
$\mathcal{P}$。
\end{enumerate}
由于这样的态射 $g$ 是开映射（见《空间的态射》引理
\ref{spaces-morphisms-lemma-fppf-open}），
可见 $W_{set}$ 是 $|Y|$ 的开子集。
把开子空间记作 $W \subset Y$，其底点集为 $W_{set}$，见
% P09-ERR-0350
《空间的性质》引理
\ref{spaces-properties-lemma-open-subspaces}。
由于 $\mathcal{P}$ 对 $\tau$ 拓扑是局部的，限制
$X_W \to W$ 具有性质 $\mathcal{P}$，因为我们有覆盖
$\{Y' \to W\}$，它是 $W$ 的一个覆盖，使各拉回具有
$\mathcal{P}$。
这证明了存在性，也证明 $W(f)$ 具有性质 (1)。
为看出性质 (2)，注意
$W(f') \supset g^{-1}(W(f))$，因为 $\mathcal{P}$
在沿平坦且局部有限展示、syntomic、光滑或平展的态射作
基变换时保持，见引理
\ref{spaces-descent-lemma-pullback-property-local-target}。
另一方面，若 $Y'' \subset Y'$ 是一个开子空间，并且
$X_{Y''} \to Y''$ 具有性质 $\mathcal{P}$，则由构造
$Y'' \to Y$ 经 $W$ 分解，即
$Y'' \subset g^{-1}(W(f))$。这证明了 (2)。
断言 (3) 由 (2) 得出，因为每个态射 $Y_i \to Y$
按 $\tau$-覆盖的定义都平坦且局部有限展示、
syntomic、光滑或平展。
\end{proof}

\begin{lemma}
\label{spaces-descent-lemma-descending-properties-morphisms}
设 $S$ 为概形。设 $\mathcal{P}$ 为 $S$ 上
代数空间态射的一个性质。假设：
\begin{enumerate}
\item 若 $X_i \to Y_i$（$i = 1, 2$）具有性质
$\mathcal{P}$，则 $X_1 \amalg X_2 \to Y_1 \amalg Y_2$
也如此；
% P09-ERR-0351
\item 代数空间态射 $f : X \to Y$ 具有性质
$\mathcal{P}$，当且仅当对每个仿射概形 $Z$
与态射 $Z \to Y$，基变换
$Z \times_Y X \to Z$，即 $f$ 的基变换，具有性质
$\mathcal{P}$；并且
\item 对仿射概形的任意满平坦态射
$Z' \to Z$ 位于 $S$ 上，以及一个态射
$f : X \to Z$ 从代数空间到 $Z$，都有
$$
f' : Z' \times_Z X \to Z'\text{ 具有 }\mathcal{P}
\Rightarrow
f\text{ 具有 }\mathcal{P}.
$$
\end{enumerate}
则 $\mathcal{P}$ 关于基是 fpqc 局部的。
\end{lemma}

\begin{proof}
若 $\mathcal{P}$ 具有性质 (2)，则它自动在任意基变换下保持。
因此，定义 \ref{spaces-descent-definition-property-morphisms-local}
中的正向蕴含成立。

\medskip\noindent
设 $\{Y_i \to Y\}_{i \in I}$ 为 $S$ 上代数空间的
一个 fpqc 覆盖。
设 $f : X \to Y$ 为 $S$ 上代数空间的一个态射。
假设每个基变换 $f_i : Y_i \times_Y X \to Y_i$
都具有性质 $\mathcal{P}$。我们的目标是证明 $f$
具有 $\mathcal{P}$。
设 $Z$ 为仿射概形，并设 $Z \to Y$ 为一个态射。
由 (2)，只须证明代数空间态射
$Z \times_Y X \to Z$ 具有 $\mathcal{P}$。
由于 $\{Y_i \to Y\}_{i \in I}$ 是 fpqc 覆盖，
我们知道存在一个标准 fpqc 覆盖
$\{Z_j \to Z\}_{j = 1, \ldots , n}$，
并且存在态射 $Z_j \to Y_{i_j}$，它们位于 $Y$ 上，
其中 $i_j \in I$ 是适当的指标。
由于 $f_{i_j}$ 具有 $\mathcal{P}$，可见
$$
Z_j \times_Y X
=
Z_j \times_{Y_{i_j}} (Y_{i_j} \times_Y X)
\longrightarrow
Z_j
$$
具有 $\mathcal{P}$，因为它是 $f_{i_j}$ 的基变换
（见证明的第一条注记）。
令 $Z' = \coprod_{j = 1, \ldots, n} Z_j$，
于是 $Z' \to Z$ 是 $S$ 上仿射概形的平坦满态射。
由 (1)，我们得出 $Z' \times_Y X \to Z'$
具有性质 $\mathcal{P}$。
由于这是态射 $Z \times_Y X \to Z$ 沿态射
$Z' \to Z$ 的基变换，我们得出
$Z \times_Y X \to Z$ 具有性质 $\mathcal{P}$，如所需。
\end{proof}


\section{fpqc 拓扑中态射性质的下降}
\label{spaces-descent-section-descending-properties-morphisms-fpqc}


\noindent
本节给出代数空间态射的大量性质，它们在 fpqc 拓扑中
关于基是局部的。关于概形态射的情形，请比较
《下降》第 \ref{descent-section-descending-properties-morphisms-fpqc} 节。

\begin{lemma}
\label{spaces-descent-lemma-descending-property-quasi-compact}
设 $S$ 为概形。性质 $\mathcal{P}(f) =$``$f$ 是拟紧的''
关于 $S$ 上代数空间的基是 fpqc 局部的。
\end{lemma}

\begin{proof}
我们用引理 \ref{spaces-descent-lemma-descending-properties-morphisms} 证明。
该引理的假设 (1) 与 (2) 由《空间的态射》引理
\ref{spaces-morphisms-lemma-quasi-compact-local} 得出。
设 $Z' \to Z$ 为 $S$ 上仿射概形的满平坦态射。
设 $f : X \to Z$ 为代数空间的态射，并假设基变换
$f' : Z' \times_Z X \to Z'$ 拟紧。我们须证明 $f$ 拟紧。
再次使用《空间的态射》引理
\ref{spaces-morphisms-lemma-quasi-compact-local}，只须证明：
对每个仿射概形 $Y$ 与态射 $Y \to Z$，纤维积
$Y \times_Z X$ 是拟紧的。图示如下：
\begin{equation}
\label{spaces-descent-equation-cube}
\vcenter{
\xymatrix{
Y \times_Z Z' \times_Z X \ar[dd] \ar[rr] \ar[rd] & &
Z' \times_Z X \ar'[d][dd]^{f'} \ar[rd] \\
& Y \times_Z X \ar[dd] \ar[rr] & & X \ar[dd]^f \\
Y \times_Z Z' \ar'[r][rr] \ar[rd] & & Z' \ar[rd] \\
& Y \ar[rr] & & Z
}
}
\end{equation}
注意所有方块均为笛卡尔方块，且底部方块由仿射概形组成。
$f'$ 拟紧这一假设，结合 $Y \times_Z Z'$ 为仿射的事实，说明
$Y \times_Z Z' \times_Z X$ 拟紧。由于
$$
Y \times_Z Z' \times_Z X \longrightarrow Y \times_Z X
$$
是 $Z' \to Z$ 的基变换，因而是满的，所以
$Y \times_Z X$ 拟紧；见《空间的态射》引理
\ref{spaces-morphisms-lemma-surjection-from-quasi-compact}。
证明完毕。
\end{proof}

\begin{lemma}
\label{spaces-descent-lemma-descending-property-quasi-separated}
设 $S$ 为概形。性质 $\mathcal{P}(f) =$``$f$ 是拟分离的''
关于 $S$ 上代数空间的基是 fpqc 局部的。
\end{lemma}

\begin{proof}
拟分离态射的任意基变换仍拟分离；见《空间的态射》引理
\ref{spaces-morphisms-lemma-base-change-separated}。
因此，定义 \ref{spaces-descent-definition-property-morphisms-local} 中的正向蕴含成立。

\medskip\noindent
设 $\{Y_i \to Y\}_{i \in I}$ 为 $S$ 上代数空间的 fpqc 覆盖。
设 $f : X \to Y$ 为 $S$ 上代数空间的一个态射。
假设每个基变换 $X_i := Y_i \times_Y X \to Y_i$ 都拟分离。
这意味着每个态射
$$
\Delta_i :
X_i
\longrightarrow
X_i \times_{Y_i} X_i = Y_i \times_Y (X \times_Y X)
$$
均拟紧。fpqc 覆盖的基变换仍为 fpqc 覆盖；见
% P09-ERR-0352
《空间上的拓扑》引理 \ref{spaces-topologies-lemma-fpqc}，
故 $\{Y_i \times_Y (X \times_Y X) \to X \times_Y X\}$
是代数空间的 fpqc 覆盖。此外，每个 $\Delta_i$ 都是态射
$\Delta : X \to X \times_Y X$ 的基变换。因此，由引理
\ref{spaces-descent-lemma-descending-property-quasi-compact}，$\Delta$ 拟紧，
即 $f$ 拟分离。
\end{proof}

\begin{lemma}
\label{spaces-descent-lemma-descending-property-universally-closed}
设 $S$ 为概形。性质 $\mathcal{P}(f) =$``$f$ 是普遍闭的''
关于 $S$ 上代数空间的基是 fpqc 局部的。
\end{lemma}

\begin{proof}
我们用引理 \ref{spaces-descent-lemma-descending-properties-morphisms} 证明。
该引理的假设 (1) 与 (2) 由《空间的态射》引理
\ref{spaces-morphisms-lemma-universally-closed-local} 得出。
设 $Z' \to Z$ 为 $S$ 上仿射概形的满平坦态射。
设 $f : X \to Z$ 为代数空间的态射，并假设基变换
$f' : Z' \times_Z X \to Z'$ 普遍闭。我们须证明 $f$ 普遍闭。
再次使用《空间的态射》引理
\ref{spaces-morphisms-lemma-universally-closed-local}，只须证明：
对每个仿射概形 $Y$ 与态射 $Y \to Z$，映射
$|Y \times_Z X| \to |Y|$ 是闭的。
考虑立方图 (\ref{spaces-descent-equation-cube})。
$f'$ 普遍闭这一假设意味着
$|Y \times_Z Z' \times_Z X| \to |Y \times_Z Z'|$ 是闭的。
由于 $Y \times_Z Z' \to Y$ 是 $Z' \to Z$ 的基变换，
它拟紧、满且平坦，所以映射
$|Y \times_Z Z'| \to |Y|$ 是商映射；见《态射》引理
\ref{morphisms-lemma-fpqc-quotient-topology}。此外，映射
$$
|Y \times_Z Z' \times_Z X|
\longrightarrow
|Y \times_Z Z'| \times_{|Y|} |Y \times_Z X|
$$
是满的；见《空间的性质》引理
\ref{spaces-properties-lemma-points-cartesian}。
由初等拓扑可知 $|Y \times_Z X| \to |Y|$ 是闭的。
\end{proof}

\begin{lemma}
\label{spaces-descent-lemma-descending-property-universally-open}
设 $S$ 为概形。性质 $\mathcal{P}(f) =$``$f$ 是普遍开的''
关于 $S$ 上代数空间的基是 fpqc 局部的。
\end{lemma}

\begin{proof}
证明与引理 \ref{spaces-descent-lemma-descending-property-universally-closed} 的证明相同。
\end{proof}

\begin{lemma}
\label{spaces-descent-lemma-descending-property-universally-submersive}
性质 $\mathcal{P}(f) =$``$f$ 是普遍商映的''
关于基是 fpqc 局部的。
\end{lemma}

\begin{proof}
证明与引理 \ref{spaces-descent-lemma-descending-property-universally-closed} 的证明相同。
\end{proof}

\begin{lemma}
\label{spaces-descent-lemma-descending-property-surjective}
性质 $\mathcal{P}(f) =$``$f$ 是满的'' 关于基是 fpqc 局部的。
\end{lemma}

\begin{proof}
略。（提示：使用《空间的性质》引理
\ref{spaces-properties-lemma-points-cartesian}。）
\end{proof}

\begin{lemma}
\label{spaces-descent-lemma-descending-property-universally-injective}
性质 $\mathcal{P}(f) =$``$f$ 是普遍单的''
关于基是 fpqc 局部的。
\end{lemma}

\begin{proof}
我们用引理 \ref{spaces-descent-lemma-descending-properties-morphisms} 证明。
该引理的假设 (1) 与 (2) 由《空间的态射》引理
% P09-ERR-0353
\ref{spaces-morphisms-lemma-universally-closed-local} 得出。
设 $Z' \to Z$ 为 $S$ 上仿射概形的平坦满态射，并设
$f : X \to Z$ 为从代数空间到 $Z$ 的态射。
假设基变换 $f' : X' \to Z'$ 普遍单。
设 $K$ 为域，并设 $a, b : \Spec(K) \to X$ 为两个态射，
使得 $f \circ a = f \circ b$。由于 $Z' \to Z$ 是满的，
存在域扩张 $K'/K$ 与态射 $\Spec(K') \to Z'$，
使下列实线图交换：
$$
\xymatrix{
\Spec(K') \ar[rrd] \ar@{-->}[rd]_{a', b'} \ar[dd] \\
 &
X' \ar[r] \ar[d] &
Z' \ar[d] \\
\Spec(K) \ar[r]^{a, b} &
X \ar[r] &
Z
}
$$
由于方块是笛卡尔的，得到使图交换的两条虚线箭头
$a'$、$b'$。因为 $X' \to Z'$ 普遍单，所以 $a' = b'$。
这迫使 $a = b$，因为 $\{\Spec(K') \to \Spec(K)\}$
是 fpqc 覆盖；见《空间的性质》命题
\ref{spaces-properties-proposition-sheaf-fpqc}。
因此 $f$ 普遍单，如所需。
\end{proof}

\begin{lemma}
\label{spaces-descent-lemma-descending-property-universal-homeomorphism}
性质 $\mathcal{P}(f) =$``$f$ 是普遍同胚''
关于基是 fpqc 局部的。
\end{lemma}

\begin{proof}
这可按与引理 \ref{spaces-descent-lemma-descending-property-universally-closed}
完全相同的方法证明。也可使用如下事实：拓扑空间之间的映射是同胚，
当且仅当它是单的、满的且开的。因此，普遍同胚等价于一个满、
普遍单且普遍开的态射。见《空间的态射》引理
\ref{spaces-morphisms-lemma-base-change-surjective} 以及定义
\ref{spaces-morphisms-definition-universally-injective}、
\ref{spaces-morphisms-definition-surjective}、
\ref{spaces-morphisms-definition-open}、
\ref{spaces-morphisms-definition-universal-homeomorphism}。
故本引理由引理 \ref{spaces-descent-lemma-descending-property-surjective}、
\ref{spaces-descent-lemma-descending-property-universally-injective} 与
\ref{spaces-descent-lemma-descending-property-universally-open} 得出。
\end{proof}

\begin{lemma}
\label{spaces-descent-lemma-descending-property-locally-finite-type}
性质 $\mathcal{P}(f) =$``$f$ 是局部有限型的''
关于基是 fpqc 局部的。
\end{lemma}

\begin{proof}
我们用引理 \ref{spaces-descent-lemma-descending-properties-morphisms} 证明。
该引理的假设 (1) 与 (2) 由《空间的态射》引理
\ref{spaces-morphisms-lemma-finite-type-local} 得出。
设 $Z' \to Z$ 为 $S$ 上仿射概形的满平坦态射。
设 $f : X \to Z$ 为代数空间的态射，并假设基变换
$f' : Z' \times_Z X \to Z'$ 局部有限型。我们须证明 $f$
局部有限型。设 $U$ 为概形，并设 $U \to X$ 为满平展态射。
再次使用《空间的态射》引理
\ref{spaces-morphisms-lemma-finite-type-local}，只须证明
$U \to Z$ 局部有限型。由于 $f'$ 局部有限型，且
$Z' \times_Z U$ 是 $Z' \times_Z X$ 上平展的概形，
由同一引理可知 $Z' \times_Z U \to Z'$ 局部有限型。
由于 $\{Z' \to Z\}$ 是 fpqc 覆盖，由《下降》引理
\ref{descent-lemma-descending-property-locally-finite-type}，
$U \to Z$ 局部有限型，如所需。
\end{proof}

\begin{lemma}
\label{spaces-descent-lemma-descending-property-locally-finite-presentation}
性质 $\mathcal{P}(f) =$``$f$ 是局部有限展示的''
关于基是 fpqc 局部的。
\end{lemma}

\begin{proof}
我们用引理 \ref{spaces-descent-lemma-descending-properties-morphisms} 证明。
该引理的假设 (1) 与 (2) 由《空间的态射》引理
\ref{spaces-morphisms-lemma-finite-presentation-local} 得出。
设 $Z' \to Z$ 为 $S$ 上仿射概形的满平坦态射。
设 $f : X \to Z$ 为代数空间的态射，并假设基变换
$f' : Z' \times_Z X \to Z'$ 局部有限展示。我们须证明 $f$
局部有限展示。设 $U$ 为概形，并设 $U \to X$ 为满平展态射。
再次使用《空间的态射》引理
\ref{spaces-morphisms-lemma-finite-presentation-local}，只须证明
$U \to Z$ 局部有限展示。由于 $f'$ 局部有限展示，且
$Z' \times_Z U$ 是 $Z' \times_Z X$ 上平展的概形，
由同一引理可知 $Z' \times_Z U \to Z'$ 局部有限展示。
由于 $\{Z' \to Z\}$ 是 fpqc 覆盖，由《下降》引理
\ref{descent-lemma-descending-property-locally-finite-presentation}，
$U \to Z$ 局部有限展示，如所需。
\end{proof}

\begin{lemma}
\label{spaces-descent-lemma-descending-property-finite-type}
性质 $\mathcal{P}(f) =$``$f$ 是有限型的''
关于基是 fpqc 局部的。
\end{lemma}

\begin{proof}
合并引理 \ref{spaces-descent-lemma-descending-property-quasi-compact} 与
\ref{spaces-descent-lemma-descending-property-locally-finite-type}。
\end{proof}

\begin{lemma}
\label{spaces-descent-lemma-descending-property-finite-presentation}
性质 $\mathcal{P}(f) =$``$f$ 是有限展示的''
关于基是 fpqc 局部的。
\end{lemma}

\begin{proof}
合并引理 \ref{spaces-descent-lemma-descending-property-quasi-compact}、
\ref{spaces-descent-lemma-descending-property-quasi-separated} 与
\ref{spaces-descent-lemma-descending-property-locally-finite-presentation}。
\end{proof}

\begin{lemma}
\label{spaces-descent-lemma-descending-property-flat}
性质 $\mathcal{P}(f) =$``$f$ 是平坦的''
关于基是 fpqc 局部的。
\end{lemma}

\begin{proof}
我们用引理 \ref{spaces-descent-lemma-descending-properties-morphisms} 证明。
该引理的假设 (1) 与 (2) 由《空间的态射》引理
\ref{spaces-morphisms-lemma-flat-local} 得出。
设 $Z' \to Z$ 为 $S$ 上仿射概形的满平坦态射。
设 $f : X \to Z$ 为代数空间的态射，并假设基变换
$f' : Z' \times_Z X \to Z'$ 平坦。我们须证明 $f$ 平坦。
设 $U$ 为概形，并设 $U \to X$ 为满平展态射。
再次使用《空间的态射》引理
\ref{spaces-morphisms-lemma-flat-local}，只须证明 $U \to Z$ 平坦。
由于 $f'$ 平坦，且 $Z' \times_Z U$ 是 $Z' \times_Z X$
上平展的概形，由同一引理可知 $Z' \times_Z U \to Z'$ 平坦。
由于 $\{Z' \to Z\}$ 是 fpqc 覆盖，由《下降》引理
\ref{descent-lemma-descending-property-flat}，$U \to Z$ 平坦，如所需。
\end{proof}

\begin{lemma}
\label{spaces-descent-lemma-descending-property-open-immersion}
性质 $\mathcal{P}(f) =$``$f$ 是开浸入''
关于基是 fpqc 局部的。
\end{lemma}

\begin{proof}
我们用引理 \ref{spaces-descent-lemma-descending-properties-morphisms} 证明。
该引理的假设 (1) 与 (2) 由《空间的态射》引理
\ref{spaces-morphisms-lemma-closed-immersion-local} 得出。
考虑如下笛卡尔图
$$
\xymatrix{
X' \ar[r] \ar[d] & X \ar[d] \\
Z' \ar[r] & Z
}
$$
它由 $S$ 上代数空间组成，其中 $Z' \to Z$ 为仿射概形的满平坦态射，
$X' \to Z'$ 为开浸入。我们须证明 $X \to Z$ 是开浸入。
注意 $|X'| \subset |Z'|$ 对应于一个开子概形
$U' \subset Z'$（同构于 $X'$），且具有性质
$\text{pr}_0^{-1}(U') = \text{pr}_1^{-1}(U')$；
这是作为 $Z' \times_Z Z'$ 的开子概形的等式。
因此存在开子概形 $U \subset Z$，使得
$X' = (Z' \to Z)^{-1}(U)$；见《下降》引理
\ref{descent-lemma-open-fpqc-covering}。
由《空间的性质》命题
\ref{spaces-properties-proposition-sheaf-fpqc}，
$X$ 对 fpqc 拓扑满足层条件。现在有 fpqc 覆盖
$\mathcal{U} = \{U' \to U\}$ 与元素
$U' \to X' \to X \in \check{H}^0(\mathcal{U}, X)$。
由层条件，得到一个态射 $U \to X$，使得
$$
\xymatrix{
U' \ar[r] \ar[d]^{\cong} \ar@/_3ex/[dd] & U \ar[d] \ar@/^3ex/[dd] \\
X' \ar[r] \ar[d] & X \ar[d] \\
Z' \ar[r] & Z
}
$$
交换。另一方面，我们知道，对任意概形 $T$（它位于 $S$ 上）
与 $T$-值点 $T \to X$，复合 $T \to X \to Z$ 具有如下性质：
$Z' \times_Z T \to Z'$ 经 $U'$ 分解。显然，这意味着
$T \to Z$ 经 $U$ 分解。换言之，层的映射 $U \to X$
是双射，结论得证。
\end{proof}

\begin{lemma}
\label{spaces-descent-lemma-descending-property-isomorphism}
性质 $\mathcal{P}(f) =$``$f$ 是同构''
关于基是 fpqc 局部的。
\end{lemma}

\begin{proof}
合并引理 \ref{spaces-descent-lemma-descending-property-surjective} 与
\ref{spaces-descent-lemma-descending-property-open-immersion}。
\end{proof}

\begin{lemma}
\label{spaces-descent-lemma-descending-property-affine}
性质 $\mathcal{P}(f) =$``$f$ 是仿射的''
关于基是 fpqc 局部的。
\end{lemma}

\begin{proof}
我们用引理 \ref{spaces-descent-lemma-descending-properties-morphisms} 证明。
该引理的假设 (1) 与 (2) 由《空间的态射》引理
\ref{spaces-morphisms-lemma-affine-local} 得出。
设 $Z' \to Z$ 为 $S$ 上仿射概形的满平坦态射。
设 $f : X \to Z$ 为代数空间的态射，并假设基变换
$f' : Z' \times_Z X \to Z'$ 仿射。
令 $X'$ 为表示 $Z' \times_Z X$ 的概形。
得到典范同构
$$
\varphi : X' \times_Z Z' \longrightarrow Z' \times_Z X'
$$
因为两个概形都表示代数空间 $Z' \times_Z Z' \times_Z X$。
这是 $X'/Z'/Z$ 的一个下降资料；见《下降》定义
\ref{descent-definition-descent-datum}
（略去验证；比较《下降》引理
\ref{descent-lemma-descent-data-sheaves}）。
由于 $X' \to Z'$ 是仿射的，该下降资料有效；见《下降》引理
\ref{descent-lemma-affine}。因此存在概形 $Y \to Z$，
它位于 $Z$ 上，并存在与下降资料相容的同构
$\psi : Z' \times_Z Y \to X'$。当然 $Y \to Z$ 是仿射的
（由构造，或由《下降》引理
\ref{descent-lemma-descending-property-affine}）。
注意 $\mathcal{Y} = \{Z' \times_Z Y \to Y\}$ 是 fpqc 覆盖；
把 $\psi$ 解释为 $X(Z' \times_Z Y)$ 的元素，可见
$\psi \in \check{H}^0(\mathcal{Y}, X)$。
由 $X$ 关于该覆盖的层条件（见《空间的性质》命题
\ref{spaces-properties-proposition-sheaf-fpqc}），得到态射 $Y \to X$。
按构造，它到 $Z'$ 的基变换是同构，因而由引理
\ref{spaces-descent-lemma-descending-property-isomorphism}，它本身也是同构。
这证明 $X$ 可由仿射概形表示，结论得证。
\end{proof}

\begin{lemma}
\label{spaces-descent-lemma-descending-property-closed-immersion}
性质 $\mathcal{P}(f) =$``$f$ 是闭浸入''
关于基是 fpqc 局部的。
\end{lemma}

\begin{proof}
我们用引理 \ref{spaces-descent-lemma-descending-properties-morphisms} 证明。
该引理的假设 (1) 与 (2) 由《空间的态射》引理
\ref{spaces-morphisms-lemma-closed-immersion-local} 得出。
考虑如下笛卡尔图
$$
\xymatrix{
X' \ar[r] \ar[d] & X \ar[d] \\
Z' \ar[r] & Z
}
$$
它由 $S$ 上代数空间组成，其中 $Z' \to Z$ 为仿射概形的
满平坦态射，$X' \to Z'$ 为闭浸入。我们须证明 $X \to Z$
是闭浸入。态射 $X' \to Z'$ 是仿射的。因此，由引理
\ref{spaces-descent-lemma-descending-property-affine}，$X$ 是概形且
$X \to Z$ 是仿射的。由《下降》引理
\ref{descent-lemma-descending-property-closed-immersion}，
$X \to Z$ 是闭浸入，如所需。
\end{proof}

\begin{lemma}
\label{spaces-descent-lemma-descending-property-separated}
性质 $\mathcal{P}(f) =$``$f$ 是分离的''
关于基是 fpqc 局部的。
\end{lemma}

\begin{proof}
分离态射的任意基变换仍分离；见《空间的态射》引理
\ref{spaces-morphisms-lemma-base-change-separated}。
因此，定义 \ref{spaces-descent-definition-property-morphisms-local} 中的正向蕴含成立。

\medskip\noindent
设 $\{Y_i \to Y\}_{i \in I}$ 为 $S$ 上代数空间的 fpqc 覆盖。
设 $f : X \to Y$ 为 $S$ 上代数空间的一个态射。
假设每个基变换 $X_i := Y_i \times_Y X \to Y_i$ 都分离。
这意味着每个态射
$$
\Delta_i :
X_i
\longrightarrow
X_i \times_{Y_i} X_i = Y_i \times_Y (X \times_Y X)
$$
都是闭浸入。fpqc 覆盖的基变换仍为 fpqc 覆盖；见
《空间上的拓扑》引理 \ref{spaces-topologies-lemma-fpqc}，
故 $\{Y_i \times_Y (X \times_Y X) \to X \times_Y X\}$
是代数空间的 fpqc 覆盖。此外，每个 $\Delta_i$ 都是态射
$\Delta : X \to X \times_Y X$ 的基变换。因此，由引理
\ref{spaces-descent-lemma-descending-property-closed-immersion}，$\Delta$ 是闭浸入，
即 $f$ 分离。
\end{proof}

\begin{lemma}
\label{spaces-descent-lemma-descending-property-proper}
性质 $\mathcal{P}(f) =$``$f$ 是固有的''
关于基是 fpqc 局部的。
\end{lemma}

\begin{proof}
本引理由引理 \ref{spaces-descent-lemma-descending-property-universally-closed}、
\ref{spaces-descent-lemma-descending-property-separated} 与
\ref{spaces-descent-lemma-descending-property-finite-type} 合并得出。
\end{proof}

\begin{lemma}
\label{spaces-descent-lemma-descending-property-quasi-affine}
性质 $\mathcal{P}(f) =$``$f$ 是拟仿射的''
关于基是 fpqc 局部的。
\end{lemma}

\begin{proof}
我们用引理 \ref{spaces-descent-lemma-descending-properties-morphisms} 证明。
该引理的假设 (1) 与 (2) 由《空间的态射》引理
\ref{spaces-morphisms-lemma-quasi-affine-local} 得出。
设 $Z' \to Z$ 为 $S$ 上仿射概形的满平坦态射。
设 $f : X \to Z$ 为代数空间的态射，并假设基变换
$f' : Z' \times_Z X \to Z'$ 拟仿射。
令 $X'$ 为表示 $Z' \times_Z X$ 的概形。
得到典范同构
$$
\varphi : X' \times_Z Z' \longrightarrow Z' \times_Z X'
$$
因为两个概形都表示代数空间 $Z' \times_Z Z' \times_Z X$。
这是 $X'/Z'/Z$ 的一个下降资料；见《下降》定义
\ref{descent-definition-descent-datum}
（略去验证；比较《下降》引理
\ref{descent-lemma-descent-data-sheaves}）。
由于 $X' \to Z'$ 是拟仿射的，该下降资料有效；见《下降》引理
\ref{descent-lemma-quasi-affine}。因此存在概形 $Y \to Z$，
它位于 $Z$ 上，并存在与下降资料相容的同构
$\psi : Z' \times_Z Y \to X'$。当然 $Y \to Z$ 是拟仿射的
（由构造，或由《下降》引理
\ref{descent-lemma-descending-property-quasi-affine}）。
注意 $\mathcal{Y} = \{Z' \times_Z Y \to Y\}$ 是 fpqc 覆盖；
把 $\psi$ 解释为 $X(Z' \times_Z Y)$ 的元素，可见
$\psi \in \check{H}^0(\mathcal{Y}, X)$。
由 $X$ 的层条件（见《空间的性质》命题
\ref{spaces-properties-proposition-sheaf-fpqc}），得到态射 $Y \to X$。
按构造，它到 $Z'$ 的基变换是同构，因而由引理
\ref{spaces-descent-lemma-descending-property-isomorphism}，它本身也是同构。
这证明 $X$ 可由拟仿射概形表示，结论得证。
\end{proof}

\begin{lemma}
\label{spaces-descent-lemma-descending-property-quasi-compact-immersion}
性质 $\mathcal{P}(f) =$``$f$ 是拟紧浸入''
关于基是 fpqc 局部的。
\end{lemma}

\begin{proof}
我们用引理 \ref{spaces-descent-lemma-descending-properties-morphisms} 证明。
该引理的假设 (1) 与 (2) 由《空间的态射》引理
\ref{spaces-morphisms-lemma-closed-immersion-local} 与
\ref{spaces-morphisms-lemma-quasi-compact-local} 得出。
考虑如下笛卡尔图
$$
\xymatrix{
X' \ar[r] \ar[d] & X \ar[d] \\
Z' \ar[r] & Z
}
$$
它由 $S$ 上代数空间组成，其中 $Z' \to Z$ 为仿射概形的
满平坦态射，$X' \to Z'$ 为拟紧浸入。我们须证明 $X \to Z$
% P09-ERR-0354
是闭浸入。态射 $X' \to Z'$ 是拟仿射的。因此，由引理
\ref{spaces-descent-lemma-descending-property-quasi-affine}，$X$ 是概形且
$X \to Z$ 是拟仿射的。由《下降》引理
\ref{descent-lemma-descending-property-quasi-compact-immersion}，
$X \to Z$ 是拟紧浸入，如所需。
\end{proof}

\begin{lemma}
\label{spaces-descent-lemma-descending-property-integral}
性质 $\mathcal{P}(f) =$``$f$ 是整的''
关于基是 fpqc 局部的。
\end{lemma}

\begin{proof}
整态射等价于仿射且普遍闭的态射；见《空间的态射》引理
\ref{spaces-morphisms-lemma-integral-universally-closed}。
故本引理由引理 \ref{spaces-descent-lemma-descending-property-universally-closed}
与 \ref{spaces-descent-lemma-descending-property-affine} 合并得出。
\end{proof}

\begin{lemma}
\label{spaces-descent-lemma-descending-property-finite}
性质 $\mathcal{P}(f) =$``$f$ 是有限的''
关于基是 fpqc 局部的。
\end{lemma}

\begin{proof}
% P09-ERR-0355
有限态射等价于局部有限型的整态射；见《空间的态射》引理
\ref{spaces-morphisms-lemma-finite-integral}。
故本引理由引理 \ref{spaces-descent-lemma-descending-property-locally-finite-type}
与 \ref{spaces-descent-lemma-descending-property-integral} 合并得出。
\end{proof}

\begin{lemma}
\label{spaces-descent-lemma-descending-property-quasi-finite}
性质 $\mathcal{P}(f) =$``$f$ 是局部拟有限的''
与性质 $\mathcal{P}(f) =$``$f$ 是拟有限的''
都关于基是 fpqc 局部的。
\end{lemma}

\begin{proof}
我们已经看到``拟紧''关于基是 fpqc 局部的；见引理
\ref{spaces-descent-lemma-descending-property-quasi-compact}。因此，只须对
``局部拟有限''证明本引理。我们用引理
\ref{spaces-descent-lemma-descending-properties-morphisms} 证明。
该引理的假设 (1) 与 (2) 由《空间的态射》引理
\ref{spaces-morphisms-lemma-quasi-finite-local} 得出。
设 $Z' \to Z$ 为 $S$ 上仿射概形的满平坦态射。
设 $f : X \to Z$ 为代数空间的态射，并假设基变换
$f' : Z' \times_Z X \to Z'$ 局部拟有限。我们须证明 $f$
局部拟有限。设 $U$ 为概形，并设 $U \to X$ 为满平展态射。
再次使用《空间的态射》引理
\ref{spaces-morphisms-lemma-quasi-finite-local}，只须证明
$U \to Z$ 局部拟有限。由于 $f'$ 局部拟有限，且
$Z' \times_Z U$ 是 $Z' \times_Z X$ 上平展的概形，
由同一引理可知 $Z' \times_Z U \to Z'$ 局部拟有限。
由于 $\{Z' \to Z\}$ 是 fpqc 覆盖，由《下降》引理
\ref{descent-lemma-descending-property-quasi-finite}，
$U \to Z$ 局部拟有限，如所需。
\end{proof}

\begin{lemma}
\label{spaces-descent-lemma-descending-property-syntomic}
性质 $\mathcal{P}(f) =$``$f$ 是 syntomic 的''
关于基是 fpqc 局部的。
\end{lemma}

\begin{proof}
我们用引理 \ref{spaces-descent-lemma-descending-properties-morphisms} 证明。
该引理的假设 (1) 与 (2) 由《空间的态射》引理
\ref{spaces-morphisms-lemma-syntomic-local} 得出。
设 $Z' \to Z$ 为 $S$ 上仿射概形的满平坦态射。
设 $f : X \to Z$ 为代数空间的态射，并假设基变换
$f' : Z' \times_Z X \to Z'$ 是 syntomic 的。我们须证明
$f$ 是 syntomic 的。设 $U$ 为概形，并设 $U \to X$
为满平展态射。再次使用《空间的态射》引理
\ref{spaces-morphisms-lemma-syntomic-local}，只须证明
$U \to Z$ 是 syntomic 的。由于 $f'$ 是 syntomic 的，且
$Z' \times_Z U$ 是 $Z' \times_Z X$ 上平展的概形，
由同一引理可知 $Z' \times_Z U \to Z'$ 是 syntomic 的。
由于 $\{Z' \to Z\}$ 是 fpqc 覆盖，由《下降》引理
\ref{descent-lemma-descending-property-syntomic}，
$U \to Z$ 是 syntomic 的，如所需。
\end{proof}

\begin{lemma}
\label{spaces-descent-lemma-descending-property-smooth}
性质 $\mathcal{P}(f) =$``$f$ 是光滑的''
关于基是 fpqc 局部的。
\end{lemma}

\begin{proof}
我们用引理 \ref{spaces-descent-lemma-descending-properties-morphisms} 证明。
该引理的假设 (1) 与 (2) 由《空间的态射》引理
\ref{spaces-morphisms-lemma-smooth-local} 得出。
设 $Z' \to Z$ 为 $S$ 上仿射概形的满平坦态射。
设 $f : X \to Z$ 为代数空间的态射，并假设基变换
$f' : Z' \times_Z X \to Z'$ 光滑。我们须证明 $f$ 光滑。
设 $U$ 为概形，并设 $U \to X$ 为满平展态射。
再次使用《空间的态射》引理
\ref{spaces-morphisms-lemma-smooth-local}，只须证明 $U \to Z$ 光滑。
由于 $f'$ 光滑，且 $Z' \times_Z U$ 是 $Z' \times_Z X$
上平展的概形，由同一引理可知 $Z' \times_Z U \to Z'$ 光滑。
由于 $\{Z' \to Z\}$ 是 fpqc 覆盖，由《下降》引理
\ref{descent-lemma-descending-property-smooth}，$U \to Z$ 光滑，如所需。
\end{proof}

\begin{lemma}
\label{spaces-descent-lemma-descending-property-unramified}
性质 $\mathcal{P}(f) =$``$f$ 是非分歧的''
关于基是 fpqc 局部的。
\end{lemma}

\begin{proof}
我们用引理 \ref{spaces-descent-lemma-descending-properties-morphisms} 证明。
该引理的假设 (1) 与 (2) 由《空间的态射》引理
\ref{spaces-morphisms-lemma-unramified-local} 得出。
设 $Z' \to Z$ 为 $S$ 上仿射概形的满平坦态射。
设 $f : X \to Z$ 为代数空间的态射，并假设基变换
$f' : Z' \times_Z X \to Z'$ 非分歧。我们须证明 $f$ 非分歧。
设 $U$ 为概形，并设 $U \to X$ 为满平展态射。
再次使用《空间的态射》引理
\ref{spaces-morphisms-lemma-unramified-local}，只须证明
$U \to Z$ 非分歧。由于 $f'$ 非分歧，且 $Z' \times_Z U$
是 $Z' \times_Z X$ 上平展的概形，由同一引理可知
$Z' \times_Z U \to Z'$ 非分歧。由于 $\{Z' \to Z\}$
是 fpqc 覆盖，由《下降》引理
\ref{descent-lemma-descending-property-unramified}，
$U \to Z$ 非分歧，如所需。
\end{proof}

\begin{lemma}
\label{spaces-descent-lemma-descending-property-etale}
性质 $\mathcal{P}(f) =$``$f$ 是平展的''
关于基是 fpqc 局部的。
\end{lemma}

\begin{proof}
我们用引理 \ref{spaces-descent-lemma-descending-properties-morphisms} 证明。
该引理的假设 (1) 与 (2) 由《空间的态射》引理
\ref{spaces-morphisms-lemma-etale-local} 得出。
设 $Z' \to Z$ 为 $S$ 上仿射概形的满平坦态射。
设 $f : X \to Z$ 为代数空间的态射，并假设基变换
$f' : Z' \times_Z X \to Z'$ 平展。我们须证明 $f$ 平展。
设 $U$ 为概形，并设 $U \to X$ 为满平展态射。
再次使用《空间的态射》引理
\ref{spaces-morphisms-lemma-etale-local}，只须证明 $U \to Z$ 平展。
由于 $f'$ 平展，且 $Z' \times_Z U$ 是 $Z' \times_Z X$
上平展的概形，由同一引理可知 $Z' \times_Z U \to Z'$ 平展。
由于 $\{Z' \to Z\}$ 是 fpqc 覆盖，由《下降》引理
\ref{descent-lemma-descending-property-etale}，$U \to Z$ 平展，如所需。
\end{proof}

\begin{lemma}
\label{spaces-descent-lemma-descending-property-finite-locally-free}
性质 $\mathcal{P}(f) =$``$f$ 是有限局部自由的''
关于基是 fpqc 局部的。
\end{lemma}

\begin{proof}
有限局部自由等价于有限、平坦且局部有限展示
（《空间的态射》引理 \ref{spaces-morphisms-lemma-finite-flat}）。
故结论由引理 \ref{spaces-descent-lemma-descending-property-finite}、
\ref{spaces-descent-lemma-descending-property-flat} 与
\ref{spaces-descent-lemma-descending-property-locally-finite-presentation} 得出。
\end{proof}

\begin{lemma}
\label{spaces-descent-lemma-descending-property-monomorphism}
性质 $\mathcal{P}(f) =$``$f$ 是单态射''
关于基是 fpqc 局部的。
\end{lemma}

\begin{proof}
设 $f : X \to Y$ 为代数空间的态射。
设 $\{Y_i \to Y\}$ 为 fpqc 覆盖，并假设每个基变换
$f_i : X_i \to Y_i$（它是 $f$ 的基变换）都是单态射。
我们须证明 $f$ 是单态射。

\medskip\noindent
第一种证明。注意 $f$ 是单态射，当且仅当
$\Delta : X \to X \times_Y X$ 是同构。把这一点应用于
$f_i$，可见每个态射
$$
\Delta_i :
X_i
\longrightarrow
X_i \times_{Y_i} X_i = Y_i \times_Y (X \times_Y X)
$$
都是同构。fpqc 覆盖的基变换仍为 fpqc 覆盖；见
《空间上的拓扑》引理 \ref{spaces-topologies-lemma-fpqc}，
故 $\{Y_i \times_Y (X \times_Y X) \to X \times_Y X\}$
是代数空间的 fpqc 覆盖。此外，每个 $\Delta_i$ 都是态射
$\Delta : X \to X \times_Y X$ 的基变换。因此，由引理
\ref{spaces-descent-lemma-descending-property-isomorphism}，$\Delta$ 是同构，
即 $f$ 是单态射。

\medskip\noindent
第二种证明。
设 $V$ 为概形，并设 $V \to Y$ 为满平展态射。
若能证明 $V \times_Y X \to V$ 是单态射，则可推出
$X \to Y$ 是单态射。事实上，给定层的任意笛卡尔图
$$
\vcenter{
\xymatrix{
\mathcal{F} \ar[r]_a \ar[d]_b & \mathcal{G} \ar[d]^c \\
\mathcal{H} \ar[r]^d & \mathcal{I}
}
}
\quad
\quad
\mathcal{F} = \mathcal{H} \times_\mathcal{I} \mathcal{G}
$$
若 $c$ 是层的满射且 $a$ 是单射，则 $d$ 也是单射。
这把问题约化到 $Y$ 为概形的情形。此外，在这种情形下，
也可假设代数空间 $Y_i$ 都是概形，因为总能加细覆盖以置身于
这种情形；见《空间上的拓扑》引理
\ref{spaces-topologies-lemma-refine-fpqc-schemes}。

\medskip\noindent
假设 $\{Y_i \to Y\}$ 是概形的 fpqc 覆盖。
设 $a, b : T \to X$ 为两个态射，使得
$f \circ a = f \circ b$。我们须证明 $a = b$。
由于 $f_i$ 是单态射，可知 $a_i = b_i$，其中
$a_i, b_i : Y_i \times_Y T \to X_i$ 为基变换。
特别地，复合 $Y_i \times_Y T \to T \to X$ 相等。
由于 $\{Y_i \times_Y T \to T\}$ 是 fpqc 覆盖，
由《空间的性质》命题
\ref{spaces-properties-proposition-sheaf-fpqc} 得到 $a = b$。
\end{proof}


\section{fppf 拓扑中态射性质的下降}
\label{spaces-descent-section-descending-properties-morphisms-fppf}

\noindent
本节给出代数空间态射的一些性质：我们（尚）未能证明它们在
fpqc 拓扑中关于基是局部的，但它们在 fppf 拓扑中关于基是局部的。

\begin{lemma}
\label{spaces-descent-lemma-descending-fppf-property-immersion}
性质 $\mathcal{P}(f) =$``$f$ 是浸入''
关于基是 fppf 局部的。
\end{lemma}

\begin{proof}
设 $f : X \to Y$ 为代数空间的态射。
设 $\{Y_i \to Y\}_{i \in I}$ 为 $Y$ 的 fppf 覆盖。
设 $f_i : X_i \to Y_i$ 为 $f$ 的基变换。

\medskip\noindent
若 $f$ 是浸入，则由《空间》引理
\ref{spaces-lemma-base-change-immersions}，每个 $f_i$ 都是浸入。
这证明了定义 \ref{spaces-descent-definition-property-morphisms-local} 中的正向蕴含。

\medskip\noindent
反之，假设每个 $f_i$ 都是浸入。由《空间的态射》引理
\ref{spaces-morphisms-lemma-immersions-monomorphisms}，
这意味着每个 $f_i$ 都分离。由《空间的态射》引理
\ref{spaces-morphisms-lemma-immersion-quasi-finite}，
这意味着每个 $f_i$ 都局部拟有限。因此，应用引理
\ref{spaces-descent-lemma-descending-property-separated} 与
\ref{spaces-descent-lemma-descending-property-quasi-finite}，可见 $f$
局部拟有限且分离。由《空间的态射》引理
\ref{spaces-morphisms-lemma-locally-quasi-finite-separated-representable}，
这意味着 $f$ 可表！

\medskip\noindent
由《空间的态射》引理
\ref{spaces-morphisms-lemma-closed-immersion-local}，只须证明：
对每个概形 $Z$ 与态射 $Z \to Y$，基变换
$Z \times_Y X \to Z$ 是浸入。由《空间上的拓扑》引理
\ref{spaces-topologies-lemma-refine-fppf-schemes}，可找到由概形组成的
fppf 覆盖 $\{Z_i \to Z\}$，它加细覆盖
$\{Y_i \to Y\}$ 到 $Z$ 上的拉回。因此可见
$Z \times_Y X \to Z$（由前一段的结论，这是概形的态射）
在拉回到 $Z$ 的一个 fppf（由概形组成）覆盖的各成员后成为浸入。
% P09-ERR-0356
故由概形情形的结论，$Z \times_Y X \to Z$ 是浸入；见《下降》引理
\ref{descent-lemma-descending-fppf-property-immersion}。
\end{proof}

\begin{lemma}
\label{spaces-descent-lemma-descending-fppf-property-locally-separated}
性质 $\mathcal{P}(f) =$``$f$ 是局部分离的''
关于基是 fppf 局部的。
\end{lemma}

\begin{proof}
局部分离态射的任意基变换仍局部分离；见《空间的态射》引理
\ref{spaces-morphisms-lemma-base-change-separated}。
因此，定义 \ref{spaces-descent-definition-property-morphisms-local} 中的正向蕴含成立。

\medskip\noindent
设 $\{Y_i \to Y\}_{i \in I}$ 为 $S$ 上代数空间的 fppf 覆盖。
设 $f : X \to Y$ 为 $S$ 上代数空间的一个态射。
假设每个基变换 $X_i := Y_i \times_Y X \to Y_i$ 都局部分离。
这意味着每个态射
$$
\Delta_i :
X_i
\longrightarrow
X_i \times_{Y_i} X_i = Y_i \times_Y (X \times_Y X)
$$
都是浸入。fppf 覆盖的基变换仍为 fppf 覆盖；见
% P09-ERR-0357
《空间上的拓扑》引理 \ref{spaces-topologies-lemma-fppf}，
故 $\{Y_i \times_Y (X \times_Y X) \to X \times_Y X\}$
是代数空间的 fppf 覆盖。此外，每个 $\Delta_i$ 都是态射
$\Delta : X \to X \times_Y X$ 的基变换。因此，由引理
\ref{spaces-descent-lemma-descending-fppf-property-immersion}，$\Delta$ 是浸入，
即 $f$ 局部分离。
\end{proof}


\section{态射性质下降的应用}
\label{spaces-descent-section-application-descending-properties-morphisms}

\noindent
本节是《下降》第
\ref{descent-section-application-descending-properties-morphisms} 节的对应版本。

\begin{lemma}
\label{spaces-descent-lemma-descending-property-ample}
设 $S$ 为概形。设 $f : X \to Y$ 为 $S$ 上代数空间的态射。
设 $\mathcal{L}$ 为可逆 $\mathcal{O}_X$-模。
设 $\{g_i : Y_i \to Y\}_{i \in I}$ 为 fpqc 覆盖。
设 $f_i : X_i \to Y_i$ 为 $f$ 的基变换，并设 $\mathcal{L}_i$
为 $\mathcal{L}$ 到 $X_i$ 上的拉回。下列条件等价：
\begin{enumerate}
\item $\mathcal{L}$ 在 $X/Y$ 上丰沛；以及
\item $\mathcal{L}_i$ 在 $X_i/Y_i$ 上丰沛，对每个 $i \in I$ 成立。
\end{enumerate}
\end{lemma}

\begin{proof}
蕴含 (1) $\Rightarrow$ (2) 由《空间上的除子》引理
\ref{spaces-divisors-lemma-ample-base-change} 得出。
假设 (2)。为检验 $\mathcal{L}$ 在 $X/Y$ 上丰沛，
可在 $Y$ 上平展局部地工作；见《空间上的除子》引理
\ref{spaces-divisors-lemma-relatively-ample-local}。
因此可假设 $Y$ 是概形，继而也可假设每个 $Y_i$ 都是概形；
见《空间上的拓扑》引理
\ref{spaces-topologies-lemma-refine-fpqc-schemes}。
换言之，可假设 $\{Y_i \to Y\}$ 是概形的 fpqc 覆盖。

\medskip\noindent
由《空间上的除子》引理
\ref{spaces-divisors-lemma-relatively-ample-properties}，
可见 $X_i \to Y_i$ 可表（即 $X_i$ 是概形）、拟紧且分离。
因此，由引理 \ref{spaces-descent-lemma-descending-property-quasi-compact} 与
\ref{spaces-descent-lemma-descending-property-separated}，$f$ 拟紧且分离。
这意味着
$\mathcal{A} = \bigoplus_{d \geq 0} f_*\mathcal{L}^{\otimes d}$
是一个拟凝聚分次 $\mathcal{O}_Y$-代数
（《空间的态射》引理 \ref{spaces-morphisms-lemma-pushforward}）。
此外，由《空间的上同调》引理
\ref{spaces-cohomology-lemma-flat-base-change-cohomology}，
$\mathcal{A}$ 的形成与平坦基变换交换。特别地，若令
$\mathcal{A}_i = \bigoplus_{d \geq 0} f_{i, *}\mathcal{L}_i^{\otimes d}$，
则有 $\mathcal{A}_i = g_i^*\mathcal{A}$。
因此，自然映射
$\psi_d : f^*\mathcal{A}_d \to \mathcal{L}^{\otimes d}$
是 $\mathcal{O}_X$-模的映射；其拉回给出自然映射
$\psi_{i, d} : f_i^*(\mathcal{A}_i)_d \to \mathcal{L}_i^{\otimes d}$，
后者是 $\mathcal{O}_{X_i}$-模的映射。
由于 $\mathcal{L}_i$ 在 $X_i/Y_i$ 上丰沛，对任意点
$x_i \in X_i$，存在 $d \geq 1$，使得
$f_i^*(\mathcal{A}_i)_d \to \mathcal{L}_i^{\otimes d}$
在 $x_i$ 处的茎上为满射。这可直接由相对丰沛模的定义得出，
也可由《态射》引理
\ref{morphisms-lemma-characterize-relatively-ample} 得出。
若 $x \in |X|$，则可选择某个 $i$ 以及点
$x_i \in X_i$，它映到 $x$。由于
$\mathcal{O}_{X, \overline{x}} \to \mathcal{O}_{X_i, \overline{x}_i}$
平坦，因而忠实平坦，
% P09-ERR-0358
我们得出：对每个 $x \in |X|$，存在 $d \geq 1$，使得
$f^*\mathcal{A}_d \to \mathcal{L}^{\otimes d}$
在 $x$ 处的茎上为满射。这意味着开子集 $U(\psi) \subset X$
（见《空间上的除子》引理
\ref{spaces-divisors-lemma-invertible-map-into-relative-proj}）
所对应的映射
$\psi : f^*\mathcal{A} \to \bigoplus_{d \geq 0} \mathcal{L}^{\otimes d}$
是分次 $\mathcal{O}_X$-代数的映射，并且该开子集等于 $X$。
考虑相应的态射
$$
r_{\mathcal{L}, \psi} : X \longrightarrow \underline{\text{Proj}}_Y(\mathcal{A})
$$
由上文显然可见，$r_{\mathcal{L}, \psi}$ 到 $Y_i$ 的基变换
是态射 $r_{\mathcal{L}_i, \psi_i}$；由《态射》引理
\ref{morphisms-lemma-characterize-relatively-ample}，它是开浸入。
因此，由引理 \ref{spaces-descent-lemma-descending-property-open-immersion}，
$r_{\mathcal{L}, \psi}$ 是开浸入。于是 $X$ 是概形，并由
《态射》引理 \ref{morphisms-lemma-characterize-relatively-ample}
得出 $\mathcal{L}$ 在 $X/Y$ 上丰沛。
\end{proof}

\begin{lemma}
\label{spaces-descent-lemma-ample-in-neighbourhood}
设 $S$ 为概形。设 $f : X \to Y$ 为 $S$ 上代数空间的固有态射。
设 $\mathcal{L}$ 为可逆 $\mathcal{O}_X$-模。
存在开子空间 $V \subset Y$，它由下列性质刻画：
代数空间的态射 $Y' \to Y$ 经 $V$ 分解，当且仅当
$\mathcal{L}'$（即 $\mathcal{L}$ 到 $X' = Y' \times_Y X$ 上的拉回）
在 $X'/Y'$ 上丰沛（其含义见《空间上的除子》定义
\ref{spaces-divisors-definition-relatively-ample}）。
\end{lemma}

\begin{proof}
先假设当 $Y$ 是概形时本引理成立。
设 $U$ 为概形，并设 $U \to Y$ 为满平展态射。
令 $R = U \times_Y U$，其投影为 $t, s : R \to U$。
记 $X_U = U \times_Y X$，并记拉回为 $\mathcal{L}_U$。
则得到引理所述的开子概形 $V' \subset U$；这是对
$(X_U \to U, \mathcal{L}_U)$ 所得到的开子概形。由函子性刻画可见
$s^{-1}(V') = t^{-1}(V')$。因此存在开子空间 $V \subset Y$，
使 $V'$ 是 $V$ 在 $U$ 中的逆像。特别地，$V' \to V$
满且平展，从而 $\mathcal{L}_V$ 在 $X_V/V$ 上丰沛
（《空间上的除子》引理
\ref{spaces-divisors-lemma-relatively-ample-local}）。
现在，若 $Y' \to Y$ 是一个态射，且 $\mathcal{L}'$
在 $X'/Y'$ 上丰沛，则 $U \times_Y Y' \to Y'$ 必经 $V'$ 分解，
由此得出 $Y' \to Y$ 经 $V$ 分解。因此 $V \subset Y$
正是引理陈述中的开子空间。这样便约化到下一段处理的情形。

\medskip\noindent
假设 $Y$ 是概形。由于问题在 $Y$ 上是局部的，可假设 $Y$
是仿射概形。我们将证明：
\begin{enumerate}
\item[(A)] 若 $\Spec(k) \to Y$ 是一个态射，使得
$\mathcal{L}_k$ 在 $X_k/k$ 上丰沛，则存在开邻域
$V \subset Y$，它是 $\Spec(k) \to Y$ 的像的邻域，并使得
$\mathcal{L}_V$ 在 $X_V/V$ 上丰沛。
\end{enumerate}
显然，(A) 蕴含本引理。

\medskip\noindent
设 $X \to Y$、$\mathcal{L}$、$\Spec(k) \to Y$ 如 (A) 所述。
由引理 \ref{spaces-descent-lemma-descending-property-ample}，可假设
$k = \kappa(y)$ 是某点 $y$ 的剩余域，该点属于 $Y$。

\medskip\noindent
由于 $Y$ 仿射，可找到一个有向集 $I$，以及代数空间态射的逆系统
$X_i \to Y_i$，其中 $Y_i$ 在 $\mathbf{Z}$ 上有限展示，
转移态射 $X_i \to X_{i'}$ 与 $Y_i \to Y_{i'}$ 均仿射，
$X_i \to Y_i$ 固有且有限展示，并且
$X \to Y = \lim (X_i \to Y_i)$。见《空间的极限》引理
\ref{spaces-limits-lemma-proper-limit-of-proper-finite-presentation-noetherian}。
缩小 $I$ 后，可假设 $Y_i$ 对所有 $i$ 都是（仿射）概形；见
《空间的极限》引理 \ref{spaces-limits-lemma-limit-is-affine}。
再次缩小 $I$ 后，可假设存在相容的可逆
$\mathcal{O}_{X_i}$-模系统 $\mathcal{L}_i$，其拉回为
$\mathcal{L}$；见《空间的极限》引理
\ref{spaces-limits-lemma-descend-invertible-modules}。
令 $y_i \in Y_i$ 为 $y$ 的像。则
$\kappa(y) = \colim \kappa(y_i)$。所以
$X_y = \lim X_{i, y_i}$；缩小 $I$ 后，可假设
$X_{i, y_i}$ 对所有 $i$ 都是概形；见《空间的极限》引理
\ref{spaces-limits-lemma-limit-is-scheme}。
因此，对某个 $i$，由《极限》引理
\ref{limits-lemma-limit-ample}，$\mathcal{L}_{i, y_i}$
在 $X_{i, y_i}$ 上丰沛。由《空间上的除子》引理
\ref{spaces-divisors-lemma-ample-in-neighbourhood}，
存在开邻域 $V_i \subset Y_i$，它包含 $y_i$，并使得
$\mathcal{L}_i$ 限制到 $f_i^{-1}(V_i)$ 后相对于 $V_i$ 丰沛。
令 $V \subset Y$ 为 $V_i$ 的逆像，即得证明
（提示：使用《态射》引理
\ref{morphisms-lemma-ample-base-change}，以及
$X \to Y \times_{Y_i} X_i$ 是仿射的事实；还要使用：
沿仿射态射拉回丰沛可逆层仍丰沛，见《态射》引理
\ref{morphisms-lemma-pullback-ample-tensor-relatively-ample}）。
\end{proof}


\section{在源上局部的态射性质}
\label{spaces-descent-section-properties-morphisms-local-source}

\noindent
本节定义代数空间态射的性质在源上局部的含义。
请比较《下降》第
\ref{descent-section-properties-morphisms-local-source} 节。

\begin{definition}
\label{spaces-descent-definition-property-morphisms-local-source}
设 $S$ 为概形。设 $\mathcal{P}$ 为 $S$ 上代数空间态射的一个性质。
设 $\tau \in \{fpqc, \linebreak[0] fppf, \linebreak[0] syntomic, \linebreak[0]
smooth, \linebreak[0] \etale\}$。称 $\mathcal{P}$
{it 在源上 $\tau$-局部}，或称它
{it 对 $\tau$-拓扑在源上局部}，如果对代数空间的任意态射
$f : X \to Y$（它位于 $S$ 上），以及代数空间的任意 $\tau$-覆盖
$\{X_i \to X\}_{i \in I}$，都有
$$
f \text{ 具有 }\mathcal{P}
\Leftrightarrow
\text{每个 }X_i \to Y\text{ 都具有 }\mathcal{P}.
$$
\end{definition}

\noindent
需要明确的是，由于同构总是覆盖，可见（或要求）性质
$\mathcal{P}$ 对 $X \to Y$ 成立，当且仅当它对任意箭头
$X' \to Y'$ 成立，只要该箭头同构于 $X \to Y$。
若一个性质在源上 $\tau$-局部，则把具有该性质的态射预复合
$\tau$-覆盖中出现的态射后，该性质仍保持。形式化陈述如下。

\begin{lemma}
\label{spaces-descent-lemma-precompose-property-local-source}
设 $S$ 为概形。设
$\tau \in \{fpqc, \linebreak[0] fppf, \linebreak[0] syntomic, \linebreak[0]
smooth, \linebreak[0] \etale\}$。
设 $\mathcal{P}$ 为 $S$ 上代数空间态射的一个性质，且它在源上
$\tau$-局部。设 $f : X \to Y$ 具有性质 $\mathcal{P}$。
对于任意态射 $a : X' \to X$，若它分别为平坦、平坦且局部有限展示、
syntomic、光滑、平展态射，则复合 $f \circ a : X' \to Y$
具有性质 $\mathcal{P}$。
\end{lemma}

\begin{proof}
这是因为可以把 $X' \to X$ 放入一个构成 $\tau$-覆盖的态射族中。
\end{proof}

\begin{lemma}
\label{spaces-descent-lemma-transfer-from-schemes}
设 $S$ 为概形。设
$\tau \in \{fpqc, \linebreak[0] fppf, \linebreak[0] syntomic, \linebreak[0]
smooth, \linebreak[0] \etale\}$。
假设 $\mathcal{P}$ 为 $S$ 上概形态射的一个性质，且它在源与靶上
% P09-ERR-0361
平展局部。把代数空间态射的相应性质记作
$\mathcal{P}_{spaces}$；这是 $S$ 上的性质，见《空间的态射》定义
\ref{spaces-morphisms-definition-P}。若 $\mathcal{P}$ 对 $\tau$-拓扑
在源上局部，则 $\mathcal{P}_{spaces}$ 对 $\tau$-拓扑也在源上局部。
\end{lemma}

\begin{proof}
设 $f : X \to Y$ 为 $S$ 上代数空间的态射。
设 $\{X_i \to X\}_{i \in I}$ 为代数空间的 $\tau$-覆盖。
选择概形 $V$ 与满平展态射 $V \to Y$。
选择概形 $U$ 与满平展态射 $U \to X \times_Y V$。
对每个 $i$，选择概形 $U_i$ 与满平展态射
$U_i \to X_i \times_X U$。

\medskip\noindent
注意 $\{X_i \times_X U \to U\}_{i \in I}$ 是 $\tau$-覆盖。
每个 $\{U_i \to X_i \times_X U\}$ 都是平展覆盖，因而是
$\tau$-覆盖。所以 $\{U_i \to U\}_{i \in I}$ 是代数空间的
$\tau$-覆盖，且位于 $S$ 上。但由于 $U$ 与每个 $U_i$ 都是概形，
% P09-ERR-0360
可见 $\{U_i \to U\}_{i \in I}$ 是概形的 $\tau$-覆盖，且位于 $S$ 上。

\medskip\noindent
现在有
\begin{align*}
f \text{ 具有 }\mathcal{P}_{spaces}
& \Leftrightarrow
U \to V \text{ 具有 }\mathcal{P} \\
& \Leftrightarrow
\text{每个 }U_i \to V \text{ 都具有 }\mathcal{P} \\
& \Leftrightarrow
\text{每个 }X_i \to Y\text{ 都具有 }\mathcal{P}_{spaces}.
\end{align*}
% P09-ERR-0359
第一与最后一个等价关系由 $\mathcal{P}_{spaces}$ 的定义得出；
中间的等价关系源于我们假设 $\mathcal{P}$ 对 $\tau$-拓扑在源上局部。
\end{proof}


\section{在源上对 fpqc 拓扑局部的态射性质}
\label{spaces-descent-section-fpqc-local-source}

\noindent
下面给出一些在源上 fpqc 局部的态射性质。

\begin{lemma}
\label{spaces-descent-lemma-flat-fpqc-local-source}
性质 $\mathcal{P}(f)=$``$f$ 是平坦的''在源上 fpqc 局部。
\end{lemma}

\begin{proof}
结论由引理 \ref{spaces-descent-lemma-transfer-from-schemes} 得出；这里使用
《空间的态射》定义 \ref{spaces-morphisms-definition-flat} 与
《下降》引理 \ref{descent-lemma-flat-fpqc-local-source}。
\end{proof}


\section{在源上对 fppf 拓扑局部的态射性质}
\label{spaces-descent-section-fppf-local-source}

\noindent
下面给出一些在源上 fppf 局部的态射性质。

\begin{lemma}
\label{spaces-descent-lemma-locally-finite-presentation-fppf-local-source}
性质 $\mathcal{P}(f)=$``$f$ 是局部有限展示的''
在源上 fppf 局部。
\end{lemma}

\begin{proof}
结论由引理 \ref{spaces-descent-lemma-transfer-from-schemes} 得出；这里使用
《空间的态射》定义
\ref{spaces-morphisms-definition-locally-finite-presentation} 与
《下降》引理
\ref{descent-lemma-locally-finite-presentation-fppf-local-source}。
\end{proof}

\begin{lemma}
\label{spaces-descent-lemma-locally-finite-type-fppf-local-source}
性质 $\mathcal{P}(f)=$``$f$ 是局部有限型的''
在源上 fppf 局部。
\end{lemma}

\begin{proof}
结论由引理 \ref{spaces-descent-lemma-transfer-from-schemes} 得出；这里使用
《空间的态射》定义
\ref{spaces-morphisms-definition-locally-finite-type} 与
《下降》引理 \ref{descent-lemma-locally-finite-type-fppf-local-source}。
\end{proof}

\begin{lemma}
\label{spaces-descent-lemma-open-fppf-local-source}
性质 $\mathcal{P}(f)=$``$f$ 是开的''在源上 fppf 局部。
\end{lemma}

\begin{proof}
结论由引理 \ref{spaces-descent-lemma-transfer-from-schemes} 得出；这里使用
《空间的态射》定义 \ref{spaces-morphisms-definition-open} 与
《下降》引理 \ref{descent-lemma-open-fppf-local-source}。
\end{proof}

\begin{lemma}
\label{spaces-descent-lemma-universally-open-fppf-local-source}
性质 $\mathcal{P}(f)=$``$f$ 是普遍开的''
在源上 fppf 局部。
\end{lemma}

\begin{proof}
结论由引理 \ref{spaces-descent-lemma-transfer-from-schemes} 得出；这里使用
《空间的态射》定义 \ref{spaces-morphisms-definition-open} 与
《下降》引理 \ref{descent-lemma-universally-open-fppf-local-source}。
\end{proof}


\section{在源上对 syntomic 拓扑局部的态射性质}
\label{spaces-descent-section-syntomic-local-source}

\noindent
下面给出一些在源上 syntomic 局部的态射性质。

\begin{lemma}
\label{spaces-descent-lemma-syntomic-syntomic-local-source}
性质 $\mathcal{P}(f)=$``$f$ 是 syntomic 的''
在源上 syntomic 局部。
\end{lemma}

\begin{proof}
结论由引理 \ref{spaces-descent-lemma-transfer-from-schemes} 得出；这里使用
《空间的态射》定义 \ref{spaces-morphisms-definition-syntomic} 与
《下降》引理 \ref{descent-lemma-syntomic-syntomic-local-source}。
\end{proof}


\section{在源上对光滑拓扑局部的态射性质}
\label{spaces-descent-section-smooth-local-source}

\noindent
下面给出一些在源上光滑局部的态射性质。

\begin{lemma}
\label{spaces-descent-lemma-smooth-smooth-local-source}
性质 $\mathcal{P}(f)=$``$f$ 是光滑的''
在源上光滑局部。
\end{lemma}

\begin{proof}
结论由引理 \ref{spaces-descent-lemma-transfer-from-schemes} 得出；这里使用
《空间的态射》定义 \ref{spaces-morphisms-definition-smooth} 与
《下降》引理 \ref{descent-lemma-smooth-smooth-local-source}。
\end{proof}


\section{在源上对平展拓扑局部的态射性质}
\label{spaces-descent-section-etale-local-source}

\noindent
下面给出一些在源上平展局部的态射性质。

\begin{lemma}
\label{spaces-descent-lemma-etale-etale-local-source}
性质 $\mathcal{P}(f)=$``$f$ 是平展的''在源上平展局部。
\end{lemma}

\begin{proof}
结论由引理 \ref{spaces-descent-lemma-transfer-from-schemes} 得出；这里使用
《空间的态射》定义 \ref{spaces-morphisms-definition-etale} 与
《下降》引理 \ref{descent-lemma-etale-etale-local-source}。
\end{proof}

\begin{lemma}
\label{spaces-descent-lemma-locally-quasi-finite-etale-local-source}
性质 $\mathcal{P}(f)=$``$f$ 是局部拟有限的''
在源上平展局部。
\end{lemma}

\begin{proof}
结论由引理 \ref{spaces-descent-lemma-transfer-from-schemes} 得出；这里使用
《空间的态射》定义
\ref{spaces-morphisms-definition-locally-quasi-finite} 与
《下降》引理
\ref{descent-lemma-locally-quasi-finite-etale-local-source}。
\end{proof}

\begin{lemma}
\label{spaces-descent-lemma-unramified-etale-local-source}
性质 $\mathcal{P}(f)=$``$f$ 是非分歧的''
在源上平展局部。
\end{lemma}

\begin{proof}
结论由引理 \ref{spaces-descent-lemma-transfer-from-schemes} 得出；这里使用
《空间的态射》定义 \ref{spaces-morphisms-definition-unramified} 与
《下降》引理 \ref{descent-lemma-unramified-etale-local-source}。
\end{proof}


\section{在源与靶上光滑局部的态射性质}
\label{spaces-descent-section-properties-local-source-target}

\noindent
设 $\mathcal{P}$ 为代数空间态射的一个性质。短语
``$\mathcal{P}$ 在源与靶上光滑局部''有一种直观含义。
然而，这个概念并不等同于同时要求 $\mathcal{P}$ 在源上光滑局部
且在靶上光滑局部。我们已在《下降》第
\ref{descent-section-properties-etale-local-source-target} 节
非常详细地讨论了类似现象（关于平展拓扑与概形范畴；快速概览见
《下降》注 \ref{descent-remark-compare-definitions}）。
不过，光滑拓扑与平展拓扑的情形之间有一个重要差别。
为看清这个差别，我们建议读者思考《下降》引理
\ref{descent-lemma-local-source-target-implies} 与引理
\ref{spaces-descent-lemma-local-source-target-implies} 之间的差别，以及《下降》引理
\ref{descent-lemma-local-source-target-characterize} 与引理
\ref{spaces-descent-lemma-local-source-target-characterize} 之间的差别。
也就是说，在平展情形下，靶的平展``覆盖''的选择无关紧要；
在光滑情形下则并非如此。

\begin{definition}
\label{spaces-descent-definition-local-source-target}
设 $S$ 为概形。设 $\mathcal{P}$ 为 $S$ 上代数空间态射的一个性质。
若满足以下条件，则称 $\mathcal{P}$ {it 在源与靶上光滑局部}：
\begin{enumerate}
\item （在与光滑映射预复合下保持）若 $f : X \to Y$ 光滑，
且 $g : Y \to Z$ 具有 $\mathcal{P}$，则 $g \circ f$ 具有
$\mathcal{P}$；
\item （在光滑基变换下保持）若 $f : X \to Y$ 具有
$\mathcal{P}$，且 $Y' \to Y$ 光滑，则基变换
$f' : Y' \times_Y X \to Y'$ 具有 $\mathcal{P}$；并且
\item （局部性）给定态射 $f : X \to Y$，下列条件等价：
\begin{enumerate}
\item $f$ 具有 $\mathcal{P}$；
\item 对每个 $x \in |X|$，存在交换图
$$
\xymatrix{
U \ar[d]_a \ar[r]_h & V \ar[d]^b \\
X \ar[r]^f & Y
}
$$
其中竖直箭头光滑，并存在 $u \in |U|$，满足
$a(u) = x$，且 $h$ 具有 $\mathcal{P}$。
\end{enumerate}
\end{enumerate}
\end{definition}

\noindent
以上作为我们的定义。在下面的引理中，我们将说明这等价于：
$\mathcal{P}$ 在靶上光滑局部、在源上光滑局部，并且在与光滑态射
后复合下保持。

\begin{lemma}
\label{spaces-descent-lemma-local-source-target-implies}
设 $S$ 为概形。设 $\mathcal{P}$ 为 $S$ 上代数空间态射的一个性质，
并且它在源与靶上光滑局部。则：
\begin{enumerate}
\item $\mathcal{P}$ 在源上光滑局部；
\item $\mathcal{P}$ 在靶上光滑局部；
\item $\mathcal{P}$ 在与光滑态射后复合下保持：若
$f : X \to Y$ 具有 $\mathcal{P}$，且 $g : Y \to Z$ 光滑，
则 $g \circ f$ 具有 $\mathcal{P}$。
\end{enumerate}
\end{lemma}

\begin{proof}
我们把细节完整写出。

\medskip\noindent
(1) 的证明。设 $f : X \to Y$ 为 $S$ 上代数空间的态射。
设 $\{X_i \to X\}_{i \in I}$ 为 $X$ 的光滑覆盖。
若每个复合 $h_i : X_i \to Y$ 都具有 $\mathcal{P}$，则对每个
% P09-ERR-0362
$|x| \in X$，可找到 $i \in I$ 与点 $x_i \in |X_i|$，
后者映到 $x$。于是 $(X_i, x_i) \to (X, x)$
是带点对之间的光滑态射，$\text{id}_Y : Y \to Y$
是光滑态射，而且 $h_i$ 如定义
\ref{spaces-descent-definition-local-source-target} 第 (3) 部分所述。
因此可见 $f$ 具有 $\mathcal{P}$。反之，若 $f$ 具有
$\mathcal{P}$，则由定义 \ref{spaces-descent-definition-local-source-target}
第 (1) 部分，每个 $X_i \to Y$ 都具有 $\mathcal{P}$。

\medskip\noindent
(2) 的证明。设 $f : X \to Y$ 为 $S$ 上代数空间的态射。
设 $\{Y_i \to Y\}_{i \in I}$ 为 $Y$ 的光滑覆盖。
记 $X_i = Y_i \times_Y X$，并记
$h_i : X_i \to Y_i$ 为 $f$ 的基变换。若每个 $h_i : X_i \to Y_i$
都具有 $\mathcal{P}$，则对每个 $x \in |X|$，选择
$i \in I$ 与点 $x_i \in |X_i|$，后者映到 $x$。
于是 $(X_i, x_i) \to (X, x)$ 是带点对之间的光滑态射，
$Y_i \to Y$ 光滑，且 $h_i$ 如定义
\ref{spaces-descent-definition-local-source-target} 第 (3) 部分所述。
因此可见 $f$ 具有 $\mathcal{P}$。反之，若 $f$ 具有
$\mathcal{P}$，则由定义 \ref{spaces-descent-definition-local-source-target}
第 (2) 部分，每个 $X_i \to Y_i$ 都具有 $\mathcal{P}$。

\medskip\noindent
(3) 的证明。假设 $f : X \to Y$ 具有 $\mathcal{P}$，且
$g : Y \to Z$ 光滑。对每个 $x \in |X|$，可把
$(X, x) \to (X, x)$ 看作带点对之间的光滑态射，
$Y \to Z$ 是光滑态射，而 $h = f$ 如定义
\ref{spaces-descent-definition-local-source-target} 第 (3) 部分所述。
因此可见 $g \circ f$ 具有 $\mathcal{P}$。
\end{proof}

\noindent
下一个引理是《态射》引理
\ref{morphisms-lemma-locally-P-characterize} 的对应版本。

\begin{lemma}
\label{spaces-descent-lemma-local-source-target-characterize}
设 $S$ 为概形。设 $\mathcal{P}$ 为 $S$ 上代数空间态射的一个性质，
并且它在源与靶上光滑局部。设 $f : X \to Y$ 为 $S$
上代数空间的一个态射。下列条件等价：
\begin{enumerate}
\item[(a)] $f$ 具有性质 $\mathcal{P}$；
\item[(b)] 对每个 $x \in |X|$，存在带点对之间的光滑态射
$a : (U, u) \to (X, x)$、光滑态射 $b : V \to Y$，以及态射
$h : U \to V$，使得 $f \circ a = b \circ h$，且 $h$ 具有
$\mathcal{P}$；
\item[(c)] 存在交换图
$$
\xymatrix{
U \ar[d]_a \ar[r]_h & V \ar[d]^b \\
X \ar[r]^f & Y
}
$$
其中 $a$、$b$ 光滑，$a$ 为满态射，且 $h$ 具有 $\mathcal{P}$；
\item[(d)] 对任意交换图
$$
\xymatrix{
U \ar[d]_a \ar[r]_h & V \ar[d]^b \\
X \ar[r]^f & Y
}
$$
若 $b$ 光滑且 $U \to X \times_Y V$ 光滑，则态射 $h$
具有 $\mathcal{P}$；
\item[(e)] 存在光滑覆盖 $\{Y_i \to Y\}_{i \in I}$，
使每个基变换 $Y_i \times_Y X \to Y_i$ 都具有 $\mathcal{P}$；
\item[(f)] 存在光滑覆盖 $\{X_i \to X\}_{i \in I}$，
使每个复合 $X_i \to Y$ 都具有 $\mathcal{P}$；
\item[(g)] 存在光滑覆盖 $\{Y_i \to Y\}_{i \in I}$，且对每个
$i \in I$，存在光滑覆盖
$\{X_{ij} \to Y_i \times_Y X\}_{j \in J_i}$，使每个态射
$X_{ij} \to Y_i$ 都具有 $\mathcal{P}$。
\end{enumerate}
\end{lemma}

\begin{proof}
(a) 与 (b) 的等价性是定义
\ref{spaces-descent-definition-local-source-target} 的一部分。
(a) 与 (e) 的等价性是引理
\ref{spaces-descent-lemma-local-source-target-implies} 第 (2) 部分。
(a) 与 (f) 的等价性是引理
\ref{spaces-descent-lemma-local-source-target-implies} 第 (1) 部分。
由于 (a) 现在分别等价于 (e) 与 (f)，所以
% P09-ERR-0363
(a) 等价于 (g)。

\medskip\noindent
显然，(c) 蕴含 (b)。若 (b) 成立，则对任意
$x \in |X|$，可选择带点对之间的光滑态射
$a_x : (U_x, u_x) \to (X, x)$、光滑态射 $b_x : V_x \to Y$，
以及态射 $h_x : U_x \to V_x$，使得
$f \circ a_x = b_x \circ h_x$，且 $h_x$ 具有 $\mathcal{P}$。
于是 $h = \coprod h_x : \coprod U_x \to \coprod V_x$，连同
$a = \coprod a_x$ 与 $b = \coprod b_x$，构成 (c) 中所述的图。
（注意，$h$ 具有性质 $\mathcal{P}$，因为
$\{V_x \to \coprod V_x\}$ 是光滑覆盖，且 $\mathcal{P}$
在靶上光滑局部。）因此 (b) 与 (c) 等价。

\medskip\noindent
现在已知 (a)、(b)、(c)、(e)、(f)、(g) 等价。
假设 (a) 成立。设 $U, V, a, b, h$ 如 (d) 所述。
$X \times_Y V \to V$ 具有 $\mathcal{P}$，因为
$\mathcal{P}$ 在光滑基变换下保持；$U \to V$ 具有
$\mathcal{P}$，因为 $\mathcal{P}$ 在与光滑态射预复合下保持。
反之，若 (d) 成立，
令 $U = X$、$V = Y$，可见 $f$ 具有 $\mathcal{P}$。
\end{proof}

\begin{lemma}
\label{spaces-descent-lemma-smooth-local-source-target}
设 $S$ 为概形。
设 $\mathcal{P}$ 为 $S$ 上代数空间态射的一个性质。假设：
\begin{enumerate}
\item $\mathcal{P}$ 在源上光滑局部；
\item $\mathcal{P}$ 在靶上光滑局部；
\item $\mathcal{P}$ 在与光滑态射后复合下保持：若
$f : X \to Y$ 具有 $\mathcal{P}$，且 $Y \to Z$ 为光滑态射，
则 $X \to Z$ 具有 $\mathcal{P}$。
\end{enumerate}
则 $\mathcal{P}$ 在源与靶上光滑局部。
\end{lemma}

\begin{proof}
设 $\mathcal{P}$ 为满足本引理条件 (1)、(2)、(3) 的代数空间态射性质。
由引理 \ref{spaces-descent-lemma-precompose-property-local-source}，可见
$\mathcal{P}$ 在与光滑态射预复合下保持。由引理
\ref{spaces-descent-lemma-pullback-property-local-target}，可见
$\mathcal{P}$ 在光滑基变换下保持。因此只需证明定义
\ref{spaces-descent-definition-local-source-target} 的第 (3) 部分成立。

\medskip\noindent
更明确地，假设 $f : X \to Y$ 是 $S$ 上代数空间的态射，且满足
定义 \ref{spaces-descent-definition-local-source-target} 第 (3)(b) 部分。换言之，对每个
% P09-ERR-0364
$x \in X$，存在光滑态射 $a_x : U_x \to X$、点
$u_x \in |U_x|$（其像为 $x$）、光滑态射 $b_x : V_x \to Y$，以及态射
$h_x : U_x \to V_x$，使得 $f \circ a_x = b_x \circ h_x$，且
$h_x$ 具有 $\mathcal{P}$。只要证明 $f$ 具有 $\mathcal{P}$，
本引理的证明便告完成。令 $U = \coprod U_x$、$a = \coprod a_x$、
$V = \coprod V_x$、$b = \coprod b_x$、$h = \coprod h_x$。得到交换图
$$
\xymatrix{
U \ar[d]_a \ar[r]_h & V \ar[d]^b \\
X \ar[r]^f & Y
}
$$
其中
% P09-ERR-0365
$a$、$b$ 光滑，而 $a$ 为满态射。注意，$h$ 具有 $\mathcal{P}$，
因为每个 $h_x$ 都如此，且 $\mathcal{P}$ 在靶上光滑局部。
由于 $a$ 为满态射且 $\mathcal{P}$ 在源上光滑局部，只需证明
$b \circ h$ 具有 $\mathcal{P}$。这由我们的假设立即得出：
$\mathcal{P}$ 在与光滑态射后复合下保持，而 $b$ 光滑。
\end{proof}

\begin{remark}
\label{spaces-descent-remark-list-local-source-target}
利用引理 \ref{spaces-descent-lemma-smooth-local-source-target} 以及本章前面各节的结果，
很容易列出在源与靶上光滑局部的态射类型。在每种情形中，我们列出
说明该性质在源上光滑局部的引理，以及说明该性质在靶上光滑局部的
引理。每种情形下，引理 \ref{spaces-descent-lemma-smooth-local-source-target}
的第三项假设都易于检验，故予以省略。清单如下：
\begin{enumerate}
\item 平坦，见引理 \ref{spaces-descent-lemma-flat-fpqc-local-source} 和
\ref{spaces-descent-lemma-descending-property-flat}；
\item 局部有限展示，见引理
\ref{spaces-descent-lemma-locally-finite-presentation-fppf-local-source} 和
\ref{spaces-descent-lemma-descending-property-locally-finite-presentation}；
% P09-ERR-0366
\item 局部有限型，见引理 \ref{spaces-descent-lemma-locally-finite-type-fppf-local-source} 和
\ref{spaces-descent-lemma-descending-property-locally-finite-type}；
\item 普遍开，见引理 \ref{spaces-descent-lemma-universally-open-fppf-local-source} 和
\ref{spaces-descent-lemma-descending-property-universally-open}；
\item syntomic，见引理 \ref{spaces-descent-lemma-syntomic-syntomic-local-source} 和
\ref{spaces-descent-lemma-descending-property-syntomic}；
\item 光滑，见引理 \ref{spaces-descent-lemma-smooth-smooth-local-source} 和
\ref{spaces-descent-lemma-descending-property-smooth}；
% P09-ERR-0367
\item 必要时可在此增补更多条目。
\end{enumerate}
\end{remark}

\section{在源与靶上平展-光滑局部的态射性质}
\label{spaces-descent-section-properties-etale-smooth-local-source-target}

\noindent
本节是在源上平展局部、在靶上光滑局部的态射性质所对应的
第 \ref{spaces-descent-section-properties-local-source-target} 节。我们特意给这种
性质取了一个长得有些夸张的名称，以免过多使用它。

\begin{definition}
\label{spaces-descent-definition-etale-smooth-local-source-target}
设 $S$ 为概形。设 $\mathcal{P}$ 为 $S$ 上代数空间态射的一个性质。
称 $\mathcal{P}$ {it 在源与靶上平展-光滑局部}，如果：
\begin{enumerate}
\item （在与平展态射预复合下保持）若 $f : X \to Y$ 平展，且
$g : Y \to Z$ 具有 $\mathcal{P}$，则 $g \circ f$ 具有 $\mathcal{P}$；
\item （在光滑基变换下保持）若 $f : X \to Y$ 具有 $\mathcal{P}$，
且 $Y' \to Y$ 光滑，则基变换 $f' : Y' \times_Y X \to Y'$ 具有
$\mathcal{P}$；
\item （局部性）给定态射 $f : X \to Y$，下列条件等价：
\begin{enumerate}
\item $f$ 具有 $\mathcal{P}$；
\item 对每个 $x \in |X|$，存在交换图
$$
\xymatrix{
U \ar[d]_a \ar[r]_h & V \ar[d]^b \\
X \ar[r]^f & Y
}
$$
其中 $b$ 光滑、$U \to X \times_Y V$ 平展，并存在
$u \in |U|$ 满足 $a(u) = x$，使得 $h$ 具有 $\mathcal{P}$。
\end{enumerate}
\end{enumerate}
\end{definition}

\noindent
以上作为我们的定义。在下面的引理中，我们将说明，这等价于
% P09-ERR-0368
$\mathcal{P}$ 在靶上平展局部、在源上光滑局部，并且在与平展态射
后复合下保持。

\begin{lemma}
\label{spaces-descent-lemma-etale-smooth-local-source-target-implies}
设 $S$ 为概形。设 $\mathcal{P}$ 为 $S$ 上代数空间态射的一个性质，
并且它在源与靶上平展-光滑局部。则：
\begin{enumerate}
\item $\mathcal{P}$ 在源上平展局部；
\item $\mathcal{P}$ 在靶上光滑局部；
\item $\mathcal{P}$ 在与平展态射后复合下保持：若
$f : X \to Y$ 具有 $\mathcal{P}$，且 $g : Y \to Z$ 平展，
则 $g \circ f$ 具有 $\mathcal{P}$；
\item $\mathcal{P}$ 具有下列保持性质：给定 $f : X \to Y$ 以及
平展态射 $g : Y \to Z$，若 $g \circ f$ 具有 $\mathcal{P}$，
则 $f$ 具有 $\mathcal{P}$。
\end{enumerate}
\end{lemma}

\begin{proof}
我们把细节完整写出。

\medskip\noindent
(1) 的证明。设 $f : X \to Y$ 为 $S$ 上代数空间的态射。
设 $\{X_i \to X\}_{i \in I}$ 为 $X$ 的平展覆盖。若每个复合
$h_i : X_i \to Y$ 都具有 $\mathcal{P}$，则对每个
% P09-ERR-0369
$|x| \in X$，可找到 $i \in I$ 与点 $x_i \in |X_i|$，
后者映到 $x$。于是 $(X_i, x_i) \to (X, x)$
是带点对之间的平展态射，$\text{id}_Y : Y \to Y$ 是光滑态射，
而且 $h_i$ 如定义
\ref{spaces-descent-definition-etale-smooth-local-source-target} 第 (3) 部分所述。
因此可见 $f$ 具有 $\mathcal{P}$。反之，若 $f$ 具有
$\mathcal{P}$，则由定义
\ref{spaces-descent-definition-etale-smooth-local-source-target} 第 (1) 部分，
每个 $X_i \to Y$ 都具有 $\mathcal{P}$。

\medskip\noindent
(2) 的证明。设 $f : X \to Y$ 为 $S$ 上代数空间的态射。
设 $\{Y_i \to Y\}_{i \in I}$ 为 $Y$ 的光滑覆盖。
记 $X_i = Y_i \times_Y X$，并记 $h_i : X_i \to Y_i$ 为
$f$ 的基变换。若每个 $h_i : X_i \to Y_i$ 都具有 $\mathcal{P}$，
则对每个 $x \in |X|$，选择 $i \in I$ 与点
$x_i \in |X_i|$，后者映到 $x$。于是
$X_i \to X \times_Y Y_i$ 是平展态射（因为它是同构），
$Y_i \to Y$ 光滑，且 $h_i$ 如定义
% P09-ERR-0370
\ref{spaces-descent-definition-local-source-target} 第 (3) 部分所述。因此可见
$f$ 具有 $\mathcal{P}$。反之，若 $f$ 具有 $\mathcal{P}$，
则由定义 \ref{spaces-descent-definition-local-source-target} 第 (2) 部分，
每个 $X_i \to Y_i$ 都具有 $\mathcal{P}$。

\medskip\noindent
(3) 的证明。假设 $f : X \to Y$ 具有 $\mathcal{P}$，且
$g : Y \to Z$ 平展。态射 $X \to Y \times_Z X$ 是平展的，
因为它是 $X$ 上两个平展代数空间之间的态射（见《空间的性质》
引理 \ref{spaces-properties-lemma-etale-permanence}）。此外，
$Y \to Z$ 平展，
% P09-ERR-0371
因而是光滑态射。因此，下图
$$
\xymatrix{
X \ar[d] \ar[r]_f & Y \ar[d] \\
X \ar[r]^{g \circ f} & Z
}
$$
对每个 $x \in |X|$ 都适用于
% P09-ERR-0372
定义 \ref{spaces-descent-definition-local-source-target} 第 (3) 部分，
由此得出 $g \circ f$ 具有 $\mathcal{P}$。

\medskip\noindent
(4) 的证明。设 $f : X \to Y$ 为态射，$g : Y \to Z$ 平展，
且 $g \circ f$ 具有 $\mathcal{P}$。由定义
\ref{spaces-descent-definition-etale-smooth-local-source-target} 第 (2) 部分，
可见 $\text{pr}_Y : Y \times_Z X \to Y$ 具有 $\mathcal{P}$。
但态射 $(f, 1) : X \to Y \times_Z X$ 是平展的，因为它是平展投影
$\text{pr}_X : Y \times_Z X \to X$ 的截面；见《空间的态射》
引理 \ref{spaces-morphisms-lemma-etale-permanence}。因此，由定义
\ref{spaces-descent-definition-etale-smooth-local-source-target} 第 (1) 部分，
$f = \text{pr}_Y \circ (f, 1)$ 具有 $\mathcal{P}$。
\end{proof}

\begin{lemma}
\label{spaces-descent-lemma-etale-smooth-local-source-target-characterize}
设 $S$ 为概形。设 $\mathcal{P}$ 为 $S$ 上代数空间态射的一个性质，
并且它在源与靶上平展-光滑局部。设 $f : X \to Y$ 为 $S$
上代数空间的一个态射。下列条件等价：
\begin{enumerate}
\item[(a)] $f$ 具有性质 $\mathcal{P}$；
\item[(b)] 对每个 $x \in |X|$，存在光滑态射 $b : V \to Y$、
平展态射 $a : U \to V \times_Y X$，以及点 $u \in |U|$，
它映到 $x$，使得 $U \to V$ 具有 $\mathcal{P}$；
\item[(c)] 存在交换图
$$
\xymatrix{
U \ar[d]_a \ar[r]_h & V \ar[d]^b \\
X \ar[r]^f & Y
}
$$
其中 $b$ 光滑、$U \to V \times_Y X$ 平展、$a$ 为满态射，
且态射 $h$ 具有 $\mathcal{P}$；
\item[(d)] 对任意交换图
$$
\xymatrix{
U \ar[d]_a \ar[r]_h & V \ar[d]^b \\
X \ar[r]^f & Y
}
$$
若 $b$ 光滑且 $U \to X \times_Y V$ 平展，则态射 $h$
具有 $\mathcal{P}$；
\item[(e)] 存在光滑覆盖 $\{Y_i \to Y\}_{i \in I}$，
使每个基变换 $Y_i \times_Y X \to Y_i$ 都具有 $\mathcal{P}$；
\item[(f)] 存在平展覆盖 $\{X_i \to X\}_{i \in I}$，
使每个复合 $X_i \to Y$ 都具有 $\mathcal{P}$；
\item[(g)] 存在光滑覆盖 $\{Y_i \to Y\}_{i \in I}$，且对每个
$i \in I$，存在平展覆盖
$\{X_{ij} \to Y_i \times_Y X\}_{j \in J_i}$，使每个态射
$X_{ij} \to Y_i$ 都具有 $\mathcal{P}$。
\end{enumerate}
\end{lemma}

\begin{proof}
(a) 与 (b) 的等价性是定义
\ref{spaces-descent-definition-etale-smooth-local-source-target} 的一部分。
(a) 与 (e) 的等价性是引理
\ref{spaces-descent-lemma-etale-smooth-local-source-target-implies} 第 (2) 部分。
(a) 与 (f) 的等价性是引理
\ref{spaces-descent-lemma-etale-smooth-local-source-target-implies} 第 (1) 部分。
由于 (a) 现在分别等价于 (e) 与 (f)，所以
% P09-ERR-0373
(a) 等价于 (g)。

\medskip\noindent
显然，(c) 蕴含 (b)。若 (b) 成立，则对任意 $x \in |X|$，
% P09-ERR-0374
可选择光滑态射 $b_x : V_x \to Y$、平展态射
$U_x \to V_x \times_Y X$，以及点 $u_x \in |U_x|$，后者映到 $x$，
使得 $U_x \to V_x$ 具有 $\mathcal{P}$。于是
% P09-ERR-0375
$h = \coprod h_x : \coprod U_x \to \coprod V_x$，连同
$a = \coprod a_x$ 与 $b = \coprod b_x$，构成 (c) 中所述的图。
（注意，$h$ 具有性质 $\mathcal{P}$，因为
$\{V_x \to \coprod V_x\}$ 是光滑覆盖，且 $\mathcal{P}$
在靶上光滑局部。）因此 (b) 与 (c) 等价。

\medskip\noindent
现在已知 (a)、(b)、(c)、(e)、(f)、(g) 等价。
假设 (a) 成立。设 $U, V, a, b, h$ 如 (d) 所述。
$X \times_Y V \to V$ 具有 $\mathcal{P}$，因为 $\mathcal{P}$
在光滑基变换下保持；由此 $U \to V$ 具有 $\mathcal{P}$，
因为 $\mathcal{P}$ 在与平展态射预复合下保持。反之，若 (d) 成立，
令 $U = X$、$V = Y$，可见 $f$ 具有 $\mathcal{P}$。
\end{proof}

\begin{lemma}
\label{spaces-descent-lemma-etale-smooth-local-source-target}
设 $S$ 为概形。设 $\mathcal{P}$ 为 $S$ 上代数空间态射的一个性质。
假设：
\begin{enumerate}
\item $\mathcal{P}$ 在源上平展局部；
\item $\mathcal{P}$ 在靶上光滑局部；
\item $\mathcal{P}$ 在与开浸入后复合下保持：若
$f : X \to Y$ 具有 $\mathcal{P}$，且 $Y \subset Z$ 为开嵌入，
则 $X \to Z$ 具有 $\mathcal{P}$。
\end{enumerate}
则 $\mathcal{P}$ 在源与靶上平展-光滑局部。
\end{lemma}

\begin{proof}
设 $\mathcal{P}$ 为满足本引理条件 (1)、(2)、(3) 的代数空间态射性质。
由引理 \ref{spaces-descent-lemma-precompose-property-local-source}，可见
$\mathcal{P}$ 在与平展态射预复合下保持。由引理
\ref{spaces-descent-lemma-pullback-property-local-target}，可见 $\mathcal{P}$
在光滑基变换下保持。因此只需证明
% P09-ERR-0376
定义 \ref{spaces-descent-definition-local-source-target} 第 (3) 部分成立。

\medskip\noindent
更明确地，假设 $f : X \to Y$ 是 $S$ 上代数空间的态射，且满足
定义 \ref{spaces-descent-definition-local-source-target} 第 (3)(b) 部分。换言之，对每个
% P09-ERR-0377
$x \in X$，存在光滑态射 $b_x : V_x \to Y$、平展态射
$U_x \to V_x \times_Y X$，以及点 $u_x \in |U_x|$，后者映到 $x$，
使得 $h_x : U_x \to V_x$ 具有 $\mathcal{P}$。只要证明
$f$ 具有 $\mathcal{P}$，本引理的证明便告完成。

\medskip\noindent
设 $a_x : U_x \to X$ 为复合 $U_x \to V_x \times_Y X \to X$。
令 $U = \coprod U_x$、$a = \coprod a_x$、$V = \coprod V_x$、
$b = \coprod b_x$、$h = \coprod h_x$。得到交换图
$$
\xymatrix{
U \ar[d]_a \ar[r]_h & V \ar[d]^b \\
X \ar[r]^f & Y
}
$$
其中 $b$ 光滑、$U \to V \times_Y X$ 平展，
% P09-ERR-0378
$a$ 为满态射。注意，$h$ 具有 $\mathcal{P}$，因为每个
$h_x$ 都如此，且 $\mathcal{P}$ 在靶上光滑局部。下一段将证明，
可以假设 $U, V, X, Y$ 都是概形；我们建议读者跳过这一段。

\medskip\noindent
设 $X, Y, U, V, a, b, f, h$ 如上一段所述。我们需要证明
$f$ 具有 $\mathcal{P}$。取满平展态射 $X' \to X$，使
% P09-ERR-0379
$X_i$ 为概形。令 $U' = X' \times_X U$。则 $U' \to X'$ 为满态射，
且 $U' \to X' \times_Y V$ 平展。由于 $\mathcal{P}$ 在源上平展局部，
可见 $U' \to V$ 具有 $\mathcal{P}$，并且只需证明
$X' \to Y$ 具有 $\mathcal{P}$。换言之，可以假设 $X$ 为概形。
接着，选择满平展态射 $Y' \to Y$，使 $Y'$ 为概形。令
$V' = V \times_Y Y'$、$X' = X \times_Y Y'$，以及
$U' = U \times_Y Y'$。则 $U' \to X'$ 为满态射，且
$U' \to X' \times_{Y'} V'$ 平展。由于 $\mathcal{P}$ 在靶上光滑局部，
可见 $U' \to V'$ 具有 $\mathcal{P}$，并且只需证明
$X' \to Y'$ 具有 $\mathcal{P}$。因此可以假设 $X$ 与 $Y$
都是概形。选择满平展态射 $V' \to V$，使 $V'$ 为概形。令
$U' = U \times_V V'$。则 $U' \to X$ 为满态射，且
$U' \to X \times_Y V'$ 平展。由于
% P09-ERR-0380
$\mathcal{P}$ 在源上光滑局部，可见 $U' \to V'$ 具有
$\mathcal{P}$。因此可以把 $U, V$ 替换为 $U', V'$，并假设
$X, Y, V$ 都是概形。最后，把 $U$ 替换为一个满平展于 $U$
之上的概形，由此可见，可以假设 $U, V, X, Y$ 全都是概形。

\medskip\noindent
若 $U, V, X, Y$ 是概形，则由《下降》引理
\ref{descent-lemma-etale-tau-local-source-target}，$f$ 具有 $\mathcal{P}$。
\end{proof}

\begin{remark}
\label{spaces-descent-remark-list-etale-smooth-local-source-target}
利用引理 \ref{spaces-descent-lemma-etale-smooth-local-source-target} 以及本章前面各节的
结果，很容易列出
% P09-ERR-0381
在源与靶上光滑局部的态射类型。在每种情形中，我们列出说明该性质
在源上平展局部的引理，以及说明该性质在靶上光滑局部的引理。
每种情形下，引理 \ref{spaces-descent-lemma-etale-smooth-local-source-target}
的第三项假设都易于检验，故予以省略。清单如下：
\begin{enumerate}
\item 平展，见引理 \ref{spaces-descent-lemma-etale-etale-local-source} 和
\ref{spaces-descent-lemma-descending-property-etale}；
\item 局部拟有限，见引理
\ref{spaces-descent-lemma-locally-quasi-finite-etale-local-source} 和
\ref{spaces-descent-lemma-descending-property-quasi-finite}；
\item 非分歧，见引理 \ref{spaces-descent-lemma-unramified-etale-local-source} 和
\ref{spaces-descent-lemma-descending-property-unramified}；
% P09-ERR-0382
\item 必要时可在此增补更多条目。
\end{enumerate}
当然，注 \ref{spaces-descent-remark-list-local-source-target} 中列出的任何性质，
更是可以列在这里的例子。
\end{remark}

\section{空间之上空间的下降资料}
\label{spaces-descent-section-descent-datum}

\noindent
本节是代数空间情形下《下降》第
\ref{descent-section-descent-datum} 节的对应版本。本节的大多数论证
都是形式的，只依赖于下降资料的定义。

\begin{definition}
\label{spaces-descent-definition-descent-datum}
设 $S$ 为概形。设 $f : Y \to X$ 为 $S$ 上代数空间的态射。
\begin{enumerate}
\item 设 $V \to Y$ 为代数空间的态射。$V/Y/X$ 的一个
{\it 下降资料}，是一个同构
$\varphi : V \times_X Y \to Y \times_X V$；它是
$Y \times_X Y$ 上代数空间的同构，并满足如下
{\it 余圈条件}：交换图
$$
\xymatrix{
V \times_X Y \times_X Y \ar[rd]^{\varphi_{01}} \ar[rr]_{\varphi_{02}} &
&
Y \times_X Y \times_X V\\
&
Y \times_X V \times_X Y \ar[ru]^{\varphi_{12}}
}
$$
交换（采用显然的记号）。
\item 也称对 $(V/Y, \varphi)$ 是相对于 $Y \to X$ 的一个
{\it 下降资料}。
\item 下降资料之间的一个
{\it 态射 $f : (V/Y, \varphi) \to (V'/Y, \varphi')$，相对于
$Y \to X$}，是态射 $f : V \to V'$；它是
$Y$ 上代数空间的态射，并使交换图
$$
\xymatrix{
V \times_X Y \ar[r]_{\varphi} \ar[d]_{f \times \text{id}_Y} &
Y \times_X V \ar[d]^{\text{id}_Y \times f} \\
V' \times_X Y \ar[r]^{\varphi'} & Y \times_X V'
}
$$
交换。
\end{enumerate}
\end{definition}

\begin{remark}
\label{spaces-descent-remark-easier}
设 $S$ 为概形。设 $Y \to X$ 为 $S$ 上代数空间的态射。
设 $(V/Y, \varphi)$ 为相对于 $Y \to X$ 的下降资料。可以把同构
$\varphi$ 看作代数空间的同构
$$
(Y \times_X Y) \times_{\text{pr}_0, Y} V
\longrightarrow
(Y \times_X Y) \times_{\text{pr}_1, Y} V
$$
它位于 $Y \times_X Y$ 之上。因此，宽泛地说，可以把 $\varphi$ 看作映射
$\varphi : \text{pr}_0^*V \to \text{pr}_1^*V$\footnote{遗憾的是，
这里的箭头方向选“反”了。按照“函数”与“空间”对偶的一般原理，
定义 \ref{spaces-descent-definition-descent-datum} 和
\ref{spaces-descent-definition-descent-datum-for-family-of-morphisms} 中的方向，
本应与定义 \ref{spaces-descent-definition-descent-datum-quasi-coherent} 中所取的方向
相反。}。此时余圈条件表示
$\text{pr}_{02}^*\varphi =
\text{pr}_{12}^*\varphi \circ \text{pr}_{01}^*\varphi$。
这样一来，它与拟凝聚层上下降资料的情形非常相似。
\end{remark}

\noindent
下面给出具有固定靶的态射族之情形下的定义。

\begin{definition}
\label{spaces-descent-definition-descent-datum-for-family-of-morphisms}
设 $S$ 为概形。设 $\{X_i \to X\}_{i \in I}$ 为 $S$
上具有固定靶 $X$ 的代数空间态射族。
\begin{enumerate}
\item 一个{\it 下降资料 $(V_i, \varphi_{ij})$，相对于族
$\{X_i \to X\}$}，由下列对象给出：
代数空间 $V_i$ 位于 $X_i$ 上（对每个 $i \in I$）；有同构
$\varphi_{ij} : V_i \times_X X_j \to X_i \times_X V_j$，它是
$X_i \times_X X_j$ 上代数空间的同构（对每一对 $(i, j) \in I^2$）；
并且对每个
指标三元组 $(i, j, k) \in I^3$，交换图
$$
\xymatrix{
V_i \times_X X_j \times_X X_k
\ar[rd]^{\text{pr}_{01}^*\varphi_{ij}}
\ar[rr]_{\text{pr}_{02}^*\varphi_{ik}} &
&
X_i \times_X X_j \times_X V_k\\
&
X_i \times_X V_j \times_X X_k
\ar[ru]^{\text{pr}_{12}^*\varphi_{jk}}
}
$$
作为 $X_i \times_X X_j \times_X X_k$ 上代数空间的图交换
（采用显然的记号）。
\item 下降资料之间的一个
{\it 态射 $\psi : (V_i, \varphi_{ij}) \to
(V'_i, \varphi'_{ij})$}，由态射族
$\psi = (\psi_i)_{i \in I}$ 给出，其中
$\psi_i : V_i \to V'_i$ 是 $X_i$ 上代数空间的态射，并且所有图
$$
\xymatrix{
V_i \times_X X_j \ar[r]_{\varphi_{ij}} \ar[d]_{\psi_i \times \text{id}} &
X_i \times_X V_j \ar[d]^{\text{id} \times \psi_j} \\
V'_i \times_X X_j \ar[r]^{\varphi'_{ij}} & X_i \times_X V'_j
}
$$
都交换。
\end{enumerate}
\end{definition}

\begin{remark}
\label{spaces-descent-remark-easier-family}
设 $S$ 为概形。设 $\{X_i \to X\}_{i \in I}$ 为 $S$
上具有固定靶 $X$ 的代数空间态射族。设
$(V_i, \varphi_{ij})$ 为相对于 $\{X_i \to X\}$ 的下降资料。
可以把各同构 $\varphi_{ij}$ 看作代数空间的同构
$$
(X_i \times_X X_j) \times_{\text{pr}_0, X_i} V_i
\longrightarrow
(X_i \times_X X_j) \times_{\text{pr}_1, X_j} V_j
$$
它位于 $X_i \times_X X_j$ 之上。因此，宽泛地说，可以把
$\varphi_{ij}$ 看作同构
$\text{pr}_0^*V_i \to \text{pr}_1^*V_j$，它位于
$X_i \times_X X_j$ 之上。此时余圈条件表示
$\text{pr}_{02}^*\varphi_{ik} =
\text{pr}_{12}^*\varphi_{jk} \circ \text{pr}_{01}^*\varphi_{ij}$。
这样一来，它与拟凝聚层上下降资料的情形非常相似。
\end{remark}

\noindent
通常采用只含一个态射的族版本，其理由是下面的引理。

\begin{lemma}
\label{spaces-descent-lemma-family-is-one}
设 $S$ 为概形。设 $\{X_i \to X\}_{i \in I}$ 为 $S$
上具有固定靶 $X$ 的代数空间态射族。令
$Y = \coprod_{i \in I} X_i$。存在范畴的典范等价
$$
\begin{matrix}
\text{下降资料范畴 } \\
\text{相对于族 } \{X_i \to X\}_{i \in I}
\end{matrix}
\longrightarrow
\begin{matrix}
\text{下降资料范畴} \\
\text{相对于 } Y/X
\end{matrix}
$$
它把 $(V_i, \varphi_{ij})$ 映到 $(V, \varphi)$，其中
$V = \coprod_{i\in I} V_i$ 且 $\varphi = \coprod \varphi_{ij}$。
\end{lemma}

\begin{proof}
注意，$Y \times_X Y = \coprod_{ij} X_i \times_X X_j$，
更高次纤维积也同样。给出态射 $V \to Y$，恰好等价于给出一个族
$V_i \to X_i$；而给出下降资料 $\varphi$，恰好等价于给出一个族
$\varphi_{ij}$。
\end{proof}

\begin{lemma}
\label{spaces-descent-lemma-pullback}
下降资料的拉回。设 $S$ 为概形。
\begin{enumerate}
\item 设
$$
\xymatrix{
Y' \ar[r]_f \ar[d]_{a'} & Y \ar[d]^a \\
X' \ar[r]^h & X
}
$$
为 $S$ 上代数空间的交换图。构造
$$
(V \to Y, \varphi) \longmapsto f^*(V \to Y, \varphi)
= (V' \to Y', \varphi')
$$
其中 $V' = Y' \times_Y V$，而 $\varphi'$ 定义为复合
% P09-ERR-0383
$$
\xymatrix{
V' \times_{X'} Y' \ar@{=}[r] &
(Y' \times_Y V) \times_{X'} Y' \ar@{=}[r] &
(Y' \times_{X'} Y') \times_{Y \times_X Y} (V \times_X Y)
\ar[d]^{\text{id} \times \varphi} \\
Y' \times_{X'} V' \ar@{=}[r] &
Y' \times_{X'} (Y' \times_Y V) &
(Y' \times_X Y') \times_{Y \times_X Y} (Y \times_X V) \ar@{=}[l]
}
$$
它定义了一个函子，从相对于 $Y \to X$ 的下降资料范畴到
相对于 $Y' \to X'$ 的下降资料范畴。
\item 给定使图交换的两个态射 $f_i : Y' \to Y$，$i = 0, 1$，
函子 $f_0^*$ 与 $f_1^*$ 典范同构。
\end{enumerate}
\end{lemma}

\begin{proof}
我们略去 (1) 的证明，但指出：按照注 \ref{spaces-descent-remark-easier}
中引入的记号，态射 $\varphi'$ 就是态射 $(f \times f)^*\varphi$。
对于 (2)，我们指出给出函子同构的态射
$f_0^*V \to f_1^*V$。具体而言，由于 $f_0$ 与 $f_1$
都可置于该交换图中，可见存在唯一态射
$r : Y' \to Y \times_X Y$，满足 $f_i = \text{pr}_i \circ r$。
于是取
\begin{eqnarray*}
f_0^*V & = &
Y' \times_{f_0, Y} V \\
& = &
Y' \times_{\text{pr}_0 \circ r, Y} V \\
& = &
Y' \times_{r, Y \times_X Y} (Y \times_X Y) \times_{\text{pr}_0, Y} V \\
& \xrightarrow{\varphi} &
Y' \times_{r, Y \times_X Y} (Y \times_X Y) \times_{\text{pr}_1, Y} V \\
& = &
Y' \times_{\text{pr}_1 \circ r, Y} V \\
& = &
Y' \times_{f_1, Y} V \\
& = & f_1^*V
\end{eqnarray*}
我们略去这一构造确实适用的验证。
\end{proof}

\begin{definition}
\label{spaces-descent-definition-pullback-functor}
设 $S, X, X', Y, Y', f, a, a', h$ 如引理 \ref{spaces-descent-lemma-pullback} 中所述。
该引理构造的函子
$$
(V, \varphi) \longmapsto f^*(V, \varphi)
$$
称为下降资料上的{\it 拉回函子}。
\end{definition}

\begin{lemma}
\label{spaces-descent-lemma-pullback-family}
设 $S$ 为概形。设
$\mathcal{U}' = \{X'_i \to X'\}_{i \in I'}$ 与
% P09-ERR-0384
$\mathcal{U} = \{X_j \to X\}_{i \in I}$ 为具有固定靶的态射族。
设 $\alpha : I' \to I$、$g : X' \to X$ 以及
$g_i : X'_i \to X_{\alpha(i)}$ 给出具有固定靶的映射族之间的态射；
见《位点》定义 \ref{sites-definition-morphism-coverings}。
\begin{enumerate}
\item 设 $(V_i, \varphi_{ij})$ 为相对于族 $\mathcal{U}$ 的下降资料。
系统
$$
\left(
g_i^*V_{\alpha(i)}, (g_i \times g_j)^*\varphi_{\alpha(i) \alpha(j)}
\right)
$$
（采用注 \ref{spaces-descent-remark-easier-family} 中的记号）是相对于
$\mathcal{U}'$ 的下降资料。
\item 此构造定义了一个函子，从相对于 $\mathcal{U}$ 的下降资料范畴
到相对于 $\mathcal{U}'$ 的下降资料范畴。
\item 给定第二个 $\beta : I' \to I$、$h : X' \to X$ 以及
$h'_i : X'_i \to X_{\beta(i)}$，它们给出具有固定靶的映射族之间的
态射；若 $g = h$，则所得的两个下降资料函子典范同构。
\item 通过引理 \ref{spaces-descent-lemma-family-is-one}，这些函子与引理
\ref{spaces-descent-lemma-pullback} 中构造的拉回函子相符。
\end{enumerate}
\end{lemma}

\begin{proof}
这由引理 \ref{spaces-descent-lemma-pullback} 通过引理
\ref{spaces-descent-lemma-family-is-one} 的对应立即得到。
\end{proof}

\begin{definition}
\label{spaces-descent-definition-pullback-functor-family}
设 $\mathcal{U}' = \{X'_i \to X'\}_{i \in I'}$、
$\mathcal{U} = \{X_i \to X\}_{i \in I}$、$\alpha : I' \to I$、
$g : X' \to X$ 以及 $g_i : X'_i \to X_{\alpha(i)}$ 如引理
\ref{spaces-descent-lemma-pullback-family} 中所述。该引理构造的函子
$$
(V_i, \varphi_{ij}) \longmapsto
(g_i^*V_{\alpha(i)}, (g_i \times g_j)^*\varphi_{\alpha(i) \alpha(j)})
$$
称为下降资料上的{\it 拉回函子}。
\end{definition}

\noindent
若 $\mathcal{U}$ 与 $\mathcal{U}'$ 具有同一个靶 $X$，并且
$\mathcal{U}'$ 加细 $\mathcal{U}$（见《位点》定义
\ref{sites-definition-morphism-coverings}），但没有明确给定一对
$(\alpha, g_i)$，我们仍可谈论拉回函子；这是因为引理
\ref{spaces-descent-lemma-pullback-family} 已说明，这一对的选择无关紧要
（相差一个典范同构）。

\begin{definition}
\label{spaces-descent-definition-effective}
设 $S$ 为概形。设 $f : Y \to X$ 为 $S$ 上代数空间的态射。
\begin{enumerate}
\item 给定代数空间 $U$ 位于 $X$ 上，有 $U$ 相对于
$\text{id} : X \to X$ 的{\it 平凡下降资料}，即 $U$ 上的恒等态射。
\item 由引理 \ref{spaces-descent-lemma-pullback}，得到 $Y \times_X U$ 上相对于
$Y \to X$ 的一个{\it 典范下降资料}，它由平凡下降资料沿 $f$
拉回而来。我们常把这个下降资料记作
% P09-ERR-0385
$(Y \times_X U, can)$。
\item 下降资料 $(V, \varphi)$ 相对于 $Y/X$，称其为{\it 有效的}，
如果 $(V, \varphi)$ 同构于典范下降资料
$(Y \times_X U, can)$，其中某个代数空间 $U$ 位于 $X$ 上。
\end{enumerate}
\end{definition}

\noindent
因此，有效性意味着存在代数空间 $U$ 位于 $X$ 上，以及同构
$\psi : V \to Y \times_X U$ 位于 $Y$ 上，使 $\varphi$ 等于复合
% P09-ERR-0386
$$
V \times_X Y \xrightarrow{\psi \times \text{id}_Y}
Y \times_X U \times_S Y =
Y \times_X Y \times_X U
\xrightarrow{\text{id}_Y \times \psi^{-1}}
Y \times_X V
$$
这里有一个小问题：当 $Y$ 与 $X$ 都是概形时，这一定义在精神上
与《下降》定义 \ref{descent-definition-effective} 有冲突。不过，
从上下文总能清楚看出我们指的是哪个版本。

\begin{definition}
\label{spaces-descent-definition-effective-family}
设 $S$ 为概形。设 $\{X_i \to X\}$ 为 $S$ 上具有固定靶 $X$
的代数空间态射族。
\begin{enumerate}
\item 给定代数空间 $U$ 位于 $X$ 上，在代数空间族
$X_i \times_X U$ 上有一个{\it 典范下降资料}；它由 $U$ 相对于
% P09-ERR-0387
$\{\text{id} : S \to S\}$ 的平凡下降资料拉回而来。把这个下降资料记作
$(X_i \times_X U, can)$。
\item 下降资料 $(V_i, \varphi_{ij})$ 相对于
% P09-ERR-0388
$\{X_i \to S\}$，称其为{\it 有效的}，如果存在代数空间 $U$
位于 $X$ 上，使 $(V_i, \varphi_{ij})$ 同构于
$(X_i \times_X U, can)$。
\end{enumerate}
\end{definition}

\section{以层表述下降资料}
\label{spaces-descent-section-descent-data-sheaves}

\noindent
本节是《下降》第 \ref{descent-section-descent-data-sheaves} 节的
对应版本。由于代数空间本身已经是层，这里的表述略有不同。

\begin{lemma}
\label{spaces-descent-lemma-descent-data-sheaves}
设 $S$ 为概形。设 $\{X_i \to X\}_{i \in I}$ 为 $S$
上代数空间的一个 fppf 覆盖（《空间上的拓扑》定义
\ref{spaces-topologies-definition-fppf-covering}）。存在范畴的等价
$$
\left\{
\begin{matrix}
\text{下降资料 }(V_i, \varphi_{ij})\\
\text{相对于 }\{X_i \to X\}
\end{matrix}
\right\}
\leftrightarrow
\left\{
\begin{matrix}
\text{层 }F\text{ 位于 }(\Sch/S)_{fppf}\text{ 上并配有}\\
\text{映射 }F \to X\text{，使每个}\\
X_i \times_X F\text{ 都是代数空间}
\end{matrix}
\right\}.
$$
此外：
\begin{enumerate}
\item 右侧的代数空间 $X_i \times_X F$ 对应于左侧的 $V_i$；
\item 层 $F$ 是代数空间\footnote{以后将看到，只要 $I$ 不太大，
这总是成立；见《自举》引理
\ref{bootstrap-lemma-descend-algebraic-space}。}，当且仅当对应的下降资料
% P09-ERR-0389
$(X_i, \varphi_{ij})$ 是有效的。
\end{enumerate}
\end{lemma}

\begin{proof}
先构造从右到左的函子。设 $F \to X$ 是
$(\Sch/S)_{fppf}$ 上层的映射，并且每个
$V_i = X_i \times_X F$ 都是代数空间。我们有投影
$V_i \to X_i$。于是 $V_i \times_X X_j$ 与
$X_i \times_X V_j$ 都表示层 $X_i \times_X F \times_X X_j$，
因而得到同构
% P09-ERR-0390
$$
\varphi_{ii'} : V_i \times_X X_j \to X_i \times_X V_j
$$
容易看出，各映射 $\varphi_{ij}$ 是 $X_i \times_X X_j$
上的态射，并满足余圈条件。从右到左的函子由此构造
$F \mapsto (V_i, \varphi_{ij})$ 给出。

\medskip\noindent
下面构造从左到右的函子。各同构 $\varphi_{ij}$ 给出同构
$$
\varphi_{ij} : V_i \times_X X_j \longrightarrow X_i \times_X V_j
$$
它们位于
% P09-ERR-0391
$X_i \times X_j$ 之上。令 $F$ 等于下图的余等化子：
% P09-ERR-0392
$$
\xymatrix{
\coprod_{i, i'} V_i \times_X X_j
\ar@<1ex>[rr]^-{\text{pr}_0}
\ar@<-1ex>[rr]_-{\text{pr}_1 \circ \varphi_{ij}}
& &
\coprod_i V_i \ar[r]
&
F
}
$$
余圈条件保证 $F$ 带有映射 $F \to X$，而且
$X_i \times_X F$ 同构于 $V_i$。从左到右的函子由此构造
$(V_i, \varphi_{ij}) \mapsto F$ 给出。

\medskip\noindent
我们略去这些构造互为拟逆函子的验证。最后的陈述 (1) 与 (2)
由这些构造得到。
\end{proof}
% Generated mechanically from the frozen spaces-perfect.tex target.
% Do not edit this generated file; edit the bound locale source instead.

% OK, start here.
%

\setcounter{chapter}{74}
\chapter{空间的导出范畴}



\phantomsection
\label{spaces-perfect-section-phantom}


\section{引言}
\label{spaces-perfect-section-introduction}

\noindent
本章讨论代数空间上模的导出范畴。似乎没有直接论述这一主题的良好入门文献；
现有文献通常通过引用处理代数栈上模的导出范畴的论文来涵盖这一主题，例如见
\cite{olsson_sheaves}。



\section{约定}
\label{spaces-perfect-section-conventions}

\noindent
若 $\mathcal{A}$ 是阿贝尔范畴，$M$ 是 $\mathcal{A}$ 的对象，则我们仍以
$M$ 表示 $K(\mathcal{A})$ 和（或）$D(\mathcal{A})$ 中的相应对象；它对应的复形
以 $M$ 为次数 $0$ 处的项，在其余所有次数处为零。

\medskip\noindent
若给定环 $A$，则 $K(A)$ 表示 $A$-模复形的同伦范畴，$D(A)$ 表示相应的
导出范畴。类似地，若给定带环空间 $(X, \mathcal{O}_X)$，则符号
$K(\mathcal{O}_X)$ 表示 $\mathcal{O}_X$-模复形的同伦范畴，而
$D(\mathcal{O}_X)$ 表示相应的导出范畴。





\section{一般结果}
\label{spaces-perfect-section-generalities}

\noindent
本节给出关于代数空间上模的无界复形之上同调的若干一般结果。

\begin{lemma}
\label{spaces-perfect-lemma-restrict-direct-image-open}
设 $S$ 为概形，$f : X \to Y$ 为 $S$ 上代数空间的态射。给定平展态射
$V \to Y$，令 $U = V \times_Y X$，并以 $g : U \to V$ 表示投影态射。
则 $(Rf_*E)|_V = Rg_*(E|_U)$ 对 $E$（$D(\mathcal{O}_X)$ 中的对象）成立。
\end{lemma}

\begin{proof}
以 $E$ 的一个 K-单射复形 $\mathcal{I}^\bullet$ 表示，其中各项均为
$\mathcal{O}_X$-模。
于是 $Rf_*(E) = f_*\mathcal{I}^\bullet$，并且由“位点上的上同调”引理
\ref{sites-cohomology-lemma-restrict-K-injective-to-open}，
$Rg_*(E|_U) = g_*(\mathcal{I}^\bullet|_U)$。
因此结论由“代数空间的性质”引理
\ref{spaces-properties-lemma-pushforward-etale-base-change-modules} 得出。
\end{proof}


\begin{definition}
\label{spaces-perfect-definition-supported-on}
设 $S$ 为概形。设 $X$ 为 $S$ 上的代数空间。
设 $E$ 为 $D(\mathcal{O}_X)$ 中的对象。
设 $T \subset |X|$ 为闭子集。
称 $E$ {\it 支撑在 $T$ 上}，如果其上同调层 $H^i(E)$ 都支撑在 $T$ 上。
\end{definition}

\section{小平展位点上拟凝聚模的导出范畴}
\label{spaces-perfect-section-derived-quasi-coherent-etale}

\noindent
设 $X$ 为概形。本节证明 $D_\QCoh(\mathcal{O}_X)$ 可以用记作
$X_\etale$ 的 $X$ 的小平展位点来定义。以 $\mathcal{O}_\etale$ 表示
$X_\etale$ 上的结构层。
考虑带环位点的态射
\begin{equation}
\label{spaces-perfect-equation-epsilon}
\epsilon :
(X_\etale, \mathcal{O}_\etale)
\longrightarrow
(X_{Zar}, \mathcal{O}_X).
\end{equation}
在“下降”引理 \ref{descent-lemma-compare-sites} 中，它记作
$\text{id}_{small, \etale, Zar}$。

\begin{lemma}
\label{spaces-perfect-lemma-epsilon-flat}
(\ref{spaces-perfect-equation-epsilon}) 中的态射 $\epsilon$ 是带环位点的平坦态射。特别地，
函子
$\epsilon^* : \textit{Mod}(\mathcal{O}_X) \to
\textit{Mod}(\mathcal{O}_\etale)$ 正合。此外，若
$\epsilon^*\mathcal{F} = 0$，则 $\mathcal{F} = 0$。
\end{lemma}

\begin{proof}
态射 $\epsilon$ 的平坦性即“下降”引理
\ref{descent-lemma-compare-etale-zariski-flat}。下面给出另一证明。我们须证明
$\mathcal{O}_\etale$ 是平坦的 $\epsilon^{-1}\mathcal{O}_X$-模。为此，只需验证态射
$\mathcal{O}_{X, x} \to \mathcal{O}_{\etale, \overline{x}}$
对任意几何点 $\overline{x}$（属于 $X$）是平坦的；见“位点上的模”引理
\ref{sites-modules-lemma-check-flat-stalks}、“位点”引理
\ref{sites-lemma-point-morphism-sites} 以及“平展上同调”评注
\ref{etale-cohomology-remarks-enough-points}。
由“平展上同调”引理
\ref{etale-cohomology-lemma-describe-etale-local-ring}，
$\mathcal{O}_{\etale, \overline{x}}$ 是 $\mathcal{O}_{X, x}$ 的严格亨泽尔化。
因此由“代数进阶”引理
\ref{more-algebra-lemma-dumb-properties-henselization}，
$\mathcal{O}_{X, x} \to \mathcal{O}_{\etale, \overline{x}}$
是忠实平坦的。

\medskip\noindent
由“位点上的模”引理 \ref{sites-modules-lemma-flat-pullback-exact}，
$\epsilon^*$ 的正合性由 $\epsilon$ 的平坦性得出。

\medskip\noindent
设 $\mathcal{F}$ 为 $\mathcal{O}_X$-模。若
$\epsilon^*\mathcal{F} = 0$，则采用上述记号，有
$$
0 = \epsilon^*\mathcal{F}_{\overline{x}} =
\mathcal{F}_x \otimes_{\mathcal{O}_{X, x}} \mathcal{O}_{\etale, \overline{x}}
$$
对所有几何点 $\overline{x}$ 都成立（“位点上的模”引理
\ref{sites-modules-lemma-pullback-stalk}）。由
$\mathcal{O}_{X, x} \to \mathcal{O}_{\etale, \overline{x}}$
的忠实平坦性，可推出 $\mathcal{F}_x = 0$ 对所有 $x \in X$ 都成立。
\end{proof}

\noindent
设 $X$ 为概形，记号同 (\ref{spaces-perfect-equation-epsilon})。回忆，由“下降”命题
\ref{descent-proposition-equivalence-quasi-coherent} 及评注
\ref{descent-remark-change-topologies-ringed-sites}，
$\epsilon^* : \QCoh(\mathcal{O}_X)
\to \QCoh(\mathcal{O}_\etale)$ 是等价。此外，由“下降”引理
\ref{descent-lemma-equivalence-quasi-coherent-limits}，
$\QCoh(\mathcal{O}_\etale)$ 构成
$\textit{Mod}(\mathcal{O}_\etale)$ 的 Serre 子范畴。因此可以把
$D_\QCoh(\mathcal{O}_\etale)$ 定义为 $D(\mathcal{O}_\etale)$ 的三角子范畴，
其对象是上同调层均拟凝聚的复形；见“导出范畴”第
\ref{derived-section-triangulated-sub} 节。函子 $\epsilon^*$ 正合（引理
\ref{spaces-perfect-lemma-epsilon-flat}），故诱导
$\epsilon^* :  D(\mathcal{O}_X) \to D(\mathcal{O}_\etale)$；
又因拟凝聚模的拉回仍拟凝聚，它还诱导
$\epsilon^* : D_\QCoh(\mathcal{O}_X) \to
D_\QCoh(\mathcal{O}_\etale)$。

\begin{lemma}
\label{spaces-perfect-lemma-derived-quasi-coherent-small-etale-site}
设 $X$ 为概形。上面定义的函子
$\epsilon^* : D_\QCoh(\mathcal{O}_X) \to
D_\QCoh(\mathcal{O}_\etale)$ 是等价。
\end{lemma}

\begin{proof}
我们证明函子
$R\epsilon_* : D(\mathcal{O}_\etale) \to D(\mathcal{O}_X)$
诱导一个拟逆。以下将直接使用如下事实：$\epsilon_*$ 由限制到
$X_{Zar} \subset X_\etale$ 给出，而“下降”引理
\ref{descent-lemma-compare-sites} 描述了
$\epsilon^* = \text{id}_{small, \etale, Zar}^*$。

\medskip\noindent
对拟凝聚 $\mathcal{O}_X$-模 $\mathcal{F}$，伴随映射
$\mathcal{F} \to \epsilon_*\epsilon^*\mathcal{F}$ 是同构，因为
$\mathcal{F}^a$（“下降”定义
\ref{descent-definition-structure-sheaf}）是层；这一点由“下降”引理
\ref{descent-lemma-sheaf-condition-holds} 证明。反过来，每个拟凝聚
$\mathcal{O}_\etale$-模 $\mathcal{H}$ 都具有形式
$\epsilon^*\mathcal{F}$，其中某个拟凝聚 $\mathcal{O}_X$-模记为
$\mathcal{F}$；见“下降”命题
\ref{descent-proposition-equivalence-quasi-coherent}。由刚才所述，
$\mathcal{F} = \epsilon_*\mathcal{H}$，从而伴随映射
$\epsilon^*\epsilon_*\mathcal{H} \to \mathcal{H}$ 对每个拟凝聚
$\mathcal{O}_\etale$-模 $\mathcal{H}$ 都是同构。

\medskip\noindent
设 $E$ 为 $D_\QCoh(\mathcal{O}_\etale)$ 的对象，并以
$\mathcal{H}^q = H^q(E)$ 表示其第 $q$ 个上同调层。令
$\mathcal{B}$ 为 $X_\etale$ 的仿射对象所成的集合。等式
$H^p(U, \mathcal{H}^q) = 0$ 对所有 $p > 0$、所有
$q \in \mathbf{Z}$ 及所有 $U \in \mathcal{B}$ 都成立；见“下降”命题
\ref{descent-proposition-same-cohomology-quasi-coherent} 和“概形的上同调”引理
\ref{coherent-lemma-quasi-coherent-affine-cohomology-zero}。
由“位点上的上同调”引理
\ref{sites-cohomology-lemma-cohomology-over-U-trivial}，这意味着
$$
H^q(U, E) = H^0(U, \mathcal{H}^q)
$$
对所有 $U \in \mathcal{B}$ 都成立。特别地，它对仿射开集
$U \subset X$ 成立。因此，第 $q$ 个上同调——对 $R\epsilon_*E$ 而言——
在 $U$ 上等于层 $\epsilon_*\mathcal{H}^q$ 在 $U$ 上的值。层化后得到
$$
H^q(R\epsilon_*E) = \epsilon_*\mathcal{H}^q
$$
这特别表明 $R\epsilon_*$ 诱导函子
$D_\QCoh(\mathcal{O}_\etale) \to D_\QCoh(\mathcal{O}_X)$。
由于 $\epsilon^*$ 正合，进而有
$H^q(\epsilon^*R\epsilon_*E) = \epsilon^*\epsilon_*\mathcal{H}^q =
\mathcal{H}^q$（由上述讨论）。故伴随映射
$\epsilon^*R\epsilon_*E \to E$ 是同构。

\medskip\noindent
反过来，对 $F \in D_\QCoh(\mathcal{O}_X)$，伴随映射
$F \to R\epsilon_*\epsilon^*F$ 由于同样的理由也是同构；也就是说，
$R\epsilon_*\epsilon^*F$ 的上同调层同构于
$\epsilon_*H^m(\epsilon^*F) = \epsilon_*\epsilon^*H^m(F) = H^m(F)$。
\end{proof}


\section{拟凝聚模的导出范畴}
\label{spaces-perfect-section-derived-quasi-coherent}

\noindent
设 $S$ 为概形。引理
\ref{spaces-perfect-lemma-derived-quasi-coherent-small-etale-site}
表明，范畴 $D_\QCoh(\mathcal{O}_S)$ 可以用 $\mathcal{O}_S$-模复形
（在概形 $S$ 上）来定义，也可以用 $\mathcal{O}$-模复形
（在 $S$ 的小平展位点上）来定义。因此下面的定义与概形情形下的定义相容。

\begin{definition}
\label{spaces-perfect-definition-derived-quasi-coherent}
设 $S$ 为概形，$X$ 为 $S$ 上的代数空间。{\it 上同调层拟凝聚的
$\mathcal{O}_X$-模的导出范畴}记作 $D_\QCoh(\mathcal{O}_X)$。
\end{definition}

\noindent
由“代数空间的性质”引理
\ref{spaces-properties-lemma-properties-quasi-coherent}
及“导出范畴”第 \ref{derived-section-triangulated-sub} 节，这一定义是有意义的。
于是得到典范函子
\begin{equation}
\label{spaces-perfect-equation-compare}
D(\QCoh(\mathcal{O}_X))
\longrightarrow
D_\QCoh(\mathcal{O}_X)
\end{equation}
；见“导出范畴”等式 (\ref{derived-equation-compare})。

\medskip\noindent
注意，代数空间的平坦态射 $f : Y \to X$ 诱导正合函子
$f^* : \textit{Mod}(\mathcal{O}_X) \to \textit{Mod}(\mathcal{O}_Y)$；
见“代数空间的态射”引理
\ref{spaces-morphisms-lemma-flat-morphism-sites} 和“位点上的模”引理
\ref{sites-modules-lemma-flat-pullback-exact}。特别地，
$Lf^* : D(\mathcal{O}_X) \to D(\mathcal{O}_Y)$ 可以在任意代表复形上计算
（“导出范畴”引理 \ref{derived-lemma-right-derived-exact-functor}）。
我们将写作 $Lf^* = f^*$，如果 $f$ 平坦；在这种情形下，
$H^i(f^*E) = f^*H^i(E)$ 对 $E$（$D(\mathcal{O}_X)$ 中的对象）成立。
当 $f$ 平展时，我们将经常使用这一点。当然，在平展情形下，拉回函子就是
限制到 $Y_\etale$；见“代数空间的性质”等式
(\ref{spaces-properties-equation-restrict-modules})。

\begin{lemma}
\label{spaces-perfect-lemma-check-quasi-coherence-on-covering}
设 $S$ 为概形，$X$ 为 $S$ 上的代数空间，$E$ 为
$D(\mathcal{O}_X)$ 的对象。下列条件等价：
\begin{enumerate}
\item $E$ 属于 $D_\QCoh(\mathcal{O}_X)$；
\item 对每个平展态射 $\varphi : U \to X$，其中 $U$ 为仿射概形，
$\varphi^*E$ 是 $D_\QCoh(\mathcal{O}_U)$ 的对象；
\item 对每个平展态射 $\varphi : U \to X$，其中 $U$ 为概形，
$\varphi^*E$ 是 $D_\QCoh(\mathcal{O}_U)$ 的对象；
\item 存在满平展态射 $\varphi : U \to X$，其中 $U$ 为概形，使得
$\varphi^*E$ 是 $D_\QCoh(\mathcal{O}_U)$ 的对象；以及
\item 存在代数空间的满平展态射 $f : Y \to X$，使得
$Lf^*E$ 是 $D_\QCoh(\mathcal{O}_Y)$ 的对象。
\end{enumerate}
\end{lemma}

\begin{proof}
这由引理前的讨论和“代数空间的性质”引理
\ref{spaces-properties-lemma-characterize-quasi-coherent} 立即得出。
\end{proof}

\begin{lemma}
\label{spaces-perfect-lemma-quasi-coherence-direct-sums}
设 $S$ 为概形，$X$ 为 $S$ 上的代数空间。则
$D_\QCoh(\mathcal{O}_X)$ 中存在直和。
\end{lemma}

\begin{proof}
由“单射对象”引理 \ref{injectives-lemma-derived-products}，导出范畴
$D(\mathcal{O}_X)$ 中存在直和，并且可在任意代表上逐项取直和来计算。
因此，直和的上同调层显然就是上同调层的直和，因为取直和是正合函子
（在任意 Grothendieck 阿贝尔范畴中）。拟凝聚层的直和仍拟凝聚，见
“代数空间的性质”引理
\ref{spaces-properties-lemma-properties-quasi-coherent}，故引理得证。
\end{proof}

\noindent
我们还需要一些关于导出极限的信息。提醒读者，在下面的引理中，导出极限通常
不是 $D_\QCoh$ 的对象。

\begin{lemma}
\label{spaces-perfect-lemma-Rlim-quasi-coherent}
设 $S$ 为概形，$X$ 为 $S$ 上的代数空间。设 $(K_n)$ 是
$D_\QCoh(\mathcal{O}_X)$ 中的逆系统，其导出极限
$K = R\lim K_n$ 位于 $D(\mathcal{O}_X)$ 中。假设映射
$H^q(K_{n + 1}) \to H^q(K_n)$ 对所有 $q \in \mathbf{Z}$ 和
$n \geq 1$ 都是满射。则
\begin{enumerate}
\item $H^q(K) = \lim H^q(K_n)$；
\item $R\lim H^q(K_n) = \lim H^q(K_n)$；以及
\item 对每个仿射开集 $U \subset X$，有
$H^p(U, \lim H^q(K_n)) = 0$（当 $p > 0$）。
\end{enumerate}
\end{lemma}

\begin{proof}
令 $\mathcal{B} \subset \Ob(X_\etale)$ 为仿射对象所成的集合。由于
$H^q(K_n)$ 拟凝聚，故有 $H^p(U, H^q(K_n)) = 0$，其中
$U \in \mathcal{B}$；见“代数空间的上同调”第
\ref{spaces-cohomology-section-higher-direct-image} 节和“概形的上同调”引理
\ref{coherent-lemma-quasi-coherent-affine-cohomology-zero}。此外，由类似的理由，映射
$H^0(U, H^q(K_{n + 1})) \to H^0(U, H^q(K_n))$ 对
$U \in \mathcal{B}$ 是满射。第 (1) 项由“位点上的上同调”引理
\ref{sites-cohomology-lemma-derived-limit-suitable-system} 得出，因为我们刚验证了
该引理的条件。第 (2)、(3) 项由“位点上的上同调”引理
\ref{sites-cohomology-lemma-inverse-limit-is-derived-limit} 得出。
\end{proof}

\begin{lemma}
\label{spaces-perfect-lemma-quasi-coherence-pullback}
设 $S$ 为概形。设 $f : Y \to X$ 为 $S$ 上代数空间的态射。
函子 $Lf^*$ 把 $D_\QCoh(\mathcal{O}_X)$ 映入
$D_\QCoh(\mathcal{O}_Y)$。
\end{lemma}

\begin{proof}
选择一个图
$$
\xymatrix{
U \ar[d]_a \ar[r]_h & V \ar[d]^b \\
X \ar[r]^f & Y
}
$$
% P09-ERR-0393
其中 $U$ 和 $V$ 是概形，竖直箭头平展，并且 $a$ 是满射。由于
$a^* \circ Lf^* = Lh^* \circ b^*$，结论由引理
\ref{spaces-perfect-lemma-check-quasi-coherence-on-covering} 以及概形情形得出；概形情形即
“概形的导出范畴”引理
\ref{perfect-lemma-quasi-coherence-pullback}。
\end{proof}

\begin{lemma}
\label{spaces-perfect-lemma-quasi-coherence-tensor-product}
设 $S$ 为概形，$X$ 为 $S$ 上的代数空间。对对象 $K, L$
（属于 $D_\QCoh(\mathcal{O}_X)$），导出张量积
$K \otimes^\mathbf{L} L$ 属于 $D_\QCoh(\mathcal{O}_X)$。
\end{lemma}

\begin{proof}
设 $\varphi : U \to X$ 为从概形 $U$ 出发的满平展态射。由于
$\varphi^*(K \otimes_{\mathcal{O}_X}^\mathbf{L} L) =
\varphi^*K \otimes_{\mathcal{O}_U}^\mathbf{L} \varphi^*L$，由引理
\ref{spaces-perfect-lemma-check-quasi-coherence-on-covering} 可知，结论归结为概形情形；后者即
“概形的导出范畴”引理
\ref{perfect-lemma-quasi-coherence-tensor-product}。
\end{proof}

\noindent
下面的引理将帮助我们在 $D_\QCoh(\mathcal{O}_X)$ 的对象上“计算”右导出函子。

\begin{lemma}
\label{spaces-perfect-lemma-nice-K-injective}
设 $S$ 为概形，$X$ 为 $S$ 上的代数空间，$E$ 为
$D_\QCoh(\mathcal{O}_X)$ 的对象。则“导出范畴”评注
\ref{derived-remark-map-into-derived-limit-truncations} 中的映射
$E \to R\lim \tau_{\geq -n}E$ 是同构\footnote{特别地，
$E$ 有 K-单射代表；见“导出范畴”引理
\ref{derived-lemma-difficulty-K-injectives}。}。
\end{lemma}

\begin{proof}
以 $\mathcal{H}^i = H^i(E)$ 表示第 $i$ 个上同调层（对 $E$ 而言）。令
$\mathcal{B}$ 为 $X_\etale$ 的仿射对象所成的集合。于是
$H^p(U, \mathcal{H}^i) = 0$ 对所有 $p > 0$、所有
$i \in \mathbf{Z}$ 以及所有 $U \in \mathcal{B}$ 都成立，因为
$U$ 是仿射概形。见“代数空间的上同调”第
\ref{spaces-cohomology-section-higher-direct-image} 节和“概形的上同调”引理
\ref{coherent-lemma-quasi-coherent-affine-cohomology-zero}。
因此，本引理由“位点上的上同调”引理
\ref{sites-cohomology-lemma-is-limit-dimension}（取 $d = 0$）得出。
\end{proof}

\begin{lemma}
\label{spaces-perfect-lemma-application-nice-K-injective}
设 $S$ 为概形，$X$ 为 $S$ 上的代数空间。设
$F : \textit{Mod}(\mathcal{O}_X) \to \textit{Ab}$
为函子，$N \geq 0$ 为整数。假设
\begin{enumerate}
\item $F$ 左正合；
\item $F$ 与可数直积可交换；
\item $R^pF(\mathcal{F}) = 0$ 对所有 $p \geq N$ 及拟凝聚
$\mathcal{F}$ 都成立。
\end{enumerate}
则对 $E \in D_\QCoh(\mathcal{O}_X)$，有
\begin{enumerate}
\item $H^i(RF(\tau_{\leq a}E) \to H^i(RF(E))$ 在
$i \leq a$ 时为同构； % P09-ERR-0394
\item $H^i(RF(E)) \to H^i(RF(\tau_{\geq b - N + 1}E))$ 在
$i \geq b$ 时为同构；
\item 若 $H^i(E) = 0$ 在 $i \not \in [a, b]$ 时成立，其中某个
$-\infty \leq a \leq b \leq \infty$，则 $H^i(RF(E)) = 0$
在 $i \not \in [a, b + N - 1]$ 时成立。
\end{enumerate}
\end{lemma}

\begin{proof}
陈述 (1) 即“导出范畴”引理 \ref{derived-lemma-negative-vanishing}。

\medskip\noindent
证明陈述 (2)。写 $E_n = \tau_{\geq -n}E$。我们有
$E = R\lim E_n$，见引理 \ref{spaces-perfect-lemma-nice-K-injective}。于是由“单射对象”引理
\ref{injectives-lemma-RF-commutes-with-Rlim}，有
$RF(E) = R\lim RF(E_n)$（在 $D(\textit{Ab})$ 中）。因此对每个
$i \in \mathbf{Z}$，存在短正合列
$$
0 \to R^1\lim H^{i - 1}(RF(E_n)) \to H^i(RF(E)) \to \lim H^i(RF(E_n)) \to 0
$$
；见“代数进阶”评注
\ref{more-algebra-remark-compare-derived-limit}。为证明 (2)，我们将证明左端项为零，
且右端项等于 $H^i(RF(E_{-b + N - 1})$；这里 $b$ 任取并满足
$i \geq b$。 % P09-ERR-0395

\medskip\noindent
对每个 $n$，有特异三角
$$
H^{-n}(E)[n] \to E_n \to E_{n - 1} \to H^{-n}(E)[n + 1]
$$
（“导出范畴”评注
\ref{derived-remark-truncation-distinguished-triangle}），它位于
$D(\mathcal{O}_X)$ 中。由于 $H^{-n}(E)$ 拟凝聚，有
$$
H^i(RF(H^{-n}(E)[n])) = R^{i + n}F(H^{-n}(E)) = 0
$$
，当 $i + n \geq N$；并且有
$$
H^i(RF(H^{-n}(E)[n + 1])) = R^{i + n + 1}F(H^{-n}(E)) = 0
$$
，当 $i + n + 1 \geq N$。由此可知
$$
H^i(RF(E_n)) \to H^i(RF(E_{n - 1}))
$$
在 $n \geq N - i$ 时是同构。因此，各系统 $H^i(RF(E_n))$ 都满足
ML 条件，而短正合列中的 $R^1\lim$ 项为零（见“代数进阶”第
\ref{more-algebra-section-Rlim} 节中的讨论）。此外，系统
$H^i(RF(E_n))$ 从 $n = N - i - 1$ 起为常值系统，正如所需。

\medskip\noindent
证明 (3)。在关于 $E$ 的假设下，有 $\tau_{\leq a - 1}E = 0$，
由 (1) 得到 $H^i(RF(E))$ 在 $i \leq a - 1$ 时消失。类似地，我们有
$\tau_{\geq b + 1}E = 0$，因而由第 (2) 项得到
$H^i(RF(E))$ 在 $i \geq b + n$ 时消失。 % P09-ERR-0396
\end{proof}


\section{总正像}
\label{spaces-perfect-section-total-direct-image}

\noindent
下面的引理是“代数空间的上同调”引理
\ref{spaces-cohomology-lemma-vanishing-higher-direct-images} 的类似形式。

\begin{lemma}
\label{spaces-perfect-lemma-quasi-coherence-direct-image}
设 $S$ 为概形，$f : X \to Y$ 为 $S$ 上代数空间的拟分离拟紧态射。
\begin{enumerate}
\item 函子 $Rf_*$ 把 $D_\QCoh(\mathcal{O}_X)$ 映入
$D_\QCoh(\mathcal{O}_Y)$。
\item 若 $Y$ 拟紧，则存在整数 $N = N(X, Y, f)$，使得对对象 $E$
（属于 $D_\QCoh(\mathcal{O}_X)$），如果 $H^m(E) = 0$ 对
$m > 0$ 成立，那么 $H^m(Rf_*E) = 0$ 对 $m \geq N$ 成立。
\item 事实上，若 $Y$ 拟紧，则可找到 $N = N(X, Y, f)$，使得对代数空间的
每个态射 $Y' \to Y$，函子 $R(f')_*$ 都满足同一结论，其中
$f' : X' \to Y'$ 是 $f$ 的基变换。
\end{enumerate}
\end{lemma}

\begin{proof}
设 $E$ 为 $D_\QCoh(\mathcal{O}_X)$ 的对象。为证明 (1)，须证明
$Rf_*E$ 的上同调层拟凝聚。这个问题在 $Y$ 上是局部的，因此可以假设
$Y$ 拟紧。取“代数空间的上同调”引理
\ref{spaces-cohomology-lemma-vanishing-higher-direct-images} 中的
$N = N(X, Y, f)$。于是 $R^pf_*\mathcal{F} = 0$ 对所有拟凝聚
$\mathcal{O}_X$-模 $\mathcal{F}$ 及所有 $p \geq N$ 都成立。
此外，$R^pf_*\mathcal{F}$ 对所有 $p$ 均拟凝聚；见“代数空间的上同调”引理
\ref{spaces-cohomology-lemma-higher-direct-image}。这些陈述在基变换后仍成立。

\medskip\noindent
首先假设 $E$ 下有界。我们将证明，对这样的 $E$，以所选的 $N$，
(1)、(2)、(3) 均成立。在这种情形下，例如可以使用谱序列
$$
R^pf_*H^q(E) \Rightarrow R^{p + q}f_*E
$$
（“导出范畴”引理 \ref{derived-lemma-two-ss-complex-functor}），并使用
$R^pf_*H^q(E)$ 的拟凝聚性以及 $R^pf_*H^q(E)$ 在 $p \geq N$ 时的消失性，
从而看出 (1)、(2)、(3) 在这种情形下成立。

\medskip\noindent
接下来证明 (2)、(3)。设 $H^m(E) = 0$ 对 $m > 0$ 成立。令
$V$ 为 $Y_\etale$ 的仿射对象。我们有
$H^p(V \times_Y X, \mathcal{F}) = 0$（当 $p \geq N$）；见
“代数空间的上同调”引理
\ref{spaces-cohomology-lemma-quasi-coherence-higher-direct-images-application}。
因此可把引理 \ref{spaces-perfect-lemma-application-nice-K-injective} 应用于函子
$\Gamma(V \times_Y X, -)$，从而看出
$$
R\Gamma(V, Rf_*E) = R\Gamma(V \times_Y X, E)
$$
在次数 $\geq N$ 处的上同调消失。由于这对所有 $V$
（$Y_\etale$ 中仿射）均成立，可得 $H^m(Rf_*E) = 0$ 对
$m \geq N$ 成立。

\medskip\noindent
最后证明一般情形下的 (1)。回忆，有特异三角
$$
\tau_{\leq -n - 1}E \to E \to \tau_{\geq -n}E \to
(\tau_{\leq -n - 1}E)[1]
$$
，它位于 $D(\mathcal{O}_X)$ 中；见“导出范畴”评注
\ref{derived-remark-truncation-distinguished-triangle}。由 (2) 可知，
$Rf_*\tau_{\leq -n - 1}E$ 在次数 $\geq -n + N$ 处的上同调层消失。
因此，给定整数 $q$，可见 $R^qf_*E$ 等于
$R^qf_*\tau_{\geq -n}E$（对某个 $n$），于是可以应用上面的结果。
\end{proof}

\begin{lemma}
\label{spaces-perfect-lemma-quasi-coherence-pushforward-direct-sums}
设 $S$ 为概形，$f : X \to Y$ 为 $S$ 上代数空间的拟分离拟紧态射。则
$Rf_* : D_\QCoh(\mathcal{O}_X) \to D_\QCoh(\mathcal{O}_Y)$
与直和可交换。
\end{lemma}

\begin{proof}
设 $E_i$ 为 $D_\QCoh(\mathcal{O}_X)$ 的一族对象，并令
$E = \bigoplus E_i$。我们要证明映射
$$
\bigoplus Rf_*E_i \longrightarrow Rf_*E
$$
是同构。我们将证明它在次数 $0$ 的上同调层上诱导同构，这便蕴含引理。
选取引理 \ref{spaces-perfect-lemma-quasi-coherence-direct-image} 中的整数 $N$。
于是由所引引理，有 $R^0f_*E = R^0f_*\tau_{\geq -N}E$ 和
$R^0f_*E_i = R^0f_*\tau_{\geq -N}E_i$。注意
$\tau_{\geq -N}E = \bigoplus \tau_{\geq -N}E_i$。
因此可以假设所有 $E_i$ 在次数 $< -N$ 处的上同调层均消失。
接着使用谱序列
$$
R^pf_*H^q(E) \Rightarrow R^{p + q}f_*E
\quad\text{以及}\quad
R^pf_*H^q(E_i) \Rightarrow R^{p + q}f_*E_i
$$
（“导出范畴”引理 \ref{derived-lemma-two-ss-complex-functor}），把问题归结为
拟凝聚层的直和情形。该情形由“代数空间的上同调”引理
\ref{spaces-cohomology-lemma-colimit-cohomology} 处理。
\end{proof}

\begin{remark}
\label{spaces-perfect-remark-match-total-direct-images}
设 $S$ 为概形，$f : X \to Y$ 为可表代数空间 $X$ 与 $Y$
（在 $S$ 上）之间的态射。
设 $f_0 : X_0 \to Y_0$ 为表示 $f$ 的概形态射（这个记号很别扭，但只是临时的）。
则图
$$
\xymatrix{
D_\QCoh(\mathcal{O}_{X_0})
\ar@{=}[rrrrrr]_{\text{引理
\ref{spaces-perfect-lemma-derived-quasi-coherent-small-etale-site}}}
& & & & & &
D_\QCoh(\mathcal{O}_X) \\
D_\QCoh(\mathcal{O}_{Y_0})
\ar[u]^{Lf^*_0}
\ar@{=}[rrrrrr]^{\text{引理
\ref{spaces-perfect-lemma-derived-quasi-coherent-small-etale-site}}}
& & & & & &
D_\QCoh(\mathcal{O}_Y) \ar[u]_{Lf^*}
}
$$
可交换（引理 \ref{spaces-perfect-lemma-quasi-coherence-pullback} 和
“概形的导出范畴”引理
\ref{perfect-lemma-quasi-coherence-pullback}）。这是因为引理
\ref{spaces-perfect-lemma-derived-quasi-coherent-small-etale-site} 中的等价
$D_\QCoh(\mathcal{O}_{X_0}) \to D_\QCoh(\mathcal{O}_X)$ 和
$D_\QCoh(\mathcal{O}_{Y_0}) \to D_\QCoh(\mathcal{O}_Y)$
来自沿带环位点的（平坦）态射
$\epsilon : X_\etale \to X_{0, Zar}$ 和
$\epsilon : Y_\etale \to Y_{0, Zar}$ 的拉回，并且带环位点的图
$$
\xymatrix{
X_{0, Zar} \ar[d]_{f_0} & X_\etale \ar[l]^\epsilon \ar[d]^f \\
Y_{0, Zar} & Y_\etale \ar[l]_\epsilon
}
$$
可交换（略去细节）。若 $f$ 拟紧且拟分离——等价地，若 $f_0$ 拟紧且
拟分离——则我们断言
$$
\xymatrix{
D_\QCoh(\mathcal{O}_{X_0})
\ar[d]_{Rf_{0, *}} \ar@{=}[rrrrrr]_{\text{引理
\ref{spaces-perfect-lemma-derived-quasi-coherent-small-etale-site}}}
& & & & & &
D_\QCoh(\mathcal{O}_X) \ar[d]^{Rf_*} \\
D_\QCoh(\mathcal{O}_{Y_0})
\ar@{=}[rrrrrr]^{\text{引理
\ref{spaces-perfect-lemma-derived-quasi-coherent-small-etale-site}}}
& & & & & &
D_\QCoh(\mathcal{O}_Y)
}
$$
也可交换（引理 \ref{spaces-perfect-lemma-quasi-coherence-direct-image} 和
“概形的导出范畴”引理
\ref{perfect-lemma-quasi-coherence-direct-image}）。这也由上面所示的位点交换图得出，
因为引理 \ref{spaces-perfect-lemma-derived-quasi-coherent-small-etale-site} 的证明表明，函子
$R\epsilon_*$ 给出等价
$D_\QCoh(\mathcal{O}_X) \to D_\QCoh(\mathcal{O}_{X_0})$ 和
$D_\QCoh(\mathcal{O}_Y) \to D_\QCoh(\mathcal{O}_{Y_0})$。
\end{remark}

\begin{lemma}
\label{spaces-perfect-lemma-affine-morphism}
设 $S$ 为概形，$f : X \to Y$ 为 $S$ 上代数空间的仿射态射。则
$Rf_* : D_\QCoh(\mathcal{O}_X) \to D_\QCoh(\mathcal{O}_Y)$
反映同构。
\end{lemma}

\begin{proof}
该陈述的含义是：态射 $\alpha : E \to F$
（属于 $D_\QCoh(\mathcal{O}_X)$）若满足 $Rf_*\alpha$ 是同构，
则其本身也是同构。
这可以在上同调层上检验。特别地，这个问题在 $Y$ 上是平展局部的。
因此可以假设 $Y$，从而 $X$，是仿射的。在这种情形下，问题归结为概形情形
（“概形的导出范畴”引理 \ref{perfect-lemma-affine-morphism}）；这一归结通过
引理 \ref{spaces-perfect-lemma-derived-quasi-coherent-small-etale-site} 和评注
\ref{spaces-perfect-remark-match-total-direct-images} 实现。
\end{proof}

\begin{lemma}
\label{spaces-perfect-lemma-affine-morphism-pull-push}
设 $S$ 为概形，$f : X \to Y$ 为 $S$ 上代数空间的仿射态射。对 $E$
（$D_\QCoh(\mathcal{O}_Y)$ 中的对象），有
$Rf_* Lf^* E = E \otimes^\mathbf{L}_{\mathcal{O}_Y} f_*\mathcal{O}_X$。
\end{lemma}

\begin{proof}
由于 $f$ 仿射，映射 $f_*\mathcal{O}_X \to Rf_*\mathcal{O}_X$
是同构（“代数空间的上同调”引理
\ref{spaces-cohomology-lemma-affine-vanishing-higher-direct-images}）。
存在典范映射
$E \otimes^\mathbf{L} f_*\mathcal{O}_X =
E \otimes^\mathbf{L} Rf_*\mathcal{O}_X \to Rf_* Lf^* E$，
它与映射
$$
Lf^*(E \otimes^\mathbf{L} Rf_*\mathcal{O}_X) =
Lf^*E \otimes^\mathbf{L} Lf^*Rf_*\mathcal{O}_X \longrightarrow
Lf^* E \otimes^\mathbf{L} \mathcal{O}_X = Lf^* E
$$
伴随；后一个映射来自 $1 : Lf^*E \to Lf^*E$ 和典范映射
$Lf^*Rf_*\mathcal{O}_X \to \mathcal{O}_X$。为检验所构造的映射是同构，
可以在 $Y$ 上局部地工作。因此可以假设 $Y$，从而 $X$，是仿射的。
在这种情形下，问题归结为概形情形（“概形的导出范畴”引理
\ref{perfect-lemma-affine-morphism-pull-push}）；这一归结通过引理
\ref{spaces-perfect-lemma-derived-quasi-coherent-small-etale-site} 和评注
\ref{spaces-perfect-remark-match-total-direct-images} 实现。
\end{proof}


\section{在基底上为固有}
\label{spaces-perfect-section-proper-over-base}

\noindent
本节是“概形的上同调”第
\ref{coherent-section-proper-over-base} 节的类似形式。和通常一样，
只要材料涉及点集上的拓扑，把它转移到代数空间时就必须谨慎。

\begin{lemma}
\label{spaces-perfect-lemma-closed-proper-over-base}
设 $S$ 为概形，$f : X \to Y$ 为 $S$ 上代数空间的局部有限型态射。
设 $T \subset |X|$ 为闭子集。下列条件等价：
\begin{enumerate}
\item 态射 $Z \to Y$ 固有，如果 $Z$ 是 $T$ 上的约化诱导代数空间结构
（“代数空间的性质”定义
\ref{spaces-properties-definition-reduced-induced-space}）；
\item 对某个闭子空间 $Z \subset X$，若 $|Z| = T$，则态射
$Z \to Y$ 固有；以及
\item 对任意闭子空间 $Z \subset X$，若 $|Z| = T$，则态射
$Z \to Y$ 固有。
\end{enumerate}
\end{lemma}

\begin{proof}
蕴含关系 (3) $\Rightarrow$ (1) 和 (1) $\Rightarrow$ (2) 是直接的。
故只需证明 (2) 蕴含 (3)。我们力劝读者自行找出这个事实的证明。
设 $Z'$ 和 $Z''$ 为闭子空间，满足 $T = |Z'| = |Z''|$，并且
$Z' \to Y$ 是代数空间的固有态射。须证明 $Z'' \to Y$ 也固有。
令 $Z''' = Z' \cup Z''$ 为概形论并；见“代数空间的态射”定义
\ref{spaces-morphisms-definition-scheme-theoretic-intersection-union}。
于是 $Z'''$ 是另一个满足 $|Z'''| = T$ 的闭子空间。例如，这由
“代数空间的态射”引理
\ref{spaces-morphisms-lemma-scheme-theoretic-union} 中概形论并的描述得出。
由于 $Z'' \to Z'''$ 是闭浸入，只需证明 $Z''' \to Y$ 固有
（见“代数空间的态射”引理
\ref{spaces-morphisms-lemma-closed-immersion-proper} 和
\ref{spaces-morphisms-lemma-composition-proper}）。态射
$Z' \to Z'''$ 是双射的闭浸入，特别地，它是满的且普遍闭的。
于是 $Z' \to Y$ 分离这一事实蕴含 $Z''' \to Y$ 分离；见
“代数空间的态射”引理
\ref{spaces-morphisms-lemma-image-universally-closed-separated}。
此外，$Z''' \to Y$ 局部有限型，因为 $X \to Y$ 局部有限型
（“代数空间的态射”引理
\ref{spaces-morphisms-lemma-immersion-locally-finite-type} 和
\ref{spaces-morphisms-lemma-composition-finite-type}）。由于
$Z' \to Y$ 拟紧，且 $Z' \to Z'''$ 是普遍同胚，可见
$Z''' \to Y$ 拟紧。最后，因为 $Z' \to Y$ 普遍闭，由
“代数空间的态射”引理
\ref{spaces-morphisms-lemma-image-proper-is-proper} 可见
$Z''' \to Y$ 也普遍闭。这就完成了证明。
\end{proof}

\begin{definition}
\label{spaces-perfect-definition-proper-over-base}
设 $S$ 为概形。设 $f : X \to Y$ 为 $S$ 上代数空间的局部有限型态射。
设 $T \subset |X|$ 为闭子集。如果引理
\ref{spaces-perfect-lemma-closed-proper-over-base} 中的等价条件成立，我们称
{\it $T$ 在 $Y$ 上固有}。
\end{definition}

\noindent
若态射 $f : X \to Y$ 不是局部有限型的，则上面定义中使用的引理不成立。
因此我们力劝读者：若 $f$ 不是局部有限型的，就不要使用这个术语。

\begin{lemma}
\label{spaces-perfect-lemma-closed-closed-proper-over-base}
设 $S$ 为概形。设 $f : X \to Y$ 为 $S$ 上代数空间的局部有限型态射。
设 $T' \subset T \subset |X|$ 为闭子集。若 $T$ 在 $Y$ 上固有，
则 $T'$ 也是如此。
\end{lemma}

\begin{proof}
略。
\end{proof}

\begin{lemma}
\label{spaces-perfect-lemma-base-change-closed-proper-over-base}
设 $S$ 为概形。考虑 $S$ 上代数空间的笛卡尔图
$$
\xymatrix{
X' \ar[d]_{f'} \ar[r]_{g'} & X \ar[d]^f \\
Y' \ar[r]^g & Y
}
$$
，其中 $f$ 局部有限型。若 $T$ 是 $|X|$ 的闭子集并在 $Y$ 上固有，
则 $|g'|^{-1}(T)$ 是 $|X'|$ 的闭子集并在 $Y'$ 上固有。
\end{lemma}

\begin{proof}
注意，由“代数空间的态射”引理
\ref{spaces-morphisms-lemma-base-change-finite-type}，$f'$ 局部有限型，
故该陈述是有意义的。设 $Z \subset X$ 为 $T$ 上的约化诱导闭子空间结构。
以 $Z' = (g')^{-1}(Z)$ 表示概形论逆像。于是
$Z' = X' \times_X Z = (Y' \times_Y X) \times_X Z = Y' \times_Y Z$
在 $Y'$ 上固有，因为它是 $Z$ 在 $Y$ 上的基变换
（“代数空间的态射”引理
\ref{spaces-morphisms-lemma-base-change-proper}）。另一方面，
我们有 $T' = |Z'|$。 % P09-ERR-0397
故引理成立。
\end{proof}

\begin{lemma}
\label{spaces-perfect-lemma-functoriality-closed-proper-over-base}
设 $S$ 为概形，$B$ 为 $S$ 上的代数空间。设 $f : X \to Y$ 为代数空间的态射，
其源与靶均在 $B$ 上局部有限型。
\begin{enumerate}
\item 若 $Y$ 在 $B$ 上分离，且 $T \subset |X|$ 是在 $B$ 上固有的闭子集，
则 $|f|(T)$ 是 $|Y|$ 的闭子集并在 $B$ 上固有。
\item 若 $f$ 普遍闭，且 $T \subset |X|$ 是在 $B$ 上固有的闭子集，
则 $|f|(T)$ 是 $Y$ 的闭子集并在 $B$ 上固有。 % P09-ERR-0398
\item 若 $f$ 固有，且 $T \subset |Y|$ 是在 $B$ 上固有的闭子集，
则 $|f|^{-1}(T)$ 是 $|X|$ 的闭子集并在 $B$ 上固有。
\end{enumerate}
\end{lemma}

\begin{proof}
证明 (1)。假设 $Y$ 在 $B$ 上分离，且 $T \subset |X|$ 是在 $B$ 上固有的
闭子集。设 $Z$ 为 $T$ 上的约化诱导闭子空间结构，并把
“代数空间的态射”引理
\ref{spaces-morphisms-lemma-scheme-theoretic-image-is-proper}
应用于 $Z \to Y$（在 $B$ 上），即得结论。

\medskip\noindent
证明 (2)。假设 $f$ 普遍闭，且 $T \subset |X|$ 是在 $B$ 上固有的闭子集。
设 $Z$ 为 $T$ 上的约化诱导闭子空间结构，$Z'$ 为 $|f|(T)$ 上的
约化诱导闭子空间结构。于是得到诱导态射 $Z \to Z'$。
以 $Z'' = f^{-1}(Z')$ 表示概形论逆像。由于 $Z'' \to Z'$ 是 $f$ 的基变换，
它普遍闭（“代数空间的态射”引理
\ref{spaces-morphisms-lemma-base-change-proper}）。因此 $Z \to Z'$ 作为闭浸入
$Z \to Z''$ 与 $Z'' \to Z'$ 的复合而普遍闭
（“代数空间的态射”引理
\ref{spaces-morphisms-lemma-closed-immersion-proper} 和
\ref{spaces-morphisms-lemma-composition-proper}）。由“代数空间的态射”引理
\ref{spaces-morphisms-lemma-image-universally-closed-separated}，可得
$Z' \to B$ 分离。由于 $Z \to B$ 拟紧且 $Z \to Z'$ 满，可见
$Z' \to B$ 拟紧。因为 $Z' \to B$ 是 $Z' \to Y$ 与 $Y \to B$ 的复合，
可见 $Z' \to B$ 局部有限型
（“代数空间的态射”引理
\ref{spaces-morphisms-lemma-immersion-locally-finite-type} 和
\ref{spaces-morphisms-lemma-composition-finite-type}）。最后，因为
$Z \to B$ 普遍闭，由“代数空间的态射”引理
\ref{spaces-morphisms-lemma-image-proper-is-proper} 可见
$Z' \to B$ 也普遍闭。这就完成了证明。

\medskip\noindent
证明 (3)。假设 $f$ 固有，且 $T \subset |Y|$ 是在 $B$ 上固有的闭子集。
设 $Z$ 为 $T$ 上的约化诱导闭子空间结构。以
$Z' = f^{-1}(Z)$ 表示概形论逆像。由于 $Z' \to Z$ 是 $f$ 的基变换，
它固有（“代数空间的态射”引理
\ref{spaces-morphisms-lemma-base-change-proper}）。从而 $Z' \to B$ 作为
$Z' \to Z$ 与 $Z \to B$ 的复合而固有
（“代数空间的态射”引理
\ref{spaces-morphisms-lemma-composition-proper}）。这就完成了证明。
\end{proof}

\begin{lemma}
\label{spaces-perfect-lemma-union-closed-proper-over-base}
设 $S$ 为概形。设 $f : X \to Y$ 为 $S$ 上代数空间的局部有限型态射。
设 $T_i \subset |X|$（$i = 1, \ldots, n$）为闭子集。若
$T_i$（$i = 1, \ldots, n$）均在 $Y$ 上固有，则
$T_1 \cup \ldots \cup T_n$ 也是如此。
\end{lemma}

\begin{proof}
设 $Z_i$ 为 $T_i$ 上的约化诱导闭子概形结构。 % P09-ERR-0399
态射
$$
Z_1 \amalg \ldots \amalg Z_n \longrightarrow X
$$
是有限的；见“代数空间的态射”引理
\ref{spaces-morphisms-lemma-closed-immersion-finite} 和
\ref{spaces-morphisms-lemma-finite-union-finite}。有限态射普遍闭
（“代数空间的态射”引理 \ref{spaces-morphisms-lemma-finite-proper}），
并且 $Z_1 \amalg \ldots \amalg Z_n$ 在 $S$ 上固有， % P09-ERR-0400
故由引理 \ref{spaces-perfect-lemma-functoriality-closed-proper-over-base} 第 (2) 项可得，
像 $Z_1 \cup \ldots \cup Z_n$ 在 $S$ 上固有。
\end{proof}

\noindent
设 $S$ 为概形。设 $f : X \to Y$ 为 $S$ 上代数空间的局部有限型态射。
设 $\mathcal{F}$ 为有限型拟凝聚 $\mathcal{O}_X$-模。于是支撑
$\text{Supp}(\mathcal{F})$（属于 $\mathcal{F}$）是 $|X|$ 的闭子集；
见“代数空间的态射”引理
\ref{spaces-morphisms-lemma-support-finite-type}。因此，说
“$\mathcal{F}$ 的支撑在 $Y$ 上固有”是有意义的。

\begin{lemma}
\label{spaces-perfect-lemma-module-support-proper-over-base}
设 $S$ 为概形，$f : X \to Y$ 为 $S$ 上代数空间的局部有限型态射。
设 $\mathcal{F}$ 为有限型拟凝聚 $\mathcal{O}_X$-模。下列条件等价：
\begin{enumerate}
\item $\mathcal{F}$ 的支撑在 $Y$ 上固有；
\item $\mathcal{F}$ 的概形论支撑（“代数空间的态射”定义
\ref{spaces-morphisms-definition-scheme-theoretic-support}）在 $Y$ 上固有；以及
\item 存在闭子空间 $Z \subset X$ 和有限型拟凝聚
$\mathcal{O}_Z$-模 $\mathcal{G}$，使得 (a) $Z \to Y$ 固有，并且 (b)
$(Z \to X)_*\mathcal{G} = \mathcal{F}$。
\end{enumerate}
\end{lemma}

\begin{proof}
支撑 $\text{Supp}(\mathcal{F})$（属于 $\mathcal{F}$）是 $|X|$ 的闭子集；
见“代数空间的态射”引理
\ref{spaces-morphisms-lemma-support-finite-type}。故可应用定义
\ref{spaces-perfect-definition-proper-over-base}。由于 $\mathcal{F}$ 的概形论支撑是闭子空间，
其底闭子集为 $\text{Supp}(\mathcal{F})$，由定义
\ref{spaces-perfect-definition-proper-over-base} 可见 (1)、(2) 等价。显然 (2) 蕴含 (3)。
反过来，若 (3) 成立，则
$\text{Supp}(\mathcal{F}) \subset |Z|$，因而例如由引理
\ref{spaces-perfect-lemma-closed-closed-proper-over-base}，$\text{Supp}(\mathcal{F})$
在 $Y$ 上固有。
\end{proof}

\begin{lemma}
\label{spaces-perfect-lemma-base-change-module-support-proper-over-base}
设 $S$ 为概形。考虑 $S$ 上代数空间的笛卡尔图
$$
\xymatrix{
X' \ar[d]_{f'} \ar[r]_{g'} & X \ar[d]^f \\
Y' \ar[r]^g & Y
}
$$
，其中 $f$ 局部有限型。设 $\mathcal{F}$ 为有限型拟凝聚
$\mathcal{O}_X$-模。若 $\mathcal{F}$ 的支撑在 $Y$ 上固有，则
$(g')^*\mathcal{F}$ 的支撑在 $Y'$ 上固有。
\end{lemma}

\begin{proof}
注意，该陈述是有意义的，因为 $(g')*\mathcal{F}$ 由“位点上的模”引理
\ref{sites-modules-lemma-local-pullback} 为有限型。 % P09-ERR-0401
由“代数空间的态射”引理
\ref{spaces-morphisms-lemma-support-finite-type}，有
$\text{Supp}((g')^*\mathcal{F}) = |g'|^{-1}(\text{Supp}(\mathcal{F}))$。
于是本引理由引理
\ref{spaces-perfect-lemma-base-change-closed-proper-over-base} 得出。
\end{proof}

\begin{lemma}
\label{spaces-perfect-lemma-cat-module-support-proper-over-base}
设 $S$ 为概形。设 $f : X \to Y$ 为 $S$ 上代数空间的局部有限型态射。
设 $\mathcal{F}$、$\mathcal{G}$ 为有限型拟凝聚 $\mathcal{O}_X$-模。
\begin{enumerate}
\item 若 $\mathcal{F}$、$\mathcal{G}$ 的支撑在 $Y$ 上固有，则下列各模也具有
这一性质：$\mathcal{F} \oplus \mathcal{G}$；
$\mathcal{G}$ 被 $\mathcal{F}$ 扩张所得的任意扩张；
$\Im(u)$ 和 $\Coker(u)$（给定任意 $\mathcal{O}_X$-模映射
$u : \mathcal{F} \to \mathcal{G}$）；以及 $\mathcal{F}$ 或
$\mathcal{G}$ 的任意拟凝聚商。
\item 若 $Y$ 局部诺特，则凝聚 $\mathcal{O}_X$-模中支撑在 $Y$ 上固有者
所成的范畴，是凝聚 $\mathcal{O}_X$-模的阿贝尔范畴的
Serre 子范畴（“同调”定义
\ref{homology-definition-serre-subcategory}）。
\end{enumerate}
\end{lemma}

\begin{proof}
证明 (1)。设 $T$、$T'$ 分别为 $\mathcal{F}$、$\mathcal{G}$ 的支撑。
则 (1) 中提到的所有层，其支撑均包含于 $T \cup T'$。因此，只要检验这些层
有限型且拟凝聚，断言本身便由引理
\ref{spaces-perfect-lemma-closed-closed-proper-over-base} 和
\ref{spaces-perfect-lemma-union-closed-proper-over-base} 直接得出。关于拟凝聚性，参见
“代数空间的性质”第 \ref{spaces-properties-section-quasi-coherent} 节。
关于“有限型”，参见“代数空间的性质”第
\ref{spaces-properties-section-properties-modules} 节。

\medskip\noindent
证明 (2)。证明与 (1) 的证明相同。注意，这些断言是有意义的，因为由
“代数空间的态射”引理
\ref{spaces-morphisms-lemma-locally-finite-type-locally-noetherian}，
$X$ 局部诺特；还要使用“代数空间的上同调”第
\ref{spaces-cohomology-section-coherent} 节中对凝聚模范畴的描述。
\end{proof}

\begin{lemma}
\label{spaces-perfect-lemma-support-proper-over-base-pushforward}
设 $S$ 为概形，$f : X \to Y$ 为 $S$ 上代数空间的态射。假设 $f$ 局部有限型，
且 $Y$ 局部诺特。设 $\mathcal{F}$ 为凝聚 $\mathcal{O}_X$-模，其支撑在
$Y$ 上固有。则 $R^pf_*\mathcal{F}$ 是凝聚 $\mathcal{O}_Y$-模，
对所有 $p \geq 0$ 都成立。
\end{lemma}

\begin{proof}
由引理 \ref{spaces-perfect-lemma-module-support-proper-over-base}，存在闭浸入
$i : Z \to X$，使得 $g = f \circ i : Z \to Y$ 固有，并且有
$\mathcal{F} = i_*\mathcal{G}$，其中 $\mathcal{G}$ 是凝聚模（在 $Z$ 上）。
由“代数空间的上同调”引理
\ref{spaces-cohomology-lemma-proper-pushforward-coherent}，可见
$R^pg_*\mathcal{G}$ 在 $S$ 上凝聚。 % P09-ERR-0402
另一方面，有 $R^qi_*\mathcal{G} = 0$（当 $q > 0$）
（“代数空间的上同调”引理
\ref{spaces-cohomology-lemma-finite-pushforward-coherent}）。
由“位点上的上同调”引理
\ref{sites-cohomology-lemma-relative-Leray}，得到
$R^pf_*\mathcal{F} = R^pg_*\mathcal{G}$，引理得证。
\end{proof}

\section{凝聚模的导出范畴}
\label{spaces-perfect-section-derived-coherent}

\noindent
设 $S$ 为概形。设 $X$ 为 $S$ 上的局部诺特代数空间。
此时，范畴
$\textit{Coh}(\mathcal{O}_X) \subset \textit{Mod}(\mathcal{O}_X)$
（即凝聚 $\mathcal{O}_X$-模所成的范畴）是一个弱 Serre 子范畴；参见“同调”第
\ref{homology-section-serre-subcategories} 节以及
“代数空间的上同调”引理
\ref{spaces-cohomology-lemma-coherent-abelian-Noetherian}。
记
$$
D_{\textit{Coh}}(\mathcal{O}_X) \subset D(\mathcal{O}_X)
$$
为上同调层均凝聚的复形所成的子范畴；参见“导出范畴”第
\ref{derived-section-triangulated-sub} 节。因此得到典范函子
\begin{equation}
\label{spaces-perfect-equation-compare-coherent}
D(\textit{Coh}(\mathcal{O}_X))
\longrightarrow
D_{\textit{Coh}}(\mathcal{O}_X)
\end{equation}
参见“导出范畴”等式 (\ref{derived-equation-compare})。

\begin{lemma}
\label{spaces-perfect-lemma-direct-image-coherent}
设 $S$ 为概形。设 $f : X \to Y$ 为 $S$ 上代数空间的态射。
假设 $f$ 局部有限型，且 $Y$ 诺特。设 $E$ 为
$D^b_{\textit{Coh}}(\mathcal{O}_X)$ 的对象，并且
$H^i(E)$ 的支撑在 $Y$ 上固有，对所有 $i$ 均成立。则 $Rf_*E$ 是
$D^b_{\textit{Coh}}(\mathcal{O}_Y)$ 的对象。
\end{lemma}

\begin{proof}
考虑谱序列
$$
R^pf_*H^q(E) \Rightarrow R^{p + q}f_*E
$$
参见“导出范畴”引理 \ref{derived-lemma-two-ss-complex-functor}。
由假设和引理 \ref{spaces-perfect-lemma-support-proper-over-base-pushforward}，
层 $R^pf_*H^q(E)$ 是凝聚的。因此
$R^{p + q}f_*E$ 凝聚，即 $E \in D_{\textit{Coh}}(\mathcal{O}_Y)$。 % P09-ERR-0403
下有界性显然。上有界性由“代数空间的上同调”引理
\ref{spaces-cohomology-lemma-vanishing-higher-direct-images} 推出，
亦可由引理 \ref{spaces-perfect-lemma-quasi-coherence-direct-image} 推出。
\end{proof}

\begin{lemma}
\label{spaces-perfect-lemma-direct-image-coherent-bdd-below}
设 $S$ 为概形。设 $f : X \to Y$ 为 $S$ 上代数空间的态射。
假设 $f$ 局部有限型，且 $Y$ 诺特。设 $E$ 为
$D^+_{\textit{Coh}}(\mathcal{O}_X)$ 的对象，并且
$H^i(E)$ 的支撑在 $S$ 上固有，对所有 $i$ 均成立。 % P09-ERR-0404
则 $Rf_*E$ 是 $D^+_{\textit{Coh}}(\mathcal{O}_Y)$ 的对象。
\end{lemma}

\begin{proof}
证明与引理 \ref{spaces-perfect-lemma-direct-image-coherent} 的证明相同。
考察正合函子 $Rf_*$ 对特异三角
$\tau_{\leq a}E \to E \to \tau_{\geq a + 1}E \to \tau_{\leq a}E[1]$
的作用，也可由引理 \ref{spaces-perfect-lemma-direct-image-coherent} 推出此结论。
\end{proof}

\begin{lemma}
\label{spaces-perfect-lemma-coherent-internal-hom}
设 $S$ 为概形。设 $X$ 为 $S$ 上的局部诺特代数空间。
若 $L$ 属于 $D^+_{\textit{Coh}}(\mathcal{O}_X)$，且
$K$ 属于 $D^-_{\textit{Coh}}(\mathcal{O}_X)$，则
$R\SheafHom(K, L)$ 属于 $D^+_{\textit{Coh}}(\mathcal{O}_X)$。
\end{lemma}

\begin{proof}
一个 $D(\mathcal{O}_X)$ 的对象是否属于
$D_{\textit{Coh}}(\mathcal{O}_X)$，可在 $X$ 上平展局部地检验；
参见“代数空间的上同调”引理
\ref{spaces-cohomology-lemma-coherent-Noetherian}。
因此，本引理由概形情形推出；参见“概形的导出范畴”引理
\ref{perfect-lemma-coherent-internal-hom}。
\end{proof}

\begin{lemma}
\label{spaces-perfect-lemma-ext-finite}
设 $A$ 为诺特环。设 $X$ 为 $A$ 上的固有代数空间。
若 $L$ 属于 $D^+_{\textit{Coh}}(\mathcal{O}_X)$，且 $K$ 属于
$D^-_{\textit{Coh}}(\mathcal{O}_X)$，则 $A$-模
$\Ext_{\mathcal{O}_X}^n(K, L)$ 均有限生成。
\end{lemma}

\begin{proof}
回顾
$$
\Ext_{\mathcal{O}_X}^n(K, L) =
H^n(X, R\SheafHom_{\mathcal{O}_X}(K, L)) =
H^n(\Spec(A), Rf_*R\SheafHom_{\mathcal{O}_X}(K, L))
$$
参见“位点上的上同调”引理
\ref{sites-cohomology-lemma-section-RHom-over-U} 以及“位点上的上同调”第
\ref{sites-cohomology-section-leray} 节。因此，结论由引理
\ref{spaces-perfect-lemma-coherent-internal-hom} 和
\ref{spaces-perfect-lemma-direct-image-coherent-bdd-below} 推出。
\end{proof}

\section{归纳原理}
\label{spaces-perfect-section-induction}

\noindent
本节讨论代数空间的一条归纳原理，它类似于概形情形的“概形的上同调”引理
\ref{coherent-lemma-induction-principle}。为表述这条原理，我们引入
{\it 初等特异方形}的概念；这一术语取自 \cite{MV}。
这里所表述的原理隐含于 Raynaud 与 Gruson 的论文 \cite{GruRay} 中。
David Rydh 的 \cite[Theorem D]{rydh_etale_devissage}
给出了代数栈的一条相关原理。

\begin{definition}
\label{spaces-perfect-definition-elementary-distinguished-square}
设 $S$ 为概形。交换图
$$
\xymatrix{
U \times_W V \ar[r] \ar[d] & V \ar[d]^f \\
U \ar[r]^j & W
}
$$
是 $S$ 上代数空间的图，称为{\it 初等特异方形}，如果
\begin{enumerate}
\item $U$ 是 $W$ 的开子空间，且 $j$ 是包含态射；
\item $f$ 平展；并且
\item 令 $T = W \setminus U$（赋予约化诱导子空间结构），则态射
$f^{-1}(T) \to T$ 是同构。
\end{enumerate}
我们将用如下说法表示这一点：“设 $(U \subset W, f : V \to W)$
为初等特异方形。”
\end{definition}

\noindent
注意，若 $(U \subset W, f : V \to W)$ 为初等特异方形，则
$W = U \cup f(V)$。因此，$\{U \to W, V \to W\}$ 是 $W$ 的平展覆盖。
事实表明，这些平展覆盖具有良好性质，而且在某种意义下它们“足够多”。

\begin{lemma}
\label{spaces-perfect-lemma-make-more-elementary-distinguished-squares}
设 $S$ 为概形。设 $(U \subset W, f : V \to W)$ 为 $S$ 上代数空间的
初等特异方形。
\begin{enumerate}
\item 若 $V' \subset V$ 和 $U \subset U' \subset W$ 是开子空间，且
$W' = U' \cup f(V')$，则
$(U' \subset W', f|_{V'} : V' \to W')$ 是初等特异方形。
\item 若 $p : W' \to W$ 是代数空间的态射，则
$(p^{-1}(U) \subset W', V \times_W W' \to W')$ 是初等特异方形。
\item 若 $S' \to S$ 是概形的态射，则
$(S' \times_S U \subset S' \times_S W, S' \times_S V \to S' \times_S W)$
是初等特异方形。
\end{enumerate}
\end{lemma}

\begin{proof}
略。
\end{proof}

\begin{lemma}
\label{spaces-perfect-lemma-induction-principle}
设 $S$ 为概形。设 $X$ 为 $S$ 上拟紧且拟分离的代数空间。设 $P$ 是
$X_{spaces, \etale}$ 的拟紧拟分离对象的一项性质。假设
\begin{enumerate}
\item $P$ 对 $X_{spaces, \etale}$ 的每个仿射对象成立；
\item 对每个初等特异方形 $(U \subset W, f : V \to W)$，若
\begin{enumerate}
\item $W$ 是 $X_{spaces, \etale}$ 的拟紧拟分离对象；
\item $U$ 拟紧；
\item $V$ 仿射；并且
\item $P$ 对 $U$、$V$ 和 $U \times_W V$ 成立，
\end{enumerate}
则 $P$ 对 $W$ 成立。
\end{enumerate}
于是，$P$ 对 $X_{spaces, \etale}$ 的每个拟紧拟分离对象成立，特别地对
$X$ 成立。
\end{lemma}

\begin{proof}
首先断言，$P$ 对 $X_{spaces, \etale}$ 的每个可表拟紧拟分离对象成立。
事实上，设 $U \to X$ 平展，且 $U$ 是拟紧拟分离概形。由假设 (1)，
性质 $P$ 对 $U$ 的每个仿射开集成立。此外，若 $W, V \subset U$
是拟紧开集，$V$ 仿射，而且 $P$ 对 $W$、$V$ 和 $W \cap V$ 成立，
则由 (2)，$P$ 对 $W \cup V$ 成立（因为
$(W \subset W \cup V, V \to W \cup V)$ 这一对是初等特异方形）。
因此，由概形的归纳原理，$P$ 对 $U$ 成立；参见“概形的上同调”引理
\ref{coherent-lemma-induction-principle}。

\medskip\noindent
为完成证明，只需证明 $P$ 对 $X$ 成立（因为可将当前的 $X$ 替换为
$X_{spaces, \etale}$ 中我们想证明结论的任意拟紧拟分离对象）。
我们将使用“良态空间”引理
\ref{decent-spaces-lemma-filter-quasi-compact-quasi-separated} 中的滤过
$$
\emptyset = U_{n + 1} \subset
U_n \subset U_{n - 1} \subset \ldots \subset U_1 = X
$$
以及态射 $f_p : V_p \to U_p$。我们将证明 $P$ 对 $U_p$ 成立，并对
$p$ 作降序归纳。注意，由 (1)，$P$ 对 $U_{n + 1}$ 成立，
因为空代数空间是仿射的。假设 $P$ 对 $U_{p + 1}$ 成立。
注意，$(U_{p + 1} \subset U_p, f_p : V_p \to U_p)$ 是初等特异方形，
但 (2) 可能不适用，因为 $V_p$ 未必仿射。然而，由于 $V_p$ 是拟紧概形，
可选取有限仿射开覆盖
$V_p = V_{p, 1} \cup \ldots \cup V_{p, m}$。令
$W_{p, 0} = U_{p + 1}$，并令
$$
W_{p, i} = U_{p + 1} \cup f_p(V_{p, 1} \cup \ldots \cup V_{p, i})
$$
其中 $i = 1, \ldots, m$。这些是 $X$ 的拟紧开子空间。于是
$$
U_{p + 1} = W_{p, 0} \subset
W_{p, 1} \subset \ldots \subset
W_{p, m} = U_p
$$
且各对
$$
(W_{p, 0} \subset W_{p, 1}, f_p|_{V_{p, 1}}),
(W_{p, 1} \subset W_{p, 2}, f_p|_{V_{p, 2}}),\ldots,
(W_{p, m - 1} \subset W_{p, m}, f_p|_{V_{p, m}})
$$
由引理 \ref{spaces-perfect-lemma-make-more-elementary-distinguished-squares}
均为初等特异方形。注意，$P$ 对每个 $V_{p, 1}$ 均成立（因为它们是
仿射概形）； % P09-ERR-0405
而 $W_{p, i} \times_{W_{p, i + 1}} V_{p, i + 1}$ 是
$V_{p, i + 1}$ 的拟紧开集，所以由本证明第一段，$P$ 对它也成立。
因此，(2) 可用于其中每一对；归纳可得 $P$ 对
$W_{p, 1}, \ldots, W_{p, m} = U_p$ 成立。
\end{proof}

\begin{lemma}
\label{spaces-perfect-lemma-induction-principle-separated}
设 $S$ 为概形。设 $X$ 为 $S$ 上拟紧且拟分离的代数空间。设
$\mathcal{B} \subset \Ob(X_{spaces, \etale})$。设 $P$ 是
$\mathcal{B}$ 中元素的一项性质。假设
\begin{enumerate}
\item 每个 $W \in \mathcal{B}$ 都拟紧且拟分离；
\item 若 $W \in \mathcal{B}$，且 $U \subset W$ 是拟紧开集，则
$U \in \mathcal{B}$；
\item 若 $V \in \Ob(X_{spaces, \etale})$ 仿射，则
(a) $V \in \mathcal{B}$，并且 (b) $P$ 对 $V$ 成立；
\item 对每个初等特异方形 $(U \subset W, f : V \to W)$，若
\begin{enumerate}
\item $W \in \mathcal{B}$；
\item $U$ 拟紧；
\item $V$ 仿射；并且
\item $P$ 对 $U$、$V$ 和 $U \times_W V$ 成立，
\end{enumerate}
则 $P$ 对 $W$ 成立。
\end{enumerate}
于是，$P$ 对每个 $W \in \mathcal{B}$ 成立。
\end{lemma}

\begin{proof}
证明与引理 \ref{spaces-perfect-lemma-induction-principle} 的证明完全相同。
（补充说明，(4)(d) 是有意义的，因为 $U \times_W V$ 是 $V$ 的拟紧开集，
所以由条件 (2) 和 (3)，它是 $\mathcal{B}$ 的元素。）
\end{proof}

\begin{remark}
\label{spaces-perfect-remark-how-to}
在引理 \ref{spaces-perfect-lemma-induction-principle-separated} 中，应如何选择集合
$\mathcal{B}$？以下是一些例子：
\begin{enumerate}
\item 若 $X$ 拟紧且分离，则可取 $\mathcal{B}$ 为
$X_{spaces, \etale}$ 的拟紧分离对象所成的集合。这时
$X \in \mathcal{B}$，且 $\mathcal{B}$ 满足 (1)、(2) 和 (3)(a)。
在这样选择 $\mathcal{B}$ 时，引理
\ref{spaces-perfect-lemma-induction-principle-separated} 重现引理
\ref{spaces-perfect-lemma-induction-principle}。
\item 若 $X$ 拟紧，且其在 $\mathbf{Z}$ 上的对角态射仿射（见“代数空间的性质”定义
\ref{spaces-properties-definition-separated}），则可取 $\mathcal{B}$ 为
$X_{spaces, \etale}$ 中拟紧且在 $\mathbf{Z}$ 上具有仿射对角态射的对象所成的集合。
同样，$X \in \mathcal{B}$，且 $\mathcal{B}$ 满足 (1)、(2) 和 (3)(a)。
\item 若 $X$ 拟紧且拟分离，则满足 (1)、(2) 和 (3)(a) 的最小子集
$\mathcal{B}$（它包含 $X$）由下述规则给出：$W \in \mathcal{B}$ 当且仅当，或者
$W$ 是 $X$ 的拟紧开子空间，或者 $W$ 是 $X_{spaces, \etale}$
某个仿射对象的拟紧开集。
\end{enumerate}
\end{remark}

\noindent
下面给出一个变式，将某个开子空间上的成立性扩张到更大的开子空间。

\begin{lemma}
\label{spaces-perfect-lemma-induction-principle-enlarge}
设 $S$ 为概形。设 $X$ 为 $S$ 上拟紧且拟分离的代数空间。
设 $W \subset X$ 为拟紧开子空间。设 $P$ 是 $X$ 的拟紧开子空间的一项性质。
假设
\begin{enumerate}
\item $P$ 对 $W$ 成立；并且
\item 对每个初等特异方形
$(W_1 \subset W_2, f : V \to W_2)$，若满足下列条件： % P09-ERR-0406
\begin{enumerate}
\item $W_1$、$W_2$ 是 $X$ 的拟紧开子空间；
\item $W \subset W_1$；
\item $V$ 仿射；并且
\item $P$ 对 $W_1$ 成立，
\end{enumerate}
则 $P$ 对 $W_2$ 成立。
\end{enumerate}
于是，$P$ 对 $X$ 成立。
\end{lemma}

\begin{proof}
这可由引理 \ref{spaces-perfect-lemma-induction-principle-separated} 推出；不过，我们将直接
重做引理 \ref{spaces-perfect-lemma-induction-principle} 的证明，给出一个直接论证。
我们将使用滤过
$$
\emptyset = U_{n + 1} \subset
U_n \subset U_{n - 1} \subset \ldots \subset U_1 = X
$$
以及“良态空间”引理
\ref{decent-spaces-lemma-filter-quasi-compact-quasi-separated} 中的态射
$f_p : V_p \to U_p$。我们将证明 $P$ 对
$W_p = W \cup U_p$ 成立，并对 $p$ 作降序归纳。由于 $W_1 = X$，这将完成证明。
注意，由 (1)，$P$ 对 $W_{n + 1} = W \cap U_{n + 1} = W$ 成立。 % P09-ERR-0407
假设 $P$ 对 $W_{p + 1}$ 成立。注意，
$W_p \setminus W_{p + 1}$（赋予约化诱导子空间结构）是
$U_p \setminus U_{p + 1}$ 的闭子空间。由于
$(U_{p + 1} \subset U_p, f_p : V_p \to U_p)$ 是初等特异方形，
$(W_{p + 1} \subset W_p, f_p : V_p \to W_p)$ 也是如此。
然而，(2) 可能不适用，因为 $V_p$ 未必仿射。
然而，由于 $V_p$ 是拟紧概形，可选取有限仿射开覆盖
$V_p = V_{p, 1} \cup \ldots \cup V_{p, m}$。令
$W_{p, 0} = W_{p + 1}$，并令
$$
W_{p, i} = W_{p + 1} \cup f_p(V_{p, 1} \cup \ldots \cup V_{p, i})
$$
其中 $i = 1, \ldots, m$。这些是 $X$ 的拟紧开子空间，并包含 $W$。于是
$$
W_{p + 1} = W_{p, 0} \subset
W_{p, 1} \subset \ldots \subset
W_{p, m} = W_p
$$
且各对
$$
(W_{p, 0} \subset W_{p, 1}, f_p|_{V_{p, 1}}),
(W_{p, 1} \subset W_{p, 2}, f_p|_{V_{p, 2}}),\ldots,
(W_{p, m - 1} \subset W_{p, m}, f_p|_{V_{p, m}})
$$
由引理 \ref{spaces-perfect-lemma-make-more-elementary-distinguished-squares}
均为初等特异方形。现在，(2) 可用于其中每一对；归纳可得 $P$ 对
$W_{p, 1}, \ldots, W_{p, m} = W_p$ 成立。
\end{proof}

\section{Mayer--Vietoris}
\label{spaces-perfect-section-mayer-vietoris}

\noindent
本节证明，一个初等特异三角会导出各种 Mayer--Vietoris 序列。 % P09-ERR-0408

\medskip\noindent
设 $S$ 为概形。设 $U \to X$ 为 $S$ 上代数空间的平展态射。在
“代数空间的性质”第 \ref{spaces-properties-section-localize} 节中已经证明，

$U_{spaces, \etale} = X_{spaces, \etale}/U$

并且与结构层相容。因此，在这种情形下，我们常把态射 $j_U : U \to X$
看作局部化态射（参见“位点上的模”定义
\ref{sites-modules-definition-localize-ringed-site}）。特别地，我们把拉回
$j_U^*$ 看作限制到 $U$，并常将其记为 ${}|_U$；这与
“代数空间的性质”等式
(\ref{spaces-properties-equation-restrict-modules}) 相容。特别地，我们看到，
\begin{equation}
\label{spaces-perfect-equation-stalk-restriction}
(\mathcal{F}|_U)_{\overline{u}} = \mathcal{F}_{\overline{x}}
\end{equation}
若 $\overline{u}$ 是 $U$ 的几何点，且 $\overline{x}$ 是
$\overline{u}$ 在 $X$ 中的像。此外，限制有一个正合左伴随 $j_{U!}$；参见“位点上的模”引理
\ref{sites-modules-lemma-extension-by-zero} 和
\ref{sites-modules-lemma-extension-by-zero-exact}。
最后，回顾若 $\mathcal{G}$ 是一个 $\mathcal{O}_X$-模，则有
\begin{equation}
\label{spaces-perfect-equation-stalk-j-shriek}
(j_{U!}\mathcal{G})_{\overline{x}} =
\bigoplus\nolimits_{\overline{u}} \mathcal{G}_{\overline{u}}
\end{equation}
对任意几何点 $\overline{x} : \Spec(k) \to X$，其中直和取遍态射
$\overline{u} : \Spec(k) \to U$，这些态射满足 $j_U \circ \overline{u} = \overline{x}$；参见“位点上的模”引理
\ref{sites-modules-lemma-stalk-j-shriek} 以及
“代数空间的性质”引理
\ref{spaces-properties-lemma-points-small-etale-site}。

\begin{lemma}
\label{spaces-perfect-lemma-exact-sequence-lower-shriek}
设 $S$ 为概形。设 $(U \subset X, V \to X)$ 为 $S$ 上代数空间的初等特异方形。
\begin{enumerate}
\item 对于 $\mathcal{O}_X$-模层 $\mathcal{F}$，有短正合列
$$
0 \to j_{U \times_X V!}\mathcal{F}|_{U \times_X V} \to
j_{U!}\mathcal{F}|_U \oplus j_{V!}\mathcal{F}|_V \to \mathcal{F} \to 0
$$
\item 对于对象 $E$（属于 $D(\mathcal{O}_X)$），有特异三角
$$
j_{U \times_X V!}E|_{U \times_X V} \to
j_{U!}E|_U \oplus j_{V!}E|_V \to E \to
j_{U \times_X V!}E|_{U \times_X V}[1]
$$
在 $D(\mathcal{O}_X)$ 中。
\end{enumerate}
\end{lemma}

\begin{proof}
为证明 (1) 中的列正合，根据“代数空间的性质”定理
\ref{spaces-properties-theorem-exactness-stalks}，可在几何点处取茎来检验。
设 $\overline{x}$ 为 $X$ 的几何点。由等式
(\ref{spaces-perfect-equation-stalk-restriction}) 和 (\ref{spaces-perfect-equation-stalk-j-shriek})，
在 $\overline{x}$ 处取茎得到序列
$$
0 \to
\bigoplus\nolimits_{(\overline{u}, \overline{v})} \mathcal{F}_{\overline{x}}
\to
\bigoplus\nolimits_{\overline{u}} \mathcal{F}_{\overline{x}}
\oplus
\bigoplus\nolimits_{\overline{v}} \mathcal{F}_{\overline{x}}
\to
\mathcal{F}_{\overline{x}} \to 0
$$
此序列正合，因为对每个 $\overline{x}$，要么恰有一个
$\overline{u}$ 映到 $\overline{x}$，要么没有 $\overline{u}$，且恰有一个
$\overline{v}$ 映到 $\overline{x}$。

\medskip\noindent
证明 (2)。我们在“位点上的上同调”第
\ref{sites-cohomology-section-properties-K-injective} 节中看到，
导出范畴上的限制函子和零延拓函子，只需将函子逐项作用于任意复形即可计算。
设 $\mathcal{E}^\bullet$ 为 $\mathcal{O}_X$-模复形，表示对象 $E$。
本引理的特异三角是下述 $\mathcal{O}_X$-模复形短正合列所关联的特异三角，
其中所用关联来自“导出范畴”第
\ref{derived-section-canonical-delta-functor} 节，尤其是引理
\ref{derived-lemma-derived-canonical-delta-functor}：
$$
0 \to j_{U \times_X V!}\mathcal{E}^\bullet|_{U \times_X V} \to
j_{U!}\mathcal{E}^\bullet|_U \oplus j_{V!}\mathcal{E}^\bullet|_V
\to \mathcal{E}^\bullet \to 0
$$
由 (1)，此列短正合。
\end{proof}

\begin{lemma}
\label{spaces-perfect-lemma-exact-sequence-j-star}
设 $S$ 为概形。设 $(U \subset X, V \to X)$ 为 $S$ 上代数空间的初等特异方形。
\begin{enumerate}
\item 对每个 $\mathcal{O}_X$-模层 $\mathcal{F}$，有短正合列
$$
0 \to \mathcal{F} \to
j_{U, *}\mathcal{F}|_U \oplus j_{V, *}\mathcal{F}|_V \to
j_{U \times_X V, *}\mathcal{F}|_{U \times_X V} \to 0
$$
\item 对于对象 $E$（属于 $D(\mathcal{O}_X)$），有特异三角
$$
E \to
Rj_{U, *}E|_U \oplus Rj_{V, *}E|_V \to
Rj_{U \times_X V, *}E|_{U \times_X V} \to
E[1]
$$
在 $D(\mathcal{O}_X)$ 中。
\end{enumerate}
\end{lemma}

\begin{proof}
设 $W$ 为 $X_\etale$ 中的对象。我们断言，序列
$$
0 \to
\mathcal{F}(W) \to
\mathcal{F}(W \times_X U) \oplus \mathcal{F}(W \times_X V) \to
\mathcal{F}(W \times_X U \times_X V)
$$
正合，并且末项中的元素可在 $W$ 上局部地提升到中间项。
由引理 \ref{spaces-perfect-lemma-make-more-elementary-distinguished-squares}，
$(W \times_X U \subset W, V \times_X W \to W)$ 是初等特异方形。
因此可设 $W = X$，只需证明
$$
0 \to
\mathcal{F}(X) \to
\mathcal{F}(U) \oplus \mathcal{F}(V) \to
\mathcal{F}(U \times_X V)
$$
我们已经在引理 \ref{spaces-perfect-lemma-exact-sequence-lower-shriek} 中看到，
$$
0 \to j_{U \times_X V!}\mathcal{O}_{U \times_X V}
\to j_{U!}\mathcal{O}_U \oplus
j_{V!}\mathcal{O}_V \to
\mathcal{O}_X \to 0
$$
是 $\mathcal{O}_X$-模的正合列，而对其施加右正合函子
$\Hom_{\mathcal{O}_X}(- , \mathcal{F})$，得到上面的序列。 % P09-ERR-0409
这也意味着，将 $s \in \mathcal{F}(U \times_X V)$ 提升为
$\mathcal{F}(U) \oplus \mathcal{F}(V)$ 中元素的障碍属于
$\Ext^1_{\mathcal{O}_X}(\mathcal{O}_X, \mathcal{F}) = H^1(X, \mathcal{F})$。
由上同调的局部性（“位点上的上同调”引理
\ref{sites-cohomology-lemma-kill-cohomology-class-on-covering}），
此障碍在 $X$ 上平展局部地消失，故 (1) 的证明完成。

\medskip\noindent
证明 (2)。选取 K-单射复形 $\mathcal{I}^\bullet$ 表示对象 $E$，其各项
$\mathcal{I}^n$ 是 $\textit{Mod}(\mathcal{O}_X)$ 中的单射对象；参见“单射”定理
\ref{injectives-theorem-K-injective-embedding-grothendieck}。
于是 $\mathcal{I}^\bullet|U$ 是 K-单射复形（“位点上的上同调”引理
\ref{sites-cohomology-lemma-restrict-K-injective-to-open}）。因此，
$Rj_{U, *}E|_U$ 由 $j_{U, *}\mathcal{I}^\bullet|_U$ 表示。
对 $V$ 及 $U \times_X V$ 也同理。因此，本引理的特异三角是下述复形短正合列
所关联的特异三角，其中所用关联来自“导出范畴”第
\ref{derived-section-canonical-delta-functor} 节，尤其是引理
\ref{derived-lemma-derived-canonical-delta-functor}：
$$
0 \to
\mathcal{I}^\bullet \to
j_{U, *}\mathcal{I}^\bullet|_U \oplus j_{V, *}\mathcal{I}^\bullet|_V \to
j_{U \times_X V, *}\mathcal{I}^\bullet|_{U \times_X V} \to
0.
$$
由 (1)，此列正合。
\end{proof}
\begin{lemma}
\label{spaces-perfect-lemma-unbounded-relative-mayer-vietoris}
设 $S$ 为概形。设 $f : X \to Y$ 为
$S$ 上代数空间之间的态射。设 $(U \subset X, V \to X)$ 为初等特异方形。
记
$a = f|_U : U \to Y$、$b = f|_V : V \to Y$，以及
$c = f|_{U \times_X V} : U \times_X V \to Y$ 为相应限制。
对于对象 $E$（属于 $D(\mathcal{O}_X)$），存在特异三角
$$
Rf_*E \to
Ra_*(E|_U) \oplus Rb_*(E|_V) \to
Rc_*(E|_{U \times_X V}) \to
Rf_*E[1]
$$
在 $D(\mathcal{O}_Y)$ 中。该三角关于 $E$ 具有函子性。
\end{lemma}

\begin{proof}
选取 K-单射复形 $\mathcal{I}^\bullet$ 表示对象 $E$。不妨设
$\mathcal{I}^n$ 都是 $\textit{Mod}(\mathcal{O}_X)$ 中的单射对象，对所有 $n$ 均成立；参见“单射”定理
\ref{injectives-theorem-K-injective-embedding-grothendieck}。
于是 $Rf_*E$ 由 $f_*\mathcal{I}^\bullet$ 计算。
对 $U$、$V$ 以及 $U \cap V$ 也同理 % P09-ERR-0410
，参见“位点上的上同调”引理
\ref{sites-cohomology-lemma-restrict-K-injective-to-open}。
因此，本引理的特异三角是下述短正合复形列所关联的特异三角，其中所用关联来自
“导出范畴”第 \ref{derived-section-canonical-delta-functor} 节以及尤其是引理
\ref{derived-lemma-derived-canonical-delta-functor}：
$$
0 \to
f_*\mathcal{I}^\bullet \to
a_*\mathcal{I}^\bullet|_U \oplus b_*\mathcal{I}^\bullet|_V \to
c_*\mathcal{I}^\bullet|_{U \times_X V} \to
0.
$$
为说明这是复形的短正合列，论证如下。取单射对象 $\mathcal{I}$（属于
$\textit{Mod}(\mathcal{O}_X)$）。对下列短正合列施加 $f_*$：
$$
0 \to \mathcal{I} \to
j_{U, *}\mathcal{I}|_U \oplus j_{V, *}\mathcal{I}|_V \to
j_{U \times_X V, *}\mathcal{I}|_{U \times_X V} \to 0
$$
该列来自引理 \ref{spaces-perfect-lemma-exact-sequence-j-star}；再利用
$R^1f_*\mathcal{I} = 0$，得到短正合列
$$
0 \to f_*\mathcal{I} \to
f_*j_{U, *}\mathcal{I}|_U \oplus f_*j_{V, *}\mathcal{I}|_V \to
f_*j_{U \times_X V, *}\mathcal{I}|_{U \times_X V} \to 0
$$
最后注意到 $a_* = f_*j_{U, *}$，对 $b_*$ 和 $c_*$ 也同样成立，即得证明。
\end{proof}

\begin{lemma}
\label{spaces-perfect-lemma-mayer-vietoris-hom}
设 $S$ 为概形。设 $(U \subset X, V \to X)$ 为 $S$ 上代数空间的初等特异方形。
对于对象 $E$、$F$（属于 $D(\mathcal{O}_X)$），有 Mayer--Vietoris 序列
$$
\xymatrix{
& \ldots \ar[r] &
\Ext^{-1}(E_{U \times_X V}, F_{U \times_X V}) \ar[lld] \\
\Hom(E, F) \ar[r] &
\Hom(E_U, F_U) \oplus
\Hom(E_V, F_V) \ar[r] &
\Hom(E_{U \times_X V}, F_{U \times_X V})
}
$$
其中下标表示限制到相应开子空间，并且各个 $\Hom$ 都是在相应导出范畴中取的。
\end{lemma}

\begin{proof}
使用引理 \ref{spaces-perfect-lemma-exact-sequence-lower-shriek} 的特异三角，
得到 $\Hom$ 的长正合序列（由“导出范畴”引理
\ref{derived-lemma-representable-homological} 得到），并利用
\[
\Hom(j_{U!}E|_U, F) = \Hom(E|_U, F|_U)
\]
以及“位点上的上同调”引理
\ref{sites-cohomology-lemma-adjoint-lower-shriek-restrict}。
\end{proof}

\begin{lemma}
\label{spaces-perfect-lemma-unbounded-mayer-vietoris}
设 $S$ 为概形。设 $(U \subset X, V \to X)$ 为 $S$ 上代数空间的初等特异方形。
对于对象 $E$（属于 $D(\mathcal{O}_X)$），有特异三角
$$
R\Gamma(X, E) \to R\Gamma(U, E) \oplus R\Gamma(V, E) \to
R\Gamma(U \times_X V, E) \to R\Gamma(X, E)[1]
$$
以及特别的长正合上同调序列
$$
\ldots \to
H^n(X, E) \to
H^n(U, E) \oplus H^n(V, E) \to
H^n(U \times_X V, E) \to
H^{n + 1}(X, E) \to \ldots
$$
特异三角和长正合序列的构造关于 $E$ 具有函子性。
\end{lemma}

\begin{proof}
选取 K-单射复形 $\mathcal{I}^\bullet$ 表示对象 $E$，其各项
$\mathcal{I}^n$ 都是 $\textit{Mod}(\mathcal{O}_X)$ 中的单射对象；参见“单射”定理
\ref{injectives-theorem-K-injective-embedding-grothendieck}。
在引理 \ref{spaces-perfect-lemma-exact-sequence-j-star} 的证明中，我们得到复形的短正合列
$$
0 \to \mathcal{I}^\bullet \to
j_{U, *}\mathcal{I}^\bullet|_U \oplus j_{V, *}\mathcal{I}^\bullet|_V \to
j_{U \times_X V, *}\mathcal{I}^\bullet|_{U \times_X V} \to 0
$$
由于 $H^1(X, \mathcal{I}^n) = 0$，可知取整体截面得到复形的正合列
$$
0 \to \Gamma(X, \mathcal{I}^\bullet) \to
\Gamma(U, \mathcal{I}^\bullet) \oplus
\Gamma(V, \mathcal{I}^\bullet) \to
\Gamma(U \times_X V, \mathcal{I}^\bullet) \to 0
$$
这些复形分别表示
$R\Gamma(X, E)$、$R\Gamma(U, E)$、$R\Gamma(V, E)$ 以及
$R\Gamma(U \times_X V, E)$，故由“导出范畴”第
\ref{derived-section-canonical-delta-functor} 节以及尤其是引理
\ref{derived-lemma-derived-canonical-delta-functor}，
得到特异三角。
\end{proof}

\begin{lemma}
\label{spaces-perfect-lemma-restrict-lower-shriek}
设 $S$ 为概形。设 $j : U \to X$ 为 $S$ 上代数空间之间的平展态射。
给定平展态射 $V \to Y$，令 $W = V \times_X U$，并记投影态射 % P09-ERR-0411
$j_W : W \to V$。则
$(j_!E)|_V = j_{W!}(E|_W)$，其中 $E$ 属于 $D(\mathcal{O}_U)$。
\end{lemma}

\begin{proof}
这是因为
$(j_!\mathcal{F})|_V = j_{W!}(\mathcal{F}|_W)$
对于 $\mathcal{O}_X$-模 $\mathcal{F}$ 成立；这直接来自函子 $j_!$ 和 $j_{W!}$ 的构造；参见“位点上的模”引理
\ref{sites-modules-lemma-extension-by-zero}。
\end{proof}

\begin{lemma}
\label{spaces-perfect-lemma-pushforward-with-support-in-open}
设 $S$ 为概形。设 $(U \subset X, j : V \to X)$ 为 $S$ 上代数空间的初等特异方形。
置 $T = |X| \setminus |U|$。
\begin{enumerate}
\item 若 $E$ 是 $D(\mathcal{O}_X)$ 中支撑在 $T$ 上的对象，则
(a) $E \to Rj_*(E|_V)$ 以及 (b) $j_!(E|_V) \to E$ 都是同构。
\item 若 $F$ 是 $D(\mathcal{O}_V)$ 中支撑在 $j^{-1}T$ 上的对象，则
(a) $F \to (j_!F)|_V$、(b) $(Rj_*F)|_V \to F$，以及 (c)
$j_!F \to Rj_*F$ 都是同构。
\end{enumerate}
\end{lemma}

\begin{proof}
设 $E$ 为 $D(\mathcal{O}_X)$ 中的对象，且其上同调层支撑在 $T$ 上。
于是 $E|_U = 0$ 且 $E|_{U \times_X V} = 0$，因为 $T$ 不与 $U$ 相交，
而 $j^{-1}T$ 不与 $U \times_X V$ 相交。因此 (1)(a) 由引理
\ref{spaces-perfect-lemma-exact-sequence-j-star} 得出。完全同样地，(1)(b) 由引理
\ref{spaces-perfect-lemma-exact-sequence-lower-shriek} 得出。

\medskip\noindent
设 $F$ 为 $D(\mathcal{O}_V)$ 中的对象，且其上同调层支撑在 $j^{-1}T$ 上。
由引理 \ref{spaces-perfect-lemma-restrict-direct-image-open}，有
$(Rj_*F)|_U = Rj_{W, *}(F|_W) = 0$，因为按假设 $F|_W = 0$。
同理，由引理 \ref{spaces-perfect-lemma-restrict-lower-shriek}，
$(j_!F)|_U = j_{W!}(F|_W) = 0$。
因此 $j_!F$ 和 $Rj_*F$ 支撑在 $T$ 上，而
$(j_!F)|_V$ 和 $(Rj_*F)|_V$ 支撑在 $j^{-1}(T)$ 上。根据“代数空间的性质”定理
\ref{spaces-properties-theorem-exactness-stalks}，要在导出范畴中检验映射
(2)(a)、(b)、(c) 为同构，只需检验它们在 $T$ 和 $j^{-1}(T)$ 的几何点处
诱导的上同调层茎上的映射为同构。 % P09-ERR-0412
为此，可作如下替换：将 $X$ 换为 $V$，将 $U$ 换为
$U \times_X V$，将 $V$ 换为 $V \times_X V$，并将 $F$ 换为
$F|_{V \times_X V}$（通过第一投影限制）；参见引理
\ref{spaces-perfect-lemma-restrict-direct-image-open}、
\ref{spaces-perfect-lemma-restrict-lower-shriek} 以及
\ref{spaces-perfect-lemma-make-more-elementary-distinguished-squares}。
由于 $V \times_X V \to V$ 有截面，这将 (2) 归结为
$j : V \to X$ 有截面的情形。

\medskip\noindent
设 $j$ 有截面 $\sigma : X \to V$。置 $V' = \sigma(X)$，这是 $V$ 的开子空间。
置 $U' = j^{-1}(U)$，这是 $V$ 的另一个开子空间。
于是 $(U' \subset V, V' \to V)$ 是初等特异方形。
注意，$F|_{U'} = 0$ 且 $F|_{V' \cap U'} = 0$，因为 $F$ 支撑在
$j^{-1}(T)$ 上。记开浸入 $j' : V' \to V$，并记
$j_{V'} : V' \to X$ 为复合态射 $V' \to V \to X$，它是 $\sigma$ 的逆。
置 $F' = \sigma^*F$。引理
\ref{spaces-perfect-lemma-exact-sequence-lower-shriek} 和
\ref{spaces-perfect-lemma-exact-sequence-j-star} 的特异三角表明
$F = j'_!(F|_{V'})$ 且 $F = Rj'_*(F|_{V'})$。
于是
$j_!F = j_!j'_!(F|_{V'}) = j_{V'!}F = F'$，
因为 $j_{V'} : V' \to X$ 是同构且是 $\sigma$ 的逆。
同理，$Rj_*F = Rj_*Rj'_*F = Rj_{V', *}F = F'$。
这证明了 (2)(c)。为证明 (2)(a) 和 (2)(b)，只需证明
$F = F'|_V$。这是显然的，因为 $F$ 和 $F'|_V$ 在 $U'$ 及
$U' \cap V'$ 上都限制为零，并且在 $V'$ 上是同一个对象。
\end{proof}

\noindent
我们可以粘合复形！

\begin{lemma}
\label{spaces-perfect-lemma-glue}
设 $S$ 为概形。设 $(U \subset X, V \to X)$ 为 $S$ 上代数空间的初等特异方形。
给定
\begin{enumerate}
\item 对象 $A$，属于 $D(\mathcal{O}_U)$，
\item 对象 $B$，属于 $D(\mathcal{O}_V)$，以及
\item 同构 $c : A|_{U \times_X V} \to B|_{U \times_X V}$。
\end{enumerate}
则存在对象 $F$，属于 $D(\mathcal{O}_X)$，以及同构
$f : F|_U \to A$、$g : F|_V \to B$，使得
$c = g|_{U \times_X V} \circ f^{-1}|_{U \times_X V}$。
此外，给定
\begin{enumerate}
\item 对象 $E$，属于 $D(\mathcal{O}_X)$，
\item 态射 $a : A \to E|_U$，属于 $D(\mathcal{O}_U)$，
\item 态射 $b : B \to E|_V$，属于 $D(\mathcal{O}_V)$，
\end{enumerate}
满足
$$
a|_{U \times_X V}  = b|_{U \times_X V} \circ c.
$$
则存在态射 $F \to E$，属于 $D(\mathcal{O}_X)$，
其限制到 $U$ 为 $a \circ f$，限制到 $V$ 为 $b \circ g$。
\end{lemma}

\begin{proof}
记 $j_U$、$j_V$、$j_{U \times_X V}$ 为指向 $X$ 的相应态射。选取特异三角
$$
F \to Rj_{U, *}A \oplus Rj_{V, *}B \to
Rj_{U \times_X V, *}(B|_{U \times_X V}) \to F[1]
$$
其中映射 $Rj_{V, *}B \to Rj_{U \times_X V, *}(B|_{U \times_X V})$
是显然的。映射
$Rj_{U, *}A \to Rj_{U \times_X V, *}(B|_{U \times_X V})$
是复合映射：先取
$Rj_{U, *}A \to Rj_{U \times_X V, *}(A|_{U \times_X V})$，
再与 $Rj_{U \times_X V, *}c$ 复合。将其限制到 $U$，得到
$$
F|_U \to A \oplus (Rj_{V, *}B)|_U \to
(Rj_{U \times_X V, *}(B|_{U \times_X V}))|_U \to F|_U[1]
$$
记 $j : U \times_X V \to U$。根据引理
\ref{spaces-perfect-lemma-restrict-direct-image-open}，限制与总正像相容，故
$(Rj_{V, *}B)|_U$ 以及
$(Rj_{U \times_X V, *}(B|_{U \times_X V}))|_U$
都典范同构于 $Rj_*(B|_{U \times_X V})$。
因此上面最后一个显示式中的第二个箭头有截面，从而
$F|_U \to A$ 是同构。

\medskip\noindent
为证明 $F|_V \to B$ 是同构，我们使用一个技巧。也就是说，选取特异三角
$$
F|_V \to B \to B' \to F[1]|_V
$$
在 $D(\mathcal{O}_V)$ 中。由于 $F|_U \to A$ 是同构，且有同构
$c : A|_{U \times_X V} \to B|_{U \times_X V}$，
$F|_V \to B$ 在 $U \times_X V$ 上的限制是同构。
因此 $B'$ 支撑在 $j_V^{-1}(T)$ 上，其中 $T = |X| \setminus |U|$。
另一方面，有特异三角的态射
$$
\xymatrix{
F \ar[r] \ar[d] &
Rj_{U, *}F|_U \oplus Rj_{V, *}F|_V \ar[r] \ar[d] &
Rj_{U \times_X V, *}F|_{U \times_X V} \ar[r] \ar[d] &
F[1] \ar[d] \\
F \ar[r] &
Rj_{U, *}A \oplus Rj_{V, *}B \ar[r] &
Rj_{U \times_X V, *}(B|_{U \times_X V}) \ar[r] &
F[1]
}
$$
该图中所有竖直映射都是同构，除了
$Rj_{V, *}F|_V \to Rj_{V, *}B$；因此由“导出范畴”引理
\ref{derived-lemma-third-isomorphism-triangle}，该映射也为同构。 % P09-ERR-0413
这意味着 $Rj_{V, *}B' = 0$。于是由引理
\ref{spaces-perfect-lemma-pushforward-with-support-in-open} 得到 $B' = 0$。

\medskip\noindent
态射 $F \to E$ 的存在性由 Mayer--Vietoris 的 $\Hom$ 序列得出，参见引理
\ref{spaces-perfect-lemma-mayer-vietoris-hom}。
\end{proof}

\section{凝聚化函子}
\label{spaces-perfect-section-coherator}

\noindent
设 $S$ 为概形。设 $X$ 为 $S$ 上的代数空间。
{\it 凝聚化函子}是函子
$$
Q_X :
\textit{Mod}(\mathcal{O}_X)
\longrightarrow
\QCoh(\mathcal{O}_X)
$$
它是包含函子
$\QCoh(\mathcal{O}_X) \to \textit{Mod}(\mathcal{O}_X)$ 的右伴随。
它对任意代数空间 $X$ 都存在；此外，伴随映射
$Q_X(\mathcal{F}) \to \mathcal{F}$ 对每个拟凝聚模 $\mathcal{F}$ 都是同构，参见“代数空间的性质”命题
\ref{spaces-properties-proposition-coherator}。
由于 $Q_X$ 是左正合的（作为右伴随），我们可以考虑其右导出延拓
$$
RQ_X :
D(\mathcal{O}_X)
\longrightarrow
D(\QCoh(\mathcal{O}_X)).
$$
由于 $Q_X$ 是包含函子
$\QCoh(\mathcal{O}_X) \to \textit{Mod}(\mathcal{O}_X)$ 的右伴随，
由“导出范畴”引理
\ref{derived-lemma-derived-adjoint-functors} 可知，$RQ_X$ 是典范函子
$D(\QCoh(\mathcal{O}_X)) \to D(\mathcal{O}_X)$ 的右伴随。

\medskip\noindent
本节研究函子 $RQ_X$。在
\ref{spaces-perfect-section-better-coherator} 节中，我们将研究包含函子
$D_\QCoh(\mathcal{O}_X) \to D(\mathcal{O}_X)$ 的（密切相关的）右伴随（若其存在）。

\begin{lemma}
\label{spaces-perfect-lemma-affine-pushforward}
设 $S$ 为概形。设 $f : X \to Y$ 为 $S$ 上代数空间的仿射态射。
则 $f_*$ 定义导出函子
$f_* : D(\QCoh(\mathcal{O}_X)) \to D(\QCoh(\mathcal{O}_Y))$。
该函子具有如下性质：
$$
\xymatrix{
D(\QCoh(\mathcal{O}_X)) \ar[d]_{f_*} \ar[r] &
D_\QCoh(\mathcal{O}_X) \ar[d]^{Rf_*} \\
D(\QCoh(\mathcal{O}_Y)) \ar[r] &
D_\QCoh(\mathcal{O}_Y)
}
$$
该图交换。
\end{lemma}

\begin{proof}
函子
$f_* : \QCoh(\mathcal{O}_X) \to \QCoh(\mathcal{O}_Y)$
是正合的，参见“空间的上同调”引理
\ref{spaces-cohomology-lemma-affine-vanishing-higher-direct-images}。
因此 $f_*$ 定义导出函子
$f_* : D(\QCoh(\mathcal{O}_X)) \to D(\QCoh(\mathcal{O}_Y))$，
只需将 $f_*$ 逐项作用于任意代表复形即可；参见“导出范畴”引理
\ref{derived-lemma-right-derived-exact-functor}。
对于任意 $\mathcal{O}_X$-模复形 $\mathcal{F}^\bullet$，有典范映射
$f_*\mathcal{F}^\bullet \to Rf_*\mathcal{F}^\bullet$。
为完成证明，我们证明当 $\mathcal{F}^\bullet$ 是每一项
$\mathcal{F}^n$ 都拟凝聚的复形时，该映射是拟同构。
该断言在 $Y$ 上平展局部成立，因此不妨设 $Y$ 仿射。
由于仿射态射是可表的，利用评注
\ref{spaces-perfect-remark-match-total-direct-images} 的相容性，可将问题归结为概形情形。
概形情形是“概形的导出范畴”引理
\ref{perfect-lemma-affine-pushforward}。
\end{proof}

\begin{lemma}
\label{spaces-perfect-lemma-flat-pushforward-coherator}
设 $S$ 为概形。设 $f : X \to Y$ 为 $S$ 上代数空间的态射。
假设 $f$ 拟紧、拟分离且平坦。记
$$
\Phi : D(\QCoh(\mathcal{O}_X)) \to D(\QCoh(\mathcal{O}_Y))
$$
为函子
$f_* : \QCoh(\mathcal{O}_X) \to \QCoh(\mathcal{O}_Y)$ 的右导出函子，
则有 $RQ_Y \circ Rf_* = \Phi \circ RQ_X$。
\end{lemma}

\begin{proof}
我们将证明 $RQ_Y \circ Rf_*$ 和 $\Phi \circ RQ_X$
是同一个函子的右伴随，该函子为
$D(\QCoh(\mathcal{O}_Y)) \to D(\mathcal{O}_X)$。

\medskip\noindent
由于 $f$ 拟紧且拟分离，$f_*$ 保持拟凝聚性，参见“空间的态射”引理
\ref{spaces-morphisms-lemma-pushforward}。
回顾 $\QCoh(\mathcal{O}_X)$ 是 Grothendieck 阿贝尔范畴（“代数空间的性质”命题
\ref{spaces-properties-proposition-coherator}）。
因此，任意对象 $K$ 属于 $D(\QCoh(\mathcal{O}_X))$，都可由
K-单射复形 $\mathcal{I}^\bullet$ 表示，该复形属于 $\QCoh(\mathcal{O}_X)$，参见“单射”定理
\ref{injectives-theorem-K-injective-embedding-grothendieck}。
于是可定义 $\Phi(K) = f_*\mathcal{I}^\bullet$。

\medskip\noindent
由于 $f$ 平坦，函子 $f^*$ 正合。因此 $f^*$ 定义
$f^* : D(\mathcal{O}_Y) \to D(\mathcal{O}_X)$ 以及
$f^* : D(\QCoh(\mathcal{O}_Y)) \to D(\QCoh(\mathcal{O}_X))$。
函子 $f^* = Lf^* : D(\mathcal{O}_Y) \to D(\mathcal{O}_X)$
是 $Rf_* : D(\mathcal{O}_X) \to D(\mathcal{O}_Y)$ 的左伴随，
参见“位点上的上同调”引理 \ref{sites-cohomology-lemma-adjoint}。
同理，函子
$f^* : D(\QCoh(\mathcal{O}_Y)) \to D(\QCoh(\mathcal{O}_X))$
是 $\Phi : D(\QCoh(\mathcal{O}_X)) \to D(\QCoh(\mathcal{O}_Y))$ 的左伴随，
参见“导出范畴”引理 \ref{derived-lemma-derived-adjoint-functors}。

\medskip\noindent
设 $A$ 为 $D(\QCoh(\mathcal{O}_Y))$ 中的对象，设 $E$ 为
$D(\mathcal{O}_X)$ 中的对象。则
\begin{align*}
\Hom_{D(\QCoh(\mathcal{O}_Y))}(A, RQ_Y(Rf_*E))
& =
\Hom_{D(\mathcal{O}_Y)}(A, Rf_*E) \\
& =
\Hom_{D(\mathcal{O}_X)}(f^*A, E) \\
& =
\Hom_{D(\QCoh(\mathcal{O}_X))}(f^*A, RQ_X(E)) \\
& =
\Hom_{D(\QCoh(\mathcal{O}_Y))}(A, \Phi(RQ_X(E)))
\end{align*}
这就证明了所需结论。
\end{proof}

\begin{lemma}
\label{spaces-perfect-lemma-affine-coherator}
设 $S$ 为概形。设 $X$ 为 $S$ 上的仿射代数空间。
置 $A = \Gamma(X, \mathcal{O}_X)$。则
\begin{enumerate}
\item $Q_X : \textit{Mod}(\mathcal{O}_X) \to \QCoh(\mathcal{O}_X)$
是如下函子：它把 $\mathcal{F}$ 送到拟凝聚 $\mathcal{O}_X$-模，
该模与 $A$-模 $\Gamma(X, \mathcal{F})$ 相关联；
\item $RQ_X : D(\mathcal{O}_X) \to D(\QCoh(\mathcal{O}_X))$
是如下函子：它把 $E$ 送到拟凝聚 $\mathcal{O}_X$-模复形，
该复形与对象 $R\Gamma(X, E)$（属于 $D(A)$）相关联；
\item 在 $D_\QCoh(\mathcal{O}_X)$ 上限制时，函子 $RQ_X$ 定义
(\ref{spaces-perfect-equation-compare}) 的拟逆。
\end{enumerate}
\end{lemma}

\begin{proof}
令 $X_0 = \Spec(A)$ 为表示 $X$ 的仿射概形。
回顾存在带环位点态射
$\epsilon : X_\etale \to X_{0, Zar}$，它诱导等价
$$
\xymatrix{
\QCoh(\mathcal{O}_X) \ar@<1ex>[r]^{{\epsilon_*}} &
\QCoh(\mathcal{O}_{X_0}) \ar@<1ex>[l]^{{\epsilon^*}}
}
$$
参见引理
\ref{spaces-perfect-lemma-derived-quasi-coherent-small-etale-site}。
因此，由伴随函子的唯一性，$Q_X = \epsilon^* \circ Q_{X_0} \circ \epsilon_*$。
于是，根据“概形的导出范畴”引理
\ref{perfect-lemma-affine-coherator} 对 $Q_{X_0}$ 的描述以及
$\Gamma(X_0, \epsilon_*\mathcal{F}) = \Gamma(X, \mathcal{F})$，得到 (1)。
断言 (2) 由 (1) 以及从 $A$-模到拟凝聚
$\mathcal{O}_X$-模的函子是正合的这一事实推出。
第三个断言现在由概形情形的结果
（“概形的导出范畴”引理 \ref{perfect-lemma-affine-coherator}）
以及引理
\ref{spaces-perfect-lemma-derived-quasi-coherent-small-etale-site} 得出。
\end{proof}

\noindent
接下来，我们证明函子
$D(\QCoh(\mathcal{O}_X)) \to D_\QCoh(\mathcal{O}_X)$
何时为等价。

\begin{lemma}
\label{spaces-perfect-lemma-argument-proves}
设 $S$ 为概形。设 $X$ 为 $S$ 上拟紧且拟分离的代数空间。
假设对于每个平展态射
$j : V \to W$，其中 $W \subset X$ 为拟紧开子空间且 $V$ 为仿射，
右导出函子
$$
\Phi : D(\QCoh(\mathcal{O}_U)) \to D(\QCoh(\mathcal{O}_W))
$$ % P09-ERR-0414
（左正合函子
$j_* : \QCoh(\mathcal{O}_V) \to \QCoh(\mathcal{O}_W)$ 的右导出函子）
都包含在如下交换图中：
$$
\xymatrix{
D(\QCoh(\mathcal{O}_V)) \ar[d]_\Phi \ar[r]_{i_V} &
D_\QCoh(\mathcal{O}_V) \ar[d]^{Rj_*} \\
D(\QCoh(\mathcal{O}_W)) \ar[r]^{i_W} &
D_\QCoh(\mathcal{O}_W)
}
$$
则函子
(\ref{spaces-perfect-equation-compare})
$$
D(\QCoh(\mathcal{O}_X))
\longrightarrow
D_\QCoh(\mathcal{O}_X)
$$
是等价，并且拟逆由 $RQ_X$ 给出。
\end{lemma}

\begin{proof}
首先使用归纳原理证明 $i_X$ 是全忠实的。
更确切地说，我们将使用引理
\ref{spaces-perfect-lemma-induction-principle-enlarge}。
设 $(U \subset W, V \to W)$ 为初等特异方形，且 $V$ 为仿射、
$U, W$ 为 $X$ 中的拟紧开子空间。假设 $i_U$ 是全忠实的。
我们必须证明 $i_W$ 是全忠实的。
可以将 $X$ 替换为 $W$，即假设 $W = X$（这样做只是为了简化记号——
注意引理陈述中的条件在此操作下保持不变）。

\medskip\noindent
设 $A, B$ 为 $D(\QCoh(\mathcal{O}_X))$ 中的对象。
我们要证明
$$
\Hom_{D(\QCoh(\mathcal{O}_X))}(A, B)
\longrightarrow
\Hom_{D(\mathcal{O}_X)}(i_X(A), i_X(B))
$$
是双射。置 $T = |X| \setminus |U|$。

\medskip\noindent
先假设 $i_X(B)$ 支撑在 $T$ 上。在此情形，映射
$$
i_X(B) \to Rj_{V, *}(i_X(B)|_V) = Rj_{V, *}(i_V(B|_V))
$$
是拟同构（引理
\ref{spaces-perfect-lemma-pushforward-with-support-in-open}）。
由假设有同构
$i_X(\Phi(B|_V)) \to Rj_{V, *}(i_V(B|_V))$，其位于
$D(\mathcal{O}_X)$ 中。
此外，$\Phi$ 与 ${-}|_V$ 是拟凝聚模导出范畴上的伴随函子（参见
“导出范畴”引理 \ref{derived-lemma-derived-adjoint-functors}）。
伴随映射 $B \to \Phi(B|_V)$ 在施加 $i_X$ 后变为同构，
因此它在 $D(\QCoh(\mathcal{O}_X))$ 中是同构。
于是
\begin{align*}
\Mor_{D(\QCoh(\mathcal{O}_X))}(A, B)
& =
\Mor_{D(\QCoh(\mathcal{O}_X))}(A, \Phi(B|_V)) \\
& =
\Mor_{D(\QCoh(\mathcal{O}_V))}(A|_V, B|_V) \\
& =
\Mor_{D(\mathcal{O}_V)}(i_V(A|_V), i_V(B|_V)) \\
& =
\Mor_{D(\mathcal{O}_X)}(i_X(A), Rj_{V, *}(i_V(B|_V))) \\
& =
\Mor_{D(\mathcal{O}_X)}(i_X(A), i_X(B))
\end{align*}
如所需。这里使用了 $i_V$ 是全忠实的这一事实
（引理 \ref{spaces-perfect-lemma-affine-coherator}）。

\medskip\noindent
一般地，选取表示 B 的复形 $\mathcal{B}^\bullet$，其各项为拟凝聚
$\mathcal{O}_X$-模，表示 $B$。接着，选取复形的任意拟同构
$s : \mathcal{B}^\bullet|_U \to \mathcal{C}^\bullet$，其中复形由
$U$ 上的拟凝聚模组成。由于 $j_U : U \to X$ 拟紧且拟分离，
函子 $j_{U, *}$ 将拟凝聚模变为拟凝聚模
（空间的态射，引理 \ref{spaces-morphisms-lemma-pushforward}）。
因此有典范映射
$\mathcal{B}^\bullet \to j_{U, *}(\mathcal{B}^\bullet|_U) \to
j_{U, *}\mathcal{C}^\bullet$，它是 $X$ 上拟凝聚模复形之间的映射。
置 $B'' = j_{U, *}\mathcal{C}^\bullet$ 于
$D(\QCoh(\mathcal{O}_X))$ 中，并选取特异三角
$$
B \to B'' \to B' \to B[1]
$$
于 $D(\QCoh(\mathcal{O}_X))$ 中。
由于该三角的第一个箭头在 $U$ 上的限制是同构，故 $B'$ 支撑在 $T$ 上。
于是，在下图中
$$
\xymatrix{
\Hom_{D(\QCoh(\mathcal{O}_X))}(A, B'[-1]) \ar[r] \ar[d] &
\Hom_{D(\mathcal{O}_X)}(i_X(A), i_X(B')[-1]) \ar[d] \\
\Hom_{D(\QCoh(\mathcal{O}_X))}(A, B) \ar[r] \ar[d] &
\Hom_{D(\mathcal{O}_X)}(i_X(A), i_X(B)) \ar[d] \\
\Hom_{D(\QCoh(\mathcal{O}_X))}(A, B'') \ar[r] \ar[d] &
\Hom_{D(\mathcal{O}_X)}(i_X(A), i_X(B'')) \ar[d] \\
\Hom_{D(\QCoh(\mathcal{O}_X))}(A, B') \ar[r] &
\Hom_{D(\mathcal{O}_X)}(i_X(A), i_X(B'))
}
$$
两列是正合的，且顶端与底端的水平箭头是双射。
最后，选取拟凝聚模复形 $\mathcal{A}^\bullet$ 表示 $A$。

\medskip\noindent
设 $\alpha : i_X(A) \to i_X(B)$ 是
$D(\mathcal{O}_X)$ 中的一个态射。 % P09-ERR-0415
限制 $\alpha|_U$ 来自 $D(\QCoh(\mathcal{O}_U))$ 中的态射，
因为 $i_U$ 是全忠实的。 % P09-ERR-0416
因此，可以选择如上的
$s : \mathcal{B}^\bullet|_U \to \mathcal{C}^\bullet$，
使得 $\alpha|_U$ 由复形的实际映射
$\mathcal{A}^\bullet|_U \to \mathcal{C}^\bullet$ 表示。
这对应于复形的映射
$\mathcal{A} \to j_{U, *}\mathcal{C}^\bullet$。
换言之，$\alpha$ 在
$\Hom_{D(\mathcal{O}_X)}(i_X(A), i_X(B''))$ 中的像
来自 $\Hom_{D(\QCoh(\mathcal{O}_X))}(A, B'')$ 中的一个元素。
通过图追逐可知，$\alpha$ 来自态射 $A \to B$，
其属于 $D(\QCoh(\mathcal{O}_X))$。
最后，设 $a : A \to B$ 是
$D(\QCoh(\mathcal{O}_X))$ 中的态射，且在 $D(\mathcal{O}_X)$ 中变为零。
适当选择 $\mathcal{B}^\bullet$ 后，可以假设 $a$ 由复形的态射
$a^\bullet : \mathcal{A}^\bullet \to \mathcal{B}^\bullet$ 表示。
由于 $i_U$ 是全忠实的，限制 $a^\bullet|_U$ 在
$D(\QCoh(\mathcal{O}_U))$ 中为零。因此，可以选择 $s$，使得
$s \circ a^\bullet|_U : \mathcal{A}^\bullet|_U \to \mathcal{C}^\bullet$
同伦于零。施加函子 $j_{U, *}$ 后，我们得到
$\mathcal{A}^\bullet \to j_{U, *}\mathcal{C}^\bullet$ 同伦于零。
因此，$a$ 在
$\Hom_{D(\QCoh(\mathcal{O}_X))}(A, B'')$ 中的像为零。
于是可以假设 $a$ 是某个元素
$b \in \Hom_{D(\QCoh(\mathcal{O}_X))}(A, B'[-1])$ 的像。
$b$ 在
$\Hom_{D(\mathcal{O}_X)}(i_X(A), i_X(B')[-1])$ 中的像
来自
$\gamma \in \Hom_{D(\mathcal{O}_X)}(A, B''[-1])$。 % P09-ERR-0417
（这是因为 $a$ 在右侧的群中映射为零。）如上所见，水平箭头是满射的，
因此 $\gamma$ 来自某个 $c$，属于
$\Hom_{D(\QCoh(\mathcal{O}_X))}(A, B''[-1])$，
从而 $a = 0$，如所需。

\medskip\noindent
此时我们知道，$i_X$ 对于原来的 $X$ 是全忠实的。
由于 $RQ_X$ 是它的右伴随，我们看到
$RQ_X \circ i_X = \text{id}$（范畴，引理
\ref{categories-lemma-adjoint-fully-faithful}）。
为完成证明，我们证明对于任意 $E$（其属于
$D_\QCoh(\mathcal{O}_X)$），映射
$i_X(RQ_X(E)) \to E$ 是同构。选取特异三角
$$
i_X(RQ_X(E)) \to E \to E' \to i_X(RQ_X(E))[1]
$$
于 $D_\QCoh(\mathcal{O}_X)$ 中。
利用上述结果的形式论证表明 $i_X(RQ_X(E')) = 0$。
因此，只需证明对于 $E \in D_\QCoh(\mathcal{O}_X)$，条件
$i_X(RQ_X(E)) = 0$ 蕴含 $E = 0$。
考虑平展态射 $j : V \to X$，其中 $V$ 为仿射。根据引理
\ref{spaces-perfect-lemma-affine-coherator} 和
\ref{spaces-perfect-lemma-flat-pushforward-coherator}，并利用我们的假设，有
$$
Rj_*(E|_V) = Rj_*(i_V(RQ_V(E|_V))) = i_X(\Phi(RQ_V(E|_V))) =
i_X(RQ_X(Rj_*(E|_V)))
$$
选取特异三角
$$
E \to Rj_*(E|_V) \to E' \to E[1]
$$
对其施加 $RQ_X$，得到特异三角
$$
0 \to RQ_X(Rj_*(E|_V)) \to RQ_X(E') \to 0[1]
$$
换言之，中间的映射是同构。
结合上面的等式串，我们发现第一个特异三角变成特异三角
$$
E \to i_X(RQ_X(E')) \to E' \to E[1]
$$
其中中间态射是伴随映射。然而，复合映射
$E \to E'$ 为零，因此由伴随性 $E \to i_X(RQ_X(E'))$ 为零！
由于该态射同构于映射
$E \to Rj_*(E|_V)$，后者是
$\text{id} : E|_V \to E|_V$ 的伴随，我们得到 $E|_V$ 为零。
这对所有平展的仿射 $V$（其位于 $X$ 上）都成立，故 $E$ 为零，如所需。
\end{proof}

\begin{proposition}
\label{spaces-perfect-proposition-quasi-compact-affine-diagonal}
设 $S$ 为概形。设 $X$ 为 $S$ 上的拟紧代数空间，且在
$\mathbf{Z}$ 上具有仿射对角态射（如“代数空间的性质”定义
\ref{spaces-properties-definition-separated} 中所述）。则函子
(\ref{spaces-perfect-equation-compare})
$$
D(\QCoh(\mathcal{O}_X))
\longrightarrow
D_\QCoh(\mathcal{O}_X)
$$
是等价，并且拟逆由 $RQ_X$ 给出。
\end{proposition}

\begin{proof}
设 $V \to W$ 为平展态射，其中 $V$ 为仿射，$W$ 为 $X$ 的拟紧开子空间。
则态射 $V \to W$ 是仿射的，因为 $W$ 在 $\mathbf{Z}$ 上具有仿射对角态射，
且 $V$ 为仿射（空间的态射，引理
\ref{spaces-morphisms-lemma-affine-permanence}）。
于是引理 \ref{spaces-perfect-lemma-affine-pushforward} 保证
引理 \ref{spaces-perfect-lemma-argument-proves} 的假设成立。
因此结论成立。
\end{proof}

\begin{lemma}
\label{spaces-perfect-lemma-direct-image-coherator}
设 $S$ 为概形，设 $f : X \to Y$ 为 $S$ 上代数空间的态射。
假设 $X$ 和 $Y$ 都拟紧，并在 $\mathbf{Z}$ 上具有仿射对角态射
（如“代数空间的性质”定义
\ref{spaces-properties-definition-separated} 中所述）。记
$$
\Phi : D(\QCoh(\mathcal{O}_X)) \to D(\QCoh(\mathcal{O}_Y))
$$
为函子
$f_* : \QCoh(\mathcal{O}_X) \to \QCoh(\mathcal{O}_Y)$ 的右导出函子，
则下图
$$
\xymatrix{
D(\QCoh(\mathcal{O}_X)) \ar[d]_\Phi \ar[r] &
D_\QCoh(\mathcal{O}_X) \ar[d]^{Rf_*} \\
D(\QCoh(\mathcal{O}_Y)) \ar[r] &
D_\QCoh(\mathcal{O}_Y)
}
$$
是交换的。 % P09-ERR-0418
\end{lemma}

\begin{proof}
注意，由命题
\ref{spaces-perfect-proposition-quasi-compact-affine-diagonal}，图中的水平箭头
都是范畴等价。因此可以将这些范畴等同（对于其他具有仿射对角态射的
拟紧代数空间也同样处理），于是引理的陈述变为：典范映射
$\Phi(K) \to Rf_*(K)$ 对所有对象 $K$ 属于
$D(\QCoh(\mathcal{O}_X))$ 都是同构。
注意，若
$K_1 \to K_2 \to K_3 \to K_1[1]$
是 $D(\QCoh(\mathcal{O}_X))$ 中的特异三角，且该陈述对其中两个成立，
则对第三个也成立。

\medskip\noindent
令 $\mathcal{B} \subset \Ob(X_{spaces, \etale})$ 为由如下对象组成的集合：
这些对象拟紧且具有仿射对角态射。
对 $U \in \mathcal{B}$ 以及任意态射 $g : U \to Z$，其中 $Z$ 是
在 $S$ 上具有仿射对角态射的拟紧代数空间，记
$$
\Phi_g : D(\QCoh(\mathcal{O}_U)) \to D(\QCoh(\mathcal{O}_Z))
$$
为 $g_*$ 的导出延拓。令
$P(U) =$ “对于任意 $K$ 属于 $D(\QCoh(\mathcal{O}_U))$，以及上述任意
$g : U \to Z$，映射 $\Phi_g(K) \to Rg_*K$ 是同构”。
根据评注 \ref{spaces-perfect-remark-how-to}，归纳原理引理
\ref{spaces-perfect-lemma-induction-principle-separated} 的条件 (1)、(2) 和 (3)(a)
成立，因此还需证明 (3)(b) 和 (4)。

\medskip\noindent
检验条件 (3)(b)。设 $U$ 为在 $X$ 上平展的仿射概形。
设 $g : U \to Z$ 如上。由于 $Z$ 的对角态射是仿射的，
态射 $g : U \to Z$ 是仿射的（空间的态射，引理
\ref{spaces-morphisms-lemma-affine-permanence}）。
因此，由引理 \ref{spaces-perfect-lemma-affine-pushforward}，$P(U)$ 成立。

\medskip\noindent
检验条件 (4)。设 $(U \subset W, V \to W)$ 为
$X_{spaces, \etale}$ 中的初等特异方形，其中 $U, W, V$ 属于
$\mathcal{B}$ 且 $V$ 为仿射。假设 $P$ 对 $U$、$V$ 以及
$U \times_W V$ 成立。
我们必须证明 $P$ 对 $W$ 成立。设 $g : W \to Z$ 为到一个
具有仿射对角态射的拟紧代数空间的态射。
设 $K$ 为 $D(\QCoh(\mathcal{O}_W))$ 中的对象。考虑特异三角
$$
K \to Rj_{U, *}K|_U \oplus Rj_{V, *}K|_V \to
Rj_{U \times_W V, *}K|_{U \times_W V} \to K[1]
$$
于 $D(\mathcal{O}_W)$ 中。由上面提到的两出三性质，
只需证明
$\Phi_g(Rj_{U, *}K|_U) \to Rg_*(Rj_{U, *}K|_U)$
是同构，并且对 $V$ 和 $U \times_W V$ 也同样成立。
这将在下一段中讨论。

\medskip\noindent
设 $j : U \to W$ 为 $X_{spaces, \etale}$ 中的态射，其中
$U, W$ 属于 $\mathcal{B}$，且 $P$ 对 $U$ 成立。 % P09-ERR-0419
设 $g : W \to Z$ 为到一个具有仿射对角态射的拟紧代数空间的态射。
为完成证明，必须证明
$\Phi_g(Rj_*K) \to Rg_*(Rj_*K)$
对于任意 $K$ 属于 $D(\QCoh(\mathcal{O}_U))$ 都是同构。
设 $\mathcal{I}^\bullet$ 为 $\QCoh(\mathcal{O}_U)$ 中表示 $K$ 的
K-单射复形。
将 $P(U)$ 应用于 $j$，可知
$j_*\mathcal{I}^\bullet$ 表示 $Rj_*K$。
由于
$j_* : \QCoh(\mathcal{O}_U) \to \QCoh(\mathcal{O}_X)$
具有正合左伴随
$j^* : \QCoh(\mathcal{O}_X) \to \QCoh(\mathcal{O}_U)$， % P09-ERR-0420
可知 $j_*\mathcal{I}^\bullet$ 是
$\QCoh(\mathcal{O}_W)$ 中的 K-单射复形；参见“导出范畴”引理
\ref{derived-lemma-adjoint-preserve-K-injectives}。
因此 $\Phi_g(Rj_*K)$ 由
$g_*j_*\mathcal{I}^\bullet = (g \circ j)_*\mathcal{I}^\bullet$ 表示。
将 $P(U)$ 应用于 $g \circ j$，可知它表示
$R_{g \circ j, *}(K) = Rg_*(Rj_*K)$。证明完成。
\end{proof}

\section{诺特空间上的凝聚化函子}
\label{spaces-perfect-section-coherator-Noetherian}

\noindent
为处理诺特代数空间的情形，我们还需要一些关于单射模的事实。

\begin{lemma}
\label{spaces-perfect-lemma-affine-injective-colimit-direct-sum-pushforwards-artin}
设 $S$ 为诺特仿射概形。$\QCoh(\mathcal{O}_S)$ 的每个单射对象都是
拟凝聚层的滤余极限 $\colim_i \mathcal{F}_i$，其中
$$
\mathcal{F}_i = (Z_i \to S)_*\mathcal{G}_i
$$
且 $Z_i$ 是某个阿廷环的谱，$\mathcal{G}_i$ 是 $Z_i$ 上的相干模。
\end{lemma}

\begin{proof}
令 $S = \Spec(A)$。设 $\mathcal{J}$ 是
$\QCoh(\mathcal{O}_S)$ 中的单射对象。由于 $\QCoh(\mathcal{O}_S)$
等价于 $A$-模的范畴，我们看到 $\mathcal{J}$ 等于某个单射模所对应的层
$\widetilde{J}$，其来源是某个单射 $A$-模 $J$。
根据《对偶化复形》命题
\ref{dualizing-proposition-structure-injectives-noetherian}，我们可以写成
$J = \bigoplus E_\alpha$，其中 $E_\alpha$ 不可分解，
因而同构于某个点处的 reside 域的单射包。 % P09-ERR-0421
因此（因为阿廷概形的有限不交并仍是阿廷的），不妨设 $J$ 是
某个素理想处 $\kappa(\mathfrak p)$ 的单射包，其中该素理想为
$\mathfrak p$，且属于 $A$。
于是 $J = \bigcup J[\mathfrak p^n]$，其中 $J[\mathfrak p^n]$ 是
$\kappa(\mathfrak p)$ 在 $A_\mathfrak/\mathfrak p^nA_\mathfrak p$ 上的单射包；见
《对偶化复形》引理 \ref{dualizing-lemma-union-artinian}。 % P09-ERR-0422
因此 $\widetilde{J}$ 是下列层的余极限：
$(Z_n \to X)_*\mathcal{G}_n$，其中 % P09-ERR-0423
$Z_n = \Spec(A_\mathfrak p/\mathfrak p^nA_\mathfrak p)$，且
$\mathfrak G_n$ 是与有限的
$A_\mathfrak/\mathfrak p^nA_\mathfrak p$-模 $J[\mathfrak p^n]$ 关联的凝聚层。
\par % P09-ERR-0424
有限性由《对偶化复形》引理
\ref{dualizing-lemma-finite} 得到。
\end{proof}

\begin{lemma}
\label{spaces-perfect-lemma-injective-colimit-direct-sum-pushforwards-artin}
设 $S$ 为仿射概形，$X$ 为 $S$ 上的诺特代数空间。$\QCoh(\mathcal{O}_X)$
中的每个单射对象都是滤余极限 $\colim_i \mathcal{F}_i$ 的直接因子，
其中这些层具有如下拟凝聚层的形式：
$$
\mathcal{F}_i = (Z_i \to X)_*\mathcal{G}_i
$$
其中 $Z_i$ 是某个阿廷环的谱，$\mathcal{G}_i$ 是 $Z_i$ 上的凝聚模。
\end{lemma}

\begin{proof}
取一个仿射概形 $U$ 以及满射平展态射 $j : U \to X$（《空间的性质》引理
\ref{spaces-properties-lemma-quasi-compact-affine-cover}）。
于是 $U$ 是诺特仿射概形。取一个单射对象 $\mathcal{J}'$，它属于
$\QCoh(\mathcal{O}_U)$，使得存在一个单射 $\mathcal{J}|_U \to \mathcal{J}'$。于是
$$
\mathcal{J} \to j_*\mathcal{J}'
$$
是 $\QCoh(\mathcal{O}_X)$ 中的单射态射，
从而将 $\mathcal{J}$ 识别为 $j_*\mathcal{J}'$ 的一个直接因子。
因此，结论由关于 $\mathcal{J}'$ 的相应结果推出；该结果已在引理
\ref{spaces-perfect-lemma-affine-injective-colimit-direct-sum-pushforwards-artin} 中证明。
\end{proof}

\begin{lemma}
\label{spaces-perfect-lemma-flat-pullback-injective-quasi-coherent}
设 $S$ 为概形。设 $f : X \to Y$ 为 $S$ 上代数空间之间的平坦、拟紧且
拟分离态射。若 $\mathcal{J}$ 是 $\QCoh(\mathcal{O}_X)$ 中的单射对象，
则 $f_*\mathcal{J}$ 是 $\QCoh(\mathcal{O}_Y)$ 中的单射对象。
\end{lemma}

\begin{proof}
由于 $f$ 拟紧且拟分离，函子 $f_*$ 将拟凝聚层变为拟凝聚层
（《空间的态射》引理 \ref{spaces-morphisms-lemma-pushforward}）。
函子 $f^*$ 是 $f_*$ 的左伴随，并且将单射变为单射。
因此结论由《同调》引理
\ref{homology-lemma-adjoint-preserve-injectives} 得到 % P09-ERR-0425
\end{proof}

\begin{lemma}
\label{spaces-perfect-lemma-injective-pushforward}
设 $S$ 为概形。设 $X$ 为 $S$ 上的诺特代数空间。若
$\mathcal{J}$ 是 $\QCoh(\mathcal{O}_X)$ 中的单射对象，则
\begin{enumerate}
\item 有 $H^p(U, \mathcal{J}|_U) = 0$，其中 $p > 0$，且 $U$ 是
在 $X$ 上平展的拟紧拟分离代数空间；
\item 对于代数空间之间的任意 $f : X \to Y$（均为 $S$ 上的代数空间），
若 $Y$ 拟分离，则有 $R^pf_*\mathcal{J} = 0$（$p > 0$）。
\end{enumerate}
\end{lemma}

\begin{proof}
第 (1) 部分的证明。将 $\mathcal{J}$ 写成
$\colim \mathcal{F}_i$ 的一个直接因子，其中
$\mathcal{F}_i = (Z_i \to X)_*\mathcal{G}_i$，
如引理 \ref{spaces-perfect-lemma-injective-colimit-direct-sum-pushforwards-artin} 所述。
显然，只需对 $\colim \mathcal{F}_i$ 证明消灭即可。
由于拉回与余极限交换，且 $U$ 拟紧拟分离，只需证明
$H^p(U, \mathcal{F}_i|_U) = 0$（$p > 0$）；见
《空间的上同调》引理 \ref{spaces-cohomology-lemma-colimits}。
注意 $Z_i \to X$ 是仿射态射；见《空间的态射》引理
\ref{spaces-morphisms-lemma-Artinian-affine}。因此
$$
\mathcal{F}_i|_U = (Z_i \times_X U \to U)_*\mathcal{G}'_i =
R(Z_i \times_X U \to U)_*\mathcal{G}'_i
$$
其中 $\mathcal{G}'_i$ 是 $\mathcal{G}_i$ 拉回到
$Z_i \times_X U$ 所得的层；见《空间的上同调》引理
\ref{spaces-cohomology-lemma-affine-base-change}。
由于 $Z_i \times_X U$ 是仿射的，我们得到
$\mathcal{G}'_i$ 在 $Z_i \times_X U$ 上没有高阶上同调。
由 Leray 谱序列可知，对 $\mathcal{F}_i|_U$ 也有同样结论
（《位点上的上同调》引理
\ref{sites-cohomology-lemma-apply-Leray}）。

\medskip\noindent
第 (2) 部分的证明。设 $f : X \to Y$ 为 $S$ 上代数空间之间的态射。
设 $V \to Y$ 为一个平展态射，且 $V$ 仿射。则
$V \times_Y X \to X$ 是平展态射，并且 $V \times_Y X$ 是一个在 $X$ 上
平展的拟紧拟分离代数空间（细节略）。因此，由第 (1) 部分
$H^p(V \times_Y X, \mathcal{J})$ 为零。
由于 $R^pf_*\mathcal{J}$ 是下述预层所关联的层：
$V \mapsto H^p(V \times_Y X, \mathcal{J})$，结论得证。
\end{proof}

\begin{lemma}
\label{spaces-perfect-lemma-Noetherian-pushforward}
设 $S$ 为概形。
设 $f : X \to Y$ 为 $S$ 上诺特代数空间之间的态射。
则拟凝聚层上的 $f_*$ 具有右导出延拓
$\Phi : D(\QCoh(\mathcal{O}_X)) \to D(\QCoh(\mathcal{O}_Y))$，
使得下图交换：
$$
\xymatrix{
D(\QCoh(\mathcal{O}_X)) \ar[d]_{\Phi} \ar[r] &
D_\QCoh(\mathcal{O}_X) \ar[d]^{Rf_*} \\
D(\QCoh(\mathcal{O}_Y)) \ar[r] &
D_\QCoh(\mathcal{O}_Y)
}
$$
\end{lemma}

\begin{proof}
由于 $X$ 和 $Y$ 是诺特的，该态射是拟紧且拟分离的（见
《空间的态射》引理
\ref{spaces-morphisms-lemma-quasi-compact-quasi-separated-permanence}）。
因此 $f_*$ 保持拟凝聚性；见《空间的态射》引理
\ref{spaces-morphisms-lemma-pushforward}。 % P09-ERR-0426
接下来，设 $K$ 是 $D(\QCoh(\mathcal{O}_X))$ 中的对象。
由于 $\QCoh(\mathcal{O}_X)$ 是 Grothendieck 阿贝尔范畴
（《空间的性质》命题 \ref{spaces-properties-proposition-coherator}），
我们可以将 $K$ 表示为一个 K-单射复形 $\mathcal{I}^\bullet$，使得每个
$\mathcal{I}^n$ 都是 $\QCoh(\mathcal{O}_X)$ 中的单射对象；见
《单射》定理 \ref{injectives-theorem-K-injective-embedding-grothendieck}。
因此，函子 $\Phi$ 定义为
$$
\Phi(K) = f_*\mathcal{I}^\bullet
$$
其中右端被视为 $D(\QCoh(\mathcal{O}_Y))$ 中的对象。为了完成引理的证明，
只需证明典范映射
$$
f_*\mathcal{I}^\bullet \longrightarrow Rf_*\mathcal{I}^\bullet
$$
是 $D(\mathcal{O}_Y)$ 中的同构。为此，只需证明该映射在上同调层上诱导同构。
任取 $m \in \mathbf{Z}$。令 $N = N(X, Y, f)$ 如引理
\ref{spaces-perfect-lemma-quasi-coherence-direct-image} 中所述。
考虑如下拟凝聚层复形的短正合序列
$$
0 \to \sigma_{\geq m - N - 1}\mathcal{I}^\bullet \to
\mathcal{I}^\bullet \to \sigma_{\leq m - N - 2}\mathcal{I}^\bullet \to 0
$$
这是 $X$ 上的复形。由
由引理 \ref{spaces-perfect-lemma-quasi-coherence-direct-image}，我们看到
$Rf_*\sigma_{\leq m - N - 2}\mathcal{I}^\bullet$ 的上同调层在
$\geq m - 1$ 的次数中为零。因此
$R^mf_*\mathcal{I}^\bullet$ 同构于
$R^mf_*\sigma_{\geq m - N - 1}\mathcal{I}^\bullet$。
换言之，我们可以假设 $\mathcal{I}^\bullet$ 是
$\QCoh(\mathcal{O}_X)$ 中单射对象构成的有下界复形。
此情形由 Leray 无环性引理（《导出范畴》引理
\ref{derived-lemma-leray-acyclicity}）以及引理
\ref{spaces-perfect-lemma-injective-pushforward} 所给出的必要消灭性推出。 % P09-ERR-0427
\end{proof}

\begin{proposition}
\label{spaces-perfect-proposition-Noetherian}
设 $S$ 为概形，$X$ 为 $S$ 上的诺特代数空间。
则函子（\ref{spaces-perfect-equation-compare}）
$$
D(\QCoh(\mathcal{O}_X))
\longrightarrow
D_\QCoh(\mathcal{O}_X)
$$
是一个等价，并且准逆由 $RQ_X$ 给出。
\end{proposition}

\begin{proof}
这立即由引理 \ref{spaces-perfect-lemma-Noetherian-pushforward} 和
\ref{spaces-perfect-lemma-argument-proves} 得到。
\end{proof}

\section{拟凝聚复形与完美复形}
\label{spaces-perfect-section-spell-out}

\noindent
本节研究《位点上的上同调》第
\ref{sites-cohomology-section-strictly-perfect}、
\ref{sites-cohomology-section-pseudo-coherent}、
\ref{sites-cohomology-section-tor} 和
\ref{sites-cohomology-section-perfect} 节中为代数空间的平展位点
定义的一般概念。特别地，我们将其与概形情形相匹配。

\medskip\noindent
首先比较概形上的拟凝聚复形与其相应小平展位点上的拟凝聚复形这一概念。

\begin{lemma}
\label{spaces-perfect-lemma-descend-finite-type}
设 $X$ 为概形，$\mathcal{F}$ 为 $\mathcal{O}_X$-模。
以下条件等价：
\begin{enumerate}
\item $\mathcal{F}$ 作为 $\mathcal{O}_X$-模是有限型的；
\item $\epsilon^*\mathcal{F}$ 作为小平展位点上的
$\mathcal{O}_\etale$-模是有限型的，且该位点属于 $X$。
\end{enumerate}
这里 $\epsilon$ 如（\ref{spaces-perfect-equation-epsilon}）中所述。
\end{lemma}

\begin{proof}
(1) $\Rightarrow$ (2) 的蕴含是一般事实，见《位点上的模》引理
\ref{sites-modules-lemma-local-pullback}。
假设 (2)。按假设，存在一个平展覆盖
$\{f_i : X_i \to X\}$，使得
$\epsilon^*\mathcal{F}|_{(X_i)_\etale}$ 由有限多个截面生成。
任取 $x \in X$。我们将证明 $\mathcal{F}$
在 $x$ 的某个邻域中由有限多个截面生成。
设 $x$ 在 $X_i \to X$ 的像中，并记 $X' = X_i$。令
$s_1, \ldots, s_n \in
\Gamma(X', \epsilon^*\mathcal{F}|_{X'_\etale})$
为生成截面。由于
$\epsilon^*\mathcal{F} =
\epsilon^{-1}\mathcal{F} \otimes_{\epsilon^{-1}\mathcal{O}_X}
\mathcal{O}_\etale$
，可以找到一个平展态射 $X'' \to X'$，使得 $x$
在 $X' \to X$ 的像中，并且
$s_i|_{X''} = \sum s_{ij} \otimes a_{ij}$，其中
$s_{ij} \in \epsilon^{-1}\mathcal{F}(X'')$，且
$a_{ij} \in \mathcal{O}_\etale(X'')$。记 $U \subset X$ 为
$X'' \to X$ 的像。这是一个开子概形，因为 $f'' : X'' \to X$ 是平展的
（《态射》引理 \ref{morphisms-lemma-etale-open}）。 % P09-ERR-0428
再将 $X''$ 缩小一些后，可假设 $s_{ij}$ 来自
$t_{ij} \in \mathcal{F}(U)$，这由逆像函子 $\epsilon^{-1}$ 的构造得到。
现在断言 $t_{ij}$ 生成 $\mathcal{F}|_U$，从而完成引理的证明。事实上，
相应的映射
$\mathcal{O}_U^{\oplus N} \to \mathcal{F}|_U$
具有如下性质：其经 $f''$ 拉回到 $X''$ 后是满射。由于
$f'' : X'' \to U$ 是概形的满射平坦态射，这意味着
$\mathcal{O}_U^{\oplus N} \to \mathcal{F}|_U$ 由考察茎并利用
$\mathcal{O}_{U, f''(z)} \to \mathcal{O}_{X'', z}$
对所有 $z \in X''$ 都忠实平坦而为满射。
\end{proof}

\noindent
在上述情形中，位点态射 $\epsilon$ 是平坦的，因而在模复形上定义拉回。

\begin{lemma}
\label{spaces-perfect-lemma-descend-pseudo-coherent}
设 $X$ 为概形，$E$ 为 $D(\mathcal{O}_X)$ 中的对象。
以下条件等价：
\begin{enumerate}
\item $E$ 是 $m$-拟凝聚的；
\item $\epsilon^*E$ 是 $m$-拟凝聚的，位于 $X$ 的小平展位点上。
\end{enumerate}
这里 $\epsilon$ 如（\ref{spaces-perfect-equation-epsilon}）中所述。
\end{lemma}

\begin{proof}
(1) $\Rightarrow$ (2) 的蕴含是一般事实，见《位点上的上同调》引理
\ref{sites-cohomology-lemma-pseudo-coherent-pullback}。
假设 $\epsilon^*E$ 是 $m$-拟凝聚的。
以下不再特别说明：$\epsilon^*$ 是正合函子，因此
$$
\epsilon^*H^i(E) = H^i(\epsilon^*E).
$$
为证明 $E$ 是 $m$-拟凝聚的，我们可以在 $X$ 上局部工作，
因而不妨设 $X$ 拟紧（例如为仿射的）。
由于 $X$ 拟紧，每个平展覆盖 $\{U_i \to X\}$
都有有限细化。因此 $\epsilon^*E$ 是
$D^{-}(\mathcal{O}_\etale)$ 中的对象；见《位点上的上同调》定义
\ref{sites-cohomology-definition-pseudo-coherent} 后的评注。
由引理 \ref{spaces-perfect-lemma-epsilon-flat} 可知 $E$ 是
$D^-(\mathcal{O}_X)$ 中的对象。

\medskip\noindent
令 $n \in \mathbf{Z}$ 为使得 $H^n(E)$ 非零的最大整数；则 $n$ 也是使得
$H^n(\epsilon^*E)$ 非零的最大整数。
我们将对 $n - m$ 作归纳来证明引理。
若 $n < m$，则引理显然成立。
若 $n \geq m$，则 $H^n(\epsilon^*E)$ 是有限的
$\mathcal{O}_\etale$-模；见《位点上的上同调》引理
\ref{sites-cohomology-lemma-finite-cohomology}。
于是 $H^n(E)$ 是有限的 $\mathcal{O}_X$-模；见引理
\ref{spaces-perfect-lemma-descend-finite-type}。
将 $X$ 换成一个开覆盖的成员后，不妨设存在满射
$\mathcal{O}_X^{\oplus t} \to H^n(E)$。
在 $X$ 上局部地可将其提升为复形映射
$\alpha : \mathcal{O}_X^{\oplus t}[-n] \to E$（细节略）。
取一个特异三角
$$
\mathcal{O}_X^{\oplus t}[-n] \to E \to C \to \mathcal{O}_X^{\oplus t}[-n + 1]
$$
则 $C$ 在次数 $\geq n$ 上的上同调消失。另一方面，复形
$\epsilon^*C$ 是 $m$-拟凝聚的，见《位点上的上同调》引理
\ref{sites-cohomology-lemma-cone-pseudo-coherent}。
因此由归纳假设，$C$ 是 $m$-拟凝聚的。再次应用《位点上的上同调》引理
\ref{sites-cohomology-lemma-cone-pseudo-coherent}，即可得出结论。
\end{proof}

\begin{lemma}
\label{spaces-perfect-lemma-descend-tor-amplitude}
设 $X$ 为概形，$E$ 为 $D(\mathcal{O}_X)$ 中的对象。则
\begin{enumerate}
\item $E$ 具有 $[a, b]$ 中的 Tor 振幅，当且仅当
$\epsilon^*E$ 具有 $[a, b]$ 中的 Tor 振幅；
\item $E$ 具有有限 Tor 维数，当且仅当 $\epsilon^*E$ 具有有限
Tor 维数。
\end{enumerate}
这里 $\epsilon$ 如（\ref{spaces-perfect-equation-epsilon}）中所述。
\end{lemma}

\begin{proof}
容易的蕴含由《位点上的上同调》引理
\ref{sites-cohomology-lemma-tor-amplitude-pullback} 得到。
反之，假设 $\epsilon^*E$ 具有 $[a, b]$ 中的 Tor 振幅。
设 $\mathcal{F}$ 为 $\mathcal{O}_X$-模。由于 $\epsilon$ 是带环位点的平坦
态射（引理 \ref{spaces-perfect-lemma-epsilon-flat}），有
$$
\epsilon^*(E \otimes^\mathbf{L}_{\mathcal{O}_X} \mathcal{F})
=
\epsilon^*E
\otimes^\mathbf{L}_{\mathcal{O}_\etale}
\epsilon^*\mathcal{F}
$$
因此，右端（按假设）上同调层的消灭性，通过引理
\ref{spaces-perfect-lemma-epsilon-flat}，推出
$E \otimes^\mathbf{L}_{\mathcal{O}_X} \mathcal{F}$ 的上同调层具有所需的消灭性。
\end{proof}

\begin{lemma}
\label{spaces-perfect-lemma-tor-dimension-rel}
设 $f : X \to Y$ 为概形的态射，$E$ 为 $D(\mathcal{O}_X)$ 中的对象。则
\begin{enumerate}
\item $E$ 作为 $D(f^{-1}\mathcal{O}_Y)$ 中的对象具有 $[a, b]$ 中的
Tor 振幅，当且仅当 $\epsilon^*E$ 具有 $[a, b]$ 中的 Tor 振幅，作为
$D(f_{small}^{-1}\mathcal{O}_{Y_\etale})$ 中的对象；
\item $E$ 作为 $D(f^{-1}\mathcal{O}_Y)$ 中的对象局部具有有限 Tor 维数，
当且仅当 $\epsilon^*E$ 作为
$D(f_{small}^{-1}\mathcal{O}_{Y_\etale})$ 中的对象局部具有有限 Tor 维数。
\end{enumerate}
这里 $\epsilon$ 如（\ref{spaces-perfect-equation-epsilon}）中所述。
\end{lemma}

\begin{proof}
(1) 中容易的方向由《位点上的上同调》引理
\ref{sites-cohomology-lemma-tor-amplitude-pullback} 得到。
设 $x \in X$ 为一点，设 $\overline{x}$ 为几何点，其位于 $x$ 之上。
令 $y = f(x)$，并记 $\overline{y}$ 为 $Y$ 中由 $\overline{x}$ 给出的几何点。
则 $(f^{-1}\mathcal{O}_Y)_x = \mathcal{O}_{Y, y}$
（《层》引理 \ref{sheaves-lemma-stalk-pullback}），并且
$$
(f_{small}^{-1}\mathcal{O}_{Y_\etale})_{\overline{x}} =
\mathcal{O}_{Y_\etale, \overline{y}} =
\mathcal{O}_{Y, y}^{sh}
$$
是严格 Hensel 化
（根据《平展上同调》引理
\ref{etale-cohomology-lemma-stalk-pullback}
和 \ref{etale-cohomology-lemma-describe-etale-local-ring}）。
由于 $\mathcal{O}_{X_\etale}$ 在 $X$ 处的茎为
$\mathcal{O}_{X, x}^{sh}$，由拉回的传递性得到 % P09-ERR-0429
$$
(\epsilon^*E)_{\overline{x}} =
E_x \otimes_{\mathcal{O}_{X, x}}^\mathbf{L} \mathcal{O}_{X, x}^{sh}
$$
若 $\epsilon^*E$ 具有 $[a, b]$ 中的 Tor 振幅，作为
$f_{small}^{-1}\mathcal{O}_{Y_\etale}$-复形，则 $(\epsilon^*E)_{\overline{x}}$
具有 $[a, b]$ 中的 Tor 振幅，作为 $\mathcal{O}_{Y, y}^{sh}$-模复形
（因为取茎就是拉回，并使用上述引理）。由《平坦性补充》引理
\ref{flat-lemma-tor-amplitude-up-down-henselization}，可知
$(\epsilon^*E)_{\overline{x}}$
作为 $\mathcal{O}_{Y, y}$-模复形的 Tor 振幅属于 $[a, b]$。
由于 $\mathcal{O}_{X, x} \to \mathcal{O}_{X, x}^{sh}$ 是忠实平坦的
（《代数补充》引理
\ref{more-algebra-lemma-dumb-properties-henselization}），且
$(\epsilon^*E)_{\overline{x}} =
E_x \otimes_{\mathcal{O}_{X, x}}^\mathbf{L} \mathcal{O}_{X, x}^{sh}$，
我们可以应用《代数补充》引理
\ref{more-algebra-lemma-no-change-tor-amplitude}，从而得到
$E_x$ 作为 $\mathcal{O}_{Y, y}$-模复形的 Tor 振幅属于 $[a, b]$。
由《上同调》引理 \ref{cohomology-lemma-tor-amplitude-stalk}，得到
$E$ 作为 $D(f^{-1}\mathcal{O}_Y)$ 中的对象具有 $[a, b]$ 中的 Tor 振幅。
这给出了 (1) 的逆向蕴含。(2) 正式地由 (1) 得到。
\end{proof}

\begin{lemma}
\label{spaces-perfect-lemma-descend-perfect}
设 $X$ 为概形，$E$ 为 $D(\mathcal{O}_X)$ 中的对象。
则 $E$ 是 $D(\mathcal{O}_X)$ 中的完美对象，当且仅当
$\epsilon^*E$ 是 $D(\mathcal{O}_\etale)$ 中的完美对象。
这里 $\epsilon$ 如（\ref{spaces-perfect-equation-epsilon}）中所述。
\end{lemma}

\begin{proof}
容易的蕴含由《位点上的上同调》引理
\ref{sites-cohomology-lemma-perfect-pullback} 中的一般结果得到。
反之，可以使用《位点上的上同调》引理
\ref{sites-cohomology-lemma-perfect} 的等价刻画，
以及关于拟凝聚复形和有限 Tor 维数复形的相应结果，即引理
\ref{spaces-perfect-lemma-descend-pseudo-coherent} 和
\ref{spaces-perfect-lemma-descend-tor-amplitude}。
略去一些细节。
\end{proof}

\begin{lemma}
\label{spaces-perfect-lemma-pseudo-coherent}
设 $S$ 为概形，$X$ 为 $S$ 上的代数空间。
若 $E$ 是 $m$-拟凝聚对象，属于 $D(\mathcal{O}_X)$，
则 $H^i(E)$ 是拟凝聚 $\mathcal{O}_X$-模，其中 $i > m$。
若 $E$ 是拟凝聚的，则 $E$ 是
$D_\QCoh(\mathcal{O}_X)$ 中的对象。
\end{lemma}

\begin{proof}
局部地，$H^i(E)$ 同构于 $H^i(\mathcal{E}^\bullet)$，
其中 $\mathcal{E}^\bullet$ 严格完美。层 $\mathcal{E}^i$
是有限自由模的直接因子，因而拟凝聚。引理得证。
\end{proof}

\begin{lemma}
\label{spaces-perfect-lemma-identify-pseudo-coherent-noetherian}
设 $S$ 为概形，$X$ 为 $S$ 上的诺特代数空间。
设 $E$ 为 $D_\QCoh(\mathcal{O}_X)$ 中的对象。对于
$m \in \mathbf{Z}$，以下条件等价：
\begin{enumerate}
\item $H^i(E)$ 对于 $i \geq m$ 是凝聚的，并且当 $i \gg 0$ 时为零；
\item $E$ 是 $m$-拟凝聚的。
\end{enumerate}
特别地，$E$ 是拟凝聚的，当且仅当 $E$ 是
$D^-_{\textit{Coh}}(\mathcal{O}_X)$ 中的对象。
\end{lemma}

\begin{proof}
由于 $X$ 拟紧，可以找到一个仿射概形 $U$ 以及满射平展态射
$U \to X$（《空间的性质》引理
\ref{spaces-properties-lemma-quasi-compact-affine-cover}）。
注意 $U$ 是诺特的。
注意，$E$ 是 $m$-拟凝聚的，当且仅当 $E|_U$ 是
$m$-拟凝聚的（这由定义或《位点上的上同调》引理
\ref{sites-cohomology-lemma-pseudo-coherent-independent-representative} 推出）。
同样，$H^i(E)$ 是凝聚的，当且仅当
$H^i(E)|_U = H^i(E|_U)$ 是凝聚的（见《空间的上同调》引理
\ref{spaces-cohomology-lemma-coherent-Noetherian}）。
因此不妨设 $X$ 可表。

\medskip\noindent
若 $X$ 可由概形 $X_0$ 表示，则（引理
\ref{spaces-perfect-lemma-derived-quasi-coherent-small-etale-site}）
可写 $E = \epsilon^*E_0$，其中 $E_0$ 是
$D_\QCoh(\mathcal{O}_{X_0})$ 中的对象，且
$\epsilon : X_\etale \to (X_0)_{Zar}$ 如
（\ref{spaces-perfect-equation-epsilon}）中所述。
在此情形中，由引理 \ref{spaces-perfect-lemma-descend-pseudo-coherent}，
$E$ 是 $m$-拟凝聚的，当且仅当 $E_0$ 也是如此。
同样，由引理 \ref{spaces-perfect-lemma-descend-finite-type}，$H^i(E_0)$
是有限型的（即凝聚的），当且仅当 $H^i(E)$ 也是如此。
最后，由引理 \ref{spaces-perfect-lemma-epsilon-flat}，$H^i(E_0) = 0$
当且仅当 $H^i(E) = 0$。
因此归约到概形情形，即《概形的导出范畴》引理
\ref{perfect-lemma-identify-pseudo-coherent-noetherian}。
\end{proof}

\begin{lemma}
\label{spaces-perfect-lemma-tor-qc-qs}
设 $S$ 为概形，$X$ 为 $S$ 上的拟分离代数空间。
设 $E$ 为 $D_\QCoh(\mathcal{O}_X)$ 中的对象。设 $a \leq b$。
以下条件等价：
\begin{enumerate}
\item $E$ 具有 $[a, b]$ 中的 Tor 振幅；
\item 对所有 $\mathcal{F}$，若其属于 $\QCoh(\mathcal{O}_X)$，则均有
$H^i(E \otimes_{\mathcal{O}_X}^\mathbf{L} \mathcal{F}) = 0$，
其中 $i \not \in [a, b]$。
\end{enumerate}
\end{lemma}

\begin{proof}
显然 (1) 蕴含 (2)。假设 (2)。设 $j : U \to X$ 为平展态射，
其中 $U$ 仿射。由于 $X$ 拟分离，$j$ 既拟紧又分离。
因此 $j_*$ 将拟凝聚模变为拟凝聚模（《空间的态射》引理
\ref{spaces-morphisms-lemma-pushforward}）。
任取拟凝聚模 $\mathcal{G}$，其定义在 $U$ 上。
由假设，
$H^i(E \otimes_{\mathcal{O}_X}^\mathbf{L} j_*\mathcal{G})$
在 $i \not \in [a, b]$ 时消失。经平坦态射 $j$ 拉回，得到
$H^i(E|_U \otimes_{\mathcal{O}_U}^\mathbf{L} j^*j_*\mathcal{G})$
在 $i \not \in [a, b]$ 时消失。由《空间的上同调》引理
\ref{spaces-cohomology-lemma-etale-pull-push-split}，
$\mathcal{G}$ 是 $j^*j_*\mathcal{G}$ 的直接因子。
因此条件 (2) 推出
$H^i(E|_U \otimes_{\mathcal{O}_U}^\mathbf{L} \mathcal{G})$
在 $i \not \in [a, b]$ 时消失，这对每个拟凝聚
$\mathcal{O}_U$-模 $\mathcal{G}$ 都成立。由于只需证明
$E|_U$ 具有 $[a, b]$ 中的 Tor 振幅，故归约到 $X$ 可表的情形。

\medskip\noindent
若 $X$ 可由概形 $X_0$ 表示，则（引理
\ref{spaces-perfect-lemma-derived-quasi-coherent-small-etale-site}）
可写 $E = \epsilon^*E_0$，其中 $E_0$ 是
$D_\QCoh(\mathcal{O}_{X_0})$ 中的对象，且
$\epsilon : X_\etale \to (X_0)_{Zar}$ 如
（\ref{spaces-perfect-equation-epsilon}）中所述。对于每个拟凝聚模
$\mathcal{F}_0$（定义在 $X_0$ 上），模 $\epsilon^*\mathcal{F}_0$ 在 $X$ 上拟凝聚，并且
$$
H^i(E \otimes_{\mathcal{O}_X}^\mathbf{L} \epsilon^*\mathcal{F}_0)
=
\epsilon^*H^i(E_0 \otimes_{\mathcal{O}_{X_0}}^\mathbf{L} \mathcal{F}_0)
$$
因为 $\epsilon$ 平坦（引理 \ref{spaces-perfect-lemma-epsilon-flat}）。
此外，这些层在 $i \not \in [a, b]$ 时消失，由同一引理可推出
$H^i(E_0 \otimes_{\mathcal{O}_{X_0}}^\mathbf{L} \mathcal{F}_0)$
也具有相同的消灭性。因此问题归约到概形情形，该情形在
《概形的导出范畴》引理 \ref{perfect-lemma-tor-qc-qs} 中处理。
\end{proof}

\begin{lemma}
\label{spaces-perfect-lemma-descend-RHom}
设 $X$ 为概形，$E, F$ 为 $D(\mathcal{O}_X)$ 中的对象。
假设下列条件之一成立：
\begin{enumerate}
\item $E$ 是拟凝聚的，且 $F$ 属于 $D^+(\mathcal{O}_X)$；
\item $E$ 是完美的，而 $F$ 任意。
\end{enumerate}
则存在典范同构
$$
\epsilon^*R\SheafHom(E, F) \longrightarrow R\SheafHom(\epsilon^*E, \epsilon^*F)
$$
这里 $\epsilon$ 如（\ref{spaces-perfect-equation-epsilon}）中所述。
\end{lemma}

\begin{proof}
回忆 $\epsilon$ 是平坦的（引理 \ref{spaces-perfect-lemma-epsilon-flat}），
因而 $\epsilon^* = L\epsilon^*$。由《位点上的上同调》评注
\ref{sites-cohomology-remark-prepare-fancy-base-change}，
存在从左端到右端的典范映射。
为了证明它是同构，可以局部工作，即不妨设 $X$ 为仿射概形。

\medskip\noindent
在情形 (1) 中，可以将 $E$ 表示为有上界复形
$\mathcal{E}^\bullet$，它由有限自由 $\mathcal{O}_X$-模构成；见《概形的导出范畴》引理
\ref{perfect-lemma-lift-pseudo-coherent}。
也可将 $F$ 表示为有下界复形 $\mathcal{F}^\bullet$，
它由 $\mathcal{O}_X$-模构成。应用《上同调》引理
\ref{cohomology-lemma-Rhom-complex-of-direct-summands-finite-free}，
可知 $R\SheafHom(E, F)$ 由各项为
$$
\bigoplus\nolimits_{n = - p + q}
\SheafHom_{\mathcal{O}_X}(\mathcal{E}^p, \mathcal{F}^q)
$$
的复形表示。应用《位点上的上同调》引理
\ref{sites-cohomology-lemma-Rhom-complex-of-direct-summands-finite-free}，
可知 $R\SheafHom(\epsilon^*E, \epsilon^*F)$ 由各项为
$$
\bigoplus\nolimits_{n = - p + q}
\SheafHom_{\mathcal{O}_\etale}
(\epsilon^*\mathcal{E}^p, \epsilon^*\mathcal{F}^q)
$$
的复形表示。因此引理的陈述归结为以下事实：典范映射
$$
\epsilon^*\SheafHom_{\mathcal{O}_X}(\mathcal{E}, \mathcal{F})
\longrightarrow
\SheafHom_{\mathcal{O}_\etale}
(\epsilon^*\mathcal{E}, \epsilon^*\mathcal{F})
$$
对于任意 $\mathcal{O}_X$-模 $\mathcal{F}$ 以及有限自由
$\mathcal{O}_X$-模 $\mathcal{E}$ 都是同构。

\medskip\noindent
在情形 (2) 中，可以将 $E$ 表示为严格完美复形
$\mathcal{E}^\bullet$，它由 $\mathcal{O}_X$-模构成；使用《概形的导出范畴》引理
\ref{perfect-lemma-affine-compare-bounded} 和
\ref{perfect-lemma-perfect-affine}，以及模的完美复形可由有限生成投射模构成的
有限复形表示这一事实。于是可以逐字使用上面的证明，只需将对《上同调》引理
\ref{cohomology-lemma-Rhom-complex-of-direct-summands-finite-free}
的引用替换为对《上同调》引理
\ref{cohomology-lemma-Rhom-strictly-perfect} 的引用。
\end{proof}

\begin{lemma}
\label{spaces-perfect-lemma-quasi-coherence-internal-hom}
设 $S$ 为概形，$X$ 为 $S$ 上的代数空间。
设 $L, K$ 为 $D(\mathcal{O}_X)$ 中的对象。
若下列条件之一成立：
\begin{enumerate}
\item $L$ 属于 $D^+_\QCoh(\mathcal{O}_X)$，且 $K$ 是拟凝聚的；
\item $L$ 属于 $D_\QCoh(\mathcal{O}_X)$，且 $K$ 是完美的，
\end{enumerate} % P09-ERR-0430
则 $R\SheafHom(K, L)$ 属于 $D_\QCoh(\mathcal{O}_X)$。
\end{lemma}

\begin{proof}
这由概形的相应结果（《概形的导出范畴》引理
\ref{perfect-lemma-quasi-coherence-internal-hom}）推出，
其中使用引理 \ref{spaces-perfect-lemma-check-quasi-coherence-on-covering} 的判别准则、
引理 \ref{spaces-perfect-lemma-descend-pseudo-coherent} 和
\ref{spaces-perfect-lemma-descend-perfect} 的判别准则，以及引理
\ref{spaces-perfect-lemma-descend-RHom} 的结果。
\end{proof}

\begin{lemma}
\label{spaces-perfect-lemma-internal-hom-evaluate-tensor-isomorphism}
设 $S$ 为概形。
设 $X$ 为 $S$ 上的代数空间。
设 $K, L, M$ 为 $D_\QCoh(\mathcal{O}_X)$ 中的对象。
《位点上的上同调》引理
\ref{sites-cohomology-lemma-internal-hom-diagonal-better} 中的映射
$$
K \otimes_{\mathcal{O}_X}^\mathbf{L} R\SheafHom(M, L)
\longrightarrow
R\SheafHom(M, K \otimes_{\mathcal{O}_X}^\mathbf{L} L)
$$
在下列情形中为同构：
\begin{enumerate}
\item $M$ 是完美的；或
\item $K$ 是完美的；或
\item $M$ 是拟凝聚的，$L \in D^+(\mathcal{O}_X)$，且 $K$ 具有有限
Tor 维数。
\end{enumerate}
\end{lemma}

\begin{proof}
映射是否为同构可在平展局部检验，因而不妨设 $X$ 为仿射概形。
由引理 \ref{spaces-perfect-lemma-derived-quasi-coherent-small-etale-site}，
可将 $K, L, M$ 写成 $\epsilon^*K_0, \epsilon^*L_0, \epsilon^*M_0$，
其中某些 $K_0, L_0, M_0$ 属于 $D_\QCoh(\mathcal{O}_X)$。
于是由引理 \ref{spaces-perfect-lemma-descend-RHom} 可见，情形 (1) 和 (3)
化归为概形的情形，即《概形的导出范畴》引理
\ref{perfect-lemma-internal-hom-evaluate-tensor-isomorphism}。
如果 $K$ 是完美的而不作其他假设，则我们不知道箭头的任一端
是否属于 $D_\QCoh(\mathcal{O}_X)$；但结论仍然成立，因为进一步局部化后，
$K$ 可由有限自由模构成的有限复形表示，此时结论显然。
\end{proof}









\section{用完美复形作逼近}
\label{spaces-perfect-section-approximation}

\noindent
本节继续《概形的导出范畴》第
\ref{perfect-section-approximation} 节中开始的讨论。

\begin{definition}
\label{spaces-perfect-definition-approximation-holds}
设 $S$ 为概形，$X$ 为 $S$ 上的代数空间。
考虑三元组 $(T, E, m)$，其中
\begin{enumerate}
\item $T \subset |X|$ 是闭子集；
\item $E$ 是 $D_\QCoh(\mathcal{O}_X)$ 中的对象；且
\item $m \in \mathbf{Z}$。
\end{enumerate}
所谓三元组 $(T, E, m)$ 的{\it 逼近成立}，是指存在完美对象 $P$，
它是 $D(\mathcal{O}_X)$ 中支撑在 $T$ 上的对象，并存在映射
$\alpha : P \to E$，它诱导同构 $H^i(P) \to H^i(E)$（对 $i > m$），
以及满射 $H^m(P) \to H^m(E)$。
\end{definition}

\noindent
逼近不可能对每个三元组都成立。其原因见《概形的导出范畴》定义
\ref{perfect-definition-approximation-holds} 后的评注。

\begin{definition}
\label{spaces-perfect-definition-approximation}
设 $S$ 为概形，$X$ 为 $S$ 上的代数空间。
所谓在 $X$ 上{\it 用完美复形作逼近成立}，是指对于任意闭子集
$T \subset |X|$，只要态射
$X \setminus T \to X$ 拟紧，就存在整数 $r$，使得对定义
\ref{spaces-perfect-definition-approximation-holds} 中任意满足下列条件的三元组
$(T, E, m)$：
\begin{enumerate}
\item $E$ 是 $(m - r)$-拟凝聚的；且
\item $H^i(E)$ 支撑在 $T$ 上（对 $i \geq m - r$），
\end{enumerate}
逼近均成立。
\end{definition}

\begin{lemma}
\label{spaces-perfect-lemma-pushforward-perfect}
设 $S$ 为概形，$(U \subset X, j : V \to X)$ 为 $S$ 上代数空间的
初等特异方形。设 $E$ 为 $D(\mathcal{O}_V)$ 中支撑在
$j^{-1}(T)$ 上的完美对象，其中 $T = |X| \setminus |U|$。
则 $Rj_*E$ 是 $D(\mathcal{O}_X)$ 中的完美对象。
\end{lemma}

\begin{proof}
完美性在 $X_\etale$ 上是局部的。因此，只需检验 $Rj_*E$
在 $U$ 和 $V$ 上的限制是完美的。由引理
\ref{spaces-perfect-lemma-pushforward-with-support-in-open}，有 $Rj_*E|_V = E$，
故其完美。又有 $Rj_*E|_U = 0$，因为
$E|_{V \setminus j^{-1}(T)} = 0$（使用引理
\ref{spaces-perfect-lemma-restrict-direct-image-open}）。
\end{proof}

\begin{lemma}
\label{spaces-perfect-lemma-open}
设 $S$ 为概形，$(U \subset X, j : V \to X)$ 为 $S$ 上代数空间的
初等特异方形。设 $T$ 为 $|X| \setminus |U|$ 的闭子集，并设
$(T, E, m)$ 为定义 \ref{spaces-perfect-definition-approximation-holds} 中的三元组。
若
\begin{enumerate}
\item 三元组 $(j^{-1}T, E|_V, m)$ 的逼近成立；且
\item 层 $H^i(E)$（对 $i \geq m$）支撑在 $T$ 上，
\end{enumerate}
则三元组 $(T, E, m)$ 的逼近成立。
\end{lemma}

\begin{proof}
设 $P \to E|_V$ 是三元组 $(j^{-1}T, E|_V, m)$ 在 $V$ 上的一个逼近。
由引理 \ref{spaces-perfect-lemma-pushforward-perfect}，$Rj_*P$ 是
$D(\mathcal{O}_X)$ 中的完美对象。另一方面，由引理
\ref{spaces-perfect-lemma-pushforward-with-support-in-open}，$Rj_*P = j_!P$。
例如由（\ref{spaces-perfect-equation-stalk-j-shriek}）可见，$j_!P$ 支撑在 $T$ 上。
因此得到逼近 $Rj_*P = j_!P \to j_!(E|_V) \to E$。
\end{proof}

\begin{lemma}
\label{spaces-perfect-lemma-approximation-affine}
设 $S$ 为概形，$X$ 为 $S$ 上可由仿射概形表示的代数空间。
若定义 \ref{spaces-perfect-definition-approximation-holds} 中的三元组
$(T, E, m)$ 满足：存在整数 $r \geq 0$ 使得
\begin{enumerate}
\item $E$ 是 $m$-拟凝聚的；
\item $H^i(E)$ 支撑在 $T$ 上（对 $i \geq m - r + 1$）；
\item $X \setminus T$ 是 $r$ 个仿射开集的并，
\end{enumerate}
则其逼近成立。特别地，仿射概形上的完美复形逼近成立。
\end{lemma}

\begin{proof}
设 $X_0$ 为表示 $X$ 的仿射概形。设 $T_0 \subset X_0$
为与 $T$ 对应的闭子集。% P09-ERR-0431
令 $\epsilon : X_\etale \to X_{0, Zar}$ 为（\ref{spaces-perfect-equation-epsilon}）
中的态射。由引理 \ref{spaces-perfect-lemma-derived-quasi-coherent-small-etale-site}，
可写成 $E = \epsilon^*E_0$，其中某个对象 $E_0$ 属于
$D_\QCoh(\mathcal{O}_{X_0})$。
由引理 \ref{spaces-perfect-lemma-descend-pseudo-coherent}，$E_0$ 是
$m$-拟凝聚的。比较上同调层的茎（见引理
\ref{spaces-perfect-lemma-epsilon-flat} 的证明），可见 $H^i(E_0)$ 支撑在 $T_0$ 上
（对 $i \geq m - r + 1$）。由《概形的导出范畴》引理
\ref{perfect-lemma-approximation-affine}，存在逼近 $P_0 \to E_0$，
它对应于三元组 $(T_0, E_0, m)$。
由引理 \ref{spaces-perfect-lemma-descend-perfect}，可见
$P = \epsilon^*P_0$ 是 $D(\mathcal{O}_X)$ 中的完美对象。
拉回后得到所需逼近 $P = \epsilon^*P_0 \to \epsilon^*E_0 = E$。
\end{proof}

\begin{lemma}
\label{spaces-perfect-lemma-induction-step}
设 $S$ 为概形，$(U \subset X, j : V \to X)$ 为 $S$ 上代数空间的
初等特异方形。假设 $U$ 拟紧、$V$ 仿射，且 $U \times_X V$ 拟紧。
若 $U$ 上用完美复形作逼近成立，则 $X$ 上用完美复形作逼近也成立。
\end{lemma}

\begin{proof}
设 $T \subset |X|$ 为闭子集，且 $X \setminus T \to X$ 拟紧。
设 $r_U$ 为定义 \ref{spaces-perfect-definition-approximation} 中适合于对子
$(U, T \cap |U|)$ 的整数。令 $T' = T \setminus |U|$，并赋予
$T'$ 诱导的既约子空间结构。由于 $|T'|$ 包含于
$|X| \setminus |U|$，可见 $j^{-1}(T') \to T'$ 是同构。此外，
$V \setminus j^{-1}(T')$ 拟紧，因为它是 $U \times_X V$ 与
$X \setminus T$ 在 $X$ 上的纤维积，并且已假设
$U \times_X V$ 拟紧且 $X \setminus T \to X$ 拟紧。
设 $r'$ 为覆盖 $V \setminus j^{-1}(T')$ 所需的仿射开集数。
我们断言 $r = \max(r_U, r')$ 对于对子 $(X, T)$ 可用。

\medskip\noindent
为证明此点，选取三元组 $(T, E, m)$，使得 $E$ 是
$(m - r)$-拟凝聚的，且 $H^i(E)$ 支撑在 $T$ 上（对
$i \geq m - r$）。设 $t$ 为使 $H^t(E)|_U$ 非零的最大整数。
（这样的整数存在，因为 $U$ 拟紧，且 $E|_U$ 是
$(m - r)$-拟凝聚的。）我们将证明 $E$ 可被逼近，并对 $t$ 作归纳。

\medskip\noindent
基本情形：$t \leq m - r'$。这意味着 $H^i(E)$ 支撑在
$T'$ 上（对 $i \geq m - r'$）。因此引理
\ref{spaces-perfect-lemma-approximation-affine} 保证存在逼近
$P \to E|_V$，它是三元组 $(T', E|_V, m)$ 在 $V$ 上的逼近。
应用引理 \ref{spaces-perfect-lemma-open} 可见 $(T', E, m)$ 可被逼近。
这样的逼近也是 $(T, E, m)$ 的逼近。

\medskip\noindent
归纳步骤。选取逼近 $P \to E|_U$，它对应于
$(T \cap |U|, E|_U, m)$。这尤其给出满射
$H^t(P) \to H^t(E|_U)$。
在证明的余下部分，我们将对可表的代数空间 $V$ 和
$U \times_X V$ 使用引理
\ref{spaces-perfect-lemma-derived-quasi-coherent-small-etale-site} 的等价
（以及评注 \ref{spaces-perfect-remark-match-total-direct-images} 的相容性）。
我们还将使用如下事实：$(m - r)$-拟凝聚性以及完美性在
Zariski 位点和平展位点上分别相同；见引理
\ref{spaces-perfect-lemma-descend-pseudo-coherent} 和
\ref{spaces-perfect-lemma-descend-perfect}。
因此，可对开浸入 $U \times_X V \subset V$ 使用
《概形的导出范畴》第 \ref{perfect-section-lift} 节的结果。
由此，《概形的导出范畴》引理
\ref{perfect-lemma-lift-perfect-complex-plus-shift-support} 表明，
存在完美对象 $Q$，它属于 $D(\mathcal{O}_V)$ 且支撑在 $j^{-1}(T)$ 上，
以及同构
$Q|_{U \times_X V} \to (P \oplus P[1])|_{U \times_X V}$。
由《概形的导出范畴》引理 \ref{perfect-lemma-lift-map}，
可将 $Q$ 替换为 $Q \otimes^\mathbf{L} I$，并且可以假设映射
$$
Q|_{U \times_X V} \longrightarrow
(P \oplus P[1])|_{U \times_X V} \longrightarrow
P|_{U \times_X V} \longrightarrow E|_{U \times_X V}
$$
可提升为 $Q \to E|_V$。由引理 \ref{spaces-perfect-lemma-glue}，得到
态射 $a : R \to E$，它是 $D(\mathcal{O}_X)$ 中的态射，并使得 $a|_U$
同构于 $P \oplus P[1] \to E|_U$，而 $a|_V$ 同构于
$Q \to E|_V$。% P09-ERR-0432
于是 $R$ 是支撑在 $T$ 上的完美对象，并且映射
$H^t(R) \to H^t(E)$ 限制到 $U$ 上是满射。选取特异三角
$$
R \to E \to E' \to R[1]
$$
则 $E'$ 是 $(m - r)$-拟凝聚的
（《位点上的上同调》引理
\ref{sites-cohomology-lemma-cone-pseudo-coherent}），
$H^i(E')|_U = 0$（对 $i \geq t$），且 $H^i(E')$ 支撑在
$T$ 上（对 $i \geq m - r$）。归纳地，得到逼近 $R' \to E'$，
它对应于三元组 $(T, E', m)$。将复合
$R' \to E' \to R[1]$ 嵌入特异三角
$R \to R'' \to R' \to R[1]$，并将态射 $R' \to E'$ 和
$R[1] \to R[1]$ 扩充成特异三角的态射
$$
\xymatrix{
R \ar[r] \ar[d] & R'' \ar[d] \ar[r] & R' \ar[d] \ar[r] & R[1] \ar[d] \\
R \ar[r] &  E \ar[r] & E' \ar[r] & R[1]
}
$$
这里使用 TR3。于是 $R''$ 是完美复形
（《位点上的上同调》引理
\ref{sites-cohomology-lemma-two-out-of-three-perfect}），且支撑在 $T$ 上。
简单的图追踪表明，$R'' \to E$ 就是所需的逼近。
\end{proof}

\begin{theorem}
\label{spaces-perfect-theorem-approximation}
设 $S$ 为概形。
设 $X$ 为 $S$ 上拟紧且拟分离的代数空间。
则 $X$ 上用完美复形作逼近成立。
\end{theorem}

\begin{proof}
这由引理 \ref{spaces-perfect-lemma-induction-principle} 的归纳原理以及引理
\ref{spaces-perfect-lemma-induction-step} 和 \ref{spaces-perfect-lemma-approximation-affine} 推出。
\end{proof}






\section{生成导出范畴}
\label{spaces-perfect-section-generating}

\noindent
本节是《概形的导出范畴》第 \ref{perfect-section-generating} 节的类似物。
不过，我们先证明下述引理；它是《概形的导出范畴》引理
\ref{perfect-lemma-lift-map-from-perfect-complex-with-support} 的类似物。

\begin{lemma}
\label{spaces-perfect-lemma-lift-map-from-perfect-complex-with-support}
设 $S$ 为概形，$X$ 为 $S$ 上拟紧且拟分离的代数空间。
设 $W \subset X$ 为拟紧开集。设 $T \subset |X|$ 为闭子集，使得
$X \setminus T \to X$ 是拟紧态射。设 $E$ 为
$D_\QCoh(\mathcal{O}_X)$ 中的对象。设 $\alpha : P \to E|_W$ 为映射，
其中 $P$ 是 $D(\mathcal{O}_W)$ 中支撑在 $T \cap W$ 上的完美对象。
则存在映射 $\beta : R \to E$，其中 $R$ 是
$D(\mathcal{O}_X)$ 中支撑在 $T$ 上的完美对象，使得 $P$ 是
$R|_W$ 的直和项（在 $D(\mathcal{O}_W)$ 中），并且与 $\alpha$ 和
$\beta|_W$ 相容。% P09-ERR-0433
\end{lemma}

\begin{proof}
我们将用引理 \ref{spaces-perfect-lemma-induction-principle-enlarge} 的归纳原理证明此事。
因此立即化归到如下情形：有初等特异方形
$(W \subset X, f : V \to X)$，其中 $V$ 仿射，且
$P \to E|_W$ 如引理陈述中所述。在证明的余下部分，我们将对可表代数空间
$V$ 和 $W \times_X V$ 使用引理
\ref{spaces-perfect-lemma-derived-quasi-coherent-small-etale-site}
（以及评注 \ref{spaces-perfect-remark-match-total-direct-images} 的相容性）。
我们还将使用如下事实：完美性在 Zariski 位点和平展位点上相同；
见引理 \ref{spaces-perfect-lemma-descend-perfect}。

\medskip\noindent
由《概形的导出范畴》引理
\ref{perfect-lemma-lift-perfect-complex-plus-shift-support}，
可选取完美对象 $Q$，它属于 $D(\mathcal{O}_V)$，支撑在 $f^{-1}T$ 上，
并且有同构
$Q|_{W \times_X V} \to (P \oplus P[1])|_{W \times_X V}$。
由《概形的导出范畴》引理 \ref{perfect-lemma-lift-map}，
可将 $Q$ 替换为 $Q \otimes^\mathbf{L} I$（仍支撑在 $f^{-1}T$ 上），
并且可以假设映射
$$
Q|_{W \times_X V} \to (P \oplus P[1])|_{W \times V} % P09-ERR-0434
\longrightarrow P|_{W \times_X V}
\longrightarrow
E|_{W \times_X V}
$$
可提升为 $Q \to E|_V$。由引理 \ref{spaces-perfect-lemma-glue}，得到态射
$a : R \to E$，它是 $D(\mathcal{O}_X)$ 中的态射，使得 $a|_W$
同构于 $P \oplus P[1] \to E|_W$，而 $a|_V$ 同构于
$Q \to E|_V$。% P09-ERR-0435
于是 $R$ 是支撑在 $T$ 上的完美对象，如所要求。
\end{proof}

\begin{remark}
\label{spaces-perfect-remark-addendum}
引理 \ref{spaces-perfect-lemma-lift-map-from-perfect-complex-with-support} 的证明表明
$$
R|_W = P \oplus P^{\oplus n_1}[1] \oplus \ldots \oplus P^{\oplus n_m}[m]
$$
其中某些 $m \geq 0$ 且 $n_j \geq 0$。因此，$R|_W$ 的最高次非零
上同调层等于 $P$ 的最高次非零上同调层。对映射
$P^{\oplus n_1}[1] \oplus \ldots \oplus P^{\oplus n_m}[m] \to R|_W$
重复该构造、取锥并使用归纳，即可使 $R|_W$ 与 $P$ 在任意给定次数
以上的上同调层相等。
\end{remark}

\begin{lemma}
\label{spaces-perfect-lemma-direct-summand-of-a-restriction}
设 $S$ 为概形。
设 $X$ 为 $S$ 上拟紧且拟分离的代数空间。
设 $W$ 为 $X$ 的拟紧开子空间。
设 $P$ 为 $D(\mathcal{O}_W)$ 中的完美对象。
则 $P$ 是 $D(\mathcal{O}_X)$ 中某个完美对象之限制的直和项。
\end{lemma}

\begin{proof}
这是引理 \ref{spaces-perfect-lemma-lift-map-from-perfect-complex-with-support} 的特殊情形。
\end{proof}

\begin{theorem}
\label{spaces-perfect-theorem-bondal-van-den-Bergh}
设 $S$ 为概形，$X$ 为 $S$ 上拟紧且拟分离的代数空间。
范畴 $D_\QCoh(\mathcal{O}_X)$ 可由单个完美对象生成。
更精确地，存在完美对象 $P$，它属于 $D(\mathcal{O}_X)$，使得对于
$E \in D_\QCoh(\mathcal{O}_X)$，下列条件等价：
\begin{enumerate}
\item $E = 0$；且
\item $\Hom_{D(\mathcal{O}_X)}(P[n], E) = 0$（对所有
$n \in \mathbf{Z}$）。
\end{enumerate}
\end{theorem}

\begin{proof}
我们将使用引理 \ref{spaces-perfect-lemma-induction-principle} 的归纳原理证明此事。

\medskip\noindent
若 $X$ 仿射，则 $\mathcal{O}_X$ 是完美生成元。这由引理
\ref{spaces-perfect-lemma-derived-quasi-coherent-small-etale-site} 和
《概形的导出范畴》引理
\ref{perfect-lemma-affine-compare-bounded} 推出。

\medskip\noindent
假设 $(U \subset X, f : V \to X)$ 是初等特异方形，其中 $U$ 拟紧，
定理对 $U$ 成立，且 $V$ 是仿射概形。
设 $P$ 为 $D(\mathcal{O}_U)$ 中的完美对象，它是
$D_\QCoh(\mathcal{O}_U)$ 的生成元。使用引理
\ref{spaces-perfect-lemma-direct-summand-of-a-restriction}，可选取完美对象
$Q$，它属于 $D(\mathcal{O}_X)$，其在 $U$ 上的限制是某个直和，
且 $P$ 是该直和的一个直和项。记 $V = \Spec(A)$。设 $Z \subset V$
为 $X \setminus U$ 的逆像所给的既约闭子概形，并同构地映到该逆像
（见定义 \ref{spaces-perfect-definition-elementary-distinguished-square}）。
这是 $V$ 的逆紧闭子集。选取 $f_1, \ldots, f_r \in A$，使得
$Z = V(f_1, \ldots, f_r)$。设 $K \in D(\mathcal{O}_V)$ 为与
$f_1, \ldots, f_r$ 在 $A$ 上的 Koszul 复形相对应的完美对象。
注意，因为 $K$ 支撑在 $Z$ 上，推前
$K' = Rf_*K$ 是 $D(\mathcal{O}_X)$ 中的完美对象，其在 $V$ 上的限制
是 $K$（见引理 \ref{spaces-perfect-lemma-pushforward-perfect} 和
\ref{spaces-perfect-lemma-pushforward-with-support-in-open}）。
我们断言 $Q \oplus K'$ 是 $D_\QCoh(\mathcal{O}_X)$ 的生成元。

\medskip\noindent
设 $E$ 为 $D_\QCoh(\mathcal{O}_X)$ 中的对象，使得从
$Q \oplus K'$ 的任意平移到 $E$ 都没有非零映射。
由引理 \ref{spaces-perfect-lemma-pushforward-with-support-in-open}，有
$K' =  f_! K$，因而
$$
\Hom_{D(\mathcal{O}_X)}(K'[n], E) = \Hom_{D(\mathcal{O}_V)}(K[n], E|_V)
$$
于是由《概形的导出范畴》引理
\ref{perfect-lemma-orthogonal-koszul-complex}
（还使用引理 \ref{spaces-perfect-lemma-derived-quasi-coherent-small-etale-site}），
这些群的消失蕴含 $E|_V$ 同构于
$R(U \times_X V \to V)_*E|_{U \times_X V}$。这又蕴含
$E = R(U \to X)_*E|_U$（略去一个小细节）。在此情形下
$$
\Hom_{D(\mathcal{O}_X)}(Q[n], E) = \Hom_{D(\mathcal{O}_U)}(Q|_U[n], E|_U)
$$
其中含有 $\Hom_{D(\mathcal{O}_U)}(P[n], E|_U)$ 作为直和项。
因此，根据对 $P$ 的选取，这些群的消失蕴含 $E|_U$ 为零，
从而 $E$ 为零。
\end{proof}












\begin{remark}
\label{spaces-perfect-remark-pullback-generator}
设 $S$ 为概形。
设 $f : X \to Y$ 为 $S$ 上拟紧且拟分离的代数空间之间的态射。
设 $E \in D_\QCoh(\mathcal{O}_Y)$ 为生成元
（见定理 \ref{spaces-perfect-theorem-bondal-van-den-Bergh}）。
则下列条件等价：
\begin{enumerate}
\item 对 $K \in D_\QCoh(\mathcal{O}_X)$，有
$Rf_*K = 0$ 当且仅当 $K = 0$；
\item $Rf_* : D_\QCoh(\mathcal{O}_X) \to D_\QCoh(\mathcal{O}_Y)$
反映同构；且
\item $Lf^*E$ 是 $D_\QCoh(\mathcal{O}_X)$ 的生成元。
\end{enumerate}
(1) 与 (2) 的等价性是如下事实的形式推论：
$Rf_* : D_\QCoh(\mathcal{O}_X) \to D_\QCoh(\mathcal{O}_Y)$ 是三角范畴的
正合函子。类似地，(1) 与 (3) 的等价性由 $Lf^*$ 是 $Rf_*$ 的左伴随
这一事实形式地推出。若 $f$ 仿射（引理 \ref{spaces-perfect-lemma-affine-morphism}），
或 $f$ 是开浸入，或 $f$ 是这类态射的复合，则这些条件成立。
\end{remark}

\noindent
下述结果是定理 \ref{spaces-perfect-theorem-bondal-van-den-Bergh} 的一个加强，
并且用完全相同的方法证明。% P09-ERR-0436
设 $T \subset |X|$ 为闭子集，其中 $X$ 是代数空间。
记 $D_T(\mathcal{O}_X)$ 为由上同调层支撑在 $T$ 上的复形所构成的
严格全、饱和三角子范畴。

\begin{lemma}
\label{spaces-perfect-lemma-generator-with-support}
设 $S$ 为概形，$X$ 为 $S$ 上拟紧且拟分离的代数空间。
设 $T \subset |X|$ 为闭子集，且 $|X| \setminus T$ 拟紧。
沿用上述记号，范畴 $D_{\QCoh, T}(\mathcal{O}_X)$ 可由单个完美对象生成。
\end{lemma}

\begin{proof}
我们将使用引理 \ref{spaces-perfect-lemma-induction-principle} 的归纳原理证明此事。
对于位点 $X_{spaces, \etale}$ 中可表、拟紧且拟分离的对象，该性质由
《概形的导出范畴》引理
\ref{perfect-lemma-generator-with-support} 成立。

\medskip\noindent
假设 $(U \subset X, f : V \to X)$ 是初等特异方形，且引理对 $U$
成立，而 $V$ 仿射。为完成证明，需说明结论对 $X$ 成立。
设 $P$ 为 $D(\mathcal{O}_U)$ 中支撑在 $T \cap U$ 上的完美对象，
它是 $D_{\QCoh, T \cap U}(\mathcal{O}_U)$ 的生成元。使用引理
\ref{spaces-perfect-lemma-lift-map-from-perfect-complex-with-support}，可选取
完美对象 $Q$，它属于 $D(\mathcal{O}_X)$ 且支撑在 $T$ 上，其在 $U$ 上的限制
是某个直和，且 $P$ 是该直和的一个直和项。
写 $V = \Spec(B)$。设 $Z = X \setminus U$。则 $f^{-1}Z$ 是 $V$ 的闭子集，
且 $V \setminus f^{-1}Z$ 拟紧。由于 $X$ 拟分离，可得
$f^{-1}Z \cap f^{-1}T = f^{-1}(Z \cap T)$ 是 $V$ 的闭子集，并且
$W = V \setminus f^{-1}(Z \cap T)$ 拟紧。因此可以选取
$g_1, \ldots, g_s \in B$，使得
$f^{-1}(Z \cap T) = V(g_1, \ldots, g_r)$。% P09-ERR-0437
设 $K \in D(\mathcal{O}_V)$ 为与 $g_1, \ldots, g_s$ 在 $B$ 上的
Koszul 复形相对应的完美对象。注意，因为 $K$ 支撑在闭子集
$f^{-1}(Z \cap T) \subset V$ 上，% P09-ERR-0438
推前
$K' = R(V \to X)_*K$ 是 $D(\mathcal{O}_X)$ 中的完美对象，
其在 $V$ 上的限制是 $K$（见引理 \ref{spaces-perfect-lemma-pushforward-perfect} 和
\ref{spaces-perfect-lemma-pushforward-with-support-in-open}）。
我们断言 $Q \oplus K'$ 是 $D_{\QCoh, T}(\mathcal{O}_X)$ 的生成元。

\medskip\noindent
设 $E$ 为 $D_{\QCoh, T}(\mathcal{O}_X)$ 中的对象，使得从
$Q \oplus K'$ 的任意平移到 $E$ 都没有非零映射。
由引理 \ref{spaces-perfect-lemma-pushforward-with-support-in-open}，有
$K' =  R(V \to X)_! K$，因而
$$
\Hom_{D(\mathcal{O}_X)}(K'[n], E) = \Hom_{D(\mathcal{O}_V)}(K[n], E|_V)
$$
于是由《概形的导出范畴》引理
\ref{perfect-lemma-orthogonal-koszul-complex}，有
$E|_V = Rj_*E|_W$，其中 $j : W \to V$ 是包含。图示为
$$
\xymatrix{
W \ar[r]_j & V & Z \cap T \ar[l] \ar[d] \\
V \setminus f^{-1}Z \ar[u]^{j'} \ar[ru]_{j''} & & Z \ar[lu]
}
$$
由于 $E$ 支撑在 $T$ 上，可见 $E|_W$ 支撑在
$f^{-1}T \cap W = f^{-1}T \cap (V \setminus f^{-1}Z)$ 上，
这是 $W$ 中的闭子集。由此得到
$$
E|_V = Rj_*(E|_W) = Rj_*(Rj'_*(E|_{U \cap V})) = Rj''_*(E|_{U \cap V})
$$
这里第二个等号是《上同调》引理
\ref{cohomology-lemma-pushforward-restriction} 的 (1)。该引理可以应用，
因为 $V$ 是概形且 $E$ 具有拟凝聚上同调层，所以沿拟紧开浸入 $j'$
的推前与底层概形上的推前一致；见评注
\ref{spaces-perfect-remark-match-total-direct-images}。
这蕴含 $E = R(U \to X)_*E|_U$（略去一个小细节）。在此情形下
$$
\Hom_{D(\mathcal{O}_X)}(Q[n], E) = \Hom_{D(\mathcal{O}_U)}(Q|_U[n], E|_U)
$$
其中含有 $\Hom_{D(\mathcal{O}_U)}(P[n], E|_U)$ 作为直和项。
因此，根据对 $P$ 的选取，这些群的消失蕴含 $E|_U$ 为零，
从而 $E$ 为零。
\end{proof}

\section{紧对象与完美对象}
\label{spaces-perfect-section-compact}

\noindent
本节是《概形的导出范畴》第 \ref{perfect-section-compact} 节的类似物。

\begin{proposition}
\label{spaces-perfect-proposition-compact-is-perfect}
设 $S$ 为概形。
设 $X$ 为 $S$ 上拟紧且拟分离的代数空间。
$D_\QCoh(\mathcal{O}_X)$ 中的对象是紧的，当且仅当它是完美的。
\end{proposition}

\begin{proof}
若 $K$ 是 $D(\mathcal{O}_X)$ 中的完美对象，其对偶为 $K^\vee$
（《位点上的上同调》引理
\ref{sites-cohomology-lemma-dual-perfect-complex}），则有
$$
\Hom_{D(\mathcal{O}_X)}(K, M) =
H^0(X, K^\vee \otimes_{\mathcal{O}_X}^\mathbf{L} M)
$$
且关于 $M$ 函子性成立。由于 $K^\vee \otimes_{\mathcal{O}_X}^\mathbf{L} -$
与直和可交换，并且由引理
\ref{spaces-perfect-lemma-quasi-coherence-pushforward-direct-sums}，$H^0(X, -)$ 在
$D_\QCoh(\mathcal{O}_X)$ 上与直和可交换，故 $K$ 在
$D_\QCoh(\mathcal{O}_X)$ 中是紧的。

\medskip\noindent
反之，设 $K$ 为 $D_\QCoh(\mathcal{O}_X)$ 中的紧对象。
为证明 $K$ 是完美的，只需证明 $K|_U$ 是完美的，其中 $U$ 是任意
仿射概形并且平展于 $X$；见《位点上的上同调》引理
\ref{sites-cohomology-lemma-perfect-independent-representative}。
注意 $j : U \to X$ 是拟紧且分离的态射。因此
$Rj_* : D_\QCoh(\mathcal{O}_U) \to D_\QCoh(\mathcal{O}_X)$
与直和可交换；见引理
\ref{spaces-perfect-lemma-quasi-coherence-pushforward-direct-sums}。
于是，限制到 $U$ 与 $Rj_*$ 的伴随性蕴含 $K|_U$ 是
$D_\QCoh(\mathcal{O}_U)$ 中的完美对象。因此化归到 $X$ 仿射的情形，
特别地，它是拟紧且拟分离的概形。通过引理
\ref{spaces-perfect-lemma-derived-quasi-coherent-small-etale-site} 和
\ref{spaces-perfect-lemma-descend-perfect}，化归到概形的情形，即《概形的导出范畴》命题
\ref{perfect-proposition-compact-is-perfect}。
\end{proof}

\begin{remark}
\label{spaces-perfect-remark-classical-generator}
设 $S$ 为概形，$X$ 为 $S$ 上拟紧且拟分离的代数空间。
设 $G$ 为 $D(\mathcal{O}_X)$ 中的完美对象，并且是
$D_\QCoh(\mathcal{O}_X)$ 的生成元。由定理
\ref{spaces-perfect-theorem-bondal-van-den-Bergh}，至少存在一个这样的对象。
结合引理 \ref{spaces-perfect-lemma-quasi-coherence-direct-sums}、命题
\ref{spaces-perfect-proposition-compact-is-perfect} 以及《导出范畴》命题
\ref{derived-proposition-generator-versus-classical-generator}，
可见 $G$ 是 $D_{perf}(\mathcal{O}_X)$ 的经典生成元。
\end{remark}

\noindent
下述结果是命题 \ref{spaces-perfect-proposition-compact-is-perfect} 的加强。
设 $T \subset |X|$ 为闭子集，其中 $X$ 是代数空间。
如前，$D_T(\mathcal{O}_X)$ 表示由上同调层支撑在 $T$ 上的复形所构成的
严格全、饱和三角子范畴。由于取直和与取上同调层可交换，
可得 $D_T(\mathcal{O}_X)$ 有直和，并且这些直和等于
$D(\mathcal{O}_X)$ 中的直和。

\begin{lemma}
\label{spaces-perfect-lemma-compact-is-perfect-with-support}
设 $S$ 为概形，$X$ 为 $S$ 上拟紧且拟分离的代数空间。
设 $T \subset |X|$ 为闭子集，且 $|X| \setminus T$ 拟紧。
$D_{\QCoh, T}(\mathcal{O}_X)$ 中的对象是紧的，当且仅当它作为
$D(\mathcal{O}_X)$ 中的对象是完美的。
\end{lemma}

\begin{proof}
由引理前的评注，$D_{\QCoh, T}(\mathcal{O}_X)$ 是有直和的三角范畴。
由命题 \ref{spaces-perfect-proposition-compact-is-perfect}，完美对象给出
$D(\mathcal{O}_X)$ 中的紧对象，% P09-ERR-0440
因而更是任意在取直和下封闭的子范畴中的紧对象。反之，我们将使用如下事实：
存在生成元 $E \in D_{\QCoh, T}(\mathcal{O}_X)$，它是
$\mathcal{O}_X$-模的完美复形；见引理
\ref{spaces-perfect-lemma-generator-with-support}。% P09-ERR-0439
于是由上文，$E$ 是紧的。因此由《导出范畴》命题
\ref{derived-proposition-generator-versus-classical-generator}，
$E$ 是 $D_{\QCoh, T}(\mathcal{O}_X)$ 的紧对象所成全子范畴的经典生成元。
因此，任意紧对象均可由 $E$ 通过有限次下列运算构造出来：
(a) 取平移，(b) 取有限直和，(c) 取锥，(d) 取直和项。
这些运算中的每一种都保持完美性，结论随之成立。
\end{proof}

\begin{remark}
\label{spaces-perfect-remark-classical-generator-with-support}
设 $S$ 为概形，$X$ 为 $S$ 上拟紧且拟分离的代数空间。
设 $T \subset |X|$ 为闭子集，且 $|X| \setminus T$ 拟紧。
设 $G$ 为 $D_{\QCoh, T}(\mathcal{O}_X)$ 中的完美对象，并且是
$D_{\QCoh, T}(\mathcal{O}_X)$ 的生成元。由引理
\ref{spaces-perfect-lemma-generator-with-support}，至少存在一个这样的对象。
结合 $D_{\QCoh, T}(\mathcal{O}_X)$ 有直和这一事实、引理
\ref{spaces-perfect-lemma-compact-is-perfect-with-support} 以及《导出范畴》命题
\ref{derived-proposition-generator-versus-classical-generator}，
可见 $G$ 是 $D_{perf, T}(\mathcal{O}_X)$ 的经典生成元。
\end{remark}

\noindent
下述引理应用了前一引理证明中所用的思想。

\begin{lemma}
\label{spaces-perfect-lemma-map-from-pseudo-coherent-to-complex-with-support}
设 $S$ 为概形，$X$ 为 $S$ 上拟紧且拟分离的代数空间。
设 $T \subset |X|$ 为闭子集，使得补集 $U \subset X$ 拟紧。
设 $\alpha : P \to E$ 为 $D_\QCoh(\mathcal{O}_X)$ 中的态射，且满足
下列条件之一：
\begin{enumerate}
\item $P$ 是完美的，且 $E$ 支撑在 $T$ 上；或
\item $P$ 是拟凝聚的，$E$ 支撑在 $T$ 上，且 $E$ 有下界。
\end{enumerate}
则存在 $\mathcal{O}_X$-模的完美复形 $I$ 以及映射
$I \to \mathcal{O}_X[0]$，使得 $I \otimes^\mathbf{L} P \to E$ 为零，
并且 $I|_U \to \mathcal{O}_U[0]$ 是同构。
\end{lemma}

\begin{proof}
置 $\mathcal{D} = D_{\QCoh, T}(\mathcal{O}_X)$。在两种情形中，复形
$K = R\SheafHom(P, E)$ 都是 $\mathcal{D}$ 中的对象。拟凝聚性见引理
\ref{spaces-perfect-lemma-quasi-coherence-internal-hom}。显然，$K$ 支撑在 $T$ 上，
因为 $R\SheafHom$ 的构造与限制到开集可交换。映射 $\alpha$ 定义了
$H^0(K) = \Hom_{D(\mathcal{O}_X)}(\mathcal{O}_X[0], K)$ 中的元素。
于是，只需对映射 $\alpha : \mathcal{O}_X[0] \to K$ 证明结论。

\medskip\noindent
设 $E \in \mathcal{D}$ 为完美生成元；见引理
\ref{spaces-perfect-lemma-generator-with-support}。如《导出范畴》引理
\ref{derived-lemma-write-as-colimit} 中那样，写成
$$
K = \text{hocolim} K_n
$$
这里使用生成元 $E$。
由于函子 $\mathcal{D} \to D(\mathcal{O}_X)$ 与直和可交换，
可见 $K = \text{hocolim} K_n$ 在 $D(\mathcal{O}_X)$ 中成立。
由于 $\mathcal{O}_X$ 是 $D(\mathcal{O}_X)$ 中的紧对象，得到某个 $n$
以及态射 $\alpha_n : \mathcal{O}_X \to K_n$，它给出 $\alpha$；
见《导出范畴》引理
\ref{derived-lemma-commutes-with-countable-sums}。
对态射 $\mathcal{O}_X[0] \to K_n$（在环境范畴
$D(\mathcal{O}_X)$ 中）应用《导出范畴》引理
\ref{derived-lemma-factor-through}，可见 $\alpha_n$ 分解为
$\mathcal{O}_X[0] \to Q \to K_n$，其中 $Q$ 是
$\langle E \rangle$ 中的对象。由此断定 $Q$ 是支撑在 $T$ 上的完美复形。

\medskip\noindent
选取特异三角
$$
I \to \mathcal{O}_X[0] \to Q \to I[1]
$$
由构造，$I$ 是完美的，映射 $I \to \mathcal{O}_X[0]$ 在 $U$ 上限制为
同构，并且复合 $I \to K$ 为零，因为 $\alpha$ 经 $Q$ 分解。
这证明了引理。
\end{proof}








\section{作为模范畴的导出范畴}
\label{spaces-perfect-section-derived-is-dga}

\noindent
本节是《概形的导出范畴》第
\ref{perfect-section-derived-is-dga} 节的类似结果。

\begin{lemma}
\label{spaces-perfect-lemma-tensor-with-QCoh-complex}
设 $S$ 为概形，$X$ 为 $S$ 上的代数空间。设 $K^\bullet$
为 $\mathcal{O}_X$-模复形，其上同调层拟凝聚。置
$(E, d) = \Hom_{\text{Comp}^{dg}(\mathcal{O}_X)}(K^\bullet, K^\bullet)$
为自同态微分分次代数。则《微分分次代数》引理
\ref{dga-lemma-tensor-with-complex-derived} 中的函子
$$
- \otimes_E^\mathbf{L} K^\bullet :
D(E, \text{d}) \longrightarrow D(\mathcal{O}_X)
$$
的像包含在 $D_\QCoh(\mathcal{O}_X)$ 中。
\end{lemma}

\begin{proof}
设 $P$ 为微分分次 $E$-模，且具有性质 $P$。% P09-ERR-0445
设 $F_\bullet$ 为《微分分次代数》第
\ref{dga-section-P-resolutions} 节中给出的 $P$ 上的滤过。
于是
$$
P \otimes_E K^\bullet = \text{hocolim}\ F_iP \otimes_E K^\bullet
$$
每个 $F_iP$ 都有有限滤过，其分次片是 $E[k]$ 的直和。
结论立即得出。
\end{proof}

\noindent
下述引理可以加强（对于所有仅在某个固定范围内具有非零上同调层的
$L$，其消失性具有一致性）。

\begin{lemma}
\label{spaces-perfect-lemma-ext-from-perfect-into-bounded-QCoh}
设 $S$ 为概形。
设 $X$ 为 $S$ 上拟紧且拟分离的代数空间。
设 $K$ 为 $D(\mathcal{O}_X)$ 中的完美对象。则
\begin{enumerate}
\item 存在整数 $a \leq b$，使得
$\Hom_{D(\mathcal{O}_X)}(K, L) = 0$，只要
$L \in D_\QCoh(\mathcal{O}_X)$ 且有 $H^i(L) = 0$，其中
$i \in [a, b]$；并且
\item 如果 $L$ 有界，则 $\Ext^n_{D(\mathcal{O}_X)}(K, L)$
除有限多个 $n$ 外均为零。
\end{enumerate}
\end{lemma}

\begin{proof}
由于 $\Ext^n_{D(\mathcal{O}_X)}(K, L) =
\Hom_{D(\mathcal{O}_X)}(K, L[n])$，第 (2) 项由第 (1) 项推出。
我们证明第 (1) 项。由于 $K$ 完美，有
$$
\Ext^i_{D(\mathcal{O}_X)}(K, L) =
H^i(X, K^\vee \otimes_{\mathcal{O}_X}^\mathbf{L} L)
$$
其中 $K^\vee$ 是 $K$ 的“对偶”完美复形；见《位点上的上同调》引理
\ref{sites-cohomology-lemma-dual-perfect-complex}。
注意，由引理 \ref{spaces-perfect-lemma-quasi-coherence-tensor-product} 和
\ref{spaces-perfect-lemma-pseudo-coherent}（后者用于说明完美复形具有拟凝聚上同调层），
$P = K^\vee \otimes_{\mathcal{O}_X}^\mathbf{L} L$
属于 $D_\QCoh(X)$。设 $K^\vee$ 的 Tor 振幅属于 $[a, b]$。
于是谱序列
$$
E_1^{p, q} = H^p(K^\vee \otimes_{\mathcal{O}_X}^\mathbf{L} H^q(L))
\Rightarrow
H^{p + q}(K^\vee \otimes_{\mathcal{O}_X}^\mathbf{L} L)
$$
表明 $H^j(K^\vee \otimes_{\mathcal{O}_X}^\mathbf{L} L)$ 为零，
只要 $H^q(L) = 0$ 对于 $q \in [j - b, j - a]$ 成立。
设 $N$ 为《空间的上同调》引理
\ref{spaces-cohomology-lemma-vanishing-quasi-separated}
中的整数 $\max(d_p + p)$。那么，
$H^0(X, K^\vee \otimes_{\mathcal{O}_X}^\mathbf{L} L)$ 为零，
只要上同调层
$$
H^{-N}(K^\vee \otimes_{\mathcal{O}_X}^\mathbf{L} L),
\ H^{-N + 1}(K^\vee \otimes_{\mathcal{O}_X}^\mathbf{L} L),
\ \ldots,
\ H^0(K^\vee \otimes_{\mathcal{O}_X}^\mathbf{L} L)
$$
都为零。事实上，由所引引理和引理
\ref{spaces-perfect-lemma-application-nice-K-injective}，有
$$
H^0(X, K^\vee \otimes_{\mathcal{O}_X}^\mathbf{L} L) =
H^0(X, \tau_{\geq -N}(K^\vee \otimes_{\mathcal{O}_X}^\mathbf{L} L))
$$
而由上同调层的消失性，它等于
$H^0(X, \tau_{\geq 1}(K^\vee \otimes_{\mathcal{O}_X}^\mathbf{L} L))$；
后者由《导出范畴》引理
\ref{derived-lemma-negative-vanishing} 为零。
因此，$\Hom_{D(\mathcal{O}_X)}(K, L)$ 为零，只要
$H^i(L) = 0$ 对于 $i \in [-b - N, -a]$ 成立。
\end{proof}

\noindent
下述定理是《概形的导出范畴》定理
\ref{perfect-theorem-DQCoh-is-Ddga} 的类似结果。

\begin{theorem}
\label{spaces-perfect-theorem-DQCoh-is-Ddga}
设 $S$ 为概形。
设 $X$ 为 $S$ 上拟紧且拟分离的代数空间。
则存在微分分次代数 $(E, \text{d})$，其非零上同调群 $H^i(E)$
只有有限多个，并且 $D_\QCoh(\mathcal{O}_X)$
等价于 $D(E, \text{d})$。
\end{theorem}

\begin{proof}
设 $K^\bullet$ 为 $\mathcal{O}$-模的 K-内射复形，% P09-ERR-0441
它是完美的并生成 $D_\QCoh(\mathcal{O}_X)$。
由定理 \ref{spaces-perfect-theorem-bondal-van-den-Bergh} 和 K-内射消解的存在性，
这样的对象存在。我们将证明，取
$$
(E, \text{d}) = \Hom_{\text{Comp}^{dg}(\mathcal{O}_X)}(K^\bullet, K^\bullet)
$$
时定理成立，其中 $\text{Comp}^{dg}(\mathcal{O}_X)$ 是
$\mathcal{O}$-模复形的微分分次范畴。参见《微分分次代数》第
\ref{dga-section-variant-base-change} 节。由于 $K^\bullet$ 是 K-内射的，
对于所有 $n \in \mathbf{Z}$，有
\begin{equation}
\label{spaces-perfect-equation-E-is-OK}
H^n(E) = \Ext^n_{D(\mathcal{O}_X)}(K^\bullet, K^\bullet)
\end{equation}
由引理 \ref{spaces-perfect-lemma-ext-from-perfect-into-bounded-QCoh}，
这些 Ext 中只有有限多个非零。考虑《微分分次代数》引理
\ref{dga-lemma-tensor-with-complex-derived} 中的函子
$$
- \otimes_E^\mathbf{L} K^\bullet :
D(E, \text{d}) \longrightarrow D(\mathcal{O}_X)
$$
由于 $K^\bullet$ 完美，由命题
\ref{spaces-perfect-proposition-compact-is-perfect}，它给出
$D(\mathcal{O}_X)$ 中的紧对象。% P09-ERR-0442
结合 (\ref{spaces-perfect-equation-E-is-OK})，由《微分分次代数》引理
\ref{dga-lemma-fully-faithful-in-compact-case}，上述函子全忠实。
由《微分分次代数》引理
\ref{dga-lemma-tensor-with-complex-hom-adjoint}，% P09-ERR-0443
它具有右伴随
$$
R\Hom(K^\bullet, - ) : D(\mathcal{O}_X) \longrightarrow D(E, \text{d})
$$
由伴随函子的一般理论，该右伴随是左拟逆函子。另一方面，由引理
\ref{spaces-perfect-lemma-tensor-with-QCoh-complex}，我们得到
$$
- \otimes_E^\mathbf{L} K^\bullet :
D(E, \text{d}) \longrightarrow D_\QCoh(\mathcal{O}_X)
$$
而由我们选取 $K^\bullet$ 为 $D_\QCoh(\mathcal{O}_X)$ 的生成元，
该伴随限制到 $D_\QCoh(\mathcal{O}_X)$ 后的核为零。
形式论证表明，这给出所需等价；见《导出范畴》引理
\ref{derived-lemma-fully-faithful-adjoint-kernel-zero}。
\end{proof}

\begin{remark}[带支撑的变式]
\label{spaces-perfect-remark-DQCoh-is-Ddga-with-support}
设 $S$ 为概形。% P09-ERR-0444
设 $X$ 为拟紧且拟分离的代数空间。
设 $T \subset |X|$ 为闭子集，且 $|X| \setminus T$ 拟紧。
定理 \ref{spaces-perfect-theorem-DQCoh-is-Ddga} 的类似结果对
$D_{\QCoh, T}(\mathcal{O}_X)$ 成立。
这由与该定理证明完全相同的论证得出，其中使用引理
\ref{spaces-perfect-lemma-generator-with-support}、\ref{spaces-perfect-lemma-compact-is-perfect-with-support}
以及引理 \ref{spaces-perfect-lemma-tensor-with-QCoh-complex} 的带支撑变式。
如果以后需要这一结果，我们将在此给出精确陈述和详细证明。
\end{remark}

\begin{remark}[微分分次代数的唯一性]
\label{spaces-perfect-remark-independence-choice}
设 $X$ 为环 $R$ 上拟紧且拟分离的代数空间。
由定理 \ref{spaces-perfect-theorem-DQCoh-is-Ddga} 证明中的构造，
存在微分分次代数 $(A, \text{d})$，它定义在 $R$ 上，使得三角范畴
$D_\QCoh(X)$ $R$-线性等价于 $D(A, \text{d})$。
可以问：$(A, \text{d})$ 在何种程度上唯一？
答案（仅仅）略好于只说 $(A, \text{d})$ 在导出等价意义下确定。
具体地，设 $(B, \text{d})$ 是第二个这样的偶。则有
$$
(A, \text{d}) = \Hom_{\text{Comp}^{dg}(\mathcal{O}_X)}(K^\bullet, K^\bullet)
$$
以及
$$
(B, \text{d}) = \Hom_{\text{Comp}^{dg}(\mathcal{O}_X)}(L^\bullet, L^\bullet)
$$
其中 $K^\bullet$ 和 $L^\bullet$ 是 $\mathcal{O}_X$-模的 K-内射复形，
分别对应于 $D_\QCoh(\mathcal{O}_X)$ 的完美生成元。置
$$
\Omega = \Hom_{\text{Comp}^{dg}(\mathcal{O}_X)}(K^\bullet, L^\bullet)
\quad
\Omega' = \Hom_{\text{Comp}^{dg}(\mathcal{O}_X)}(L^\bullet, K^\bullet)
$$
则 $\Omega$ 是微分分次 $B^{opp} \otimes_R A$-模，
而 $\Omega'$ 是微分分次 $A^{opp} \otimes_R B$-模。
此外，等价
$$
D(A, \text{d}) \to D_\QCoh(\mathcal{O}_X) \to
D(B, \text{d})
$$
由函子 $- \otimes_A^\mathbf{L} \Omega'$ 给出，拟逆也类似。
于是我们处于《微分分次代数》评注
\ref{dga-remark-hochschild-cohomology} 的情形。
如果以后需要这一评注，我们将在此给出精确陈述和详细证明。
\end{remark}











\section{刻画拟凝聚复形，I}
\label{spaces-perfect-section-pseudo-coherent-hocolim}

\noindent
这一材料将在《空间态射进阶》第
\ref{spaces-more-morphisms-section-characterize-pseudo-coherent} 节中继续讨论。
我们可以把拟凝聚对象刻画为完美对象的逼近所成的导出同伦极限。
% P09-ERR-0446

\begin{lemma}
\label{spaces-perfect-lemma-pseudo-coherent-hocolim}
设 $S$ 为概形。
设 $X$ 为 $S$ 上拟紧且拟分离的代数空间。
设 $K \in D(\mathcal{O}_X)$。下列条件等价：
\begin{enumerate}
\item $K$ 是拟凝聚的；并且
\item $K = \text{hocolim} K_n$，其中
$K_n$ 完美，并且映射
$\tau_{\geq -n}K_n \to \tau_{\geq -n}K$ 对所有 $n$ 都是同构。
\end{enumerate}
\end{lemma}

\begin{proof}
蕴含 (2) $\Rightarrow$ (1) 在任意带环位点上都成立。
事实上，假设 (2) 成立。回忆导出范畴中的完美对象是拟凝聚的；见
《位点上的上同调》引理 \ref{sites-cohomology-lemma-perfect}。
于是由定义，$\tau_{\geq -n}K_n$ 是 $(-n + 1)$-拟凝聚的，
从而 $\tau_{\geq -n}K$ 是 $(-n + 1)$-拟凝聚的，
因而 $K$ 是 $(-n + 1)$-拟凝聚的。这对所有 $n$ 都成立，
所以 $K$ 是拟凝聚的；见《位点上的上同调》定义
\ref{sites-cohomology-definition-pseudo-coherent}。

\medskip\noindent
假设 (1) 成立。先选取逼近 $K_1 \to K$；这是
$(X, K, -2)$ 的逼近，并由完美复形 $K_1$ 给出；见定义
\ref{spaces-perfect-definition-approximation-holds}、\ref{spaces-perfect-definition-approximation}
和定理 \ref{spaces-perfect-theorem-approximation}。
归纳地假设已有
$$
K_1 \to K_2 \to \ldots \to K_n \to K
$$
其中 $K_i$ 完美，并且
$\tau_{\geq -i}K_i \to \tau_{\geq -i}K$ 是同构，
对所有 $1 \leq i \leq n$ 都成立。% P09-ERR-0447
于是选取 $a \leq b$，即引理
\ref{spaces-perfect-lemma-ext-from-perfect-into-bounded-QCoh} 中用于完美对象 $K_n$
的整数。
选取逼近 $K_{n + 1} \to K$，即
$(X, K, \min(a - 1, -n - 1))$ 的逼近。选取特异三角
$$
K_{n + 1} \to K \to C \to K_{n + 1}[1]
$$
于是可见 $C \in D_\QCoh(\mathcal{O}_X)$ 满足
$H^i(C) = 0$，对所有 $i \geq a$ 都成立。
因此，由 $a, b$ 的选取可知
$\Hom_{D(\mathcal{O}_X)}(K_n, C) = 0$。
故复合 $K_n \to K \to C$ 为零。因此，由《导出范畴》引理
\ref{derived-lemma-representable-homological}，
我们可以将 $K_n \to K$ 经 $K_{n + 1}$ 分解，从而完成归纳步骤。

\medskip\noindent
还须证明 $K = \text{hocolim} K_n$。
这由把《导出范畴》引理
\ref{derived-lemma-cohomology-of-hocolim} 应用于函子
$H^i( - ) : D(\mathcal{O}_X) \to \textit{Mod}(\mathcal{O}_X)$
以及我们对 $K_n$ 的选取得出。
\end{proof}

\begin{lemma}
\label{spaces-perfect-lemma-pseudo-coherent-hocolim-with-support}
设 $X$ 为拟紧且拟分离的概形。
设 $T \subset X$ 为闭子集，且 $X \setminus T$ 拟紧。
设 $K \in D(\mathcal{O}_X)$ 支撑在 $T$ 上。下列条件等价：
\begin{enumerate}
\item $K$ 是拟凝聚的；并且
\item $K = \text{hocolim} K_n$，其中
$K_n$ 完美、支撑在 $T$ 上，并且
$\tau_{\geq -n}K_n \to \tau_{\geq -n}K$ 对所有 $n$ 都是同构。
\end{enumerate}
\end{lemma}

\begin{proof}
本引理的证明与引理
\ref{spaces-perfect-lemma-pseudo-coherent-hocolim} 的证明完全相同，
唯一差别是在选取逼近时使用三元组 $(T, K, m)$。
\end{proof}











\section{再论凝聚化函子}
\label{spaces-perfect-section-better-coherator}

\noindent
在第 \ref{spaces-perfect-section-coherator} 节中，我们构造并研究了右伴随 $RQ_X$；
它所伴随的典范函子是
$D(\QCoh(\mathcal{O}_X)) \to D(\mathcal{O}_X)$。
它被构造成凝聚化函子
$Q_X : \textit{Mod}(\mathcal{O}_X) \to \QCoh(\mathcal{O}_X)$
的右导出延拓。本节研究包含函子
$$
D_\QCoh(\mathcal{O}_X) \longrightarrow D(\mathcal{O}_X)
$$
何时具有右伴随。如果这个右伴随存在，我们将把它
记作\footnote{这大概不是标准记号。不过，我们已经用 $Q_X$
表示凝聚化函子，并用 $RQ_X$ 表示它的导出延拓。}
$$
DQ_X :
D(\mathcal{O}_X) \longrightarrow D_\QCoh(\mathcal{O}_X)
$$
事实证明，拟紧且拟分离的代数空间具有这样的右伴随。

\begin{lemma}
\label{spaces-perfect-lemma-better-coherator}
设 $S$ 为概形。
设 $X$ 为 $S$ 上拟紧且拟分离的代数空间。
包含函子 $D_\QCoh(\mathcal{O}_X) \to D(\mathcal{O}_X)$ 具有右伴随。
\end{lemma}

\begin{proof}[第一种证明]
我们将使用引理 \ref{spaces-perfect-lemma-induction-principle} 中的归纳原理证明此事。
如果 $D(\QCoh(\mathcal{O}_X)) \to D_\QCoh(\mathcal{O}_X)$ 是等价，
则本引理成立，因为第 \ref{spaces-perfect-section-coherator} 节的函子 $RQ_X$
是函子 $D(\QCoh(\mathcal{O}_X)) \to D(\mathcal{O}_X)$ 的右伴随。
特别地，本引理对仿射代数空间成立；见引理
\ref{spaces-perfect-lemma-affine-coherator}。因此，只须证明如下命题：
如果 $(U \subset X, f : V \to X)$ 是初等特异方形，其中 $U$ 拟紧、
$V$ 仿射，并且本引理对 $U$、$V$ 和 $U \times_X V$ 成立，
那么本引理对 $X$ 成立。

\medskip\noindent
右伴随存在，当且仅当对每个对象 $K$，只要它属于
$D(\mathcal{O}_X)$，
我们都能找到特异三角 % P09-ERR-0448
$$
E' \to E \to K \to E'[1]
$$
它位于 $D(\mathcal{O}_X)$ 中，并且 $E'$ 属于
$D_\QCoh(\mathcal{O}_X)$，同时 $\Hom(M, K) = 0$
对所有 $M$（属于 $D_\QCoh(\mathcal{O}_X)$）成立。见《导出范畴》引理
\ref{derived-lemma-right-adjoint}。考虑引理
\ref{spaces-perfect-lemma-exact-sequence-j-star} 的特异三角
$$
E \to Rj_{U, *}E|_U \oplus Rj_{V, *}E|_V \to
Rj_{U \times_X V, *}E|_{U \times_X V} \to E[1]
$$
它位于 $D(\mathcal{O}_X)$ 中。
由《导出范畴》引理 \ref{derived-lemma-prepare-adjoint}，只须为
$Rj_{U, *}E|_U$、$Rj_{V, *}E|_V$ 和
$Rj_{U \times_X V, *}E|_{U \times_X V}$
构造所需的特异三角。这把我们归约到下一段讨论的命题。

\medskip\noindent
设 $j : U \to X$ 为一个对应于如下情形的 \'etale 态射：
$U$ 拟紧且拟分离，并且本引理对 $U$ 成立。% P09-ERR-0449
设 $L$ 为 $D(\mathcal{O}_U)$ 的对象。
则存在特异三角
$$
E' \to Rj_*L \to K \to E'[1]
$$
它位于 $D(\mathcal{O}_X)$ 中，使得 $E'$ 属于
$D_\QCoh(\mathcal{O}_X)$，并且 $\Hom(M, K) = 0$ 对所有
$M$（属于 $D_\QCoh(\mathcal{O}_X)$）成立。
为说明这一点，选取特异三角
$$
L' \to L \to Q \to L'[1]
$$
它位于 $D(\mathcal{O}_U)$ 中，
使得 $L'$ 属于 $D_\QCoh(\mathcal{O}_U)$，并且
$\Hom(N, Q) = 0$ 对所有 $N$（属于
$D_\QCoh(\mathcal{O}_U)$）成立。
这是可能的，因为《导出范畴》引理
\ref{derived-lemma-right-adjoint} 中的陈述是一个当且仅当命题。
我们得到特异三角
$$
Rj_*L' \to Rj_*L \to Rj_*Q \to Rj_*L'[1]
$$
它位于 $D(\mathcal{O}_X)$ 中。
注意，由引理 \ref{spaces-perfect-lemma-quasi-coherence-direct-image}，
$Rj_*L'$ 属于 $D_\QCoh(\mathcal{O}_X)$。
另一方面，如果 $M$ 属于 $D_\QCoh(\mathcal{O}_X)$，% P09-ERR-0450
则
$$
\Hom(M, Rj_*Q) = \Hom(Lj^*M, Q) = 0
$$
因为由引理 \ref{spaces-perfect-lemma-quasi-coherence-pullback}，
$Lj^*M$ 属于 $D_\QCoh(\mathcal{O}_U)$。
这完成了证明。
\end{proof}

\begin{proof}[第二种证明]
由《导出范畴》命题 \ref{derived-proposition-brown}，该伴随存在。
其假设均满足：首先，注意 $D_\QCoh(\mathcal{O}_X)$ 有直和，
并且取直和与包含函子可交换（引理
\ref{spaces-perfect-lemma-quasi-coherence-direct-sums}）。
另一方面，$D_\QCoh(\mathcal{O}_X)$ 是紧生成的，因为它有完美生成元；
见定理 \ref{spaces-perfect-theorem-bondal-van-den-Bergh}，% P09-ERR-0451
并且由命题 \ref{spaces-perfect-proposition-compact-is-perfect}，完美对象是紧的。
\end{proof}

\begin{lemma}
\label{spaces-perfect-lemma-pushforward-better-coherator}
设 $S$ 为概形。
设 $f : X \to Y$ 为 $S$ 上代数空间的拟紧且拟分离态射。
如果右伴随 $DQ_X$ 和 $DQ_Y$——即包含函子
$D_\QCoh \to D$ 对 $X$ 和 $Y$ 的右伴随——存在，则
$$
Rf_* \circ DQ_X = DQ_Y \circ Rf_*
$$
\end{lemma}

\begin{proof}
该陈述有意义，因为由引理
\ref{spaces-perfect-lemma-quasi-coherence-direct-image}，$Rf_*$
把 $D_\QCoh(\mathcal{O}_X)$ 映入 $D_\QCoh(\mathcal{O}_Y)$。
该陈述成立，是因为 $Lf^*$ 同样把
$D_\QCoh(\mathcal{O}_Y)$ 映入 $D_\QCoh(\mathcal{O}_X)$
（引理 \ref{spaces-perfect-lemma-quasi-coherence-pullback}），
因而 $Rf_* \circ DQ_X$ 和 $DQ_Y \circ Rf_*$
都是 $Lf^* : D_\QCoh(\mathcal{O}_Y) \to D(\mathcal{O}_X)$ 的右伴随。
\end{proof}

\begin{remark}
\label{spaces-perfect-remark-explain-consequence}
设 $S$ 为概形。设 $(U \subset X, f : V \to X)$ 为
$S$ 上代数空间的初等特异方形。
假设 $X$、$U$、$V$ 拟紧且拟分离。
由引理 \ref{spaces-perfect-lemma-better-coherator}，函子
$DQ_X$、$DQ_U$、$DQ_V$、$DQ_{U \times_X V}$ 存在。
此外，有典范特异三角
$$
DQ_X(K) \to Rj_{U, *}DQ_U(K|_U) \oplus Rj_{V, *}DQ_V(K|_V)
\to Rj_{U \times_X V, *}DQ_{U \times_X V}(K|_{U \times_X V}) \to
$$
对任意 $K \in D(\mathcal{O}_X)$ 都成立。
% P09-ERR-0452
这是把正合函子 $DQ_X$ 应用于引理
\ref{spaces-perfect-lemma-exact-sequence-j-star} 的特异三角，
并使用三次引理 \ref{spaces-perfect-lemma-pushforward-better-coherator} 得出的。
\end{remark}

\begin{lemma}
\label{spaces-perfect-lemma-boundedness-better-coherator}
设 $S$ 为概形。
设 $X$ 为 $S$ 上拟紧且拟分离的代数空间。
引理 \ref{spaces-perfect-lemma-better-coherator} 的函子 $DQ_X$
具有如下有界性：存在整数 $N = N(X)$，使得，如果
$K$ 属于 $D(\mathcal{O}_X)$，并且 $H^i(U, K) = 0$，
其中 $U$ 是在 $X$ 上仿射且 \'etale 的，而
$i \not \in [a, b]$，那么上同调层 $H^i(DQ_X(K))$ 为零，
只要 $i \not \in [a, b + N]$。
\end{lemma}

\begin{proof}
我们将使用引理 \ref{spaces-perfect-lemma-induction-principle} 的归纳原理证明此事。

\medskip\noindent
如果 $X$ 仿射，则本引理对 $N = 0$ 成立，因为此时
$RQ_X = DQ_X$ 由取与 $R\Gamma(X, K)$ 相伴的拟凝聚层复形给出。
见引理 \ref{spaces-perfect-lemma-affine-coherator}。

\medskip\noindent
设 $(U \subset W, f : V \to W)$ 为初等特异方形，其中
$W$ 拟紧且拟分离，$U \subset W$ 为拟紧开子空间，$V$ 仿射，
并且本引理对 $U$、$V$ 和 $U \times_W V$ 成立。
设相应整数为 $N(U)$、$N(V)$ 和 $N(U \times_W V)$。
现在假设 $K$ 属于 $D(\mathcal{O}_X)$，并且
$H^i(W, K) = 0$，其中 $W$ 是在 $X$ 上仿射且 \'etale 的，
而 $i \not \in [a, b]$。% P09-ERR-0453
则 $K|_U$、$K|_V$、$K|_{U \times_W V}$ 具有相同性质。
因此可见 $RQ_U(K|_U)$、$RQ_V(K|_V)$ 和 % P09-ERR-0454
$RQ_{U \cap V}(K|_{U \times_W V})$ % P09-ERR-0455
在区间
$[a, b + \max(N(U), N(V), N(U \times_W V))$ % P09-ERR-0456
之外具有消失的上同调层。
由于函子 $Rj_{U, *}$、$Rj_{V, *}$、$Rj_{U \times_W V, *}$
由引理 \ref{spaces-perfect-lemma-quasi-coherence-direct-image} 在
$D_\QCoh$ 上具有有限上同调维数，可见存在 $N$，使得
$Rj_{U, *}DQ_U(K|_U)$、$Rj_{V, *}DQ_V(K|_V)$ 和
$Rj_{U \cap V, *}DQ_{U \times_W V}(K|_{U \times_W V})$
在区间 $[a, b + N]$ 之外具有消失的上同调层。
最后，由评注 \ref{spaces-perfect-remark-explain-consequence} 的特异三角得出结论。
\end{proof}

\begin{example}
\label{spaces-perfect-example-inverse-limit-quasi-coherent}
设 $S$ 为概形。
设 $X$ 为 $S$ 上拟紧且拟分离的代数空间。
设 $(\mathcal{F}_n)$ 为 $X$ 上拟凝聚层的逆系统。
由于 $DQ_X$ 是右伴随，它与积可交换，因而与导出极限可交换。
因此
$$
DQ_X(R\lim \mathcal{F}_n) =
(R\lim\text{ 于 }D_\QCoh(\mathcal{O}_X))(\mathcal{F}_n)
$$
其中第一个 $R\lim$ 取于 $D(\mathcal{O}_X)$ 中。
事实上，为此记 $K = R\lim \mathcal{F}_n$。
对任意 $U$，只要它在 $X$ 上仿射且 \'etale，就有
$$
H^i(U, K) =
H^i(R\Gamma(U, R\lim \mathcal{F}_n)) =
H^i(R\lim R\Gamma(U, \mathcal{F}_n)) =
H^i(R\lim \Gamma(U, \mathcal{F}_n))
$$
因为上同调与导出极限可交换，并且拟凝聚层
$\mathcal{F}_n$ 在仿射对象上没有高阶上同调。
由阿贝尔群范畴中对 $R\lim$ 的计算可见，
$H^i(U, K) = 0$，除非 $i \in [0, 1]$。
最后断定，$R\lim$（取于 $D_\QCoh(\mathcal{O}_X)$ 中），
即上文的 $DQ_X(K)$，属于 $D^b_\QCoh(\mathcal{O}_X)$，
并且由引理 \ref{spaces-perfect-lemma-boundedness-better-coherator}，
其负次数上同调层为零。
\end{example}










\section{上同调与基变换，IV}
\label{spaces-perfect-section-cohomology-and-base-change-perfect}

\noindent
本节是《概形的导出范畴》第
\ref{perfect-section-cohomology-and-base-change-perfect} 节的类似结果。

\begin{lemma}
\label{spaces-perfect-lemma-cohomology-base-change}
设 $S$ 为概形。设 $f : X \to Y$ 为 $S$ 上代数空间的
拟紧且拟分离态射。对 $E$ 属于
$D_\QCoh(\mathcal{O}_X)$ 且 $K$ 属于
$D_\QCoh(\mathcal{O}_Y)$，有
$$
Rf_*(E) \otimes_{\mathcal{O}_Y}^\mathbf{L} K =
Rf_*(E \otimes_{\mathcal{O}_X}^\mathbf{L} Lf^*K)
$$
\end{lemma}

\begin{proof}
不加任何假设时，存在映射
$Rf_*(E) \otimes_{\mathcal{O}_Y}^\mathbf{L} K \to
Rf_*(E \otimes_{\mathcal{O}_X}^\mathbf{L} Lf^*K)$。
具体地，它是典范映射
$$
Lf^*(Rf_*(E) \otimes_{\mathcal{O}_Y}^\mathbf{L} K) =
Lf^*(Rf_*(E)) \otimes_{\mathcal{O}_X}^\mathbf{L} Lf^*K
\longrightarrow
E \otimes_{\mathcal{O}_X}^\mathbf{L} Lf^*K
$$
的伴随映射，而该映射来自 $Lf^*Rf_*E \to E$。见《位点上的上同调》引理
\ref{sites-cohomology-lemma-pullback-tensor-product} 和
\ref{sites-cohomology-lemma-adjoint}。
为检验它是同构，可以在 $Y$ 上 \'etale 局部地工作。
因此归约到 $Y$ 为仿射概形的情形。

\medskip\noindent
假设 $K = \bigoplus K_i$ 是若干复形
$K_i \in D_\QCoh(\mathcal{O}_Y)$ 的直和。
如果该陈述对每个 $K_i$ 成立，那么它对 $K$ 成立。
事实上，函子 $Lf^*$ 和 $\otimes^\mathbf{L}$ 由构造保持直和，
而由引理 \ref{spaces-perfect-lemma-quasi-coherence-pushforward-direct-sums}，
$Rf_*$ 与直和可交换（对具有拟凝聚上同调层的复形）。
此外，假设 $K \to L \to M \to K[1]$ 是
$D_\QCoh(Y)$ 中的特异三角。那么，如果本引理的陈述对
$K, L, M$ 中两个成立，它就对第三个成立
（因为所涉函子都是三角范畴的正合函子）。

\medskip\noindent
假设 $Y$ 仿射，记 $Y = \Spec(A)$。函子
$\widetilde{\ } : D(A) \to D_\QCoh(\mathcal{O}_Y)$ 是等价；见引理
\ref{spaces-perfect-lemma-derived-quasi-coherent-small-etale-site} 和
《概形的导出范畴》引理
\ref{perfect-lemma-affine-compare-bounded}。
令 $T$ 表示关于 $K \in D(A)$ 的如下性质：
本引理的陈述对 $\widetilde{K}$ 成立。
上面的讨论和《代数进阶》评注
\ref{more-algebra-remark-P-resolution} 表明，% P09-ERR-0457
只须证明 $T$ 对 $A[k]$ 成立。这完成了证明，因为本引理的陈述
对结构层的平移显然成立。
\end{proof}

\begin{definition}
\label{spaces-perfect-definition-tor-independent}
设 $S$ 为概形，$B$ 为 $S$ 上的代数空间。
设 $X$、$Y$ 为 $B$ 上的代数空间。称 $X$ 和
$Y$ 在 $B$ 上 {\it Tor 独立}，当且仅当对每个交换图
$$
\xymatrix{
\Spec(k) \ar[d]_{\overline{y}} \ar[dr]_{\overline{b}} \ar[r]_-{\overline{x}} &
X \ar[d] \\
Y \ar[r] & B
}
$$
的几何点，环 $\mathcal{O}_{X, \overline{x}}$ 和
$\mathcal{O}_{Y, \overline{y}}$ 在
$\mathcal{O}_{B, \overline{b}}$ 上 Tor 独立（见《代数进阶》定义
\ref{more-algebra-definition-tor-independent}）。
\end{definition}

\noindent
下述引理特别表明：在代数空间可表的情形中，本定义与我们已有的定义一致。

\begin{lemma}
\label{spaces-perfect-lemma-tor-independent}
设 $S$ 为概形，$B$ 为 $S$ 上的代数空间。
设 $X$、$Y$ 为 $B$ 上的代数空间。下列条件等价：
\begin{enumerate}
\item $X$ 和 $Y$ 在 $B$ 上 Tor 独立；
\item 对每个交换图
$$
\xymatrix{
U \ar[d] \ar[r] & W \ar[d] & V \ar[d] \ar[l] \\
X \ar[r] & B & Y \ar[l]
}
$$
若竖直箭头为 \'etale，则 $U$ 和 $V$ 在 $W$ 上 Tor 独立；% P09-ERR-0458
\item 对某个如 (2) 中的交换图，满足 (a) $W \to B$ 为
\'etale 满射，(b) $U \to X \times_B W$ 为 \'etale 满射，
(c) $V \to Y \times_B W$ 为 \'etale 满射，空间 $U$ 和 $V$
在 $W$ 上 Tor 独立；并且
\item 对某个如 (3) 中的交换图，其中 $U$、$V$、$W$ 是概形，
概形 $U$ 和 $V$ 在 $W$ 上 Tor 独立，其意义见《概形的导出范畴》定义
\ref{perfect-definition-tor-independent}。
\end{enumerate}
\end{lemma}

\begin{proof}
对代数空间的 \'etale 态射 $\varphi : U \to X$ 以及几何点
$\overline{u}$，局部环映射
$\mathcal{O}_{X, \varphi(\overline{u})} \to
\mathcal{O}_{U, \overline{u}}$ 是同构。
因此 (1) 与 (2) 的等价性成立，蕴含 (1) $\Rightarrow$ (3)
也成立。假设 (3)，并选取一个如定义
\ref{spaces-perfect-definition-tor-independent} 中的几何点图。
这些假设表明，我们可以先把 $\overline{b}$ 提升到几何点
$\overline{w}$（属于 $W$），再把几何点
$(\overline{x}, \overline{b})$ 提升到几何点
$\overline{u}$（属于 $U$），最后把几何点
$(\overline{y}, \overline{b})$ 提升到几何点
$\overline{v}$（属于 $V$）。使用《空间的性质》引理
\ref{spaces-properties-lemma-geometric-lift-to-usual} 找到这些提升。
利用上面关于局部环的评注，可断定给定的图满足定义中的条件。

\medskip\noindent
作出这些初步选择后，显然，(4) 归结为如下陈述：
当 $X$、$Y$、$B$ 是概形时，定义
\ref{spaces-perfect-definition-tor-independent} 与《概形的导出范畴》定义
\ref{perfect-definition-tor-independent} 一致。

\medskip\noindent
设 $\overline{x}, \overline{b}, \overline{y}$ 如定义
\ref{spaces-perfect-definition-tor-independent} 中那样，位于点
$x, y, b$ 之上。回忆
$\mathcal{O}_{X, \overline{x}} = \mathcal{O}_{X, x}^{sh}$
（《空间的性质》引理
\ref{spaces-properties-lemma-describe-etale-local-ring}），其余两个也类似。
由《代数》引理
\ref{algebra-lemma-strictly-henselian-functorial-improve} 可见，
$\mathcal{O}_{X, \overline{x}}$ 是
$\mathcal{O}_{X, x} \otimes_{\mathcal{O}_{B, b}}
\mathcal{O}_{B, \overline{b}}$ 的严格亨泽尔化。
特别地，环映射
$$
\mathcal{O}_{X, x} \otimes_{\mathcal{O}_{B, b}} \mathcal{O}_{B, \overline{b}}
\longrightarrow
\mathcal{O}_{X, \overline{x}}
$$
是平坦的（《代数进阶》引理
\ref{more-algebra-lemma-dumb-properties-henselization}）。
由《代数进阶》引理
\ref{more-algebra-lemma-tor-independent-flat} 可见
$$
\text{Tor}_i^{\mathcal{O}_{B, b}}(\mathcal{O}_{X, x}, \mathcal{O}_{Y, y})
\otimes_{\mathcal{O}_{X, x} \otimes_{\mathcal{O}_{B, b}} \mathcal{O}_{Y, y}}
(\mathcal{O}_{X, \overline{x}} \otimes_{\mathcal{O}_{B, \overline{b}}}
\mathcal{O}_{Y, \overline y})
=
\text{Tor}_i^{\mathcal{O}_{B, \overline{b}}}(
\mathcal{O}_{X, \overline{x}}, \mathcal{O}_{Y, \overline{y}})
$$
因此，如果作为概形的 $X$ 和 $Y$ 在 $B$ 上 Tor 独立，
那么作为代数空间的 $X$ 和 $Y$ 在 $B$ 上 Tor 独立。

\medskip\noindent
反之，可以假设 $X$、$Y$、$B$ 仿射。注意，环映射
$$
\mathcal{O}_{X, x} \otimes_{\mathcal{O}_{B, b}} \mathcal{O}_{Y, y}
\longrightarrow
\mathcal{O}_{X, \overline{x}} \otimes_{\mathcal{O}_{B, \overline{b}}}
\mathcal{O}_{Y, \overline y}
$$
由上面的观察是平坦的。此外，该映射在谱上的像包含所有素理想
$$
\mathfrak s \subset
\mathcal{O}_{X, x} \otimes_{\mathcal{O}_{B, b}} \mathcal{O}_{Y, y}
$$
只要它们位于 $\mathfrak m_x$ 和 $\mathfrak m_y$ 之上。
因此，由此及上面显示的 Tor 公式可见，如果作为代数空间的
$X$ 和 $Y$ 在 $B$ 上 Tor 独立，那么
$$
\text{Tor}_i^{\mathcal{O}_{B, b}}
(\mathcal{O}_{X, x}, \mathcal{O}_{Y, y})_\mathfrak s = 0
$$
对所有 $i > 0$ 以及上面所有的 $\mathfrak s$ 成立。
把《代数进阶》引理
\ref{more-algebra-lemma-tor-independent} 应用于环映射
$\Gamma(B, \mathcal{O}_B) \to \Gamma(X, \mathcal{O}_X)$
以及
$\Gamma(B, \mathcal{O}_B) \to \Gamma(X, \mathcal{O}_X)$，% P09-ERR-0459
可知这蕴含 $X$ 和 $Y$ 在 $B$ 上 Tor 独立。
\end{proof}

\begin{lemma}
\label{spaces-perfect-lemma-compare-base-change}
设 $S$ 为概形。设 $g : Y' \to Y$ 为 $S$ 上代数空间的态射。
设 $f : X \to Y$ 为 $S$ 上代数空间的拟紧且拟分离态射。
考虑基变换图
$$
\xymatrix{
X' \ar[r]_{g'} \ar[d]_{f'} &
X \ar[d]^f \\
Y' \ar[r]^g &
Y
}
$$
如果 $X$ 和 $Y'$ 在 $Y$ 上 Tor 独立，那么对所有
$E \in D_\QCoh(\mathcal{O}_X)$，有
$Rf'_*L(g')^*E = Lg^*Rf_*E$。
\end{lemma}

\begin{proof}
对任意对象 $E$（属于 $D(\mathcal{O}_X)$），可以使用
《位点上的上同调》评注
\ref{sites-cohomology-remark-base-change} 得到典范基变换映射
$Lg^*Rf_*E \to Rf'_*L(g')^*E$。为检验它是同构，可以在
$Y'$ 上 \'etale 局部地工作。因此可以假设
$g : Y' \to Y$ 是仿射概形的态射。特别地，$g$ 是仿射的，
并且只须证明
$$
Rg_*Lg^*Rf_*E \to Rg_*Rf'_*L(g')^*E = Rf_*(Rg'_* L(g')^* E)
$$
是同构；见引理 \ref{spaces-perfect-lemma-affine-morphism}
（并使用引理 \ref{spaces-perfect-lemma-quasi-coherence-pullback}、
\ref{spaces-perfect-lemma-quasi-coherence-tensor-product} 和
\ref{spaces-perfect-lemma-quasi-coherence-direct-image} 来说明对象
$Rf'_*L(g')^*E$ 和 $Lg^*Rf_*E$ 具有拟凝聚上同调层）。
注意，$g'$ 也是仿射的（《空间的态射》引理
\ref{spaces-morphisms-lemma-base-change-affine}）。
由引理 \ref{spaces-perfect-lemma-affine-morphism-pull-push}，该映射变为映射
$$
Rf_*E \otimes_{\mathcal{O}_Y}^\mathbf{L} g_*\mathcal{O}_{Y'}
\longrightarrow
Rf_*(E \otimes_{\mathcal{O}_X}^\mathbf{L} g'_*\mathcal{O}_{X'})
$$
注意 $g'_*\mathcal{O}_{X'} = f^*g_*\mathcal{O}_{Y'}$。
因此由引理 \ref{spaces-perfect-lemma-cohomology-base-change}，只须证明
$Lf^*g_*\mathcal{O}_{Y'} = f^*g_*\mathcal{O}_{Y'}$。
这由假设 $X$ 和 $Y'$ 在 $Y$ 上 Tor 独立得出。
事实上，为检验这一点，可以在 $X$ 上 \'etale 局部地工作，
因此还可以假设 $X$ 仿射。记
$X = \Spec(A)$、$Y = \Spec(R)$ 且 $Y' = \Spec(R')$。
我们的假设蕴含 $A$ 和 $R'$ 在 $R$ 上 Tor 独立（见引理
\ref{spaces-perfect-lemma-tor-independent} 和《代数进阶》引理
\ref{more-algebra-lemma-tor-independent}），即
$\text{Tor}_i^R(A, R') = 0$ 对 $i > 0$ 成立。换言之，
$A \otimes_R^\mathbf{L} R' = A \otimes_R R'$，
这恰好意味着
$Lf^*g_*\mathcal{O}_{Y'} = f^*g_*\mathcal{O}_{Y'}$。
\end{proof}

\noindent
下述引理将在对偶化复形一章中使用。

\begin{lemma}
\label{spaces-perfect-lemma-affine-morphism-and-hom-out-of-perfect}
设 $g : S' \to S$ 为仿射概形的态射。考虑笛卡尔方形
$$
\xymatrix{
X' \ar[r]_{g'} \ar[d]_{f'} & X \ar[d]^f \\
S' \ar[r]^g & S
}
$$
其中各代数空间拟紧且拟分离。假设 $g$ 和 $f$ Tor 独立。
写作 $S = \Spec(R)$ 且 $S' = \Spec(R')$。对
$M, K \in D(\mathcal{O}_X)$，典范映射
$$
R\Hom_X(M, K) \otimes^\mathbf{L}_R R'
\longrightarrow
R\Hom_{X'}(L(g')^*M, L(g')^*K)
$$
在 $D(R')$ 中，于下列两种情形为同构：
\begin{enumerate}
\item $M \in D(\mathcal{O}_X)$ 完美，且 $K \in D_\QCoh(X)$；或
\item $M \in D(\mathcal{O}_X)$ 拟凝聚，
$K \in D_\QCoh^+(X)$，且 $R'$ 在 $R$ 上具有有限 Tor 维数。
\end{enumerate}
\end{lemma}

\begin{proof}
存在典范映射
$R\Hom_X(M, K) \to R\Hom_{X'}(L(g')^*M, L(g')^*K)$，
它位于全局 Hom 复形的 $D(\Gamma(X, \mathcal{O}_X))$ 中；见
《位点上的上同调》第
\ref{sites-cohomology-section-global-RHom} 节。
限制标量后，可以把它看作 $D(R)$ 中的映射。
然后使用限制标量与 $- \otimes_R^\mathbf{L} R'$ 的伴随性，
得到本引理中显示的映射。定义了该映射后，只须证明它在
阿贝尔群的导出范畴中是同构。

\medskip\noindent
右端等于
$$
R\Hom_X(M, R(g')_*L(g')^*K) =
R\Hom_X(M, K \otimes_{\mathcal{O}_X}^\mathbf{L} g'_*\mathcal{O}_{X'})
$$
这是引理 \ref{spaces-perfect-lemma-affine-morphism-pull-push} 的结论。
在两种情形中，由引理 \ref{spaces-perfect-lemma-quasi-coherence-internal-hom}，
复形 $R\SheafHom(M, K)$ 都是
$D_\QCoh(\mathcal{O}_X)$ 的对象。存在自然映射
$$
R\SheafHom(M, K) \otimes_{\mathcal{O}_X}^\mathbf{L} g'_*\mathcal{O}_{X'}
\longrightarrow
R\SheafHom(M, K \otimes_{\mathcal{O}_X}^\mathbf{L} g'_*\mathcal{O}_{X'})
$$
它在两种情形中都是同构；引理
\ref{spaces-perfect-lemma-internal-hom-evaluate-tensor-isomorphism}。% P09-ERR-0461
为说明该引理适用于情形 (2)，注意
$g'_*\mathcal{O}_{X'} = Rg'_*\mathcal{O}_{X'} =
Lf^*g_*\mathcal{O}_X$，其中第二个等号由引理
\ref{spaces-perfect-lemma-compare-base-change} 给出。% P09-ERR-0462
使用《概形的导出范畴》引理
\ref{perfect-lemma-tor-dimension-affine}、引理
\ref{spaces-perfect-lemma-descend-tor-amplitude} 和《位点上的上同调》引理
\ref{sites-cohomology-lemma-tor-amplitude-pullback}，可断定
$g'_*\mathcal{O}_{X'}$ 具有有限 Tor 维数。
因此，在两种情形中，把 $K$ 替换为 $R\SheafHom(M, K)$，
我们都归约到证明
$$
R\Gamma(X, K) \otimes^\mathbf{L}_A A' \longrightarrow
R\Gamma(X, K \otimes^\mathbf{L}_{\mathcal{O}_X} g'_*\mathcal{O}_{X'})
$$
是同构。% P09-ERR-0463
注意，由引理 \ref{spaces-perfect-lemma-affine-morphism-pull-push}，左端等于
$R\Gamma(X', L(g')^*K)$。因此结论由引理
\ref{spaces-perfect-lemma-compare-base-change} 得出。
\end{proof}

\begin{remark}
\label{spaces-perfect-remark-multiplication-map}
沿用引理 \ref{spaces-perfect-lemma-affine-morphism-and-hom-out-of-perfect} 的记号。% P09-ERR-0460
下图交换：
$$
\xymatrix{
R\Hom_X(M, Rg'_*L) \otimes_R^\mathbf{L} R' \ar[r] \ar[d]_\mu &
R\Hom_{X'}(L(g')^*M, L(g')^*Rg'_*L) \ar[d]^a \\
R\Hom_X(M, R(g')_*L) \ar@{=}[r] &
R\Hom_{X'}(L(g')^*M, L)
}
$$
其中上方水平箭头是该引理中的映射，$\mu$ 是乘法映射，
而 $a$ 来自伴随映射 $L(g')^*Rg'_*L \to L$。
乘法映射是伴随映射
$K' \otimes_R^\mathbf{L} R' \to K'$，其中
$K' \in D(R')$ 任意。
\end{remark}

\begin{lemma}
\label{spaces-perfect-lemma-tor-independence-and-tor-amplitude}
设 $S$ 为概形。考虑代数空间的笛卡尔方形
$$
\xymatrix{
X' \ar[r]_{g'} \ar[d]_{f'} & X \ar[d]^f \\
Y' \ar[r]^g & Y
}
$$
它位于 $S$ 上。
假设 $g$ 和 $f$ Tor 独立。
\begin{enumerate}
\item 如果 $E \in D(\mathcal{O}_X)$ 具有 $[a, b]$ 中的
Tor 振幅（把它视为 $f^{-1}\mathcal{O}_Y$-模复形），
那么 $L(g')^*E$ 具有 $[a, b]$ 中的 Tor 振幅
（把它视为 $f^{-1}\mathcal{O}_{Y'}$-模复形）。% P09-ERR-0464
\item 如果 $\mathcal{G}$ 是 $\mathcal{O}_X$-模且在
$Y$ 上平坦，那么
$L(g')^*\mathcal{G} = (g')^*\mathcal{G}$。
\end{enumerate}
\end{lemma}

\begin{proof}
可以在茎上计算 Tor 维数；见《位点上的上同调》引理
\ref{sites-cohomology-lemma-tor-amplitude-stalk} 和《空间的性质》定理
\ref{spaces-properties-theorem-exactness-stalks}。
如果 $\overline{x}'$ 是 $X'$ 的几何点，且 $\overline{x}$
是它在 $X$ 中的像，那么
$$
(L(g')^*E)_{\overline{x}'} =
E_{\overline{x}}
\otimes_{\mathcal{O}_{X, \overline{x}}}^\mathbf{L}
\mathcal{O}_{X', \overline{x}'}
$$
设 $\overline{y}'$（属于 $Y'$）和 $\overline{y}$（属于 $Y$）
分别是 $\overline{x}'$ 和 $\overline{x}$ 的像。
由于 $X$ 和 $Y'$ 在 $Y$ 上 Tor 独立，可以应用《代数进阶》引理
\ref{more-algebra-lemma-base-change-comparison}，看出所显示公式的右端等于
$E_{\overline{x}}
\otimes_{\mathcal{O}_{Y, \overline{y}}}^\mathbf{L}
\mathcal{O}_{Y', \overline{y}'}$
（在 $D(\mathcal{O}_{Y', \overline{y}'})$ 中）。
因此 (1) 由《代数进阶》引理
\ref{more-algebra-lemma-pull-tor-amplitude} 得出。
为证明 (2)，注意 $\mathcal{G}$ 的平坦性等价于
$\mathcal{G}[0]$ 具有 $[0, 0]$ 中的 Tor 振幅。
应用 (1) 即得结论。
\end{proof}










\section{上同调与基变换，V}
\label{spaces-perfect-section-cohomology-base-change}

\noindent
本节是《概形的导出范畴》第
\ref{perfect-section-cohomology-base-change} 节的类似结果。
在第 \ref{spaces-perfect-section-cohomology-and-base-change-perfect} 节中，
我们看到，当态射 Tor 独立时，基变换定理成立。
即使在仿射情形中，若无这种条件，也不可能有基变换定理；见
《代数进阶》第 \ref{more-algebra-section-tor-independence} 节。
本节分析何时可以“逐个复形”得到基变换结果。

\medskip\noindent
为此，设 $S$ 为基概形，并假设我们有交换图
$$
\xymatrix{
X' \ar[r]_{g'} \ar[d]_{f'} &
X \ar[d]^f \\
Y' \ar[r]^g &
Y
}
$$
它由 $S$ 上的代数空间组成（通常我们会假设它是笛卡尔图）。
设 $K \in D_\QCoh(\mathcal{O}_X)$，并设
$L(g')^*K \to K'$ 是 $D_\QCoh(\mathcal{O}_{X'})$ 中的映射。
对几何点 $\overline{x}'$（属于 $X'$），考虑几何点
$\overline{x} = g'(\overline{x}')$、
$\overline{y}' = f'(\overline{x}')$、
$\overline{y} = f(\overline{x}) = g(\overline{y}')$，
它们分别属于 $X$、$Y'$、$Y$。
于是可以考虑映射
$$
K_{\overline{x}}
\otimes_{\mathcal{O}_{Y, \overline{y}}}^\mathbf{L}
\mathcal{O}_{Y', \overline{y}'}
\to
K_{\overline{x}}
\otimes_{\mathcal{O}_{X, \overline{x}}}^\mathbf{L}
\mathcal{O}_{X', \overline{x}'} \to
K'_{\overline{x}'}
$$
其中第一个箭头是《代数进阶》方程
(\ref{more-algebra-equation-comparison-map})，第二个箭头来自
$$
(L(g')^*K)_{\overline{x}'} =
K_{\overline{x}} \otimes_{\mathcal{O}_{X, \overline{x}}}^\mathbf{L}
\mathcal{O}_{X', \overline{x}'}
$$
以及给定映射 $L(g')^*K \to K'$。
对每个 $i \in \mathbf{Z}$，得到一个
$\mathcal{O}_{X, \overline{x}} \otimes_{\mathcal{O}_{Y, \overline{y}}}
\mathcal{O}_{Y', \overline{y}'}$-模结构，其承载模为
$H^i(K_{\overline{x}}
\otimes_{\mathcal{O}_{Y, \overline{y}}}^\mathbf{L}
\mathcal{O}_{Y', \overline{y}'})$。
把所有这些放在一起，得到典范映射
\begin{equation}
\label{spaces-perfect-equation-bc}
H^i(K_{\overline{x}}
\otimes_{\mathcal{O}_{Y, \overline{y}}}^\mathbf{L}
\mathcal{O}_{Y', \overline{y}'})
\otimes_{(\mathcal{O}_{X, \overline{x}}
\otimes_{\mathcal{O}_{Y, \overline{y}}}
\mathcal{O}_{Y', \overline{y}'})}
\mathcal{O}_{X', \overline{x}'}
\longrightarrow
H^i(K'_{\overline{x}'})
\end{equation}
它是 $\mathcal{O}_{X', \overline{x}'}$-模的映射。

\begin{lemma}
\label{spaces-perfect-lemma-single-complex-base-change-condition}
设 $S$ 为概形。设
$$
\xymatrix{
X' \ar[r]_{g'} \ar[d]_{f'} &
X \ar[d]^f \\
Y' \ar[r]^g &
Y
}
$$
为 $S$ 上代数空间的笛卡尔图。
设 $K \in D_\QCoh(\mathcal{O}_X)$，并设 $L(g')^*K \to K'$
是 $D_\QCoh(\mathcal{O}_{X'})$ 中的映射。下列条件等价：
\begin{enumerate}
\item 对任意 $x' \in X'$ 和 $i \in \mathbf{Z}$，% P09-ERR-0465
映射 (\ref{spaces-perfect-equation-bc}) 是同构；
\item 对任意交换图
$$
\xymatrix{
& U \ar[d] \ar[rd]^a \\
V' \ar[r] \ar[rd]^c & V \ar[rd]^b & X \ar[d]^f \\
& Y' \ar[r]^g & Y
}
$$
若 $a, b, c$ 为 \'etale，$U, V, V'$ 为概形，并且
$U' = V' \times_V U$，则《概形的导出范畴》引理
\ref{perfect-lemma-single-complex-base-change-condition} 的等价条件
对 $(U \to X)^*K$ 和 $(U' \to X')^*K'$ 成立；并且
\item 存在一个如 (2) 中的图，使得 $U' \to X'$ 满射。
\end{enumerate}
\end{lemma}

\begin{proof}
注意，(1) 在 $X'$ 上是 \'etale 局部的。
逐一考察已知条件之间的形式蕴含，可见只须证明：当
$X, Y, Y', X'$ 分别由概形 $X_0, Y_0, Y'_0, X'_0$ 表示时，
本引理的条件 (1) 等价于《概形的导出范畴》引理
\ref{perfect-lemma-single-complex-base-change-condition} 的条件 (1)。
用 $f_0, g_0, g'_0, f'_0$ 表示这些概形之间分别对应于
$f, g, g', f'$ 的态射。% P09-ERR-0466
可以假设 $K = \epsilon^*K_0$ 且 $K' = \epsilon^*K'_0$，
其中 $K_0 \in D_\QCoh(\mathcal{O}_{X_0})$ 且
$K'_0 \in D_\QCoh(\mathcal{O}_{X'_0})$；见引理
\ref{spaces-perfect-lemma-derived-quasi-coherent-small-etale-site}。
此外，映射 $Lg^*K \to K'$ 是映射 % P09-ERR-0467
$L(g_0)^*K_0 \to K'_0$ 的拉回；记号见评注
\ref{spaces-perfect-remark-match-total-direct-images}。
回忆 $\mathcal{O}_{X, \overline{x}}$ 是
$\mathcal{O}_{X, x}$ 的严格亨泽尔化
（《空间的性质》引理
\ref{spaces-properties-lemma-describe-etale-local-ring}），并且
$$
K_{\overline{x}} =
K_{0, x} \otimes_{\mathcal{O}_{X, x}}^\mathbf{L} \mathcal{O}_{X, \overline{x}}
\quad\text{且}\quad
K'_{\overline{x}'} =
K'_{0, x'} \otimes_{\mathcal{O}_{X', x'}}^\mathbf{L}
\mathcal{O}_{X', \overline{x}'}
$$
（类似于《空间的性质》引理
\ref{spaces-properties-lemma-stalk-quasi-coherent}）。
考虑交换图
$$
\xymatrix{
H^i(K_{\overline{x}}
\otimes_{\mathcal{O}_{Y, \overline{y}}}^\mathbf{L}
\mathcal{O}_{Y', \overline{y}'})
\otimes_{(\mathcal{O}_{X, \overline{x}}
\otimes_{\mathcal{O}_{Y, \overline{y}}}
\mathcal{O}_{Y', \overline{y}'})}
\mathcal{O}_{X', \overline{x}'}
\ar[r] &
H^i(K'_{\overline{x}'}) \\
H^i(K_{0, x} \otimes_{\mathcal{O}_{Y, y}}^\mathbf{L} \mathcal{O}_{Y', y'})
\otimes_{(\mathcal{O}_{X, x} \otimes_{\mathcal{O}_{Y, y}} \mathcal{O}_{Y', y'})}
\mathcal{O}_{X', x'}
\ar[u] \ar[r] &
H^i(K'_{0, x'}) \ar[u]
}
$$
必须证明，下方水平箭头是同构当且仅当上方水平箭头是同构。
由于 $\mathcal{O}_{X', x'} \to \mathcal{O}_{X', \overline{x}'}$
是忠实平坦的（《代数进阶》引理
\ref{more-algebra-lemma-dumb-properties-henselization}），
只须证明上方箭头是下方箭头沿此映射的基变换。
对于右侧对象，这立即由上面给出的茎之间的关系得出。
对于左侧对象，只须证明
\begin{align*}
&
H^i\left(
(K_{0, x}
\otimes_{\mathcal{O}_{X, x}}^\mathbf{L} \mathcal{O}_{X, \overline{x}})
\otimes_{\mathcal{O}_{Y, \overline{y}}}^\mathbf{L}
\mathcal{O}_{Y', \overline{y}'}\right) \\
& =
H^i(K_{0, x} \otimes_{\mathcal{O}_{Y, y}}^\mathbf{L} \mathcal{O}_{Y', y'})
\otimes_{(\mathcal{O}_{X, x} \otimes_{\mathcal{O}_{Y, y}} \mathcal{O}_{Y', y'})}
(\mathcal{O}_{X, \overline{x}}
\otimes_{\mathcal{O}_{Y, \overline{y}}}
\mathcal{O}_{Y', \overline{y}'})
\end{align*}
这由《代数进阶》引理
\ref{more-algebra-lemma-lemma-tor-independent-flat-compare} 得出。
该引理的平坦性假设由上文以及《代数》引理
\ref{algebra-lemma-strictly-henselian-functorial-improve} 成立；
后一个引理说明 $\mathcal{O}_{X, \overline{x}}$ 是
$\mathcal{O}_{X, x} \otimes_{\mathcal{O}_{Y, y}}
\mathcal{O}_{Y, \overline{y}}$ 的严格亨泽尔化，并且
$\mathcal{O}_{Y', \overline{y}'}$ 是
$\mathcal{O}_{Y', y'} \otimes_{\mathcal{O}_{Y, y}}
\mathcal{O}_{Y, \overline{y}}$ 的严格亨泽尔化。
\end{proof}

\begin{lemma}
\label{spaces-perfect-lemma-single-complex-base-change}
设 $S$ 为概形。设
$$
\xymatrix{
X' \ar[r]_{g'} \ar[d]_{f'} &
X \ar[d]^f \\
Y' \ar[r]^g &
Y
}
$$
为 $S$ 上代数空间的笛卡尔图。
设 $K \in D_\QCoh(\mathcal{O}_X)$，并设 $L(g')^*K \to K'$
是 $D_\QCoh(\mathcal{O}_{X'})$ 中的映射。如果
\begin{enumerate}
\item 引理 \ref{spaces-perfect-lemma-single-complex-base-change-condition} 的等价条件成立；并且
\item $f$ 拟紧且拟分离，
\end{enumerate}
那么复合映射 $Lg^*Rf_*K \to Rf'_*L(g')^*K \to Rf'_*K'$
是同构。
\end{lemma}

\begin{proof}
为检验该映射是同构，可以在 $Y'$ 上 \'etale 局部地工作。
因此可以假设 $g : Y' \to Y$ 是仿射概形的态射。
在这种情形中，将使用引理 \ref{spaces-perfect-lemma-induction-principle} 的归纳原理，
证明：对于拟紧且拟分离代数空间 $U$（它在 $X$ 上 \'etale），
类似构造的映射
$Lg^*R(U \to Y)_*K|_U \to R(U' \to Y')_*K'|_{U'}$
是同构。这里
$U' = X' \times_{g', X} U = Y' \times_{g, Y} U$。

\medskip\noindent
如果 $U$ 是概形（例如仿射概形），那么该结果成立。
事实上，此时 $Y, Y', U, U'$ 是概形；由引理
\ref{spaces-perfect-lemma-derived-quasi-coherent-small-etale-site}，$K$ 和 $K'$
来自底层概形导出范畴中的对象；而《概形的导出范畴》引理
\ref{perfect-lemma-single-complex-base-change-condition} 的条件
对这些复形成立，这是引理
\ref{spaces-perfect-lemma-single-complex-base-change-condition} 的结论。
因此（利用评注
\ref{spaces-perfect-remark-match-total-direct-images} 中说明的相容性），
可以应用概形情形下的结果，即《概形的导出范畴》引理
\ref{perfect-lemma-single-complex-base-change}。

\medskip\noindent
下面是归纳步骤。设 $(U \subset W, V \to W)$ 为初等特异方形，
其中 $W$ 是在 $X$ 上 \'etale 的拟紧且拟分离代数空间，
$U$ 拟紧，$V$ 仿射，并且结果对 $U$、$V$ 和
$U \times_W V$ 成立。为简化记号，在此把 $W$ 替换为 $X$
（这一步在此处是允许的）。% P09-ERR-0468
用 $a : U \to Y$、$b : V \to Y$ 和
$c : U \times_X V \to Y$ 表示显然的态射。% P09-ERR-0469
设 $a' : U' \to Y'$、$b' : V' \to Y'$ 和
$c' : U' \times_{X'} V' \to Y'$ 分别为 $a$、$b$ 和 $c$ 的基变换。
使用相对 Mayer--Vietoris 的特异三角
（引理 \ref{spaces-perfect-lemma-unbounded-relative-mayer-vietoris}），
得到交换图
$$
\xymatrix{
Lg^*Rf_*K \ar[r] \ar[d] & Rf'_*K' \ar[d] \\
Lg^*Ra_*K|_U \oplus Lg^*Rb_*K|_V \ar[r] \ar[d] &
Ra'_* K'|_{U'} \oplus Rb'_* K'|_{V'} \ar[d] \\
Lg^*Rc_*K|_{U \times_X V} \ar[r] \ar[d] &
Rc'_*K'|_{U' \times_{X'} V'} \ar[d] \\
Lg^*Rf_* K[1] \ar[r] &
Rf'_* K'[1]
}
$$
由于第二和第三个水平箭头是同构，第一个也是同构
（《导出范畴》引理 \ref{derived-lemma-third-isomorphism-triangle}），
从而本引理的证明完成。
\end{proof}

\begin{lemma}
\label{spaces-perfect-lemma-single-complex-base-change-condition-inherited}
设 $S$ 为概形。设
$$
\xymatrix{
X' \ar[r]_{g'} \ar[d]_{f'} &
X \ar[d]^f \\
S' \ar[r]^g &
S
}
$$
为 $S$ 上代数空间的笛卡尔图。
设 $K \in D_\QCoh(\mathcal{O}_X)$，并设 $L(g')^*K \to K'$
是 $D_\QCoh(\mathcal{O}_{X'})$ 中的映射。
如果引理 \ref{spaces-perfect-lemma-single-complex-base-change-condition} 的等价条件成立，
那么
\begin{enumerate}
\item 对 $E \in D_\QCoh(\mathcal{O}_X)$，
引理 \ref{spaces-perfect-lemma-single-complex-base-change-condition} 的等价条件
对映射
$L(g')^*(E \otimes^\mathbf{L} K) \to L(g')^*E \otimes^\mathbf{L} K'$
成立；
\item 如果 $E$ 属于 $D(\mathcal{O}_X)$ 且完美，% P09-ERR-0470
那么引理 \ref{spaces-perfect-lemma-single-complex-base-change-condition} 的等价条件
对映射
$L(g')^*R\SheafHom(E, K) \to R\SheafHom(L(g')^*E, K')$ 成立；并且
\item 如果 $K$ 有下界，且 $E$ 属于 $D(\mathcal{O}_X)$
并且拟凝聚，那么引理 % P09-ERR-0471
\ref{spaces-perfect-lemma-single-complex-base-change-condition} 的等价条件对映射
$L(g')^*R\SheafHom(E, K) \to R\SheafHom(L(g')^*E, K')$ 成立。
\end{enumerate}
\end{lemma}

\begin{proof}
由引理
\ref{spaces-perfect-lemma-quasi-coherence-pullback}、
\ref{spaces-perfect-lemma-quasi-coherence-tensor-product} 和
\ref{spaces-perfect-lemma-quasi-coherence-internal-hom}，以及《位点上的上同调》引理
\ref{sites-cohomology-lemma-pseudo-coherent-pullback} 和
\ref{sites-cohomology-lemma-perfect-pullback}，
所涉复形具有拟凝聚上同调层，因而该陈述有意义。
说明这一点后，在情形 (1) 中，利用张量积的传递性计算
(\ref{spaces-perfect-equation-bc}) 的源和靶，即可检验映射
(\ref{spaces-perfect-equation-bc}) 是同构；见《代数进阶》引理
\ref{more-algebra-lemma-triple-tensor-product}。
情形 (2) 和 (3) 中的映射是复合映射
$$
L(g')^*R\SheafHom(E, K) \to R\SheafHom(L(g')^*E, L(g')^*K)
\to R\SheafHom(L(g')^*E, K')
$$
其中第一个箭头是《位点上的上同调》评注
\ref{sites-cohomology-remark-prepare-fancy-base-change} 中的箭头，
第二个箭头来自给定映射 $L(g')^*K \to K'$。
为证明映射 (\ref{spaces-perfect-equation-bc}) 是同构，在情形 (2) 中，
把 $E_x$ 表示为有限生成投射 $\mathcal{O}_{X. x}$-模的有界复形；% P09-ERR-0472
在情形 (3) 中，则用有限秩自由模的有上界复形表示它；
然后计算该箭头的源和靶。
略去若干细节。
\end{proof}

\begin{lemma}
\label{spaces-perfect-lemma-base-change-tensor}
设 $S$ 为概形。设 $f : X \to Y$ 为 $S$ 上代数空间的
拟紧且拟分离态射。设 $E \in D_\QCoh(\mathcal{O}_X)$。
设 $\mathcal{G}^\bullet$ 是由拟凝聚 $\mathcal{O}_X$-模组成的
有上界复形，并且每一项在 $Y$ 上平坦。那么
$$
Rf_*(E \otimes^\mathbf{L}_{\mathcal{O}_X} \mathcal{G}^\bullet)
$$
的形成与任意基变换可交换（精确陈述见证明）。
\end{lemma}

\begin{proof}
该陈述的意义如下。设 $g : Y' \to Y$ 为代数空间的态射，
并考虑基变换图
$$
\xymatrix{
X' \ar[r]_{g'} \ar[d]_{f'} &
X \ar[d]^f \\
Y' \ar[r]^g &
Y
}
$$
换言之，$X' = Y' \times_Y X$。本引理断言
$$
Lg^*Rf_*(E \otimes^\mathbf{L}_{\mathcal{O}_X} \mathcal{G}^\bullet)
\longrightarrow
Rf'_*(L(g')^*E \otimes^\mathbf{L}_{\mathcal{O}_{X'}} (g')^*\mathcal{G}^\bullet)
$$
是同构。注意，在右端对 $\mathcal{G}^\bullet$ {\bf 不}使用导出拉回。
为证明这一点，应用引理 \ref{spaces-perfect-lemma-single-complex-base-change} 和
\ref{spaces-perfect-lemma-single-complex-base-change-condition-inherited}，可见只须证明典范映射
$$
L(g')^*\mathcal{G}^\bullet \to (g')^*\mathcal{G}^\bullet
$$
满足引理 \ref{spaces-perfect-lemma-single-complex-base-change-condition} 的等价条件。
这可通过在茎上检验条件得出；在茎上，它立即源于如下事实：
$$
\mathcal{G}^\bullet_{\overline{x}}
\otimes_{\mathcal{O}_{Y, \overline{y}}}
\mathcal{O}_{Y', \overline{y}'}
$$
由我们对复形 $\mathcal{G}^\bullet$ 的假设计算导出张量积。
\end{proof}

\begin{lemma}
\label{spaces-perfect-lemma-base-change-RHom}
设 $S$ 为概形。设 $f : X \to Y$ 为 $S$ 上代数空间的
拟紧且拟分离态射。设 $E$ 是 $D(\mathcal{O}_X)$ 的对象。
设 $\mathcal{G}^\bullet$ 为拟凝聚 $\mathcal{O}_X$-模复形。如果
\begin{enumerate}
\item $E$ 完美，$\mathcal{G}^\bullet$ 有上界，% P09-ERR-0473
并且 $\mathcal{G}^n$ 在 $Y$ 上平坦；或
\item $E$ 拟凝聚，$\mathcal{G}^\bullet$ 有界，
并且 $\mathcal{G}^n$ 在 $Y$ 上平坦，
\end{enumerate}
那么
$$
Rf_*R\SheafHom(E, \mathcal{G}^\bullet)
$$
的形成与任意基变换可交换（精确陈述见证明）。
\end{lemma}

\begin{proof}
该陈述的意义如下。设 $g : Y' \to Y$ 为代数空间的态射，
并考虑基变换图
$$
\xymatrix{
X' \ar[r]_h \ar[d]_{f'} & % P09-ERR-0474
X \ar[d]^f \\
Y' \ar[r]^g &
Y
}
$$
换言之，$X' = Y' \times_Y X$。本引理断言
$$
Lg^*Rf_*R\SheafHom(E, \mathcal{G}^\bullet)
\longrightarrow
R(f')_*R\SheafHom(L(g')^*E, (g')^*\mathcal{G}^\bullet)
$$
是同构。注意，在右端对 $\mathcal{G}^\bullet$ {\bf 不}使用导出拉回。
为证明这一点，应用引理 \ref{spaces-perfect-lemma-single-complex-base-change} 和
\ref{spaces-perfect-lemma-single-complex-base-change-condition-inherited}，可见只须证明典范映射
$$
L(g')^*\mathcal{G}^\bullet \to (g')^*\mathcal{G}^\bullet
$$
满足引理 \ref{spaces-perfect-lemma-single-complex-base-change-condition} 的等价条件。
这一点已在引理 \ref{spaces-perfect-lemma-base-change-tensor} 的证明中得到证明。
\end{proof}















\section{构造完美复形}
\label{spaces-perfect-section-producing-perfect}

\noindent
下述引理是我们构造完美复形的主要技术工具。
这个结果后面的版本将通过诺特逼近归约到它。

\begin{lemma}
\label{spaces-perfect-lemma-perfect-direct-image}
设 $S$ 为概形。设 $Y$ 为 $S$ 上的诺特代数空间。
设 $f : X \to Y$ 为局部有限型且拟分离的代数空间态射。
设 $E \in D(\mathcal{O}_X)$，满足
\begin{enumerate}
\item $E \in D^b_{\textit{Coh}}(\mathcal{O}_X)$；
\item $H^i(E)$ 的支撑在 $Y$ 上对所有 $i$ 都适当；
\item $E$ 作为 $D(f^{-1}\mathcal{O}_Y)$ 的对象具有有限 Tor 维数。
\end{enumerate}
那么 $Rf_*E$ 是 $D(\mathcal{O}_Y)$ 的完美对象。
\end{lemma}

\begin{proof}
由引理 \ref{spaces-perfect-lemma-direct-image-coherent}，$Rf_*E$ 是
$D^b_{\textit{Coh}}(\mathcal{O}_Y)$ 的对象。因此 $Rf_*E$ 拟凝聚
（引理 \ref{spaces-perfect-lemma-identify-pseudo-coherent-noetherian}）。
所以只须证明 $Rf_*E$ 具有有限 Tor 维数；见《位点上的上同调》引理
\ref{sites-cohomology-lemma-perfect}。
由引理 \ref{spaces-perfect-lemma-tor-qc-qs}，只须检验
$Rf_*(E) \otimes_{\mathcal{O}_Y}^\mathbf{L} \mathcal{F}$
对所有拟凝聚层 $\mathcal{F}$（位于 $Y$ 上）
都具有一致有界的上同调。它有上界是显然的，因为 $Rf_*(E)$ 有上界。
设 $T \subset |X|$ 为 $H^i(E)$ 的支撑之并，其中 $i$ 取遍所有整数。
由假设 (1)、(2) 和引理 \ref{spaces-perfect-lemma-union-closed-proper-over-base}，
$T$ 在 $Y$ 上适当。特别地，存在拟紧开子空间
$X' \subset X$，它包含 $T$。置 $f' = f|_{X'}$。我们有
$Rf_*(E) = Rf'_*(E|_{X'})$，因为 $E$ 在
$X \setminus T$ 上限制为零。
因此可以把 $X$ 换成 $X'$，并假设 $f$ 拟紧。
已经假设 $f$ 拟分离。因此
$$
Rf_*(E) \otimes_{\mathcal{O}_Y}^\mathbf{L} \mathcal{F} =
Rf_*\left(E \otimes_{\mathcal{O}_X}^\mathbf{L} Lf^*\mathcal{F}\right) =
Rf_*\left(E \otimes_{f^{-1}\mathcal{O}_Y}^\mathbf{L} f^{-1}\mathcal{F}\right)
$$
这由引理 \ref{spaces-perfect-lemma-cohomology-base-change} 和《位点上的上同调》引理
\ref{sites-cohomology-lemma-variant-derived-pullback} 得出。
由假设 (3)，复形
$E \otimes_{f^{-1}\mathcal{O}_Y}^\mathbf{L} f^{-1}\mathcal{F}$
的上同调层位于某个给定有限范围中，记为 $[a, b]$。
于是它的 $Rf_*$ 的上同调位于 $[a, \infty)$ 中，结论得证。
\end{proof}

\begin{lemma}
\label{spaces-perfect-lemma-tensor-perfect}
设 $S$ 为概形。设 $B$ 为 $S$ 上的诺特代数空间。
设 $f : X \to B$ 为局部有限型且拟分离的代数空间态射。
设 $E \in D(\mathcal{O}_X)$ 完美。
设 $\mathcal{G}^\bullet$ 为由拟凝聚 $\mathcal{O}_X$-模组成的有界复形，
每一项在 $B$ 上平坦，且支撑在 $B$ 上适当。那么
$K = Rf_*(E \otimes^\mathbf{L}_{\mathcal{O}_X} \mathcal{G}^\bullet)$
是 $D(\mathcal{O}_B)$ 的完美对象。
\end{lemma}

\begin{proof}
由引理 \ref{spaces-perfect-lemma-perfect-direct-image}，对象 $K$ 完美。
下面检验该引理适用：局部地，$E$ 同构于有限秩自由
$\mathcal{O}_X$-模的有限复形。因此局部地，
$E \otimes^\mathbf{L}_{\mathcal{O}_X} \mathcal{G}^\bullet$
同构于一个有限复形，其各项具有形式
$$
\bigoplus\nolimits_{i = a, \ldots, b} (\mathcal{G}^i)^{\oplus r_i}
$$
其中 $a, b, r_a, \ldots, r_b$ 是整数。这立即蕴含上同调层
$H^i(E \otimes^\mathbf{L}_{\mathcal{O}_X} \mathcal{G})$ 是凝聚的。% P09-ERR-0475
Tor 维数的假设也成立，因为 $\mathcal{G}^i$ 在
$f^{-1}\mathcal{O}_Y$ 上平坦。% P09-ERR-0476
\end{proof}

\begin{lemma}
\label{spaces-perfect-lemma-ext-perfect}
设 $S$ 为概形。设 $B$ 为 $S$ 上的诺特代数空间。
设 $f : X \to B$ 为局部有限型且拟分离的代数空间态射。
设 $E \in D(\mathcal{O}_X)$ 完美。
设 $\mathcal{G}^\bullet$ 为由拟凝聚 $\mathcal{O}_X$-模组成的有界复形，
每一项在 $B$ 上平坦，且支撑在 $B$ 上适当。那么
$K = Rf_*R\SheafHom(E, \mathcal{G})$ 是 % P09-ERR-0477
$D(\mathcal{O}_B)$ 的完美对象。
\end{lemma}

\begin{proof}
由于 $E$ 是完美复形，存在对偶完美复形 $E^\vee$；见《位点上的上同调》引理
\ref{sites-cohomology-lemma-dual-perfect-complex}。
注意
$R\SheafHom(E, \mathcal{G}^\bullet) =
E^\vee \otimes^\mathbf{L}_{\mathcal{O}_X} \mathcal{G}^\bullet$。
因此，由引理 \ref{spaces-perfect-lemma-tensor-perfect}，$K$ 是完美的。
\end{proof}






\section{Ext 的投影公式}
\label{spaces-perfect-section-ext}

\noindent
引理 \ref{spaces-perfect-lemma-compute-ext}（或文献中的类似结果）
有时可用于验证障碍理论的性质，而这些性质是验证 Quot 函子、
Hilbert 概形及其他模问题的某个 Artin 判据所需的。
假设 $f : X \to Y$ 是代数空间的适当、平坦、有限表现态射，
并且 $E \in D(\mathcal{O}_X)$ 完美。这里该引理说明
$$
\Ext^i_X(E, f^*\mathcal{F}) =
\Ext^i_Y((Rf_*E^\vee)^\vee, \mathcal{F})
$$
其中 $\mathcal{F}$ 是 $Y$ 上的拟凝聚层。
这样写使它看起来像 Ext 的投影公式；事实上，该结果很容易由引理
\ref{spaces-perfect-lemma-cohomology-base-change} 得出。

\begin{lemma}
\label{spaces-perfect-lemma-compute-tensor-perfect}
假设和记号如引理 \ref{spaces-perfect-lemma-tensor-perfect}。那么存在函子性同构
$$
H^i(B, K \otimes^\mathbf{L}_{\mathcal{O}_B} \mathcal{F})
\longrightarrow
H^i(X, E \otimes^\mathbf{L}_{\mathcal{O}_X}
(\mathcal{G}^\bullet \otimes_{\mathcal{O}_X} f^*\mathcal{F}))
$$
其中 $\mathcal{F}$ 是 $B$ 上的拟凝聚层；这些同构与边界映射相容
（见证明）。
\end{lemma}

\begin{proof}
我们有
$$
\mathcal{G}^\bullet \otimes_{\mathcal{O}_X}^\mathbf{L} Lf^*\mathcal{F} =
\mathcal{G}^\bullet \otimes_{f^{-1}\mathcal{O}_B}^\mathbf{L} f^{-1}\mathcal{F} =
\mathcal{G}^\bullet \otimes_{f^{-1}\mathcal{O}_B} f^{-1}\mathcal{F} =
\mathcal{G}^\bullet \otimes_{\mathcal{O}_X} f^*\mathcal{F}
$$
第一个等号由《位点上的上同调》引理
\ref{sites-cohomology-lemma-variant-derived-pullback} 得出，% P09-ERR-0478
第二个等号源于 $\mathcal{G}^n$ 是平坦
$f^{-1}\mathcal{O}_B$-模，第三个等号源于拉回的定义。
因此得到
\begin{align*}
H^i(X, E \otimes^\mathbf{L}_{\mathcal{O}_X}
(\mathcal{G}^\bullet \otimes_{\mathcal{O}_X} f^*\mathcal{F}))
& =
H^i(X, E \otimes^\mathbf{L}_{\mathcal{O}_X} \mathcal{G}^\bullet
\otimes_{\mathcal{O}_X}^\mathbf{L} Lf^*\mathcal{F}) \\
& =
H^i(B,
Rf_*(E \otimes^\mathbf{L}_{\mathcal{O}_X} \mathcal{G}^\bullet
\otimes^\mathbf{L}_{\mathcal{O}_X} Lf^*\mathcal{F})) \\
& =
H^i(B,
Rf_*(E \otimes^\mathbf{L}_{\mathcal{O}_X} \mathcal{G}^\bullet)
\otimes^\mathbf{L}_{\mathcal{O}_B} \mathcal{F}) \\
& =
H^i(B, K \otimes^\mathbf{L}_{\mathcal{O}_B} \mathcal{F})
\end{align*}
第一个等号由上文得出，第二个等号由 Leray 得出
（《位点上的上同调》评注 \ref{sites-cohomology-remark-before-Leray}），
第三个等号由引理 \ref{spaces-perfect-lemma-cohomology-base-change} 得出。
关于边界映射的陈述含义如下：给定短正合列
$0 \to \mathcal{F}_1 \to \mathcal{F}_2 \to \mathcal{F}_3 \to 0$，
这些同构嵌入交换图
$$
\xymatrix{
H^i(B, K \otimes^\mathbf{L}_{\mathcal{O}_B} \mathcal{F}_3)
\ar[r] \ar[d]_\delta &
H^i(X, E \otimes^\mathbf{L}_{\mathcal{O}_X}
(\mathcal{G}^\bullet \otimes_{\mathcal{O}_X} f^*\mathcal{F}_3)) \ar[d]^\delta \\
H^{i + 1}(B, K \otimes^\mathbf{L}_{\mathcal{O}_B} \mathcal{F}_1)
\ar[r] &
H^{i + 1}(X, E \otimes^\mathbf{L}_{\mathcal{O}_X}
(\mathcal{G}^\bullet \otimes_{\mathcal{O}_X} f^*\mathcal{F}_1))
}
$$
其中边界映射来自特异三角
$$
K \otimes^\mathbf{L}_{\mathcal{O}_B} \mathcal{F}_1 \to
K \otimes^\mathbf{L}_{\mathcal{O}_B} \mathcal{F}_2 \to
K \otimes^\mathbf{L}_{\mathcal{O}_B} \mathcal{F}_3 \to
K \otimes^\mathbf{L}_{\mathcal{O}_B} \mathcal{F}_1[1]
$$
以及 $D(\mathcal{O}_X)$ 中与复形的短正合列
$$
0 \to
\mathcal{G}^\bullet \otimes_{\mathcal{O}_X} f^*\mathcal{F}_1 \to
\mathcal{G}^\bullet \otimes_{\mathcal{O}_X} f^*\mathcal{F}_2 \to
\mathcal{G}^\bullet \otimes_{\mathcal{O}_X} f^*\mathcal{F}_3 \to 0
$$
相伴的特异三角。这个序列正合，因为 $\mathcal{G}^n$ 在 $B$ 上平坦。
略去对所显示图之交换性的验证。
\end{proof}

\begin{lemma}
\label{spaces-perfect-lemma-compute-ext-perfect}
假设和记号如引理 \ref{spaces-perfect-lemma-ext-perfect}。% P09-ERR-0479
那么存在函子性同构
$$
H^i(B, K \otimes^\mathbf{L}_{\mathcal{O}_B} \mathcal{F})
\longrightarrow
\Ext^i_{\mathcal{O}_X}(E,
\mathcal{G}^\bullet \otimes_{\mathcal{O}_X} f^*\mathcal{F})
$$
其中 $\mathcal{F}$ 是 $B$ 上的拟凝聚层；这些同构与边界映射相容
（见证明）。
\end{lemma}

\begin{proof}
如引理 \ref{spaces-perfect-lemma-ext-perfect} 的证明中那样，令
$E^\vee$ 为对偶完美复形，并回忆
$K = Rf_*(E^\vee \otimes_{\mathcal{O}_X}^\mathbf{L} \mathcal{G}^\bullet)$。
此外，我们还有
$$
\Ext^i_{\mathcal{O}_X}(E,
\mathcal{G}^\bullet \otimes_{\mathcal{O}_X} f^*\mathcal{F})
=
H^i(X, E^\vee \otimes^\mathbf{L}_{\mathcal{O}_X}
(\mathcal{G}^\bullet \otimes_{\mathcal{O}_X} f^*\mathcal{F}))
$$
这是 $E^\vee$ 的构造所给出的。所以这些同构的存在性由引理
\ref{spaces-perfect-lemma-compute-tensor-perfect} 应用于 $E^\vee$ 和
$\mathcal{G}^\bullet$ 得出。
关于边界映射的陈述含义如下：给定短正合列
$0 \to \mathcal{F}_1 \to \mathcal{F}_2 \to \mathcal{F}_3 \to 0$，
这些同构组成交换图
$$
\xymatrix{
H^i(B, K \otimes^\mathbf{L}_{\mathcal{O}_B} \mathcal{F}_3)
\ar[r] \ar[d]_\delta &
\Ext^i_{\mathcal{O}_X}(E,
\mathcal{G}^\bullet \otimes_{\mathcal{O}_X} f^*\mathcal{F}_3) \ar[d]^\delta \\
H^{i + 1}(B, K \otimes^\mathbf{L}_{\mathcal{O}_B} \mathcal{F}_1)
\ar[r] &
\Ext^{i + 1}_{\mathcal{O}_X}(E,
\mathcal{G}^\bullet \otimes_{\mathcal{O}_X} f^*\mathcal{F}_1)
}
$$
其中边界映射来自特异三角
$$
K \otimes^\mathbf{L}_{\mathcal{O}_B} \mathcal{F}_1 \to
K \otimes^\mathbf{L}_{\mathcal{O}_B} \mathcal{F}_2 \to
K \otimes^\mathbf{L}_{\mathcal{O}_B} \mathcal{F}_3 \to
K \otimes^\mathbf{L}_{\mathcal{O}_B} \mathcal{F}_1[1]
$$
以及 $D(\mathcal{O}_X)$ 中与复形的短正合列
$$
0 \to
\mathcal{G}^\bullet \otimes_{\mathcal{O}_X} f^*\mathcal{F}_1 \to
\mathcal{G}^\bullet \otimes_{\mathcal{O}_X} f^*\mathcal{F}_2 \to
\mathcal{G}^\bullet \otimes_{\mathcal{O}_X} f^*\mathcal{F}_3 \to 0
$$
相伴的特异三角。这个序列正合，因为 $\mathcal{G}^n$ 在 $B$ 上平坦。
略去对所显示图之交换性的验证。
\end{proof}

\begin{lemma}
\label{spaces-perfect-lemma-compute-ext}
设 $S$ 为概形。设 $f : X \to B$ 为 $S$ 上代数空间的态射，
$E \in D(\mathcal{O}_X)$，并设 $\mathcal{F}^\bullet$ 为 % P09-ERR-0480
$\mathcal{O}_X$-模复形。假设
\begin{enumerate}
\item $B$ 是诺特的；
\item $f$ 局部有限型且拟分离；
\item $E \in D^-_{\textit{Coh}}(\mathcal{O}_X)$；
\item $\mathcal{G}^\bullet$ 是由凝聚 $\mathcal{O}_X$-模组成的
有界复形，每一项在 $B$ 上平坦，且支撑在 $B$ 上适当。% P09-ERR-0481
\end{enumerate}
那么下列两个陈述成立：
\begin{enumerate}
\item[(A)] 对每个 $m \in \mathbf{Z}$，存在完美对象 $K$
（属于 $D(\mathcal{O}_B)$）和函子性映射
$$
\alpha^i_\mathcal{F} :
\Ext^i_{\mathcal{O}_X}(E,
\mathcal{G}^\bullet \otimes_{\mathcal{O}_X} f^*\mathcal{F})
\longrightarrow
H^i(B, K \otimes^\mathbf{L}_{\mathcal{O}_B} \mathcal{F})
$$
其中 $\mathcal{F}$ 是 $B$ 上的拟凝聚层；这些映射与边界映射相容
（见证明），并且 $\alpha^i_\mathcal{F}$ 在 $i \leq m$ 时是同构；
\item[(B)] 存在拟凝聚对象 $L \in D(\mathcal{O}_B)$ 和函子性同构
$$
\Ext^i_{\mathcal{O}_B}(L, \mathcal{F}) \longrightarrow
\Ext^i_{\mathcal{O}_X}(E,
\mathcal{G}^\bullet \otimes_{\mathcal{O}_X} f^*\mathcal{F})
$$
其中 $\mathcal{F}$ 是 $B$ 上的拟凝聚层；这些同构与边界映射相容。
\end{enumerate}
\end{lemma}

\begin{proof}
(A) 的证明。假设 $\mathcal{G}^i$ 仅当 $i \in [a, b]$ 时非零。
可以把 $X$ 换成 $\mathcal{G}^i$ 的支撑之并的拟紧开邻域。
因此可以假设 $X$ 是诺特的。在这种情形中，$X$ 和 $f$
都是拟紧且拟分离的。选取逼近 $P \to E$，其中 $P$ 是
$(X, E, -m - 1 + a)$ 的完美复形逼近
（这由定理 \ref{spaces-perfect-theorem-approximation} 保证存在）。
于是诱导映射
$$
\Ext^i_{\mathcal{O}_X}(E,
\mathcal{G}^\bullet \otimes_{\mathcal{O}_X} f^*\mathcal{F})
\longrightarrow
\Ext^i_{\mathcal{O}_X}(P,
\mathcal{G}^\bullet \otimes_{\mathcal{O}_X} f^*\mathcal{F})
$$
当 $i \leq m$ 时是同构。事实上，该映射的核和余核分别是
下列两个群的商模和子模：
$$
\Ext^i_{\mathcal{O}_X}(C,
\mathcal{G}^\bullet \otimes_{\mathcal{O}_X} f^*\mathcal{F})
\quad\text{分别}\quad
\Ext^{i + 1}_{\mathcal{O}_X}(C,
\mathcal{G}^\bullet \otimes_{\mathcal{O}_X} f^*\mathcal{F})
$$
其中 $C$ 是 $P \to E$ 的锥。由于 $C$ 在次数
$\geq -m - 1 + a$ 的上同调层消失，由《导出范畴》引理
\ref{derived-lemma-negative-exts}，这些 $\Ext$ 群对
$i \leq m + 1$ 为零。
这把问题归约到 $E$ 为完美复形的情形，也就是引理
\ref{spaces-perfect-lemma-compute-ext-perfect}。% P09-ERR-0482
关于边界的陈述已在引理
\ref{spaces-perfect-lemma-compute-ext-perfect} 的证明中说明。

\medskip\noindent
(B) 的证明。如 (A) 的证明中那样，可以假设 $X$ 是诺特的。
由引理 \ref{spaces-perfect-lemma-identify-pseudo-coherent-noetherian}，$E$ 拟凝聚。
由引理 \ref{spaces-perfect-lemma-pseudo-coherent-hocolim}，可以写成
$E = \text{hocolim} E_n$，其中 $E_n$ 完美，并且 $E_n \to E$
在截断 $\tau_{\geq -n}$ 上诱导同构。设 $E_n^\vee$ 为对偶完美复形
（《位点上的上同调》引理
\ref{sites-cohomology-lemma-dual-perfect-complex}）。
得到完美对象的逆系统
$\ldots \to E_3^\vee \to E_2^\vee \to E_1^\vee$。
这又给出逆系统
$$
\ldots \to K_3 \to K_2 \to K_1\quad\text{其中}\quad
K_n = Rf_*(E_n^\vee \otimes_{\mathcal{O}_X}^\mathbf{L} \mathcal{G}^\bullet)
$$
它在 $Y$ 上完美；见引理 \ref{spaces-perfect-lemma-tensor-perfect}。% P09-ERR-0483
由引理 \ref{spaces-perfect-lemma-compute-ext-perfect} 及其证明，以及上一段的论证
（取 $P = E_n$），对任意拟凝聚层 $\mathcal{F}$（位于 $Y$ 上），
存在函子性典范映射
$$
\xymatrix{
& \Ext^i_{\mathcal{O}_X}(E,
\mathcal{G}^\bullet \otimes_{\mathcal{O}_X} f^*\mathcal{F})
\ar[ld] \ar[rd] \\
H^i(Y, K_{n + 1} \otimes_{\mathcal{O}_Y}^\mathbf{L} \mathcal{F})
\ar[rr] & &
H^i(Y, K_n \otimes_{\mathcal{O}_Y}^\mathbf{L} \mathcal{F})
}
$$
它们对 $i \leq n + a$ 是同构。
设 $L_n = K_n^\vee$ 为对偶完美复形。
于是 $L_1 \to L_2 \to L_3 \to \ldots$ 是
$D(\mathcal{O}_Y)$ 中的完美对象系统，使得对于任意拟凝聚层
$\mathcal{F}$（位于 $Y$ 上），映射
$$
\Ext^i_{\mathcal{O}_Y}(L_{n + 1}, \mathcal{F})
\longrightarrow
\Ext^i_{\mathcal{O}_Y}(L_n, \mathcal{F})
$$
对 $i \leq n + a - 1$ 是同构。这蕴含 $L_n \to L_{n + 1}$
在截断 $\tau_{\geq -n - a + 2}$ 上诱导同构
（提示：取 $L_n \to L_{n + 1}$ 的锥，并考察它最后一个非零上同调层）。
因此 $L = \text{hocolim} L_n$ 拟凝聚；见引理
\ref{spaces-perfect-lemma-pseudo-coherent-hocolim}。同伦余极限的映射性质给出
$\Ext^i_{\mathcal{O}_Y}(L, \mathcal{F}) =
\Ext^i_{\mathcal{O}_Y}(L_n, \mathcal{F})$
对 $i \leq n + a - 3$ 成立，从而证明完成。
\end{proof}

\begin{remark}
\label{spaces-perfect-remark-base-change-of-L}
拟凝聚复形 $L$（即引理 \ref{spaces-perfect-lemma-compute-ext} 的 (B) 部分中的复形）
典范地由这一情形确定。例如，$L$ 按 (B) 中方式的构造与基变换相容。
换言之，给定概形的笛卡儿图 % P09-ERR-0484
$$
\xymatrix{
X' \ar[r]_{g'} \ar[d]_{f'} &
X \ar[d]^f \\
Y' \ar[r]^g &
Y
}
$$
我们有函子性的典范同构
$$
\Ext^i_{\mathcal{O}_{Y'}}(Lg^*L, \mathcal{F}') \longrightarrow
\Ext^i_{\mathcal{O}_X}(L(g')^*E,
(g')^*\mathcal{G}^\bullet \otimes_{\mathcal{O}_{X'}} (f')^*\mathcal{F}')
$$
% P09-ERR-0485
其中 $\mathcal{F}'$ 是 $Y'$ 上的拟凝聚层。注意，在右端我们\textbf{没有}对 % P09-ERR-0486
$\mathcal{G}^\bullet$ 使用导出拉回。若将来需要这种结果，我们会在此
表述一个精确结果并给出详细证明。
\end{remark}

\section{极限与导出范畴}
\label{spaces-perfect-section-limits}

\noindent
本节汇集关于如下代数空间之导出范畴的一些结果：该代数空间是
一个代数空间逆系统的极限。更确切地说，我们将在下列情形中工作。

\begin{situation}
\label{spaces-perfect-situation-descent}
设 $S$ 为概形。设 $X = \lim_{i \in I} X_i$ 是 $S$ 上代数空间的
有向逆系统之极限，其转移态射
$f_{i'i} : X_{i'} \to X_i$ 是仿射的。把投影记为 % P09-ERR-0487
$f_i : X \to X_i$。假设 $X_i$ 对所有 $i \in I$ 都拟紧且拟分离。
还选定一个元素 $0 \in I$。
\end{situation}

\begin{lemma}
\label{spaces-perfect-lemma-descend-homomorphisms}
在情形 \ref{spaces-perfect-situation-descent} 中，令 $E_0$ 和 $K_0$ 为
$D(\mathcal{O}_{X_0})$ 的对象。置
$E_i = Lf_{i0}^*E_0$ 和 $K_i = Lf_{i0}^*K_0$（$i \geq 0$），
并置 $E = Lf_0^*E_0$ 和 $K = Lf_0^*K_0$。那么映射
$$
\colim_{i \geq 0} \Hom_{D(\mathcal{O}_{X_i})}(E_i, K_i)
\longrightarrow
\Hom_{D(\mathcal{O}_X)}(E, K)
$$
在下列任一情形中为同构：
\begin{enumerate}
\item $E_0$ 完美且 $K_0 \in D_\QCoh(\mathcal{O}_{X_0})$；或
\item $E_0$ 拟凝聚，且
$K_0 \in D_\QCoh(\mathcal{O}_{X_0})$ 具有有限 Tor 维数。
\end{enumerate}
\end{lemma}

\begin{proof}
对每个拟紧拟分离对象 $U_0$（属于 $(X_0)_{spaces, \etale}$），
考虑如下条件 $P$：典范映射
$$
\colim_{i \geq 0} \Hom_{D(\mathcal{O}_{U_i})}(E_i|_{U_i}, K_i|_{U_i})
\longrightarrow
\Hom_{D(\mathcal{O}_U)}(E|_U, K|_U)
$$
是同构，其中 $U = X \times_{X_0} U_0$ 且
$U_i = X_i \times_{X_0} U_0$。我们将用引理
\ref{spaces-perfect-lemma-induction-principle} 的归纳原理证明 $P$ 对每个 $U_0$ 成立。
该引理的条件 (2) 立即由导出范畴中 Hom 的 Mayer--Vietoris 性质得出；
见引理 \ref{spaces-perfect-lemma-mayer-vietoris-hom}。
因此只需在 $X_0$ 仿射时证明本引理。

\medskip\noindent
若 $X_0$ 仿射，则结果由概形情形得出；见《概形的导出范畴》引理
\ref{perfect-lemma-descend-homomorphisms}。
为此，使用引理 \ref{spaces-perfect-lemma-derived-quasi-coherent-small-etale-site}
的等价，并使用引理
\ref{spaces-perfect-lemma-descend-pseudo-coherent}、
\ref{spaces-perfect-lemma-descend-tor-amplitude} 和
\ref{spaces-perfect-lemma-descend-perfect} 中说明的性质对应。
\end{proof}

\begin{lemma}
\label{spaces-perfect-lemma-perfect-on-limit}
在情形 \ref{spaces-perfect-situation-descent} 中，$D(\mathcal{O}_X)$ 的完美对象范畴
是 $D(\mathcal{O}_{X_i})$ 的完美对象范畴之余极限。
\end{lemma}

\begin{proof}
对每个拟紧拟分离对象 $U_0$（属于 $(X_0)_{spaces, \etale}$），
考虑条件 $P$：函子
$$
\colim_{i \geq 0} D_{perf}(\mathcal{O}_{U_i})
\longrightarrow
D_{perf}(\mathcal{O}_U)
$$
是等价；其中 ${}_{perf}$ 表示由完美对象构成的充满子范畴，而
$U = X \times_{X_0} U_0$ 且
$U_i = X_i \times_{X_0} U_0$。我们将用引理
\ref{spaces-perfect-lemma-induction-principle} 的归纳原理证明 $P$ 对每个 $U_0$ 成立。
首先，由引理 \ref{spaces-perfect-lemma-descend-homomorphisms}，已经知道该函子全忠实。
因此只需证明本质满性。

\medskip\noindent
先验证归纳原理的条件 (2)。因而，设有初等特异方形
$(U_0 \subset X_0, V_0 \to X_0)$，并且 $P$ 对
$U_0$、$V_0$ 和 $U_0 \times_{X_0} V_0$ 成立。设 $E$ 是
$D(\mathcal{O}_X)$ 的完美对象。可以找到 $i \geq 0$、完美对象
$E_{U, i}$（位于 $U_i$ 上）和完美对象 $E_{V, i}$（位于 $V_i$ 上），
它们到 $U$ 和 $V$ 的拉回分别同构于
$E|_U$ 和 $E|_V$。记
$$
a : E_{U, i} \to (R(X \to X_i)_*E)|_{U_i}
\quad\text{且}\quad
b : E_{V, i} \to (R(X \to X_i)_*E)|_{V_i}
$$
为与同构 $L(U \to U_i)^*E_{U, i} \to E|_U$
和 $L(V \to V_i)^*E_{V, i} \to E|_V$ 伴随的映射。
由全忠实性，% P09-ERR-0488
在增大 $i$ 后，可以找到同构
$c : E_{U, i}|_{U_i \times_{X_i} V_i} \to E_{V, i}|_{U_i \times_{X_i} V_i}$，
其拉回给出下列识别：
$$
L(U \to U_i)^*E_{U, i}|_{U \times_X V} \to E|_{U \times_X V} \to
L(V \to V_i)^*E_{V, i}|_{U \times_X V}.
$$
应用引理 \ref{spaces-perfect-lemma-glue}，得到对象 $E_i$（位于 $X_i$ 上）和映射
$d : E_i \to R(X \to X_i)_*E$；其限制是映射 $a$ 和 $b$，
分别位于 $U_i$ 和 $V_i$ 上。于是显然 $E_i$ 完美，并且
$d$ 与同构 $L(X \to X_i)^*E_i \to E$ 伴随。

\medskip\noindent
最后验证归纳原理的条件 (1)，换言之，验证本引理在 $X_0$ 仿射时成立。
这由概形情形得出；见《概形的导出范畴》引理
\ref{perfect-lemma-descend-perfect}。
为此，使用引理
\ref{spaces-perfect-lemma-derived-quasi-coherent-small-etale-site} 的等价，
并使用引理 \ref{spaces-perfect-lemma-descend-perfect} 中说明的性质对应。
\end{proof}

\section{上同调与基变换，VI}
\label{spaces-perfect-section-cohomology-and-base-change-final}

\noindent
这是关于上同调与基变换的最后一节，继续第
\ref{spaces-perfect-section-cohomology-and-base-change-perfect}、
\ref{spaces-perfect-section-cohomology-base-change} 和
\ref{spaces-perfect-section-producing-perfect} 节的讨论。
备注 \ref{spaces-perfect-remark-explain-perfect-direct-image} 给出了一个易于把握的特例。

\begin{lemma}
\label{spaces-perfect-lemma-base-change-tensor-perfect}
设 $S$ 为概形。设 $f : X \to Y$ 是 $S$ 上代数空间之间的
有限表现态射。设 $E \in D(\mathcal{O}_X)$ 为完美对象。
设 $\mathcal{G}^\bullet$ 是由有限表现 $\mathcal{O}_X$-模组成的
有界复形，它在 $Y$ 上平坦，且支撑在 $Y$ 上适当。那么
$$
K = Rf_*(E \otimes_{\mathcal{O}_X}^\mathbf{L} \mathcal{G}^\bullet)
$$
是 $D(\mathcal{O}_Y)$ 的完美对象，并且它的构造与任意基变换可交换。
\end{lemma}

\begin{proof}
关于基变换的陈述是引理 \ref{spaces-perfect-lemma-base-change-tensor}。
因此只需证明 $K$ 是完美对象。若 $Y$ 诺特，则这由引理
\ref{spaces-perfect-lemma-tensor-perfect} 得出。
我们将用诺特逼近把问题归约到这种情形。
建议读者跳过本证明的其余部分。

\medskip\noindent
问题在 $Y$ 上是局部的，故可假设 $Y$ 仿射。
设 $Y = \Spec(R)$。把 $R = \colim R_i$ 写成诺特环
$R_i$ 的滤过余极限。由《空间的极限》引理
\ref{spaces-limits-lemma-descend-finite-presentation}，
存在 $i$ 和代数空间 $X_i$，它在 $R_i$ 上有限表现，且到
$R$ 的基变换是 $X$。由《空间的极限》引理
\ref{spaces-limits-lemma-descend-modules-finite-presentation}，
增大 $i$ 后，可以假设存在由有限表现
$\mathcal{O}_{X_i}$-模组成的有界复形 $\mathcal{G}_i^\bullet$，
它到 $X$ 的拉回是 $\mathcal{G}^\bullet$。增大 $i$ 后，
可以假设 $\mathcal{G}_i^n$ 在 $R_i$ 上平坦；见《空间的极限》引理
\ref{spaces-limits-lemma-descend-flat}。
增大 $i$ 后，可以假设 $\mathcal{G}_i^n$ 的支撑在 $R_i$ 上适当；
见《空间的极限》引理
\ref{spaces-limits-lemma-eventually-proper-support}。
最后，由引理 \ref{spaces-perfect-lemma-perfect-on-limit}，增大 $i$ 后，
可以假设存在完美对象 $E_i$（属于 $D(\mathcal{O}_{X_i})$），
它到 $X$ 的拉回是 $E$。由引理 \ref{spaces-perfect-lemma-tensor-perfect}，
$$
K_i =
Rf_{i, *}(E_i \otimes_{\mathcal{O}_{X_i}}^\mathbf{L} \mathcal{G}_i^\bullet)
$$
在 $\Spec(R_i)$ 上完美，其中 $f_i : X_i \to \Spec(R_i)$ 是结构态射。
由基变换结果（引理 \ref{spaces-perfect-lemma-base-change-tensor}），$K_i$ 到
$Y = \Spec(R)$ 的拉回是 $K$，故结论成立。
\end{proof}

\begin{remark}
\label{spaces-perfect-remark-explain-perfect-direct-image}
设 $R$ 为环。设 $X$ 为 $R$ 上有限表现的代数空间。
设 $\mathcal{G}$ 为有限表现 $\mathcal{O}_X$-模，它在 $R$ 上平坦，
且支撑在 $R$ 上适当。由引理
\ref{spaces-perfect-lemma-base-change-tensor-perfect}，存在由有限投射
$R$-模组成的有限复形 $M^\bullet$，使得
$$
R\Gamma(X_{R'}, \mathcal{G}_{R'}) = M^\bullet \otimes_R R'
$$
对 $R$-代数 $R'$ 函子性成立。
\end{remark}

\begin{lemma}
\label{spaces-perfect-lemma-base-change-tensor-pseudo-coherent}
设 $S$ 为概形。
设 $f : X \to Y$ 是 $S$ 上代数空间之间的有限表现态射。
设 $E \in D(\mathcal{O}_X)$ 为拟凝聚对象。
设 $\mathcal{G}^\bullet$ 是由有限表现 $\mathcal{O}_X$-模组成的
上有界复形，它在 $Y$ 上平坦，且支撑在 $Y$ 上适当。那么
$$
K = Rf_*(E \otimes_{\mathcal{O}_X}^\mathbf{L} \mathcal{G}^\bullet)
$$
是 $D(\mathcal{O}_Y)$ 的拟凝聚对象，并且它的构造与任意基变换可交换。
\end{lemma}

\begin{proof}
关于基变换的陈述是引理 \ref{spaces-perfect-lemma-base-change-tensor}。
因此只需证明 $K$ 是拟凝聚对象。
这将由引理 \ref{spaces-perfect-lemma-base-change-tensor-perfect}
通过完美复形逼近得出。建议读者跳过本证明的其余部分。

\medskip\noindent
问题在 $Y$ 上是平展局部的，故可假设 $Y$ 仿射。
于是 $X$ 拟紧且拟分离。此外，存在整数 $N$，使得全正像
$Rf_* : D_\QCoh(\mathcal{O}_X) \to D_\QCoh(\mathcal{O}_Y)$
具有上同调维数 $N$；见引理
\ref{spaces-perfect-lemma-quasi-coherence-direct-image} 中的说明。
选取整数 $b$，使得 $\mathcal{G}^i = 0$ 对 $i > b$ 成立。
只需证明 $K$ 对每个 $m$ 都是 $m$-拟凝聚的。
选取逼近 $P \to E$，其中 $P$ 是完美复形，且该逼近属于
$(X, E, m - N - 1 - b)$。这由定理
\ref{spaces-perfect-theorem-approximation} 保证存在。
选取特异三角
$$
P \to E \to C \to P[1]
$$
，它位于 $D_\QCoh(\mathcal{O}_X)$ 中。$C$ 的上同调层在次数
$\geq m - N - 1 - b$ 为零。
因此 $C \otimes^\mathbf{L} \mathcal{G}^\bullet$ 的上同调层
在次数 $\geq m - N - 1$ 为零。
从而 $Rf_*(C \otimes^\mathbf{L} \mathcal{G})$ 的上同调层 % P09-ERR-0489
在次数 $\geq m - 1$ 为零。因此
$$
Rf_*(P \otimes^\mathbf{L} \mathcal{G}) \to
Rf_*(E \otimes^\mathbf{L} \mathcal{G})
$$
在次数 $\geq m$ 的上同调层上是同构。
其次，假设 $H^i(P) = 0$ 对 $i > a$ 成立。那么
$
P \otimes^\mathbf{L} \sigma_{\geq m - N - 1 - a}\mathcal{G}^\bullet
\longrightarrow
P \otimes^\mathbf{L} \mathcal{G}^\bullet
$
在次数 $\geq m - N - 1$ 的上同调层上是同构。
因而再次得到
$$
Rf_*(P \otimes^\mathbf{L} \sigma_{\geq m - N - 1 - a}\mathcal{G}^\bullet) \to
Rf_*(P \otimes^\mathbf{L} \mathcal{G}^\bullet)
$$
在次数 $\geq m$ 的上同调层上是同构。
由引理 \ref{spaces-perfect-lemma-base-change-tensor-perfect}，源是完美复形。
因此 $K$ 如所需地为 $m$-拟凝聚。
\end{proof}

\begin{lemma}
\label{spaces-perfect-lemma-flat-proper-perfect-direct-image-general}
设 $S$ 为概形。设 $f : X \to Y$ 是 $S$ 上代数空间之间的
适当有限表现态射。
\begin{enumerate}
\item 设 $E \in D(\mathcal{O}_X)$ 完美且 $f$ 平坦。那么
$Rf_*E$ 是 $D(\mathcal{O}_Y)$ 的完美对象，并且它的构造
与任意基变换可交换。
\item 设 $\mathcal{G}$ 为有限表现 $\mathcal{O}_X$-模，
且在 $S$ 上平坦。那么 $Rf_*\mathcal{G}$ 是 % P09-ERR-0490
$D(\mathcal{O}_Y)$ 的完美对象，并且它的构造与任意基变换可交换。
\end{enumerate}
\end{lemma}

\begin{proof}
这是引理 \ref{spaces-perfect-lemma-base-change-tensor-perfect} 的特例，分别取
(1) $\mathcal{G}^\bullet$ 等于 $\mathcal{O}_X$（置于次数 $0$），以及
(2) $E = \mathcal{O}_X$，并令 $\mathcal{G}^\bullet$
由 $\mathcal{G}$ 组成（置于次数 $0$）。
\end{proof}

\begin{lemma}
\label{spaces-perfect-lemma-flat-proper-pseudo-coherent-direct-image-general}
设 $S$ 为概形。设 $f : X \to Y$ 是 $S$ 上代数空间之间的
平坦适当有限表现态射。设 $E \in D(\mathcal{O}_X)$ 拟凝聚。
那么 $Rf_*E$ 是 $D(\mathcal{O}_Y)$ 的拟凝聚对象，
并且它的构造与任意基变换可交换。
\end{lemma}

\noindent
更一般地，若 $f : X \to Y$ 适当，而 $E$（位于 $X$ 上）相对于
$Y$ 拟凝聚（《空间态射进阶》定义
\ref{spaces-more-morphisms-definition-relative-pseudo-coherence}），
则 $Rf_*E$ 拟凝聚
（但在这种一般性下，其构造不与基变换可交换）。
概形的这种情形在 \cite{Kiehl} 中得到证明。

\begin{proof}
这是引理 \ref{spaces-perfect-lemma-base-change-tensor-pseudo-coherent} 的特例，
取 $\mathcal{G} = \mathcal{O}_X$。
\end{proof}

\begin{lemma}
\label{spaces-perfect-lemma-pullback-and-limits}
设 $R$ 为环。设 $X$ 为代数空间，并设
$f : X \to \Spec(R)$ 适当、平坦且有限表现。
设 $(M_n)$ 为具有满射转移映射的 $R$-模逆系统。
那么典范映射
$$
\mathcal{O}_X \otimes_R (\lim M_n)
\longrightarrow
\lim \mathcal{O}_X \otimes_R M_n
$$
诱导从源到对目标应用 $DQ_X$ 所得对象的同构。
\end{lemma}

\begin{proof}
该陈述意指：对任意对象 $E$（属于
$D_\QCoh(\mathcal{O}_X)$），诱导映射
$$
\Hom(E, \mathcal{O}_X \otimes_R (\lim M_n))
\longrightarrow
\Hom(E, \lim \mathcal{O}_X \otimes_R M_n)
$$
是同构。由于 $D_\QCoh(\mathcal{O}_X)$ 有完美生成元
（定理 \ref{spaces-perfect-theorem-bondal-van-den-Bergh}），只需对完美 $E$ 验证。
由引理 \ref{spaces-perfect-lemma-Rlim-quasi-coherent}，
$\lim \mathcal{O}_X \otimes_R M_n = R\lim \mathcal{O}_X \otimes_R M_n$。
《位点上的上同调》第 \ref{sites-cohomology-section-global-RHom} 节的
正合函子
$R\Hom_X(E, -) : D_\QCoh(\mathcal{O}_X) \to D(R)$
与积可交换，因而与导出极限可交换，所以
$$
R\Hom_X(E, \lim \mathcal{O}_X \otimes_R M_n) =
R\lim R\Hom_X(E, \mathcal{O}_X \otimes_R M_n)
$$
令 $E^\vee$ 为对偶完美复形；见《位点上的上同调》引理
\ref{sites-cohomology-lemma-dual-perfect-complex}。我们有
$$
R\Hom_X(E, \mathcal{O}_X \otimes_R M_n) =
R\Gamma(X, E^\vee \otimes_{\mathcal{O}_X}^\mathbf{L} Lf^*M_n) =
R\Gamma(X, E^\vee) \otimes_R^\mathbf{L} M_n
$$
，这由引理 \ref{spaces-perfect-lemma-cohomology-base-change} 得出。
由引理 \ref{spaces-perfect-lemma-flat-proper-perfect-direct-image-general} 可知
$R\Gamma(X, E^\vee)$ 是 $R$-模的完美复形。
特别地，它是拟凝聚复形；由《代数进阶》引理
\ref{more-algebra-lemma-pseudo-coherent-tensor-limit}，得到
$$
R\lim R\Gamma(X, E^\vee) \otimes_R^\mathbf{L} M_n =
R\Gamma(X, E^\vee) \otimes_R^\mathbf{L} \lim M_n
$$
，正如所需。
\end{proof}

\begin{lemma}
\label{spaces-perfect-lemma-perfect-enough}
设 $A$ 为环。设 $X$ 为 $A$ 上拟紧拟分离的代数空间。
设 $K \in D^-_\QCoh(\mathcal{O}_X)$。若
$R\Gamma(X, E \otimes^\mathbf{L} K)$ 在 $D(A)$ 中拟凝聚，
对每个完美 $E$（属于 $D(\mathcal{O}_X)$）都如此，那么
$R\Gamma(X, E \otimes^\mathbf{L} K)$ 在 $D(A)$ 中拟凝聚，
对每个拟凝聚 $E$（属于 $D(\mathcal{O}_X)$）都如此。
\end{lemma}

\begin{proof}
存在整数 $N$，使得
$R\Gamma(X, -) : D_\QCoh(\mathcal{O}_X) \to D(A)$
具有上同调维数 $N$；见引理
\ref{spaces-perfect-lemma-quasi-coherence-direct-image} 中的说明。
取 $b \in \mathbf{Z}$，使得 $H^i(K) = 0$ 对 $i > b$ 成立。
设 $E$ 在 $X$ 上拟凝聚。
只需证明 $R\Gamma(X, E \otimes^\mathbf{L} K)$
对每个 $m$ 都是 $m$-拟凝聚的。
选取逼近 $P \to E$，其中 $P$ 是完美复形，且该逼近属于
$(X, E, m - N - 1 - b)$。这由定理
\ref{spaces-perfect-theorem-approximation} 保证存在。
选取特异三角
$$
P \to E \to C \to P[1]
$$
，它位于 $D_\QCoh(\mathcal{O}_X)$ 中。$C$ 的上同调层
在次数 $\geq m - N - 1 - b$ 为零。因此
$C \otimes^\mathbf{L} K$ 的上同调层在次数
$\geq m - N - 1$ 为零。
从而 $R\Gamma(X, C \otimes^\mathbf{L} K)$ 的上同调 % P09-ERR-0491
在次数 $\geq m - 1$ 为零。因此
$$
R\Gamma(X, P \otimes^\mathbf{L} K) \to R\Gamma(X, E \otimes^\mathbf{L} K)
$$
在次数 $\geq m$ 的上同调上是同构。
根据假设，源拟凝聚。
因此 $R\Gamma(X, E \otimes^\mathbf{L} K)$
如所需地为 $m$-拟凝聚。
\end{proof}

\begin{lemma}
\label{spaces-perfect-lemma-base-change-RHom-perfect}
设 $S$ 为概形。
设 $f : X \to Y$ 是 $S$ 上代数空间之间的有限表现态射。
设 $E \in D(\mathcal{O}_X)$ 为完美对象。设 $\mathcal{G}^\bullet$
为由有限表现 $\mathcal{O}_X$-模组成的有界复形，
它在 $Y$ 上平坦，且支撑在 $Y$ 上适当。那么
$$
K = Rf_*R\SheafHom(E, \mathcal{G}^\bullet)
$$
是 $D(\mathcal{O}_Y)$ 的完美对象，并且它的构造与任意基变换可交换。
\end{lemma}

\begin{proof}
关于基变换的陈述是引理 \ref{spaces-perfect-lemma-base-change-RHom}。
因此只需证明 $K$ 是完美对象。若 $Y$ 诺特，则这由引理
\ref{spaces-perfect-lemma-ext-perfect} 得出。
我们将用诺特逼近把问题归约到这种情形。
建议读者跳过本证明的其余部分。

\medskip\noindent
问题在 $Y$ 上是局部的，故可假设 $Y$ 仿射。
设 $Y = \Spec(R)$。把 $R = \colim R_i$ 写成诺特环
$R_i$ 的滤过余极限。由《空间的极限》引理
\ref{spaces-limits-lemma-descend-finite-presentation}，
存在 $i$ 和代数空间 $X_i$，它在 $R_i$ 上有限表现，且到
$R$ 的基变换是 $X$。由《空间的极限》引理
\ref{spaces-limits-lemma-descend-modules-finite-presentation}，
增大 $i$ 后，可以假设存在由有限表现
$\mathcal{O}_{X_i}$-模组成的有界复形 $\mathcal{G}_i^\bullet$，% P09-ERR-0492
它到 $X$ 的拉回是 $\mathcal{G}$。增大 $i$ 后，% P09-ERR-0493
可以假设 $\mathcal{G}_i^n$ 在 $R_i$ 上平坦；见《空间的极限》引理
\ref{spaces-limits-lemma-descend-flat}。
增大 $i$ 后，可以假设 $\mathcal{G}_i^n$ 的支撑在 $R_i$ 上适当；
见《空间的极限》引理
\ref{spaces-limits-lemma-eventually-proper-support}。
最后，由引理 \ref{spaces-perfect-lemma-descend-perfect}，增大 $i$ 后，
可以假设存在完美对象 $E_i$（属于 $D(\mathcal{O}_{X_i})$），
它到 $X$ 的拉回是 $E$。把引理
\ref{spaces-perfect-lemma-compute-ext-perfect} 应用于
$X_i \to \Spec(R_i)$、$E_i$、$\mathcal{G}_i^\bullet$，
并使用已经证明的基变换性质，即得结果。
\end{proof}

\section{完美复形}
\label{spaces-perfect-section-perfect-complexes}

\noindent
先讨论完美复形的 Betti 数之跳跃轨迹。
首先必须定义 Betti 数。

\medskip\noindent
设 $S$ 为概形。设 $X$ 为 $S$ 上的代数空间。
设 $E$ 为 $D(\mathcal{O}_X)$ 的对象。
设 $x \in |X|$。我们要定义
$\beta_i(x) \in \{0, 1, 2, \ldots \} \cup \{\infty\}$。
为此，选取态射 $f : \Spec(k) \to X$，它属于 $x$ 的等价类。
于是 $Lf^*E$ 是 $D(\Spec(k)_\etale, \mathcal{O})$ 的对象。
由《平展上同调》引理
\ref{etale-cohomology-lemma-all-modules-quasi-coherent} 和定理
\ref{etale-cohomology-theorem-quasi-coherent}，得到
$D(\Spec(k)_\etale, \mathcal{O}) = D(k)$，即 $k$-向量空间的导出范畴。
因此 $Lf^*E$ 是 $k$-向量空间的复形，可以取
$\beta_i(x) = \dim_k H^i(Lf^*E)$。容易看出，这不依赖于
$x$ 中代表元的选取。此外，若 $X$ 是概形，这与《概形的导出范畴》第
\ref{perfect-section-perfect-complexes} 节中使用的概念相同。

\begin{lemma}
\label{spaces-perfect-lemma-jump-loci}
设 $S$ 为概形。设 $X$ 为 $S$ 上的代数空间。
设 $E \in D(\mathcal{O}_X)$ 拟凝聚（例如完美）。
对任意 $i \in \mathbf{Z}$，考虑上面定义的函数
$$
\beta_i : |X| \longrightarrow \{0, 1, 2, \ldots\}
$$
。那么
\begin{enumerate}
\item $\beta_i$ 的构造与任意基变换可交换；
\item 函数 $\beta_i$ 上半连续；且
\item $\beta_i$ 的水平集在平展局部可构造。
\end{enumerate}
\end{lemma}

\begin{proof}
选取概形 $U$ 和满射平展态射 $\varphi : U \to X$。
那么 $L\varphi^*E$ 是概形 $U$ 上的拟凝聚复形（使用引理
\ref{spaces-perfect-lemma-descend-pseudo-coherent}），从而可以应用概形情形的结果；
见《概形的导出范畴》引理 \ref{perfect-lemma-jump-loci}。
第 (3) 部分意指这些水平集到 $U$ 的逆像局部可构造；见
《空间的性质》定义
\ref{spaces-properties-definition-locally-constructible}。
\end{proof}

\begin{lemma}
\label{spaces-perfect-lemma-jump-loci-geometric}
设 $Y$ 为概形，并设 $X$ 为 $Y$ 上的代数空间，使得结构态射
$f : X \to Y$ 平坦、适当且有限表现。
设 $\mathcal{F}$ 为有限表现 $\mathcal{O}_X$-模，且在 $Y$ 上平坦。
固定 $i \in \mathbf{Z}$，考虑函数
$$
\beta_i : |Y| \to \{0, 1, 2, \ldots\},\quad
y \longmapsto \dim_{\kappa(y)} H^i(X_y, \mathcal{F}_y)
$$
。那么
\begin{enumerate}
\item $\beta_i$ 的构造与任意基变换可交换；
\item 函数 $\beta_i$ 上半连续；且
\item $\beta_i$ 的水平集在 $Y$ 中局部可构造。
\end{enumerate}
\end{lemma}

\begin{proof}
由上同调与基变换（更确切地说，由引理
\ref{spaces-perfect-lemma-flat-proper-perfect-direct-image-general}），对象
$K = Rf_*\mathcal{F}$ 是 $Y$ 的导出范畴中的完美对象，
并且它的构造与任意基变换可交换。特别地，
$$
H^i(X_y, \mathcal{F}_y) = H^i(K \otimes_{\mathcal{O}_Y}^\mathbf{L} \kappa(y))
$$
。因此本引理由引理 \ref{spaces-perfect-lemma-jump-loci} 得出。
\end{proof}

\begin{lemma}
\label{spaces-perfect-lemma-chi-locally-constant}
设 $S$ 为概形。设 $X$ 为 $S$ 上的代数空间。
设 $E \in D(\mathcal{O}_X)$ 完美。函数
$$
\chi_E : |X| \longrightarrow \mathbf{Z},\quad
x \longmapsto \sum (-1)^i \beta_i(x)
$$
在 $X$ 上局部常值。
\end{lemma}

\begin{proof}
略。提示：
由概形情形通过平展局部化得出。见《概形的导出范畴》引理
\ref{perfect-lemma-chi-locally-constant}。
\end{proof}

\begin{lemma}
\label{spaces-perfect-lemma-open-where-cohomology-in-degree-i-rank-r}
设 $S$ 为概形。设 $X$ 为 $S$ 上的代数空间。
设 $E \in D(\mathcal{O}_X)$ 完美。
给定 $i, r \in \mathbf{Z}$，存在由下列性质刻画的开子空间
$U \subset X$：
\begin{enumerate}
\item $E|_U \cong H^i(E|_U)[-i]$，且 $H^i(E|_U)$ 是局部自由
$\mathcal{O}_U$-模，秩为 $r$；
\item 态射 $f : Y \to X$ 经由 $U$ 分解，当且仅当
$Lf^*E$ 同构于局部自由模，秩为 $r$，并置于次数 $i$。
\end{enumerate}
\end{lemma}

\begin{proof}
略。提示：
由概形情形通过平展局部化得出。见《概形的导出范畴》引理
\ref{perfect-lemma-open-where-cohomology-in-degree-i-rank-r}。
\end{proof}

\begin{lemma}
\label{spaces-perfect-lemma-open-where-cohomology-in-degree-i-rank-r-geometric}
设 $S$ 为概形。
设 $f : X \to Y$ 是 $S$ 上代数空间的态射，且适当、平坦、有限表现。
设 $\mathcal{F}$ 为有限表现 $\mathcal{O}_X$-模，且在 $Y$ 上平坦。
固定 $i, r \in \mathbf{Z}$。那么存在开子空间
$V \subset Y$，具有如下性质：
态射 $T \to Y$ 经由 $V$ 分解，当且仅当
$Rf_{T, *}\mathcal{F}_T$ 同构于有限局部自由模，秩为 $r$，
并置于次数 $i$。
\end{lemma}

\begin{proof}
由上同调与基变换（引理
\ref{spaces-perfect-lemma-flat-proper-perfect-direct-image-general}），对象
$K = Rf_*\mathcal{F}$ 是 $Y$ 的导出范畴中的完美对象，
并且它的构造与任意基变换可交换。因此本引理由引理
\ref{spaces-perfect-lemma-open-where-cohomology-in-degree-i-rank-r} 立即得出。
\end{proof}

\begin{lemma}
\label{spaces-perfect-lemma-locally-closed-where-H0-locally-free}
设 $S$ 为概形。设 $X$ 为 $S$ 上的代数空间。
设 $E \in D(\mathcal{O}_X)$ 完美，且 Tor 振幅在
$[a, b]$ 中，其中 $a, b \in \mathbf{Z}$。
设 $r \geq 0$。
那么存在由下列性质刻画的局部闭子空间 $j : Z \to X$：
\begin{enumerate}
\item $H^a(Lj^*E)$ 是局部自由 $\mathcal{O}_Z$-模，秩为 $r$；且
\item 态射 $f : Y \to X$ 经由 $Z$ 分解，当且仅当对所有态射
$g : Y' \to Y$，$\mathcal{O}_{Y'}$-模
$H^a(L(f \circ g)^*E)$ 都局部自由且秩为 $r$。
\end{enumerate}
此外，$j : Z \to X$ 是有限表现的，并且
\begin{enumerate}
\item[(3)] 若 $f : Y \to X$ 分解为
$Y \xrightarrow{g} Z \to X$，则
$H^a(Lf^*E) = g^*H^a(Lj^*E)$；
\item[(4)] 若 $\beta_a(x) \leq r$ 对所有 $x \in |X|$ 成立，
则 $j$ 是闭浸入；并且给定 $f : Y \to X$，下列条件等价：
\begin{enumerate}
\item $f : Y \to X$ 经由 $Z$ 分解；
\item $H^0(Lf^*E)$ 是局部自由 $\mathcal{O}_Y$-模，秩为 $r$；% P09-ERR-0494
\end{enumerate}
而当 $r = 1$ 时，它们还等价于
\begin{enumerate}
\item[(c)] $\mathcal{O}_Y \to
\SheafHom_{\mathcal{O}_Y}(H^0(Lf^*E), H^0(Lf^*E))$
为单射。
\end{enumerate}
\end{enumerate}
\end{lemma}

\begin{proof}
略。提示：
由概形情形通过平展局部化得出。见《概形的导出范畴》引理
\ref{perfect-lemma-locally-closed-where-H0-locally-free}。
\end{proof}

\begin{lemma}
\label{spaces-perfect-lemma-proper-flat-h0}
设 $S$ 为概形。
设 $f : X \to Y$ 是 $S$ 上代数空间的态射。假设
\begin{enumerate}
\item $f$ 适当、平坦且有限表现；并且
\item 对于一个态射 $\Spec(k) \to Y$，其中 $k$ 是域，有 % P09-ERR-0495
$k = H^0(X_k, \mathcal{O}_{X_k})$。
\end{enumerate}
那么
\begin{enumerate}
\item[(a)] $f_*\mathcal{O}_X = \mathcal{O}_Y$，且这在任意基变换后仍成立；
\item[(b)] 在 $Y$ 的平展局部有
$$
Rf_*\mathcal{O}_X = \mathcal{O}_Y \oplus P
$$
（在 $D(\mathcal{O}_Y)$ 中），其中 $P$ 完美且 Tor 振幅在
$[1, \infty)$ 中。
\end{enumerate}
\end{lemma}

\begin{proof}
只需在 $Y$ 的平展局部证明 (a) 和 (b)，故可以并且确实假设
$Y$ 是仿射概形。
由上同调与基变换
（引理 \ref{spaces-perfect-lemma-flat-proper-perfect-direct-image-general}），复形
$E = Rf_*\mathcal{O}_X$
完美，并且它的构造与任意基变换可交换。
特别地，对 $y \in Y$ 有
$H^0(E \otimes^\mathbf{L} \kappa(y)) =
H^0(X_y, \mathcal{O}_{X_y}) = \kappa(y)$。
因此，用引理 \ref{spaces-perfect-lemma-jump-loci} 的记号，
$\beta_0(y) \leq 1$ 对所有 $y \in Y$ 成立。应用引理
\ref{spaces-perfect-lemma-locally-closed-where-H0-locally-free}，取
$a = 0$ 和 $r = 1$。得到泛闭子概形
$j : Z \to Y$，其中 $H^0(Lj^*E)$ 可逆，并由该引理
(4)(a)、(b)、(c) 的等价性刻画。
由于 $E$ 的构造与基变换可交换，有
$$
Lf^*E = R\text{pr}_{1, *}\mathcal{O}_{X \times_Y X}
$$
。态射 $\text{pr}_1 : X \times_Y X$ 有截面，
即对角态射 $\Delta$（$X$ 在 $Y$ 上）。得到映射
$$
\mathcal{O}_X \longrightarrow R\text{pr}_{1, *}\mathcal{O}_{X \times_Y X}
\longrightarrow \mathcal{O}_X
$$
（在 $D(\mathcal{O}_X)$ 中），其复合是恒等态射。因此
$R\text{pr}_{1, *}\mathcal{O}_{X \times_Y X} = \mathcal{O}_X \oplus E'$
（在 $D(\mathcal{O}_X)$ 中）。从而 $\mathcal{O}_X$ 是
$H^0(Lf^*E)$ 的直和因子；由上述引理 (4)(c) 与 (4)(a) 的等价性，
得出 $X \to Y$ 经由 $Z$ 分解。
由于 $\{X \to Y\}$ 是 fppf 覆盖，有 $Z = Y$。
于是 $f_*\mathcal{O}_X$ 是可逆 $\mathcal{O}_Y$-模。
我们断定 $\mathcal{O}_Y \to f_*\mathcal{O}_X$ 是同构，因为若环同态
$A \to B$ 使 $B$ 作为 $A$-模可逆，则该同态是同构。
由于假设在基变换下保持，(a) 得证。

\medskip\noindent
(b) 的证明。上面已经看到，对每个 $y \in Y$，映射
$\mathcal{O}_Y \to H^0(E \otimes^\mathbf{L} \kappa(y))$ 满射。
因此可以应用《代数进阶》引理
\ref{more-algebra-lemma-better-cut-complex-in-two}，得出在 $y$ 的
一个开邻域上有分解
$Rf_*\mathcal{O}_X = \mathcal{O}_Y \oplus P$。
\end{proof}

\begin{lemma}
\label{spaces-perfect-lemma-proper-flat-geom-red-connected}
设 $S$ 为概形。
设 $f : X \to Y$ 是 $S$ 上代数空间的态射。假设
\begin{enumerate}
\item $f$ 适当、平坦且有限表现；并且
\item $f$ 的几何纤维是既约且连通的。
\end{enumerate}
那么 $f_*\mathcal{O}_X = \mathcal{O}_Y$，且这在任意基变换后仍成立。
\end{lemma}

\begin{proof}
由引理 \ref{spaces-perfect-lemma-proper-flat-h0}，只需证明
$k = H^0(X_k, \mathcal{O}_{X_k})$ 对所有态射
$\Spec(k) \to Y$（其中 $k$ 是域）成立。这由《域上的空间》引理
\ref{spaces-over-fields-lemma-proper-geometrically-reduced-global-sections}
以及 $X_k$ 几何连通且几何既约这一事实得出。
\end{proof}

\section{其他应用}
\label{spaces-perfect-section-other-applications}

\noindent
本节陈述并证明一些可由上述理论推出的结果。

\begin{lemma}
\label{spaces-perfect-lemma-countable-cohomology}
设 $S$ 为概形。设 $X$ 为 $S$ 上拟紧拟分离的代数空间。
设 $K$ 为 $D_\QCoh(\mathcal{O}_X)$ 的对象，并且上同调层
$H^i(K)$ 在平展于 $X$ 的仿射概形上的截面集合可数。
那么对任意拟紧拟分离平展态射 $U \to X$ 和任意完美对象
$E$（属于 $D(\mathcal{O}_X)$），集合
$$
H^i(U, K \otimes^\mathbf{L} E),\quad \Ext^i(E|_U, K|_U)
$$
都可数。
\end{lemma}

\begin{proof}
由《位点上的上同调》引理
\ref{sites-cohomology-lemma-dual-perfect-complex}，
只需对群 $H^i(U, K \otimes^\mathbf{L} E)$ 证明结果。
我们将使用归纳原理证明本引理；见引理
\ref{spaces-perfect-lemma-induction-principle}。

\medskip\noindent
当 $U = \Spec(A)$ 仿射时，结果由概形情形得出；
见《概形的导出范畴》引理
\ref{perfect-lemma-countable-cohomology}。

\medskip\noindent
为完成证明，只需说明：若 $(U \subset W, V \to W)$
是初等特异三角，% P09-ERR-0496; same source-defect class as P09-ERR-0408
并且结果对 $U$、$V$ 和 $U \times_W V$ 成立，则结果对 $W$ 成立。
这立即由 Mayer--Vietoris 序列得出；见引理
\ref{spaces-perfect-lemma-unbounded-mayer-vietoris}。
\end{proof}

\begin{lemma}
\label{spaces-perfect-lemma-countable}
设 $S$ 为概形。
设 $X$ 为 $S$ 上拟紧拟分离的代数空间。
假设 $\mathcal{O}_X$ 在平展于 $X$ 的仿射概形上的截面集合可数。
设 $K$ 为 $D_\QCoh(\mathcal{O}_X)$ 的对象。下列条件等价：
\begin{enumerate}
\item $K = \text{hocolim} E_n$，其中 $E_n$ 是
$D(\mathcal{O}_X)$ 的完美对象；且
\item 上同调层 $H^i(K)$ 在平展于 $X$ 的仿射概形上的截面集合可数。
\end{enumerate}
\end{lemma}

\begin{proof}
若 (1) 成立，则 (2) 成立，因为同伦余极限与取上同调层可交换，% P09-ERR-0497
（由《导出范畴》引理
\ref{derived-lemma-cohomology-of-hocolim}），并且完美复形局部同构于
有限自由 $\mathcal{O}_X$-模的有限复形，所以根据关于 $X$ 的假设满足 (2)。

\medskip\noindent
假设 (2)。
选取 K-内射复形 $\mathcal{K}^\bullet$ 来代表 $K$。
选取完美生成元 $E$（属于 $D_\QCoh(\mathcal{O}_X)$），并用 K-内射复形
$\mathcal{I}^\bullet$ 代表它。根据定理
\ref{spaces-perfect-theorem-DQCoh-is-Ddga} 及其证明，存在三角范畴等价
$F : D_\QCoh(\mathcal{O}_X) \to D(A, \text{d})$，
其中 $(A, \text{d})$ 是微分分次代数
$$
(A, \text{d}) =
\Hom_{\text{Comp}^{dg}(\mathcal{O}_X)}
(\mathcal{I}^\bullet, \mathcal{I}^\bullet)
$$
，且它把 $K$ 映到微分分次模
$$
M = \Hom_{\text{Comp}^{dg}(\mathcal{O}_X)}
(\mathcal{I}^\bullet, \mathcal{K}^\bullet)
$$
。注意 $H^i(A) = \Ext^i(E, E)$ 且
$H^i(M) = \Ext^i(E, K)$。
此外，由于 $F$ 是等价，它及其拟逆都与同伦余极限可交换。
因此，只需把 $M$ 写成 $D(A, \text{d})$ 的紧对象之同伦余极限。
由《微分分次代数》引理 \ref{dga-lemma-countable}，只需证明 % P09-ERR-0498
$\Ext^i(E, E)$ 和 $\Ext^i(E, K)$ 对每个 $i$ 都可数。
这由引理 \ref{spaces-perfect-lemma-countable-cohomology} 得出。
\end{proof}

\begin{lemma}
\label{spaces-perfect-lemma-computing-sections-as-colim}
设 $A$ 为环。设 $f : U \to X$ 是 $A$ 上代数空间之间的
平坦有限表现态射。那么
\begin{enumerate}
\item 存在由完美对象 $L_n$（属于 $D(\mathcal{O}_X)$）构成的逆系统，
使得
$$
R\Gamma(U, Lf^*K) = \text{hocolim}\ R\Hom_X(L_n, K)
$$
在 $D(A)$ 中对 $K$（属于 $D_\QCoh(\mathcal{O}_X)$）函子性成立；且
\item 存在由完美对象 $E_n$（属于 $D(\mathcal{O}_X)$）构成的系统，
使得
$$
R\Gamma(U, Lf^*K) = \text{hocolim}\ R\Gamma(X, E_n \otimes^\mathbf{L} K)
$$
在 $D(A)$ 中对 $K$（属于 $D_\QCoh(\mathcal{O}_X)$）函子性成立。
\end{enumerate}
\end{lemma}

\begin{proof}
由引理 \ref{spaces-perfect-lemma-cohomology-base-change}，
$$
R\Gamma(U, Lf^*K) = R\Gamma(X, Rf_*\mathcal{O}_U \otimes^\mathbf{L} K)
$$
对 $K$ 函子性成立。注意，$R\Gamma(X, -)$ 与同伦余极限可交换，
因为由引理 \ref{spaces-perfect-lemma-quasi-coherence-pushforward-direct-sums}，
它与直和可交换。类似地，$- \otimes^\mathbf{L} K$
与导出余极限可交换，因为 $- \otimes^\mathbf{L} K$ 与直和可交换
（因为 $D(\mathcal{O}_X)$ 中的直和由代表复形的直和给出）。
因此，为证明 (2)，只需把
$Rf_*\mathcal{O}_U = \text{hocolim} E_n$ 写成由完美对象
$E_n$（属于 $D(\mathcal{O}_X)$）组成的系统。
做到这一点后，置 $L_n = E_n^\vee$ 即得 (1)；见《位点上的上同调》引理
\ref{sites-cohomology-lemma-dual-perfect-complex}。

\medskip\noindent
把 $A = \colim A_i$ 写成环 $A_i$ 的余极限，它们
在 $\mathbf{Z}$ 上有限型。由《空间的极限》引理
\ref{spaces-limits-lemma-descend-finite-presentation}，
可以找到 $i$ 和有限表现态射
$U_i \to X_i \to \Spec(A_i)$，其到 $\Spec(A)$ 的基变换恢复
$U \to X \to \Spec(A)$。
增大 $i$ 后，可以假设 $f_i : U_i \to X_i$ 平坦；见《空间的极限》引理
\ref{spaces-limits-lemma-descend-flat}。
由引理 \ref{spaces-perfect-lemma-compare-base-change}，对象
$Rf_{i, *}\mathcal{O}_{U_i}$ 沿 $g : X \to X_i$ 的导出拉回等于
$Rf_*\mathcal{O}_U$。
由于 $Lg^*$ 与导出余极限可交换，只需对 $f_i$ 证明所需结论。
因此可以假设 $U$ 和 $X$ 在 $\mathbf{Z}$ 上有限型。

\medskip\noindent
假设 $f : U \to X$ 是 $\mathbf{Z}$ 上有限型代数空间的态射。
为完成证明，我们将说明 $Rf_*\mathcal{O}_U$ 是完美复形的同伦余极限。
为此应用引理 \ref{spaces-perfect-lemma-countable}。所以只需证明
$R^if_*\mathcal{O}_U$ 在平展于 $X$ 的仿射概形上的截面集合可数。
这由引理 \ref{spaces-perfect-lemma-countable-cohomology} 应用于结构层得出。
\end{proof}









\section{分解性质}
\label{spaces-perfect-section-resolution-property}

\noindent
本节是《概形的导出范畴》第
\ref{perfect-section-resolution-property} 节在代数空间情形的对应版本；
请先阅读该节。目前尚不清楚域上的光滑固有代数空间是否总具有分解性质，
抑或并非如此。如知此问题的答案，请发送邮件至
\href{mailto:stacks.project@gmail.com}{stacks.project@gmail.com}。

\medskip\noindent
尽管对一般代数空间考虑如下概念几乎没有意义，我们仍可作出这个定义。

\begin{definition}
\label{spaces-perfect-definition-resolution-property}
设 $S$ 为概形，$X$ 为 $S$ 上的代数空间。称 $X$ 具有
{\it 分解性质}，如果每个有限型拟凝聚 $\mathcal{O}_X$-模
都是某个有限局部自由 $\mathcal{O}_X$-模的商。
\end{definition}

\noindent
若 $X$ 是拟紧拟分离代数空间，则只需检验每个有限表示的
$\mathcal{O}_X$-模（自动拟凝聚）都是某个有限局部自由
$\mathcal{O}_X$-模的商；见《空间的极限》引理 % P09-ERR-0499
\ref{spaces-limits-lemma-finite-directed-colimit-surjective-maps}。
若 $X$ 是诺特代数空间，则有限型拟凝聚模恰为凝聚
$\mathcal{O}_X$-模；见《空间的上同调》引理
\ref{spaces-cohomology-lemma-coherent-Noetherian}。

\begin{lemma}
\label{spaces-perfect-lemma-resolution-property-ample-relative}
设 $S$ 为概形。设 $f : X \to Y$ 是 $S$ 上代数空间之间的态射。假设
\begin{enumerate}
\item $Y$ 拟紧、拟分离且具有分解性质；
\item 存在一个 $f$-丰沛可逆模，其定义在 $X$ 上
（《空间上的除子》定义
\ref{spaces-divisors-definition-relatively-ample}）。
\end{enumerate}
则 $X$ 具有分解性质。
\end{lemma}

\begin{proof}
令 $\mathcal{F}$ 为有限型拟凝聚 $\mathcal{O}_X$-模，
$\mathcal{L}$ 为 $f$-丰沛可逆模。取仿射概形 $V$ 以及满平展态射
$V \to Y$。置 $U = V \times_Y X$，则 $\mathcal{L}|_U$ 在 $U$ 上丰沛。
由《性质》命题
\ref{properties-proposition-characterize-ample}，存在有限多个映射 % P09-ERR-0500
$s_i : \mathcal{L}^{\otimes n_i}|_U \to \mathcal{F}|_U$，
它们联合为满。考虑拟凝聚 $\mathcal{O}_Y$-模
$$
\mathcal{H}_n =
f_*(\mathcal{F} \otimes_{\mathcal{O}_X} \mathcal{L}^{\otimes n})
$$
可将 $s_i$ 视为 $V$ 上的一个截面；这里所取的层是
$\mathcal{H}_{-n_i}$。
假设可以找到有限局部自由 $\mathcal{O}_Y$-模 $\mathcal{E}_i$ 和映射
$\mathcal{E}_i \to \mathcal{H}_{-n_i}$，使 $s_i$ 属于其像。
那么相应的映射
$$
f^*\mathcal{E}_i
\otimes_{\mathcal{O}_X}
\mathcal{L}^{\otimes n_i} \longrightarrow \mathcal{F}
$$
将联合为满，于是引理得证。由《空间的极限》引理
\ref{spaces-limits-lemma-directed-colimit-finite-type}，
对每个 $i$ 都能找到有限型拟凝聚子模
$\mathcal{H}'_i \subset  \mathcal{H}_{-n_i}$，
它包含截面 $s_i$（定义在 $V$ 上）。
因此，$Y$ 的分解性质给出满射
$\mathcal{E}_i \to \mathcal{H}'_i$，结论随即成立。
\end{proof}

\begin{lemma}
\label{spaces-perfect-lemma-resolution-property-goes-up-affine}
设 $S$ 为概形。设 $f : X \to Y$ 是 $S$ 上代数空间之间的仿射或拟仿射
态射，并且 $Y$ 拟紧拟分离。若 $Y$ 具有分解性质，则 $X$ 亦然。
\end{lemma}

\begin{proof}
由《空间上的除子》引理
\ref{spaces-divisors-lemma-quasi-affine-O-ample}，
这是引理 \ref{spaces-perfect-lemma-resolution-property-ample-relative} 的特殊情形。
\end{proof}

\noindent
下面给出一个可以证明分解性质向下传递的情形。

\begin{lemma}
\label{spaces-perfect-lemma-resolution-property-goes-down-finite-flat}
设 $S$ 为概形。设 $f : X \to Y$ 是 $S$ 上代数空间之间的满有限局部自由
态射。若 $X$ 具有分解性质，则 $Y$ 亦然。
\end{lemma}

\begin{proof}
该条件意味着 $f$ 仿射，且 $f_*\mathcal{O}_X$ 是秩为正的有限局部自由
$\mathcal{O}_Y$-模。令 $\mathcal{G}$ 为有限型拟凝聚
$\mathcal{O}_Y$-模。由假设，存在满射
$\mathcal{E} \to f^*\mathcal{G}$，其源是有限局部自由
$\mathcal{O}_X$-模 $\mathcal{E}$。由于 $f_*$ 正合
（《空间的上同调》第 \ref{spaces-cohomology-section-finite-morphisms} 节），
我们得到满射
$$
f_*\mathcal{E}
\longrightarrow
f_*f^*\mathcal{G} = \mathcal{G} \otimes_{\mathcal{O}_Y} f_*\mathcal{O}_X
$$
取对偶后，得到满射 % P09-ERR-0501
$$
f_*\mathcal{E}
\otimes_{\mathcal{O}_Y}
\SheafHom_{\mathcal{O}_Y}(f_*\mathcal{O}_X, \mathcal{O}_Y)
\longrightarrow
\mathcal{G}
$$
由于 $f_*\mathcal{E}$ 有限局部自由，结论成立。
\end{proof}

\noindent
关于代数空间分解性质的更多内容，见《空间态射进阶》第
\ref{spaces-more-morphisms-section-resolution-property} 节。













\section{判定有界性}
\label{spaces-perfect-section-detecting-boundedness}

\noindent
本节将说明，第 \ref{spaces-perfect-section-generating} 节所构造的拟紧拟分离概形之
$D_\QCoh$ 的紧生成元具有一种特殊性质。% P09-ERR-0502
建议先阅读该节，因为它与本节非常相似。

\begin{lemma}
\label{spaces-perfect-lemma-bounded-truncation}
设 $S$ 为概形。设 $X$ 为 $S$ 上的拟紧拟分离代数空间。
设 $P \in D_{perf}(\mathcal{O}_X)$ 且
$E \in D_{\QCoh}(\mathcal{O}_X)$。设 $a \in \mathbf{Z}$。
以下条件等价：
\begin{enumerate}
\item $\Hom_{D(\mathcal{O}_X)}(P[-i], E) = 0$
对 $i \gg 0$ 成立；且
\item $\Hom_{D(\mathcal{O}_X)}(P[-i], \tau_{\geq a} E) = 0$
对 $i \gg 0$ 成立。
\end{enumerate}
\end{lemma}

\begin{proof}
利用三角 $ \tau_{< a} E \to E \to \tau_{\geq a} E \to$，
可知只要能证明
$$
\Hom_{D(\mathcal{O}_X)}(P[-i], \tau_{< a} E) =
\Hom_{D(\mathcal{O}_X)}(P, (\tau_{< a} E)[i]) = 0
$$
对 $i \gg 0$ 成立，所述等价性就成立。由于 $P$ 完美，
这由引理 \ref{spaces-perfect-lemma-ext-from-perfect-into-bounded-QCoh} 得出。
\end{proof}

\begin{lemma}
\label{spaces-perfect-lemma-bounded-below-truncation}
设 $S$ 为概形。设 $X$ 为 $S$ 上的拟紧拟分离代数空间。
设 $P \in D_{perf}(\mathcal{O}_X)$ 且
$E \in D_{\QCoh}(\mathcal{O}_X)$。设 $a \in \mathbf{Z}$。
以下条件等价：
\begin{enumerate}
\item $\Hom_{D(\mathcal{O}_X)}(P[-i], E) = 0$
对 $i \ll 0$ 成立；且
\item $\Hom_{D(\mathcal{O}_X)}(P[-i], \tau_{\leq a} E) = 0$
对 $i \ll 0$ 成立。
\end{enumerate}
\end{lemma}

\begin{proof}
利用三角 $ \tau_{\leq a} E \to E \to \tau_{> a} E \to$，
可知只要能证明
$$
\Hom_{D(\mathcal{O}_X)}(P[-i], \tau_{> a} E) =
\Hom_{D(\mathcal{O}_X)}(P, (\tau_{> a} E)[i]) = 0
$$
对 $i \ll 0$ 成立，所述等价性就成立。由于 $P$ 完美，
这由引理 \ref{spaces-perfect-lemma-ext-from-perfect-into-bounded-QCoh} 得出。
\end{proof}

\begin{proposition}
\label{spaces-perfect-proposition-detecting-bounded-above}
设 $S$ 为概形。设 $X$ 为 $S$ 上的拟紧拟分离代数空间。
设 $G \in D_{perf}(\mathcal{O}_X)$ 是生成
$D_\QCoh (\mathcal{O}_X)$ 的完美复形。设
$E \in D_\QCoh (\mathcal{O}_X)$。以下条件等价：
\begin{enumerate}
\item $E \in D^-_\QCoh (\mathcal{O}_X)$；
\item $\Hom_{D(\mathcal{O}_X)}(G[-i], E) = 0$
对 $i \gg 0$ 成立；
\item $\Ext^i_X(G, E) = 0$ 对 $i \gg 0$ 成立；
\item $R\Hom_X(G, E)$ 属于 $D^-(\mathbf{Z})$；
\item $H^i(X, G^\vee \otimes_{\mathcal{O}_X}^\mathbf{L} E) = 0$
对 $i \gg 0$ 成立；
\item $R\Gamma(X, G^\vee \otimes_{\mathcal{O}_X}^\mathbf{L} E)$
属于 $D^-(\mathbf{Z})$；
\item 对每个完美对象 $P$（属于 $D(\mathcal{O}_X)$），
\begin{enumerate}
\item 将 $G$ 换成 $P$ 后，断言 (2)、(3)、(4) 成立；且
\item $H^i(X, P \otimes_{\mathcal{O}_X}^\mathbf{L} E) = 0$
对 $i \gg 0$ 成立；
\item $R\Gamma(X, P \otimes_{\mathcal{O}_X}^\mathbf{L} E)$
属于 $D^-(\mathbf{Z})$。
\end{enumerate}
\end{enumerate}
\end{proposition}

\begin{proof}
假设 (1) 成立。由于
$\Hom_{D(\mathcal{O}_X)}(G[-i], E) = \Hom_{D(\mathcal{O}_X)}(G, E[i])$，
由引理 \ref{spaces-perfect-lemma-ext-from-perfect-into-bounded-QCoh}，
这在 $i \gg 0$ 时为零。这证明了 (1) 推出 (2)。

\medskip\noindent
由《位点上的上同调》第 \ref{sites-cohomology-section-global-RHom} 节
的讨论，(2)、(3)、(4) 等价。按定义，
$H^i(X, -) = H^i(R\Gamma(X, -))$，故 (5) 和 (6) 等价。
由《位点上的上同调》引理
\ref{sites-cohomology-lemma-dual-perfect-complex} 以及等式
$(G[-i])^\vee = G^\vee[i]$，等价条件 (2)、(3)、(4)
又等价于等价条件 (5)、(6)。

\medskip\noindent
显然 (7) 推出 (2)。反过来，我们证明等价条件 (2)--(6) 推出 (7)。
回顾注 \ref{spaces-perfect-remark-classical-generator}：$G$ 是
$D_{perf}(\mathcal{O}_X)$ 的经典生成元。对
$P \in D_{perf}(\mathcal{O}_X)$，令 $T(P)$ 表示断言
$R\Hom_X(P, E)$ 属于 $D^-(\mathbf{Z})$。
显然，$T$ 对直和、直和项和移位封闭，并且对特异三角满足两出三性质。
因此，由《导出范畴》注 \ref{derived-remark-check-on-generator}，
(4) 推出 $T$ 在整个 $D_{perf}(\mathcal{O}_X)$ 上成立。
同一论证适用于所有其他性质；但对性质 (7)(b) 和 (7)(c)，% P09-ERR-0503
还要用到 $P \mapsto P^\vee$ 是 $D_{perf}(\mathcal{O}_X)$ 的自等价。
略去这个小细节。

\medskip\noindent
我们将利用引理 \ref{spaces-perfect-lemma-induction-principle} 的归纳原理，
证明等价条件 (2)--(7) 推出 (1)。

\medskip\noindent
首先，我们证明 (2)--(7) $\Rightarrow$ (1) 在 $X$ 仿射时成立。
这由概形的情形得出；见《概形的导出范畴》命题
\ref{perfect-proposition-detecting-bounded-above}。

\medskip\noindent
现在假设 $(U \subset X, j : V \to X)$ 是 $S$ 上拟紧拟分离
代数空间的初等特异方形，并且已知蕴涵 (2)--(7) $\Rightarrow$ (1)
对 $U$、$V$ 和 $U \times_X V$ 成立。
为完成证明，必须说明蕴涵 (2)--(7) $\Rightarrow$ (1) 对 $X$ 成立。设
$E \in D_\QCoh(\mathcal{O}_X)$ 满足 (2)--(7)。
由引理 \ref{spaces-perfect-lemma-direct-summand-of-a-restriction} 和定理
\ref{spaces-perfect-theorem-bondal-van-den-Bergh}，存在完美复形 $Q$（在 $X$ 上），
使 $Q|_U$ 生成 $D_\QCoh (\mathcal{O}_U)$。

\medskip\noindent
写 $V = \Spec(A)$。令 $Z \subset V$ 为 $X \setminus U$ 的逆像
所对应的约化闭子概形；它同构地映到该逆像
（见定义 \ref{spaces-perfect-definition-elementary-distinguished-square}）。
这是 $V$ 的逆紧闭子集。取 $f_1, \ldots, f_r \in A$，使得
$Z = V(f_1, \ldots, f_r)$。令 $K \in D(\mathcal{O}_V)$ 为
由 $f_1, \ldots, f_r$ 在 $A$ 上的 Koszul 复形所对应的完美对象。
注意，由于 $K$ 支撑在 $Z$ 上，其正像
$K' = Rj_*K$ 是 $D(\mathcal{O}_X)$ 的完美对象，并且它限制到
$V$ 后是 $K$（见引理 \ref{spaces-perfect-lemma-pushforward-perfect} 和
\ref{spaces-perfect-lemma-pushforward-with-support-in-open}）。
由假设，$R\Hom_{\mathcal{O}_X}(Q, E)$ 和
$R\Hom_{\mathcal{O}_X}(K', E)$ 都有上界。

\medskip\noindent
由引理 \ref{spaces-perfect-lemma-pushforward-with-support-in-open}，有
$K' =  j_!K$，因而
$$
\Hom_{D(\mathcal{O}_X)}(K'[-i], E) = \Hom_{D(\mathcal{O}_V)}(K[-i], E|_V) = 0
$$
对 $i \gg 0$ 成立。因此，可将《概形的导出范畴》引理
\ref{perfect-lemma-orthogonal-koszul-first-variant} 应用于 $E|_V$，
得到整数 $a$，使得
$\tau_{\geq a}(E|_V) =
\tau_{\geq a} R (U \times_X V \to V)_* (E|_{U \times_X V})$。
于是 $\tau_{\geq a} E = \tau_{\geq a} R (U \to X)_* (E |_U)$
（限制到 $U$ 和 $V$ 后检验典范映射是同构）。
因而两次使用引理 \ref{spaces-perfect-lemma-bounded-truncation} 可得
$$
\Hom_{D(\mathcal{O}_U)}(Q|_U [-i], E|_U) =
\Hom_{D(\mathcal{O}_X)}(Q[-i], R (U \to X)_* (E|_U)) = 0
$$
对 $i \gg 0$ 成立。因为该命题对 $U$ 和生成元 $Q|_U$ 成立，
故 $E|_U \in D^-_\QCoh(\mathcal{O}_U)$。又因为函子
$R (U \to X)_*$ 保持 $D^-_\QCoh$
（由引理 \ref{spaces-perfect-lemma-quasi-coherence-direct-image}），所以
$\tau_{\geq a}E \in D^-_\QCoh(\mathcal{O}_X)$。因此
$E \in D^-_\QCoh (\mathcal{O}_X)$。
\end{proof}

\begin{proposition}
\label{spaces-perfect-proposition-detecting-bounded-below}
设 $S$ 为概形。设 $X$ 为 $S$ 上的拟紧拟分离代数空间。
设 $G \in D_{perf}(\mathcal{O}_X)$ 是生成
$D_\QCoh (\mathcal{O}_X)$ 的完美复形。设
$E \in D_\QCoh (\mathcal{O}_X)$。以下条件等价：
\begin{enumerate}
\item $E \in D^+_\QCoh (\mathcal{O}_X)$；
\item $\Hom_{D(\mathcal{O}_X)}(G[-i], E) = 0$
对 $i \ll 0$ 成立；
\item $\Ext^i_X(G, E) = 0$ 对 $i \ll 0$ 成立；
\item $R\Hom_X(G, E)$ 属于 $D^+(\mathbf{Z})$；
\item $H^i(X, G^\vee \otimes_{\mathcal{O}_X}^\mathbf{L} E) = 0$
对 $i \ll 0$ 成立；
\item $R\Gamma(X, G^\vee \otimes_{\mathcal{O}_X}^\mathbf{L} E)$
属于 $D^+(\mathbf{Z})$；
\item 对每个完美对象 $P$（属于 $D(\mathcal{O}_X)$），
\begin{enumerate}
\item 将 $G$ 换成 $P$ 后，断言 (2)、(3)、(4) 成立；且
\item $H^i(X, P \otimes_{\mathcal{O}_X}^\mathbf{L} E) = 0$
对 $i \ll 0$ 成立；
\item $R\Gamma(X, P \otimes_{\mathcal{O}_X}^\mathbf{L} E)$
属于 $D^+(\mathbf{Z})$。
\end{enumerate}
\end{enumerate}
\end{proposition}

\begin{proof}
假设 (1) 成立。由于
$\Hom_{D(\mathcal{O}_X)}(G[-i], E) = \Hom_{D(\mathcal{O}_X)}(G, E[i])$，
由引理 \ref{spaces-perfect-lemma-ext-from-perfect-into-bounded-QCoh}，
这在 $i \ll 0$ 时为零。这证明了 (1) 推出 (2)。

\medskip\noindent
由《位点上的上同调》第 \ref{sites-cohomology-section-global-RHom} 节
的讨论，(2)、(3)、(4) 等价。按定义，
$H^i(X, -) = H^i(R\Gamma(X, -))$，故 (5) 和 (6) 等价。
由《位点上的上同调》引理
\ref{sites-cohomology-lemma-dual-perfect-complex} 以及等式
$(G[-i])^\vee = G^\vee[i]$，等价条件 (2)、(3)、(4)
又等价于等价条件 (5)、(6)。

\medskip\noindent
显然 (7) 推出 (2)。反过来，我们证明等价条件 (2)--(6) 推出 (7)。
回顾注 \ref{spaces-perfect-remark-classical-generator}：$G$ 是
$D_{perf}(\mathcal{O}_X)$ 的经典生成元。对
$P \in D_{perf}(\mathcal{O}_X)$，令 $T(P)$ 表示断言
$R\Hom_X(P, E)$ 属于 $D^+(\mathbf{Z})$。
显然，$T$ 对直和、直和项和移位封闭，并且对特异三角满足两出三性质。
因此，由《导出范畴》注 \ref{derived-remark-check-on-generator}，
(4) 推出 $T$ 在整个 $D_{perf}(\mathcal{O}_X)$ 上成立。
同一论证适用于所有其他性质；但对性质 (7)(b) 和 (7)(c)，% P09-ERR-0503
还要用到 $P \mapsto P^\vee$ 是 $D_{perf}(\mathcal{O}_X)$ 的自等价。
略去这个小细节。

\medskip\noindent
我们将利用引理 \ref{spaces-perfect-lemma-induction-principle} 的归纳原理，
证明等价条件 (2)--(7) 推出 (1)。

\medskip\noindent
首先，我们证明 (2)--(7) $\Rightarrow$ (1) 在 $X$ 仿射时成立。
这由概形的情形得出；见《概形的导出范畴》命题
\ref{perfect-proposition-detecting-bounded-below}。

\medskip\noindent
现在假设 $(U \subset X, j : V \to X)$ 是 $S$ 上拟紧拟分离
代数空间的初等特异方形，并且已知蕴涵 (2)--(7) $\Rightarrow$ (1)
对 $U$、$V$ 和 $U \times_X V$ 成立。
为完成证明，必须说明蕴涵 (2)--(7) $\Rightarrow$ (1) 对 $X$ 成立。设
$E \in D_\QCoh(\mathcal{O}_X)$ 满足 (2)--(7)。
由引理 \ref{spaces-perfect-lemma-direct-summand-of-a-restriction} 和定理
\ref{spaces-perfect-theorem-bondal-van-den-Bergh}，存在完美复形 $Q$（在 $X$ 上），
使 $Q|_U$ 生成 $D_\QCoh (\mathcal{O}_U)$。

\medskip\noindent
写 $V = \Spec(A)$。令 $Z \subset V$ 为 $X \setminus U$ 的逆像
所对应的约化闭子概形；它同构地映到该逆像
（见定义 \ref{spaces-perfect-definition-elementary-distinguished-square}）。
这是 $V$ 的逆紧闭子集。取 $f_1, \ldots, f_r \in A$，使得
$Z = V(f_1, \ldots, f_r)$。令 $K \in D(\mathcal{O}_V)$ 为
由 $f_1, \ldots, f_r$ 在 $A$ 上的 Koszul 复形所对应的完美对象。
注意，由于 $K$ 支撑在 $Z$ 上，其正像
$K' = Rj_*K$ 是 $D(\mathcal{O}_X)$ 的完美对象，并且它限制到
$V$ 后是 $K$（见引理 \ref{spaces-perfect-lemma-pushforward-perfect} 和
\ref{spaces-perfect-lemma-pushforward-with-support-in-open}）。
由假设，$R\Hom_{\mathcal{O}_X}(Q, E)$ 和
$R\Hom_{\mathcal{O}_X}(K', E)$ 都有下界。

\medskip\noindent
由引理 \ref{spaces-perfect-lemma-pushforward-with-support-in-open}，有
$K' =  j_!K$，因而
$$
\Hom_{D(\mathcal{O}_X)}(K'[-i], E) = \Hom_{D(\mathcal{O}_V)}(K[-i], E|_V) = 0
$$
对 $i \ll 0$ 成立。因此，可将《概形的导出范畴》引理
\ref{perfect-lemma-orthogonal-koszul-second-variant} 应用于 $E|_V$，
得到整数 $a$，使得
$\tau_{\leq a}(E|_V) =
\tau_{\leq a} R (U \times_X V \to V)_* (E|_{U \times_X V})$。
于是 $\tau_{\leq a} E = \tau_{\leq a} R (U \to X)_* (E |_U)$
（限制到 $U$ 和 $V$ 后检验典范映射是同构）。
因而两次使用引理 \ref{spaces-perfect-lemma-bounded-below-truncation} 可得
$$
\Hom_{D(\mathcal{O}_U)}(Q|_U [-i], E|_U) =
\Hom_{D(\mathcal{O}_X)}(Q[-i], R (U \to X)_* (E|_U)) = 0
$$
对 $i \ll 0$ 成立。因为该命题对 $U$ 和生成元 $Q|_U$ 成立，
故 $E|_U \in D^+_\QCoh(\mathcal{O}_U)$。又因为函子
$R (U \to X)_*$ 保持 $D^+_\QCoh$
（由引理 \ref{spaces-perfect-lemma-quasi-coherence-direct-image}），所以
$\tau_{\leq a}E \in D^+_\QCoh(\mathcal{O}_X)$。因此
$E \in D^+_\QCoh (\mathcal{O}_X)$。
\end{proof}











\section{导出范畴中的拟凝聚对象}
\label{spaces-perfect-section-QC}

\noindent
设 $S$ 为概形，$X$ 为 $S$ 上的代数空间。回顾，
$X_{affine, \etale}$ 表示 $X_\etale$ 的仿射对象所成的范畴，
其拓扑由标准平展覆盖给出；见《空间的性质》定义
\ref{spaces-properties-definition-affine-etale-site}。
提醒读者，$X_{affine, \etale}$ 的拓扑斯就是 $X$ 的小平展拓扑斯；
见《空间的性质》引理
\ref{spaces-properties-lemma-alternative}。位点 $X_\etale$
带有结构层 $\mathcal{O}_X$；它限制到 $X_{affine, \etale}$ 后，
我们仍将其记作 $\mathcal{O}_X$。于是存在带环拓扑斯的等价
$$
(\Sh(X_{affine, \etale}), \mathcal{O}_X)
\longrightarrow
(\Sh(X_\etale), \mathcal{O}_X)
$$
见《空间上的下降》公式
(\ref{spaces-descent-equation-alternative-small-ringed})
以及《空间上的下降》第
\ref{spaces-descent-section-alternative-quasi-coherent} 节中的讨论。

\medskip\noindent
本节以 $X_{affine}$ 表示 $X_{affine, \etale}$ 的底层范畴，
并赋予它混沌拓扑，即层与预层相同的拓扑。特别地，结构层
$\mathcal{O}_X$ 也成为 $X_{affine}$ 上的层。
我们得到带环位点的态射 % P09-ERR-0504
$$
\epsilon :
(X_{affine, \etale}, \mathcal{O}_X)
\longrightarrow
(X_{affine}, \mathcal{O}_X)
$$
，如《位点上的上同调》第 \ref{sites-cohomology-section-compare} 节。
本节将把 $D_\QCoh(\mathcal{O}_X)$ 等同于
《位点上的上同调》第 \ref{sites-cohomology-section-modules-cohomology} 节
中引入的范畴 $\mathit{QC}(X_{affine}, \mathcal{O}_X)$。

\begin{lemma}
\label{spaces-perfect-lemma-DQCoh-alternative-small}
在上述情形中，下列三角范畴之间存在典范正合等价：
\begin{enumerate}
\item $D_\QCoh(\mathcal{O}_X)$；
\item $D_\QCoh(X_{affine, \etale}, \mathcal{O}_X)$；
\item $D_\QCoh(X_{affine}, \mathcal{O}_X)$；以及
\item $\mathit{QC}(X_{affine}, \mathcal{O}_X)$。
\end{enumerate}
\end{lemma}

\begin{proof}
若 $U \to V \to X$ 是平展态射且 $U$ 和 $V$ 仿射，
则环同态 $\mathcal{O}_X(V) \to \mathcal{O}_X(U)$ 平坦。
因此，(3) 与 (4) 之间的等价是《位点上的上同调》引理
\ref{sites-cohomology-lemma-cartesian-flat-quasi-coherent} 的特殊情形
（其证明也阐明了该陈述）。

\medskip\noindent
引理之前的讨论表明，我们有带环拓扑斯的等价
$(\Sh(X_{affine, \etale}), \mathcal{O}_X) \to (\Sh(X_\etale), \mathcal{O}_X)$，
从而有模的阿贝尔范畴之间的等价。由于拟凝聚模的概念是内蕴的
（《位点上的模》引理 \ref{sites-modules-lemma-special-locally-free}），
可知该等价保持拟凝聚模的子范畴。因而我们得到 (1) 与 (2)
中三角范畴之间的典范正合等价。

\medskip\noindent
为得到 (2) 与 (3) 中三角范畴之间的正合等价，我们把《位点上的上同调》
引理 \ref{sites-cohomology-lemma-compare-topologies-derived-adequate-modules}
应用于上面的态射 $\epsilon : (X_{affine, \etale}, \mathcal{O}_X) \to
(X_{affine}, \mathcal{O}_X)$。取
$\mathcal{B} = \Ob(X_{affine})$，并取
$\mathcal{A} \subset \textit{PMod}(X_{affine}, \mathcal{O})$
为满足下述条件的预层 $\mathcal{F}$ 所成的全子范畴：
$\mathcal{F}(V) \otimes_{\mathcal{O}_X(V)} \mathcal{O}_X(U)
\to \mathcal{F}(U)$ 是同构。注意，由《空间上的下降》引理
\ref{spaces-descent-lemma-quasi-coherent-alternative-small}，
$\mathcal{A}$ 的对象恰为平展拓扑中那些在
$(X_{affine, \etale}, \mathcal{O}_X)$ 上是拟凝聚模的层。
另一方面，由《位点上的模》引理
\ref{sites-modules-lemma-chaotic-quasi-coherent}，
$\mathcal{A}$ 的对象恰为 $(X_{affine}, \mathcal{O}_X)$ 上，
即混沌拓扑中的拟凝聚模。因此，若能说明《位点上的上同调》引理
\ref{sites-cohomology-lemma-compare-topologies-derived-adequate-modules}
适用，那么使用 $\epsilon^*$ 和 $R\epsilon_*$，确实得到
类别 (2) 与 (3) 之间的典范等价。

\medskip\noindent
我们必须验证四个条件：
\begin{enumerate}
\item $\mathcal{A}$ 的每个对象都是平展拓扑的层。这一点已在上面看到。
\item $\mathcal{A}$ 是
$\textit{Mod}(X_{affine, \etale}, \mathcal{O}_X)$ 的弱 Serre 子范畴。
上面已经看到
$\mathcal{A} = \QCoh(X_{affine, \etale}, \mathcal{O}_X)$，
并且也已看到，在等价
$\textit{Mod}(X_{affine, \etale}, \mathcal{O}) =
\textit{Mod}(X_\etale, \mathcal{O}_X)$ 下，它们对应于 $X$ 上的拟凝聚模。
所以结论由《空间的性质》引理
\ref{spaces-properties-lemma-properties-quasi-coherent} 和《同调》引理
\ref{homology-lemma-characterize-weak-serre-subcategory} 得出。% P09-ERR-0505
\item $X_{affine}$ 的每个对象在混沌拓扑中都有一个覆盖，其成员属于
$\mathcal{B}$。这是因为 $\mathcal{B}$ 包含所有对象。
\item 对每个对象 $U$（作为 $X_{affine}$ 的对象）和每个
$\mathcal{F}$（属于 $\mathcal{A}$），有
$H^p_{Zar}(U, \mathcal{F}) = 0$，其中 $p > 0$。这由仿射概形上拟凝聚模的上同调消失得出；
见《空间的上同调》第 \ref{spaces-cohomology-section-higher-direct-image} 节
中的讨论以及《概形的上同调》引理
\ref{coherent-lemma-quasi-coherent-affine-cohomology-zero}。
\end{enumerate}
证明完成。
\end{proof}

\begin{remark}
\label{spaces-perfect-remark-QC-compare}
设 $S$ 为概形，$X$ 为 $S$ 上的代数空间。稍后我们将证明
$\mathit{QC}((\textit{Aff}/X), \mathcal{O})$
也典范等价于 $D_\QCoh(\mathcal{O}_X)$。
见《栈上的层》命题 \ref{stacks-sheaves-proposition-QC-compare}。
\end{remark}
% Generated mechanically from the frozen spaces-more-morphisms.tex target.
% Do not edit this generated file; edit the bound locale source instead.

% OK, start here.
%

\setcounter{chapter}{75}
\chapter{空间态射进阶}



\phantomsection
\label{spaces-more-morphisms-section-phantom}


\section{引言}
\label{spaces-more-morphisms-section-introduction}

\noindent
本章继续研究代数空间态射的性质。一个基本参考文献是 \cite{Kn}。





\section{约定}
\label{spaces-more-morphisms-section-conventions}

\noindent
我们始终假设所有概形均包含在一个大 fppf 位点 $\Sch_{fppf}$ 中，
并且所考虑的所有环 $A$ 都满足：$\Spec(A)$（同构于）该大位点的一个对象。

\medskip\noindent
设 $S$ 为概形，$X$ 为 $S$ 上的代数空间。在本章及下一章中，
我们将以 $X \times_S X$ 表示 $X$ 与自身的乘积（在 $S$ 上代数空间的范畴中），
而不用 $X \times X$。







\section{根态射}
\label{spaces-more-morphisms-section-radicial}

\noindent
与概形态射的情形不同，根态射并不等同于普遍单射态射。
事实上，前者略强一些。

\begin{definition}
\label{spaces-more-morphisms-definition-radicial}
设 $S$ 为概形，$f : X \to Y$ 为 $S$ 上代数空间的态射。
若 $f$ 满足下述条件，则称其为{\it 根态射}：对任意态射
$\Spec(K) \to Y$（其中 $K$ 为域），约化
$(\Spec(K) \times_Y X)_{red}$ 或为空，或可由 $K$ 的某个
纯不可分域扩张之谱表示。
\end{definition}

\begin{lemma}
\label{spaces-more-morphisms-lemma-radicial-implies-universally-injective}
代数空间的根态射是普遍单射的。
\end{lemma}

\begin{proof}
设 $S$ 为概形，$f : X \to Y$ 为 $S$ 上代数空间的根态射。
由定义显然可知，给定态射 $\Spec(K) \to Y$，它至多有一个
到 $X$ 的提升。因此，由《空间的态射》引理
\ref{spaces-morphisms-lemma-universally-injective}，
$f$ 是普遍单射的。
\end{proof}

\begin{example}
\label{spaces-more-morphisms-example-universally-injective-not-radicial}
普遍单射与根态不再等价。例如，《空间》例
\ref{spaces-example-Qbar} 中的态射
$$
X = [\Spec(\overline{\mathbf{Q}})/
\text{Gal}(\overline{\mathbf{Q}}/\mathbf{Q})]
\longrightarrow
S = \Spec(\mathbf{Q})
$$
是普遍单射的，但不是上述意义下的根态射。
\end{example}

\noindent
不过，逆向蕴含在许多情形下仍然成立。

\begin{lemma}
\label{spaces-more-morphisms-lemma-when-universally-injective-radicial}
设 $S$ 为概形，$f : X \to Y$ 为 $S$ 上代数空间的普遍单射态射。
\begin{enumerate}
\item 若 $f$ 良态，则 $f$ 为根态射。
\item 若 $f$ 拟分离，则 $f$ 为根态射。
\item 若 $f$ 局部分离，则 $f$ 为根态射。
\end{enumerate}
\end{lemma}

\begin{proof}
设 $\mathcal{P}$ 为代数空间态射的一个性质，它在基变换与复合下稳定，
且对闭浸入成立。假设 $f : X \to Y$ 具有 $\mathcal{P}$ 且普遍单射。
那么，在定义 \ref{spaces-more-morphisms-definition-radicial} 的情形中，态射
$(\Spec(K) \times_Y X)_{red} \to \Spec(K)$
普遍单射且具有 $\mathcal{P}$。因此，为证明
$$
\mathcal{P} + \text{普遍单射}
\Rightarrow
\text{根态}
$$
只需证明：任意非空约化代数空间 $X$，若其在域上的结构态射
$X \to \Spec(K)$ 普遍单射且具有 $\mathcal{P}$，则该空间可由某个域的谱表示。
% P09-ERR-0506
确实，这时 $X \to \Spec(K)$ 将是概形态射，而概形态射的根态与普遍单射等价；
见《态射》引理 \ref{morphisms-lemma-universally-injective}。

\medskip\noindent
先证明 (1)。假设 $f$ 良态且普遍单射。由《良态空间》诸引理
\ref{decent-spaces-lemma-base-change-relative-conditions}、
\ref{decent-spaces-lemma-composition-relative-conditions} 及
\ref{decent-spaces-lemma-properties-trivial-implications}
（最后一个引理说明浸入是良态的），第一段的讨论适用。
设 $X$ 为域 $K$ 上非空、良态、约化且普遍单射的代数空间。
特别地，$|X|$ 是单点集。由《良态空间》引理
\ref{decent-spaces-lemma-when-field}，可知
$X \cong \Spec(L)$，其中 $K \subset L$ 是某个扩张，如所需。

\medskip\noindent
拟分离态射是良态的，见《良态空间》引理
\ref{decent-spaces-lemma-properties-trivial-implications}。
故 (1) 蕴含 (2)。

\medskip\noindent
下面证明 (3)。回顾：分离性公理在基变换与复合下稳定，并且闭浸入是分离的；
见《空间的态射》诸引理
\ref{spaces-morphisms-lemma-base-change-separated}、
\ref{spaces-morphisms-lemma-composition-separated} 及
\ref{spaces-morphisms-lemma-immersions-monomorphisms}。
因此，证明第一段的讨论适用。
设 $X$ 为域 $K$ 上约化、普遍单射且局部分离的代数空间。
特别地，$|X|$ 是单点集，因而 $X$ 拟紧；见《空间的性质》引理
\ref{spaces-properties-lemma-quasi-compact-space}。
由《空间的性质》引理
\ref{spaces-properties-lemma-quasi-compact-affine-cover}，
可找到满平展态射 $U \to X$，其中 $U$ 仿射。考虑概形态射
$$
j :
U \times_X U
\longrightarrow
U \times_{\Spec(K)} U
$$
由于 $X \to \Spec(K)$ 普遍单射，$j$ 是满射；
由于 $X \to \Spec(K)$ 局部分离，$j$ 是浸入。
满浸入是闭浸入，见《概形》引理
\ref{schemes-lemma-immersion-when-closed}。
因此 $R = U \times_X U$ 作为仿射概形的闭子概形而仿射，
特别地，$R$ 拟紧。于是 $X = U/R$ 拟分离，结论由 (2) 得出。
\end{proof}

\begin{remark}
\label{spaces-more-morphisms-remark-weakly-radicial}
设 $X \to Y$ 为代数空间的态射。对根态射的某些应用，
只需要求对每个态射 $\Spec(K) \to Y$（其中 $K$ 为域）：
\begin{enumerate}
\item 空间 $|\Spec(K) \times_Y X|$ 为单点集；
\item 存在单态射 $\Spec(L) \to \Spec(K) \times_Y X$；以及
\item $K \subset L$ 纯不可分。
\end{enumerate}
若以后有需要，我们可以称这样的态射为{\it 弱根态射}。% P09-ERR-0507
例如，若 $X \to Y$ 是满的弱根态射，则映射
$X(k) \to Y(k)$ 对每个代数闭域 $k$ 都是满射。注意，例
\ref{spaces-more-morphisms-example-universally-injective-not-radicial} 中态射的基变换
$X_{\overline{\mathbf{Q}}} \to \Spec(\overline{\mathbf{Q}})$
是弱根态射，但不是根态射。引理
\ref{spaces-more-morphisms-lemma-when-universally-injective-radicial} 的对应结论是：
若 $X \to Y$ 具有性质 ($\beta$) 且普遍单射，则它是弱根态射
（证明略）。
\end{remark}

\begin{lemma}
\label{spaces-more-morphisms-lemma-check-universally-injective}
设 $S$ 为概形，$f : X \to Y$ 为 $S$ 上代数空间的态射。假设
\begin{enumerate}
\item $f$ 局部有限型；
\item 对每个平展态射 $V \to Y$，映射
$|X \times_Y V| \to |V|$ 是单射。
\end{enumerate}
则 $f$ 普遍单射。
\end{lemma}

\begin{proof}
由《空间的态射》引理
\ref{spaces-morphisms-lemma-universally-injective-local}，
该问题在 $Y$ 上是平展局部的。因此可假设 $Y$ 是概形。
于是 $Y$ 特别是良态的；由《良态空间》引理
\ref{decent-spaces-lemma-conditions-on-point-in-fibre-and-qf}，
$f$ 局部拟有限。设 $y \in Y$ 为一点，并令 $X_y$ 为概形论纤维。
假设 $X_y$ 非空。由《域上的空间》引理
\ref{spaces-over-fields-lemma-locally-quasi-finite-over-field}，
$X_y$ 是在 $\kappa(y)$ 上局部拟有限的概形。
由于 $|X_y| \subset |X|$ 是 $|X| \to |Y|$ 在 $y$ 上的纤维，
$X_y$ 有唯一一点 $x$。空间
$X_y \times_{\Spec(\kappa(y))} \Spec(k)$
对任意有限可分扩张 $k/\kappa(y)$ 亦如此，
因为可将 $k$ 实现为位于 $y$ 上方的、某个 $Y$ 上平展概形中的一点的剩余域；
见《态射进阶》引理
\ref{more-morphisms-lemma-realize-prescribed-residue-field-extension-etale}。
于是由《簇》引理
\ref{varieties-lemma-characterize-geometrically-disconnected}，
$X_y$ 几何连通。这蕴含有限扩张 $\kappa(x)/\kappa(y)$
是纯不可分的。

\medskip\noindent
因此（在 $Y$ 为概形的情形），对每个 $y \in Y$，
纤维 $X_y$ 或为空，或
$(X_y)_{red} = \Spec(\kappa(x))$，其中
$\kappa(y) \subset \kappa(x)$ 纯不可分。
故 $f$ 为根态射（略去若干细节），从而由引理
\ref{spaces-more-morphisms-lemma-radicial-implies-universally-injective} 也普遍单射。
\end{proof}




\section{单态射}
\label{spaces-more-morphisms-section-monomorphisms}

\noindent
本节续接《空间的态射》第
\ref{spaces-morphisms-section-monomorphisms} 节。
我们想知道代数空间的每个单态射是否都可表。如果你能证明这一点，
或能给出反例，请发邮件至
\href{mailto:stacks.project@gmail.com}{stacks.project@gmail.com}。
目前已知下列情形成立：
\begin{enumerate}
\item 对局部有限型的单态射成立
（更一般地，由《空间的态射》引理
\ref{spaces-morphisms-lemma-locally-quasi-finite-separated-representable}，
任意分离且局部拟有限的态射都可表；又由《空间的态射》引理
\ref{spaces-morphisms-lemma-monomorphism-loc-finite-type-loc-quasi-finite}，
局部有限型的单态射局部拟有限）；
\item 若靶是若干零维局部环之谱的不交并，则成立
（《良态空间》引理
\ref{decent-spaces-lemma-monomorphism-toward-disjoint-union-dim-0-rings}）；
\item 对平坦单态射成立（见下文）。
\end{enumerate}

\begin{lemma}[David Rydh]
\label{spaces-more-morphisms-lemma-flat-case}
代数空间的平坦单态射可由概形表示。
\end{lemma}

\begin{proof}
设 $f : X \to Y$ 为代数空间的平坦单态射。
要证明 $f$ 可表，必须证明 $X \times_Y V$ 对每个概形
$V$（它映到 $Y$）都是概形。由于成为概形是局部性质（《空间的性质》引理
\ref{spaces-properties-lemma-subscheme}），可假设 $V$ 仿射。
因此可假设 $Y = \Spec(B)$ 是仿射概形。
接着，以一个拟紧开子空间替换 $X$，可假设 $X$ 拟紧。
空间 $X$ 是分离的，因为 $X \to X \times_{\Spec(B)} X$ 是同构。
应用《空间的极限》引理 \ref{spaces-limits-lemma-enough-local}，
可归约到下述情形：$B$ 是局部环，$X \to \Spec(B)$ 是平坦单态射，
并且存在点 $x \in X$ 映到 $\Spec(B)$ 的闭点。
由于分离代数空间的平坦态射可提升一般化（见《良态空间》引理
\ref{decent-spaces-lemma-generalizations-lift-flat}），
$X \to \Spec(B)$ 是满射。因此
$\{X \to \Spec(B)\}$ 是 fpqc 覆盖。
态射 $X \to \Spec(B)$ 沿 $X \to \Spec(B)$ 基变换后成为同构，
故由 fpqc 下降，它本身是同构；见《空间上的下降》引理
\ref{spaces-descent-lemma-descending-property-isomorphism}。
\end{proof}

\noindent
下面是上述引理的一个（某种意义下的）变体。

\begin{lemma}
\label{spaces-more-morphisms-lemma-ui-case}
设 $S$ 为概形，$f : X \to Y$ 为代数空间的拟紧单态射，
并且对每个 $T \to Y$，映射
$$
\mathcal{O}_T \to f_{T,*}\mathcal{O}_{X \times_Y T}
$$
都是单射。则 $f$ 是同构（因而可由概形表示）。
\end{lemma}

\begin{proof}
该问题在 $Y$ 上是平展局部的，故可假设
$Y = \Spec(A)$ 仿射。于是 $X$ 拟紧；可选取仿射概形
$U = \Spec(B)$ 及满平展态射 $U \to X$
（《空间的性质》引理
\ref{spaces-properties-lemma-quasi-compact-affine-cover}）。
注意 $U \times_X U = \Spec(B \otimes_A B)$。
因此，拟凝聚 $\mathcal{O}_X$-模的范畴等价于
模下降数据范畴 $DD_{B/A}$，其中下降数据相对于环同态 $A \to B$。
见《空间的性质》命题
\ref{spaces-properties-proposition-quasi-coherent}、
《下降》定义 \ref{descent-definition-descent-datum-modules} 及
《下降》小节 \ref{descent-subsection-descent-modules-morphisms}。
另一方面，
$$
A \to B
$$
是普遍单射的环同态。事实上，给定一个 $A$-模 $M$，
由引理的假设可知
$A \oplus M \to B \otimes_A (A \oplus M)$ 是单射。
因此，由《下降》定理 \ref{descent-theorem-descent}，
$DD_{B/A}$ 等价于 $A$-模范畴。于是沿
$f : X \to \Spec(A)$ 的拉回给出拟凝聚模范畴的等价。
特别地，$f^*$ 在拟凝聚模上正合，因而 $f$ 平坦
（略去一个小细节）。此外，$f$ 显然满射
（例如，因为 $\Spec(B) \to \Spec(A)$ 满射）。
故 $\{X \to \Spec(A)\}$ 是 fpqc 覆盖。
态射 $X \to \Spec(A)$ 沿 $X \to \Spec(A)$ 基变换后成为同构，
所以由 fpqc 下降，它本身是同构；见《空间上的下降》引理
\ref{spaces-descent-lemma-descending-property-isomorphism}。
\end{proof}

\begin{lemma}
\label{spaces-more-morphisms-lemma-flat-surjective-monomorphism}
代数空间的拟紧、平坦、满的单态射是同构。
\end{lemma}

\begin{proof}
这样的态射满足引理 \ref{spaces-more-morphisms-lemma-ui-case} 的假设。
\end{proof}







\section{浸入的余法层}
\label{spaces-more-morphisms-section-conormal-sheaf}

\noindent
设 $S$ 为概形，$i : Z \to X$ 为 $S$ 上代数空间的闭浸入。
令 $\mathcal{I} \subset \mathcal{O}_X$ 为对应的拟凝聚理想层；
见《空间的态射》引理
\ref{spaces-morphisms-lemma-closed-immersion-ideals}。
考虑短正合列
$$
0 \to \mathcal{I}^2 \to \mathcal{I} \to \mathcal{I}/\mathcal{I}^2 \to 0
$$
，它由 $X$ 上的拟凝聚层组成。
由于层 $\mathcal{I}/\mathcal{I}^2$ 被 $\mathcal{I}$ 零化，
由《空间的态射》引理
\ref{spaces-morphisms-lemma-i-star-equivalence}，它对应于 $Z$ 上的一个层。
这个拟凝聚 $\mathcal{O}_Z$-模称为{\it $Z$ 在 $X$ 中的余法层}。
沿用《空间的态射》第
\ref{spaces-morphisms-section-closed-immersions-quasi-coherent} 节中所述的
记号滥用，通常仍将它记作 $\mathcal{I}/\mathcal{I}^2$。

\medskip\noindent
若 $i : Z \to X$ 是（局部闭）浸入，则把 $i$ 的余法层定义为闭浸入
$i : Z \to X \setminus \partial Z$ 的余法层；见《空间的态射》注
\ref{spaces-morphisms-remark-immersion}。
它通常记作 $\mathcal{I}/\mathcal{I}^2$，其中 $\mathcal{I}$
是闭浸入 $i : Z \to X \setminus \partial Z$ 的理想层。

\begin{definition}
\label{spaces-more-morphisms-definition-conormal-sheaf}
设 $i : Z \to X$ 为浸入。称 $\mathcal{C}_{Z/X}$ 为
{\it $Z$ 在 $X$ 中的余法层}，亦称{\it $i$ 的余法层}；它是上述拟凝聚
$\mathcal{O}_Z$-模 $\mathcal{I}/\mathcal{I}^2$。
\end{definition}

\noindent
在 \cite[IV Definition 16.1.2]{EGA} 中，该层记作
$\mathcal{N}_{Z/X}$。我们不采用这一约定，因为希望把记号
$\mathcal{N}_{Z/X}$ 留给{\it 浸入的法层}。若余法层有限表示，
则法层定义为
$$
\mathcal{N}_{Z/X} =
\SheafHom_{\mathcal{O}_Z}(\mathcal{C}_{Z/X}, \mathcal{O}_Z) =
\SheafHom_{\mathcal{O}_Z}(\mathcal{I}/\mathcal{I}^2, \mathcal{O}_Z)
$$
（否则法层甚至可能不是拟凝聚的）。稍后将回到法层
（此处插入未来引用）。% P09-ERR-0508

\begin{lemma}
\label{spaces-more-morphisms-lemma-etale-conormal}
设 $S$ 为概形，$i : Z \to X$ 为浸入。设
$\varphi : U \to X$ 为平展态射，其中 $U$ 是概形。
令 $Z_U = U \times_X Z$；它是 $U$ 的局部闭子概形。则典范地有
$$
\mathcal{C}_{Z/X}|_{Z_U} = \mathcal{C}_{Z_U/U}
$$
并且该等式关于 $U$ 具有函子性。
\end{lemma}

\begin{proof}
取闭子空间 $T \subset X$，使 $i$ 定义到 $X \setminus T$ 中的闭浸入。
令 $\mathcal{I}$ 为 $X \setminus T$ 上定义 $Z$ 的拟凝聚理想层。
则引理只是说，$\mathcal{I}|_{U \setminus \varphi^{-1}(T)}$
是浸入 $Z_U \to U \setminus \varphi^{-1}(T)$ 的理想层。
这由《空间的态射》引理
\ref{spaces-morphisms-lemma-closed-immersion-ideals} 中
$\mathcal{I}$ 的构造立即可见。
\end{proof}

\begin{lemma}
\label{spaces-more-morphisms-lemma-conormal-functorial}
设 $S$ 为概形，并设
$$
\xymatrix{
Z \ar[r]_i \ar[d]_f & X \ar[d]^g \\
Z' \ar[r]^{i'} & X'
}
$$
为 $S$ 上代数空间的交换图。假设 $i$、$i'$ 为浸入。% P09-ERR-0509
则存在典范的
$\mathcal{O}_Z$-模映射
$$
f^*\mathcal{C}_{Z'/X'}
\longrightarrow
\mathcal{C}_{Z/X}
$$
\end{lemma}

\begin{proof}
先找开子空间 $U' \subset X'$ 和 $U \subset X$，使
$g(U) \subset U'$，且 $i(Z) \subset U$ 和
$i(Z') \subset U'$ 都是闭的（略去存在性的证明）。% P09-ERR-0510
将 $X$ 替换为 $U$，将 $X'$ 替换为 $U'$，可假设 $i$ 与 $i'$
都是闭浸入。令 $\mathcal{I}' \subset \mathcal{O}_{X'}$ 和
$\mathcal{I} \subset \mathcal{O}_X$ 为分别对应于 $i'$ 与 $i$
的拟凝聚理想层；见《空间的态射》引理
\ref{spaces-morphisms-lemma-closed-immersion-ideals}。
考虑复合
$$
g^{-1}\mathcal{I}' \to g^{-1}\mathcal{O}_{X'}
\xrightarrow{g^\sharp} \mathcal{O}_X \to
\mathcal{O}_X/\mathcal{I} = i_*\mathcal{O}_Z
$$
由于 $g(i(Z)) \subset Z'$，可知该复合为零（见《空间的态射》引理
\ref{spaces-morphisms-lemma-closed-immersion-ideals} 中关于分解的陈述）。
因此得到交换图
$$
\xymatrix{
0 \ar[r] &
\mathcal{I} \ar[r] &
\mathcal{O}_X \ar[r] &
i_*\mathcal{O}_Z \ar[r] &
0 \\
0 \ar[r] &
g^{-1}\mathcal{I}' \ar[r] \ar[u] &
g^{-1}\mathcal{O}_{X'} \ar[r] \ar[u] &
g^{-1}i'_*\mathcal{O}_{Z'} \ar[r] \ar[u] &
0
}
$$
由于 $g^{-1}$ 是正合函子，下行正合。由正合性还可看出
$(g^{-1}\mathcal{I}')^2 = g^{-1}((\mathcal{I}')^2)$。
故该图诱导映射
$g^{-1}(\mathcal{I}'/(\mathcal{I}')^2) \to \mathcal{I}/\mathcal{I}^2$。
将其拉回（例如使用 $i^{-1}$）到 $Z$，得到
$i^{-1}g^{-1}(\mathcal{I}'/(\mathcal{I}')^2) \to \mathcal{C}_{Z/X}$。
由于 $i^{-1}g^{-1} = f^{-1}(i')^{-1}$，这给出映射
$f^{-1}\mathcal{C}_{Z'/X'} \to \mathcal{C}_{Z/X}$，
进而诱导所需映射。
\end{proof}

\begin{lemma}
\label{spaces-more-morphisms-lemma-conormal-functorial-more}
设 $S$ 为概形。定义 \ref{spaces-more-morphisms-definition-conormal-sheaf} 的余法层及
引理 \ref{spaces-more-morphisms-lemma-conormal-functorial} 的函子性满足下列性质：% P09-ERR-0511
\begin{enumerate}
\item 若 $Z \to X$ 是 $S$ 上概形的浸入，则该余法层与
《态射》定义 \ref{morphisms-definition-conormal-sheaf} 中的余法层相同。
\item 若引理 \ref{spaces-more-morphisms-lemma-conormal-functorial} 中所有空间都是概形，则映射
$f^*\mathcal{C}_{Z'/X'} \to \mathcal{C}_{Z/X}$
与《态射》引理
\ref{morphisms-lemma-conormal-functorial} 中构造的映射相同。
\item 给定交换图
$$
\xymatrix{
Z \ar[r]_i \ar[d]_f & X \ar[d]^g \\
Z' \ar[r]^{i'} \ar[d]_{f'} & X' \ar[d]^{g'} \\
Z'' \ar[r]^{i''} & X''
}
$$
则映射 $(f' \circ f)^*\mathcal{C}_{Z''/X''} \to \mathcal{C}_{Z/X}$
等于映射 $f^*\mathcal{C}_{Z'/X'} \to \mathcal{C}_{Z/X}$
与下述映射沿 $f$ 的拉回之复合：
$(f')^*\mathcal{C}_{Z''/X''} \to \mathcal{C}_{Z'/X'}$。
\end{enumerate}
\end{lemma}

\begin{proof}
略。注意，(1) 是引理 \ref{spaces-more-morphisms-lemma-etale-conormal} 的特殊情形。
\end{proof}

\begin{lemma}
\label{spaces-more-morphisms-lemma-conormal-functorial-flat}
设 $S$ 为概形，并设
$$
\xymatrix{
Z \ar[r]_i \ar[d]_f & X \ar[d]^g \\
Z' \ar[r]^{i'} & X'
}
$$
为 $S$ 上代数空间的纤维积图。假设 $i$、$i'$ 为浸入。% P09-ERR-0509
则引理 \ref{spaces-more-morphisms-lemma-conormal-functorial} 的典范映射
$f^*\mathcal{C}_{Z'/X'} \to \mathcal{C}_{Z/X}$ 是满射。
若 $g$ 平坦，则它是同构。
\end{lemma}

\begin{proof}
选择交换图
$$
\xymatrix{
U \ar[r] \ar[d] & X \ar[d] \\
U' \ar[r] & X'
}
$$
其中 $U$、$U'$ 是概形，且水平箭头满且平展；见《空间》引理
\ref{spaces-lemma-lift-morphism-presentations}。
利用引理 \ref{spaces-more-morphisms-lemma-etale-conormal} 和
\ref{spaces-more-morphisms-lemma-conormal-functorial-more}，问题归约为概形态射的情形。
对概形而言，这就是《态射》引理
\ref{morphisms-lemma-conormal-functorial-flat}。
\end{proof}

\begin{lemma}
\label{spaces-more-morphisms-lemma-transitivity-conormal}
设 $S$ 为概形，$Z \to Y \to X$ 为代数空间的浸入。
则存在典范正合列
$$
i^*\mathcal{C}_{Y/X} \to
\mathcal{C}_{Z/X} \to
\mathcal{C}_{Z/Y} \to 0
$$
其中各映射来自引理 \ref{spaces-more-morphisms-lemma-conormal-functorial}，
而 $i : Z \to Y$ 是第一个态射。
\end{lemma}

\begin{proof}
设 $U$ 为概形，$U \to X$ 为满平展态射。
借助引理 \ref{spaces-more-morphisms-lemma-etale-conormal} 和
\ref{spaces-more-morphisms-lemma-conormal-functorial-more}，上述序列的正合性立即转化为
概形浸入
$Z \times_X U \to Y \times_X U \to U$
对应序列的正合性。因此，引理由《态射》引理
\ref{morphisms-lemma-transitivity-conormal} 得出。
\end{proof}








\section{浸入的法锥}
\label{spaces-more-morphisms-section-normal-cone}

\noindent
设 $S$ 为概形，$i : Z \to X$ 为 $S$ 上代数空间的闭浸入。
令 $\mathcal{I} \subset \mathcal{O}_X$ 为对应的拟凝聚理想层；
见《空间的态射》引理
\ref{spaces-morphisms-lemma-closed-immersion-ideals}。
考虑拟凝聚分次 $\mathcal{O}_X$-代数层
$\bigoplus_{n \geq 0} \mathcal{I}^n/\mathcal{I}^{n + 1}$。
由于每个层 $\mathcal{I}^n/\mathcal{I}^{n + 1}$ 都被
$\mathcal{I}$ 零化，由《空间的态射》引理
\ref{spaces-morphisms-lemma-i-star-equivalence}，
该分次代数对应于一个拟凝聚分次 $\mathcal{O}_Z$-代数层。
这个拟凝聚分次 $\mathcal{O}_Z$-代数称为{\it $Z$ 在 $X$ 中的余法代数}。
沿用《空间的态射》第
\ref{spaces-morphisms-section-closed-immersions-quasi-coherent} 节中所述的
记号滥用，通常简称为
$\bigoplus_{n \geq 0} \mathcal{I}^n/\mathcal{I}^{n + 1}$。

\medskip\noindent
若 $i : Z \to X$ 是（局部闭）浸入，则把 $i$ 的余法代数定义为闭浸入
$i : Z \to X \setminus \partial Z$ 的余法代数；见《空间的态射》注
\ref{spaces-morphisms-remark-immersion}。
它通常记作
$\bigoplus_{n \geq 0} \mathcal{I}^n/\mathcal{I}^{n + 1}$，
其中 $\mathcal{I}$ 是闭浸入
$i : Z \to X \setminus \partial Z$ 的理想层。

\begin{definition}
\label{spaces-more-morphisms-definition-conormal-algebra}
设 $i : Z \to X$ 为浸入。称 $\mathcal{C}_{Z/X, *}$ 为
{\it $Z$ 在 $X$ 中的余法代数}，亦称{\it $i$ 的余法代数}；
它是上述拟凝聚分次 $\mathcal{O}_Z$-代数层
$\bigoplus_{n \geq 0} \mathcal{I}^n/\mathcal{I}^{n + 1}$。
\end{definition}

\noindent
因此 $\mathcal{C}_{Z/X, 1} = \mathcal{C}_{Z/X}$ 是该浸入的余法层。
此外，$\mathcal{C}_{Z/X, 0} = \mathcal{O}_Z$，而
$\mathcal{C}_{Z/X, n}$ 是由下述性质刻画的拟凝聚
$\mathcal{O}_Z$-模：
\begin{equation}
\label{spaces-more-morphisms-equation-conormal-in-degree-n}
i_*\mathcal{C}_{Z/X, n} = \mathcal{I}^n/\mathcal{I}^{n + 1}
\end{equation}
其中 $i : Z \to X \setminus \partial Z$，而 $\mathcal{I}$ 是上述
$i$ 的理想层。最后，注意存在典范满射
\begin{equation}
\label{spaces-more-morphisms-equation-conormal-algebra-quotient}
\text{Sym}^*(\mathcal{C}_{Z/X}) \longrightarrow \mathcal{C}_{Z/X, *}
\end{equation}
，它是拟凝聚分次 $\mathcal{O}_Z$-代数的映射，并且在次数
$0$ 和 $1$ 上是同构。

\begin{lemma}
\label{spaces-more-morphisms-lemma-etale-conormal-algebra}
设 $S$ 为概形，$i : Z \to X$ 为 $S$ 上代数空间的浸入。
设 $\varphi : U \to X$ 为平展态射，其中 $U$ 是概形。
令 $Z_U = U \times_X Z$；它是 $U$ 的局部闭子概形。则典范地有
$$
\mathcal{C}_{Z/X, *}|_{Z_U} = \mathcal{C}_{Z_U/U, *}
$$
并且该等式关于 $U$ 具有函子性。
\end{lemma}

\begin{proof}
取闭子空间 $T \subset X$，使 $i$ 定义到 $X \setminus T$ 中的闭浸入。
令 $\mathcal{I}$ 为 $X \setminus T$ 上定义 $Z$ 的拟凝聚理想层。
该引理由如下事实得出：
$\mathcal{I}|_{U \setminus \varphi^{-1}(T)}$ 是浸入
$Z_U \to U \setminus \varphi^{-1}(T)$ 的理想层。
这由《空间的态射》引理
\ref{spaces-morphisms-lemma-closed-immersion-ideals} 中
$\mathcal{I}$ 的构造立即可见。
\end{proof}

\begin{lemma}
\label{spaces-more-morphisms-lemma-conormal-algebra-functorial}
设 $S$ 为概形，并设
$$
\xymatrix{
Z \ar[r]_i \ar[d]_f & X \ar[d]^g \\
Z' \ar[r]^{i'} & X'
}
$$
为 $S$ 上代数空间的交换图。假设 $i$、$i'$ 为浸入。% P09-ERR-0512
则存在典范的分次 $\mathcal{O}_Z$-代数映射
$$
f^*\mathcal{C}_{Z'/X', *}
\longrightarrow
\mathcal{C}_{Z/X, *}
$$
\end{lemma}

\begin{proof}
先找开子空间 $U' \subset X'$ 和 $U \subset X$，使
$g(U) \subset U'$，且 $i(Z) \subset U$ 和
$i(Z') \subset U'$ 都是闭的（略去存在性的证明）。% P09-ERR-0513
将 $X$ 替换为 $U$，将 $X'$ 替换为 $U'$，可假设 $i$ 与 $i'$
都是闭浸入。令 $\mathcal{I}' \subset \mathcal{O}_{X'}$ 和
$\mathcal{I} \subset \mathcal{O}_X$ 为分别对应于 $i'$ 与 $i$
的拟凝聚理想层；见《空间的态射》引理
\ref{spaces-morphisms-lemma-closed-immersion-ideals}。
考虑复合
$$
g^{-1}\mathcal{I}' \to g^{-1}\mathcal{O}_{X'}
\xrightarrow{g^\sharp} \mathcal{O}_X \to
\mathcal{O}_X/\mathcal{I} = i_*\mathcal{O}_Z
$$
由于 $g(i(Z)) \subset Z'$，可知该复合为零（见《空间的态射》引理
\ref{spaces-morphisms-lemma-closed-immersion-ideals} 中关于分解的陈述）。
因此得到交换图
$$
\xymatrix{
0 \ar[r] &
\mathcal{I} \ar[r] &
\mathcal{O}_X \ar[r] &
i_*\mathcal{O}_Z \ar[r] &
0 \\
0 \ar[r] &
g^{-1}\mathcal{I}' \ar[r] \ar[u] &
g^{-1}\mathcal{O}_{X'} \ar[r] \ar[u] &
g^{-1}i'_*\mathcal{O}_{Z'} \ar[r] \ar[u] &
0
}
$$
由于 $g^{-1}$ 是正合函子，下行正合。由正合性还可看出，
$(g^{-1}\mathcal{I}')^n = g^{-1}((\mathcal{I}')^n)$
对所有 $n \geq 1$ 成立。
故该图诱导映射
$g^{-1}((\mathcal{I}')^n/(\mathcal{I}')^{n + 1}) \to
\mathcal{I}^n/\mathcal{I}^{n + 1}$。
将其拉回（例如使用 $i^{-1}$）到 $Z$，得到
$i^{-1}g^{-1}((\mathcal{I}')^n/(\mathcal{I}')^{n + 1}) \to
\mathcal{C}_{Z/X, n}$。
由于 $i^{-1}g^{-1} = f^{-1}(i')^{-1}$，这给出各映射
$f^{-1}\mathcal{C}_{Z'/X', n} \to \mathcal{C}_{Z/X, n}$，
它们诱导所需映射。
\end{proof}

\begin{lemma}
\label{spaces-more-morphisms-lemma-conormal-algebra-functorial-flat}
设 $S$ 为概形，并设
$$
\xymatrix{
Z \ar[r]_i \ar[d]_f & X \ar[d]^g \\
Z' \ar[r]^{i'} & X'
}
$$
为 $S$ 上代数空间的笛卡尔方块，其中
$i$、$i'$ 为浸入。% P09-ERR-0514
则引理 \ref{spaces-more-morphisms-lemma-conormal-algebra-functorial} 的典范映射
$f^*\mathcal{C}_{Z'/X', *} \to \mathcal{C}_{Z/X, *}$ 是满射。
若 $g$ 平坦，则它是同构。
\end{lemma}

\begin{proof}
该陈述可在对 $X'$ 作平展局部化后检验。
在此情形可假设 $X' \to X$ 是概形的态射；% P09-ERR-0515
于是 $Z$ 和 $Z'$ 是概形，结论由概形的情形得出；
见《除子》引理
\ref{divisors-lemma-conormal-algebra-functorial-flat}。
\end{proof}

\noindent
对于代数空间上的锥与向量丛，我们采用与概形情形相同的约定
（概形情形采用 EGA 的约定）；见《构造》第
\ref{constructions-section-cone} 节和第
\ref{constructions-section-vector-bundle} 节。
特别地，向量丛是一个非常一般的对象
（并不要求局部同构于仿射空间丛）。

\begin{definition}
\label{spaces-more-morphisms-definition-normal-cone}
设 $S$ 为概形，$i : Z \to X$ 为 $S$ 上代数空间的浸入。
称{\it 法锥 $C_ZX$} 为 $Z$ 在 $X$ 中的如下锥：
$$
C_ZX = \underline{\Spec}_Z(\mathcal{C}_{Z/X, *})
$$
，见《空间的态射》定义
\ref{spaces-morphisms-definition-relative-spec}。
{\it $Z$ 在 $X$ 中的法丛}是向量丛
$$
N_ZX = \underline{\Spec}_Z(\text{Sym}(\mathcal{C}_{Z/X}))
$$
\end{definition}

\noindent
因此 $C_ZX \to Z$ 是 $Z$ 上的锥，而
$N_ZX \to Z$ 是 $Z$ 上的向量丛。此外，分次代数的典范满射
(\ref{spaces-more-morphisms-equation-conormal-algebra-quotient}) 定义了典范闭浸入
\begin{equation}
\label{spaces-more-morphisms-equation-normal-cone-in-normal-bundle}
C_ZX \longrightarrow N_ZX
\end{equation}
，它是 $Z$ 上锥的态射。










\section{态射的微分层}
\label{spaces-more-morphisms-section-sheaf-differentials}

\noindent
建议读者参阅交换代数章中的相应一节
（《代数》第 \ref{algebra-section-differentials} 节）、
概形态射章中的相应一节
（《态射》第 \ref{morphisms-section-sheaf-differentials} 节），
以及《位点上的模》第
\ref{sites-modules-section-differentials} 节。
我们先证明：概形态射的微分层概念与相应的小平展（带环）位点态射
给出的概念一致。

\medskip\noindent
为清楚陈述下一个引理，我们暂时重新以 $\mathcal{F}^a$
表示一个 $\mathcal{O}_{X_\etale}$-模层，它伴随于拟凝聚
$\mathcal{O}_X$-模 $\mathcal{F}$，而后者定义在概形 $X$ 上；见
《下降》定义 \ref{descent-definition-structure-sheaf}。

\begin{lemma}
\label{spaces-more-morphisms-lemma-match-modules-differentials}
设 $f : X \to Y$ 为概形的态射。设
$f_{small} : X_\etale \to Y_\etale$ 为相伴的小平展位点态射，见
《下降》注 \ref{descent-remark-change-topologies-ringed}。
则存在与泛导子相容的典范同构
$$
(\Omega_{X/Y})^a = \Omega_{X_\etale/Y_\etale}
$$
。这里第一个模是 $X_\etale$ 上的层，它伴随于拟凝聚
$\mathcal{O}_X$-模 $\Omega_{X/Y}$，见
《态射》定义 \ref{morphisms-definition-sheaf-differentials}；
第二个模则来自《位点上的模》定义
\ref{sites-modules-definition-module-differentials}。
\end{lemma}

\begin{proof}
设 $h : U \to X$ 为平展态射。此时自然映射
$h^*\Omega_{X/Y} \to \Omega_{U/Y}$ 是同构，见
《关于态射的更多内容》引理
\ref{more-morphisms-lemma-sheaf-differentials-etale-localization}。
这意味着存在自然的 $\mathcal{O}_{Y_\etale}$-导子
$$
\text{d}^a : \mathcal{O}_{X_\etale} \longrightarrow (\Omega_{X/Y})^a
$$
，因为我们刚刚看到，$(\Omega_{X/Y})^a$ 在任意对象
$U$（它属于 $X_\etale$）上的值典范地等同于
$\Gamma(U, \Omega_{U/Y})$。由
$\text{d}_{X/Y} :
\mathcal{O}_{X_\etale}
\to
\Omega_{X_\etale/Y_\etale}$
的泛性质，存在唯一的 $\mathcal{O}_{X_\etale}$-线性映射
$c : \Omega_{X_\etale/Y_\etale} \to (\Omega_{X/Y})^a$，使得
$\text{d}^a = c \circ \text{d}_{X/Y}$。

\medskip\noindent
反过来，设 $\mathcal{F}$ 为
$\mathcal{O}_{X_\etale}$-模，
而 $D : \mathcal{O}_{X_\etale} \to \mathcal{F}$ 为
$\mathcal{O}_{Y_\etale}$-导子。只需把
$D$ 限制到小 Zariski 位点 $X_{Zar}$，即 $X$ 的相应位点。
由于 $X_{Zar}$ 上的层与 $X$ 上的层一致，见
《下降》注 \ref{descent-remark-Zariski-site-space}，
可见 $D|_{X_{Zar}} : \mathcal{O}_X \to \mathcal{F}|_{X_{Zar}}$
正是一个“通常的” $Y$-导子。因此得到映射
$\psi : \Omega_{X/Y} \longrightarrow \mathcal{F}|_{X_{Zar}}$，
使得 $D|_{X_{Zar}} = \psi \circ \text{d}$。特别地，对
$\mathcal{F} = \Omega_{X_\etale/Y_\etale}$ 应用此结论，得到映射
$$
c' :
\Omega_{X/Y}
\longrightarrow
\Omega_{X_\etale/Y_\etale}|_{X_{Zar}}
$$
考虑《下降》注
\ref{descent-remark-change-topologies-ringed} 和引理
\ref{descent-lemma-compare-sites} 中讨论的带环位点态射
$\text{id}_{small, \etale, Zar} : X_\etale \to X_{Zar}$。
由于限制函子 $\mathcal{F} \mapsto \mathcal{F}|_{X_{Zar}}$
等于 $\text{id}_{small, \etale, Zar, *}$，又由于
$\text{id}_{small, \etale, Zar}^*$ 是
$\text{id}_{small, \etale, Zar, *}$ 的左伴随，且
$(\Omega_{X/Y})^a = \text{id}_{small, \etale, Zar}^*\Omega_{X/Y}$，
可见 $c'$ 伴随于映射
$$
c'' :
(\Omega_{X/Y})^a
\longrightarrow
\Omega_{X_\etale/Y_\etale}.
$$
我们断言 $c''$ 与 $c'$ 互逆。% P09-ERR-0516
该断言即完成引理的证明。
为验证这一点，只需说明：
$c''(\text{d}(f)) = \text{d}_{X/Y}(f)$ 且
$c(\text{d}_{X/Y}(f)) = \text{d}(f)$，其中 $f$ 是
$\mathcal{O}_X$ 在 $X$ 的一个开集上的局部截面。
验证略去。
\end{proof}

\noindent
这样便可作如下定义。另一种做法见注 \ref{spaces-more-morphisms-remark-alternative}。

\begin{definition}
\label{spaces-more-morphisms-definition-sheaf-differentials}
设 $S$ 为概形。设 $f : X \to Y$ 为 $S$ 上代数空间的态射。
{\it 微分层 $\Omega_{X/Y}$}，即 $X$ 在 $Y$ 上的微分层，% P09-ERR-0517
是带环拓扑斯态射
$$
(f_{small}, f^\sharp) :
(X_\etale, \mathcal{O}_X)
\to
(Y_\etale, \mathcal{O}_Y)
$$
的微分层（《位点上的模》定义
\ref{sites-modules-definition-sheaf-differentials}）；该态射见
《空间的性质》引理
\ref{spaces-properties-lemma-morphism-ringed-topoi}。
{\it 泛 $Y$-导子}记为
$\text{d}_{X/Y} : \mathcal{O}_X \to \Omega_{X/Y}$。
\end{definition}

\noindent
由引理 \ref{spaces-more-morphisms-lemma-match-modules-differentials}，
当 $X$ 与 $Y$ 可表时，这一定义不与已有概念冲突。
此后若 $X$ 与 $Y$ 可表，我们不再区分上述微分层与
《态射》定义 \ref{morphisms-definition-sheaf-differentials}
中的微分层。我们希望把它同概形态射通常的微分模联系起来。
关键引理如下。

\begin{lemma}
\label{spaces-more-morphisms-lemma-localize-differentials}
设 $S$ 为概形。设 $f : X \to Y$ 为 $S$ 上代数空间的态射。
考虑任意交换图
$$
\xymatrix{
U \ar[d]_a \ar[r]_\psi & V \ar[d]^b \\
X \ar[r]^f & Y
}
$$
，其中竖直箭头是代数空间的平展态射。则
$$
\Omega_{X/Y}|_{U_\etale} = \Omega_{U/V}
$$
特别地，若 $U$、$V$ 为概形，则它等于概形态射
$U \to V$ 通常的微分层。
\end{lemma}

\begin{proof}
由《空间的性质》引理
\ref{spaces-properties-lemma-etale-morphism-topoi}
和公式 (\ref{spaces-properties-equation-restrict})，
可把 $X_\etale$ 上层到 $U_\etale$ 的限制视为由
$a_{small}$ 所作的拉回。对 $b$ 亦然。由
《位点上的模》引理
\ref{sites-modules-lemma-localize-differentials}，有
$$
\Omega_{X/Y}|_{U_\etale} =
\Omega_{\mathcal{O}_{U_\etale}/
a_{small}^{-1}f_{small}^{-1}\mathcal{O}_{Y_\etale}}
$$
由于
$a_{small}^{-1}f_{small}^{-1}\mathcal{O}_{Y_\etale}
= \psi_{small}^{-1}b_{small}^{-1}\mathcal{O}_{Y_\etale}
= \psi_{small}^{-1}\mathcal{O}_{V_\etale}$，
引理得证。
\end{proof}

\begin{lemma}
\label{spaces-more-morphisms-lemma-module-differentials-quasi-coherent}
设 $S$ 为概形。设 $f : X \to Y$ 为 $S$ 上代数空间的态射。
则 $\Omega_{X/Y}$ 是拟凝聚 $\mathcal{O}_X$-模。
\end{lemma}

\begin{proof}
选择引理 \ref{spaces-more-morphisms-lemma-localize-differentials} 中的一个图，
使 $a$ 与 $b$ 都满，且 $U$ 与 $V$ 都是概形。
于是 $\Omega_{X/Y}|_U = \Omega_{U/V}$，它是拟凝聚的
（例如见《态射》引理
\ref{morphisms-lemma-differentials-diagonal}）。
故由《空间的性质》引理
\ref{spaces-properties-lemma-characterize-quasi-coherent}，
$\Omega_{X/Y}$ 是拟凝聚的。
\end{proof}

\begin{remark}
\label{spaces-more-morphisms-remark-alternative}
既然已知 $\Omega_{X/Y}$ 拟凝聚，就可以尝试用另一种方式构造它。
例如可利用《空间的性质》第
\ref{spaces-properties-section-quasi-coherent-presentation} 节的结果，
通过粘合来构造微分层。
例如，若 $Y$ 是概形，而 $U \to X$ 是从概形到 $X$ 的满平展态射，
则 $\Omega_{U/Y}$ 是拟凝聚 $\mathcal{O}_U$-模；又因
$s, t : R \to U$ 都平展，得到同构
$$
\alpha : s^*\Omega_{U/Y} \to \Omega_{R/Y} \to t^*\Omega_{U/Y}
$$
，这里使用了《态射》引理
\ref{morphisms-lemma-triangle-differentials-smooth}。
可以检验它满足余圈条件，于是构造完成。
若 $Y$ 不是概形，则把 $\Omega_{U/Y}$ 定义为映射
$(U \to Y)^*\Omega_{Y/S} \to \Omega_{U/S}$ 的余核，
再按上述方式进行。这种两步过程略显笨拙。% P09-ERR-0518
另一种可能是对引理 \ref{spaces-more-morphisms-lemma-localize-differentials}
中任意一个图粘合各层 $\Omega_{U/V}$，但这也不很优雅。
不过两种方法都可行，并会给出微分层略为初等的构造。
\end{remark}

\begin{lemma}
\label{spaces-more-morphisms-lemma-functoriality-differentials}
设 $S$ 为概形。设
$$
\xymatrix{
X' \ar[d] \ar[r]_f & X \ar[d] \\
Y' \ar[r] & Y
}
$$
为代数空间的交换图。映射
$f^\sharp : \mathcal{O}_X \to f_*\mathcal{O}_{X'}$
同映射
$f_*\text{d}_{X'/Y'} : f_*\mathcal{O}_{X'} \to f_*\Omega_{X'/Y'}$
的复合是一个 $Y$-导子。因此得到典范的
$\mathcal{O}_X$-模映射
$\Omega_{X/Y} \to f_*\Omega_{X'/Y'}$；再由
$f_*$ 与 $f^*$ 的伴随性，得到典范的
$\mathcal{O}_{X'}$-模同态
$$
c_f : f^*\Omega_{X/Y} \longrightarrow \Omega_{X'/Y'}.
$$
它由下述性质唯一刻画：
$f^*\text{d}_{X/Y}(t)$ 映到
$\text{d}_{X'/Y'}(f^* t)$，
其中 $t$ 是 $\mathcal{O}_X$ 的任意局部截面。
\end{lemma}

\begin{proof}
这是《位点上的模》引理
\ref{sites-modules-lemma-functoriality-differentials-ringed-topoi}
的特例。
\end{proof}

\begin{lemma}
\label{spaces-more-morphisms-lemma-check-functoriality-differentials}
设 $S$ 为概形。设
$$
\xymatrix{
X'' \ar[d] \ar[r]_g & X' \ar[d] \ar[r]_f & X \ar[d] \\
Y'' \ar[r] & Y' \ar[r] & Y
}
$$
为 $S$ 上代数空间的交换图。则有
$$
c_{f \circ g} = c_g \circ g^* c_f
$$
，作为映射 $(f \circ g)^*\Omega_{X/Y} \to \Omega_{X''/Y''}$。
\end{lemma}

\begin{proof}
略。提示：使用 $c_f, c_g, c_{f \circ g}$
对局部截面的作用来刻画这些映射。
\end{proof}

\begin{lemma}
\label{spaces-more-morphisms-lemma-triangle-differentials}
设 $S$ 为概形。
设 $f : X \to Y$、$g : Y \to B$ 为 $S$ 上代数空间的态射。
则存在典范正合列
$$
f^*\Omega_{Y/B} \to \Omega_{X/B} \to \Omega_{X/Y} \to 0
$$
，其中各映射由应用引理
\ref{spaces-more-morphisms-lemma-functoriality-differentials} 得到。
\end{lemma}

\begin{proof}
这由概形情形的相应结果经平展局部化得出；见
《态射》引理 \ref{morphisms-lemma-triangle-differentials}
以及引理 \ref{spaces-more-morphisms-lemma-localize-differentials}。
\end{proof}

\begin{lemma}
\label{spaces-more-morphisms-lemma-immersion-differentials}
设 $S$ 为概形。若 $X \to Y$ 是 $S$ 上代数空间的浸入，
则 $\Omega_{X/S}$ 为零。% P09-ERR-0519
\end{lemma}

\begin{proof}
这由概形情形的相应结果经平展局部化得出；见
《态射》引理 \ref{morphisms-lemma-immersion-differentials}
以及引理 \ref{spaces-more-morphisms-lemma-localize-differentials}。
\end{proof}

\begin{lemma}
\label{spaces-more-morphisms-lemma-differentials-relative-immersion}
设 $S$ 为概形。设 $B$ 为 $S$ 上的代数空间。
设 $i : Z \to X$ 为 $B$ 上代数空间的浸入。
存在典范正合列
$$
\mathcal{C}_{Z/X} \to i^*\Omega_{X/B} \to \Omega_{Z/B} \to 0
$$
，其中第一箭头由 $\text{d}_{X/B}$ 诱导，
第二箭头来自引理 \ref{spaces-more-morphisms-lemma-functoriality-differentials}。
\end{lemma}

\begin{proof}
这是《态射》引理
\ref{morphisms-lemma-differentials-relative-immersion}
的代数空间版本；由该引理经平展局部化即可推出，见引理
\ref{spaces-more-morphisms-lemma-localize-differentials} 和
\ref{spaces-more-morphisms-lemma-etale-conormal}。
不过，我们应当确认第一箭头可在整体上定义。
因此在此说明“由 $\text{d}_{X/B}$ 诱导”的含义。
也就是说，可假设 $i$ 是闭浸入，方法是把 $X$ 替换为一个开子空间。
设 $\mathcal{I} \subset \mathcal{O}_X$
为对应于 $Z \subset X$ 的拟凝聚理想层。
则 $\text{d}_{X/S} : \mathcal{I} \to \Omega_{X/S}$% P09-ERR-0520
把子层 $\mathcal{I}^2 \subset \mathcal{I}$ 映到
$\mathcal{I}\Omega_{X/S}$。因此它诱导映射
$\mathcal{I}/\mathcal{I}^2 \to
\Omega_{X/S}/\mathcal{I}\Omega_{X/S}$，
该映射是 $\mathcal{O}_X/\mathcal{I}$-线性的。
由《空间的态射》引理
\ref{spaces-morphisms-lemma-i-star-equivalence}，
这对应于所需的映射
$\mathcal{C}_{Z/X} \to i^*\Omega_{X/S}$。
\end{proof}

\begin{lemma}
\label{spaces-more-morphisms-lemma-differentials-relative-immersion-section}
设 $S$ 为概形。设 $B$ 为 $S$ 上的代数空间。
设 $i : Z \to X$ 为 $B$ 上代数空间的浸入，并假设
$i$（平展局部地）有左逆。则典范列
$$
0 \to \mathcal{C}_{Z/X} \to i^*\Omega_{X/B} \to \Omega_{Z/B} \to 0
$$
（平展局部地）分裂正合；该列来自引理
\ref{spaces-more-morphisms-lemma-differentials-relative-immersion}。
\end{lemma}

\begin{proof}
说明如下：我们断言，若 $g : X \to Z$ 是 $i$
在 $B$ 上的左逆，则 $i^*c_g$ 是映射
$i^*\Omega_{X/B} \to \Omega_{Z/B}$ 的右逆。
说明这一点后，结论由概形态射的相应结果经平展局部化得出，见引理
\ref{spaces-more-morphisms-lemma-localize-differentials} 和
\ref{spaces-more-morphisms-lemma-etale-conormal}。
\end{proof}

\begin{lemma}
\label{spaces-more-morphisms-lemma-base-change-differentials}
设 $S$ 为概形。
设 $X \to Y$ 为 $S$ 上代数空间的态射。
设 $g : Y' \to Y$ 为 $S$ 上代数空间的态射。
设 $X' = X_{Y'}$ 为 $X$ 的基变换。
以 $g' : X' \to X$ 表示投影。
则映射
$$
(g')^*\Omega_{X/Y} \to \Omega_{X'/Y'}
$$
是同构；该映射来自引理
\ref{spaces-more-morphisms-lemma-functoriality-differentials}。
\end{lemma}

\begin{proof}
这由概形情形的相应结果与平展局部化得出；见
《态射》引理 \ref{morphisms-lemma-base-change-differentials}
以及引理 \ref{spaces-more-morphisms-lemma-localize-differentials}。
\end{proof}

\begin{lemma}
\label{spaces-more-morphisms-lemma-differential-product}
设 $S$ 为概形。
设 $f : X \to B$ 和 $g : Y \to B$ 为 $S$ 上具有相同靶的
代数空间态射。
设 $p : X \times_B Y \to X$ 和 $q : X \times_B Y \to Y$
为投影态射。引理
\ref{spaces-more-morphisms-lemma-functoriality-differentials} 中的映射
$$
p^*\Omega_{X/B} \oplus q^*\Omega_{Y/B}
\longrightarrow
\Omega_{X \times_B Y/B}
$$
给出同构。
\end{lemma}

\begin{proof}
这由概形情形的相应结果与平展局部化得出；见
《态射》引理 \ref{morphisms-lemma-differential-product}
以及引理 \ref{spaces-more-morphisms-lemma-localize-differentials}。
\end{proof}

\begin{lemma}
\label{spaces-more-morphisms-lemma-finite-type-differentials}
设 $S$ 为概形。
设 $f : X \to Y$ 为 $S$ 上代数空间的态射。
若 $f$ 局部有限型，则 $\Omega_{X/Y}$
是有限型 $\mathcal{O}_X$-模。
\end{lemma}

\begin{proof}
这由概形情形的相应结果与平展局部化得出；见
《态射》引理 \ref{morphisms-lemma-finite-type-differentials}
以及引理 \ref{spaces-more-morphisms-lemma-localize-differentials}。
\end{proof}

\begin{lemma}
\label{spaces-more-morphisms-lemma-finite-presentation-differentials}
设 $S$ 为概形。
设 $f : X \to Y$ 为 $S$ 上代数空间的态射。
若 $f$ 局部有限表示，则 $\Omega_{X/Y}$
是有限表示的 $\mathcal{O}_X$-模。
\end{lemma}

\begin{proof}
这由概形情形的相应结果与平展局部化得出；见
《态射》引理 \ref{morphisms-lemma-finite-presentation-differentials}
以及引理 \ref{spaces-more-morphisms-lemma-localize-differentials}。
\end{proof}

\begin{lemma}
\label{spaces-more-morphisms-lemma-smooth-omega-finite-locally-free}
设 $S$ 为概形。
设 $f : X \to Y$ 为 $S$ 上代数空间的光滑态射。
则微分模 $\Omega_{X/Y}$ 是有限局部自由的。
\end{lemma}

\begin{proof}
由引理 \ref{spaces-more-morphisms-lemma-localize-differentials}，该陈述在
$X$ 与 $Y$ 上都是平展局部的。
故它由概形情形得出，见《态射》引理
\ref{morphisms-lemma-smooth-omega-finite-locally-free}。
\end{proof}













\section{平展位点的拓扑不变性}
\label{spaces-more-morphisms-section-topological-invariance}

\noindent
我们证明位点 $X_{spaces, \etale}$ 是一个“拓扑不变量”。
于是 $X_\etale$ 也是拓扑不变量；它由
$X_{spaces, \etale}$ 中的可表对象组成。见引理
\ref{spaces-more-morphisms-lemma-topological-invariance}。

\begin{theorem}
\label{spaces-more-morphisms-theorem-topological-invariance}
设 $S$ 为概形。
设 $f : X \to Y$ 为 $S$ 上代数空间的态射。
假设 $f$ 整、普遍单射且满。函子
$$
V \longmapsto V_X = X \times_Y V
$$
定义范畴等价
$Y_{spaces, \etale} \to X_{spaces, \etale}$。
\end{theorem}

\begin{proof}
态射 $f$ 可表且为普遍同胚，见《空间的态射》第
\ref{spaces-morphisms-section-universal-homeomorphisms} 节。

\medskip\noindent
先证明该函子忠实。
设 $V', V$ 为 $Y_{spaces, \etale}$ 的对象，而
$a, b : V' \to V$ 为 $Y$ 上两个不同的态射。
由于 $V', V$ 在 $Y$ 上平展，如下等化子
$$
E =  V' \times_{(a, b), V \times_Y V, \Delta_{V/Y}} V
$$
（即 $a, b$ 的等化子）也在 $Y$ 上平展。因此 $E \to V'$ 是平展单态射
（即开浸入）；它为同构当且仅当它满。
由于 $X \to Y$ 是普遍同胚，可见此条件成立当且仅当
$E_X = V'_X$，即当且仅当 $a_X = b_X$。

\medskip\noindent
其次证明该函子全忠实。
设 $V', V$ 为 $Y_{spaces, \etale}$ 的对象，而
$c : V'_X \to V_X$ 为 $X$ 上的态射。我们要构造
态射 $a : V' \to V$，它定义在 $Y$ 上，并使得 $a_X = c$。
取满平展态射 $a' : V'' \to V'$，使 $V''$ 为分离代数空间。
若能构造态射 $a'' : V'' \to V$，使得
$a''_X = c \circ a'_X$，则两个复合
$$
V'' \times_{V'} V'' \xrightarrow{\text{pr}_i} V'' \xrightarrow{a''} V
$$
由第一段已经证明的函子忠实性而相等。
因此 $a''$ 唯一地经态射 $a : V' \to V$ 分解，
因为作为层的 $V'$ 是 $V''$ 关于等价关系
$V'' \times_{V'} V''$ 的商。因此可以假设 $V'$ 分离。
此时图
$$
\Gamma_c \subset (V' \times_Y V)_X
$$
既开又闭（细节略）。由于 $X \to Y$ 是普遍同胚，存在既开又闭的子空间
$\Gamma \subset V' \times_Y V$，使得 $\Gamma_X = \Gamma_c$。
投影 $\Gamma \to V'$ 是平展态射，其到 $X$ 的基变换是同构。
因此 $\Gamma \to V'$ 平展、普遍单射且满，故由
《空间的态射》引理
\ref{spaces-morphisms-lemma-etale-universally-injective-open}
可知它是同构。于是 $\Gamma$ 正是所需态射
$a : V' \to V$ 的图。

\medskip\noindent
最后证明该函子本质满。
设 $U$ 为 $X_{spaces, \etale}$ 的对象。
需要找到对象 $V$，它属于 $Y_{spaces, \etale}$，
使得 $V_X \cong U$。取满平展态射 $U' \to U$，
使得 $U' \cong V'_X$ 且
$U' \times_U U' \cong V''_X$，其中 $V'', V'$ 是
$Y_{spaces, \etale}$ 的某些对象。
于是由函子的全忠实性，得到态射% P09-ERR-0521
$s, t : V'' \to V'$，其中
$t_X = \text{pr}_0$ 且 $s_X = \text{pr}_1$；这里它们作为态射
$U' \times_U U' \to U'$ 来看。
利用
$(\text{pr}_0, \text{pr}_1) : U' \times_U U' \to U' \times_S U'$
是平展等价关系，以及 $U' \to V'$ 和
$U' \times_U U' \to V''$ 都普遍单射且满，可推出
$(t, s) : V'' \to V' \times_S V'$ 是平展等价关系。
于是商 $V = V'/V''$（见《空间》定理
\ref{spaces-theorem-presentation}）
是代数空间 $V$，它位于 $Y$ 上。存在态射
$V' \to V$，使得 $V'' = V' \times_V V'$。
因而得到态射 $V \to Y$（见《空间上的下降》引理
\ref{spaces-descent-lemma-fpqc-universal-effective-epimorphisms}）。
基变换到 $X$ 后，可见存在态射 $U' \to V_X$
以及相容同构
$U' \times_{V_X} U' = U' \times_U U'$；
这蕴含 $V_X \cong U$（再次使用刚才引用的引理）。

\medskip\noindent
选择概形 $W$ 和满平展态射 $W \to Y$。
选择概形 $U'$ 和满平展态射 $U' \to U \times_X W_X$。
注意，$U'$ 与 $U' \times_U U'$ 都是在 $X$ 上平展的概形，
且它们到 $X$ 的结构态射经概形 $W_X$ 分解。
因此由《平展上同调》定理
\ref{etale-cohomology-theorem-topological-invariance}，
存在概形 $V', V''$，它们在 $W$ 上平展，且到
$W_X$ 的基变换分别同构于 $U'$ 和 $U' \times_U U'$。
这就完成了证明。
\end{proof}

\begin{lemma}
\label{spaces-more-morphisms-lemma-topological-invariance}
在定理 \ref{spaces-more-morphisms-theorem-topological-invariance}
的假设和记号下，范畴等价
$Y_{spaces, \etale} \to X_{spaces, \etale}$
限制为范畴等价
$Y_\etale \to X_\etale$ 和
$Y_{affine, \etale} \to X_{affine, \etale}$。
\end{lemma}

\begin{proof}
这不过是下述陈述：给定对象
$V \in Y_{spaces, \etale}$，$V$ 为（仿射）概形，
当且仅当 $V \times_Y X$ 为（仿射）概形。
由于 $V \times_Y X \to V$ 整、普遍单射且满
（它是 $X \to Y$ 的基变换），这由《空间的极限》引理
\ref{spaces-limits-lemma-integral-universally-bijective-scheme}
和命题 \ref{spaces-limits-proposition-affine} 得出。
\end{proof}

\begin{remark}
\label{spaces-more-morphisms-remark-topological-invariance-etale-site}
\begin{reference}
Lenny Taelman 于 2016 年 5 月 1 日所发电子邮件。
\end{reference}
代数空间的普遍同胚未必可表，见《空间的态射》例
\ref{spaces-morphisms-example-universal-homeomorphism}。
事实上，定理 \ref{spaces-more-morphisms-theorem-topological-invariance}
对于一般普遍同胚并不成立。为说明这一点，设 $k$ 为特征 $0$
的代数闭域，并令
$$
\mathbf{A}^1 \to X \to \mathbf{A}^1
$$
如《空间的态射》例
\ref{spaces-morphisms-example-universal-homeomorphism}
所述。回忆第一个态射平展，并把
$t$ 与 $-t$（其中 $t \in \mathbf{A}^1_k \setminus \{0\}$）
等同；第二个态射则是我们的普遍同胚。
由于 $\mathbf{A}^1_k$ 没有非平凡的连通有限平展覆盖
（因为 $k$ 是特征零的代数闭域；细节略），
只需构造一个非平凡连通有限平展覆盖
$Y \to X$。为此，令 $Y$ 为原点加倍的仿射直线
（《概形》例 \ref{schemes-example-affine-space-zero-doubled}）。
于是 $Y = Y_1 \cup Y_2$，其中
$Y_i = \mathbf{A}^1_k$，并沿
$\mathbf{A}^1_k \setminus \{0\}$ 粘合。
为定义态射 $Y \to X$，使用态射
$$
Y_1 \xrightarrow{1} \mathbf{A}^1_k \to X
\quad\text{且}\quad
Y_2 \xrightarrow{-1} \mathbf{A}^1_k \to X.
$$
它们在 $Y_1 \cap Y_2$ 上粘合，这是由 $X$ 的构造保证的，
因而定义态射 $Y \to X$。事实上，我们断言
$$
\xymatrix{
Y \ar[d] & Y_1 \amalg Y_2 \ar[l] \ar[d] \\
X & \mathbf{A}^1_k \ar[l]
}
$$
是笛卡尔方块。细节略去；例如可使用《群胚》引理
\ref{groupoids-lemma-criterion-fibre-product}。
由于 $\mathbf{A}^1_k \to X$ 平展且满，
这证明 $Y \to X$ 是次数为 $2$ 的有限平展态射，
从而给出所需例子。

\medskip\noindent
还可用如下更简单的论证。概形 $Y$ 上有群
$G = \{+1, -1\}$ 的自由作用，其中 $-1$
通过交换 $Y_1$ 与 $Y_2$ 并改变坐标符号而作用。
于是 $X = Y/G$（见《空间》定义
\ref{spaces-definition-quotient}），故 $Y \to X$ 有限平展。
还可直接证明存在普遍同胚 $X \to \mathbf{A}^1_k$，
方法是在仿射空间上使用 $t \mapsto t^2$。
事实上，这个 $X$ 就是上面的 $X$。
\end{remark}








\section{加厚}
\label{spaces-more-morphisms-section-thickenings}

\noindent
下述术语未必完全标准，但使用起来很方便。

\begin{definition}
\label{spaces-more-morphisms-definition-thickening}
加厚。设 $S$ 为概形。
\begin{enumerate}
\item 称代数空间 $X'$ 是代数空间 $X$ 的一个{\it 加厚}，
若 $X$ 是 $X'$ 的闭子空间且相伴的拓扑空间相等。
\item 称 $X'$ 是 $X$ 的一个{\it 一阶加厚}，
若 $X$ 是 $X'$ 的闭子空间，并且拟凝聚理想层
$\mathcal{I} \subset \mathcal{O}_{X'}$ 定义 $X$ 且平方为零。
\item 称 $X'$ 是 $X$ 的一个{\it 有限阶加厚}，
若 $X$ 是 $X$ 的闭子空间，% P09-ERR-0522
并且拟凝聚理想层
$\mathcal{I} \subset \mathcal{O}_{X'}$ 定义 $X$ 且幂零；也就是说，
存在整数 $n \geq 0$ 使得 $\mathcal{I}^{n + 1} = 0$。
\item 称 $X'$ 为一个{\it $n$ 阶加厚}（关于 $X$），
若 $X$ 是 $X$ 的闭子空间，且
$\mathcal{I}^{n + 1} = 0$，其中
$\mathcal{I} \subset \mathcal{O}_{X'}$ 是定义 $X$ 的拟凝聚理想层。
% P09-ERR-0522
\item 给定两个加厚 $X \subset X'$ 和 $Y \subset Y'$，
若态射 $f' : X' \to Y'$ 满足
$f(X) \subset Y$，则称它为{\it 加厚的态射}；% P09-ERR-0523
也就是说，$f'|_X$ 经闭子空间 $Y$ 分解。在此情形令
$f = f'|_X : X \to Y$，并称
$(f, f') : (X \subset X') \to (Y \subset Y')$
为加厚的态射。
\item 设 $B$ 为代数空间。类似地定义
{\it $B$ 上的加厚}和{\it $B$ 上加厚的态射}。
这表示上述空间 $X, X', Y, Y'$ 都是赋予了到 $B$
的结构态射的代数空间，而态射
$X \to X'$、$Y \to Y'$ 和 $f' : X' \to Y'$
都是 $B$ 上的态射。
\end{enumerate}
\end{definition}

\noindent
基本等价。注意，若 $X \subset X'$ 是加厚，则
$X \to X'$ 整且普遍双射。这蕴含
\begin{equation}
\label{spaces-more-morphisms-equation-equivalence-etale-spaces}
X_{spaces, \etale} = X'_{spaces, \etale}
\end{equation}
，其中等价由拉回函子给出；见定理
\ref{spaces-more-morphisms-theorem-topological-invariance}。
因此可把 $\mathcal{O}_{X'}$ 看作
$X_{spaces, \etale}$ 上的层。于是有局部带环拓扑斯的典范等价
\begin{equation}
\label{spaces-more-morphisms-equation-fundamental-equivalence}
(\Sh(X'_{spaces, \etale}), \mathcal{O}_{X'})
\cong
(\Sh(X_{spaces, \etale}), \mathcal{O}_{X'})
\end{equation}
。下文将经常把它同《空间的性质》定理
\ref{spaces-properties-theorem-fully-faithful}
的全忠实性结果结合使用。% P09-ERR-0525
例如，闭浸入 $i_X : X \to X'$ 对应于满射
$i_X^\sharp : \mathcal{O}_{X'} \to \mathcal{O}_X$。

\medskip\noindent
设 $S$ 为概形，并设 $B$ 为 $S$ 上的代数空间。
设 $(f, f') : (X \subset X') \to (Y \subset Y')$
为 $B$ 上加厚的态射。注意，连续函子图
$$
\xymatrix{
X_{spaces, \etale} &
Y_{spaces, \etale} \ar[l] \\
X'_{spaces, \etale} \ar[u] &
Y'_{spaces, \etale} \ar[u] \ar[l]
}
$$
交换，且竖直箭头是等价。因此
$f_{spaces, \etale}$、$f_{small}$、
$f'_{spaces, \etale}$ 和 $f'_{small}$
都定义同一个拓扑斯态射。因此可把
$$
(f')^\sharp :
f_{spaces, \etale}^{-1}\mathcal{O}_{Y'}
\longrightarrow
\mathcal{O}_{X'}
$$
看作 $\mathcal{O}_B$-代数层的映射，并且它嵌入交换图
$$
\xymatrix{
f_{spaces, \etale}^{-1}\mathcal{O}_Y \ar[r]_-{f^\sharp} \ar[r] &
\mathcal{O}_X \\
f_{spaces, \etale}^{-1}\mathcal{O}_{Y'} \ar[r]^-{(f')^\sharp}
\ar[u]^{i_Y^\sharp} &
\mathcal{O}_{X'} \ar[u]_{i_X^\sharp}
}
$$
。这里 $i_X : X \to X'$ 和 $i_Y : Y \to Y'$
是给定闭浸入的名称。

\begin{lemma}
\label{spaces-more-morphisms-lemma-first-order-thickening-maps}
设 $S$ 为概形。设 $B$ 为 $S$ 上的代数空间。
设 $X \subset X'$ 和 $Y \subset Y'$ 为 $B$ 上代数空间的加厚。
设 $f : X \to Y$ 为 $B$ 上代数空间的态射。给定任意
$\mathcal{O}_B$-代数映射
$$
\alpha : f_{spaces, \etale}^{-1}\mathcal{O}_{Y'} \to \mathcal{O}_{X'}
$$
，使得图
$$
\xymatrix{
f_{spaces, \etale}^{-1}\mathcal{O}_Y \ar[r]_-{f^\sharp} \ar[r] &
\mathcal{O}_X \\
f_{spaces, \etale}^{-1}\mathcal{O}_{Y'} \ar[r]^-\alpha
\ar[u]^{i_Y^\sharp} &
\mathcal{O}_{X'} \ar[u]_{i_X^\sharp}
}
$$
交换，则存在唯一的 $(f, f')$ 的态射% P09-ERR-0524
作为 $B$ 上加厚的态射，使得 $\alpha = (f')^\sharp$。
\end{lemma}

\begin{proof}
为求得 $f'$，由《空间的性质》定理
\ref{spaces-properties-theorem-fully-faithful}，
只需证明带环拓扑斯态射
$$
(f_{spaces, \etale}, \alpha) :
(\Sh(X_{spaces, \etale}), \mathcal{O}_{X'})
\longrightarrow
(\Sh(Y_{spaces, \etale}), \mathcal{O}_{Y'})
$$
是局部带环拓扑斯的态射。这直接来自局部带环拓扑斯态射的定义
（《位点上的模》定义
\ref{sites-modules-definition-morphism-locally-ringed-topoi}）、
$(f, f^\sharp)$ 是局部带环拓扑斯的态射这一事实
（《空间的性质》引理
\ref{spaces-properties-lemma-morphism-locally-ringed}）、
$\alpha$ 嵌入给定交换图这一事实，以及
$i_X^\sharp$ 与 $i_Y^\sharp$ 的核局部幂零这一事实。
最后，$f' \circ i_X = i_Y \circ f$ 来自该图的交换性和
再次应用《空间的性质》定理
\ref{spaces-properties-theorem-fully-faithful}。
关于 $f'$ 是 $B$ 上态射的验证略去。
\end{proof}

\begin{lemma}
\label{spaces-more-morphisms-lemma-open-subspace-thickening}
设 $S$ 为概形。设 $X \subset X'$ 为 $S$ 上代数空间的加厚。
对任意开子空间 $U \subset X$，存在唯一的开子空间
$U' \subset X'$，使得
$U = X \times_{X'} U'$。
\end{lemma}

\begin{proof}
令 $U' \to X'$ 为 $X'_{spaces, \etale}$ 中的对象，
它经由 (\ref{spaces-more-morphisms-equation-equivalence-etale-spaces}) 对应于对象
$U \to X$，后者属于 $X_{spaces, \etale}$。
态射 $U' \to X'$ 平展且普遍单射，因而为开浸入，见
《空间的态射》引理
\ref{spaces-morphisms-lemma-etale-universally-injective-open}。
\end{proof}

\noindent
设 $i_X : X \to X'$ 为代数空间的加厚。
核 $\mathcal{I} = \Ker(i_X^\sharp) \subset \mathcal{O}_{X'}$
的任意局部截面都局部幂零。这对结构层
$\mathcal{O}_{X'}$ 给出一定控制；但从技术上说，
有限阶加厚这一类 $X \subset X'$ 容易处理得多。
具体而言，若 $X'$ 是一个 $n$ 阶加厚，则有过滤
$$
0 \subset \mathcal{I}^n \subset \mathcal{I}^{n - 1} \subset
\ldots \subset \mathcal{I} \subset \mathcal{O}_{X'}
$$
，并可见 $X'$ 被闭子空间过滤：
$$
X = X_0 \subset X_1 \subset \ldots \subset X_{n - 1} \subset X_{n + 1} = X'
$$
，其中每一对 $X_i \subset X_{i + 1}$ 都是 $B$ 上的一阶加厚。
% P09-ERR-0526 P09-ERR-0527
利用简单的归纳论证，许多关于一阶加厚证明的结果都可改述为
有限阶加厚的结果。

\begin{lemma}
\label{spaces-more-morphisms-lemma-first-order-thickening-surjective}
设 $S$ 为概形。设 $X \subset X'$ 为 $S$ 上代数空间的加厚。
设 $U$ 为 $X_{spaces, \etale}$ 的仿射对象。则
$$
\Gamma(U, \mathcal{O}_{X'}) \to \Gamma(U, \mathcal{O}_X)
$$
满；这里通过 (\ref{spaces-more-morphisms-equation-fundamental-equivalence})
把 $\mathcal{O}_{X'}$ 看作 $X_{spaces, \etale}$ 上的层。
\end{lemma}

\begin{proof}
设 $U' \to X'$ 为代数空间的平展态射，使得
$U = X \times_{X'} U'$；见定理
\ref{spaces-more-morphisms-theorem-topological-invariance}。
由《空间的极限》引理 \ref{spaces-limits-lemma-affine}，
$U'$ 是仿射概形。因此
$\Gamma(U, \mathcal{O}_{X'}) = \Gamma(U', \mathcal{O}_{U'}) \to
\Gamma(U, \mathcal{O}_U)$
为满射，因为 $U \to U'$ 是仿射概形的闭浸入。
下面给出有限阶加厚的直接证明；这是实践中最常使用的情形。
\end{proof}

\begin{proof}[有限阶加厚情形的证明]
由上述原理，可假设 $X \subset X'$ 是一阶加厚。
以 $\mathcal{I}$ 表示满射
$\mathcal{O}_{X'} \to \mathcal{O}_X$ 的核。
由于 $\mathcal{I}$ 是拟凝聚
$\mathcal{O}_{X'}$-模，且一阶加厚的定义给出
$\mathcal{I}^2 = 0$，可应用《空间的态射》引理
\ref{spaces-morphisms-lemma-i-star-equivalence}，
得到 $\mathcal{I}$ 是拟凝聚 $\mathcal{O}_X$-模。
因此，引理由短正合列
$$
0 \to \mathcal{I} \to \mathcal{O}_{X'} \to \mathcal{O}_X \to 0
$$
所关联的上同调长正合列，以及
$H^1_\etale(U, \mathcal{I}) = 0$ 这一事实得出；
后者成立是因为 $\mathcal{I}$ 拟凝聚，见《下降》命题
\ref{descent-proposition-same-cohomology-quasi-coherent}
和《概形的上同调》引理
\ref{coherent-lemma-quasi-coherent-affine-cohomology-zero}。
\end{proof}

\begin{lemma}
\label{spaces-more-morphisms-lemma-thickening-scheme}
设 $S$ 为概形。设 $X \subset X'$ 为 $S$ 上代数空间的加厚。
若 $X$（可由）概形表示，则 $X'$ 亦然。
\end{lemma}

\begin{proof}
注意 $X'_{red} = X_{red}$。因此若 $X$ 是概形，
则 $X'_{red}$ 是概形。于是结论由
《空间的极限》引理
\ref{spaces-limits-lemma-reduction-scheme} 得出。
下面给出有限阶加厚的直接证明；这是实践中最常使用的情形。
\end{proof}

\begin{proof}[有限阶加厚情形的证明]
只需在 $X'$ 是 $X$ 的一阶加厚时证明。
由《空间的性质》引理
\ref{spaces-properties-lemma-subscheme}，
$X'$ 有一个最大的概形开子空间。因此只需证明，
每个点 $x$（在 $|X'| = |X|$ 中）都包含在 $X'$ 的某个概形开子空间中。
使用引理 \ref{spaces-more-morphisms-lemma-open-subspace-thickening}，
可把 $X \subset X'$ 替换为 $U \subset U'$，其中
$x \in U$ 且 $U$ 是仿射概形。因此可假设 $X$ 仿射。
这样就化归到下一段讨论的情形。

\medskip\noindent
假设 $X \subset X'$ 是一阶加厚，其中 $X$ 是仿射概形。
令 $A = \Gamma(X, \mathcal{O}_X)$ 且
$A' = \Gamma(X', \mathcal{O}_{X'})$。
由引理 \ref{spaces-more-morphisms-lemma-first-order-thickening-surjective}，
映射 $A' \to A$ 满。核 $I$ 是平方零理想。
由《空间的性质》引理
\ref{spaces-properties-lemma-morphism-to-affine-scheme}，
得到典范态射 $f : X' \to \Spec(A')$，它嵌入交换图
$$
\xymatrix{
X \ar@{=}[d] \ar[r] &  X' \ar[d]^f \\
\Spec(A) \ar[r] & \Spec(A')
}
$$
。因为水平箭头都是加厚，显然 $f$ 普遍单射且满。
因此只需证明 $f$ 平展；因为此时《空间的态射》引理
\ref{spaces-morphisms-lemma-etale-universally-injective-open}
将蕴含 $f$ 为同构。

\medskip\noindent
为证明 $f$ 平展，选择仿射概形 $U'$ 以及平展态射
$U' \to X'$。只需证明复合
$U' \to X' \to \Spec(A')$ 平展，见《空间的性质》定义
\ref{spaces-properties-definition-etale}。
写 $U' = \Spec(B')$。令 $U = X \times_{X'} U'$。
由于 $U$ 是 $U'$ 的闭子空间，它是闭子概形；故
$U = \Spec(B)$，其中 $B' \to B$ 满。
记 $J = \Ker(B' \to B)$，并注意
$J = \Gamma(U, \mathcal{I})$，其中
$\mathcal{I} = \Ker(\mathcal{O}_{X'} \to \mathcal{O}_X)$
是 $X_{spaces, \etale}$ 上的层，如引理
\ref{spaces-more-morphisms-lemma-first-order-thickening-surjective} 的证明中所述。
态射 $U' \to X' \to \Spec(A')$ 诱导交换图
$$
\xymatrix{
0 \ar[r] &
J \ar[r] &
B' \ar[r] &
B \ar[r] & 0 \\
0 \ar[r] &
I \ar[r] \ar[u] &
A' \ar[r] \ar[u] &
A \ar[r] \ar[u] & 0
}
$$
。现在，由于 $\mathcal{I}$ 是拟凝聚 $\mathcal{O}_X$-模，
有 $\mathcal{I} = (\widetilde I)^a$；记号见《下降》定义
\ref{descent-definition-structure-sheaf}，
这一事实的理由见《下降》命题
\ref{descent-proposition-equivalence-quasi-coherent}。
因此 $J = I \otimes_A B$。
最后，$A \to B$ 平展，因为 $U \to X$ 平展，
而后者是平展态射 $U' \to X'$ 的基变换。
由《代数》引理
\ref{algebra-lemma-lift-etale-infinitesimal}，
可断定 $A' \to B'$ 平展。
\end{proof}

\begin{lemma}
\label{spaces-more-morphisms-lemma-thickening-equivalence}
设 $S$ 为概形。设 $X \subset X'$ 为 $S$ 上代数空间的加厚。
函子
$$
V' \longmapsto V = X \times_{X'} V'
$$
定义范畴等价
$X'_\etale \to X_\etale$。
\end{lemma}

\begin{proof}
函子 $V' \mapsto V$ 定义范畴等价
$X'_{spaces, \etale} \to X_{spaces, \etale}$，见定理
\ref{spaces-more-morphisms-theorem-topological-invariance}。
因此只需证明 $V$ 是概形当且仅当 $V'$ 是概形。
这正是引理 \ref{spaces-more-morphisms-lemma-thickening-scheme} 的内容。
\end{proof}

\noindent
一阶加厚可描述如下。% P09-ERR-0528

\begin{lemma}
\label{spaces-more-morphisms-lemma-first-order-thickening}
设 $S$ 为概形。
设 $f : X \to B$ 为 $S$ 上代数空间的态射。
考虑短正合列
$$
0 \to \mathcal{I} \to \mathcal{A} \to \mathcal{O}_X \to 0
$$
，它由 $X_\etale$ 上的层组成，其中 $\mathcal{A}$ 是
$f^{-1}\mathcal{O}_B$-代数层，
$\mathcal{A} \to \mathcal{O}_X$ 是
$f^{-1}\mathcal{O}_B$-代数层的满射，而
$\mathcal{I}$ 是其核。若
\begin{enumerate}
\item $\mathcal{I}$ 是 $\mathcal{A}$ 中的平方零理想；且
\item $\mathcal{I}$ 作为 $\mathcal{O}_X$-模是拟凝聚的，
\end{enumerate}
则存在一阶加厚 $X \subset X'$（在 $B$ 上），以及同构
$\mathcal{O}_{X'} \to \mathcal{A}$；它是
$f^{-1}\mathcal{O}_B$-代数同构，并与到 $\mathcal{O}_X$ 的满射相容。
\end{lemma}

\begin{proof}
在本证明中，我们重做引理
\ref{spaces-more-morphisms-lemma-first-order-thickening-surjective} 和
\ref{spaces-more-morphisms-lemma-thickening-scheme} 的证明里使用过的一些论证。
先处理 $B = S = \Spec(\mathbf{Z})$ 的情形。
设 $U$ 为仿射概形，并设 $U \to X$ 平展。则
$$
0 \to \mathcal{I}(U) \to \mathcal{A}(U) \to \mathcal{O}_X(U) \to 0
$$
正合，因为 $H^1(U_\etale, \mathcal{I}) = 0$；
这是由于 $\mathcal{I}$ 拟凝聚，见《下降》命题
\ref{descent-proposition-same-cohomology-quasi-coherent}
和《概形的上同调》引理
\ref{coherent-lemma-quasi-coherent-affine-cohomology-zero}。
若 $V \to U$ 是 $X_{spaces, \etale}$ 的仿射对象之间的态射，则
$$
\mathcal{I}(V) = \mathcal{I}(U) \otimes_{\mathcal{O}_X(U)} \mathcal{O}_X(V)
$$
，因为 $\mathcal{I}$ 是拟凝聚 $\mathcal{O}_X$-模，见
《下降》命题
\ref{descent-proposition-equivalence-quasi-coherent}。
因此 $\mathcal{A}(U) \to \mathcal{A}(V)$ 是平展环同态，见
《代数》引理 \ref{algebra-lemma-lift-etale-infinitesimal}。
故可见
$$
U \longmapsto U' = \Spec(\mathcal{A}(U))
$$
是从 $X_{affine, \etale}$ 到以仿射概形为对象、平展态射为态射的范畴的函子。
事实上，我们断言该函子可扩张为函子 $U \mapsto U'$，
其定义域为整个 $X_\etale$。为说明这一点，若 $U$ 是 $X_\etale$ 的对象，注意
$$
0 \to \mathcal{I}|_{U_{Zar}} \to \mathcal{A}|_{U_{Zar}} \to
\mathcal{O}_X|_{U_{Zar}} \to 0
$$
且 $\mathcal{I}|_{U_{Zar}}$ 是 $U$ 上的拟凝聚层，见《下降》命题
\ref{descent-proposition-equivalence-quasi-coherent-functorial}。
因此由《关于态射的更多内容》引理
\ref{more-morphisms-lemma-first-order-thickening}，
得到概形的一阶加厚 $U \subset U'$，使得
$\mathcal{O}_{U'}$ 同构于 $\mathcal{A}|_{U_{Zar}}$。
显然，此构造与上面对仿射对象的构造相容。

\medskip\noindent
选择表示 $X = U/R$，见《空间》定义
\ref{spaces-definition-presentation}，使
$s, t : R \to U$ 定义平展等价关系。
应用上述函子，得到概形中的平展等价关系
$s', t' : R' \to U'$。考虑代数空间
$X' = U'/R'$（见《空间》定理
\ref{spaces-theorem-presentation}）。
态射 $X = U/R \to U'/R' = X'$ 是一阶加厚。
把 $\mathcal{O}_{X'}$ 看作 $X_\etale$ 上的层。
由构造，有同构
$$
\gamma :
\mathcal{O}_{X'}|_{U_\etale}
\longrightarrow
\mathcal{A}|_{U_\etale}
$$
，使得 $s^{-1}\gamma$ 与 $t^{-1}\gamma$ 在 $R_\etale$ 上一致。
因此由《空间的性质》引理
\ref{spaces-properties-lemma-descent-sheaf}，
$\gamma$ 来自唯一的所需同构
$\mathcal{O}_{X'} \to \mathcal{A}$。

\medskip\noindent
为处理一般基代数空间 $B$ 的情形，先按上法把
$X'$ 构造为 $\mathbf{Z}$ 上的代数空间。
然后用同构 $\mathcal{O}_{X'} \to \mathcal{A}$ 定义
$f^{-1}\mathcal{O}_B \to \mathcal{O}_{X'}$。
依引理 \ref{spaces-more-morphisms-lemma-first-order-thickening-maps}，
这定义了态射 $X' \to B$，它与给定态射
$X \to B$ 相容，证明完毕。
\end{proof}

\begin{lemma}
\label{spaces-more-morphisms-lemma-base-change-thickening}
设 $S$ 为概形。设 $Y \subset Y'$ 为 $S$ 上代数空间的加厚。
设 $X' \to Y'$ 为态射，并令 $X = Y \times_{Y'} X'$。
则
$(X \subset X') \to (Y \subset Y')$
是加厚的态射。若 $Y \subset Y'$ 是一阶
（相应地，有限阶）加厚，则 $X \subset X'$ 是一阶
（相应地，有限阶）加厚。
\end{lemma}

\begin{proof}
略。
\end{proof}

\begin{lemma}
\label{spaces-more-morphisms-lemma-composition-thickening}
设 $S$ 为概形。若 $X \subset X'$ 和 $X' \subset X''$
是 $S$ 上代数空间的加厚，则 $X \subset X''$ 亦然。
\end{lemma}

\begin{proof}
略。
\end{proof}

\begin{lemma}
\label{spaces-more-morphisms-lemma-descending-property-thickening}
“是加厚”这一性质是 fpqc 局部的。
一阶加厚亦然。
\end{lemma}

\begin{proof}
该陈述的含义如下：设 $S$ 为概形，并设
$X \to X'$ 为 $S$ 上代数空间的态射。
设 $\{g_i : X'_i \to X'\}$ 为代数空间的 fpqc 覆盖，
使基变换 $X_i \to X'_i$ 对所有 $i$ 都是加厚。
则 $X \to X'$ 是加厚。由于各态射 $g_i$ 联合满，
可知 $X \to X'$ 满。由《空间上的下降》引理
\ref{spaces-descent-lemma-descending-property-closed-immersion}，
可知 $X \to X'$ 是闭浸入。
因此 $X \to X'$ 是加厚。一阶加厚情形的证明略去。
\end{proof}


\section{加厚的态射}
\label{spaces-more-morphisms-section-morphisms-thickenings}

\noindent
若 $(f, f') : (X \subset X') \to (Y \subset Y')$ 是代数空间加厚的态射，
则态射 $f$ 的性质往往也由 $f'$ 继承。这里有若干变体。

\begin{lemma}
\label{spaces-more-morphisms-lemma-thicken-property-morphisms}
设 $S$ 为概形。设
$(f, f') : (X \subset X') \to (Y \subset Y')$
为 $S$ 上代数空间加厚的态射。则
\begin{enumerate}
\item $f$ 是仿射态射当且仅当 $f'$ 是仿射态射；
\item $f$ 是满态射当且仅当 $f'$ 是满态射；
\item $f$ 拟紧当且仅当 $f'$ 拟紧；% P09-ERR-0529
\item $f$ 普遍闭当且仅当 $f'$ 普遍闭；
\item $f$ 是整态射当且仅当 $f'$ 是整态射；
\item $f$（拟）分离当且仅当 $f'$（拟）分离；
\item $f$ 普遍单射当且仅当 $f'$ 普遍单射；
\item $f$ 普遍开当且仅当 $f'$ 普遍开；
\item $f$ 可表当且仅当 $f'$ 可表；且
\item 在此添加更多内容。% P09-ERR-0530
\end{enumerate}
\end{lemma}

\begin{proof}
注意 $Y \to Y'$ 和 $X \to X'$ 都是整的普遍同胚。
这立即推出第 (2)、(3)、(4)、(7) 和 (8) 项。
第 (1) 项由《空间的极限》命题
\ref{spaces-limits-proposition-affine} 得出；该命题告诉我们，
平展于 $X$ 与 $X'$ 的仿射概形之间有一一对应，
平展于 $Y$ 与 $Y'$ 的仿射概形之间亦然。
第 (5) 项由第 (1)、(4) 项及《空间的态射》引理
\ref{spaces-morphisms-lemma-integral-universally-closed} 得出。
最后，注意
$$
X \times_Y X = X \times_{Y'} X \to X \times_{Y'} X' \to X' \times_{Y'} X'
$$
是加厚（由引理 \ref{spaces-more-morphisms-lemma-base-change-thickening}，两个箭头都是加厚）。
因此，对加厚的态射
$(X \subset X') \to (X \times_Y X \to X' \times_{Y'} X')$
应用第 (3)、(4) 项，即得第 (6) 项。% P09-ERR-0531
最后，第 (9) 项来自如下事实：概形的代数空间加厚仍为概形，见引理
\ref{spaces-more-morphisms-lemma-thickening-scheme}。
\end{proof}

\begin{lemma}
\label{spaces-more-morphisms-lemma-thicken-property-morphisms-cartesian}
设 $S$ 为概形。设
$(f, f') : (X \subset X') \to (Y \subset Y')$
为 $S$ 上代数空间加厚的态射，并满足
$X = Y \times_{Y'} X'$。若 $X \subset X'$ 是有限阶加厚，则
\begin{enumerate}
\item $f$ 是闭浸入当且仅当 $f'$ 是闭浸入；
\item $f$ 局部有限型当且仅当 $f'$ 局部有限型；
\item $f$ 局部拟有限当且仅当 $f'$ 局部拟有限；
\item $f$ 局部有限型且相对维数为 $d$ 当且仅当
$f'$ 局部有限型且相对维数为 $d$；
\item $\Omega_{X/Y} = 0$ 当且仅当 $\Omega_{X'/Y'} = 0$；
\item $f$ 非分歧当且仅当 $f'$ 非分歧；
\item $f$ 固有当且仅当 $f'$ 固有；
\item $f$ 是有限态射当且仅当 $f'$ 是有限态射；% P09-ERR-0532
\item $f$ 是单态射当且仅当 $f'$ 是单态射；
\item $f$ 是浸入当且仅当 $f'$ 是浸入；且
\item 在此添加更多内容。% P09-ERR-0533
\end{enumerate}
\end{lemma}

\begin{proof}
选择概形 $V'$ 及满平展态射 $V' \to Y'$。
选择概形 $U'$ 及满平展态射
$U' \to X' \times_{Y'} V'$。令 $V = Y \times_{Y'} V'$ 且
$U = X \times_{X'} U'$。于是，对态射的平展局部性质，
可化归到概形加厚的态射
$(U \subset U') \to (V \subset V')$，并应用《关于态射的更多内容》引理
\ref{more-morphisms-lemma-thicken-property-morphisms-cartesian}。
这证明了第 (2)、(3)、(4)、(5) 和 (6) 项。

\medskip\noindent
第 (1)、(7)、(8)、(9)、(10) 项中的态射性质在基变换下稳定；
因此，若 $f'$ 具有性质 $\mathcal{P}$，则 $f$ 也具有它。见
《空间》引理 \ref{spaces-lemma-base-change-immersions}，以及
《空间的态射》引理
\ref{spaces-morphisms-lemma-base-change-proper}、
\ref{spaces-morphisms-lemma-base-change-integral} 和
\ref{spaces-morphisms-lemma-base-change-monomorphism}。

\medskip\noindent
第 (1)、(7)、(8)、(9)、(10) 项中值得讨论的方向，是假设
$f$ 具有该性质并推出 $f'$ 也具有它。
对加厚的阶数作归纳，可假设 $Y \subset Y'$ 是一阶加厚，
见上面对有限阶加厚的讨论。

\medskip\noindent
第 (1) 项的证明。选择概形 $V'$ 及满平展态射
$V' \to Y'$。令 $V = Y \times_{Y'} V'$、
$U' = X' \times_{Y'} V'$ 且 $U = X \times_Y V$。
于是 $U \to V$ 是闭浸入，这推出 $U$ 是概形，
进而推出 $U'$ 是概形（引理 \ref{spaces-more-morphisms-lemma-thickening-scheme}）。
故可对
$(U \subset U') \to (V \subset V')$
应用概形情形的引理
（《关于态射的更多内容》引理
\ref{more-morphisms-lemma-thicken-property-morphisms-cartesian}），
从而得到结论。

\medskip\noindent
第 (7) 项的证明。把第 (2) 项与引理
\ref{spaces-more-morphisms-lemma-thicken-property-morphisms} 的结果结合，
并使用如下事实：固有等于拟紧 $+$ 分离 $+$ 局部有限型
$+$ 普遍闭。

\medskip\noindent
第 (8) 项的证明。把第 (2) 项与引理
\ref{spaces-more-morphisms-lemma-thicken-property-morphisms} 的结果结合，
并使用如下事实：有限等于整 $+$ 局部有限型
（《态射》引理 \ref{morphisms-lemma-finite-integral}）。

\medskip\noindent
第 (9) 项的证明。由于 $f$ 是单态射，有 $X = X \times_Y X$。
可把目前已证的结果应用于加厚的态射
$(X \subset X') \to (X \times_Y X \subset X' \times_{Y'} X')$。
由第 (1) 项可知 $X' \to X' \times_{Y'} X'$ 是闭浸入。
事实上，它是一阶加厚，因为定义闭浸入
$X' \to X' \times_{Y'} X'$ 的理想包含在理想
$\mathcal{I} \subset \mathcal{O}_{Y'}$ 的拉回中；该理想切出
$Y$（在 $Y'$ 中）。确实，
$X = X \times_Y X = (X' \times_{Y'} X') \times_{Y'} Y$
包含在 $X'$ 中。闭浸入
$\Delta : X' \to X' \times_{Y'} X'$ 的余法层等于
$\Omega_{X'/Y'}$（这是《态射》引理
\ref{morphisms-lemma-differentials-diagonal} 对代数空间的类似陈述；
它可由平展局部化推出，也可由引理
\ref{spaces-more-morphisms-lemma-differentials-relative-immersion-section} 和
\ref{spaces-more-morphisms-lemma-differential-product} 结合推出；略去若干细节）。
因此只需证明 $\Omega_{X'/Y'} = 0$；这由第 (5) 项及
$X/Y$ 的相应陈述得出。

\medskip\noindent
第 (10) 项的证明。若 $f : X \to Y$ 是浸入，则它分解为
$X \to V \to Y$，其中 $V \to Y$ 是开子空间，而
$X \to V$ 是闭浸入，见《空间的态射》评注
\ref{spaces-morphisms-remark-immersion}。
设 $V' \subset Y'$ 为开子空间，使其底拓扑空间
$|V'|$ 与 $|V| \subset |Y| = |Y'|$ 相同。
于是 $X' \to Y'$ 通过 $V'$ 分解；由第 (1) 项可知
$X' \to V'$ 是闭浸入。证明完毕。
\end{proof}

\noindent
下面的引理是前一引理的一个变体。我们不再假设所涉及的加厚为
有限阶加厚（该假设使局部有限型这一性质可从 $f$ 传递到 $f'$），
而是直接假设 $f$ 和 $f'$ 都局部有限型。

\begin{lemma}
\label{spaces-more-morphisms-lemma-properties-that-extend-over-thickenings}
设 $S$ 为概形。设
$(f, f') : (X \subset X') \to (Y \to Y')$
为 $S$ 上代数空间加厚的态射。% P09-ERR-0534
假设 $f$ 和 $f'$ 局部有限型，并且
$X = Y \times_{Y'} X'$。则
\begin{enumerate}
\item $f$ 局部拟有限当且仅当 $f'$ 局部拟有限；
\item $f$ 有限当且仅当 $f'$ 有限；
\item $f$ 是闭浸入当且仅当 $f'$ 是闭浸入；
\item $\Omega_{X/Y} = 0$ 当且仅当 $\Omega_{X'/Y'} = 0$；
\item $f$ 非分歧当且仅当 $f'$ 非分歧；
\item $f$ 是单态射当且仅当 $f'$ 是单态射；
\item $f$ 是浸入当且仅当 $f'$ 是浸入；
\item $f$ 固有当且仅当 $f'$ 固有；且
\item 在此添加更多内容。% P09-ERR-0535
\end{enumerate}
\end{lemma}

\begin{proof}
选择概形 $V'$ 及满平展态射 $V' \to Y'$。
选择概形 $U'$ 及满平展态射
$U' \to X' \times_{Y'} V'$。令 $V = Y \times_{Y'} V'$ 且
$U = X \times_{X'} U'$。于是，对态射的平展局部性质，
可化归到概形加厚的态射
$(U \subset U') \to (V \subset V')$，并应用
《关于态射的更多内容》引理
\ref{more-morphisms-lemma-properties-that-extend-over-thickenings}。
这证明了第 (1)、(4) 和 (5) 项。

\medskip\noindent
第 (2)、(3)、(6)、(7)、(8) 项中的性质在基变换下稳定；
因此，若 $f'$ 具有性质 $\mathcal{P}$，则 $f$ 也具有它。见
《空间》引理 \ref{spaces-lemma-base-change-immersions}，以及
《空间的态射》引理
\ref{spaces-morphisms-lemma-base-change-proper}、
\ref{spaces-morphisms-lemma-base-change-integral} 和
\ref{spaces-morphisms-lemma-base-change-monomorphism}。
因此在每一种情形中，只需证明：若 $f$ 具有所需性质，
则 $f'$ 也具有它。

\medskip\noindent
第 (2) 种情形由引理 \ref{spaces-more-morphisms-lemma-thicken-property-morphisms}
的第 (5) 种情形，以及如下事实得出：有限态射恰为
局部有限型的整态射
（《空间的态射》引理 \ref{spaces-morphisms-lemma-finite-integral}）。

\medskip\noindent
第 (3) 种情形。这立即由《空间的极限》引理
\ref{spaces-limits-lemma-check-closed-infinitesimally} 得出。

\medskip\noindent
第 (6) 项的证明。由于 $f$ 是单态射，有 $X = X \times_Y X$。
可把目前已证的结果应用于加厚的态射
$(X \subset X') \to (X \times_Y X \subset X' \times_{Y'} X')$。
由第 (3) 项可知
$\Delta_{X'/Y'} : X' \to X' \times_{Y'} X'$ 是闭浸入。
事实上，$\Delta_{X'/Y'}$ 诱导双射
$|X'| \to |X' \times_{Y'} X'|$，故 $\Delta_{X'/Y'}$ 是加厚。
另一方面，由《空间的态射》引理
\ref{spaces-morphisms-lemma-diagonal-morphism-finite-type}，
$\Delta_{X'/Y'}$ 局部有限表示。换言之，
$\Delta_{X'/Y'}(X')$ 由有限型拟凝聚理想层
$\mathcal{J} \subset \mathcal{O}_{X' \times_{Y'} X'}$ 切出。
由于第 (5) 项给出 $\Omega_{X'/Y'} = 0$，可知闭浸入
$X' \to X' \times_{Y'} X'$ 的余法层为零。
（闭浸入 $\Delta_{X'/Y'}$ 的余法层等于
$\Omega_{X'/Y'}$；这是《态射》引理
\ref{morphisms-lemma-differentials-diagonal} 对代数空间的类似陈述；
它可由平展局部化推出，也可由引理
\ref{spaces-more-morphisms-lemma-differentials-relative-immersion-section} 和
\ref{spaces-more-morphisms-lemma-differential-product} 结合推出；略去若干细节。）
换言之，$\mathcal{J}/\mathcal{J}^2 = 0$。
这推出 $\Delta_{X'/Y'}$ 是同构，例如可用《代数》引理
\ref{algebra-lemma-ideal-is-squared-union-connected}。

\medskip\noindent
第 (7) 项的证明。若 $f : X \to Y$ 是浸入，则它分解为
$X \to V \to Y$，其中 $V \to Y$ 是开子空间，而
$X \to V$ 是闭浸入，见《空间的态射》评注
\ref{spaces-morphisms-remark-immersion}。
设 $V' \subset Y'$ 为开子空间，使其底拓扑空间
$|V'|$ 与 $|V| \subset |Y| = |Y'|$ 相同。
于是 $X' \to Y'$ 通过 $V'$ 分解；由第 (3) 项可知
$X' \to V'$ 是闭浸入。

\medskip\noindent
第 (8) 种情形由引理 \ref{spaces-more-morphisms-lemma-thicken-property-morphisms}
以及固有态射的定义得出：它们是局部有限型、拟紧、
普遍闭且分离的态射。
\end{proof}


\section{加厚的 Picard 群}
\label{spaces-more-morphisms-section-picard-group-thickening}

\noindent
这里给出一些关于加厚的 Picard 群的材料。

\begin{lemma}
\label{spaces-more-morphisms-lemma-picard-group-first-order-thickening}
设 $S$ 为概形。设 $X \subset X'$ 为 $S$ 上代数空间的一阶加厚，
其理想层为 $\mathcal{I}$。则有典范正合列
$$
\xymatrix{
0 \ar[r] &
H^0(X, \mathcal{I}) \ar[r] &
H^0(X', \mathcal{O}_{X'}^*) \ar[r] &
H^0(X, \mathcal{O}^*_X) \ar `r[d] `d[l] `l[llld] `d[dll] [dll] \\
& H^1(X, \mathcal{I}) \ar[r] &
\Pic(X') \ar[r] &
\Pic(X) \ar `r[d] `d[l] `l[llld] `d[dll] [dll] \\
& H^2(X, \mathcal{I}) \ar[r] & \ldots \ar[r] & \ldots
}
$$
，其中各项均为阿贝尔群。
\end{lemma}

\begin{proof}
回忆 $X_\etale = X'_\etale$，见引理
\ref{spaces-more-morphisms-lemma-thickening-equivalence}，以及更一般地见第
\ref{spaces-more-morphisms-section-thickenings} 节中的讨论。
引理中的序列是与阿贝尔群层的短正合列
$$
0 \to \mathcal{I} \to \mathcal{O}_{X'}^* \to \mathcal{O}_X^* \to 0
$$
相伴的长正合上同调列；该短正合列位于 $X_\etale$ 上，
其中第一个映射把局部截面 $f$（属于 $\mathcal{I}$）
送到可逆截面 $1 + f$（属于 $\mathcal{O}_{X'}$）。
还要使用带环位点的 Picard 群与可逆函数层的一阶上同调群之间的同一，
见《位点上的上同调》引理
\ref{sites-cohomology-lemma-h1-invertible}。
\end{proof}


\section{无穷小邻域}
\label{spaces-more-morphisms-section-first-order-infinitesimal-neighbourhood}

\noindent
有限阶加厚的一种自然构造如下。假设
$i : Z \to X$ 是代数空间的浸入。% P09-ERR-0536
选择开子空间 $U \subset X$，使 $i$ 把 $Z$ 同一为闭子空间
$Z \subset U$（见《空间的态射》评注
\ref{spaces-morphisms-remark-immersion}）。
设 $\mathcal{I} \subset \mathcal{O}_U$ 为定义 $Z$（在 $U$ 中）的
拟凝聚理想层，见《空间的态射》引理
\ref{spaces-morphisms-lemma-closed-immersion-ideals}。
对 $n \geq 1$，可考虑闭子空间 $Z_n \subset U$，
它由拟凝聚理想层 $\mathcal{I}^{n + 1}$ 定义。

\begin{definition}
\label{spaces-more-morphisms-definition-first-order-infinitesimal-neighbourhood}
设 $i : Z \to X$ 为代数空间的浸入。
\begin{enumerate}
\item $Z$ 在 $X$ 中的{\it 一阶无穷小邻域}，是上面描述的
一阶加厚 $Z \subset Z_1$（在 $X$ 上）。
\item {\it $n$ 阶无穷小邻域}，即 $Z$ 在 $X$ 中的这一邻域，
是上面描述的 $n$ 阶加厚 $Z \subset Z_n$（在 $X$ 上）。
\end{enumerate}
\end{definition}

\noindent
该加厚具有如下泛性质（这也消除了上述构造依赖于开集
$U$ 的任何疑虑）。

\begin{lemma}
\label{spaces-more-morphisms-lemma-first-order-infinitesimal-neighbourhood}
设 $i : Z \to X$ 为代数空间的浸入。
\begin{enumerate}
\item 一阶无穷小邻域 $Z'$（即 $Z$ 在 $X$ 中的这一邻域）
具有如下泛性质：
给定任意交换图
$$
\xymatrix{
Z \ar[d]_i & T \ar[l]^a \ar[d] \\
X & T' \ar[l]_b
}
$$
其中 $T \subset T'$ 是 $X$ 上的一阶加厚，则存在唯一的加厚的态射
$(a', a) : (T \subset T') \to (Z \subset Z')$（在 $X$ 上）。% P09-ERR-0537
\item 对 $n \geq 1$，$n$ 阶无穷小邻域 $Z_n$
（即 $Z$ 在 $X$ 中的这一邻域）具有如下泛性质：给定任意交换图
$$
\xymatrix{
Z \ar[d]_i & T \ar[l]^a \ar[d] \\
X & T' \ar[l]_b
}
$$
其中 $T \subset T'$ 是 $n$ 阶加厚（在 $X$ 上），则存在唯一的加厚的态射
$(a', a) : (T \subset T') \to (Z \subset Z_n)$（在 $X$ 上）。% P09-ERR-0537
\end{enumerate}
\end{lemma}

\begin{proof}
只证明第 (1) 项。
设 $U \subset X$ 为构造 $Z'$ 时使用的开子空间；换言之，
该开集把 $Z$ 同一为 $U$ 中由拟凝聚理想层
$\mathcal{I}$ 切出的闭子空间。
由于 $|T| = |T'|$，可见 $|b|(|T'|) \subset |U|$。
因此可把 $b$ 看作到 $U$ 的态射，见《空间的性质》引理
\ref{spaces-properties-lemma-factor-through-open-subspace}。
设 $\mathcal{J} \subset \mathcal{O}_{T'}$ 为切出 $T$ 的
平方零拟凝聚理想层。由图的交换性，有
$b|_T = i \circ a$，其中 $i : Z \to U$ 是闭浸入。
由《空间的态射》引理
\ref{spaces-morphisms-lemma-closed-immersion-ideals}，可知
$b^\sharp(b^{-1}\mathcal{I}) \subset \mathcal{J}$。
由于 $T'$ 是 $T$ 的一阶加厚，有 $\mathcal{J}^2 = 0$，
从而 $b^\sharp(b^{-1}(\mathcal{I}^2)) = 0$。
由《空间的态射》引理
\ref{spaces-morphisms-lemma-closed-immersion-ideals}，
这意味着 $b$ 通过 $Z'$ 分解。令
$a' : T' \to Z'$ 为该分解，即得结论。
\end{proof}

\begin{lemma}
\label{spaces-more-morphisms-lemma-infinitesimal-neighbourhood-conormal}
设 $i : Z \to X$ 为代数空间的浸入。
设 $Z \subset Z'$ 为 $Z$ 在 $X$ 中的一阶无穷小邻域。则图
$$
\xymatrix{
Z \ar[r] \ar[d] & Z' \ar[d] \\
Z \ar[r] & X
}
$$
由引理 \ref{spaces-more-morphisms-lemma-conormal-functorial} 诱导余法层之间的映射
$\mathcal{C}_{Z/X} \to \mathcal{C}_{Z/Z'}$。
该映射是同构。
\end{lemma}

\begin{proof}
这由上述 $Z'$ 的构造显然可见。
\end{proof}


\section{形式光滑、形式平展与形式非分歧的变换}
\label{spaces-more-morphisms-section-formally-smooth-etale-unramified}

\noindent
回忆：环同态 $R \to A$ 称为{\it 形式光滑}、相应地称为
{\it 形式平展}、相应地称为{\it 形式非分歧}
（见《代数》定义 \ref{algebra-definition-formally-smooth}、
相应地定义 \ref{algebra-definition-formally-etale}、
相应地定义 \ref{algebra-definition-formally-unramified}），
如果对每个实线部分交换的图
$$
\xymatrix{
A \ar[r] \ar@{-->}[rd] & B/I \\
R \ar[r] \ar[u] & B \ar[u]
}
$$
（其中 $I \subset B$ 为平方零理想），存在、相应地存在唯一、
相应地至多存在一个虚线箭头使该图交换。
这启发了代数空间态射以及更一般的函子的如下类似概念。

\begin{definition}
\label{spaces-more-morphisms-definition-formally-smooth-etale-unramified}
设 $S$ 为概形。设 $a : F \to G$ 为函子
$F, G : (\Sch/S)_{fppf}^{opp} \to \textit{Sets}$ 之间的变换。
考虑如下实线部分交换的图：
$$
\xymatrix{
F \ar[d]_a & T \ar[d]^i \ar[l] \\
G & T' \ar[l] \ar@{-->}[lu]
}
$$
其中 $T$ 和 $T'$ 为仿射概形，而 $i$ 是由平方零理想定义的闭浸入。
\begin{enumerate}
\item 若对任意上述实线图，均存在使图交换的虚线箭头，
则称 $a$ {\it 形式光滑}\footnote{这只是这里可以采用的一种定义。
另一稍弱的条件，是要求虚线箭头在 $T'$ 上 fppf 局部存在。
在某种意义下，这一较弱概念具有更好的形式性质。}。
\item 若对任意上述实线图，恰有一个使图交换的虚线箭头，
则称 $a$ {\it 形式平展}。
\item 若对任意上述实线图，至多有一个使图交换的虚线箭头，
则称 $a$ {\it 形式非分歧}。
\end{enumerate}
\end{definition}

\begin{lemma}
\label{spaces-more-morphisms-lemma-formally-etale-is-combination}
设 $S$ 为概形。设 $a : F \to G$ 为函子
$F, G : (\Sch/S)_{fppf}^{opp} \to \textit{Sets}$ 之间的变换。
则 $a$ 形式平展当且仅当 $a$ 既形式光滑又形式非分歧。
\end{lemma}

\begin{proof}
这由定义形式地得出。
\end{proof}

\begin{lemma}
\label{spaces-more-morphisms-lemma-composition-formally-smooth-etale-unramified}
关于复合，有：
\begin{enumerate}
\item 函子之间的形式光滑变换的复合仍形式光滑；
\item 函子之间的形式平展变换的复合仍形式平展；
\item 函子之间的形式非分歧变换的复合仍形式非分歧。
\end{enumerate}
\end{lemma}

\begin{proof}
这是形式的。
\end{proof}

\begin{lemma}
\label{spaces-more-morphisms-lemma-base-change-formally-smooth-etale-unramified}
设 $S$ 为包含在 $\Sch_{fppf}$ 中的概形。
设 $F, G, H : (\Sch/S)_{fppf}^{opp} \to \textit{Sets}$。
设 $a : F \to G$、$b : H \to G$ 为函子之间的变换。
考虑纤维积图
$$
\xymatrix{
H \times_{b, G, a} F \ar[r]_-{b'} \ar[d]_{a'} & F \ar[d]^a \\
H \ar[r]^b & G
}
$$
\begin{enumerate}
\item 若 $a$ 形式光滑，则基变换 $a'$ 形式光滑；
\item 若 $a$ 形式平展，则基变换 $a'$ 形式平展；
\item 若 $a$ 形式非分歧，则基变换 $a'$ 形式非分歧。
\end{enumerate}
\end{lemma}

\begin{proof}
这是形式的。
\end{proof}

\begin{lemma}
\label{spaces-more-morphisms-lemma-representable-property-formally-property}
设 $S$ 为概形。
设 $F, G : (\Sch/S)_{fppf}^{opp} \to \textit{Sets}$。
设 $a : F \to G$ 为可表的函子之间的变换。
\begin{enumerate}
\item 若 $a$ 光滑，则 $a$ 形式光滑；
\item 若 $a$ 平展，则 $a$ 形式平展；
\item 若 $a$ 非分歧，则 $a$ 形式非分歧。
\end{enumerate}
\end{lemma}

\begin{proof}
考虑实线部分交换的图
$$
\xymatrix{
F \ar[d]_a & T \ar[d]^i \ar[l] \\
G & T' \ar[l] \ar@{-->}[lu]
}
$$
，如定义 \ref{spaces-more-morphisms-definition-formally-smooth-etale-unramified} 中所述。
则 $F \times_G T'$ 是 $T'$ 上光滑（相应地，平展；相应地，非分歧）的概形。
因此由《关于态射的更多内容》引理
\ref{more-morphisms-lemma-smooth-formally-smooth}
（相应地，引理 \ref{more-morphisms-lemma-etale-formally-etale}；
相应地，引理 \ref{more-morphisms-lemma-unramified-formally-unramified}），
可以补入（相应地，唯一地补入；相应地，至多以一种方式补入）
下图中的虚线箭头：
$$
\xymatrix{
F \times_G T' \ar[d] & T \ar[d]^i \ar[l] \\
T' & T' \ar[l] \ar@{-->}[lu]
}
$$
，并因而也得到第一个图中的相应断言。% P09-ERR-0538
\end{proof}

\begin{lemma}
\label{spaces-more-morphisms-lemma-etale-on-top}
设 $S$ 为包含在 $\Sch_{fppf}$ 中的概形。
设 $F, G, H : (\Sch/S)_{fppf}^{opp} \to \textit{Sets}$。
设 $a : F \to G$、$b : G \to H$ 为函子之间的变换。
假设 $a$ 可表、满且平展。
\begin{enumerate}
\item 若 $b$ 形式光滑，则 $b \circ a$ 形式光滑；
\item 若 $b$ 形式平展，则 $b \circ a$ 形式平展；
\item 若 $b$ 形式非分歧，则 $b \circ a$ 形式非分歧。
\end{enumerate}
反过来，考虑实线部分交换的图
$$
\xymatrix{
G \ar[d]_b & T \ar[d]^i \ar[l] \\
H & T' \ar[l] \ar@{-->}[lu]
}
$$
，其中 $T'$ 是 $S$ 上的仿射概形，而
$i : T \to T'$ 是由平方零理想定义的闭浸入。
\begin{enumerate}
\item[(4)] 若 $b \circ a$ 形式光滑，则对每个 $t \in T$，
存在仿射概形之间的平展态射 $U' \to T'$ 及态射
$U' \to G$，使得图
$$
\xymatrix{
G \ar[d]_b & T \ar[l] & T \times_{T'} U' \ar[d] \ar[l]\\
H & T' \ar[l] & U' \ar[llu] \ar[l]
}
$$
交换，且 $t$ 位于 $U' \to T'$ 的像中。
\item[(5)] 若 $b \circ a$ 形式非分歧，则上述图中至多有一个
虚线箭头；即 $b$ 形式非分歧。
\item[(6)] 若 $b \circ a$ 形式平展，则上述图中恰有一个
虚线箭头；即 $b$ 形式平展。
\end{enumerate}
\end{lemma}

\begin{proof}
假设 $b$ 形式光滑（相应地，形式平展；相应地，形式非分歧）。
由于平展态射既光滑又非分歧，可知 $a$ 可表且光滑
（相应地，平展；相应地，非分歧）。因此第 (1)、(2)、(3) 项
由引理 \ref{spaces-more-morphisms-lemma-representable-property-formally-property} 与引理
\ref{spaces-more-morphisms-lemma-composition-formally-smooth-etale-unramified} 结合得出。

\medskip\noindent
假设 $b \circ a$ 形式光滑。考虑引理陈述中的图。令
$W = F \times_G T$。由假设，$W$ 是 $T$ 上满平展的概形。
由《平展态射》定理 \ref{etale-theorem-remarkable-equivalence}，
存在概形 $W'$，它平展于 $T'$，使得
$W = T \times_{T'} W'$。选择仿射开子概形
$U' \subset W'$，使 $t$ 位于 $U' \to T'$ 的像中。
因为 $b \circ a$ 形式光滑，可知存在态射 $U' \to F$，% P09-ERR-0539
使图
$$
\xymatrix{
F \ar[d]_{b \circ a} & W \ar[l] & T \times_{T'} U' \ar[d] \ar[l]\\
H & T' \ar[l] & U' \ar[llu] \ar[l]
}
$$
交换。取复合 $U' \to F \to G$，便得到引理第 (5) 项中的映射。
% P09-ERR-0540

\medskip\noindent
假设 $f, g : T' \to G$ 是两个适合引理中图的虚线箭头。
令 $W = F \times_G T$。由假设，$W$ 是 $T$ 上满平展的概形。
由《平展态射》定理 \ref{etale-theorem-remarkable-equivalence}，
存在概形 $W'$，它平展于 $T'$，使得
$W = T \times_{T'} W'$。由于 $a$ 形式平展，复合
$$
W' \to T' \xrightarrow{f} G
\quad\text{和}\quad
W' \to T' \xrightarrow{g} G
$$
可提升为态射 $f', g' : W' \to F$
（在仿射开集上提升，并由唯一性粘合）。
现在若 $b \circ a : F \to H$ 形式非分歧，则
$f' = g'$，故 $f = g$，因为 $W' \to T'$ 是平展覆盖。
这证明了引理的第 (6) 项。% P09-ERR-0541

\medskip\noindent
假设 $b \circ a$ 形式平展。由第 (4) 项，可在 $T'$ 上平展局部地
找到适合该图的虚线箭头；由第 (5) 项，该虚线箭头唯一。
因此可粘合局部解，得到断言 (6)。略去若干细节。
\end{proof}

\begin{remark}
\label{spaces-more-morphisms-remark-tempting}
在引理 \ref{spaces-more-morphisms-lemma-etale-on-top} 的情形中，很容易猜想
``$b$ 形式光滑'' $\Leftrightarrow$ ``$b \circ a$ 形式光滑''。
然而，这在一般情形下很可能不成立。
\end{remark}

\begin{lemma}
\label{spaces-more-morphisms-lemma-formally-permanence}
设 $S$ 为概形。
设 $F, G, H : (\Sch/S)_{fppf}^{opp} \to \textit{Sets}$。
设 $a : F \to G$、$b : G \to H$ 为函子之间的变换。
假设 $b$ 形式非分歧。
\begin{enumerate}
\item 若 $b \circ a$ 形式非分歧，则 $a$ 形式非分歧；
\item 若 $b \circ a$ 形式平展，则 $a$ 形式平展；
\item 若 $b \circ a$ 形式光滑，则 $a$ 形式光滑。
\end{enumerate}
\end{lemma}

\begin{proof}
设 $T \subset T'$ 为由平方零理想定义的仿射概形闭浸入。
给定 $g' : T' \to G$ 和 $f : T \to F$，使得
$g'|_T = a \circ f$。因为 $b$ 形式非分歧，下列两个集合之间
存在一一对应：
$$
\{f' : T' \to F \mid f = f'|_T\text{ 且 }a \circ f' = g'\}
$$
与
$$
\{f' : T' \to F \mid f = f'|_T\text{ 且 }b \circ a \circ f' = b \circ g'\}.
$$
引理由此形式地得出。
\end{proof}








\section{形式非分歧态射}
\label{spaces-more-morphisms-section-formally-unramified}

\noindent
本节具体说明代数空间态射形式非分歧的含义。

\begin{definition}
\label{spaces-more-morphisms-definition-formally-unramified}
设 $S$ 为概形。设 $f : X \to Y$ 为代数空间的态射，并且它在 $S$ 上。
若它作为函子间的变换是形式非分歧的，则称它为
{\it 形式非分歧的}；其含义见定义
\ref{spaces-more-morphisms-definition-formally-smooth-etale-unramified}。
\end{definition}

\noindent
我们不再重述在函子间形式非分歧变换这一更一般的框架中证明的结果，
参见第 \ref{spaces-more-morphisms-section-formally-smooth-etale-unramified} 节。
这一性质可用相对微分模的消失刻画，见引理
\ref{spaces-more-morphisms-lemma-characterize-formally-unramified}。

\begin{lemma}
\label{spaces-more-morphisms-lemma-formally-unramified}
设 $S$ 为概形。设 $f : X \to Y$ 为 $S$ 上代数空间的态射。
下列条件等价：
\begin{enumerate}
\item $f$ 形式非分歧；
\item 对每个图
$$
\xymatrix{
U \ar[d] \ar[r]_\psi & V \ar[d] \\
X \ar[r]^f & Y
}
$$
其中 $U$、$V$ 为概形且竖直箭头平展，概形态射 $\psi$
是形式非分歧的（参见《态射进阶》定义
\ref{more-morphisms-definition-formally-unramified}）；
\item 对于一个这样的图，若竖直箭头满，则 $\psi$ 形式非分歧。
\end{enumerate}
\end{lemma}

\begin{proof}
假设 $f$ 形式非分歧。由引理
\ref{spaces-more-morphisms-lemma-representable-property-formally-property}，
态射 $U \to X$ 和 $V \to Y$ 形式非分歧。因此，由引理
\ref{spaces-more-morphisms-lemma-composition-formally-smooth-etale-unramified}，
复合 $U \to Y$ 形式非分歧。再由引理
\ref{spaces-more-morphisms-lemma-formally-permanence}，可知 $U \to V$ 形式非分歧。
故 (1) 蕴含 (2)，而 (2) 显然蕴含 (3)。% P09-ERR-0542

\medskip\noindent
假设给定 (3) 中的图。由引理
\ref{spaces-more-morphisms-lemma-representable-property-formally-property}，
态射 $V \to Y$ 形式非分歧。因此，由引理
\ref{spaces-more-morphisms-lemma-composition-formally-smooth-etale-unramified}，
复合 $U \to Y$ 形式非分歧。再由引理
\ref{spaces-more-morphisms-lemma-etale-on-top}，可知 $X \to Y$ 形式非分歧，
即 (1) 成立。
\end{proof}

\begin{lemma}
\label{spaces-more-morphisms-lemma-formally-unramified-not-affine}
设 $S$ 为概形。
若 $f : X \to Y$ 是 $S$ 上代数空间的形式非分歧态射，则对于任意实线交换图
$$
\xymatrix{
X \ar[d]_f & T \ar[d]^i \ar[l] \\
Y & T' \ar[l] \ar@{-->}[lu]
}
$$
其中 $T \subset T'$ 是 $S$ 上代数空间的一阶加厚，至多存在一个
使该图交换的虚线箭头。换言之，在定义
\ref{spaces-more-morphisms-definition-formally-unramified} 中，可去掉 $T$ 为仿射概形的条件。
\end{lemma}

\begin{proof}
这是因为存在满平展态射 $U' \to T'$，其中 $U'$ 是仿射概形的不交并
（见《空间的性质》引理
\ref{spaces-properties-lemma-cover-by-union-affines}），
而态射 $T' \to X$ 由其在 $U'$ 上的限制唯一确定。
\end{proof}

\begin{lemma}
\label{spaces-more-morphisms-lemma-composition-formally-unramified}
形式非分歧态射的复合仍形式非分歧。
\end{lemma}

\begin{proof}
这是形式地成立的。
\end{proof}

\begin{lemma}
\label{spaces-more-morphisms-lemma-base-change-formally-unramified}
形式非分歧态射的基变换仍形式非分歧。
\end{lemma}

\begin{proof}
这是形式地成立的。
\end{proof}

\begin{lemma}
\label{spaces-more-morphisms-lemma-characterize-formally-unramified}
设 $S$ 为概形。设 $f : X \to Y$ 为 $S$ 上代数空间的态射。
下列条件等价：
\begin{enumerate}
\item $f$ 形式非分歧；
\item $\Omega_{X/Y} = 0$。
\end{enumerate}
\end{lemma}

\begin{proof}
这是下列结果的组合：引理
\ref{spaces-more-morphisms-lemma-formally-unramified}、《态射进阶》引理
\ref{more-morphisms-lemma-formally-unramified-differentials}
以及引理 \ref{spaces-more-morphisms-lemma-localize-differentials}。
\end{proof}

\begin{lemma}
\label{spaces-more-morphisms-lemma-unramified-formally-unramified}
设 $S$ 为概形。
设 $f : X \to Y$ 为 $S$ 上代数空间的态射。
下列条件等价：
\begin{enumerate}
\item 态射 $f$ 非分歧；
\item 态射 $f$ 局部有限型且 $\Omega_{X/Y} = 0$；
\item 态射 $f$ 局部有限型且形式非分歧。
\end{enumerate}
\end{lemma}

\begin{proof}
选择图
$$
\xymatrix{
U \ar[d] \ar[r]_\psi & V \ar[d] \\
X \ar[r]^f & Y
}
$$
其中 $U$、$V$ 为概形，且竖直箭头平展并且满。于是有
\begin{align*}
f\text{ 非分歧}
& \Leftrightarrow
\psi\text{ 非分歧} \\
& \Leftrightarrow
\psi\text{ 局部有限型且 }\Omega_{U/V} = 0 \\
& \Leftrightarrow
f\text{ 局部有限型且 }\Omega_{X/Y} = 0 \\
& \Leftrightarrow
f\text{ 局部有限型且形式非分歧}
\end{align*}
% P09-ERR-0543
这里使用了《态射》引理
\ref{morphisms-lemma-unramified-omega-zero}
以及引理 \ref{spaces-more-morphisms-lemma-characterize-formally-unramified}。
\end{proof}

\begin{lemma}
\label{spaces-more-morphisms-lemma-universally-injective-unramified}
设 $S$ 为概形。
设 $f : X \to Y$ 为 $S$ 上代数空间的态射。
下列条件等价：
\begin{enumerate}
\item $f$ 非分歧且为单态射；
\item $f$ 非分歧且普遍单射；
\item $f$ 局部有限型且为单态射；
\item $f$ 普遍单射、局部有限型且形式非分歧。
\end{enumerate}
此外，在这种情形下，$f$ 还可表、分离且局部拟有限。
\end{lemma}

\begin{proof}
由引理 \ref{spaces-more-morphisms-lemma-unramified-formally-unramified}，
形式非分歧且局部有限型等价于非分歧。因此 (4) 等价于 (2)。
单态射当然形式非分歧，故 (3) 蕴含 (4)。% P09-ERR-0544
显然 (1) 蕴含 (3)。最后，若 (2) 成立，则
$\Delta : X \to X \times_Y X$ 既是开浸入
（《空间的态射》引理
\ref{spaces-morphisms-lemma-diagonal-unramified-morphism}），
又是满的（《空间的态射》引理
\ref{spaces-morphisms-lemma-universally-injective}），
因而是同构，即 $f$ 为单态射。由此可见 (2) 蕴含 (1)。
最后，由《空间的态射》诸引理
\ref{spaces-morphisms-lemma-monomorphism-loc-finite-type-loc-quasi-finite}
和
\ref{spaces-morphisms-lemma-locally-quasi-finite-separated-representable}，
可知 $f$ 可表、分离且局部拟有限。
\end{proof}

\begin{lemma}
\label{spaces-more-morphisms-lemma-characterize-closed-immersion}
设 $S$ 为概形。
设 $f : X \to Y$ 为 $S$ 上代数空间的态射。
下列条件等价：
\begin{enumerate}
\item $f$ 是闭浸入；
\item $f$ 普遍闭、非分歧且为单态射；
\item $f$ 普遍闭、非分歧且普遍单射；
\item $f$ 普遍闭、局部有限型且为单态射；
\item $f$ 普遍闭、普遍单射、局部有限型且形式非分歧。
\end{enumerate}
\end{lemma}

\begin{proof}
(2) -- (5) 的等价性立即由引理
\ref{spaces-more-morphisms-lemma-universally-injective-unramified} 得出。
此外，若 (2) -- (5) 成立，则 $f$ 可表。
同样，若 (1) 成立，则 $f$ 可表。
故结论归结为概形情形；见《平展态射》引理
\ref{etale-lemma-characterize-closed-immersion}。
\end{proof}

\section{泛一阶加厚}
\label{spaces-more-morphisms-section-universal-thickening}

\noindent
设 $S$ 为概形。
设 $h : Z \to X$ 为 $S$ 上代数空间的态射。
$Z$ 在 $X$ 上的一个{\it 泛一阶加厚}，是一个一阶加厚
$Z \subset Z'$（在 $X$ 上），使得给定任意一阶加厚
$T \subset T'$（在 $X$ 上）及实线交换图
\begin{equation}
\label{spaces-more-morphisms-equation-universal-first-order-thickening}
\vcenter{
\xymatrix{
& Z \ar[ld] & & T \ar[rd] \ar[ll]^a \\
Z' \ar[rrd] & & & & T' \ar@{..>}[llll]_{a'} \ar[lld]^b \\
 & & X
}
}
\end{equation}
都存在唯一的虚线箭头使该图交换。注意，在这种情形下，
$(a, a') : (T \subset T') \to (Z \subset Z')$
是 $X$ 上加厚之间的态射。因此，若泛一阶加厚存在，
则它在唯一同构意义下唯一。一般而言，泛一阶加厚未必存在；
但若 $h$ 形式非分歧，则它确实存在。在证明这一点之前，
先说明概形范畴中的泛一阶加厚也是代数空间范畴中的泛一阶加厚。

\begin{lemma}
\label{spaces-more-morphisms-lemma-check-universal-first-order-thickening}
设 $S$ 为概形。
设 $h : Z \to X$ 为 $S$ 上代数空间的态射。
设 $Z \subset Z'$ 为 $X$ 上的一阶加厚。
下列条件等价：% P09-ERR-0545
\begin{enumerate}
\item $Z \subset Z'$ 是泛一阶加厚；
\item 对任意图 (\ref{spaces-more-morphisms-equation-universal-first-order-thickening})，
若 $T'$ 为概形，则存在唯一的虚线箭头使该图交换；
\item 对任意图 (\ref{spaces-more-morphisms-equation-universal-first-order-thickening})，
若 $T'$ 为仿射概形，则存在唯一的虚线箭头使该图交换。
\end{enumerate}
\end{lemma}

\begin{proof}
蕴含关系 (1) $\Rightarrow$ (2) $\Rightarrow$ (3) 是形式的。
假设 (3)，并给定任意一个图
(\ref{spaces-more-morphisms-equation-universal-first-order-thickening})。% P09-ERR-0546
选择展示 $T' = U'/R'$，见《空间》定义
\ref{spaces-definition-presentation}。
可假设 $U' = \coprod U'_i$ 是仿射概形的不交并，于是
$R' = U' \times_{T'} U' = \coprod_{i, j} U'_i \times_T' U'_j$。% P09-ERR-0547
对每一对 $(i, j)$，选择仿射开覆盖
$U'_i \times_T' U'_j = \bigcup_k R'_{ijk}$。
用 $U_i, R_{ijk}$ 表示它们与 $T$ 在 $T'$ 上的纤维积。
于是每个 $U_i \subset U'_i$ 和 $R_{ijk} \subset R'_{ijk}$
都是仿射概形的一阶加厚。用 $a_i : U_i \to Z$，相应地用
$a_{ijk} : R_{ijk} \to Z$，表示 $a : T \to Z$ 分别与态射
$U_i \to T$、$R_{ijk} \to T$ 的复合。
将 (3) 应用于 $a_i : U_i \to Z$，得到唯一态射
$a'_i : U'_i \to Z'$。将 (3) 应用于 $a_{ijk}$，可知两个复合
$R'_{ijk} \to R'_i \to Z'$ 和 $R'_{ijk} \to R'_j \to Z'$ 相等。% P09-ERR-0548
因此 $a' = \coprod a'_i : U' = \coprod U'_i \to Z'$
下降到商层 $T' = U'/R'$，结论成立。
\end{proof}

\begin{lemma}
\label{spaces-more-morphisms-lemma-universal-thickening-over-formally-etale}
设 $S$ 为概形。
设 $Z \to Y \to X$ 为 $S$ 上代数空间的态射。
若 $Z \subset Z'$ 是 $Z$ 在 $Y$ 上的泛一阶加厚，且
$Y \to X$ 形式平展，则 $Z \subset Z'$ 是 $Z$ 在 $X$ 上的
泛一阶加厚。
\end{lemma}

\begin{proof}
这是形式的。事实上，由引理
\ref{spaces-more-morphisms-lemma-check-universal-first-order-thickening}，
只需考虑 $T'$ 为仿射概形的实线交换图
(\ref{spaces-more-morphisms-equation-universal-first-order-thickening})。
复合 $T \to Z \to Y$ 唯一地提升为 $T' \to Y$，
因为假设 $Y \to X$ 形式平展。于是 $Z \subset Z'$ 是 $Y$ 上的
泛一阶加厚这一事实给出所需态射 $a' : T' \to Z'$。
\end{proof}

\begin{lemma}
\label{spaces-more-morphisms-lemma-etale-morphism-of-universal-thickenings}
设 $S$ 为概形。
设 $Z \to Y \to X$ 为 $S$ 上代数空间的态射。
假设 $Z \to Y$ 平展。
\begin{enumerate}
\item 若 $Y \subset Y'$ 是 $Y$ 在 $X$ 上的泛一阶加厚，
则唯一平展态射 $Z' \to Y'$，即满足 $Z = Y \times_{Y'} Z'$ 的那个态射，
（见定理 \ref{spaces-more-morphisms-theorem-topological-invariance}）
是 $Z$ 在 $X$ 上的泛一阶加厚。
\item 若 $Z \to Y$ 满，并且
$(Z \subset Z') \to (Y \subset Y')$ 是 $X$ 上一阶加厚之间的
平展态射，而 $Z'$ 是 $Z$ 在 $X$ 上的泛一阶加厚，
则 $Y'$ 是 $Y$ 在 $X$ 上的泛一阶加厚。
\end{enumerate}
\end{lemma}

\begin{proof}
先证 (1)。由引理
\ref{spaces-more-morphisms-lemma-check-universal-first-order-thickening}，
只需考虑 $T'$ 为仿射概形的实线交换图
(\ref{spaces-more-morphisms-equation-universal-first-order-thickening})。
复合 $T \to Z \to Y$ 唯一地提升为 $T' \to Y'$，
因为 $Y'$ 是泛一阶加厚。然后 $Z' \to Y'$ 平展这一事实
意味着（见引理 \ref{spaces-more-morphisms-lemma-representable-property-formally-property}）
$T' \to Y'$ 可提升为所需态射 $a' : T' \to Z'$。

\medskip\noindent
再证 (2)。设 $T \subset T'$ 为 $X$ 上的一阶加厚，
并设 $a : T \to Y$ 为态射。令 $W = T \times_Y Z$，
并以 $c : W \to Z$ 表示投影。% P09-ERR-0549
设 $W' \to T'$ 为满足 $W = T \times_{T'} W'$ 的唯一平展态射，
见定理 \ref{spaces-more-morphisms-theorem-topological-invariance}。
注意，$W' \to T'$ 满，因为 $Z \to Y$ 满。
由假设，得到唯一态射 $c' : W' \to Z'$，它在 $X$ 上，
且其限制是 $c$（在 $W$ 上）。由唯一性，$c'$ 在 $W' \times_{T'} W'$ 上的
两个限制相等（因为 $c$ 在 $W \times_T W$ 上的两个限制相等）。
因此 $c'$ 下降为唯一态射 $a' : T' \to Y'$，结论成立。
\end{proof}

\begin{lemma}
\label{spaces-more-morphisms-lemma-universal-thickening}
设 $S$ 为概形。
设 $h : Z \to X$ 为 $S$ 上代数空间的形式非分歧态射。
存在泛一阶加厚 $Z \subset Z'$，它是 $Z$ 在 $X$ 上的加厚。
\end{lemma}

\begin{proof}
任取交换图
$$
\xymatrix{
V \ar[d] \ar[r] & U \ar[d] \\
Z \ar[r] & X
}
$$
其中 $V$、$U$ 为概形，且竖直箭头平展。
注意，$V \to U$ 是概形的形式非分歧态射，见引理
\ref{spaces-more-morphisms-lemma-formally-unramified}。
合并引理 \ref{spaces-more-morphisms-lemma-check-universal-first-order-thickening}
与《态射进阶》引理
\ref{more-morphisms-lemma-universal-thickening}，
可知存在泛一阶加厚 $V \subset V'$，它是 $V$ 在 $U$ 上的加厚。
由引理 \ref{spaces-more-morphisms-lemma-universal-thickening-over-formally-etale} 第 (1) 项，% P09-ERR-0550
$V'$ 是 $V$ 在 $X$ 上的泛一阶加厚。

\medskip\noindent
固定概形 $U$ 及满平展态射 $U \to X$。上述论证表明，
对于任意平展态射 $V \to Z$，其中 $V$ 为概形，
若 $V \to Z \to X$ 分解通过 $U$，则存在泛一阶加厚
$V \subset V'$，它是 $V$ 在 $X$ 上的加厚（但它不依赖于所选的
$V \to X$ 通过 $U$ 的分解）。现在可以选择 $V$，使得
$V \to Z$ 满平展（见《空间》引理
\ref{spaces-lemma-lift-morphism-presentations}）。
于是 $R = V \times_Z V$ 是一个平展于 $Z$ 的概形，
且 $R \to X$ 也分解通过 $U$。% P09-ERR-0551
因此得到泛一阶加厚 $V \subset V'$ 和 $R \subset R'$，
它们都在 $X$ 上。由于 $V \subset V'$
是泛一阶加厚，两个投影 $s, t : R \to V$ 提升为态射
$s', t': R' \to V'$。由引理
\ref{spaces-more-morphisms-lemma-etale-morphism-of-universal-thickenings}，
因为 $R'$ 是 $R$ 在 $X$ 上的泛一阶加厚，这些态射平展。
于是 $(t', s') : R' \to V'$ 是平展等价关系，% P09-ERR-0552
可以令 $Z' = V'/R'$。由于 $V' \to Z'$ 满平展，且
$v'$ 是 $V$ 在 $X$ 上的泛一阶加厚，% P09-ERR-0553
由引理 \ref{spaces-more-morphisms-lemma-universal-thickening-over-formally-etale} 第 (2) 项，% P09-ERR-0554
可知 $Z'$ 是 $Z$ 在 $X$ 上的泛一阶加厚。
\end{proof}

\begin{definition}
\label{spaces-more-morphisms-definition-universal-thickening}
设 $S$ 为概形。
设 $h : Z \to X$ 为 $S$ 上代数空间的形式非分歧态射。
\begin{enumerate}
\item $Z$ 在 $X$ 上的{\it 泛一阶加厚}，是引理
\ref{spaces-more-morphisms-lemma-universal-thickening} 中构造的加厚 $Z \subset Z'$。
\item $Z$ 在 $X$ 上的{\it 余法层}，是 $Z$ 在其泛一阶加厚
$Z'$ 中的余法层；该加厚定义在 $X$ 上。
\end{enumerate}
在这种情形下，我们常将余法层记为 $\mathcal{C}_{Z/X}$。
\end{definition}

\noindent
因此，有层的短正合列
$$
0 \to \mathcal{C}_{Z/X} \to \mathcal{O}_{Z'} \to \mathcal{O}_Z \to 0
$$
定义在 $Z_\etale$ 上，且 $\mathcal{C}_{Z/X}$ 是拟凝聚
$\mathcal{O}_Z$-模。
下列引理证明，当 $Z \to X$ 是浸入时，本定义与相应定义并不冲突。

\begin{lemma}
\label{spaces-more-morphisms-lemma-immersion-universal-thickening}
设 $S$ 为概形。
设 $i : Z \to X$ 为 $S$ 上代数空间的浸入。则
\begin{enumerate}
\item $i$ 形式非分歧；
\item $Z$ 在 $X$ 上的泛一阶加厚是定义
\ref{spaces-more-morphisms-definition-first-order-infinitesimal-neighbourhood}
中的 $Z$ 在 $X$ 中的一阶无穷小邻域；
\item 定义 \ref{spaces-more-morphisms-definition-conormal-sheaf} 意义下 $i$ 的余法层
与定义 \ref{spaces-more-morphisms-definition-universal-thickening} 意义下 $i$ 的余法层相同。
\end{enumerate}
\end{lemma}

\begin{proof}
代数空间的浸入按定义是可表态射。因此，由《态射》诸引理
\ref{morphisms-lemma-open-immersion-unramified} 和
\ref{morphisms-lemma-closed-immersion-unramified}，
通过《空间》引理
\ref{spaces-lemma-representable-transformations-property-implication}
这一抽象原理，浸入是非分歧的。故由引理
\ref{spaces-more-morphisms-lemma-unramified-formally-unramified}，它形式非分歧。
其余断言由引理
\ref{spaces-more-morphisms-lemma-first-order-infinitesimal-neighbourhood}、
\ref{spaces-more-morphisms-lemma-infinitesimal-neighbourhood-conormal} 与各定义合并得出。
\end{proof}

\begin{lemma}
\label{spaces-more-morphisms-lemma-universal-thickening-unramified}
设 $S$ 为概形。
设 $Z \to X$ 为 $S$ 上代数空间的形式非分歧态射。
则泛一阶加厚 $Z'$ 在 $X$ 上形式非分歧。
\end{lemma}

\begin{proof}
设 $T \subset T'$ 为 $X$ 上仿射概形的一阶加厚。设
$$
\xymatrix{
Z' \ar[d] & T \ar[l]^c \ar[d] \\
X & T' \ar[l] \ar[lu]^{a, b}
}
$$
为交换图。令 $T_0 = c^{-1}(Z) \subset T$，并令
$T'_a = a^{-1}(Z)$（概形论意义下）。
由于 $Z'$ 是 $Z$ 的一阶加厚，可知 $T'$ 是 $T'_a$ 的一阶加厚。
又因 $c = a|_T$，可知 $T_0 = T \cap T'_a$（概形论意义下）。
由于 $T'$ 是 $T$ 的一阶加厚，故 $T'_a$ 是 $T_0$ 的一阶加厚。
现在 $a|_{T'_a}$ 和 $b|_{T'_a}$ 是从 $T'_a$ 到 $Z'$
的态射，并且在 $X$ 上；它们在 $T_0$ 上相同（作为到 $Z$ 的态射）。
因此，由 $Z'$ 的泛性质，得到
$a|_{T'_a} = b|_{T'_a}$。于是 $a$、$b$ 是从一阶加厚
$T'$ 出发的态射；这个加厚是 $T'_a$ 的，并且它们在 $T'_a$
上的限制作为到 $Z$ 的态射相同。% P09-ERR-0555
再次使用 $Z'$ 的泛性质，得到 $a = b$。换言之，
形式非分歧态射的定义性质对 $Z' \to X$ 成立，正如所需。
\end{proof}

\begin{lemma}
\label{spaces-more-morphisms-lemma-universal-thickening-functorial}
设 $S$ 为概形。% P09-ERR-0556
考虑 $S$ 上代数空间的交换图
$$
\xymatrix{
Z \ar[r]_h \ar[d]_f & X \ar[d]^g \\
W \ar[r]^{h'} & Y
}
$$
其中 $h$ 和 $h'$ 形式非分歧。设 $Z \subset Z'$ 是 $Z$ 在 $X$
上的泛一阶加厚，设 $W \subset W'$ 是 $W$ 在 $Y$ 上的泛一阶加厚。
存在加厚之间的典范态射
$(f, f') : (Z, Z') \to (W, W')$，它在 $Y$ 上，
并嵌入下列交换图：
$$
\xymatrix{
& & & Z' \ar[ld] \ar[d]^{f'} \\
Z \ar[rr] \ar[d]_f \ar[rrru] & & X \ar[d] & W' \ar[ld] \\
W \ar[rrru]|!{[rr];[rruu]}\hole \ar[rr] & & Y
}
$$
特别地，加厚之间的态射 $(f, f')$ 诱导余法层之间的态射
$f^*\mathcal{C}_{W/Y} \to \mathcal{C}_{Z/X}$。
\end{lemma}

\begin{proof}
第一项断言由 $W'$ 的泛性质显然成立。
余法层上的诱导映射，是将引理
\ref{spaces-more-morphisms-lemma-conormal-functorial} 应用于
$(Z \subset Z') \to (W \subset W')$ 所得的映射。
\end{proof}

\begin{lemma}
\label{spaces-more-morphisms-lemma-universal-thickening-fibre-product}
设 $S$ 为概形。设
$$
\xymatrix{
Z \ar[r]_h \ar[d]_f & X \ar[d]^g \\
W \ar[r]^{h'} & Y
}
$$
为 $S$ 上代数空间的纤维积图，其中 $h'$ 形式非分歧。
则 $h$ 形式非分歧；并且若 $W \subset W'$ 是 $W$ 在 $Y$ 上的
泛一阶加厚，则
$Z = X \times_Y W \subset X \times_Y W'$ 是 $Z$ 在 $X$ 上的
泛一阶加厚。特别地，引理
\ref{spaces-more-morphisms-lemma-universal-thickening-functorial} 中的典范映射
$f^*\mathcal{C}_{W/Y} \to \mathcal{C}_{Z/X}$ 是满射。
\end{lemma}

\begin{proof}
由引理 \ref{spaces-more-morphisms-lemma-base-change-formally-unramified}，态射 $h$
形式非分歧。显然，$X \times_Y W'$ 是一阶加厚。
由 $W'$ 的泛性质（以及纤维积的映射性质），直接验证可知
它具有该泛性质。关于这为何蕴含余法层之间的映射为满射，
见引理 \ref{spaces-more-morphisms-lemma-conormal-functorial-flat}。
\end{proof}

\begin{lemma}
\label{spaces-more-morphisms-lemma-universal-thickening-fibre-product-flat}
设 $S$ 为概形。设
$$
\xymatrix{
Z \ar[r]_h \ar[d]_f & X \ar[d]^g \\
W \ar[r]^{h'} & Y
}
$$
为 $S$ 上代数空间的纤维积图，其中 $h'$ 形式非分歧且 $g$ 平坦。
在这种情形下，相应的泛一阶加厚之间的映射 $Z' \to W'$ 平坦，
且 $f^*\mathcal{C}_{W/Y} \to \mathcal{C}_{Z/X}$ 是同构。
\end{lemma}

\begin{proof}
平坦性在基变换下保持，见《空间的态射》引理
\ref{spaces-morphisms-lemma-base-change-flat}。
因此，第一项断言由引理
\ref{spaces-more-morphisms-lemma-universal-thickening-fibre-product} 中对 $W'$ 的描述得出。
显然，$X \times_Y W'$ 是一阶加厚。
由 $W'$ 的泛性质（以及纤维积的映射性质），直接验证可知
它具有该泛性质。关于这为何蕴含余法层之间的映射为同构，
见引理 \ref{spaces-more-morphisms-lemma-conormal-functorial-flat}。
\end{proof}

\begin{lemma}
\label{spaces-more-morphisms-lemma-universal-thickening-localize}
取泛一阶加厚的操作与平展局部化交换。更准确地说，
设 $h : Z \to X$ 为基概形 $S$ 上代数空间的形式非分歧态射。
设
$$
\xymatrix{
V \ar[d] \ar[r] & U \ar[d] \\
Z \ar[r] & X
}
$$
为竖直箭头平展的交换图。
设 $Z'$ 为 $Z$ 在 $X$ 上的泛一阶加厚。
则 $V \to U$ 形式非分歧，并且泛一阶加厚 $V'$（即 $V$
在 $U$ 上的那个）平展于 $Z'$。特别地，
$\mathcal{C}_{Z/X}|_V = \mathcal{C}_{V/U}$。
\end{lemma}

\begin{proof}
第一项断言就是引理 \ref{spaces-more-morphisms-lemma-formally-unramified}。
泛一阶加厚的相容性由引理
\ref{spaces-more-morphisms-lemma-universal-thickening-over-formally-etale} 和
\ref{spaces-more-morphisms-lemma-etale-morphism-of-universal-thickenings} 得出。
\end{proof}

\begin{lemma}
\label{spaces-more-morphisms-lemma-differentials-universally-unramified}
设 $S$ 为概形，$B$ 为 $S$ 上的代数空间。
设 $h : Z \to X$ 为 $B$ 上代数空间的形式非分歧态射。
设 $Z \subset Z'$ 为 $Z$ 在 $X$ 上的泛一阶加厚，
其结构态射为 $h' : Z' \to X$。典范映射
$$
\text{d}h' : (h')^*\Omega_{X/B} \to \Omega_{Z'/B}
$$
% P09-ERR-0557
诱导同构
$h^*\Omega_{X/B} \to \Omega_{Z'/B} \otimes \mathcal{O}_Z$。
\end{lemma}

\begin{proof}
映射 $c_{h'}$ 是引理
\ref{spaces-more-morphisms-lemma-functoriality-differentials} 中定义的映射。
若 $i : Z \to Z'$ 是给定闭浸入，则 $i^*c_{h'}$ 是映射
$h^*\Omega_{X/S} \to \Omega_{Z'/S} \otimes \mathcal{O}_Z$。% P09-ERR-0558
通过平展局部化，验证它是同构可归结为概形情形，见引理
\ref{spaces-more-morphisms-lemma-universal-thickening-localize} 和
\ref{spaces-more-morphisms-lemma-localize-differentials}。
在这种情形下，结论就是《态射进阶》引理
\ref{more-morphisms-lemma-differentials-universally-unramified}。
\end{proof}

\begin{lemma}
\label{spaces-more-morphisms-lemma-universally-unramified-differentials-sequence}
设 $S$ 为概形，$B$ 为 $S$ 上的代数空间。
设 $h : Z \to X$ 为 $B$ 上代数空间的形式非分歧态射。
有典范正合列
$$
\mathcal{C}_{Z/X} \to h^*\Omega_{X/B} \to \Omega_{Z/B} \to 0.
$$
第一支箭头由 $\text{d}_{Z'/B}$ 诱导，其中 $Z'$ 是 $Z$ 在 $X$
上的泛一阶邻域。% P09-ERR-0559
\end{lemma}

\begin{proof}
已知有典范正合列
$$
\mathcal{C}_{Z/Z'} \to
\Omega_{Z'/S} \otimes \mathcal{O}_Z \to
\Omega_{Z/S} \to 0.
$$
% P09-ERR-0560
见引理 \ref{spaces-more-morphisms-lemma-differentials-relative-immersion}。
因此，应用引理 \ref{spaces-more-morphisms-lemma-differentials-universally-unramified}
即得结论。
\end{proof}

\begin{lemma}
\label{spaces-more-morphisms-lemma-two-unramified-morphisms}
设 $S$ 为概形。设
$$
\xymatrix{
Z \ar[r]_i \ar[rd]_j & X \ar[d] \\
& Y
}
$$
为 $S$ 上代数空间的交换图，其中 $i$、$j$ 形式非分歧。
则有典范正合列
$$
\mathcal{C}_{Z/Y} \to
\mathcal{C}_{Z/X} \to
i^*\Omega_{X/Y} \to 0
$$
其中第一支箭头来自引理
\ref{spaces-more-morphisms-lemma-universal-thickening-functorial}，
第二支箭头来自引理
\ref{spaces-more-morphisms-lemma-universally-unramified-differentials-sequence}。
\end{lemma}

\begin{proof}
这些映射已经定义，故通过平展局部化，验证该序列正合
可归结为概形情形，见引理
\ref{spaces-more-morphisms-lemma-universal-thickening-localize} 和
\ref{spaces-more-morphisms-lemma-localize-differentials}。
在这种情形下，结论就是《态射进阶》引理
\ref{more-morphisms-lemma-two-unramified-morphisms}。
\end{proof}

\begin{lemma}
\label{spaces-more-morphisms-lemma-transitivity-conormal-unramified}
设 $S$ 为概形。
设 $Z \to Y \to X$ 为 $S$ 上代数空间的形式非分歧态射。
\begin{enumerate}
\item 若 $Z \subset Z'$ 是 $Z$ 在 $X$ 上的泛一阶加厚，
$Y \subset Y'$ 是 $Y$ 在 $X$ 上的泛一阶加厚，则存在态射
$Z' \to Y'$，并且 $Y \times_{Y'} Z'$ 是 $Z$ 在 $Y$ 上的
泛一阶加厚。
\item 有典范正合列
$$
i^*\mathcal{C}_{Y/X} \to
\mathcal{C}_{Z/X} \to
\mathcal{C}_{Z/Y} \to 0
$$
其中映射来自引理 \ref{spaces-more-morphisms-lemma-universal-thickening-functorial}，
而 $i : Z \to Y$ 是第一个态射。
\end{enumerate}
\end{lemma}

\begin{proof}
(1) 中的映射 $h : Z' \to Y'$ 来自引理
\ref{spaces-more-morphisms-lemma-universal-thickening-functorial}。
$Y \times_{Y'} Z'$ 是 $Z$ 在 $Y$ 上的泛一阶加厚这一断言，
由 $Z'$ 和 $Y'$ 的泛性质显然成立。由引理
\ref{spaces-more-morphisms-lemma-transitivity-conormal}，有正合列
$$
(i')^*\mathcal{C}_{Y \times_{Y'} Z'/Z'} \to
\mathcal{C}_{Z/Z'} \to
\mathcal{C}_{Z/Y \times_{Y'} Z'} \to 0
$$
其中 $i' : Z \to Y \times_{Y'} Z'$ 是给定态射。由引理
\ref{spaces-more-morphisms-lemma-conormal-functorial-flat}，存在满射
$h^*\mathcal{C}_{Y/Y'} \to \mathcal{C}_{Y \times_{Y'} Z'/Z'}$。
结合等式
$\mathcal{C}_{Y/Y'} = \mathcal{C}_{Y/X}$、
$\mathcal{C}_{Z/Z'} = \mathcal{C}_{Z/X}$ 以及
$\mathcal{C}_{Z/Y \times_{Y'} Z'} = \mathcal{C}_{Z/Y}$，
即得引理。
\end{proof}

\section{形式平展态射}
\label{spaces-more-morphisms-section-formally-etale}

\noindent
本节具体说明代数空间态射形式平展的含义。

\begin{definition}
\label{spaces-more-morphisms-definition-formally-etale}
设 $S$ 为概形。设 $f : X \to Y$ 为代数空间的态射，并且它在 $S$ 上。
若它作为函子间的变换是形式平展的，则称它为{\it 形式平展的}；
其含义见定义 \ref{spaces-more-morphisms-definition-formally-smooth-etale-unramified}。
\end{definition}

\noindent
我们不再重述在函子间形式平展变换这一更一般的框架中证明的结果，
参见第 \ref{spaces-more-morphisms-section-formally-smooth-etale-unramified} 节。

\begin{lemma}
\label{spaces-more-morphisms-lemma-formally-etale}
设 $S$ 为概形。设 $f : X \to Y$ 为代数空间的态射，并且它在
$S$ 上。下列条件等价：
\begin{enumerate}
\item $f$ 形式平展；
\item 对每个图
$$
\xymatrix{
U \ar[d] \ar[r]_\psi & V \ar[d] \\
X \ar[r]^f & Y
}
$$
其中 $U$、$V$ 为概形，且竖直箭头平展，概形态射 $\psi$
是形式平展的（参见《态射进阶》定义
\ref{more-morphisms-definition-formally-etale}）；
\item 对于一个这样的图，若竖直箭头满，则 $\psi$ 形式平展。
\end{enumerate}
\end{lemma}

\begin{proof}
假设 $f$ 形式平展。由引理
\ref{spaces-more-morphisms-lemma-representable-property-formally-property}，
态射 $U \to X$ 和 $V \to Y$ 形式平展。因此，由引理
\ref{spaces-more-morphisms-lemma-composition-formally-smooth-etale-unramified}，
复合 $U \to Y$ 形式平展。再由引理
\ref{spaces-more-morphisms-lemma-formally-permanence}，可知 $U \to V$ 形式平展。
故 (1) 蕴含 (2)，而 (2) 显然蕴含 (3)。% P09-ERR-0561

\medskip\noindent
假设给定 (3) 中的图。由引理
\ref{spaces-more-morphisms-lemma-representable-property-formally-property}，
态射 $V \to Y$ 形式平展。因此，由引理
\ref{spaces-more-morphisms-lemma-composition-formally-smooth-etale-unramified}，
复合 $U \to Y$ 形式平展。再由引理
\ref{spaces-more-morphisms-lemma-etale-on-top}，可知 $X \to Y$ 形式平展，即 (1) 成立。
\end{proof}

\begin{lemma}
\label{spaces-more-morphisms-lemma-formally-etale-not-affine}
设 $S$ 为概形。
设 $f : X \to Y$ 为代数空间的形式平展态射，并且它在 $S$ 上。
则给定任意实线交换图
$$
\xymatrix{
X \ar[d]_f & T \ar[d]^i \ar[l]_a \\
Y & T' \ar[l] \ar@{-->}[lu]
}
$$
其中 $T \subset T'$ 是 $Y$ 上代数空间的一阶加厚，
恰有一个虚线箭头使该图交换。换言之，在定义
\ref{spaces-more-morphisms-definition-formally-etale} 中，可去掉 $T$ 为仿射概形的条件。
\end{lemma}

\begin{proof}
设 $U' \to T'$ 为满平展态射，其中
$U' = \coprod U'_i$ 是仿射概形的不交并。令
$U_i = T \times_{T'} U'_i$。于是得到态射
$a'_i : U'_i \to X$，使得 $a'_i|_{U_i}$ 等于复合
$U_i \to T \to X$。由唯一性（见引理
\ref{spaces-more-morphisms-lemma-formally-unramified-not-affine}），
可知 $a'_i$ 和 $a'_j$ 在纤维积
$U'_i \times_{T'} U'_j$ 上相同。因此
$\coprod a'_i : U' \to X$ 下降为唯一态射 $a' : T' \to X$。
\end{proof}

\begin{lemma}
\label{spaces-more-morphisms-lemma-composition-formally-etale}
形式平展态射的复合仍形式平展。
\end{lemma}

\begin{proof}
这是形式地成立的。
\end{proof}

\begin{lemma}
\label{spaces-more-morphisms-lemma-base-change-formally-etale}
形式平展态射的基变换仍形式平展。
\end{lemma}

\begin{proof}
这是形式地成立的。
\end{proof}

\begin{lemma}
\label{spaces-more-morphisms-lemma-characterize-formally-etale}
设 $S$ 为概形。
设 $f : X \to Y$ 为代数空间的态射，并且它在 $S$ 上。% P09-ERR-0562
下列条件等价：
\begin{enumerate}
\item $f$ 形式平展；
\item $f$ 形式非分歧，且 $X$ 在 $Y$ 上的泛一阶加厚等于 $X$；
\item $f$ 形式非分歧，且 $\mathcal{C}_{X/Y} = 0$；
\item $\Omega_{X/Y} = 0$ 且 $\mathcal{C}_{X/Y} = 0$。
\end{enumerate}
\end{lemma}

\begin{proof}
事实上，最后一项断言之所以有意义，只是因为
$\Omega_{X/Y} = 0$ 蕴含 $\mathcal{C}_{X/Y}$ 可通过引理
\ref{spaces-more-morphisms-lemma-characterize-formally-unramified} 和定义
\ref{spaces-more-morphisms-definition-universal-thickening} 加以定义。% P09-ERR-0563
这也说明 (3) 和 (4) 等价。

\medskip\noindent
(1)、(2)、(3) 中任一假设都蕴含 $f$ 形式非分歧。% P09-ERR-0564
故可假设 $f$ 形式非分歧。(1)、(2)、(3) 的等价性由
泛一阶加厚 $X'$（即 $X$ 在 $S$ 上的那个）的泛性质，% P09-ERR-0565 % P09-ERR-0566
以及下述事实得出：
$X = X' \Leftrightarrow \mathcal{C}_{X/Y} = 0$；毕竟按定义，
$\mathcal{C}_{X/Y} = \mathcal{C}_{X/X'}$ 是 $X$ 在 $X'$ 中的理想层。
\end{proof}

\begin{lemma}
\label{spaces-more-morphisms-lemma-unramified-flat-formally-etale}
非分歧平坦态射是形式平展的。
\end{lemma}

\begin{proof}
这由概形情形、《态射进阶》引理
\ref{more-morphisms-lemma-unramified-flat-formally-etale}，
以及平展局部化得出；对于后者，见引理
\ref{spaces-more-morphisms-lemma-formally-unramified}、\ref{spaces-more-morphisms-lemma-formally-etale}
以及《空间的态射》引理
\ref{spaces-morphisms-lemma-flat-local}。
\end{proof}

\begin{lemma}
\label{spaces-more-morphisms-lemma-etale-formally-etale}
设 $S$ 为概形。
设 $f : X \to Y$ 为 $S$ 上代数空间的态射。
下列条件等价：
\begin{enumerate}
\item 态射 $f$ 平展；
\item 态射 $f$ 局部有限表示且形式平展。
\end{enumerate}
\end{lemma}

\begin{proof}
这由概形情形、《态射进阶》引理
\ref{more-morphisms-lemma-etale-formally-etale}，
以及平展局部化得出；对于后者，见引理
\ref{spaces-more-morphisms-lemma-formally-etale}
以及《空间的态射》诸引理
\ref{spaces-morphisms-lemma-finite-presentation-local} 和
\ref{spaces-morphisms-lemma-etale-local}。
\end{proof}

\section{映射的无穷小形变}
\label{spaces-more-morphisms-section-action-by-derivations}

\noindent
本节说明如何用导子无穷小地移动一个映射。全节都使用如下事实：
加厚 $X'$（它是 $X$ 的加厚）上的层可视为 $X$ 上的层，见等式
(\ref{spaces-more-morphisms-equation-equivalence-etale-spaces}) 和
(\ref{spaces-more-morphisms-equation-fundamental-equivalence})。

\begin{lemma}
\label{spaces-more-morphisms-lemma-difference-derivation}
设 $S$ 为概形，$B$ 为 $S$ 上的代数空间。
设 $X \subset X'$ 和 $Y \subset Y'$ 是 $B$ 上代数空间的两个一阶加厚。
设
$(a, a'), (b, b') : (X \subset X') \to (Y \subset Y')$
是 $B$ 上加厚之间的两个态射。假设
\begin{enumerate}
\item $a = b$；
\item 两个映射 $a^*\mathcal{C}_{Y/Y'} \to \mathcal{C}_{X/X'}$
（引理 \ref{spaces-more-morphisms-lemma-conormal-functorial}）相等。
\end{enumerate}
则映射 $(a')^\sharp - (b')^\sharp$ 分解为
$$
\mathcal{O}_{Y'} \to \mathcal{O}_Y \xrightarrow{D}
a_*\mathcal{C}_{X/X'} \to a_*\mathcal{O}_{X'}
$$
其中 $D$ 是 $\mathcal{O}_B$-导子。
\end{lemma}

\begin{proof}
我们不在 $Y$ 上工作，而在 $X$ 上工作，其优点是拉回函子
$a^{-1}$ 正合。使用 (1) 和 (2)，得到行正合的交换图
$$
\xymatrix{
0 \ar[r] &
\mathcal{C}_{X/X'} \ar[r] &
\mathcal{O}_{X'} \ar[r] &
\mathcal{O}_X \ar[r] & 0 \\
0 \ar[r] &
a^{-1}\mathcal{C}_{Y/Y'} \ar[r] \ar[u] &
a^{-1}\mathcal{O}_{Y'}
\ar[r] \ar@<1ex>[u]^{(a')^\sharp} \ar@<-1ex>[u]_{(b')^\sharp} &
a^{-1}\mathcal{O}_Y \ar[r] \ar[u] & 0
}
$$
一般而言，在这种情形下，$\mathcal{O}_B$-代数映射
$(a')^\sharp$ 与 $(b')^\sharp$ 之差，是一个 $\mathcal{O}_B$-导子，
即从 $a^{-1}\mathcal{O}_Y$ 到 $\mathcal{C}_{X/X'}$ 的导子。
由函子 $a^{-1}$ 与 $a_*$ 的伴随性，这等同于一个
$\mathcal{O}_B$-导子，即从 $\mathcal{O}_Y$ 到
$a_*\mathcal{C}_{X/X'}$ 的导子。略去若干细节。
\end{proof}

\noindent
注意，在上述引理的情形中，可将 $D$ 写成
\begin{equation}
\label{spaces-more-morphisms-equation-D}
D = \text{d}_{Y/B} \circ \theta
\end{equation}
% P09-ERR-0567
其中 $\theta$ 是 $\mathcal{O}_Y$-线性映射
$\theta : \Omega_{Y/B} \to a_*\mathcal{C}_{X/X'}$。
当然，再次利用伴随性，也可将 $\theta$ 视为 $\mathcal{O}_X$-线性映射
$\theta : a^*\Omega_{Y/B} \to \mathcal{C}_{X/X'}$。

\begin{lemma}
\label{spaces-more-morphisms-lemma-action-by-derivations}
设 $S$ 为概形，$B$ 为 $S$ 上的代数空间。
设 $(a, a') : (X \subset X') \to (Y \subset Y')$
为 $B$ 上一阶加厚之间的态射。设
$$
\theta : a^*\Omega_{Y/B} \to \mathcal{C}_{X/X'}
$$
为 $\mathcal{O}_X$-线性映射。则存在唯一的序对态射
$(b, b') : (X \subset X') \to (Y \subset Y')$，
使得引理 \ref{spaces-more-morphisms-lemma-difference-derivation} 的 (1)、(2) 成立，
并且导子 $D$ 与 $\theta$ 由等式 (\ref{spaces-more-morphisms-equation-D}) 联系。
\end{lemma}

\begin{proof}
考虑映射
$$
\alpha = (a')^\sharp + D : a^{-1}\mathcal{O}_{Y'} \to \mathcal{O}_{X'}
$$
其中 $D$ 如等式 (\ref{spaces-more-morphisms-equation-D}) 所示。因为 $D$ 是
$\mathcal{O}_B$-导子，所以 $\alpha$ 是 $\mathcal{O}_B$-代数层之间的映射。
由构造，$i_X^\sharp \circ \alpha = a^\sharp \circ i_Y^\sharp$，其中
$i_X : X \to X'$ 和 $i_Y : Y \to Y'$ 是给定闭浸入。由引理
\ref{spaces-more-morphisms-lemma-first-order-thickening-maps}，得到加厚之间的唯一态射
$(a, b') : (X  \subset X') \to (Y \subset Y')$，它在 $B$ 上，
使得 $\alpha = (b')^\sharp$。令 $b = a$ 即得结论。
\end{proof}

\begin{remark}
\label{spaces-more-morphisms-remark-action-by-derivations}
沿用引理 \ref{spaces-more-morphisms-lemma-action-by-derivations} 的假设与记号。
局部截面 $\theta$ 在 $a'$ 上的作用有时记作
$\theta \cdot a'$。这无非意味着 $(a')^\sharp$ 与
$(\theta \cdot a')^\sharp$ 相差一个导子 $D$，
而该导子与 $\theta$ 由等式 (\ref{spaces-more-morphisms-equation-D}) 联系。
\end{remark}

\begin{lemma}
\label{spaces-more-morphisms-lemma-sheaf}
设 $S$ 为概形，$B$ 为 $S$ 上的代数空间。
设 $X \subset X'$ 和 $Y \subset Y'$ 是 $B$ 上的一阶加厚。
假设给定态射 $a : X \to Y$ 以及映射
$A : a^*\mathcal{C}_{Y/Y'} \to \mathcal{C}_{X/X'}$；
这是一个 $\mathcal{O}_X$-模的映射。对对象 $U'$（它属于
$(X')_{spaces, \etale}$），令
$U = X \times_{X'} U'$，并考虑满足下列条件的态射
$a' : U' \to Y'$：
\begin{enumerate}
\item $a'$ 是 $B$ 上的态射；
\item $a'|_U = a|_U$；
\item 诱导映射
$a^*\mathcal{C}_{Y/Y'}|_U \to \mathcal{C}_{X/X'}|_U$
是 $A$ 在 $U$ 上的限制。
\end{enumerate}
则规则
\begin{equation}
\label{spaces-more-morphisms-equation-sheaf}
U' \mapsto
\{a' : U' \to Y'\text{ 满足 (1)、(2)、(3)。}\}
\end{equation}
定义 $(X')_{spaces, \etale}$ 上的集合层。
\end{lemma}

\begin{proof}
用 $\mathcal{F}$ 表示引理中的规则。映射
$\mathcal{F}(U') \to \mathcal{F}(V')$（其中
$V' \subset U' \subset X'$）作为 $\mathcal{F}$ 的限制映射，实际上就是
$a' \mapsto a'|_{V'}$。在此定义下，$\mathcal{F}$ 显然是层，
因为代数空间的态射满足平展下降；见《空间上的下降》引理
\ref{spaces-descent-lemma-fpqc-universal-effective-epimorphisms}。
\end{proof}

\begin{lemma}
\label{spaces-more-morphisms-lemma-action-sheaf}
沿用引理 \ref{spaces-more-morphisms-lemma-sheaf} 的记号与假设。
我们识别 $X$ 和 $X'$ 上的层，所用等价见
(\ref{spaces-more-morphisms-equation-equivalence-etale-spaces})。
层
$$
\SheafHom_{\mathcal{O}_X}(a^*\Omega_{Y/B}, \mathcal{C}_{X/X'})
$$
作用在层 (\ref{spaces-more-morphisms-equation-sheaf}) 上。此外，对于任意对象
$U'$（它属于 $(X')_{spaces, \etale}$），若层
(\ref{spaces-more-morphisms-equation-sheaf}) 在其上有截面，则该作用简便传递。
\end{lemma}

\begin{proof}
这是下列引理的组合：
\ref{spaces-more-morphisms-lemma-difference-derivation}、
\ref{spaces-more-morphisms-lemma-action-by-derivations} 和 \ref{spaces-more-morphisms-lemma-sheaf}。
\end{proof}

\begin{remark}
\label{spaces-more-morphisms-remark-special-case}
引理 \ref{spaces-more-morphisms-lemma-difference-derivation}、
\ref{spaces-more-morphisms-lemma-action-by-derivations}、
\ref{spaces-more-morphisms-lemma-sheaf} 和
\ref{spaces-more-morphisms-lemma-action-sheaf} 的一个特殊情形是 $Y = Y'$。
此时映射 $A$ 总为零。引理 \ref{spaces-more-morphisms-lemma-sheaf} 中的层就是由规则
$$
U' \mapsto
\{a' : U' \to Y\text{ 在 }B\text{ 上，且 } a'|_U = a|_U\}
$$
给出的层，并由层
$\SheafHom_{\mathcal{O}_X}(a^*\Omega_{Y/B}, \mathcal{C}_{X/X'})$
作用于其上。
\end{remark}

\begin{remark}
\label{spaces-more-morphisms-remark-another-special-case}
引理 \ref{spaces-more-morphisms-lemma-difference-derivation}、
\ref{spaces-more-morphisms-lemma-action-by-derivations}、
\ref{spaces-more-morphisms-lemma-sheaf} 和
\ref{spaces-more-morphisms-lemma-action-sheaf} 的另一个特殊情形是：$B$ 本身是加厚
$Z \subset Z' = B$，且 $Y = Z \times_{Z'} Y'$。图示为
$$
\xymatrix{
(X \subset X') \ar@{..>}[rr]_{(a, ?)} \ar[rd]_{(g, g')} & &
(Y \subset Y') \ar[ld]^{(h, h')} \\
& (Z \subset Z')
}
$$
此时映射 $A : a^*\mathcal{C}_{Y/Y'} \to \mathcal{C}_{X/X'}$
由 $a$ 决定：映射
$h^*\mathcal{C}_{Z/Z'} \to \mathcal{C}_{Y/Y'}$ 是满射（因为假设
$Y = Z \times_{Z'} Y'$），所以拉回
$g^*\mathcal{C}_{Z/Z'} = a^*h^*\mathcal{C}_{Z/Z'} \to
a^*\mathcal{C}_{Y/Y'}$ 是满射，而复合
$g^*\mathcal{C}_{Z/Z'} \to a^*\mathcal{C}_{Y/Y'} \to \mathcal{C}_{X/X'}$
必须是由 $g'$ 诱导的典范映射。因此，引理
\ref{spaces-more-morphisms-lemma-sheaf} 中的层就是由规则
$$
U' \mapsto
\{a' : U' \to Y'\text{ 在 }Z'\text{ 上，且 } a'|_U = a|_U\}
$$
给出的层，并由层
$\SheafHom_{\mathcal{O}_X}(a^*\Omega_{Y/Z}, \mathcal{C}_{X/X'})$
作用于其上。
\end{remark}

\begin{lemma}
\label{spaces-more-morphisms-lemma-action-by-derivations-etale-localization}
设 $S$ 为概形。考虑一阶加厚的交换图
$$
\vcenter{
\xymatrix{
(T_2 \subset T_2') \ar[d]_{(h, h')} \ar[rr]_{(a_2, a_2')} & &
(X_2 \subset X_2') \ar[d]^{(f, f')} \\
(T_1 \subset T_1') \ar[rr]^{(a_1, a_1')} & &
(X_1 \subset X_1')
}
}
\quad
\begin{matrix}
\text{以及交换} \\
\text{图}
\end{matrix}
\quad
\vcenter{
\xymatrix{
X_2' \ar[r] \ar[d] & B_2 \ar[d] \\
X_1' \ar[r] & B_1
}
}
$$
它们都是 $S$ 上代数空间的图，并且 $X_2 \to X_1$ 和
$B_2 \to B_1$ 平展。对任意 $\mathcal{O}_{T_1}$-线性映射
$\theta_1 : a_1^*\Omega_{X_1/B_1} \to \mathcal{C}_{T_1/T'_1}$，
令 $\theta_2$ 为复合
$$
\xymatrix{
a_2^*\Omega_{X_2/B_2} \ar@{=}[r] &
h^*a_1^*\Omega_{X_1/B_1} \ar[r]^-{h^*\theta_1} &
h^*\mathcal{C}_{T_1/T'_1} \ar[r] &
\mathcal{C}_{T_2/T'_2}
}
$$
（等号在证明中说明）。则图
$$
\xymatrix{
T_2' \ar[rr]_{\theta_2 \cdot a_2'} \ar[d] & & X'_2 \ar[d] \\
T_1' \ar[rr]^{\theta_1 \cdot a_1'} & & X'_1
}
$$
交换，其中作用 $\theta_2 \cdot a_2'$ 和 $\theta_1 \cdot a_1'$
如评注 \ref{spaces-more-morphisms-remark-action-by-derivations} 所述。
\end{lemma}

\begin{proof}
等号来自下述识别：
$f^*\Omega_{X_1/S_1} = \Omega_{X_2/S_2}$；该识别之所以成立，
是因为微分层的构造与平展局部化（在源与靶两边）相容，见引理
\ref{spaces-more-morphisms-lemma-localize-differentials}。% P09-ERR-0568
事实上，由此有
$a_2^*\Omega_{X_2/S_2} = a_2^*f^*\Omega_{X_1/S_1} =
h^*a_1^*\Omega_{X_1/S_1}$，因为 $f \circ a_2 = a_1 \circ h$。% P09-ERR-0569
说明这一点后，可平展局部地检验该图的交换性。% P09-ERR-0570
因此可假设 $T'_i$、$X'_i$、$B_2$ 和 $B_1$ 都是概形；
在这种情形下，引理由《态射进阶》引理
\ref{more-morphisms-lemma-action-by-derivations-etale-localization} 得出。
另一种证明：由引理 \ref{spaces-more-morphisms-lemma-first-order-thickening-maps}，
只需证明 $X_1'$ 上环层的某个图交换；然后完全仿照上述
《态射进阶》引理
\ref{more-morphisms-lemma-action-by-derivations-etale-localization}
的证明，即可看出确实如此。
\end{proof}
\section{代数空间的无穷小形变}
\label{spaces-more-morphisms-section-deform}

\noindent
下面这个简单引理往往是检验映射的无穷小形变是否平坦的便利工具。

\begin{lemma}
\label{spaces-more-morphisms-lemma-deform}
设 $S$ 为概形。设
$(f, f') : (X \subset X') \to (Y \subset Y')$
是 $S$ 上代数空间一阶加厚的态射，并假设 $f$ 平坦。则下列条件等价：
\begin{enumerate}
\item $f'$ 平坦且 $X = Y \times_{Y'} X'$；
\item 典范映射 $f^*\mathcal{C}_{Y/Y'} \to \mathcal{C}_{X/X'}$
是同构。
\end{enumerate}
\end{lemma}

\begin{proof}
选取概形 $V'$ 和满平展态射 $V' \to Y'$。
选取概形 $U'$ 和满平展态射
$U' \to X' \times_{Y'} V'$。令 $U = X \times_{X'} U'$ 且
$V = Y \times_{Y'} V'$。由代数空间平坦态射的定义可知，诱导映射
$g : U \to V$ 是概形的平坦态射，并且 $f'$ 平坦当且仅当相应态射
$g' : U' \to V'$ 平坦。此外，$X = Y \times_{Y'} X'$ 当且仅当
$U = V \times_{V'} V'$。% P09-ERR-0571
最后，映射
$f^*\mathcal{C}_{Y/Y'} \to \mathcal{C}_{X/X'}$
是同构当且仅当
$g^*\mathcal{C}_{V/V'} \to \mathcal{C}_{U/U'}$ 是同构。
因此，引理由概形态射的相应结论得出；见《态射进阶》引理
\ref{more-morphisms-lemma-deform}。
\end{proof}

\noindent
下面这个引理是“按纤维判定平坦性的准则”的“幂零”版本；见
第 \ref{spaces-more-morphisms-section-criterion-flat-fibres} 节。

\begin{lemma}
\label{spaces-more-morphisms-lemma-flatness-morphism-thickenings}
设 $S$ 为概形。考虑交换图
$$
\xymatrix{
(X \subset X') \ar[rr]_{(f, f')} \ar[rd] & & (Y \subset Y') \ar[ld] \\
& (B \subset B')
}
$$
，它是 $S$ 上代数空间加厚的交换图。假设
\begin{enumerate}
\item $X'$ 在 $B'$ 上平坦；
\item $f$ 平坦；
\item $B \subset B'$ 是有限阶加厚；
\item $X = B \times_{B'} X'$ 且 $Y = B \times_{B'} Y'$。
\end{enumerate}
则 $f'$ 平坦，并且 $Y'$ 在 $B'$ 上平坦，至少在 $f'$ 的像中每一点处如此。
\end{lemma}

\begin{proof}
选取概形 $U'$ 和满平展态射 $U' \to B'$。
选取概形 $V'$ 和满平展态射
$V' \to U' \times_{B'} Y'$。
选取概形 $W'$ 和满平展态射
$W' \to V' \times_{Y'} X'$。令 $U, V, W$ 分别为
$U', V', W'$ 沿 $B \to B'$ 的基变换。于是 $f'$ 的平坦性等价于
$W' \to V'$ 的平坦性，而已知 $W \to V$ 平坦。因此可以把概形情形的
引理应用于概形加厚的图
$$
\xymatrix{
(W \subset W') \ar[rr] \ar[rd] & & (V \subset V') \ar[ld] \\
& (U \subset U')
}
$$
；见《态射进阶》引理
\ref{more-morphisms-lemma-flatness-morphism-thickenings}。
关于 $Y'/B'$ 在 $f'$ 像中各点处平坦的断言也以同样方式得出。
\end{proof}

\noindent
概形态射的许多性质在平坦形变下保持不变。

\begin{lemma}
\label{spaces-more-morphisms-lemma-deform-property}
设 $S$ 为概形。考虑交换图
$$
\xymatrix{
(X \subset X') \ar[rr]_{(f, f')} \ar[rd] & & (Y \subset Y') \ar[ld] \\
& (B \subset B')
}
$$
，它是 $S$ 上代数空间加厚的交换图。假设
$B \subset B'$ 是有限阶加厚，$X'$ 在 $B'$ 上平坦，
$X = B \times_{B'} X'$，且 $Y = B \times_{B'} Y'$。则
\begin{enumerate}
\item $f$ 可表当且仅当 $f'$ 可表，
\label{spaces-more-morphisms-item-representable}
\item $f$ 平坦当且仅当 $f'$ 平坦，
\label{spaces-more-morphisms-item-flat}
\item $f$ 是同构当且仅当 $f'$ 是同构，
\label{spaces-more-morphisms-item-isomorphism}
\item $f$ 是开浸入当且仅当 $f'$ 是开浸入，
\label{spaces-more-morphisms-item-open-immersion}
\item $f$ 拟紧当且仅当 $f'$ 拟紧，
\label{spaces-more-morphisms-item-quasi-compact}
\item $f$ 普遍闭当且仅当 $f'$ 普遍闭，
\label{spaces-more-morphisms-item-universally-closed}
\item $f$（拟）分离当且仅当 $f'$（拟）分离，
\label{spaces-more-morphisms-item-separated}
\item $f$ 是单态射当且仅当 $f'$ 是单态射，
\label{spaces-more-morphisms-item-monomorphism}
\item $f$ 满射当且仅当 $f'$ 满射，
\label{spaces-more-morphisms-item-surjective}
\item $f$ 普遍单射当且仅当 $f'$ 普遍单射，
\label{spaces-more-morphisms-item-universally-injective}
\item $f$ 仿射当且仅当 $f'$ 仿射，
\label{spaces-more-morphisms-item-affine}
\item
\label{spaces-more-morphisms-item-finite-type}
$f$ 局部有限型当且仅当 $f'$ 局部有限型，
\item $f$ 局部拟有限当且仅当 $f'$ 局部拟有限，
\label{spaces-more-morphisms-item-quasi-finite}
\item
\label{spaces-more-morphisms-item-finite-presentation}
$f$ 局部有限表示当且仅当 $f'$ 局部有限表示，
\item
\label{spaces-more-morphisms-item-relative-dimension-d}
$f$ 是相对维数 $d$ 的局部有限型态射当且仅当
$f'$ 是相对维数 $d$ 的局部有限型态射，
\item $f$ 普遍开当且仅当 $f'$ 普遍开，
\label{spaces-more-morphisms-item-universally-open}
\item $f$ 是 syntomic 态射当且仅当 $f'$ 是 syntomic 态射，
\label{spaces-more-morphisms-item-syntomic}
\item $f$ 光滑当且仅当 $f'$ 光滑，
\label{spaces-more-morphisms-item-smooth}
\item $f$ 非分歧当且仅当 $f'$ 非分歧，
\label{spaces-more-morphisms-item-unramified}
\item $f$ 平展当且仅当 $f'$ 平展，
\label{spaces-more-morphisms-item-etale}
\item $f$ 固有当且仅当 $f'$ 固有，
\label{spaces-more-morphisms-item-proper}
\item $f$ 整当且仅当 $f'$ 整，
\label{spaces-more-morphisms-item-integral}
\item $f$ 有限当且仅当 $f'$ 有限，
\label{spaces-more-morphisms-item-finite}
\item
\label{spaces-more-morphisms-item-finite-locally-free}
$f$ 有限局部自由（秩为 $d$）当且仅当 $f'$ 有限局部自由
（秩为 $d$）；
\item 在此添加更多情形。% P09-ERR-0572
\end{enumerate}
\end{lemma}

\begin{proof}
情形 (\ref{spaces-more-morphisms-item-representable}) 由引理
\ref{spaces-more-morphisms-lemma-thicken-property-morphisms} 得出。

\medskip\noindent
选取概形 $U'$ 和满平展态射 $U' \to B'$。
选取概形 $V'$ 和满平展态射
$V' \to U' \times_{B'} Y'$。
选取概形 $W'$ 和满平展态射
$W' \to V' \times_{Y'} X'$。令 $U, V, W$ 分别为
$U', V', W'$ 沿 $B \to B'$ 的基变换。考虑概形加厚的图
$$
\xymatrix{
(W \subset W') \ar[rr] \ar[rd] & & (V \subset V') \ar[ld] \\
& (U \subset U')
}
$$
。对于任何对源与靶平展局部的性质，把概形加厚态射的相应结论
应用于上图便立即得到所需结论。因此，情形
(\ref{spaces-more-morphisms-item-flat})、(\ref{spaces-more-morphisms-item-finite-type})、
(\ref{spaces-more-morphisms-item-quasi-finite})、(\ref{spaces-more-morphisms-item-finite-presentation})、
(\ref{spaces-more-morphisms-item-relative-dimension-d})、(\ref{spaces-more-morphisms-item-syntomic})、
(\ref{spaces-more-morphisms-item-smooth})、(\ref{spaces-more-morphisms-item-unramified})、(\ref{spaces-more-morphisms-item-etale})
由《态射进阶》引理
\ref{more-morphisms-lemma-deform-property} 的相应情形得出。

\medskip\noindent
由于 $X \to X'$ 和 $Y \to Y'$ 都是普遍同胚，关于映射
$X \to Y$ 和 $X' \to Y'$ 的任何拓扑问题都有相同答案。因此情形
(\ref{spaces-more-morphisms-item-quasi-compact})、(\ref{spaces-more-morphisms-item-universally-closed})、
(\ref{spaces-more-morphisms-item-surjective})、(\ref{spaces-more-morphisms-item-universally-injective}) 和
(\ref{spaces-more-morphisms-item-universally-open}) 成立。

\medskip\noindent
在其余每种情形中，我们只证明蕴含
$f\text{ 具有 }P \Rightarrow f'\text{ 具有 }P$，因为反向蕴含由
$P$ 在基变换下稳定得出；见《空间》引理
\ref{spaces-lemma-base-change-immersions} 以及《空间的态射》引理
\ref{spaces-morphisms-lemma-base-change-separated}、
\ref{spaces-morphisms-lemma-base-change-monomorphism}、
\ref{spaces-morphisms-lemma-base-change-affine}、
\ref{spaces-morphisms-lemma-base-change-proper}、
\ref{spaces-morphisms-lemma-base-change-integral} 和
\ref{spaces-morphisms-lemma-base-change-finite-locally-free}。

\medskip\noindent
情形 (\ref{spaces-more-morphisms-item-open-immersion})。假设 $f$ 是开浸入。
由 (\ref{spaces-more-morphisms-item-etale})，$f'$ 平展；又由
(\ref{spaces-more-morphisms-item-universally-injective}) 普遍单射，故 $f'$ 是开浸入；
见《空间的态射》引理
\ref{spaces-morphisms-lemma-etale-universally-injective-open}。
若先用 (\ref{spaces-more-morphisms-item-representable}) 约化到概形情形，则也可以不用这个引理。

\medskip\noindent
情形 (\ref{spaces-more-morphisms-item-isomorphism})。它由情形
(\ref{spaces-more-morphisms-item-open-immersion}) 和 (\ref{spaces-more-morphisms-item-surjective}) 得出。

\medskip\noindent
情形 (\ref{spaces-more-morphisms-item-separated})。见引理
\ref{spaces-more-morphisms-lemma-thicken-property-morphisms}。

\medskip\noindent
情形 (\ref{spaces-more-morphisms-item-monomorphism})。假设 $f$ 是单态射。
考虑对角态射 $\Delta_{X'/Y'} : X' \to X' \times_{Y'} X'$。
$\Delta_{X'/Y'}$ 沿 $B \to B'$ 的基变换是 $\Delta_{X/Y}$，
后者依假设是同构。由 (\ref{spaces-more-morphisms-item-isomorphism}) 可知
$\Delta_{X'/Y'}$ 是同构，因而 $f'$ 是单态射。

\medskip\noindent
情形 (\ref{spaces-more-morphisms-item-affine})。见引理 \ref{spaces-more-morphisms-lemma-thicken-property-morphisms}。

\medskip\noindent
情形 (\ref{spaces-more-morphisms-item-proper})。见引理
\ref{spaces-more-morphisms-lemma-thicken-property-morphisms-cartesian}。

\medskip\noindent
情形 (\ref{spaces-more-morphisms-item-integral})。见引理
\ref{spaces-more-morphisms-lemma-thicken-property-morphisms}。

\medskip\noindent
情形 (\ref{spaces-more-morphisms-item-finite})。见引理
\ref{spaces-more-morphisms-lemma-thicken-property-morphisms-cartesian}。

\medskip\noindent
情形 (\ref{spaces-more-morphisms-item-finite-locally-free})。假设 $f$ 有限局部自由。
由 (\ref{spaces-more-morphisms-item-finite}) 可知 $f'$ 有限。
由 (\ref{spaces-more-morphisms-item-flat}) 可知 $f'$ 平坦。
由 (\ref{spaces-more-morphisms-item-finite-presentation})，$f'$ 局部有限表示。
因此，由《空间的态射》引理
\ref{spaces-morphisms-lemma-finite-flat}，$f'$ 有限局部自由。
\end{proof}

\noindent
下面这个引理是“按纤维判定平坦性的准则”的“局部幂零”版本；见
第 \ref{spaces-more-morphisms-section-criterion-flat-fibres} 节。

\begin{lemma}
\label{spaces-more-morphisms-lemma-flatness-morphism-thickenings-fp-over-ft}
设 $S$ 为概形。考虑交换图
$$
\xymatrix{
(X \subset X') \ar[rr]_{(f, f')} \ar[rd] & & (Y \subset Y') \ar[ld] \\
& (B \subset B')
}
$$
，它是 $S$ 上代数空间加厚的交换图。假设
\begin{enumerate}
\item $Y' \to B'$ 局部有限型；
\item $X' \to B'$ 平坦且局部有限表示；
\item $f$ 平坦；
\item $X = B \times_{B'} X'$ 且 $Y = B \times_{B'} Y'$。
\end{enumerate}
则 $f'$ 平坦，并且对于所有 $y' \in |Y'|$，只要它在
$|f'|$ 的像中，态射 $Y' \to B'$ 就在 $y'$ 处平坦。
\end{lemma}

\begin{proof}
选取概形 $U'$ 和满平展态射 $U' \to B'$。
选取概形 $V'$ 和满平展态射
$V' \to U' \times_{B'} Y'$。
选取概形 $W'$ 和满平展态射
$W' \to V' \times_{Y'} X'$。令 $U, V, W$ 分别为
$U', V', W'$ 沿 $B \to B'$ 的基变换。于是 $f'$ 的平坦性等价于
$W' \to V'$ 的平坦性，而已知 $W \to V$ 平坦。因此可以把概形情形的
引理应用于概形加厚的图
$$
\xymatrix{
(W \subset W') \ar[rr] \ar[rd] & & (V \subset V') \ar[ld] \\
& (U \subset U')
}
$$
；见《态射进阶》引理
\ref{more-morphisms-lemma-flatness-morphism-thickenings-fp-over-ft}。
关于 $Y'/B'$ 在 $f'$ 像中各点处平坦的断言也以同样方式得出。
\end{proof}

\noindent
概形态射的许多性质在上述引理中的平坦形变下保持不变。

\begin{lemma}
\label{spaces-more-morphisms-lemma-deform-property-fp-over-ft}
设 $S$ 为概形。考虑交换图
$$
\xymatrix{
(X \subset X') \ar[rr]_{(f, f')} \ar[rd] & & (Y \subset Y') \ar[ld] \\
& (B \subset B')
}
$$
，它是 $S$ 上代数空间加厚的交换图。假设
$Y' \to B'$ 局部有限型，
$X' \to B'$ 平坦且局部有限表示，
$X = B \times_{B'} X'$，且 $Y = B \times_{B'} Y'$。则
\begin{enumerate}
\item $f$ 可表当且仅当 $f'$ 可表，
\label{spaces-more-morphisms-item-representable-fp-over-ft}
\item $f$ 平坦当且仅当 $f'$ 平坦，
\label{spaces-more-morphisms-item-flat-fp-over-ft}
\item $f$ 是同构当且仅当 $f'$ 是同构，
\label{spaces-more-morphisms-item-isomorphism-fp-over-ft}
\item $f$ 是开浸入当且仅当 $f'$ 是开浸入，
\label{spaces-more-morphisms-item-open-immersion-fp-over-ft}
\item $f$ 拟紧当且仅当 $f'$ 拟紧，
\label{spaces-more-morphisms-item-quasi-compact-fp-over-ft}
\item $f$ 普遍闭当且仅当 $f'$ 普遍闭，
\label{spaces-more-morphisms-item-universally-closed-fp-over-ft}
\item $f$（拟）分离当且仅当 $f'$（拟）分离，
\label{spaces-more-morphisms-item-separated-fp-over-ft}
\item $f$ 是单态射当且仅当 $f'$ 是单态射，
\label{spaces-more-morphisms-item-monomorphism-fp-over-ft}
\item $f$ 满射当且仅当 $f'$ 满射，
\label{spaces-more-morphisms-item-surjective-fp-over-ft}
\item $f$ 普遍单射当且仅当 $f'$ 普遍单射，
\label{spaces-more-morphisms-item-universally-injective-fp-over-ft}
\item $f$ 仿射当且仅当 $f'$ 仿射，
\label{spaces-more-morphisms-item-affine-fp-over-ft}
\item $f$ 局部拟有限当且仅当 $f'$ 局部拟有限，
\label{spaces-more-morphisms-item-quasi-finite-fp-over-ft}
\item
\label{spaces-more-morphisms-item-relative-dimension-d-fp-over-ft}
$f$ 是相对维数 $d$ 的局部有限型态射当且仅当
$f'$ 是相对维数 $d$ 的局部有限型态射，
\item $f$ 普遍开当且仅当 $f'$ 普遍开，
\label{spaces-more-morphisms-item-universally-open-fp-over-ft}
\item $f$ 是 syntomic 态射当且仅当 $f'$ 是 syntomic 态射，
\label{spaces-more-morphisms-item-syntomic-fp-over-ft}
\item $f$ 光滑当且仅当 $f'$ 光滑，
\label{spaces-more-morphisms-item-smooth-fp-over-ft}
\item $f$ 非分歧当且仅当 $f'$ 非分歧，
\label{spaces-more-morphisms-item-unramified-fp-over-ft}
\item $f$ 平展当且仅当 $f'$ 平展，
\label{spaces-more-morphisms-item-etale-fp-over-ft}
\item $f$ 固有当且仅当 $f'$ 固有，
\label{spaces-more-morphisms-item-proper-fp-over-ft}
\item $f$ 有限当且仅当 $f'$ 有限，
\label{spaces-more-morphisms-item-finite-fp-over-ft}
\item
\label{spaces-more-morphisms-item-finite-locally-free-fp-over-ft}
$f$ 有限局部自由（秩为 $d$）当且仅当 $f'$ 有限局部自由
（秩为 $d$）；
\item 在此添加更多情形。% P09-ERR-0573
\end{enumerate}
\end{lemma}

\begin{proof}
情形 (\ref{spaces-more-morphisms-item-representable-fp-over-ft}) 由引理
\ref{spaces-more-morphisms-lemma-thicken-property-morphisms} 得出。

\medskip\noindent
选取概形 $U'$ 和满平展态射 $U' \to B'$。
选取概形 $V'$ 和满平展态射
$V' \to U' \times_{B'} Y'$。
选取概形 $W'$ 和满平展态射
$W' \to V' \times_{Y'} X'$。令 $U, V, W$ 分别为
$U', V', W'$ 沿 $B \to B'$ 的基变换。考虑概形加厚的图
$$
\xymatrix{
(W \subset W') \ar[rr] \ar[rd] & & (V \subset V') \ar[ld] \\
& (U \subset U')
}
$$
。对于任何对源与靶平展局部的性质，把概形加厚态射的相应结论
应用于上图便立即得到所需结论。因此，情形
(\ref{spaces-more-morphisms-item-flat-fp-over-ft})、
(\ref{spaces-more-morphisms-item-quasi-finite-fp-over-ft})、
(\ref{spaces-more-morphisms-item-relative-dimension-d-fp-over-ft})、
(\ref{spaces-more-morphisms-item-syntomic-fp-over-ft})、
(\ref{spaces-more-morphisms-item-smooth-fp-over-ft})、
(\ref{spaces-more-morphisms-item-unramified-fp-over-ft})、
(\ref{spaces-more-morphisms-item-etale-fp-over-ft})
由《态射进阶》引理
\ref{more-morphisms-lemma-deform-property-fp-over-ft} 的相应情形得出。

\medskip\noindent
由于 $X \to X'$ 和 $Y \to Y'$ 都是普遍同胚，关于映射
$X \to Y$ 和 $X' \to Y'$ 的任何拓扑问题都有相同答案。因此情形
(\ref{spaces-more-morphisms-item-quasi-compact-fp-over-ft})、
(\ref{spaces-more-morphisms-item-universally-closed-fp-over-ft})、
(\ref{spaces-more-morphisms-item-surjective-fp-over-ft})、
(\ref{spaces-more-morphisms-item-universally-injective-fp-over-ft}) 和
(\ref{spaces-more-morphisms-item-universally-open-fp-over-ft}) 成立。

\medskip\noindent
在其余每种情形中，我们只证明蕴含
$f\text{ 具有 }P \Rightarrow f'\text{ 具有 }P$，因为反向蕴含由
$P$ 在基变换下稳定得出；见《空间》引理
\ref{spaces-lemma-base-change-immersions} 以及《空间的态射》引理
\ref{spaces-morphisms-lemma-base-change-separated}、
\ref{spaces-morphisms-lemma-base-change-monomorphism}、
\ref{spaces-morphisms-lemma-base-change-affine}、
\ref{spaces-morphisms-lemma-base-change-proper}、
\ref{spaces-morphisms-lemma-base-change-integral} 和
\ref{spaces-morphisms-lemma-base-change-finite-locally-free}。

\medskip\noindent
情形 (\ref{spaces-more-morphisms-item-open-immersion-fp-over-ft})。
假设 $f$ 是开浸入。
由 (\ref{spaces-more-morphisms-item-etale-fp-over-ft})，$f'$ 平展；又由
(\ref{spaces-more-morphisms-item-universally-injective-fp-over-ft}) 普遍单射，故
$f'$ 是开浸入；见《空间的态射》引理
\ref{spaces-morphisms-lemma-etale-universally-injective-open}。
若先用 (\ref{spaces-more-morphisms-item-representable-fp-over-ft}) 约化到概形情形，
则也可以不用这个引理。

\medskip\noindent
情形 (\ref{spaces-more-morphisms-item-isomorphism-fp-over-ft})。它由情形
(\ref{spaces-more-morphisms-item-open-immersion-fp-over-ft}) 和
(\ref{spaces-more-morphisms-item-surjective-fp-over-ft}) 得出。

\medskip\noindent
情形 (\ref{spaces-more-morphisms-item-separated-fp-over-ft})。见引理
\ref{spaces-more-morphisms-lemma-thicken-property-morphisms}。

\medskip\noindent
情形 (\ref{spaces-more-morphisms-item-monomorphism-fp-over-ft})。假设 $f$ 是单态射。
考虑对角态射 $\Delta_{X'/Y'} : X' \to X' \times_{Y'} X'$。
$\Delta_{X'/Y'}$ 沿 $B \to B'$ 的基变换是 $\Delta_{X/Y}$，
后者依假设是同构。由 (\ref{spaces-more-morphisms-item-isomorphism-fp-over-ft}) 可知
$\Delta_{X'/Y'}$ 是同构，因而 $f'$ 是单态射。

\medskip\noindent
情形 (\ref{spaces-more-morphisms-item-affine-fp-over-ft})。
见引理 \ref{spaces-more-morphisms-lemma-thicken-property-morphisms}。

\medskip\noindent
情形 (\ref{spaces-more-morphisms-item-proper-fp-over-ft})。见引理
\ref{spaces-more-morphisms-lemma-properties-that-extend-over-thickenings}。

\medskip\noindent
情形 (\ref{spaces-more-morphisms-item-finite-fp-over-ft})。见引理
\ref{spaces-more-morphisms-lemma-properties-that-extend-over-thickenings}。

\medskip\noindent
情形 (\ref{spaces-more-morphisms-item-finite-locally-free-fp-over-ft})。
假设 $f$ 有限局部自由。
由 (\ref{spaces-more-morphisms-item-finite-fp-over-ft}) 可知 $f'$ 有限。
由 (\ref{spaces-more-morphisms-item-flat-fp-over-ft}) 可知 $f'$ 平坦。
此外，由《空间的态射》引理% P09-ERR-0574
\ref{spaces-morphisms-lemma-finite-presentation-permanence}，
$f'$ 局部有限表示。
因此，由《空间的态射》引理
\ref{spaces-morphisms-lemma-finite-flat}，$f'$ 有限局部自由。
\end{proof}













\section{形式光滑态射}
\label{spaces-more-morphisms-section-formally-smooth}

\noindent
本节引入代数空间形式光滑态射 $X \to Y$ 的概念。这类态射的特征是：
只要 $T$ 仿射，$X$ 的 $T$-值点就能提升到 $T$ 的无穷小加厚。
主要结论是，形式光滑且局部有限表示的态射是光滑的；见引理
\ref{spaces-more-morphisms-lemma-smooth-formally-smooth}。事实证明，这个判据往往比雅可比判据
更便于使用。

\begin{definition}
\label{spaces-more-morphisms-definition-formally-smooth}
设 $S$ 为概形。若代数空间的态射 $f : X \to Y$（它在 $S$ 上）作为函子变换
在定义 \ref{spaces-more-morphisms-definition-formally-smooth-etale-unramified} 的意义下形式光滑，
则称它为{\it 形式光滑的}。
\end{definition}

\noindent
对于形式非分歧和形式平展态射，可以去掉 $T'$ 仿射这一条件；见引理
\ref{spaces-more-morphisms-lemma-formally-unramified-not-affine} 和
\ref{spaces-more-morphisms-lemma-formally-etale-not-affine}。对于形式光滑态射，这不再成立。
事实上，稍微自然一些的条件是要求能够在 $T'$ 上平展局部地补出虚线箭头。
进一步分析引理 \ref{spaces-more-morphisms-lemma-smooth-formally-smooth} 的证明可知，这与目前的
定义等价。要求虚线箭头在 $T'$ 上 fppf 局部地存在也足够，不过证明起来
稍微困难一些。

\medskip\noindent
第 \ref{spaces-more-morphisms-section-formally-smooth-etale-unramified} 节已经在形式光滑函子变换的
更一般背景下证明了若干结果，这里不再重述。

\begin{lemma}
\label{spaces-more-morphisms-lemma-composition-formally-smooth}
形式光滑态射的复合仍形式光滑。
\end{lemma}

\begin{proof}
略。
\end{proof}

\begin{lemma}
\label{spaces-more-morphisms-lemma-base-change-formally-smooth}
形式光滑态射的基变换仍形式光滑。
\end{lemma}

\begin{proof}
略；代数版本见《代数》引理
\ref{algebra-lemma-base-change-fs}。
\end{proof}

\begin{lemma}
\label{spaces-more-morphisms-lemma-formally-etale-unramified-smooth}
设 $f : X \to S$ 为概形态射。则 $f$ 形式平展当且仅当
$f$ 形式光滑且形式非分歧。
\end{lemma}

\begin{proof}
略。
\end{proof}

\noindent
下面是一个辅助引理，稍后将由引理 \ref{spaces-more-morphisms-lemma-formally-smooth} 取代。

\begin{lemma}
\label{spaces-more-morphisms-lemma-helper-formally-smooth}
设 $S$ 为概形，并设
$$
\xymatrix{
U \ar[d] \ar[r]_\psi & V \ar[d] \\
X \ar[r]^f & Y
}
$$
为 $S$ 上代数空间态射的交换图。若竖直箭头平展且 $f$ 形式光滑，
则 $\psi$ 形式光滑。
\end{lemma}

\begin{proof}
由引理 \ref{spaces-more-morphisms-lemma-representable-property-formally-property}，态射
$U \to X$ 和 $V \to Y$ 形式平展。由引理
\ref{spaces-more-morphisms-lemma-composition-formally-smooth-etale-unramified}，复合
$U \to Y$ 形式光滑。再由引理 \ref{spaces-more-morphisms-lemma-formally-permanence} 可知
$\psi : U \to V$ 形式光滑。
\end{proof}

\noindent
下面的引理是本节的主要结果。结合《空间的极限》命题
\ref{spaces-limits-proposition-characterize-locally-finite-presentation}，
它说明可以判定代数空间态射 $f : X \to Y$ 是否光滑，而只需考察函子变换
$X \to Y$ 的“简单”性质。

\begin{lemma}[无穷小提升判据]
\label{spaces-more-morphisms-lemma-smooth-formally-smooth}
设 $S$ 为概形。设 $f : X \to Y$ 为 $S$ 上代数空间的态射。
下列条件等价：
\begin{enumerate}
\item 态射 $f$ 光滑。
\item 态射 $f$ 局部有限表示且形式光滑。
\end{enumerate}
\end{lemma}

\begin{proof}
假设 $f : X \to S$ 局部有限表示且形式光滑。% P09-ERR-0575
考虑交换图
$$
\xymatrix{
U \ar[d] \ar[r]_\psi & V \ar[d] \\
X \ar[r]^f & Y
}
$$
，其中 $U$ 和 $V$ 是概形，竖直箭头平展且满。由引理
\ref{spaces-more-morphisms-lemma-helper-formally-smooth} 可知 $\psi : U \to V$ 形式光滑。
由《空间的态射》引理
\ref{spaces-morphisms-lemma-finite-presentation-local}，态射 $\psi$ 局部有限表示。
因此，由概形情形可知态射 $\psi$ 光滑；见《态射进阶》引理
\ref{more-morphisms-lemma-smooth-formally-smooth}。于是 $f$ 光滑；见
《空间的态射》引理
\ref{spaces-morphisms-lemma-smooth-local}。

\medskip\noindent
反过来，假设 $f : X \to Y$ 光滑。考虑定义
\ref{spaces-more-morphisms-definition-formally-smooth} 中实线部分交换的图
$$
\xymatrix{
X \ar[d]_f & T \ar[d]^i \ar[l]^a \\
Y & T' \ar[l] \ar@{-->}[lu]
}
$$
。我们将证明虚线箭头存在，从而证明 $f$ 形式光滑。
令 $\mathcal{F}$ 为引理 \ref{spaces-more-morphisms-lemma-sheaf} 在评注
\ref{spaces-more-morphisms-remark-special-case} 所述特殊情形下定义于
$(T')_{spaces, \etale}$ 上的集合层。令
$$
\mathcal{H} =
\SheafHom_{\mathcal{O}_T}(a^*\Omega_{X/Y}, \mathcal{C}_{T/T'})
$$
为 $\mathcal{O}_T$-模层（定义于 $T_{spaces, \etale}$ 上），并按引理
\ref{spaces-more-morphisms-lemma-action-sheaf} 取作用
$\mathcal{H} \times \mathcal{F} \to \mathcal{F}$。
这个作用 $\mathcal{H} \times \mathcal{F} \to \mathcal{F}$ 使
$\mathcal{F}$ 成为伪 $\mathcal{H}$-挠子；见《位点上的上同调》定义
\ref{sites-cohomology-definition-torsor}。目标是证明 $\mathcal{F}$ 是平凡的
$\mathcal{H}$-挠子。分两步：(I) 为证明 $\mathcal{F}$ 是挠子，必须证明
$\mathcal{F}$ 在平展局部有截面；(II) 为证明 $\mathcal{F}$ 是平凡挠子，
只需证明 $H^1(T_\etale, \mathcal{H}) = 0$；见《位点上的上同调》引理
\ref{sites-cohomology-lemma-torsors-h1}。

\medskip\noindent
先证明 (I)。为此，选取交换图
$$
\xymatrix{
U \ar[d] \ar[r]_\psi & V \ar[d] \\
X \ar[r]^f & Y
}
$$
，其中 $U$ 和 $V$ 是概形，竖直箭头平展且满。由于假设 $f$ 光滑，
可知 $\psi$ 光滑，因而由引理
\ref{spaces-more-morphisms-lemma-representable-property-formally-property} 可知它形式光滑。
由同一引理，态射 $V \to Y$ 形式平展。因此，由引理
\ref{spaces-more-morphisms-lemma-composition-formally-smooth-etale-unramified}，复合
$U \to Y$ 形式光滑。于是 (I) 由引理 \ref{spaces-more-morphisms-lemma-etale-on-top} 的第 (4) 部分得出。

\medskip\noindent
最后证明 (II)。由引理
\ref{spaces-more-morphisms-lemma-finite-presentation-differentials} 可知
$\Omega_{X/S}$ 有限表示。% P09-ERR-0576
因此 $a^*\Omega_{X/S}$ 有限表示（见《空间的性质》第
\ref{spaces-properties-section-properties-modules} 节）。故层
$\mathcal{H} =
\SheafHom_{\mathcal{O}_T}(a^*\Omega_{X/Y}, \mathcal{C}_{T/T'})$
拟凝聚；见《空间的性质》引理
\ref{spaces-properties-lemma-properties-quasi-coherent}。因而由《下降》命题
\ref{descent-proposition-same-cohomology-quasi-coherent} 和《概形的上同调》引理
\ref{coherent-lemma-quasi-coherent-affine-cohomology-zero}，有
$$
H^1(T_{spaces, \etale}, \mathcal{H}) =
H^1(T_\etale, \mathcal{H}) =
H^1(T, \mathcal{H}) = 0
$$
，正合所需。
\end{proof}

\noindent
光滑态射满足强局部提升性质；见引理 \ref{spaces-more-morphisms-lemma-smooth-strong-lift}。
若在该引理中假设 $T'$ 仿射，我们不知道是否有必要取平展覆盖。
更确切地说，若有代数空间的交换图
$$
\xymatrix{
X \ar[d] & T \ar[l] \ar[d] \\
Y & T' \ar[l] \ar@{..>}[lu]
}
$$
，其中 $X \to Y$ 光滑且 $T \to T'$ 是仿射概形的加厚，那么使图交换的
虚线箭头是否总存在？% P09-ERR-0577
若你知道答案或掌握参考文献，请发送邮件至
\href{mailto:stacks.project@gmail.com}{stacks.project@gmail.com}。

\begin{lemma}
\label{spaces-more-morphisms-lemma-smooth-strong-lift}
设 $S$ 为概形。考虑交换图
$$
\xymatrix{
X \ar[d] & T \ar[l] \ar[d] \\
Y & T' \ar[l]
}
$$
，它是 $S$ 上代数空间的交换图，其中 $X \to Y$ 光滑且
$T \to T'$ 是加厚。则存在平展覆盖
$\{T'_i \to T'\}$，使得能在图
$$
\xymatrix{
X \ar[d] & T \ar[l] \ar[d] & T \times_{T'} T'_i \ar[l] \ar[d] \\
Y & T' \ar[l] & T'_i \ar[l] \ar@{..>}[llu]
}
$$
中找到使图交换的虚线箭头（对所有 $i$）。
\end{lemma}

\begin{proof}
选取平展覆盖 $\{Y_i \to Y\}$，其中每个 $Y_i$ 都仿射。
用诱导的平展覆盖替换 $T'$ 后，可假设 $Y$ 仿射。

\medskip\noindent
假设 $Y$ 仿射。选取平展覆盖 $\{X_i \to X\}$。
这给出 $T$ 的平展覆盖。由定理 \ref{spaces-more-morphisms-theorem-topological-invariance}
（另见第 \ref{spaces-more-morphisms-section-thickenings} 节的讨论），$T$ 的这个平展覆盖来自
$T'$ 的某个平展覆盖。因此可假设 $X$ 仿射。

\medskip\noindent
假设 $X$ 和 $Y$ 都仿射。可再对 $T'$ 取一次平展覆盖并假设 $T'$ 仿射。
在这种情形下，引理由《代数》引理
\ref{algebra-lemma-smooth-strong-lift} 得出。
\end{proof}







\noindent
下面再做一些工作，证明形式光滑性对源平展局部。首先证明形式光滑态射
具有良好的微分层。局部投射拟凝聚模的概念定义于《空间的性质》第
\ref{spaces-properties-section-locally-projective} 节。

\begin{lemma}
\label{spaces-more-morphisms-lemma-formally-smooth-sheaf-differentials}
设 $S$ 为概形。设 $f : X \to Y$ 为 $S$ 上代数空间的形式光滑态射。
则 $\Omega_{X/Y}$ 在 $X$ 上局部投射。
\end{lemma}

\begin{proof}
选取图
$$
\xymatrix{
U \ar[d] \ar[r]_\psi & V \ar[d] \\
X \ar[r]^f & Y
}
$$
，其中 $U$ 和 $V$ 是仿射（！）概形，竖直箭头平展。由引理
\ref{spaces-more-morphisms-lemma-helper-formally-smooth} 可知 $\psi : U \to V$ 形式光滑。
因此
$\Gamma(V, \mathcal{O}_V) \to \Gamma(U, \mathcal{O}_U)$
是形式光滑环同态；见《态射进阶》引理
\ref{more-morphisms-lemma-affine-formally-smooth}。于是由《代数》引理
\ref{algebra-lemma-characterize-formally-smooth-again}，
$\Gamma(U, \mathcal{O}_U)$-模
$\Omega_{\Gamma(U, \mathcal{O}_U)/\Gamma(V, \mathcal{O}_V)}$
是投射模。因此 $\Omega_{U/V}$ 局部投射；见《性质》第
\ref{properties-section-locally-projective} 节。由于
$\Omega_{X/Y}|_U = \Omega_{U/V}$，可知 $\Omega_{X/Y}$ 也局部投射。
（因为可以找到 $X$ 的平展覆盖，其成员是上述图中的仿射 $U$；细节从略。）
\end{proof}

\begin{lemma}
\label{spaces-more-morphisms-lemma-h1-is-zero}
设 $T$ 为仿射概形。设 $\mathcal{F}$、$\mathcal{G}$ 是拟凝聚
$\mathcal{O}_T$-模（定义于 $T_\etale$ 上）。考虑内 Hom 层
$\mathcal{H} = \SheafHom_{\mathcal{O}_T}(\mathcal{F}, \mathcal{G})$，
它定义于 $T_\etale$ 上。
若 $\mathcal{F}$ 局部投射，则 $H^1(T_\etale, \mathcal{H}) = 0$。
\end{lemma}

\begin{proof}
由代数空间上局部投射层的定义（见《空间的性质》定义
\ref{spaces-properties-definition-locally-projective}），可知
$\mathcal{F}_{Zar} = \mathcal{F}|_{T_{Zar}}$ 是概形 $T$ 上的局部投射层。
因此 $\mathcal{F}_{Zar}$ 是某个自由 $\mathcal{O}_{T_{Zar}}$-模的直和项。
进而（因为 $\mathcal{F} = (\mathcal{F}_{Zar})^a$，见《下降》命题
\ref{descent-proposition-equivalence-quasi-coherent}），可知 $\mathcal{F}$ 是
某个自由 $\mathcal{O}_T$-模的直和项（定义于 $T_\etale$ 上）。故可假设
$\mathcal{F} = \bigoplus_{i \in I} \mathcal{O}_T$ 是自由模。
在这种情形下，$\mathcal{H} = \prod_{i \in I} \mathcal{G}$ 是拟凝聚模的积。
由《位点上的上同调》引理
\ref{sites-cohomology-lemma-cohomology-products} 可得 $H^1 = 0$，因为
仿射概形上拟凝聚层的上同调为零；见《下降》命题
\ref{descent-proposition-same-cohomology-quasi-coherent} 和《概形的上同调》引理
\ref{coherent-lemma-quasi-coherent-affine-cohomology-zero}。
\end{proof}

\begin{lemma}
\label{spaces-more-morphisms-lemma-formally-smooth}
设 $S$ 为概形。设 $f : X \to Y$ 为 $S$ 上代数空间的态射。
下列条件等价：
\begin{enumerate}
\item $f$ 形式光滑；
\item 对于每个图
$$
\xymatrix{
U \ar[d] \ar[r]_\psi & V \ar[d] \\
X \ar[r]^f & Y
}
$$
，其中 $U$ 和 $V$ 是概形且竖直箭头平展，概形态射 $\psi$ 都形式光滑
（即《态射进阶》定义
\ref{more-morphisms-definition-formally-unramified} 的意义）；% P09-ERR-0578
\item 对于一个竖直箭头为满射的这类图，态射 $\psi$ 形式光滑。
\end{enumerate}
\end{lemma}

\begin{proof}
引理 \ref{spaces-more-morphisms-lemma-helper-formally-smooth} 已经说明 (1) 蕴含 (2) 和 (3)。
假设 (3)。证明 $f$ 形式光滑的过程与引理
\ref{spaces-more-morphisms-lemma-smooth-formally-smooth} 中 (1) $\Rightarrow$ (2) 的证明完全类似。

\medskip\noindent
考虑定义 \ref{spaces-more-morphisms-definition-formally-smooth} 中实线部分交换的图
$$
\xymatrix{
X \ar[d]_f & T \ar[d]^i \ar[l]^a \\
Y & T' \ar[l] \ar@{-->}[lu]
}
$$
。我们将证明虚线箭头存在，从而证明 $f$ 形式光滑。
令 $\mathcal{F}$ 为引理 \ref{spaces-more-morphisms-lemma-sheaf} 在评注
\ref{spaces-more-morphisms-remark-special-case} 所述特殊情形下定义于
$(T')_{spaces, \etale}$ 上的集合层。令
$$
\mathcal{H} =
\SheafHom_{\mathcal{O}_T}(a^*\Omega_{X/Y}, \mathcal{C}_{T/T'})
$$
为 $\mathcal{O}_T$-模层（定义于 $T_{spaces, \etale}$ 上），并按引理
\ref{spaces-more-morphisms-lemma-action-sheaf} 取作用
$\mathcal{H} \times \mathcal{F} \to \mathcal{F}$。
这个作用 $\mathcal{H} \times \mathcal{F} \to \mathcal{F}$ 使
$\mathcal{F}$ 成为伪 $\mathcal{H}$-挠子；见《位点上的上同调》定义
\ref{sites-cohomology-definition-torsor}。目标是证明 $\mathcal{F}$ 是平凡的
$\mathcal{H}$-挠子。分两步：(I) 为证明 $\mathcal{F}$ 是挠子，必须证明
$\mathcal{F}$ 在平展局部有截面；(II) 为证明 $\mathcal{F}$ 是平凡挠子，
只需证明 $H^1(T_\etale, \mathcal{H}) = 0$；见《位点上的上同调》引理
\ref{sites-cohomology-lemma-torsors-h1}。

\medskip\noindent
先证明 (I)。为此考虑图（它因假设 (3) 而存在）
$$
\xymatrix{
U \ar[d] \ar[r]_\psi & V \ar[d] \\
X \ar[r]^f & Y
}
$$
，其中 $U$ 和 $V$ 是概形，竖直箭头平展且满，并且 $\psi$ 形式光滑。
由引理 \ref{spaces-more-morphisms-lemma-representable-property-formally-property}，态射
$V \to Y$ 形式平展。因此，由引理
\ref{spaces-more-morphisms-lemma-composition-formally-smooth-etale-unramified}，复合
$U \to Y$ 形式光滑。于是 (I) 由引理 \ref{spaces-more-morphisms-lemma-etale-on-top} 的第 (4) 部分得出。

\medskip\noindent
最后证明 (II)。由引理
\ref{spaces-more-morphisms-lemma-formally-smooth-sheaf-differentials} 可知
$\Omega_{U/V}$ 局部投射。% P09-ERR-0579
因此 $\Omega_{X/Y}$ 局部投射；见《空间上的下降》引理
\ref{spaces-descent-lemma-locally-projective-descends}。
故 $a^*\Omega_{X/Y}$ 局部投射；见《空间的性质》引理
\ref{spaces-properties-lemma-locally-projective-pullback}。因此
$$
H^1(T_\etale, \mathcal{H}) =
H^1(T_\etale,
\SheafHom_{\mathcal{O}_T}(a^*\Omega_{X/Y}, \mathcal{C}_{T/T'}) = 0
$$
% P09-ERR-0580
由引理 \ref{spaces-more-morphisms-lemma-h1-is-zero} 得到，正合所需。
\end{proof}

\begin{lemma}
\label{spaces-more-morphisms-lemma-descending-property-formally-smooth}
性质 $\mathcal{P}(f) =$“$f$ 形式光滑”在基底上是 fpqc 局部的。
\end{lemma}

\begin{proof}
设 $f : X \to Y$ 是某概形 $S$ 上代数空间的态射。
选取指标集 $I$ 和图
$$
\xymatrix{
U_i \ar[d] \ar[r]_{\psi_i} & V_i \ar[d] \\
X \ar[r]^f & Y
}
$$
，其中竖直箭头平展，且 $U_i$、$V_i$ 是仿射概形。此外，假设
$\coprod U_i \to X$ 和 $\coprod V_i \to Y$ 满射；见《空间的性质》引理
\ref{spaces-properties-lemma-cover-by-union-affines}。由引理
\ref{spaces-more-morphisms-lemma-formally-smooth} 可知，$f$ 形式光滑当且仅当每个态射
$\psi_i$ 都形式光滑。% P09-ERR-0581
因此约化到仿射概形的态射。在这种情形下，结论由《代数》引理
\ref{algebra-lemma-descent-formally-smooth} 得出。部分细节从略。
\end{proof}

\begin{lemma}
\label{spaces-more-morphisms-lemma-triangle-differentials-formally-smooth}
设 $S$ 为概形。设 $f : X \to Y$、$g : Y \to Z$ 是 $S$ 上代数空间的态射。
假设 $f$ 形式光滑。则
$$
0 \to f^*\Omega_{Y/Z} \to \Omega_{X/Z} \to \Omega_{X/Y} \to 0
$$
在引理 \ref{spaces-more-morphisms-lemma-triangle-differentials} 中是短正合的。% P09-ERR-0582
\end{lemma}

\begin{proof}
由概形情形和平展局部化得出；前者见《态射进阶》引理
\ref{more-morphisms-lemma-triangle-differentials-formally-smooth}，
后者见引理 \ref{spaces-more-morphisms-lemma-formally-smooth} 和
\ref{spaces-more-morphisms-lemma-localize-differentials}。
\end{proof}

\begin{lemma}
\label{spaces-more-morphisms-lemma-differentials-formally-unramified-formally-smooth}
设 $S$ 为概形，$B$ 为 $S$ 上的代数空间。设 $h : Z \to X$ 是
$B$ 上代数空间的形式非分歧态射。假设 $Z$ 在 $B$ 上形式光滑。
则引理 \ref{spaces-more-morphisms-lemma-universally-unramified-differentials-sequence} 中的
典范正合序列
$$
0 \to \mathcal{C}_{Z/X} \to h^*\Omega_{X/B} \to \Omega_{Z/B} \to 0
$$
是短正合的。
\end{lemma}

\begin{proof}
设 $Z \to Z'$ 是 $Z$ 在 $X$ 上的泛一阶加厚。由引理
\ref{spaces-more-morphisms-lemma-universally-unramified-differentials-sequence} 的证明可知，
我们的序列与下列序列相识别：
$$
\mathcal{C}_{Z/Z'} \to \Omega_{Z'/B} \otimes \mathcal{O}_Z \to
\Omega_{Z/B} \to 0.
$$
由于 $Z \to S$ 形式光滑，% P09-ERR-0583
可以在 $Z'$ 上平展局部地找到左逆 $Z' \to Z$；它在 $B$ 上左逆于
含入映射 $Z \to Z'$。因此该序列平展局部分裂；见引理
\ref{spaces-more-morphisms-lemma-differentials-relative-immersion-section}。
\end{proof}

\begin{lemma}
\label{spaces-more-morphisms-lemma-two-unramified-morphisms-formally-smooth}
设 $S$ 为概形。设
$$
\xymatrix{
Z \ar[r]_i \ar[rd]_j & X \ar[d]^f \\
& Y
}
$$
为 $S$ 上代数空间的交换图，其中 $i$ 和 $j$ 形式非分歧，且 $f$ 形式光滑。
则引理 \ref{spaces-more-morphisms-lemma-two-unramified-morphisms} 中的典范正合序列
$$
0 \to
\mathcal{C}_{Z/Y} \to
\mathcal{C}_{Z/X} \to
i^*\Omega_{X/Y} \to 0
$$
正合且局部分裂。
\end{lemma}

\begin{proof}
记 $Z \to Z'$ 为 $Z$ 在 $X$ 上的泛一阶加厚。
记 $Z \to Z''$ 为 $Z$ 在 $Y$ 上的泛一阶加厚。由引理
\ref{spaces-more-morphisms-lemma-universally-unramified-differentials-sequence}，% P09-ERR-0584
这里有典范态射 $Z' \to Z''$，% P09-ERR-0585
从而得到交换图
$$
\xymatrix{
Z \ar[r]_{i'} \ar[rd]_{j'} & Z' \ar[r]_a \ar[d]^k & X \ar[d]^f \\
& Z'' \ar[r]^b & Y
}
$$
。上述序列借助形式非分歧态射余法层的定义与序列
$$
\mathcal{C}_{Z/Z''} \to
\mathcal{C}_{Z/Z'} \to
(i')^*\Omega_{Z'/Z''} \to 0
$$
相识别。设 $U'' \to Z''$ 是平展态射且 $U''$ 仿射。
记 $U \to Z$ 和 $U' \to Z'$ 为相应的仿射概形，它们分别在
$Z$ 和 $Z'$ 上平展。由于 $f$ 形式光滑，存在态射 $h : U'' \to X$，
它与 $i$ 一致（限制到 $U$ 后），并且 $f \circ h$ 等于 $b|_{U''}$。
由于 $Z'$ 是泛一阶加厚，得到唯一态射 $g : U'' \to Z'$，使得
$g = a \circ h$。% P09-ERR-0586
$Z''$ 的泛性质蕴含 $k \circ g$ 是含入映射 $U'' \to Z''$。
因此 $g$ 是 $k$ 的左逆。图示为
$$
\xymatrix{
U \ar[d] \ar[r] & Z' \ar[d]^k \\
U'' \ar[r] \ar[ru]^g & Z''
}
$$
。于是 $g$ 诱导映射
$\mathcal{C}_{Z/Z'}|_U \to \mathcal{C}_{Z/Z''}|_U$，它是映射
$\mathcal{C}_{Z/Z''} \to \mathcal{C}_{Z/Z'}$ 在 $U$ 上的左逆。
\end{proof}











\section{诺特基底上的光滑性}
\label{spaces-more-morphisms-section-smooth-Noetherian}

\noindent
本节对应于《态射进阶》第
\ref{more-morphisms-section-smooth-Noetherian} 节。

\begin{lemma}
\label{spaces-more-morphisms-lemma-lifting-along-artinian-at-point}
设 $S$ 为概形。设 $f : X \to Y$ 是 $S$ 上代数空间的态射。
设 $x \in |X|$。假设 $Y$ 局部诺特且 $f$ 局部有限型。
下列条件等价：
\begin{enumerate}
\item $f$ 在 $x$ 处光滑；
\item 对于每个实线部分交换的图
$$
\xymatrix{
X \ar[d]_f & \Spec(B) \ar[d]^i \ar[l]^-\alpha \\
Y & \Spec(B') \ar[l]_-{\beta} \ar@{-->}[lu]
}
$$
，其中 $B' \to B$ 是局部环的满射，$\Ker(B' \to B)$ 的平方为零，
并且 $\alpha$ 把 $\Spec(B)$ 的闭点映到 $x$，都存在使图交换的虚线箭头；
\item 与 (2) 相同，但让 $B' \to B$ 遍历小扩张（见《代数》定义
\ref{algebra-definition-small-extension}）。
\end{enumerate}
\end{lemma}

\begin{proof}
条件 (1) 意味着存在开子空间 $X' \subset X$，使得 $X' \to Y$ 光滑。
因此，由引理 \ref{spaces-more-morphisms-lemma-smooth-formally-smooth}，(1) 蕴含 (2) 和 (3)。
条件 (2) 显然蕴含条件 (3)。假设 (3)。选取交换图
$$
\xymatrix{
X \ar[d] & U \ar[l] \ar[d] \\
Y & V \ar[l]
}
$$
，其中 $U$ 和 $V$ 仿射，水平箭头平展，并且存在点
$u \in U$ 映到 $x$。接着，考虑 (3) 中的图
$$
\xymatrix{
X \ar[d] & U \ar[l] \ar[d] & \Spec(B) \ar[d]^i \ar[l]^-\alpha  \\
Y & V \ar[l] & \Spec(B') \ar[l]_-{\beta}
}
$$
，不过现针对 $u \in U \to V$。令 $\gamma : \Spec(B') \to X$
为假设 (3) 对 $X$ 成立所给出的箭头。因为 $U \to X$ 平展，故形式平展
（引理 \ref{spaces-more-morphisms-lemma-etale-formally-etale}），态射 $\gamma$ 有唯一提升到 $U$，
并与 $\alpha$ 相容。又因为 $V \to Y$ 平展，因而形式平展，该提升与
$\beta$ 相容。因此 (3) 对 $u \in U \to V$ 成立，由《态射进阶》引理
\ref{more-morphisms-lemma-lifting-along-artinian-at-point} 可知
$U \to V$ 在 $u$ 处光滑。这证明 $X \to Y$ 在 $x$ 处光滑，证明完毕。
\end{proof}

\noindent
有时，只需对以有限型点为“中心”的小扩张检验提升判据这一事实很有用
（见《空间的态射》第
\ref{spaces-morphisms-section-points-finite-type} 节）。

\begin{lemma}
\label{spaces-more-morphisms-lemma-lifting-along-artinian}
设 $S$ 为概形。设 $f : X \to Y$ 是 $S$ 上代数空间的态射。
假设 $Y$ 局部诺特且 $f$ 局部有限型。下列条件等价：
\begin{enumerate}
\item $f$ 光滑；
\item 对于每个实线部分交换的图
$$
\xymatrix{
X \ar[d]_f & \Spec(B) \ar[d]^i \ar[l]^-\alpha \\
Y & \Spec(B') \ar[l]_-{\beta} \ar@{-->}[lu]
}
$$
，其中 $B' \to B$ 是阿廷局部环的小扩张，且 $\beta$ 有限型（！），
都存在使图交换的虚线箭头。
\end{enumerate}
\end{lemma}

\begin{proof}
若 $f$ 光滑，则无穷小提升判据（引理
\ref{spaces-more-morphisms-lemma-smooth-formally-smooth}）说明 $f$ 形式光滑，从而 (2) 成立。

\medskip\noindent
假设 $f$ 不光滑。点 $x \in X$ 中使 $f$ 不光滑者构成闭子集
$T$，它位于 $|X|$ 中。由《空间的态射》引理
\ref{spaces-morphisms-lemma-enough-finite-type-points}，存在点
$x \in T \subset X$，且 $x \in X_{\text{ft-pts}}$。选取交换图
$$
\xymatrix{
X \ar[d] & U \ar[l] \ar[d] & u \ar@{|->}[d] \\
Y & V \ar[l] & v
}
$$
，其中 $U$ 和 $V$ 仿射，水平箭头平展，并且存在点
$u \in U$ 映到 $x$。于是 $u$ 是 $U$ 的有限型点。因为
$U \to V$ 在点 $u$ 处不光滑，由《态射进阶》引理
\ref{more-morphisms-lemma-lifting-along-artinian-at-point}，存在图
$$
\xymatrix{
X \ar[d] & U \ar[l] \ar[d] & \Spec(B) \ar[d]^i \ar[l]^-\alpha  \\
Y & V \ar[l] & \Spec(B') \ar[l]_-{\beta} \ar@{-->}[lu]
}
$$
，其中 $B' \to B$ 是（阿廷）局部环的小扩张，$B$ 的剩余域等于
$\kappa(v)$，并且虚线箭头不存在。由于 $U \to V$ 有限型，
可知 $v$ 是 $V$ 的有限型点。由《态射》引理
\ref{morphisms-lemma-artinian-finite-type}，态射 $\beta$ 有限型，
因而复合 $\Spec(B) \to Y$ 也有限型。% P09-ERR-0587
完全仿照引理 \ref{spaces-more-morphisms-lemma-lifting-along-artinian-at-point} 的证明
（利用 $U \to X$ 和 $V \to Y$ 平展，因而形式平展），可知不存在
箭头 $\Spec(B) \to X$ 嵌入最后一个图的外矩形。% P09-ERR-0588
换言之，(2) 不成立，证明完毕。
\end{proof}

\noindent
下面是一个有用的应用。

\begin{lemma}
\label{spaces-more-morphisms-lemma-check-smoothness-on-infinitesimal-nbhds}
设 $S$ 为概形。设 $f : X \to Y$ 是 $S$ 上代数空间的态射。
假设 $f$ 局部有限型且 $Y$ 局部诺特。设 $Z \subset Y$ 为闭子空间，
其 $n$ 阶无穷小邻域为 $Z_n \subset Y$。令
$X_n = Z_n \times_Y X$。
\begin{enumerate}
\item 若 $X_n \to Z_n$ 对所有 $n$ 都光滑，则 $f$ 在
$f^{-1}(Z)$ 的每一点处光滑。
\item 若 $X_n \to Z_n$ 对所有 $n$ 都平展，则 $f$ 在
$f^{-1}(Z)$ 的每一点处平展。
\end{enumerate}
\end{lemma}

\begin{proof}
假设 $X_n \to Z_n$ 对所有 $n$ 都光滑。设 $x \in X$ 位于 $Z$ 的某点上方。
给定引理 \ref{spaces-more-morphisms-lemma-lifting-along-artinian-at-point} 第 (3) 部分中的
小扩张 $B' \to B$ 和态射 $\alpha$、$\beta$，$B'$ 的极大理想幂零
（因为 $B'$ 是阿廷环），所以态射 $\beta$ 通过 $Z_n$ 分解，而 $\alpha$
通过 $X_n$ 分解；对某个适当的 $n$ 会出现这种情形。于是
$X_n \to Z_n$ 的提升性质
给出图中所需的虚线箭头。这证明 (1)。第 (2) 部分由 (1) 以及下述事实得出：
态射平展当且仅当它光滑且相对维数为 $0$。
\end{proof}











\section{朴素余切复形}
\label{spaces-more-morphisms-section-netherlander}

\noindent
本节接续《位点上的模》第 \ref{sites-modules-section-netherlander} 节，
后者又接续《代数》第 \ref{algebra-section-netherlander} 节中的讨论。

\begin{definition}
\label{spaces-more-morphisms-definition-netherlander}
设 $S$ 为概形。设 $f : X \to Y$ 是 $S$ 上代数空间的态射。
{\it $f$ 的朴素余切复形}是《位点上的模》定义
\ref{sites-modules-definition-cotangent-complex-morphism-ringed-topoi}
对带环拓扑斯态射 $f_{small}$ 所定义的复形；该态射位于 $X$ 与 $Y$ 的
小平展位点之间；
见《空间的性质》引理
\ref{spaces-properties-lemma-morphism-ringed-topoi}。记作 $\NL_f$ 或
$\NL_{X/Y}$。
\end{definition}

\noindent
下面几个引理说明，这一定义与环同态和概形的定义相容，并且
$\NL_{X/Y}$ 是 $D_\QCoh(\mathcal{O}_X)$ 的对象。

\begin{lemma}
\label{spaces-more-morphisms-lemma-NL-etale-localization}
设 $S$ 为概形。考虑交换图
$$
\xymatrix{
U \ar[d]_p \ar[r]_g & V \ar[d]^q \\
X \ar[r]^f & Y
}
$$
，它是 $S$ 上代数空间的图，且 $p$ 和 $q$ 平展。则
有典范识别 $\NL_{X/Y}|_{U_\etale} = \NL_{U/V}$，它位于
$D(\mathcal{O}_U)$ 中。
\end{lemma}

\begin{proof}
朴素余切复形的形成与拉回交换（《位点上的模》引理
\ref{sites-modules-lemma-pullback-NL}），并且有
$p_{small}^{-1}\mathcal{O}_X = \mathcal{O}_U$ 以及
$g_{small}^{-1}\mathcal{O}_{V_\etale} =
p_{small}^{-1}f_{small}^{-1}\mathcal{O}_{Y_\etale}$，
因为由《空间的性质》引理
\ref{spaces-properties-lemma-etale-exact-pullback} 有
$q_{small}^{-1}\mathcal{O}_{Y_\etale} =
\mathcal{O}_{V_\etale}$。逐一定义追踪可得
$\NL_{X/Y}|_{U_\etale} = \NL_{U/V}$。
\end{proof}

\begin{lemma}
\label{spaces-more-morphisms-lemma-NL-compare-spaces-schemes}
设 $S$ 为概形。设 $f : X \to Y$ 是 $S$ 上代数空间的态射。
假设 $X$ 和 $Y$ 分别由概形 $X_0$ 和 $Y_0$ 表示。% P09-ERR-0589
则有典范识别 $\NL_{X/Y} = \epsilon^*\NL_{X_0/Y_0}$；它位于
$D(\mathcal{O}_X)$ 中，其中 $\epsilon$ 如《空间的导出范畴》
第 \ref{spaces-perfect-section-derived-quasi-coherent-etale} 节所述，
而 $\NL_{X_0/Y_0}$ 如《态射进阶》定义
\ref{more-morphisms-definition-netherlander} 所述。
\end{lemma}

\begin{proof}
设 $f_0 : X_0 \to Y_0$ 为与 $f$ 对应的概形态射。有典范映射
$\epsilon^{-1}f_0^{-1}\mathcal{O}_{Y_0} \to f_{small}^{-1}\mathcal{O}_Y$，
它与
$\epsilon^\sharp : \epsilon^{-1}\mathcal{O}_{X_0} \to \mathcal{O}_X$
相容，因为有交换图
$$
\xymatrix{
X_{0, Zar} \ar[d]_{f_0} & X_\etale \ar[l]^\epsilon \ar[d]^f \\
Y_{0, Zar} & Y_\etale \ar[l]_\epsilon
}
$$
；见《空间的导出范畴》评注
\ref{spaces-perfect-remark-match-total-direct-images}。因此由朴素余切复形的
函子性得到典范映射
$$
\epsilon^{-1}\NL_{X_0/Y_0} =
\epsilon^{-1}\NL_{\mathcal{O}_{X_0}/f_0^{-1}\mathcal{O}_{Y_0}} =
\NL_{\epsilon^{-1}\mathcal{O}_{X_0}/\epsilon^{-1}f_0^{-1}\mathcal{O}_{Y_0}}
\to
\NL_{\mathcal{O}_X/f^{-1}_{small}\mathcal{O}_Y} = \NL_{X/Y}
$$
。
为看出诱导映射 $\epsilon^*\NL_{X_0/Y_0} \to \NL_{X/Y}$ 在
$D(\mathcal{O}_X)$ 中是同构，可在几何点的茎上检验（《空间的性质》定理
\ref{spaces-properties-theorem-exactness-stalks}）。设
$\overline{x} : \Spec(k) \to X_0$ 是位于 $x \in X_0$ 上方的几何点，
并令 $\overline{y} = f \circ \overline{x}$ 位于 $y \in Y_0$ 上方。则
$$
\NL_{X/Y, \overline{x}} =
\NL_{\mathcal{O}_{X, \overline{x}}/\mathcal{O}_{Y, \overline{y}}}
$$
。
这是因为在 $\overline{x}$ 处取茎等同于经
$\overline{x} : \Spec(k) \to X$ 取逆像，因而可以应用《位点上的模》引理
\ref{sites-modules-lemma-pullback-NL}。另一方面，有
$$
(\epsilon^*\NL_{X_0/Y_0})_{\overline{x}} =
\NL_{X_0/Y_0, x} \otimes_{\mathcal{O}_{X_0, x}}
\mathcal{O}_{X, \overline{x}} =
\NL_{\mathcal{O}_{X_0, x}/\mathcal{O}_{Y_0, y}}
\otimes_{\mathcal{O}_{X_0, x}} \mathcal{O}_{X, \overline{x}}
$$
。
略去一些细节（提示：使用拉回的茎等于像点处的茎这一事实，见《位点》引理
\ref{sites-lemma-point-morphism-sites}；也使用相应的模版本，见《位点上的模》
引理 \ref{sites-modules-lemma-pullback-stalk}）。注意
$\mathcal{O}_{X, \overline{x}}$ 是 $\mathcal{O}_{X_0, x}$ 的严格亨泽尔化，
$\mathcal{O}_{Y, \overline{y}}$ 也类似（《空间的性质》引理
\ref{spaces-properties-lemma-describe-etale-local-ring}）。因此结论由《代数进阶》
引理 \ref{more-algebra-lemma-henselization-NL} 得出。
\end{proof}

\begin{lemma}
\label{spaces-more-morphisms-lemma-netherlander-quasi-coherent}
设 $S$ 为概形。设 $f : X \to Y$ 是 $S$ 上代数空间的态射。
复形 $\NL_{X/Y}$ 的上同调层拟凝聚，在次数 $-1$、$0$ 之外为零，
并且等于 $\Omega_{X/Y}$（在次数 $0$）。
\end{lemma}

\begin{proof}
由《位点上的模》第 \ref{sites-modules-section-netherlander} 节中朴素余切复形的
构造可知，$\NL_{X/Y}$ 是位于次数 $-1$、$0$ 的复形，其次数 $0$ 上同调为
$\Omega_{X/Y}$（根据第 \ref{spaces-more-morphisms-section-sheaf-differentials} 节中
$\Omega_{X/Y}$ 的构造）。微分层拟凝聚（引理
\ref{spaces-more-morphisms-lemma-module-differentials-quasi-coherent}）。为完成证明，只需说明
$H^{-1}(\NL_{X/Y})$ 拟凝聚。这可通过平展局部检验得出：引理
\ref{spaces-more-morphisms-lemma-NL-etale-localization} 和《空间的性质》引理
\ref{spaces-properties-lemma-characterize-quasi-coherent} 允许这样做；
再由引理 \ref{spaces-more-morphisms-lemma-NL-compare-spaces-schemes} 约化到概形情形；
最后使用概形情形的结果，即《态射进阶》引理
\ref{more-morphisms-lemma-netherlander-quasi-coherent}。
\end{proof}

\begin{lemma}
\label{spaces-more-morphisms-lemma-netherlander-fp}
设 $S$ 为概形。设 $f : X \to Y$ 是 $S$ 上代数空间的态射。
若 $f$ 局部有限表示，则 $\NL_{X/Y}$ 在 $X$ 上平展局部拟同构于复形
$$
\ldots \to 0 \to \mathcal{F}^{-1} \to \mathcal{F}^0 \to 0 \to \ldots
$$
，其各项为拟凝聚 $\mathcal{O}_X$-模，并且 $\mathcal{F}^0$ 有限表示，
$\mathcal{F}^{-1}$ 有限型。
\end{lemma}

\begin{proof}
由引理 \ref{spaces-more-morphisms-lemma-NL-etale-localization}，朴素余切复形的形成与平展局部化交换。
再由引理 \ref{spaces-more-morphisms-lemma-NL-compare-spaces-schemes} 约化到概形情形。
概形情形的结论是《态射进阶》引理
\ref{more-morphisms-lemma-netherlander-fp}。
\end{proof}

\begin{lemma}
\label{spaces-more-morphisms-lemma-NL-formally-smooth}
设 $S$ 为概形。设 $f : X \to Y$ 是 $S$ 上代数空间的态射。
下列条件等价：
\begin{enumerate}
\item $f$ 形式光滑；
\item $H^{-1}(\NL_{X/Y}) = 0$，并且
$H^0(\NL_{X/Y}) = \Omega_{X/Y}$ 局部投射。
\end{enumerate}
\end{lemma}

\begin{proof}
这由引理 \ref{spaces-more-morphisms-lemma-formally-smooth}、
引理 \ref{spaces-more-morphisms-lemma-NL-etale-localization}、
引理 \ref{spaces-more-morphisms-lemma-NL-compare-spaces-schemes} 以及概形情形得出；
后者即《态射进阶》引理
\ref{more-morphisms-lemma-NL-formally-smooth}。
\end{proof}

\begin{lemma}
\label{spaces-more-morphisms-lemma-NL-formally-etale}
设 $f : X \to Y$ 为概形态射。% P09-ERR-0590
下列条件等价：
\begin{enumerate}
\item $f$ 形式平展；
\item $H^{-1}(\NL_{X/Y}) = H^0(\NL_{X/Y}) = 0$。
\end{enumerate}
\end{lemma}

\begin{proof}
假设 (1)。形式平展态射是形式光滑态射。因此，由引理
\ref{spaces-more-morphisms-lemma-NL-formally-smooth}，$H^{-1}(\NL_{X/Y}) = 0$。
另一方面，形式平展态射是形式非分歧的，% P09-ERR-0591
故由引理 \ref{spaces-more-morphisms-lemma-characterize-formally-unramified} 有
$\Omega_{X/Y} = 0$。反过来，若 (2) 成立，则由引理
\ref{spaces-more-morphisms-lemma-NL-formally-smooth}，$f$ 形式光滑；由引理
\ref{spaces-more-morphisms-lemma-characterize-formally-unramified}，它形式非分歧；
因而由引理 \ref{spaces-more-morphisms-lemma-formally-etale-unramified-smooth}，它形式平展。% P09-ERR-0592
\end{proof}

\begin{lemma}
\label{spaces-more-morphisms-lemma-NL-smooth}
设 $f : X \to Y$ 为概形态射。% P09-ERR-0593
下列条件等价：
\begin{enumerate}
\item $f$ 光滑；
\item $f$ 局部有限表示，$H^{-1}(\NL_{X/Y}) = 0$，并且
$H^0(\NL_{X/Y}) = \Omega_{X/Y}$ 有限局部自由。
\end{enumerate}
\end{lemma}

\begin{proof}
这由引理 \ref{spaces-more-morphisms-lemma-formally-smooth}、
引理 \ref{spaces-more-morphisms-lemma-NL-etale-localization}、
引理 \ref{spaces-more-morphisms-lemma-NL-compare-spaces-schemes} 以及概形情形得出；
后者即《态射进阶》引理
\ref{more-morphisms-lemma-NL-smooth}。
\end{proof}











\section{平坦轨迹的开性}
\label{spaces-more-morphisms-section-open-flat}

\noindent
本节对应于《态射进阶》第
\ref{more-morphisms-section-open-flat} 节。% P09-ERR-0594
注意，《空间的态射》第 \ref{spaces-morphisms-section-flat-modules} 节已经定义了
代数空间上拟凝聚模的平坦性。

\begin{theorem}
\label{spaces-more-morphisms-theorem-openness-flatness}
设 $S$ 为概形。设 $f : X \to Y$ 是 $S$ 上代数空间的态射。
设 $\mathcal{F}$ 为 $X$ 上的拟凝聚层。假设 $f$ 局部有限表示，
且 $\mathcal{F}$ 是局部有限表示的 $\mathcal{O}_X$-模。则
$$
\{x \in |X| : \mathcal{F}\text{ is flat over }Y\text{ 在 }x\}
$$
在 $|X|$ 中开。
\end{theorem}

\begin{proof}
选取交换图
$$
\xymatrix{
U \ar[d]_p \ar[r]_\alpha &
V \ar[d]^q \\
X \ar[r]^a & Y
}
$$
，其中 $U$、$V$ 是概形，而 $p$、$q$ 如《空间》引理
\ref{spaces-lemma-lift-morphism-presentations} 中那样满且平展。
由《态射进阶》定理
\ref{more-morphisms-theorem-openness-flatness}，集合
$U' = \{u \in |U| : p^*\mathcal{F}\text{ is flat over }V\text{ 在 }u\}$
在 $U$ 中开。由《空间的态射》定义
\ref{spaces-morphisms-definition-flat-module}，$U'$ 在 $|X|$ 中的像就是定理中的
集合。最后，映射 $|U| \to |X|$ 是开映射，故结论成立；见《空间的性质》引理
\ref{spaces-properties-lemma-topology-points}。
\end{proof}

\begin{lemma}
\label{spaces-more-morphisms-lemma-flat-locus-base-change}
设 $S$ 为概形。设
$$
\xymatrix{
X' \ar[r]_{g'} \ar[d]_{f'} & X \ar[d]^f \\
Y' \ar[r]^g & Y
}
$$
是 $S$ 上代数空间的笛卡尔图。设 $\mathcal{F}$ 是拟凝聚
$\mathcal{O}_X$-模。假设 $g$ 平坦，$f$ 局部有限表示，并且
$\mathcal{F}$ 局部有限表示。则
$$
\{x' \in |X'| : (g')^*\mathcal{F}\text{ is flat over }Y'\text{ 在 }x'\}
$$
是定理 \ref{spaces-more-morphisms-theorem-openness-flatness} 中开子集在连续映射
$|g'| : |X'| \to |X|$ 下的逆像。
\end{lemma}

\begin{proof}
这由《空间的态射》引理
\ref{spaces-morphisms-lemma-base-change-module-flat} 得出。
\end{proof}












\section{平坦性的纤维判据}
\label{spaces-more-morphisms-section-criterion-flat-fibres}

\noindent
设 $S$ 为概形。考虑 $S$ 上代数空间的交换图
$$
\xymatrix{
X \ar[rr]_f \ar[dr]_g & & Y \ar[dl]^h \\
& Z
}
$$
以及拟凝聚 $\mathcal{O}_X$-模 $\mathcal{F}$。
给定一点 $x \in |X|$，我们考察 $\mathcal{F}$ 是否相对于
$Y$ 在 $x$ 处平坦。若 $\mathcal{F}$ 相对于 $Z$ 在 $x$ 处平坦，
则下述定理说明，
这个问题与如下问题密切相关：$\mathcal{F}$ 在 $X \to Z$ 位于
$g(x)$ 上的纤维上的限制，是否相对于 $Y \to Z$ 位于 $g(x)$ 上的纤维
平坦。为使这一说法有意义，我们先给出如下引理。

\begin{lemma}
\label{spaces-more-morphisms-lemma-flat-on-fibres-at-point}
在上述情形中，下列条件等价：
\begin{enumerate}
\item 选取几何点 $\overline{x}$；它是 $X$ 的几何点并位于 $x$ 上方。
令 $\overline{y} = f \circ \overline{x}$ 且
$\overline{z} = g \circ \overline{x}$。则模
$\mathcal{F}_{\overline{x}}/
\mathfrak m_{\overline{z}}\mathcal{F}_{\overline{x}}$
在
$\mathcal{O}_{Y, \overline{y}}/
\mathfrak m_{\overline{z}}\mathcal{O}_{Y, \overline{y}}$
上平坦。
\item 选取态射 $x : \Spec(K) \to X$；它位于 $x$ 的等价类中。
令 $z = g \circ x$、$X_z = \Spec(K) \times_{z, Z} X$、
$Y_z = \Spec(K) \times_{z, Z} Y$，并令 $\mathcal{F}_z$ 为
$\mathcal{F}$ 到 $X_z$ 的拉回。则 $\mathcal{F}_z$ 在 $x$ 处相对于
$Y_z$ 平坦（如《空间的态射》定义
\ref{spaces-morphisms-definition-flat-module} 所定义）。
\item 选取交换图
% P09-ERR-0595
$$
\xymatrix{
& & & U \ar[llld]_a \ar[rr] \ar[dr] & & V \ar[llld]_>>>>>>>b \ar[dl] \\
X \ar[rr]_f \ar[dr]_g & & Y \ar[dl]^h &  & W \ar[llld]_c \\
& Z
}
$$
，其中 $U, V, W$ 是概形，$a, b, c$ 是平展的；并选取一点
$u \in U$，它映到 $x$。令 $w \in W$ 为 $u$ 的像。令
$\mathcal{F}_w$ 为 $\mathcal{F}$ 到纤维 $U_w$ 上的拉回；这是
$U \to W$ 在 $w$ 处的纤维。则 $\mathcal{F}_w$ 相对于 $V_w$
在 $u$ 处平坦。
\end{enumerate}
\end{lemma}

\begin{proof}
注意，在 (2) 中，态射 $x : \Spec(K) \to X$ 给出一个
$K$-有理点，它属于 $X_z$，故该陈述有意义。此外，(2) 中的条件与
$\Spec(K) \to X$ 的选择无关，只要该选择属于 $x$ 的等价类
（略去细节；这也将由下述论证推出，因为
其余条件不依赖于这个选择）。还要注意，我们总能如下选取 (3) 中的图：
首先选取概形 $W$ 以及满平展态射 $W \to Z$，再选取概形 $V$ 以及
满平展态射 $V \to W \times_Z Y$，最后选取概形 $U$ 以及满平展态射
$U \to V \times_Y X$。作出这些选择后，令 $U \to W$ 等于复合
$U \to V \to W$；可以选取一点 $u \in U$ 映到 $x$，因为态射
$U \to X$ 是满的。

\medskip\noindent
假设同时给定 (3) 中的图以及 (1) 中的几何点
$\overline{x} : \Spec(k) \to X$。由《空间的性质》引理
\ref{spaces-properties-lemma-geometric-lift-to-usual}，
可选取几何点 $\overline{u} : \Spec(k) \to U$，它位于 $u$ 上方，
使得 $\overline{x} = a \circ \overline{u}$。
用 $\overline{v} : \Spec(k) \to V$ 和
$\overline{w} : \Spec(k) \to W$ 表示由此诱导的 $V$ 和 $W$ 的几何点。
在此情形下，我们知道
$\mathcal{O}_{X, \overline{x}} = \mathcal{O}_{U, u}^{sh}$，
对 $Y$ 和 $Z$ 亦然；见《空间的性质》引理
\ref{spaces-properties-lemma-describe-etale-local-ring}。
同样有
$$
\mathcal{F}_{\overline{x}} =
(a^*\mathcal{F})_u \otimes_{\mathcal{O}_{U, u}}
\mathcal{O}_{U, u}^{sh}
$$
；见《空间的性质》引理
\ref{spaces-properties-lemma-stalk-quasi-coherent}。
注意，$\mathcal{F}_w$ 在 $u$ 处的茎为
$$
(\mathcal{F}_w)_u = (a^*\mathcal{F})_u/\mathfrak m_w(a^*\mathcal{F})_u
$$
，而 $V_w$ 在 $v$ 处的局部环为
$$
\mathcal{O}_{V_w, v} = \mathcal{O}_{V, v}/\mathfrak m_w\mathcal{O}_{V, v}.
$$
由于 $\mathfrak m_{\overline{z}} =
\mathfrak m_w \mathcal{O}_{Z, \overline{z}} =
\mathfrak m_w \mathcal{O}_{W, w}^{sh}$，可知
\begin{align*}
\mathcal{F}_{\overline{x}}/
\mathfrak m_{\overline{z}}\mathcal{F}_{\overline{x}} & =
(a^*\mathcal{F})_u \otimes_{\mathcal{O}_{U, u}}
\mathcal{O}_{X, \overline{x}}/
\mathfrak m_{\overline{z}}\mathcal{O}_{X, \overline{x}} \\
& =
(\mathcal{F}_w)_u \otimes_{\mathcal{O}_{U_w, u}}
\mathcal{O}_{U, u}^{sh}/\mathfrak m_w\mathcal{O}_{U, u}^{sh} \\
& = (\mathcal{F}_w)_u \otimes_{\mathcal{O}_{U_w, u}}
\mathcal{O}_{U_w, \overline{u}}^{sh} \\
& = (\mathcal{F}_w)_{\overline{u}}
\end{align*}
；倒数第二个等号由《代数》引理
\ref{algebra-lemma-quotient-strict-henselization} 得出，
最后一个等号由《空间的性质》引理
\ref{spaces-properties-lemma-stalk-quasi-coherent} 得出。
对 $V$ 和 $Y$ 的结构层应用同样的论证，得到
$$
\mathcal{O}_{V_w, \overline{v}}^{sh} =
\mathcal{O}_{V, v}^{sh}/\mathfrak m_w \mathcal{O}_{V, v}^{sh} =
\mathcal{O}_{Y, \overline{y}}/
\mathfrak m_{\overline{z}}\mathcal{O}_{Y, \overline{y}}.
$$
现在应用《空间的态射》引理
\ref{spaces-morphisms-lemma-flat-at-point}，可知 (1) 与 (3) 等价。

\medskip\noindent
最后证明 (2) 与 (3) 等价。
为此，选取域扩张 $\tilde K$；它扩张 $K$。再选取态射
$\tilde x : \Spec(\tilde K) \to U$，它位于 $u$ 上方（这是可能的，因为
$u \times_{X, x} \Spec(K)$ 是非空概形）。令
$\tilde z : \Spec(\tilde K) \to U \to W$ 为该复合。% P09-ERR-0596
我们得到交换图
% P09-ERR-0595
$$
\xymatrix{
& & & U_w \times_w \tilde z \ar[llld]_a \ar[rr] \ar[dr] & &
V_w \times_w \tilde z \ar[llld]_>>>>>>>b \ar[dl] \\
X_z \ar[rr]_f \ar[dr]_g & & Y_z \ar[dl]^h &  & \tilde z \ar[llld]_c \\
& z
}
$$
，其中 $z = \Spec(K)$ 且 $w = \Spec(\kappa(w))$。显然，
$\mathcal{F}_w$ 与 $\mathcal{F}_z$ 在 $U_w \times_w \tilde z$ 上的拉回
是同一个模。由此得到交换图
$$
\xymatrix{
X_z \ar[d] & U_w \times_w \tilde z \ar[l] \ar[d] \ar[r] & U_w \ar[d] \\
Y_z & V_w \times_w \tilde z \ar[l] \ar[r] & V_w
}
$$
；它的两个方块都是笛卡尔的，且底部两个水平箭头都是平坦的：
左下水平箭头是态射
$Y \times_Z \tilde z \to Y \times_Z z = Y_z$（平坦态射的基变换）、
平展态射 $V \times_Z \tilde z \to Y \times_Z \tilde z$ 以及
平展态射 $V \times_W \tilde z \to V \times_Z \tilde z$ 的复合。
因此，由《空间的态射》引理
\ref{spaces-morphisms-lemma-base-change-module-flat}，有
$$
\mathcal{F}_z\text{ flat at }x\text{ 相对于 }Y_z
\Leftrightarrow
\mathcal{F}|_{U_w \times_w \tilde z}
\text{ flat at }\tilde x\text{ 相对于 }V_w \times_w \tilde z
\Leftrightarrow
\mathcal{F}_w\text{ flat at }u\text{ 相对于 }V_w
$$
，于是结论成立。
\end{proof}

\begin{definition}
\label{spaces-more-morphisms-definition-module-flat-on-fibre}
设 $S$ 为概形。设 $X \to Y \to Z$ 是 $S$ 上代数空间的态射。
设 $\mathcal{F}$ 为拟凝聚 $\mathcal{O}_X$-模。
设 $x \in |X|$ 为一点，并用 $z \in |Z|$ 表示它的像。
\begin{enumerate}
\item 我们称 {\it $\mathcal{F}$ 到其位于 $z$ 上的纤维的限制在 $x$ 处
相对于 $Y$ 位于 $z$ 上的纤维平坦}，如果引理
\ref{spaces-more-morphisms-lemma-flat-on-fibres-at-point} 的等价条件成立。
\item 我们称 {\it $X$ 位于 $z$ 上的纤维在 $x$ 处相对于 $Y$ 位于
$z$ 上的纤维平坦}，如果引理 \ref{spaces-more-morphisms-lemma-flat-on-fibres-at-point}
的等价条件在 $\mathcal{F} = \mathcal{O}_X$ 时成立。
\item 我们称 {\it $X$ 位于 $z$ 上的纤维相对于 $Y$ 位于 $z$ 上的纤维
平坦}，如果对所有 $x \in |X|$，只要它位于 $z$ 上方，$X$ 位于
$z$ 上的纤维在 $x$ 处相对于 $Y$ 位于 $z$ 上的纤维就平坦
% P09-ERR-0597
\end{enumerate}
\end{definition}

\noindent
有了这个定义，可以如下陈述该判据的一个版本。诺特情形见第
\ref{spaces-more-morphisms-section-flat-over-Noetherian} 节。

\begin{theorem}
\label{spaces-more-morphisms-theorem-criterion-flatness-fibre}
设 $S$ 为概形。
设 $f : X \to Y$ 和 $Y \to Z$ 是 $S$ 上代数空间的态射。
设 $\mathcal{F}$ 为拟凝聚 $\mathcal{O}_X$-模。
假设
\begin{enumerate}
\item $X$ 在 $Z$ 上局部有限表示；
\item $\mathcal{F}$ 是有限表示的 $\mathcal{O}_X$-模；并且 % P09-ERR-0598
\item $Y$ 在 $Z$ 上局部有限型。
\end{enumerate}
设 $x \in |X|$，并令 $y \in |Y|$ 和 $z \in |Z|$ 为 $x$ 的像。
若 $\mathcal{F}_{\overline{x}} \not = 0$，则下列条件等价：% P09-ERR-0599
\begin{enumerate}
\item $\mathcal{F}$ 相对于 $Z$ 在 $x$ 处平坦，且
$\mathcal{F}$ 到其位于 $z$ 上的纤维的限制在 $x$ 处相对于
$Y$ 位于 $z$ 上的纤维平坦；
\item $Y$ 相对于 $Z$ 在 $y$ 处平坦，且 $\mathcal{F}$ 相对于 $Y$
在 $x$ 处平坦。
\end{enumerate}
此外，(1) 和 (2) 成立的点 $x$ 所成的集合在
$\text{Supp}(\mathcal{F})$ 中开。
\end{theorem}

\begin{proof}
选取引理 \ref{spaces-more-morphisms-lemma-flat-on-fibres-at-point} 第 (3) 项中的图。
由定义可知，这把问题化归为概形态射 $U \to V \to W$、
拟凝聚层 $a^*\mathcal{F}$ 以及点 $u \in U$ 的相应定理。
因此，该定理由概形的相应结果，即《态射进阶》定理
\ref{more-morphisms-theorem-criterion-flatness-fibre} 得出。
\end{proof}

\begin{lemma}
\label{spaces-more-morphisms-lemma-morphism-between-flat}
设 $S$ 为概形。
设 $f : X \to Y$ 和 $Y \to Z$ 是 $S$ 上代数空间的态射。% P09-ERR-0600
假设
\begin{enumerate}
\item $X$ 在 $Z$ 上局部有限表示；
\item $X$ 在 $Z$ 上平坦；
\item 对每个 $z \in |Z|$，$X$ 位于 $z$ 上的纤维相对于
$Y$ 位于 $z$ 上的纤维平坦；并且
\item $Y$ 在 $Z$ 上局部有限型。
\end{enumerate}
则 $f$ 平坦。若 $f$ 还是满的，则 $Y$ 在 $Z$ 上平坦。
\end{lemma}

\begin{proof}
这是定理 \ref{spaces-more-morphisms-theorem-criterion-flatness-fibre} 的特殊情形。
\end{proof}

\begin{lemma}
\label{spaces-more-morphisms-lemma-base-change-criterion-flatness-fibre}
设 $S$ 为概形。设 $f : X \to Y$ 和 $Y \to Z$ 是 $S$ 上代数空间的
态射。设 $\mathcal{F}$ 为拟凝聚 $\mathcal{O}_X$-模。
假设
\begin{enumerate}
\item $X$ 在 $Z$ 上局部有限表示；
\item $\mathcal{F}$ 是有限表示的 $\mathcal{O}_X$-模；% P09-ERR-0598
\item $\mathcal{F}$ 在 $Z$ 上平坦；并且
\item $Y$ 在 $Z$ 上局部有限型。
\end{enumerate}
则集合
% P09-ERR-0601
$$
A = \{x \in |X| : \mathcal{F} \text{ flat at }x \text{ 相对于 }Y\}.
$$
在 $|X|$ 中开，且它的形成与任意基变换可交换：
若 $Z' \to Z$ 是代数空间的态射，并且 $A'$ 是
$X' = X \times_Z Z'$ 中使
$\mathcal{F}' = \mathcal{F} \times_Z Z'$ 在
$Y' = Y \times_Z Z'$ 上平坦的点所成的集合，% P09-ERR-0602
则 $A'$ 是 $A$ 在连续映射 $|X'| \to |X|$ 下的逆像。
\end{lemma}

\begin{proof}
一种证明方式是把《态射进阶》引理
\ref{more-morphisms-lemma-morphism-between-flat} 中的证明
翻译到代数空间的范畴中。这里我们改为把它化归到概形情形。
具体而言，选取引理 \ref{spaces-more-morphisms-lemma-flat-on-fibres-at-point} 第 (3) 项中的图，
使得 $a$、$b$、$c$ 都是满的。
由定义可知，这把问题化归为概形态射 $U \to V \to W$、
拟凝聚层 $a^*\mathcal{F}$ 以及点 $u \in U$ 的相应定理。
唯一需要说明的小点是：给定代数空间的态射 $Z' \to Z$，
选取概形 $W'$ 以及满平展态射 $W' \to W \times_Z Z'$。
然后令 $U' = W' \times_W U$ 且 $V' = W' \times_W V$。
用 $a', b', c'$ 表示从 $U', V', W'$ 到 $X', Y', Z'$ 的态射。
在这种情形下，$A$ 和 $A'$ 分别是 $U$ 和 $U'$ 中与
$a^*\mathcal{F}$ 和 $(a')^*\mathcal{F}'$ 相关的开子集的像。
这确实把引理化归为《态射进阶》引理
\ref{more-morphisms-lemma-morphism-between-flat}。
\end{proof}

\begin{lemma}
\label{spaces-more-morphisms-lemma-base-change-flatness-fibres}
设 $S$ 为概形。
设 $f : X \to Y$ 和 $Y \to Z$ 是 $S$ 上代数空间的态射。% P09-ERR-0600
假设
\begin{enumerate}
\item $X$ 在 $Z$ 上局部有限表示；
\item $X$ 在 $Z$ 上平坦；并且
\item $Y$ 在 $Z$ 上局部有限型。
\end{enumerate}
则集合
% P09-ERR-0601
$$
\{x \in |X| : X\text{ flat at }x \text{ 相对于 }Y\}.
$$
在 $|X|$ 中开，且它的形成与任意基变换 $Z' \to Z$ 可交换。
\end{lemma}

\begin{proof}
这是引理 \ref{spaces-more-morphisms-lemma-base-change-criterion-flatness-fibre} 的特殊情形。
\end{proof}

\begin{lemma}
\label{spaces-more-morphisms-lemma-flat-and-free-at-point-fibre}
设 $S$ 为概形。设 $f : X \to Y$ 是 $S$ 上代数空间的态射，
并且局部有限表示。
设 $\mathcal{F}$ 是有限表示的 $\mathcal{O}_X$-模。
设 $x \in |X|$，其像为 $y \in |Y|$。若 $\mathcal{F}$ 在 $x$ 处
相对于 $Y$ 平坦，则下列条件等价：% P09-ERR-0603
\begin{enumerate}
\item $(\mathcal{F}_{\overline{y}})_{\overline{x}}$ 是平坦的
$\mathcal{O}_{X_{\overline{y}}, \overline{x}}$-模；
\item $(\mathcal{F}_{\overline{y}})_{\overline{x}}$ 是自由的
$\mathcal{O}_{X_{\overline{y}}, \overline{x}}$-模；
\item $\mathcal{F}_{\overline{y}}$ 在 $\overline{x}$ 的某个平展邻域上
有限自由；该邻域位于 $X_{\overline{y}}$ 中；并且
\item $\mathcal{F}$ 在 $x$ 的某个平展邻域上有限自由；该邻域位于
$X$ 中。
\end{enumerate}
这里 $\overline{x}$ 是 $X$ 的一个位于 $x$ 上方的几何点，
且 $\overline{y} = f \circ \overline{x}$。
\end{lemma}

\begin{proof}
选取交换图
$$
\xymatrix{
U \ar[d] \ar[r] & V \ar[d] \\
X \ar[r] & Y
}
$$
，其中 $U$ 和 $V$ 是概形，竖直箭头是平展的，且存在一点
$u \in U$ 映到 $x$。令 $v \in V$ 为 $u$ 的像。
把引理 \ref{spaces-more-morphisms-lemma-flat-on-fibres-at-point} 应用于
$\text{id} : X \to X$（底为 $Y$），可知 (1) 转化为条件
``$\mathcal{F}|_{U_v}$ 相对于 $U_v$ 在 $u$ 处平坦''。
换言之，(1) 等价于 $(\mathcal{F}|_{U_v})_u$ 是平坦的
$\mathcal{O}_{U_v, u}$-模。
由概形情形（《态射进阶》引理
\ref{more-morphisms-lemma-flat-and-free-at-point-fibre}），
这推出 $\mathcal{F}|_U$ 在 $u$ 的某个开邻域上有限自由。
由此可见，(1) 推出 (4)。
(4) $\Rightarrow$ (3) 和 (2) $\Rightarrow$ (1) 是立即的。
为证明 (3) $\Rightarrow$ (2)，使用《空间的性质》诸引理
\ref{spaces-properties-lemma-describe-etale-local-ring} 和
\ref{spaces-properties-lemma-stalk-quasi-coherent}
对局部环和茎的描述。
\end{proof}

\begin{lemma}
\label{spaces-more-morphisms-lemma-finite-free-open}
设 $S$ 为概形。设 $f : X \to Y$ 是 $S$ 上代数空间的态射，
并且局部有限表示。
设 $\mathcal{F}$ 是有限表示的 $\mathcal{O}_X$-模，并且在 $Y$ 上平坦。
则集合
$$
\{x \in |X| : \mathcal{F}\text{ free in an \'etale neighbourhood of }x\}
$$
在 $|X|$ 中开，且它的形成与任意基变换 $Y' \to Y$ 可交换。
\end{lemma}

\begin{proof}
开性是显然的。设 $Y' \to Y$ 是代数空间的态射，
令 $X' = Y' \times_Y X$，
并设 $x' \in |X'|$ 是位于 $x \in |X|$ 上方的一点。
由引理 \ref{spaces-more-morphisms-lemma-flat-and-free-at-point-fibre}，
$x$ 属于我们的集合，当且仅当
$(\mathcal{F}_{\overline{y}})_{\overline{x}}$ 是平坦的
$\mathcal{O}_{X_{\overline{y}}, \overline{x}}$-模。
同样，$x'$ 属于拉回 $\mathcal{F}'$ 所对应的类似集合；这里该模是
$\mathcal{F}$ 到 $X'$ 的拉回，当且仅当
$(\mathcal{F}'_{\overline{y}'})_{\overline{x}'}$ 是平坦的
$\mathcal{O}_{X'_{\overline{y}'}, \overline{x}'}$-模
（符号含义显然）。把引理 \ref{spaces-more-morphisms-lemma-flat-on-fibres-at-point} 应用于
态射 $\text{id} : X \to X$（底为 $Y$），可知这两个断言等价。
因此，关于基变换的陈述成立。
\end{proof}







\section{诺特基底上的平坦性}
\label{spaces-more-morphisms-section-flat-over-Noetherian}

\noindent
下面给出诺特情形的``平坦性的纤维判据''。

\begin{theorem}
\label{spaces-more-morphisms-theorem-criterion-flatness-fibre-Noetherian}
设 $S$ 为概形。
设 $f : X \to Y$ 和 $Y \to Z$ 是 $S$ 上代数空间的态射。
设 $\mathcal{F}$ 为拟凝聚 $\mathcal{O}_X$-模。
假设
\begin{enumerate}
\item $X$、$Y$、$Z$ 都局部诺特；并且 % P09-ERR-0604
\item $\mathcal{F}$ 是凝聚 $\mathcal{O}_X$-模。% P09-ERR-0604
\end{enumerate}
设 $x \in |X|$，并令 $y \in |Y|$ 和 $z \in |Z|$ 为 $x$ 的像。
若 $\mathcal{F}_{\overline{x}} \not = 0$，则下列条件等价：% P09-ERR-0605
\begin{enumerate}
\item $\mathcal{F}$ 相对于 $Z$ 在 $x$ 处平坦，且
$\mathcal{F}$ 到其位于 $z$ 上的纤维的限制在 $x$ 处相对于
$Y$ 位于 $z$ 上的纤维平坦；
\item $Y$ 相对于 $Z$ 在 $y$ 处平坦，且 $\mathcal{F}$ 相对于 $Y$
在 $x$ 处平坦。
\end{enumerate}
\end{theorem}

\begin{proof}
选取引理 \ref{spaces-more-morphisms-lemma-flat-on-fibres-at-point} 第 (3) 项中的图。
由定义可知，这把问题化归为概形态射 $U \to V \to W$、
拟凝聚层 $a^*\mathcal{F}$ 以及点 $u \in U$ 的相应定理。
因此，该定理由概形的相应结果，即《态射进阶》定理
\ref{more-morphisms-theorem-criterion-flatness-fibre-Noetherian} 得出。
\end{proof}

\begin{lemma}
\label{spaces-more-morphisms-lemma-morphism-between-flat-Noetherian}
设 $S$ 为概形。
设 $f : X \to Y$ 和 $Y \to Z$ 是 $S$ 上代数空间的态射。% P09-ERR-0606
假设
\begin{enumerate}
\item $X$、$Y$、$Z$ 都局部诺特；% P09-ERR-0604
\item $X$ 在 $Z$ 上平坦；
\item 对每个 $z \in |Z|$，$X$ 位于 $z$ 上的纤维相对于
$Y$ 位于 $z$ 上的纤维平坦。
\end{enumerate}
则 $f$ 平坦。若 $f$ 还是满的，则 $Y$ 在 $Z$ 上平坦。
\end{lemma}

\begin{proof}
这是定理 \ref{spaces-more-morphisms-theorem-criterion-flatness-fibre-Noetherian} 的特殊情形。
\end{proof}

\noindent
正如检验光滑性一样，若基底诺特，只需在阿廷环上检验平坦性。
下面给出一个示例性陈述。

\begin{lemma}
\label{spaces-more-morphisms-lemma-flatness-over-Noetherian-ring}
设 $A$ 为诺特环。设 $I \subset A$ 为理想。
设 $X$ 为在 $S = \Spec(A)$ 上局部有限表示的代数空间。
对 $n \geq 1$，令 $S_n = \Spec(A/I^n)$ 且
$X_n = S_n \times_S X$。设 $\mathcal{F}$ 为凝聚
$\mathcal{O}_X$-模。% P09-ERR-0607
若对每个 $n \geq 1$，$\mathcal{F}_n$ 是 $\mathcal{F}$ 到 $X$ 的拉回，% P09-ERR-0608
并且在 $S_n$ 上平坦，则 $\mathcal{F}$ 在 $X$ 上平坦的（开）轨迹 % P09-ERR-0609
包含 $V(I)$ 在 $X \to S$ 下的逆像。
\end{lemma}

\begin{proof}
$\mathcal{F}$ 在 $S$ 上平坦的轨迹由定理
\ref{spaces-more-morphisms-theorem-openness-flatness} 在 $|X|$ 中开。
把 $X$ 替换为一个平展覆盖的成员不影响该陈述，故可假设 $X$ 是仿射概形。
在这种情形下，结论立即由《代数》引理
\ref{algebra-lemma-flat-module-powers} 得出。
略去一些细节。
\end{proof}







\section{再论正规化}
\label{spaces-more-morphisms-section-normalization}

\noindent
正规化与光滑基变换可交换。

\begin{lemma}
\label{spaces-more-morphisms-lemma-integral-closure-smooth-pullback}
设 $S$ 为概形。设 $f : Y \to X$ 为 $S$ 上代数空间的光滑态射。
设 $\mathcal{A}$ 为拟凝聚 $\mathcal{O}_X$-代数层。则
$\mathcal{O}_Y$ 在 $f^*\mathcal{A}$ 中的整闭包等于 $f^*\mathcal{A}'$，
其中 $\mathcal{A}' \subset \mathcal{A}$ 是 $\mathcal{O}_X$ 在
$\mathcal{A}$ 中的整闭包。
\end{lemma}

\begin{proof}
由整闭包的构造（见《空间的态射》定义
\ref{spaces-morphisms-definition-integral-closure}），
这立即化归为 $X$ 和 $Y$ 均仿射的情形。
在这种情形下，结论就是《代数》引理
\ref{algebra-lemma-integral-closure-commutes-smooth}。
\end{proof}

\begin{lemma}[正规化与光滑基变换可交换]
\label{spaces-more-morphisms-lemma-normalization-smooth-localization}
设 $S$ 为概形。设
$$
\xymatrix{
Y_2 \ar[r] \ar[d] & Y_1 \ar[d]^f \\
X_2 \ar[r]^\varphi & X_1
}
$$
是 $S$ 上代数空间的笛卡尔方块。假设 $f$ 拟紧且拟分离，
并且 $\varphi$ 光滑。
设 $Y_i \to X_i' \to X_i$ 为 $X_i$ 在 $Y_i$ 中的正规化。
则 $X_2' \cong X_2 \times_{X_1} X_1'$。
\end{lemma}

\begin{proof}
把分解 $Y_1 \to X_1' \to X_1$ 基变换到 $X_2$，得到分解
$Y_2 \to X_2 \times_{X_1} X_1' \to X_1$，且
$X_2 \times_{X_1} X_1' \to X_1$ 是整态射 % P09-ERR-0610
（《空间的态射》引理 \ref{spaces-morphisms-lemma-base-change-integral}）。
因此，由《空间的态射》引理
\ref{spaces-morphisms-lemma-characterize-normalization} 的泛性质，
得到态射 $h : X_2' \to X_2 \times_{X_1} X_1'$。
注意，$X_2'$ 是 $\mathcal{O}_{X_2}$ 在
$f_{2, *}\mathcal{O}_{Y_2}$ 中的整闭包的相对谱。
若 $\mathcal{A}' \subset f_{1, *}\mathcal{O}_{Y_1}$ 表示
$\mathcal{O}_{X_2}$ 的整闭包，% P09-ERR-0611
则 $X_2 \times_{X_1} X_1'$ 是 $\varphi^*\mathcal{A}'$ 的相对谱，
因为相对谱的构造与任意基变换可交换。由《空间的上同调》引理
\ref{spaces-cohomology-lemma-flat-base-change-cohomology}，
我们知道 $f_{2, *}\mathcal{O}_{Y_2} = \varphi^*f_{1, *}\mathcal{O}_{Y_1}$。
因此，结论由引理 \ref{spaces-more-morphisms-lemma-integral-closure-smooth-pullback} 得出。
\end{proof}









\section{Cohen--Macaulay 态射}
\label{spaces-more-morphisms-section-CM}

\noindent
本节对应于《态射进阶》第 \ref{more-morphisms-section-CM} 节。

\begin{lemma}
\label{spaces-more-morphisms-lemma-CM-local-ring-fibre}
概形芽的态射的性质
\begin{align*}
& \mathcal{P}((X, x) \to (S, s)) = \\
& \text{纤维的局部环 }
\mathcal{O}_{X_s, x}
\text{ 是诺特且 Cohen--Macaulay 的}
\end{align*}
在源和目标上都是平展局部的（《下降》定义
\ref{descent-definition-local-source-target-at-point}）。
\end{lemma}

\begin{proof}
给定《下降》定义
\ref{descent-definition-local-source-target-at-point} 中的图，
我们得到纤维的平展态射 $U'_{v'} \to U_v$，它把 $u'$ 映到 $u$；
见《下降》引理 \ref{descent-lemma-etale-on-fiber}。
因此，局部环 $\mathcal{O}_{U'_{v'}, u'}$ 和
$\mathcal{O}_{U_v, u}$ 的严格亨泽尔化相同。结论由《代数进阶》引理
\ref{more-algebra-lemma-henselization-CM} 得出。
\end{proof}

\begin{definition}
\label{spaces-more-morphisms-definition-CM}
设 $S$ 为概形。
设 $f : X \to Y$ 为 $S$ 上代数空间的态射。
假设 $f$ 的纤维局部诺特（《空间上的除子》定义
\ref{spaces-divisors-definition-locally-Noetherian-fibre}）。
\begin{enumerate}
\item 设 $x \in |X|$ 且 $y = f(x)$。我们称 $f$
{\it 在 $x$ 处为 Cohen--Macaulay}，如果 $f$ 在 $x$ 处平坦，并且
《空间的态射》引理
\ref{spaces-morphisms-lemma-local-source-target-at-point} 中的等价条件
对于引理 \ref{spaces-more-morphisms-lemma-CM-local-ring-fibre} 所述性质 $\mathcal{P}$ 成立。
\item 我们称 $f$ 为 {\it Cohen--Macaulay 态射}，如果 $f$ 在 $X$
的每一点处都是 Cohen--Macaulay 的。
\end{enumerate}
\end{definition}

\noindent
下面给出相应的转述。

\begin{lemma}
\label{spaces-more-morphisms-lemma-CM}
设 $S$ 为概形。设 $f : X \to Y$ 为 $S$ 上代数空间的态射。
假设 $f$ 的纤维局部诺特。下列条件等价：% P09-ERR-0612
\begin{enumerate}
\item $f$ 是 Cohen--Macaulay 的；
\item $f$ 平坦，并且存在满平展态射 $V \to Y$，其中 $V$ 是概形，
使得 $X_V \to V$ 的纤维都是 Cohen--Macaulay 代数空间；并且
\item $f$ 平坦，并且对任意平展态射 $V \to Y$，其中 $V$ 是概形，
$X_V \to V$ 的纤维都是 Cohen--Macaulay 代数空间。
\end{enumerate}
给定 $x \in |X|$，其像为 $y \in |Y|$，下列条件等价：% P09-ERR-0612
\begin{enumerate}
\item[(a)] $f$ 在 $x$ 处为 Cohen--Macaulay；
\item[(b)] $\mathcal{O}_{Y, \overline{y}} \to \mathcal{O}_{X, \overline{x}}$
平坦，且
$\mathcal{O}_{X, \overline{x}}/
\mathfrak m_{\overline{y}}\mathcal{O}_{X, \overline{x}}$
是 Cohen--Macaulay 的。% P09-ERR-0613
\end{enumerate}
\end{lemma}

\begin{proof}
给定平展态射 $V \to Y$，其中 $V$ 是概形；选取概形 $U$ 以及
满平展态射 $U \to X \times_Y V$。考虑交换图
$$
\xymatrix{
U \ar[d] \ar[r] & V \ar[d] \\
X \ar[r] & Y
}
$$
设 $u \in U$，其像为 $x \in |X|$、$y \in |Y|$ 和 $v \in V$。
则 $f$ 在 $x$ 处为 Cohen--Macaulay，当且仅当 $U \to V$ 在 $u$ 处为
Cohen--Macaulay（由定义）。此外，态射
$U_v \to X_v = (X_V)_v$ 是满平展的。因此，概形 $U_v$ 是
Cohen--Macaulay 的，当且仅当代数空间 $X_v$ 是 Cohen--Macaulay 的。
于是，(1)、(2)、(3) 的等价性由概形态射的相应等价性通过形式论证得出；
见《态射进阶》引理 \ref{more-morphisms-lemma-CM}。

\medskip\noindent
证明 (a) 与 (b) 等价。平坦性的相应等价性即《空间的态射》引理
\ref{spaces-morphisms-lemma-flat-at-point-etale-local-rings}。
因此，在证明该等价性时，可以假设 $f$ 在 $x$ 处平坦。
考虑上述图以及 $x, y, u, v$。则
$\mathcal{O}_{Y, \overline{y}} \to \mathcal{O}_{X, \overline{x}}$
等于局部环严格亨泽尔化之间的映射
$\mathcal{O}_{V, v}^{sh} \to \mathcal{O}_{U, u}^{sh}$；
见《空间的性质》引理
\ref{spaces-properties-lemma-describe-etale-local-ring}。
因此有
$$
\mathcal{O}_{X, \overline{x}}/
\mathfrak m_{\overline{y}}\mathcal{O}_{X, \overline{x}} =
(\mathcal{O}_{U, u}/\mathfrak m_v \mathcal{O}_{U, u})^{sh}
$$
；这是《代数》引理
\ref{algebra-lemma-quotient-strict-henselization}。
故只需证明诺特局部环
$\mathcal{O}_{U, u}/\mathfrak m_v \mathcal{O}_{U, u}$
为 Cohen--Macaulay，当且仅当其严格亨泽尔化为 Cohen--Macaulay。
这就是《代数进阶》引理 \ref{more-algebra-lemma-henselization-CM}。
\end{proof}

\begin{lemma}
\label{spaces-more-morphisms-lemma-composition-CM}
设 $S$ 为概形。
设 $f : X \to Y$ 和 $g : Y \to Z$ 是 $S$ 上代数空间的态射。
假设 $f$、$g$ 和 $g \circ f$ 的纤维局部诺特。
设 $x \in |X|$，其像为 $y \in |Y|$ 和 $z \in |Z|$。
\begin{enumerate}
\item 若 $f$ 在 $x$ 处为 Cohen--Macaulay，且 $g$ 在 $f(x)$ 处为
Cohen--Macaulay，则 $g \circ f$ 在 $x$ 处为 Cohen--Macaulay。
\item 若 $f$ 和 $g$ 都是 Cohen--Macaulay 的，则 $g \circ f$ 也是
Cohen--Macaulay 的。
\item 若 $g \circ f$ 在 $x$ 处为 Cohen--Macaulay，且 $f$ 在 $x$ 处平坦，
则 $f$ 在 $x$ 处为 Cohen--Macaulay，且 $g$ 在 $f(x)$ 处为
Cohen--Macaulay。
\item 若 $f \circ g$ 是 Cohen--Macaulay 的，且 $f$ 平坦，% P09-ERR-0614
则 $f$ 是 Cohen--Macaulay 的，且 $g$ 在 $f$ 的像中的每一点处都是
Cohen--Macaulay 的。
\end{enumerate}
\end{lemma}

\begin{proof}
在平展局部上工作，结论由概形的相应结果得出；见《态射进阶》引理
\ref{more-morphisms-lemma-composition-CM}。
或者，可以使用引理 \ref{spaces-more-morphisms-lemma-CM} 中 (a) 与 (b) 的等价性。
因此，考虑诺特局部环的局部同态
$$
\mathcal{O}_{Y, \overline{y}}/
\mathfrak m_{\overline{z}}\mathcal{O}_{Y, \overline{y}}
\longrightarrow
\mathcal{O}_{X, \overline{x}}/
\mathfrak m_{\overline{z}}\mathcal{O}_{X, \overline{x}}
$$
，其纤维为
$$
\mathcal{O}_{X, \overline{x}}/
\mathfrak m_{\overline{y}}\mathcal{O}_{X, \overline{x}}
$$
；然后应用《代数》引理 \ref{algebra-lemma-CM-goes-up}。
\end{proof}

\begin{lemma}
\label{spaces-more-morphisms-lemma-flat-morphism-from-CM}
设 $S$ 为概形。
设 $f : X \to Y$ 为 $S$ 上局部诺特代数空间之间的平坦态射。
若 $X$ 是 Cohen--Macaulay 的，则 $f$ 是 Cohen--Macaulay 的，并且
$\mathcal{O}_{Y, f(\overline{x})}$ 对所有 $x \in |X|$ 都是
Cohen--Macaulay 的。% P09-ERR-0615
\end{lemma}

\begin{proof}
利用引理 \ref{spaces-more-morphisms-lemma-CM} 转译为代数（与引理
\ref{spaces-more-morphisms-lemma-composition-CM} 的证明比较）后，结论由《代数》引理
\ref{algebra-lemma-CM-goes-up} 得出。
\end{proof}

\begin{lemma}
\label{spaces-more-morphisms-lemma-base-change-CM}
设 $S$ 为概形。
设 $f : X \to Y$ 为 $S$ 上代数空间的态射。
假设 $f$ 的纤维局部诺特。
设 $Y' \to Y$ 局部有限型。设 $f' : X' \to Y'$ 为 $f$ 的基变换。
设 $x' \in |X'|$ 为一点，其像为 $x \in |X|$。
\begin{enumerate}
\item 若 $f$ 在 $x$ 处为 Cohen--Macaulay，则
$f' : X' \to Y'$ 在 $x'$ 处为 Cohen--Macaulay。
\item 若 $f$ 在 $x$ 处平坦，且 $f'$ 在 $x'$ 处为 Cohen--Macaulay，
则 $f$ 在 $x$ 处为 Cohen--Macaulay。
\item 若 $Y' \to Y$ 在 $f'(x')$ 处平坦，且 $f'$ 在 $x'$ 处为
Cohen--Macaulay，则 $f$ 在 $x$ 处为 Cohen--Macaulay。
\end{enumerate}
\end{lemma}

\begin{proof}
用 $y \in |Y|$ 和 $y' \in |Y'|$ 表示 $x'$ 的像。% P09-ERR-0616
选取满平展态射 $V \to Y$，其中 $V$ 是概形。
选取满平展态射 $U \to X \times_Y V$，其中 $U$ 是概形。
选取满平展态射 $V' \to Y' \times_Y V$，其中 $V'$ 是概形。% P09-ERR-0617
则 $U' = U \times_V V'$ 是概形，并带有满平展态射 $U' \to X'$。
选取 $u' \in U'$ 映到 $x'$。用 $u \in U$ 表示 $u'$ 的像。
于是，引理由 $U \to V$、其基变换 $U' \to V'$ 以及点 $u'$、$u$
的引理得出（这由定义可知）。
故引理由概形情形得出；见《态射进阶》引理
\ref{more-morphisms-lemma-base-change-CM}。
\end{proof}

\begin{lemma}
\label{spaces-more-morphisms-lemma-flat-finite-presentation-CM-open}
设 $S$ 为概形。
设 $f : X \to Y$ 为 $S$ 上代数空间的态射，并且平坦且局部有限表示。
令
$$
W = \{x \in |X| : f\text{ 为 Cohen--Macaulay 于点 }x\}
$$
则 $W$ 在 $|X|$ 中开，且 $W$ 的形成与 $f$ 的任意基变换可交换：
对任意态射 $g : Y' \to Y$，考虑基变换 $f' : X' \to Y'$（它是
$f$ 的基变换）以及投影 $g' : X' \to X$。
则相应集合 $W'$（针对态射 $f'$）等于 $W' = (g')^{-1}(W)$。% P09-ERR-0618
\end{lemma}

\begin{proof}
选取交换图
$$
\xymatrix{
U \ar[d] \ar[r] & V \ar[d] \\
X \ar[r] & Y
}
$$
，其中竖直箭头平展，而 $U$ 和 $V$ 是概形。
设 $u \in U$，其像为 $x \in |X|$。
则 $f$ 在 $x$ 处为 Cohen--Macaulay，当且仅当 $U \to V$ 在 $u$ 处为
Cohen--Macaulay（由定义）。
因而化归为态射 $U \to V$ 的情形。见《态射进阶》引理
\ref{more-morphisms-lemma-flat-finite-presentation-CM-open}。
\end{proof}

\begin{lemma}
\label{spaces-more-morphisms-lemma-lfp-CM-relative-dimension}
设 $S$ 为概形。设 $f : X \to Y$ 为 $S$ 上代数空间的态射。
假设 $f$ 局部有限表示且为 Cohen--Macaulay。
则存在开闭子概形 $X_d \subset X$，% P09-ERR-0619
使得 $X = \coprod_{d \geq 0} X_d$，且 $f|_{X_d} : X_d \to Y$
具有相对维数 $d$。
\end{lemma}

\begin{proof}
选取交换图
$$
\xymatrix{
U \ar[d] \ar[r] & V \ar[d] \\
X \ar[r] & Y
}
$$
，其中竖直箭头平展，而 $U$ 和 $V$ 是概形。
则 $U \to V$ 局部有限表示且为 Cohen--Macaulay（由定义立即得出）。
因此有分解 $U = \coprod_{d \geq 0} U_d$，其中每个分量都是开闭子概形，
且 $f|_{U_d} : U_d \to V$ 具有相对维数 $d$；% P09-ERR-0620
见《态射》引理
\ref{morphisms-lemma-flat-finite-presentation-CM-fibres-relative-dimension}。
设 $u \in U$，其像为 $x \in |X|$。则 $f$ 具有相对维数 $d$
在 $x$ 处成立，当且仅当 $U \to V$ 具有相对维数 $d$ 在 $u$ 处成立
（这由定义得出）。
由此可见，$U_d$ 是子集 $X_d \subset |X|$ 的逆像，
而该子集必为开闭的。用 $X_d$ 表示 $X$ 的相应开闭代数子空间，
即可看出引理成立。
\end{proof}







\section{Gorenstein 态射}
\label{spaces-more-morphisms-section-gorenstein}

\noindent
本节对应于《概形的对偶性》第
\ref{duality-section-gorenstein-morphisms} 节。

\begin{lemma}
\label{spaces-more-morphisms-lemma-gorenstein-local-ring-fibre}
概形芽的态射的下述性质
\begin{align*}
& \mathcal{P}((X, x) \to (S, s)) = \\
& \text{纤维的局部环 }
\mathcal{O}_{X_s, x}
\text{ 是诺特且 Gorenstein 的}
\end{align*}
在源和靶上为平展局部的（《下降》定义
\ref{descent-definition-local-source-target-at-point}）。
\end{lemma}

\begin{proof}
给定《下降》定义
\ref{descent-definition-local-source-target-at-point} 中的图，
我们得到纤维的平展态射 $U'_{v'} \to U_v$，它把 $u'$ 映到 $u$；
见《下降》引理 \ref{descent-lemma-etale-on-fiber}。
因此，$\mathcal{O}_{U_v, u} \to \mathcal{O}_{U'_{v'}, u'}$
是一个平展环映射的局部化。故前者是诺特的，当且仅当后者是诺特的；
见《代数进阶》引理 \ref{more-algebra-lemma-Noetherian-etale-extension}。
又因为
$\mathcal{O}_{U'_{v'}, u'}/\mathfrak m_u \mathcal{O}_{U'_{v'}, u'}
= \kappa(u')$（《代数》引理 \ref{algebra-lemma-etale-at-prime}）
是 Gorenstein 环，所以由《对偶化复形》引理
\ref{dualizing-lemma-flat-under-gorenstein} 可见，
$\mathcal{O}_{U_v, u}$ 是 Gorenstein 的，当且仅当
$\mathcal{O}_{U'_{v'}, u'}$ 是 Gorenstein 的。
\end{proof}

\begin{definition}
\label{spaces-more-morphisms-definition-gorenstein}
设 $S$ 为概形。
设 $f : X \to Y$ 为 $S$ 上代数空间的态射。
假设 $f$ 的纤维局部诺特（《空间上的除子》定义
\ref{spaces-divisors-definition-locally-Noetherian-fibre}）。
\begin{enumerate}
\item 设 $x \in |X|$ 且 $y = f(x)$。我们称 $f$
{\it 在 $x$ 处为 Gorenstein}，如果 $f$ 在 $x$ 处平坦，并且
《空间的态射》引理
\ref{spaces-morphisms-lemma-local-source-target-at-point} 中的等价条件
对于引理 \ref{spaces-more-morphisms-lemma-gorenstein-local-ring-fibre} 所述性质
$\mathcal{P}$ 成立。
\item 我们称 $f$ 为 {\it Gorenstein 态射}，如果 $f$ 在 $X$
的每一点处都是 Gorenstein 的。
\end{enumerate}
\end{definition}

\noindent
下面给出相应的转述。

\begin{lemma}
\label{spaces-more-morphisms-lemma-gorenstein}
设 $S$ 为概形。设 $f : X \to Y$ 为 $S$ 上代数空间的态射。
假设 $f$ 的纤维局部诺特。下列条件等价：% P09-ERR-0621
\begin{enumerate}
\item $f$ 是 Gorenstein 的；
\item $f$ 平坦，并且存在满平展态射 $V \to Y$，其中 $V$ 是概形，
使得 $X_V \to V$ 的纤维都是 Gorenstein 代数空间；并且
\item $f$ 平坦，并且对任意平展态射 $V \to Y$，其中 $V$ 是概形，
$X_V \to V$ 的纤维都是 Gorenstein 代数空间。
\end{enumerate}
给定 $x \in |X|$，其像为 $y \in |Y|$，下列条件等价：% P09-ERR-0621
\begin{enumerate}
\item[(a)] $f$ 在 $x$ 处为 Gorenstein；
\item[(b)] $\mathcal{O}_{Y, \overline{y}} \to \mathcal{O}_{X, \overline{x}}$
平坦，并且
$\mathcal{O}_{X, \overline{x}}/
\mathfrak m_{\overline{y}}\mathcal{O}_{X, \overline{x}}$
是 Gorenstein 的。% P09-ERR-0622
\end{enumerate}
\end{lemma}

\begin{proof}
给定平展态射 $V \to Y$，其中 $V$ 是概形；选取概形 $U$ 以及
满平展态射 $U \to X \times_Y V$。考虑交换图
$$
\xymatrix{
U \ar[d] \ar[r] & V \ar[d] \\
X \ar[r] & Y
}
$$
设 $u \in U$，其像为 $x \in |X|$、$y \in |Y|$ 和 $v \in V$。
则 $f$ 在 $x$ 处为 Gorenstein，当且仅当 $U \to V$ 在 $u$ 处为
Gorenstein（由定义）。此外，态射
$U_v \to X_v = (X_V)_v$ 是满平展的。因此，概形 $U_v$ 是
Gorenstein 的，当且仅当代数空间 $X_v$ 是 Gorenstein 的。
于是，(1)、(2) 与 (3) 的等价性由概形态射的相应等价性形式地得出；
见《概形的对偶性》引理 \ref{duality-lemma-gorenstein}。

\medskip\noindent
证明 (a) 与 (b) 等价。关于平坦性的相应等价性是《空间的态射》引理
\ref{spaces-morphisms-lemma-flat-at-point-etale-local-rings}。
因此，在证明该等价性时可以假设 $f$ 在 $x$ 处平坦。
考虑上述图以及 $x, y, u, v$。于是
$\mathcal{O}_{Y, \overline{y}} \to \mathcal{O}_{X, \overline{x}}$
等于局部环的严格亨泽尔化之间的映射
$\mathcal{O}_{V, v}^{sh} \to \mathcal{O}_{U, u}^{sh}$；
见《空间的性质》引理
\ref{spaces-properties-lemma-describe-etale-local-ring}。
故由《代数》引理 \ref{algebra-lemma-quotient-strict-henselization}，有
$$
\mathcal{O}_{X, \overline{x}}/
\mathfrak m_{\overline{y}}\mathcal{O}_{X, \overline{x}} =
(\mathcal{O}_{U, u}/\mathfrak m_v \mathcal{O}_{U, u})^{sh}
$$
因此，只须证明诺特局部环
$\mathcal{O}_{U, u}/\mathfrak m_v \mathcal{O}_{U, u}$
为 Gorenstein，当且仅当其严格亨泽尔化为 Gorenstein。
这立即由《对偶化复形》引理
\ref{dualizing-lemma-completion-henselization-dualizing}
以及 Gorenstein 局部环的定义得出：它是一个以自身为其对偶化复形的诺特局部环。
\end{proof}

\begin{lemma}
\label{spaces-more-morphisms-lemma-composition-gorenstein}
设 $S$ 为概形。
设 $f : X \to Y$ 和 $g : Y \to Z$ 是 $S$ 上代数空间的态射。
假设 $f$、$g$ 和 $g \circ f$ 的纤维局部诺特。
设 $x \in |X|$，其像为 $y \in |Y|$ 和 $z \in |Z|$。
\begin{enumerate}
\item 若 $f$ 在 $x$ 处为 Gorenstein，且 $g$ 在 $f(x)$ 处为
Gorenstein，则 $g \circ f$ 在 $x$ 处为 Gorenstein。
\item 若 $f$ 和 $g$ 都是 Gorenstein 的，则 $g \circ f$ 也是
Gorenstein 的。
\item 若 $g \circ f$ 在 $x$ 处为 Gorenstein，且 $f$ 在 $x$ 处平坦，
则 $f$ 在 $x$ 处为 Gorenstein，且 $g$ 在 $f(x)$ 处为 Gorenstein。
\item 若 $f \circ g$ 是 Gorenstein 的，且 $f$ 平坦，% P09-ERR-0623
则 $f$ 是 Gorenstein 的，且 $g$ 在 $f$ 的像中的每一点处都是
Gorenstein 的。
\end{enumerate}
\end{lemma}

\begin{proof}
在平展局部上工作，结论由概形的相应结果得出；见《概形的对偶性》引理
\ref{duality-lemma-composition-gorenstein}。
或者，可以使用引理 \ref{spaces-more-morphisms-lemma-gorenstein} 中 (a) 与 (b) 的等价性。
因此考虑诺特局部环的局部同态
$$
\mathcal{O}_{Y, \overline{y}}/
\mathfrak m_{\overline{z}}\mathcal{O}_{Y, \overline{y}}
\longrightarrow
\mathcal{O}_{X, \overline{x}}/
\mathfrak m_{\overline{z}}\mathcal{O}_{X, \overline{x}}
$$
其纤维为
$$
\mathcal{O}_{X, \overline{x}}/
\mathfrak m_{\overline{y}}\mathcal{O}_{X, \overline{x}}
$$
；然后应用《对偶化复形》引理
\ref{dualizing-lemma-flat-under-gorenstein}。
\end{proof}

\begin{lemma}
\label{spaces-more-morphisms-lemma-flat-morphism-from-gorenstein}
设 $S$ 为概形。
设 $f : X \to Y$ 为 $S$ 上局部诺特代数空间之间的平坦态射。
若 $X$ 是 Gorenstein 的，则 $f$ 是 Gorenstein 的，并且
$\mathcal{O}_{Y, f(\overline{x})}$ 对所有 $x \in |X|$ 都是
Gorenstein 的。% P09-ERR-0624
\end{lemma}

\begin{proof}
利用引理 \ref{spaces-more-morphisms-lemma-gorenstein} 转译为代数（与引理
\ref{spaces-more-morphisms-lemma-composition-gorenstein} 的证明比较）后，结论由
《对偶化复形》引理 \ref{dualizing-lemma-flat-under-gorenstein} 得出。
\end{proof}

\begin{lemma}
\label{spaces-more-morphisms-lemma-base-change-gorenstein}
设 $S$ 为概形。
设 $f : X \to Y$ 为 $S$ 上代数空间的态射。
假设 $f$ 的纤维局部诺特。
设 $Y' \to Y$ 局部有限型。设 $f' : X' \to Y'$ 为 $f$ 的基变换。
设 $x' \in |X'|$ 为一点，其像为 $x \in |X|$。
\begin{enumerate}
\item 若 $f$ 在 $x$ 处为 Gorenstein，则
$f' : X' \to Y'$ 在 $x'$ 处为 Gorenstein。
\item 若 $f$ 在 $x$ 处平坦，且 $f'$ 在 $x'$ 处为 Gorenstein，
则 $f$ 在 $x$ 处为 Gorenstein。
\item 若 $Y' \to Y$ 在 $f'(x')$ 处平坦，且 $f'$ 在 $x'$ 处为
Gorenstein，则 $f$ 在 $x$ 处为 Gorenstein。
\end{enumerate}
\end{lemma}

\begin{proof}
用 $y \in |Y|$ 和 $y' \in |Y'|$ 表示 $x'$ 的像。% P09-ERR-0625
选取满平展态射 $V \to Y$，其中 $V$ 是概形。
选取满平展态射 $U \to X \times_Y V$，其中 $U$ 是概形。
选取满平展态射 $V' \to Y' \times_Y V$，其中 $V'$ 是概形。% P09-ERR-0626
则 $U' = U \times_V V'$ 是概形，并带有满平展态射 $U' \to X'$。
选取 $u' \in U'$ 映到 $x'$。用 $u \in U$ 表示 $u'$ 的像。
于是，引理由 $U \to V$、其基变换 $U' \to V'$ 以及点 $u'$、$u$
的引理得出（这由定义可知）。
故引理由概形情形得出；见《概形的对偶性》引理
\ref{duality-lemma-base-change-gorenstein}。
\end{proof}

\begin{lemma}
\label{spaces-more-morphisms-lemma-flat-finite-presentation-gorenstein-open}
设 $S$ 为概形。
设 $f : X \to Y$ 为 $S$ 上代数空间的态射，并且平坦且局部有限表示。
令
$$
W = \{x \in |X| : f\text{ 为 Gorenstein 于点 }x\}
$$
则 $W$ 在 $|X|$ 中开，且 $W$ 的形成与 $f$ 的任意基变换可交换：
对任意态射 $g : Y' \to Y$，考虑基变换 $f' : X' \to Y'$（它是
$f$ 的基变换）以及投影 $g' : X' \to X$。
则相应集合 $W'$（针对态射 $f'$）等于 $W' = (g')^{-1}(W)$。% P09-ERR-0627
\end{lemma}

\begin{proof}
选取交换图% P09-ERR-0628
$$
\xymatrix{
U \ar[d] \ar[r] & V \ar[d] \\
X \ar[r] & Y
}
$$
设 $u \in U$，其像为 $x \in |X|$。
则 $f$ 在 $x$ 处为 Gorenstein，当且仅当 $U \to V$ 在 $u$ 处为
Gorenstein（由定义）。
因而化归为概形态射 $U \to V$ 的情形。开性由
《概形的对偶性》引理
\ref{duality-lemma-flat-finite-presentation-characterize-gorenstein} 证明，
与基变换的相容性由《概形的对偶性》引理
\ref{duality-lemma-flat-lft-base-change-gorenstein} 证明。
\end{proof}






\section{Cohen--Macaulay 态射的切片}
\label{spaces-more-morphisms-section-slice}

\noindent
设 $S$ 为概形。设 $X$ 为 $S$ 上的代数空间。
设 $f_1, \ldots, f_r \in \Gamma(X, \mathcal{O}_X)$。在此情形下，
用 $V(f_1, \ldots, f_r)$ 表示 {\it $X$ 中由 $f_1, \ldots, f_r$
切出的闭子空间}。更确切地，可以把 $V(f_1, \ldots, f_r)$
定义为 $X$ 的闭子空间，它对应于由 $f_1, \ldots, f_r$ 生成的
拟凝聚理想层；见《空间的态射》引理
\ref{spaces-morphisms-lemma-closed-immersion-ideals}。
也可以选取表示 $X = U/R$，并考虑闭子概形 $Z \subset U$，它由
$f_1|U, \ldots, f_r|_U$ 切出。显然，$Z$ 是 $R$-不变闭子概形
（见《广群》定义 \ref{groupoids-definition-invariant-open}），并可令
$V(f_1, \ldots, f_r) = Z/R_Z$。

\begin{lemma}
\label{spaces-more-morphisms-lemma-slice}
设 $S$ 为概形。考虑笛卡尔图
$$
\xymatrix{
X \ar[d] & F \ar[l]^p \ar[d] \\
Y & \Spec(k) \ar[l]
}
$$
其中 $X \to Y$ 是 $S$ 上代数空间的平坦且局部有限表示的态射，
而 $k$ 是 $S$ 上的域。设
$f_1, \ldots, f_r \in \Gamma(X, \mathcal{O}_X)$，且 $z \in |F|$，
使得 $f_1, \ldots, f_r$ 在局部环
$\mathcal{O}_{F, \overline{z}}$ 中的像构成正则序列。% P09-ERR-0629
则以 $X$ 的一个包含 $p(z)$ 的开子空间替换之以后，态射
$$
V(f_1, \ldots, f_r) \longrightarrow Y
$$
平坦且局部有限表示。
\end{lemma}

\begin{proof}
令 $Z = V(f_1, \ldots, f_r)$。显然，$Z \to X$ 局部有限表示，
故复合 $Z \to Y$ 局部有限表示；见《空间的态射》引理
\ref{spaces-morphisms-lemma-composition-finite-presentation}。
因此，只须证明 $Z \to Y$ 在 $p(z)$ 的一个邻域中平坦。
设 $k'/k$ 为域扩张。则
$F' = F \times_{\Spec(k)} \Spec(k')$ 在 $F$ 上满且平坦，% P09-ERR-0630
故可找到点 $z' \in |F'|$，它映到 $z$，且局部环映射
$\mathcal{O}_{F, \overline{z}} \to \mathcal{O}_{F', \overline{z}'}$
平坦；见《空间的态射》引理
\ref{spaces-morphisms-lemma-flat-at-point-etale-local-rings}。
因此，$f_1, \ldots, f_r$ 在
$\mathcal{O}_{F', \overline{z}'}$ 中的像也构成正则序列；
见《代数》引理 \ref{algebra-lemma-flat-increases-depth}。
所以在证明过程中可以用扩域替换 $k$。特别地，可以假设
$z \in |F|$ 来自截面 $z : \Spec(k) \to F$，它是结构态射
$F \to \Spec(k)$ 的截面。

\medskip\noindent
选取概形 $V$ 以及满平展态射 $V \to Y$。选取概形 $U$ 以及
满平展态射 $U \to X \times_Y V$。必要时再次扩大 $k$，可以假设
$\Spec(k) \to F \to X$ 通过 $U$ 分解（因为 $U \to X$ 是满的）。
设 $u : \Spec(k) \to U$ 为这样的分解，并用 $v \in V$ 表示 $u$ 的像。
注意，态射
$$
U_v \times_{\Spec(\kappa(v))} \Spec(k) =
U \times_V \Spec(k) \to U \times_Y \Spec(k) \to F
$$
都是平展的（第一个是 $V \to V \times_Y V$ 的基变换，
第二个是 $U \to X$ 的基变换）。此外，依构造，点
$u : \Spec(k) \to U$ 给出最左空间中的一点，它在右端映到 $z$。% P09-ERR-0631
因此，$f_1, \ldots, f_r$ 在下述映射右端的局部环中的像构成正则序列：
$$
\mathcal{O}_{U_v, u}
\longrightarrow
\mathcal{O}_{U_v \times_{\Spec(\kappa(v)} \Spec(k), \overline{u}}
=
\mathcal{O}_{U \times_V \Spec(k), \overline{u}}.
$$
% P09-ERR-0632
但是，该显示箭头平坦（合用《平坦性进阶》引理
\ref{flat-lemma-flat-up-down-henselization} 与《空间的态射》引理
\ref{spaces-morphisms-lemma-flat-at-point-etale-local-rings}），
故由《代数》引理 \ref{algebra-lemma-flat-increases-depth} 可见，
$f_1, \ldots, f_r$ 在 $\mathcal{O}_{U_v, u}$ 中映成正则序列。% P09-ERR-0633
由《态射进阶》引理 \ref{more-morphisms-lemma-slice-given-elements}，
概形态射
$$
V(f_1, \ldots, f_r) \times_X U = V(f_1|_U, \ldots, f_r|_U) \to V
$$
在开邻域 $U'$ 中平坦，该邻域含 $u$。设 $X' \subset X$ 为与
$|U'| \to |X|$ 的像对应的开子空间（见《空间的性质》引理
\ref{spaces-properties-lemma-topology-points} 和
\ref{spaces-properties-lemma-open-subspaces}）。
由此，$V(f_1, \ldots, f_r) \cap X' \to Y$ 平坦
（见《空间的态射》定义 \ref{spaces-morphisms-definition-flat}），
因为有交换图
$$
\xymatrix{
V(f_1, \ldots, f_r) \times_X U' \ar[d]_a \ar[r] & V \ar[d]^b \\
V(f_1, \ldots, f_r) \cap X' \ar[r] & Y
}
$$
，其中 $a, b$ 平展，且 $a$ 满。
\end{proof}






\section{约化纤维}
\label{spaces-more-morphisms-section-reduced}

\noindent
本节对应于《态射进阶》第
\ref{more-morphisms-section-reduced} 节。

\begin{lemma}
\label{spaces-more-morphisms-lemma-geometrically-reduced-fibre}
设 $S$ 为概形。设 $f : X \to Y$ 为 $S$ 上代数空间的态射。
设 $y \in |Y|$。下列条件等价：% P09-ERR-0634
\begin{enumerate}
\item 对于某个态射 $\Spec(k) \to Y$，它位于 $y$ 的等价类中，
代数空间 $X_k$ 在 $k$ 上几何约化；
\item 对于每个态射 $\Spec(k) \to Y$，它位于 $y$ 的等价类中，
代数空间 $X_k$ 在 $k$ 上几何约化；
\item 对于每个态射 $\Spec(k) \to Y$，它位于 $y$ 的等价类中，
代数空间 $X_k$ 约化。
\end{enumerate}
\end{lemma}

\begin{proof}
这立即由《域上的空间》引理
\ref{spaces-over-fields-lemma-geometrically-reduced-upstairs}
以及《空间的性质》第 \ref{spaces-properties-section-points} 节中给出的、
定义 $|X|$ 的等价关系得出。% P09-ERR-0635
\end{proof}

\begin{definition}
\label{spaces-more-morphisms-definition-geometrically-reduced-fibre}
设 $S$ 为概形。设 $f : X \to Y$ 为 $S$ 上代数空间的态射。
设 $y \in |Y|$。若引理
\ref{spaces-more-morphisms-lemma-geometrically-reduced-fibre} 的等价条件成立，
则称 {\it $f : X \to Y$ 在 $y$ 处的纤维几何约化}。
\end{definition}

\noindent
下面给出必需的引理。

\begin{lemma}
\label{spaces-more-morphisms-lemma-base-change-fibres-geometrically-reduced}
设 $S$ 为概形。设 $f : X \to Y$ 和 $g : Y' \to Y$ 为 $S$ 上
代数空间的态射。用 $f' : X' \to Y'$ 表示 $f$ 沿 $g$ 的基变换。则
\begin{align*}
\{y' \in |Y'| :
\text{态射 }f' : X' \to Y'\text{ 在点 }y'
\text{ 的纤维几何约化}\} \\
= g^{-1}(\{y \in |Y| :
\text{态射 }f : X \to Y\text{ 在点 }y
\text{ 的纤维几何约化}\}).
\end{align*}
\end{lemma}

\begin{proof}
对 $y' \in |Y'|$，选取态射 $\Spec(k) \to Y'$，它位于 $y'$
的等价类中。则 $g(y')$ 由复合
$\Spec(k) \to Y' \to Y$ 表示。于是
$X' \times_{Y'} \Spec(k) = X \times_Y \Spec(k)$，
结论由定义得出。
\end{proof}

\begin{lemma}
\label{spaces-more-morphisms-lemma-geometrically-reduced-constructible}
设 $S$ 为概形。设 $f : X \to Y$ 为 $S$ 上代数空间的拟紧且
局部有限表示的态射。则集合
$$
E = \{y \in |Y| : \text{态射 }f : X \to Y\text{ 在点 }y
\text{ 的纤维几何约化}\}
$$
平展局部可构造。
\end{lemma}

\begin{proof}
选取仿射概形 $V$ 以及平展态射 $V  \to Y$。
陈述的含义是 $E$ 在 $|V|$ 中的逆像可构造。由引理
\ref{spaces-more-morphisms-lemma-base-change-fibres-geometrically-reduced}，
可以替换 $Y$，以 $V$ 代之，即可以假设 $Y$ 为仿射概形。
于是 $X$ 拟紧。选取仿射概形 $U$ 以及满平展态射 $U \to X$。
对于态射 $\Spec(k) \to Y$，纤维之间的态射
$U_k \to X_k$ 满且平展。因此，$U_k$ 在 $k$ 上几何约化，
当且仅当 $X_k$ 在 $k$ 上几何约化；见《域上的空间》引理
\ref{spaces-over-fields-lemma-geometrically-reduced-etale-local}。
所以，集合 $E$（针对 $X \to Y$）与集合 $E$（针对 $U \to Y$）相同。
由此可见，引理由概形情形得出；见《态射进阶》引理
\ref{more-morphisms-lemma-geometrically-reduced-constructible}。
\end{proof}

\begin{lemma}
\label{spaces-more-morphisms-lemma-proper-flat-over-dvr-reduced-fibre}
设 $X$ 为离散赋值环 $R$ 上的代数空间，其结构态射
$X \to \Spec(R)$ 固有且平坦。若特殊纤维约化，则
$X$ 和一般纤维 $X_\eta$ 都约化。
\end{lemma}

\begin{proof}
选取平展态射 $U \to X$，其中 $U$ 为仿射概形。
则 $U$ 在 $R$ 上有限型。设 $u \in U$ 位于特殊纤维中。
局部环 $A = \mathcal{O}_{U, u}$ 在 $R$ 上本质有限型，因而是诺特的。
设 $\pi \in R$ 为一致化元。由于 $X$ 在 $R$ 上平坦，
可见 $\pi \in \mathfrak m_A$ 在 $A$ 上为非零因子；又由于 $X$
的特殊纤维约化，$A/\pi A$ 约化。
若 $a \in A$ 且 $a \not = 0$，则存在 $n \geq 0$ 以及元素
$a' \in A$，使得 $a = \pi^n a'$ 且 $a' \not \in \pi A$。
这由 Krull 交定理得出（《代数》引理
\ref{algebra-lemma-intersect-powers-ideal-module-zero}）。
若 $a$ 幂零，则 $a'$ 也幂零，因为 $\pi$ 是非零因子。
但 $a'$ 在约化环 $A/\pi A$ 中的像非零，这是不可能的。
故 $A$ 约化。因此，存在 $u$ 的一个约化开邻域，它位于 $U$ 中
（略去一个小细节；使用 $U$ 为诺特的这一事实）。% P09-ERR-0636
于是可以找到平展态射 $U \to X$，其中 $U$ 是约化概形，
使得 $X$ 的特殊纤维的每一点都在其像中。由于 $X$ 在 $R$ 上固有，
可知 $U \to X$ 满。因此 $X$ 约化。由于 $U \to \Spec(R)$ 的
一般纤维也约化（在仿射片上，它通过取局部化计算），
一般纤维也具有同样的性质。
\end{proof}

\begin{lemma}
\label{spaces-more-morphisms-lemma-geometrically-reduced-open}
设 $S$ 为概形。设 $f : X \to Y$ 为 $S$ 上代数空间的态射。
若 $f$ 平坦、固有且有限表示，则集合
$$
E = \{y \in |Y| : \text{态射 }f : X \to Y\text{ 在点 }y
\text{ 的纤维几何约化}\}
$$
在 $|Y|$ 中开。
\end{lemma}

\begin{proof}
由引理 \ref{spaces-more-morphisms-lemma-base-change-fibres-geometrically-reduced}，
$E$ 的形成与基变换可交换。为检验 $|Y|$ 的子集为开集，
可以用一个平展覆盖的各成员替换 $Y$。故可假设 $Y$ 仿射。
于是 $Y$ 是有限型仿射 $\mathbf{Z}$-概形的余滤极限。
因此，可以假设 $X \to Y$ 是 $X_0 \to Y_0$ 的基变换，其中
$Y_0$ 是有限型 $\mathbf{Z}$-代数的谱，而 $X_0 \to Y_0$ 平坦且固有。
见《空间的极限》引理
\ref{spaces-limits-lemma-descend-finite-presentation}、
\ref{spaces-limits-lemma-descend-flat} 和
\ref{spaces-limits-lemma-eventually-proper}。由于 $E$ 的形成与基变换
可交换（见上），可以假设基底为诺特的。

\medskip\noindent
假设 $Y$ 诺特。由引理
\ref{spaces-more-morphisms-lemma-geometrically-reduced-constructible}，该集合可构造。
因此，只须证明该集合在一般化下稳定（《拓扑》引理
\ref{topology-lemma-characterize-closed-Noetherian}）。由《性质》引理
\ref{properties-lemma-locally-Noetherian-specialization-dvr}，
约化到如下情形：$Y = \Spec(R)$，$R$ 为离散赋值环，且闭纤维
$X_y$ 几何约化。要证明：一般纤维 $X_\eta$ 几何约化。

\medskip\noindent
若不然，则存在有限扩张 $L$，它扩张 $R$ 的分式域，且使 $X_L$ 非约化；
见《域上的空间》引理
\ref{spaces-over-fields-lemma-perfect-reduced}（特征零）和
\ref{spaces-over-fields-lemma-geometrically-reduced-positive-characteristic}
（正特征）。存在离散赋值环 $R' \subset L$，其分式域为 $L$ 且支配
$R$；见《代数》引理 \ref{algebra-lemma-integral-closure-Dedekind}。
替换 $R$，以 $R'$ 代之以后，约化到引理
\ref{spaces-more-morphisms-lemma-proper-flat-over-dvr-reduced-fibre}。
\end{proof}







\section{纤维的连通分支}
\label{spaces-more-morphisms-section-connected}

\noindent
本节对应于《态射进阶》第
\ref{more-morphisms-section-connected} 节。

\begin{lemma}
\label{spaces-more-morphisms-lemma-base-change-fibres-nr-geometrically-connected-components}
设 $S$ 为概形。
设 $f : X \to Y$ 为 $S$ 上代数空间的态射。设
$$
n_{X/Y} : |Y| \to \{0, 1, 2, 3, \ldots, \infty\}
$$
为如下函数：它把 $y \in Y$ 映到 $X_k$ 的连通分支数，% P09-ERR-0637
其中 $\Spec(k) \to Y$ 位于 $y$ 的等价类中，且 $k$ 代数闭。
该函数良定义；若 $g : Y' \to Y$ 为态射，则
$$
n_{X'/Y'} = n_{X/Y} \circ g
% P09-ERR-0645
$$
，其中 $X' \to Y'$ 是 $f$ 的基变换。
\end{lemma}

\begin{proof}
假设 $y' \in Y'$ 的像为 $y \in Y$。% P09-ERR-0638
设 $\Spec(k') \to Y'$ 位于 $y'$ 的等价类中，且 $k'$ 代数闭。
则可以选取交换图
$$
\xymatrix{
\Spec(K) \ar[r] \ar[rd] &
\Spec(k') \ar[r] & Y' \ar[d] \\
& \Spec(k) \ar[r] & Y
}
$$
，其中 $K$ 为代数闭域。结论成立，因为概形态射
$$
\xymatrix{
X'_{k'} & (X'_{k'})_K = (X_k)_K \ar[l] \ar[r] & X_k
}
$$
在连通分支之间诱导双射；见《域上的空间》引理
\ref{spaces-over-fields-lemma-separably-closed-field-connected-components}。
为用此证明函数良定义，取 $Y' = Y$。
\end{proof}







\section{纤维的维数}
\label{spaces-more-morphisms-section-dimension}

\noindent
本节对应于《态射进阶》第
\ref{more-morphisms-section-dimension} 节。

\begin{lemma}
\label{spaces-more-morphisms-lemma-dimension-fibre}
设 $S$ 为概形。设 $f : X \to Y$ 为 $S$ 上代数空间的有限型态射。
设 $y \in |Y|$。下列量相同：% P09-ERR-0639
\begin{enumerate}
\item 令 $d = -\infty$（如果 $y$ 不在 $|f|$ 的像中）；否则，$d$ 是满足
如下条件的最小整数：$f$ 在所有下述点处的相对维数均 $\leq d$：
$x \in |X|$ 且其像为 $y$；
\item 代数空间 $X_k = \Spec(k) \times_Y X$ 的维数，其中
$\Spec(k) \to Y$ 是定义 $y$ 的等价类中的任意态射。
\end{enumerate}
\end{lemma}

\begin{proof}
为理解此陈述，须参阅《空间的态射》定义
\ref{spaces-morphisms-definition-dimension-fibre}、
《空间的性质》定义 \ref{spaces-properties-definition-dimension}，以及
《空间的性质》定义 \ref{spaces-properties-definition-dimension-at-point}。
对于一个固定态射 $\Spec(k) \to Y$，我们将证明 (1) 与 (2) 中的数相等。
选取平展态射 $V \to Y$，其中 $V$ 为仿射概形；并选取点
$v \in V$，其像为 $y$。由于 $V \times_Y \Spec(k) \to \Spec(k)$ 满且平展
（由《空间的性质》引理 \ref{spaces-properties-lemma-points-cartesian}），
可以找到有限可分扩张 $k'/k$（由《态射》引理
\ref{morphisms-lemma-etale-over-field}）以及交换图
$$
\xymatrix{
\Spec(k') \ar[r] \ar[d] & V \ar[d] \\
\Spec(k) \ar[r] & Y
}
$$
可以把 $X \to Y$ 替换为 $V \times_Y X \to V$，并把
$X_k$ 替换为 $X_{k'} = \Spec(k') \times_V (V \times_Y X)$，
因为由《空间的性质》引理
\ref{spaces-properties-lemma-dimension-decent-invariant-under-etale}
和《空间的态射》引理
\ref{spaces-morphisms-lemma-dimension-fibre-after-base-change}，
这不改变所论维数。因此，可以假设 $Y$ 为仿射概形。
在此情形下，可以假设 $k = \kappa(y)$，因为
$X_{\kappa(y)}$ 与 $X_k$ 的维数相同；这由前述《空间的态射》引理
\ref{spaces-morphisms-lemma-dimension-fibre-after-base-change}
以及如下事实得出：对于代数空间 $Z$ 及其底域 $K$，$Z$ 在点
$z \in |Z|$ 处的相对维数依定义等于 $\dim_z(Z)$。
假设 $Y$ 仿射且 $k = \kappa(y)$。则 $X$ 拟紧；% P09-ERR-0640
可以选取仿射概形 $U$ 以及满平展态射 $U \to X$。% P09-ERR-0641
于是 $\dim(X_k) = \dim(U_k) = \max \dim_u(U_k)$，% P09-ERR-0642
它依定义等于 (1) 中给出的数。
\end{proof}

\begin{lemma}
\label{spaces-more-morphisms-lemma-base-change-dimension-fibres}
设 $S$ 为概形。设 $f : X \to Y$ 为 $S$ 上代数空间的有限型态射。设
$$
n_{X/Y} : |Y| \to \{-\infty, 0, 1, 2, 3, \ldots\}
$$
为把 $y \in |Y|$ 映到引理 \ref{spaces-more-morphisms-lemma-dimension-fibre} 中所论整数的函数。
若 $g : Y' \to Y$ 为态射，则
$$
n_{X'/Y'} = n_{X/Y} \circ |g|
$$
，其中 $X' \to Y'$ 是 $f$ 的基变换。
\end{lemma}

\begin{proof}
这立即由引理 \ref{spaces-more-morphisms-lemma-dimension-fibre} 得出。
\end{proof}

\begin{lemma}
\label{spaces-more-morphisms-lemma-dimension-fibres-flat}
设 $S$ 为概形。
设 $f : X \to Y$ 为 $S$ 上代数空间的平坦有限表示态射。
设 $n_{X/Y}$ 为引理 \ref{spaces-more-morphisms-lemma-base-change-dimension-fibres}
中引入的、在 $Y$ 上给出 $f$ 的纤维维数的函数。
则 $n_{X/Y}$ 下半连续。
\end{lemma}

\begin{proof}
设 $V \to Y$ 为满平展态射，其中 $V$ 是概形。
若能证明复合 $n_{X/Y} \circ |g|$ 下半连续，% P09-ERR-0643
则因为 $|g|$ 开，引理成立。故可以假设 $Y$ 为概形。
在局部上工作，可以假设 $V$ 为仿射概形。% P09-ERR-0644
于是可以选取仿射概形 $U$ 以及满平展态射 $U \to X$。
于是 $n_{X/Y} = n_{U/Y}$。故可以假设 $X$ 和 $Y$ 都是概形。
在此情形下，引理由《态射进阶》引理
\ref{more-morphisms-lemma-dimension-fibres-flat} 得出。
\end{proof}

\begin{lemma}
\label{spaces-more-morphisms-lemma-dimension-fibres-proper}
设 $S$ 为概形。
设 $f : X \to Y$ 为 $S$ 上代数空间的固有态射。
设 $n_{X/Y}$ 为引理 \ref{spaces-more-morphisms-lemma-base-change-dimension-fibres}
中引入的、在 $Y$ 上给出 $f$ 的纤维维数的函数。
则 $n_{X/Y}$ 上半连续。
\end{lemma}

\begin{proof}
令 $Z_d = \{x \in |X| :
\text{态射 }f\text{ 在点 }x\text{ 的纤维维数 }> d\}$。
由《空间的态射》引理
\ref{spaces-morphisms-lemma-openness-bounded-dimension-fibres}，
$Z_d$ 是 $|X|$ 的闭子集。由于 $f$ 固有，% P09-ERR-0646
$f(Z_d)$ 在 $|Y|$ 中闭。
由于 $y \in f(Z_d) \Leftrightarrow n_{X/Y}(y) > d$，引理成立。
\end{proof}

\begin{lemma}
\label{spaces-more-morphisms-lemma-dimension-fibres-proper-flat}
设 $S$ 为概形。设 $f : X \to Y$ 为 $S$ 上代数空间的固有、平坦且
有限表示的态射。
设 $n_{X/Y}$ 为引理 \ref{spaces-more-morphisms-lemma-base-change-dimension-fibres}
中引入的、在 $Y$ 上给出 $f$ 的纤维维数的函数。
则 $n_{X/Y}$ 局部常值。
\end{lemma}

\begin{proof}
这是引理 \ref{spaces-more-morphisms-lemma-dimension-fibres-flat} 和
\ref{spaces-more-morphisms-lemma-dimension-fibres-proper} 的直接推论。
\end{proof}





\section{catenary 代数空间}
\label{spaces-more-morphisms-section-catenary}

\noindent
本节继续《良态空间》第 \ref{decent-spaces-section-catenary} 节中开始的讨论。
下一个引理的证明将使用如下引理。

\begin{lemma}
\label{spaces-more-morphisms-lemma-construct-glueing}
设 $S$ 为概形。设 $f : X \to Y$ 为 $S$ 上代数空间的整态射。
设 $y \in |Y|$ 为可由闭浸入 $y : \Spec(k) \to Y$ 表示的点。
则存在分解 $X \to X' \to Y$，其复合为 $f$，并且
\begin{enumerate}
\item $X' \to Y$ 为整态射；
\item $X \to X'$ 在 $X' \setminus X'_y$ 上为同构；
\item $X'_y$ 有唯一点 $x'$，且 $\kappa(x') = k$。
\end{enumerate}
此外，若 $f$ 有限且 $Y$ 局部 Noetherian，则 $X' \to Y$ 有限。
\end{lemma}

\begin{proof}
由《空间的态射》引理 \ref{spaces-morphisms-lemma-pushforward}，层
$f_*\mathcal{O}_X$、$(X_y \to Y)_*\mathcal{O}_{X_y}$ 和
$y_*\mathcal{O}_{\Spec(k)}$ 都是拟凝聚 $\mathcal{O}_Y$-代数层。考虑映射
$$
f_*\mathcal{O}_Y \longrightarrow
% P09-ERR-0647
(X_y \to Y)_*\mathcal{O}_{X_y} \longleftarrow
y_*\mathcal{O}_{\Spec(k)}
$$
其纤维积是拟凝聚 $\mathcal{O}_Y$-代数层 $\mathcal{A}'$。
可以把 $X' \to Y$ 定义为 $\mathcal{A}'$ 在 $Y$ 上的相对谱；见
《态射》引理 \ref{morphisms-lemma-affine-equivalence-algebras}。
此构造与任意基变换可交换。特别地，显然在开子空间
$|Y| \setminus \{y\}$ 上态射 $X \to X'$ 是同构，而在
$|Y| \setminus \{y\}$ 上态射 $X' \to Y$ 是整态射。
尚须在 $y$ 的一个小邻域内证明这些陈述。
选取仿射概形 $V = \Spec(R)$ 以及平展态射
$\varphi : V \to Y$，使得 $y$ 位于 $\varphi$ 的像中。
则 $V_y$ 是 $V$ 的闭子概形，并且在 $k$ 上平展，因而由有限多个点组成，
且每个点的剩余域都是 $k$ 的可分扩张
（见《良态空间》评注 \ref{decent-spaces-remark-recall}）。
缩小 $V$ 后，可以假设存在唯一的闭点
$v = \Spec(l) \to V$ 映到 $y$，其中 $l/k$ 有限可分。
可以写成 $V \times_Y X = \Spec(C)$，其中 $R \to C$ 是整环同态。
上述与基变换的相容性给出
$U \times_X Y' = \Spec(C')$，其中 % P09-ERR-0648
$$
C' = C \times_{C \otimes_R l} l
$$
由于 $R \to l$ 满，$C \to C \otimes_R l$ 也满，因而这是
《代数进阶》情形 \ref{more-algebra-situation-module-over-fibre-product}
中研究的那类纤维积
（其中 $A, A', B, B'$ 分别对应于 $C \otimes_R l, C, l, C'$）。
注意到 $C'$ 是一个 $R$-子代数，包含于 $C$ 中，因而在 $R$ 上整；
这证明了 (1)。
最后，《代数进阶》引理
\ref{more-algebra-lemma-points-of-fibre-product} 表明
$V \times_X Y' = \Spec(C')$ % P09-ERR-0649
有唯一的点 $y''$ 位于 $v$ 上方，其剩余域为 $l$
（它对应于显然的满映射 $C' \to l$）。
因此 $X_y \times_{\Spec(k)} \Spec(l)$ 有唯一的点，% P09-ERR-0650
其剩余域为 $l$。由于 $l/k$ 有限可分，这蕴含 $X'_y$
有唯一的点，且其剩余域为 $k$，即 (3) 成立。

\medskip\noindent
为证明最后一个陈述，注意到若 $Y$ 局部 Noetherian，则 $R$ 是
Noetherian 环；若 $f$ 有限，则 $R \to C$ 有限。于是由《代数进阶》引理
\ref{more-algebra-lemma-fibre-product-finite-type}，$C'$ 是有限型
$R$-代数。这证明了 $X' \to Y$ 有限。
\end{proof}

\begin{lemma}
\label{spaces-more-morphisms-lemma-universally-catenary-dimension-function}
设 $S$ 为概形。设 $B$ 为 $S$ 上的代数空间。设
$\delta : |B| \to \mathbf{Z}$ 为函数。假设 $B$ 良态、局部 Noetherian
且普遍 catenary，并且 $\delta$ 是维数函数。若 $X$ 是 $B$ 上的良态
代数空间，其结构态射 $f : X \to B$ 局部有限型，则按下述规则定义
$\delta_X : |X| \to \mathbf{Z}$：
$$
\delta_X(x) = \delta(f(x)) + \text{超越次数 }x/f(x)
% P09-ERR-0651
$$
（《空间的态射》定义
\ref{spaces-morphisms-definition-dimension-fibre}）。
则 $\delta_X$ 是维数函数。
\end{lemma}

\begin{proof}
问题在 $B$ 上是局部的，故可以假设 $B$ 拟紧。
由《良态空间》引理
\ref{decent-spaces-lemma-locally-Noetherian-decent-quasi-separated}，
$B$ 拟分离。由《空间的极限》命题
\ref{spaces-limits-proposition-there-is-a-scheme-finite-over}，
可以选取有限满态射 $\pi : Y \to X$，其中 $Y$ 是概形。% P09-ERR-0652
断言：$\delta_Y$ 是维数函数。

\medskip\noindent
该断言蕴含引理。对引理中的 $X \to B$，令
$Z = Y \times_B X$，并记投影为 $p : Z \to Y$ 和 $q : Z \to X$。
则
$$
\delta_Z(z) = \delta_Y(p(z)) + \text{超越次数 }z/p(z)
% P09-ERR-0651
$$
且 $\delta_Z(z) = \delta_X(q(z))$。这由《空间的态射》引理
\ref{spaces-morphisms-lemma-dimension-fibre-at-a-point-additive}
以及这些超越次数对于有限态射为零这一事实得出。由《良态空间》引理
\ref{decent-spaces-lemma-scheme-with-dimension-function} 和该断言，
可知 $\delta_Z$ 是维数函数。再由《良态空间》引理
\ref{decent-spaces-lemma-check-dimension-function-finite-cover}，
可知 $\delta_X$ 是维数函数。

\medskip\noindent
证明该断言。考虑特化 $y \leadsto y'$，其中 $y \not = y'$，
二者都是 Noetherian 概形 $Y$ 的点。于是
$\delta_Y(y) > \delta_Y(y')$，因为 $Y$ 的纤维中的点之间不存在特化
（见《良态空间》引理
\ref{decent-spaces-lemma-conditions-on-fibre-and-qf}）。
把它用于一条特化链，得到
$$
\delta_Y(y) - \delta_Y(y') \geq
\text{余维}(\overline{\{y'\}}, \overline{\{y\}})
$$
我们的任务是证明等号成立。由《性质》引理
\ref{properties-lemma-locally-Noetherian-closed-point}，
可以选取一个特化 $y' \leadsto y_0$。只须证明
$\delta_Y(y) - \delta_Y(y_0) =
\text{余维}(\overline{\{y_0\}}, \overline{\{y\}})$，
因为这将同时蕴含 $y \leadsto y'$ 和 $y' \leadsto y_0$ 的等式。

\medskip\noindent
选取极大特化链
$y = y_c \leadsto y_{c - 1} \leadsto \ldots \leadsto y_0$，它位于 $Y$ 中。
令 $b = \pi(y)$ 且 $b_0 = \pi(y_0)$。
选取极大特化链
$b = b_e \leadsto b_{e - 1} \leadsto \ldots \leadsto b_0$，它位于 $|B|$ 中。
须证明 $e = c$。由于 $\pi$ 闭
（《空间的态射》引理 \ref{spaces-morphisms-lemma-finite-proper}），
可以找到特化序列
$y = y'_e \leadsto y'_{e - 1} \leadsto \ldots \leadsto y'_0$，
它映到
$b = b_e \leadsto b_{e - 1} \leadsto \ldots \leadsto b_0$。
注意 $y'_e \leadsto y'_{e - 1} \leadsto \ldots \leadsto y'_0$
也是极大链。若 $y_0 = y'_0$，则因 $Y$ catenary，得到所需的
$e = c$。下一段用一个小技巧约化到此情形，并以同样方式得出结论。

\medskip\noindent
由于 $\pi$ 闭，$b_0$ 是 $|B|$ 的闭点。
由《良态空间》引理 \ref{decent-spaces-lemma-decent-space-closed-point}，
可以把 $b_0$ 表示为闭浸入 $b_0 : \Spec(k) \to B$。
由引理 \ref{spaces-more-morphisms-lemma-construct-glueing}，可以找到分解
$$
Y \to Y' \to X
% P09-ERR-0653
$$
，其中 $\pi' : Y' \to X$ 有限，而 $Y \to Y'$ 是把 % P09-ERR-0653
$y_0$ 和 $y'_0$ 映到同一点、且在 $b_0$ 的逆像之外为同构的态射。
% P09-ERR-0654
（当然，$Y'$ 不会是概形，但这不影响下面的论证。）
显然，极大特化链
$y_c \leadsto y_{c - 1} \leadsto \ldots \leadsto y_0$ 和
$y'_e \leadsto y'_{e - 1} \leadsto \ldots \leadsto y'_0$
映为 $Y'$ 中有相同起点和终点的极大特化链。
由于 $B$ 普遍 catenary，$|Y'|$ catenary，结论同前。
\end{proof}





\section{态射的平展局部化}
\label{spaces-more-morphisms-section-etale-localization}

\noindent
本节对应于《态射进阶》第
\ref{more-morphisms-section-etale-localization} 节。

\begin{lemma}
\label{spaces-more-morphisms-lemma-etale-splits-off-quasi-finite-part}
设 $S$ 为概形。
设 $f : X \to Y$ 为 $S$ 上代数空间的态射。
设 $y \in |Y|$。设 $x_1, \ldots, x_n \in |X|$ 均映到 $y$。
假设
\begin{enumerate}
\item $f$ 局部有限型；
\item $f$ 分离；
\item $f$ 在 $x_1, \ldots, x_n$ 处拟有限；并且
\item $f$ 拟紧，或 $Y$ 良态。
\end{enumerate}
则存在带基点代数空间的平展态射 $(U, u) \to (Y, y)$ 以及分解
$$
U \times_Y X = W \amalg V
$$
为既开又闭子空间，使得态射 $V \to U$ 有限；
$|V| \to |U|$ 在 $u$ 上方纤维的每个点都映到某个 $x_i$；
而 $|W| \to |U|$ 在 $u$ 上方的纤维不含映到任何 $x_i$ 的点。
\end{lemma}

\begin{proof}
设 $(U, u) \to (Y, y)$ 为代数空间的平展态射，并考虑满足如下条件的
$w \in |U \times_Y X|$ 之集合：它映到 $u \in |U|$ 以及某个
$x_i \in |X|$。由《良态空间》引理
\ref{decent-spaces-lemma-qf-and-qc-finite-fibre}（若 $f$ 有限型）
或《良态空间》引理
\ref{decent-spaces-lemma-decent-finite-fibre}（若 $Y$ 良态），
此集合有限。因此可以把 $f$ 替换为基变换 $U \times_Y X \to U$，
并把 $x_1, \ldots, x_n$ 替换为这些 $w$ 之集合。
特别地，可以并且确实假设 $Y$ 为仿射概形；于是 $X$ 是分离代数空间。

\medskip\noindent
选取仿射概形 $Z$ 以及平展态射 $Z \to X$，使
$x_1, \ldots, x_n$ 位于 $|Z| \to |X|$ 的像中。
$|Z| \to |X|$ 的纤维有限；见《空间的性质》引理
\ref{spaces-properties-lemma-finite-fibres-presentation}
（或《良态空间》第
\ref{decent-spaces-section-reasonable-decent} 节中更一般的讨论）。
令 $\{z_1, \ldots, z_m\} \subset |Z|$ 为
$\{x_1, \ldots, x_n\}$ 的逆像。由《态射进阶》引理
\ref{more-morphisms-lemma-etale-splits-off-quasi-finite-part-technical}，
存在平展态射 $(U, u) \to (Y, y)$，使得
$U \times_Y Z = Z_1 \amalg Z_2$，其中 $Z_1 \to U$ 有限，且
$(Z_1)_y = \{z_1, \ldots, z_m\}$。% P09-ERR-0655
可以假设 $U$ 仿射，因而 $Z_1$ 也仿射。

\medskip\noindent
由于 $f$ 分离，像 $V$（即 $Z_1 \to X$ 的像）既开又闭
（《空间的态射》引理
\ref{spaces-morphisms-lemma-universally-closed-permanence}）。% P09-ERR-0656
令 $W = X \setminus V$，便得到引理所述的分解。% P09-ERR-0656
为完成证明，须证明 $V \to U$ 有限。
由于 $Z_1 \to V$ 满且平展，$V$ 是 $Z_1$ 关于平展等价关系
$R = Z_1 \times_V Z_1$ 的商；见《空间》引理
\ref{spaces-lemma-space-presentation}。
由于 $f$ 分离，$V \to U$ 分离，且 $R$ 在
$Z_1 \times_U Z_1$ 中闭。由于 $Z_1 \to U$ 有限，投影
$s, t : R \to Z_1$ 有限。因此，由《群胚》命题
\ref{groupoids-proposition-finite-flat-equivalence}，$V$ 是仿射概形。
由《态射》引理 \ref{morphisms-lemma-image-proper-is-proper}，
可知 $V \to U$ 固有；再由《态射》引理
\ref{morphisms-lemma-finite-proper}，可知 $V \to U$ 有限，
从而完成证明。
\end{proof}

\begin{lemma}
\label{spaces-more-morphisms-lemma-etale-splits-off-just-one-quasi-finite-part}
设 $S$ 为概形。
设 $f : X \to Y$ 为 $S$ 上代数空间的态射。
设 $x \in |X|$，其像为 $y \in |Y|$。假设
\begin{enumerate}
\item $f$ 局部有限型；
\item $f$ 分离；并且
\item $f$ 在 $x$ 处拟有限。
\end{enumerate}
则存在带基点代数空间的平展态射 $(U, u) \to (Y, y)$ 以及分解
$$
U \times_Y X = W \amalg V
$$
为既开又闭子空间，使得态射 $V \to U$ 有限，并且存在点
$v \in |V|$，它映到 $x$（在 $|X|$ 中）以及 $u$（在 $|U|$ 中）。
\end{lemma}

\begin{proof}
选取概形 $U$、点 $u \in U$ 以及平展态射
$U \to Y$，该态射把 $u$ 映到 $y$。存在点
$x' \in |U \times_Y X|$，它映到 $x$（在 $|X|$ 中）以及
$u$（在 $|U|$ 中）（《空间的性质》引理
\ref{spaces-properties-lemma-points-cartesian}）。
为完成证明，把引理 \ref{spaces-more-morphisms-lemma-etale-splits-off-quasi-finite-part}
应用于态射 $U \times_Y X \to U$ 和点 $x'$。
它适用，因为 $U$ 是概形，因而 $u$ 来自单态射
$\Spec(\kappa(u)) \to U$。
\end{proof}





\section{Zariski 主定理}
\label{spaces-more-morphisms-section-structure-quasi-finite}

\noindent
本节应用上一节的结果，证明代数空间态射的 Zariski 主定理。
本节对应于《态射进阶》第
\ref{more-morphisms-section-application-etale-neighbourhoods} 节。

\begin{lemma}
\label{spaces-more-morphisms-lemma-finite-type-separated}
设 $S$ 为概形。设 $f : X \to Y$ 为 $S$ 上代数空间的有限型分离态射。
设 $Y'$ 为 $Y$ 在 $X$ 中的正规化。图示为：
$$
\xymatrix{
X \ar[rd]_f \ar[rr]_{f'} & & Y' \ar[ld]^\nu \\
& Y &
}
$$
则存在开子空间 $U' \subset Y'$，使得
\begin{enumerate}
\item $(f')^{-1}(U') \to U'$ 为同构；并且
\item $(f')^{-1}(U') \subset X$ 恰为 $f$ 拟有限的点之集合。
\end{enumerate}
\end{lemma}

\begin{proof}
由《空间的态射》引理
\ref{spaces-morphisms-lemma-locally-finite-type-quasi-finite-part}，
存在开子空间 $U \subset X$，它对应于 $|X|$ 中 $f$ 拟有限的点。
须证明
\begin{enumerate}
\item[(a)] $|U| \to |Y'|$ 的像为 $|U'|$，其中 $U'$ 是
$Y'$ 的某个开子空间；
\item[(b)] $U = f^{-1}(U')$；并且 % P09-ERR-0657
\item[(c)] $U \to U'$ 为同构。
\end{enumerate}
由于 $U$ 的形成与任意基变换可交换
（《空间的态射》引理
\ref{spaces-morphisms-lemma-locally-finite-type-quasi-finite-part}），
由于正规化 $Y'$ 的形成与光滑基变换可交换
（引理 \ref{spaces-more-morphisms-lemma-normalization-smooth-localization}），
由于平展态射是开态射，并且由于“为同构”在底上是 fpqc 局部的
（《空间上的下降》引理
\ref{spaces-descent-lemma-descending-property-isomorphism}），
只须在 $Y$ 上平展局部地证明 (a)、(b)、(c)
（省略若干细节）。% P09-ERR-0658
因此，可以假设 $Y$ 为仿射概形。这蕴含 $Y'$ 也是（仿射）概形。

\medskip\noindent
设 $x \in |U|$。断言：存在开邻域 $f'(x) \in V \subset Y'$，
使得 $(f')^{-1}V \to V$ 为同构。先证明该断言蕴含引理。
事实上，此时 $(f')^{-1}V \cong V$ 是概形（因为它是 $Y'$ 的开集），
在 $Y$ 上局部有限型（因为它是 $X$ 的开子空间）；而对于 $v \in V$，
剩余域扩张 $\kappa(v)/\kappa(\nu(v))$ 是代数的（因为
$V \subset Y'$ 且 $Y'$ 在 $Y$ 上整）。因此 $V \to Y$ 的纤维离散
（《态射》引理
\ref{morphisms-lemma-algebraic-residue-field-extension-closed-point-fibre}），
且 $(f')^{-1}V \to Y$ 局部拟有限
（《态射》引理 \ref{morphisms-lemma-locally-quasi-finite-fibres}）。
这蕴含 $(f')^{-1}V \subset U$ 且 $V \subset U'$。% P09-ERR-0659
由于 $x$ 任意，可知 (a)、(b)、(c) 成立。

\medskip\noindent
令 $y = f(x) \in |Y|$。设 $(T, t) \to (Y, y)$ 为带基点概形的
平展态射。用下标 $ {}_T$ 表示到 $T$ 的基变换。
设 $z \in X_T$ 为纤维 $X_t$ 中位于 $x$ 上方的点。% P09-ERR-0660
注意 $U_T \subset X_T$ 是 $f_T$ 拟有限的点之集合；见
《空间的态射》引理
\ref{spaces-morphisms-lemma-locally-finite-type-quasi-finite-part}。
注意
$$
X_T \xrightarrow{f'_T} Y'_T \xrightarrow{\nu_T} T
$$
是 $T$ 在 $X_T$ 中的正规化；见引理
\ref{spaces-more-morphisms-lemma-normalization-smooth-localization}。
假设断言对 $z \in U_T \subset X_T \to Y'_T \to T$ 成立，% P09-ERR-0660
即假设可以找到开邻域 $f'_T(z) \in V' \subset Y'_T$，使得
$(f'_T)^{-1}V' \to V'$ 为同构。态射 $Y'_T \to Y'$ 平展，
故像 $V \subset Y'$（即 $V'$ 的像）是开集。注意 $f'(x) \in V$，
因为 $f'_T(z) \in V'$。还要注意
$$
\xymatrix{
(f'_T)^{-1}V' \ar[r] \ar[d] & (f')^{-1}(V) \ar[d] \\
V' \ar[r] & V
}
$$
是纤维方块（因为 $Y'_T \times_{Y'} X = X_T$）。
由于左侧竖直箭头是同构，且 $\{V' \to V\}$ 是平展覆盖，
由《空间上的下降》引理
\ref{spaces-descent-lemma-descending-property-isomorphism}，
右侧竖直箭头也是同构。换言之，断言对
$x \in U \subset X \to Y' \to Y$ 成立。% P09-ERR-0660

\medskip\noindent
由上一段的结果，为证明断言对 $x \in |U|$ 成立，可以把 $Y$
替换为平展邻域 $T$，它是 $y = f(x)$ 的邻域；并把 $x$ 替换为
位于 $x$ 上方、属于 $T \times_Y X$ 的任意点。% P09-ERR-0660
因此，可以假设存在分解
$$
X = V \amalg W
$$
为既开又闭子空间，其中 $V \to Y$ 有限且 $x \in V$；
见引理 \ref{spaces-more-morphisms-lemma-etale-splits-off-quasi-finite-part}。
由于 $X$ 是 $V$ 与 $W$ 在 $Y$ 上的不交并，并且 $V \to Y$ 有限，
可知 $Y$ 在 $X$ 中的正规化是态射
$$
X = V \amalg W \longrightarrow V \amalg W' \longrightarrow S
% P09-ERR-0661
$$
，其中 $W'$ 是 $Y$ 在 $W$ 中的正规化；见《空间的态射》引理
\ref{spaces-morphisms-lemma-normalization-in-disjoint-union}、
\ref{spaces-morphisms-lemma-finite-integral} 和
\ref{spaces-morphisms-lemma-normalization-in-integral}。
断言成立，证明完成。
\end{proof}

\noindent
下一个引理与《空间的态射》引理
\ref{spaces-morphisms-lemma-quasi-finite-separated-quasi-affine}
重复。保留同一引理的两个副本，是因为证明有所不同。下面给出的证明
依赖于上面所述不可表示代数空间态射的 Zariski 主定理，而《空间的态射》
引理 \ref{spaces-morphisms-lemma-quasi-finite-separated-quasi-affine}
的证明则依赖于《空间的态射》命题
\ref{spaces-morphisms-proposition-locally-quasi-finite-separated-over-scheme}，
从而约化到概形态射的情形。

\begin{lemma}
\label{spaces-more-morphisms-lemma-quasi-finite-separated-quasi-affine}
设 $S$ 为概形。
设 $f : X \to Y$ 为 $S$ 上代数空间的态射。
假设 $f$ 拟有限且分离。设 $Y'$ 为 $Y$ 在 $X$ 中的正规化。
图示为：
$$
\xymatrix{
X \ar[rd]_f \ar[rr]_{f'} & & Y' \ar[ld]^\nu \\
& Y &
}
$$
则 $f'$ 是拟紧开浸入，且 $\nu$ 是整态射。特别地，$f$ 拟仿射。
\end{lemma}

\begin{proof}
这由引理 \ref{spaces-more-morphisms-lemma-finite-type-separated} 得出。事实上，由该引理，
存在开子空间 $U' \subset Y'$，使得
$(f')^{-1}(U') = X$（！）且 $X \to U'$ 是同构！换言之，$f'$ 是开浸入。
注意 $f'$ 拟紧，因为 $f$ 拟紧，且 $\nu : Y' \to Y$ 分离
（《空间的态射》引理
\ref{spaces-morphisms-lemma-quasi-compact-permanence}）。
因此，对于每个仿射概形 $Z$ 和态射 $Z \to Y$，纤维积
$Z \times_Y X$ 是仿射概形 $Z \times_Y Y'$ 的拟紧开子概形。
故依定义，$f$ 拟仿射。
\end{proof}

\begin{lemma}[Zariski 主定理]
\label{spaces-more-morphisms-lemma-quasi-finite-separated-pass-through-finite}
设 $S$ 为概形。
设 $f : X \to Y$ 为 $S$ 上代数空间的态射。
假设 $f$ 拟有限且分离，并假设 $Y$ 拟紧且拟分离。则存在分解
$$
\xymatrix{
X \ar[rd]_f \ar[rr]_j & & T \ar[ld]^\pi \\
& Y &
}
$$
，其中 $j$ 是拟紧开浸入，且 $\pi$ 有限。
\end{lemma}

\begin{proof}
设 $X \to Y' \to Y$ 如引理
\ref{spaces-more-morphisms-lemma-quasi-finite-separated-quasi-affine} 的结论所述。
由《空间的极限》引理
\ref{spaces-limits-lemma-integral-algebra-directed-colimit-finite}，
可以把 $\nu_*\mathcal{O}_{Y'} = \colim_{i \in I} \mathcal{A}_i$
写成有限拟凝聚 $\mathcal{O}_X$-代数 % P09-ERR-0662
$\mathcal{A}_i \subset \nu_*\mathcal{O}_{Y'}$ 的有向余极限。于是
$\pi_i : T_i = \underline{\Spec}_Y(\mathcal{A}_i) \to Y$
对于每个 $i$ 都是有限态射。注意转移态射 $T_{i'} \to T_i$
都是仿射态射，并且
$Y' = \lim T_i$。

\medskip\noindent
由《空间的极限》引理 \ref{spaces-limits-lemma-descend-opens}，
存在 $i$ 以及拟紧开集 $U_i \subset T_i$，其在 $Y'$ 中的逆像等于
$f'(X)$。对于 $i' \geq i$，令 $U_{i'}$ 为 $U_i$ 在 $T_{i'}$
中的逆像。于是
$X \cong f'(X) = \lim_{i' \geq i} U_{i'}$；见《空间的极限》引理
\ref{spaces-limits-lemma-directed-inverse-system-has-limit}。
由《空间的极限》引理
\ref{spaces-limits-lemma-finite-type-eventually-closed}，可知
$X \to U_{i'}$ 对于某个 $i' \geq i$ 是闭浸入。
（事实上，$X \cong U_{i'}$ 对充分大的 $i'$ 成立，但这里不需要。）
因此 $X \to T_{i'}$ 是浸入。由《空间的态射》引理
\ref{spaces-morphisms-lemma-factor-the-other-way}，可以将其分解为
$X \to T \to T_{i'}$，其中第一个箭头是开浸入，第二个是闭浸入。
证明完成。
\end{proof}

\begin{lemma}
\label{spaces-more-morphisms-lemma-quasi-finite-separated-pass-through-finite-addendum}
采用引理 \ref{spaces-more-morphisms-lemma-quasi-finite-separated-pass-through-finite}
的记号和假设。再假设 $f$ 局部有限表示。则可以选取分解，使
$T$ 在 $Y$ 上有限且有限表示。
\end{lemma}

\begin{proof}
由《空间的极限》引理
\ref{spaces-limits-lemma-finite-in-finite-and-finite-presentation}，
可以写成 $T = \lim T_i$，其中所有 $T_i$ 在 $Y$ 上有限且有限表示，
而转移态射 $T_{i'} \to T_i$ 都是闭浸入。由《空间的极限》引理
\ref{spaces-limits-lemma-descend-opens}，存在 $i$ 以及开子概形
$U_i \subset T_i$，其在 $T$ 中的逆像是 $X$。% P09-ERR-0663
由《空间的极限》引理
\ref{spaces-limits-lemma-finite-type-eventually-closed}，可知
$X \cong U_i$ 对充分大的 $i$ 成立。把 $T$ 替换为 $T_i$ 即完成证明。
\end{proof}





\section{Zariski 主定理的应用 I}
\label{spaces-more-morphisms-section-applications-zmt}

\noindent
第一个应用是把有限态射刻画为具有有限纤维的固有态射。

\begin{lemma}
\label{spaces-more-morphisms-lemma-characterize-finite}
设 $S$ 为概形。设 $f : X \to Y$ 为 $S$ 上代数空间的态射。
下列条件等价：
\begin{enumerate}
\item $f$ 有限；
\item $f$ 固有且局部拟有限；
\item $f$ 固有，并且 $|X_k|$ 对于每个态射
$\Spec(k) \to Y$ 都是离散空间，其中 $k$ 是域；
\item $f$ 万有闭、分离、局部有限型，并且 $|X_k|$ 对于每个态射
$\Spec(k) \to Y$ 都是离散空间，其中 $k$ 是域。
\end{enumerate}
\end{lemma}

\begin{proof}
由《空间的态射》引理
\ref{spaces-morphisms-lemma-finite-proper}、
\ref{spaces-morphisms-lemma-finite-quasi-finite}，有 (1) $\Rightarrow$ (2)。
由《空间的态射》引理
\ref{spaces-morphisms-lemma-locally-quasi-finite}，有 (2) $\Rightarrow$ (3)。
依定义，(3) 蕴含 (4)。

\medskip\noindent
假设 (4)。由于 $f$ 万有闭，它拟紧
（《空间的态射》引理
\ref{spaces-morphisms-lemma-universally-closed-quasi-compact}）。
选取点 $y$，它属于 $|Y|$。选取一个表示 $y$ 的态射
$\Spec(k) \to Y$。
注意 $|X_k|$ 作为拟紧离散空间是有限离散空间。
映射 $|X_k| \to |X|$ 满射到 $|X| \to |Y|$ 在 $y$ 上方的纤维
（《空间的性质》引理 \ref{spaces-properties-lemma-points-cartesian}）。
由《空间的态射》引理
\ref{spaces-morphisms-lemma-quasi-finite-at-point}，可知
$X \to Y$ 在 $|X| \to |Y|$ 位于 $y$ 上方的纤维的所有点处拟有限。
选取基本平展邻域 $(U, u) \to (Y, y)$ 以及分解
$X_U = V \amalg W$，如引理
\ref{spaces-more-morphisms-lemma-etale-splits-off-quasi-finite-part} 所述，并使其适合
$|X|$ 中位于 $y$ 上方的所有点。注意 $W_u = \emptyset$，
因为采用了 $|X| \to |Y|$ 位于 $y$ 上方纤维中的所有点。
由于 $f$ 万有闭，$|W|$ 在 $|U|$ 中的像是不含 $u$ 的闭集。
缩小 $U$ 后，可以假设 $W = \emptyset$。换言之，
$X_U = V$ 在 $U$ 上有限。由于 $y \in |Y|$ 任意，这意味着存在
平展态射族 $\{U_i \to Y\}$，其像覆盖 $Y$，并且基变换
$X_{U_i} \to U_i$ 都有限。由《空间的态射》引理
\ref{spaces-morphisms-lemma-integral-local}，可知 $f$ 有限。
\end{proof}

\noindent
作为推论，得到如下有用结果。

\begin{lemma}
\label{spaces-more-morphisms-lemma-proper-finite-fibre-finite-in-neighbourhood}
设 $S$ 为概形。设 $f : X \to Y$ 为 $S$ 上代数空间的态射。
设 $y \in |Y|$。假设
\begin{enumerate}
\item $f$ 固有；并且
\item $f$ 在所有 $x \in |X|$ 处拟有限，且这些点位于 $y$ 上方
（《良态空间》引理
\ref{decent-spaces-lemma-conditions-on-fibre-and-qf}）。
\end{enumerate}
则存在开邻域 $V \subset Y$，它是 $y$ 的邻域，并使得
$f|_{f^{-1}(V)} : f^{-1}(V) \to V$ 有限。
\end{lemma}

\begin{proof}
由《空间的态射》引理
\ref{spaces-morphisms-lemma-locally-finite-type-quasi-finite-part}，
$f$ 拟有限的点之集合是开集 $U \subset X$。
令 $Z = X \setminus U$。则 $y \not \in f(Z)$。由于 $f$ 固有，
$f(Z) \subset Y$ 是闭集。选取任意开邻域 $V \subset Y$，它含 $y$
且满足 $Z \cap V = \emptyset$。% P09-ERR-0664
于是 $f^{-1}(V) \to V$ 局部拟有限且固有。因此由引理
\ref{spaces-more-morphisms-lemma-characterize-finite}，$f^{-1}(V) \to V$ 有限。
\end{proof}

\begin{lemma}
\label{spaces-more-morphisms-lemma-flat-proper-family-cannot-collapse-fibre}
\begin{slogan}
压缩固有族的一个纤维，会迫使附近的纤维也被压缩。
\end{slogan}
设 $S$ 为概形。设
$$
\xymatrix{
X \ar[rr]_h \ar[rd]_f & & Y \ar[ld]^g \\
& B
}
$$
为 $S$ 上代数空间态射的交换图。% P09-ERR-0665
设 $b \in B$，并设 $\Spec(k) \to B$ 为 $b$ 的等价类中的态射。
% P09-ERR-0666
假设
\begin{enumerate}
\item $X \to B$ 是固有态射；
\item $Y \to B$ 分离且局部有限型；
\item 下列条件之一成立：
\begin{enumerate}
\item $|X_k| \to |Y_k|$ 的像有限；
\item $|f|^{-1}(\{b\})$ 在 $|Y|$ 中的像有限，且 $B$ 良态。
\end{enumerate}
\end{enumerate}
则存在开子空间 $B' \subset B$，它包含 $b$，并使得
$X_{B'} \to Y_{B'}$ 分解经过闭子空间 $Z \subset Y_{B'}$，
且后者在 $B'$ 上有限。
\end{lemma}

\begin{proof}
令 $Z \subset Y$ 为 $h$ 的概形论像；见《空间的态射》第
\ref{spaces-morphisms-section-scheme-theoretic-image} 节。
由《空间的态射》引理
\ref{spaces-morphisms-lemma-scheme-theoretic-image-is-proper}，
态射 $X \to Z$ 满，且 $Z \to B$ 固有。因此
$$
\{x \in |X|\text{ 映到 }b\} \to
\{z \in |Z|\text{ 映到 }b\}
$$
以及 $|X_k| \to |Z_k|$ 都满。可见，(3)(a) 或 (3)(b) 都蕴含
$Z \to B$ 在 $|Z|$ 中位于 $b$ 上方的所有点处拟有限。% P09-ERR-0667
这由《良态空间》引理
\ref{decent-spaces-lemma-conditions-on-fibre-and-qf} 得出。
因此，由引理
\ref{spaces-more-morphisms-lemma-proper-finite-fibre-finite-in-neighbourhood}，
$Z \to B$ 在 $b$ 的一个开邻域上有限。
\end{proof}







\section{Stein 分解}
\label{spaces-more-morphisms-section-stein-factorization}

\noindent
Stein 分解是如下陈述：若固有态射 $f : X \to S$ 满足
$f_*\mathcal{O}_X = \mathcal{O}_S$，则它的纤维连通。

\begin{lemma}
\label{spaces-more-morphisms-lemma-stein-universally-closed}
设 $S$ 为概形。设 $f : X \to Y$ 为 $S$ 上代数空间之间的普遍闭
拟分离态射。存在分解
$$
\xymatrix{
X \ar[rr]_{f'} \ar[rd]_f & & Y' \ar[dl]^\pi \\
& Y &
}
$$
满足下列性质：
\begin{enumerate}
\item 态射 $f'$ 普遍闭、拟紧、拟分离且满；
\item 态射 $\pi : Y' \to Y$ 为整态射；
\item 有 $f'_*\mathcal{O}_X = \mathcal{O}_{Y'}$；
\item 有 $Y' = \underline{\Spec}_Y(f_*\mathcal{O}_X)$；
\item $Y'$ 是《空间的态射》定义
\ref{spaces-morphisms-definition-normalization-X-in-Y}
所定义的 $Y$ 在 $X$ 中的正规化。
\end{enumerate}
分解 $f = \pi \circ f'$ 的形成与平坦基变换可交换。
\end{lemma}

\begin{proof}
由《空间的态射》引理
\ref{spaces-morphisms-lemma-universally-closed-quasi-compact}，
态射 $f$ 拟紧。我们直接把 $Y'$ 定义为 $Y$ 在 $X$ 中的正规化，
故 (5) 与 (2) 自动成立。由《空间的态射》引理
\ref{spaces-morphisms-lemma-normalization-in-universally-closed}，
可知 (4) 成立。由《空间的态射》引理
\ref{spaces-morphisms-lemma-universally-closed-permanence}，
态射 $f'$ 普遍闭。由《空间的态射》引理
\ref{spaces-morphisms-lemma-quasi-compact-permanence}，它拟紧；
由《空间的态射》引理
\ref{spaces-morphisms-lemma-compose-after-separated}，它拟分离。

\medskip\noindent
为证明其余陈述，可以假设基底 $Y$ 仿射
（因为取正规化与平展局部化可交换）。设 $Y = \Spec(R)$。
则 $Y' = \Spec(A)$，其中
$A = \Gamma(X, \mathcal{O}_X)$ 是整 $R$-代数。
于是显然 $f'_*\mathcal{O}_X$ 就是 $\mathcal{O}_{Y'}$
（因为由《空间的态射》引理
\ref{spaces-morphisms-lemma-pushforward}，
$f'_*\mathcal{O}_X$ 拟凝聚，从而等于 $\widetilde{A}$）。这证明了 (3)。

\medskip\noindent
下面证明 $f'$ 满。由于 $f'$ 普遍闭（见上文），$f'$ 的像是闭子集
$V(I) \subset Y' = \Spec(A)$。取 $h \in I$。则
$h|_X = f^\sharp(h)$ 是 $X$ 的结构层的整体截面，并且在每一点消失。
% P09-ERR-0684
由于 $X$ 拟紧，这意味着 $h|_X$ 是幂零截面，即
$h^n|X = 0$ 对某个 $n > 0$ 成立。% P09-ERR-0668
但 $A = \Gamma(X, \mathcal{O}_X)$，故 $h^n = 0$。
换言之，$I$ 包含于 $A$ 的 Jacobson 根中；由此得到
$V(I) = Y'$，如所需。% P09-ERR-0669
\end{proof}

\begin{lemma}
\label{spaces-more-morphisms-lemma-stein-universally-closed-residue-fields}
在引理 \ref{spaces-more-morphisms-lemma-stein-universally-closed} 中，进一步假设
$f$ 局部有限型且 $Y$ 仿射。则对 $y \in Y$，纤维
$\pi^{-1}(\{y\}) = \{y_1, \ldots, y_n\}$ 有限，并且域扩张
$\kappa(y_i)/\kappa(y)$ 有限。% P09-ERR-0670
\end{lemma}

\begin{proof}
回忆，$\pi^{-1}(\{y\})$ 的各点之间不存在特殊化；见《代数》引理
\ref{algebra-lemma-integral-no-inclusion}。
由于 $f'$ 满，可知 $|X_y| \to \pi^{-1}(\{y\})$ 满。
注意，$X_y$ 是域上有限型的拟分离代数空间
（拟紧性已在所引引理的证明中证得）。因此 $|X_y|$ 是诺特拓扑空间
（《空间的态射》引理
\ref{spaces-morphisms-lemma-finite-presentation-noetherian}）。
现在一个拓扑论证（略）表明 $\pi^{-1}(\{y\})$ 有限。% P09-ERR-0671
对每个 $i$，可以选取有限型点 $x_i \in |X_y|$，它映到 $y_i$
（《空间的态射》引理
\ref{spaces-morphisms-lemma-enough-finite-type-points}）。
由此可知 $\kappa(y_i)/\kappa(y)$ 有限：$x_i$ 可由有限型态射
$\Spec(k_i) \to X_y$ 表示（根据有限型点的定义），因而
$\Spec(k_i) \to y = \Spec(\kappa(y))$ 是有限型态射
（作为有限型态射的复合），% P09-ERR-0672
所以 $k_i/\kappa(y)$ 有限（《态射》引理
\ref{morphisms-lemma-point-finite-type}）。
\end{proof}

\noindent
设 $f : X \to Y$ 为代数空间的态射，并设
$\overline{y} : \Spec(k) \to Y$ 为几何点。则 $f$ 在
$\overline{y}$ 上方的{\it 纤维}是代数空间
$X_{\overline{y}} = X \times_{Y, \overline{y}} \Spec(k)$，它定义在 $k$ 上。
若 $Y$ 是概形且 $y \in Y$ 是一点，则仍照常记
$X_y = X \times_Y \Spec(\kappa(y))$ 为该纤维。% P09-ERR-0673

\begin{lemma}
\label{spaces-more-morphisms-lemma-characterize-geometrically-connected-fibres}
设 $S$ 为概形。设 $f : X \to Y$ 为 $S$ 上代数空间的态射。
设 $\overline{y}$ 为 $Y$ 的几何点。则 $X_{\overline{y}}$ 连通，当且仅当：
对每个平展邻域 $(V, \overline{v}) \to (Y, \overline{y})$，其中
$V$ 为概形，基变换 $X_V \to V$ 的纤维 $X_v$ 都连通。
\end{lemma}

\begin{proof}
由于 $\overline{y}$ 的平展邻域所成范畴是余滤过的，并包含由概形组成的
余终子集（《空间的性质》引理
\ref{spaces-properties-lemma-cofinal-etale}），
可以用其中一个邻域替换 $Y$，从而假设 $Y$ 为概形。
设 $y \in Y$ 为对应于 $\overline{y}$ 的点。则 $X_y$ 在
$\kappa(y)$ 上几何连通，当且仅当 $X_{\overline{y}}$ 连通；这又当且仅当
$(X_y)_{k'}$ 对每个有限可分扩张 $k'$（它扩张 $\kappa(y)$）都连通。
见《域上空间》第
\ref{spaces-over-fields-section-geometrically-connected} 节，尤其是引理
\ref{spaces-over-fields-lemma-characterize-geometrically-disconnected}。
由《态射进阶》引理
\ref{more-morphisms-lemma-realize-prescribed-residue-field-extension-etale}，
存在仿射平展邻域 $(V, v) \to (Y, y)$，使得
$\kappa(s) \subset \kappa(u)$ 与 $\kappa(s) \subset k'$
对任意给定的有限可分扩张相识别。% P09-ERR-0674
引理得证。
\end{proof}

\begin{theorem}[Stein 分解；诺特情形]
\label{spaces-more-morphisms-theorem-stein-factorization-Noetherian}
设 $S$ 为概形。设 $f : X \to Y$ 为 $S$ 上代数空间的固有态射，
并且 $Y$ 局部诺特。存在分解
$$
\xymatrix{
X \ar[rr]_{f'} \ar[rd]_f & & Y' \ar[dl]^\pi \\
& Y &
}
$$
满足下列性质：
\begin{enumerate}
\item 态射 $f'$ 固有且具有连通的几何纤维；
\item 态射 $\pi : Y' \to Y$ 有限；
\item 有 $f'_*\mathcal{O}_X = \mathcal{O}_{Y'}$；
\item 有 $Y' = \underline{\Spec}_Y(f_*\mathcal{O}_X)$；
\item $Y'$ 是 $Y$ 在 $X$ 中的正规化；见《态射》定义
\ref{morphisms-definition-normalization-X-in-Y}。
\end{enumerate}
\end{theorem}

\begin{proof}
令 $f = \pi \circ f'$ 为引理
\ref{spaces-more-morphisms-lemma-stein-universally-closed} 给出的分解。注意，除引理
\ref{spaces-more-morphisms-lemma-stein-universally-closed} 的结论外，$f'$ 还分离
（《空间的态射》引理
\ref{spaces-morphisms-lemma-compose-after-separated}）且有限型
（《空间的态射》引理
\ref{spaces-morphisms-lemma-permanence-finite-type}）。
因此 $f'$ 固有。由《空间的上同调》引理
\ref{spaces-cohomology-lemma-proper-pushforward-coherent}，
可知 $f_*\mathcal{O}_X$ 是凝聚 $\mathcal{O}_Y$-模。
因此 $\pi$ 有限，即 (2) 成立。

\medskip\noindent
这证明了除最有意义的断言外的所有断言；该断言是 $f'$ 的几何纤维连通。
由上面的讨论显然可以把 $Y$ 替换为 $Y'$。于是 $Y$ 局部诺特，
$f : X \to Y$ 固有，并且 $f_*\mathcal{O}_X = \mathcal{O}_Y$。
设 $\overline{y}$ 为 $Y$ 的一个几何点。此时应用形式函数定理，更确切地说，
应用《空间的上同调》引理
\ref{spaces-cohomology-lemma-formal-functions-stalk}。它告诉我们
$$
\mathcal{O}^\wedge_{Y, \overline{y}} =
\lim_n H^0(X_n, \mathcal{O}_{X_n})
$$
其中 $X_n =
\Spec(\mathcal{O}_{Y, \overline{y}}/\mathfrak m_{\overline{y}}^n) \times_Y X$。
注意，$X_1 = X_{\overline{y}} \to X_n$ 是（有限阶）加厚，
因而 $X_n$ 的底拓扑空间等于 $X_{\overline{y}}$ 的底拓扑空间。
于是，若 $X_{\overline{y}} = T_1 \amalg T_2$ 是两个非空开闭子空间的
不交并，则同样有 $X_n = T_{1, n} \amalg T_{2, n}$，且这对所有
$n$ 成立。这又意味着
$H^0(X_n, \mathcal{O}_{X_n})$ 含有非平凡幂等元 $e_{1, n}$，
即恒取值 $1$ 于 $T_{1, n}$、且恒取值 $0$ 于 $T_{2, n}$ 的函数。
显然，$e_{1, n + 1}$ 限制为 $e_{1, n}$ 于 $X_n$ 上。
因此 $e_1 = \lim e_{1, n}$ 是该极限中的非平凡幂等元。
这与 $\mathcal{O}^\wedge_{Y, \overline{y}}$ 是局部环矛盾。
故假设错误，即 $X_{\overline{y}}$ 连通，如所需。
\end{proof}

\begin{theorem}[Stein 分解；一般情形]
\label{spaces-more-morphisms-theorem-stein-factorization-general}
设 $S$ 为概形。设 $f : X \to Y$ 为 $S$ 上代数空间的固有态射。
存在分解
$$
\xymatrix{
X \ar[rr]_{f'} \ar[rd]_f & & Y' \ar[dl]^\pi \\
& Y &
}
$$
满足下列性质：
\begin{enumerate}
\item 态射 $f'$ 固有且具有连通的几何纤维；
\item 态射 $\pi : Y' \to Y$ 为整态射；
\item 有 $f'_*\mathcal{O}_X = \mathcal{O}_{Y'}$；
\item 有 $Y' = \underline{\Spec}_Y(f_*\mathcal{O}_X)$；
\item $Y'$ 是 $Y$ 在 $X$ 中的正规化（《空间的态射》定义
\ref{spaces-morphisms-definition-normalization-X-in-Y}）。
\end{enumerate}
\end{theorem}

\begin{proof}
可以应用引理 \ref{spaces-more-morphisms-lemma-stein-universally-closed} 得到态射
$f' : X \to Y'$。注意，除引理
\ref{spaces-more-morphisms-lemma-stein-universally-closed} 的结论外，$f'$ 还分离
（《空间的态射》引理
\ref{spaces-morphisms-lemma-compose-after-separated}）且有限型
（《空间的态射》引理
\ref{spaces-morphisms-lemma-permanence-finite-type}）。
因此 $f'$ 固有。至此，除 $f'$ 具有连通几何纤维这一陈述外，
所有陈述都已得到证明。

\medskip\noindent
由讨论可知，显然可以把 $Y$ 替换为 $Y'$。于是
$f : X \to Y$ 固有且 $f_*\mathcal{O}_X = \mathcal{O}_Y$。
注意，这些条件在平坦基变换下保持不变
（《空间的态射》引理 \ref{spaces-morphisms-lemma-base-change-proper}
以及《空间的上同调》引理
\ref{spaces-cohomology-lemma-flat-base-change-cohomology}）。
设 $\overline{y}$ 为 $Y$ 的几何点。由引理
\ref{spaces-more-morphisms-lemma-characterize-geometrically-connected-fibres}
及刚才的评注，问题约化到如下情形：$Y$ 是概形，$y \in Y$ 是一点，
% P09-ERR-0675
$f : X \to Y$ 是 $Y$ 上的固有代数空间，且
$f_*\mathcal{O}_X = \mathcal{O}_Y$；需要证明纤维 $X_y$ 连通。
把 $Y$ 替换为 $y$ 的仿射邻域，可以假设 $Y = \Spec(R)$ 仿射。
于是 $f_*\mathcal{O}_X = \mathcal{O}_Y$ 意味着环同态
$R \to \Gamma(X, \mathcal{O}_X)$ 是双射。

\medskip\noindent
由《空间的极限》引理
\ref{spaces-limits-lemma-proper-limit-of-proper-finite-presentation-noetherian}，
可以写成 $(X \to Y) = \lim (X_i \to Y_i)$，其中
$X_i \to Y_i$ 固有且有限表示，$Y_i$ 诺特。对充分大的 $i$，
$Y_i$ 仿射（《空间的极限》引理
\ref{spaces-limits-lemma-limit-is-affine}）。设 $Y_i = \Spec(R_i)$，
并令 $R'_i = \Gamma(X_i, \mathcal{O}_{X_i})$。
注意，有环同态 $R_i \to R_i' \to R$：第一个来自 $X_i$ 是
$R_i$ 上的代数空间；第二个来自 $X \to X_i$ 以及
$R = \Gamma(X, \mathcal{O}_X)$。由《空间的极限》引理
\ref{spaces-limits-lemma-descend-section}，注意到 $R = \colim R'_i$。
于是
$$
\xymatrix{
X \ar[d]  \ar[r] & X_i \ar[d] \\
Y \ar[r] & Y'_i \ar[r] & Y_i
}
$$
交换，其中 $Y'_i = \Spec(R'_i)$。设 $y'_i \in Y'_i$ 为 $y$ 的像。
% P09-ERR-0676
有 $X_y = \lim X_{i, y'_i}$，因为 $X = \lim X_i$、
$Y = \lim Y'_i$ 且 $\kappa(y) = \colim \kappa(y'_i)$。现在设
$X_y = U \amalg V$，其中 $U$ 与 $V$ 都既开又闭。
则 $U, V$ 是开集 $U_i, V_i$ 的逆像，而这些开集位于 $X_{i, y'_i}$ 中
（《空间的极限》引理 \ref{spaces-limits-lemma-descend-opens}）。
由定理 \ref{spaces-more-morphisms-theorem-stein-factorization-Noetherian}，
$X_i \to Y'_i$ 的纤维连通，故 $U$ 或 $V$ 之一为空。
证明完毕。
\end{proof}

\noindent
这里给出一个应用。

\begin{lemma}
\label{spaces-more-morphisms-lemma-geometrically-connected-fibres-towards-normal}
设 $S$ 为概形。设 $f : X \to Y$ 为 $S$ 上代数空间的态射。假设
\begin{enumerate}
\item $f$ 固有；
\item $Y$ 是整的（《域上空间》定义
\ref{spaces-over-fields-definition-integral-algebraic-space}），
其一般点为 $\xi$；
\item $Y$ 正规；
\item $X$ 约化；
\item $|X|$ 的每个不可约分支的一般点都映到 $\xi$；
\item 有 $H^0(X_\xi, \mathcal{O}) = \kappa(\xi)$。
\end{enumerate}
则 $f_*\mathcal{O}_X = \mathcal{O}_Y$，并且 $f$ 具有几何连通纤维。
\end{lemma}

\begin{proof}
应用定理 \ref{spaces-more-morphisms-theorem-stein-factorization-general} 得到分解
$X \to Y' \to Y$。只须证明 $Y' = Y$。
只须证明 $Y' \times_Y V \to V$ 是同构，其中 $V \to Y$ 为平展态射，
且 $V$ 是仿射整概形；见《域上空间》引理
\ref{spaces-over-fields-lemma-normal-integral-cover-by-affines}。
$Y'$ 的形成与平展基变换可交换；见《空间的态射》引理
\ref{spaces-morphisms-lemma-properties-normalization}。
$X \times_Y V$ 的一般点位于 $X$ 的一般点上方
（《良态空间》引理 \ref{decent-spaces-lemma-decent-generic-points}），
故由假设 (5)，它们映到 $V$ 的一般点。此外，条件 (6) 在由
$V \to Y$ 所作的基变换下保持不变，例如这可由平坦基变换得出
（《空间的上同调》引理
\ref{spaces-cohomology-lemma-flat-base-change-cohomology}）。
因此只须在 $Y$ 是正规整仿射概形的情形证明引理。

\medskip\noindent
假设 $Y$ 是正规整仿射概形。我们将应用《态射》引理
\ref{morphisms-lemma-finite-birational-over-normal}
证明 $Y' \to Y$ 是同构。事实上，$Y'$ 约化，因为 $X$ 约化
（《空间的态射》引理
\ref{spaces-morphisms-lemma-normalization-in-reduced}）。
由上面所引定理，态射 $Y' \to Y$ 为整态射。
由于 $Y$ 良态且 $X \to Y$ 分离，可知 $X$ 也良态；为此可用
《良态空间》引理
\ref{decent-spaces-lemma-properties-trivial-implications} 和
\ref{decent-spaces-lemma-property-over-property}。由假设 (5)、
《空间的态射》引理 \ref{spaces-morphisms-lemma-normalization-generic}
以及《良态空间》引理
\ref{decent-spaces-lemma-decent-generic-points}，可知
$|Y'|$ 的每个不可约分支的一般点都映到 $\xi$。另一方面，
由于 $Y'$ 是 $f_*\mathcal{O}_X$ 的相对谱，可知概形论纤维
$Y'_\xi$ 是 $H^0(X_\xi, \mathcal{O})$ 的谱；由假设，它等于
$\kappa(\xi)$。故 $Y'$ 是整概形，并且其函数域等于 $Y$ 的函数域。
证明完毕。
\end{proof}

\noindent
这里再给出一个应用。

\begin{lemma}
\label{spaces-more-morphisms-lemma-proper-flat-nr-geom-conn-comps-lower-semicontinuous}
设 $S$ 为概形。设 $X \to Y$ 为 $S$ 上代数空间的态射。
若 $f$ 固有、平坦且有限表示，则函数
% P09-ERR-0677
$n_{X/Y} : |Y| \to \mathbf{Z}$ 是下半连续的；该函数计数 $f$ 的纤维的
几何连通分支数（引理
\ref{spaces-more-morphisms-lemma-base-change-fibres-nr-geometrically-connected-components}）。
\end{lemma}

\begin{proof}
问题在 $Y$ 上是平展局部的，故可以且确实假设 $Y$ 是仿射概形。
设 $y \in Y$。令 $n = n_{X/S}(y)$。% P09-ERR-0678
注意 $n < \infty$，因为 $X \to Y$ 在 $y$ 处的几何纤维是域上的
固有代数空间，故诺特，因而只有有限个连通分支。
需要找出开邻域 $V$，它含 $y$，使得 $n_{X/S}|_V \geq n$。
% P09-ERR-0679
令 $X \to Y' \to Y$ 为定理
\ref{spaces-more-morphisms-theorem-stein-factorization-general} 中的 Stein 分解。
由引理 \ref{spaces-more-morphisms-lemma-stein-universally-closed-residue-fields}，
存在有限多个点 $y'_1, \ldots, y'_m \in Y'$，它们位于 $y$ 上方，
且扩张
$\kappa(y'_i)/\kappa(y)$ 有限。
《态射进阶》引理
\ref{more-morphisms-lemma-etale-makes-integral-split}
告诉我们，在把 $Y$ 替换为 $y$ 的一个平展邻域后，可以假设
$Y' = V_1 \amalg \ldots \amalg V_m$ 是概形，且
$y'_i \in V_i$ 以及 $\kappa(y'_i)/\kappa(y)$ 纯不可分。
于是代数空间 $X_{y_i'}$ 在 $\kappa(y)$ 上几何连通，故 $m = n$。
代数空间 $X_i = (f')^{-1}(V_i)$（$i = 1, \ldots, n$）
在 $Y$ 上平坦且有限表示。因此 $X_i \to Y$ 的像是开集
（《空间的态射》引理 \ref{spaces-morphisms-lemma-fppf-open}）。
故在 $y$ 的一个邻域内，$n_{X/Y}$ 至少为 $n$。
\end{proof}

\begin{lemma}
\label{spaces-more-morphisms-lemma-proper-flat-geom-red}
设 $S$ 为概形。设 $f : X \to Y$ 为 $S$ 上代数空间的态射。假设
\begin{enumerate}
\item $f$ 固有、平坦且有限表示；
\item $f$ 的几何纤维约化。
\end{enumerate}
则函数 $n_{X/S} : |Y| \to \mathbf{Z}$ 是局部常值的；该函数计数
% P09-ERR-0680
$f$ 的纤维的几何连通分支数
（引理 \ref{spaces-more-morphisms-lemma-base-change-fibres-nr-geometrically-connected-components}）。
\end{lemma}

\begin{proof}
由引理 \ref{spaces-more-morphisms-lemma-proper-flat-nr-geom-conn-comps-lower-semicontinuous}，
函数 $n_{X/Y}$ 下半连续。% P09-ERR-0681
故只须证明它上半连续。为此，可以在 $Y$ 上平展局部地工作，
因而可以假设 $Y$ 是仿射概形。对 $y \in Y$，考虑
$\kappa(y)$-代数
$$
A = H^0(X_y, \mathcal{O}_{X_y})
$$
由《域上空间》引理
\ref{spaces-over-fields-lemma-proper-geometrically-reduced-global-sections}
以及 $X_y$ 几何约化这一事实，$A$ 是 $\kappa(y)$ 的有限可分扩张的
有限乘积。因此 $A \otimes_{\kappa(y)} \kappa(\overline{y})$ 是由
$\beta_0(y) = \dim_{\kappa(y)} A$ 个 $\kappa(\overline{y})$ 副本组成的积。
故 $X_{\overline{y}}$ 有 $\beta_0(y)$ 个连通分支。
换言之，有 $n_{X/S} = \beta_0$，这是 $Y$ 上函数的等式。
% P09-ERR-0682
因此，由《空间的导出范畴》引理
\ref{spaces-perfect-lemma-jump-loci-geometric}，$n_{X/Y}$ 上半连续。
证明完毕。
\end{proof}

\begin{lemma}
\label{spaces-more-morphisms-lemma-stein-factorization-etale}
设 $S$ 为概形。设 $f : X \to Y$ 为 $S$ 上代数空间的固有态射。
设 $X \to Y' \to Y$ 为 $f$ 的 Stein 分解
（定理 \ref{spaces-more-morphisms-theorem-stein-factorization-general}）。
若 $f$ 有限表示、平坦且具有几何约化纤维
（定义 \ref{spaces-more-morphisms-definition-geometrically-reduced-fibre}），
则 $Y' \to Y$ 有限平展。
\end{lemma}

\begin{proof}
Stein 分解的形成与平坦基变换可交换；见引理
\ref{spaces-more-morphisms-lemma-stein-universally-closed}。因此可以在 $Y$ 上平展局部地工作，
并可假设 $Y$ 是仿射概形。此时 $Y'$ 是仿射概形，且
$Y' \to Y$ 为整态射。

\medskip\noindent
设 $y \in Y$。令 $n$ 为几何纤维 $X_{\overline{y}}$ 的连通分支数。
注意 $n < \infty$，因为 $X \to Y$ 在 $y$ 处的几何纤维是域上的
固有代数空间，故诺特，因而只有有限个连通分支。
由引理 \ref{spaces-more-morphisms-lemma-stein-universally-closed-residue-fields}，
存在有限多个点 $y'_1, \ldots, y'_m \in Y'$，它们位于 $y$ 上方；
并且对每个 $i$，可以选取有限型点
$x_i \in |X_y|$ 映到 $y'_i$，且扩张
$\kappa(y'_i)/\kappa(y)$ 有限。% P09-ERR-0683
因此《态射进阶》引理
\ref{more-morphisms-lemma-etale-makes-integral-split}
告诉我们，在把 $Y$ 替换为 $y$ 的一个平展邻域后，可以假设
$Y' = V_1 \amalg \ldots \amalg V_m$ 是概形，且
$y'_i \in V_i$ 以及 $\kappa(y'_i)/\kappa(y)$ 纯不可分。
此时，代数空间 $X_{y_i'}$ 在 $\kappa(y)$ 上几何连通，故 $m = n$。
代数空间 $X_i = (f')^{-1}(V_i)$（$i = 1, \ldots, n$）
在 $Y$ 上固有、平坦、有限表示，且具有几何约化纤维。
只须对每个态射 $X_i \to Y$ 证明引理。这把我们约化到
$X_{\overline{y}}$ 连通的情形。

\medskip\noindent
假设 $X_{\overline{y}}$ 连通。由引理
\ref{spaces-more-morphisms-lemma-proper-flat-geom-red}，可知 $X \to Y$ 在 $y$ 的一个邻域内
具有几何连通纤维。因此可以假设 $X \to Y$ 的纤维都几何连通。
于是由《空间的导出范畴》引理
\ref{spaces-perfect-lemma-proper-flat-geom-red-connected}，有
$f_*\mathcal{O}_X = \mathcal{O}_Y$，这就完成了证明。
\end{proof}

\noindent
下面引理的证明使用了概形的 Stein 分解，这就是它出现在本节的原因。

\begin{lemma}
\label{spaces-more-morphisms-lemma-split-off-proper-part-henselian}
设 $(A, I)$ 为亨泽尔对。设 $X$ 为 $A$ 上分离有限型的代数空间。令
$X_0 = X \times_{\Spec(A)} \Spec(A/I)$。
设 $Y \subset X_0$ 为既开又闭子空间，并且
$Y \to \Spec(A/I)$ 固有。则存在既开又闭子空间 $W \subset X$，
它在 $A$ 上固有，且
$W \times_{\Spec(A)} \Spec(A/I) = Y$。
\end{lemma}

\begin{proof}
我们用 $T \mapsto T_0$ 表示由
$\Spec(A/I) \to \Spec(A)$ 所作的基变换。
由 Chow 引理的一个弱版本（其形式见《空间的上同调》引理
\ref{spaces-cohomology-lemma-weak-chow}），存在满固有态射
$\varphi : X' \to X$，使得 $X'$ 可浸入 $\mathbf{P}^n_A$。
令 $Y' = \varphi^{-1}(Y)$。这是 $X'_0$ 的开闭子概形。
由《态射进阶》引理
\ref{more-morphisms-lemma-split-off-proper-part-henselian}，
引理对 $(X', Y')$ 成立。设 $W' \subset X'$ 为 $A$ 上固有的开闭子概形，
使得 $Y' = W'_0$。由《空间的态射》引理
\ref{spaces-morphisms-lemma-universally-closed-permanence}，
$Q_1 = \varphi(|W'|) \subset |X|$ 和
$Q_2 = \varphi(|X' \setminus W'|) \subset |X|$
是闭子集；又由《空间的态射》引理
\ref{spaces-morphisms-lemma-image-proper-is-proper}，
$Q_1$ 上的任意闭子空间结构在 $A$ 上固有。
$Q_1 \cap Q_2$ 在 $\Spec(A)$ 中的像是闭集。
由于 $(A, I)$ 是亨泽尔对，若 $Q_1 \cap Q_2$ 非空，则可知
$Q_1 \cap Q_2$ 有一点位于 $\Spec(A/I)$ 上方。
这是不可能的，因为 $W'_0 = Y' = \varphi^{-1}(Y)$。
由此得到 $Q_1$ 在 $|X|$ 中既开又闭。
令 $W \subset X$ 为对应的既开又闭子空间。则 $W$ 在 $A$ 上固有，
且 $W_0 = Y$。
\end{proof}










\section{从开子空间延拓性质}
\label{spaces-more-morphisms-section-extending-properties}

\noindent
本节汇集若干如下形式的结果：若 $f : X \to Y$ 是代数空间的平坦态射，
并且 $f$ 在 $Y$ 的一个稠密开子空间上满足某个性质，则 $f$ 在整个
$Y$ 上满足同一性质。

\begin{lemma}
\label{spaces-more-morphisms-lemma-flat-finite-type-finitely-presented-over-dense-open}
设 $S$ 为概形。设 $f : X \to Y$ 为 $S$ 上代数空间的态射。
设 $\mathcal{F}$ 为拟凝聚 $\mathcal{O}_X$-模。
设 $V \subset Y$ 为开子空间。假设
\begin{enumerate}
\item $f$ 局部有限表示；
\item $\mathcal{F}$ 有限型且在 $Y$ 上平坦；
\item $V \to Y$ 拟紧且概形论稠密；
\item $\mathcal{F}|_{f^{-1}V}$ 有限表示。
\end{enumerate}
则 $\mathcal{F}$ 有限表示。
\end{lemma}

\begin{proof}
只须证明 $\mathcal{F}$ 到某个概形上的拉回有限表示，该概形满且平展地
映到 $X$。故可假设 $X$ 为概形。% P09-ERR-0685
类似地，可以用一个概形替换 $Y$，该概形满且平展地映到 $Y$
（$V$ 在此概形中的逆像概形论稠密；见《空间的态射》第
\ref{spaces-morphisms-section-scheme-theoretic-closure} 节）。
% P09-ERR-0686
于是问题约化到概形情形，即《平坦性进阶》引理
\ref{flat-lemma-flat-finite-type-finitely-presented-over-dense-open}。
\end{proof}

\begin{lemma}
\label{spaces-more-morphisms-lemma-flat-finite-type-finitely-presented-over-dense-open-X}
设 $S$ 为概形。设 $f : X \to Y$ 为 $S$ 上代数空间的态射。
设 $V \subset Y$ 为开子空间。假设
\begin{enumerate}
\item $f$ 局部有限型且平坦；
\item $V \to Y$ 拟紧且概形论稠密；
\item $f|_{f^{-1}V} : f^{-1}V \to V$ 局部有限表示。
\end{enumerate}
则 $f$ 局部有限表示。% P09-ERR-0687
\end{lemma}

\begin{proof}
证明与引理
\ref{spaces-more-morphisms-lemma-flat-finite-type-finitely-presented-over-dense-open}
的证明相同，只须改用《平坦性进阶》引理
\ref{flat-lemma-flat-finite-type-finitely-presented-over-dense-open-X}。
\end{proof}

\begin{lemma}
\label{spaces-more-morphisms-lemma-flat-fp-dimension-over-dense-open}
设 $S$ 为概形。设 $f : X \to Y$ 为 $S$ 上代数空间的态射，
它平坦且局部有限型。设 $V \subset Y$ 为开子空间，满足
$|V| \subset |Y|$ 稠密且 $X_V \to V$ 的相对维数 $\leq d$。
若还满足下列条件之一：
\begin{enumerate}
\item $f$ 局部有限表示；
\item $V \to Y$ 拟紧，
\end{enumerate}
则 $f : X \to Y$ 的相对维数 $\leq d$。
\end{lemma}

\begin{proof}
可以把 $Y$ 替换为其约化，因而可以假设 $Y$ 约化。
于是 $V$ 在 $Y$ 中概形论稠密；见《空间的态射》引理
\ref{spaces-morphisms-lemma-quasi-compact-immersion}。
根据定义，相对维数 $\leq d$ 这一性质可以在平展覆盖上检验；见
《空间的态射》第 \ref{spaces-morphisms-section-relative-dimension} 节。
% P09-ERR-0688
因此，只须证明 $f$ 的相对维数 $\leq d$；为此可把 $X$ 替换为
满且平展地映到 $X$ 的一个概形。类似地，可以用一个概形替换 $Y$，
该概形满且平展地映到 $Y$。% P09-ERR-0689
$V$ 在此概形中的逆像概形论稠密；见《空间的态射》第
\ref{spaces-morphisms-section-scheme-theoretic-closure} 节。
由于概形的概形论稠密开子概形尤其稠密，问题约化到概形情形，即
《平坦性进阶》引理
\ref{flat-lemma-flat-finite-presentation-dimension-over-dense-open}。
\end{proof}

\begin{lemma}
\label{spaces-more-morphisms-lemma-proper-flat-finite-over-dense-open}
设 $S$ 为概形。设 $f : X \to Y$ 为 $S$ 上代数空间的态射，
它平坦且固有。设 $V \to Y$ 为开子空间，满足
$|V| \subset |Y|$ 稠密且 $X_V \to V$ 有限。% P09-ERR-0690
若还满足 $f$ 局部有限表示或 $V \to Y$ 拟紧之一，则 $f$ 有限。
% P09-ERR-0691
\end{lemma}

\begin{proof}
由引理 \ref{spaces-more-morphisms-lemma-flat-fp-dimension-over-dense-open}，$f$ 的纤维维数为零。
由《空间的态射》引理
\ref{spaces-morphisms-lemma-locally-quasi-finite-rel-dimension-0}，
这蕴含 $f$ 局部拟有限。由《空间的态射》引理
\ref{spaces-morphisms-lemma-locally-quasi-finite-separated-representable}，
这蕴含 $f$ 可表。可以检验 $f$ 是否有限，而这一检验在 $Y$ 上平展局部；
故可假设 $Y$ 是概形。由于 $f$ 可表，问题约化到概形情形，即
《平坦性进阶》引理 \ref{flat-lemma-proper-flat-finite-over-dense-open}。
\end{proof}

\begin{lemma}
\label{spaces-more-morphisms-lemma-zariski}
设 $S$ 为概形。设 $f : X \to Y$ 为 $S$ 上代数空间的态射。
设 $V \subset Y$ 为开子空间。若
\begin{enumerate}
\item $f$ 分离、局部有限型且平坦；
\item $f^{-1}(V) \to V$ 是同构；
\item $V \to Y$ 拟紧且概形论稠密，
\end{enumerate}
则 $f$ 是开浸入。
\end{lemma}

\begin{proof}
应用引理
\ref{spaces-more-morphisms-lemma-flat-finite-type-finitely-presented-over-dense-open-X}，
可知 $f$ 局部有限表示。应用引理
\ref{spaces-more-morphisms-lemma-flat-fp-dimension-over-dense-open}，可知 $f$ 的相对维数
$\leq 0$。由《空间的态射》引理
\ref{spaces-morphisms-lemma-locally-quasi-finite-rel-dimension-0}，
这蕴含 $f$ 局部拟有限。由《空间的态射》引理
\ref{spaces-morphisms-lemma-locally-quasi-finite-separated-representable}，
这蕴含 $f$ 可表。由《空间上的下降》引理
\ref{spaces-descent-lemma-descending-property-open-immersion}，
可以检验 $f$ 是否为开浸入，而这一检验在 $Y$ 上平展局部。
故可假设 $Y$ 是概形。由于 $f$ 可表，问题约化到概形情形，即
《平坦性进阶》引理 \ref{flat-lemma-zariski}。
\end{proof}

\section{爆破与平坦性}
\label{spaces-more-morphisms-section-blowup-flat}

\noindent
我们不再重复《平坦性进阶》第 \ref{flat-section-blowup-flat} 节中的工作，
而是证明《平坦性进阶》引理 \ref{flat-lemma-push-ideal} 的一个类似结果；
它说明，寻找合适爆破的问题通常在基底上是平展局部的。

\begin{lemma}
\label{spaces-more-morphisms-lemma-push-ideal}
设 $S$ 为概形。设 $X$ 为 $S$ 上拟紧且拟分离的代数空间。
设 $\varphi : W \to X$ 为拟紧、分离的平展态射。
设 $U \subset X$ 为拟紧开子空间。
设 $\mathcal{I} \subset \mathcal{O}_W$ 为有限型拟凝聚理想层，满足
$V(\mathcal{I}) \cap \varphi^{-1}(U) = \emptyset$。
则存在有限型拟凝聚理想层
$\mathcal{J} \subset \mathcal{O}_X$，使得
\begin{enumerate}
\item $V(\mathcal{J}) \cap U = \emptyset$；
\item $\varphi^{-1}(\mathcal{J})\mathcal{O}_W = \mathcal{I} \mathcal{I}'$，
其中 $\mathcal{I}' \subset \mathcal{O}_W$ 是某个有限型拟凝聚理想。
\end{enumerate}
\end{lemma}

\begin{proof}
选择分解 $W \to Y \to X$，其中 $j : W \to Y$ 是拟紧开浸入，
$\pi : Y \to X$ 是有限表示的有限态射
（引理 \ref{spaces-more-morphisms-lemma-quasi-finite-separated-pass-through-finite-addendum}）。
令 $V = j(W) \cup \pi^{-1}(U) \subset Y$。注意，
$\mathcal{I}$ 在 $W \cong j(W)$ 上，它与另一开子空间上的
$\mathcal{O}_{\pi^{-1}(U)}$ 可粘合成有限型拟凝聚理想层
$\mathcal{I}_1 \subset \mathcal{O}_V$。由《空间的极限》引理
\ref{spaces-limits-lemma-extend}，存在有限型拟凝聚理想层
$\mathcal{I}_2 \subset \mathcal{O}_Y$，使得
$\mathcal{I}_2|_V = \mathcal{I}_1$。换言之，
$\mathcal{I}_2 \subset \mathcal{O}_Y$ 是有限型拟凝聚理想层，满足
$V(\mathcal{I}_2)$ 与 $\pi^{-1}(U)$ 不交，并且
$j^{-1}\mathcal{I}_2 = \mathcal{I}$。以 $i : Z \to Y$
表示相应的闭浸入；它是有限表示的（《空间的态射》引理
\ref{spaces-morphisms-lemma-closed-immersion-finite-presentation}）。
% P09-ERR-0692
特别地，复合 $\tau = \pi \circ i : Z \to X$ 是有限且有限表示的
（《空间的态射》引理
\ref{spaces-morphisms-lemma-composition-finite-presentation} 和
\ref{spaces-morphisms-lemma-composition-integral}）。

\medskip\noindent
令 $\mathcal{F} = \tau_*\mathcal{O}_Z$，并把它视为拟凝聚
$\mathcal{O}_X$-模。由《空间上的下降》引理
\ref{spaces-descent-lemma-finite-finitely-presented-module}，
可知 $\mathcal{F}$ 是有限表示的 $\mathcal{O}_X$-模。
令 $\mathcal{J} = \text{Fit}_0(\mathcal{F})$。
（此处插入带环拓扑斯上模的 Fitting 理想的参考文献。）
% P09-ERR-0693
这是 $X$ 上的有限型拟凝聚理想层（因为 $\mathcal{F}$ 有限表示；见
《代数进阶》引理 \ref{more-algebra-lemma-fitting-ideal-basics}）。
$|\tau|(|Z|) \cap |U| = \emptyset$；这是由
$\mathcal{I}_2$ 的选取而得。而零模的第 $0$ 个 Fitting 理想
等于结构层，故引理的第 (1) 部分成立。为证明 (2)，注意
$\varphi^{-1}(\mathcal{J})\mathcal{O}_W = \text{Fit}_0(\varphi^*\mathcal{F})$，
因为取 Fitting 理想与基变换可交换。另一方面，由于
$\varphi : W \to X$ 分离且平展，可知
$(1, j) : W \to W \times_X Y$ 是开闭浸入。
故有 $W \times_Y Z = V(\mathcal{I}) \amalg Z'$，其中
$\tau' : Z' \to W$ 是某个有限且有限表示的代数空间态射。
% P09-ERR-0694
因此
\begin{align*}
\text{Fit}_0(\varphi^*\mathcal{F}) & =
\text{Fit}_0((W \times_Y Z \to W)_*\mathcal{O}_{W \times_Y Z}) \\
& =
\text{Fit}_0(\mathcal{O}_W/\mathcal{I}) \cdot
\text{Fit}_0(\tau'_*\mathcal{O}_{Z'}) \\
& =
\mathcal{I} \cdot \text{Fit}_0(\tau'_*\mathcal{O}_{Z'})
\end{align*}
其中第二个等式由《代数进阶》引理
\ref{more-algebra-lemma-fitting-ideal-basics}
在带环拓扑斯上的层中转述而得。% P09-ERR-0695
令 $\mathcal{I}' = \text{Fit}_0(\tau'_*\mathcal{O}_{Z'})$，
即完成引理的证明。
\end{proof}

\begin{theorem}
\label{spaces-more-morphisms-theorem-flatten-module}
设 $S$ 为概形。设 $B$ 为 $S$ 上拟紧且拟分离的代数空间。
设 $X$ 为 $B$ 上的代数空间。设 $\mathcal{F}$ 为 $X$ 上的拟凝聚模。
设 $U \subset B$ 为拟紧开子空间。假设
\begin{enumerate}
\item $X$ 拟紧；
\item $X$ 在 $B$ 上局部有限表示；
\item $\mathcal{F}$ 是有限型模；
\item $\mathcal{F}_U$ 有限表示；
\item $\mathcal{F}_U$ 在 $U$ 上平坦。
\end{enumerate}
则存在一个 $U$-容许爆破 $B' \to B$，使得严格变换
$\mathcal{F}'$（即 $\mathcal{F}$ 的严格变换）是有限表示的
$\mathcal{O}_{X \times_B B'}$-模，并且在 $B'$ 上平坦。
\end{theorem}

\begin{proof}
选择仿射概形 $V$ 以及满平展态射 $V \to X$。
由于严格变换与平展局部化可交换
（《空间上的除子》引理 \ref{spaces-divisors-lemma-strict-transform-local}），
只须证明把 $X$ 替换为 $V$ 后所得的结论。故可假设
$X \to B$ 可表（此外仍满足引理的各项假设）。% P09-ERR-0696

\medskip\noindent
假设 $X \to B$ 可表。选择仿射概形 $W$ 以及满平展态射
$\varphi : W \to B$。注意 $X \times_B W$ 是概形。
由概形情形（《平坦性进阶》定理 \ref{flat-theorem-flatten-module}），
可以找到有限型拟凝聚理想层
$\mathcal{I} \subset \mathcal{O}_W$，使得
(a) $|V(\mathcal{I})| \cap |\varphi^{-1}(U)| = \emptyset$，并且 (b)
$\mathcal{F}|_{X \times_B W}$ 关于爆破 $W' \to W$
（其中心为 $\mathcal{I}$）的严格变换在 $W'$ 上平坦且是有限表示模。
按引理 \ref{spaces-more-morphisms-lemma-push-ideal} 选择有限型理想层
$\mathcal{J} \subset \mathcal{O}_B$。令 $B' \to B$ 为
$\mathcal{J}$ 的爆破。我们断言这一爆破满足要求。事实上，显然
$B' \to B$ 是 $U$-容许的；这是由 $\mathcal{J}$ 的选取而得。
此外，基变换 $B' \times_B W \to W$ 是 $W$ 沿
$\varphi^{-1}\mathcal{J} = \mathcal{I}\mathcal{I}'$ 的爆破
（爆破与平坦基变换的相容性；见《空间上的除子》引理
\ref{spaces-divisors-lemma-flat-base-change-blowing-up}）。
因而存在分解
$$
W \times_B B' \to W' \to W
$$
其中第一个态射也是爆破；见《空间上的除子》引理
\ref{spaces-divisors-lemma-blowing-up-two-ideals}。
$\mathcal{F}'$（它位于 $B' \times_B X$ 上）到
$W \times_B B' \times_B X$ 的限制，是
$\mathcal{F}|_{X \times_B W}$ 的严格变换
（《空间上的除子》引理 \ref{spaces-divisors-lemma-strict-transform-local}），
因而也是 $\mathcal{F}|_{X \times_B W}$ 经上述两个爆破所得的
两次迭代严格变换（《空间上的除子》引理
\ref{spaces-divisors-lemma-strict-transform-composition-blowups}）。
第一次爆破后，所得层按构造已经在基底上平坦且有限表示。
因此第二次严格变换后仍然如此（因为由《空间上的除子》引理
\ref{spaces-divisors-lemma-strict-transform-flat}，该变换是拉回）。
于是可知，$\mathcal{F}'$ 到 $B' \times_B X$ 的一个平展覆盖上的限制
具有所需性质，定理得证。
\end{proof}

\section{应用}
\label{spaces-more-morphisms-section-applications-flattening-by-blowing-up}

\noindent
本节应用通过爆破实现平坦化的结果。

\begin{lemma}
\label{spaces-more-morphisms-lemma-flat-after-blowing-up}
设 $S$ 为概形。设 $f : X \to B$ 为 $S$ 上代数空间的态射。
设 $U \subset B$ 为开子空间。假设
\begin{enumerate}
\item $B$ 拟紧且拟分离；
\item $U$ 拟紧；
\item $f : X \to B$ 有限型且拟分离；
\item $f^{-1}(U) \to U$ 平坦且局部有限表示。
\end{enumerate}
则存在一个 $U$-容许爆破 $B' \to B$，使得严格变换
$X'$（即 $X$ 的严格变换）在 $B'$ 上平坦且有限表示。
\end{lemma}

\begin{proof}
设 $B' \to B$ 为一个 $U$-容许爆破。注意，$X$ 的严格变换在
$B'$ 上拟紧且拟分离，因为 $X$ 在 $B$ 上拟紧且拟分离。
故只须寻找一个 $U$-容许爆破，使得严格变换平坦且局部有限表示。
不能直接应用定理 \ref{spaces-more-morphisms-theorem-flatten-module}，因为 $X$ 在 $B$
上不局部有限表示。

\medskip\noindent
选择仿射概形 $V$ 以及满平展态射 $V \to X$。
（这是可能的，因为 $X$ 作为拟紧空间 $B$ 上的有限型空间而拟紧。）
于是只须对态射 $V \to B$ 证明结论（因为严格变换与平展局部化可交换；
见《空间上的除子》引理
\ref{spaces-divisors-lemma-strict-transform-local}）。
故可假设 $X \to B$ 除有限型外还分离。
在此情形下，可以找到闭浸入 $i : X \to Y$，其中 $Y \to B$
分离且有限表示；见《空间的极限》命题
\ref{spaces-limits-proposition-separated-closed-in-finite-presentation}。

\medskip\noindent
对 $\mathcal{F} = i_*\mathcal{O}_X$（视为 $Y/B$ 上的模）应用定理
\ref{spaces-more-morphisms-theorem-flatten-module}。得到一个 $U$-容许爆破
$B' \to B$，使得 $\mathcal{F}$ 的严格变换在 $B'$ 上平坦且有限表示。
设 $X'$ 为 $X$ 在爆破 $B' \to B$ 下的严格变换。
设 $i' : X' \to Y \times_B B'$ 为诱导态射。
由于取严格变换与沿仿射态射的推前可交换（《空间上的除子》引理
\ref{spaces-divisors-lemma-strict-transform-affine}），
可知 $i'_*\mathcal{O}_{X'}$ 在 $B'$ 上平坦，并且作为
$\mathcal{O}_{Y \times_B B'}$-模有限表示。
因此 $X' \to B'$ 平坦且局部有限表示。
结合前述讨论即得引理。
\end{proof}

\begin{lemma}
\label{spaces-more-morphisms-lemma-finite-after-blowing-up}
设 $S$ 为概形。设 $f : X \to B$ 为 $S$ 上代数空间的态射。
设 $U \subset B$ 为开子空间。假设
\begin{enumerate}
\item $B$ 拟紧且拟分离；
\item $U$ 拟紧；
\item $f : X \to B$ 固有；
\item $f^{-1}(U) \to U$ 有限局部自由。
\end{enumerate}
则存在一个 $U$-容许爆破 $B' \to B$，使得严格变换
$X'$（即 $X$ 的严格变换）在 $B'$ 上有限局部自由。
\end{lemma}

\begin{proof}
由引理 \ref{spaces-more-morphisms-lemma-flat-after-blowing-up}，可假设
$X \to B$ 平坦且有限表示。必要时把 $B$ 替换为一个
$U$-容许爆破，可以假设 $U \subset B$ 概形论稠密。
于是由引理 \ref{spaces-more-morphisms-lemma-proper-flat-finite-over-dense-open}，$f$ 有限。
故由《空间的态射》引理
\ref{spaces-morphisms-lemma-finite-flat}，$f$ 有限局部自由。
\end{proof}

\begin{lemma}
\label{spaces-more-morphisms-lemma-isomorphism-after-blowing-up}
设 $S$ 为概形。设 $f : X \to B$ 为 $S$ 上代数空间的态射。
设 $U \subset B$ 为开子空间。假设
\begin{enumerate}
\item $B$ 拟紧且拟分离；
\item $U$ 拟紧；
\item $f : X \to B$ 固有；
\item $f^{-1}(U) \to U$ 是同构。
\end{enumerate}
则存在一个 $U$-容许爆破 $B' \to B$，使得严格变换
$X'$（即 $X$ 的严格变换）同构地映到 $B'$。
\end{lemma}

\begin{proof}
由引理 \ref{spaces-more-morphisms-lemma-flat-after-blowing-up}，可假设
$X \to B$ 平坦且有限表示。必要时把 $B$ 替换为一个
$U$-容许爆破，可以假设 $U \subset B$ 概形论稠密。
于是由引理 \ref{spaces-more-morphisms-lemma-proper-flat-finite-over-dense-open}，$f$ 有限；
由引理 \ref{spaces-more-morphisms-lemma-zariski}，它是开浸入。
故 $f$ 是像为闭集且包含稠密开子空间 $U$ 的开浸入，
从而 $f$ 是同构。
\end{proof}

\begin{lemma}
\label{spaces-more-morphisms-lemma-zariski-after-blowup}
设 $S$ 为概形。设 $f : X \to B$ 为 $S$ 上代数空间的态射。
设 $U \subset B$ 为开子空间。假设
\begin{enumerate}
\item $B$ 拟紧且拟分离；% P09-ERR-0698
\item $U$ 拟紧；
\item $f$ 有限型；
\item $f^{-1}(U) \to U$ 是同构。
\end{enumerate}
则存在一个 $U$-容许爆破 $B' \to B$，使 $U$ 在 $B'$ 中概形论稠密，
并且严格变换 $X'$（即 $X$ 的严格变换）同构地映到 $B'$ 的一个开子空间。
\end{lemma}

\begin{proof}
本引理是引理 \ref{spaces-more-morphisms-lemma-isomorphism-after-blowing-up} 的推广。
由于 $U$-容许爆破的复合仍是 $U$-容许的（《空间上的除子》引理
\ref{spaces-divisors-lemma-composition-admissible-blowups}），
可以分阶段进行。选取有限型拟凝聚理想层
$\mathcal{I} \subset \mathcal{O}_B$，使得
$|B| \setminus |U| = |V(\mathcal{I})|$。把 $B$ 替换为 $B$ 沿
$\mathcal{I}$ 的爆破，并把 $X$ 替换为 $X$ 的严格变换。
替换后，$B \setminus U$ 是某个有效 Cartier 除子 $D$ 的支撑
（《空间上的除子》引理
\ref{spaces-divisors-lemma-blowing-up-gives-effective-Cartier-divisor}）。
特别地，$U$ 在 $B$ 中概形论稠密（《空间上的除子》引理
\ref{spaces-divisors-lemma-complement-effective-Cartier-divisor}）。
接着再作一次 $U$-容许爆破，使得 $X \to B$ 平坦且有限表示；
见引理 \ref{spaces-more-morphisms-lemma-flat-after-blowing-up}。
注意 $U$ 在 $B$ 中仍概形论稠密。因此由引理
\ref{spaces-more-morphisms-lemma-zariski}，$X \to B$ 是开浸入。
\end{proof}

\noindent
下述引理说明，一个修改可被某个爆破支配。

\begin{lemma}
\label{spaces-more-morphisms-lemma-dominate-modification-by-blowup}
设 $S$ 为概形。设 $f : X \to B$ 为 $S$ 上代数空间的态射。
设 $U \subset B$ 为开子空间。假设
\begin{enumerate}
\item $B$ 拟紧且拟分离；
\item $U$ 拟紧；
\item $f : X \to B$ 固有；
\item $f^{-1}(U) \to U$ 是同构。% P09-ERR-0699
\end{enumerate}
则存在一个 $U$-容许爆破 $B' \to B$，它支配 $X$；
换言之，爆破态射存在分解 $B' \to X \to B$。
\end{lemma}

\begin{proof}
由引理 \ref{spaces-more-morphisms-lemma-isomorphism-after-blowing-up}，可找到一个
$U$-容许爆破 $B' \to B$，使得严格变换 $X'$ 同构地映到
$B'$。于是可用 $B' = X' \to X$ 作为该分解。
\end{proof}

\begin{lemma}
\label{spaces-more-morphisms-lemma-get-section-after-blowup}
设 $S$ 为概形。设 $X$、$Y$ 为 $S$ 上的代数空间。
设 $U \subset W \subset Y$ 为开子空间。
设 $f : X \to W$ 与 $s : U \to X$ 为态射，并满足
$f \circ s = \text{id}_U$。假设
\begin{enumerate}
\item $f$ 固有；
\item $Y$ 拟紧且拟分离；
\item $U$ 与 $W$ 拟紧。
\end{enumerate}
则存在一个 $U$-容许爆破 $b : Y' \to Y$ 以及态射
$s' : b^{-1}(W) \to X$；它延拓 $s$，并满足
$f \circ s' = b|_{b^{-1}(W)}$。
\end{lemma}

\begin{proof}
可以并且将把 $X$ 替换为 $s$ 的概形论像。
此时 $X \to W$ 在 $U$ 上是同构；见《空间的态射》引理
\ref{spaces-morphisms-lemma-scheme-theoretic-image-of-partial-section}。
由引理 \ref{spaces-more-morphisms-lemma-dominate-modification-by-blowup}，存在一个
$U$-容许爆破 $W' \to W$ 以及态射 $W' \to X$；后者延拓 $s$。
最后应用《空间上的除子》引理
\ref{spaces-divisors-lemma-extend-admissible-blowups}，
把 $W' \to W$ 延拓为一个 $U$-容许爆破，其靶为 $Y$，即完成证明。
\end{proof}
 
\section{Chow 引理}
\label{spaces-more-morphisms-section-chow}

\noindent
本节证明 Chow 引理
（引理 \ref{spaces-more-morphisms-lemma-chow-noetherian-separated}）。
我们建议读者参阅《空间的上同调》第
\ref{spaces-cohomology-section-weak-chow} 节，那里给出了
Chow 引理的一个容易证明并足以用于许多应用的弱版本。

\medskip\noindent
由于我们尚未定义代数空间的射影态射，
各引理的陈述（例如见引理
\ref{spaces-more-morphisms-lemma-blowup-to-find-embedding}）将涉及可表固有态射，
而不是射影态射。

\begin{lemma}
\label{spaces-more-morphisms-lemma-find-common-blowups}
设 $S$ 为概形。设 $Y$ 为 $S$ 上拟紧且拟分离的代数空间。
设 $U \to X_1$ 与 $U \to X_2$ 是 $Y$ 上代数空间的开浸入，
并假设 $U$、$X_1$、$X_2$ 在 $Y$ 上都是有限型且分离的。% P09-ERR-0700
则存在代数空间的交换图
$$
\xymatrix{
X_1' \ar[d] \ar[r] & X & X_2' \ar[l] \ar[d] \\
X_1 & U \ar[l] \ar[lu] \ar[u] \ar[ru] \ar[r] & X_2
}
$$
它在 $Y$ 上，其中 $X_i' \to X_i$ 是 $U$-容许爆破，
$X_i' \to X$ 是开浸入，
且 $X$ 在 $Y$ 上分离且有限型。
\end{lemma}

\begin{proof}
在整个证明中，所有代数空间都在 $Y$ 上分离且有限型。% P09-ERR-0701
特别地，这蕴含这些代数空间拟紧且拟分离，并且它们之间的态射
拟紧且分离。见《空间的态射》第
\ref{spaces-morphisms-section-separation-axioms} 节和第
\ref{spaces-morphisms-section-quasi-compact} 节。
我们将使用如下事实：若 $U \to W$ 是这类 $Y$ 上空间的浸入，
则概形论像 $Z$，即 $U$ 在 $W$ 中的概形论像，是 $W$ 的闭子空间，
$U \to Z$ 是开浸入，$U \subset Z$ 概形论稠密，并且
$|U| \subset |Z|$ 稠密。见《空间的态射》引理
\ref{spaces-morphisms-lemma-quasi-compact-immersion}。

\medskip\noindent
设 $X_{12} \subset X_1 \times_Y X_2$ 为
$U \to X_1 \times_Y X_2$ 的概形论像。由《空间的态射》引理
\ref{spaces-morphisms-lemma-scheme-theoretic-image-of-partial-section}，
投影 $p_i : X_{12} \to X_i$ 诱导同构 $p_i^{-1}(U) \to U$。
选取一个 $U$-容许爆破 $X_i^i \to X_i$，使得严格变换
$X_{12}^i$，即 $X_{12}$ 的严格变换，同构于 $X_i^i$ 的一个开子空间；见引理
\ref{spaces-more-morphisms-lemma-zariski-after-blowup}。
设 $\mathcal{I}_i \subset \mathcal{O}_{X_i}$ 为相应的有限型拟凝聚理想层。
回忆 $X_{12}^i \to X_{12}$ 是沿
$p_i^{-1}\mathcal{I}_i \mathcal{O}_{X_{12}}$ 的爆破；见
《空间上的除子》引理 \ref{spaces-divisors-lemma-strict-transform}。
设 $X_{12}'$ 为 $X_{12}$ 沿
$p_1^{-1}\mathcal{I}_1 p_2^{-1}\mathcal{I}_2 \mathcal{O}_{X_{12}}$
的爆破；关于这意味着什么，见《空间上的除子》引理
\ref{spaces-divisors-lemma-blowing-up-two-ideals}。我们得到交换图
$$
\xymatrix{
X_{12}' \ar[d] \ar[r] & X_{12}^2 \ar[d] \\
X_{12}^1 \ar[r] & X_{12}
}
$$
其中所有态射都是 $U$-容许爆破。由于
$X_{12}^i \subset X_i^i$ 是开子空间，可以选取一个
$U$-容许爆破 $X_i' \to X_i^i$，其限制为
$X_{12}' \to X_{12}^i$；见《空间上的除子》引理
\ref{spaces-divisors-lemma-extend-admissible-blowups}。
于是 $X_{12}' \subset X_i'$ 是开子空间，并且图
$$
\xymatrix{
X_{12}' \ar[d] \ar[r] & X_i' \ar[d] \\
X_{12}^i \ar[r] & X_i^i
}
$$
交换，其中竖直箭头是爆破，水平箭头是开浸入。% P09-ERR-0702
注意 $X'_{12} \to X_1' \times_Y X_2'$ 是浸入且固有
（使用 $X'_{12} \to X_{12}$ 固有、
$X_{12} \to X_1 \times_Y X_2$ 为闭态射，且
$X_1' \times_Y X_2' \to X_1 \times_Y X_2$ 分离，再应用
《空间的态射》引理
\ref{spaces-morphisms-lemma-universally-closed-permanence}）。
因此 $X'_{12} \to  X_1' \times_Y X_2'$ 是闭浸入。
若把 $X$ 定义为将 $X_1'$ 与 $X_2'$ 沿公共开子空间
$X_{12}'$ 黏合所得的空间，则 $X \to Y$ 有限型且分离
\footnote{因为可以检验对角态射 $X \to X \times_Y X$ 在四个开部分
$X'_i \times_Y X'_j$ 上的闭性；这些开部分覆盖 $X \times_Y X$，
而在其上闭性是显然的。}。
由于 $U$-容许爆破的复合仍为 $U$-容许爆破
（《空间上的除子》引理
\ref{spaces-divisors-lemma-composition-admissible-blowups}），
引理得证。
\end{proof}

\begin{lemma}
\label{spaces-more-morphisms-lemma-blowup-to-find-embedding}
设 $S$ 为概形。设 $f : X \to Y$ 为 $S$ 上代数空间的态射。
设 $U \subset X$ 为开子空间。假设
\begin{enumerate}
\item $U$ 拟紧；
\item $Y$ 拟紧且拟分离；
\item 存在浸入 $U \to \mathbf{P}^n_Y$，它在 $Y$ 上；
\item $f$ 有限型且分离。
\end{enumerate}
则存在交换图
$$
\xymatrix{
& U \ar[ld] \ar[d] \ar[rd] \ar[rrd] \\
X \ar[rd] & X' \ar[l] \ar[d] \ar[r] & Z' \ar[ld] \ar[r] & Z \ar[ld] \\
& Y & \mathbf{P}^n_Y \ar[l]
}
$$
其中以 $U$ 为源的箭头是开浸入，
$X' \to X$ 是 $U$-容许爆破，
$X' \to Z'$ 是开浸入，
$Z' \to Y$ 是代数空间的固有可表态射。
更精确地，$Z' \to Z$ 是 $U$-容许爆破，
而 $Z \to \mathbf{P}^n_Y$ 是闭浸入。
\end{lemma}

\begin{proof}
设 $Z \subset \mathbf{P}^n_Y$ 为浸入
$U \to \mathbf{P}^n_Y$ 的概形论像。由于
$U \to \mathbf{P}^n_Y$ 拟紧，可知 $U \subset Z$ 是
（概形论）稠密开子空间（《空间的态射》引理
\ref{spaces-morphisms-lemma-quasi-compact-immersion}）。
应用引理 \ref{spaces-more-morphisms-lemma-find-common-blowups} 得到图
$$
\xymatrix{
X' \ar[d] \ar[r] & \overline{X}' & Z' \ar[l] \ar[d] \\
X & U \ar[l] \ar[lu] \ar[u] \ar[ru] \ar[r] & Z
}
$$
它具有该引理陈述中列出的性质。
由于 $X' \to X$ 与 $Z' \to Z$ 都是 $U$-容许爆破，
可知 $U$ 是 $X'$ 与 $Z'$ 二者的概形论稠密开子空间
（见《空间上的除子》引理
\ref{spaces-divisors-lemma-blowing-up-gives-effective-Cartier-divisor} 和
\ref{spaces-divisors-lemma-complement-effective-Cartier-divisor}）。
由于 $Z' \to Z \to Y$ 固有，可知
$Z' \subset \overline{X}'$ 是闭子空间（见《空间的态射》引理
\ref{spaces-morphisms-lemma-universally-closed-permanence}）。
于是 $X' \subset Z'$（在概形论意义下），因而 $X'$ 是 $Z'$ 的
开子空间（略去一个小细节），引理得证。% P09-ERR-0703
\end{proof}

\begin{lemma}
\label{spaces-more-morphisms-lemma-chow-noetherian}
设 $S$ 为概形。设 $f : X \to Y$ 为 $S$ 上代数空间的态射。
假设 $f$ 分离且有限型，并且 $Y$ 诺特。% P09-ERR-0704
则存在稠密开子空间 $U \subset X$ 以及交换图
$$
\xymatrix{
& U \ar[ld] \ar[d] \ar[rd] \ar[rrd] \\
X \ar[rd] & X' \ar[l] \ar[d] \ar[r] & Z' \ar[ld] \ar[r] & Z \ar[ld] \\
& Y & \mathbf{P}^n_Y \ar[l]
}
$$
其中以 $U$ 为源的箭头是开浸入，
$X' \to X$ 是 $U$-容许爆破，
$X' \to Z'$ 是开浸入，
$Z' \to Y$ 是代数空间的固有可表态射。
更精确地，$Z' \to Z$ 是 $U$-容许爆破，
而 $Z \to \mathbf{P}^n_Y$ 是闭浸入。
\end{lemma}

\begin{proof}
由《空间的极限》引理
\ref{spaces-limits-lemma-embedding-into-affine-over-qs}，存在稠密开子空间
$U \subset X$ 以及浸入 $U \to \mathbf{A}^n_Y$，它在 $Y$ 上。
将其与开浸入 $\mathbf{A}^n_Y \to \mathbf{P}^n_Y$ 复合，
便得到引理 \ref{spaces-more-morphisms-lemma-blowup-to-find-embedding} 的情形，结论随之成立。
\end{proof}

\begin{remark}
\label{spaces-more-morphisms-remark-chow-Noetherian}
在引理 \ref{spaces-more-morphisms-lemma-blowup-to-find-embedding} 和
\ref{spaces-more-morphisms-lemma-chow-noetherian} 中，态射 $g : Z' \to Y$ 是射影态射的复合。
推测（由《态射》引理 \ref{morphisms-lemma-ample-composition}
对代数空间的类似结论），存在一个 $g$-丰沛可逆层，它位于 $Z'$ 上。
若以后用到这一点，我们将在此处陈述并证明它。
\end{remark}

\noindent
下述结果是 \cite[IV Theorem 3.1]{Kn}。注意浸入
$X' \to \mathbf{P}^n_Y$ 拟紧，因而可分解为
$X' \to Z' \to \mathbf{P}^n_Y$，其中第一个态射为开浸入，
第二个态射为闭浸入（《空间的态射》引理
\ref{spaces-morphisms-lemma-quasi-compact-immersion}）。

\begin{lemma}[Chow 引理]
\label{spaces-more-morphisms-lemma-chow-noetherian-separated}
\begin{reference}
\cite[IV Theorem 3.1]{Kn}
\end{reference}
设 $S$ 为概形。设 $f : X \to Y$ 为 $S$ 上代数空间的态射。
假设 $f$ 分离且有限型，并且 $Y$ 分离且诺特。% P09-ERR-0705
则存在交换图
$$
\xymatrix{
X \ar[rd] & X' \ar[l] \ar[d] \ar[r] & \mathbf{P}^n_Y \ar[ld] \\
& Y
}
$$
其中 $X' \to X$ 是 $U$-容许爆破，这里 $U \subset X$
是某个稠密开子空间，并且态射 $X' \to \mathbf{P}^n_Y$ 是浸入。
\end{lemma}

\begin{proof}
在证明的第一段中，我们把引理约化到 $Y$ 在
$\Spec(\mathbf{Z})$ 上有限型的情形。
我们可以并且将把基概形 $S$ 替换为 $\Spec(\mathbf{Z})$。
可将 $Y = \lim Y_i$ 写成 $\Spec(\mathbf{Z})$ 上有限型分离代数空间的
有向逆极限；见《空间的极限》命题
\ref{spaces-limits-proposition-approximate} 和引理
\ref{spaces-limits-lemma-descend-separated}。
对所有充分大的 $i$，可以找到有限型分离态射
$X_i \to Y_i$，使得 $X = Y \times_{Y_i} X_i$；见《空间的极限》引理
\ref{spaces-limits-lemma-descend-finite-presentation} 和
\ref{spaces-limits-lemma-descend-separated-morphism}。
设 $\eta_1, \ldots, \eta_n$ 为 $|X|$ 的不可约分支的一般点
（$X$ 是诺特代数空间 $Y$ 上有限型分离的代数空间，因而诺特，
所以 $|X|$ 是诺特拓扑空间）。
由《空间的极限》引理 \ref{spaces-limits-lemma-topology-limit}，
可知 $\eta_1, \ldots, \eta_n$ 在 $|X_i|$ 中的像彼此不同，
只要 $i$ 足够大。可以把 $X_i$ 替换为态射 $X \to X_i$
（它拟紧，事实上是仿射的）的概形论像。
替换后可知，$\eta_1, \ldots, \eta_n$ 在 $|X_i|$ 中的像是
$|X_i|$ 的不可约分支的一般点；见《空间的态射》引理
\ref{spaces-morphisms-lemma-quasi-compact-scheme-theoretic-image}。
在作出这些说明后，假设能找到图
$$
\xymatrix{
X_i \ar[rd] & X_i' \ar[l] \ar[d] \ar[r] & \mathbf{P}^n_{Y_i} \ar[ld] \\
& Y % P09-ERR-0706
}
$$
其中 $X_i' \to X_i$ 是 $U_i$-容许爆破，这里
$U_i \subset X_i$ 是某个稠密开子空间，并且态射
$X_i' \to \mathbf{P}^n_{Y_i}$ 是浸入。
则严格变换 $X' \to X$，即 $X$ 关于 $X_i' \to X_i$ 的严格变换，
是 $U$-容许爆破，其中 $U \subset X$ 是 $U_i$ 在 $X$ 中的逆像。
由于我们精心选取了指标 $i$，可得
$\eta_1, \ldots, \eta_n \in |U|$，且 $U \subset X$ 稠密。
此外，$X' \to \mathbf{P}^n_Y$ 是浸入，因为 $X'$ 在
$X_i' \times_{X_i} X = X_i' \times_{Y_i} Y$ 中为闭子空间，
而后者带有到 $\mathbf{P}^n_Y$ 的浸入。因此我们已约化到下一段的情形。

\medskip\noindent
假设 $Y$ 在 $\Spec(\mathbf{Z})$ 上分离且有限型。
于是 $X \to \Spec(\mathbf{Z})$ 也分离且有限型。
把引理 \ref{spaces-more-morphisms-lemma-chow-noetherian} 应用于
$X \to \Spec(\mathbf{Z})$，得到稠密开子空间 $U \subset X$
以及交换图
$$
\xymatrix{
& U \ar[ld] \ar[d] \ar[rd] \ar[rrd] \\
X \ar[rd] & X' \ar[l] \ar[d] \ar[r] & Z' \ar[ld] \ar[r] & Z \ar[ld] \\
& \Spec(\mathbf{Z}) & \mathbf{P}^n_\mathbf{Z} \ar[l]
}
$$
它具有该引理中列出的全部性质。
注意 $Z$ 有丰沛可逆层，即
$\mathcal{O}_{\mathbf{P}^n}(1)|_Z$。因此由《态射》引理
\ref{morphisms-lemma-projective-over-quasi-projective-is-H-projective}，
$Z' \to Z$ 是 H-射影态射。
由《态射》引理 \ref{morphisms-lemma-H-projective-composition}，
可知 $Z' \to \Spec(\mathbf{Z})$ 是 H-射影态射。
因此存在闭浸入 $Z' \to \mathbf{P}^m_{\Spec(\mathbf{Z})}$，
其中 $m \geq 0$。
于是对角态射
$$
X' \to Y \times \mathbf{P}^m_\mathbf{Z} = \mathbf{P}^m_Y
$$
是浸入（因为它与到 $\mathbf{P}^m_\mathbf{Z}$ 的投影之复合是浸入），
证明即告完成。
\end{proof}

\section{Chow 引理的变体}
\label{spaces-more-morphisms-section-chows-lemma}

\noindent
本节证明 Chow 引理的若干变体，它们处理非诺特代数空间之间的态射。
相应的诺特版本是引理 \ref{spaces-more-morphisms-lemma-chow-noetherian} 和引理
\ref{spaces-more-morphisms-lemma-chow-noetherian-separated}。

\begin{lemma}
\label{spaces-more-morphisms-lemma-chow-finite-type}
设 $S$ 为概形。设 $Y$ 为 $S$ 上拟紧且拟分离的代数空间。
设 $f : X \to Y$ 为有限型分离态射。则存在交换图
$$
\xymatrix{
X \ar[rd] & X' \ar[l] \ar[d] \ar[r] & \overline{X}' \ar[ld] \\
& Y
}
$$
其中 $X' \to X$ 固有且满，
$X' \to \overline{X}'$ 是开浸入，并且
$\overline{X}' \to Y$ 是代数空间的固有可表态射。% P09-ERR-0707
\end{lemma}

\begin{proof}
由《空间的极限》命题
\ref{spaces-limits-proposition-separated-closed-in-finite-presentation}，
可以找到闭浸入 $X \to X_1$，其中 $X_1$ 在 $Y$ 上分离且有限表示。
显然，若对 $X_1 \to Y$ 证明了断言，则 $X$ 的结论随之成立。
故可假设 $X$ 在 $Y$ 上有限表示。

\medskip\noindent
我们可以并且将把基概形 $S$ 替换为 $\Spec(\mathbf{Z})$。
把 $Y = \lim_i Y_i$ 写成 $\Spec(\mathbf{Z})$ 上有限型拟分离
代数空间的有向逆极限；见《空间的极限》命题
\ref{spaces-limits-proposition-approximate}。
由《空间的极限》引理
\ref{spaces-limits-lemma-descend-finite-presentation}，可以找到指标
$i \in I$ 以及有限表示的概形 $X_i \to Y_i$，% P09-ERR-0708
使得 $X = Y \times_{Y_i} X_i$。
由《空间的极限》引理
\ref{spaces-limits-lemma-descend-separated-morphism}，
可以假设 $X_i \to Y_i$ 分离。
显然，若对 $X_i$（在 $Y_i$ 上）证明了断言，则 $X$ 的断言也成立。
$X_i \to Y_i$ 的情形由引理 \ref{spaces-more-morphisms-lemma-chow-noetherian} 处理。
\end{proof}

\begin{lemma}
\label{spaces-more-morphisms-lemma-chow-finite-type-separated}
设 $S$ 为概形。设 $f : X \to Y$ 为 $S$ 上代数空间的态射。
假设 $f$ 分离且有限型，并且 $Y$ 分离且拟紧。% P09-ERR-0709
则存在交换图
$$
\xymatrix{
X \ar[rd] & X' \ar[l] \ar[d] \ar[r] & \mathbf{P}^n_Y \ar[ld] \\
& Y
}
$$
其中 $X' \to X$ 是固有满态射，% P09-ERR-0710
并且态射 $X' \to \mathbf{P}^n_Y$ 是浸入。
\end{lemma}

\begin{proof}
由《空间的极限》命题
\ref{spaces-limits-proposition-separated-closed-in-finite-presentation}，
可以找到闭浸入 $X \to X_1$，其中 $X_1$ 在 $Y$ 上分离且有限表示。
显然，若对 $X_1 \to Y$ 证明了断言，则 $X$ 的结论随之成立。
故可假设 $X$ 在 $Y$ 上有限表示。

\medskip\noindent
我们可以并且将把基概形 $S$ 替换为 $\Spec(\mathbf{Z})$。
把 $Y = \lim_i Y_i$ 写成 $\Spec(\mathbf{Z})$ 上有限型拟分离
代数空间的有向逆极限；见《空间的极限》命题
\ref{spaces-limits-proposition-approximate}。
由《空间的极限》引理 \ref{spaces-limits-lemma-descend-separated}，
可以假设 $Y_i$ 对所有 $i$ 都分离。
由《空间的极限》引理
\ref{spaces-limits-lemma-descend-finite-presentation}，可以找到指标
$i \in I$ 以及有限表示的概形 $X_i \to Y_i$，% P09-ERR-0711
使得 $X = Y \times_{Y_i} X_i$。
由《空间的极限》引理
\ref{spaces-limits-lemma-descend-separated-morphism}，
可以假设 $X_i \to Y_i$ 分离。
显然，若对 $X_i$（在 $Y_i$ 上）证明了断言，则 $X$ 的断言也成立。
$X_i \to Y_i$ 的情形由引理
\ref{spaces-more-morphisms-lemma-chow-noetherian-separated} 处理。
\end{proof}

\section{Grothendieck 存在定理}
\label{spaces-more-morphisms-section-existence-proper}

\noindent
本节讨论代数空间的 Grothendieck 存在定理。我们不发展“形式代数空间”
的理论，而是暂时建立一些术语，以替代“诺特 adic 形式空间上的凝聚模”
这一概念。

\medskip\noindent
设 $S$ 为概形。设 $X$ 为 $S$ 上的诺特代数空间。
设 $\mathcal{I} \subset \mathcal{O}_X$ 为拟凝聚理想层。
下面考虑逆系统 $(\mathcal{F}_n)$，其各项为凝聚
$\mathcal{O}_X$-模，并且满足
\begin{enumerate}
\item $\mathcal{F}_n$ 被 $\mathcal{I}^n$ 零化；
\item 转移态射诱导同构
$\mathcal{F}_{n + 1}/\mathcal{I}^n\mathcal{F}_{n + 1} \to \mathcal{F}_n$。
\end{enumerate}
这类逆系统之间的态射
$\alpha : (\mathcal{F}_n) \to (\mathcal{G}_n)$，就是相容的态射系统
$\alpha_n : \mathcal{F}_n \to \mathcal{G}_n$。
把这些逆系统所成的范畴记为 $\textit{Coh}(X, \mathcal{I})$。
我们将建立关于这些系统的一些理论，它与概形情形中的相应结果平行；
见《概形的上同调》第 \ref{coherent-section-existence}、
\ref{coherent-section-existence-proper}、
\ref{coherent-section-existence-proper-support} 和
\ref{coherent-section-algebraization} 节。% P09-ERR-0712

\medskip\noindent
函子性。设 $f : X \to Y$ 为某概形 $S$ 上诺特代数空间的态射，
并设 $\mathcal{J} \subset \mathcal{O}_Y$ 为拟凝聚理想层。
令 $\mathcal{I} = f^{-1}\mathcal{J}\mathcal{O}_X$。
在此情形下有函子
$$
f^* : \textit{Coh}(Y, \mathcal{J}) \longrightarrow \textit{Coh}(X, \mathcal{I})
$$
它把 $(\mathcal{G}_n)$ 映到 $(f^*\mathcal{G}_n)$。
参见《概形的上同调》引理
\ref{coherent-lemma-inverse-systems-pullback}。
若 $f$ 平展，则可把它简单地看作将系统限制到 $X$；
见《空间的性质》公式
\ref{spaces-properties-equation-restrict-modules}。

\medskip\noindent
平展下降。设 $S$ 为概形。设 $U_0 \to X$ 为诺特代数空间的满平展态射。
令 $U_1 = U_0 \times_X U_0$ 以及
$U_2 = U_0 \times_X U_0 \times_X U_0$。
设 $\mathcal{I} \subset \mathcal{O}_{X}$ 为拟凝聚理想层。
令 $\mathcal{I}_i = \mathcal{I}|_{U_i}$。在此情形下得到范畴图
$$
\xymatrix{
\textit{Coh}(X, \mathcal{I}) \ar[r] &
\textit{Coh}(U_0, \mathcal{I}_0) \ar@<0.5ex>[r] \ar@<-0.5ex>[r] &
\textit{Coh}(U_1, \mathcal{I}_1) \ar@<1ex>[r] \ar[r] \ar@<-1ex>[r] &
\textit{Coh}(U_2, \mathcal{I}_2)
}
$$
并且第一个箭头把 $\textit{Coh}(X, \mathcal{I})$ 展示为该图右部的
同伦极限。% P09-ERR-0713
更精确地，给定一个{\it 下降资料}，即一对
$((\mathcal{G}_n), \varphi)$，其中 $(\mathcal{G}_n)$ 是
$\textit{Coh}(U_0, \mathcal{I}_0)$ 的对象，且
$\varphi : \text{pr}_0^*(\mathcal{G}_n) \to \text{pr}_1^*(\mathcal{G}_n)$
是 $\textit{Coh}(U_1, \mathcal{I}_1)$ 中的同构，并在
$\textit{Coh}(U_2, \mathcal{I}_2)$ 中满足余圈条件，则存在唯一对象
$(\mathcal{F}_n)$ 属于 $\textit{Coh}(X, \mathcal{I})$，
其相伴的典范下降资料同构于 $((\mathcal{G}_n), \varphi)$。
参见《空间上的下降》定义
\ref{spaces-descent-definition-descent-datum-effective-quasi-coherent}。
对下降资料 $(\mathcal{G}_n, \varphi_n)$（其中 $n$ 变化）应用
《空间上的下降》命题
\ref{spaces-descent-proposition-fpqc-descent-quasi-coherent}，
便立即得到这一陈述的证明。

\begin{lemma}
\label{spaces-more-morphisms-lemma-inverse-systems-abelian}
设 $S$ 为概形。设 $X$ 为 $S$ 上的诺特代数空间，
并设 $\mathcal{I} \subset \mathcal{O}_X$ 为拟凝聚理想层。
\begin{enumerate}
\item 范畴 $\textit{Coh}(X, \mathcal{I})$ 是阿贝尔范畴。
\item $\textit{Coh}(X, \mathcal{I})$ 中的正合性可在平展局部检验。
\item 对诺特代数空间的任意平坦态射 $f : X' \to X$，函子
$f^* : \textit{Coh}(X, \mathcal{I}) \to
\textit{Coh}(X', f^{-1}\mathcal{I}\mathcal{O}_{X'})$ 是正合的。
\end{enumerate}
\end{lemma}

\begin{proof}
先证 (1)。选择仿射概形 $U_0$ 以及满平展态射
$U_0 \to X$。如同上述平展下降的讨论，令
$U_1 = U_0 \times_X U_0$ 以及
$U_2 = U_0 \times_X U_0 \times_X U_0$。
各范畴 $\textit{Coh}(U_i, \mathcal{I}_i)$ 都是阿贝尔范畴
（《概形的上同调》引理
\ref{coherent-lemma-inverse-systems-abelian}），并且拉回函子
$\textit{Coh}(U_0, \mathcal{I}_0) \to \textit{Coh}(U_1, \mathcal{I}_1)$
以及
$\textit{Coh}(U_1, \mathcal{I}_1) \to \textit{Coh}(U_2, \mathcal{I}_2)$
都是正合函子（《概形的上同调》引理
\ref{coherent-lemma-inverse-systems-pullback}）。
于是，由把 $\textit{Coh}(X, \mathcal{I})$ 描述成下降资料范畴，
形式地得到引理。略去一些细节；参见《群胚》引理
\ref{groupoids-lemma-abelian} 的证明。% P09-ERR-0714

\medskip\noindent
第 (2) 部分立即由上一段的讨论得到。在 (3) 的情形中，
选择交换图
$$
\xymatrix{
U' \ar[d] \ar[r] & U \ar[d] \\
X' \ar[r] & X
}
$$
其中 $U'$ 与 $U$ 是仿射概形，竖直态射满且平展。
于是 $U' \to U$ 是诺特概形的平坦态射（《空间的态射》引理
\ref{spaces-morphisms-lemma-flat-local}），故拉回函子
$\textit{Coh}(U, \mathcal{I}\mathcal{O}_U) \to
\textit{Coh}(U', \mathcal{I}\mathcal{O}_{U'})$
由《概形的上同调》引理
\ref{coherent-lemma-inverse-systems-pullback} 而正合。
由于可以检验 $\textit{Coh}(X, \mathcal{O}_X)$ 中的正合性，
而这一检验可在 $U$ 上进行，并且对 $X', U'$ 也类似，
断言成立。% P09-ERR-0715
\end{proof}

\begin{lemma}
\label{spaces-more-morphisms-lemma-inverse-systems-surjective}
设 $S$ 为概形。设 $X$ 为 $S$ 上的诺特代数空间，
并设 $\mathcal{I} \subset \mathcal{O}_X$ 为拟凝聚理想层。
态射 $(\mathcal{F}_n) \to (\mathcal{G}_n)$ 在
$\textit{Coh}(X, \mathcal{I})$ 中为满，当且仅当
$\mathcal{F}_1 \to \mathcal{G}_1$ 为满。
\end{lemma}

\begin{proof}
由引理 \ref{spaces-more-morphisms-lemma-inverse-systems-abelian}，可以在 $X$ 的某个仿射
平展覆盖上检验。因而问题约化到概形情形，即《概形的上同调》引理
\ref{coherent-lemma-inverse-systems-surjective}。
\end{proof}

\noindent
设 $S$ 为概形。设 $X$ 为 $S$ 上的诺特代数空间，
并设 $\mathcal{I} \subset \mathcal{O}_X$ 为拟凝聚理想层。
存在函子
\begin{equation}
\label{spaces-more-morphisms-equation-completion-functor}
\textit{Coh}(\mathcal{O}_X) \longrightarrow \textit{Coh}(X, \mathcal{I}), \quad
\mathcal{F} \longmapsto \mathcal{F}^\wedge
\end{equation}
它把凝聚 $\mathcal{O}_X$-模 $\mathcal{F}$ 映到对象
$\mathcal{F}^\wedge = (\mathcal{F}/\mathcal{I}^n\mathcal{F})$，
该对象属于 $\textit{Coh}(X, \mathcal{I})$。

\begin{lemma}
\label{spaces-more-morphisms-lemma-exact}
函子 (\ref{spaces-more-morphisms-equation-completion-functor}) 是正合的。
\end{lemma}

\begin{proof}
只须在 $X$ 上平展局部检验这一点；见引理
\ref{spaces-more-morphisms-lemma-inverse-systems-abelian}。
因而问题约化到概形情形，即《概形的上同调》引理
\ref{coherent-lemma-exact}。
\end{proof}

\begin{lemma}
\label{spaces-more-morphisms-lemma-completion-internal-hom}
设 $S$ 为概形。设 $X$ 为 $S$ 上的诺特代数空间，
并设 $\mathcal{I} \subset \mathcal{O}_X$ 为拟凝聚理想层。
设 $\mathcal{F}$、$\mathcal{G}$ 为凝聚 $\mathcal{O}_X$-模。令
$\mathcal{H} = \SheafHom_{\mathcal{O}_X}(\mathcal{F}, \mathcal{G})$。
则
$$
\lim H^0(X, \mathcal{H}/\mathcal{I}^n\mathcal{H}) =
\Mor_{\textit{Coh}(X, \mathcal{I})}
(\mathcal{F}^\wedge, \mathcal{G}^\wedge).
$$
\end{lemma}

\begin{proof}
由于 $\mathcal{H}$ 是 $X_\etale$ 上的层，而且
$\textit{Coh}(X, \mathcal{I})$ 的对象具有平展下降，
只须在平展局部证明。因而问题约化到概形情形，即《概形的上同调》引理
\ref{coherent-lemma-completion-internal-hom}。
\end{proof}

\noindent
现在引入本节其余部分始终关注的情形。

\begin{situation}
\label{spaces-more-morphisms-situation-existence}
这里 $A$ 是关于理想 $I$ 完备的诺特环。此外，
$f : X \to \Spec(A)$ 是代数空间的有限型分离态射，且
$\mathcal{I} = I\mathcal{O}_X$。
\end{situation}

\noindent
在此情形中，记
$$
\textit{Coh}_{\text{支撑固有于 } A}(\mathcal{O}_X)
$$
为 $\textit{Coh}(\mathcal{O}_X)$ 的全子范畴，% P09-ERR-0716
其对象是凝聚 $\mathcal{O}_X$-模，而且其支撑在 $\Spec(A)$ 上固有；
等价地，其概形论支撑在 $\Spec(A)$ 上固有。见
《空间的导出范畴》引理
\ref{spaces-perfect-lemma-module-support-proper-over-base}。
类似地，令
$$
\textit{Coh}_{\text{支撑固有于 } A}(X, \mathcal{I})
$$
为 $\textit{Coh}(X, \mathcal{I})$ 的全子范畴，
它由满足下述条件的对象 $(\mathcal{F}_n)$ 组成：
$\mathcal{F}_1$ 的支撑在 $\Spec(A)$ 上固有。
由于商模的支撑包含在原模的支撑中，函子
(\ref{spaces-more-morphisms-equation-completion-functor}) 诱导函子
\begin{equation}
\label{spaces-more-morphisms-equation-completion-functor-proper-over-A}
\textit{Coh}_{\text{支撑固有于 }A}(\mathcal{O}_X)
\longrightarrow
\textit{Coh}_{\text{支撑固有于 }A}(X, \mathcal{I})
\end{equation}
第一个结果是此函子全忠实。

\begin{lemma}
\label{spaces-more-morphisms-lemma-fully-faithful}
处于情形 \ref{spaces-more-morphisms-situation-existence}。
设 $\mathcal{F}$、$\mathcal{G}$ 为凝聚 $\mathcal{O}_X$-模。
假设 $\mathcal{F}$ 与 $\mathcal{G}$ 的支撑之交在
$\Spec(A)$ 上固有。则由 (\ref{spaces-more-morphisms-equation-completion-functor}) 得到的映射
$$
\Mor_{\textit{Coh}(\mathcal{O}_X)}(\mathcal{F}, \mathcal{G})
\longrightarrow
\Mor_{\textit{Coh}(X, \mathcal{I})}
(\mathcal{F}^\wedge, \mathcal{G}^\wedge)
$$
是双射。特别地，(\ref{spaces-more-morphisms-equation-completion-functor-proper-over-A})
全忠实。
\end{lemma}

\begin{proof}
令 $\mathcal{H} = \SheafHom_{\mathcal{O}_X}(\mathcal{G}, \mathcal{F})$。
这是凝聚 $\mathcal{O}_X$-模，因为它限制到平展映到 $X$ 的概形后，
由《模》引理
\ref{modules-lemma-internal-hom-locally-kernel-direct-sum} 是凝聚的。
% P09-ERR-0717 P09-ERR-0718
由引理 \ref{spaces-more-morphisms-lemma-completion-internal-hom}，映射
$$
\lim_n H^0(X, \mathcal{H}/\mathcal{I}^n\mathcal{H})
\to
\Mor_{\textit{Coh}(X, \mathcal{I})}
(\mathcal{G}^\wedge, \mathcal{F}^\wedge)
$$
是双射。设 $i : Z \to X$ 为 $\mathcal{H}$ 的概形论支撑。
显然，$Z$ 是闭子空间，且 $|Z|$ 包含在 $\mathcal{F}$ 与
$\mathcal{G}$ 的支撑之交中。因此由假设，$Z \to \Spec(A)$ 固有
（见《空间的导出范畴》第
\ref{spaces-perfect-section-proper-over-base} 节）。
写成 $\mathcal{H} = i_*\mathcal{H}'$，其中
$\mathcal{O}_Z$-模 $\mathcal{H}'$ 是凝聚的。我们有
$i_*(\mathcal{H}'/I^n\mathcal{H}') = \mathcal{H}/I^n\mathcal{H}$。
因此得到
\begin{align*}
\lim_n H^0(X, \mathcal{H}/\mathcal{I}^n\mathcal{H})
& =
\lim_n H^0(Z, \mathcal{H}'/\mathcal{I}^n\mathcal{H}') \\
& =
H^0(Z, \mathcal{H}') \\
& =
H^0(X, \mathcal{H}) \\
& = 
\Mor_{\textit{Coh}(\mathcal{O}_X)}(\mathcal{F}, \mathcal{G})
\end{align*}
其中第二个等式由形式函数定理得到（《空间的上同调》引理
\ref{spaces-cohomology-lemma-spell-out-theorem-formal-functions}）。
这就证明了引理。
\end{proof}

\begin{remark}
\label{spaces-more-morphisms-remark-inverse-systems-kernel-cokernel-annihilated-by}
设 $S$ 为概形。设 $X$ 为 $S$ 上的诺特代数空间，并设
$\mathcal{I}, \mathcal{K} \subset \mathcal{O}_X$ 为拟凝聚理想层。
设 $\alpha : (\mathcal{F}_n) \to (\mathcal{G}_n)$ 为
$\textit{Coh}(X, \mathcal{I})$ 的态射。
给定仿射概形 $U = \Spec(A)$ 以及满平展态射 $U \to X$，
以 $I, K \subset A$ 表示与限制
$\mathcal{I}|_U, \mathcal{K}|_U$ 对应的理想。
以 $\alpha_U : M \to N$ 表示有限 $A^\wedge$-模之间的态射，
它经由《概形的上同调》引理
\ref{coherent-lemma-inverse-systems-affine} 对应于 $\alpha|_U$。
% P09-ERR-0719
我们断言下列条件等价：
\begin{enumerate}
\item 存在整数 $t \geq 1$，使得 $\Ker(\alpha_n)$ 与
$\Coker(\alpha_n)$ 都被 $\mathcal{K}^t$ 零化，对所有 $n \geq 1$；
\item 对任意（或某个）如上的仿射开集
$\Spec(A) = U \subset X$，模 $\Ker(\alpha_U)$ 与
$\Coker(\alpha_U)$ 都被 $K^t$ 零化，其中某个整数 $t \geq 1$。
% P09-ERR-0720
\end{enumerate}
若这些等价条件成立，就称 $\alpha$ 是一个
{\it 其核与余核被 $\mathcal{K}$ 的某次幂零化的映射}。
关于等价性的证明，参见《概形的上同调》评注
\ref{coherent-remark-inverse-systems-kernel-cokernel-annihilated-by}。
\end{remark}

\begin{lemma}
\label{spaces-more-morphisms-lemma-existence-easy}
设 $S$ 为概形。设 $X$ 为 $S$ 上的诺特代数空间，并设
$\mathcal{I} \subset \mathcal{O}_X$ 为拟凝聚理想层。
设 $\mathcal{G}$ 为凝聚 $\mathcal{O}_X$-模，$(\mathcal{F}_n)$ 为
$\textit{Coh}(X, \mathcal{I})$ 的对象，并设
$\alpha : (\mathcal{F}_n) \to \mathcal{G}^\wedge$ 是一个其核与余核
被 $\mathcal{I}$ 的某次幂零化的映射。
则存在唯一的（精确到唯一同构）三元组
$(\mathcal{F}, a, \beta)$，其中
\begin{enumerate}
\item $\mathcal{F}$ 是凝聚 $\mathcal{O}_X$-模；
\item $a : \mathcal{F} \to \mathcal{G}$ 是 $\mathcal{O}_X$-模态射，
其核与余核被 $\mathcal{I}$ 的某次幂零化；
\item $\beta : (\mathcal{F}_n) \to \mathcal{F}^\wedge$ 是同构；
\item $\alpha = a^\wedge \circ \beta$。
\end{enumerate}
\end{lemma}

\begin{proof}
$\textit{Coh}(X, \mathcal{I})$ 和 $\textit{Coh}(\mathcal{O}_X)$
的对象具有唯一性和平展下降，这蕴含只须构造
$(\mathcal{F}, a, \beta)$，而这一构造可在 $X$ 上平展局部进行。
% P09-ERR-0721
因而问题约化到概形情形，即《概形的上同调》引理
\ref{coherent-lemma-existence-easy}。
\end{proof}

\begin{lemma}
\label{spaces-more-morphisms-lemma-existence-tricky}
处于情形 \ref{spaces-more-morphisms-situation-existence}。设
$\mathcal{K} \subset \mathcal{O}_X$ 为拟凝聚理想层。
设 $X_e \subset X$ 为由 $\mathcal{K}^e$ 切出的闭子空间。
令 $\mathcal{I}_e = \mathcal{I}\mathcal{O}_{X_e}$。
设 $(\mathcal{F}_n)$ 为
$\textit{Coh}_{\text{支撑固有于 } A}(X, \mathcal{I})$ 的对象。假设
\begin{enumerate}
\item 函子
$\textit{Coh}_{\text{支撑固有于 } A}(\mathcal{O}_{X_e})
\to \textit{Coh}_{\text{支撑固有于 } A}(X_e, \mathcal{I}_e)$
对所有 $e \geq 1$ 都是范畴等价；
\item 存在对象 $\mathcal{H}$，它属于
$\textit{Coh}_{\text{支撑固有于 } A}(\mathcal{O}_X)$，以及映射
$\alpha : (\mathcal{F}_n) \to \mathcal{H}^\wedge$，其核与余核
被 $\mathcal{K}$ 的某次幂零化。
\end{enumerate}
则 $(\mathcal{F}_n)$ 属于
(\ref{spaces-more-morphisms-equation-completion-functor-proper-over-A}) 的本质像。
\end{lemma}

\begin{proof}
在本证明中，我们将不再另行说明地使用下述事实：对闭浸入
$i : Z \to X$，函子 $i_*$ 给出 $Z$ 上凝聚模的范畴与 $X$ 上
被 $Z$ 的理想层零化的凝聚模范畴之间的等价；见
《空间的上同调》引理
\ref{spaces-cohomology-lemma-i-star-equivalence}。
特别地，我们把
$$
\textit{Coh}_{\text{支撑固有于 } A}(\mathcal{O}_{X_e})
\subset
\textit{Coh}_{\text{支撑固有于 } A}(\mathcal{O}_X)
$$
看成由被 $\mathcal{K}^e$ 零化的模组成的全子范畴，并把
$$
\textit{Coh}_{\text{支撑固有于 } A}(X_e, \mathcal{I}_e)
\subset
\textit{Coh}_{\text{支撑固有于 } A}(X, \mathcal{I})
$$
看成由被 $\mathcal{K}^e$ 零化的对象组成的全子范畴。% P09-ERR-0722
此外，(1) 告诉我们，这两个范畴在完备化函子
(\ref{spaces-more-morphisms-equation-completion-functor-proper-over-A}) 下等价。

\medskip\noindent
应用此等价，得到凝聚 $\mathcal{O}_X$-模 $\mathcal{G}_e$；
它被 $\mathcal{K}^e$ 零化，并对应于系统
$(\mathcal{F}_n/\mathcal{K}^e\mathcal{F}_n)$；该系统属于
$\textit{Coh}_{\text{支撑固有于 } A}(X, \mathcal{I})$。态射
$\mathcal{F}_n/\mathcal{K}^{e + 1}\mathcal{F}_n \to
\mathcal{F}_n/\mathcal{K}^e\mathcal{F}_n$ 对应于典范态射
$\mathcal{G}_{e + 1} \to \mathcal{G}_e$，后者诱导同构
$\mathcal{G}_{e + 1}/\mathcal{K}^e\mathcal{G}_{e + 1} \to \mathcal{G}_e$。
我们得到对象 $(\mathcal{G}_e)$，它属于范畴
$\textit{Coh}_{\text{支撑固有于 } A}(X, \mathcal{K})$。
映射 $\alpha$ 诱导态射系统
$$
\mathcal{F}_n/\mathcal{K}^e\mathcal{F}_n
\longrightarrow
\mathcal{H}/(\mathcal{I}^n + \mathcal{K}^e)\mathcal{H}
$$
从而再次由范畴等价得到态射
$\mathcal{G}_e \to \mathcal{H}/\mathcal{K}^e\mathcal{H}$。
由假设 (2)，存在整数 $t \geq 1$，使得 $\mathcal{K}^t$ 零化所有态射
$\mathcal{F}_n \to \mathcal{H}/\mathcal{I}^n\mathcal{H}$ 的核与余核。
于是 $\mathcal{K}^{2t}$ 零化态射
$\mathcal{F}_n/\mathcal{K}^e\mathcal{F}_n \to
\mathcal{H}/(\mathcal{I}^n + \mathcal{K}^e)\mathcal{H}$
的核与余核（略去细节；见《概形的上同调》评注
\ref{coherent-remark-inverse-systems-kernel-cokernel-annihilated-by}）。
进而可知 $\mathcal{K}^{4t}$ 零化下列态射的核与余核：
$$
\mathcal{G}_e
\longrightarrow
\mathcal{H}/\mathcal{K}^e\mathcal{H},
$$
（略去细节；见《概形的上同调》评注
\ref{coherent-remark-inverse-systems-kernel-cokernel-annihilated-by}）。
% P09-ERR-0723
应用引理 \ref{spaces-more-morphisms-lemma-existence-easy}，得到凝聚
$\mathcal{O}_X$-模 $\mathcal{F}$、态射
$a : \mathcal{F} \to \mathcal{H}$ 以及同构
$\beta : (\mathcal{G}_e) \to (\mathcal{F}/\mathcal{K}^e\mathcal{F})$，
后者位于 $\textit{Coh}(X, \mathcal{K})$ 中。逆向追溯可知，
对给定的 $n$，三元组
$(\mathcal{F}/\mathcal{I}^n\mathcal{F}, a \bmod \mathcal{I}^n, \beta
\bmod \mathcal{I}^n)$ 是上述引理中对应于态射
$\alpha_n \bmod \mathcal{K}^e :
(\mathcal{F}_n/\mathcal{K}^e\mathcal{F}_n) \to
(\mathcal{H}/(\mathcal{I}^n + \mathcal{K}^e)\mathcal{H})$
的三元组，该态射位于 $\textit{Coh}(X, \mathcal{K})$ 中。
因此，引理 \ref{spaces-more-morphisms-lemma-existence-easy} 中的唯一性给出典范同构
$\mathcal{F}/\mathcal{I}^n\mathcal{F} \to \mathcal{F}_n$，
它与这里涉及的所有态射相容。

\medskip\noindent
为完成引理的证明，还须说明 $\mathcal{F}$ 的支撑在 $A$ 上固有。
按构造，$a : \mathcal{F} \to \mathcal{H}$ 的核被
$\mathcal{K}$ 的某次幂零化。因此，此核的支撑包含在
$\mathcal{G}_1$ 的支撑中。由于 $\mathcal{G}_1$ 是
$\textit{Coh}_{\text{支撑固有于 } A}(\mathcal{O}_{X_1})$ 的对象，
可知此支撑在 $A$ 上固有。再结合 $\mathcal{H}$ 的支撑在 $A$
上固有这一事实，由《空间的导出范畴》引理
\ref{spaces-perfect-lemma-union-closed-proper-over-base}，
可知 $\mathcal{F}$ 的支撑在 $A$ 上固有。
\end{proof}

\begin{lemma}
\label{spaces-more-morphisms-lemma-inverse-systems-push-pull}
设 $S$ 为概形。设 $f : X \to Y$ 为 $S$ 上诺特代数空间的
可表固有态射。设
$\mathcal{J}, \mathcal{K} \subset \mathcal{O}_Y$ 为拟凝聚理想层。
假设 $f$ 在 $V = Y \setminus V(\mathcal{K})$ 上是同构。
令 $\mathcal{I} = f^{-1}\mathcal{J} \mathcal{O}_X$。
设 $(\mathcal{G}_n)$ 为 $\textit{Coh}(Y, \mathcal{J})$ 的对象，
设 $\mathcal{F}$ 为凝聚 $\mathcal{O}_X$-模，并设
$\beta : (f^*\mathcal{G}_n)  \to \mathcal{F}^\wedge$ 是
$\textit{Coh}(X, \mathcal{I})$ 中的同构。则存在映射
$$
\alpha :
(\mathcal{G}_n)
\longrightarrow
(f_*\mathcal{F})^\wedge
$$
它位于 $\textit{Coh}(Y, \mathcal{J})$ 中，并且其核与余核
被 $\mathcal{K}$ 的某次幂零化。
\end{lemma}

\begin{proof}
由于 $f$ 是固有态射，$f_*\mathcal{F}$ 是凝聚
$\mathcal{O}_Y$-模（《空间的上同调》引理
\ref{spaces-cohomology-lemma-proper-pushforward-coherent}）。
故引理的陈述有意义。考虑复合
$$
\gamma_n : \mathcal{G}_n \to
f_*f^*\mathcal{G}_n \to
f_*(\mathcal{F}/\mathcal{I}^n\mathcal{F}).
$$
这里第一个映射是伴随映射，第二个是 $f_*\beta_n$。
我们断言，存在唯一的引理所述 $\alpha$，使得复合
$$
\mathcal{G}_n \xrightarrow{\alpha_n}
f_*\mathcal{F}/\mathcal{J}^nf_*\mathcal{F} \to
f_*(\mathcal{F}/\mathcal{I}^n\mathcal{F})
$$
都等于 $\gamma_n$，对所有 $n$。由唯一性以及
$\textit{Coh}(Y, \mathcal{J})$ 的平展下降，只须在 $Y$ 上
平展局部证明。故可假设 $Y$ 是某个诺特环的谱。
由于 $f$ 可表，可知 $X$ 也是概形。这样便约化到概形情形；
见《概形的上同调》引理
\ref{coherent-lemma-inverse-systems-push-pull} 的证明。
\end{proof}

\begin{theorem}[Grothendieck 存在定理]
\label{spaces-more-morphisms-theorem-grothendieck-existence}
在情形 \ref{spaces-more-morphisms-situation-existence} 中，函子
(\ref{spaces-more-morphisms-equation-completion-functor-proper-over-A}) 是范畴等价。
\end{theorem}

\begin{proof}
在定理的证明中，我们将不再另行说明地使用《空间的上同调》引理
\ref{spaces-cohomology-lemma-i-star-equivalence} 的范畴等价。
由引理 \ref{spaces-more-morphisms-lemma-fully-faithful}，该函子全忠实。
因此只须证明该函子本质满。

\medskip\noindent
考虑集合 $\Xi$；其元素为拟凝聚理想层
$\mathcal{K} \subset \mathcal{O}_X$，并满足：断言对每个对象
$(\mathcal{F}_n)$ 都成立，其中该对象属于
$\textit{Coh}_{\text{支撑固有于 }A}(X, \mathcal{I})$，并被
$\mathcal{K}$ 零化。
我们要证明 $(0)$ 属于 $\Xi$。若不然，由于 $X$ 诺特，存在
极大拟凝聚理想层 $\mathcal{K}$，它不属于 $\Xi$；见《空间的上同调》引理
\ref{spaces-cohomology-lemma-acc-coherent}。
把 $X$ 替换为 $X$ 中与 $\mathcal{K}$ 对应的闭子概形后，% P09-ERR-0724
可以假设每个非零 $\mathcal{K}$ 都属于 $\Xi$。
设 $(\mathcal{F}_n)$ 为
$\textit{Coh}_{\text{支撑固有于 }A}(X, \mathcal{I})$ 的对象。
我们将证明此对象属于本质像，从而完成定理的证明。

\medskip\noindent
应用 Chow 引理（引理 \ref{spaces-more-morphisms-lemma-chow-noetherian-separated}），
找到固有满态射 $f : Y \to X$；它在某个稠密开子空间
$U \subset X$ 上是同构，并且 $Y$ 在 $A$ 上 H-拟射影。
注意 $Y$ 是概形，且 $f$ 可表。选取开浸入
$j : Y \to Y'$，其中 $Y'$ 在 $A$ 上射影；见《态射》引理
\ref{morphisms-lemma-H-quasi-projective-open-H-projective}。
设 $T_n$ 为 $\mathcal{F}_n$ 的概形论支撑。
注意 $|T_n| = |T_1|$，故 $T_n$ 在 $A$ 上固有，对所有 $n$
（《空间的态射》引理
\ref{spaces-morphisms-lemma-image-proper-is-proper}）。
于是 $f^*\mathcal{F}_n$ 支撑在闭子概形 $f^{-1}T_n$ 上；
后者在 $A$ 上固有（由《空间的态射》引理
\ref{spaces-morphisms-lemma-composition-proper} 以及 $f$ 的固有性）。
特别地，复合 $f^{-1}T_n \to Y \to Y'$ 是闭态射
（《态射》引理 \ref{morphisms-lemma-image-proper-scheme-closed}）。
设 $T'_n \subset Y'$ 为相应的闭子概形；它包含在开子概形 $Y$ 中，
并且等于 $f^{-1}T_n$，这里把二者视为 $Y$ 的闭子概形。
设 $\mathcal{F}_n'$ 为凝聚 $\mathcal{O}_{Y'}$-模，它对应于
$f^*\mathcal{F}_n$；后者被视为 $Y'$ 上的凝聚模，所用闭浸入为
$f^{-1}T_n = T'_n \subset Y'$。
于是 $(\mathcal{F}_n')$ 是
$\textit{Coh}(Y', I\mathcal{O}_{Y'})$ 的对象。
由射影情形的 Grothendieck 存在定理（《概形的上同调》引理
\ref{coherent-lemma-existence-projective}），存在凝聚
$\mathcal{O}_{Y'}$-模 $\mathcal{F}'$ 以及同构
$(\mathcal{F}')^\wedge \cong (\mathcal{F}'_n)$，
该同构位于 $\textit{Coh}(Y', I\mathcal{O}_{Y'})$ 中。
设 $Z' \subset Y'$ 为 $\mathcal{F}'$ 的概形论支撑。
由于 $\mathcal{F}'/I\mathcal{F}' = \mathcal{F}'_1$，可知
$Z' \cap V(I\mathcal{O}_{Y'}) = T'_1$（作为点集）。
结构态射 $p' : Y' \to \Spec(A)$ 固有，故
$p'(Z' \cap (Y' \setminus Y))$ 在 $\Spec(A)$ 中闭。
若它非空，则它会包含 $V(I)$ 的一点，因为 $I$ 包含在 $A$ 的
Jacobson 根中（《代数》引理
\ref{algebra-lemma-radical-completion}）。
但上面已经看到
$Z' \cap (p')^{-1}V(I) = T'_1 \subset Y$，故可知
$Z' \subset Y$。于是 $\mathcal{F}'|_Y$ 支撑在 $Y$ 的某个
闭子概形上，而该闭子概形在 $A$ 上固有。

\medskip\noindent
设 $\mathcal{K}$ 为切出约化补集 $X \setminus U$ 的拟凝聚理想层。
由《空间的上同调》引理
\ref{spaces-cohomology-lemma-proper-pushforward-coherent}，
$\mathcal{O}_X$-模 $\mathcal{H} = f_*\mathcal{F}'$ 是凝聚的；
并且由引理 \ref{spaces-more-morphisms-lemma-inverse-systems-push-pull}，存在态射
$\alpha : (\mathcal{F}_n) \to \mathcal{H}^\wedge$；它位于范畴
$\textit{Coh}_{\text{支撑固有于 } A}(X, \mathcal{I})$ 中，且其核与余核
被 $\mathcal{K}$ 的某次幂零化。% P09-ERR-0725
设 $Z_0 \subset X$ 为 $\mathcal{H}$ 的概形论支撑。
显然 $|Z_0| \subset f(|Z'|)$。因此
$Z_0 \to \Spec(A)$ 固有（《空间的态射》引理
\ref{spaces-morphisms-lemma-image-proper-is-proper}）。
所以 $\mathcal{H}$ 是
$\textit{Coh}_{\text{支撑固有于 } A}(\mathcal{O}_X)$ 的对象。
由于每个理想层 $\mathcal{K}^e$ 都属于 $\Xi$，
引理 \ref{spaces-more-morphisms-lemma-existence-tricky} 的假设成立，结论随之得到。
\end{proof}

\begin{remark}[展开 Grothendieck 存在定理]
\label{spaces-more-morphisms-remark-reformulate-existence-theorem}
设 $A$ 为关于理想 $I$ 完备的诺特环。
写成 $S = \Spec(A)$ 以及 $S_n = \Spec(A/I^n)$。
设 $X \to S$ 为分离且有限型的代数空间态射。
对 $n \geq 1$，令 $X_n = X \times_S S_n$。图示如下：
$$
\xymatrix{
X_1 \ar[r]_{i_1} \ar[d] & X_2 \ar[r]_{i_2} \ar[d] & X_3 \ar[r] \ar[d] &
\ldots & X \ar[d] \\
S_1 \ar[r] & S_2 \ar[r] & S_3 \ar[r] & \ldots & S
}
$$
在此情形下，考虑系统 $(\mathcal{F}_n, \varphi_n)$，其中
\begin{enumerate}
\item $\mathcal{F}_n$ 是凝聚 $\mathcal{O}_{X_n}$-模；
\item $\varphi_n : i_n^*\mathcal{F}_{n + 1} \to \mathcal{F}_n$
是同构；
\item $\text{Supp}(\mathcal{F}_1)$ 在 $S_1$ 上固有。
\end{enumerate}
定理 \ref{spaces-more-morphisms-theorem-grothendieck-existence} 说明完备化函子
$$
\begin{matrix}
\text{凝聚 }\mathcal{O}_X\text{-模 }\mathcal{F} \\
\text{其支撑固有于 }A
\end{matrix}
\quad
\longrightarrow
\quad
\begin{matrix}
\text{系统 }(\mathcal{F}_n) \\
\text{如上}
\end{matrix}
$$
是范畴等价。在 $X$ 在 $A$ 上固有这一特殊情形中，可以省略
关于支撑的条件。
\end{remark}

\section{Grothendieck 代数化定理}
\label{spaces-more-morphisms-section-algebraization}

\noindent
本节是《概形的上同调》第
\ref{coherent-section-algebraization} 节的类似版本。
不过，本节没有给出带丰沛可逆层的固有代数空间之形变的代数化结果，
因为一个带丰沛可逆层的固有代数空间已经是概形。
我们确实给出了相对维数为 $1$ 的固有代数空间的代数化结果。
第一个结果以闭子概形和有限态射的语言转述 Grothendieck 存在定理。
% P09-ERR-0726

\begin{lemma}
\label{spaces-more-morphisms-lemma-algebraize-formal-closed-subscheme}
设 $A$ 为关于理想 $I$ 完备的诺特环。
写成 $S = \Spec(A)$ 以及 $S_n = \Spec(A/I^n)$。
设 $X \to S$ 为分离且有限型的代数空间态射。
对 $n \geq 1$，令 $X_n = X \times_S S_n$。
设给定代数空间的交换图
$$
\xymatrix{
Z_1 \ar[r] \ar[d] & Z_2 \ar[r] \ar[d] & Z_3 \ar[r] \ar[d] & \ldots \\
X_1 \ar[r]^{i_1} & X_2 \ar[r]^{i_2} & X_3 \ar[r] & \ldots
}
$$
且其中各方块为笛卡尔方块。假设
\begin{enumerate}
\item $Z_1 \to X_1$ 是闭浸入；
\item $Z_1 \to S_1$ 固有。
\end{enumerate}
则存在代数空间的闭浸入 $Z \to X$，使得
$Z_n = Z \times_S S_n$，对所有 $n \geq 1$。
此外，$Z$ 在 $S$ 上固有。
\end{lemma}

\begin{proof}
把竖直态射记为 $j_n : Z_n \to X_n$。
由于陈述中的方块为笛卡尔方块，$j_n$ 到 $X_1$ 的基变换为
$j_1$。因此《空间的极限》引理
\ref{spaces-limits-lemma-check-closed-infinitesimally} 表明
$j_n$ 是闭浸入。
令 $\mathcal{F}_n = j_{n, *}\mathcal{O}_{Z_n}$，于是
$j_n^\sharp$ 是满射 $\mathcal{O}_{X_n} \to \mathcal{F}_n$。
再次使用方块为笛卡尔方块这一事实，可知
$\mathcal{F}_{n + 1}$ 到 $X_n$ 的拉回是 $\mathcal{F}_n$。
因此，由评注 \ref{spaces-more-morphisms-remark-reformulate-existence-theorem} 中转述的
Grothendieck 存在定理，存在映射 $\mathcal{O}_X \to \mathcal{F}$；
它是凝聚 $\mathcal{O}_X$-模之间的映射，并且限制到 $X_n$ 后恢复
$\mathcal{O}_{X_n} \to \mathcal{F}_n$。
此外，$\mathcal{F}$ 的支撑在 $S$ 上固有。
由于完备化函子正合（引理 \ref{spaces-more-morphisms-lemma-exact}），可知
$\mathcal{O}_X \to \mathcal{F}$ 是满射。
因此有 $\mathcal{F} = \mathcal{O}_X/\mathcal{J}$，其中
$\mathcal{J}$ 是某个拟凝聚理想层。
令 $Z = V(\mathcal{J})$，证明即告完成。
\end{proof}

\begin{lemma}
\label{spaces-more-morphisms-lemma-algebraize-formal-algebraic-space-finite-over-proper}
设 $A$ 为关于理想 $I$ 完备的诺特环。
写成 $S = \Spec(A)$ 以及 $S_n = \Spec(A/I^n)$。
设 $X \to S$ 为分离且有限型的代数空间态射。
对 $n \geq 1$，令 $X_n = X \times_S S_n$。
设给定代数空间的交换图
$$
\xymatrix{
Y_1 \ar[r] \ar[d] & Y_2 \ar[r] \ar[d] & Y_3 \ar[r] \ar[d] & \ldots \\
X_1 \ar[r]^{i_1} & X_2 \ar[r]^{i_2} & X_3 \ar[r] & \ldots
}
$$
且其中各方块为笛卡尔方块。假设
\begin{enumerate}
\item $Y_1 \to X_1$ 是有限态射；
\item $Y_1 \to S_1$ 固有。
\end{enumerate}
则存在代数空间的有限态射 $Y \to X$，使得
$Y_n = Y \times_S S_n$，对所有 $n \geq 1$。
此外，$Y$ 在 $S$ 上固有。
\end{lemma}

\begin{proof}
把竖直态射记为 $f_n : Y_n \to X_n$。
由于陈述中的方块为笛卡尔方块，$f_n$ 到 $X_1$ 的基变换为
$f_1$。因此引理 \ref{spaces-more-morphisms-lemma-thicken-property-morphisms-cartesian}
表明 $f_n$ 是有限态射。
令 $\mathcal{F}_n = f_{n, *}\mathcal{O}_{Y_n}$。
使用方块为笛卡尔方块这一事实，可知
$\mathcal{F}_{n + 1}$ 到 $X_n$ 的拉回是 $\mathcal{F}_n$。
因此，由评注 \ref{spaces-more-morphisms-remark-reformulate-existence-theorem} 中转述的
Grothendieck 存在定理，存在凝聚 $\mathcal{O}_X$-模
$\mathcal{F}$，它限制到 $X_n$ 后恢复 $\mathcal{F}_n$。
此外，$\mathcal{F}$ 的支撑在 $S$ 上固有。
由于完备化函子全忠实（定理 \ref{spaces-more-morphisms-theorem-grothendieck-existence}），
可知乘法映射
$\mathcal{F}_n \otimes_{\mathcal{O}_{X_n}} \mathcal{F}_n \to
\mathcal{F}_n$ 相容地给出 $\mathcal{F}$ 上的代数结构。
% P09-ERR-0727
令 $Y = \underline{\Spec}_X(\mathcal{F})$，证明即告完成。
\end{proof}

\begin{lemma}
\label{spaces-more-morphisms-lemma-algebraize-morphism}
设 $A$ 为关于理想 $I$ 完备的诺特环。
写成 $S = \Spec(A)$ 以及 $S_n = \Spec(A/I^n)$。
设 $X$、$Y$ 为 $S$ 上的代数空间。对 $n \geq 1$，令
$X_n = X \times_S S_n$ 以及 $Y_n = Y \times_S S_n$。
设给定相容的交换图系统
$$
\xymatrix{
& & X_{n + 1} \ar[rd] \ar[rr]_{g_{n + 1}} & & Y_{n + 1} \ar[ld] \\
X_n \ar[rru] \ar[rd] \ar[rr]_{g_n} & & Y_n \ar[rru] \ar[ld] & S_{n + 1} \\
& S_n \ar[rru]
}
$$
假设
\begin{enumerate}
\item $X \to S$ 固有；
\item $Y \to S$ 分离且有限型。% P09-ERR-0728
\end{enumerate}
则存在唯一的代数空间态射 $g : X \to Y$（在 $S$ 上），
使得 $g_n$ 是 $g$ 到 $S_n$ 的基变换。
\end{lemma}

\begin{proof}
态射 $(1, g_n) : X_n \to X_n \times_S Y_n$ 是闭浸入，
因为 $Y_n \to S_n$ 分离（《空间的态射》引理
\ref{spaces-morphisms-lemma-section-immersion}）。
因此由引理 \ref{spaces-more-morphisms-lemma-algebraize-formal-closed-subscheme}，
存在闭子空间 $Z \subset X \times_S Y$，它在 $S$ 上固有，
且到 $S_n$ 的基变换恢复
$X_n \subset X_n \times_S Y_n$。第一投影 $p : Z \to X$
是固有态射（因为 $Z$ 在 $S$ 上固有；见《空间的态射》引理
\ref{spaces-morphisms-lemma-universally-closed-permanence}），
并且它到 $S_n$ 的基变换对所有 $n$ 都是同构。
特别地，$p : Z \to X$ 在 $Z$ 的一个开子空间上拟有限，
该开子空间包含 $Z_0$ 的每一点；例如由《空间的态射》引理
\ref{spaces-morphisms-lemma-locally-finite-type-quasi-finite-part} 可得。
% P09-ERR-0729
由于 $Z$ 在 $S$ 上固有，此开邻域就是整个 $Z$。% P09-ERR-0730
因此由 Zariski 主定理可知 $p : Z \to X$ 有限
（例如应用引理 \ref{spaces-more-morphisms-lemma-quasi-finite-separated-pass-through-finite}，
并利用 $Z$ 在 $X$ 上的固有性看出该浸入是闭浸入）。
应用定理 \ref{spaces-more-morphisms-theorem-grothendieck-existence} 的范畴等价，
可知 $p_*\mathcal{O}_Z = \mathcal{O}_X$，因为这在模 $I^n$
之后对所有 $n$ 都成立。因此 $p$ 是同构，并且得到态射
$g$，即复合 $X \cong Z \to Y$。
我们略去唯一性的证明。% P09-ERR-0731
\end{proof}

\begin{remark}
\label{spaces-more-morphisms-remark-weaken-separation-axioms-question}
可以问，在 Grothendieck 代数化定理（采用引理
\ref{spaces-more-morphisms-lemma-algebraize-morphism} 的形式）中，能否对靶采用更弱的分离公理。
更精确地说，设 $A$、$I$、$S$、$S_n$、$X$、$Y$、$X_n$、
$Y_n$ 和 $g_n$ 如引理 \ref{spaces-more-morphisms-lemma-algebraize-morphism} 的陈述，并假设
\begin{enumerate}
\item $X \to S$ 固有；
\item $Y \to S$ 局部有限型。
\end{enumerate}
是否存在代数空间态射 $g : X \to Y$（在 $S$ 上），使得
$g_n$ 是 $g$ 到 $S_n$ 的基变换？一般情形下我们不知道答案；
若你知道，请发送电子邮件至
\href{mailto:stacks.project@gmail.com}{stacks.project@gmail.com}。
若 $Y \to S$ 分离，则由该引理结论成立（可立即约化到
$X$ 在 $S$ 上有限型的情形，具体是选择一个包含 $g_1$ 的像的拟紧开集）。
% P09-ERR-0732
若只假设 $Y \to S$ 拟分离，结论也成立。首先，如前一样，
可假设 $Y$ 既拟紧又拟分离。然后可使用 \cite{Bhatt-Algebraize}
或从 \cite{Hall-Rydh-coherent} 将 $(g_n)$ 代数化。% P09-ERR-0733
具体而言，为应用第一篇参考文献，使用
$$
D_{perf}(X) \to \lim D_{perf}(X_n)
\xrightarrow{\lim Lg_n^*}
\lim D_{perf}(Y_n) = D_{perf}(Y)
$$
其中最后一步使用固有代数空间 $Y$ 在 $R$ 上的导出范畴的
Grothendieck 存在结果（参见《空间上的平坦性》评注
\ref{spaces-flat-remark-correct-generality}）。
所引论文表明，此箭头按所需确定态射 $Y \to X$。% P09-ERR-0734
为应用第二篇参考文献，对凝聚模使用相同论证：
$$
\textit{Coh}(\mathcal{O}_X) \to \lim \textit{Coh}(\mathcal{O}_{X_n})
\xrightarrow{\lim g_n^*}
\lim \textit{Coh}(\mathcal{O}_{Y_n}) = \textit{Coh}(\mathcal{O}_Y)
$$
其中最后一个等式是 Grothendieck 存在定理（定理
\ref{spaces-more-morphisms-theorem-grothendieck-existence}）的推论。
第二篇参考文献说明，此函子对应于态射 $Y \to X$（在 $R$ 上）。
% P09-ERR-0735
若以后需要这一推广，我们将在此处精确陈述并仔细证明结果。
\end{remark}

\begin{lemma}
\label{spaces-more-morphisms-lemma-formal-algebraic-space-proper-reldim-1}
设 $(A, \mathfrak m, \kappa)$ 为完备局部诺特环。
令 $S = \Spec(A)$ 以及 $S_n = \Spec(A/\mathfrak m^n)$。
考虑代数空间的交换图
$$
\xymatrix{
X_1 \ar[r]_{i_1} \ar[d] & X_2 \ar[r]_{i_2} \ar[d] & X_3 \ar[r] \ar[d] &
\ldots \\
S_1 \ar[r] & S_2 \ar[r] & S_3 \ar[r] & \ldots
}
$$
且其中各方块为笛卡尔方块。若 $\dim(X_1) \leq 1$，
则存在概形的射影态射 $X \to S$ 以及同构
$X_n \cong X \times_S S_n$，且这些同构与 $i_n$ 相容。
% P09-ERR-0736
\end{lemma}

\begin{proof}
由《域上的空间》引理
\ref{spaces-over-fields-lemma-codim-1-point-in-schematic-locus}，
代数空间 $X_1$ 是概形。因此，$X_1$ 是维数
$\leq 1$ 的固有概形（在 $\kappa$ 上）。由《簇》引理
\ref{varieties-lemma-dim-1-proper-projective}，可知 $X_1$ 在
$\kappa$ 上 H-射影。设 $\mathcal{L}_1$ 为 $X_1$ 上的丰沛可逆层。

\medskip\noindent
我们将证明 $\mathcal{L}_1$ 可提升为可逆层的相容系统
$\{\mathcal{L}_n\}$；这些层定义在 $\{X_n\}$ 上。注意，由引理 \ref{spaces-more-morphisms-lemma-thickening-scheme}，
$X_n$ 也是概形。回顾 $X_1 \to X_n$ 诱导底拓扑空间之间的同胚。
在证明的其余部分，我们不区分 $X_n$ 上的层和 $X_1$ 上的层。
假设给定 $\mathcal{L}_n$ 到 $X_n$ 的一个提升。% P09-ERR-0737
我们考虑正合序列
$$
1 \to
(1 + \mathfrak m^n\mathcal{O}_{X_{n + 1}})^* \to
\mathcal{O}_{X_{n + 1}}^* \to \mathcal{O}_{X_n}^* \to 1
$$
它由 $X_{n + 1}$ 上的层组成。
$\mathcal{L}_n$ 在 $H^1(X_n, \mathcal{O}_{X_n}^*)$ 中的类（见
《上同调》引理 \ref{cohomology-lemma-h1-invertible}）可提升为
$H^1(X_{n + 1}, \mathcal{O}_{X_{n + 1}}^*)$ 的元素，当且仅当
$H^2(X_{n + 1}, (1 + \mathfrak m^n\mathcal{O}_{X_{n + 1}})^*)$
中的障碍为零。由于 $X_1$ 是维数 $\leq 1$ 的诺特概形，
此上同调群消失（《上同调》命题
\ref{cohomology-proposition-vanishing-Noetherian}）。

\medskip\noindent
由 Grothendieck 代数化定理（《概形的上同调》定理
\ref{coherent-theorem-algebraization}），我们得到概形的射影态射
$X \to S = \Spec(A)$ 以及相容的同构系统
$X_n = S_n \times_S X$。
\end{proof}

\begin{lemma}
\label{spaces-more-morphisms-lemma-projective-over-complete}
设 $(A, \mathfrak m, \kappa)$ 为完备诺特局部环。
设 $X$ 为 $\Spec(A)$ 上的代数空间。若
$X \to \Spec(A)$ 固有且 $\dim(X_\kappa) \leq 1$，则
$X$ 是 $A$ 上的射影概形。
\end{lemma}

\begin{proof}
令 $X_n = X \times_{\Spec(A)} \Spec(A/\mathfrak m^n)$。
由引理 \ref{spaces-more-morphisms-lemma-formal-algebraic-space-proper-reldim-1}，
存在射影态射 $Y \to \Spec(A)$ 以及相容的同构
$Y \times_{\Spec(A)} \Spec(A/\mathfrak m^n) \cong
X \times_{\Spec(A)} \Spec(A/\mathfrak m^n)$。
由引理 \ref{spaces-more-morphisms-lemma-algebraize-morphism}，可知 $X \cong Y$，
证明即告完成。
\end{proof}

\section{正则浸入}
\label{spaces-more-morphisms-section-regular-immersions}

\noindent
本节是《除子》第 \ref{divisors-section-regular-immersions} 节
针对代数空间态射的类似版本。建议读者先在该节了解概形的正则浸入。

\medskip\noindent
在《除子》第 \ref{divisors-section-regular-immersions} 节中，
我们为概形态射定义了四类正则浸入。据我们所知，其中只有三类
关于靶在平展拓扑下局部；通常意义下的正则浸入仍然不是。
正因如此，我们实际上无法为代数空间态射定义正则浸入。
（这类麻烦促使 Grothendieck 及其学派以 Koszul-正则浸入取代
正则浸入的原始概念；见 \cite[Exposee VII, Definition 1.4]{SGA6}。）
% P09-ERR-0738 P09-ERR-0739
不过，我们可以定义 Koszul-正则、$H_1$-正则和拟正则浸入。
还要指出，由于 Koszul-正则浸入不在任意基变换下保持，
我们不能采用《空间的态射》第
\ref{spaces-morphisms-section-representable} 节的策略来定义它们。
类似地，由于 Koszul-正则浸入不关于源平展局部，我们也不能使用
《空间的态射》引理
\ref{spaces-morphisms-lemma-local-source-target} 来定义它们。
我们改以下面的引理代替该引理。

\begin{lemma}
\label{spaces-more-morphisms-lemma-representable-etale-local-target}
设 $\mathcal{P}$ 为概形态射的一项性质，且它关于靶平展局部。
设 $S$ 为概形。设 $f : X \to Y$ 为 $S$ 上代数空间的可表态射。
考虑交换图
$$
\xymatrix{
X \times_Y V \ar[d] \ar[r] & V \ar[d] \\
X \ar[r]^f & Y
}
$$
其中 $V$ 为概形，且 $V \to Y$ 平展。以下条件等价：
\begin{enumerate}
\item 对任意如上图，投影 $X \times_Y V \to V$ 具有性质
$\mathcal{P}$；
\item 对某个如上图，若 $V \to Y$ 满，则投影
$X \times_Y V \to V$ 具有性质 $\mathcal{P}$。
\end{enumerate}
若 $X$ 和 $Y$ 可表，则这些条件还等价于 $f$（作为概形态射）
具有性质 $\mathcal{P}$。
\end{lemma}

\begin{proof}
先证明 (1) 与 (2) 等价。蕴含 (1) $\Rightarrow$ (2) 是立即的。
假设
$$
\xymatrix{
X \times_Y V \ar[d] \ar[r] & V \ar[d] \\
X \ar[r]^f & Y
}
\quad\quad
\xymatrix{
X \times_Y V' \ar[d] \ar[r] & V' \ar[d] \\
X \ar[r]^f & Y
}
$$
是引理中的两个图。假设 $V \to Y$ 满，且
$X \times_Y V \to V$ 具有性质 $\mathcal{P}$。
为证明 (2) 蕴含 (1)，我们必须证明
$X \times_Y V' \to V'$ 具有 $\mathcal{P}$。为此，考虑图
$$
\xymatrix{
X \times_Y V \ar[d] &
(X \times_Y V) \times_X (X \times_Y V') \ar[l] \ar[d] \ar[r] &
X \times_Y V' \ar[d] \\
V &
V \times_Y V' \ar[l] \ar[r] &
V'
}
$$
由我们关于 $\mathcal{P}$ 在源上平展局部的假设，% P09-ERR-0740
可知 $\mathcal{P}$ 在平展基变换下保持；见《下降》引理
\ref{descent-lemma-pullback-property-local-target}。
因此，若左边的竖直箭头具有 $\mathcal{P}$，则中间的竖直箭头也如此。
% P09-ERR-0741
由于 $U \times_X U' \to U'$ 满且平展（因而给出 $U'$ 的平展覆盖），
% P09-ERR-0742
又因为假设 $\mathcal{P}$ 关于靶在平展拓扑下局部，
故这蕴含左边的竖直箭头具有 $\mathcal{P}$。% P09-ERR-0743

\medskip\noindent
若 $X$ 和 $Y$ 可表，则可取 $\text{id}_Y : Y \to Y$ 作为平展覆盖，
从而看出引理的最后一个陈述成立。
\end{proof}

\noindent
注意，“是 Koszul-正则（相应地，$H_1$-正则、拟正则）浸入”
是概形态射的一项性质，并且它关于靶 fpqc 局部；见《下降》引理
\ref{descent-lemma-descending-property-regular-immersion}。
因此，下面的定义是有意义的。

\begin{definition}
\label{spaces-more-morphisms-definition-regular-immersion}
设 $S$ 为概形。设 $i : X \to Y$ 为 $S$ 上代数空间的态射。
\begin{enumerate}
\item 称 $i$ 为{\it Koszul-正则浸入}，如果 $i$ 可表，并且引理
\ref{spaces-more-morphisms-lemma-representable-etale-local-target} 的等价条件对
$\mathcal{P}(f) =$“$f$ 是 Koszul-正则浸入”成立。
\item 称 $i$ 为{\it $H_1$-正则浸入}，如果 $i$ 可表，并且引理
\ref{spaces-more-morphisms-lemma-representable-etale-local-target} 的等价条件对
$\mathcal{P}(f) =$“$f$ 是 $H_1$-正则浸入”成立。
\item 称 $i$ 为{\it 拟正则浸入}，如果 $i$ 可表，并且引理
\ref{spaces-more-morphisms-lemma-representable-etale-local-target} 的等价条件对
$\mathcal{P}(f) =$“$f$ 是拟正则浸入”成立。
\end{enumerate}
\end{definition}

\begin{lemma}
\label{spaces-more-morphisms-lemma-regular-quasi-regular-immersion}
设 $S$ 为概形。设 $i : Z \to X$ 为 $S$ 上代数空间的浸入。
我们有以下蕴含：
$i$ 是 Koszul-正则 $\Rightarrow$
$i$ 是 $H_1$-正则 $\Rightarrow$
$i$ 是拟正则。
\end{lemma}

\begin{proof}
由定义，本引理立即约化为《除子》引理
\ref{divisors-lemma-regular-quasi-regular-immersion}。
\end{proof}

\begin{lemma}
\label{spaces-more-morphisms-lemma-regular-immersion-noetherian}
设 $S$ 为概形。设 $i : Z \to X$ 为 $S$ 上代数空间的浸入。
假设 $X$ 局部诺特。则
$i$ 是 Koszul-正则 $\Leftrightarrow$
$i$ 是 $H_1$-正则 $\Leftrightarrow$
$i$ 是拟正则。
\end{lemma}

\begin{proof}
由定义 \ref{spaces-more-morphisms-definition-regular-immersion}
（以及《空间的性质》第
\ref{spaces-properties-section-types-properties} 节中局部诺特代数空间的定义），
这立即转化为概形情形，即《除子》引理
\ref{divisors-lemma-regular-immersion-noetherian}。
\end{proof}

\begin{lemma}
\label{spaces-more-morphisms-lemma-flat-base-change-regular-immersion}
\begin{slogan}
正则浸入在平坦基变换下稳定。
\end{slogan}
设 $S$ 为概形。设 $i : Z \to X$ 为代数空间的
Koszul-正则、$H_1$-正则或拟正则浸入（在 $S$ 上）。
设 $X' \to X$ 为 $S$ 上代数空间的平坦态射。
则基变换 $i' : Z \times_X X' \to X'$ 是 Koszul-正则、
$H_1$-正则或拟正则浸入。
\end{lemma}

\begin{proof}
由定义 \ref{spaces-more-morphisms-definition-regular-immersion}
（以及《空间的态射》第
\ref{spaces-morphisms-section-flat} 节中代数空间平坦态射的定义），
本引理约化为概形情形；见《除子》引理
\ref{divisors-lemma-flat-base-change-regular-immersion}。
\end{proof}

\begin{lemma}
\label{spaces-more-morphisms-lemma-quasi-regular-immersion}
设 $S$ 为概形。设 $i : Z \to X$ 为 $S$ 上代数空间的浸入。
则 $i$ 是拟正则浸入，当且仅当以下条件成立：
\begin{enumerate}
\item $i$ 局部有限表示；
\item 余法层 $\mathcal{C}_{Z/X}$ 有限局部自由；
\item 映射 (\ref{spaces-more-morphisms-equation-conormal-algebra-quotient}) 是同构。
\end{enumerate}
\end{lemma}

\begin{proof}
通过平展局部化，这由概形情形推出（《除子》引理
\ref{divisors-lemma-quasi-regular-immersion}）；使用定义
\ref{spaces-more-morphisms-definition-regular-immersion} 和引理
\ref{spaces-more-morphisms-lemma-etale-conormal-algebra}。
\end{proof}

\begin{lemma}
\label{spaces-more-morphisms-lemma-transitivity-conormal-quasi-regular}
设 $S$ 为概形。设 $Z \to Y \to X$ 为 $S$ 上代数空间的浸入。
假设 $Z \to Y$ 是 $H_1$-正则的。则引理
\ref{spaces-more-morphisms-lemma-transitivity-conormal} 的典范序列
$$
0 \to i^*\mathcal{C}_{Y/X} \to
\mathcal{C}_{Z/X} \to
\mathcal{C}_{Z/Y} \to 0
$$
正合，并且（平展）局部分裂。
\end{lemma}

\begin{proof}
由于 $\mathcal{C}_{Z/Y}$ 有限局部自由（见引理
\ref{spaces-more-morphisms-lemma-quasi-regular-immersion} 和引理
\ref{spaces-more-morphisms-lemma-regular-quasi-regular-immersion}），只须证明该序列正合。
因为该序列一般已经右正合，只须说明第一个映射是单射。
在 $X$ 上作平展局部化后，这约化为概形情形；见《除子》引理
\ref{divisors-lemma-transitivity-conormal-quasi-regular}。
\end{proof}

\noindent
拟正则浸入的复合未必拟正则；见《代数》评注
\ref{algebra-remark-join-quasi-regular-sequences}。
其他类型的正则浸入在复合下保持。

\begin{lemma}
\label{spaces-more-morphisms-lemma-composition-regular-immersion}
设 $S$ 为概形。设 $i : Z \to Y$ 和 $j : Y \to X$ 为
$S$ 上代数空间的浸入。
\begin{enumerate}
\item 若 $i$ 和 $j$ 都是 Koszul-正则浸入，则 $j \circ i$ 也是。
\item 若 $i$ 和 $j$ 都是 $H_1$-正则浸入，则 $j \circ i$ 也是。
\item 若 $i$ 是 $H_1$-正则浸入且 $j$ 是拟正则浸入，
则 $j \circ i$ 是拟正则浸入。
\end{enumerate}
\end{lemma}

\begin{proof}
这由概形情形立即推出；见《除子》引理
\ref{divisors-lemma-composition-regular-immersion}。
\end{proof}

\begin{lemma}
\label{spaces-more-morphisms-lemma-permanence-regular-immersion}
设 $S$ 为概形。设 $i : Z \to Y$ 和 $j : Y \to X$ 为
$S$ 上代数空间的浸入。假设 $j$ 局部有限表示，并且引理
\ref{spaces-more-morphisms-lemma-transitivity-conormal} 的序列
$$
0 \to i^*\mathcal{C}_{Y/X} \to
\mathcal{C}_{Z/X} \to
\mathcal{C}_{Z/Y} \to 0
$$
正合且局部分裂。
\begin{enumerate}
\item 若 $j \circ i$ 是拟正则浸入，则 $i$ 也是。
\item 若 $j \circ i$ 是 $H_1$-正则浸入，则 $i$ 也是。
\item 若 $j$ 和 $j \circ i$ 都是 Koszul-正则浸入，则 $i$ 也是。
\end{enumerate}
\end{lemma}

\begin{proof}
这由概形情形立即推出；见《除子》引理
\ref{divisors-lemma-permanence-regular-immersion}。
\end{proof}

\begin{lemma}
\label{spaces-more-morphisms-lemma-extra-permanence-regular-immersion-noetherian}
设 $S$ 为概形。设 $i : Z \to Y$ 和 $j : Y \to X$ 为
$S$ 上代数空间的浸入。假设 $X$ 局部诺特。以下条件等价：
\begin{enumerate}
\item $i$ 和 $j$ 是 Koszul 正则浸入；
\item $i$ 和 $j \circ i$ 是 Koszul 正则浸入；
\item $j \circ i$ 是 Koszul 正则浸入，且余法序列
$$
0 \to i^*\mathcal{C}_{Y/X} \to
\mathcal{C}_{Z/X} \to
\mathcal{C}_{Z/Y} \to 0
$$
正合且局部分裂。
\end{enumerate}
% P09-ERR-0744
\end{lemma}

\begin{proof}
这由概形情形立即推出；见《除子》引理
\ref{divisors-lemma-extra-permanence-regular-immersion-noetherian}。
\end{proof}

\section{相对拟凝聚性}
\label{spaces-more-morphisms-section-relative-pseudo-coherence}

\noindent
本节是《态射进阶》第
\ref{more-morphisms-section-relative-pseudo-coherence} 节的类似版本。
不过，在处理代数空间的这些材料时，我们决定只使用导出范畴中
上同调层拟凝聚的对象。理由有二：(1) 这大大简化了叙述；
(2) 我们目前用不到更一般的概念。

\begin{remark}
\label{spaces-more-morphisms-remark-match-relative-pseudo-coherence}
设 $S$ 为概形。设 $f : X \to Y$ 为 $S$ 上可表代数空间的
局部有限型态射。设 $f_0 : X_0 \to Y_0$ 为表示 $f$ 的概形态射
（这一记号略显别扭，但只是暂用）。则 $f_0$ 局部有限型。
若 $E$ 是 $D_\QCoh(\mathcal{O}_X)$ 的对象，则 $E$ 是
唯一对象 $E_0$ 的拉回，该对象属于 $D_\QCoh(\mathcal{O}_{X_0})$；见
《空间的导出范畴》引理
\ref{spaces-perfect-lemma-derived-quasi-coherent-small-etale-site}。
在此情形下，“$E$ 是 $m$-拟凝聚的（相对于 $Y$）”意指
“$E_0$ 是 $m$-拟凝聚的（相对于 $Y_0$）”，其定义见《态射进阶》第
\ref{more-morphisms-section-relative-pseudo-coherence} 节。
\end{remark}

\begin{lemma}
\label{spaces-more-morphisms-lemma-qcoh-relative-pseudo-coherence-characterize}
设 $S$ 为概形。设 $f : X \to Y$ 为 $S$ 上代数空间的
局部有限型态射。设 $m \in \mathbf{Z}$。设
$E \in D_\QCoh(\mathcal{O}_X)$。采用评注
\ref{spaces-more-morphisms-remark-match-relative-pseudo-coherence} 中说明的记号，
以下条件等价：
\begin{enumerate}
\item 对每个交换图
$$
\xymatrix{
U \ar[d] \ar[r] & V \ar[d] \\
X \ar[r] & Y
}
$$
其中 $U$、$V$ 为概形且竖直箭头平展，复形 $E|_U$ 是
$m$-拟凝聚的（相对于 $V$）；
\item 对某个如 (1) 的交换图，若 $U \to X$ 满，则复形
$E|_U$ 是 $m$-拟凝聚的（相对于 $V$）；
\item 对每个如 (1) 的交换图，若 $U$ 和 $V$ 仿射，则复形
$R\Gamma(U, E)$（作为 $\mathcal{O}_X(U)$-模）是
$m$-拟凝聚的（相对于 $\mathcal{O}_Y(V)$）。
\end{enumerate}
\end{lemma}

\begin{proof}
由《态射进阶》引理
\ref{more-morphisms-lemma-qcoh-relative-pseudo-coherence-characterize}，
(1) 蕴含 (3)。

\medskip\noindent
假设 (3)。任取一个如 (1) 的交换图，其中 $U \to X$ 满。
选取仿射开覆盖 $V = \bigcup V_j$，并选取仿射开覆盖
$(U \to V)^{-1}(V_j) = \bigcup U_{ij}$。由 (3) 及《态射进阶》引理
\ref{more-morphisms-lemma-qcoh-relative-pseudo-coherence-characterize}，
可知 $E|_U$ 是 $m$-拟凝聚的（相对于 $V$）。因此 (3) 蕴含 (2)。

\medskip\noindent
假设 (2)。选取交换图
$$
\xymatrix{
U \ar[d] \ar[r] & V \ar[d] \\
X \ar[r] & Y
}
$$
其中 $U$、$V$ 为概形，竖直箭头平展，态射 $U \to X$ 满，
且 $E|_U$ 是 $m$-拟凝聚的（相对于 $V$）。接着，设给定第二个交换图
$$
\xymatrix{
U' \ar[d] \ar[r] & V' \ar[d] \\
X \ar[r] & Y
}
$$
其中竖直箭头平展，且 $U', V'$ 为概形。我们要证明
$E|_{U'}$ 是 $m$-拟凝聚的（相对于 $V'$）。
态射 $U'' = U \times_X U' \to U'$ 满且平展，而
$U'' \to V'$ 分解通过 $V'' = V' \times_Y V$；后者在 $V'$ 上平展。
因此只须证明 $E|_{U''}$ 是 $m$-拟凝聚的（相对于 $V''$）；见
《态射进阶》引理
\ref{more-morphisms-lemma-relative-pseudo-coherent-descends-fppf} 和
\ref{more-morphisms-lemma-relative-pseudo-coherent-post-compose}。
再使用一次第二个引理，只须证明 $E|_{U''}$ 是 $m$-拟凝聚的
（相对于 $V$）。由《态射进阶》引理
\ref{more-morphisms-lemma-pull-relative-pseudo-coherent}，
以及概形的平展态射拟凝聚这一事实，此断言成立；后者见《态射进阶》引理
\ref{more-morphisms-lemma-flat-finite-presentation-pseudo-coherent}。
\end{proof}

\begin{definition}
\label{spaces-more-morphisms-definition-relative-pseudo-coherence}
设 $S$ 为概形。设 $f : X \to Y$ 为 $S$ 上代数空间的
局部有限型态射。设 $E$ 为 $D_\QCoh(\mathcal{O}_X)$ 的对象。
设 $\mathcal{F}$ 为拟凝聚 $\mathcal{O}_X$-模。固定
$m \in \mathbf{Z}$。
\begin{enumerate}
\item 称 $E$ {\it 为 $m$-拟凝聚的（相对于 $Y$）}，如果引理
\ref{spaces-more-morphisms-lemma-qcoh-relative-pseudo-coherence-characterize} 的等价条件成立。
\item 称 $E$ {\it 相对于 $Y$ 拟凝聚}，如果 $E$ 是
$m$-拟凝聚的（相对于 $Y$），对所有 $m \in \mathbf{Z}$。
\item 称 $\mathcal{F}$ {\it 为 $m$-拟凝聚的（相对于 $Y$）}，
如果把 $\mathcal{F}$ 视为 $D_\QCoh(\mathcal{O}_X)$ 的对象时，
它是 $m$-拟凝聚的（相对于 $Y$）。
\item 称 $\mathcal{F}$ {\it 相对于 $Y$ 拟凝聚}，
如果把 $\mathcal{F}$ 视为 $D_\QCoh(\mathcal{O}_X)$ 的对象时，
它相对于 $Y$ 拟凝聚。
\end{enumerate}
\end{definition}

\noindent
复形相对于基拟凝聚的大多数性质将立即由概形情形的相应性质推出。
我们会在需要时在此加入相关引理。

\begin{lemma}
\label{spaces-more-morphisms-lemma-relative-pseudo-coherent-is-moot}
设 $S$ 为概形。设 $f : X \to Y$ 为 $S$ 上代数空间的态射。
设 $E$ 为 $D_\QCoh(\mathcal{O}_X)$ 的对象。% P09-ERR-0745
若 $f$ 平坦且局部有限表示，则以下条件等价：
\begin{enumerate}
\item $E$ 相对于 $Y$ 拟凝聚；
\item $E$ 在 $X$ 上拟凝聚。
\end{enumerate}
\end{lemma}

\begin{proof}
由平展局部化和定义，可以假设 $X$ 和 $Y$ 是概形。
对概形情形，此结论由《态射进阶》引理
\ref{more-morphisms-lemma-check-relative-pseudo-coherence-on-charts} 推出。
\end{proof}

\section{拟凝聚态射}
\label{spaces-more-morphisms-section-pseudo-coherent}

\noindent
本节是《态射进阶》第
\ref{more-morphisms-section-pseudo-coherent} 节针对概形态射的类似版本。
建议读者先在该节了解概形的拟凝聚态射。

\medskip\noindent
概形态射的“拟凝聚”性质关于源和靶平展局部。为此，使用
《态射进阶》引理
\ref{more-morphisms-lemma-descending-property-pseudo-coherent} 和
\ref{more-morphisms-lemma-pseudo-coherent-fppf-local-source}，
以及《下降》引理 \ref{descent-lemma-etale-local-source-target}。
由《空间的态射》引理
\ref{spaces-morphisms-lemma-local-source-target}，
可以如下定义代数空间拟凝聚态射；当所涉及的代数空间可表时，
它与《态射进阶》第
\ref{more-morphisms-section-pseudo-coherent} 节已经定义的概念一致。

\begin{definition}
\label{spaces-more-morphisms-definition-pseudo-coherent}
设 $S$ 为概形。设 $f : X \to Y$ 为 $S$ 上代数空间的态射。
\begin{enumerate}
\item 称 $f$ {\it 拟凝聚}，如果《空间的态射》引理
\ref{spaces-morphisms-lemma-local-source-target} 的等价条件对
$\mathcal{P} =$“拟凝聚”成立。
\item 设 $x \in |X|$。称 $f$ {\it 在 $x$ 处拟凝聚}，如果存在
开邻域 $X' \subset X$（它含有 $x$），使得
$f|_{X'} : X' \to Y$ 拟凝聚。
\end{enumerate}
\end{definition}

\noindent
注意，拟凝聚态射的基变换一般并不拟凝聚。

\begin{lemma}
\label{spaces-more-morphisms-lemma-flat-base-change-pseudo-coherent}
拟凝聚态射的平坦基变换是拟凝聚的。
\end{lemma}

\begin{proof}
从略。提示：使用本引理的概形版本；见《态射进阶》引理
\ref{more-morphisms-lemma-flat-base-change-pseudo-coherent}。
\end{proof}

\begin{lemma}
\label{spaces-more-morphisms-lemma-composition-pseudo-coherent}
拟凝聚态射的复合是拟凝聚的。
\end{lemma}

\begin{proof}
从略。提示：使用本引理的概形版本；见《态射进阶》引理
\ref{more-morphisms-lemma-composition-pseudo-coherent}。
\end{proof}

\begin{lemma}
\label{spaces-more-morphisms-lemma-pseudo-coherent-finite-presentation}
拟凝聚态射局部有限表示。
\end{lemma}

\begin{proof}
这由定义立即得到。
\end{proof}

\begin{lemma}
\label{spaces-more-morphisms-lemma-flat-finite-presentation-pseudo-coherent}
局部有限表示的平坦态射是拟凝聚的。
\end{lemma}

\begin{proof}
从略。提示：使用本引理的概形版本；见《态射进阶》引理
\ref{more-morphisms-lemma-flat-finite-presentation-pseudo-coherent}。
\end{proof}

\begin{lemma}
\label{spaces-more-morphisms-lemma-permanence-pseudo-coherent}
设 $f : X \to Y$ 为代数空间的态射，并且它在某个基代数空间
$B$ 上拟凝聚。则 $f$ 拟凝聚。
\end{lemma}

\begin{proof}
从略。提示：使用本引理的概形版本；见《态射进阶》引理
\ref{more-morphisms-lemma-permanence-pseudo-coherent}。
\end{proof}

\begin{lemma}
\label{spaces-more-morphisms-lemma-Noetherian-pseudo-coherent}
设 $S$ 为概形。设 $f : X \to Y$ 为 $S$ 上代数空间的态射。
若 $Y$ 局部诺特，则 $f$ 拟凝聚，当且仅当 $f$ 局部有限型。
\end{lemma}

\begin{proof}
从略。提示：使用本引理的概形版本；见《态射进阶》引理
\ref{more-morphisms-lemma-Noetherian-pseudo-coherent}。
\end{proof}

\section{完美态射}
\label{spaces-more-morphisms-section-perfect}

\noindent
本节是《态射进阶》第
\ref{more-morphisms-section-perfect} 节关于概形态射的内容
在代数空间情形的对应版本。建议读者先阅读该节关于概形完美态射的内容。

\medskip\noindent
概形态射的“完美”性质在源与靶上平展局部。为看出这一点，使用
《态射进阶》引理
\ref{more-morphisms-lemma-descending-property-perfect} 和
\ref{more-morphisms-lemma-perfect-fppf-local-source}
以及《下降》引理
\ref{descent-lemma-etale-local-source-target}。
由《空间的态射》引理
\ref{spaces-morphisms-lemma-local-source-target}，
可以如下定义代数空间的完美态射；当所涉及的代数空间可表时，
它与《态射进阶》第 \ref{more-morphisms-section-perfect} 节
已经定义的概念一致。

\begin{definition}
\label{spaces-more-morphisms-definition-perfect}
设 $S$ 为概形。设 $f : X \to Y$ 为 $S$ 上代数空间的态射。
\begin{enumerate}
\item 称 $f$ {\it 完美}，如果《空间的态射》引理
\ref{spaces-morphisms-lemma-local-source-target} 的等价条件对
$\mathcal{P} =$“完美”成立。
\item 设 $x \in |X|$。称 $f$ {\it 在 $x$ 处完美}，如果存在
开邻域 $X' \subset X$，它含有 $x$，并且使得
$f|_{X'} : X' \to Y$ 完美。
\end{enumerate}
\end{definition}

\noindent
注意，完美态射是拟凝聚的，因而局部有限表示。要注意，
完美态射的基变换一般并不完美。

\begin{lemma}
\label{spaces-more-morphisms-lemma-flat-base-change-perfect}
完美态射的平坦基变换是完美的。
\end{lemma}

\begin{proof}
从略。提示：使用本引理的概形版本；见《态射进阶》引理
\ref{more-morphisms-lemma-flat-base-change-perfect}。
\end{proof}

\begin{lemma}
\label{spaces-more-morphisms-lemma-composition-perfect}
完美态射的复合是完美的。
\end{lemma}

\begin{proof}
从略。提示：使用本引理的概形版本；见《态射进阶》引理
\ref{more-morphisms-lemma-composition-perfect}。
\end{proof}

\begin{lemma}
\label{spaces-more-morphisms-lemma-flat-finite-presentation-perfect}
设 $S$ 为概形。设 $f : X \to Y$ 为 $S$ 上代数空间的态射。
下列条件等价：
\begin{enumerate}
\item $f$ 平坦且完美；
\item $f$ 平坦且局部有限表示。
\end{enumerate}
\end{lemma}

\begin{proof}
从略。提示：使用本引理的概形版本；见《态射进阶》引理
\ref{more-morphisms-lemma-flat-finite-presentation-perfect}。
\end{proof}

\begin{lemma}
\label{spaces-more-morphisms-lemma-perfect-proper-perfect-direct-image}
设 $S$ 为概形。设 $Y$ 为 $S$ 上的诺特代数空间。
设 $f : X \to Y$ 为代数空间的完美固有态射。
设 $E \in D(\mathcal{O}_X)$ 为完美的。则
$Rf_*E$ 是 $D(\mathcal{O}_Y)$ 的完美对象。
\end{lemma}

\begin{proof}
我们断言《空间的导出范畴》引理
\ref{spaces-perfect-lemma-perfect-direct-image} 适用。
条件 (1) 和 (2) 是立即的。条件 (3) 在 $X$ 上是局部的。
因此可以假设 $X$ 和 $Y$ 仿射，并且 $E$ 由一个
$\mathcal{O}_X$-模的严格完美复形表示。
所以只需证明：在平展位点上，$\mathcal{O}_X$ 作为
$f^{-1}\mathcal{O}_Y$-模层具有有限 Tor 维数。由《空间的导出范畴》引理
\ref{spaces-perfect-lemma-tor-dimension-rel}，只需在 Zariski 位点上检验这一点。
对于概形的有限型态射，由《态射进阶》引理
\ref{more-morphisms-lemma-check-perfect-stalks}，
这等价于完美性。
\end{proof}

\section{局部完全交态射}
\label{spaces-more-morphisms-section-lci}

\noindent
本节是《态射进阶》第 \ref{more-morphisms-section-lci} 节
关于概形态射的内容在代数空间情形的对应版本。
建议读者先阅读该节关于概形局部完全交态射的内容。

\medskip\noindent
概形态射的“为局部完全交态射”这一性质在源与靶上平展局部。
为看出这一点，使用《态射进阶》引理
\ref{more-morphisms-lemma-descending-property-lci} 和
\ref{more-morphisms-lemma-lci-syntomic-local-source}
以及《下降》引理 \ref{descent-lemma-etale-local-source-target}。
由《空间的态射》引理
\ref{spaces-morphisms-lemma-local-source-target}，
可以如下定义代数空间的局部完全交态射；当所涉及的代数空间可表时，
它与《态射进阶》第 \ref{more-morphisms-section-lci} 节
已经定义的概念一致。

\begin{definition}
\label{spaces-more-morphisms-definition-lci}
设 $S$ 为概形。设 $f : X \to Y$ 为 $S$ 上代数空间的态射。
\begin{enumerate}
\item 称 $f$ 是一个 {\it Koszul 态射}，或称 $f$ 是一个
{\it 局部完全交态射}，如果《空间的态射》引理
\ref{spaces-morphisms-lemma-local-source-target} 的等价条件对
$\mathcal{P}(f) =$“$f$ 是一个局部完全交态射”成立。
\item 设 $x \in |X|$。称 $f$ {\it 在 $x$ 处为 Koszul}，如果存在
开邻域 $X' \subset X$，它含有 $x$，并且使得
$f|_{X'} : X' \to Y$ 是局部完全交态射。
\end{enumerate}
\end{definition}

\noindent
在某种意义下，下面的引理给出了局部完全交态射的定义性质。

\begin{lemma}
\label{spaces-more-morphisms-lemma-lci}
设 $S$ 为概形。设 $f : X \to Y$ 为 $S$ 上代数空间的
局部完全交态射。设 $P$ 为在 $Y$ 上光滑的代数空间。
设 $U \to X$ 为代数空间的平展态射，并设
% P09-ERR-0746
$i : U \to P$ 为 $Y$ 上代数空间的浸入。图示如下：
$$
\xymatrix{
X \ar[rd] & U \ar[l] \ar[d] \ar[r]_i & P \ar[ld] \\
& Y
}
$$
则 $i$ 是代数空间的 Koszul-正则浸入。
\end{lemma}

\begin{proof}
取概形 $V$ 和满平展态射 $V \to Y$。
取概形 $W$ 和满平展态射 $W \to P \times_Y V$。
令 $U' = U \times_P W$；它是一个在 $U$ 上平展的概形。
我们必须证明 $U' \to W$ 是概形的 Koszul-正则浸入；见
定义 \ref{spaces-more-morphisms-definition-regular-immersion}。
由上面的定义 \ref{spaces-more-morphisms-definition-lci}，概形态射 $U' \to V$
是局部完全交态射。因此结论由《态射进阶》引理
\ref{more-morphisms-lemma-lci} 得到。
\end{proof}

\noindent
在此汇集所有 Koszul 态射共有的一些性质看来是合适的。

\begin{lemma}
\label{spaces-more-morphisms-lemma-lci-properties}
设 $S$ 为概形。设 $f : X \to Y$ 为 $S$ 上代数空间的
局部完全交态射。则
\begin{enumerate}
\item $f$ 局部有限表示；
\item $f$ 拟凝聚；
\item $f$ 完美。
\end{enumerate}
\end{lemma}

\begin{proof}
从略。提示：使用本引理的概形版本；见《态射进阶》引理
\ref{more-morphisms-lemma-lci-properties}。
\end{proof}

\noindent
要注意，Koszul 态射的基变换一般并非 Koszul。

\begin{lemma}
\label{spaces-more-morphisms-lemma-flat-base-change-lci}
局部完全交态射的平坦基变换是局部完全交态射。
\end{lemma}

\begin{proof}
从略。提示：使用本引理的概形版本；见《态射进阶》引理
\ref{more-morphisms-lemma-flat-base-change-lci}。
\end{proof}

\begin{lemma}
\label{spaces-more-morphisms-lemma-composition-lci}
局部完全交态射的复合是局部完全交态射。
\end{lemma}

\begin{proof}
从略。提示：使用本引理的概形版本；见《态射进阶》引理
\ref{more-morphisms-lemma-composition-lci}。
\end{proof}

\begin{lemma}
\label{spaces-more-morphisms-lemma-flat-lci}
\begin{slogan}
Syntomic 等于平坦加 lci（对代数空间而言）。
\end{slogan}
设 $S$ 为概形。设 $f : X \to Y$ 为 $S$ 上代数空间的态射。
下列条件等价：
\begin{enumerate}
\item $f$ 平坦且为局部完全交态射；
\item $f$ 为 syntomic。
\end{enumerate}
\end{lemma}

\begin{proof}
从略。提示：使用本引理的概形版本；见《态射进阶》引理
\ref{more-morphisms-lemma-flat-lci}。
\end{proof}

\begin{lemma}
\label{spaces-more-morphisms-lemma-regular-immersion-lci}
设 $S$ 为概形。$S$ 上代数空间的 Koszul-正则浸入是局部完全交态射。
\end{lemma}

\begin{proof}
设 $i : X \to Y$ 为 $S$ 上代数空间的 Koszul-正则浸入。
按定义，存在满平展态射 $V \to Y$，其中 $V$ 为概形，
使得 $X \times_Y V$ 是概形，并且基变换
$X \times_Y V \to V$ 是概形的 Koszul-正则浸入。
由《态射进阶》引理
\ref{more-morphisms-lemma-regular-immersion-lci} 可知，
$X \times_Y V \to V$ 是局部完全交态射。
由定义 \ref{spaces-more-morphisms-definition-lci} 得到 $i$ 是代数空间的局部完全交态射。
\end{proof}

\begin{lemma}
\label{spaces-more-morphisms-lemma-lci-permanence}
设 $S$ 为概形。设
$$
\xymatrix{
X \ar[rr]_f \ar[rd] & & Y \ar[ld] \\
& Z
}
$$
为 $S$ 上代数空间态射的交换图。假设 $Y \to Z$ 光滑，
并且 $X \to Z$ 是局部完全交态射。则
$f : X \to Y$ 是局部完全交态射。
\end{lemma}

\begin{proof}
取概形 $W$ 和满平展态射 $W \to Z$。
取概形 $V$ 和满平展态射 $V \to W \times_Z Y$。
取概形 $U$ 和满平展态射 $U \to V \times_Y X$。
于是 $U \to W$ 是概形的局部完全交态射，且
$V \to W$ 是概形的光滑态射。由概形情形的结论
（《态射进阶》引理 \ref{more-morphisms-lemma-lci-permanence}），
得到 $U \to V$ 是局部完全交态射。按定义，这意味着
$f$ 是局部完全交态射。
\end{proof}

\begin{lemma}
\label{spaces-more-morphisms-lemma-descending-property-lci}
性质 $\mathcal{P}(f) =$“$f$ 是局部完全交态射”在基底上 fpqc 局部。
\end{lemma}

\begin{proof}
设 $S$ 为概形。设 $f : X \to Y$ 为 $S$ 上代数空间的态射。
设 $\{Y_i \to Y\}$ 为 fpqc 覆盖（《空间上的拓扑》定义
\ref{spaces-topologies-definition-fpqc-covering}）。
设 $f_i : X_i \to Y_i$ 为 $f$ 沿 $Y_i \to Y$ 的基变换。
如果 $f$ 是局部完全交态射，则由引理
\ref{spaces-more-morphisms-lemma-flat-base-change-lci}，每个 $f_i$ 都是局部完全交态射。

\medskip\noindent
反过来，假设每个 $f_i$ 都是局部完全交态射。
可以用一个加细替换该覆盖（这仍是因为平坦基变换保持
为局部完全交态射这一性质）。因此可以假设 $Y_i$ 对每个 $i$ 都是概形；见
《空间上的拓扑》引理 \ref{spaces-topologies-lemma-refine-fpqc-schemes}。
取概形 $V$ 和满平展态射 $V \to Y$。
取概形 $U$ 和满平展态射 $U \to V \times_Y X$。
我们必须证明 $U \to V$ 是概形的局部完全交态射。
由《空间上的拓扑》引理
\ref{spaces-topologies-lemma-recognize-fpqc-covering}，
$\{Y_i \times_Y V \to V\}$ 是概形的 fpqc 覆盖。
由概形情形（《态射进阶》引理
\ref{more-morphisms-lemma-descending-property-lci}），只需证明
% P09-ERR-0747
$$
U \times_Y Y_i = U \times_V (V \times_Y Y_i) \longrightarrow V
$$
这一 $U \to V$ 沿 $V \times_Y Y_i \to V$ 的基变换
是局部完全交态射。可以把它写成如下复合：
$$
U \times_Y Y_i \longrightarrow
(V \times_Y X) \times_Y Y_i = V \times_Y X_i \longrightarrow
V \times_Y Y_i
$$
第一支箭头是概形的平展态射（因为它是
$U \to V \times_Y X$ 的基变换），第二支箭头是概形的
局部完全交态射，因为它是 $f_i$ 的平坦基变换。
结论成立，因为为局部完全交态射这一性质在源上 syntomic 局部，
并且平展态射是 syntomic（《态射进阶》引理
\ref{more-morphisms-lemma-lci-syntomic-local-source}
及《态射》引理 \ref{morphisms-lemma-etale-syntomic}）。
\end{proof}

\begin{lemma}
\label{spaces-more-morphisms-lemma-lci-syntomic-local-source}
性质 $\mathcal{P}(f) =$“$f$ 是局部完全交态射”在源上 syntomic 局部。
\end{lemma}

\begin{proof}
这由《空间上的下降》引理
\ref{spaces-descent-lemma-transfer-from-schemes}
和《态射进阶》引理
\ref{more-morphisms-lemma-lci-syntomic-local-source} 得到。
\end{proof}

\begin{lemma}
\label{spaces-more-morphisms-lemma-base-change-lci-fibres}
设 $S$ 为概形。考虑代数空间的交换图
$$
\xymatrix{
X \ar[rr]_f \ar[rd]_p & & Y \ar[ld]^q \\
& Z
}
$$
其中各态射在 $S$ 上。
假设 $p$ 和 $q$ 都平坦且局部有限表示。
则存在开子空间 $U(f) \subset X$，使得 $|U(f)| \subset |X|$
是 $f$ 为 Koszul 之点的集合。此外，对代数空间的任意态射
$Z' \to Z$，如果 $f' : X' \to Y'$ 是 $f$ 沿 $Z' \to Z$
的基变换，则 $U(f')$ 是 $U(f)$ 在投影 $X' \to X$ 下的逆像。
\end{lemma}

\begin{proof}
本引理是《态射进阶》引理
\ref{more-morphisms-lemma-base-change-lci-fibres} 的对应版本，
事实上我们将由该引理推出本引理。按定义 \ref{spaces-more-morphisms-definition-lci}，
集合 $\{x \in |X| : f \text{ 为 Koszul 于点 }x\}$
在 $|X|$ 中开，因而由《空间的性质》引理
\ref{spaces-properties-lemma-open-subspaces}，
它对应于开子空间 $U(f)$；这是 $X$ 的一个开子空间。
所以只需证明最后一个断言。

\medskip\noindent
取概形 $W$ 和满平展态射 $W \to Z$。
取概形 $V$ 和满平展态射 $V \to W \times_Z Y$。
取概形 $U$ 和满平展态射 $U \to V \times_Y X$。
最后，取概形 $W'$ 和满平展态射 $W' \to W \times_Z Z'$。
令 $V' = W' \times_W V$ 且 $U' = W' \times_W U$，
从而得到满平展态射 $V' \to Y'$ 和 $U' \to X'$。
以下将不再说明地使用这一事实：代数空间的平展态射在相伴拓扑空间上
诱导开映射（见《空间的性质》引理
\ref{spaces-properties-lemma-etale-open}）。
注意，按定义，$U(f)$ 在 $|X|$ 中表示集合 $T$ 的像；
该集合由 $U$ 中使概形态射 $U \to V$ 为 Koszul 的那些点组成。
类似地，$U(f')$ 在 $|X'|$ 中表示集合 $T'$ 的像；
该集合由 $U'$ 中使概形态射 $U' \to V'$ 为 Koszul 的那些点组成。
现在，按构造，图
$$
\xymatrix{
U' \ar[r] \ar[d] & U \ar[d] \\
V' \ar[r] & V
}
$$
是笛卡尔的（在概形范畴中）。因此，上述《态射进阶》引理
\ref{more-morphisms-lemma-base-change-lci-fibres} 适用，并说明
$T'$ 是 $T$ 的逆像。由于 $|U'| \to |X'|$ 满，这就推出了本引理。
\end{proof}

\begin{lemma}
\label{spaces-more-morphisms-lemma-unramified-lci}
设 $S$ 为概形。设 $f : X \to Y$ 为 $S$ 上代数空间的
局部完全交态射。则 $f$ 非分歧，当且仅当 $f$ 形式非分歧；
在这种情况下，余法层 $\mathcal{C}_{X/Y}$ 在 $X$ 上有限局部自由。
\end{lemma}

\begin{proof}
这由概形态射的相应结论经平展局部化得到；见《态射进阶》引理
\ref{more-morphisms-lemma-unramified-lci} 和引理
\ref{spaces-more-morphisms-lemma-universal-thickening-localize}。
（注意，在该引理的情形中，态射 $V \to U$
按定义非分歧且为局部完全交态射。）
\end{proof}

\begin{lemma}
\label{spaces-more-morphisms-lemma-transitivity-conormal-lci}
设 $S$ 为概形。设 $Z \to Y \to X$ 为 $S$ 上代数空间的
形式非分歧态射。假设 $Z \to Y$ 是局部完全交态射。
% P09-ERR-0748
引理 \ref{spaces-more-morphisms-lemma-transitivity-conormal} 的正合序列
$$
0 \to i^*\mathcal{C}_{Y/X} \to
\mathcal{C}_{Z/X} \to
\mathcal{C}_{Z/Y} \to 0
$$
是短正合的。
\end{lemma}

\begin{proof}
取概形 $U$ 和满平展态射 $U \to X$。
取概形 $V$ 和满平展态射 $V \to U \times_X Y$。
取概形 $W$ 和满平展态射 $W \to V \times_Y Z$。
由引理 \ref{spaces-more-morphisms-lemma-universal-thickening-localize}，
态射 $W \to V$ 和 $V \to U$ 都形式非分歧。此外，序列
$i^*\mathcal{C}_{Y/X} \to \mathcal{C}_{Z/X} \to \mathcal{C}_{Z/Y} \to 0$
限制为相应序列
$i^*\mathcal{C}_{V/U} \to \mathcal{C}_{W/U} \to \mathcal{C}_{W/V} \to 0$；
这是针对 $W \to V \to U$ 而言的。
所以，由概形情形的结论（《态射进阶》引理
\ref{more-morphisms-lemma-transitivity-conormal-lci}）可得所求，
因为按定义，态射 $W \to V$ 是局部完全交态射。
\end{proof}

\section{一个态射何时是同构？}
\label{spaces-more-morphisms-section-when-isomorphism}

\noindent
更一般地，可以问：“一个态射何时具有性质 $\mathcal{P}$？”
更精确的问题如下。假设给定代数空间的交换图
$$
\xymatrix{
X \ar[rr]_f \ar[rd]_p & & Y \ar[ld]^q \\
& Z
}
$$
是否存在代数空间的单态射 $W \to Z$，使得下列两个性质成立：
\begin{enumerate}
\item 基变换 $f_W : X_W \to Y_W$ 具有性质 $\mathcal{P}$；
\item 代数空间的任意态射 $Z' \to Z$ 经由 $W$ 分解，
当且仅当基变换 $f_{Z'} : X_{Z'} \to Y_{Z'}$ 具有性质
$\mathcal{P}$。
\end{enumerate}
在许多情形中，如果 $W \to Z$ 存在，那么它是浸入、开浸入或闭浸入。

\medskip\noindent
这一问题的答案可能依赖于态射 $f$、$p$ 和 $q$ 的辅助性质。
例如 $\mathcal{P}(f) =$“$f$ 平坦”；在 $Y = S$ 的情形中，
我们已经在“平坦性进阶”一章对概形态射详细讨论过这一例子，
起点是《平坦性进阶》第 \ref{flat-section-flattening-functors} 节。

\begin{lemma}
\label{spaces-more-morphisms-lemma-where-unramified}
考虑代数空间的交换图
$$
\xymatrix{
X \ar[rr]_f \ar[rd]_p & & Y \ar[ld]^q \\
& Z
}
$$
假设 $p$ 局部有限型且为闭映射。则存在开子空间 $W \subset Z$，
使得态射 $Z' \to Z$ 经由 $W$ 分解，当且仅当基变换
$f_{Z'} : X_{Z'} \to Y_{Z'}$ 非分歧。
\end{lemma}

\begin{proof}
由《空间的态射》引理
\ref{spaces-morphisms-lemma-where-unramified}，
存在开子空间 $U(f) \subset X$，它是使 $f$ 非分歧的点的集合。
此外，$U(f)$ 的构造与任意基变换交换。令 $W \subset Z$
为如下开子空间（见《空间的性质》引理
\ref{spaces-properties-lemma-open-subspaces}），其底层点集为
$$
|W| = |Z| \setminus |p|\left(|X| \setminus |U(f)|\right)
$$
也就是说，$z \in |Z|$ 是 $W$ 的点，当且仅当 $f$
在 $X$ 中位于 $z$ 上方的每一点处都非分歧。
由于假设 $p$ 为闭映射，该集合为开。因为 $U(f)$ 的构造与任意
基变换交换，立即可见（使用《空间的性质》引理
\ref{spaces-properties-lemma-factor-through-open-subspace}）
$W$ 具有所需的泛性质。
\end{proof}

\begin{lemma}
\label{spaces-more-morphisms-lemma-where-unramified-universally-injective}
考虑代数空间的交换图
$$
\xymatrix{
X \ar[rr]_f \ar[rd]_p & & Y \ar[ld]^q \\
& Z
}
$$
假设
\begin{enumerate}
\item $p$ 局部有限型；
\item $p$ 为闭映射；
\item $p_2 : X \times_Y X \to Z$ 为闭映射。
\end{enumerate}
则存在开子空间 $W \subset Z$，使得态射 $Z' \to Z$ 经由 $W$
分解，当且仅当基变换 $f_{Z'} : X_{Z'} \to Y_{Z'}$
非分歧且普遍单射。
\end{lemma}

\begin{proof}
用引理 \ref{spaces-more-morphisms-lemma-where-unramified} 中得到的开子空间替换 $Z$ 后，
可以假设 $f$ 已经非分歧；注意，这不会破坏假设 (2) 或 (3)。
由《空间的态射》引理
\ref{spaces-morphisms-lemma-diagonal-unramified-morphism}，
可知 $\Delta_{X/Y} : X \to X \times_Y X$ 是开浸入。
任意基变换后这一点仍成立。因此，由《空间的态射》引理
\ref{spaces-morphisms-lemma-universally-injective}，
$f_{Z'}$ 普遍单射，当且仅当对角线的基变换
$X_{Z'} \to (X \times_Y X)_{Z'}$ 是同构。令 $W \subset Z$
为如下开子空间（见《空间的性质》引理
\ref{spaces-properties-lemma-open-subspaces}），其底层点集为
$$
|W| = |Z| \setminus
|p_2|\left(|X \times_Y X| \setminus \Im(|\Delta_{X/Y}|)\right)
$$
也就是说，$z \in |Z|$ 是 $W$ 的点，当且仅当
$|X \times_Y X| \to |Z|$ 在 $z$ 上的纤维包含于
$|X| \to |X \times_Y X|$ 的像中。由上面的讨论立即可知，
限制 $p^{-1}(W) \to q^{-1}(W)$（即 $f$ 的限制）非分歧且普遍单射。

\medskip\noindent
反过来，假设 $f_{Z'}$ 非分歧且普遍单射。为了证明
$Z' \to Z$ 经由 $W$ 分解，只需证明 $|Z'| \to |Z|$
的像包含于 $|W|$；见《空间的性质》引理
\ref{spaces-properties-lemma-factor-through-open-subspace}。
因此，只需证明 $Z'$ 是一个域的谱时的结论。
% P09-ERR-0749
记 $z \in |Z|$ 为 $|Z'| \to |Z|$ 的像。上面的讨论说明
$$
|X_{Z'}| \longrightarrow |(X \times_Y X)_{Z'}|
$$
是满的。由《空间的性质》引理
\ref{spaces-properties-lemma-points-cartesian}，在交换图
$$
\xymatrix{
|X_{Z'}| \ar[d] \ar[r] &
|(X \times_Y X)_{Z'}| \ar[d] \\
|p|^{-1}(\{z\}) \ar[r] &
|p_2|^{-1}(\{z\})
}
$$
中，竖直箭头均为满的。于是得到所需的 $z \in |W|$。
\end{proof}

\begin{lemma}
\label{spaces-more-morphisms-lemma-where-closed-immersion}
考虑代数空间的交换图
$$
\xymatrix{
X \ar[rr]_f \ar[rd]_p & & Y \ar[ld]^q \\
& Z
}
$$
假设
\begin{enumerate}
\item $p$ 局部有限型；
\item $p$ 普遍闭；
\item $q : Y \to Z$ 可分。
\end{enumerate}
则存在开子空间 $W \subset Z$，使得态射 $Z' \to Z$ 经由 $W$
分解，当且仅当基变换 $f_{Z'} : X_{Z'} \to Y_{Z'}$ 是闭浸入。
\end{lemma}

\begin{proof}
我们将使用闭浸入作为普遍闭、非分歧且普遍单射态射的刻画；见引理
\ref{spaces-more-morphisms-lemma-characterize-closed-immersion}。
首先，注意到 $p$ 普遍闭且 $q$ 可分，所以 $f$ 普遍闭；见
《空间的态射》引理
\ref{spaces-morphisms-lemma-universally-closed-permanence}。
于是 $f$ 的任意基变换都普遍闭；见《空间的态射》引理
\ref{spaces-morphisms-lemma-base-change-universally-closed}。
因此，为完成本引理的证明，只需证明引理
\ref{spaces-more-morphisms-lemma-where-unramified-universally-injective} 的假设成立。
投影 $\text{pr}_0 : X \times_Y X \to X$ 作为 $f$ 的基变换而普遍闭；
见《空间的态射》引理
\ref{spaces-morphisms-lemma-base-change-universally-closed}。
所以 $X \times_Y X \to Z$ 作为普遍闭态射的复合而普遍闭
（见《空间的态射》引理
\ref{spaces-morphisms-lemma-composition-universally-closed}）。
这就完成了本引理的证明。
\end{proof}

\begin{lemma}
\label{spaces-more-morphisms-lemma-where-flat}
考虑代数空间的交换图
$$
\xymatrix{
X \ar[rr]_f \ar[rd]_p & & Y \ar[ld]^q \\
& Z
}
$$
假设
\begin{enumerate}
\item $p$ 局部有限表示；
\item $p$ 平坦；
\item $p$ 为闭映射；
\item $q$ 局部有限型。
\end{enumerate}
则存在开子空间 $W \subset Z$，使得态射 $Z' \to Z$ 经由 $W$
分解，当且仅当基变换 $f_{Z'} : X_{Z'} \to Y_{Z'}$ 平坦。
\end{lemma}

\begin{proof}
由引理 \ref{spaces-more-morphisms-lemma-base-change-flatness-fibres}，集合
$$
A = \{x \in |X| : X\text{ 在点 }x \text{ 处平坦于 }Y\}.
$$
在 $|X|$ 中开，并且它的构造与任意基变换交换。令
$W \subset Z$ 为如下开子空间（见《空间的性质》引理
\ref{spaces-properties-lemma-open-subspaces}），其底层点集为
$$
|W| = |Z| \setminus |p|\left(|X| \setminus A\right)
$$
也就是说，$z \in |Z|$ 是 $W$ 的点，当且仅当
$|X| \to |Z|$ 在 $z$ 上的整个纤维包含于 $A$。
由于 $p$ 为闭映射，该集合为开。因为 $A$ 的构造与任意
基变换交换，所以 $W$ 符合要求。
\end{proof}

\begin{lemma}
\label{spaces-more-morphisms-lemma-where-surjective-flat}
考虑代数空间的交换图
$$
\xymatrix{
X \ar[rr]_f \ar[rd]_p & & Y \ar[ld]^q \\
& Z
}
$$
假设
\begin{enumerate}
\item $p$ 局部有限表示；
\item $p$ 平坦；
\item $p$ 为闭映射；
\item $q$ 局部有限型；
\item $q$ 为闭映射。
\end{enumerate}
则存在开子空间 $W \subset Z$，使得态射 $Z' \to Z$ 经由 $W$
分解，当且仅当基变换 $f_{Z'} : X_{Z'} \to Y_{Z'}$ 满且平坦。
\end{lemma}

\begin{proof}
由引理 \ref{spaces-more-morphisms-lemma-where-flat}，可以假设 $f$ 平坦。
由《空间的态射》引理
\ref{spaces-morphisms-lemma-finite-presentation-permanence}，
注意到 $f$ 局部有限表示。所以 $f$ 为开映射；见《空间的态射》引理
\ref{spaces-morphisms-lemma-fppf-open}。令 $W \subset Z$
为如下开子空间（见《空间的性质》引理
\ref{spaces-properties-lemma-open-subspaces}），其底层点集为
$$
|W| = |Z| \setminus |q|\left(|Y| \setminus |f|(|X|)\right).
$$
换言之，对 $z \in |Z|$，有 $z \in |W|$，当且仅当
$|Y| \to |Z|$ 在 $z$ 上的整个纤维包含于
$|X| \to |Y|$ 的像。由于 $q$ 为闭映射，该集合在 $|Z|$ 中开。
按构造，态射 $X_W \to Y_W$ 是满的。最后，假设
$X_{Z'} \to Y_{Z'}$ 是满的。为了证明 $Z' \to Z$ 经由 $W$ 分解，
只需证明 $|Z'| \to |Z|$ 的像包含于 $|W|$；见《空间的性质》引理
\ref{spaces-properties-lemma-factor-through-open-subspace}。
因此，只需证明 $Z'$ 是一个域的谱时的结论。
% P09-ERR-0750
记 $z \in |Z|$ 为 $|Z'| \to |Z|$ 的像。由《空间的性质》引理
\ref{spaces-properties-lemma-points-cartesian}，在交换图
$$
\xymatrix{
|X_{Z'}| \ar[d] \ar[r] &
|Y_{Z'}| \ar[d] \\
|p|^{-1}(\{z\}) \ar[r] &
|q|^{-1}(\{z\})
}
$$
中，竖直箭头均为满的。于是得到所需的 $z \in |W|$。
\end{proof}

\begin{lemma}
\label{spaces-more-morphisms-lemma-where-isomorphism}
考虑代数空间的交换图
$$
\xymatrix{
X \ar[rr]_f \ar[rd]_p & & Y \ar[ld]^q \\
& Z
}
$$
假设
\begin{enumerate}
\item $p$ 局部有限表示；
\item $p$ 平坦；
\item $p$ 普遍闭；
\item $q$ 局部有限型；
\item $q$ 为闭映射；
\item $q$ 可分。
\end{enumerate}
则存在开子空间 $W \subset Z$，使得态射 $Z' \to Z$ 经由 $W$
分解，当且仅当基变换 $f_{Z'} : X_{Z'} \to Y_{Z'}$ 是同构。
\end{lemma}

\begin{proof}
由引理 \ref{spaces-more-morphisms-lemma-where-surjective-flat}，存在开子空间
$W_1 \subset Z$，使得 $f_{Z'}$ 满且平坦，当且仅当
$Z' \to Z$ 经由 $W_1$ 分解。由引理
\ref{spaces-more-morphisms-lemma-where-closed-immersion}，存在开子空间
$W_2 \subset Z$，使得 $f_{Z'}$ 是闭浸入，当且仅当
$Z' \to Z$ 经由 $W_2$ 分解。我们断言 $W = W_1 \cap W_2$ 符合要求。
当然，如果 $f_{Z'}$ 是同构，那么 $Z' \to Z$ 经由 $W$ 分解。
所以只需证明 $f_W$ 是同构。按构造，$f_W$ 是满、平坦的闭浸入。
特别地，$f_W$ 可表。由于概形的满平坦闭浸入是同构
（见《态射》引理
\ref{morphisms-lemma-characterize-flat-closed-immersions}），结论成立。
（注意，事实上 $f_W$ 局部有限表示，从而为开映射；因此若愿意，
可以避免使用这个引理。）
\end{proof}

\begin{lemma}
\label{spaces-more-morphisms-lemma-where-lci}
考虑代数空间的交换图
$$
\xymatrix{
X \ar[rr]_f \ar[rd]_p & & Y \ar[ld]^q \\
& Z
}
$$
假设
\begin{enumerate}
\item $p$ 平坦且局部有限表示；
\item $p$ 为闭映射；
\item $q$ 平坦且局部有限表示。
\end{enumerate}
则存在开子空间 $W \subset Z$，使得态射 $Z' \to Z$ 经由 $W$
分解，当且仅当基变换 $f_{Z'} : X_{Z'} \to Y_{Z'}$
是局部完全交态射。
\end{lemma}

\begin{proof}
由引理 \ref{spaces-more-morphisms-lemma-base-change-lci-fibres}，存在开子空间
$U(f) \subset X$，它是使 $f$ 为 Koszul 的点的集合。
此外，$U(f)$ 的构造与任意基变换交换。令 $W \subset Z$
为如下开子空间（见《空间的性质》引理
\ref{spaces-properties-lemma-open-subspaces}），其底层点集为
$$
|W| = |Z| \setminus |p|\left(|X| \setminus |U(f)|\right)
$$
也就是说，$z \in |Z|$ 是 $W$ 的点，当且仅当 $f$
在 $X$ 中位于 $z$ 上方的每一点处都为 Koszul。
由于假设 $p$ 为闭映射，该集合为开。因为 $U(f)$ 的构造与任意
基变换交换，立即可见（使用《空间的性质》引理
\ref{spaces-properties-lemma-factor-through-open-subspace}）
$W$ 具有所需的泛性质。
\end{proof}

\section{微分层与余法层的正合列}
\label{spaces-more-morphisms-section-exact}

\noindent
本节汇集关于余法层和微分层的正合列的一些结果。
在某种意义下，这些都是与代数空间的可复合态射相伴的
余切复形三角形的实现。

\medskip\noindent
% P09-ERR-0751
下面各个序列中的每个映射都是在引理
\ref{spaces-more-morphisms-lemma-functoriality-differentials} 或引理
\ref{spaces-more-morphisms-lemma-universal-thickening-functorial} 中构造的。
设 $S$ 为概形。设 $g : Z \to Y$ 和 $f : Y \to X$
为 $S$ 上代数空间的态射。
\begin{enumerate}
\item 有典范正合列
$$
g^*\Omega_{Y/X} \to \Omega_{Z/X} \to \Omega_{Z/Y} \to 0,
$$
见引理 \ref{spaces-more-morphisms-lemma-triangle-differentials}。
如果 $g : Z \to Y$ 形式光滑，那么这个序列是短正合列；见引理
\ref{spaces-more-morphisms-lemma-triangle-differentials-formally-smooth}。
\item 如果 $g$ 形式非分歧，那么有典范正合列
$$
\mathcal{C}_{Z/Y} \to g^*\Omega_{Y/X} \to \Omega_{Z/X} \to 0,
$$
见引理 \ref{spaces-more-morphisms-lemma-universally-unramified-differentials-sequence}。
如果 $f \circ g : Z \to X$ 形式光滑，那么这个序列是短正合列；见引理
\ref{spaces-more-morphisms-lemma-differentials-formally-unramified-formally-smooth}。
\item 如果 $g$ 和 $f \circ g$ 都形式非分歧，那么有典范正合列
$$
\mathcal{C}_{Z/X} \to \mathcal{C}_{Z/Y} \to g^*\Omega_{Y/X} \to 0,
$$
见引理 \ref{spaces-more-morphisms-lemma-two-unramified-morphisms}。
如果 $f : Y \to X$ 形式光滑，那么这个序列是短正合列；见引理
\ref{spaces-more-morphisms-lemma-two-unramified-morphisms-formally-smooth}。
\item 如果 $g$ 和 $f$ 都形式非分歧，那么有典范正合列
$$
g^*\mathcal{C}_{Y/X} \to \mathcal{C}_{Z/X} \to \mathcal{C}_{Z/Y} \to 0.
$$
见引理 \ref{spaces-more-morphisms-lemma-transitivity-conormal-unramified}。
如果 $g : Z \to Y$ 是局部完全交态射，那么这个序列是短正合列；
见引理 \ref{spaces-more-morphisms-lemma-transitivity-conormal-lci}。
\end{enumerate}

\section{拟凝聚复形的刻画，II}
\label{spaces-more-morphisms-section-characterize-pseudo-coherent}

\noindent
本节用上同调讨论拟凝聚复形的一种刻画。关于代数空间上拟凝聚复形的
早先材料可见《空间的导出范畴》第
\ref{spaces-perfect-section-spell-out} 节和《空间的导出范畴》第
\ref{spaces-perfect-section-pseudo-coherent-hocolim} 节。
本节在概形情形的对应版本是《态射进阶》第
\ref{more-morphisms-section-characterize-pseudo-coherent} 节。
一个基本工具是利用 Chow 引理的导出版本约化到射影空间的情形；
见引理 \ref{spaces-more-morphisms-lemma-derived-chow}。

\begin{lemma}
\label{spaces-more-morphisms-lemma-case-of-tor-independence}
设 $S$ 为概形。考虑代数空间的交换图
$$
\xymatrix{
Z' \ar[d] \ar[r] & Y' \ar[d] \\
X' \ar[r] & B'
}
$$
它在 $S$ 上。设 $B \to B'$ 为态射。记 $X$ 和 $Y$ 分别为
$X'$ 和 $Y'$ 到 $B$ 的基变换。假设 $Y' \to B'$ 和
$Z' \to X'$ 平坦。则 $X \times_B Y$ 和 $Z'$
在 $X' \times_{B'} Y'$ 上 Tor 独立。
\end{lemma}

\begin{proof}
由《空间的导出范畴》引理
\ref{spaces-perfect-lemma-tor-independent}，可以在
$X \times_B Y$ 和 $Z'$ 上平展局部地检验 Tor 独立性。
这\footnote{更详细的论证如下。取满平展态射 $W' \to B'$，
其中 $W'$ 为概形。取满平展态射 $W \to B \times_{B'} W'$，
其中 $W$ 为概形。取满平展态射
$U' \to X' \times_{B'} W'$，其中 $U'$ 为概形。取满平展态射
$V' \to Y' \times_{B'} W'$，其中 $V'$ 为概形。
注意，$U' \times_{W'} V' \to X' \times_{B'} Y'$ 满平展。
取满平展态射
$T' \to Z' \times_{X' \times_{B'} Y'} U' \times_{W'} V'$，
其中 $T'$ 为概形。
% P09-ERR-0752
记 $U$ 和 $V$ 为 $U'$ 和 $V'$ 到 $W$ 的基变换。
于是本引理断言：作为代数空间，$X \times_B Y$ 和 $Z'$
在 $X' \times_{B'} Y'$ 上 Tor 独立，当且仅当作为概形，
$U \times_W V$ 和 $T'$ 在 $U' \times_{W'} V'$ 上 Tor 独立。
因此，只需对以 $T', U', V', W'$ 为四个角的方块及基变换
$W \to W'$ 证明本引理。$Y' \to B'$ 和 $Z' \to X'$ 的平坦性
蕴含 $V' \to W'$ 和 $T' \to U'$ 的平坦性。}
把本引理约化到概形的情形，而该情形就是《态射进阶》引理
\ref{more-morphisms-lemma-case-of-tor-independence}。
\end{proof}

\begin{lemma}[导出 Chow 引理]
\label{spaces-more-morphisms-lemma-derived-chow}
设 $A$ 为环。设 $X$ 为 $A$ 上可分且有限表示的代数空间。
设 $x \in |X|$。则存在一个 $n \geq 0$、一个闭子空间
$Z \subset X \times_A \mathbf{P}^n_A$、一个点 $z \in |Z|$、
一个开集 $V \subset \mathbf{P}^n_A$，以及一个对象
$E$（它属于 $D(\mathcal{O}_{X \times_A \mathbf{P}^n_A})$），使得
\begin{enumerate}
\item $Z \to X \times_A \mathbf{P}^n_A$ 有限表示；
\item $c : Z \to \mathbf{P}^n_A$ 在 $V$ 上是闭浸入；
令 $W = c^{-1}(V)$；
\item $b : Z \to X$ 到 $W$ 的限制平展，
$z \in W$，并且 $b(z) = x$；
\item $E|_{X \times_A V} \cong
(b, c)_*\mathcal{O}_Z|_{X \times_A V}$；
\item $E$ 拟凝聚且支撑在 $Z$ 上。
\end{enumerate}
\end{lemma}

\begin{proof}
可以找到有限型 $\mathbf{Z}$-子代数 $A' \subset A$
以及代数空间 $X'$，它在 $A'$ 上可分且有限表示，并使它到 $A$
的基变换为 $X$。见《空间的极限》引理
\ref{spaces-limits-lemma-descend-finite-presentation} 和
\ref{spaces-limits-lemma-descend-separated-morphism}。
设 $x' \in |X'|$ 为 $x$ 的像。如果能对 $(X'/A', x')$
证明本引理，那么对 $(X/A, x)$ 的结论随即成立。
确实，如果 $n', Z', z', V', E'$ 给出 $(X'/A', x')$ 的解，
那么可以令 $n = n'$；令 $Z \subset X \times \mathbf{P}^n$
为 $Z'$ 的逆像；
% P09-ERR-0753
令 $z \in Z$ 为映到 $x$ 的唯一点；令
$V \subset \mathbf{P}^n_A$ 为 $V'$ 的逆像；并令 $E$
为 $E'$ 的导出拉回。注意，由《位点上的上同调》引理
\ref{sites-cohomology-lemma-pseudo-coherent-pullback}，$E$ 拟凝聚。
% P09-ERR-0754
只需检验 (5)。为此，令 $W = c^{-1}(V)$ 和
$W' = (c')^{-1}(V')$，并考虑笛卡尔方块
$$
\xymatrix{
W \ar[d]_{(b, c)} \ar[r] & W' \ar[d]^{(b', c')} \\
X \times_A V \ar[r] & X' \times_{A'} V'
}
$$
由引理 \ref{spaces-more-morphisms-lemma-case-of-tor-independence}，$X \times_A V$ 和 $W'$
在 $X' \times_{A'} V'$ 上 Tor 独立。因此，由《空间的导出范畴》引理
\ref{spaces-perfect-lemma-compare-base-change}，
$(b', c')_*\mathcal{O}_{W'}$ 到 $X \times_A V$ 的导出拉回是
$(b, c)_*\mathcal{O}_W$。这里还使用了
$R(b', c')_*\mathcal{O}_{Z'} = (b', c')_*\mathcal{O}_{Z'}$，
因为 $(b', c')$ 是闭浸入；对 $(b, c)_*\mathcal{O}_Z$ 也同样如此。
% P09-ERR-0755
由于 $E'|_{U' \times_{A'} V'} =
(b', c')_*\mathcal{O}_{W'}$，得到
$E|_{U \times_A V} = (b, c)_*\mathcal{O}_W$，从而 (5) 成立。
这把我们约化到下一段所述的情形。

\medskip\noindent
假设 $A$ 在 $\mathbf{Z}$ 上有限型。取平展态射 $U \to X$，
其中 $U$ 为仿射概形，并取一点 $u \in U$，它映到 $x$。
于是 $U$ 在 $A$ 上有限型。取闭浸入 $U \to \mathbf{A}^n_A$，并记
% P09-ERR-0756
$j : U \to \mathbf{P}^n_A$ 为将其与开浸入
$\mathbf{A}^n_A \to \mathbf{P}^n_A$ 复合而得到的浸入。
令 $Z$ 为
$$
(\text{id}_U, j) : U \longrightarrow X \times_A \mathbf{P}^n_A
$$
的概形论闭包。设 $z \in Z$ 为 $u$ 的像。
设 $Y \subset \mathbf{P}^n_A$ 为 $j$ 的概形论闭包。
于是显然，$Z \subset X \times_A Y$ 是
$(\text{id}_U, j) : U \to X \times_A Y$ 的概形论闭包。
由于 $X$ 可分，态射 $X \times_A Y \to Y$ 也可分。
因此，由《空间的态射》引理
\ref{spaces-morphisms-lemma-scheme-theoretic-image-of-partial-section}，
可见 $Z \to Y$ 在开子概形 $j(U) \subset Y$ 上是同构。
取开集 $V \subset \mathbf{P}^n_A$，使 $V \cap Y = j(U)$。
于是 (2) 成立，且 $W = (\text{id}_U, j)(U)$，从而 (3) 成立。
(1) 成立是因为 $A$ 诺特。

\medskip\noindent
因为 $A$ 诺特，所以 $X$ 和 $X \times_A \mathbf{P}^n_A$
都是诺特代数空间。因此在这种情形可以取
$E = (b, c)_*\mathcal{O}_Z$：(4) 是显然的；关于 (5)，见
《空间的导出范畴》引理
\ref{spaces-perfect-lemma-identify-pseudo-coherent-noetherian}。
这就完成了证明。
\end{proof}

\begin{lemma}
\label{spaces-more-morphisms-lemma-compute-Fourier-Mukai-for-derived-chow}
设 $X/A$、$x \in |X|$ 以及 $n, Z, z, V, E$
如引理 \ref{spaces-more-morphisms-lemma-derived-chow} 中所述。对任意
$K \in D_\QCoh(\mathcal{O}_X)$，有
$$
Rq_*(Lp^*K \otimes^\mathbf{L} E)|_V = R(W \to V)_*K|_W
$$
其中 $p : X \times_A \mathbf{P}^n_A \to X$ 和
$q : X \times_A \mathbf{P}^n_A \to \mathbf{P}^n_A$ 是投影，
而态射 $W \to V$ 是有限表示闭浸入 $c|_W : W \to V$。
\end{lemma}

\begin{proof}
由于 $W = c^{-1}(V)$，并且 $c$ 在 $V$ 上是闭浸入，
可知 $c|_W$ 是闭浸入。它有限表示，因为 $W$ 和 $V$
都在 $A$ 上有限表示；见《空间的态射》引理
\ref{spaces-morphisms-lemma-finite-presentation-permanence}。
首先有
$$
Rq_*(Lp^*K \otimes^\mathbf{L} E)|_V =
Rq'_*\left((Lp^*K \otimes^\mathbf{L} E)|_{X \times_A V}\right)
$$
其中 $q' : X \times_A V \to V$ 是投影，因为总正像的构造与局部化交换。
% P09-ERR-0757
记 $i = (b, c)|_W : W \to X \times_A V$ 为给定的闭浸入。则
$$
Rq'_*\left((Lp^*K \otimes^\mathbf{L} E)|_{X \times_A V}\right) =
Rq'_*(Lp^*K|_{X \times_A V} \otimes^\mathbf{L} i_*\mathcal{O}_W)
$$
这由性质 (5) 得到。由于 $i$ 是闭浸入，有
$i_*\mathcal{O}_W = Ri_*\mathcal{O}_W$。使用《空间的导出范畴》引理
\ref{spaces-perfect-lemma-cohomology-base-change}，可以把它改写为
$$
Rq'_* Ri_* Li^* Lp^*K|_{X \times_A V} =
R(q' \circ i)_* Lb^*K|_W =
R(W \to V)_* K|_W
$$
这正是所要证明的。（注意，限制到 $W$ 与经由 $W \to X$
导出拉回是同一件事，因为 $W$ 在 $X$ 上平展。）
\end{proof}

\begin{lemma}
\label{spaces-more-morphisms-lemma-characterize-pseudo-coherent}
设 $A$ 为环。设 $X$ 为 $A$ 上可分且有限表示的代数空间。
设 $K \in D_\QCoh(\mathcal{O}_X)$。如果
$R\Gamma(X, E \otimes^\mathbf{L} K)$ 是 $D(A)$ 中的拟凝聚对象，
并且这对每个拟凝聚的 $E$（它属于 $D(\mathcal{O}_X)$）都成立，
那么 $K$ 相对于 $A$ 拟凝聚
（定义 \ref{spaces-more-morphisms-definition-relative-pseudo-coherence}）。
\end{lemma}

\begin{proof}
假设 $K \in D_\QCoh(\mathcal{O}_X)$，并且
$R\Gamma(X, E \otimes^\mathbf{L} K)$ 是 $D(A)$ 中的拟凝聚对象，
其中 $E$ 是 $D(\mathcal{O}_X)$ 中任意拟凝聚对象。
设 $x \in |X|$。我们将证明 $K$ 相对于 $A$ 在 $x$
的某个平展邻域中拟凝聚。由相对拟凝聚性的定义，这将证明本引理。

\medskip\noindent
取 $n, Z, z, V, E$ 如引理 \ref{spaces-more-morphisms-lemma-derived-chow} 中所述。
记 $p : X \times \mathbf{P}^n \to X$ 和
$q : X \times \mathbf{P}^n \to \mathbf{P}^n_A$ 为投影。
于是对任意 $i \in \mathbf{Z}$，有
\begin{align*}
& R\Gamma(\mathbf{P}^n_A,
Rq_*(Lp^*K \otimes^\mathbf{L} E)
\otimes^\mathbf{L}
\mathcal{O}_{\mathbf{P}^n_A}(i)) \\
& =
R\Gamma(X \times \mathbf{P}^n,
Lp^*K \otimes^\mathbf{L} E \otimes^\mathbf{L}
Lq^*\mathcal{O}_{\mathbf{P}^n_A}(i)) \\
& =
R\Gamma(X,
K \otimes^\mathbf{L} Rp_*(E \otimes^\mathbf{L}
Lq^*\mathcal{O}_{\mathbf{P}^n_A}(i)))
\end{align*}
这由《空间的导出范畴》引理
\ref{spaces-perfect-lemma-cohomology-base-change} 得到。
由《空间的导出范畴》引理
\ref{spaces-perfect-lemma-flat-proper-pseudo-coherent-direct-image-general}，
复形 $Rp_*(E \otimes^\mathbf{L} Lq^*\mathcal{O}_{\mathbf{P}^n_A}(i))$
在 $X$ 上拟凝聚。因此，假设告诉我们，上面显示公式中的表达式
是 $D(A)$ 中的拟凝聚对象。由《概形的导出范畴》引理
\ref{perfect-lemma-pseudo-coherent-on-projective-space}，
得到 $Rq_*(Lp^*K \otimes^\mathbf{L} E)$ 在
$\mathbf{P}^n_A$ 上拟凝聚。由引理
\ref{spaces-more-morphisms-lemma-compute-Fourier-Mukai-for-derived-chow}，有
% P09-ERR-0758
$$
Rq_*(Lp^*K \otimes^\mathbf{L} E)|_{X \times_A V} =
R(W \to V)_*K|_W
$$
由于 $W \to V$ 是到 $\mathbf{P}^n_A$ 的开子概形中的闭浸入，
这意味着 $K|_W$ 相对于 $A$ 拟凝聚；例如可由《态射进阶》引理
\ref{more-morphisms-lemma-check-relative-pseudo-coherence-on-charts} 看出。
\end{proof}

\begin{lemma}
\label{spaces-more-morphisms-lemma-characterize-pseudo-coh-improved}
设 $A$ 为环。设 $X$ 为 $A$ 上可分且有限表示的代数空间。
设 $K \in D_\QCoh(\mathcal{O}_X)$。如果
$R \Gamma (X, E \otimes ^{\mathbf{L}} K)$ 是 $D(A)$ 中的拟凝聚对象，
并且这对每个完美的 $E \in D(\mathcal{O}_X)$ 都成立，
那么 $K$ 相对于 $A$ 拟凝聚。
\end{lemma}

\begin{proof}
鉴于引理 \ref{spaces-more-morphisms-lemma-characterize-pseudo-coherent}，只需证明
$R \Gamma (X, E \otimes ^{\mathbf{L}} K)$ 是 $D(A)$ 中的拟凝聚对象，
并且这对每个拟凝聚的 $E \in D(\mathcal{O}_X)$ 都成立。
由《空间的导出范畴》命题
\ref{spaces-perfect-proposition-detecting-bounded-above}，
得到 $K \in D^-_\QCoh (\mathcal{O}_X)$。现在结论由
《空间的导出范畴》引理
\ref{spaces-perfect-lemma-perfect-enough} 得到。
\end{proof}

\section{相对完美对象}
\label{spaces-more-morphisms-section-relatively-perfect}

\noindent
本节引入 \cite{lieblich-complexes} 中的一个概念。
《概形的导出范畴》第 \ref{perfect-section-relatively-perfect} 节
已经对概形态射讨论过这个概念。

\begin{definition}
\label{spaces-more-morphisms-definition-relatively-perfect}
设 $S$ 为概形。设 $f : X \to Y$ 为 $S$ 上代数空间的态射，
并且它平坦且局部有限表示。称对象 $E$（它属于 $D(\mathcal{O}_X)$）
{\it 相对于 $Y$ 完美}，或称它 {\it $Y$-完美}，如果 $E$ 拟凝聚
（《位点上的上同调》定义
\ref{sites-cohomology-definition-pseudo-coherent}），并且 $E$
作为 $D(f^{-1}\mathcal{O}_Y)$ 的对象局部具有有限 Tor 维数
（《位点上的上同调》定义
\ref{sites-cohomology-definition-tor-amplitude}）。
\end{definition}

\noindent
有关讨论请见《概形的导出范畴》评注
\ref{perfect-remark-discuss-rel-perfect}；这里仅指出：由引理
\ref{spaces-more-morphisms-lemma-relative-pseudo-coherent-is-moot}，$E$ 拟凝聚等价于
$E$ 相对于 $Y$ 拟凝聚。此外，$E$ 的拟凝聚性蕴含
$E \in D_\QCoh(\mathcal{O}_X)$；见《空间的导出范畴》引理
\ref{spaces-perfect-lemma-pseudo-coherent}。

\begin{example}
\label{spaces-more-morphisms-example-relatively-perfect-field}
设 $k$ 为域。设 $X$ 为 $k$ 上有限表示的代数空间
（特别地，$X$ 拟紧）。则对象 $E$（它属于 $D(\mathcal{O}_X)$）
为 $k$-完美，当且仅当它有界且拟凝聚（按定义），也就是说，
当且仅当它属于 $D^b_{\textit{Coh}}(X)$（由《空间的导出范畴》引理
\ref{spaces-perfect-lemma-identify-pseudo-coherent-noetherian}）。
所以相对完美{\bf 并不}意味着“在纤维上完美”。
\end{example}

\noindent
相应的代数概念在《代数进阶》第
\ref{more-algebra-section-relatively-perfect} 节中研究。
可以如下把代数空间的概念与代数概念联系起来。

\begin{lemma}
\label{spaces-more-morphisms-lemma-affine-locally-rel-perfect}
设 $S$ 为概形。设 $f : X \to Y$ 为 $S$ 上代数空间的态射，
并且它平坦且局部有限表示。设
$E \in D_\QCoh(\mathcal{O}_X)$。下列条件等价：
\begin{enumerate}
\item $E$ 是 $Y$-完美的；
\item 对每个交换图
$$
\xymatrix{
U \ar[d] \ar[r]_g & V \ar[d] \\
X \ar[r]^f & Y
}
$$
其中 $U$、$V$ 是概形且竖直箭头平展，复形 $E|_U$
在《概形的导出范畴》定义
\ref{perfect-definition-relatively-perfect} 的意义下是 $V$-完美的；
\item 对某个如 (2) 的交换图，其中 $U \to X$ 为满态射，
复形 $E|_U$ 在《概形的导出范畴》定义
\ref{perfect-definition-relatively-perfect} 的意义下是 $V$-完美的；
\item 对每个如 (2) 的交换图，其中 $U$ 和 $V$ 仿射，
复形 $R\Gamma(U, E)$ 是 $\mathcal{O}_Y(V)$-完美的。
\end{enumerate}
\end{lemma}

\begin{proof}
为使本引理的 (2)、(3)、(4) 有意义，注意对象
$E|_U$（它属于 $D_\QCoh(\mathcal{O}_U)$）对应于对象
$E_0$（它属于 $D_\QCoh(\mathcal{O}_{U_0})$），其中 $U_0$
表示 $U$ 的底层概形；见《空间的导出范畴》引理
\ref{spaces-perfect-lemma-derived-quasi-coherent-small-etale-site}。
此外，在这种情形中，$E_0$ 拟凝聚，当且仅当 $E|_U$ 拟凝聚；
见《空间的导出范畴》引理
\ref{spaces-perfect-lemma-descend-pseudo-coherent}。
同样，由《空间的导出范畴》引理
\ref{spaces-perfect-lemma-tor-dimension-rel}，$E|_U$ 在
$f^{-1}\mathcal{O}_Y|_U = g^{-1}\mathcal{O}_V$ 上局部具有有限 Tor 维数，
当且仅当 $E_0$ 在 $g_0^{-1}\mathcal{O}_{V_0}$ 上局部具有有限 Tor 维数。
这里 $g_0 : U_0 \to V_0$ 是表示 $g : U \to V$ 的概形态射
（记号如《空间的导出范畴》评注
\ref{spaces-perfect-remark-match-total-direct-images}）。
最后，注意“拟凝聚”是平展局部性质，当然，“局部具有有限 Tor 维数”
也是平展局部性质。因此，只需平展局部地检验 $Y$-完美性；
由上面的讨论，(1) 蕴含 (2)，而 (3) 蕴含 (1)。由于
《概形的导出范畴》引理
\ref{perfect-lemma-affine-locally-rel-perfect} 说明 (4) 与 (2)、(3) 等价，
证明完毕。
\end{proof}

\begin{lemma}
\label{spaces-more-morphisms-lemma-triangulated}
设 $S$ 为概形。设 $f : X \to Y$ 为 $S$ 上代数空间的态射，
并且它平坦且局部有限表示。$D(\mathcal{O}_X)$ 中由
$Y$-完美对象组成的全子范畴是饱和的\footnote{《导出范畴》定义
\ref{derived-definition-saturated}。}三角子范畴。
\end{lemma}

\begin{proof}
这由《位点上的上同调》引理
\ref{sites-cohomology-lemma-cone-pseudo-coherent}、
\ref{sites-cohomology-lemma-summands-pseudo-coherent}、
\ref{sites-cohomology-lemma-cone-tor-amplitude} 和
\ref{sites-cohomology-lemma-summands-tor-amplitude} 得到。
\end{proof}

\begin{lemma}
\label{spaces-more-morphisms-lemma-perfect-relatively-perfect}
设 $S$ 为概形。设 $f : X \to Y$ 为 $S$ 上代数空间的态射，
并且它平坦且局部有限表示。$D(\mathcal{O}_X)$ 的完美对象是
$Y$-完美的。如果 $K, M \in D(\mathcal{O}_X)$，则
$K \otimes_{\mathcal{O}_X}^\mathbf{L} M$ 为 $Y$-完美，只要
$K$ 完美且 $M$ 为 $Y$-完美。
\end{lemma}

\begin{proof}
使用引理 \ref{spaces-more-morphisms-lemma-affine-locally-rel-perfect} 约化到概形的情形，
然后应用《概形的导出范畴》引理
\ref{perfect-lemma-perfect-relatively-perfect}。
\end{proof}

\begin{lemma}
\label{spaces-more-morphisms-lemma-base-change-relatively-perfect}
设 $S$ 为概形。设 $f : X \to Y$ 为 $S$ 上代数空间的态射，
并且它平坦且局部有限表示。设 $g : Y' \to Y$
为 $S$ 上代数空间的态射。令 $X' = Y' \times_Y X$，并记
% P09-ERR-0759
$g' : X' \to X$ 为投影。如果 $K \in D(\mathcal{O}_X)$
为 $Y$-完美的，那么 $L(g')^*K$ 为 $Y'$-完美的。
\end{lemma}

\begin{proof}
使用引理 \ref{spaces-more-morphisms-lemma-affine-locally-rel-perfect} 约化到概形的情形，
然后应用《概形的导出范畴》引理
\ref{perfect-lemma-base-change-relatively-perfect}。
\end{proof}

\begin{situation}
\label{spaces-more-morphisms-situation-relative-descent}
设 $S$ 为概形。设 $Y = \lim_{i \in I} Y_i$ 为 $S$ 上
代数空间有向系的极限，其仿射转移态射为
$g_{i'i} : Y_{i'} \to Y_i$。假设 $Y_i$ 对每个 $i \in I$
都拟紧且拟分离。记 $g_i : Y \to Y_i$ 为投影。固定元素
$0 \in I$ 和有限表示平坦态射 $X_0 \to Y_0$。
令 $X_i = Y_i \times_{Y_0} X_0$ 和 $X = Y \times_{Y_0} X_0$，
并记转移态射 $f_{i'i} : X_{i'} \to X_i$ 以及投影
$f_i : X \to X_i$。
\end{situation}

\begin{lemma}
\label{spaces-more-morphisms-lemma-relative-descend-homomorphisms}
在情形 \ref{spaces-more-morphisms-situation-relative-descent} 中，设 $K_0$ 和 $L_0$
为 $D(\mathcal{O}_{X_0})$ 的对象。令
$K_i = Lf_{i0}^*K_0$ 和 $L_i = Lf_{i0}^*L_0$（其中 $i \geq 0$），并令
$K = Lf_0^*K_0$ 和 $L = Lf_0^*L_0$。则映射
$$
\colim_{i \geq 0} \Hom_{D(\mathcal{O}_{X_i})}(K_i, L_i)
\longrightarrow
\Hom_{D(\mathcal{O}_X)}(K, L)
$$
是同构，条件是 $K_0$ 拟凝聚，并且
$L_0 \in D_\QCoh(\mathcal{O}_{X_0})$ 作为
$D((X_0 \to Y_0)^{-1}\mathcal{O}_{Y_0})$ 的对象
（局部）具有有限 Tor 维数。
% P09-ERR-0760
\end{lemma}

\begin{proof}
对每个拟紧拟分离对象 $U_0$（它属于 $(X_0)_{spaces, \etale}$），
考虑条件 $P$：
$$
\colim_{i \geq 0} \Hom_{D(\mathcal{O}_{U_i})}(K_i|_{U_i}, L_i|_{U_i})
\longrightarrow
\Hom_{D(\mathcal{O}_U)}(K|_U, L|_U)
$$
是同构，其中 $U = X \times_{X_0} U_0$ 且
$U_i = X_i \times_{X_0} U_0$。我们将证明 $P$ 对每个 $U_0$ 成立。

\medskip\noindent
假设 $(U_0 \subset W_0, V_0 \to W_0)$ 是
$(X_0)_{spaces, \etale}$ 中的初等特异方形，并且 $P$ 对
$U_0, V_0, U_0 \times_{W_0} V_0$ 成立。
则由导出范畴中 Hom 的 Mayer--Vietoris 性质，$P$ 对 $W_0$ 成立；
见《空间的导出范畴》引理
\ref{spaces-perfect-lemma-mayer-vietoris-hom}。

\medskip\noindent
首先考虑 $U_0 = W_0 \times_{Y_0} X_0$，其中 $W_0$ 是
$(Y_0)_{spaces, \etale}$ 的拟紧拟分离对象。
把《空间的导出范畴》引理
\ref{spaces-perfect-lemma-induction-principle} 的归纳原理
应用于这些 $W_0$ 以及上一段，便发现只需证明 $P$，
而只考虑 $U_0 = W_0 \times_{Y_0} X_0$ 且 $W_0$ 仿射的情形。
换言之，已经约化到 $Y_0$ 仿射的情形。接着再次应用归纳原理，
这次应用于 $X_0$ 的所有拟紧拟分离开集，从而进一步约化到
$X_0$ 也仿射的情形。

\medskip\noindent
如果 $X_0$ 和 $Y_0$ 仿射，就回到了概形情形；该情形已在
《概形的导出范畴》引理
\ref{perfect-lemma-relative-descend-homomorphisms} 中证明。
读者可以使用《空间的导出范畴》引理
\ref{spaces-perfect-lemma-pseudo-coherent}、
\ref{spaces-perfect-lemma-derived-quasi-coherent-small-etale-site}、
\ref{spaces-perfect-lemma-descend-pseudo-coherent} 和
\ref{spaces-perfect-lemma-tor-dimension-rel}，把陈述转换为仅涉及
概形及概形上模的导出范畴的陈述。
\end{proof}

\begin{lemma}
\label{spaces-more-morphisms-lemma-descend-relatively-perfect}
在情形 \ref{spaces-more-morphisms-situation-relative-descent} 中，
由 $Y$-完美对象组成的范畴（这些对象属于 $D(\mathcal{O}_X)$）
是由 $Y_i$-完美对象组成的范畴（这些对象属于
$D(\mathcal{O}_{X_i})$）的余极限。
\end{lemma}

\begin{proof}
对每个拟紧拟分离对象 $U_0$（它属于 $(X_0)_{spaces, \etale}$），
考虑条件 $P$：函子
$$
\colim_{i \geq 0} D_{Y_i\text{-完美}}(\mathcal{O}_{U_i})
\longrightarrow
D_{Y\text{-完美}}(\mathcal{O}_U)
$$
是等价，其中 $U = X \times_{X_0} U_0$ 且
$U_i = X_i \times_{X_0} U_0$。由引理
\ref{spaces-more-morphisms-lemma-relative-descend-homomorphisms}，已知这个函子全忠实。
因此只需证明本质满。

\medskip\noindent
假设 $(U_0 \subset W_0, V_0 \to W_0)$ 是
$(X_0)_{spaces, \etale}$ 中的初等特异方形，并且 $P$ 对
$U_0, V_0, U_0 \times_{W_0} V_0$ 成立。我们断言 $P$ 对 $W_0$ 成立。
将使用记号 $U_i = X_i \times_{X_0} U_0$、
$U = X \times_{X_0} U_0$，并对 $V_0$ 和 $W_0$ 使用类似记号。
我们将滥用符号 $f_i$ 来表示所有态射 $U \to U_i$、$V \to V_i$、
$U \times_W V \to U_i \times_{W_i} V_i$ 和 $W \to W_i$。
% P09-ERR-0761
设 $E$ 是 $Y$-完美对象，它属于 $D(\mathcal{O}_W)$。
目标：证明 $E$ 在该函子的本质像中。由假设，可以找到
$i \geq 0$、$Y_i$-完美对象 $E_{U, i}$（在 $U_i$ 上）、
$Y_i$-完美对象 $E_{V, i}$（在 $V_i$ 上），以及同构
$Lf_i^*E_{U, i} \to E|_U$ 和 $Lf_i^*E_{V, i} \to E|_V$。
令
% P09-ERR-0762
$$
a : E_{U, i} \to (Rf_{i, *}E)|_{U_i}
\quad\text{且}\quad
b : E_{V, i} \to (Rf_{i, *}E)|_{V_i}
$$
为与同构 $Lf_i^*E_{U, i} \to E|_U$ 和
$Lf_i^*E_{V, i} \to E|_V$ 伴随的映射。
% P09-ERR-0763
由全忠实性，增大 $i$ 后，可以找到同构
$c : E_{U, i}|_{U_i \times_{W_i} V_i} \to E_{V, i}|_{U_i \times_{W_i} V_i}$，
其拉回为如下复合识别：
$$
Lf_i^*E_{U, i}|_{U \times_W V} \to E|_{U \times_W V} \to
Lf_i^*E_{V, i}|_{U \times_W V}.
$$
应用《空间的导出范畴》引理
\ref{spaces-perfect-lemma-glue}，得到对象 $E_i$（在 $W_i$ 上）
以及映射 $d : E_i \to Rf_{i, *}E$，它限制为映射 $a$ 和 $b$，
分别在 $U_i$ 和 $V_i$ 上。于是显然，$E_i$ 是 $Y_i$-完美的
（因为相对完美是平展局部性质），而 $d$ 与同构
$Lf_i^*E_i \to E$ 伴随。

\medskip\noindent
完全照搬引理 \ref{spaces-more-morphisms-lemma-relative-descend-homomorphisms} 的证明中
使用归纳原理（《空间的导出范畴》引理
\ref{spaces-perfect-lemma-induction-principle}）的论证，约化到
$X_0$ 和 $Y_0$ 都仿射的情形：先处理
$(Y_0)_{spaces, \etale}$ 中的拟紧拟分离对象，把问题约化到
$Y_0$ 仿射；再处理 $(X_0)_{spaces, \etale}$ 中的拟紧拟分离对象，
% P09-ERR-0764
把问题约化到 $X_0$ 仿射。在仿射情形，结论由概形情形得到，
即《概形的导出范畴》引理
\ref{perfect-lemma-descend-relatively-perfect}。
到概形情形的转换由引理
\ref{spaces-more-morphisms-lemma-affine-locally-rel-perfect} 完成。
\end{proof}

\begin{lemma}
\label{spaces-more-morphisms-lemma-derived-pushforward-rel-perfect}
设 $S$ 为概形。设 $f : X \to Y$ 为 $S$ 上代数空间的态射，
并且它平坦、固有且有限表示。设 $E \in D(\mathcal{O}_X)$
为 $Y$-完美的。则 $Rf_*E$ 是 $D(\mathcal{O}_Y)$ 的完美对象，
并且它的构造与任意基变换交换。
\end{lemma}

\begin{proof}
关于基变换的陈述是《空间的导出范畴》引理
\ref{spaces-perfect-lemma-base-change-tensor}（其中
$\mathcal{G}^\bullet$ 等于 $\mathcal{O}_X$，并集中在次数 $0$）。
所以只需证明 $Rf_*E$ 是完美对象。我们将通过极限论证约化到
$Y$ 为诺特仿射的情形。

\medskip\noindent
这个问题在 $Y$ 上平展局部，故可假设 $Y$ 仿射。
记 $Y = \Spec(R)$。把 $R = \colim R_i$ 写成诺特环 $R_i$
的滤过余极限。由《空间的极限》引理
\ref{spaces-limits-lemma-descend-finite-presentation}，存在 $i$
和代数空间 $X_i$，它在 $R_i$ 上有限表示，并使它到 $R$
的基变换为 $X$。由《空间的极限》引理
\ref{spaces-limits-lemma-eventually-proper} 和
\ref{spaces-limits-lemma-descend-flat}，可以假设 $X_i$
在 $R_i$ 上固有且平坦。由引理
\ref{spaces-more-morphisms-lemma-descend-relatively-perfect}，可以假设存在
% P09-ERR-0765
$R_i$-完美对象 $E_i$（它属于 $D(\mathcal{O}_{X_i})$），
其到 $X$ 的拉回为 $E$。把《空间的导出范畴》引理
\ref{spaces-perfect-lemma-perfect-direct-image} 应用于
$X_i \to \Spec(R_i)$ 和 $E_i$，再使用已经证明的基变换性质，
便得到结论。
\end{proof}

\begin{lemma}
\label{spaces-more-morphisms-lemma-compute-ext-rel-perfect}
设 $S$ 为概形。设 $f : X \to Y$ 为 $S$ 上代数空间的态射。
设 $E, K \in D(\mathcal{O}_X)$。假设
\begin{enumerate}
\item $Y$ 拟紧且拟分离；
\item $f$ 固有、平坦且有限表示；
\item $E$ 是 $Y$-完美的；
\item $K$ 拟凝聚。
\end{enumerate}
则存在拟凝聚的 $L \in D(\mathcal{O}_Y)$，使得
$$
Rf_*R\SheafHom(K, E) = R\SheafHom(L, \mathcal{O}_Y)
$$
而且任意基变换后同样成立：给定
$$
\vcenter{
\xymatrix{
X' \ar[r]_{g'} \ar[d]_{f'} &
X \ar[d]^f \\
Y' \ar[r]^g &
Y
}
}
\quad\quad
\begin{matrix}
\text{若该图为笛卡尔图，则有} \\
Rf'_*R\SheafHom(L(g')^*K, L(g')^*E) \\
= R\SheafHom(Lg^*L, \mathcal{O}_{Y'})
\end{matrix}
$$
\end{lemma}

\begin{proof}
由于 $Y$ 拟紧且拟分离，$X$ 也如此。
由《空间的导出范畴》引理
\ref{spaces-perfect-lemma-pseudo-coherent-hocolim}，可以写成
$K = \text{hocolim} K_n$，其中 $K_n$ 完美，并且 $K_n \to K$
在截断 $\tau_{\geq -n}$ 上诱导同构。设 $K_n^\vee$ 为对偶完美复形
（《位点上的上同调》引理
\ref{sites-cohomology-lemma-dual-perfect-complex}）。得到完美对象的逆系
$\ldots \to K_3^\vee \to K_2^\vee \to K_1^\vee$。
由引理 \ref{spaces-more-morphisms-lemma-perfect-relatively-perfect}，
$K_n^\vee \otimes_{\mathcal{O}_X} E$ 是 $Y$-完美的。
因此，可以把引理 \ref{spaces-more-morphisms-lemma-derived-pushforward-rel-perfect}
应用于 $K_n^\vee \otimes_{\mathcal{O}_X} E$，得到逆系
$$
\ldots \to M_3 \to M_2 \to M_1
$$
它由 $Y$ 上的完美复形组成，其中
$$
M_n = Rf_*(K_n^\vee \otimes_{\mathcal{O}_X}^\mathbf{L} E) =
Rf_*R\SheafHom(K_n, E)
$$
此外，这些复形的构造与任意基变换交换；具体地，
$Lg^*M_n =
Rf'_*((L(g')^*K_n)^\vee \otimes_{\mathcal{O}_{X'}}^\mathbf{L} L(g')^*E) =
Rf'_*R\SheafHom(L(g')^*K_n, L(g')^*E)$。

\medskip\noindent
由于 $K_n \to K$ 在 $\tau_{\geq -n}$ 上诱导同构，可知
$K_n \to K_{n + 1}$ 在 $\tau_{\geq -n}$ 上诱导同构。
因此，$K_{n + 1}^\vee \to K_n^\vee$ 在 $\tau_{\leq n}$ 上诱导同构，
因为 $K_n^\vee = R\SheafHom(K_n, \mathcal{O}_X)$。
假设 $E$ 的 Tor 振幅在 $[a, b]$ 中，这里把它看作
$f^{-1}\mathcal{O}_Y$-模复形。则任意基变换后也是如此；见《空间的导出范畴》引理
\ref{spaces-perfect-lemma-tor-independence-and-tor-amplitude}。
可以看出映射
$K_{n + 1}^\vee \otimes_{\mathcal{O}_X} E \to
K_n^\vee \otimes_{\mathcal{O}_X} E$
在 $\tau_{\leq n + a}$ 上诱导同构，任意基变换后也同样成立。
应用右导出函子 $Rf_*$，得到映射 $M_{n + 1} \to M_n$
在 $\tau_{\leq n + a}$ 上诱导同构，任意基变换后也同样成立。
取特异三角
$$
M_{n + 1} \to M_n \to C_n \to M_{n + 1}[1]
$$
取 $y \in |Y|$。选择初等平展邻域
$(U, u) \to (Y, y)$；这是可行的，见《良好空间》引理
\ref{decent-spaces-lemma-decent-space-elementary-etale-neighbourhood}。
取 $Y'$ 为 $u$ 处剩余域的谱。拉回后可见
$C_n|_U \otimes_{\mathcal{O}_U}^\mathbf{L} \kappa(u)$
仅在次数 $\geq n + a$ 具有非零上同调。由《代数进阶》引理
\ref{more-algebra-lemma-lift-perfect-from-residue-field}，
可见完美复形 $C_n|_U$ 的 Tor 振幅在 $[n + a, m_n]$ 中，
其中 $m_n$ 为某个整数；必要时再缩小 $U$。
所以 $C_n$ 的 Tor 振幅在 $[n + a, m_n]$ 中，其中 $m_n$
为某个整数（因为 $Y$ 拟紧）。特别地，对偶完美复形
$C_n^\vee$ 的 Tor 振幅在 $[-m_n, -n - a]$ 中。

\medskip\noindent
令 $L_n = M_n^\vee$ 为对偶完美复形。上一段讨论的结论是，
$L_n \to L_{n + 1}$ 在 $\tau_{\geq -n - a}$ 上诱导同构。
所以 $L = \text{hocolim} L_n$ 拟凝聚；见《空间的导出范畴》引理
\ref{spaces-perfect-lemma-pseudo-coherent-hocolim}。由于
$$
R\SheafHom(K, E) = R\SheafHom(\text{hocolim} K_n, E) =
R\lim R\SheafHom(K_n, E) = R\lim K_n^\vee \otimes_{\mathcal{O}_X} E
$$
（《位点上的上同调》引理
\ref{sites-cohomology-lemma-colim-and-lim-of-duals}），并且
$R\lim$ 与 $Rf_*$ 交换，得到
$$
Rf_*R\SheafHom(K, E) = R\lim M_n = R\lim R\SheafHom(L_n, \mathcal{O}_Y) =
R\SheafHom(L, \mathcal{O}_Y)
$$
这证明了 $Y$ 上的公式。由于 $M_n$ 的构造与
% P09-ERR-0766
基变换相容，任意基变换后公式仍成立。
\end{proof}

\begin{remark}
\label{spaces-more-morphisms-remark-compare-L}
读者可能已经注意到引理 \ref{spaces-more-morphisms-lemma-compute-ext-rel-perfect}
与《空间的导出范畴》引理
\ref{spaces-perfect-lemma-compute-ext} 之间的相似性。
事实上，引理 \ref{spaces-more-morphisms-lemma-compute-ext-rel-perfect} 的拟凝聚复形 $L$
可刻画为 $Y$ 上唯一的拟凝聚复形，使得存在函子性同构
$$
\Ext^i_{\mathcal{O}_Y}(L, \mathcal{F}) \longrightarrow
\Ext^i_{\mathcal{O}_X}(K,
E \otimes_{\mathcal{O}_X}^\mathbf{L} Lf^*\mathcal{F})
$$
当 $\mathcal{F}$ 遍历 $\QCoh(\mathcal{O}_Y)$ 时，这些同构与边界映射相容。
如果以后需要这一点，我们将在此陈述精确结果并给出详细证明。
\end{remark}

\begin{lemma}
\label{spaces-more-morphisms-lemma-bounded-on-fibres}
设 $S$ 为概形。设 $X$ 为 $S$ 上代数空间，使得结构态射
$f : X \to S$ 平坦且局部有限表示。设 $E$ 为
$D(\mathcal{O}_X)$ 的拟凝聚对象。下列条件等价：
\begin{enumerate}
\item $E$ 是 $S$-完美的；
\item $E$ 局部下有界，并且对每个点 $s \in S$，
对象 $L(X_s \to X)^*E$（它属于 $D(\mathcal{O}_{X_s})$）局部下有界。
\end{enumerate}
\end{lemma}

\begin{proof}
由于一切都是局部的，立即约化到 $X$ 和 $S$ 仿射的情形；
见引理 \ref{spaces-more-morphisms-lemma-affine-locally-rel-perfect}。这一情形由
《概形的导出范畴》引理 \ref{perfect-lemma-bounded-on-fibres} 处理。
\end{proof}

\begin{lemma}
\label{spaces-more-morphisms-lemma-characterize-relatively-perfect}
设 $A$ 为环。设 $X$ 为在 $A$ 上可分、有限表示且平坦的代数空间。
设 $K \in D_\QCoh(\mathcal{O}_X)$。如果
$R \Gamma (X, E \otimes^\mathbf{L} K)$ 是 $D(A)$ 中的完美对象，
并且这对每个完美的 $E \in D(\mathcal{O}_X)$ 都成立，
那么 $K$ 是 $\Spec(A)$-完美的。
\end{lemma}

\begin{proof}
由引理 \ref{spaces-more-morphisms-lemma-characterize-pseudo-coh-improved}，$K$
相对于 $A$ 拟凝聚。由引理
\ref{spaces-more-morphisms-lemma-relative-pseudo-coherent-is-moot}，$K$
在 $D(\mathcal{O}_X)$ 中拟凝聚。由《空间的导出范畴》命题
\ref{spaces-perfect-proposition-detecting-bounded-below}，
可见 $K$ 属于 $D^-(\mathcal{O}_X)$。设 $\mathfrak{p}$ 为 $A$ 的素理想，并记
% P09-ERR-0767
$i : Y \to X$ 为 $\mathfrak{p}$ 上概形论纤维的包含；也就是说，
% P09-ERR-0768
$Y$ 是 $\kappa(\mathfrak p)$ 上的概形。由引理
\ref{spaces-more-morphisms-lemma-bounded-on-fibres}，只要能证明 $Li^*(K)$ 下有界，
结论就成立。设 $G \in D_{perf} (\mathcal{O}_X)$ 为生成
$D_\QCoh (\mathcal{O}_X)$ 的完美复形；见《空间的导出范畴》定理
\ref{spaces-perfect-theorem-bondal-van-den-Bergh}。有
\begin{align*}
R\Hom _{\mathcal{O}_Y}(Li^*(G), Li^*(K))
& =
R\Gamma(Y, Li^*(G ^\vee \otimes ^\mathbf{L} K)) \\
& =
R\Gamma(X, G^\vee \otimes ^{\mathbf{L}} K)
\otimes^\mathbf{L}_A \kappa(\mathfrak{p})
\end{align*}
第一个等式使用了 $Li^*$ 保持完美对象和对偶这一事实以及
《位点上的上同调》引理
\ref{sites-cohomology-lemma-dual-perfect-complex}；略去一些细节。
第二个等式由《空间的导出范畴》引理
\ref{spaces-perfect-lemma-compare-base-change} 得到，因为 $X$
在 $A$ 上平坦。由假设，这给出 $D(\kappa(\mathfrak{p}))$
中的完美对象。对象 $Li^*(G) \in D_{perf}(\mathcal{O}_Y)$
由《空间的导出范畴》评注
\ref{spaces-perfect-remark-pullback-generator} 生成
$D_\QCoh(\mathcal{O}_Y)$。于是《空间的导出范畴》命题
\ref{spaces-perfect-proposition-detecting-bounded-below}
蕴含 $Li^*(K)$ 下有界，从而结论成立。
\end{proof}

\section{立方定理}
\label{spaces-more-morphisms-section-theorem-cube}

\noindent
本节是《态射进阶》第
\ref{more-morphisms-section-theorem-cube} 节的对应版本。
下面的引理说明，在引理所作的假设下，Picard 函子的对角线
可由局部闭浸入表示。

\begin{lemma}
\label{spaces-more-morphisms-lemma-diagonal-picard-flat-proper}
设 $S$ 为概形。设 $f : X \to Y$ 为 $S$ 上代数空间的
平坦有限表示固有态射。设 $\mathcal{E}$ 为有限局部自由
$\mathcal{O}_X$-模。对态射 $g : Y' \to Y$，考虑基变换图
$$
\xymatrix{
X' \ar[d]_{f'} \ar[r]_{g'} & X \ar[d]^f \\
Y' \ar[r]^g & Y
}
$$
假设 $\mathcal{O}_{Y'} \to f'_*\mathcal{O}_{X'}$
对所有 $g : Y' \to Y$ 都是同构。
则存在有限表示浸入 $j : Z \to Y$，使得态射
$g : Y' \to Y$ 经由 $Z$ 分解，当且仅当存在有限局部自由
$\mathcal{O}_{Y'}$-模 $\mathcal{N}$，满足
% P09-ERR-0769
$(f')^*\mathcal{N} \cong (g')^*\mathcal{L}$。
\end{lemma}

\begin{proof}
设 $y : \Spec(k) \to Y$ 为域值点。纤维 $X_y$
（即 $f$ 在 $y$ 处的纤维）连通，这是因为我们的假设
$H^0(X_y, \mathcal{O}_{X_y}) = k$。所以 $\mathcal{E}$ 的秩在各纤维上恒定。
由于 $f$ 为开映射（《空间的态射》引理
\ref{spaces-morphisms-lemma-fppf-open}）且为闭映射，得到分解
$Y = \coprod Y_r$，把 $Y$ 分解为开闭子空间，使得
$\mathcal{E}$ 恒秩 $r$，这是在 $Y_r$ 的逆像上而言的。
所以可以假设 $\mathcal{E}$ 恒秩 $r$。记
% P09-ERR-0770
$\mathcal{E}^\vee = \SheafHom(\mathcal{E}, \mathcal{O}_X)$
为秩 $r$ 的对偶模。

\medskip\noindent
由上同调与基变换（更精确地，由《空间的导出范畴》引理
\ref{spaces-perfect-lemma-flat-proper-perfect-direct-image-general}），
可见 $E = Rf_*\mathcal{E}$ 是 $Y$ 的导出范畴中的完美对象，
并且它的构造与任意基变换交换。对
$E' = Rf_*\mathcal{E}^\vee$ 也同样如此。
由于在次数 $< 0$ 从来没有上同调，$E$ 和 $E'$ 的 Tor 振幅
（局部）包含在 $[0, b]$ 中，其中 $b$ 是某个数。
注意，对任意 $g : Y' \to Y$，有
$f'_*((g')^*\mathcal{E}) = H^0(Lg^*E)$ 和
$f'_*((g')^*\mathcal{E}^\vee) = H^0(Lg^*E')$。
设 $j : Z \to Y$ 和 $j' : Z' \to Y$ 为《空间的导出范畴》引理
\ref{spaces-perfect-lemma-locally-closed-where-H0-locally-free}
对 $E$ 和 $E'$（取 $a = 0$ 和 $r = r$）构造的局部闭浸入；
它们由如下性质刻画：
% P09-ERR-0771
$H^0(Lj^*E)$ 和 $H^0((j')^*E')$ 是秩 $r$ 的局部自由模，
并且与拉回相容。

\medskip\noindent
设 $g : Y' \to Y$ 为态射。如果存在引理中所述的
$\mathcal{N}$，那么使用投影公式《位点上的上同调》引理
\ref{sites-cohomology-lemma-projection-formula}，可见模
$$
f'_*((g')^*\mathcal{E}) \cong
f'_*((f')^*\mathcal{N}) \cong
\mathcal{N} \otimes_{\mathcal{O}_{Y'}} f'_*\mathcal{O}_{X'} \cong
\mathcal{N}\quad\text{且类似地}\quad
f'_*((g')^*\mathcal{E}^\vee) \cong \mathcal{N}^\vee
$$
局部自由且秩为 $r$，并且仍局部自由且秩为 $r$，
即使经过任意进一步的基变换 $Y'' \to Y'$ 也是如此。
所以在这种情形，$g : Y' \to Y$ 经由 $j$ 和 $j'$ 分解。
因此，可以替换 $Y$，所用空间为 $Z \times_Y Z'$，并假设
$f_*\mathcal{E}$ 和 $f_*\mathcal{E}^\vee$ 是局部自由
$\mathcal{O}_Y$-模且秩为 $r$，其构造与任意基变换交换。

\medskip\noindent
在这种情形下，如果 $g : Y' \to Y$ 为态射，并且存在引理中所述的
$\mathcal{N}$，那么映射（次数 $0$ 的杯积）
$$
f'_*((g')^*\mathcal{E})
\otimes_{\mathcal{O}_{Y'}}
f'_*((g')^*\mathcal{E}^\vee)
\longrightarrow \mathcal{O}_{Y'}
$$
是完美配对。反过来，如果这个杯积映射是完美配对，
那么在 $Y'$ 上局部地可找到截面的基
% P09-ERR-0769
$\sigma_1, \ldots, \sigma_r$（属于 $f'_*((g')^*\mathcal{L})$）以及
$\tau_1, \ldots, \tau_r$（属于
$f'_*((g')^*\mathcal{E}^\vee)$），其乘积满足
$\sigma_i \tau_j = \delta_{ij}$。把 $\sigma_i$ 看作
$(g')^*\mathcal{L}$ 的截面（在 $X'$ 上），并把 $\tau_j$ 看作
$(g')^*\mathcal{L}^\vee$ 的截面（在 $X'$ 上），得到
$$
\sigma_1, \ldots, \sigma_r :
\mathcal{O}_{X'}^{\oplus r}
\longrightarrow
(g')^*\mathcal{E}
$$
是同构，其逆由
$$
\tau_1, \ldots, \tau_r :
(g')^*\mathcal{E}
\longrightarrow
\mathcal{O}_{X'}^{\oplus r}
$$
给出。换言之，
$(f')^*f'_*(g')^*\mathcal{E} \cong (g')^*\mathcal{E}$。
而杯积非退化这一条件切出一个逆紧开子概形
（即某个适当行列式非零的轨迹），证明完毕。
\end{proof}

\section{复形有限性性质的下降}
\label{spaces-more-morphisms-section-descent-finiteness}

\noindent
本节是《态射进阶》第
\ref{more-morphisms-section-descent-finiteness} 节和
《概形的导出范畴》第
\ref{perfect-section-descent-finiteness} 节的对应版本。

\begin{lemma}
\label{spaces-more-morphisms-lemma-pseudo-coherent-descends-fpqc}
设 $S$ 为概形。设 $\{f_i : X_i \to X\}$ 为 $S$ 上代数空间的
fpqc 覆盖。设 $E \in D_\QCoh(\mathcal{O}_X)$，并设
$m \in \mathbf{Z}$。则 $E$ 是 $m$-拟凝聚的，当且仅当每个
$Lf_i^*E$ 都是 $m$-拟凝聚的。
\end{lemma}

\begin{proof}
拉回总保持 $m$-拟凝聚性，见《位点上的上同调》引理
\ref{sites-cohomology-lemma-pseudo-coherent-pullback}。
所以只需假设 $Lf_i^*E$ 是 $m$-拟凝聚的，并证明 $E$
是 $m$-拟凝聚的。首先，可以假设 $X_i$ 对所有 $i$ 都是概形，见
《空间上的拓扑》引理
\ref{spaces-topologies-lemma-refine-fpqc-schemes}。
% P09-ERR-0772
接着，选取满 étale 态射 $U \to X$，其中 $U$ 是概形。
则 $U_i = U \times_X X_i$ 是概形，并得到概形的 fpqc 覆盖
$\{U_i \to U\}$，见《空间上的拓扑》引理
\ref{spaces-topologies-lemma-recognize-fpqc-covering}。
由概形情形可知结论对
$\{U_i \to U\}_{i \in I}$ 成立，见《概形的导出范畴》引理
\ref{perfect-lemma-pseudo-coherent-descends-fpqc}。
另一方面，限制 $E|_U$ 来自
$D_\QCoh(\mathcal{O}_U)$ 中的一个对象（这里使用 Zariski
拓扑和 $U$ 的“通常”结构层来定义），见《空间的导出范畴》引理
\ref{spaces-perfect-lemma-derived-quasi-coherent-small-etale-site}。
由于两种拟凝聚概念（étale 与 Zariski）由《空间的导出范畴》引理
\ref{spaces-perfect-lemma-descend-pseudo-coherent}
可知相同，故引理得证。
\end{proof}

\begin{lemma}
\label{spaces-more-morphisms-lemma-tor-amplitude-descends-fppf}
设 $S$ 为概形。设 $\{g_i : Y_i \to Y\}$ 为 $S$ 上代数空间的
fpqc 覆盖。设 $f : X \to Y$ 为代数空间的态射，并置
$X_i = Y_i \times_Y X$，投影记为 $f_i : X_i \to Y_i$
和 $g'_i : X_i \to X$。设 $E \in D_\QCoh(\mathcal{O}_X)$，
并设 $a, b \in \mathbf{Z}$。则下列条件等价：
\begin{enumerate}
\item $E$ 具有包含在 $[a, b]$ 中的 Tor 振幅，这是作为
$D(f^{-1}\mathcal{O}_Y)$ 的对象而言的；以及
% P09-ERR-0773
\item $L(g'_i)^*E$ 具有包含在 $[a, b]$ 中的 Tor 振幅，
这是作为 $D(f_i^{-1}\mathcal{O}_{Y_i})$ 的对象而言的，
并且对所有 $i$ 都如此。
\end{enumerate}
把“具有包含在 $[a, b]$ 中的 Tor 振幅”替换为
“局部具有有限 Tor 维数”，结论也成立。
\end{lemma}

\begin{proof}
% P09-ERR-0774
由《空间的导出范畴》引理
\ref{spaces-perfect-lemma-tor-independence-and-tor-amplitude}
可知，拉回保持“具有包含在 $[a, b]$ 中的 Tor 振幅”。
注意，$Y_i$ 和 $X$ 在 $Y$ 上 Tor 无关，因为
$Y_i \to Y$ 平坦。假设 (2)，来证明 (1)。Tor 维数可在茎处计算，见
《位点上的上同调》引理
\ref{sites-cohomology-lemma-tor-amplitude-stalk}
和《空间的性质》定理
\ref{spaces-properties-theorem-exactness-stalks}。
设 $\overline{x}$ 为 $X$ 的几何点。选取 $i$ 和几何点
$\overline{x}_i$（它位于 $X_i$ 中），使其像为
$\overline{x}$（在 $X$ 中）。则
% P09-ERR-0775
$$
(L(g_i')^*E)_{\overline{x}_i} =
E_{\overline{x}}
\otimes_{\mathcal{O}_{X, \overline{x}}}^\mathbf{L}
\mathcal{O}_{X_, \overline{x}_i}
$$
设 $\overline{y}_i$（在 $Y_i$ 中）和 $\overline{y}$（在 $Y$ 中）
分别为 $\overline{x}_i$ 和 $\overline{x}$ 的像。
由于 $X$ 和 $Y_i$ 在 $Y$ 上 Tor 无关，可应用《代数进阶》引理
\ref{more-algebra-lemma-base-change-comparison}，从而上式右端等于
$E_{\overline{x}}
\otimes_{\mathcal{O}_{Y, \overline{y}}}^\mathbf{L}
\mathcal{O}_{Y_i, \overline{y}_i}$，这是在
$D(\mathcal{O}_{Y_i, \overline{y}_i})$ 中而言的。
% P09-ERR-0776
由于已假设它的 Tor 振幅包含在 $[a, b]$ 中，由《代数进阶》引理
\ref{more-algebra-lemma-flat-descent-tor-amplitude}
可知 $E_{\overline{x}}$ 在
$D(\mathcal{O}_{Y, \overline{y}})$ 中的 Tor 振幅包含在
$[a, b]$ 中。因此 (1) 成立。

\medskip\noindent
再使用一些初等拓扑，“局部有限 Tor 维数”的情形也随之成立。
\end{proof}

\noindent
下面的引理其实并不属于本节。

\begin{lemma}
\label{spaces-more-morphisms-lemma-thickening-pseudo-coherent}
设 $S$ 为概形。设 $i : X \to X'$ 为代数空间的有限阶加厚。
设 $K' \in D(\mathcal{O}_{X'})$ 为对象，并且
$K = Li^*K'$ 是拟凝聚的。则 $K'$ 是拟凝聚的。
\end{lemma}

\begin{proof}
先证明 $K'$ 具有拟凝聚上同调层；建议读者跳过这一部分。
为此，由第 \ref{spaces-more-morphisms-section-thickenings} 节可约化到一阶加厚的情形。
设 $\mathcal{I} \subset \mathcal{O}_{X'}$ 为切出 $X$
的拟凝聚理想层。把短正合列
$$
0 \to \mathcal{I} \to \mathcal{O}_{X'} \to i_*\mathcal{O}_X \to 0
$$
与 $K'$ 作张量，得到特异三角
$$
K' \otimes_{\mathcal{O}_{X'}}^\mathbf{L} \mathcal{I}
\to K' \to
K' \otimes_{\mathcal{O}_{X'}}^\mathbf{L} i_*\mathcal{O}_X
\to
(K' \otimes_{\mathcal{O}_{X'}}^\mathbf{L} \mathcal{I})[1]
$$
由于 $i_* = Ri_*$，并且可以把 $\mathcal{I}$ 看作拟凝聚
$\mathcal{O}_X$-模（因为这里是一阶加厚），可将此改写为
$$
i_*(K \otimes_{\mathcal{O}_X}^\mathbf{L} \mathcal{I})
\to K' \to
i_*K \to
i_*(K \otimes_{\mathcal{O}_X}^\mathbf{L} \mathcal{I})[1]
$$
这些项的识别请使用《空间的上同调》引理
\ref{spaces-cohomology-lemma-projection-formula-finite}。
由于 $K$ 属于 $D_\QCoh(\mathcal{O}_X)$，可知
$K'$ 属于 $D_\QCoh(\mathcal{O}_{X'})$；这里使用《空间的导出范畴》引理
\ref{spaces-perfect-lemma-pseudo-coherent}、
\ref{spaces-perfect-lemma-quasi-coherence-tensor-product} 和
\ref{spaces-perfect-lemma-quasi-coherence-direct-image}。

\medskip\noindent
假设 $K'$ 属于 $D_\QCoh(\mathcal{O}_{X'})$。这个问题在 $X'$
上是 étale 局部的，所以可以假设 $X'$ 仿射。
在这种情形，结论由概形情形得出（《态射进阶》引理
\ref{more-morphisms-lemma-thickening-pseudo-coherent}）。
将其转写为概形语言时，使用《空间的导出范畴》引理
\ref{spaces-perfect-lemma-derived-quasi-coherent-small-etale-site} 和
\ref{spaces-perfect-lemma-descend-pseudo-coherent}，以及注记
\ref{spaces-perfect-remark-match-total-direct-images}。
\end{proof}

\begin{lemma}
\label{spaces-more-morphisms-lemma-thickening-relatively-perfect}
设 $S$ 为概形。考虑笛卡儿图
$$
\xymatrix{
X \ar[r]_i \ar[d]_f & X' \ar[d]^{f'} \\
Y \ar[r]^j & Y'
}
$$
这是 $S$ 上代数空间的图。假设 $X' \to Y'$ 平坦且局部有限表示，而 $Y \to Y'$
为有限阶加厚。设 $E' \in D(\mathcal{O}_{X'})$。如果
$E = Li^*(E')$ 是 $Y$-完美的，则 $E'$ 是 $Y'$-完美的。
\end{lemma}

\begin{proof}
回忆一下，$Y$-完美性对于 $E$ 意味着 $E$ 拟凝聚，并且作为
$f^{-1}\mathcal{O}_Y$-模的复形在局部具有有限 Tor 维数
（定义 \ref{spaces-more-morphisms-definition-relatively-perfect}）。
由引理 \ref{spaces-more-morphisms-lemma-thickening-pseudo-coherent}，$E'$ 拟凝聚。
特别地，$E'$ 属于 $D_\QCoh(\mathcal{O}_{X'})$，见
《空间的导出范畴》引理
\ref{spaces-perfect-lemma-pseudo-coherent}。
由引理 \ref{spaces-more-morphisms-lemma-affine-locally-rel-perfect}，这把问题约化到概形情形。
概形情形是《态射进阶》引理
\ref{more-morphisms-lemma-thickening-relatively-perfect}。
\end{proof}

\begin{lemma}
\label{spaces-more-morphisms-lemma-henselian-relatively-perfect}
设 $(R, I)$ 为由环和包含在 Jacobson 根中的理想 $I$ 所组成的偶。
置 $S = \Spec(R)$ 和 $S_0 = \Spec(R/I)$。设 $X$ 为 $R$
上的代数空间，其结构态射 $f : X \to S$ 固有、平坦且有限表示。
% P09-ERR-0777
记 $X_0 = S_0 \times_S X$。设
$E \in D(\mathcal{O}_X)$ 拟凝聚。如果导出限制 $E_0$
（即 $E$ 到 $X_0$ 的导出限制）是 $S_0$-完美的，
则 $E$ 是 $S$-完美的。
\end{lemma}

\begin{proof}
选取满 étale 态射 $U \to X$，其中 $U$ 仿射。
选取闭浸入 $U \to \mathbf{A}^d_S$。置
$U_0 = S_0 \times_S U$。复形 $E_0|_{U_0}$
的 Tor 振幅包含在某个区间 $[a, b]$ 中，其中
$a, b \in \mathbf{Z}$。设 $\overline{x}$ 为 $X$ 的几何点。
我们将证明 $E_{\overline{x}}$ 在 $R$ 上的 Tor 振幅包含在
$[a - d, b]$ 中。这样即可完成证明，因为由《位点上的上同调》引理
\ref{sites-cohomology-lemma-tor-amplitude-stalk}
和《空间的性质》定理
\ref{spaces-properties-theorem-exactness-stalks}，
Tor 振幅可从茎读出。

\medskip\noindent
设 $x \in |X|$ 为由 $\overline{x}$ 确定的点。
回忆 $|X| \to |S|$ 为闭映射（根据固有态射的定义）。
由于 $I$ 包含在 Jacobson 根中，$S$ 的每个非空闭子集
都含有闭子概形 $S_0$ 的一个点。因此，可找到特化
$x \leadsto x_0$（在 $|X|$ 中），其中 $x_0 \in |X_0|$。
选取 $u_0 \in U_0$，它映到 $x_0$。由《适度空间》引理
\ref{decent-spaces-lemma-generalizations-lift-flat}
（或由直接适用于 étale 态射的《适度空间》引理
\ref{decent-spaces-lemma-specialization}），可找到特化
$u \leadsto u_0$（在 $U$ 中），使得 $u$ 映到 $x$。可以把
$\overline{x}$ 提升为几何点 $\overline{u}$，它是 $U$ 的点且位于
$u$ 上方。
于是 $E_{\overline{x}} = (E|_U)_{\overline{u}}$。

\medskip\noindent
% P09-ERR-0778
写成 $U = \Spec(A)$。则 $A$ 是平坦有限表示的 $R$-代数，
并且是多项式 $R$-代数的商，此多项式代数含 $d$ 个变量。
限制 $E|_U$（由《空间的导出范畴》引理
\ref{spaces-perfect-lemma-pseudo-coherent}、
\ref{spaces-perfect-lemma-derived-quasi-coherent-small-etale-site} 和
\ref{spaces-perfect-lemma-descend-pseudo-coherent}，
以及《概形的导出范畴》引理
\ref{perfect-lemma-affine-compare-bounded} 和
\ref{perfect-lemma-pseudo-coherent-affine}）对应于拟凝聚对象 $K$，
它属于 $D(A)$。
注意，$E_0$ 对应于 $K \otimes_A^\mathbf{L} A/IA$。
设 $\mathfrak q \subset \mathfrak q_0 \subset A$
为对应于 $u \leadsto u_0$ 的素理想。则
$$
E_{\overline{x}} =
(E|_U)_{\overline{u}} =
E_u \otimes_{\mathcal{O}_{U, u}}^\mathbf{L} \mathcal{O}_{U, \overline{u}} =
K_{\mathfrak q} \otimes_{A_\mathfrak q}^\mathbf{L} A_{\mathfrak q}^{sh}
$$
（略去了一些细节）。由于
$A_\mathfrak q \to A_\mathfrak q^{sh}$ 平坦，它作为
$R$-模的 Tor 振幅等于 $K_\mathfrak q$ 作为 $R$-模的
Tor 振幅（《代数进阶》引理
\ref{more-algebra-lemma-no-change-tor-amplitude}）。
此外，$K_{\mathfrak q}$ 是 $K_{\mathfrak q_0}$ 的局部化。
所以只需证明 $K_{\mathfrak q_0}$ 具有包含在
$[a - d, b]$ 中的 Tor 振幅，这是作为 $R$-模的复形而言的。

\medskip\noindent
设 $I \subset \mathfrak p_0 \subset R$ 为对应于
$f(x_0)$ 的素理想。则有
\begin{align*}
K \otimes_R^\mathbf{L} \kappa(\mathfrak p_0)
& =
(K \otimes_R^\mathbf{L} R/I) \otimes_{R/I}^\mathbf{L}
\kappa(\mathfrak p_0) \\
& =
(K \otimes_A^\mathbf{L} A/IA) \otimes_{R/I}^\mathbf{L} \kappa(\mathfrak p_0)
\end{align*}
% P09-ERR-0779
第二个等式成立是因为 $R \to A$ 平坦。
根据对 $a,b$ 的选择，这个复形仅在区间 $[a,b]$
内的次数上具有上同调。最后，将《代数进阶》引理
\ref{more-algebra-lemma-lift-from-fibre-relatively-perfect}
应用于 $R \to A$、$\mathfrak q_0$、$\mathfrak p_0$ 和 $K$，
即可得到结论。
\end{proof}

\section{结点曲线的范族}
\label{spaces-more-morphisms-section-families-nodal}

\noindent
本节接续《代数曲线》第
\ref{curves-section-families-nodal} 节。
关于我们的术语选择，也请参阅该节。

\medskip\noindent
概形态射的“相对维数为 $1$ 且至多带结点奇性”这一性质，
在源和目标上都是 étale 局部的，见《下降》引理
\ref{descent-lemma-etale-local-source-target}，以及
《代数曲线》引理
\ref{curves-lemma-nodal-family-etale-local-source}、
\ref{curves-lemma-nodal-family-fpqc-local-target} 和
\ref{curves-lemma-nodal-family-postcompose-etale}。
它还在基变换下稳定，并且在目标上是 fpqc 局部的，见
《代数曲线》引理 \ref{curves-lemma-base-change-nodal-family} 和
\ref{curves-lemma-nodal-family-fpqc-local-target}。因此，根据
《空间的态射》引理 \ref{spaces-morphisms-lemma-local-source-target}，
可如下为代数空间定义相对维数为 $1$ 且至多带结点奇性的态射；
当该态射可表时，这与《空间的态射》第
\ref{spaces-morphisms-section-representable} 节中已经定义的概念一致。

\begin{definition}
\label{spaces-more-morphisms-definition-nodal-family}
设 $S$ 是概形。
设 $f : X \to Y$ 是 $S$ 上代数空间的态射。
称 $f$ 是{\it 相对维数为 $1$ 且至多带结点奇性的}，
如果《空间的态射》引理
\ref{spaces-morphisms-lemma-local-source-target} 的等价条件对
$\mathcal{P} =$“相对维数为 $1$ 且至多带结点奇性”成立。
\end{definition}

\begin{lemma}
\label{spaces-more-morphisms-lemma-base-change-nodal}
相对维数为 $1$ 且至多带结点奇性这一性质在基变换下保持。
\end{lemma}

\begin{proof}
见《空间的态射》评注 \ref{spaces-morphisms-remark-base-change-P}
和《代数曲线》引理 \ref{curves-lemma-base-change-nodal-family}。
\end{proof}

\begin{lemma}
\label{spaces-more-morphisms-lemma-nodal-local}
设 $S$ 是概形。
设 $f : X \to Y$ 是 $S$ 上代数空间的态射。
下列条件等价：
\begin{enumerate}
\item $f$ 相对维数为 $1$ 且至多带结点奇性；
\item 对每个概形 $Z$ 和任意态射 $Z \to Y$，态射
$Z \times_Y X \to Z$ 都相对维数为 $1$ 且至多带结点奇性；
\item 对每个仿射概形 $Z$ 和任意态射
$Z \to Y$，态射 $Z \times_Y X \to Z$ 都
相对维数为 $1$ 且至多带结点奇性；
\item 存在概形 $V$ 和满 étale 态射
$V \to Y$，使得 $V \times_Y X \to V$
相对维数为 $1$ 且至多带结点奇性；
\item 存在概形 $U$ 和满 étale 态射
$\varphi : U \to X$，使得复合态射 $f \circ \varphi$
相对维数为 $1$ 且至多带结点奇性；
\item 对每个交换图
$$
\xymatrix{
U \ar[d] \ar[r] & V \ar[d] \\
X \ar[r] & Y
}
$$
% P09-ERR-0780
其中 $U$、$V$ 是概形，竖直箭头均为 étale 态射，则上方的
水平箭头相对维数为 $1$ 且至多带结点奇性；
\item 存在交换图
$$
\xymatrix{
U \ar[d] \ar[r] & V \ar[d] \\
X \ar[r] & Y
}
$$
% P09-ERR-0781
其中 $U$、$V$ 是概形，竖直箭头均为 étale 态射，
$U \to X$ 为满射，并且上方的水平箭头
相对维数为 $1$ 且至多带结点奇性；以及
\item 存在 Zariski 覆盖 $Y = \bigcup_{i \in I} Y_i$，
% P09-ERR-0782
以及 $f^{-1}(Y_i) = \bigcup X_{ij}$，使得每个态射
$X_{ij} \to Y_i$ 都相对维数为 $1$ 且至多带结点奇性。
\end{enumerate}
\end{lemma}

\begin{proof}
% P09-ERR-0783
略去。
\end{proof}

\noindent
下面的引理说明，是否相对维数为 $1$ 且至多带结点奇性，
可以在纤维上检验。

\begin{lemma}
\label{spaces-more-morphisms-lemma-locus-where-nodal}
设 $S$ 是概形。
设 $f : X \to Y$ 是 $S$ 上代数空间的态射，并且它平坦且局部有限表示。
则存在最大的开子空间 $X' \subset X$，使得
$f|_{X'} : X' \to Y$ 相对维数为 $1$ 且至多带结点奇性。
此外，$X'$ 的构造与任意基变换交换。
\end{lemma}

\begin{proof}
选取交换图
$$
\xymatrix{
U \ar[d] \ar[r]_h & V \ar[d] \\
X \ar[r]^f & Y
}
$$
其中 $U$、$V$ 是概形，竖直箭头均为 étale 态射，并且
$U \to X$ 是满射。由概形情形的引理
（《代数曲线》引理 \ref{curves-lemma-locus-where-nodal}），
可找到最大的开子概形 $U' \subset U$，使得
$h|_{U'} : U' \to V$ 相对维数为 $1$ 且至多带结点奇性，
并且 $U'$ 的构造与基变换交换。
令 $X' \subset X$ 为其点对应于开子集
$\Im(|U'| \to |X|)$ 的开子空间。
由引理 \ref{spaces-more-morphisms-lemma-nodal-local}，$X' \to Y$
相对维数为 $1$ 且至多带结点奇性，而且 $X'$ 是具有此性质的
最大开子空间（这也推出 $U'$ 是 $X'$ 在 $U$ 中的逆像，但我们
不需要这一点）。因为基变换后同样如此，证明完成。
\end{proof}




\section{分解性质}
\label{spaces-more-morphisms-section-resolution-property}

\noindent
我们继续《空间的导出范畴》第
\ref{spaces-perfect-section-resolution-property} 节中的讨论。

\begin{situation}
\label{spaces-more-morphisms-situation-descend-resolution-property}
设 $S$ 是概形。设 $X$ 是 $S$ 上拟紧且拟分离的代数空间。
设 $V \to X$ 是满 étale 态射，其中 $V$ 是仿射概形
（这样的态射由《空间的性质》引理
\ref{spaces-properties-lemma-quasi-compact-affine-cover} 保证存在）。
选取交换图
$$
\xymatrix{
V \ar[rd]_\varphi \ar[rr]_j & & Y \ar[ld]^\pi \\
& X
}
$$
其中 $j$ 是开浸入，而 $\pi$ 是代数空间的有限态射
（这样的图由引理
\ref{spaces-more-morphisms-lemma-quasi-finite-separated-pass-through-finite} 保证存在）。
设 $\mathcal{I} \subset \mathcal{O}_Y$ 是 $Y$ 上有限型拟凝聚理想层，满足
$V(\mathcal{I}) = Y \setminus j(V)$（这样的理想层由《空间的极限》引理
\ref{spaces-limits-lemma-quasi-coherent-finite-type-ideals} 保证存在）。
\end{situation}

\begin{lemma}
\label{spaces-more-morphisms-lemma-surjection-in-Noetherian-case}
在情形 \ref{spaces-more-morphisms-situation-descend-resolution-property} 中，假设
$X$ 是诺特的。那么对任意凝聚 $\mathcal{O}_X$-模
$\mathcal{F}$，存在 $r \geq 0$、整数
$n_1, \ldots, n_r \geq 0$，以及满射
$$
\bigoplus\nolimits_{i = 1, \ldots, r} \pi_*(\mathcal{I}^{n_i})
\longrightarrow \mathcal{F}
$$
这是 $\mathcal{O}_X$-模意义下的满射。
\end{lemma}

\begin{proof}
% P09-ERR-0784
以 $\omega_{Y/X}$ 表示如下凝聚 $\mathcal{O}_Y$-模：存在
同构
$$
\pi_*\omega_{Y/X} \cong
\SheafHom_{\mathcal{O}_X}(\pi_*\mathcal{O}_Y, \mathcal{O}_X)
$$
它是 $\pi_*\mathcal{O}_Y$-模的同构，见《空间的态射》引理
\ref{spaces-morphisms-lemma-affine-equivalence-modules} 和
《空间上的下降》引理
\ref{spaces-descent-lemma-finite-over-finite-module}。
典范映射 $\mathcal{O}_X \to \pi_*\mathcal{O}_Y$ 给出典范映射
$$
\text{Tr}_\pi : \pi_*\omega_{Y/X} \longrightarrow \mathcal{O}_X
$$
% P09-ERR-0785
因为 $V$ 是诺特仿射概形，我们可以选取截面
$$
s_1, \ldots, s_r \in
\Gamma(V, \pi^*\mathcal{F} \otimes_{\mathcal{O}_Y} \omega_{Y/X})
$$
使它们生成凝聚模
% P09-ERR-0786
$\pi^*\mathcal{F} \otimes_{\mathcal{O}_X} \omega_{Y/X}$
在 $V$ 上的限制。
由《空间的上同调》引理
\ref{spaces-cohomology-lemma-homs-over-open}，可以选取整数
$n_i \geq 0$，使得 $s_i$ 延拓为映射
$s_i' : \mathcal{I}^{n_i} \to
\pi^*\mathcal{F} \otimes_{\mathcal{O}_Y} \omega_{Y/X}$。
将其推前到 $X$，得到映射
$$
\sigma_i :
\pi_*\mathcal{I}^{n_i}
\xrightarrow{\pi_*s'_i}
\pi_*(\pi^*\mathcal{F} \otimes_{\mathcal{O}_Y} \omega_{Y/X}) =
\mathcal{F} \otimes_{\mathcal{O}_X} \pi_*\omega_{Y/X}
\xrightarrow{\text{Tr}_\pi} \mathcal{F}
$$
其中等号来自《空间的上同调》引理
\ref{spaces-cohomology-lemma-finite-rings}。
为完成证明，我们将说明这些映射之和是满射。

\medskip\noindent
设 $x \in |X|$ 是 $X$ 的一点。设 $v \in |V|$ 是映到 $x$ 的一点。
可以选取 étale 邻域 $(U, u) \to (X, x)$，使得
$$
U \times_X Y = W \coprod W'
$$
（代数空间的不交并），其中 $W \to U$ 是同构，并且
唯一一点 $w \in W$ 位于 $u$ 上方，并映到一点 $v$；
后者属于 $V \subset Y$。
为说明这一点，可使用引理
\ref{spaces-more-morphisms-lemma-etale-splits-off-just-one-quasi-finite-part} 和
《Étale 态射》引理 \ref{etale-lemma-etale-etale-local}。
必要时进一步缩小 $U$，可假设 $W$ 通过投影映入
$V \subset Y$。由于 $\omega_{Y/X}$ 的构造与 étale 局部化交换，
可见
$$
\pi_*\omega_{Y/X}|_U = (\pi|_W)_*\omega_{W/U} \oplus (\pi|_{W'})_*\omega_{W'/U}
$$
我们有 $(\pi|_W)_*\omega_{W/U} = \mathcal{O}_U$，而这个同构由
迹映射 $\text{Tr}_\pi|_U$ 在上述分解的第一个直和项上的限制给出。
因为 $W$ 映入 $V$，可见
$\mathcal{I}^{n_i}|_W = \mathcal{O}_W$。因此
$$
\pi_*(\mathcal{I}^{n_i})|_U = \mathcal{O}_U \oplus
(W' \to U)_*(\mathcal{I}^{n_i}|_{W'})
$$
追踪交换图可知（略去细节），$\sigma_i|_U$ 在直和项
$\mathcal{O}_U$ 上就是映射 $\mathcal{O}_U \to \mathcal{F}$，
它对应于截面
$$
s_i|_W  \in
\Gamma(W,\pi^*\mathcal{F} \otimes_{\mathcal{O}_Y} \omega_{Y/X})
=
\Gamma(W, \mathcal{F}|_W \otimes_{\mathcal{O}_W} \omega_{W/U})
=
\Gamma(U, \mathcal{F})
$$
因为截面 $s_i$ 生成模
$\pi^*\mathcal{F} \otimes_{\mathcal{O}_Y} \omega_{Y/X}$
在 $V$ 上的限制，
并且 $W$ 映入 $V$，所以
$\bigoplus \sigma_i$ 在 $U$ 上的限制是满射。
由于 $x$ 是任意一点，证明完成。
\end{proof}

\begin{lemma}
\label{spaces-more-morphisms-lemma-resolution-property-in-Noetherian-case}
在情形 \ref{spaces-more-morphisms-situation-descend-resolution-property} 中，假设
$X$ 是诺特的。那么 $X$ 具有分解性质，当且仅当
$\pi_*\mathcal{I}$ 是某个有限秩局部自由
$\mathcal{O}_X$-模的商。
\end{lemma}

\begin{proof}
由《空间的上同调》引理
\ref{spaces-cohomology-lemma-finite-pushforward-coherent}，
模 $\pi_*\mathcal{I}$ 是凝聚的。因此，如果 $X$ 具有分解性质，
则 $\pi_*\mathcal{I}$ 是某个有限秩局部自由
$\mathcal{O}_X$-模的商。
反过来，假设给定满射 $\mathcal{E} \to \pi_*\mathcal{I}$，
其源是一个有限秩局部自由 $\mathcal{O}_X$-模 $\mathcal{E}$。
注意，对所有 $n \geq 1$，都有满射
$$
\pi_*\mathcal{I} \otimes_{\mathcal{O}_X} \pi_*\mathcal{I}^n
\longrightarrow
\pi_*\mathcal{I}^{n + 1}
$$
因此，$\mathcal{E}^{\otimes n}$ 满射到
$\pi_*\mathcal{I}^n$，而且这对所有 $n \geq 1$ 成立。将此与引理
\ref{spaces-more-morphisms-lemma-surjection-in-Noetherian-case} 的结论结合，
便知 $X$ 具有分解性质。
\end{proof}

\begin{lemma}
\label{spaces-more-morphisms-lemma-resolution-property-one-sheaf}
在情形 \ref{spaces-more-morphisms-situation-descend-resolution-property} 中，
代数空间 $X$ 具有分解性质，当且仅当
$\pi_*\mathcal{I}$ 是某个有限秩局部自由
$\mathcal{O}_X$-模的商。
\end{lemma}

\begin{proof}
对于推前 $\pi_*\mathcal{G}$，若这里的源是有限型拟凝聚
$\mathcal{O}_Y$-模 $\mathcal{G}$，则由《空间上的下降》引理
\ref{spaces-descent-lemma-finite-over-finite-module}，
可知它是有限型拟凝聚 $\mathcal{O}_X$-模。
特别地，如果 $X$ 具有分解性质，则由《空间的导出范畴》定义
\ref{spaces-perfect-definition-resolution-property}，
$\pi_*\mathcal{I}$ 是某个有限秩局部自由
$\mathcal{O}_X$-模的商。

\medskip\noindent
假设有满射 $\mathcal{E} \to \pi_*\mathcal{I}$，
其源是有限秩局部自由 $\mathcal{O}_X$-模 $\mathcal{E}$。
在证明的其余部分中，我们把问题约化到引理
\ref{spaces-more-morphisms-lemma-resolution-property-in-Noetherian-case}
已经处理的诺特情形，以说明 $X$ 具有分解性质。
我们建议读者跳过证明的其余部分。

\medskip\noindent
第一个约化是，可以把 $X$ 视为
$\Spec(\mathbf{Z})$ 上的代数空间，见
《空间》定义 \ref{spaces-definition-base-change}。
（这不影响引理的条件或结论。）

\medskip\noindent
由《空间的极限》引理
\ref{spaces-limits-lemma-finite-in-finite-and-finite-presentation}，
可以写成 $Y = \lim Y_i$，其中 $Y_i$ 在 $X$ 上有限且有限表示，
而转移映射是闭浸入。考虑闭子空间
$Z = V(\mathcal{I})$，它位于 $Y$ 中。由于 $\mathcal{I}$ 是有限型的，
态射 $Z \to Y$ 是有限表示的。因此，由《空间的极限》引理
\ref{spaces-limits-lemma-descend-finite-presentation}，
可以找到 $i$ 和有限表示态射 $Z_i \to Y_i$，其到 $Y$ 的
基变换是 $Z \to Y$。
对 $i' \geq i$，记
$Z_{i'} = Z_i \times_{Y_i} Y_{i'}$。
增大 $i$ 后，由《空间的极限》引理
\ref{spaces-limits-lemma-descend-closed-immersion}，
可假设 $Z_i \to Y_i$ 是（有限表示的）闭浸入。
% P09-ERR-0787
以 $\mathcal{I}_i \subset \mathcal{O}_{Y_i}$ 表示
$Z_i$ 的理想层，并以 $\pi_i : Y_i \to X$ 表示结构态射。
对 $i' \geq i$ 也同样记号。由于
$Z = \lim_{i' \geq i} Z_{i'}$，我们有
$$
\pi_*\mathcal{I} = \colim \pi_{i', *}\mathcal{I}_{i'}
$$
% P09-ERR-0788
这个系统中的转移映射全都是满射；这由映射
$\pi_{i, *}\mathcal{O}_{Y_i} \to \pi_{i', *}\mathcal{O}_{Y_{i'}}$
的满射性以及事实
$Z_{i'} = Z_i \times_{Y_i} Y_{i'}$ 推出。
由《空间的上同调》引理
\ref{spaces-cohomology-lemma-finite-presentation-quasi-compact-colimit}，
对某个 $i' \geq i$，映射
$\mathcal{E} \to \pi_*\mathcal{I}$ 可提升为映射
$\mathcal{E} \to \pi_{i', *}\mathcal{I}_{i'}$。
增大 $i'$ 后，该映射
$\mathcal{E} \to \pi_{i', *}\mathcal{I}_{i'}$
成为满射（否则，具有满转移映射的余核系统之余极限非零）。
这把我们约化到下一段讨论的情形。

\medskip\noindent
假设 $X$ 是 $\mathbf{Z}$ 上的代数空间，并且
$Y \to X$ 是有限表示的。由绝对诺特逼近，可以把
$X = \lim X_i$ 写成有向极限，其中每个 $X_i$ 都是
$\mathbf{Z}$ 上有限型的拟分离代数空间，而转移态射是仿射的，
见《空间的极限》命题 \ref{spaces-limits-proposition-approximate}。
由于 $\pi : Y \to X$ 是有限表示的，由《空间的极限》引理
\ref{spaces-limits-lemma-descend-finite-presentation}，
可以找到 $i$ 和有限表示态射
$\pi_i : Y_i \to X_i$，其到 $X$ 的基变换是 $\pi$。
增大 $i$ 后，由《空间的极限》引理
\ref{spaces-limits-lemma-descend-finite}，
可假设 $\pi_i$ 是有限的。
接着，由《空间的极限》引理
\ref{spaces-limits-lemma-descend-invertible-modules}，
可假设存在有限秩局部自由
$\mathcal{O}_{X_i}$-模 $\mathcal{E}_i$，
其到 $X$ 的拉回是 $\mathcal{E}$。
还可假设存在映射
$\mathcal{E}_i \to \pi_{i, *}\mathcal{O}_{Y_i}$，
其到 $X$ 的拉回是复合映射
$\mathcal{E} \to \pi_*\mathcal{I} \to \pi_*\mathcal{O}_Y$，
见《空间的极限》引理
\ref{spaces-limits-lemma-descend-modules-finite-presentation}。
余核
$$
\mathcal{E}_i \to \pi_{i, *}\mathcal{O}_{Y_i} \to \mathcal{Q}_i \to 0
$$
% P09-ERR-0789
是凝聚的 $\mathcal{O}_{Y_i}$-模，其到 $X$ 的拉回是（有限表示）余核
$\mathcal{Q}$；这个余核来自映射
$\mathcal{E} \to \pi_*\mathcal{O}_Y$。换言之，我们有
$\mathcal{Q} = \pi_*(\mathcal{O}_Y/\mathcal{I})$。考虑映射
$$
\mathcal{E}_i
\otimes_{\mathcal{O}_{X_i}}
\pi_{i, *} \mathcal{O}_{Y_i}
\longrightarrow
\pi_{i, *}\mathcal{O}_{Y_i}
\otimes_{\mathcal{O}_{X_i}}
\pi_{i, *} \mathcal{O}_{Y_i}
\to
\pi_{i, *} \mathcal{O}_{Y_i}
\to
\mathcal{Q}_i
$$
其中第二个箭头由
$\pi_{i, *}\mathcal{O}_{Y_i}$ 的代数结构给出。
% P09-ERR-0790
这个映射到 $Y$ 的拉回为零，因为
$\mathcal{E} \to \pi_*\mathcal{O}_Y$ 的像是理想
$\pi_*\mathcal{I}$。因此，由《空间的极限》引理
\ref{spaces-limits-lemma-descend-modules-finite-presentation}，
增大 $i$ 后，可假设所显示的复合映射为零。
% P09-ERR-0791
这恰好意味着 $\mathcal{E}_i \to \pi_{i, *}\mathcal{O}_{Y_i}$
的像具有形式 $\pi_{i, *}\mathcal{I}_i$，
其中 $\mathcal{I}_i \subset \mathcal{O}_{Y_i}$ 是某个凝聚理想层。
由于 $\mathcal{E}_i \to \pi_{i, *}\mathcal{O}_{Y_i}$
拉回为 $\mathcal{E} \to \pi_*\mathcal{O}_Y$，可见
$\mathcal{I}_i$ 到 $Y$ 的拉回生成 $\mathcal{I}$。
% P09-ERR-0792
以 $V_i \subset Y_i$ 表示其补集为
$V(\mathcal{I}_i) \subset Y_i$ 的开子空间。
由上述评注，$V$ 是 $V_i$ 的逆像。
增大 $i$ 后，可假设 $V_i$ 是仿射的，并且
$\pi_i|_{V_i} : V_i \to X_i$ 是 étale 的，见
《空间的极限》引理
\ref{spaces-limits-lemma-limit-is-affine} 和
\ref{spaces-limits-lemma-descend-etale}。
至此，可应用引理
\ref{spaces-more-morphisms-lemma-resolution-property-in-Noetherian-case}
断定 $X_i$ 具有分解性质。由于 $X \to X_i$ 是仿射的，
由《空间的导出范畴》引理
\ref{spaces-perfect-lemma-resolution-property-goes-up-affine}，
可断定 $X$ 也具有分解性质。
\end{proof}

\begin{lemma}
\label{spaces-more-morphisms-lemma-resolution-property-descends}
设 $S$ 是概形。
设 $X = \lim X_i$ 是 $S$ 上拟紧且拟分离代数空间的一个
直系统的极限，其转移态射都是仿射的。
% P09-ERR-0793
那么 $X$ 具有分解性质，当且仅当 $X_i$ 对某个 $i$ 具有分解性质。
\end{lemma}

\begin{proof}
如果 $X_i$ 具有分解性质，那么由《空间的导出范畴》引理
\ref{spaces-perfect-lemma-resolution-property-goes-up-affine}，
$X$ 也具有分解性质。
假设 $X$ 具有分解性质。
选取 $i \in I$。可以选取仿射概形 $V_i$ 和满 étale 态射
$V_i \to X_i$
（《空间的性质》引理
\ref{spaces-properties-lemma-quasi-compact-affine-cover}）。
可以选取嵌入 $j : V_i \to Y_i$，其中 $Y_i$ 在 $X_i$ 上
有限且有限表示
（引理 \ref{spaces-more-morphisms-lemma-quasi-finite-separated-pass-through-finite-addendum}）。
可以选取有限型拟凝聚理想
$\mathcal{I}_i \subset \mathcal{O}_{Y_i}$，使得
$V_i = Y_i \setminus V(\mathcal{I}_i)$
（《空间的极限》引理
\ref{spaces-limits-lemma-quasi-coherent-finite-type-ideals}）。
% P09-ERR-0794
以 $V \to Y \to X$ 表示 $V_i \to Y_i \to X_i$ 到 $X$ 的基变换。
以 $\mathcal{I} \subset \mathcal{O}_Y$ 表示理想
$\mathcal{I}_i$ 的拉回。由引理
\ref{spaces-more-morphisms-lemma-resolution-property-one-sheaf} 的容易方向，
存在有限秩局部自由 $\mathcal{O}_X$-模
$\mathcal{E}$ 和满射
$\mathcal{E} \to \pi_*\mathcal{I}$。
注意，由于 $\pi_i : Y_i \to X_i$ 是有限且有限表示的，
$\pi : Y \to X$ 也有限且有限表示，并且
$\mathcal{O}_{X_i}$-模
$\pi_{i, *}\mathcal{O}_{Y_i}$ 和
$\pi_{i, *}(\mathcal{O}_{Y_i}/\mathcal{I}_i)$
% P09-ERR-0795
都是有限表示的；它们拉回到 $X$ 后分别给出
$\pi_*\mathcal{O}_Y$ 和
$\pi_*(\mathcal{O}_Y/\mathcal{I})$。因此，由《空间的极限》引理
\ref{spaces-limits-lemma-descend-modules-finite-presentation}，
增大 $i$ 后，可以找到有限秩局部自由
$\mathcal{O}_{X_i}$-模 $\mathcal{E}_i$ 和映射
$\mathcal{E}_i \to \pi_{i, *}\mathcal{O}_{Y_i}$，
其到 $X$ 的基变换恢复复合映射
$\mathcal{E} \to \pi_*\mathcal{I} \to \pi_*\mathcal{O}_Y$。
有限表示 $\mathcal{O}_{X_i}$-模
$\Coker(\mathcal{E}_i \to \pi_{i, *}\mathcal{O}_{Y_i})$ 和
$\pi_{i, *}(\mathcal{O}_{Y_i}/\mathcal{I}_i)$
到 $X$ 的拉回作为 $\pi_*\mathcal{O}_Y$ 的商彼此相同。
因此，由《空间的极限》引理
\ref{spaces-limits-lemma-descend-modules-finite-presentation}，
可以假设二者相同；换言之，
% P09-ERR-0796
$\mathcal{E}_i \to \pi_{i, *}\mathcal{O}_{X_i}$ 的像等于
$\pi_{i, *}\mathcal{I}_i$。
于是，由引理 \ref{spaces-more-morphisms-lemma-resolution-property-one-sheaf}，
可断定 $X_i$ 具有分解性质。
\end{proof}

\begin{lemma}
\label{spaces-more-morphisms-lemma-resolution-property-affine-diagonal}
设 $S$ 是概形。
设 $X$ 是具有分解性质的拟紧且拟分离代数空间。
那么 $X$ 在 $\mathbf{Z}$ 上具有仿射对角线
（定义见《空间的性质》定义
\ref{spaces-properties-definition-separated}）。
\end{lemma}

\begin{proof}
我们可以像概形情形那样证明，但这里改为从概形情形推出本引理。
首先，可以并且确实假设
$S = \Spec(\mathbf{Z})$。
接着，选取概形 $Y$ 和满整态射
$f : Y \to X$，见《合宜空间》引理
\ref{decent-spaces-lemma-there-is-a-scheme-integral-over}。
于是 $f$ 是仿射的，故由《空间的导出范畴》引理
\ref{spaces-perfect-lemma-resolution-property-goes-up-affine}，
$Y$ 具有分解性质。因此，由概形情形，
概形 $Y$ 具有仿射对角线，见《概形的导出范畴》引理
\ref{perfect-lemma-resolution-property-affine-diagonal}。
接着，考虑交换图
$$
\xymatrix{
Y \ar[d] \ar[rr]_{\Delta_Y} & & Y \times_{\mathbf{Z}} Y \ar[d] \\
X \ar[rr]^{\Delta_X} & & X \times_{\mathbf{Z}} X
}
$$
注意，右侧竖直箭头是整的，因而是仿射的。
设 $W \to X \times_{\mathbf{Z}} X$ 是态射，其中 $W$ 是仿射的。
于是可见
$$
Y \times_{X \times_{\mathbf{Z}} X} W =
Y \times_{\Delta_Y, Y \times_{\mathbf{Z}} Y}
(Y \times_{\mathbf{Z}} Y) \times_{X \times_{\mathbf{Z}} X} W
$$
是仿射的。另一方面，$Y \to X$ 是整且满的，因此
$$
Y \times_{X \times_{\mathbf{Z}} X} W
\longrightarrow
X \times_{X \times_{\mathbf{Z}} X} W
$$
% P09-ERR-0797
是整满态射，因为它是 $Y \to X$ 到 $W$ 的基变换。
由《空间的极限》命题
\ref{spaces-limits-proposition-affine}，
可断定这个箭头的目标是仿射的。
于是 $\Delta_X$ 如所需是仿射的。
\end{proof}






\section{爆破与分解性质}
\label{spaces-more-morphisms-section-resolution-property-by-blowing-up}

\noindent
我们证明，作一次爆破后便满足分解性质。

\begin{lemma}
\label{spaces-more-morphisms-lemma-blow-up-resolution-property}
设 $S$ 为概形。设 $X$ 为 $S$ 上拟紧且拟分离的
代数空间。假设 $|X|$ 只有有限多个不可约分支。
则存在稠密拟紧开子空间 $U \subset X$ 和一个
$U$-容许爆破 $X' \to X$，使得代数空间
$X'$ 具有分解性质。
\end{lemma}

\begin{proof}
由《空间的极限》引理
\ref{spaces-limits-lemma-there-is-a-scheme-finite-over-refined}，
存在满、有限且有限表示的态射
$f : Y \to X$，其中 $Y$ 是概形；
% P09-ERR-0798
还存在拟紧稠密开子空间 $U \subset X$，使得
$f^{-1}(U) \to U$ 有限平展。
由《态射详论》引理 \ref{more-morphisms-lemma-blow-up-ample-family}，
存在拟紧稠密开子空间 $V \subset Y$ 和一个
$V$-容许爆破 $Y' \to Y$，使得 $Y'$ 具有丰沛可逆模族。
缩小 $U$ 后，可以假设 $f^{-1}(U) \subset V$
（细节略）。因此 $f' : Y' \to X$ 在 $U$ 上有限平展；
特别地，态射 $(f')^{-1}(U) \to U$ 有限局部自由。
由引理 \ref{spaces-more-morphisms-lemma-finite-after-blowing-up}，存在一个
$U$-容许爆破 $X' \to X$，使得严格变换 $Y''$
（即 $Y'$ 的严格变换）在 $X'$ 上有限局部自由。
% P09-ERR-0799
图示如下：
$$
\xymatrix{
Y'' \ar[d] \ar[r]_g &
Y' \ar[r] &
Y \ar[d] \\
X' \ar[rr] & & X
}
$$
由于 $g : Y'' \to Y'$ 是一次爆破
（《空间上的除子》引理 \ref{spaces-divisors-lemma-strict-transform}），
其中心是 $X' \to X$ 的爆破中心之逆像，
可知 $g : Y'' \to Y'$ 是射影态射，并且存在某个
$g$-丰沛可逆模于 $Y''$ 上。因此，由《态射详论》引理
\ref{more-morphisms-lemma-ample-family-ample-relative}，
可知 $Y''$ 具有丰沛可逆模族。
所以 $Y''$ 具有分解性质；见《概形的导出范畴》引理
\ref{perfect-lemma-resolution-property-ample-family}。
最后，由《空间的导出范畴》引理
\ref{spaces-perfect-lemma-resolution-property-goes-down-finite-flat}，
可断定 $X'$ 具有分解性质。
\end{proof}

\begin{lemma}
\label{spaces-more-morphisms-lemma-blow-up-resolution-property-filtered}
设 $S$ 为概形。设 $X$ 为 $S$ 上拟紧且拟分离的
代数空间。则存在 $t \geq 0$ 及闭子空间
$$
X \supset Z_0 \supset Z_1 \supset \ldots \supset Z_t = \emptyset
$$
使得 $Z_i \to X$ 有限表示，
$Z_0 \subset X$ 是加厚，并且对于每个
% P09-ERR-0800
$i = 0, \ldots t - 1$，存在一个
% P09-ERR-0801
$(Z_i \setminus Z_{i - 1})$-容许爆破
$Z'_i \to Z_i$，使得 $Z'_i$ 具有分解性质。
\end{lemma}

\begin{proof}
本段使用绝对诺特逼近，把问题约化到在
$\Spec(\mathbf{Z})$ 上有限表示的代数空间情形。
我们可以把 $X$ 看成 $\Spec(\mathbf{Z})$ 上的代数空间；见
《空间》定义 \ref{spaces-definition-base-change} 和
《空间的性质》定义 \ref{spaces-properties-definition-separated}。
因此可以应用《空间的极限》命题
\ref{spaces-limits-proposition-approximate}。
由此可找到仿射态射 $X \to X_0$，其中 $X_0$
在 $\mathbf{Z}$ 上有限表示。
如果能对 $X_0$ 证明本引理，那么把分层和各爆破中心拉回到 $X$，
就得到 $X$ 的结论；这里用到分解性质沿仿射态射上升
（《空间的导出范畴》引理
\ref{spaces-perfect-lemma-resolution-property-goes-up-affine}），
以及仿射态射的严格变换仍为仿射态射——细节略。
这就把问题约化到下一段讨论的情形。

\medskip\noindent
假设 $X$ 在 $\mathbf{Z}$ 上有限表示。
那么 $X$ 是诺特的，并且 $|X|$ 是有限维诺特拓扑空间
（具有有限多个不可约分支）。
因此可以对 $\dim(|X|)$ 作归纳。
由引理 \ref{spaces-more-morphisms-lemma-blow-up-resolution-property}，
存在稠密开子空间 $U \subset X$ 和一个 $U$-容许爆破
$X' \to X$，使得 $X'$ 具有分解性质。
令 $Z_0 = X$，并令 $Z_1 \subset X$ 为满足
$|Z_1| = |X| \setminus |U|$ 的约化闭子空间。
由归纳假设，可找到整数 $t \geq 0$ 及由闭子空间组成的滤过
$$
Z_1 \supset Z_{1, 0} \supset Z_{1, 1} \supset \ldots
\supset Z_{1, t} = \emptyset
$$
其中 $Z_{1, 0} \to Z_1$ 是加厚，并且存在
$(Z_{1, i} \setminus Z_{1, i + 1})$-容许爆破
$Z'_{1, i} \to Z_{1, i}$，使得 $Z'_{1, i}$
具有分解性质。由于 $Z_1$ 是约化的，故 $Z_1 = Z_{1, 0}$。
因此，可以令 $Z_i = Z_{1, i - 1}$ 和
$Z'_i = Z'_{1, i - 1}$（$i \geq 1$）；引理得证。
\end{proof}
% Generated mechanically from the frozen spaces-flat.tex target.
% Do not edit this generated file; edit the bound locale source instead.

% OK, start here.
%

\setcounter{chapter}{76}
\chapter{代数空间上的平坦性}


\phantomsection
\label{spaces-flat-section-phantom}




\section{引言}
\label{spaces-flat-section-introduction}

\noindent
本章在代数空间的框架下讨论平坦模与平坦态射的一些进阶结果。
我们强烈建议读者先阅读概形框架下的对应章节，见
《平坦性进阶》第 \ref{flat-section-introduction} 节。
一个参考文献是 Raynaud 与 Gruson 的论文 
\cite{GruRay}。






\section{不纯性}
\label{spaces-flat-section-impure}

\noindent
本节是《平坦性进阶》第 \ref{flat-section-impure} 节的对应版本。

\begin{situation}
\label{spaces-flat-situation-pre-pure}
设 $S$ 为概形。设 $f : X \to Y$ 为有限型良态态射\footnote{
拟分离态射是良态的，见《良态空间》引理
\ref{decent-spaces-lemma-properties-trivial-implications}。
对任意态射 $\Spec(k) \to Y$，其中 $k$ 是域，代数空间
$X_k$ 在 $k$ 上有限表示；这是因为它在 $k$ 上有限型，并且由
《良态空间》引理
\ref{decent-spaces-lemma-locally-Noetherian-decent-quasi-separated}
知它是拟分离的。}，它是 $S$ 上代数空间之间的态射。
另设 $\mathcal{F}$ 为有限型拟凝聚 $\mathcal{O}_X$-模。
最后，设 $y \in |Y|$ 为 $Y$ 的一点。
\end{situation}

\noindent
在此情形下，考虑概形 $T$、态射 $g : T \to Y$、
一点 $t \in T$（满足 $g(t)=y$）、特化 $t' \leadsto t$（在 $T$ 中），
以及一点 $\xi \in |X_T|$，它位于 $t'$ 上方。这里
$X_T = T \times_Y X$。
% P09-ERR-0802
图示为
\begin{equation}
\label{spaces-flat-equation-impurity}
\vcenter{
\xymatrix{
\xi \ar@{|->}[d] & \\
t' \ar@{~>}[r] & t \ar@{|->}[r] \ar[r] & y
}
}
\quad\quad
\vcenter{
\xymatrix{
X_T \ar[d]_{f_T} \ar[r] & X \ar[d]^f \\
T \ar[r]^g & Y
}
}
\end{equation}
再用 $\mathcal{F}_T$ 表示 $\mathcal{F}$ 在 $X_T$ 上的拉回。

\begin{definition}
\label{spaces-flat-definition-impurity}
在情形 \ref{spaces-flat-situation-pre-pure} 中，称图
(\ref{spaces-flat-equation-impurity}) 定义了
{\it $\mathcal{F}$ 在 $y$ 上方的一个不纯性}，如果
$\xi \in \text{Ass}_{X_T/T}(\mathcal{F}_T)$ 且
$t \not \in f_T(\overline{\{\xi\}})$。
我们把这表述为：“令 $(g : T \to Y, t' \leadsto t, \xi)$
为 $\mathcal{F}$ 在 $y$ 上方的一个不纯性。”
\end{definition}

\noindent
换一种说法：$(g : T \to Y, t' \leadsto t, \xi)$ 是
$\mathcal{F}$ 在 $y$ 上方的一个不纯性，当且仅当不存在特化
$\xi \leadsto \theta$（在拓扑空间 $|X_T|$ 中）满足
$f_T(\theta)=t$。
非良态代数空间中的特化表现不佳。如果态射 $f$ 是良态的，
则 $X_T$ 都是良态代数空间（对上述每个态射 $g : T \to Y$）；见
《良态空间》定义 \ref{decent-spaces-definition-relative-conditions}。

\begin{lemma}
\label{spaces-flat-lemma-impure-limit}
在情形 \ref{spaces-flat-situation-pre-pure} 中，令
% P09-ERR-0803
$(g : T \to S, t' \leadsto t, \xi)$ 为
$\mathcal{F}$ 在 $y$ 上方的一个不纯性。假设
$T = \lim_{i \in I} T_i$ 是 $Y$ 上仿射概形的有向极限。
那么对某个 $i$，三元组
$(T_i \to Y, t'_i \leadsto t_i, \xi_i)$ 是
$\mathcal{F}$ 在 $y$ 上方的一个不纯性。
\end{lemma}

\begin{proof}
陈述中的记号含义如下。令 $p_i : T \to T_i$ 为投影态射，
令 $t_i=p_i(t)$ 和 $t'_i=p_i(t')$。
最后，$\xi_i \in |X_{T_i}|$ 是 $\xi$ 的像。由
《空间上的除子》引理
\ref{spaces-divisors-lemma-base-change-relative-assassin}，有
$\xi_i \in \text{Ass}_{X_{T_i}/T_i}(\mathcal{F}_{T_i})$。
因此只须证明
$t_i \not \in f_{T_i}(\overline{\{\xi_i\}})$ 对某个 $i$ 成立。

\medskip\noindent
令 $Z_i \subset X_{T_i}$ 为
$\overline{\{\xi_i\}} \subset |X_{T_i}|$ 上的约化诱导概形结构，
并令 $Z \subset X_T$ 为
$\overline{\{\xi\}} \subset |X_T|$ 上的约化诱导概形结构。
于是由《空间的极限》引理
\ref{spaces-limits-lemma-inverse-limit-irreducibles}，有
$Z=\lim Z_i$（该引理适用，因为每个 $X_{T_i}$ 都是良态的）。
选取域 $k$ 及态射 $\Spec(k) \to T$，其像为 $t$。则
$$
\emptyset =
Z \times_T \Spec(k) = (\lim Z_i) \times_{(\lim T_i)} \Spec(k)
= \lim Z_i \times_{T_i} \Spec(k)
$$
因为极限与纤维积交换（即极限彼此交换）。
每个 $Z_i \times_{T_i} \Spec(k)$ 都拟紧，因为
$X_{T_i} \to T_i$ 有限型，从而 $Z_i \to T_i$ 有限型。
故 $Z_i \times_{T_i} \Spec(k)$ 对某个 $i$ 为空；
这是由《空间的极限》引理 \ref{spaces-limits-lemma-limit-nonempty} 得出的。
由于复合态射 $\Spec(k) \to T \to T_i$ 的像是 $t_i$，
这就给出所需结论。
\end{proof}


\noindent
不纯性沿平坦基变换上升。

\begin{lemma}
\label{spaces-flat-lemma-flat-ascent-impurity}
在情形 \ref{spaces-flat-situation-pre-pure} 中，设
$(Y_1,y_1) \to (Y,y)$ 为 $S$ 上带点代数空间之间的态射。
假设 $Y_1 \to Y$ 在 $y_1$ 处平坦。如果
$(T \to Y,t' \leadsto t,\xi)$ 是
$\mathcal{F}$ 在 $y$ 上方的一个不纯性，那么存在一个不纯性
$(T_1 \to Y_1,t_1' \leadsto t_1,\xi_1)$：它是模
$\mathcal{F}_1$ 的一个不纯性；这个模是 $\mathcal{F}$ 在
$X_1=Y_1\times_YX$ 上的拉回，并且它位于 $y_1$ 上方；
其中 $T_1$ 在 $Y_1\times_YT$ 上平展。
\end{lemma}

\begin{proof}
选取平展态射 $T_1 \to Y_1\times_YT$，其中 $T_1$ 是概形，
并令 $t_1\in T_1$ 为映到 $y_1$ 和 $t$ 的一点。
由《空间的性质》引理
\ref{spaces-properties-lemma-points-cartesian}，
可以找到这样的二元组 $(T_1,t_1)$。
概形态射 $T_1\to T$ 在 $t_1$ 处平坦
（使用《空间的态射》引理
\ref{spaces-morphisms-lemma-base-change-flat}
以及代数空间平坦态射的定义）。
% P09-ERR-0804
于是存在特化 $t'_1\leadsto t_1$，它位于
$t'\leadsto t$ 上方；见《态射》引理
\ref{morphisms-lemma-generalizations-lift-flat}。
选取一点 $\xi_1\in |X_{T_1}|$，它映到 $t'_1$ 和 $\xi$，
并且 $\xi_1\in
\text{Ass}_{X_{T_1}/T_1}(\mathcal{F}_{T_1})$；
% P09-ERR-0805
它对应于 $\Spec(\kappa(t'_1)\otimes_{\kappa(t')}\kappa(\xi))$
的一点。由《空间上的除子》引理
\ref{spaces-divisors-lemma-base-change-relative-assassin}，
这样的选取是可能的。
由于闭包 $Z_1$（即 $\{\xi_1\}$ 在 $|X_{T_1}|$ 中的闭包）
映入 $\{\xi\}$ 在 $|X_T|$ 中的闭包，可知 $Z_1$
在 $|T_1|$ 中的像不可能包含 $t_1$。
因此，$(T_1\to Y_1,t'_1\leadsto t_1,\xi_1)$
是 $\mathcal{F}_1$ 在
% P09-ERR-0806
$Y_1$ 上方的一个不纯性。
\end{proof}

\begin{lemma}
\label{spaces-flat-lemma-pure-along-X-y}
在情形 \ref{spaces-flat-situation-pre-pure} 中，设 $\overline{y}$ 为位于
$y$ 上方的几何点。令
$\mathcal{O}=\mathcal{O}_{Y,\overline{y}}$
为 $Y$ 在 $\overline{y}$ 处的平展局部环。记
$Y^{sh}=\Spec(\mathcal{O})$、$X^{sh}=X\times_YY^{sh}$，
并以 $\mathcal{F}^{sh}$ 表示 $\mathcal{F}$ 在 $X^{sh}$ 上的拉回。
下列条件等价：
\begin{enumerate}
\item 存在不纯性
$(Y^{sh}\to Y,y'\leadsto\overline{y},\xi)$，它是
$\mathcal{F}$ 在 $y$ 上方的不纯性；
% P09-ERR-0807
\item $\text{Ass}_{X^{sh}/Y^{sh}}(\mathcal{F}^{sh})$
中的每一点都特化到闭纤维 $X_{\overline{y}}$ 的一点；
\item 存在不纯性 $(T\to Y,t'\leadsto t,\xi)$，它是
$\mathcal{F}$ 在 $y$ 上方的不纯性，并且
$(T,t)\to(Y,y)$ 是平展邻域；
\item 存在不纯性 $(T\to Y,t'\leadsto t,\xi)$，它是
$\mathcal{F}$ 在 $y$ 上方的不纯性，并且 $T\to Y$ 在 $t$ 处拟有限。
\end{enumerate}
\end{lemma}

\begin{proof}
(1) 与 (2) 等价由定义立即可得。

\medskip\noindent
回忆 $\mathcal{O}=\mathcal{O}_{Y,\overline{y}}$
是 $\mathcal{O}(V)$ 的滤过余极限，其中
$(V,\overline{v})\to(Y,\overline{y})$
遍历所有平展邻域（《空间的性质》引理
\ref{spaces-properties-lemma-cofinal-etale}）。
此外，只须考虑仿射平展邻域 $V$。因此
$Y^{sh}=\Spec(\mathcal{O})=\lim\Spec(\mathcal{O}(V))=\lim V$。
于是由引理 \ref{spaces-flat-lemma-impure-limit}，(1) 蕴含 (3)。

\medskip\noindent
由于平展态射局部拟有限（《空间的态射》引理
\ref{spaces-morphisms-lemma-etale-locally-quasi-finite}），
可知 (3) 蕴含 (4)。

\medskip\noindent
最后，假设 (4) 成立。将 $T$ 替换为一个包含 $t$ 的开邻域后，
可以假设 $T\to Y$ 局部拟有限。由引理
\ref{spaces-flat-lemma-flat-ascent-impurity}，可找到不纯性
$(T_1\to Y^{sh},t_1'\leadsto t_1,\xi_1)$，
其中 $T_1\to T\times_YY^{sh}$ 平展。
由于平展态射局部拟有限，并使用《空间的态射》引理
\ref{spaces-morphisms-lemma-base-change-quasi-finite}
及《态射》引理
\ref{morphisms-lemma-composition-quasi-finite}，
可知 $T_1\to Y^{sh}$ 局部拟有限。
由于 $\mathcal{O}$ 严格亨泽尔，可以应用《态射详论》引理
\ref{more-morphisms-lemma-etale-makes-quasi-finite-finite-at-point}；
将 $T_1$ 替换为一个包含 $t_1$ 的既开又闭邻域后，
可以假设 $T_1\to Y^{sh}=\Spec(\mathcal{O})$ 有限。
令 $\theta\in|X^{sh}|$ 为 $\xi_1$ 的像，令
$y'\in\Spec(\mathcal{O})$ 为 $t_1'$ 的像。
由《空间上的除子》引理
\ref{spaces-divisors-lemma-base-change-relative-assassin}，有
$\theta\in\text{Ass}_{X^{sh}/Y^{sh}}(\mathcal{F}^{sh})$。
由于 $\pi:X_{T_1}\to X^{sh}$ 有限，它诱导闭映射
$|X_{T_1}|\to|X^{sh}|$。所以
$\overline{\{\xi_1\}}$ 的像是 $\overline{\{\theta\}}$。
由此，$(Y^{sh}\to Y,y'\leadsto\overline{y},\theta)$
是 $\mathcal{F}$ 在 $y$ 上方的一个不纯性，证明完成。
\end{proof}


\section{相对纯模}
\label{spaces-flat-section-pure}

\noindent
本节是《平坦性进阶》第 \ref{flat-section-pure} 节的对应版本。

\begin{definition}
\label{spaces-flat-definition-pure}
在情形 \ref{spaces-flat-situation-pre-pure} 中：
\begin{enumerate}
\item 称 $\mathcal{F}$ 在 $y$ 上方{\it 纯}，如果引理
\ref{spaces-flat-lemma-pure-along-X-y} 中的等价条件{\bf 均不}成立。
\item 称 $\mathcal{F}$ 在 $y$ 上方{\it 普遍纯}，如果不存在
$\mathcal{F}$ 在 $y$ 上方的任何不纯性。
\item 称 $X$ 在 $y$ 上方{\it 纯}，如果 $\mathcal{O}_X$
在 $y$ 上方纯。
\item 称 $\mathcal{F}$ {\it 普遍 $Y$-纯}，或称它
{\it 相对于 $Y$ 普遍纯}，如果 $\mathcal{F}$ 在 $y$ 上方普遍纯，
并且这对每个 $y \in |Y|$ 均成立。
\item 称 $\mathcal{F}$ {\it $Y$-纯}，或称它
{\it 相对于 $Y$ 纯}，如果 $\mathcal{F}$ 在 $y$ 上方纯，
并且这对每个 $y \in |Y|$ 均成立。
\item 称 $X$ {\it $Y$-纯}，或称它{\it 相对于 $Y$ 纯}，
如果 $\mathcal{O}_X$ 相对于 $Y$ 纯。
\end{enumerate}
\end{definition}

\noindent
下面给出一些必需的引理。

\begin{lemma}
\label{spaces-flat-lemma-base-change-universally}
在情形 \ref{spaces-flat-situation-pre-pure} 中，下列条件等价：
% P09-ERR-0808
\begin{enumerate}
\item $\mathcal{F}$ 在 $y$ 上方普遍纯；
\item 对带点代数空间的每个态射
$(Y', y') \to (Y, y)$，拉回 $\mathcal{F}_{Y'}$
在 $y'$ 上方纯。
\end{enumerate}
特别地，$\mathcal{F}$ 相对于 $Y$ 普遍纯，当且仅当
每个基变换 $\mathcal{F}_{Y'}$（由 $\mathcal{F}$ 得到）
都相对于 $Y'$ 纯。
\end{lemma}

\begin{proof}
这是形式上的。
\end{proof}

\begin{lemma}
\label{spaces-flat-lemma-quasi-finite-base-change}
在情形 \ref{spaces-flat-situation-pre-pure} 中，设
$(Y', y') \to (Y, y)$ 为带点代数空间之间的态射。
如果 $Y' \to Y$ 在 $y'$ 处拟有限，并且 $\mathcal{F}$
在 $y$ 上方纯，那么 $\mathcal{F}_{Y'}$ 在 $y'$ 上方纯。
\end{lemma}

\begin{proof}
% P09-ERR-0809
如果 $(T \to Y', t' \leadsto t, \xi)$ 是
$\mathcal{F}_{Y'}$ 在 $y'$ 上方的一个不纯性，且
$T \to Y'$ 在 $t$ 处拟有限，那么
% P09-ERR-0810
$(T \to Y, t' \to t, \xi)$ 是 $\mathcal{F}$
在 $y$ 上方的一个不纯性，且 $T \to Y$ 在 $t$ 处拟有限；见
《空间的态射》引理
\ref{spaces-morphisms-lemma-composition-quasi-finite}。
因此，该引理由纯性的定义立即得到。
\end{proof}

\noindent
纯性满足平坦下降。

\begin{lemma}
\label{spaces-flat-lemma-flat-descend-pure}
在情形 \ref{spaces-flat-situation-pre-pure} 中，设
$(Y_1, y_1) \to (Y, y)$ 为带点代数空间之间的态射。
假设 $Y_1 \to Y$ 在 $y_1$ 处平坦。
\begin{enumerate}
\item 如果 $\mathcal{F}_{Y_1}$ 在 $y_1$ 上方纯，
那么 $\mathcal{F}$ 在 $y$ 上方纯。
\item 如果 $\mathcal{F}_{Y_1}$ 在 $y_1$ 上方普遍纯，
那么 $\mathcal{F}$ 在 $y$ 上方普遍纯。
\end{enumerate}
\end{lemma}

\begin{proof}
这是因为不纯性沿平坦基变换上升；见引理
\ref{spaces-flat-lemma-flat-ascent-impurity}。例如，对 (1)，任取不纯性
% P09-ERR-0811
$(T \to Y, t' \leadsto t, \xi)$，它是 $\mathcal{F}$
在 $y$ 上方的不纯性，且 $T \to Y$ 在 $t$ 处拟有限；
上述引理给出不纯性
$(T_1 \to Y_1, t_1' \leadsto t_1, \xi_1)$，它是模
$\mathcal{F}_1$ 的一个不纯性；该模是 $\mathcal{F}$ 在
$X_1 = Y_1 \times_Y X$ 上的拉回，并且这个不纯性位于
$y_1$ 上方，其中 $T_1$ 在 $Y_1 \times_Y T$ 上平展。
所以 $T_1 \to Y_1$ 在 $t_1$ 处拟有限，因为平展态射局部拟有限，
而局部拟有限态射的复合仍局部拟有限
（《空间的态射》引理
\ref{spaces-morphisms-lemma-etale-locally-quasi-finite} 和
\ref{spaces-morphisms-lemma-composition-quasi-finite}）。
对 (2) 的论证相同。
\end{proof}

\begin{lemma}
\label{spaces-flat-lemma-supported-on-closed}
在情形 \ref{spaces-flat-situation-pre-pure} 中，设 $i : Z \to X$ 为闭浸入，
并假设 $\mathcal{F} = i_*\mathcal{G}$，其中 $\mathcal{G}$ 是
$Z$ 上某个有限型拟凝聚层。那么 $\mathcal{G}$ 在 $y$ 上方
（普遍）纯，当且仅当 $\mathcal{F}$ 在 $y$ 上方（普遍）纯。
\end{lemma}

\begin{proof}
这由《空间上的除子》引理
\ref{spaces-divisors-lemma-relative-weak-assassin-finite} 得出。
\end{proof}

\begin{lemma}
\label{spaces-flat-lemma-proper-pure}
在情形 \ref{spaces-flat-situation-pre-pure} 中：
\begin{enumerate}
\item 如果 $\mathcal{F}$ 的支集在 $Y$ 上固有，那么
$\mathcal{F}$ 相对于 $Y$ 普遍纯。
\item 如果 $f$ 固有，那么 $\mathcal{F}$ 相对于 $Y$ 普遍纯。
\item 如果 $f$ 固有，那么 $X$ 相对于 $Y$ 普遍纯。
\end{enumerate}
\end{lemma}

\begin{proof}
先把 (1) 归约到 (2)。具体地，令 $Z \subset X$ 为
$\mathcal{F}$ 的概形论支集（《空间的态射》定义
\ref{spaces-morphisms-definition-scheme-theoretic-support}）。
令 $i : Z \to X$ 为相应的闭浸入，并写成
$\mathcal{F} = i_*\mathcal{G}$，其中存在某个有限型拟凝聚
$\mathcal{O}_Z$-模 $\mathcal{G}$。在情形 (1) 中，$Z \to Y$ 由假设是固有的。
因此由引理 \ref{spaces-flat-lemma-supported-on-closed}，情形 (1) 归约到情形 (2)。

\medskip\noindent
假设 $f$ 固有。令 $(g : T \to Y, t' \leadsto t, \xi)$ 为
$\mathcal{F}$ 在 $y$ 上方的一个不纯性。由于 $f$ 固有，
它是普遍闭的。因此 $f_T : X_T \to T$ 是闭映射。
由于 $f_T(\xi) = t'$，这蕴含
% P09-ERR-0812
$t \in f(\overline{\{\xi\}})$，矛盾。
\end{proof}







\section{有限型平坦模}
\label{spaces-flat-section-finite-type-flat}

\noindent
请与《平坦性进阶》第
\ref{flat-section-finite-type-flat-I}、
\ref{flat-section-finite-type-flat-II} 和
\ref{flat-section-finite-type-flat-III} 节比较。
% P09-ERR-0813
其中多数结果经平展局部化，立即给出代数空间情形的推论。

\begin{lemma}
\label{spaces-flat-lemma-existence-complete}
设 $S$ 为概形。设 $X \to Y$ 为 $S$ 上代数空间之间的有限型态射。
设 $\mathcal{F}$ 为有限型拟凝聚 $\mathcal{O}_X$-模。
设 $y \in |Y|$ 为一点。存在平展态射
$(Y', y') \to (Y, y)$，其中 $Y'$ 是仿射概形；还存在平展态射
$h_i : W_i \to X_{Y'}$，$i = 1, \ldots, n$，使得对每个
$i$，$\mathcal{F}_i/W_i/Y'$ 在 $y'$ 上存在一个完全逐层剥离，
其中 $\mathcal{F}_i$ 是 $\mathcal{F}$ 在 $W_i$ 上的拉回，并且
$|(X_{Y'})_{y'}| \subset \bigcup h_i(W_i)$。
\end{lemma}

\begin{proof}
所论问题在 $Y$ 上是平展局部的，因而可以假设 $Y$ 是仿射概形。
于是 $X$ 拟紧，所以可以选取仿射概形 $X'$ 以及满平展态射
$X' \to X$。然后把《平坦性进阶》引理
\ref{flat-lemma-existence-complete} 应用于 $X' \to Y$、
$(X' \to Y)^*\mathcal{F}$ 和 $y$，即得所需结论。
\end{proof}

\begin{lemma}
\label{spaces-flat-lemma-open-in-fibre-where-flat}
设 $S$ 为概形。设 $f : X \to Y$ 为 $S$ 上局部有限型的
代数空间态射。设 $\mathcal{F}$ 为有限型拟凝聚
$\mathcal{O}_X$-模。令 $y \in |Y|$ 且
$F = f^{-1}(\{y\}) \subset |X|$。则集合
$$
\{x \in F \mid \mathcal{F} \text{ flat over }Y\text{ 在 }x\}
$$
在 $F$ 中开。
\end{lemma}

\begin{proof}
选取概形 $V$、一点 $v \in V$ 以及平展态射 $V \to Y$，
此态射把 $v$ 映到 $y$。再选取概形 $U$ 以及满平展态射
$U \to V \times_Y X$。于是 $|U_v| \to F$ 是拓扑空间的开连续映射，
因为 $|U| \to |X|$ 连续且开。因此结论由概形情形得到；后者即
《平坦性进阶》引理 \ref{flat-lemma-open-in-fibre-where-flat}。
\end{proof}

\begin{lemma}
\label{spaces-flat-lemma-bourbaki-finite-type-general-base-at-point}
设 $S$ 为概形。设 $f : X \to Y$ 为 $S$ 上局部有限型的
代数空间态射。设 $x \in |X|$，其像为 $y \in |Y|$。
设 $\mathcal{F}$ 为 $X$ 上有限型拟凝聚层，设 $\mathcal{G}$
为 $Y$ 上拟凝聚层。如果 $\mathcal{F}$ 在 $x$ 处关于 $Y$ 平坦，
那么
$$
x \in \text{WeakAss}_X(\mathcal{F} \otimes_{\mathcal{O}_X} f^*\mathcal{G})
\Leftrightarrow
y \in \text{WeakAss}_Y(\mathcal{G})
\text{ 且 }
x \in \text{Ass}_{X/Y}(\mathcal{F}).
$$
\end{lemma}

\begin{proof}
选取交换图
$$
\xymatrix{
U \ar[d] \ar[r]_g & V \ar[d] \\
X \ar[r]^f & Y
}
$$
其中 $U$ 和 $V$ 是概形，且竖直箭头是满平展态射。
选取 $u \in U$ 映到 $x$。令 $\mathcal{E} = \mathcal{F}|_U$，
并令 $\mathcal{H} = \mathcal{G}|_V$。令 $v \in V$ 为 $u$ 的像。于是
$x \in \text{WeakAss}_X(\mathcal{F} \otimes_{\mathcal{O}_X} f^*\mathcal{G})$
当且仅当
% P09-ERR-0814
$u \in \text{WeakAss}_X(\mathcal{E} \otimes_{\mathcal{O}_X} g^*\mathcal{H})$；
这是由《空间上的除子》定义
\ref{spaces-divisors-definition-weakly-associated} 得出的。
类似地，$y \in \text{WeakAss}_Y(\mathcal{G})$ 当且仅当
$v \in \text{WeakAss}_V(\mathcal{H})$。
最后，由《空间上的除子》定义
\ref{spaces-divisors-definition-relative-weak-assassin}，
$x \in \text{Ass}_{X/Y}(\mathcal{F})$ 当且仅当
$u \in \text{Ass}_{U_v}(\mathcal{E}|_{U_v})$。
注意，$\mathcal{F}$ 在 $x$ 处平坦等价于 $\mathcal{E}$ 在 $u$ 处平坦；见
《空间的态射》定义 \ref{spaces-morphisms-definition-flat-module}。
对 $g : U \to V$、$\mathcal{E}$、$\mathcal{H}$、$u$ 和 $v$ 的等价性，
就是《平坦性进阶》引理
\ref{flat-lemma-bourbaki-finite-type-general-base-at-point}。
\end{proof}

\begin{lemma}
\label{spaces-flat-lemma-bourbaki-finite-type-general-base}
设 $S$ 为概形。设 $f : X \to Y$ 为 $S$ 上局部有限型的
代数空间态射。设 $\mathcal{F}$ 为 $X$ 上有限型拟凝聚层，
并且它在 $Y$ 上平坦。设 $\mathcal{G}$ 为 $Y$ 上拟凝聚层。则
$$
\text{WeakAss}_X(\mathcal{F} \otimes_{\mathcal{O}_X} f^*\mathcal{G}) =
\text{Ass}_{X/Y}(\mathcal{F}) \cap
|f|^{-1}(\text{WeakAss}_Y(\mathcal{G}))
$$
\end{lemma}

\begin{proof}
这是引理 \ref{spaces-flat-lemma-bourbaki-finite-type-general-base-at-point}
的直接推论。
\end{proof}

\begin{theorem}
\label{spaces-flat-theorem-finite-type-flat}
设 $S$ 为概形。设 $f : X \to Y$ 为 $S$ 上代数空间之间的态射。
设 $\mathcal{F}$ 为拟凝聚 $\mathcal{O}_X$-模。假设
\begin{enumerate}
\item $X \to Y$ 局部有限表示；
\item $\mathcal{F}$ 是有限型 $\mathcal{O}_X$-模；
\item $Y$ 的弱伴随点集在 $Y$ 中局部有限。
\end{enumerate}
那么 $U = \{x \in |X| : \mathcal{F}\text{ flat at }x\text{ 相对于 }Y\}$
在 $X$ 中开，且 $\mathcal{F}|_U$ 是有限表示
$\mathcal{O}_U$-模并在 $Y$ 上平坦。
\end{theorem}

\begin{proof}
条件 (3) 的含义如下：如果 $V \to Y$ 是满平展态射且 $V$ 是概形，
那么 $V$ 的弱伴随点集在概形 $V$ 上局部有限。（回忆，由
《空间上的除子》定义
\ref{spaces-divisors-definition-weakly-associated}，$V$ 的弱伴随点集
恰为 $Y$ 的弱伴随点集的逆像。）说明这一点之后，所论问题在
$X$ 和 $Y$ 上都是平展局部的，因而可以假设 $X$ 和 $Y$ 是概形。
于是结论由《平坦性进阶》定理
\ref{flat-theorem-finite-type-flat} 得到。
\end{proof}

\begin{lemma}
\label{spaces-flat-lemma-finite-type-flat-along-fibre-free-variant}
设 $S$ 为概形。设 $f : X \to Y$ 为 $S$ 上代数空间之间的态射。
设 $\mathcal{F}$ 为 $X$ 上拟凝聚层。令 $y \in |Y|$，并置
$F = f^{-1}(\{y\}) \subset |X|$。假设
\begin{enumerate}
\item $f$ 有限型；
\item $\mathcal{F}$ 有限型；
\item $\mathcal{F}$ 关于 $Y$ 平坦，并且这在所有 $x \in F$ 处均成立。
\end{enumerate}
那么存在平展态射 $(Y', y') \to (Y, y)$，其中 $Y'$ 是概形，
并且存在代数空间的交换图
$$
\xymatrix{
X \ar[d] & X' \ar[l]^g \ar[d] \\
Y & \Spec(\mathcal{O}_{Y', y'}) \ar[l]
}
$$
使得 $X' \to X \times_Y \Spec(\mathcal{O}_{Y', y'})$ 平展，
$|X'_{y'}| \to F$ 满射，$X'$ 仿射，并且
$\Gamma(X', g^*\mathcal{F})$ 是自由 $\mathcal{O}_{Y', y'}$-模。
\end{lemma}

\begin{proof}
选取平展态射 $(Y', y') \to (Y, y)$，其中 $Y'$ 是仿射概形。
于是 $X \times_Y Y'$ 拟紧。选取仿射概形 $X'$ 以及满平展态射
$X' \to X \times_Y Y'$。
% P09-ERR-0815
图示为
$$
\xymatrix{
X \ar[d] & X' \ar[l]^g \ar[d] \\
Y & Y' \ar[l]
}
$$
于是 $\mathcal{F}' = g^*\mathcal{F}$ 在 $Y'$ 上、在
$X'_{y'}$ 的所有点处平坦；见《空间的态射》引理
\ref{spaces-morphisms-lemma-base-change-module-flat}。
因此可以把概形情形的引理（《平坦性进阶》引理
\ref{flat-lemma-finite-type-flat-along-fibre-free-variant}）应用于态射
$X' \to Y'$、拟凝聚层 $g^*\mathcal{F}$ 和点 $y'$。
这给出平展态射 $(Y'', y'') \to (Y', y')$ 以及交换图
$$
\xymatrix{
X \ar[d] & X' \ar[l]^g \ar[d] & X'' \ar[l]^{g'} \ar[d] \\
Y & Y' \ar[l] & \Spec(\mathcal{O}_{Y'', y''}) \ar[l]
}
$$
为得到所需结论，取 $(Y'', y'') \to (Y, y)$ 以及
$g \circ g' : X'' \to X$ 即可。
\end{proof}

\begin{theorem}
\label{spaces-flat-theorem-check-flatness-at-associated-points}
设 $S$ 为概形。设 $f : X \to Y$ 为 $S$ 上局部有限型的
代数空间态射。设 $\mathcal{F}$ 为有限型拟凝聚
$\mathcal{O}_X$-模。设 $x \in |X|$，其像为 $y \in |Y|$。
置 $F = f^{-1}(\{y\}) \subset |X|$。考虑下列条件：
\begin{enumerate}
\item $\mathcal{F}$ 在 $x$ 处关于 $Y$ 平坦；
\item 对每个 $x' \in F \cap \text{Ass}_{X/Y}(\mathcal{F})$，
如果它特化到 $x$，那么
$\mathcal{F}$ 在 $x'$ 处关于 $Y$ 平坦。
\end{enumerate}
总有 (2) $\Rightarrow$ (1)。如果 $X$ 和 $Y$ 良态，
那么 (1) $\Rightarrow$ (2)。
\end{theorem}

\begin{proof}
假设 (2)。选取概形 $V$ 以及满平展态射 $V \to Y$。
选取概形 $U$ 以及满平展态射 $U \to V \times_Y X$。
选取 $u \in U$ 映到 $x$。令 $v \in V$ 为 $u$ 的像。
我们将由 $\mathcal{F}|_U = (U \to X)^*\mathcal{F}$、点 $u$
和相应的概形情形推出结论。
% P09-ERR-0816
$U_v$。这是可行的，因为
$\text{Ass}_{U/V}(\mathcal{F}|_U) \cap |U_v|$
等于 $\text{Ass}_{U_v}(\mathcal{F}|_{U_v})$，也等于
$F \cap \text{Ass}_{X/Y}(\mathcal{F})$ 的逆像。
由于映射 $|U_v| \to F$ 连续，$|U_v|$ 中的特化映到
$F$ 中的特化。因此条件 (2) 由 $U \to V$、
$\mathcal{F}|_U$ 和点 $u$ 继承。于是《平坦性进阶》定理
\ref{flat-theorem-check-flatness-at-associated-points} 适用，
从而 (1) 成立。

\medskip\noindent
如果 $Y$ 良态，那么点 $y$ 可以由拟紧单态射 $\Spec(k) \to Y$
表示（这由良态空间的定义得到，见《良态空间》定义
\ref{decent-spaces-definition-very-reasonable}）。于是
$F = |X_k|$；见《良态空间》引理
\ref{decent-spaces-lemma-topology-fibre}。
如果再假设 $X$ 良态（更一般地，如果 $f$ 良态；见
《良态空间》定义 \ref{decent-spaces-definition-relative-conditions}
和《良态空间》引理
\ref{decent-spaces-lemma-property-for-morphism-out-of-property}），
那么 $X_y$ 也是良态空间。此外，$F$ 中的特化可以提升为
$U_v \to X_y$ 中的特化；见《良态空间》引理
\ref{decent-spaces-lemma-decent-specialization}。
据此，反向蕴含显然成立，因为它在概形情形中成立。
\end{proof}

\begin{lemma}
\label{spaces-flat-lemma-check-along-closed-fibre}
设 $S$ 为局部概形，闭点为 $s$。设 $f : X \to S$ 为从代数空间
$X$ 到 $S$ 的局部有限型态射。设 $\mathcal{F}$ 为有限型拟凝聚
$\mathcal{O}_X$-模。假设
\begin{enumerate}
\item $\text{Ass}_{X/S}(\mathcal{F})$ 中每一点都特化到闭纤维
$X_s$ 的一点\footnote{例如，当 $f$ 有限型且 $\mathcal{F}$
沿 $X_s$ 纯，或当 $f$ 固有时，该条件成立。}；
\item $\mathcal{F}$ 关于 $S$ 平坦，并且这在 $X_s$ 的每一点处都成立。
\end{enumerate}
那么 $\mathcal{F}$ 在 $S$ 上平坦。
\end{lemma}

\begin{proof}
由定理 \ref{spaces-flat-theorem-check-flatness-at-associated-points}，
只须在 $\mathcal{F}$ 关于 $S$ 的相对伴随点处检验平坦性；
结论立即得出。
\end{proof}




\section{平坦有限表示模}
\label{spaces-flat-section-finitely-presented-flat}

\noindent
本节是《平坦性进阶》第
\ref{flat-section-finitely-presented-flat} 节的对应版本。

\begin{proposition}
\label{spaces-flat-proposition-finite-presentation-flat-at-point}
设 $S$ 为概形。设 $f : X \to Y$ 为 $S$ 上代数空间之间的态射。
设 $\mathcal{F}$ 为 $X$ 上拟凝聚层。设 $x \in |X|$，其像为
$y \in |Y|$。假设
\begin{enumerate}
\item $f$ 局部有限表示；
\item $\mathcal{F}$ 有限表示；
\item $\mathcal{F}$ 在 $x$ 处关于 $Y$ 平坦。
\end{enumerate}
那么存在带点概形的交换图
$$
\xymatrix{
(X, x) \ar[d] & (X', x') \ar[l]^g \ar[d] \\
(Y, y) & (Y', y') \ar[l]
}
$$
其水平箭头平展，$X'$、$Y'$ 仿射，并且
$\Gamma(X', g^*\mathcal{F})$ 是投射
$\Gamma(Y', \mathcal{O}_{Y'})$-模。
\end{proposition}

\begin{proof}
按此表述，该命题立即归约到概形情形；后者即
《平坦性进阶》命题
\ref{flat-proposition-finite-presentation-flat-at-point}。
\end{proof}

\begin{lemma}
\label{spaces-flat-lemma-finite-presentation-flat-along-fibre}
设 $S$ 为概形。设 $f : X \to Y$ 为 $S$ 上代数空间之间的态射。
设 $\mathcal{F}$ 为 $X$ 上拟凝聚层。令 $y \in |Y|$，并置
$F = f^{-1}(\{y\}) \subset |X|$。假设
\begin{enumerate}
\item $f$ 有限表示；
\item $\mathcal{F}$ 有限表示；
\item $\mathcal{F}$ 关于 $Y$ 平坦，并且这在所有 $x \in F$ 处均成立。
\end{enumerate}
那么存在代数空间的交换图
$$
\xymatrix{
X \ar[d] & X' \ar[l]^g \ar[d] \\
Y & Y' \ar[l]_h
}
$$
使得 $h$ 和 $g$ 平展，存在一点 $y' \in |Y'|$ 映到 $y$，
有 $F \subset g(|X'|)$，代数空间 $X'$、$Y'$ 仿射，并且
$\Gamma(X', g^*\mathcal{F})$ 是投射
$\Gamma(Y', \mathcal{O}_{Y'})$-模。
\end{lemma}

\begin{proof}
按此表述，该引理立即归约到概形情形；后者即
《平坦性进阶》引理
\ref{flat-lemma-finite-presentation-flat-along-fibre}。
\end{proof}




\section{纯性判据}
\label{spaces-flat-section-criterion-purity}

\noindent
本节是《平坦性进阶》第
\ref{flat-section-criterion-purity} 节的对应版本。

\begin{lemma}
\label{spaces-flat-lemma-associated-point-specializes}
设 $S$ 为概形，设 $X$ 为 $S$ 上局部有限型的良态代数空间。
设 $\mathcal{F}$ 为有限型拟凝聚 $\mathcal{O}_X$-模。
% P09-ERR-0817
设 $s \in S$，并且 $\mathcal{F}$ 关于 $S$ 平坦，
这一点在 $X_s$ 的所有点处都成立。
设 $x' \in \text{Ass}_{X/S}(\mathcal{F})$。如果 $\{x'\}$
在 $|X|$ 中的闭包与 $|X_s|$ 相交，那么该闭包与
$\text{Ass}_{X/S}(\mathcal{F}) \cap |X_s|$ 相交。
\end{lemma}

\begin{proof}
注意，$|X_s| \subset |X|$ 是 $|X|$ 中位于 $s \in S$ 上方的点集；见
《良态空间》引理 \ref{decent-spaces-lemma-topology-fibre}。
令 $t \in |X_s|$ 为 $x'$ 在 $|X|$ 中的一个特化。
选取仿射概形 $U$、一点 $u \in U$，以及平展态射
$\varphi : U \to X$；该态射把 $u$ 映到 $t$。由《良态空间》引理
\ref{decent-spaces-lemma-decent-specialization}，可以选取特化
$u' \leadsto u$，其中 $u'$ 映到 $x'$。置
% P09-ERR-0818
$g = f \circ \varphi$。注意，$s' = g(u') = f(x')$ 特化到 $s$。
按 $\text{Ass}_{X/S}(\mathcal{F})$ 的定义，有
$u' \in \text{Ass}_{U/S}(\varphi^*\mathcal{F})$。
由该引理的概形版本（《平坦性进阶》引理
\ref{flat-lemma-associated-point-specializes}），可知存在特化
$u' \leadsto u$，且
$u \in \text{Ass}_{U_s}(\varphi^*\mathcal{F}_s) =
\text{Ass}_{U/S}(\varphi^*\mathcal{F}) \cap U_s$。
因此 $x = \varphi(u) \in \text{Ass}_{X/S}(\mathcal{F})$
位于 $s$ 上方，引理得证。
\end{proof}

\begin{lemma}
\label{spaces-flat-lemma-contains-relative-ass-after-base-change}
设 $Y$ 为概形 $S$ 上的代数空间。设 $g : X' \to X$ 为
$Y$ 上代数空间之间的态射，其中 $X$ 在 $Y$ 上局部有限型。
设 $\mathcal{F}$ 为拟凝聚 $\mathcal{O}_X$-模。如果
$\text{Ass}_{X/Y}(\mathcal{F}) \subset g(|X'|)$，那么对任意态射
$Z \to Y$，都有
$\text{Ass}_{X_Z/Z}(\mathcal{F}_Z) \subset g_Z(|X'_Z|)$。
\end{lemma}

\begin{proof}
由《空间的性质》引理
\ref{spaces-properties-lemma-points-cartesian}，映射
$|X'_Z| \to |X_Z| \times_{|X|} |X'|$ 满射，因为
$X'_Z$ 等于 $X_Z \times_X X'$。由《空间上的除子》引理
\ref{spaces-divisors-lemma-base-change-relative-assassin}，映射
$|X_Z| \to |X|$ 把
$\text{Ass}_{X_Z/Z}(\mathcal{F}_Z)$ 映入
$\text{Ass}_{X/Y}(\mathcal{F})$。引理得证。
\end{proof}

\begin{lemma}
\label{spaces-flat-lemma-pure-on-top}
设 $Y$ 为概形 $S$ 上的代数空间。设 $g : X' \to X$ 为
$Y$ 上代数空间之间的平展态射。假设结构态射
$X' \to Y$ 和 $X \to Y$ 良态且有限型。
设 $\mathcal{F}$ 为有限型拟凝聚 $\mathcal{O}_X$-模。
设 $y \in |Y|$。置
% P09-ERR-0819
$F = f^{-1}(\{y\}) \subset |X|$。
\begin{enumerate}
\item 如果 $\text{Ass}_{X/Y}(\mathcal{F}) \subset g(|X'|)$，
且 $g^*\mathcal{F}$ 在 $y$ 上方（普遍）纯，那么
$\mathcal{F}$ 在 $y$ 上方（普遍）纯。
\item 如果 $\mathcal{F}$ 在 $y$ 上方纯，$g(|X'|)$ 包含 $F$，
且 $Y$ 是以 $y$ 为闭点的仿射局部概形，那么
$\text{Ass}_{X/Y}(\mathcal{F}) \subset g(|X'|)$。
\item 如果 $\mathcal{F}$ 在 $y$ 上方纯，$\mathcal{F}$
在 $F$ 的所有点处平坦，$g(|X'|)$ 包含
$\text{Ass}_{X/Y}(\mathcal{F}) \cap F$，且 $Y$ 是以
$y$ 为闭点的仿射局部概形，那么
$\text{Ass}_{X/Y}(\mathcal{F}) \subset g(|X'|)$。
% P09-ERR-0820
\item 此处添加更多内容。
\end{enumerate}
\end{lemma}

\begin{proof}
关于 $X \to Y$ 和 $X' \to Y$ 的假设保证我们可以把第
\ref{spaces-flat-section-impure} 节和第 \ref{spaces-flat-section-pure} 节的材料
应用于这些态射以及层 $\mathcal{F}$ 和 $g^*\mathcal{F}$。
由于 $g$ 平展，可知
$\text{Ass}_{X'/Y}(g^*\mathcal{F})$ 是
$\text{Ass}_{X/Y}(\mathcal{F})$ 的逆像，并且基变换后仍然如此。

\medskip\noindent
证明 (1)。假设
$\text{Ass}_{X/Y}(\mathcal{F}) \subset g(|X'|)$。
设 $(T \to Y, t' \leadsto t, \xi)$ 是 $\mathcal{F}$
在 $y$ 上方的一个不纯性。由引理
\ref{spaces-flat-lemma-contains-relative-ass-after-base-change}，有
$\text{Ass}_{X_T/T}(\mathcal{F}_T) \subset g_T(|X'_T|)$，
所以可以选取一点 $\xi' \in |X'_T|$ 映到 $\xi$。由上可知
$(T \to Y, t' \leadsto t, \xi')$ 是
$g^*\mathcal{F}$ 在
% P09-ERR-0821
$y'$ 上方的一个不纯性。这就证明了 (1)。

\medskip\noindent
证明 (2)。这是因为 $g(|X'|)$ 在 $|X|$ 中开，并且由纯性，
$\text{Ass}_{X/Y}(\mathcal{F})$ 中每一点都特化到 $F$ 的一点。

\medskip\noindent
证明 (3)。这是因为 $g(|X'|)$ 在 $|X|$ 中开，并且结合纯性与
引理 \ref{spaces-flat-lemma-associated-point-specializes}，
$\text{Ass}_{X/Y}(\mathcal{F})$ 中每一点都特化到
$\text{Ass}_{X/Y}(\mathcal{F}) \cap F$ 的一点。
\end{proof}

\begin{lemma}
\label{spaces-flat-lemma-finite-type-flat-pure-along-fibre-is-universal}
设 $S$ 为概形。设 $f : X \to Y$ 为 $S$ 上代数空间之间的态射。
设 $\mathcal{F}$ 为拟凝聚 $\mathcal{O}_X$-模。设 $y \in |Y|$。
假设
\begin{enumerate}
\item $f$ 良态且有限型；
\item $\mathcal{F}$ 有限型；
\item $\mathcal{F}$ 关于 $Y$ 平坦，并且这在所有位于 $y$ 上方的点处成立；
\item $\mathcal{F}$ 在 $y$ 上方纯。
\end{enumerate}
那么 $\mathcal{F}$ 在 $y$ 上方普遍纯。
\end{lemma}

\begin{proof}
考虑态射
% P09-ERR-0822
$\Spec(\mathcal{O}_{Y, \overline{y}}) \to Y$。
这是一个从严格亨泽尔局部环之谱出发的平坦态射，并把闭点映到
$y$。由引理 \ref{spaces-flat-lemma-flat-descend-pure}，可归约到下一段所述情形。

\medskip\noindent
假设 $Y$ 是严格亨泽尔局部环 $R$ 的谱，其闭点为 $y$。
由引理 \ref{spaces-flat-lemma-finite-type-flat-along-fibre-free-variant}，
存在平展态射 $g : X' \to X$，使得
$g(|X'|) \supset |X_y|$、$X'$ 仿射，并且
$\Gamma(X', g^*\mathcal{F})$ 是自由 $R$-模。
于是 $g^*\mathcal{F}$ 相对于 $Y$ 普遍纯；见《平坦性进阶》引理
\ref{flat-lemma-affine-locally-projective-pure}。
因此，由引理 \ref{spaces-flat-lemma-pure-on-top} 的 (1)，只须证明
$g(|X'|)$ 包含 $\text{Ass}_{X/Y}(\mathcal{F})$。
而这又由引理 \ref{spaces-flat-lemma-pure-on-top} 的 (2) 得出。
\end{proof}

\begin{lemma}
\label{spaces-flat-lemma-finite-type-flat-pure-is-universal}
设 $S$ 为概形。设 $f : X \to Y$ 为 $S$ 上代数空间之间的
良态有限型态射。设 $\mathcal{F}$ 为有限型拟凝聚
$\mathcal{O}_X$-模。假设 $\mathcal{F}$ 在 $Y$ 上平坦。
在这种情形下，$\mathcal{F}$ 相对于 $Y$ 纯，当且仅当
$\mathcal{F}$ 相对于 $Y$ 普遍纯。
\end{lemma}

\begin{proof}
这是引理
\ref{spaces-flat-lemma-finite-type-flat-pure-along-fibre-is-universal}
与各定义的直接推论。
\end{proof}

\begin{lemma}
\label{spaces-flat-lemma-universally-separating}
设 $Y$ 为概形 $S$ 上的代数空间。设 $g : X' \to X$ 为
$Y$ 上代数空间之间的平坦态射，其中 $X$ 在 $Y$ 上局部有限型。
设 $\mathcal{F}$ 为有限型拟凝聚 $\mathcal{O}_X$-模，并且它在
$Y$ 上平坦。如果
$\text{Ass}_{X/Y}(\mathcal{F}) \subset g(|X'|)$，那么典范映射
$$
\mathcal{F} \longrightarrow g_*g^*\mathcal{F}
$$
单射，并且在任意基变换后仍为单射。
\end{lemma}

\begin{proof}
最后一个断言的含义是：映射
$\mathcal{F}_Z \to (g_Z)_*g_Z^*\mathcal{F}_Z$
对任意态射 $Z \to Y$ 都单射。
由于关于相对伴随点集的假设在基变换下保持（引理
\ref{spaces-flat-lemma-contains-relative-ass-after-base-change}），
只须证明所显示的箭头单射。

\medskip\noindent
令 $\mathcal{K} = \Ker(\mathcal{F} \to g_*g^*\mathcal{F})$。
目标是证明 $\mathcal{K} = 0$。为此，只须证明
$\text{WeakAss}_X(\mathcal{K}) = \emptyset$；见《空间上的除子》引理
\ref{spaces-divisors-lemma-weakly-ass-zero}。由《空间上的除子》引理
\ref{spaces-divisors-lemma-ses-weakly-ass}，有
$\text{WeakAss}_X(\mathcal{K}) \subset \text{WeakAss}_X(\mathcal{F})$。
由于 $\mathcal{F}$ 平坦，引理
\ref{spaces-flat-lemma-bourbaki-finite-type-general-base} 给出
$\text{WeakAss}_X(\mathcal{F}) \subset \text{Ass}_{X/Y}(\mathcal{F})$。
按假设，任一点 $x$（属于 $\text{Ass}_{X/Y}(\mathcal{F})$）
都是某个 $x' \in |X'|$ 的像。由于 $g$ 平坦，局部环映射
$\mathcal{O}_{X, \overline{x}} \to \mathcal{O}_{X', \overline{x}'}$
忠实平坦，因而映射
$$
\mathcal{F}_{\overline{x}}
\longrightarrow
(g^*\mathcal{F})_{\overline{x}'} =
\mathcal{F}_{\overline{x}} \otimes_{\mathcal{O}_{X, \overline{x}}}
\mathcal{O}_{X', \overline{x}'}
$$
单射（见《代数》引理
\ref{algebra-lemma-faithfully-flat-universally-injective}）。
由于所显示的箭头经由
$\mathcal{F}_{\overline{x}} \to (g_*g^*\mathcal{F})_{\overline{x}}$
分解，可知 $\mathcal{K}_{\overline{x}} = 0$。因此 $x$
不可能是 $\mathcal{K}$ 的弱伴随点，结论得证。
\end{proof}







\section{平坦化函子}
\label{spaces-flat-section-F-zero}

\noindent
本节是《平坦性进阶》第
\ref{flat-section-flattening-functors} 节的对应版本。
我们建议读者初次阅读时跳过本节。

\begin{situation}
\label{spaces-flat-situation-iso}
设 $S$ 为概形。设 $f : X \to B$ 为 $S$ 上代数空间之间的态射。
设 $u : \mathcal{F} \to \mathcal{G}$ 为拟凝聚
$\mathcal{O}_X$-模之间的同态。对任意概形 $T$（在 $B$ 上），以
$u_T : \mathcal{F}_T \to \mathcal{G}_T$ 表示 $u$ 到 $T$ 的基变换；
换言之，$u_T$ 是 $u$ 经投影态射
$X_T = X \times_B T \to X$ 的拉回。在此情形下，可以考虑函子
\begin{equation}
\label{spaces-flat-equation-iso}
F_{iso} : (\Sch/B)^{opp} \longrightarrow \textit{Sets}, \quad
T \longrightarrow \left\{
\begin{matrix}
\{*\} & \text{如果 }u_T\text{ 是同构}, \\
\emptyset & \text{否则。}
\end{matrix}
\right.
\end{equation}
还有变体 $F_{inj}$、$F_{surj}$、$F_{zero}$，其中分别要求
$u_T$ 为单射、满射或零映射。
\end{situation}

\noindent
在情形 \ref{spaces-flat-situation-iso} 中，有时把函子
$F_{iso}$、$F_{inj}$、$F_{surj}$ 和 $F_{zero}$ 看作从
$(\Sch/S)^{opp} \to \textit{Sets}$ 的函子，并分别带有态射
$F_{iso} \to B$、$F_{inj} \to B$、$F_{surj} \to B$ 和
$F_{zero} \to B$。具体地，如果 $T$ 是 $S$ 上概形，
那么元素 $h \in F_{iso}(T)$ 是态射 $h : T \to B$，使得
$u$ 经 $h$ 的基变换是同构。特别地，当我们说 $F_{iso}$
是代数空间时，是指相应函子
$(\Sch/S)^{opp} \to \textit{Sets}$ 是代数空间。

\begin{lemma}
\label{spaces-flat-lemma-iso-sheaf}
在情形 \ref{spaces-flat-situation-iso} 中，函子
$F_{iso}$、$F_{inj}$、$F_{surj}$、$F_{zero}$
中的每一个都满足 fpqc 拓扑的层性质。
\end{lemma}

\begin{proof}
设 $\{T_i \to T\}_{i \in I}$ 为 $B$ 上概形的 fpqc 覆盖。置
% P09-ERR-0823
$X_i = X_{T_i} = X \times_S T_i$ 及 $u_i = u_{T_i}$。
注意，$\{X_i \to X_T\}_{i \in I}$ 是 $X_T$ 的 fpqc 覆盖；见
《空间上的拓扑》引理 \ref{spaces-topologies-lemma-fpqc}。
特别地，对每个 $x \in |X_T|$，存在 $i \in I$ 以及
$x_i \in |X_i|$ 映到 $x$。由于
% P09-ERR-0824
$\mathcal{O}_{X_T, \overline{x}} \to \mathcal{O}_{X_i, \overline{x_i}}$
平坦，因而忠实平坦（见《空间的态射》第
\ref{spaces-morphisms-section-flat} 节），可知 $(u_i)_{x_i}$
为单射、满射、双射或零映射，当且仅当 $(u_T)_x$
相应地为单射、满射、双射或零映射。引理得证。
\end{proof}

\begin{lemma}
\label{spaces-flat-lemma-iso-go-up}
在情形 \ref{spaces-flat-situation-iso} 中，设 $X' \to X$ 为代数空间的平坦态射。
以 $u' : \mathcal{F}' \to \mathcal{G}'$ 表示 $u$ 到 $X'$ 的拉回。
以 $F'_{iso}$、$F'_{inj}$、$F'_{surj}$、$F'_{zero}$ 表示
$\Sch/B$ 上由 $u'$ 得到的函子。
\begin{enumerate}
\item 如果 $\mathcal{G}$ 有限型，且 $|X'| \to |X|$ 的像包含
$\mathcal{G}$ 的支集，那么 $F_{surj} = F'_{surj}$ 且
$F_{zero} = F'_{zero}$。
\item 如果 $\mathcal{F}$ 有限型，且 $|X'| \to |X|$ 的像包含
$\mathcal{F}$ 的支集，那么 $F_{inj} = F'_{inj}$ 且
$F_{zero} = F'_{zero}$。
\item 如果 $\mathcal{F}$ 和 $\mathcal{G}$ 都有限型，且
$|X'| \to |X|$ 的像包含 $\mathcal{F}$ 与 $\mathcal{G}$ 的支集，
那么 $F_{iso} = F'_{iso}$。
\end{enumerate}
\end{lemma}

\begin{proof}
设 $v : \mathcal{H} \to \mathcal{E}$ 为代数空间 $Y$ 上拟凝聚模
之间的映射，设 $\varphi : Y' \to Y$ 为代数空间的满平坦态射。
那么 $v$ 为同构、单射、满射或零映射，当且仅当 $\varphi^*v$
相应地为同构、单射、满射或零映射。确实，对每个 $y \in |Y|$，
存在 $y' \in |Y'|$，并且局部环映射
$\mathcal{O}_{Y, \overline{y}} \to \mathcal{O}_{Y', \overline{y'}}$
忠实平坦（见《空间的态射》第
\ref{spaces-morphisms-section-flat} 节）。当然，检验单射性或零性时，
只须考察 $\mathcal{H}$ 的支集中的点；检验满射性时，只须考察
$\mathcal{E}$ 的支集中的点。此外，在引理陈述中的有限型假设下，
取支集与基变换交换；见《空间的态射》引理
\ref{spaces-morphisms-lemma-support-finite-type}。因此引理成立。
\end{proof}

\noindent
回忆，我们已在《空间的态射》定义
\ref{spaces-morphisms-definition-scheme-theoretic-support}
中定义了有限型拟凝聚模的概形论支集。

\begin{lemma}
\label{spaces-flat-lemma-iso-limits}
在情形 \ref{spaces-flat-situation-iso} 中：
\begin{enumerate}
\item 如果 $\mathcal{G}$ 有限型，且 $\mathcal{G}$ 的概形论支集在
$B$ 上拟紧，那么 $F_{surj}$ 保极限。
% P09-ERR-0825
\item 如果 $\mathcal{F}$ 有限型，且 $\mathcal{F}$ 的概形论支集在
$B$ 上拟紧，那么 $F_{zero}$ 保极限。
\item 如果 $\mathcal{F}$ 有限型、$\mathcal{G}$ 有限表示，且
$\mathcal{F}$ 与 $\mathcal{G}$ 的概形论支集在 $B$ 上拟紧，
那么 $F_{iso}$ 保极限。
\end{enumerate}
\end{lemma}

\begin{proof}
证明 (1)。令 $i : Z \to X$ 为 $\mathcal{G}$ 的概形论支集，
并把 $\mathcal{G}$ 看作 $Z$ 上有限型拟凝聚模。可以把 $X$ 替换为
$Z$，并把 $u$ 替换为映射 $i^*\mathcal{F} \to \mathcal{G}$
（略去细节）。因此可以假设 $f$ 拟紧且 $\mathcal{G}$ 有限型。
设 $T = \lim_{i \in I} T_i$ 为仿射 $B$-概形的有向极限，并假设
$u_T$ 满射。置
% P09-ERR-0823
$X_i = X_{T_i} = X \times_S T_i$ 以及
$u_i = u_{T_i} : \mathcal{F}_i = \mathcal{F}_{T_i}
\to \mathcal{G}_i = \mathcal{G}_{T_i}$。
为证明 (1)，须说明 $u_i$ 对某个 $i$ 满射。选取 $0 \in I$，并把
$I$ 替换为 $\{i \mid i \geq 0\}$。由于 $f$ 拟紧，$X_0$ 拟紧。
因此可以选取满平展态射 $\varphi_0 : W_0 \to X_0$，其中
$W_0$ 为仿射概形。置 $W = W_0 \times_{T_0} T$，并置
$W_i = W_0 \times_{T_0} T_i$（$i \geq 0$）。这些都是仿射概形，
% P09-ERR-0832
并带有满平展态射 $\varphi : W \to X_T$ 和
$\varphi_i : W_i \to X_i$。注意 $W = \lim W_i$。
于是 $\varphi^*u_T$ 满射，而只须证明 $\varphi_i^*u_i$ 对某个
$i$ 满射。这样便归约到仿射情形，即《代数》引理
\ref{algebra-lemma-module-map-property-in-colimit} 的 (2)。

\medskip\noindent
证明 (2)。假设 $\mathcal{F}$ 有限型，并且其概形论支集
% P09-ERR-0826
$Z \subset B$ 在 $B$ 上拟紧。设
$T = \lim_{i \in I} T_i$ 为仿射 $B$-概形的有向极限，并假设
$u_T$ 为零映射。置 $X_i = T_i \times_B X$，并以
$u_i : \mathcal{F}_i \to \mathcal{G}_i$ 表示拉回。选取
$0 \in I$，并把 $I$ 替换为 $\{i \mid i \geq 0\}$。置
$Z_0 = Z \times_X X_0$。由《空间的态射》引理
\ref{spaces-morphisms-lemma-support-finite-type}，
% P09-ERR-0827
$\mathcal{F}_i$ 的支集是 $|Z_0|$。由于 $|Z_0|$ 拟紧，
可以找到仿射概形 $W_0$ 及平展态射 $W_0 \to X_0$，使得
$|Z_0| \subset \Im(|W_0| \to |X_0|)$。置
$W = W_0 \times_{T_0} T$，并置
$W_i = W_0 \times_{T_0} T_i$（$i \geq 0$）。这些是带有平展态射
$\varphi : W \to X_T$ 和 $\varphi_i : W_i \to X_i$ 的仿射概形。
注意 $W = \lim W_i$，并且 $\mathcal{F}_T$ 与 $\mathcal{F}_i$
的支集分别包含在 $|W| \to |X_T|$ 与 $|W_i| \to |X_i|$ 的像中。
% P09-ERR-0828
现在 $\varphi^*u_T$ 单射，而只须证明 $\varphi_i^*u_i$ 对某个
$i$ 单射。这样便归约到仿射情形，即《代数》引理
\ref{algebra-lemma-module-map-property-in-colimit} 的 (1)。

\medskip\noindent
证明 (3)。可以完全类似于前两段，使用《代数》引理
\ref{algebra-lemma-module-map-property-in-colimit} 的 (3) 来证明。
也可以如下由 (1) 和 (2) 推出。设
$T = \lim_{i \in I} T_i$ 为仿射 $B$-概形的有向极限，并假设
$u_T$ 是同构。由 (1)，存在 $0 \in I$，使得 $u_{T_0}$ 满射。
置 $\mathcal{K} = \Ker(u_{T_0})$，并考虑拟凝聚模的映射
$v : \mathcal{K} \to \mathcal{F}_{T_0}$。对 $i \geq 0$，
基变换 $v_{T_i}$ 为零映射，当且仅当 $u_i$ 是同构。此外，
$v_T$ 为零映射。由于 $\mathcal{G}_{T_0}$ 有限表示、
$\mathcal{F}_{T_0}$ 有限型且 $u_{T_0}$ 满射，可知
$\mathcal{K}$ 有限型（《位形上的模》引理
\ref{sites-modules-lemma-kernel-surjection-finite-onto-finite-presentation}）。
显然，$\mathcal{K}$ 的支集包含在 $\mathcal{F}_{T_0}$ 的支集中，
而后者在 $T_0$ 上拟紧。因此可以应用 (2)，得 $v_{T_i}$ 对某个
$i$ 为零映射。
\end{proof}

\begin{lemma}
\label{spaces-flat-lemma-relate-zero-iso}
在情形 \ref{spaces-flat-situation-iso} 中，假设给定正合列
$$
\mathcal{F} \xrightarrow{u} \mathcal{G} \xrightarrow{v} \mathcal{H} \to 0
$$
那么按显然的记号，有 $F_{v, iso} = F_{u, zero}$。
\end{lemma}

\begin{proof}
由于拉回右正合，可见
$\mathcal{F}_T \to \mathcal{G}_T \to \mathcal{H}_T \to 0$
对每个概形 $T$（在 $B$ 上）正合。因此
% P09-ERR-0829
$u_T$ 满射，当且仅当 $v_T$ 是同构。
\end{proof}

\begin{lemma}
\label{spaces-flat-lemma-relate-zero-affine}
在情形 \ref{spaces-flat-situation-iso} 中，假设给定仿射态射 $i : Z \to X$
以及拟凝聚 $\mathcal{O}_Z$-模 $\mathcal{H}$，使得
$\mathcal{G} = i_*\mathcal{H}$。设
$v : i^*\mathcal{F} \to \mathcal{H}$ 为与 $u$ 伴随的映射。那么
\begin{enumerate}
\item $F_{v, zero} = F_{u, zero}$；
\item 如果 $i$ 是闭浸入，那么 $F_{v, surj} = F_{u, surj}$。
\end{enumerate}
\end{lemma}

\begin{proof}
设 $T$ 为 $B$ 上概形。以 $i_T : Z_T \to X_T$ 表示 $i$ 的基变换，
以 $\mathcal{H}_T$ 表示 $\mathcal{H}$ 到 $Z_T$ 的拉回。注意
$(i^*\mathcal{F})_T = i_T^*\mathcal{F}_T$ 且
$i_{T, *}\mathcal{H}_T = (i_*\mathcal{H})_T$。
第一条等式来自拉回的交换性，第二条来自《空间的上同调》引理
\ref{spaces-cohomology-lemma-affine-base-change}。因此 $u_T$ 与 $v_T$
仍是伴随映射。故 $u_T = 0$ 当且仅当 $v_T = 0$，这证明了 (1)。
在情形 (2) 中，$u_T$ 满射当且仅当 $v_T$ 满射，因为 $u_T$ 分解为
$$
\mathcal{F}_T \to
i_{T, *}i_T^*\mathcal{F}_T \xrightarrow{i_{T, *}v_T} i_{T, *}\mathcal{H}_T
$$
并且 $i_{T, *}$ 是正合函子，把 $Z_T$ 上拟凝聚模范畴完全忠实地
嵌入拟凝聚 $\mathcal{O}_{X_T}$-模范畴。见《空间的态射》引理
\ref{spaces-morphisms-lemma-i-star-equivalence}。
\end{proof}

\begin{lemma}
\label{spaces-flat-lemma-relate-zero-affine-push}
在情形 \ref{spaces-flat-situation-iso} 中，假设给定仿射态射
$g : X \to X'$。置
% P09-ERR-0830
$u' = f_*u : f_*\mathcal{F} \to f_*\mathcal{G}$。那么
$F_{u, iso} = F_{u', iso}$、$F_{u, inj} = F_{u', inj}$、
$F_{u, surj} = F_{u', surj}$ 且 $F_{u, zero} = F_{u', zero}$。
\end{lemma}

\begin{proof}
由《空间的上同调》引理
\ref{spaces-cohomology-lemma-affine-base-change}，有
$g_{T, *}u_T = u'_T$。此外，
% P09-ERR-0831
$g_{T, *} : \QCoh(\mathcal{O}_{X_T}) \to \QCoh(\mathcal{O}_X)$
是忠实正合函子，并且反映同构、单射和满射。
\end{proof}

\begin{situation}
\label{spaces-flat-situation-flat}
设 $S$ 为概形。设 $f : X \to Y$ 为 $S$ 上代数空间之间的态射。
设 $\mathcal{F}$ 为拟凝聚 $\mathcal{O}_X$-模。对任意概形
$T$（在 $Y$ 上），以 $\mathcal{F}_T$ 表示 $\mathcal{F}$ 到 $T$ 的基变换；
换言之，$\mathcal{F}_T$ 是 $\mathcal{F}$ 经投影态射
$X_T = X \times_Y T \to X$ 的拉回。由于平坦模的基变换仍平坦，
得到函子
\begin{equation}
\label{spaces-flat-equation-flat}
F_{flat} : (\Sch/Y)^{opp} \longrightarrow \textit{Sets}, \quad
T \longrightarrow \left\{
\begin{matrix}
\{*\} & \text{如果 } \mathcal{F}_T \text{ 平坦于 }T, \\
\emptyset & \text{否则。}
\end{matrix}
\right.
\end{equation}
\end{situation}

\noindent
在情形 \ref{spaces-flat-situation-flat} 中，有时把 $F_{flat}$ 看作函子
$(\Sch/S)^{opp} \to \textit{Sets}$，并赋有态射
$F_{flat} \to Y$。具体地，如果 $T$ 是 $S$ 上概形，
那么元素 $h \in F_{flat}(T)$ 是态射 $h : T \to Y$，使得
$\mathcal{F}$ 经 $h$ 的基变换在 $T$ 上平坦。特别地，当我们说
$F_{flat}$ 是代数空间时，是指相应函子
$(\Sch/S)^{opp} \to \textit{Sets}$ 是代数空间。

\begin{lemma}
\label{spaces-flat-lemma-flat}
在情形 \ref{spaces-flat-situation-flat} 中：
\begin{enumerate}
\item 函子 $F_{flat}$ 满足 fpqc 拓扑的层性质。
\item 如果 $f$ 拟紧且局部有限表示，并且 $\mathcal{F}$ 有限表示，
那么函子 $F_{flat}$ 保极限。
\end{enumerate}
\end{lemma}

\begin{proof}
(1) 由下述事实得到：如果 $T' \to T$ 是 $Y$ 上代数空间的
满平坦态射，那么 $\mathcal{F}_{T'}$ 在 $T'$ 上平坦，当且仅当
$\mathcal{F}_T$ 在 $T$ 上平坦；见《空间的态射》引理
\ref{spaces-morphisms-lemma-base-change-module-flat}。
如果 $f$ 还拟分离（即 $f$ 有限表示），则 (2) 由《空间的极限》引理
\ref{spaces-limits-lemma-descend-flat} 得出。一般情形下，先归约到基
仿射的情形，再用有限多个仿射开覆盖 $X$，归约到拟分离情形。
略去细节。
\end{proof}


\section{使映射为零}
\label{spaces-flat-section-zero-map}

\noindent
本节在关于概形的对应章节中没有类似内容。

\begin{situation}
\label{spaces-flat-situation-somewhat-closed}
设 $S = \Spec(R)$ 为仿射概形。设 $X$ 为 $S$ 上代数空间。
设 $u : \mathcal{F} \to \mathcal{G}$ 为拟凝聚
$\mathcal{O}_X$-模之间的映射。
% P09-ERR-0833
假设 $\mathcal{G}$ 在 $S$ 上平坦。
\end{situation}

\begin{lemma}
\label{spaces-flat-lemma-F-zero-somewhat-closed}
在情形 \ref{spaces-flat-situation-somewhat-closed} 中，设 $T \to S$ 为概形的
拟紧态射，并且基变换 $u_T$ 为零映射。
% P09-ERR-0834
那么存在闭子概形
$Z \subset S$，使得 (a) $T \to S$ 分解经过 $Z$，且
(b) 基变换 $u_Z$ 为零映射。
如果 $\mathcal{F}$ 是有限型 $\mathcal{O}_X$-模，且
$\mathcal{F}$ 的概形论支集拟紧，那么可以取 $Z \to S$ 为有限表示。
\end{lemma}

\begin{proof}
设 $U \to X$ 为代数空间的满平展态射，其中
$U = \coprod U_i$ 是仿射概形的不交并（见《空间的性质》引理
\ref{spaces-properties-lemma-cover-by-union-affines}）。
由引理 \ref{spaces-flat-lemma-iso-go-up}，可以把 $X$ 替换为 $U$。换言之，
可以假设 $X = \coprod X_i$ 是仿射概形 $X_i$ 的不交并。
假设对 $u_i = u|_{X_i}$ 可以证明本引理。于是得到闭子概形
$Z_i \subset S$，使得 $T \to S$ 分解经过 $Z_i$，并且
$u_{i, Z_i}$ 为零映射。如果
$Z_i = \Spec(R/I_i) \subset \Spec(R) = S$，那么取
$Z = \Spec(R/\sum I_i)$ 即可。因此可以假设
$X = \Spec(A)$ 为仿射概形。

\medskip\noindent
选取有限仿射开覆盖 $T = T_1 \cup \ldots \cup T_m$。
显然，可以把 $T$ 替换为 $\coprod_{j = 1, \ldots, m} T_j$。
因此可以假设 $T$ 仿射，记 $T = \Spec(R')$。
% P09-ERR-0835
设 $u : M \to N$ 为 $A$-模同态，它对应于
$u : \mathcal{F} \to \mathcal{G}$。
于是 $N$ 是平坦 $R$-模，因为 $\mathcal{G}$ 在 $S$ 上平坦。
引理的假设意味着映射
$$
M \otimes_R R' \to N \otimes_R R'
$$
为零。设 $z \in M$。由 Lazard 定理（《代数》定理
\ref{algebra-theorem-lazard}）以及 $\otimes$ 与余极限交换这一事实，
可以找到自由 $R$-模 $F_z$、元素 $\tilde z \in F_z$ 以及映射
$F_z \to N$，使得 $u(z)$ 是 $\tilde z$ 的像，并且 $\tilde z$
在 $F_z \otimes_R R'$ 中的像为零。选取一组基
$\{e_{z, \alpha}\}$ 作为 $F_z$ 的基，并写成
$\tilde z = \sum f_{z, \alpha} e_{z, \alpha}$，其中
$f_{z, \alpha} \in R$。令 $I \subset R$ 为由元素
$f_{z, \alpha}$ 生成的理想，其中 $z$ 遍历 $M$ 的所有元素。
由构造，$I$ 在 $R'$ 中的像为零，而各元素 $\tilde z$ 在
$F_z/IF_z$ 中、从而在 $N/IN$ 中的像为零。因此在这种情形下，
$Z = \Spec(R/I)$ 是所求解答。

\medskip\noindent
假设 $\mathcal{F}$ 有限型，并且其概形论支集拟紧。写
$Z = \Spec(R/I)$。把 $I = \bigcup I_\lambda$ 写成有限生成理想的
滤过并。置 $Z_\lambda = \Spec(R/I_\lambda)$，于是
% P09-ERR-0836
$Z = \colim Z_\lambda$。由于 $u_Z$ 为零映射，由引理
\ref{spaces-flat-lemma-iso-limits} 可知，$u_{Z_\lambda}$ 对某个 $\lambda$
为零映射。这就完成了引理的证明。
\end{proof}

\begin{lemma}
\label{spaces-flat-lemma-F-zero-module-map}
设 $A$ 为环，设 $u : M \to N$ 为 $A$-模之间的映射。
如果 $N$ 作为 $A$-模是投射的，那么存在理想 $I \subset A$，
使得对任意环同态 $\varphi : A \to B$，下列条件等价：
\begin{enumerate}
\item $u \otimes 1 : M \otimes_A B \to N \otimes_A B$ 为零映射；
\item $\varphi(I) = 0$。
\end{enumerate}
\end{lemma}

\begin{proof}
由于 $N$ 投射，可以找到投射 $A$-模 $C$，使得
% P09-ERR-0837
$F = N \oplus C$ 是自由 $R$-模。把 $u$ 替换为
% P09-ERR-0838
$u \oplus 1 : F = M \oplus C \to N \oplus C$，
可知可以假设 $N$ 自由。在此情形下，令 $I$ 为 $A$ 的理想：
它由 $\Im(u)$ 中所有元素相对于 $N$ 的某组（固定）基的系数生成。
\end{proof}

\begin{lemma}
\label{spaces-flat-lemma-F-zero-somewhat-closed-points}
在情形 \ref{spaces-flat-situation-somewhat-closed} 中，设 $T \subset S$ 为子集，
设 $s \in S$ 位于 $T$ 的闭包中。对 $t \in T$，令 $u_t$ 为 $u$
在 $X_t$ 上的拉回，并令 $u_s$ 为 $u$ 在 $X_s$ 上的拉回。
如果 $X$ 在 $S$ 上局部有限表示、$\mathcal{G}$ 有限表示\footnote{
只要求 $X$ 在 $S$ 上局部有限型且 $\mathcal{G}$ 相对于 $S$
有限表示便已足够，但在代数空间框架下尚未定义这一概念。
关于概形的定义见《态射进阶》第
\ref{more-morphisms-section-finite-type-finite-presentation} 节。}，
并且都有 $u_t = 0$（对所有 $t \in T$），那么 $u_s = 0$。
\end{lemma}

\begin{proof}
% P09-ERR-0839
检验 $u_s$ 是否为零在纤维 $X_s$ 上是平展局部的。因此可以选取
一点 $x \in |X_s| \subset |X|$，并在一个平展邻域中检验。选取
$$
\xymatrix{
(X, x) \ar[d] & (X', x') \ar[l]^g \ar[d] \\
(S, s) & (S', s') \ar[l]
}
$$
如命题 \ref{spaces-flat-proposition-finite-presentation-flat-at-point} 中所示。
令 $T' \subset S'$ 为 $T$ 的逆像。注意，$s'$ 位于 $T'$ 的闭包中，
因为 $S' \to S$ 是开映射。因此归约到下一段所述的代数问题。

\medskip\noindent
有一个 $R$-模映射 $u : M \to N$，其中 $N$ 作为 $R$-模投射，
并且映射
$u_t : M \otimes_R \kappa(t) \to N \otimes_R \kappa(t)$
对每个 $t \in T$ 都为零。问题是证明 $u_s = 0$。令 $I \subset R$
为引理 \ref{spaces-flat-lemma-F-zero-module-map} 中定义的理想。那么 $I$ 在
$\kappa(t)$ 中的像对所有 $t \in T$ 都为零。因此 $T \subset V(I)$。
由于 $s$ 位于 $T$ 的闭包中，可知 $s \in V(I)$，故 $u_s = 0$。
\end{proof}

\noindent
若能为引理 \ref{spaces-flat-lemma-F-zero-closed-pure} 或引理
\ref{spaces-flat-lemma-F-zero-closed-proper} 找到一个“简单”的直接证明，
并使用引理 \ref{spaces-flat-lemma-F-zero-somewhat-closed} 和
\ref{spaces-flat-lemma-F-zero-somewhat-closed-points} 中所用的论证，
将会很有意思。当 $f : X \to B$ 是投射态射且 $B$ 是诺特概形时，
此引理的一个“经典”证明如下：(a) 选取相对丰沛可逆层
$\mathcal{O}_X(1)$；(b) 置
$u_n : f_*\mathcal{F}(n) \to f_*\mathcal{G}(n)$；
(c) 注意，$f_*\mathcal{G}(n)$ 对所有 $n \gg 0$ 都是有限局部自由层；
(d) $F_{zero}$ 由 $u_n$ 的零化轨迹表示，其中某个 $n \gg 0$。

\begin{lemma}
\label{spaces-flat-lemma-F-zero-closed-pure}
在情形 \ref{spaces-flat-situation-iso} 中，假设：
\begin{enumerate}
\item $f$ 有限表示；
\item $\mathcal{G}$ 有限表示、在 $B$ 上平坦且相对于 $B$ 纯。
\end{enumerate}
那么 $F_{zero}$ 是代数空间，并且 $F_{zero} \to B$ 是闭浸入。
如果 $\mathcal{F}$ 有限型，那么 $F_{zero} \to B$ 有限表示。
\end{lemma}

\begin{proof}
由引理 \ref{spaces-flat-lemma-finite-type-flat-pure-is-universal}，模
$\mathcal{G}$ 相对于 $B$ 普遍纯。为证明 $F_{zero}$ 是代数空间，
只须说明 $F_{zero} \to B$ 可表；见《空间》引理
\ref{spaces-lemma-representable-over-space}。设 $B' \to B$ 为态射，
其中 $B'$ 是概形；设 $u' : \mathcal{F}' \to \mathcal{G}'$ 为 $u$
在 $X' = X_{B'}$ 上的拉回。那么相应函子 $F'_{zero}$ 等于
$F_{zero} \times_B B'$。这把问题归约到 $B$ 是概形的情形。

\medskip\noindent
假设 $B$ 是概形。我们将说明 $F_{zero}$ 可由 $B$ 的闭子概形表示。
由引理 \ref{spaces-flat-lemma-iso-sheaf} 以及《下降》引理
\ref{descent-lemma-closed-immersion} 和
\ref{descent-lemma-descent-data-sheaves}，这个问题对于 $B$ 上的
平展拓扑是局部的。设 $b \in B$。先把 $B$ 替换为 $b$ 的一个仿射邻域。
选取图表
$$
\xymatrix{
X \ar[d] & X' \ar[l]^g \ar[d] \\
B & B' \ar[l]
}
$$
以及 $b' \in B'$，它映到 $b \in B$，如引理
\ref{spaces-flat-lemma-finite-presentation-flat-along-fibre} 中所示。
由于我们在平展局部工作，可以把 $B$ 替换为 $B'$，并假设有图表
$$
\xymatrix{
X \ar[rd] & & X' \ar[ll]^g \ar[ld] \\
& B
}
$$
其中 $B$ 和 $X'$ 都仿射，$\Gamma(X', g^*\mathcal{G})$ 是投射
$\Gamma(B, \mathcal{O}_B)$-模，并且
$g(|X'|) \supset |X_b|$。令 $U \subset X$ 为满足
$|U| = g(|X'|)$ 的开子空间。由《空间上的除子》引理
\ref{spaces-divisors-lemma-relative-assassin-constructible}，集合
$$
E = \{t \in B : \text{Ass}_{X_t}(\mathcal{G}_t) \subset |U_t|\} =
\{t \in B : \text{Ass}_{X/B}(\mathcal{G}) \cap |X_t| \subset |U_t|\}
$$
在 $B$ 中可构。由引理 \ref{spaces-flat-lemma-pure-on-top} 的 (2)，可知 $E$
包含 $\Spec(\mathcal{O}_{B, b})$。再由《态射》引理
\ref{morphisms-lemma-constructible-containing-open}，可知 $E$ 包含
$b$ 的一个开邻域。因此把 $B$ 替换为 $b$ 的一个更小仿射邻域后，
可以假设 $\text{Ass}_{X/B}(\mathcal{G}) \subset g(|X'|)$。

\medskip\noindent
由引理 \ref{spaces-flat-lemma-universally-separating} 可知，
% P09-ERR-0840
$u : \mathcal{F} \to \mathcal{G}$ 单射，当且仅当
$g^*u : g^*\mathcal{F} \to g^*\mathcal{G}$ 单射；而且任意基变换后
仍然如此。因此已归约到如下情形：除了
% P09-ERR-0841
定理中的假设外，$X \to B$ 还是仿射概形之间的态射，并且
$\Gamma(X, \mathcal{G})$ 是投射
$\Gamma(B, \mathcal{O}_B)$-模。此情形立即由引理
\ref{spaces-flat-lemma-F-zero-module-map} 得出。

\medskip\noindent
还须说明：$F_{zero} \to B$ 在 $\mathcal{F}$ 有限型时有限表示。
这由引理 \ref{spaces-flat-lemma-iso-limits} 与《空间的极限》命题
\ref{spaces-limits-proposition-characterize-locally-finite-presentation}
合并得出。
\end{proof}

\begin{lemma}
\label{spaces-flat-lemma-F-zero-closed-proper}
在情形 \ref{spaces-flat-situation-iso} 中，假设：
\begin{enumerate}
\item $f$ 局部有限表示；
\item $\mathcal{G}$ 是有限表示的 $\mathcal{O}_X$-模，并且在
$B$ 上平坦；
\item $\mathcal{G}$ 的支集在 $B$ 上是固有的。
\end{enumerate}
那么函子 $F_{zero}$ 是代数空间，并且 $F_{zero} \to B$ 是闭浸入。
如果 $\mathcal{F}$ 有限型，那么 $F_{zero} \to B$ 有限表示。
\end{lemma}

\begin{proof}
如果 $f$ 有限表示，那么结论立即由引理
\ref{spaces-flat-lemma-F-zero-closed-pure} 和 \ref{spaces-flat-lemma-proper-pure} 得出。
这是唯一有意义的情形；我们建议读者跳过证明的其余部分。其余部分
处理如下可能性（本引理的假设允许这一可能性）：$f$ 不拟分离或不拟紧。

\medskip\noindent
设 $i : Z \to X$ 为由 $\mathcal{G}$ 的第零 Fitting 理想截出的
闭子空间（《空间上的除子》第
\ref{spaces-divisors-section-fitting-ideals} 节）。由假设，$Z \to B$
是固有的（见《空间的导出范畴》第
\ref{spaces-perfect-section-proper-over-base} 节）。另一方面，$i$
有限表示（《空间上的除子》引理
\ref{spaces-divisors-lemma-fitting-ideal-of-finitely-presented} 以及
《空间的态射》引理
\ref{spaces-morphisms-lemma-closed-immersion-finite-presentation}）。
存在有限型拟凝聚 $\mathcal{O}_Z$-模 $\mathcal{H}$，使得
$i_*\mathcal{H} = \mathcal{G}$（《空间上的除子》引理
\ref{spaces-divisors-lemma-on-subscheme-cut-out-by-Fit-0}）。事实上，
由《代数》引理 \ref{algebra-lemma-finitely-presented-over-subring}，
$\mathcal{H}$ 作为 $\mathcal{O}_Z$-模有限表示（略去细节）。
于是 $F_{zero}$ 与函子 $F_{zero}$ 相同；后一个函子对应于映射
$i^*\mathcal{F} \to \mathcal{H}$，该映射与 $u$ 伴随。见引理
\ref{spaces-flat-lemma-relate-zero-affine}。层 $\mathcal{H}$ 相对于 $B$ 平坦，
因为 $\mathcal{G}$ 具有同一性质（在茎上检验；略去细节）。此外，
注意，如果 $\mathcal{F}$ 有限型，那么 $i^*\mathcal{F}$ 有限型。
因此已把本引理归约到证明第一段讨论的情形。
\end{proof}









\section{映射的平坦化}
\label{spaces-flat-section-flattening-map}

\noindent
本节是《平坦性进阶》第 \ref{flat-section-flattening-map} 节的对应版本。
特别地，下述结果是《平坦性进阶》定理
\ref{flat-theorem-flattening-map} 的一个变体。

\begin{theorem}
\label{spaces-flat-theorem-flattening-map}
在情形 \ref{spaces-flat-situation-iso} 中，假设：
\begin{enumerate}
\item $f$ 有限表示；
\item $\mathcal{F}$ 有限表示、在 $B$ 上平坦且相对于 $B$ 纯；
\item $u$ 满射。
\end{enumerate}
那么 $F_{iso}$ 可由闭浸入 $Z \to B$ 表示。此外，
% P09-ERR-0842
$Z \to S$ 在 $\mathcal{G}$ 有限表示时有限表示。
\end{theorem}

\begin{proof}
令 $\mathcal{K} = \Ker(u)$，并以
$v : \mathcal{K} \to \mathcal{F}$ 表示包含映射。由引理
\ref{spaces-flat-lemma-relate-zero-iso}，有 $F_{u, iso} = F_{v, zero}$。
把引理 \ref{spaces-flat-lemma-F-zero-closed-pure} 应用于 $v$，可知
$F_{u, iso} = F_{v, zero}$ 可由 $B$ 的闭子空间表示。注意，
$\mathcal{K}$ 在 $\mathcal{G}$ 有限表示时有限型；见《位形上的模》引理
\ref{sites-modules-lemma-kernel-surjection-finite-onto-finite-presentation}。
% P09-ERR-0843
因此也得到引理的最后一个断言。
\end{proof}

\begin{lemma}
\label{spaces-flat-lemma-F-iso-closed}
在情形 \ref{spaces-flat-situation-iso} 中，假设：
\begin{enumerate}
\item $f$ 局部有限表示；
\item $\mathcal{F}$ 局部有限表示且在 $B$ 上平坦；
\item $\mathcal{F}$ 的支集在 $B$ 上固有；
\item $u$ 满射。
\end{enumerate}
那么函子 $F_{iso}$ 是代数空间，并且 $F_{iso} \to B$ 是闭浸入。
如果 $\mathcal{G}$ 有限表示，那么 $F_{iso} \to B$ 有限表示。
\end{lemma}

\begin{proof}
令 $\mathcal{K} = \Ker(u)$，并以
$v : \mathcal{K} \to \mathcal{F}$ 表示包含映射。由引理
\ref{spaces-flat-lemma-relate-zero-iso}，有 $F_{u, iso} = F_{v, zero}$。
把引理 \ref{spaces-flat-lemma-F-zero-closed-proper} 应用于 $v$，可知
$F_{u, iso} = F_{v, zero}$ 可由 $B$ 的闭子空间表示。注意，
$\mathcal{K}$ 在 $\mathcal{G}$ 有限表示时有限型；见《位形上的模》引理
\ref{sites-modules-lemma-kernel-surjection-finite-onto-finite-presentation}。
因此也得到引理的最后一个断言。
\end{proof}

\noindent
讨论 Quot 函子时，我们将使用下述（容易的）结果。

\begin{lemma}
\label{spaces-flat-lemma-F-surj-open}
在情形 \ref{spaces-flat-situation-iso} 中，假设：
\begin{enumerate}
\item $f$ 局部有限表示；
\item $\mathcal{G}$ 有限型；
\item $\mathcal{G}$ 的支集在 $B$ 上固有。
\end{enumerate}
那么 $F_{surj}$ 是代数空间，并且 $F_{surj} \to B$ 是开浸入。
\end{lemma}

\begin{proof}
考虑 $\Coker(u)$。注意，$\Coker(u_T) = \Coker(u)_T$ 对任意
$T/B$ 成立。有限型拟凝聚模的支集的形成与拉回交换
（《空间的态射》引理 \ref{spaces-morphisms-lemma-support-covering}）。
因此 $F_{surj}$ 可由 $B$ 的如下开子空间表示：它对应于开集
$$
|B| \setminus |f|(\text{Supp}(\Coker(u)))
$$
见《空间的性质》引理 \ref{spaces-properties-lemma-open-subspaces}。
% P09-ERR-0844
这是开集，因为 $|f|$ 在 $\text{Supp}(\mathcal{G})$ 上是闭映射，
而 $\text{Supp}(\Coker(u))$ 是 $\text{Supp}(\mathcal{G})$ 的闭子集。
\end{proof}






\section{局部情形下的平坦化}
\label{spaces-flat-section-flattening-local}

\noindent
本节是《平坦性进阶》第 \ref{flat-section-flattening-local} 节的对应版本。

\begin{lemma}
\label{spaces-flat-lemma-freebie}
设 $S$ 是亨泽尔局部环的谱，其闭点为 $s$。设 $X \to S$ 为
局部有限型的代数空间态射。设 $\mathcal{F}$ 为有限型拟凝聚
$\mathcal{O}_X$-模。设 $E \subset |X_s|$ 为子集。存在闭子概形
$Z \subset S$，具有下述性质：对任意带点概形的态射
$(T, t) \to (S, s)$，下列条件等价：
\begin{enumerate}
\item $\mathcal{F}_T$ 在 $T$ 上平坦，所指的是在 $|X_t|$ 中
所有映到 $E \subset |X_s|$ 中一点的点处平坦；
\item $\Spec(\mathcal{O}_{T, t}) \to S$ 分解经过 $Z$。
\end{enumerate}
此外，如果 $X \to S$ 局部有限表示、$\mathcal{F}$ 有限表示，
并且 $E \subset |X_s|$ 闭且拟紧，那么 $Z \to S$ 有限表示。
\end{lemma}

\begin{proof}
选取概形 $U$ 和平展态射 $\varphi : U \to X$。令
$E' \subset |U_s|$ 为 $E$ 的逆像。如果 $E' \to E$ 满射，那么
条件 (1) 等价于：$(\varphi^*\mathcal{F})_T$ 在 $T$ 上平坦，
所指的是在 $|U_t|$ 中所有映到
% P09-ERR-0845
$E' \subset |U_t|$ 中一点的点处平坦。选取 $\varphi$ 为满射，
% P09-ERR-0846
我们便归约到概形的情形，即《平坦性进阶》引理
\ref{flat-lemma-freebie}。如果 $E$ 闭且拟紧，那么可以选取 $U$ 仿射，
使得 $E' \to E$ 满射。此时 $E'$ 闭且拟紧，最后一个断言由
《平坦性进阶》引理 \ref{flat-lemma-freebie} 的最后一个断言得出。
\end{proof}









\section{普遍平坦化}
\label{spaces-flat-section-flattening-final}

\noindent
本节是《平坦性进阶》第 \ref{flat-section-flattening-final} 节的对应版本。
我们的主要目标是证明引理 \ref{spaces-flat-lemma-when-universal-flattening}。
不过，我们没有看到直接从概形的对应结果推出这一结果的方法，
因此必须在这里重新建立部分材料。但先给出一个定义。

\begin{definition}
\label{spaces-flat-definition-flattening}
设 $S$ 为概形。设 $X \to Y$ 为 $S$ 上代数空间之间的态射。
设 $\mathcal{F}$ 为拟凝聚 $\mathcal{O}_X$-模。我们称
{\it $\mathcal{F}$ 的普遍平坦化存在}，如果情形
\ref{spaces-flat-situation-flat} 中定义的函子 $F_{flat}$ 是代数空间。
我们称{\it $X$ 的普遍平坦化存在}，如果
$\mathcal{O}_X$ 的普遍平坦化存在。
\end{definition}

\noindent
这有些不令人满意，因为我们不知道代数空间的每个单态射是否可表
（《空间的态射进阶》第
\ref{spaces-more-morphisms-section-monomorphisms} 节），所以此处的
普遍平坦化定义与概形情形所用的定义并不一致。希望这永远不会引起混淆。

\begin{lemma}
\label{spaces-flat-lemma-pre-flat-dimension-n}
设 $S$ 为概形。设 $f : X \to Y$ 为局部有限型的代数空间态射。
设 $\mathcal{F}$ 为有限型拟凝聚 $\mathcal{O}_X$-模，设
$n \geq 0$。下列条件等价：
\begin{enumerate}
\item 对某个交换图
$$
\xymatrix{
U \ar[d]_\varphi \ar[r] & V \ar[d] \\
X \ar[r] & Y
}
$$
其竖直箭头满且平展，$U$ 和 $V$ 为概形，并且层
$\varphi^*\mathcal{F}$ 在 $V$ 上于维数 $\geq n$ 处平坦
（《平坦性进阶》定义 \ref{flat-definition-flat-dimension-n}）；
\item 对每个交换图
$$
\xymatrix{
U \ar[d]_\varphi \ar[r] & V \ar[d] \\
X \ar[r] & Y
}
$$
其竖直箭头平展，$U$ 和 $V$ 为概形，并且层
$\varphi^*\mathcal{F}$ 在 $V$ 上于维数 $\geq n$ 处平坦；
\item 对任意 $x \in |X|$，若 $\mathcal{F}$ 在 $x$ 处不在 $Y$ 上
平坦，则 $x/f(x)$ 的超越次数 $< n$（《空间的态射》定义
\ref{spaces-morphisms-definition-dimension-fibre}）。
\end{enumerate}
如果这些条件成立，那么任意基变换 $Y' \to Y$ 后仍成立。
\end{lemma}

\begin{proof}
假设有 (1) 中的图。由《平坦性进阶》引理
\ref{flat-lemma-pre-flat-dimension-n} 中各条件的等价性，可知
(1) 与 (3) 等价。而条件 (3) 由 $\varphi^*\mathcal{F}$ 继承；
这对 (2) 中任意 $U \to V$ 都适用。因此再次使用概形情形的结果，
可知 (3) 蕴含 (2)。概形情形的结果也蕴含关于基变换的断言。
\end{proof}

\begin{definition}
\label{spaces-flat-definition-flat-dimension-n}
设 $S$ 为概形。设 $f : X \to Y$ 为 $S$ 上局部有限型的
代数空间态射。设 $\mathcal{F}$ 为有限型拟凝聚
$\mathcal{O}_X$-模。设 $n \geq 0$。如果引理
\ref{spaces-flat-lemma-pre-flat-dimension-n} 中的等价条件成立，则称
{\it $\mathcal{F}$ 在 $Y$ 上于维数 $\geq n$ 处平坦}。
\end{definition}

\begin{situation}
\label{spaces-flat-situation-flat-dimension-n}
设 $S$ 为概形。设 $f : X \to Y$ 为 $S$ 上局部有限型的
代数空间态射。设 $\mathcal{F}$ 为有限型拟凝聚
$\mathcal{O}_X$-模。对任意概形 $T$（在 $Y$ 上），以
$\mathcal{F}_T$ 表示 $\mathcal{F}$ 到 $T$ 的基变换；换言之，
$\mathcal{F}_T$ 是 $\mathcal{F}$ 经投影态射
$X_T = X \times_Y T \to X$ 的拉回。注意，
% P09-ERR-0847
$f_T : X_T \to T$ 有限型，并且 $\mathcal{F}_T$ 是有限型
$\mathcal{O}_{X_T}$-模（《空间的态射》引理
\ref{spaces-morphisms-lemma-base-change-finite-type} 和
《位形上的模》引理 \ref{sites-modules-lemma-local-pullback}）。
设 $n \geq 0$。由定义 \ref{spaces-flat-definition-flat-dimension-n} 和引理
\ref{spaces-flat-lemma-pre-flat-dimension-n}，得到函子
\begin{equation}
\label{spaces-flat-equation-flat-dimension-n}
F_n : (\Sch/Y)^{opp} \longrightarrow \textit{Sets}, \quad
T \longrightarrow \left\{
\begin{matrix}
\{*\} & \text{如果 }\mathcal{F}_T\text{ 对 }T\text{ 而言在 }\dim \geq n, \\
\emptyset & \text{否则。}
\end{matrix}
\right.
\end{equation}
\end{situation}

\noindent
在情形 \ref{spaces-flat-situation-flat-dimension-n} 中，有时把 $F_n$ 看作函子
$(\Sch/S)^{opp} \to \textit{Sets}$，并赋有态射 $F_n \to Y$。
具体地，如果 $T$ 是 $S$ 上概形，那么元素 $h \in F_n(T)$ 是态射
$h : T \to Y$，使得 $\mathcal{F}$ 经 $h$ 的基变换在 $T$ 上于
$\dim \geq n$ 处平坦。特别地，当我们说 $F_n$ 是代数空间时，
是指相应函子 $(\Sch/S)^{opp} \to \textit{Sets}$ 是代数空间。

\begin{lemma}
\label{spaces-flat-lemma-flat-dimension-n}
在情形 \ref{spaces-flat-situation-flat-dimension-n} 中：
\begin{enumerate}
\item 函子 $F_n$ 满足 fpqc 拓扑的层性质。
\item 如果 $f$ 拟紧且局部有限表示，并且 $\mathcal{F}$ 有限表示，
那么函子 $F_n$ 保极限。
\end{enumerate}
\end{lemma}

\begin{proof}
证明 (1)。假设 $\{T_i \to T\}$ 是概形 $T$（在 $Y$ 上）的 fpqc 覆盖。
须说明：如果 $F_n(T_i)$ 对所有 $i$ 都非空，那么 $F_n(T)$ 非空。
选取引理 \ref{spaces-flat-lemma-pre-flat-dimension-n} 的 (1) 中的图。以 $F'_n$
表示 $\varphi^*\mathcal{F}$ 和态射 $U \to V$ 所对应的函子。
由《平坦性进阶》引理 \ref{flat-lemma-flat-dimension-n}，$F'_n$
具有层性质。因此 $F_n$ 具有层性质，因为对 $T \to Y$，由引理
\ref{spaces-flat-lemma-pre-flat-dimension-n} 有
$F_n(T) = F'_n(V \times_Y T)$，并且
$\{V \times_Y T_i \to V \times_Y T\}$ 是 fpqc 覆盖。

\medskip\noindent
证明 (2)。假设 $T = \lim_{i \in I} T_i$ 是仿射概形
$T_i$（在 $Y$ 上）的滤过极限，并假设 $F_n(T)$ 非空。须说明
$F_n(T_i)$ 对某个 $i$ 非空。选取引理 \ref{spaces-flat-lemma-pre-flat-dimension-n} 的 (1)
中的图。固定 $i \in I$，并选取仿射开集
$W_i \subset V \times_Y T_i$，它满射到 $T_i$。对 $i' \geq i$，
令 $W_{i'}$ 为 $W_i$ 在 $V \times_Y T_{i'}$ 中的逆像；令
$W \subset V \times_Y T$ 为 $W_i$ 的逆像。于是
% P09-ERR-0848
$W = \lim_{i' \geq i} W_i$ 是 $V$ 上仿射概形的滤过极限。
再次由引理 \ref{spaces-flat-lemma-pre-flat-dimension-n}，只须说明
$F'_n(W_{i'})$ 对某个 $i' \geq i$ 非空。但 $F'_n(W)$ 非空，因为假设
$F_n(T) = F'_n(V \times_Y T)$ 非空。因此可以应用《平坦性进阶》引理
\ref{flat-lemma-flat-dimension-n} 得出结论。
\end{proof}

\begin{lemma}
\label{spaces-flat-lemma-localize-flat-dimension-n}
在情形 \ref{spaces-flat-situation-flat-dimension-n} 中，设
$h : X' \to X$ 为平展态射。置 $\mathcal{F}' = h^*\mathcal{F}$ 和
$f' = f \circ h$。令 $F_n'$ 为与
$(f' : X' \to Y, \mathcal{F}')$ 相联系的
(\ref{spaces-flat-equation-flat-dimension-n})。那么 $F_n$ 是 $F_n'$ 的子函子；
如果 $h(X') \supset \text{Ass}_{X/Y}(\mathcal{F})$，那么
$F_n = F'_n$。
\end{lemma}

\begin{proof}
选取引理 \ref{spaces-flat-lemma-pre-flat-dimension-n} 的 (1) 中的
$U \to X$、$V \to Y$、$U \to V$。选取满平展态射
$U' \to U \times_X X'$，其中 $U'$ 是概形。于是，对于由
两个函子 $F_{U, n}$ 与 $F_{U', n}$，本引理成立；这两个函子由
$U' \to U$ 和 $\mathcal{F}|_U$（在 $V$ 上）确定；见《平坦性进阶》引理
\ref{flat-lemma-localize-flat-dimension-n}。另一方面，引理
\ref{spaces-flat-lemma-pre-flat-dimension-n} 告诉我们：给定 $T \to Y$，有
$F_n(T) = F_{U, n}(V \times_Y T)$ 以及
$F'_n(T) = F_{U', n}(V \times_Y T)$。这证明了引理。
\end{proof}

\begin{theorem}
\label{spaces-flat-theorem-flat-dimension-n-representable}
在情形 \ref{spaces-flat-situation-flat-dimension-n} 中。进一步假设 $f$ 有限表示，
$\mathcal{F}$ 是有限表示的 $\mathcal{O}_X$-模，并且
$\mathcal{F}$ 相对于 $Y$ 纯。那么 $F_n$ 是代数空间，并且
$F_n \to Y$ 是有限表示的单态射。
\end{theorem}

\begin{proof}
由引理 \ref{spaces-flat-lemma-flat-dimension-n}，函子 $F_n$ 是 fppf 拓扑的层。
因为 $F_n \to Y$ 是位点 $(\Sch/S)_{fppf}$ 上层的单态射，所以
$\Delta : F_n \to F_n \times F_n$ 是对角态射
$\Delta_Y : Y \to Y \times_S Y$ 的拉回。因此，$\Delta_Y$ 的可表性
蕴含 $F_n$ 的同一性质。故只须证明：存在概形 $W$（在 $S$ 上）以及
满平展态射 $W \to F_n$。

\medskip\noindent
为构造 $W \to F_n$，选取平展覆盖 $\{Y_i \to Y\}$，其中 $Y_i$
是概形。令 $X_i = X \times_Y Y_i$，并令 $\mathcal{F}_i$ 为
$\mathcal{F}$ 到 $X_i$ 的拉回。于是，或者由定义，或者由引理
\ref{spaces-flat-lemma-quasi-finite-base-change}，$\mathcal{F}_i$ 相对于 $Y_i$ 纯。
定理的其他假设也都保持。最后，$F_n$ 在 $Y_i$ 上的限制就是函子
$F_n$；后者对应于 $X_i \to Y_i$ 和 $\mathcal{F}_i$。因此只须证明：
给定定理陈述中的 $\mathcal{F}$ 和 $f : X \to Y$，其中 $Y$ 是概形，
函子 $F_n$ 可由概形 $Z_n$ 表示，并且 $Z_n \to Y$ 是有限表示的
单态射。

\medskip\noindent
注意，有限表示的单态射是分离且拟有限的（《态射》引理
\ref{morphisms-lemma-monomorphism-loc-finite-type-loc-quasi-finite}）。
因此，结合《下降》引理 \ref{descent-lemma-descent-data-sheaves}、
《态射进阶》引理
\ref{more-morphisms-lemma-separated-locally-quasi-finite-morphisms-fppf-descend}
% P09-ERR-0849
以及《下降》引理 \ref{descent-lemma-descending-property-monomorphism} 和
\ref{descent-lemma-descending-property-finite-presentation}，可知这个问题
对于 $Y$ 上的平展拓扑是局部的。

\medskip\noindent
特别地，这一情形对于 $Y$ 上的 Zariski 拓扑是局部的，因而可以假设
$Y$ 仿射。此时 $f$ 的纤维维数有上界，故可知 $F_n$ 在 $n$ 足够大时
可表。因此可以对 $n$ 作降序归纳。假设已知 $F_{n + 1}$ 可由有限表示
的单态射 $Z_{n + 1} \to Y$ 表示。考虑基变换
$X_{n + 1} = Z_{n + 1} \times_Y X$，并令 $\mathcal{F}_{n + 1}$ 为
$\mathcal{F}$ 到 $X_{n + 1}$ 的拉回。态射 $Z_{n + 1} \to Y$ 是有限
表示的单态射，因而拟有限；故引理 \ref{spaces-flat-lemma-quasi-finite-base-change}
蕴含 $\mathcal{F}_{n + 1}$ 相对于 $Z_{n + 1}$ 纯。因为 $F_n$ 是
$F_{n + 1}$ 的子函子，所以为了对 $F_n$ 证明结果，只须对情形
$\mathcal{F}_{n + 1}/X_{n + 1}/Z_{n + 1}$ 所对应的函子证明结果。
这样便把问题归约为对 $F_n$ 证明结果，其中
% P09-ERR-0850
$Y_{n + 1} = Y$；换言之，可以假设
$\mathcal{F}$ 于维数 $\geq n + 1$ 处在 $Y$ 上平坦。

\medskip\noindent
固定 $n$，并假设 $\mathcal{F}$ 于维数 $\geq n + 1$ 处在仿射概形
$Y$ 上平坦。为完成证明，须说明 $F_n$ 可由有限表示的单态射
% P09-ERR-0851
$Z_n \to S$ 表示。由于这个问题对于 $Y$ 上的平展拓扑是局部的，
只须说明：对每个 $y \in Y$，都存在平展邻域
$(Y', y') \to (Y, y)$，使得基变换到 $Y'$ 后结果成立。因此，由引理
\ref{spaces-flat-lemma-existence-complete}，可以假设存在平展态射
$h_j : W_j \to X$，$j = 1, \ldots, m$，使得对每个 $j$，
$\mathcal{F}_j/W_j/Y$ 在 $y$ 上都有完全逐层剥离，其中
$\mathcal{F}_j$ 是 $\mathcal{F}$ 到 $W_j$ 的拉回，并且
$|X_y| \subset \bigcup h_j(W_j)$。由于 $h_j$ 平展，由引理
\ref{spaces-flat-lemma-pre-flat-dimension-n}，
% P09-ERR-0852
层 $\mathcal{F}_j$ 仍于维数 $\geq n + 1$ 处在 $Y$ 上平坦。
置 $W = \bigcup h_j(W_j)$，这是 $X$ 的拟紧开集。由于
$\mathcal{F}$ 沿 $X_y$ 纯，可知
% P09-ERR-0853
$$
E = \{t \in |Y| : \text{Ass}_{X_t}(\mathcal{F}_t) \subset W \}.
$$
包含 $y$ 的所有一般化。由《空间上的除子》引理
\ref{spaces-divisors-lemma-relative-assassin-constructible}，$E$ 是
$Y$ 的可构子集。我们已经看到
$\Spec(\mathcal{O}_{Y, y}) \subset E$。由《态射》引理
\ref{morphisms-lemma-constructible-containing-open}，可知 $E$ 包含
$y$ 的一个开邻域。因此缩小 $Y$ 后，可以假设
% P09-ERR-0854
$E = Y$。由引理 \ref{spaces-flat-lemma-localize-flat-dimension-n}，只须对函子
$F_n$ 证明本引理；该函子与 $X = \coprod W_j$ 和
$\mathcal{F} = \coprod \mathcal{F}_j$ 相联系。若 $F_{j, n}$ 表示 $W_j \to Y$ 和层
$\mathcal{F}_j$ 所对应的函子，则 $F_n = \prod F_{j, n}$。因此只须证明
每个 $F_{j, n}$ 都可由某个有限表示的单态射 $Z_{j, n} \to Y$ 表示，
因为此时
% P09-ERR-0855
$$
Z_n = Z_{1, n} \times_Y \ldots \times_Y Z_{m, n}
$$
这样便把定理归约到《平坦性进阶》引理
\ref{flat-lemma-flat-dimension-n-representable} 所处理的特殊情形。
\end{proof}

\noindent
于是终于得到所需结果。

\begin{lemma}
\label{spaces-flat-lemma-when-universal-flattening}
设 $S$ 为概形。设 $f : X \to Y$ 为 $S$ 上代数空间之间的态射。
设 $\mathcal{F}$ 为拟凝聚 $\mathcal{O}_X$-模。
\begin{enumerate}
\item 如果 $f$ 有限表示，$\mathcal{F}$ 是有限表示的
$\mathcal{O}_X$-模，并且 $\mathcal{F}$ 相对于 $Y$ 纯，那么存在
普遍平坦化 $Y' \to Y$，其对象是 $\mathcal{F}$。此外，$Y' \to Y$ 是有限表示
的单态射。
\item 如果 $f$ 有限表示，并且 $X$ 相对于 $Y$ 纯，那么存在普遍平坦化
$Y' \to Y$，其对象是 $X$。此外，$Y' \to Y$ 是有限表示的单态射。
\item 如果 $f$ 固有且有限表示，并且 $\mathcal{F}$ 是有限表示的
$\mathcal{O}_X$-模，那么存在普遍平坦化 $Y' \to Y$，其对象是
$\mathcal{F}$。此外，$Y' \to Y$ 是有限表示的单态射。
\item 如果 $f$ 固有且有限表示，那么存在普遍平坦化 $Y' \to Y$，
其对象是 $X$。
\end{enumerate}
\end{lemma}

\begin{proof}
把定理 \ref{spaces-flat-theorem-flat-dimension-n-representable} 应用于
$F_0 = F_{flat}$，再利用 $f$ 固有时 $\mathcal{F}$ 自动相对于基纯这一
事实，便立即得到这些陈述；见引理 \ref{spaces-flat-lemma-proper-pure}。
\end{proof}






\section{Grothendieck 存在定理}
\label{spaces-flat-section-existence}

\noindent
本节是《平坦性进阶》第 \ref{flat-section-existence} 节的对应版本，
并继续《空间的态射进阶》第
\ref{spaces-more-morphisms-section-existence-proper} 节的讨论。
我们将在下述情形中工作。

\begin{situation}
\label{spaces-flat-situation-existence}
这里有一个环的逆系统 $(A_n)$，其转移映射满且核局部幂零。置
$A = \lim A_n$。有一个代数空间 $X$，它在 $A$ 上分离且有限表示。
置 $X_n = X \times_{\Spec(A)} \Spec(A_n)$，并把它视为 $X$ 的闭子空间。
进一步假设给定系统 $(\mathcal{F}_n, \varphi_n)$，其中
$\mathcal{F}_n$ 是有限表示的 $\mathcal{O}_{X_n}$-模，在 $A_n$ 上平坦，
其支撑在 $A_n$ 上固有，并且
$$
\varphi_n :
\mathcal{F}_n \otimes_{\mathcal{O}_{X_n}} \mathcal{O}_{X_{n - 1}}
\longrightarrow
\mathcal{F}_{n - 1}
$$
是同构（这里采用《空间的态射》引理
\ref{spaces-morphisms-lemma-i-star-equivalence} 中等价性所给出的记号）。
\end{situation}

\noindent
我们的目标是考察能否找到拟凝聚层 $\mathcal{F}$（在 $X$ 上），使得
$\mathcal{F}_n = \mathcal{F} \otimes_{\mathcal{O}_X} \mathcal{O}_{X_n}$
对所有 $n$ 都成立。

\begin{lemma}
\label{spaces-flat-lemma-compute-what-it-should-be}
在情形 \ref{spaces-flat-situation-existence} 中，考虑
$$
K = R\lim_{D_\QCoh(\mathcal{O}_X)}(\mathcal{F}_n) =
DQ_X(R\lim_{D(\mathcal{O}_X)}\mathcal{F}_n)
$$
那么 $K$ 属于 $D^b_{\QCoh}(\mathcal{O}_X)$；事实上，$K$ 只有次数
$\geq 0$ 的上同调层可能非零。
\end{lemma}

\begin{proof}
这是《空间的导出范畴》例
\ref{spaces-perfect-example-inverse-limit-quasi-coherent} 的特殊情形。
\end{proof}

\begin{lemma}
\label{spaces-flat-lemma-compute-against-perfect}
在情形 \ref{spaces-flat-situation-existence} 中，设 $K$ 如引理
\ref{spaces-flat-lemma-compute-what-it-should-be} 所示。对任意完美对象 $E$（属于
$D(\mathcal{O}_X)$），有：
\begin{enumerate}
\item $M = R\Gamma(X, K \otimes^\mathbf{L} E)$ 是 $D(A)$ 的完美对象，
并且有典范同构
$R\Gamma(X_n, \mathcal{F}_n \otimes^\mathbf{L} E|_{X_n}) =
M \otimes_A^\mathbf{L} A_n$，它位于 $D(A_n)$ 中；
\item $N = R\Hom_X(E, K)$ 是 $D(A)$ 的完美对象，并且有典范同构
$R\Hom_{X_n}(E|_{X_n}, \mathcal{F}_n) = N \otimes_A^\mathbf{L} A_n$，
它位于 $D(A_n)$ 中。
\end{enumerate}
两个陈述中的 $E|_{X_n}$ 都表示 $E$ 到 $X_n$ 的导出拉回。
\end{lemma}

\begin{proof}
证明 (2)。记 $E_n = E|_{X_n}$，并置
$N_n = R\Hom_{X_n}(E_n, \mathcal{F}_n)$。回顾
$R\Hom_{X_n}(-, -)$ 等于 $R\Gamma(X_n, R\SheafHom(-, -))$；见
《位形上的上同调》第 \ref{sites-cohomology-section-global-RHom} 节。
因此，由《空间的导出范畴》引理
\ref{spaces-perfect-lemma-base-change-RHom-perfect}，可知 $N_n$ 是
$D(A_n)$ 的完美对象，其构造与基变换交换。故映射
$N_n \otimes_{A_n}^\mathbf{L} A_{n - 1} \to N_{n - 1}$（由
$\varphi_n$ 得到）都是同构。
由《代数进阶》引理 \ref{more-algebra-lemma-Rlim-perfect-gives-perfect}，
可知 $R\lim N_n$ 完美，并且把它基变换回 $A_n$ 即恢复 $N_n$。
另一方面，三角范畴的正合函子
$R\Hom_X(E, -) : D_\QCoh(\mathcal{O}_X) \to D(A)$ 与积交换，因而也与
导出极限交换，所以
$$
R\Hom_X(E, K) =
R\lim R\Hom_X(E, \mathcal{F}_n) =
R\lim R\Hom_X(E_n, \mathcal{F}_n) =
R\lim N_n
$$
这证明了 (2)。利用《位形上的上同调》引理
\ref{sites-cohomology-lemma-dual-perfect-complex} 把 (1) 转化为 (2)，
便可看出 (1) 成立。
\end{proof}

\begin{lemma}
\label{spaces-flat-lemma-relative-pseudo-coherence}
在情形 \ref{spaces-flat-situation-existence} 中，设 $K$ 如引理
\ref{spaces-flat-lemma-compute-what-it-should-be} 所示。那么 $K$ 相对于 $A$ 伪凝聚。
\end{lemma}

\begin{proof}
% P09-ERR-0856
结合引理 \ref{spaces-flat-lemma-compute-against-perfect} 和《空间的导出范畴》引理
\ref{spaces-perfect-lemma-perfect-enough}，可知
$R\Gamma(X, K \otimes^\mathbf{L} E)$ 是 $D(A)$ 中的伪凝聚对象；
这对所有伪凝聚对象 $E$（属于 $D(\mathcal{O}_X)$）都成立。因此，
本引理由《空间的态射进阶》引理
\ref{spaces-more-morphisms-lemma-characterize-pseudo-coherent} 得出。
\end{proof}

\begin{lemma}
\label{spaces-flat-lemma-compute-over-affine}
在情形 \ref{spaces-flat-situation-existence} 中，设 $K$ 如引理
\ref{spaces-flat-lemma-compute-what-it-should-be} 所示。对任意平展态射 $U \to X$，
其中 $U$ 拟紧且拟分离，有
$$
R\Gamma(U, K) \otimes_A^\mathbf{L} A_n =
R\Gamma(U_n, \mathcal{F}_n)
$$
这是 $D(A_n)$ 中的等式，其中 $U_n = U \times_X X_n$。
\end{lemma}

\begin{proof}
固定 $n$。由《空间的导出范畴》引理
\ref{spaces-perfect-lemma-computing-sections-as-colim}，存在完美复形系统
$E_m$（在 $X$ 上），使得
$R\Gamma(U, K) = \text{hocolim} R\Gamma(X, K \otimes^\mathbf{L} E_m)$。
事实上，这个公式不只对 $K$ 成立，而且对
$D_\QCoh(\mathcal{O}_X)$ 的每个对象都成立。把它应用于
$\mathcal{F}_n$，得到
\begin{align*}
R\Gamma(U_n, \mathcal{F}_n)
& =
R\Gamma(U, \mathcal{F}_n) \\
& =
\text{hocolim}_m R\Gamma(X, \mathcal{F}_n \otimes^\mathbf{L} E_m) \\
& =
\text{hocolim}_m R\Gamma(X_n, \mathcal{F}_n \otimes^\mathbf{L} E_m|_{X_n})
\end{align*}
利用引理 \ref{spaces-flat-lemma-compute-against-perfect} 以及
$- \otimes_A^\mathbf{L} A_n$ 与同伦余极限交换这一事实，便得到结果。
\end{proof}

\begin{lemma}
\label{spaces-flat-lemma-finitely-presented}
在情形 \ref{spaces-flat-situation-existence} 中，设 $K$ 如引理
\ref{spaces-flat-lemma-compute-what-it-should-be} 所示。以 $X_0 \subset |X|$ 表示由
位于 $\Spec(A_1) = \Spec(A_2) = \ldots$ 这个 $\Spec(A)$ 的闭子集上方
的点组成的闭子集。存在开子空间 $W \subset X$，它包含 $X_0$，并使得
\begin{enumerate}
\item $H^i(K)|_W$ 为零，除非 $i = 0$；
\item $\mathcal{F} = H^0(K)|_W$ 有限表示；
\item $\mathcal{F}_n = \mathcal{F} \otimes_{\mathcal{O}_X} \mathcal{O}_{X_n}$。
\end{enumerate}
\end{lemma}

\begin{proof}
固定 $n \geq 1$。由构造，有典范映射 $K \to \mathcal{F}_n$（在
$D_\QCoh(\mathcal{O}_X)$ 中），从而有拟凝聚层的典范映射
$H^0(K) \to \mathcal{F}_n$。这说明了 (3) 的含义。

\medskip\noindent
设 $x \in X_0$ 为一点。我们将找出开邻域 $W$，它包含 $x$，并使 (1)、(2)、
(3) 成立。由于 $X_0$ 拟紧，这将证明引理。设 $U \to X$ 为平展态射，
其中 $U$ 仿射，且 $u \in U$ 映到 $x$。由于 $|U| \to |X|$ 是开映射，
只须找出 $u$ 的开邻域（在 $U$ 中），使 (1)、(2)、(3) 成立。记
$U = \Spec(B)$。选取满射 $P \to B$，其中 $P$ 在 $A$ 上光滑。
由引理 \ref{spaces-flat-lemma-relative-pseudo-coherence} 和相对伪凝聚性的定义，
存在上有界复形 $F^\bullet$，它由有限自由 $P$-模组成，并且表示
% P09-ERR-0857
$Ri_*K$；这里 $i : U \to \Spec(P)$ 是该表示所诱导的闭浸入。
令 $M_n$ 为 $B$-模；它对应于 $\mathcal{F}_n|_U$。由引理
\ref{spaces-flat-lemma-compute-over-affine}，
$$
H^i(F^\bullet \otimes_A A_n) =
\left\{
\begin{matrix}
0 & \text{如果} & i \not = 0 \\
M_n & \text{如果} & i = 0
\end{matrix}
\right.
$$
设 $i$ 为使 $F^i$ 非零的最大指标。如果 $i \leq 0$，那么 (1)、(2)、
(3) 成立。否则 $i > 0$，而且映射
$$
F^{i - 1} \to F^i
$$
% P09-ERR-0858
在点 $u$ 处的秩最大。因此，在 $u$ 的一个开邻域上（该邻域取在 $\Spec(P)$ 中），
其秩最大。故把 $P$ 替换为一个主局部化后，可以假设所显示的映射满。
由于 $F^i$ 有限自由，可以选取分裂 $F^{i - 1} = F' \oplus F^i$。
于是可以把 $F^\bullet$ 替换为复形
$$
\ldots \to F^{i - 2} \to F' \to 0 \to \ldots
$$
再对 $i$ 作归纳即得结论。
\end{proof}

\begin{lemma}
\label{spaces-flat-lemma-proper-support}
在情形 \ref{spaces-flat-situation-existence} 中，设 $K$ 如引理
\ref{spaces-flat-lemma-compute-what-it-should-be} 所示。设 $W \subset X$ 如引理
\ref{spaces-flat-lemma-finitely-presented} 所示。置 $\mathcal{F} = H^0(K)|_W$。
那么，必要时缩小开集 $W$ 后，$\mathcal{F}$ 的支撑在 $A$ 上固有。
\end{lemma}

\begin{proof}
固定 $n \geq 1$。令 $I_n = \Ker(A \to A_n)$。由《代数进阶》引理
\ref{more-algebra-lemma-limit-henselian}，对 $(A, I_n)$ 是 Hensel 的。
令 $Z \subset W$ 为 $\mathcal{F}$ 的概形论支撑。由于 $\mathcal{F}$
有限表示，这是闭子空间。由引理 \ref{spaces-flat-lemma-finitely-presented} 的 (3)，
可知 $Z \times_{\Spec(A)} \Spec(A_n)$ 等于 $\mathcal{F}_n$ 的支撑，
因而在
% P09-ERR-0859
$\Spec(A/I)$ 上固有。由《空间的态射进阶》引理
\ref{spaces-more-morphisms-lemma-split-off-proper-part-henselian}，可以写成
$Z = Z_1 \amalg Z_2$，其中 $Z_1, Z_2$ 在 $Z$ 中开且闭，$Z_1$ 在 $A$
上固有，并且 $Z_1 \times_{\Spec(A)} \Spec(A/I_n)$ 等于
$\mathcal{F}_n$ 的支撑。换言之，$|Z_2|$ 不与 $X_0$ 相交。因此把
$W$ 替换为 $W \setminus Z_2$ 后，便得到引理。
\end{proof}

\begin{theorem}[Grothendieck 存在定理]
\label{spaces-flat-theorem-existence}
在情形 \ref{spaces-flat-situation-existence} 中，存在有限表示的
$\mathcal{O}_X$-模 $\mathcal{F}$，它在 $A$ 上平坦，支撑在 $A$ 上固有，
并且有
$\mathcal{F}_n = \mathcal{F} \otimes_{\mathcal{O}_X} \mathcal{O}_{X_n}$；
对所有 $n$，这些等式均与映射 $\varphi_n$ 相容。
\end{theorem}

\begin{proof}
应用引理 \ref{spaces-flat-lemma-compute-what-it-should-be}、
\ref{spaces-flat-lemma-compute-against-perfect}、
\ref{spaces-flat-lemma-relative-pseudo-coherence}、
\ref{spaces-flat-lemma-compute-over-affine}、
\ref{spaces-flat-lemma-finitely-presented} 和
\ref{spaces-flat-lemma-proper-support}，得到开子空间 $W \subset X$，它包含所有位于
$\Spec(A_n)$ 上方的点，还得到有限表示的 $\mathcal{O}_W$-模
$\mathcal{F}$，其支撑在 $A$ 上固有，并且有
$\mathcal{F}_n = \mathcal{F} \otimes_{\mathcal{O}_W} \mathcal{O}_{X_n}$；
这个等式对所有 $n \geq 1$ 成立。
（这是有意义的，因为 $X_n \subset W$。）由引理
\ref{spaces-flat-lemma-proper-pure}，可知 $\mathcal{F}$ 相对于 $\Spec(A)$ 普遍纯。
由定理 \ref{spaces-flat-theorem-flat-dimension-n-representable}（解释见引理
\ref{spaces-flat-lemma-when-universal-flattening}），存在普遍平坦化
$S' \to \Spec(A)$，其所针对的模为 $\mathcal{F}$；此外，态射
$S' \to \Spec(A)$ 是有限表示的单态射。
特别地，$S'$ 是概形（这可由定理的证明得出；也可
% P09-ERR-0860
事后由《空间的态射》命题
\ref{spaces-morphisms-proposition-locally-quasi-finite-separated-over-scheme}
得出）。由于 $\mathcal{F}$ 到 $\Spec(A_n)$ 的基变换是
$\mathcal{F}_n$，可知态射 $\Spec(A_n) \to \Spec(A)$（唯一地）分解经过
$S'$；这对每个 $n$ 成立。由《平坦性进阶》引理
\ref{flat-lemma-monomorphism-isomorphism}，可知 $S' = \Spec(A)$。
这意味着 $\mathcal{F}$ 在 $A$ 上平坦。最后，令 $Z$ 为 $\mathcal{F}$ 的
概形论支撑。由于这一支撑在 $\Spec(A)$ 上固有，态射 $Z \to X$ 是闭的。因此，
推前 $(W \to X)_*\mathcal{F}$ 支撑在 $W$ 上，并具有全部所需性质。
\end{proof}







\section{Grothendieck 存在定理，再论}
\label{spaces-flat-section-existence-derived}

\noindent
% P09-ERR-0861
本节沿用第 \ref{spaces-flat-section-existence} 节中处理拟凝聚模的方法，
证明导出范畴中的 Grothendieck 存在定理之对应形式。本节是《平坦性进阶》
第 \ref{flat-section-existence-derived} 节对代数空间的对应形式。
经典情形（对代数空间而言）见《空间的态射进阶》第
\ref{spaces-more-morphisms-section-existence-proper} 节。我们将在下述情形中工作。

\begin{situation}
\label{spaces-flat-situation-existence-derived}
这里有环的逆系统 $(A_n)$，其过渡映射为满射，且各核局部幂零。
置 $A = \lim A_n$。我们有代数空间 $X$，它在 $A$ 上固有、平坦且有限表示。
置 $X_n = X \times_{\Spec(A)} \Spec(A_n)$，并把它看作 $X$ 的闭子空间。
再假设给定系统 $(K_n, \varphi_n)$，其中 $K_n$ 是
$D(\mathcal{O}_{X_n})$ 的伪凝聚对象，而
$$
\varphi_n : K_n \longrightarrow K_{n - 1}
$$
是 $D(\mathcal{O}_{X_n})$ 中的映射，它诱导同构
$K_n \otimes_{\mathcal{O}_{X_n}}^\mathbf{L} \mathcal{O}_{X_{n - 1}}
\to K_{n - 1}$；该同构位于 $D(\mathcal{O}_{X_{n - 1}})$ 中。
\end{situation}

\noindent
更精确地说，应当写成
$\varphi_n : K_n \to Ri_{n - 1, *}K_{n - 1}$，
其中 $i_{n - 1} : X_{n - 1} \to X_n$ 是包含态射；在这种记号下，
条件是伴随映射 $Li_{n - 1}^*K_n \to K_{n - 1}$ 为同构。
我们的目标是找出伪凝聚对象 $K \in D(\mathcal{O}_X)$，使得
$K_n = K \otimes_{\mathcal{O}_X}^\mathbf{L} \mathcal{O}_{X_n}$
对所有 $n$ 成立（沿用同样的记号滥用）。

\begin{lemma}
\label{spaces-flat-lemma-compute-what-it-should-be-derived}
在情形 \ref{spaces-flat-situation-existence-derived} 中，考虑
$$
K = R\lim_{D_\QCoh(\mathcal{O}_X)}(K_n) =
DQ_X(R\lim_{D(\mathcal{O}_X)} K_n)
$$
那么 $K$ 属于 $D^-_{\QCoh}(\mathcal{O}_X)$。
\end{lemma}

\begin{proof}
函子 $DQ_X$ 存在，因为 $X$ 拟紧且拟分离，见《空间的导出范畴》引理
\ref{spaces-perfect-lemma-better-coherator}。
由于 $DQ_X$ 是右伴随，它与积交换，因而也与导出极限交换。
这便给出引理陈述中的等式。

\medskip\noindent
由《空间的导出范畴》引理
\ref{spaces-perfect-lemma-boundedness-better-coherator}，函子 $DQ_X$
具有有界的上同调维数。因此，只须证明
$R\lim K_n \in D^-(\mathcal{O}_X)$。为此，设 $U \to X$ 为平展态射，
其中 $U$ 仿射。由《位形上的上同调》引理
\ref{sites-cohomology-lemma-RGamma-commutes-with-Rlim}，有典范正合序列
$$
0 \to
R^1\lim H^{m - 1}(U, K_n) \to H^m(U, R\lim K_n) \to
\lim H^m(U, K_n) \to 0
$$
由于 $U$ 仿射，且 $K_n$ 伪凝聚（从而其上同调层拟凝聚，见
《空间的导出范畴》引理 \ref{spaces-perfect-lemma-pseudo-coherent}），
由《概形的导出范畴》引理 \ref{perfect-lemma-affine-compare-bounded}，有
$H^m(U, K_n) = H^m(K_n)(U)$。因此，只须证明 $K_n$ 的上有界界限
与 $n$ 无关。

\medskip\noindent
由于 $K_n$ 伪凝聚，有 $K_n \in D^-(\mathcal{O}_{X_n})$。
% P09-ERR-0862
假设 $a_n$ 是使 $H^{a_n}(K_n)$ 非零的最大整数。当然，
$a_1 \leq a_2 \leq a_3 \leq \ldots$。注意，$H^{a_n}(K_n)$ 是有限表示的
$\mathcal{O}_{X_n}$-模（《位形上的上同调》引理
\ref{sites-cohomology-lemma-finite-cohomology}）。我们有
$H^{a_n}(K_{n - 1}) =
H^{a_n}(K_n) \otimes_{\mathcal{O}_{X_n}} \mathcal{O}_{X_{n - 1}}$。
由于 $X_{n - 1} \to X_n$ 是加厚，由 Nakayama 引理（《代数》引理
\ref{algebra-lemma-NAK}）可知，如果
$H^{a_n}(K_n) \otimes_{\mathcal{O}_{X_n}} \mathcal{O}_{X_{n - 1}}$
为零，那么 $H^{a_n}(K_n)$ 也为零（例如可在茎上检验；略去这一小细节）。
% P09-ERR-0863
所以 $a_{n - 1} = a_n$ 对所有 $n$ 都成立，结论随即成立。
\end{proof}

\begin{lemma}
\label{spaces-flat-lemma-compute-against-perfect-derived}
在情形 \ref{spaces-flat-situation-existence-derived} 中，设 $K$ 如引理
\ref{spaces-flat-lemma-compute-what-it-should-be-derived} 所示。对任意完美对象
$E$（属于 $D(\mathcal{O}_X)$），导出上同调对象
$$
M = R\Gamma(X, K \otimes^\mathbf{L} E)
$$
是 $D(A)$ 的伪凝聚对象，并且有典范同构
$$
R\Gamma(X_n, K_n \otimes^\mathbf{L} E|_{X_n}) = M \otimes_A^\mathbf{L} A_n
$$
这是 $D(A_n)$ 中的同构。这里 $E|_{X_n}$ 表示 $E$ 到 $X_n$ 的导出拉回。
\end{lemma}

\begin{proof}
记 $E_n = E|_{X_n}$，并记
$M_n = R\Gamma(X_n, K_n \otimes^\mathbf{L} E|_{X_n})$。
由《空间的导出范畴》引理
\ref{spaces-perfect-lemma-flat-proper-pseudo-coherent-direct-image-general}，
可知 $M_n$ 是 $D(A_n)$ 的伪凝聚对象，且其构造与基变换交换。因此，映射
$M_n \otimes_{A_n}^\mathbf{L} A_{n - 1} \to M_{n - 1}$
（由 $\varphi_n$ 得到）都是同构。
由《代数进阶》引理
\ref{more-algebra-lemma-Rlim-pseudo-coherent-gives-pseudo-coherent}，
可知 $R\lim M_n$ 伪凝聚，并且把它基变换回 $A_n$ 即恢复 $M_n$。
另一方面，三角范畴的正合函子
$R\Gamma(X, -) : D_\QCoh(\mathcal{O}_X) \to D(A)$ 与积交换，
因而也与导出极限交换，所以
$$
R\Gamma(X, E \otimes^\mathbf{L} K) =
R\lim R\Gamma(X,  E \otimes^\mathbf{L} K_n) =
R\lim R\Gamma(X_n, E_n \otimes^\mathbf{L} K_n) =
R\lim M_n
$$
这正是所需结论。
\end{proof}

\begin{lemma}
\label{spaces-flat-lemma-relative-pseudo-coherence-derived}
在情形 \ref{spaces-flat-situation-existence-derived} 中，设 $K$ 如引理
\ref{spaces-flat-lemma-compute-what-it-should-be-derived} 所示。那么 $K$ 在 $X$ 上伪凝聚。
\end{lemma}

\begin{proof}
% P09-ERR-0864
结合引理 \ref{spaces-flat-lemma-compute-against-perfect-derived} 和《空间的导出范畴》引理
\ref{spaces-perfect-lemma-perfect-enough}，可知
$R\Gamma(X, K \otimes^\mathbf{L} E)$ 是 $D(A)$ 中的伪凝聚对象；
这对所有伪凝聚对象 $E$（属于 $D(\mathcal{O}_X)$）均成立。
因此，由《空间的态射进阶》引理
\ref{spaces-more-morphisms-lemma-characterize-pseudo-coherent}，
可知 $K$ 相对于 $A$ 伪凝聚。
% P09-ERR-0865
由于 $X$ 在 $A$ 上平坦且有限表示，这等价于在 $X$ 上伪凝聚，见
《空间的态射进阶》引理
\ref{spaces-more-morphisms-lemma-relative-pseudo-coherent-is-moot}。
\end{proof}

\begin{lemma}
\label{spaces-flat-lemma-compute-over-affine-derived}
在情形 \ref{spaces-flat-situation-existence-derived} 中，设 $K$ 如引理
\ref{spaces-flat-lemma-compute-what-it-should-be-derived} 所示。对任意平展态射
$U \to X$，其中 $U$ 拟紧且拟分离，有
$$
R\Gamma(U, K) \otimes_A^\mathbf{L} A_n =
R\Gamma(U_n, K_n)
$$
这是 $D(A_n)$ 中的等式，其中 $U_n = U \times_X X_n$。
\end{lemma}

\begin{proof}
固定 $n$。由《空间的导出范畴》引理
\ref{spaces-perfect-lemma-computing-sections-as-colim}，存在完美复形系统
$E_m$（在 $X$ 上），使得
$R\Gamma(U, K) = \text{hocolim} R\Gamma(X, K \otimes^\mathbf{L} E_m)$。
事实上，这个公式不只对 $K$ 成立，而且对
$D_\QCoh(\mathcal{O}_X)$ 的每个对象都成立。把它应用于 $K_n$，得到
\begin{align*}
R\Gamma(U_n, K_n)
& =
R\Gamma(U, K_n) \\
& =
\text{hocolim}_m R\Gamma(X, K_n \otimes^\mathbf{L} E_m) \\
& =
\text{hocolim}_m R\Gamma(X_n, K_n \otimes^\mathbf{L} E_m|_{X_n})
\end{align*}
利用引理 \ref{spaces-flat-lemma-compute-against-perfect-derived} 以及
$- \otimes_A^\mathbf{L} A_n$ 与同伦余极限交换这一事实，便得到结果。
\end{proof}

\begin{theorem}[导出 Grothendieck 存在定理]
\label{spaces-flat-theorem-existence-derived}
在情形 \ref{spaces-flat-situation-existence-derived} 中，存在
伪凝聚对象 $K$（属于 $D(\mathcal{O}_X)$），使得
$K_n = K \otimes_{\mathcal{O}_X}^\mathbf{L} \mathcal{O}_{X_n}$；
对所有 $n$，这些等式均与映射 $\varphi_n$ 相容。
\end{theorem}

\begin{proof}
应用引理 \ref{spaces-flat-lemma-compute-what-it-should-be-derived}、
\ref{spaces-flat-lemma-compute-against-perfect-derived}、
\ref{spaces-flat-lemma-relative-pseudo-coherence-derived}，得到
伪凝聚对象 $K$（属于 $D(\mathcal{O}_X)$）。在引理
\ref{spaces-flat-lemma-compute-over-affine-derived} 中选取仿射的 $U$，立即可知
$K$ 的限制为 $K_n$（这里限制到 $X_n$）。
\end{proof}

\begin{remark}
\label{spaces-flat-remark-correct-generality}
本节的结果可以推广。很可能只须假设 $X \to \Spec(A)$ 分离且有限表示，
并假设 $K_n$ 相对于 $A_n$ 伪凝聚，且支撑在 $X_n$ 的某个在 $A_n$ 上
固有的闭子集上，结论便仍然正确。所得对象将是 $K$，它相对于 $A$
伪凝聚，并支撑在某个在 $A$ 上固有的闭子集上。若今后需要这一结果，
我们将在此给出精确陈述并加以证明。
\end{remark}
% Generated mechanically from the frozen spaces-groupoids.tex target.
% Do not edit this generated file; edit the bound locale source instead.

% OK, start here.
%

\setcounter{chapter}{77}
\chapter{代数空间中的群胚}



\phantomsection
\label{spaces-groupoids-section-phantom}


\section{引言}
\label{spaces-groupoids-section-introduction}

\noindent
本章专论代数空间中群胚的一般理论。我们推荐阅读 Keel 和 Mori 的优美论文
\cite{K-M}。

\medskip\noindent
这里的很多内容重复了群胚概形一章中已经讨论过的内容，见《群胚》第
\ref{groupoids-section-introduction} 节。这里关于商栈的讨论则是新的。


\section{约定}
\label{spaces-groupoids-section-conventions}

\noindent
我们始终假设所有概形都包含在一个大 fppf 位形 $\Sch_{fppf}$ 中，
并且所有考虑的环 $A$ 都具有如下性质：$\Spec(A)$（在同构意义下）
是这个大位形的对象。

\medskip\noindent
设 $S$ 为概形，$X$ 为 $S$ 上的代数空间。在本章及下一章中，我们用
$X \times_S X$ 表示 $X$ 与自身的积（取在 $S$ 上代数空间的范畴中），
而不用 $X \times X$。

\medskip\noindent
我们继续沿用从指标 $0$ 开始标记投影映射的约定，因而有
$\text{pr}_0 : X \times_S Y \to X$ 和
$\text{pr}_1 : X \times_S Y \to Y$。




\section{记号}
\label{spaces-groupoids-section-notation}

\noindent
设 $S$ 为概形；它将是我们的基概形，所有代数空间都取在 $S$ 上。
设 $B$ 为 $S$ 上的代数空间；它将是我们的基代数空间，其他代数空间以及
概形往往取在 $B$ 上。若称 $X$ 是 $B$ 上的代数空间，意思是 $X$ 是
$S$ 上的代数空间，并配备结构态射 $X \to B$。此外，我们尽量保留字母
$T$ 来表示 $B$ 上的“测试”概形。换言之，$T$ 是配备结构态射
$T \to B$ 的概形。在这种情形下，我们用 $X(T)$ 表示 $T$-值点所成的集合；
这些点取值于 $X$ 且位于 $B$ 上。用公式写就是
$$
X(T) = \Mor_B(T, X).
$$
类似地，给定另一个代数空间 $Y$（在 $B$ 上），置
$$
X(Y) = \Mor_B(Y, X).
$$
假设如上给定代数空间 $X$、$Y$（在 $B$ 上），以及态射
$f : X \to Y$（在 $B$ 上）。对任意概形 $T$（在 $B$ 上），得到诱导的集合映射
$$
f : X(T) \longrightarrow Y(T)
$$
而且它对概形 $T$（在 $B$ 上）具有函子性。由于 $f$ 是
$(\Sch/S)_{fppf}$ 上、位于层 $B$ 之上的层映射，显然 $f$ 决定这个规则，
也由这个规则决定。更一般地，我们对纤维积之间的映射采用相同记号。
例如，若 $X$、$Y$、$Z$ 是 $B$ 上的代数空间，而
$m : X \times_B Y \to Z \times_B Z$ 是 $B$ 上代数空间的态射，
那么可把 $m$ 看作对应于如下 $T$-值点之间的一族映射：
$$
X(T) \times Y(T) \longrightarrow Z(T) \times Z(T).
$$
其余情形依此类推。

\medskip\noindent
最后，给定代数空间的两个映射 $f, g : X \to Y$（在 $B$ 上），若诱导映射
$f, g : X(T) \to Y(T)$ 对每个概形 $T$（在 $B$ 上）都相等，那么
$f = g$；因而映射 $f, g : X(Z) \to Y(Z)$ 对每个第三个代数空间
$Z$（在 $B$ 上）也相等。例如，要检验群代数空间 $G$（在 $B$ 上）的
% P09-ERR-0866
公理，只须在 $T$-值点上检验图表的交换性；这里 $T$ 是 $B$ 上的概形，
% P09-ERR-0867
具体做法如下面的定义 \ref{spaces-groupoids-definition-group-space} 所示。


\section{等价关系}
\label{spaces-groupoids-section-equivalence-relations}

\noindent
记号参见《群胚》第 \ref{groupoids-section-equivalence-relations} 节。

\begin{definition}
\label{spaces-groupoids-definition-equivalence-relation}
设 $B \to S$ 如第 \ref{spaces-groupoids-section-notation} 节所示。设 $U$ 为代数空间（在
$B$ 上）。
\begin{enumerate}
\item $U$ 在 $B$ 上的一个\emph{预关系}，是任意态射
$j : R \to U \times_B U$（它是 $B$ 上代数空间的态射）。在这种情形下，置
$t = \text{pr}_0 \circ j$ 和 $s = \text{pr}_1 \circ j$，从而
$j = (t, s)$。
\item $U$ 在 $B$ 上的一个\emph{关系}，是单态射
$j : R \to U \times_B U$（它是 $B$ 上代数空间的态射）。
\item 一个\emph{预等价关系}，是预关系 $j : R \to U \times_B U$，
它满足：映射 $j : R(T) \to U(T) \times U(T)$ 的像对所有概形
$T$（在 $B$ 上）都是等价关系。
\item 称态射 $R \to U \times_B U$（它是 $B$ 上代数空间的态射）是
\emph{$U$ 在 $B$ 上的等价关系}，当且仅当对每个 $T$（在 $B$ 上），
相应的 $T$-值点（属于 $R$）在相应的 $T$-值点集（属于 $U$）上定义一个
等价关系。
\end{enumerate}
\end{definition}

\noindent
换言之，等价关系就是满足 $j$ 为关系的预等价关系。

\begin{lemma}
\label{spaces-groupoids-lemma-restrict-relation}
设 $B \to S$ 如第 \ref{spaces-groupoids-section-notation} 节所示。设 $U$ 为代数空间（在
$B$ 上）。设 $j : R \to U \times_B U$ 为预关系。设
$g : U' \to U$ 为代数空间的态射（在 $B$ 上）。最后，置
$$
R' = (U' \times_B U')\times_{U \times_B U} R
\xrightarrow{j'}
U' \times_B U'
$$
那么 $j'$ 是 $U'$ 在 $B$ 上的预关系。如果 $j$ 是关系，那么 $j'$ 也是关系。
如果 $j$ 是预等价关系，那么 $j'$ 也是预等价关系。如果 $j$ 是等价关系，
那么 $j'$ 也是等价关系。
\end{lemma}

\begin{proof}
略。
\end{proof}







\begin{definition}
\label{spaces-groupoids-definition-restrict-relation}
设 $B \to S$ 如第 \ref{spaces-groupoids-section-notation} 节所示。设 $U$ 为代数空间（在
$B$ 上）。设 $j : R \to U \times_B U$ 为预关系。设
$g : U' \to U$ 为代数空间的态射（在 $B$ 上）。引理
\ref{spaces-groupoids-lemma-restrict-relation} 中的预关系
$j' : R' \to U' \times_B U'$ 称为预关系 $j$ 到 $U'$ 的\emph{限制}或
\emph{拉回}。在这种情形下，有时写成 $R' = R|_{U'}$。
\end{definition}

\begin{lemma}
\label{spaces-groupoids-lemma-pre-equivalence-equivalence-relation-points}
设 $B \to S$ 如第 \ref{spaces-groupoids-section-notation} 节所示。设
$j : R \to U \times_B U$ 为代数空间的预关系（在 $B$ 上）。考虑由下述规则
在 $|U|$ 上定义的关系：
$$
x \sim y
\Leftrightarrow
\exists\ r \in |R| :
t(r) = x,
s(r) = y.
$$
如果 $j$ 是预等价关系，那么这就是等价关系。
\end{lemma}

\begin{proof}
假设 $x \sim y$ 且 $y \sim z$。选取 $r \in |R|$，使得
$t(r) = x$、$s(r) = y$；再选取 $r' \in |R|$，使得
$t(r') = y$、$s(r') = z$。可以选取域 $K$，使得 $r$ 和 $r'$ 可由态射
$r, r' : \Spec(K) \to R$ 表示，并且 $s \circ r = t \circ r'$。
记 $x = t \circ r$、$y = s \circ r = t \circ r'$ 和
$z = s \circ r'$，于是 $x, y, z : \Spec(K) \to U$。由构造，
$(x, y) \in j(R(K))$ 且 $(y, z) \in j(R(K))$。由于 $j$ 是预等价关系，
可知 $(x, z) \in j(R(K))$ 也成立。这显然推出 $x \sim z$。

\medskip\noindent
略去 $\sim$ 具有自反性和对称性的证明。
\end{proof}
\section{群代数空间}
\label{spaces-groupoids-section-group-spaces}

\noindent
记号参见《群胚》第 \ref{groupoids-section-group-schemes} 节。

\begin{definition}
\label{spaces-groupoids-definition-group-space}
设 $B \to S$ 如第 \ref{spaces-groupoids-section-notation} 节所示。
\begin{enumerate}
\item 一个\emph{$B$ 上的群代数空间}是偶对 $(G, m)$，其中 $G$ 是
$B$ 上的代数空间，而 $m : G \times_B G \to G$ 是 $B$ 上代数空间的态射，
并具有如下性质：对每个概形 $T$（在 $B$ 上），偶对 $(G(T), m)$ 是群。
\item 一个\emph{态射 $\psi : (G, m) \to (G', m')$（它是 $B$ 上
群代数空间的态射）}，是态射 $\psi : G \to G'$（它是 $B$ 上代数空间的态射），
并且对每个 $T/B$，诱导映射
$\psi : G(T) \to G'(T)$ 都是群同态。
\end{enumerate}
\end{definition}

\noindent
设 $(G, m)$ 为代数空间 $B$ 上的群代数空间。根据《群胚》第
\ref{groupoids-section-group-schemes} 节的讨论，得到 $B$ 上代数空间的态射
（单位元）$e : B \to G$ 和（逆元）$i : G \to G$，使得对每个 $T$，
四元组 $(G(T), m, e, i)$ 满足群公理。

\medskip\noindent
设 $(G, m)$、$(G', m')$ 为 $B$ 上的群代数空间。设
$f : G \to G'$ 为 $B$ 上代数空间的态射。由定义，$f$ 是 $B$ 上群代数空间的
态射，当且仅当下图交换：
$$
\xymatrix{
G \times_B G \ar[r]_-{f \times f} \ar[d]_m &
% P09-ERR-0868
G' \times_B G' \ar[d]^m \\
G \ar[r]^f & G'
}
$$

\begin{lemma}
\label{spaces-groupoids-lemma-base-change-group-space}
设 $B \to S$ 如第 \ref{spaces-groupoids-section-notation} 节所示。设 $(G, m)$ 为 $B$ 上的
群代数空间。设 $B' \to B$ 为代数空间的态射。拉回
$(G_{B'}, m_{B'})$ 是 $B'$ 上的群代数空间。
\end{lemma}

\begin{proof}
略。
\end{proof}







\section{群代数空间的性质}
\label{spaces-groupoids-section-properties-group-spaces}

\noindent
本节汇集群代数空间在任意基底上成立的一些简单性质。

\begin{lemma}
\label{spaces-groupoids-lemma-group-scheme-separated}
设 $S$ 为概形，$B$ 为 $S$ 上的代数空间，$G$ 为 $B$ 上的群代数空间。
那么 $G \to B$ 分离（相应地，拟分离；局部分离），当且仅当单位元态射
$e : B \to G$ 是闭浸入（相应地，拟紧；浸入）。
\end{lemma}

\begin{proof}
回忆《空间的态射》引理 \ref{spaces-morphisms-lemma-section-immersion}：
$e$ 在 $G \to B$ 分离（相应地，拟分离；局部分离）时是闭浸入
（相应地，拟紧；浸入）。为证明逆命题，考虑图表
$$
\xymatrix{
G \ar[r]_-{\Delta_{G/B}} \ar[d] &
G \times_B G \ar[d]^{(g, g') \mapsto m(i(g), g')} \\
B \ar[r]^e & G
}
$$
运用代数几何中的函子点观点，不难验证这个图表是笛卡尔图表。换言之，
$\Delta_{G/B}$ 是 $e$ 的基变换。因此，如果 $e$ 是闭浸入（相应地，拟紧；
浸入），那么 $\Delta_{G/B}$ 也具有相应性质；参见《空间》引理
\ref{spaces-lemma-base-change-immersions}（相应地，《空间的态射》引理
\ref{spaces-morphisms-lemma-base-change-quasi-compact}；《空间》引理
\ref{spaces-lemma-base-change-immersions}）。
\end{proof}

\begin{lemma}
\label{spaces-groupoids-lemma-group-scheme-unramified-or-lqf}
设 $S$ 为概形，$B$ 为 $S$ 上的代数空间，$G$ 为 $B$ 上的群代数空间。
假设 $G \to B$ 局部有限型。那么 $G \to B$ 非分歧（相应地，局部拟有限），
当且仅当 $G \to B$ 在 $e(b)$ 处非分歧（相应地，拟有限）；这对所有
$b \in |B|$ 都成立。
\end{lemma}

\begin{proof}
由《空间的态射》引理 \ref{spaces-morphisms-lemma-where-unramified}
（相应地，《空间的态射》引理
\ref{spaces-morphisms-lemma-base-change-quasi-finite-locus}），存在最大的
开子空间 $U \subset G$，使得 $U \to B$ 非分歧（相应地，局部拟有限），
且 $U$ 的构造与基变换交换。因此，可归约到 $B = \Spec(k)$ 是域的谱这一情形。
设 $g \in G(K)$ 为取值于扩域 $K/k$ 的点。为了检验 $g$ 是否属于 $U$，
可以把基底变换到 $K$。所以只须证明
$$
G \to \Spec(k)\text{ 非分歧于 }e
\Leftrightarrow
G \to \Spec(k)\text{ 非分歧于 }g
$$
其中在 $k$ 上，$g$ 是有理点（相应地，在 $g$ 与 $e$ 处拟有限的情形也一样）。
由于沿 $g$ 的平移是
$G$ 在 $k$ 上的自同构，这一点很清楚。
\end{proof}

\begin{lemma}
\label{spaces-groupoids-lemma-open-over-which-unramified-or-lqf}
设 $S$ 为概形，$B$ 为 $S$ 上的代数空间，$G$ 为 $B$ 上的群代数空间。
假设 $G \to B$ 局部有限型。
\begin{enumerate}
\item 存在最大的开子空间 $U \subset B$，使得 $G_U \to U$ 非分歧，
且 $U$ 的构造与基变换交换。
\item 存在最大的开子空间 $U \subset B$，使得 $G_U \to U$ 局部拟有限，
且 $U$ 的构造与基变换交换。
\end{enumerate}
\end{lemma}

\begin{proof}
由《空间的态射》引理 \ref{spaces-morphisms-lemma-where-unramified}
（相应地，《空间的态射》引理
\ref{spaces-morphisms-lemma-base-change-quasi-finite-locus}），存在最大的
开子空间 $W \subset G$，使得 $W \to B$ 非分歧（相应地，局部拟有限）。
此外，$W$ 的构造与基变换交换。由引理
\ref{spaces-groupoids-lemma-group-scheme-unramified-or-lqf}，可知两种情形下都有
$U = e^{-1}(W)$。
\end{proof}
\section{群代数空间的例}
\label{spaces-groupoids-section-examples-group-spaces}

\noindent
如果 $G \to S$ 是基概形 $S$ 上的群概形，那么由引理
\ref{spaces-groupoids-lemma-base-change-group-space}，基变换 $G_B$ 是群代数空间；这里 $B$ 是
$S$ 上任意代数空间，而所得群代数空间取在 $B$ 上。
% P09-ERR-0869
我们将在下面的例中频繁使用这一点。

\begin{example}[乘法群代数空间]
\label{spaces-groupoids-example-multiplicative-group}
设 $B \to S$ 如第 \ref{spaces-groupoids-section-notation} 节所示。考虑如下函子：它把任意概形
$T$（在 $B$ 上）映到结构层整体截面中单位元所成的群
$\Gamma(T, \mathcal{O}_T^*)$。这个函子由群代数空间
$$
\mathbf{G}_{m, B} = B \times_S \mathbf{G}_{m, S}
$$
表示；该空间取在 $B$ 上。这里 $\mathbf{G}_{m, S}$ 是 $S$ 上的乘法群概形，见《群胚》例
\ref{groupoids-example-multiplicative-group}。
\end{example}

\begin{example}[作为群代数空间的单位根]
\label{spaces-groupoids-example-roots-of-unity}
设 $B \to S$ 如第 \ref{spaces-groupoids-section-notation} 节所示。设 $n \in \mathbf{N}$。
考虑如下函子：它把任意概形 $T$（在 $B$ 上）映到
$\Gamma(T, \mathcal{O}_T^*)$ 中由 $n$ 次单位根组成的子群。这个函子由
群代数空间
$$
\mu_{n, B} = B \times_S \mu_{n, S}
$$
表示；该空间取在 $B$ 上。这里 $\mu_{n, S}$ 是 $n$ 次单位根群概形（在 $S$ 上），见《群胚》例
\ref{groupoids-example-roots-of-unity}。
\end{example}

\begin{example}[加法群代数空间]
\label{spaces-groupoids-example-additive-group}
设 $B \to S$ 如第 \ref{spaces-groupoids-section-notation} 节所示。考虑如下函子：它把任意概形
$T$（在 $B$ 上）映到结构层整体截面所成的群
$\Gamma(T, \mathcal{O}_T)$。这个函子由群代数空间
$$
\mathbf{G}_{a, B} = B \times_S \mathbf{G}_{a, S}
$$
表示；该空间取在 $B$ 上。这里 $\mathbf{G}_{a, S}$ 是 $S$ 上的加法群概形，见《群胚》例
\ref{groupoids-example-additive-group}。
\end{example}

\begin{example}[一般线性群代数空间]
\label{spaces-groupoids-example-general-linear-group}
设 $B \to S$ 如第 \ref{spaces-groupoids-section-notation} 节所示。设 $n \geq 1$。
考虑如下函子：它把任意概形 $T$（在 $B$ 上）映到群
$$
\text{GL}_n(\Gamma(T, \mathcal{O}_T))
$$
即结构层整体截面上的可逆 $n \times n$ 矩阵群。这个函子由群代数空间
$$
\text{GL}_{n, B} = B \times_S \text{GL}_{n, S}
$$
表示；该空间取在 $B$ 上。这里
% P09-ERR-0870
$\mathbf{G}_{m, S}$ 是 $S$ 上的一般线性群概形，见《群胚》例
\ref{groupoids-example-general-linear-group}。
\end{example}

\begin{example}
\label{spaces-groupoids-example-determinant}
设 $B \to S$ 如第 \ref{spaces-groupoids-section-notation} 节所示。设 $n \geq 1$。
行列式定义群代数空间的态射
$$
\det : \text{GL}_{n, B} \longrightarrow \mathbf{G}_{m, B}
$$
这是 $B$ 上的态射。它是《群胚》例
\ref{groupoids-example-determinant} 中 $S$ 上行列式态射的基变换。
\end{example}

\begin{example}[常值群代数空间]
\label{spaces-groupoids-example-constant-group}
设 $B \to S$ 如第 \ref{spaces-groupoids-section-notation} 节所示。设 $G$ 为抽象群。
考虑如下函子：它把任意概形 $T$（在 $B$ 上）映到局部常值映射
$T \to G$ 所成的群（这里 $T$ 取 Zariski 拓扑，$G$ 取离散拓扑）。
这个函子由群代数空间
$$
G_B = B \times_S G_S
$$
表示；该空间取在 $B$ 上。这里 $G_S$ 是《群胚》例 \ref{groupoids-example-constant-group}
中引入的常值群概形。
\end{example}
\section{群代数空间的作用}
\label{spaces-groupoids-section-action-group-space}

\noindent
记号参见《群胚》第 \ref{groupoids-section-action-group-scheme} 节。

\begin{definition}
\label{spaces-groupoids-definition-action-group-space}
设 $B \to S$ 如第 \ref{spaces-groupoids-section-notation} 节所示。设 $(G, m)$ 为 $B$ 上的
群代数空间。设 $X$ 为 $B$ 上的代数空间。
\begin{enumerate}
\item \emph{$G$ 在代数空间 $X/B$ 上的作用}，是态射
$a : G \times_B X \to X$（在 $B$ 上），它满足：对每个概形 $T$（在 $B$ 上），映射
$a : G(T) \times X(T) \to X(T)$ 定义一个 $G(T)$-集结构，其底层集合为
$X(T)$。
\item 假设 $X$、$Y$ 是 $B$ 上的代数空间，并且各自配备 $G$ 的作用。
一个\emph{等变}态射，更精确地说，一个\emph{$G$-等变态射}
$\psi : X \to Y$，是 $B$ 上代数空间的态射，它满足：对每个 $T$（在
$B$ 上），映射 $\psi : X(T) \to Y(T)$ 是 $G(T)$-集的态射。
\end{enumerate}
\end{definition}

\noindent
在情形 (1) 中，这意味着图表
\begin{equation}
\label{spaces-groupoids-equation-action}
\xymatrix{
G \times_B G \times_B X \ar[r]_-{1_G \times a} \ar[d]_{m \times 1_X} &
G \times_B X \ar[d]^a \\
G \times_B X \ar[r]^a & X
}
\quad
\xymatrix{
G \times_B X \ar[r]_-a & X \\
X\ar[u]^{e \times 1_X} \ar[ru]_{1_X}
}
\end{equation}
交换。在情形 (2) 中，这仅意味着图表
$$
\xymatrix{
% P09-ERR-0871
G \times_B X \ar[r]_-{\text{id} \times f} \ar[d]_a &
G \times_B Y \ar[d]^a \\
X \ar[r]^f & Y
}
$$
交换。

\begin{definition}
\label{spaces-groupoids-definition-free-action}
设 $B \to S$、$G \to B$ 和 $X \to B$ 如定义
\ref{spaces-groupoids-definition-action-group-space} 所示。设
$a : G \times_B X \to X$ 为 $G$ 在 $X/B$ 上的作用。如果对每个概形
$T$（在 $B$ 上），作用 $a : G(T) \times X(T) \to X(T)$ 都是群
$G(T)$ 在集合 $X(T)$ 上的自由作用，就称该作用\emph{自由}。
\end{definition}

\begin{lemma}
\label{spaces-groupoids-lemma-free-action}
情形如定义 \ref{spaces-groupoids-definition-free-action} 所示。
% P09-ERR-0872
作用 $a$ 自由，当且仅当
$$
G \times_B X \to X \times_B X, \quad (g, x) \mapsto (a(g, x), x)
$$
是代数空间的单态射。
\end{lemma}

\begin{proof}
由定义立即得到。
\end{proof}
\section{主齐性空间}
\label{spaces-groupoids-section-principal-homogeneous}

\noindent
本节对应于《群胚》第
\ref{groupoids-section-principal-homogeneous} 节。我们建议先阅读该节。

\begin{definition}
\label{spaces-groupoids-definition-pseudo-torsor}
设 $S$ 为概形。设 $B$ 为 $S$ 上的代数空间。设 $(G, m)$ 为 $B$ 上的
群代数空间。设 $X$ 为 $B$ 上的代数空间，并设
$a : G \times_B X \to X$ 为 $G$ 在 $X$ 上的作用。
\begin{enumerate}
\item 称 $X$ 为一个\emph{伪 $G$-挠子}，或称 $X$
\emph{在 $G$ 作用下形式主齐性}，如果诱导态射
$G \times_B X \to X \times_B X$，
$(g, x) \mapsto (a(g, x), x)$ 是同构。
\item 一个伪 $G$-挠子 $X$ 称为\emph{平凡的}，如果存在
% P09-ERR-0873
一个 $G$-等变同构 $G \to X$（在 $B$ 上），其中 $G$ 通过左乘作用在
$G$ 上。
\end{enumerate}
\end{definition}

\noindent
显然，如果 $B' \to B$ 是代数空间的态射，那么伪挠子的拉回
$X_{B'}$（原挠子为伪 $G$-挠子，取在 $B$ 上）是伪
$G_{B'}$-挠子（取在 $B'$ 上）。

\begin{lemma}
\label{spaces-groupoids-lemma-characterize-trivial-pseudo-torsors}
情形如定义 \ref{spaces-groupoids-definition-pseudo-torsor} 所示。
\begin{enumerate}
\item 代数空间 $X$ 是伪 $G$-挠子，当且仅当对每个概形 $T$（在 $B$
上），集合 $X(T)$ 或者为空，或者群 $G(T)$ 在 $X(T)$ 上的作用单传递。
\item 伪 $G$-挠子 $X$ 平凡，当且仅当态射 $X \to B$ 有截面。
\end{enumerate}
\end{lemma}

\begin{proof}
略。
\end{proof}

\begin{definition}
\label{spaces-groupoids-definition-principal-homogeneous-space}
设 $S$ 为概形。设 $B$ 为 $S$ 上的代数空间。设 $(G, m)$ 为 $B$ 上的
群代数空间。设 $X$ 为伪 $G$-挠子（在 $B$ 上）。
\begin{enumerate}
\item 称 $X$ 为一个\emph{主齐性空间}，更精确地说，称它为一个
\emph{主齐性 $G$-空间（在 $B$ 上）}，如果存在一个 fpqc 覆盖\footnote{《群胚》定义
\ref{groupoids-definition-principal-homogeneous-space} 中默认的挠子类型，
是在某个 fpqc 覆盖上平凡的伪挠子。由于
% P09-ERR-0874
$G$ 作为代数空间可以视为群层，那里已经有一个 $G$-挠子的概念；这个概念
% P09-ERR-0875
对应于 fppf-挠子，见引理 \ref{spaces-groupoids-lemma-torsor}。因此，对于 fpqc 局部平凡
的伪挠子，我们使用“主齐性空间”这一名称，并尽量避免在这种情形下使用
“挠子”一词。}
$\{B_i \to B\}_{i \in I}$，使得每个 $X_{B_i} \to B_i$ 都有截面
（即它是平凡的伪 $G_{B_i}$-挠子）。
\item 设 $\tau \in \{Zariski, \etale, smooth, syntomic, fppf\}$。
称 $X$ 为一个\emph{$G$-挠子（取 $\tau$ 拓扑）}，或称为一个
\emph{$\tau$ $G$-挠子}，或简称为一个\emph{$\tau$-挠子}，如果存在
一个 $\tau$ 覆盖 $\{B_i \to B\}_{i \in I}$，使得每个
$X_{B_i} \to B_i$ 都有截面。
\item 如果 $X$ 是主齐性 $G$-空间（在 $B$ 上），并且它是
\'etale 拓扑中的挠子，就称它\emph{拟等平凡}。
\item 如果 $X$ 是主齐性 $G$-空间（在 $B$ 上），并且它是 Zariski
拓扑中的挠子，就称它\emph{局部平凡}。
\end{enumerate}
\end{definition}

\noindent
我们有时说“设 $X$ 为 $G$-主齐性空间（在 $B$ 上）”，意思是说 $X$ 是
$B$ 上配备了 $G$ 的一个作用的代数空间，并且这个作用使它成为 $B$ 上的
主齐性空间。下面我们证明，在两种记号都适用时，它们相互一致。

\begin{lemma}
\label{spaces-groupoids-lemma-torsor}
设 $S$ 为概形。设 $(G, m)$ 为 $S$ 上的群代数空间。设 $X$ 为 $S$
上的代数空间，并设 $a : G \times_S X \to X$ 为 $G$ 在 $X$ 上的
作用。那么，按定义 \ref{spaces-groupoids-definition-principal-homogeneous-space}，
$X$ 是 $G$-挠子（取 $fppf$-拓扑），当且仅当按《位点上的上同调》定义
\ref{sites-cohomology-definition-torsor}，$X$ 是 $G$-挠子（在
$(\Sch/S)_{fppf}$ 上）。
\end{lemma}

\begin{proof}
略。
\end{proof}

\begin{lemma}
\label{spaces-groupoids-lemma-pseudo-torsor-implications}
设 $S$ 为概形。设 $B$ 为 $S$ 上的代数空间。设 $G$ 为 $B$ 上的群
代数空间。设 $X$ 为伪 $G$-挠子（在 $B$ 上）。假设
% P09-ERR-0876
$G$ 和 $X$ 在 $B$ 上局部有限型。
\begin{enumerate}
\item 如果 $G \to B$ 非分歧，那么 $X \to B$ 非分歧。
\item 如果 $G \to B$ 局部拟有限，那么 $X \to B$ 局部拟有限。
\end{enumerate}
\end{lemma}

\begin{proof}
(1) 的证明。由《空间的态射》引理
\ref{spaces-morphisms-lemma-where-unramified}，我们约化到 $B$ 是域的
谱的情形。如果 $X$ 为空，则结论成立。如果 $X$ 非空，那么在扩大该域之后，
可以假设 $X$ 有一个点。于是 $G \cong X$，结论成立。

\medskip\noindent
(2) 的证明完全相同，只需使用《空间的态射》引理
\ref{spaces-morphisms-lemma-base-change-quasi-finite-locus}。
\end{proof}
\section{等变拟凝聚层}
\label{spaces-groupoids-section-equivariant}

\noindent
请与《群胚》第 \ref{groupoids-section-equivariant} 节比较。

\begin{definition}
\label{spaces-groupoids-definition-equivariant-module}
设 $B \to S$ 如第 \ref{spaces-groupoids-section-notation} 节所示。设 $(G, m)$ 为 $B$ 上的
群代数空间，并设 $a : G \times_B X \to X$ 为 $G$ 在代数空间 $X$ 上的
作用，该空间取在 $B$ 上。一个\emph{$G$-等变拟凝聚
$\mathcal{O}_X$-模}，或简称一个\emph{等变拟凝聚 $\mathcal{O}_X$-模}，
是一个偶 $(\mathcal{F}, \alpha)$，其中 $\mathcal{F}$ 是拟凝聚
$\mathcal{O}_X$-模，而 $\alpha$ 是 $\mathcal{O}_{G \times_B X}$-模映射
$$
\alpha : a^*\mathcal{F} \longrightarrow \text{pr}_1^*\mathcal{F}
$$
这里 $\text{pr}_1 : G \times_B X \to X$ 是投影，并且满足：
\begin{enumerate}
\item 图表
$$
\xymatrix{
% P09-ERR-0877
(1_G \times a)^*\text{pr}_2^*\mathcal{F} \ar[r]_-{\text{pr}_{12}^*\alpha} &
\text{pr}_2^*\mathcal{F} \\
(1_G \times a)^*a^*\mathcal{F} \ar[u]^{(1_G \times a)^*\alpha} \ar@{=}[r] &
(m \times 1_X)^*a^*\mathcal{F} \ar[u]_{(m \times 1_X)^*\alpha}
}
$$
% P09-ERR-0878
在 $\mathcal{O}_{G \times_B G \times_B X}$-模范畴中交换；
\item 拉回
$$
(e \times 1_X)^*\alpha : \mathcal{F} \longrightarrow \mathcal{F}
$$
是恒等映射。
\end{enumerate}
解释可与式 (\ref{spaces-groupoids-equation-action}) 的有关图表比较。
\end{definition}

\noindent
注意，第一个图表的交换性保证 $(e \times 1_X)^*\alpha$ 是
$\mathcal{F}$ 上的幂等算子；因此条件 (2) 恰好是说它为同构。

\begin{lemma}
\label{spaces-groupoids-lemma-pullback-equivariant}
设 $B \to S$ 如第 \ref{spaces-groupoids-section-notation} 节所示。设 $G$ 为 $B$ 上的
群代数空间。设 $f : X \to Y$ 为一个 $G$-等变态射，它连接两个在 $B$
上且配备 $G$-作用的代数空间。
% P09-ERR-0879
拉回 $f^*$ 由
$(\mathcal{F}, \alpha) \mapsto (f^*\mathcal{F}, (1_G \times f)^*\alpha)$
给出，它定义一个从拟凝聚 $G$-等变层（在 $Y$ 上）的范畴到拟凝聚
$G$-等变层（在 $X$ 上）的范畴的函子。
\end{lemma}

\begin{proof}
略。
\end{proof}
\section{代数空间中的群胚}
\label{spaces-groupoids-section-groupoids}

\noindent
记号参见《群胚》第 \ref{groupoids-section-groupoids} 节。

\begin{definition}
\label{spaces-groupoids-definition-groupoid}
设 $B \to S$ 如第 \ref{spaces-groupoids-section-notation} 节所示。
\begin{enumerate}
\item 一个\emph{$B$ 上代数空间中的群胚}，是一个五元组
$(U, R, s, t, c)$，其中 $U$ 和 $R$ 是 $B$ 上的代数空间，而
$s, t : R \to U$ 和 $c : R \times_{s, U, t} R \to R$ 是 $B$ 上
代数空间的态射，并具有如下性质：对每个概形 $T$（在 $B$ 上），五元组
$$
(U(T), R(T), s, t, c)
$$
是群胚范畴。
\item 代数空间中的群胚的一个态射
$f : (U, R, s, t, c) \to (U', R', s', t', c')$
（取在 $B$ 上），由代数空间的态射 $f : U \to U'$ 和
$f : R \to R'$（在 $B$ 上）给出，并具有如下性质：对每个概形
$T$（在 $B$ 上），这些映射 $f$ 定义一个从群胚范畴
$(U(T), R(T), s, t, c)$ 到群胚范畴
$(U'(T), R'(T), s', t', c')$ 的函子。
\end{enumerate}
\end{definition}

\noindent
设 $(U, R, s, t, c)$ 为代数空间中的群胚（在 $B$ 上）。注意，存在唯一的
代数空间态射 $e : U \to R$ 和 $i : R \to R$（在 $B$ 上），使得对
每个概形 $T$（在 $B$ 上），诱导映射
% P09-ERR-0880
$e : U(T) \to R(T)$ 是恒等，而
% P09-ERR-0881
$i : R(T) \to R(T)$ 是该群胚范畴的逆。七元组
$(U, R, s, t, c, e, i)$ 满足若干交换图表，它们分别对应于《群胚》第
\ref{groupoids-section-groupoids} 节的公理 (1)、(2)(a)、(2)(b)、(3)(a)
和 (3)(b)。反过来，给定具有这些性质的七元组，五元组
$(U, R, s, t, c)$ 就是代数空间中的群胚（在 $B$ 上）。注意，$i$ 是
同构，而 $e$ 同时是 $s$ 和 $t$ 的截面。此外，对一个 $B$ 上代数空间
中的群胚，我们记
$$
j = (t, s) : R \longrightarrow U \times_B U
$$
这与上面第 \ref{spaces-groupoids-section-equivalence-relations} 节的约定相容。
我们有时说“设 $(U, R, s, t, c, e, i)$ 为代数空间中的群胚（在 $B$
上）”，以强调恒等态射和逆态射的存在。

\begin{lemma}
\label{spaces-groupoids-lemma-groupoid-pre-equivalence}
设 $B \to S$ 如第 \ref{spaces-groupoids-section-notation} 节所示。给定代数空间中的
群胚 $(U, R, s, t, c)$（在 $B$ 上），态射
$j : R \to U \times_B U$ 是预等价关系。
\end{lemma}

\begin{proof}
略。这是一个很好的定义练习。
\end{proof}

\begin{lemma}
\label{spaces-groupoids-lemma-equivalence-groupoid}
设 $B \to S$ 如第 \ref{spaces-groupoids-section-notation} 节所示。给定等价关系
$j : R \to U \times_B U$（在 $B$ 上），存在唯一的方式将它扩充成
代数空间中的群胚 $(U, R, s, t, c)$（在 $B$ 上）。
\end{lemma}

\begin{proof}
略。这是一个很好的定义练习。
\end{proof}

\begin{lemma}
\label{spaces-groupoids-lemma-diagram}
设 $B \to S$ 如第 \ref{spaces-groupoids-section-notation} 节所示。设
$(U, R, s, t, c)$ 为代数空间中的群胚（在 $B$ 上）。在交换图表
$$
\xymatrix{
& U & \\
R \ar[d]_s \ar[ru]^t &
R \times_{s, U, t} R
\ar[l]^-{\text{pr}_0} \ar[d]^{\text{pr}_1} \ar[r]_-c &
R \ar[d]^s \ar[lu]_t \\
U & R \ar[l]_t \ar[r]^s & U
}
$$
中，下方两个方块都是纤维积方块。此外，上方的三角形（实际上是方块）也是
笛卡尔的。
\end{lemma}

\begin{proof}
略。这是一个关于定义以及代数几何中函子点观点的练习。
\end{proof}

\begin{lemma}
\label{spaces-groupoids-lemma-diagram-pull}
设 $B \to S$ 如第 \ref{spaces-groupoids-section-notation} 节所示。设
$(U, R, s, t, c, e, i)$ 为代数空间中的群胚（在 $B$ 上）。图表
\begin{equation}
\label{spaces-groupoids-equation-pull}
\xymatrix{
R \times_{t, U, t} R
\ar@<1ex>[r]^-{\text{pr}_1} \ar@<-1ex>[r]_-{\text{pr}_0}
\ar[d]_{\text{pr}_0 \times c \circ (i, 1)} &
R \ar[r]^t \ar[d]^{\text{id}_R} &
U \ar[d]^{\text{id}_U} \\
R \times_{s, U, t} R
\ar@<1ex>[r]^-c \ar@<-1ex>[r]_-{\text{pr}_0} \ar[d]_{\text{pr}_1} &
R \ar[r]^t \ar[d]^s &
U \\
R \ar@<1ex>[r]^s \ar@<-1ex>[r]_t &
U
}
\end{equation}
交换。上方两行通过给出的竖直映射彼此同构。左下方两个方块是笛卡尔的。
\end{lemma}

\begin{proof}
图表的交换性由群胚的公理得到。注意，用群胚的语言说，左上方的竖直箭头把
具有相同目标的一对态射 $(\alpha, \beta)$ 映到一对态射
$(\alpha, \alpha^{-1} \circ \beta)$。在任意群胚中，这定义了
$$
\text{Arrows} \times_{t, \text{Ob}, t} \text{Arrows}
$$
与
$$
\text{Arrows} \times_{s, \text{Ob}, t} \text{Arrows}
$$
之间的双射。因此得到引理的第二个断言。最后一个断言由引理
\ref{spaces-groupoids-lemma-diagram} 得到。
\end{proof}

\begin{lemma}
\label{spaces-groupoids-lemma-base-change-groupoid}
设 $B \to S$ 如第 \ref{spaces-groupoids-section-notation} 节所示。设
$(U, R, s, t, c)$ 为代数空间中的群胚（在 $B$ 上）。设
$B' \to B$ 为代数空间的态射。那么基变换 $U' = B' \times_B U$、
$R' = B' \times_B R$，连同基变换 $s'$、$t'$、$c'$（分别来自态射
$s, t, c$），构成代数空间中的群胚 $(U', R', s', t', c')$（在 $B'$ 上）；
并且这些投影确定代数空间中的群胚的一个态射
$$
(U', R', s', t', c') \to (U, R, s, t, c)
$$
（取在 $B$ 上）。
\end{lemma}

\begin{proof}
略。提示：
$R' \times_{s', U', t'} R' = B' \times_B (R \times_{s, U, t} R)$。
\end{proof}
\section{群胚上的拟凝聚层}
\label{spaces-groupoids-section-groupoids-quasi-coherent}

\noindent
请与《群胚》第 \ref{groupoids-section-groupoids-quasi-coherent} 节比较。

\begin{definition}
\label{spaces-groupoids-definition-groupoid-module}
设 $B \to S$ 如第 \ref{spaces-groupoids-section-notation} 节所示。设
$(U, R, s, t, c)$ 为代数空间中的群胚（在 $B$ 上）。一个
\emph{群胚 $(U, R, s, t, c)$ 上的拟凝聚模}，是一个偶
$(\mathcal{F}, \alpha)$，其中 $\mathcal{F}$ 是拟凝聚
$\mathcal{O}_U$-模，而 $\alpha$ 是 $\mathcal{O}_R$-模映射
$$
\alpha : t^*\mathcal{F} \longrightarrow s^*\mathcal{F}
$$
并且满足：
\begin{enumerate}
\item 图表
$$
\xymatrix{
& \text{pr}_1^*t^*\mathcal{F} \ar[r]_-{\text{pr}_1^*\alpha} &
\text{pr}_1^*s^*\mathcal{F} \ar@{=}[rd] & \\
\text{pr}_0^*s^*\mathcal{F} \ar@{=}[ru] & & & c^*s^*\mathcal{F} \\
& \text{pr}_0^*t^*\mathcal{F} \ar[lu]^{\text{pr}_0^*\alpha} \ar@{=}[r] &
c^*t^*\mathcal{F} \ar[ru]_{c^*\alpha}
}
$$
% P09-ERR-0882
在 $\mathcal{O}_{R \times_{s, U, t} R}$-模范畴中交换；
\item 拉回
$$
e^*\alpha : \mathcal{F} \longrightarrow \mathcal{F}
$$
是恒等映射。
\end{enumerate}
请与引理 \ref{spaces-groupoids-lemma-diagram} 的交换图表比较。
\end{definition}

\noindent
第一个图表的交换性迫使算子 $e^*\alpha$ 为幂等。因此第二个条件可以改述为
$e^*\alpha$ 是同构。事实上，该条件蕴含 $\alpha$ 是同构。

\begin{lemma}
\label{spaces-groupoids-lemma-isomorphism}
设 $B \to S$ 如第 \ref{spaces-groupoids-section-notation} 节所示。设
$(U, R, s, t, c)$ 为代数空间中的群胚（在 $B$ 上）。如果
$(\mathcal{F}, \alpha)$ 是群胚 $(U, R, s, t, c)$ 上的拟凝聚模，
那么 $\alpha$ 是同构。
\end{lemma}

\begin{proof}
把定义 \ref{spaces-groupoids-definition-groupoid-module} 中的交换图表沿态射
$(i, 1) : R \to R \times_{s, U, t} R$ 拉回。于是可见
$i^*\alpha \circ \alpha = s^*e^*\alpha$。沿态射 $(1, i)$ 拉回，
得到关系
$\alpha \circ i^*\alpha = t^*e^*\alpha$。由第二个假设，这些态射
是恒等态射。因此 $i^*\alpha$ 是 $\alpha$ 的逆。
\end{proof}

\begin{lemma}
\label{spaces-groupoids-lemma-pullback}
设 $B \to S$ 如第 \ref{spaces-groupoids-section-notation} 节所示。考虑一个态射
$f : (U, R, s, t, c) \to (U', R', s', t', c')$，
% P09-ERR-0883
它是代数空间中的群胚的态射（在 $B$ 上）。
% P09-ERR-0884
拉回 $f^*$ 由
$$
(\mathcal{F}, \alpha) \mapsto (f^*\mathcal{F}, f^*\alpha)
$$
给出，并定义一个从群胚 $(U', R', s', t', c')$ 上的拟凝聚层范畴到群胚
$(U, R, s, t, c)$ 上的拟凝聚层范畴的函子。
\end{lemma}

\begin{proof}
略。
\end{proof}

\begin{lemma}
\label{spaces-groupoids-lemma-pushforward}
设 $B \to S$ 如第 \ref{spaces-groupoids-section-notation} 节所示。考虑代数空间中的群胚的
一个态射
$f : (U, R, s, t, c) \to (U', R', s', t', c')$（在 $B$ 上）。
假设：
\begin{enumerate}
\item $f : U \to U'$ 拟紧且拟分离；
\item 方块
$$
\xymatrix{
R \ar[d]_t \ar[r]_f & R' \ar[d]^{t'} \\
U \ar[r]^f & U'
}
$$
是笛卡尔的；
\item $s'$ 和 $t'$ 平坦。
\end{enumerate}
% P09-ERR-0885
前推 $f_*$ 由
$$
(\mathcal{F}, \alpha) \mapsto (f_*\mathcal{F}, f_*\alpha)
$$
给出，并定义一个从群胚 $(U, R, s, t, c)$ 上的拟凝聚层范畴到群胚
$(U', R', s', t', c')$ 上的拟凝聚层范畴的函子；这个函子是引理
\ref{spaces-groupoids-lemma-pullback} 所定义拉回的右伴随。
\end{lemma}

\begin{proof}
由于 $U \to U'$ 拟紧且拟分离，可知 $f_*$ 把拟凝聚层变为拟凝聚层
（《空间的态射》引理 \ref{spaces-morphisms-lemma-pushforward}）。
此外，由于方块
$$
\vcenter{
\xymatrix{
R \ar[d]_t \ar[r]_f & R' \ar[d]^{t'} \\
U \ar[r]^f & U'
}
}
\quad\text{和}\quad
\vcenter{
\xymatrix{
R \ar[d]_s \ar[r]_f & R' \ar[d]^{s'} \\
U \ar[r]^f & U'
}
}
$$
是笛卡尔的，得到
$(t')^*f_*\mathcal{F} = f_*t^*\mathcal{F}$ 和
$(s')^*f_*\mathcal{F} = f_*s^*\mathcal{F}$，见《空间的上同调》引理
\ref{spaces-cohomology-lemma-flat-base-change-cohomology}。因此，把
$f_*\alpha$ 看作映射
$(t')^*f_*\mathcal{F} \to (s')^*f_*\mathcal{F}$ 是有意义的。
类似的论证表明 $f_*\alpha$ 满足余圈条件。这个函子与拉回函子互为伴随，
因为环化空间上模的拉回与前推互为伴随。略去一些细节。
\end{proof}

\begin{lemma}
\label{spaces-groupoids-lemma-colimits}
设 $B \to S$ 如第 \ref{spaces-groupoids-section-notation} 节所示。设
$(U, R, s, t, c)$ 为代数空间中的群胚（在 $B$ 上）。群胚
$(U, R, s, t, c)$ 上的拟凝聚模范畴有余极限。
\end{lemma}

\begin{proof}
设 $i \mapsto (\mathcal{F}_i, \alpha_i)$ 为以范畴 $\mathcal{I}$ 为
指标的图表。可以形成余极限
$\mathcal{F} = \colim \mathcal{F}_i$；它是 $U$ 上的拟凝聚层，见
《空间的性质》引理
\ref{spaces-properties-lemma-properties-quasi-coherent}。由于余极限与
拉回交换，可知
$s^*\mathcal{F} = \colim s^*\mathcal{F}_i$，类似地
$t^*\mathcal{F} = \colim t^*\mathcal{F}_i$。因此可以置
$\alpha = \colim \alpha_i$。我们略去如下断言的证明：
$(\mathcal{F}, \alpha)$ 是拟凝聚模范畴中该图表的余极限，该范畴取在群胚
$(U, R, s, t, c)$ 上。
\end{proof}

\begin{lemma}
\label{spaces-groupoids-lemma-abelian}
设 $B \to S$ 如第 \ref{spaces-groupoids-section-notation} 节所示。设
$(U, R, s, t, c)$ 为代数空间中的群胚（在 $B$ 上）。如果 $s$、$t$
平坦，那么群胚 $(U, R, s, t, c)$ 上的拟凝聚模范畴是 Abel 范畴。
\end{lemma}

\begin{proof}
设 $\varphi : (\mathcal{F}, \alpha) \to (\mathcal{G}, \beta)$ 是
群胚 $(U, R, s, t, c)$ 上拟凝聚模的同态。由于 $s$ 平坦，可知
$$
0 \to s^*\Ker(\varphi)
\to s^*\mathcal{F} \to s^*\mathcal{G} \to s^*\Coker(\varphi) \to 0
$$
正合；沿 $t$ 拉回时也一样。因此 $\alpha$ 和 $\beta$ 诱导同构
$\kappa : t^*\Ker(\varphi) \to s^*\Ker(\varphi)$ 和
$\lambda : t^*\Coker(\varphi) \to s^*\Coker(\varphi)$，它们满足
余圈条件。于是直接验证可知 $(\Ker(\varphi), \kappa)$ 和
$(\Coker(\varphi), \lambda)$ 分别是拟凝聚模范畴中的核和余核，该范畴
取在群胚 $(U, R, s, t, c)$ 上。此外，由于等式
$\Coim(\varphi) = \Im(\varphi)$ 在 $U$ 上成立，所以这里也成立。
\end{proof}
\section{拟凝聚模的余极限}
\label{spaces-groupoids-section-colimits}

\noindent
本节对应于《群胚》第 \ref{groupoids-section-colimits} 节。

\begin{lemma}
\label{spaces-groupoids-lemma-construct-quasi-coherent}
设 $B \to S$ 如第 \ref{spaces-groupoids-section-notation} 节所示。设
$(U, R, s, t, c)$ 为代数空间中的群胚（在 $B$ 上）。假设 $s, t$ 平坦、
拟紧且拟分离。对任意拟凝聚模 $\mathcal{G}$（在 $U$ 上），存在典范同构
$\alpha : t^*s_*t^*\mathcal{G} \to s^*s_*t^*\mathcal{G}$，它使
$(s_*t^*\mathcal{G}, \alpha)$ 成为群胚 $(U, R, s, t, c)$ 上的
拟凝聚模。这个构造定义函子
$$
\QCoh(\mathcal{O}_U) \longrightarrow \QCoh(U, R, s, t, c)
$$
它是遗忘函子 $(\mathcal{F}, \beta) \mapsto \mathcal{F}$ 的右伴随。
\end{lemma}

\begin{proof}
沿拟紧拟分离态射前推拟凝聚模，所得仍为拟凝聚模，见《空间的态射》引理
\ref{spaces-morphisms-lemma-pushforward}。因此 $s_*t^*\mathcal{G}$ 是
拟凝聚的。使用引理 \ref{spaces-groupoids-lemma-diagram} 的记号，有
$$
t^*s_*t^*\mathcal{G} =
\text{pr}_{1, *}\text{pr}_0^*t^*\mathcal{G} =
\text{pr}_{1, *}c^*t^*\mathcal{G} =
s^*s_*t^*\mathcal{G}
$$
中间的等号来自等式
$t \circ c = t \circ \text{pr}_0$，这是态射
$R \times_{s, U, t} R \to U$ 的等式；第一个和最后一个等号来自这些步骤中
基变换与前推交换，见《空间的上同调》引理
\ref{spaces-cohomology-lemma-flat-base-change-cohomology}。

\medskip\noindent
为了验证定义 \ref{spaces-groupoids-definition-groupoid-module} 中 $\alpha$ 的余圈条件以及
伴随性质，我们用另一种方式描述构造
$\mathcal{G} \mapsto (s_*t^*\mathcal{G}, \alpha)$。考虑
% P09-ERR-0886
代数空间中的群胚
% P09-ERR-0888
$(R, R \times_{t, U, t} R, \text{pr}_0, \text{pr}_1, \text{pr}_{02})$；
它来自等价关系 $R \times_{t, U, t} R$（定义在 $R$ 上），见引理
\ref{spaces-groupoids-lemma-equivalence-groupoid}。存在态射
$$
f :
(R, R \times_{t, U, t} R, \text{pr}_1, \text{pr}_0, \text{pr}_{02})
\longrightarrow
(U, R, s, t, c)
$$
% P09-ERR-0887
这是代数空间中的群胚的态射，由 $s : R \to U$ 和映射
$R \times_{t, U, t} R \to R$ 给出，后一个映射为
$(r_0, r_1) \mapsto r_0^{-1} \circ r_1$；我们略去所需图表交换性的验证。
由于 $t, s : R \to U$ 拟紧、拟分离且平坦，并且有笛卡尔方块
$$
\xymatrix{
R \times_{t, U, t} R \ar[d]_{\text{pr}_0}
\ar[rr]_-{(r_0, r_1) \mapsto r_0^{-1} \circ r_1} & & R \ar[d]^t \\
R \ar[rr]^s & & U
}
$$
由引理 \ref{spaces-groupoids-lemma-diagram-pull} 可知，引理 \ref{spaces-groupoids-lemma-pushforward}
适用于 $f$。因此，沿 $f$ 前推和拉回拟凝聚模是互为伴随的函子。为了完成证明，
我们把这些函子与上面描述的函子等同起来。为此注意，函子
$$
t^* :
\QCoh(\mathcal{O}_U)
\longrightarrow
\QCoh(R, R \times_{t, U, t} R, \text{pr}_1, \text{pr}_0, \text{pr}_{02})
$$
是等价，因为根据拟凝聚层的下降理论，
$\{t : R \to U\}$ 是 fpqc 覆盖；见《空间上的下降》命题
\ref{spaces-descent-proposition-fpqc-descent-quasi-coherent}。

\medskip\noindent
沿 $f$ 前推并预合成等价 $t^*$，把 $\mathcal{G}$ 映到
$(s_*t^*\mathcal{G}, \alpha)$；我们略去如下验证：以这种方式得到的同构
$\alpha$ 与上面构造的同构相同。

\medskip\noindent
沿 $f$ 拉回并后合成等价 $t^*$ 的逆，把 $(\mathcal{F}, \beta)$ 映到
相对于 $\{t : R \to U\}$ 的下降；被下降的模为 $s^*\mathcal{F}$，配备
下降资料 $\gamma$；该资料定义在 $R \times_{t, U, t} R$ 上，并且是
$\beta$ 沿映射
$(r_0, r_1) \mapsto r_0^{-1} \circ r_1$ 的拉回。考虑同构
$\beta : t^*\mathcal{F} \to s^*\mathcal{F}$。$t^*\mathcal{F}$ 上相对于
$\{t : R \to U\}$ 的典范下降资料（《空间上的下降》
定义 \ref{spaces-descent-definition-descent-datum-effective-quasi-coherent}）
经由 $\beta$ 变成映射
$$
\text{pr}_0^*s^*\mathcal{F}
\xrightarrow{\text{pr}_0^*\beta^{-1}}
\text{pr}_0^*t^*\mathcal{F}
\xrightarrow{can}
\text{pr}_1^*t^*\mathcal{F}
\xrightarrow{\text{pr}_1^*\beta}
\text{pr}_1^*s^*\mathcal{F}
$$
由于 $\beta$ 满足余圈条件，这等于 $\beta$ 沿映射
$(r_0, r_1) \mapsto r_0^{-1} \circ r_1$ 的拉回。为说明这一点，取定义
\ref{spaces-groupoids-definition-groupoid-module} 中实际的余圈关系，并沿态射
$(\text{pr}_0, c \circ (i, 1)) : R \times_{t, U, t} R \to R \times_{s, U, t} R$
将它拉回；该态射也出现在引理 \ref{spaces-groupoids-lemma-diagram-pull} 的交换图表中。
由此可知 $(s^*\mathcal{F}, \gamma)$ 同构于
$(t^*\mathcal{F}, can)$。综上，沿 $f$ 拉回并后合成等价 $t^*$ 的逆所得的
函子，同构于遗忘函子 $(\mathcal{F}, \beta) \mapsto \mathcal{F}$。
\end{proof}

\begin{remark}
\label{spaces-groupoids-remark-adjunction-map}
在引理 \ref{spaces-groupoids-lemma-construct-quasi-coherent} 的情形下，记
$$
F : \QCoh(U, R, s, t, c) \to \QCoh(\mathcal{O}_U),\quad
(\mathcal{F}, \beta) \mapsto \mathcal{F}
$$
为遗忘函子，并记
$$
G : \QCoh(\mathcal{O}_U) \to \QCoh(U, R, s, t, c),\quad
\mathcal{G} \mapsto (s_*t^*\mathcal{G}, \alpha)
$$
为引理中构造的右伴随。那么，该伴随的单位
$\eta : \text{id} \to G \circ F$ 在 $(\mathcal{F}, \beta)$ 上的取值由
映射
$$
\mathcal{F} \to s_*s^*\mathcal{F} \xrightarrow{\beta^{-1}} s_*t^*\mathcal{F}
$$
给出。我们略去验证。
\end{remark}

\begin{lemma}
\label{spaces-groupoids-lemma-push-pull}
设 $S$ 为概形。设 $f : Y \to X$ 为 $S$ 上代数空间的态射。设
$\mathcal{F}$ 为拟凝聚 $\mathcal{O}_X$-模，设 $\mathcal{G}$ 为拟凝聚
$\mathcal{O}_Y$-模，并设
$\varphi : \mathcal{G} \to f^*\mathcal{F}$ 为模映射。假设：
\begin{enumerate}
\item $\varphi$ 是单射；
\item $f$ 拟紧、拟分离、平坦且满；
\item $X$、$Y$ 局部 Noether；
\item $\mathcal{G}$ 是凝聚 $\mathcal{O}_Y$-模。
\end{enumerate}
那么 $\mathcal{F} \cap f_*\mathcal{G}$ 由拉回
$$
\xymatrix{
\mathcal{F} \ar[r] & f_*f^*\mathcal{F} \\
\mathcal{F} \cap f_*\mathcal{G} \ar[u] \ar[r] &
f_*\mathcal{G} \ar[u]
}
$$
定义，并且它是凝聚 $\mathcal{O}_X$-模。
\end{lemma}

\begin{proof}
我们将自由使用《空间的上同调》引理
\ref{spaces-cohomology-lemma-coherent-Noetherian} 对凝聚模的刻画，以及
凝聚模构成 $\QCoh(\mathcal{O}_X)$ 的 Serre 子范畴这一事实；见
《空间的上同调》引理
\ref{spaces-cohomology-lemma-coherent-Noetherian-quasi-coherent-sub-quotient}。
如果 $f$ 有截面 $\sigma$，那么可见 $\mathcal{F} \cap f_*\mathcal{G}$
包含在
$\sigma^*\mathcal{G} \to \sigma^*f^*\mathcal{F} = \mathcal{F}$
的像中，因而是凝聚的。一般地，为证明
$\mathcal{F} \cap f_*\mathcal{G}$ 凝聚，
% P09-ERR-0889
只需证明 $f^*(\mathcal{F} \cap f_*\mathcal{G})$ 凝聚（见
《空间上的下降》引理 \ref{spaces-descent-lemma-finite-type-descends}）。
由于 $f$ 平坦，这等于
$f^*\mathcal{F} \cap f^*f_*\mathcal{G}$。由于 $f$ 平坦、拟紧且拟分离，
可知 $f^*f_*\mathcal{G} = p_*q^*\mathcal{G}$，其中
$p, q : Y \times_X Y \to Y$ 是投影；见《空间的上同调》引理
\ref{spaces-cohomology-lemma-flat-base-change-cohomology}。由于 $p$
有截面，结论成立。
\end{proof}

\noindent
设 $B \to S$ 如第 \ref{spaces-groupoids-section-notation} 节所示。设
$(U, R, s, t, c)$ 为代数空间中的群胚（在 $B$ 上）。假设 $U$ 局部
Noether。在下面的引理中，如果拟凝聚层 $(\mathcal{F}, \alpha)$（定义在群胚
$(U, R, s, t, c)$ 上）满足 $\mathcal{F}$ 是凝聚
$\mathcal{O}_U$-模，就称它\emph{凝聚}。

\begin{lemma}
\label{spaces-groupoids-lemma-colimit-coherent}
设 $B \to S$ 如第 \ref{spaces-groupoids-section-notation} 节所示。设
$(U, R, s, t, c)$ 为代数空间中的群胚（在 $B$ 上）。假设：
\begin{enumerate}
\item $U$、$R$ 是 Noether 的；
\item $s, t$ 平坦、拟紧且拟分离。
\end{enumerate}
那么每个拟凝聚模 $(\mathcal{F}, \alpha)$（定义在群胚
$(U, R, s, t, c)$ 上）都是凝聚模的滤过余极限。
\end{lemma}

\begin{proof}
我们将使用《空间的上同调》引理
\ref{spaces-cohomology-lemma-coherent-Noetherian} 对局部 Noether
代数空间上凝聚模的刻画，而不再提及。可以把
$\mathcal{F} = \colim \mathcal{H}_i$ 写成凝聚子模
$\mathcal{H}_i \subset \mathcal{F}$ 的滤过余极限，见《空间的上同调》
引理 \ref{spaces-cohomology-lemma-directed-colimit-coherent}。给定拟凝聚层
$\mathcal{H}$（在 $U$ 上），记
$(s_*t^*\mathcal{H}, \alpha)$ 为引理
\ref{spaces-groupoids-lemma-construct-quasi-coherent} 所给出的群胚
$(U, R, s, t, c)$ 上的拟凝聚层。考虑伴随映射
$(\mathcal{F}, \beta) \to (s_*t^*\mathcal{F}, \alpha)$（在
$\QCoh(U, R, s, t, c)$ 中），见注
\ref{spaces-groupoids-remark-adjunction-map}。置
$$
(\mathcal{F}_i, \beta_i) =
(\mathcal{F}, \beta)
\times_{(s_*t^*\mathcal{F}, \alpha)}
(s_*t^*\mathcal{H}_i, \alpha)
$$
于 $\QCoh(U, R, s, t, c)$ 中。由引理 \ref{spaces-groupoids-lemma-abelian} 的证明，
限制到 $U$ 在 $\QCoh(U, R, s, t, c)$ 上是正合函子，因此得到拉回图表
$$
\xymatrix{
\mathcal{F} \ar[r] & s_*t^*\mathcal{F} \\
\mathcal{F}_i \ar[r] \ar[u] & s_*t^*\mathcal{H}_i \ar[u]
}
$$
换言之，$\mathcal{F}_i = \mathcal{F} \cap s_*t^*\mathcal{H}_i$。
根据注 \ref{spaces-groupoids-remark-adjunction-map} 中伴随映射的描述，该图表同构于图表
$$
\xymatrix{
\mathcal{F} \ar[r] & s_*s^*\mathcal{F} \\
\mathcal{F}_i \ar[r] \ar[u] & s_*t^*\mathcal{H}_i \ar[u]
}
$$
其中右侧竖直箭头是把 $s_*$ 施于映射
$$
t^*\mathcal{H}_i \to t^*\mathcal{F} \xrightarrow{\beta} s^*\mathcal{F}
$$
所得的结果。该箭头是单射，因为 $t$ 是平坦态射。由引理
\ref{spaces-groupoids-lemma-push-pull} 可知 $\mathcal{F}_i$ 凝聚。最后，因为 $s$ 拟紧且
拟分离，可知 $s_*$ 与余极限交换（见《概形的上同调》引理
\ref{coherent-lemma-colimit-cohomology}）。因此
$s_*t^*\mathcal{F} = \colim s_*t^*\mathcal{H}_i$，从而如期有
$(\mathcal{F}, \beta) = \colim (\mathcal{F}_i, \beta_i)$。
\end{proof}

\section{拟凝聚层中的晶体}
\label{spaces-groupoids-section-crystals}

\noindent
设 $(I, \Phi, j)$ 由集合 $I$ 和预关系
$j : \Phi \to I \times I$ 组成。假设对每个 $i \in I$ 给定概形 $X_i$，
并对每个 $\phi \in \Phi$ 给定概形态射
$f_\phi : X_{i'} \to X_i$，其中 $j(\phi) = (i, i')$。置
$X = (\{X_i\}_{i \in I}, \{f_\phi\}_{\phi \in \Phi})$。
所谓 $X$ 上拟凝聚模的\emph{晶体}，是指如下规则：它对每个
% P09-ERR-0890
$i \in \Ob(\mathcal{I})$ 指定拟凝聚层 $\mathcal{F}_i$（在 $X_i$ 上），
并对每个 $\phi \in \Phi$（满足 $j(\phi) = (i, i')$）指定拟凝聚层同构
$$
\alpha_\phi : f_\phi^*\mathcal{F}_i \longrightarrow \mathcal{F}_{i'}
$$
（定义在 $X_{i'}$ 上）。
这些拟凝聚模晶体构成加性范畴 $\textit{CQC}(X)$\footnote{可以选出一族三元组
$\phi, \phi', \phi'' \in \Phi$，它们满足
$j(\phi) = (i, i')$、$j(\phi') = (i', i'')$ 和
$j(\phi'') = (i, i'')$，并且
$f_{\phi''} = f_\phi \circ f_{\phi'}$；然后要求这些三元组满足
$\alpha_{\phi'} \circ f_{\phi'}^*\alpha_\phi = \alpha_{\phi''}$。
这会定义一个加性子范畴。例如，资料 $(I, \Phi)$ 可以是某个指标范畴的
对象集与箭头集，而 $X$ 可以是该指标范畴上的概形图表。引理
\ref{spaces-groupoids-lemma-crystals-in-quasi-coherent-modules} 的结论立即给出相应子范畴中的结论。}。
这个范畴有余极限（证明与引理 \ref{spaces-groupoids-lemma-colimits} 的证明相同）。
如果所有态射 $f_\phi$ 都平坦，那么 $\textit{CQC}(X)$ 是 Abel 范畴
（证明与引理 \ref{spaces-groupoids-lemma-abelian} 的证明相同）。设 $\kappa$ 为基数。
称拟凝聚模晶体 $\mathcal{F}$（在 $X$ 上）是\emph{$\kappa$-生成的}，
如果它的每个 $\mathcal{F}_i$ 都是 $\kappa$-生成的（见《性质》定义
\ref{properties-definition-kappa-generated}）。

\begin{lemma}
\label{spaces-groupoids-lemma-crystals-in-quasi-coherent-modules}
在上述情形中，如果所有态射 $f_\phi$ 都平坦，那么存在基数 $\kappa$，
使得每个对象
$(\{\mathcal{F}_i\}_{i \in I}, \{\alpha_\phi\}_{\phi \in \Phi})$
（属于 $\textit{CQC}(X)$）
都是其 $\kappa$-生成子模的有向余极限。
\end{lemma}

\begin{proof}
在本引理及其证明中，
$(\{\mathcal{F}_i\}_{i \in I}, \{\alpha_\phi\}_{\phi \in \Phi})$
的\emph{子模}，是指给定拟凝聚子模
$\mathcal{G}_i \subset \mathcal{F}_i$（对所有 $i$），使得
$\alpha_\phi(f_\phi^*\mathcal{G}_i) = \mathcal{G}_{i'}$
作为 $\mathcal{F}_{i'}$ 的子层，并且这对所有 $\phi \in \Phi$ 成立。
这个定义有意义，因为 $f_\phi$ 平坦时，
拉回 $f^*_\phi$ 是正合的，也就是说它保持子层。
% P09-ERR-0891
本证明是《性质》引理 \ref{properties-lemma-colimit-kappa} 的证明的一个变体。
建议读者先阅读该证明。

\medskip\noindent
我们声称，只需在所有概形 $X_i$ 都仿射的情形下证明本引理。为说明这一点，令
$$
J = \coprod\nolimits_{i \in I} \{U \subset X_i\text{ 为仿射开集}\}
$$
并令
% P09-ERR-0892
\begin{align*}
\Psi = & \coprod\nolimits_{\phi \in \Phi}
\{
(U, V) \mid
U \subset X_i, V \subset X_{i'}\text{ 为仿射开集且 } f_\phi(U) \subset V
\} \\
&
\amalg \coprod\nolimits_{i \in I}
\{
(U, U') \mid
U, U' \subset X_i\text{ 为仿射开集且 } U \subset U'
\}.
\end{align*}
给它配备显然的映射 $\Psi \to J \times J$。于是我们的
$(\mathcal{F}, \alpha)$ 诱导拟凝聚层晶体
$(\{\mathcal{H}_j\}_{j \in J}, \{\beta_\psi\}_{\psi \in \Psi})$
（在 $Y = (J, \Psi)$ 上）：置
$\mathcal{H}_{(i, U)} = \mathcal{F}_i|_U$（对 $(i, U) \in J$），并置
$\beta_\psi$（对 $\psi \in \Psi$）如下：
% P09-ERR-0893
它等于 $\alpha_\phi$ 在 $U$ 上的限制（若
$\psi = (\phi, U, V)$），并等于
$\text{id} : (\mathcal{F}_i|_{U'})|_U \to \mathcal{F}_i|_U$
（当 $\psi = (i, U, U')$）。此外，
$(\{\mathcal{H}_j\}_{j \in J}, \{\beta_\psi\}_{\psi \in \Psi})$
的子模按 $1$ 对 $1$ 地对应于
$(\{\mathcal{F}_i\}_{i \in I}, \{\alpha_\phi\}_{\phi \in \Phi})$
的子模。我们略去证明（提示：使用《层》第
\ref{sheaves-section-bases} 节）。此外，显然，如果 $\kappa$ 对 $Y$ 可行，
那么同一个 $\kappa$ 对 $X$ 也可行（由 $\kappa$-生成模的定义）。
% P09-ERR-0894
因此，只需对 $Y$ 上拟凝聚层晶体证明本引理。

\medskip\noindent
假设所有概形 $X_i$ 都仿射。取无限基数 $\kappa$，使其大于 $I$ 或
$\Phi$ 的基数。设
$(\{\mathcal{F}_i\}_{i \in I}, \{\alpha_\phi\}_{\phi \in \Phi})$
为 $\textit{CQC}(X)$ 的对象。对每个 $i$ 写成
$X_i = \Spec(A_i)$ 和 $M_i = \Gamma(X_i, \mathcal{F}_i)$。
对每个 $\phi \in \Phi$（满足 $j(\phi) = (i, i')$），映射
$\alpha_\phi$ 转化为 $A_{i'}$-模同构
$$
\alpha_\phi : M_i \otimes_{A_i} A_{i'} \longrightarrow M_{i'}
$$
利用选择公理，选取规则
% P09-ERR-0895
$$
(\phi, m) \longmapsto S(\phi, m')
$$
其中源是满足如下条件的偶 $(\phi, m')$ 的集合：
$\phi \in \Phi$、$j(\phi) = (i, i')$ 且 $m' \in M_{i'}$；其输出是
有限子集 $S(\phi, m') \subset M_i$，使得
$$
m' = \alpha_\phi\left(\sum\nolimits_{m \in S(\phi, m')} m \otimes a'_m\right)
$$
对某些 $a'_m \in A_{i'}$ 成立。

\medskip\noindent
作出这些选择后，我们声称，任意 $\mathcal{F}_i$ 在任意 $X_i$ 上的任意
截面都属于某个 $\kappa$-生成子模。为说明这一点，假设给定子集族
$\mathcal{S} = \{S_i\}_{i \in I}$，其中每个
$S_i \subset M_i$ 的基数至多为 $\kappa$。定义新子集族
$\mathcal{S}' = \{S'_i\}_{i \in I}$，其中
$$
S'_i = S_i \cup
\bigcup\nolimits_{(\phi, m'),\ j(\phi) = (i, i'),\ m' \in S_{i'}} S(\phi, m')
$$
注意每个 $S'_i$ 的基数仍至多为 $\kappa$。置
$\mathcal{S}^{(0)} = \mathcal{S}$、$\mathcal{S}^{(1)} = \mathcal{S}'$，
并递归地置
$\mathcal{S}^{(n + 1)} = (\mathcal{S}^{(n)})'$。再置
$S_i^{(\infty)} = \bigcup_{n \geq 0} S_i^{(n)}$ 和
$\mathcal{S}^{(\infty)} = \{S_i^{(\infty)}\}_{i \in I}$。
根据构造，对每个 $\phi \in \Phi$（满足 $j(\phi) = (i, i')$）及每个
$m' \in S^{(\infty)}_{i'}$，都能把 $m'$ 写成像
$\alpha_\phi(m \otimes 1)$ 的有限线性组合，其中
$m \in S_i^{(\infty)}$。因此，令 $N_i$ 为 $A_i$-子模；它是在
$M_i$ 中由 $S_i^{(\infty)}$ 生成的。则相应的拟凝聚子模
% P09-ERR-0896
$\widetilde{N_i} \subset \mathcal{F}_i$ 构成一个 $\kappa$-生成子模。
这就完成了证明。
\end{proof}

\begin{lemma}
\label{spaces-groupoids-lemma-set-generators}
设 $B \to S$ 如第 \ref{spaces-groupoids-section-notation} 节所示。设
$(U, R, s, t, c)$ 为代数空间中的群胚（在 $B$ 上）。如果 $s$、$t$
平坦，那么存在集合 $T$ 和一族对象
$(\mathcal{F}_t, \alpha_t)_{t \in T}$（属于 $\QCoh(U, R, s, t, c)$），使得每个对象
$(\mathcal{F}, \alpha)$ 都是其如下子模的有向余极限：这些子模同构于
对象 $(\mathcal{F}_t, \alpha_t)$ 中的某一个。
\end{lemma}

\begin{proof}
本引理推广了《群胚》引理 \ref{groupoids-lemma-colimit-kappa}；
后者处理概形中的群胚情形。我们不能直接使用同一个论证，因此使用上面发展的
“拟凝聚层晶体”材料。

\medskip\noindent
选取概形 $W$ 和满 \'etale 态射 $W \to U$。选取概形 $V$ 和满
\'etale 态射 $V \to W \times_{U, s} R$。选取概形 $V'$ 和满
\'etale 态射 $V' \to R \times_{t, U} W$。考虑概形族
$$
I = \{W, W \times_U W, V, V', V \times_R V'\}
$$
以及概形态射集
$$
\Phi = \{\text{pr}_i : W \times_U W \to W, V \to W, V' \to W,
V \times_R V' \to V, V \times_R V' \to V'\}
$$
置 $X = (I, \Phi)$。回忆我们已定义了范畴
$\textit{CQC}(X)$，即 $X$ 上拟凝聚层晶体的范畴。存在函子
$$
\QCoh(U, R, s, t, c) \longrightarrow \textit{CQC}(X)
$$
该函子把 $(\mathcal{F}, \alpha)$ 分别映到：层
$\mathcal{F}|_W$（在 $W$ 上）；层
$\mathcal{F}|_{W \times_U W}$（在 $W \times_U W$ 上）；
$\mathcal{F}$ 经由 $V \to W \times_{U, s} R \to W \to U$
的拉回（在 $V$ 上）；$\mathcal{F}$ 经由
$V' \to R \times_{t, U} W \to W \to U$ 的拉回（在 $V'$ 上）；
最后是 $\mathcal{F}$ 经由
$V \times_R V' \to V \to W \times_{U, s} R \to W \to U$
的拉回（在 $V \times_R V'$ 上）。比较映射
$\{\alpha_\phi\}_{\phi \in \Phi}$ 除一个以外都取显然的映射
（来自拉回的结合性）；对映射
$\phi = \text{pr}_{V'} : V \times_R V' \to V'$，则使用
$\alpha : t^*\mathcal{F} \to s^*\mathcal{F}$ 在
$V \times_R V'$ 上的拉回。下列交换图表说明这个定义有意义：
$$
\xymatrix{
& V \times_R V' \ar[ld] \ar[rd] \\
V \ar[rd] \ar[dd] & & V' \ar[ld] \ar[dd] \\
& R \ar@<-1ex>[dd]_s \ar@<1ex>[dd]^t \\
W \ar[rd] & & W \ar[ld] \\
& U
}
$$
上面显示的函子不是范畴等价。不过，因为 $W \to U$ 是满
\'etale 态射，它是忠实的\footnote{事实上该函子是全忠实的，但我们不需要这一点。}。
由于上图中的所有态射都平坦，可见它是 Abel 范畴间的正合函子。此外，我们声称，
% P09-ERR-0897
给定 $(\mathcal{F}, \alpha)$，其像为
$(\{\mathcal{F}_i\}_{i \in I}, \{\alpha_\phi\}_{\phi \in \Phi})$，
则存在 $1$ 对 $1$ 的对应：$(\mathcal{F}, \alpha)$ 的拟凝聚子模与
$(\{\mathcal{F}_i\}_{i \in I}, \{\alpha_\phi\}_{\phi \in \Phi})$
的子模相对应。具体而言，给定
$(\{\mathcal{F}_i\}_{i \in I}, \{\alpha_\phi\}_{\phi \in \Phi})$
的子模，它在 $W$ 上的部分与投影映射
$W \times_U W \to W$ 相容，这保证该子模来自 $\mathcal{F}$ 的一个
拟凝聚子模（由《空间的性质》命题
\ref{spaces-properties-proposition-quasi-coherent}）；它与
$\alpha_{\text{pr}_{V'}}$ 相容，这保证该子层与 $\alpha$ 相容
（细节略）。

\medskip\noindent
选取基数 $\kappa$，如引理
\ref{spaces-groupoids-lemma-crystals-in-quasi-coherent-modules} 对系统 $X = (I, \Phi)$ 所给。
由《性质》引理 \ref{properties-lemma-set-of-iso-classes} 可直接看出，
$\kappa$-生成的拟凝聚层晶体（在 $X$ 上）的同构类构成一个集合。因此结论显然。
\end{proof}






\section{群胚与群空间}
\label{spaces-groupoids-section-groupoids-group-spaces}

\noindent
请与《群胚》第 \ref{groupoids-section-groupoids-group-schemes} 节比较。

\begin{lemma}
\label{spaces-groupoids-lemma-groupoid-from-action}
设 $B \to S$ 如第 \ref{spaces-groupoids-section-notation} 节所示。设 $(G, m)$ 为
$B$ 上的群代数空间，其单位元为 $e_G$、逆元为 $i_G$。设 $X$ 为
$B$ 上的代数空间，并设 $a : G \times_B X \to X$ 为 $G$ 在
$X$ 上的作用（在 $B$ 上）。那么按如下方式得到代数空间中的群胚
$(U, R, s, t, c, e, i)$（在 $B$ 上）：
\begin{enumerate}
\item 置 $U = X$ 和 $R = G \times_B X$。
\item 置 $s : R \to U$ 等于 $(g, x) \mapsto x$。
\item 置 $t : R \to U$ 等于 $(g, x) \mapsto a(g, x)$。
\item 置 $c : R \times_{s, U, t} R \to R$ 等于
$((g, x), (g', x')) \mapsto (m(g, g'), x')$。
\item 置 $e : U \to R$ 等于 $x \mapsto (e_G(x), x)$。
\item 置 $i : R \to R$ 等于 $(g, x) \mapsto (i_G(g), a(g, x))$。
\end{enumerate}
\end{lemma}

\begin{proof}
略。提示：只需在集合层次证明上述构造可行。为此，使用引理上方的描述：
$g$ 是从 $v$ 到 $a(g, v)$ 的箭头。
\end{proof}

\begin{lemma}
\label{spaces-groupoids-lemma-action-groupoid-modules}
设 $B \to S$ 如第 \ref{spaces-groupoids-section-notation} 节所示。设 $(G, m)$ 为
$B$ 上的群代数空间。设 $X$ 为 $B$ 上的代数空间，并设
$a : G \times_B X \to X$ 为 $G$ 在 $X$ 上的作用（在 $B$ 上）。设
$(U, R, s, t, c)$ 为引理 \ref{spaces-groupoids-lemma-groupoid-from-action} 中构造的
代数空间中的群胚。规则
$(\mathcal{F}, \alpha) \mapsto (\mathcal{F}, \alpha)$ 给出如下两个
范畴之间的等价：$G$-等变 $\mathcal{O}_X$-模的范畴，与群胚
$(U, R, s, t, c)$ 上拟凝聚模的范畴。
\end{lemma}

\begin{proof}
该断言有意义，因为 $t = a$ 且 $s = \text{pr}_1$（作为态射
$R = G \times_B X \to X$）；见定义 \ref{spaces-groupoids-definition-equivariant-module}
和 \ref{spaces-groupoids-definition-groupoid-module}。利用引理
\ref{spaces-groupoids-lemma-groupoid-from-action} 中的转译，两个定义中的交换性要求
恰好相符。
\end{proof}






\section{稳定子群代数空间}
\label{spaces-groupoids-section-stabilizer}

\noindent
请与《群胚》第 \ref{groupoids-section-stabilizer} 节比较。给定代数空间中的群胚，
可如下得到一个群代数空间。

\begin{lemma}
\label{spaces-groupoids-lemma-groupoid-stabilizer}
设 $B \to S$ 如第 \ref{spaces-groupoids-section-notation} 节所示。设
$(U, R, s, t, c)$ 为代数空间中的群胚（在 $B$ 上）。代数空间 $G$ 由笛卡儿方块
$$
\xymatrix{
G \ar[r] \ar[d] & R \ar[d]^{j = (t, s)} \\
U \ar[r]^-{\Delta} & U \times_B U
}
$$
定义。它是 $U$ 上的群代数空间，其复合法则 $m$ 由复合法则
$c$ 诱导。
\end{lemma}

\begin{proof}
这是因为，在群胚范畴中，任一对象的自映射集构成一个群。
\end{proof}

\noindent
由于 $\Delta$ 是单态射，可知 $G = j^{-1}(\Delta_{U/B})$ 是 $R$ 的子层。
按这种方式看它时，结构态射 $G = j^{-1}(\Delta_{U/B}) \to U$ 由 $s$ 或
$t$ 诱导（两者相同），而 $m$ 由 $c$ 诱导。

\begin{definition}
\label{spaces-groupoids-definition-stabilizer-groupoid}
设 $B \to S$ 如第 \ref{spaces-groupoids-section-notation} 节所示。设
$(U, R, s, t, c)$ 为代数空间中的群胚（在 $B$ 上）。群代数空间
$j^{-1}(\Delta_{U/B}) \to U$ 称为emph{代数空间中的群胚
$(U, R, s, t, c)$ 的稳定子}。
\end{definition}

\noindent
文献中常用 $S$ 表示稳定子群代数空间（推测是因为 stabilizer 一词以 ``s''
开头）；我们不能这样做，因为已经用 $S$ 表示基概形。

\begin{lemma}
\label{spaces-groupoids-lemma-groupoid-action-stabilizer}
设 $B \to S$ 如第 \ref{spaces-groupoids-section-notation} 节所示。设
$(U, R, s, t, c)$ 为代数空间中的群胚（在 $B$ 上），并设 $G/U$ 为其稳定子。
用 $R_t/U$ 表示代数空间 $R$；它被看作 $U$ 上的代数空间，其结构态射为
$t : R \to U$。存在典范左作用
$$
a : G \times_U R_t \longrightarrow R_t
$$
它由复合法则 $c$ 诱导。
\end{lemma}

\begin{proof}
用 $T/B$ 上的点表示时，定义 $a(g, r) = c(g, r)$。
\end{proof}






\section{限制群胚}
\label{spaces-groupoids-section-restrict-groupoid}

\noindent
记号请参见《群胚》第 \ref{groupoids-section-restrict-groupoid} 节。

\begin{lemma}
\label{spaces-groupoids-lemma-restrict-groupoid}
设 $B \to S$ 如第 \ref{spaces-groupoids-section-notation} 节所示。设
$(U, R, s, t, c)$ 为代数空间中的群胚（在 $B$ 上）。设
$g : U' \to U$ 为代数空间的态射。考虑图表
$$
\xymatrix{
R' \ar[d] \ar[r] \ar@/_3pc/[dd]_{t'} \ar@/^1pc/[rr]^{s'}&
R \times_{s, U} U' \ar[r] \ar[d] &
U' \ar[d]^g \\
U' \times_{U, t} R \ar[d] \ar[r] &
R \ar[r]^s \ar[d]_t &
U \\
U' \ar[r]^g &
U
}
$$
其中所有方块都是纤维积方块。那么存在典范复合法则
$c' : R' \times_{s', U', t'} R' \to R'$，使得
$(U', R', s', t', c')$ 是代数空间中的群胚（在 $B$ 上），并且
$U' \to U$、$R' \to R$ 定义代数空间中的群胚的态射
$(U', R', s', t', c') \to (U, R, s, t, c)$（在 $B$ 上）。此外，对
任意概形 $T$（在 $B$ 上），群胚函子
$$
(U'(T), R'(T), s', t', c') \to (U(T), R(T), s, t, c)
$$
是 $(U(T), R(T), s, t, c)$ 经映射 $U'(T) \to U(T)$ 所作的限制
（见《群胚》第 \ref{groupoids-section-restrict-groupoid} 节）。
\end{lemma}

\begin{proof}
略。
\end{proof}

\begin{definition}
\label{spaces-groupoids-definition-restrict-groupoid}
设 $B \to S$ 如第 \ref{spaces-groupoids-section-notation} 节所示。设
$(U, R, s, t, c)$ 为代数空间中的群胚（在 $B$ 上）。设
$g : U' \to U$ 为 $B$ 上代数空间的态射。引理
\ref{spaces-groupoids-lemma-restrict-groupoid} 中构造的代数空间中的群胚的态射
$(U', R', s', t', c') \to (U, R, s, t, c)$ 称为
\emph{$(U, R, s, t, c)$ 到 $U'$ 的限制}。
% P09-ERR-0898
在这种情形下，有时使用记号 $R' = R|_{U'}$。
\end{definition}

\begin{lemma}
\label{spaces-groupoids-lemma-restrict-groupoid-relation}
定义 \ref{spaces-groupoids-definition-restrict-groupoid} 和
\ref{spaces-groupoids-definition-restrict-relation} 中所定义的限制群胚与限制（预）等价关系的概念，
通过引理 \ref{spaces-groupoids-lemma-groupoid-pre-equivalence} 和
\ref{spaces-groupoids-lemma-equivalence-groupoid} 的构造相符。
\end{lemma}

\begin{proof}
这里所说的是，引理 \ref{spaces-groupoids-lemma-restrict-groupoid} 中的 $R'$ 也等于
$$
R' = (U' \times_B U')\times_{U \times_B U} R
\longrightarrow
U' \times_B U'
$$
事实上，这也许是陈述该引理的一种更清楚的方式。
\end{proof}






\section{不变子空间}
\label{spaces-groupoids-section-invariant}

\noindent
本节简要讨论不变子空间的概念。

\begin{definition}
\label{spaces-groupoids-definition-invariant-open}
设 $B \to S$ 如第 \ref{spaces-groupoids-section-notation} 节所示。设
$(U, R, s, t, c)$ 为代数空间中的群胚（在基底 $B$ 上）。
\begin{enumerate}
\item 称开子空间 $W \subset U$ 是\emph{$R$-不变的}，如果
$t(s^{-1}(W)) \subset W$。
\item 若局部闭子空间 $Z \subset U$ 作为 $R$ 的局部闭子空间满足
$t^{-1}(Z) = s^{-1}(Z)$，则称它是\emph{$R$-不变的}。
\item 若代数空间的单态射 $T \to U$ 作为 $R$ 上的代数空间满足
$T \times_{U, t} R = R \times_{s, U} T$，则称它是\emph{$R$-不变的}。
\end{enumerate}
\end{definition}

\noindent
对于开子空间 $W \subset U$，$R$-不变性也等价于要求
$s^{-1}(W) = t^{-1}(W)$。如果 $W \subset U$ 是 $R$-不变的，那么
$R$ 到 $W$ 的限制就是 $R_W = s^{-1}(W) = t^{-1}(W)$。类似地，如果
$Z \subset U$ 是 $R$-不变局部闭子空间，那么 $R$ 到 $Z$ 的限制就是
$R_Z = s^{-1}(Z) = t^{-1}(Z)$。

\begin{lemma}
\label{spaces-groupoids-lemma-constructing-invariant-opens}
设 $B \to S$ 如第 \ref{spaces-groupoids-section-notation} 节所示。设
$(U, R, s, t, c)$ 为代数空间中的群胚（在 $B$ 上）。
\begin{enumerate}
\item 如果 $s$ 和 $t$ 都是开态射，那么对每个开集 $W \subset U$，开集
$s(t^{-1}(W))$ 是 $R$-不变的。
\item 如果 $s$ 和 $t$ 都开且拟紧，那么 $U$ 有一个开覆盖，其成员都是
$R$-不变拟紧开子空间。
\end{enumerate}
\end{lemma}

\begin{proof}
假设 $s$ 和 $t$ 都开，且 $W \subset U$ 开。由于 $s$ 开，可知
$W' = s(t^{-1}(W))$ 是 $U$ 的开子空间。
% P09-ERR-0899
从函子观点很容易看出这是 $R$-不变开子集（在 $U$ 中）；不过，我们认为直接使用一些
图表来论证颇有启发性。注意 $t^{-1}(W')$ 是态射
$$
A := t^{-1}(W) \times_{s|_{t^{-1}(W)}, U, t} R
\xrightarrow{\text{pr}_1} R
$$
的像，而 $s^{-1}(W')$ 是态射
$$
B := R \times_{s, U, s|_{t^{-1}(W)}} t^{-1}(W)
\xrightarrow{\text{pr}_0} R.
$$
上述箭头左侧的代数空间 $A$、$B$ 分别是
$R \times_{s, U, t} R$ 和 $R \times_{s, U, s} R$ 的开子空间。由引理
\ref{spaces-groupoids-lemma-diagram}，图表
$$
\xymatrix{
R \times_{s, U, t} R \ar[rd]_{\text{pr}_1} \ar[rr]_{(\text{pr}_1, c)} & &
R \times_{s, U, s} R \ar[ld]^{\text{pr}_0} \\
& R &
}
$$
交换，且水平箭头是同构。此外，显然
$(\text{pr}_1, c)(A) = B$。因此得到
$s^{-1}(W') = t^{-1}(W')$，并且 $W'$ 是 $R$-不变的。这证明了 (1)。

\medskip\noindent
现在假设 $s$、$t$ 都开且拟紧。那么，如果 $W \subset U$ 是拟紧开集，
$W' = s(t^{-1}(W))$ 也是拟紧开集，并且由上面的讨论可知它不变。让 $W$
遍历所有 \'etale 于 $U$ 的仿射概形的像，便得到 (2)。
\end{proof}






\section{商层}
\label{spaces-groupoids-section-quotient-sheaves}

\noindent
设 $S$ 为概形，$B$ 为 $S$ 上的代数空间。设
$j : R \to U \times_B U$ 为 $B$ 上的预关系。对每个概形 $S'$（在 $S$ 上），
可以取等价关系 $\sim_{S'}$，它由
$j(S') : R(S') \to U(S') \times U(S')$ 的像生成。于是得到预层
\begin{equation}
\label{spaces-groupoids-equation-quotient-presheaf}
\begin{matrix}
(\Sch/S)^{opp}_{fppf} &
\longrightarrow &
\textit{Sets}, \\
S' &
\longmapsto  &
U(S')/\sim_{S'}
\end{matrix}
\end{equation}
注意，因为 $j$ 是 $B$ 上代数空间的态射，且取值于 $U \times_B U$，
所以存在从预层（\ref{spaces-groupoids-equation-quotient-presheaf}）到 $B$ 的预层的典范变换。

\begin{definition}
\label{spaces-groupoids-definition-quotient-sheaf}
设 $B \to S$ 以及预关系 $j : R \to U \times_B U$ 如上。在此情形下，
\emph{商层 $U/R$}（与 $j$ 相伴）是预层
（\ref{spaces-groupoids-equation-quotient-presheaf}）在 $(\Sch/S)_{fppf}$ 上的层化。如果
$j : R \to U \times_B U$ 来自群代数空间 $G$（在 $B$ 上）对 $U$ 的作用，
如引理 \ref{spaces-groupoids-lemma-groupoid-from-action} 所述，那么把商层记为 $U/G$。
\end{definition}

\noindent
这恰好意味着图表
$$
\xymatrix{
R \ar@<1ex>[r] \ar@<-1ex>[r] &
U \ar[r] &
U/R
}
$$
是 $(\Sch/S)_{fppf}$ 上集合层范畴中的余等化子图表。仍然存在层的典范映射
$U/R \to B$，因为 $j$ 是 $B$ 上代数空间的态射，且取值于
$U \times_B U$。

\begin{remark}
\label{spaces-groupoids-remark-quotient-variant}
上述构造的一种变体是对函子
$$
\begin{matrix}
(\textit{Spaces}/B)^{opp}_{fppf} &
\longrightarrow &
\textit{Sets}, \\
X &
\longmapsto  &
U(X)/\sim_X
\end{matrix}
$$
作层化；这里 $\sim_X \subset U(X) \times U(X)$ 是由
$j : R(X) \to U(X) \times U(X)$ 的像生成的等价关系。当然，
$U(X) = \Mor_B(X, U)$ 且 $R(X) = \Mor_B(X, R)$。事实上，通过如下恒等，
所得结果相同：
% P09-ERR-0900
（在此插入《空间的拓扑》中的未来引用）。
\end{remark}

\begin{definition}
\label{spaces-groupoids-definition-representable-quotient}
% P09-ERR-0901
在定义 \ref{spaces-groupoids-definition-quotient-sheaf} 的情形下，称预关系 $j$ 有一个
\emph{可由代数空间表示的商}，如果层 $U/R$ 是代数空间。称预关系 $j$
有一个\emph{可表示商}，如果层 $U/R$ 可由概形表示。如果群胚
$(U, R, s, t, c)$ 是代数空间中的群胚（在 $B$ 上），并且取
商 $U/R$（其中 $j = (t, s)$）可表示（相应地，是代数空间），则称它有一个
\emph{可表示商}（相应地，\emph{可由代数空间表示的商}）。
\end{definition}

\noindent
如果商 $U/R$ 可由 $M$ 表示（表示对象或为概形，或为 $S$ 上的代数空间），
那么如上所见，它带有典范结构态射 $M \to B$。

\medskip\noindent
下一个引理刻画表示该商的 $M$。例如，当 $U \to M$ 平坦、有限表现且满，
并且 $R \cong U \times_M U$ 时，可以应用该引理。

\begin{lemma}
\label{spaces-groupoids-lemma-criterion-quotient-representable}
在定义 \ref{spaces-groupoids-definition-quotient-sheaf} 的情形下，假设存在代数空间
$M$（在 $S$ 上）和态射 $U \to M$，使得
\begin{enumerate}
\item 态射 $U \to M$ 等化 $s,t$，
\item 映射 $U \to M$ 是层的满射，并且
\item 诱导映射 $(t, s) : R \to U \times_M U$ 是层的满射。
\end{enumerate}
在这种情形下，$M$ 表示商层 $U/R$。
\end{lemma}

\begin{proof}
条件 (1) 说明 $U \to M$ 分解经过 $U/R$。条件 (2) 说明
$U/R \to M$ 作为层的映射是满射。条件 (3) 说明
$U/R \to M$ 作为层的映射是单射。因此引理成立。
\end{proof}

\noindent
如果不要求 $j$ 为预等价关系（例如，只要求它为预关系），下一个引理就是错的。

\begin{lemma}
\label{spaces-groupoids-lemma-quotient-pre-equivalence}
设 $S$ 为概形。设 $B$ 为 $S$ 上的代数空间。设
$j : R \to U \times_B U$ 为 $B$ 上的预等价关系。对概形 $S'$（在 $S$ 上）
和 $a, b \in U(S')$，下列条件等价：
\begin{enumerate}
\item $a$ 和 $b$ 映到 $(U/R)(S')$ 的同一个元素，并且
\item 存在 fppf 覆盖 $\{f_i : S_i \to S'\}$（覆盖 $S'$）和态射
$r_i : S_i \to R$，使得
$a \circ f_i = s \circ r_i$ 且 $b \circ f_i = t \circ r_i$。
\end{enumerate}
换言之，在此情形下，层的映射
$$
R \longrightarrow U \times_{U/R} U
$$
是满射。
\end{lemma}

\begin{proof}
略。提示：这个结论成立的原因是，在此情形下预层
（\ref{spaces-groupoids-equation-quotient-presheaf}）实际上由
$T \mapsto U(T)/j(R(T))$ 给出，因为
$j(R(T)) \subset U(T) \times U(T)$ 是等价关系；见定义
\ref{spaces-groupoids-definition-equivalence-relation}。
\end{proof}

\begin{lemma}
\label{spaces-groupoids-lemma-quotient-pre-equivalence-relation-restrict}
设 $S$ 为概形。设 $B$ 为 $S$ 上的代数空间。设
$j : R \to U \times_B U$ 为 $B$ 上的预等价关系，并设
$g : U' \to U$ 为 $B$ 上代数空间的态射。设
$j' : R' \to U' \times_B U'$ 为 $j$ 到 $U'$ 的限制。商层的映射
$$
U'/R' \longrightarrow U/R
$$
是单射。如果 $U' \to U$ 作为层的映射是满射，例如，如果
$\{g : U' \to U\}$ 是 fppf 覆盖（见《空间上的拓扑》定义
\ref{spaces-topologies-definition-fppf-covering}），那么
$U'/R' \to U/R$ 是层的同构。
\end{lemma}

\begin{proof}
假设截面 $\xi, \xi' \in (U'/R')(S')$ 映到 $U/R$ 的同一个截面。那么可以找到
fppf 覆盖 $\mathcal{S} = \{S_i \to S'\}$（覆盖 $S'$），使得
$\xi|_{S_i}, \xi'|_{S_i}$ 由 $a_i, a_i' \in U'(S_i)$ 给出。由引理
\ref{spaces-groupoids-lemma-quotient-pre-equivalence} 和位点公理，在作加细后可以假设
% P09-ERR-0902
$\mathcal{T}$ 存在态射 $r_i : S_i \to R$，使得
$g \circ a_i = s \circ r_i$、$g \circ a_i' = t \circ r_i$。由于根据构造
$R' = R \times_{U \times_S U} (U' \times_S U')$，可知
$(r_i, (a_i, a_i')) \in R'(S_i)$；这说明 $a_i$ 和 $a_i'$ 定义
$U'/R'$ 的同一个截面（在 $S_i$ 上）。由层条件，这推出 $\xi = \xi'$。

\medskip\noindent
如果 $U' \to U$ 是层的满射，那么 $U'/R' \to U/R$ 也是满射。最后，如果
$\{g : U' \to U\}$ 是 fppf 覆盖，那么层的映射 $U' \to U$ 是满射；见
《空间上的拓扑》引理 \ref{spaces-topologies-lemma-fppf-covering-surjective}。
\end{proof}

\begin{lemma}
\label{spaces-groupoids-lemma-quotient-groupoid-restrict}
设 $S$ 为概形。设 $B$ 为 $S$ 上的代数空间。设
$(U, R, s, t, c)$ 为代数空间中的群胚（在 $B$ 上）。
% P09-ERR-0904
设 $g : U' \to U$ 为 $B$ 上代数空间的态射。设
$(U', R', s', t', c')$ 为 $(U, R, s, t, c)$ 到 $U'$ 的限制。商层的映射
$$
U'/R' \longrightarrow U/R
$$
是单射。如果复合
$$
\xymatrix{
U' \times_{g, U, t} R \ar[r]_-{\text{pr}_1} \ar@/^3ex/[rr]^h
& R \ar[r]_s & U
}
$$
是 fppf 层的满射，那么该映射是双射。例如，当
$\{h : U' \times_{g, U, t} R \to U\}$ 是 $fppf$-覆盖，或
$U' \to U$ 是层的满射，或 $\{g : U' \to U\}$ 是 fppf 拓扑中的覆盖时，
这个条件成立。
\end{lemma}

\begin{proof}
单射性由引理 \ref{spaces-groupoids-lemma-groupoid-pre-equivalence} 和
\ref{spaces-groupoids-lemma-quotient-pre-equivalence-relation-restrict} 合并得到。为说明满射性
（层的满射映射的刻画见《位点》第 \ref{sites-section-sheaves-injective} 节），
按如下方式论证。假设 $T$ 为概形且 $\sigma \in U/R(T)$。存在覆盖
$\{T_i \to T\}$，使得 $\sigma|_{T_i}$ 是某个元素 $f_i \in U(T_i)$ 的像。
% P09-ERR-0905
因此，可以假设 $\sigma$ 是 $f \in U(T)$ 的像。由 $h$ 是层的满射这一假设，
可以找到 fppf 覆盖 $\{\varphi_i : T_i \to T\}$ 和态射
$f_i : T_i \to U' \times_{g, U, t} R$，使得
$f \circ \varphi_i = h \circ f_i$。记
$f'_i = \text{pr}_0 \circ f_i : T_i \to U'$。于是可见，
$f'_i \in U'(T_i)$ 映到 $g \circ f'_i \in U(T_i)$，并且
% P09-ERR-0906
$g \circ f'_i \sim_{T_i} h \circ f_i = f \circ \varphi_i$；记号如
（\ref{spaces-groupoids-equation-quotient-presheaf}）所示。给出该关系的 $R(T_i)$ 中元素是
$\text{pr}_1 \circ f_i$。这意味着，$\sigma$ 到 $T_i$ 的限制属于映射
% P09-ERR-0907
$U'/R'(T_i) \to U/R(T_i)$ 的像，正如所需。

\medskip\noindent
如果 $\{h\}$ 是 fppf 覆盖，那么它诱导层的满射；见《空间上的拓扑》引理
\ref{spaces-topologies-lemma-fppf-covering-surjective}。如果 $U' \to U$ 满，
那么 $h$ 也满，因为 $s$ 有一个截面
% P09-ERR-0908
（即群胚概形的中性元 $e$）。
\end{proof}






\section{商叠}
\label{spaces-groupoids-section-stacks}

\noindent
本节和随后的几节描述上面第 \ref{spaces-groupoids-section-quotient-sheaves} 节及《群胚》第
\ref{groupoids-section-quotient-sheaves} 节的一种推广。区别如下：我们将取商叠，
而不是商层。

\medskip\noindent
假设有概形 $S$、代数空间 $B$（在 $S$ 上），以及代数空间中的群胚
$(U, R, s, t, c)$（在 $B$ 上）。给定这些数据，考虑函子
\begin{equation}
\label{spaces-groupoids-equation-quotient-stack}
\begin{matrix}
(\Sch/S)_{fppf}^{opp} &
\longrightarrow &
\textit{Groupoids} \\
S' &
\longmapsto &
(U(S'), R(S'), s, t, c)
\end{matrix}
\end{equation}
由《范畴》例 \ref{categories-example-functor-groupoids}，这个“群胚预层”对应于
$(\Sch/S)_{fppf}$ 上的群胚纤维化范畴。在本章中把它记为
$$
[U/_{\!p}R] \to (\Sch/S)_{fppf}
$$
其中下标 ${}_p$ 用来与商叠相区别。

\begin{definition}
\label{spaces-groupoids-definition-quotient-stack}
商叠。设 $B \to S$ 如上。
\begin{enumerate}
\item 设 $(U, R, s, t, c)$ 为代数空间中的群胚（在 $B$ 上）。
\emph{商叠}
$$
p : [U/R] \longrightarrow (\Sch/S)_{fppf}
$$
属于 $(U, R, s, t, c)$；它是群胚纤维化范畴
$[U/_{\!p}R]$（在 $(\Sch/S)_{fppf}$ 上）的叠化（见《叠》引理
\ref{stacks-lemma-stackify-groupoids}），该范畴与
（\ref{spaces-groupoids-equation-quotient-stack}）相伴。
\item 设 $(G, m)$ 为 $B$ 上的群代数空间。设
% P09-ERR-0909
$a : G \times_B X \to X$ 为 $G$ 在 $B$ 上某个代数空间上的作用。
\emph{商叠}
$$
p : [X/G] \longrightarrow (\Sch/S)_{fppf}
$$
是与引理 \ref{spaces-groupoids-lemma-groupoid-from-action} 中代数空间中的群胚
$(X, G \times_B X, s, t, c)$（在 $B$ 上）相伴的商叠。
\end{enumerate}
\end{definition}

\noindent
因此 $[U/R]$ 和 $[X/G]$ 是 $(\Sch/S)_{fppf}$ 上的群胚叠。这些叠以后非常重要，
因而值得详细描述。回忆：给定代数空间 $X$（在 $S$ 上），我们用记号
$\mathcal{S}_X \to (\Sch/S)_{fppf}$ 表示与层 $X$ 相伴的集合叠；见《范畴》引理
\ref{categories-lemma-2-category-fibred-sets} 和《叠》引理
\ref{stacks-lemma-when-stack-in-sets}。

\begin{lemma}
\label{spaces-groupoids-lemma-quotient-stack-arrows}
假设 $B \to S$ 和 $(U, R, s, t, c)$ 如定义
\ref{spaces-groupoids-definition-quotient-stack} (1) 所示。存在典范 $1$-态射
$\pi : \mathcal{S}_U \to [U/R]$ 和 $[U/R] \to \mathcal{S}_B$，它们是
$(\Sch/S)_{fppf}$ 上群胚叠的态射。复合
$\mathcal{S}_U \to \mathcal{S}_B$ 是 $1$-态射，即与结构态射 $U \to B$ 相伴者。
\end{lemma}

\begin{proof}
在本证明中，用 $[U/_{\!p}R]$ 表示与群胚预层
（\ref{spaces-groupoids-equation-quotient-stack}）相伴的群胚纤维化范畴。由叠化的构造，存在
$1$-态射 $[U/_{\!p}R] \to [U/R]$。$1$-态射
$\mathcal{S}_U \to [U/R]$ 就是复合
$\mathcal{S}_U \to [U/_{\!p}R] \to [U/R]$；其中第一个箭头把概形
$S'/S$ 以及态射 $x : S' \to U$（在 $S$ 上）映到对象
$x \in U(S')$，它属于 $[U/_{\!p}R]$ 在 $S'$ 上的纤维范畴。

\medskip\noindent
为构造 $1$-态射 $[U/R] \to \mathcal{S}_B$，只需构造 $1$-态射
$[U/_{\!p}R] \to \mathcal{S}_B$；见《叠》引理
\ref{stacks-lemma-stackify-groupoids-universal-property}。在 $S'/S$ 上的对象处，
只使用映射
$$
U(S') \longrightarrow B(S')
$$
它来自结构态射 $U \to B$。显然，如果 $a \in R(S')$ 是一个“箭头”，其源为
$s(a) \in U(S')$、靶为 $t(a) \in U(S')$，那么因为 $s$ 和 $t$ 都是
\emph{$B$ 上}的态射，两者都映到同一个元素 $\overline{a}$（属于 $B(S')$）。
因此可将箭头 $a \in R(S')$ 映到 $\overline{a}$ 的恒等态射。（这是合适的，
因为纤维范畴 $(\mathcal{S}_B)_{S'}$ 只包含恒等态射。）我们略去验证这条规则与
这些分裂纤维化范畴上的拉回相容，因而如所需定义了 $1$-态射
$[U/_{\!p}R] \to \mathcal{S}_B$。

\medskip\noindent
略去最后一个断言的验证。
\end{proof}

\begin{lemma}
\label{spaces-groupoids-lemma-quotient-stack-2-arrow}
假设与记号如引理 \ref{spaces-groupoids-lemma-quotient-stack-arrows} 所示。存在典范 $2$-态射
$\alpha : \pi \circ s \to \pi \circ t$，它使图表
$$
\xymatrix{
\mathcal{S}_R \ar[r]_s \ar[d]_t & \mathcal{S}_U \ar[d]^\pi \\
\mathcal{S}_U \ar[r]^-\pi & [U/R]
}
$$
$2$-交换。
\end{lemma}

\begin{proof}
设 $S'$ 为 $S$ 上的概形。设 $r : S' \to R$ 为 $S$ 上的态射。那么
$r \in R(S')$ 是对象 $s \circ r, t \circ r \in U(S')$ 之间的同构。此外，
这个构造与拉回相容。这给出典范 $2$-态射
$\alpha_p : \pi_p \circ s \to \pi_p \circ t$，其中
$\pi_p : \mathcal{S}_U \to [U/_{\!p}R]$ 如引理
\ref{spaces-groupoids-lemma-quotient-stack-arrows} 的证明所示。因此，甚至图表
$$
\xymatrix{
\mathcal{S}_R \ar[r]_s \ar[d]^t & \mathcal{S}_U \ar[d]^{\pi_p} \\
\mathcal{S}_U \ar[r]^-{\pi_p} & [U/_{\!p}R]
}
$$
也是 $2$-交换的。因此更不必说，引理中的图表是 $2$-交换的。
\end{proof}

\begin{remark}
\label{spaces-groupoids-remark-fundamental-square}
在以后的各章中，我们将使用一种有歧义的记号：与其用
$\mathcal{S}_X$ 表示与 $X$ 相伴的集合叠，不如直接写成 $X$。使用这个记号，
引理 \ref{spaces-groupoids-lemma-quotient-stack-2-arrow} 的图表变成大家熟悉的图表
$$
\xymatrix{
R \ar[r]_s \ar[d]_t & U \ar[d]^\pi \\
U \ar[r]^-\pi & [U/R]
}
$$
在随后各节中，我们将证明这个图表有许多良好性质。特别地，将证明它是
$2$-纤维积（第 \ref{spaces-groupoids-section-quotient-stack-2-cartesian} 节），并且它几乎是
$2$-余等化子（对 $s$ 和 $t$ 而言；见第
\ref{spaces-groupoids-section-quotient-stacks-2-coequalize} 节）。
\end{remark}





\section{商叠的函子性}
\label{spaces-groupoids-section-functoriality-quotient-stacks}

\noindent
代数空间中的群胚的态射给出相伴的商叠态射。

\begin{lemma}
\label{spaces-groupoids-lemma-quotient-stack-functorial}
设 $S$ 为概形。设 $B$ 为 $S$ 上的代数空间。设
$f : (U, R, s, t, c) \to (U', R', s', t', c')$ 为代数空间中的群胚的态射
（在 $B$ 上）。那么 $f$ 诱导商叠的典范 $1$-态射
$$
[f] : [U/R] \longrightarrow [U'/R'].
$$
\end{lemma}

\begin{proof}
用 $[U/_{\!p}R]$ 和 $[U'/_{\!p}R']$ 表示与函子
（\ref{spaces-groupoids-equation-quotient-stack}）相伴的群胚纤维化范畴；它们取在基位点
$(\Sch/S)_{fppf}$ 上。显然，$f$ 定义 $1$-态射
$[U/_{\!p}R] \to [U'/_{\!p}R']$；把它与 $[U'/R']$ 的叠化函子复合，得到
$[U/_{\!p}R] \to [U'/R']$。
% P09-ERR-0910
再由叠化函子 $[U/_{\!p}R] \to [U/R]$ 的泛性质（见《叠》引理
\ref{stacks-lemma-stackify-groupoids-universal-property}），得到
$[U/R] \to [U'/R']$。
\end{proof}

\noindent
设 $B \to S$ 和 $f : (U, R, s, t, c) \to (U', R', s', t', c')$ 如引理
\ref{spaces-groupoids-lemma-quotient-stack-functorial} 所示。在此情形下，如下定义第三个代数空间中的
群胚（在 $B$ 上）；我们使用 $T$-值点的语言，其中 $T$ 是 $B$ 上变动的概形：
\begin{enumerate}
\item $U'' = U \times_{f, U', t'} R'$，所以一个 $T$-值点是偶
$(u, r')$，满足 $f(u) = t'(r')$；
\item $R'' = R \times_{f \circ s, U', t'} R'$，所以一个 $T$-值点是偶
$(r, r')$，满足 $f(s(r)) = t'(r')$；
\item $s'' : R'' \to U''$ 由 $s''(r, r') = (s(r), r')$ 给出；
\item $t'' : R'' \to U''$ 由 $t''(r, r') = (t(r), c'(f(r), r'))$ 给出；
\item $c'' : R'' \times_{s'', U'', t''} R'' \to R''$ 由
$c''((r_1, r'_1), (r_2, r'_2)) = (c(r_1, r_2), r'_2)$ 给出。
\end{enumerate}
$c''$ 的公式有意义，因为 $s''(r_1, r'_1) = t''(r_2, r'_2)$。显然
$c''$ 满足结合律。恒等元 $e''$ 由
% P09-ERR-0914
$e''(u, r) = (e(u), r)$ 给出。$(r, r')$ 的逆元由
$(i(r), c'(f(r), r'))$ 给出。因此，确实得到代数空间中的群胚（在 $B$ 上）。

\medskip\noindent
显然，映射 $U'' \to U$ 和 $R'' \to R$ 定义代数空间中的群胚的态射
$g : (U'', R'', s'', t'', c'') \to (U, R, s, t, c)$（在 $B$ 上）。此外，
映射 $U'' \to U'$、$(u, r') \mapsto s'(r')$ 以及
$R'' \to U'$、$(r, r') \mapsto s'(r')$ 表明，事实上
$(U'', R'', s'', t'', c'')$ 是代数空间中的群胚（在 $U'$ 上）。

\begin{lemma}
\label{spaces-groupoids-lemma-cartesian-square-of-morphism}
记号与假设如引理 \ref{spaces-groupoids-lemma-quotient-stack-functorial} 所示。设
$(U'', R'', s'', t'', c'')$ 为上面构造的代数空间中的群胚（在 $B$ 上）。
存在 $2$-交换方块
$$
\xymatrix{
[U''/R''] \ar[d] \ar[r]_{[g]} & [U/R] \ar[d]^{[f]} \\
\mathcal{S}_{U'} \ar[r] & [U'/R']
}
$$
它把 $[U''/R'']$ 识别为该 $2$-纤维积。
\end{lemma}

\begin{proof}
映射 $[f]$ 和 $[g]$ 来自应用引理
\ref{spaces-groupoids-lemma-quotient-stack-functorial}，另两个映射来自引理
\ref{spaces-groupoids-lemma-quotient-stack-arrows}（以及
$(U'', R'', s'', t'', c'')$ 取在 $U'$ 上这一事实）。为证明 $2$-纤维积性质，
只需对群胚纤维化范畴的图表
$$
\xymatrix{
[U''/_{\!p}R''] \ar[d] \ar[r]_{[g]} & [U/_{\!p}R] \ar[d]^{[f]} \\
\mathcal{S}_{U'} \ar[r] & [U'/_{\!p}R']
}
$$
证明引理；见《叠》引理
\ref{stacks-lemma-stackification-fibre-product-categories-fibred-in-groupoids}。
换言之，只需证明 $2$-纤维积的一个对象
% P09-ERR-0911
$\mathcal{S}_U \times_{[U'/_{\!p}R']} [U/_{\!p}R]$（在 $T$ 上）对应于
一个 $T$-值点（属于 $U''$），并对态射作类似证明。当然，这恰好是我们一开始构造
$U''$ 和 $R''$ 的方式。

\medskip\noindent
具体地，$\mathcal{S}_U \times_{[U'/_{\!p}R']} [U/_{\!p}R]$ 在 $T$ 上的
一个对象是三元组 $(u', u, r')$，其中 $u'$ 是 $T$-值点（属于 $U'$），$u$ 是
$T$-值点（属于 $U$），而 $r'$ 是从 $u'$ 到 $f(u)$ 的态射
% P09-ERR-0915
（在 $[U'/R']_T$ 中）；即
% P09-ERR-0912
$r'$ 是一个 $T$-值点（属于 $R$），满足 $s'(r') = u'$ 且
$t'(r') = f(u)$。显然，可以忘掉 $u'$ 而不丢失信息，并且可见这些对象与
% P09-ERR-0913
$T$-值点（属于 $R''$）一一对应。

\medskip\noindent
态射的情形类似：设 $(u'_1, u_1, r'_1)$ 和 $(u'_2, u_2, r'_2)$ 是该纤维积
在 $T$ 上的两个对象。那么，从 $(u'_2, u_2, r'_2)$ 到
$(u'_1, u_1, r'_1)$ 的态射由 $(1, r)$ 给出，其中
$1 : u'_1 \to u'_2$ 仅仅意味着 $u'_1 = u'_2$（这是因为
$\mathcal{S}_U$ 纤维化于集合），而 $r$ 是 $T$-值点（属于 $R$），满足
$s(r) = u_2$、$t(r) = u_1$，并且还有
$c'(f(r), r'_2) = r'_1$。因此箭头
$$
(1, r) : (u'_2, u_2, r'_2) \to (u'_1, u_1, r'_1)
$$
完全由偶 $(r, r'_2)$ 决定。于是箭头函子由 $R''$ 表示，而且态射
$s''$、$t''$、$c''$ 显然分别对应于 $2$-纤维积
$\mathcal{S}_U \times_{[U'/_{\!p}R']} [U/_{\!p}R]$ 中的源、靶与复合。
\end{proof}






\section{商叠的 2-笛卡儿方块}
\label{spaces-groupoids-section-quotient-stack-2-cartesian}

\noindent
本节计算商叠的 $\mathit{Isom}$-层，并由此推出商叠的定义图表是
$2$-纤维积。

\begin{lemma}
\label{spaces-groupoids-lemma-quotient-stack-morphisms}
假设 $B \to S$、$(U, R, s, t, c)$ 和
$\pi : \mathcal{S}_U \to [U/R]$ 如引理
\ref{spaces-groupoids-lemma-quotient-stack-arrows} 所示。设 $S'$ 为 $S$ 上的概形。设
$x, y \in \Ob([U/R]_{S'})$ 为商叠在 $S'$ 上的对象。如果
$x = \pi(x')$ 且 $y = \pi(y')$，其中 $x', y' : S' \to U$ 是某些态射，那么
$$
\mathit{Isom}(x, y) = S' \times_{(y', x'), U \times_S U} R
$$
作为 $S'$ 上的层。
\end{lemma}

\begin{proof}
设 $[U/_{\!p}R]$ 为与群胚预层（\ref{spaces-groupoids-equation-quotient-stack}）相伴的
群胚纤维化范畴，如引理 \ref{spaces-groupoids-lemma-quotient-stack-arrows} 的证明所述。由构造，
层 $\mathit{Isom}(x, y)$ 是与预层 $\mathit{Isom}(x', y')$ 相伴的层。另一方面，
由 $[U/_{\!p}R]$ 中态射的定义，有
$$
\mathit{Isom}(x', y') = S' \times_{(y', x'), U \times_S U} R
$$
而右侧是代数空间，因而是层。
\end{proof}

\begin{lemma}
\label{spaces-groupoids-lemma-quotient-stack-2-cartesian}
假设 $B \to S$、$(U, R, s, t, c)$ 和
$\pi : \mathcal{S}_U \to [U/R]$ 如引理
\ref{spaces-groupoids-lemma-quotient-stack-arrows} 所示。引理
\ref{spaces-groupoids-lemma-quotient-stack-2-arrow} 的 $2$-交换方块
$$
\xymatrix{
\mathcal{S}_R \ar[r]_s \ar[d]_t & \mathcal{S}_U \ar[d]^\pi \\
\mathcal{S}_U \ar[r]^-\pi & [U/R]
}
$$
是一个 $2$-纤维积（在 $(\Sch/S)_{fppf}$ 的群胚叠中）。
\end{lemma}

\begin{proof}
由《叠》引理 \ref{stacks-lemma-2-product-stacks-in-groupoids}，该引理的断言
有意义。它还说明，我们必须证明函子
$$
\mathcal{S}_R \longrightarrow \mathcal{S}_U \times_{[U/R]} \mathcal{S}_U
$$
是等价；该函子把 $r : T \to R$ 映到
$(T, t(r), s(r), \alpha(r))$，其中右侧是《范畴》引理
\ref{categories-lemma-2-product-categories-over-C} 中所述的 $2$-纤维积。
把定义展开以后，这恰好是引理
\ref{spaces-groupoids-lemma-quotient-stack-morphisms} 的内容。（另一种证明：
% P09-ERR-0916
在此情形下展开引理 \ref{spaces-groupoids-lemma-cartesian-square-of-morphism} 的含义，
也会得到该结论。）
\end{proof}

\begin{lemma}
\label{spaces-groupoids-lemma-quotient-stack-isom}
假设 $B \to S$ 和 $(U, R, s, t, c)$ 如定义
\ref{spaces-groupoids-definition-quotient-stack} (1) 所示。对任意概形 $T$（在 $S$ 上）以及
对象 $x,y$（属于 $[U/R]$，在 $T$ 上），层
$\mathit{Isom}(x, y)$（在 $(\Sch/T)_{fppf}$ 上）具有如下性质：存在 fppf 覆盖
$\{T_i \to T\}_{i \in I}$，使得
$\mathit{Isom}(x, y)|_{(\Sch/T_i)_{fppf}}$ 可由代数空间表示。
\end{lemma}

\begin{proof}
这直接由引理 \ref{spaces-groupoids-lemma-quotient-stack-morphisms} 以及如下事实得到：根据商叠的
构造，$x$ 和 $y$ 在 fppf 拓扑下局部地都来自 $\mathcal{S}_U$ 的对象。
\end{proof}






\section{商叠的 2-余等化子性质}
\label{spaces-groupoids-section-quotient-stacks-2-coequalize}

\noindent
在群胚上有复合，它为上一个引理中的典范 $2$-态射导出上闭链条件。为了给出精确
表述，我们使用《范畴》第 \ref{categories-section-formal-cat-cat} 节和第
\ref{categories-section-2-categories} 节中引入的记号。

\begin{lemma}
\label{spaces-groupoids-lemma-quotient-stack-cocycle}
假设与记号如引理 \ref{spaces-groupoids-lemma-quotient-stack-arrows} 和
\ref{spaces-groupoids-lemma-quotient-stack-2-arrow} 所示。下图中两个 2-态射的竖直复合
$$
\xymatrix@C=15pc{
\mathcal{S}_{R \times_{s, U, t} R}
\ruppertwocell^{\pi \circ s \circ \text{pr}_1 = \pi \circ s \circ c}{\ \ \ \ \ \ \alpha \star \text{id}_{\text{pr}_1}}
\ar[r]_(.3){\pi \circ t \circ \text{pr}_1 = \pi \circ s \circ \text{pr}_0}
\rlowertwocell_{\pi \circ t \circ \text{pr}_0 = \pi \circ t \circ c}{\ \ \ \ \ \ \alpha \star \text{id}_{\text{pr}_0}}
&
[U/R]
}
$$
是 $2$-态射 $\alpha \star \text{id}_c$。用公式表示，
$\alpha \star \text{id}_c =
(\alpha \star \text{id}_{\text{pr}_0})
\circ
(\alpha \star \text{id}_{\text{pr}_1})
$。
\end{lemma}

\begin{proof}
作两点说明：
\begin{enumerate}
\item 公式
$\alpha \star \text{id}_c = (\alpha \star \text{id}_{\text{pr}_0}) \circ
(\alpha \star \text{id}_{\text{pr}_1})$ 仅在实现下列\emph{等式}后才有意义：
$\pi \circ s \circ \text{pr}_1 = \pi \circ s \circ c$、
$\pi \circ t \circ \text{pr}_1 = \pi \circ s \circ \text{pr}_0$ 和
$\pi \circ t \circ \text{pr}_0 = \pi \circ t \circ c$。确切地说，第二个等式
保证竖直复合 $\circ$ 有意义，另两个等式则保证公式两侧都是有相同源与靶的
$2$-态射。
\item 引理成立的原因是，与群胚预层（\ref{spaces-groupoids-equation-quotient-stack}）相伴的
群胚纤维化范畴 $[U/_{\!p}R]$ 中的复合来自复合法则
$c : R \times_{s, U, t} R \to R$。
\end{enumerate}
略去引理的证明。
\end{proof}

\noindent
注意，在引理的情形下，实际上已有等式
$s \circ \text{pr}_1 = s \circ c$、
$t \circ \text{pr}_1 = s \circ \text{pr}_0$ 和
$t \circ \text{pr}_0 = t \circ c$；这些等式是在与 $\pi$ 复合之前成立的。因此，下一个引理中的公式与上一个引理中的
公式以完全相同的方式有意义。

\begin{lemma}
\label{spaces-groupoids-lemma-quotient-stack-2-coequalizer}
假设与记号如引理 \ref{spaces-groupoids-lemma-quotient-stack-arrows} 和
\ref{spaces-groupoids-lemma-quotient-stack-2-arrow} 所示。引理
\ref{spaces-groupoids-lemma-quotient-stack-2-arrow} 的 $2$-交换图表在如下意义下是
$2$-余等化子：给定
\begin{enumerate}
\item 群胚叠 $\mathcal{X}$（在 $(\Sch/S)_{fppf}$ 上）；
\item $1$-态射 $f : \mathcal{S}_U \to \mathcal{X}$；以及
\item $2$-箭头 $\beta : f \circ s \to f \circ t$，
\end{enumerate}
使得
$$
\beta \star \text{id}_c
=
(\beta \star \text{id}_{\text{pr}_0})
\circ
(\beta \star \text{id}_{\text{pr}_1})
$$
那么存在 $1$-态射 $[U/R] \to \mathcal{X}$，它使图表
$$
\xymatrix{
\mathcal{S}_R \ar[r]_s \ar[d]^t & \mathcal{S}_U \ar[d] \ar[ddr]^f \\
\mathcal{S}_U \ar[r] \ar[rrd]_f & [U/R] \ar[rd] \\
& & \mathcal{X}
}
$$
$2$-交换。
\end{lemma}

\begin{proof}
假设给定引理中的 $\mathcal{X}$、$f$ 和 $\beta$。由《叠》引理
\ref{stacks-lemma-stackify-groupoids-universal-property}，只需构造 $1$-态射
$g : [U/_{\!p}R] \to \mathcal{X}$。首先注意，$1$-态射
$\mathcal{S}_U \to [U/_{\!p}R]$ 在对象上是双射。因此在对象上可以置
$g(x) = f(x)$，其中
$x \in \Ob(\mathcal{S}_U) = \Ob([U/_{\!p}R])$。态射
$\varphi : x \to y$（属于 $[U/_{\!p}R]$）来自交换图表
$$
\xymatrix{
S_2 \ar[dd]_h \ar[r]_x \ar[dr]_\varphi & U \\
& R \ar[u]_s \ar[d]^t \\
S_1 \ar[r]^y & U.
}
$$
于是可以把 $g(\varphi)$ 置为复合
$$
\xymatrix{
f(x) \ar@{=}[r] \ar[rrrrrd] &
f(s \circ \varphi) \ar@{=}[r] &
(f \circ s)(\varphi) \ar[r]^\beta &
(f \circ t)(\varphi) \ar@{=}[r] &
f(t \circ \varphi) \ar@{=}[r] &
f(y \circ h) \ar[d] \\
& & & & & f(y).
}
$$
竖直箭头是把函子 $f$ 应用于典范态射
$y \circ h \to y$（在 $\mathcal{S}_U$ 中）所得的结果（即提升 $h$ 且以
$y$ 为靶的强笛卡儿态射）。
% P09-ERR-0917
我们验证如此定义的 $f$ 与复合相容，至少在纤维范畴上如此。设 $S'$ 为 $S$ 上的
概形，并设 $a : S' \to R \times_{s, U, t} R$ 为态射。在此情形下，置
$x = s \circ \text{pr}_1 \circ a = s \circ c \circ a$、
$y = t \circ \text{pr}_1 \circ a = s \circ \text{pr}_0 \circ a$，以及
% P09-ERR-0918
$z = t \circ \text{pr}_0 \circ a = t \circ \text{pr}_0 \circ c$，从而得到交换图表
$$
\xymatrix{
x \ar[rr]_{c \circ a} \ar[rd]_{\text{pr}_1 \circ a} & & z \\
& y \ar[ru]_{\text{pr}_0 \circ a}
}
$$
它位于纤维范畴 $[U/_{\!p}R]_{S'}$ 中。
此外，该纤维范畴中的每个交换三角形都具有这种形式。于是由上面的定义可知，
$f$ 把它映到交换图表，当且仅当图表
$$
\xymatrix{
& (f \circ s)(c \circ a) \ar[r]_-{\beta} &
(f \circ t)(c \circ a) \ar@{=}[rd] & \\
(f \circ s)(\text{pr}_1 \circ a) \ar[rd]^\beta \ar@{=}[ru] & & &
(f \circ t)(\text{pr}_0 \circ a) \\
& (f \circ t)(\text{pr}_1 \circ a) \ar@{=}[r] &
(f \circ s)(\text{pr}_0 \circ a) \ar[ru]^\beta
}
$$
交换；这恰好是引理中公式所表达的条件。略去验证 $f$ 把恒等态射映到恒等态射，
并且与任意态射的复合相容。
\end{proof}

\section{商叠的显式描述}
\label{spaces-groupoids-section-explicit-quotient-stacks}

\noindent
为陈述结果，我们需要引入一些记号。假设 $B \to S$ 以及
$(U, R, s, t, c)$ 如定义 \ref{spaces-groupoids-definition-quotient-stack} (1) 中所述。
设 $T$ 为 $S$ 上的概形。设
$\mathcal{T} = \{T_i \to T\}_{i \in I}$ 为一个 fppf 覆盖。一个
{\it $[U/R]$-下降资料}（相对于 $\mathcal{T}$）由系统
$(u_i, r_{ij})$ 给出，其中
\begin{enumerate}
\item 对每个 $i$，有一个态射 $u_i : T_i \to U$，并且
\item 对每个 $i, j$，有一个态射 $r_{ij} : T_i \times_T T_j \to R$，
\end{enumerate}
使得
\begin{enumerate}
\item[(a)] 作为态射 $T_i \times_T T_j \to U$，有
$$
s \circ r_{ij} = u_i \circ \text{pr}_0
\quad\text{和}\quad
t \circ r_{ij} = u_j \circ \text{pr}_1,
$$
\item[(b)] 作为态射 $T_i \times_T T_j \times_T T_k \to R$，有
$$
c \circ (r_{jk} \circ \text{pr}_{12}, r_{ij} \circ \text{pr}_{01})
= r_{ik} \circ \text{pr}_{02}.
$$
\end{enumerate}
{\it 态射 $(u_i, r_{ij}) \to (u'_i, r'_{ij})$} 是两个
$[U/R]$-下降资料在同一个覆盖 $\mathcal{T}$ 上的态射，指一个族
$(r_i : T_i \to R)$，使得
\begin{enumerate}
\item[] $(\alpha)$ 作为态射 $T_i \to U$，有
$$
u_i = s \circ r_i
\quad\text{和}\quad
u'_i = t \circ r_i
$$
\item[] $(\beta)$ 作为态射 $T_i \times_T T_j \to R$，有
$$
c \circ (r'_{ij}, r_i \circ \text{pr}_0)
=
c \circ (r_j \circ \text{pr}_1, r_{ij}).
$$
\end{enumerate}
相对于一个固定覆盖的下降资料态射上有自然的复合法则，由此得到下降资料范畴。
这个范畴是一个群胚。最后，若
$\mathcal{T}' = \{T'_j \to T\}_{j \in J}$
是加细 $\mathcal{T}$ 的第二个 fppf 覆盖，则下降资料存在拉回的概念。在这种情形下，
它尤其容易显式描述。确切地说，若 $\alpha : J \to I$ 且
% P09-ERR-0919
$\varphi_j : T'_j \to T_{\alpha(i)}$ 是覆盖的态射，则下降资料
$(u_i, r_{ii'})$ 的拉回就是
$$
(u_{\alpha(i)} \circ \varphi_j,
r_{\alpha(j)\alpha(j')} \circ \varphi_j \times \varphi_{j'}).
$$
% P09-ERR-0920
以这种方式定义的拉回给出一个从 $\mathcal{T}$ 上的下降资料范畴到
$\mathcal{T}'$ 上的下降资料范畴的函子。

\begin{lemma}
\label{spaces-groupoids-lemma-quotient-stack-objects}
假设 $B \to S$ 以及 $(U, R, s, t, c)$ 如
定义 \ref{spaces-groupoids-definition-quotient-stack} (1) 中所述。
设 $\pi : \mathcal{S}_U \to [U/R]$ 如
引理 \ref{spaces-groupoids-lemma-quotient-stack-arrows} 中所述。
设 $T$ 为 $S$ 上的概形。
\begin{enumerate}
\item 对每个对象 $x$，若它属于纤维范畴 $[U/R]_T$，则存在一个 fppf 覆盖
$\{f_i : T_i \to T\}_{i \in I}$，使得
$f_i^*x \cong \pi(u_i)$，其中 $u_i \in U(T_i)$；
\item 下列同构的复合
$$
\pi(u_i \circ \text{pr}_0)
=
\text{pr}_0^*\pi(u_i)
\cong
\text{pr}_0^*f_i^*x
\cong
\text{pr}_1^*f_j^*x
\cong
\text{pr}_1^*\pi(u_j)
=
\pi(u_j \circ \text{pr}_1)
$$
具有 $\pi(r_{ij})$ 的形式，其中
$r_{ij} : T_i \times_T T_j \to R$ 是某些态射；
\item 系统 $(u_i, r_{ij})$ 构成如上定义的一个 $[U/R]$-下降资料；
\item 任意 $[U/R]$-下降资料 $(u_i, r_{ij})$ 都以这种方式产生；
\item 若 $x$ 如上对应于 $(u_i, r_{ij})$，而
$y \in \Ob([U/R]_T)$ 对应于 $(u'_i, r'_{ij})$，则有典范双射
$$
\Mor_{[U/R]_T}(x, y)
\longleftrightarrow
\left\{
\begin{matrix}
\text{态射 }(u_i, r_{ij}) \to (u'_i, r'_{ij})\\
\text{作为 }[U/R]\text{-下降资料}
\end{matrix}
\right\}
$$
\item 这个对应与 fppf 覆盖的加细相容。
\end{enumerate}
\end{lemma}

\begin{proof}
% P09-ERR-0921
断言 (1) 是叠化构造的一部分。
第 (2) 部分由引理 \ref{spaces-groupoids-lemma-quotient-stack-morphisms} 得出。
略去对 (3) 的验证。第 (4) 部分是叠中所有下降资料均有效这一事实的改述。
略去对 (5) 和 (6) 的验证。
\end{proof}

\section{限制与商叠}
\label{spaces-groupoids-section-quotient-stack-restrict}

\noindent
本节研究取限制时商叠会发生什么。

\begin{lemma}
\label{spaces-groupoids-lemma-quotient-stack-restrict}
记号与假设如引理 \ref{spaces-groupoids-lemma-quotient-stack-functorial} 中所述。
商叠的态射
$$
[f] : [U/R] \longrightarrow [U'/R']
$$
全忠实，当且仅当 $R$ 是 $R'$ 通过态射 $f : U \to U'$ 得到的限制。
\end{lemma}

\begin{proof}
设 $x, y$ 是 $[U/R]$ 在概形 $T/S$ 上的对象。设 $x', y'$ 是 $x, y$
在范畴 $[U'/R']_T$ 中的像。函子 $[f]$ 全忠实，当且仅当层的态射
$$
\mathit{Isom}(x, y) \longrightarrow \mathit{Isom}(x', y')
$$
对每个 $T, x, y$ 都是同构。可以在 $T$ 上局部地（就 fppf 拓扑而言）检验这一点。
因此，由引理 \ref{spaces-groupoids-lemma-quotient-stack-objects}，可以假设 $x, y$ 来自
$a, b \in U(T)$。在这种情形下，可见 $x', y'$ 对应于
$f \circ a, f \circ b$。由引理 \ref{spaces-groupoids-lemma-quotient-stack-morphisms}，
在这种情形下，上面显示的层态射成为
$$
T \times_{(a, b), U \times_B U} R
\longrightarrow
T \times_{f \circ a, f \circ b, U' \times_B U'} R'.
$$
若 $R$ 是该限制，这就是同构，因为在这种情形下
$R = (U \times_B U) \times_{U' \times_B U'} R'$；参见
引理 \ref{spaces-groupoids-lemma-restrict-groupoid-relation} 及其证明。反过来，若最后显示的态射
对所有 $T, a, b$ 都是同构，则可推出
$R = (U \times_B U) \times_{U' \times_B U'} R'$，即
$R$ 是 $R'$ 的限制。
\end{proof}

\begin{lemma}
\label{spaces-groupoids-lemma-quotient-stack-restrict-equivalence}
记号与假设如引理 \ref{spaces-groupoids-lemma-quotient-stack-functorial} 中所述。
商叠的态射
$$
[f] : [U/R] \longrightarrow [U'/R']
$$
是等价，当且仅当
\begin{enumerate}
\item $(U, R, s, t, c)$ 是 $(U', R', s', t', c')$ 通过
$f : U \to U'$ 得到的限制，并且
\item 映射
$$
\xymatrix{
U \times_{f, U', t'} R' \ar[r]_-{\text{pr}_1} \ar@/^3ex/[rr]^h
& R' \ar[r]_{s'} & U'
}
$$
是层的满射。
\end{enumerate}
例如，若 $\{h : U \times_{f, U', t'} R' \to U'\}$ 是 fppf 覆盖，
或 $f : U \to U'$ 是层的满射，或 $\{f : U \to U'\}$ 是 fppf 覆盖，
则第 (2) 部分成立。
\end{lemma}

\begin{proof}
% P09-ERR-0922
由引理 \ref{spaces-groupoids-lemma-quotient-stack-restrict}，我们已经知道第 (1) 部分等价于
全忠实性。因此可以假设 (1) 成立且 $[f]$ 全忠实。我们的目标是在这些假设下证明：
$[f]$ 是等价，当且仅当 (2) 成立。可以使用刻画等价的
Stacks，引理 \ref{stacks-lemma-characterize-essentially-surjective-when-ff}。

\medskip\noindent
假设 (2)。我们将使用
Stacks，引理 \ref{stacks-lemma-characterize-essentially-surjective-when-ff}
来证明 $[f]$ 是等价。设 $T$ 是概形且
$x' \in \Ob([U'/R']_T)$。存在一个覆盖 $\{g_i : T_i \to T\}$，
使得 $g_i^*x'$ 是某个元素 $a'_i \in U'(T_i)$ 的像；参见
引理 \ref{spaces-groupoids-lemma-quotient-stack-objects}。因此可以假设 $x'$ 是
$a' \in U'(T)$ 的像。由 $h$ 是层的满射这一假设，可以找到一个 fppf 覆盖
$\{\varphi_i : T_i \to T\}$ 和态射
% P09-ERR-0923
$b_i : T_i \to U \times_{g, U', t'} R'$，使得
$a' \circ \varphi_i = h \circ b_i$。记
$a_i = \text{pr}_0 \circ b_i : T_i \to U$。于是可见
$a_i \in U(T_i)$ 映到 $f \circ a_i \in U'(T_i)$，并且
$f \circ a_i \cong_{T_i} h \circ b_i = a' \circ \varphi_i$，
其中 $\cong_{T_i}$ 表示纤维范畴 $[U'/R']_{T_i}$ 中的同构。确切地说，
给出该同构的 $R'(T_i)$ 中元素是 $\text{pr}_1 \circ b_i$。
% P09-ERR-0924
这意味着 $x$ 到 $T_i$ 的限制位于函子
$[U/R]_{T_i} \to [U'/R']_{T_i}$ 的本质像中，正如所需。

\medskip\noindent
假设 $[f]$ 是等价。令 $\xi' \in [U'/R']_{U'}$ 表示对应于 $U'$
的恒等态射的对象。应用
Stacks，引理 \ref{stacks-lemma-characterize-essentially-surjective-when-ff}，
可见存在一个 fppf 覆盖
$\mathcal{U}' = \{g'_i : U'_i \to U'\}$，使得
$(g'_i)^*\xi' \cong [f](\xi_i)$，其中某个 $\xi_i$ 位于
$[U/R]_{U'_i}$ 中。加细覆盖 $\mathcal{U}'$ 之后（使用
引理 \ref{spaces-groupoids-lemma-quotient-stack-objects}），可以假设 $\xi_i$ 来自态射
$a_i : U'_i \to U$。事实
$[f](\xi_i) \cong (g'_i)^*\xi'$ 意味着，在必要时再次加细覆盖
$\mathcal{U}'$ 后，存在态射 $r'_i : U'_i \to R'$，满足
$t' \circ r'_i = f \circ a_i$ 且
$s' \circ r'_i = \text{id}_{U'} \circ g'_i$。图示如下：
$$
\xymatrix{
U \ar[d]^f & & U'_i \ar[ll]^{a_i} \ar[ld]_{r'_i} \ar[d]^{g'_i} \\
U' & R' \ar[l]_{t'} \ar[r]^{s'} & U'
}
$$
因此 $(a_i, r'_i) : U'_i \to U \times_{g, U', t'} R'$ 是态射，
满足 $h \circ (a_i, r'_i) = g'_i$；由此可知
$\{h : U \times_{g, U', t'} R' \to U'\}$ 可由 fppf 覆盖
$\mathcal{U}'$ 加细，这意味着 $h$ 诱导层的满射；参见
代数空间上的拓扑，引理
\ref{spaces-topologies-lemma-fppf-covering-surjective}。

\medskip\noindent
若 $\{h\}$ 是 fppf 覆盖，则它诱导层的满射；参见
代数空间上的拓扑，引理
\ref{spaces-topologies-lemma-fppf-covering-surjective}。
% P09-ERR-0925
若 $U' \to U$ 是满射，则 $h$ 也是满射，因为 $s$ 有一个截面
（即代数空间群胚的中性元 $e$）。
\end{proof}

\begin{lemma}
\label{spaces-groupoids-lemma-criterion-fibre-product}
记号与假设如引理 \ref{spaces-groupoids-lemma-quotient-stack-functorial} 中所述。
假设
$$
\xymatrix{
R \ar[d]_s \ar[r]_f & R' \ar[d]^{s'} \\
U \ar[r]^f & U'
}
$$
是笛卡儿的。则
$$
\xymatrix{
\mathcal{S}_U \ar[d] \ar[r] & [U/R] \ar[d]^{[f]} \\
\mathcal{S}_{U'} \ar[r] & [U'/R']
}
$$
是一个 $2$-纤维积方块。
\end{lemma}

\begin{proof}
把逆同构 $i : R \to R$ 和 $i' : R' \to R'$ 应用于引理断言中的
（第一个）笛卡儿图表，可见
$$
\xymatrix{
R \ar[d]_t \ar[r]_f & R' \ar[d]^{t'} \\
U \ar[r]^f & U'
}
$$
也是笛卡儿的。由引理 \ref{spaces-groupoids-lemma-cartesian-square-of-morphism}，有一个
$2$-纤维积方块
$$
\xymatrix{
[U''/R''] \ar[d] \ar[r] & [U/R] \ar[d] \\
\mathcal{S}_{U'} \ar[r] & [U'/R']
}
$$
其中 $U'' = U \times_{f, U', t'} R'$ 且
$R'' = R \times_{f \circ s, U', t'} R'$。由上可见
$(t, f) : R \to U''$ 是同构，并且
% P09-ERR-0926
$$
R'' =
R \times_{f \circ s, U', t'} R' =
R \times_{s, U} U \times_{f, U', t'} R' =
R \times_{s, U, t} \times R.
$$
显式地说，同构 $R \times_{s, U, t} R \to R''$ 由法则
$(r_0, r_1) \mapsto (r_0, f(r_1))$ 给出。此外，$s'', t'', c''$
转化为映射
$$
R \times_{s, U, t} R \to R,
\quad
s''(r_0, r_1) = r_1, \quad t''(r_0, r_1) = c(r_0, r_1)
$$
以及
$$
\begin{matrix}
c'' : &
(R \times_{s, U, t} R) \times_{s'', R, t''} (R \times_{s, U, t} R)
&
\longrightarrow
&
R \times_{s, U, t} R, \\
&
((r_0, r_1), (r_2, r_3)) &
\longmapsto & (c(r_0, r_2), r_3).
\end{matrix}
$$
预复合以同构
$$
R \times_{s, U, s} R \longrightarrow R \times_{s, U, t} R,
\quad
(r_0, r_1) \longmapsto (c(r_0, i(r_1)), r_1)
$$
后，可见 $t''$ 和 $s''$ 转化为 $\text{pr}_0$ 和 $\text{pr}_1$，
而 $c''$ 转化为
$\text{pr}_{02} :
R \times_{s, U, s} R \times_{s, U, s} R \to R \times_{s, U, s} R$。
因此可见有同构
$[U''/R''] \cong [R/R \times_{s, U, s} R]$，其中，作为代数空间中的群胚，
$(R, R \times_{s, U, s} R, s'', t'', c'')$ 是平凡群胚
$(U, U, \text{id}, \text{id}, \text{id})$ 通过 $s : R \to U$ 得到的限制。
由于 $s : R \to U$ 是 fppf 层的满射（因为它有右逆），态射
$$
[U''/R''] \cong [R/R \times_{s, U, s} R]
\longrightarrow
[U/U] = \mathcal{S}_U
$$
由引理 \ref{spaces-groupoids-lemma-quotient-stack-restrict-equivalence} 是等价。这就证明了引理。
\end{proof}

\section{惯性与商叠}
\label{spaces-groupoids-section-inertia}

\noindent
群胚叠的（相对）惯性叠定义于
Stacks，第 \ref{stacks-section-the-inertia-stack} 节。
% P09-ERR-0927
在纤维化范畴的设定中，实际构造及其若干性质见
范畴，第 \ref{categories-section-inertia} 节。

\begin{lemma}
\label{spaces-groupoids-lemma-presentation-inertia}
假设 $B \to S$ 以及 $(U, R, s, t, c)$ 如
定义 \ref{spaces-groupoids-definition-quotient-stack} (1) 中所述。设 $G/U$ 为群胚
$(U, R, s, t, c, e, i)$ 的稳定子群代数空间；参见
定义 \ref{spaces-groupoids-definition-stabilizer-groupoid}。置 $R' = R \times_{s, U} G$，并置
\begin{enumerate}
\item $s' : R' \to G$，$(r, g) \mapsto g$；
\item $t' : R' \to G$，$(r, g) \mapsto c(r, c(g, i(r)))$；
% P09-ERR-0928
\item $c' : R' \times_{s', G, t'} R' \to R'$，
$((r_1, g_1), (r_2, g_2) \mapsto (c(r_1, r_2), g_1)$。
\end{enumerate}
则 $(G, R', s', t', c')$ 是 $B$ 上代数空间中的群胚，并且
$$
\mathcal{I}_{[U/R]} = [G/ R'].
$$
即，相伴商叠是 $[U/R]$ 的惯性叠。
\end{lemma}

\begin{proof}
由 Stacks，引理 \ref{stacks-lemma-stackification-inertia}，只需证明
$\mathcal{I}_{[U/_{\!p}R]} = [G/_{\!p} R']$。设 $T$ 为 $S$ 上的概形。
回顾，$[U/_{\!p}R]$ 在 $T$ 上的惯性纤维化范畴的一个对象由一对
$(x, g)$ 给出，其中 $x$ 是 $[U/_{\!\!p}R]$ 在 $T$ 上的一个对象，而
$g$ 是 $x$ 在 $T$ 上其纤维范畴中的一个自同构。换言之，
$x : T \to U$ 且 $g : T \to R$，满足
$x = s \circ g = t \circ g$。这恰好意味着 $g : T \to G$。
惯性纤维化范畴中从 $(x, g) \to (y, h)$ 的一个态射（在 $T$ 上）由
$r : T \to R$ 给出，满足 $s(r) = x$、$t(r) = y$ 且
$c(r, g) = c(h, r)$；参见范畴，引理
\ref{categories-lemma-inertia-fibred-category} 中的交换图表。写成公式为
$$
h = c(r, c(g, i(r))) = c(c(r, g), i(r)).
$$
% P09-ERR-0930
记号 $s(r)$ 等是 $s \circ r$ 等的简写。
% P09-ERR-0929
$r_1 : (x_2, g_2) \to (x_1, g_1)$ 与
$r_2 : (x_1, g_1) \to (x_2, g_2)$ 的复合是
$c(r_1, r_2) : (x_1, g_1) \to (x_3, g_3)$。

\medskip\noindent
注意，在上面，可以用 $g$ 而不用 $(x, g)$ 表示
$\mathcal{I}_{[U/_{\!p}R]}$ 在 $T$ 上的一个对象，因为 $x$ 是 $g$
在结构态射 $G \to U$ 下的像。于是，态射 $g \to h$ 若视为
$\mathcal{I}_{[U/_{\!p}R]}$ 在 $T$ 上的态射，则恰好对应于态射
$r' : T \to R'$，满足
$s'(r') = g$ 且 $t'(r') = h$。此外，复合对应于 (3) 中所述的法则。
引理得证。
\end{proof}

\begin{lemma}
\label{spaces-groupoids-lemma-2-cartesian-inertia}
假设 $B \to S$ 以及 $(U, R, s, t, c)$ 如
定义 \ref{spaces-groupoids-definition-quotient-stack} (1) 中所述。设 $G/U$ 为群胚
$(U, R, s, t, c, e, i)$ 的稳定子群代数空间；参见
定义 \ref{spaces-groupoids-definition-stabilizer-groupoid}。有一个典范 $2$-笛卡儿图表
$$
\xymatrix{
\mathcal{S}_G \ar[r] \ar[d] & \mathcal{S}_U \ar[d] \\
\mathcal{I}_{[U/R]} \ar[r] & [U/R]
}
$$
它由 $(\Sch/S)_{fppf}$ 上的群胚叠组成。
\end{lemma}

\begin{proof}
由引理 \ref{spaces-groupoids-lemma-criterion-fibre-product}，只需证明
% P09-ERR-0931
引理 \ref{spaces-groupoids-lemma-presentation-inertia} 的态射 $s' : R' \to G$
同构于 $s$ 沿结构态射 $G \to U$ 的基变换。这个基变换性质由
$s'$ 的构造显然可见。
\end{proof}

\section{Gerbe 与商叠}
\label{spaces-groupoids-section-gerbes}

\noindent
本节把商叠与 Stacks，第 \ref{stacks-section-gerbes} 节中的讨论联系起来，
特别是与 Stacks，定义
\ref{stacks-definition-gerbe-over-stack-in-groupoids} 中定义的 gerbe 联系起来。
一般而言，本节出现的群胚叠并不是代数叠！

\begin{lemma}
\label{spaces-groupoids-lemma-when-gerbe}
记号与假设同引理 \ref{spaces-groupoids-lemma-quotient-stack-functorial}。商叠的态射
$$
[f] : [U/R] \longrightarrow [U'/R']
$$
$[U/R]$ 成为 $[U'/R']$ 上 gerbe 的充分条件是 $f : U \to U'$ 与
$R \to R'|_U$ 均为 fppf 层的满射。这里 $R'|_U$ 表示 $R'$ 到 $U$
的限制，它是经 $f : U \to U'$ 得到的。
\end{lemma}

\begin{proof}
我们将验证 Stacks，引理 \ref{stacks-lemma-when-gerbe} 的性质
(2) (a) 与 (2) (b) 成立。性质 (2)(a) 成立，因为 $U \to U'$ 是层的满射
（利用引理 \ref{spaces-groupoids-lemma-quotient-stack-objects} 可知，$[U'/R']$ 中的对象局部来自
$U'$）。为证明 (2)(b)，令 $x, y$ 为 $[U/R]$ 的对象，它们位于概形 $T/S$ 上。
% P09-ERR-0932
令 $x', y'$ 为 $x, y$ 在范畴 $[U'/'R]_T$ 中的像。条件 (2)(b) 要求我们检验层的映射
$$
\mathit{Isom}(x, y) \longrightarrow \mathit{Isom}(x', y')
$$
在 $(\Sch/T)_{fppf}$ 上是满射。为此，可以在 $T$ 上 fppf 局部地工作，
% P09-ERR-0933
并假设它们来自 $a, b \in U(T)$。在这种情形下，$x', y'$ 分别对应于
$f \circ a, f \circ b$。由引理 \ref{spaces-groupoids-lemma-quotient-stack-morphisms}，
在这种情形下，上述层的映射变为
$$
T \times_{(a, b), U \times_B U} R
\longrightarrow
T \times_{f \circ a, f \circ b, U' \times_B U'} R' =
T \times_{(a, b), U \times_B U} R'|_U.
$$
因此，$R \to R'|_U$ 是 $(\Sch/S)_{fppf}$ 上 fppf 层的满射这一假设
蕴含所需的满射性。
\end{proof}

\begin{lemma}
\label{spaces-groupoids-lemma-group-quotient-gerbe}
设 $S$ 为概形，$B$ 为 $S$ 上的代数空间，$G$ 为 $B$ 上的群代数空间。
赋予 $B$ 以 $G$ 的平凡作用。态射
$$
[B/G] \longrightarrow \mathcal{S}_B
$$
（引理 \ref{spaces-groupoids-lemma-quotient-stack-arrows}）使 $[B/G]$ 成为 $B$ 上的 gerbe。
\end{lemma}

\begin{proof}
由引理 \ref{spaces-groupoids-lemma-when-gerbe} 立即可得，因为态射 $B \to B$ 与
$B \times_B G \to B$ 作为层的态射都是满射。
\end{proof}

\section{商叠与大位点的变换}
\label{spaces-groupoids-section-bigger-site}

\noindent
我们建议初次阅读时跳过本节。叠的拉回定义于
Stacks，第 \ref{stacks-section-inverse-image} 节。

\begin{lemma}
\label{spaces-groupoids-lemma-quotient-stack-change-big-site}
设给定大位点 $\Sch_{fppf}$ 与 $\Sch'_{fppf}$。假设
$\Sch_{fppf}$ 包含于 $\Sch'_{fppf}$ 中；参见拓扑，第
\ref{topologies-section-change-alpha} 节。设
$S \in \Ob(\Sch_{fppf})$。设
$B, U, R \in \Sh((\Sch/S)_{fppf})$ 为代数空间，并设
$(U, R, s, t, c)$ 为 $B$ 上代数空间中的群胚。
% P09-ERR-0934
设 $f : (\Sch'/S)_{fppf} \to (\Sch/S)_{fppf}$ 为与包含函子
$u : \Sch_{fppf} \to \Sch'_{fppf}$ 对应的位点态射。则有典范等价
$$
[f^{-1}U/f^{-1}R]
\longrightarrow
f^{-1}[U/R]
$$
它由 $(\Sch'/S)_{fppf}$ 上的群胚叠组成。
\end{lemma}

\begin{proof}
注意，$f^{-1}B, f^{-1}U, f^{-1}R \in \Sh((\Sch'/S)_{fppf})$
由空间，引理 \ref{spaces-lemma-change-big-site} 是代数空间，因而
$(f^{-1}U, f^{-1}R, f^{-1}s, f^{-1}t, f^{-1}c)$ 是
$f^{-1}B$ 上代数空间中的群胚。因此该陈述有意义。

\medskip\noindent
% P09-ERR-0935
范畴 $u_p[U/_{\!p}R]$ 是范畴 $u_{pp}[U/_{\!p}R]$ 关于态射的
右乘法系 $I$ 的局部化。$u_{pp}[U/_{\!p}R]$ 的一个对象是三元组
$$
(T', \phi : T' \to T, x)
$$
其中
$T' \in \Ob((\Sch'/S)_{fppf})$，
$T \in \Ob((\Sch/S)_{fppf})$，$\phi$ 是 $S$ 上概形的态射，
而 $x : T \to U$ 是 $(\Sch/S)_{fppf}$ 上的层态射。注意，概形态射
$\phi : T' \to T$ 等同于态射 $\phi : T' \to u(T)$；又由于
$u(T)$ 表示 $f^{-1}T$，这也等同于态射 $T' \to f^{-1}T$。
此外，由于 $f^{-1}$ 在代数空间上全忠实（参见空间，引理
\ref{spaces-lemma-fully-faithful}），我们也可以把 $x$ 视为态射
$x : f^{-1}T \to f^{-1}U$。下文将不再说明地作这些等同。态射
$$
(a, a', \alpha) :
(T'_1, \phi_1 : T'_1 \to T_1, x_1)
\longrightarrow
(T'_2, \phi_2 : T'_2 \to T_2, x_2)
$$
若是 $u_{pp}[U/_{\!p}R]$ 的态射，就是一个交换图表
$$
\xymatrix{
& & U \\
T'_1 \ar[d]_{a'} \ar[r]_{\phi_1} &
T_1 \ar[d]_a \ar[ru]^{x_1} \ar[r]_\alpha &
R \ar[d]^t \ar[u]_s \\
T'_2 \ar[r]^{\phi_2} &
T_2 \ar[r]^{x_2} &
U
}
$$
而且这样的态射属于 $I$ 当且仅当 $T'_1 = T'_2$ 且
$a' = \text{id}$。我们按如下法则定义函子
$$
u_{pp}[U/_{\!p}R] \longrightarrow [f^{-1}U/_{\!p}f^{-1}R]
$$
它在对象上的法则为
$$
(T', \phi : T' \to T, x) \longmapsto (x \circ \phi : T' \to f^{-1}U)
$$
在上述态射上的法则为
$$
(a, a', \alpha) \longmapsto (\alpha \circ \phi_1 : T'_1 \to f^{-1}R)
$$
显然，$I$ 的元素被变换为同构，因为
$(f^{-1}U, f^{-1}R, f^{-1}s, f^{-1}t, f^{-1}c)$ 是
$f^{-1}B$ 上代数空间中的群胚。因此，这个函子典范地通过函子
$$
u_p[U/_{\!p}R] \longrightarrow [f^{-1}U/_{\!p}f^{-1}R]
$$
分解。应用叠化，可得叠函子
$$
f^{-1}[U/R] \longrightarrow [f^{-1}U/f^{-1}R]
$$
它在 $(\Sch'/S)_{fppf}$ 上；这是因为由 Stacks，引理
\ref{stacks-lemma-technical-up}，
叠 $f^{-1}[U/R]$ 是 $u_p[U/_{\!p}R]$ 的叠化。

\medskip\noindent
至此我们已有叠的态射。要验证它是等价，只需证明它全忠实且对象局部地
位于本质像中；参见 Stacks，引理 \ref{stacks-lemma-characterize-ff} 与
\ref{stacks-lemma-characterize-essentially-surjective-when-ff}。
% P09-ERR-0936
关于对象的陈述成立，因为 $f^{-1}R$ 有满平展态射
$f^{-1}W \to f^{-1}R$，其中 $W$ 是 $(\Sch/S)_{fppf}$ 的某个对象。
要证明函子“满”，只需证明态射局部地位于该函子的像中；这成立是因为
$f^{-1}U$ 有满平展态射 $f^{-1}W \to f^{-1}U$，其中 $W$ 是
$(\Sch/S)_{fppf}$ 的某个对象。略去函子忠实性的证明。
\end{proof}

\section{分离性条件}
\label{spaces-groupoids-section-separation}

\noindent
这实际上意指态射 $j : R \to U \times_B U$ 所满足的条件；这里给定一个
代数空间中的群胚 $(U, R, s, t, c)$，它在 $B$ 上。与上一节一样，
我们先写出相应的图表。

\begin{lemma}
\label{spaces-groupoids-lemma-diagram-diagonal}
设 $B \to S$ 如第 \ref{spaces-groupoids-section-notation} 节所示。设
$(U, R, s, t, c)$ 为 $B$ 上代数空间中的群胚。设
$G \to U$ 为稳定子群代数空间。交换图表
$$
\xymatrix{
R \ar[d]^{\Delta_{R/U \times_B U}} \ar[rrr]_{f \mapsto (f, s(f))} & & &
R \times_{s, U} U \ar[d] \ar[r] & U \ar[d] \\
R \times_{(U \times_B U)} R \ar[rrr]^{(f, g) \mapsto (f, f^{-1} \circ g)} & & &
R \times_{s, U} G \ar[r] & G
}
$$
% P09-ERR-0937
中的左边两条水平箭头是同构，而右方块是纤维积方块。
\end{lemma}

\begin{proof}
略。这是关于定义以及代数几何中函子观点的练习。
\end{proof}

\begin{lemma}
\label{spaces-groupoids-lemma-diagonal}
设 $B \to S$ 如第 \ref{spaces-groupoids-section-notation} 节所示。设
$(U, R, s, t, c)$ 为 $B$ 上代数空间中的群胚。设
$G \to U$ 为稳定子群代数空间。
\begin{enumerate}
\item 下列条件等价：
\begin{enumerate}
\item $j : R \to U \times_B U$ 是分离的；
\item $G \to U$ 是分离的；以及
\item $e : U \to G$ 是闭浸入。
\end{enumerate}
\item 下列条件等价：
\begin{enumerate}
\item $j : R \to U \times_B U$ 是局部分离的；
\item $G \to U$ 是局部分离的；以及
\item $e : U \to G$ 是浸入。
\end{enumerate}
\item 下列条件等价：
\begin{enumerate}
\item $j : R \to U \times_B U$ 是拟分离的；
\item $G \to U$ 是拟分离的；以及
\item $e : U \to G$ 是拟紧的。
\end{enumerate}
\end{enumerate}
\end{lemma}

\begin{proof}
群代数空间 $G \to U$ 是 $R \to U \times_B U$ 沿对角态射
$U \to U \times_B U$ 的基变换；参见引理
\ref{spaces-groupoids-lemma-groupoid-stabilizer}。因此，如果 $j$ 是分离的（相应地，
局部分离的、拟分离的），那么 $G \to U$ 是分离的（相应地，
局部分离的、拟分离的）。参见空间的态射，引理
\ref{spaces-morphisms-lemma-base-change-separated}。因此在 (1)、(2) 与
(3) 中，(a) $\Rightarrow$ (b)。

\medskip\noindent
反之，如果 $G \to U$ 是分离的（相应地，局部分离的、拟分离的），
那么态射 $e : U \to G$ 作为结构态射 $G \to U$ 的截面，是闭浸入
（相应地，浸入、拟紧的）；参见空间的态射，引理
\ref{spaces-morphisms-lemma-section-immersion}。因此在 (1)、(2) 与
(3) 中，(b) $\Rightarrow$ (c)。

\medskip\noindent
如果 $e$ 是闭浸入（相应地，浸入、拟紧的），那么由引理
\ref{spaces-groupoids-lemma-diagram-diagonal} 的结论（以及空间，引理
\ref{spaces-lemma-base-change-immersions} 和空间的态射，引理
\ref{spaces-morphisms-lemma-base-change-quasi-compact}），可知
$\Delta_{R/U \times_B U}$ 是闭浸入（相应地，浸入、拟紧的）。
因此在 (1)、(2) 与 (3) 中，(c) $\Rightarrow$ (a)。
\end{proof}
% Generated mechanically from the frozen spaces-more-groupoids.tex target.
% Do not edit this generated file; edit the bound locale source instead.

% OK, start here.
%

\setcounter{chapter}{78}
\chapter{代数空间中的群胚进阶}



\phantomsection
\label{spaces-more-groupoids-section-phantom}


\section{引言}
\label{spaces-more-groupoids-section-introduction}

\noindent
本章专论代数空间中群胚的若干进阶主题。虽然各项结果均以代数空间中的
群胚来表述，读者仍应始终记住下面的 $2$-笛卡尔图：
\begin{equation}
\label{spaces-more-groupoids-equation-quotient-stack}
\vcenter{
\xymatrix{
R \ar[r] \ar[d] & U \ar[d] \\
U \ar[r] & [U/R]
}
}
\end{equation}
其中 $[U/R]$ 是商叠，见《代数空间中的群胚》注
\ref{spaces-groupoids-remark-fundamental-square}。本章许多结果的动机都来自
对这张图的思考。例如可参见 Keel 和 Mori 的优美论文 \cite{K-M}。





\section{记号}
\label{spaces-more-groupoids-section-notation}

\noindent
我们继续遵循《代数空间中的群胚》第
\ref{spaces-groupoids-section-notation} 节中引入的约定和记号。





\section{常用图表}
\label{spaces-more-groupoids-section-diagrams}

\noindent
为便于本章引用，我们简要重述《代数空间中的群胚》引理
\ref{spaces-groupoids-lemma-diagram} 和
\ref{spaces-groupoids-lemma-diagram-pull} 的结果。
设 $S$ 为概形，$B$ 为 $S$ 上的代数空间。
设 $(U, R, s, t, c)$ 为 $B$ 上代数空间中的群胚。在交换图
\begin{equation}
\label{spaces-more-groupoids-equation-diagram}
\vcenter{
\xymatrix{
& U & \\
R \ar[d]_s \ar[ru]^t &
R \times_{s, U, t} R
\ar[l]^-{\text{pr}_0} \ar[d]^{\text{pr}_1} \ar[r]_-c &
R \ar[d]^s \ar[lu]_t \\
U & R \ar[l]_t \ar[r]^s & U
}
}
\end{equation}
中，下方两个方块都是纤维积方块。此外，上方三角形（其实是一个方块）
也是笛卡尔的。

\medskip\noindent
图
\begin{equation}
\label{spaces-more-groupoids-equation-pull}
\vcenter{
\xymatrix{
R \times_{t, U, t} R
\ar@<1ex>[r]^-{\text{pr}_1} \ar@<-1ex>[r]_-{\text{pr}_0}
\ar[d]_{\text{pr}_0 \times c \circ (i, 1)} &
R \ar[r]^t \ar[d]^{\text{id}_R} &
U \ar[d]^{\text{id}_U} \\
R \times_{s, U, t} R
\ar@<1ex>[r]^-c \ar@<-1ex>[r]_-{\text{pr}_0} \ar[d]_{\text{pr}_1} &
R \ar[r]^t \ar[d]^s &
U \\
R \ar@<1ex>[r]^s \ar@<-1ex>[r]_t &
U
}
}
\end{equation}
是交换的。上面两行通过所给竖直映射同构。左下方两个方块都是笛卡尔的。






\section{局部结构}
\label{spaces-more-groupoids-section-local}

\noindent
设 $S$ 为概形，设 $(U, R, s, t, c, e, i)$ 为 $S$ 上代数空间中的群胚。
设 $\overline{u}$ 为 $U$ 的一个几何点。本节说明我们在局部环
（《代数空间的性质》定义
\ref{spaces-properties-definition-etale-local-rings}）
$$
A = \mathcal{O}_{U, \overline{u}}
\quad\text{且}\quad
B = \mathcal{O}_{R, e(\overline{u})}
$$
上得到怎样的结构。我们约定，由态射 $s,t,c,e,i$ 诱导的局部环同态仍以
相应字母表示。特别地，有局部环的交换图
$$
\xymatrix{
A \ar[rd]_t \ar[rrd]^1 \\
& B \ar[r]^e & A \\
A \ar[ru]^s \ar[rru]_1
}
$$
因此，若 $I \subset B$ 表示 $e : B \to A$ 的核，则
$B = s(A) \oplus I = t(A) \oplus I$。记
$$
C = \mathcal{O}_{R \times_{s, U, t} R, (e, e)(\overline{u})}
$$
于是
$$
C =
(B \otimes_{s, A, t} B)_{\mathfrak m_B \otimes B + B \otimes \mathfrak m_B}^h
$$
因为局部化
$(B \otimes_{s, A, t} B)_{\mathfrak m_B \otimes B + B \otimes \mathfrak m_B}$
的剩余域是可分闭的。设 $J \subset C$ 为 $C$ 的一个理想，它由
$I \otimes B + B \otimes I$ 生成。于是 $J$ 也正是局部环同态
$$
(e, e) : C \longrightarrow A
$$
的核。合成律 $c : R \times_{s, U, t} R \to R$ 对应于环映射
$$
c : B \longrightarrow C
$$
它把 $I$ 映入 $J$。

\begin{lemma}
\label{spaces-more-groupoids-lemma-first-order-structure-c}
映射 $I/I^2 \to J/J^2$ 由 $c$ 诱导，它是复合
$$
I/I^2 \xrightarrow{(1, 1)} I/I^2 \oplus I/I^2 \to J/J^2
$$
其中第二个箭头来自等式 $J = (I \otimes B + B \otimes I)C$。
映射 $i : B \to B$ 诱导映射 $-1 : I/I^2 \to I/I^2$。
\end{lemma}

\begin{proof}
要描述从 $C$ 到另一个亨泽尔局部环的局部同态，只须说明形如
$b_1 \otimes b_2$ 的元素被映到哪里；例如可用《代数》引理
\ref{algebra-lemma-henselian-functorial}。记住这一点，我们有两个典范映射
$$
e_2 : C \to B,\ b_1 \otimes b_2 \mapsto b_1s(e(b_2)),\quad
e_1 : C \to B,\ b_1 \otimes b_2 \mapsto t(e(b_1))b_2
$$
它们对应于嵌入 $R \to R \times_{s, U, t} R$，分别由
$r \mapsto (r, e(s(r)))$ 和 $r \mapsto (e(t(r)), r)$ 给出。
这些映射定义了映射 $J/J^2 \to I/I^2$；二者合在一起给出引理中映射
$I/I^2 \oplus I/I^2 \to J/J^2$ 的逆。因此，要证明该陈述，只须说明
% P09-ERR-0938
$e_1 \circ c : B \to B$ 和 $e_2 \circ c : B \to B$ 都是恒等映射。
这来自如下事实：两个复合 $R \to R \times_{s, U, t} R \to R$
都是恒等映射。

\medskip\noindent
关于 $i$ 的陈述来自关于 $c$ 的陈述以及等式
$c \circ (1, i) = e \circ t$。略去若干细节。
\end{proof}










\section{截面群胚}
\label{spaces-more-groupoids-section-groupoid-sections}

\noindent
假设给定一个群胚
$(\text{Ob}, \text{Arrows}, s, t, c, e, i)$。我们可以构造一个幺半群
$\Gamma$；它的元素是映射 $\delta : \text{Ob} \to \text{Arrows}$，满足
$s \circ \delta = \text{id}_{\text{Ob}}$，其合成定义为
$$
\delta_1 \circ \delta_2 = c(\delta_1 \circ t \circ \delta_2, \delta_2)
$$
换言之，$\Gamma$ 的元素是一条规则 $\delta$，它从每个对象出发指定一个箭头，
而合成就是自然的做法。例如，用显然的记号有
$$
\vcenter{
\xymatrix{
& \bullet \ar[dl] \\
\bullet \ar@(ul, dl)[] & \bullet \ar[d] \\
& \bullet \ar[lu]
}
}
\quad\circ\quad
\vcenter{
\xymatrix{
& \bullet \ar[d] \\
\bullet \ar[ru] & \bullet \ar[d] \\
& \bullet \ar[lu]
}
}
\quad = \quad\quad
\vcenter{
\xymatrix{
& \bullet \ar@/^/[dd] \\
\bullet \ar@(ul, dl)[] & \bullet \ar[l] \\
& \bullet \ar[lu]
}
}
$$

\medskip\noindent
同样的过程可用于代数空间中的群胚 $(U, R, s, t, c, e, i)$；
这里群胚取在概形 $S$ 上。具体地，令 $\Gamma$ 的元素组成集合
$$
\Gamma = \{\delta : U \to R \mid s \circ \delta = \text{id}_U\}
$$
而合成 $\circ : \Gamma \times \Gamma \to \Gamma$ 由上面的规则给出：
\begin{equation}
\label{spaces-more-groupoids-equation-composition}
\delta_1 \circ \delta_2 = c(\delta_1 \circ t \circ \delta_2, \delta_2).
\end{equation}
单位元由 $e \in \Gamma$ 给出。
% P09-ERR-0939
群胚 $\Gamma$ 一般并不是群，因为可能存在没有逆元的元素
$\delta \in \Gamma$。显然，$\delta \in \Gamma$ 有逆元当且仅当
$t \circ \delta$ 是 $U$ 的自同构；在这种情形下，
$\delta^{-1} = i \circ \delta \circ (t \circ \delta)^{-1}$。

\medskip\noindent
为后文使用，我们讨论子群胚 $\Gamma_0$；它是 $\Gamma$ 中由与单位元 $e$
无穷小接近的截面构成。更确切地，假设给定一个 $R$-不变闭子空间
$U_0 \subset U$，并且 $U$ 是 $U_0$ 的一阶加厚。记
$R_0 = s^{-1}(U_0) = t^{-1}(U_0)$，并令
$(U_0, R_0, s_0, t_0, c_0, e_0, i_0)$ 为相应的代数空间中的群胚。置
$$
\Gamma_0 = \{\delta \in \Gamma \mid \delta|_{U_0} = e_0\}
$$
若 $s$ 和 $t$ 平坦，则 $\Gamma_0$ 的每个元素都可逆。事实上，
$t \circ \delta$ 是态射 $U \to U$，并且在
$\mathcal{O}_{U_0}$ 以及 $\mathcal{C}_{U_0/U}$ 上都诱导恒等映射
（引理 \ref{spaces-more-groupoids-lemma-idenity-on-conormal}）；再利用短正合列
$0 \to \mathcal{C}_{U_0/U} \to \mathcal{O}_U \to \mathcal{O}_{U_0} \to 0$
即可得到结论。

\begin{lemma}
\label{spaces-more-groupoids-lemma-idenity-on-conormal}
在本节所讨论的情形下，设 $\delta \in \Gamma_0$，并令
$f = t \circ \delta : U \to U$。若 $s,t$ 平坦，则典范映射
$\mathcal{C}_{U_0/U} \to \mathcal{C}_{U_0/U}$（它由 $f$ 诱导；见
《代数空间态射进阶》引理
\ref{spaces-more-morphisms-lemma-conormal-functorial}）是恒等映射。
\end{lemma}

\begin{proof}
为说明这一点，把图 (\ref{spaces-more-groupoids-equation-pull}) 的下部延拓如下：
$$
\xymatrix{
Y \ar[r] \ar[d] &
R \times_{s, U, t} R
\ar@<1ex>[r]^-c \ar@<-1ex>[r]_-{\text{pr}_0} \ar[d]_{\text{pr}_1} &
R \ar[r]^t \ar[d]^s &
U \\
U \ar[r]_\delta &
R \ar@<1ex>[r]^s \ar@<-1ex>[r]_t &
U
}
$$
其中左方块是笛卡尔的，这也就是我们对 $Y$ 的定义；我们不需要知道
$Y$ 的更多性质。把各处基变换到 $U_0$，得到具有同样性质的类似图。
我们要证明 $\text{id}_U = s \circ \delta$ 和
$f = t \circ \delta$ 在余法层上诱导相同映射。由于 $s$ 平坦且满，
只须对沿顶行的两个复合 $a,b : Y \to R$ 证明同样的事情。
注意 $a_0=b_0$，并且 $a$、$b$ 中有一个是同构，因为已知
$s \circ \delta$ 是同构。因此，两个态射 $a,b : Y \to R$ 都是
$U$ 上平坦代数空间之间的态射（分别通过态射 $t : R \to U$
和态射 $t \circ a = t \circ b : Y \to U$）。
这蕴含所需结论。事实上，由《代数空间态射进阶》引理
\ref{spaces-more-morphisms-lemma-conormal-functorial-more}
中与合成的相容性可知，两个映射
$a_0^*\mathcal{C}_{R_0/R} \to \mathcal{C}_{Y_0/Y}$
都嵌入交换图
$$
\xymatrix{
a_0^*\mathcal{C}_{R_0/R} \ar[rr] & & \mathcal{C}_{Y_0/Y} \\
a_0^*t_0^*\mathcal{C}_{U_0/U} \ar[u] \ar@{=}[rr] & &
(t_0 \circ a_0)^*\mathcal{C}_{U_0/U} \ar[u]
}
$$
中；由《代数空间态射进阶》引理
\ref{spaces-more-morphisms-lemma-deform}，其中竖直箭头都是同构。
因此引理成立。
\end{proof}

\noindent
下面辨认群 $\Gamma_0$。把《代数空间态射进阶》注
\ref{spaces-more-morphisms-remark-action-by-derivations} 和
\ref{spaces-more-morphisms-remark-another-special-case} 中的讨论应用于图
$$
\xymatrix{
(U_0 \subset U) \ar@{..>}[rr]_{(e_0, \delta)}
\ar[rd]_{(\text{id}_{U_0}, \text{id}_U)} & &
(R_0 \subset R) \ar[ld]^{(s_0, s)} \\
& (U_0 \subset U)
}
$$
可见 $\delta = \theta \cdot e$；这里相对于 $\mathcal{O}_{U_0}$，
$\theta : e_0^*\Omega_{R_0/U_0} \to \mathcal{C}_{U_0/U}$ 是唯一的线性映射。
于是，应用《代数空间态射进阶》引理
\ref{spaces-more-morphisms-lemma-action-sheaf}，我们得到双射
\begin{equation}
\label{spaces-more-groupoids-equation-isomorphism}
\Hom_{\mathcal{O}_{U_0}}(e_0^*\Omega_{R_0/U_0}, \mathcal{C}_{U_0/U})
\longrightarrow
\Gamma_0.
\end{equation}

\begin{lemma}
\label{spaces-more-groupoids-lemma-composition-is-addition}
双射 (\ref{spaces-more-groupoids-equation-isomorphism}) 是群同构。
\end{lemma}

\begin{proof}
设 $\delta_1, \delta_2 \in \Gamma_0$ 分别对应于如上的
$\theta_1,\theta_2$，并设合成 $\delta = \delta_1 \circ \delta_2$
在 $\Gamma_0$ 中对应于 $\theta$。我们要证明
$\theta = \theta_1 + \theta_2$。回忆《代数空间态射进阶》引理
\ref{spaces-more-morphisms-lemma-action-by-derivations}：
$\theta_1,\theta_2,\theta$ 分别对应于导子
$D_1,D_2,D : e_0^{-1}\mathcal{O}_{R_0} \to \mathcal{C}_{U_0/U}$，
其中 $D_1 = \theta_1 \circ \text{d}_{R_0/U_0}$，其余依此类推。
只须检验 $D=D_1+D_2$。

\medskip\noindent
这一等式可在茎上检验。设 $\overline{u}$ 为 $U$ 的一个几何点，并使用
第 \ref{spaces-more-groupoids-section-local} 节中引入的局部环 $A,B,C$。
态射 $\delta_i$ 对应于环映射 $\delta_i:B\to A$。设 $K\subset A$ 为
平方为零的理想，且 $A/K=\mathcal{O}_{U_0,\overline{u}}$。换言之，
$K$ 是 $\mathcal{C}_{U_0/U}$ 在 $\overline{u}$ 处的茎。
$\delta_i\in\Gamma_0$ 这一事实恰好意味着 $\delta_i(I)\subset K$。
导子 $D_i$ 就是映射 $\delta_i-e:B\to A$。由于
$B=s(A)\oplus I$，可见 $D_i$ 由它在 $I$ 上的限制决定，而这个限制正是
$\delta_i|_I$。此外，$D_i$ 乃至 $\delta_i$ 都消去 $I^2$，因为
% P09-ERR-0940
$I = \Ker(I)$。

\medskip\noindent
为完成证明，注意 $\delta$ 对应于复合
$$
B \to C =
(B \otimes_{s, A, t} B)^h_{\mathfrak m_B \otimes B + B \otimes \mathfrak m_B}
\to A
$$
其中第一个箭头是 $c$，而由 (\ref{spaces-more-groupoids-equation-composition})，第二个箭头由规则
$b_1 \otimes b_2 \mapsto \delta_2(t(\delta_1(b_1))) \delta_2(b_2)$
决定。由引理 \ref{spaces-more-groupoids-lemma-first-order-structure-c}，元素 $\zeta$ 属于 $I$，并被映到
$\zeta\otimes 1+1\otimes\zeta$ 加上高阶项。因此
$$
D(\zeta) = (\delta_2 \circ t)\left(D_1(\zeta)\right) + D_2(\zeta)
$$
然而，由引理 \ref{spaces-more-groupoids-lemma-idenity-on-conormal}，
$\delta_2\circ t$ 在 $K=\mathcal{C}_{U_0/U,\overline{u}}$ 上的作用是恒等，
结论即得。
\end{proof}

\section{群胚的性质}
\label{spaces-more-groupoids-section-technical-lemma}

\noindent
本节是《群胚进阶》第 \ref{more-groupoids-section-technical-lemma} 节的类似版本。
强烈建议读者先阅读该节。

\medskip\noindent
下面的引理是《群胚进阶》引理
\ref{more-groupoids-lemma-property-invariant} 的类似版本。

\begin{lemma}
\label{spaces-more-groupoids-lemma-property-invariant}
设 $B \to S$ 如第 \ref{spaces-more-groupoids-section-notation} 节中所述。
设 $(U, R, s, t, c)$ 为 $B$ 上代数空间中的群胚。设
$\tau \in \{fppf, \linebreak[0] \etale, \linebreak[0]
smooth, \linebreak[0] syntomic\}$。
设 $\mathcal{P}$ 为代数空间态射的一个性质，并且它在目标上对
$\tau$ 是局部的（《代数空间的下降》定义
\ref{spaces-descent-definition-property-morphisms-local}）。
假设 $\{s : R \to U\}$ 和 $\{t : R \to U\}$ 是
$\tau$-拓扑的覆盖。设 $W\subset U$ 为最大的开子空间，使得
$s^{-1}(W)\to W$ 具有性质 $\mathcal{P}$。那么 $W$ 是 $R$-不变的
（《代数空间中的群胚》定义
\ref{spaces-groupoids-definition-invariant-open}）。
\end{lemma}

\begin{proof}
开子空间 $W\subset U$ 的存在性及其性质见《代数空间的下降》引理
\ref{spaces-descent-lemma-largest-open-of-the-base}。
在图 (\ref{spaces-more-groupoids-equation-diagram}) 中，令 $W_1\subset R$ 为最大的开子概形，
使得其上的态射
$\text{pr}_1:R\times_{s,U,t}R\to R$ 具有性质 $\mathcal{P}$。
% P09-ERR-0943
由上述《代数空间的下降》引理
\ref{spaces-descent-lemma-largest-open-of-the-base} 以及
$\{s:R\to U\}$ 和 $\{t:R\to U\}$ 是 $\tau$-拓扑覆盖这一假设，得到
$t^{-1}(W)=W_1=s^{-1}(W)$，正是所需结论。
\end{proof}

\begin{lemma}
\label{spaces-more-groupoids-lemma-property-G-invariant}
设 $B\to S$ 如第 \ref{spaces-more-groupoids-section-notation} 节中所述。
设 $(U,R,s,t,c)$ 为 $B$ 上代数空间中的群胚。
设 $G\to U$ 为它的稳定子群代数空间。设
$\tau \in \{fppf, \linebreak[0] \etale, \linebreak[0]
smooth, \linebreak[0] syntomic\}$。
设 $\mathcal{P}$ 为代数空间态射的一个性质，并且它在目标上对
$\tau$ 是局部的。假设 $\{s:R\to U\}$ 和 $\{t:R\to U\}$
是 $\tau$-拓扑的覆盖。设 $W\subset U$ 为最大的开子空间，使得
$G_W\to W$ 具有性质 $\mathcal{P}$。那么 $W$ 是 $R$-不变的（见
《代数空间中的群胚》定义
\ref{spaces-groupoids-definition-invariant-open}）。
\end{lemma}

\begin{proof}
开子空间 $W\subset U$ 的存在性及其性质见《代数空间的下降》引理
\ref{spaces-descent-lemma-largest-open-of-the-base}。态射
$$
G \times_{U, t} R \longrightarrow R \times_{s, U} G, \quad
(g, r) \longmapsto (r, r^{-1} \circ g \circ r)
$$
是 $R$ 上代数空间的同构（其中 $\circ$ 表示群胚中的合成）。因此，
有 $s^{-1}(W)=t^{-1}(W)$；这是由上述《代数空间的下降》引理
\ref{spaces-descent-lemma-largest-open-of-the-base} 中证明的 $W$ 的性质得到的。
\end{proof}




\section{比较纤维}
\label{spaces-more-groupoids-section-fibres}

\noindent
本节是《群胚进阶》第 \ref{more-groupoids-section-fibres} 节的类似版本。
强烈建议读者先阅读该节。

\begin{lemma}
\label{spaces-more-groupoids-lemma-two-fibres}
设 $B\to S$ 如第 \ref{spaces-more-groupoids-section-notation} 节中所述。
设 $(U,R,s,t,c)$ 为 $B$ 上代数空间中的群胚。设 $K$ 为域，并设
$r,r':\Spec(K)\to R$ 为满足
$t\circ r=t\circ r':\Spec(K)\to U$ 的态射。
置 $u=s\circ r$、$u'=s\circ r'$，并以
$F_u=\Spec(K)\times_{u,U,s}R$ 和
$F_{u'}=\Spec(K)\times_{u',U,s}R$ 表示相应纤维积。
那么 $F_u\cong F_{u'}$；这是 $K$ 上代数空间的同构。
\end{lemma}

\begin{proof}
我们使用图 (\ref{spaces-more-groupoids-equation-diagram}) 的存在性和性质。
存在态射 $\xi:\Spec(K)\to R\times_{s,U,t}R$，使得
$\text{pr}_0\circ\xi=r$ 且 $c\circ\xi=r'$。
令 $\tilde r=\text{pr}_1\circ\xi:\Spec(K)\to R$。
考察图 (\ref{spaces-more-groupoids-equation-diagram}) 下方的两个方块，可见 $F_u$ 和 $F_{u'}$
都与代数空间
$\Spec(K)\times_{\tilde r,R,\text{pr}_1}(R\times_{s,U,t}R)$ 等同。
\end{proof}

\noindent
实际上，在引理的情形下，成对态射 $s:(R,r)\to(U,u)$ 和
$s:(R,r')\to(U,u')$ 在 $\tau$-拓扑中局部同构，只要
$\{s:R\to U\}$ 是一个 $\tau$-覆盖。
% P09-ERR-0941
如果有需要，我们会在这里加入一个精确陈述。







\section{限制群胚}
\label{spaces-more-groupoids-section-restricting-groupoids}

\noindent
本节汇集一批关于群胚性质的引理；这些性质会遗传给限制。
其中大多数引理都可通过考察限制的定义图
\begin{equation}
\label{spaces-more-groupoids-equation-restriction}
\vcenter{
\xymatrix{
R' \ar[d] \ar[r] \ar@/_3pc/[dd]_{t'} \ar@/^1pc/[rr]^{s'}&
R \times_{s, U} U' \ar[r] \ar[d] &
U' \ar[d]^g \\
U' \times_{U, t} R \ar[d] \ar[r] &
R \ar[r]^s \ar[d]_t &
U \\
U' \ar[r]^g &
U
}
}
\end{equation}
得到。见《代数空间中的群胚》引理
\ref{spaces-groupoids-lemma-restrict-groupoid}。

\begin{lemma}
\label{spaces-more-groupoids-lemma-restrict-preserves-type}
设 $S$ 为概形，$B$ 为 $S$ 上的代数空间。
设 $(U,R,s,t,c)$ 为 $B$ 上代数空间中的群胚。
设 $g:U'\to U$ 为 $B$ 上代数空间的态射。
设 $(U',R',s',t',c')$ 为 $(U,R,s,t,c)$ 经 $g$ 的限制。
\begin{enumerate}
\item 若 $s,t$ 局部有限型且 $g$ 局部有限型，则 $s',t'$ 局部有限型。
\item 若 $s,t$ 局部有限表示且 $g$ 局部有限表示，则 $s',t'$ 局部有限表示。
\item 若 $s,t$ 平坦且 $g$ 平坦，则 $s',t'$ 平坦。
\item % P09-ERR-0942
在这里补充更多内容。
\end{enumerate}
\end{lemma}

\begin{proof}
局部有限型这一性质在合成和任意基变换下稳定，见《代数空间的态射》引理
\ref{spaces-morphisms-lemma-composition-finite-type} 和
\ref{spaces-morphisms-lemma-base-change-finite-type}。
因此，由图 (\ref{spaces-more-groupoids-equation-restriction}) 可知 (1) 是显然的。
其他情形见《代数空间的态射》引理
\ref{spaces-morphisms-lemma-composition-finite-presentation}、
\ref{spaces-morphisms-lemma-base-change-finite-presentation}、
\ref{spaces-morphisms-lemma-composition-flat} 和
\ref{spaces-morphisms-lemma-base-change-flat}。
\end{proof}


\section{域上的群与域上的群胚的性质}
\label{spaces-more-groupoids-section-properties-groupoids-on-fields}

\noindent
建议读者先考察群胚概形的相应各节，见《群胚》第
\ref{groupoids-section-properties-group-schemes-field} 节和
《群胚进阶》第
\ref{more-groupoids-section-properties-groupoids-on-fields} 节。

\begin{situation}
\label{spaces-more-groupoids-situation-group-over-field}
这里 $S$ 是概形，$k$ 是 $S$ 上的域，而 $(G, m)$ 是
$\Spec(k)$ 上的群代数空间。
\end{situation}

\begin{situation}
\label{spaces-more-groupoids-situation-groupoid-on-field}
这里 $S$ 是概形，$B$ 是代数空间，而 $(U, R, s, t, c)$ 是
$B$ 上代数空间中的群胚，其中 $U = \Spec(k)$，这里 $k$ 是某个域。
\end{situation}

\noindent
注意，在情形 \ref{spaces-more-groupoids-situation-group-over-field} 中，我们得到代数空间中的群胚
\begin{equation}
\label{spaces-more-groupoids-equation-groupoid-from-group}
(\Spec(k), G, p, p, m)
\end{equation}
其中 $p : G \to \Spec(k)$ 是 $G$ 的结构态射，见
《空间中的群胚》引理
\ref{spaces-groupoids-lemma-groupoid-from-action}。
这是情形 \ref{spaces-more-groupoids-situation-groupoid-on-field} 中的情形。
本节下文将不再说明地使用这一点。

\begin{lemma}
\label{spaces-more-groupoids-lemma-groupoid-on-field-open-multiplication}
在情形 \ref{spaces-more-groupoids-situation-groupoid-on-field} 中，合成态射
$c : R \times_{s, U, t} R \to R$ 是平坦且普遍开的。
在情形 \ref{spaces-more-groupoids-situation-group-over-field} 中，群律
$m : G \times_k G \to G$ 是平坦且普遍开的。
\end{lemma}

\begin{proof}
由图 (\ref{spaces-more-groupoids-equation-pull})，合成同构于投影映射
$\text{pr}_1 : R \times_{t, U, t} R \to R$。
该投影作为平坦态射 $t$ 的基变换是平坦的；由《代数空间的态射》引理
\ref{spaces-morphisms-lemma-space-over-field-universally-open}，它也是开的。
第二个断言立即由第一个断言得到，因为在
(\ref{spaces-more-groupoids-equation-groupoid-from-group}) 中 $m$ 与 $c$ 相符。
\end{proof}

\noindent
注意，以下引理特别适用于拟分离或局部分离的代数空间
（《良态空间》引理
\ref{decent-spaces-lemma-locally-separated-decent}）。

\begin{lemma}
\label{spaces-more-groupoids-lemma-group-scheme-over-field-separated}
在情形 \ref{spaces-more-groupoids-situation-groupoid-on-field} 中，假设 $R$ 是良态空间。
那么 $R$ 是分离代数空间。在情形 \ref{spaces-more-groupoids-situation-group-over-field} 中，
假设 $G$ 是良态代数空间。那么 $G$ 是分离代数空间。
% P09-ERR-0945
\end{lemma}

\begin{proof}
我们先证明第二个断言。由《空间中的群胚》引理
\ref{spaces-groupoids-lemma-group-scheme-separated}，
我们必须证明 $e : S \to G$ 是闭浸入。
% P09-ERR-0946
这由《良态空间》引理
\ref{decent-spaces-lemma-finite-residue-field-extension-finite} 得到。

\medskip\noindent
接着证明第一个断言。为此，可以把 $B$ 替换为 $S$。
由上段，稳定子群概形 $G \to U$ 是分离的。
% P09-ERR-0947
由《空间中的群胚》引理
\ref{spaces-groupoids-lemma-diagonal}，态射
$j = (t, s) : R \to U \times_S U$ 是分离的。
由于 $U$ 是域的谱，概形 $U \times_S U$ 是仿射的
（由《概形》第 \ref{schemes-section-fibre-products} 节中纤维积的构造）。
因此 $R$ 是分离的，见《代数空间的态射》引理
\ref{spaces-morphisms-lemma-separated-over-separated}。
\end{proof}

\begin{lemma}
\label{spaces-more-groupoids-lemma-restrict-groupoid-on-field}
在情形 \ref{spaces-more-groupoids-situation-groupoid-on-field} 中。
% P09-ERR-0948
设 $k'/k$ 为域扩张，$U' = \Spec(k')$，并设
$(U', R', s', t', c')$ 是 $(U, R, s, t, c)$ 经 $U' \to U$ 的限制。
在定义图
$$
\xymatrix{
R' \ar[d] \ar[r] \ar@/_3pc/[dd]_{t'} \ar@/^1pc/[rr]^{s'} \ar@{..>}[rd] &
R \times_{s, U} U' \ar[r] \ar[d] &
U' \ar[d] \\
U' \times_{U, t} R \ar[d] \ar[r] &
R \ar[r]^s \ar[d]_t &
U \\
U' \ar[r] &
U
}
$$
中，所有态射都是满、平坦且普遍开的。此外，虚线箭头
$R' \to R$ 是仿射的。
\end{lemma}

\begin{proof}
态射 $U' \to U$ 等于 $\Spec(k') \to \Spec(k)$，
因而是仿射、满且平坦的。由《态射》引理
\ref{morphisms-lemma-scheme-over-field-universally-open}，
态射 $s, t : R \to U$ 和态射 $U' \to U$ 是普遍开的。
% P09-ERR-0949
由于 $R$ 非空且 $U$ 是域的谱，态射 $s, t : R \to U$ 是满且平坦的。
然后使用《代数空间的态射》引理
\ref{spaces-morphisms-lemma-base-change-surjective}、
\ref{spaces-morphisms-lemma-composition-surjective}、
\ref{spaces-morphisms-lemma-composition-open}、
\ref{spaces-morphisms-lemma-base-change-affine}、
\ref{spaces-morphisms-lemma-composition-affine}、
\ref{spaces-morphisms-lemma-base-change-flat} 和
\ref{spaces-morphisms-lemma-composition-flat} 即可得出结论。
\end{proof}

\begin{lemma}
\label{spaces-more-groupoids-lemma-groupoid-on-field-explain-points}
在情形 \ref{spaces-more-groupoids-situation-groupoid-on-field} 中。
% P09-ERR-0948
对任意点 $r \in |R|$，存在
\begin{enumerate}
\item 域扩张 $k'/k$，其中 $k'$ 代数闭；
\item 点 $r' : \Spec(k') \to R'$，其中
$(U', R', s', t', c')$ 是 $(U, R, s, t, c)$ 经
$\Spec(k') \to \Spec(k)$ 的限制，
\end{enumerate}
使得
\begin{enumerate}
\item 点 $r'$ 映到 $r$，所用态射为 $R' \to R$；且
\item 映射
$s' \circ r', t' \circ r' : \Spec(k') \to \Spec(k')$
是自同构。
\end{enumerate}
\end{lemma}

\begin{proof}
把 $r$ 表示为态射 $r : \Spec(K) \to R$，其中 $K$ 是某个域。
为证明该引理，必须找到一个代数闭域 $k'$ 和交换图
$$
\xymatrix{
k' & k' \ar[l]^1 & \\
k' \ar[u]^\tau & K \ar[lu]^\sigma & k \ar[l]^-s \ar[lu]_i \\
& k \ar[lu]^i \ar[u]_t
}
$$
其中 $s, t : k \to K$ 是来自 $s \circ r$ 和 $t \circ r$ 的域映射。
《群胚进阶》引理
\ref{more-groupoids-lemma-groupoid-on-field-explain-points} 的证明说明了
如何构造这样的图。
\end{proof}

\begin{lemma}
\label{spaces-more-groupoids-lemma-groupoid-on-field-move-point}
在情形 \ref{spaces-more-groupoids-situation-groupoid-on-field} 中。
% P09-ERR-0948
若 $r : \Spec(k) \to R$ 是使得
$s \circ r, t \circ r$ 为 $\Spec(k)$ 的自同构的态射，则映射
$$
R \longrightarrow R, \quad
x \longmapsto c(r, x)
$$
是自同构 $R \to R$，它把 $e$ 映到 $r$。
\end{lemma}

\begin{proof}
证明与《群胚进阶》引理
\ref{more-groupoids-lemma-groupoid-on-field-move-point} 的证明相同。
\end{proof}

\begin{lemma}
\label{spaces-more-groupoids-lemma-groupoid-on-field-geometrically-irreducible}
% P09-ERR-0950
在情形 \ref{spaces-more-groupoids-situation-groupoid-on-field} 中，代数空间 $R$ 是几何单支的。
在情形 \ref{spaces-more-groupoids-situation-group-over-field} 中，代数空间 $G$ 是几何单支的。
\end{lemma}

\begin{proof}
设 $r \in |R|$。我们必须证明 $R$ 在 $r$ 处几何单支。
结合引理 \ref{spaces-more-groupoids-lemma-restrict-groupoid-on-field} 与
《代数空间上的下降》引理
\ref{spaces-descent-lemma-descend-unibranch}，可见只需证明如下情形：
$k$ 是代数闭域，并且 $r$ 来自态射 $r : \Spec(k) \to R$，使得
$s \circ r$ 和 $t \circ r$ 都是 $\Spec(k)$ 的自同构。
由引理 \ref{spaces-more-groupoids-lemma-groupoid-on-field-move-point}，可将问题约化到
$r = e$ 是 $R$ 的单位元且 $k$ 代数闭的情形。

\medskip\noindent
假设 $r = e$ 且 $k$ 代数闭。令
$A = \mathcal{O}_{R, e}$ 为 $R$ 在 $e$ 处的 \'etale 局部环，并令
$C = \mathcal{O}_{R \times_{s, U, t} R, (e, e)}$
为 $R \times_{s, U, t} R$ 在 $(e, e)$ 处的 \'etale 局部环。
由《代数进阶》引理
\ref{more-algebra-lemma-minimal-primes-tensor-strictly-henselian}，
极小素理想 $\mathfrak q$（在 $C$ 中）与极小素理想对
$1$ 对 $1$ 地对应于 $\mathfrak p, \mathfrak p' \subset A$。
另一方面，合成律诱导平坦环同态
$$
\xymatrix{
A \ar[r]_{c^\sharp} & C & \mathfrak q \\
& A \otimes_{s^\sharp, k, t^\sharp} A \ar[u] &
\mathfrak p \otimes A + A \otimes \mathfrak p' \ar@{|}[u]
}
$$
注意，$(c^\sharp)^{-1}(\mathfrak q)$ 同时包含
$\mathfrak p$ 和 $\mathfrak p'$，因为由 (
\ref{spaces-more-groupoids-equation-diagram}) 可知图
$$
\xymatrix{
A \ar[r]_{c^\sharp} & C \\
A \otimes_{s^\sharp, k} k \ar[u] &
A \otimes_{s^\sharp, k, t^\sharp} A \ar[l]_{1 \otimes e^\sharp} \ar[u]
}
\quad\quad
\xymatrix{
A \ar[r]_{c^\sharp} & C \\
k \otimes_{k, t^\sharp} A \ar[u] &
A \otimes_{s^\sharp, k, t^\sharp} A \ar[l]_{e^\sharp \otimes 1} \ar[u]
}
$$
交换。
由于 $c^\sharp$ 平坦（因为由引理
\ref{spaces-more-groupoids-lemma-groupoid-on-field-open-multiplication}，$c$ 是平坦态射），
可见 $(c^\sharp)^{-1}(\mathfrak q)$ 是 $A$ 的极小素理想。
因此 $\mathfrak p = (c^\sharp)^{-1}(\mathfrak q) = \mathfrak p'$。
\end{proof}

\noindent
以下引理采用《代数空间的性质》第
\ref{spaces-properties-section-dimension} 节中定义的代数空间（在一点处）的维数。
我们还采用《代数空间的性质》第
\ref{spaces-properties-section-dimension-local-ring} 节中定义的局部环维数，
以及点的超越次数，见《代数空间的态射》第
\ref{spaces-morphisms-section-relative-dimension} 节。

\begin{lemma}
\label{spaces-more-groupoids-lemma-groupoid-on-field-locally-finite-type-dimension}
在情形 \ref{spaces-more-groupoids-situation-groupoid-on-field} 中，假设 $s, t$ 局部有限型。
对所有 $r \in |R|$，有
\begin{enumerate}
\item $\dim(R) = \dim_r(R)$；
\item $r$ 相对于 $\Spec(k)$ 的超越次数（经 $s$），等于
$r$ 相对于 $\Spec(k)$ 的超越次数（经 $t$）；且
\item 若 (2) 中所述的超越次数为 $0$，则
$\dim(R) = \dim(\mathcal{O}_{R, \overline{r}})$。
\end{enumerate}
\end{lemma}

\begin{proof}
设 $r \in |R|$。以 $\text{trdeg}(r/_{\!\! s}k)$ 表示
$r$ 相对于 $\Spec(k)$ 的超越次数（经 $s$）。
选取 \'etale 态射 $\varphi : V \to R$，其中 $V$ 是概形，
$v \in V$ 映到 $r$。使用引理上方提到的定义，可见
$$
\dim_r(R) = \dim_v(V) =
\dim(\mathcal{O}_{V, v}) + \text{trdeg}_{s(k)}(\kappa(v)) =
\dim(\mathcal{O}_{R, \overline{r}}) + \text{trdeg}(r/_{\!\! s}k)
$$
对 $t$ 也同样如此（第二个等号由《态射》引理
\ref{morphisms-lemma-dimension-fibre-at-a-point} 得到）。
因此可见
$\text{trdeg}(r/_{\!\! s}k) = \text{trdeg}(r/_{\!\! t}k)$，
即 (2) 成立。

\medskip\noindent
设 $k'/k$ 为域扩张。注意，限制 $R'$，即 $R$ 经 $\Spec(k')$ 的限制（见引理
\ref{spaces-more-groupoids-lemma-restrict-groupoid-on-field}）由 $R$ 经两个域态射的基变换得到。
因此，《代数空间的态射》引理
\ref{spaces-morphisms-lemma-dimension-fibre-after-base-change}
说明这一操作不改变 $R$ 在一点处的维数。
于是，为证明 (1)，由引理
\ref{spaces-more-groupoids-lemma-groupoid-on-field-explain-points}，可以假设 $r$ 由态射
$r : \Spec(k) \to R$ 表示，并且 $s \circ r$ 和 $t \circ r$
都是 $\Spec(k)$ 的自同构。
在此情形下，存在自同构 $R \to R$，它把 $r$ 映到 $e$
（引理 \ref{spaces-more-groupoids-lemma-groupoid-on-field-move-point}）。
因此有 $\dim_r(R) = \dim_e(R)$，对任意 $r$ 都成立。
按照定义，这意味着 $\dim_r(R) = \dim(R)$。

\medskip\noindent
(3) 是上面所得结果的形式推论。
\end{proof}

\begin{lemma}
\label{spaces-more-groupoids-lemma-group-over-field-locally-finite-type-dimension}
在情形 \ref{spaces-more-groupoids-situation-group-over-field} 中，假设 $G$ 局部有限型。
% P09-ERR-0951
对所有 $g \in |G|$，有
\begin{enumerate}
\item $\dim(G) = \dim_g(G)$；
\item 若 $g$ 相对于 $k$ 的超越次数为 $0$，则
$\dim(G) = \dim(\mathcal{O}_{G, \overline{g}})$。
\end{enumerate}
\end{lemma}

\begin{proof}
这立即由引理
\ref{spaces-more-groupoids-lemma-groupoid-on-field-locally-finite-type-dimension}
经由 (\ref{spaces-more-groupoids-equation-groupoid-from-group}) 得到。
\end{proof}

\begin{lemma}
\label{spaces-more-groupoids-lemma-groupoid-on-field-dimension-equal-stabilizer}
在情形 \ref{spaces-more-groupoids-situation-groupoid-on-field} 中，假设 $s, t$ 局部有限型。
令
$$
G = \Spec(k)
\times_{\Delta, \Spec(k) \times_B \Spec(k), t \times s} R
$$
为稳定子群代数空间。那么 $\dim(R) = \dim(G)$。
\end{lemma}

\begin{proof}
由于 $G$ 和 $R$ 都等维（见引理
\ref{spaces-more-groupoids-lemma-groupoid-on-field-locally-finite-type-dimension} 和
\ref{spaces-more-groupoids-lemma-group-over-field-locally-finite-type-dimension}），
只需证明 $\dim_e(R) = \dim_e(G)$。
设 $V$ 为仿射概形，$v \in V$，并设
$\varphi : V \to R$ 为概形的 \'etale 态射，使得 $\varphi(v) = e$。
注意，$V$ 是 Noether 概形，因为 $s \circ \varphi$ 作为局部有限型态射的合成
是局部有限型的，并且 $V$ 拟紧（使用《代数空间的态射》引理
\ref{spaces-morphisms-lemma-composition-finite-type}、
\ref{spaces-morphisms-lemma-etale-locally-finite-presentation} 和
\ref{spaces-morphisms-lemma-finite-presentation-finite-type}，以及
《态射》引理 \ref{morphisms-lemma-finite-type-noetherian}）。
因此 $V$ 局部连通（见《性质》引理
\ref{properties-lemma-Noetherian-topology} 和
《拓扑学》引理
\ref{topology-lemma-locally-Noetherian-locally-connected}）。
于是可以把 $V$ 替换为包含 $v$ 的连通分支（它仍是仿射的，
因为它是 $V$ 的既开又闭子概形）。令 $T = V_{red}$ 为 $V$ 的约化。
考察两个态射 $a, b : T \to \Spec(k)$，其中
$a = s \circ \varphi|_T$，$b = t \circ \varphi|_T$。
注意，$a, b$ 诱导同一个域映射 $k \to \kappa(v)$，因为
$\varphi(v) = e$！令
$k_a \subset \Gamma(T, \mathcal{O}_T)$ 为
$a^\sharp(k) \subset \Gamma(T, \mathcal{O}_T)$ 的整闭包。
类似地，令 $k_b \subset \Gamma(T, \mathcal{O}_T)$ 为
$b^\sharp(k) \subset \Gamma(T, \mathcal{O}_T)$ 的整闭包。
由《代数簇》命题
\ref{varieties-proposition-unique-base-field}，可见 $k_a = k_b$。
因而得到交换图
$$
\xymatrix{
k \ar[rd]^a \ar[rrrd] \\
& k_a = k_b \ar[r] & \Gamma(T, \mathcal{O}_T) \ar[r] & \kappa(v) \\
k \ar[ru]_b \ar[rrru]
}
$$
如上所述，两条长箭头相同。
由于 $k_a = k_b \to \kappa(v)$ 是单射，遂有两个态射 $a$ 和 $b$ 相同。
因此 $T \to R$ 分解经过 $G$。于是有 $R_{red} = G_{red}$
（在 $e$ 的一个开邻域中），这当然蕴含 $\dim_e(R) = \dim_e(G)$。
\end{proof}






\section{域上的群代数空间}
\label{spaces-more-groupoids-section-group-algebraic-spaces-over-fields}

\noindent
域上存在非分离的群代数空间，即特征零域上的
$\mathbf{G}_a/\mathbf{Z}$，见《例》第
\ref{examples-section-non-separated-group-space} 节。
事实上，域上的任意群概形都是分离的
（引理 \ref{spaces-more-groupoids-lemma-group-scheme-over-field-separated}），
所以域上的每个非分离群代数空间都不可表。
另一方面，域上的群代数空间一旦是良态的便是分离的，见引理
\ref{spaces-more-groupoids-lemma-group-scheme-over-field-separated}。
本节将证明，域上的分离群代数空间是可表的，也就是说，它是概形。

\begin{lemma}
\label{spaces-more-groupoids-lemma-group-space-scheme-over-kbar}
设 $k$ 为域，其代数闭包为 $\overline{k}$。
设 $G$ 为 $k$ 上的分离群代数空间\footnote{只需假设 $G$ 良态；
例如，由引理 \ref{spaces-more-groupoids-lemma-group-scheme-over-field-separated}，
局部分离或拟分离便足够。}。
那么 $G_{\overline{k}}$ 是概形。
\end{lemma}

\begin{proof}
由《域上的空间》引理
\ref{spaces-over-fields-lemma-when-scheme-after-base-change}，
只需证明 $G_K$ 是概形，其中 $K/k$ 是某个域扩张。
以 $G_K' \subset G_K$ 表示 $G_K$ 的概形轨迹，见《代数空间的性质》引理
\ref{spaces-properties-lemma-subscheme}。
由《代数空间的性质》命题
\ref{spaces-properties-proposition-locally-quasi-separated-open-dense-scheme}，
可见 $G_K' \subset G_K$ 是稠密开集，特别地，它非空。
选取概形 $U$ 和满 \'etale 态射 $U \to G$。
由《代数簇》引理 \ref{varieties-lemma-make-Jacobson}，
若 $K$ 是超越次数充分大的代数闭域，则 $U_K$ 是 Jacobson 概形，
且 $U_K$ 的每个闭点都是 $K$-有理点。
因此 $G_K'$ 有一个 $K$-有理点，并且只需证明每个 $K$-有理点
（属于 $G_K$）都属于 $G_K'$。
若 $g \in G_K(K)$ 是 $K$-有理点，而 $g' \in G_K'(K)$ 是概形轨迹中的
$K$-有理点，则可见 $g$ 属于 $G_K'$ 在自同构
$$
G_K \longrightarrow G_K,\quad
h \longmapsto g(g')^{-1}h
$$
下的像；这里该映射是 $G_K$ 的自同构。由于 $G_K$
作为代数空间的自同构保持 $G_K'$，
遂有 $g \in G_K'$，如所需。
\end{proof}

\begin{lemma}
\label{spaces-more-groupoids-lemma-group-space-scheme-locally-finite-type-over-k}
设 $k$ 为域，$G$ 为 $k$ 上的群代数空间。
若 $G$ 在 $k$ 上分离且局部有限型，则 $G$ 是概形。
\end{lemma}

\begin{proof}
这由引理 \ref{spaces-more-groupoids-lemma-group-space-scheme-over-kbar}、
《群胚》引理 \ref{groupoids-lemma-points-in-affine} 和
《域上的空间》引理
\ref{spaces-over-fields-lemma-scheme-over-algebraic-closure-enough-affines}
得到。
\end{proof}

\begin{proposition}
\label{spaces-more-groupoids-proposition-group-space-scheme-over-field}
设 $k$ 为域，$G$ 为 $k$ 上的群代数空间。
若 $G$ 分离，则 $G$ 是概形。
\end{proposition}

\begin{proof}
本引理推广了引理
\ref{spaces-more-groupoids-lemma-group-space-scheme-locally-finite-type-over-k}
（它涵盖了实践中所关心的全部情形）。
% P09-ERR-0952
该证明与《域上的空间》引理
\ref{spaces-over-fields-lemma-scheme-over-algebraic-closure-enough-affines}
的证明非常相似；后者用在引理
\ref{spaces-more-groupoids-lemma-group-space-scheme-locally-finite-type-over-k} 的证明中。
我们建议读者先阅读那个证明。

\medskip\noindent
由引理 \ref{spaces-more-groupoids-lemma-group-space-scheme-over-kbar}，基变换
$G_{\overline{k}}$ 是概形。
设 $K/k$ 为超越次数非常大的纯超越扩张。
由《域上的空间》引理
\ref{spaces-over-fields-lemma-scheme-after-purely-transcendental-base-change}，
只需证明 $G_K$ 是概形。
令 $K^{perf}$ 为 $K$ 的完美闭包。由《域上的空间》引理
\ref{spaces-over-fields-lemma-scheme-after-purely-inseparable-base-change}，
只需证明 $G_{K^{perf}}$ 是概形。
令 $K \subset K^{perf} \subset \overline{K}$ 为 $K$ 的代数闭包。
% P09-ERR-0953
可以选取嵌入 $\overline{k} \to \overline{K}$（在 $k$ 上），使得
$G_{\overline{K}}$ 是概形 $G_{\overline{k}}$ 经
$\overline{k} \to \overline{K}$ 的基变换。
由《代数簇》引理 \ref{varieties-lemma-make-Jacobson}，
可见 $G_{\overline{K}}$ 是 Jacobson 概形，且它的每个闭点的剩余域都是
$\overline{K}$。

\medskip\noindent
由于 $G_{\overline{K}} \to G_{K^{perf}}$ 是满的，只需证明
像点 $g \in |G_{K^{perf}}|$（它来自 $G_{\overline{K}}$ 的任意闭点）
属于 $G_K$ 的概形轨迹。
% P09-ERR-0954
特别地，可以把 $g$ 表示为态射 $g : \Spec(L) \to G_{K^{perf}}$，
其中 $L/K^{perf}$ 是可分代数扩张（例如可以取 $L = \overline{K}$）。
于是概形
\begin{align*}
T & = \Spec(L) \times_{G_{K^{perf}}} G_{\overline{K}} \\
& =
\Spec(L) \times_{\Spec(K^{perf})} \Spec(\overline{K}) \\
& =
\Spec(L \otimes_{K^{perf}} \overline{K})
\end{align*}
是某个 $\overline{K}$-代数的谱；该代数是有限个
$\overline{K}$ 的副本之积所成代数的滤过余极限。
因此，由《群胚》引理 \ref{groupoids-lemma-compact-set-in-affine}，
可以找到仿射开集 $W \subset G_{\overline{K}}$，它包含
$g_{\overline{K}} : T \to G_{\overline{K}}$ 的像。

\medskip\noindent
选取拟紧开集 $V \subset G_{K^{perf}}$，使其包含 $W$ 的像。
由《域上的空间》引理
\ref{spaces-over-fields-lemma-when-scheme-after-base-change}，
可见 $V_{K'}$ 是概形，其中 $K'/K^{perf}$ 是某个有限扩张。
扩大 $K'$ 后，可以假设存在仿射开集
$U' \subset V_{K'} \subset G_{K'}$，其到 $\overline{K}$ 的基变换恢复 $W$
（使用 $V_{\overline{K}}$ 是概形 $V_{K''}$ 的极限，其中
$K' \subset K'' \subset \overline{K}$ 是有限扩张，并使用《极限》引理
\ref{limits-lemma-descend-opens} 和
\ref{limits-lemma-limit-affine}）。可以假设 $K'/K^{perf}$ 是 Galois 扩张
（取正规闭包《域》引理 \ref{fields-lemma-normal-closure}，并使用
$K^{perf}$ 是完美的这一事实）。
% P09-ERR-0955
置 $H = \text{Gal}(K'/K^{perf})$。
按照构造，$H$-不变闭子概形
$\Spec(L) \times_{G_{K^{perf}}} G_{K'}$ 包含在 $U'$ 中。
由《域上的空间》引理
\ref{spaces-over-fields-lemma-base-change-by-Galois} 和
\ref{spaces-over-fields-lemma-when-quotient-scheme-at-point}，结论得证。
\end{proof}






\section{群上不存在有理曲线}
\label{spaces-more-groupoids-section-no-rational-curves}

\noindent
本节证明，从 $\mathbf{P}^1$ 到域上局部有限型群代数空间的态射
不存在非常值的。

\begin{lemma}
\label{spaces-more-groupoids-lemma-factor-through-over-open}
设 $S$ 为概形，$B$ 为 $S$ 上的代数空间。
设 $f : X \to Y$ 和 $g : X \to Z$ 为 $B$ 上代数空间的态射。假设
\begin{enumerate}
\item $Y \to B$ 分离；
\item $g$ 满、平坦且局部有限表示；
\item 存在概形论稠密开集 $V \subset Z$，使得
$f|_{g^{-1}(V)} : g^{-1}(V) \to Y$ 分解经过 $V$。
\end{enumerate}
那么 $f$ 分解经过 $g$。
\end{lemma}

\begin{proof}
置 $R = X \times_Z X$。由 (2)，可见作为层有 $Z = X/R$。
此外，(2) 蕴含 $V$ 在 $R$ 中的逆像在 $R$ 中概形论稠密
（《代数空间的态射》引理
\ref{spaces-morphisms-lemma-flat-morphism-scheme-theoretically-dense-open}）。
可见两个合成 $R \to X \to Y$ 相等，这是由《代数空间的态射》引理
\ref{spaces-morphisms-lemma-equality-of-morphisms} 得到的。
% P09-ERR-0956
引理得证。
\end{proof}

\begin{lemma}
\label{spaces-more-groupoids-lemma-quotient-power-P1}
\begin{slogan}
从域上非空个射影直线的乘积到域上分离有限型代数空间的态射，
经投影到某个射影直线的乘积后，分解为有限态射。
\end{slogan}
设 $k$ 为域，$n \geq 1$，并设 $(\mathbf{P}^1_k)^n$
为 $n$ 重自乘积（在 $\Spec(k)$ 上）。设
$f : (\mathbf{P}^1_k)^n \to Z$ 为 $k$ 上代数空间的态射。
若 $Z$ 在 $k$ 上分离且有限型，则 $f$ 分解为
$$
(\mathbf{P}^1_k)^n \xrightarrow{projection}
(\mathbf{P}^1_k)^m \xrightarrow{finite} Z.
$$
\end{lemma}

\begin{proof}
可以假设 $k$ 代数闭（略去细节）；这样做只是为了可以使用有理点来论证，
但读者若愿意，也可以绕开这一点。证明中的乘积都在 $k$ 上取。
$(\mathbf{P}^1_k)^n$ 的自同构群代数空间包含
$G = (\text{GL}_{2, k})^n$。若 $C \subset (\mathbf{P}^1_k)^n$
是映到一点的闭子簇（特别地，它在 $k$ 上不可约），则可将
《代数空间的态射进阶》引理
\ref{spaces-more-morphisms-lemma-flat-proper-family-cannot-collapse-fibre}
应用于态射
$$
G \times C \to G \times Z,\quad (g, c) \mapsto (g, f(g \cdot c))
$$
（在 $G$ 上）。因此，$g(C)$ 映到一点，只要
$g \in G(k)$ 位于某个 Zariski 开集 $U \subset G$ 中。假设
$x = (x_1, \ldots, x_n)$、$y = (y_1, \ldots, y_n)$
是 $k$-值点，且属于 $(\mathbf{P}^1_k)^n$。令
$I \subset \{1, \ldots, n\}$ 为指标 $i$ 之集，其中 $x_i = y_i$。
那么
$$
\{g(x) \mid g(y) = y,\ g \in U(k)\}
$$
在投影的纤维中 Zariski 稠密，其中该投影为
$\pi_I : (\mathbf{P}^1_k)^n \to \prod_{i \in I} \mathbf{P}^1_k$
（练习）。因此，若不同的 $x, y \in C(k)$，则可得 $f$ 把
$\pi_I$ 的、包含 $x, y$ 的整个纤维映到同一点。
此外，$U(k)$ 作用下 $C$ 的轨道与 $\pi_I$ 的某个 Zariski 开集中的纤维相交。
由引理 \ref{spaces-more-groupoids-lemma-factor-through-over-open}，态射 $f$ 分解经过 $\pi_I$。
有限次重复这一过程后，便达到 $f$ 在 $k$-点上方的全部纤维都是有限集的阶段。
此时，$f$ 是有限态射，这是由《代数空间的态射进阶》引理
\ref{spaces-more-morphisms-lemma-proper-finite-fibre-finite-in-neighbourhood}
以及 $k$-点在 $Z$ 中稠密这一事实（《域上的空间》引理
\ref{spaces-over-fields-lemma-smooth-separable-closed-points-dense}）得到的。
\end{proof}

\begin{lemma}
\label{spaces-more-groupoids-lemma-no-nonconstant-morphism-from-P1-to-group}
\begin{slogan}
群上不存在完备有理曲线。
\end{slogan}
设 $k$ 为域，$G$ 为 $k$ 上分离且局部有限型的群代数空间。
不存在非常值态射 $f : \mathbf{P}^1_k \to G$（在 $\Spec(k)$ 上）。
\end{lemma}

\begin{proof}
假设 $f$ 非常值。考察态射
$$
\mathbf{P}^1_k \times_{\Spec(k)} \ldots \times_{\Spec(k)} \mathbf{P}^1_k
\longrightarrow G,
\quad (t_1, \ldots, t_n) \longmapsto f(g_1) \ldots f(g_n)
$$
% P09-ERR-0957
其中右端使用群中的乘法。由引理 \ref{spaces-more-groupoids-lemma-quotient-power-P1}
及 $f$ 非常值这一假设，此态射在其像上是有限的。
因此 $\dim(G) \geq n$ 对所有 $n$ 都成立；这不可能，因为由引理
\ref{spaces-more-groupoids-lemma-group-over-field-locally-finite-type-dimension}，
$G$ 在 $k$ 上局部有限型。
\end{proof}




\section{态射的有限部分}
\label{spaces-more-groupoids-section-finite}

\noindent
设 $S$ 为概形。
设 $f : X \to Y$ 为 $S$ 上代数空间的态射。
对代数空间或概形 $T$（在 $S$ 上），考察偶对 $(a, Z)$，其中
\begin{equation}
\label{spaces-more-groupoids-equation-finite-conditions}
\begin{matrix}
a : T \to Y\text{ 是态射（相对于 }S, \\
Z \subset T \times_Y X\text{ 是开子空间} \\
\text{并且 }\text{pr}_0|_Z : Z \to T\text{ 是有限态射。}
\end{matrix}
\end{equation}
假设 $h : T' \to T$ 为 $S$ 上代数空间的态射，
而 $(a, Z)$ 是 $T$ 上满足 (
\ref{spaces-more-groupoids-equation-finite-conditions}) 的偶对。置
$a' = a \circ h$ 和
$Z' = (h \times \text{id}_X)^{-1}(Z) = T' \times_T Z$。
那么 $(a', Z')$ 是 $T'$ 上满足 (
\ref{spaces-more-groupoids-equation-finite-conditions}) 的偶对。
这是因为有限态射在基变换下保持，见《代数空间的态射》引理
\ref{spaces-morphisms-lemma-base-change-integral}。
由此得到函子
\begin{equation}
\label{spaces-more-groupoids-equation-finite}
\begin{matrix}
(X/Y)_{fin} : &
(\Sch/S)^{opp} &
\longrightarrow &
\textit{Sets} \\
& T & \longmapsto &
\{(a, Z)\text{ 如上}\}
\end{matrix}
\end{equation}
在应用中，我们主要关心函子 $(X/Y)_{fin}$ 在 $f$ 分离且局部有限型
时的情形。为了理解这里的主题，可以先看注
\ref{spaces-more-groupoids-remark-finite-quasi-finite-separated-morphism-schemes}。

\begin{lemma}
\label{spaces-more-groupoids-lemma-finite-sheaf}
设 $S$ 为概形。
设 $f : X \to Y$ 为 $S$ 上代数空间的态射。那么
\begin{enumerate}
\item 预层 $(X/Y)_{fin}$ 对 fppf 拓扑满足层条件。
\item 若 $T$ 为 $S$ 上的代数空间，则存在典范双射
$$
\Mor_{\Sh((\Sch/S)_{fppf})}(T, (X/Y)_{fin})
=
\{(a, Z)\text{ 满足 \ref{spaces-more-groupoids-equation-finite-conditions}}\}
$$
\end{enumerate}
\end{lemma}

\begin{proof}
设 $T$ 为 $S$ 上的代数空间。
设 $\{T_i \to T\}$ 为 fppf 覆盖（由代数空间组成）。
设 $s_i = (a_i, Z_i)$ 为 $T_i$ 上满足
\ref{spaces-more-groupoids-equation-finite-conditions} 的偶对，并且
$s_i|_{T_i \times_T T_j} = s_j|_{T_i \times_T T_j}$。
首先，这特别蕴含 $a_i$ 和 $a_j$ 定义同一个态射
$T_i \times_T T_j \to Y$。由《代数空间上的下降》引理
\ref{spaces-descent-lemma-fpqc-universal-effective-epimorphisms}，
可得存在唯一态射 $a : T \to Y$，使得 $a_i$ 等于合成
$T_i \to T \to Y$。
其次，这蕴含 $Z_i \subset T_i \times_Y X$ 是开子空间，并且它们在
$(T_i \times_T T_j) \times_Y X$ 中的逆像相等。
由于 $\{T_i \times_Y X \to T \times_Y X\}$ 是 fppf 覆盖，
可得存在唯一开子空间 $Z \subset T \times_Y X$，其限制回
$Z_i$（在 $T_i$ 上），见《代数空间上的下降》引理
\ref{spaces-descent-lemma-open-fpqc-covering}。
我们断言投影 $Z \to T$ 是有限态射。
这是因为有限性对 fpqc 拓扑是局部的，见《代数空间上的下降》引理
\ref{spaces-descent-lemma-descending-property-finite}。

\medskip\noindent
注意，上一段的结果特别蕴含 (1)。

\medskip\noindent
设 $T$ 为 $S$ 上的代数空间。为证明 (2)，我们将在所显示的集合之间
构造互逆映射。下文所谓“偶对”均指满足条件
\ref{spaces-more-groupoids-equation-finite-conditions} 的偶对。

\medskip\noindent
设 $v : T \to (X/Y)_{fin}$ 为自然变换。
选取概形 $U$ 和满 \'etale 态射 $p : U \to T$。
那么 $v(p) \in (X/Y)_{fin}(U)$ 对应于偶对 $(a_U, Z_U)$（在 $U$ 上）。
令 $R = U \times_T U$，投影为 $t, s : R \to U$。
由于 $v$ 是函子的变换，可见 $(a_U, Z_U)$ 经 $s$ 和 $t$ 的拉回相同。
因此，因为 $\{U \to T\}$ 是 fppf 覆盖，可以应用第一段的结果，
从而推出存在唯一偶对 $(a, Z)$（在 $T$ 上）。
% P09-ERR-0958

\medskip\noindent
反之，设 $(a, Z)$ 为 $T$ 上的偶对。
令 $U \to T$、$R = U \times_T U$ 和 $t, s : R \to U$ 如上。
那么限制 $(a, Z)|_U$ 由 Yoneda 引理给出函子变换
$v : h_U \to (X/Y)_{fin}$
（《范畴》引理 \ref{categories-lemma-yoneda}）。
由于两个拉回 $s^*(a, Z)|_U$ 和 $t^*(a, Z)|_U$ 相等，
可见 $v$ 使两个映射 $h_t, h_s : h_R \to h_U$ 余等化。
由《代数空间》引理 \ref{spaces-lemma-space-presentation}，
$T = U/R$ 是 fppf 商层；又由 (1)，$(X/Y)_{fin}$ 是 fppf 层，
故 $v$ 分解经过映射 $T \to (X/Y)_{fin}$。

\medskip\noindent
略去验证上述两个构造互逆的过程。
\end{proof}

\begin{lemma}
\label{spaces-more-groupoids-lemma-finite-open}
设 $S$ 为概形。考察交换图
$$
\xymatrix{
X' \ar[rr]_j \ar[rd] & & X \ar[ld] \\
& Y
}
$$
，其中所有代数空间都在 $S$ 上。若 $j$ 是开浸入，则存在层的典范单射
$j : (X'/Y)_{fin} \to (X/Y)_{fin}$。
\end{lemma}

\begin{proof}
若 $(a, Z)$ 是 $T$ 上的偶对（对于 $X'/Y$），则
$(a, j(Z))$ 是 $T$ 上的偶对（对于 $X/Y$）。
\end{proof}

\begin{lemma}
\label{spaces-more-groupoids-lemma-finite-lives-on-locally-quasi-finite-part}
设 $S$ 为概形。
设 $f : X \to Y$ 为 $S$ 上局部有限型的代数空间态射。
设 $X' \subset X$ 为使 $f$ 局部拟有限的最大开子空间，见
《代数空间的态射》引理
\ref{spaces-morphisms-lemma-locally-finite-type-quasi-finite-part}。
那么 $(X/Y)_{fin} = (X'/Y)_{fin}$。
\end{lemma}

\begin{proof}
引理 \ref{spaces-more-groupoids-lemma-finite-open} 给出单射
$(X'/Y)_{fin} \to (X/Y)_{fin}$。
《代数空间的态射》引理
\ref{spaces-morphisms-lemma-locally-finite-type-quasi-finite-part}
保证 $X'$ 的形成与基变换交换。
因此，一切归结为证明：若 $Z \subset X$ 是开子空间，并且
$f|_Z : Z \to Y$ 是有限态射，则 $Z \subset X'$。
这成立是因为有限态射局部拟有限，见《代数空间的态射》引理
\ref{spaces-morphisms-lemma-finite-quasi-finite}。
\end{proof}

\begin{lemma}
\label{spaces-more-groupoids-lemma-finite-separated}
设 $S$ 为概形。
设 $f : X \to Y$ 为 $S$ 上代数空间的态射。
设 $T$ 为 $S$ 上的代数空间，并设 $(a, Z)$ 为
\ref{spaces-more-groupoids-equation-finite-conditions} 中的偶对。
若 $f$ 分离，则 $Z$ 在 $T \times_Y X$ 中闭。
\end{lemma}

\begin{proof}
由《代数空间的态射》引理
\ref{spaces-morphisms-lemma-finite-proper}，
代数空间的有限态射是普遍闭的。
由于 $f$ 分离，态射 $T \times_Y X \to T$ 也分离，见
《代数空间的态射》引理
\ref{spaces-morphisms-lemma-base-change-separated}。
因此，$Z$ 的闭性由《代数空间的态射》引理
\ref{spaces-morphisms-lemma-universally-closed-permanence} 得到。
\end{proof}

\begin{remark}
\label{spaces-more-groupoids-remark-finite-monoid}
设 $f : X \to Y$ 为代数空间的分离态射。
层 $(X/Y)_{fin}$ 带有自然映射 $(X/Y)_{fin} \to Y$：它把
偶对 $(a, Z) \in (X/Y)_{fin}(T)$ 映到元素 $a \in Y(T)$。
可以使用引理 \ref{spaces-more-groupoids-lemma-finite-separated} 定义运算
$$
\star_i : (X/Y)_{fin} \times_Y (X/Y)_{fin} \longrightarrow (X/Y)_{fin}
$$
，规则为
\begin{align*}
\star_1 : ((a, Z_1), (a, Z_2)) & \longmapsto (a, Z_1 \cup Z_2) \\
\star_2 : ((a, Z_1), (a, Z_2)) & \longmapsto (a, Z_1 \cap Z_2) \\
\star_3 : ((a, Z_1), (a, Z_2)) & \longmapsto (a, Z_1 \setminus Z_2) \\
\star_4 : ((a, Z_1), (a, Z_2)) & \longmapsto (a, Z_2 \setminus Z_1).
\end{align*}
这可行的原因是 $Z_1 \cap Z_2$ 在 $Z_1$ 和 $Z_2$ 中都既开又闭
（这也蕴含 $Z_1 \cup Z_2$ 是另外三部分的不交并）。
因此，可以把 $(X/Y)_{fin}$ 看成无幺 $\mathbf{F}_2$-代数（在 $Y$ 上），
其乘法由 $ss' = \star_2(s, s')$ 给出，而加法由
$$
s + s' = \star_1(\star_3(s, s'), \star_4(s, s'))
$$
给出；这无非是取对称差。
注意，在这个代数层中 $0 = (1_Y, \emptyset)$，并且
$s + s = 0$ 对任意局部截面 $s$ 都成立。
若 $f : X \to Y$ 有限，则该代数有单位元，即
$1 = (1_Y, X)$，并且 $\star_3(s, s') = s(1 + s')$，
$\star_4(s, s') = (1 + s)s'$。
\end{remark}

\begin{remark}
\label{spaces-more-groupoids-remark-finite-quasi-finite-separated-morphism-schemes}
设 $f : X \to Y$ 为概形的分离、局部拟有限态射。
在此情形下，层 $(X/Y)_{fin}$ 与层 $f_!\mathbf{F}_2$ 密切相关
（在 $Y_\etale$ 上；在这里插入未来的参考文献）。
% P09-ERR-0959
具体地，若 $V \to Y$ 是 \'etale 的，且
$s \in \Gamma(V, f_!\mathbf{F}_2)$，则
$s \in \Gamma(V \times_Y X, \mathbf{F}_2)$ 是一个截面，
其支集 $Z = \text{Supp}(s)$ 在 $V$ 上真。
由于 $f$ 还局部拟有限，可见投影 $Z \to V$ 实际上是有限态射。
由于常值阿贝尔层的截面的支集是开的，可见偶对
$(V \to Y, \text{Supp}(s))$ 满足
\ref{spaces-more-groupoids-equation-finite-conditions}。
事实上，在此情形下有
$f_!\mathbf{F}_2 \cong (X/Y)_{fin}|_{Y_\etale}$；
这也解释了注 \ref{spaces-more-groupoids-remark-finite-monoid} 中引入的
$\mathbf{F}_2$-代数结构。
\end{remark}

\begin{lemma}
\label{spaces-more-groupoids-lemma-finite-diagonal}
设 $S$ 为概形。
设 $f : X \to Y$ 为 $S$ 上代数空间的态射。
$(X/Y)_{fin} \to Y$ 的对角态射
$$
(X/Y)_{fin} \longrightarrow (X/Y)_{fin} \times_Y (X/Y)_{fin}
$$
可表（由概形），且为开浸入；而“绝对”对角态射
$$
(X/Y)_{fin} \longrightarrow (X/Y)_{fin} \times (X/Y)_{fin}
$$
可表（由概形）。
\end{lemma}

\begin{proof}
第二个断言由第一个断言得到，因为绝对对角态射是相对对角态射与
$Y$ 的对角态射（可由概形表）的一个基变换之合成，见《代数空间》第
\ref{spaces-section-representable} 节。
为证明第一个断言，必须证明如下事实：给定概形 $T$，以及
具有相同第一分量的两个偶对 $(a, Z_1)$ 和 $(a, Z_2)$（在 $T$ 上），
它们满足 \ref{spaces-more-groupoids-equation-finite-conditions}，则存在开子概形
$V \subset T$，具有如下性质：对任意概形态射 $h : T' \to T$，有
$$
h(T') \subset V \Leftrightarrow
\Big(T' \times_T Z_1 = T' \times_T Z_2
\text{，视为以下空间的子空间：}T' \times_Y X\Big)
$$
现在构造 $V$。注意，$Z_1 \cap Z_2$ 在 $Z_1$ 中开，也在 $Z_2$ 中开。
由于 $\text{pr}_0|_{Z_i} : Z_i \to T$ 有限，故为真态射（见
《代数空间的态射》引理 \ref{spaces-morphisms-lemma-finite-proper}），
可见
$$
E =
\text{pr}_0|_{Z_1}\left(Z_1 \setminus Z_1 \cap Z_2)\right)
\cup
\text{pr}_0|_{Z_2}\left(Z_2 \setminus Z_1 \cap Z_2)\right)
$$
% P09-ERR-0961
在 $T$ 中闭。现在显然 $V = T \setminus E$ 满足要求。
\end{proof}

\begin{lemma}
\label{spaces-more-groupoids-lemma-finite-criterion-etale}
设 $S$ 为概形。
设 $f : X \to Y$ 为 $S$ 上代数空间的态射。
假设 $U$ 是概形，$U \to Y$ 是 \'etale 态射，并且
$Z \subset U \times_Y X$ 是在 $U$ 上有限的开子空间。
那么诱导态射 $U \to (X/Y)_{fin}$ 是 \'etale 的。
\end{lemma}

\begin{proof}
这可由引理 \ref{spaces-more-groupoids-lemma-finite-diagonal} 中对角态射的描述形式地得到，
但由于这是理论发展中的重要一步，我们把它写出来。
必须验证：对任意概形 $T$（在 $S$ 上）和态射
$T \to (X/Y)_{fin}$，投影映射
$$
T \times_{(X/Y)_{fin}} U \longrightarrow T
$$
是 \'etale 的。注意
$$
T \times_{(X/Y)_{fin}} U
=
(X/Y)_{fin} \times_{((X/Y)_{fin} \times_Y (X/Y)_{fin})} (T \times_Y U)
$$
应用引理 \ref{spaces-more-groupoids-lemma-finite-diagonal} 的结果，可见
$T \times_{(X/Y)_{fin}} U$ 由 $T \times_Y U$ 的一个开子概形表示。
由于投影 $T \times_Y U \to T$ 由《代数空间的态射》引理
\ref{spaces-morphisms-lemma-base-change-etale} 是 \'etale 的，结论得证。
\end{proof}

\begin{lemma}
\label{spaces-more-groupoids-lemma-finite-pullback}
设 $S$ 为概形。设
$$
\xymatrix{
X' \ar[d] \ar[r] & X \ar[d] \\
Y' \ar[r] & Y
}
$$
为 $S$ 上代数空间的纤维积方块。那么
$$
\xymatrix{
(X'/Y')_{fin} \ar[d] \ar[r] & (X/Y)_{fin} \ar[d] \\
Y' \ar[r] & Y
}
$$
是 $(\Sch/S)_{fppf}$ 上层的纤维积方块。
\end{lemma}

\begin{proof}
由定义立即可知，层 $(X'/Y')_{fin}$ 等于层
$Y' \times_Y (X/Y)_{fin}$。
\end{proof}

\begin{lemma}
\label{spaces-more-groupoids-lemma-finite-surjective-etale-cover}
设 $S$ 为概形。
设 $f : X \to Y$ 为 $S$ 上代数空间的态射。
若 $f$ 分离且局部拟有限，则存在概形 $U$（在 $Y$ 上 \'etale），
以及满 \'etale 态射 $U \to (X/Y)_{fin}$（在 $Y$ 上）。
\end{lemma}

\begin{proof}
由引理 \ref{spaces-more-groupoids-lemma-finite-diagonal} 关于 $(X/Y)_{fin}$ 的对角态射的结果，
这个断言是有意义的，见《代数空间》引理
\ref{spaces-lemma-representable-diagonal}。
设 $V$ 为概形，并设 $V \to Y$ 为满 \'etale 态射。
由引理 \ref{spaces-more-groupoids-lemma-finite-pullback}，态射
$(V \times_Y X/V)_{fin} \to (X/Y)_{fin}$ 是映射 $V \to Y$ 的基变换，
因而满且 \'etale，见《代数空间》引理
\ref{spaces-lemma-base-change-representable-transformations-property}。
因此，只需对 $(V \times_Y X/V)_{fin}$ 证明该引理。
（这里隐式使用如下事实：函子的可表、满且 \'etale 的变换之合成
仍然可表、满且 \'etale；见《代数空间》引理
\ref{spaces-lemma-composition-representable-transformations} 和
\ref{spaces-lemma-composition-representable-transformations-property}，以及
《态射》引理 \ref{morphisms-lemma-composition-surjective} 和
\ref{morphisms-lemma-composition-etale}。）
注意，分离和局部拟有限这两个性质在基变换下保持，见
《代数空间的态射》引理
\ref{spaces-morphisms-lemma-base-change-separated} 和
\ref{spaces-morphisms-lemma-base-change-quasi-finite}。
因此，$V \times_Y X \to V$ 也分离且局部拟有限，并且由
《代数空间的态射》命题
\ref{spaces-morphisms-proposition-locally-quasi-finite-separated-over-scheme}，
可见 $V \times_Y X$ 也是概形。
于是可以假设 $f : X \to Y$ 是概形的分离、局部拟有限态射。

\medskip\noindent
选取点 $y \in Y$。选取点 $x_1, \ldots, x_n \in X$；它们位于
$y$ 上方。选取 \'etale 邻域
$a : (U, u) \to (Y, y)$ 和分解
$$
U \times_S X =
W \amalg
\ \coprod\nolimits_{i = 1, \ldots, n}
\ \coprod\nolimits_{j = 1, \ldots, m_j}
V_{i, j}
$$
，如《态射进阶》引理
\ref{more-morphisms-lemma-etale-splits-off-quasi-finite-part-technical-variant}
中所述。
% P09-ERR-0960
任取子集
$$
I \subset \{(i, j) \mid 1 \leq i \leq n, \ 1 \leq j \leq m_i\}.
$$
给定这些选择，得到偶对 $(a, Z)$，其中
$Z = \bigcup_{(i, j) \in I} V_{i, j}$，
它满足条件 \ref{spaces-more-groupoids-equation-finite-conditions}。换言之，我们得到态射
$U \to (X/Y)_{fin}$。此态射的构造依赖于上面选取的全部数据，
所以实际上应当写成
$$
U(y, n, x_1, \ldots, x_n, a, I) \longrightarrow (X/Y)_{fin}
$$
由引理 \ref{spaces-more-groupoids-lemma-finite-criterion-etale}，此态射是 \'etale 的。

\medskip\noindent
断言：所有这些态射的不交并满射到 $(X/Y)_{fin}$。
显然，若断言成立，则引理成立。

\medskip\noindent
为证明满性，必须证明如下事实（见《代数空间》注
\ref{spaces-remark-warning}）：给定概形 $T$，并设它在 $S$ 上；再给定点 $t \in T$
以及映射 $T \to (X/Y)_{fin}$，可以找到如上的数据
$(y, n, x_1, \ldots, x_n, a, I)$，使得 $t$ 属于投影映射
$$
U(y, n, x_1, \ldots, x_n, a, I) \times_{(X/Y)_{fin}} T \longrightarrow T.
$$
的像。为证明这一点，显然可以把 $T$ 替换为
$\Spec(\overline{\kappa(t)})$，并把 $T \to (X/Y)_{fin}$ 替换为合成
$\Spec(\overline{\kappa(t)}) \to T \to (X/Y)_{fin}$。
换言之，可以假设 $T$ 是代数闭域的谱。

\medskip\noindent
设 $T = \Spec(k)$ 为代数闭域 $k$ 的谱。
态射 $T \to (X/Y)_{fin}$ 由满足条件
\ref{spaces-more-groupoids-equation-finite-conditions} 的偶对 $(T \to Y, Z)$ 给出。
图示如下：
$$
\xymatrix{
& Z \ar[d] \ar[r] & X \ar[d] \\
\Spec(k) \ar@{=}[r] & T \ar[r] & Y
}
$$
设 $y \in Y$ 为 $T \to Y$ 的像点。
由于 $Z$ 在 $k$ 上有限，它只有有限多个点。
因此存在有限多个点 $x_1, \ldots, x_n \in X$，使得 $Z$ 在 $X$ 中的像
包含在 $\{x_1, \ldots, x_n\}$ 中。
如上选取 $a : (U, u) \to (Y, y)$，使其适合于 $y$ 和
$x_1, \ldots, x_n$，得到图
$$
\xymatrix{
W \amalg
\ \coprod\nolimits_{i = 1, \ldots, n}
\ \coprod\nolimits_{j = 1, \ldots, m_j}
V_{i, j} \ar[d] \ar[r] &
X \ar[d] \\
U \ar[r] & Y.
}
$$
由于 $k$ 代数闭且 $\kappa(y) \subset \kappa(u)$ 是有限可分扩张，
可以使态射 $T = \Spec(k) \to Y$ 分解经过态射
$u = \Spec(\kappa(u)) \to \Spec(\kappa(y)) = y \subset Y$。
在这一选择下得到交换图
$$
\xymatrix{
Z \ar[d] \ar[r] &
W \amalg
\ \coprod\nolimits_{i = 1, \ldots, n}
\ \coprod\nolimits_{j = 1, \ldots, m_j}
V_{i, j} \ar[d] \ar[r] &
X \ar[d] \\
\Spec(k) \ar[r] &
U \ar[r] & Y
}
$$
已知左上箭头的像落在 $\coprod V_{i, j}$ 中。
还要回忆，$Z$ 是 $\Spec(k) \times_Y X$ 的开子概形；
这是由 $(X/Y)_{fin}$ 的定义给出的，并且右侧方块是纤维积方块。
因此可见
$$
Z \subset
\coprod\nolimits_{i = 1, \ldots, n}\ \coprod\nolimits_{j = 1, \ldots, m_j}
\Spec(k) \times_U V_{i, j}
$$
是开子概形。按照构造（见《态射进阶》引理
\ref{more-morphisms-lemma-etale-splits-off-quasi-finite-part-technical-variant}），
每个 $V_{i, j}$ 都有唯一一个点 $v_{i, j}$，该点位于 $u$ 上方，
其剩余域扩张 $\kappa(v_{i, j})/\kappa(u)$ 是纯不可分的。
因此，每个概形 $\Spec(k) \times_U V_{i, j}$ 恰有一个点。
所以可见
$$
Z = \coprod\nolimits_{(i, j) \in I} \Spec(k) \times_U V_{i, j}
$$
，其中唯一确定的子集为
$I \subset \{(i, j) \mid 1 \leq i \leq n, \ 1 \leq j \leq m_i\}$。
展开定义可知
$$
U(y, n, x_1, \ldots, x_n, a, I) \times_{(X/Y)_{fin}} T
$$
对上面找到的 $I$ 非空，如所需。
\end{proof}

\begin{proposition}
\label{spaces-more-groupoids-proposition-finite-algebraic-space}
设 $S$ 为概形。
设 $f : X \to Y$ 为 $S$ 上分离且局部有限型的代数空间态射。
那么 $(X/Y)_{fin}$ 是代数空间。此外，态射
$(X/Y)_{fin} \to Y$ 是 \'etale 的。
\end{proposition}

\begin{proof}
由引理 \ref{spaces-more-groupoids-lemma-finite-lives-on-locally-quasi-finite-part}，
可以把 $X$ 替换为在 $Y$ 上局部拟有限的开子概形。
% P09-ERR-0963
因此，可以假设 $f$ 分离且局部拟有限。
我们将验证《代数空间》定义
\ref{spaces-definition-algebraic-space} 的三个条件。
条件 (1) 由引理 \ref{spaces-more-groupoids-lemma-finite-sheaf} 得到。
条件 (2) 由引理 \ref{spaces-more-groupoids-lemma-finite-diagonal} 得到。
最后，条件 (3) 由引理 \ref{spaces-more-groupoids-lemma-finite-surjective-etale-cover} 得到。
所以 $(X/Y)_{fin}$ 是代数空间。
此外，该引理说明存在交换图
$$
\xymatrix{
U \ar[rr] \ar[rd] & & (X/Y)_{fin} \ar[ld] \\
& Y
}
$$
，其中水平箭头满且 \'etale，而东南箭头 \'etale。
由《代数空间的性质》引理
\ref{spaces-properties-lemma-etale-local}，这蕴含西南箭头也 \'etale。
% P09-ERR-0964
\end{proof}

\begin{remark}
\label{spaces-more-groupoids-remark-warning}
命题 \ref{spaces-more-groupoids-proposition-finite-algebraic-space} 中 $f$ 分离的条件不能删去。
一个例子是取 $X$ 为原点加倍的仿射直线，见《概形》例
\ref{schemes-example-affine-space-zero-doubled}；取
$Y = \mathbf{A}^1_k$ 为仿射直线，并取 $X \to Y$ 为显然的映射。
回忆在 $0 \in Y$ 上方有两个点 $0_1$ 和 $0_2$ 位于 $X$ 中。
因此 $(X/Y)_{fin}$ 在 $0$ 上方有四个点，即
$\emptyset, \{0_1\}, \{0_2\}, \{0_1, 0_2\}$。
在这四个点中，只有三个能提升为 $U \times_Y X$ 的开子概形，
该开子概形在 $U$ 上有限，其中 $U \to Y$ 是 \'etale 的；这三个点是
$\emptyset, \{0_1\}, \{0_2\}$。
这说明，若 $(X/Y)_{fin}$ 可由代数空间表示，则它在 $Y$ 上并非 \'etale。
类似论证说明 $(X/Y)_{fin}$ 确实不是代数空间。略去细节。
\end{remark}

\begin{remark}
\label{spaces-more-groupoids-remark-not-scheme}
设 $Y = \mathbf{A}^1_{\mathbf{R}}$ 为实数域上的仿射直线，
并设 $X = \Spec(\mathbf{C})$ 映到 $\mathbf{R}$-有理点 $0$；
该点位于 $Y$ 中。在此情形下，态射 $f : X \to Y$ 是有限的，
但 $(X/Y)_{fin}$ 不是概形。事实上，可以证明此时的代数空间
$(X/Y)_{fin}$ 同构于《代数空间》例
\ref{spaces-example-non-representable-descent} 中与扩张
$\mathbf{R} \subset \mathbf{C}$ 相联系的代数空间。
因此，为表示层 $(X/Y)_{fin}$，确实有必要离开概形范畴；
即便 $f$ 是有限态射，情况也是如此。
\end{remark}

\begin{lemma}
\label{spaces-more-groupoids-lemma-finite-separated-flat-locally-finite-presentation}
设 $S$ 为概形。
设 $f : X \to Y$ 为 $S$ 上分离、平坦且局部有限表示的代数空间态射。
在此情形下：
\begin{enumerate}
\item $(X/Y)_{fin} \to Y$ 是分离的、可表的且 \'etale 的；
\item 若 $Y$ 是概形，则 $(X/Y)_{fin}$（可由）概形（表示）。
\end{enumerate}
\end{lemma}

\begin{proof}
特别地，$f$ 分离且局部有限型（见《代数空间的态射》引理
\ref{spaces-morphisms-lemma-finite-presentation-finite-type}），
所以由命题 \ref{spaces-more-groupoids-proposition-finite-algebraic-space}，
$(X/Y)_{fin}$ 是代数空间。
为了证明 $(X/Y)_{fin} \to Y$ 分离，必须证明如下事实：
给定概形 $T$ 以及两个偶对 $(a, Z_1)$ 和 $(a, Z_2)$；
它们在 $T$ 上、首分量相同，并满足 \ref{spaces-more-groupoids-equation-finite-conditions}。
存在闭子概形 $V \subset T$，具有如下性质：对任意概形态射
$h : T' \to T$，有
$$
h \text{ 分解经过 } V \Leftrightarrow
\Big(T' \times_T Z_1 = T' \times_T Z_2
\text{，作为子空间位于 }T' \times_Y X\Big)
$$
在引理 \ref{spaces-more-groupoids-lemma-finite-diagonal} 的证明中，已经看到
$V = T' \setminus E$ 是 $T'$ 的开子概形，其闭补为
% P09-ERR-0965
$$
E =
\text{pr}_0|_{Z_1}\left(Z_1 \setminus Z_1 \cap Z_2)\right)
\cup
\text{pr}_0|_{Z_2}\left(Z_2 \setminus Z_1 \cap Z_2)\right).
$$
所以一切都归结为证明 $E$ 也是开的。由引理
\ref{spaces-more-groupoids-lemma-finite-separated} 可见，$Z_1$ 和 $Z_2$ 在
$T' \times_Y X$ 中闭。因此 $Z_1 \setminus Z_1 \cap Z_2$ 在 $Z_1$ 中开。
% P09-ERR-0966
由于 $f$ 平坦且局部有限表示，$\text{pr}_0|_{Z_1}$ 也如此。
这是因为 $Z_1$ 是基变换 $T' \times_Y X$ 的开子空间，并且可应用
《代数空间的态射》引理
\ref{spaces-morphisms-lemma-base-change-finite-presentation} 和引理
\ref{spaces-morphisms-lemma-base-change-flat}。
因此 $\text{pr}_0|_{Z_1}$ 是开映射，见《代数空间的态射》引理
\ref{spaces-morphisms-lemma-fppf-open}。
于是 $\text{pr}_0|_{Z_1}\left(Z_1 \setminus Z_1 \cap Z_2)\right)$ 是开的，
从而 $E$ 如所需是开的。

\medskip\noindent
由命题 \ref{spaces-more-groupoids-proposition-finite-algebraic-space}，已经知道
$(X/Y)_{fin} \to Y$ 是 \'etale 的。
因此现在知道它局部拟有限（见《代数空间的态射》引理
\ref{spaces-morphisms-lemma-etale-locally-quasi-finite}）且分离，
从而由《代数空间的态射》引理
\ref{spaces-morphisms-lemma-locally-quasi-finite-separated-representable}
可表。最后一个断言是清楚的（也可以使用《代数空间的态射》命题
\ref{spaces-morphisms-proposition-locally-quasi-finite-separated-over-scheme}）。
\end{proof}

\noindent
变体：设 $S$ 为概形。
设 $f : X \to Y$ 为 $S$ 上的代数空间态射。
设 $\sigma : Y \to X$ 为 $f$ 的截面。
对于代数空间或概形 $T$，设其在 $S$ 上，并考虑偶对 $(a, Z)$，其中
\begin{equation}
\label{spaces-more-groupoids-equation-finite-conditions-variant}
\begin{matrix}
a : T \to Y\text{ 是相对于下列基底的态射：}S, \\
Z \subset T \times_Y X\text{ 是开子空间}, \\
\text{使得 }\text{pr}_0|_Z : Z \to T\text{ 是有限的，并且} \\
(1_T, \sigma \circ a) : T \to T \times_Y X\text{ 分解经过 }Z.
\end{matrix}
\end{equation}
用 $(X/Y, \sigma)_{fin}$ 表示 $(X/Y)_{fin}$ 中参数化这些偶对的子函子。

\begin{lemma}
\label{spaces-more-groupoids-lemma-finite-plus-section}
设 $S$ 为概形。
设 $f : X \to Y$ 为 $S$ 上的代数空间态射。
设 $\sigma : Y \to X$ 为 $f$ 的截面。考虑上面定义的函子变换
$$
t : (X/Y, \sigma)_{fin} \longrightarrow (X/Y)_{fin}.
$$
则：
\begin{enumerate}
\item $t$ 可由开浸入表示；
\item 若 $f$ 分离，则 $t$ 可由既开又闭的浸入表示；
\item 若 $(X/Y)_{fin}$ 是代数空间，则 $(X/Y, \sigma)_{fin}$ 是代数空间，
并且是 $(X/Y)_{fin}$ 的开子空间；
\item 若 $(X/Y)_{fin}$ 是概形，则 $(X/Y, \sigma)_{fin}$ 是其开子概形。
\end{enumerate}
\end{lemma}

\begin{proof}
略。提示：给定偶对 $(a, Z)$；它在 $T$ 上，并且如式
(\ref{spaces-more-groupoids-equation-finite-conditions}) 所示。$Z$ 在
$(1_T, \sigma \circ a) : T \to T \times_Y X$ 下的逆像
就是所求的 $T$ 的开子概形。
\end{proof}





\section{箭头的有限集合}
\label{spaces-more-groupoids-section-finite-set-arrows}

\noindent
设 $\mathcal{C}$ 为群胚，见《范畴》定义
\ref{categories-definition-groupoid}。
如《群胚》第 \ref{groupoids-section-groupoids} 节所述，
这对应于七元组 $(\text{Ob}, \text{Arrows}, s, t, c, e, i)$。

\medskip\noindent
利用这些数据，可以如下构造另一个群胚 $\mathcal{C}_{fin}$：
\begin{enumerate}
\item $\mathcal{C}_{fin}$ 的一个对象由有限子集
$Z \subset \text{Arrows}$ 构成，并满足如下性质：
\begin{enumerate}
\item $s(Z) = \{u\}$ 是单点集；
\item $e(u) \in Z$。
\end{enumerate}
\item $\mathcal{C}_{fin}$ 的一个态射由偶对 $(Z, z)$ 构成，其中
$Z$ 是 $\mathcal{C}_{fin}$ 的对象，且 $z \in Z$。
\item $(Z, z)$ 的源为 $Z$。
\item $(Z, z)$ 的靶为
$t(Z, z) = \{z' \circ z^{-1}; z' \in Z\}$。
\item 给定 $(Z_1, z_1)$、$(Z_2, z_2)$，若
$s(Z_1, z_1) = t(Z_2, z_2)$，则复合
$(Z_1, z_1) \circ (Z_2, z_2)$ 为 $(Z_2, z_1 \circ z_2)$。
\end{enumerate}
略去关于这确实定义了一个群胚的验证。
从图像上看，$\mathcal{C}_{fin}$ 的一个对象可视为图
$$
\xymatrix{
& \bullet \\
\bullet \ar@(ul, dl)[]_e \ar[ru] \ar[r] \ar[rd] & \bullet \\
& \bullet
}
$$
为了构造 $\mathcal{C}_{fin}$ 的一个态射，选取其中一条箭头，
并把其余箭头同它的逆作前复合。例如，若选取中间的水平箭头，
则靶为下图：
$$
\xymatrix{
& \bullet \\
\bullet & \bullet \ar[l] \ar[u] \ar@(dr, ur)[]_e \ar[d] \\
& \bullet
}
$$
注意，$s(Z, z)$ 和 $t(Z, z)$ 的基数相等。
所以 $\mathcal{C}_{fin}$ 确实是群胚的可数不交并。




\section{群胚的有限部分}
\label{spaces-more-groupoids-section-finite-part-groupoid}

\noindent
本节将使用第 \ref{spaces-more-groupoids-section-finite-set-arrows} 节所述的想法，
取代数空间中一个群胚的有限部分。

\medskip\noindent
设 $S$ 为概形。
设 $B$ 为 $S$ 上的代数空间。
设 $(U, R, s, t, c, e, i)$ 为 $B$ 上代数空间中的群胚。
假设：态射 $s, t$ 分离且局部有限型。
本节始终固定这些记号和假设。
% P09-ERR-0967

\medskip\noindent
将 $R_s$ 记为代数空间 $R$；这里把它看成 $U$ 上的代数空间，
其结构态射为 $s$。
令 $U' = (R_s/U, e)_{fin}$。由于 $s$ 分离且局部有限型，
由命题 \ref{spaces-more-groupoids-proposition-finite-algebraic-space} 和引理
\ref{spaces-more-groupoids-lemma-finite-plus-section}，可见 $U'$ 是一个代数空间，
并带有 \'etale 态射 $g : U' \to U$。此外，由引理
\ref{spaces-more-groupoids-lemma-finite-sheaf}，存在万有开子空间
$Z_{univ} \subset R \times_{s, U, g} U'$，它在 $U'$ 上有限，
并使得 $(1_{U'}, e \circ g) : U' \to R \times_{s, U, g} U'$
分解经过 $Z_{univ}$。再由引理
\ref{spaces-more-groupoids-lemma-finite-separated}，开子空间 $Z_{univ}$ 在
$R \times_{s, U', g} U$ 中也是闭的。
% P09-ERR-0968
目前的图示如下：
$$
\xymatrix{
Z_{univ} \ar[d] \ar[rd] & \\
R \times_{s, U, g} U' \ar[d] \ar[r] & U' \ar[d]^g \\
R \ar[r]^s & U
}
$$
设 $T$ 为 $B$ 上的概形。可见，一个 $T$-值的 $Z_{univ}$ 点
可视为三元组 $(u, Z, z)$，其中：
\begin{enumerate}
\item $u : T \to U$ 是一个 $T$-值点，取值于 $U$；
\item $Z \subset R \times_{s, U, u} T$ 是在 $T$ 上有限的既开又闭子空间，
并且 $(e \circ u, 1_T)$ 分解经过它；
\item $z : T \to R$ 是一个 $T$-值点，取值于 $R$，满足 $s \circ z = u$，
并且 $(z, 1_T)$ 分解经过 $Z$。
\end{enumerate}
说到这里，由第 \ref{spaces-more-groupoids-section-finite-set-arrows} 节的讨论，
在观念上已经清楚，可以把 $(Z_{univ}, U')$ 变成
$B$ 上代数空间中的群胚。为确保无误，我们将用函子观点逐一定义态射
$s', t', c', e', i'$。
% P09-ERR-0969
（请先阅读并理解第 \ref{spaces-more-groupoids-section-finite-set-arrows} 节中的简单构造，
再阅读这里。）

\medskip\noindent
态射 $s' : Z_{univ} \to U'$ 对应于规则
$$
s' : (u, Z, z) \mapsto (u, Z).
$$
态射 $t' : Z_{univ} \to U'$ 由规则
$$
t' : (u, Z, z) \mapsto (t \circ z, c(Z, i \circ z)).
$$
给出。项 $c(Z, i \circ z)$ 有意义，因为映射
$c(-, i \circ z) : R \times_{s, U, u} T \to
R \times_{s, U, t \circ z} T$
是同构，其逆为 $c(-, z)$。
态射 $e' : U' \to Z_{univ}$ 由规则
$$
e' : (u, Z) \mapsto (u, Z, (e \circ u, 1_T)).
$$
给出。注意，由 $(e \circ u, 1_T)$ 分解经过 $Z$ 的要求，
这一定义有意义。
态射 $i' : Z_{univ} \to Z_{univ}$ 由规则
$$
i' : (u, Z, z) \mapsto (t \circ z, c(Z, i \circ z), i \circ z).
$$
给出。最后，复合由规则
$$
c' : ((u_1, Z_1, z_1), (u_2, Z_2, z_2)) \mapsto
(u_2, Z_2, z_1 \circ z_2).
$$
定义。略去关于 $(U', Z_{univ}, s', t', c', e', i')$
满足代数空间中群胚公理的验证。

\medskip\noindent
最后一项信息是，存在群胚的典范态射
$$
(U', Z_{univ}, s', t', c', e', i')
\longrightarrow
(U, R, s, t, c, e, i)
$$
事实上，态射 $U' \to U$ 是态射 $g : U' \to U$，
它由规则 $(u, Z) \mapsto u$ 定义。态射 $Z_{univ} \to R$
由规则 $(u, Z, z) \mapsto z$ 定义。这就完成了构造。
现将所得结果概括如下。

\begin{lemma}
\label{spaces-more-groupoids-lemma-finite-part-groupoid}
设 $S$ 为概形。
设 $B$ 为 $S$ 上的代数空间。
设 $(U, R, s, t, c, e, i)$ 为 $B$ 上代数空间中的群胚。
假设态射 $s, t$ 分离且局部有限型。存在典范态射
$$
(U', Z_{univ}, s', t', c', e', i')
\longrightarrow
(U, R, s, t, c, e, i)
$$
它是 $B$ 上代数空间中群胚的态射，其中：
\begin{enumerate}
\item $g : U' \to U$ 与 $(R_s/U, e)_{fin} \to U$ 等同；
\item $Z_{univ} \subset R \times_{s, U, g} U'$ 是在 $U'$ 上有限、
包含单位元 $e$ 的基变换的万有开（且闭）子空间。
\end{enumerate}
\end{lemma}

\begin{proof}
见上面的讨论。
\end{proof}









\section{群胚概形的 \'etale 局部化}
\label{spaces-more-groupoids-section-etale-localize}

\noindent
本节证明与 \cite[Proposition 4.2]{K-M} 类似的结果。
我们尝试稍微更一般一些，并用态射的有限部分来避免使用 Hilbert 概形。
目标是在 \'etale 局部化之后，在一点上“分裂”代数空间中的群胚。
定义如下（与 \cite[Definition 4.1]{K-M} 非常相似）。

\begin{definition}
\label{spaces-more-groupoids-definition-split-at-point}
设 $S$ 为概形。设 $B$ 为 $S$ 上的代数空间
% P09-ERR-0970
设 $(U, R, s, t, c)$ 为 $B$ 上代数空间中的群胚。
设 $u \in |U|$ 为一点。
\begin{enumerate}
\item 称 $R$ 在 $u$ 上\emph{强分裂}，如果存在开子空间
$P \subset R$，使得：
\begin{enumerate}
\item $(U, P, s|_P, t|_P, c|_{P \times_{s, U, t} P})$
是 $B$ 上代数空间中的群胚；
\item $s|_P$、$t|_P$ 有限；
\item $\{r \in |R| : s(r) = u, t(r) = u\} \subset |P|$，
\end{enumerate}
这样的 $P$ 的选择称为 $R$ 在 $u$ 上的\emph{强分裂结构}。
\item 称 $R$ 在 $u$ 上\emph{分裂}，如果存在开子空间
$P \subset R$，使得：
\begin{enumerate}
\item $(U, P, s|_P, t|_P, c|_{P \times_{s, U, t} P})$
是 $B$ 上代数空间中的群胚；
\item $s|_P$、$t|_P$ 有限；
\item $\{g \in |G| : g\text{ 映到 }u\} \subset |P|$，
其中 $G \to U$ 为稳定子，
\end{enumerate}
这样的 $P$ 的选择称为 $R$ 在 $u$ 上的\emph{分裂结构}。
\item 称 $R$ 在 $u$ 上\emph{拟分裂}，如果存在开子空间
$P \subset R$，使得：
\begin{enumerate}
\item $(U, P, s|_P, t|_P, c|_{P \times_{s, U, t} P})$
是 $B$ 上代数空间中的群胚；
\item $s|_P$、$t|_P$ 有限；
\item $e(u) \in |P|$\footnote{此条件由 (a) 蕴含。}，
\end{enumerate}
这样的 $P$ 的选择称为 $R$ 在 $u$ 上的\emph{拟分裂结构}。
\end{enumerate}
\end{definition}

\noindent
注意，关于 $P$ 的这些条件与式
(\ref{spaces-more-groupoids-equation-finite-conditions}) 中偶对的条件相似。
特别地，若 $s, t$ 分离，则 $P$ 在 $R$ 中也闭
（见引理 \ref{spaces-more-groupoids-lemma-finite-separated}）。

\medskip\noindent
假设从代数空间中的群胚 $(U, R, s, t, c)$ 出发，设它在 $B$ 上；
再取一点 $u \in |U|$。
由于目标是在 \'etale 局部化之后分裂群胚，不妨用仿射概形替换 $U$
（意思是，对任何可能的应用，这样做都无害）。
此外，我们将不得不施加的附加假设会迫使 $R$ 至少在
$\{r \in |R| : s(r) = u, t(r) = u\}$ 或 $e(u)$ 的一个邻域中成为概形。
因此，我们从下面所述的群胚概形开始。
然而，我们的证明方法会把我们带出概形范畴；正因如此，
上面把分裂表述成了代数空间中群胚的情形。
另一方面，我们只知道态射 $s$、$t$ 还平坦且有限表示这一情形的应用；
在此情形下，最终又回到概形范畴中。

\begin{situation}[强分裂]
\label{spaces-more-groupoids-situation-strong-splitting}
设 $S$ 为概形。
设 $(U, R, s, t, c)$ 为 $S$ 上的群胚概形。
设 $u \in U$ 为一点。假设：
\begin{enumerate}
\item $s, t : R \to U$ 分离；
\item $s$、$t$ 局部有限型；
\item 集合 $\{r \in R : s(r) = u, t(r) = u\}$ 有限；
\item $s$ 在 (3) 中集合的每一点处拟有限。
\end{enumerate}
注意，若纤维 $s^{-1}(\{u\})$ 有限，则假设 (3) 和 (4) 均成立，
见《态射》引理 \ref{morphisms-lemma-finite-fibre}。
\end{situation}

\begin{situation}[分裂]
\label{spaces-more-groupoids-situation-splitting}
设 $S$ 为概形。
设 $(U, R, s, t, c)$ 为 $S$ 上的群胚概形。
设 $u \in U$ 为一点。假设：
\begin{enumerate}
\item $s, t : R \to U$ 分离；
\item $s$、$t$ 局部有限型；
\item 集合 $\{g \in G : g\text{ 映到 }u\}$ 有限，其中
$G \to U$ 是稳定子；
\item $s$ 在 (3) 中集合的每一点处拟有限。
\end{enumerate}
\end{situation}

\begin{situation}[拟分裂]
\label{spaces-more-groupoids-situation-quasi-splitting}
设 $S$ 为概形。
设 $(U, R, s, t, c)$ 为 $S$ 上的群胚概形。
设 $u \in U$ 为一点。假设：
\begin{enumerate}
\item $s, t : R \to U$ 分离；
\item $s$、$t$ 局部有限型；
\item $s$ 在 $e(u)$ 处拟有限。
\end{enumerate}
\end{situation}

\noindent
对于代数空间存在性定理的应用，拟分裂情形已经足够。
此外，拟分裂情形将使我们能够证明拟-DM 叠的 \'etale 局部结构定理。
分裂情形将用于证明 Keel--Mori 定理的一个版本。
强分裂情形可给出对角拟紧的拟-DM 代数叠的 \'etale 局部结构定理。

\begin{lemma}[强分裂的存在性]
\label{spaces-more-groupoids-lemma-strong-splitting}
在情形 \ref{spaces-more-groupoids-situation-strong-splitting} 中，存在代数空间 $U'$、
\'etale 态射 $U' \to U$，以及点
$u' : \Spec(\kappa(u)) \to U'$；它位于
$u : \Spec(\kappa(u)) \to U$ 上方，使得限制
$R' = R|_{U'}$，即 $R$ 在 $U'$ 上的限制，在 $u'$ 上强分裂。
\end{lemma}

\begin{proof}
设 $f : (U', Z_{univ}, s', t', c') \to (U, R, s, t, c)$
为引理 \ref{spaces-more-groupoids-lemma-finite-part-groupoid} 中构造的态射。
回忆
$R' = R \times_{(U \times_S U)} (U' \times_S U')$。
因此得到代数空间中群胚的态射
$(f, t', s') : Z_{univ} \to R'$：
$$
(U', Z_{univ}, s', t', c') \to (U', R', s', t', c')
$$
（滥用记号，以相同符号表示两个群胚中的态射。）
现在，由于 $Z_{univ} \subset R \times_{s, U, g} U'$ 是开的，
而 $R' \to R \times_{s, U, g} U'$ 是 \'etale 的
（它是 $U' \to U$ 的基变换），可见
$Z_{univ} \to R'$ 是开浸入。按照构造，态射
$s', t' : Z_{univ} \to U'$ 有限。还须找到点 $u'$，它属于 $U'$。

\medskip\noindent
如引理陈述那样，把 $u$ 看成态射
$\Spec(\kappa(u)) \to U$。令
$F_u = R \times_{s, U} \Spec(\kappa(u))$。
由假设，集合 $\{r \in R : s(r) = u, t(r) = u\}$ 有限，
而 $F_u \to \Spec(\kappa(u))$ 在其每个元素处拟有限。
因此，可以找到既开又闭子概形的分解
$$
F_u = Z_u \amalg Rest
$$
，其中概形 $Z_u$ 在 $\kappa(u)$ 上有限，其支集为
$\{r \in R : s(r) = u, t(r) = u\}$。注意 $e(u) \in Z_u$。
因此，由第 \ref{spaces-more-groupoids-section-finite-part-groupoid} 节中 $U'$ 的构造，
$(u, Z_u)$ 定义了一个 $\Spec(\kappa(u))$-值点 $u'$，
该点属于 $U'$。

\medskip\noindent
还须证明集合
$\{r' \in |R'| : s'(r') = u', t'(r') = u'\}$
包含在 $|Z_{univ}|$ 中。任取此集合中的一点 $r'$，
并用态射 $z' : \Spec(k) \to R'$ 表示它。将 $z : \Spec(k) \to R$
记为 $z'$ 与映射 $R' \to R$ 的合成。
显然，$z$ 定义了集合
$\{r \in R : s(r) = u, t(r) = u\}$ 的一个元素。
此外，合成 $s \circ z, t \circ z : \Spec(k) \to U$
都分解经过 $u$，所以可把 $s \circ z, t \circ z$
看成态射 $\Spec(k) \to \Spec(\kappa(u))$。于是
$$
z' = (z, u' \circ t \circ z, u'\circ s \circ u)
$$
作为到 $R' = R \times_{(U \times_S U)} (U' \times_S U')$ 的态射。
% P09-ERR-0971
考虑三元组
$$
(s \circ z, Z_u \times_{\Spec(\kappa(u)), s \circ z} \Spec(k), z)
$$
其中 $Z_u$ 如上。它定义了一个 $\Spec(k)$-值点，取值于
$Z_{univ}$；其在 $s', t'$ 下的 $U'$ 中之像为 $u'$，
并且它在 $Z_{univ} \to R'$ 下的像就是点 $r'$；
这是由上面所述 $z$ 与 $z'$ 的关系得到的。
这就完成了证明。
\end{proof}

\begin{lemma}[分裂的存在性]
\label{spaces-more-groupoids-lemma-splitting}
在情形 \ref{spaces-more-groupoids-situation-splitting} 中，存在代数空间 $U'$、
\'etale 态射 $U' \to U$，以及点
$u' : \Spec(\kappa(u)) \to U'$；它位于
$u : \Spec(\kappa(u)) \to U$ 上方，使得限制
$R' = R|_{U'}$，即 $R$ 在 $U'$ 上的限制，在 $u'$ 上分裂。
\end{lemma}

\begin{proof}
设 $f : (U', Z_{univ}, s', t', c') \to (U, R, s, t, c)$
为引理 \ref{spaces-more-groupoids-lemma-finite-part-groupoid} 中构造的态射。
回忆 $R' = R \times_{(U \times_S U)} (U' \times_S U')$。
因此得到代数空间中群胚的态射
$(f, t', s') : Z_{univ} \to R'$：
$$
(U', Z_{univ}, s', t', c') \to (U', R', s', t', c')
$$
（滥用记号，以相同符号表示两个群胚中的态射。）
现在，由于 $Z_{univ} \subset R \times_{s, U, g} U'$ 是开的，
而 $R' \to R \times_{s, U, g} U'$ 是 \'etale 的
（它是 $U' \to U$ 的基变换），可见
$Z_{univ} \to R'$ 是开浸入。按照构造，态射
$s', t' : Z_{univ} \to U'$ 有限。还须找到点 $u'$，它属于 $U'$。

\medskip\noindent
如引理陈述那样，把 $u$ 看成态射
$\Spec(\kappa(u)) \to U$。令
$F_u = R \times_{s, U} \Spec(\kappa(u))$。
设 $G_u \subset F_u$ 为稳定子 $G \to U$ 在 $u$ 上的概形论纤维。
由假设，$G_u$ 有限，并且 $F_u \to \Spec(\kappa(u))$
在 $G_u$ 的每一点处拟有限。因此，可以找到既开又闭子概形的分解
$$
F_u = Z_u \amalg Rest
$$
，其中概形 $Z_u$ 在 $\kappa(u)$ 上有限，其支集为 $G_u$。
注意 $e(u) \in Z_u$。因此，由第
\ref{spaces-more-groupoids-section-finite-part-groupoid} 节中 $U'$ 的构造，
$(u, Z_u)$ 定义了一个 $\Spec(\kappa(u))$-值点 $u'$，
该点属于 $U'$。

\medskip\noindent
还须证明集合 $\{g' \in |G'| : g'\text{ 映到 }u'\}$
包含在 $|Z_{univ}|$ 中。
% P09-ERR-0972
任取此集合中的一点 $g'$，并用态射
$z' : \Spec(k) \to G'$ 表示它。将
$z : \Spec(k) \to G$ 记为 $z'$ 与映射 $G' \to G$ 的合成。
显然，$z$ 定义了 $G_u$ 的一点。事实上，将相应的映射记作
$\tilde u : \Spec(k) \to u \to U$；这是到 $u$ 或 $U$ 的映射。
考虑三元组
$$
(\tilde u, Z_u \times_{u, \tilde u} \Spec(k), z)
$$
，其中 $Z_u$ 如上。它定义了一个 $\Spec(k)$-值点，取值于
$Z_{univ}$；其在 $s', t'$ 下的 $U'$ 中之像为 $u'$，
并且它在 $Z_{univ} \to R'$ 下的像就是点 $z'$
（因为它在 $R$ 中的像为 $z$）。
这就完成了证明。
\end{proof}

\begin{lemma}[拟分裂的存在性]
\label{spaces-more-groupoids-lemma-quasi-splitting}
在情形 \ref{spaces-more-groupoids-situation-quasi-splitting} 中，存在代数空间 $U'$、
\'etale 态射 $U' \to U$，以及点
$u' : \Spec(\kappa(u)) \to U'$；它位于
$u : \Spec(\kappa(u)) \to U$ 上方，使得限制
$R' = R|_{U'}$，即 $R$ 在 $U'$ 上的限制，在 $u'$ 上拟分裂。
\end{lemma}

\begin{proof}
设 $f : (U', Z_{univ}, s', t', c') \to (U, R, s, t, c)$
为引理 \ref{spaces-more-groupoids-lemma-finite-part-groupoid} 中构造的态射。
回忆 $R' = R \times_{(U \times_S U)} (U' \times_S U')$。
因此得到代数空间中群胚的态射
$(f, t', s') : Z_{univ} \to R'$：
$$
(U', Z_{univ}, s', t', c') \to (U', R', s', t', c')
$$
（滥用记号，以相同符号表示两个群胚中的态射。）
现在，由于 $Z_{univ} \subset R \times_{s, U, g} U'$ 是开的，
而 $R' \to R \times_{s, U, g} U'$ 是 \'etale 的
（它是 $U' \to U$ 的基变换），可见
$Z_{univ} \to R'$ 是开浸入。按照构造，态射
$s', t' : Z_{univ} \to U'$ 有限。还须找到点 $u'$，它属于 $U'$。

\medskip\noindent
如引理陈述那样，把 $u$ 看成态射
$\Spec(\kappa(u)) \to U$。令
$F_u = R \times_{s, U} \Spec(\kappa(u))$。
由假设，态射 $F_u \to \Spec(\kappa(u))$ 在 $e(u)$ 处拟有限。
因此，可以找到既开又闭子概形的分解
$$
F_u = Z_u \amalg Rest
$$
，其中概形 $Z_u$ 在 $\kappa(u)$ 上有限，其支集为 $e(u)$。
因此，由第 \ref{spaces-more-groupoids-section-finite-part-groupoid} 节中 $U'$ 的构造，
$(u, Z_u)$ 定义了一个 $\Spec(\kappa(u))$-值点 $u'$，
该点属于 $U'$。为完成证明，须证明
$e'(u') \in Z_{univ}$，这是清楚的。
\end{proof}

\noindent
最后，加入附加假设后可得到概形。

\begin{lemma}
\label{spaces-more-groupoids-lemma-strong-splitting-scheme}
在情形 \ref{spaces-more-groupoids-situation-strong-splitting} 中，再假设
$s, t$ 平坦且局部有限表示。则存在概形 $U'$、分离的 \'etale 态射
$U' \to U$，以及点 $u' \in U'$；它位于 $u$ 上方，并满足
$\kappa(u) = \kappa(u')$，使得限制
$R' = R|_{U'}$，即 $R$ 在 $U'$ 上的限制，在 $u'$ 上强分裂。
\end{lemma}

\begin{proof}
这由引理 \ref{spaces-more-groupoids-lemma-strong-splitting} 的证明中 $U'$ 的构造得到，
因为在此情形下，$U' = (R_s/U, e)_{fin}$ 是在 $U$ 上分离的概形；
这里使用引理
\ref{spaces-more-groupoids-lemma-finite-separated-flat-locally-finite-presentation} 和
\ref{spaces-more-groupoids-lemma-finite-plus-section}。
\end{proof}

\begin{lemma}
\label{spaces-more-groupoids-lemma-splitting-scheme}
在情形 \ref{spaces-more-groupoids-situation-splitting} 中，再假设
$s, t$ 平坦且局部有限表示。则存在概形 $U'$、分离的 \'etale 态射
$U' \to U$，以及点 $u' \in U'$；它位于 $u$ 上方，并满足
$\kappa(u) = \kappa(u')$，使得限制
$R' = R|_{U'}$，即 $R$ 在 $U'$ 上的限制，在 $u'$ 上分裂。
\end{lemma}

\begin{proof}
这由引理 \ref{spaces-more-groupoids-lemma-splitting} 的证明中 $U'$ 的构造得到，
因为在此情形下，$U' = (R_s/U, e)_{fin}$ 是在 $U$ 上分离的概形；
这里使用引理
\ref{spaces-more-groupoids-lemma-finite-separated-flat-locally-finite-presentation} 和
\ref{spaces-more-groupoids-lemma-finite-plus-section}。
\end{proof}

\begin{lemma}
\label{spaces-more-groupoids-lemma-quasi-splitting-scheme}
在情形 \ref{spaces-more-groupoids-situation-quasi-splitting} 中，再假设
$s, t$ 平坦且局部有限表示。则存在概形 $U'$、分离的 \'etale 态射
$U' \to U$，以及点 $u' \in U'$；它位于 $u$ 上方，并满足
$\kappa(u) = \kappa(u')$，使得限制
$R' = R|_{U'}$，即 $R$ 在 $U'$ 上的限制，在 $u'$ 上拟分裂。
\end{lemma}

\begin{proof}
这由引理 \ref{spaces-more-groupoids-lemma-quasi-splitting} 的证明中 $U'$ 的构造得到，
因为在此情形下，$U' = (R_s/U, e)_{fin}$ 是在 $U$ 上分离的概形；
这里使用引理
\ref{spaces-more-groupoids-lemma-finite-separated-flat-locally-finite-presentation} 和
\ref{spaces-more-groupoids-lemma-finite-plus-section}。
\end{proof}

\noindent
事实上，可以应用前面关于有限局部自由群胚的结果来得到仿射概形。

\begin{lemma}
\label{spaces-more-groupoids-lemma-strong-splitting-affine-scheme}
在情形 \ref{spaces-more-groupoids-situation-strong-splitting} 中，再假设
$s, t$ 平坦且局部有限表示，并且 $U$ 仿射。
则存在仿射概形 $U'$、\'etale 态射 $U' \to U$，
以及点 $u' \in U'$；它位于 $u$ 上方，并满足
$\kappa(u) = \kappa(u')$，使得限制
$R' = R|_{U'}$，即 $R$ 在 $U'$ 上的限制，在 $u'$ 上强分裂。
\end{lemma}

\begin{proof}
设 $U' \to U$ 和 $u' \in U'$ 为引理
\ref{spaces-more-groupoids-lemma-strong-splitting-scheme} 中找到的概形之间的分离
\'etale 态射及其点。设 $P \subset R'$ 为 $R'$ 在 $u'$ 上的强分裂结构。
由《群胚进阶》引理
\ref{more-groupoids-lemma-restrict-preserves-type}，态射
$s', t' : R' \to U'$ 平坦且局部有限表示。由假设，它们有限。
% P09-ERR-0973
因此 $s', t'$ 有限局部自由，见《态射》引理
\ref{morphisms-lemma-finite-flat}。
特别地，$t(s^{-1}(u'))$ 是点的有限集合
$\{u'_1, u'_2, \ldots, u'_n\}$，这些点位于 $U'$ 中。
选取拟紧开子集 $W \subset U'$，使它包含每个 $u'_i$。
由于 $U$ 仿射，态射 $W \to U$ 拟紧
（见《概形》引理 \ref{schemes-lemma-quasi-compact-affine}）。
态射 $W \to U$ 还局部拟有限（见《态射》引理
\ref{morphisms-lemma-etale-locally-quasi-finite}）且分离。
因此，由《态射进阶》引理
\ref{more-morphisms-lemma-quasi-finite-separated-quasi-affine}
（Zariski 主定理的一个版本），可知 $W$ 拟仿射。
由《性质》引理 \ref{properties-lemma-ample-finite-set-in-affine}，
可见 $\{u'_1, \ldots, u'_n\}$ 包含在 $U'$ 的一个仿射开子集中。
于是可以应用《群胚》引理
\ref{groupoids-lemma-find-invariant-affine}，推出存在仿射的
$P$-不变开子集 $U'' \subset U'$，它包含 $u'$。

\medskip\noindent
为完成证明，将 $R'' = R|_{U''}$ 记为 $R$ 在 $U''$ 上的限制。
这也等于 $R'$ 在 $U''$ 上的限制。由于
$P \subset R'$ 是既开又闭子概形，$P|_{U''} \subset R''$ 也如此。
按照构造，开子概形 $U'' \subset U'$ 是 $P$-不变的，这意味着
$$
P|_{U''} = (s'|_P)^{-1}(U'') = (t'|_P)^{-1}(U'')
$$
（见《群胚》第 \ref{groupoids-section-invariant} 节中的讨论），
所以 $s''$ 和 $t''$ 在 $P|_{U''}$ 上的限制仍然有限。
子群胚概形 $P|_{U''}$ 仍然是 $R''$ 在 $u''$ 上的强分裂结构；
% P09-ERR-0974
上面已经验证了 (a)、(b)，而 (c) 成立是因为显然有
$$
\{r' \in R': t'(r') = u', s'(r') = u'\} =
\{r'' \in R'': t''(r'') = u', s''(r'') = u'\}
$$
引理得证。
\end{proof}

\begin{lemma}
\label{spaces-more-groupoids-lemma-splitting-affine-scheme}
在情形 \ref{spaces-more-groupoids-situation-splitting} 中，再假设
$s, t$ 平坦且局部有限表示，并且 $U$ 仿射。
则存在仿射概形 $U'$、\'etale 态射 $U' \to U$，
以及点 $u' \in U'$；它位于 $u$ 上方，并满足
$\kappa(u) = \kappa(u')$，使得限制
$R' = R|_{U'}$，即 $R$ 在 $U'$ 上的限制，在 $u'$ 上分裂。
\end{lemma}

\begin{proof}
此引理的证明逐字等同于引理
\ref{spaces-more-groupoids-lemma-strong-splitting-affine-scheme} 的证明，
只须将“强分裂”替换为“分裂”（共 2 处），并把对引理
\ref{spaces-more-groupoids-lemma-strong-splitting-scheme} 的引用替换为对引理
\ref{spaces-more-groupoids-lemma-splitting-scheme} 的引用。
\end{proof}

\begin{lemma}
\label{spaces-more-groupoids-lemma-quasi-splitting-affine-scheme}
在情形 \ref{spaces-more-groupoids-situation-quasi-splitting} 中，再假设
$s, t$ 平坦且局部有限表示，并且 $U$ 仿射。
则存在仿射概形 $U'$、\'etale 态射 $U' \to U$，
以及点 $u' \in U'$；它位于 $u$ 上方，并满足
$\kappa(u) = \kappa(u')$，使得限制
$R' = R|_{U'}$，即 $R$ 在 $U'$ 上的限制，在 $u'$ 上拟分裂。
\end{lemma}

\begin{proof}
此引理的证明逐字等同于引理
\ref{spaces-more-groupoids-lemma-strong-splitting-affine-scheme} 的证明，
只须将“强分裂”替换为“拟分裂”（共 2 处），并把对引理
\ref{spaces-more-groupoids-lemma-strong-splitting-scheme} 的引用替换为对引理
\ref{spaces-more-groupoids-lemma-quasi-splitting-scheme} 的引用。
\end{proof}
% Generated mechanically from the frozen bootstrap.tex target.
% Do not edit this generated file; edit the bound locale source instead.

% OK, start here.
%

\setcounter{chapter}{79}
\chapter{自举}


\phantomsection
\label{bootstrap-section-phantom}





\section{引言}
\label{bootstrap-section-introduction}

\noindent
本章利用前面各章的材料，给出概形范畴上的集合值预层成为代数空间的判据。
其中一部分内容源自 Artin 的工作，见 \cite{ArtinI}、\cite{ArtinII}、
\cite{Artin-Theorem-Representability},
\cite{Artin-Construction-Techniques},
\cite{Artin-Algebraic-Spaces},
\cite{Artin-Algebraic-Approximation},
\cite{Artin-Implicit-Function},
以及 \cite{ArtinVersal}。
不过，我们的方法将尽可能采用与 Keel 和 Mori 的论文中相近的论证，见
\cite{K-M}。

\section{约定}
\label{bootstrap-section-conventions}

\noindent
始终假定所有概形都包含在一个大 fppf 位点 $\Sch_{fppf}$ 中。
并且所考虑的每个环 $A$ 都满足：$\Spec(A)$（同构于）该大位点中的一个对象。

\medskip\noindent
设 $S$ 为概形，$X$ 为 $S$ 上的代数空间。
在本章及以下各章中，我们把 $X$ 与自身的积（在 $S$ 上代数空间的范畴中）
记作 $X \times_S X$，而不记作 $X \times X$。




\section{可由代数空间表示的态射}
\label{bootstrap-section-morphism-representable-by-spaces}

\noindent
这里定义一个预层相对于另一个预层可由代数空间表示的概念，并证明这一概念的若干性质。

\begin{definition}
\label{bootstrap-definition-morphism-representable-by-spaces}
设 $S$ 为包含在 $\Sch_{fppf}$ 中的概形，
$F$、$G$ 为 $\Sch_{fppf}/S$ 上的预层。
若对每个 $U \in \Ob((\Sch/S)_{fppf})$ 以及任意
$\xi : U \to G$，纤维积 $U \times_{\xi, G} F$ 都是代数空间，
则称态射 $a : F \to G$ {\it 可由代数空间表示}。
\end{definition}

\noindent
先作一项基本检验。

\begin{lemma}
\label{bootstrap-lemma-morphism-spaces-is-representable-by-spaces}
设 $S$ 为概形，$f : X \to Y$ 为 $S$ 上代数空间的态射。
则 $f$ 可由代数空间表示。
\end{lemma}

\begin{proof}
这是一个形式推论。它只用到 $S$ 上代数空间的范畴具有纤维积这一事实，见
《空间》章，引理 \ref{spaces-lemma-fibre-product-spaces}。
\end{proof}

\begin{lemma}
\label{bootstrap-lemma-base-change-transformation}
\begin{slogan}
预层之间可由代数空间表示的态射经基变换后仍可由代数空间表示。
\end{slogan}
设 $S$ 为概形。设
$$
\xymatrix{
G' \times_G F \ar[r] \ar[d]^{a'} & F \ar[d]^a \\
G' \ar[r] & G
}
$$
是 $(\Sch/S)_{fppf}$ 上预层的纤维方块。
若 $a$ 可由代数空间表示，则 $a'$ 也可由代数空间表示。
\end{lemma}

\begin{proof}
略。提示：这是一个形式推论。
\end{proof}

\begin{lemma}
\label{bootstrap-lemma-representable-by-spaces-transformation-to-sheaf}
设 $S$ 为包含在 $\Sch_{fppf}$ 中的概形，并设
$F, G : (\Sch/S)_{fppf}^{opp} \to \textit{Sets}$。
设 $a : F \to G$ 可由代数空间表示。
若 $G$ 是层，则 $F$ 也是层。
\end{lemma}

\begin{proof}
（与《空间》章，引理
\ref{spaces-lemma-representable-transformation-to-sheaf} 的证明相同。）
设 $\{\varphi_i : T_i \to T\}$ 为位点 $(\Sch/S)_{fppf}$ 的一个覆盖，
并设满足层条件的 $s_i \in F(T_i)$。
于是 $\sigma_i = a(s_i) \in G(T_i)$ 也满足层条件。
因此存在唯一的 $\sigma \in G(T)$，使得
$\sigma_i = \sigma|_{T_i}$。由假设，
$F' = h_T \times_{\sigma, G, a} F$ 是层。
又有 $(\varphi_i, s_i) \in F'(T_i)$ 也满足层条件，
因而来自唯一的 $(\text{id}_T, s) \in F'(T)$。
显然，$s$ 正是所求的 $F$ 的截面。
\end{proof}

\begin{lemma}
\label{bootstrap-lemma-representable-by-spaces-transformation-diagonal}
设 $S$ 为包含在 $\Sch_{fppf}$ 中的概形，并设
$F, G : (\Sch/S)_{fppf}^{opp} \to \textit{Sets}$。
设 $a : F \to G$ 可由代数空间表示。
则 $\Delta_{F/G} : F \to F \times_G F$ 可由代数空间表示。
\end{lemma}

\begin{proof}
（与《空间》章，引理
\ref{spaces-lemma-representable-transformation-diagonal} 的证明相同。）
设 $U$ 为概形，并设
$\xi = (\xi_1, \xi_2) \in (F \times_G F)(U)$。
令 $\xi' = a(\xi_1) = a(\xi_2) \in G(U)$。
由假设，存在代数空间 $V$ 及态射 $V \to U$，表示纤维积
$U \times_{\xi', G} F$。
特别地，元素 $\xi_1,\xi_2$ 给出 $U$ 上的态射
$f_1, f_2 : U \to V$。由于 $V$ 表示纤维积
$U \times_{\xi', G} F$，且
$\xi' = a \circ \xi_1 = a \circ \xi_2$
，可知对任意态射 $g : U' \to U$，
$$
g^*\xi_1 = g^*\xi_2
\Leftrightarrow
f_1 \circ g = f_2 \circ g.
$$
换言之，$U \times_{\xi, F \times_G F} F$ 由
$V \times_{\Delta, V \times V, (f_1, f_2)} U$ 表示，
而后者是代数空间。
\end{proof}

\noindent
下面引理 \ref{bootstrap-lemma-representable-by-spaces-over-space} 的证明
其实略有微妙之处。具体而言，我们不能采用《空间》章，引理
\ref{spaces-lemma-representable-over-space} 的证明论证，
因为此时尚不知道两个可由代数空间表示的态射之复合仍可由代数空间表示。
事实上，我们将用这个引理来证明该命题。

\begin{lemma}
\label{bootstrap-lemma-representable-by-spaces-over-space}
设 $S$ 为包含在 $\Sch_{fppf}$ 中的概形，并设
$F, G : (\Sch/S)_{fppf}^{opp} \to \textit{Sets}$。
设 $a : F \to G$ 可由代数空间表示。
若 $G$ 是代数空间，则 $F$ 也是代数空间。
\end{lemma}

\begin{proof}
由引理 \ref{bootstrap-lemma-representable-by-spaces-transformation-to-sheaf}
可知 $F$ 是层。

\medskip\noindent
设 $U$ 为概形，且 $U \to G$ 为满的平展态射。
此时 $U \times_G F$ 是代数空间。再设 $W$ 为概形，且
$W \to U \times_G F$ 为满的平展态射。

\medskip\noindent
首先断言 $W \to F$ 是可表的。
为此，设 $X$ 为概形，并设 $X \to F$ 为态射。则
$$
W \times_F X = W \times_{U \times_G F} U \times_G F \times_F X
= W \times_{U \times_G F} (U \times_G X)
$$
由于 $U \times_G F$ 和 $G$ 都是代数空间，可知这是一个概形。

\medskip\noindent
其次断言 $W \to F$ 是满的且为平展态射（既已知道它可表，这一说法便有意义）。
这由上式得出，因为 $W \to U \times_G F$ 与 $U \to G$
均为满的平展态射，故
$W \times_{U \times_G F} (U \times_G X) \to U \times_G X$ 与
$U \times_G X \to X$ 均为满的平展态射，而满的平展态射之复合仍是满的平展态射。

\medskip\noindent
令 $R = W \times_F W$。由上可知 $R$ 是概形，且投影
$t, s : R \to W$ 均为平展态射。
显然，$R$ 是等价关系，并且 $W \to F$ 是层的满射。
因此 $R$ 是平展等价关系且 $F=W/R$。
由《空间》章，定理 \ref{spaces-theorem-presentation}，
$F$ 是代数空间。
\end{proof}

\begin{lemma}
\label{bootstrap-lemma-representable-by-spaces}
设 $S$ 为概形，$a : F \to G$ 为 $(\Sch/S)_{fppf}$ 上预层的态射。
假设 $a : F \to G$ 可由代数空间表示。
若 $X$ 是 $S$ 上的代数空间，且 $X \to G$ 是预层的态射，
则 $X \times_G F$ 是代数空间。
\end{lemma}

\begin{proof}
由引理 \ref{bootstrap-lemma-base-change-transformation}，态射
$X \times_G F \to X$ 可由代数空间表示。
因而由引理 \ref{bootstrap-lemma-representable-by-spaces-over-space}，
$X \times_G F$ 是代数空间。
\end{proof}

\begin{lemma}
\label{bootstrap-lemma-composition-transformation}
设 $S$ 为概形。设
$$
\xymatrix{
F \ar[r]^a & G \ar[r]^b & H
}
$$
是 $(\Sch/S)_{fppf}$ 上预层的态射。
若 $a$ 与 $b$ 均可由代数空间表示，则 $b \circ a$ 也可由代数空间表示。
\end{lemma}

\begin{proof}
设 $T$ 为 $S$ 上的概形，且 $T \to H$ 为态射。
由假设，$T \times_H G$ 是代数空间。
于是由引理 \ref{bootstrap-lemma-representable-by-spaces} 可知
$T \times_H F = (T \times_H G) \times_G F$ 也是代数空间。
\end{proof}

\begin{lemma}
\label{bootstrap-lemma-product-transformations}
设 $S$ 为概形。令
$F_i, G_i : (\Sch/S)_{fppf}^{opp} \to \textit{Sets}$，$i=1,2$。
设 $a_i : F_i \to G_i$（$i=1,2$）可由代数空间表示。则
$$
a_1 \times a_2 : F_1 \times F_2 \longrightarrow G_1 \times G_2
$$
可由代数空间表示。
\end{lemma}

\begin{proof}
把 $a_1 \times a_2$ 写成复合
$F_1 \times F_2 \to G_1 \times F_2 \to G_1 \times G_2$。
第一支箭头是 $a_1$ 沿态射 $G_1 \times F_2 \to G_1$ 的基变换，
第二支箭头是 $a_2$ 沿态射 $G_1 \times G_2 \to G_2$ 的基变换。
因此，本引理是引理 \ref{bootstrap-lemma-composition-transformation}
与 \ref{bootstrap-lemma-base-change-transformation} 的形式推论。
\end{proof}

\begin{lemma}
\label{bootstrap-lemma-representable-by-spaces-permanence}
设 $S$ 为概形。设 $a : F \to G$ 与 $b : G \to H$ 是函子
$(\Sch/S)_{fppf}^{opp} \to \textit{Sets}$ 之间的态射。假设
\begin{enumerate}
\item $\Delta : G \to G \times_H G$ 可由代数空间表示；并且
\item $b \circ a : F \to H$ 可由代数空间表示。
\end{enumerate}
则 $a$ 可由代数空间表示。
\end{lemma}

\begin{proof}
设 $U$ 为 $S$ 上的概形，且 $\xi \in G(U)$。则
$$
U \times_{\xi, G, a} F =
(U \times_{b(\xi), H, b \circ a} F) \times_{(\xi, a), (G \times_H G), \Delta} G
$$
由引理 \ref{bootstrap-lemma-representable-by-spaces} 即得结论。
\end{proof}

\begin{lemma}
\label{bootstrap-lemma-glueing-sheaves}
设 $S \in \Ob(\Sch_{fppf})$，并设 $F$ 为
$(\Sch/S)_{fppf}$ 上的集合值预层。假设
\begin{enumerate}
\item $F$ 是 $(\Sch/S)_{fppf}$ 上 Zariski 拓扑的层；
\item 存在指标集 $I$ 及子函子 $F_i \subset F$，使得
\begin{enumerate}
\item 每个 $F_i$ 都是 fppf 层；
\item 每个 $F_i \to F$ 都可由代数空间表示；
\item $\coprod F_i \to F$ 经 fppf 层化后成为满射。
\end{enumerate}
\end{enumerate}
则 $F$ 是 fppf 层。
\end{lemma}

\begin{proof}
设 $T \in \Ob((\Sch/S)_{fppf})$ 且 $s \in F(T)$。
由 (2)(c)，存在 fppf 覆盖 $\{T_j \to T\}$，使得对某个
$\alpha(j) \in I$，$s|_{T_j}$ 是 $F_{\alpha(j)}$ 的截面。
令 $W_j \subset T$ 为 $T_j \to T$ 的像；由《态射》章，引理
\ref{morphisms-lemma-fppf-open}，它是开子概形。
由 (2)(b)，
$F_{\alpha(j)} \times_{F, s|_{W_j}} W_j \to W_j$
是代数空间的单态射，且 $T_j$ 经其分解。
由于 $\{T_j \to W_j\}$ 是 fppf 覆盖，可得
$F_{\alpha(j)} \times_{F, s|_{W_j}} W_j = W_j$；换言之，
$s|_{W_j} \in F_{\alpha(j)}(W_j)$。
因此 $\coprod F_i \to F$ 对 Zariski 拓扑是满的。

\medskip\noindent
设 $\{T_j \to T\}$ 是 $(\Sch/S)_{fppf}$ 中的 fppf 覆盖。
设 $s,s' \in F(T)$，且对所有 $j$ 都有
$s|_{T_j}=s'|_{T_j}$。我们要证明 $s=s'$。
由 (1)，$F$ 是 Zariski 层，所以可以在 $T$ 上作 Zariski 局部讨论。
根据上一段的结论，可以假定存在 $i$ 使 $s\in F_i(T)$。
于是 $s'|_{T_j}$ 是 $F_i$ 的截面。由 (2)(b)，
$F_i \times_{F,s'}T\to T$ 是代数空间的单态射，且所有 $T_j$ 均经其分解。
因此 $s'\in F_i(T)$。由于 $F_i$ 是 fppf 拓扑的层，遂有 $s=s'$。

\medskip\noindent
设 $\{T_j \to T\}$ 是 $(\Sch/S)_{fppf}$ 中的 fppf 覆盖，并设
$s_j \in F(T_j)$ 满足
$s_j|_{T_j \times_T T_{j'}} = s_{j'}|_{T_j \times_T T_{j'}}$。
由假设 (2)(c)，可将覆盖加细，并假定对某个 $\alpha(j)\in I$ 有
$s_j\in F_{\alpha(j)}(T_j)$。
令 $W_j\subset T$ 为 $T_j\to T$ 的像；由《态射》章，引理
\ref{morphisms-lemma-fppf-open}，它是开子概形。
于是 $\{T_j\to W_j\}$ 是 fppf 覆盖。
由于 $F_{\alpha(j)}$ 是 $F$ 的子预层，$s_j$ 在
$T_j\times_{W_j}T_j$ 上的两个限制作为
$F_{\alpha(j)}(T_j\times_{W_j}T_j)$ 的元素相等。
因此由 $F_{\alpha(j)}$ 的层条件，存在
$s'_j\in F_{\alpha(j)}(W_j)$，其在 $T_j$ 上的限制为 $s_j$。
对任意一对指标 $j$、$j'$，由上一段的结论，$F$ 的截面
$s'_j|_{W_j\cap W_{j'}}$ 与
$s'_{j'}|_{W_j\cap W_{j'}}$ 相等。
最后利用 $F$ 是 Zariski 层，即得证明。
\end{proof}





\section{可由代数空间表示的预层态射的性质}
\label{bootstrap-section-representable-by-spaces-properties}

\noindent
下面的定义使这一做法得以成立。

\begin{definition}
\label{bootstrap-definition-property-transformation}
设 $S$ 为概形，$a : F \to G$ 为 $(\Sch/S)_{fppf}$ 上预层的态射，
且该态射可由代数空间表示。
设 $\mathcal{P}$ 为代数空间态射的一项性质，并满足
\begin{enumerate}
\item 在任意基变换下保持；并且
\item 相对于基底在 fppf 意义下是局部的，见《空间上的下降》章，
定义 \ref{spaces-descent-definition-property-morphisms-local}。
\end{enumerate}
此时，若对每个概形 $U$ 及每个 $\xi : U \to G$，所得代数空间态射
$U \times_G F \to U$ 都具有性质 $\mathcal{P}$，
则称 $a$ {\it 具有性质 $\mathcal{P}$}。
\end{definition}

\noindent
需要强调的是，我们只对在基变换下稳定且相对于基底在 fppf 拓扑下局部的
态射性质使用这一定义。这并不是因为在其他情形下定义没有意义；
而是因为对所考虑的具体性质，我们可能希望采用另一个更合适的定义。

\medskip\noindent
上述定义例如适用于\footnote{基变换下保持性由《空间的态射》章的下列引理给出：
\ref{spaces-morphisms-lemma-base-change-surjective},
\ref{spaces-morphisms-lemma-base-change-quasi-compact},
\ref{spaces-morphisms-lemma-base-change-etale},
\ref{spaces-morphisms-lemma-base-change-smooth},
\ref{spaces-morphisms-lemma-base-change-flat},
\ref{spaces-morphisms-lemma-base-change-separated},
\ref{spaces-morphisms-lemma-base-change-finite-type},
\ref{spaces-morphisms-lemma-base-change-quasi-finite},
\ref{spaces-morphisms-lemma-base-change-finite-presentation},
\ref{spaces-morphisms-lemma-base-change-affine},
\ref{spaces-morphisms-lemma-base-change-proper}，
以及《空间》章，引理
\ref{spaces-lemma-base-change-immersions}。
相对于基底的 fppf 局部性由《空间上的下降》章的下列引理给出：
\ref{spaces-descent-lemma-descending-property-surjective},
\ref{spaces-descent-lemma-descending-property-quasi-compact},
\ref{spaces-descent-lemma-descending-property-etale},
\ref{spaces-descent-lemma-descending-property-smooth},
\ref{spaces-descent-lemma-descending-property-flat},
\ref{spaces-descent-lemma-descending-property-separated},
\ref{spaces-descent-lemma-descending-property-finite-type},
\ref{spaces-descent-lemma-descending-property-quasi-finite},
\ref{spaces-descent-lemma-descending-property-locally-finite-presentation},
\ref{spaces-descent-lemma-descending-property-affine},
\ref{spaces-descent-lemma-descending-property-proper}，
以及
\ref{spaces-descent-lemma-descending-property-closed-immersion}.
}
“满射”“拟紧”“平展”“光滑”“平坦”“分离”“（局部）有限型”
“（局部）拟有限”“（局部）有限呈示”“仿射”“固有”以及“闭浸入”等性质。
换言之，称 $a$ {\it 为满射}
（相应地，{\it 拟紧}、{\it 平展}、{\it 光滑}、{\it 平坦}、
{\it 分离}、{\it （局部）有限型}、{\it （局部）拟有限}、
{\it （局部）有限呈示}、{\it 固有}、{\it 闭浸入}），
是指对每个概形 $T$ 及态射 $\xi : T \to G$，
代数空间态射 $T \times_{\xi, G} F \to T$ 为满射
（相应地，拟紧、平展、平坦、分离、（局部）有限型、
（局部）拟有限、（局部）有限呈示、固有、闭浸入）。

\medskip\noindent
下面检验该定义与已有概念相容。
由引理 \ref{bootstrap-lemma-morphism-spaces-is-representable-by-spaces}，
$S$ 上任意两个代数空间之间的态射都可由代数空间表示。
再由《空间的态射》章，引理
\ref{spaces-morphisms-lemma-surjective-local}
(resp.\ \ref{spaces-morphisms-lemma-quasi-compact-local},
\ref{spaces-morphisms-lemma-etale-local},
\ref{spaces-morphisms-lemma-smooth-local},
\ref{spaces-morphisms-lemma-flat-local},
\ref{spaces-morphisms-lemma-separated-local},
\ref{spaces-morphisms-lemma-finite-type-local},
\ref{spaces-morphisms-lemma-quasi-finite-local},
\ref{spaces-morphisms-lemma-finite-presentation-local},
\ref{spaces-morphisms-lemma-affine-local},
\ref{spaces-morphisms-lemma-proper-local},
\ref{spaces-morphisms-lemma-closed-immersion-local})
，上述满射（相应地，拟紧、平展、光滑、平坦、分离、
（局部）有限型、（局部）拟有限、（局部）有限呈示、仿射、固有、闭浸入）
的定义与代数空间态射的已有定义一致。

\medskip\noindent
下面给出一些形式性的引理。

\begin{lemma}
\label{bootstrap-lemma-base-change-transformation-property}
设 $S$ 为概形。设 $\mathcal{P}$ 为定义
\ref{bootstrap-definition-property-transformation} 中的性质。设
$$
\xymatrix{
G' \times_G F \ar[r] \ar[d]^{a'} & F \ar[d]^a \\
G' \ar[r] & G
}
$$
是 $(\Sch/S)_{fppf}$ 上预层的纤维方块。
若 $a$ 可由代数空间表示且具有性质 $\mathcal{P}$，
则 $a'$ 也如此。
\end{lemma}

\begin{proof}
略。提示：这是一个形式推论。
\end{proof}

\begin{lemma}
\label{bootstrap-lemma-composition-transformation-property}
设 $S$ 为概形。设 $\mathcal{P}$ 为定义
\ref{bootstrap-definition-property-transformation} 中的性质，
并假设 $\mathcal{P}$ 在复合下稳定。设
$$
\xymatrix{
F \ar[r]^a & G \ar[r]^b & H
}
$$
是 $(\Sch/S)_{fppf}$ 上预层的态射。
若 $a$、$b$ 均可由代数空间表示且具有性质 $\mathcal{P}$，
则 $b \circ a$ 也如此。
\end{lemma}

\begin{proof}
略。提示：参见引理 \ref{bootstrap-lemma-composition-transformation}，
并使用复合下的稳定性。
\end{proof}

\begin{lemma}
\label{bootstrap-lemma-product-transformations-property}
设 $S$ 为概形。令
$F_i, G_i : (\Sch/S)_{fppf}^{opp} \to \textit{Sets}$，$i=1,2$。
设 $a_i : F_i \to G_i$（$i=1,2$）可由代数空间表示。
设 $\mathcal{P}$ 为定义 \ref{bootstrap-definition-property-transformation}
中的性质，并且在复合下稳定。
若 $a_1$ 和 $a_2$ 都具有性质 $\mathcal{P}$，则
$a_1 \times a_2 : F_1 \times F_2 \longrightarrow G_1 \times G_2$
也具有性质 $\mathcal{P}$。
\end{lemma}

\begin{proof}
由引理 \ref{bootstrap-lemma-product-transformations}，本引理的陈述确有意义。
证明略。
\end{proof}

\begin{lemma}
\label{bootstrap-lemma-transformations-property-implication}
设 $S$ 为概形，并设
$F, G : (\Sch/S)_{fppf}^{opp} \to \textit{Sets}$。
设 $a : F \to G$ 为可由代数空间表示的函子态射。
设 $\mathcal{P}$、$\mathcal{P}'$ 为定义
\ref{bootstrap-definition-property-transformation} 中的性质。
假设对 $S$ 上代数空间的任意态射 $f:X\to Y$ 都有
$\mathcal{P}(f) \Rightarrow \mathcal{P}'(f)$。
若 $a$ 具有性质 $\mathcal{P}$，则 $a$ 具有性质 $\mathcal{P}'$。
\end{lemma}

\begin{proof}
这是形式推论。
\end{proof}

\begin{lemma}
\label{bootstrap-lemma-surjective-flat-locally-finite-presentation}
设 $S$ 为概形，并设
$F,G:(\Sch/S)_{fppf}^{opp}\to\textit{Sets}$ 为层。
设 $a:F\to G$ 可由代数空间表示、平坦、局部有限呈示且为满射。
则作为层的态射，$a:F\to G$ 是满射。
\end{lemma}

\begin{proof}
设 $T$ 为 $S$ 上的概形，$g:T\to G$ 为 $G$ 的一个 $T$ 值点。
由假设，$T'=F\times_G T$ 是代数空间，且态射 $T'\to T$
是平坦、局部有限呈示的满态射。
设 $U\to T'$ 为满的平展态射，其中 $U$ 是概形。
由代数空间平坦态射的定义，概形态射 $U\to T$ 是平坦的；
“局部有限呈示”也同理。态射 $U\to T$ 还是满的，见《空间的态射》章，
引理 \ref{spaces-morphisms-lemma-surjective-local}。
因此 $\{U\to T\}$ 是一个 fppf 覆盖，并且
$g|_U\in G(U)$ 来自 $F(U)$ 的一个元素，即态射 $U\to T'\to F$。
这证明该态射作为层的态射是满射，见《位点》章，定义
\ref{sites-definition-sheaves-injective-surjective}。
\end{proof}




\section{对角态射的自举}
\label{bootstrap-section-bootstrap-diagonal}

\noindent
本节证明：对 $(\Sch/S)_{fppf}$ 上的层 $F$，只要存在一个由概形或
代数空间给出的 $F$ 的“fppf 覆盖”，其对角态射便是可表的，见
引理 \ref{bootstrap-lemma-bootstrap-diagonal}。

\begin{lemma}
\label{bootstrap-lemma-representable-diagonal}
\begin{slogan}
一个预层的对角态射可由代数空间表示，当且仅当从任意概形到该预层的
每个态射均可由代数空间表示。
\end{slogan}
设 $S$ 为概形，$F$ 为 $(\Sch/S)_{fppf}$ 上的预层。
下列条件等价：
\begin{enumerate}
\item $\Delta_F : F \to F \times F$ 可由代数空间表示；
\item 对每个概形 $T$，任意态射 $T \to F$ 均可由代数空间表示；
\item 对每个代数空间 $X$，任意态射 $X \to F$ 均可由代数空间表示。
\end{enumerate}
\end{lemma}

\begin{proof}
假设 (1) 成立。取 (3) 中的 $X\to F$。设 $T$ 为概形，且
$T\to F$ 为态射。则
$$
T \times_F X = (T \times_S X) \times_{F \times F, \Delta} F
$$
由引理 \ref{bootstrap-lemma-representable-by-spaces} 及 (1)，这是一个代数空间。
因此 $X\to F$ 是可表的，即 (3) 成立。
蕴含 (3) $\Rightarrow$ (2) 是显然的。现假设 (2) 成立。
设 $T$ 为概形，且 $(a,b):T\to F\times F$ 为态射。则
$$
F \times_{\Delta_F, F \times F} T =
(T \times_{a, F, b} T) \times_{T \times T, \Delta_T} T
$$
由假设，这是一个代数空间。因此 $\Delta_F$ 可由代数空间表示，
即 (1) 成立。
\end{proof}

\noindent
特别地，若预层 $F$ 满足该引理中的等价条件，则对任意态射
$X\to F$（其中 $X$ 是代数空间），可以按照定义
\ref{bootstrap-definition-property-transformation} 有意义地说 $X\to F$
是满射（相应地，平展、平坦、局部有限呈示）。

\medskip\noindent
在实际进行自举之前，先证明一个有趣的引理。

\begin{lemma}
\label{bootstrap-lemma-after-fppf-sep-lqf}
设 $S$ 为概形。设
$$
\xymatrix{
E \ar[r]_a \ar[d]_f & F \ar[d]^g \\
H \ar[r]^b & G
}
$$
是 $(\Sch/S)_{fppf}$ 上层的笛卡儿图，因而
$E=H\times_G F$。若
\begin{enumerate}
\item $g$ 可由代数空间表示、为满射、平坦且局部有限呈示；
\item $a$ 可由代数空间表示、分离且局部拟有限，
\end{enumerate}
则 $b$ 可表（由概形表示），并且分离且局部拟有限。
\end{lemma}

\begin{proof}
设 $T$ 为概形，且 $T\to G$ 为态射。需要证明 $T\times_GH$ 是概形，
并且态射 $T\times_GH\to T$ 分离且局部拟有限。
因此，可将整个图基变换到 $T$，并假设 $G$ 是概形；此时 $F$ 是代数空间。
设 $U$ 为概形，且 $U\to F$ 为满的平展态射。
由《空间的态射》章，引理 \ref{spaces-morphisms-lemma-etale-flat}
及 \ref{spaces-morphisms-lemma-etale-locally-finite-presentation}，
$U\to F$ 可表、为满射、平坦且局部有限呈示。
由引理 \ref{bootstrap-lemma-composition-transformation}，
$U\to G$ 也为满射、平坦且局部有限呈示。
注意，$a$ 的基变换 $E\times_FU\to U$ 仍然分离且局部拟有限
（由引理 \ref{bootstrap-lemma-base-change-transformation-property}）。
因此，可以用 $E\times_FU\to U$ 替换引理之图的上半部分。
换言之，可以假定 $F\to G$ 是概形之间局部有限呈示的平坦满态射。
特别地，$\{F\to G\}$ 是概形的一个 fppf 覆盖。
由《空间的态射》章，命题
\ref{spaces-morphisms-proposition-locally-quasi-finite-separated-over-scheme}，
可知 $E$ 也是概形。
由《下降》章，引理 \ref{descent-lemma-descent-data-sheaves}，
等式 $E=H\times_GF$ 表明，相对于 fppf 覆盖 $\{F\to G\}$，
在 $E$ 上得到一个下降数据。
由《态射进阶》章，引理
\ref{more-morphisms-lemma-separated-locally-quasi-finite-morphisms-fppf-descend}，
该下降数据是有效的。
再次应用《下降》章，引理 \ref{descent-lemma-descent-data-sheaves}，
可知 $H$ 是概形。
最后由《下降》章，引理 \ref{descent-lemma-descending-property-separated}
及 \ref{descent-lemma-descending-property-quasi-finite}，
得到 $b$ 分离且局部拟有限。
\end{proof}

\noindent
下面给出本节标题所指的结果。

\begin{lemma}
\label{bootstrap-lemma-bootstrap-diagonal}
设 $S$ 为概形，并设
$F:(\Sch/S)_{fppf}^{opp}\to\textit{Sets}$ 为函子。假设
\begin{enumerate}
\item 预层 $F$ 是层；
\item 存在代数空间 $X$ 及态射 $X\to F$，该态射可由代数空间表示、
为满射、平坦且局部有限呈示。
\end{enumerate}
则 $\Delta_F$ 可表（由概形表示）。
\end{lemma}

\begin{proof}
设 $U\to X$ 是从概形 $U$ 到 $X$ 的满平展态射。
由《空间的态射》章，引理 \ref{spaces-morphisms-lemma-etale-flat}
及 \ref{spaces-morphisms-lemma-etale-locally-finite-presentation}，
$U\to X$ 可表、为满射、平坦且局部有限呈示。
由引理 \ref{bootstrap-lemma-composition-transformation-property}，
复合 $U\to F$ 也可由代数空间表示、为满射、平坦且局部有限呈示。
因此 $R=U\times_FU$ 是代数空间，见引理
\ref{bootstrap-lemma-representable-by-spaces}。
代数空间态射 $R\to U\times_SU$ 是单态射，故为分离态射
（因为单态射的对角态射是同构，见《空间的态射》章，引理
\ref{spaces-morphisms-lemma-monomorphism}）。
由于 $U\to F$ 局部有限呈示，两个态射 $R\to U$ 均局部有限呈示，
见引理 \ref{bootstrap-lemma-base-change-transformation-property}。
因此 $R\to U\times_SU$ 局部有限型（使用《空间的态射》章，引理
\ref{spaces-morphisms-lemma-finite-presentation-finite-type} 及
\ref{spaces-morphisms-lemma-permanence-finite-type}）。
综上，$R\to U\times_SU$ 是局部有限型的单态射，因而是分离且局部拟有限的态射，
见《空间的态射》章，引理
\ref{spaces-morphisms-lemma-monomorphism-loc-finite-type-loc-quasi-finite}。

\medskip\noindent
现在来证明 $\Delta_F$ 可表。
设 $T$ 为概形，且 $(a,b):T\to F\times F$ 为态射。令
$$
T' = (U \times_S U) \times_{F \times F} T.
$$
由引理 \ref{bootstrap-lemma-product-transformations-property}，
$U\times_SU\to F\times F$ 可由代数空间表示、为满射、平坦且局部有限呈示。
因此 $T'$ 是代数空间，投影态射 $T'\to T$ 为满射、平坦且局部有限呈示。
考虑 $Z=T\times_{F\times F}F$（这是一个层）以及
$$
Z' = T' \times_{U \times_S U} R
= T' \times_T Z.
$$
可知 $Z'$ 是代数空间。由本证明第一段中关于 $R$ 是代数空间且
态射 $R\to U\times_SU$ 具有相应性质的讨论，
$Z'\to T'$ 分离且局部拟有限。
因此可对下图应用引理 \ref{bootstrap-lemma-after-fppf-sep-lqf}：
$$
\xymatrix{
Z' \ar[r] \ar[d] & T' \ar[d] \\
Z \ar[r] & T
}
$$
从而得到结论。
\end{proof}

\noindent
下面是上述结果的一个变式。

\begin{lemma}
\label{bootstrap-lemma-bootstrap-locally-quasi-finite}
设 $S$ 为概形，$F:(\Sch/S)_{fppf}^{opp}\to\textit{Sets}$ 为函子。
设 $X$ 为概形，且 $X\to F$ 可由代数空间表示并局部拟有限。
则 $X\to F$ 可表（由概形表示）。
\end{lemma}

\begin{proof}
设 $T$ 为概形，且 $T\to F$ 为态射。需要证明代数空间
$X\times_FT$ 可由概形表示。考虑态射
$$
X \times_F T  \longrightarrow X \times_{\Spec(\mathbf{Z})} T
$$
由于 $X\times_FT\to T$ 局部拟有限，所显示的箭头也局部拟有限
（《空间的态射》章，引理
\ref{spaces-morphisms-lemma-permanence-quasi-finite}）。
另一方面，所显示的箭头是单态射，因而分离
（《空间的态射》章，引理
\ref{spaces-morphisms-lemma-monomorphism-separated}）。
因此由《空间的态射》章，命题
\ref{spaces-morphisms-proposition-locally-quasi-finite-separated-over-scheme}，
$X\times_FT$ 是概形。
\end{proof}











\section{自举}
\label{bootstrap-section-bootstrap}

\noindent
先提醒读者，本节的结果将被更强的定理
\ref{bootstrap-theorem-final-bootstrap} 所取代。
另一方面，本节定理的证明容易得多，同时仍能充分揭示其中的机制，
尤其适合主要关注 Deligne--Mumford 叠的读者。

\medskip\noindent
在《空间》章，第 \ref{spaces-section-algebraic-spaces} 节中，
代数空间被定义为 fppf 拓扑下具有可表对角态射的层，
并且存在从某个概形到它的满平展态射。
本节证明：若 fppf 拓扑下的层具有可由代数空间表示的对角态射，
并具有由某个代数空间给出的满平展覆盖，则它也是代数空间。
换言之，代数空间范畴是对概形范畴的扩充，加入了那些具有可表对角态射
以及概形平展覆盖的 fppf 层 $F$。
本节结果表明，从代数空间范畴出发再次执行同样的过程，并不会产生又一个新范畴。

\medskip\noindent
本节内容的另一动机是，它将在后文保证：惯性叠平凡的 Deligne--Mumford 叠
等价于一个代数空间，见《代数叠》章，引理
\ref{algebraic-lemma-algebraic-stack-no-automorphisms}。

\medskip\noindent
下面是本节的主要结果（如上所述，它将被更强的定理
\ref{bootstrap-theorem-final-bootstrap} 所取代）。

\begin{theorem}
\label{bootstrap-theorem-bootstrap}
设 $S$ 为概形，$F:(\Sch/S)_{fppf}^{opp}\to\textit{Sets}$ 为函子。
假设
\begin{enumerate}
\item 预层 $F$ 是层；
\item 对角态射 $F\to F\times F$ 可由代数空间表示；并且
\item 存在代数空间 $X$ 及满平展态射 $X\to F$。
\end{enumerate}
或者假设
\begin{enumerate}
\item[(a)] 预层 $F$ 是层；并且
\item[(b)] 存在代数空间 $X$ 及态射 $X\to F$，
该态射可由代数空间表示、为满射且平展。
\end{enumerate}
则 $F$ 是代数空间。
\end{theorem}

\begin{proof}
以下将使用定义 \ref{bootstrap-definition-property-transformation}
紧随其后的说明，而不再逐一指出。

\medskip\noindent
假设 (1)、(2)、(3) 成立，并取 (3) 中的 $X\to F$。
由引理 \ref{bootstrap-lemma-representable-diagonal}，态射
$X\to F$ 可由代数空间表示。因此 (a)、(b) 成立。

\medskip\noindent
假设 (a)、(b) 成立，并取 (b) 中的 $X\to F$。
设 $U\to X$ 是从概形 $U$ 到 $X$ 的满平展态射。
由引理 \ref{bootstrap-lemma-composition-transformation}，复合态射
$U\to F$ 可由代数空间表示、为满射且平展。
因此，要证明 $F$ 是代数空间，只须证明 $\Delta_F$ 可表
（《空间》章，定义 \ref{spaces-definition-algebraic-space}）。
这立即由引理 \ref{bootstrap-lemma-bootstrap-diagonal} 得出。
另一方面，也可以绕开该引理，按下一段直接证明 $F$ 是代数空间。

\medskip\noindent
具体地，设 $U$ 为概形，且 $U\to F$ 可由代数空间表示、为满射且平展。
考虑纤维积 $R=U\times_FU$。两个投影 $R\to U$ 均可由代数空间表示、
为满射且平展（引理 \ref{bootstrap-lemma-base-change-transformation-property}）。
特别地，由引理 \ref{bootstrap-lemma-representable-by-spaces-over-space}，
$R$ 是代数空间。
代数空间态射 $R\to U\times_SU$ 是单态射，因而分离
（因为单态射的对角态射是同构）。
由于 $R\to U$ 平展，可知 $R\to U$ 局部拟有限，见《空间的态射》章，
引理 \ref{spaces-morphisms-lemma-etale-locally-quasi-finite}。
再由《空间的态射》章，引理
\ref{spaces-morphisms-lemma-permanence-quasi-finite}，
$R\to U\times_SU$ 也局部拟有限。
于是可应用《空间的态射》章，命题
\ref{spaces-morphisms-proposition-locally-quasi-finite-separated-over-scheme}，
从而 $R$ 是概形。
由引理 \ref{bootstrap-lemma-surjective-flat-locally-finite-presentation}，
$U\to F$ 是层的满射。因此 $F=U/R$。
最后由《空间》章，定理 \ref{spaces-theorem-presentation}，
$F$ 是代数空间。
\end{proof}









\section{寻找开子空间}
\label{bootstrap-section-finding-opens}


\medskip\noindent
首先证明一个引理；它把《空间》章，引理
\ref{spaces-lemma-finding-opens} 略加改进，并推广到群胚所关联的商层。

\begin{lemma}
\label{bootstrap-lemma-better-finding-opens}
设 $S$ 为概形，$(U,R,s,t,c)$ 为 $S$ 上的群胚概形，
$g:U'\to U$ 为态射。假设
\begin{enumerate}
\item 复合
$$
\xymatrix{
U' \times_{g, U, t} R \ar[r]_-{\text{pr}_1} \ar@/^3ex/[rr]^h
& R \ar[r]_s & U
}
$$
的像是开子概形 $W\subset U$；并且
\item 所得态射 $h:U'\times_{g,U,t}R\to W$
在 fppf 拓扑下给出层的满射。
\end{enumerate}
令 $R'=R|_{U'}$ 为 $R$ 在 $U'$ 上的限制。则 fppf 拓扑下的商层态射
$$
U'/R' \to U/R
$$
是可表的，并且是开浸入。
\end{lemma}

\begin{proof}
注意，$W$ 是 $U$ 的 $R$-不变开子概形。
这是因为，$W$ 的点恰好是 $U$ 中按《群胚》章，引理
\ref{groupoids-lemma-pre-equivalence-equivalence-relation-points}
之意义与 $g(U')\subset U$ 中某点等价的那些点
（该引理适用，是因为由《群胚》章，引理
\ref{groupoids-lemma-groupoid-pre-equivalence}，
$j:R\to U\times_SU$ 是预等价关系）。
此外，$g:U'\to U$ 经 $W$ 分解。
令 $R|_W$ 为 $R$ 在 $W$ 上的限制。
于是 $R'$ 也是 $R|_W$ 在 $U'$ 上的限制。
因此，引理中的层态射可分解为
$$
U'/R' \longrightarrow W/R|_W \longrightarrow U/R
$$
由《群胚》章，引理 \ref{groupoids-lemma-quotient-groupoid-restrict}，
第一支箭头是层的同构。
所以只须在 $g$ 是某个 $R$-不变开子概形到 $U$ 的浸入时证明本引理。

\medskip\noindent
假设 $U'\subset U$ 是 $R$-不变开子概形，且 $g$ 为包含态射。
令 $F=U/R$、$F'=U'/R'$。由《群胚》章，引理
\ref{groupoids-lemma-quotient-pre-equivalence-relation-restrict}
或 \ref{groupoids-lemma-quotient-groupoid-restrict}，
态射 $F'\to F$ 是单射。设 $\xi\in F(T)$。
需要证明 $T\times_{\xi,F}F'$ 可由 $T$ 的开子概形表示。
存在 fppf 覆盖 $\{f_i:T_i\to T\}$，使得 $\xi|_{T_i}$
是某态射 $a_i:T_i\to U$ 在 $U\to U/R$ 下的像。
令 $V_i=a_i^{-1}(U')$。
断言作为 $T_i\times_TT_j$ 的开子概形，有
$V_i\times_TT_j=T_i\times_TV_j$。

\medskip\noindent
由于 $a_i\circ\text{pr}_0$ 与 $a_j\circ\text{pr}_1$ 是
$T_i\times_TT_j\to U$ 的态射，且二者均映到截面
$\xi|_{T_i\times_TT_j}\in F(T_i\times_TT_j)$，
可以找到 fppf 覆盖
$\{f_{ijk}:T_{ijk}\to T_i\times_TT_j\}$ 及态射
$r_{ijk}:T_{ijk}\to R$，使得
$$
a_i \circ \text{pr}_0 \circ f_{ijk} = s \circ r_{ijk},
\quad
a_j \circ \text{pr}_1 \circ f_{ijk} = t \circ r_{ijk},
$$
见《群胚》章，引理 \ref{groupoids-lemma-quotient-pre-equivalence}。
由于 $U'$ 是 $R$-不变的，有 $s^{-1}(U')=t^{-1}(U')$，因而
$f_{ijk}^{-1}(V_i\times_TT_j)=f_{ijk}^{-1}(T_i\times_TV_j)$。
由于 $\{f_{ijk}\}$ 是满射，这便证明了上述断言。
于是由《下降》章，引理 \ref{descent-lemma-open-fpqc-covering}，
存在开子概形 $V\subset T$ 使得 $f_i^{-1}(V)=V_i$。
断言 $V$ 表示 $T\times_{\xi,F}F'$。

\medskip\noindent
第一步证明 $\xi|_V\in F'(V)\subset F(V)$。
态射族 $\{V_i\to V\}$ 是 fppf 覆盖，并且由构造有
$\xi|_{V_i}\in F'(V_i)$。所以由 $F'$ 的层性质，
$\xi|_V\in F'(V)$。
最后，设 $T'\to T$ 为概形态射，并且 $\xi|_{T'}\in F'(T')$。
为完成证明，需要说明 $T'\to T$ 经 $V$ 分解。
可以找到 fppf 覆盖 $\{T'_j\to T'\}_{j\in J}$ 及态射
$b_j:T'_j\to U'$，使得 $\xi|_{T'_j}$ 是 $b_j$ 在
$U'\to U/R$ 下的像。显然，只须证明复合 $T'_j\to T$ 经 $V$ 分解。
因此，可以假定 $\xi|_{T'}$ 是某态射 $b:T'\to U'$ 的像。
此时只须证明 $T'\times_TT_i\to T_i$ 经 $V_i$ 分解。
在概形 $T'\times_TT_i$ 上，$\xi$ 的限制是
$(U/R)(T'\times_TT_i)$ 中两个元素的像，即
$a_i\circ\text{pr}_1$ 与 $b\circ\text{pr}_0$；
后者经 $R$-不变开子概形 $U'$ 分解。
由《群胚》章，引理 \ref{groupoids-lemma-quotient-pre-equivalence}，
存在覆盖 $\{h_k:Z_k\to T'\times_TT_i\}$ 及态射 $r_k:Z_k\to R$，
使得 $a_i\circ\text{pr}_1\circ h_k=s\circ r_k$ 且
$b\circ\text{pr}_0\circ h_k=t\circ r_k$。
由于 $U'$ 是 $R$-不变开子概形，而 $b$ 的像包含在 $U'$ 中，
每个 $a_i\circ\text{pr}_1\circ h_k$ 的像也包含在 $U'$ 中。
按 $V_i$ 的定义，这说明 $T'\times_TT_i\to T_i$ 的像包含在 $V_i$ 中，
从而完成证明。
\end{proof}












\section{等价关系的切片}
\label{bootstrap-section-slicing}

\noindent
本节说明如何通过切片来“改进”一个给定的等价关系。
这不是读者可能熟悉的某种“平展切片”，而是一种粗糙得多的切片。


\begin{lemma}
\label{bootstrap-lemma-slice-equivalence-relation}
设 $S$ 为概形，$j:R\to U\times_SU$ 为 $S$ 上概形的等价关系。
假设 $s,t:R\to U$ 平坦且局部有限呈示。
则存在 $S$ 上概形的等价关系 $j':R'\to U'\times_SU'$，以及同构
$$
U'/R' \longrightarrow U/R
$$
它由一个把 $R'$ 映入 $R$ 的态射 $U'\to U$ 诱导，并且
$s',t':R\to U$ 平坦、局部有限呈示且局部拟有限。
\end{lemma}

\begin{proof}
分若干步证明本引理。以下将不再逐一说明地使用这些事实：
等价关系给出一个群胚概形，而等价关系的限制仍是等价关系。见《群胚》章，引理
\ref{groupoids-lemma-restrict-relation},
\ref{groupoids-lemma-equivalence-groupoid} 以及
\ref{groupoids-lemma-restrict-groupoid-relation}.

\medskip\noindent
步骤 1：可以假定 $s,t:R\to U$ 是局部有限呈示的 Cohen--Macaulay 态射。
事实上，如《群胚进阶》章，引理 \ref{more-groupoids-lemma-make-CM}，
取开子概形 $g:U'\to U$，使得 $t^{-1}(U')\subset R$ 是
$s:R\to U$ 在其上为 Cohen--Macaulay 态射的最大开集，
并以 $R'$ 表示 $R$ 在 $U'$ 上的限制。由上述引理可知
$$
\xymatrix{
t^{-1}(U') \ar@{=}[r] &
U' \times_{g, U, t} R \ar[r]_-{\text{pr}_1} \ar@/^3ex/[rr]^h &
R \ar[r]_s &
U
}
$$
是满射。由于 $h$ 平坦且局部有限呈示，$\{h\}$ 是 fppf 覆盖。
所以由《群胚》章，引理 \ref{groupoids-lemma-quotient-groupoid-restrict}，
$U'/R'\to U/R$ 是同构。由 $U'$ 的构造，$s',t'$ 为
Cohen--Macaulay 且局部有限呈示。

\medskip\noindent
步骤 2：假设 $s,t$ 为 Cohen--Macaulay 且局部有限呈示。
设 $u\in U$ 为有限型点。由《群胚进阶》章，引理
\ref{more-groupoids-lemma-max-slice-quasi-finite}，
存在仿射概形 $U'$ 及态射 $g:U'\to U$，使得
\begin{enumerate}
\item $g$ 是浸入；
\item $u\in U'$；
\item $g$ 局部有限呈示；
\item $h$ 平坦、局部有限呈示且局部拟有限；
\item 态射 $s',t':R'\to U'$ 平坦、局部有限呈示且局部拟有限。
\end{enumerate}
这里使用了《群胚进阶》章，情形
\ref{more-groupoids-situation-slice} 中引入的记号。

\medskip\noindent
步骤 3：对每个有限型点 $u\in U$，选择步骤 2 中的
$g_u:U'_u\to U$，并以 $R'_u$ 表示 $R$ 在 $U'_u$ 上的限制。
记 $h_u=s\circ\text{pr}_1:U'_u\times_{g_u,U,t}R\to U$。令
$U'=\coprod_{u\in U}U'_u$ 且 $g=\coprod g_u$。
如上令 $R'$ 为 $R$ 在 $U'$ 上的限制。断言二元组 $(U',g)$ 满足要求
\footnote{这里应检验 $U'$ 不会过大，即它同构于范畴
$\Sch_{fppf}$ 中的一个对象；见第 \ref{bootstrap-section-conventions} 节。
这是纯粹的集合论问题。使用《集合》章，第
\ref{sets-section-categories-schemes} 节引入的概形大小概念。
每个 $U'_u$ 的大小至多为 $U$ 的大小，而指标集的基数至多为
$|U|$ 的基数，后者又以 $U$ 的大小为界。
因此由《集合》章，引理 \ref{sets-lemma-what-is-in-it} 的 (6)，
$U'$ 同构于 $\Sch_{fppf}$ 中的一个对象。}。注意
\begin{align*}
R' = &
\coprod\nolimits_{u_1, u_2 \in U}
(U'_{u_1} \times_{g_{u_1}, U, t} R)
\times_R
(R \times_{s, U, g_{u_2}} U'_{u_2}) \\
= &
\coprod\nolimits_{u_1, u_2 \in U}
(U'_{u_1} \times_{g_{u_1}, U, t} R) \times_{h_{u_1}, U, g_{u_2}} U'_{u_2}
\end{align*}
投影 $s':R'\to U'=\coprod U'_{u_2}$ 是 $\coprod h_{u_1}$ 的基变换，
因而平坦、局部有限呈示且局部拟有限。
最后，由构造，态射 $h:U'\times_{g,U,t}R\to U$ 等于
$\coprod h_u$，所以其像包含 $U$ 的所有有限型点。
由于每个 $h_u$ 都平坦且局部有限呈示，$h$ 也平坦且局部有限呈示。
特别地，$h$ 的像是开的（见《态射》章，引理
\ref{morphisms-lemma-fppf-open}）；又因有限型点的集合稠密（见《态射》章，
引理 \ref{morphisms-lemma-enough-finite-type-points}），
可知 $h$ 的像是 $U$。因此 $\{h\}$ 是 fppf 覆盖。
由《群胚》章，引理 \ref{groupoids-lemma-quotient-groupoid-restrict}，
这意味着 $U'/R'\to U/R$ 是同构，从而完成引理的证明。
\end{proof}










\section{对子群胚取商}
\label{bootstrap-section-dividing}

\noindent
在进行最终自举之前，还需要一个引理。
在着手处理概形论版本之前，先用“普通”群胚来说明其中的情形。

\medskip\noindent
设 $\mathcal{C}$ 为群胚，见《范畴》章，定义
\ref{categories-definition-groupoid}。
如《群胚》章，第 \ref{groupoids-section-groupoids} 节所述，
它对应于五元组 $(\text{Ob},\text{Arrows},s,t,c)$。
假设给定子集 $P\subset\text{Arrows}$，使得
$(\text{Ob},P,s|_P,t|_P,c|_P)$ 也是群胚，且 $P$ 中没有非平凡自同构。
于是可以如下构造商群胚
$(\overline{\text{Ob}}, \overline{\text{Arrows}}, \overline{s},
\overline{t}, \overline{c})$
：
\begin{enumerate}
\item $\overline{\text{Ob}}=\text{Ob}/P$ 是 $P$-同构类的集合；
\item $\overline{\text{Arrows}}=P\backslash\text{Arrows}/P$
是 $\mathcal{C}$ 中的箭头模去与 $P$ 中箭头的前复合和后复合所得的集合；
\item 源映射和靶映射
$\overline{s},\overline{t}:P\backslash\text{Arrows}/P\to\text{Ob}/P$
由 $s,t$ 诱导；
\item 复合由良定义的规则
$\overline{c}(\overline{a},\overline{b})=\overline{c(a,b)}$ 定义。
\end{enumerate}
事实上，原群胚 $(\text{Ob},\text{Arrows},s,t,c)$ 典范同构于群胚
$(\overline{\text{Ob}},\overline{\text{Arrows}},\overline{s},
\overline{t},\overline{c})$ 沿商映射
$g:\text{Ob}\to\overline{\text{Ob}}$ 的限制
（见《群胚》章，第 \ref{groupoids-section-restrict-groupoid} 节的讨论）。
回忆这意味着
$$
\text{Arrows} =
\text{Ob}
\times_{g, \overline{\text{Ob}}, \overline{t}}
\overline{\text{Arrows}}
\times_{\overline{s}, \overline{\text{Ob}}, g}
\text{Ob}
$$
该等式成立是因为 $P$ 没有非平凡自同构。细节从略。

\medskip\noindent
下面的引理在远为一般的情形下仍然成立，但这里给出的是最终自举证明中
所用的版本（此后可以更容易地证明该引理的更一般版本）。

\begin{lemma}
\label{bootstrap-lemma-divide-subgroupoid}
设 $S$ 为概形，$(U,R,s,t,c)$ 为 $S$ 上的群胚概形，
$P\to R$ 为概形的单态射。假设
\begin{enumerate}
\item $(U,P,s|_P,t|_P,c|_{P\times_{s,U,t}P})$ 是群胚概形；
\item $s|_P,t|_P:P\to U$ 有限局部自由；
\item $j|_P:P\to U\times_SU$ 是单态射；
\item $U$ 是仿射的；并且
\item $j:R\to U\times_SU$ 分离且局部拟有限，
\end{enumerate}
则 $U/P$ 可由仿射概形 $\overline{U}$ 表示，商态射
$U\to\overline{U}$ 有限局部自由，且
$P=U\times_{\overline{U}}U$。此外，$R$ 是 $\overline{U}$ 上群胚概形
$(\overline{U},\overline{R},\overline{s},\overline{t},\overline{c})$
沿商态射 $U\to\overline{U}$ 的限制。
\end{lemma}

\begin{proof}
由条件 (1)、(2)、(3)、(4) 以及《群胚》章，命题
\ref{groupoids-proposition-finite-flat-equivalence}，
存在表示 $U/P$ 的仿射概形 $\overline{U}$，态射
$U\to\overline{U}$ 有限局部自由，且
$P=U\times_{\overline{U}}U$。此等同满足
$t|_P=\text{pr}_0$、$s|_P=\text{pr}_1$，而复合等于
$\text{pr}_{02}:U\times_{\overline{U}}U\times_{\overline{U}}U
\to U\times_{\overline{U}}U$。
有限局部自由态射的积仍有限局部自由（见《空间》章，引理
\ref{spaces-lemma-product-representable-transformations-property}，
以及《态射》章，引理 \ref{morphisms-lemma-base-change-finite-locally-free}
与 \ref{morphisms-lemma-composition-finite-locally-free}）。
为得到 $\overline{R}$，要沿有限局部自由态射
$U\times_SU\to\overline{U}\times_S\overline{U}$ 对概形 $R$ 作下降。
注意
$$
(U \times_S U)
\times_{(\overline{U} \times_S \overline{U})}
(U \times_S U)
=
P \times_S P
$$
，其中使用了上述结论。因此，给出
$R/U\times_SU/\overline{U}\times_S\overline{U}$ 的下降数据
（见《下降》章，定义 \ref{descent-definition-descent-datum}），
就是给出同构
$$
\varphi :
R \times_{(U \times_S U), t \times t} (P \times_S P)
\longrightarrow
(P \times_S P) \times_{s \times s, (U \times_S U)} R
$$
，它位于 $P\times_SP$ 上并满足余圈条件。
在 $T$ 值点上按如下规则定义 $\varphi$：
$$
\varphi : (r, (p, p')) \longmapsto ((p, p'), p^{-1} \circ r \circ p')
$$
其中复合取自在群胚范畴 $(U(T),R(T),s,t,c)$ 中的复合。
这是有意义的，因为 $(r,(p,p'))$ 要成为 $\varphi$ 的源的 $T$ 值点，
必须满足 $t(r)=t(p)$ 且 $s(r)=t(p')$。
注意，该映射是同构，其逆映射为
$((p,p'),r')\mapsto(p\circ r'\circ(p')^{-1},(p,p'))$。
为检验余圈条件，需要验证
$\varphi_{02} = \varphi_{12} \circ \varphi_{01}$
作为下列空间上的映射成立：
$$
(U \times_S U)
\times_{(\overline{U} \times_S \overline{U})} (U \times_S U)
\times_{(\overline{U} \times_S \overline{U})} (U \times_S U) =
(P \times_S P) \times_{s \times s, (U \times_S U), t \times t} (P \times_S P)
$$
直接计算可得
$$
\begin{matrix}
\varphi_{02} & (r, (p_1, p_1'), (p_2, p_2')) & \mapsto &
((p_1, p_1'), (p_2, p_2'),
(p_1 \circ p_2)^{-1} \circ r \circ (p_1' \circ p_2')) \\
\varphi_{01} & (r, (p_1, p_1'), (p_2, p_2')) & \mapsto &
((p_1, p_1'), p_1^{-1} \circ r \circ p_1', (p_2, p_2')) \\
\varphi_{12} & ((p_1, p_1'), r, (p_2, p_2')) & \mapsto &
((p_1, p_1'), (p_2, p_2'), p_2^{-1} \circ r \circ p_2')
\end{matrix}
$$
（记号含义显然），这便给出所需结论。
由 (5)，$j$ 分离且局部拟有限，因此可以应用《态射进阶》章，引理
\ref{more-morphisms-lemma-separated-locally-quasi-finite-morphisms-fppf-descend}，
得到概形 $\overline{R}\to\overline{U}\times_S\overline{U}$ 及同构
$$
R \to \overline{R} \times_{(\overline{U} \times_S \overline{U})} (U \times_S U)
$$
，它把下降数据 $\varphi$ 与
$\overline{R}\times_{(\overline{U}\times_S\overline{U})}(U\times_SU)$
上的典范下降数据等同；见《下降》章，定义
\ref{descent-definition-effective}。

\medskip\noindent
由于 $U\times_SU\to\overline{U}\times_S\overline{U}$ 有限局部自由，
作为其基变换，$R\to\overline{R}$ 也有限局部自由。
因此作为 $(\Sch/S)_{fppf}$ 上层的态射，
$R\to\overline{R}$ 是满射。
对给定的 $T$ 值点 $r,r'\in R(T)$，我们对 $\varphi$ 的选择表明：
二者在 $\overline{R}$ 中有相同的像，当且仅当对某些
$p,p'\in P(T)$ 有 $p^{-1}\circ r\circ p'$。因此
$\overline{R}$ 表示层
$$
T \longmapsto  \overline{R(T)} = P(T)\backslash R(T)/P(T)
$$
其中记号沿用引理之前的讨论。
于是可以完全按照“普通”群胚情形中的讨论，在
$(\overline{U}=U/P,\overline{R}=P\backslash R/P)$ 上定义群胚结构。
由此可知，$(U,R,s,t,c)$ 是该群胚结构沿态射
$U\to\overline{U}$ 的拉回，从而完成证明。
\end{proof}















\section{最终自举}
\label{bootstrap-section-final-bootstrap}

\noindent
下面的结果比先前各结果强得多。

\begin{theorem}
\label{bootstrap-theorem-final-bootstrap}
设 $S$ 为概形，$F:(\Sch/S)_{fppf}^{opp}\to\textit{Sets}$ 为函子。
下列任一条件均蕴含 $F$ 是代数空间：
\begin{enumerate}
\item $F=U/R$，其中 $(U,R,s,t,c)$ 是 $S$ 上代数空间中的群胚，
$s,t$ 平坦且局部有限呈示，而
$j=(t,s):R\to U\times_SU$ 是等价关系；
\item $F=U/R$，其中 $(U,R,s,t,c)$ 是 $S$ 上的群胚概形，
$s,t$ 平坦且局部有限呈示，而
$j=(t,s):R\to U\times_SU$ 是等价关系；
\item $F$ 是层，并且存在代数空间 $U$ 及态射 $U\to F$，
该态射可由代数空间表示、为满射、平坦且局部有限呈示；
\item $F$ 是层，并且存在概形 $U$ 及态射 $U\to F$，
该态射可由代数空间或概形表示、为满射、平坦且局部有限呈示；
\item $F$ 是层，$\Delta_F$ 可由代数空间表示，并且存在代数空间 $U$
及满、平坦、局部有限呈示的态射 $U\to F$；或者
\item $F$ 是层，$\Delta_F$ 可表，并且存在概形 $U$
及满、平坦、局部有限呈示的态射 $U\to F$。
\end{enumerate}
\end{theorem}

\begin{proof}
先作两个显然的观察：(6) 是 (5) 的特例，(4) 是 (3) 的特例。
首先证明情形 (5) 和 (3) 可归约到情形 (1)。
由对角态射的自举引理 \ref{bootstrap-lemma-bootstrap-diagonal}，
(3) 蕴含 (5)。在情形 (5) 中，令 $R=U\times_FU$；
由假设，这是代数空间。此外，两个投影 $s,t:R\to U$
均为满射、平坦且局部有限呈示。
态射 $j:R\to U\times_SU$ 显然是等价关系。
由引理 \ref{bootstrap-lemma-surjective-flat-locally-finite-presentation}，
$U\to F$ 是层的满射。因此 $F=U/R$，从而归约到情形 (1)。

\medskip\noindent
下面证明 (1) 可归约到 (2)。
设 $(U,R,s,t,c)$ 是 $S$ 上代数空间中的群胚，
$s,t$ 平坦且局部有限呈示，而
$j=(t,s):R\to U\times_SU$ 是等价关系。
选择概形 $U'$ 及满平展态射 $U'\to U$。
令 $R'=R|_{U'}$ 为 $R$ 在 $U'$ 上的限制。
由《空间中的群胚》章，引理
\ref{spaces-groupoids-lemma-quotient-pre-equivalence-relation-restrict}，
$U/R=U'/R'$。由于 $s',t':R'\to U'$ 也平坦且局部有限呈示
（见《空间中的群胚进阶》章，引理
\ref{spaces-more-groupoids-lemma-restrict-preserves-type}），
可归约到 $U$ 为概形的情形。
因为 $j$ 是等价关系，所以 $j$ 是单态射。
因为 $s:R\to U$ 局部有限呈示，态射
$j:R\to U\times_SU$ 局部有限型；见《空间的态射》章，引理
\ref{spaces-morphisms-lemma-permanence-finite-type}。
由《空间的态射》章，引理
\ref{spaces-morphisms-lemma-monomorphism-loc-finite-type-loc-quasi-finite}，
$j$ 局部拟有限且分离。
因此，当 $U$ 为概形时，由《空间的态射》章，命题
\ref{spaces-morphisms-proposition-locally-quasi-finite-separated-over-scheme}，
$R$ 也是概形。于是归约到在情形 (2) 中证明定理。

\medskip\noindent
假设 $F=U/R$，其中 $(U,R,s,t,c)$ 是 $S$ 上的群胚概形，
$s,t$ 平坦且局部有限呈示，而
$j=(t,s):R\to U\times_SU$ 是等价关系。
由引理 \ref{bootstrap-lemma-slice-equivalence-relation}，
可归约到 $s,t$ 平坦、局部有限呈示且局部拟有限的情形。
设 $U=\bigcup_{i\in I}U_i$ 为仿射开覆盖
（为避免后面的集合论问题，指标集 $I$ 的基数 $\leq$ $U$ 的大小；
大多数读者可以放心忽略这条说明）。
令 $(U_i,R_i,s_i,t_i,c_i)$ 为 $R$ 在 $U_i$ 上的限制。
由于 $R_i$ 是 $R$ 的开子概形
$s^{-1}(U_i)\cap t^{-1}(U_i)$，而 $s_i,t_i$ 是 $s,t$ 在该开集上的限制，
显然 $s_i,t_i$ 仍然平坦、局部有限呈示且局部拟有限。
由引理 \ref{bootstrap-lemma-better-finding-opens}
（或较简单的《空间》章，引理 \ref{spaces-lemma-finding-opens}），
态射 $U_i/R_i\to U/R$ 可由开浸入表示。
因此，若能证明 $F_i=U_i/R_i$ 是代数空间，则由《空间》章，引理
\ref{spaces-lemma-coproduct-algebraic-spaces}，
$\coprod_{i\in I}F_i$ 是代数空间。
由于 $U=\bigcup U_i$ 是开覆盖，$\coprod F_i\to F$ 显然是满射。
于是由《空间》章，引理 \ref{spaces-lemma-glueing-algebraic-spaces}，
$U/R$ 是代数空间。
这样便归约到下述情形：$U$ 是仿射的，$s,t$ 平坦、局部有限呈示且局部拟有限，
而 $j$ 是等价关系。

\medskip\noindent
假设 $(U,R,s,t,c)$ 是 $S$ 上的群胚概形，$U$ 仿射，
$s,t$ 平坦、局部有限呈示且局部拟有限，而 $j$ 是等价关系。
选择 $u\in U$。对 $u\in U,R,s,t,c$ 应用《空间中的群胚进阶》章，
引理 \ref{spaces-more-groupoids-lemma-quasi-splitting-affine-scheme}。
得到仿射概形 $U'$、平展态射 $g:U'\to U$ 以及点
$u'\in U'$，其中 $\kappa(u)=\kappa(u')$，使得限制
$R'=R|_{U'}$ 在 $u'$ 上拟分裂。
由于 $g$ 平展，其像 $g(U')$ 是包含 $u$ 的开集。
因此，反复应用该引理，可以找到有限多个点
$u_i\in U$（$i=1,\ldots,n$）、仿射概形 $U'_i$、
平展态射 $g_i:U_i'\to U$ 以及点 $u'_i\in U'_i$，满足
$g(u'_i)=u_i$，并且：(a) 每个限制 $R'_i$ 在 $U'_i$ 中某点上拟分裂；
(b) $U=\bigcup_{i=1,\ldots,n}g_i(U'_i)$。
现在重复上一段论证的最后部分。
由引理 \ref{bootstrap-lemma-better-finding-opens}
（或较简单的《空间》章，引理 \ref{spaces-lemma-finding-opens}），
态射 $U'_i/R'_i\to U/R$ 可由开浸入表示。
若能证明 $F_i=U'_i/R'_i$ 是代数空间，则由《空间》章，引理
\ref{spaces-lemma-coproduct-algebraic-spaces}，
$\coprod_{i\in I}F_i$ 是代数空间。
由于 $\{g_i:U'_i\to U\}$ 是平展覆盖，$\coprod F_i\to F$ 显然是满射。
于是由《空间》章，引理 \ref{spaces-lemma-glueing-algebraic-spaces}，
$U/R$ 是代数空间。
这样便归约到下述情形：$U$ 仿射，$s,t$ 平坦、局部有限呈示且局部拟有限，
$j$ 是等价关系，并且对某个 $u\in U$，$R$ 在 $u$ 上拟分裂。

\medskip\noindent
假设 $(U,R,s,t,c)$ 是 $S$ 上的群胚概形，$U$ 仿射，$u\in U$，
$s,t$ 平坦、局部有限呈示且局部拟有限，
$j=(t,s):R\to U\times_SU$ 是等价关系，并且 $R$ 在 $u$ 上拟分裂。
设 $P\subset R$ 为 $R$ 在 $u$ 上的一个拟分裂结构。
由引理 \ref{bootstrap-lemma-divide-subgroupoid}，
$(U,R,s,t,c)$ 是群胚
$(\overline{U},\overline{R},\overline{s},\overline{t},\overline{c})$
沿满的有限局部自由态射 $U\to\overline{U}$ 的限制，并且
$P=U\times_{\overline{U}}U$。注意，$s$ 有分解
$$
R = U \times_{\overline{U}, \overline{t}} \overline{R}
\times_{\overline{s}, \overline{U}} U
\xrightarrow{\text{pr}_{23}}
\overline{R} \times_{\overline{s}, \overline{U}} U
\xrightarrow{\text{pr}_2} U
$$
态射 $\text{pr}_2$ 是 $\overline{s}$ 的基变换，而
$\text{pr}_{23}$ 是满的有限局部自由态射
$U\to\overline{U}$ 的基变换。
由于 $s$ 平坦、局部有限呈示且局部拟有限，而 $\text{pr}_{23}$
满且有限局部自由（它是此类态射的基变换），由《下降》章，引理
\ref{descent-lemma-flat-fpqc-local-source}、
\ref{descent-lemma-locally-finite-presentation-fppf-local-source}
以及《态射》章，引理 \ref{morphisms-lemma-quasi-finite-local-source}，
可知 $\text{pr}_2$ 平坦、局部有限呈示且局部拟有限。
由于 $\text{pr}_2$ 是 $\overline{s}$ 沿 $U\to\overline{U}$ 的基变换，
而 $\{U\to\overline{U}\}$ 是 fppf 覆盖，
由《下降》章，引理 \ref{descent-lemma-descending-property-flat}、
\ref{descent-lemma-descending-property-locally-finite-presentation}
及 \ref{descent-lemma-descending-property-quasi-finite}，
可知 $\overline{s}$ 平坦、局部有限呈示且局部拟有限。
$\overline{t}$ 亦然。考虑交换图
$$
\xymatrix{
U \times_{\overline{U}} U \ar@{=}[r] \ar[rd] & P \ar[r] \ar[d] & R \ar[d] \\
& \overline{U} \ar[r]^{\overline{e}} & \overline{R}
}
$$
关于限制的一般事实表明，外侧四个角构成笛卡儿图。
由等式可知内方块也是笛卡儿的。
由于 $P$ 在 $R$ 中开（由拟分裂结构的定义），
由《下降》章，引理 \ref{descent-lemma-descending-property-open-immersion}，
$\overline{e}$ 是开浸入。
应用《群胚》章，引理
\ref{groupoids-lemma-quotient-pre-equivalence-relation-restrict}，
得到 $U/R=\overline{U}/\overline{R}$。
所以已经归约到如下情形：$(U,R,s,t,c)$ 是 $S$ 上的群胚概形，
$U$ 仿射，$u\in U$，$s,t$ 平坦、局部有限呈示且局部拟有限，
$j=(t,s):R\to U\times_SU$ 是等价关系，而 $e:U\to R$ 是开浸入！

\medskip\noindent
当然，若 $e$ 是开浸入，而 $s,t$ 平坦且局部有限呈示，
则态射 $t,s$ 都是平展的。
例如，应用《群胚进阶》章，引理
\ref{more-groupoids-lemma-sheaf-differentials} 可知
$\Omega_{R/U}=0$，进而 $s,t:R\to U$ 是 G-非分歧的
（见《态射》章，引理 \ref{morphisms-lemma-unramified-omega-zero}），
再进而 $s,t$ 平展（见《态射》章，引理
\ref{morphisms-lemma-flat-unramified-etale}）。
最后，若 $s,t$ 平展，则由《空间》章，定理
\ref{spaces-theorem-presentation}，$U/R$ 是代数空间。
\end{proof}





\section{应用}
\label{bootstrap-section-applications}

\noindent
第一个应用给出如下基本事实：
$$
\fbox{在 fppf 局部为代数空间的层本身是代数空间。}
$$
这就是下面引理的内容。
注意，假设 (2) 等价于 $F|_{(\Sch/S_i)_{fppf}}$ 是代数空间这一条件，
见《空间》章，引理 \ref{spaces-lemma-rephrase}。
假设 (3) 是集合论条件；不关心集合论问题的读者可以忽略它。

\begin{lemma}
\label{bootstrap-lemma-locally-algebraic-space}
\begin{slogan}
代数空间的定义在 fppf 意义下是局部的。
\end{slogan}
设 $S$ 为概形，$F:(\Sch/S)_{fppf}^{opp}\to\textit{Sets}$ 为函子，
$\{S_i\to S\}_{i\in I}$ 为 $(\Sch/S)_{fppf}$ 的覆盖。假设
\begin{enumerate}
\item $F$ 是层；
\item 每个 $F_i=h_{S_i}\times F$ 都是代数空间；并且
\item $\coprod_{i\in I}F_i$ 是代数空间（见《空间》章，引理
\ref{spaces-lemma-coproduct-algebraic-spaces}）。
\end{enumerate}
则 $F$ 是代数空间。
\end{lemma}

\begin{proof}
考虑态射 $\coprod F_i\to F$。它是 $\coprod S_i\to S$
沿 $F\to S$ 的基变换。因此，由 fppf 覆盖的定义及引理
\ref{bootstrap-lemma-base-change-transformation-property}，
该态射可表、局部有限呈示、平坦且为满射。
于是应用定理 \ref{bootstrap-theorem-final-bootstrap} 可知 $F$ 是代数空间。
\end{proof}

\noindent
下面是引理 \ref{bootstrap-lemma-locally-algebraic-space} 的一个特例，
在其中不必担心集合论问题。

\begin{lemma}
\label{bootstrap-lemma-locally-algebraic-space-finite-type}
设 $S$ 为概形，$F:(\Sch/S)_{fppf}^{opp}\to\textit{Sets}$ 为函子，
$\{S_i\to S\}_{i\in I}$ 为 $(\Sch/S)_{fppf}$ 的覆盖。假设
\begin{enumerate}
\item $F$ 是层；
\item 每个 $F_i=h_{S_i}\times F$ 都是代数空间；并且
\item 态射 $F_i\to S_i$ 为有限型。
\end{enumerate}
则 $F$ 是代数空间。
\end{lemma}

\begin{proof}
将使用上面的引理 \ref{bootstrap-lemma-locally-algebraic-space}。
为此，要利用 $F_i$ 在 $S_i$ 上有限型这一假设，证明该引理中的集合论条件成立
（可能先稍微加细给定的 $S$ 的覆盖）。
建议读者略过证明的其余部分。

\medskip\noindent
若 $S'_i\to S_i$ 是概形态射，则
$$
h_{S'_i} \times F =
h_{S'_i} \times_{h_{S_i}} h_{S_i} \times F =
h_{S'_i} \times_{h_{S_i}} F_i
$$
是 $S'_i$ 上有限型的代数空间；见《空间》章，引理
\ref{spaces-lemma-fibre-product-spaces} 以及《空间的态射》章，引理
\ref{spaces-morphisms-lemma-base-change-finite-type}。
因此可以加细给定的覆盖。此后可假定：
(a) 每个 $S_i$ 都仿射；(b) $I$ 的基数至多为 $S$ 的点集的基数。
（因为要覆盖整个 $S$，只须每个点都属于某个 $S_i\to S$ 的像。）

\medskip\noindent
由于每个 $S_i$ 都仿射，而每个 $F_i$ 在 $S_i$ 上有限型，
$F_i$ 是拟紧的。于是由《空间的性质》章，引理
\ref{spaces-properties-lemma-quasi-compact-affine-cover}，
可以找到仿射的 $U_i\in\Ob((\Sch/S)_{fppf})$ 及满平展态射
$U_i\to F_i$。由 $F_i\to S_i$ 局部有限型可知
$U_i\to S_i$ 局部有限型，特别地 $U_i\to S$ 局部有限型。
由《集合》章，引理 \ref{sets-lemma-bound-finite-type}，
$\text{size}(U_i)\leq\text{size}(S)$。
又有 $|I|\leq\text{size}(S)$，故由《集合》章，引理
\ref{sets-lemma-bound-size} 及 $\Sch$ 的构造，
$\coprod_{i\in I}U_i$ 同构于 $(\Sch/S)_{fppf}$ 的一个对象。
再由《空间》章，引理 \ref{spaces-lemma-coproduct-algebraic-spaces}，
$\coprod F_i$ 是代数空间，证毕。
\end{proof}

\noindent
第二个应用给出
$$
\fbox{代数空间的任意 fppf 下降数据都是有效的。}
$$
撇开集合论上的困难，该命题成立；下面给出一个示例性结果。

\begin{lemma}
\label{bootstrap-lemma-descend-algebraic-space}
\begin{slogan}
代数空间的 fppf 下降数据是有效的。
\end{slogan}
设 $S$ 为概形，$\{X_i\to X\}_{i\in I}$ 为 $S$ 上代数空间的 fppf 覆盖。
\begin{enumerate}
\item 若 $I$ 可数\footnote{不关心集合论问题的读者可以忽略可数性限制。
若能控制被下降的代数空间的大小，也可以允许更大的指标集；例如见引理
\ref{bootstrap-lemma-locally-algebraic-space-finite-type}。}，
则相对于 $\{X_i\to X\}$ 的任意代数空间下降数据都是有效的。
\item 相对于 $\{X_i\to X\}_{i\in I}$ 的任意下降数据
$(Y_i,\varphi_{ij})$（《空间上的下降》章，定义
\ref{spaces-descent-definition-descent-datum-for-family-of-morphisms}），
若 $Y_i\to X_i$ 为有限型，则该下降数据有效。
\end{enumerate}
\end{lemma}

\begin{proof}
证明 (1)。由《空间上的下降》章，引理
\ref{spaces-descent-lemma-descent-data-sheaves}，
问题转化为如下命题：若带有态射 $F\to X$ 的 fppf 层 $F$
满足每个 $F\times_XX_i$ 都是代数空间，则 $F$ 是代数空间。
对 $I$ 基数的限制蕴含，以 $I$ 为指标的代数空间余积仍是代数空间；
见《空间》章，引理 \ref{spaces-lemma-coproduct-algebraic-spaces}
及《集合》章，引理 \ref{sets-lemma-what-is-in-it}。
态射
$$
\coprod F \times_X X_i \longrightarrow F
$$
可由代数空间表示（它是 $\coprod X_i\to X$ 的基变换，见引理
\ref{bootstrap-lemma-base-change-transformation}），并且为满射、平坦且局部有限呈示
（同样因为它是 $\coprod X_i\to X$ 的基变换，见引理
\ref{bootstrap-lemma-base-change-transformation-property}）。
所以 (1) 由定理 \ref{bootstrap-theorem-final-bootstrap} 得出。

\medskip\noindent
证明 (2)。首先应用《空间上的下降》章，引理
\ref{spaces-descent-lemma-descent-data-sheaves}，
得到带有态射 $F\to X$ 的 fppf 层 $F$，使得对所有 $i\in I$ 都有
$F\times_XX_i=Y_i$。目标是证明 $F$ 是代数空间。
选择概形 $U$ 及满平展态射 $U\to X$。
作为 $U\to X$ 的基变换，
$F'=U\times_XF\to F$ 可表、为满射且平展。
由定理 \ref{bootstrap-theorem-final-bootstrap}，只须证明
$F'=U\times_XF$ 是代数空间。
可以选择 fppf 覆盖 $\{U_j\to U\}_{j\in J}$，
其中 $U_j$ 是概形，并且该覆盖加细 fppf 覆盖
$\{X_i\times_XU\to U\}_{i\in I}$；见《空间上的拓扑》章，引理
\ref{spaces-topologies-lemma-refine-fppf-schemes}。
于是得到映射 $a:J\to I$，并且对每个 $j$，得到 $X$ 上的态射
$U_j\to X_{a(j)}$。
此时 $U_j\times_UF'=U_j\times_{X_{a(j)}}Y_{a(j)}$
在 $U_j$ 上有限型。因此由引理
\ref{bootstrap-lemma-locally-algebraic-space-finite-type}，$F'$ 是代数空间。
\end{proof}

\noindent
下面给出另一类应用。

\begin{lemma}
\label{bootstrap-lemma-representable-by-spaces-cover}
设 $S$ 为概形。设 $a:F\to G$ 与 $b:G\to H$ 是函子
$(\Sch/S)_{fppf}^{opp}\to\textit{Sets}$ 之间的态射。假设
\begin{enumerate}
\item $F,G,H$ 都是层；
\item $a:F\to G$ 可由代数空间表示、平坦、局部有限呈示且为满射；并且
\item $b\circ a:F\to H$ 可由代数空间表示。
\end{enumerate}
则 $b$ 可由代数空间表示。
\end{lemma}

\begin{proof}
设 $U$ 为 $S$ 上的概形，且 $\xi\in H(U)$。需要证明
$U\times_{\xi,H}G$ 是代数空间。
另一方面，已知 $U\times_{\xi,H}F$ 是代数空间，并且作为态射 $a$ 的基变换，
$U\times_{\xi,H}F\to U\times_{\xi,H}G$ 可由代数空间表示、
平坦、局部有限呈示且为满射（见引理
\ref{bootstrap-lemma-base-change-transformation-property}）。
结论遂由定理 \ref{bootstrap-theorem-final-bootstrap} 得出。
\end{proof}

\begin{lemma}
\label{bootstrap-lemma-quotient-stack-isom}
设 $B\to S$ 与 $(U,R,s,t,c)$ 如《空间中的群胚》章，定义
\ref{spaces-groupoids-definition-quotient-stack} (1) 所述。
对任意 $S$ 上概形 $T$ 以及 $[U/R]$ 在 $T$ 上的对象 $x,y$，
$(\Sch/T)_{fppf}$ 上的层 $\mathit{Isom}(x,y)$ 是代数空间。
\end{lemma}

\begin{proof}
由《空间中的群胚》章，引理
\ref{spaces-groupoids-lemma-quotient-stack-isom}，
存在 fppf 覆盖 $\{T_i\to T\}_{i\in I}$，使得对每个 $i$，
$\mathit{Isom}(x,y)|_{(\Sch/T_i)_{fppf}}$ 都是代数空间。
由《空间》章，引理 \ref{spaces-lemma-rephrase}，
这意味着每个 $F_i=h_{S_i}\times\mathit{Isom}(x,y)$ 都是代数空间。
所以，为证明本引理，只须检验上面引理
\ref{bootstrap-lemma-locally-algebraic-space} 中的集合论条件，即
$\coprod F_i$ 是代数空间。为此使用《空间》章，引理
\ref{spaces-lemma-coproduct-algebraic-spaces}；它要求证明
$I$ 与各 $F_i$ 都不会“过大”。建议读者略过证明的其余部分。

\medskip\noindent
选择 $U'\in\Ob(\Sch/S)_{fppf}$ 及满平展态射 $U'\to U$。
令 $R'$ 为 $R$ 在 $U'$ 上的限制。
由于 $[U/R]=[U'/R']$，以 $U'$ 替换 $U$ 后，可以假定 $U$ 是概形。
（作这一步是为了使下面的纤维积位于一个概形上。）

\medskip\noindent
注意，加细覆盖 $\{T_i\to T\}$ 后，每个 $F_i$ 仍是代数空间。
因此可假定每个 $T_i$ 都仿射。
由于 $T_i\to T$ 局部有限呈示，由《集合》章，引理
\ref{sets-lemma-bound-finite-type}，
$\text{size}(T_i)\leq\text{size}(T)$。
还可假定指标集 $I$ 的基数至多为 $T$ 的点集的基数，
因为要得到覆盖，只须检验 $T$ 的每个点都属于某个像。
故 $|I|\leq\text{size}(T)$。
选择 $W\in\Ob((\Sch/S)_{fppf})$ 及满平展态射 $W\to R$。
注意，在《空间中的群胚》章，引理
\ref{spaces-groupoids-lemma-quotient-stack-isom} 的证明中，
我们已证明，对某些 $x_i,y_i:T_i\to U$，$F_i$ 可由
$T_i\times_{(y_i,x_i),U\times_BU}R$ 表示。
因此
$V_i=T_i\times_{(y_i,x_i),U\times_BU}W$ 是概形，
并带有平展满射 $V_i\to F_i$。
由《集合》章，引理 \ref{sets-lemma-bound-size-fibre-product}，
$$
\text{size}(V_i) \leq \max\{\text{size}(T_i), \text{size}(W)\}
\leq \max\{\text{size}(T), \text{size}(W)\}
$$
于是由《集合》章，引理 \ref{sets-lemma-bound-size}，
$$
\text{size}(\coprod\nolimits_{i \in I} V_i)
\leq \max\{|I|, \text{size}(T), \text{size}(W)\}.
$$
因此由 $\Sch$ 的构造，$\coprod_{i\in I}V_i$
同构于 $(\Sch/S)_{fppf}$ 中的一个对象 $V$。
这验证了《空间》章，引理
\ref{spaces-lemma-coproduct-algebraic-spaces} 的假设，证毕。
\end{proof}

\begin{lemma}
\label{bootstrap-lemma-covering-quotient}
设 $S$ 为概形。考虑形如 $F=U/R$ 的代数空间 $F$，
其中 $(U,R,s,t,c)$ 是 $S$ 上代数空间中的群胚，
$s,t$ 平坦且局部有限呈示，而
$j=(t,s):R\to U\times_SU$ 是等价关系。
则 $U\to F$ 为满射、平坦且局部有限呈示。
\end{lemma}

\begin{proof}
这几乎是显然的，但还不完全是。
由《空间中的群胚》章，引理
\ref{spaces-groupoids-lemma-quotient-pre-equivalence}
以及 $j$ 是单态射这一事实，可知 $R=U\times_FU$。
选择概形 $W$ 及满平展态射 $W\to F$。
由于 $U\to F$ 是层的满射，可以找到 fppf 覆盖
$\{W_i\to W\}$ 及态射 $W_i\to U$，提升各态射 $W_i\to F$。
于是
$$
W_i \times_F U = W_i \times_U U \times_F U = W_i \times_{U, t} R
$$
而投影 $W_i\times_FU\to W_i$ 是 $t:R\to U$ 的基变换，
因而平坦且局部有限呈示；见《空间的态射》章，引理
\ref{spaces-morphisms-lemma-base-change-flat} 及
\ref{spaces-morphisms-lemma-base-change-finite-presentation}。
因此由《空间上的下降》章，引理
\ref{spaces-descent-lemma-descending-property-flat} 及
\ref{spaces-descent-lemma-descending-property-locally-finite-presentation}，
$U\to F$ 平坦且局部有限呈示。
由《空间》章，注 \ref{spaces-remark-warning}，它还是满射。
\end{proof}

\begin{lemma}
\label{bootstrap-lemma-quotient-free-action}
设 $S$ 为概形，$X\to B$ 为 $S$ 上代数空间的态射，
$G$ 为 $B$ 上的群代数空间，并设
$a:G\times_BX\to X$ 为 $G$ 在 $X$ 上相对于 $B$ 的作用。若
\begin{enumerate}
\item $a$ 是自由作用；并且
\item $G\to B$ 平坦且局部有限呈示，
\end{enumerate}
则 $X/G$（见《空间中的群胚》章，定义
\ref{spaces-groupoids-definition-quotient-sheaf}）是代数空间，
态射 $X\to X/G$ 为满射、平坦且局部有限呈示，
并且 $X$ 是 $X/G$ 上的 fppf $G$-挠子。
\end{lemma}

\begin{proof}
由定理 \ref{bootstrap-theorem-final-bootstrap} 及各定义，立即可知
$X/G$ 是代数空间。事实上，$X/G=X/R$，其中 $R=G\times_BX$。
态射 $s,t:G\times_BX\to X$ 平坦且局部有限呈示
（对 $s$，这是因为它是 $G\to B$ 的基变换；利用逆元及对称性，
对 $t$ 也成立）。由于作用自由，由《空间中的群胚》章，引理
\ref{spaces-groupoids-lemma-free-action}，
态射 $j:G\times_BX\to X\times_BX$ 是单态射。
由引理 \ref{bootstrap-lemma-covering-quotient}，
$X\to X/G$ 为满射、平坦且局部有限呈示。
要证明 $X\to X/G$ 是 fppf $G$-挠子
（《空间中的群胚》章，定义
\ref{spaces-groupoids-definition-principal-homogeneous-space}），
需要证明 $G\times_SX\to X\times_{X/G}X$ 是同构，
并且 $X\to X/G$ 在 fppf 局部具有截面。
第二点由已经证明的 $X\to X/G$ 的性质显然成立。
由于作用自由，态射 $G\times_SX\to X\times_{X/G}X$
作为 fppf 层的态射是单射。
最后，由《空间中的群胚》章，引理
\ref{spaces-groupoids-lemma-quotient-pre-equivalence}，
该态射作为层的态射也是满射，从而完成证明。
\end{proof}

\begin{lemma}
\label{bootstrap-lemma-descent-torsor}
设 $\{S_i\to S\}_{i\in I}$ 为 $(\Sch/S)_{fppf}$ 的覆盖，
$G$ 为 $S$ 上的群代数空间，并以 $G_i=G_{S_i}$ 表示各基变换。
假设给定
\begin{enumerate}
\item 对每个 $i\in I$，一个 $S_i$ 上的 fppf $G_i$-挠子 $X_i$；以及
\item 对每个 $i,j\in I$，一个 $G_{S_i\times_SS_j}$-等变同构
$\varphi_{ij}:X_i\times_SS_j\to S_i\times_SX_j$，
它在每个 $S_i\times_SS_j\times_SS_j$ 上满足余圈条件。
\end{enumerate}
则存在 $S$ 上的 fppf $G$-挠子 $X$，其到 $S_i$ 的基变换同构于 $X_i$，
并恢复下降数据 $\varphi_{ij}$。
\end{lemma}

\begin{proof}
可以把 $X_i$ 看作 $(\Sch/S_i)_{fppf}$ 上的层；见《空间》章，第
\ref{spaces-section-change-base-scheme} 节。
由《位点》章，第 \ref{sites-section-glueing-sheaves} 节，
下降数据 $(X_i,\varphi_{ij})$ 是有效的，其含义是：
存在 $(\Sch/S)_{fppf}$ 上唯一的层 $X$，限制回
$(\Sch/S_i)_{fppf}$ 后恢复代数空间 $X_i$。
因此 $X_i=h_{S_i}\times X$。
由引理 \ref{bootstrap-lemma-locally-algebraic-space}，$X$ 是代数空间，
只须验证 $\coprod X_i$ 是代数空间；这将在证明末尾完成。
由《位点》章，引理 \ref{sites-lemma-mapping-property-glue}
中的范畴等价，各作用态射
$G_i\times_{S_i}X_i\to X_i$ 可黏合成态射
$a:G\times_SX\to X$。
现在需要证明 $a$ 是作用，$X$ 是伪挠子，并且在 fppf 局部平凡
（见《空间中的群胚》章，定义
\ref{spaces-groupoids-definition-principal-homogeneous-space}）。
这些性质都可以在 fppf 局部检验，因而由各作用
$G_i\times_{S_i}X_i\to X_i$ 的相应性质得出。本引理因此成立。

\medskip\noindent
建议读者略过证明的其余部分，它完全是集合论的。
选择 $(\Sch/S)_{fppf}$ 的覆盖
$\{S_{ij}\to S_j\}_{j\in J_i}$，使 $G_i$-挠子 $X_i$ 平凡化
（这可由假设以及《拓扑》章，引理
\ref{topologies-lemma-fppf-induced} (1) 实现）。
于是 $\{S_{ij}\to S\}_{i\in I,j\in J_i}$ 是
$(\Sch/S)_{fppf}$ 的覆盖，故可以假定每个 $X_i$ 都是平凡挠子！
当然还可进一步加细覆盖，所以可假定每个 $S_i$ 都仿射，并且指标集 $I$
的基数以 $S$ 的点集的基数为界。
选择 $U\in\Ob((\Sch/S)_{fppf})$ 及满平展态射 $U\to G$。
于是 $U_i=U\times_SS_i$ 带有到 $X_i\cong G_i$ 的平展满射。
由《集合》章，引理 \ref{sets-lemma-bound-size-fibre-product}，
$\text{size}(U_i)\leq\max\{\text{size}(U),\text{size}(S_i)\}$。
由《集合》章，引理 \ref{sets-lemma-bound-finite-type}，
$\text{size}(S_i)\leq\text{size}(S)$。
所以对所有 $i\in I$，
$\text{size}(U_i)\leq\max\{\text{size}(U),\text{size}(S)\}$。
结合上面得到的 $|I|$ 的界，再由《集合》章，引理
\ref{sets-lemma-bound-size}，
$\text{size}(\coprod U_i)\leq
\max\{\text{size}(U),\text{size}(S)\}$。
因此可以应用《空间》章，引理
\ref{spaces-lemma-coproduct-algebraic-spaces}，得到
$\coprod X_i$ 是代数空间；这正是所要证明的。
\end{proof}




\section{平展拓扑中的代数空间}
\label{bootstrap-section-spaces-etale}

\noindent
设 $S$ 为概形。除了研究大 fppf 位点 $(\Sch/S)_{fppf}$ 上的层，
也可以研究大平展位点 $(\Sch/S)_\etale$ 上的层。
《代数空间》章，第 \ref{spaces-section-representable} 节及
第 \ref{spaces-section-representable-properties} 节的全部内容，
对 $(\Sch/S)_\etale$ 上的层仍然有意义。
因此，在平展拓扑中工作会得到第二种代数空间概念。
先验地说，它弱于《代数空间》章，定义
\ref{spaces-definition-algebraic-space} 中引入的概念，
因为 fppf 拓扑下的层当然也是平展拓扑下的层。
不过，下面的引理表明这两个概念等价。

\begin{lemma}
\label{bootstrap-lemma-spaces-etale}
以 $\Sch_\alpha$ 表示 $\Sch_{fppf}$ 与 $\Sch_\etale$
共同的底层范畴（见《拓扑》章，注
\ref{topologies-remark-choice-sites}）。
设 $S$ 为 $\Sch_\alpha$ 的对象，并设
$$
F : (\Sch_\alpha/S)^{opp} \longrightarrow \textit{Sets}
$$
为满足下列性质的预层：
\begin{enumerate}
\item $F$ 是平展拓扑下的层；
\item 对角态射 $\Delta:F\to F\times F$ 可表；并且
\item 存在 $U\in\Ob(\Sch_\alpha/S)$ 及满平展态射 $U\to F$。
\end{enumerate}
则 $F$ 是《代数空间》章，定义
\ref{spaces-definition-algebraic-space} 意义下的代数空间。
\end{lemma}

\begin{proof}
注意，本引理中的性质 (2)、(3) 以及《代数空间》章，定义
\ref{spaces-definition-algebraic-space} 中相应的性质 (2)、(3)
都与拓扑无关。这是因为这些性质只涉及预层的纤维积、预层态射、
可表函子态射的概念，以及这类态射为满射和平展的含义。
因此，只须证明具有性质 (2)、(3) 的平展层 $F$ 也是 fppf 层。

\medskip\noindent
为此，令 $R=U\times_FU$。由 (2)，预层 $R$ 可由概形表示；
由 (3)，投影 $R\to U$ 平展。因此
$j:R\to U\times_SU$ 是平展等价关系。
此外，$U\to F$ 把 $F$ 等同为平展拓扑下 $U$ 对 $R$ 的商：
(a) 若 $T\to F$ 为态射，则 $\{T\times_FU\to T\}$ 是平展覆盖，
所以 $U\to F$ 是平展拓扑下层的满射；
(b) 若 $a,b:T\to U$ 映到 $F$ 的同一截面，则
$(a,b):T\to R$，所以 $a,b$ 在平展拓扑下 $U$ 对 $R$ 的商中有相同的像。
接着，以 $U/R$ 表示 fppf 拓扑下的商层；由《空间》章，定理
\ref{spaces-theorem-presentation}，它是代数空间。
因此有态射（函子之间的态射）
$$
U \to F \to U/R.
$$
由上述《空间》章，定理 \ref{spaces-theorem-presentation}，
该复合可表、为满射且平展。
因此，对任意概形 $T$ 及态射 $T\to U/R$，纤维积
$V=T\times_{U/R}U$ 是在 $T$ 上满且平展的概形。
换言之，$\{V\to U\}$ 是平展覆盖。
这证明 $U\to U/R$ 作为平展拓扑下层的态射是满射，
进而 $F\to U/R$ 作为平展拓扑下层的态射是满射。
另一方面，再次由《空间》章，定理
\ref{spaces-theorem-presentation}，有 $R=U\times_{U/R}U$，
所以 $F\to U/R$ 作为预层态射是单射。
因此 $F\to U/R$ 是平展层的同构；见《位点》章，引理
\ref{sites-lemma-mono-epi-sheaves}。证明完成。
\end{proof}

\noindent
还有《空间》章，引理
\ref{spaces-lemma-etale-locally-representable-gives-space} 的对应结果。

\begin{lemma}
\label{bootstrap-lemma-spaces-etale-locally-representable}
以 $\Sch_\alpha$ 表示 $\Sch_{fppf}$ 与 $\Sch_\etale$
共同的底层范畴（见《拓扑》章，注
\ref{topologies-remark-choice-sites}）。
设 $S$ 为 $\Sch_\alpha$ 的对象，并设
$$
F : (\Sch_\alpha/S)^{opp} \longrightarrow \textit{Sets}
$$
为满足下列性质的预层：
\begin{enumerate}
\item $F$ 是平展拓扑下的层；
\item 存在 $S$ 上的代数空间 $U$ 及态射 $U\to F$，
该态射可由代数空间表示、为满射且平展。
\end{enumerate}
则 $F$ 是《代数空间》章，定义
\ref{spaces-definition-algebraic-space} 意义下的代数空间。
\end{lemma}

\begin{proof}
令 $R=U\times_FU$。由于假设 $U\to F$ 可由代数空间表示，
$R$ 是代数空间。由于假设 $U\to F$ 平展，投影
$s,t:R\to U$ 是代数空间的平展态射。
因为 $R=U\times_FU$，态射
$j=(t,s):R\to U\times_SU$ 是单态射和等价关系。
由定理 \ref{bootstrap-theorem-final-bootstrap}，fppf 商层
$F'=U/R$ 是代数空间。
由引理 \ref{bootstrap-lemma-covering-quotient}，
态射 $U\to F'$ 为满射、平坦且局部有限呈示。
由《空间中的群胚》章，引理
\ref{spaces-groupoids-lemma-quotient-pre-equivalence}，
$R\to U\times_{F'}U$ 作为 fppf 层的态射是满射；
又因 $j$ 是单态射，它是同构。
所以 $U\to F'$ 沿自身的基变换是平展的，
再由《空间上的下降》章，引理
\ref{spaces-descent-lemma-descending-property-etale}，
$U\to F'$ 平展。
因此 $U\to F'$ 作为平展层的态射是满射。
这意味着 $F'$ 等于平展拓扑下的商层 $U/R$（略去一项简单检验）。
于是得到典范分解 $U\to F'\to F$，且 $F'\to F$ 是层的单射。
另一方面，$U\to F$ 作为平展层的态射是满射，故 $F'\to F$ 也是满射。
所以 $F'=F$，证明完成。
\end{proof}

\noindent
事实上，只需有概形给出的光滑覆盖，并且只需假设对角态射可由代数空间表示。

\begin{lemma}
\label{bootstrap-lemma-spaces-etale-smooth-cover}
以 $\Sch_\alpha$ 表示 $\Sch_{fppf}$ 与 $\Sch_\etale$
共同的底层范畴（见《拓扑》章，注
\ref{topologies-remark-choice-sites}）。设 $S$ 为 $\Sch_\alpha$ 的对象。
$$
F : (\Sch_\alpha/S)^{opp} \longrightarrow \textit{Sets}
$$
为满足下列性质的预层：
\begin{enumerate}
\item $F$ 是平展拓扑下的层；
\item 对角态射 $\Delta:F\to F\times F$ 可由代数空间表示；并且
\item 存在 $U\in\Ob(\Sch_\alpha/S)$ 及满光滑态射 $U\to F$。
\end{enumerate}
则 $F$ 是《代数空间》章，定义
\ref{spaces-definition-algebraic-space} 意义下的代数空间。
\end{lemma}

\begin{proof}
证明仿照引理 \ref{bootstrap-lemma-spaces-etale} 的证明。
令 $R=U\times_FU$。由 (2)，预层 $R$ 是代数空间；
由 (3)，投影 $R\to U$ 光滑且为满射。
以 $(U,R,s,t,c)$ 表示等价关系 $j:R\to U\times_SU$
所关联的群胚（见《空间中的群胚》章，引理
\ref{spaces-groupoids-lemma-equivalence-groupoid}）。
由定理 \ref{bootstrap-theorem-final-bootstrap}，
$X=U/R$（fppf 拓扑下的商）是代数空间。
光滑拓扑与平展拓扑具有相同的层（由《态射进阶》章，引理
\ref{more-morphisms-lemma-etale-dominates-smooth}），
故态射 $U\to F$ 把 $F$ 等同为光滑拓扑下 $U$ 对 $R$ 的商
（细节从略）。因此有态射（函子之间的态射）
$$
U \to F \to X.
$$
由引理 \ref{bootstrap-lemma-covering-quotient}，
$U\to X$ 为满射、平坦且局部有限呈示。
由《空间中的群胚》章，引理
\ref{spaces-groupoids-lemma-quotient-pre-equivalence}
（以及 $j$ 是单态射这一事实），有 $R=U\times_XU$。
由于投影 $R\to U$ 光滑且为满射，而 $\{U\to X\}$ 是 fppf 覆盖，
由《空间上的下降》章，引理
\ref{spaces-descent-lemma-descending-property-smooth}，
$U\to X$ 光滑且为满射。
因此，对任意概形 $T$ 及态射 $T\to X$，
纤维积 $T\times_XU$ 是在 $T$ 上光滑且为满射的代数空间。
选择概形 $V$ 及满平展态射 $V\to T\times_XU$。
于是 $\{V\to T\}$ 是光滑覆盖，并且 $V\to T\to X$
可提升为态射 $V\to U$。
这证明 $U\to X$ 作为光滑拓扑下层的态射是满射，
进而 $F\to X$ 作为光滑拓扑下层的态射是满射。
另一方面，由于 $R=U\times_XU$，$F\to X$ 作为预层态射是单射。
因此 $F\to X$ 是光滑（即平展）层的同构；见《位点》章，引理
\ref{sites-lemma-mono-epi-sheaves}。证明完成。
\end{proof}

\noindent
最后，给出《空间》章，引理
\ref{spaces-lemma-etale-locally-representable-gives-space} 的对应结果，
其中用光滑态射覆盖空间。

\begin{lemma}
\label{bootstrap-lemma-spaces-smooth-locally-representable}
以 $\Sch_\alpha$ 表示 $\Sch_{fppf}$ 与 $\Sch_\etale$
共同的底层范畴（见《拓扑》章，注
\ref{topologies-remark-choice-sites}）。
设 $S$ 为 $\Sch_\alpha$ 的对象，并设
$$
F : (\Sch_\alpha/S)^{opp} \longrightarrow \textit{Sets}
$$
为满足下列性质的预层：
\begin{enumerate}
\item $F$ 是平展拓扑下的层；
\item 存在 $S$ 上的代数空间 $U$ 及态射 $U\to F$，
该态射可由代数空间表示、为满射且光滑。
\end{enumerate}
则 $F$ 是《代数空间》章，定义
\ref{spaces-definition-algebraic-space} 意义下的代数空间。
\end{lemma}

\begin{proof}
证明与引理 \ref{bootstrap-lemma-spaces-etale-locally-representable} 的证明相同。
令 $R=U\times_FU$。由于假设 $U\to F$ 可由代数空间表示，
$R$ 是代数空间。由于假设 $U\to F$ 光滑，投影
$s,t:R\to U$ 是代数空间的光滑态射。
因为 $R=U\times_FU$，态射
$j=(t,s):R\to U\times_SU$ 是单态射和等价关系。
由定理 \ref{bootstrap-theorem-final-bootstrap}，
fppf 商层 $F'=U/R$ 是代数空间。
由引理 \ref{bootstrap-lemma-covering-quotient}，
态射 $U\to F'$ 为满射、平坦且局部有限呈示。
由《空间中的群胚》章，引理
\ref{spaces-groupoids-lemma-quotient-pre-equivalence}，
$R\to U\times_{F'}U$ 作为 fppf 层的态射是满射；
又因 $j$ 是单态射，它是同构。
所以 $U\to F'$ 沿自身的基变换是光滑的，
再由《空间上的下降》章，引理
\ref{spaces-descent-lemma-descending-property-smooth}，
$U\to F'$ 光滑。
因此 $U\to F'$ 作为平展层的态射是满射
（由《态射进阶》章，引理
\ref{more-morphisms-lemma-etale-dominates-smooth}，
光滑拓扑与平展拓扑具有相同的层）。
这意味着 $F'$ 等于平展拓扑下的商层 $U/R$（略去一项简单检验）。
于是得到典范分解 $U\to F'\to F$，且 $F'\to F$ 是层的单射。
另一方面，$U\to F$ 作为平展层的态射是满射
（因为光滑拓扑与平展拓扑具有相同的层），故 $F'\to F$ 也是满射。
所以 $F'=F$，证明完成。
\end{proof}
% Generated mechanically from the frozen spaces-pushouts.tex target.
% Do not edit this generated file; edit the bound locale source instead.

% OK, start here.
%

\setcounter{chapter}{80}
\chapter{代数空间的推出}



\phantomsection
\label{spaces-pushouts-section-phantom}


\section{引言}
\label{spaces-pushouts-section-introduction}

\noindent
本章旨在讨论代数空间范畴中的推出。可以在不同的假设下进行这种构造。
文献 \cite{Temkin-Tyomkin} 给出了一种相当一般的推出构造：
其中一个态射是仿射的，另一个态射是闭浸入。
我们在第 \ref{spaces-pushouts-section-pushouts} 节讨论它的一个特殊情形，
即假设其中一个态射是仿射的，另一个态射是加厚；
这种情形在形变理论中经常出现。

\medskip\noindent
在第 \ref{spaces-pushouts-section-formal-glueing} 节和第
\ref{spaces-pushouts-section-formal-glueing-spaces} 节中，我们讨论图表
$$
\xymatrix{
f^{-1}(X \setminus Z) \ar[r] \ar[d] & Y \ar[d]^f \\
X \setminus Z \ar[r] & X
}
$$
其中 $f$ 是代数空间的拟紧且拟分离态射，
$Z \to X$ 是有限呈示的闭浸入，映射 $f^{-1}(Z) \to Z$ 是同构，
并且 $f$ 沿 $f^{-1}(Z)$ 平坦。在这种情形下，我们把
$X \setminus Z$ 和 $Y$ 上的拟凝聚模粘合为 $X$ 上的拟凝聚模
（第 \ref{spaces-pushouts-section-formal-glueing} 节），并把
$X \setminus Z$ 和 $Y$ 上的代数空间粘合为 $X$ 上的代数空间
（第 \ref{spaces-pushouts-section-formal-glueing-spaces} 节）。

\medskip\noindent
在第 \ref{spaces-pushouts-section-coequalizer-glue} 节中，我们讨论诺特代数空间的
固有双有理态射如何在某种意义下给出代数空间中的余等化子图表。

\medskip\noindent
在第 \ref{spaces-pushouts-section-compactifications} 节中，我们利用第
\ref{spaces-pushouts-section-elementary-dsitnguished-squares} 节对初等特异方形的构造，
证明代数空间情形下的 Nagata 紧化定理。






\section{约定}
\label{spaces-pushouts-section-conventions}

\noindent
我们始终假设所有概形都包含在一个大 fppf 位点 $\Sch_{fppf}$ 中。
此外，所考虑的每个环 $A$ 都满足：$\Spec(A)$（同构于）
这个大位点中的一个对象。

\medskip\noindent
设 $S$ 是概形，$X$ 是 $S$ 上的代数空间。
在本章及以下各章中，我们把 $X$ 与自身的乘积
（在 $S$ 上代数空间的范畴中）写成 $X \times_S X$，
而不写成 $X \times X$。




\section{代数空间的余极限}
\label{spaces-pushouts-section-pushouts-generalities}

\noindent
我们简要讨论代数空间的余极限。设 $S$ 是概形。
令 $\mathcal{I} \to (\Sch/S)_{fppf}$，$i \mapsto X_i$
为一个图表（见《范畴》第 \ref{categories-section-limits} 节）。
对每个 $i$，可以考虑小平展位点 $X_{i, \etale}$；
其对象是 $X_i$ 上平展的概形，见《空间的性质》第
\ref{spaces-properties-section-etale-site} 节。
对 $\mathcal{I}$ 的每个态射 $i \to j$，都有态射 $X_i \to X_j$，
从而有拉回函子 $X_{j, \etale} \to X_{i, \etale}$。
因此得到一个从 $\mathcal{I}^{opp}$ 到范畴的 $2$-范畴的伪函子。记
$$
\lim_i X_{i, \etale}
$$
为这个 $2$-极限（此处以后补入参引）。具体而言，这是什么意思？
这个极限的一个对象是 $\mathcal{I}$ 上的一族平展态射
$U_i \to X_i$，使得对 $\mathcal{I}$ 中的每个 $i \to j$，图表
$$
\xymatrix{
U_i \ar[r] \ar[d] & U_j \ar[d] \\
X_i \ar[r] & X_j
}
$$
是笛卡尔的。对象之间的态射按显然的方式定义。
设 $f_i : X_i \to T$ 是一族态射，并且对每个 $i \to j$，
复合 $X_i \to X_j \to T$ 等于 $f_i$。
于是得到函子 $T_\etale \to \lim X_{i, \etale}$。
有了这些记号，我们可以陈述下面的引理。

\begin{lemma}
\label{spaces-pushouts-lemma-colimit-agrees}
设 $S$ 是概形。令 $\mathcal{I} \to (\Sch/S)_{fppf}$，$i \mapsto X_i$
为如上的 $S$ 上概形图表。假设
\begin{enumerate}
\item 在概形范畴中存在 $X = \colim X_i$；
\item $\coprod X_i \to X$ 是满射；
\item 若 $U \to X$ 平展且 $U_i = X_i \times_X U$，则在概形范畴中
$U = \colim U_i$；以及
\item $\lim X_{i, \etale}$ 中每个满足 $U_i \to X_i$ 分离的对象
$(U_i \to X_i)$，都属于函子
$X_\etale \to \lim X_{i, \etale}$ 的本质像。
\end{enumerate}
则在 $S$ 上代数空间的范畴中也有 $X = \colim X_i$。
\end{lemma}

\begin{proof}
设 $Z$ 是 $S$ 上的代数空间。设 $f_i : X_i \to Z$ 是一族态射，
使得对每个 $i \to j$，复合 $X_i \to X_j \to Z$ 等于 $f_i$。
我们必须构造代数空间的态射 $f : X \to Z$，使得 $f_i$ 可由复合
$X_i \to X \to Z$ 得到。令 $W \to Z$ 是从一个概形到 $Z$ 的满平展态射。
可以假设 $W$ 是仿射概形的不交并，特别地，可以假设 $W \to Z$ 分离。
对每个 $i$，令 $U_i = W \times_{Z, f_i} X_i$，并把投影记作
$h_i : U_i \to W$。于是 $U_i \to X_i$ 构成
$\lim X_{i, \etale}$ 的一个对象，且 $U_i \to X_i$ 分离。
由假设 (4)，可以找到平展态射 $U \to X$ 以及（函子性的）同构
$U_i = X_i \times_X U$。由假设 (3)，存在态射 $h : U \to W$，
使复合 $U_i \to U \to W$ 都是 $h_i$。
令 $g : U \to Z$ 为 $h$ 与映射 $W \to Z$ 的复合。
为完成证明，必须说明 $g : U \to Z$ 可下降为态射 $X \to Z$。
为此，考虑态射
$(h, h) : U \times_X U \to W \times_S W$.
将它与 $U_i \times_{X_i} U_i \to U \times_X U$ 复合，所得态射
$(h_i, h_i)$ 通过 $W \times_Z W$ 分解。由 (3)，$U \times_X U$
是概形 $U_i \times_{X_i} U_i$ 的余极限，故 $(h,h)$ 通过
$W \times_Z W$ 分解。因此两个复合
$U \times_X U \to U \to W \to Z$ 相等。
由于每个 $U_i \to X_i$ 都是满射，再由假设 (2)，可知 $U \to X$ 是满射。
由于 $Z$ 是平展拓扑下的层，便得到 $g : U \to Z$ 按所需下降为
$f : X \to Z$。
\end{proof}

\noindent
可以在余锥上（fpqc）局部地检验一个余锥是否为余极限。

\begin{lemma}
\label{spaces-pushouts-lemma-pushout-fpqc-local}
设 $S$ 是概形，$B$ 是 $S$ 上的代数空间。
令 $\mathcal{I} \to (\Sch/S)_{fppf}$，$i \mapsto X_i$
为 $B$ 上代数空间的一个图表。令 $(X, X_i \to X)$ 为此图表在
$B$ 上代数空间范畴中的一个余锥
（《范畴》注 \ref{categories-remark-cones-and-cocones}）。
若存在一个 fpqc 覆盖 $\{U_a \to X\}_{a \in A}$，使得
\begin{enumerate}
\item 对所有 $a \in A$，在 $B$ 上代数空间范畴中有
$U_a = \colim X_i \times_X U_a$
；以及
\item 对所有 $a,b \in A$，在 $B$ 上代数空间范畴中有
$U_a \times_X U_b = \colim X_i \times_X U_a \times_X U_b$
，
\end{enumerate}
则在 $B$ 上代数空间的范畴中有 $X = \colim X_i$。
\end{lemma}

\begin{proof}
事实上，对 $B$ 上的代数空间 $Y$，一个 $B$ 上的态射 $X \to Y$
等价于一族态射 $U_a \to Y$，它们对所有 $a,b \in A$ 都在交叠
$U_a \times_X U_b$ 上相合；见《空间上的下降》引理
\ref{spaces-descent-lemma-fpqc-universal-effective-epimorphisms}。
\end{proof}

\noindent
下面给出引理 \ref{spaces-pushouts-lemma-colimit-agrees} 和
\ref{spaces-pushouts-lemma-pushout-fpqc-local} 的一个共同的部分推广；
特别地，它可以把余极限构造约化到全体代数空间范畴的一个子范畴中。

\medskip\noindent
设 $S$ 是概形，$B$ 是 $S$ 上的代数空间。令 $\mathcal{I}$ 为指标范畴，
并令 $i \mapsto X_i$ 为 $B$ 上代数空间范畴中的一个图表；
见《范畴》第 \ref{categories-section-limits} 节。
对每个 $i$，可以考虑小平展位点 $X_{i, spaces, \etale}$；
其对象是 $X_i$ 上平展的代数空间，见《空间的性质》第
\ref{spaces-properties-section-etale-site} 节。
对 $\mathcal{I}$ 的每个态射 $i \to j$，都有态射 $X_i \to X_j$，
从而有拉回函子
$X_{j, spaces, \etale} \to X_{i, spaces, \etale}$。
因此得到一个从 $\mathcal{I}^{opp}$ 到范畴的 $2$-范畴的伪函子。记
$$
\lim_i X_{i, spaces, \etale}
$$
为这个 $2$-极限（此处以后补入参引）。具体而言，这是什么意思？
这个极限的一个对象，是 $B$ 上代数空间的箭头范畴中的一个图表
$i \mapsto (U_i \to X_i)$，使得对 $\mathcal{I}$ 中每个 $i \to j$，图表
$$
\xymatrix{
U_i \ar[r] \ar[d] & U_j \ar[d] \\
X_i \ar[r] & X_j
}
$$
是笛卡尔的。对象之间的态射按显然的方式定义。
设 $f_i : X_i \to Z$ 是 $B$ 上代数空间的一族态射，
使得对每个 $i \to j$，复合 $X_i \to X_j \to Z$ 等于 $f_i$。
于是得到函子
$Z_{spaces, \etale} \to \lim X_{i, spaces, \etale}$。
有了这些记号，我们可以陈述下一个引理。

\begin{lemma}
\label{spaces-pushouts-lemma-colimit-check-etale-locally}
设 $S$ 是概形，$B$ 是 $S$ 上的代数空间。
令 $\mathcal{I} \to (\Sch/S)_{fppf}$，$i \mapsto X_i$
为 $B$ 上代数空间的一个图表。令 $(X, X_i \to X)$ 为此图表在
$B$ 上代数空间范畴中的一个余锥
（《范畴》注 \ref{categories-remark-cones-and-cocones}）。假设
\begin{enumerate}
\item 基变换函子
$X_{spaces, \etale} \to \lim X_{i, spaces, \etale}$,% P10-E0035
把 $U$ 送到 $U_i = X_i \times_X U$，并且它是一个等价；
\item 给定
\begin{enumerate}
\item $B'$ 在 $B$ 上仿射且平展；
\item $Z$ 是 $B'$ 上的仿射概形；
\item $U \to X \times_B B'$ 是代数空间的平展态射，且 $U$ 仿射；
\item $f_i : U_i \to Z$ 是图表
$i \mapsto U_i = U \times_X X_i$ 在 $B'$ 上的一个余锥；
\end{enumerate}
则存在唯一的 $B'$ 上态射 $f : U \to Z$，使得 $f_i$ 等于复合
$U_i \to U \to Z$。
\end{enumerate}
则在全体 $B$ 上代数空间的范畴中有 $X = \colim X_i$。
\end{lemma}

\begin{proof}
本段把问题约化到 $B$ 为仿射概形的情形。
令 $B' \to B$ 是代数空间的平展态射。注意，若用
$B'$、$X_i \times_B B'$、$X \times_B B'$ 分别替换 $B$、$X_i$、$X$，
则条件 (1) 和 (2) 保持不变。
令 $\{B_a \to B\}_{a \in A}$ 为一个平展覆盖，其中各 $B_a$ 仿射；
见《空间的性质》引理
\ref{spaces-properties-lemma-cover-by-union-affines}。
对 $a \in A$，分别以 $X_a$、$X_{a,i}$ 表示 $X$ 和该图表到 $B_a$ 的基变换。
对 $a,b \in A$，分别以 $X_{a,b}$、$X_{a,b,i}$ 表示 $X$ 和该图表到
$B_a \times_B B_b$ 的基变换。
由引理 \ref{spaces-pushouts-lemma-pushout-fpqc-local}，只需证明
$X_a = \colim X_{a,i}$ 和 $X_{a,b} = \colim X_{a,b,i}$。
这把问题约化到 $B=B_a$（仿射概形）或
$B=B_a \times_B B_b$（分离概形）的情形。再重复一次这个论证，
便可假设 $B$ 是仿射概形（这里用到了分离概形中两个仿射开集的交仍为仿射）。

\medskip\noindent
假设 $B$ 是仿射概形。令 $Z$ 为 $B$ 上的代数空间。我们必须证明
$$
\Mor_B(X, Z) \longrightarrow \lim \Mor_B(X_i, Z)
$$
是双射。

\medskip\noindent
证明单射性。令 $f,g : X \to Z$ 为态射，并且对所有 $i$，复合
$f_i,g_i : X_i \to Z$ 都相同。选取仿射概形 $Z'$ 和平展态射
$Z' \to Z$。由《空间的性质》引理
\ref{spaces-properties-lemma-cover-by-union-affines}，可知这样的仿射概形
能够覆盖 $Z$。令 $U = X \times_{f,Z} Z'$、$U'=X\times_{g,Z}Z'$，
并把投影记作 $p:U\to X$ 和 $p':U'\to X$。
由于对所有 $i$ 都有 $f_i=g_i$，可知
$$
U_i = X_i \times_{f_i, Z} Z' = X_i \times_{g_i, Z} Z' = U'_i
$$
且它们与转移态射相容。由 (1)，存在唯一的 $X$ 上代数空间同构
$\epsilon:U\to U'$，即满足 $p=p'\circ\epsilon$，并与上述等同相容。
选取一个平展覆盖 $\{h_a:U_a\to U\}$，其中各 $U_a$ 仿射。
由 (2)，有
$f\circ p\circ h_a=g\circ p'\circ\epsilon\circ h_a
=g\circ p\circ h_a$。
由于 $\{h_a:U_a\to U\}$ 是平展覆盖，得到 $f\circ p=g\circ p$。
按这种方式得到的所有态射 $p:U\to X$ 构成一个平展覆盖，故 $f=g$。

\medskip\noindent
证明满射性。令 $f_i:X_i\to Z$ 为证明首段所示箭头右端的一个元素。
只需找到一个平展覆盖 $\{U_c\to X\}_{c\in C}$，使得各族
$f_{c,i}\in\lim_i\Mor_B(X_i\times_XU_c,Z)$ 来自态射
$f_c:U_c\to Z$。事实上，由上面证明的唯一性，诸 $f_c$ 在
$U_c\times_XU_b$ 上相合，因而可下降为所需态射 $f:X\to Z$。
为找到这个覆盖，先选取一个平展覆盖
$\{g_a:Z_a\to Z\}_{a\in A}$，其中各 $Z_a$ 仿射。然后令
$U_{a,i}=X_i\times_{f_i,Z}Z_a$。由 (1)，对某个在 $X$ 上平展的
代数空间 $U_a$，有 $U_{a,i}=X_i\times_XU_a$。
再选取平展覆盖 $\{U_{a,b}\to U_a\}_{b\in B_a}$，
其中各 $U_{a,b}$ 仿射，并考虑态射
$$
U_{a, b, i} = X_i \times_X U_{a, b} \to
X_i \times_X U_a = X_i \times_{f_i, Z} Z_a \to Z_a
$$
由 (2)，得到与这些态射相容的态射 $f_{a,b}:U_{a,b}\to Z_a$。
令 $C=\coprod_{a\in A}B_a$；若 $c\in C$ 对应于 $b\in B_a$，
则令 $U_c=U_{a,b}$ 且
$f_c=g_a\circ f_{a,b}:U_c\to Z$，结论随即得到。
\end{proof}

\noindent
下面应用这些思路，把一般情形约化到分离代数空间的情形。

\begin{lemma}
\label{spaces-pushouts-lemma-colimit-separated-enough}
设 $S$ 是概形，$B$ 是 $S$ 上的代数空间。
令 $\mathcal{I}\to(\Sch/S)_{fppf}$，$i\mapsto X_i$
为 $B$ 上代数空间的一个图表。假设
\begin{enumerate}
\item 每个 $X_i$ 都在 $B$ 上分离；
\item 在 $B$ 上分离代数空间的范畴中存在 $X=\colim X_i$；
\item $\coprod X_i\to X$ 是满射；
\item 若 $U\to X$ 是代数空间的平展分离态射且
$U_i=X_i\times_XU$，则在 $B$ 上分离代数空间的范畴中
$U=\colim U_i$；以及
\item $\lim X_{i,spaces,\etale}$ 中每个满足 $U_i\to X_i$ 分离的对象
$(U_i\to X_i)$，都可写成 $U_i=X_i\times_XU$，其中
$U\to X$ 是某个代数空间的平展分离态射。
\end{enumerate}
则在全体 $B$ 上代数空间的范畴中有 $X=\colim X_i$。
\end{lemma}

\begin{proof}
我们建议读者改看引理 \ref{spaces-pushouts-lemma-colimit-check-etale-locally} 及其证明。

\medskip\noindent
设 $Z$ 是 $B$ 上的代数空间。设 $f_i:X_i\to Z$ 是一族态射，
使得对每个 $i\to j$，复合 $X_i\to X_j\to Z$ 等于 $f_i$。
我们必须构造一个 $B$ 上代数空间的态射 $f:X\to Z$，
使 $f_i$ 可由复合 $X_i\to X\to Z$ 得到。
令 $W\to Z$ 是从一个概形到 $Z$ 的满平展态射。
可以假设 $W$ 是仿射概形的不交并；特别地，可以假设
$W\to Z$ 分离且 $W$ 在 $B$ 上分离。对每个 $i$，令
$U_i=W\times_{Z,f_i}X_i$，并把投影记作 $h_i:U_i\to W$。
于是 $U_i\to X_i$ 构成 $\lim X_{i,spaces,\etale}$ 的一个对象，
且 $U_i\to X_i$ 分离。由假设 (5)，可以找到代数空间的平展分离态射
$U\to X$ 以及（函子性的）同构 $U_i=X_i\times_XU$。
由假设 (4)，存在 $B$ 上态射 $h:U\to W$，使复合
$U_i\to U\to W$ 都是 $h_i$。令 $g:U\to Z$ 为 $h$ 与映射
$W\to Z$ 的复合。为完成证明，必须说明 $g:U\to Z$
可下降为态射 $X\to Z$。为此，考虑态射
$(h, h) : U \times_X U \to W \times_S W$.
将它与 $U_i\times_{X_i}U_i\to U\times_XU$ 复合，所得态射
$(h_i,h_i)$ 通过 $W\times_ZW$ 分解。由 (4)，$U\times_XU$
是诸代数空间 $U_i\times_{X_i}U_i$ 在 $B$ 上分离代数空间范畴中的
余极限，故 $(h,h)$ 通过 $W\times_ZW$ 分解。因此两个复合
$U\times_XU\to U\to W\to Z$ 相等。
由于每个 $U_i\to X_i$ 都是满射，再由假设 (2)，可知 $U\to X$ 是满射。
由于 $Z$ 是平展拓扑下的层，便得到 $g:U\to Z$ 按所需下降为
$f:X\to Z$。
\end{proof}





\section{平展层的下降}
\label{spaces-pushouts-section-glueing-etale}

\noindent
本节是《平展上同调》第 \ref{etale-cohomology-section-glueing-etale} 节
在代数空间情形下的类似版本。

\medskip\noindent
为了方便地表述结果，需要先引入一些记号。
设 $S$ 是概形。令 $\mathcal{U}=\{f_i:X_i\to X\}$
为 $S$ 上代数空间的一族具有固定靶的态射。
关于 $\mathcal{U}$ 的平展层{\it 下降数据}，是一个族
$((\mathcal{F}_i)_{i \in I}, (\varphi_{ij})_{i, j \in I})$，其中
\begin{enumerate}
\item $\mathcal{F}_i$ 属于 $\Sh(X_{i, \etale})$；并且
\item $\varphi_{ij} :
\text{pr}_{0, small}^{-1} \mathcal{F}_i
\longrightarrow
\text{pr}_{1, small}^{-1} \mathcal{F}_j$
是 $\Sh((X_i \times_X X_j)_\etale)$ 中的同构。
\end{enumerate}
并且满足{\it 上链条件}：图表
$$
\xymatrix{
\text{pr}_{0, small}^{-1}\mathcal{F}_i
\ar[dr]_{\text{pr}_{02, small}^{-1}\varphi_{ik}}
\ar[rr]^{\text{pr}_{01, small}^{-1}\varphi_{ij}} & &
\text{pr}_{1, small}^{-1}\mathcal{F}_j
\ar[dl]^{\text{pr}_{12, small}^{-1}\varphi_{jk}} \\
& \text{pr}_{2, small}^{-1}\mathcal{F}_k
}
$$
在 $\Sh((X_i \times_X X_j \times_X X_k)_\etale)$ 中交换。
存在显然的{\it 下降数据态射}概念，由此得到下降数据的范畴。
称下降数据
$((\mathcal{F}_i)_{i \in I}, (\varphi_{ij})_{i, j \in I})$
是{\it 有效的}，如果存在 $\Sh(X_\etale)$ 中的 $\mathcal{F}$ 以及同构
$\varphi_i : f_{i, small}^{-1} \mathcal{F} \to \mathcal{F}_i$
它们在 $\Sh(X_{i,\etale})$ 中与 $\varphi_{ij}$ 相容；也就是说，满足
$$
\varphi_{ij} =
\text{pr}_{1, small}^{-1} (\varphi_j) \circ
\text{pr}_{0, small}^{-1} (\varphi_i^{-1})
$$
也可以表述如下。给定 $\Sh(X_\etale)$ 的一个对象 $\mathcal{F}$，
得到{\it 典范下降数据} $(f_{i, small}^{-1}\mathcal{F}_i,c_{ij})$，
其中 $c_{ij}$ 是典范同构
$$
c_{ij} : \text{pr}_{0, small}^{-1} f_{i, small}^{-1}\mathcal{F}
\longrightarrow
\text{pr}_{1, small}^{-1} f_{j, small}^{-1}\mathcal{F}
$$
下降数据
$((\mathcal{F}_i)_{i \in I}, (\varphi_{ij})_{i, j \in I})$
有效，当且仅当它同构于与 $\Sh(X_\etale)$ 中某个 $\mathcal{F}$
相伴的典范下降数据。

\medskip\noindent
如果该族只含一个态射 $\{X\to Y\}$，则把下降数据看作一对
$(\mathcal{F},\varphi)$，其中 $\mathcal{F}$ 是 $\Sh(X_\etale)$ 的对象，
而 $\varphi$ 是 $\Sh((X\times_YX)_\etale)$ 中的同构
$$
\text{pr}_{0, small}^{-1} \mathcal{F}
\longrightarrow
\text{pr}_{1, small}^{-1} \mathcal{F}
$$
，并满足上链条件：
$$
\xymatrix{
\text{pr}_{0, small}^{-1}\mathcal{F}
\ar[dr]_{\text{pr}_{02, small}^{-1}\varphi}
\ar[rr]^{\text{pr}_{01, small}^{-1}\varphi} & &
\text{pr}_{1, small}^{-1}\mathcal{F}
\ar[dl]^{\text{pr}_{12, small}^{-1}\varphi} \\
& \text{pr}_{2, small}^{-1}\mathcal{F}
}
$$
该图表在 $\Sh((X\times_YX\times_YX)_\etale)$ 中交换。
下降数据态射与有效性的概念和前面完全相同。

\begin{lemma}
\label{spaces-pushouts-lemma-glue-etale-sheaf-etale}
设 $S$ 是概形。令 $\{f_i:X_i\to X\}$ 为代数空间的一个平展覆盖。函子
$$
\Sh(X_\etale)
\longrightarrow
\text{关于下列族的平展层下降数据 }\{f_i : X_i \to X\}
$$
是范畴等价。
\end{lemma}

\begin{proof}
在《空间的性质》第 \ref{spaces-properties-section-etale-site} 节中，
我们定义了一个位点 $X_{spaces,\etale}$：其对象是在 $X$ 上平展的
代数空间，其覆盖为平展覆盖。此外，我们有等同
$\Sh(X_\etale)=\Sh(X_{spaces,\etale})$，它与代数空间的态射相容，
即与正像和逆像相容。因此，本引理来自《位点》第
\ref{sites-section-glueing-sheaves} 节中更为一般的讨论。
\end{proof}

\begin{lemma}
\label{spaces-pushouts-lemma-reduce-to-scheme-base}
设 $S$ 是概形。令 $f:X\to Y$ 为 $S$ 上代数空间的态射，
并令 $\{Y_i\to Y\}_{i\in I}$ 为代数空间的一个平展覆盖。
如果对每个 $i\in I$，函子
$$
\Sh(Y_{i, \etale})
\longrightarrow
\text{关于下列族的平展层下降数据 }\{X \times_Y Y_i \to Y_i\}
$$
是范畴等价，并且对每个 $i,j\in I$，函子
$$
\Sh((Y_i \times_Y Y_j)_\etale)
\longrightarrow
\text{关于下列族的平展层下降数据 }
\{X \times_Y Y_i \times_Y Y_j \to Y_i \times_Y Y_j\}
$$
也是范畴等价，那么
$$
\Sh(Y_\etale)
\longrightarrow
\text{关于下列族的平展层下降数据 }\{X \to Y\}
$$
是范畴等价。
\end{lemma}

\begin{proof}
这是引理 \ref{spaces-pushouts-lemma-glue-etale-sheaf-etale} 和定义的形式推论。
\end{proof}

\begin{lemma}
\label{spaces-pushouts-lemma-representable-case}
设 $S$ 是概形。令 $f:X\to Y$ 为 $S$ 上代数空间的态射。
假设 $f$（由概形）可表，并且 $f$ 具有下列性质之一：
满且整；满且固有；或者满、平坦且局部有限呈示。那么
$$
\Sh(Y_\etale)
\longrightarrow
\text{关于下列族的平展层下降数据 }\{X \to Y\}
$$
是范畴等价。
\end{lemma}

\begin{proof}
引理陈述中提到的代数空间态射的每一种性质都在任意基变换下保持；
见《空间》第 \ref{spaces-section-lists} 节中的列表。
因此可以应用引理 \ref{spaces-pushouts-lemma-reduce-to-scheme-base}，
从而在 $Y$ 上平展局部地工作。这样便约化到 $Y$ 是概形的情形；
略去一些细节。此时 $X$ 也是概形，结论由《平展上同调》引理
\ref{etale-cohomology-lemma-glue-etale-sheaf-integral-surjective},
\ref{etale-cohomology-lemma-glue-etale-sheaf-proper-surjective}, or
\ref{etale-cohomology-lemma-glue-etale-sheaf-fppf-cover}.
\end{proof}

\begin{lemma}
\label{spaces-pushouts-lemma-reduce-to-scheme-source}
设 $S$ 是概形。令 $f:X\to Y$ 为 $S$ 上代数空间的态射，
并令 $\pi:X'\to X$ 为代数空间的态射。假设
\begin{enumerate}
\item $f\circ\pi$（由概形）可表；
\item $f\circ\pi$ 具有下列性质之一：满且整；满且固有；
或者满、平坦且局部有限呈示。
\end{enumerate}
那么
$$
\Sh(Y_\etale)
\longrightarrow
\text{关于下列族的平展层下降数据 }\{X \to Y\}
$$
是范畴等价。
\end{lemma}

\begin{proof}
这是引理 \ref{spaces-pushouts-lemma-representable-case} 和《叠》引理
\ref{stacks-lemma-compare-descent-condition} 的形式推论。
\end{proof}

\begin{lemma}
\label{spaces-pushouts-lemma-glue-etale-sheaf-proper-surjective}
设 $S$ 是概形。令 $f:X\to Y$ 为 $S$ 上代数空间的态射，
并且具有下列性质之一：满且整；满且固有；或者满、平坦且局部有限呈示。
则函子
$$
\Sh(Y_\etale)
\longrightarrow
\text{关于下列族的平展层下降数据 }\{X \to Y\}
$$
是范畴等价。
\end{lemma}

\begin{proof}
注意，固有满态射的基变换仍然固有且满；见《空间的态射》引理
\ref{spaces-morphisms-lemma-base-change-proper} 和
\ref{spaces-morphisms-lemma-base-change-surjective}。
因此，由引理 \ref{spaces-pushouts-lemma-reduce-to-scheme-base}，
可以在 $Y$ 上平展局部地工作。于是约化到 $Y$ 为仿射概形的情形；
略去一些细节。

\medskip\noindent
假设 $Y$ 仿射。由引理 \ref{spaces-pushouts-lemma-reduce-to-scheme-source}，
只需找到态射 $X'\to X$，其中 $X'$ 是概形，且 $X'\to Y$
满足下列性质之一：满且整；满且固有；或者满、平坦且局部有限呈示。

\medskip\noindent
当 $X\to Y$ 整且满时，可以取 $X'=X$，因为整态射可表。

\medskip\noindent
如果 $f$ 固有且满，那么代数空间 $X$ 拟紧且分离；
见《空间的态射》第 \ref{spaces-morphisms-section-quasi-compact} 节和引理
\ref{spaces-morphisms-lemma-separated-over-separated}。
选取一个概形 $X'$ 和一个满有限态射 $X'\to X$；见《空间的极限》命题
\ref{spaces-limits-proposition-there-is-a-scheme-finite-over}。
于是 $X'\to Y$ 满且固有。

\medskip\noindent
最后，如果 $X\to Y$ 满、平坦且局部有限呈示，
则可取一个仿射平展覆盖 $\{U_i\to X\}$，并令
$X'$ 等于不交并 $\coprod U_i$。
\end{proof}

\begin{lemma}
\label{spaces-pushouts-lemma-glue-etale-sheaf-fppf}
设 $S$ 是概形。令 $\{f_i:X_i\to X\}$ 为 $S$ 上代数空间的一个
fppf 覆盖。函子
$$
\Sh(X_\etale)
\longrightarrow
\text{关于下列族的平展层下降数据 }\{f_i : X_i \to X\}
$$
是范畴等价。
\end{lemma}

\begin{proof}
对态射 $f:\coprod X_i\to X$，可以应用引理
\ref{spaces-pushouts-lemma-glue-etale-sheaf-proper-surjective}。
然后，一个形式论证表明，关于 $f$ 的下降数据与关于该覆盖的下降数据
是一回事；比较《下降》引理 \ref{descent-lemma-family-is-one}。
略去细节。
\end{proof}

\begin{lemma}
\label{spaces-pushouts-lemma-glue-etale-sheaf-modification}
设 $S$ 是概形。令 $f:Y'\to Y$ 为 $S$ 上代数空间的固有态射，
并令 $i:Z\to Y$ 为闭浸入。置 $E=Z\times_YY'$。图表为
$$
\xymatrix{
E \ar[d]_g \ar[r]_j & Y' \ar[d]^f \\
Z \ar[r]^i & Y
}
$$
如果 $f$ 在 $Y\setminus Z$ 上是同构，那么函子
$$
\Sh(Y_\etale)
\longrightarrow
\Sh(Y'_\etale) \times_{\Sh(E_\etale)} \Sh(Z_\etale)
$$
是范畴等价。
\end{lemma}

\begin{proof}
注意，$X=Y'\coprod Z\to Y$ 是固有满态射。因此，只需构造范畴等价
$$
\Sh(Y'_\etale) \times_{\Sh(E_\etale)} \Sh(Z_\etale)
\longrightarrow
\text{关于下列族的平展层下降数据 }\{X \to Y\}
$$
并要求它与来自 $Y$ 的拉回函子相容；这样就可以用引理
\ref{spaces-pushouts-lemma-glue-etale-sheaf-proper-surjective} 得出结论。
因此，令 $(\mathcal{G}',\mathcal{G},\alpha)$ 为
$\Sh(Y'_\etale)\times_{\Sh(E_\etale)}\Sh(Z_\etale)$ 的一个对象，
记号如《范畴》例 \ref{categories-example-2-fibre-product-categories}。
可以考虑 $X$ 上的层 $\mathcal{F}$：在分支 $Y'$ 上取
$\mathcal{G}'$，在分支 $Z$ 上取 $\mathcal{G}$。我们有
$$
X \times_Y X = Y' \times_Y Y' \amalg
Y' \times_Y Z \amalg Z \times_Y Y' \amalg Z \times_Y Z =
Y' \times_Y Y' \amalg E \amalg E \amalg Z
$$
在分支 $E$、$E$、$Z$ 上，$\mathcal{F}$ 到这个代数空间的两个拉回之间
的同构是显然的。证明中有实质的部分，是在 $Y'\times_YY'$ 上找到同构
$\text{pr}_{0, small}^{-1}\mathcal{G}' \to
\text{pr}_{1, small}^{-1}\mathcal{G}'$
并使它满足上链条件。然而，$Y'\to Y$ 在 $Y\setminus Z$ 上为同构的假设蕴含
$$
h : Y \coprod E \times_Z E \longrightarrow Y' \times_Y Y'
$$
是满固有态射。（事实上它是有限态射，因为它是两个闭浸入的不交并。）
因此，只需在由 $h_{small}$ 拉回
$\text{pr}_{0,small}^{-1}\mathcal{G}'$ 和
$\text{pr}_{1,small}^{-1}\mathcal{G}'$ 后，构造满足某个上链条件的同构。
在对角线上，如何构造是清楚的。对于到 $E\times_ZE$ 的拉回，
我们利用这样一个事实：两个层都拉回为 $\mathcal{G}$ 经态射
$E\times_ZE\to Z$ 的拉回。略去细节。
\end{proof}






\section{代数空间平展态射的下降}
\label{spaces-pushouts-section-descending-etale}

\noindent
本节把第 \ref{spaces-pushouts-section-glueing-etale} 节中关于平展层的粘合结果
与代数空间的灵活性结合起来，得到代数空间平展态射的一些下降陈述。

\begin{lemma}
\label{spaces-pushouts-lemma-descend-etale-proper-surjective}
设 $S$ 是概形。令 $f:X\to Y$ 为 $S$ 上代数空间的固有满态射。
关于 $f$ 的任一下降数据 $(U/X,\varphi)$
（《空间上的下降》定义 \ref{spaces-descent-definition-descent-datum}），
若 $U$ 在 $X$ 上平展，则它是有效的
（《空间上的下降》定义 \ref{spaces-descent-definition-effective}）。
更确切地说，存在代数空间的平展态射 $V\to Y$，
其相应的典范下降数据同构于 $(U/X,\varphi)$。
\end{lemma}

\begin{proof}
回忆 $U$ 给出 $\Sh(X_{spaces,\etale})=\Sh(X_\etale)$ 中的可表层
$\mathcal{F}=h_U$；见《空间的性质》第
\ref{spaces-properties-section-etale-site} 节。
$U$ 关于 $f$ 的下降数据恰好给出平展层关于 $\{X\to Y\}$ 的下降数据
$(\mathcal{F},\varphi)$。由引理
\ref{spaces-pushouts-lemma-glue-etale-sheaf-proper-surjective}，这个下降数据是有效的。
令 $\mathcal{G}$ 为 $Y_\etale$ 上相应的层。
由《空间的性质》引理 \ref{spaces-properties-lemma-sheaf-gives-space}，
得到对应于 $\mathcal{G}$ 的代数空间平展态射 $V\to Y$；
我们略去对集合论条件的验证\footnote{这由 $\mathcal{F}$ 满足相应条件这一事实推出。}。
给定的同构 $\mathcal{F}\to f_{small}^{-1}\mathcal{G}$
对应于与下降数据相容的同构 $U\to V\times_YX$。
\end{proof}

\begin{lemma}
\label{spaces-pushouts-lemma-glue-etale-space-modification}
设 $S$ 是概形。令 $f:Y'\to Y$ 为 $S$ 上代数空间的固有态射，
令 $i:Z\to Y$ 为闭浸入，并置 $E=Z\times_YY'$。图表为
$$
\xymatrix{
E \ar[d]_g \ar[r]_j & Y' \ar[d]^f \\
Z \ar[r]^i & Y
}
$$
如果 $f$ 在 $Y\setminus Z$ 上是同构，那么函子
$$
Y_{spaces, \etale}
\longrightarrow
Y'_{spaces, \etale} \times_{E_{spaces, \etale}} Z_{spaces, \etale}
$$
是范畴等价。
\end{lemma}

\begin{proof}
令 $(V'\to Y',W\to Z,\alpha)$ 为右端的一个对象。
回忆 $V'$ 和 $W$ 分别给出可表层
$\mathcal{G}'=h_{V'}$（属于
$\Sh(Y'_{spaces,\etale})=\Sh(Y'_\etale)$）和
$\mathcal{G}=h_W$（属于 $\Sh(Z_{spaces,\etale})=\Sh(Z_\etale)$）；
见《空间的性质》第 \ref{spaces-properties-section-etale-site} 节。
同构 $\alpha:V'\times_{Y'}E\to W\times_ZE$ 决定 $E$ 上诸层的同构
$j_{small}^{-1}\mathcal{G}'\to g_{small}^{-1}\mathcal{G}$。
由引理 \ref{spaces-pushouts-lemma-glue-etale-sheaf-modification}，
得到 $Y$ 上的唯一层 $\mathcal{F}$，它以与该同构相容的方式分别拉回为
$\mathcal{G}'$ 和 $\mathcal{G}$。% P10-E0011
由《空间的性质》引理 \ref{spaces-properties-lemma-sheaf-gives-space}，
得到对应于 $\mathcal{F}$ 的代数空间平展态射 $V\to Y$；
我们略去对集合论条件的验证\footnote{这由 $\mathcal{G}$ 和
$\mathcal{G}'$ 都满足相应条件这一事实推出。}。% P10-E0012
给定的同构 $\mathcal{G}'\to f_{small}^{-1}\mathcal{F}$ 和
$\mathcal{G}\to i_{small}^{-1}\mathcal{F}$ 分别对应于同构
$V'\to V\times_YY'$ 和 $W\to V\times_YZ$；
它们按所需与 $\alpha$ 相容。
\end{proof}









\section{沿加厚与仿射态射的推出}
\label{spaces-pushouts-section-pushouts}

\noindent
本节类似于《态射续篇》第 \ref{more-morphisms-section-pushouts} 节。

\begin{lemma}
\label{spaces-pushouts-lemma-pushout-along-thickening-schemes}
设 $S$ 是概形。令 $X\to X'$ 为 $S$ 上概形的加厚，
令 $X\to Y$ 为 $S$ 上概形的仿射态射。
令 $Y'=Y\amalg_XX'$ 为概形范畴中的推出（见《态射续篇》引理
\ref{more-morphisms-lemma-pushout-along-thickening}）。
那么 $Y'$ 在 $S$ 上代数空间的范畴中也是推出。
\end{lemma}

\begin{proof}
这是引理 \ref{spaces-pushouts-lemma-colimit-agrees} 和《态射续篇》引理
\ref{more-morphisms-lemma-pushout-along-thickening},
\ref{more-morphisms-lemma-equivalence-categories-schemes-over-pushout} 和
\ref{more-morphisms-lemma-equivalence-categories-schemes-over-pushout-flat}.
\end{proof}

\begin{lemma}
\label{spaces-pushouts-lemma-pushout-along-thickening}
设 $S$ 是概形。令 $X\to X'$ 为 $S$ 上代数空间的加厚，
令 $X\to Y$ 为 $S$ 上代数空间的仿射态射。那么在 $S$ 上代数空间的
范畴中存在推出
$$
\xymatrix{
X \ar[r] \ar[d]_f
&
X' \ar[d]^{f'}
\\
Y \ar[r]
&
Y \amalg_X X'
}
$$
。此外，$Y'=Y\amalg_XX'$ 是 $Y$ 的加厚，并且
$$
\mathcal{O}_{Y'} = \mathcal{O}_Y \times_{f_*\mathcal{O}_X} f'_*\mathcal{O}_{X'}
$$
作为 $Y_\etale=(Y')_\etale$ 上的层，有上述等式。
\end{lemma}

\begin{proof}
选取概形 $V$ 和满平展态射 $V\to Y$。置 $U=V\times_YX$。
这是一个在 $V$ 上仿射的概形，并且有满平展态射 $U\to X$。
由《空间的态射续篇》引理
\ref{spaces-more-morphisms-lemma-thickening-equivalence}，
存在满平展态射 $U'\to X'$，满足 $U=U'\times_{X'}X$。
特别地，概形态射 $U\to U'$ 也是加厚。应用《态射续篇》引理
\ref{more-morphisms-lemma-pushout-along-thickening}，
在概形范畴中得到推出 $V'=V\amalg_UU'$。

\medskip\noindent
重复这一过程，在概形范畴中构造推出
$$
\xymatrix{
U \times_X U \ar[d] \ar[r] & U' \times_{X'} U' \ar[d] \\
V \times_Y V \ar[r] & R'
}
$$
。考虑态射
$$
U \times_X U \to U \to V',\quad
U' \times_{X'} U' \to U' \to V',\quad
V \times_Y V \to V \to V'
$$
其中每种情形都使用第一个投影。显然，这些态射可粘合成态射
$t':R'\to V'$；由《态射续篇》引理
\ref{more-morphisms-lemma-equivalence-categories-schemes-over-pushout-flat}.
可知它平展。类似地，得到平展态射 $s':R'\to V'$。
态射 $j'=(t',s'):R'\to V'\times_SV'$ 是非分歧的
（因为 $t'$ 平展），且限制到闭子概形 $V\times_YV\subset R'$ 后是单态射。
由于 $V\times_YV\subset R'$ 是加厚，故 $j'$ 本身也是单态射。
最后，$j'$ 是等价关系：可以利用概形推出的函子性构造态射
$c':R'\times_{s',V',t'}R'\to R'$（略去细节）。
此时置 $Y'=V'/R'$；见《空间》定理 \ref{spaces-theorem-presentation}。

\medskip\noindent
我们有态射 $X'=U'/U'\times_{X'}U'\to V'/R'=Y'$ 和
$Y=V/V\times_YV\to V'/R'=Y'$。按构造，它们组成交换图表
$$
\xymatrix{
X \ar[r] \ar[d]_f & X' \ar[d]^{f'} \\
Y \ar[r] & Y'
}
$$
由于 $Y\to Y'$ 是加厚，有 $Y_\etale=(Y')_\etale$；
见《空间的态射续篇》引理
\ref{spaces-more-morphisms-lemma-thickening-equivalence}。
图表的交换性给出这个位点上的层态射
$$
\mathcal{O}_{Y'}
\longrightarrow
\mathcal{O}_Y \times_{f_*\mathcal{O}_X} f'_*\mathcal{O}_{X'}
$$
。由《态射续篇》引理
\ref{more-morphisms-lemma-pushout-along-thickening}
可知，将此映射限制到概形 $V'$ 后是同构，因此它本身也是同构。

\medskip\noindent
为完成证明，我们说明上述图表在代数空间范畴中是推出。
为此，令 $Z$ 为代数空间，并令 $a':X'\to Z$ 和 $b:Y\to Z$
为代数空间的态射。由引理
\ref{spaces-pushouts-lemma-pushout-along-thickening-schemes}，
得到唯一的态射 $h:V'\to Z$，使下列图表交换：
$$
\vcenter{
\xymatrix{
U' \ar[d] \ar[r] & V' \ar[d]^h \\
X' \ar[r]^{a'} & Z
}
}
\quad\text{和}\quad
\vcenter{
\xymatrix{
V \ar[r] \ar[d] & V' \ar[d]^h \\
Y \ar[r]^b & Z
}
}
$$
唯一性表明 $h\circ t'=h\circ s'$。因此 $h$ 唯一地分解为
$V'\to Y'\to Z$，证明完毕。
\end{proof}

\noindent
在下面的引理中，我们使用《范畴》例
\ref{categories-example-2-fibre-product-categories} 所定义的范畴纤维积。

\begin{lemma}
\label{spaces-pushouts-lemma-categories-spaces-over-pushout}
设 $S$ 是基概形。令 $X\to X'$ 为 $S$ 上代数空间的加厚，
令 $X\to Y$ 为 $S$ 上代数空间的仿射态射。
令 $Y'=Y\amalg_XX'$ 为推出（见引理
\ref{spaces-pushouts-lemma-pushout-along-thickening}）。基变换给出函子
$$
F :
(\textit{Spaces}/Y')
\longrightarrow
(\textit{Spaces}/Y) \times_{(\textit{Spaces}/Y')} (\textit{Spaces}/X')
$$
它由 $V'\longmapsto(V'\times_{Y'}Y,V'\times_{Y'}X',1)$ 给出，
并把 $(\Sch/Y')$ 映入
$(\Sch/Y)\times_{(\Sch/Y')}(\Sch/X')$。函子 $F$ 有左伴随
$$
G :
(\textit{Spaces}/Y) \times_{(\textit{Spaces}/Y')} (\textit{Spaces}/X')
\longrightarrow
(\textit{Spaces}/Y')
$$
它把三元组 $(V,U',\varphi)$ 映为 $S$ 上代数空间范畴中的推出
$V\amalg_{(V\times_YX)}U'$。函子 $G$ 把
$(\Sch/Y)\times_{(\Sch/Y')}(\Sch/X')$ 映入 $(\Sch/Y')$。
\end{lemma}

\begin{proof}
证明完全是形式的。由于态射 $X\to X'$ 和 $X\to Y$ 可表，
显然 $F$ 把 $(\Sch/Y')$ 映入
$(\Sch/Y)\times_{(\Sch/Y')}(\Sch/X')$。

\medskip\noindent
来构造 $G$。令 $(V,U',\varphi)$ 为纤维积范畴的一个对象，
并置 $U=U'\times_{X'}X$。注意 $U\to U'$ 是加厚。
由于 $\varphi:V\times_YX\to U'\times_{X'}X=U$ 是同构，
我们得到 $X\to Y$ 上的态射 $U\to V$，它把 $U$ 等同于纤维积
$X\times_YV$。特别地，$U\to V$ 仿射；见《空间的态射》引理
\ref{spaces-morphisms-lemma-base-change-affine}。
因此可以应用引理 \ref{spaces-pushouts-lemma-pushout-along-thickening}，
得到推出 $V'=V\amalg_UU'$。由于 $V'$ 是推出，且给定的态射
$V\to Y$ 与 $U'\to X'$ 作为到 $Y'$ 的态射在 $U$ 上相合，
我们得到态射 $V'\to Y'$。令 $G(V,U',\varphi)=V'$，便得到函子 $G$。

\medskip\noindent
如果 $(V,U',\varphi)$ 是
$(\Sch/Y)\times_{(\Sch/Y')}(\Sch/X')$ 的对象，
那么 $U=U'\times_{X'}X$ 也是概形，并且由《态射续篇》引理
\ref{more-morphisms-lemma-pushout-along-thickening}，
可以在概形范畴中作推出 $V'=V\amalg_UU'$。
由引理 \ref{spaces-pushouts-lemma-pushout-along-thickening-schemes}，
这在概形范畴中也是推出，故 $G$ 把
$(\Sch/Y)\times_{(\Sch/Y')}(\Sch/X')$ 映入 $(\Sch/Y')$。% P10-E0014

\medskip\noindent
来证明 $G$ 是 $F$ 的左伴随。令 $Z$ 为 $Y'$ 上的代数空间。
我们必须证明
$$
\Mor(V', Z) = \Mor((V, U', \varphi), F(Z))
$$
其中态射集取自各自的范畴。令 $g':V'\to Z$ 为一个态射。
分别把 $g'$ 与态射 $V\to V'$、$U'\to V'$ 复合，所得复合记作
$\tilde g$、$\tilde f'$。分别沿 $Y\to Y'$、$X'\to Y'$ 对
$\tilde g$、$\tilde f'$ 作基变换，得到态射
$g:V\to Z\times_{Y'}Y$ 和 $f':U'\to Z\times_{Y'}X'$。
于是 $(g,f')$ 是上述等式右端的元素（略去细节）。
反过来，设 $(g,f'):(V,U',\varphi)\to F(Z)$ 是右端的一个元素。
可以把 $g$、$f$ 分别与
$Z\times_{Y'}X'\to Z$、$Z\times_{Y'}Y\to Z$ 复合，
得到复合 $\tilde g:V\to Z$、$\tilde f':U'\to Z$。% P10-E0015
$\tilde g$ 与 $\tilde f'$ 作为从 $U$ 到 $Z$ 的态射相合。
由推出的泛性质，得到态射 $g':V'\to Z$，即等式左端的一个元素。
略去对这两个构造互逆的验证。
\end{proof}

\begin{lemma}
\label{spaces-pushouts-lemma-diagram}
设 $S$ 是概形。令
$$
\xymatrix{
A \ar[r] \ar[d] & C \ar[d] \ar[r] & E \ar[d] \\
B \ar[r] & D \ar[r] & F
}
$$
为 $S$ 上代数空间的交换图表。假设 $A,B,C,D$ 和 $A,B,E,F$
分别组成笛卡尔方形，并且 $B\to D$ 满且平展。
那么 $C,D,E,F$ 组成笛卡尔方形。
\end{lemma}

\begin{proof}
这是形式的。
\end{proof}

\begin{lemma}
\label{spaces-pushouts-lemma-equivalence-categories-spaces-over-pushout}
在引理 \ref{spaces-pushouts-lemma-categories-spaces-over-pushout} 的情形下，
函子 $F\circ G$ 同构于恒等函子。
\end{lemma}

\begin{proof}
我们把这个断言约化为《态射续篇》引理
\ref{more-morphisms-lemma-equivalence-categories-schemes-over-pushout}.
中的相应断言，从而证明 $F\circ G$ 同构于恒等函子。

\medskip\noindent
选取概形 $Y_1$ 和满平展态射 $Y_1\to Y$。置 $X_1=Y_1\times_YX$。
这是一个在 $Y_1$ 上仿射的概形，并且有满平展态射 $X_1\to X$。
由《空间的态射续篇》引理
\ref{spaces-more-morphisms-lemma-thickening-equivalence}，
存在满平展态射 $X'_1\to X'$，满足 $X_1=X'_1\times_{X'}X$。
特别地，概形态射 $X_1\to X'_1$ 也是加厚。应用《态射续篇》引理
\ref{more-morphisms-lemma-pushout-along-thickening}，
在概形范畴中得到推出 $Y'_1=Y_1\amalg_{X_1}X'_1$。
在引理 \ref{spaces-pushouts-lemma-pushout-along-thickening} 的证明中，
我们把 $Y'$ 构造为 $Y'_1$ 上某个平展等价关系的商，使得得到交换图表
\begin{equation}
\label{spaces-pushouts-equation-cube}
\vcenter{
\xymatrix{
& X \ar[rr] \ar'[d][dd] & & X' \ar[dd] \\
X_1 \ar[rr] \ar[dd] \ar[ru] & & X_1' \ar[dd] \ar[ru] & \\
& Y \ar'[r][rr] & & Y' \\
Y_1 \ar[rr] \ar[ru] & & Y_1' \ar[ru]
}
}
\end{equation}
其中除前、后方形外，所有方形都是笛卡尔的
（前、后方形是推出），东北方向的箭头都是满平展态射。
把《态射续篇》引理
\ref{more-morphisms-lemma-equivalence-categories-schemes-over-pushout}
中对前方形构造的函子记作 $F_1$、$G_1$。那么范畴图表
$$
\xymatrix{
(\Sch/Y_1') \ar@<-1ex>[r]_-{F_1} \ar[d] &
(\Sch/Y_1) \times_{(\Sch/Y_1')} (\Sch/X_1') \ar[d] \ar@<-1ex>[l]_-{G_1} \\
(\textit{Spaces}/Y') \ar@<-1ex>[r]_-F &
(\textit{Spaces}/Y) \times_{(\textit{Spaces}/Y')} (\textit{Spaces}/X')
\ar@<-1ex>[l]_-G
}
$$
是交换的；这来自对基变换函子的简单考察，以及引理
\ref{spaces-pushouts-lemma-pushout-along-thickening-schemes} 中概形推出与代数空间推出的一致性。

\medskip\noindent
令 $(V,U',\varphi)$ 为
$(\textit{Spaces}/Y)\times_{(\textit{Spaces}/Y')}(\textit{Spaces}/X')$
的一个对象。记 $U=U'\times_{X'}X$，从而
$G(V,U',\varphi)=V\amalg_UU'$。选取概形 $V_1$ 和满平展态射
$V_1\to Y_1\times_YV$。置 $U_1=V_1\times_YX$。于是
$$
U_1 = V_1 \times_Y X
\longrightarrow
(Y_1 \times_Y V) \times_Y X =
X_1 \times_Y V = X_1 \times_X X \times_Y V = X_1 \times_X U
$$
也满且平展。由《空间的态射续篇》引理
\ref{spaces-more-morphisms-lemma-thickening-equivalence}
存在加厚 $U_1\to U'_1$ 和满平展态射
$U'_1\to X'_1\times_{X'}U'$，其到 $X_1\times_XU$ 的基变换就是上述态射。
此时 $(V_1,U'_1,\varphi_1)$ 是
$(\Sch/Y_1)\times_{(\Sch/Y'_1)}(\Sch/X'_1)$ 的对象。
在引理 \ref{spaces-pushouts-lemma-pushout-along-thickening} 的证明中，
我们把 $G(V,U',\varphi)=V\amalg_UU'$ 构造为
$G_1(V_1,U'_1,\varphi_1)=V_1\amalg_{U_1}U'_1$
上某个平展等价关系的商，使得得到交换图表
\begin{equation}
\label{spaces-pushouts-equation-cube-over}
\vcenter{
\xymatrix{
& U \ar[rr] \ar'[d][dd] & & U' \ar[dd] \\
U_1 \ar[rr] \ar[dd] \ar[ru] & & U_1' \ar[dd] \ar[ru] & \\
& V \ar'[r][rr] & & G(V, U', \varphi) \\
V_1 \ar[rr] \ar[ru] & & G_1(V_1, U_1', \varphi_1) \ar[ru]
}
}
\end{equation}
其中除前、后方形外，所有方形都是笛卡尔的
（前、后方形是推出），东北方向的箭头都是满平展态射。特别地，
$$
G_1(V_1, U_1', \varphi_1) \to G(V, U', \varphi)
$$
满且平展。

\medskip\noindent
最后来证明本引理。我们必须说明伴随映射
$(V,U',\varphi)\to F(G(V,U',\varphi))$ 是同构。
由《态射续篇》引理
\ref{more-morphisms-lemma-equivalence-categories-schemes-over-pushout}.
，可知 $(V_1,U'_1,\varphi_1)\to F_1(G_1(V_1,U'_1,\varphi_1))$
是同构。回忆 $F$ 和 $F_1$ 都由基变换给出。
利用 (\ref{spaces-pushouts-equation-cube-over}) 的性质和引理 \ref{spaces-pushouts-lemma-diagram}，可知
$V \to G(V, U', \varphi) \times_{Y'} Y$ 和
$U'\to G(V,U',\varphi)\times_{Y'}X'$ 都是同构；也就是说，
$(V,U',\varphi)\to F(G(V,U',\varphi))$ 是同构。
\end{proof}

\begin{lemma}
\label{spaces-pushouts-lemma-space-over-pushout-flat-modules}
设 $S$ 是基概形。令 $X\to X'$ 为 $S$ 上代数空间的加厚，
令 $X\to Y$ 为 $S$ 上代数空间的仿射态射。
令 $Y'=Y\amalg_XX'$ 为推出（见引理
\ref{spaces-pushouts-lemma-pushout-along-thickening}）。
令 $V'\to Y'$ 为 $S$ 上代数空间的态射，并置
$V=Y\times_{Y'}V'$、$U'=X'\times_{Y'}V'$、$U=X\times_{Y'}V'$。
下列两个范畴之间存在等价：
\begin{enumerate}
\item 在 $Y'$ 上平坦的拟凝聚 $\mathcal{O}_{V'}$-模；以及
\item 三元组 $(\mathcal{G},\mathcal{F}',\varphi)$ 的范畴，其中
\begin{enumerate}
\item $\mathcal{G}$ 是在 $Y$ 上平坦的拟凝聚 $\mathcal{O}_V$-模；
\item $\mathcal{F}'$ 是在 $X$ 上平坦的拟凝聚 $\mathcal{O}_{U'}$-模；以及
\item $\varphi : (U \to V)^*\mathcal{G} \to (U \to U')^*\mathcal{F}'$
是 $\mathcal{O}_U$-模同构。
\end{enumerate}
\end{enumerate}
这个等价把 $\mathcal{G}'$ 映为
$((V\to V')^*\mathcal{G}',(U'\to V')^*\mathcal{G}',can)$。
设 $\mathcal{G}'$ 对应于三元组 $(\mathcal{G},\mathcal{F}',\varphi)$。那么
\begin{enumerate}
\item[(a)] $\mathcal{G}'$ 是有限型 $\mathcal{O}_{V'}$-模，当且仅当
$\mathcal{G}$ 和 $\mathcal{F}'$ 分别是有限型
$\mathcal{O}_Y$-模和 $\mathcal{O}_{U'}$-模；
\item[(b)] 如果 $V'\to Y'$ 局部有限呈示，那么 $\mathcal{G}'$
是有限呈示 $\mathcal{O}_{V'}$-模，当且仅当 $\mathcal{G}$ 和
$\mathcal{F}'$ 分别是有限呈示 $\mathcal{O}_Y$-模和
$\mathcal{O}_{U'}$-模。
\end{enumerate}
\end{lemma}

\begin{proof}
一个拟逆函子把三元组 $(\mathcal{G},\mathcal{F}',\varphi)$ 映为纤维积
$$
(V \to V')_*\mathcal{G}
\times_{(U \to V')_*\mathcal{F}}
(U' \to V')_*\mathcal{F}'
$$
其中 $\mathcal{F}=(U\to U')^*\mathcal{F}'$。这是可行的，因为在
$V'$ 和 $Y'$ 上平展的仿射概形上，我们恢复了《代数续篇》引理
\ref{more-algebra-lemma-relative-flat-module-over-fibre-product}.
中的等价。略去细节。

\medskip\noindent
通过平展局部化（《空间的性质》第
\ref{spaces-properties-section-properties-modules} 节），
(a) 和 (b) 约化到 $V'$ 和 $Y'$ 都仿射的情形；此时结论来自
《代数续篇》引理
\ref{more-algebra-lemma-relative-finite-module-over-fibre-product} 和
\ref{more-algebra-lemma-relative-finitely-presented-module-over-fibre-product}.
\end{proof}


\begin{lemma}
\label{spaces-pushouts-lemma-equivalence-categories-spaces-pushout-flat}
在引理 \ref{spaces-pushouts-lemma-equivalence-categories-spaces-over-pushout} 的情形下，
如果对某个三元组 $(V,U',\varphi)$ 有 $V'=G(V,U',\varphi)$，那么
\begin{enumerate}
\item $V'\to Y'$ 局部有限型，当且仅当 $V\to Y$ 和 $U'\to X'$
都局部有限型；
\item $V'\to Y'$ 平坦，当且仅当 $V\to Y$ 和 $U'\to X'$ 都平坦；
\item $V'\to Y'$ 平坦且局部有限呈示，当且仅当 $V\to Y$ 和
$U'\to X'$ 都平坦且局部有限呈示；
\item $V'\to Y'$ 光滑，当且仅当 $V\to Y$ 和 $U'\to X'$ 都光滑；
\item $V'\to Y'$ 平展，当且仅当 $V\to Y$ 和 $U'\to X'$ 都平展；以及
\item 可按需要在此增添更多性质。
\end{enumerate}
如果 $W'$ 在 $Y'$ 上平坦，那么伴随映射 $G(F(W'))\to W'$ 是同构。
因此，$F$ 和 $G$ 在以下两个范畴之间定义互为拟逆的函子：
$Y'$ 上平坦的空间的范畴；以及满足 $V\to Y$ 和 $U'\to X'$ 平坦的
三元组 $(V,U',\varphi)$ 的范畴。
\end{lemma}

\begin{proof}
选取图表 (\ref{spaces-pushouts-equation-cube})，如引理
\ref{spaces-pushouts-lemma-equivalence-categories-spaces-over-pushout} 的证明那样。

\medskip\noindent
证明 (1)--(5)。令 $(V,U',\varphi)$ 为
$(\textit{Spaces}/Y)\times_{(\textit{Spaces}/Y')}(\textit{Spaces}/X')$
的一个对象。构造图表 (\ref{spaces-pushouts-equation-cube-over})，如引理
\ref{spaces-pushouts-lemma-equivalence-categories-spaces-over-pushout} 的证明那样。那么
$G(V,U',\varphi)\to Y'$ 到 $Y'_1$ 的基变换是
$G_1(V_1,U'_1,\varphi_1)\to Y'_1$。因此 (1)--(5) 立即来自
《态射续篇》引理
\ref{more-morphisms-lemma-equivalence-categories-schemes-over-pushout-flat}
中关于概形的相应断言。

\medskip\noindent
假设 $W'\to Y'$ 平坦。选取概形 $W'_1$ 和满平展态射
$W'_1\to Y'_1\times_{Y'}W'$。注意，$W'_1\to W'$ 是满平展态射，
因为它是满平展态射的复合。由《态射续篇》引理
\ref{more-morphisms-lemma-equivalence-categories-schemes-over-pushout-flat}
将其应用于 $Y'_1$ 上的 $W'_1$ 和图表的前面（函子 $G_1$、$F_1$
如引理 \ref{spaces-pushouts-lemma-equivalence-categories-spaces-over-pushout} 的证明），
可知 $G_1(F_1(W'_1))\to W'_1$ 是同构。
于是 $G(F(W'))$ 的构造（即作为引理
\ref{spaces-pushouts-lemma-pushout-along-thickening} 中构造的推出）表明
$G_1(F_1(W'_1))\to G(F(W))$ 满且平展。% P10-E0018
从而 $G(F(W))\to W$ 平展；例如见《空间的性质》引理
\ref{spaces-properties-lemma-etale-local}。
但按构造，$G(F(W))\to W$ 在底层约化代数空间上是同构，故它是同构。
\end{proof}







\section{沿闭浸入与整态射的推出}
\label{spaces-pushouts-section-pushouts-II}

\noindent
本节类似于《态射续篇》第 \ref{more-morphisms-section-pushouts-II} 节。

\begin{lemma}
\label{spaces-pushouts-lemma-pushout-along-closed-immersion-and-integral}
在《态射续篇》情形
\ref{more-morphisms-situation-pushout-along-closed-immersion-and-integral}
中，令 $Y\amalg_ZX$ 为概形范畴中的推出
（《态射续篇》命题
\ref{more-morphisms-proposition-pushout-along-closed-immersion-and-integral}）。
那么 $Y\amalg_ZX$ 在 $S$ 上代数空间的范畴中也是推出。
\end{lemma}

\begin{proof}
这是引理 \ref{spaces-pushouts-lemma-colimit-agrees}、本引理中提到的命题以及
《态射续篇》引理
\ref{more-morphisms-lemma-pushout-functor} 和
\ref{more-morphisms-lemma-pushout-functor-equivalence-flat}.
立即得到引理 \ref{spaces-pushouts-lemma-colimit-agrees} 的条件 (1) 和 (2)。
为验证 (3) 和 (4)，注意平展态射局部拟有限，并利用《态射续篇》引理
\ref{more-morphisms-lemma-pushout-functor-equivalence-flat}
中的范畴等价是用《态射续篇》引理
\ref{more-morphisms-lemma-pushout-functor} 的推出构造得到的。
略去次要细节。
\end{proof}








\section{推出与导出范畴}
\label{spaces-pushouts-section-pushouts-derived}

\noindent
本节讨论模的导出范畴在推出下的表现。

\begin{lemma}
\label{spaces-pushouts-lemma-pushout-along-thickening-derived}
设 $S$ 是概形。考虑 $S$ 上代数空间范畴中的推出
$$
\xymatrix{
X \ar[r]_i \ar[d]_f & X' \ar[d]^{f'}
\\
Y \ar[r]^j & Y'
}
$$
，它如引理 \ref{spaces-pushouts-lemma-pushout-along-thickening} 中所述。
假设 $i$ 是加厚。那么函子\footnote{所有函子都由导出拉回给出。}
的本质像
$$
D(\mathcal{O}_{Y'}) \longrightarrow
D(\mathcal{O}_Y) \times_{D(\mathcal{O}_X)} D(\mathcal{O}_{X'})
$$
包含每个三元组 $(M,K',\alpha)$，其中
$M\in D(\mathcal{O}_Y)$ 和 $K'\in D(\mathcal{O}_{X'})$ 都是伪凝聚的。
\end{lemma}

\begin{proof}
令 $(M,K',\alpha)$ 为本引理函子靶范畴的一个对象。
这里 $\alpha:Lf^*M\to Li^*K'$ 是同构，它伴随于映射
$\beta:M\to Rf_*Li^*K'$。于是得到映射
$$
Rj_*M \xrightarrow{Rj_*\beta}
Rj_*Rf_*Li^*K' = Rf'_*Ri_*Li^*K' \leftarrow Rf'_*K'
$$
其中指向左的箭头来自 $K'\to Ri_*Li^*K'$。选取 $D(\mathcal{O}_{Y'})$
中的特异三角
$$
M' \to Rj_*M \oplus Rf'_*K' \to Rj_*Rf_*Li^*K' \to M'[1]
$$
。第一个箭头定义与 $\alpha$ 相容的典范映射
$Lj^*M'\to M$ 和 $L(f')^*M'\to K'$。
因此，只需说明映射 $Lj^*M'\to M$ 和 $L(f')^*M'\to K$ 是同构。% P10-E0019
这可以在 $Y'$ 上平展局部地检验，因而可以假设 $Y'$ 是平展的。% P10-E0020

\medskip\noindent
假设 $Y'$ 仿射，且 $M\in D(\mathcal{O}_Y)$ 和
$K'\in D(\mathcal{O}_{X'})$ 都伪凝聚。设我们的推出对应于环的纤维积
$$
\xymatrix{
B & B' \ar[l] \\
A \ar[u] & A' \ar[l] \ar[u]
}
$$
，其中 $B'\to B$ 是满射，其核 $I$ 局部幂零
（因而 $A'\to A$ 也是满射，且其核 $I$ 局部幂零）。
关于 $M$ 和 $K'$ 的假设蕴含：$M$ 来自 $D(A)$ 的一个伪凝聚对象，
$K'$ 来自 $D(B')$ 的一个伪凝聚对象；见《空间的导出范畴》引理
\ref{spaces-perfect-lemma-pseudo-coherent},
\ref{spaces-perfect-lemma-derived-quasi-coherent-small-etale-site} 和
\ref{spaces-perfect-lemma-descend-pseudo-coherent}
以及《概形的导出范畴》引理
\ref{perfect-lemma-affine-compare-bounded} 和
\ref{perfect-lemma-pseudo-coherent-affine}.
此外，正像和导出拉回与模的导出范畴上的相应运算一致；
见《空间的导出范畴》注
\ref{spaces-perfect-remark-match-total-direct-images}
以及《概形的导出范畴》引理
\ref{perfect-lemma-quasi-coherence-pushforward} 和
\ref{perfect-lemma-quasi-coherence-pullback}.
这把问题约化为下一段所表述的断言。
（确切地说，这些参引表明对象 $M'$ 属于
$D_\QCoh(\mathcal{O}_{Y'})$，因为它是 $D(\mathcal{O}_{Y'})$
的三角子范畴。）

\medskip\noindent
给定如上的环图表以及三元组 $(M,K',\alpha)$，其中
$M\in D(A)$、$K'\in D(B')$ 伪凝聚，且
$\alpha:M\otimes_A^\mathbf{L}B\to K'\otimes_{B'}^\mathbf{L}B$ 是同构。
设有 $D(A')$ 中的特异三角
$$
M' \to M \oplus K' \to K' \otimes_{B'}^\mathbf{L} B \to M'[1]
$$
。目标是证明诱导映射 $M'\otimes_{A'}^\mathbf{L}A\to M$ 和
$M'\otimes_{A'}^\mathbf{L}B'\to K'$ 都是同构。
为此，选取一个表示 $M$ 的有限自由 $A$-模的有上界复形 $E^\bullet$。
由于 $(B',I)$ 是 Hensel 对
（《代数续篇》引理 \ref{more-algebra-lemma-locally-nilpotent-henselian}），
且 $B=B'/I$，可以应用《代数续篇》引理
\ref{more-algebra-lemma-lift-complex-finite-projectives}
，得到自由 $B'$-模的有上界复形 $P^\bullet$，使 $\alpha$ 由同构
$E^\bullet\otimes_AB\cong P^\bullet\otimes_{B'}B$ 表示。
然后可以考虑 $B'$-模复形的短正合列
$$
0 \to L^\bullet \to
E^\bullet \oplus P^\bullet \to P^\bullet \otimes_{B'} B \to 0
$$
。《代数续篇》引理
\ref{more-algebra-lemma-finitely-presented-module-over-fibre-product}
蕴含 $L^\bullet$ 是有限投射 $A'$-模的有上界复形
（事实上，在我们的情形中直接证明 $L^n$ 有限自由相当容易），并且有
$L^\bullet \otimes_{A'} A = E^\bullet$ 以及
$L^\bullet \otimes_{A'} B' = P^\bullet$.
该短正合列在 $D(A')$ 中给出特异三角
$$
L^\bullet \to M \oplus K' \to K' \otimes_{B'}^\mathbf{L} B \to (L^\bullet)[1]
$$
（《导出范畴》第 \ref{derived-section-canonical-delta-functor} 节）。
由三角范畴的一般性质（《导出范畴》第
\ref{derived-section-elementary-results} 节），它同构于给定的特异三角。
换言之，$L^\bullet$ 以与给定映射相容的方式表示 $M'$。
因此，映射 $M'\otimes_{A'}^\mathbf{L}A\to M$ 和
$M'\otimes_{A'}^\mathbf{L}B'\to K'$ 都是同构，
因为刚刚已经看到 $L^\bullet$ 的相应断言成立。
\end{proof}







\section{构造初等特异方形}
\label{spaces-pushouts-section-elementary-dsitnguished-squares}

\noindent
初等特异方形定义于《空间的导出范畴》第
\ref{spaces-perfect-section-induction} 节。

\begin{lemma}
\label{spaces-pushouts-lemma-elementary-distinguished-square-pushout}
设 $S$ 是概形。令 $(U\subset W,f:V\to W)$ 为初等特异方形。那么
$$
\xymatrix{
U \times_W V \ar[r] \ar[d] &
V \ar[d]^f \\
U \ar[r] & W
}
$$
是 $S$ 上代数空间范畴中的推出。
\end{lemma}

\begin{proof}
注意，$U\amalg V\to W$ 是满平展态射。纤维积
$$
(U \amalg V) \times_W (U \amalg V)
$$
是四个部分的不交并，即 $U=U\times_WU$、$U\times_WV$、
$V\times_WU$ 和 $V\times_WV$。存在满平展态射
$$
V \amalg (U \times_W V) \times_U (U \times_W V) \longrightarrow V \times_W V
$$
，因为 $f$ 在 $W\setminus U$ 上诱导同构
（这是初等特异方形定义的一部分）。
令 $B$ 为 $S$ 上的代数空间，并令 $g:V\to B$ 和 $h:U\to B$
为 $S$ 上态射，它们限制到 $U\times_WV$ 后相合。
于是，上面对 $(U\amalg V)\times_W(U\amalg V)$ 的描述表明，
$h\amalg g:U\amalg V\to B$ 使两个投影相等。
由于 $B$ 是平展拓扑下的层，按所需得到 $h\amalg g$
通过 $W$ 的唯一分解。
\end{proof}

\begin{lemma}
\label{spaces-pushouts-lemma-construct-elementary-distinguished-square}
设 $S$ 是概形，$V$、$U$ 是 $S$ 上的代数空间。
令 $V'\subset V$ 为开子空间，令 $f':V'\to U$ 为 $S$ 上代数空间的
分离平展态射。那么在 $S$ 上代数空间的范畴中存在推出
$$
\xymatrix{
V' \ar[r] \ar[d] &
V \ar[d]^f \\
U \ar[r] & W
}
$$
，而且 $(U\subset W,f:V\to W)$ 是初等特异方形。
\end{lemma}

\begin{proof}
我们把 $W$ 构造为 $U\amalg V$ 上某个平展等价关系 $R$ 的商。
例如，由《自举》定理 \ref{bootstrap-theorem-final-bootstrap}，
这样的商是代数空间。此外，引理
\ref{spaces-pushouts-lemma-elementary-distinguished-square-pushout} 的证明告诉我们应取
$$
R = U \amalg V' \amalg V' \amalg V \amalg
(V' \times_U V' \setminus \Delta_{V'/U}(V'))
$$
由于假设 $V'\to U$ 分离，$\Delta_{V'/U}$ 的像是闭的，
因而其补集是开子空间。态射
$j:R\to(U\amalg V)\times_S(U\amalg V)$ 由下式给出：
$$
u,\ v',\ v',\ v,\ (v'_1, v'_2) \mapsto
(u, u),\ (f'(v'), v'),\ (v', f'(v')),\ (v, v),\ (v'_1, v'_2)
$$
记号含义显然。立即可验证这是一个单态射和等价关系，
且诱导态射 $s,t:R\to U\amalg V$ 平展。
令 $W=(U\amalg V)/R$ 为商代数空间。
得到本引理陈述中的交换图表。为完成证明，只需说明该图表是初等特异方形；
因为这样一来，引理 \ref{spaces-pushouts-lemma-elementary-distinguished-square-pushout}
就蕴含它是推出。因此必须说明 $U\to W$ 是开态射，
$f$ 平展且在 $W\setminus U$ 上是同构。
这些都由 $R$ 的选取推出；略去细节。
\end{proof}






\section{拟凝聚模的形式粘合}
\label{spaces-pushouts-section-formal-glueing}

\noindent
本节类似于《代数续篇》第 \ref{more-algebra-section-formal-glueing} 节。
对于概形态射，相关结果见 Joyet 的论文 \cite{Joyet}；
这是一个很好的入门起点。关于它在叠的下降问题中的应用，
见 Moret-Bailly 的论文 \cite{MB}。对于概形的仿射态射，
论文 \cite{Ferrand-Raynaud} 的附录中有一个陈述，
但必须补充闭子概形由有限生成理想截出的假设
（如 Joyet 的论文中那样），否则结论并不成立。
这部分内容向（高阶）导出范畴的推广，以及它对非平坦情形的潜在应用，
见 \cite[第 5 节]{Bhatt-Algebraize}。

\medskip\noindent
先从一个关于支撑在闭子集上的阿贝尔层的引理开始。

\begin{lemma}
\label{spaces-pushouts-lemma-stalk-pushforward-with-support}
设 $S$ 是概形。令 $f:Y\to X$ 为 $S$ 上代数空间的态射。
令 $Z\subset X$ 为闭子空间，使 $f^{-1}Z\to Z$ 整且万有单射。
令 $\overline y$ 为 $Y$ 的几何点，并令
$\overline x=f(\overline y)$。那么
$$
(Rf_*Q)_{\overline{x}} = Q_{\overline{y}}
$$
对 $D(Y_\etale)$ 中任一支撑在 $|f^{-1}Z|$ 上的对象 $Q$，
上述等式在 $D(\textit{Ab})$ 中成立。
\end{lemma}

\begin{proof}
考虑代数空间的交换图表
$$
\xymatrix{
f^{-1}Z \ar[r]_{i'} \ar[d]_{f'} & Y \ar[d]_f \\
Z \ar[r]^i & X
}
$$
由《空间的上同调》引理
\ref{spaces-cohomology-lemma-complexes-with-support-on-closed}，
对 $D(f^{-1}Z_\etale)$ 中某个对象 $K'$，可以写成 $Q=Ri'_*K'$。
由《空间的态射》引理
\ref{spaces-morphisms-lemma-integral-universally-injective-push-pull}，
有 $K'=(f')^{-1}K$，其中 $K=Rf'_*K'$。
于是 $Rf_*Q=Rf_*Ri'_*K'=Ri_*Rf'_*K'=Ri_*K$。
令 $\overline z$ 为 $Z$ 中对应于 $\overline x$ 的几何点，
令 $\overline z'$ 为 $f^{-1}Z$ 中对应于 $\overline y$ 的几何点。
如下得到本引理的结论：
$$
Q_{\overline{y}} = (Ri'_*K')_{\overline{y}} = K'_{\overline{z}'} =
(f')^{-1}K_{\overline{z}'} = K_{\overline{z}} = Ri_*K_{\overline{x}} =
Rf_*Q_{\overline{x}}
$$
中间的等式成立，是因为《空间的性质》引理
\ref{spaces-properties-lemma-stalk-pullback} 给出了拉回茎的描述。
\end{proof}

\begin{lemma}
\label{spaces-pushouts-lemma-stalk-formal-glueing}
设 $S$ 是概形。令 $f:Y\to X$ 为 $S$ 上代数空间的态射。
令 $Z\subset X$ 为闭子空间，使 $f^{-1}Z\to Z$ 整且万有单射。
令 $\overline y$ 为 $Y$ 的几何点，并令
$\overline x=f(\overline y)$。令 $\mathcal{G}$ 为 $Y$ 上的阿贝尔层。
那么二项复形之间的映射
$$
\left(f_*\mathcal{G}_{\overline{x}} \to
(f \circ j')_*(\mathcal{G}|_V)_{\overline{x}}\right)
\longrightarrow
\left(\mathcal{G}_{\overline{y}} \to j'_*(\mathcal{G}|_V)_{\overline{y}}\right)
$$
在核上诱导同构，在余核上诱导单射。这里
$V=Y\setminus f^{-1}Z$，而 $j':V\to Y$ 是包含映射。
\end{lemma}

\begin{proof}
选取 $D(Y_\etale)$ 中的特异三角
$$
\mathcal{G} \to Rj'_*\mathcal{G}|_V \to Q \to \mathcal{G}[1]
$$
。% P10-E0021
$Q$ 的上同调层支撑在 $|f^{-1}Z|$ 上。应用 $Rf_*$，得到
$$
Rf_*\mathcal{G} \to Rf_*Rj'_*\mathcal{G}|_V \to Rf_*Q
\to Rf_*\mathcal{G}[1]
$$
在 $\overline x$ 处取茎，得到正合列
$$
0 \to
(R^{-1}f_*Q)_{\overline{x}} \to
f_*\mathcal{G}_{\overline{x}} \to
(f \circ j')_*(\mathcal{G}|_V)_{\overline{x}} \to
(R^0f_*Q)_{\overline{x}}
$$
可以把它与正合列
$$
0 \to
H^{-1}(Q)_{\overline{y}} \to
\mathcal{G}_{\overline{y}} \to
j'_*(\mathcal{G}|_V)_{\overline{y}} \to
H^0(Q)_{\overline{y}}
$$
相比较。由引理 \ref{spaces-pushouts-lemma-stalk-pushforward-with-support}，
有 $Q_{\overline y}=Rf_*Q_{\overline x}$，故本引理成立。
\end{proof}

\begin{lemma}
\label{spaces-pushouts-lemma-stalk-of-pushforward}
设 $S$ 是概形，$X$ 是 $S$ 上的代数空间。
令 $f:Y\to X$ 为拟紧且拟分离的态射。
令 $\overline x$ 为 $X$ 的几何点，并令
$\Spec(\mathcal{O}_{X,\overline x})\to X$ 为典范态射。
对 $Y$ 上的拟凝聚模 $\mathcal{G}$，有
$$
f_*\mathcal{G}_{\overline{x}} =
\Gamma(Y \times_X \Spec(\mathcal{O}_{X, \overline{x}}), p^*\mathcal{F})
$$
其中 $p:Y\times_X\Spec(\mathcal{O}_{X,\overline x})\to Y$ 是投影。% P10-E0022
\end{lemma}

\begin{proof}
注意，$f_*\mathcal{G}_{\overline x}=
\Gamma(\Spec(\mathcal{O}_{X,\overline x}),h^*f_*\mathcal{G})$，
其中 $h:\Spec(\mathcal{O}_{X,\overline x})\to X$。
由于 $h$ 平坦，可以应用《空间的上同调》引理
\ref{spaces-cohomology-lemma-flat-base-change-cohomology}，故结论成立。
\end{proof}

\begin{lemma}
\label{spaces-pushouts-lemma-stalk-of-module-with-support}
设 $S$ 是概形，$X$ 是 $S$ 上的代数空间。
令 $i:Z\to X$ 为有限呈示的闭浸入，
令 $Q\in D_\QCoh(\mathcal{O}_X)$ 支撑在 $|Z|$ 上。
令 $\overline x$ 为 $X$ 的几何点，令
$I_{\overline x}\subset\mathcal{O}_{X,\overline x}$ 为 $Z$ 的理想层的茎。
那么上同调模 $H^n(Q_{\overline x})$ 是 $I_{\overline x}$-幂挠的
（见《代数续篇》定义 \ref{more-algebra-definition-f-power-torsion}）。
\end{lemma}

\begin{proof}
选取仿射概形 $U$ 和平展态射 $U\to X$，使 $\overline x$
提升为 $U$ 的几何点 $\overline u$。于是可以分别用
$U$、$U\times_XZ$、限制 $Q|_U$、$\overline u$ 替换
$X$、$Z$、$Q$、$\overline x$。因此可以假设 $X=\Spec(A)$ 仿射。
令 $I\subset A$ 为定义 $Z$ 的理想。由于 $i:Z\to X$ 有限呈示，
理想 $I=(f_1,\ldots,f_r)$ 有限生成。
对象 $Q$ 来自一个 $A$-模复形 $M^\bullet$；见《空间的导出范畴》引理
\ref{spaces-perfect-lemma-derived-quasi-coherent-small-etale-site}
以及《概形的导出范畴》引理
\ref{perfect-lemma-affine-compare-bounded}.
由于 $Q$ 的上同调层支撑在 $Z$ 上，可知对每个 $f\in I$，
局部化 $M^\bullet_f$ 都无上同调。取 $x\in H^p(M^\bullet)$。
由上可找到 $n_i$，使得对每个 $i$，在 $H^p(M^\bullet)$ 中有
$f_i^{n_i}x=0$。令 $n=\sum n_i$，便可知 $I^n$ 消去 $x$。
因此 $H^p(M^\bullet)$ 是 $I$-幂挠的。由于环映射
$A\to\mathcal{O}_{X,\overline x}$ 平坦，并且
$I_{\overline x}=I\mathcal{O}_{X,\overline x}$，结论随即得到。
\end{proof}

\begin{lemma}
\label{spaces-pushouts-lemma-formal-glueing-on-closed}
设 $S$ 是概形。令 $f:Y\to X$ 为 $S$ 上代数空间的态射，
令 $Z\subset X$ 为闭子空间。假设 $f^{-1}Z\to Z$ 是同构，
且 $f$ 在 $f^{-1}Z$ 的每一点都平坦。
对 $D_\QCoh(\mathcal{O}_Y)$ 中任一支撑在 $|f^{-1}Z|$ 上的 $Q$，有
$Lf^*Rf_*Q=Q$。
\end{lemma}

\begin{proof}
在 $\overline y$ 处检验茎，以证明典范映射 $Lf^*Rf_*Q\to Q$ 是同构。
如果 $\overline y$ 不在 $f^{-1}Z$ 中，那么两端都为零，结论成立。
假设 $\overline y$ 的像 $\overline x$ 属于 $Z$。
由引理 \ref{spaces-pushouts-lemma-stalk-pushforward-with-support}，有
$Rf_*Q_{\overline x}=Q_{\overline y}$；又因为 $f$ 在 $\overline y$ 处平坦，故
$$
(Lf^*Rf_*Q)_{\overline{y}} =
(Rf_*Q)_{\overline{x}}
\otimes_{\mathcal{O}_{X, \overline{x}}}
\mathcal{O}_{Y, \overline{y}} =
Q_{\overline{y}} \otimes_{\mathcal{O}_{X, \overline{x}}}
\mathcal{O}_{Y, \overline{y}}
$$
因此只需验证典范映射
$$
Q_{\overline{y}} \otimes_{\mathcal{O}_{X, \overline{x}}}
\mathcal{O}_{Y, \overline{y}}
\longrightarrow Q_{\overline{y}}
$$
在导出范畴中是同构。令
$I_{\overline x}\subset\mathcal{O}_{X,\overline x}$ 为定义 $Z$ 的理想层的茎。
由于 $Z\to X$ 局部有限呈示，这个理想有限生成；
把引理 \ref{spaces-pushouts-lemma-stalk-of-module-with-support} 应用于 $Y$ 上的 $Q$，
可知 $Q_{\overline y}$ 的上同调群是
$I_{\overline y}=I_{\overline x}\mathcal{O}_{Y,\overline y}$-幂挠的。
因此它们也是 $I_{\overline x}$-幂挠的。环映射
$\mathcal{O}_{X,\overline x}\to\mathcal{O}_{Y,\overline y}$ 平坦；
又因为假设 $f^{-1}Z\to Z$ 是同构，所以分别模去
$I_{\overline x}$ 和 $I_{\overline y}$ 后，它诱导同构。因此，由
《代数续篇》引理 \ref{more-algebra-lemma-neighbourhood-isomorphism}，
可知
$Q_{\overline{y}} \otimes_{\mathcal{O}_{X, \overline{x}}}
\mathcal{O}_{Y, \overline{y}}$
的上同调模等于 $Q_{\overline y}$ 的上同调模，证明完成。
\end{proof}

\begin{situation}
\label{spaces-pushouts-situation-formal-glueing}
这里 $S$ 是基概形，$f:Y\to X$ 是 $S$ 上代数空间的拟紧且拟分离态射，
$Z\to X$ 是有限呈示的闭浸入。假设 $f^{-1}(Z)\to Z$ 是同构，
且 $f$ 在每一点 $x\in|f^{-1}Z|$ 都平坦。
置 $U=X\setminus Z$ 和 $V=Y\setminus f^{-1}(Z)$。图表为
$$
\xymatrix{
V \ar[r]_{j'} \ar[d]_{f|_V} & Y \ar[d]^f \\
U \ar[r]^j & X
}
$$
\end{situation}

\noindent
在情形 \ref{spaces-pushouts-situation-formal-glueing} 中，定义
$\textit{QCoh}(Y\to X,Z)$ 为三元组
$(\mathcal{H},\mathcal{G},\varphi)$ 的范畴，其中
$\mathcal{H}$ 是拟凝聚 $\mathcal{O}_U$-模层，
$\mathcal{G}$ 是拟凝聚 $\mathcal{O}_Y$-模层，而
$\varphi:f^*\mathcal{H}\to\mathcal{G}|_V$ 是 $\mathcal{O}_V$-模同构。
存在典范函子
\begin{equation}
\label{spaces-pushouts-equation-formal-glueing-modules}
\QCoh(\mathcal{O}_X) \longrightarrow \textit{QCoh}(Y \to X, Z)
\end{equation}
它把 $\mathcal{F}$ 映为系统 $(\mathcal{F}|_U,f^*\mathcal{F},can)$。
类比仿射情形下给出的证明，我们构造反方向的函子。
对对象 $(\mathcal{H},\mathcal{G},\varphi)$，赋予 $\mathcal{O}_X$-模
\begin{equation}
\label{spaces-pushouts-equation-reverse}
\Ker(j_*\mathcal{H} \oplus f_*\mathcal{G} \to (f \circ j')_*\mathcal{G}|_V)
\end{equation}
注意，由于 $Z\to X$ 有限呈示，$j$ 和 $j'$ 都是拟紧态射。
因此，$f_*$、$j_*$ 和 $(f\circ j')_*$ 把拟凝聚模映为拟凝聚模
（《空间的态射》引理 \ref{spaces-morphisms-lemma-pushforward}）。
故模 (\ref{spaces-pushouts-equation-reverse}) 拟凝聚。

\begin{lemma}
\label{spaces-pushouts-lemma-adjoint}
在情形 \ref{spaces-pushouts-situation-formal-glueing} 中，函子
(\ref{spaces-pushouts-equation-reverse}) 是函子
(\ref{spaces-pushouts-equation-formal-glueing-modules}) 的右伴随。
\end{lemma}

\begin{proof}
这很容易从 $f^*$ 与 $f_*$、$j^*$ 与 $j_*$ 的伴随性推出。略去细节。
\end{proof}

\begin{lemma}
\label{spaces-pushouts-lemma-reverse-commutes-with-flat-base-change}
在情形 \ref{spaces-pushouts-situation-formal-glueing} 中，令 $X'\to X$ 为代数空间的平坦态射。
置 $Z'=X'\times_XZ$ 和 $Y'=X'\times_XY$。
拉回 $\QCoh(\mathcal{O}_X)\to\QCoh(\mathcal{O}_{X'})$ 和
$\QCoh(Y\to X,Z)\to\QCoh(Y'\to X',Z')$ 与函子
(\ref{spaces-pushouts-equation-reverse}) 和
\ref{spaces-pushouts-equation-formal-glueing-modules}) 相容。% P10-E0023
\end{lemma}

\begin{proof}
这是因为拉回与拉回交换，而平坦拉回与沿拟紧且拟分离态射的正像交换；
见《空间的上同调》引理
\ref{spaces-cohomology-lemma-flat-base-change-cohomology}.
\end{proof}

\begin{proposition}
\label{spaces-pushouts-proposition-formal-glueing-modules}
在情形 \ref{spaces-pushouts-situation-formal-glueing} 中，函子
(\ref{spaces-pushouts-equation-formal-glueing-modules}) 是等价，
其拟逆由 (\ref{spaces-pushouts-equation-reverse}) 给出。
\end{proposition}

\begin{proof}
先处理 $X$ 和 $Y$ 都是仿射概形且态射 $f$ 平坦的特殊情形。
设 $X=\Spec(R)$、$Y=\Spec(S)$。那么 $f$ 对应于平坦环映射 $R\to S$。
此外，$Z\subset X$ 由有限生成理想 $I\subset R$ 截出。
选取生成元 $f_1,\ldots,f_t\in I$。
由用模描述拟凝聚模的方式
（《概形》第 \ref{schemes-section-quasi-coherent-affine} 节），
可知范畴 $\textit{QCoh}(Y\to X,Z)$ 典范等价于《代数续篇》注
\ref{more-algebra-remark-glueing-data} 中的范畴
$\text{Glue}(R\to S,f_1,\ldots,f_t)$，并且函子
(\ref{spaces-pushouts-equation-formal-glueing-modules}) 和 (\ref{spaces-pushouts-equation-reverse})
分别对应于函子 $\text{Can}$ 和 $H^0$。
因此，在这种情形下，结论来自《代数续篇》命题
\ref{more-algebra-proposition-equivalence}。

\medskip\noindent
回到一般情形。令 $\mathcal{F}$ 为 $X$ 上的拟凝聚模。我们将证明
$$
\alpha :
\mathcal{F}
\longrightarrow
\Ker\left(j_*\mathcal{F}|_U \oplus f_*f^*\mathcal{F} \to
(f \circ j')_*f^*\mathcal{F}|_V\right)
$$
是同构。令 $(\mathcal{H},\mathcal{G},\varphi)$ 为
$\QCoh(Y\to X,Z)$ 的一个对象。我们将证明
$$
\beta :
f^*\Ker\left(
j_*\mathcal{H} \oplus f_*\mathcal{G} \to (f \circ j')_*\mathcal{G}|_V
\right)
\longrightarrow
\mathcal{G}
$$
以及
$$
\gamma :
j^*\Ker\left(
j_*\mathcal{H} \oplus f_*\mathcal{G} \to (f \circ j')_*\mathcal{G}|_V
\right)
\longrightarrow
\mathcal{H}
$$
都是同构。为证明这些断言，只需考察茎。
令 $\overline y$ 为 $Y$ 的几何点，它映到 $X$ 的几何点 $\overline x$。

\medskip\noindent
固定 $\QCoh(Y\to X,Z)$ 的一个对象
$(\mathcal{H},\mathcal{G},\varphi)$。由引理
\ref{spaces-pushouts-lemma-stalk-formal-glueing} 和一次图追逐（略去），典范映射
$$
\Ker(j_*\mathcal{H} \oplus f_*\mathcal{G} \to
(f \circ j')_*\mathcal{G}|_V)_{\overline{x}}
\longrightarrow
\Ker(
j_*\mathcal{H}_{\overline{x}} \oplus \mathcal{G}_{\overline{y}}
\to
j'_*\mathcal{G}_{\overline{y}}
)
$$
是同构。

\medskip\noindent
特别地，如果 $\overline y$ 是 $V$ 的几何点，那么
$j'_*\mathcal{G}_{\overline y}=\mathcal{G}_{\overline y}$，
从而这个核等于 $\mathcal{H}_{\overline x}$。
这很容易蕴含在此情形下 $\alpha_{\overline x}$、
$\beta_{\overline x}$ 和 $\beta_{\overline y}$ 都是同构。% P10-E0024

\medskip\noindent
接着，假设 $\overline y$ 是 $f^{-1}Z$ 的一点。
分别令 $I_{\overline x}\subset\mathcal{O}_{X,\overline x}$、
$I_{\overline y}\subset\mathcal{O}_{Y,\overline y}$ 为截出
$Z$、$f^{-1}Z$ 的理想的茎。那么 $I_{\overline x}$ 是有限生成理想，
$I_{\overline y}=I_{\overline x}\mathcal{O}_{Y,\overline y}$，而
$\mathcal{O}_{X,\overline x}\to\mathcal{O}_{Y,\overline y}$
是平坦局部同态，并诱导同构
$\mathcal{O}_{X,\overline x}/I_{\overline x}=
\mathcal{O}_{Y,\overline y}/I_{\overline y}$。
此时可以利用范畴图表进行自举：
$$
\xymatrix{
\QCoh(\mathcal{O}_X) \ar[r]_-{(\ref{spaces-pushouts-equation-formal-glueing-modules})} \ar[d] &
\QCoh(Y \to X, Z) \ar[d] \ar@/_2pc/[l]^{(\ref{spaces-pushouts-equation-reverse})} \\
\text{Mod}_{\mathcal{O}_{X, \overline{x}}} \ar[r]^-{\text{Can}} &
\text{Glue}(\mathcal{O}_{X, \overline{x}} \to \mathcal{O}_{Y, \overline{y}},
f_1, \ldots, f_t) \ar@/^2pc/[l]_{H^0}
}
$$
也就是说，如证明第一段中那样，我们等同
$$
\text{Glue}(\mathcal{O}_{X, \overline{x}} \to \mathcal{O}_{Y, \overline{y}},
f_1, \ldots, f_t)
=
\QCoh(\Spec(\mathcal{O}_{Y, \overline{y}}) \to
\Spec(\mathcal{O}_{X, \overline{x}}), V(I_{\overline{x}}))
$$
右侧竖直函子由拉回给出，显然内方形交换。
把证明第三段中对核之茎的计算与引理
\ref{spaces-pushouts-lemma-stalk-of-pushforward} 结合起来，可知外方形
（使用弯曲箭头）交换。因此，可用证明第一段中已经处理的
仿射概形平坦态射情形得出结论。
\end{proof}

\begin{lemma}
\label{spaces-pushouts-lemma-derived-equivalent}
在情形 \ref{spaces-pushouts-situation-formal-glueing} 中，函子 $Rf_*$ 诱导
$D_{\QCoh,|f^{-1}Z|}(\mathcal{O}_Y)$ 与
$D_{\QCoh,|Z|}(\mathcal{O}_X)$ 之间的等价，其拟逆由 $Lf^*$ 给出。
\end{lemma}

\begin{proof}
由于 $f$ 拟紧且拟分离，$Rf_*$ 定义从
$D_{\QCoh,|f^{-1}Z|}(\mathcal{O}_Y)$ 到
$D_{\QCoh,|Z|}(\mathcal{O}_X)$ 的函子；见《空间的导出范畴》引理
\ref{spaces-perfect-lemma-quasi-coherence-direct-image}.
由《空间的导出范畴》引理
\ref{spaces-perfect-lemma-quasi-coherence-pullback}
，可知 $Lf^*$ 把 $D_{\QCoh,|Z|}(\mathcal{O}_X)$ 映入
$D_{\QCoh,|f^{-1}Z|}(\mathcal{O}_Y)$。
在引理 \ref{spaces-pushouts-lemma-formal-glueing-on-closed} 中已经看到，
对 $D_{\QCoh,|f^{-1}Z|}(\mathcal{O}_Y)$ 中的 $Q$，有
$Lf^*Rf_*Q=Q$。由《导出范畴》引理
\ref{derived-lemma-fully-faithful-adjoint-kernel-zero}
的对偶，为完成证明，只需说明：对
$D_{\QCoh,|Z|}(\mathcal{O}_X)$ 中的 $K$，$Lf^*K=0$ 蕴含 $K=0$。
这由 $f$ 在 $f^{-1}Z$ 的所有点都平坦以及
$f^{-1}Z\to Z$ 为满射这两个事实推出。
\end{proof}

\begin{lemma}
\label{spaces-pushouts-lemma-dominate-by-fpqc-covering}
在情形 \ref{spaces-pushouts-situation-formal-glueing} 中，存在一个加细族
$\{U\to X,Y\to X\}$ 的 fpqc 覆盖 $\{X_i\to X\}_{i\in I}$。
\end{lemma}

\begin{proof}
关于 fpqc 覆盖的定义和一般性质，参见《拓扑》第
\ref{topologies-section-fpqc} 节。特别地，可以先选取一个平展覆盖
$\{X_i\to X\}$，其中各 $X_i$ 仿射；把 $Y$、$Z$、$U$
基变换到每个 $X_i$，便约化到 $X$ 仿射的情形。
此时 $U$ 拟紧，因而是仿射开集的有限并
$U=U_1\cup\ldots\cup U_n$。于是 $Z$ 拟紧，故 $f^{-1}Z$ 也拟紧。
因此可以选取仿射概形 $W$ 和平展态射 $h:W\to Y$，使
$h^{-1}f^{-1}Z\to f^{-1}Z$ 为满射。
设 $W=\Spec(B)$ 且 $h^{-1}f^{-1}Z=V(J)$，其中 $J\subset B$
是有限型理想。由《Pro-平展上同调》引理
\ref{proetale-lemma-localization}，存在局部化 $B\to B'$，
使 $\Spec(B')$ 的点恰好对应于 $W=\Spec(B)$ 中特化到
$h^{-1}f^{-1}Z=V(J)$ 的点。由于假设 $f:Y\to X$ 在
$f^{-1}Z$ 的所有点都平坦，所以复合
$\Spec(B')\to\Spec(B)=W\to Y\to X$ 平坦。于是由《拓扑》引理
\ref{topologies-lemma-recognize-fpqc-covering}，
$\{\Spec(B')\to X,U_1\to X,\ldots,U_n\to X\}$ 是 fpqc 覆盖。
\end{proof}




\section{代数空间的形式粘合}
\label{spaces-pushouts-section-formal-glueing-spaces}

\noindent
在情形 \ref{spaces-pushouts-situation-formal-glueing} 中，考虑由如下形状的
$S$ 上代数空间交换图表组成的范畴 $\textit{Spaces}(Y\to X,Z)$：
$$
\xymatrix{
U' \ar[d] & V' \ar[l] \ar[d] \ar[r] & Y' \ar[d] \\
U & V \ar[l] \ar[r] & Y
}
$$
其中两个方形都是笛卡尔的。存在典范函子
\begin{equation}
\label{spaces-pushouts-equation-formal-glueing-spaces}
\textit{Spaces}/X \longrightarrow \textit{Spaces}(Y \to X, Z)
\end{equation}
它把 $X'\to X$ 映为诸态射
$U\times_XX'\leftarrow V\times_XX'\rightarrow Y\times_XX'$。

\begin{lemma}
\label{spaces-pushouts-lemma-equivalence-on-affine}
在情形 \ref{spaces-pushouts-situation-formal-glueing} 中，函子
(\ref{spaces-pushouts-equation-formal-glueing-spaces}) 限制为下列等价：
\begin{enumerate}
\item 从 $X$ 上仿射的代数空间范畴，到
$\textit{Spaces}(Y\to X,Z)$ 的满子范畴；后者由满足
$U'\to U$、$V'\to V$、$Y'\to Y$ 都仿射的
$(U'\leftarrow V'\rightarrow Y')$ 组成；
\item 从闭浸入 $X'\to X$ 的范畴，到
$\textit{Spaces}(Y\to X,Z)$ 的满子范畴；后者由满足
$U'\to U$、$V'\to V$、$Y'\to Y$ 都是闭浸入的
$(U'\leftarrow V'\rightarrow Y')$ 组成；以及
\item 对有限态射成立与 (2) 相同的陈述。
\end{enumerate}
\end{lemma}

\begin{proof}
在 $X$ 上仿射的代数空间范畴，等价于拟凝聚 $\mathcal{O}_X$-代数层
$\mathcal{A}$ 的范畴。$\textit{Spaces}(Y\to X,Z)$ 中由满足
$U'\to U$、$V'\to V$、$Y'\to Y$ 都仿射的
$(U'\leftarrow V'\rightarrow Y')$ 所组成的满子范畴，
等价于 $\QCoh(Y\to X,Z)$ 中代数对象的范畴。
在两种情形下，这都来自《空间的态射》引理
\ref{spaces-morphisms-lemma-affine-equivalence-algebras}
；拟逆由相对谱构造给出（《空间的态射》定义
\ref{spaces-morphisms-definition-relative-spec}），
而相对谱构造与任意基变换交换。因此本引理的 (1) 来自命题
\ref{spaces-pushouts-proposition-formal-glueing-modules}。

\medskip\noindent
第 (2) 部分的全忠实性由 (1) 推出。% P10-E0025
对于本质满射性，由 (1)，问题约化为证明：$X'\to X$ 是闭浸入，
当且仅当 $U\times_XX'\to U$ 和 $Y\times_XX'\to Y$ 都是闭浸入。
由引理 \ref{spaces-pushouts-lemma-dominate-by-fpqc-covering}，
$\{U\to X,Y\to X\}$ 可由某个 fpqc 覆盖加细。因此结论来自
《空间上的下降》引理
\ref{spaces-descent-lemma-descending-property-closed-immersion}.

\medskip\noindent
对 (3)，使用证明 (2) 的论证以及《空间上的下降》引理
\ref{spaces-descent-lemma-descending-property-finite}.
\end{proof}

\begin{lemma}
\label{spaces-pushouts-lemma-reflects-isomorphisms}
在情形 \ref{spaces-pushouts-situation-formal-glueing} 中，函子
(\ref{spaces-pushouts-equation-formal-glueing-spaces}) 反映同构。
\end{lemma}

\begin{proof}
通过关于基变换的形式论证，这约化为如下问题：
若代数空间态射 $a:X'\to X$ 使 $U\times_XX'\to U$ 和
$Y\times_XX'\to Y$ 都是同构，则该态射是同构。
由引理 \ref{spaces-pushouts-lemma-dominate-by-fpqc-covering}，
族 $\{U\to X,Y\to X\}$ 可由某个 fpqc 覆盖加细。
因此结论来自《空间上的下降》引理
\ref{spaces-descent-lemma-descending-property-isomorphism}.
\end{proof}

\begin{lemma}
\label{spaces-pushouts-lemma-fully-faithful-on-separated}
在情形 \ref{spaces-pushouts-situation-formal-glueing} 中，函子
(\ref{spaces-pushouts-equation-formal-glueing-spaces}) 在 $X$ 上分离的代数空间上全忠实。
更确切地说，只要 $X'_2\to X$ 分离，它就诱导双射
$$
\Mor_X(X'_1, X'_2)
\longrightarrow
\Mor_{\textit{Spaces}(Y \to X, Z)}(F(X'_1), F(X'_2))
$$
。
\end{lemma}

\begin{proof}
由于 $X'_2\to X$ 分离，$X$ 上态射 $X'_1\to X'_2$ 的图
$i:X'_1\to X'_1\times_XX'_2$ 是闭浸入；见《空间的态射》引理
\ref{spaces-morphisms-lemma-semi-diagonal}。
此外，闭浸入 $i:T\to X'_1\times_XX'_2$ 是某个态射的图，
当且仅当 $\text{pr}_1\circ i$ 是同构。相同的断言也适用于
\begin{enumerate}
\item $U$ 上态射 $U\times_XX'_1\to U\times_XX'_2$ 的图；
\item $V$ 上态射 $V\times_XX'_1\to V\times_XX'_2$ 的图；以及
\item $Y$ 上态射 $Y\times_XX'_1\to Y\times_XX'_2$ 的图。
\end{enumerate}
此外，如果 (1)、(2)、(3) 中的态射相互配合，组成范畴
$\textit{Spaces}(Y\to X,Z)$ 中的一个态射，那么这些图也相互配合，
给出下列范畴的一个对象：
$\textit{Spaces}(Y \times_X (X'_1 \times_X X'_2) \to X'_1 \times_X X'_2,
Z \times_X (X'_1 \times_X X'_2))$
其三个态射都是闭浸入。应用引理
\ref{spaces-pushouts-lemma-equivalence-on-affine} 和
\ref{spaces-pushouts-lemma-reflects-isomorphisms} 即完成证明。
\end{proof}









\section{粘合与 Beauville--Laszlo 定理}
\label{spaces-pushouts-section-glueing-beauville-laszlo}

\noindent
令 $R\to R'$ 为环同态，$f\in R$ 为一个元素，使得
$$
0 \to R \to R_f \oplus R' \to R'_f \to 0
$$
是短正合列。这蕴含对所有 $n$ 都有 $R/f^nR\cong R'/f^nR'$，
并且 $(R\to R',f)$ 是《代数续篇》第
\ref{more-algebra-section-beauville-laszlo} 节意义下的粘合对。
置 $X=\Spec(R)$、$U=\Spec(R_f)$、$X'=\Spec(R')$、
$U'=\Spec(R'_f)$。图表为
$$
\xymatrix{
U' \ar[r] \ar[d] & X' \ar[d] \\
U \ar[r] & X
}
$$
在这种情形下，可以考虑范畴
$\textit{Spaces}(U\leftarrow U'\to X')$，其对象为交换图表
$$
\xymatrix{
V \ar[d] & V' \ar[l] \ar[d] \ar[r] & Y' \ar[d] \\
U & U' \ar[l] \ar[r] & X'
}
$$
，其中各项为代数空间，两个方形都是笛卡尔的，态射按显然方式定义。% P10-E0026
这个范畴的对象记作 $(V,V',Y')$，省略记号中的箭头。存在函子
\begin{equation}
\label{spaces-pushouts-equation-beauville-laszlo-glueing-spaces}
\textit{Spaces}/X
\longrightarrow
\textit{Spaces}(U \leftarrow U' \to X')
\end{equation}
由基变换给出：$Y\mapsto(U\times_XY,U'\times_XY,X'\times_XY)$。

\medskip\noindent
我们已在《代数续篇》第 \ref{more-algebra-section-beauville-laszlo} 节看到，
并非每个 $R$-模 $M$ 都能从其粘合数据恢复。类似地，函子
(\ref{spaces-pushouts-equation-beauville-laszlo-glueing-spaces})
在全体 $X$ 上空间的范畴上并不全忠实。
为挑出 $X$ 上代数空间的一个适当子范畴，需要下面的引理。

\begin{lemma}
\label{spaces-pushouts-lemma-glueable}
令 $(R\to R',f)$ 为如上的粘合对，$Y$ 为 $X$ 上的代数空间。
下列条件等价：
\begin{enumerate}
\item 存在平展覆盖 $\{Y_i\to Y\}_{i\in I}$，其中各 $Y_i$ 仿射，
且 $\Gamma(Y_i,\mathcal{O}_{Y_i})$ 作为 $R$-模可粘合；
\item 对每个平展态射 $W\to Y$，只要 $W$ 仿射，
$\Gamma(W,\mathcal{O}_W)$ 作为 $R$-模就可粘合。
\end{enumerate}
\end{lemma}

\begin{proof}
显然 (2) 蕴含 (1)。假设 $\{Y_i\to Y\}$ 如 (1) 所述，
且 $W\to Y$ 如 (2) 所述。那么
$\{Y_i\times_YW\to W\}_{i\in I}$ 是平展覆盖，
可由一个各 $W_j$ 仿射的平展覆盖
$\{W_j\to W\}_{j=1,\ldots,m}$ 加细
（《拓扑》引理 \ref{topologies-lemma-etale-affine}）。
因此，为完成证明，只需证明下列三个代数断言：
\begin{enumerate}
\item 如果 $R\to A\to B$ 是环映射，$A\to B$ 平展，
且 $A$ 作为 $R$-模可粘合，那么 $B$ 作为 $R$-模也可粘合；
\item 可粘合 $R$-模的有限积可粘合；
\item 如果 $R\to A\to B$ 是环映射，$A\to B$ 忠实平展，
且 $B$ 作为 $R$-模可粘合，那么 $A$ 作为 $R$-模也可粘合。
\end{enumerate}
具体地，第一个断言蕴含 $\Gamma(W_j,\mathcal{O}_{W_j})$
是可粘合 $R$-模，第二个断言蕴含
$\prod\Gamma(W_j,\mathcal{O}_{W_j})$ 是可粘合 $R$-模，
第三个断言则蕴含 $\Gamma(W,\mathcal{O}_W)$ 是可粘合 $R$-模。

\medskip\noindent
考虑平展 $R$-代数同态 $A\to B$。置
$A'=A\otimes_RR'$ 和 $B'=B\otimes_RR'=A'\otimes_AB$。
此时，断言 (1) 和 (3) 来自下列事实：
(a) $A$（分别地，$B$）可粘合，当且仅当序列
$$
0 \to A \to A_f \oplus A' \to A'_f \to 0,
\quad\text{分别}\quad
0 \to B \to B_f \oplus B' \to B'_f \to 0,
$$
正合；(b) 第二个序列由函子 $-\otimes_AB$ 作用于第一个序列得到；
(c) $A\to B$（忠实）平坦。略去对 (2) 的证明。
\end{proof}

\noindent
令 $(R\to R',f)$ 为如上的粘合对。如果引理
\ref{spaces-pushouts-lemma-glueable} 的等价条件得到满足，
就称 $X=\Spec(R)$ 上的代数空间 $Y$
{\it 对 $(R\to R',f)$ 可粘合}。

\begin{lemma}
\label{spaces-pushouts-lemma-glueing-affines}
令 $(R\to R',f)$ 为如上的粘合对。函子
(\ref{spaces-pushouts-equation-beauville-laszlo-glueing-spaces}) 限制为下列两个范畴之间的等价：
对 $(R\to R',f)$ 可粘合的仿射 $Y/X$ 的范畴；以及
$\textit{Spaces}(U\leftarrow U'\to X')$ 中由 $V$、$V'$、$Y'$
都仿射的对象 $(V,V',Y')$ 所组成的满子范畴。
\end{lemma}

\begin{proof}
令 $(V,V',Y')$ 为 $\textit{Spaces}(U\leftarrow U'\to X')$
的一个对象，并假设 $V$、$V'$、$Y'$ 都仿射。
写成 $V=\Spec(A_1)$ 和 $Y'=\Spec(A')$。
由范畴 $\textit{Spaces}(U\leftarrow U'\to X')$ 的定义，
可知 $V'$ 既是 $A_1\otimes_{R_f}R'_f=A_1\otimes_RR'$ 的谱，
也是 $A'_f$ 的谱。因此得到 $R'_f$-代数同构
$\varphi:A'_f\to A_1\otimes_RR'$。
由《代数续篇》定理 \ref{more-algebra-theorem-BL-glueing}，
存在唯一的可粘合 $R$-模 $A$，以及与 $\varphi$ 相容的模同构
$A_f\to A_1$ 和 $A\otimes_RR'\to A'$。由于序列
$$
0 \to A \to A_1 \oplus A' \to A'_f \to 0
$$
短正合，$A_1$ 和 $A'$ 上的乘法在 $A$ 上定义唯一的 $R$-代数结构，
使映射 $A\to A_1$ 和 $A\to A'$ 是环同态。
略去如下验证：此构造给出函子
(\ref{spaces-pushouts-equation-beauville-laszlo-glueing-spaces})
限制到本引理所述子范畴后的拟逆。
\end{proof}

\begin{lemma}
\label{spaces-pushouts-lemma-glueing-affines-etale}
令 $P$ 为态射的下列性质之一：``有限''、``闭浸入''、``平坦''、
``有限型''、``平坦且有限呈示''、``平展''。
在引理 \ref{spaces-pushouts-lemma-glueing-affines} 的等价下，
具有性质 $P$ 的态射对应于各分量都具有性质 $P$ 的三元组态射。
\end{lemma}

\begin{proof}
令 $P'$ 为环同态的下列性质之一：``有限''、``满射''、``平坦''、
``有限型''、``平坦且有限呈示''、``平展''。
把陈述翻译为代数语言，就是：如果 $A\to B$ 是 $R$-代数同态，
且 $A$、$B$ 对 $(R\to R',f)$ 可粘合，那么
$A_f\to B_f$ 和 $A\otimes_RR'\to B\otimes_RR'$ 具有性质 $P'$，
当且仅当 $A\to B$ 具有性质 $P'$。

\medskip\noindent
由《代数续篇》引理 \ref{more-algebra-lemma-faithful-descent} 和
\ref{more-algebra-lemma-BL-flat}，当 $P'$ 为``有限''或``平坦''时，
这个代数断言成立。

\medskip\noindent
如果 $A_f\to B_f$ 和 $A\otimes_RR'\to B\otimes_RR'$ 都是满射，
那么 $N=B/A$ 是满足 $N_f=0$ 且 $N\otimes_RR'=0$ 的 $R$-模，
因而由《代数续篇》引理 \ref{more-algebra-lemma-BL-faithful}，该模为零。
所以 $A\to B$ 是满射。

\medskip\noindent
如果 $A_f\to B_f$ 和 $A\otimes_RR'\to B\otimes_RR'$ 都有限型，
则可选取 $A$-代数同态 $A[x_1,\ldots,x_n]\to B$，使
$A_f[x_1,\ldots,x_n]\to B_f$ 和
$(A\otimes_RR')[x_1,\ldots,x_n]\to B\otimes_RR'$ 都是满射
（略去一个小细节）。由前面的结果，可知
$A[x_1,\ldots,x_n]\to B$ 是满射。因此 $A\to B$ 有限型。

\medskip\noindent
如果 $A_f\to B_f$ 和 $A\otimes_RR'\to B\otimes_RR'$ 都平坦且有限呈示，
那么由已经证明的结论，$A\to B$ 平坦且有限型。
选取满射 $A[x_1,\ldots,x_n]\to B$，并把核记作 $I$。
由 $B$ 在 $A$ 上的平坦性，可知 $I_f$ 是
$A_f[x_1,\ldots,x_n]\to B_f$ 的核，而 $I\otimes_RR'$ 是
$A\otimes_RR'[x_1,\ldots,x_n]\to B\otimes_RR'$ 的核。
因此，$I_f$ 是有限 $A_f[x_1,\ldots,x_n]$-模，
$I\otimes_RR'$ 是有限 $(A\otimes_RR')[x_1,\ldots,x_n]$-模。
把《代数续篇》引理 \ref{more-algebra-lemma-faithful-descent}
应用于视作 $A[x_1,\ldots,x_n]$-模的 $I$，可知 $I$ 是有限生成理想，
故 $A\to B$ 平坦且有限呈示。

\medskip\noindent
如果 $A_f\to B_f$ 和 $A\otimes_RR'\to B\otimes_RR'$ 都平展，
那么由已经证明的结论，$A\to B$ 平坦且有限呈示。
由于 $\Spec(B)\to\Spec(A)$ 的纤维同构于
$\Spec(B_f)\to\Spec(A_f)$ 或
$\Spec(B/fB)\to\Spec(A/fA)$ 的纤维，可知 $A\to B$ 非分歧；
见《态射》引理 \ref{morphisms-lemma-unramified-over-field} 和
\ref{morphisms-lemma-unramified-etale-fibres}。
例如再由《态射》引理 \ref{morphisms-lemma-flat-unramified-etale}，
得到 $A\to B$ 平展。
\end{proof}

\begin{lemma}
\label{spaces-pushouts-lemma-glueing-f}
令 $(R\to R',f)$ 为如上的粘合对。函子
(\ref{spaces-pushouts-equation-beauville-laszlo-glueing-spaces})
在对 $(R\to R',f)$ 可粘合的代数空间 $Y/X$ 所组成的满子范畴上忠实。
\end{lemma}

\begin{proof}
令 $f,g:Y\to Z$ 为 $X$ 上代数空间的两个态射，
其中 $Y$、$Z$ 对 $(R\to R',f)$ 可粘合；并假设 $f$、$g$ 被映为范畴
$\textit{Spaces}(U\leftarrow U'\to X')$ 中的同一个态射。
必须证明 $f$ 与 $g$ 的等化子 $E\to Y$ 是同构。
在 $Y$ 上平展局部地工作，可以假设 $Y$ 是仿射概形。
那么 $E$ 是概形，且态射 $E\to Y$ 是单态射并局部拟有限；
见《空间的态射》引理 \ref{spaces-morphisms-lemma-properties-diagonal}。
此外，$E\to Y$ 到 $U$ 和 $X'$ 的基变换都是同构。
由于 $Y$ 是仿射开集 $V=U\times_XY$ 与仿射闭集
$V(f)\times_XY$ 的不交并，可知 $E$ 是它们的同构逆像的不交并。
特别地，$E$ 拟紧。由 Zariski 主定理（《态射续篇》引理
\ref{more-morphisms-lemma-quasi-finite-separated-pass-through-finite}），
可知 $E$ 拟仿射。
置 $B=\Gamma(E,\mathcal{O}_E)$ 和 $A=\Gamma(Y,\mathcal{O}_Y)$，
从而有 $R$-代数同态 $A\to B$。由于 $E\to Y$ 基变换到 $U$ 和
$X'$ 后都是同构，得到环映射 $B\to A_f$ 和
$B\to A\otimes_RR'$，它们作为到 $A\otimes_RR'_f$ 的映射相合。
由于 $A$ 对 $(R\to R',f)$ 可粘合，得到环映射 $B\to A$，
它是映射 $A\to B$ 的左逆。相应的态射
$Y=\Spec(A)\to\Spec(B)$ 逐点映入开子概形 $E\subset\Spec(B)$，
因为基变换到 $U$ 和 $X'$ 后确是如此。因此得到 $Y$ 上态射 $Y\to E$。
由于 $E\to Y$ 是单态射，% P10-E0027
可知 $Y\to E$ 按所需是同构。
\end{proof}

\begin{lemma}
\label{spaces-pushouts-lemma-glueing-ff}
令 $(R\to R',f)$ 为如上的粘合对。函子
(\ref{spaces-pushouts-equation-beauville-laszlo-glueing-spaces}) 在由如下代数空间
$Y/X$ 所组成的满子范畴上全忠实：(a) 对 $(R\to R',f)$ 可粘合；
(b) 对角态射 $Y\to Y\times_XY$ 仿射。
\end{lemma}

\begin{proof}
令 $Y,Z$ 为 $X$ 上的两个代数空间，二者都对 $(R\to R',f)$ 可粘合，
并假设 $Z$ 的对角态射仿射。令
$a:U\times_XY\to U\times_XZ$ 为 $U$ 上态射，
$b:X'\times_XY\to X'\times_XZ$ 为 $X'$ 上态射；
假设它们在 $U'$ 上诱导同一个态射
$c:U'\times_XY\to U'\times_XZ$。
我们要构造 $X$ 上态射 $f:Y\to Z$，使其基变换到 $U$、$X'$ 后
分别给出 $a$、$b$。由引理 \ref{spaces-pushouts-lemma-glueing-f} 的忠实性，
只需在 $Y$ 上平展局部地构造 $f$（略去细节）。
因此可以并且确实假设 $Y$ 仿射。

\medskip\noindent
令 $y\in|Y|$ 为一点。如果 $y$ 映入开集 $U\subset X$，
那么 $U\times_XY$ 是 $Y$ 的开集，态射 $f$ 已在其上定义
（直接取 $a$ 即可）。因此可以假设 $y$ 映入 $X$ 的闭子集 $V(f)$。
由于 $R/fR=R'/fR'$，存在唯一的点
$y'\in|X'\times_XY|$ 映到 $y$。
令 $z'=b(y')\in|X'\times_XZ|$，并令 $z\in|Z|$ 为 $y'$ 的像。
选取一个平展邻域 $(W,w)\to(Z,z)$，其中 $W$ 仿射。注意，
$$
(U \times_X W) \times_{U \times_X Z, a} (U \times_X Y),\quad
(U' \times_X W) \times_{U' \times_X Z, c} (U' \times_X Y),
$$
以及
$$
(X' \times_X W) \times_{X' \times_X Z, b} (X' \times_X Y)
$$
组成 $\textit{Spaces}(U\leftarrow U'\to X')$ 的一个对象，
且各部分都仿射（这里用到了 $Z$ 具有仿射对角态射）。
因此，由引理 \ref{spaces-pushouts-lemma-glueing-affines}，
存在唯一的、对 $(R\to R',f)$ 可粘合的仿射概形 $V$，使得
$$
(U \times_X V, U' \times_X V, X' \times_X V)
$$
就是上述三元组。由仿射情形的全忠实性
（引理 \ref{spaces-pushouts-lemma-glueing-affines}），得到唯一的态射 $V\to W$ 和
$V\to Y$，它们分别与上述构造中在 $U$、$X'$ 上的第一、第二投影态射相合。
由引理 \ref{spaces-pushouts-lemma-glueing-affines-etale}，态射 $V\to Y$ 平展。
为完成证明，只需说明存在一点 $v\in|V|$ 映到 $y$
（因为这样一来，$f$ 就在 $y$ 的一个平展邻域——即 $V$——上定义）。
存在唯一的点 $w'\in|X'\times_XW|$ 映到 $w$。
由唯一性，$w'$ 在映射 $|X'\times_XW|\to|X'\times_XZ|$ 下映到 $z'$。
于是考虑笛卡尔图表
$$
\xymatrix{
X' \times_X V \ar[r] \ar[d] & X' \times_X W \ar[d] \\
X' \times_X Y \ar[r] & X' \times_X Z
}
$$
由此可知存在一点 $v'\in|X'\times_XV|$ 映到 $y'$ 和 $w'$；
见《空间的性质》引理 \ref{spaces-properties-lemma-points-cartesian}。
当然，$v'$ 在 $|V|$ 中的像 $v$ 映到 $y$，证明完成。
\end{proof}

\begin{lemma}
\label{spaces-pushouts-lemma-glueing-quasi-affines}
令 $(R\to R',f)$ 为如上的粘合对。
$\textit{Spaces}(U\leftarrow U'\to X')$ 中任一满足
$V$、$V'$、$Y'$ 都拟仿射的对象 $(V,V',Y')$，
都同构于某个 $X$ 上分离代数空间 $Y$ 在函子
(\ref{spaces-pushouts-equation-beauville-laszlo-glueing-spaces}) 下的像。
\end{lemma}

\begin{proof}
如《性质》引理 \ref{properties-lemma-quasi-affine-presentation} 中那样，
选取 $n'$、$T'\to Y'$ 以及 $n_1$、$T_1\to V$。图表为
$$
\xymatrix{
& &
T_1 \times_V V' \times_Y T' \ar[ld] \ar[rd] \\
T_1 \ar[d] &
T_1 \times_V V' \ar[l] \ar[dr] & &
V' \times_{Y'} T' \ar[r] \ar[dl] &
T' \ar[d] \\
V & &
V' \ar[rr] \ar[ll] & &
Y'
}
$$
注意，$T_1\times_VV'$ 和 $V'\times_{Y'}T'$ 仿射
（事实上，态射 $V'\to V$ 和 $V'\to Y'$ 分别是仿射态射
$U'\to U$ 和 $U'\to X'$ 的基变换，因而仿射）。按构造，有
$$
\mathbf{A}^{n'}_{T_1 \times_V V'} \cong
T_1 \times_V V' \times_{Y'} T' \cong
\mathbf{A}^{n_1}_{V' \times_{Y'} T'}
$$
换言之，仿射概形 $\mathbf{A}^{n'}_{T_1}$ 和
$\mathbf{A}^{n_1}_{T'}$ 是一个三元组的组成部分，该三元组构成
$\textit{Spaces}(U\leftarrow U'\to X')$ 的仿射对象。
由引理 \ref{spaces-pushouts-lemma-glueing-affines}，存在仿射概形态射 $T\to X$
以及同构 $U\times_XT\cong\mathbf{A}^{n'}_{T_1}$ 和
$X'\times_XT\cong\mathbf{A}^{n_1}_{T'}$，它们与上述同构相容。
这些同构给出态射
$$
U \times_X T \longrightarrow V
\quad\text{和}\quad
X' \times_X T \longrightarrow Y'
$$
它们在 $n=n'+n_1$ 时满足《性质》引理
\ref{properties-lemma-quasi-affine-presentation} 的性质，
并且还在范畴 $\textit{Spaces}(U\leftarrow U'\to X')$ 中定义从三元组
$(U\times_XT,U'\times_XT,X'\times_XT)$ 到三元组
$(V,V',Y')$ 的态射。

\medskip\noindent
由引理 \ref{spaces-pushouts-lemma-glueing-affines}，存在仿射概形 $W$，
它在 $\textit{Spaces}(U\leftarrow U'\to X')$ 中的像同构于三元组
$$
((U \times_X T) \times_V (U \times_X T),
(U' \times_X T) \times_{V'} (U' \times_X T),
(X' \times_X T) \times_{Y'} (X' \times_X T))
$$
由这个构造的全忠实性，得到两个映射 $p_0,p_1:W\to T$，
它们到 $U,U',X'$ 的基变换都是投影态射。
由引理 \ref{spaces-pushouts-lemma-glueing-affines-etale}，态射 $p_0,p_1$
平坦且有限呈示，而态射 $(p_0,p_1):W\to T\times_XT$ 是闭浸入。
事实上，$W\to T\times_XT$ 是等价关系：由上面使用的诸引理，
可以在基变换到 $U$ 和 $X'$ 后检验对称性、自反性和传递性；
在那里这些性质是显然的（略去细节）。因此商层
$$
Y = T/W
$$
是代数空间；例如可用《自举》定理
\ref{bootstrap-theorem-final-bootstrap}。
由于显然 $Y/X$ 被映到三元组 $(V,V',Y')$。% P10-E0028
沿拟紧满平坦态射 $T\times_XT\to Y\times_XY$ 对对角态射
$\Delta:Y\to Y\times_XY$ 作基变换，得到闭浸入
$W\to T\times_XT$。因此，由《空间上的下降》引理
\ref{spaces-descent-lemma-descending-property-closed-immersion}，
$\Delta$ 是闭浸入。故代数空间 $Y$ 分离，证明完成。
\end{proof}








\section{余等化子与粘合}
\label{spaces-pushouts-section-coequalizer-glue}

\noindent
令 $X$ 为诺特代数空间，$Z\to X$ 为闭子空间，
$X'\to X$ 为沿 $Z$ 的爆破。本节说明，可以从 $X'$、$Z_n$ 和粘合数据
恢复 $X$；这里 $Z_n$ 是 $Z$ 在 $X$ 中的第 $n$ 阶无穷小邻域。

\begin{lemma}
\label{spaces-pushouts-lemma-coequalizer}
设 $S$ 是概形。令
$$
g : Y \longrightarrow X
$$
为 $S$ 上代数空间的态射。假设 $X$ 局部诺特，且 $g$ 固有。
令 $R=Y\times_XY$，其投影态射为 $t,s:R\to Y$。
在 $S$ 上代数空间范畴中，$s,t:R\to Y$ 存在余等化子 $X'$。此外，
\begin{enumerate}
\item 态射 $X'\to X$ 有限；
\item 态射 $Y\to X'$ 固有；
\item 态射 $Y\to X'$ 为满射；
\item 态射 $X'\to X$ 万有单射；
\item 如果 $g$ 是满射，则态射 $X'\to X$ 是万有同胚。
\end{enumerate}
\end{lemma}

\begin{proof}
记 $h:R\to X$ 为 $s$ 或 $t$ 与 $g$ 的复合。% P10-E0029
由《空间的态射》引理
\ref{spaces-morphisms-lemma-base-change-proper} 和
\ref{spaces-morphisms-lemma-composition-proper}.
，$h$ 固有。层
$$
g_*\mathcal{O}_Y
\quad\text{和}\quad
h_*\mathcal{O}_R
$$
由《空间的上同调》引理
\ref{spaces-cohomology-lemma-proper-pushforward-coherent}
可知是凝聚 $\mathcal{O}_X$-代数。$X$-态射 $s$、$t$ 诱导从前者到后者的
$\mathcal{O}_X$-代数映射 $s^\sharp,t^\sharp$。置
$$
\mathcal{A} = \text{等化子}\left(s^\sharp, t^\sharp :
g_*\mathcal{O}_Y \longrightarrow h_*\mathcal{O}_R\right)
$$
那么 $\mathcal{A}$ 是凝聚 $\mathcal{O}_X$-代数，可以定义
$$
X' = \underline{\Spec}_X(\mathcal{A})
$$
，如《空间的态射》定义 \ref{spaces-morphisms-definition-relative-spec}
中那样。由《空间的态射》注
\ref{spaces-morphisms-remark-factorization-quasi-compact-quasi-separated}
以及 $\underline{\Spec}$ 构造的函子性，存在分解
$$
Y \longrightarrow X' \longrightarrow X
$$
，且态射 $g':Y\to X'$ 使 $s$ 与 $t$ 相等。

\medskip\noindent
在证明 $X'$ 是 $s$ 与 $t$ 的余等化子之前，先证明 $Y\to X'$ 和
$X'\to X$ 具有所需性质。由于 $\mathcal{A}$ 是凝聚
$\mathcal{O}_X$-模，显然 $X'\to X$ 是代数空间的有限态射。这证明了 (1)。
由《空间的态射》引理
\ref{spaces-morphisms-lemma-universally-closed-permanence}，
态射 $Y\to X'$ 固有。这证明了 (2)。
以 $Y\to Y'\to X$ 表示 $g$ 的 Stein 分解，其中
$Y'=\underline{\Spec}_X(g_*\mathcal{O}_Y)$；见《空间的态射续篇》定理
\ref{spaces-more-morphisms-theorem-stein-factorization-Noetherian}.
当然，得到与上面研究的态射相配合的态射 $Y\to Y'\to X'\to X$。
由于 $\mathcal{O}_{X'}\subset g_*\mathcal{O}_Y$ 是有限扩张，
可知 $Y'\to X'$ 有限且满。略去一些细节；提示：使用《代数》引理
\ref{algebra-lemma-integral-overring-surjective}，并通过平展局部化约化到仿射情形。
由于 $Y\to Y'$ 为满射（且具有几何连通的纤维），可知 $Y\to X'$ 为满射。
这证明了 (3)。为证明 $X'\to X$ 万有单射，必须证明
$X'\to X'\times_XX'$ 为满射；见《空间的态射》定义
\ref{spaces-morphisms-definition-universally-injective}
和引理 \ref{spaces-morphisms-lemma-universally-injective}。
由于 $Y\to X'$ 为满射（见上），而且由《空间的态射》引理
\ref{spaces-morphisms-lemma-base-change-surjective} 和
\ref{spaces-morphisms-lemma-composition-surjective}
，满射的基变换和复合仍是满射，故
$Y\times_XY\to X'\times_XX'$ 为满射。然而，由于 $Y\to X'$
使 $s$ 与 $t$ 相等，可知 $Y\times_XY\to X'\times_XX'$ 通过
$X'\to X'\times_XX'$ 分解，因而后一个映射为满射。这证明了 (4)。
最后，如果 $g$ 为满射，那么由于 $g$ 通过 $X'\to X$ 分解，
可知 $X'\to X$ 为满射。满、万有单射的有限态射是万有同胚
（因为它万有双射且万有闭），故 (5) 得证。

\medskip\noindent
在证明的其余部分，说明 $Y\to X'$ 是 $s$ 与 $t$ 在 $S$ 上代数空间范畴中的
余等化子。注意，$X'$ 局部诺特
（《空间的态射》引理
\ref{spaces-morphisms-lemma-locally-finite-type-locally-noetherian}）。
还要注意，由于 $Y\to X'$ 使 $s$ 与 $t$ 相等，
$Y\times_{X'}Y\to Y\times_XY$ 是同构（这是一个范畴论陈述）。
因此，为证明 $Y\to X'$ 是 $s$ 与 $t$ 的余等化子，
可以并且确实假设 $X=X'$。换言之，$\mathcal{O}_X$ 是映射
$s^\sharp,t^\sharp:g_*\mathcal{O}_Y\to h_*\mathcal{O}_R$ 的等化子。

\medskip\noindent
令 $X_1\to X$ 为 $S$ 上代数空间的平坦态射，且 $X_1$ 局部诺特。
分别以 $g_1:Y_1\to X_1$、$h_1:R_1\to X_1$ 和
$s_1,t_1:R_1\to Y_1$ 表示 $g,h,s,t$ 到 $X_1$ 的基变换。
当然，$g_1$ 固有，且 $R_1=Y_1\times_{X_1}Y_1$。
由于拟凝聚模的正像满足平坦基变换（《空间的上同调》引理
\ref{spaces-cohomology-lemma-flat-base-change-cohomology}），
可知 $\mathcal{O}_{X_1}$ 是映射
$s_1^\sharp,t_1^\sharp:g_{1,*}\mathcal{O}_{Y_1}\to
h_{1,*}\mathcal{O}_{R_1}$ 的等化子。因此我们的所有假设都在此基变换下保持。

\medskip\noindent
现在验证引理 \ref{spaces-pushouts-lemma-colimit-check-etale-locally} 的条件 (1) 和 (2)。
条件 (1) 来自引理 \ref{spaces-pushouts-lemma-descend-etale-proper-surjective}
以及 $g$ 固有且满这一事实（因为 $X=X'$）。
为验证条件 (2)，由上面关于基变换的说明，
约化为下一段讨论并证明的断言。

\medskip\noindent
假设 $S=\Spec(A)$ 是仿射概形，$X=X'$ 是仿射概形，
而 $Z$ 是 $S$ 上的仿射概形。必须证明
$$
\Mor_S(X, Z) \longrightarrow
\text{等化子}(s, t : \Mor_S(Y, Z) \to \Mor_S(R, Z))
$$
是双射。然而，这从 $X=X'$ 这一事实即可看出；该事实蕴含
$\mathcal{O}_X$ 是映射
$s^\sharp,t^\sharp:g_*\mathcal{O}_Y\to h_*\mathcal{O}_R$ 的等化子，
进而蕴含
$$
\Gamma(X, \mathcal{O}_X) =
\text{等化子}\left(
s^\sharp, t^\sharp : \Gamma(Y, \mathcal{O}_Y) \to
\Gamma(R, \mathcal{O}_R)
\right)
$$
事实上，我们有
$$
\Mor_S(X, Z) = \Hom_A(\Gamma(Z, \mathcal{O}_Z), \Gamma(X, \mathcal{O}_X))
$$
，对 $Y$ 和 $R$ 也有类似等式；见《空间的性质》引理
\ref{spaces-properties-lemma-morphism-to-affine-scheme}。
\end{proof}

\noindent
下面将在如下情形中工作。

\begin{situation}
\label{spaces-pushouts-situation-coequalizer-glue}
设 $S$ 是概形，$X$ 是 $S$ 上局部诺特的代数空间。
令 $Z\to X$ 为闭浸入，$U\subset X$ 为其补开子空间。
最后，令 $f:X'\to X$ 为代数空间的固有态射，
使 $f^{-1}(U)\to U$ 是同构。
\end{situation}

\begin{lemma}
\label{spaces-pushouts-lemma-coequalizer-glue}
在情形 \ref{spaces-pushouts-situation-coequalizer-glue} 中，令 $Y=X'\amalg Z$，
$R=Y\times_XY$，投影为 $t,s:R\to Y$。
在 $S$ 上代数空间的范畴中，$s,t:R\to Y$ 存在余等化子 $X_1$。
态射 $X_1\to X$ 是有限万有同胚，在 $U$ 上为同构，
并且 $Z\to X$ 可提升到 $X_1$。
\end{lemma}

\begin{proof}
$X_1$ 的存在性以及 $X_1\to X$ 是有限万有同胚这一事实，
是引理 \ref{spaces-pushouts-lemma-coequalizer} 的特殊情形。
$X_1$ 的形成与 $X$ 上的平展局部化交换
（见引理 \ref{spaces-pushouts-lemma-coequalizer} 的证明）。
因此态射 $X_1\to X$ 在 $U$ 上为同构。
从构造立即可知 $Z\to X$ 可提升到 $X_1$。
\end{proof}

\noindent
在情形 \ref{spaces-pushouts-situation-coequalizer-glue} 中，对 $n\geq1$，令
$Z_n\subset X$ 为 $Z$ 在 $X$ 中的第 $n$ 阶无穷小邻域，
即由截出 $Z$ 的理想层的第 $n$ 次幂所定义的闭子概形。
考虑 $Y_n=X'\amalg Z_n$、$R_n=Y_n\times_XY_n$ 以及引理
\ref{spaces-pushouts-lemma-coequalizer-glue} 中的余等化子
$$
\xymatrix{
R_n \ar@<1ex>[r] \ar@<-1ex>[r] & Y_n \ar[r] & X_n \ar[r] & X
}
$$
。映射 $Y_n\to Y_{n+1}$ 和 $R_n\to R_{n+1}$ 诱导态射
\begin{equation}
\label{spaces-pushouts-equation-system-coequalizers}
X_1 \to X_2 \to X_3 \to \ldots \to X
\end{equation}
由于态射 $X_n\to X$ 是万有同胚，这些态射中的每一个也都是万有同胚。

\begin{lemma}
\label{spaces-pushouts-lemma-essentially-constant}
在情形 \ref{spaces-pushouts-situation-coequalizer-glue} 中，假设 $X$ 拟紧。
在 (\ref{spaces-pushouts-equation-system-coequalizers}) 中，对所有充分大的 $n$，
存在 $m$，使 $X_n\to X_{n+m}$ 通过闭浸入 $X\to X_{n+m}$ 分解。
\end{lemma}

\begin{proof}
更仔细地考察 $X_n$ 的构造以及它随 $n$ 增大而发生的变化。我们有
$X_n=\underline{\Spec}(\mathcal{A}_n)$，其中 $\mathcal{A}_n$ 是从
$g_{n,*}\mathcal{O}_{Y_n}$ 到 $h_{n,*}\mathcal{O}_{R_n}$ 的映射
$s_n^\sharp$ 与 $t_n^\sharp$ 的等化子。
这里 $g_n:Y_n=X'\amalg Z_n\to X$ 和
$h_n:R_n=Y_n\times_XY_n\to X$ 是给定态射。
令 $\mathcal{I}\subset\mathcal{O}_X$ 为对应于 $Z$ 的凝聚理想层。那么
$$
g_{n, *}\mathcal{O}_{Y_n} =
f_*\mathcal{O}_{X'} \times \mathcal{O}_X/\mathcal{I}^n
$$
类似地，有分解
$$
R_n = X' \times_X X' \amalg X' \times_X Z_n \amalg
Z_n \times_X X' \amalg Z_n \times_X Z_n
$$
由于 $Z_n\to X$ 是单态射，可知
$X'\times_XZ_n=Z_n\times_XX'$，而这个等同与到 $X$ 的两个态射、
到 $X'$ 的两个态射以及到 $Z_n$ 的两个态射都相容。
把到 $X$ 的态射记作 $f_n:X'\times_XZ_n\to X$。记
$$
\mathcal{A} = \text{等化子}(
\xymatrix{
f_*\mathcal{O}_{X'} \ar@<1ex>[r] \ar@<-1ex>[r] &
(f \times f)_*\mathcal{O}_{X' \times_X X'}
}
)
$$
由上面的说明可知
$$
\mathcal{A}_n =
\text{等化子}(
\xymatrix{
\mathcal{A} \times \mathcal{O}_X/\mathcal{I}^n \ar@<1ex>[r] \ar@<-1ex>[r] &
f_{n, *}\mathcal{O}_{X' \times_X Z_n}
}
)
$$
我们有凝聚 $\mathcal{O}_X$-代数的典范映射
$$
\mathcal{O}_X \to \ldots \to \mathcal{A}_3 \to \mathcal{A}_2 \to \mathcal{A}_1
$$
。本引理的陈述意味着：对充分大的 $n$，存在 $m\geq0$，
使 $\mathcal{A}_{n+m}\to\mathcal{A}_n$ 的像同构于 $\mathcal{O}_X$。
这可以在 $X$ 上平展局部地检验。因此，由《空间的性质》引理
\ref{spaces-properties-lemma-quasi-compact-affine-cover}，
可以假设 $X$ 是仿射诺特概形。

\medskip\noindent
由于 $X_n\to X$ 在 $U$ 上为同构，可知
$\mathcal{O}_X\to\mathcal{A}_n$ 的核支撑在 $|Z|$ 上。
由于 $X$ 诺特，核序列
$\mathcal{J}_n=\Ker(\mathcal{O}_X\to\mathcal{A}_n)$ 稳定
（《空间的上同调》引理 \ref{spaces-cohomology-lemma-acc-coherent}）。
设 $\mathcal{J}_{n_0}=\mathcal{J}_{n_0+1}=\ldots=\mathcal{J}$。
由《空间的上同调》引理
\ref{spaces-cohomology-lemma-power-ideal-kills-sheaf}，
对某个 $t\geq0$，有 $\mathcal{I}^t\mathcal{J}=0$。
另一方面，存在 $\mathcal{O}_X$-代数映射
$\mathcal{A}_n\to\mathcal{O}_X/\mathcal{I}^n$，
故对所有 $n$ 都有 $\mathcal{J}\subset\mathcal{I}^n$。
由 Artin--Rees（《空间的上同调》引理
\ref{spaces-cohomology-lemma-Artin-Rees}），存在 $c\geq0$，使对所有
$n\gg0$ 都有
$\mathcal{J}\cap\mathcal{I}^n\subset\mathcal{I}^{n-c}\mathcal{J}$。
因此 $\mathcal{J}=0$。

\medskip\noindent
如上一段那样，取 $n\geq n_0$。那么
$\mathcal{O}_X\to\mathcal{A}_n$ 是单射。因此，现在只需找到
$m\geq0$，使 $\mathcal{A}_{n+m}\to\mathcal{A}_n$ 的像等于
$\mathcal{O}_X$ 的像。注意，$\mathcal{A}_n$ 位于短正合列
$$
0 \to \Ker(\mathcal{A} \to f_{n, *}\mathcal{O}_{X' \times_X Z_n})
\to \mathcal{A}_n \to \mathcal{O}_X/\mathcal{I}^n \to 0
$$
中；对 $\mathcal{A}_{n+m}$ 也类似。因此，只需证明
$$
\Ker(\mathcal{A} \to f_{n + m, *}\mathcal{O}_{X' \times_X Z_{n + m}})
\subset
\Im(\mathcal{I}^n \to \mathcal{A})
$$
对某个 $m\geq0$ 成立。为此，可以在 $X$ 上平展局部地工作；
又因为 $X$ 诺特，可以假设 $X$ 是诺特仿射概形。
设 $X=\Spec(R)$，而 $\mathcal{I}$ 对应于理想 $I\subset R$。
对某个有限 $R$-代数 $A$，令 $\mathcal{A}=\widetilde A$；
对某个有限 $R$-代数 $B$，令 $f_*\mathcal{O}_{X'}=\widetilde B$。
那么 $R\to A\subset B$，且反演 $I$ 中任一元素后，这些映射都成为同构。

\medskip\noindent
注意，在《空间的上同调》第
\ref{spaces-cohomology-section-theorem-formal-functions} 节所用的记号下，
$f_{n,*}\mathcal{O}_{X'\times_XZ_n}$ 等于
$f_*(\mathcal{O}_{X'}/I^n\mathcal{O}_{X'})$。
由《空间的上同调》引理
\ref{spaces-cohomology-lemma-ML-cohomology-powers-ideal}
，可知存在 $c\geq0$，使得
$$
\Ker(B \to \Gamma(X, f_*(\mathcal{O}_{X'}/I^{n + m + c}\mathcal{O}_{X'}))
$$
包含于 $I^{n+m}B$。% P10-E0030
另一方面，由于 $R\to B$ 有限，且反演 $I$ 中任一元素后成为同构，
可知当 $m$ 充分大时
$I^{n+m}B\subset\Im(I^n\to B)$（可以选得与 $n$ 无关）。
由于 $A\subset B$，这就完成了证明。
\end{proof}

\begin{remark}
\label{spaces-pushouts-remark-essentially-constant}
引理 \ref{spaces-pushouts-lemma-essentially-constant} 的含义是：系统
$X_1\to X_2\to X_3\to\ldots$ 本质上是常值为 $X$ 的。
见《范畴》定义 \ref{categories-definition-essentially-constant-diagram}。
\end{remark}






\section{紧化}
\label{spaces-pushouts-section-compactifications}

\noindent
本节类似于《平坦性续篇》第 \ref{flat-section-compactifications} 节。
本节的定理是 \cite{CLO} 中的主定理。

\medskip\noindent
令 $B$ 为某个基概形 $S$ 上拟紧且拟分离的代数空间。
如果存在从它到某个在 $B$ 上固有的代数空间 $\overline X$
的拟紧开浸入 $X\to\overline X$，就称 $B$ 上的代数空间 $X$
{\it 在 $B$ 上有紧化}或{\it 在 $B$ 上可紧化}。
如果 $X$ 在 $B$ 上有紧化，那么 $X\to B$ 分离且有限型。
本节的主定理断言其逆命题也成立。

\begin{lemma}
\label{spaces-pushouts-lemma-check-separated}
设 $S$ 是概形，$X\to Y$ 是 $S$ 上代数空间的态射。
如果 $(U\subset X,f:V\to X)$ 是初等特异方形，
$U\to Y$ 和 $V\to Y$ 分离，且
$U\times_XV\to U\times_YV$ 是闭态射，那么 $X\to Y$ 分离。
\end{lemma}

\begin{proof}
必须验证 $\Delta:X\to X\times_YX$ 是闭浸入。
$X\times_YX$ 有一个由四部分
$U\times_YU$、$U\times_YV$、$V\times_YU$、$V\times_YV$ 给出的平展覆盖。
注意，
$(U \times_Y U) \times_{(X \times_Y X), \Delta} X  = U$,
$(U \times_Y V) \times_{(X \times_Y X), \Delta} X = U \times_X V$,
$(V \times_Y U) \times_{(X \times_Y X), \Delta} X = V \times_X U$，以及
$(V \times_Y V) \times_{(X \times_Y X), \Delta} X = V$.
因此本引理的假设恰好说明 $\Delta$ 是闭浸入。
\end{proof}

\begin{lemma}
\label{spaces-pushouts-lemma-separate-disjoint-locally-closed-by-blowing-up}
设 $S$ 是概形，$X$ 是 $S$ 上拟紧且拟分离的代数空间，
$U\subset X$ 是拟紧开子空间。
\begin{enumerate}
\item 如果 $Z_1,Z_2\subset X$ 是有限呈示的闭子空间，且
$Z_1\cap Z_2\cap U=\emptyset$，那么存在 $U$-容许爆破 $X'\to X$，
使 $Z_1$ 与 $Z_2$ 的严格变换不交；
\item 如果 $T_1,T_2\subset|U|$ 是不交的可构造闭子集，
那么存在 $U$-容许爆破 $X'\to X$，使 $T_1$ 与 $T_2$ 的闭包不交。
\end{enumerate}
\end{lemma}

\begin{proof}
证明 (1)。$Z_i\to X$ 有限呈示这一假设意味着 $Z_i$ 的拟凝聚理想层
$\mathcal{I}_i$ 有限型；见《空间的态射》引理
\ref{spaces-morphisms-lemma-closed-immersion-finite-presentation}.
以 $Z\subset X$ 表示由乘积 $\mathcal{I}_1\mathcal{I}_2$ 截出的闭子空间。
注意，$Z\cap U$ 是 $Z_1\cap U$ 与 $Z_2\cap U$ 的不交并。
由《空间上的除子》引理
\ref{spaces-divisors-lemma-separate-disjoint-opens-by-blowing-up}
，存在 $U\cap Z$-容许爆破 $Z'\to Z$，
使 $Z_1$ 与 $Z_2$ 的严格变换不交。
以 $Y\subset Z$ 表示这个爆破的中心。那么 $Y\to X$ 是有限呈示的闭浸入，
因为它是 $Y\to Z$ 与 $Z\to X$ 的复合
（《空间上的除子》定义
\ref{spaces-divisors-definition-admissible-blowup} 和
《空间的态射》引理
\ref{spaces-morphisms-lemma-composition-finite-presentation}).
）。因此，沿 $Y$ 的爆破 $X'\to X$ 是 $U$-容许爆破。
由严格变换的一般性质，$Z_1$、$Z_2$ 关于 $X'\to X$ 的严格变换，
与它们关于 $Z'\to Z$ 的严格变换相同；见《空间上的除子》引理
\ref{spaces-divisors-lemma-strict-transform}。这证明了 (1)。

\medskip\noindent
证明 (2)。由《空间的极限》引理
\ref{spaces-limits-lemma-quasi-coherent-finite-type-ideals}
，存在有限型拟凝聚理想层
$\mathcal{J}_i\subset\mathcal{O}_U$，使得
$T_i=V(\mathcal{J}_i)$（作为集合）。
由《空间的极限》引理 \ref{spaces-limits-lemma-extend}，
存在有限型拟凝聚理想层 $\mathcal{I}_i\subset\mathcal{O}_X$，
其到 $U$ 的限制为 $\mathcal{J}_i$。
把 (1) 的结果应用于闭子空间 $Z_i=V(\mathcal{I}_i)$，即得结论。
\end{proof}

\begin{lemma}
\label{spaces-pushouts-lemma-blowup-iso-along}
设 $S$ 是概形，$f:X\to Y$ 是 $S$ 上拟紧且拟分离代数空间的固有态射。
令 $V\subset Y$ 为拟紧开子空间，$U=f^{-1}(V)$。
令 $T\subset|V|$ 为闭子集，使 $f|_U:U\to V$ 在 $V$ 中 $T$
的某个开邻域上为同构。那么存在 $V$-容许爆破 $Y'\to Y$，
使 $f$ 的严格变换 $f':X'\to Y'$ 在 $T$ 于 $|Y'|$ 中的闭包的
某个开邻域上为同构。
\end{lemma}

\begin{proof}
令 $T'\subset|V|$ 为 $f|_U$ 在其上为同构的最大开集的补集。
那么 $T',T$ 在 $|V|$ 中闭，且 $T\cap T'=\emptyset$。
由于 $|V|$ 是谱拓扑空间（《空间的性质》引理
\ref{spaces-properties-lemma-quasi-compact-quasi-separated-spectral})
），可以在 $|V|$ 中找到可构造闭子集 $T_c,T'_c$，使
$T\subset T_c$、$T'\subset T'_c$ 且 $T_c\cap T'_c=\emptyset$
（选取拟紧开集 $W$，它在 $|V|$ 中包含 $T'$ 且与 $T$ 不交，
置 $T_c=|V|\setminus W$；再选取拟紧开集 $W'$，它在 $|V|$ 中
包含 $T_c$ 且与 $T'$ 不交，置 $T'_c=|V|\setminus W'$）。
由引理 \ref{spaces-pushouts-lemma-separate-disjoint-locally-closed-by-blowing-up}，
在用一个 $V$-容许爆破替换 $Y$ 后，可以假设 $T_c$ 与 $T'_c$
在 $|Y|$ 中的闭包不交。令 $Y_0$ 为 $Y$ 中对应于开集
$|Y|\setminus\overline{T}'_c$ 的开子空间，并置
$V_0=V\cap Y_0$、$U_0=U\times_VV_0$、$X_0=X\times_YY_0$。
由于 $U_0\to V_0$ 是同构，可以找到 $V_0$-容许爆破
$Y'_0\to Y_0$，使 $X_0$ 的严格变换 $X'_0$ 同构地映到 $Y'_0$；
见《空间的态射续篇》引理
\ref{spaces-more-morphisms-lemma-zariski-after-blowup}.
由《空间上的除子》引理
\ref{spaces-divisors-lemma-extend-admissible-blowups}
，存在 $V$-容许爆破 $Y'\to Y$，其到 $Y_0$ 的限制为 $Y'_0\to Y_0$。
如果以 $f':X'\to Y'$ 表示 $f$ 的严格变换，
那么由于 $f'$ 限制到 $Y'_0$ 上为同构，所需结论成立。
\end{proof}

\begin{lemma}
\label{spaces-pushouts-lemma-blowup-etale-along}
设 $S$ 是概形。考虑 $S$ 上拟紧且拟分离代数空间的图表
$$
\xymatrix{
X \ar[d]_f & U \ar[l] \ar[d]_{f|_U} & A \ar[d] \ar[l] \\
Y & V \ar[l] & B \ar[l]
}
$$
。假设
\begin{enumerate}
\item $f$ 固有；
\item $V$ 是 $Y$ 的拟紧开子空间，$U=f^{-1}(V)$；
\item $B\subset V$ 和 $A\subset U$ 是闭子空间；
\item $f|_A:A\to B$ 是同构，且 $f$ 在 $A$ 的每一点都平展。
\end{enumerate}
那么存在 $V$-容许爆破 $Y'\to Y$，使严格变换 $f':X'\to Y'$ 满足：
对 $|A|$ 在 $|X'|$ 中的闭包的每个几何点 $\overline a$，
都存在商 $\mathcal{O}_{X',\overline a}\to\mathcal{O}$，使得
$\mathcal{O}_{Y',f'(\overline a)}\to\mathcal{O}$ 有限平坦。
\end{lemma}

\noindent
从证明可以看出，实际上有更强的结论；但这个陈述已经足够长，
而且它足以满足后面的需要。

\begin{proof}
令 $T'\subset|U|$ 为 $f|_U$ 在其上平展的最大开集的补集。
那么 $T'$ 在 $|U|$ 中闭，并且与 $|A|$ 不交。
由于 $|U|$ 是谱拓扑空间（《空间的性质》引理
\ref{spaces-properties-lemma-quasi-compact-quasi-separated-spectral})
），可以找到 $|U|$ 的可构造闭子集 $T_c,T'_c$，使
$|A|\subset T_c$、$T'\subset T'_c$ 且 $T_c\cap T'_c=\emptyset$
（见引理 \ref{spaces-pushouts-lemma-blowup-iso-along} 的证明）。
由引理 \ref{spaces-pushouts-lemma-separate-disjoint-locally-closed-by-blowing-up}，
存在 $U$-容许爆破 $X_1\to X$，使 $T_c$ 与 $T'_c$ 在 $|X_1|$
中的闭包不交。令 $X_{1,0}$ 为 $X_1$ 中对应于开集
$|X_1|\setminus\overline{T}'_c$ 的开子空间，并置
$U_0=U\cap X_{1,0}$。注意，按构造，$A$ 的概形论像
$\overline A_1\subset X_1$ 包含于 $X_{1,0}$。

\medskip\noindent
在用一个 $V$-容许爆破替换 $Y$ 并取严格变换后，
可以假设 $X_{1,0}\to Y$ 平坦、拟有限且有限呈示；
见《空间的态射续篇》引理
\ref{spaces-more-morphisms-lemma-flat-after-blowing-up} 和
\ref{spaces-more-morphisms-lemma-flat-fp-dimension-over-dense-open}.
考虑交换图表
$$
\vcenter{
\xymatrix{
X_1 \ar[rr] \ar[rd] & & X \ar[ld] \\
& Y
}
}
\quad\text{以及图表}\quad
\vcenter{
\xymatrix{
\overline{A}_1 \ar[rr] \ar[rd] & & \overline{A} \ar[ld] \\
& \overline{B}
}
}
$$
，以及概形论像的图表。态射 $\overline A_1\to\overline A$ 是满射：
它固有，因而 $\overline A_1\to\overline A$ 的概形论像必等于
$\overline A$，然后可以使用《空间的态射》引理
\ref{spaces-morphisms-lemma-scheme-theoretic-image-is-proper}.
。关于平展局部环的陈述可如下得到：把几何点 $\overline a$
提升为 $\overline A_1$ 的几何点 $\overline a_1$，并置
$\mathcal{O}=\mathcal{O}_{X_1,\overline a_1}$。
事实上，由于 $X_1\to Y$ 在
$X_{1,0}\supset\overline A_1$ 上平坦且拟有限，映射
$\mathcal{O}_{Y',f'(\overline a)}\to
\mathcal{O}_{X_1,\overline a_1}$ 有限平坦；见《代数》引理
\ref{algebra-lemma-quasi-finite-strict-henselization}
和 \ref{algebra-lemma-characterize-henselian}。
\end{proof}

\begin{lemma}
\label{spaces-pushouts-lemma-replaced-by-strict-transform}
设 $S$ 是概形，$X\to B$ 和 $Y\to B$ 是 $S$ 上代数空间的态射，
$U\subset X$ 是开子空间。令 $V\to X\times_BY$ 为拟紧态射，
它与第一个投影的复合映入 $U$。令 $Z\subset X\times_BY$ 为
$V\to X\times_BY$ 的概形论像。令 $X'\to X$ 为 $U$-容许爆破。
那么 $V\to X'\times_BY$ 的概形论像，就是 $Z$ 关于该爆破的严格变换。
\end{lemma}

\begin{proof}
以 $Z'\to Z$ 表示严格变换。态射 $Z'\to X'$ 诱导闭浸入
$Z'\to X'\times_BY$（因为按定义，$Z'$ 是 $X'\times_XZ$ 的闭子空间）。
因此，为完成证明，只需说明 $V\to Z'$ 的概形论像 $Z''$ 就是 $Z'$。
注意，$Z''\subset Z'$ 是闭子空间，且 $V\to Z'$ 通过 $Z''$ 分解。
由于 $V\to X\times_BY$ 和 $V\to X'\times_BY$ 都拟紧
（对后者，这来自《空间的态射》引理
\ref{spaces-morphisms-lemma-quasi-compact-permanence}
以及如下事实：$X'\times_BY\to X\times_BY$ 是固有态射的基变换，因而分离），
由《空间的态射》引理
\ref{spaces-morphisms-lemma-quasi-compact-scheme-theoretic-image}
，可知 $Z\cap(U\times_BY)=Z''\cap(U\times_BY)$。
因此，包含态射 $Z''\to Z'$ 在 $Z'\to Z$ 的例外除子 $E$ 之外是同构。
然而，$Z'$ 的结构层没有任何支撑在 $E$ 上的非零截面
（这是严格变换的定义），所以满射
$\mathcal{O}_{Z'}\to\mathcal{O}_{Z''}$ 必为同构。
\end{proof}

\begin{lemma}
\label{spaces-pushouts-lemma-compactification-dominates}
设 $S$ 是概形，$B$ 是 $S$ 上拟紧且拟分离的代数空间，
$U$ 是 $B$ 上有限型且分离的代数空间，$V\to U$ 是平展态射。
如果 $V$ 在 $B$ 上有紧化 $V\subset Y$，那么存在 $V$-容许爆破
$Y'\to Y$ 和开子空间 $V\subset V'\subset Y'$，
使 $V\to U$ 延拓为固有态射 $V'\to U$。
\end{lemma}

\begin{proof}
考虑``对角''态射 $V\to Y\times_BU$ 的概形论像
$Z\subset Y\times_BU$。如果用一个 $V$-容许爆破替换 $Y$，
那么 $Z$ 被替换为关于该爆破的严格变换；见引理
\ref{spaces-pushouts-lemma-replaced-by-strict-transform}。因此，由《空间的态射续篇》引理
\ref{spaces-more-morphisms-lemma-zariski-after-blowup}
，可以假设 $Z\to Y$ 是开浸入。若以 $V'\subset Y$ 表示其像，
则诱导态射 $V'\to U$ 固有，因为投影 $Y\times_BU\to U$ 固有，
而 $V'\cong Z$ 是 $Y\times_BU$ 的闭子空间。
\end{proof}

\noindent
下面的引理针对某个 $\mathbf Z$ 上有限型代数空间上的有限型分离代数空间来表述。
拟紧且拟分离代数空间的版本也成立（证明本质上相同），
但它将立即由本节的主定理推出。我们强烈建议读者先阅读本引理在概形情形下的证明。

\begin{lemma}
\label{spaces-pushouts-lemma-two-compactifications}
令 $B$ 为 $\mathbf Z$ 上有限型的代数空间，
$U$ 为 $B$ 上有限型且分离的代数空间。
令 $(U_2\subset U,f:U_1\to U)$ 为初等特异方形。
假设 $U_1$ 和 $U_2$ 在 $B$ 上都有紧化，且
$U_1\times_UU_2\to U$ 的像稠密。那么 $U$ 在 $B$ 上有紧化。
\end{lemma}

\begin{proof}
对 $i=1,2$，选取 $B$ 上的紧化 $U_i\subset X_i$。
可以假设 $U_i$ 在 $X_i$ 中概形论稠密。由引理
\ref{spaces-pushouts-lemma-compactification-dominates}，可以假设存在开子空间
$V_i\subset X_i$ 和固有态射 $\psi_i:V_i\to U$，延拓 $U_i\to U$。
图表为
$$
\xymatrix{
U_i \ar[r] \ar[d] & V_i \ar[r] \ar[dl]^{\psi_i} & X_i \\
U
}
$$
以 $Z_1\subset U$ 表示对应于闭子集 $|U|\setminus|U_2|$ 的约化闭子空间。
回忆 $f^{-1}Z_1$ 是 $U_1$ 的闭子空间，并同构地映到 $Z_1$。
以 $Z_2\subset U$ 表示对应于闭子集
$|U|\setminus\Im(|f|)=
|U_2|\setminus\Im(|U_1\times_UU_2|\to|U_2|)$ 的约化闭子空间。
于是，作为集合有
$$
U = U_2 \amalg Z_1 = Z_2 \amalg \Im(f) =
Z_2 \amalg \Im(U_1 \times_U U_2 \to U_2) \amalg Z_1
$$
。以 $Z_{i,i}\subset V_i$ 表示 $Z_i$ 在 $\psi_i$ 下的逆像。
注意，$\psi_2$ 在 $Z_2$ 的某个开邻域上为同构。又注意，对某个与
$f^{-1}Z_1$ 不交的闭子空间 $T\subset V_1$，有
$Z_{1,1}=\psi_1^{-1}Z_1=f^{-1}Z_1\amalg T$，
而且 $\psi_1$ 沿 $f^{-1}Z_1$ 平展。
以 $Z_{i,j}\subset V_i$ 表示 $Z_j$ 在 $\psi_i$ 下的逆像。
注意，$\psi_i:Z_{i,j}\to Z_j$ 是固有态射。
由于 $Z_i$ 与 $Z_j$ 是 $U$ 的不交闭子空间，
$Z_{i,i}$ 与 $Z_{i,j}$ 是 $V_i$ 的不交闭子空间。

\medskip\noindent
以 $\overline Z_{i,i}$ 和 $\overline Z_{i,j}$ 表示
$Z_{i,i}$ 和 $Z_{i,j}$ 在 $X_i$ 中的概形论像。
回忆 $|Z_{i,j}|$ 在 $|\overline Z_{i,j}|$ 中稠密；
见《空间的态射》引理
\ref{spaces-morphisms-lemma-quasi-compact-immersion}.
用一个 $V_i$-容许爆破替换 $X_i$ 后，可以假设
$\overline Z_{i,i}$ 与 $\overline Z_{i,j}$ 不交；见引理
\ref{spaces-pushouts-lemma-separate-disjoint-locally-closed-by-blowing-up}。
假设这一点对 $X_1$ 和 $X_2$ 都成立。
注意，若再用一个 $V_i$-容许爆破替换 $X_i$，此性质仍保持。
因此，可以用另一个 $V_1$-容许爆破替换 $X_1$，并假设
$|\overline Z_{1,1}|$ 是 $|T|$ 与 $|f^{-1}Z_1|$ 在 $|X_1|$
中的闭包的不交并。

\medskip\noindent
置 $V_{12}=V_1\times_UV_2$。我们有浸入
$V_{12}\to X_1\times_BX_2$，它是闭浸入
$V_{12}=V_1\times_UV_2\to V_1\times_BV_2$
（《空间的态射》引理
\ref{spaces-morphisms-lemma-fibre-product-after-map})
）与开浸入 $V_1\times_BV_2\to X_1\times_BX_2$ 的复合。
令 $X_{12}\subset X_1\times_BX_2$ 为
$V_{12}\to X_1\times_BX_2$ 的概形论像。投影态射
$$
p_1 : X_{12} \to X_1
\quad\text{和}\quad
p_2 : X_{12} \to X_2
$$
固有，因为 $X_1$ 和 $X_2$ 在 $B$ 上固有。如果用一个
$V_1$-容许爆破替换 $X_1$，那么 $X_{12}$ 被替换为关于这个爆破的严格变换；
见引理 \ref{spaces-pushouts-lemma-replaced-by-strict-transform}。

\medskip\noindent
以 $\psi:V_{12}\to U$ 表示复合
$\psi=\psi_1\circ p_1|_{V_{12}}=\psi_2\circ p_2|_{V_{12}}$。
考虑闭子空间
$$
Z_{12, 2} =
(p_1|_{V_{12}})^{-1}Z_{1, 2} =
(p_2|_{V_{12}})^{-1}Z_{2, 2} =
\psi^{-1}Z_2 \subset V_{12}
$$
态射 $p_1|_{V_{12}}:V_{12}\to V_1$ 在 $Z_{1,2}$ 的某个开邻域上为同构，
因为 $\psi_2:V_2\to U$ 在 $Z_2$ 的某个开邻域上为同构，
且 $V_{12}=V_1\times_UV_2$。由引理 \ref{spaces-pushouts-lemma-blowup-iso-along}，
存在 $V_1$-容许爆破 $X'_1\to X_1$，使 $p_1$ 的严格变换
$p'_1:X'_{12}\to X'_1$ 在 $|Z_{1,2}|$ 于 $|X'_1|$ 中的闭包的
某个开邻域上为同构。用 $X'_1$ 替换 $X_1$、用 $X'_{12}$ 替换
$X_{12}$ 后，可以假设 $p_1$ 在 $|\overline Z_{1,2}|$
的某个开邻域上为同构。

\medskip\noindent
上一段的结果告诉我们
$$
X_{12} \cap (\overline{Z}_{1, 2} \times_B \overline{Z}_{2, 1}) = \emptyset
$$
，其中交取在 $X_1\times_BX_2$ 中。事实上，$X_{12}$ 中的逆像
$p_1^{-1}\overline Z_{1,2}$ 同构地映到 $\overline Z_{1,2}$。
特别地，$|Z_{12,2}|$ 在 $|p_1^{-1}\overline Z_{1,2}|$ 中稠密。
因此，$p_2$ 把 $|p_1^{-1}\overline Z_{1,2}|$ 映入
$|\overline Z_{2,2}|$。由于
$|\overline Z_{2,2}|\cap|\overline Z_{2,1}|=\emptyset$，结论随即得到。

\medskip\noindent
事实表明，在结束论证前还需要再作一次爆破。令
$V_2\subset W_2\subset X_2$ 为具有如下底层拓扑空间的开子空间：
$$
|W_2| =
|V_2| \cup (|X_2| \setminus |\overline{Z}_{2, 1}|) =
|X_2| \setminus \left(|\overline{Z}_{2, 1}| \setminus |Z_{2, 1}|\right)
$$
由于 $p_2(p_1^{-1}\overline Z_{1,2})$ 包含于 $W_2$（见上），
用一个 $W_2$-容许爆破替换 $X_2$，并用相应的严格变换替换 $X_{21}$，% P10-E0031
仍会保持 $p_1$ 在 $\overline Z_{1,2}$ 的某个开邻域上为同构这一性质。
由于 $\overline Z_{2,1}\cap W_2=\overline Z_{2,1}\cap V_2=Z_{2,1}$，
可知 $Z_{2,1}$ 是 $W_2$ 和 $V_2$ 的闭子空间。
注意，作为 $X_{12}$ 的开子空间，有
$V_{12}=V_1\times_UV_2=p_1^{-1}(V_1)=p_2^{-1}(V_2)$，
因为它是 $X_{12}$ 中态射 $\psi:V_{12}\to U$ 可延拓到的最大开子空间；
略去细节\footnote{事实上，$V_1\times_UV_2$ 在 $U$ 上固有，
所以如果 $\psi$ 可延拓到 $X_{12}$ 的一个更大开集，那么由《空间的态射》引理
\ref{spaces-morphisms-lemma-universally-closed-permanence}.
，$V_1\times_UV_2$ 在这个开集中将是闭的。
又因为 $V_{12}\subset X_{12}$ 稠密，所以得到相等。}。
我们有 $V_{12}$ 的闭子空间之间的下列等式：
$$
p_2^{-1}Z_{2, 1} = p_2^{-1} \psi_2^{-1} Z_1 =
p_1^{-1} \psi_1^{-1} Z_1= p_1^{-1}Z_{1, 1} =
p_1^{-1}f^{-1}Z_1 \amalg p_1^{-1}T
$$
这里以及下文中，我们略微滥用记号，以 $p_2$ 表示 $p_2$ 到% P10-E0032
$V_{12}$ 的限制，等等。由于 $Z_{2,1}$ 是 $W_2$ 的闭子空间，
$p_2^{-1}(Z_{2,1})$ 是 $p_2^{-1}(W_2)$ 的闭子空间；因此
$p_1^{-1}f^{-1}Z_1$ 也是 $p_2^{-1}(W_2)$ 的闭子空间。
最后，因为 $\psi_1$ 沿 $f^{-1}Z_1$ 平展，且
$V_{12}=V_1\times_UV_2$，态射 $p_2:X_{12}\to X_2$
在 $p_1^{-1}f^{-1}Z_1$ 的各点平展。
于是可以把引理 \ref{spaces-pushouts-lemma-blowup-etale-along} 应用于态射
$p_2:X_{12}\to X_2$、开子空间 $W_2$、闭子空间
$Z_{2,1}\subset W_2$ 以及闭子空间
$p_1^{-1}f^{-1}Z_1\subset p_2^{-1}(W_2)$。
因此，用一个 $W_2$-容许爆破替换 $X_2$，并用相应的严格变换替换
$X_{12}$ 后，对 $|p_1^{-1}f^{-1}Z_1|$ 的闭包中的每个几何点
$\overline y$，都得到局部环映射
$\mathcal{O}_{X_{12},\overline y}\to\mathcal{O}$，使
$\mathcal{O}_{X_2,p_2(\overline y)}\to\mathcal{O}$ 有限平坦。

\medskip\noindent
考虑由推出得到的代数空间
$$
W_2 = U \coprod\nolimits_{U_2} (X_2 \setminus \overline{Z}_{2, 1}),
$$
，以及以第一段中的 $T\subset V_1$ 定义的代数空间
$$
W_1 = U \coprod\nolimits_{U_1}
(X_1 \setminus \overline{Z}_{1, 2} \cup \overline{T}),
$$
；见引理 \ref{spaces-pushouts-lemma-construct-elementary-distinguished-square}。
应用引理 \ref{spaces-pushouts-lemma-check-separated} 来说明 $W_i\to B$ 分离。
首先，$U\to B$ 和 $X_i\to B$ 分离。利用赋值判据，来验证拟紧浸入
$U_i\to U\times_B(X_i\setminus\overline Z_{i,j})$ 是闭的；
见《空间的态射》引理
\ref{spaces-morphisms-lemma-quasi-compact-existence-universally-closed}.
选取 $B$ 上的赋值环 $A$，其分式域为 $K$，并选取相容态射
$(u,x_i):\Spec(A)\to U\times_BX_i$ 和
$u_i:\Spec(K)\to U_i$。由于 $\psi_i$ 固有，可以找到与 $u$、$u_i$
相容的唯一态射 $v_i:\Spec(A)\to V_i$。
由于 $X_i$ 在 $B$ 上固有，可知 $x_i=v_i$。
如果 $v_i$ 不通过 $U_i\subset V_i$ 分解，那么可知 $x_i$
把 $\Spec(A)$ 的闭点映入 $Z_{i,j}$，或者在 $i=1$ 时映入 $T$。
这就完成了证明，因为在构造 $W_i$ 时移除了
$\overline Z_{i,j}$ 和 $\overline T$。% P10-E0033

\medskip\noindent
另一方面，对任一 $B$ 上赋值环 $A$（其分式域为 $K$）以及任一 $B$ 上态射
$$
\gamma : \Spec(K) \to \Im(U_1 \times_U U_2 \to U)
$$
，我们声称：用赋值环的某个扩张替换 $A$ 后，存在某个 $i$，
使 $\gamma$ 延拓为态射 $h_i:\Spec(A)\to W_i$。
事实上，先利用固有性的赋值判据把 $\gamma$ 延拓为态射
$g_2:\Spec(A)\to X_2$。如果 $g_2$ 的像不与
$\overline Z_{2,1}$ 相交，就得到所需的到 $W_2$ 的态射。
否则，令 $\overline z\in\overline Z_{2,1}$ 为位于 $g_2$
作用下闭点之像上方的一个几何点。可以把它提升为 $X_{12}$ 的一个几何点
$\overline y$，该点属于 $|p_1^{-1}f^{-1}Z_1|$ 的闭包；
这是因为空间映射
$|p_1^{-1}f^{-1}Z_1|\to|\overline Z_{2,1}|$ 是闭映射，
且其像包含稠密开集 $|Z_{2,1}|$。用 $A$ 的严格 Hensel 化替换该环
（《代数续篇》引理 \ref{more-algebra-lemma-henselization-valuation-ring}）后，
得到图表
$$
\xymatrix{
A \ar@{..>}[rr] & & A' \\
\mathcal{O}_{X_2, \overline{z}} \ar[r] \ar[u] &
\mathcal{O}_{X_{12}, \overline{y}} \ar[r] &
\mathcal{O} \ar@{..>}[u]
}
$$
，其中 $\mathcal{O}_{X_{12},\overline y}\to\mathcal{O}$ 是证明第五段中
找到的映射。由于水平复合有限且平坦，可以找到赋值环的扩张
$A'/A$ 和使图表交换的虚线箭头。用 $A'$ 替换 $A$ 后，
这意味着得到提升 $g_{12}:\Spec(A)\to X_{12}$，
其闭点映入 $|p_1^{-1}f^{-1}Z_1|$ 的闭包。
于是 $g_1=p_1\circ g_{12}:\Spec(A)\to X_1$ 的闭点映入
$|f^{-1}Z_1|$ 的闭包。由于 $|f^{-1}Z_1|$ 的闭包与 $|T|$ 的闭包不交，
并包含于 $|\overline Z_{1,1}|$，而后者与
$|\overline Z_{1,2}|$ 不交，所以 $g_1$ 按所需定义态射
$h_1:\Spec(A)\to W_1$。

\medskip\noindent
考虑图表
$$
\xymatrix{
W_1' \ar[d] \ar[r] & W & W_2' \ar[l] \ar[d] \\
W_1 & U \ar[l] \ar[lu] \ar[u] \ar[ru] \ar[r] & W_2
}
$$
，如《空间的态射续篇》引理
\ref{spaces-more-morphisms-lemma-find-common-blowups}.
中所述。由上一段，对每个实线图表
$$
\xymatrix{
\Spec(K) \ar[r]_\gamma  \ar[d] & W \ar[d] \\
\Spec(A) \ar@{..>}[ru] \ar[r] & B
}
$$
，若 $\Im(\gamma)\subset\Im(U_1\times_UU_2\to U)$，
那么在必要时用赋值环的某个扩张替换 $A$ 后，存在某个 $i$ 以及
$\gamma$ 的延拓 $h_i:\Spec(A)\to W_i$。
再利用 $W'_i\to W_i$ 的固有性赋值判据，可以把 $h_i$ 提升为
$h'_i:\Spec(A)\to W'_i$。因此，在必要时扩张 $A$ 后，
图表中的虚线箭头存在。由于 $W$ 在 $B$ 上分离，
可知其实不需要选取扩张，而且该箭头也是唯一的；
见《空间的态射》引理 \ref{spaces-morphisms-lemma-usual-enough} 和
\ref{spaces-morphisms-lemma-separated-implies-valuative}。
最后，由《空间的态射》引理
\ref{spaces-morphisms-lemma-refined-valuative-criterion-universally-closed}，
虚线箭头的存在性蕴含 $W\to B$ 万有闭。
由于 $W\to B$ 已经有限型且分离，结论得证。
\end{proof}

\begin{lemma}
\label{spaces-pushouts-lemma-filter-Noetherian-space}
设 $S$ 是概形，$X$ 是 $S$ 上的诺特代数空间，
$U\subset X$ 是真稠密开子空间。那么存在仿射概形 $V$ 和平展态射
$V\to X$，使得
\begin{enumerate}
\item 开子空间 $W=U\cup\Im(V\to X)$ 严格大于 $U$；
\item $(U\subset W,V\to W)$ 是特异方形；并且
\item $U\times_WV\to U$ 的像稠密。
\end{enumerate}
\end{lemma}

\begin{proof}
选取一个分层
$$
\emptyset = U_{n + 1} \subset
U_n \subset U_{n - 1} \subset \ldots \subset U_1 = X
$$
以及如《合宜空间》引理
\ref{decent-spaces-lemma-filter-quasi-compact-quasi-separated}.
中那样的态射 $f_p:V_p\to U_p$。
令 $p$ 为满足 $U_p\not\subset U$ 的最小整数
（由于 $U \not = X$，这是可能的）。选取仿射开集 $V\subset V_p$，
使平展态射 $f_p|_V:V\to X$ 不通过 $U$ 分解。
考虑开集 $W=U\cup\Im(V\to X)$，以及满足
$|Z|=|W|\setminus|U|$ 的约化闭子空间 $Z\subset W$。
那么 $f^{-1}Z\to Z$ 是同构，% P10-E0034
因为对态射 $f_p$ 有相应性质；见上引理。因此
$(U\subset W,f:V\to W)$ 是特异方形。
开集 $I=\Im(U\times_WV\to U)$ 未必在 $U$ 中稠密。
底层集合为 $|U|\setminus\overline{|I|}$ 的代数空间 $U'\subset U$
是诺特的，因而可以找到稠密开子概形 $U''\subset U'$；
例如见《空间的性质》命题
\ref{spaces-properties-proposition-locally-quasi-separated-open-dense-scheme}.
然后可以找到稠密仿射开集 $U'''\subset U''$；见《性质》引理
\ref{properties-lemma-Noetherian-irreducible-components} 和
\ref{properties-lemma-maximal-points-affine}。
用 $V\amalg U'''\to X$ 替换 $f$ 后，一切都很清楚。
\end{proof}

\begin{theorem}
\label{spaces-pushouts-theorem-nagata}
\begin{reference}
\cite{CLO}
\end{reference}
设 $S$ 是概形，$B$ 是 $S$ 上拟紧且拟分离的代数空间，
$X\to B$ 是分离有限型态射。那么 $X$ 在 $B$ 上有紧化。
\end{theorem}

\begin{proof}
先约化到诺特情形。我们强烈建议读者跳过本段。
首先，可以用 $\Spec(\mathbf Z)$ 替换 $S$；见《空间》第
\ref{spaces-section-change-base-scheme} 节和《空间的性质》定义
\ref{spaces-properties-definition-separated}。
存在闭浸入 $X\to X'$，且 $X'\to B$ 有限呈示并分离；
见《空间的极限》命题
\ref{spaces-limits-proposition-separated-closed-in-finite-presentation}.
如果找到 $X'$ 在 $B$ 上的紧化，那么取其中 $X$ 的概形论闭包，
就得到 $X$ 在 $B$ 上的紧化。因此，可以假设 $X\to B$ 分离且有限呈示。
可以把 $B=\lim B_i$ 写成一个系统的有向极限，其中各项是
$\Spec(\mathbf Z)$ 上有限型的诺特代数空间，转移态射仿射；
见《空间的极限》命题 \ref{spaces-limits-proposition-approximate}。
可以选取某个 $i$ 和有限呈示态射 $X_i\to B_i$，
使它到 $B$ 的基变换为 $X\to B$；见《空间的极限》引理
\ref{spaces-limits-lemma-descend-finite-presentation}。
增大 $i$ 后，可以假设 $X_i\to B_i$ 分离；见《空间的极限》引理
\ref{spaces-limits-lemma-descend-separated-morphism}.
如果能找到 $X_i$ 在 $B_i$ 上的紧化，那么它到 $B$ 的基变换就是
$X$ 在 $B$ 上的紧化。这把问题约化到下一段讨论的情形。

\medskip\noindent
除拟紧且拟分离外，再假设 $B$ 在 $\mathbf Z$ 上有限型。
令 $U\to X$ 为代数空间的平展态射，并假设 $U$ 在
$\Spec(\mathbf Z)$ 上有紧化 $Y$。态射
$$
U \longrightarrow B \times_{\Spec(\mathbf{Z})} Y
$$
由《空间的态射》引理
\ref{spaces-morphisms-lemma-monomorphism-loc-finite-type-loc-quasi-finite}
可知分离且拟有限（所示态射分解为一个浸入，因而是单态射）。
所以由 Zariski 主定理（《空间的态射续篇》引理
\ref{spaces-more-morphisms-lemma-quasi-finite-separated-pass-through-finite})
），存在 $U$ 到某个代数空间 $Y'$ 的开浸入，且后者在
$B\times_{\Spec(\mathbf Z)}Y$ 上有限。那么 $Y'\to B$ 固有，
因为它是两个固有态射的复合
$Y'\to B\times_{\Spec(\mathbf Z)}Y\to B$
（使用《空间的态射》引理 \ref{spaces-morphisms-lemma-finite-proper}、
\ref{spaces-morphisms-lemma-composition-proper} 和
\ref{spaces-morphisms-lemma-base-change-proper}）。
）。因此 $U$ 在 $B$ 上有紧化。

\medskip\noindent
存在一个为概形的稠密开子空间 $U\subset X$
（《空间的性质》命题
\ref{spaces-properties-proposition-locally-quasi-separated-open-dense-scheme}).
）。事实上，可以选取 $U$ 为仿射概形
（《性质》引理 \ref{properties-lemma-Noetherian-irreducible-components} 和
\ref{properties-lemma-maximal-points-affine}）。
因此 $U$ 在 $\Spec(\mathbf Z)$ 上有紧化；
这可以直接轻易证明，也可由概形的相应定理推出，
见《平坦性续篇》定理 \ref{flat-theorem-nagata}。
由上一段，$U$ 在 $B$ 上有紧化。
通过诺特归纳，可以找到一个在 $B$ 上有紧化的极大稠密开子空间
$U\subset X$。我们将说明假设 $U \not = X$ 会导致矛盾。
事实上，由引理 \ref{spaces-pushouts-lemma-filter-Noetherian-space}，
可以找到严格更大的开子空间 $U\subset W\subset X$ 和特异方形
$(U\subset W,f:V\to W)$，其中 $V$ 仿射，且
$U\times_WV$ 在 $U$ 中的像稠密。
由于 $V$ 仿射，它和前面一样在 $B$ 上有紧化。
因此可以应用引理 \ref{spaces-pushouts-lemma-two-compactifications}，
得到 $W$ 在 $B$ 上有紧化，这正是所需矛盾。
\end{proof}
% Generated mechanically from the frozen spaces-chow.tex target.
% Do not edit this generated file; edit the bound locale source instead.

% OK, start here.
%

\setcounter{chapter}{81}
\chapter{代数空间的周群}


\phantomsection
\label{spaces-chow-section-phantom}




\section{引言}
\label{spaces-chow-section-introduction}

\noindent
本章先讨论代数空间的周群。定义这些群之后，我们把向量丛的陈类
定义为作用于这些周群的算子。所用策略与概形情形完全相同。因此，
我们建议读者先阅读概形相应章节的引言
（《Chow 同调》第 \ref{chow-section-introduction} 节）。

\medskip\noindent
相关论文有 \cite{edidin-graham} 和 \cite{kresch_cycle}。



\section{设定}
\label{spaces-chow-section-setup}

\noindent
我们先固定所要使用的代数空间范畴。请在本章始终注意，根据
《合宜空间》引理
\ref{decent-spaces-lemma-locally-Noetherian-decent-quasi-separated}，
“合宜 $+$ 局部诺特”等价于“拟分离 $+$ 局部诺特”。

\begin{situation}
\label{spaces-chow-situation-setup}
这里 $S$ 是概形，$B$ 是 $S$ 上的代数空间。假设 $B$ 拟分离、
局部诺特且普遍链状（《合宜空间》定义
\ref{decent-spaces-definition-universally-catenary}）。此外，假设给定维数函数
$\delta : |B| \longrightarrow \mathbf{Z}$。
若 $X$ 是 $B$ 上的代数空间，且其结构态射 $f:X\to B$ 拟分离并局部有限型，
则称 $X/B$ 是{\it 良好的}。在这种情况下，定义
$$
\delta = \delta_{X/B} : |X| \longrightarrow \mathbf{Z}
$$
为把 $x$ 映到 $\delta(f(x))$ 与 $x/f(x)$ 的超越次数之和的映射
（《空间的态射》定义
\ref{spaces-morphisms-definition-dimension-fibre}）。由《空间的态射续篇》引理
\ref{spaces-more-morphisms-lemma-universally-catenary-dimension-function}，
这是一个维数函数。
\end{situation}

\noindent
一个特殊情形是 $S=B$ 为概形，且 $(S,\delta)$ 如《Chow 同调》情形
\ref{chow-situation-setup} 中所述。因此，$B$ 可以是一个域的谱
（《Chow 同调》例 \ref{chow-example-field}），也可以是
$B=\Spec(\mathbf{Z})$
（《Chow 同调》例 \ref{chow-example-domain-dimension-1}）。

\medskip\noindent
许多关于代数空间的引理、命题、定理和定义，在情形
\ref{spaces-chow-situation-setup} 的设定下更容易处理，因为我们所用的代数空间拟分离
（从而更是合宜的）且局部诺特。本章将穿插如下评注来指出这一点。% P10-E0036

\begin{remark}
\label{spaces-chow-remark-sober}
在情形 \ref{spaces-chow-situation-setup} 中，若 $X/B$ 良好，则 $|X|$ 是清醒拓扑空间；
见《空间的性质》引理 \ref{spaces-properties-lemma-quasi-separated-sober}
或《合宜空间》命题 \ref{decent-spaces-proposition-reasonable-sober}。
下文选取 $|X|$ 中不可约闭子集的泛点时，将直接使用这一事实而不再说明。
\end{remark}

\begin{remark}
\label{spaces-chow-remark-integral}
在情形 \ref{spaces-chow-situation-setup} 中，若 $X/B$ 良好，则 $X$ 为整代数空间
（《域上的空间》定义
\ref{spaces-over-fields-definition-integral-algebraic-space}）当且仅当 $X$ 既约且
$|X|$ 不可约。此外，对任意点 $\xi\in|X|$，存在唯一的整闭子空间
$Z\subset X$，使得 $\xi$ 是闭子集 $|Z|\subset|X|$ 的泛点；见
《域上的空间》引理
\ref{spaces-over-fields-lemma-decent-irreducible-closed}。
\end{remark}

\noindent
若 $B$ 是 Jacobson 的，并且 $\delta$ 把闭点映到零，则 $\delta$ 正是把每一点
映到其闭包之维数的函数。

\begin{lemma}
\label{spaces-chow-lemma-delta-is-dimension}
在情形 \ref{spaces-chow-situation-setup} 中，假设 $B$ 是 Jacobson 的，并且对 $|B|$ 的
每个闭点 $b$ 都有 $\delta(b)=0$。令 $X/B$ 良好。若 $Z\subset X$ 是
泛点为 $\xi\in|Z|$ 的整闭子空间，则下列整数相等：
\begin{enumerate}
\item $\delta(\xi) = \delta_{X/B}(\xi)$,
\item $\dim(|Z|)$,
\item $\text{codim}(\{z\}, |Z|)$ for $z \in |Z|$ closed,
\item 对任意闭点 $z\in|Z|$，$Z$ 在 $z$ 处局部环的维数；以及
\item $\dim(\mathcal{O}_{Z, \overline{z}})$ for $z \in |Z|$ closed.
\end{enumerate}
\end{lemma}

\begin{proof}
令 $X,Z,\xi$ 如引理所述。由于 $X$ 在 $B$ 上局部有限型，$X$ 是 Jacobson 的；
见《合宜空间》引理
\ref{decent-spaces-lemma-Jacobson-universally-Jacobson}。
于是，由《合宜空间》引理
\ref{decent-spaces-lemma-decent-Jacobson-ft-pts}，
$X_{\text{ft-pts}}\subset|X|$ 是闭点集。给定 $|Z|$ 中不可约闭子集的链
$T_0\supset\ldots\supset T_e$，由《空间的态射》引理
\ref{spaces-morphisms-lemma-enough-finite-type-points}，
$T_e\cap X_{\text{ft-pts}}$ 非空。因此，总可以假设这样的链终止于
$T_e=\{z\}$，其中 $z\in|Z|$ 为闭点。于是
$\dim(Z)=\sup_z\text{codim}(\{z\},|Z|)$，其中 $z$ 遍历 $|Z|$ 的闭点。
由《拓扑学》引理 \ref{topology-lemma-dimension-function-catenary}，有
$\text{codim}(\{z\},Z)=\delta(\xi)-\delta(z)$。
由《空间的态射》引理
\ref{spaces-morphisms-lemma-finite-type-points-morphism}，
$z$ 的像是 $B$ 的有限型点，即 $|B|$ 的闭点。由《空间的态射》引理
\ref{spaces-morphisms-lemma-jacobson-finite-type-points}，
$z/b$ 的超越次数为 $0$。结合假设，得到
$\delta(z)=\delta(b)=0$。因此有等式
$$
\dim(|Z|) = \text{codim}(\{z\}, Z) = \delta(\xi)
$$
对每个闭点 $z\in|Z|$ 都成立。最后，由《合宜空间》引理
\ref{decent-spaces-lemma-codimension-local-ring}，
$\text{codim}(\{z\},Z)$ 等于 $Z$ 在 $z$ 处局部环的维数；再由
《空间的性质》引理 \ref{spaces-properties-lemma-dimension-local-ring}，
后者等于 $\dim(\mathcal{O}_{Z,\overline{z}})$。
\end{proof}

\noindent
在上面引理的情形中，$\delta$ 在一个不可约闭子集之泛点处的值，
就是该不可约闭子集的维数。这引出如下定义。

\begin{definition}
\label{spaces-chow-definition-delta-dimension}
在情形 \ref{spaces-chow-situation-setup} 中，对任意良好的 $X/B$ 以及任意不可约闭子集
$T\subset|X|$，定义
$$
\dim_\delta(T) = \delta(\xi)
$$
其中 $\xi\in T$ 是 $T$ 的泛点。称它为 $T$ 的{\it $\delta$-维数}。
若 $T\subset|X|$ 是任意闭子集，则定义 $\dim_\delta(T)$ 为 $T$ 的各个
不可约分支之 $\delta$-维数的上确界。若 $Z$ 是 $X$ 的闭子空间，则置
$\dim_\delta(Z) = \dim_\delta(|Z|)$.
\end{definition}

\noindent
当然，这恰好意味着
$\dim_\delta(T) = \sup \{\delta(t) \mid t \in T\}$.







\section{代数闭链}
\label{spaces-chow-section-cycles}

\noindent
本节是《Chow 同调》第 \ref{chow-section-cycles} 节的代数空间版本。

\medskip\noindent
由于并未假设所考虑的空间拟紧，定义闭链时必须稍加小心。例如，
一个有理函数可能有无穷多个极点，因此必须允许无穷和。无论如何，
若 $X$ 拟紧，则闭链仍像通常那样是有限和。

\begin{definition}
\label{spaces-chow-definition-cycles}
在情形 \ref{spaces-chow-situation-setup} 中，令 $X/B$ 良好，并令
$k\in\mathbf{Z}$。
\begin{enumerate}
\item $X$ 上的{\it 闭链}是形式和
$$
\alpha = \sum n_Z [Z]
$$
其中求和遍历整闭子空间 $Z\subset X$，每个 $n_Z\in\mathbf{Z}$，并且
$\{|Z|;n_Z\not=0\}$ 是 $|X|$ 的局部有限子集族
（《拓扑学》定义 \ref{topology-definition-locally-finite}）。
\item $X$ 上的{\it $k$ 维闭链}是形如
$$
\alpha = \sum n_Z [Z]
$$
的闭链，其中 $n_Z\not=0\Rightarrow\dim_\delta(Z)=k$。
\item $X$ 上全体 $k$ 维闭链组成的阿贝尔群记为 $Z_k(X)$。
\end{enumerate}
\end{definition}

\noindent
换言之，$X$ 上的 $k$ 维闭链，是由 $\delta$-维数为 $k$ 的整闭子空间
（评注 \ref{spaces-chow-remark-integral}）组成的局部有限形式
$\mathbf{Z}$-线性组合。$k$ 维闭链 $\alpha=\sum n_Z[Z]$ 与
$\beta=\sum m_Z[Z]$ 的和定义为
$$
\alpha + \beta = \sum (n_Z + m_Z)[Z],
$$
即逐项相加系数。




\section{重数}
\label{spaces-chow-section-multiplicities}

\noindent
本节给出几个关于长度与重数的简单结果。

\begin{lemma}
\label{spaces-chow-lemma-length}
设 $S$ 为概形，$X$ 为 $S$ 上的代数空间，$\mathcal{F}$ 为拟凝聚
$\mathcal{O}_X$-模。令 $x\in|X|$，并令
$d\in\{0,1,2,\ldots,\infty\}$。下列条件等价：
\begin{enumerate}
\item
$\text{length}_{\mathcal{O}_{X, \overline{x}}} \mathcal{F}_{\overline{x}} = d$
\item 存在平展态射 $U\to X$，其中 $U$ 是概形，并存在映到 $x$ 的
$u\in U$，使得
$\text{length}_{\mathcal{O}_{U, u}} (\mathcal{F}|_U)_u = d$
\item 对任意平展态射 $U\to X$，其中 $U$ 是概形，以及任意映到 $x$ 的
$u\in U$，都有
$\text{length}_{\mathcal{O}_{U, u}} (\mathcal{F}|_U)_u = d$
\end{enumerate}
\end{lemma}

\begin{proof}
令 $U\to X$ 与 $u\in U$ 如 (2) 或 (3) 所述。于是
$\mathcal{O}_{X,\overline{x}}$ 是 $\mathcal{O}_{U,u}$ 的严格 Hensel 化，并且
$$
\mathcal{F}_{\overline{x}} =
(\mathcal{F}|_U)_u \otimes_{\mathcal{O}_{U, u}} \mathcal{O}_{X, \overline{x}}
$$
见《空间的性质》引理
\ref{spaces-properties-lemma-describe-etale-local-ring} 和
\ref{spaces-properties-lemma-stalk-quasi-coherent}。
由《代数学》引理 \ref{algebra-lemma-pullback-module}、
$\mathcal{O}_{U,u}\to\mathcal{O}_{X,\overline{x}}$ 的平坦性，以及
$\mathcal{O}_{X,\overline{x}}/\mathfrak m_u\mathcal{O}_{X,\overline{x}}$
等于 $\mathcal{O}_{X,\overline{x}}$ 的剩余域，便得到所需的长度相等。
关于严格 Hensel 化的这些事实见《代数续篇》引理
\ref{more-algebra-lemma-dumb-properties-henselization}。
\end{proof}

\begin{definition}
\label{spaces-chow-definition-length-at-x}
设 $S$ 为概形，$X$ 为 $S$ 上的代数空间，$\mathcal{F}$ 为拟凝聚
$\mathcal{O}_X$-模。令 $x\in|X|$，并令
$d\in\{0,1,2,\ldots,\infty\}$。若引理 \ref{spaces-chow-lemma-length} 中的
等价条件成立，则称{\it $\mathcal{F}$ 在 $x$ 处的长度为 $d$}。
\end{definition}

\begin{lemma}
\label{spaces-chow-lemma-length-closed-immersion}
设 $S$ 为概形，$i:Y\to X$ 为 $S$ 上代数空间的闭浸入，
$\mathcal{G}$ 为拟凝聚 $\mathcal{O}_Y$-模。令 $y\in|Y|$ 的像为
$x\in|X|$，并令 $d\in\{0,1,2,\ldots,\infty\}$。下列条件等价：
\begin{enumerate}
\item $\mathcal{G}$ 在 $y$ 处的长度为 $d$；
\item $i_*\mathcal{G}$ 在 $x$ 处的长度为 $d$。
\end{enumerate}
\end{lemma}

\begin{proof}
选取平展态射 $f:U\to X$，其中 $U$ 为概形，并选取映到 $x$ 的
$u\in U$。置 $V=Y\times_XU$，并把投影记为 $g:V\to Y$ 和
$j:V\to U$。那么 $j:V\to U$ 是闭浸入，并存在唯一的点 $v\in V$
同时映到 $y\in|Y|$ 与 $u\in U$（使用《空间的性质》引理
\ref{spaces-properties-lemma-points-cartesian} 和《空间》引理
\ref{spaces-lemma-base-change-immersions}）。作为概形 $V$ 上的模，% P10-E0037
有 $j_*(\mathcal{G}|_V)=(i_*\mathcal{G})|_U$，而 $j_*$ 是概形态射
$j$ 的“通常”模推出；见《空间的上同调》式
(\ref{spaces-cohomology-equation-representable-higher-direct-image})
附近的讨论。这样便约化到概形情形：若 $i:Y\to X$ 是概形的闭浸入，则
$$
(i_*\mathcal{G})_x = \mathcal{G}_y
$$
这是 $\mathcal{O}_{X,x}$-模的等式，其中右端的模结构由满射
$i_y^\sharp:\mathcal{O}_{X,x}\to\mathcal{O}_{Y,y}$ 给出。
于是由《代数学》引理 \ref{algebra-lemma-length-independent} 得到长度相等。
\end{proof}

\begin{lemma}
\label{spaces-chow-lemma-length-finite}
设 $S$ 为概形，$X$ 为 $S$ 上局部诺特的代数空间，$\mathcal{F}$ 为凝聚
$\mathcal{O}_X$-模，并令 $x\in|X|$。下列条件等价：
\begin{enumerate}
\item 存在平展态射 $U\to X$，其中 $U$ 为概形，并存在映到 $x$ 的
$u\in U$，使得 $u$ 是 $\text{Supp}(\mathcal{F}|_U)$ 的某个不可约分支的泛点；
\item 对任意平展态射 $U\to X$，其中 $U$ 为概形，以及任意映到 $x$ 的
$u\in U$，点 $u$ 都是 $\text{Supp}(\mathcal{F}|_U)$ 的某个不可约分支的泛点；
\item $\mathcal{F}$ 在 $x$ 处的长度有限且非零。
\end{enumerate}
若 $X$ 合宜（等价地，拟分离），则这些条件还等价于
\begin{enumerate}
\item[(4)] $x$ 是 $\text{Supp}(\mathcal{F})$ 的某个不可约分支的泛点。
\end{enumerate}
\end{lemma}

\begin{proof}
设 $f:U\to X$ 是平展态射，$U$ 为概形，并且 $u\in U$ 映到 $x$。
那么 $\mathcal{F}|_U=f^*\mathcal{F}$ 是局部诺特概形 $U$ 上的凝聚
$\mathcal{O}_U$-模，特别地，$(\mathcal{F}|_U)_u$ 是有限
$\mathcal{O}_{U,u}$-模；见《空间的上同调》引理
\ref{spaces-cohomology-lemma-coherent-Noetherian} 与《概形的上同调》引理
\ref{coherent-lemma-coherent-Noetherian}。回忆 $\mathcal{F}|_U$ 的支集是
$U$ 的闭子集（《态射》引理 \ref{morphisms-lemma-support-finite-type}），而
$(\mathcal{F}|_U)_u$ 的支集是 $\mathcal{F}|_U$ 的支集沿态射
$\Spec(\mathcal{O}_{U,u})\to U$ 的拉回。因此，$u$ 是
$\text{Supp}(\mathcal{F}|_U)$ 的某个不可约分支的泛点，当且仅当
$(\mathcal{F}|_U)_u$ 的支集等于 $\mathcal{O}_{U,u}$ 的极大理想。
现在，(1)、(2)、(3) 的等价性由《代数学》引理
\ref{algebra-lemma-support-point} 得出。% P10-E0038

\medskip\noindent
若 $X$ 合宜，则选取平展态射 $f:U\to X$ 以及映到 $x$ 的点
$u\in U$。$\mathcal{F}$ 的支集拉回为 $\mathcal{F}|_U$ 的支集；见
《空间的态射》引理 \ref{spaces-morphisms-lemma-support-finite-type}。
此外，$|X|$ 中的特化 $x'\leadsto x$ 可提升为 $U$ 中的特化
$u'\leadsto u$，而 $U$ 中任意非平凡特化 $u'\leadsto u$ 都映到
$|X|$ 中的非平凡特化 $f(u')\leadsto f(u)$；见《合宜空间》引理
\ref{decent-spaces-lemma-decent-specialization} 和
\ref{decent-spaces-lemma-decent-no-specializations-map-to-same-point}。
利用 $|X|$ 和 $U$ 都是清醒拓扑空间这一事实（《合宜空间》命题
\ref{decent-spaces-proposition-reasonable-sober} 与《概形》引理
\ref{schemes-lemma-scheme-sober}），可知 $x$ 是 $\mathcal{F}$ 支集的泛点，
当且仅当 $u$ 是 $\mathcal{F}|_U$ 支集的泛点。因此，(4) 等价于 (1)。

\medskip\noindent
括号中的陈述来自《合宜空间》引理
\ref{decent-spaces-lemma-locally-Noetherian-decent-quasi-separated}.
\end{proof}

\begin{lemma}
\label{spaces-chow-lemma-point-of-max-dimension}
在情形 \ref{spaces-chow-situation-setup} 中，令 $X/B$ 良好，$T\subset|X|$ 为闭子集，
且 $t\in T$。若 $\dim_\delta(T)\leq k$ 且 $\delta(t)=k$，则 $t$ 是
$T$ 的某个不可约分支的泛点。
\end{lemma}

\begin{proof}
点 $t$ 包含在某个不可约分支 $T'\subset T$ 中。令 $t'\in T'$ 为其泛点。
那么 $k\geq\delta(t')\geq\delta(t)$。由于 $\delta$ 是维数函数，
必有 $t=t'$。
\end{proof}



\section{闭子空间的伴随闭链}
\label{spaces-chow-section-cycle-of-closed-subscheme}

\noindent
本节是《Chow 同调》第
\ref{chow-section-cycle-of-closed-subscheme} 节的代数空间版本。

\begin{remark}
\label{spaces-chow-remark-irreducible-component}
在情形 \ref{spaces-chow-situation-setup} 中，令 $X/B$ 良好，并令 $Y\subset X$ 为
闭子空间。由评注 \ref{spaces-chow-remark-sober} 和 \ref{spaces-chow-remark-integral}，
下列三者之间存在一一对应：
\begin{enumerate}
\item $|Y|$ 的不可约分支 $T$；
\item $|Y|$ 的不可约分支的泛点；
\item 整闭子空间 $Z\subset Y$，满足 $|Z|$ 是 $|Y|$ 的不可约分支。
\end{enumerate}
本章称 (3) 中的 $Z$ 为 $Y$ 的{\it 不可约分支}，并称
$\xi\in|Z|$ 为该分支的{\it 泛点}。
\end{remark}

\begin{definition}
\label{spaces-chow-definition-cycle-associated-to-closed-subscheme}
在情形 \ref{spaces-chow-situation-setup} 中，令 $X/B$ 良好，并令
$Y\subset X$ 为闭子空间。
\begin{enumerate}
\item 对泛点为 $\xi$ 的不可约分支 $Z\subset Y$，称 $\mathcal{O}_Y$
在 $\xi$ 处的长度（定义 \ref{spaces-chow-definition-length-at-x}）为
{\it $Z$ 在 $Y$ 中的重数}。把引理 \ref{spaces-chow-lemma-length-finite}
应用于 $Y$ 上的 $\mathcal{O}_Y$，可知这是正整数。
\item 假设 $\dim_\delta(Y)\leq k$。定义{\it $Y$ 的伴随 $k$ 维闭链}为
$$
[Y]_k = \sum m_{Z, Y}[Z]
$$
其中求和遍历 $Y$ 的所有 $\delta$-维数为 $k$ 的不可约分支 $Z$，而
$m_{Z,Y}$ 是 $Z$ 在 $Y$ 中的重数。由《域上的空间》引理
\ref{spaces-over-fields-lemma-components-locally-finite}，
这是一个 $k$ 维闭链。
\end{enumerate}
\end{definition}

\noindent
务必注意，只有当 $Y$ 的 $\delta$-维数不超过 $k$ 时，我们才定义
$[Y]_k$。换言之，按照约定，一旦写出 $[Y]_k$，就隐含
$\dim_\delta(Y)\leq k$。







\section{凝聚层的伴随闭链}
\label{spaces-chow-section-cycle-of-coherent-sheaf}

\noindent
本节是《Chow 同调》第
\ref{chow-section-cycle-of-coherent-sheaf} 节的代数空间版本。

\begin{definition}
\label{spaces-chow-definition-cycle-associated-to-coherent-sheaf}
在情形 \ref{spaces-chow-situation-setup} 中，令 $X/B$ 良好，并令 $\mathcal{F}$ 为
凝聚 $\mathcal{O}_X$-模。
\begin{enumerate}
\item 设整闭子空间 $Z\subset X$ 的泛点为 $\xi$，并且 $|Z|$ 是
$\text{Supp}(\mathcal{F})$ 的不可约分支。称 $\mathcal{F}$ 在 $\xi$ 处的长度
（定义 \ref{spaces-chow-definition-length-at-x}）为{\it $Z$ 在 $\mathcal{F}$ 中的重数}。
由引理 \ref{spaces-chow-lemma-length-finite}，这是正整数。
\item 假设 $\dim_\delta(\text{Supp}(\mathcal{F}))\leq k$。定义
{\it $\mathcal{F}$ 的伴随 $k$ 维闭链}为
$$
[\mathcal{F}]_k = \sum m_{Z, \mathcal{F}}[Z]
$$
其中求和遍历整闭子空间 $Z\subset X$，它们对应于
$\text{Supp}(\mathcal{F})$ 中 $\delta$-维数为 $k$ 的不可约分支，而
$m_{Z,\mathcal{F}}$ 是 $Z$ 在 $\mathcal{F}$ 中的重数。由《域上的空间》引理
\ref{spaces-over-fields-lemma-components-locally-finite}，
这是一个 $k$ 维闭链。
\end{enumerate}
\end{definition}

\noindent
务必注意，只有当 $\mathcal{F}$ 凝聚且
$\text{Supp}(\mathcal{F})$ 的 $\delta$-维数不超过 $k$ 时，我们才定义
$[\mathcal{F}]_k$。换言之，按照约定，一旦写出 $[\mathcal{F}]_k$，就隐含
$\mathcal{F}$ 是 $X$ 上的凝聚模，并且
$\dim_\delta(\text{Supp}(\mathcal{F}))\leq k$。

\begin{lemma}
\label{spaces-chow-lemma-reformulate-coeff-coherent}
在情形 \ref{spaces-chow-situation-setup} 中，令 $X/B$ 良好。设 $\mathcal{F}$ 为凝聚
$\mathcal{O}_X$-模，且
$\dim_\delta(\text{Supp}(\mathcal{F}))\leq k$。令 $Z$ 为 $X$ 的整闭子空间，
满足 $\dim_\delta(Z)=k$，并令 $\xi\in|Z|$ 为其泛点。那么，
$[\mathcal{F}]_k$ 中 $Z$ 的系数等于 $\mathcal{F}$ 在 $\xi$ 处的长度。
\end{lemma}

\begin{proof}
注意，$|Z|$ 是 $\text{Supp}(\mathcal{F})$ 的不可约分支，当且仅当
$\xi\in\text{Supp}(\mathcal{F})$；见引理
\ref{spaces-chow-lemma-point-of-max-dimension}。此外，若
$\xi\not \in\text{Supp}(\mathcal{F})$，则 $\mathcal{F}$ 在 $\xi$ 处的长度为零。
结合定义 \ref{spaces-chow-definition-cycle-associated-to-coherent-sheaf} 即得结论。
\end{proof}

\begin{lemma}
\label{spaces-chow-lemma-cycle-closed-coherent}
在情形 \ref{spaces-chow-situation-setup} 中，令 $X/B$ 良好，并令 $Y\subset X$ 为
闭子空间。若 $\dim_\delta(Y)\leq k$，则
$[Y]_k=[i_*\mathcal{O}_Y]_k$，其中 $i:Y\to X$ 为包含态射。
\end{lemma}

\begin{proof}
令 $Z$ 为 $X$ 的整闭子空间，且 $\dim_\delta(Z)=k$。若
$Z\not\subset Y$，则 $Z$ 在 $[Y]_k$ 和 $[i_*\mathcal{O}_Y]_k$ 中的系数
都为零。% P10-E0039
若 $Z\subset Y$，则可把 $Z$ 的泛点视为 $y\in|Y|$，其像为
$x\in|X|$。于是，$[Y]_k$ 中 $Z$ 的系数是 $\mathcal{O}_Y$ 在 $y$ 处的
长度，而 $[i_*\mathcal{O}_Y]_k$ 中 $Z$ 的系数是
$i_*\mathcal{O}_Y$ 在 $x$ 处的长度。系数相等便由引理
\ref{spaces-chow-lemma-length-closed-immersion} 得出。
\end{proof}

\begin{lemma}
\label{spaces-chow-lemma-additivity-sheaf-cycle}
在情形 \ref{spaces-chow-situation-setup} 中，令 $X/B$ 良好。设
$0\to\mathcal{F}\to\mathcal{G}\to\mathcal{H}\to0$ 为凝聚
$\mathcal{O}_X$-模的短正合列。假设 $\mathcal{F},\mathcal{G},\mathcal{H}$
的支集之 $\delta$-维数都不超过 $k$。那么
$[\mathcal{G}]_k=[\mathcal{F}]_k+[\mathcal{H}]_k$。
\end{lemma}

\begin{proof}
令 $Z$ 为 $X$ 的整闭子空间，且 $\dim_\delta(Z)=k$。只需证明 $Z$ 在
$[\mathcal{G}]_k,[\mathcal{F}]_k,[\mathcal{H}]_k$ 中的系数满足相应的
可加性。由引理 \ref{spaces-chow-lemma-reformulate-coeff-coherent}，只需证明
$$
\mathcal{G}\text{ 在 }x\text{ 处的长度} =
\mathcal{F}\text{ 在 }x\text{ 处的长度} +
\mathcal{H}\text{ 在 }x\text{ 处的长度}
$$
对任意 $x\in|X|$ 成立。根据定义 \ref{spaces-chow-definition-length-at-x}，
这立即来自长度的可加性；见《代数学》引理
\ref{algebra-lemma-length-additive}。
\end{proof}





\section{固有推出的准备}
\label{spaces-chow-section-preparation-pushforward}

\noindent
本节是《Chow 同调》第
\ref{chow-section-preparation-pushforward} 节的代数空间版本。

\begin{lemma}
\label{spaces-chow-lemma-proper-image}
在情形 \ref{spaces-chow-situation-setup} 中，令 $X,Y/B$ 良好，并令 $f:X\to Y$
为 $B$ 上的态射。若 $Z\subset X$ 为整闭子空间，则存在唯一的整闭子空间
$Z'\subset Y$，使得存在交换图表
$$
\xymatrix{
Z \ar[r] \ar[d] & X \ar[d]^f \\
Z' \ar[r] & Y
}
$$
且 $Z\to Z'$ 支配。若 $f$ 固有，则 $Z\to Z'$ 固有且满射。
\end{lemma}

\begin{proof}
令 $\xi\in|Z|$ 为泛点，并令 $Z'\subset Y$ 为泛点
$\xi'=f(\xi)$ 的整闭子空间；见评注 \ref{spaces-chow-remark-integral}。由
《空间的性质》引理 \ref{spaces-properties-lemma-points-cartesian}，有
$\xi\in|f^{-1}(Z')|=|f|^{-1}(|Z'|)$；又由于 $Z$ 既约且
$|Z|=\overline{\{\xi\}}$，% P10-E0040
可知 $Z\subset f^{-1}(Z')$ 是 $X$ 的闭子空间（见《空间的性质》引理
\ref{spaces-properties-lemma-map-into-reduction}）。于是得到态射 $Z\to Z'$。
该态射支配，因为 $Z$ 的泛点映到 $Z'$ 的泛点。$Z'$ 的唯一性显然。
若 $f$ 固有，则 $Z\to Y$ 作为固有态射的复合而固有（《空间的态射》引理
\ref{spaces-morphisms-lemma-base-change-proper} 和
\ref{spaces-morphisms-lemma-closed-immersion-proper}）。再由
《空间的态射》引理
\ref{spaces-morphisms-lemma-universally-closed-permanence}，
$Z\to Z'$ 固有。由于固有态射的像是闭集，随即得到满射性。
\end{proof}

\begin{remark}
\label{spaces-chow-remark-residue-field}
在情形 \ref{spaces-chow-situation-setup} 中，令 $X/B$ 良好。每个 $x\in|X|$ 都可由
一个（唯一的）单态射 $\Spec(k)\to X$ 表示，其中 $k$ 为域；见
《合宜空间》引理 \ref{decent-spaces-lemma-decent-points-monomorphism}。
此时，$k$ 称为 $x$ 的{\it 剩余域}，记为 $\kappa(x)$。回忆 $X$ 有一个
稠密开子概形 $U\subset X$（《空间的性质》命题
\ref{spaces-properties-proposition-locally-quasi-separated-open-dense-scheme}）。
若 $x\in U$，则 $\kappa(x)$ 与把 $U$ 视为概形时 $x$ 的剩余域一致。
见《合宜空间》第
\ref{decent-spaces-section-residue-fields-henselian-local-rings} 节。
\end{remark}

\begin{remark}
\label{spaces-chow-remark-function-field}
在情形 \ref{spaces-chow-situation-setup} 中，令 $X/B$ 良好，并假设 $X$ 为整空间。
在这种情况下，$X$ 的{\it 函数域} $R(X)$ 有定义，并等于 $X$ 在其泛点处的
剩余域。
见《域上的空间》定义
\ref{spaces-over-fields-definition-function-field}。
结合评注 \ref{spaces-chow-remark-integral} 可知，对任意 $x\in X$，剩余域
$\kappa(x)$ 是以 $x$ 为泛点的唯一整闭子空间 $Z\subset X$ 的函数域。
\end{remark}

\begin{lemma}
\label{spaces-chow-lemma-equal-dimension}
在情形 \ref{spaces-chow-situation-setup} 中，令 $X,Y/B$ 良好，并令 $f:X\to Y$
为 $B$ 上的态射。假设 $X,Y$ 都是整空间且
$\dim_\delta(X)=\dim_\delta(Y)$。那么，或者 $f$ 通过 $Y$ 的某个真闭子空间
分解，或者 $f$ 支配且函数域扩张 $R(X)/R(Y)$ 有限。
\end{lemma}

\begin{proof}
由引理 \ref{spaces-chow-lemma-proper-image}，存在唯一的整闭子空间 $Z\subset Y$，
使得 $f$ 通过支配态射 $X\to Z$ 分解。此时，$Z=Y$ 当且仅当
$\dim_\delta(Z)=\dim_\delta(Y)$。另一方面，由维数函数的构造（见情形
\ref{spaces-chow-situation-setup}），有
$\dim_\delta(X)=\dim_\delta(Z)+r$，其中 $r$ 是扩张 $R(X)/R(Z)$ 的
超越次数。% P10-E0041
结合《域上的空间》引理 \ref{spaces-over-fields-lemma-finite-degree} 即得结论。
\end{proof}

\begin{lemma}
\label{spaces-chow-lemma-quasi-compact-locally-finite}
在情形 \ref{spaces-chow-situation-setup} 中，令 $X,Y/B$ 良好，并令 $f:X\to Y$
为 $B$ 上的态射。假设 $f$ 拟紧，且 $\{T_i\}_{i\in I}$ 是 $|X|$ 中闭子集的
局部有限族。那么，$\{\overline{|f|(T_i)}\}_{i\in I}$ 是 $|Y|$ 中闭子集的
局部有限族。
\end{lemma}

\begin{proof}
令 $V\subset|Y|$ 为拟紧开子集。由《空间的态射》引理
\ref{spaces-morphisms-lemma-quasi-compact-is-quasi-compact}，
$|f|^{-1}(V)\subset|X|$ 拟紧。因此，集合
$\{i \in I : T_i \cap |f|^{-1}(V) \not = \emptyset \}$
由一个略去的简单拓扑论证可知是有限的。由于它与下列集合相同：
$$
\{i \in I : |f|(T_i) \cap V \not = \emptyset \} =
\{i \in I : \overline{|f|(T_i)} \cap V \not = \emptyset \}
$$
引理得证。
\end{proof}










\section{固有推出}
\label{spaces-chow-section-proper-pushforward}

\noindent
本节是《Chow 同调》第
\ref{chow-section-proper-pushforward} 节的代数空间版本。

\begin{definition}
\label{spaces-chow-definition-proper-pushforward}
在情形 \ref{spaces-chow-situation-setup} 中，令 $X,Y/B$ 良好，并令 $f:X\to Y$
为 $B$ 上的态射。假设 $f$ 固有。
\begin{enumerate}
\item 令 $Z\subset X$ 为整闭子空间，且 $\dim_\delta(Z)=k$。
令 $Z'\subset Y$ 为引理 \ref{spaces-chow-lemma-proper-image} 意义下 $Z$ 的像。定义
$$
f_*[Z] =
\left\{
\begin{matrix}
0 & \text{若} & \dim_\delta(Z')< k, \\
\deg(Z/Z') [Z'] & \text{若} & \dim_\delta(Z') = k.
\end{matrix}
\right.
$$
由引理 \ref{spaces-chow-lemma-equal-dimension} 和《域上的空间》定义
\ref{spaces-over-fields-definition-degree}，若
$\dim_\delta(Z')=\dim_\delta(Z)$，则 $Z$ 在 $Z'$ 上的次数有定义且有限。
\item 令 $\alpha=\sum n_Z[Z]$ 为 $X$ 上的 $k$ 维闭链。定义 $\alpha$ 的
{\it 推出}为下列和：% P10-E0042
$$
f_* \alpha = \sum n_Z f_*[Z]
$$
其中每个 $f_*[Z]$ 均如上定义。由引理
\ref{spaces-chow-lemma-quasi-compact-locally-finite}，该和局部有限。
\end{enumerate}
\end{definition}

\noindent
按照定义，闭链的固有推出
$$
f_* : Z_k(X) \longrightarrow Z_k(Y)
$$
是阿贝尔群同态。它使 $X\mapsto Z_k(X)$ 成为如下范畴上的协变函子：
对象是 $B$ 上良好的代数空间，态射是 $B$ 上的固有态射。% P10-E0043

\begin{lemma}
\label{spaces-chow-lemma-compose-pushforward}
在情形 \ref{spaces-chow-situation-setup} 中，令 $X,Y,Z/B$ 良好，并令
$f:X\to Y$ 与 $g:Y\to Z$ 为 $B$ 上的固有态射。那么，作为映射
$Z_k(X)\to Z_k(Z)$，有 $g_*\circ f_*=(g\circ f)_*$。
\end{lemma}

\begin{proof}
令 $W\subset X$ 为维数 $k$ 的整闭子空间。先把引理
\ref{spaces-chow-lemma-proper-image} 应用于 $f$ 与 $W$，再应用于 $g$ 与所得的 $W'$，
得到整闭子空间 $W'\subset Y$ 和 $W''\subset Z$。那么
$W\to W'$ 与 $W'\to W''$ 都是固有满射。需要证明
$g_*(f_*[W])=(f\circ g)_*[W]$。% P10-E0044
若 $\dim_\delta(W'')<k$，则两端都为零。若
$\dim_\delta(W'')=k$，则 $W\to W'$ 和 $W'\to W''$ 都满足引理
\ref{spaces-chow-lemma-equal-dimension} 的假设。因此
$$
g_*(f_*[W]) = \deg(W/W')\deg(W'/W'')[W''],
\quad
(f \circ g)_*[W] = \deg(W/W'')[W''].
$$
应用《域上的空间》引理
\ref{spaces-over-fields-lemma-degree-composition} 即得结论。
\end{proof}

\begin{lemma}
\label{spaces-chow-lemma-cycle-push-sheaf}
在情形 \ref{spaces-chow-situation-setup} 中，令 $f:X\to Y$ 为 $B$ 上良好代数空间的
固有态射。
\begin{enumerate}
\item 令 $Z\subset X$ 为闭子空间，且 $\dim_\delta(Z)\leq k$。那么
$$
f_*[Z]_k = [f_*{\mathcal O}_Z]_k.
$$
\item 令 $\mathcal{F}$ 为 $X$ 上的凝聚层，且
$\dim_\delta(\text{Supp}(\mathcal{F}))\leq k$。那么
$$
f_*[\mathcal{F}]_k = [f_*{\mathcal F}]_k.
$$
\end{enumerate}
注意，该陈述有意义，因为由《空间的上同调》引理
\ref{spaces-cohomology-lemma-proper-pushforward-coherent}，
$f_*\mathcal{F}$ 和 $f_*\mathcal{O}_Z$ 都是凝聚 $\mathcal{O}_Y$-模。
\end{lemma}

\begin{proof}
(1) 由 (2) 与引理 \ref{spaces-chow-lemma-cycle-closed-coherent} 得出。令
$\mathcal{F}$ 为 $X$ 上的凝聚层，并假设
$\dim_\delta(\text{Supp}(\mathcal{F}))\leq k$。由《空间的上同调》引理
\ref{spaces-cohomology-lemma-coherent-support-closed}，
存在闭浸入 $i:Z\to X$ 和凝聚 $\mathcal{O}_Z$-模 $\mathcal{G}$，使得
$i_*\mathcal{G}\cong\mathcal{F}$，且 $\mathcal{F}$ 的支集为 $Z$。
令 $Z'\subset Y$ 为 $f|_Z:Z\to Y$ 的概形论像；见《空间的态射》定义
\ref{spaces-morphisms-definition-scheme-theoretic-image}。考虑交换图表
$$
\xymatrix{
Z \ar[r]_i \ar[d]_{f|_Z} &
X \ar[d]^f \\
Z' \ar[r]^{i'} & Y
}
$$
这是 $B$ 上代数空间的图表。注意，$f|_Z$ 满射（由《空间的态射》引理
\ref{spaces-morphisms-lemma-quasi-compact-scheme-theoretic-image}
及 $|f|$ 为闭映射得出）且固有（由《空间的态射》引理
\ref{spaces-morphisms-lemma-base-change-proper}、
\ref{spaces-morphisms-lemma-closed-immersion-proper} 和
\ref{spaces-morphisms-lemma-universally-closed-permanence} 得出）。
沿图表的两条路径计算，得到
$f_*\mathcal{F}=f_*i_*\mathcal{G}=i'_*(f|_Z)_*\mathcal{G}$。
假设已经知道结论对闭浸入以及对 $f|_Z$ 成立。那么
$$
f_*[\mathcal{F}]_k = f_*i_*[\mathcal{G}]_k
= (i')_*(f|_Z)_*[\mathcal{G}]_k =
(i')_*[(f|_Z)_*\mathcal{G}]_k =
[(i')_*(f|_Z)_*\mathcal{G}]_k = [f_*\mathcal{F}]_k
$$
正是所需等式。闭浸入情形由引理
\ref{spaces-chow-lemma-length-closed-immersion} 和各定义得出。因此，问题约化为
$\dim_\delta(X)\leq k$ 且 $f:X\to Y$ 为固有满射的情形。

\medskip\noindent
假设 $\dim_\delta(X)\leq k$，且 $f:X\to Y$ 固有满射。对每个泛点为
$\eta$ 的不可约分支 $Z\subset Y$，存在 $\xi\in X$ 使得
$f(\xi)=\eta$。因此 $\delta(\eta)\leq\delta(\xi)\leq k$。所以在表达式
$$
f_*[\mathcal{F}]_k = \sum n_Z[Z],
\quad
\text{且}
\quad
[f_*\mathcal{F}]_k = \sum m_Z[Z].
$$
中，只要 $n_Z\not=0$ 或 $m_Z\not=0$，整闭子空间 $Z$ 实际上就是 $Y$
的一个 $\delta$-维数为 $k$ 的不可约分支（见引理
\ref{spaces-chow-lemma-point-of-max-dimension}）。选取这样的整闭子空间
$Z\subset Y$，并把其泛点记为 $\eta$。注意，对任意满足
$f(\xi)=\eta$ 的 $\xi\in X$，都有 $\delta(\xi)\geq k$，从而 $\xi$ 也是
$X$ 的某个 $\delta$-维数为 $k$ 的不可约分支的泛点（见引理
\ref{spaces-chow-lemma-point-of-max-dimension}）。由《域上的空间》引理
\ref{spaces-over-fields-lemma-finite-in-codim-1}，存在开子空间
$\eta\in V\subset Y$，使得 $f^{-1}(V)\to V$ 有限。由于 $\eta$ 是
$|Y|$ 某个不可约分支的泛点，可以假设 $V$ 为仿射概形；见
《空间的性质》命题
\ref{spaces-properties-proposition-locally-quasi-separated-open-dense-scheme}。
以 $V$ 代替 $Y$、以 $f^{-1}(V)$ 代替 $X$，便约化到 $Y$ 仿射且 $f$ 有限的
情形。特别地，$X$ 与 $Y$ 都是概形，于是约化到概形的相应结果；见
《Chow 同调》引理 \ref{chow-lemma-cycle-push-sheaf}
（取 $S=Y$）。
\end{proof}













\section{平坦拉回的准备}
\label{spaces-chow-section-preparation-flat-pullback}

\noindent
本节对应于《Chow 同调》第
\ref{chow-section-preparation-flat-pullback} 节。

\medskip\noindent
回顾：若在源与目标上平展局部地得到一个相对维数为 $r$ 的概形态射，
则称一个代数空间态射具有相对维数 $r$。精确定义与此等价，但实际上稍有不同；
见《空间的态射》定义
\ref{spaces-morphisms-definition-relative-dimension}。

\begin{lemma}
\label{spaces-chow-lemma-flat-inverse-image-dimension}
在情形 \ref{spaces-chow-situation-setup} 中，设 $X,Y/B$ 是良好的。
设 $f:X\to Y$ 是 $B$ 上的态射，并假设 $f$ 平坦且相对维数为 $r$。
对任意闭子集 $T\subset |Y|$，只要 $|f|^{-1}(T)$ 非空，就有
$$
\dim_\delta(|f|^{-1}(T)) = \dim_\delta(T) + r.
$$
若 $Z\subset Y$ 是整闭子概形，且 $Z'\subset f^{-1}(Z)$ 是一个不可约分支，
则 $Z'$ 支配 $Z$，并且
$\dim_\delta(Z')=\dim_\delta(Z)+r$。
\end{lemma}

\begin{proof}
闭子集的 $\delta$-维数是其各不可约分支的 $\delta$-维数的上确界，
故只需证明最后一个断言。可以用整闭子概形 $Z$ 代替 $Y$，并用
$f^{-1}(Z)=Z\times_YX$ 代替 $X$。因而可以假设 $Z=Y$ 是整的，
且 $f$ 是相对维数为 $r$ 的平坦态射。由于 $Y$ 局部诺特，局部有限型的
态射 $f$ 实际上局部有限呈示。因此《空间的态射》引理
\ref{spaces-morphisms-lemma-fppf-open} 适用，从而 $f$ 是开的。

设 $\xi\in X$ 是 $X$ 的某个不可约分支的泛点。由 $f$ 的开性可知，
$f(\xi)$ 是 $Z=Y$ 的泛点 $\eta$。因此 $Z'$ 支配 $Z=Y$。
最后，由《空间的性质》命题
\ref{spaces-properties-proposition-locally-quasi-separated-open-dense-scheme}
可知 $\xi$ 与 $\eta$ 分别位于 $X$ 与 $Y$ 的概形轨迹中。
由于 $\xi$ 是 $X$ 的一个泛点，
$\mathcal{O}_{X,\xi}=\mathcal{O}_{X_\eta,\xi}$ 只有一个素理想，
故其维数为 $0$（由于 $\xi$ 与 $\eta$ 位于 $X$ 与 $Y$ 的概形轨迹中，
这里可以使用通常的局部环）。于是由《空间的态射》引理
\ref{spaces-morphisms-lemma-dimension-fibre-at-a-point}
（以及给定相对维数之态射的定义）可知，$\kappa(\xi)$ 在
$\kappa(\eta)$ 上的超越次数为 $r$。换言之，
$\delta(\xi)=\delta(\eta)+r$，这正是所需结论。
\end{proof}

\noindent
下面的引理将用来证明：闭子概形的局部有限族在平坦拉回后仍局部有限。

\begin{lemma}
\label{spaces-chow-lemma-inverse-image-locally-finite}
在情形 \ref{spaces-chow-situation-setup} 中，设 $X,Y/B$ 是良好的，
且 $f:X\to Y$ 是 $B$ 上的态射。假设
$\{T_i\}_{i\in I}$ 是 $|Y|$ 的闭子集的一个局部有限族。
则 $\{|f|^{-1}(T_i)\}_{i\in I}$ 是 $X$ 的闭子集的一个局部有限族。
\end{lemma}

\begin{proof}
设 $U\subset |X|$ 是拟紧开子集。由于像
$|f|(U)\subset |Y|$ 是拟紧子集，存在拟紧开集 $V\subset |Y|$，
使得 $|f|(U)\subset V$。注意
$$
\{i \in I : |f|^{-1}(T_i) \cap U \not = \emptyset \}
\subset
\{i \in I : T_i \cap V \not = \emptyset \}.
$$
依假设，右端集合是有限的，故得证。
\end{proof}

















\section{平坦拉回}
\label{spaces-chow-section-flat-pullback}

\noindent
本节对应于《Chow 同调》第 \ref{chow-section-flat-pullback} 节。

\medskip\noindent
设 $S$ 是概形，且 $f:X\to Y$ 是 $S$ 上代数空间的态射。
设 $Z\subset Y$ 是闭子空间。本章有时把《空间的态射》定义
\ref{spaces-morphisms-definition-inverse-image-closed-subspace}
中构造的 $Z$ 的逆像 $f^{-1}(Z)$ 称为
{\it 概形论逆像}。这一概形论逆像是纤维积
$$
\xymatrix{
f^{-1}(Z) \ar[r] \ar[d] & X \ar[d] \\
Z \ar[r] & Y
}
$$
若 $\mathcal{I}\subset\mathcal{O}_Y$ 是 $Y$ 中与 $Z$ 对应的拟凝聚理想层，
则 $f^{-1}(\mathcal{I})\mathcal{O}_X$ 是 $X$ 中与 $f^{-1}(Z)$ 对应的
拟凝聚理想层。

\begin{definition}
\label{spaces-chow-definition-flat-pullback}
在情形 \ref{spaces-chow-situation-setup} 中，设 $X,Y/B$ 是良好的，
且 $f:X\to Y$ 是 $B$ 上的态射。假设 $f$ 平坦且相对维数为 $r$。
\begin{enumerate}
\item 设 $Z\subset Y$ 是 $\delta$-维数为 $k$ 的整闭子空间。
把 $f^*[Z]$ 定义为 $X$ 上与概形论逆像相伴的 $(k+r)$-闭链：
$$
f^*[Z] = [f^{-1}(Z)]_{k+r}.
$$
由引理 \ref{spaces-chow-lemma-flat-inverse-image-dimension}，
$\dim_\delta(f^{-1}(Z))=k+r$，故此定义有意义。
\item 设 $\alpha=\sum n_i[Z_i]$ 是 $Y$ 上的 $k$-闭链。
把和式
$$
f^* \alpha = \sum n_i f^*[Z_i]
$$
称为 $\alpha$ 沿 $f$ 的{\it 平坦拉回}，其中每个 $f^*[Z_i]$ 如上定义。
由引理 \ref{spaces-chow-lemma-inverse-image-locally-finite}，这个和式局部有限。
\item 由此得到的阿贝尔群同态记作
$f^*:Z_k(Y)\to Z_{k+r}(X)$。
\end{enumerate}
\end{definition}

\noindent
开浸入是平坦的。这是平坦态射的一个虽平凡却重要的特例。
若 $U\subset X$ 是开的，则闭链沿 $j:U\to X$ 的拉回有时称为该闭链
在 $U$ 上的{\it 限制}。注意在此情形中，所有映射
$$
j^* : Z_k(X) \longrightarrow Z_k(U)
$$
都是{\it 满射}。事实上，给定任意整闭子空间 $Z'\subset U$，
可以取 $Z'$ 在 $X$ 中的闭包 $Z$，并把它看作 $X$ 的既约闭子空间
（见《空间的性质》定义
\ref{spaces-properties-definition-reduced-induced-space}）。
% P10-E0046
显然 $Z\cap U=Z'$，换言之 $j^*[Z]=[Z']$，从而得到满射性。
事实上还有稍强的结论。

\begin{lemma}
\label{spaces-chow-lemma-exact-sequence-open}
在情形 \ref{spaces-chow-situation-setup} 中，设 $X/B$ 是良好的，且 $U\subset X$
是开子空间。设 $Y$ 是满足 $|Y|=|X|\setminus |U|$ 的 $X$ 的既约闭子空间，
并以 $i:Y\to X$ 表示包含态射。对每个 $k\in\mathbf{Z}$，序列
$$
\xymatrix{
Z_k(Y) \ar[r]^{i_*} & Z_k(X) \ar[r]^{j^*} & Z_k(U) \ar[r] & 0
}
$$
是阿贝尔群的正合复形。
\end{lemma}

\begin{proof}
上面已经证明了 $j^*$ 的满射性。先假设 $X$ 拟紧。
此时 $Z_k(X)$ 是自由 $\mathbf{Z}$-模，其基由元素 $[Z]$ 给出，
其中 $Z\subset X$ 是 $\delta$-维数为 $k$ 的整闭子空间。
这样的基元素或者映到 $Z_k(U)$ 的基元素 $[Z\cap U]$，
或者在 $Z\subset Y$ 时映到零。因此在这种情形中引理显然成立。
一般情形类似，略去证明。
\end{proof}

\begin{lemma}
\label{spaces-chow-lemma-etale-pullback}
在情形 \ref{spaces-chow-situation-setup} 中，设 $f:X\to Y$ 是 $B$ 上良好代数空间之间的
平展态射。若 $Z\subset Y$ 是整闭子空间，则
$f^*[Z]=\sum[Z']$，其中求和遍历 $f^{-1}(Z)$ 的各不可约分支
（注 \ref{spaces-chow-remark-irreducible-component}）。
\end{lemma}

\begin{proof}
引理的含义是 $[Z']$ 的系数为 $1$。这是因为 $f^{-1}(Z)$ 在整代数空间
$Z$ 上平展，因而是既约代数空间。
\end{proof}

\begin{lemma}
\label{spaces-chow-lemma-compose-flat-pullback}
在情形 \ref{spaces-chow-situation-setup} 中，设 $X,Y,Z/B$ 是良好的。
设 $f:X\to Y$ 与 $g:Y\to Z$ 是 $B$ 上相对维数分别为 $r$ 与 $s$
的平坦态射。则 $g\circ f$ 平坦且相对维数为 $r+s$，并且
$$
f^* \circ g^* = (g \circ f)^*
$$
作为映射 $Z_k(Z)\to Z_{k+r+s}(X)$ 相等。
\end{lemma}

\begin{proof}
由《空间的态射》引理
\ref{spaces-morphisms-lemma-dimension-fibre-at-a-point-additive} 与
\ref{spaces-morphisms-lemma-composition-flat}，复合态射平坦且相对维数为
$r+s$。假设
\begin{enumerate}
\item $A\subset Z$ 是 $\delta$-维数为 $k$ 的整闭子空间；
\item $A'\subset Y$ 是 $\delta$-维数为 $k+s$ 的整闭子空间，且
$A'\subset g^{-1}(A)$；以及
\item $A''\subset Y$ 是 $\delta$-维数为 $k+s+r$ 的整闭子空间，且
$A''\subset f^{-1}(W')$。
% P10-E0047
% P10-E0048
\end{enumerate}
必须证明 $[A'']$ 在 $(g\circ f)^*[A]$ 中的系数 $n$ 等于
$[A'']$ 在 $f^*(g^*[A])$ 中的系数 $m$。可以选取交换图
$$
\xymatrix{
U \ar[d] \ar[r] & V \ar[d] \ar[r] & W \ar[d] \\
X \ar[r] & Y \ar[r] & Z
}
$$
其中 $U,V,W$ 是概形，竖直箭头均平展，并且存在点
$u\in U$、$v\in V$、$w\in W$，使得 $u\mapsto v\mapsto w$，
且 $u,v,w$ 分别映到 $A'',A',A$ 的泛点。（略去细节。）
于是有平坦局部环同态
% P10-E0049
$\mathcal{O}_{W, w} \to \mathcal{O}_{V, v}$,
$\mathcal{O}_{V, v} \to \mathcal{O}_{U, u}$,
并且反复使用引理 \ref{spaces-chow-lemma-length} 得到
$$
n = \text{length}_{\mathcal{O}_{U, u}}(
\mathcal{O}_{U, u}/\mathfrak m_w\mathcal{O}_{U, u})
$$
以及
$$
m =
\text{length}_{\mathcal{O}_{V, v}}(
\mathcal{O}_{V, v}/\mathfrak m_w\mathcal{O}_{V, v})
\text{length}_{\mathcal{O}_{U, u}}(
\mathcal{O}_{U, u}/\mathfrak m_v\mathcal{O}_{U, u})
$$
因此所需等式由《代数》引理
\ref{algebra-lemma-pullback-transitive} 得出。
\end{proof}

\begin{lemma}
\label{spaces-chow-lemma-pullback-coherent}
在情形 \ref{spaces-chow-situation-setup} 中，设 $X,Y/B$ 是良好的，并设
$f:X\to Y$ 是相对维数为 $r$ 的平坦态射。
\begin{enumerate}
\item 设 $Z\subset Y$ 是满足 $\dim_\delta(Z)\leq k$ 的闭子空间。
则 $\dim_\delta(f^{-1}(Z))\leq k+r$，且在 $Z_{k+r}(X)$ 中有
$[f^{-1}(Z)]_{k+r}=f^*[Z]_k$。
\item 设 $\mathcal{F}$ 是 $Y$ 上满足
$\dim_\delta(\text{Supp}(\mathcal{F}))\leq k$ 的凝聚层。
则 $\dim_\delta(\text{Supp}(f^*\mathcal{F}))\leq k+r$，且
$$
f^*[{\mathcal F}]_k = [f^*{\mathcal F}]_{k+r}
$$
在 $Z_{k+r}(X)$ 中成立。
\end{enumerate}
\end{lemma}

\begin{proof}
由引理 \ref{spaces-chow-lemma-cycle-closed-coherent} 以及
$f^*\mathcal{O}_Z=\mathcal{O}_{f^{-1}(Z)}$，(1) 可由 (2) 推出。

\medskip\noindent
证明 (2)。由于 $X,Y$ 局部诺特，可以应用《空间的上同调》引理
\ref{spaces-cohomology-lemma-coherent-Noetherian} 得知 $\mathcal{F}$ 有限型；
于是 $f^*\mathcal{F}$ 有限型（《位点上的模》引理
\ref{sites-modules-lemma-local-pullback}），从而 $f^*\mathcal{F}$ 凝聚
（再次使用《空间的上同调》引理
\ref{spaces-cohomology-lemma-coherent-Noetherian}）。
因此引理中的表达式确有意义。设 $W\subset Y$ 是 $\delta$-维数为 $k$
的整闭子空间，并设 $W'\subset X$ 是维数为 $k+r$ 且在 $f$ 下映入
$W$ 的整闭子空间。必须证明 $[W']$ 在 $f^*[{\mathcal F}]_k$ 中的系数
$n$ 等于 $[W']$ 在 $[f^*{\mathcal F}]_{k+r}$ 中的系数 $m$。可以选取交换图
$$
\xymatrix{
U \ar[d] \ar[r] & V \ar[d] \\
X \ar[r] & Y
}
$$
其中 $U,V$ 是概形，竖直箭头均平展，并且存在点 $u\in U$、$v\in V$，
使得 $u\mapsto v$，且 $u,v$ 分别映到 $W',W$ 的泛点。（略去细节。）
把茎 $M=(\mathcal{F}|_V)_v$ 看作 $\mathcal{O}_{V,v}$-模。
（注意，由维数假设，$M$ 的长度有限，但实际上无须验证这一点；
见引理 \ref{spaces-chow-lemma-length-finite}。）有
$(f^*\mathcal{F}|_U)_u = \mathcal{O}_{U, u} \otimes_{\mathcal{O}_{V, v}} M$.
因此
$$
n = \text{length}_{\mathcal{O}_{U, u}}
(\mathcal{O}_{U, u} \otimes_{\mathcal{O}_{V, v}} M)
\quad
\text{且}
\quad
m = \text{length}_{\mathcal{O}_{V, v}}(M)
\text{length}_{\mathcal{O}_{V, v}}(
\mathcal{O}_{U, u}/\mathfrak m_v \mathcal{O}_{U, u})
$$
故所需等式由《代数》引理 \ref{algebra-lemma-pullback-module} 得出。
\end{proof}









\section{推出与拉回}
\label{spaces-chow-section-push-pull}

\noindent
本节对应于《Chow 同调》第 \ref{chow-section-push-pull} 节。

\medskip\noindent
本节验证：在有意义时，固有推出与平坦拉回相容。由上面的工作，
这是上同调与基变换的一个推论。

\begin{lemma}
\label{spaces-chow-lemma-flat-pullback-proper-pushforward}
在情形 \ref{spaces-chow-situation-setup} 中，设
$$
\xymatrix{
X' \ar[r]_{g'} \ar[d]_{f'} & X \ar[d]^f \\
Y' \ar[r]^g & Y
}
$$
是 $B$ 上良好代数空间的纤维积图。假设 $f:X\to Y$ 固有，
且 $g:Y'\to Y$ 平坦、相对维数为 $r$。则 $f'$ 也固有，
$g'$ 也平坦且相对维数为 $r$。对 $X$ 上任意 $k$-闭链 $\alpha$，有
$$
g^*f_*\alpha = f'_*(g')^*\alpha
$$
在 $Z_{k+r}(Y')$ 中成立。
\end{lemma}

\begin{proof}
关于 $f'$ 固有的断言由《空间的态射》引理
\ref{spaces-morphisms-lemma-base-change-proper} 得出。
关于 $g'$ 平坦且相对维数为 $r$ 的断言由《空间的态射》引理
\ref{spaces-morphisms-lemma-dimension-fibre-after-base-change}
与 \ref{spaces-morphisms-lemma-base-change-flat} 得出。
只需在 $\alpha=[W]$ 时证明闭链的等式，其中 $W\subset X$ 是
$\delta$-维数为 $k$ 的某个整闭子空间。注意在此情形中
$\alpha=[\mathcal{O}_W]_k$；见引理 \ref{spaces-chow-lemma-cycle-closed-coherent}。
因此由引理 \ref{spaces-chow-lemma-cycle-push-sheaf} 与
\ref{spaces-chow-lemma-pullback-coherent}，只需证明
$f'_*(g')^*\mathcal{O}_W$ 与 $g^*f_*\mathcal{O}_W$ 同构。
这由上同调与基变换得出；见《空间的上同调》引理
\ref{spaces-cohomology-lemma-flat-base-change-cohomology}。
\end{proof}

\begin{lemma}
\label{spaces-chow-lemma-finite-flat}
在情形 \ref{spaces-chow-situation-setup} 中，设 $X,Y/B$ 是良好的。
设 $f:X\to Y$ 是次数为 $d$ 的有限局部自由态射（见《空间的态射》定义
\ref{spaces-morphisms-definition-finite-locally-free}）。
则 $f$ 既固有又平坦，且相对维数为 $0$；并且对每个
$\alpha\in Z_k(Y)$ 有
$$
f_*f^*\alpha = d\alpha
$$
\end{lemma}

\begin{proof}
由《空间的态射》引理 \ref{spaces-morphisms-lemma-finite-flat}，
有限局部自由态射平坦且有限；又由《空间的态射》引理
\ref{spaces-morphisms-lemma-finite-proper}，有限态射固有。
略去有限态射相对维数为 $0$ 的证明。因此公式有意义。
为证明该公式，设 $Z\subset Y$ 是 $\delta$-维数为 $k$ 的整闭子概形。
只需对 $\alpha=[Z]$ 证明公式。由于有限局部自由态射的基变换仍有限局部自由
（《空间的态射》引理
\ref{spaces-morphisms-lemma-base-change-finite-locally-free}），
$f_*f^*\mathcal{O}_Z$ 是 $Z$ 上秩为 $d$ 的有限局部自由层。
故 $f_*f^*\mathcal{O}_Z$ 在 $Z$ 的泛点处显然具有长度 $d$。因此
$$
f_*f^*[Z] = f_*f^*[\mathcal{O}_Z]_k =
[f_*f^*\mathcal{O}_Z]_k = d[Z]
$$
其中使用了引理 \ref{spaces-chow-lemma-pullback-coherent} 与
\ref{spaces-chow-lemma-cycle-push-sheaf}。
\end{proof}










\section{主因子的准备}
\label{spaces-chow-section-preparation-principal-divisors}

\noindent
本节对应于《Chow 同调》第
\ref{chow-section-preparation-principal-divisors} 节。
本节的一些材料与《域上的空间》第
\ref{spaces-over-fields-section-Weil-divisors} 节中的讨论部分重叠。

\begin{lemma}
\label{spaces-chow-lemma-divisor-delta-dimension}
在情形 \ref{spaces-chow-situation-setup} 中，设 $X/B$ 是良好的，并假设 $X$ 整。
\begin{enumerate}
\item 若 $Z\subset X$ 是整闭子空间，则下列条件等价：
\begin{enumerate}
\item $Z$ 是素因子；
\item $|Z|$ 在 $|X|$ 中的余维为 $1$；
\item $\dim_\delta(Z)=\dim_\delta(X)-1$。
\end{enumerate}
\item 若 $Z$ 是 $X$ 上某个有效 Cartier 因子的不可约分支，则
$\dim_\delta(Z)=\dim_\delta(X)-1$。
\end{enumerate}
\end{lemma}

\begin{proof}
(1) 由素因子的定义（《域上的空间》定义
\ref{spaces-over-fields-definition-Weil-divisor}）、《合宜空间》引理
\ref{decent-spaces-lemma-codimension-local-ring} 以及维数函数的定义
（《拓扑学》定义 \ref{topology-definition-dimension-function}）得出。

\medskip\noindent
设 $D\subset X$ 是有效 Cartier 因子，$Z\subset D$ 是一个不可约分支，
且 $\xi\in|Z|$ 是泛点。选取平展邻域
$(U,u)\to(X,\xi)$，其中 $U=\Spec(A)$，且 $D\times_XU$ 由非零因子
$f\in A$ 截出；见《空间上的因子》引理
\ref{spaces-divisors-lemma-characterize-effective-Cartier-divisor}。
由《合宜空间》引理 \ref{decent-spaces-lemma-decent-generic-points}，
$u$ 是 $V(f)$ 的泛点。因此由 Krull 主理想定理（《代数》引理
\ref{algebra-lemma-minimal-over-1}），$\mathcal{O}_{U,u}$ 的维数为 $1$。
所以 $\xi$ 是 $X$ 上余维为 $1$ 的点，$Z$ 是素因子，正合所需。
\end{proof}












\section{主因子}
\label{spaces-chow-section-principal-divisors}

\noindent
本节对应于《Chow 同调》第 \ref{chow-section-principal-divisors} 节。
下面的定义是在当前设定下与《域上的空间》定义
\ref{spaces-over-fields-definition-principal-divisor} 对应的版本。

\begin{definition}
\label{spaces-chow-definition-principal-divisor}
在情形 \ref{spaces-chow-situation-setup} 中，设 $X/B$ 是良好的，并假设 $X$ 整且
$\dim_\delta(X)=n$。设 $f\in R(X)^*$。与 $f$ 相伴的{\it 主因子}
是 $(n-1)$-闭链
$$
\text{div}(f) = \text{div}_X(f) = \sum \text{ord}_Z(f) [Z]
$$
（定义见《域上的空间》定义
\ref{spaces-over-fields-definition-principal-divisor}）。
由引理 \ref{spaces-chow-lemma-divisor-delta-dimension}，素因子的 $\delta$-维数为
$n-1$，故此定义有意义。
\end{definition}

\noindent
在上述定义的情形中，对 $f,g\in R(X)^*$ 有
$$
\text{div}_X(fg) = \text{div}_X(f) + \text{div}_X(g)
$$
在 $Z_{n-1}(X)$ 中成立；见《域上的空间》引理
\ref{spaces-over-fields-lemma-div-additive}。
下面的引理使我们能把关于主因子的陈述约化到概形情形。

\begin{lemma}
\label{spaces-chow-lemma-etale-pullback-principal-divisor}
在情形 \ref{spaces-chow-situation-setup} 中，设 $f:X\to Y$ 是 $B$ 上良好代数空间之间的
平展态射，并假设 $Y$ 整。设 $g\in R(Y)^*$。作为 $X$ 上的闭链，有
$$
f^*(\text{div}_Y(g)) =
\sum\nolimits_{X'} (X' \to X)_*\text{div}_{X'}(g \circ f|_{X'})
$$
其中求和遍历 $X$ 的各不可约分支
（注 \ref{spaces-chow-remark-irreducible-component}）。
\end{lemma}

\begin{proof}
映射 $|X|\to|Y|$ 是开的。$|X|$ 的不可约分支族在 $|X|$ 中局部有限。
因此对每个不可约分支 $X'\subset X$，$f|_{X'}:X'\to Y$ 都是支配的。
于是 $g\circ f|_{X'}$ 有定义（《空间的态射》第
\ref{spaces-morphisms-section-rational-maps} 节），从而
$\text{div}_{X'}(g\circ f|_{X'})$ 有定义。此外，该和式局部有限，
所以右端确实是 $X$ 上的闭链。左端由定义
\ref{spaces-chow-definition-flat-pullback} 以及平展态射平坦且相对维数为 $0$
这一事实定义。

\medskip\noindent
由于 $f$ 平展，对所有 $x\in|X|$ 有
$\delta_X(x)=\delta_y(f(x))$。因此若 $\dim_\delta(Y)=n$，则 $X$ 的每个
不可约分支 $X'$ 都满足 $\dim_\delta(X')=n$（因为 $X$ 的各泛点映到
$Y$ 的泛点，见上）。所以等式两端都是 $(n-1)$-闭链。

\medskip\noindent
设 $Z\subset X$ 是满足 $\dim_\delta(Z)=n-1$ 的整闭子空间。
为证明等式，需要证明 $Z$ 的系数相同。设 $Z'\subset Y$ 是引理
\ref{spaces-chow-lemma-proper-image} 中构造的整闭子空间。此时也有
$\dim_\delta(Z')=n-1$。设 $\xi\in|Z|$ 是泛点，则
$\xi'=f(\xi)\in|Z'|$ 是泛点。考虑交换图
$$
\xymatrix{
\Spec(\mathcal{O}_{X, \xi}^h) \ar[r] \ar[d] & X \ar[d] \\
\Spec(\mathcal{O}_{Y, \xi'}^h) \ar[r] & Y
}
$$
（见《合宜空间》注
\ref{decent-spaces-remark-functoriality-henselian-local-ring}）。
需要稍加小心，因为既约诺特局部环
$\mathcal{O}_{X,\xi}^h$ 与 $\mathcal{O}_{Y,\xi'}^h$ 未必是整环。
因此我们使用全分式环 $Q(-)$ 而不是分式域。依定义，为得到 $Z'$ 在
$\text{div}_Y(g)$ 中的系数，把 $g$ 在
$Q(\mathcal{O}_{Y,\xi'}^h)$ 中的像写成 $a/b$，其中
$a,b\in\mathcal{O}_{Y,\xi'}^h$ 是非零因子，并取
$$
\text{ord}_{Z'}(g) =
\text{length}_{\mathcal{O}_{Y, \xi'}^h}
(\mathcal{O}_{Y, \xi'}^h/a \mathcal{O}_{Y, \xi'}^h) -
\text{length}_{\mathcal{O}_{Y, \xi'}^h}
(\mathcal{O}_{Y, \xi'}^h/b \mathcal{O}_{Y, \xi'}^h)
$$
注意，$Z$ 在 $f^*\text{div}_Y(G)$ 中的系数是同一个整数；见引理
\ref{spaces-chow-lemma-etale-pullback}。
% P10-E0050
假设 $\xi\in X'$。于是可以考虑映射
$$
\mathcal{O}_{Y, \xi'}^h \to
\mathcal{O}_{X, \xi}^h \to
\mathcal{O}_{X', \xi}^h
$$
第一个箭头平坦，第二个箭头是维数为 $1$ 的既约诺特局部环之间的满射。
因此这两个映射都把非零因子映到非零因子，从而 $Z'$ 在
$\text{div}_{X'}(g\circ f|_{X'})$ 中的系数为
% P10-E0051
$$
\text{ord}_Z(g \circ f|_{X'}) =
\text{length}_{\mathcal{O}_{X', \xi}^h}
(\mathcal{O}_{X', \xi}^h/a \mathcal{O}_{X', \xi}^h) -
\text{length}_{\mathcal{O}_{Y, \xi'}^h}
(\mathcal{O}_{X', \xi}^h/b \mathcal{O}_{X', \xi}^h)
$$
其定义方式与上面相同。
% P10-E0052
因此只需证明
$$
\text{length}_{\mathcal{O}_{Y, \xi'}^h}
(\mathcal{O}_{Y, \xi'}^h/a \mathcal{O}_{Y, \xi'}^h) =
\sum\nolimits_{\xi \in |X'|}
\text{length}_{\mathcal{O}_{X', \xi}^h}
(\mathcal{O}_{X', \xi}^h/a \mathcal{O}_{X', \xi}^h)
$$
% P10-E0053
首先，由于环同态 $\mathcal{O}_{Y,\xi'}^h\to
\mathcal{O}_{X,\xi}^h$ 平坦且非分歧，故
$$
\text{length}_{\mathcal{O}_{Y, \xi'}^h}
(\mathcal{O}_{Y, \xi'}^h/a \mathcal{O}_{Y, \xi'}^h) =
\text{length}_{\mathcal{O}_{X, \xi}^h}
(\mathcal{O}_{X, \xi}^h/a \mathcal{O}_{X, \xi}^h)
$$
（由《代数》引理 \ref{algebra-lemma-pullback-module}）。
设 $\mathfrak q_1,\ldots,\mathfrak q_t$ 是
$\mathcal{O}_{X,\xi}^h$ 的非极大素理想，并令
$R_j=\mathcal{O}_{X,\xi}^h/\mathfrak q_j$。
对上述 $X'$，以 $J(X')\subset\{1,\ldots,t\}$ 表示满足下列条件的指标集：
$\mathfrak q_j$ 对应于 $X'$ 的一个点；亦即在满射
$\mathcal{O}_{X,\xi}^h\to\mathcal{O}_{X',\xi}$ 下，
素理想 $\mathfrak q_j$ 对应于 $\mathcal{O}_{X',\xi}$ 的一个素理想。
% P10-E0054
由《Chow 同调》引理 \ref{chow-lemma-additivity-divisors-restricted} 得到
$$
\text{length}_{\mathcal{O}_{X, \xi}^h}
(\mathcal{O}_{X, \xi}^h/a \mathcal{O}_{X, \xi}^h) =
\sum\nolimits_j \text{length}_{R_j}(R_j/a R_j)
$$
and
$$
\text{length}_{\mathcal{O}_{X', \xi}^h}
(\mathcal{O}_{X', \xi}^h/a \mathcal{O}_{X', \xi}^h) =
\sum\nolimits_{j \in J(X')} \text{length}_{R_j}(R_j/a R_j)
$$
故引理的结论成立，因为 $\{1,\ldots,t\}$ 是各集合 $J(X')$ 的不交并：
$X$ 上每个余维为 $0$ 的点都位于唯一的 $X'$ 上。
\end{proof}







\section{主因子与推出}
\label{spaces-chow-section-two-fun}

\noindent
本节对应于《Chow 同调》第 \ref{chow-section-two-fun} 节。

\begin{lemma}
\label{spaces-chow-lemma-proper-pushforward-alteration}
在情形 \ref{spaces-chow-situation-setup} 中，设 $X,Y/B$ 是良好的。
假设 $X,Y$ 整，且 $n=\dim_\delta(X)=\dim_\delta(Y)$。
设 $p:X\to Y$ 是支配固有态射，并设 $f\in R(X)^*$。令
$$
g = \text{Nm}_{R(X)/R(Y)}(f).
$$
则 $p_*\text{div}(f)=\text{div}(g)$。
\end{lemma}

\begin{proof}
我们将通过平展局部化，从概形情形推出此结论。设 $Z\subset Y$ 是
$\delta$-维数为 $n-1$ 的整闭子空间。要证明 $[Z]$ 在
$p_*\text{div}(f)$ 与 $\text{div}(g)$ 中的系数相等。
把《域上的空间》引理 \ref{spaces-over-fields-lemma-finite-in-codim-1}
应用于态射 $p:X\to Y$ 及泛点 $\xi\in|Z|$。由此可以用包含 $\xi$
的开子空间代替 $Y$，并假设 $p:X\to Y$ 有限。选取平展邻域
$(V,v)\to(Y,\xi)$，其中 $V$ 是仿射概形。由引理
\ref{spaces-chow-lemma-etale-pullback}，只需在拉回到 $V$ 后证明闭链相等。
令 $U=V\times_YX$，并考虑交换图
$$
\xymatrix{
U \ar[r]_a \ar[d]_{p'} & X \ar[d]^p \\
V \ar[r]^b & Y
}
$$
设 $V_j\subset V$（$j=1,\ldots,m$）是 $V$ 的不可约分支。
对每个 $i$，设 $U_{j,i}$（$i=1,\ldots,n_j$）是 $U$ 中支配
$V_j$ 的各不可约分支。
% P10-E0055
以 $p'_{j,i}:U_{j,i}\to V_j$ 表示 $p':U\to V$ 的限制。
由概形情形（《Chow 同调》引理
\ref{chow-lemma-proper-pushforward-alteration}）可知
$$
p'_{j, i, *}\text{div}_{U_{j, i}}(f_{j, i}) = \text{div}_{V_j}(g_{j, i})
$$
其中 $f_{j,i}$ 是 $f$ 在 $U_{j,i}$ 上的限制，而 $g_{j,i}$ 是
$f_{j,i}$ 沿有限扩张 $R(U_{j,i})/R(V_j)$ 的范数。有
\begin{align*}
b^* p_*\text{div}_X(f)
& =
p'_* a^* \text{div}_X(f) \\
& =
p'_*\left(\sum\nolimits_{j, i}
(U_{j, i} \to U)_*\text{div}_{U_{j, i}}(f_{j, i})\right) \\
& =
\sum\nolimits_{j, i}
(V_j \to V)_*p'_{j, i, *}\text{div}_{U_{j, i}}(f_{j, i}) \\
& =
\sum\nolimits_j
(V_j \to V)_*\left(\sum\nolimits_i \text{div}_{V_j}(g_{j, i})\right) \\
& =
\sum\nolimits_j (V_j \to V)_*\text{div}_{V_j}(\prod\nolimits_i g_{j, i})
\end{align*}
这里使用了引理
\ref{spaces-chow-lemma-flat-pullback-proper-pushforward},
\ref{spaces-chow-lemma-etale-pullback-principal-divisor} 以及
\ref{spaces-chow-lemma-compose-pushforward}。
为完成证明，再次使用引理 \ref{spaces-chow-lemma-etale-pullback-principal-divisor}，
只需证明
$$
g \circ b|_{V_j} = \prod\nolimits_i g_{j, i}
$$
作为 $V_j$ 的函数域中的元素相等。用域来表述，这就是下列命题：
设 $L/K$ 是有限扩张，$M/K$ 是有限可分扩张，并写成
$M\otimes_KL=\prod M_i$。若 $t\in L$ 在各 $M_i$ 中的像为
$t_i$，则 $\text{Norm}_{L/K}(t)$ 在 $M$ 中的像为
$\prod\text{Norm}_{M_i/M}(t_i)$。略去证明。
\end{proof}







\section{有理等价}
\label{spaces-chow-section-rational-equivalence}

\noindent
本节对应于《Chow 同调》第 \ref{chow-section-rational-equivalence} 节。
本节定义 $k$-闭链上的{\it 有理等价}。我们允许主因子在闭浸入下之像的
局部有限和。这会导致一些颇为奇特的现象（见概形一章中的例子）。
然而，若不允许这些和式，我们便不知道如何证明与线丛的陈类作帽积
可通过有理等价分解。

\begin{definition}
\label{spaces-chow-definition-rational-equivalence}
在情形 \ref{spaces-chow-situation-setup} 中，设 $X/B$ 是良好的，并设
$k\in\mathbf{Z}$。
\begin{enumerate}
\item 给定任意局部有限族 $\{W_j\subset X\}$，其中各 $W_j$ 是满足
$\dim_\delta(W_j)=k+1$ 的整闭子空间，并给定任意
$f_j\in R(W_j)^*$，可以考虑
$$
\sum (i_j)_*\text{div}(f_j) \in Z_k(X)
$$
其中 $i_j:W_j\to X$ 是包含态射。由于态射
$\coprod i_j:\coprod W_j\to X$ 固有，这一定义有意义。
\item 若 $\alpha\in Z_k(X)$ 是上述形式的闭链，则称 $\alpha$
{\it 有理等价于零}。
\item 若 $\alpha-\beta$ 有理等价于零，则称
$\alpha,\beta\in Z_k(X)$ {\it 有理等价}，并记作
$\alpha\sim_{rat}\beta$。
\item 定义
$$
\CH_k(X) = Z_k(X) / \sim_{rat}
$$
为{\it $X$ 上 $k$-闭链的周群}。它有时也称为
{\it $X$ 上模有理等价的 $k$-闭链周群}。
\end{enumerate}
\end{definition}

\noindent
还有许多其他有趣的等价关系。有理等价是其中最粗的。
下面是一个十分简单但很重要的引理。

\begin{lemma}
\label{spaces-chow-lemma-restrict-to-open}
在情形 \ref{spaces-chow-situation-setup} 中，设 $X/B$ 是良好的。
设 $U\subset X$ 是开子空间，$Y$ 是满足
$|Y|=|X|\setminus|U|$ 的 $X$ 的既约闭子空间，并以
$i:Y\to X$ 表示包含态射。设 $k\in\mathbf{Z}$，并假设
$\alpha,\beta\in Z_k(X)$。若 $\alpha|_U\sim_{rat}\beta|_U$，
则存在闭链 $\gamma\in Z_k(Y)$，使得
$$
\alpha \sim_{rat} \beta + i_*\gamma.
$$
换言之，序列
$$
\xymatrix{
\CH_k(Y) \ar[r]^{i_*} & \CH_k(X) \ar[r]^{j^*} & \CH_k(U) \ar[r] & 0
}
$$
是阿贝尔群的正合复形。
\end{lemma}

\begin{proof}
设 $\{W_j\}_{j\in J}$ 是 $U$ 中 $\delta$-维数为 $k+1$ 的整闭子空间
的局部有限族，并设 $f_j\in R(W_j)^*$ 满足定义中的等式
$(\alpha-\beta)|_U=\sum(i_j)_*\text{div}(f_j)$。
设 $W_j'\subset X$ 是相应的整闭子空间，即它与 $W_j$ 有同一个泛点。
假设 $V\subset X$ 是拟紧开集。由于 $V$ 诺特，$V\cap U$ 也是 $U$
中的拟紧开集。于是，由 $\{W_j\}$ 的局部有限性，集合
$\{j\in J\mid W_j\cap V\not=\emptyset\}
=\{j\in J\mid W'_j\cap V\not=\emptyset\}$ 有限。
换言之，$\{W'_j\}$ 也局部有限。由于 $R(W_j)=R(W'_j)$，可知
$$
\alpha - \beta - \sum (i'_j)_*\text{div}(f_j)
$$
是 $X$ 上的闭链，且其在 $U$ 上的限制为零。
应用引理 \ref{spaces-chow-lemma-exact-sequence-open} 即得结论。
\end{proof}

\begin{remark}
\label{spaces-chow-remark-infinite-sums-rational-equivalences}
在情形 \ref{spaces-chow-situation-setup} 中，设 $X/B$ 是良好的。
假设有 $X$ 上 $k$-闭链的无限族
$\alpha_i,\beta_i\in Z_k(X)$（$i\in I$）。假设 $\alpha_i$ 与
$\beta_i$ 的支撑分别构成 $X$ 的闭子集的局部有限族，从而
$\sum\alpha_i$ 与 $\sum\beta_i$ 都定义为闭链。再假设对每个 $i$
都有 $\alpha_i\sim_{rat}\beta_i$。这时并不显然有
$\sum\alpha_i\sim_{rat}\sum\beta_i$。问题在于，这些有理等价可能由
局部有限族
$\{W_{i,j},f_{i,j}\in R(W_{i,j})^*\}_{j\in J_i}$ 给出，
但其并族 $\{W_{i,j}\}_{i\in I,j\in J_i}$ 未必局部有限。

\medskip\noindent
实践中经常有 $|X|$ 的闭子集的局部有限族
$\{T_i\}_{i\in I}$，使得 $\alpha_i,\beta_i$ 支撑在 $T_i$ 上，
并在 $T_i$ “上”满足 $\alpha_i\sim_{rat}\beta_i$。
更精确地说，各族
$\{W_{i,j},f_{i,j}\in R(W_{i,j})^*\}_{j\in J_i}$
由满足 $|W_{i,j}|\subset T_i$ 的整闭子空间 $W_{i,j}$ 构成。
在这种情形中，$X$ 上确实有
$\sum\alpha_i\sim_{rat}\sum\beta_i$，因为族
$\{W_{i,j}\}_{i\in I,j\in J_i}$ 此时自动局部有限。
\end{remark}











\section{有理等价与推出、拉回}
\label{spaces-chow-section-properties-rational-equivalence}

\noindent
本节对应于《Chow 同调》第
\ref{chow-section-properties-rational-equivalence} 节。
本节证明平坦拉回与固有推出同有理等价相容。

\begin{lemma}
\label{spaces-chow-lemma-prepare-flat-pullback-rational-equivalence}
在情形 \ref{spaces-chow-situation-setup} 中，设 $X,Y/B$ 是良好的。
假设 $Y$ 整且 $\dim_\delta(Y)=k$。设 $f:X\to Y$ 是相对维数为
$r$ 的平坦态射。则对 $g\in R(Y)^*$ 有
$$
f^*\text{div}_Y(g) =
\sum m_{X', X} (X' \to X)_*\text{div}_{X'}(g \circ f|_{X'})
$$
作为 $X$ 上的 $(k+r-1)$-闭链相等，其中求和遍历 $X$ 的各不可约分支
$X'$，而 $m_{X',X}$ 是 $X'$ 在 $X$ 中的重数。
\end{lemma}

\begin{proof}
注意，$X$ 的任意不可约分支都支配 $Y$
（引理 \ref{spaces-chow-lemma-flat-inverse-image-dimension}），因而复合
$g\circ f|_{X'}$ 有定义（《空间的态射》第
\ref{spaces-morphisms-section-rational-maps} 节）。
我们把问题约化到概形情形。选取概形 $V$ 和满平展态射 $V\to Y$，
再选取概形 $U$ 和满平展态射 $U\to V\times_YX$。图示为
$$
\xymatrix{
U \ar[r]_a \ar[d]_h & X \ar[d]^f \\
V \ar[r]^b & Y
}
$$
由于 $a$ 满且平展，由引理 \ref{spaces-chow-lemma-etale-pullback}，只需在沿 $a$
拉回后证明闭链相等。由引理
\ref{spaces-chow-lemma-etale-pullback-principal-divisor}，可以写成
$$
b^*\text{div}_Y(g) = \sum (V' \to V)_*\text{div}_{V'}(g \circ b|_{V'})
$$
其中求和遍历 $V$ 的各不可约分支 $V'$。使用引理
\ref{spaces-chow-lemma-flat-pullback-proper-pushforward} 得
$$
h^*b^*\text{div}_Y(g) =
\sum (V' \times_V U \to U)_*(h')^*\text{div}_{V'}(g \circ b|_{V'})
$$
其中 $h':V'\times_VU\to V'$ 是投影。把概形情形中的引理
（《Chow 同调》引理
\ref{chow-lemma-prepare-flat-pullback-rational-equivalence}）
应用于态射 $h':V'\times_VU\to V'$，得到
$$
(h')^*\text{div}_{V'}(g \circ b|_{V'}) =
\sum
m_{U', V' \times_V U}
(U' \to V' \times_V U)_*\text{div}_{U'}(g \circ b|_{V'} \circ h'|_{U'})
$$
其中求和遍历 $V'\times_VU$ 的各不可约分支 $U'$。
此和式中出现的每个 $U'$ 都是 $U$ 的不可约分支；反过来，$U$ 的每个
不可约分支 $U'$ 都是某个唯一不可约分支 $V'\subset V$ 所对应的
$V'\times_VU$ 的不可约分支。
给定不可约分支 $U'\subset U$，以 $\overline{a(U')}\subset X$
表示其在 $X$ 中的“像”（引理 \ref{spaces-chow-lemma-proper-image}）；
例如由引理 \ref{spaces-chow-lemma-flat-inverse-image-dimension}，它是 $X$
的不可约分支。重数 $m_{U',V'\times_VU}$ 等于重数
$m_{\overline{a(U')},X}$。
% P10-E0056
这由等式 $h^*a^*[Y]=b^*f^*[Y]$（引理
\ref{spaces-chow-lemma-compose-flat-pullback}）、定义以及引理
\ref{spaces-chow-lemma-etale-pullback} 得出。综合以上结论，得到
$$
a^*f^*\text{div}_Y(g) =
h^*b^*\text{div}_Y(g) =
\sum m_{\overline{a(U')}, X}
(U' \to U)_*\text{div}_{U'}(g \circ (f \circ a)|_{U'})
$$
接下来分析引理陈述中公式右端沿 $a$ 拉回后的情形。
首先使用引理 \ref{spaces-chow-lemma-flat-pullback-proper-pushforward} 得到
$$
a^*\sum m_{X', X} (X' \to X)_*\text{div}_{X'}(g \circ f|_{X'}) =
\sum m_{X', X} (X' \times_X U \to U)_*(a')^*\text{div}_{X'}(g \circ f|_{X'})
$$
其中 $a':X'\times_XU\to X'$ 是投影。由引理
\ref{spaces-chow-lemma-etale-pullback-principal-divisor} 得
$$
(a')^*\text{div}_{X'}(g \circ f|_{X'}) =
\sum (U' \to X' \times_X U)_*\text{div}_{U'}(g \circ (f \circ a)|_{U'})
$$
其中求和遍历 $X'\times_XU$ 的各不可约分支 $U'$。
这些 $U'$ 是 $U$ 的不可约分支，并且恰好是满足
$\overline{a(U')}=X'$ 的那些 $U$ 的不可约分支。
与上面所得公式比较即得结论。
\end{proof}

\begin{lemma}
\label{spaces-chow-lemma-flat-pullback-rational-equivalence}
在情形 \ref{spaces-chow-situation-setup} 中，设 $X,Y/B$ 是良好的。
设 $f:X\to Y$ 是相对维数为 $r$ 的平坦态射，并设
$\alpha\sim_{rat}\beta$ 是 $Y$ 上有理等价的 $k$-闭链。
则作为 $X$ 上的 $(k+r)$-闭链，有
$f^*\alpha\sim_{rat}f^*\beta$。
\end{lemma}

\begin{proof}
需要证明什么？假设给定一族闭浸入
$$
i_j : W_j \longrightarrow Y
$$
其中每个 $W_j$ 整且 $\delta$-维数为 $k+1$，并给定有理函数
$g_j\in R(W_j)^*$。再假设族
$\{|i_j|(|W_j|)\}_{j\in J}$ 在 $|Y|$ 中局部有限。
那么需要证明
$$
f^*(\sum i_{j, *}\text{div}(g_j)) = \sum f^*i_{j, *}\text{div}(g_j)
$$
在 $X$ 上有理等价于零。由引理
\ref{spaces-chow-lemma-inverse-image-locally-finite}，右端的和有意义。

\medskip\noindent
考虑纤维积
$$
i'_j : W'_j = W_j \times_Y X \longrightarrow X.
$$
并以 $f_j:W'_j\to W_j$ 表示第一投影。
由引理 \ref{spaces-chow-lemma-flat-pullback-proper-pushforward}，上述和式可写成
$$
\sum i'_{j, *}(f_j^*\text{div}(g_j))
$$
由引理 \ref{spaces-chow-lemma-prepare-flat-pullback-rational-equivalence}，
每个 $f_j^*\text{div}(g_j)$ 都在 $W'_j$ 上有理等价于零。
因此每个 $i'_{j,*}(f_j^*\text{div}(g_j))$ 都有理等价于零。
再由注 \ref{spaces-chow-remark-infinite-sums-rational-equivalences} 中的讨论，
所显示的和式也有理等价于零。
\end{proof}

\begin{lemma}
\label{spaces-chow-lemma-proper-pushforward-rational-equivalence}
在情形 \ref{spaces-chow-situation-setup} 中，设 $X,Y/B$ 是良好的。
设 $p:X\to Y$ 是固有态射。假设
$\alpha,\beta\in Z_k(X)$ 有理等价。则 $p_*\alpha$ 与 $p_*\beta$
有理等价。
\end{lemma}

\begin{proof}
需要证明什么？假设给定一族闭浸入
$$
i_j : W_j \longrightarrow X
$$
其中每个 $W_j$ 整且 $\delta$-维数为 $k+1$，并给定有理函数
$f_j\in R(W_j)^*$。再假设族 $\{i_j(W_j)\}_{j\in J}$ 在 $X$
上局部有限。那么需要证明
$$
p_*\left(\sum i_{j, *}\text{div}(f_j)\right)
$$
在 $X$ 上有理等价于零。
% P10-E0057

\medskip\noindent
注意该和式等于
$$
\sum p_*i_{j, *}\text{div}(f_j).
$$
设 $W'_j\subset Y$ 是 $p\circ i_j$ 的像这一整闭子空间；见引理
\ref{spaces-chow-lemma-proper-image}。由引理
\ref{spaces-chow-lemma-quasi-compact-locally-finite}，族 $\{W'_j\}$ 在 $Y$
中局部有限。因此，只需对给定的 $j$ 证明：
$p_*i_{j,*}\text{div}(f_j)=0$，或者它等于某个
$g_j\in R(W'_j)^*$ 所对应的 $i'_{j,*}\text{div}(g_j)$。

\medskip\noindent
因此上述论证把我们约化到单个 $\delta$-维数为 $k+1$ 的整闭子空间
$W\subset X$ 的情形。
% P10-E0058
设 $f\in R(W)^*$，并令 $W'=p(W)$ 如上。得到态射的交换图
$$
\xymatrix{
W \ar[r]_i \ar[d]_{p'} & X \ar[d]^p \\
W' \ar[r]^{i'} & Y
}
$$
由引理 \ref{spaces-chow-lemma-compose-pushforward}，注意
$p_*i_*\text{div}(f)=i'_*(p')_*\text{div}(f)$。
如上所述，必须证明 $(p')_*\text{div}(f)$ 是 $W'$ 上某个有理函数的
因子，或者为零。分三种情形。

\medskip\noindent
情形 $\dim_\delta(W')<k$。此时自动有
$(p')_*\text{div}(f)=0$，无需证明。

\medskip\noindent
情形 $\dim_\delta(W')=k$。我们证明此时
$(p')_*\text{div}(f)=0$。由于 $(p')_*\text{div}(f)$ 是 $k$-闭链，
存在 $n\in\mathbf{Z}$ 使得 $(p')_*\text{div}(f)=n[W']$。
为证明 $n=0$，可以用非空开子空间代替 $W'$。特别地，可以并且确实
假设 $W'$ 是概形。设 $\eta\in W'$ 是泛点，并令函数域
$K=\kappa(\eta)=R(W')$。考虑基变换图
$$
\xymatrix{
W_\eta \ar[r] \ar[d]_c & W \ar[d]^{p'} \\
\Spec(K) \ar[r]^\eta & W'
}
$$
注意 $c$ 固有，而且 $|W_\eta|$ 的维数为 $1$：使用《合宜空间》引理
\ref{decent-spaces-lemma-topology-fibre}，把 $|W_\eta|$ 识别为
$|W|$ 中映到 $\eta$ 的点所成的子空间；又由于
$\dim_\delta(W)=k+1$ 且 $\delta(\eta)=k$，$W_\eta$ 中的点的
$\delta$-值只能是 $k$ 或 $k+1$。于是由《合宜空间》引理
\ref{decent-spaces-lemma-codimension-local-ring}，各局部环的维数
$\leq1$。由《域上的空间》引理
\ref{spaces-over-fields-lemma-codim-1-point-in-schematic-locus}，
$W_\eta$ 是概形。由于 $\Spec(K)$ 是 $W'$ 的各非空仿射开子概形的极限，
由《空间的极限》引理 \ref{spaces-limits-lemma-limit-is-scheme}，
可以假设 $W$ 是概形。最后，由概形情形（《Chow 同调》引理
\ref{chow-lemma-proper-pushforward-rational-equivalence}）得 $n=0$。

\medskip\noindent
情形 $\dim_\delta(W')=k+1$。此时引理
\ref{spaces-chow-lemma-proper-pushforward-alteration} 适用，因而确有某个
$g\in R(W')^*$ 使得
$p'_*\text{div}(f)=\text{div}(g)$，正合所需。
\end{proof}












\section{与可逆层相伴的因子}
\label{spaces-chow-section-divisor-invertible-sheaf}

\noindent
本节对应于《Chow 同调》第
\ref{chow-section-divisor-invertible-sheaf} 节。
下面的定义是在当前设定下与《域上的空间》定义
\ref{spaces-over-fields-definition-divisor-invertible-sheaf} 对应的版本。

\begin{definition}
\label{spaces-chow-definition-divisor-invertible-sheaf}
在情形 \ref{spaces-chow-situation-setup} 中，设 $X/B$ 是良好的。
假设 $X$ 整且 $n=\dim_\delta(X)$。设 $\mathcal{L}$ 是可逆
$\mathcal{O}_X$-模。
\begin{enumerate}
\item 对 $\mathcal{L}$ 的任意非零亚纯截面 $s$，定义
{\it 与 $s$ 相伴的 Weil 因子}为 $(n-1)$-闭链
% P10-E0059
$$
\text{div}_\mathcal{L}(s) =
\sum \text{ord}_{Z, \mathcal{L}}(s) [Z]
$$
（定义见《域上的空间》定义
\ref{spaces-over-fields-definition-divisor-invertible-sheaf}）。
由引理 \ref{spaces-chow-lemma-divisor-delta-dimension}，Weil 因子的
$\delta$-维数为 $n-1$，故此定义有意义。
\item 定义{\it 与 $\mathcal{L}$ 相伴的 Weil 因子}为
$$
c_1(\mathcal{L}) \cap [X] =
\text{div}_\mathcal{L}(s)\text{ 的类} \in \CH_{n - 1}(X)
$$
其中 $s$ 是 $\mathcal{L}$ 在 $X$ 上的任意非零亚纯截面。
由《域上的空间》引理
\ref{spaces-over-fields-lemma-divisor-meromorphic-well-defined}，
这一定义良好。
\end{enumerate}
\end{definition}

\noindent
非零截面的零概形是有效 Cartier 因子，其 Weil 因子类给出与该可逆模
相伴的 Weil 因子。

\begin{lemma}
\label{spaces-chow-lemma-compute-c1}
在情形 \ref{spaces-chow-situation-setup} 中，设 $X/B$ 是良好的。
假设 $X$ 整且 $n=\dim_\delta(X)$。设 $\mathcal{L}$ 是可逆
$\mathcal{O}_X$-模，并设 $s\in\Gamma(X,\mathcal{L})$ 是非零整体截面。
则
$$
\text{div}_\mathcal{L}(s) = [Z(s)]_{n - 1}
$$
在 $Z_{n-1}(X)$ 中成立，并且
$$
c_1(\mathcal{L}) \cap [X] = [Z(s)]_{n - 1}
$$
在 $\CH_{n-1}(X)$ 中成立。
\end{lemma}

\begin{proof}
设 $Z\subset X$ 是 $\delta$-维数为 $n-1$ 的整闭子空间，
并设 $\xi\in|Z|$ 是其泛点。为证明第一个等式，比较两端 $Z$ 的系数。
选取初等平展邻域 $(U,u)\to(X,\xi)$；见《合宜空间》第
\ref{decent-spaces-section-residue-fields-henselian-local-rings} 节。
回顾在此情形中
$\mathcal{O}_{X,\xi}^h=\mathcal{O}_{U,u}^h$。
用 $u$ 的一个开邻域代替 $U$ 后，可以假设 $\mathcal{L}|_U$
有平凡化截面 $s_U$。对某个 $f\in\Gamma(U,\mathcal{O}_U)$ 写成
$s|_U=fs_U$。则作为 $U$ 的闭子概形，$Z\times_XU$ 等于 $V(f)$；
见《空间上的因子》定义
\ref{spaces-divisors-definition-zero-scheme-s}。
如《域上的空间》第 \ref{spaces-over-fields-section-c1} 节，以
$\mathcal{L}_\xi$ 表示 $\mathcal{L}$ 沿典范态射
$c_\xi:\Spec(\mathcal{O}_{X,\xi}^h)\to X$ 的拉回。
以 $s_\xi$ 表示 $s_U$ 的拉回；它是 $\mathcal{L}_\xi$ 的一个平凡化。
于是 $c_\xi^*(s)=fs_\xi$。依定义，$Z$ 在
$[Z(s)]_{n-1}$ 中的系数为
$$
\text{length}_{\mathcal{O}_{U, u}}(\mathcal{O}_{U, u}/f\mathcal{O}_{U, u})
$$
由于 $\mathcal{O}_{U,u}\to\mathcal{O}_{X,\xi}^h$ 平坦且诱导剩余域
同构，这等于
$$
\text{length}_{\mathcal{O}_{X, \xi}^h}
(\mathcal{O}_{X, \xi}^h/f\mathcal{O}_{X, \xi}^h)
$$
（由《代数》引理 \ref{algebra-lemma-pullback-module}）。
由《域上的空间》定义
\ref{spaces-over-fields-definition-order-vanishing-meromorphic}，
最后这个量等于 $\text{ord}_{Z,\mathcal{L}}(s)$，亦即 $Z$ 在
$\text{div}_\mathcal{L}(s)$ 中的系数，正合所需。
\end{proof}

\begin{lemma}
\label{spaces-chow-lemma-Gm-torsor}
在情形 \ref{spaces-chow-situation-setup} 中，设 $X/B$ 是良好的，且
$\mathcal{L}$ 是可逆 $\mathcal{O}_X$-模。态射
$$
q :
T = \underline{\Spec}\left(
\bigoplus\nolimits_{n \in \mathbf{Z}} \mathcal{L}^{\otimes n}\right)
\longrightarrow
X
$$
具有下列性质：
\begin{enumerate}
\item $q$ 满、光滑、仿射，且相对维数为 $1$；
\item 存在同构 $\alpha:q^*\mathcal{L}\cong\mathcal{O}_T$；
\item $(q:T\to X,\alpha)$ 的构造与基变换相容；
\item $q^*:Z_k(X)\to Z_{k+1}(T)$ 是单射；
\item 若 $Z\subset X$ 是整闭子空间，则
$q^{-1}(Z)\subset T$ 是整闭子空间；
\item 若 $Z\subset X$ 是 $\delta$-维数 $\leq k$ 的闭子空间，
则 $q^{-1}(Z)$ 是 $T$ 中 $\delta$-维数 $\leq k+1$ 的闭子空间，且
$q^*[Z]_k=[q^{-1}(Z)]_{k+1}$；
\item 若 $\xi'\in|T|$ 是 $|T|\to|X|$ 在 $\xi$ 上的纤维的泛点，
则环同态 $\mathcal{O}_{X,\xi}^h\to\mathcal{O}_{T,\xi'}^h$ 平坦，
有 $\mathfrak m_{\xi'}^h=\mathfrak m_\xi^h\mathcal{O}_{T,\xi'}^h$，
且剩余域扩张是超越次数为 $1$ 的纯超越扩张；以及
\item 按需在此增补。
% P10-E0061
\end{enumerate}
\end{lemma}

\begin{proof}
设 $U\to X$ 是使 $\mathcal{L}|_U$ 平凡的平展态射。
则 $T\times_XU\to U$ 同构于投影态射
$\mathbf{G}_m\times U\to U$，其中 $\mathbf{G}_m$ 是乘法群概形；
见《群胚》例 \ref{groupoids-example-multiplicative-group}。
% P10-E0060
因此 (1) 显然。

\medskip\noindent
为证明 (2)，注意
$q_*q^*\mathcal{L} =
\bigoplus_{n \in \mathbf{Z}} \mathcal{L}^{\otimes n + 1}$.
因此存在显然的 $q_*\mathcal{O}_T$-模同构
$q_*q^*\mathcal{L}\to q_*\mathcal{O}_T$。由《空间的态射》引理
\ref{spaces-morphisms-lemma-affine-equivalence-modules}，
它确定一个同构 $q^*\mathcal{L}\to\mathcal{O}_T$。

\medskip\noindent
(3) 成立，因为相对谱的构造与任意基变换相容，而同构 $\alpha$
显然也有同样的性质。

\medskip\noindent
(4) 立即由 (1) 与定义得出。

\medskip\noindent
(5) 由下述事实得出：若 $Z$ 是整代数空间，则
$\mathbf{G}_m\times Z$ 是整代数空间。

\medskip\noindent
(6) 由长度保持这一事实得出：若 $(A,\mathfrak m)$ 是局部环，
$B=A[x]_{\mathfrak m A[x]}$，且 $M$ 是 $A$-模，则
$\text{length}_A(M) = \text{length}_B(M \otimes_A B)$.
这蕴含：若 $\mathcal{F}$ 是凝聚 $\mathcal{O}_X$-模，
$\xi\in|X|$，且 $\xi'\in|T|$ 是 $\xi$ 上纤维的泛点，
则 $\mathcal{F}$ 在 $\xi$ 处的长度等于 $q^*\mathcal{F}$
在 $\xi'$ 处的长度。追溯各定义即可得到 (6)，以及更多结论。

\medskip\noindent
(7) 中的映射来自《合宜空间》注
\ref{decent-spaces-remark-functoriality-henselian-local-ring}。
不过在本情形中有
$$
\Spec(\mathcal{O}_{X, \xi}^h) \times_X T =
\mathbf{G}_m \times \Spec(\mathcal{O}_{X, \xi}^h) =
\Spec(\mathcal{O}_{X, \xi}^h[t, t^{-1}])
$$
且 $\xi'$ 对应于它在 $\Spec(\mathcal{O}_{X,\xi}^h)$ 上特殊纤维的泛点。
因此 $\mathcal{O}_{T,\xi'}^h$ 是
$\mathcal{O}_{X,\xi}^h[t,t^{-1}]$ 在相应素理想处的局部化之 Hensel 化。
(7) 由此及一些交换代数结论得出；略去细节。
\end{proof}

\begin{lemma}
\label{spaces-chow-lemma-Gm-torsor-divisor-meromorphic-section}
在情形 \ref{spaces-chow-situation-setup} 中，设 $X/B$ 是良好的，且
$\mathcal{L}$ 是可逆 $\mathcal{O}_X$-模。假设 $X$ 整，并设 $s$
是 $\mathcal{L}$ 的非零亚纯截面。设 $q:T\to X$ 是引理
\ref{spaces-chow-lemma-Gm-torsor} 中的态射。则
$$
q^*\text{div}_\mathcal{L}(s) = \text{div}_T(q^*(s))
$$
其中利用同构 $q^*\mathcal{L}\to\mathcal{O}_T$，把拉回 $q^*(s)$
视为 $T$ 上的非零亚纯函数。
\end{lemma}

\begin{proof}
由《域上的空间》第
\ref{spaces-over-fields-section-Weil-divisors} 节与
\ref{spaces-over-fields-section-c1} 节中各构造的相容性，注意
$\text{div}_T(q^*(s))=\text{div}_{\mathcal{O}_T}(q^*(s))$。
本证明将说明它与 $\text{div}_{\mathcal{O}_T}(q^*(s))$ 一致。
下文将直接使用引理 \ref{spaces-chow-lemma-Gm-torsor} 中所述的
$q:T\to X$ 的全部性质，不再另行说明。设 $Z\subset T$ 是素因子。
那么或者 $Z\to X$ 支配，或者对某个素因子 $Z'\subset X$ 有
$Z=q^{-1}(Z')$。若 $Z\to X$ 支配，则 $Z$ 在引理等式两端的系数
均为零。因此可以假设 $Z=q^{-1}(Z')$，其中 $Z'\subset X$ 是素因子。
设 $\xi'\in|Z'|$ 与 $\xi\in|Z|$ 分别是泛点。于是得到交换图
$$
\xymatrix{
\Spec(\mathcal{O}_{T, \xi}^h) \ar[r]_-{c_\xi} \ar[d]_h & T \ar[d]^q \\
\Spec(\mathcal{O}_{X, \xi'}^h) \ar[r]^-{c_{\xi'}} & X
}
$$
；见《合宜空间》注
\ref{decent-spaces-remark-functoriality-henselian-local-ring}。
选取 $\mathcal{L}_{\xi'}=c_{\xi'}^*\mathcal{L}$ 的平凡化
$s_{\xi'}$。于是可以用 $s_{\xi'}$ 沿 $h$ 的拉回 $s_\xi$
作为 $\mathcal{L}_\xi=c_\xi^*q^*\mathcal{L}$ 的平凡化。
对非零因子 $a,b\in\mathcal{O}_{X,\xi'}$，写成
$s/s_{\xi'}=a/b$。依定义，$Z'$ 在
$\text{div}_\mathcal{L}(s)$ 中的系数为
$$
\text{length}_{\mathcal{O}_{X, \xi'}^h}(
\mathcal{O}_{X, \xi'}^h/a \mathcal{O}_{X, \xi'}^h)
-
\text{length}_{\mathcal{O}_{X, \xi'}^h}(
\mathcal{O}_{X, \xi'}^h/b \mathcal{O}_{X, \xi'}^h)
$$
由于 $Z=q^{-1}(Z')$，这也是 $Z$ 在
$q^*\text{div}_\mathcal{L}(s)$ 中的系数。由于
$\mathcal{O}_{X,\xi'}^h\to\mathcal{O}_{T,\xi}^h$ 平坦，元素 $a,b$
在 $\mathcal{O}_{T,\xi}^h$ 中的像是非零因子。因此在
$\mathcal{O}_{T,\xi}^h$ 的全商环中有
$q^*(s)/s_\xi=a/b$。依定义，$Z$ 在
$\text{div}_T(q^*(s))$ 中的系数为
$$
\text{length}_{\mathcal{O}_{T, \xi}^h}(
\mathcal{O}_{T, \xi}^h/a \mathcal{O}_{T, \xi}^h)
-
\text{length}_{\mathcal{O}_{T, \xi}^h}(
\mathcal{O}_{T, \xi}^h/b \mathcal{O}_{T, \xi}^h)
$$
由《代数》引理 \ref{algebra-lemma-pullback-module} 以及引理
\ref{spaces-chow-lemma-Gm-torsor} 中证明的
$\mathfrak m_\xi^h=\mathfrak m_{\xi'}^h\mathcal{O}_{T,\xi}^h$，
这些长度与前面相同，证明完成。
\end{proof}









\section{与可逆层相交}
\label{spaces-chow-section-intersecting-with-divisors}

\noindent
本节对应于《Chow 同调》第
\ref{chow-section-intersecting-with-divisors} 节。
本节研究下列构造。

\begin{definition}
\label{spaces-chow-definition-cap-c1}
在情形 \ref{spaces-chow-situation-setup} 中，设 $X/B$ 是良好的，且
$\mathcal{L}$ 是可逆 $\mathcal{O}_X$-模。对每个整数 $k$，定义运算
$$
c_1(\mathcal{L}) \cap - :
Z_{k + 1}(X) \to \CH_k(X)
$$
称为{\it 与 $\mathcal{L}$ 的第一陈类相交}。
\begin{enumerate}
\item 给定满足 $\dim_\delta(W)=k+1$ 的整闭子空间
$i:W\to X$，定义
$$
c_1(\mathcal{L}) \cap [W] = i_*(c_1({i^*\mathcal{L}}) \cap [W])
$$
其中右端定义见定义 \ref{spaces-chow-definition-divisor-invertible-sheaf}。
\item 对一般的 $(k+1)$-闭链 $\alpha=\sum n_i[W_i]$，令
$$
c_1(\mathcal{L}) \cap \alpha = \sum n_i c_1(\mathcal{L}) \cap [W_i]
$$
\end{enumerate}
\end{definition}

\noindent
把每个 $c_1(\mathcal{L})\cap W_i=\sum_jn_{i,j}[Z_{i,j}]$ 写出，
其中 $\{Z_{i,j}\}_j$ 是 $W_i$
的整闭子空间的局部有限和。
% P10-E0062
由于 $\{W_i\}$ 是 $X$ 上整闭子空间的局部有限族，容易看出
$\{Z_{i,j}\}_{i,j}$ 是 $X$ 的闭子空间的局部有限族。因此
$c_1(\mathcal{L})\cap\alpha=\sum n_in_{i,j}[Z_{i,j}]$ 是闭链。
% P10-E0063
另一种通常更方便的理解方式是注意态射 $\coprod W_i\to X$ 固有。
所以 $c_1(\mathcal{L})\cap\alpha$ 可以看作
$\CH_k(\coprod W_i)=\prod\CH_k(W_i)$ 中某个类的推出。
这也解释了为什么所得结果在 $X$ 上模有理等价是良定义的。

\medskip\noindent
接下来几节的主要目标是证明，与 $c_1(\mathcal{L})$ 相交可通过
有理等价分解。这并不是显然的。

\begin{lemma}
\label{spaces-chow-lemma-c1-cap-additive}
在情形 \ref{spaces-chow-situation-setup} 中，设 $X/B$ 是良好的。
设 $\mathcal{L},\mathcal{N}$ 是 $X$ 上的可逆层。则
% P10-E0065
$$
c_1(\mathcal{L}) \cap \alpha  + c_1(\mathcal{N}) \cap \alpha =
c_1(\mathcal{L} \otimes_{\mathcal{O}_X} \mathcal{N}) \cap \alpha
$$
对每个 $\alpha\in Z_{k-1}(X)$ 在 $\CH_k(X)$ 中成立。此外，
对所有 $\alpha$ 都有 $c_1(\mathcal{O}_X)\cap\alpha=0$。
% P10-E0064
\end{lemma}

\begin{proof}
可加性直接由《域上的空间》引理
\ref{spaces-over-fields-lemma-c1-additive} 与各定义得出。
为证明 $c_1(\mathcal{O}_X)\cap\alpha=0$，考虑截面
$1\in\Gamma(X,\mathcal{O}_X)$。它在任意整闭子空间
$W\subset X$ 上的限制处处非零。因此
$c_1(\mathcal{O}_X)\cap[W]=0$，正合所需。
\end{proof}

\noindent
回顾：$Z(s)\subset X$ 表示代数空间 $X$ 上可逆层的整体截面 $s$
的零概形；见《空间上的因子》定义
\ref{spaces-divisors-definition-zero-scheme-s}。

\begin{lemma}
\label{spaces-chow-lemma-prepare-geometric-cap}
在情形 \ref{spaces-chow-situation-setup} 中，设 $Y/B$ 是良好的，且
$\mathcal{L}$ 是可逆 $\mathcal{O}_Y$-模。设
$s\in\Gamma(Y,\mathcal{L})$ 是正则截面，并假设
$\dim_\delta(Y)\leq k+1$。写成
$[Y]_{k+1}=\sum n_i[Y_i]$，其中 $Y_i\subset Y$ 是 $Y$ 中
$\delta$-维数为 $k+1$ 的各不可约分支。令
$s_i=s|_{Y_i}\in\Gamma(Y_i,\mathcal{L}|_{Y_i})$。则
\begin{equation}
\label{spaces-chow-equation-equal-as-cycles}
[Z(s)]_k =  \sum n_i[Z(s_i)]_k
\end{equation}
作为 $Y$ 上的 $k$-闭链相等。
\end{lemma}

\begin{proof}
设 $\varphi:V\to Y$ 是满平展态射，其中 $V$ 是概形。
只需在沿 $\varphi$ 拉回后证明等式；见引理
\ref{spaces-chow-lemma-etale-pullback}。同一引理还给出
$\varphi^*[Y_i]=[\varphi^{-1}(Y_i)]=\sum[V_{i,j}]$，
其中 $V_{i,j}$ 是 $V$ 中位于 $Y_i$ 上方的各不可约分支。
因此，先把概形情形（《Chow 同调》引理
\ref{chow-lemma-prepare-geometric-cap}）应用于
$Y_i\times_YV$ 上的 $\varphi^*s_i$，得到
$\varphi^*[Z(s_i)]_k=[Z(\varphi^*s_i)]=\sum[Z(s_{i,j})]_k$，
其中 $s_{i,j}$ 是 $s$ 在 $V_{i,j}$ 上的拉回。
再把概形情形应用于 $\varphi^*s$，得到
$$
\varphi^*[Z(s)]_k =
[Z(\varphi^*s)]_k =
\sum n_i[Z(s_{i, j})]_k
$$
（使用了上面关于重数的说明）。综合以上各式即完成证明。
\end{proof}

\noindent
下面的引理是计算可逆层的 $c_1$ 与闭子概形所伴闭链之交积的一个
有用结果。回顾：$Z(s)\subset X$ 表示概形 $X$ 上可逆层的整体截面
$s$ 的零概形；见《因子》定义
\ref{divisors-definition-zero-scheme-s}。

\begin{lemma}
\label{spaces-chow-lemma-geometric-cap}
在情形 \ref{spaces-chow-situation-setup} 中，设 $X/B$ 是良好的，且
$\mathcal{L}$ 是可逆 $\mathcal{O}_X$-模。设 $Y\subset X$
是满足 $\dim_\delta(Y)\leq k+1$ 的闭子概形，并设
$s\in\Gamma(Y,\mathcal{L}|_Y)$ 是正则截面。则
$$
c_1(\mathcal{L}) \cap [Y]_{k + 1} = [Z(s)]_k
$$
在 $\CH_k(X)$ 中成立。
\end{lemma}

\begin{proof}
写成
$$
[Y]_{k + 1} = \sum n_i[Y_i]
$$
其中 $Y_i\subset Y$ 是 $Y$ 中 $\delta$-维数为 $k+1$ 的各不可约分支，
且 $n_i>0$。依假设，限制
$s_i=s|_{Y_i}\in\Gamma(Y_i,\mathcal{L}|_{Y_i})$ 非零，因而是正则截面。
由引理 \ref{spaces-chow-lemma-compute-c1}，$[Z(s_i)]_k$ 表示
$c_1(\mathcal{L}|_{Y_i})$。所以依定义
$$
c_1(\mathcal{L}) \cap [Y]_{k + 1} = \sum n_i[Z(s_i)]_k
$$
结论遂由引理 \ref{spaces-chow-lemma-prepare-geometric-cap} 得出。
\end{proof}









\section{与可逆层相交及推出、拉回}
\label{spaces-chow-section-intersecting-with-divisors-push-pull}

\noindent
本节对应于《Chow 同调》第
\ref{chow-section-intersecting-with-divisors-push-pull} 节。
本节证明运算 $c_1(\mathcal{L})\cap-$ 与平坦拉回及固有推出相容。

\begin{lemma}
\label{spaces-chow-lemma-prepare-flat-pullback-cap-c1}
在情形 \ref{spaces-chow-situation-setup} 中，设 $X,Y/B$ 是良好的。
设 $f:X\to Y$ 是相对维数为 $r$ 的平坦态射，且 $\mathcal{L}$
是 $Y$ 上的可逆层。假设 $Y$ 整且 $n=\dim_\delta(Y)$。
设 $s$ 是 $\mathcal{L}$ 的非零亚纯截面。则
$$
f^*\text{div}_\mathcal{L}(s) = \sum n_i\text{div}_{f^*\mathcal{L}|_{X_i}}(s_i)
$$
在 $Z_{n+r-1}(X)$ 中成立。这里求和遍历 $X$ 中 $\delta$-维数为
$n+r$ 的各不可约分支 $X_i\subset X$；截面
$s_i=f|_{X_i}^*(s)$ 是 $s$ 的拉回；而
$n_i=m_{X_i,X}$ 是 $X_i$ 在 $X$ 中的重数。
\end{lemma}

\begin{proof}
我们用一个技巧从引理
\ref{spaces-chow-lemma-prepare-flat-pullback-rational-equivalence} 推出结论。
（另一种做法是在可逆模的亚纯截面的设定下重做该引理的证明。）
具体地，设 $q:T\to Y$ 是用 $Y$ 上的 $\mathcal{L}$ 按引理
\ref{spaces-chow-lemma-Gm-torsor} 构造的态射。我们将使用该引理中所述的
$T$ 的全部性质。考虑纤维积图
$$
\xymatrix{
T' \ar[r]_{q'} \ar[d]_h & X \ar[d]^f \\
T \ar[r]^q & Y
}
$$
则 $q':T'\to X$ 是用 $X$ 上的 $f^*\mathcal{L}$ 构造的态射。
只需证明
$$
(q')^*f^*\text{div}_\mathcal{L}(s) =
\sum n_i (q')^*\text{div}_{f^*\mathcal{L}|_{X_i}}(s_i)
$$
注意 $T'_i=q^{-1}(X_i)$ 是 $T'$ 的各不可约分支，且 $n_i$ 是
$T'_i$ 在 $T'$ 中的重数。
% P10-E0066
左端等于
$$
h^*q^*\text{div}_\mathcal{L}(s) = h^*\text{div}_T(q^*(s))
$$
（由引理 \ref{spaces-chow-lemma-Gm-torsor-divisor-meromorphic-section}
以及引理 \ref{spaces-chow-lemma-compose-flat-pullback}）。
另一方面，以 $q'_i:T'_i\to X_i$ 表示 $q'$ 的限制，则引理
\ref{spaces-chow-lemma-Gm-torsor-divisor-meromorphic-section} 还说明右端等于
$$
\sum n_i \text{div}_{T_i}((q'_i)^*(s_i))
$$
% P10-E0067
在这两个公式中，表达式 $q^*(s)$ 与 $(q'_i)^*(s_i)$ 分别表示：
通过同构 $\alpha:q^*\mathcal{L}\to\mathcal{O}_T$ 及其到 $T$
上方各空间的拉回，与 $q^*\mathcal{L}$ 和
$(q'_i)^*f^*\mathcal{L}|_{X_i}$ 的拉回亚纯截面相对应的有理函数。
按此约定，$(q'_i)^*(s_i)$ 显然是 $T$ 上的有理函数 $q^*(s)$
与态射 $h|_{T'_i}:T'_i\to T$ 的复合。因此引理
\ref{spaces-chow-lemma-prepare-flat-pullback-rational-equivalence} 恰好说明
$$
h^*\text{div}_T(q^*(s)) = \sum n_i \text{div}_{T_i}((q'_i)^*(s_i))
$$
正合所需。
\end{proof}

\begin{lemma}
\label{spaces-chow-lemma-flat-pullback-cap-c1}
在情形 \ref{spaces-chow-situation-setup} 中，设 $X,Y/B$ 是良好的。
设 $f:X\to Y$ 是相对维数为 $r$ 的平坦态射，$\mathcal{L}$ 是
$Y$ 上的可逆层，且 $\alpha$ 是 $Y$ 上的 $k$-闭链。则
$$
f^*(c_1(\mathcal{L}) \cap \alpha) = c_1(f^*\mathcal{L}) \cap f^*\alpha
$$
在 $\CH_{k+r-1}(X)$ 中成立。
\end{lemma}

\begin{proof}
写成 $\alpha=\sum n_i[W_i]$。我们将通过在 $X$ 的闭子空间
$f^{-1}(W_i)$ 上构造有理等价，证明
$$
f^*(c_1(\mathcal{L}) \cap [W_i]) = c_1(f^*\mathcal{L}) \cap f^*[W_i]
$$
在 $\CH_{k+r-1}(X)$ 中成立。由注
\ref{spaces-chow-remark-infinite-sums-rational-equivalences} 中的讨论，
这便能证明引理中的等式。

\medskip\noindent
设 $W\subset Y$ 是 $\delta$-维数为 $k$ 的整闭子空间。
考虑闭子空间 $W'=f^{-1}(W)=W\times_YX$，于是有纤维积图
$$
\xymatrix{
W' \ar[r] \ar[d]_h & X \ar[d]^f \\
W \ar[r] & Y
}
$$
必须证明
$f^*(c_1(\mathcal{L})\cap[W])
=c_1(f^*\mathcal{L})\cap f^*[W]$。
选取 $\mathcal{L}|_W$ 的非零亚纯截面 $s$。设 $W'_i\subset W'$
是 $\delta$-维数为 $k+r$ 的各不可约分支。写成
$[W']_{k+r}=\sum n_i[W'_i]$，其中依定义 $n_i$ 是 $W'_i$
在 $W'$ 中的重数。因此在 $Z_{k+r}(X)$ 中
$f^*[W]=\sum n_i[W'_i]$。由于每个 $W'_i\to W$ 都支配，
对每个 $i$，$s_i=s|_{W'_i}$ 都是非零亚纯截面。
由引理 \ref{spaces-chow-lemma-prepare-flat-pullback-cap-c1}，有闭链等式
$$
h^*\text{div}_{\mathcal{L}|_W}(s) =
\sum n_i\text{div}_{f^*\mathcal{L}|_{W'_i}}(s_i)
$$
在 $Z_{k+r-1}(W')$ 中成立。左端是 $W'$ 上的闭链，其推出在
$\CH_{k+r-1}(X)$ 中为 $f^*(c_1(\mathcal{L})\cap[W])$；
右端是 $W'$ 上的闭链，其推出在 $\CH_{k+r-1}(X)$ 中为
$c_1(f^*\mathcal{L})\cap f^*[W]$。故证明完成。
\end{proof}

\begin{lemma}
\label{spaces-chow-lemma-equal-c1-as-cycles}
在情形 \ref{spaces-chow-situation-setup} 中，设 $X,Y/B$ 是良好的。
设 $f:X\to Y$ 是固有态射，且 $\mathcal{L}$ 是 $Y$ 上的可逆层。
假设 $X,Y$ 整、$f$ 支配且 $\dim_\delta(X)=\dim_\delta(Y)$。
设 $s$ 是 $\mathcal{L}$ 在 $Y$ 上的非零亚纯截面。则
$$
f_*\left(\text{div}_{f^*\mathcal{L}}(f^*s)\right) =
[R(X) : R(Y)]\text{div}_\mathcal{L}(s).
$$
作为 $Y$ 上的闭链相等。特别地，
$$
f_*(c_1(f^*\mathcal{L}) \cap [X]) = c_1(\mathcal{L}) \cap f_*[Y].
$$
% P10-E0068
\end{lemma}

\begin{proof}
依定义 $f_*[X]=[R(X):R(Y)][Y]$，故最后一个等式由第一个等式得出。
现在证明第一个等式。
% P10-E0069
设 $q:T\to Y$ 是用 $Y$ 上的 $\mathcal{L}$ 按引理
\ref{spaces-chow-lemma-Gm-torsor} 构造的态射。我们将使用该引理中所述的
$T$ 的全部性质。考虑纤维积图
$$
\xymatrix{
T' \ar[r]_{q'} \ar[d]_h & X \ar[d]^f \\
T \ar[r]^q & Y
}
$$
则 $q':T'\to X$ 是用 $X$ 上的 $f^*\mathcal{L}$ 构造的态射。
只需在拉回到 $T'$ 后证明等式。
% P10-E0070
左端拉回为
\begin{align*}
q^*f_*\left(\text{div}_{f^*\mathcal{L}}(f^*s)\right)
& =
h_*(q')^*\text{div}_{f^*\mathcal{L}}(f^*s) \\
& =
h_*\text{div}_{(q')^*f^*\mathcal{L}}((q')^*f^*s) \\
& =
h_*\text{div}_{h^*q^*\mathcal{L}}(h^*q^*s)
\end{align*}
第一个等式由引理 \ref{spaces-chow-lemma-flat-pullback-proper-pushforward} 得出。
第二个等式使用 $T'$ 整这一事实及引理
\ref{spaces-chow-lemma-prepare-flat-pullback-cap-c1}。第三个等式来自所示图的交换性。
右端拉回为
$$
[R(X) : R(Y)]q^*\text{div}_\mathcal{L}(s) =
[R(T') : R(T)]\text{div}_{q^*\mathcal{L}}(q^*s)
$$
这由引理 \ref{spaces-chow-lemma-prepare-flat-pullback-cap-c1}、$T$ 整这一事实，
以及等式 $[R(T'):R(T)]=[R(X):R(Y)]$ 得出；略去后一等式的证明
（它可由引理 \ref{spaces-chow-lemma-flat-pullback-proper-pushforward} 得出，
但那将是一种很笨拙的证明方式）。
因此只需对 $h:T'\to T$、可逆模 $q^\mathcal{L}$ 及截面 $q^*s$
证明引理。
% P10-E0071
由于 $q^*\mathcal{L}=\mathcal{O}_T$，问题约化到下一段讨论的
$\mathcal{L}\cong\mathcal{O}$ 情形。

\medskip\noindent
假设 $\mathcal{L}=\mathcal{O}_Y$。此时 $s$ 对应于有理函数
$g\in R(Y)$。利用嵌入 $R(Y)\subset R(X)$，可以把 $g$ 看作 $X$
上的有理函数，而我们只需证明
$$
f_*\left(\text{div}_X(g)\right) = [R(X) : R(Y)]\text{div}_Y(g).
$$
与引理 \ref{spaces-chow-lemma-proper-pushforward-alteration} 的结论比较可知，
这确实成立，因为对 $g\in R(Y)^*$ 有
$\text{Nm}_{R(X)/R(Y)}(g)=g^{[R(X):R(Y)]}$。
\end{proof}

\begin{lemma}
\label{spaces-chow-lemma-pushforward-cap-c1}
在情形 \ref{spaces-chow-situation-setup} 中，设 $X,Y/B$ 是良好的。
设 $p:X\to Y$ 是固有态射，$\alpha\in Z_{k+1}(X)$，且
$\mathcal{L}$ 是 $Y$ 上的可逆层。则
$$
p_*(c_1(p^*\mathcal{L}) \cap \alpha) = c_1(\mathcal{L}) \cap p_*\alpha
$$
在 $\CH_k(Y)$ 中成立。
\end{lemma}

\begin{proof}
先假设 $p$ 具有下列性质：对每个整闭子空间 $W\subset X$，
映射 $p|_W:W\to Y$ 都是闭浸入。此时由与 $c_1(\mathcal{L})$
作帽积的定义，引理成立。

\medskip\noindent
我们用这个说明约化到一个特殊情形。写成
$\alpha=\sum n_i[W_i]$，其中 $n_i\not=0$，且各 $W_i$ 两两不同。
设 $W'_i\subset Y$ 是引理 \ref{spaces-chow-lemma-proper-image} 意义下
$W_i$ 的“像”。考虑图
$$
\xymatrix{
X' = \coprod W_i \ar[r]_-q \ar[d]_{p'} & X \ar[d]^p \\
Y' = \coprod W'_i \ar[r]^-{q'} & Y.
}
$$
由于 $\{W_i\}$ 在 $X$ 上局部有限且 $p$ 固有，可知
$\{W'_i\}$ 在 $Y$ 上局部有限，并且 $q,q',p'$ 也都是固有态射。
还可以把 $\sum n_i[W_i]$ 看作 $k$-闭链
$\alpha'\in Z_k(X')$。显然 $q_*\alpha'=\alpha$。
% P10-E0072
我们有
$q_*(c_1(q^*p^*\mathcal{L}) \cap \alpha')
= c_1(p^*\mathcal{L}) \cap q_*\alpha'$
以及
$(q')_*(c_1((q')^*\mathcal{L}) \cap p'_*\alpha') =
c_1(\mathcal{L}) \cap q'_*p'_*\alpha'$，
这是由证明开头的说明得出的。因此，只需对态射 $p'$ 与闭链
$\sum n_i[W_i]$ 证明引理。显然，这意味着可以假设 $X,Y$ 整、
$f:X\to Y$ 支配且 $\alpha=[X]$。
% P10-E0073
此时结论由引理 \ref{spaces-chow-lemma-equal-c1-as-cycles} 得出。
\end{proof}













\section{关键公式}
\label{spaces-chow-section-key}

\noindent
本节对应于《Chow 同调》第 \ref{chow-section-key} 节。
我们强烈建议读者先阅读该处的证明。

\medskip\noindent
在情形 \ref{spaces-chow-situation-setup} 中，设 $X/B$ 是良好的。假设 $X$ 整且
$\dim_\delta(X)=n$。设 $\mathcal{L}$ 与 $\mathcal{N}$ 是可逆
$\mathcal{O}_X$-模，$s$ 是 $\mathcal{L}$ 的非零亚纯截面，而
$t$ 是 $\mathcal{N}$ 的非零亚纯截面。设 $Z\subset X$ 是素因子，
其泛点为 $\xi\in|Z|$。考虑《域上的空间》第
\ref{spaces-over-fields-section-c1} 节中使用的态射
$$
c_\xi : \Spec(\mathcal{O}_{X, \xi}^h) \longrightarrow X
$$
。以 $\mathcal{L}_\xi$ 与 $\mathcal{N}_\xi$ 分别表示
$\mathcal{L}$ 与 $\mathcal{N}$ 沿 $c_\xi$ 的拉回；我们常把
$\mathcal{L}_\xi$ 与 $\mathcal{N}_\xi$ 看作由此得到的秩 $1$ 自由
$\mathcal{O}_{X,\xi}^h$-模。注意 $s$（相应地 $t$）的拉回是
$\mathcal{L}_\xi$（相应地 $\mathcal{N}_\xi$）的正则亚纯截面。

\medskip\noindent
设 $Z_i\subset X$（$i\in I$）是满足下列性质的素因子的局部有限集：
若 $Z\not \in\{Z_i\}$，则 $s$ 是 $\mathcal{L}_\xi$ 的生成元，
且 $t$ 是 $\mathcal{N}_\xi$ 的生成元。由《域上的空间》引理
\ref{spaces-over-fields-lemma-divisor-meromorphic-locally-finite}，
这样的集合存在。于是
$$
\text{div}_\mathcal{L}(s) = \sum \text{ord}_{Z_i, \mathcal{L}}(s) [Z_i]
$$
，类似地
$$
\text{div}_\mathcal{N}(t) = \sum \text{ord}_{Z_i, \mathcal{N}}(t) [Z_i]
$$
进一步展开定义，对每个 $i$ 选取生成元
$s_i\in\mathcal{L}_{\xi_i}$ 与 $t_i\in\mathcal{N}_{\xi_i}$，
其中 $\xi_i$ 是 $Z_i$ 的泛点。于是可以写成
$$
s = f_i s_i
\quad\text{且}\quad
t = g_i t_i
$$
其中 $f_i,g_i$ 是全分式环
$Q(\mathcal{O}_{X,\xi_i}^h)$ 中的可逆元素。简记
$B_i=\mathcal{O}_{X,\xi_i}^h$。记
$$
\text{ord}_{B_i} : Q(B_i)^* \longrightarrow \mathbf{Z},\quad
a/b \longmapsto
\text{length}_{B_i}(B_i/aB_i) - \text{length}_{B_i}(B_i/bB_i)
$$
。换言之，暂把《代数》定义 \ref{algebra-definition-ord}
推广到这些维数为 $1$ 的既约诺特局部环。于是依定义
$$
\text{ord}_{Z_i, \mathcal{L}}(s) = \text{ord}_{B_i}(f_i)
\quad\text{且}\quad
\text{ord}_{Z_i, \mathcal{N}}(t) = \text{ord}_{B_i}(g_i)
$$
。由于 $\xi_i$ 是 $Z_i$ 的泛点，剩余域 $\kappa(\xi_i)$ 是 $Z_i$
的函数域。此外，$\kappa(\xi_i)$ 是 $B_i$ 的剩余域；见《合宜空间》引理
\ref{decent-spaces-lemma-residue-field-henselian-local-ring}。
由于 $t_i$ 是 $\mathcal{N}_{\xi_i}$ 的生成元，它在纤维
$\mathcal{N}_{\xi_i}\otimes_{B_i}\kappa(\xi_i)$ 中的像是
$\mathcal{N}|_{Z_i}$ 的非零亚纯截面。把这个像记作 $t_i|_{Z_i}$。
由定义可得
$$
c_1(\mathcal{N}) \cap \text{div}_\mathcal{L}(s) =
\sum \text{ord}_{B_i}(f_i)
(Z_i \to X)_*\text{div}_{\mathcal{N}|_{Z_i}}(t_i|_{Z_i})
$$
以及类似的等式
$$
c_1(\mathcal{L}) \cap \text{div}_\mathcal{N}(t) =
\sum \text{ord}_{B_i}(g_i)
(Z_i \to X)_*\text{div}_{\mathcal{L}|_{Z_i}}(s_i|_{Z_i})
$$
，它们位于 $\CH_{n-2}(X)$ 中。我们将找出这两个闭链之间的有理等价。
为此考虑驯顺符号
$$
\partial_{B_i}(f_i, g_i) \in \kappa(\xi_i)^* = R(Z_i)^*
$$
；见《Chow 同调》第 \ref{chow-section-tame-symbol} 节。

\begin{lemma}[关键公式]
\label{spaces-chow-lemma-key-formula}
在上述情形中，闭链
$$
\sum
(Z_i \to X)_*\left(
\text{ord}_{B_i}(f_i) \text{div}_{\mathcal{N}|_{Z_i}}(t_i|_{Z_i}) -
\text{ord}_{B_i}(g_i) \text{div}_{\mathcal{L}|_{Z_i}}(s_i|_{Z_i}) \right)
$$
等于闭链
$$
\sum (Z_i \to X)_*\text{div}(\partial_{B_i}(f_i, g_i))
$$
\end{lemma}

\begin{proof}
证明策略如下：先约化到 $\mathcal{L}$ 与 $\mathcal{N}$ 都是平凡可逆模
的情形；再改变局部平凡化的选取；最后用平展局部化约化到概形情形
\footnote{也许可以通过立即应用平展局部化给出更短的证明。}。

\medskip\noindent
第一步。设 $q:T\to X$ 是引理 \ref{spaces-chow-lemma-Gm-torsor} 中构造的态射。
下文将直接使用该引理所述的全部性质，不再另行说明。特别地，只需证明
闭链沿 $q$ 拉回后相等。以 $s'$ 与 $t'$ 分别表示 $s$ 与 $t$
拉回得到的 $q^*\mathcal{L}$ 与 $q^*\mathcal{N}$ 的亚纯截面。
记 $Z'_i=q^{-1}(Z_i)$；以 $\xi'_i\in|Z'_i|$ 表示泛点；
记 $B'_i=\mathcal{O}_{T,\xi'_i}^h$；并以
$\mathcal{L}_{\xi'_i}$ 与 $\mathcal{N}_{\xi'_i}$ 表示
$\mathcal{L}$ 与 $\mathcal{N}$ 到 $\Spec(B'_i)$ 的拉回。
回顾有交换图
$$
\xymatrix{
\Spec(B'_i) \ar[r]_-{c_{\xi'_i}} \ar[d] & T \ar[d]^q \\
\Spec(B_i) \ar[r]^-{c_{\xi_i}} & X
}
$$
；见《合宜空间》注
\ref{decent-spaces-remark-functoriality-henselian-local-ring}。
以 $s'_i$ 与 $t'_i$ 表示 $s_i$ 与 $t_i$ 的拉回；它们分别是
$\mathcal{L}_{\xi'_i}$ 与 $\mathcal{N}_{\xi'_i}$ 的生成元。于是
$$
s' = f'_i s'_i
\quad\text{且}\quad
t' = g'_i t'_i
$$
其中 $f'_i$ 与 $g'_i$ 是 $f_i,g_i$ 在 $B_i\to B'_i$ 所诱导的
映射 $Q(B_i)\to Q(B'_i)$ 下的像。由《代数》引理
\ref{algebra-lemma-pullback-module} 有
$$
\text{ord}_{B_i}(f_i) = \text{ord}_{B'_i}(f'_i)
\quad\text{且}\quad
\text{ord}_{B_i}(g_i) = \text{ord}_{B'_i}(g'_i)
$$
把引理 \ref{spaces-chow-lemma-prepare-flat-pullback-cap-c1} 应用于
$q:Z'_i\to Z_i$，得到
$$
q^*\text{div}_{\mathcal{N}|_{Z_i}}(t_i|_{Z_i}) =
\text{div}_{q^*\mathcal{N}|_{Z'_i}}(t'_i|_{Z'_i})
\quad\text{且}\quad
q^*\text{div}_{\mathcal{L}|_{Z_i}}(s_i|_{Z_i}) =
\text{div}_{q^*\mathcal{L}|_{Z'_i}}(s'_i|_{Z'_i})
$$
这已经说明，引理陈述中的第一个闭链拉回为与
$s',t',Z'_i,s'_i,t'_i$ 相对应的闭链。为说明第二个闭链也如此，
注意由《Chow 同调》引理
\ref{chow-lemma-tame-symbol-formally-smooth} 有
$$
\partial_{B_i}(f_i, g_i) \mapsto
\partial_{B'_i}(f'_i, g'_i)
\quad\text{经由}\quad
\kappa(\xi_i) \to \kappa(\xi'_i)
$$
因此与前面相同的引理说明
$$
q^*\text{div}(\partial_{B_i}(f_i, g_i)) =
\text{div}(\partial_{B'_i}(f'_i, g'_i))
$$
由于 $q^*\mathcal{L}\cong\mathcal{O}_T$，只需在 $\mathcal{L}$
平凡时证明等式。交换 $\mathcal{L}$ 与 $\mathcal{N}$ 的角色，
同理还可假设 $\mathcal{N}$ 平凡。第一步完成。

\medskip\noindent
第二步。假设 $\mathcal{L}=\mathcal{O}_X$ 且
$\mathcal{N}=\mathcal{O}_X$。以 $1$ 表示 $\mathcal{L}$ 的平凡化截面。
则对某个单位 $u\in B_i$ 有 $s_i=u\cdot1$。考察用 $1$ 代替
$s_i$ 后的情形。此时 $f_i$ 被 $uf_i$ 代替。因此，引理第一个表达式
的第一部分不变，而在第二部分中增加
$$
\text{ord}_{B_i}(g_i)\text{div}(u|_{Z_i})
$$
，其中由《域上的空间》引理
\ref{spaces-over-fields-lemma-divisor-meromorphic-well-defined}，
$u|_{Z_i}$ 是 $u$ 在剩余域中的像；而在第二个表达式中增加
$$
\text{div}(\partial_{B_i}(u, g_i))
$$
，这是由驯顺符号的双线性得到的。这两项由《Chow 同调》公式
(\ref{chow-item-normalization}) 中给出的驯顺符号性质而相等。

\medskip\noindent
设 $Y\subset X$ 是满足 $\dim_\delta(Y)=n-2$ 的整闭子空间。
为说明 $Y$ 在引理两个闭链中的系数相同，可以像上一段一样用 $1$
代替 $s_i$。完全相同的论证说明也可以用 $1$ 代替 $t_i$。
% P10-E0074
由于只有有限多个 $Z_i$ 满足 $Y\subset Z_i$，可以对所有这些
$Z_i$ 假设 $s_i=1$ 且 $t_i=1$。

\medskip\noindent
第三步。设 $Y$ 是满足 $\dim_\delta(Y)=n-2$ 的整闭子空间，并且
$\mathcal{L}=\mathcal{O}_X$、$\mathcal{N}=\mathcal{O}_X$，而对所有
满足 $Y\subset Z_i$ 的 $Z_i$ 都有 $s_i=1$、$t_i=1$。
我们证明 $Y$ 在引理两个闭链中的系数相等。用更小的开子空间代替
$X$ 后，事实上可以假设对所有 $i$，$s_i$ 与 $t_i$ 都等于 $1$。
此时第一个闭链为零。我们的任务是证明 $Y$ 在第二个闭链中的系数也为零。

\medskip\noindent
首先，由于 $\mathcal{L}=\mathcal{O}_X$ 且
$\mathcal{N}=\mathcal{O}_X$，可以并且确实把 $s,t$ 看作 $X$ 上的
有理函数 $f,g$。由于 $s_i,t_i$ 都等于 $1$，对所有 $i$，
$f_i$（相应地 $g_i$）就是 $f$（相应地 $g$）在 $Q(B_i)$ 中的像。
设 $\zeta\in|Y|$ 是泛点。选取平展邻域
$$
(U, u) \longrightarrow (X, \zeta)
$$
并记 $Y'=\overline{\{u\}}\subset U$。由于平展态射平坦，可以把
$f,g$ 拉回为 $U$ 上的正则亚纯函数，仍分别记作 $f,g$。
对每个素因子 $Y\subset Z\subset X$，概形 $Z\times_XU$
是 $U$ 的若干素因子的并。反过来，给定素因子
$Y'\subset Z'\subset U$，存在素因子 $Y\subset Z\subset X$，
使得 $Z'$ 是 $Z\times_XU$ 的一个分支。对这样的偶 $(Z,Z')$，
环同态
$$
\mathcal{O}_{X, \xi}^h \to \mathcal{O}_{U, \xi'}^h
$$
平展（事实上有限平展）。因此由《Chow 同调》引理
\ref{chow-lemma-tame-symbol-formally-smooth} 得
$$
\partial_{\mathcal{O}_{X, \xi}^h}(f, g) \mapsto
\partial_{\mathcal{O}_{U, \xi'}^h}(f, g)
\quad\text{经由}\quad
\kappa(\xi) \to \kappa(\xi')
$$
于是引理 \ref{spaces-chow-lemma-etale-pullback-principal-divisor} 适用，并给出
$$
(Z \times_X U \to Z)^*\text{div}_Z(\partial_{\mathcal{O}_{X, \xi}^h}(f, g))
=
\sum\nolimits_{Z' \subset Z \times_X U}
\text{div}_{Z'}(\partial_{\mathcal{O}_{U, \xi'}^h}(f, g))
$$
由于平坦拉回与沿闭浸入的推出相容（引理
\ref{spaces-chow-lemma-flat-pullback-proper-pushforward}），只需证明
$Y'$ 在
$$
\sum\nolimits_{Z' \subset U} 
(Z' \to U)_*\text{div}_{Z'}(\partial_{\mathcal{O}_{U, \xi'}^h}(f, g))
$$
中的系数为零。

\medskip\noindent
令 $A=\mathcal{O}_{U,u}$。则 $f,g\in Q(A)^*$。
因而可写成 $f=a/b$、$g=c/d$，其中 $a,b,c,d\in A$ 都是非零因子。
上述表达式中 $Y'$ 的系数为
$$
\sum\nolimits_{\mathfrak q \subset A\text{ 高度 }1}
\text{ord}_{A/\mathfrak q}(\partial_{A_\mathfrak q}(f, g))
$$
% P10-E0075
由 $\partial_A$ 的双线性，只需证明
$$
\sum\nolimits_{\mathfrak q \subset A\text{ 高度 }1}
\text{ord}_{A/\mathfrak q}(\partial_{A_\mathfrak q}(a, c))
$$
为零；对其他各偶 $(a,d)$、$(b,c)$ 与 $(b,d)$ 也同样如此。
这由《Chow 同调》引理
\ref{chow-lemma-key-nonzerodivisors} 得出。
\end{proof}











\section{与可逆层相交及有理等价}
\label{spaces-chow-section-commutativity}

\noindent
本节对应于《Chow 同调》第 \ref{chow-section-commutativity} 节。
应用关键引理，我们得到与可逆层相交的基本性质。特别地，将会看到
$c_1(\mathcal{L})\cap-$ 可通过有理等价分解，并且不同可逆层所对应的
这些运算彼此交换。

\begin{lemma}
\label{spaces-chow-lemma-commutativity-on-integral}
在情形 \ref{spaces-chow-situation-setup} 中，设 $X/B$ 是良好的。
假设 $X$ 整且 $\dim_\delta(X)=n$。设
$\mathcal{L},\mathcal{N}$ 是 $X$ 上的可逆层。选取
$\mathcal{L}$ 的非零亚纯截面 $s$ 与 $\mathcal{N}$ 的非零亚纯截面
$t$。令 $\alpha=\text{div}_\mathcal{L}(s)$，
$\beta=\text{div}_\mathcal{N}(t)$。则
$$
c_1(\mathcal{N}) \cap \alpha
=
c_1(\mathcal{L}) \cap \beta
$$
在 $\CH_{n-2}(X)$ 中成立。
\end{lemma}

\begin{proof}
立即由关键引理 \ref{spaces-chow-lemma-key-formula} 及其前面的讨论得出。
\end{proof}

\begin{lemma}
\label{spaces-chow-lemma-factors}
在情形 \ref{spaces-chow-situation-setup} 中，设 $X/B$ 是良好的，且
$\mathcal{L}$ 是 $X$ 上的可逆层。运算
$\alpha\mapsto c_1(\mathcal{L})\cap\alpha$ 可通过有理等价分解，
从而给出运算
$$
c_1(\mathcal{L}) \cap - : \CH_{k + 1}(X) \to \CH_k(X)
$$
\end{lemma}

\begin{proof}
设 $\alpha\in Z_{k+1}(X)$ 且 $\alpha\sim_{rat}0$。
必须证明定义 \ref{spaces-chow-definition-cap-c1} 中所定义的
$c_1(\mathcal{L})\cap\alpha$ 为零。由定义
\ref{spaces-chow-definition-rational-equivalence}，存在满足
$\dim_\delta(W_j)=k+2$ 的整闭子空间的局部有限族 $\{W_j\}$
及有理函数 $f_j\in R(W_j)^*$，使得
$$
\alpha = \sum (i_j)_*\text{div}_{W_j}(f_j)
$$
注意 $p:\coprod W_j\to X$ 是固有态射，故
$\alpha=p_*\alpha'$，其中 $\alpha'\in Z_{k+1}(\coprod W_j)$
是各主因子 $\text{div}_{W_j}(f_j)$ 的和。由引理
\ref{spaces-chow-lemma-pushforward-cap-c1}，
$c_1(\mathcal{L})\cap\alpha
=p_*(c_1(p^*\mathcal{L})\cap\alpha')$。
所以只需证明每个
$c_1(\mathcal{L}|_{W_j})\cap\text{div}_{W_j}(f_j)$ 都为零。
换言之，可以假设 $X$ 整且
$\alpha=\text{div}_X(f)$，其中 $f\in R(X)^*$。

\medskip\noindent
假设 $X$ 整，且对某个 $f\in R(X)^*$ 有
$\alpha=\text{div}_X(f)$。可以把 $f$ 看作可逆层
$\mathcal{N}=\mathcal{O}_X$ 的正则亚纯截面。选取 $\mathcal{L}$
的一个亚纯截面 $s$，并记
$\beta=\text{div}_\mathcal{L}(s)$。
% P10-E0077
由引理 \ref{spaces-chow-lemma-commutativity-on-integral} 得
$$
c_1(\mathcal{L}) \cap \alpha = c_1(\mathcal{O}_X) \cap \beta.
$$
但由引理 \ref{spaces-chow-lemma-c1-cap-additive}，右端在 $\CH_k(X)$ 中为零，
正合所需。
\end{proof}

\noindent
在情形 \ref{spaces-chow-situation-setup} 中，设 $X/B$ 是良好的，且
$\mathcal{L}$ 是 $X$ 上的可逆层。记
$$
c_1(\mathcal{L})^s \cap - : \CH_{k + s}(X) \to \CH_k(X)
$$
为运算 $c_1(\mathcal{L})\cap-$。由引理 \ref{spaces-chow-lemma-factors}，
这有意义。对所有 $s\geq0$，以
$c_1(\mathcal{L}^s\cap-$ 表示该运算的 $s$ 次迭代。
% P10-E0076

\begin{lemma}
\label{spaces-chow-lemma-cap-commutative}
在情形 \ref{spaces-chow-situation-setup} 中，设 $X/B$ 是良好的。
设 $\mathcal{L},\mathcal{N}$ 是 $X$ 上的可逆层。
对任意 $\alpha\in\CH_{k+2}(X)$，有
$$
c_1(\mathcal{L}) \cap c_1(\mathcal{N}) \cap \alpha
=
c_1(\mathcal{N}) \cap c_1(\mathcal{L}) \cap \alpha
$$
作为 $\CH_k(X)$ 的元素相等。
\end{lemma}

\begin{proof}
对某个满足 $\dim_\delta(Z_j)=k+2$ 的整闭子空间
$Z_j\subset X$ 的局部有限族，写成
$\alpha=\sum m_j[Z_j]$。考虑固有态射
$p:\coprod Z_j\to X$。把
$\alpha'=\sum m_j[Z_j]$ 看作 $\coprod Z_j$ 上的 $(k+2)$-闭链。
多次应用引理 \ref{spaces-chow-lemma-pushforward-cap-c1}，可知
$c_1(\mathcal{L}) \cap c_1(\mathcal{N}) \cap \alpha
= p_*(c_1(p^*\mathcal{L}) \cap c_1(p^*\mathcal{N}) \cap \alpha')$
以及
$c_1(\mathcal{N}) \cap c_1(\mathcal{L}) \cap \alpha
= p_*(c_1(p^*\mathcal{N}) \cap c_1(p^*\mathcal{L}) \cap \alpha')$.
因此只需在 $X$ 整且 $\alpha=[X]$ 时证明公式。
此时结论由引理 \ref{spaces-chow-lemma-commutativity-on-integral} 与各定义得出。
\end{proof}

















\section{与有效 Cartier 因子相交}
\label{spaces-chow-section-intersecting-effective-Cartier}

\noindent
本节对应于《Chow 同调》第
\ref{chow-section-intersecting-effective-Cartier} 节。
请阅读该节的导言以了解动机。% P10-E0078

\medskip\noindent
回顾，有效 Cartier 因子与二元组 $(\mathcal{L},s)$ 的同构类一一对应，
其中 $\mathcal{L}$ 是可逆层，$s$ 是整体截面；见《空间上的因子》引理
\ref{spaces-divisors-lemma-characterize-OD}。若 $D$ 对应于
$(\mathcal{L},s)$，则 $\mathcal{L}=\mathcal{O}_X(D)$。
阅读本节时请始终记住这一点。

\begin{definition}
\label{spaces-chow-definition-gysin-homomorphism}
在情形 \ref{spaces-chow-situation-setup} 中，设 $X/B$ 是良好的。
设 $(\mathcal{L},s)$ 是由可逆层与整体截面
$s\in\Gamma(X,\mathcal{L})$ 组成的二元组。
令 $D=Z(s)$ 为 $s$ 的零化轨迹，并以 $i:D\to X$ 表示闭浸入。
对每个整数 $k$，定义一个（精细）{\it Gysin 同态}
$$
i^* : Z_{k + 1}(X) \to \CH_k(D).
$$
规则如下：
\begin{enumerate}
\item 给定整闭子空间 $W\subset X$，其中
$\dim_\delta(W)=k+1$，定义
\begin{enumerate}
\item 若 $W\not\subset D$，则令 $i^*[W]=[D\cap W]_k$，
视为 $D$ 上的 $k$-闭链；
\item 若 $W\subset D$，则令
$i^*[W] = i'_*(c_1(\mathcal{L}|_W) \cap [W])$,
其中 $i':W\to D$ 是诱导的闭浸入。
\end{enumerate}
\item 对一般的 $(k+1)$-闭链 $\alpha=\sum n_j[W_j]$，令
$$
i^*\alpha = \sum n_j i^*[W_j]
$$
\item 若 $D$ 是有效 Cartier 因子，则以
$D\cdot\alpha=i_*i^*\alpha$ 表示把该类推出为 $X$ 上的类。
\end{enumerate}
\end{definition}

\noindent
事实上，后文将看到，这个 Gysin 同态 $i^*$ 可视为非平坦拉回的一例。
因此，我们有时会非正式地把类 $i^*\alpha$ 称为类 $\alpha$ 的
{\it 拉回}。

\begin{remark}
\label{spaces-chow-remark-on-cycles}
设 $S$、$B$、$X$、$\mathcal{L}$、$s$、$i:D\to X$ 如定义
\ref{spaces-chow-definition-gysin-homomorphism}，并假设
$\mathcal{L}|_D\cong\mathcal{O}_D$。此时可以在闭链上定义典范映射
$i^*:Z_{k+1}(X)\to Z_k(D)$，方法是要求每当 $W\subset D$ 时都有
$i^*[W]=0$。以后会用到这种做法。
\end{remark}

\begin{remark}
\label{spaces-chow-remark-pullback-pairs}
在情形 \ref{spaces-chow-situation-setup} 中，设 $f:X'\to X$ 是 $B$ 上良好代数空间
之间的态射。设 $(\mathcal{L},s,i:D\to X)$ 是定义
\ref{spaces-chow-definition-gysin-homomorphism} 中的三元组。令
$\mathcal{L}'=f^*\mathcal{L}$、$s'=f^*s$ 以及
$D'=X'\times_XD=Z(s')$，便得到交换图
$$
\xymatrix{
D' \ar[d]_g \ar[r]_{i'} & X' \ar[d]^f \\
D \ar[r]^i & X
}
$$
于是可以讨论 $i^*$ 与 $(i')^*$ 之间的各种相容性。
\end{remark}

\begin{lemma}
\label{spaces-chow-lemma-support-cap-effective-Cartier}
在情形 \ref{spaces-chow-situation-setup} 中，设 $X/B$ 是良好的。设
$(\mathcal{L},s,i:D\to X)$ 如定义
\ref{spaces-chow-definition-gysin-homomorphism}，并设 $\alpha$ 是 $X$ 上的
$(k+1)$-闭链。则在 $\CH_k(X)$ 中有
$i_*i^*\alpha=c_1(\mathcal{L})\cap\alpha$。特别地，若 $D$ 是有效
Cartier 因子，则
$D\cdot\alpha=c_1(\mathcal{O}_X(D))\cap\alpha$。
\end{lemma}

\begin{proof}
写成 $\alpha=\sum n_j[W_j]$，其中 $i_j:W_j\to X$ 是
$\dim_\delta(W_j)=k$ 的整闭子空间。% P10-E0079
由于 $D$ 是 $s$ 的零化轨迹，$D\cap W_j$ 是限制
$s|_{W_j}$ 的零化轨迹。因此，对每个满足 $W_j\not\subset D$ 的 $j$，
由引理 \ref{spaces-chow-lemma-geometric-cap} 有
$c_1(\mathcal{L})\cap[W_j]=[D\cap W_j]_k$。从而
% P10-E0080
$$
c_1(\mathcal{L}) \cap \alpha
=
\sum\nolimits_{W_j \not \subset D} n_j[D \cap W_j]_k
+
\sum\nolimits_{W_j \subset D}
n_j i_{j, *}(c_1(\mathcal{L})|_{W_j}) \cap [W_j])
$$
在 $\CH_k(X)$ 中成立，这是由定义 \ref{spaces-chow-definition-cap-c1} 得到的。
右端逐项等于定义 \ref{spaces-chow-definition-gysin-homomorphism} 中 $D$ 上的类
$i^*\alpha$ 的推出。因此结论成立。
\end{proof}

\begin{lemma}
\label{spaces-chow-lemma-closed-in-X-gysin}
在情形 \ref{spaces-chow-situation-setup} 中，设 $f:X'\to X$ 是 $B$ 上良好代数空间
之间的固有态射。设 $(\mathcal{L},s,i:D\to X)$ 如定义
\ref{spaces-chow-definition-gysin-homomorphism}。作图
$$
\xymatrix{
D' \ar[d]_g \ar[r]_{i'} & X' \ar[d]^f \\
D \ar[r]^i & X
}
$$
如注 \ref{spaces-chow-remark-pullback-pairs}。对 $X'$ 上任意 $(k+1)$-闭链
$\alpha'$，在 $\CH_k(D)$ 中有
$i^*f_*\alpha'=g_*(i')^*\alpha'$（这是有意义的，因为 $f_*$
在闭链层次上已有定义）。
\end{lemma}

\begin{proof}
设 $\alpha=[W']$，其中 $W'\subset X'$ 是某个整闭子空间。% P10-E0081
令 $W\subset X$ 为引理 \ref{spaces-chow-lemma-proper-image} 意义下 $W'$ 的“像”。
若 $W'\not\subset D'$，则 $W\not\subset D$，并且
$$
[W' \cap D']_k = \text{div}_{\mathcal{L}'|_{W'}}({s'|_{W'}})
\quad\text{且}\quad
[W \cap D]_k = \text{div}_{\mathcal{L}|_W}(s|_W)
$$
故由引理 \ref{spaces-chow-lemma-equal-c1-as-cycles}，第一个闭链在 $f_*$ 下的像
等于第二个闭链。因此等式在闭链层次成立。若 $W'\subset D'$，则
$W\subset D$，并且由引理 \ref{spaces-chow-lemma-equal-c1-as-cycles} 的第二个断言，
$f_*(c_1(\mathcal{L}|_{W'})\cap[W'])$ 在 $\CH_k(W)$ 中等于
$c_1(\mathcal{L}|_W)\cap[W]$。
由注 \ref{spaces-chow-remark-infinite-sums-rational-equivalences}，结论对一般的
$\alpha'$ 也成立。
\end{proof}

\begin{lemma}
\label{spaces-chow-lemma-gysin-flat-pullback}
在情形 \ref{spaces-chow-situation-setup} 中，设 $f:X'\to X$ 是 $B$ 上良好代数空间
之间相对维数为 $r$ 的平坦态射。设
$(\mathcal{L},s,i:D\to X)$ 如定义
\ref{spaces-chow-definition-gysin-homomorphism}。作图
$$
\xymatrix{
D' \ar[d]_g \ar[r]_{i'} & X' \ar[d]^f \\
D \ar[r]^i & X
}
$$
如注 \ref{spaces-chow-remark-pullback-pairs}。对 $X$ 上任意 $(k+1)$-闭链
$\alpha$，有
% P10-E0082
% P10-E0083
在 $\CH_{k + r}(D)$ 中有
$(i')^*f^*\alpha = g^*i^*\alpha'$
（这是有意义的，因为 $f^*$ 在闭链层次上已有定义）。
\end{lemma}

\begin{proof}
设 $\alpha=[W]$，其中 $W\subset X$ 是某个整闭子空间。
令 $W'=f^{-1}(W)\subset X'$。若 $W\not\subset D$，则
$W'\not\subset D'$，并且
$$
W' \cap D' = g^{-1}(W \cap D)
$$
作为 $D'$ 的闭子空间相等。因此由引理
\ref{spaces-chow-lemma-pullback-coherent}，等式在闭链层次成立。
若 $W\subset D$，则 $W'\subset D'$ 且 $W'=g^{-1}(W)$，并有
$[W']_{k+1+r}=g^*[W]$；由引理 \ref{spaces-chow-lemma-flat-pullback-cap-c1}，
等式在 $\CH_{k+r}(D')$ 中成立。
由注 \ref{spaces-chow-remark-infinite-sums-rational-equivalences}，结论对一般的
$\alpha'$ 也成立。% P10-E0082
\end{proof}

\begin{lemma}
\label{spaces-chow-lemma-easy-gysin}
在情形 \ref{spaces-chow-situation-setup} 中，设 $X/B$ 是良好的。设
$(\mathcal{L},s,i:D\to X)$ 如定义
\ref{spaces-chow-definition-gysin-homomorphism}。设 $Z\subset X$ 是闭子概形，
满足 $\dim_\delta(Z)\leq k+1$，且 $D\cap Z$ 是 $Z$ 上的有效
Cartier 因子。则
$i^*([Z]_{k+1})=[D\cap Z]_k$。
\end{lemma}

\begin{proof}
该假设意味着 $s|_Z$ 是 $\mathcal{L}|_Z$ 的正则截面。
因此 $D\cap Z=Z(s)$，并有
$$
[D \cap Z]_k = \sum n_i [Z(s_i)]_k
$$
作为闭链相等；这里 $s_i=s|_{Z_i}$，$Z_i$ 是 $\delta$-维数为
$k+1$ 的不可约分支，并且 $[Z]_{k+1}=\sum n_i[Z_i]$。
见引理 \ref{spaces-chow-lemma-prepare-geometric-cap}。又有
$D\cap Z_i=Z(s_i)$。与 Gysin 映射的定义比较即得结论。
\end{proof}













\section{Gysin 同态}
\label{spaces-chow-section-gysin}

\noindent
本节对应于《Chow 同调》第 \ref{chow-section-gysin} 节。
本节利用关键公式证明 Gysin 同态可通过有理等价因子化。% P10-E0084

\begin{lemma}
\label{spaces-chow-lemma-gysin-factors-general}
在情形 \ref{spaces-chow-situation-setup} 中，设 $X/B$ 是良好的。
假设 $X$ 整且 $n=\dim_\delta(X)$。设 $i:D\to X$ 是有效
Cartier 因子。设 $\mathcal{N}$ 是可逆 $\mathcal{O}_X$-模，
$t$ 是 $\mathcal{N}$ 的非零亚纯截面。则在 $\CH_{n-2}(D)$ 中有
$i^*\text{div}_\mathcal{N}(t)=c_1(\mathcal{N})\cap[D]_{n-1}$。
\end{lemma}

\begin{proof}
写成
$\text{div}_\mathcal{N}(t)=\sum\text{ord}_{Z_i,\mathcal{N}}(t)[Z_i]$，
其中 $Z_i\subset X$ 是一些 $\delta$-维数为 $n-1$ 的整闭子空间。
不妨假设族 $\{Z_i\}$ 局部有限，$t\in\Gamma(U,\mathcal{N}|_U)$
是生成元，其中 $U=X\setminus\bigcup Z_i$，并且 $D$ 的每个不可约分支
都是某个 $Z_i$；见《域上的空间》引理
\ref{spaces-over-fields-lemma-components-locally-finite}、
\ref{spaces-over-fields-lemma-divisor-locally-finite} 以及
\ref{spaces-over-fields-lemma-divisor-meromorphic-locally-finite}。

\medskip\noindent
令 $\mathcal{L}=\mathcal{O}_X(D)$，并以
$s\in\Gamma(X,\mathcal{O}_X(D))=\Gamma(X,\mathcal{L})$
表示典范截面。我们把第 \ref{spaces-chow-section-key} 节的讨论用于当前情形。
对每个 $i$，令 $\xi_i\in|Z_i|$ 为其一般点，并令
$B_i=\mathcal{O}_{X,\xi_i}^h$。对每个 $i$，选取
$\mathcal{L}_{\xi_i}$ 的生成元 $s_i$ 和 $\mathcal{N}_{\xi_i}$ 的
生成元 $t_i$；二者均作为 $B_i$-模的生成元，但当
$Z_i\not\subset D$ 时，要求选取 $s_i=s$。写成
$s=f_is_i$、$t=g_it_i$，其中 $f_i,g_i\in B_i$。% P10-E0085
于是 $\text{ord}_{Z_i,\mathcal{N}}(t)=\text{ord}_{B_i}(g_i)$。
另一方面，$f_i\in B_i$，并且
$$
[D]_{n - 1} = \sum \text{ord}_{B_i}(f_i)[Z_i]
$$
这是由 $s_i$ 的选取方式得到的。我们断言
$$
i^*\text{div}_\mathcal{N}(t) =
\sum \text{ord}_{B_i}(g_i) \text{div}_{\mathcal{L}|_{Z_i}}(s_i|_{Z_i})
$$
在闭链层次成立。更确切地说，右端是代表左端的一个闭链。
这由 $\text{div}_\mathcal{N}(t)$ 的公式以及以下事实立即得到：当
$Z_i\not\subset D$ 时，
$\text{div}_{\mathcal{L}|_{Z_i}}(s_i|_{Z_i}) = [Z(s_i|_{Z_i})]_{n - 2} =
[Z_i \cap D]_{n - 2}$，因为此时 $s_i|_{Z_i}=s|_{Z_i}$ 是正则截面；
见引理 \ref{spaces-chow-lemma-compute-c1}。类似地，
$$
c_1(\mathcal{N}) \cap [D]_{n - 1} =
\sum \text{ord}_{B_i}(f_i) \text{div}_{\mathcal{N}|_{Z_i}}(t_i|_{Z_i})
$$
关键公式（引理 \ref{spaces-chow-lemma-key-formula}）给出闭链等式
$$
\sum \left(
\text{ord}_{B_i}(f_i) \text{div}_{\mathcal{N}|_{Z_i}}(t_i|_{Z_i}) -
\text{ord}_{B_i}(g_i) \text{div}_{\mathcal{L}|_{Z_i}}(s_i|_{Z_i}) \right) =
\sum \text{div}_{Z_i}(\partial_{B_i}(f_i, g_i))
$$
若 $Z_i\not\subset D$，则 $f_i=1$，从而
$\text{div}_{Z_i}(\partial_{B_i}(f_i,g_i))=0$。因此，在 $D$ 上，
代表 $i^*\text{div}_\mathcal{N}(t)$ 与
$c_1(\mathcal{N})\cap[D]_{n-1}$ 的上述具体闭链彼此有理等价。
证明完毕。
\end{proof}

\begin{lemma}
\label{spaces-chow-lemma-gysin-factors}
在情形 \ref{spaces-chow-situation-setup} 中，设 $X/B$ 是良好的。设
$(\mathcal{L},s,i:D\to X)$ 如定义
\ref{spaces-chow-definition-gysin-homomorphism}。Gysin 同态通过有理等价因子化，
从而给出映射 $i^*:\CH_{k+1}(X)\to\CH_k(D)$。
\end{lemma}

\begin{proof}
设 $\alpha\in Z_{k+1}(X)$，并假设 $\alpha\sim_{rat}0$。
这意味着存在一族局部有限的、$\delta$-维数为 $k+2$ 的整闭子空间
$W_j\subset X$，以及 $f_j\in R(W_j)^*$，使得
$\alpha=\sum i_{j,*}\text{div}_{W_j}(f_j)$。令
$X'=\coprod W_i$，并考虑图% P10-E0086
$$
\xymatrix{
D' \ar[d]_q \ar[r]_{i'} & X' \ar[d]^p \\
D \ar[r]^i & X
}
$$
它来自注 \ref{spaces-chow-remark-pullback-pairs}。由于 $X'\to X$ 固有，由引理
\ref{spaces-chow-lemma-closed-in-X-gysin} 有 $i^*p_*=q_*(i')^*$。
又已知 $q_*$ 通过有理等价因子化（引理
\ref{spaces-chow-lemma-proper-pushforward-rational-equivalence}），故只需对 $X'$ 上
$\alpha'=\sum\text{div}_{W_j}(f_j)$ 证明结论。显然，这把我们约化到
$X$ 整且 $\alpha=\text{div}(f)$ 的情形，其中 $f\in R(X)^*$。

\medskip\noindent
假设 $X$ 整，且对某个 $f\in R(X)^*$ 有
$\alpha=\text{div}(f)$。若 $X=D$，则 $i^*\alpha$ 等于
$c_1(\mathcal{L})\cap\alpha$；由引理 \ref{spaces-chow-lemma-factors}，它有理等价于零。
若 $D\not = X$，则由引理 \ref{spaces-chow-lemma-gysin-factors-general}，
$i^*\text{div}_X(f)$ 在 $\CH_k(D)$ 中等于
$c_1(\mathcal{O}_D)\cap[D]_{n-1}$。% P10-E0087
当然，与 $c_1(\mathcal{O}_D)$ 相交是零映射。
\end{proof}

\begin{lemma}
\label{spaces-chow-lemma-gysin-commutes-cap-c1}
在情形 \ref{spaces-chow-situation-setup} 中，设 $X/B$ 是良好的。设
$(\mathcal{L},s,i:D\to X)$ 是定义
\ref{spaces-chow-definition-gysin-homomorphism} 中的三元组。
设 $\mathcal{N}$ 是可逆 $\mathcal{O}_X$-模。则对所有
$\alpha\in\CH_k(Z)$，在 $\CH_{k-2}(D)$ 中有% P10-E0088
$i^*(c_1(\mathcal{N})\cap\alpha)=c_1(i^*\mathcal{N})\cap i^*\alpha$。
\end{lemma}

\begin{proof}
采用与引理 \ref{spaces-chow-lemma-gysin-factors} 完全相同的证明，由引理
\ref{spaces-chow-lemma-pushforward-cap-c1}、
\ref{spaces-chow-lemma-cap-commutative} 以及
\ref{spaces-chow-lemma-gysin-factors-general} 即得结论。
\end{proof}

\begin{lemma}
\label{spaces-chow-lemma-gysin-commutes-gysin}
在情形 \ref{spaces-chow-situation-setup} 中，设 $X/B$ 是良好的。设
$(\mathcal{L},s,i:D\to X)$ 和
$(\mathcal{L}',s',i':D'\to X)$ 是定义
\ref{spaces-chow-definition-gysin-homomorphism} 中的两个三元组。则图
$$
\xymatrix{
\CH_k(X) \ar[r]_{i^*} \ar[d]_{(i')^*} & \CH_{k - 1}(D) \ar[d] \\
\CH_{k - 1}(D') \ar[r] & \CH_{k - 2}(D \cap D')
}
$$
交换，其中每个映射都是 Gysin 映射。
\end{lemma}

\begin{proof}
以 $j:D\cap D'\to D$ 和 $j':D\cap D'\to D'$ 表示分别对应于
$(\mathcal{L}|_{D'},s|_{D'}$ 与
$(\mathcal{L}'_D,s|_D)$ 的闭浸入。% P10-E0089
需要证明，对所有 $\alpha\in\CH_k(X)$ 都有
$(j')^*i^*\alpha=j^*(i')^*\alpha$。设 $W\subset X$ 是维数为 $k$ 的
整闭子概形。% P10-E0090
我们先对 $\alpha=[W]$ 证明该等式。一般情形随后由注
\ref{spaces-chow-remark-infinite-sums-rational-equivalences} 中的观察
（以及下文构造的有理等价的具体形式）得出。
我们将从关键公式推出 $\alpha=[W]$ 时的等式。

\medskip\noindent
令 $\sigma$ 是 $\mathcal{L}|_W$ 的非零亚纯截面；当
$W\not\subset D$ 时，要求它等于 $s|_W$。令 $\sigma'$ 是
$\mathcal{L}'|_W$ 的非零亚纯截面；当 $W\not\subset D'$ 时，
要求它等于 $s'|_W$。写成
$$
\text{div}_{\mathcal{L}|_W}(\sigma) =
\sum \text{ord}_{Z_i, \mathcal{L}|_W}(\sigma)[Z_i] = \sum n_i[Z_i]
$$
以及类似的
$$
\text{div}_{\mathcal{L}'|_W}(\sigma') =
\sum \text{ord}_{Z_i, \mathcal{L}'|_W}(\sigma')[Z_i] = \sum n'_i[Z_i]
$$
如第 \ref{spaces-chow-section-key} 节的讨论。于是，当 $n_i\not =0$ 时有
$Z_i\subset D$；当 $n'_i\not =0$ 时有 $Z'_i\subset D'$。% P10-E0091
对每个 $i$，令 $\xi_i\in|Z_i|$ 为一般点。如第
\ref{spaces-chow-section-key} 节那样，对每个 $i$ 选取元素
$\sigma_i\in\mathcal{L}_{\xi_i}$ 和
$\sigma'_i\in\mathcal{L}'_{\xi_i}$，它们分别生成
$B_i=\mathcal{O}_{W,\xi_i}^h$ 上的相应模；并且当
$Z_i\not\subset D$ 时，前者等于 $s$ 的像，而当
$Z_i\not\subset D'$ 时，后者等于 $s'$ 的像。
写成 $\sigma=f_i\sigma_i$、$\sigma'=f'_i\sigma'_i$，于是
$n_i=\text{ord}_{B_i}(f_i)$ 且
$n'_i=\text{ord}_{B_i}(f'_i)$。由定义可得闭链等式
$$
(j')^*i^*[W] =
\sum \text{ord}_{B_i}(f_i) \text{div}_{\mathcal{L}'|_{Z_i}}(\sigma'_i|_{Z_i})
$$
以及
$$
j^*(i')^*[W] =
\sum \text{ord}_{B_i}(f'_i) \text{div}_{\mathcal{L}|_{Z_i}}(\sigma_i|_{Z_i})
$$
关键公式（引理 \ref{spaces-chow-lemma-key-formula}）现在给出闭链等式
$$
\sum \left(
\text{ord}_{B_i}(f_i) \text{div}_{\mathcal{L}'|_{Z_i}}(\sigma'_i|_{Z_i}) -
\text{ord}_{B_i}(f'_i) \text{div}_{\mathcal{L}|_{Z_i}}(\sigma_i|_{Z_i})
\right) =
\sum \text{div}_{Z_i}(\partial_{B_i}(f_i, f'_i))
$$
注意，若 $Z_i\not\subset D\cap D'$，则
$\text{div}_{Z_i}(\partial_{B_i}(f_i,f'_i))=0$，因为此时
$f_i=1$ 或 $f'_i=1$。因此，在 $D\cap D'\cap W$ 上，代表
$(j')^*i^*[W]$ 与 $j^*(i')^*[W]$ 的上述具体闭链彼此有理等价。
\end{proof}










\section{相对有效 Cartier 因子}
\label{spaces-chow-section-relative-effective-cartier}

\noindent
本节对应于《Chow 同调》第
\ref{chow-section-relative-effective-cartier} 节。
相对有效 Cartier 因子定义于《空间上的因子》第
\ref{spaces-divisors-section-effective-Cartier-morphisms} 节。
为了建立向量丛陈类的基本结果，我们只需要环境概形与有效 Cartier 因子
都在基上平坦的情形。% P10-E0092

\begin{lemma}
\label{spaces-chow-lemma-relative-effective-cartier}
在情形 \ref{spaces-chow-situation-setup} 中，设 $X,Y/B$ 是良好的。设
$p:X\to Y$ 是相对维数为 $r$ 的平坦态射。设 $i:D\to X$ 是相对
有效 Cartier 因子（《空间上的因子》定义
\ref{spaces-divisors-definition-relative-effective-Cartier-divisor}）。
令 $\mathcal{L}=\mathcal{O}_X(D)$。对任意
$\alpha\in\CH_{k+1}(Y)$，在 $\CH_{k+r}(D)$ 中有
$$
i^*p^*\alpha = (p|_D)^*\alpha
$$
并且在 $\CH_{k+r}(X)$ 中有
$$
c_1(\mathcal{L}) \cap p^*\alpha = i_* ((p|_D)^*\alpha)
$$
\end{lemma}

\begin{proof}
设 $W\subset Y$ 是 $\delta$-维数为 $k+1$ 的整闭子空间。
由《空间上的因子》引理
\ref{spaces-divisors-lemma-relative-Cartier}，
$D\cap p^{-1}W$ 是 $p^{-1}W$ 上的有效 Cartier 因子。
由引理 \ref{spaces-chow-lemma-easy-gysin}，下式的第一个等号成立：
$$
i^*[p^{-1}W]_{k + r + 1} =
[D \cap p^{-1}W]_{k + r} =
[(p|_D)^{-1}(W)]_{k + r}.
$$
第二个等号成立，是因为作为代数空间有
$D\cap p^{-1}(W)=(p|_D)^{-1}(W)$。由定义，
$p^*[W]=[p^{-1}W]_{k+r+1}$，故在闭链层次有
$i^*p^*[W]=(p|_D)^*[W]$。若
$\alpha=\sum m_j[W_j]$ 是一般的 $(k+1)$-闭链，则在闭链层次有
$i^*\alpha=\sum m_j i^*p^*[W_j]=\sum m_j(p|_D)^*[W_j]$。% P10-E0093
这证明了第一个等式。% P10-E0094
对第一个等式应用引理 \ref{spaces-chow-lemma-support-cap-effective-Cartier}，
即可推出第二个等式。
\end{proof}












\section{仿射丛}
\label{spaces-chow-section-affine-vector}

\noindent
本节对应于《Chow 同调》第 \ref{chow-section-affine-vector} 节。
对仿射丛而言，周群上的拉回映射是满射。

\begin{lemma}
\label{spaces-chow-lemma-pullback-affine-fibres-surjective}
在情形 \ref{spaces-chow-situation-setup} 中，设 $X,Y/B$ 是良好的。设
$f:X\to Y$ 是 $B$ 上相对维数为 $r$ 的拟紧平坦态射。
假设对每个 $y\in Y$ 都有 $X_y\cong\mathbf{A}^r_{\kappa(y)}$。
则对所有 $k\in\mathbf{Z}$，
$f^*:\CH_k(Y)\to\CH_{k+r}(X)$ 都是满射。
\end{lemma}

\begin{proof}
设 $\alpha\in\CH_{k+r}(X)$。写成 $\alpha=\sum m_j[W_j]$，其中
$m_j\not =0$，且 $W_j$ 是两两不同的、$\delta$-维数为 $k+r$ 的
整闭子空间。于是族 $\{W_j\}$ 在 $X$ 中局部有限。
如引理 \ref{spaces-chow-lemma-proper-image}，令 $Z_j\subset Y$ 为使我们得到
支配态射 $W_j\to Z_j$ 的整闭子空间。对 $Y$ 中任意拟紧开集
$V\subset Y$，只有有限多个 $j$ 使 $f^{-1}(V)\cap W_j$ 非空。
因此，作为这些像之闭包的 $Z_j$ 构成 $Y$ 中一族局部有限的整闭子空间。

\medskip\noindent
考虑纤维积图
$$
\xymatrix{
f^{-1}(Z_j) \ar[r] \ar[d]_{f_j} & X \ar[d]^f \\
Z_j \ar[r] & Y
}
$$
假设对某个 $k$-闭链 $\beta_j\in\CH_k(Z_j)$，
$[W_j]\in Z_{k+r}(f^{-1}(Z_j))$ 有理等价于 $f_j^*\beta_j$。
那么 $\beta=\sum m_j\beta_j$ 是 $Y$ 上的 $k$-闭链，且
$f^*\beta=\sum m_jf_j^*\beta_j$ 与 $\alpha$ 有理等价
（见注 \ref{spaces-chow-remark-infinite-sums-rational-equivalences}）。
这把我们约化到 $Y$ 整，且 $\alpha=[W]$ 的情形，其中 $W$ 是
$X$ 中支配 $Y$ 的某个整闭子概形。% P10-E0095
特别地，不妨假设 $d=\dim_\delta(Y)<\infty$。

\medskip\noindent
因此可以对 $d=\dim_\delta(Y)$ 作归纳。若 $d<k$，则
$\CH_{k+r}(X)=0$，引理成立；这是归纳的起始情形。
取非空开集 $V\subset Y$。假设可以证明，对某个
$\beta\in Z_k(V)$ 有 $\alpha|_{f^{-1}(V)}=f^*\beta$。
由引理 \ref{spaces-chow-lemma-exact-sequence-open}，存在
$\beta'\in Z_k(Y)$ 使得 $\beta=\beta'|_V$。
由引理 \ref{spaces-chow-lemma-restrict-to-open} 的正合列
$\CH_k(f^{-1}(Y\setminus V))\to\CH_k(X)\to\CH_k(f^{-1}(V))$，% P10-E0096
$\alpha-f^*\beta'$ 来自某个闭链
$\alpha'\in\CH_{k+r}(f^{-1}(Y\setminus V))$。
由于 $\dim_\delta(Y\setminus V)<d$，对 $d$ 使用归纳假设即得结论。

\medskip\noindent
特别地，以适当开子空间替换 $Y$ 后，可以假设 $Y$ 是一般点为
$\eta$ 的概形。同构 $Y_\eta\cong\mathbf{A}^r_\eta$ 可延拓为
$Y$ 中某个非空开集 $V$ 上的同构；见《空间的极限》引理
\ref{spaces-limits-lemma-descend-finite-presentation}。% P10-E0097
这把我们约化到概形情形，即《Chow 同调》引理
\ref{chow-lemma-pullback-affine-fibres-surjective}。
\end{proof}

\begin{lemma}
\label{spaces-chow-lemma-linebundle}
在情形 \ref{spaces-chow-situation-setup} 中，设 $X/B$ 是良好的。
设 $\mathcal{L}$ 是可逆 $\mathcal{O}_X$-模。令
$$
p :
L = \underline{\Spec}(\text{Sym}^*(\mathcal{L}))
\longrightarrow
X
$$
为 $X$ 上的相伴向量丛。则对所有 $k$，
$p^*:\CH_k(X)\to\CH_{k+1}(L)$ 都是同构。
\end{lemma}

\begin{proof}
满射性见引理 \ref{spaces-chow-lemma-pullback-affine-fibres-surjective}。
令 $o:X\to L$ 为 $L\to X$ 的零截面，即对应于满射
$\text{Sym}^*(\mathcal{L})\to\mathcal{O}_X$ 的态射；该满射把所有
$n>0$ 的 $\mathcal{L}^{\otimes n}$ 都映为零。于是
$p\circ o=\text{id}_X$，且 $o(X)$ 是 $L$ 上的有效 Cartier 因子。
因此由引理 \ref{spaces-chow-lemma-relative-effective-cartier}，
$o^*\circ p^*=\text{id}$，从而 $p^*$ 也是单射。
\end{proof}























\section{双变相交理论}
\label{spaces-chow-section-bivariant}

\noindent
本节对应于《Chow 同调》第 \ref{chow-section-bivariant} 节。
为了恰当地讨论向量丛的高阶陈类，我们依照 \cite{FM} 引入下述概念。
由 \cite[Theorem 17.1]{F} 可知，除去我们所处背景不同这一点外，
这里的定义与 \cite{F} 中的定义一致。

\begin{definition}
\label{spaces-chow-definition-bivariant-class}
\begin{reference}
类似于 \cite[Definition 17.1]{F}
\end{reference}
在情形 \ref{spaces-chow-situation-setup} 中，设 $f:X\to Y$ 是 $B$ 上良好代数空间
之间的态射，并设 $p\in\mathbf{Z}$。一个{\it 关于 $f$ 的 $p$ 次
双变类 $c$}，是这样一条规则：对 $B$ 上良好代数空间的每个态射
$Y'\to Y$ 及每个 $k$，指定一个映射
$$
c \cap - : \CH_k(Y') \longrightarrow \CH_{k - p}(X')
$$
其中 $X'=Y'\times_YX$，并满足下列条件：
\begin{enumerate}
\item 若 $Y''\to Y'$ 是固有态射，则
$c \cap (Y'' \to Y')_*\alpha'' = (X'' \to X')_*(c \cap \alpha'')$
对 $Y''$ 上所有 $\alpha''$ 都成立；
\item 若 $Y''\to Y'$ 是 $B$ 上良好代数空间之间相对维数为 $r$ 的
平坦态射，% P10-E0098
则
$c \cap (Y'' \to Y')^*\alpha' = (X'' \to X')^*(c \cap \alpha')$
对 $Y'$ 上所有 $\alpha'$ 都成立；
\item 若 $(\mathcal{L}',s',i':D'\to Y')$ 如定义
\ref{spaces-chow-definition-gysin-homomorphism}，而其到 $X'$ 的拉回为
$(\mathcal{N}',t',j':E'\to X')$，则
$c\cap(i')^*\alpha'=(j')^*(c\cap\alpha')$
对 $Y'$ 上所有 $\alpha'$ 都成立。
\end{enumerate}
关于 $f$ 的所有 $p$ 次双变类的集合记作 $A^p(X\to Y)$。
\end{definition}

\noindent
在情形 \ref{spaces-chow-situation-setup} 中，设 $X\to Y$ 和 $Y\to Z$ 是 $B$ 上
良好代数空间之间的态射，并设 $p\in\mathbf{Z}$。显然，
$A^p(X\to Y)$ 是阿贝尔群。此外，显然存在双线性复合
$$
A^p(X \to Y) \times A^q(Y \to Z) \to A^{p + q}(X \to Z)
$$
且它满足结合律。我们最关注 $A^p(X)=A^p(X\to X)$；它始终表示关于
$\text{id}_X$ 的双变上同调类。陈类正是位于这里。

\begin{definition}
\label{spaces-chow-definition-chow-cohomology}
在情形 \ref{spaces-chow-situation-setup} 中，设 $X/B$ 是良好的。$X$ 的
{\it Chow 上同调}是分次 $\mathbf{Z}$-代数 $A^*(X)$，其 $p$ 次分量为
$A^p(X\to X)$。
\end{definition}

\noindent
警告：尚不清楚 $A^*(X)$ 的 $\mathbf{Z}$-代数结构是否交换，
但我们将看到陈类位于它的中心。

\begin{remark}
\label{spaces-chow-remark-pullback-cohomology}
在情形 \ref{spaces-chow-situation-setup} 中，设 $f:X\to Y$ 是 $B$ 上良好代数空间
之间的态射。则存在典范的 $\mathbf{Z}$-代数映射
$A^*(Y)\to A^*(X)$。具体地，给定 $c\in A^p(Y)$ 与 $X'\to X$，
把 $X'$ 视为 $Y$ 上的代数空间，由映射
$c\cap-:\CH_k(X')\to\CH_{k-p}(X')$ 定义 $f^*c$。
\end{remark}

\begin{lemma}
\label{spaces-chow-lemma-cap-c1-bivariant}
在情形 \ref{spaces-chow-situation-setup} 中，设 $X/B$ 是良好的。
设 $\mathcal{L}$ 是可逆 $\mathcal{O}_X$-模。对 $f:X'\to X$ 指定
$c_1(f^*\mathcal{L})\cap-:\CH_k(X')\to\CH_{k-1}(X')$
的规则是 $1$ 次双变类。
\end{lemma}

\begin{proof}
这由引理 \ref{spaces-chow-lemma-factors}、
\ref{spaces-chow-lemma-pushforward-cap-c1}、
\ref{spaces-chow-lemma-flat-pullback-cap-c1} 以及
\ref{spaces-chow-lemma-gysin-commutes-cap-c1} 得到。
\end{proof}

\begin{lemma}
\label{spaces-chow-lemma-flat-pullback-bivariant}
在情形 \ref{spaces-chow-situation-setup} 中，设 $f:X\to Y$ 是 $B$ 上良好代数空间
之间相对维数为 $r$ 的平坦态射。对 $Y'\to Y$ 指定
$(f')^*:\CH_k(Y')\to\CH_{k+r}(X')$ 的规则是 $-r$ 次双变类，
其中 $X'=X\times_YY'$。
\end{lemma}

\begin{proof}
这由引理 \ref{spaces-chow-lemma-flat-pullback-rational-equivalence}、
\ref{spaces-chow-lemma-compose-flat-pullback}、
\ref{spaces-chow-lemma-flat-pullback-proper-pushforward} 以及
\ref{spaces-chow-lemma-gysin-flat-pullback} 得到。
\end{proof}

\begin{lemma}
\label{spaces-chow-lemma-gysin-bivariant}
在情形 \ref{spaces-chow-situation-setup} 中，设 $X/B$ 是良好的。设
$(\mathcal{L},s,i:D\to X)$ 是定义
\ref{spaces-chow-definition-gysin-homomorphism} 中的三元组。对 $f:X'\to X$
指定 $(i')^*:\CH_k(X')\to\CH_{k-1}(D')$ 的规则是 $1$ 次双变类，
其中 $D'=D\times_XX'$。
\end{lemma}

\begin{proof}
这由引理 \ref{spaces-chow-lemma-gysin-factors}、
\ref{spaces-chow-lemma-closed-in-X-gysin}、
\ref{spaces-chow-lemma-gysin-flat-pullback} 以及
\ref{spaces-chow-lemma-gysin-commutes-gysin} 得到。
\end{proof}

\begin{lemma}
\label{spaces-chow-lemma-push-proper-bivariant}
在情形 \ref{spaces-chow-situation-setup} 中，设 $f:X\to Y$ 和 $g:Y\to Z$ 是
$B$ 上良好代数空间之间的态射。设 $c\in A^p(X\to Z)$，并假设
$f$ 固有。对 $X'\to X$ 指定% P10-E0099
$\alpha\longmapsto f_*(c\cap\alpha)$ 的规则是 $p$ 次双变类。
\end{lemma}

\begin{proof}
这由引理 \ref{spaces-chow-lemma-compose-pushforward}、
\ref{spaces-chow-lemma-flat-pullback-proper-pushforward} 以及
\ref{spaces-chow-lemma-closed-in-X-gysin} 得到。
\end{proof}

\noindent
下面说明 $c_1(\mathcal{L})$ 位于 $A^*(X)$ 的中心。

\begin{lemma}
\label{spaces-chow-lemma-c1-center}
在情形 \ref{spaces-chow-situation-setup} 中，设 $X/B$ 是良好的。
设 $\mathcal{L}$ 是可逆 $\mathcal{O}_X$-模。则
$c_1(\mathcal{L})\in A^1(X)$ 与每个元素 $c\in A^p(X)$ 交换。
\end{lemma}

\begin{proof}
令 $p:L\to X$ 如引理 \ref{spaces-chow-lemma-linebundle}，并令 $o:X\to L$ 为零截面。
注意，$p^*\mathcal{L}^{\otimes-1}$ 有一个典范截面，其零化轨迹恰为
有效 Cartier 因子 $o(X)$。设 $\alpha\in\CH_k(X)$。则
$$
p^*(c_1(\mathcal{L}^{\otimes -1}) \cap \alpha) =
c_1(p^*\mathcal{L}^{\otimes -1}) \cap p^*\alpha =
o_* o^* p^*\alpha
$$
这是由引理 \ref{spaces-chow-lemma-flat-pullback-cap-c1} 和
\ref{spaces-chow-lemma-relative-effective-cartier} 得到的。由于 $c$ 是双变类，故
\begin{align*}
p^*(c \cap c_1(\mathcal{L}^{\otimes -1}) \cap \alpha)
& =
c \cap p^*(c_1(\mathcal{L}^{\otimes -1}) \cap \alpha) \\
& =
c \cap o_* o^* p^*\alpha \\
& =
o_* o^* p^*(c \cap \alpha) \\
& =
p^*(c_1(\mathcal{L}^{\otimes -1}) \cap c \cap \alpha)
\end{align*}
（最后一个等号来自把上式用于 $c\cap\alpha$）。
由上文所引引理，$p^*$ 是单射，故
$c_1(\mathcal{L}^{\otimes-1})$ 位于 $A^*(X)$ 的中心。引理得证。
\end{proof}

\noindent
下面给出双变类为零的一个判据。% P10-E0100

\begin{lemma}
\label{spaces-chow-lemma-bivariant-zero}
在情形 \ref{spaces-chow-situation-setup} 中，设 $X/B$ 是良好的，并设
$c\in A^p(X)$。则 $c=0$，当且仅当对每个在 $X$ 上局部有限型的
整代数空间 $Y$，都有 $c\cap[Y]=0$ 于 $\CH_*(Y)$ 中。
\end{lemma}

\begin{proof}
一个方向是显然的。反之，假设对每个在 $X$ 上局部有限型的整代数空间
$Y$，都有 $c\cap[Y]=0$ 于 $\CH_*(Y)$ 中。
设 $X'\to X$ 局部有限型，并设 $\alpha\in\CH_k(X')$。
写成 $\alpha=\sum n_i[Y_i]$，其中 $Y_i\subset X'$ 是一族局部有限的、
$\delta$-维数为 $k$ 的整闭子概形。% P10-E0101
于是，$\alpha$ 是 $X''=\coprod Y_i$ 上闭链
$\alpha'=\sum n_i[Y_i]$ 沿固有态射 $X''\to X'$ 的推出。
由双变类的性质，只需证明
$c\cap\alpha'=0$ 于 $\CH_{k-p}(X'')$ 中。
由定义立即有
$\CH_{k-p}(X'')=\prod\CH_{k-p}(Y_i)$。投影映射
$\CH_{k-p}(X'')\to\CH_{k-p}(Y_i)$ 由平坦拉回给出。
由于与 $c$ 相交同平坦拉回交换，只需证明
$c\cap[Y_i]$ 在 $\CH_{k-p}(Y_i)$ 中为零，而这正是我们的假设。
\end{proof}















\section{射影空间丛公式}
\label{spaces-chow-section-projective-space-bundle-formula}

\noindent
在情形 \ref{spaces-chow-situation-setup} 中，设 $X/B$ 是良好的。
取秩为 $r$ 的有限局部自由 $\mathcal{O}_X$-模 $\mathcal{E}$。
我们的约定是，{\it 与 $\mathcal{E}$ 相伴的射影丛}指态射
$$
\xymatrix{
\mathbf{P}(\mathcal{E}) =
\underline{\text{Proj}}_X(\text{Sym}^*(\mathcal{E}))
\ar[r]^-\pi
& X
}
$$
并把 $\mathcal{O}_{\mathbf{P}(\mathcal{E})}(1)$ 规范化为满足
$\pi_*(\mathcal{O}_{\mathbf{P}(\mathcal{E})}(1))=\mathcal{E}$。
特别地，存在满射
$\pi^*\mathcal{E}\to\mathcal{O}_{\mathbf{P}(\mathcal{E})}(1)$。
我们将非正式地用“设 $(\pi:P\to X,\mathcal{O}_P(1))$ 是与
$\mathcal{E}$ 相伴的射影丛”表示如下情形：
$P=\mathbf{P}(\mathcal{E})$ 且
$\mathcal{O}_P(1)=\mathcal{O}_{\mathbf{P}(\mathcal{E})}(1)$。

\begin{lemma}
\label{spaces-chow-lemma-cap-projective-bundle}
在情形 \ref{spaces-chow-situation-setup} 中，设 $X/B$ 是良好的。
设 $\mathcal{E}$ 是秩为 $r$ 的有限局部自由 $\mathcal{O}_X$-模
$\mathcal{E}$。% P10-E0102
设 $(\pi:P\to X,\mathcal{O}_P(1))$ 是与 $\mathcal{E}$ 相伴的射影丛。
对任意 $\alpha\in\CH_k(X)$，元素
$$
\pi_*\left(
c_1(\mathcal{O}_P(1))^s \cap \pi^*\alpha
\right)
\in
\CH_{k + r - 1 - s}(X)
$$
在 $s<r-1$ 时为 $0$，在 $s=r-1$ 时等于 $\alpha$。
\end{lemma}

\begin{proof}
设 $Z\subset X$ 是 $\delta$-维数为 $k$ 的整闭子空间。
我们对 $\alpha=[Z]$ 证明引理。省略由这一特殊情形推出一般情形的论证；
提示：仿照注 \ref{spaces-chow-remark-infinite-sums-rational-equivalences} 论证。

\medskip\noindent
令 $P_Z=P\times_XZ$ 为基变换。显然，
$\pi_Z:P_Z\to Z$ 是与 $\mathcal{E}|_Z$ 相伴的射影丛，而
$\mathcal{O}_P(1)$ 拉回为 $P_Z$ 上相应的可逆模。
由引理 \ref{spaces-chow-lemma-cap-c1-bivariant} 和
\ref{spaces-chow-lemma-flat-pullback-bivariant}，
$c_1(\mathcal{O}_P(1)\cap-$ 与 $\pi^*$ 都是双变类，% P10-E0103
故
$$
\pi_*\left(
c_1(\mathcal{O}_P(1))^s \cap \pi^*[Z]
\right)
=
(Z \to X)_*\pi_{Z, *}\left(
c_1(\mathcal{O}_{P_Z}(1))^s \cap \pi_Z^*[Z]
\right)
$$
因此只需在 $X$ 整且 $\alpha=[X]$ 时证明引理。

\medskip\noindent
假设 $X$ 整、$\dim_\delta(X)=k$ 且 $\alpha=[X]$。
注意，$\pi^*[X]=[P]$，因为 $P$ 整且 $\delta$-维数为 $r-1$。% P10-E0104
若 $s<r-1$，则依构造，
$c_1(\mathcal{O}_P(1))^s\cap[P]$ 一个
$(k+r-1-s)$-闭链。% P10-E0105
因此，由维数原因，该闭链的推出为零。

\medskip\noindent
令 $s=r-1$。由上述论证，对某个 $n\in\mathbf{Z}$ 有
$\pi_*(c_1(\mathcal{O}_P(1))^s\cap[P])=n[X]$。我们要证明 $n=1$。
由与上面相同的维数理由，只需在以某个稠密开子空间替换 $X$ 后证明。
因此可以假设 $X$ 是概形，此时结论由《Chow 同调》引理
\ref{chow-lemma-cap-projective-bundle} 得到。
\end{proof}

\begin{lemma}[射影空间丛公式]
\label{spaces-chow-lemma-chow-ring-projective-bundle}
设 $(S,\delta)$ 如情形 \ref{spaces-chow-situation-setup}，并设 $X$ 在 $S$ 上
局部有限型。设 $\mathcal{E}$ 是秩为 $r$ 的有限局部自由
$\mathcal{O}_X$-模 $\mathcal{E}$。% P10-E0106
设 $(\pi:P\to X,\mathcal{O}_P(1))$ 是与 $\mathcal{E}$ 相伴的射影丛。
映射
$$
\bigoplus\nolimits_{i = 0}^{r - 1}
\CH_{k + i}(X)
\longrightarrow
\CH_{k + r - 1}(P),
$$
$$
(\alpha_0, \ldots, \alpha_{r-1})
\longmapsto
\pi^*\alpha_0 +
c_1(\mathcal{O}_P(1)) \cap \pi^*\alpha_1
+ \ldots +
c_1(\mathcal{O}_P(1))^{r - 1} \cap \pi^*\alpha_{r-1}
$$
是同构。
\end{lemma}

\begin{proof}
固定 $k\in\mathbf{Z}$。先证明该映射是单射。
假设左端的元素 $(\alpha_0,\ldots,\alpha_{r-1})$ 映为零。
由引理 \ref{spaces-chow-lemma-cap-projective-bundle}，
$$
0 = \pi_*(\pi^*\alpha_0 +
c_1(\mathcal{O}_P(1)) \cap \pi^*\alpha_1
+ \ldots +
c_1(\mathcal{O}_P(1))^{r - 1} \cap \pi^*\alpha_{r-1})
= \alpha_{r - 1}
$$
其次，
$$
0 = \pi_*(c_1(\mathcal{O}_P(1)) \cap (\pi^*\alpha_0 +
c_1(\mathcal{O}_P(1)) \cap \pi^*\alpha_1
+ \ldots +
c_1(\mathcal{O}_P(1))^{r - 2} \cap \pi^*\alpha_{r - 2}))
= \alpha_{r - 2}
$$
依此类推。因此该映射是单射。

\medskip\noindent
为证明该映射是满射，我们将完全仿照引理
\ref{spaces-chow-lemma-pullback-affine-fibres-surjective} 的证明，约化到概形情形。
我们建议读者跳过这个证明。

\medskip\noindent
设 $\beta\in\CH_{k+r-1}(P)$。写成 $\beta=\sum m_j[W_j]$，其中
$m_j\not =0$，且 $W_j$ 是两两不同的、$\delta$-维数为 $k+r$ 的
整闭子空间。% P10-E0107
于是族 $\{W_j\}$ 在 $P$ 中局部有限。令 $Z_j\subset X$ 为引理
\ref{spaces-chow-lemma-proper-image} 意义下 $W_j$ 的“像”。对任意拟紧开集
$U\subset X$，只有有限多个 $j$ 使 $\pi^{-1}(U)\cap W_j$ 非空。
因此，像 $Z_j$ 构成 $X$ 中一族局部有限的整闭子空间。

\medskip\noindent
考虑纤维积图
$$
\xymatrix{
P_j \ar[r] \ar[d]_{\pi_j} & P \ar[d]^\pi \\
Z_j \ar[r] & X
}
$$
假设 $[W_j]\in Z_{k+r-1}(P_j)$ 有理等价于
$$
\pi_j^*\alpha_{j, 0} +
c_1(\mathcal{O}(1)) \cap \pi_j^*\alpha_{j, 1} +
\ldots +
c_1(\mathcal{O}(1))^{r - 1} \cap \pi_j^*\alpha_{j, r - 1}
$$
其中 $\alpha_{j,i}\in\CH_{k+i}(Z_j)$ 是某个 $(k+i)$-闭链。
那么 $\alpha_i=\sum m_j\beta_{j,i}$ 是 $X$ 上的 $(k+i)$-闭链，% P10-E0108
并且
$$
\pi^*\alpha_0 +
c_1(\mathcal{O}(1)) \cap \pi^*\alpha_1 +
\ldots +
c_1(\mathcal{O}(1))^{r - 1} \cap \pi^*\alpha_{r - 1}
$$
与 $\beta$ 有理等价（见注
\ref{spaces-chow-remark-infinite-sums-rational-equivalences}）。
这把我们约化到 $X$ 整且 $\alpha=[W]$ 的情形，% P10-E0109
其中 $W$ 是 $P$ 中支配 $X$ 的某个整闭子概形。% P10-E0110
特别地，不妨假设 $d=\dim_\delta(X)<\infty$。

\medskip\noindent
因此可以对 $d=\dim_\delta(X)$ 作归纳。若 $d<k$，则
$\CH_{k+r-1}(X)=0$，引理成立；这是归纳的起始情形。% P10-E0111
取非空开集 $U\subset X$。假设可以证明
$$
\beta|_{\pi^{-1}(U)} =
\pi^*\alpha_0 +
c_1(\mathcal{O}(1)) \cap \pi^*\alpha_1 +
\ldots +
c_1(\mathcal{O}(1))^{r - 1} \cap \pi^*\alpha_{r - 1}
$$
其中 $\alpha_i\in Z_{k+i}(U)$。由引理
\ref{spaces-chow-lemma-exact-sequence-open}，存在
$\alpha'_i\in Z_{k+i}(X)$ 使得 $\alpha_i=\alpha'_i|_U$。
由引理 \ref{spaces-chow-lemma-restrict-to-open} 的正合列
% P10-E0112
$\CH_{k + i}(\pi^{-1}(X \setminus U)) \to \CH_{k + i}(P) \to
\CH_{k + i}(\pi^{-1}(U))$
可知
$$
\beta -
\left(\pi^*\alpha'_0 +
c_1(\mathcal{O}(1)) \cap \pi^*\alpha'_1 +
\ldots +
c_1(\mathcal{O}(1))^{r - 1} \cap \pi^*\alpha'_{r - 1}\right)
$$
来自某个闭链
$\beta'\in\CH_{k+r}(\pi^{-1}(X\setminus U))$。% P10-E0113
由于 $\dim_\delta(X\setminus U)<d$，对 $d$ 使用归纳假设即得结论。

\medskip\noindent
特别地，以适当开子空间替换 $X$ 后，可以假设 $X$ 是概形，
从而已把问题约化到《Chow 同调》引理
\ref{chow-lemma-chow-ring-projective-bundle}。
\end{proof}

\begin{lemma}
\label{spaces-chow-lemma-vectorbundle}
在情形 \ref{spaces-chow-situation-setup} 中，设 $X/B$ 是良好的。
设 $\mathcal{E}$ 是 $X$ 上秩为 $r$ 的有限局部自由层。令
$$
p :
E = \underline{\Spec}(\text{Sym}^*(\mathcal{E}))
\longrightarrow
X
$$
为 $X$ 上的相伴向量丛。则对所有 $k$，
$p^*:\CH_k(X)\to\CH_{k+r}(E)$ 都是同构。
\end{lemma}

\begin{proof}
（线丛情形见引理 \ref{spaces-chow-lemma-linebundle}。）
满射性见引理 \ref{spaces-chow-lemma-pullback-affine-fibres-surjective}。
设 $(\pi:P\to X,\mathcal{O}_P(1))$ 是与有限局部自由层
$\mathcal{E}\oplus\mathcal{O}_X$ 相伴的射影空间丛。
设 $s\in\Gamma(P,\mathcal{O}_P(1))$ 对应于整体截面
$(0,1)\in\Gamma(X,\mathcal{E}\oplus\mathcal{O}_X)$。
令 $D=Z(s)\subset P$。注意，
$(\pi|_D : D \to X , \mathcal{O}_P(1)|_D)$
是与 $\mathcal{E}$ 相伴的射影空间丛。记
$\pi_D=\pi|_D$ 以及 $\mathcal{O}_D(1)=\mathcal{O}_P(1)|_D$。
此外，$D$ 是 $P$ 上的有效 Cartier 因子。因此
$\mathcal{O}_P(D)=\mathcal{O}_P(1)$（见《空间上的因子》引理
\ref{spaces-divisors-lemma-characterize-OD}）。
还有同构 $E\cong P\setminus D$。以 $j:E\to P$ 表示相应的开浸入。
为证明单射性，我们使用如下事实：映射
$$
j^* :
\CH_{k + r}(P)
\longrightarrow
\CH_{k + r}(E)
$$
的核是支集于有效 Cartier 因子 $D$ 的闭链；见引理
\ref{spaces-chow-lemma-restrict-to-open}。所以若 $p^*\alpha=0$，则对某个
$\beta\in\CH_{k+r}(D)$ 有
$\pi^*\alpha=i_*\beta$。% P10-E0114
由引理 \ref{spaces-chow-lemma-chow-ring-projective-bundle}，可写成
$$
\beta = \pi_D^*\beta_0 +
\ldots + c_1(\mathcal{O}_D(1))^{r - 1} \cap \pi_D^* \beta_{r - 1}.
$$
其中 $\beta_i\in\CH_{k+i}(X)$。% P10-E0115
由引理 \ref{spaces-chow-lemma-relative-effective-cartier} 和
\ref{spaces-chow-lemma-pushforward-cap-c1}，这推出
$$
\pi^*\alpha = i_*\beta =
c_1(\mathcal{O}_P(1)) \cap \pi^*\beta_0 +
\ldots +
c_1(\mathcal{O}_D(1))^r \cap \pi^*\beta_{r - 1}.
$$
% P10-E0116
由于 $\mathcal{E}\oplus\mathcal{O}_X$ 的秩为 $r+1$，除非所有
$\alpha$ 和所有 $\beta_i$ 都为零，否则这与引理
\ref{spaces-chow-lemma-pushforward-cap-c1} 矛盾。% P10-E0117
\end{proof}








\section{向量丛的陈类}
\label{spaces-chow-section-chern-classes-vector-bundles}

\noindent
本节对应于《Chow 同调》第
\ref{chow-section-chern-classes-vector-bundles} 节和第
\ref{chow-section-intersecting-chern-classes} 节。
不过，与那里的做法不同，我们直接把向量丛的陈类定义为双变类。
这样可以省去大量工作。

\begin{lemma}
\label{spaces-chow-lemma-segre-classes}
在情形 \ref{spaces-chow-situation-setup} 中，设 $X/B$ 是良好的。
设 $\mathcal{E}$ 是 $X$ 上秩为 $r$ 的有限局部自由层，并设
$(\pi:P\to X,\mathcal{O}_P(1))$ 是与 $\mathcal{E}$ 相伴的射影空间丛。
对 $B$ 上良好代数空间的每个态射 $X'\to X$，存在唯一的映射
$$
c_i(\mathcal{E}) \cap - : \CH_k(X') \longrightarrow \CH_{k - i}(X'),\quad
i = 0, \ldots, r
$$
使得对 $\alpha\in\CH_k(X')$ 有
$c_0(\mathcal{E})\cap\alpha=\alpha$，并且
% P10-E0118
$$
\sum\nolimits_{i = 0, \ldots, r}
(-1)^i c_1(\mathcal{O}_{P'}(1))^i \cap
(\pi')^*\left(c_{r - i}(\mathcal{E}) \cap \alpha\right) = 0
$$
其中 $\pi':P'\to X'$ 是 $\pi$ 的基变换。此外，这些映射定义了
$X$ 上的 $i$ 次双变类 $c_i(\mathcal{E})$。
\end{lemma}

\begin{proof}
映射 $c_i(\mathcal{E})\cap-$ 的唯一性与存在性，立即由引理
\ref{spaces-chow-lemma-chow-ring-projective-bundle} 以及给定的
$c_0(\mathcal{E})$ 描述得到。% P10-E0122
对每个 $i\in\mathbf{Z}$，给 $B$ 上良好代数空间的每个态射
$X'\to X$ 指定映射
% P10-E0119
$$
t_i(\mathcal{E}) \cap - :
\CH_k(X') \longrightarrow \CH_{k - i}(X'),\quad
\alpha \longmapsto
\pi'_*(c_1(\mathcal{O}_{P'}(1))^{r - 1 + i} \cap (\pi')^*\alpha)
$$
的规则是双变类\footnote{不计符号时，它们就是 $\mathcal{E}$ 的
Segre 类。}；这是由引理 \ref{spaces-chow-lemma-cap-c1-bivariant}、
\ref{spaces-chow-lemma-flat-pullback-bivariant} 以及
\ref{spaces-chow-lemma-push-proper-bivariant} 得到的。
由引理 \ref{spaces-chow-lemma-cap-projective-bundle}，当 $i<0$ 时有
$t_i(\mathcal{E})=0$，且 $t_0(\mathcal{E})=1$。
对引理陈述中的等式施行推出，再由引理
\ref{spaces-chow-lemma-cap-projective-bundle} 得到
$$
(-1)^r t_1(\mathcal{E}) + (-1)^{r - 1}c_1(\mathcal{E}) = 0
$$
特别地，$c_1(\mathcal{E})$ 是双变类。若把引理陈述中的等式乘以
$c_1(\mathcal{O}_{P'}(1))$，再把结果推出到 $X'$，便得到
$$
(-1)^r t_2(\mathcal{E}) +
(-1)^{r - 1} t_1(\mathcal{E}) \cap c_1(\mathcal{E}) +
(-1)^{r - 2} c_2(\mathcal{E}) = 0
$$
与前面一样，由此可知 $c_2(\mathcal{E})$ 是双变类。依此类推。
\end{proof}

\begin{definition}
\label{spaces-chow-definition-chern-classes}
在情形 \ref{spaces-chow-situation-setup} 中，设 $X/B$ 是良好的。
设 $\mathcal{E}$ 是 $X$ 上秩为 $r$ 的有限局部自由层。
对 $i=0,\ldots,r$，$\mathcal{E}$ 的{\it 第 $i$ 陈类}是引理
\ref{spaces-chow-lemma-segre-classes} 中构造的 $i$ 次双变类
$c_i(\mathcal{E})\in A^i(X)$。$\mathcal{E}$ 的{\it 总陈类}是形式和
$$
c(\mathcal{E}) =
c_0(\mathcal{E}) + c_1(\mathcal{E}) + \ldots + c_r(\mathcal{E})
$$
并把它视为 $X$ 上的非齐次双变类。
\end{definition}

\noindent
为方便起见，当 $i>r$ 或 $i<0$ 时，我们常令
$c_i(\mathcal{E})=0$。依定义，
$c_0(\mathcal{E})=1\in A^0(X)$。下面作一个相容性检验。

\begin{lemma}
\label{spaces-chow-lemma-first-chern-class}
在情形 \ref{spaces-chow-situation-setup} 中，设 $X/B$ 是良好的。
设 $\mathcal{L}$ 是可逆 $\mathcal{O}_X$-模。定义
\ref{spaces-chow-definition-chern-classes} 中 $\mathcal{L}$ 在 $X$ 上的第一陈类，
等于引理 \ref{spaces-chow-lemma-cap-c1-bivariant} 中的双变类。
\end{lemma}

\begin{proof}
事实上，在此情形下，由我们对射影丛所作的规范化，
$P=\mathbf{P}(\mathcal{L})=X$ 且
$\mathcal{O}_P(1)=\mathcal{L}$；见第
\ref{spaces-chow-section-projective-space-bundle-formula} 节。
因此，引理 \ref{spaces-chow-lemma-segre-classes} 中的等式变为
$$
(-1)^0 c_1(\mathcal{L})^0 \cap c^{new}_1(\mathcal{L}) \cap \alpha +
(-1)^1 c_1(\mathcal{L})^1 \cap c^{new}_0(\mathcal{L}) \cap \alpha = 0
$$
其中 $c_i^{new}(\mathcal{L})$ 如定义
\ref{spaces-chow-definition-chern-classes}。由于
$c_0^{new}(\mathcal{L})=1$ 且 $c_1(\mathcal{L})^0=1$，结论成立。
\end{proof}

\noindent
下面说明陈类位于双变 Chow 上同调环 $A^*(X)$ 的中心。

\begin{lemma}
\label{spaces-chow-lemma-cap-commutative-chern}
在情形 \ref{spaces-chow-situation-setup} 中，设 $X/B$ 是良好的。
设 $\mathcal{E}$ 是秩为 $r$ 的局部自由 $\mathcal{O}_X$-模。
则 $c_j(\mathcal{L})\in A^j(X)$ 与每个元素 $c\in A^p(X)$ 交换。% P10-E0120
特别地，若 $\mathcal{F}$ 是 $X$ 上另一个秩为 $s$ 的局部自由
$\mathcal{O}_X$-模，则对所有 $\alpha\in\CH_k(X)$，
$$
c_i(\mathcal{E}) \cap c_j(\mathcal{F}) \cap \alpha
=
c_j(\mathcal{F}) \cap c_i(\mathcal{E}) \cap \alpha
$$
在 $\CH_{k-i-j}(X)$ 中成立。
\end{lemma}

\begin{proof}
设 $X'\to X$ 是 $B$ 上良好代数空间之间的态射，并设
$\alpha\in\CH_k(X')$。写成
$\alpha_j=c_j(\mathcal{E})\cap\alpha$，于是 $\alpha_0=\alpha$。
由引理 \ref{spaces-chow-lemma-segre-classes}，
$$
\sum\nolimits_{i = 0}^r
(-1)^i c_1(\mathcal{O}_{P'}(1))^i \cap
(\pi')^*(\alpha_{r - i}) = 0
$$
在与 $(X'\to X)^*\mathcal{E}$ 相伴的射影丛
$(\pi':P'\to X',\mathcal{O}_{P'}(1))$ 的周群中成立。
施行 $c\cap-$，再利用引理 \ref{spaces-chow-lemma-c1-center} 与双变类的性质，得到
% P10-E0121
$$
\sum\nolimits_{i = 0}^r
(-1)^i c_1(\mathcal{O}_{P'}(1))^i \cap
\pi^*(c \cap \alpha_{r - i}) = 0
$$
在 $P'$ 的周群中成立。因此，由引理 \ref{spaces-chow-lemma-segre-classes} 中的
唯一性，$c\cap\alpha_j$ 等于
$c_j(\mathcal{E})\cap(c\cap\alpha)$。引理得证。
\end{proof}

\begin{remark}
\label{spaces-chow-remark-extend-to-finite-locally-free}
在情形 \ref{spaces-chow-situation-setup} 中，设 $X/B$ 是良好的。
设 $\mathcal{E}$ 是有限局部自由 $\mathcal{O}_X$-模。
若 $\mathcal{E}$ 的秩不是常数，仍可定义 $\mathcal{E}$ 的陈类。
具体地，此时可以写成
$$
X = X_0 \amalg X_1 \amalg X_2 \amalg \ldots
$$
其中 $X_r\subset X$ 是 $\mathcal{E}$ 的秩为 $r$ 的开闭子空间。
若 $X'\to X$ 是 $B$ 上良好代数空间之间的态射，则拉回给出
$X'$ 的相应分解，并且由我们的定义有
$$
\CH_*(X') = \prod\nolimits_{r \geq 0} \CH_*(X'_r)
$$
于是把 $c_i(\mathcal{E})$ 定义为保持这些直积分解的双变类，
并使它在各因子上的作用为已经定义的运算
$c_i(\mathcal{E}|_{X_r})\cap-$。注意，在这一背景下，可能有无穷多个
$i$ 使 $c_i(\mathcal{E})$ 非零。
\end{remark}












\section{陈类之间的多项式关系}
\label{spaces-chow-section-relations-chern-classes}

\noindent
在情形 \ref{spaces-chow-situation-setup} 中，设 $X/B$ 是良好的。
设 $\mathcal{E}_i$ 是有限多个有限局部自由 $\mathcal{O}_X$-模。
由引理 \ref{spaces-chow-lemma-cap-commutative-chern}，陈类
$$
c_j(\mathcal{E}_i) \in A^*(X)
$$
生成 Chow 上同调 $A^*(X)$ 的一个交换（甚至是中心的）
$\mathbf{Z}$-子代数。因此，可以谈论这些陈类的一个多项式为零，
或两个多项式相等的含义。例如，等式
$c_1(\mathcal{E}_1)^5 + c_2(\mathcal{E}_2)c_3(\mathcal{E}_3) = 0$
表示对 $B$ 上良好代数空间的所有态射 $f:Y\to X$，运算
$$
\CH_k(Y) \longrightarrow \CH_{k - 5}(Y), \quad
\alpha \longmapsto
c_1(\mathcal{E}_1)^5 \cap \alpha +
c_2(\mathcal{E}_2) \cap c_3(\mathcal{E}_3) \cap \alpha
$$
均为零。由引理 \ref{spaces-chow-lemma-bivariant-zero}，这等价于要求：
给定任意态射 $f:Y\to X$，其中 $Y$ 是在 $X$ 上局部有限型的
整代数空间，闭链
$$
c_1(\mathcal{E}_1)^5 \cap [Y] +
c_2(\mathcal{E}_2) \cap c_3(\mathcal{E}_3) \cap [Y]
$$
在 $\CH_{\dim(Y)-5}(Y)$ 中为零。% P10-E0123

\medskip\noindent
一个具体例子是关系
$$
c_1(\mathcal{L} \otimes_{\mathcal{O}_X} \mathcal{N})
=
c_1(\mathcal{L}) + c_1(\mathcal{N})
$$
它已在引理 \ref{spaces-chow-lemma-c1-cap-additive} 中得到证明。
更一般地，任意局部自由层与一个可逆层作张量积时有如下结果。

\begin{lemma}
\label{spaces-chow-lemma-chern-classes-E-tensor-L}
在情形 \ref{spaces-chow-situation-setup} 中，设 $X/B$ 是良好的。
设 $\mathcal{E}$ 是 $X$ 上秩为 $r$ 的有限局部自由层，
$\mathcal{L}$ 是 $X$ 上的可逆层。则在 $A^*(X)$ 中有
\begin{equation}
\label{spaces-chow-equation-twist}
c_i({\mathcal E} \otimes {\mathcal L})
=
\sum\nolimits_{j = 0}^i
\binom{r - i + j}{j} c_{i - j}({\mathcal E}) c_1({\mathcal L})^j
\end{equation}
\end{lemma}

\begin{proof}
证明与《Chow 同调》引理
\ref{chow-lemma-chern-classes-E-tensor-L} 的证明相同，
只需把其中使用的引理换成引理 \ref{spaces-chow-lemma-bivariant-zero} 和
\ref{spaces-chow-lemma-segre-classes}。
\end{proof}












\section{陈类的可加性}
\label{spaces-chow-section-additivity-chern-classes}

\noindent
本节对应于《Chow 同调》第
\ref{chow-section-additivity-chern-classes} 节。

\begin{lemma}
\label{spaces-chow-lemma-get-rid-of-trivial-subbundle}
在情形 \ref{spaces-chow-situation-setup} 中，设 $X/B$ 是良好的。
设 $\mathcal{E}$、$\mathcal{F}$ 分别是 $X$ 上秩为 $r$、$r-1$ 的
有限局部自由层，并且它们嵌入短正合列
$$
0 \to \mathcal{O}_X \to \mathcal{E} \to \mathcal{F} \to 0
$$
则在 $A^*(X)$ 中有
$$
c_r(\mathcal{E}) = 0, \quad
c_j(\mathcal{E}) = c_j(\mathcal{F}), \quad j = 0, \ldots, r - 1
$$
\end{lemma}

\begin{proof}
证明与《Chow 同调》引理
\ref{chow-lemma-get-rid-of-trivial-subbundle} 的证明相同，
只需把其中使用的引理换成引理
\ref{spaces-chow-lemma-bivariant-zero}、
\ref{spaces-chow-lemma-relative-effective-cartier}、
\ref{spaces-chow-lemma-pushforward-cap-c1} 以及
\ref{spaces-chow-lemma-segre-classes}。
\end{proof}

\begin{lemma}
\label{spaces-chow-lemma-additivity-invertible-subsheaf}
在情形 \ref{spaces-chow-situation-setup} 中，设 $X/B$ 是良好的。
设 $\mathcal{E}$、$\mathcal{F}$ 分别是 $X$ 上秩为 $r$、$r-1$ 的
有限局部自由层，并且它们嵌入短正合列
$$
0 \to \mathcal{L} \to \mathcal{E} \to \mathcal{F} \to 0
$$
其中 $\mathcal{L}$ 是可逆层。则在 $A^*(X)$ 中有
$$
c(\mathcal{E}) = c(\mathcal{L}) c(\mathcal{F})
$$
\end{lemma}

\begin{proof}
证明与《Chow 同调》引理
\ref{chow-lemma-additivity-invertible-subsheaf} 的证明相同，
只需把其中使用的引理换成引理
\ref{spaces-chow-lemma-get-rid-of-trivial-subbundle} 和
\ref{spaces-chow-lemma-chern-classes-E-tensor-L}。
\end{proof}

\begin{lemma}
\label{spaces-chow-lemma-additivity-chern-classes}
在情形 \ref{spaces-chow-situation-setup} 中，设 $X/B$ 是良好的。
假设 $\mathcal{E}$ 位于有限局部自由层的正合列
$$
0
\to
\mathcal{E}_1
\to
\mathcal{E}
\to
\mathcal{E}_2
\to
0
$$
中，其中 $\mathcal{E}_i$ 的秩为 $r_i$。则总陈类在 $A^*(X)$ 中满足
$$
c(\mathcal{E}) = c(\mathcal{E}_1) c(\mathcal{E}_2)
$$
\end{lemma}

\begin{proof}
证明与《Chow 同调》引理
\ref{chow-lemma-additivity-chern-classes} 的证明相同，
只需把其中使用的引理换成引理
\ref{spaces-chow-lemma-bivariant-zero}、
\ref{spaces-chow-lemma-additivity-invertible-subsheaf} 以及
\ref{spaces-chow-lemma-segre-classes}。
\end{proof}

\begin{lemma}
\label{spaces-chow-lemma-chern-filter-by-linebundles}
在情形 \ref{spaces-chow-situation-setup} 中，设 $X/B$ 是良好的。
设 ${\mathcal L}_i$（$i=1,\ldots,r$）是可逆
$\mathcal{O}_X$-模。设 $\mathcal{E}$ 是一个局部自由秩
$\mathcal{O}_X$-模，并带有滤过% P10-E0124
$$
0 = \mathcal{E}_0 \subset \mathcal{E}_1 \subset \mathcal{E}_2
\subset \ldots \subset \mathcal{E}_r = \mathcal{E}
$$
使得 $\mathcal{E}_i/\mathcal{E}_{i-1}\cong\mathcal{L}_i$。
令 $c_1({\mathcal L}_i)=x_i$。则在 $A^*(X)$ 中有
$$
c(\mathcal{E})
=
\prod\nolimits_{i = 1}^r (1 + x_i)
$$
\end{lemma}

\begin{proof}
应用引理 \ref{spaces-chow-lemma-additivity-invertible-subsheaf} 并作归纳。
\end{proof}



















\section{分裂原理}
\label{spaces-chow-section-splitting-principle}

\noindent
本节对应于《Chow 同调》第
\ref{chow-section-additivity-chern-classes} 节。

\begin{lemma}
\label{spaces-chow-lemma-splitting-principle}
在情形 \ref{spaces-chow-situation-setup} 中，设 $X/B$ 是良好的。
设 $\mathcal{E}_i$ 是有限多个秩为 $r_i$ 的局部自由
$\mathcal{O}_X$-模。存在相对维数为 $d$ 的射影平坦态射
$\pi:P\to X$，使得
\begin{enumerate}
\item 对 $B$ 上良好代数空间的任意态射 $f:Y\to X$，映射
$\pi_Y^*:\CH_*(Y)\to\CH_{*+d}(Y\times_XP)$ 是单射；
\item 每个 $\pi^*\mathcal{E}_i$ 都有一个滤过，其逐次商
$\mathcal{L}_{i,1},\ldots,\mathcal{L}_{i,r_i}$ 是可逆
${\mathcal O}_P$-模。
\end{enumerate}
\end{lemma}

\begin{proof}
对整数 $r=\sum r_i$ 作归纳。若 $r=0$，可取
$\pi=\text{id}_X$。若所有 $i$ 都有 $r_i=1$，也可取
$\pi=\text{id}_X$。假设对某个 $i_0$ 有 $r_{i_0}>1$。
设 $(\pi:P\to X,\mathcal{O}_P(1))$ 是与
$\mathcal{E}_{i_0}$ 相伴的射影丛。典范映射
$\pi^*\mathcal{E}_{i_0}\to\mathcal{O}_P(1)$ 是满射，故其核
$\mathcal{E}'_{i_0}$ 是秩为 $r_{i_0}-1$ 的有限局部自由层。
注意，对 $B$ 上良好代数空间的任意态射 $f:Y\to X$，
$\pi_Y^*$ 都是单射；见引理
\ref{spaces-chow-lemma-chow-ring-projective-bundle}。因此，只需对 $P$ 以及局部自由层
$\pi^*\mathcal{E}_i$ 证明引理。不过，由于存在具有可逆商的子丛
$\mathcal{E}_{i_0}\subset\pi^*\mathcal{E}_{i_0}$，% P10-E0125
现在只需对族
$\{\mathcal{E}_i\}_{i\not = i_0}\cup\{\mathcal{E}'_{i_0}\}$
证明引理。这样使 $r$ 减少 $1$，归纳假设给出结论。
\end{proof}

\noindent
与其解释分裂原理的含义，不如直接在若干引理的证明中使用它。

\begin{lemma}
\label{spaces-chow-lemma-chern-classes-dual}
在情形 \ref{spaces-chow-situation-setup} 中，设 $X/B$ 是良好的。
设 $\mathcal{E}$ 是有限局部自由 $\mathcal{O}_X$-模，
其对偶为 $\mathcal{E}^\vee$。则在 $A^i(X)$ 中有
$$
c_i(\mathcal{E}^\vee) = (-1)^i c_i(\mathcal{E})
$$
\end{lemma}

\begin{proof}
选取引理 \ref{spaces-chow-lemma-splitting-principle} 中的态射
$\pi:P\to X$。由于 $\pi^*$（在任意基变换后）是单射，
只需在拉回到 $P$ 后证明 $\mathcal{E}$ 与 $\mathcal{E}^\vee$
的陈类之间的关系。因此，可以假设存在可逆
$\mathcal{O}_X$-模 ${\mathcal L}_i$（$i=1,\ldots,r$）以及滤过
$$
0 = \mathcal{E}_0 \subset \mathcal{E}_1 \subset \mathcal{E}_2
\subset \ldots \subset \mathcal{E}_r = \mathcal{E}
$$
使得 $\mathcal{E}_i/\mathcal{E}_{i-1}\cong\mathcal{L}_i$。
于是得到对偶滤过
% P10-E0126
$$
0 = \mathcal{E}_r^\perp \subset \mathcal{E}_1^\perp \subset \mathcal{E}_2^\perp
\subset \ldots \subset \mathcal{E}_0^\perp = \mathcal{E}^\vee
$$
满足 $\mathcal{E}_{i-1}^\perp/\mathcal{E}_i^\perp\cong
\mathcal{L}_i^{\otimes-1}$。令 $x_i=c_1(\mathcal{L}_i)$。
由引理 \ref{spaces-chow-lemma-c1-cap-additive}，
$c_1(\mathcal{L}_i^{\otimes-1})=-x_i$。
由引理 \ref{spaces-chow-lemma-chern-filter-by-linebundles}，在 $A^*(X)$ 中有
$$
c(\mathcal{E}) = \prod\nolimits_{i = 1}^r (1 + x_i)
\quad\text{且}\quad
c(\mathcal{E}^\vee) = \prod\nolimits_{i = 1}^r (1 - x_i)
$$
结论由一个形式计算得出，略去该计算。
\end{proof}

\begin{lemma}
\label{spaces-chow-lemma-chern-classes-tensor-product}
在情形 \ref{spaces-chow-situation-setup} 中，设 $X/B$ 是良好的。
设 $\mathcal{E}$ 和 $\mathcal{F}$ 是秩分别为 $r$ 和 $s$ 的一个有限
局部自由 $\mathcal{O}_X$-模。% P10-E0127
则有
$$
c_1(\mathcal{E} \otimes \mathcal{F})
=
r c_1(\mathcal{F}) + s c_1(\mathcal{E})
$$
$$
c_2(\mathcal{E} \otimes \mathcal{F})
=
r^2 c_2(\mathcal{F}) +
rs c_1(\mathcal{F})c_1(\mathcal{E}) +
s^2 c_2(\mathcal{E})
$$
等等（见证明）。% P10-E0128
\end{lemma}

\begin{proof}
完全仿照引理 \ref{spaces-chow-lemma-chern-classes-dual} 的证明，可以假设有可逆
$\mathcal{O}_X$-模
${\mathcal L}_i$（$i=1,\ldots,r$）、
${\mathcal N}_i$（$i=1,\ldots,s$）以及滤过
$$
0 = \mathcal{E}_0 \subset \mathcal{E}_1 \subset \mathcal{E}_2
\subset \ldots \subset \mathcal{E}_r = \mathcal{E}
\quad\text{且}\quad
0 = \mathcal{F}_0 \subset \mathcal{F}_1 \subset \mathcal{F}_2
\subset \ldots \subset \mathcal{F}_s = \mathcal{F}
$$
使得 $\mathcal{E}_i/\mathcal{E}_{i-1}\cong\mathcal{L}_i$，且
$\mathcal{F}_j/\mathcal{F}_{j-1}\cong\mathcal{N}_j$。
按字典序排列二元组 $(i,j)$，得到滤过
$$
0 \subset \ldots \subset
\mathcal{E}_i \otimes \mathcal{F}_j
+
\mathcal{E}_{i - 1} \otimes \mathcal{F}
\subset \ldots \subset \mathcal{E} \otimes \mathcal{F}
$$
其逐次商为
$$
\mathcal{L}_1 \otimes \mathcal{N}_1,
\mathcal{L}_1 \otimes \mathcal{N}_2,
\ldots,
\mathcal{L}_1 \otimes \mathcal{N}_s,
\mathcal{L}_2 \otimes \mathcal{N}_1,
\ldots,
\mathcal{L}_r \otimes \mathcal{N}_s
$$
由引理 \ref{spaces-chow-lemma-chern-filter-by-linebundles}，在 $A^*(X)$ 中有
% P10-E0129
$$
c(\mathcal{E}) = \prod (1 + x_i),
\quad
c(\mathcal{F}) = \prod (1 + y_j),
\quad\text{且}\quad
c(\mathcal{F}) = \prod (1 + x_i + y_j),
$$
结论由一个形式计算得出，略去该计算。
\end{proof}











\section{零维闭链的次数}
\label{spaces-chow-section-degree-zero-cycles}

\noindent
本节对应于《Chow 同调》第 \ref{chow-section-degree-zero-cycles} 节。
先定义域上固有代数空间中零维闭链的次数。

\begin{definition}
\label{spaces-chow-definition-degree-zero-cycle}
设 $k$ 是域，$p:X\to\Spec(k)$ 是代数空间的固有态射。
$X$ 上{\it 零维闭链的次数}由固有推出
$$
p_* : \CH_0(X) \longrightarrow \CH_0(\Spec(k)) \longrightarrow \mathbf{Z}
$$
（引理 \ref{spaces-chow-lemma-proper-pushforward-rational-equivalence}）与典范同构
$\CH_0(\Spec(k))\to\mathbf{Z}$ 的复合给出；该同构把
$[\Spec(k)]$ 映为 $1$。记作 $\deg(\alpha)$。
\end{definition}

\noindent
下面把定义写得更明确。

\begin{lemma}
\label{spaces-chow-lemma-spell-out-degree-zero-cycle}
设 $k$ 是域，$X$ 是 $k$ 上的固有代数空间。
设 $\alpha=\sum n_i[Z_i]\in Z_0(X)$。则
$$
\deg(\alpha) = \sum n_i\deg(Z_i)
$$
其中 $\deg(Z_i)$ 是 $Z_i\to\Spec(k)$ 的次数，即
$\deg(Z_i)=\dim_k\Gamma(Z_i,\mathcal{O}_{Z_i})$。
\end{lemma}

\begin{proof}
这正是固有推出的定义（定义 \ref{spaces-chow-definition-proper-pushforward}）。
\end{proof}

\begin{lemma}
\label{spaces-chow-lemma-degrees-and-numerical-intersections}
设 $k$ 是域，$X$ 是 $k$ 上的固有代数空间。
设 $Z\subset X$ 是维数为 $d$ 的闭子空间，并设
$\mathcal{L}_1,\ldots,\mathcal{L}_d$ 是可逆 $\mathcal{O}_X$-模。则
% P10-E0130
$$
(\mathcal{L}_1 \cdots \mathcal{L}_d \cdot Z) =
\deg(
c_1(\mathcal{L}_1) \cap \ldots \cap c_1(\mathcal{L}_1) \cap [Z]_d)
$$
其中左端定义于《域上的空间》定义
\ref{spaces-over-fields-definition-intersection-number}。
\end{lemma}

\begin{proof}
设 $Z_i\subset Z$（$i=1,\ldots,t$）是维数为 $d$ 的不可约分支，
$m_i$ 是 $Z_i$ 在 $Z$ 中的重数。则
$[Z]_d=\sum m_i[Z_i]$，并且
$c_1(\mathcal{L}_1) \cap \ldots \cap c_1(\mathcal{L}_d) \cap [Z]_d$
是闭链
$m_i c_1(\mathcal{L}_1) \cap \ldots \cap c_1(\mathcal{L}_d) \cap [Z_i]$.
之和。由《域上的空间》引理
\ref{spaces-over-fields-lemma-numerical-polynomial-leading-term}，
$(\mathcal{L}_1\cdots\mathcal{L}_d\cdot Z)$ 也有类似分解，
故只需在 $Z=X$ 是 $k$ 上固有整代数空间时证明引理。

\medskip\noindent
由 Chow 引理，存在固有态射 $f:X'\to X$，它在某个稠密开子空间
$U\subset X$ 上是同构，且 $X'$ 是概形。见《空间态射进阶》引理
\ref{spaces-more-morphisms-lemma-chow-noetherian-separated}。
于是 $X'$ 是 $k$ 上的固有概形。以 $f^{-1}(U)$ 的概形论闭包替换
$X'$ 后，可以假设 $X'$ 整。此时
$$
(f^*\mathcal{L}_1 \cdots f^*\mathcal{L}_d \cdot X') =
(\mathcal{L}_1 \cdots \mathcal{L}_d \cdot X)
$$
这是由《域上的空间》引理
\ref{spaces-over-fields-lemma-intersection-number-and-pullback} 得到的，
并且
% P10-E0131
$$
f_*(c_1(f^*\mathcal{L}_1) \cap \ldots \cap c_1(f^*\mathcal{L}_d) \cap [Y]) =
c_1(\mathcal{L}_1) \cap \ldots \cap c_1(\mathcal{L}_d) \cap [X]
$$
这是由引理 \ref{spaces-chow-lemma-pushforward-cap-c1} 得到的。因此可以用 $X'$
替换 $X$，并假设 $X$ 是 $k$ 上的固有概形。此情形已在
《Chow 同调》引理 \ref{chow-lemma-degrees-and-numerical-intersections}
中得到证明。
\end{proof}
% Generated mechanically from the frozen groupoids-quotients.tex target.
% Do not edit this generated file; edit the bound locale source instead.

% OK, start here.
%

\setcounter{chapter}{82}
\chapter{群胚的商}



\phantomsection
\label{groupoids-quotients-section-phantom}


\section{引言}
\label{groupoids-quotients-section-introduction}

\noindent
本章讨论群胚及其商（当商存在时）的一般理论。
关于这个主题有大量文献，例如参见
\cite{GIT}, \cite{seshadri_quotients}, \cite{KollarQuotients},
\cite{K-M}, \cite{KollarFinite} 以及更多文献。





\section{约定与记号}
\label{groupoids-quotients-section-conventions-notation}

\noindent
本章采用《空间中的群胚》第
\ref{spaces-groupoids-section-conventions} 节和第
\ref{spaces-groupoids-section-notation} 节中引入的约定与记号。


\section{不变态射}
\label{groupoids-quotients-section-invariant}

\begin{definition}
\label{groupoids-quotients-definition-invariant}
设 $S$ 为概形，$B$ 为 $S$ 上的代数空间。
设 $j = (t, s) : R \to U \times_B U$ 为 $B$ 上代数空间的一个预关系。
若下图交换，则称 $B$ 上代数空间的态射
$\phi : U \to X$ 是{\it $R$-不变的}：
$$
\xymatrix{
R \ar[r]_s \ar[d]_t & U \ar[d]^\phi \\
U \ar[r]^\phi & X
}
$$
若 $j : R \to U \times_B U$ 来自群代数空间 $G$ 在 $B$ 上对 $U$
的作用，如《空间中的群胚》引理
\ref{spaces-groupoids-lemma-groupoid-from-action} 所述，
则称 $\phi$ 是{\it $G$-不变的}。
\end{definition}

\noindent
换言之，态射 $U \to X$ 是 $R$-不变的，当且仅当它使 $s$ 与 $t$
相等。用相应的商层可将这一条件改述如下。

\begin{lemma}
\label{groupoids-quotients-lemma-invariant}
设 $S$ 为概形，$B$ 为 $S$ 上的代数空间。
设 $j = (t, s) : R \to U \times_B U$ 为 $B$ 上代数空间的一个预关系。
代数空间的态射 $\phi : U \to X$ 是 $R$-不变的，当且仅当它分解为
$U \to U/R \to X$.
\end{lemma}

\begin{proof}
这由《空间中的群胚》第
\ref{spaces-groupoids-section-quotient-sheaves} 节中商层的定义立即得到。
\end{proof}

\begin{lemma}
\label{groupoids-quotients-lemma-base-change-on-invariant}
设 $S$ 为概形，$B$ 为 $S$ 上的代数空间。
设 $j = (t, s) : R \to U \times_B U$ 为 $B$ 上代数空间的一个预关系。
设 $U \to X$ 为 $B$ 上代数空间的一个 $R$-不变态射，
并设 $X' \to X$ 为代数空间的任意态射。
\begin{enumerate}
\item 令 $U' = X' \times_X U$、$R' = X' \times_X R$，则得到一个预关系
$j' : R' \to U' \times_B U'$。
\item 若 $j$ 是关系，则 $j'$ 是关系。
\item 若 $j$ 是预等价关系，则 $j'$ 是预等价关系。
\item 若 $j$ 是等价关系，则 $j'$ 是等价关系。
\item 若 $j$ 来自 $B$ 上代数空间中的群胚
$(U, R, s, t, c)$，则
\begin{enumerate}
\item $(U, R, s, t, c)$ 是 $X$ 上代数空间中的群胚；并且
\item $j'$ 来自此群胚到 $X'$ 的基变换
$(U', R', s', t', c')$，参见《空间中的群胚》引理
\ref{spaces-groupoids-lemma-base-change-groupoid}.
\end{enumerate}
\item 若 $j$ 来自群代数空间 $G/B$ 在 $U$ 上的作用，如《空间中的群胚》引理
\ref{spaces-groupoids-lemma-groupoid-from-action}
所述，则 $j'$ 来自 $G$ 在 $U'$ 上的诱导作用。
\end{enumerate}
\end{lemma}

\begin{proof}
略。提示：结合下图采用函子观点：
$$
\xymatrix{
R' = X' \times_X R \ar[dd] \ar[rr] \ar[rd] & &
X' \times_X U = U' \ar'[d][dd] \ar[rd] \\
& R \ar[dd] \ar[rr] & & U \ar[dd] \\
U' = X' \times_X U \ar'[r][rr] \ar[rd] & & X' \ar[rd] \\
& U \ar[rr] & & X
}
$$
\end{proof}

\begin{definition}
\label{groupoids-quotients-definition-base-change}
在引理 \ref{groupoids-quotients-lemma-base-change-on-invariant} 的情形下，
称 $j' : R' \to U' \times_B U'$ 为预关系 $j$ 到 $X'$ 的{\it 基变换}。
若 $X' \to X$ 是代数空间的平坦态射，则称其为{\it 平坦基变换}。
\end{definition}

\noindent
这种基变换与取商层和商叠相容良好。

\begin{lemma}
\label{groupoids-quotients-lemma-base-change-quotient-sheaf}
在引理 \ref{groupoids-quotients-lemma-base-change-on-invariant} 的情形下，存在层的同构
$$
U'/R' = X' \times_X U/R
$$
关于商层的构造，参见《空间中的群胚》第
\ref{spaces-groupoids-section-quotient-sheaves} 节。
\end{lemma}

\begin{proof}
由于 $U \to X$ 是 $R$-不变的，显然映射 $U \to X$ 通过商层 $U/R$ 分解。
回顾按定义，
$$
\xymatrix{
R \ar@<1ex>[r] \ar@<-1ex>[r] &
U \ar[r] &
U/R
}
$$
是 $(\Sch/S)_{fppf}$ 上集合层范畴 $\Sh$ 中的余等化子图。
事实上，它还是逗号范畴 $\Sh/X$ 中的余等化子图。
由于基变换函子
$X' \times_X - : \Sh/X \to \Sh/X'$ 是正合的（这在任意拓扑斯中都成立），
结论随即得到。
\end{proof}

\begin{lemma}
\label{groupoids-quotients-lemma-base-change-quotient-stack}
设 $S$ 为概形，$B$ 为 $S$ 上的代数空间。
设 $(U, R, s, t, c)$ 为 $B$ 上代数空间中的群胚。
设 $U \to X$ 为 $B$ 上代数空间的一个 $R$-不变态射，
$g : X' \to X$ 为 $B$ 上代数空间的态射，并令
$(U', R', s', t', c')$ 为引理
\ref{groupoids-quotients-lemma-base-change-on-invariant} 中的基变换。则
$$
\xymatrix{
[U'/R'] \ar[r] \ar[d] & [U/R] \ar[d] \\
\mathcal{S}_{X'} \ar[r] & \mathcal{S}_X
}
$$
是 $(\Sch/S)_{fppf}$ 上群胚叠的一个 $2$-纤维积。
关于商叠以及图中态射的构造，参见《空间中的群胚》第
\ref{spaces-groupoids-section-stacks} 节。
\end{lemma}

\begin{proof}
我们使用《空间中的群胚》引理
\ref{spaces-groupoids-lemma-quotient-stack-objects}
给出的商叠之显式描述来证明这一点。
不过，我们强烈建议读者自行找出证明。
首先，可将 $(U, R, s, t, c)$ 视为 $X$ 上代数空间中的群胚，
从而得到映射 $f : [U/R] \to \mathcal{S}_X$，参见《空间中的群胚》引理
\ref{spaces-groupoids-lemma-quotient-stack-arrows}。
类似地，有 $f' : [U'/R'] \to X'$。

\medskip\noindent
% P10-E0133
$S$ 上概形 $T$ 上的 $2$-纤维积
$\mathcal{S}_{X'} \times_{\mathcal{S}_X} [U/R]$ 的一个对象，
等价于一个态射 $x' : T \to X'$ 和 $T$ 上 $[U/R]$ 的一个对象 $y$，
满足复合 $g \circ x'$ 等于 $f(y)$。
这是有意义的，因为 $T$ 上 $\mathcal{S}_X$ 的对象就是态射 $T \to X$。
由《空间中的群胚》引理
\ref{spaces-groupoids-lemma-quotient-stack-objects}
可设 $y$ 由相对于一个 fppf 覆盖 $\{T_i \to T\}$ 的
$[U/R]$-下降数据 $(u_i, r_{ij})$ 给出。
等式 $g \circ x' = f(y)$ 意味着下列图表
$$
\vcenter{
\xymatrix{
T_i \ar[rr]_{u_i} \ar[d] & & U \ar[d] \\
T \ar[r]^{x'} & X' \ar[r]^g & X
}
}
\quad\text{以及}\quad
\vcenter{
\xymatrix{
T_i \times_T T_j \ar[rr]_{r_{ij}} \ar[d] & & R \ar[d] \\
T \ar[r]^{x'} & X' \ar[r]^g & X
}
}
$$
均交换。

\medskip\noindent
另一方面，由《空间中的群胚》引理
\ref{spaces-groupoids-lemma-quotient-stack-objects}
可知，$S$ 上概形 $T$ 上 $[U'/R']$ 的一个对象 $y'$，
由相对于一个 fppf 覆盖 $\{T_i \to T\}$ 的
$[U'/R']$-下降数据 $(u'_i, r'_{ij})$ 给出。
令 $f'(y') = x' : T \to X'$，可见下列图表
% P10-E0135
$$
\vcenter{
\xymatrix{
T_i \ar[r]_{u'_i} \ar[d] & U' \ar[d] \\
T \ar[r]^{x'} & X'
}
}
\quad\text{以及}\quad
\vcenter{
\xymatrix{
T_i \times_T T_j \ar[r]_{r'_{ij}} \ar[d] & U' \ar[d] \\
T \ar[r]^{x'} & X'
}
}
$$
均交换。

\medskip\noindent
有了这些记号，定义函子
$$
[U'/R'] \longrightarrow \mathcal{S}_{X'} \times_{\mathcal{S}_X} [U/R]
$$
如下：将上述 $y' = (u'_i, r'_{ij})$ 送到对象
$(x', (u_i, r_{ij}))$，其中 $x' = f'(y')$，
$u_i$ 是复合 $T_i \to U' \to U$，
$r_{ij}$ 是复合 $T_i \times_T T_j \to R' \to R$。
% P10-E0134
反之，给定右侧的对象 $(x', (u_i, r_{ij})$，
可将其送到左侧的对象 $((x', u_i), (x', r_{ij}))$。
关于如何处理态射的讨论从略（其处理方式完全相同）。
\end{proof}









\section{范畴商}
\label{groupoids-quotients-section-categorical}

\noindent
这是所能考虑的最基本的一类商。

\begin{definition}
\label{groupoids-quotients-definition-categorical}
设 $S$ 为概形，$B$ 为 $S$ 上的代数空间。
设 $j = (t, s) : R \to U \times_B U$ 为 $B$ 上代数空间中的预关系。
\begin{enumerate}
% P10-E0132
\item 若 $B$ 上代数空间的态射 $\phi : U \to X$ 是 $R$-不变的，
且对 $B$ 上代数空间的每个 $R$-不变态射 $\psi : U \to Y$，
都存在唯一态射 $\chi : X \to Y$，使得
$\psi = \phi \circ \chi$.
则称其为{\it 范畴商}。
\item 设 $\mathcal{C}$ 为 $B$ 上代数空间范畴的一个全子范畴，
并设 $U,R$ 为 $\mathcal{C}$ 的对象。在这种情形下，若
$X \in \Ob(\mathcal{C})$、$\phi$ 是 $R$-不变的，且对每个满足
$Y \in \Ob(\mathcal{C})$ 的 $R$-不变态射
$\psi : U \to Y$，都存在唯一态射 $\chi : X \to Y$，使得
% P10-E0132
$\psi = \phi \circ \chi$，则称 $B$ 上代数空间的态射
$\phi : U \to X$ 为{\it $\mathcal{C}$ 中的范畴商}。
\item 若 $B=S$ 且 $\mathcal{C}$ 是 $S$ 上概形的范畴，
则称 $U\to X$ 为{\it 概形范畴中的范畴商}，或简称为
{\it 概形中的范畴商}。
\end{enumerate}
\end{definition}

\noindent
我们常用某种分离性公理从 $B$ 上代数空间中选出一个范畴
$\mathcal{C}$；若干标准情形见例 \ref{groupoids-quotients-example-categories}。
注意，$\phi : U \to X$ 是范畴商，当且仅当 $U \to X$
是该范畴中态射 $t,s:R\to U$ 的余等化子。
因此立即得到下列引理。

\begin{lemma}
\label{groupoids-quotients-lemma-categorical-unique}
设 $S$ 为概形，$B$ 为 $S$ 上的代数空间，并设
$j : R \to U \times_B U$ 为 $B$ 上代数空间中的预关系。
若 $B$ 上代数空间范畴中的范畴商存在，则它在唯一同构意义下唯一。
$\textit{Spaces}/B$ 的全子范畴中的范畴商亦然。
\end{lemma}

\begin{proof}
参见《范畴》第 \ref{categories-section-coequalizers} 节。
\end{proof}

\begin{example}
\label{groupoids-quotients-example-categories}
设 $S$ 为概形，$B$ 为 $S$ 上的代数空间。
% P10-E0136
应用定义 \ref{groupoids-quotients-definition-categorical} 时常遇到如下范畴
$\mathcal{C}$ 的标准例子：
\begin{enumerate}
\item $\mathcal{C}$ 是所有 $B$ 上代数空间的范畴；
\item $B$ 是分离的，且 $\mathcal{C}$ 是所有 $B$ 上分离代数空间的范畴；
\item $B$ 是拟分离的，且 $\mathcal{C}$ 是所有 $B$ 上拟分离代数空间的范畴；
\item $B$ 是局部分离的，且 $\mathcal{C}$ 是所有 $B$ 上局部分离代数空间的范畴；
\item $B$ 是合宜的，且 $\mathcal{C}$ 是所有 $B$ 上合宜代数空间的范畴；以及
\item $S=B$，且 $\mathcal{C}$ 是 $S$ 上概形的范畴。
\end{enumerate}
在这些情形中，若 $\phi : U \to X$ 是范畴商，则分别称
$U \to X$ 为
(1) {\it 范畴商}，
(2) {\it 分离代数空间中的范畴商}，
(3) {\it 拟分离代数空间中的范畴商}，
(4) {\it 局部分离代数空间中的范畴商}，
(5) {\it 合宜代数空间中的范畴商}，
(6) {\it 概形中的范畴商}。
\end{example}

\begin{definition}
\label{groupoids-quotients-definition-universal-categorical}
设 $S$ 为概形，$B$ 为 $S$ 上的代数空间。
设 $\mathcal{C}$ 为 $B$ 上代数空间范畴中对纤维积封闭的全子范畴。
% P10-E0137
设 $j = (t, s) : R \to U \times_B U$ 为 $\mathcal{C}$ 中的预关系，
并设 $U \to X$ 为一个 $R$-不变态射，其中
$X \in \Ob(\mathcal{C})$。
\begin{enumerate}
\item 若对 $\mathcal{C}$ 中每个态射 $X'\to X$，态射
$U'=X'\times_XU\to X'$ 都是 $j$ 的基变换
% P10-E0138
$j':R'\to U'$ 在 $\mathcal{C}$ 中的范畴商，
则称 $U\to X$ 为 $\mathcal{C}$ 中的{\it 泛范畴商}。
\item 若对 $\mathcal{C}$ 中每个平坦态射 $X'\to X$，态射
$U'=X'\times_XU\to X'$ 都是 $j$ 的基变换
% P10-E0138
$j':R'\to U'$ 在 $\mathcal{C}$ 中的范畴商，
则称 $U\to X$ 为 $\mathcal{C}$ 中的{\it 一致范畴商}。
\end{enumerate}
\end{definition}

\begin{lemma}
\label{groupoids-quotients-lemma-categorical-reduced}
在定义 \ref{groupoids-quotients-definition-categorical} 的情形下，
若 $\phi : U \to X$ 是范畴商且 $U$ 是既约的，则 $X$ 是既约的。
对于例 \ref{groupoids-quotients-example-categories} 所列空间范畴 $\mathcal{C}$
中的范畴商，同样成立。
\end{lemma}

\begin{proof}
设 $X_{red}$ 为代数空间 $X$ 的既约化。
由于 $U$ 既约，态射 $\phi : U \to X$ 通过
$i : X_{red} \to X$ 分解（《空间的性质》引理
\ref{spaces-properties-lemma-map-into-reduction}）。
将所得态射记为 $\phi_{red} : U \to X_{red}$。
由 $\phi \circ s = \phi \circ t$，又因
$i : X_{red} \to X$ 是单态射，可知
$\phi_{red} \circ s = \phi_{red} \circ t$。
因此由 $\phi$ 的泛性质，存在态射 $\chi : X \to X_{red}$，使得
% P10-E0132
$\phi_{red} = \phi \circ \chi$。由唯一性，
$i \circ \chi = \text{id}_X$ 且
$\chi \circ i = \text{id}_{X_{red}}$。
故 $i$ 是同构，$X$ 既约。

\medskip\noindent
要说明此论证在范畴 $\mathcal{C}$ 中成立，只需说明
$\mathcal{C}$ 中对象的既约化仍是 $\mathcal{C}$ 的对象。
对每个标准例子的验证从略。
\end{proof}





\section{作为轨道空间的商}
\label{groupoids-quotients-section-orbits}

\noindent
设 $j = (t, s) : R \to U \times_B U$ 为预关系。
若 $j$ 是预等价关系，则粗略地说，$R$ 在 $U$ 上的“轨道”
就是 $U$ 的子集 $t(s^{-1}(\{u\}))$。
然而，若 $j$ 只是预关系，就需要取由 $R$ 生成的等价关系。

\begin{definition}
\label{groupoids-quotients-definition-orbit}
设 $S$ 为概形，$B$ 为 $S$ 上的代数空间，并设
$j : R \to U \times_B U$ 为 $B$ 上的预关系。
若 $u \in |U|$，则 $u$ 的{\it 轨道}，更确切地说其
{\it $R$-轨道}，定义为
$$
O_u =
\left\{
u' \in |U|\ :
\begin{matrix}
 \exists n \geq 1, \ \exists u_0, \ldots, u_n \in |U|\text{，使得 }
 u_0 = u \text{ 且 } u_n = u' \\
 \text{并且对所有 }i \in \{0, \ldots, n - 1\}\text{，下列情形之一成立：}
 u_i = u_{i + 1}\text{，或 } \\
 \exists r \in |R|, \ s(r) = u_i, t(r) = u_{i + 1}
 \text{，或 } \\
\exists r \in |R|, \ t(r) = u_i, s(r) = u_{i + 1}
\end{matrix}
\right\}
$$
\end{definition}

\noindent
显然，这些集合是某个等价关系的等价类；也就是说，
$u' \in O_u$ 当且仅当 $u \in O_{u'}$。
下列引理是《空间中的群胚》引理
\ref{spaces-groupoids-lemma-pre-equivalence-equivalence-relation-points}
的改述。

\begin{lemma}
\label{groupoids-quotients-lemma-pre-equivalence-equivalence-relation-points}
设 $B\to S$ 如第 \ref{groupoids-quotients-section-conventions-notation} 节所述，
并设 $j : R \to U \times_B U$ 为 $B$ 上代数空间的预等价关系。则
$$
O_u =
\{u' \in |U| \text{，使得 } \exists r \in |R|, \ s(r) = u, \ t(r) = u'\}.
$$
\end{lemma}

\begin{proof}
由上述《空间中的群胚》引理
\ref{spaces-groupoids-lemma-pre-equivalence-equivalence-relation-points}
可知，引理中定义的轨道 $O_u$ 给出 $|U|$ 的不交并分解。
因此它们等于定义 \ref{groupoids-quotients-definition-orbit} 中定义的轨道。
\end{proof}

\begin{lemma}
\label{groupoids-quotients-lemma-invariant-map-constant-on-orbit}
在定义 \ref{groupoids-quotients-definition-orbit} 的情形下，
设 $\phi : U \to X$ 为 $B$ 上代数空间的一个 $R$-不变态射。
则 $|\phi| : |U| \to |X|$ 在每条轨道上为常值。
\end{lemma}

\begin{proof}
只需对所有满足如下条件的 $u,u'\in|U|$ 证明
$\phi(u)=\phi(u')$：存在 $r\in|R|$，使得
$s(r)=u$ 且 $t(r)=u'$。由于 $\phi$ 使 $s$ 与 $t$ 相等，这是显然的。
\end{proof}

\noindent
若试图用轨道 $O_u\subset|U|$ 刻画商映射的性质，会遇到若干问题。
其中一个如下。设 $\Spec(k)\to B$ 为 $B$ 的几何点，考虑典范映射
$$
U(k) \longrightarrow |U|.
$$
通常，由 $j(R(k)) \subset U(k) \times U(k)$ 生成的等价关系之等价类，
并不是轨道 $O_u\subset|U|$ 的逆像。
一个简单得近乎可笑的例子是取
$S=B=\Spec(\mathbf{Z})$、$U=R=\Spec(k)$，且
$s=t=\text{id}_k$。此时 $|U|=|R|$ 只有一个点，
但 $U(k)/R(k)$ 却极其庞大。
一个更有意思的例子是取 $S=B=\Spec(\mathbf{Q})$，
% P10-E0139
选择数域 $K\subset L$，令 $U=\Spec(L)$、
$R=\Spec(L\otimes_KL)$，并取显然的映射 $s,t:R\to U$。
此时 $|U|$ 仍只有一个点，但商
$$
U(k)/R(k) = \Hom(K, k)
$$
含有不止一个元素。由这两个例子可见，若 $U\to X$ 是
$R$-不变映射，而我们希望它“分离轨道”，那么考虑诱导映射
$U(k)\to X(k)$ 并要求这些映射分离轨道，会得到强得多且更有意义的概念。
% P10-E0140

\medskip\noindent
这也有一个问题。具体而言，设
$S=B=\Spec(\mathbf{R})$、$U=\Spec(\mathbf{C})$，并设
$R=\Spec(\mathbf{C})\amalg\Spec(K)$，
其中 $\sigma:\mathbf{C}\to K$ 是某个域扩张。
令映射 $s,t$ 在分支 $\Spec(\mathbf{C})$ 上均由恒等映射给出，
而在第二个分支上分别由 $\sigma,\sigma\circ\tau$ 给出，
其中 $\tau$ 是复共轭。若 $K$ 是 $\mathbf{C}$ 的非平凡扩张，
则两点 $1,\tau\in U(\mathbf{C})$ 在
$j(R(\mathbf{C}))$ 下并不等价。
然而，选取基数足够大的扩张 $\mathbf{C}\subset\Omega$
（例如其基数大于 $K$ 的基数）后，
$1,\tau\in U(\mathbf{C})$ 在 $U(\Omega)$ 中的像却会变得等价！
直观上很清楚，这要么是因为 $s,t:R\to U$ 不是局部有限型的，
要么是因为域 $k$ 的基数不够大。

\medskip\noindent
基于这一点，作如下定义。

\begin{definition}
\label{groupoids-quotients-definition-geometric-orbits}
设 $S$ 为概形，$B$ 为 $S$ 上的代数空间，
$j:R\to U\times_BU$ 为 $B$ 上的预关系，并设
$\Spec(k)\to B$ 为 $B$ 的几何点。
\begin{enumerate}
\item 若 $\overline{u},\overline{u}'\in U(k)$ 在由关系
$j(R(k))\subset U(k)\times U(k)$ 生成的等价关系下属于同一等价类，
则称二者{\it 弱 $R$-等价}。
\item 若存在上域 $k\subset\Omega$，使二者在 $U(\Omega)$ 中的像
弱 $R$-等价，则称 $\overline{u},\overline{u}'\in U(k)$
{\it $R$-等价}。
\item $\overline{u}\in U(k)$ 的{\it 弱轨道}，更确切地说其
{\it 弱 $R$-轨道}，是 $U(k)$ 中所有与 $\overline{u}$
弱 $R$-等价的元素所成的集合。
\item $\overline{u}\in U(k)$ 的{\it 轨道}，更确切地说其
{\it $R$-轨道}，是 $U(k)$ 中所有与 $\overline{u}$
$R$-等价的元素所成的集合。
\end{enumerate}
% P10-E0141
\end{definition}

\noindent
在良好情形下，轨道与弱轨道实际上相同，见引理
\ref{groupoids-quotients-lemma-geometric-orbits}。下列引理在预等价关系这一特殊情形中
说明二者的差异。

\begin{lemma}
\label{groupoids-quotients-lemma-weak-orbit-pre-equivalence}
设 $S$ 为概形，$B$ 为 $S$ 上的代数空间，
$\Spec(k)\to B$ 为 $B$ 的几何点，并设
$j:R\to U\times_BU$ 为 $B$ 上的预等价关系。
此时 $\overline{u}\in U(k)$ 的弱轨道就是
$$
\{
\overline{u}' \in U(k)
\text{，使得 }
\exists \overline{r} \in R(k),
\ s(\overline{r}) = \overline{u},
\ t(\overline{r}) = \overline{u}'
\}
$$
而 $\overline{u}\in U(k)$ 的轨道为
$$
\{
\overline{u}' \in U(k) :
\exists\text{ 域扩张 }K/k, \ \exists\ r \in R(K),
\ s(r) = \overline{u}, \ t(r) = \overline{u}'\}
$$
\end{lemma}

\begin{proof}
这是因为按照预等价关系的定义，像
$j(R(k))\subset U(k)\times U(k)$ 是等价关系。
\end{proof}

\noindent
下面描述将任意预关系变成预等价关系的方法。我们将使用态射
\begin{equation}
\label{groupoids-quotients-equation-list}
\begin{matrix}
j_{diag} &
: &
U &
\longrightarrow &
U \times_B U, &
u &
\longmapsto &
(u, u) \\
j_{flip} &
: &
R &
\longrightarrow &
U \times_B U, &
r &
\longmapsto &
(s(r), t(r)) \\
j_{comp} &
: &
R \times_{s, U, t} R &
\longrightarrow &
U \times_B U, &
(r, r') &
\longmapsto &
(t(r), s(r'))
\end{matrix}
\end{equation}
定义 $j_1 = (t_1, s_1) : R_1 \to U \times_B U$ 为态射
$$
j \amalg j_{diag} \amalg j_{flip} :
R \amalg U \amalg R
\longrightarrow
U \times_B U
$$
其中记号如方程 (\ref{groupoids-quotients-equation-list}) 所示。对 $n>1$，令
$$
j_n = (t_n, s_n) :
R_n = R_1 \times_{s_1, U, t_{n - 1}} R_{n - 1} \longrightarrow U \times_B U
$$
其中 $t_n$ 是 $t_1$ 与到 $R_1$ 的投影之预复合，
$s_n$ 是 $s_{n-1}$ 与到 $R_{n-1}$ 的投影之预复合。最后记
$$
j_\infty = (t_\infty, s_\infty) :
R_\infty = \coprod\nolimits_{n \geq 1} R_n
\longrightarrow
U \times_B U.
$$

\begin{lemma}
\label{groupoids-quotients-lemma-make-pre-equivalence}
设 $S$ 为概形，$B$ 为 $S$ 上的代数空间，
并设 $j:R\to U\times_BU$ 为 $B$ 上的预关系。
则 $j_\infty:R_\infty\to U\times_BU$ 是 $B$ 上的预等价关系。此外，
\begin{enumerate}
\item $\phi:U\to X$ 是 $R$-不变的，当且仅当它是
$R_\infty$-不变的；
\item 商层的典范映射 $U/R\to U/R_\infty$（见《空间中的群胚》第
\ref{spaces-groupoids-section-quotient-sheaves} 节）是同构；
\item 弱 $R$-轨道与弱 $R_\infty$-轨道相同；
\item $R$-轨道与 $R_\infty$-轨道相同；
\item 若 $s,t$ 是局部有限型的，则 $s_\infty,t_\infty$
是局部有限型的；
\item 按需在此增补。
% P10-E0142
\end{enumerate}
\end{lemma}

\begin{proof}
略。对 (5) 的提示：$s,t$ 的任何满足下列条件的性质都会由
$s_\infty,t_\infty$ 继承：对复合稳定、对基变换稳定，
且在源上是 Zariski 局部的。
\end{proof}

\begin{lemma}
\label{groupoids-quotients-lemma-geometric-orbits}
设 $S$ 为概形，$B$ 为 $S$ 上的代数空间，
$j:R\to U\times_BU$ 为 $B$ 上的预关系，并设
$\Spec(k)\to B$ 为 $B$ 的几何点。
\begin{enumerate}
\item 若 $s,t:R\to U$ 是局部有限型的，则 $U(k)$ 上的弱
$R$-等价与 $R$-等价相同，且 $U(k)$ 上的弱 $R$-轨道与
$R$-轨道相同。
\item 若 $k$ 的基数足够大，则 $U(k)$ 上的弱 $R$-等价与
$R$-等价相同，且 $U(k)$ 上的弱 $R$-轨道与 $R$-轨道相同。
\end{enumerate}
\end{lemma}

\begin{proof}
先证明 (1)。设 $s,t$ 局部有限型。由引理
\ref{groupoids-quotients-lemma-make-pre-equivalence}，可设 $R$ 是预等价关系。
设 $k$ 为 $B$ 上的代数闭域，并设
$\overline{u},\overline{u}'\in U(k)$ 是 $R$-等价的。
则对某个扩域 $\Omega/k$，存在一点
$\overline{r}\in R(\Omega)$ 映到
$(\overline{u},\overline{u}')\in(U\times_BU)(\Omega)$，
见引理 \ref{groupoids-quotients-lemma-weak-orbit-pre-equivalence}。因此
$$
Z = R \times_{j, U \times_B U, (\overline{u}, \overline{u}')} \Spec(k)
$$
非空。由于 $s$ 局部有限型，$j$ 也局部有限型，见《空间的态射》引理
\ref{spaces-morphisms-lemma-permanence-finite-type}。
这说明 $Z$ 是代数闭域 $k$ 上非空且局部有限型的代数空间
（使用《空间的态射》引理
\ref{spaces-morphisms-lemma-base-change-finite-type}）。
于是 $Z$ 有一个 $k$-值点，见《空间的态射》引理
\ref{spaces-morphisms-lemma-locally-finite-type-surjective-geometric-points}。
故存在 $\overline{r}\in R(k)$，使
$j(\overline{r})=(\overline{u},\overline{u}')$，从而
% P10-E0143
$\overline{u},\overline{u}'$ 如所需地是 $R$-等价的。

\medskip\noindent
第 (2) 部分的证明相同，只需用《空间的态射》引理
\ref{spaces-morphisms-lemma-large-enough}
代替《空间的态射》引理
\ref{spaces-morphisms-lemma-locally-finite-type-surjective-geometric-points}。
这说明，只要 $|k|>\lambda(R)$，断言即成立；其中
$\lambda(R)$ 是在《空间的态射》引理
\ref{spaces-morphisms-lemma-locally-finite-type-surjective-geometric-points}
之前引入的。
\end{proof}

\noindent
在下述定义中，术语“$k$ 是 $B$ 上的域”是指
$\Spec(k)$ 配备了一个态射 $\Spec(k)\to B$。

\begin{definition}
\label{groupoids-quotients-definition-set-theoretically-invariant}
设 $S$ 为概形，$B$ 为 $S$ 上的代数空间，并设
$j:R\to U\times_BU$ 为 $B$ 上的预关系。
\begin{enumerate}
\item 称 $\phi:U\to X$ {\it 在集合论意义下 $R$-不变}，
当且仅当对 $B$ 上每个代数闭域 $k$，映射
$U(k)\to X(k)$ 都使两个映射 $s,t:R(k)\to U(k)$ 相等。
\item 若 $\phi:U\to X$ 在集合论意义下 $R$-不变，
且对 $B$ 上每个代数闭域 $k$，等式
$\phi(\overline{u})=\phi(\overline{u}')$ 在 $X(k)$ 中成立时，
$\overline{u},\overline{u}'\in U(k)$ 均位于同一轨道，
则称该态射{\it 分离轨道}，或{\it 分离 $R$-轨道}。
\end{enumerate}
\end{definition}

\noindent
例 \ref{groupoids-quotients-example-not-invariant} 将说明，在代数空间范畴中，
集合论意义下的不变性是一个“过弱”的概念。
关于集合论意义下不变或分离轨道之含义的更几何化改述，见引理
\ref{groupoids-quotients-lemma-separates-orbits}。

\begin{lemma}
\label{groupoids-quotients-lemma-set-theoretic-invariant}
在定义 \ref{groupoids-quotients-definition-set-theoretically-invariant} 的情形下，
态射 $\phi:U\to X$ 在集合论意义下 $R$-不变，当且仅当对
$B$ 上任意代数闭域 $k$，映射 $U(k)\to X(k)$ 在每条轨道上为常值。
\end{lemma}

\begin{proof}
这是因为该条件要求对 $B$ 上所有代数闭域成立。
\end{proof}

\begin{lemma}
\label{groupoids-quotients-lemma-invariant-set-theoretically-invariant}
在定义 \ref{groupoids-quotients-definition-set-theoretically-invariant} 的情形下，
不变态射在集合论意义下不变。
\end{lemma}

\begin{proof}
这由定义立即得到。
\end{proof}

\begin{lemma}
\label{groupoids-quotients-lemma-set-theoretically-invariant-invariant-when-reduced}
在定义 \ref{groupoids-quotients-definition-set-theoretically-invariant} 的情形下，
设 $\phi:U\to X$ 为 $B$ 上代数空间的态射。假设
\begin{enumerate}
\item $\phi$ 在集合论意义下 $R$-不变；
\item $R$ 既约；并且
\item $X$ 在 $B$ 上局部分离。
\end{enumerate}
则 $\phi$ 是 $R$-不变的。
\end{lemma}

\begin{proof}
考虑等化子代数空间
$$
Z = R \times_{(\phi, \phi) \circ j, X \times_B X, \Delta_{X/B}} X
$$
。由假设 (3)，$Z\to R$ 是浸入。由假设 (1)，
$|Z|\to|R|$ 是满射。这说明 $Z\to R$ 是双射闭浸入
（使用《概形》引理 \ref{schemes-lemma-immersion-when-closed}）；
再由假设 (2) 得 $Z=R$。
\end{proof}

\begin{example}
\label{groupoids-quotients-example-not-invariant}
存在既约拟分离代数空间 $X,Y$ 以及一对态射
$a,b:Y\to X$，它们在所有 $k$-值点上一致，但彼此不相等。
为得到一个例子，取 $Y=\Spec(k[[x]])$ 以及
$$
X = \mathbf{A}^1_k \Big/ \big(\Delta \amalg \{(x, -x) \mid x \not = 0\}\big)
$$
即《空间》例 \ref{spaces-example-affine-line-involution}
中的代数空间。两个态射 $a,b:Y\to X$ 来自
$Y$ 到 $\mathbf{A}^1_k=\Spec(k[x])$ 的两个映射
$x\mapsto x$ 和 $x\mapsto -x$。在一般点处，这两个映射相同，
因为在空间 $X$ 的开部分 $x\not =0$ 上，函数 $x$ 与 $-x$ 相等。
在闭点处，两个映射显然也相同。另一方面，$a\not =b$。
这说明，若把引理
\ref{groupoids-quotients-lemma-set-theoretically-invariant-invariant-when-reduced}
中的假设 (3) 换成 $X$ 拟分离，该引理便不成立。具体地，考虑图表
$$
\xymatrix{
Y \ar[d]_{-1} \ar[r]_1 & Y \ar[d]^a \\
Y \ar[r]^a & X
}
$$
则复合 $a\circ(-1)=b$。因此可令 $R=Y$、$U=Y$、
$s=1$、$t=-1$、$\phi=a$，从而得到一个在集合论意义下不变、
但并非不变的态射之例。
\end{example}

\noindent
上述例子很有启发性，因为映射 $Y\to X$ 甚至分离轨道。
它表明代数空间范畴中集合论意义下不变的态射实在太多。
下面定义 $R$ 为集合论等价关系的含义；同时要记住，
为使其正确成立，必须允许域扩张。

\begin{definition}
\label{groupoids-quotients-definition-set-theoretic-equivalence}
设 $S$ 为概形，$B$ 为 $S$ 上的代数空间，并设
$j:R\to U\times_BU$ 为 $B$ 上的预关系。
\begin{enumerate}
\item 若对 $B$ 上所有代数闭域 $k$，由下式定义的 $U(k)$ 上关系
$\sim_R$
$$
\overline{u} \sim_R \overline{u}'
\Leftrightarrow
\begin{matrix}
\exists\text{ 域扩张 }K/k, \ \exists\ r \in R(K), \\
s(r) = \overline{u}, \ t(r) = \overline{u}'
\end{matrix}
$$
都是等价关系，则称 $j$ 为{\it 集合论预等价关系}。
\item 若 $j$ 是泛单射且是集合论预等价关系，
则称 $j$ 为{\it 集合论等价关系}。
\end{enumerate}
\end{definition}

\noindent
下面用更几何的语言改述这一点。

\begin{lemma}
\label{groupoids-quotients-lemma-set-theoretic-pre-equivalence-geometric}
在定义 \ref{groupoids-quotients-definition-set-theoretic-equivalence} 的情形下，
下列条件等价：
\begin{enumerate}
\item 态射 $j$ 是集合论预等价关系。
\item 子集 $j(|R|)\subset|U\times_BU|$ 包含方程
(\ref{groupoids-quotients-equation-list}) 中任一态射 $j'$ 的 $|j'|$ 之像。
\item 对 $B$ 上每个基数足够大的代数闭域 $k$，子集
$j(R(k))\subset U(k)\times U(k)$ 是等价关系。
\end{enumerate}
若 $s,t$ 局部有限型，则这些条件还等价于
\begin{enumerate}
\item[(4)] 对 $B$ 上每个代数闭域 $k$，子集
$j(R(k))\subset U(k)\times U(k)$ 是等价关系。
\end{enumerate}
\end{lemma}

\begin{proof}
假设 (2)。设 $k$ 为 $B$ 上的代数闭域。我们将证明
$\sim_R$ 是等价关系。设
$\overline{u}_i:\Spec(k)\to U$（$i=1,2$）为 $U$ 的 $k$-值点，
并设 $(\overline{u}_1,\overline{u}_2)$ 是某个
$K$-值点 $r\in R(K)$ 的像。考虑实线部分交换的图表
$$
\xymatrix{
\Spec(K') \ar@{..>}[r] \ar@{..>}[d]
&
\Spec(k) \ar[d]_{(\overline{u}_2, \overline{u}_1)} &
\Spec(K) \ar[d] \ar[l] \\
R \ar[r]^-j &
U \times_B U &
R \ar[l]_-{j_{flip}}
}
$$
仍以 $r\in|R|$ 表示 $r$ 的像。依假设，
$|j_{flip}|$ 的像包含于 $|j|$ 的像；换言之，存在
$r'\in|R|$，使 $|j|(r')=|j_{flip}|(r)$。
注意，由图表实线部分的交换性，
$(\overline{u}_2,\overline{u}_1)$ 位于定义 $|j|(r')$ 的等价类中。
这意味着存在域扩张 $K'/k$ 和态射
% P10-E0144
$r':\Spec(K)\to R$（也滥用记号写成 $r'$），满足
$j\circ r'=(\overline{u}_2,\overline{u}_1)\circ i$，
其中 $i:\Spec(K')\to\Spec(K)$ 是显然的映射。
换言之，图表虚线部分交换。这证明 $\sim_R$ 是 $U(k)$ 上的对称关系。
类似地，利用 $|j_{diag}|$ 的像包含于 $|j|$ 的像，
可见 $\sim_R$ 是自反的（细节从略）。

\medskip\noindent
为证明 $\sim_R$ 的传递性，设给定
$\overline{u}_i:\Spec(k)\to U$（$i=1,2,3$）、
域扩张 $K_i/k$ 以及点
$r_i:\Spec(K_i)\to R$（$i=1,2$），使得
$j(r_1)=(\overline{u}_1,\overline{u}_2)$ 且
% P10-E0145
$j(r_1)=(\overline{u}_2,\overline{u}_3)$。于是可选择域的交换图表
$$
\xymatrix{
K & K_2 \ar[l] \\
K_1 \ar[u] & k \ar[l] \ar[u]
}
$$
并可视为 $r_1,r_2\in R(K)$。考虑实线部分交换的图表
$$
\xymatrix{
\Spec(K') \ar@{..>}[r] \ar@{..>}[d]
&
\Spec(k) \ar[d]_{(\overline{u}_1, \overline{u}_3)} &
\Spec(K) \ar[d]^{(r_1, r_2)} \ar[l]
\\
R \ar[r]^-j &
U \times_B U &
R \times_{s, U, t} R \ar[l]_-{j_{comp}}
}
$$
采用与证明第一部分完全相同的论证，但这次使用
$|j_{comp}|((r_1,r_2))$ 位于 $|j|$ 的像中，可得存在域 $K'$
和虚线箭头，使图表交换。这证明 $\sim_R$ 是传递的，
也就完成了 (2) 推出 (1) 的证明。

\medskip\noindent
假设 (1)，并设 $k$ 为 $B$ 上基数大于 $\lambda(R)$ 的代数闭域，
见《空间的态射》引理 \ref{spaces-morphisms-lemma-large-enough}。
设 $\overline{u}\sim_R\overline{u}'$，其中
$\overline{u},\overline{u}'\in U(k)$。依假设，存在 $|R|$ 中一点
映到 $(\overline{u},\overline{u}')\in|U\times_BU|$。
因此由《空间的态射》引理 \ref{spaces-morphisms-lemma-large-enough}，
存在 $\overline{r}\in R(k)$，使
$j(\overline{r})=(\overline{u},\overline{u}')$。
由此可见 (1) 推出 (3)。

\medskip\noindent
假设 (3)。来证明
$\Im(|j_{comp}|)\subset\Im(|j|)$。任取一点
$c\in|R\times_{s,U,t}R|$。可用态射
$\overline{c}:\Spec(k)\to R\times_{s,U,t}R$ 表示它，
其中 $k$ 是 $B$ 上基数足够大的域。依假设，
$j_{comp}(\overline{c})\in U(k)\times U(k)=(U\times_BU)(k)$
也可写成某个 $\overline{r}\in R(k)$ 的像 $j(\overline{r})$。
故如所需，$j_{comp}(c)=j(r)$ 在 $|U\times_BU|$ 中成立
（其中 $r\in|R|$ 是 $\overline{r}$ 的等价类）。
同一论证还说明
$\Im(|j_{diag}|)\subset\Im(|j|)$ 以及
$\Im(|j_{flip}|)\subset\Im(|j|)$（细节从略）。
由此可见 (3) 推出 (2)。至此已证明 (1)、(2)、(3) 相互等价。

\medskip\noindent
显然 (4) 推出 (3)（无需对 $s,t$ 作任何假设）。
为完成引理的证明，来说明当 $s,t$ 局部有限型时 (1) 推出 (4)。
设 $k$ 为 $B$ 上的代数闭域，并设
$\overline{u}\sim_R\overline{u}'$，其中
$\overline{u},\overline{u}'\in U(k)$。依假设，代数空间
$Z=R\times_{j,U\times_BU,(\overline{u},\overline{u}')}\Spec(k)$
非空。另一方面，由于 $j=(t,s)$ 局部有限型，态射
$Z\to\Spec(k)$ 也局部有限型（使用《空间的态射》引理
\ref{spaces-morphisms-lemma-permanence-finite-type}
和 \ref{spaces-morphisms-lemma-base-change-finite-type}）。
因此由《空间的态射》引理
\ref{spaces-morphisms-lemma-locally-finite-type-surjective-geometric-points}，
$Z$ 有一个 $k$-点，从而如所需地有
$(\overline{u},\overline{u}')\in j(R(k))$。
引理得证。
\end{proof}

\begin{lemma}
\label{groupoids-quotients-lemma-set-theoretic-equivalence-geometric}
在定义 \ref{groupoids-quotients-definition-set-theoretic-equivalence} 的情形下，
下列条件等价：
\begin{enumerate}
\item 态射 $j$ 是集合论等价关系。
\item 态射 $j$ 是泛单射，且
$j(|R|)\subset|U\times_BU|$ 包含方程
(\ref{groupoids-quotients-equation-list}) 中任一态射 $j'$ 的 $|j'|$ 之像。
\item 对 $B$ 上每个基数足够大的代数闭域 $k$，映射
$j:R(k)\to U(k)\times U(k)$ 是单射，且其像是等价关系。
\end{enumerate}
若 $j$ 合宜、局部分离或拟分离，则这些条件还等价于
\begin{enumerate}
\item[(4)] 对 $B$ 上每个代数闭域 $k$，映射
$j:R(k)\to U(k)\times U(k)$ 是单射，且其像是等价关系。
\end{enumerate}
\end{lemma}

\begin{proof}
(1) $\Rightarrow$ (2) 和 (2) $\Rightarrow$ (3) 由引理
\ref{groupoids-quotients-lemma-set-theoretic-pre-equivalence-geometric} 及定义得到。
同一引理说明 (3) 蕴含 $j$ 是集合论预等价关系。
当然，条件 (3) 还蕴含 $j$ 泛单射，见《空间的态射》引理
\ref{spaces-morphisms-lemma-universally-injective}，
故 $j$ 的确是集合论等价关系。至此已知 (1)、(2)、(3) 相互等价。

\medskip\noindent
无需对 $j$ 作任何进一步假设，条件 (4) 就推出 (3)。
设 $j$ 合宜、局部分离或拟分离，且等价条件 (1)、(2)、(3) 成立。
由《空间态射进阶》引理
\ref{spaces-more-morphisms-lemma-when-universally-injective-radicial}
可知 $j$ 是纯不可分态射。任取 $B$ 上代数闭域 $k$，并设
$\overline{u},\overline{u}'\in U(k)$ 满足
$\overline{u}\sim_R\overline{u}'$。可见
$R \times_{U \times_B U, (\overline{u}, \overline{u}')} \Spec(k)$
非空。由于 $j$ 是纯不可分态射，其既约化是 $k$ 的一个纯不可分扩域之谱。
又因 $k=\overline{k}$，它就是 $k$ 的谱。因此如所需，存在一点
$\overline{r}\in R(k)$，使
$t(\overline{r})=\overline{u}$ 且
$s(\overline{r})=\overline{u}'$。
\end{proof}

\begin{lemma}
\label{groupoids-quotients-lemma-set-theoretic-equivalence}
设 $S$ 为概形，$B$ 为 $S$ 上的代数空间，并设
$j:R\to U\times_BU$ 为 $B$ 上的预关系。
\begin{enumerate}
\item 若 $j$ 是预等价关系，则 $j$ 是集合论预等价关系。
特别地，当 $j$ 来自代数空间中的群胚，或来自群代数空间在 $U$
上的作用时，此结论成立。
\item 若 $j$ 是等价关系，则 $j$ 是集合论等价关系。
\end{enumerate}
\end{lemma}

\begin{proof}
略。
\end{proof}

\begin{lemma}
\label{groupoids-quotients-lemma-separates-orbits}
设 $B\to S$ 如第 \ref{groupoids-quotients-section-conventions-notation} 节所述，
$j:R\to U\times_BU$ 为预关系，并设
$\phi:U\to X$ 为 $B$ 上代数空间的态射。考虑图表
$$
\xymatrix{
(U \times_X U) \times_{(U \times_B U)} R \ar[d]^q \ar[r]_-p & R \ar[d]^j \\
U \times_X U \ar[r]^c & U \times_B U
}
$$
则有：
\begin{enumerate}
\item 态射 $\phi$ 在集合论意义下不变，当且仅当 $p$ 满射。
\item 若 $j$ 是集合论预等价关系，则 $\phi$ 分离轨道，
当且仅当 $p$ 与 $q$ 都满射。
\item 若 $p$ 与 $q$ 都满射，则 $j$ 是集合论预等价关系
（且 $\phi$ 分离轨道）。
\item 若 $\phi$ 是 $R$-不变的，且 $j$ 是集合论预等价关系，
则 $\phi$ 分离轨道，当且仅当诱导态射
$R\to U\times_XU$ 满射。
\end{enumerate}
\end{lemma}

\begin{proof}
假设 $\phi$ 在集合论意义下不变。这意味着对 $B$ 上任意代数闭域
$k$ 及任意 $\overline{r}\in R(k)$，都有
$\phi(s(\overline{r}))=\phi(t(\overline{r}))$。因此
% P10-E0146
$((\phi(t(\overline{r})),\phi(s(\overline{r}))),\overline{r})$
定义了纤维积中的一点，并经 $p$ 映到 $\overline{r}$。
这说明 $p$ 满射。反之，设 $p$ 满射。取
$\overline{r}\in R(k)$。由 $p$ 满射，可找到域扩张 $K/k$
以及纤维积的一个 $K$-值点 $\tilde r$，使
$p(\tilde r)=\overline{r}$。于是
$q(\tilde r)\in U\times_XU$ 在 $U\times_BU$ 中映到
$(t(\overline{r}),s(\overline{r}))$，从而
$\phi(s(\overline{r}))=\phi(t(\overline{r}))$。
这证明 $\phi$ 在集合论意义下不变。

\medskip\noindent
(2)、(3)、(4) 的证明从略。提示：设 $k$ 是 $B$ 上基数很大的代数闭域。
考虑相应的集合图表
$$
\xymatrix{
(U(k) \times_{X(k)} U(k)) \times_{U(k) \times U(k)} R(k) \ar[d]^q \ar[r]_-p &
R(k) \ar[d]^j \\
U(k) \times_{X(k)} U(k) \ar[r]^c & U(k) \times U(k)
}
$$
由上述引理，(2)、(3)、(4) 中的等价性都化为与刚才图表有关的
集合论问题；这里使用《空间的态射》引理
\ref{spaces-morphisms-lemma-large-enough}，
它将满射性转化为 $k$-值点上的满射性。
\end{proof}

\noindent
上文已看到，在代数空间范畴中，集合论意义下不变态射的概念相当弱。
因此如下定义预关系的轨道空间。

\begin{definition}
\label{groupoids-quotients-definition-orbit-space}
设 $B\to S$ 如第 \ref{groupoids-quotients-section-conventions-notation} 节所述，
并设 $j:R\to U\times_BU$ 为预关系。若
\begin{enumerate}
\item $\phi$ 是 $R$-不变的；
\item $\phi$ 分离 $R$-轨道；并且
\item $\phi$ 满射，
\end{enumerate}
则称 $\phi:U\to X$ 为{\it $R$ 的轨道空间}。
\end{definition}

\noindent
分离 $R$-轨道的定义涉及取值于代数闭域的点。
但如前所见，在许多情形中，这只对应于某些典范相伴的代数空间态射之满射性。
下面用轨道空间的刻画概括上文的部分讨论。

\begin{lemma}
\label{groupoids-quotients-lemma-orbit-space}
设 $B\to S$ 如第 \ref{groupoids-quotients-section-conventions-notation} 节所述，
并设 $j:R\to U\times_BU$ 为集合论预等价关系。
态射 $\phi:U\to X$ 是 $R$ 的轨道空间，当且仅当
\begin{enumerate}
\item $\phi\circ s=\phi\circ t$，即 $\phi$ 不变；
\item 诱导态射 $(t,s):R\to U\times_XU$ 满射；并且
\item 态射 $\phi:U\to X$ 满射。
\end{enumerate}
例如，当 $j$ 是预等价关系，或来自 $B$ 上代数空间中的群胚，
或来自 $B$ 上群代数空间在 $U$ 上的作用时，此刻画均适用。
\end{lemma}

\begin{proof}
由引理 \ref{groupoids-quotients-lemma-separates-orbits} 的第 (4) 部分立即得到。
\end{proof}

\noindent
在下列引理中，仅假设态射 $s,t$ 局部有限型（大概）还不够。
原因是某个映射 $\phi:U\to X$ 可能是轨道空间，却并非局部有限型。
在这种情形下，对 $B$ 上所有代数闭域 $k$，映射
$U(k)\to X(k)$ 未必都满射。

\begin{lemma}
\label{groupoids-quotients-lemma-orbit-space-locally-finite-type-over-base}
设 $B\to S$ 如第 \ref{groupoids-quotients-section-conventions-notation} 节所述，
$j=(t,s):R\to U\times_BU$ 为预关系，并假设 $R,U$
在 $B$ 上局部有限型。设 $\phi:U\to X$ 为 $B$ 上代数空间的
$R$-不变态射。则 $\phi$ 是 $R$ 的轨道空间，当且仅当自然映射
$$
U(k)/\big(\text{由 }j(R(k))\text{ 生成的等价关系}\big)
\longrightarrow
X(k)
$$
对 $B$ 上所有代数闭域 $k$ 都是双射。
\end{lemma}

\begin{proof}
注意，由于 $U,R$ 在 $B$ 上局部有限型，所有态射
$s,t,j,\phi$ 都局部有限型，见《空间的态射》引理
\ref{spaces-morphisms-lemma-permanence-finite-type}。
下文还将默认使用《空间的态射》引理
\ref{spaces-morphisms-lemma-locally-finite-type-surjective-geometric-points}。
假设 $\phi$ 是轨道空间。任取 $B$ 上代数闭域 $k$ 及
$\overline{x}\in X(k)$。考虑
$U \times_{\phi, X, \overline{x}} \Spec(k)$.
这是 $k$ 上非空且局部有限型的代数空间，故它有一个 $k$-值点。
这说明引理中的映射满射。设
$\overline{u},\overline{u}'\in U(k)$ 映到 $X(k)$ 中同一元素。
由定义 \ref{groupoids-quotients-definition-set-theoretically-invariant}，
这意味着 $\overline{u},\overline{u}'$ 位于同一 $R$-轨道。
由引理 \ref{groupoids-quotients-lemma-geometric-orbits}，这又意味着二者在由
$j(R(k))$ 生成的等价关系下等价。因此所显示的态射是单射。

\medskip\noindent
反之，假设对 $B$ 上所有代数闭域 $k$，所显示的映射都是双射。
该条件显然蕴含 $\phi$ 满射。我们已经假设 $\phi$ 是 $R$-不变的。
最后，所有这些映射的单射性蕴含 $\phi$ 分离轨道。
故 $\phi$ 是轨道空间。
\end{proof}










\section{粗商}
\label{groupoids-quotients-section-coarse}

\noindent
这里只是补入这一概念，以便以后说明粗商对应于粗模空间（或模概形）。

\begin{definition}
\label{groupoids-quotients-definition-coarse}
设 $S$ 为概形，$B$ 为 $S$ 上的代数空间，
并设 $j:R\to U\times_BU$ 为预关系。若 $B$ 上代数空间的态射
$\phi:U\to X$ 满足
\begin{enumerate}
\item $\phi$ 是范畴商；并且
\item $\phi$ 是轨道空间，
\end{enumerate}
则称其为{\it 粗商}。若 $S=B$ 且 $U,R$ 均为概形，
那么若概形态射 $\phi:U\to X$ 满足
\begin{enumerate}
\item $\phi$ 是概形中的范畴商；并且
\item $\phi$ 是轨道空间，
\end{enumerate}
则称其为{\it 概形中的粗商}。
\end{definition}

\noindent
在许多情形下，代数空间 $R,U$ 在 $B$ 上局部有限型，
而轨道空间条件仅意味着
$$
U(k)/\big(\text{由 }j(R(k))\text{ 生成的等价关系}\big)
\cong
X(k)
$$
对所有代数闭域 $k$ 成立。参见引理
\ref{groupoids-quotients-lemma-orbit-space-locally-finite-type-over-base}。
若 $j$ 还是（集合论）预等价关系，则此条件仅等价于：
对所有代数闭域 $k$，$U(k)/j(R(k))\to X(k)$ 是双射。









\section{拓扑性质}
\label{groupoids-quotients-section-topological}

\noindent
设 $S$ 为概形，$B$ 为 $S$ 上的代数空间，并设
$j:R\to U\times_BU$ 为预关系。若作为 $|R|$ 的子集有
$s^{-1}(T)=t^{-1}(T)$，则称子集 $T\subset|U|$ 为
{\it $R$-不变的}。注意，若 $T$ 闭，则 $U$ 中相应的既约闭子空间
未必是 $R$-不变的（如《空间中的群胚》定义
\ref{spaces-groupoids-definition-invariant-open} 所述），
因为拉回 $s^{-1}(T),t^{-1}(T)$ 未必既约。
对不变态射 $\phi:U\to X$，可考虑如下条件。

\begin{definition}
\label{groupoids-quotients-definition-topological}
设 $S$ 为概形，$B$ 为 $S$ 上的代数空间，
$j:R\to U\times_BU$ 为预关系，并设
$\phi:U\to X$ 为 $B$ 上代数空间的 $R$-不变态射。
\begin{enumerate}
\item
\label{groupoids-quotients-item-submersive}
态射 $\phi$ 是下沉的。
\item
\label{groupoids-quotients-item-invariant-closed}
对任意 $R$-不变闭子集 $Z\subset|U|$，像 $\phi(Z)$ 在 $|X|$ 中闭。
\item
\label{groupoids-quotients-item-intersect-invariant-closed}
条件 (\ref{groupoids-quotients-item-invariant-closed}) 成立，且对任意一对
$R$-不变闭子集 $Z_1,Z_2\subset|U|$，有
$$
\phi(Z_1 \cap Z_2) = \phi(Z_1) \cap \phi(Z_2)
$$
\item 态射 $(t,s):R\to U\times_XU$ 是泛下沉的。
\label{groupoids-quotients-item-strong}
\end{enumerate}
还可对这些性质逐一要求其在任意平坦基变换后成立，
或在任意基变换后成立，见定义 \ref{groupoids-quotients-definition-base-change}。
在这些情形下，分别称条件
(\ref{groupoids-quotients-item-submersive}),
(\ref{groupoids-quotients-item-invariant-closed}),
(\ref{groupoids-quotients-item-intersect-invariant-closed})，或
(\ref{groupoids-quotients-item-strong}) {\it 一致地}成立或{\it 泛地}成立。
\end{definition}








\section{不变函数}
\label{groupoids-quotients-section-functions}

\noindent
在某些情形下，要求每个不变函数都是商之结构层的局部截面，
可以方便地确定商的结构层。

\begin{definition}
\label{groupoids-quotients-definition-functions}
设 $S$ 为概形，$B$ 为 $S$ 上的代数空间，
$j:R\to U\times_BU$ 为预关系，并设
$\phi:U\to X$ 为 $R$-不变态射。记
$\phi'=\phi\circ s=\phi\circ t:R\to X$。
\begin{enumerate}
\item 以 $(\phi_*\mathcal{O}_U)^R$ 表示 $\phi_*\mathcal{O}_U$ 的
$\mathcal{O}_X$-子代数，它是下列两个映射的等化子：
$$
\xymatrix{
\phi_*\mathcal{O}_U
\ar@<1ex>[rr]^{\phi_*s^\sharp}
\ar@<-1ex>[rr]_{\phi_*t^\sharp}
& &
\phi'_*\mathcal{O}_R
}
$$
这里的映射定义在 $X_\etale$ 上。有时称其为
{\it $X$ 上的 $R$-不变函数层}。
\item 若自然映射
$\mathcal{O}_X\to(\phi_*\mathcal{O}_U)^R$ 是同构，
则称{\it $X$ 上的函数就是 $U$ 上的 $R$-不变函数}。
\end{enumerate}
\end{definition}

\noindent
当然，可以要求此性质在任意（平坦或一般）基变换后成立，
从而得到（一致或）泛的概念。为得到（更为）唯一的商，
常把此条件与其他条件一并加入。整个主题的很大一部分动机，
当然来自如下特殊情形：$U=\Spec(A)$ 是域
$S=B=\Spec(k)$ 上的仿射概形，$R=G\times U$，
其中 $G$ 是 $k$ 上的仿射群概形。此时可取如下商：
$$
X = \Spec(A^G)
$$
这样至少满足上述定义中的条件。虽然这是一个不错的构造，
但它往往不是正确的商；例如，若 $U=\text{GL}_{n,k}$ 且
$G$ 为上三角矩阵群，则上述构造给出 $X=\Spec(k)$，
而事实上存在好得多的商，即旗簇。








\section{良商}
\label{groupoids-quotients-section-good}

\noindent
特别是在取群作用的商时，下列定义很有用。

\begin{definition}
\label{groupoids-quotients-definition-good}
设 $S$ 为概形，$B$ 为 $S$ 上的代数空间，
并设 $j:R\to U\times_BU$ 为预关系。
若 $B$ 上代数空间的态射 $\phi:U\to X$ 满足
\begin{enumerate}
\item $\phi$ 不变；
\item $\phi$ 是仿射的；
\item $\phi$ 满射；
\item 条件 (\ref{groupoids-quotients-item-intersect-invariant-closed}) 泛地成立；并且
\item $X$ 上的函数就是 $U$ 上的 $R$-不变函数，
\end{enumerate}
则称其为{\it 良商}。
\end{definition}

\noindent
Seshadri 在 \cite{seshadri_quotients} 中给出几乎相同的定义；
区别是他用“条件 (\ref{groupoids-quotients-item-intersect-invariant-closed}) 成立”
代替 (4)，并不要求该条件泛地成立。







\section{几何商}
\label{groupoids-quotients-section-geometric}

\noindent
这是 Mumford 对几何商的定义（至少是 GIT 第一版中的定义；
据我们所知，后来的版本把“泛下沉”改成了“下沉”）。

\begin{definition}
\label{groupoids-quotients-definition-geometric}
设 $S$ 为概形，$B$ 为 $S$ 上的代数空间，
并设 $j:R\to U\times_BU$ 为预关系。
若 $B$ 上代数空间的态射 $\phi:U\to X$ 满足
\begin{enumerate}
\item $\phi$ 是轨道空间；
\item 条件 (\ref{groupoids-quotients-item-submersive}) 泛地成立，即 $\phi$ 泛下沉；并且
\item $X$ 上的函数就是 $U$ 上的 $R$-不变函数，
\end{enumerate}
则称其为{\it 几何商}。
\end{definition}
% Generated mechanically from the frozen spaces-more-cohomology.tex target.
% Do not edit this generated file; edit the bound locale source instead.

% OK, start here.
%

\setcounter{chapter}{83}
\chapter{空间上同调续论}


\phantomsection
\label{spaces-more-cohomology-section-phantom}





\section{引言}
\label{spaces-more-cohomology-section-introduction}

\noindent
% P10-E0147
本章继续《空间的上同调》第
\ref{spaces-cohomology-section-introduction} 节中开始的讨论。
也可以把本章视为概形平展上同调一章在代数空间情形下的对应版本，
参见《平展上同调》第
\ref{etale-cohomology-section-introduction} 节。

\medskip\noindent
事实上，本章主要旨在把已经为概形证明的结果转述为代数空间的语言。
其中一些结果可见 \cite{Kn}。





\section{约定}
\label{spaces-more-cohomology-section-conventions}

\noindent
始终假设所有概形都包含在一个大 fppf 站点 $\Sch_{fppf}$ 中。
并且所考虑的每个环 $A$ 都满足：$\Spec(A)$（同构于）该大站点的一个对象。

\medskip\noindent
设 $S$ 为概形，$X$ 为 $S$ 上的代数空间。
在本章及后续各章中，我们用 $X \times_S X$ 表示 $X$ 与自身的积
（在 $S$ 上代数空间的范畴中），而不用 $X \times X$。







\section{从概形迁移结果}
\label{spaces-more-cohomology-section-api}

\noindent
本节简要说明，概形的结果如何推出（可表）代数空间以及代数空间的
（可表）态射的相应结果。
对拟凝聚模还可得到更强的结论
（因为概形上拟凝聚模的平展上同调与扎里斯基上同调相同）；
《空间的上同调》第
\ref{spaces-cohomology-section-higher-direct-image} 节已经讨论过这一点。

\medskip\noindent
设 $S$ 为概形，$X$ 为 $S$ 上的代数空间。
现假设 $X$ 可由概形 $X_0$ 表示
（这是一个略显笨拙但只是临时采用的记号；通常我们直接说“$X$
是概形”）。在这种情形下，$X$ 与 $X_0$ 具有相同的小平展站点：
$$
X_\etale = (X_0)_\etale
$$
这在《空间的性质》第 \ref{spaces-properties-section-etale-site} 节中已有说明。
此外，若 $f : X \to Y$ 是 $S$ 上可表代数空间之间的态射，
而 $f_0 : X_0 \to Y_0$ 是表示 $f$ 的概形态射，
则所诱导的小平展拓扑斯态射彼此相同：
$$
\xymatrix{
\Sh(X_\etale) \ar[rr]_{f_{small}} \ar@{=}[d] & &
\Sh(Y_\etale) \ar@{=}[d] \\
\Sh((X_0)_\etale) \ar[rr]^{(f_0)_{small}} & &
\Sh((Y_0)_\etale)
}
$$
参见《空间的性质》引理
\ref{spaces-properties-lemma-functoriality-etale-site} 和
《拓扑》引理 \ref{topologies-lemma-morphism-big-small-etale}。

\medskip\noindent
因此，概形的平展上同调与相应代数空间的平展上同调完全没有区别。
沿概形态射的高阶正像也是如此。
事实上，若 $f : X \to Y$ 是 $S$ 上代数空间的一个（由概形）可表态射，
则 $X_\etale$ 上层 $\mathcal{F}$ 的高阶正像 $R^if_*\mathcal{F}$
可以在 $Y$ 上平展局部地计算
（《站点上的上同调》引理 \ref{sites-cohomology-lemma-higher-direct-images}），
故计算与证明常可化归到 $Y$ 和 $X$ 都是概形的情形。

\medskip\noindent
本章将直接使用上述事实而不再特别说明。
对其他拓扑同样成立；这里把关于上同调的结论明确表述为一个引理。

\begin{lemma}
\label{spaces-more-cohomology-lemma-compare-cohomology-other-topologies}
设 $S$ 为概形。设 $\tau \in \{\etale, fppf, ph\}$（此处再添加更多）。
% P10-E0148
包含函子
$$
(\Sch/S)_\tau \longrightarrow (\textit{Spaces}/S)_\tau
$$
是特殊余连续函子
（《站点》定义 \ref{sites-definition-special-cocontinuous-functor}），
因而给出拓扑斯的等同。
\end{lemma}

\begin{proof}
《站点》引理 \ref{sites-lemma-equivalence} 的条件立即得到验证，
因为该函子是全忠实的，并且每个代数空间都有由概形组成的平展覆盖。
\end{proof}







\section{真基变换}
\label{spaces-more-cohomology-section-proper-base-change}

\noindent
稍加论证，即可由概形的真基变换定理和 Chow 引理推出代数空间的真基变换定理。

\begin{lemma}
\label{spaces-more-cohomology-lemma-surjective-proper}
设 $S$ 为概形。设 $f : Y \to X$ 为 $S$ 上代数空间的满真态射。
设 $\mathcal{F}$ 为 $X_\etale$ 上的层。则
$\mathcal{F} \to f_*f^{-1}\mathcal{F}$ 是单射，其像是两个映射
$f_*f^{-1}\mathcal{F} \to g_*g^{-1}\mathcal{F}$ 的等化子，
其中 $g$ 是结构态射 $g : Y \times_X Y \to X$。
\end{lemma}

\begin{proof}
对 $S$ 上代数空间的任意满态射 $f : Y \to X$，映射
$\mathcal{F} \to f_*f^{-1}\mathcal{F}$ 都是单射。
事实上，若 $\overline{x}$ 是 $X$ 的几何点，则选择 $Y$ 中位于
$\overline{x}$ 上方的几何点 $\overline{y}$，并考虑
$$
\mathcal{F}_{\overline{x}} \to
(f_*f^{-1}\mathcal{F})_{\overline{x}} \to
(f^{-1}\mathcal{F})_{\overline{y}} = \mathcal{F}_{\overline{x}}
$$
最后一个等式参见《空间的性质》引理
\ref{spaces-properties-lemma-stalk-pullback}。

\medskip\noindent
第二个陈述在 $X$ 上对平展拓扑是局部的，故可以并且确实假设
$Y$ 是仿射概形。

\medskip\noindent
选择满真态射 $Z \to Y$，其中 $Z$ 是概形；参见《空间的上同调》引理
\ref{spaces-cohomology-lemma-weak-chow}。
$Z \to X$ 的结论蕴含 $Y \to X$ 的结论。
由于 $Z \to X$ 是概形的满真态射，因而是 ph 覆盖
（《拓扑》引理 \ref{topologies-lemma-surjective-proper-ph}），
$Z \to X$ 的结论可由《平展上同调》引理
\ref{etale-cohomology-lemma-describe-pullback-pi-ph} 推出
（事实上，它在某种意义下与该引理等价）。
\end{proof}

\begin{lemma}
\label{spaces-more-cohomology-lemma-h0-proper-over-henselian-pair}
设 $(A, I)$ 为亨泽尔对。设 $X$ 为 $A$ 上的代数空间，且结构态射
$f : X \to \Spec(A)$ 为真态射。令 $i : X_0 \to X$ 为
$X \times_{\Spec(A)} \Spec(A/I)$ 的包含映射。
则对 $X_\etale$ 上的任意层 $\mathcal{F}$，有
$\Gamma(X, \mathcal{F}) = \Gamma(X_0, i^{-1}\mathcal{F})$。
\end{lemma}

\begin{proof}
选择满真态射 $Y \to X$，其中 $Y$ 是概形；参见《空间的上同调》引理
\ref{spaces-cohomology-lemma-weak-chow}。考虑图表
$$
\xymatrix{
\Gamma(X_0, \mathcal{F}_0) \ar[r] &
\Gamma(Y_0, \mathcal{G}_0) \ar@<1ex>[r] \ar@<-1ex>[r] &
\Gamma((Y \times_X Y)_0, \mathcal{H}_0) \\
\Gamma(X, \mathcal{F}) \ar[r] \ar[u] &
\Gamma(Y, \mathcal{G}) \ar@<1ex>[r] \ar@<-1ex>[r] \ar[u] &
\Gamma(Y \times_X Y, \mathcal{H}) \ar[u]
}
$$
% P10-E0149
这里 $\mathcal{G}$ 和 $\mathcal{H}$ 分别是 $\mathcal{F}$ 在 $Y$ 和
$Y \times_X Y$ 上的拉回，下标 $0$ 表示基变换到 $\Spec(A/I)$。
由概形的情形（《平展上同调》引理
\ref{etale-cohomology-lemma-h0-proper-over-henselian-pair}），
中间和右侧的竖直箭头均为双射。
再由引理 \ref{spaces-more-cohomology-lemma-surjective-proper}，左侧箭头也是双射。
\end{proof}

\begin{lemma}
\label{spaces-more-cohomology-lemma-h0-proper-over-henselian-local}
设 $A$ 为亨泽尔局部环。设 $X$ 为 $A$ 上的代数空间，且
$f : X \to \Spec(A)$ 是真态射。令 $X_0 \subset X$ 为 $f$ 在闭点上方的纤维。
则对 $X_\etale$ 上的任意层 $\mathcal{F}$，有
$\Gamma(X, \mathcal{F}) = \Gamma(X_0, \mathcal{F}|_{X_0})$。
\end{lemma}

\begin{proof}
这是引理 \ref{spaces-more-cohomology-lemma-h0-proper-over-henselian-pair} 的特例。
\end{proof}

\begin{lemma}
\label{spaces-more-cohomology-lemma-proper-base-change-f-star}
% P10-E0150
设 $S$ 为概形。设 $f : X \to Y$ 和 $g : Y' \to Y$ 为
$S$ 上代数空间的态射，并假设 $f$ 为真态射。
令 $X' = Y' \times_Y X$，投影为 $f' : X' \to Y'$ 和 $g' : X' \to X$。
设 $\mathcal{F}$ 为 $X_\etale$ 上的任意层。则
$g^{-1}f_*\mathcal{F} = f'_*(g')^{-1}\mathcal{F}$。
\end{lemma}

\begin{proof}
该问题在 $Y'$ 上对平展拓扑是局部的。选择概形 $V$ 和满平展态射
$V \to Y$。再选择概形 $V'$ 和满平展态射
$V' \to V \times_Y Y'$。于是可分别用 $V'$ 和 $V$ 替换 $Y'$ 和 $Y$，
故可假设 $Y$ 和 $Y'$ 都是概形。随后可在 $Y$ 和 $Y'$ 上扎里斯基局部地工作，
因而还可假设 $Y$ 和 $Y'$ 都是仿射概形。

\medskip\noindent
假设 $Y$ 和 $Y'$ 是仿射概形。选择满真态射
$h_1 : X_1 \to X$，其中 $X_1$ 是概形；参见《空间的上同调》引理
\ref{spaces-cohomology-lemma-weak-chow}。
令 $X_2 = X_1 \times_X X_1$，并将结构态射记为
$h_2 : X_2 \to X$。注意它是概形。
由概形的情形（《平展上同调》引理
\ref{etale-cohomology-lemma-proper-base-change-f-star}），
可知对下列笛卡尔图表该引理成立：
$$
\vcenter{
\xymatrix{
X'_1 \ar[r] \ar[d] & X_1 \ar[d] \\
Y' \ar[r] & Y
}
}
\quad\text{和}\quad
\vcenter{
\xymatrix{
X'_2 \ar[r] \ar[d] & X_2 \ar[d] \\
Y' \ar[r] & Y
}
}
$$
以及层 $\mathcal{F}_i = (X_i \to X)^{-1}\mathcal{F}$。
由引理 \ref{spaces-more-cohomology-lemma-surjective-proper}，有正合序列
$0 \to \mathcal{F} \to h_{1, *}\mathcal{F}_1 \to h_{2, *}\mathcal{F}_2$。
由于 $X'_2 = X'_1 \times_{X'} X'_1$，对 $(g')^{-1}\mathcal{F}$ 也同样成立。
% P10-E0151
故可断定该引理成立（略去一些细节）。
\end{proof}

\noindent
设 $S$ 为概形。设 $f : Y \to X$ 为 $S$ 上代数空间的态射。
% P10-E0152
设 $\overline{x} : \Spec(k) \to S$ 为几何点。
$f$ 在 $\overline{x}$ 处的纤维是 $\Spec(k)$ 上的代数空间
$Y_{\overline{x}} = \Spec(k) \times_{\overline{x}, X} Y$。
若 $\mathcal{F}$ 是 $Y_\etale$ 上的层，则用
$\mathcal{F}_{\overline{x}} = p^{-1}\mathcal{F}$
表示 $\mathcal{F}$ 在 $(Y_{\overline{x}})_\etale$ 上的拉回，
其中 $p : Y_{\overline{x}} \to Y$ 是投影。
下文将考察集合 $\Gamma(Y_{\overline{x}}, \mathcal{F}_{\overline{x}})$。

\begin{lemma}
\label{spaces-more-cohomology-lemma-proper-pushforward-stalk}
设 $S$ 为概形。设 $f : Y \to X$ 为 $S$ 上代数空间的真态射，
并设 $\overline{x} \to X$ 为几何点。
对 $Y_\etale$ 上的任意层 $\mathcal{F}$，典范映射
$$
(f_*\mathcal{F})_{\overline{x}} \longrightarrow
\Gamma(Y_{\overline{x}}, \mathcal{F}_{\overline{x}})
$$
是双射。
\end{lemma}

\begin{proof}
这是引理 \ref{spaces-more-cohomology-lemma-proper-base-change-f-star} 的特例。
\end{proof}

\begin{theorem}
\label{spaces-more-cohomology-theorem-proper-base-change}
设 $S$ 为概形。设
$$
\xymatrix{
X' \ar[r]_{g'} \ar[d]_{f'} & X \ar[d]^f \\
Y' \ar[r]^g & Y
}
$$
为 $S$ 上代数空间的笛卡尔方块。假设 $f$ 为真态射。
设 $\mathcal{F}$ 为 $X_\etale$ 上的阿贝尔挠层。则基变换映射
$$
g^{-1}Rf_*\mathcal{F} \longrightarrow Rf'_*(g')^{-1}\mathcal{F}
$$
是同构。
\end{theorem}

\begin{proof}
本证明重复概形真基变换定理证明中的若干论证。
更多细节参见《平展上同调》第
\ref{etale-cohomology-section-proper-base-change} 节。

\medskip\noindent
该陈述在 $Y'$ 和 $Y$ 上对平展拓扑是局部的，故可假设
$Y$ 和 $Y'$ 都是仿射概形。特别地，这证明了 $f$ 可表时的定理
（下文将使用这一点）。

\medskip\noindent
对每个 $n \geq 1$，令 $\mathcal{F}[n]$ 为 $\mathcal{F}$ 中被 $n$
零化的截面所成的子层。则 $\mathcal{F} = \colim \mathcal{F}[n]$。
由《空间的上同调》引理
\ref{spaces-cohomology-lemma-colimit-cohomology}，函子
$g^{-1}R^pf_*$ 和 $R^pf'_*(g')^{-1}$ 与滤过余极限可交换。
因此只需在 $\mathcal{F}$ 被某个 $n$ 零化时证明该定理。

\medskip\noindent
取由 $\mathbf{Z}/n\mathbf{Z}$-模的内射层组成的解消
$\mathcal{F} \to \mathcal{I}^\bullet$。注意，由引理
$g^{-1}f_*\mathcal{I}^\bullet = f'_*(g')^{-1}\mathcal{I}^\bullet$
\ref{spaces-more-cohomology-lemma-proper-base-change-f-star} 有上述等式。
应用 Leray 无环性引理（《导出范畴》引理
\ref{derived-lemma-leray-acyclicity}），可知只需证明：
对 $p > 0$ 和 $m \in \mathbf{Z}$，有
$R^pf'_*(g')^{-1}\mathcal{I}^m = 0$。

\medskip\noindent
选择满真态射 $h : Z \to X$，其中 $Z$ 是概形；
参见《空间的上同调》引理 \ref{spaces-cohomology-lemma-weak-chow}。
选择单射 $h^{-1}\mathcal{I}^m \to \mathcal{J}$，其中 $\mathcal{J}$
是 $Z_\etale$ 上 $\mathbf{Z}/n\mathbf{Z}$-模的内射层。
由于 $h$ 是满射，映射 $\mathcal{I}^m \to h_*\mathcal{J}$ 是单射
（参见引理 \ref{spaces-more-cohomology-lemma-surjective-proper}）。
又因 $\mathcal{I}^m$ 内射，$\mathcal{I}^m$ 是 $h_*\mathcal{J}$ 的直和项。
所以只需对 $h_*\mathcal{J}$ 证明所需的消失性。

\medskip\noindent
% P10-E0153
用 $h'$ 表示 h 沿 $g$ 的基变换，并用 $g'' : Z' \to Z$ 表示投影。
存在谱序列
$$
E_2^{p, q} = R^pf'_* R^qh'_* (g'')^{-1}\mathcal{J}
$$
收敛到 $R^{p + q}(f' \circ h')_*(g'')^{-1}\mathcal{J}$。
由于 $h$ 和 $f \circ h$ 都（由概形）可表，我们知道所需结果对它们成立。
因此在该谱序列中，当 $q > 0$ 时 $E_2^{p, q} = 0$，并且当
$p + q > 0$ 时
$R^{p + q}(f' \circ h')_*(g'')^{-1}\mathcal{J} = 0$。
由此当 $p > 0$ 时 $E_2^{p, 0} = 0$。现在
$$
E_2^{p, 0} = R^pf'_* h'_* (g'')^{-1}\mathcal{J} =
R^pf'_* (g')^{-1}h_*\mathcal{J}
$$
其中等式由引理 \ref{spaces-more-cohomology-lemma-proper-base-change-f-star} 得到。证明完成。
\end{proof}

\begin{lemma}
\label{spaces-more-cohomology-lemma-proper-base-change}
设 $S$ 为概形。设
$$
\xymatrix{
X' \ar[r]_{g'} \ar[d]_{f'} & X \ar[d]^f \\
Y' \ar[r]^g & Y
}
$$
为 $S$ 上代数空间的笛卡尔方块。假设 $f$ 为真态射。
设 $E \in D^+(X_\etale)$ 的上同调层均为挠层。
则基变换映射 $g^{-1}Rf_*E \to Rf'_*(g')^{-1}E$ 是同构。
\end{lemma}

\begin{proof}
利用谱序列
$$
E_2^{p, q} = R^pf_*H^q(E)
\quad\text{和}\quad
{E'}_2^{p, q} = R^pf'_*(g')^{-1}H^q(E)
$$
它们分别收敛到 $R^nf_*E$ 和 $R^nf'_*(g')^{-1}E$，
即可由真基变换定理（定理 \ref{spaces-more-cohomology-theorem-proper-base-change}）简单推出结论。
这些谱序列构造于《导出范畴》引理
\ref{derived-lemma-two-ss-complex-functor}。略去一些细节。
\end{proof}

\begin{lemma}
\label{spaces-more-cohomology-lemma-proper-base-change-stalk}
设 $S$ 为概形。设 $f : X \to Y$ 为代数空间的真态射，
并设 $\overline{y} \to Y$ 为几何点。
\begin{enumerate}
\item 对 $X_\etale$ 上的挠阿贝尔层 $\mathcal{F}$，有
$(R^nf_*\mathcal{F})_{\overline{y}} =
H^n_\etale(X_{\overline{y}}, \mathcal{F}_{\overline{y}})$.
\item 对上同调层均为挠层的 $E \in D^+(X_\etale)$，有
$(R^nf_*E)_{\overline{y}} = H^n_\etale(X_{\overline{y}}, E_{\overline{y}})$.
\end{enumerate}
\end{lemma}

\begin{proof}
在陈述中，$\mathcal{F}_{\overline{y}}$ 表示 $\mathcal{F}$ 在
$X_{\overline{y}} = \overline{y} \times_Y X$ 上的拉回。
由于沿 $\overline{y} \to Y$ 拉回会给出 $\mathcal{F}$ 的茎，
引理的第一个陈述是定理 \ref{spaces-more-cohomology-theorem-proper-base-change} 的特例，
第二个陈述是引理 \ref{spaces-more-cohomology-lemma-proper-base-change} 的特例。
\end{proof}

\begin{lemma}
\label{spaces-more-cohomology-lemma-base-change-separably-closed}
设 $k'/k$ 为可分闭域的扩张。设 $X$ 为 $k$ 上的真代数空间，
$\mathcal{F}$ 为 $X$ 上的挠阿贝尔层。则对 $q \geq 0$，映射
$H^q_\etale(X, \mathcal{F}) \to
H^q_\etale(X_{k'}, \mathcal{F}|_{X_{k'}})$ 是同构。
\end{lemma}

\begin{proof}
这是定理 \ref{spaces-more-cohomology-theorem-proper-base-change} 的特例。
\end{proof}










\section{比较大拓扑斯与小拓扑斯}
\label{spaces-more-cohomology-section-compare}

\noindent
设 $S$ 为概形，$X$ 为 $S$ 上的代数空间。
在《空间上的拓扑》引理
\ref{spaces-topologies-lemma-at-the-bottom-etale} 中，我们引入了比较态射
$\pi_X : (\textit{Spaces}/X)_\etale \to X_{spaces, \etale}$ 以及
$i_X : \Sh(X_\etale) \to \Sh((\textit{Spaces}/X)_\etale)$
作为拓扑斯态射，有 $\pi_X \circ i_X = \text{id}$，并且
$\pi_{X, *} = i_X^{-1}$。
更一般地，若 $f : Y \to X$ 是 $(\textit{Spaces}/X)_\etale$ 的对象，
则存在态射
$i_f : \Sh(Y_\etale) \to \Sh((\textit{Spaces}/X)_\etale)$
使得 $f_{small} = \pi_X \circ i_f$；参见《空间上的拓扑》引理
\ref{spaces-topologies-lemma-put-in-T-etale} 和
\ref{spaces-topologies-lemma-morphism-big-small-etale}。
在《空间上的拓扑》评注
\ref{spaces-topologies-remark-change-topologies-ringed} 中，
我们把它们扩张为赋环站点的态射
$$
\pi_X :
((\textit{Spaces}/X)_\etale, \mathcal{O})
\to
(X_{spaces, \etale}, \mathcal{O}_X)
$$
以及赋环拓扑斯的态射
$$
i_X :
(\Sh(X_\etale), \mathcal{O}_X)
\to
(\Sh((\textit{Spaces}/X)_\etale), \mathcal{O})
$$
和
$$
i_f :
(\Sh(Y_\etale), \mathcal{O}_Y)
\to
(\Sh((\textit{Spaces}/X)_\etale, \mathcal{O}))
$$
注意，限制 $i_X^{-1} = \pi_{X, *}$（参见《拓扑》定义
\ref{topologies-definition-restriction-small-etale}）把 $\mathcal{O}$
变为 $\mathcal{O}_X$。类似地，$i_f^{-1}$ 把 $\mathcal{O}$ 变为
$\mathcal{O}_Y$。参见《空间上的拓扑》评注
\ref{spaces-topologies-remark-change-topologies-ringed}。
因此，对 $(\textit{Spaces}/X)_\etale$ 上的任意 $\mathcal{O}$-模
$\mathcal{F}$，有 $i_X^*\mathcal{F} = i_X^{-1}\mathcal{F}$ 和
$i_f^*\mathcal{F} = i_f^{-1}\mathcal{F}$。特别地，$i_X^*$ 和 $i_f^*$
都是正合函子。函子 $i_X^*$ 常记作
$\mathcal{F} \mapsto \mathcal{F}|_{X_\etale}$
（这与《空间上的拓扑》定义
\ref{spaces-topologies-definition-restriction-small-etale} 中的记号不冲突）。

\begin{lemma}
\label{spaces-more-cohomology-lemma-describe-pullback}
设 $S$ 为概形，$X$ 为 $S$ 上的代数空间。
设 $\mathcal{F}$ 为 $X_\etale$ 上的层。则 $\pi_X^{-1}\mathcal{F}$
由下列规则给出：
$$
(\pi_X^{-1}\mathcal{F})(Y) = \Gamma(Y_\etale, f_{small}^{-1}\mathcal{F})
$$
其中 $f : Y \to X$ 是 $(\textit{Spaces}/X)_\etale$ 的对象。
% P10-E0154
此外，$\pi_Y^{-1}\mathcal{F}$ 对光滑、syntomic、fppf、fpqc 和 ph 覆盖
满足层条件。
\end{lemma}

\begin{proof}
由于拉回具有传递性且 $f_{small} = \pi_X \circ i_f$（见上文），有
$i_f^{-1} \pi_X^{-1}\mathcal{F} = f_{small}^{-1}\mathcal{F}$。
这说明 $\pi_X^{-1}$ 具有引理中给出的描述。

\medskip\noindent
为证明 $\pi_X^{-1}\mathcal{F}$ 对 ph 拓扑是层，
由《空间上的拓扑》引理
\ref{spaces-topologies-lemma-characterize-sheaf}，只需证明：
对 $X$ 上代数空间的满真态射 $V \to U$，
$(\pi_X^{-1}\mathcal{F})(U)$ 是两个映射
$(\pi_X^{-1}\mathcal{F})(V) \to
(\pi_X^{-1}\mathcal{F})(V \times_U V)$ 的等化子。
这已由引理 \ref{spaces-more-cohomology-lemma-surjective-proper} 得到。

\medskip\noindent
光滑、syntomic 和 fppf 覆盖的情形，由 ph 覆盖的情形和
《空间上的拓扑》引理
\ref{spaces-topologies-lemma-zariski-etale-smooth-syntomic-fppf-ph} 推出。

\medskip\noindent
设 $\mathcal{U} = \{U_i \to U\}_{i \in I}$ 为 $X$ 上代数空间的 fpqc 覆盖。
设截面 $s_i \in (\pi_X^{-1}\mathcal{F})(U_i)$ 在
$U_i \times_U U_j$ 上彼此相同。需要证明存在唯一的
$s \in (\pi_X^{-1}\mathcal{F})(U)$，其在 $U_i$ 上限制为 $s_i$。
情形 I：$U$ 和 $U_i$ 都是概形。此时结论由《平展上同调》引理
\ref{etale-cohomology-lemma-describe-pullback} 得到。
情形 II：$U$ 是概形。选择满平展态射 $T_i \to U_i$，其中 $T_i$ 是概形。
则 $\mathcal{T} = \{T_i \to U\}$ 是由概形组成的 fpqc 覆盖，
由情形 I，结论对 $\mathcal{T}$ 成立。略去由此推出
$\mathcal{U}$ 的结论的验证。
情形 III：一般情形。取满平展态射 $W \to U$，其中 $W$ 是概形。
则 $\mathcal{W} = \{U_i \times_U W \to W\}$ 是概形 $W$ 的
fpqc 覆盖（覆盖对象为代数空间）。
% P10-E0155
由情形 II，结论对 $\mathcal{W}$ 成立。略去由此推出
$\mathcal{U}$ 的结论的验证。
\end{proof}

\begin{lemma}
\label{spaces-more-cohomology-lemma-compare-injectives}
设 $S$ 为概形。设 $Y \to X$ 为 $(\textit{Spaces}/S)_\etale$ 中的态射。
\begin{enumerate}
\item 若 $\mathcal{I}$ 在
$\textit{Ab}((\textit{Spaces}/X)_\etale)$ 中内射，则
\begin{enumerate}
\item $i_f^{-1}\mathcal{I}$ 在 $\textit{Ab}(Y_\etale)$ 中内射；
\item $\mathcal{I}|_{X_\etale}$ 在 $\textit{Ab}(X_\etale)$ 中内射。
\end{enumerate}
\item 若 $\mathcal{I}^\bullet$ 是 $\textit{Ab}((\textit{Spaces}/X)_\etale)$
中的 K-内射复形，则
\begin{enumerate}
\item $i_f^{-1}\mathcal{I}^\bullet$ 是 $\textit{Ab}(Y_\etale)$ 中的
K-内射复形；
\item $\mathcal{I}^\bullet|_{X_\etale}$ 是 $\textit{Ab}(X_\etale)$ 中的
K-内射复形。
\end{enumerate}
\end{enumerate}
相应的模版本不成立。
\end{lemma}

\begin{proof}
由限制函子 $\pi_{X, *} = i_X^{-1}$ 是正合函子 $\pi_X^{-1}$ 的右伴随，
形式地推出 (1)(b) 和 (2)(b)。参见《同调代数》引理
\ref{homology-lemma-adjoint-preserve-injectives} 和《导出范畴》引理
\ref{derived-lemma-adjoint-preserve-K-injectives}。

\medskip\noindent
可用两种方法证明 (1)(a) 和 (2)(a)。第一种证明：利用
$i_f^{-1}$ 是正合函子 $i_{f, !}$ 的右伴随。
对集合层，该函子构造于《拓扑》引理
\ref{topologies-lemma-put-in-T-etale}；对阿贝尔层，则构造于
《站点上的模》引理 \ref{sites-modules-lemma-g-shriek-adjoint}。
《站点上的模》引理
\ref{sites-modules-lemma-exactness-lower-shriek} 证明了它是正合的。
第二种证明：利用《拓扑》引理
\ref{topologies-lemma-morphism-big-small-etale} 所证明的
$i_f = i_Y \circ f_{big}$。由于 $f_{big}$ 是局部化，沿它的拉回保持内射对象和
K-内射对象；参见《站点上的上同调》引理
\ref{sites-cohomology-lemma-cohomology-of-open} 和
\ref{sites-cohomology-lemma-restrict-K-injective-to-open}。
然后把已经证明的 (1)(b) 和 (2)(b) 应用于函子 $i_Y^{-1}$ 即得结论。

\medskip\noindent
模情形的反例参见《平展上同调》引理
\ref{etale-cohomology-lemma-compare-injectives}。
\end{proof}

\noindent
设 $S$ 为概形。设 $f : Y \to X$ 为 $S$ 上代数空间的态射。
《空间上的拓扑》引理
\ref{spaces-topologies-lemma-morphism-big-small-etale} (3) 的交换图
诱导出赋环站点的交换图
$$
\xymatrix{
(Y_{spaces, \etale}, \mathcal{O}_Y) \ar[d]_{f_{spaces, \etale}} &
((\textit{Spaces}/Y)_\etale, \mathcal{O}) \ar[d]^{f_{big}} \ar[l]^{\pi_Y} \\
(X_{spaces, \etale}, \mathcal{O}_X) &
((\textit{Spaces}/X)_\etale, \mathcal{O}) \ar[l]_{\pi_X}
}
$$
将 $f_{small}^\sharp$、$f_{big}^\sharp$、$\pi_X^\sharp$ 和
$\pi_Y^\sharp$ 的定义展开即可容易看出这一点。特别地，这意味着
\begin{equation}
\label{spaces-more-cohomology-equation-compare-big-small}
(f_{big, *}\mathcal{F})|_{X_\etale} =
f_{small, *}(\mathcal{F}|_{Y_\etale})
\end{equation}
对 $(\textit{Spaces}/Y)_\etale$ 上的任意层 $\mathcal{F}$ 都成立；
若 $\mathcal{F}$ 是 $\mathcal{O}$-模层，则
(\ref{spaces-more-cohomology-equation-compare-big-small}) 是 $X_\etale$ 上
$\mathcal{O}_X$-模的同构。

\begin{lemma}
\label{spaces-more-cohomology-lemma-compare-higher-direct-image}
设 $S$ 为概形。设 $f : Y \to X$ 为 $S$ 上代数空间的态射。
\begin{enumerate}
\item 对 $D((\textit{Spaces}/Y)_\etale)$ 中的 $K$，有
$
(Rf_{big, *}K)|_{X_\etale} = Rf_{small, *}(K|_{Y_\etale})
$
在 $D(X_\etale)$ 中成立。
\item 对 $D((\textit{Spaces}/Y)_\etale, \mathcal{O})$ 中的 $K$，有
$
(Rf_{big, *}K)|_{X_\etale} = Rf_{small, *}(K|_{Y_\etale})
$
在 $D(\textit{Mod}(X_\etale, \mathcal{O}_X))$ 中成立。
\end{enumerate}
更一般地，设 $g : X' \to X$ 为 $(\textit{Spaces}/X)_\etale$ 的对象。
考虑纤维积
$$
\xymatrix{
Y' \ar[r]_{g'} \ar[d]_{f'} & Y \ar[d]^f \\
X' \ar[r]^g & X
}
$$
则
\begin{enumerate}
\item[(3)] 对 $D((\textit{Spaces}/Y)_\etale)$ 中的 $K$，有
$i_g^{-1}(Rf_{big, *}K) = Rf'_{small, *}(i_{g'}^{-1}K)$
在 $D(X'_\etale)$ 中成立。
\item[(4)] 对 $D((\textit{Spaces}/Y)_\etale, \mathcal{O})$ 中的 $K$，有
$i_g^*(Rf_{big, *}K) = Rf'_{small, *}(i_{g'}^*K)$
在 $D(\textit{Mod}(X'_\etale, \mathcal{O}_{X'}))$ 中成立。
\item[(5)] 对 $D((\textit{Spaces}/Y)_\etale)$ 中的 $K$，有
$g_{big}^{-1}(Rf_{big, *}K) = Rf'_{big, *}((g'_{big})^{-1}K)$
在 $D((\textit{Spaces}/X')_\etale)$ 中成立。
\item[(6)] 对 $D((\textit{Spaces}/Y)_\etale, \mathcal{O})$ 中的 $K$，有
$g_{big}^*(Rf_{big, *}K) = Rf'_{big, *}((g'_{big})^*K)$
在 $D(\textit{Mod}(X'_\etale, \mathcal{O}_{X'}))$ 中成立。
\end{enumerate}
\end{lemma}

\begin{proof}
选择表示 $K$ 的阿贝尔层 K-内射复形后，(1) 由引理
\ref{spaces-more-cohomology-lemma-compare-injectives} 和
(\ref{spaces-more-cohomology-equation-compare-big-small}) 推出。

\medskip\noindent
选择表示 $K$ 的阿贝尔层 K-内射复形后，(3) 由引理
\ref{spaces-more-cohomology-lemma-compare-injectives} 和《拓扑》引理
\ref{topologies-lemma-morphism-big-small-cartesian-diagram-etale} 推出。

\medskip\noindent
(5) 即《站点上的上同调》引理
\ref{sites-cohomology-lemma-localize-cartesian-square}。

\medskip\noindent
(6) 即《站点上的上同调》引理
\ref{sites-cohomology-lemma-localize-cartesian-square-modules}。

\medskip\noindent
(2) 可作如下证明。上文已看到，作为赋环站点的态射，有
$\pi_X \circ f_{big} = f_{small} \circ \pi_Y$。因此由
$R\pi_{X, *} \circ Rf_{big, *} = Rf_{small, *} \circ R\pi_{Y, *}$
《站点上的上同调》引理
\ref{sites-cohomology-lemma-derived-pushforward-composition} 得到上述等式。
由于限制函子 $\pi_{X, *}$ 和 $\pi_{Y, *}$ 都正合，结论成立。

\medskip\noindent
(4) 由 (6) 以及把 (2) 应用于 $f' : Y' \to X'$ 得到。
\end{proof}

\noindent
设 $S$ 为概形，$X$ 为 $S$ 上的代数空间。
设 $\mathcal{H}$ 为 $(\textit{Spaces}/X)_\etale$ 上的阿贝尔层。
回忆，$H^n_\etale(U, \mathcal{H})$ 表示 $\mathcal{H}$ 在
$(\textit{Spaces}/X)_\etale$ 的对象 $U$ 上的上同调。

\begin{lemma}
\label{spaces-more-cohomology-lemma-compare-cohomology}
设 $S$ 为概形。设 $f : Y \to X$ 为 $S$ 上代数空间的态射。则
\begin{enumerate}
\item 对 $D(X_\etale)$ 中的 $K$，有
$H^n_\etale(X, \pi_X^{-1}K) = H^n(X_\etale, K)$.
\item 对 $D(X_\etale, \mathcal{O}_X)$ 中的 $K$，有
$H^n_\etale(X, L\pi_X^*K) = H^n(X_\etale, K)$.
\item 对 $D(X_\etale)$ 中的 $K$，有
$H^n_\etale(Y, \pi_X^{-1}K) = H^n(Y_\etale, f_{small}^{-1}K)$.
\item 对 $D(X_\etale, \mathcal{O}_X)$ 中的 $K$，有
$H^n_\etale(Y, L\pi_X^*K) = H^n(Y_\etale, Lf_{small}^*K)$.
\item 对 $D((\textit{Spaces}/X)_\etale)$ 中的 $M$，有
$H^n_\etale(Y, M) = H^n(Y_\etale, i_f^{-1}M)$.
\item 对 $D((\textit{Spaces}/X)_\etale, \mathcal{O})$ 中的 $M$，有
$H^n_\etale(Y, M) = H^n(Y_\etale, i_f^*M)$.
\end{enumerate}
\end{lemma}

\begin{proof}
为证明 (5)，用阿贝尔层的 K-内射复形表示 $M$，
应用引理 \ref{spaces-more-cohomology-lemma-compare-injectives} 并展开定义即可。
由于 $i_f^{-1}\pi_X^{-1} = f_{small}^{-1}$，(3) 由此推出；
(1) 是 (3) 的特例。

\medskip\noindent
(6) 由十分一般的《站点上的上同调》引理
\ref{sites-cohomology-lemma-pullback-same-cohomology} 推出。
又因 $Lf_{small}^* = i_f^* \circ L\pi_X^*$，(4) 随之成立；
(2) 是 (4) 的特例。
\end{proof}

\begin{lemma}
\label{spaces-more-cohomology-lemma-cohomological-descent-etale}
设 $S$ 为概形，$X$ 为 $S$ 上的代数空间。
对 $K \in D(X_\etale)$，映射
$$
K \longrightarrow R\pi_{X, *}\pi_X^{-1}K
$$
是同构，其中 $\pi_X : \Sh((\textit{Spaces}/X)_\etale) \to
\Sh(X_\etale)$ 如上所述。
\end{lemma}

\begin{proof}
这是因为 $\pi_X^{-1}$ 和 $\pi_{X, *} = i_X^{-1}$ 都是正合函子，
且复合 $\pi_{X, *} \circ \pi_X^{-1}$ 是恒等函子。
\end{proof}

\begin{lemma}
\label{spaces-more-cohomology-lemma-compare-higher-direct-image-proper}
设 $S$ 为概形。设 $f : Y \to X$ 为 $S$ 上代数空间的真态射。则有
\begin{enumerate}
\item $\pi_X^{-1} \circ f_{small, *} = f_{big, *} \circ \pi_Y^{-1}$
作为函子 $\Sh(Y_\etale) \to \Sh((\textit{Spaces}/X)_\etale)$ 的等式；
\item $\pi_X^{-1}Rf_{small, *}K = Rf_{big, *}\pi_Y^{-1}K$
对上同调层均为挠层的 $D^+(Y_\etale)$ 中对象 $K$ 成立；以及
\item $\pi_X^{-1}Rf_{small, *}K = Rf_{big, *}\pi_Y^{-1}K$
若 $f$ 有限，则对 $D(Y_\etale)$ 中所有 $K$ 成立。
\end{enumerate}
\end{lemma}

\begin{proof}
(1) 的证明。设 $\mathcal{F}$ 为 $Y_\etale$ 上的层。
设 $g : X' \to X$ 为 $(\textit{Spaces}/X)_\etale$ 的对象。
考虑纤维积
$$
\xymatrix{
Y' \ar[r]_{f'} \ar[d]_{g'} & X' \ar[d]^g \\
Y \ar[r]^f & X
}
$$
则有
$$
(f_{big, *}\pi_Y^{-1}\mathcal{F})(X') =
(\pi_Y^{-1}\mathcal{F})(Y') =
((g'_{small})^{-1}\mathcal{F})(Y')  =
(f'_{small, *}(g'_{small})^{-1}\mathcal{F})(X')
$$
其中第二个等式由引理 \ref{spaces-more-cohomology-lemma-describe-pullback} 得到。另一方面，
$$
(\pi_X^{-1}f_{small, *}\mathcal{F})(X') =
(g_{small}^{-1}f_{small, *}\mathcal{F})(X')
$$
这同样由引理 \ref{spaces-more-cohomology-lemma-describe-pullback} 得到。
因此，由集合层的真基变换（引理
\ref{spaces-more-cohomology-lemma-proper-base-change-f-star}），这两个集合典范同构。
该同构与限制映射相容，并给出同构
$\pi_X^{-1}f_{small, *}\mathcal{F} = f_{big, *}\pi_Y^{-1}\mathcal{F}$，
从而给出函子同构
$\pi_X^{-1} \circ f_{small, *} = f_{big, *} \circ \pi_Y^{-1}$。

\medskip\noindent
(2) 的证明。对 $D(Y_\etale)$ 中任意 $K$，存在典范基变换映射
$\pi_X^{-1}Rf_{small, *}K \to Rf_{big, *}\pi_Y^{-1}K$
参见《站点上的上同调》评注
\ref{sites-cohomology-remark-base-change}。
% P10-E0157
为证明它是同构，只需证明：对 $(\Sch/X)_\etale$ 中任意对象
$g : X' \to X$，沿
$i_g : \Sh(X'_\etale) \to \Sh((\Sch/X)_\etale)$ 拉回该基变换映射后所得映射
是同构。
% P10-E0156
令 $T', g', f'$ 如上一段所述。该基变换映射的拉回为
\begin{align*}
g_{small}^{-1}Rf_{small, *}K
& =
i_g^{-1}\pi_X^{-1}Rf_{small, *}K \\
& \to
i_g^{-1}Rf_{big, *}\pi_Y^{-1}K \\
& =
Rf'_{small, *}(i_{g'}^{-1}\pi_Y^{-1}K) \\
& =
Rf'_{small, *}((g'_{small})^{-1}K)
\end{align*}
这里使用了 $\pi_X \circ i_g = g_{small}$、
$\pi_Y \circ i_{g'} = g'_{small}$ 以及引理
\ref{spaces-more-cohomology-lemma-compare-higher-direct-image}。
若 $K$ 下有界且 $K$ 的上同调层均为挠层，则由真基变换定理
（引理 \ref{spaces-more-cohomology-lemma-proper-base-change}），该映射是同构。

\medskip\noindent
(3) 的证明。若 $f$ 有限，则函子 $f_{small, *}$ 和 $f_{big, *}$ 都正合。
对 $f_{small}$，这由《空间的上同调》引理
\ref{spaces-cohomology-lemma-finite-higher-direct-image-zero} 得到。
由于 $f$ 的任意基变换 $f'$ 也有限，引理
\ref{spaces-more-cohomology-lemma-compare-higher-direct-image} 的 (3) 表明 $f_{big, *}$ 也正合
（因为高阶导出函子为零）。故此情形由 (1) 推出。
\end{proof}







\section{比较 fppf 拓扑与平展拓扑}
\label{spaces-more-cohomology-section-fppf-etale}

\noindent
本节是《平展上同调》第
\ref{etale-cohomology-section-fppf-etale} 节的对应版本。

\medskip\noindent
设 $S$ 为概形，$X$ 为 $S$ 上的代数空间。
在范畴 $\textit{Spaces}/X$ 上考虑 fppf 拓扑和平展拓扑。恒等函子
$(\textit{Spaces}/X)_\etale \to (\textit{Spaces}/X)_{fppf}$
是连续的，并由《站点》命题
\ref{sites-proposition-get-morphism} 定义站点态射
$$
\epsilon_X :
(\textit{Spaces}/X)_{fppf} \longrightarrow (\textit{Spaces}/X)_\etale
$$
注意，$\epsilon_{X, *}$ 在底层预层上是恒等函子，而
$\epsilon_X^{-1}$ 把一个平展层送到其 fppf 层化。
考虑比较大、小平展站点的站点态射
$$
\pi_X : (\textit{Spaces}/X)_\etale \longrightarrow X_{spaces, \etale}
$$
参见第 \ref{spaces-more-cohomology-section-compare} 节。复合给出站点态射
$$
a_X = \pi_X \circ \epsilon_X :
(\textit{Spaces}/X)_{fppf}
\longrightarrow
X_{spaces, \etale}
$$
若 $\mathcal{H}$ 是 $(\textit{Spaces}/X)_{fppf}$ 上的阿贝尔层，
则用 $H^n_{fppf}(U, \mathcal{H})$ 表示 $\mathcal{H}$ 在
$(\textit{Spaces}/X)_{fppf}$ 的对象 $U$ 上的上同调。

\begin{lemma}
\label{spaces-more-cohomology-lemma-comparison-fppf-etale}
设 $S$ 为概形，$X$ 为 $S$ 上的代数空间。
\begin{enumerate}
\item 对 $\mathcal{F} \in \Sh(X_\etale)$，有
$\epsilon_{X, *}a_X^{-1}\mathcal{F} = \pi_X^{-1}\mathcal{F}$
以及 $a_{X, *}a_X^{-1}\mathcal{F} = \mathcal{F}$。
\item 对 $\mathcal{F} \in \textit{Ab}(X_\etale)$，当 $i > 0$ 时有
$R^i\epsilon_{X, *}(a_X^{-1}\mathcal{F}) = 0$。
\end{enumerate}
\end{lemma}

\begin{proof}
有 $a_X^{-1}\mathcal{F} = \epsilon_X^{-1} \pi_X^{-1}\mathcal{F}$。
由引理 \ref{spaces-more-cohomology-lemma-describe-pullback}，平展层
$\pi_X^{-1}\mathcal{F}$ 对 fppf 拓扑也是层，因此等于
$a_X^{-1}\mathcal{F}$（因为沿 $\epsilon_X$ 拉回就是 fppf 层化）。
还要回忆，$\epsilon_{X, *}$ 在底层预层上是恒等函子。
现在由引理 \ref{spaces-more-cohomology-lemma-describe-pullback} 中 $\pi_X^{-1}$ 的显式描述，
(1) 立即成立。

\medskip\noindent
我们将把 (2) 化归到概形情形来证明；参见《平展上同调》引理
\ref{etale-cohomology-lemma-V-C-all-n-etale-fppf} 的 (1)。
由于每个代数空间在平展局部都是概形，这“显然可行”。
下文给出细节，不过建议读者跳过本证明。

\medskip\noindent
对 $(\textit{Spaces}/X)_{fppf}$ 上的阿贝尔层 $\mathcal{H}$，
高阶正像 $R^p\epsilon_{X, *}\mathcal{H}$ 是
$(\textit{Spaces}/X)_\etale$ 上预层
$U \mapsto H^p_{fppf}(U, \mathcal{H})$ 的伴随层。
参见《站点上的上同调》引理
\ref{sites-cohomology-lemma-higher-direct-images}。
由于 $(\textit{Spaces}/X)_\etale$ 的每个对象都有由概形组成的覆盖，
只需证明：给定一个概形 $U/X$ 和
$\xi \in H^p_{fppf}(U, a_X^{-1}\mathcal{F})$，可找到平展覆盖
$\{U_i \to U\}$，使得 $\xi$ 在 $U_i$ 上的限制为零。我们有
\begin{align*}
H^p_{fppf}(U, a_X^{-1}\mathcal{F})
& =
H^p((\textit{Spaces}/U)_{fppf}, (a_X^{-1}\mathcal{F})|_{\textit{Spaces}/U}) \\
& =
H^p((\Sch/U)_{fppf}, (a_X^{-1}\mathcal{F})|_{\Sch/U})
\end{align*}
其中第二个等同由引理
\ref{spaces-more-cohomology-lemma-compare-cohomology-other-topologies} 得到，第一个则是关于限制的一般事实
（《站点上的上同调》引理
\ref{sites-cohomology-lemma-cohomology-of-open}）。
考察第一段以及概形情形的对应结果（《平展上同调》引理
\ref{etale-cohomology-lemma-describe-pullback-pi-fppf}），可知层
$(a_X^{-1}\mathcal{F})|_{\Sch/U}$ 与沿“概形版本的 $a_U$”拉回所得的层相同。
因此可找到平展覆盖 $\{U_i \to U\}$，使得对每个 $i$，该类在
$H^p((\Sch/U_i)_{fppf}, (a_X^{-1}\mathcal{F})|_{\Sch/U_i})$ 中消失；
参见《平展上同调》引理
\ref{etale-cohomology-lemma-V-C-all-n-etale-fppf}
（这里应使用的准确陈述是：$V_n$ 对所有 $n$ 成立，
这正是概形情形下 (2) 的陈述）。
再反向迁移（对 $U_i$ 使用与上文相同的公式），即可断定
$\xi$ 在 $U_i$ 上的限制为零，如所需。
\end{proof}

\noindent
概形情形中已经完成的艰苦工作现在告诉我们：
对于来自小平展站点的层，平展上同调与 fppf 上同调相同。

\begin{lemma}
\label{spaces-more-cohomology-lemma-cohomological-descent-etale-fppf}
设 $S$ 为概形，$X$ 为 $S$ 上的代数空间。
对 $K \in D^+(X_\etale)$，映射
$$
\pi_X^{-1}K \longrightarrow R\epsilon_{X, *}a_X^{-1}K
\quad\text{和}\quad
K \longrightarrow Ra_{X, *}a_X^{-1}K
$$
均为同构，其中
$a_X : \Sh((\textit{Spaces}/X)_{fppf}) \to \Sh(X_\etale)$ 如上所述。
\end{lemma}

\begin{proof}
只证明第二个陈述；第一个更容易，并可用完全相同的方法证明。
立即可化归到 $K$ 由单个阿贝尔层给出的情形。事实上，用下有界复形
$\mathcal{F}^\bullet$ 表示 $K$。由层的情形可知
$\mathcal{F}^n = a_{X, *} a_X^{-1} \mathcal{F}^n$，
且层 $R^qa_{X, *}a_X^{-1}\mathcal{F}^n$ 在 $q > 0$ 时为零。
把 Leray 无环性引理（《导出范畴》引理
\ref{derived-lemma-leray-acyclicity}）应用于
$a_X^{-1}\mathcal{F}^\bullet$ 和函子 $a_{X, *}$，即得结论。
以下假设 $K = \mathcal{F}$。

\medskip\noindent
由引理 \ref{spaces-more-cohomology-lemma-comparison-fppf-etale}，有
$a_{X, *}a_X^{-1}\mathcal{F} = \mathcal{F}$。因此只需证明：当 $q > 0$ 时
$R^qa_{X, *}a_X^{-1}\mathcal{F} = 0$。
为此可使用
% P10-E0159
$a_X = \epsilon_X \circ \pi_X$ 以及 Leray 谱序列
（《站点上的上同调》引理
\ref{sites-cohomology-lemma-relative-Leray}）。
由引理 \ref{spaces-more-cohomology-lemma-comparison-fppf-etale}，当 $i > 0$ 时有
$R^i\epsilon_{X, *}(a_X^{-1}\mathcal{F}) = 0$。还有
$\epsilon_{X, *}a_X^{-1}\mathcal{F} = \pi_X^{-1}\mathcal{F}$
并且由引理 \ref{spaces-more-cohomology-lemma-cohomological-descent-etale}，当 $j > 0$ 时有
$R^j\pi_{X, *}(\pi_X^{-1}\mathcal{F}) = 0$。证明完成。
\end{proof}

\begin{lemma}
\label{spaces-more-cohomology-lemma-compare-cohomology-etale-fppf}
设 $S$ 为概形，$X$ 为 $S$ 上的代数空间，并令
$a_X : \Sh((\textit{Spaces}/X)_{fppf}) \to \Sh(X_\etale)$ 如上。则：
\begin{enumerate}
\item $H^q(X_\etale, \mathcal{F}) = H^q_{fppf}(X, a_X^{-1}\mathcal{F})$
其中 $\mathcal{F}$ 是 $X_\etale$ 上的阿贝尔层；
\item 对 $K \in D^+(X_\etale)$，有
$H^q(X_\etale, K) = H^q_{fppf}(X, a_X^{-1}K)$。
\end{enumerate}
例如，若 $A$ 是阿贝尔群，则
$H^q_\etale(X, \underline{A}) = H^q_{fppf}(X, \underline{A})$。
\end{lemma}

\begin{proof}
由引理 \ref{spaces-more-cohomology-lemma-cohomological-descent-etale-fppf} 和
《站点上的上同调》评注 \ref{sites-cohomology-remark-before-Leray} 即得。
\end{proof}

\begin{lemma}
\label{spaces-more-cohomology-lemma-push-pull-fppf-etale}
设 $S$ 为概形。设 $f : X \to Y$ 为 $S$ 上代数空间的态射。
则有拓扑斯的交换图
$$
\xymatrix{
\Sh((\textit{Spaces}/X)_{fppf}) \ar[rr]_{f_{big, fppf}} \ar[d]_{\epsilon_X} & &
\Sh((\textit{Spaces}/Y)_{fppf}) \ar[d]^{\epsilon_Y} \\
\Sh((\textit{Spaces}/X)_\etale) \ar[rr]^{f_{big, \etale}} & &
\Sh((\textit{Spaces}/Y)_\etale)
}
$$
以及
$$
\xymatrix{
\Sh((\textit{Spaces}/X)_{fppf}) \ar[rr]_{f_{big, fppf}} \ar[d]_{a_X} & &
\Sh((\textit{Spaces}/Y)_{fppf}) \ar[d]^{a_Y} \\
\Sh(X_\etale) \ar[rr]^{f_{small}} & &
\Sh(Y_\etale)
}
$$
其中 $a_X = \pi_X \circ \epsilon_X$，并且
% P10-E0160
$a_Y = \pi_X \circ \epsilon_X$。
\end{lemma}

\begin{proof}
展开所涉及态射的定义即可立即得到结论；参见《空间上的拓扑》第
\ref{spaces-topologies-section-fppf} 节以及本章第 \ref{spaces-more-cohomology-section-compare} 节。
\end{proof}

\begin{lemma}
\label{spaces-more-cohomology-lemma-proper-push-pull-fppf-etale}
在引理 \ref{spaces-more-cohomology-lemma-push-pull-fppf-etale} 中，若 $f$ 为真态射，则有
\begin{enumerate}
\item $a_Y^{-1} \circ f_{small, *} = f_{big, fppf, *} \circ a_X^{-1}$；以及
\item
$a_Y^{-1}(Rf_{small, *}K) = Rf_{big, fppf, *}(a_X^{-1}K)$
对任意上同调层均为挠层的 $K$ 属于 $D^+(X_\etale)$ 成立。
\end{enumerate}
\end{lemma}

\begin{proof}
证明 (1)。重复引理 \ref{spaces-more-cohomology-lemma-compare-higher-direct-image-proper}
第 (1) 部分的证明即可；这里我们改为由该结果推出本结论。
由于 $\epsilon_{Y, *}$ 在底层预层上是恒等函子，
它反映同构。引理 \ref{spaces-more-cohomology-lemma-comparison-fppf-etale} 表明
$\epsilon_{Y, *} \circ a_Y^{-1} = \pi_Y^{-1}$，
对 $X$ 亦然。为了证明典范映射
$a_Y^{-1}f_{small, *}\mathcal{F} \to f_{big, fppf, *}a_X^{-1}\mathcal{F}$
是同构，只须证明
\begin{align*}
\pi_Y^{-1}f_{small, *}\mathcal{F}
& =
\epsilon_{Y, *}a_Y^{-1}f_{small, *}\mathcal{F} \\
& \to 
\epsilon_{Y, *}f_{big, fppf, *}a_X^{-1}\mathcal{F} \\
& =
f_{big, \etale, *} \epsilon_{X, *}a_X^{-1}\mathcal{F} \\
& =
f_{big, \etale, *}\pi_X^{-1}\mathcal{F}
\end{align*}
是同构。这正是引理 \ref{spaces-more-cohomology-lemma-compare-higher-direct-image-proper}
第 (1) 部分。

\medskip\noindent
为证明 (2)，我们使用
\begin{align*}
R\epsilon_{Y, *}Rf_{big, fppf, *}a_X^{-1}K
& =
Rf_{big, \etale, *}R\epsilon_{X, *}a_X^{-1}K \\
& =
Rf_{big, \etale, *}\pi_X^{-1}K \\
& =
\pi_Y^{-1}Rf_{small, *}K \\
& =
R\epsilon_{Y, *} a_Y^{-1}Rf_{small, *}K
\end{align*}
% P10-E0158
第一个等式来自引理 \ref{spaces-more-cohomology-lemma-push-pull-fppf-etale} 中的交换图以及
《站点上的上同调》引理
\ref{sites-cohomology-lemma-derived-pushforward-composition}.
% P10-E0161
第二个等式是引理 \ref{spaces-more-cohomology-lemma-cohomological-descent-etale-fppf}。
第三个等式是引理 \ref{spaces-more-cohomology-lemma-compare-higher-direct-image-proper} 第 (2) 部分。
第四个等式再次使用引理 \ref{spaces-more-cohomology-lemma-cohomological-descent-etale-fppf}。
因此，基变换映射
$a_Y^{-1}(Rf_{small, *}K) \to Rf_{big, fppf, *}(a_X^{-1}K)$
诱导同构
$$
R\epsilon_{Y, *}a_Y^{-1}Rf_{small, *}K \to
R\epsilon_{Y, *}Rf_{big, fppf, *}a_X^{-1}K
$$
最后用以下观察完成证明：设 $L$ 属于 $D^+(Y_\etale)$、
$M$ 属于 $D^+((\textit{Spaces}/Y)_{fppf})$，且映射
$\alpha : a_Y^{-1}L \to M$ 使得 $R\epsilon_{Y, *}\alpha$ 为同构，
则该映射本身也是同构。事实上，我们对 $i$ 归纳证明
$H^i(\alpha)$ 为同构。对所有充分小的 $i$，这显然成立。
若对 $i \leq i_0$ 成立，则由引理
\ref{spaces-more-cohomology-lemma-comparison-fppf-etale} 以及在此范围内
$H^i(M) = a_Y^{-1}H^i(L)$，可知
对 $j > 0$ 且 $i \leq i_0$，有
$R^j\epsilon_{Y, *}H^i(M) = 0$。于是谱序列论证给出
$\epsilon_{Y, *}H^{i_0 + 1}(M) = H^{i_0 + 1}(R\epsilon_{Y, *}M)$。
从而
$\epsilon_{Y, *}H^{i_0 + 1}(M) = \pi_Y^{-1}H^{i_0 + 1}(L) =
\epsilon_{Y, *}a_Y^{-1}H^{i_0 + 1}(L)$.
这蕴含 $H^{i_0 + 1}(\alpha)$ 是同构，因为 $\epsilon_{Y, *}$
在底层预层上恒等，因而反映同构。证毕。
\end{proof}

\begin{lemma}
\label{spaces-more-cohomology-lemma-finite-push-pull-fppf-etale}
在引理 \ref{spaces-more-cohomology-lemma-push-pull-fppf-etale} 中，若 $f$ 为有限态射，则
$a_Y^{-1}(Rf_{small, *}K) = Rf_{big, fppf, *}(a_X^{-1}K)$
对任意属于 $D^+(X_\etale)$ 的 $K$ 成立。
\end{lemma}

\begin{proof}
取满平展态射 $V \to Y$，其中 $V$ 为概形。
只须证明限制到 $V$ 后基变换映射为同构，故可设 $Y$ 为概形。
由于该态射有限，因而可表，我们进一步可设 $X$ 和 $Y$ 均为概形。
此时，结论由概形的情形（《平展上同调》引理
\ref{etale-cohomology-lemma-V-C-all-n-etale-fppf} 第 (2) 部分）得出；
这里使用第 \ref{spaces-more-cohomology-section-api} 节所讨论的拓扑斯比较，
尤其是引理 \ref{spaces-more-cohomology-lemma-compare-cohomology-other-topologies}。
略去若干细节。
\end{proof}

\begin{lemma}
\label{spaces-more-cohomology-lemma-descent-sheaf-fppf-etale}
在引理 \ref{spaces-more-cohomology-lemma-push-pull-fppf-etale} 中，设
$f$ 平坦、局部有限表示且满。则函子
$$
\Sh(Y_\etale) \longrightarrow
\left\{
(\mathcal{G}, \mathcal{H}, \alpha)
\middle|
\begin{matrix}
\mathcal{G} \in \Sh(X_\etale),\ \mathcal{H} \in \Sh((\Sch/Y)_{fppf}), \\
\alpha : a_X^{-1}\mathcal{G} \to f_{big, fppf}^{-1}\mathcal{H}
\text{ 为同构}
\end{matrix}
\right\}
$$
把 $\mathcal{F}$ 映到
$(f_{small}^{-1}\mathcal{F}, a_Y^{-1}\mathcal{F}, can)$，它是范畴等价。
\end{lemma}

\begin{proof}
函子 $a_X^{-1}$ 全忠实（由引理 \ref{spaces-more-cohomology-lemma-comparison-fppf-etale}，
$a_{X, *}a_X^{-1} = \text{id}$）。因此，遗忘函子
$(\mathcal{G}, \mathcal{H}, \alpha) \mapsto \mathcal{H}$
把三元组范畴认同为 $\Sh((\Sch/Y)_{fppf})$ 的一个全子范畴。
此外，函子 $a_Y^{-1}$ 全忠实，故引理中的函子也全忠实。

\medskip\noindent
设有平展覆盖 $\{Y_i \to Y\}$。令 $f_i : X_i \to Y_i$ 为 $f$ 的基变换，
并记
$f_{ij} = f_i \times f_j : X_i \times_X X_j  \to Y_i \times_Y Y_j$。
断言：若对所有 $i,j$，引理对 $f_i$ 和 $f_{ij}$ 均成立，则它对 $f$ 也成立。
为此注意，给定平展覆盖在四个站点
$Y_\etale, X_\etale, (\Sch/Y)_{fppf}, (\Sch/X)_{fppf}$
中的每一个上都给出终对象的平展覆盖。因此在四种情形中，层范畴都等价于
该覆盖的粘合数据范畴（《站点论》引理
\ref{sites-lemma-mapping-property-glue}）。随后，一个巨大的范畴交换图即完成
该断言的证明；细节从略。该断言表明可以在 $Y$ 上平展局部地工作，
特别地可设 $Y$ 为概形。

\medskip\noindent
设 $Y$ 为概形。取概形 $X'$ 及满平展态射 $s : X' \to X$。
令 $f' = f \circ s : X' \to Y$；注意 $f'$ 满、局部有限表示且平坦。
断言：若引理对 $f'$ 成立，则对 $f$ 也成立。
事实上，给定 $f$ 的三元组 $(\mathcal{G}, \mathcal{H}, \alpha)$，
沿 $s$ 拉回便得到 $f'$ 的三元组
$(s_{small}^{-1}\mathcal{G}, \mathcal{H}, s_{big, fppf}^{-1}\alpha)$。
该三元组的一个解给出 $Y_\etale$ 上的层 $\mathcal{F}$，满足
$a_Y^{-1}\mathcal{F} = \mathcal{H}$。由证明第一段，这意味着原三元组位于
本质像中。于是问题化归为 $X$、$Y$ 均为概形的情形。
此情形由《平展上同调》引理
\ref{etale-cohomology-lemma-descent-sheaf-fppf-etale} 得出；
这里利用第 \ref{spaces-more-cohomology-section-api} 节的讨论，尤其是引理
\ref{spaces-more-cohomology-lemma-compare-cohomology-other-topologies}。
\end{proof}



\section{比较 fppf 拓扑与平展拓扑：模}
\label{spaces-more-cohomology-section-fppf-etale-modules}

\noindent
我们继续第 \ref{spaces-more-cohomology-section-fppf-etale} 节的讨论，
但本节简要考察模层的情形。

\medskip\noindent
设 $S$ 为概形，$X$ 为 $S$ 上的代数空间。
第 \ref{spaces-more-cohomology-section-fppf-etale} 节引入的站点态射
$\epsilon_X$、$\pi_X$ 及其复合 $a_X$
都自然提升为赋环站点的态射。第一个写作
$$
\epsilon_X :
((\textit{Spaces}/X)_{fppf}, \mathcal{O})
\longrightarrow
((\textit{Spaces}/X)_\etale, \mathcal{O})
$$
注意，我们确实可以对结构层使用同一符号，因为这些层有相同的底层预层。
第二个是
$$
\pi_X :
((\textit{Spaces}/X)_\etale, \mathcal{O})
\longrightarrow
(X_\etale, \mathcal{O}_X)
$$
第三个态射是
$$
a_X :
((\textit{Spaces}/X)_{fppf}, \mathcal{O})
\longrightarrow
(X_\etale, \mathcal{O}_X)
$$
下面回顾关于这些站点上拟凝聚模的已有结论。

\begin{lemma}
\label{spaces-more-cohomology-lemma-review-quasi-coherent}
设 $S$ 为概形，$X$ 为 $S$ 上的代数空间，
$\mathcal{F}$ 为拟凝聚 $\mathcal{O}_X$-模。
\begin{enumerate}
\item 规则
$$
\mathcal{F}^a : (\textit{Spaces}/X)_\etale \longrightarrow \textit{Ab},\quad
(f : Y \to X) \longmapsto \Gamma(Y, f^*\mathcal{F})
$$
对 fpqc 覆盖满足层条件，因而尤其对 fppf 覆盖和平展覆盖满足层条件；
\item 在 $(\textit{Spaces}/X)_\etale$ 上，
$\mathcal{F}^a = \pi_X^*\mathcal{F}$；
\item 在 $(\textit{Spaces}/X)_{fppf}$ 上，
$\mathcal{F}^a = a_X^*\mathcal{F}$；
\item 规则 $\mathcal{F} \mapsto \mathcal{F}^a$ 给出拟凝聚
$\mathcal{O}_X$-模与
$((\textit{Spaces}/X)_\etale, \mathcal{O})$ 上拟凝聚模之间的等价；
\item 规则 $\mathcal{F} \mapsto \mathcal{F}^a$ 给出拟凝聚
$\mathcal{O}_X$-模与
$((\textit{Spaces}/X)_{fppf}, \mathcal{O})$ 上拟凝聚模之间的等价；
\item 有 $\epsilon_{X, *}a_X^*\mathcal{F} = \pi_X^*\mathcal{F}$
及 $a_{X, *}a_X^*\mathcal{F} = \mathcal{F}$；
\item 对 $i > 0$，有 $R^i\epsilon_{X, *}(a_X^*\mathcal{F}) = 0$
及 $R^ia_{X, *}(a_X^*\mathcal{F}) = 0$。
\end{enumerate}
\end{lemma}

\begin{proof}
第 (1) 部分是拟凝聚模 fppf 下降的推论。
具体地，设 $\{f_i : U_i \to U\}$ 是
$(\textit{Spaces}/X)_\etale$ 中的 fpqc 覆盖，
记 $g : U \to X$ 为结构态射。设有一族截面
$s_i \in \Gamma(U_i , f_i^*g^*\mathcal{F})$，满足
$s_i|_{U_i \times_U U_j} = s_j|_{U_i \times_U U_j}$。
% P10-E0163
我们要找相应截面 $s \in \Gamma(U, g^*\mathcal{F})$。
可以把 $s_i$ 重新解释为一族映射
$\varphi_i : f_i^*\mathcal{O}_U = \mathcal{O}_{U_i} \to f_i^*g^*\mathcal{F}$
并且它们与 $U$ 上拟凝聚层 $\mathcal{O}_U$ 和 $g^*\mathcal{F}$
所带的典范下降数据相容。因此由《代数空间的下降》命题
\ref{spaces-descent-proposition-fpqc-descent-quasi-coherent}，
可将它们唯一地下降为映射 $\mathcal{O}_U \to g^*\mathcal{F}$，
它给出所求截面 $s$。

\medskip\noindent
我们将由概形的相应结论推出 (2)--(7)。
取平展覆盖 $\{X_i \to X\}_{i \in I}$，其中每个 $X_i$ 都是概形；
注意 $X_i \times_X X_j$ 也是概形。此覆盖在三个站点
$(\textit{Spaces}/X)_{fppf}$、$(\textit{Spaces}/X)_\etale$ 和 $X_\etale$
的每一个上都诱导终对象的覆盖。因此，这些站点上的层范畴等价于相应覆盖的
下降数据范畴；参见《站点论》引理 \ref{sites-lemma-mapping-property-glue}。
第 (2)、(3) 部分是局部的，因为已有粘合结论；拟凝聚性是局部性质，
故第 (4)、(5) 部分也是局部的。显然，第 (6)、(7) 部分也是局部的。
所以只须在 $X$ 为概形时证明引理的第 (2)--(7) 部分。

\medskip\noindent
设 $X$ 为概形。嵌入
$(\Sch/X)_\etale \subset (\textit{Spaces}/X)_\etale$ 和
$(\Sch/X)_{fppf} \subset (\textit{Spaces}/X)_{fppf}$
由引理 \ref{spaces-more-cohomology-lemma-compare-cohomology-other-topologies} 给出赋环拓扑斯的等价。
于是 (2)--(7) 由概形情形即《平展上同调》引理
\ref{etale-cohomology-lemma-review-quasi-coherent} 得出。
为经由此等价迁移拟凝聚性，使用拟凝聚性是模的内蕴性质这一事实；
参见《站点上的模》第 \ref{sites-modules-section-local} 节。
略去若干次要细节。
\end{proof}

\begin{lemma}
\label{spaces-more-cohomology-lemma-cohomological-descent-etale-fppf-modules}
设 $S$ 为概形，$X$ 为 $S$ 上的代数空间。
对拟凝聚 $\mathcal{O}_X$-模 $\mathcal{F}$，映射
$$
\pi_X^*\mathcal{F} \longrightarrow R\epsilon_{X, *}(a_X^*\mathcal{F})
\quad\text{和}\quad
\mathcal{F} \longrightarrow Ra_{X, *}(a_X^*\mathcal{F})
$$
均为同构。
\end{lemma}

\begin{proof}
这是引理 \ref{spaces-more-cohomology-lemma-review-quasi-coherent} 第 (6)、(7) 部分的直接推论。
\end{proof}

\begin{lemma}
\label{spaces-more-cohomology-lemma-vanishing-adequate}
设 $S$ 为概形，$X$ 为 $S$ 上的代数空间。
设 $\mathcal{F}_1 \to \mathcal{F}_2 \to \mathcal{F}_3$
为拟凝聚 $\mathcal{O}_X$-模的复形。令
$$
\mathcal{H}_\etale =
\Ker(\pi_X^*\mathcal{F}_2 \to \pi_X^*\mathcal{F}_3)/
\Im(\pi_X^*\mathcal{F}_1 \to \pi_X^*\mathcal{F}_2)
$$
于 $(\textit{Spaces}/X)_\etale$ 上，并令
$$
\mathcal{H}_{fppf} =
\Ker(a_X^*\mathcal{F}_2 \to a_X^*\mathcal{F}_3)/
\Im(a_X^*\mathcal{F}_1 \to a_X^*\mathcal{F}_2)
$$
于 $(\textit{Spaces}/X)_{fppf}$ 上。则
$\mathcal{H}_\etale = \epsilon_{X, *}\mathcal{H}_{fppf}$，并且
$$
H^p_\etale(U, \mathcal{H}_\etale) = H^p_{fppf}(U, \mathcal{H}_{fppf}) = 0
$$
对 $p > 0$ 以及 $(\textit{Spaces}/X)_\etale$ 的任意仿射对象 $U$ 成立。
\end{lemma}

\noindent
事实上还有更强的结论：$(\textit{Spaces}/X)_{fppf}$ 上那些在 fppf 局部
具有引理中形式的模称为适足模。
% P10-E0164
它们构成全体 $\mathcal{O}$-模范畴的一个弱 Serre 子范畴；
其上同调见《适足模》第 \ref{adequate-section-adequate} 节。

\begin{proof}
对 $(\textit{Spaces}/X)_\etale$ 的任意对象 $f : U \to X$，
考虑经由第 \ref{spaces-more-cohomology-section-compare} 节讨论的函子 $i_f^* = i_f^{-1}$
把 $\mathcal{H}_\etale$ 限制到 $U_\etale$ 所得的
$\mathcal{H}_\etale|_{U_\etale}$。层 $\mathcal{H}_\etale|_{U_\etale}$
等于复形 $f^*\mathcal{F}_\bullet$ 的第 $1$ 次同调。
这是因为作为赋环站点 $U_\etale \to X_\etale$ 的态射有
$i_f \circ \pi_X = f$。特别地，
$\mathcal{H}_\etale|_{U_\etale}$ 是拟凝聚 $\mathcal{O}_U$-模。
其次，设 $g : V \to U$ 是 $(\textit{Spaces}/X)_\etale$ 中的平坦态射。
由于
$$
i_{f \circ g}^* \circ \pi_X^* = (f \circ g)^* = g^* \circ f^*
$$
作为站点 $V_\etale \to X_\etale$ 的态射成立，且 $g$ 平坦，
所以 $g^*$ 正合，从而得到
$$
\mathcal{H}_\etale|_{V_\etale} =
g^*\left(\mathcal{H}_\etale|_{U_\etale}\right)
$$
有了这些准备，即可证明引理。

\medskip\noindent
设 $\mathcal{U} = \{g_i : U_i \to U\}_{i \in I}$ 为 fppf 覆盖，
其中 $f : U \to X$ 如上。把引理 \ref{spaces-more-cohomology-lemma-review-quasi-coherent}
第 (1) 部分应用于 $\mathcal{H}_\etale|_{U_\etale}$，再结合上述结论，
可知 $\mathcal{H}_\etale$ 对覆盖 $\mathcal{U}$ 满足层条件。
因此 $\mathcal{H}_\etale$ 已经是 fppf 层，这意味着作为预层，
$\mathcal{H}_{fppf}$ 等于 $\mathcal{H}_\etale$。特别地，
$\mathcal{H}_\etale = \epsilon_{X, *}\mathcal{H}_{fppf}$。

\medskip\noindent
最后，为证明消没性，我们使用《站点上的上同调》引理
\ref{sites-cohomology-lemma-cech-vanish-collection}。
令 $\mathcal{B}$ 为 $(\textit{Spaces}/X)_{fppf}$ 中的仿射对象，
令 $\text{Cov}$ 为有限 fppf 覆盖
$\mathcal{U} = \{U_i \to U\}_{i = 1, \ldots, n}$ 的集合，
其中 $U$ 和 $U_i$ 均仿射。我们有
$$
{\check H}^p(\mathcal{U}, \mathcal{H}_\etale) =
{\check H}^p(\mathcal{U}, \left(\mathcal{H}_\etale|_{U_\etale}\right)^a)
$$
这是因为由第一段，$\mathcal{H}_\etale$ 在那些平坦于 $U$ 的仿射概形
$U_{i_0} \times_U \ldots \times_U U_{i_p}$ 上的取值，
与拟凝聚模 $\mathcal{H}_\etale|_{U_\etale}$ 的拉回取值相同。
于是由《下降》引理 \ref{descent-lemma-standard-covering-Cech-quasi-coherent}
得到消没性。这就完成了证明。
\end{proof}

\begin{lemma}
\label{spaces-more-cohomology-lemma-cohomological-descent-etale-fppf-modules-unbounded}
设 $S$ 为概形，$X$ 为 $S$ 上的代数空间。
对 $K \in D_\QCoh(\mathcal{O}_X)$，映射
$$
L\pi_X^*K \longrightarrow R\epsilon_{X, *}(La_X^*K)
\quad\text{和}\quad
K \longrightarrow Ra_{X, *}(La_X^*K)
$$
均为同构。这里
$a_X : \Sh((\textit{Spaces}/X)_{fppf}) \to \Sh(X_\etale)$ 如上所述。
\end{lemma}

\begin{proof}
该问题在 $X$ 上是平展局部的，故可设 $X$ 仿射。
设 $X = \Spec(A)$。由《空间的导出范畴》引理
\ref{spaces-perfect-lemma-derived-quasi-coherent-small-etale-site}
及《概形的导出范畴》引理 \ref{perfect-lemma-affine-compare-bounded}，
有 $D_\QCoh(\mathcal{O}_X) = D(A)$。
% P10-E0165
因此可选取一个 $A$-模的 K-平坦复形 $K^\bullet$，
使其对应的拟凝聚 $\mathcal{O}_X$-模复形 $\mathcal{K}^\bullet$ 表示 $K$。
我们断言 $\mathcal{K}^\bullet$ 是 $\mathcal{O}_X$-模的 K-平坦复形。

\medskip\noindent
证明该断言。由《概形的导出范畴》引理
\ref{perfect-lemma-affine-K-flat}
可知 $\widetilde{K}^\bullet$ 在概形
$(\Spec(A), \mathcal{O}_{\Spec(A)})$ 上 K-平坦。
其次注意 $\mathcal{K}^\bullet = \epsilon^*\widetilde{K}^\bullet$，
其中 $\epsilon$ 如《空间的导出范畴》引理
\ref{spaces-perfect-lemma-derived-quasi-coherent-small-etale-site} 所述。
再由《站点上的上同调》引理
\ref{sites-cohomology-lemma-pullback-K-flat-points}，以及概形的平展站点有足够多点
这一事实（《平展上同调》评注 \ref{etale-cohomology-remarks-enough-points}），
可知 $\mathcal{K}^\bullet$ K-平坦。

\medskip\noindent
由该断言可知
$La_X^*K = a_X^*\mathcal{K}^\bullet$ 和
$L\pi_X^*K = \pi_X^*\mathcal{K}^\bullet$。
证明第一部分表明，拟凝聚模的拉回 $a_X^*\mathcal{K}^n$
分别对 $\epsilon_{X, *}$ 和 $a_{X, *}$ 无上同调；
那么证明必然由 Leray 无上同调性引理完成了吧？其实并非如此，
因为 Leray 无上同调性引理只适用于下有界复形。
不过下一段将说明，本结论确实可由下有界情形推出，
因为我们的复形是拟凝聚模下有界复形的导出逆极限。

\medskip\noindent
由引理 \ref{spaces-more-cohomology-lemma-vanishing-adequate}，
$\pi_X^*\mathcal{K}^\bullet$ 和 $a_X^*\mathcal{K}^\bullet$ 的上同调层
在 $(\textit{Spaces}/X)_\etale$ 的仿射对象上有消没的高阶上同调群。
因此有
$$
L\pi_X^*K = R\lim \tau_{\geq -n}(L\pi_X^*K)
\quad\text{和}\quad
La_X^*K = R\lim \tau_{\geq -n}(La_X^*K)
$$
这是《站点上的上同调》引理
\ref{sites-cohomology-lemma-is-limit-dimension} 的结论。

\medskip\noindent
% P10-E0166
证明 $L\pi_X^*K = R\epsilon_{X, *}(La_X^*\mathcal{F})$。
由上可得
$$
R\epsilon_{X, *}La_X^*K =
R\lim R\epsilon_{X, *}(\tau_{\geq -n}(La_X^*K))
$$
这是《站点上的上同调》引理
\ref{sites-cohomology-lemma-Rf-commutes-with-Rlim}。
注意 $\tau_{\geq -n}(La_X^*K)$ 由
$\tau_{\geq -n}(a_X^*\mathcal{K}^\bullet)$ 表示，
它未必等同于 $a_X^*(\tau_{\geq -n}\mathcal{K}^\bullet)$。
但以下两个系统显然
$$
\{\tau_{\geq -n}(a_X^*\mathcal{K}^\bullet)\}_{n \geq 1}
\quad\text{和}\quad
\{a_X^*(\tau_{\geq -n}\mathcal{K}^\bullet)\}_{n \geq 1}
$$
作为 pro-系统同构。由 Leray 无上同调性引理
（《导出范畴》引理 \ref{derived-lemma-leray-acyclicity}）
以及本引理第一部分，可知
$$
R\epsilon_{X, *}(a_X^*(\tau_{\geq -n}\mathcal{K}^\bullet)) =
\pi_X^*(\tau_{\geq -n}\mathcal{K}^\bullet)
$$
然后使用以下两个系统
$$
\{\tau_{\geq -n}(\pi_X^*\mathcal{K}^\bullet)\}_{n \geq 1}
\quad\text{和}\quad
\{\pi_X^*(\tau_{\geq -n}\mathcal{K}^\bullet)\}_{n \geq 1}
$$
作为 pro-系统同构这一事实。最后汇总如下：
\begin{align*}
R\epsilon_{X, *}La_X^*K
& =
R\epsilon_{X, *} (R\lim \tau_{\geq -n}(La_X^*K)) \\
& =
R\lim R\epsilon_{X, *}(\tau_{\geq -n}(La_X^*K)) \\
& =
R\lim R\epsilon_{X, *}(\tau_{\geq -n}(a_X^*\mathcal{K}^\bullet)) \\
& =
R\lim R\epsilon_{X, *}(a_X^*(\tau_{\geq -n}\mathcal{K}^\bullet)) \\
& =
R\lim \pi_X^*(\tau_{\geq -n}\mathcal{K}^\bullet) \\
& =
R\lim \tau_{\geq -n}(\pi_X^*\mathcal{K}^\bullet) \\
& =
R\lim \tau_{\geq -n}(L\pi_X^*K) \\
& =
L\pi_X^*K
\end{align*}
在第四个和第六个等式中，我们使用了同构的 pro-系统具有相同的
$R\lim$（略去一个小细节）。也可以利用引理
\ref{spaces-more-cohomology-lemma-vanishing-adequate} 所证明的、关于复形
$\tau_{\geq -n}a_X^*\mathcal{K}^\bullet$ 各项上同调的更多信息来避开这一步；
这会直接证明
$R\epsilon_{X, *}(\tau_{\geq -n}(a_X^*\mathcal{K}^\bullet)) =
\tau_{\geq -n}(\pi_X^*\mathcal{K}^\bullet)$。

\medskip\noindent
% P10-E0167
等式 $K = Ra_{X, *}(La_X^*\mathcal{F})$ 的证明完全相同；
最后一步使用《空间的导出范畴》引理
\ref{spaces-perfect-lemma-nice-K-injective} 给出的
$K = R\lim \tau_{\geq -n}K$。
\end{proof}








\section{比较 ph 拓扑与平展拓扑}
\label{spaces-more-cohomology-section-ph-etale}

\noindent
本节对应于《平展上同调》第
\ref{etale-cohomology-section-ph-etale} 节。

\medskip\noindent
设 $S$ 为概形，$X$ 为 $S$ 上的代数空间。
在范畴 $\textit{Spaces}/X$ 上考虑 ph 拓扑和平展拓扑。
恒等函子
$(\textit{Spaces}/X)_\etale \to (\textit{Spaces}/X)_{ph}$
连续，因为由《空间上的拓扑》引理
\ref{spaces-topologies-lemma-zariski-etale-smooth-syntomic-fppf-ph}，
每个平展覆盖都是 ph 覆盖。因此它定义站点态射
$$
\epsilon_X :
(\textit{Spaces}/X)_{ph} \longrightarrow (\textit{Spaces}/X)_\etale
$$
这里使用了《站点论》命题 \ref{sites-proposition-get-morphism}。
注意 $\epsilon_{X, *}$ 在底层预层上是恒等函子，而
$\epsilon_X^{-1}$ 把一个平展层映到其 ph 层化。
再考虑比较大、小平展站点的站点态射
$$
\pi_X : (\textit{Spaces}/X)_\etale \longrightarrow X_{spaces, \etale}
$$
参见第 \ref{spaces-more-cohomology-section-compare} 节。其复合给出站点态射
$$
a_X = \pi_X \circ \epsilon_X :
(\textit{Spaces}/X)_{ph}
\longrightarrow
X_{spaces, \etale}
$$
若 $\mathcal{H}$ 是 $(\textit{Spaces}/X)_{ph}$ 上的阿贝尔层，
则以 $H^n_{ph}(U, \mathcal{H})$ 表示 $\mathcal{H}$ 在
$(\textit{Spaces}/X)_{ph}$ 的对象 $U$ 上的上同调。

\begin{lemma}
\label{spaces-more-cohomology-lemma-comparison-ph-etale}
设 $S$ 为概形，$X$ 为 $S$ 上的代数空间。
\begin{enumerate}
\item 对 $\mathcal{F} \in \Sh(X_\etale)$，有
$\epsilon_{X, *}a_X^{-1}\mathcal{F} = \pi_X^{-1}\mathcal{F}$
及 $a_{X, *}a_X^{-1}\mathcal{F} = \mathcal{F}$。
\item 对挠层 $\mathcal{F} \in \textit{Ab}(X_\etale)$，
当 $i > 0$ 时有 $R^i\epsilon_{X, *}(a_X^{-1}\mathcal{F}) = 0$。
\end{enumerate}
\end{lemma}

\begin{proof}
有 $a_X^{-1}\mathcal{F} = \epsilon_X^{-1} \pi_X^{-1}\mathcal{F}$。
由引理 \ref{spaces-more-cohomology-lemma-describe-pullback}，平展层
$\pi_X^{-1}\mathcal{F}$ 对 ph 拓扑也是层，故等于
$a_X^{-1}\mathcal{F}$，因为沿 $\epsilon_X$ 拉回就是进行 ph 层化。
还要回忆，$\epsilon_{X, *}$ 在底层预层上是恒等函子。
现在，第 (1) 部分立即由引理 \ref{spaces-more-cohomology-lemma-describe-pullback}
中对 $\pi_X^{-1}$ 的显式描述得出。

\medskip\noindent
我们把第 (2) 部分化归为概形情形来证明；参见《平展上同调》引理
\ref{etale-cohomology-lemma-V-C-all-n-etale-ph} 第 (1) 部分。
由于每个代数空间平展局部都是概形，这一方法“显然可行”。
下面给出细节，不过我们建议读者跳过此证明。

\medskip\noindent
对 $(\textit{Spaces}/X)_{ph}$ 上的阿贝尔层 $\mathcal{H}$，
高阶直像 $R^p\epsilon_{X, *}\mathcal{H}$ 是
$(\textit{Spaces}/X)_\etale$ 上预层
$U \mapsto H^p_{ph}(U, \mathcal{H})$ 所关联的层；
参见《站点上的上同调》引理
\ref{sites-cohomology-lemma-higher-direct-images}。
由于 $(\textit{Spaces}/X)_\etale$ 的每个对象都有由概形组成的覆盖，
只须证明：给定概形 $U/X$ 及
$\xi \in H^p_{ph}(U, a_X^{-1}\mathcal{F})$，可以找到平展覆盖
$\{U_i \to U\}$，使 $\xi$ 在每个 $U_i$ 上的限制均为零。我们有
\begin{align*}
H^p_{ph}(U, a_X^{-1}\mathcal{F})
& =
H^p((\textit{Spaces}/U)_{ph}, (a_X^{-1}\mathcal{F})|_{\textit{Spaces}/U}) \\
& =
H^p((\Sch/U)_{ph}, (a_X^{-1}\mathcal{F})|_{\Sch/U})
\end{align*}
其中第二个认同来自引理 \ref{spaces-more-cohomology-lemma-compare-cohomology-other-topologies}，
第一个则是关于限制的一般事实（《站点上的上同调》引理
\ref{sites-cohomology-lemma-cohomology-of-open}）。
考察第一段以及概形情形的对应结果（《平展上同调》引理
\ref{etale-cohomology-lemma-describe-pullback-pi-ph}），可知层
$(a_X^{-1}\mathcal{F})|_{\Sch/U}$ 与沿“概形版 $a_U$”的拉回一致。
因此可以找到平展覆盖 $\{U_i \to U\}$，使得对每个 $i$，我们的类在
$H^p((\Sch/U_i)_{ph}, (a_X^{-1}\mathcal{F})|_{\Sch/U_i})$ 中消失；
参见《平展上同调》引理
\ref{etale-cohomology-lemma-V-C-all-n-etale-ph}。
这里应当使用的精确陈述是：$V_n$ 对所有 $n$ 成立，
这就是概形情形的第 (2) 部分。再按上述相同公式对 $U_i$ 迁移回来，
便得到 $\xi$ 在 $U_i$ 上的限制为零，如所需。
\end{proof}

\noindent
概形情形中完成的艰苦工作现在告诉我们：
对于来自小平展站点的挠阿贝尔层，平展上同调与 ph 上同调一致。

\begin{lemma}
\label{spaces-more-cohomology-lemma-cohomological-descent-etale-ph}
设 $S$ 为概形，$X$ 为 $S$ 上的代数空间。
对上同调层均为挠层的 $K \in D^+(X_\etale)$，映射
$$
\pi_X^{-1}K \longrightarrow R\epsilon_{X, *}a_X^{-1}K
\quad\text{和}\quad
K \longrightarrow Ra_{X, *}a_X^{-1}K
$$
均为同构，其中
$a_X : \Sh((\textit{Spaces}/X)_{ph}) \to \Sh(X_\etale)$ 如上所述。
\end{lemma}

\begin{proof}
只证明第二个陈述；第一个更容易，且证明方法完全相同。
可以化归到 $K$ 由单个挠阿贝尔层给出的情形。
具体地，用挠阿贝尔层的下有界复形 $\mathcal{F}^\bullet$ 表示 $K$；
《站点上的上同调》引理 \ref{sites-cohomology-lemma-torsion}
保证可以如此选择。由单个层的情形可知
$\mathcal{F}^n = a_{X, *} a_X^{-1} \mathcal{F}^n$
且当 $q > 0$ 时，层 $R^qa_{X, *}a_X^{-1}\mathcal{F}^n$ 为零。
把 Leray 无上同调性引理（《导出范畴》引理
\ref{derived-lemma-leray-acyclicity}）应用于
$a_X^{-1}\mathcal{F}^\bullet$ 和函子 $a_{X, *}$，即得结论。
以下设 $K = \mathcal{F}$，其中 $\mathcal{F}$ 为挠阿贝尔层。

\medskip\noindent
由引理 \ref{spaces-more-cohomology-lemma-comparison-ph-etale}，
$a_{X, *}a_X^{-1}\mathcal{F} = \mathcal{F}$，故只须证明当 $q > 0$ 时
$R^qa_{X, *}a_X^{-1}\mathcal{F} = 0$。
% P10-E0168
为此可使用 $a_X = \epsilon_X \circ \pi_X$ 以及 Leray 谱序列
（《站点上的上同调》引理 \ref{sites-cohomology-lemma-relative-Leray}）。
由引理 \ref{spaces-more-cohomology-lemma-comparison-ph-etale}，当 $i > 0$ 时有
$R^i\epsilon_{X, *}(a_X^{-1}\mathcal{F}) = 0$。又有
$\epsilon_{X, *}a_X^{-1}\mathcal{F} = \pi_X^{-1}\mathcal{F}$
且由引理 \ref{spaces-more-cohomology-lemma-cohomological-descent-etale}，当 $j > 0$ 时有
$R^j\pi_{X, *}(\pi_X^{-1}\mathcal{F}) = 0$。证明完成。
\end{proof}

\begin{lemma}
\label{spaces-more-cohomology-lemma-compare-cohomology-etale-ph}
设 $S$ 为概形，$X$ 为 $S$ 上的代数空间。
沿用上述 $a_X : \Sh((\textit{Spaces}/X)_{ph}) \to \Sh(X_\etale)$，则：
\begin{enumerate}
\item $H^q(X_\etale, \mathcal{F}) = H^q_{ph}(X, a_X^{-1}\mathcal{F})$
对 $X_\etale$ 上的挠阿贝尔层 $\mathcal{F}$ 成立；
\item 对上同调层均为挠层的 $K \in D^+(X_\etale)$，有
$H^q(X_\etale, K) = H^q_{ph}(X, a_X^{-1}K)$。
\end{enumerate}
例如，若 $A$ 为挠阿贝尔群，则
$H^q_\etale(X, \underline{A}) = H^q_{ph}(X, \underline{A})$。
\end{lemma}

\begin{proof}
由引理 \ref{spaces-more-cohomology-lemma-cohomological-descent-etale-ph} 以及
《站点上的上同调》评注 \ref{sites-cohomology-remark-before-Leray} 即得。
\end{proof}

\begin{lemma}
\label{spaces-more-cohomology-lemma-push-pull-ph-etale}
设 $S$ 为概形，$f : X \to Y$ 为 $S$ 上代数空间的态射。
则有拓扑斯交换图
$$
\xymatrix{
\Sh((\textit{Spaces}/X)_{ph}) \ar[rr]_{f_{big, ph}} \ar[d]_{\epsilon_X} & &
\Sh((\textit{Spaces}/Y)_{ph}) \ar[d]^{\epsilon_Y} \\
\Sh((\textit{Spaces}/X)_\etale) \ar[rr]^{f_{big, \etale}} & &
\Sh((\textit{Spaces}/Y)_\etale)
}
$$
以及
$$
\xymatrix{
\Sh((\textit{Spaces}/X)_{ph}) \ar[rr]_{f_{big, ph}} \ar[d]_{a_X} & &
\Sh((\textit{Spaces}/Y)_{ph}) \ar[d]^{a_Y} \\
\Sh(X_\etale) \ar[rr]^{f_{small}} & &
\Sh(Y_\etale)
}
$$
% P10-E0169
其中 $a_X = \pi_X \circ \epsilon_X$ 且 $a_Y = \pi_X \circ \epsilon_X$。
\end{lemma}

\begin{proof}
展开所涉及态射的定义即可立即得到结论；参见《空间上的拓扑》第
\ref{spaces-topologies-section-ph} 节以及本章第 \ref{spaces-more-cohomology-section-compare} 节。
\end{proof}

\begin{lemma}
\label{spaces-more-cohomology-lemma-proper-push-pull-ph-etale}
在引理 \ref{spaces-more-cohomology-lemma-push-pull-ph-etale} 中，若 $f$ 为真态射，则有
\begin{enumerate}
\item $a_Y^{-1} \circ f_{small, *} = f_{big, ph, *} \circ a_X^{-1}$；以及
\item
$a_Y^{-1}(Rf_{small, *}K) = Rf_{big, ph, *}(a_X^{-1}K)$
对任意上同调层均为挠层的 $K$ 属于 $D^+(X_\etale)$ 成立。
\end{enumerate}
\end{lemma}

\begin{proof}
证明 (1)。重复引理 \ref{spaces-more-cohomology-lemma-compare-higher-direct-image-proper}
第 (1) 部分的证明即可；这里我们改为由该结果推出本结论。
由于 $\epsilon_{Y, *}$ 在底层预层上是恒等函子，它反映同构。
引理 \ref{spaces-more-cohomology-lemma-comparison-ph-etale} 表明
$\epsilon_{Y, *} \circ a_Y^{-1} = \pi_Y^{-1}$，对 $X$ 亦然。
为了证明典范映射
$a_Y^{-1}f_{small, *}\mathcal{F} \to f_{big, ph, *}a_X^{-1}\mathcal{F}$
是同构，只须证明
\begin{align*}
\pi_Y^{-1}f_{small, *}\mathcal{F}
& =
\epsilon_{Y, *}a_Y^{-1}f_{small, *}\mathcal{F} \\
& \to 
\epsilon_{Y, *}f_{big, ph, *}a_X^{-1}\mathcal{F} \\
& =
f_{big, \etale, *} \epsilon_{X, *}a_X^{-1}\mathcal{F} \\
& =
f_{big, \etale, *}\pi_X^{-1}\mathcal{F}
\end{align*}
是同构。这正是引理 \ref{spaces-more-cohomology-lemma-compare-higher-direct-image-proper}
第 (1) 部分。

\medskip\noindent
为证明 (2)，我们使用
\begin{align*}
R\epsilon_{Y, *}Rf_{big, ph, *}a_X^{-1}K
& =
Rf_{big, \etale, *}R\epsilon_{X, *}a_X^{-1}K \\
& =
Rf_{big, \etale, *}\pi_X^{-1}K \\
& =
\pi_Y^{-1}Rf_{small, *}K \\
& =
R\epsilon_{Y, *} a_Y^{-1}Rf_{small, *}K
\end{align*}
% P10-E0170
第一个等式来自引理 \ref{spaces-more-cohomology-lemma-push-pull-ph-etale} 中的交换图以及
《站点上的上同调》引理
\ref{sites-cohomology-lemma-derived-pushforward-composition}.
% P10-E0171
第二个等式是引理 \ref{spaces-more-cohomology-lemma-cohomological-descent-etale-ph}。
第三个等式是引理 \ref{spaces-more-cohomology-lemma-compare-higher-direct-image-proper} 第 (2) 部分。
第四个等式再次使用引理 \ref{spaces-more-cohomology-lemma-cohomological-descent-etale-ph}。
因此，基变换映射
$a_Y^{-1}(Rf_{small, *}K) \to Rf_{big, ph, *}(a_X^{-1}K)$
诱导同构
$$
R\epsilon_{Y, *}a_Y^{-1}Rf_{small, *}K \to
R\epsilon_{Y, *}Rf_{big, ph, *}a_X^{-1}K
$$
最后用以下观察完成证明：考虑映射
$\alpha : a_Y^{-1}L \to M$，其中 $L$ 属于 $D^+(Y_\etale)$，
其上同调层均为挠层，且 $M$ 属于 $D^+((\textit{Spaces}/Y)_{ph})$。
若 $R\epsilon_{Y, *}\alpha$ 是同构，则 $\alpha$ 也是同构。
事实上，我们对 $i$ 归纳证明 $H^i(\alpha)$ 是同构。
对所有充分小的 $i$，这显然成立。若对 $i \leq i_0$ 成立，
则由引理 \ref{spaces-more-cohomology-lemma-comparison-ph-etale} 以及在此范围内
$H^i(M) = a_Y^{-1}H^i(L)$，可知
对 $j > 0$ 且 $i \leq i_0$，有
$R^j\epsilon_{Y, *}H^i(M) = 0$。于是谱序列论证给出
$\epsilon_{Y, *}H^{i_0 + 1}(M) = H^{i_0 + 1}(R\epsilon_{Y, *}M)$。
从而
$\epsilon_{Y, *}H^{i_0 + 1}(M) = \pi_Y^{-1}H^{i_0 + 1}(L) =
\epsilon_{Y, *}a_Y^{-1}H^{i_0 + 1}(L)$.
这蕴含 $H^{i_0 + 1}(\alpha)$ 是同构，因为 $\epsilon_{Y, *}$
在底层预层上恒等，因而反映同构。证毕。
\end{proof}
% Generated mechanically from the frozen spaces-simplicial.tex target.
% Do not edit this generated file; edit the bound locale source instead.

% OK, start here.
%

\setcounter{chapter}{84}
\chapter{单纯空间}



\phantomsection
\label{spaces-simplicial-section-phantom}


\section{引言}
\label{spaces-simplicial-section-introduction}

\noindent
本章发展单纯拓扑空间、单纯赋环空间、单纯概形以及单纯代数空间的若干理论。
借助单纯空间理论，有时可以证明一些用其他方法看来难以证明的
局部到整体原理。若干应用实例见论文
\cite{faltings_finiteness}、\cite{Kiehl} 和 \cite{HodgeIII}。

\medskip\noindent
全文假定读者熟悉《单纯方法》一章的基本概念与结果；参见
《单纯对象》第 \ref{simplicial-section-introduction} 节。
特别地，对于单纯对象，我们仍记作 $X$ 而不记作 $X_\bullet$。








\section{单纯拓扑空间}
\label{spaces-simplicial-section-simplicial-top}

\noindent
所谓{\it 单纯空间}，是以拓扑空间为对象、以连续映射为态射的范畴中的
单纯对象。（我们用“单纯代数空间”指代数空间范畴中的单纯对象。）
可以把单纯空间 $X$ 描绘如下：
$$
\xymatrix{
X_2
\ar@<2ex>[r]
\ar@<0ex>[r]
\ar@<-2ex>[r]
&
X_1
\ar@<1ex>[r]
\ar@<-1ex>[r]
\ar@<1ex>[l]
\ar@<-1ex>[l]
&
X_0
\ar@<0ex>[l]
}
$$
这里有两个态射 $d^1_0, d^1_1 : X_1 \to X_0$，以及一个态射
$s^0_0 : X_0 \to X_1$，等等。务须牢记，应把
$d^n_i : X_n \to X_{n - 1}$ 看成“忘掉第 $i$ 个坐标的投影”，
并把 $s^n_j : X_n \to X_{n + 1}$ 看成重复第 $j$ 个坐标的对角映射。

\medskip\noindent
设 $X$ 为单纯空间。按如下方式为 $X$ 配置一个站点
$X_{Zar}$\footnote{此记号与《站点》例
\ref{sites-example-site-topological} 及《拓扑》第
\ref{topologies-definition-big-small-Zariski} 定义中的记号相似。}：
\begin{enumerate}
\item $X_{Zar}$ 的对象是某个 $X_n$ 的开集 $U$，其中 $n$ 为某个整数；
\item $X_{Zar}$ 的态射 $U \to V$ 由一个
$\varphi : [m] \to [n]$ 给出，其中 $n, m$ 满足
$U \subset X_n$、$V \subset X_m$，且 $\varphi$ 满足
$X(\varphi)(U) \subset V$；
\item $X_{Zar}$ 中的覆盖 $\{U_i \to U\}$ 意味着
$U, U_i \subset X_n$ 均为开集，映射 $U_i \to U$ 由
$\text{id} : [n] \to [n]$ 给出，并且 $U = \bigcup U_i$。
\end{enumerate}
特别注意，若 $X_{Zar}$ 的态射 $U \to V$ 由 $\varphi$ 给出，
则 $X(\varphi) : X_n \to X_m$ 确实诱导拓扑空间的连续映射 $U \to V$。

\noindent
显然，上述构造是为任意拓扑空间图式配置站点这一构造的特例。
我们陈述相应的例行引理。

\begin{lemma}
\label{spaces-simplicial-lemma-simplicial-site}
设 $X$ 为单纯空间。则上述定义的 $X_{Zar}$ 是一个站点。
\end{lemma}

\begin{proof}
略。
\end{proof}

\noindent
设 $X$ 为单纯空间，$\mathcal{F}$ 为 $X_{Zar}$ 上的层。
由覆盖的定义可知，把 $\mathcal{F}$ 限制到 $X_n$ 的开集上，
便得到拓扑空间 $X_n$ 上的层 $\mathcal{F}_n$。
对每个 $\varphi : [m] \to [n]$，当 $U \subset X_n$、
$V \subset X_m$ 且 $X(\varphi)(U) \subset V$ 时，
$\mathcal{F}$ 的限制映射定义一个 $X(\varphi)$-映射
$\mathcal{F}(\varphi) : \mathcal{F}_m \to \mathcal{F}_n$；参见
《层》第 \ref{sheaves-definition-f-map} 定义。
此外，给定 $\varphi : [m] \to [n]$ 和 $\psi : [l] \to [m]$，有
$$
\mathcal{F}(\varphi) \circ \mathcal{F}(\psi) =
\mathcal{F}(\varphi \circ \psi)
$$
(左边使用 $f$-映射的复合；参见《层》第
\ref{sheaves-definition-composition-f-maps} 定义）。
显然，逆命题也成立：若有如上的系统
$(\{\mathcal{F}_n\}_{n \geq 0},
\{\mathcal{F}(\varphi)\}_{\varphi \in \text{箭头}(\Delta)})$
满足所示等式，则可得到 $X_{Zar}$ 上的一个层。

\begin{lemma}
\label{spaces-simplicial-lemma-describe-sheaves-simplicial-site}
设 $X$ 为单纯空间。下列两个范畴等价：
\begin{enumerate}
\item $\Sh(X_{Zar})$；
\item 上述系统 $(\mathcal{F}_n, \mathcal{F}(\varphi))$ 的范畴。
\end{enumerate}
\end{lemma}

\begin{proof}
见上文讨论。
\end{proof}

\begin{lemma}
\label{spaces-simplicial-lemma-simplicial-space-site-functorial}
设 $f : Y \to X$ 为单纯空间的态射。函子 $u : X_{Zar} \to Y_{Zar}$
把开集 $U \subset X_n$ 映到开集 $f_n^{-1}(U) \subset Y_n$；
它定义站点态射 $f_{Zar} : Y_{Zar} \to X_{Zar}$。
\end{lemma}

\begin{proof}
显然 $u$ 是连续函子，故得到函子 $f_{Zar, *} = u^s$ 和
$f_{Zar}^{-1} = u_s$；参见《站点》第 \ref{sites-section-morphism-sites} 节。
为证明所得对象是站点态射，须证明 $u_s$ 正合。我们用《站点》引理
\ref{sites-lemma-directed-morphism} 来证明这一点。
设 $V \subset Y_n$ 为开子集。范畴 $\mathcal{I}_V^u$（参见《站点》第
\ref{sites-section-functoriality-PSh} 节）的对象是二元组 $(U, \varphi)$，
其中 $\varphi : [m] \to [n]$，$U \subset X_m$ 为开集，且
$Y(\varphi)(V) \subset f_m^{-1}(U)$。态射
$(U, \varphi) \to (U', \varphi')$ 则由满足
$X(\psi)(U) \subset U'$ 和 $\varphi \circ \psi = \varphi'$ 的
$\psi : [m'] \to [m]$ 给出。我们要证明 $\mathcal{I}_V^u$ 是余滤范畴。

\medskip\noindent
验证《范畴》第 \ref{categories-definition-codirected} 定义中的条件。
由于 $(X_n, \text{id}_{[n]})$ 是对象，条件 (1) 成立。
设 $(U, \varphi)$ 为对象。条件
$Y(\varphi)(V) \subset f_m^{-1}(U)$ 等价于
$V \subset f_n^{-1}(X(\varphi)^{-1}(U))$。因此，取
$\psi = \varphi$ 便得到态射
$(X(\varphi)^{-1}(U), \text{id}_{[n]}) \to (U, \varphi)$。
再者，对任意两个形如 $(U, \text{id}_{[n]})$ 和
$(U', \text{id}_{[n]})$ 的对象，对象
$(U \cap U', \text{id}_{[n]})$ 同时映到二者，故条件 (2) 成立。
为验证条件 (3)，设有两个由 $\psi, \psi' : [m'] \to [m]$ 给出的态射
$a, a': (U, \varphi) \to (U', \varphi')$。
把二者预复合以由 $\varphi$ 给出的态射
$(X(\varphi)^{-1}(U), \text{id}_{[n]}) \to (U, \varphi)$，
便使 $a,a'$ 相等，因为
$\varphi \circ \psi = \varphi' = \varphi \circ \psi'$。
证毕。
\end{proof}

\begin{lemma}
\label{spaces-simplicial-lemma-describe-functoriality}
设 $f : Y \to X$ 为单纯空间的态射。用引理
\ref{spaces-simplicial-lemma-describe-sheaves-simplicial-site} 对层的描述，
引理 \ref{spaces-simplicial-lemma-simplicial-space-site-functorial} 中的态射
$f_{Zar}$ 可描述如下：
\begin{enumerate}
\item 若 $\mathcal{G}$ 是 $Y$ 上的层，则
$(f_{Zar, *}\mathcal{G})_n = f_{n, *}\mathcal{G}_n$。
\item 若 $\mathcal{F}$ 是 $X$ 上的层，则
$(f_{Zar}^{-1}\mathcal{F})_n = f_n^{-1}\mathcal{F}_n$。
\end{enumerate}
\end{lemma}

\begin{proof}
第一部分由定义立即得到。对于第二部分，注意在引理
\ref{spaces-simplicial-lemma-simplicial-space-site-functorial} 的证明中已经证明：
对开集 $V \subset Y_n$，范畴 $(\mathcal{I}_V^u)^{opp}$ 包含一个余终子范畴，
其对象是满足 $f_n^{-1}(U) \supset V$ 的开集 $U \subset X_n$，
态射由包含映射给出。因此，把 $u_p\mathcal{F}$ 限制到 $Y_n$ 的开集上，
所得预层就是《层》引理 \ref{sheaves-lemma-pullback-presheaves}
中定义的 $f_{n, p}\mathcal{F}_n$。
由于 $f_{Zar}^{-1}\mathcal{F} = u_s\mathcal{F}$ 是
$u_p\mathcal{F}$ 的层化，层化只使用覆盖，而 $Y_{Zar}$ 中的覆盖只使用
同一个 $Y_n$ 上开集之间的包含映射，结论便由
$f_n^{-1}\mathcal{F}_n$ 相应地是 $f_{n, p}\mathcal{F}_n$ 的层化而得；
参见《层》第 \ref{sheaves-section-presheaves-functorial} 节。
\end{proof}

\noindent
设 $X$ 为拓扑空间。在《站点》例
\ref{sites-example-site-topological} 中，我们用 $X_{Zar}$ 表示如下站点：
对象为 $X$ 的开集，态射为包含映射，覆盖为开覆盖。
我们把拓扑斯 $\Sh(X_{Zar})$ 与 $X$ 上的层范畴等同。

\begin{lemma}
\label{spaces-simplicial-lemma-restriction-to-components}
设 $X$ 为单纯空间。函子
$X_{n, Zar} \to X_{Zar}$，$U \mapsto U$，既连续又余连续。
相应的拓扑斯态射 $g_n : \Sh(X_n) \to \Sh(X_{Zar})$ 满足：
\begin{enumerate}
\item $g_n^{-1}$ 把 $X$ 上的层 $\mathcal{F}$ 映到 $X_n$ 上的层
$\mathcal{F}_n$；
\item $g_n^{-1} : \Sh(X_{Zar}) \to \Sh(X_n)$ 有左伴随 $g^{Sh}_{n!}$；
\item $g^{Sh}_{n!}$ 与有限连通极限交换；
\item $g_n^{-1} : \textit{Ab}(X_{Zar}) \to \textit{Ab}(X_n)$
有左伴随 $g_{n!}$；
\item $g_{n!}$ 正合。
\end{enumerate}
\end{lemma}

\begin{proof}
除命题中所述性质外，范畴 $X_{n, Zar}$ 还有纤维积与等化子，
且该函子与它们交换（注意 $X_{Zar}$ 并非具有所有纤维积）。
因此，本引理由《站点》第 \ref{sites-section-cocontinuous-functors}、
\ref{sites-section-cocontinuous-morphism-topoi}
节以及《站点上的模》第
\ref{sites-modules-section-exactness-lower-shriek} 节的讨论得到。
更确切地说，使用《站点》引理
\ref{sites-lemma-cocontinuous-morphism-topoi}、
\ref{sites-lemma-when-shriek}、
\ref{sites-lemma-preserve-equalizers}，以及《站点上的模》引理
\ref{sites-modules-lemma-g-shriek-adjoint}、
\ref{sites-modules-lemma-exactness-lower-shriek}。
\end{proof}

\begin{lemma}
\label{spaces-simplicial-lemma-restriction-injective-to-component}
设 $X$ 为单纯空间。若 $\mathcal{I}$ 是 $X_{Zar}$ 上的内射阿贝尔层，
则 $\mathcal{I}_n$ 是 $X_n$ 上的内射阿贝尔层。
\end{lemma}

\begin{proof}
这由《同调代数》引理 \ref{homology-lemma-adjoint-preserve-injectives}
和引理 \ref{spaces-simplicial-lemma-restriction-to-components} 得到。
\end{proof}

\begin{lemma}
\label{spaces-simplicial-lemma-restriction-to-components-functorial}
设 $f : Y \to X$ 为单纯空间的态射。则
$$
\xymatrix{
\Sh(Y_n) \ar[d] \ar[r]_{f_n} & \Sh(X_n) \ar[d] \\
\Sh(Y_{Zar}) \ar[r]^{f_{Zar}} & \Sh(X_{Zar})
}
$$
是拓扑斯的交换图。
\end{lemma}

\begin{proof}
直接由引理 \ref{spaces-simplicial-lemma-describe-functoriality} 和
\ref{spaces-simplicial-lemma-restriction-to-components} 中对拉回函子的描述得到。
\end{proof}

\begin{lemma}
\label{spaces-simplicial-lemma-augmentation}
设 $Y$ 为单纯空间，$a : Y \to X$ 为一个增广
（《单纯对象》定义 \ref{simplicial-definition-augmentation}）。
设 $a_n : Y_n \to X$ 为相应的拓扑空间态射。存在典范拓扑斯态射
$$
a : \Sh(Y_{Zar}) \to \Sh(X)
$$
满足下列性质：
\begin{enumerate}
\item $a^{-1}\mathcal{F}$ 限制到 $Y_n$ 后是 $a_n^{-1}\mathcal{F}$；
\item 对所有 $\varphi : [m] \to [n]$，有
$a_m \circ Y(\varphi) = a_n$；
\item 作为拓扑斯态射有 $a \circ g_n = a_n$，其中 $g_n$ 如引理
\ref{spaces-simplicial-lemma-restriction-to-components}；
\item 对 $\mathcal{G} \in \Sh(Y_{Zar})$，$a_*\mathcal{G}$ 是两个映射
$a_{0, *}\mathcal{G}_0 \to a_{1, *}\mathcal{G}_1$ 的等化子。
\end{enumerate}
\end{lemma}

\begin{proof}
在任意范畴中，单纯对象的增广均满足 (2)。因此
$Y(\varphi)^{-1} a_m^{-1} \mathcal{F} = a_n^{-1}\mathcal{F}$，
这定义一个从 $a_m^{-1}\mathcal{F}$ 到 $a_n^{-1}\mathcal{F}$ 的
$Y(\varphi)$-映射。于是可用 (1) 定义 $a^{-1}\mathcal{F}$
（使用引理 \ref{spaces-simplicial-lemma-describe-sheaves-simplicial-site}），并用 (4)
定义 $a_*$。若这定义拓扑斯态射，则由构造有
$g_n^{-1} \circ a^{-1} = a_n^{-1}$，从而 (3) 成立。
为验证 $a$ 是拓扑斯态射，须证明 $a^{-1}$ 左伴随于 $a_*$，并证明
$a^{-1}$ 正合。后一事实由各函子 $a_n^{-1}$ 的正合性立即得到。

\medskip\noindent
设 $\mathcal{F}$ 是 $\Sh(X)$ 的对象，$\mathcal{G}$ 是
$\Sh(Y_{Zar})$ 的对象。给定
$\beta : a^{-1}\mathcal{F} \to \mathcal{G}$，考察其分量
$\beta_n : a_n^{-1}\mathcal{F} \to \mathcal{G}_n$。
这些映射伴随于映射
$\beta_n : \mathcal{F} \to a_{n, *}\mathcal{G}_n$。
与单纯结构的相容性说明 $\beta_0$ 的像落在 $a_*\mathcal{G}$ 中。
反之，设给定映射 $\alpha : \mathcal{F} \to a_*\mathcal{G}$。
对任意 $n$，选择 $\varphi : [0] \to [n]$，并考察复合
$$
\mathcal{F} \xrightarrow{\alpha} a_*\mathcal{G}
\to a_{0, *}\mathcal{G}_0 \xrightarrow{\mathcal{G}(\varphi)}
a_{n, *}\mathcal{G}_n
$$
这些复合伴随于映射 $a_n^{-1}\mathcal{F} \to \mathcal{G}_n$，
后者定义层态射 $a^{-1}\mathcal{F} \to \mathcal{G}$。
略去验证上述构造给出互逆双射
$$
\Mor_{\Sh(Y_{Zar})}(a^{-1}\mathcal{F}, \mathcal{G}) =
\Mor_{\Sh(X)}(\mathcal{F}, a_*\mathcal{G})
$$
由此证毕。值得注意的是，这个拓扑斯态射并不对应于定义这些拓扑斯的
站点之间任何显然的几何函子。
\end{proof}

\begin{lemma}
\label{spaces-simplicial-lemma-simplicial-resolution-Z}
设 $X$ 为单纯拓扑空间。$X_{Zar}$ 上的阿贝尔预层复形
$$
\ldots \to \mathbf{Z}_{X_2} \to \mathbf{Z}_{X_1} \to \mathbf{Z}_{X_0}
$$
以 $\sum (-1)^i d^n_i$ 为边界；它是常值预层 $\mathbf{Z}$ 的消解。
\end{lemma}

\begin{proof}
设 $U \subset X_m$ 为 $X_{Zar}$ 的对象。上述复形在 $U$ 上的值是
阿贝尔群复形
$$
\ldots \to
\mathbf{Z}[\Mor_\Delta([2], [m])] \to
\mathbf{Z}[\Mor_\Delta([1], [m])] \to
\mathbf{Z}[\Mor_\Delta([0], [m])]
$$
换言之，这是单纯集 $\Delta[m]$ 上自由阿贝尔群所伴随的复形；参见
《单纯对象》例 \ref{simplicial-example-simplex-simplicial-set}。
由于 $\Delta[m]$ 与 $\Delta[0]$ 同伦等价（参见《单纯对象》例
\ref{simplicial-example-simplex-contractible}），且“取自由阿贝尔群”是函子，
上述复形与 $\Delta[0]$ 上的自由阿贝尔群同伦等价
（《单纯对象》评注 \ref{simplicial-remark-homotopy-better} 及引理
\ref{simplicial-lemma-homotopy-equivalence-s-N}）。
该复形在正次数无环，在次数 $0$ 等于 $\mathbf{Z}$。
\end{proof}

\begin{lemma}
\label{spaces-simplicial-lemma-simplicial-sheaf-cohomology}
设 $X$ 为单纯拓扑空间，$\mathcal{F}$ 为 $X$ 上的阿贝尔层。
存在谱序列 $(E_r, d_r)_{r \geq 0}$，其中
$$
E_1^{p, q} = H^q(X_p, \mathcal{F}_p)
$$
它收敛到 $H^{p + q}(X_{Zar}, \mathcal{F})$，并且关于
$\mathcal{F}$ 具有函子性。
\end{lemma}

\begin{proof}
设 $\mathcal{F} \to \mathcal{I}^\bullet$ 为内射消解。考察项为
$$
A^{p, q} = \mathcal{I}^q(X_p)
$$
的双复形，其第一微分是沿映射 $d^{p + 1}_i$ 所给映射
$\mathcal{I}_p^q \to \mathcal{I}_{p + 1}^q$ 的交错和；参见引理
\ref{spaces-simplicial-lemma-describe-sheaves-simplicial-site}。注意
$$
A^{p, q} = \Gamma(X_p, \mathcal{I}_p^q) =
\Mor_{\textit{PSh}}(h_{X_p}, \mathcal{I}^q) =
\Mor_{\textit{PAb}}(\mathbf{Z}_{X_p}, \mathcal{I}^q)
$$
因此，由引理 \ref{spaces-simplicial-lemma-simplicial-resolution-Z} 和《站点上的上同调》引理
\ref{sites-cohomology-lemma-injective-abelian-sheaf-injective-presheaf}，
双复形的各行在正次数正合，并在次数 $0$ 给出
$\Gamma(X_{Zar}, \mathcal{I}^q)$。另一方面，由于限制函子正合
（引理 \ref{spaces-simplicial-lemma-restriction-to-components}），映射
$$
\mathcal{F}_p \to \mathcal{I}_p^\bullet
$$
是消解。各层 $\mathcal{I}_p^q$ 是 $X_p$ 上的内射阿贝尔层
（引理 \ref{spaces-simplicial-lemma-restriction-injective-to-component}）。
故各列的上同调计算群 $H^q(X_p, \mathcal{F}_p)$。
最后应用《同调代数》引理 \ref{homology-lemma-first-quadrant-ss} 和
\ref{homology-lemma-double-complex-gives-resolution} 即得结论。
\end{proof}

\begin{lemma}
\label{spaces-simplicial-lemma-augmentation-pushforward-higher-direct-image}
设 $X$ 为单纯空间，$a : X \to Y$ 为增广，并设 $\mathcal{F}$ 为
$X_{Zar}$ 上的阿贝尔层。则 $R^na_*\mathcal{F}$ 是预层
$$
V \longmapsto H^n((X \times_Y V)_{Zar}, \mathcal{F}|_{(X \times_Y V)_{Zar}})
$$
所伴随的层。
\end{lemma}

\begin{proof}
这是《上同调》引理 \ref{cohomology-lemma-describe-higher-direct-images}
或《站点上的上同调》引理
\ref{sites-cohomology-lemma-higher-direct-images} 的类似结果；
我们强烈建议读者跳过此证明。选取 $X_{Zar}$ 上
$\mathcal{F}$ 的内射消解并使用定义，只须证明：
(1) 把 $X_{Zar}$ 上的内射阿贝尔层限制到
$(X \times_Y V)_{Zar}$ 后仍是内射阿贝尔层；以及
(2) $a_*\mathcal{F}$ 等于规则
$$
V \longmapsto H^0((X \times_Y V)_{Zar}, \mathcal{F}|_{(X \times_Y V)_{Zar}})
$$
(2) 由下列事实得到：
\begin{enumerate}
\item[(2a)] 由引理 \ref{spaces-simplicial-lemma-augmentation}，$a_*\mathcal{F}$ 是两个映射
$a_{0, *}\mathcal{F}_0 \to a_{1, *}\mathcal{F}_1$
的等化子；
\item[(2b)]
$a_{0, *}\mathcal{F}_0(V) =
H^0(a_0^{-1}(V), \mathcal{F}_0)$ 且
$a_{1, *}\mathcal{F}_1(V) = H^0(a_1^{-1}(V), \mathcal{F}_1)$；
\item[(2c)] $X_0 \times_Y V = a_0^{-1}(V)$ 且
$X_1 \times_Y V = a_1^{-1}(V)$；
\item[(2d)]
$H^0((X \times_Y V)_{Zar}, \mathcal{F}|_{(X \times_Y V)_{Zar}})$
是两个映射
$H^0(X_0 \times_Y V, \mathcal{F}_0) \to H^0(X_1 \times_Y V, \mathcal{F}_1)$
的等化子，例如可由引理 \ref{spaces-simplicial-lemma-simplicial-sheaf-cohomology} 得到。
\end{enumerate}
定义限制函子
$\textit{Ab}(X_{Zar}) \to \textit{Ab}((X \times_Y V)_{Zar})$
的正合左伴随
$j_! : \textit{Ab}((X \times_Y V)_{Zar}) \to \textit{Ab}(X_{Zar})$
（零延拓），再应用《同调代数》引理
\ref{homology-lemma-adjoint-preserve-injectives}，即可得到 (1)。
\end{proof}

\noindent
设 $X$ 为拓扑空间。用 $X_\bullet$ 表示取常值 $X$ 的常值单纯拓扑空间。
由引理 \ref{spaces-simplicial-lemma-describe-sheaves-simplicial-site}，
$X_{\bullet, Zar}$ 上的层等同于 $X$ 上层范畴中的余单纯对象。

\begin{lemma}
\label{spaces-simplicial-lemma-constant-simplicial-space}
设 $X$ 为拓扑空间，$X_\bullet$ 为取常值 $X$ 的常值单纯拓扑空间。函子
$$
X_{\bullet, Zar} \longrightarrow X_{Zar},\quad
U \longmapsto U
$$
既连续又余连续，并定义拓扑斯态射
$g : \Sh(X_{\bullet, Zar}) \to \Sh(X)$，且 $g^{-1}$ 有左伴随 $g_!$。
有：
\begin{enumerate}
\item $g^{-1}$ 把 $X$ 上的层映到 $X$ 上的常值余单纯层；
\item $g_!$ 把 $X_{\bullet, Zar}$ 上的层 $\mathcal{F}$ 映到层
$\mathcal{F}_0$；
\item $g_*$ 把 $X_{\bullet, Zar}$ 上的层 $\mathcal{F}$ 映到两个映射
$\mathcal{F}_0 \to \mathcal{F}_1$ 的等化子。
\end{enumerate}
\end{lemma}

\begin{proof}
关于该函子的陈述可直接验证。$g$ 与 $g_!$ 的存在性由《站点》引理
\ref{sites-lemma-cocontinuous-morphism-topoi} 和
\ref{sites-lemma-when-shriek} 得到。$g^{-1}$ 的描述直接来自《站点》引理
\ref{sites-lemma-when-shriek}。所给函子分别是 $g^{-1}$ 的右、左伴随，
因而得到 $g_*$ 与 $g_!$ 的描述。
\end{proof}








\section{单纯站点与拓扑斯}
\label{spaces-simplicial-section-simplicial-sites}

\noindent
自然的做法似乎是把{\it 单纯站点}定义为如下（大）范畴中的单纯对象：
其对象是站点，态射是站点态射。参见《站点》定义
\ref{sites-definition-site}、\ref{sites-definition-morphism-sites}；
态射的复合如《站点》引理 \ref{sites-lemma-composition-morphisms-sites}。
不过，也可考虑下列变体：(a) 改用站点之间的余连续函子
（参见《站点》第 \ref{sites-section-cocontinuous-functors}、
\ref{sites-section-cocontinuous-morphism-topoi} 节）；
(b) 在适当的站点 $2$-范畴中工作，并引入站点态射之间的
$2$-态射概念；(c) 在由余连续函子构造的 $2$-范畴中工作。
我们不选取其中一种变体作为定义，而只按需要发展理论。

\medskip\noindent
当然，{\it 单纯拓扑斯}大概应定义为从 $\Delta^{opp}$ 到拓扑斯
$2$-范畴的伪函子。参见《范畴》定义
\ref{categories-definition-functor-into-2-category} 以及《站点》第
\ref{sites-section-topoi}、\ref{sites-section-2-category} 节。
我们将尽量避免直接处理这种庞然大物。

\medskip\noindent
{\bf 情形 A。}
设 $\mathcal{C}$ 是如下范畴中的单纯对象：其对象为站点，态射为站点态射。
这意味着，对 $\Delta$ 的每个态射 $\varphi : [m] \to [n]$，
都有站点态射 $f_\varphi : \mathcal{C}_n \to \mathcal{C}_m$。
该态射由反方向的连续函子给出，记为
$u_\varphi : \mathcal{C}_m \to \mathcal{C}_n$。

\begin{lemma}
\label{spaces-simplicial-lemma-simplicial-site-site}
设 $\mathcal{C}$ 为站点范畴中的单纯对象。沿用上述记号，按如下方式构造站点
$\mathcal{C}_{total}$：
\begin{enumerate}
\item $\mathcal{C}_{total}$ 的对象是某个 $\mathcal{C}_n$ 的对象 $U$，
其中 $n$ 为某个整数；
\item $\mathcal{C}_{total}$ 的态射 $(\varphi, f) : U \to V$ 由映射
$\varphi : [m] \to [n]$ 以及 $\mathcal{C}_n$ 中的态射
$f : U \to u_\varphi(V)$ 给出，其中
$U \in \Ob(\mathcal{C}_n)$、$V \in \Ob(\mathcal{C}_m)$；
\item $\mathcal{C}_{total}$ 中的覆盖
$\{(\text{id}, f_i) :  U_i \to U\}$ 由某个 $n$ 以及
$\mathcal{C}_n$ 中的覆盖 $\{f_i : U_i \to U\}$ 给出。
\end{enumerate}
\end{lemma}

\begin{proof}
$(\varphi, f) : U \to V$ 与 $(\psi, g) : V \to W$ 的复合为
$(\varphi \circ \psi, u_\varphi(g) \circ f)$。
这里使用了 $u_\varphi \circ u_\psi = u_{\varphi \circ \psi}$。

\medskip\noindent
设 $\{(\text{id}, f_i) :  U_i \to U\}$ 为 (3) 中的覆盖，
$(\varphi, g) : W \to U$ 为态射，其中
$W \in \Ob(\mathcal{C}_m)$。我们断言
$$
W \times_{(\varphi, g), U, (\text{id}, f_i)} U_i =
W \times_{g, u_\varphi(U), u_\varphi(f_i)} u_\varphi(U_i)
$$
在范畴 $\mathcal{C}_{total}$ 中成立。此式确有意义：按我们的站点态射定义，
$u_\varphi$ 把覆盖变成覆盖，故所需的纤维积在 $\mathcal{C}_m$ 中存在。
同样的推理证明该断言（略去细节）。因此覆盖族在基变换下稳定；
站点的其他公理立即成立。
\end{proof}

\noindent
{\bf 情形 B。}
设 $\mathcal{C}$ 是如下范畴中的单纯对象：其对象为站点，态射为余连续函子。
这意味着，对 $\Delta$ 的每个态射 $\varphi : [m] \to [n]$，
都有余连续函子 $u_\varphi : \mathcal{C}_n \to \mathcal{C}_m$。
相应的拓扑斯态射记为
$f_\varphi : \Sh(\mathcal{C}_n) \to \Sh(\mathcal{C}_m)$。

\begin{lemma}
\label{spaces-simplicial-lemma-simplicial-cocontinuous-site}
设 $\mathcal{C}$ 为以站点为对象、以余连续函子为态射的范畴中的单纯对象。
沿用上述记号，假设函子
$u_\varphi : \mathcal{C}_n \to \mathcal{C}_m$ 具有《站点》评注
\ref{sites-remark-cartesian-cocontinuous} 的性质 $P$。
则可按如下方式构造站点 $\mathcal{C}_{total}$：
\begin{enumerate}
\item $\mathcal{C}_{total}$ 的对象是某个 $\mathcal{C}_n$ 的对象 $U$，
其中 $n$ 为某个整数；
\item $\mathcal{C}_{total}$ 的态射 $(\varphi, f) : U \to V$ 由映射
$\varphi : [m] \to [n]$ 以及 $\mathcal{C}_m$ 中的态射
$f : u_\varphi(U) \to V$ 给出，其中
$U \in \Ob(\mathcal{C}_n)$、$V \in \Ob(\mathcal{C}_m)$；
\item $\mathcal{C}_{total}$ 中的覆盖
$\{(\text{id}, f_i) :  U_i \to U\}$ 由某个 $n$ 以及
$\mathcal{C}_n$ 中的覆盖 $\{f_i : U_i \to U\}$ 给出。
\end{enumerate}
\end{lemma}

\begin{proof}
$(\varphi, f) : U \to V$ 与 $(\psi, g) : V \to W$ 的复合为
$(\varphi \circ \psi, g \circ u_\psi(f))$。
这里使用了 $u_\psi \circ u_\varphi = u_{\varphi \circ \psi}$。

\medskip\noindent
设 $\{(\text{id}, f_i) :  U_i \to U\}$ 为 (3) 中的覆盖，
$(\varphi, g) : W \to U$ 为态射，其中
$W \in \Ob(\mathcal{C}_m)$。我们断言
$$
W \times_{(\varphi, g), U, (\text{id}, f_i)} U_i =
W \times_{g, U, f_i} U_i
$$
在范畴 $\mathcal{C}_{total}$ 中成立；右边是《站点》评注
\ref{sites-remark-cartesian-cocontinuous} 所定义的
$\mathcal{C}_m$ 的对象，并由性质 $P$ 保证存在。
这类纤维积与函子复合的相容性推出该断言（略去细节）。
由性质 $P$，族 $\{W \times_{g, U, f_i} U_i \to W\}$ 是
$\mathcal{C}_m$ 的覆盖，故覆盖族在基变换下稳定；
站点的其他公理立即成立。
\end{proof}

\begin{situation}
\label{spaces-simplicial-situation-simplicial-site}
此处考虑下列两种情形之一：
\begin{enumerate}
\item[(A)] $\mathcal{C}$ 是以站点为对象、以站点态射为态射的范畴中的
单纯对象。对 $\Delta$ 的每个态射 $\varphi : [m] \to [n]$，
都有由连续函子 $u_\varphi : \mathcal{C}_m \to \mathcal{C}_n$ 给出的
站点态射 $f_\varphi : \mathcal{C}_n \to \mathcal{C}_m$。
\item[(B)] $\mathcal{C}$ 是以站点为对象、以满足《站点》评注
\ref{sites-remark-cartesian-cocontinuous} 的性质 $P$ 的余连续函子为态射的
范畴中的单纯对象。对 $\Delta$ 的每个态射
$\varphi : [m] \to [n]$，都有余连续函子
$u_\varphi : \mathcal{C}_n \to \mathcal{C}_m$，它诱导拓扑斯态射
$f_\varphi : \Sh(\mathcal{C}_n) \to \Sh(\mathcal{C}_m)$。
\end{enumerate}
照例，用 $f_\varphi^{-1}$ 和 $f_{\varphi, *}$ 分别表示拉回与推前。
用 $\mathcal{C}_{total}$ 表示引理 \ref{spaces-simplicial-lemma-simplicial-site-site}
（情形 A）或引理 \ref{spaces-simplicial-lemma-simplicial-cocontinuous-site}
（情形 B）中定义的站点。
\end{situation}

\noindent
设 $\mathcal{C}$ 如情形 \ref{spaces-simplicial-situation-simplicial-site}，
$\mathcal{F}$ 为 $\mathcal{C}_{total}$ 上的层。
由覆盖的定义可知，把 $\mathcal{F}$ 限制到 $\mathcal{C}_n$ 的对象上，
便定义站点 $\mathcal{C}_n$ 上的层 $\mathcal{F}_n$。
对每个 $\varphi : [m] \to [n]$，$\mathcal{F}$ 沿态射
$(\varphi, f) : U \to V$ 的限制映射定义下列集合中的元素
$\mathcal{F}(\varphi)$，其中 $U \in \Ob(\mathcal{C}_n)$、
$V \in \Ob(\mathcal{C}_m)$：
$$
\Mor_{\Sh(\mathcal{C}_m)}(\mathcal{F}_m, f_{\varphi, *}\mathcal{F}_n) =
\Mor_{\Sh(\mathcal{C}_n)}(f_\varphi^{-1}\mathcal{F}_m, \mathcal{F}_n)
$$
此外，给定 $\varphi : [m] \to [n]$ 和 $\psi : [l] \to [m]$，下列图交换：
$$
\vcenter{
\xymatrix{
\mathcal{F}_l \ar[rr]_{\mathcal{F}(\varphi \circ \psi)}
\ar[rd]_{\mathcal{F}(\psi)}
& & f_{\varphi \circ \psi, *} \mathcal{F}_n \\
& f_{\psi, *}\mathcal{F}_m \ar[ur]_{f_{\psi, *}\mathcal{F}(\varphi)}
}
}
\quad\text{且}\quad
\vcenter{
\xymatrix{
f_{\varphi \circ \psi}^{-1}\mathcal{F}_l
\ar[rr]_{\mathcal{F}(\varphi \circ \psi)}
\ar[rd]_{f_\varphi^{-1}\mathcal{F}(\psi)}
& & \mathcal{F}_n \\
& f_\varphi^{-1}\mathcal{F}_m \ar[ur]_{\mathcal{F}(\varphi)}
}
}
$$
显然，逆命题也成立：若有系统
$(\{\mathcal{F}_n\}_{n \geq 0},
\{\mathcal{F}(\varphi)\}_{\varphi \in \text{箭头}(\Delta)})$
满足上述交换性约束，则可得到 $\mathcal{C}_{total}$ 上的一个层。

\begin{lemma}
\label{spaces-simplicial-lemma-describe-sheaves-simplicial-site-site}
在情形 \ref{spaces-simplicial-situation-simplicial-site} 中，下列两个范畴等价：
\begin{enumerate}
\item $\Sh(\mathcal{C}_{total})$；
\item 上述系统 $(\mathcal{F}_n, \mathcal{F}(\varphi))$ 的范畴。
\end{enumerate}
特别地，拓扑斯 $\Sh(\mathcal{C}_{total})$ 仅依赖于各拓扑斯
$\Sh(\mathcal{C}_n)$ 及拓扑斯态射 $f_\varphi$。
\end{lemma}

\begin{proof}
见上文讨论。
\end{proof}

\begin{lemma}
\label{spaces-simplicial-lemma-restriction-to-components-site}
在情形 \ref{spaces-simplicial-situation-simplicial-site} 中，函子
$\mathcal{C}_n \to \mathcal{C}_{total}$，$U \mapsto U$，
既连续又余连续。相应的拓扑斯态射
$g_n : \Sh(\mathcal{C}_n) \to \Sh(\mathcal{C}_{total})$ 满足：
\begin{enumerate}
\item $g_n^{-1}$ 把 $\mathcal{C}_{total}$ 上的层 $\mathcal{F}$ 映到
$\mathcal{C}_n$ 上的层 $\mathcal{F}_n$；
\item $g_n^{-1} : \Sh(\mathcal{C}_{total}) \to \Sh(\mathcal{C}_n)$
有左伴随 $g^{Sh}_{n!}$；
\item 对 $\Sh(\mathcal{C}_n)$ 中的 $\mathcal{G}$，
$g_{n!}^{Sh}\mathcal{G}$ 在 $\mathcal{C}_m$ 上的限制为
$\coprod\nolimits_{\varphi : [n] \to [m]} f_\varphi^{-1}\mathcal{G}$；
\item $g_{n!}^{Sh}$ 与有限连通极限交换；
\item $g_n^{-1} : \textit{Ab}(\mathcal{C}_{total}) \to
\textit{Ab}(\mathcal{C}_n)$ 有左伴随 $g_{n!}$；
\item 对 $\textit{Ab}(\mathcal{C}_n)$ 中的 $\mathcal{G}$，
$g_{n!}\mathcal{G}$ 在 $\mathcal{C}_m$ 上的限制为
$\bigoplus\nolimits_{\varphi : [n] \to [m]} f_\varphi^{-1}\mathcal{G}$；
\item $g_{n!}$ 正合。
\end{enumerate}
\end{lemma}

\begin{proof}
情形 A。若 $\{U_i \to U\}_{i \in I}$ 是 $\mathcal{C}_n$ 中的覆盖，
则按定义（引理 \ref{spaces-simplicial-lemma-simplicial-site-site}），其像
$\{U_i \to U\}_{i \in I}$ 是 $\mathcal{C}_{total}$ 中的覆盖。
对 $\mathcal{C}_n$ 的态射 $V \to U$，$\mathcal{C}_n$ 中的纤维积
$V \times_U U_i$ 也是 $\mathcal{C}_{total}$ 中的纤维积
（由引理 \ref{spaces-simplicial-lemma-simplicial-site-site} 证明中的断言）。
因此该函子连续。另一方面，该函子在 $\mathcal{C}_n$ 中 $U$ 的覆盖
与 $\mathcal{C}_{total}$ 中 $U$ 的覆盖之间给出双射，故它也余连续。

\medskip\noindent
情形 B。若 $\{U_i \to U\}_{i \in I}$ 是 $\mathcal{C}_n$ 中的覆盖，
则按定义（引理 \ref{spaces-simplicial-lemma-simplicial-cocontinuous-site}），其像
$\{U_i \to U\}_{i \in I}$ 是 $\mathcal{C}_{total}$ 中的覆盖。
对 $\mathcal{C}_n$ 的态射 $V \to U$，$\mathcal{C}_n$ 中的纤维积
$V \times_U U_i$ 也是 $\mathcal{C}_{total}$ 中的纤维积
（由引理 \ref{spaces-simplicial-lemma-simplicial-cocontinuous-site} 证明中的断言）。
因此该函子连续。另一方面，该函子在 $\mathcal{C}_n$ 中 $U$ 的覆盖
与 $\mathcal{C}_{total}$ 中 $U$ 的覆盖之间给出双射，故它也余连续。

\medskip\noindent
至此，在情形 A 与 B 中，(1) 以及 $g^{Sh}_{n!}$、$g_{n!}$ 的存在性
由《站点》引理 \ref{sites-lemma-cocontinuous-morphism-topoi}、
\ref{sites-lemma-when-shriek} 和《站点上的模》引理
\ref{sites-modules-lemma-g-shriek-adjoint} 得到。

\medskip\noindent
(3) 的证明。设 $\mathcal{G}$ 为 $\mathcal{C}_n$ 上的层。
考察 $\mathcal{C}_{total}$ 上的层 $\mathcal{H}$，其次数 $m$ 部分是层
$$
\mathcal{H}_m = \coprod\nolimits_{\varphi : [n] \to [m]}
f_\varphi^{-1}\mathcal{G}
$$
即引理陈述 (3) 中给出的层。给定映射 $\psi : [m] \to [m']$，
映射 $\mathcal{H}(\psi) : f_\psi^{-1}\mathcal{H}_m \to \mathcal{H}_{m'}$
在各分量上由下列等同给出：
$$
f_\psi^{-1} f_\varphi^{-1} \mathcal{G} \to
f_{\psi \circ \varphi}^{-1}\mathcal{G}
$$
注意，给定 $\mathcal{C}_{total}$ 上层的映射
$\alpha : \mathcal{H} \to \mathcal{F}$，可得到映射
$\mathcal{G} \to \mathcal{F}_n$；它对应于把 $\alpha_n$ 限制到
$\mathcal{H}_n$ 中的分量 $\mathcal{G}$。反之，给定
$\mathcal{C}_n$ 上层的映射 $\beta : \mathcal{G} \to \mathcal{F}_n$，
可定义 $\alpha : \mathcal{H} \to \mathcal{F}$：令 $\alpha_m$
在各分量
$$
f_\varphi^{-1}\mathcal{G} \to \mathcal{F}_m
$$
上使用伴随于 $\mathcal{F}(\varphi) \circ f_\varphi^{-1}\beta$ 的映射。
略去验证这两个构造给出互逆映射
$$
\Mor_{\Sh(\mathcal{C}_n)}(\mathcal{G}, \mathcal{F}_n) =
\Mor_{\Sh(\mathcal{C}_{total})}(\mathcal{H}, \mathcal{F})
$$
故 $\mathcal{H} = g^{Sh}_{n!}\mathcal{G}$，如所求。

\medskip\noindent
(4) 的证明。若 $\mathcal{G}$ 是 $\mathcal{C}_n$ 上的阿贝尔层，
则完全仿照上文，只是把 $\mathcal{H}$ 定义为
$\mathcal{C}_{total}$ 上的阿贝尔层，其次数 $m$ 部分为
$$
\bigoplus\nolimits_{\varphi : [n] \to [m]} f_\varphi^{-1}\mathcal{G}
$$
转移映射完全如上定义。双射
$$
\Mor_{\textit{Ab}(\mathcal{C}_n)}(\mathcal{G}, \mathcal{F}_n) =
\Mor_{\textit{Ab}(\mathcal{C}_{total})}(\mathcal{H}, \mathcal{F})
$$
的证明与上文完全相同。故 $\mathcal{H} = g_{n!}\mathcal{G}$，如所求。

\medskip\noindent
$g^{Sh}_{n!}$ 与 $g_{n!}$ 的正合性由这些函子的所给公式得到。
\end{proof}

\begin{lemma}
\label{spaces-simplicial-lemma-restriction-injective-to-component-site}
\begin{slogan}
单纯站点上的内射阿贝尔层在每个分量上均内射
\end{slogan}
在情形 \ref{spaces-simplicial-situation-simplicial-site} 中，若 $\mathcal{I}$ 在
$\textit{Ab}(\mathcal{C}_{total})$ 中内射，则 $\mathcal{I}_n$ 在
$\textit{Ab}(\mathcal{C}_n)$ 中内射。若 $\mathcal{I}^\bullet$ 是
$\textit{Ab}(\mathcal{C}_{total})$ 中的 K-内射复形，则
$\mathcal{I}_n^\bullet$ 是 $\textit{Ab}(\mathcal{C}_n)$ 中的 K-内射复形。
\end{lemma}

\begin{proof}
第一条陈述由《同调代数》引理
\ref{homology-lemma-adjoint-preserve-injectives} 和引理
\ref{spaces-simplicial-lemma-restriction-to-components-site} 得到。
第二条陈述由《导出范畴》引理
\ref{derived-lemma-adjoint-preserve-K-injectives} 和引理
\ref{spaces-simplicial-lemma-restriction-to-components-site} 得到。
\end{proof}







\section{单纯站点的增广}
\label{spaces-simplicial-section-augmentation-simplicial-sites}

\noindent
我们继续采用第 \ref{spaces-simplicial-section-simplicial-sites} 节所述方法，
阐明该节情形 A 与 B 中增广的含义。

\begin{remark}
\label{spaces-simplicial-remark-augmentation-site}
在情形 \ref{spaces-simplicial-situation-simplicial-site} 中，所谓{\it 指向站点
$\mathcal{D}$ 的增广 $a_0$}，意指：
\begin{enumerate}
\item[(A)] $a_0 : \mathcal{C}_0 \to \mathcal{D}$ 是由连续函子
$u_0 : \mathcal{D} \to \mathcal{C}_0$ 给出的站点态射，且对所有
$\varphi, \psi : [0] \to [n]$，有
$u_\varphi \circ u_0 = u_\psi \circ u_0$；
\item[(B)] $a_0 : \Sh(\mathcal{C}_0) \to \Sh(\mathcal{D})$ 是由余连续函子
$u_0 : \mathcal{C}_0 \to \mathcal{D}$ 给出的拓扑斯态射，且对所有
$\varphi, \psi : [0] \to [n]$，有
$u_0 \circ u_\varphi = u_0 \circ u_\psi$。
\end{enumerate}
\end{remark}

\begin{lemma}
\label{spaces-simplicial-lemma-augmentation-site}
在情形 \ref{spaces-simplicial-situation-simplicial-site} 中，设 $a_0$ 为评注
\ref{spaces-simplicial-remark-augmentation-site} 所述指向站点 $\mathcal{D}$ 的增广。
则 $a_0$ 诱导：
\begin{enumerate}
\item 对所有 $n \geq 0$，拓扑斯态射
$a_n : \Sh(\mathcal{C}_n) \to \Sh(\mathcal{D})$；
\item 拓扑斯态射 $a : \Sh(\mathcal{C}_{total}) \to \Sh(\mathcal{D})$。
\end{enumerate}
它们满足：
\begin{enumerate}
\item 对所有 $\varphi : [m] \to [n]$，有
$a_m \circ f_\varphi = a_n$；
\item 若 $g_n : \Sh(\mathcal{C}_n) \to \Sh(\mathcal{C}_{total})$
如引理 \ref{spaces-simplicial-lemma-restriction-to-components-site}，则
$a \circ g_n = a_n$；
\item 对 $\mathcal{F} \in \Sh(\mathcal{C}_{total})$，
$a_*\mathcal{F}$ 是两个映射
$a_{0, *}\mathcal{F}_0 \to a_{1, *}\mathcal{F}_1$ 的等化子。
\end{enumerate}
\end{lemma}

\begin{proof}
情形 A。令 $u_n : \mathcal{D} \to \mathcal{C}_n$ 为诸函子
$u_\varphi \circ u_0$（$\varphi : [0] \to [n]$）的共同值。
则 $u_n$ 对应于站点态射 $a_n : \mathcal{C}_n \to \mathcal{D}$；
参见《站点》引理 \ref{sites-lemma-composition-morphisms-sites}。
同一引理说明，对所有 $\varphi : [m] \to [n]$，有
$a_m \circ f_\varphi = a_n$。

\medskip\noindent
情形 B。令 $u_n : \mathcal{C}_n \to \mathcal{D}$ 为诸函子
$u_0 \circ u_\varphi$（$\varphi : [0] \to [n]$）的共同值。
则 $u_n$ 余连续，因而定义拓扑斯态射
% P10-E0172
$a_n : \Sh(\mathcal{C}_n) \to \Sh(\mathcal{D)}$；参见
《站点》引理 \ref{sites-lemma-composition-cocontinuous}。
同一引理说明，对所有 $\varphi : [m] \to [n]$，有
$a_m \circ f_\varphi = a_n$。

\medskip\noindent
考察函子 $a^{-1} : \Sh(\mathcal{D}) \to \Sh(\mathcal{C}_{total})$。
它把集合层 $\mathcal{G}$ 映到层 $\mathcal{F} = a^{-1}\mathcal{G}$；
后者的分量为 $a_n^{-1}\mathcal{G}$，转移映射
$\mathcal{F}(\varphi)$ 为下列等同：
$$
f_\varphi^{-1}\mathcal{F}_m =
f_\varphi^{-1} a_m^{-1}\mathcal{G} =
a_n^{-1}\mathcal{G} =
\mathcal{F}_n
$$
其中 $\varphi : [m] \to [n]$；参见引理
\ref{spaces-simplicial-lemma-describe-sheaves-simplicial-site-site} 对
$\Sh(\mathcal{C}_{total})$ 的描述。由于各函子 $a_n^{-1}$ 正合，
$a^{-1}$ 也是正合函子。最后，把
$a_* : \Sh(\mathcal{C}_{total}) \to \Sh(\mathcal{D})$ 取为如下函子：
% P10-E0173
它把 $\Sh(\mathcal{D})$ 上的层 $\mathcal{F}$ 映到
$$
\xymatrix{
a_*\mathcal{F} \ar@{=}[r] &
\text{等化子}(a_{0, *}\mathcal{F}_0
\ar@<1ex>[r] \ar@<-1ex>[r] &
a_{1, *}\mathcal{F}_1)
}
$$
这里两个映射由两个映射 $\varphi : [0] \to [1]$ 经下式给出：
$$
a_{0, *}\mathcal{F}_0 \to
a_{0, *}f_{\varphi, *} f_\varphi^{-1}\mathcal{F}_0
\xrightarrow{\mathcal{F}(\varphi)}
a_{0, *}f_{\varphi, *} \mathcal{F}_1 = a_{1, *}\mathcal{F}_1
$$
其中第一个箭头来自 $1 \to f_{\varphi, *} f_\varphi^{-1}$。
% P10-E0174
用 $\mathcal{G}_\bullet$ 表示取常值 $\mathcal{G}$ 的常值余单纯层，
并用 $a_{\bullet, *}\mathcal{F}$ 表示次数 $n$ 分量为
$a_{n, *}\mathcal{F}_n$ 的余单纯层。
% P10-E0175
由拓扑斯态射 $a_n$ 的通常伴随关系可知，映射
$a^{-1}\mathcal{G} \to \mathcal{F}$ 等同于余单纯层的映射
$$
\mathcal{G}_\bullet \longrightarrow a_{\bullet, *}\mathcal{F}
$$
由《单纯对象》引理 \ref{simplicial-lemma-augmentation-howto} 的对偶，
这又等同于映射 $\mathcal{G} \to a_*\mathcal{F}$。
故 $a^{-1}$ 与 $a_*$ 互为伴随函子，从而得到拓扑斯态射
$a$\footnote{在情形 B 中，态射 $a$ 对应于余连续函子
$\mathcal{C}_{total} \to \mathcal{D}$；它把
$\mathcal{C}_n$ 中的 $U$ 映到 $u_n(U)$。}。
% P10-E0176
等式 $a \circ g_n = f_n$ 由定义立即得到。
\end{proof}



\section{单纯站点的态射}
\label{spaces-simplicial-section-morphism-simplicial-sites}

\noindent
我们继续采用第 \ref{spaces-simplicial-section-simplicial-sites} 节所述方法，
阐明该节情形 A 与 B 中单纯站点态射的含义。

\begin{remark}
\label{spaces-simplicial-remark-morphism-simplicial-sites}
设 $\mathcal{C}_n, f_\varphi, u_\varphi$ 和
$\mathcal{C}'_n, f'_\varphi, u'_\varphi$ 如情形
\ref{spaces-simplicial-situation-simplicial-site}。所谓{\it 单纯站点之间的态射 $h$}，意指：
\begin{enumerate}
\item[(A)] 站点态射 $h_n : \mathcal{C}_n \to \mathcal{C}'_n$，
使得对所有 $\varphi : [m] \to [n]$，作为站点态射有
$f'_\varphi \circ h_n = h_m \circ f_\varphi$；
\item[(B)] 余连续函子 $v_n : \mathcal{C}_n \to \mathcal{C}'_n$，
它诱导拓扑斯态射
$h_n : \Sh(\mathcal{C}_n) \to \Sh(\mathcal{C}'_n)$，
使得对所有 $\varphi : [m] \to [n]$，作为函子有
$u'_\varphi \circ v_n = v_m \circ u_\varphi$。
\end{enumerate}
在两种情形中，作为拓扑斯态射均有
$f'_\varphi \circ h_n = h_m \circ f_\varphi$；
情形 B 参见《站点》引理
\ref{sites-lemma-composition-cocontinuous}，情形 A 参见《站点》定义
\ref{sites-definition-composition-morphisms-sites}。
\end{remark}

\begin{lemma}
\label{spaces-simplicial-lemma-morphism-simplicial-sites}
设 $\mathcal{C}_n, f_\varphi, u_\varphi$ 和
$\mathcal{C}'_n, f'_\varphi, u'_\varphi$ 如情形
\ref{spaces-simplicial-situation-simplicial-site}，并设 $h$ 为评注
\ref{spaces-simplicial-remark-morphism-simplicial-sites} 所述的单纯站点之间的态射。
则得到拓扑斯态射
$$
h_{total} : \Sh(\mathcal{C}_{total}) \to \Sh(\mathcal{C}'_{total})
$$
以及交换图
$$
\xymatrix{
\Sh(\mathcal{C}_n) \ar[d]_{g_n} \ar[r]_{h_n} &
\Sh(\mathcal{C}'_n) \ar[d]^{g'_n} \\
\Sh(\mathcal{C}_{total}) \ar[r]^{h_{total}} &
\Sh(\mathcal{C}'_{total})
}
$$
此外，有 $(g'_n)^{-1} \circ h_{total, *} = h_{n, *} \circ g_n^{-1}$。
\end{lemma}

\begin{proof}
情形 A。设 $h_n$ 对应于连续函子
$v_n : \mathcal{C}'_n \to \mathcal{C}_n$。在次数 $n$ 使用 $v_n$，
即可定义函子 $v_{total} : \mathcal{C}'_{total} \to \mathcal{C}_{total}$。
它显然连续（参见引理 \ref{spaces-simplicial-lemma-simplicial-site-site} 中覆盖的定义）。令
$h_{total}^{-1} = v_{total, s} :
\Sh(\mathcal{C}'_{total}) \to \Sh(\mathcal{C}_{total})$，以及
$h_{total, *} = v_{total}^s = v_{total}^p :
\Sh(\mathcal{C}_{total}) \to \Sh(\mathcal{C}'_{total})$
为《站点》第 \ref{sites-section-continuous-functors}、
\ref{sites-section-morphism-sites} 节所构造并研究的一对伴随函子。
为证明 $h_{total}$ 是拓扑斯态射，还须验证 $h_{total}^{-1}$ 正合。
先注意
$(g'_n)^{-1} \circ h_{total, *} = h_{n, *} \circ g_n^{-1}$；
这可通过计算 $\mathcal{C}'_n$ 的对象 $U$ 上的截面立即看出。
因此，若按引理 \ref{spaces-simplicial-lemma-describe-sheaves-simplicial-site-site}
把 $\mathcal{C}_{total}$ 上的层 $\mathcal{F}$ 看成系统
$(\mathcal{F}_n, \mathcal{F}(\varphi))$，则
$h_{total, *}\mathcal{F}$ 对应于系统
$(h_{n, *}\mathcal{F}_n, h_{n, *}\mathcal{F}(\varphi))$。
显然，函子
$(\mathcal{F}'_n, \mathcal{F}'(\varphi)) \to
(h_n^{-1}\mathcal{F}'_n, h_n^{-1}\mathcal{F}'(\varphi))$
是它的左伴随。由伴随的唯一性，$h_{total}^{-1}$ 在系统上由此规则给出。
特别地，由该引理对 $\mathcal{C}_{total}$ 上层的描述以及各函子
$h_n^{-1}$ 的正合性，$h_{total}^{-1}$ 正合，从而得到所需拓扑斯态射。
最后还有 $g_n^{-1} \circ h_{total}^{-1} =
h_n^{-1} \circ (g'_n)^{-1}$，这证明引理所示图交换。

\medskip\noindent
情形 B。在次数 $n$ 使用 $v_n$，得到函子
$v_{total} : \mathcal{C}_{total} \to \mathcal{C}'_{total}$。
它显然余连续（参见引理 \ref{spaces-simplicial-lemma-simplicial-cocontinuous-site}
中覆盖的定义）。令 $h_{total}$ 为 $v_{total}$ 所伴随的拓扑斯态射。
引理所示图的交换性直接来自《站点》引理
\ref{sites-lemma-composition-cocontinuous}。
取左伴随后，引理最后的等式化为
$$
h_{total}^{-1} \circ (g'_n)^{Sh}_! = g^{Sh}_{n!} \circ h_n^{-1}
$$
这直接由引理 \ref{spaces-simplicial-lemma-restriction-to-components-site} 中对函子
$(g'_n)^{Sh}_!$ 与 $g^{Sh}_{n!}$ 的显式描述、对
$\varphi : [m] \to [n]$ 成立的等式
$h_n^{-1} \circ (f'_\varphi)^{-1} =
f_\varphi^{-1} \circ h_m^{-1}$，以及已知 $h_{total}^{-1}$
与限制到单纯站点次数 $n$ 部分交换这一事实得到。
\end{proof}

\begin{lemma}
\label{spaces-simplicial-lemma-direct-image-morphism-simplicial-sites}
采用引理 \ref{spaces-simplicial-lemma-morphism-simplicial-sites} 的记号与假设。
对 $K \in D(\mathcal{C}_{total})$，有
$(g'_n)^{-1}Rh_{total, *}K = Rh_{n, *}g_n^{-1}K$。
\end{lemma}

\begin{proof}
设 $\mathcal{I}^\bullet$ 是 $\mathcal{C}_{total}$ 上表示 $K$ 的
K-内射复形。则 $g_n^{-1}K$ 由
$g_n^{-1}\mathcal{I}^\bullet = \mathcal{I}_n^\bullet$ 表示；
由引理 \ref{spaces-simplicial-lemma-restriction-injective-to-component-site}，后者是 K-内射的。
% P10-E0177
我们有 $(g'_n)^{-1}h_{total, *}\mathcal{I}^\bullet =
h_{n, *}g_n^{-1}\mathcal{I}_n^\bullet$；
由引理 \ref{spaces-simplicial-lemma-morphism-simplicial-sites} 即得所需等式。
\end{proof}

\begin{remark}
\label{spaces-simplicial-remark-morphism-augmentation-simplicial-sites}
设 $\mathcal{C}_n, f_\varphi, u_\varphi$ 和
$\mathcal{C}'_n, f'_\varphi, u'_\varphi$ 如情形
\ref{spaces-simplicial-situation-simplicial-site}。设 $a_0$ 和 $a'_0$ 分别为评注
\ref{spaces-simplicial-remark-augmentation-site} 所述指向站点 $\mathcal{D}$ 和
$\mathcal{D}'$ 的增广。设 $h$ 为评注
\ref{spaces-simplicial-remark-morphism-simplicial-sites} 所述的单纯站点之间的态射。
若拓扑斯态射 $h_{-1} : \Sh(\mathcal{D}) \to \Sh(\mathcal{D}')$
满足下列条件，则称它{\it 与 $h$、$a_0$、$a'_0$ 相容}：
\begin{enumerate}
\item[(A)] $h_{-1}$ 来自站点态射
$h_{-1} : \mathcal{D} \to \mathcal{D}'$，且作为站点态射有
$a'_0 \circ h_0 = h_{-1} \circ a_0$；
\item[(B)] $h_{-1}$ 来自余连续函子
$v_{-1} : \mathcal{D} \to \mathcal{D}'$，且作为函子有
$u'_0 \circ v_0 = v_{-1} \circ u_0$。
\end{enumerate}
在两种情形中，作为拓扑斯态射均有
$a'_0 \circ h_0 = h_{-1} \circ a_0$；情形 B 参见《站点》引理
\ref{sites-lemma-composition-cocontinuous}，情形 A 参见《站点》定义
\ref{sites-definition-composition-morphisms-sites}。
\end{remark}

\begin{lemma}
\label{spaces-simplicial-lemma-morphism-augmentation-simplicial-sites}
设 $\mathcal{C}_n, f_\varphi, u_\varphi, \mathcal{D}, a_0$，
$\mathcal{C}'_n, f'_\varphi, u'_\varphi, \mathcal{D}', a'_0$，以及
$h_n$（$n \geq -1$）如评注
\ref{spaces-simplicial-remark-morphism-augmentation-simplicial-sites}。
则得到交换图
$$
\xymatrix{
\Sh(\mathcal{C}_{total}) \ar[d]_a \ar[r]_{h_{total}} &
\Sh(\mathcal{C}'_{total}) \ar[d]^{a'} \\
\Sh(\mathcal{D}) \ar[r]^{h_{-1}} &
\Sh(\mathcal{D}')
}
$$
\end{lemma}

\begin{proof}
态射 $h$ 由引理 \ref{spaces-simplicial-lemma-morphism-simplicial-sites} 定义，
态射 $a$ 与 $a'$ 由引理 \ref{spaces-simplicial-lemma-augmentation-site} 定义。
故只须证明该图交换。为此，证明
$a^{-1} \circ h_{-1}^{-1} = h_{total}^{-1} \circ (a')^{-1}$。
由引理 \ref{spaces-simplicial-lemma-morphism-simplicial-sites} 中的交换图，
以及引理 \ref{spaces-simplicial-lemma-describe-sheaves-simplicial-site-site}
借助分量对 $\Sh(\mathcal{C}_{total})$ 和
$\Sh(\mathcal{C}'_{total})$ 的描述，只须证明
$$
\xymatrix{
\Sh(\mathcal{C}_n) \ar[d]_{a_n} \ar[r]_{h_n} &
\Sh(\mathcal{C}'_n) \ar[d]^{a'_n} \\
\Sh(\mathcal{D}) \ar[r]^{h_{-1}} &
\Sh(\mathcal{D}')
}
$$
对所有 $n$ 均交换。$n=0$ 的情形是评注
\ref{spaces-simplicial-remark-morphism-augmentation-simplicial-sites} 中的假设。
当 $n>0$ 时，选取 $\varphi : [0] \to [n]$；所需交换性由
$n=0$ 的情形、关系 $a_n = a_0 \circ f_\varphi$ 与
$a'_n = a'_0 \circ f'_\varphi$，以及交换关系
$f'_\varphi \circ h_n = h_0 \circ f_\varphi$ 得到。
\end{proof}




\section{赋环单纯站点}
\label{spaces-simplicial-section-simplicial-sites-modules}

\noindent
现在为单纯拓扑斯赋予一个环层。

\begin{lemma}
\label{spaces-simplicial-lemma-restriction-module-to-components-site}
在情形 \ref{spaces-simplicial-situation-simplicial-site} 中，设 $\mathcal{O}$ 为
$\mathcal{C}_{total}$ 上的环层。存在典范赋环拓扑斯态射
$g_n : (\Sh(\mathcal{C}_n), \mathcal{O}_n) \to
(\Sh(\mathcal{C}_{total}), \mathcal{O})$
其底层拓扑斯态射与引理 \ref{spaces-simplicial-lemma-restriction-to-components-site}
中的 $g_n$ 一致。函子
$g_n^* : \textit{Mod}(\mathcal{O}) \to \textit{Mod}(\mathcal{O}_n)$
有左伴随 $g_{n!}$。对 $\textit{Mod}(\mathcal{O}_n)$ 中的
$\mathcal{G}$，$g_{n!}\mathcal{G}$ 在 $\mathcal{C}_m$ 上的限制为
$$
\bigoplus\nolimits_{\varphi : [n] \to [m]} f_\varphi^*\mathcal{G}
$$
其中 $f_\varphi : (\Sh(\mathcal{C}_m), \mathcal{O}_m) \to
(\Sh(\mathcal{C}_n), \mathcal{O}_n)$ 是赋环拓扑斯态射；
其底层拓扑斯态射与先前定义的 $f_\varphi$ 一致，并在环层上使用映射
$\mathcal{O}(\varphi) : f_\varphi^{-1}\mathcal{O}_n \to \mathcal{O}_m$
。
\end{lemma}

\begin{proof}
由引理 \ref{spaces-simplicial-lemma-restriction-to-components-site}，有
$g_n^{-1}\mathcal{O} = \mathcal{O}_n$，从而得到上述赋环拓扑斯态射。
由《站点上的模》引理 \ref{sites-modules-lemma-lower-shriek-modules}，
得到伴随 $g_{n!}$。为证明 $g_{n!}$ 的公式，先定义
$\mathcal{C}_{total}$ 上的 $\mathcal{O}$-模层 $\mathcal{H}$，
其次数 $m$ 分量为 $\mathcal{O}_m$-模
$$
\mathcal{H}_m =
\bigoplus\nolimits_{\varphi : [n] \to [m]} f_\varphi^*\mathcal{G}
$$
给定映射 $\psi : [m] \to [m']$，映射
$\mathcal{H}(\psi) : f_\psi^{-1}\mathcal{H}_m \to \mathcal{H}_{m'}$
在各分量上由下式给出：
$$
f_\psi^{-1} f_\varphi^*\mathcal{G} \to
f_\psi^* f_\varphi^*\mathcal{G} \to
f_{\psi \circ \varphi}^*\mathcal{G}
$$
由于映射 $f_\psi^{-1}\mathcal{H}_m \to \mathcal{H}_{m'}$
关于 $\mathcal{O}(\psi) : f_\psi^{-1}\mathcal{O}_m \to \mathcal{O}_{m'}$
是半线性的，这确实定义一个 $\mathcal{O}$-模
（使用引理 \ref{spaces-simplicial-lemma-describe-sheaves-simplicial-site-site}）。
随后直接证明
$$
\Mor_{\mathcal{O}_n}(\mathcal{G}, \mathcal{F}_n) =
\Mor_{\mathcal{O}}(\mathcal{H}, \mathcal{F})
$$
其论证同引理 \ref{spaces-simplicial-lemma-restriction-to-components-site} 的证明。
故 $\mathcal{H} = g_{n!}\mathcal{G}$，如所求。
\end{proof}

\begin{lemma}
\label{spaces-simplicial-lemma-restriction-injective-to-component-limp}
在情形 \ref{spaces-simplicial-situation-simplicial-site} 中，设 $\mathcal{O}$ 为
$\mathcal{C}_{total}$ 上的环层。若 $\mathcal{I}$ 在
$\textit{Mod}(\mathcal{O})$ 中内射，则 $\mathcal{I}_n$ 是
$\mathcal{C}_n$ 上的完全无环层。
\end{lemma}

\begin{proof}
把《站点上的上同调》引理
\ref{sites-cohomology-lemma-pullback-injective-limp} 应用于包含函子
$\mathcal{C}_n \to \mathcal{C}_{total}$，并使用引理
\ref{spaces-simplicial-lemma-restriction-to-components-site} 所证明的该函子性质，即得结论。
\end{proof}

\begin{lemma}
\label{spaces-simplicial-lemma-exactness-g-shriek-modules}
在引理 \ref{spaces-simplicial-lemma-restriction-module-to-components-site} 的假设下，
若对所有 $\varphi : [n] \to [m]$，映射
$f_\varphi^{-1}\mathcal{O}_n \to \mathcal{O}_m$ 均平坦，则函子
$g_{n!} : \textit{Mod}(\mathcal{O}_n) \to \textit{Mod}(\mathcal{O})$
正合。
\end{lemma}

\begin{proof}
回忆 $g_{n!}\mathcal{G}$ 是如下 $\mathcal{O}$-模：其次数 $m$ 部分是
$\mathcal{O}_m$-模
$$
\bigoplus\nolimits_{\varphi : [n] \to [m]} f_\varphi^*\mathcal{G}
$$
这里赋环拓扑斯态射
$f_\varphi : (\Sh(\mathcal{C}_m), \mathcal{O}_m) \to
(\Sh(\mathcal{C}_n), \mathcal{O}_n)$ 使用引理陈述中的映射
$f_\varphi^{-1}\mathcal{O}_n \to \mathcal{O}_m$。
若这些映射平坦，则 $f_\varphi^*$ 正合
（《站点上的模》引理 \ref{sites-modules-lemma-flat-pullback-exact}）。
由站点 $\mathcal{C}_{total}$ 的定义可知，这些函子具有所需正合性，故得结论。
\end{proof}

\begin{lemma}
\label{spaces-simplicial-lemma-restriction-injective-to-component-site-module}
在情形 \ref{spaces-simplicial-situation-simplicial-site} 中，设 $\mathcal{O}$ 为
$\mathcal{C}_{total}$ 上的环层，并假设对所有
$\varphi : [n] \to [m]$，映射
$f_\varphi^{-1}\mathcal{O}_n \to \mathcal{O}_m$ 均平坦。
若 $\mathcal{I}$ 在 $\textit{Mod}(\mathcal{O})$ 中内射，
则 $\mathcal{I}_n$ 在 $\textit{Mod}(\mathcal{O}_n)$ 中内射。
\end{lemma}

\begin{proof}
这由《同调代数》引理 \ref{homology-lemma-adjoint-preserve-injectives}
和引理 \ref{spaces-simplicial-lemma-exactness-g-shriek-modules} 得到。
\end{proof}







\section{赋环单纯站点的态射}
\label{spaces-simplicial-section-morphism-simplicial-sites-modules}

\noindent
我们继续第 \ref{spaces-simplicial-section-morphism-simplicial-sites} 节的讨论。

\begin{remark}
\label{spaces-simplicial-remark-morphism-simplicial-sites-modules}
设 $\mathcal{C}_n, f_\varphi, u_\varphi$ 与
$\mathcal{C}'_n, f'_\varphi, u'_\varphi$ 如
情形 \ref{spaces-simplicial-situation-simplicial-site} 所述。
设 $\mathcal{O}$ 与 $\mathcal{O}'$
分别为 $\mathcal{C}_{total}$ 与 $\mathcal{C}'_{total}$ 上的环层。
若 $h$ 是注 \ref{spaces-simplicial-remark-morphism-simplicial-sites}
所述的单纯站点间的态射，且
$h^\sharp : h_{total}^{-1}\mathcal{O}' \to \mathcal{O}$，
等价地，$h^\sharp : \mathcal{O}' \to h_{total, *}\mathcal{O}$
是环层的同态，则称 $(h, h^\sharp)$ 是
{\it 赋环单纯站点间的态射}。
\end{remark}

\begin{lemma}
\label{spaces-simplicial-lemma-morphism-simplicial-sites-modules}
设 $\mathcal{C}_n, f_\varphi, u_\varphi$ 与
$\mathcal{C}'_n, f'_\varphi, u'_\varphi$ 如
情形 \ref{spaces-simplicial-situation-simplicial-site} 所述。
设 $\mathcal{O}$ 与 $\mathcal{O}'$
分别为 $\mathcal{C}_{total}$ 与 $\mathcal{C}'_{total}$ 上的环层。
设 $(h, h^\sharp)$ 是注
\ref{spaces-simplicial-remark-morphism-simplicial-sites-modules} 所述的单纯站点间的态射。
则我们得到赋环拓扑斯的态射
$$
h_{total} :
(\Sh(\mathcal{C}_{total}), \mathcal{O})
\to
(\Sh(\mathcal{C}'_{total}), \mathcal{O}')
$$
以及交换图
$$
\xymatrix{
(\Sh(\mathcal{C}_n), \mathcal{O}_n) \ar[d]_{g_n} \ar[r]_{h_n} &
(\Sh(\mathcal{C}'_n), \mathcal{O}'_n) \ar[d]^{g'_n} \\
(\Sh(\mathcal{C}_{total}), \mathcal{O}) \ar[r]^{h_{total}} &
(\Sh(\mathcal{C}'_{total}), \mathcal{O}')
}
$$
；这里它们是赋环拓扑斯的图，而 $g_n$ 与 $g'_n$ 如
引理 \ref{spaces-simplicial-lemma-restriction-module-to-components-site} 所述。
此外，作为函子 $\textit{Mod}(\mathcal{O}) \to
\textit{Mod}(\mathcal{O}'_n)$，我们有
$(g'_n)^* \circ h_{total, *} = h_{n, *} \circ g_n^*$
。
\end{lemma}

\begin{proof}
由引理 \ref{spaces-simplicial-lemma-morphism-simplicial-sites} 与
\ref{spaces-simplicial-lemma-restriction-module-to-components-site}，并追踪环层即可得出。
有一个小细节：为了将 $h_n$ 定义为赋环拓扑斯的态射，我们令
$h_n^\sharp = g_n^{-1}h^\sharp :
g_n^{-1}h_{total}^{-1}\mathcal{O}' \to g_n^{-1}\mathcal{O}$
；这是有意义的，因为
$g_n^{-1}h_{total}^{-1}\mathcal{O}' = h_n^{-1}(g'_n)^{-1}\mathcal{O}' =
h_n^{-1}\mathcal{O}'_n$ 且 $g_n^{-1}\mathcal{O} = \mathcal{O}_n$。
注意，对任意 $\mathcal{O}$-模层 $\mathcal{F}$，均有
$g_n^*\mathcal{F} = g_n^{-1}\mathcal{F}$；
对 $g'_n$ 亦然。这也有助于说明，为什么
$(g'_n)^* \circ h_{total, *} = h_{n, *} \circ g_n^*$
可由引理 \ref{spaces-simplicial-lemma-morphism-simplicial-sites} 中的相应陈述推出。
\end{proof}

\begin{lemma}
\label{spaces-simplicial-lemma-direct-image-morphism-simplicial-sites-modules}
采用引理 \ref{spaces-simplicial-lemma-morphism-simplicial-sites-modules} 中的记号与假设。
对 $K \in D(\mathcal{O})$，有
$(g'_n)^*Rh_{total, *}K = Rh_{n, *}g_n^*K$。
\end{lemma}

\begin{proof}
回忆一下，由引理 \ref{spaces-simplicial-lemma-restriction-module-to-components-site}
中的构造，$g_n^{-1}\mathcal{O} = \mathcal{O}_n$，故
$g_n^* = g_n^{-1}$。特别地，$g_n^*$ 是正合的，而 $Lg_n^*$
可通过将 $g_n^*$ 作用于任意一个模复形代表来给出。对 $g'_n$ 亦然。
存在典范基变换映射
$(g'_n)^*Rh_{total, *}K \to Rh_{n, *}g_n^*K$，见
《站点上的上同调》，注 \ref{sites-cohomology-remark-base-change}。
由《站点上的上同调》，引理
\ref{sites-cohomology-lemma-modules-abelian-unbounded}
可知，它在 $D(\mathcal{C}'_n)$ 中的像为映射
$(g'_n)^{-1}Rh_{total, *}K_{ab} \to Rh_{n, *}g_n^{-1}K_{ab}$
，其中 $K_{ab}$ 是 $K$ 在 $D(\mathcal{C}_{total})$ 中的像。
引理 \ref{spaces-simplicial-lemma-direct-image-morphism-simplicial-sites}
已证明此映射是同构，故结论成立。
\end{proof}








\section{单纯站点上的上同调}
\label{spaces-simplicial-section-cohomology-simplicial-sites}

\noindent
设 $\mathcal{C}$ 如情形 \ref{spaces-simplicial-situation-simplicial-site} 所述。
在以下各引理的陈述中，我们以
$g_n : \Sh(\mathcal{C}_n) \to \Sh(\mathcal{C}_{total})$ 表示
引理 \ref{spaces-simplicial-lemma-restriction-to-components-site} 中的拓扑斯态射。
若 $\varphi : [m] \to [n]$ 是 $\Delta$ 的一个态射，则拓扑斯图
$$
\xymatrix{
\Sh(\mathcal{C}_n) \ar[rd]_{g_n} \ar[rr]_{f_\varphi} & &
\Sh(\mathcal{C}_m) \ar[ld]^{g_m} \\
& \Sh(\mathcal{C}_{total})
}
$$
并不交换，但存在来自映射
$\mathcal{F}(\varphi) : f_\varphi^{-1}\mathcal{F}_m \to \mathcal{F}_n$
的 $2$-态射 $g_n \to g_m \circ f_\varphi$。
参见《站点》第 \ref{sites-section-2-category} 节。

\begin{lemma}
\label{spaces-simplicial-lemma-simplicial-resolution-Z-site}
在情形 \ref{spaces-simplicial-situation-simplicial-site} 中，并采用上述记号，存在复形
$$
\ldots \to g_{2!}\mathbf{Z} \to g_{1!}\mathbf{Z} \to g_{0!}\mathbf{Z}
$$
；它由 $\mathcal{C}_{total}$ 上的阿贝尔层组成，并构成
$\mathcal{C}_{total}$ 上取值为 $\mathbf{Z}$ 的常值层的一个消解。
\end{lemma}

\begin{proof}
以下不再说明地使用引理 \ref{spaces-simplicial-lemma-restriction-to-components-site}
中对函子 $g_{n!}$ 的描述。
取 $\sum (-1)^i d^n_i$ 为该复形的映射，其中
$d^n_i : g_{n!}\mathbf{Z} \to g_{n - 1!}\mathbf{Z}$ 是映射 $\mathbf{Z} \to
\bigoplus_{[n - 1] \to [n]} \mathbf{Z} = g_n^{-1}g_{n - 1!}\mathbf{Z}$
的伴随；这里该映射对应于标记为
$\delta^n_i : [n - 1] \to [n]$ 的因子。
于是将 $g_m^{-1}$ 作用于该复形，得到 $\mathcal{C}_m$ 上的复形
$$
\ldots \to
\bigoplus\nolimits_{\alpha \in \Mor_\Delta([2], [m])]} \mathbf{Z} \to
\bigoplus\nolimits_{\alpha \in \Mor_\Delta([1], [m])]} \mathbf{Z} \to
\bigoplus\nolimits_{\alpha \in \Mor_\Delta([0], [m])]} \mathbf{Z}
$$
。
换言之，这是与单纯集合 $\Delta[m]$ 上的自由阿贝尔层相关联的复形；
参见《单纯对象》例 \ref{simplicial-example-simplex-simplicial-set}。
由于 $\Delta[m]$ 与 $\Delta[0]$ 同伦等价（参见《单纯对象》例
\ref{simplicial-example-simplex-contractible}），且“取其上的自由阿贝尔层”
是一个函子，故上述复形与 $\Delta[0]$ 上的自由阿贝尔层同伦等价
（《单纯对象》注 \ref{simplicial-remark-homotopy-better} 及
引理 \ref{simplicial-lemma-homotopy-equivalence-s-N}）。
此复形在正次数上无环，并在次数 $0$ 上等于 $\mathbf{Z}$。
\end{proof}

\begin{lemma}
\label{spaces-simplicial-lemma-cech-complex}
在情形 \ref{spaces-simplicial-situation-simplicial-site} 中，设 $\mathcal{F}$ 是
$\mathcal{C}_{total}$ 上的阿贝尔层。存在典范复形
$$
0 \to \Gamma(\mathcal{C}_{total}, \mathcal{F}) \to
\Gamma(\mathcal{C}_0, \mathcal{F}_0) \to
\Gamma(\mathcal{C}_1, \mathcal{F}_1) \to
\Gamma(\mathcal{C}_2, \mathcal{F}_2) \to \ldots
$$
；它在次数 $-1, 0$ 上正合；若 $\mathcal{F}$ 是内射的，
则它处处正合。
\end{lemma}

\begin{proof}
注意
$\Hom(\mathbf{Z}, \mathcal{F}) = \Gamma(\mathcal{C}_{total}, \mathcal{F})$
且
$\Hom(g_{n!}\mathbf{Z}, \mathcal{F}) = \Gamma(\mathcal{C}_n, \mathcal{F}_n)$.
因此，由引理 \ref{spaces-simplicial-lemma-simplicial-resolution-Z-site}
以及当 $\mathcal{F}$ 内射时 $\Hom(-, \mathcal{F})$ 正合这一事实，
立即得到本引理。
\end{proof}

\begin{lemma}
\label{spaces-simplicial-lemma-simplicial-sheaf-cohomology-site}
在情形 \ref{spaces-simplicial-situation-simplicial-site} 中，对属于
$D^+(\mathcal{C}_{total})$ 的 $K$，存在谱序列
$(E_r, d_r)_{r \geq 0}$，其中
$$
E_1^{p, q} = H^q(\mathcal{C}_p, K_p),\quad
d_1^{p, q} : E_1^{p, q} \to E_1^{p + 1, q}
$$
，它收敛到 $H^{p + q}(\mathcal{C}_{total}, K)$。
此谱序列关于 $K$ 是函子性的。
\end{lemma}

\begin{proof}
设 $\mathcal{I}^\bullet$ 是表示 $K$ 的下有界内射复形。
考虑各项为
$$
A^{p, q} = \Gamma(\mathcal{C}_p, \mathcal{I}^q_p)
$$
的双复形；其水平箭头来自引理 \ref{spaces-simplicial-lemma-cech-complex}，
竖直箭头来自复形 $\mathcal{I}^\bullet$ 的微分。
该双复形的各行在正次数上正合，并在次数 $0$ 上给出
$\Gamma(\mathcal{C}_{total}, \mathcal{I}^q)$。
另一方面，由于到 $\mathcal{C}_p$ 的限制是正合的
（引理 \ref{spaces-simplicial-lemma-restriction-to-components-site}），复形
$\mathcal{I}_p^\bullet$ 在 $D(\mathcal{C}_p)$ 中表示 $K_p$。
诸层 $\mathcal{I}_p^q$ 是 $\mathcal{C}_p$ 上的内射阿贝尔层
（引理 \ref{spaces-simplicial-lemma-restriction-injective-to-component-site}）。
因此，各列的上同调计算群 $H^q(\mathcal{C}_p, K_p)$。
应用《同调代数》引理 \ref{homology-lemma-first-quadrant-ss} 与
\ref{homology-lemma-double-complex-gives-resolution} 即得结论。
\end{proof}

\begin{remark}
\label{spaces-simplicial-remark-simplicial-cohomology-site}
采用引理 \ref{spaces-simplicial-lemma-simplicial-sheaf-cohomology-site}
中的假设与记号，但不要求 $D(\mathcal{C}_{total})$ 中的 $K$ 下有界。
我们断言，此时也存在自然谱序列。具体而言，设
$\mathcal{I}^\bullet$ 是 $\mathcal{C}_{total}$ 上表示 $K$
且各项内射的层 K-内射复形。则
\begin{align*}
R\Gamma(\mathcal{C}_{total}, K)
& =
R\Hom(\mathbf{Z}, K) \\
& =
R\Hom(
\ldots \to g_{2!}\mathbf{Z} \to g_{1!}\mathbf{Z} \to g_{0!}\mathbf{Z}, K) \\
& =
\Gamma(\mathcal{C}_{total},
\SheafHom^\bullet(
\ldots \to g_{2!}\mathbf{Z} \to g_{1!}\mathbf{Z} \to g_{0!}\mathbf{Z},
\mathcal{I}^\bullet)) \\
& =
\text{Tot}_\pi(A^{\bullet, \bullet})
\end{align*}
其中 $A^{\bullet, \bullet}$ 是各项为
$A^{p, q} = \Gamma(\mathcal{C}_p, \mathcal{I}^q_p)$ 的双复形，
$\text{Tot}_\pi$ 表示该双复形的乘积全化。
这里，第一个等式在任意站点上成立。第二个等式由引理
\ref{spaces-simplicial-lemma-simplicial-resolution-Z-site} 成立。
第三个等式成立是因为 $\mathcal{I}^\bullet$ 是 K-内射的；参见
《站点上的上同调》第 \ref{sites-cohomology-section-hom-complexes} 节
与第 \ref{sites-cohomology-section-internal-hom} 节。
最后一个等式由 $\SheafHom^\bullet$ 的构造以及等式
$\Hom(g_{p!}\mathbf{Z}, \mathcal{I}^q) =
\Gamma(\mathcal{C}_p, \mathcal{I}^q_p)$ 成立。
现在将 $\text{Tot}_\pi(A^{\bullet, \bullet})$ 看成滤过复形，其中
$F^i\text{Tot}^n_\pi(A^{\bullet, \bullet}) =
\prod_{p + q = n,\ p \geq i} A^{p, q}$，由此得到所需谱序列。
除正则性、有界性、收敛性与极限项等问题外，我们得到的谱序列
与《同调代数》第 \ref{homology-section-double-complex} 节中的谱序列
$({}'E_r, {}'d_r)_{r \geq 0}$ 具有同样的表现
（该节讨论的是直和全化的情形）。特别地，正如引理
\ref{spaces-simplicial-lemma-simplicial-sheaf-cohomology-site} 中一样，我们得到
$E_1^{p, q} = H^q(\mathcal{C}_p, K_p)$。
\end{remark}

\begin{lemma}
\label{spaces-simplicial-lemma-simplicial-sheaf-cohomology-site-zero}
在情形 \ref{spaces-simplicial-situation-simplicial-site} 中，设 $K$ 是
$D(\mathcal{C}_{total})$ 的一个对象。
\begin{enumerate}
\item 若对所有 $p \geq 0$ 都有 $H^{-p}(\mathcal{C}_p, K_p) = 0$，则
$H^0(\mathcal{C}_{total}, K) = 0$。
\item 若对所有 $p \geq 0$ 都有 $R\Gamma(\mathcal{C}_p, K_p) = 0$，
则 $R\Gamma(\mathcal{C}_{total}, K) = 0$。
\end{enumerate}
\end{lemma}

\begin{proof}
采用注 \ref{spaces-simplicial-remark-simplicial-cohomology-site} 中的记号，
可见 $R\Gamma(\mathcal{C}_{total}, K)$ 由
$\text{Tot}_\pi(A^{\bullet, \bullet})$ 表示。
(1) 中的假设说明 $H^{-p}(A^{p, \bullet}) = 0$。
于是 (1) 中的消失由《代数续论》引理
\ref{more-algebra-lemma-vanishing-coh-prod-totalization} 得出。
(2) 由 (1) 及平移得到。
\end{proof}

\begin{lemma}
\label{spaces-simplicial-lemma-sanity-check}
设 $\mathcal{C}$ 如情形 \ref{spaces-simplicial-situation-simplicial-site} 所述，
$U \in \Ob(\mathcal{C}_n)$，且
$\mathcal{F} \in \textit{Ab}(\mathcal{C}_{total})$。
则 $H^p(U, \mathcal{F}) = H^p(U, g_n^{-1}\mathcal{F})$，
其中在左端，$U$ 被视为 $\mathcal{C}_{total}$ 的一个对象。
\end{lemma}

\begin{proof}
注意，“将 $U$ 视为 $\mathcal{C}_{total}$ 的对象”这一说法，
在情形 (A) 下由引理 \ref{spaces-simplicial-lemma-simplicial-site-site}
中的 $\mathcal{C}_{total}$ 构造得到解释，在情形 (B) 下则由引理
\ref{spaces-simplicial-lemma-simplicial-cocontinuous-site} 中的构造得到解释。
于是，该等式由引理 \ref{spaces-simplicial-lemma-restriction-injective-to-component-site}
及上同调的定义得出。
\end{proof}






\section{单纯站点的上同调与增广}
\label{spaces-simplicial-section-cohomology-augmentation-simplicial-sites}

\noindent
考虑情形 \ref{spaces-simplicial-situation-simplicial-site} 所述的单纯站点 $\mathcal{C}$。
设 $a_0$ 是注 \ref{spaces-simplicial-remark-augmentation-site}
所述的朝向站点 $\mathcal{D}$ 的增广。
由引理 \ref{spaces-simplicial-lemma-augmentation-site}，我们得到拓扑斯的态射
$$
a : \Sh(\mathcal{C}_{total}) \longrightarrow \Sh(\mathcal{D})
$$
以及拓扑斯的态射
$g_n : \Sh(\mathcal{C}_n) \to \Sh(\mathcal{C}_{total})$
，后者如引理 \ref{spaces-simplicial-lemma-restriction-to-components-site} 所述。
将复合 $a \circ g_n$ 记为
$a_n : \Sh(\mathcal{C}_n) \to \Sh(\mathcal{D})$。
此外，单纯结构给出拓扑斯的态射
$f_\varphi : \Sh(\mathcal{C}_n) \to \Sh(\mathcal{C}_m)$
，使得对所有 $\varphi : [m] \to [n]$，均有
% P10-E0180
$a_n \circ f_\varphi = a_m$。

\begin{lemma}
\label{spaces-simplicial-lemma-simplicial-resolution-augmentation}
在情形 \ref{spaces-simplicial-situation-simplicial-site} 中，设
$a_0$ 是注 \ref{spaces-simplicial-remark-augmentation-site}
所述的朝向站点 $\mathcal{D}$ 的增广。
对 $\mathcal{D}$ 上任意阿贝尔层 $\mathcal{G}$，存在正合复形
$$
\ldots \to
g_{2!}(a_2^{-1}\mathcal{G}) \to
g_{1!}(a_1^{-1}\mathcal{G}) \to
g_{0!}(a_0^{-1}\mathcal{G}) \to
a^{-1}\mathcal{G} \to 0
$$
；它由 $\mathcal{C}_{total}$ 上的阿贝尔层组成。
\end{lemma}

\begin{proof}
建议读者先阅读引理 \ref{spaces-simplicial-lemma-simplicial-resolution-Z-site} 的证明。
以下不再说明地使用引理 \ref{spaces-simplicial-lemma-augmentation-site}，
以及引理 \ref{spaces-simplicial-lemma-restriction-to-components-site}
中对函子 $g_{n!}$ 的描述。
特别地，$g_{n!}(a_n^{-1}\mathcal{G})$ 是
$\mathcal{C}_{total}$ 上的层，其到 $\mathcal{C}_m$ 的限制为层
$$
\bigoplus\nolimits_{\varphi : [n] \to [m]} f_\varphi^{-1}a_n^{-1}\mathcal{G} =
\bigoplus\nolimits_{\varphi : [n] \to [m]} a_m^{-1}\mathcal{G}
$$
。取 $\sum (-1)^i d^n_i$ 为该复形的映射，其中
$d^n_i : g_{n!}(a_n^{-1}\mathcal{G}) \to g_{n - 1!}(a_{n - 1}^{-1}\mathcal{G})$
是映射
$a_n^{-1}\mathcal{G} \to \bigoplus_{[n - 1] \to [n]} a_n^{-1}\mathcal{G} =
g_n^{-1}g_{n - 1!}(a_{n - 1}^{-1}\mathcal{G})$
的伴随；这里该映射对应于标记为
$\delta^n_i : [n - 1] \to [n]$ 的因子。
映射 $g_{0!}(a_0^{-1}\mathcal{G}) \to a^{-1}\mathcal{G}$
与 $a_0^{-1}\mathcal{G}$ 的恒等映射伴随。
于是，将 $g_m^{-1}$ 作用于次数为 $\ldots, 2, 1, 0$ 的链复形，
得到 $\mathcal{C}_m$ 上的复形
$$
\ldots \to
\bigoplus\nolimits_{\alpha \in \Mor_\Delta([2], [m])]} a_m^{-1}\mathcal{G} \to
\bigoplus\nolimits_{\alpha \in \Mor_\Delta([1], [m])]} a_m^{-1}\mathcal{G} \to
\bigoplus\nolimits_{\alpha \in \Mor_\Delta([0], [m])]} a_m^{-1}\mathcal{G}
$$
。这等于 $a_m^{-1}\mathcal{G}$ 在常值层 $\mathbf{Z}$ 上与复形
$$
\ldots \to
\bigoplus\nolimits_{\alpha \in \Mor_\Delta([2], [m])]} \mathbf{Z} \to
\bigoplus\nolimits_{\alpha \in \Mor_\Delta([1], [m])]} \mathbf{Z} \to
\bigoplus\nolimits_{\alpha \in \Mor_\Delta([0], [m])]} \mathbf{Z}
$$
的张量积；该复形已在引理 \ref{spaces-simplicial-lemma-simplicial-resolution-Z-site}
的证明中讨论过。那里已经看到，此复形同伦等价于置于次数 $0$
的 $\mathbf{Z}$，证明完毕。
\end{proof}

\begin{lemma}
\label{spaces-simplicial-lemma-augmentation-cech-complex}
在情形 \ref{spaces-simplicial-situation-simplicial-site} 中，设
$a_0$ 是注 \ref{spaces-simplicial-remark-augmentation-site}
所述的朝向站点 $\mathcal{D}$ 的增广。
对 $\mathcal{C}_{total}$ 上的阿贝尔层 $\mathcal{F}$，
存在 $\mathcal{D}$ 上的典范复形
$$
0 \to a_*\mathcal{F} \to a_{0, *}\mathcal{F}_0 \to a_{1, *}\mathcal{F}_1 \to
a_{2, *}\mathcal{F}_2 \to \ldots
$$
；它在次数 $-1, 0$ 上正合；若 $\mathcal{F}$ 是内射的，
则它处处正合。
\end{lemma}

\begin{proof}
为构造该复形，由 Yoneda 引理，只需对 $\mathcal{D}$ 上任意
阿贝尔层 $\mathcal{G}$ 构造复形
$$
0 \to \Hom(\mathcal{G}, a_*\mathcal{F}) \to
\Hom(\mathcal{G}, a_{0, *}\mathcal{F}_0) \to
\Hom(\mathcal{G}, a_{1, *}\mathcal{F}_1) \to \ldots
$$
，并使其关于 $\mathcal{G}$ 是函子性的。为此，将
$\Hom(-, \mathcal{F})$ 作用于引理
\ref{spaces-simplicial-lemma-simplicial-resolution-augmentation} 的正合复形，
并使用拉回与推出的伴随性。
由于 $\Hom(-, \mathcal{F})$ 左正合，次数 $-1, 0$
上的正合性由该构造得出。若 $\mathcal{F}$ 是内射阿贝尔层，
则 $\Hom(-, \mathcal{F})$ 正合，故此复形正合。
\end{proof}

\begin{lemma}
\label{spaces-simplicial-lemma-augmentation-spectral-sequence}
在情形 \ref{spaces-simplicial-situation-simplicial-site} 中，设
$a_0$ 是注 \ref{spaces-simplicial-remark-augmentation-site}
所述的朝向站点 $\mathcal{D}$ 的增广。
对任意属于 $D^+(\mathcal{C}_{total})$ 的 $K$，存在谱序列
$(E_r, d_r)_{r \geq 0}$，其中
$$
E_1^{p, q} = R^qa_{p, *} K_p,\quad
d_1^{p, q} : E_1^{p, q} \to E_1^{p + 1, q}
$$
；它收敛到 $R^{p + q}a_*K$。此谱序列关于 $K$ 是函子性的。
\end{lemma}

\begin{proof}
设 $\mathcal{I}^\bullet$ 是表示 $K$ 的下有界内射复形。
考虑各项为
$$
A^{p, q} = a_{p, *}\mathcal{I}^q_p
$$
的双复形；其水平箭头来自引理
\ref{spaces-simplicial-lemma-augmentation-cech-complex}，
竖直箭头来自复形 $\mathcal{I}^\bullet$ 的微分。
该双复形的各行在正次数上正合，并在次数 $0$ 上给出
$a_*\mathcal{I}^q$。另一方面，由于到 $\mathcal{C}_p$ 的限制正合
（引理 \ref{spaces-simplicial-lemma-restriction-to-components-site}），复形
$\mathcal{I}_p^\bullet$ 在 $D(\mathcal{C}_p)$ 中表示 $K_p$。
诸层 $\mathcal{I}_p^q$ 是 $\mathcal{C}_p$ 上的内射阿贝尔层
（引理 \ref{spaces-simplicial-lemma-restriction-injective-to-component-site}）。
因此，各列的上同调计算 $R^qa_{p, *}K_p$。
应用《同调代数》引理 \ref{homology-lemma-first-quadrant-ss} 与
\ref{homology-lemma-double-complex-gives-resolution} 即得结论。
\end{proof}





\section{赋环单纯站点上的上同调}
\label{spaces-simplicial-section-cohomology-simplicial-sites-modules}

\noindent
本节是第 \ref{spaces-simplicial-section-cohomology-simplicial-sites} 节
对于模层的类似版本。

\medskip\noindent
在情形 \ref{spaces-simplicial-situation-simplicial-site} 中，设 $\mathcal{O}$ 是
$\mathcal{C}_{total}$ 上的环层。
在以下各引理的陈述中，我们以
$g_n : (\Sh(\mathcal{C}_n), \mathcal{O}_n) \to
(\Sh(\mathcal{C}_{total}), \mathcal{O})$
表示引理 \ref{spaces-simplicial-lemma-restriction-module-to-components-site}
中的赋环拓扑斯态射。
若 $\varphi : [m] \to [n]$ 是 $\Delta$ 的一个态射，则赋环拓扑斯图
$$
\xymatrix{
(\Sh(\mathcal{C}_n), \mathcal{O}_n) \ar[rd]_{g_n} \ar[rr]_{f_\varphi} & &
(\Sh(\mathcal{C}_m), \mathcal{O}_m) \ar[ld]^{g_m} \\
& (\Sh(\mathcal{C}_{total}), \mathcal{O})
}
$$
并不交换，但存在来自映射
$\mathcal{F}(\varphi) : f_\varphi^{-1}\mathcal{F}_m \to \mathcal{F}_n$
的 $2$-态射 $g_n \to g_m \circ f_\varphi$。
参见《站点》第 \ref{sites-section-2-category} 节。

\begin{lemma}
\label{spaces-simplicial-lemma-simplicial-resolution-ringed}
在情形 \ref{spaces-simplicial-situation-simplicial-site} 中，设 $\mathcal{O}$ 是
$\mathcal{C}_{total}$ 上的环层。存在 $\mathcal{O}$-模复形
$$
\ldots \to g_{2!}\mathcal{O}_2 \to g_{1!}\mathcal{O}_1 \to g_{0!}\mathcal{O}_0
$$
；它构成 $\mathcal{O}$ 的一个消解。
这里 $g_{n!}$ 如引理
\ref{spaces-simplicial-lemma-restriction-module-to-components-site} 所述。
\end{lemma}

\begin{proof}
我们将使用引理 \ref{spaces-simplicial-lemma-restriction-to-components-site}
中对 $g_{n!}$ 的描述。取 $\sum (-1)^i d^n_i$ 为该复形的映射，
其中 $d^n_i : g_{n!}\mathcal{O}_n \to g_{n - 1!}\mathcal{O}_{n - 1}$
是映射
$\mathcal{O}_n \to \bigoplus_{[n - 1] \to [n]} \mathcal{O}_n =
g_n^*g_{n - 1!}\mathcal{O}_{n - 1}$
的伴随；这里该映射对应于标记为
$\delta^n_i : [n - 1] \to [n]$ 的因子。
于是将 $g_m^{-1}$ 作用于该复形，得到 $\mathcal{C}_m$ 上的复形
$$
\ldots \to
\bigoplus\nolimits_{\alpha \in \Mor_\Delta([2], [m])]} \mathcal{O}_m \to
\bigoplus\nolimits_{\alpha \in \Mor_\Delta([1], [m])]} \mathcal{O}_m \to
\bigoplus\nolimits_{\alpha \in \Mor_\Delta([0], [m])]} \mathcal{O}_m
$$
。
换言之，这是与单纯集合 $\Delta[m]$ 上的自由
$\mathcal{O}_m$-模相关联的复形；参见《单纯对象》例
\ref{simplicial-example-simplex-simplicial-set}。
由于 $\Delta[m]$ 与 $\Delta[0]$ 同伦等价（参见《单纯对象》例
\ref{simplicial-example-simplex-contractible}），
% P10-E0178
且“取其上的自由阿贝尔层”是一个函子，故上述复形同伦等价于
$\Delta[0]$ 上的自由阿贝尔层
（《单纯对象》注 \ref{simplicial-remark-homotopy-better} 及
引理 \ref{simplicial-lemma-homotopy-equivalence-s-N}）。
此复形在正次数上无环，并在次数 $0$ 上等于 $\mathcal{O}_m$。
\end{proof}

\begin{lemma}
\label{spaces-simplicial-lemma-cech-complex-modules}
在情形 \ref{spaces-simplicial-situation-simplicial-site} 中，设 $\mathcal{O}$ 是环层，
$\mathcal{F}$ 是 $\mathcal{O}$-模层。存在典范复形
$$
0 \to \Gamma(\mathcal{C}_{total}, \mathcal{F}) \to
\Gamma(\mathcal{C}_0, \mathcal{F}_0) \to
\Gamma(\mathcal{C}_1, \mathcal{F}_1) \to
\Gamma(\mathcal{C}_2, \mathcal{F}_2) \to \ldots
$$
；它在次数 $-1, 0$ 上正合；若 $\mathcal{F}$ 是内射
$\mathcal{O}$-模，则它处处正合。
\end{lemma}

\begin{proof}
注意
$\Hom(\mathcal{O}, \mathcal{F}) = \Gamma(\mathcal{C}_{total}, \mathcal{F})$
且
$\Hom(g_{n!}\mathcal{O}_n, \mathcal{F}) = \Gamma(\mathcal{C}_n, \mathcal{F}_n)$.
因此，由引理 \ref{spaces-simplicial-lemma-simplicial-resolution-ringed}
以及当 $\mathcal{F}$ 内射时 $\Hom(-, \mathcal{F})$ 正合这一事实，
立即得到本引理。
\end{proof}

\begin{lemma}
\label{spaces-simplicial-lemma-simplicial-module-cohomology-site}
在情形 \ref{spaces-simplicial-situation-simplicial-site} 中，设 $\mathcal{O}$ 是环层。
对属于 $D^+(\mathcal{O})$ 的 $K$，存在谱序列
$(E_r, d_r)_{r \geq 0}$，其中
$$
E_1^{p, q} = H^q(\mathcal{C}_p, K_p),\quad
d_1^{p, q} : E_1^{p, q} \to E_1^{p + 1, q}
$$
；它收敛到 $H^{p + q}(\mathcal{C}_{total}, K)$。
此谱序列关于 $K$ 是函子性的。
\end{lemma}

\begin{proof}
设 $\mathcal{I}^\bullet$ 是表示 $K$ 的下有界内射
$\mathcal{O}$-模复形。考虑各项为
$$
A^{p, q} = \Gamma(\mathcal{C}_p, \mathcal{I}^q_p)
$$
的双复形；其水平箭头来自引理
\ref{spaces-simplicial-lemma-cech-complex-modules}，竖直箭头来自复形
$\mathcal{I}^\bullet$ 的微分。注意
% P10-E0179
$\Gamma(\mathcal{D}, -) =
\Hom_{\mathcal{O}_\mathcal{D}}(\mathcal{O}_\mathcal{D}, -)$
在 $\textit{Mod}(\mathcal{O}_\mathcal{D})$ 上成立。因此，该引理说明
双复形的各行在正次数上正合，并在次数 $0$ 上给出
$\Gamma(\mathcal{C}_{total}, \mathcal{I}^q)$。
于是，由《同调代数》引理
\ref{homology-lemma-double-complex-gives-resolution}，
与该双复形相关联的总复形计算
$R\Gamma(\mathcal{C}_{total}, K)$。
另一方面，由于到 $\mathcal{C}_p$ 的限制正合
（引理 \ref{spaces-simplicial-lemma-restriction-to-components-site}），复形
$\mathcal{I}_p^\bullet$ 在 $D(\mathcal{C}_p)$ 中表示 $K_p$。
诸层 $\mathcal{I}_p^q$ 在 $\mathcal{C}_p$ 上完全无环
（引理 \ref{spaces-simplicial-lemma-restriction-injective-to-component-limp}）。
因此，由 Leray 无环性引理（《导出范畴》引理
\ref{derived-lemma-leray-acyclicity}）以及《站点上的上同调》引理
\ref{sites-cohomology-lemma-limp-acyclic}，各列的上同调计算群
$H^q(\mathcal{C}_p, K_p)$。
应用《同调代数》引理 \ref{homology-lemma-first-quadrant-ss}
即得结论。
\end{proof}

\begin{lemma}
\label{spaces-simplicial-lemma-sanity-check-modules}
在情形 \ref{spaces-simplicial-situation-simplicial-site} 中，设 $\mathcal{O}$ 是环层，
$U \in \Ob(\mathcal{C}_n)$，且
$\mathcal{F} \in \textit{Mod}(\mathcal{O})$。
则 $H^p(U, \mathcal{F}) = H^p(U, g_n^*\mathcal{F})$，
其中在左端，$U$ 被视为 $\mathcal{C}_{total}$ 的一个对象。
\end{lemma}

\begin{proof}
注意，“将 $U$ 视为 $\mathcal{C}_{total}$ 的对象”这一说法，
在情形 (A) 下由引理 \ref{spaces-simplicial-lemma-simplicial-site-site}
中的 $\mathcal{C}_{total}$ 构造得到解释，在情形 (B) 下则由引理
\ref{spaces-simplicial-lemma-simplicial-cocontinuous-site} 中的构造得到解释。
在两种情形下，函子 $\mathcal{C}_n \to \mathcal{C}$
均连续且余连续；参见引理 \ref{spaces-simplicial-lemma-restriction-to-components-site}。
此外，按定义有 $g_n^{-1}\mathcal{O} = \mathcal{O}_n$。
因此，该结论是《站点上的上同调》引理
\ref{sites-cohomology-lemma-pullback-same-cohomology} 的特例。
\end{proof}






\section{赋环单纯站点的上同调与增广}
\label{spaces-simplicial-section-cohomology-augmentation-ringed-simplicial-sites}

\noindent
本节是第 \ref{spaces-simplicial-section-cohomology-augmentation-simplicial-sites} 节
对于模层的类似版本。

\medskip\noindent
考虑情形 \ref{spaces-simplicial-situation-simplicial-site} 所述的单纯站点 $\mathcal{C}$。
设 $a_0$ 是注 \ref{spaces-simplicial-remark-augmentation-site}
所述的朝向站点 $\mathcal{D}$ 的增广。
设 $\mathcal{O}$ 是 $\mathcal{C}_{total}$ 上的环层，
$\mathcal{O}_\mathcal{D}$ 是 $\mathcal{D}$ 上的环层。
假设给定态射
$$
a^\sharp : \mathcal{O}_\mathcal{D} \longrightarrow a_*\mathcal{O}
$$
，其中 $a$ 如引理 \ref{spaces-simplicial-lemma-augmentation-site} 所述。
从而，我们得到赋环拓扑斯的态射
$$
a :
(\Sh(\mathcal{C}_{total}), \mathcal{O})
\longrightarrow
(\Sh(\mathcal{D}), \mathcal{O}_\mathcal{D})
$$
我们按照引理 \ref{spaces-simplicial-lemma-restriction-module-to-components-site}，
将 $g_n : (\Sh(\mathcal{C}_n), \mathcal{O}_n) \to
(\Sh(\mathcal{C}_{total}), \mathcal{O})$ 视为赋环拓扑斯的态射。
取赋环拓扑斯态射的复合 $a_n = a \circ g_n$
（引理 \ref{spaces-simplicial-lemma-augmentation-site}），便得到
$$
a_n :
(\Sh(\mathcal{C}_n), \mathcal{O}_n)
\longrightarrow
(\Sh(\mathcal{D}), \mathcal{O}_\mathcal{D})
$$
利用转移映射 $f_\varphi^{-1}\mathcal{O}_m \to \mathcal{O}_n$，
我们得到赋环拓扑斯的态射
$$
f_\varphi : (\Sh(\mathcal{C}_n), \mathcal{O}_n) \to
(\Sh(\mathcal{C}_m), \mathcal{O}_m)
$$
，使得对所有 $\varphi : [m] \to [n]$，作为赋环拓扑斯的态射均有
$a_m \circ f_\varphi = a_n$。

\begin{lemma}
\label{spaces-simplicial-lemma-flat-augmentation-modules}
采用上述记号。态射
$a : (\Sh(\mathcal{C}_{total}), \mathcal{O}) \to
(\Sh(\mathcal{D}), \mathcal{O}_\mathcal{D})$
平坦，当且仅当对 $n \geq 0$，态射
$a_n : (\Sh(\mathcal{C}_n), \mathcal{O}_n) \to
(\Sh(\mathcal{D}), \mathcal{O}_\mathcal{D})$
平坦。
\end{lemma}

\begin{proof}
由于 $g_n : (\Sh(\mathcal{C}_n), \mathcal{O}_n) \to
(\Sh(\mathcal{C}_{total}), \mathcal{O})$ 平坦，可见若 $a$ 平坦，
则作为复合的 $a_n = a \circ g_n$ 也平坦。
反之，假设对所有 $n$，$a_n$ 都平坦。
我们须验证 $\mathcal{O}$ 作为 $a^{-1}\mathcal{O}_\mathcal{D}$-模层是平坦的。
设 $\mathcal{F} \to \mathcal{G}$ 是
$a^{-1}\mathcal{O}_\mathcal{D}$-模的单射。我们须证明
$$
\mathcal{F} \otimes_{a^{-1}\mathcal{O}_\mathcal{D}} \mathcal{O}
\to
\mathcal{G} \otimes_{a^{-1}\mathcal{O}_\mathcal{D}} \mathcal{O}
$$
是单射。这可以在 $\mathcal{C}_n$ 上检验，即在施用
$g_n^{-1}$ 之后检验。由于 $g_n^{-1}\mathcal{O} = \mathcal{O}_n$，
故 $g_n^* = g_n^{-1}$，于是得到
$$
g_n^{-1}\mathcal{F} \otimes_{g_n^{-1}a^{-1}\mathcal{O}_\mathcal{D}}
\mathcal{O}_n
\to
g_n^{-1}\mathcal{G} \otimes_{g_n^{-1}a^{-1}\mathcal{O}_\mathcal{D}}
\mathcal{O}_n
$$
；它是单射，因为
$g_n^{-1}a^{-1}\mathcal{O}_\mathcal{D} = a_n^{-1}\mathcal{O}_\mathcal{D}$
，且我们假设了 $a_n$ 平坦。
\end{proof}

\begin{lemma}
\label{spaces-simplicial-lemma-simplicial-resolution-augmentation-modules}
采用上述记号。对 $\mathcal{O}_\mathcal{D}$-模 $\mathcal{G}$，
存在正合复形
$$
\ldots \to
g_{2!}(a_2^*\mathcal{G}) \to
g_{1!}(a_1^*\mathcal{G}) \to
g_{0!}(a_0^*\mathcal{G}) \to
a^*\mathcal{G} \to 0
$$
；它由 $\mathcal{C}_{total}$ 上的 $\mathcal{O}$-模层组成。
这里 $g_{n!}$ 如引理
\ref{spaces-simplicial-lemma-restriction-module-to-components-site} 所述。
\end{lemma}

\begin{proof}
注意，$a^*\mathcal{G}$ 是 $\mathcal{C}_{total}$ 上的
$\mathcal{O}$-模，其到 $\mathcal{C}_m$ 的限制为
$\mathcal{O}_m$-模 $a_m^*\mathcal{G}$。
引理 \ref{spaces-simplicial-lemma-restriction-module-to-components-site}
中对模上的函子 $g_{n!}$ 的描述表明，
$g_{n!}(a_n^*\mathcal{G})$ 是 $\mathcal{C}_{total}$ 上的
$\mathcal{O}$-模，其到 $\mathcal{C}_m$ 的限制为 $\mathcal{O}_m$-模
$$
\bigoplus\nolimits_{\varphi : [n] \to [m]} f_\varphi^*a_n^*\mathcal{G} =
\bigoplus\nolimits_{\varphi : [n] \to [m]} a_m^*\mathcal{G}
$$
。其余证明与引理 \ref{spaces-simplicial-lemma-simplicial-resolution-augmentation}
的证明完全相同，只需将 $a_m^{-1}\mathcal{G}$ 替换为
$a_m^*\mathcal{G}$。
\end{proof}

\begin{lemma}
\label{spaces-simplicial-lemma-augmentation-cech-complex-modules}
采用上述记号。对 $\mathcal{C}_{total}$ 上的
$\mathcal{O}$-模 $\mathcal{F}$，存在典范复形
$$
0 \to a_*\mathcal{F} \to a_{0, *}\mathcal{F}_0 \to a_{1, *}\mathcal{F}_1 \to
a_{2, *}\mathcal{F}_2 \to \ldots
$$
；它由 $\mathcal{O}_\mathcal{D}$-模组成，并在次数 $-1, 0$ 上正合。
若 $\mathcal{F}$ 是内射 $\mathcal{O}$-模，则该复形在所有次数上正合，
且对任意 $\mathcal{O}_\mathcal{D}$-模 $\mathcal{G}$，
施用函子 $\Hom_{\mathcal{O}_\mathcal{D}}(\mathcal{G}, -)$ 后仍正合。
\end{lemma}

\begin{proof}
为构造该复形，由 Yoneda 引理，只需对 $\mathcal{D}$ 上任意
$\mathcal{O}_\mathcal{D}$-模 $\mathcal{G}$ 构造复形
$$
0 \to \Hom_{\mathcal{O}_\mathcal{D}}(\mathcal{G}, a_*\mathcal{F}) \to
\Hom_{\mathcal{O}_\mathcal{D}}(\mathcal{G}, a_{0, *}\mathcal{F}_0) \to
\Hom_{\mathcal{O}_\mathcal{D}}(\mathcal{G}, a_{1, *}\mathcal{F}_1) \to \ldots
$$
，并使其关于 $\mathcal{G}$ 是函子性的。为此，将
$\Hom_\mathcal{O}(-, \mathcal{F})$ 作用于引理
\ref{spaces-simplicial-lemma-simplicial-resolution-augmentation-modules} 的正合复形，
并使用拉回与推出的伴随性。
由于 $\Hom_\mathcal{O}(-, \mathcal{F})$ 左正合，次数 $-1, 0$
上的正合性由该构造得出。若 $\mathcal{F}$ 是内射
$\mathcal{O}$-模，则 $\Hom_\mathcal{O}(-, \mathcal{F})$ 正合，
故此复形正合。
\end{proof}

\begin{lemma}
\label{spaces-simplicial-lemma-augmentation-spectral-sequence-modules}
采用上述记号。对任意属于 $D^+(\mathcal{O})$ 的 $K$，
在 $\textit{Mod}(\mathcal{O}_\mathcal{D})$ 中存在谱序列
$(E_r, d_r)_{r \geq 0}$，其中
$$
E_1^{p, q} = R^qa_{p, *} K_p
$$
；它收敛到 $R^{p + q}a_*K$。此谱序列关于 $K$ 是函子性的。
\end{lemma}

\begin{proof}
设 $\mathcal{I}^\bullet$ 是表示 $K$ 的下有界内射
$\mathcal{O}$-模复形。考虑各项为
$$
A^{p, q} = a_{p, *}\mathcal{I}^q_p
$$
的双复形；其水平箭头来自引理
\ref{spaces-simplicial-lemma-augmentation-cech-complex-modules}，
竖直箭头来自复形 $\mathcal{I}^\bullet$ 的微分。
该引理说明，双复形的各行在正次数上正合，并在次数 $0$ 上给出
$a_*\mathcal{I}^q$。于是，由《同调代数》引理
\ref{homology-lemma-double-complex-gives-resolution}，
与该双复形相关联的总复形计算 $Ra_*K$。
另一方面，由于到 $\mathcal{C}_p$ 的限制正合
（引理 \ref{spaces-simplicial-lemma-restriction-to-components-site}），复形
$\mathcal{I}_p^\bullet$ 在 $D(\mathcal{C}_p)$ 中表示 $K_p$。
诸层 $\mathcal{I}_p^q$ 在 $\mathcal{C}_p$ 上完全无环
（引理 \ref{spaces-simplicial-lemma-restriction-injective-to-component-limp}）。
因此，由 Leray 无环性引理（《导出范畴》引理
\ref{derived-lemma-leray-acyclicity}）以及《站点上的上同调》引理
\ref{sites-cohomology-lemma-limp-acyclic}，各列的上同调为层
$R^qa_{p, *}K_p$。应用《同调代数》引理
\ref{homology-lemma-first-quadrant-ss} 即得结论。
\end{proof}





\section{笛卡尔层与模}
\label{spaces-simplicial-section-cartesian}

\noindent
定义如下。

\begin{definition}
\label{spaces-simplicial-definition-cartesian-sheaf}
在情形 \ref{spaces-simplicial-situation-simplicial-site} 中。
\begin{enumerate}
\item 若对所有 $\varphi : [m] \to [n]$，映射
$\mathcal{F}(\varphi) : f_\varphi^{-1}\mathcal{F}_m \to \mathcal{F}_n$
都是同构，则称 $\mathcal{C}_{total}$ 上的集合层或阿贝尔群层
$\mathcal{F}$ 是{\it 笛卡尔的}。
\item 若 $\mathcal{O}$ 是 $\mathcal{C}_{total}$ 上的环层，
且对所有 $\varphi : [m] \to [n]$，映射
$f_\varphi^*\mathcal{F}_m \to \mathcal{F}_n$ 都是同构，
则称 $\mathcal{O}$-模层 $\mathcal{F}$ 是{\it 笛卡尔的}。
\item 若对所有 $\varphi : [m] \to [n]$，映射
$f_\varphi^{-1}K_m \to K_n$
都是同构，则称 $D(\mathcal{C}_{total})$ 的对象 $K$
是{\it 笛卡尔的}。
\item 若 $\mathcal{O}$ 是 $\mathcal{C}_{total}$ 上的环层，
且对所有 $\varphi : [m] \to [n]$，映射
$Lf_\varphi^*K_m \to K_n$
都是同构，则称 $D(\mathcal{O})$ 的对象 $K$ 是{\it 笛卡尔的}。
\end{enumerate}
\end{definition}

\noindent
当然，纤维化范畴的笛卡尔截面有一个一般概念，而以上仅是其若干例子。
% P10-E0181
关于拉回的这一性质只需对退化映射检验。

\begin{lemma}
\label{spaces-simplicial-lemma-check-cartesian-module}
在情形 \ref{spaces-simplicial-situation-simplicial-site} 中。
\begin{enumerate}
\item 集合层或阿贝尔群层 $\mathcal{F}$ 是笛卡尔的，当且仅当映射
$(f_{\delta^n_j})^{-1}\mathcal{F}_{n - 1} \to \mathcal{F}_n$
是同构。
\item $D(\mathcal{C}_{total})$ 的对象 $K$ 是笛卡尔的，
当且仅当映射
$(f_{\delta^n_j})^{-1}K_{n - 1} \to K_n$
是同构。
\item 若 $\mathcal{O}$ 是 $\mathcal{C}_{total}$ 上的环层，
则 $\mathcal{O}$-模层 $\mathcal{F}$ 是笛卡尔的，当且仅当映射
$(f_{\delta^n_j})^*\mathcal{F}_{n - 1} \to \mathcal{F}_n$
是同构。
\item 若 $\mathcal{O}$ 是 $\mathcal{C}_{total}$ 上的环层，
则 $D(\mathcal{O})$ 的对象 $K$ 是笛卡尔的，当且仅当映射
$L(f_{\delta^n_j})^*K_{n - 1} \to K_n$
是同构。
% P10-E0182
\item 在此补充更多内容。
\end{enumerate}
\end{lemma}

\begin{proof}
在每一种情形下，关键都是拉回函子的复合仍为拉回函子；
关于 (4)，参见《站点上的上同调》引理
\ref{sites-cohomology-lemma-derived-pullback-composition}。
我们说明 (1) 中的论证，并略去其余情形的证明。
% P10-E0183
范畴 $\Delta$ 由态射 $\delta^n_j$ 与 $\sigma^n_j$ 生成；
参见《单纯对象》引理 \ref{simplicial-lemma-face-degeneracy}。
因此，只需检验映射
$(f_{\delta^n_j})^{-1}\mathcal{F}_{n - 1} \to \mathcal{F}_n$
与 $(f_{\sigma^n_j})^{-1}\mathcal{F}_{n + 1} \to \mathcal{F}_n$
是同构；关于记号，参见《单纯对象》引理
\ref{simplicial-lemma-characterize-simplicial-object}。
由于 $\sigma^n_j \circ \delta_j^{n + 1} = \text{id}_{[n]}$，复合
$$
\mathcal{F}_n =
(f_{\sigma^n_j})^{-1}
(f_{\delta_j^{n + 1}})^{-1}
\mathcal{F}_n \to
(f_{\sigma^n_j})^{-1}
\mathcal{F}_{n + 1} \to
\mathcal{F}_n
$$
是恒等映射。因此，对 $\delta^{n + 1}_j$ 的结论蕴含
对 $\sigma^n_j$ 的结论。
\end{proof}

\begin{lemma}
\label{spaces-simplicial-lemma-augmentation-cartesian-module}
在情形 \ref{spaces-simplicial-situation-simplicial-site} 中，设 $a_0$ 是注
\ref{spaces-simplicial-remark-augmentation-site} 所述的朝向站点 $\mathcal{D}$ 的增广。
\begin{enumerate}
\item $\mathcal{D}$ 上集合层或阿贝尔群层的拉回
$a^{-1}\mathcal{G}$ 是笛卡尔的。
\item $D(\mathcal{D})$ 的对象 $K$ 的拉回 $a^{-1}K$ 是笛卡尔的。
\end{enumerate}
设 $\mathcal{O}$ 是 $\mathcal{C}_{total}$ 上的环层，
$\mathcal{O}_\mathcal{D}$ 是 $\mathcal{D}$ 上的环层，且
$a^\sharp : \mathcal{O}_\mathcal{D} \to a_*\mathcal{O}$ 是第
\ref{spaces-simplicial-section-cohomology-augmentation-ringed-simplicial-sites} 节所述的态射。
\begin{enumerate}
\item[(3)] $\mathcal{O}_\mathcal{D}$-模层的拉回
$a^*\mathcal{F}$ 是笛卡尔的。
\item[(4)] $D(\mathcal{O}_\mathcal{D})$ 的对象 $K$
的导出拉回 $La^*K$ 是笛卡尔的。
\end{enumerate}
\end{lemma}

\begin{proof}
这立即由对所有 $\varphi : [m] \to [n]$ 成立的等式
$a_m \circ f_\varphi = a_n$ 得出。参见引理
\ref{spaces-simplicial-lemma-augmentation-site} 以及第
\ref{spaces-simplicial-section-cohomology-augmentation-ringed-simplicial-sites} 节中的讨论。
\end{proof}

\begin{lemma}
\label{spaces-simplicial-lemma-characterize-cartesian}
在情形 \ref{spaces-simplicial-situation-simplicial-site} 中，
笛卡尔集合层（相应地，笛卡尔阿贝尔群层）范畴等价于
二元组 $(\mathcal{F}, \alpha)$ 的范畴，其中 $\mathcal{F}$ 是
$\mathcal{C}_0$ 上的集合层（相应地，阿贝尔群层），且
$$
\alpha :
(f_{\delta_1^1})^{-1}\mathcal{F}
\longrightarrow (f_{\delta_0^1})^{-1}\mathcal{F}
$$
是 $\mathcal{C}_1$ 上集合层（相应地，阿贝尔群层）的同构，满足
$(f_{\delta^2_1})^{-1}\alpha =
(f_{\delta^2_0})^{-1}\alpha \circ (f_{\delta^2_2})^{-1}\alpha$
；这里等式两端都是 $\mathcal{C}_2$ 上层的映射。
\end{lemma}

\begin{proof}
简记
$d^n_j = f_{\delta^n_j} : \Sh(\mathcal{C}_n) \to \Sh(\mathcal{C}_{n - 1})$.
引理陈述中关于 $\alpha$ 的条件是有意义的，因为
$$
d^1_1 \circ d^2_2 = d^1_1 \circ d^2_1, \quad
d^1_1 \circ d^2_0 = d^1_0 \circ d^2_2, \quad
d^1_0 \circ d^2_0 = d^1_0 \circ d^2_1
$$
是拓扑斯态射 $\Sh(\mathcal{C}_2) \to \Sh(\mathcal{C}_0)$ 的等式；
参见《单纯对象》注 \ref{simplicial-remark-relations}。
因此，可将这些映射表示为
$$
\xymatrix{
& (d^2_0)^{-1}(d^1_1)^{-1}\mathcal{F} \ar[r]_-{(d^2_0)^{-1}\alpha} &
(d^2_0)^{-1}(d^1_0)^{-1}\mathcal{F} \ar@{=}[rd] & \\
(d^2_2)^{-1}(d^1_0)^{-1}\mathcal{F} \ar@{=}[ru] & & &
(d^2_1)^{-1}(d^1_0)^{-1}\mathcal{F} \\
& (d^2_2)^{-1}(d^1_1)^{-1}\mathcal{F} \ar[lu]^{(d^2_2)^{-1}\alpha} \ar@{=}[r] &
(d^2_1)^{-1}(d^1_1)^{-1}\mathcal{F} \ar[ru]_{(d^2_1)^{-1}\alpha}
}
$$
，而该条件表示此图交换。显然，给定 $\mathcal{C}_{total}$ 上的
笛卡尔集合层（相应地，阿贝尔群层）$\mathcal{G}$，
可以令 $\mathcal{F} = \mathcal{G}_0$，并令 $\alpha$ 等于复合
% P10-E0184
$$
(d_1^1)^{-1}\mathcal{G}_0 \to \mathcal{G}_1
\leftarrow (d_1^0)^{-1}\mathcal{G}_0
$$
；由于 $\mathcal{G}$ 是笛卡尔的，图中箭头均可逆。
为证明此函子是等价，我们构造一个拟逆。
拟逆的构造类似于《下降》第 \ref{descent-section-descent-modules} 节
讨论的构造；我们借用其中的记号
$\tau^n_i : [0] \to [n]$, $0 \mapsto i$ 以及
$\tau^n_{ij} : [1] \to [n]$, $0 \mapsto i$, $1 \mapsto j$。
具体而言，给定引理所述的二元组 $(\mathcal{F}, \alpha)$，令
$\mathcal{G}_n = (f_{\tau^n_n})^{-1}\mathcal{F}$。
给定 $\varphi : [n] \to [m]$，用下图定义
$\mathcal{G}(\varphi) : (f_\varphi)^{-1}\mathcal{G}_n \to \mathcal{G}_m$：
$$
\xymatrix{
(f_\varphi)^{-1}\mathcal{G}_n \ar@{=}[r] &
(f_\varphi)^{-1}(f_{\tau^n_n})^{-1}\mathcal{F} \ar@{=}[r] &
(f_{\tau^m_{\varphi(n)}})^{-1}\mathcal{F} \ar@{=}[r] &
(f_{\tau^m_{\varphi(n)m}})^{-1}(d^1_1)^{-1}\mathcal{F}
\ar[d]^{(f_{\tau^m_{\varphi(n)m}})^{-1}\alpha} \\
&
\mathcal{G}_m \ar@{=}[r] &
(f_{\tau^m_m})^{-1}\mathcal{F} \ar@{=}[r] &
(f_{\tau^m_{\varphi(n)m}})^{-1}(d^1_0)^{-1}\mathcal{F}
}
$$
我们略去如下验证：上图的交换性蕴含这些映射能够正确复合，
因而给出 $\mathcal{C}_{total}$ 上的层；参见引理
\ref{spaces-simplicial-lemma-describe-sheaves-simplicial-site-site}。
我们也略去这两个函子互为拟逆的验证。
\end{proof}

\begin{lemma}
\label{spaces-simplicial-lemma-characterize-cartesian-modules}
在情形 \ref{spaces-simplicial-situation-simplicial-site} 中，设 $\mathcal{O}$ 是
$\mathcal{C}_{total}$ 上的环层。
笛卡尔 $\mathcal{O}$-模范畴等价于二元组
$(\mathcal{F}, \alpha)$ 的范畴，其中 $\mathcal{F}$ 是
$\mathcal{O}_0$-模，且
$$
\alpha :
(f_{\delta_1^1})^*\mathcal{F}
\longrightarrow (f_{\delta_0^1})^*\mathcal{F}
$$
是 $\mathcal{O}_1$-模的同构，满足
$(f_{\delta^2_1})^*\alpha =
(f_{\delta^2_0})^*\alpha \circ (f_{\delta^2_2})^*\alpha$
；这里等式两端都是 $\mathcal{O}_2$-模映射。
\end{lemma}

\begin{proof}
证明与引理 \ref{spaces-simplicial-lemma-characterize-cartesian} 的证明相同，
只需将阿贝尔群层的拉回替换为模的拉回。
\end{proof}

\begin{lemma}
\label{spaces-simplicial-lemma-Serre-subcat-cartesian-modules}
在情形 \ref{spaces-simplicial-situation-simplicial-site} 中。
\begin{enumerate}
\item 笛卡尔阿贝尔层的满子范畴构成
$\textit{Ab}(\mathcal{C}_{total})$ 的弱 Serre 子范畴。
笛卡尔阿贝尔层系统的余极限仍是笛卡尔的。
\item 设 $\mathcal{O}$ 是 $\mathcal{C}_{total}$ 上的环层，
并且态射
$$
f_{\delta^n_j} : (\Sh(\mathcal{C}_n), \mathcal{O}_n)
\to (\Sh(\mathcal{C}_{n - 1}), \mathcal{O}_{n - 1})
$$
平坦。则笛卡尔 $\mathcal{O}$-模的满子范畴构成
$\textit{Mod}(\mathcal{O})$ 的弱 Serre 子范畴。
笛卡尔 $\mathcal{O}$-模系统的余极限仍是笛卡尔的。
\end{enumerate}
\end{lemma}

\begin{proof}
为说明 (1) 中得到弱 Serre 子范畴，我们检验《同调代数》引理
\ref{homology-lemma-characterize-weak-serre-subcategory} 所列的条件。
首先，若 $\varphi : \mathcal{F} \to \mathcal{G}$ 是笛卡尔阿贝尔层间的映射，
则 $\Ker(\varphi)$ 与 $\Coker(\varphi)$ 也是笛卡尔的，
因为限制函子
$\Sh(\mathcal{C}_{total}) \to \Sh(\mathcal{C}_n)$
以及函子 $f_\varphi^{-1}$ 都正合。类似地，若
$$
0 \to \mathcal{F} \to \mathcal{H} \to \mathcal{G} \to 0
$$
是 $\mathcal{C}_{total}$ 上阿贝尔层的短正合列，且
$\mathcal{F}$ 与 $\mathcal{G}$ 都是笛卡尔的，则由五引理可知
$\mathcal{H}$ 也是笛卡尔的。关于余极限的性质，使用如下事实：
拉回是左伴随，故余极限与拉回交换。
在模的情形下，采用同样的论证，并使用平坦拉回的正合性
（《站点上的模》引理 \ref{sites-modules-lemma-flat-pullback-exact}），
以及只需对 $f_{\delta^n_j}$ 检验该条件这一事实；
参见引理 \ref{spaces-simplicial-lemma-check-cartesian-module}。
\end{proof}

\begin{remark}[警告]
\label{spaces-simplicial-remark-warning-cartesian-modules}
尽管有引理 \ref{spaces-simplicial-lemma-Serre-subcat-cartesian-modules}，
笛卡尔 $\mathcal{O}$-模范畴仍可能是阿贝尔范畴，却并非
$\textit{Mod}(\mathcal{O})$ 的 Serre 子范畴。
具体而言，假设只知道 $f_{\delta_1^1}$ 与 $f_{\delta_0^1}$ 平坦。
由引理 \ref{spaces-simplicial-lemma-characterize-cartesian-modules} 很容易推出，
笛卡尔 $\mathcal{O}$-模范畴是阿贝尔范畴。
但若（例如）$f_{\delta_0^2}$ 不平坦，则没有理由认为从
笛卡尔 $\mathcal{O}$-模范畴到全体 $\mathcal{O}$-模范畴的包含函子正合。
\end{remark}

\begin{lemma}
\label{spaces-simplicial-lemma-derived-cartesian-modules}
在情形 \ref{spaces-simplicial-situation-simplicial-site} 中。
\begin{enumerate}
\item $D(\mathcal{C}_{total})$ 的对象 $K$ 是笛卡尔的，
当且仅当对所有 $q$，$H^q(K)$ 都是笛卡尔阿贝尔层。
\item 设 $\mathcal{O}$ 是 $\mathcal{C}_{total}$ 上的环层，
并且态射
$f_{\delta^n_j} : (\Sh(\mathcal{C}_n), \mathcal{O}_n)
\to (\Sh(\mathcal{C}_{n - 1}), \mathcal{O}_{n - 1})$ 平坦。
则 $D(\mathcal{O})$ 的对象 $K$ 是笛卡尔的，
当且仅当对所有 $q$，$H^q(K)$ 都是笛卡尔 $\mathcal{O}$-模。
\end{enumerate}
\end{lemma}

\begin{proof}
(1) 成立是因为拉回函子 $(f_\varphi)^{-1}$ 正合。
(2) 由引理 \ref{spaces-simplicial-lemma-check-cartesian-module} 中的刻画，
以及平坦性所给出的
$L(f_{\delta^n_j})^* = (f_{\delta^n_j})^*$ 得出。
\end{proof}

\begin{lemma}
\label{spaces-simplicial-lemma-derived-cartesian-shriek}
在情形 \ref{spaces-simplicial-situation-simplicial-site} 中。
\begin{enumerate}
\item $D(\mathcal{C}_{total})$ 的对象 $K$ 是笛卡尔的，
当且仅当对所有 $n$，典范映射
$$
g_{n!}K_n \longrightarrow
g_{n!}\mathbf{Z} \otimes^\mathbf{L}_\mathbf{Z} K
$$
是同构。
\item 设 $\mathcal{O}$ 是 $\mathcal{C}_{total}$ 上的环层，
并且对所有 $\varphi : [n] \to [m]$，态射
$f_\varphi^{-1}\mathcal{O}_n \to \mathcal{O}_m$ 都平坦。
则 $D(\mathcal{O})$ 的对象 $K$ 是笛卡尔的，
当且仅当对所有 $n$，典范映射
$$
g_{n!}K_n \longrightarrow
g_{n!}\mathcal{O}_n \otimes^\mathbf{L}_\mathcal{O} K
$$
是同构。
\end{enumerate}
\end{lemma}

\begin{proof}
证明 (1)。由于 $g_{n!}$ 正合，它在导出范畴上诱导一个与
$g_n^{-1}$ 伴随的函子。该映射是映射
$K_n \to (g_n^{-1}g_{n!}\mathbf{Z}) \otimes^\mathbf{L}_\mathbf{Z} K_n$
的伴随；后者对应于 $\mathbf{Z} \to g_n^{-1}g_{n!}\mathbf{Z}$，
而这一映射又与
$\text{id} : g_{n!}\mathbf{Z} \to g_{n!}\mathbf{Z}$ 伴随。
利用引理 \ref{spaces-simplicial-lemma-restriction-to-components-site}
中对 $g_{n!}$ 的描述，可见该映射到 $\mathcal{C}_m$ 的限制为
$$
\bigoplus\nolimits_{\varphi : [n] \to [m]} f_\varphi^{-1}K_n
\longrightarrow
\bigoplus\nolimits_{\varphi : [n] \to [m]} K_m
$$
因此，该陈述是清楚的。

\medskip\noindent
证明 (2)。由于 $g_{n!}$ 正合
（引理 \ref{spaces-simplicial-lemma-exactness-g-shriek-modules}），
它在导出范畴上诱导一个与同样正合的 $g_n^*$ 伴随的函子。
该映射是映射
$K_n \to (g_n^*g_{n!}\mathcal{O}_n) \otimes^\mathbf{L}_{\mathcal{O}_n} K_n$
的伴随；后者对应于 $\mathcal{O}_n \to g_n^*g_{n!}\mathcal{O}_n$，
而这一映射又与
$\text{id} : g_{n!}\mathcal{O}_n \to g_{n!}\mathcal{O}_n$ 伴随。
利用引理 \ref{spaces-simplicial-lemma-restriction-module-to-components-site}
中对 $g_{n!}$ 的描述，可见该映射到 $\mathcal{C}_m$ 的限制为
$$
\bigoplus\nolimits_{\varphi : [n] \to [m]} f_\varphi^*K_n
\longrightarrow
\bigoplus\nolimits_{\varphi : [n] \to [m]}
f_\varphi^*\mathcal{O}_n \otimes_{\mathcal{O}_m} K_m =
\bigoplus\nolimits_{\varphi : [n] \to [m]} K_m
$$
因此，该陈述是清楚的。
\end{proof}

\begin{lemma}
\label{spaces-simplicial-lemma-quasi-coherent-sheaf}
在情形 \ref{spaces-simplicial-situation-simplicial-site} 中，设 $\mathcal{O}$ 是
$\mathcal{C}_{total}$ 上的环层，$\mathcal{F}$ 是 $\mathcal{O}$-模层。
则 $\mathcal{F}$ 在《站点上的模》定义
\ref{sites-modules-definition-site-local} 的意义下拟凝聚，
当且仅当 $\mathcal{F}$ 是笛卡尔的，且对所有 $n$，
$\mathcal{F}_n$ 都是拟凝聚 $\mathcal{O}_n$-模。
\end{lemma}

\begin{proof}
假设 $\mathcal{F}$ 拟凝聚。由于拟凝聚模的拉回仍拟凝聚
（《站点上的模》引理 \ref{sites-modules-lemma-local-pullback}），
可见对所有 $n$，$\mathcal{F}_n$ 都是拟凝聚
$\mathcal{O}_n$-模。为证明 $\mathcal{F}$ 是笛卡尔的，
设 $U$ 是 $\mathcal{C}_n$ 的对象，其中 $n$ 为某个整数，并将 $U$
视为 $\mathcal{C}_{total}$ 的对象。由于 $\mathcal{F}$ 拟凝聚，
存在覆盖 $\{U_i \to U\}$，并且对每个 $i$，存在表现
$$
\bigoplus\nolimits_{j \in J_i} \mathcal{O}_{\mathcal{C}_{total}/U_i} \to
\bigoplus\nolimits_{k \in K_i} \mathcal{O}_{\mathcal{C}_{total}/U_i} \to
\mathcal{F}|_{\mathcal{C}_{total}/U_i} \to 0
$$
。由站点 $\mathcal{C}_{total}$ 的构造可知，
$\{U_i \to U\}$ 是 $\mathcal{C}_n$ 中的覆盖。
接着，设 $V$ 是 $\mathcal{C}_m$ 的对象，其中 $m$ 为某个整数，并设
$V \to U$ 是 $\mathcal{C}_{total}$ 中位于
$\varphi : [n] \to [m]$ 之上的态射。纤维积
$V_i = V \times_U U_i$ 存在，并在 $\mathcal{C}_m$ 中得到诱导覆盖
$\{V_i \to V\}$。将上述表现限制到站点
$\mathcal{C}_n/U_i$ 与 $\mathcal{C}_m/V_i$，得到表现
% P10-E0185
$$
\bigoplus\nolimits_{j \in J_i} \mathcal{O}_{\mathcal{C}_m/U_i} \to
\bigoplus\nolimits_{k \in K_i} \mathcal{O}_{\mathcal{C}_m/U_i} \to
\mathcal{F}_n|_{\mathcal{C}_n/U_i} \to 0
$$
以及
$$
\bigoplus\nolimits_{j \in J_i} \mathcal{O}_{\mathcal{C}_m/V_i} \to
\bigoplus\nolimits_{k \in K_i} \mathcal{O}_{\mathcal{C}_m/V_i} \to
\mathcal{F}_m|_{\mathcal{C}_m/V_i} \to 0
$$
。这些表现与映射
$\mathcal{F}(\varphi) : f_\varphi^*\mathcal{F}_n \to \mathcal{F}_m$
相容（因为该映射利用 $\mathcal{F}$ 沿
$\mathcal{C}_{total}$ 中位于 $\varphi$ 之上的态射的限制映射来定义）。
由此可知 $\mathcal{F}(\varphi)|_{\mathcal{C}_m/V_i}$ 是同构。
由于 $\{V_i \to V\}$ 是覆盖，可知
$\mathcal{F}(\varphi)|_{\mathcal{C}_m/V}$ 是同构。
因 $V$ 与 $U$ 任意，这证明了 $\mathcal{F}$ 是笛卡尔的。
（在情形 A 下使用《站点》引理
\ref{sites-lemma-morphism-of-sites-covering}。）

\medskip\noindent
反之，假设对所有 $n$，$\mathcal{F}_n$ 都拟凝聚，
且 $\mathcal{F}$ 是笛卡尔的。于是，对任意 $n$ 及
$\mathcal{C}_n$ 的对象 $U$，可以选取 $\mathcal{C}_n$ 中的覆盖
$\{U_i \to U\}$，并对每个 $i$ 选取表现
% P10-E0186
$$
\bigoplus\nolimits_{j \in J_i} \mathcal{O}_{\mathcal{C}_m/U_i} \to
\bigoplus\nolimits_{k \in K_i} \mathcal{O}_{\mathcal{C}_m/U_i} \to
\mathcal{F}_n|_{\mathcal{C}_n/U_i} \to 0
$$
。将其拉回到 $\mathcal{C}_{total}/U_i$，得到模复形
$$
\bigoplus\nolimits_{j \in J_i} \mathcal{O}_{\mathcal{C}_{total}/U_i} \to
\bigoplus\nolimits_{k \in K_i} \mathcal{O}_{\mathcal{C}_{total}/U_i} \to
\mathcal{F}|_{\mathcal{C}_{total}/U_i} \to 0
$$
。这是 $\mathcal{C}_{total}/U_i$ 上的模复形。
$\mathcal{F}$ 的笛卡尔性蕴含此复形正合。细节从略。
\end{proof}











\section{导出范畴的单纯系统}
\label{spaces-simplicial-section-glueing}

\noindent
本节将在阿贝尔层导出范畴的框架下证明
\cite[Proposition 3.2.9]{BBD} 的一个特例。
模的情形在第 \ref{spaces-simplicial-section-glueing-modules} 节中讨论。

\begin{definition}
\label{spaces-simplicial-definition-cartesian-derived}
在情形 \ref{spaces-simplicial-situation-simplicial-site} 中，
一个{\it 导出范畴的单纯系统}由下列数据组成：
\begin{enumerate}
\item 对每个 $n$，给定 $D(\mathcal{C}_n)$ 的对象 $K_n$；
\item 对每个 $\varphi : [m] \to [n]$，给定
$D(\mathcal{C}_n)$ 中的映射
$K_\varphi : f_\varphi^{-1}K_m \to K_n$。
\end{enumerate}
这些数据须满足：
$$
K_{\varphi \circ \psi} = K_\varphi \circ f_\varphi^{-1}K_\psi :
f_{\varphi \circ \psi}^{-1}K_l = f_\varphi^{-1} f_\psi^{-1}K_l
\longrightarrow
K_n
$$
对 $\Delta$ 的任意态射 $\varphi : [m] \to [n]$ 与
$\psi : [l] \to [m]$，上述等式均成立。
若对所有 $\varphi$，映射 $K_\varphi$ 都是同构，
则称该单纯系统是{\it 笛卡尔的}。
对两个导出范畴的单纯系统，显然可以定义
{\it 导出范畴的单纯系统的态射}。
\end{definition}

\noindent
我们有意给这个概念取了一个长得荒唐的名字。
目标是证明：在某些假设下，导出范畴的单纯系统来自
$D(\mathcal{C}_{total})$ 的一个对象。

\begin{lemma}
\label{spaces-simplicial-lemma-cartesian-objects-derived}
在情形 \ref{spaces-simplicial-situation-simplicial-site} 中，若
$K \in D(\mathcal{C}_{total})$，则
$(K_n, K(\varphi))$ 是导出范畴的单纯系统。
若 $K$ 是笛卡尔的，则该系统也是笛卡尔的。
\end{lemma}

\begin{proof}
这是显然的。
\end{proof}

\begin{lemma}
\label{spaces-simplicial-lemma-construct-cartesian-objects-derived}
在情形 \ref{spaces-simplicial-situation-simplicial-site} 中，假设给定
$K_0 \in D(\mathcal{C}_0)$ 以及同构
$$
\alpha :
f_{\delta_1^1}^{-1}K_0
\longrightarrow
f_{\delta_0^1}^{-1}K_0
$$
，它满足余圈条件。令
$\tau^n_i : [0] \to [n]$, $0 \mapsto i$，并令
$K_n = f_{\tau^n_n}^{-1}K_0$。则诸 $K_n$
构成导出范畴的一个笛卡尔单纯系统。
\end{lemma}

\begin{proof}
请与引理 \ref{spaces-simplicial-lemma-characterize-cartesian} 及其证明比较
（其中也明确写出了余圈条件）。
该构造类似于《下降》第 \ref{descent-section-descent-modules} 节
讨论的构造；我们借用其中的记号
$\tau^n_i : [0] \to [n]$, $0 \mapsto i$ 以及
$\tau^n_{ij} : [1] \to [n]$, $0 \mapsto i$, $1 \mapsto j$。
给定 $\varphi : [n] \to [m]$，用下图定义
$K_\varphi : f_\varphi^{-1}K_n \to K_m$：
$$
\xymatrix{
f_\varphi^{-1}K_n \ar@{=}[r] &
f_\varphi^{-1} f_{\tau^n_n}^{-1}K_0 \ar@{=}[r] &
f_{\tau^m_{\varphi(n)}}^{-1}K_0 \ar@{=}[r] &
f_{\tau^m_{\varphi(n)m}}^{-1}f_{\delta^1_1}^{-1}K_0
\ar[d]_{f_{\tau^m_{\varphi(n)m}}^{-1}\alpha} \\
&
K_m \ar@{=}[r] &
f_{\tau^m_m}^{-1}K_0 \ar@{=}[r] &
f_{\tau^m_{\varphi(n)m}}^{-1}f_{\delta^1_0}^{-1}K_0
}
$$
我们略去如下验证：余圈条件蕴含这些映射
（在各自的导出范畴中）能够正确复合，因而给出
导出范畴的一个单纯系统。
\end{proof}

\begin{lemma}
\label{spaces-simplicial-lemma-abelian-postnikov}
在情形 \ref{spaces-simplicial-situation-simplicial-site} 中，设 $K$ 是
$D(\mathcal{C}_{total})$ 的对象。令
$$
X_n = (g_{n!}\mathbf{Z})
\otimes^\mathbf{L}_\mathbf{Z} K
\quad\text{且}\quad
Y_n =
(g_{n!}\mathbf{Z} \to \ldots \to g_{0!}\mathbf{Z})[-n]
\otimes^\mathbf{L}_\mathbf{Z} K
$$
为 $D(\mathcal{C}_{total})$ 的对象，其中的映射如引理
\ref{spaces-simplicial-lemma-simplicial-resolution-Z-site} 所述。
配备显然的典范映射 $Y_n \to X_n$ 以及
$Y_0 \to Y_1[1] \to Y_2[2] \to \ldots$ 后，有：
\begin{enumerate}
\item 区别三角 $Y_n \to X_n \to Y_{n - 1} \to Y_n[1]$
为 $\ldots \to X_2 \to X_1 \to X_0$ 定义一个 Postnikov 系统
（《导出范畴》定义 \ref{derived-definition-postnikov-system}）；
\item 在 $D(\mathcal{C}_{total})$ 中，
$K = \text{hocolim} Y_n[n]$。
\end{enumerate}
\end{lemma}

\begin{proof}
首先，若 $K = \mathbf{Z}$，则这正是将《导出范畴》例
\ref{derived-example-key-postnikov} 的构造应用于
$\textit{Ab}(\mathcal{C}_{total})$ 中的复形
$$
\ldots \to
g_{2!}\mathbf{Z} \to
g_{1!}\mathbf{Z} \to
g_{0!}\mathbf{Z}
$$
，再结合如下事实：由引理 \ref{spaces-simplicial-lemma-simplicial-resolution-Z-site}，
此复形在 $D(\mathcal{C}_{total})$ 中表示 $K = \mathbf{Z}$。
一般情形由此以及下列两个事实得出：正合函子
$- \otimes^\mathbf{L}_\mathbf{Z} K$ 将 Postnikov 系统映为
Postnikov 系统；函子
$- \otimes^\mathbf{L}_\mathbf{Z} K$ 与同伦余极限交换。
\end{proof}

\begin{lemma}
\label{spaces-simplicial-lemma-nullity-cartesian-objects-derived}
在情形 \ref{spaces-simplicial-situation-simplicial-site} 中。
% P10-E0187
设 $K, K' \in D(\mathcal{C}_{total})$。假设
\begin{enumerate}
\item $K$ 是笛卡尔的；
\item 当 $i > 0$ 时，$\Hom(K_i[i], K'_i) = 0$；
\item 当 $i \geq 0$ 时，$\Hom(K_i[i + 1], K'_i) = 0$。
\end{enumerate}
则任何诱导零映射 $K_0 \to K'_0$ 的映射 $K \to K'$ 都是零映射。
\end{lemma}

\begin{proof}
考虑引理 \ref{spaces-simplicial-lemma-abelian-postnikov} 中与 $K$ 相关联的对象
$X_n$ 及 Postnikov 系统 $Y_n$。
由于 $K = \text{hocolim} Y_n[n]$，映射 $K \to K'$
诱导相容的态射族 $Y_n[n] \to K'$。
由 (1) 与引理 \ref{spaces-simplicial-lemma-derived-cartesian-shriek}，有
$X_n = g_{n!}K_n$。由于 $Y_0 = X_0$，可见
$K_0 \to K'_0$ 为零蕴含 $Y_0 \to K'$ 为零。
假设已经证明对某个 $n \geq 0$，映射 $Y_n[n] \to K'$ 为零。
由区别三角
$$
Y_n[n] \to Y_{n + 1}[n + 1] \to X_{n + 1}[n + 1] \to Y_n[n + 1]
$$
得到正合列
$$
\Hom(X_{n + 1}[n + 1], K') \to
\Hom(Y_{n + 1}[n + 1], K') \to
\Hom(Y_n[n], K')
$$
由于 $X_{n + 1}[n + 1] = g_{n + 1!}K_{n + 1}[n + 1]$，
第一个群等于
$$
\Hom(K_{n + 1}[n + 1], K'_{n + 1})
$$
，它由假设 (2) 为零。归纳可知所有映射
$Y_n[n] \to K'$ 都为零。考虑同伦余极限的定义区别三角
$$
\bigoplus Y_n[n] \to
\bigoplus Y_n[n] \to
K \to
(\bigoplus Y_n[n])[1]
$$
。如上论证，可见只需证明
$$
\Hom((\bigoplus Y_n[n])[1], K') = \prod \Hom(Y_n[n + 1], K')
$$
对所有 $n \geq 0$ 均为零。为此，再次如上论证，只需证明
$$
\Hom(K_n[n + 1], K'_n)  = 0
$$
对所有 $n \geq 0$ 成立；这由条件 (3) 得出。
\end{proof}

\begin{lemma}
\label{spaces-simplicial-lemma-hom-cartesian-objects-derived}
在情形 \ref{spaces-simplicial-situation-simplicial-site} 中。
% P10-E0188
设 $K, K' \in D(\mathcal{C}_{total})$。假设
\begin{enumerate}
\item $K$ 是笛卡尔的；
\item 当 $i > 1$ 时，$\Hom(K_i[i - 1], K'_i) = 0$。
\end{enumerate}
则 $K$ 与 $K'$ 所关联的单纯系统之间的任何映射
$\{K_n \to K'_n\}$，均来自 $D(\mathcal{C}_{total})$
中的一个映射 $K \to K'$。
\end{lemma}

\begin{proof}
设 $\{K_n \to K'_n\}_{n \geq 0}$ 是导出范畴单纯系统的一个态射。
考虑引理 \ref{spaces-simplicial-lemma-abelian-postnikov} 中与 $K$ 相关联的对象
$X_n$ 及 Postnikov 系统 $Y_n$。
由 (1) 与引理 \ref{spaces-simplicial-lemma-derived-cartesian-shriek}，有
$X_n = g_{n!}K_n$。特别地，映射 $K_0 \to K'_0$
诱导态射 $X_0 \to K'$。由于 $\{K_n \to K'_n\}$
是系统的态射，一个从略的计算表明复合
$$
X_1 \to X_0 \to K'
$$
为零。由于 $Y_0 = X_0$，且 $Y_1$ 位于区别三角
$$
Y_1 \to X_1 \to Y_0 \to Y_1[1]
$$
中，可知存在态射 $Y_1[1] \to K'$，它与
$X_0 = Y_0 \to Y_1[1]$ 的复合正是上述态射 $X_0 \to K'$。
假设已对 $n \geq 1$ 给定映射 $Y_n[n] \to K'$。
由区别三角
$$
X_{n + 1}[n] \to Y_n[n] \to Y_{n + 1}[n + 1] \to X_{n + 1}[n + 1]
$$
得到正合列
$$
\Hom(Y_{n + 1}[n + 1], K') \to
\Hom(Y_n[n], K') \to
\Hom(X_{n + 1}[n], K')
$$
由于 $X_{n + 1}[n] = g_{n + 1!}K_{n + 1}[n]$，最后一个群等于
$$
\Hom(K_{n + 1}[n], K'_{n + 1})
$$
，它由假设 (2) 为零。归纳得到与转移映射相容的映射系统
$Y_n[n] \to K'$，且它在 $Y_0$ 上约化为给定映射。
这产生映射
$$
\gamma :
K = \text{hocolim} Y_n[n]
\longrightarrow
K'
$$
无论如何，此映射使图
$$
\xymatrix{
X_0 \ar[rd] \ar[r] &
K \ar[d]^\gamma \\
& K'
}
$$
交换。限制到 $\mathcal{C}_0$，可知映射
$\gamma_0 : K_0 \to K'_0$ 与该单纯系统态射的第一个映射
$K_0 \to K'_0$ 相同。由于 $K$ 是笛卡尔的，由此容易看出
$\{\gamma_n\}$ 正是我们最初给定的单纯系统态射。
\end{proof}

\begin{lemma}
\label{spaces-simplicial-lemma-cartesian-object-derived-from-simplicial}
在情形 \ref{spaces-simplicial-situation-simplicial-site} 中，设
$(K_n, K_\varphi)$ 是导出范畴的一个单纯系统。假设
\begin{enumerate}
\item $(K_n, K_\varphi)$ 是笛卡尔的；
\item 当 $i \geq 0$ 且 $t > 0$ 时，
$\Hom(K_i[t], K_i) = 0$。
\end{enumerate}
则存在 $D(\mathcal{C}_{total})$ 的一个笛卡尔对象 $K$，
其所关联的单纯系统同构于 $(K_n, K_\varphi)$。
\end{lemma}

\begin{proof}
在 $D(\mathcal{C}_{total})$ 中令 $X_n = g_{n!}K_n$。
对每个 $n \geq 1$，有
$$
\Hom(X_n, X_{n - 1}) =
\Hom(K_n, g_n^{-1}g_{n - 1!}K_{n - 1}) =
\bigoplus\nolimits_{\varphi : [n - 1] \to [n]}
\Hom(K_n, f_\varphi^{-1}K_{n - 1})
$$
因此，我们得到映射 $X_n \to X_{n - 1}$，它对应于映射
$K_\varphi^{-1} : K_n \to f_\varphi^{-1}K_{n - 1}$
在 $\varphi$ 遍历 $\delta^n_0, \ldots, \delta^n_n$ 时的交错和。
这可以进行，因为由假设 (1)，$K_\varphi$ 可逆。
请注意，这与引理 \ref{spaces-simplicial-lemma-simplicial-resolution-Z-site}
证明中映射的定义相似。由此得到
$D(\mathcal{C}_{total})$ 中的复形
$$
\ldots \to X_2 \to X_1 \to X_0
$$
。我们略去说明各复合为零的计算。由《导出范畴》引理
\ref{derived-lemma-existence-postnikov-system}，若有
$$
\Hom(X_i[i - j - 2], X_j) = 0\text{，当 }i > j + 2
$$
，则可将此复形延拓为一个 Postnikov 系统。该群等于
$$
\Hom(K_i[i - j - 2], g_i^{-1}g_{j!}K_j)
$$
再次利用 $(K_n, K_\varphi)$ 是笛卡尔的，可见
$g_i^{-1}g_{j!}K_j$ 同构于有限个 $K_i$ 副本的直和。
因此，该群由假设 (2) 消失。
令 Postnikov 系统由 $Y_0 = X_0$ 以及对 $n \geq 1$
的区别三角 $Y_n \to X_n \to Y_{n - 1} \to Y_n[1]$ 给出。
令
$$
K = \text{hocolim} Y_n[n]
$$
为完成证明，须证明对所有 $m$，$g_m^{-1}K$ 同构于 $K_m$，
且这些同构与映射 $K_\varphi$ 相容。注意
$$
g_m^{-1} K = \text{hocolim} g_m^{-1}Y_n[n]
$$
，且 $g_m^{-1}Y_n[n]$ 是 $g_m^{-1}X_n$ 的一个 Postnikov 系统。
考虑同构
$$
g_m^{-1}X_n =
\bigoplus\nolimits_{\varphi : [n] \to [m]} f_\varphi^{-1}K_n
\xrightarrow{\bigoplus K_\varphi}
\bigoplus\nolimits_{\varphi : [n] \to [m]} K_m
$$
这些映射定义复形的同构
$$
\xymatrix{
\ldots \ar[r] &
g_m^{-1}X_2 \ar[r] \ar[d] &
g_m^{-1}X_1 \ar[r] \ar[d] &
g_m^{-1}X_0 \ar[d] \\
\ldots \ar[r] &
\bigoplus\nolimits_{\varphi : [2] \to [m]} K_m \ar[r] &
\bigoplus\nolimits_{\varphi : [1] \to [m]} K_m \ar[r] &
\bigoplus\nolimits_{\varphi : [0] \to [m]} K_m
}
$$
；它位于 $D(\mathcal{C}_m)$ 中，而下行的箭头如引理
\ref{spaces-simplicial-lemma-simplicial-resolution-Z-site} 的证明所述。
由于我们对复形 $\ldots \to X_2 \to X_1 \to X_0$
中箭头的选取，各方块交换；计算从略。
下行复形有由下式给出的 Postnikov 塔：
$$
Y'_{m, n} =
\left(\bigoplus\nolimits_{\varphi : [n] \to [m]} \mathbf{Z} \to
\ldots \to
\bigoplus\nolimits_{\varphi : [0] \to [m]} \mathbf{Z}\right)[-n]
\otimes^\mathbf{L}_\mathbf{Z} K_m
$$
，并且 $\text{hocolim} Y'_{m, n} = K_m$
（请与引理 \ref{spaces-simplicial-lemma-abelian-postnikov} 的证明以及《导出范畴》例
\ref{derived-example-key-postnikov} 比较）。
应用《导出范畴》引理
\ref{derived-lemma-existence-postnikov-system} 的第二部分，
只要有
$$
\Hom(g_m^{-1}X_i[i - j - 1], \bigoplus\nolimits_{\varphi : [j] \to [m]} K_m)
= 0\text{，当 }i > j + 1
$$
，大图中的竖直映射便可延拓为 Postnikov 系统的同构。
% P10-E0189
若当 $i > j + 1$ 时
$\Hom(K_m[i - j - 1], K_m) = 0$，上述消失成立；
而这由假设 (2) 得到。在 $D(\mathcal{C}_m)$ 中，选取由
$\gamma_{m, n} : g_m^{-1}Y_n \to Y'_{m, n}$ 给出的
Postnikov 系统同构。由同伦余极限的唯一性，可以找到同构
$$
g_m^{-1} K = \text{hocolim} g_m^{-1}Y_n[n]
\xrightarrow{\gamma_m}
\text{hocolim} Y'_{m, n} = K_m
$$
，它与 $\gamma_{m, n}$ 相容。

\medskip\noindent
还须证明，对每个 $\varphi : [m] \to [n]$，映射 $\gamma_m$
构成交换图
$$
\xymatrix{
f_\varphi^{-1}g_m^{-1}K \ar[d]_{f_\varphi^{-1}\gamma_m} \ar[r]_{K(\varphi)} &
g_n^{-1}K \ar[d]^{\gamma_n} \\
f_\varphi^{-1}K_m \ar[r]^{K_\varphi} &
K_n
}
$$
。考虑图
$$
\xymatrix{
f_\varphi^{-1}(\bigoplus_{\psi : [0] \to [m]} f_\psi^{-1}K_0)
\ar@{=}[r] \ar[d]_{f_\varphi^{-1}(\bigoplus K_\psi)} &
f_\varphi^{-1}g_m^{-1}X_0 \ar[d] \ar[r]_{X_0(\varphi)} &
g_n^{-1}X_0 \ar[d] &
\bigoplus_{\chi : [0] \to [n]} f_\chi^{-1}K_0
\ar@{=}[l] \ar[d]^{\bigoplus K_\chi} \\
f_\varphi^{-1}(\bigoplus_{\psi : [0] \to [m]} K_m) \ar@{=}[d] &
f_\varphi^{-1}g_m^{-1}K \ar[d]_{f_\varphi^{-1}\gamma_m} \ar[r]_{K(\varphi)} &
g_n^{-1}K \ar[d]^{\gamma_n} &
\bigoplus_{\chi : [0] \to [n]} K_n \ar@{=}[d] \\
f_\varphi^{-1}Y'_{0, m} \ar[r] &
f_\varphi^{-1}K_m \ar[r]^{K_\varphi} &
K_n &
Y'_{0, n} \ar[l]
}
$$
由于 $X_0 \to K$ 是单纯对象的态射，上方中间的方块交换。
由于 $\gamma_m$（相应地，$\gamma_n$）与
$\gamma_{0, m}$（相应地，$\gamma_{0, n}$）相容，
左矩形（相应地，右矩形）交换；后两者正是图中的箭头
$\bigoplus K_\psi$ 与 $\bigoplus K_\chi$。
% P10-E0190
由于 $(K_n, K_\varphi)$ 是单纯系统，且映射 $X_0(\varphi)$
由显然的恒等式
$f_\varphi^{-1}f_\psi^{-1}K_0 = f_{\varphi \circ \psi}^{-1}K_0$
给出，故沿该图的外矩形绕行是交换的。
注意，箭头 $\bigoplus_\psi K_m \to Y'_{0, m} \to K_m$
在任意一个直和分量上诱导同构
（这是因为上面对 Postnikov 系统 $Y'_{i, j}$ 的显式构造）。
% P10-E0191
因此，若取左上角的一个直和分量，则由于 $\gamma_m$ 是同构，
该分量同构地映到 $f_\varphi^{-1}g_m^{-1}K$。
展开上述陈述，但只考察这个直和分量，即可推出下方中间的方块也交换。
证明完毕。
\end{proof}









\section{导出范畴的单纯系统：模}
\label{spaces-simplicial-section-glueing-modules}

\noindent
本节将在 $\mathcal{O}$-模导出范畴的框架下证明
\cite[Proposition 3.2.9]{BBD} 的一个特例。
（稍微）容易一些的阿贝尔层情形在第
\ref{spaces-simplicial-section-glueing} 节中讨论。

\begin{definition}
\label{spaces-simplicial-definition-cartesian-derived-modules}
在情形 \ref{spaces-simplicial-situation-simplicial-site} 中，设 $\mathcal{O}$ 是
$\mathcal{C}_{total}$ 上的环层。一个
{\it 模导出范畴的单纯系统}由下列数据组成：
\begin{enumerate}
\item 对每个 $n$，给定 $D(\mathcal{O}_n)$ 的对象 $K_n$；
\item 对每个 $\varphi : [m] \to [n]$，给定
$D(\mathcal{O}_n)$ 中的映射
$K_\varphi : Lf_\varphi^*K_m \to K_n$。
\end{enumerate}
这些数据须满足：
$$
K_{\varphi \circ \psi} = K_\varphi \circ Lf_\varphi^*K_\psi :
Lf_{\varphi \circ \psi}^*K_l = Lf_\varphi^* Lf_\psi^*K_l
\longrightarrow
K_n
$$
对 $\Delta$ 的任意态射 $\varphi : [m] \to [n]$ 与
$\psi : [l] \to [m]$，上述等式均成立。
若对所有 $\varphi$，映射 $K_\varphi$ 都是同构，
则称该单纯系统是{\it 笛卡尔的}。
对两个导出范畴的单纯系统，显然可以定义
{\it 模导出范畴的单纯系统的态射}。
\end{definition}

\noindent
我们有意给这个概念取了一个长得荒唐的名字。
目标是证明：在某些假设下，模导出范畴的单纯系统来自
$D(\mathcal{O})$ 的一个对象。

\begin{lemma}
\label{spaces-simplicial-lemma-cartesian-objects-derived-modules}
在情形 \ref{spaces-simplicial-situation-simplicial-site} 中，设 $\mathcal{O}$ 是
$\mathcal{C}_{total}$ 上的环层。
若 $K \in D(\mathcal{O})$，则 $(K_n, K(\varphi))$
是模导出范畴的一个单纯系统。
若 $K$ 是笛卡尔的，则该系统也是笛卡尔的。
\end{lemma}

\begin{proof}
这立即由定义得出。
\end{proof}

\begin{lemma}
\label{spaces-simplicial-lemma-construct-cartesian-objects-derived-modules}
在情形 \ref{spaces-simplicial-situation-simplicial-site} 中，设 $\mathcal{O}$ 是
$\mathcal{C}_{total}$ 上的环层。假设给定
$K_0 \in D(\mathcal{O}_0)$ 以及同构
$$
\alpha :
L(f_{\delta_1^1})^*K_0
\longrightarrow
L(f_{\delta_0^1})^*K_0
$$
，它满足余圈条件。令 $\tau^n_i : [0] \to [n]$, $0 \mapsto i$，
并令 $K_n = Lf_{\tau^n_n}^*K_0$。诸对象 $K_n$
构成模导出范畴的一个笛卡尔单纯系统的成员。
\end{lemma}

\begin{proof}
请与引理 \ref{spaces-simplicial-lemma-construct-cartesian-objects-derived}、
引理 \ref{spaces-simplicial-lemma-characterize-cartesian} 及其证明比较
（其中也明确写出了余圈条件）。
该构造类似于《下降》第 \ref{descent-section-descent-modules} 节
讨论的构造；我们借用其中的记号
$\tau^n_i : [0] \to [n]$, $0 \mapsto i$ 以及
$\tau^n_{ij} : [1] \to [n]$, $0 \mapsto i$, $1 \mapsto j$。
给定 $\varphi : [n] \to [m]$，用下图定义
$K_\varphi : L(f_\varphi)^*K_n \to K_m$：
$$
\xymatrix{
L(f_\varphi)^*K_n \ar@{=}[r] &
L(f_\varphi)^* L(f_{\tau^n_n})^*K_0 \ar@{=}[r] &
L(f_{\tau^m_{\varphi(n)}})^*K_0 \ar@{=}[r] &
L(f_{\tau^m_{\varphi(n)m}})^* L(f_{\delta^1_1})^*K_0
\ar[d]_{L(f_{\tau^m_{\varphi(n)m}})^*\alpha} \\
&
K_m \ar@{=}[r] &
L(f_{\tau^m_m})^*K_0 \ar@{=}[r] &
L(f_{\tau^m_{\varphi(n)m}})^* L(f_{\delta^1_0})^*K_0
}
$$
% P10-E0192
我们略去如下验证：余圈条件蕴含这些映射
（在各自的导出范畴中）能够正确复合，因而给出
模导出范畴的一个单纯系统。
\end{proof}

\begin{lemma}
\label{spaces-simplicial-lemma-modules-postnikov}
在情形 \ref{spaces-simplicial-situation-simplicial-site} 中，设 $\mathcal{O}$ 是
$\mathcal{C}_{total}$ 上的环层。
% P10-E0193
设 $K$ 是 $D(\mathcal{C}_{total})$ 的对象。令
$$
X_n = (g_{n!}\mathcal{O}_n)
\otimes^\mathbf{L}_\mathcal{O} K
\quad\text{且}\quad
Y_n =
(g_{n!}\mathcal{O}_n \to \ldots \to g_{0!}\mathcal{O}_0)[-n]
\otimes^\mathbf{L}_\mathcal{O} K
$$
为 $D(\mathcal{O})$ 的对象，其中的映射如
% P10-E0194
引理 \ref{spaces-simplicial-lemma-simplicial-resolution-Z-site} 所述。
配备显然的典范映射 $Y_n \to X_n$ 以及
$Y_0 \to Y_1[1] \to Y_2[2] \to \ldots$ 后，有：
\begin{enumerate}
\item 区别三角 $Y_n \to X_n \to Y_{n - 1} \to Y_n[1]$
为 $\ldots \to X_2 \to X_1 \to X_0$ 定义一个 Postnikov 系统
（《导出范畴》定义 \ref{derived-definition-postnikov-system}）；
\item 在 $D(\mathcal{O})$ 中，$K = \text{hocolim} Y_n[n]$。
\end{enumerate}
\end{lemma}

\begin{proof}
首先，若 $K = \mathcal{O}$，则这正是将《导出范畴》例
\ref{derived-example-key-postnikov} 的构造应用于
$\textit{Ab}(\mathcal{C}_{total})$ 中的复形
$$
\ldots \to
g_{2!}\mathcal{O}_2 \to
g_{1!}\mathcal{O}_1 \to
g_{0!}\mathcal{O}_0
$$
，再结合如下事实：由引理 \ref{spaces-simplicial-lemma-simplicial-resolution-ringed}，
此复形在 $D(\mathcal{C}_{total})$ 中表示 $K = \mathcal{O}$。
一般情形由此以及下列两个事实得出：正合函子
$- \otimes^\mathbf{L}_\mathcal{O} K$ 将 Postnikov 系统映为
Postnikov 系统；函子
$- \otimes^\mathbf{L}_\mathcal{O} K$ 与同伦余极限交换。
\end{proof}

\begin{lemma}
\label{spaces-simplicial-lemma-nullity-cartesian-modules-derived}
在情形 \ref{spaces-simplicial-situation-simplicial-site} 中，设 $\mathcal{O}$ 是
$\mathcal{C}_{total}$ 上的环层。
% P10-E0195
设 $K, K' \in D(\mathcal{O})$。假设
\begin{enumerate}
% P10-E0196
\item 对 $\varphi : [m] \to [n]$，
$f_\varphi^{-1}\mathcal{O}_n \to \mathcal{O}_m$ 平坦；
\item $K$ 是笛卡尔的；
\item 当 $i > 0$ 时，$\Hom(K_i[i], K'_i) = 0$；
\item 当 $i \geq 0$ 时，$\Hom(K_i[i + 1], K'_i) = 0$。
\end{enumerate}
则任何诱导零映射 $K_0 \to K'_0$ 的映射 $K \to K'$ 都是零映射。
\end{lemma}

\begin{proof}
证明与引理 \ref{spaces-simplicial-lemma-nullity-cartesian-objects-derived} 的证明完全相同，
只需用引理 \ref{spaces-simplicial-lemma-modules-postnikov} 代替
引理 \ref{spaces-simplicial-lemma-abelian-postnikov}。
\end{proof}

\begin{lemma}
\label{spaces-simplicial-lemma-hom-cartesian-modules-derived}
在情形 \ref{spaces-simplicial-situation-simplicial-site} 中，设 $\mathcal{O}$ 是
$\mathcal{C}_{total}$ 上的环层。
% P10-E0197
设 $K, K' \in D(\mathcal{O})$。假设
\begin{enumerate}
% P10-E0196
\item 对 $\varphi : [m] \to [n]$，
$f_\varphi^{-1}\mathcal{O}_n \to \mathcal{O}_m$ 平坦；
\item $K$ 是笛卡尔的；
\item 当 $i > 1$ 时，$\Hom(K_i[i - 1], K'_i) = 0$。
\end{enumerate}
则 $K$ 与 $K'$ 所关联的单纯系统之间的任何映射
$\{K_n \to K'_n\}$，均来自 $D(\mathcal{O})$ 中的一个映射
$K \to K'$。
\end{lemma}

\begin{proof}
证明与引理 \ref{spaces-simplicial-lemma-hom-cartesian-objects-derived} 的证明完全相同，
只需用引理 \ref{spaces-simplicial-lemma-modules-postnikov} 代替
引理 \ref{spaces-simplicial-lemma-abelian-postnikov}。
\end{proof}

\begin{lemma}
\label{spaces-simplicial-lemma-cartesian-module-derived-from-simplicial}
在情形 \ref{spaces-simplicial-situation-simplicial-site} 中，设 $\mathcal{O}$ 是
$\mathcal{C}_{total}$ 上的环层，且
$(K_n, K_\varphi)$ 是模导出范畴的一个单纯系统。假设
\begin{enumerate}
% P10-E0196
\item 对 $\varphi : [m] \to [n]$，
$f_\varphi^{-1}\mathcal{O}_n \to \mathcal{O}_m$ 平坦；
\item $(K_n, K_\varphi)$ 是笛卡尔的；
\item 当 $i \geq 0$ 且 $t > 0$ 时，
$\Hom(K_i[t], K_i) = 0$。
\end{enumerate}
则存在 $D(\mathcal{O})$ 的一个笛卡尔对象 $K$，
其所关联的单纯系统同构于 $(K_n, K_\varphi)$。
\end{lemma}

\begin{proof}
证明与引理 \ref{spaces-simplicial-lemma-cartesian-object-derived-from-simplicial}
的证明完全相同，但作如下改动：
\begin{enumerate}
\item 处处使用 $g_n^* = Lg_n^*$ 代替 $g_n^{-1}$；
\item 处处使用 $f_\varphi^* = Lf_\varphi^*$ 代替 $f_\varphi^{-1}$；
\item 引用引理 \ref{spaces-simplicial-lemma-simplicial-resolution-ringed}，
而非引理 \ref{spaces-simplicial-lemma-simplicial-resolution-Z-site}；
\item 在 $Y'_{m, n}$ 的构造中使用 $\mathcal{O}_m$，
而非 $\mathbf{Z}$；
\item 与引理 \ref{spaces-simplicial-lemma-modules-postnikov} 的证明比较，
而非与引理 \ref{spaces-simplicial-lemma-abelian-postnikov} 的证明比较。
\end{enumerate}
证明完毕。
\end{proof}







\section{与半可表示对象关联的站点}
\label{spaces-simplicial-section-semi-representable}

\noindent
设 $\mathcal{C}$ 是一个站点。回忆，$\mathcal{C}$ 的一个
{\it 半可表示对象}，就是 $\mathcal{C}$ 的对象族
$\{U_i\}_{i \in I}$。一个
{\it 半可表示对象的态射
$\{U_i\}_{i \in I} \to \{V_j\}_{j \in J}$}
由映射 $\alpha : I \to J$ 以及对每个 $i \in I$
给定的 $\mathcal{C}$ 中的态射
$f_i : U_i \to V_{\alpha(i)}$ 组成。
$\mathcal{C}$ 的半可表示对象范畴记为
$\text{SR}(\mathcal{C})$。更多信息参见《超覆盖》定义
\ref{hypercovering-definition-SR} 及其所在一节。

\medskip\noindent
对 $\mathcal{C}$ 的半可表示对象 $K = \{U_i\}_{i \in I}$，令
$$
\mathcal{C}/K = \coprod\nolimits_{i \in I} \mathcal{C}/U_i
$$
为 $\mathcal{C}$ 在诸 $U_i$ 处的局部化之不交并。
此范畴上有自然的站点结构，其覆盖继承自各局部化
$\mathcal{C}/U_i$。站点 $\mathcal{C}/K$ 称为
{\it $\mathcal{C}$ 在 $K$ 处的局部化}。
注意，$\mathcal{C}/K$ 上的层等同于各
$\mathcal{C}/U_i$ 上的层 $\mathcal{F}_i$ 所成的族，即
$$
\Sh(\mathcal{C}/K) = \prod\nolimits_{i \in I} \Sh(\mathcal{C}/U_i)
$$
这一描述有时有助于理解所发生的事情。

\medskip\noindent
设 $\mathcal{C}$ 是一个站点，且
$K = \{U_i\}_{i \in I}$ 是 $\text{SR}(\mathcal{C})$ 的对象。
存在连续且余连续的局部化函子
$j : \mathcal{C}/K \to \mathcal{C}$；它是诸局部化函子的积
% P10-E0198
$j_i : \mathcal{C}/V_i \to \mathcal{C}$.
正如《站点》第 \ref{sites-section-localize} 节中一样，
我们得到函子 $j_!$、$j^{-1}$、$j_*$。
用积分解
$\Sh(\mathcal{C}/K) = \prod\nolimits_{i \in I} \Sh(\mathcal{C}/U_i)$
表示时，有
$$
\begin{matrix}
j_! & : &
(\mathcal{F}_i)_{i \in I} &
\longmapsto &
\coprod j_{i, !}\mathcal{F}_i \\
j^{-1} & : &
\mathcal{G} &
\longmapsto &
(j_i^{-1}\mathcal{G})_{i \in I} \\
j_* & : &
(\mathcal{F}_i)_{i \in I} &
\longmapsto &
\prod j_{i, *}\mathcal{F}_i
\end{matrix}
$$
，读者容易验证这些公式。

\medskip\noindent
设 $f : K \to L$ 是 $\text{SR}(\mathcal{C})$ 的一个态射。
则逐分量应用《站点》引理 \ref{sites-lemma-relocalize}
的构造，得到连续且余连续的函子
$$
v : \mathcal{C}/K \longrightarrow \mathcal{C}/L
$$
。更准确地，假设 $f = (\alpha, f_i)$，其中
$K = \{U_i\}_{i \in I}$、$L = \{V_j\}_{j \in J}$、
$\alpha : I \to J$，且 $f_i : U_i \to V_{\alpha(i)}$。
则函子 $v$ 通过上述引理的构造，将分量
$\mathcal{C}/U_i$ 映入分量 $\mathcal{C}/V_{\alpha(i)}$。
特别地，得到拓扑斯的态射
$$
f : \Sh(\mathcal{C}/K) \to \Sh(\mathcal{C}/L)
$$
。用积分解
$\Sh(\mathcal{C}/K) = \prod\nolimits_{i \in I} \Sh(\mathcal{C}/U_i)$ and
$\Sh(\mathcal{C}/L) = \prod\nolimits_{j \in J} \Sh(\mathcal{C}/V_j)$
表示时，读者可验证
$$
\begin{matrix}
f_! & : &
(\mathcal{F}_i)_{i \in I} &
\longmapsto &
(\coprod\nolimits_{i \in I, \alpha(i) = j} f_{i, !}\mathcal{F}_i)_{j \in J} \\
f^{-1} & : &
(\mathcal{G}_j)_{j \in J} &
\longmapsto &
(f_i^{-1}\mathcal{G}_{\alpha(i)})_{i \in I} \\
f_* & : &
(\mathcal{F}_i)_{i \in I} &
\longmapsto &
(\prod\nolimits_{i \in I, \alpha(i) = j} f_{i, *}\mathcal{F}_i)_{j \in J}
\end{matrix}
$$
，其中 $f_i : \Sh(\mathcal{C}/U_i) \to
\Sh(\mathcal{C}/V_{\alpha(i)})$ 是与
$f_i : U_i \to V_{\alpha(i)}$ 所对应的局部化函子
$\mathcal{C}/U_i \to \mathcal{C}/V_{\alpha(i)}$ 相关联的态射。

\begin{lemma}
\label{spaces-simplicial-lemma-has-P}
设 $\mathcal{C}$ 是一个站点。
\begin{enumerate}
\item 对 $\text{SR}(\mathcal{C})$ 中的 $K$，函子
$j : \mathcal{C}/K \to \mathcal{C}$ 连续、余连续，并具有
《站点》注 \ref{sites-remark-cartesian-cocontinuous} 的性质 P。
\item 对 $\text{SR}(\mathcal{C})$ 中的 $f : K \to L$，
上述函子 $v : \mathcal{C}/K \to \mathcal{C}/L$
连续、余连续，并具有《站点》注
\ref{sites-remark-cartesian-cocontinuous} 的性质 P。
\end{enumerate}
\end{lemma}

\begin{proof}
证明 (2)。采用引理前讨论中的记号。由《站点》第
\ref{sites-section-localize} 节，局部化函子
$\mathcal{C}/U_i \to \mathcal{C}/V_{\alpha(i)}$
连续且余连续；由《站点》注
\ref{sites-remark-localization-cartesian-cocontinuous}，它们满足 $P$。
由此形式地推出 $v$ 连续、余连续并具有 $P$。细节从略。
(1) 的证明也从略。
\end{proof}

\begin{lemma}
\label{spaces-simplicial-lemma-push-pull-localization}
设 $\mathcal{C}$ 是一个站点，且
$K$ 属于 $\text{SR}(\mathcal{C})$。对
$\Sh(\mathcal{C})$ 中的 $\mathcal{F}$，有
$$
j_*j^{-1}\mathcal{F} = \SheafHom(F(K)^\#, \mathcal{F})
$$
，其中 $F$ 如《超覆盖》定义
\ref{hypercovering-definition-SR-F} 所述。
\end{lemma}

\begin{proof}
设 $K = \{U_i\}_{i \in I}$。
利用上面对函子 $j^{-1}$ 与 $j_*$ 的描述，可见
$$
j_*j^{-1}\mathcal{F} =
\prod\nolimits_{i \in I} j_{i, *}(\mathcal{F}|_{\mathcal{C}/U_i}) =
\prod\nolimits_{i \in I} \SheafHom(h_{U_i}^\#, \mathcal{F})
$$
% P10-E0199
第二个等式由《站点》引理 \ref{sites-lemma-hom-sheaf-hU} 得出。
% P10-E0200
由于在 $\textit{PSh}(\mathcal{C}$ 中
$F(K) = \coprod h_{U_i}$，故在 $\Sh(\mathcal{C})$ 中
$F(K)^\# = \coprod h_{U_i}^\#$。又因
$\SheafHom(-, \mathcal{F})$ 将余积变为积
（这立即由《站点》第 \ref{sites-section-glueing-sheaves} 节
的构造得出），结论成立。
\end{proof}

\begin{lemma}
\label{spaces-simplicial-lemma-localize-compare}
设 $\mathcal{C}$ 是一个站点。
\begin{enumerate}
\item 对 $\text{SR}(\mathcal{C})$ 中的 $K$，函子 $j_!$ 给出等价
$\Sh(\mathcal{C}/K) \to \Sh(\mathcal{C})/F(K)^\#$，
其中 $F$ 如《超覆盖》定义
\ref{hypercovering-definition-SR-F} 所述。
\item 经由 (1) 中的等同，函子
$j^{-1} : \Sh(\mathcal{C}) \to \Sh(\mathcal{C}/K)$ 对应于
$\mathcal{F} \mapsto (\mathcal{F} \times F(K)^\# \to F(K)^\#)$.
\item 对 $\text{SR}(\mathcal{C})$ 中的 $f : K \to L$，
经由 (1) 中的等同，函子 $f^{-1}$ 对应于函子
$\Sh(\mathcal{C})/F(L)^\# \to \Sh(\mathcal{C})/F(K)^\#$,
$(\mathcal{G} \to F(L)^\#) \mapsto
(\mathcal{G} \times_{F(L)^\#} F(K)^\# \to F(K)^\#)$.
\end{enumerate}
\end{lemma}

\begin{proof}
注意，若 $K = \{U_i\}_{i \in I}$，则范畴
$\Sh(\mathcal{C}/K)$ 分解为诸范畴
$\Sh(\mathcal{C}/U_i)$ 的积。又注意
$F(K)^\# = \coprod_{i \in I} h_{U_i}^\#$（层范畴中的余积）。
因此，$\Sh(\mathcal{C})/F(K)^\#$ 是诸范畴
$\Sh(\mathcal{C})/h_{U_i}^\#$ 的积。
于是，(1) 与 (2) 由每个 $i$ 对应的陈述得出；参见《站点》引理
\ref{sites-lemma-essential-image-j-shriek} 与
\ref{sites-lemma-compute-j-shriek-restrict}。
类似地，若 $L = \{V_j\}_{j \in J}$，且 $f$ 由
$\alpha : I \to J$ 及 $f_i : U_i \to V_{\alpha(i)}$ 给出，
则对每个再局部化态射
$\mathcal{C}/U_i \to \mathcal{C}/V_{\alpha(i)}$
应用《站点》引理 \ref{sites-lemma-relocalize-explicit}，即得 (3)。
\end{proof}

\begin{lemma}
\label{spaces-simplicial-lemma-localize-injective}
设 $\mathcal{C}$ 是一个站点。对
$K$ 属于 $\text{SR}(\mathcal{C})$，函子 $j^{-1}$ 将内射阿贝尔层
映为内射阿贝尔层。类似地，函子 $j^{-1}$ 将阿贝尔层的
K-内射复形映为阿贝尔层的 K-内射复形。
\end{lemma}

\begin{proof}
第一个陈述是《站点上的上同调》引理
\ref{sites-cohomology-lemma-cohomology-of-open}
向半可表示对象的自然推广。事实上，利用
$\Sh(\mathcal{C}/K)$ 的积分解以及上面对函子 $j^{-1}$ 的描述，
它由该引理得出。第二个陈述是《站点上的上同调》引理
\ref{sites-cohomology-lemma-restrict-K-injective-to-open}
的自然推广，并由拓扑斯的积分解从该引理得出。

\medskip\noindent
另一种证明：由引理 \ref{spaces-simplicial-lemma-localize-compare} 的 (1)，
$j$ 诱导拓扑斯的一个局部化，故通过扩大站点，结论也立即由
《站点上的上同调》引理
\ref{sites-cohomology-lemma-cohomology-of-open} 与
\ref{sites-cohomology-lemma-restrict-K-injective-to-open} 得出；
在内射层的情形下，可与《站点上的上同调》引理
\ref{sites-cohomology-lemma-cohomology-on-sheaf-sets} 的证明比较。
\end{proof}

\begin{remark}[对象之上的变体]
\label{spaces-simplicial-remark-semi-representable-over-object}
设 $\mathcal{C}$ 是一个站点，且 $X \in \Ob(\mathcal{C})$。
{\it $X$ 之上的半可表示对象}范畴
$\text{SR}(\mathcal{C}, X)$ 由公式
$\text{SR}(\mathcal{C}, X) = \text{SR}(\mathcal{C}/X)$ 定义。
参见《超覆盖》定义 \ref{hypercovering-definition-SR}。
因此，可以将上述讨论应用于站点 $\mathcal{C}/X$。
简言之，上述构造给出：
\begin{enumerate}
\item 对 $\text{SR}(\mathcal{C}, X)$ 中的 $K$，给出站点
$\mathcal{C}/K$；
\item 若 $K = \{U_i/X\}$，给出分解
$\Sh(\mathcal{C}/K) = \prod \Sh(\mathcal{C}/U_i)$；
\item 给出局部化函子 $j : \mathcal{C}/K \to \mathcal{C}/X$；
\item 对 $\text{SR}(\mathcal{C}, X)$ 中的 $f : K \to L$，
给出态射 $f : \Sh(\mathcal{C}/K) \to \Sh(\mathcal{C}/L)$。
\end{enumerate}
将 $\mathcal{C}$ 处处替换为 $\mathcal{C}/X$，本节所有结果
在此情形下均成立。
\end{remark}

\begin{remark}[赋环变体]
\label{spaces-simplicial-remark-semi-representable-ringed}
设 $\mathcal{C}$ 是一个站点，且 $\mathcal{O}_\mathcal{C}$
是 $\mathcal{C}$ 上的环层。在此情形下，对 $\mathcal{C}$
的任意半可表示对象 $K$，站点 $\mathcal{C}/K$ 是赋环站点，
其环层为 $\mathcal{O}_K = j^{-1}\mathcal{O}_\mathcal{C}$。
上述构造给出：
\begin{enumerate}
\item 对 $\text{SR}(\mathcal{C})$ 中的 $K$，给出赋环站点
$(\mathcal{C}/K, \mathcal{O}_K)$；
\item 若 $K = \{U_i\}$，给出分解
$\textit{Mod}(\mathcal{O}_K) =
\prod \textit{Mod}(\mathcal{O}_{U_i})$；
\item 给出赋环拓扑斯的局部化态射
$j : (\Sh(\mathcal{C}/K), \mathcal{O}_K) \to
(\Sh(\mathcal{C}), \mathcal{O}_\mathcal{C})$
；
\item 对 $\text{SR}(\mathcal{C})$ 中的 $f : K \to L$，
给出赋环拓扑斯的态射
$f : (\Sh(\mathcal{C}/K), \mathcal{O}_K) \to
(\Sh(\mathcal{C}/L), \mathcal{O}_L)$。
\end{enumerate}
上述许多结果在此框架下成立。例如，函子 $j^*$ 有正合左伴随
$$
j_! : \textit{Mod}(\mathcal{O}_K) \to \textit{Mod}(\mathcal{O}_\mathcal{C}),
$$
；用 (2) 的积分解表示时，它将
$(\mathcal{F}_i)_{i \in I}$ 映到
$\bigoplus j_{i, !}\mathcal{F}_i$。
类似地，给定上述 $f : K \to L$，函子 $f^*$ 有正合左伴随
$f_! : \textit{Mod}(\mathcal{O}_K) \to \textit{Mod}(\mathcal{O}_L)$。
因此，函子 $j^*$ 与 $f^*$ 正合，即 $j$ 与 $f$
是赋环拓扑斯的平坦态射（这也由等式
$\mathcal{O}_K = j^{-1}\mathcal{O}_\mathcal{C}$ 与
$\mathcal{O}_K = f^{-1}\mathcal{O}_L$ 得出）。
\end{remark}

\begin{remark}[对象之上的赋环变体]
\label{spaces-simplicial-remark-semi-representable-ringed-over-object}
设 $\mathcal{C}$ 是一个站点，$\mathcal{O}_\mathcal{C}$
是 $\mathcal{C}$ 上的环层，且 $X \in \Ob(\mathcal{C})$。
% P10-E0201
记 $\mathcal{O}_X = \mathcal{O}_\mathcal{C}|_{\mathcal{C}/U}$。
结合注 \ref{spaces-simplicial-remark-semi-representable-over-object} 与
\ref{spaces-simplicial-remark-semi-representable-ringed} 中的构造，得到：
\begin{enumerate}
\item 对 $\text{SR}(\mathcal{C}, X)$ 中的 $K$，给出赋环站点
$(\mathcal{C}/K, \mathcal{O}_K)$；
\item 若 $K = \{U_i\}$，给出分解
$\textit{Mod}(\mathcal{O}_K) =
\prod \textit{Mod}(\mathcal{O}_{U_i})$；
\item 给出赋环拓扑斯的局部化态射
$j : (\Sh(\mathcal{C}/K), \mathcal{O}_K) \to
(\Sh(\mathcal{C}/X), \mathcal{O}_X)$
；
\item 对 $\text{SR}(\mathcal{C}, X)$ 中的 $f : K \to L$，
给出赋环拓扑斯的态射
$f : (\Sh(\mathcal{C}/K), \mathcal{O}_K) \to
(\Sh(\mathcal{C}/L), \mathcal{O}_L)$。
\end{enumerate}
当然，注 \ref{spaces-simplicial-remark-semi-representable-ringed}
提到的所有结果在此框架下也成立。
\end{remark}





% P10-E0202
\section{与单纯半可表示对象关联的站点}
\label{spaces-simplicial-section-simplicial-semi-representable}

\noindent
设 $\mathcal{C}$ 是一个站点，且 $K$ 是
$\text{SR}(\mathcal{C})$ 的单纯对象。照例，令
$K_n = K([n])$，并以 $K(\varphi) : K_n \to K_m$
表示与 $\varphi : [m] \to [n]$ 相关联的态射。
由第 \ref{spaces-simplicial-section-semi-representable} 节中的构造，
在以站点为对象、以余连续函子为态射的范畴中，
我们得到单纯对象 $n \mapsto \mathcal{C}/K_n$。
换言之，得到第 \ref{spaces-simplicial-section-simplicial-sites} 节情形 B 所述的资料。
由引理 \ref{spaces-simplicial-lemma-has-P}，这些函子满足性质 P。
因此，可应用引理 \ref{spaces-simplicial-lemma-simplicial-cocontinuous-site}，
得到站点 $(\mathcal{C}/K)_{total}$。

\medskip\noindent
站点 $(\mathcal{C}/K)_{total}$ 可显式描述如下。
设 $K_n = \{U_{n, i}\}_{i \in I_n}$。对
$\varphi : [m] \to [n]$，态射
$K(\varphi) : K_n \to K_m$ 由映射
$\alpha(\varphi) : I_n \to I_m$ 以及对每个 $i \in I_n$
给定的态射
$f_{\varphi, i} : U_{n, i} \to U_{m, \alpha(\varphi)(i)}$ 所组成。
于是：
\begin{enumerate}
\item $(\mathcal{C}/K)_{total}$ 的对象对应于某个 $n$
及某个 $i \in I_n$ 下 $\mathcal{C}/U_{n, i}$ 的对象
$(U/U_{n, i})$；
\item $U/U_{n, i}$ 与 $V/U_{m, j}$ 之间的态射是二元组
$(\varphi, f)$，其中 $\varphi : [m] \to [n]$、
$j = \alpha(\varphi)(i)$，且 $f : U \to V$ 是
$\mathcal{C}$ 的态射，使得图
$$
\vcenter{
\xymatrix{
U \ar[r]_f \ar[d] & V \ar[d] \\
U_{n, i} \ar[r]^-{f_{\varphi, i}} &
U_{m, j}
}
}
$$
交换；
\item 对象 $U/U_{n, i}$ 的覆盖按如下方式构造：
从 $\mathcal{C}$ 中的覆盖 $\{f_j : U_j \to U\}$ 出发，
并令 $\{(\text{id}, f_j) : U_j/U_{n, i} \to U/U_{n, i}\}$
为 $(\mathcal{C}/K)_{total}$ 中的覆盖。
\end{enumerate}
我们为单纯站点建立的全部一般理论均适用于
$(\mathcal{C}/K)_{total}$。注意，显然的遗忘函子
$$
j_{total} : (\mathcal{C}/K)_{total} \longrightarrow \mathcal{C}
$$
连续且余连续。事实表明，与之关联的拓扑斯态射来自一个
（显然的）增广。

\begin{lemma}
\label{spaces-simplicial-lemma-augmentation-simplicial-semi-representable}
设 $\mathcal{C}$ 是一个站点，且 $K$ 是
$\text{SR}(\mathcal{C})$ 的单纯对象。局部化函子
$j_0 : \mathcal{C}/K_0 \to \mathcal{C}$ 定义增广
$a_0 : \Sh(\mathcal{C}/K_0) \to \Sh(\mathcal{C})$，
如注 \ref{spaces-simplicial-remark-augmentation-site} 的情形 (B) 所述。
引理 \ref{spaces-simplicial-lemma-augmentation-site} 中相应的拓扑斯态射
$$
a_n : \Sh(\mathcal{C}/K_n) \longrightarrow \Sh(\mathcal{C}),\quad
a : \Sh((\mathcal{C}/K)_{total}) \longrightarrow \Sh(\mathcal{C})
$$
等于与下列连续且余连续的局部化函子相关联的拓扑斯态射：
$j_n : \mathcal{C}/K_n \to \mathcal{C}$ 以及
$j_{total} : (\mathcal{C}/K)_{total} \to \mathcal{C}$。
\end{lemma}

\begin{proof}
逐一展开定义即可立即得出。关于 $a$ 与 $j_{total}$ 的关系，
特别参见引理 \ref{spaces-simplicial-lemma-augmentation-site} 证明中的脚注。
\end{proof}

\begin{lemma}
\label{spaces-simplicial-lemma-comparison}
采用引理 \ref{spaces-simplicial-lemma-augmentation-simplicial-semi-representable}
中的假设与记号，具有下列性质：
\begin{enumerate}
\item 存在函子
$a^{Sh}_! : \Sh((\mathcal{C}/K)_{total}) \to \Sh(\mathcal{C})$
，它左伴随于
$a^{-1} : \Sh(\mathcal{C}) \to \Sh((\mathcal{C}/K)_{total})$；
\item 存在函子
$a_! : \textit{Ab}((\mathcal{C}/K)_{total}) \to \textit{Ab}(\mathcal{C})$
，它左伴随于
$a^{-1} : \textit{Ab}(\mathcal{C}) \to
\textit{Ab}((\mathcal{C}/K)_{total})$；
\item 函子 $a^{-1}$ 将 $\Sh(\mathcal{C})$ 中的 $\mathcal{F}$
映到 $(\mathcal{C}/K)_{total}$ 上的层，该层在次数 $n$
上等于 $a_n^{-1}\mathcal{F}$；
\item 函子 $a_*$ 将
$\textit{Ab}((\mathcal{C}/K)_{total})$ 中的 $\mathcal{G}$
映到两个映射
$j_{0, *}\mathcal{G}_0 \to j_{1, *}\mathcal{G}_1$ 的等化子。
\end{enumerate}
\end{lemma}

\begin{proof}
(3) 与 (4) 对单纯站点的任意增广均成立；参见引理
\ref{spaces-simplicial-lemma-augmentation-site}。
由于 $j_{total}$ 连续且余连续，(1) 与 (2) 成立。
函子 $a^{Sh}_!$ 构造于《站点》引理
\ref{sites-lemma-when-shriek}，函子 $a_!$ 构造于
《站点上的模》引理
\ref{sites-modules-lemma-g-shriek-adjoint}。
\end{proof}

\begin{lemma}
\label{spaces-simplicial-lemma-sanity-check-simplicial-semi-representable}
设 $\mathcal{C}$ 是一个站点，$K$ 是
$\text{SR}(\mathcal{C})$ 的单纯对象，$U/U_{n, i}$ 是
$\mathcal{C}/K_n$ 的对象，且
$\mathcal{F} \in \textit{Ab}((\mathcal{C}/K)_{total})$。则
$$
H^p(U, \mathcal{F}) = H^p(U, \mathcal{F}_{n, i})
$$
，其中
\begin{enumerate}
% P10-E0203
\item 左端的 $U$ 被视为 $\mathcal{C}_{total}$ 的对象；
\item 右端的 $\mathcal{F}_{n, i}$ 是层 $\mathcal{F}_n$
在第 \ref{spaces-simplicial-section-semi-representable} 节的分解
$\Sh(\mathcal{C}/K_n) = \prod \Sh(\mathcal{C}/U_{n, i})$
中的第 $i$ 个分量；这里该层是
$\mathcal{C}/K_n$ 上的层。
\end{enumerate}
\end{lemma}

\begin{proof}
这立即由引理 \ref{spaces-simplicial-lemma-sanity-check} 以及第
\ref{spaces-simplicial-section-semi-representable} 节的积分解得出。
\end{proof}

\begin{remark}[对象之上的变体]
\label{spaces-simplicial-remark-augmentation-over-object}
设 $\mathcal{C}$ 是一个站点，且 $X \in \Ob(\mathcal{C})$。
回忆，我们有 $X$ 之上的半可表示对象范畴
$\text{SR}(\mathcal{C}, X) = \text{SR}(\mathcal{C}/X)$
；参见注 \ref{spaces-simplicial-remark-semi-representable-over-object}。
可以将上述讨论应用于站点 $\mathcal{C}/X$。
简言之，上述构造给出：
\begin{enumerate}
\item 对 $\text{SR}(\mathcal{C}, X)$ 的单纯对象 $K$，
给出站点 $(\mathcal{C}/K)_{total}$；
\item 给出局部化函子
$j_{total} : (\mathcal{C}/K)_{total} \to \mathcal{C}/X$；
\item 给出局部化函子 $j_n : \mathcal{C}/K_n \to \mathcal{C}/X$；
\item 给出拓扑斯态射
$a : \Sh((\mathcal{C}/K)_{total}) \to \Sh(\mathcal{C}/X)$；
\item 给出拓扑斯态射
$a_n : \Sh(\mathcal{C}/K_n) \to \Sh(\mathcal{C}/X)$；
\item 给出函子
$a^{Sh}_! : \Sh((\mathcal{C}/K)_{total}) \to \Sh(\mathcal{C}/X)$
，它左伴随于 $a^{-1}$；
\item 给出函子
$a_! : \textit{Ab}((\mathcal{C}/K)_{total}) \to \textit{Ab}(\mathcal{C}/X)$
，它左伴随于 $a^{-1}$。
\end{enumerate}
本节所有结果在此框架下均成立。为证明这一点，
只需将站点 $\mathcal{C}$ 处处替换为 $\mathcal{C}/X$。
\end{remark}

\begin{remark}[赋环变体]
\label{spaces-simplicial-remark-augmentation-ringed}
设 $\mathcal{C}$ 是一个站点，且 $\mathcal{O}_\mathcal{C}$ 是环层。
给定 $\mathcal{C}$ 的单纯半可表示对象 $K$，令
$\mathcal{O} = a^{-1}\mathcal{O}_\mathcal{C}$，其中 $a$ 如引理
\ref{spaces-simplicial-lemma-augmentation-simplicial-semi-representable} 与
\ref{spaces-simplicial-lemma-comparison} 所述。
如注 \ref{spaces-simplicial-remark-semi-representable-ringed} 中一样追踪环层，
上述构造给出：
\begin{enumerate}
\item 对 $\text{SR}(\mathcal{C})$ 的单纯对象 $K$，
给出赋环站点 $((\mathcal{C}/K)_{total}, \mathcal{O})$；
\item 给出赋环拓扑斯的态射
$a : (\Sh((\mathcal{C}/K)_{total}), \mathcal{O}) \to
(\Sh(\mathcal{C}), \mathcal{O}_\mathcal{C})$；
\item 给出赋环拓扑斯的态射
$a_n : (\Sh(\mathcal{C}/K_n), \mathcal{O}_n) \to
(\Sh(\mathcal{C}), \mathcal{O}_\mathcal{C})$；
\item 给出函子
$a_! : \textit{Mod}(\mathcal{O}) \to \textit{Mod}(\mathcal{O}_\mathcal{C})$
，它左伴随于 $a^*$。
\end{enumerate}
函子 $a_!$ 存在（但一般并不正合），因为
$a^{-1}\mathcal{O}_\mathcal{C} = \mathcal{O}$，且可以把
《站点上的模》引理 \ref{sites-modules-lemma-g-shriek-adjoint}
在引理 \ref{spaces-simplicial-lemma-comparison} 证明中的使用，替换为
《站点上的模》引理
\ref{sites-modules-lemma-lower-shriek-modules}。
如注 \ref{spaces-simplicial-remark-semi-representable-ringed} 所讨论的，
存在正合函子
$a_{n!} : \textit{Mod}(\mathcal{O}_n) \to
\textit{Mod}(\mathcal{O}_\mathcal{C})$
，它左伴随于 $a_n^*$。从而，态射 $a$ 与 $a_n$ 平坦。
注 \ref{spaces-simplicial-remark-semi-representable-ringed} 蕴含：对
$\varphi : [m] \to [n]$，赋环拓扑斯的态射
$f_\varphi : (\Sh(\mathcal{C}/K_n), \mathcal{O}_n) \to
(\Sh(\mathcal{C}/K_m), \mathcal{O}_m)$
平坦，且存在正合函子
$f_{\varphi !} : \textit{Mod}(\mathcal{O}_n) \to \textit{Mod}(\mathcal{O}_m)$
，它左伴随于 $f_\varphi^*$。这又蕴含：对赋环拓扑斯的平坦态射
$g_n : (\Sh(\mathcal{C}/K_n), \mathcal{O}_n) \to
(\Sh((\mathcal{C}/K)_{total}), \mathcal{O})$
，左伴随于 $g_n^*$ 的函子
$g_{n!} : \textit{Mod}(\mathcal{O}_n) \to
\textit{Mod}(\mathcal{O})$ 正合；参见引理
\ref{spaces-simplicial-lemma-exactness-g-shriek-modules}。
\end{remark}

\begin{remark}[对象之上的赋环变体]
\label{spaces-simplicial-remark-augmentation-ringed-over-object}
设 $\mathcal{C}$ 是一个站点，$\mathcal{O}_\mathcal{C}$ 是环层，
且 $X \in \Ob(\mathcal{C})$。记
$\mathcal{O}_X = \mathcal{O}_\mathcal{C}|_{\mathcal{C}/X}$。
结合注 \ref{spaces-simplicial-remark-augmentation-over-object} 与
\ref{spaces-simplicial-remark-augmentation-ringed} 中的构造，得到：
\begin{enumerate}
\item 对 $\text{SR}(\mathcal{C}, X)$ 的单纯对象 $K$，
给出赋环站点 $((\mathcal{C}/K)_{total}, \mathcal{O})$；
\item 给出赋环拓扑斯的态射
$a : (\Sh((\mathcal{C}/K)_{total}), \mathcal{O}) \to
(\Sh(\mathcal{C}/X), \mathcal{O}_X)$；
\item 给出赋环拓扑斯的态射
$a_n : (\Sh(\mathcal{C}/K_n), \mathcal{O}_n) \to
(\Sh(\mathcal{C}/X), \mathcal{O}_X)$；
\item 给出函子
$a_! : \textit{Mod}(\mathcal{O}) \to \textit{Mod}(\mathcal{O}_X)$
，它左伴随于 $a^*$。
\end{enumerate}
当然，注 \ref{spaces-simplicial-remark-augmentation-ringed}
提到的所有结果在此框架下也成立。
\end{remark}






\section{超覆盖的上同调下降}
\label{spaces-simplicial-section-cohomological-descent-hypercoverings}

\noindent
设 $\mathcal{C}$ 是一个站点。本节假设 $\mathcal{C}$
有等化子与纤维积。设 $K$ 是《超覆盖》定义
\ref{hypercovering-definition-hypercovering-variant}
所定义的超覆盖。我们将研究第
\ref{spaces-simplicial-section-simplicial-semi-representable} 节中的增广
$$
a : \Sh((\mathcal{C}/K)_{total}) \longrightarrow \Sh(\mathcal{C})
$$
。

\begin{lemma}
\label{spaces-simplicial-lemma-hypercovering-descent-sheaves}
设 $\mathcal{C}$ 是具有等化子与纤维积的站点，且 $K$ 是超覆盖。则
\begin{enumerate}
\item $a^{-1} : \Sh(\mathcal{C}) \to \Sh((\mathcal{C}/K)_{total})$
满忠实，其本质像为笛卡尔集合层；
\item $a^{-1} : \textit{Ab}(\mathcal{C}) \to
\textit{Ab}((\mathcal{C}/K)_{total})$
满忠实，其本质像为笛卡尔阿贝尔群层。
\end{enumerate}
在两种情形下，$a_*$ 都给出拟逆函子。
\end{lemma}

\begin{proof}
由于函子 $a^{-1}$ 与积交换，阿贝尔层的情形立即由集合层的情形得出。
在其余证明中，我们只处理集合层。
由引理 \ref{spaces-simplicial-lemma-augmentation-cartesian-module}，对
$\Sh(\mathcal{C})$ 中的 $\mathcal{F}$，$a^{-1}\mathcal{F}$ 是笛卡尔的。
只需证明：伴随映射 $\mathcal{F} \to a_*a^{-1}\mathcal{F}$
% P10-E0204
对 $\Sh(\mathcal{C})$ 中的 $\mathcal{F}$ 是同构；并且对
$(\mathcal{C}/K)_{total}$ 上的笛卡尔层 $\mathcal{G}$，
伴随映射 $a^{-1}a_*\mathcal{G} \to \mathcal{G}$ 是同构。

\medskip\noindent
设 $\mathcal{F}$ 是 $\mathcal{C}$ 上的层。
回忆，$a_*a^{-1}\mathcal{F}$ 是两个映射
$a_{0, *}a_0^{-1}\mathcal{F} \to a_{1, *}a_1^{-1}\mathcal{F}$
的等化子；参见引理 \ref{spaces-simplicial-lemma-comparison}。
由引理 \ref{spaces-simplicial-lemma-push-pull-localization}，
$$
a_{0, *}a_0^{-1}\mathcal{F} = \SheafHom(F(K_0)^\#, \mathcal{F})
\quad\text{且}\quad
a_{1, *}a_1^{-1}\mathcal{F} = \SheafHom(F(K_1)^\#, \mathcal{F})
$$
另一方面，我们知道
$$
\xymatrix{
F(K_1)^\# \ar@<1ex>[r] \ar@<-1ex>[r] &
F(K_0)^\# \ar[r] & \Sh(\mathcal{C})\text{ 的终对象 }*
}
$$
按超覆盖的定义，这是集合层范畴中的余等化子图。
因此，只需证明 $\SheafHom(-, \mathcal{F})$ 将余等化子变为等化子；
这立即由《站点》第 \ref{sites-section-glueing-sheaves} 节
的构造得出。

\medskip\noindent
设 $\mathcal{G}$ 是 $(\mathcal{C}/K)_{total}$ 上的笛卡尔层。
我们将证明，存在 $\mathcal{C}$ 上的层 $\mathcal{F}$，使得
$\mathcal{G} = a^{-1}\mathcal{F}$。这将完成证明，因为由上一段的结果，
此时 $a^{-1}a_*\mathcal{G} = a^{-1}a_*a^{-1}\mathcal{F} =
a^{-1}\mathcal{F} = \mathcal{G}$。
对 $n \geq 0$，置 $\mathcal{K}_n = F(K_n)^\#$。于是有层的映射
$$
\xymatrix{
\mathcal{K}_2
\ar@<1ex>[r]
\ar@<0ex>[r]
\ar@<-1ex>[r]
&
\mathcal{K}_1
\ar@<0.5ex>[r]
\ar@<-0.5ex>[r]
&
\mathcal{K}_0
}
$$
它们来自 $K$ 是单纯半可表示对象这一事实。
$K$ 是超覆盖这一条件意味着
$$
\mathcal{K}_1 \to \mathcal{K}_0 \times \mathcal{K}_0
\quad\text{且}\quad
\mathcal{K}_2 \to
\left(\text{cosk}_1(
\xymatrix{
\mathcal{K}_1
\ar@<0.5ex>[r]
\ar@<-0.5ex>[r]
&
\mathcal{K}_0 \ar[l]
})\right)_2
$$
是层的满射。利用引理 \ref{spaces-simplicial-lemma-characterize-cartesian}
对 $(\mathcal{C}/K)_{total}$ 上笛卡尔层的描述，以及引理
\ref{spaces-simplicial-lemma-localize-compare} 对 $\Sh(\mathcal{C}/K_n)$ 的描述，
可见我们的问题完全可以用如下两项来表述\footnote{虽然具体表述并不重要，
我们仍将其写出：问题是证明，给定
$\Sh(\mathcal{C})/\mathcal{K}_0$ 的对象
$\mathcal{G}_0/\mathcal{K}_0$ 以及同构
$$
\alpha :
\mathcal{G}_0 \times_{\mathcal{K}_0, \mathcal{K}(\delta^1_1)} \mathcal{K}_1 \to
\mathcal{G}_0 \times_{\mathcal{K}_0, \mathcal{K}(\delta^1_0)} \mathcal{K}_1
$$
它位于 $\mathcal{K}_1$ 上方并在
$\Sh(\mathcal{C})/\mathcal{K}_2$ 中满足余圈条件；则存在
$\Sh(\mathcal{C})$ 中的 $\mathcal{F}$ 和 $\mathcal{K}_0$ 上方与
$\alpha$ 相容的同构
$\mathcal{F} \times \mathcal{K}_0 \to \mathcal{G}_0$。}：
\begin{enumerate}
\item 拓扑斯 $\Sh(\mathcal{C})$；
\item $\Sh(\mathcal{C})$ 中以 $\mathcal{K}_n$ 为各项的单纯对象
$\mathcal{K}$。
\end{enumerate}
因此，按《站点》引理 \ref{sites-lemma-topos-good-site}，以另一个站点
$\mathcal{C}'$ 替换 $\mathcal{C}$ 后，可以假设 $\mathcal{C}$
具有所有有限极限，$\mathcal{C}$ 上的拓扑是亚典范的；并且
$\mathcal{C}$ 的一族态射 $\{V_j \to V\}$ 是覆盖，当且仅当
$\coprod h_{V_j} \to V$ 是满射；此外，$\mathcal{C}$ 中存在单纯对象
$U$，使得作为单纯层有 $\mathcal{K}_n = h_{U_n}$。
沿上述等价反向追溯，可以假设对所有 $n$ 都有 $K_n = \{U_n\}$。

\medskip\noindent
设 $X$ 是 $\mathcal{C}$ 的终对象。于是
$\{U_0 \to X\}$ 是覆盖，$\{U_1 \to U_0 \times U_0\}$ 是覆盖，
且 $\{U_2 \to (\text{cosk}_1 \text{sk}_1 U)_2\}$ 是覆盖。
沿用《单纯对象》定义 \ref{simplicial-definition-face-degeneracy}
的记号，以 $d^n_i : U_n \to U_{n - 1}$ 和
$s^n_j : U_n \to U_{n + 1}$ 分别表示对应于
$\delta^n_i$ 和 $\sigma^n_j$ 的态射。
稍微滥用记号，对 $\mathcal{C}$ 的态射 $c : V \to W$，
仍以同一字母表示拓扑斯态射
$c : \Sh(\mathcal{C}/V) \to \Sh(\mathcal{C}/W)$。
现在，$\mathcal{G}$ 由 $\mathcal{C}/U_0$ 上的层 $\mathcal{G}_0$
和同构
$\alpha : (d^1_1)^{-1}\mathcal{G}_0 \to (d^1_0)^{-1}\mathcal{G}_0$
给出；该同构满足引理 \ref{spaces-simplicial-lemma-characterize-cartesian}
在 $\mathcal{C}/U_2$ 上表述的余圈条件。
由于 $\{U_2 \to (\text{cosk}_1 \text{sk}_1 U)_2\}$ 是覆盖，
层上的相应拉回函子是忠实的（略去一个小细节）。
因此，可以用 $\text{cosk}_1 \text{sk}_1 U$ 替换 $U$，因为这会以
$(\text{cosk}_1 \text{sk}_1 U)_2$ 替换 $U_2$，而保持
$U_1$ 和 $U_0$ 不变。于是
$$
(d^2_0, d^2_1, d^2_2) : U_2 \to U_1 \times U_1 \times U_1
$$
% P10-E0205
是单态射，其在 $T$-值点上的像由《单纯对象》引理
\ref{simplicial-lemma-work-out} 描述。特别地，存在态射 $c$
嵌入如下交换图：
$$
\xymatrix{
U_1 \times_{(d^1_1, d^1_0), U_0 \times U_0, (d^1_1, d^1_0)} U_1
\ar[d] \ar[rr]_c & & U_2 \ar[d] \\
U_1 \times U_1
\ar[rr]^{(\text{pr}_1, \text{pr}_2, s^0_0 \circ d^1_1 \circ \text{pr}_1)} & &
U_1 \times U_1 \times U_1
}
$$
因为沿图的另一方向复合会给出 $U_2$ 的一个点。
% P10-E0206
将 $U_2$ 上 $\alpha$ 的余圈条件拉回，得到如下条件：
经两个投影拉回到
$U_1 \times_{(d^1_1, d^1_0), U_0 \times U_0, (d^1_1, d^1_0)} U_1$
的 $\alpha$ 的拉回，与经
$s^0_0 \circ d^1_1 \circ \text{pr}_1$ 拉回的 $\alpha$ 相同，
而后者为恒同映射（即 $\alpha$ 经 $s^0_0$ 的拉回是恒同映射）。
由《站点》引理 \ref{sites-lemma-glue-maps}，这意味着
$\alpha$ 来自 $\mathcal{C}/U_0 \times U_0$ 上层的同构
$$
\alpha' : \text{pr}_1^{-1}\mathcal{G}_0 \to \text{pr}_2^{-1}\mathcal{G}_0
$$
。
最后，利用 $U_1 \to U_0 \times U_0$ 的满射性和上面对 $U_2$
的描述，容易看出态射 $U_2 \to U_0 \times U_0 \times U_0$
在伴随层上是满射。因此，$\alpha'$ 在
$U_0 \times U_0 \times U_0$ 上满足余圈条件。
将《站点》引理 \ref{sites-lemma-mapping-property-glue}
应用于覆盖 $\{U_0 \to X\}$，即可完成证明。
\end{proof}

\begin{lemma}
\label{spaces-simplicial-lemma-hypercovering-cech-complex}
设 $\mathcal{C}$ 是具有等化子与纤维积的站点，且 $K$ 是超覆盖。
引理 \ref{spaces-simplicial-lemma-augmentation-cech-complex} 中与
$a^{-1}\mathcal{F}$ 关联的 {\v C}ech 复形
$$
a_{0, *}a_0^{-1}\mathcal{F} \to a_{1, *}a_1^{-1}\mathcal{F} \to
a_{2, *}a_2^{-1}\mathcal{F} \to \ldots
$$
等于复形 $\SheafHom(s(\mathbf{Z}_{F(K)}^\#), \mathcal{F})$。
这里的 $s(\mathbf{Z}_{F(K)}^\#)$ 如《超覆盖》定义
\ref{hypercovering-definition-homology} 所述。
\end{lemma}

\begin{proof}
由引理 \ref{spaces-simplicial-lemma-push-pull-localization}，有
$$
a_{n, *}a_n^{-1}\mathcal{F} = \SheafHom'(F(K_n)^\#, \mathcal{F})
$$
其中 $\SheafHom'$ 如《站点》第 \ref{sites-section-glueing-sheaves} 节所述。
引理 \ref{spaces-simplicial-lemma-augmentation-cech-complex} 的复形中的边界映射
来自单纯结构。因此，复形的相等性来自对
$\Sh(\mathcal{C})$ 中 $
\mathcal{G}$ 的典范恒等式
$\SheafHom'(\mathcal{G}, \mathcal{F}) =
\SheafHom(\mathbf{Z}_\mathcal{G}, \mathcal{F})$。
\end{proof}

\begin{lemma}
\label{spaces-simplicial-lemma-hypercovering-descent-bounded-abelian}
设 $\mathcal{C}$ 是具有等化子与纤维积的站点，且 $K$ 是超覆盖。
对 $E \in D(\mathcal{C})$，映射
$$
E \longrightarrow Ra_*a^{-1}E
$$
是同构。
\end{lemma}

\begin{proof}
首先，设 $\mathcal{I}$ 是 $\mathcal{C}$ 上的内射阿贝尔层。
引理 \ref{spaces-simplicial-lemma-augmentation-spectral-sequence}
对层 $a^{-1}\mathcal{I}$ 给出的谱序列退化，因为由引理
\ref{spaces-simplicial-lemma-localize-injective}，
$(a^{-1}\mathcal{I})_p = a_p^{-1}\mathcal{I}$ 是内射的。因此，复形
$$
a_{0, *}a_0^{-1}\mathcal{I} \to
a_{1, *}a_1^{-1}\mathcal{I} \to
a_{2, *}a_2^{-1}\mathcal{I} \to\ldots
$$
计算 $Ra_*a^{-1}\mathcal{I}$。由引理
\ref{spaces-simplicial-lemma-hypercovering-cech-complex}，它等于复形
$\SheafHom(s(\mathbf{Z}_{F(K)}^\#), \mathcal{I})$。
由于 $K$ 是超覆盖，将《超覆盖》引理
\ref{hypercovering-lemma-acyclic-hypercover-sheaves}
应用于单纯预层 $F(K)$，可知 $s(\mathbf{Z}_{F(K)}^\#)$
在次数 $> 0$ 处正合。由于 $\mathcal{I}$ 是内射的，函子
$\SheafHom(-, \mathcal{I})$ 正合，从而
$\SheafHom(s(\mathbf{Z}_{F(K)}^\#), \mathcal{I})$
在正次数处正合。因此，对 $p > 0$，有
$R^pa_*a^{-1}\mathcal{I} = 0$。另一方面，由引理
\ref{spaces-simplicial-lemma-hypercovering-descent-sheaves}，有
$\mathcal{I} = a_*a^{-1}\mathcal{I}$。

\medskip\noindent
下有界情形。设 $E \in D^+(\mathcal{C})$。
选取表示 $E$ 的下有界内射复形 $\mathcal{I}^\bullet$。
由第一段的结果和 Leray 无环性引理
（《导出范畴》引理 \ref{derived-lemma-leray-acyclicity}），
$Ra_*a^{-1}\mathcal{I}^\bullet$ 由复形
$a_*a^{-1}\mathcal{I}^\bullet = \mathcal{I}^\bullet$
计算。因此，本情形下引理成立。

\medskip\noindent
无界情形。我们建议读者略过这一部分，因为论证与上文相同，
只是为绕开收敛性问题，我们使用双复形的显式表示。
设 $E \in D(\mathcal{C})$。为了证明映射
$E \to Ra_*a^{-1}E$ 是同构，只需证明对
$\mathcal{C}$ 的每个对象 $U$ 都有
$$
R\Gamma(U, E) = R\Gamma(U, Ra_*a^{-1}E)
$$
我们将计算两边，并证明映射 $E \to Ra_*a^{-1}E$ 诱导同构。
选取表示 $E$ 的 K-内射复形 $\mathcal{I}^\bullet$。
再选取拟同构
$a^{-1}\mathcal{I}^\bullet \to \mathcal{J}^\bullet$，
其中 $\mathcal{J}^\bullet$ 是 $(\mathcal{C}/K)_{total}$
上的某个 K-内射复形。我们有
$$
R\Gamma(U, E) = R\Hom(\mathbf{Z}_U^\#, E)
$$
以及
$$
R\Gamma(U, Ra_*a^{-1}E) = R\Hom(\mathbf{Z}_U^\#, Ra_*a^{-1}E) =
R\Hom(a^{-1}\mathbf{Z}_U^\#, a^{-1}E)
$$
由引理 \ref{spaces-simplicial-lemma-simplicial-resolution-augmentation}，有拟同构
$$
\Big(\ldots \to
g_{2!}(a_2^{-1}\mathbf{Z}_U^\#) \to
g_{1!}(a_1^{-1}\mathbf{Z}_U^\#) \to
g_{0!}(a_0^{-1}\mathbf{Z}_U^\#)\Big)
\longrightarrow
a^{-1}\mathbf{Z}_U^\#
$$
因此，$R\Hom(a^{-1}\mathbf{Z}_U^\#, a^{-1}E)$ 等于
$$
R\Gamma((\mathcal{C}/K)_{total},
R\SheafHom(
\ldots \to
g_{2!}(a_2^{-1}\mathbf{Z}_U^\#) \to
g_{1!}(a_1^{-1}\mathbf{Z}_U^\#) \to
g_{0!}(a_0^{-1}\mathbf{Z}_U^\#),
\mathcal{J}^\bullet))
$$
根据《站点上的上同调》第
\ref{sites-cohomology-section-internal-hom} 节中的构造，并利用
$\mathcal{J}^\bullet$ 是 K-内射的，可知它由具有如下各项的
阿贝尔群复形表示：
$$
\prod\nolimits_{p + q = n}
\Hom(g_{p!}(a_p^{-1}\mathbf{Z}_U^\#), \mathcal{J}^q) =
\prod\nolimits_{p + q = n}
\Hom(a_p^{-1}\mathbf{Z}_U^\#, g_p^{-1}\mathcal{J}^q)
$$
更多细节见《站点上的上同调》引理
\ref{sites-cohomology-lemma-RHom-into-K-injective} 和
\ref{sites-cohomology-lemma-section-RHom-over-U}。
因此，$R\Gamma(U, Ra_*a^{-1}E)$ 由乘积全复形
$\text{Tot}_\pi(B^{\bullet, \bullet})$ 计算，其中
$B^{p, q} = \Hom(a_p^{-1}\mathbf{Z}_U^\#, g_p^{-1}\mathcal{J}^q)$。
对另一边作类似论证。首先注意到
$$
s(\mathbf{Z}_{F(K)}^\#) \longrightarrow \mathbf{Z}
$$
由《超覆盖》引理 \ref{hypercovering-lemma-acyclic-hypercover-sheaves}，
这是 $\mathcal{C}$ 上复形的拟同构。由于 $\mathbf{Z}_U^\#$
是平坦的 $\mathbf{Z}$-模层，可知
$$
s(\mathbf{Z}_{F(K)}^\#) \otimes_\mathbf{Z} \mathbf{Z}_U^\#
\longrightarrow
\mathbf{Z}_U^\#
$$
是拟同构。因此，$R\Hom(\mathbf{Z}_U^\#, E)$ 等于
$$
R\Gamma(\mathcal{C}, R\SheafHom(
s(\mathbf{Z}_{F(K)}^\#) \otimes_\mathbf{Z} \mathbf{Z}_U^\#,
\mathcal{I}^\bullet))
$$
根据 $R\SheafHom$ 的构造，并利用 $\mathcal{I}^\bullet$
是 K-内射的，它由具有如下各项的阿贝尔群复形表示：
$$
\prod\nolimits_{p + q = n}
\Hom(\mathbf{Z}^\#_{K_p} \otimes_\mathbf{Z} \mathbf{Z}_U^\#, \mathcal{I}^q)
=
\prod\nolimits_{p + q = n}
\Hom(a_p^{-1}\mathbf{Z}_U^\#, a_p^{-1}\mathcal{I}^q)
$$
各项相等是因为由《站点上的模》评注
\ref{sites-modules-remark-j-shriek-tensor}，有
$\mathbf{Z}^\#_{K_p} \otimes_\mathbf{Z} \mathbf{Z}_U^\# =
a_{p!}a_p^{-1}\mathbf{Z}_U^\#$。
因此，$R\Gamma(U, E)$ 由乘积全复形
$\text{Tot}_\pi(A^{\bullet, \bullet})$ 计算，其中
$A^{p, q} = \Hom(a_p^{-1}\mathbf{Z}_U^\#, a_p^{-1}\mathcal{I}^q)$。

\medskip\noindent
由于 $\mathcal{I}^\bullet$ 是 K-内射的，由引理
\ref{spaces-simplicial-lemma-localize-injective}，$a_p^{-1}\mathcal{I}^\bullet$
是 K-内射的。由于 $\mathcal{J}^\bullet$ 是 K-内射的，由引理
\ref{spaces-simplicial-lemma-restriction-injective-to-component-site}，
$g_p^{-1}\mathcal{J}^\bullet$ 是 K-内射的。二者都表示对象
$a_p^{-1}E$。因此，对每个 $p \geq 0$，由给定映射
$a^{-1}\mathcal{I}^\bullet \to \mathcal{J}^\bullet$
在 $g_p^{-1}$ 下的像诱导的复形映射
$$
A^{p, \bullet} = \Hom(a_p^{-1}\mathbf{Z}_U^\#, a_p^{-1}\mathcal{I}^\bullet)
\longrightarrow
\Hom(a_p^{-1}\mathbf{Z}_U^\#, g_p^{-1}\mathcal{J}^\bullet) = B^{p, \bullet}
$$
% P10-E0207
是拟同构，因为这两个复形都计算
$$
R\Hom(a_p^{-1}\mathbf{Z}_U^\#, a_p^{-1}E)
$$
由《代数进阶》引理 \ref{more-algebra-lemma-prod-qis-gives-qis}，
可知如下交换图中的右侧竖直箭头是拟同构：
$$
\xymatrix{
R\Gamma(U, E) \ar[r] \ar[d] &
\text{Tot}_\pi(A^{\bullet, \bullet}) \ar[d] \\
R\Gamma(U, Ra_*a^{-1}E) \ar[r] &
\text{Tot}_\pi(B^{\bullet, \bullet})
}
$$
由于上面已经看到水平箭头是拟同构，左侧竖直箭头也是拟同构。
\end{proof}

\begin{lemma}
\label{spaces-simplicial-lemma-compare-cohomology-hypercovering}
设 $\mathcal{C}$ 是具有等化子与纤维积的站点，且 $K$ 是超覆盖。
则有典范同构
$$
R\Gamma(\mathcal{C}, E) =
R\Gamma((\mathcal{C}/K)_{total}, a^{-1}E)
$$
其中 $E \in D(\mathcal{C})$。
\end{lemma}

\begin{proof}
这由引理 \ref{spaces-simplicial-lemma-hypercovering-descent-bounded-abelian} 得出，
因为按《站点上的上同调》评注
\ref{sites-cohomology-remark-before-Leray}，有
$R\Gamma((\mathcal{C}/K)_{total}, -) =
R\Gamma(\mathcal{C}, -) \circ Ra_*$。
\end{proof}

\begin{lemma}
\label{spaces-simplicial-lemma-hypercovering-equivalence-bounded}
设 $\mathcal{C}$ 是具有等化子与纤维积的站点，且 $K$ 是超覆盖。
以 $\mathcal{A} \subset \textit{Ab}((\mathcal{C}/K)_{total})$
表示笛卡尔阿贝尔层组成的弱 Serre 子范畴。则函子 $a^{-1}$ 给出等价
$$
D^+(\mathcal{C}) \longrightarrow D_\mathcal{A}^+((\mathcal{C}/K)_{total})
$$
，其拟逆为 $Ra_*$。
\end{lemma}

\begin{proof}
由引理 \ref{spaces-simplicial-lemma-Serre-subcat-cartesian-modules}，
$\mathcal{A}$ 是弱 Serre 子范畴。该等价是迄今所得结果的形式推论。
应用引理 \ref{spaces-simplicial-lemma-hypercovering-descent-sheaves}、
\ref{spaces-simplicial-lemma-hypercovering-descent-bounded-abelian} 以及
《站点上的上同调》引理
\ref{sites-cohomology-lemma-equivalence-bounded}。
\end{proof}

\noindent
我们建议读者略过下面的评注。

\begin{remark}
\label{spaces-simplicial-remark-compare-cohomology-hypercovering-presheaf}
设 $\mathcal{C}$ 是站点，$\mathcal{G}$ 是 $\mathcal{C}$
上的集合预层。如果 $\mathcal{C}$ 具有等化子与纤维积，
那么我们已在《超覆盖》定义
\ref{hypercovering-definition-hypercovering-variant}
中定义了 $\mathcal{G}$ 的超覆盖。我们断言，本节的所有结果
在这一框架下都有有效的对应版本。为此，完全仿照《站点》引理
\ref{sites-lemma-localize-topos-site}，定义 $\mathcal{C}$
在 $\mathcal{G}$ 处的局部化 $\mathcal{C}/\mathcal{G}$
（该引理只对层表述；拓扑斯 $\Sh(\mathcal{C}/\mathcal{G})$
等于拓扑斯 $\Sh(\mathcal{C})$ 在层 $\mathcal{G}^\#$ 处的局部化）。
读者容易证明，站点 $\mathcal{C}/\mathcal{G}$ 具有纤维积与等化子，
并且 $\mathcal{C}$ 中 $\mathcal{G}$ 的超覆盖与站点
$\mathcal{C}/\mathcal{G}$ 的超覆盖是同一回事。因此，在上面关于
超覆盖的引理中用 $\mathcal{C}/\mathcal{G}$ 替换站点 $\mathcal{C}$，
便得到关于 $\mathcal{G}$ 的超覆盖的相应结果的证明。
例如，对 $\mathcal{G}$ 的超覆盖 $K$，有
$$
R\Gamma(\mathcal{C}/\mathcal{G}, E) =
R\Gamma((\mathcal{C}/K)_{total}, a^{-1}E)
$$
其中 $E \in D^+(\mathcal{C}/\mathcal{G})$，而
$a : \Sh((\mathcal{C}/K)_{total}) \to \Sh(\mathcal{C}/\mathcal{G})$
是典范增广。这就是引理 \ref{spaces-simplicial-lemma-compare-cohomology-hypercovering}。
把 $R\Gamma(\mathcal{G}, -) : D(\mathcal{C}) \to D(\textit{Ab})$
定义为函子 $H^0(\mathcal{G}, -) = H^0(\mathcal{G}^\#, -)$
的导出函子；该函子见《超覆盖》第
\ref{hypercovering-section-hypercoverings-verdier} 节和
《站点上的上同调》第 \ref{sites-cohomology-section-limp} 节。我们有
$$
R\Gamma(\mathcal{G}, E) = R\Gamma(\mathcal{C}/\mathcal{G}, j^{-1}E)
$$
% P10-E0208
这是《站点上的上同调》引理
\ref{sites-cohomology-lemma-cohomology-of-open}
对局部化函子 $j : \mathcal{C}/\mathcal{G} \to \mathcal{C}$
的类似版本。综合以上各点，得到
$$
R\Gamma(\mathcal{G}, E) =
R\Gamma((\mathcal{C}/K)_{total}, a^{-1}j^{-1}E) =
R\Gamma((\mathcal{C}/K)_{total}, g^{-1}E)
$$
其中 $E \in D^+(\mathcal{C})$，而
$g : \Sh((\mathcal{C}/K)_{total}) \to \Sh(\mathcal{C})$
是 $a$ 与 $j$ 的复合。
\end{remark}







\section{超覆盖的上同调下降：模}
\label{spaces-simplicial-section-cohomological-descent-hypercoverings-modules}

\noindent
设 $\mathcal{C}$ 是站点，$\mathcal{O}_\mathcal{C}$ 是环层。
假设 $\mathcal{C}$ 具有等化子与纤维积，并设 $K$ 是《超覆盖》定义
\ref{hypercovering-definition-hypercovering-variant}
所定义的超覆盖。我们将研究评注 \ref{spaces-simplicial-remark-augmentation-ringed}
中的增广
$$
a :
(\Sh((\mathcal{C}/K)_{total}), \mathcal{O})
\longrightarrow
(\Sh(\mathcal{C}), \mathcal{O}_\mathcal{C})
$$
的上同调下降。

\begin{lemma}
\label{spaces-simplicial-lemma-hypercovering-descent-modules}
设 $\mathcal{C}$ 是具有等化子与纤维积的站点，
$\mathcal{O}_\mathcal{C}$ 是环层，且 $K$ 是超覆盖。
采用上述记号，
$$
a^* : \textit{Mod}(\mathcal{O}_\mathcal{C}) \to \textit{Mod}(\mathcal{O})
$$
满忠实，其本质像为笛卡尔 $\mathcal{O}$-模。
函子 $a_*$ 给出拟逆。
\end{lemma}

\begin{proof}
由于 $a^{-1}\mathcal{O}_\mathcal{C} = \mathcal{O}$，有
$a^* = a^{-1}$。因此，本引理立即由引理
\ref{spaces-simplicial-lemma-hypercovering-descent-sheaves} 得出。
\end{proof}

\begin{lemma}
\label{spaces-simplicial-lemma-hypercovering-descent-bounded-modules}
设 $\mathcal{C}$ 是具有等化子与纤维积的站点，
$\mathcal{O}_\mathcal{C}$ 是环层，且 $K$ 是超覆盖。
对 $E \in D(\mathcal{O}_\mathcal{C})$，映射
$$
E \longrightarrow Ra_*La^*E
$$
是同构。
\end{lemma}

\begin{proof}
由于 $a^{-1}\mathcal{O}_\mathcal{C} = \mathcal{O}$，有
$La^* = a^* = a^{-1}$。此外，模上的 $Ra_*$ 与阿贝尔层上的
% P10-E0209
$Ra_*$ 一致；参见《站点上的上同调》引理
\ref{sites-cohomology-lemma-modules-abelian-unbounded}。
因此，本引理立即由引理
\ref{spaces-simplicial-lemma-hypercovering-descent-bounded-abelian} 得出。
\end{proof}

\begin{lemma}
\label{spaces-simplicial-lemma-compare-cohomology-hypercovering-modules}
设 $\mathcal{C}$ 是具有等化子与纤维积的站点，
$\mathcal{O}_\mathcal{C}$ 是环层，且 $K$ 是超覆盖。
则有典范同构
$$
R\Gamma(\mathcal{C}, E) =
R\Gamma((\mathcal{C}/K)_{total}, La^*E)
$$
其中 $E \in D(\mathcal{O}_\mathcal{C})$。
\end{lemma}

\begin{proof}
这由引理 \ref{spaces-simplicial-lemma-hypercovering-descent-bounded-modules} 得出，
因为按《站点上的上同调》评注
\ref{sites-cohomology-remark-before-Leray}，或按《站点上的上同调》引理
\ref{sites-cohomology-lemma-Leray-unbounded}，有
$R\Gamma((\mathcal{C}/K)_{total}, -) =
R\Gamma(\mathcal{C}, -) \circ Ra_*$。
\end{proof}

\begin{lemma}
\label{spaces-simplicial-lemma-hypercovering-equivalence-bounded-modules}
设 $\mathcal{C}$ 是具有等化子与纤维积的站点，
$\mathcal{O}_\mathcal{C}$ 是环层，且 $K$ 是超覆盖。
以 $\mathcal{A} \subset \textit{Mod}(\mathcal{O})$
表示笛卡尔 $\mathcal{O}$-模组成的弱 Serre 子范畴。
则函子 $La^*$ 给出等价
$$
D^+(\mathcal{O}_\mathcal{C}) \longrightarrow D_\mathcal{A}^+(\mathcal{O})
$$
，其拟逆为 $Ra_*$。
\end{lemma}

\begin{proof}
由引理 \ref{spaces-simplicial-lemma-Serre-subcat-cartesian-modules}，
$\mathcal{A}$ 是弱 Serre 子范畴（评注
\ref{spaces-simplicial-remark-augmentation-ringed} 中的讨论表明所需假设成立）。
该等价是迄今所得结果的形式推论。应用引理
\ref{spaces-simplicial-lemma-hypercovering-descent-modules}、
\ref{spaces-simplicial-lemma-hypercovering-descent-bounded-modules} 以及
《站点上的上同调》引理 \ref{sites-cohomology-lemma-equivalence-bounded}。
\end{proof}






\section{对象的超覆盖之上同调下降}
\label{spaces-simplicial-section-cohomological-descent-hypercoverings-X}

\noindent
本节假设 $\mathcal{C}$ 具有纤维积，并设 $X \in \Ob(\mathcal{C})$。
设 $K$ 是《超覆盖》定义
\ref{hypercovering-definition-hypercovering} 所定义的 $X$ 的超覆盖。
我们将研究评注 \ref{spaces-simplicial-remark-augmentation-over-object} 中的增广
$$
a : \Sh((\mathcal{C}/K)_{total}) \longrightarrow \Sh(\mathcal{C}/X)
$$
。注意，$\mathcal{C}/X$ 是具有等化子与纤维积的站点；并且
$K$ 是站点 $\mathcal{C}/X$ 的超覆盖\footnote{逆命题未必成立。
也就是说，若 $K$ 是
$\text{SR}(\mathcal{C}, X) = \text{SR}(\mathcal{C}/X)$
的单纯对象，并按《超覆盖》定义
\ref{hypercovering-definition-hypercovering-variant}
定义了站点 $\mathcal{C}/X$ 的超覆盖，则 $K$ 未必定义 $X$
的超覆盖。例如，若 $K_0 = \{U_{0, i}\}_{i \in I_0}$，
则后一个条件保证 $\{U_{0, i} \to X\}$ 是 $\mathcal{C}$ 的覆盖，
而前一个条件仅要求
$\coprod h_{U_{0, i}}^\# \to h_X^\#$ 是层的满射。}；
这由《超覆盖》引理 \ref{hypercovering-lemma-hypercovering-F} 得出。
这意味着，第 \ref{spaces-simplicial-section-cohomological-descent-hypercoverings}
节对超覆盖证明的每个结果，在本节的情形下都有直接的对应版本。

\begin{lemma}
\label{spaces-simplicial-lemma-hypercovering-X-descent-sheaves}
设 $\mathcal{C}$ 是具有纤维积的站点，$X \in \Ob(\mathcal{C})$，
且 $K$ 是 $X$ 的超覆盖。则
\begin{enumerate}
\item $a^{-1} : \Sh(\mathcal{C}/X) \to \Sh((\mathcal{C}/K)_{total})$
满忠实，其本质像为笛卡尔集合层；
\item $a^{-1} : \textit{Ab}(\mathcal{C}/X) \to
\textit{Ab}((\mathcal{C}/K)_{total})$
满忠实，其本质像为笛卡尔阿贝尔群层。
\end{enumerate}
在两种情形下，$a_*$ 都给出拟逆函子。
\end{lemma}

\begin{proof}
由评注 \ref{spaces-simplicial-remark-semi-representable-over-object}、
\ref{spaces-simplicial-remark-augmentation-over-object} 以及本节引言中的讨论，
这由引理 \ref{spaces-simplicial-lemma-hypercovering-descent-sheaves} 得出。
\end{proof}

\begin{lemma}
\label{spaces-simplicial-lemma-hypercovering-X-descent-bounded-abelian}
% P10-E0210
设 $\mathcal{C}$ 是具有纤维积的站点，$X \in \Ob(\mathcal{C})$，
且 $K$ 是 $X$ 的超覆盖。对 $E \in D(\mathcal{C}/X)$，映射
$$
E \longrightarrow Ra_*a^{-1}E
$$
是同构。
\end{lemma}

\begin{proof}
由评注 \ref{spaces-simplicial-remark-semi-representable-over-object}、
\ref{spaces-simplicial-remark-augmentation-over-object} 以及本节引言中的讨论，
这由引理 \ref{spaces-simplicial-lemma-hypercovering-descent-bounded-abelian} 得出。
\end{proof}

\begin{lemma}
\label{spaces-simplicial-lemma-compare-cohomology-hypercovering-X}
设 $\mathcal{C}$ 是具有纤维积的站点，$X \in \Ob(\mathcal{C})$，
且 $K$ 是 $X$ 的超覆盖。则有典范同构
$$
R\Gamma(X, E) = R\Gamma((\mathcal{C}/K)_{total}, a^{-1}E)
$$
其中 $E \in D(\mathcal{C}/X)$。
\end{lemma}

\begin{proof}
由评注 \ref{spaces-simplicial-remark-semi-representable-over-object} 和
\ref{spaces-simplicial-remark-augmentation-over-object}，这由引理
\ref{spaces-simplicial-lemma-compare-cohomology-hypercovering} 得出。
\end{proof}

\begin{lemma}
\label{spaces-simplicial-lemma-hypercovering-X-equivalence-bounded}
设 $\mathcal{C}$ 是具有纤维积的站点，$X \in \Ob(\mathcal{C})$，
且 $K$ 是 $X$ 的超覆盖。以
$\mathcal{A} \subset \textit{Ab}((\mathcal{C}/K)_{total})$
表示笛卡尔阿贝尔层组成的弱 Serre 子范畴。
则函子 $a^{-1}$ 给出等价
$$
D^+(\mathcal{C}/X) \longrightarrow D_\mathcal{A}^+((\mathcal{C}/K)_{total})
$$
，其拟逆为 $Ra_*$。
\end{lemma}

\begin{proof}
由评注 \ref{spaces-simplicial-remark-semi-representable-over-object} 和
\ref{spaces-simplicial-remark-augmentation-over-object}，这由引理
\ref{spaces-simplicial-lemma-hypercovering-equivalence-bounded} 得出。
\end{proof}








\section{对象的超覆盖之上同调下降：模}
\label{spaces-simplicial-section-cohomological-descent-hypercoverings-X-modules}

\noindent
本节假设 $\mathcal{C}$ 具有纤维积，并设 $X \in \Ob(\mathcal{C})$。
设 $K$ 是《超覆盖》定义
\ref{hypercovering-definition-hypercovering} 所定义的 $X$ 的超覆盖。
设 $\mathcal{O}_\mathcal{C}$ 是 $\mathcal{C}$ 上的环层，并置
$\mathcal{O}_X = \mathcal{O}_\mathcal{C}|_{\mathcal{C}/X}$。
我们将研究评注 \ref{spaces-simplicial-remark-augmentation-ringed-over-object} 中的增广
$$
a :
(\Sh((\mathcal{C}/K)_{total}), \mathcal{O})
\longrightarrow
(\Sh(\mathcal{C}/X), \mathcal{O}_X)
$$
。注意，$\mathcal{C}/X$ 是具有等化子与纤维积的站点，
而 $K$ 是站点 $\mathcal{C}/X$ 的超覆盖。因此，本节的结果
立即由第 \ref{spaces-simplicial-section-cohomological-descent-hypercoverings-modules}
节中的相应结果得出。

\begin{lemma}
\label{spaces-simplicial-lemma-hypercovering-X-descent-modules}
设 $\mathcal{C}$ 是具有纤维积的站点，$X \in \Ob(\mathcal{C})$，
$\mathcal{O}_\mathcal{C}$ 是环层，且 $K$ 是 $X$ 的超覆盖。
采用上述记号，
$$
a^* : \textit{Mod}(\mathcal{O}_X) \to \textit{Mod}(\mathcal{O})
$$
满忠实，其本质像为笛卡尔 $\mathcal{O}$-模。
函子 $a_*$ 给出拟逆。
\end{lemma}

\begin{proof}
由评注 \ref{spaces-simplicial-remark-semi-representable-ringed-over-object}、
\ref{spaces-simplicial-remark-augmentation-ringed-over-object} 以及本节引言中的讨论，
这由引理 \ref{spaces-simplicial-lemma-hypercovering-descent-modules} 得出。
\end{proof}

\begin{lemma}
\label{spaces-simplicial-lemma-hypercovering-X-descent-bounded-modules}
设 $\mathcal{C}$ 是具有纤维积的站点，$X \in \Ob(\mathcal{C})$，
$\mathcal{O}_\mathcal{C}$ 是环层，且 $K$ 是 $X$ 的超覆盖。
对 $E \in D(\mathcal{O}_X)$，映射
$$
E \longrightarrow Ra_*La^*E
$$
是同构。
\end{lemma}

\begin{proof}
由评注 \ref{spaces-simplicial-remark-semi-representable-ringed-over-object}、
\ref{spaces-simplicial-remark-augmentation-ringed-over-object} 以及本节引言中的讨论，
这由引理 \ref{spaces-simplicial-lemma-hypercovering-descent-bounded-modules} 得出。
\end{proof}

\begin{lemma}
\label{spaces-simplicial-lemma-compare-cohomology-hypercovering-X-modules}
设 $\mathcal{C}$ 是具有纤维积的站点，$X \in \Ob(\mathcal{C})$，
$\mathcal{O}_\mathcal{C}$ 是环层，且 $K$ 是 $X$ 的超覆盖。
则有典范同构
$$
R\Gamma(X, E) = R\Gamma((\mathcal{C}/K)_{total}, La^*E)
$$
% P10-E0211
其中 $E \in D(\mathcal{O}_\mathcal{C})$。
\end{lemma}

\begin{proof}
由评注 \ref{spaces-simplicial-remark-semi-representable-ringed-over-object}、
\ref{spaces-simplicial-remark-augmentation-ringed-over-object} 以及本节引言中的讨论，
这由引理 \ref{spaces-simplicial-lemma-compare-cohomology-hypercovering-modules} 得出。
\end{proof}

\begin{lemma}
\label{spaces-simplicial-lemma-hypercovering-X-equivalence-bounded-modules}
设 $\mathcal{C}$ 是具有纤维积的站点，$X \in \Ob(\mathcal{C})$，
$\mathcal{O}_\mathcal{C}$ 是环层，且 $K$ 是 $X$ 的超覆盖。
以 $\mathcal{A} \subset \textit{Mod}(\mathcal{O})$
表示笛卡尔 $\mathcal{O}$-模组成的弱 Serre 子范畴。
则函子 $La^*$ 给出等价
$$
D^+(\mathcal{O}_X) \longrightarrow D_\mathcal{A}^+(\mathcal{O})
$$
，其拟逆为 $Ra_*$。
\end{lemma}

\begin{proof}
由评注 \ref{spaces-simplicial-remark-semi-representable-ringed-over-object}、
\ref{spaces-simplicial-remark-augmentation-ringed-over-object} 以及本节引言中的讨论，
这由引理 \ref{spaces-simplicial-lemma-hypercovering-equivalence-bounded-modules} 得出。
\end{proof}










\section{站点的单纯对象给出的超覆盖}
\label{spaces-simplicial-section-hypercovering}

\noindent
设 $\mathcal{C}$ 是具有纤维积的站点，且 $X \in \Ob(\mathcal{C})$。
本节将阐明第 \ref{spaces-simplicial-section-cohomological-descent-hypercoverings-X}
节在超覆盖由站点的单纯对象给出时的结果。
设 $U$ 是 $\mathcal{C}$ 的单纯对象。照例，记
$U_n = U([n])$，并以 $f_\varphi : U_n \to U_m$ 表示对应于
$\varphi : [m] \to [n]$ 的态射 $f_\varphi = U(\varphi)$。
假设有增广
$$
a : U \to X
$$
由此得到单纯站点 $(\mathcal{C}/U)_{total}$ 和增广态射
$$
a : \Sh((\mathcal{C}/U)_{total}) \longrightarrow \Sh(\mathcal{C}/X)
$$
具体而言，由 $U$ 得到 $\text{SR}(\mathcal{C}, X)$ 的单纯对象
% P10-E0212
$K$，其次数 $n$ 部分为 $K_n = \{U_n \to X\}$；于是可以应用
评注 \ref{spaces-simplicial-remark-augmentation-over-object} 中的构造。
更明确地说，站点 $(\mathcal{C}/U)_{total}$ 的对象由
% P10-E0213
$V/U_n$ 给出；态射 $(\varphi, f) : V/U_n \to W/U_m$
由 $\Delta$ 中的态射 $\varphi : [m] \to [n]$ 和态射
$f : V \to W$ 给出，并要求下图交换：
$$
\xymatrix{
V \ar[r]_f \ar[d] & W \ar[d] \\
U_n \ar[r]^{f_\varphi} & U_m
}
$$
拓扑斯态射 $a$ 由余连续函子 $V/U_n \mapsto V/X$ 给出。
仅此而已！

\medskip\noindent
本节称增广 $a : U \to X$ 是
{\it $\mathcal{C}$ 中 $X$ 的超覆盖}，如果满足下列条件：
\begin{enumerate}
\item $\{U_0 \to X\}$ 是 $\mathcal{C}$ 的覆盖；
\item $\{U_1 \to U_0 \times_X U_0\}$ 是 $\mathcal{C}$ 的覆盖；
\item $\{U_{n + 1} \to (\text{cosk}_n\text{sk}_n U)_{n + 1}\}$
对 $n \geq 1$ 是 $\mathcal{C}$ 的覆盖。
\end{enumerate}
这等价于上述 $K$ 是 $X$ 的超覆盖这一条件；参见《超覆盖》例
\ref{hypercovering-example-hypercovering-in-C}。

\begin{lemma}
\label{spaces-simplicial-lemma-hypercovering-X-simple-descent-sheaves}
% P10-E0210
设 $\mathcal{C}$ 是具有纤维积的站点，$X \in \Ob(\mathcal{C})$，
且 $a : U \to X$ 是上述意义下 $\mathcal{C}$ 中 $X$ 的超覆盖。则
\begin{enumerate}
\item $a^{-1} : \Sh(\mathcal{C}/X) \to \Sh((\mathcal{C}/U)_{total})$
满忠实，其本质像为笛卡尔集合层；
\item $a^{-1} : \textit{Ab}(\mathcal{C}/X) \to
\textit{Ab}((\mathcal{C}/U)_{total})$
满忠实，其本质像为笛卡尔阿贝尔群层。
\end{enumerate}
在两种情形下，$a_*$ 都给出拟逆函子。
\end{lemma}

\begin{proof}
这是引理 \ref{spaces-simplicial-lemma-hypercovering-X-descent-sheaves} 的特殊情形。
\end{proof}

\begin{lemma}
\label{spaces-simplicial-lemma-hypercovering-X-simple-descent-bounded-abelian}
% P10-E0210
设 $\mathcal{C}$ 是具有纤维积的站点，$X \in \Ob(\mathcal{C})$，
且 $a : U \to X$ 是上述意义下 $\mathcal{C}$ 中 $X$ 的超覆盖。
对 $E \in D(\mathcal{C}/X)$，映射
$$
E \longrightarrow Ra_*a^{-1}E
$$
是同构。
\end{lemma}

\begin{proof}
这是引理 \ref{spaces-simplicial-lemma-hypercovering-X-descent-bounded-abelian}
的特殊情形。
\end{proof}

\begin{lemma}
\label{spaces-simplicial-lemma-compare-cohomology-hypercovering-X-simple}
设 $\mathcal{C}$ 是具有纤维积的站点，$X \in \Ob(\mathcal{C})$，
且 $a : U \to X$ 是上述意义下 $\mathcal{C}$ 中 $X$ 的超覆盖。
则有典范同构
$$
R\Gamma(X, E) = R\Gamma((\mathcal{C}/U)_{total}, a^{-1}E)
$$
其中 $E \in D(\mathcal{C}/X)$。
\end{lemma}

\begin{proof}
这是引理 \ref{spaces-simplicial-lemma-compare-cohomology-hypercovering-X} 的特殊情形。
\end{proof}

\begin{lemma}
\label{spaces-simplicial-lemma-hypercovering-X-simple-equivalence-bounded}
% P10-E0210
设 $\mathcal{C}$ 是具有纤维积的站点，$X \in \Ob(\mathcal{C})$，
且 $a : U \to X$ 是上述意义下 $\mathcal{C}$ 中 $X$ 的超覆盖。
以 $\mathcal{A} \subset \textit{Ab}((\mathcal{C}/U)_{total})$
表示笛卡尔阿贝尔层组成的弱 Serre 子范畴。
则函子 $a^{-1}$ 给出等价
$$
D^+(\mathcal{C}/X) \longrightarrow D_\mathcal{A}^+((\mathcal{C}/U)_{total})
$$
，其拟逆为 $Ra_*$。
\end{lemma}

\begin{proof}
% P10-E0214
这是引理 \ref{spaces-simplicial-lemma-hypercovering-X-equivalence-bounded} 的特殊情形。
\end{proof}

\begin{lemma}
\label{spaces-simplicial-lemma-sr-when-fibre-products}
设 $U$ 是具有纤维积的站点 $\mathcal{C}$ 的单纯对象。
\begin{enumerate}
\item $\mathcal{C}/U$ 在以站点为对象、以站点态射为态射的范畴中
具有单纯对象的结构；
\item 将引理 \ref{spaces-simplicial-lemma-simplicial-site-site} 的构造应用于 (1)
中的结构，重新得到上述站点 $(\mathcal{C}/U)_{total}$；
\item 若 $a : U \to X$ 是增广，则
$a_0 : \mathcal{C}/U_0 \to \mathcal{C}/X$ 是评注
\ref{spaces-simplicial-remark-augmentation-site} (A) 部分意义下的增广，并给出与上面相同的
拓扑斯态射
$a : \Sh((\mathcal{C}/U)_{total}) \to \Sh(\mathcal{C}/X)$
。
\end{enumerate}
\end{lemma}

\begin{proof}
给定 $\mathcal{C}$ 的对象之间的态射 $V \to W$，局部化态射
$j : \mathcal{C}/V \to \mathcal{C}/W$ 是基变换函子
$\mathcal{C}/W \to \mathcal{C}/V$ 的左伴随。基变换函子连续，
并诱导与 $j$ 相同的拓扑斯态射。参见《站点》引理
\ref{sites-lemma-relocalize-given-fibre-products}。这证明了 (1)。

\medskip\noindent
(2) 成立，因为在引理 \ref{spaces-simplicial-lemma-simplicial-site-site}
所构造的范畴中，态射 $V/U_n \to W/U_m$ 是 $U_n$ 上方的态射
$V \to W \times_{U_m, f_\varphi} U_n$；这等价于位于态射
$f_\varphi : U_n \to U_m$ 上方的态射 $f : V \to W$，
也就是上面定义的范畴 $(\mathcal{C}/U)_{total}$ 中的态射。
覆盖族的集合相等立即由定义得出。

\medskip\noindent
略去 (3) 的证明。
\end{proof}







\section{站点的单纯对象给出的超覆盖：模}
\label{spaces-simplicial-section-hypercovering-modules}

\noindent
设 $\mathcal{C}$ 是具有纤维积的站点，$X \in \Ob(\mathcal{C})$，
而 $\mathcal{O}_\mathcal{C}$ 是 $\mathcal{C}$ 上的环层。
设 $U \to X$ 是第 \ref{spaces-simplicial-section-hypercovering} 节所定义的
$\mathcal{C}$ 中 $X$ 的超覆盖。本节研究如下增广：
$$
a :
(\Sh((\mathcal{C}/U)_{total}), \mathcal{O})
\longrightarrow
(\Sh(\mathcal{C}/X), \mathcal{O}_X)
$$
把 $U$ 视为 $\mathcal{C}/X$ 的单纯半可表示对象，
% P10-E0212
其次数 $n$ 部分是单元素族 $\{U_n/X\}$，再应用评注
\ref{spaces-simplicial-remark-augmentation-ringed-over-object} 中的构造，便得到该增广。
因此，本节的所有结果立即由第
\ref{spaces-simplicial-section-cohomological-descent-hypercoverings-X-modules}
节中的相应结果得出。

\begin{lemma}
\label{spaces-simplicial-lemma-hypercovering-X-simple-descent-modules}
设 $\mathcal{C}$ 是具有纤维积的站点，$X \in \Ob(\mathcal{C})$，
$\mathcal{O}_\mathcal{C}$ 是环层，且 $U$ 是 $\mathcal{C}$ 中
$X$ 的超覆盖。采用上述记号，
$$
a^* : \textit{Mod}(\mathcal{O}_X) \to \textit{Mod}(\mathcal{O})
$$
满忠实，其本质像为笛卡尔 $\mathcal{O}$-模。
函子 $a_*$ 给出拟逆。
\end{lemma}

\begin{proof}
这是引理 \ref{spaces-simplicial-lemma-hypercovering-X-descent-modules} 的特殊情形。
\end{proof}

\begin{lemma}
\label{spaces-simplicial-lemma-hypercovering-X-simple-descent-bounded-modules}
设 $\mathcal{C}$ 是具有纤维积的站点，$X \in \Ob(\mathcal{C})$，
$\mathcal{O}_\mathcal{C}$ 是环层，且 $U$ 是 $\mathcal{C}$ 中
$X$ 的超覆盖。对 $E \in D(\mathcal{O}_X)$，映射
$$
E \longrightarrow Ra_*La^*E
$$
是同构。
\end{lemma}

\begin{proof}
这是引理 \ref{spaces-simplicial-lemma-hypercovering-X-descent-bounded-modules}
的特殊情形。
\end{proof}

\begin{lemma}
\label{spaces-simplicial-lemma-compare-cohomology-hypercovering-X-simple-modules}
设 $\mathcal{C}$ 是具有纤维积的站点，$X \in \Ob(\mathcal{C})$，
$\mathcal{O}_\mathcal{C}$ 是环层，且 $U$ 是 $\mathcal{C}$ 中
$X$ 的超覆盖。则有典范同构
$$
R\Gamma(X, E) = R\Gamma((\mathcal{C}/U)_{total}, La^*E)
$$
% P10-E0211
其中 $E \in D(\mathcal{O}_\mathcal{C})$。
\end{lemma}

\begin{proof}
这是引理 \ref{spaces-simplicial-lemma-compare-cohomology-hypercovering-X-modules}
的特殊情形。
\end{proof}

\begin{lemma}
\label{spaces-simplicial-lemma-hypercovering-X-simple-equivalence-bounded-modules}
设 $\mathcal{C}$ 是具有纤维积的站点，$X \in \Ob(\mathcal{C})$，
$\mathcal{O}_\mathcal{C}$ 是环层，且 $U$ 是 $\mathcal{C}$ 中
$X$ 的超覆盖。以 $\mathcal{A} \subset \textit{Mod}(\mathcal{O})$
表示笛卡尔 $\mathcal{O}$-模组成的弱 Serre 子范畴。
则函子 $La^*$ 给出等价
$$
D^+(\mathcal{O}_X) \longrightarrow D_\mathcal{A}^+(\mathcal{O})
$$
，其拟逆为 $Ra_*$。
\end{lemma}

\begin{proof}
这是引理 \ref{spaces-simplicial-lemma-hypercovering-X-equivalence-bounded-modules}
的特殊情形。
\end{proof}







\section{超覆盖的无界上同调下降}
\label{spaces-simplicial-section-unbounded-cohomological-descent}

\noindent
本节讨论无界上同调下降。结果本身将立即由前几节关于有界上同调下降
的结果，以及《站点上的上同调》引理
\ref{sites-cohomology-lemma-equivalence-unbounded-one} 和（或）
\ref{sites-cohomology-lemma-equivalence-unbounded-two} 得出；
真正的工作在于建立记号并选择适当的假设。
我们的讨论受 \cite{six-I} 中讨论的启发，不过细节颇为不同。

\medskip\noindent
设 $(\mathcal{C}, \mathcal{O}_\mathcal{C})$ 是赋环站点。
假设对 $\mathcal{C}$ 的每个对象 $U$，给定弱 Serre 子范畴
$\mathcal{A}_U \subset \textit{Mod}(\mathcal{O}_U)$，
并满足下列性质：
\begin{enumerate}
\item
\label{spaces-simplicial-item-restriction}
给定 $\mathcal{C}$ 的态射 $U \to V$，限制函子
$\textit{Mod}(\mathcal{O}_V) \to \textit{Mod}(\mathcal{O}_U)$
将 $\mathcal{A}_V$ 映入 $\mathcal{A}_U$；
\item
\label{spaces-simplicial-item-local}
给定 $\mathcal{C}$ 的覆盖 $\{U_i \to U\}_{i \in I}$，
$\textit{Mod}(\mathcal{O}_U)$ 的对象 $\mathcal{F}$ 属于
$\mathcal{A}_U$，当且仅当对所有 $i \in I$，
$\mathcal{F}$ 在 $\mathcal{C}/U_i$ 上的限制属于
$\mathcal{A}_{U_i}$；
\item
\label{spaces-simplicial-item-bounded-dimension}
存在子集 $\mathcal{B} \subset \Ob(\mathcal{C})$，使得
\begin{enumerate}
\item $\mathcal{C}$ 的每个对象都有一个覆盖，其各成员属于
$\mathcal{B}$；
\item 对每个 $V \in \mathcal{B}$，存在整数 $d_V$ 和 $V$
的覆盖组成的余终系 $\text{Cov}_V$，使得
$$
H^p(V_i, \mathcal{F}) = 0 \text{，当 }
\{V_i \to V\} \in \text{Cov}_V,\ p > d_V, \text{且 }
\mathcal{F} \in \Ob(\mathcal{A}_V)
$$
\end{enumerate}
\end{enumerate}
注意，我们要求这一点对 $\mathcal{A}_V$ 中的 $\mathcal{F}$ 成立，
而不仅仅对“整体”对象成立（因此它强于《站点上的上同调》情形
\ref{sites-cohomology-situation-olsson-laszlo} 中施加的条件）。
在这种情形下，存在弱 Serre 子范畴
$\mathcal{A} \subset \textit{Mod}(\mathcal{O}_\mathcal{C})$，
它由如下对象组成：对所有 $U \in \Ob(\mathcal{C})$，
其在 $\mathcal{C}/U$ 上的限制属于 $\mathcal{A}_U$。
此外，还有导出范畴 $D_\mathcal{A}(\mathcal{O}_\mathcal{C})$
和 $D_{\mathcal{A}_U}(\mathcal{O}_U)$，而限制函子在它们之间映射。

\begin{example}
\label{spaces-simplicial-example-quasi-coherent-spaces-etale}
设 $S$ 是概形，$X$ 是 $S$ 上的代数空间。令
$\mathcal{C} = X_{spaces, \etale}$ 为在 $X$ 上平展的代数空间范畴
上的平展站点；参见《空间的性质》定义
\ref{spaces-properties-definition-spaces-etale-site}。
以 $\mathcal{O}_\mathcal{C}$ 表示结构层，即由规则
$U \mapsto \Gamma(U, \mathcal{O}_U)$ 给出的层。
以 $\mathcal{A}_U$ 表示拟凝聚 $\mathcal{O}_U$-模的范畴。
令 $\mathcal{B} = \Ob(\mathcal{C})$；对 $V \in \mathcal{B}$，
置 $d_V = 0$，并以 $\text{Cov}_V$ 表示对所有 $i$，
$V_i$ 均为仿射的覆盖 $\{V_i \to V\}$。
此时假设 (1)、(2)、(3) 均成立。性质 (1) 和 (2) 见
《空间的性质》引理
\ref{spaces-properties-lemma-pullback-quasi-coherent} 和
\ref{spaces-properties-lemma-properties-quasi-coherent}；
(3) 中的消灭性由《概形的上同调》引理
\ref{coherent-lemma-quasi-coherent-affine-cohomology-zero}
以及《空间的上同调》第
\ref{spaces-cohomology-section-higher-direct-image} 节中的讨论得出。
\end{example}

\begin{example}
\label{spaces-simplicial-example-etale}
设 $S$ 是下列类型之一的概形：
\begin{enumerate}
\item 有限域的谱；
\item 可分闭域的谱；
\item 严格 Hensel 的 Noether 局部环的谱；
\item 剩余域有限的 Hensel Noether 局部环的谱；
\item % P10-E0215
在此添加更多情形。
\end{enumerate}
设 $\Lambda$ 是有限环，其阶在 $S$ 上可逆。令
$\mathcal{C} \subset (\Sch/S)_\etale$ 为由 $S$ 上局部有限型概形
组成的全子范畴，并赋予平展拓扑。令
$\mathcal{O}_\mathcal{C} = \underline{\Lambda}$ 为常值层。
置 $\mathcal{A}_U = \textit{Mod}(\mathcal{O}_U)$；换言之，
我们考虑所有 $\Lambda$-模的平展层。令
$\mathcal{B} \subset \Ob(\mathcal{C})$ 为拟紧对象的集合。
对 $V \in \mathcal{B}$，置
$$
d_V = 1 + 2\dim(S) +
\sup\nolimits_{v \in V}(\text{trdeg}_{\kappa(s)}(\kappa(v)) +
2 \dim \mathcal{O}_{V, v})
$$
并以 $\text{Cov}_V$ 表示对所有 $i$，$V_i$ 均拟紧的平展覆盖
$\{V_i \to V\}$。界 $d_V$ 的这一选择来自 Gabber
关于上同调维数的定理。要看出在此选择下条件 (3) 成立，
可使用 \cite[Expos\'e VIII-A, Corollary 1.2 and Lemma 2.2]{Traveaux}
以及关于域的上同调维数的初等论证。公式中加 $1$，
是因为我们的列表包含允许 $S$ 具有有限剩余域的情形。
稍后还会回到这个例子（% P10-E0216
插入将来的引用）。
\end{example}

\noindent
设 $(\mathcal{C}, \mathcal{O}_\mathcal{C})$ 是赋环站点。
假设给定满足条件 (\ref{spaces-simplicial-item-restriction}) 的弱 Serre 子范畴
$\mathcal{A}_U \subset \textit{Mod}(\mathcal{O}_U)$。则
\begin{enumerate}
\item 给定半可表示对象 $K = \{U_i\}_{i \in I}$，取
$\prod \mathcal{A}_{U_i} \subset
\prod \textit{Mod}(\mathcal{O}_{U_i}) = \textit{Mod}(\mathcal{O}_K)$，
便得到弱 Serre 子范畴
$\mathcal{A}_K \subset \textit{Mod}(\mathcal{O}_K)$；
\item 给定半可表示对象之间的态射 $f : K \to L$，拉回映射
% P10-E0217
$f^* : \textit{Mod}(\mathcal{O}_L) \to \textit{Mod}(\mathcal{O}_L)$
将 $\mathcal{A}_L$ 映入 $\mathcal{A}_K$。
\end{enumerate}
记号和说明见评注 \ref{spaces-simplicial-remark-semi-representable-ringed}。
特别地，给定单纯半可表示对象 $K$，评注
\ref{spaces-simplicial-remark-augmentation-ringed} 中 $\textit{Mod}(\mathcal{O})$
的对象 $\mathcal{F}$ 的限制 $\mathcal{F}_n$ 对所有 $n$
均属于 $\mathcal{A}_{K_n}$，其含义是无歧义的。

\begin{lemma}
\label{spaces-simplicial-lemma-hypercovering-equivalence-modules}
设 $(\mathcal{C}, \mathcal{O}_\mathcal{C})$ 是赋环站点。
假设给定满足上述条件 (\ref{spaces-simplicial-item-restriction})、
(\ref{spaces-simplicial-item-local}) 和 (\ref{spaces-simplicial-item-bounded-dimension}) 的弱 Serre 子范畴
$\mathcal{A}_U \subset \textit{Mod}(\mathcal{O}_U)$。
假设 $\mathcal{C}$ 具有等化子与纤维积，并设 $K$ 是超覆盖。
令 $((\mathcal{C}/K)_{total}, \mathcal{O})$ 如评注
\ref{spaces-simplicial-remark-augmentation-ringed} 所述。以
$\mathcal{A}_{total} \subset \textit{Mod}(\mathcal{O})$
表示由如下笛卡尔 $\mathcal{O}$-模 $\mathcal{F}$ 组成的弱 Serre
子范畴：对所有 $n$，其限制 $\mathcal{F}_n$ 属于
$\mathcal{A}_{K_n}$（含义如上所定义）。则函子 $La^*$ 给出等价
$$
D_\mathcal{A}(\mathcal{O}_\mathcal{C})
\longrightarrow
D_{\mathcal{A}_{total}}(\mathcal{O})
$$
，其拟逆为 $Ra_*$。
\end{lemma}

\begin{proof}
由引理 \ref{spaces-simplicial-lemma-Serre-subcat-cartesian-modules}，
笛卡尔 $\mathcal{O}$-模组成弱 Serre 子范畴
（评注 \ref{spaces-simplicial-remark-augmentation-ringed} 中的讨论表明所需假设成立）。
由于限制函子
% P10-E0219
$g_n^* : \textit{Mod}(\mathcal{O}) \to \textit{Mod}(\mathcal{O}_n)$
正合，$\mathcal{A}_{total}$ 也是弱 Serre 子范畴。

\medskip\noindent
我们来证明 $a^* : \mathcal{A} \to \mathcal{A}_{total}$ 是范畴的等价，
% P10-E0218
其逆由 $La_*$ 给出。由有界版本（引理
\ref{spaces-simplicial-lemma-hypercovering-equivalence-bounded-modules}），
我们已经知道 $La_*a^*\mathcal{F} = \mathcal{F}$。
对 $\mathcal{A}$ 中的 $\mathcal{F}$，显然
$a^*\mathcal{F}$ 属于 $\mathcal{A}_{total}$。反之，假设
$\mathcal{G} \in \mathcal{A}_{total}$。因为 $\mathcal{G}$ 是笛卡尔的，
由引理 \ref{spaces-simplicial-lemma-hypercovering-descent-modules}，存在某个
$\mathcal{O}_\mathcal{C}$-模 $\mathcal{F}$，使得
$\mathcal{G} = a^*\mathcal{F}$。我们要证明 $\mathcal{F}$ 属于
$\mathcal{A}$。取 $U \in \Ob(\mathcal{C})$。须证明
$\mathcal{F}$ 在 $\mathcal{C}/U$ 上的限制属于 $\mathcal{A}_U$。
照例，写 $K_0 = \{U_{0, i}\}_{i \in I_0}$。由于 $K$ 是超覆盖，
映射 $\coprod_{i \in I_0} h_{U_{0, i}} \to *$ 在层化后成为满射。
这意味着存在覆盖 $\{U_j \to U\}_{j \in J}$、映射
$\tau : J \to I_0$，并且对每个 $j \in J$，存在态射
$\varphi_j : U_j \to U_{0, \tau(j)}$。
由于 $\mathcal{G}_0 = a_0^*\mathcal{F}$，可知
$\mathcal{F}$ 在 $\mathcal{C}/U_j$ 上的限制，等于
$\mathcal{G}_0$ 的第 $\tau(j)$ 个分量经态射
% P10-E0220
$\varphi_j : U_j \to U_{0, \tau(i)}$ 在 $\mathcal{C}/U_j$
上的限制。因此，由 (\ref{spaces-simplicial-item-restriction})，
$\mathcal{F}|_{\mathcal{C}/U_j}$ 属于 $\mathcal{A}_{U_j}$；
进而由 (\ref{spaces-simplicial-item-local})，$\mathcal{F}|_{\mathcal{C}/U}$
属于 $\mathcal{A}_U$。

\medskip\noindent
特别地，引理的陈述是有意义的。现在，引理由《站点上的上同调》引理
\ref{sites-cohomology-lemma-equivalence-unbounded-one} 得出。
假设 (1) 显然成立（见评注 \ref{spaces-simplicial-remark-augmentation-ringed}）。
假设 (2) 和 (3) 已在上一段证明。假设 (4) 立即由
(\ref{spaces-simplicial-item-bounded-dimension}) 得出。对假设 (5)，令
$\mathcal{B}_{total}$ 为站点 $(\mathcal{C}/K)_{total}$ 中满足
$U \in \mathcal{B}$ 的对象 $U/U_{n, i}$ 的集合，其中
$\mathcal{B}$ 如 (\ref{spaces-simplicial-item-bounded-dimension}) 所述。
这里使用第 \ref{spaces-simplicial-section-simplicial-semi-representable} 节对站点
$(\mathcal{C}/K)_{total}$ 的描述。此外，置
$\text{Cov}_{U/U_{n, i}}$ 等于 $\text{Cov}_U$，并置
$d_{U/U_{n, i}}$ 等于 $d_U$，其中 $\text{Cov}_U$ 与 $d_U$
由 (\ref{spaces-simplicial-item-bounded-dimension}) 给出。我们断言，在这些选择下
条件 (5) 成立。这立即由引理
\ref{spaces-simplicial-lemma-sanity-check-simplicial-semi-representable}
和如下事实得出：$\mathcal{F} \in \mathcal{A}_{total}$ 蕴含
$\mathcal{F}_n \in \mathcal{A}_{K_n}$，从而
$\mathcal{F}_{n, i} \in \mathcal{A}_{U_{n, i}}$。
（担心阿贝尔层上同调与模层上同调之间差别的读者，可参阅
《站点上的上同调》引理
\ref{sites-cohomology-lemma-cohomology-modules-abelian-agree}。）
\end{proof}










\section{粘合复形}
\label{spaces-simplicial-section-glueing-complexes}

\noindent
本节是《上同调》第 \ref{cohomology-section-glueing-complexes} 节的续篇。
目标是证明 \cite[Theorem 3.2.4]{BBD} 的一个轻微推广。
我们的方法略有不同，因为会使用第 \ref{spaces-simplicial-section-glueing} 节和第
\ref{spaces-simplicial-section-glueing-modules} 节的材料。我们还将按 \cite{six-I}
中的方法重新证明无界版本。

\medskip\noindent
给读者的建议：不妨先看引理 \ref{spaces-simplicial-lemma-bbd-glueing-cech}
的陈述及其第二个证明。

\medskip\noindent
下面给出我们所关心的情形。

\begin{situation}
\label{spaces-simplicial-situation-locally-given}
设 $(\mathcal{C}, \mathcal{O}_\mathcal{C})$ 是赋环站点。给定
\begin{enumerate}
\item 范畴 $\mathcal{B}$ 和函子 $u : \mathcal{B} \to \mathcal{C}$；
\item 对 $U \in \Ob(\mathcal{B})$，对象
$E_U$ 属于 $D(\mathcal{O}_{u(U)})$；
\item 对 $\mathcal{B}$ 中的 $a : V \to U$，同构
$\rho_a : E_U|_{\mathcal{C}/u(V)} \to E_V$
属于 $D(\mathcal{O}_{u(V)})$。
\end{enumerate}
并且每当 $\mathcal{B}$ 中有可复合箭头
$b : W \to V$ 和 $a : V \to U$ 时，都有
$\rho_{a \circ b} = \rho_b \circ \rho_a|_{\mathcal{C}/u(W)}$。
\end{situation}

\noindent
若不再加假设，我们无法由此证明任何结论。一个重要情形是：
$\mathcal{B}$ 是全子范畴，并且 $\mathcal{C}$ 的每个对象都有
一个由 $\mathcal{B}$ 的对象组成的覆盖（\cite{BBD} 考虑的正是此情形）。
对我们而言，也必须允许并非如此的情形；主要的另一种情形是：
有站点态射 $f : \mathcal{C} \to \mathcal{D}$，而 $\mathcal{B}$
是 $\mathcal{D}$ 的全子范畴，使得 $\mathcal{D}$ 的每个对象
都有一个由 $\mathcal{B}$ 的对象组成的覆盖。

\medskip\noindent
在情形 \ref{spaces-simplicial-situation-locally-given} 中，{\it 解}是二元组
$(E, \rho_U)$，其中 $E$ 是 $D(\mathcal{O}_\mathcal{C})$ 的对象，
而对 $U \in \Ob(\mathcal{B})$，
$\rho_U : E|_{\mathcal{C}/u(U)} \to E_U$ 是同构；并且对
$\mathcal{B}$ 中的 $a : V \to U$，有
$\rho_a \circ \rho_U|_{\mathcal{C}/u(V)} = \rho_V$。

\begin{lemma}
\label{spaces-simplicial-lemma-prepare-bbd-glueing}
处于情形 \ref{spaces-simplicial-situation-locally-given}。假设 $E_U$ 在
$D(\mathcal{O}_{u(U)})$ 中所有负次数自 Ext 均为零。
设 $L$ 是 $\text{SR}(\mathcal{B})$ 的单纯对象。
考虑 $\text{SR}(\mathcal{C})$ 的单纯对象 $K = u(L)$，
并令 $((\mathcal{C}/K)_{total}, \mathcal{O})$ 如评注
\ref{spaces-simplicial-remark-augmentation-ringed} 所述。存在 $D(\mathcal{O})$
的笛卡尔对象 $E$，使得写 $L_n = \{U_{n, i}\}_{i \in I_n}$ 时，
$E$ 在 $D(\mathcal{O}_{\mathcal{C}/u(U_{n, i})})$ 上的限制
相容地等于 $E_{U_{n, i}}$（细节见证明）。此外，$E$
在唯一同构的意义下唯一。
\end{lemma}

\begin{proof}
回忆
$\Sh(\mathcal{C}/K_n) = \prod_{i \in I_n} \Sh(\mathcal{C}/u(U_{n, i}))$
，模范畴也有类似分解。模层的导出范畴也继承这一乘积分解；
而且该分解与单纯半可表示对象 $K$ 中的态射相容。
参见第 \ref{spaces-simplicial-section-semi-representable} 节。因此，可以在
$D(\mathcal{O}_n)$ 中置
$E_n = \prod_{i \in I_n} E_{U_{n, i}}$（“形式”乘积）。
取情形 \ref{spaces-simplicial-situation-locally-given} 中映射 $\rho_a$ 的形式乘积，
得到同构 $E_\varphi : f_\varphi^*E_n \to E_m$。
% P10-E0221
关于映射 $\rho_a$ 之复合的假设立即蕴含 $(E_n, E_\varphi)$
按定义 \ref{spaces-simplicial-definition-cartesian-derived-modules} 构成模的导出范畴的
单纯系统。引理中假设的负次数 Ext 消灭性蕴含：对
$n \geq 0$ 和 $t > 0$，有 $\Hom(E_n[t], E_n) = 0$。
因此，由引理 \ref{spaces-simplicial-lemma-cartesian-module-derived-from-simplicial}，
得到 $E$。在唯一同构的意义下的唯一性由引理
\ref{spaces-simplicial-lemma-nullity-cartesian-modules-derived} 和
\ref{spaces-simplicial-lemma-hom-cartesian-modules-derived} 得出。
\end{proof}

\begin{lemma}[BBD 粘合引理]
\label{spaces-simplicial-lemma-bbd-glueing}
处于情形 \ref{spaces-simplicial-situation-locally-given}。假设
\begin{enumerate}
\item $\mathcal{C}$ 具有等化子与纤维积；
\item 存在由连续函子 $u : \mathcal{D} \to \mathcal{C}$
给出的站点态射 $f : \mathcal{C} \to \mathcal{D}$，使得
\begin{enumerate}
\item $\mathcal{D}$ 具有等化子与纤维积，且 $u$ 与它们交换；
\item $\mathcal{B}$ 是 $\mathcal{D}$ 的全子范畴，而
$u : \mathcal{B} \to \mathcal{C}$ 是 $u$ 的限制；
\item $\mathcal{D}$ 的每个对象都有一个由 $\mathcal{B}$
的对象组成的覆盖；
\end{enumerate}
\item 对 $\mathcal{B}$ 中的所有 $U$，$E_U$ 在
$D(\mathcal{O}_{u(U)})$ 中所有负次数自 Ext 均为零；
\item 存在 $t \in \mathbf{Z}$，使得当 $i < t$ 且
$U \in \Ob(\mathcal{B})$ 时，$H^i(E_U) = 0$。
\end{enumerate}
则存在一个解，并且它在唯一同构的意义下唯一。
\end{lemma}

\begin{proof}
由《超覆盖》引理 \ref{hypercovering-lemma-hypercovering-site}，
存在站点 $\mathcal{D}$ 的超覆盖 $L$，使得
$L_n = \{U_{n, i}\}_{i \in I_n}$，其中
% P10-E0222
$U_{i, n} \in \Ob(\mathcal{B})$。置 $K = u(L)$。
应用引理 \ref{spaces-simplicial-lemma-prepare-bbd-glueing}，得到站点
$(\mathcal{C}/K)_{total}$ 上 $D(\mathcal{O})$ 的笛卡尔对象 $E$，
它在 $\mathcal{C}/u(U_{n, i})$ 上相容地限制为 $E_{U_{n, i}}$。
关于 $t$ 的假设蕴含 $E \in D^+(\mathcal{O})$。
由《超覆盖》引理
\ref{hypercovering-lemma-hypercovering-morphism-sites}，
$K$ 也是超覆盖。由引理
\ref{spaces-simplicial-lemma-hypercovering-equivalence-bounded-modules}，
存在 $D^+(\mathcal{O}_\mathcal{C})$ 中的 $F$，使得 $E = a^*F$。

\medskip\noindent
为了证明 $F$ 是解，我们将使用《超覆盖》引理
\ref{hypercovering-lemma-hypercovering-site} 的证明中给出的
$L_0$ 和 $L_1$ 的构造。（这略显不雅，但似乎没有完全直接的绕行办法。）

\medskip\noindent
具体地，$I_0 = \Ob(\mathcal{B})$，所以
$L_0 = \{U\}_{U \in \Ob(\mathcal{B})}$。因此，按我们对 $E$
的选择，将同构 $a^*F \to E$ 限制到 $\mathcal{C}/K_0$
的分量 $\mathcal{C}/u(U)$，便对 $U \in \Ob(\mathcal{B})$
定义了同构 $\rho_U : F|_{\mathcal{C}/u(U)} \to E_U$。

\medskip\noindent
% P10-E0224
为了证明 $\rho_U$ 满足与情形 \ref{spaces-simplicial-situation-locally-given}
中的映射 $\rho_a$ 相容这一要求，我们使用 $I_1$ 包含集合
$$
\Omega =
\{(U, V, W, a, b) \mid U, V, W \in \mathcal{B}, a : U \to V, b : U \to W\}
$$
并且对 $\Omega$ 中的 $i = (U, V, W, a, b)$，有
$U_{1, i} = U$。此外，两个态射 $K_1 \to K_0$ 的分量映射
$f_{\delta^1_0, i}$ 和 $f_{\delta^1_1, i}$ 是态射
% P10-E0223
$$
a : U \to V \quad\text{且}\quad b : U \to V
$$
因此，引理 \ref{spaces-simplicial-lemma-prepare-bbd-glueing} 中所述的相容性给出
$$
\rho_a \circ \rho_V|_{\mathcal{C}/u(U)} = \rho_U
\quad\text{且}\quad
\rho_b \circ \rho_W|_{\mathcal{C}/u(U)} = \rho_U
$$
例如取 $i = (U, V, U, a, \text{id}_U) \in \Omega$，
就得到所需的相容性。$F$ 的唯一性由前一引理中 $E$
的唯一性得出（略去一个小细节）。
\end{proof}

\begin{lemma}[无界 BBD 粘合引理]
\label{spaces-simplicial-lemma-bbd-unbounded-glueing}
处于情形 \ref{spaces-simplicial-situation-locally-given}。假设
\begin{enumerate}
\item $\mathcal{C}$ 具有等化子与纤维积；
\item 存在由连续函子 $u : \mathcal{D} \to \mathcal{C}$
给出的站点态射 $f : \mathcal{C} \to \mathcal{D}$，使得
\begin{enumerate}
\item $\mathcal{D}$ 具有等化子与纤维积，且 $u$ 与它们交换；
\item $\mathcal{B}$ 是 $\mathcal{D}$ 的全子范畴，而
$u : \mathcal{B} \to \mathcal{C}$ 是 $u$ 的限制；
\item $\mathcal{D}$ 的每个对象都有一个由 $\mathcal{B}$
的对象组成的覆盖；
\end{enumerate}
\item $E_U$ 在 $D(\mathcal{O}_{u(U)})$ 中所有负次数自 Ext 均为零；
\item 对所有 $U \in \Ob(\mathcal{C})$，存在满足条件
(\ref{spaces-simplicial-item-restriction})、(\ref{spaces-simplicial-item-local}) 和
(\ref{spaces-simplicial-item-bounded-dimension}) 的弱 Serre 子范畴
$\mathcal{A}_U \subset \textit{Mod}(\mathcal{O}_U)$；
\item % P10-E0225
$E_U \in D_{\mathcal{A}_U}(\mathcal{O}_U)$。
\end{enumerate}
则存在一个解，并且它在唯一同构的意义下唯一。
\end{lemma}

\begin{proof}
证明与引理 \ref{spaces-simplicial-lemma-bbd-glueing} 的证明{\bf 完全}相同。
唯一的改变是 $E$ 为 $D_{\mathcal{A}_{total}}(\mathcal{O})$ 的对象，
所以我们使用引理 \ref{spaces-simplicial-lemma-hypercovering-equivalence-modules}
得到满足 $E = a^*F$ 的 $F$，而不是使用引理
\ref{spaces-simplicial-lemma-hypercovering-equivalence-bounded-modules}。
\end{proof}

\noindent
下面给出上述一般理论的一个应用实例。

\begin{lemma}
\label{spaces-simplicial-lemma-bbd-glueing-cech}
\begin{reference}
Martin Olsson 于 2021 年 9 月 9 日发出的电子邮件。
\end{reference}
设 $(\mathcal{C}, \mathcal{O}_\mathcal{C})$ 是赋环站点，
并假设 $\mathcal{C}$ 具有纤维积。设
$\{U_i \to X\}_{i \in I}$ 是 $\mathcal{C}$ 中的覆盖。
对 $i \in I$，设 $E_i$ 是 $D(\mathcal{O}_{U_i})$ 的对象；
对 $i, j \in I$，设
$$
\rho_{ij} : E_i|_{\mathcal{C}/U_{ij}} \longrightarrow E_j|_{\mathcal{C}/U_{ij}}
$$
为 $D(\mathcal{O}_{U_{ij}})$ 中的同构，其中
$U_{ij} = U_i \times_X U_j$。假设
\begin{enumerate}
\item 对所有 $i, j, k \in I$，$\rho_{ij}$ 在
$U_i \times_X U_j \times_X U_k$ 上满足余圈条件；
\item $\SheafExt^p_{\mathcal{O}_{U_i}}(E_i, E_i) = 0$
对所有 $p < 0$ 和 $i \in I$ 成立；
\item 存在 $t \in \mathbf{Z}$，使得对 $p < t$ 及所有
$i \in I$，都有 $H^p(E_i) = 0$。
\end{enumerate}
则存在唯一的二元组 $(E, \rho_i)$，其中 $E$ 是
$D(\mathcal{O}_X)$ 的对象，而
$\rho_i : E|_{U_i} \to E_i$ 是 $D(\mathcal{O}_{U_i})$
中与 $\rho_{ij}$ 相容的同构。
\end{lemma}

\begin{proof}[第一个证明]
本证明由十分一般的引理 \ref{spaces-simplicial-lemma-bbd-glueing} 推出该引理。
% P10-E0226
我们建议读者改看第二个证明。

\medskip\noindent
可以用 $\mathcal{C}/X$ 替换 $\mathcal{C}$。
因此，我们可以并且确实假设 $X$ 是 $\mathcal{C}$ 的终对象，
且 $\mathcal{C}$ 具有所有有限极限。

\medskip\noindent
令 $\mathcal{B}$ 为 $\mathcal{C}$ 的全子范畴，其对象是满足下列条件的
$U \in \Ob(\mathcal{C})$：存在 $i(U) \in I$ 和态射
$a_U : U \to U_{i(U)}$。在 $D(\mathcal{O}_U)$ 中，以
$E_U = a_U^*E_{i(U)}$ 表示 $E_i$ 经 $a_U$ 的拉回（限制）。
% P10-E0227
给定 $\mathcal{B}$ 的态射 $a : U \to U'$，得到态射
$(a_{U'} \circ a, a_U) : U \to U_{i(U')} \times_X U_{i(U)} = U_{i(U')i(U)}$
，从而得到 $D(\mathcal{O}_U)$ 中的同构
$$
\rho_a :
a^*E_{U'} = a^*a_{U'}^*E_{i(U')}
\xrightarrow{(a_{U'} \circ a, a_U)^*\rho_{i(U')i(U)}}
a_{U}^*E_{i(U)} = E_{U}
$$
。数据 $\mathcal{B}, E_U, \rho_a$ 如情形
\ref{spaces-simplicial-situation-locally-given} 所述；同构 $\rho_a$
恰因引理陈述中的条件 (1) 而满足余圈条件（略去细节）。

\medskip\noindent
我们将把引理 \ref{spaces-simplicial-lemma-bbd-glueing} 应用于上述
$\mathcal{B}, E_U, \rho_a$，并取 $\mathcal{D} = \mathcal{C}$、
$f : \mathcal{C} \to \mathcal{D}$ 为恒同态射。
引理 \ref{spaces-simplicial-lemma-bbd-glueing} 的假设 (1) 和 (2)(a) 已在上面看到。
引理 \ref{spaces-simplicial-lemma-bbd-glueing} 的假设 (2)(b) 显然成立。
引理 \ref{spaces-simplicial-lemma-bbd-glueing} 的假设 (2)(c) 成立，因为
$\{U_i \to X\}$ 是覆盖\footnote{事实上，只需映射
$\coprod_{i \in I} h_{U_i} \to h_X$ 在层化后成为满射；
在这种情形下，该引理以同一证明成立。}。
引理 \ref{spaces-simplicial-lemma-bbd-glueing} 的假设 (3) 成立，因为我们已假设
$E_i$ 的所有负次数 Ext 层均消灭；这当然蕴含对位于 $U_i$
上方的任意对象 $U$，$E_i|_U$ 的负次数自 Ext 均为零。
引理 \ref{spaces-simplicial-lemma-bbd-glueing} 的假设 (4) 成立，因为我们已假设
每个 $E_i$ 的上同调层在 $t$ 的左侧为零。

\medskip\noindent
我们得到唯一的解 $(E, \rho_U)$。置 $\rho_i = \rho_{U_i}$，
便得到引理。
\end{proof}

\begin{proof}[第二个证明]
我们概述一个更直接的证明。以 $K$ 表示与覆盖
$\{U_i \to X\}_{i \in I}$ 关联的 $X$ 的 {\v C}ech 超覆盖；
参见《超覆盖》例 \ref{hypercovering-example-cech}。
例如，$K_0 = \{U_i \to X\}_{i \in I}$，而
$K_1 = \{U_i \times_X U_j \to X\}_{i, j \in I}$，依此类推。
令 $((\mathcal{C}/K)_{total}, \mathcal{O})$、$a$、$a_n$
如评注 \ref{spaces-simplicial-remark-augmentation-ringed-over-object} 所述。
对象 $E_i$ 在 $D(\mathcal{O}_0) = \prod D(\mathcal{O}_{U_i})$
中确定对象 $M_0$。类似地，同构 $\rho_{ij}$ 确定同构
$$
\alpha :
L(f_{\delta_1^1})^*M_0
\longrightarrow
L(f_{\delta_0^1})^*M_0
$$
，它属于 $D(\mathcal{O}_1)$ 并满足余圈条件。
由引理 \ref{spaces-simplicial-lemma-construct-cartesian-objects-derived-modules}，
得到导出范畴的笛卡尔单纯系统 $(M_n)$。
由所假设的负次数 Ext 层消灭性，可知对象 $M_n$
的负次数自 Ext 均消灭。因此，由引理
\ref{spaces-simplicial-lemma-cartesian-module-derived-from-simplicial}，得到
$D(\mathcal{O})$ 的笛卡尔对象 $M$，其关联单纯系统同构于 $(M_n)$。
由于 $M$ 的上同调层在次数 $< t$ 处为零，由引理
\ref{spaces-simplicial-lemma-hypercovering-X-equivalence-bounded-modules}，存在
$D(\mathcal{O}_X)$ 中的 $E$，使得 $M = La^*E$。
将同构 $La^*E \to M$ 限制到 $\mathcal{C}/U_i$，
便得到同构 $\rho_i$。略去它们与同构 $\rho_{ij}$ 相容性的验证。
\end{proof}












\section{拓扑中的真超覆盖}
\label{spaces-simplicial-section-proper-hypercovering}

\noindent
在由 Hausdorff 局部拟紧拓扑空间与连续映射组成的范畴
$\textit{LC}$ 中工作；参见《站点上的上同调》第
\ref{sites-cohomology-section-cohomology-LC} 节。
设 $X$ 是 $\textit{LC}$ 的对象，$U$ 是 $\textit{LC}$
的单纯对象。假设有增广
$$
a : U \to X
$$
如果满足下列条件，就称 $U$ 是 $X$ 的{\it 真超覆盖}：
\begin{enumerate}
\item $U_0 \to X$ 是真满射；
\item $U_1 \to U_0 \times_X U_0$ 是真满射；
\item $U_{n + 1} \to (\text{cosk}_n\text{sk}_n U)_{n + 1}$
对 $n \geq 1$ 是真满射。
\end{enumerate}
$\textit{LC}$ 具有所有有限极限，所以上述表述中使用的余骨架存在。
$$
\fbox{原理：可以用真超覆盖计算上同调。}
$$
证明该原理的一个关键思想，是在 $\textit{LC}$ 上找到一个比通常拓扑
更细的拓扑，使得：(a) 真满射定义覆盖；(b) 通常层相对于这个更细
拓扑的上同调与通常上同调一致。qc 拓扑满足性质 (a) 和 (b)；
参见《站点上的上同调》第
\ref{sites-cohomology-section-cohomology-LC} 节。
一旦有了 (a) 和 (b)，便可由本章前面的工作推出该原理。

\begin{lemma}
\label{spaces-simplicial-lemma-compare-simplicial-objects}
设 $U$ 是 $\textit{LC}$ 的单纯对象，$a : U \to X$ 是增广。
则有交换图
$$
\xymatrix{
\Sh((\textit{LC}_{qc}/U)_{total}) \ar[r]_-h \ar[d]_{a_{qc}} &
\Sh(U_{Zar}) \ar[d]^a \\
\Sh(\textit{LC}_{qc}/X) \ar[r]^-{h_{-1}} &
\Sh(X)
}
$$
其中左侧竖直箭头定义于第 \ref{spaces-simplicial-section-hypercovering} 节，
右侧竖直箭头定义于引理 \ref{spaces-simplicial-lemma-augmentation}。
\end{lemma}

\begin{proof}
写 $\Sh(X) = \Sh(X_{Zar})$。注意，
$(\textit{LC}_{qc}/U)_{total}$ 与 $U_{Zar}$ 都属于情形
\ref{spaces-simplicial-situation-simplicial-site} 的 A 类。对 $U_{Zar}$，
这立即由第 \ref{spaces-simplicial-section-simplicial-top} 节中的构造得出；
对 $(\textit{LC}_{qc}/U)_{total}$，则由引理
\ref{spaces-simplicial-lemma-sr-when-fibre-products} 得出。接着考虑函子
$U_{n, Zar} \to \textit{LC}_{qc}/U_n$, $U \mapsto U/U_n$
以及
$X_{Zar} \to \textit{LC}_{qc}/X$, $U \mapsto U/X$.
我们已在《站点上的上同调》第
\ref{sites-cohomology-section-cohomology-LC} 节看到，
它们定义站点态射。因此，得到评注
\ref{spaces-simplicial-remark-morphism-augmentation-simplicial-sites}
意义下与增广相容的单纯站点态射；应用引理
\ref{spaces-simplicial-lemma-morphism-augmentation-simplicial-sites} 即得结论。
\end{proof}

\begin{lemma}
\label{spaces-simplicial-lemma-descent-sheaves-for-proper-hypercovering}
设 $U$ 是 $\textit{LC}$ 的单纯对象，$a : U \to X$ 是增广。
若 $a : U \to X$ 给出 $X$ 的真超覆盖，则
$$
a^{-1} : \Sh(X) \to \Sh(U_{Zar})
\quad\text{且}\quad
a^{-1} : \textit{Ab}(X) \to \textit{Ab}(U_{Zar})
$$
满忠实，其本质像为笛卡尔层，拟逆由 $a_*$ 给出。
这里 $a : \Sh(U_{Zar}) \to \Sh(X)$ 如引理
\ref{spaces-simplicial-lemma-augmentation} 所述。
\end{lemma}

\begin{proof}
我们证明集合层的陈述。它几乎是已建立结果的形式推论。
考虑引理 \ref{spaces-simplicial-lemma-compare-simplicial-objects} 的图。
由《站点上的上同调》引理
\ref{sites-cohomology-lemma-describe-pullback-pi}，函子
$(h_{-1})^{-1}$ 满忠实，其拟逆为 $h_{-1, *}$。
对 $h$ 的分量 $h_n$ 也同样成立。由引理
\ref{spaces-simplicial-lemma-morphism-simplicial-sites} 对函子 $h^{-1}$ 和 $h_*$
的描述，可知 $h^{-1}$ 满忠实，其拟逆为 $h_*$。
$U$ 是 $\textit{LC}_{qc}$ 中 $X$ 的超覆盖
（定义见第 \ref{spaces-simplicial-section-hypercovering} 节）；这由《站点上的上同调》引理
\ref{sites-cohomology-lemma-proper-surjective-is-qc-covering} 得出。
由引理 \ref{spaces-simplicial-lemma-hypercovering-X-simple-descent-sheaves}，
$a_{qc}^{-1}$ 满忠实，其拟逆为 $a_{qc, *}$，本质像为
$(\textit{LC}_{qc}/U)_{total}$ 上的笛卡尔层。
现在，一个形式论证（沿图追逐）表明 $a^{-1}$ 满忠实。

\medskip\noindent
最后，假设 $\mathcal{G}$ 是 $U_{Zar}$ 上的笛卡尔层。
则 $h^{-1}\mathcal{G}$ 是 $\textit{LC}_{qc}/U$ 上的笛卡尔层。
因此，对 $\textit{LC}_{qc}/X$ 上的某个层 $\mathcal{H}$，有
$h^{-1}\mathcal{G} = a_{qc}^{-1}\mathcal{H}$。计算得
\begin{align*}
(h_{-1})^{-1}(a_*\mathcal{G})
& =
(h_{-1})^{-1}
\text{等化子}(
\xymatrix{
a_{0, *}\mathcal{G}_0
\ar@<1ex>[r] \ar@<-1ex>[r] &
a_{1, *}\mathcal{G}_1
}
) \\
& =
\text{等化子}(
\xymatrix{
(h_{-1})^{-1}a_{0, *}\mathcal{G}_0
\ar@<1ex>[r] \ar@<-1ex>[r] &
(h_{-1})^{-1}a_{1, *}\mathcal{G}_1
}
) \\
& =
\text{等化子}(
\xymatrix{
a_{qc, 0, *}h_0^{-1}\mathcal{G}_0
\ar@<1ex>[r] \ar@<-1ex>[r] &
a_{qc, 1, *}h_1^{-1}\mathcal{G}_1
}
) \\
& =
\text{等化子}(
\xymatrix{
a_{qc, 0, *}a_{qc, 0}^{-1}\mathcal{H}
\ar@<1ex>[r] \ar@<-1ex>[r] &
a_{qc, 1, *}a_{qc, 1}^{-1}\mathcal{H}
}
) \\
& =
a_{qc, *}a_{qc}^{-1}\mathcal{H} \\
& =
\mathcal{H}
\end{align*}
第一个等号由引理 \ref{spaces-simplicial-lemma-augmentation} 得出；第二个等号来自
$(h_{-1})^{-1}$ 是正合函子；第三个等号由《站点上的上同调》引理
\ref{sites-cohomology-lemma-push-pull-LC} 得出
（此处使用 $a_0 : U_0 \to X$ 和 $a_1: U_1 \to X$ 为真态射）；
第四个等号来自 $a_{qc}^{-1}\mathcal{H} = h^{-1}\mathcal{G}$；
第五个等号来自引理 \ref{spaces-simplicial-lemma-augmentation-site}；
第六个等号已在上面看到。由于
$a_{qc}^{-1}\mathcal{H} = h^{-1}\mathcal{G}$，可得
$h^{-1}\mathcal{G} \cong h^{-1}a^{-1}a_*\mathcal{G}$；
再利用 $h^{-1}$ 的满忠实性，证明即告完成。
\end{proof}

\begin{lemma}
\label{spaces-simplicial-lemma-cohomological-descent-for-proper-hypercovering}
设 $U$ 是 $\textit{LC}$ 的单纯对象，$a : U \to X$ 是增广。
若 $a : U \to X$ 给出 $X$ 的真超覆盖，则对 $K \in D^+(X)$，
$$
K \to Ra_*(a^{-1}K)
$$
是同构，其中 $a : \Sh(U_{Zar}) \to \Sh(X)$ 如引理
\ref{spaces-simplicial-lemma-augmentation} 所述。
\end{lemma}

\begin{proof}
考虑引理 \ref{spaces-simplicial-lemma-compare-simplicial-objects} 的图。
由《站点上的上同调》引理
\ref{sites-cohomology-lemma-cohomological-descent-LC}，
$Rh_{n, *}h_n^{-1}$ 是 $D^+(U_n)$ 上的恒同函子。
因此，由引理 \ref{spaces-simplicial-lemma-direct-image-morphism-simplicial-sites}，
$Rh_*h^{-1}$ 是 $D^+(U_{Zar})$ 上的恒同函子。我们有
\begin{align*}
Ra_*(a^{-1}K)
& =
Ra_*Rh_*h^{-1}a^{-1}K \\
& =
Rh_{-1, *}Ra_{qc, *}a_{qc}^{-1}(h_{-1})^{-1}K \\
& =
Rh_{-1, *}(h_{-1})^{-1}K \\
& =
K
\end{align*}
% P10-E0228
第一个等号来自上述讨论；第二个等号来自引理
\ref{spaces-simplicial-lemma-compare-simplicial-objects} 中图的交换性；
第三个等号来自引理
\ref{spaces-simplicial-lemma-hypercovering-X-simple-descent-bounded-abelian}
（由《站点上的上同调》引理
\ref{sites-cohomology-lemma-proper-surjective-is-qc-covering}，
$U$ 是 $\textit{LC}_{qc}$ 中 $X$ 的超覆盖）；最后一个等号
来自已经使用过的《站点上的上同调》引理
\ref{sites-cohomology-lemma-cohomological-descent-LC}。
\end{proof}

\begin{lemma}
\label{spaces-simplicial-lemma-compute-via-proper-hypercovering}
设 $U$ 是 $\textit{LC}$ 的单纯对象，$a : U \to X$ 是增广。
若 $U$ 是 $X$ 的真超覆盖，则
$$
R\Gamma(X, K) = R\Gamma(U_{Zar}, a^{-1}K)
$$
其中 $K \in D^+(X)$，而 $a : \Sh(U_{Zar}) \to \Sh(X)$
如引理 \ref{spaces-simplicial-lemma-augmentation} 所述。
\end{lemma}

\begin{proof}
这由引理 \ref{spaces-simplicial-lemma-cohomological-descent-for-proper-hypercovering} 得出，
因为按《站点上的上同调》评注
\ref{sites-cohomology-remark-before-Leray}，有
$R\Gamma(U_{Zar}, -) = R\Gamma(X, -) \circ Ra_*$。
\end{proof}

\begin{lemma}
\label{spaces-simplicial-lemma-proper-hypercovering-equivalence-bounded}
设 $U$ 是 $\textit{LC}$ 的单纯对象，$a : U \to X$ 是增广。
以 $\mathcal{A} \subset \textit{Ab}(U_{Zar})$
表示笛卡尔阿贝尔层组成的弱 Serre 子范畴。
若 $U$ 是 $X$ 的真超覆盖，则函子 $a^{-1}$ 给出等价
$$
D^+(X) \longrightarrow D_\mathcal{A}^+(U_{Zar})
$$
，其拟逆为 $Ra_*$；这里 $a : \Sh(U_{Zar}) \to \Sh(X)$
如引理 \ref{spaces-simplicial-lemma-augmentation} 所述。
\end{lemma}

\begin{proof}
由引理 \ref{spaces-simplicial-lemma-Serre-subcat-cartesian-modules}，
$\mathcal{A}$ 是弱 Serre 子范畴。该等价是迄今所得结果的形式推论。
应用引理 \ref{spaces-simplicial-lemma-descent-sheaves-for-proper-hypercovering}、
\ref{spaces-simplicial-lemma-cohomological-descent-for-proper-hypercovering} 以及
《站点上的上同调》引理 \ref{sites-cohomology-lemma-equivalence-bounded}。
\end{proof}

\begin{lemma}
\label{spaces-simplicial-lemma-spectral-sequence-proper-hypercovering}
设 $U$ 是 $\textit{LC}$ 的单纯对象，$a : U \to X$ 是增广。
设 $\mathcal{F}$ 是 $X$ 上的阿贝尔层，$\mathcal{F}_n$
是它到 $U_n$ 的拉回。若 $U$ 是 $X$ 的真超覆盖，
则存在典范谱序列
$$
E_1^{p, q} = H^q(U_p, \mathcal{F}_p)
$$
，它收敛到 $H^{p + q}(X, \mathcal{F})$。
\end{lemma}

\begin{proof}
这是引理 \ref{spaces-simplicial-lemma-compute-via-proper-hypercovering} 和
\ref{spaces-simplicial-lemma-simplicial-sheaf-cohomology} 的立即推论。
\end{proof}




\section{单纯概形}
\label{spaces-simplicial-section-simplicial}

\noindent
{\it 单纯概形}是概形范畴中的单纯对象；参见《单纯对象》定义
\ref{simplicial-definition-simplicial-object}。回忆，单纯概形形如
$$
\xymatrix{
X_2
\ar@<2ex>[r]
\ar@<0ex>[r]
\ar@<-2ex>[r]
&
X_1
\ar@<1ex>[r]
\ar@<-1ex>[r]
\ar@<1ex>[l]
\ar@<-1ex>[l]
&
X_0
\ar@<0ex>[l]
}
$$
这里有两个态射 $d^1_0, d^1_1 : X_1 \to X_0$ 和一个态射
$s^0_0 : X_0 \to X_1$，依此类推。这些态射满足一些必需关系，
例如 $d^1_0 \circ s^0_0 = \text{id}_{X_0} = d^1_1 \circ s^0_0$；
参见《单纯对象》引理
\ref{simplicial-lemma-characterize-simplicial-object}。
把 $d^n_i : X_n \to X_{n - 1}$ 看作“忘掉第 $i$ 个坐标的投影”，
把 $s^n_j : X_n \to X_{n + 1}$ 看作“重复第 $j$ 个坐标的对角映射”，
会很有帮助。

\medskip\noindent
{\it 单纯概形的态射} $h : X \to Y$，就是概形范畴中单纯对象的态射；
参见《单纯对象》定义 \ref{simplicial-definition-simplicial-object}。
因此，$h$ 由概形态射 $h_n : X_n \to Y_n$ 组成，并且只要等式有意义，
就有 $h_{n - 1} \circ d^n_j = d^n_j \circ h_n$ 和
$h_{n + 1} \circ s^n_j = s^n_j \circ h_n$。

\medskip\noindent
单纯概形 $X$ 的{\it 增广}是概形态射 $a_0 : X_0 \to S$，
满足 $a_0 \circ d^1_0 = a_0 \circ d^1_1$。
参见《单纯对象》第 \ref{simplicial-section-augmentation} 节。

\medskip\noindent
设 $X$ 是单纯概形。将第 \ref{spaces-simplicial-section-simplicial-top} 节的构造
应用于其底单纯拓扑空间，得到站点 $X_{Zar}$。另一方面，
对每个 $n$ 有小 Zariski 站点 $X_{n, Zar}$（《拓扑》定义
\ref{topologies-definition-big-small-Zariski}）；对每个态射
$\varphi : [m] \to [n]$，有与概形态射
$X(\varphi) : X_n \to X_m$ 关联的站点态射
$f_\varphi = X(\varphi)_{small} : X_{n, Zar} \to X_{m, Zar}$
（《拓扑》引理 \ref{topologies-lemma-morphism-big-small}）。
这给出站点范畴中的单纯对象 $\mathcal{C}$。在引理
\ref{spaces-simplicial-lemma-simplicial-site-site} 中，我们构造了关联站点
$\mathcal{C}_{total}$。把开浸入映到其像，定义了等价
$\mathcal{C}_{total} \to X_{Zar}$，并在此等价下识别层；即
$\Sh(\mathcal{C}_{total}) = \Sh(X_{Zar})$。
$\mathcal{C}_{total}$ 与 $X_{Zar}$ 之间的差别，类似于小 Zariski
站点 $S_{Zar}$ 与 $S$ 的底拓扑空间之间的差别。下文默认识别这些站点。

\medskip\noindent
设 $X_{Zar}$ 是与单纯概形 $X$ 关联的站点。
$X_{Zar}$ 上存在环层 $\mathcal{O}$，它在 $X_n$ 上的限制
是结构层 $\mathcal{O}_{X_n}$。这由引理
\ref{spaces-simplicial-lemma-describe-sheaves-simplicial-site} 或引理
\ref{spaces-simplicial-lemma-describe-sheaves-simplicial-site-site} 得出。
我们称{\it $\mathcal{O}$ 是单纯概形 $X$ 的结构层}。
至此，为单纯（赋环）站点发展的一切材料都适用；参见第
\ref{spaces-simplicial-section-simplicial-sites}、
\ref{spaces-simplicial-section-augmentation-simplicial-sites}、
\ref{spaces-simplicial-section-morphism-simplicial-sites}、
\ref{spaces-simplicial-section-simplicial-sites-modules}、
\ref{spaces-simplicial-section-cohomology-simplicial-sites}、
\ref{spaces-simplicial-section-cohomology-augmentation-simplicial-sites}、
\ref{spaces-simplicial-section-cohomology-simplicial-sites-modules}、
\ref{spaces-simplicial-section-cohomology-augmentation-ringed-simplicial-sites}、
\ref{spaces-simplicial-section-cartesian}、
\ref{spaces-simplicial-section-glueing} 和
\ref{spaces-simplicial-section-glueing-modules} 节。

\medskip\noindent
设 $X$ 是结构层为 $\mathcal{O}$ 的单纯概形。
与任何赋环拓扑斯上一样，可以定义
{\it $X_{Zar}$ 上的拟凝聚 $\mathcal{O}$-模}；参见
《站点上的模》定义 \ref{sites-modules-definition-site-local}。
不过，$X_{Zar}$ 上的拟凝聚 $\mathcal{O}$-模，恰是如下笛卡尔
$\mathcal{O}$-模 $\mathcal{F}$：其限制 $\mathcal{F}_n$
在 $X_n$ 上拟凝聚；参见引理 \ref{spaces-simplicial-lemma-quasi-coherent-sheaf}。

\medskip\noindent
设 $h : X \to Y$ 是单纯概形的态射。由引理
\ref{spaces-simplicial-lemma-simplicial-space-site-functorial}，或由引理
\ref{spaces-simplicial-lemma-morphism-simplicial-sites} 的证明，得到站点态射
$h_{Zar} : X_{Zar} \to Y_{Zar}$。回忆，$h_{Zar}^{-1}$ 和
$h_{Zar, *}$ 可用各分量简单描述；参见引理
\ref{spaces-simplicial-lemma-describe-functoriality} 或引理
\ref{spaces-simplicial-lemma-morphism-simplicial-sites}。
以 $\mathcal{O}_X$、$\mathcal{O}_Y$ 分别表示 $X$、$Y$ 的结构层。
定义
% P10-E0229
$h_{Zar}^\sharp : h_{Zar, *}\mathcal{O}_X \to \mathcal{O}_Y$
为 $Y_{Zar}$ 上的环层映射，它在 $Y_n$ 上由
$h_n^\sharp : h_{n, *}\mathcal{O}_{X_n} \to \mathcal{O}_{Y_n}$ 给出。
由此得到赋环站点的态射
$$
h_{Zar} : (X_{Zar}, \mathcal{O}_X) \longrightarrow (Y_{Zar}, \mathcal{O}_Y)
$$

\medskip\noindent
设 $X$ 是结构层为 $\mathcal{O}$ 的单纯概形。设 $S$ 是概形，
$a_0 : X_0 \to S$ 是 $X$ 的增广。由引理
\ref{spaces-simplicial-lemma-augmentation} 或引理 \ref{spaces-simplicial-lemma-augmentation-site}，
得到相应拓扑斯态射 $a : \Sh(X_{Zar}) \to \Sh(S)$。
注意，$a^{-1}\mathcal{G}$ 是 $X_{Zar}$ 上分量为
$a_n^{-1}\mathcal{G}$ 的层。因此，可以利用映射
$a_n^\sharp : a_n^{-1}\mathcal{O}_S \to \mathcal{O}_{X_n}$
定义映射 $a^\sharp : a^{-1}\mathcal{O}_S \to \mathcal{O}$；
等价地，由伴随定义映射
$a^\sharp : \mathcal{O}_S \to a_*\mathcal{O}$
（照例仍用同一名称）。这正处于第
\ref{spaces-simplicial-section-cohomology-augmentation-ringed-simplicial-sites}
节所讨论的情形。因此，得到赋环拓扑斯的态射
$$
a : (\Sh(X_{Zar}), \mathcal{O}) \longrightarrow (\Sh(S), \mathcal{O}_S)
$$

\medskip\noindent
最后作如下观察。设给定单纯概形 $X$ 与 $Y$ 的态射
$h : X \to Y$，其结构层分别为 $\mathcal{O}_X$、$\mathcal{O}_Y$；
并给定增广 $a_0 : X_0 \to X_{-1}$、$b_0 : Y_0 \to Y_{-1}$
和态射 $h_{-1} : X_{-1} \to Y_{-1}$，使得
$$
\xymatrix{
X_0 \ar[r]_{h_0} \ar[d]_{a_0} & Y_0 \ar[d]^{b_0} \\
X_{-1} \ar[r]^{h_{-1}} & Y_{-1}
}
$$
交换。由上面阐明的构造，得到如下赋环拓扑斯态射的交换图：
$$
\xymatrix{
(\Sh(X_{Zar}), \mathcal{O}_X) \ar[r]_{h_{Zar}} \ar[d]_a &
(\Sh(Y_{Zar}), \mathcal{O}_Y) \ar[d]^b \\
(\Sh(X_{-1}), \mathcal{O}_{X_{-1}}) \ar[r]^{h_{-1}} &
(\Sh(Y_{-1}), \mathcal{O}_{Y_{-1}})
}
$$








\section{用单纯概形表述下降}
\label{spaces-simplicial-section-simplicial-descent}

\noindent
笛卡尔态射定义如下。

\begin{definition}
\label{spaces-simplicial-definition-cartesian-morphism}
设 $a : Y \to X$ 是单纯概形的态射。如果对 $\Delta$
的每个态射 $\varphi : [n] \to [m]$，相应的图
$$
\xymatrix{
Y_m \ar[r]_a \ar[d]_{Y(\varphi)} & X_m \ar[d]^{X(\varphi)}\\
Y_n \ar[r]^{a} & X_n
}
$$
是概形范畴中的纤维方块，就称 $a$ 是{\it 笛卡尔的}，
或称{\it $Y$ 在 $X$ 上是笛卡尔的}。
\end{definition}

\noindent
笛卡尔态射与下降数据有关。首先证明一个一般引理，
描述固定单纯概形上的笛卡尔单纯概形之范畴。在该引理中，
以 $f^* : \Sch/X \to \Sch/Y$ 表示与概形态射
$f :Y \to X$ 关联的基变换函子。

\begin{lemma}
\label{spaces-simplicial-lemma-characterize-cartesian-schemes}
设 $X$ 是单纯概形。在 $X$ 上笛卡尔的单纯概形之范畴，
等价于二元组 $(V, \varphi)$ 的范畴，其中 $V$ 是 $X_0$ 上的概形，而
$$
\varphi :
V \times_{X_0, d^1_1} X_1
\longrightarrow
X_1 \times_{d^1_0, X_0} V
$$
是 $X_1$ 上的同构，满足 $(s_0^0)^*\varphi = \text{id}_V$，且
$$
(d^2_1)^*\varphi = (d^2_0)^*\varphi \circ (d^2_2)^*\varphi
$$
作为 $X_2$ 上的概形态射成立。
\end{lemma}

\begin{proof}
所显示等式的陈述有意义，因为
$d^1_1 \circ d^2_2 = d^1_1 \circ d^2_1$,
$d^1_1 \circ d^2_0 = d^1_0 \circ d^2_2$, and
$d^1_0 \circ d^2_0 = d^1_0 \circ d^2_1$ 作为态射
$X_2 \to X_0$ 成立；参见《单纯对象》评注
\ref{simplicial-remark-relations}。因此，可以把这些映射画成
$$
\xymatrix{
&
X_2 \times_{d^1_1 \circ d^2_0, X_0} V
\ar[r]_-{(d^2_0)^*\varphi} &
X_2 \times_{d^1_0 \circ d^2_0, X_0} V
\ar@{=}[rd] & \\
X_2 \times_{d^1_0 \circ d^2_2, X_0} V
\ar@{=}[ru] & & &
X_2 \times_{d^1_0 \circ d^2_1, X_0} V \\
&
X_2 \times_{d^1_1 \circ d^2_2, X_0} V
\ar[lu]^{(d^2_2)^*\varphi} \ar@{=}[r] &
X_2 \times_{d^1_1 \circ d^2_1, X_0} V
\ar[ru]_{(d^2_1)^*\varphi}
}
$$
，而该条件表示此图交换。显然，给定在 $X$ 上笛卡尔的单纯概形
$Y$，可以置 $V = Y_0$，并令 $\varphi$ 等于如下由笛卡尔结构
给出的恒等式之复合：
$$
V \times_{X_0, d^1_1} X_1 =
Y_0 \times_{X_0, d^1_1} X_1 = Y_1 =
X_1 \times_{X_0, d^1_0} Y_0 =
X_1 \times_{X_0, d^1_0} V
$$
为证明该函子是等价，构造一个拟逆。拟逆的构造类似于《下降》第
\ref{descent-section-descent-modules} 节所讨论的构造；我们从那里借用
记号 $\tau^n_i : [0] \to [n]$, $0 \mapsto i$ 和
$\tau^n_{ij} : [1] \to [n]$, $0 \mapsto i$, $1 \mapsto j$。
具体地，给定引理中的二元组 $(V, \varphi)$，置
$Y_n = X_n \times_{X(\tau^n_n), X_0} V$。给定
$\beta : [n] \to [m]$，定义
% P10-E0230
$V(\beta) : Y_m \to Y_n$ 为映射 $\varphi$ 与投影
$X_m \times_{X(\beta), X_n} Y_n \to Y_n$ 复合后，经
$X(\tau^m_{\beta(n)m})$ 的拉回。这是有意义的，因为
$$
X_m \times_{X(\tau^m_{\beta(n)m}), X_1} X_1 \times_{d^1_1, X_0} V
=
X_m \times_{X(\tau^m_m), X_0} V = Y_m
$$
以及
$$
X_m \times_{X(\tau^m_{\beta(n)m}), X_1} X_1 \times_{d^1_0, X_0} V =
X_m \times_{X(\tau^m_{\beta(n)}), X_0} V =
X_m \times_{X(\beta), X_n} Y_n.
$$
略去如下验证：上面所显示图的交换性蕴含这些映射正确复合。
也略去两个函子彼此为拟逆的验证。
\end{proof}

\begin{definition}
\label{spaces-simplicial-definition-fibre-products-simplicial-scheme}
设 $f : X \to S$ 是概形态射。与 $f$ {\it 关联的单纯概形}，
记为 $(X/S)_\bullet$，是函子
$\Delta^{opp} \to \Sch$, $[n] \mapsto X \times_S \ldots \times_S X$
，其描述见《单纯对象》例
\ref{simplicial-example-fibre-products-simplicial-object}。
\end{definition}

\noindent
因此，$(X/S)_n$ 是 $X$ 在 $S$ 上的 $(n + 1)$ 重纤维积。
态射 $d^1_0 : X \times_S X \to X$ 是映射
$(x_0, x_1) \mapsto x_1$，态射 $d^1_1$ 是另一个投影。
态射 $s^0_0$ 是对角态射 $X \to X \times_S X$。

\begin{lemma}
\label{spaces-simplicial-lemma-cartesian-over}
设 $f : X \to S$ 是概形态射，
$\pi : Y \to (X/S)_\bullet$ 是单纯概形的笛卡尔态射。
置 $V = Y_0$，并把它视为 $X$ 上的概形。态射
$d^1_0, d^1_1 : Y_1 \to Y_0$ 和态射
$\pi_1 : Y_1 \to X \times_S X$ 诱导同构
$$
\xymatrix{
V \times_S X & &
Y_1 \ar[ll]_-{(d^1_1, \text{pr}_1 \circ \pi_1)}
\ar[rr]^-{(\text{pr}_0 \circ \pi_1, d^1_0)} & &
X \times_S V.
}
$$
以 $\varphi : V \times_S X \to X \times_S V$ 表示由此得到的同构。
则二元组 $(V, \varphi)$ 是相对于 $X \to S$ 的下降数据。
\end{lemma}

\begin{proof}
这是引理 \ref{spaces-simplicial-lemma-characterize-cartesian-schemes}（一部分）
的特殊情形，因为该引理所显示的等式等价于《下降》定义
\ref{descent-definition-descent-datum} 的余圈条件。
\end{proof}

\begin{lemma}
\label{spaces-simplicial-lemma-cartesian-equivalent-descent-datum}
设 $f : X \to S$ 是概形态射。引理
\ref{spaces-simplicial-lemma-cartesian-over} 的构造
$$
\begin{matrix}
\text{笛卡尔概形的范畴} \\
\text{位于 } (X/S)_\bullet \text{ 上}
\end{matrix}
\longrightarrow
\begin{matrix}
\text{下降数据的范畴} \\
\text{相对于 } X/S
\end{matrix}
$$
是范畴的等价。
\end{lemma}

\begin{proof}
从左到右的函子由引理 \ref{spaces-simplicial-lemma-cartesian-over} 给出。
因此，这是引理 \ref{spaces-simplicial-lemma-characterize-cartesian-schemes} 的特殊情形。
\end{proof}

\noindent
可以如下重新解释《下降》引理 \ref{descent-lemma-pullback}
中的拉回。设给定单纯概形的态射 $f : X' \to X$ 和单纯概形的
笛卡尔态射 $Y \to X$。则单纯概形的纤维积（视为“拉回”）
$$
f^*Y = Y \times_X X'
$$
是在 $X'$ 上笛卡尔的单纯概形。设给定概形态射的交换图
$$
\xymatrix{
X' \ar[r]_f \ar[d] & X \ar[d] \\
S' \ar[r] & S.
}
$$
这给出单纯概形的态射
$$
f_\bullet : (X'/S')_\bullet \longrightarrow (X/S)_\bullet.
$$
我们断言，沿态射
$f_\bullet : (X'/S')_\bullet \to (X/S)_\bullet$ 的“拉回”
$f_\bullet^*$，经引理
\ref{spaces-simplicial-lemma-cartesian-equivalent-descent-datum} 的对应，
正是上述《下降》引理 \ref{descent-lemma-pullback}
中用下降数据定义的拉回。







\section{单纯概形上的拟凝聚模}
\label{spaces-simplicial-section-modules-simplicial}

\begin{lemma}
\label{spaces-simplicial-lemma-pullback-cartesian-module}
设 $f : V \to U$ 是单纯概形的态射。给定 $U_{Zar}$ 上的拟凝聚模
$\mathcal{F}$，其拉回 $f^*\mathcal{F}$ 是 $V_{Zar}$ 上的拟凝聚模。
\end{lemma}

\begin{proof}
回忆，$\mathcal{F}$ 是笛卡尔的，且 $\mathcal{F}_n$ 拟凝聚；
参见引理 \ref{spaces-simplicial-lemma-quasi-coherent-sheaf}。由引理
\ref{spaces-simplicial-lemma-describe-functoriality}，有
$(f^*\mathcal{F})_n = f_n^*\mathcal{F}_n$（略去一些细节）。
因此，$(f^*\mathcal{F})_n$ 拟凝聚。该事实与 $\mathcal{F}$
的笛卡尔性质共同蕴含 $f^*\mathcal{F}$ 的笛卡尔性质。
% P10-E0231
于是，再由引理 \ref{spaces-simplicial-lemma-quasi-coherent-sheaf}，
$\mathcal{F}$ 拟凝聚。
\end{proof}

\begin{lemma}
\label{spaces-simplicial-lemma-pushforward-cartesian-module}
设 $f : V \to U$ 是单纯概形的笛卡尔态射。
假设态射 $d^n_j : U_n \to U_{n - 1}$ 平坦，而态射
$V_n \to U_n$ 拟紧且拟分离。对 $V_{Zar}$ 上的拟凝聚模
$\mathcal{G}$，推前 $f_*\mathcal{G}$ 是 $U_{Zar}$ 上的拟凝聚模。
\end{lemma}

\begin{proof}
若 $\mathcal{F} = f_* \mathcal{G}$，则由引理
\ref{spaces-simplicial-lemma-describe-functoriality}，有
$\mathcal{F}_n = f_{n , *}\mathcal{G}_n$。
映射 $\mathcal{F}(\varphi)$ 用基变换映射定义；参见《上同调》第
\ref{cohomology-section-base-change-map} 节。由《概形》引理
\ref{schemes-lemma-push-forward-quasi-coherent}，以及由引理
\ref{spaces-simplicial-lemma-quasi-coherent-sheaf} 得知的 $\mathcal{G}_n$ 拟凝聚性，
层 $\mathcal{F}_n$ 拟凝聚。由《概形的上同调》引理
\ref{coherent-lemma-flat-base-change-cohomology}，以及由引理
\ref{spaces-simplicial-lemma-quasi-coherent-sheaf} 得知的 $\mathcal{G}$ 的笛卡尔性质，
% P10-E0232
沿退化态射 $d^n_j$ 的基变换映射是同构。
因此，由引理 \ref{spaces-simplicial-lemma-check-cartesian-module}，
$\mathcal{F}$ 是笛卡尔的；再由引理
\ref{spaces-simplicial-lemma-quasi-coherent-sheaf}，$\mathcal{F}$ 拟凝聚。
\end{proof}

\begin{lemma}
\label{spaces-simplicial-lemma-adjoint-functors-cartesian-modules}
设 $f : V \to U$ 是单纯概形的笛卡尔态射。
假设态射 $d^n_j : U_n \to U_{n - 1}$ 平坦，而态射
$V_n \to U_n$ 拟紧且拟分离。则 $f^*$ 与 $f_*$
构成 $U_{Zar}$ 和 $V_{Zar}$ 上拟凝聚模范畴之间的一对伴随函子。
\end{lemma}

\begin{proof}
引理 \ref{spaces-simplicial-lemma-pullback-cartesian-module} 和
\ref{spaces-simplicial-lemma-pushforward-cartesian-module} 已表明此陈述有意义。
伴随性立即由每个 $f_n^*$ 都左伴随于 $f_{n, *}$ 这一事实得出。
\end{proof}

\begin{lemma}
\label{spaces-simplicial-lemma-cartesian-modules-with-section}
设 $f : X \to S$ 是具有截面的概形态射\footnote{事实上，
只需假设 $f$ 在 $S$ 上 fpqc 局部地具有截面即可，因为由《下降》第
\ref{descent-section-fpqc-descent-quasi-coherent} 节，
拟凝聚模满足下降。}。令 $(X/S)_\bullet$ 为与 $X \to S$
关联的单纯概形；参见定义
\ref{spaces-simplicial-definition-fibre-products-simplicial-scheme}。
则拉回定义了拟凝聚 $\mathcal{O}_S$-模范畴与
$((X/S)_\bullet)_{Zar}$ 上拟凝聚模范畴之间的等价。
\end{lemma}

\begin{proof}
设 $\sigma : S \to X$ 是 $f$ 的截面。设 $(\mathcal{F}, \alpha)$
是引理 \ref{spaces-simplicial-lemma-characterize-cartesian-modules} 中的二元组。
置 $\mathcal{G} = \sigma^*\mathcal{F}$。考虑图
$$
\xymatrix{
X \ar[r]_-{(\sigma \circ f, 1)} \ar[d]_f &
X \times_S X \ar[d]^{\text{pr}_0} \ar[r]_-{\text{pr}_1} & X \\
S \ar[r]^\sigma & X
}
$$
注意，$\text{pr}_0 = d^1_1$ 且 $\text{pr}_1 = d^1_0$。
因此，$(\sigma \circ f, 1)^*\alpha$ 定义同构
$$
f^*\mathcal{G} = (\sigma \circ f, 1)^*\text{pr}_0^*\mathcal{F}
\longrightarrow
(\sigma \circ f, 1)^*\text{pr}_1^*\mathcal{F} = \mathcal{F}
$$
略去如下验证：该同构与 $\alpha$ 以及典范同构
$\text{pr}_0^*f^*\mathcal{G} \to \text{pr}_1^*f^*\mathcal{G}$ 相容。
\end{proof}






\section{群胚与单纯概形}
\label{spaces-simplicial-section-groupoids-simplicial}

\noindent
给定概形中的群胚，可以构造单纯概形。我们将看到，
群胚上的拟凝聚层范畴等价于关联单纯概形上的笛卡尔拟凝聚层范畴。

\begin{lemma}
\label{spaces-simplicial-lemma-groupoid-simplicial}
设 $(U, R, s, t, c, e, i)$ 是 $S$ 上的群胚概形。
则存在 $S$ 上的单纯概形 $X$，满足下列性质：
\begin{enumerate}
\item $X_0 = U$, $X_1 = R$, $X_2 = R \times_{s, U, t} R$；
\item $s_0^0 = e : X_0 \to X_1$；
\item $d^1_0 = s : X_1 \to X_0$, $d^1_1 = t : X_1 \to X_0$；
\item $s_0^1 = (e \circ t, 1) : X_1 \to X_2$,
$s_1^1 = (1, e \circ t) : X_1 \to X_2$；% P10-E0233
\item $d^2_0 = \text{pr}_1 : X_2 \to X_1$,
$d^2_1 = c : X_2 \to X_1$,
$d^2_2 = \text{pr}_0$；
\item $X = \text{cosk}_2 \text{sk}_2 X$.
\end{enumerate}
对所有 $n$，$X_n = R \times_{s, U, t} \ldots \times_{s, U, t} R$，
其中有 $n$ 个因子。映射 $d^n_j : X_n \to X_{n - 1}$
在点函子上由下式给出：
$$
(r_1, \ldots, r_n) \longmapsto (r_1, \ldots, c(r_j, r_{j + 1}), \ldots, r_n)
$$
这里 $1 \leq j \leq n - 1$；而
$d^n_0(r_1, \ldots, r_n) = (r_2, \ldots, r_n)$，且
$d^n_n(r_1, \ldots, r_n) = (r_1, \ldots, r_{n - 1})$。
\end{lemma}

\begin{proof}
只需验证 (1)、(2)、(3)、(4)、(5) 中规定的规则定义了
$S$ 上的 $2$-截断单纯概形 $U'$，因为此时 (6) 允许置
$X = \text{cosk}_2 U'$；参见《单纯对象》引理
\ref{simplicial-lemma-existence-cosk}。采用点函子方法，
只需验证：若 $(\text{对象}, \text{箭头}, s, t, c, e, i)$ 是群胚，则
$$
\xymatrix{
\text{箭头} \times_{s, \text{对象}, t} \text{箭头}
\ar@<8ex>[d]^{\text{pr}_0}
\ar@<0ex>[d]_c
\ar@<-8ex>[d]_{\text{pr}_1}
\\
\text{箭头}
\ar@<4ex>[d]^t
\ar@<-4ex>[d]_s
\ar@<4ex>[u]^{1, e}
\ar@<-4ex>[u]_{e, 1}
\\
\text{对象}
\ar@<0ex>[u]_e
}
$$
是 $2$-截断单纯集。略去细节。

\medskip\noindent
最后，对 $n > 2$ 的 $X_n$ 的描述，由 $X_0$、$X_1$、$X_2$
的描述以及《单纯对象》评注
\ref{simplicial-remark-inductive-coskeleton} 和引理
\ref{simplicial-lemma-work-out} 归纳得出。也可以证明，
将 $\text{cosk}_2$ 应用于上面所显示的 $2$-截断单纯集，
得到的单纯集之第 $n$ 项等于
$\text{箭头} \times_{s, \text{对象}, t} \ldots
\times_{s, \text{对象}, t} \text{箭头}$，有 $n$ 个因子，
而退化映射如引理所述。略去一些细节。
\end{proof}

\begin{lemma}
\label{spaces-simplicial-lemma-quasi-coherent-groupoid-simplicial}
设 $S$ 是概形，$(U, R, s, t, c)$ 是 $S$ 上的群胚概形，
$X$ 是引理 \ref{spaces-simplicial-lemma-groupoid-simplicial} 构造的 $S$ 上单纯概形。
则 $(U, R, s, t, c)$ 上的拟凝聚模范畴，
等价于 $X_{Zar}$ 上的拟凝聚模范畴。
\end{lemma}

\begin{proof}
这由引理 \ref{spaces-simplicial-lemma-quasi-coherent-sheaf}、
\ref{spaces-simplicial-lemma-characterize-cartesian-modules} 和《群胚》定义
\ref{groupoids-definition-groupoid-module} 显然成立。
\end{proof}

\noindent
在下一个引理中，将使用定义 \ref{spaces-simplicial-definition-cartesian-morphism}
所定义的单纯概形之笛卡尔态射 $V \to U$ 的概念。

\begin{lemma}
\label{spaces-simplicial-lemma-quasi-coherent-groupoid-R-cartesian}
设 $(U, R, s, t, c)$ 是概形 $S$ 上的群胚概形，
$X$ 是引理 \ref{spaces-simplicial-lemma-groupoid-simplicial} 构造的 $S$ 上单纯概形。
令 $(R/U)_\bullet$ 为与 $s : R \to U$ 关联的单纯概形；
参见定义 \ref{spaces-simplicial-definition-fibre-products-simplicial-scheme}。
则存在单纯概形的笛卡尔态射
$t_\bullet : (R/U)_\bullet \to X$，其低次数态射由下图给出：
$$
\xymatrix{
R \times_{s, U, s} R \times_{s, U, s} R
\ar@<3ex>[r]_-{\text{pr}_{12}}
\ar@<0ex>[r]_-{\text{pr}_{02}}
\ar@<-3ex>[r]_-{\text{pr}_{01}}
\ar[dd]_{(r_0, r_1, r_2) \mapsto (r_0 \circ r_1^{-1}, r_1 \circ r_2^{-1})} &
R \times_{s, U, s} R
\ar@<1ex>[r]_-{\text{pr}_1} \ar@<-2ex>[r]_-{\text{pr}_0}
\ar[dd]_{(r_0, r_1) \mapsto r_0 \circ r_1^{-1}} &
R \ar[dd]^t
\\
\\
R \times_{s, U, t} R
\ar@<3ex>[r]_{\text{pr}_1}
\ar@<0ex>[r]_c
\ar@<-3ex>[r]_{\text{pr}_0} &
R \ar@<1ex>[r]_s \ar@<-2ex>[r]_t &
U
}
$$
\end{lemma}

\begin{proof}
对任意 $n$，按规则
$$
(r_0, \ldots, r_n)
\longrightarrow
(r_0 \circ r_1^{-1}, \ldots, r_{n - 1} \circ r_n^{-1})
$$
定义 $(R/U)_\bullet \to X_n$。由引理
\ref{spaces-simplicial-lemma-groupoid-simplicial} 对退化映射的描述，
它与退化映射的相容性显然成立。
% P10-E0234
略去这些映射保持态射 $s^n_j$ 的验证。《群胚》引理
\ref{groupoids-lemma-diagram-pull}（交换 $s$ 与 $t$ 的角色）
表明右侧两个方块是笛卡尔的。完全类似地，可证所有其他方块
也都是笛卡尔的。因此，该态射是笛卡尔的。
\end{proof}




\section{下降数据给出等价关系}
\label{spaces-simplicial-section-equivalence-relation}

\noindent
在第 \ref{spaces-simplicial-section-simplicial-descent} 节中，我们看到如何用
$(X/S)_\bullet$ 上的笛卡尔单纯概形，表述相对于 $X \to S$
的下降数据。现在把它与等价关系联系起来。

\begin{lemma}
\label{spaces-simplicial-lemma-equivalence-relation}
设 $f : X \to S$ 是概形态射，
$\pi : Y \to (X/S)_\bullet$ 是单纯概形的笛卡尔态射；
参见定义 \ref{spaces-simplicial-definition-cartesian-morphism} 和
\ref{spaces-simplicial-definition-fibre-products-simplicial-scheme}。则态射
$$
j = (d^1_1, d^1_0) : Y_1 \to Y_0 \times_S Y_0
$$
定义 $S$ 上 $Y_0$ 的等价关系；参见《群胚》定义
\ref{groupoids-definition-equivalence-relation}。
\end{lemma}

\begin{proof}
注意，$j$ 是单态射。事实上，由于 $\pi$ 是笛卡尔的，
复合 $Y_1 \to Y_0 \times_S Y_0 \to Y_0 \times_S X$ 是同构。

\medskip\noindent
考虑态射
$$
(d^2_2, d^2_0) : Y_2 \to Y_1 \times_{d^1_0, Y_0, d^1_1} Y_1.
$$
这是可行的，因为 $d_0 \circ d_2 = d_1 \circ d_0$；
参见《单纯对象》评注 \ref{simplicial-remark-relations}。
此外，它是 $(X/S)_2$ 上的态射。由于
$Y \to (X/S)_\bullet$ 是笛卡尔的，它是同构。
例如，右边同构于
$Y_0 \times_{\pi_0, X, \text{pr}_1} (X \times_S X \times_S X) =
X \times_S Y_0 \times_S X$
，因为 $\pi$ 是笛卡尔的。略去细节。

\medskip\noindent
仿照《群胚》定义 \ref{groupoids-definition-equivalence-relation}，
记 $t = \text{pr}_0 \circ j = d^1_1$ 和
$s = \text{pr}_1 \circ j = d^1_0$。上述同构与态射
$d^2_1 : Y_2 \to Y_1$ 共同给出复合态射
$$
c : Y_1 \times_{s, Y_0, t} Y_1 \longrightarrow Y_1
$$
，它位于 $Y_0 \times_S Y_0$ 上方。这立即蕴含：对任意概形
$T/S$，关系 $Y_1(T) \subset Y_0(T) \times Y_0(T)$ 具有传递性。

\medskip\noindent
自反性来自如下事实：态射 $j$ 在对角
$\Delta : X \to X \times_S X$ 上的限制是同构
（再次使用 $\pi$ 的笛卡尔性质）。

\medskip\noindent
为证明对称性，考虑态射
$$
(d^2_2, d^2_1) : Y_2 \to Y_1 \times_{d^1_1, Y_0, d^1_1} Y_1.
$$
这是可行的，因为 $d_1 \circ d_2 = d_1 \circ d_1$；
参见《单纯对象》评注 \ref{simplicial-remark-relations}。
由于 $Y \to (X/S)_\bullet$ 是笛卡尔的，它是同构。
例如，右边同构于
$Y_0 \times_{\pi_0, X, \text{pr}_0} (X \times_S X \times_S X) =
Y_0 \times_S X \times_S X$
，因为 $\pi$ 是笛卡尔的。略去细节。

\medskip\noindent
设 $T/S$ 是概形。对 $a, b \in Y_0(T)$，规定 $a \sim b$
与 $(a, b) \in Y_1(T)$ 同义。上述同构 $(d^2_2, d^2_1)$
蕴含：若 $a \sim b$ 且 $a \sim c$，则 $b \sim c$。
结合自反性可知，$\sim$ 是等价关系。
\end{proof}







\section{一个示例情形}
\label{spaces-simplicial-section-example}

\noindent
本节证明，Artin 环的谱之不交并可以沿概形的拟紧平坦满射下降。

\begin{lemma}
\label{spaces-simplicial-lemma-equivalence-classes-points}
设 $X \to S$ 是概形态射，并假设
$Y \to (X/S)_\bullet$ 是单纯概形的笛卡尔态射。
对点 $y \in Y_0$，定义 $Y_0$ 的子集
$$
T_y = \{y' \in Y_0 \mid \exists\ y_1 \in Y_1:
d^1_1(y_1) = y, d^1_0(y_1) = y'\}
$$
。则 $y \in T_y$，且
$T_y \cap T_{y'} \not = \emptyset \Rightarrow T_y = T_{y'}$。
\end{lemma}

\begin{proof}
合用引理 \ref{spaces-simplicial-lemma-equivalence-relation} 和《群胚》引理
\ref{groupoids-lemma-pre-equivalence-equivalence-relation-points}。
\end{proof}

\begin{lemma}
\label{spaces-simplicial-lemma-quasi-compact}
设 $X \to S$ 是概形态射，并假设
$Y \to (X/S)_\bullet$ 是单纯概形的笛卡尔态射。
设 $y \in Y_0$ 是点。若 $X \to S$ 拟紧，则
$$
T_y = \{y' \in Y_0 \mid \exists\ y_1 \in Y_1:
d^1_1(y_1) = y, d^1_0(y_1) = y'\}
$$
是 $Y_0$ 的拟紧子集。
\end{lemma}

\begin{proof}
设 $F_y$ 是 $d^1_1 : Y_1 \to Y_0$ 在 $y$ 处的概形论纤维。
则 $T_y$ 是如下态射的像：
$$
\xymatrix{
F_y \ar[r] \ar[d] &
Y_1 \ar[r]^{d^1_0} \ar[d]^{d^1_1} &
Y_0 \\
y \ar[r] &
Y_0 &
}
$$
注意，$F_y$ 拟紧。这证明了引理。
\end{proof}

\begin{lemma}
\label{spaces-simplicial-lemma-descent-disjoint-union-Artinian-along-fields}
设 $X \to S$ 是拟紧平坦满射，$(V, \varphi)$ 是相对于
$X \to S$ 的下降数据。若 $V$ 是 Artin 环的谱之不交并，
则 $(V, \varphi)$ 有效。
\end{lemma}

\begin{proof}
由引理 \ref{spaces-simplicial-lemma-cartesian-equivalent-descent-datum}，
设 $Y \to (X/S)_\bullet$ 为对应于 $(V, \varphi)$ 的单纯概形之
笛卡尔态射。注意，$Y_0 = V$。写
$V = \coprod_{i \in I} \Spec(A_i)$，其中每个 $A_i$ 是局部 Artin 环。
此外，对所有 $i \in I$，设 $v_i \in V$ 为 $\Spec(A_i)$
的唯一闭点。采用上面引理
\ref{spaces-simplicial-lemma-equivalence-classes-points} 的记号，规定
$i \sim j$ 当且仅当 $v_i \in T_{v_j}$。由引理
\ref{spaces-simplicial-lemma-equivalence-classes-points} 和 \ref{spaces-simplicial-lemma-quasi-compact}，
这是等价类均有限的等价关系。令 $\overline{I} = I/\sim$。
于是可以写成
$V = \coprod_{\overline{i} \in \overline{I}} V_{\overline{i}}$
，其中
$V_{\overline{i}} = \coprod_{i \in \overline{i}} \Spec(A_i)$.
由构造，$\varphi : V \times_S X \to X \times_S V$
把开闭子空间 $V_{\overline{i}} \times_S X$ 映入开闭子空间
$X \times_S V_{\overline{i}}$。换言之，得到下降数据
$(V_{\overline{i}}, \varphi_{\overline{i}})$，而 $(V, \varphi)$
是它们在下降数据范畴中的余积。由于每个 $V_{\overline{i}}$
都是局部 Artin 环的谱的有限并，态射 $V_{\overline{i}} \to X$
是仿射的；参见《态射》引理
\ref{morphisms-lemma-Artinian-affine}。由于 $\{X \to S\}$
是 fpqc 覆盖，由《下降》引理 \ref{descent-lemma-affine}，
所有下降数据 $(V_{\overline{i}}, \varphi_{\overline{i}})$ 都有效。
\end{proof}

\noindent
诚然，上述引理的适用范围十分有限！









\section{单纯代数空间}
\label{spaces-simplicial-section-simplicial-algebraic-spaces}

\noindent
设 $S$ 是概形。{\it 单纯代数空间}是 $S$ 上代数空间范畴中的单纯对象；
参见《单纯对象》定义 \ref{simplicial-definition-simplicial-object}。
回忆，单纯代数空间形如
$$
\xymatrix{
X_2
\ar@<2ex>[r]
\ar@<0ex>[r]
\ar@<-2ex>[r]
&
X_1
\ar@<1ex>[r]
\ar@<-1ex>[r]
\ar@<1ex>[l]
\ar@<-1ex>[l]
&
X_0
\ar@<0ex>[l]
}
$$
这里有两个态射 $d^1_0, d^1_1 : X_1 \to X_0$ 和一个态射
$s^0_0 : X_0 \to X_1$，依此类推。这些态射满足一些必需关系，
例如 $d^1_0 \circ s^0_0 = \text{id}_{X_0} = d^1_1 \circ s^0_0$；
参见《单纯对象》引理
\ref{simplicial-lemma-characterize-simplicial-object}。
把 $d^n_i : X_n \to X_{n - 1}$ 看作“忘掉第 $i$ 个坐标的投影”，
把 $s^n_j : X_n \to X_{n + 1}$ 看作“重复第 $j$ 个坐标的对角映射”，
会很有帮助。

\medskip\noindent
{\it 单纯代数空间的态射} $h : X \to Y$，就是 $S$ 上代数空间范畴中
单纯对象的态射；参见《单纯对象》定义
\ref{simplicial-definition-simplicial-object}。因此，$h$
由代数空间态射 $h_n : X_n \to Y_n$ 组成，并且只要等式有意义，
就有 $h_{n - 1} \circ d^n_j = d^n_j \circ h_n$ 和
$h_{n + 1} \circ s^n_j = s^n_j \circ h_n$。

\medskip\noindent
单纯代数空间 $X$ 的{\it 增广} $a : X \to X_{-1}$，
由满足 $a_0 \circ d^1_0 = a_0 \circ d^1_1$ 的代数空间态射
$a_0 : X_0 \to X_{-1}$ 给出。参见《单纯对象》第
\ref{simplicial-section-augmentation} 节。在这种情形下，
对 $n \geq 0$，总以 $a_n : X_n \to X_{-1}$ 表示所诱导的态射。

\medskip\noindent
设 $X$ 是单纯代数空间。对每个 $n$，有站点
$X_{n, spaces, \etale}$（《空间的性质》定义
\ref{spaces-properties-definition-spaces-etale-site}）；
对每个态射 $\varphi : [m] \to [n]$，有站点态射
$$
f_\varphi = X(\varphi)_{spaces, \etale} :
X_{n, spaces, \etale} \to X_{m, spaces, \etale},
$$
，它与代数空间态射 $X(\varphi) : X_n \to X_m$ 关联
（《空间的性质》引理
\ref{spaces-properties-lemma-functoriality-etale-site}）。
这给出站点范畴中的单纯对象。在引理
\ref{spaces-simplicial-lemma-simplicial-site-site} 中，我们构造了关联站点，
记为 $X_{spaces, \etale}$。站点 $X_{spaces, \etale}$ 的对象，
% P10-E0235
是对某个 $n$ 在 $X_n$ 上平展的代数空间 $U$；态射
$(\varphi, f) : U/X_n \to V/X_m$ 由 $\Delta$ 中的态射
$\varphi : [m] \to [n]$ 和代数空间态射 $f : U \to V$ 给出，
并要求下图交换：
$$
\xymatrix{
U \ar[r]_f \ar[d] & V \ar[d] \\
X_n \ar[r]^{f_\varphi} & X_m
}
$$
考虑全子范畴
$$
X_{affine, \etale} \subset X_\etale \subset X_{spaces, \etale}
$$
，其对象是 $U/X_n$，其中 $U$ 分别为仿射的或为概形。
给这些范畴赋予其自然拓扑（参见《空间的性质》引理
\ref{spaces-properties-lemma-alternative}、定义
\ref{spaces-properties-definition-etale-site} 和引理
\ref{spaces-properties-lemma-compare-etale-sites}），
这些包含函子定义拓扑斯等价
$$
\Sh(X_{affine, \etale}) = \Sh(X_\etale) = \Sh(X_{spaces, \etale})
$$
下文默认识别这些拓扑斯。称 $X_\etale$ 为
{\it $X$ 的小平展站点}，称其拓扑斯为{\it $X$ 的小平展拓扑斯}。

\medskip\noindent
设 $X_\etale$ 是单纯代数空间 $X$ 的小平展站点。
$X_\etale$ 上存在环层 $\mathcal{O}$，它在 $X_n$ 上的限制
是结构层 $\mathcal{O}_{X_n}$。这由引理
\ref{spaces-simplicial-lemma-describe-sheaves-simplicial-site-site} 得出。
我们称{\it $\mathcal{O}$ 是单纯代数空间 $X$ 的结构层}。
至此，为单纯（赋环）站点发展的一切材料都适用；参见第
\ref{spaces-simplicial-section-simplicial-sites}、
\ref{spaces-simplicial-section-augmentation-simplicial-sites}、
\ref{spaces-simplicial-section-morphism-simplicial-sites}、
\ref{spaces-simplicial-section-simplicial-sites-modules}、
\ref{spaces-simplicial-section-cohomology-simplicial-sites}、
\ref{spaces-simplicial-section-cohomology-augmentation-simplicial-sites}、
\ref{spaces-simplicial-section-cohomology-simplicial-sites-modules}、
\ref{spaces-simplicial-section-cohomology-augmentation-ringed-simplicial-sites}、
\ref{spaces-simplicial-section-cartesian}、
\ref{spaces-simplicial-section-glueing} 和
\ref{spaces-simplicial-section-glueing-modules} 节。

\medskip\noindent
设 $X$ 是结构层为 $\mathcal{O}$ 的单纯代数空间。
与任何赋环拓扑斯上一样，可以定义
{\it $X_\etale$ 上的拟凝聚 $\mathcal{O}$-模}；参见
《站点上的模》定义 \ref{sites-modules-definition-site-local}。
不过，$X_\etale$ 上的拟凝聚 $\mathcal{O}$-模，恰是如下笛卡尔
$\mathcal{O}$-模 $\mathcal{F}$：其限制 $\mathcal{F}_n$
在 $X_n$ 上拟凝聚；参见引理 \ref{spaces-simplicial-lemma-quasi-coherent-sheaf}。

\medskip\noindent
设 $h : X \to Y$ 是 $S$ 上单纯代数空间的态射。将引理
\ref{spaces-simplicial-lemma-morphism-simplicial-sites} 应用于站点态射
% P10-E0237
$(h_n)_{spaces, \etale} : X_{spaces, \etale} \to Y_{spaces, \etale}$
（《空间的性质》引理
\ref{spaces-properties-lemma-functoriality-etale-site}），
得到小平展拓扑斯的态射
$h_\etale : \Sh(X_\etale) \to \Sh(Y_\etale)$。
回忆，$h_\etale^{-1}$ 和 $h_{\etale, *}$ 可用各分量简单描述；
参见引理 \ref{spaces-simplicial-lemma-morphism-simplicial-sites}。
以 $\mathcal{O}_X$、$\mathcal{O}_Y$ 分别表示 $X$、$Y$ 的结构层。
定义
% P10-E0229
$h_\etale^\sharp : h_{\etale, *}\mathcal{O}_X \to \mathcal{O}_Y$
为 $Y_\etale$ 上的环层映射，它在 $Y_n$ 上由
$h_n^\sharp : h_{n, *}\mathcal{O}_{X_n} \to \mathcal{O}_{Y_n}$ 给出。
由此得到赋环拓扑斯的态射
$$
h_\etale :
(\Sh(X_\etale), \mathcal{O}_X)
\longrightarrow
(\Sh(Y_\etale), \mathcal{O}_Y)
$$

\medskip\noindent
设 $X$ 是结构层为 $\mathcal{O}$ 的单纯代数空间。
设 $X_{-1}$ 是 $S$ 上的代数空间，$a_0 : X_0 \to X_{-1}$
是 $X$ 的增广。将引理 \ref{spaces-simplicial-lemma-augmentation-site}
应用于站点态射
$(a_0)_{spaces, \etale} :
X_{0, spaces, \etale} \to X_{-1, spaces, \etale}$
，得到相应拓扑斯态射
$a : \Sh(X_\etale) \to \Sh(X_{-1, \etale})$。
注意，$a^{-1}\mathcal{G}$ 是 $X_\etale$ 上分量为
$a_n^{-1}\mathcal{G}$ 的层。因此，可以利用映射
$a_n^\sharp : a_n^{-1}\mathcal{O}_{X_{-1}} \to \mathcal{O}_{X_n}$
定义映射 $a^\sharp : a^{-1}\mathcal{O}_{X_{-1}} \to \mathcal{O}$；
等价地，由伴随定义映射
$a^\sharp : \mathcal{O}_{X_{-1}} \to a_*\mathcal{O}$
（照例仍用同一名称）。这正处于第
\ref{spaces-simplicial-section-cohomology-augmentation-ringed-simplicial-sites}
节所讨论的情形。因此，得到赋环拓扑斯的态射
$$
a :
(\Sh(X_\etale), \mathcal{O})
\longrightarrow
(\Sh(X_{-1}), \mathcal{O}_{X_{-1}})
$$

\medskip\noindent
最后作如下观察。设给定单纯代数空间 $X$ 与 $Y$ 的态射
$h : X \to Y$，其结构层分别为 $\mathcal{O}_X$、$\mathcal{O}_Y$；
并给定增广 $a_0 : X_0 \to X_{-1}$、$b_0 : Y_0 \to Y_{-1}$
和态射 $h_{-1} : X_{-1} \to Y_{-1}$，使得
$$
\xymatrix{
X_0 \ar[r]_{h_0} \ar[d]_{a_0} & Y_0 \ar[d]^{b_0} \\
X_{-1} \ar[r]^{h_{-1}} & Y_{-1}
}
$$
交换。由上面阐明的构造，得到如下赋环拓扑斯态射的交换图：
% P10-E0236
$$
\xymatrix{
(\Sh(X_\etale), \mathcal{O}_X) \ar[r]_{h_\etale} \ar[d]_a &
(\Sh(Y_\etale), \mathcal{O}_Y) \ar[d]^b \\
(\Sh(X_{-1}), \mathcal{O}_{X_{-1}}) \ar[r]^{h_{-1}} &
(\Sh(Y_{-1}), \mathcal{O}_{Y_{-1}})
}
$$




\section{代数空间的 fppf 超覆盖}
\label{spaces-simplicial-section-fppf-hypercovering}

\noindent
本节是第 \ref{spaces-simplicial-section-proper-hypercovering} 节在代数空间和 fppf
超覆盖情形下的类似版本。有意如此做的读者，可以处处把“代数空间”
替换为“概形”，得到同样有效的结果。这样做的优点是，可以把对
《空间的上同调进阶》第
\ref{spaces-more-cohomology-section-fppf-etale} 节的引用，
替换为对《平展上同调》第
\ref{etale-cohomology-section-fppf-etale} 节的引用。

\medskip\noindent
固定基概形 $S$。设 $X$ 是 $S$ 上的代数空间，
$U$ 是 $S$ 上的单纯代数空间。假设有增广
$$
a : U \to X
$$
参见第 \ref{spaces-simplicial-section-simplicial-algebraic-spaces} 节。
如果满足下列条件，就称 $U$ 是 $X$ 的{\it fppf 超覆盖}：
\begin{enumerate}
\item $U_0 \to X$ 平坦、局部有限呈示且满射；
\item $U_1 \to U_0 \times_X U_0$ 平坦、局部有限呈示且满射；
\item $U_{n + 1} \to (\text{cosk}_n\text{sk}_n U)_{n + 1}$
对 $n \geq 1$ 平坦、局部有限呈示且满射。
\end{enumerate}
$S$ 上的代数空间范畴具有所有有限极限，所以上述表述中使用的余骨架存在。
$$
\fbox{原理：可以用 fppf 超覆盖计算平展上同调。}
$$
证明该原理的关键思想，是比较范畴 $\textit{Spaces}/S$ 上的 fppf
拓扑与平展拓扑。具体地，fppf 拓扑比平展拓扑更细，并且：
(a) 平坦、局部有限呈示的满射定义 fppf 覆盖；
(b) 从小平展站点拉回的层之 fppf 上同调与平展上同调一致，
正如《空间的上同调进阶》第
\ref{spaces-more-cohomology-section-fppf-etale} 节所述。

\begin{lemma}
\label{spaces-simplicial-lemma-compare-simplicial-objects-fppf-etale}
设 $S$ 是概形，$X$ 是 $S$ 上的代数空间，
$U$ 是 $S$ 上的单纯代数空间，$a : U \to X$ 是增广。
则有交换图
$$
\xymatrix{
\Sh((\textit{Spaces}/U)_{fppf, total}) \ar[r]_-h \ar[d]_{a_{fppf}} &
\Sh(U_\etale) \ar[d]^a \\
\Sh((\textit{Spaces}/X)_{fppf}) \ar[r]^-{h_{-1}} &
\Sh(X_\etale)
}
$$
其中左侧竖直箭头定义于第 \ref{spaces-simplicial-section-hypercovering} 节，
右侧竖直箭头定义于第
\ref{spaces-simplicial-section-simplicial-algebraic-spaces} 节。
\end{lemma}

\begin{proof}
记号 $(\textit{Spaces}/U)_{fppf, total}$ 表示，我们对站点
$(\textit{Spaces}/S)_{fppf}$ 及其单纯对象 $U$ 使用第
\ref{spaces-simplicial-section-hypercovering} 节的构造\footnote{也可以使用平展拓扑，
此时记为 $(\textit{Spaces}/U)_{\etale, total}$。}。
对右边的拓扑斯，我们将使用站点 $X_{spaces, \etale}$ 和
$U_{spaces, \etale}$；这是允许的，见第
\ref{spaces-simplicial-section-simplicial-algebraic-spaces} 节中的讨论。

\medskip\noindent
注意，$(\textit{Spaces}/U)_{fppf, total}$ 与
$U_{spaces, \etale}$ 都属于情形
\ref{spaces-simplicial-situation-simplicial-site} 的 A 类。对 $U_\etale$，
这立即由第 \ref{spaces-simplicial-section-simplicial-algebraic-spaces} 节中的构造得出；
对 $(\textit{Spaces}/U)_{fppf, total}$，则由引理
\ref{spaces-simplicial-lemma-sr-when-fibre-products} 得出。接着考虑函子
$U_{n, spaces, \etale} \to (\textit{Spaces}/U_n)_{fppf}$, $U \mapsto U/U_n$
以及
$X_{spaces, \etale} \to (\textit{Spaces}/X)_{fppf}$, $U \mapsto U/X$.
我们已在《空间的上同调进阶》第
\ref{spaces-more-cohomology-section-fppf-etale} 节看到，
它们定义站点态射；在那里分别记为
$a_{U_n} = \epsilon_{U_n} \circ \pi_{u_n}$ 和
$a_X = \epsilon_X \circ \pi_X$。因此，得到评注
\ref{spaces-simplicial-remark-morphism-augmentation-simplicial-sites}
意义下与增广相容的单纯站点态射；应用引理
\ref{spaces-simplicial-lemma-morphism-augmentation-simplicial-sites} 即得结论。
\end{proof}

\begin{lemma}
\label{spaces-simplicial-lemma-descent-sheaves-for-fppf-hypercovering}
设 $S$ 是概形，$X$ 是 $S$ 上的代数空间，
$U$ 是 $S$ 上的单纯代数空间，$a : U \to X$ 是增广。
若 $a : U \to X$ 是 $X$ 的 fppf 超覆盖，则
$$
a^{-1} : \Sh(X_\etale) \to \Sh(U_\etale)
\quad\text{且}\quad
a^{-1} : \textit{Ab}(X_\etale) \to \textit{Ab}(U_\etale)
$$
满忠实，其本质像为笛卡尔层，拟逆由 $a_*$ 给出。
这里 $a : \Sh(U_\etale) \to \Sh(X_\etale)$ 如第
\ref{spaces-simplicial-section-simplicial-algebraic-spaces} 节所述。
\end{lemma}

\begin{proof}
我们证明集合层的陈述。它几乎是已建立结果的形式推论。
考虑引理 \ref{spaces-simplicial-lemma-compare-simplicial-objects-fppf-etale} 的图。
该引理的证明表明，$h_{-1}$ 是《空间的上同调进阶》第
\ref{spaces-more-cohomology-section-fppf-etale} 节中的态射 $a_X$。
因此，由《空间的上同调进阶》引理
\ref{spaces-more-cohomology-lemma-comparison-fppf-etale}，
$(h_{-1})^{-1}$ 满忠实，其拟逆为 $h_{-1, *}$。
对 $h$ 的分量 $h_n$ 也同样成立。由引理
\ref{spaces-simplicial-lemma-morphism-simplicial-sites} 对函子 $h^{-1}$ 和 $h_*$
的描述，可知 $h^{-1}$ 满忠实，其拟逆为 $h_*$。
$U$ 是 $(\textit{Spaces}/S)_{fppf}$ 中 $X$ 的超覆盖，
其定义见第 \ref{spaces-simplicial-section-hypercovering} 节。
由引理 \ref{spaces-simplicial-lemma-hypercovering-X-simple-descent-sheaves}，
$a_{fppf}^{-1}$ 满忠实，其拟逆为 $a_{fppf, *}$，本质像为
$(\textit{Spaces}/U)_{fppf, total}$ 上的笛卡尔层。
现在，一个形式论证（沿图追逐）表明 $a^{-1}$ 满忠实。

\medskip\noindent
最后，假设 $\mathcal{G}$ 是 $U_\etale$ 上的笛卡尔层。
则 $h^{-1}\mathcal{G}$ 是 $(\textit{Spaces}/U)_{fppf, total}$
上的笛卡尔层。因此，对 $(\textit{Spaces}/X)_{fppf}$ 上的某个层
$\mathcal{H}$，有
$h^{-1}\mathcal{G} = a_{fppf}^{-1}\mathcal{H}$。
特别地，$h_0^{-1}\mathcal{G}_0 =
(a_{0, big, fppf})^{-1}\mathcal{H}$。回忆 $h_0 = a_{U_0}$，
且 $U_0 \to X$ 平坦、局部有限呈示并且满射；由
《空间的上同调进阶》引理
\ref{spaces-more-cohomology-lemma-descent-sheaf-fppf-etale}，
存在 $X_\etale$ 上的层 $\mathcal{F}$ 及同构
$\mathcal{H} = (h_{-1})^{-1}\mathcal{F}$。
由于 $a_{fppf}^{-1}\mathcal{H} = h^{-1}\mathcal{G}$，可得
$h^{-1}\mathcal{G} \cong h^{-1}a^{-1}\mathcal{F}$。
由 $h^{-1}$ 的满忠实性，得到
$a^{-1}\mathcal{F} \cong \mathcal{G}$。

\medskip\noindent
固定同构 $\theta : a^{-1}\mathcal{F} \to \mathcal{G}$。
为完成证明，须证明 $\mathcal{G} = a^{-1}a_*\mathcal{G}$
（目的是证明拟逆由 $a_*$ 给出；其余内容均已在上面证明）。
由于 $a^{-1}$ 满忠实，由《范畴》引理
\ref{categories-lemma-adjoint-fully-faithful}，有
$\text{id} \cong a_*a^{-1}$。因此，
$\mathcal{F} \cong a_*a^{-1}\mathcal{F}$ 与
$a_*\theta : a_*a^{-1}\mathcal{F} \to a_*\mathcal{G}$
合成同构 $\mathcal{F} \to a_*\mathcal{G}$。
将其沿 $a$ 拉回并与 $\theta^{-1}$ 预复合，便得到所需同构。
\end{proof}

\begin{lemma}
\label{spaces-simplicial-lemma-cohomological-descent-for-fppf-hypercovering}
设 $S$ 是概形，$X$ 是 $S$ 上的代数空间，
$U$ 是 $S$ 上的单纯代数空间，$a : U \to X$ 是增广。
若 $a : U \to X$ 是 $X$ 的 fppf 超覆盖，则对
$K \in D^+(X_\etale)$，
$$
K \to Ra_*(a^{-1}K)
$$
是同构。这里 $a : \Sh(U_\etale) \to \Sh(X_\etale)$ 如第
\ref{spaces-simplicial-section-simplicial-algebraic-spaces} 节所述。
\end{lemma}

\begin{proof}
考虑引理 \ref{spaces-simplicial-lemma-compare-simplicial-objects-fppf-etale} 的图。
由《空间的上同调进阶》引理
\ref{spaces-more-cohomology-lemma-cohomological-descent-etale-fppf}，
$Rh_{n, *}h_n^{-1}$ 是 $D^+(U_{n, \etale})$ 上的恒同函子。
因此，由引理 \ref{spaces-simplicial-lemma-direct-image-morphism-simplicial-sites}，
$Rh_*h^{-1}$ 是 $D^+(U_\etale)$ 上的恒同函子。我们有
\begin{align*}
Ra_*(a^{-1}K)
& =
Ra_*Rh_*h^{-1}a^{-1}K \\
& =
Rh_{-1, *}Ra_{fppf, *}a_{fppf}^{-1}(h_{-1})^{-1}K \\
& =
Rh_{-1, *}(h_{-1})^{-1}K \\
& =
K
\end{align*}
% P10-E0228
第一个等号来自上述讨论；第二个等号来自
% P10-E0238
引理 \ref{spaces-simplicial-lemma-compare-simplicial-objects} 中图的交换性；
第三个等号来自引理
\ref{spaces-simplicial-lemma-hypercovering-X-simple-descent-bounded-abelian}，
因为 $U$ 是 $(\textit{Spaces}/S)_{fppf}$ 中 $X$ 的超覆盖；
最后一个等号来自已经使用过的《空间的上同调进阶》引理
\ref{spaces-more-cohomology-lemma-cohomological-descent-etale-fppf}。
\end{proof}

\begin{lemma}
\label{spaces-simplicial-lemma-compute-via-fppf-hypercovering}
设 $S$ 是概形，$X$ 是 $S$ 上的代数空间，
$U$ 是 $S$ 上的单纯代数空间，$a : U \to X$ 是增广。
若 $a : U \to X$ 是 $X$ 的 fppf 超覆盖，则
$$
R\Gamma(X_\etale, K) = R\Gamma(U_\etale, a^{-1}K)
$$
其中 $K \in D^+(X_\etale)$。这里
$a : \Sh(U_\etale) \to \Sh(X_\etale)$ 如第
\ref{spaces-simplicial-section-simplicial-algebraic-spaces} 节所述。
\end{lemma}

\begin{proof}
这由引理 \ref{spaces-simplicial-lemma-cohomological-descent-for-fppf-hypercovering} 得出，
因为按《站点上的上同调》评注
\ref{sites-cohomology-remark-before-Leray}，有
$R\Gamma(U_\etale, -) = R\Gamma(X_\etale, -) \circ Ra_*$。
\end{proof}

\begin{lemma}
\label{spaces-simplicial-lemma-fppf-hypercovering-equivalence-bounded}
设 $S$ 是概形，$X$ 是 $S$ 上的代数空间，
$U$ 是 $S$ 上的单纯代数空间，$a : U \to X$ 是增广。
以 $\mathcal{A} \subset \textit{Ab}(U_\etale)$
表示笛卡尔阿贝尔层组成的弱 Serre 子范畴。
若 $U$ 是 $X$ 的 fppf 超覆盖，则函子 $a^{-1}$ 给出等价
$$
D^+(X_\etale) \longrightarrow D_\mathcal{A}^+(U_\etale)
$$
，其拟逆为 $Ra_*$。这里 $a : \Sh(U_\etale) \to \Sh(X_\etale)$
如第 \ref{spaces-simplicial-section-simplicial-algebraic-spaces} 节所述。
\end{lemma}

\begin{proof}
由引理 \ref{spaces-simplicial-lemma-Serre-subcat-cartesian-modules}，
$\mathcal{A}$ 是弱 Serre 子范畴。该等价是迄今所得结果的形式推论。
应用引理 \ref{spaces-simplicial-lemma-descent-sheaves-for-fppf-hypercovering}、
\ref{spaces-simplicial-lemma-cohomological-descent-for-fppf-hypercovering} 以及
《站点上的上同调》引理 \ref{sites-cohomology-lemma-equivalence-bounded}。
\end{proof}

\begin{lemma}
\label{spaces-simplicial-lemma-spectral-sequence-fppf-hypercovering}
设 $S$ 是概形，$X$ 是 $S$ 上的代数空间，
$U$ 是 $S$ 上的单纯代数空间，$a : U \to X$ 是增广。
设 $\mathcal{F}$ 是 $X_\etale$ 上的阿贝尔层，$\mathcal{F}_n$
是它到 $U_{n, \etale}$ 的拉回。若 $U$ 是 $X$ 的 fppf 超覆盖，
则存在典范谱序列
$$
E_1^{p, q} = H^q_\etale(U_p, \mathcal{F}_p)
$$
，它收敛到 $H^{p + q}_\etale(X, \mathcal{F})$。
\end{lemma}

\begin{proof}
这是引理 \ref{spaces-simplicial-lemma-compute-via-fppf-hypercovering} 和
\ref{spaces-simplicial-lemma-simplicial-sheaf-cohomology-site} 的立即推论。
\end{proof}








\section{代数空间的 fppf 超覆盖：模}
\label{spaces-simplicial-section-fppf-hypercovering-modules}

\noindent
我们继续第 \ref{spaces-simplicial-section-fppf-hypercovering} 节所开始的关于 fppf
超覆盖的（上同调）下降之讨论，但在本节中我们讨论模层的情形。
我们主要讨论拟凝聚模；事实证明，对它们可以进行无界上同调下降。

\begin{lemma}
\label{spaces-simplicial-lemma-compare-simplicial-objects-fppf-etale-modules}
设 $S$ 是概形。设 $X$ 是 $S$ 上的代数空间。
设 $U$ 是 $S$ 上的单纯代数空间。
设 $a : U \to X$ 是增广。存在如下赋环拓扑斯交换图
$$
\xymatrix{
(\Sh((\textit{Spaces}/U)_{fppf, total}), \mathcal{O}_{big, total})
\ar[r]_-h \ar[d]_{a_{fppf}} &
(\Sh(U_\etale), \mathcal{O}_U) \ar[d]^a \\
(\Sh((\textit{Spaces}/X)_{fppf}), \mathcal{O}_{big}) \ar[r]^-{h_{-1}} &
(\Sh(X_\etale), \mathcal{O}_X)
}
$$
其中左侧竖直箭头定义于第 \ref{spaces-simplicial-section-hypercovering-modules} 节，
右侧竖直箭头定义于第 \ref{spaces-simplicial-section-simplicial-algebraic-spaces} 节。
\end{lemma}

\begin{proof}
关于底层拓扑斯图，参见引理
\ref{spaces-simplicial-lemma-compare-simplicial-objects-fppf-etale} 的证明中的讨论。
层 $\mathcal{O}_U$ 是第 \ref{spaces-simplicial-section-simplicial-algebraic-spaces} 节
所定义的单纯代数空间 $U$ 的结构层。层 $\mathcal{O}_X$
是代数空间 $X$ 的通常结构层。环层 $\mathcal{O}_{big, total}$
和 $\mathcal{O}_{big}$ 以第 \ref{spaces-simplicial-section-hypercovering-modules} 节
所说明的方式来自 $(\textit{Spaces}/S)_{fppf}$ 上的结构层；
该节还将 $a_{fppf}$ 构造为赋环拓扑斯的态射。
由空间上上同调进阶的第
\ref{spaces-more-cohomology-section-fppf-etale-modules} 节，
分量态射 $h_n = a_{U_n}$ 和 $h_{-1} = a_X$
是赋环拓扑斯的态射。最后，用于定义 $h$ 的连续函子
$u : U_{spaces, \etale} \to (\textit{Spaces}/U)_{fppf, total}$
\footnote{这是在引理
\ref{spaces-simplicial-lemma-compare-simplicial-objects-fppf-etale} 的证明中，
通过应用引理 \ref{spaces-simplicial-lemma-morphism-augmentation-simplicial-sites}
完成的。}
由 $V/U_n \mapsto V/U_n$ 给出，故
$h_*\mathcal{O}_{big, total} = \mathcal{O}_U$。
我们据此按注 \ref{spaces-simplicial-remark-morphism-simplicial-sites-modules}
赋予 $h$ 赋环单纯站点之态射的结构。
于是由引理 \ref{spaces-simplicial-lemma-morphism-simplicial-sites-modules}，
得到作为赋环拓扑斯态射的 $h$。请注意，态射 $h_n$
确实与上面所述的态射 $a_{U_n}$ 一致。
我们略去该图作为赋环拓扑斯之图的交换性验证
（我们已经知道它作为拓扑斯之图是交换的）。
\end{proof}

\begin{lemma}
\label{spaces-simplicial-lemma-descent-qcoh-for-fppf-hypercovering}
设 $S$ 是概形，$X$ 是 $S$ 上的代数空间。
设 $U$ 是 $S$ 上的单纯代数空间，$a : U \to X$ 是增广。
若 $a : U \to X$ 是 $X$ 的 fppf 超覆盖，则
$$
a^* : \QCoh(\mathcal{O}_X) \to \QCoh(\mathcal{O}_U)
$$
% P10-E0239
是一个全忠实的等价，其拟逆由 $a_*$ 给出。
这里 $a : \Sh(U_\etale) \to \Sh(X_\etale)$
如第 \ref{spaces-simplicial-section-simplicial-algebraic-spaces} 节所述。
\end{lemma}

\begin{proof}
考虑引理
\ref{spaces-simplicial-lemma-compare-simplicial-objects-fppf-etale-modules} 中的图。
在该引理的证明中，我们已看到 $h_{-1}$ 是空间上上同调进阶的
第 \ref{spaces-more-cohomology-section-fppf-etale-modules} 节
中的态射 $a_X$。因此，由空间上上同调进阶的引理
\ref{spaces-more-cohomology-lemma-review-quasi-coherent}，
$$
(h_{-1})^* :
\QCoh(\mathcal{O}_X)
\longrightarrow
\QCoh(\mathcal{O}_{big})
$$
是以 $h_{-1, *}$ 为拟逆的等价。对 $h$ 的分量 $h_n$
同样成立。回忆 $\QCoh(\mathcal{O}_U)$ 和
$\QCoh(\mathcal{O}_{big, total})$ 由各分量拟凝聚的笛卡尔模组成，
见引理 \ref{spaces-simplicial-lemma-quasi-coherent-sheaf}。
由于引理 \ref{spaces-simplicial-lemma-morphism-simplicial-sites-modules} 中的函子
$h^*$ 和 $h_*$ 在各分量上与函子 $h_n^*$ 和 $h_{n, *}$ 一致，
我们得到
$$
h^* :
\QCoh(\mathcal{O}_U)
\longrightarrow 
\QCoh(\mathcal{O}_{big, total})
$$
是以 $h_*$ 为拟逆的等价。注意，按第
\ref{spaces-simplicial-section-hypercovering} 节的定义，$U$ 是
$(\textit{Spaces}/S)_{fppf}$ 中 $X$ 的超覆盖。
由引理 \ref{spaces-simplicial-lemma-hypercovering-X-simple-descent-modules}，
$a_{fppf}^*$ 是全忠实的，拟逆为 $a_{fppf, *}$，
且其本质像为 $\mathcal{O}_{fppf, total}$-模的笛卡尔层。
因此，由引理 \ref{spaces-simplicial-lemma-quasi-coherent-sheaf} 对
$\QCoh(\mathcal{O}_{big})$ 和
$\QCoh(\mathcal{O}_{big, total})$ 的描述，我们得到等价
$$
a_{fppf}^* :
\QCoh(\mathcal{O}_{big})
\longrightarrow
\QCoh(\mathcal{O}_{big, total})
$$
，其拟逆由 $a_{fppf, *}$ 给出。现在，一个形式论证
（沿图追逐）表明，$a^*$ 在 $\QCoh(\mathcal{O}_X)$ 上全忠实，
且其像包含在 $\QCoh(\mathcal{O}_U)$ 中。

\medskip\noindent
最后，设 $\mathcal{G}$ 属于 $\QCoh(\mathcal{O}_U)$。
则 $h^*\mathcal{G}$ 属于 $\QCoh(\mathcal{O}_{big, total})$。
因此 $h^*\mathcal{G} = a_{fppf}^*\mathcal{H}$，其中
$\mathcal{H} = a_{fppf, *}h^*\mathcal{G}$
属于 $\QCoh(\mathcal{O}_{big})$（见上文）。
进而，$\mathcal{H} = (h_{-1})^*\mathcal{F}$，其中
$\mathcal{F} = h_{-1, *}\mathcal{H}$ 属于
$\QCoh(\mathcal{O}_X)$。沿该图追逐可得
$h^*\mathcal{G} \cong h^*a^*\mathcal{F}$。
由 $h^*$ 的全忠实性，我们得到
$a^*\mathcal{F} \cong \mathcal{G}$。
又因为 $\mathcal{F} = h_{-1, *}a_{fppf, *}h^*\mathcal{G} =
a_*h_*h^*\mathcal{G} = a_*\mathcal{G}$，
所以也得到拟逆由 $a_*$ 给出的断言。
\end{proof}

\begin{lemma}
\label{spaces-simplicial-lemma-cohomological-descent-qcoh-for-fppf-hypercovering}
设 $S$ 是概形，$X$ 是 $S$ 上的代数空间。
设 $U$ 是 $S$ 上的单纯代数空间，$a : U \to X$ 是增广。
若 $a : U \to X$ 是 $X$ 的 fppf 超覆盖，则对任意拟凝聚
$\mathcal{O}_X$-模 $\mathcal{F}$，映射
$$
\mathcal{F} \to Ra_*(a^*\mathcal{F})
$$
是同构。这里 $a : \Sh(U_\etale) \to \Sh(X_\etale)$
如第 \ref{spaces-simplicial-section-simplicial-algebraic-spaces} 节所述。
\end{lemma}

\begin{proof}
% P10-E0238
考虑引理 \ref{spaces-simplicial-lemma-compare-simplicial-objects-fppf-etale} 中的图。
令 $\mathcal{F}_n = a_n^*\mathcal{F}$ 为
$a^*\mathcal{F}$ 的第 $n$ 个分量。它是拟凝聚
$\mathcal{O}_{U_n}$-模。由空间上上同调进阶的引理
\ref{spaces-more-cohomology-lemma-cohomological-descent-etale-fppf-modules}，
有 $\mathcal{F}_n = Rh_{n, *}h_n^*\mathcal{F}_n$。
因而由引理 \ref{spaces-simplicial-lemma-direct-image-morphism-simplicial-sites-modules}，
$a^*\mathcal{F} = Rh_*h^*a^*\mathcal{F}$。我们有
\begin{align*}
Ra_*(a^*\mathcal{F})
& =
Ra_*Rh_*h^*a^*\mathcal{F} \\
& =
Rh_{-1, *}Ra_{fppf, *}a_{fppf}^*(h_{-1})^*\mathcal{F} \\
& =
Rh_{-1, *}(h_{-1})^*\mathcal{F} \\
& =
\mathcal{F}
\end{align*}
% P10-E0228
第一个等号来自上面的讨论；第二个等号来自
引理 \ref{spaces-simplicial-lemma-compare-simplicial-objects} 中图的交换性；
第三个等号来自引理
\ref{spaces-simplicial-lemma-hypercovering-X-simple-descent-bounded-modules}，
因为 $U$ 是 $(\textit{Spaces}/S)_{fppf}$ 中 $X$ 的超覆盖，
且 $a_{fppf}$ 平坦，从而
$La_{fppf}^* = a_{fppf}^*$
（即 $a_{fppf}^{-1}\mathcal{O}_{big} =
\mathcal{O}_{big, total}$，见注 \ref{spaces-simplicial-remark-augmentation-ringed}）；
最后一个等号来自已经用过的空间上上同调进阶的引理
\ref{spaces-more-cohomology-lemma-cohomological-descent-etale-fppf-modules}。
\end{proof}

\begin{lemma}
\label{spaces-simplicial-lemma-coh-descent-qcoh-unbounded-for-fppf-hypercovering}
设 $S$ 是概形，$X$ 是 $S$ 上的代数空间。
设 $U$ 是 $S$ 上的单纯代数空间，$a : U \to X$ 是增广。
假设 $a : U \to X$ 是 $X$ 的 fppf 超覆盖。
则 $\QCoh(\mathcal{O}_U)$ 是
$\textit{Mod}(\mathcal{O}_U)$ 的弱 Serre 子范畴，且
$$
a^* : D_\QCoh(\mathcal{O}_X) \longrightarrow D_\QCoh(\mathcal{O}_U)
$$
是范畴等价，其拟逆由 $Ra_*$ 给出。
这里 $a : \Sh(U_\etale) \to \Sh(X_\etale)$
如第 \ref{spaces-simplicial-section-simplicial-algebraic-spaces} 节所述。
\end{lemma}

\begin{proof}
首先注意，由超覆盖的注 \ref{hypercovering-remark-P-covering}，
映射 $a_n : U_n \to X$ 和 $d^n_i : U_n \to U_{n - 1}$
是平坦、局部有限呈示且满的。

\medskip\noindent
回忆，$\mathcal{O}_U$-模 $\mathcal{F}$ 拟凝聚，当且仅当
它是笛卡尔的，并且对所有 $n$，$\mathcal{F}_n$ 都拟凝聚；
见引理 \ref{spaces-simplicial-lemma-quasi-coherent-sheaf}。
由引理 \ref{spaces-simplicial-lemma-Serre-subcat-cartesian-modules}
（以及上面所示映射 $d^n_i : U_n \to U_{n - 1}$ 的平坦性），
% P10-E0240
笛卡尔模构成 $\textit{Mod}(\mathcal{O}_U)$ 的弱 Serre 子范畴。
另一方面，对每个 $n$，
$\QCoh(\mathcal{O}_{U_n}) \subset \textit{Mod}(\mathcal{O}_{U_n})$
是弱 Serre 子范畴（空间的性质，引理
\ref{spaces-properties-lemma-properties-quasi-coherent}）。
合并二者可见
$\QCoh(\mathcal{O}_U) \subset \textit{Mod}(\mathcal{O}_U)$
是弱 Serre 子范畴。

\medskip\noindent
为完成证明，我们逐一验证站点上的上同调的引理
\ref{sites-cohomology-lemma-equivalence-unbounded-one}
中的条件 (1) -- (5)。

\medskip\noindent
% P10-E0241
关于 (1)。由引理 \ref{spaces-simplicial-lemma-flat-augmentation-modules}，
$a_n$ 平坦（见上文）蕴含 $a$ 平坦，故该条件成立。

\medskip\noindent
关于 (2)。这正是引理
\ref{spaces-simplicial-lemma-descent-qcoh-for-fppf-hypercovering} 的内容。

\medskip\noindent
关于 (3)。这正是引理
\ref{spaces-simplicial-lemma-cohomological-descent-qcoh-for-fppf-hypercovering} 的内容。

\medskip\noindent
关于 (4)。回忆，我们既可用站点 $U_\etale$，也可用
$U_{spaces, \etale}$ 来定义小平展拓扑斯
$\Sh(U_\etale)$；见第 \ref{spaces-simplicial-section-simplicial-algebraic-spaces} 节。
站点上的上同调的情形
\ref{sites-cohomology-situation-olsson-laszlo} 的假设对三元组
$(U_{spaces, \etale}, \mathcal{O}_U, \QCoh(\mathcal{O}_U))$
成立；同理也对三元组
$(U_\etale, \mathcal{O}_U, \QCoh(\mathcal{O}_U))$.
成立。事实上，取
$$
\mathcal{B} \subset \Ob(U_\etale) \subset \Ob(U_{spaces, \etale})
$$
为仿射对象之集。对 $V/U_n \in \mathcal{B}$，
取 $d_{V/U_n} = 0$，并令 $\text{Cov}_{V/U_n}$ 为
$V_i$ 仿射的平展覆盖 $\{V_i \to V\}$ 之集。
于是得到所需的消失，因为对
$\mathcal{F} \in \QCoh(\mathcal{O}_U)$
和任意 $V/U_n \in \mathcal{B}$，由引理
\ref{spaces-simplicial-lemma-sanity-check-modules} 有
$$
H^p(V/U_n, \mathcal{F}) = H^p(V, \mathcal{F}_n)
$$
。右端是 $U_n$ 上拟凝聚层 $\mathcal{F}_n$
在 $U_{n, \etale}$ 的仿射对象 $V$ 上的上同调。
由空间的上同调的第
\ref{spaces-cohomology-section-higher-direct-image} 节中的讨论
以及概形的上同调的引理
\ref{coherent-lemma-quasi-coherent-affine-cohomology-zero}，
当 $p > 0$ 时它消失。

\medskip\noindent
关于 (5)。取 $\mathcal{B} \subset \Ob(X_{spaces, \etale})$
为仿射对象之集，并应用上面给出的文献，即得该条件。
\end{proof}

\begin{lemma}
\label{spaces-simplicial-lemma-compute-via-fppf-hypercovering-modules}
设 $S$ 是概形，$X$ 是 $S$ 上的代数空间。
设 $U$ 是 $S$ 上的单纯代数空间，$a : U \to X$ 是增广。
若 $a : U \to X$ 是 $X$ 的 fppf 超覆盖，则
$$
R\Gamma(X_\etale, K) = R\Gamma(U_\etale, a^*K)
$$
对 $K \in D_\QCoh(\mathcal{O}_X)$ 成立。
这里 $a : \Sh(U_\etale) \to \Sh(X_\etale)$
如第 \ref{spaces-simplicial-section-simplicial-algebraic-spaces} 节所述。
\end{lemma}

\begin{proof}
这由引理
\ref{spaces-simplicial-lemma-coh-descent-qcoh-unbounded-for-fppf-hypercovering} 得出，
因为由站点上的上同调的注
\ref{sites-cohomology-remark-before-Leray}，
$R\Gamma(U_\etale, -) = R\Gamma(X_\etale, -) \circ Ra_*$。
\end{proof}

\begin{lemma}
\label{spaces-simplicial-lemma-spectral-sequence-fppf-hypercovering-modules}
设 $S$ 是概形，$X$ 是 $S$ 上的代数空间。
设 $U$ 是 $S$ 上的单纯代数空间，$a : U \to X$ 是增广。
% P10-E0242
设 $\mathcal{F}$ 是拟凝聚 $\mathcal{O}_X$-模，
并设 $\mathcal{F}_n$ 是它在 $U_{n, \etale}$ 上的拉回。
若 $U$ 是 $X$ 的 fppf 超覆盖，则存在典范谱序列
$$
E_1^{p, q} = H^q_\etale(U_p, \mathcal{F}_p)
$$
，它收敛于 $H^{p + q}_\etale(X, \mathcal{F})$。
\end{lemma}

\begin{proof}
这是引理 \ref{spaces-simplicial-lemma-compute-via-fppf-hypercovering-modules}
和 \ref{spaces-simplicial-lemma-simplicial-module-cohomology-site} 的直接推论。
\end{proof}







\section{复形的 fppf 下降}
\label{spaces-simplicial-section-fppf-descent-derived}

\noindent
在本节中，我们把前面证明的若干结果汇集起来，
用于代数空间的 fppf 覆盖和拟凝聚模的导出范畴。

\begin{lemma}
\label{spaces-simplicial-lemma-fppf-neg-ext-zero-hom}
设 $X$ 是概形 $S$ 上的代数空间。
设 $K, E \in D_\QCoh(\mathcal{O}_X)$。
设 $a : U \to X$ 是 fppf 超覆盖。
假设对所有 $n \geq 0$ 都有
$$
\Ext_{\mathcal{O}_{U_n}}^i(La_n^*K, La_n^*E) = 0
\text{，当 } i < 0
$$
。则
\begin{enumerate}
\item 当 $i < 0$ 时，$\Ext_{\mathcal{O}_X}^i(K, E) = 0$；并且
\item 存在正合序列
$$
0
\to
\Hom_{\mathcal{O}_X}(K, E)
\to
\Hom_{\mathcal{O}_{U_0}}(La_0^*K, La_0^*E)
\to
\Hom_{\mathcal{O}_{U_1}}(La_1^*K, La_1^*E)
$$
\end{enumerate}
\end{lemma}

\begin{proof}
记 $K_n = La_n^*K$ 且 $E_n = La_n^*E$。
那么，它们分别是与 $La^*K$ 和 $La^*E$ 关联的模之导出范畴的
单纯系统（定义 \ref{spaces-simplicial-definition-cartesian-derived-modules} 和
引理 \ref{spaces-simplicial-lemma-cartesian-objects-derived-modules}），
其中 $a : U_\etale \to X_\etale$
如第 \ref{spaces-simplicial-section-simplicial-algebraic-spaces} 节所述。
先证明 (2)。由引理
\ref{spaces-simplicial-lemma-coh-descent-qcoh-unbounded-for-fppf-hypercovering}，
$$
\Hom_{\mathcal{O}_X}(K, E) =
\Hom_{\mathcal{O}_U}(La^*K, La^*E)
$$
。因此，该序列具有如下形式：
$$
0
\to
\Hom_{\mathcal{O}_U}(La^*K, La^*E)
\to
\Hom_{\mathcal{O}_{U_0}}(K_0, E_0)
\to
\Hom_{\mathcal{O}_{U_1}}(K_1, E_1)
$$
由引理 \ref{spaces-simplicial-lemma-nullity-cartesian-modules-derived}，
第一个箭头是单射。由引理
\ref{spaces-simplicial-lemma-hom-cartesian-modules-derived}，
该箭头的像是第二个箭头的核。这就完成了 (2) 的证明。
对 $i > 0$ 将 (2) 应用于 $K[i]$ 和 $E$，即得 (1)。
\end{proof}

\begin{lemma}
\label{spaces-simplicial-lemma-fppf-glue-neg-ext-zero}
设 $X$ 是概形 $S$ 上的代数空间。
设 $a : U \to X$ 是 fppf 超覆盖。
假设给定 $K_0 \in D_\QCoh(U_0)$ 和同构
$$
\alpha :
L(f_{\delta_1^1})^*K_0
\longrightarrow
L(f_{\delta_0^1})^*K_0
$$
% P10-E0243
，它在 $U_1$ 上满足余圈条件。令
$\tau^n_i : [0] \to [n]$，$0 \mapsto i$，并令
$K_n = Lf_{\tau^n_n}^*K_0$。
假设当 $i < 0$ 时
$\Ext^i_{\mathcal{O}_{U_n}}(K_n, K_n) = 0$。
% P10-E0244
则存在对象 $K \in D_\QCoh(\mathcal{O}_X)$
以及与 $\alpha$ 相容的同构 $La_0^*K \to K$。
\end{lemma}

\begin{proof}
由引理 \ref{spaces-simplicial-lemma-construct-cartesian-objects-derived-modules}，
对象 $K_n$ 构成模之导出范畴的一个单纯系统的各成员。
于是得到对象
$K' \in D_\QCoh(\mathcal{O}_{U_\etale})$，
使得 $(K_n, K_\varphi)$ 是由 $K'$ 推出的系统；见引理
\ref{spaces-simplicial-lemma-cartesian-module-derived-from-simplicial}。
最后，应用引理
\ref{spaces-simplicial-lemma-coh-descent-qcoh-unbounded-for-fppf-hypercovering}，
可见对某个 $K \in D_\QCoh(\mathcal{O}_X)$，
$K' = La^*K$，这正是所需的。
\end{proof}











\section{代数空间的真超覆盖}
\label{spaces-simplicial-section-proper-hypercovering-spaces}

\noindent
本节给出代数空间情形下第
\ref{spaces-simplicial-section-proper-hypercovering} 节的对应版本。
愿意如此处理的读者可以在各处把“代数空间”替换为“概形”，
所得结果同样有效。这样做的好处是，可以把对空间上上同调进阶的
第 \ref{spaces-more-cohomology-section-ph-etale} 节的引用，
替换为对平展上同调的第
\ref{etale-cohomology-section-ph-etale} 节的引用。

\medskip\noindent
固定一个基概形 $S$。设 $X$ 是 $S$ 上的代数空间，
$U$ 是 $S$ 上的单纯代数空间。假设有增广
$$
a : U \to X
$$
；见第 \ref{spaces-simplicial-section-simplicial-algebraic-spaces} 节。
如果满足下列条件，就称 $U$ 是 $X$ 的{\it 真超覆盖}：
\begin{enumerate}
\item $U_0 \to X$ 是真且满的；
\item $U_1 \to U_0 \times_X U_0$ 是真且满的；
\item $U_{n + 1} \to (\text{cosk}_n\text{sk}_n U)_{n + 1}$
对 $n \geq 1$ 是真且满的。
\end{enumerate}
$S$ 上代数空间的范畴具有所有有限极限，因而上述表述中使用的
余骨架存在。
$$
\fbox{原理：真超覆盖可用于计算平展上同调。}
$$
证明该原理的关键思想，是比较范畴 $\textit{Spaces}/S$
上的 ph 拓扑与平展拓扑。具体地说，ph 拓扑强于平展拓扑，并且
(a) 真满映射给出 ph 覆盖；以及
(b) 从小平展站点拉回的层的 ph 上同调与平展上同调一致，
正如我们在空间上上同调进阶的第
\ref{spaces-more-cohomology-section-ph-etale} 节中所见。

\medskip\noindent
本节的所有结果都可推广到 $U \to X$ 仅为“ph 超覆盖”的情形，
即按第 \ref{spaces-simplicial-section-hypercovering} 节的定义，
它是站点 $(\textit{Spaces}/S)_{ph}$ 中 $X$ 的超覆盖。
若将来需要这一点，我们会在此处精确表述并证明它。

\begin{lemma}
\label{spaces-simplicial-lemma-compare-simplicial-objects-ph-etale}
设 $S$ 是概形，$X$ 是 $S$ 上的代数空间。
设 $U$ 是 $S$ 上的单纯代数空间。
设 $a : U \to X$ 是增广。存在交换图
$$
\xymatrix{
\Sh((\textit{Spaces}/U)_{ph, total}) \ar[r]_-h \ar[d]_{a_{ph}} &
\Sh(U_\etale) \ar[d]^a \\
\Sh((\textit{Spaces}/X)_{ph}) \ar[r]^-{h_{-1}} &
\Sh(X_\etale)
}
$$
，其中左侧竖直箭头定义于第 \ref{spaces-simplicial-section-hypercovering} 节，
右侧竖直箭头定义于第
\ref{spaces-simplicial-section-simplicial-algebraic-spaces} 节。
\end{lemma}

\begin{proof}
记号 $(\textit{Spaces}/U)_{ph, total}$ 表示我们把
第 \ref{spaces-simplicial-section-hypercovering} 节的构造用于站点
$(\textit{Spaces}/S)_{ph}$ 及其单纯对象 $U$
\footnote{这是为了与第 \ref{spaces-simplicial-section-fppf-hypercovering} 节中
用 fppf 拓扑定义的
$(\textit{Spaces}/U)_{fppf, total}$ 相区别。}。
对右侧的拓扑斯，我们将使用站点
$X_{spaces, \etale}$ 和 $U_{spaces, \etale}$；
这是容许的，见第 \ref{spaces-simplicial-section-simplicial-algebraic-spaces} 节中的讨论。

\medskip\noindent
注意，$(\textit{Spaces}/U)_{ph, total}$ 和
$U_{spaces, \etale}$ 都属于情形
\ref{spaces-simplicial-situation-simplicial-site} 的情形 A。
对 $U_\etale$，这直接来自第
\ref{spaces-simplicial-section-simplicial-algebraic-spaces} 节中的构造；
对 $(\textit{Spaces}/U)_{ph, total}$，则由引理
\ref{spaces-simplicial-lemma-sr-when-fibre-products} 得出。
接下来考虑函子
$U_{n, spaces, \etale} \to (\textit{Spaces}/U_n)_{ph}$, $U \mapsto U/U_n$
以及
$X_{spaces, \etale} \to (\textit{Spaces}/X)_{ph}$, $U \mapsto U/X$.
。我们已在空间上上同调进阶的第
\ref{spaces-more-cohomology-section-ph-etale} 节中看到，
它们定义站点的态射；在那里这些态射记为
% P10-E0168
% P10-E0245
$a_{U_n} = \epsilon_{U_n} \circ \pi_{u_n}$
和 $a_X = \epsilon_X \circ \pi_X$。
因此，如注 \ref{spaces-simplicial-remark-morphism-augmentation-simplicial-sites}，
我们得到与增广相容的单纯站点态射；
应用引理 \ref{spaces-simplicial-lemma-morphism-augmentation-simplicial-sites}
即可得出结论。
\end{proof}

\begin{lemma}
\label{spaces-simplicial-lemma-descent-sheaves-for-ph-hypercovering}
设 $S$ 是概形，$X$ 是 $S$ 上的代数空间。
设 $U$ 是 $S$ 上的单纯代数空间，$a : U \to X$ 是增广。
若 $a : U \to X$ 是 $X$ 的真超覆盖，则
$$
a^{-1} : \Sh(X_\etale) \to \Sh(U_\etale)
\quad\text{且}\quad
a^{-1} : \textit{Ab}(X_\etale) \to \textit{Ab}(U_\etale)
$$
是全忠实的，其本质像为笛卡尔层，拟逆由 $a_*$ 给出。
这里 $a : \Sh(U_\etale) \to \Sh(X_\etale)$
如第 \ref{spaces-simplicial-section-simplicial-algebraic-spaces} 节所述。
\end{lemma}

\begin{proof}
我们对集合层证明该断言。它几乎是已经建立的结果的形式推论。
考虑引理 \ref{spaces-simplicial-lemma-compare-simplicial-objects-ph-etale} 中的图。
在该引理的证明中，我们已看到 $h_{-1}$ 是空间上上同调进阶的
第 \ref{spaces-more-cohomology-section-ph-etale} 节中的态射 $a_X$。
因此，由空间上上同调进阶的引理
\ref{spaces-more-cohomology-lemma-comparison-ph-etale}，
$(h_{-1})^{-1}$ 是全忠实的，拟逆为 $h_{-1, *}$。
对 $h$ 的分量 $h_n$ 同样成立。由引理
\ref{spaces-simplicial-lemma-morphism-simplicial-sites} 对函子 $h^{-1}$ 和 $h_*$
的描述，我们得到 $h^{-1}$ 全忠实，拟逆为 $h_*$。
注意，按第 \ref{spaces-simplicial-section-hypercovering} 节的定义，
$U$ 是 $(\textit{Spaces}/S)_{ph}$ 中 $X$ 的超覆盖，
因为由空间上的拓扑的引理
\ref{spaces-topologies-lemma-surjective-proper-ph}，
真满态射给出 ph 覆盖。
由引理 \ref{spaces-simplicial-lemma-hypercovering-X-simple-descent-sheaves}，
$a_{ph}^{-1}$ 是全忠实的，拟逆为 $a_{ph, *}$，
且其本质像为 $(\textit{Spaces}/U)_{ph, total}$ 上的笛卡尔层。
现在，一个形式论证（沿图追逐）表明 $a^{-1}$ 全忠实。

\medskip\noindent
最后，设 $\mathcal{G}$ 是 $U_\etale$ 上的笛卡尔层。
则 $h^{-1}\mathcal{G}$ 是
$(\textit{Spaces}/U)_{ph, total}$ 上的笛卡尔层。
因而对 $(\textit{Spaces}/X)_{ph}$ 上的某个层
$\mathcal{H}$，有
$h^{-1}\mathcal{G} = a_{ph}^{-1}\mathcal{H}$。
用稍显繁琐的记号计算如下：
\begin{align*}
(h_{-1})^{-1}(a_*\mathcal{G})
& =
(h_{-1})^{-1}
\text{等化子}(
\xymatrix{
a_{0, small, *}\mathcal{G}_0
\ar@<1ex>[r] \ar@<-1ex>[r] &
a_{1, small, *}\mathcal{G}_1
}
) \\
& =
\text{等化子}(
\xymatrix{
(h_{-1})^{-1}a_{0, small, *}\mathcal{G}_0
\ar@<1ex>[r] \ar@<-1ex>[r] &
(h_{-1})^{-1}a_{1, small, *}\mathcal{G}_1
}
) \\
& =
\text{等化子}(
\xymatrix{
a_{0, big, ph, *}h_0^{-1}\mathcal{G}_0
\ar@<1ex>[r] \ar@<-1ex>[r] &
a_{1, big, ph, *}h_1^{-1}\mathcal{G}_1
}
) \\
& =
\text{等化子}(
\xymatrix{
a_{0, big, ph, *}(a_{0, big, ph})^{-1}\mathcal{H}
\ar@<1ex>[r] \ar@<-1ex>[r] &
a_{1, big, ph, *}(a_{1, big, ph})^{-1}\mathcal{H}
}
) \\
& =
a_{ph, *}a_{ph}^{-1}\mathcal{H} \\
& =
\mathcal{H}
\end{align*}
这里，第一个等号来自引理 \ref{spaces-simplicial-lemma-augmentation-site}；
第二个等号来自 $(h_{-1})^{-1}$ 是正合函子；
第三个等号来自空间上上同调进阶的引理
\ref{spaces-more-cohomology-lemma-proper-push-pull-ph-etale}
（这里用到 $a_0 : U_0 \to X$ 和 $a_1: U_1 \to X$ 为真）；
第四个等号来自
$a_{ph}^{-1}\mathcal{H} = h^{-1}\mathcal{G}$；
第五个等号来自引理 \ref{spaces-simplicial-lemma-augmentation-site}；
第六个等号已在上面见到。由于
$a_{ph}^{-1}\mathcal{H} = h^{-1}\mathcal{G}$，
我们推出
$h^{-1}\mathcal{G} \cong h^{-1}a^{-1}a_*\mathcal{G}$；
再由 $h^{-1}$ 的全忠实性，证明完毕。
\end{proof}

\begin{lemma}
\label{spaces-simplicial-lemma-cohomological-descent-for-ph-hypercovering}
设 $S$ 是概形，$X$ 是 $S$ 上的代数空间。
设 $U$ 是 $S$ 上的单纯代数空间，$a : U \to X$ 是增广。
若 $a : U \to X$ 是 $X$ 的真超覆盖，则对
$K \in D^+(X_\etale)$，
$$
K \to Ra_*(a^{-1}K)
$$
是同构。这里 $a : \Sh(U_\etale) \to \Sh(X_\etale)$
如第 \ref{spaces-simplicial-section-simplicial-algebraic-spaces} 节所述。
\end{lemma}

\begin{proof}
考虑引理 \ref{spaces-simplicial-lemma-compare-simplicial-objects-ph-etale} 中的图。
由空间上上同调进阶的引理
\ref{spaces-more-cohomology-lemma-cohomological-descent-etale-ph}，
$Rh_{n, *}h_n^{-1}$ 是 $D^+(U_{n, \etale})$ 上的恒等函子。
因此，由引理 \ref{spaces-simplicial-lemma-direct-image-morphism-simplicial-sites}，
$Rh_*h^{-1}$ 是 $D^+(U_\etale)$ 上的恒等函子。我们有
\begin{align*}
Ra_*(a^{-1}K)
& =
Ra_*Rh_*h^{-1}a^{-1}K \\
& =
Rh_{-1, *}Ra_{ph, *}a_{ph}^{-1}(h_{-1})^{-1}K \\
& =
Rh_{-1, *}(h_{-1})^{-1}K \\
& =
K
\end{align*}
% P10-E0228
第一个等号来自上面的讨论；第二个等号来自
% P10-E0238
引理 \ref{spaces-simplicial-lemma-compare-simplicial-objects} 中图的交换性；
第三个等号来自引理
\ref{spaces-simplicial-lemma-hypercovering-X-simple-descent-bounded-abelian}，
因为由空间上的拓扑的引理
\ref{spaces-topologies-lemma-surjective-proper-ph}，
$U$ 是 $(\textit{Spaces}/S)_{ph}$ 中 $X$ 的超覆盖；
最后一个等号来自已经用过的空间上上同调进阶的引理
\ref{spaces-more-cohomology-lemma-cohomological-descent-etale-ph}。
\end{proof}

\begin{lemma}
\label{spaces-simplicial-lemma-compute-via-ph-hypercovering}
设 $S$ 是概形，$X$ 是 $S$ 上的代数空间。
设 $U$ 是 $S$ 上的单纯代数空间，$a : U \to X$ 是增广。
若 $a : U \to X$ 是 $X$ 的真超覆盖，则
$$
R\Gamma(X_\etale, K) = R\Gamma(U_\etale, a^{-1}K)
$$
对 $K \in D^+(X_\etale)$ 成立。
这里 $a : \Sh(U_\etale) \to \Sh(X_\etale)$
如第 \ref{spaces-simplicial-section-simplicial-algebraic-spaces} 节所述。
\end{lemma}

\begin{proof}
这由引理 \ref{spaces-simplicial-lemma-cohomological-descent-for-ph-hypercovering} 得出，
因为由站点上的上同调的注
\ref{sites-cohomology-remark-before-Leray}，
$R\Gamma(U_\etale, -) = R\Gamma(X_\etale, -) \circ Ra_*$。
\end{proof}

\begin{lemma}
\label{spaces-simplicial-lemma-ph-hypercovering-equivalence-bounded}
设 $S$ 是概形，$X$ 是 $S$ 上的代数空间。
设 $U$ 是 $S$ 上的单纯代数空间，$a : U \to X$ 是增广。
以 $\mathcal{A} \subset \textit{Ab}(U_\etale)$
表示笛卡尔阿贝尔层所成的弱 Serre 子范畴。
若 $U$ 是 $X$ 的真超覆盖，则函子 $a^{-1}$ 定义等价
$$
D^+(X_\etale) \longrightarrow D_\mathcal{A}^+(U_\etale)
$$
，其拟逆为 $Ra_*$。这里
$a : \Sh(U_\etale) \to \Sh(X_\etale)$
如第 \ref{spaces-simplicial-section-simplicial-algebraic-spaces} 节所述。
\end{lemma}

\begin{proof}
由引理 \ref{spaces-simplicial-lemma-Serre-subcat-cartesian-modules}，
$\mathcal{A}$ 是弱 Serre 子范畴。
该等价是目前所得结果的形式推论。应用引理
\ref{spaces-simplicial-lemma-descent-sheaves-for-ph-hypercovering}、
\ref{spaces-simplicial-lemma-cohomological-descent-for-ph-hypercovering}
以及站点上的上同调的引理
\ref{sites-cohomology-lemma-equivalence-bounded}。
\end{proof}

\begin{lemma}
\label{spaces-simplicial-lemma-spectral-sequence-ph-hypercovering}
设 $S$ 是概形，$X$ 是 $S$ 上的代数空间。
设 $U$ 是 $S$ 上的单纯代数空间，$a : U \to X$ 是增广。
设 $\mathcal{F}$ 是 $X_\etale$ 上的阿贝尔层，
并设 $\mathcal{F}_n$ 是它在 $U_{n, \etale}$ 上的拉回。
若 $U$ 是 $X$ 的 ph 超覆盖，则存在典范谱序列
$$
E_1^{p, q} = H^q_\etale(U_p, \mathcal{F}_p)
$$
，它收敛于 $H^{p + q}_\etale(X, \mathcal{F})$。
\end{lemma}

\begin{proof}
这是引理 \ref{spaces-simplicial-lemma-compute-via-ph-hypercovering}
和 \ref{spaces-simplicial-lemma-simplicial-sheaf-cohomology-site} 的直接推论。
\end{proof}
% Generated mechanically from the frozen spaces-duality.tex target.
% Do not edit this generated file; edit the bound locale source instead.

% OK, start here.
%

\setcounter{chapter}{85}
\chapter{空间的对偶性}



\phantomsection
\label{spaces-duality-section-phantom}


\section{引言}
\label{spaces-duality-section-introduction}

\noindent
本章是概形情形相应章节的类似版本，见《概形的对偶性》第
\ref{duality-section-introduction} 节。这里的论述与下列论文中的论述相近：
\cite{Neeman-Grothendieck}, \cite{LN},
\cite{Lipman-notes}、\cite{Neeman-improvement}。




\section{代数空间上的对偶复形}
\label{spaces-duality-section-dualizing-spaces}

\noindent
设 $U$ 是局部 Noether 概形。以 $\mathcal{O}_\etale$ 表示 $U$ 上的
结构层（定义在 $U$ 的小平展站点上）。若存在《概形的对偶性》第
\ref{duality-section-dualizing-schemes} 节意义下的对偶复形
$\omega_U^\bullet$，使得
$K = \epsilon^*(\omega_U^\bullet)$，则称对象
$K \in D_\QCoh(\mathcal{O}_\etale)$ 是 $U$ 上的对偶复形。
这里，$\epsilon^* : D_\QCoh(\mathcal{O}_U) \to
D_\QCoh(\mathcal{O}_\etale)$ 是《空间的导出范畴》引理
\ref{spaces-perfect-lemma-derived-quasi-coherent-small-etale-site}
中的等价。《概形的对偶性》第
\ref{duality-section-dualizing-schemes} 节中研究的
$\omega_U^\bullet$ 的大部分性质，经《空间的导出范畴》第
\ref{spaces-perfect-section-derived-quasi-coherent-etale} 与
\ref{spaces-perfect-section-spell-out} 节的讨论由 $K$ 继承。

\medskip\noindent
我们把局部 Noether 代数空间上的对偶复形定义为这样一种复形：
它在平展局部来自相应概形上的对偶复形。

\begin{lemma}
\label{spaces-duality-lemma-equivalent-definitions}
设 $S$ 是概形，$X$ 是 $S$ 上的局部 Noether 代数空间，
$K$ 是 $D_\QCoh(\mathcal{O}_X)$ 的对象。下列条件等价：
\begin{enumerate}
\item 对每个以概形 $U$ 为源的平展态射 $U \to X$，限制
$K|_U$ 都是 $U$ 上的对偶复形（含义如上）。
\item 存在以概形 $U$ 为源的满平展态射 $U \to X$，使得
$K|_U$ 是 $U$ 上的对偶复形。
\end{enumerate}
\end{lemma}

\begin{proof}
设 $U \to X$ 是以概形 $U$ 为源的满平展态射，并设
$V \to X$ 是以概形 $V$ 为源的平展态射。则
$$
U \leftarrow U \times_X V \rightarrow V
$$
是概形间的平展态射，其中指向 $V$ 的箭头是满的。因此，由
《概形的对偶性》引理 \ref{duality-lemma-descent-ascent} 可知：
若 $K|_U$ 是 $U$ 上的对偶复形，则 $K|_V$ 是 $V$ 上的对偶复形。
\end{proof}

\begin{definition}
\label{spaces-duality-definition-dualizing-scheme}
设 $S$ 是概形，$X$ 是 $S$ 上的局部 Noether 代数空间。
若 $D_\QCoh(\mathcal{O}_X)$ 的对象 $K$ 满足引理
\ref{spaces-duality-lemma-equivalent-definitions} 中的等价条件，则称 $K$ 为
{\it 对偶复形}。
\end{definition}

\begin{lemma}
\label{spaces-duality-lemma-affine-duality}
设 $A$ 是 Noether 环，$X = \Spec(A)$。以 $\mathcal{O}_\etale$
表示 $X$ 上的结构层（定义在 $X$ 的小平展站点上）。设 $K,L$ 是
$D(A)$ 的对象。
若 $K \in D_{\textit{Coh}}(A)$ 且 $L$ 具有有限内射维数，则
$$
\epsilon^*\widetilde{R\Hom_A(K, L)} =
R\SheafHom_{\mathcal{O}_\etale}(\epsilon^*\widetilde{K},
\epsilon^*\widetilde{L})
$$
在 $D(\mathcal{O}_\etale)$ 中成立，其中
$\epsilon : (X_\etale, \mathcal{O}_\etale) \to (X, \mathcal{O}_X)$
如《空间的导出范畴》第
\ref{spaces-perfect-section-derived-quasi-coherent-etale} 节所述。
\end{lemma}

\begin{proof}
由《概形的对偶性》引理 \ref{duality-lemma-affine-duality}，有典范同构
$$
\widetilde{R\Hom_A(K, L)} =
R\SheafHom_{\mathcal{O}_X}(\widetilde{K}, \widetilde{L})
$$
于 $D(\mathcal{O}_X)$ 中。又有典范映射
$$
\epsilon^*R\Hom_{\mathcal{O}_X}(\widetilde{K}, \widetilde{L})
\longrightarrow
R\SheafHom_{\mathcal{O}_\etale}(\epsilon^*\widetilde{K},
\epsilon^*\widetilde{L})
$$
于 $D(\mathcal{O}_\etale)$ 中，见《站点上的上同调》评注
\ref{sites-cohomology-remark-prepare-fancy-base-change}。
我们将证明该箭头两端具有同构的上同调层，但略去验证这一同构
确由该箭头给出。

\medskip\noindent
可以假设 $L$ 由内射 $A$-模组成的有限复形 $I^\bullet$ 给出。
对 $I^\bullet$ 的长度作归纳，并利用这些构造与特异三角的相容性，
可归约到 $L=I[0]$ 的情形，其中 $I$ 是内射 $A$-模。回忆
% P10-E0246
$R\SheafHom_{\mathcal{O}_\etale}(\epsilon^*\widetilde{K},
\epsilon^*\widetilde{L}))$
的上同调层是下述预层的层化：它把 $X$ 上平展的 $U$ 映到
$\epsilon^*\widetilde{K}$ 与 $\epsilon^*\widetilde{L}$ 在
$U_\etale$ 上的限制之间的第 $i$ 个 Ext 群。见《站点上的上同调》
引理 \ref{sites-cohomology-lemma-section-RHom-over-U}。
若 $U=\Spec(B)$ 是仿射的，则由《空间的导出范畴》引理
\ref{spaces-perfect-lemma-derived-quasi-coherent-small-etale-site} 及
《概形的导出范畴》引理
\ref{perfect-lemma-affine-compare-bounded}
给出的等价，该 Ext 群等于
$\text{Ext}^i_B(K \otimes_A B, L \otimes_A B)$（这里还使用了
《空间的导出范畴》评注
\ref{spaces-perfect-remark-match-total-direct-images} 中详述的相容性）。
由于 $A\to B$ 是平展的，由《对偶复形》引理
\ref{dualizing-lemma-injective-goes-up} 可知
$I\otimes_A B$ 是内射 $B$-模。因此
\begin{align*}
\Ext^n_B(K \otimes_A B, I \otimes_A B)
& =
\Hom_B(H^{-n}(K \otimes_A B), I \otimes_A B) \\
& =
% P10-E0247
\Hom_{A_f}(H^{-n}(K) \otimes_A B, I \otimes_A B) \\
& =
\Hom_A(H^{-n}(K), I) \otimes_A B \\
& =
\text{Ext}^n_A(K, I) \otimes_A B
\end{align*}
% P10-E0248
倒数第二个等号成立，是因为 $H^{-n}(K)$ 是有限 $A$-模；见
《代数详论》引理
\ref{more-algebra-lemma-pseudo-coherence-and-base-change-ext}。
因此，引理中等式两端的上同调层相同。
\end{proof}

\begin{lemma}
\label{spaces-duality-lemma-dualizing-spaces}
设 $S$ 是概形，$X$ 是 $S$ 上的局部 Noether 代数空间，
$K$ 是 $X$ 上的对偶复形。则 $K$ 是
$D_{\textit{Coh}}(\mathcal{O}_X)$ 的对象，并且
$D = R\SheafHom_{\mathcal{O}_X}(-, K)$ 诱导反等价
$$
D :
D_{\textit{Coh}}(\mathcal{O}_X)
\longrightarrow
D_{\textit{Coh}}(\mathcal{O}_X)
$$
，并配有典范同构 $\text{id}\to D\circ D$。若 $X$ 拟紧，则
$D$ 交换 $D^+_{\textit{Coh}}(\mathcal{O}_X)$ 与
$D^-_{\textit{Coh}}(\mathcal{O}_X)$，并诱导等价
$D^b_{\textit{Coh}}(\mathcal{O}_X) \to
D^b_{\textit{Coh}}(\mathcal{O}_X)$。
\end{lemma}

\begin{proof}
取以仿射概形 $U$ 为源的平展态射 $U\to X$。设 $U=\Spec(A)$，
并令 $\omega_A^\bullet$ 为与 $K|_U$ 对应的 $A$ 的对偶复形，
其含义如引理 \ref{spaces-duality-lemma-equivalent-definitions} 及《概形的对偶性》
引理 \ref{duality-lemma-equivalent-definitions} 所述。由引理
\ref{spaces-duality-lemma-affine-duality}，下图交换：
$$
\xymatrix{
D_{\textit{Coh}}(A) \ar[r] \ar[d]_{R\Hom_A(-, \omega_A^\bullet)} &
D_{\textit{Coh}}(\mathcal{O}_\etale)
\ar[d]^{R\SheafHom_{\mathcal{O}_\etale}(-, K|_U)} \\
D_{\textit{Coh}}(A) \ar[r] &
D(\mathcal{O}_\etale)
}
$$
其中 $\mathcal{O}_\etale$ 是 $U$ 的小平展站点上的结构层。
由于 $R\SheafHom$ 的形成与限制相交换，可知 $D$ 把
$D_{\textit{Coh}}(\mathcal{O}_X)$ 映入
$D_{\textit{Coh}}(\mathcal{O}_X)$。此外，典范映射
$$
L \longrightarrow
R\SheafHom_{\mathcal{O}_X}(R\SheafHom_{\mathcal{O}_X}(L, K), K)
$$
（《站点上的上同调》引理
\ref{sites-cohomology-lemma-internal-hom-evaluate}）
对 $D_{\textit{Coh}}(\mathcal{O}_X)$ 中每个 $L$ 都是同构；这是因为，
由《对偶复形》引理 \ref{dualizing-lemma-dualizing}，在上述每个 $U$
上这一结论都成立。拟紧情形下关于函子 $D$ 的有界性性质，也由
《对偶复形》引理 \ref{dualizing-lemma-dualizing} 中的相应结论推出。
\end{proof}

\noindent
设 $(\mathcal{C},\mathcal{O})$ 是赋环站点。回忆：若
$D(\mathcal{O})$ 的对象 $L$ 对于由导出张量积给出的
% P10-E0249
$D(\mathcal{O}_X)$ 上的对称幺半结构是可逆对象，则称它
{\it 可逆}。由《站点上的上同调》引理
\ref{sites-cohomology-lemma-invertible-derived} 可知，
% P10-E0250
这等价于 $L$ 是完美的；并且若 $(\mathcal{C},\mathcal{O})$
是局部赋环站点，则对 $\mathcal{C}$ 的每个对象 $U$，存在
$\mathcal{C}$ 中 $U$ 的覆盖 $\{U_i\to U\}$，使得对某些整数 $n_i$ 有
$L|_{U_i}\cong\mathcal{O}_{U_i}[-n_i]$。

\medskip\noindent
设 $S$ 是概形，$X$ 是 $S$ 上的代数空间。若
$D(\mathcal{O}_X)$ 中的 $L$ 可逆，则存在不交并分解
$X=\coprod_{n\in\mathbf{Z}}X_n$，使得 $L|_{X_n}$ 是位于次数 $n$
的可逆模。特别地，$L=\bigoplus H^n(L)[-n]$；这给出了一个表示 $L$
的良定义 $\mathcal{O}_X$-模复形（其微分均为零）。

\begin{lemma}
\label{spaces-duality-lemma-dualizing-unique-spaces}
设 $S$ 是概形，$X$ 是 $S$ 上的局部 Noether 代数空间。
若 $K$ 与 $K'$ 是 $X$ 上的对偶复形，则 $K'$ 同构于
$K\otimes_{\mathcal{O}_X}^\mathbf{L}L$，其中 $L$ 是
$D(\mathcal{O}_X)$ 中的某个可逆对象。
\end{lemma}

\begin{proof}
置
$$
L = R\SheafHom_{\mathcal{O}_X}(K, K')
$$
它是 $D(\mathcal{O}_X)$ 中的可逆对象，因为这一点在仿射局部成立；
可使用引理 \ref{spaces-duality-lemma-affine-duality} 以及《对偶复形》引理
\ref{dualizing-lemma-dualizing-unique} 及其证明。由同样的理由，
求值映射 $L\otimes_{\mathcal{O}_X}^\mathbf{L}K\to K'$ 是同构。
\end{proof}

\begin{lemma}
\label{spaces-duality-lemma-dimension-function-scheme}
设 $S$ 是概形，$X$ 是 $S$ 上局部 Noether 且拟分离的代数空间，
$\omega_X^\bullet$ 是 $X$ 上的对偶复形。则对 $X$，函数
% P10-E0251
$|X| \to \mathbf{Z}$ 由下式定义：
$$
x \longmapsto \delta(x)\text{，使得 }
\omega_{X, \overline{x}}^\bullet[-\delta(x)]
\text{ 是下列环上的正规化对偶复形：}
\mathcal{O}_{X, \overline{x}}
$$
该函数是 $|X|$ 上的维数函数。
\end{lemma}

\begin{proof}
取概形 $U$ 及满平展态射 $U\to X$。令 $\omega_U^\bullet$ 为 $U$ 上
与 $\omega_X^\bullet|_U$ 对应的对偶复形。若 $u\in U$ 映到
$x\in|X|$，则 $\mathcal{O}_{X,\overline{x}}$ 是
$\mathcal{O}_{U,u}$ 的严格 Hensel 化。由《对偶复形》引理
\ref{dualizing-lemma-flat-unramified} 可知，若 $\omega^\bullet$ 是
$\mathcal{O}_{U,u}$ 的正规化对偶复形，则
$\omega^\bullet\otimes_{\mathcal{O}_{U,u}}
\mathcal{O}_{X,\overline{x}}$ 是 $\mathcal{O}_{X,\overline{x}}$
的正规化对偶复形。因此，《概形的对偶性》引理
\ref{duality-lemma-dimension-function-scheme} 对概形 $U$ 与复形
$\omega_U^\bullet$ 给出的维数函数 $U\to\mathbf{Z}$，等于
$U\to|X|$ 与 $\delta$ 的复合。利用 $|X|$ 中的特化可提升为
$U$ 中的特化，且 $U$ 中的非平凡特化映到 $X$ 中的非平凡特化
（《适度空间》引理 \ref{decent-spaces-lemma-decent-specialization} 及
\ref{decent-spaces-lemma-decent-no-specializations-map-to-same-point}），
一个简单的拓扑论证表明 $\delta$ 是 $|X|$ 上的维数函数。
\end{proof}






\section{推前的右伴随}
\label{spaces-duality-section-twisted-inverse-image}

\noindent
这是《概形的对偶性》第 \ref{duality-section-twisted-inverse-image} 节
的类似版本。

\begin{lemma}
\label{spaces-duality-lemma-twisted-inverse-image}
\begin{reference}
这几乎就是 \cite[Example 4.2]{Neeman-Grothendieck}。
\end{reference}
设 $S$ 是概形，$f:X\to Y$ 是 $S$ 上拟分离且拟紧的代数空间之间
的态射。函子 $Rf_*:D_\QCoh(X)\to D_\QCoh(Y)$ 有右伴随。
\end{lemma}

\begin{proof}
我们验证《导出范畴》命题 \ref{derived-proposition-brown} 的假设，
由此证明右伴随存在。首先，范畴 $D_\QCoh(\mathcal{O}_X)$ 有直和，
见《空间的导出范畴》引理
\ref{spaces-perfect-lemma-quasi-coherence-direct-sums}；范畴
$D_\QCoh(\mathcal{O}_X)$ 还是紧生成的，
见《空间的导出范畴》定理
\ref{spaces-perfect-theorem-bondal-van-den-Bergh}。
由于 $X$ 与 $Y$ 拟紧且拟分离，$f$ 也如此；见《空间的态射》引理
\ref{spaces-morphisms-lemma-compose-after-separated} 及
\ref{spaces-morphisms-lemma-quasi-compact-permanence}。
因而函子 $Rf_*$ 与直和交换，见《空间的导出范畴》引理
\ref{spaces-perfect-lemma-quasi-coherence-pushforward-direct-sums}。
证明完毕。
\end{proof}

\begin{lemma}
\label{spaces-duality-lemma-twisted-inverse-image-bounded-below}
记号与假设同引理 \ref{spaces-duality-lemma-twisted-inverse-image}。设
$a:D_\QCoh(\mathcal{O}_Y)\to D_\QCoh(\mathcal{O}_X)$ 是 $Rf_*$
的右伴随。则 $a$ 把 $D^+_\QCoh(\mathcal{O}_Y)$ 映入
$D^+_\QCoh(\mathcal{O}_X)$。事实上，存在整数 $N$，使得
$i\leq c$ 时 $H^i(K)=0$ 蕴含 $i\leq c-N$ 时 $H^i(a(K))=0$。
\end{lemma}

\begin{proof}
由《空间的导出范畴》引理
\ref{spaces-perfect-lemma-quasi-coherence-direct-image}，函子 $Rf_*$
具有有限上同调维数。换言之，存在整数 $N$，使得若 $i\geq c$ 时
$H^i(L)=0$，则 $i\geq N+c$ 时 $H^i(Rf_*L)=0$。
设 $K\in D^+_\QCoh(\mathcal{O}_Y)$ 满足 $i\leq c$ 时
$H^i(K)=0$。则
$$
\Hom_{D(\mathcal{O}_X)}(\tau_{\leq c - N}a(K), a(K)) =
\Hom_{D(\mathcal{O}_Y)}(Rf_*\tau_{\leq c - N}a(K), K) = 0
$$
，其中最后一个等号由上述结论得到。显然，这蕴含在 $i\leq c-N$ 时
$H^i(a(K))=0$。
\end{proof}

\noindent
设 $S$ 是概形，$f:X\to Y$ 是 $S$ 上拟分离且拟紧的代数空间之间
的态射。以 $a$ 表示
$Rf_*:D_\QCoh(\mathcal{O}_X)\to D_\QCoh(\mathcal{O}_Y)$ 的右伴随。
对每个 $K\in D_\QCoh(\mathcal{O}_Y)$ 与
$L\in D_\QCoh(\mathcal{O}_X)$，得到典范映射
\begin{equation}
\label{spaces-duality-equation-sheafy-trace}
Rf_*R\SheafHom_{\mathcal{O}_X}(L, a(K))
\longrightarrow
R\SheafHom_{\mathcal{O}_Y}(Rf_*L, K)
\end{equation}
具体而言，该映射构造为复合
$$
Rf_*R\SheafHom_{\mathcal{O}_X}(L, a(K)) \to
R\SheafHom_{\mathcal{O}_Y}(Rf_*L, Rf_*a(K)) \to
R\SheafHom_{\mathcal{O}_Y}(Rf_*L, K)
$$
，其中第一支箭头来自《站点上的上同调》评注
\ref{sites-cohomology-remark-projection-formula-for-internal-hom}，
第二支箭头是该伴随的余单位 $Rf_*a(K)\to K$。

\begin{lemma}
\label{spaces-duality-lemma-iso-on-RSheafHom}
设 $S$ 是概形，$f:X\to Y$ 是 $S$ 上拟紧且拟分离的代数空间之间
的态射，$a$ 是
$Rf_*:D_\QCoh(\mathcal{O}_X)\to D_\QCoh(\mathcal{O}_Y)$ 的右伴随。
设 $L\in D_\QCoh(\mathcal{O}_X)$ 且
$K\in D_\QCoh(\mathcal{O}_Y)$。则映射
(\ref{spaces-duality-equation-sheafy-trace})
$$
Rf_*R\SheafHom_{\mathcal{O}_X}(L, a(K))
\longrightarrow
R\SheafHom_{\mathcal{O}_Y}(Rf_*L, K)
$$
在施加《空间的导出范畴》第
\ref{spaces-perfect-section-better-coherator} 节讨论的函子
$DQ_Y:D(\mathcal{O}_Y)\to D_\QCoh(\mathcal{O}_Y)$ 后成为同构。
\end{lemma}

\begin{proof}
由《空间的导出范畴》引理
\ref{spaces-perfect-lemma-better-coherator}，$DQ_Y$ 存在，故该陈述有意义。
由于 $DQ_Y$ 是包含函子
$D_\QCoh(\mathcal{O}_Y)\to D(\mathcal{O}_Y)$ 的右伴随，为证明引理，
只需证明对任意 $M\in D_\QCoh(\mathcal{O}_Y)$，映射
(\ref{spaces-duality-equation-sheafy-trace}) 诱导双射
% P10-E0252
$$
\Hom_Y(M, Rf_*R\SheafHom_{\mathcal{O}_X}(L, a(K)))
\longrightarrow
\Hom_Y(M, R\SheafHom_{\mathcal{O}_Y}(Rf_*L, K))
$$
。为此，使用下列一串等式：
\begin{align*}
\Hom_Y(M, Rf_*R\SheafHom_{\mathcal{O}_X}(L, a(K)))
& =
\Hom_X(Lf^*M, R\SheafHom_{\mathcal{O}_X}(L, a(K))) \\
& =
\Hom_X(Lf^*M \otimes_{\mathcal{O}_X}^\mathbf{L} L, a(K)) \\
& =
\Hom_Y(Rf_*(Lf^*M \otimes_{\mathcal{O}_X}^\mathbf{L} L), K) \\
& =
\Hom_Y(M \otimes_{\mathcal{O}_Y}^\mathbf{L} Rf_*L, K) \\
& =
\Hom_Y(M, R\SheafHom_{\mathcal{O}_Y}(Rf_*L, K))
\end{align*}
第一个等号由《站点上的上同调》引理
\ref{sites-cohomology-lemma-adjoint} 得到。
% P10-E0253
第二个等号由《站点上的上同调》引理
\ref{sites-cohomology-lemma-internal-hom} 得到；第三个等号来自 $a$ 的构造；
第四个等号由《空间的导出范畴》引理
\ref{spaces-perfect-lemma-cohomology-base-change} 得到（这是关键一步）；
第五个等号由《站点上的上同调》引理
\ref{sites-cohomology-lemma-internal-hom} 得到。
\end{proof}

\begin{example}
\label{spaces-duality-example-iso-on-RSheafHom}
若不施加“拟凝聚化子”$DQ_Y$，引理 \ref{spaces-duality-lemma-iso-on-RSheafHom}
的陈述并不成立。见《概形的对偶性》例
\ref{duality-example-iso-on-RSheafHom}。
\end{example}

\begin{remark}
\label{spaces-duality-remark-iso-on-RSheafHom}
在引理 \ref{spaces-duality-lemma-iso-on-RSheafHom} 的情形中，由《空间的导出范畴》
引理 \ref{spaces-perfect-lemma-pushforward-better-coherator} 有
$$
DQ_Y(Rf_*R\SheafHom_{\mathcal{O}_X}(L, a(K))) =
Rf_* DQ_X(R\SheafHom_{\mathcal{O}_X}(L, a(K)))
$$
。因此，若
$R\SheafHom_{\mathcal{O}_X}(L,a(K))\in D_\QCoh(\mathcal{O}_X)$，
便可从箭头左端“消去”$DQ_Y$。另一方面，若已知
$R\SheafHom_{\mathcal{O}_Y}(Rf_*L,K)\in D_\QCoh(\mathcal{O}_Y)$，
便可从箭头右端“消去”$DQ_Y$。若两者均成立，则
(\ref{spaces-duality-equation-sheafy-trace}) 是同构。结合《空间的导出范畴》引理
\ref{spaces-perfect-lemma-quasi-coherence-internal-hom} 可知，在下列任一
情形中，$Rf_*R\SheafHom_{\mathcal{O}_X}(L,a(K))\to
R\SheafHom_{\mathcal{O}_Y}(Rf_*L,K)$ 是同构：
\begin{enumerate}
\item $L$ 与 $Rf_*L$ 都是完美的；或
\item $K$ 下有界，且 $L$ 与 $Rf_*L$ 都是伪凝聚的。
\end{enumerate}
对于 (2)，使用如下事实：若 $K$ 下有界，则 $a(K)$ 下有界；见引理
\ref{spaces-duality-lemma-twisted-inverse-image-bounded-below}。
\end{remark}

\begin{example}
\label{spaces-duality-example-iso-on-RSheafHom-noetherian}
设 $S$ 是概形，$f:X\to Y$ 是 $S$ 上 Noether 代数空间之间的真态射，
$L\in D^-_{\textit{Coh}}(X)$，且
$K\in D^+_{\QCoh}(\mathcal{O}_Y)$。则映射
$Rf_*R\SheafHom_{\mathcal{O}_X}(L,a(K))\to
R\SheafHom_{\mathcal{O}_Y}(Rf_*L,K)$ 是同构。事实上，由《空间的
导出范畴》引理
\ref{spaces-perfect-lemma-identify-pseudo-coherent-noetherian} 及
\ref{spaces-perfect-lemma-direct-image-coherent}，复形 $L$ 与 $Rf_*L$
都是伪凝聚的，故可应用评注 \ref{spaces-duality-remark-iso-on-RSheafHom} 中的讨论。
\end{example}

\begin{lemma}
\label{spaces-duality-lemma-iso-global-hom}
设 $S$ 是概形，$f:X\to Y$ 是 $S$ 上拟分离且拟紧的代数空间之间
的态射。对所有 $L\in D_\QCoh(\mathcal{O}_X)$ 与
$K\in D_\QCoh(\mathcal{O}_Y)$，(\ref{spaces-duality-equation-sheafy-trace}) 诱导
整体导出 Hom 的同构
$R\Hom_X(L,a(K))\to R\Hom_Y(Rf_*L,K)$。
\end{lemma}

\begin{proof}
由构造（《站点上的上同调》第
\ref{sites-cohomology-section-global-RHom} 节），复形
$$
R\Hom_X(L, a(K)) =
R\Gamma(X, R\SheafHom_{\mathcal{O}_X}(L, a(K))) =
R\Gamma(Y, Rf_*R\SheafHom_{\mathcal{O}_X}(L, a(K)))
$$
与
% P10-E0254
$$
R\Hom_Y(Rf_*L, K) = R\Gamma(Y, R\SheafHom_{\mathcal{O}_X}(Rf_*L, a(K)))
$$
。因此，本引理由引理 \ref{spaces-duality-lemma-iso-on-RSheafHom} 推出。
具体而言，若 $D(\mathcal{O}_Y)$ 中的映射 $E\to E'$ 诱导同构
$DQ_Y(E)\to DQ_Y(E')$，则它诱导拟同构
$R\Gamma(Y,E)\to R\Gamma(Y,E')$。事实上，
$H^i(Y, E) = \Ext^i_Y(\mathcal{O}_Y, E) = \Hom(\mathcal{O}_Y[-i], E) =
\Hom(\mathcal{O}_Y[-i], DQ_Y(E))$。这是因为 $\mathcal{O}_Y[-i]$
属于 $D_\QCoh(\mathcal{O}_Y)$，且 $DQ_Y$
是包含函子 $D_\QCoh(\mathcal{O}_Y)\to D(\mathcal{O}_Y)$ 的右伴随。
\end{proof}









\section{推前的右伴随与基变换（一）}
\label{spaces-duality-section-base-change-map}

\noindent
我们来定义推前函子的右伴随之间的基变换映射。设 $S$ 是概形，
考虑笛卡尔图
\begin{equation}
\label{spaces-duality-equation-base-change}
\vcenter{
\xymatrix{
X' \ar[r]_{g'} \ar[d]_{f'} & X \ar[d]^f \\
Y' \ar[r]^g & Y
}
}
\end{equation}
，其中 $Y'$ 与 $X$ 在 $Y$ 上 {\bf Tor 独立}。以
$$
a  : D_\QCoh(\mathcal{O}_Y) \to D_\QCoh(\mathcal{O}_X)
\quad\text{且}\quad
a' : D_\QCoh(\mathcal{O}_{Y'}) \to D_\QCoh(\mathcal{O}_{X'})
$$
表示 $Rf_*$ 与 $Rf'_*$ 的右伴随（引理
\ref{spaces-duality-lemma-twisted-inverse-image}）。《站点上的上同调》评注
\ref{sites-cohomology-remark-base-change} 的基变换映射给出函子变换
$$
Lg^* \circ Rf_* \longrightarrow Rf'_* \circ L(g')^*
$$
于具有拟凝聚上同调的层的导出范畴上。因此，得到反向的右伴随之间
的变换
$$
a \circ Rg_* \longleftarrow Rg'_* \circ a'
$$

\begin{lemma}
\label{spaces-duality-lemma-flat-precompose-pus}
在图 (\ref{spaces-duality-equation-base-change}) 中，映射
$a\circ Rg_*\leftarrow Rg'_*\circ a'$ 是同构。
\end{lemma}

\begin{proof}
由《空间的导出范畴》引理
\ref{spaces-perfect-lemma-compare-base-change}，对
$D_\QCoh(\mathcal{O}_X)$ 中每个 $K$，基变换映射
$Lg^*\circ Rf_*K\to Rf'_*\circ L(g')^*K$ 都是同构（这里使用了
Tor 独立的假设）。因此，伴随函子之间相应的变换也是同构。
\end{proof}

\noindent
于是可以考虑由下列复合给出的函子
$D_\QCoh(\mathcal{O}_Y)\to D_\QCoh(\mathcal{O}_{X'})$ 的态射：
\begin{equation}
\label{spaces-duality-equation-base-change-map}
L(g')^* \circ a \to
L(g')^* \circ a \circ Rg_* \circ Lg^* \leftarrow
L(g')^* \circ Rg'_* \circ a' \circ Lg^* \to a' \circ Lg^*
\end{equation}
第一支箭头来自伴随映射 $\text{id}\to Rg_*Lg^*$，最后一支箭头来自
伴随映射 $L(g')^*Rg'_*\to\text{id}$。为逆转中间的箭头，需要 Tor
独立的假设，见引理 \ref{spaces-duality-lemma-flat-precompose-pus}。或者，利用
$L(g')^*$ 与 $R(g')_*$ 的伴随性，可以把
(\ref{spaces-duality-equation-base-change-map}) 看作自然变换
$$
a \to a \circ Rg_* \circ Lg^* \leftarrow Rg'_* \circ a' \circ Lg^*
$$
% P10-E0255
，其中第二支箭头仍然可逆。若 $M\in D_\QCoh(\mathcal{O}_X)$ 且
$K\in D_\QCoh(\mathcal{O}_Y)$，则该映射在 Yoneda 函子上由下式给出：
\begin{align*}
\Hom_X(M, a(K))
& =
\Hom_Y(Rf_*M, K) \\
& \to
\Hom_Y(Rf_*M, Rg_* Lg^*K) \\
& =
\Hom_{Y'}(Lg^*Rf_*M, Lg^*K) \\
& \leftarrow
\Hom_{Y'}(Rf'_* L(g')^*M, Lg^*K) \\
& =
\Hom_{X'}(L(g')^*M, a'(Lg^*K)) \\
& =
\Hom_X(M, Rg'_*a'(Lg^*K))
\end{align*}
% P10-E0256
（其中向左的箭头由《空间的导出范畴》引理
\ref{spaces-perfect-lemma-compare-base-change}
给出的基变换定理可知是可逆的），这使构造更为明确。

\medskip\noindent
本节先证明，该基变换映射对于方块的纵横叠合满足若干自然相容性；
这类似于通常基变换映射的《站点上的上同调》评注
\ref{sites-cohomology-remark-compose-base-change} 与
\ref{sites-cohomology-remark-compose-base-change-horizontal}。
建议读者初读时跳过本节余下部分。

\begin{lemma}
\label{spaces-duality-lemma-compose-base-change-maps}
设 $S$ 是概形。考虑交换图
$$
\xymatrix{
X' \ar[r]_k \ar[d]_{f'} & X \ar[d]^f \\
Y' \ar[r]^l \ar[d]_{g'} & Y \ar[d]^g \\
Z' \ar[r]^m & Z
}
$$
，其中各对象均为 $S$ 上拟紧且拟分离的代数空间，两个方块均为
笛卡尔的，并且 $f$ 与 $l$、$g$ 与 $m$ 分别 Tor 独立。则两个方块的
映射 (\ref{spaces-duality-equation-base-change-map}) 复合后给出外侧矩形的基变换映射
（精确陈述见证明）。
\end{lemma}

\begin{proof}
由假设可知 $g\circ f$ 与 $m$ Tor 独立（略去细节），故该陈述有意义。
本证明以 $k^*$ 代替 $Lk^*$，以 $f_*$ 代替 $Rf_*$。令 $a,b,c$
分别为引理 \ref{spaces-duality-lemma-twisted-inverse-image} 对 $f,g,g\circ f$ 给出的
右伴随；带撇版本的记号类似。上方方块对应的箭头是复合
$$
\gamma_{top} :
k^* \circ a \to k^* \circ a \circ l_* \circ l^*
\xleftarrow{\xi_{top}} k^* \circ k_* \circ a' \circ l^* \to a' \circ l^*
$$
，其中 $\xi_{top}:k_*\circ a'\to a\circ l_*$ 是同构（因而可逆），
并且是基变换映射 $l^*\circ f_*\to f'_*\circ k^*$ 的“对偶”箭头。
两端的箭头来自典范映射 $1\to l_*\circ l^*$ 与
$k^*\circ k_*\to1$。类似地，对下方方块有
$$
\gamma_{bot} :
l^* \circ b \to l^* \circ b \circ m_* \circ m^*
\xleftarrow{\xi_{bot}} l^* \circ l_* \circ b' \circ m^* \to b' \circ m^*
$$
对于外侧矩形，得到
$$
\gamma_{rect} :
k^* \circ c \to k^* \circ c \circ m_* \circ m^*
\xleftarrow{\xi_{rect}} k^* \circ k_* \circ c' \circ m^* \to c' \circ m^*
$$
由于 $(g\circ f)_*=g_*\circ f_*$，故 $c=a\circ b$；类似地，
$c'=a'\circ b'$。引理的陈述是，$\gamma_{rect}$ 等于复合
$$
k^* \circ c = k^* \circ a \circ b \xrightarrow{\gamma_{top}}
a' \circ l^* \circ b \xrightarrow{\gamma_{bot}}
a' \circ b' \circ m^* = c' \circ m^*
$$
为证明这一点，考察下图：
$$
\xymatrix{
& & k^* \circ a \circ b \ar[d] \ar[lldd] \\
& & k^* \circ a \circ l_* \circ l^* \circ b \ar[ld] \\
k^* \circ a \circ b \circ m_* \circ m^* \ar[r] &
k^* \circ a \circ l_* \circ l^* \circ b \circ m_* \circ m^* &
k^* \circ k_* \circ a' \circ l^* \circ b \ar[u]_{\xi_{top}} \ar[d] \ar[ld] \\
& k^*\circ k_* \circ a' \circ l^* \circ b \circ m_* \circ m^*
\ar[u]_{\xi_{top}} \ar[rd] &
a' \circ l^* \circ b \ar[d] \\
k^* \circ k_* \circ a' \circ b' \circ m^* \ar[uu]_{\xi_{rect}} \ar[ddrr] &
k^*\circ k_* \circ a' \circ l^* \circ l_* \circ b' \circ m^*
\ar[u]_{\xi_{bot}} \ar[l] \ar[dr] &
a' \circ l^* \circ b \circ m_* \circ m^* \\
& & a' \circ l^* \circ l_* \circ b' \circ m^* \ar[u]_{\xi_{bot}} \ar[d] \\
& & a' \circ b' \circ m^*
}
$$
沿右侧向下得到上述复合，沿左侧向下得到 $\gamma_{rect}$。
由《范畴》引理 \ref{categories-lemma-properties-2-cat-cats}，或者更直接地
由《范畴》定义 \ref{categories-definition-horizontal-composition} 之前的
讨论，此图右侧的所有四边形均交换。因此，只需证明下图
$$
\xymatrix{
a \circ l_* \circ l^* \circ b \circ m_* &
a \circ b \circ m_* \ar[l] \\
k_* \circ a' \circ l^* \circ b \circ m_* \ar[u]_{\xi_{top}} & \\
k_* \circ a' \circ l^* \circ l_* \circ b' \ar[u]_{\xi_{bot}} \ar[r] &
k_* \circ a' \circ b' \ar[uu]_{\xi_{rect}}
}
$$
在逆转箭头 $\xi_{top},\xi_{bot},\xi_{rect}$ 后交换（注意，这不同于
要求原图本身交换）。另一方面，下图
$$
\xymatrix{
& a \circ l_* \circ l^* \circ b \circ m_* \\
a \circ l_* \circ l^* \circ l_* \circ b'
\ar[ru]^{\xi_{bot}} & &
k_* \circ a' \circ l^* \circ b \circ m_* \ar[ul]_{\xi_{top}} \\
& k_* \circ a' \circ l^* \circ l_* \circ b'
\ar[ul]^{\xi_{top}} \ar[ur]_{\xi_{bot}}
}
$$
由《范畴》引理 \ref{categories-lemma-properties-2-cat-cats} 可知交换。
又因为下列两图
% P10-E0257
$$
\vcenter{
\xymatrix{
a \circ l_* \circ l^* \circ b \circ m_* &
a \circ b \circ m \ar[l] \\
a \circ l_* \circ l^* \circ l_* \circ b' \ar[u] &
a \circ l_* \circ b' \ar[l] \ar[u]
}
}
\quad\text{且}\quad
\vcenter{
\xymatrix{
a \circ l_* \circ l^* \circ l_* \circ b' \ar[r] & a \circ l_* \circ b' \\
k_* \circ a' \circ l^* \circ l_* \circ b' \ar[u] \ar[r] &
k_* \circ a' \circ b' \ar[u]
}
}
$$
交换（见所引文献），并且复合
$l_*\to l_*\circ l^*\circ l_*\to l_*$ 是恒等映射，所以只需证明
% P10-E0258
$$
k \circ a' \circ b' \xrightarrow{\xi_{bot}} a \circ l_* \circ b
\xrightarrow{\xi_{top}} a \circ b \circ m_*
$$
等于 $\xi_{rect}$（通过恒等 $a\circ b=c$ 与 $a'\circ b'=c'$）。
这正是《站点上的上同调》评注
\ref{sites-cohomology-remark-compose-base-change} 的对偶陈述，证明完毕。
\end{proof}

\begin{lemma}
\label{spaces-duality-lemma-compose-base-change-maps-horizontal}
设 $S$ 是概形。考虑交换图
$$
\xymatrix{
X'' \ar[r]_{g'} \ar[d]_{f''} & X' \ar[r]_g \ar[d]_{f'} & X \ar[d]^f \\
Y'' \ar[r]^{h'} & Y' \ar[r]^h & Y
}
$$
，其中各对象均为 $S$ 上拟紧且拟分离的代数空间，两个方块均为
笛卡尔的，并且 $f$ 与 $h$、$f'$ 与 $h'$ 分别 Tor 独立。则两个方块的
映射 (\ref{spaces-duality-equation-base-change-map}) 复合后给出外侧矩形的基变换映射
（精确陈述见证明）。
\end{lemma}

\begin{proof}
由假设可知 $f$ 与 $h\circ h'$ Tor 独立（略去细节），故该陈述有意义。
本证明以 $g^*$ 代替 $Lg^*$，以 $f_*$ 代替 $Rf_*$。令 $a,a',a''$
分别为引理 \ref{spaces-duality-lemma-twisted-inverse-image} 对 $f,f',f''$ 给出的
右伴随。右侧方块对应的箭头是复合
$$
\gamma_{right} :
g^* \circ a \to g^* \circ a \circ h_* \circ h^*
\xleftarrow{\xi_{right}} g^* \circ g_* \circ a' \circ h^* \to a' \circ h^*
$$
，其中 $\xi_{right}:g_*\circ a'\to a\circ h_*$ 是同构（因而可逆），
并且是基变换映射 $h^*\circ f_*\to f'_*\circ g^*$ 的“对偶”箭头。
两端的箭头来自典范映射 $1\to h_*\circ h^*$ 与
$g^*\circ g_*\to1$。类似地，对左侧方块有
$$
\gamma_{left} :
(g')^* \circ a' \to (g')^* \circ a' \circ (h')_* \circ (h')^*
\xleftarrow{\xi_{left}}
(g')^* \circ (g')_* \circ a'' \circ (h')^* \to a'' \circ (h')^*
$$
对于外侧矩形，得到
$$
\gamma_{rect} :
k^* \circ a \to
k^* \circ a \circ m_* \circ m^* \xleftarrow{\xi_{rect}}
k^* \circ k_* \circ a'' \circ m^* \to
a'' \circ m^*
$$
，其中 $k=g\circ g'$ 且 $m=h\circ h'$。有
$k^*=(g')^*\circ g^*$ 与 $m^*=(h')^*\circ h^*$。
引理的陈述是，$\gamma_{rect}$ 等于复合
$$
k^* \circ a =
(g')^* \circ g^* \circ a \xrightarrow{\gamma_{right}}
(g')^* \circ a' \circ h^* \xrightarrow{\gamma_{left}}
a'' \circ (h')^* \circ h^* = a'' \circ m^*
$$
为证明这一点，考察下图：
$$
\xymatrix{
& (g')^* \circ g^* \circ a \ar[d] \ar[ddl] \\
& (g')^* \circ g^* \circ a \circ h_* \circ h^* \ar[ld] \\
(g')^* \circ g^* \circ a \circ h_* \circ (h')_* \circ (h')^* \circ h^* &
(g')^* \circ g^* \circ g_* \circ a' \circ h^*
\ar[u]_{\xi_{right}} \ar[d] \ar[ld] \\
(g')^* \circ g^* \circ g_* \circ a' \circ (h')_* \circ (h')^* \circ h^*
\ar[u]_{\xi_{right}} \ar[dr] &
(g')^* \circ a' \circ h^* \ar[d] \\
(g')^* \circ g^* \circ g_* \circ (g')_* \circ a'' \circ (h')^* \circ h^*
\ar[u]_{\xi_{left}} \ar[ddr] \ar[dr] &
(g')^* \circ a' \circ (h')_* \circ (h')^* \circ h^* \\
& (g')^*\circ (g')_* \circ a'' \circ (h')^* \circ h^*
\ar[u]_{\xi_{left}} \ar[d] \\
& a'' \circ (h')^* \circ h^*
}
$$
沿右侧向下得到上述复合，沿左侧向下得到 $\gamma_{rect}$。
由《范畴》引理 \ref{categories-lemma-properties-2-cat-cats}，或者更直接地
由《范畴》定义 \ref{categories-definition-horizontal-composition} 之前的
讨论，此图右侧的所有四边形均交换。因此，只需证明
$$
g_* \circ (g')_* \circ a'' \xrightarrow{\xi_{left}}
g_* \circ a' \circ (h')_* \xrightarrow{\xi_{right}}
a \circ h_* \circ (h')_*
$$
等于 $\xi_{rect}$。这正是《上同调》评注
\ref{cohomology-remark-compose-base-change-horizontal} 的对偶陈述，
证明完毕。
\end{proof}

\begin{remark}
\label{spaces-duality-remark-going-around}
设 $S$ 是概形。考虑交换图
$$
\xymatrix{
X'' \ar[r]_{k'} \ar[d]_{f''} & X' \ar[r]_k \ar[d]_{f'} & X \ar[d]^f \\
Y'' \ar[r]^{l'} \ar[d]_{g''} & Y' \ar[r]^l \ar[d]_{g'} & Y \ar[d]^g \\
Z'' \ar[r]^{m'} & Z' \ar[r]^m & Z
}
$$
，其中各对象均为 $S$ 上拟紧且拟分离的代数空间，所有方块均为
笛卡尔的，并且 $(f,l),(g,m),(f',l'),(g',m')$ 都是 Tor 独立的
态射对。令 $a,a',a'',b,b',b''$ 分别为引理
\ref{spaces-duality-lemma-twisted-inverse-image} 对 $f,f',f'',g,g',g''$ 给出的右伴随。
如下把图中的方块标记为 $A,B,C,D$：
$$
\begin{matrix}
A & B \\
C & D
\end{matrix}
$$
于是，各方块的映射 (\ref{spaces-duality-equation-base-change-map}) 如下（这里使用
$k^*=Lk^*$ 等简写）：
$$
\begin{matrix}
\gamma_A : (k')^* \circ a' \to a'' \circ (l')^* &
\gamma_B : k^* \circ a \to a' \circ l^* \\
\gamma_C : (l')^* \circ b' \to b'' \circ (m')^* &
\gamma_D : l^* \circ b \to b' \circ m^*
\end{matrix}
$$
对于 $2\times1$ 与 $1\times2$ 矩形，另有四个基变换映射：
% P10-E0259
$$
\begin{matrix}
\gamma_{A + B} : (k \circ k')^* \circ a \to a'' \circ (l \circ l')^* \\
\gamma_{C + D} : (l \circ l')^* \circ b \to b'' \circ (m \circ m')^* \\
\gamma_{A + C} : (k')^* \circ (a' \circ b') \to (a'' \circ b'') \circ (m')^* \\
\gamma_{A + C} : k^* \circ (a \circ b) \to (a' \circ b') \circ m^*
\end{matrix}
$$
由引理 \ref{spaces-duality-lemma-compose-base-change-maps-horizontal}，有
$$
\gamma_{A + B} = \gamma_A \circ \gamma_B, \quad
\gamma_{C + D} = \gamma_C \circ \gamma_D
$$
；由引理 \ref{spaces-duality-lemma-compose-base-change-maps}，有
$$
\gamma_{A + C} = \gamma_C \circ \gamma_A, \quad
\gamma_{B + D} = \gamma_D \circ \gamma_B
$$
这里，更准确的写法应为
$\gamma_{A + B} = (\gamma_A \star \text{id}_{l^*}) \circ
(\text{id}_{(k')^*} \star \gamma_B)$，记号同《范畴》第
\ref{categories-section-formal-cat-cat} 节；其他情形类似。不过，我们
继续沿用引理 \ref{spaces-duality-lemma-compose-base-change-maps} 与
\ref{spaces-duality-lemma-compose-base-change-maps-horizontal} 的证明中的记号滥用，
省略与恒等变换的 $\star$ 乘积，因为只要已知变换的源与目标，
% P10-E0260
便可判断应补入哪些恒等变换。综上，先验地得到两个变换
$$
(k')^* \circ k^* \circ a \circ b
\longrightarrow
a'' \circ b'' \circ (m')^* \circ m^*
$$
，即
$$
\gamma_C \circ \gamma_A \circ \gamma_D \circ \gamma_B =
\gamma_{A + C} \circ \gamma_{B + D}
$$
以及
$$
\gamma_C \circ \gamma_D \circ \gamma_A \circ \gamma_B =
\gamma_{C + D} \circ \gamma_{A + B}
$$
本评注的要点是说明这两个变换相等。具体而言，只需证明下图交换：
$$
\xymatrix{
(k')^* \circ a' \circ l^* \circ b \ar[r]_{\gamma_D} \ar[d]_{\gamma_A} &
(k')^* \circ a' \circ b' \circ m^* \ar[d]^{\gamma_A} \\
a'' \circ (l')^* \circ l^* \circ b \ar[r]^{\gamma_D} &
a'' \circ (l')^* \circ b' \circ m^*
}
$$
由《范畴》引理 \ref{categories-lemma-properties-2-cat-cats}，或者更直接地
由《范畴》定义 \ref{categories-definition-horizontal-composition} 之前的
讨论，这确实成立。
\end{remark}








\section{推前的右伴随与基变换（二）}
\label{spaces-duality-section-base-change-II}

\noindent
本节证明，第 \ref{spaces-duality-section-base-change-map} 节的基变换映射在某些情形下
是同构。

\begin{lemma}
\label{spaces-duality-lemma-more-base-change}
在图 (\ref{spaces-duality-equation-base-change}) 中，进一步假设
$g:Y'\to Y$ 是仿射概形之间的态射，且 $f:X\to Y$ 是真的。
则基变换映射 (\ref{spaces-duality-equation-base-change-map}) 在下列情形中诱导同构
$$
L(g')^*a(K) \longrightarrow a'(Lg^*K)
$$
\begin{enumerate}
\item 若 $f$ 平坦且有限表示，则对所有
% P10-E0261
$K\in D_\QCoh(\mathcal{O}_X)$ 均成立；
\item 若 $f$ 完美且 $Y$ 是 Noether 的，则对所有
$K\in D_\QCoh(\mathcal{O}_X)$ 均成立；
\item 若 $g$ 具有有限 Tor 维数且 $Y$ 是 Noether 的，则对
$K\in D_\QCoh^+(\mathcal{O}_X)$ 成立。
\end{enumerate}
\end{lemma}

\begin{proof}
写成 $Y=\Spec(A)$ 与 $Y'=\Spec(A')$。作为仿射态射的基变换，$g'$
是仿射的。取 $D_\QCoh(\mathcal{O}_X)$ 的完美生成元 $M$，见
《空间的导出范畴》定理
\ref{spaces-perfect-theorem-bondal-van-den-Bergh}。则 $L(g')^*M$ 是
$D_\QCoh(\mathcal{O}_{X'})$ 的生成元，见《空间的导出范畴》评注
\ref{spaces-perfect-remark-pullback-generator}。因此，只需证明
(\ref{spaces-duality-equation-base-change-map}) 诱导整体 Hom 复形的同构
\begin{equation}
\label{spaces-duality-equation-iso}
R\Hom_{X'}(L(g')^*M, L(g')^*a(K))
\longrightarrow
R\Hom_{X'}(L(g')^*M, a'(Lg^*K))
\end{equation}
；见《站点上的上同调》第
\ref{sites-cohomology-section-global-RHom} 节，因为这将蕴含
$L(g')^*a(K)\to a'(Lg^*K)$ 的锥为零。证明的结构如下：先证明这些
Hom 复形同构，再在证明的最后部分说明该同构由
(\ref{spaces-duality-equation-iso}) 诱导。

\medskip\noindent
先看左端。由于 $M$ 完美，典范映射
$$
R\Hom_X(M, a(K)) \otimes^\mathbf{L}_A A'
\longrightarrow
R\Hom_{X'}(L(g')^*M, L(g')^*a(K))
$$
由《空间的导出范畴》引理
\ref{spaces-perfect-lemma-affine-morphism-and-hom-out-of-perfect}
可知是同构。再结合引理 \ref{spaces-duality-lemma-iso-global-hom} 的同构
$R\Hom_Y(Rf_*M, K) = R\Hom_X(M, a(K))$
，可知左端等于 $R\Hom_Y(Rf_*M,K)\otimes_A^\mathbf{L}A'$。

\medskip\noindent
再看右端。先使用引理 \ref{spaces-duality-lemma-iso-global-hom} 的同构
$$
R\Hom_{X'}(L(g')^*M, a'(Lg^*K)) = R\Hom_{Y'}(Rf'_*L(g')^*M, Lg^*K)
$$
。由于 $f$ 与 $g$ Tor 独立，由《空间的导出范畴》引理
\ref{spaces-perfect-lemma-compare-base-change}，基变换映射
$Lg^*Rf_*M\to Rf'_*L(g')^*M$ 是同构。因此可将右端改写为
$R\Hom_{Y'}(Lg^*Rf_*M,Lg^*K)$。由于 $Y,Y'$ 仿射，且 $K,Rf_*M$
属于 $D_\QCoh(\mathcal{O}_Y)$（《空间的导出范畴》引理
\ref{spaces-perfect-lemma-quasi-coherence-direct-image}），有典范映射
$$
\beta :
R\Hom_Y(Rf_*M, K) \otimes^\mathbf{L}_A A'
\longrightarrow
R\Hom_{Y'}(Lg^*Rf_*M, Lg^*K)
$$
于 $D(A')$ 中。这就是《代数详论》等式
(\ref{more-algebra-equation-base-change-RHom}) 的箭头；这里使用了
《概形的导出范畴》引理 \ref{perfect-lemma-affine-compare-bounded} 与
\ref{perfect-lemma-quasi-coherence-internal-hom} 在代数表述之间来回转换。
\begin{enumerate}
\item 若 $f$ 平坦且有限表示，则由《空间的导出范畴》引理
\ref{spaces-perfect-lemma-flat-proper-perfect-direct-image-general}，复形
$Rf_*M$ 在 $Y$ 上完美；由《代数详论》引理
\ref{more-algebra-lemma-base-change-RHom} 的 (1)，$\beta$ 是同构。
\item 若 $f$ 完美且 $Y$ 是 Noether 的，则由《空间的态射详论》引理
\ref{spaces-more-morphisms-lemma-perfect-proper-perfect-direct-image}，
复形 $Rf_*M$ 在 $Y$ 上完美；同上，$\beta$ 是同构。
\item 若 $g$ 具有有限 Tor 维数且 $Y$ 是 Noether 的，则复形 $Rf_*M$
在 $Y$ 上伪凝聚（《空间的导出范畴》引理
\ref{spaces-perfect-lemma-direct-image-coherent} 与
\ref{spaces-perfect-lemma-identify-pseudo-coherent-noetherian}）；由
《代数详论》引理 \ref{more-algebra-lemma-base-change-RHom} 的 (4)，
$\beta$ 是同构。
\end{enumerate}
因此，得到与上一段相同的结果。

\medskip\noindent
在证明余下部分，我们说明第二、三段给出的
(\ref{spaces-duality-equation-iso}) 左右两端的认同，事实上正由
(\ref{spaces-duality-equation-iso}) 给出。为使公式易于处理，记
$(-,-)_X=R\Hom_X(-,-)$，以 $-\otimes A'$ 代替
% P10-E0262
$-\otimes_A^\mathbf{L}A'$，并简记 $g^*=Lg^*$、$f_*=Rf_*$。
考虑下列交换图：
% P10-E0263
$$
\xymatrix{
((g')^*M, (g')^*a(K))_{X'} \ar[d] &
(M, a(K))_X \otimes A' \ar[l]^-\alpha \ar[d] &
(f_*M, K)_Y \otimes A' \ar@{=}[l] \ar[d] \\
((g')^*M, (g')^*a(g_*g^*K))_{X'} &
(M, a(g_*g^*K))_X \otimes A' \ar[l]^-\alpha &
(f_*M, g_*g^*K)_Y \otimes A' \ar@{=}[l] \ar@/_4pc/[dd]_{\mu'} \\
((g')^*M, (g')^*g'_*a'(g^*K))_{X'} \ar[u] \ar[d] &
(M, g'_*a'(g^*K))_X \otimes A' \ar[u] \ar[l]^-\alpha \ar[ld]^\mu &
(f_*M, K) \otimes A' \ar[d]^\beta \\
((g')^*M, a'(g^*K))_{X'} &
(f'_*(g')^*M, g^*K)_{Y'} \ar@{=}[l] \ar[r] &
(g^*f_*M, g^*K)_{Y'}
}
$$
标有 $\alpha$ 的箭头，是《空间的导出范畴》引理
\ref{spaces-perfect-lemma-affine-morphism-and-hom-out-of-perfect}
对以 $X',X,Y',Y$ 为顶点之图所给出的映射。由于水平箭头对各变量
具有函子性，图的上部交换。中间的竖直箭头来自引理
\ref{spaces-duality-lemma-flat-precompose-pus} 的可逆变换
$g'_*\circ a'\to a\circ g_*$，因而中间方块交换。沿左侧向下正是
(\ref{spaces-duality-equation-iso})。上方的水平箭头给出证明第二段使用的认同；
包括 $\beta$ 在内的下方水平箭头给出证明第三段使用的认同。
给定 $E\in D(A)$、$E'\in D(A')$ 以及 $D(A)$ 中的
$c:E\to E'$，以 $\mu_c:E\otimes A'\to E'$ 表示由 $c$ 以及限制与
基变换的伴随性所诱导的映射。若 $c$ 明确，便写成 $\mu=\mu_c$，
即从记号中略去 $c$。图中的映射 $\mu$ 属于这种形式，其中 $c$
由下列认同给出：
% P10-E0264
$(M, g'_*a(g^*K))_X = ((g')^*M, a'(g^*K))_{X'}$
。由《空间的导出范畴》评注
\ref{spaces-perfect-remark-multiplication-map}，含 $\mu$ 的三角形交换。

\medskip\noindent
注意，图
$$
\xymatrix{
(M, a(g_*g^*K))_X &
(f_*M, g_* g^*K)_Y \ar@{=}[l] &
(g^*f_*M, g^*K)_{Y'} \ar@{=}[l] \\
(M, g'_* a'(g^*K))_X \ar[u] &
((g')^*M, a'(g^*K))_{X'} \ar@{=}[l] &
(f'_*(g')^*M, g^*K)_{Y'} \ar@{=}[l] \ar[u]
}
$$
由变换 $g'_*\circ a'\to a\circ g_*$ 的定义即知交换。令 $\mu'$ 如上，
对应于认同
% P10-E0265
$(f_*M, g_*g^*K)_X = (g^*f_*M, g^*K)_{Y'}$
，则六边形也交换。因此，只需证明 $\beta$ 等于下列映射与 $\mu'$
的复合：
% P10-E0266
$(f_*M, K)_Y \otimes A' \to (f_*M, g_*g^*K)_X \otimes A'$
。为此，只需证明所诱导的两个映射
$(f_*M,K)_Y\to(g^*f_*M,g^*K)_{Y'}$ 相同。换言之，只需证明下图交换：
$$
\xymatrix{
R\Hom_A(E, K) \ar[rr]_{\text{由 }\beta\text{ 诱导}} \ar[rd] & &
R\Hom_{A'}(E \otimes_A^\mathbf{L} A', K \otimes_A^\mathbf{L} A') \\
& R\Hom_A(E, K \otimes_A^\mathbf{L} A') \ar[ru]
}
$$
对所有 $E,K\in D(A)$ 都成立。而《代数详论》第
\ref{more-algebra-section-base-change-RHom} 节正是如此构造 $\beta$ 的，
故证明完毕。
\end{proof}









\section{推前的右伴随与迹映射}
\label{spaces-duality-section-trace}

\noindent
设 $S$ 是概形，$f:X\to Y$ 是 $S$ 上拟紧且拟分离的代数空间之间
的态射。令 $a:D_\QCoh(\mathcal{O}_Y)\to D_\QCoh(\mathcal{O}_X)$
为引理 \ref{spaces-duality-lemma-twisted-inverse-image} 中的右伴随。由《范畴》第
\ref{categories-section-adjoint} 节，得到函子变换
$$
\text{Tr}_f : Rf_* \circ a \longrightarrow \text{id}
$$
对 $K\in D_\QCoh(\mathcal{O}_Y)$，相应的映射
$\text{Tr}_{f,K}:Rf_*a(K)\longrightarrow K$ 有时称为{\it 迹映射}。
它具有如下性质：刻画右伴随的双射
$$
\Hom_X(L, a(K)) \longrightarrow \Hom_Y(Rf_*L, K)
$$
（其中 $L\in D_\QCoh(\mathcal{O}_X)$）由下式给出：
$$
\varphi \longmapsto \text{Tr}_{f, K} \circ Rf_*\varphi
$$
典范映射 (\ref{spaces-duality-equation-sheafy-trace})
$$
Rf_*R\SheafHom_{\mathcal{O}_X}(L, a(K))
\longrightarrow
R\SheafHom_{\mathcal{O}_Y}(Rf_*L, K)
$$
由与 $\text{Tr}_{f,K}$ 复合得到。本节将考察的每个迹映射都是该迹映射
的特殊情形。在讨论这些特殊情形之前，先证明迹映射的形成与基变换
相交换。

\begin{lemma}[迹映射与基变换]
\label{spaces-duality-lemma-trace-map-and-base-change}
设有图 (\ref{spaces-duality-equation-base-change})。则映射
$1 \star \text{Tr}_f : Lg^* \circ Rf_* \circ a \to Lg^*$ 与
$\text{Tr}_{f'} \star 1 : Rf'_* \circ a' \circ Lg^* \to Lg^*$
通过下列基变换映射相一致：
$\beta : Lg^* \circ Rf_* \to Rf'_* \circ L(g')^*$
（《站点上的上同调》评注 \ref{sites-cohomology-remark-base-change}）
以及 $\alpha : L(g')^* \circ a \to a' \circ Lg^*$
（\ref{spaces-duality-equation-base-change-map}）。
更确切地说，下列函子变换之图交换：
$$
\xymatrix{
Lg^* \circ Rf_* \circ a
\ar[d]_{\beta \star 1} \ar[r]_-{1 \star \text{Tr}_f} &
Lg^* \\
Rf'_* \circ L(g')^* \circ a \ar[r]^{1 \star \alpha} &
Rf'_* \circ a' \circ Lg^* \ar[u]_{\text{Tr}_{f'} \star 1}
}
$$
\end{lemma}

\begin{proof}
本证明以 $f_*$ 代替 $Rf_*$，以 $g^*$ 代替 $Lg^*$，并省略与恒等变换
的 $\star$ 乘积，因为只要已知变换的源与目标，便可判断应补入哪些
恒等变换。% P10-E0268
回忆，$\beta:g^*\circ f_*\to f'_*\circ(g')^*$ 是同构，而 $\alpha$
用 $\beta$ 的伴随同构
$\beta^\vee:g'_*\circ a'\to a\circ g_*$ 定义；见引理
\ref{spaces-duality-lemma-flat-precompose-pus} 及其证明。首先注意，引理之图上方的
水平箭头等于复合
$$
g^* \circ f_* \circ a \to
g^* \circ f_* \circ a \circ g_* \circ g^* \to
g^* \circ g_* \circ g^* \to g^*
$$
，其中第一支箭头是 $(g^*,g_*)$ 的单位，第二支箭头是
$\text{Tr}_f$，第三支箭头是 $(g^*,g_*)$ 的余单位。这直接源于如下
事实：单位与余单位的复合 $g^*\to g^*\circ g_*\circ g^*\to g^*$
是恒等变换。考虑图
$$
\xymatrix{
& g^* \circ f_* \circ a \ar[ld]_\beta \ar[d] \ar[r]_{\text{Tr}_f} & g^* \\
f'_* \circ (g')^* \circ a \ar[dr] &
g^* \circ f_* \circ a \circ g_* \circ g^* \ar[d]_\beta \ar[ru] &
g^* \circ f_* \circ g'_* \circ a' \circ g^* \ar[l]_{\beta^\vee} \ar[d]_\beta &
f'_* \circ a' \circ g^* \ar[lu]_{\text{Tr}_{f'}} \\
& f'_* \circ (g')^* \circ a \circ g_* \circ g^* &
f'_* \circ (g')^* \circ g'_* \circ a' \circ g^* \ar[ru] \ar[l]_{\beta^\vee}
}
$$
% P10-E0267
在此图中，由《范畴》引理 \ref{categories-lemma-properties-2-cat-cats}，
或者更直接地由《范畴》定义
\ref{categories-definition-horizontal-composition} 之前的讨论，两个方块
均交换。由上述讨论，三角形交换。由《范畴》引理
\ref{categories-lemma-transformation-between-functors-and-adjoints}
，下方方块
$$
\xymatrix{
g^* \circ f_* \circ g'_* \circ a' \ar[d]_{\beta^\vee} \ar[r]_-\beta &
f'_* \circ (g')^* \circ g'_* \circ a' \ar[d] \\
g^* \circ f_* \circ a \circ g_* \ar[r] &
\text{id}
}
$$
交换，从而大图中的五边形交换。由于 $\beta$ 与 $\beta^\vee$ 都是
同构，并且按定义沿大图外侧得到
% P10-E0269
$\text{Tr}_f\circ\alpha\circ\beta$，引理得证。
\end{proof}

\noindent
设 $S$ 是概形，$f:X\to Y$ 是 $S$ 上拟紧且拟分离的代数空间之间
的态射。令 $a:D_\QCoh(\mathcal{O}_Y)\to D_\QCoh(\mathcal{O}_X)$
为引理 \ref{spaces-duality-lemma-twisted-inverse-image} 中 $Rf_*$ 的右伴随。由
《范畴》第 \ref{categories-section-adjoint} 节，得到函子变换
$$
\eta_f : \text{id} \to  a \circ Rf_*
$$
，称为该伴随的单位。

\begin{lemma}
\label{spaces-duality-lemma-unit-and-base-change}
设有图 (\ref{spaces-duality-equation-base-change})。则映射
$1 \star \eta_f : L(g')^* \to L(g')^* \circ a \circ Rf_*$ 与
$\eta_{f'} \star 1 : L(g')^* \to a' \circ Rf'_* \circ L(g')^*$
通过下列基变换映射相一致：
$\beta : Lg^* \circ Rf_* \to Rf'_* \circ L(g')^*$
（《站点上的上同调》评注 \ref{sites-cohomology-remark-base-change}）
以及 $\alpha : L(g')^* \circ a \to a' \circ Lg^*$
（\ref{spaces-duality-equation-base-change-map}）。
更确切地说，下列函子变换之图交换：
$$
\xymatrix{
L(g')^* \ar[r]_-{1 \star \eta_f} \ar[d]_{\eta_{f'} \star 1} &
L(g')^* \circ a \circ Rf_* \ar[d]^\alpha \\
a' \circ Rf'_* \circ L(g')^* &
a' \circ Lg^* \circ Rf_* \ar[l]_-\beta
}
$$
\end{lemma}

\begin{proof}
本证明与引理 \ref{spaces-duality-lemma-trace-map-and-base-change} 的证明对偶。
本证明以 $f_*$ 代替 $Rf_*$，以 $g^*$ 代替 $Lg^*$，并省略与恒等变换
的 $\star$ 乘积，因为只要已知变换的源与目标，便可判断应补入哪些
恒等变换。% P10-E0268
回忆，$\beta:g^*\circ f_*\to f'_*\circ(g')^*$ 是同构，而 $\alpha$
用 $\beta$ 的伴随同构
$\beta^\vee:g'_*\circ a'\to a\circ g_*$ 定义；见引理
\ref{spaces-duality-lemma-flat-precompose-pus} 及其证明。首先注意，引理之图左侧的
竖直箭头等于复合
$$
(g')^* \to (g')^* \circ g'_* \circ (g')^* \to
(g')^* \circ g'_* \circ a' \circ f'_* \circ (g')^* \to
a' \circ f'_* \circ (g')^*
$$
，其中第一支箭头是 $((g')^*,g'_*)$ 的单位，第二支箭头是
$\eta_{f'}$，第三支箭头是 $((g')^*,g'_*)$ 的余单位。这直接源于如下
事实：单位与余单位的复合
$(g')^*\to(g')^*\circ(g')_*\circ(g')^*\to(g')^*$ 是恒等变换。
考虑图
$$
\xymatrix{
& (g')^* \circ a \circ f_* \ar[r] &
(g')^* \circ a \circ g_* \circ g^* \circ f_*
\ar[ld]_\beta \\
(g')^* \ar[ru]^{\eta_f} \ar[dd]_{\eta_{f'}} \ar[rd] &
(g')^* \circ a \circ g_* \circ f'_* \circ (g')^* &
(g')^* \circ g'_* \circ a' \circ g^* \circ f_*
\ar[u]_{\beta^\vee} \ar[ld]_\beta \ar[d] \\
& (g')^* \circ g'_* \circ a' \circ f'_* \circ (g')^*
\ar[ld] \ar[u]_{\beta^\vee} &
a' \circ g^* \circ f_* \ar[lld]^\beta \\
a' \circ f'_* \circ (g')^*
}
$$
% P10-E0267
在此图中，由《范畴》引理 \ref{categories-lemma-properties-2-cat-cats}，
或者更直接地由《范畴》定义
\ref{categories-definition-horizontal-composition} 之前的讨论，两个方块
均交换。由上述讨论，三角形交换。由《范畴》引理
\ref{categories-lemma-transformation-between-functors-and-adjoints}
的对偶，下方方块
$$
\xymatrix{
\text{id} \ar[r] \ar[d] &
g'_* \circ a' \circ g^* \circ f_* \ar[d]^\beta \\
g'_* \circ a' \circ g^* \circ f_* \ar[r]^{\beta^\vee} &
a \circ g_* \circ f'_* \circ (g')^*
}
$$
交换，从而大图中的五边形交换。由于 $\beta$ 与 $\beta^\vee$ 都是
同构，并且按定义沿大图外侧得到
$\beta\circ\alpha\circ\eta_f$，引理得证。
\end{proof}





\section{推前的右伴随与拉回}
\label{spaces-duality-section-compare-with-pullback}

\noindent
设 $S$ 是概形，$f:X\to Y$ 是 $S$ 上拟紧且拟分离的代数空间之间
的态射。令 $a$ 为引理 \ref{spaces-duality-lemma-twisted-inverse-image} 中推前的右伴随。
对 $K,L\in D_\QCoh(\mathcal{O}_Y)$，有典范映射
$$
Lf^*K \otimes^\mathbf{L}_{\mathcal{O}_X} a(L)
\longrightarrow
a(K \otimes_{\mathcal{O}_Y}^\mathbf{L} L)
$$
具体而言，该映射伴随于映射
$$
Rf_*(Lf^*K \otimes^\mathbf{L}_{\mathcal{O}_X} a(L)) =
K \otimes^\mathbf{L}_{\mathcal{O}_Y} Rf_*(a(L))
\longrightarrow
K \otimes^\mathbf{L}_{\mathcal{O}_Y} L
$$
（等号由《空间的导出范畴》引理
\ref{spaces-perfect-lemma-cohomology-base-change} 给出），其中使用迹映射
$Rf_*a(L)\to L$。当 $L=\mathcal{O}_Y$ 时，得到映射
\begin{equation}
\label{spaces-duality-equation-compare-with-pullback}
Lf^*K \otimes^\mathbf{L}_{\mathcal{O}_X} a(\mathcal{O}_Y) \longrightarrow a(K)
\end{equation}
；它对 $K$ 具有函子性，并与特异三角相容。

\begin{lemma}
\label{spaces-duality-lemma-compare-with-pullback-perfect}
设 $S$ 是概形，$f:X\to Y$ 是 $S$ 上拟紧且拟分离的代数空间之间
的态射。上面为 $K,L\in D_\QCoh(\mathcal{O}_Y)$ 定义的映射
$Lf^*K \otimes^\mathbf{L}_{\mathcal{O}_X} a(L) \to
a(K \otimes_{\mathcal{O}_Y}^\mathbf{L} L)$
在 $K$ 完美时是同构。特别地，若 $K$ 完美，则
(\ref{spaces-duality-equation-compare-with-pullback}) 是同构。
\end{lemma}

\begin{proof}
令 $K^\vee$ 为 $K$ 的“对偶”，见《站点上的上同调》引理
\ref{sites-cohomology-lemma-dual-perfect-complex}。对
$M\in D_\QCoh(\mathcal{O}_X)$，有
\begin{align*}
\Hom_{D(\mathcal{O}_Y)}(Rf_*M, K \otimes^\mathbf{L}_{\mathcal{O}_Y} L)
& =
\Hom_{D(\mathcal{O}_Y)}(
Rf_*M \otimes^\mathbf{L}_{\mathcal{O}_Y} K^\vee, L) \\
& =
\Hom_{D(\mathcal{O}_X)}(
M \otimes^\mathbf{L}_{\mathcal{O}_X} Lf^*K^\vee, a(L)) \\
& =
\Hom_{D(\mathcal{O}_X)}(M,
Lf^*K \otimes^\mathbf{L}_{\mathcal{O}_X} a(L))
\end{align*}
% P10-E0270
第二个等号由 $a$ 的定义与投影公式（《站点上的上同调》引理
\ref{sites-cohomology-lemma-projection-formula}）得到，也可使用更一般的
《空间的导出范畴》引理
\ref{spaces-perfect-lemma-cohomology-base-change}。结论遂由 Yoneda 引理
推出。
\end{proof}

\begin{lemma}
\label{spaces-duality-lemma-restriction-compare-with-pullback}
设有图 (\ref{spaces-duality-equation-base-change})，并设
$K\in D_\QCoh(\mathcal{O}_Y)$。则下图交换：
$$
\xymatrix{
L(g')^*(Lf^*K \otimes^\mathbf{L}_{\mathcal{O}_X} a(\mathcal{O}_Y))
\ar[r] \ar[d] & L(g')^*a(K) \ar[d] \\
L(f')^*Lg^*K \otimes_{\mathcal{O}_{X'}}^\mathbf{L} a'(\mathcal{O}_{Y'})
\ar[r] & a'(Lg^*K)
}
$$
其中水平箭头分别是 $K$ 与 $Lg^*K$ 的映射
(\ref{spaces-duality-equation-compare-with-pullback})，竖直箭头则使用《站点上的上同调》
评注 \ref{sites-cohomology-remark-base-change} 及
(\ref{spaces-duality-equation-base-change-map}) 构造。
\end{lemma}

\begin{proof}
本证明以 $f_*$ 代替 $Rf_*$，以 $f^*$ 代替 $Lf^*$ 等，并以
$\otimes$ 代替 $\otimes^\mathbf{L}_{\mathcal{O}_X}$ 等。把
(\ref{spaces-duality-equation-compare-with-pullback}) 写成复合
\begin{align*}
f^*K \otimes a(\mathcal{O}_Y)
& \to
a(f_*(f^*K \otimes a(\mathcal{O}_Y))) \\
& \leftarrow
% P10-E0271
a(K \otimes f_*a(\mathcal{O}_K)) \\
& \to
a(K \otimes \mathcal{O}_Y) \\
& \to
a(K)
\end{align*}
。这里，第一支箭头是单位 $\eta_f$；第二支箭头是把 $a$ 施加于
《站点上的上同调》等式
(\ref{sites-cohomology-equation-projection-formula-map}) 所得，而由
《空间的导出范畴》引理
\ref{spaces-perfect-lemma-cohomology-base-change}，该箭头是同构；
第三支箭头是把 $a$ 施加于 $\text{id}_K\otimes\text{Tr}_f$ 所得；
第四支箭头是把 $a$ 施加于同构 $K\otimes\mathcal{O}_Y=K$ 所得。
为证明引理，只需说明这些映射中的每一个都产生引理陈述中的交换方块。
对 $\eta_f$ 与 $\text{Tr}_f$，这分别由引理
\ref{spaces-duality-lemma-unit-and-base-change} 与
\ref{spaces-duality-lemma-trace-map-and-base-change} 给出。对使用《站点上的上同调》等式
(\ref{sites-cohomology-equation-projection-formula-map}) 的箭头，这由
《站点上的上同调》评注
\ref{sites-cohomology-remark-compatible-with-diagram} 给出。
对乘法映射则显然。证明完毕。
\end{proof}











\section{真平坦态射推前的右伴随}
\label{spaces-duality-section-proper-flat}

\noindent
对于拟紧且拟分离的代数空间之间真、平坦且有限表示的态射，推前的
右伴随具有若干显著性质。

\begin{lemma}
\label{spaces-duality-lemma-proper-flat}
设 $S$ 是概形，$Y$ 是 $S$ 上拟紧且拟分离的代数空间，
$f:X\to Y$ 是真、平坦且有限表示的代数空间态射。令 $a$ 为引理
\ref{spaces-duality-lemma-twisted-inverse-image} 中
$Rf_*:D_\QCoh(\mathcal{O}_X)\to D_\QCoh(\mathcal{O}_Y)$ 的右伴随。
则 $a$ 与直和交换。
\end{lemma}

\begin{proof}
设 $P$ 是 $D(\mathcal{O}_X)$ 的完美对象。由《空间的导出范畴》引理
\ref{spaces-perfect-lemma-flat-proper-perfect-direct-image-general}，
复形 $Rf_*P$ 在 $Y$ 上完美。设 $K_i$ 是
$D_\QCoh(\mathcal{O}_Y)$ 中的一族对象。则
\begin{align*}
\Hom_{D(\mathcal{O}_X)}(P, a(\bigoplus K_i))
& =
\Hom_{D(\mathcal{O}_Y)}(Rf_*P, \bigoplus K_i) \\
& =
\bigoplus \Hom_{D(\mathcal{O}_Y)}(Rf_*P, K_i) \\
& =
\bigoplus \Hom_{D(\mathcal{O}_X)}(P, a(K_i))
\end{align*}
，因为完美对象是紧的（《空间的导出范畴》命题
\ref{spaces-perfect-proposition-compact-is-perfect}）。由于
$D_\QCoh(\mathcal{O}_X)$ 有完美生成元（《空间的导出范畴》定理
\ref{spaces-perfect-theorem-bondal-van-den-Bergh}），可知映射
$\bigoplus a(K_i)\to a(\bigoplus K_i)$ 是同构，即 $a$ 与直和交换。
\end{proof}

\begin{lemma}
\label{spaces-duality-lemma-compare-with-pullback-flat-proper}
设 $S$ 是概形，$Y$ 是 $S$ 上拟紧且拟分离的代数空间，
$f:X\to Y$ 是真、平坦且有限表示的代数空间态射。对
$D_\QCoh(\mathcal{O}_Y)$ 的每个对象 $K$，映射
(\ref{spaces-duality-equation-compare-with-pullback}) 都是同构。
\end{lemma}

\begin{proof}
由引理 \ref{spaces-duality-lemma-proper-flat}，$a$ 与直和交换。因此，使
(\ref{spaces-duality-equation-compare-with-pullback}) 成为同构的
$D_\QCoh(\mathcal{O}_Y)$ 对象，构成 $D_\QCoh(\mathcal{O}_Y)$ 的一个
严格满、饱和的三角子范畴，
并且在取直和下保持。由于 $D_\QCoh(\mathcal{O}_Y)$ 是由单个完美对象
生成的模范畴（《空间的导出范畴》定理
\ref{spaces-perfect-theorem-DQCoh-is-Ddga} 与
\ref{spaces-perfect-theorem-bondal-van-den-Bergh}），可如《代数详论》
评注 \ref{more-algebra-remark-P-resolution} 那样论证，只需对一个完美
对象证明 (\ref{spaces-duality-equation-compare-with-pullback}) 是同构。然而，对完美
对象该结论已经成立，见引理
\ref{spaces-duality-lemma-compare-with-pullback-perfect}。
\end{proof}

\begin{lemma}
\label{spaces-duality-lemma-properties-relative-dualizing}
设 $Y$ 是仿射概形，$f:X\to Y$ 是真、平坦且有限表示的代数空间态射。
令 $a$ 为引理 \ref{spaces-duality-lemma-twisted-inverse-image} 中
$Rf_*:D_\QCoh(\mathcal{O}_X)\to D_\QCoh(\mathcal{O}_Y)$ 的右伴随。
则
\begin{enumerate}
\item $a(\mathcal{O}_Y)$ 是 $D(\mathcal{O}_X)$ 中的 $Y$-完美对象；
\item $Rf_*a(\mathcal{O}_Y)$ 的正次数上同调层均为零；
\item $\mathcal{O}_X \to
R\SheafHom_{\mathcal{O}_X}(a(\mathcal{O}_Y), a(\mathcal{O}_Y))$
是同构。
\end{enumerate}
\end{lemma}

\begin{proof}
% P10-E0272
对 $D(\mathcal{O}_X)$ 的完美对象 $E$，有
\begin{align*}
Rf_*(E \otimes_{\mathcal{O}_X}^\mathbf{L} \omega_{X/Y}^\bullet)
& =
Rf_*R\SheafHom_{\mathcal{O}_X}(E^\vee, \omega_{X/Y}^\bullet) \\
& =
R\SheafHom_{\mathcal{O}_Y}(Rf_*E^\vee, \mathcal{O}_Y) \\
& =
(Rf_*E^\vee)^\vee
\end{align*}
第一个等号见《站点上的上同调》引理
\ref{sites-cohomology-lemma-dual-perfect-complex}。第二个等号见引理
\ref{spaces-duality-lemma-iso-on-RSheafHom}、评注 \ref{spaces-duality-remark-iso-on-RSheafHom} 及
《空间的导出范畴》引理
\ref{spaces-perfect-lemma-flat-proper-perfect-direct-image-general}。
第三个等号是对偶的定义。特别地，这些引文还表明所得对象是
$D(\mathcal{O}_Y)$ 的完美对象。于是由《空间的态射详论》引理
\ref{spaces-more-morphisms-lemma-characterize-relatively-perfect}，
$\omega_{X/Y}^\bullet$ 是 $Y$-完美的。这证明了 (1)。

\medskip\noindent
设 $M$ 是 $D_\QCoh(\mathcal{O}_Y)$ 的对象。则
\begin{align*}
\Hom_Y(M, Rf_*a(\mathcal{O}_Y)) & =
\Hom_X(Lf^*M, a(\mathcal{O}_Y)) \\
& =
\Hom_Y(Rf_*Lf^*M, \mathcal{O}_Y) \\
& =
% P10-E0274
\Hom_Y(M \otimes_{\mathcal{O}_Y}^\mathbf{L} Rf_*\mathcal{O}_Y, \mathcal{O}_Y)
\end{align*}
第一个等号由《站点上的上同调》引理
\ref{sites-cohomology-lemma-adjoint} 得到。
% P10-E0273
第二个等号由 $a$ 的构造得到；第三个等号由《空间的导出范畴》引理
\ref{spaces-perfect-lemma-cohomology-base-change} 得到。回忆：对某个 $N$，
$Rf_*\mathcal{O}_X$ 是 Tor 振幅位于 $[0,N]$ 的完美复形；见
《空间的导出范畴》引理
\ref{spaces-perfect-lemma-flat-proper-perfect-direct-image-general}。
因此，可以用位于次数 $[0,N]$ 的有限投射模复形表示
$Rf_*\mathcal{O}_X$（使用《代数详论》引理
\ref{more-algebra-lemma-perfect} 以及 $Y$ 仿射这一事实）。故若对某个
$i>0$ 有 $M=\mathcal{O}_Y[-i]$，则最后一个群为零。由于 $Y$ 仿射，
可知当 $i>0$ 时 $H^i(Rf_*a(\mathcal{O}_Y))=0$。这证明了 (2)。

\medskip\noindent
设 $E$ 是 $D_\QCoh(\mathcal{O}_X)$ 的完美对象。则
\begin{align*}
% P10-E0275
\Hom_X(E, R\SheafHom_{\mathcal{O}_X}(a(\mathcal{O}_Y), a(\mathcal{O}_Y))
& =
\Hom_X(E \otimes_{\mathcal{O}_X}^\mathbf{L} a(\mathcal{O}_Y),
a(\mathcal{O}_Y)) \\
& =
\Hom_Y(Rf_*(E \otimes_{\mathcal{O}_X}^\mathbf{L} a(\mathcal{O}_Y)),
\mathcal{O}_Y) \\
& =
\Hom_Y(Rf_*(R\SheafHom_{\mathcal{O}_X}(E^\vee, a(\mathcal{O}_Y))),
\mathcal{O}_Y) \\
& =
\Hom_Y(R\SheafHom_{\mathcal{O}_Y}(Rf_*E^\vee, \mathcal{O}_Y),
\mathcal{O}_Y) \\
& =
R\Gamma(Y, Rf_*E^\vee) \\
& =
\Hom_X(E, \mathcal{O}_X)
\end{align*}
第一个等号由《站点上的上同调》引理
\ref{sites-cohomology-lemma-internal-hom} 得到。第二个等号是 $a$ 的定义。
第三个等号来自对偶完美复形 $E^\vee$ 的构造，见《站点上的上同调》
引理 \ref{sites-cohomology-lemma-dual-perfect-complex}。第四个等号来自
证明第一段所得的等式
$Rf_*R\SheafHom_{\mathcal{O}_X}(E^\vee, \omega_{X/Y}^\bullet) =
R\SheafHom_{\mathcal{O}_Y}(Rf_*E^\vee, \mathcal{O}_Y)$
。第五个等号由完美复形的二次对偶性（《站点上的上同调》引理
\ref{sites-cohomology-lemma-dual-perfect-complex}）以及下列事实得到：
% P10-E0276
由《空间的导出范畴》引理
\ref{spaces-perfect-lemma-flat-proper-perfect-direct-image-general}
$Rf_*E$ 是完美的。% P10-E0277
最后一个等号是 $f$ 的 Leray 等式。这一串等式实质上通过 Yoneda 引理
证明了 (3)。具体而言，由《空间的导出范畴》引理
\ref{spaces-perfect-lemma-quasi-coherence-internal-hom}，对象
$R\SheafHom(a(\mathcal{O}_Y),a(\mathcal{O}_Y))$ 属于
$D_\QCoh(\mathcal{O}_X)$。在上式中取 $E=\mathcal{O}_X$，得到映射
$\alpha : \mathcal{O}_X \to
R\SheafHom_{\mathcal{O}_X}(a(\mathcal{O}_Y), a(\mathcal{O}_Y))$
，它对应于
$\text{id}_{\mathcal{O}_X} \in \Hom_X(\mathcal{O}_X, \mathcal{O}_X)$.
。由于上述所有同构对 $E$ 具有函子性，$\alpha$ 的锥是
$D_\QCoh(\mathcal{O}_X)$ 中的对象 $C$，并且对所有完美 $E$ 都有
$\Hom(E,C)=0$。由于完美对象生成该范畴（《空间的导出范畴》定理
\ref{spaces-perfect-theorem-bondal-van-den-Bergh}），可知 $\alpha$
是同构。
\end{proof}










\section{真平坦态射的相对对偶复形}
\label{spaces-duality-section-relative-dualizing-proper-flat}

\noindent
受《概形的对偶性》第 \ref{duality-section-proper-flat}、
\ref{duality-section-relative-dualizing-complexes} 节及本章第
\ref{spaces-duality-section-proper-flat} 节内容的启发，作如下定义。

\begin{definition}
\label{spaces-duality-definition-relative-dualizing-proper-flat}
设 $S$ 是概形，$f:X\to Y$ 是 $S$ 上真、平坦且有限表示的代数空间
态射。$X/Y$ 的{\it 相对对偶复形}是一个二元组
$(\omega_{X/Y}^\bullet,\tau)$，它由 $D(\mathcal{O}_X)$ 中的
$Y$-完美对象 $\omega_{X/Y}^\bullet$ 与映射
$$
\tau : Rf_*\omega_{X/Y}^\bullet \longrightarrow \mathcal{O}_Y
$$
组成，并且对任意笛卡尔方块
$$
\xymatrix{
X' \ar[r]_{g'} \ar[d]_{f'} & X \ar[d]^f \\
Y' \ar[r]^g & Y
}
$$
，若 $Y'$ 是仿射概形，则二元组
$(L(g')^*\omega_{X/Y}^\bullet, Lg^*\tau)$
同构于二元组
$(a'(\mathcal{O}_{Y'}), \text{Tr}_{f', \mathcal{O}_{Y'}})$；该二元组见第
\ref{spaces-duality-section-twisted-inverse-image},
\ref{spaces-duality-section-base-change-map},
\ref{spaces-duality-section-base-change-II},
\ref{spaces-duality-section-trace},
\ref{spaces-duality-section-compare-with-pullback} 与
\ref{spaces-duality-section-proper-flat} 节。
\end{definition}

\noindent
这里需要作几点说明。
\begin{enumerate}
\item 在定义 \ref{spaces-duality-definition-relative-dualizing-proper-flat} 中，可以删去
$\omega_{X/Y}^\bullet$ 是 $Y$-完美的这一假设。事实上，让 $Y'$ 遍历
$Y$ 的一个仿射平展覆盖的成员；由引理
\ref{spaces-duality-lemma-properties-relative-dualizing} 可知，
$\omega_{X/Y}^\bullet$ 在 $X$ 的一个平展覆盖的各成员上的限制都是
$Y$-完美的。这蕴含 $\omega_{X/Y}^\bullet$ 是 $Y$-完美的；见
《空间的态射详论》第
\ref{spaces-more-morphisms-section-relatively-perfect} 节。
\item 考虑相对对偶复形 $(\omega_{X/Y}^\bullet,\tau)$ 与定义
\ref{spaces-duality-definition-relative-dualizing-proper-flat} 中的笛卡尔方块。我们把
同构
$(L(g')^*\omega_{X/Y}^\bullet, Lg^*\tau) \cong
(a'(\mathcal{O}_{Y'}), \text{Tr}_{f', \mathcal{O}_{Y'}})$
的存在理解为：对任意 $M'\in D_\QCoh(\mathcal{O}_{X'})$，映射
$$
\Hom_{X'}(M', L(g')^*\omega_{X/Y}^\bullet)
\longrightarrow
\Hom_{Y'}(Rf'_*M', \mathcal{O}_{Y'}),\quad
\varphi' \longmapsto Lg^*\tau \circ Rf'_*\varphi'
$$
是同构。这由 $a'$ 的定义及第 \ref{spaces-duality-section-trace} 节的讨论得到。特别地，
Yoneda 引理保证该同构唯一。
\item 若 $Y$ 本身仿射，则相对对偶复形
$(\omega_{X/Y}^\bullet,\tau)$ 存在，并典范同构于
$(a(\mathcal{O}_Y),\text{Tr}_{f,\mathcal{O}_Y})$，其中 $a$ 是引理
\ref{spaces-duality-lemma-twisted-inverse-image} 中 $Rf_*$ 的右伴随，$\text{Tr}_f$
如第 \ref{spaces-duality-section-trace} 节所述。事实上，给定定义中的图，由引理
\ref{spaces-duality-lemma-more-base-change} 得到同构
$L(g')^*a(\mathcal{O}_Y)\to a'(\mathcal{O}_{Y'})$；由引理
\ref{spaces-duality-lemma-trace-map-and-base-change}，它与迹映射相容。
\end{enumerate}
这些信息恰好足以把局部给出的相对对偶复形粘合成整体对象。建议读者
跳过以下各引理的证明。

\begin{lemma}
\label{spaces-duality-lemma-relative-dualizing-RHom}
% P10-E0278
设 $S$ 是概形，设 $X\to Y$ 是真、平坦且有限表示的代数空间态射。
若 $(\omega_{X/Y}^\bullet,\tau)$ 是相对对偶复形，则
$\mathcal{O}_X \to
R\SheafHom_{\mathcal{O}_X}(\omega_{X/Y}^\bullet, \omega_{X/Y}^\bullet)$
是同构，并且 $Rf_*\omega_{X/Y}^\bullet$ 的正次数上同调层均为零。
\end{lemma}

\begin{proof}
只需在用一个平展于 $Y$ 的仿射概形作基变换后证明；此时结论由引理
\ref{spaces-duality-lemma-properties-relative-dualizing} 得到。
\end{proof}

\begin{lemma}
\label{spaces-duality-lemma-uniqueness-relative-dualizing}
设 $S$ 是概形，设 $X\to Y$ 是真、平坦且有限表示的代数空间态射。
若 $(\omega_j^\bullet,\tau_j)$（$j=1,2$）是 $X/Y$ 上的两个相对对偶
复形，则存在唯一同构
$(\omega_1^\bullet, \tau_1) \to (\omega_2^\bullet, \tau_2)$。
\end{lemma}

\begin{proof}
取平展态射 $g:Y'\to Y$，其中 $Y'$ 是仿射概形，并记基变换为
$X'=Y'\times_YX$。由定义
\ref{spaces-duality-definition-relative-dualizing-proper-flat} 及其后的讨论，存在唯一同构
$\iota : (\omega_1^\bullet|_{X'}, \tau_1|_{Y'}) \to
(\omega_2^\bullet|_{X'}, \tau_2|_{Y'})$。若 $Y''\to Y'$ 还是仿射
概形间的平展态射，且 $X''=Y''\times_YX$，则 $\iota|_{X''}$ 是唯一同构
$(\omega_1^\bullet|_{X''}, \tau_1|_{Y''}) \to
(\omega_2^\bullet|_{X''}, \tau_2|_{Y''})$（由唯一性）。此外，
$$
\text{Ext}^p_{X'}(\omega_1^\bullet|_{X'}, \omega_2^\bullet|_{X'}) = 0,
\quad p < 0
$$
，因为由引理 \ref{spaces-duality-lemma-relative-dualizing-RHom}，
$\mathcal{O}_{X'} \cong
R\SheafHom_{\mathcal{O}_{X'}}(\omega_1^\bullet|_{X'}, \omega_1^\bullet|_{X'})
\cong
R\SheafHom_{\mathcal{O}_{X'}}(\omega_1^\bullet|_{X'}, \omega_2^\bullet|_{X'})$。

\medskip\noindent
选取平展超覆盖 $b:V\to Y$，使得每个
$V_n=\coprod_{i\in I_n}Y_{n,i}$，其中各 $Y_{n,i}$ 均仿射。这可由
《超覆盖》引理 \ref{hypercovering-lemma-hypercovering-object} 与评注
\ref{hypercovering-remark-take-unions-hypercovering-X} 做到（把前一引理
产生的超覆盖换成每一次数上都是不交并的超覆盖）。记
$X_{n,i}=Y_{n,i}\times_YX$ 与 $U_n=V_n\times_YX$，于是得到平展超覆盖
$a:U\to X$（《超覆盖》引理
\ref{hypercovering-lemma-hypercovering-morphism-sites}），且
$U_n=\coprod X_{n,i}$。对 $a:U\to X$ 以及复形
$\omega_1^\bullet$、$\omega_2^\bullet$，均满足《单纯空间》引理
\ref{spaces-simplicial-lemma-fppf-neg-ext-zero-hom} 的假设。因此得到唯一态射
$\iota : \omega_1^\bullet \to \omega_2^\bullet$
，它在 $X_{0,i}$ 上的限制是唯一同构
$(\omega_1^\bullet|_{X_{0, i}}, \tau_1|_{Y_{0, i}}) \to
(\omega_2^\bullet|_{X_{0, i}}, \tau_2|_{Y_{0, i}})$。% P10-E0280
还需证明下图交换：
% P10-E0279
$$
\xymatrix{
Rf_*\omega_1^\bullet \ar[rd]_{\tau_1} \ar[rr]_{Rf_*\iota} & &
Rf_*\omega_1^\bullet \ar[ld]^{\tau_2} \\
& \mathcal{O}_Y
}
$$
不过，$Rf_*\omega_1^\bullet$ 与 $Rf_*\omega_2^\bullet$ 的正次数上同调层
均为零（引理 \ref{spaces-duality-lemma-relative-dualizing-RHom}），故只需限制到各仿射
概形 $Y_{0,i}$ 后证明交换性；在那里它依构造成立。
\end{proof}

\begin{lemma}
\label{spaces-duality-lemma-covering-enough}
% P10-E0281
设 $S$ 是概形，设 $X\to Y$ 是真、平坦且有限表示的代数空间态射。
设 $(\omega^\bullet,\tau)$ 是由 $D(\mathcal{O}_X)$ 中的 $Y$-完美对象与映射
$\tau : Rf_*\omega^\bullet \to \mathcal{O}_Y$。
组成的二元组。设有笛卡尔图
$$
\xymatrix{
X_i \ar[r]_{g_i'} \ar[d]_{f_i} & X \ar[d]^f \\
Y_i \ar[r]^{g_i} & Y
}
$$
，其中 $Y_i$ 仿射，$\{g_i:Y_i\to Y\}$ 是平展覆盖，并且有定义
\ref{spaces-duality-definition-relative-dualizing-proper-flat} 所述的二元组同构
$(\omega^\bullet|_{X_i},\tau|_{Y_i})
\to(a_i(\mathcal{O}_{Y_i}),\text{Tr}_{f_i,\mathcal{O}_{Y_i}})$。
则 $(\omega^\bullet,\tau)$ 是 $X$ 在 $Y$ 上的相对对偶复形。
\end{lemma}

\begin{proof}
令 $g:Y'\to Y$ 以及 $X',f',g',a'$ 如定义
\ref{spaces-duality-definition-relative-dualizing-proper-flat} 所述。置
$((\omega')^\bullet,\tau')=(L(g')^*\omega^\bullet,Lg^*\tau)$。
可以找到一个由仿射概形组成的有限平展覆盖 $\{Y'_j\to Y'\}$，它加细
$\{Y_i\times_YY'\to Y'\}$（《拓扑》引理
\ref{topologies-lemma-etale-affine}）。因此，对每个 $j$，存在 $i_j$ 以及
$Y$ 上的态射 $k_j:Y'_j\to Y_{i_j}$。考虑纤维积
$$
\xymatrix{
X'_j \ar[r]_{h_j'} \ar[d]_{f'_j} &
X' \ar[d]^{f'} \\
Y'_j \ar[r]^{h_j} & Y'
}
$$
。记 $k'_j:X'_j\to X_{i_j}$ 为诱导态射（即 $k_j$ 沿 $f_{i_j}$ 的基变换）。
经态射 $k'_j$ 把给定同构限制到 $Y'_j$，得到二元组同构
% P10-E0282
$((\omega')^\bullet|_{X'_j}, \tau'|_{Y'_j})
\to (a_j(\mathcal{O}_{Y'_j}), \text{Tr}_{f'_j, \mathcal{O}_{Y'_j}})$。
以 $f':X'\to Y'$ 代替 $f:X\to Y$ 后，便归结为下一段所解决的问题。

\medskip\noindent
设 $Y$ 仿射。问题是证明 $(\omega^\bullet,\tau)$ 同构于
$(\omega_{X/Y}^\bullet,\text{Tr})=
(a(\mathcal{O}_Y),\text{Tr}_{f,\mathcal{O}_Y})$。
可以假设覆盖 $\{Y_i\to Y\}$ 由仿射概形间的一个满平展态射
$\{g:Y'\to Y\}$ 给出。事实上，可以先用一个有限子覆盖替换
$\{g_i:Y_i\to Y\}$，再令
$g=\coprod g_i:Y'=\coprod Y_i\to Y$；略去若干细节。置
$X'=Y'\times_YX$，并令 $f',g'$ 如定义
\ref{spaces-duality-definition-relative-dualizing-proper-flat} 所述。此时给定的信息仅为同构
$$
(\omega^\bullet|_{X'}, \tau|_{Y'}) \to
(a'(\mathcal{O}_{Y'}), \text{Tr}_{f', \mathcal{O}_{Y'}})
$$
。由于 $(\omega_{X/Y}^\bullet,\text{Tr})$ 是相对对偶复形（见定义
\ref{spaces-duality-definition-relative-dualizing-proper-flat} 后的讨论），故有唯一同构
$$
(\omega_{X/Y}^\bullet|_{X'}, \text{Tr}|_{Y'}) \to
(a'(\mathcal{O}_{Y'}), \text{Tr}_{f', \mathcal{O}_{Y'}})
$$
。例如，唯一性由引理 \ref{spaces-duality-lemma-uniqueness-relative-dualizing} 得到。
合并以上两个同构，得到同构
$$
\alpha :
(\omega^\bullet|_{X'}, \tau|_{Y'}) \to
(\omega_{X/Y}^\bullet|_{X'}, \text{Tr}|_{Y'})
$$
。% P10-E0283
置 $Y''=Y'\times_YY'$ 与 $X''=Y''\times_YX$。由唯一性，$\alpha$ 到
$X''$ 的两个拉回必相同。又因为 $X'$ 上
$\omega_{X'/Y'}^\bullet$ 的负次数自 Ext 均为零（引理
\ref{spaces-duality-lemma-relative-dualizing-RHom}），而且沿任一投影
$Y'\times_Y\ldots\times_YY'\to Y'$ 拉回后仍然如此（略去一个小细节；
比较引理 \ref{spaces-duality-lemma-uniqueness-relative-dualizing} 的证明），故由
《单纯空间》引理 \ref{spaces-simplicial-lemma-fppf-neg-ext-zero-hom}，
$\alpha$ 下降为 $X$ 上的同构
$\omega^\bullet\to\omega_{X/Y}^\bullet$。
\end{proof}

\begin{lemma}
\label{spaces-duality-lemma-existence-relative-dualizing}
% P10-E0284
设 $S$ 是概形，设 $X\to Y$ 是真、平坦且有限表示的代数空间态射。
存在相对对偶复形 $(\omega_{X/Y}^\bullet,\tau)$。
\end{lemma}

\begin{proof}
选取平展超覆盖 $b:V\to Y$，使每个
$V_n=\coprod_{i\in I_n}Y_{n,i}$，其中各 $Y_{n,i}$ 均仿射。这可由
《超覆盖》引理 \ref{hypercovering-lemma-hypercovering-object} 与评注
\ref{hypercovering-remark-take-unions-hypercovering-X} 做到（把前一引理
产生的超覆盖换成每一次数上都是不交并的超覆盖）。记
$X_{n,i}=Y_{n,i}\times_YX$ 与 $U_n=V_n\times_YX$，于是得到平展超覆盖
$a:U\to X$（《超覆盖》引理
\ref{hypercovering-lemma-hypercovering-morphism-sites}），且
$U_n=\coprod X_{n,i}$。对每个 $n,i$，$X_{n,i}/Y_{n,i}$ 上都存在相对
对偶复形 $(\omega_{n,i}^\bullet,\tau_{n,i})$；见定义
\ref{spaces-duality-definition-relative-dualizing-proper-flat} 后的讨论。
对 $\varphi:[m]\to[n]$ 与 $i\in I_n$，考虑给定超覆盖结构中的态射
% P10-E0285
$g_{\varphi, i} : Y_{n, i} \to Y_{m, \alpha(\varphi)}$ 与
$g'_{\varphi, i} : X_{n, i} \to X_{m, \alpha(\varphi)}$
（《超覆盖》第 \ref{hypercovering-section-hypercovering-sites} 节）。于是
有唯一的二元组同构% P10-E0286
% P10-E0287
$$
\iota_{n, i, \varphi} :
(L(g'_{n, i})^*\omega_{n, i}^\bullet, Lg_{n, i}^*\tau_{n, i})
\longrightarrow
(\omega_{m, \alpha(\varphi)(i)}^\bullet, \tau_{m, \alpha(\varphi)(i)})
$$
；见定义 \ref{spaces-duality-definition-relative-dualizing-proper-flat} 后的讨论。由引理
\ref{spaces-duality-lemma-relative-dualizing-RHom}，$X_{n,i}$ 上
$\omega_{n,i}^\bullet$ 的负次数自 Ext 均为零。记
$(\omega_n^\bullet,\tau_n)$ 为用 $i\in I_n$ 的各二元组
$(\omega_{n,i}^\bullet,\tau_{n,i})$ 构造出的 $U_n/V_n$ 上二元组。
对 $\varphi:[m]\to[n]$ 与 $i\in I_n$，考虑单纯代数空间 $V$ 与 $U$
的结构态射 $g_\varphi:V_n\to V_m$ 与 $g'_\varphi:U_n\to U_m$。
此时有由各分片上的同构构造出的唯一二元组同构
$$
\iota_\varphi :
(L(g'_\varphi)^*\omega_n^\bullet, Lg_\varphi^*\tau_n)
\longrightarrow
(\omega_m^\bullet, \tau_m)
$$
。唯一性保证这些同构满足《单纯空间》定义
\ref{spaces-simplicial-definition-cartesian-derived-modules} 所表述的传递性
条件。对 $a:U\to X$、复形 $\omega_n^\bullet$ 及同构
$\iota_\varphi$，均满足《单纯空间》引理
\ref{spaces-simplicial-lemma-fppf-glue-neg-ext-zero} 的假设\footnote{%
% P10-E0288
该引理只用到 $\omega_0^\bullet$ 以及两个映射
$\delta_1^1,\delta_0^1:[1]\to[0]$。读者可以跳过所引引理证明的开头
几行，因为这里实际上已经给出了模导出范畴中的单纯系统。}。因此得到
$D_\QCoh(\mathcal{O}_X)$ 中的对象 $\omega^\bullet$，以及与两个同构
$\iota_{\delta^1_0}$、$\iota_{\delta^1_1}$ 相容的同构
$\iota_0 : \omega^\bullet|_{U_0} \to \omega_0^\bullet$
。最后应用《单纯空间》引理
\ref{spaces-simplicial-lemma-fppf-neg-ext-zero-hom}，得到唯一态射
$$
\tau : Rf_*\omega^\bullet \longrightarrow \mathcal{O}_Y
$$
，其在 $V_0$ 上的限制与 $\tau_0$ 一致。略去若干细节；例如，可与引理
\ref{spaces-duality-lemma-uniqueness-relative-dualizing} 证明的末尾比较，便可看出为何所需
的负次数 Ext 消失。% P10-E0289
由引理 \ref{spaces-duality-lemma-covering-enough}，二元组
$(\omega^\bullet,\tau)$ 是相对对偶复形，证明完成。
\end{proof}

\begin{lemma}
\label{spaces-duality-lemma-base-change-relative-dualizing}
设 $S$ 是概形。考虑 $S$ 上代数空间的笛卡尔方块
$$
\xymatrix{
X' \ar[d]_{f'} \ar[r]_{g'} & X \ar[d]^f \\
Y' \ar[r]^g & Y
}
$$
。设 $X\to Y$ 真、平坦且有限表示，并设
$(\omega_{X/Y}^\bullet,\tau)$ 是 $f$ 的相对对偶复形。则
$(L(g')^*\omega_{X/Y}^\bullet,Lg^*\tau)$ 是 $f'$ 的相对对偶复形。
\end{lemma}

\begin{proof}
由《空间的态射详论》引理
\ref{spaces-more-morphisms-lemma-base-change-relatively-perfect}，
$L(g')^*\omega_{X/Y}^\bullet$ 是 $Y'$-完美的。定义
\ref{spaces-duality-definition-relative-dualizing-proper-flat} 中的另一个条件由纤维积的
传递性成立。
\end{proof}







\section{与概形情形的比较}
\label{spaces-duality-section-comparison}

\noindent
本节还应补充大量内容。

\begin{lemma}
\label{spaces-duality-lemma-compare}
设 $S$ 是概形，设 $f:X\to Y$ 是 $S$ 上拟紧拟分离代数空间之间的
态射。设 $X$ 与 $Y$ 均可表，并设概形态射 $f_0:X_0\to Y_0$ 表示
$f$（这一记号别扭，但只是临时使用）。令
$a:D_\QCoh(\mathcal{O}_Y)\to D_\QCoh(\mathcal{O}_X)$ 为引理
\ref{spaces-duality-lemma-twisted-inverse-image} 中 $Rf_*$ 的右伴随；令
$a_0:D_\QCoh(\mathcal{O}_{Y_0})\to D_\QCoh(\mathcal{O}_{X_0})$ 为
《概形的对偶性》引理 \ref{duality-lemma-twisted-inverse-image} 中
$Rf_*$ 的右伴随。则下图交换：
$$
\xymatrix{
D_\QCoh(\mathcal{O}_{X_0})
\ar@{=}[rrrrrr]_{\text{《空间的导出范畴》引理
\ref{spaces-perfect-lemma-derived-quasi-coherent-small-etale-site}}}
& & & & & &
D_\QCoh(\mathcal{O}_X) \\
D_\QCoh(\mathcal{O}_{Y_0}) \ar[u]^{a_0}
\ar@{=}[rrrrrr]^{\text{《空间的导出范畴》引理
\ref{spaces-perfect-lemma-derived-quasi-coherent-small-etale-site}}}
& & & & & &
D_\QCoh(\mathcal{O}_Y) \ar[u]_a
}
$$
\end{lemma}

\begin{proof}
这由伴随的唯一性以及《空间的导出范畴》评注
\ref{spaces-perfect-remark-match-total-direct-images} 中的相容性得到。
\end{proof}
% Generated mechanically from the frozen formal-spaces.tex target.
% Do not edit this generated file; edit the bound locale source instead.

% OK, start here.
%

\setcounter{chapter}{86}
\chapter{形式代数空间}



\phantomsection
\label{formal-spaces-section-phantom}


\section{引言}
\label{formal-spaces-section-introduction}

\noindent
形式概形由 \cite{EGA} 引入。\cite{McQuillan} 引入了形式概形的一个更一般的版本，
\cite{Yasuda} 又引入了另一个版本。形式代数空间由 \cite{Kn} 引入。
相关材料以及更多内容可见 \cite{Abbes} 和 \cite{Fujiwara-Kato}。
本章介绍我们将采用的形式代数空间概念。我们的定义足够一般，
可以把文献中大多数类别的形式概形/空间作为全子范畴包括进来。

\medskip\noindent
虽然我们确实讨论了其中若干替代理论与本理论的比较，
但在理论的逻辑发展并不需要时，我们并不总是给出全部细节。

\medskip\noindent
除引入形式代数空间外，我们还证明若干非常基本的性质，并讨论几类态射。









\section{EGA 意义下的形式概形}
\label{formal-spaces-section-formal-schemes-EGA}

\noindent
本节回顾 \cite{EGA} 中形式概形的构造。这个概念虽然在代数几何中非常有用，
却未必总是应当考虑的正确概念。或许更准确的说法是：建立该理论时作出了若干选择，
而对于不同目的，其他选择可能更合适。事实上，文献中可以找到许多彼此密切相关、
并针对作者所考虑问题而调整的不同理论。不过，这里所概述的理论有一个主要优点：
我们处理的是确定的几何对象。

\medskip\noindent
开始之前，我们应指出拓扑环层，或者更一般地拓扑空间层的层条件所涉及的一个问题。
具体而言，下列大范畴
\begin{enumerate}
\item 拓扑空间范畴，
\item 拓扑群范畴，
\item 拓扑环范畴，
\item 给定拓扑环上的拓扑模范畴，
\end{enumerate}
连同它们到 $\textit{Sets}$ 的自然遗忘函子，并不是
《层》之第 \ref{sheaves-section-algebraic-structures} 节所定义的代数结构类型的实例。
因此，我们不能贸然对它们应用该章建立的机制。另一方面，上列每个范畴都有极限和等化子，
并且到集合、群、环、模的遗忘函子与这些构造交换
（见《拓扑》之引理 \ref{topology-lemma-limits}、
\ref{topology-lemma-topological-group-limits}、
\ref{topology-lemma-topological-ring-limits} 和
\ref{topology-lemma-topological-module-limits}）。
因此，我们可以像《层》之定义
\ref{sheaves-definition-sheaf-values-in-category} 那样定义层的概念，
而其底层的集合、群、环或模预层确实是层。关键区别在于：对于开覆盖
$U = \bigcup_{i \in I} U_i$，图表
$$
\xymatrix{
\mathcal{F}(U) \ar[r]
&
\prod\nolimits_{i\in I}
\mathcal{F}(U_i)
\ar@<1ex>[r] \ar@<-1ex>[r]
&
\prod\nolimits_{(i_0, i_1) \in I \times I}
\mathcal{F}(U_{i_0} \cap U_{i_1})
}
$$
必须是拓扑空间、拓扑群、拓扑环或拓扑模范畴中的等化子图表；
也就是说，第一个映射把 $\mathcal{F}(U)$ 识别为
$\prod_{i \in I} \mathcal{F}(U_i)$ 的一个子空间，
而后者带有乘积拓扑。

\medskip\noindent
拓扑空间、拓扑群、拓扑环或拓扑模层 $\mathcal{F}$ 在点 $x \in X$ 处的茎
$\mathcal{F}_x$，定义为相应范畴中对开邻域所取的余极限
$$
\mathcal{F}_x = \colim_{x\in U} \mathcal{F}(U)
$$
这等同于在集合、群、环或模层面取余极限
（见《拓扑》之引理 \ref{topology-lemma-colimits}、
\ref{topology-lemma-topological-group-colimits}、
\ref{topology-lemma-topological-ring-colimits} 和
\ref{topology-lemma-topological-module-colimits}），
但它还带有一个拓扑。警告：所得拓扑取决于所处的范畴，见
《例》之第 \ref{examples-section-colimit-topology} 节。
我们可以对拓扑空间、拓扑群、拓扑环或拓扑模的预层作层化，
而取茎与此操作交换，见注 \ref{formal-spaces-remark-sheafification-of-presheaves-in-top}。

\medskip\noindent
设 $f : X \to Y$ 是拓扑空间的连续映射。从拓扑空间、拓扑群、拓扑环或拓扑模层的范畴，
到 $Y$ 上相应层范畴有一个函子 $f_*$，照常由
$f_*\mathcal{F}(V) = \mathcal{F}(f^{-1}V)$ 定义。
（这个情形中的拉回留待后文讨论。）我们定义 $Y$ 上拓扑空间层 $\mathcal{G}$
与 $X$ 上拓扑空间层 $\mathcal{F}$ 之间的 $f$-映射
$\xi : \mathcal{G} \to \mathcal{F}$，其方式与《层》之定义
\ref{sheaves-definition-f-map} 完全相同，但另加条件：
对每个开集 $V \subset Y$，
$\xi_V : \mathcal{G}(V) \to \mathcal{F}(f^{-1}V)$ 都连续。与
《层》之引理 \ref{sheaves-lemma-f-map} 一样，有
$$
\{\mathcal{G}\text{ 到 }\mathcal{F}\text{ 的 }f\text{-映射}\} =
\Mor_{\Sh(Y, \textit{Top})}(\mathcal{G}, f_*\mathcal{F})
$$
拓扑群、拓扑环和拓扑模层的情形也相同。最后，设
$\xi : \mathcal{G} \to \mathcal{F}$ 是上述 $f$-映射。
给定 $x \in X$，令其像为 $y = f(x)$，则有茎之间的连续映射
$$
\xi_x : \mathcal{G}_y \longrightarrow \mathcal{F}_x
$$
其定义方式与《层》之定义
\ref{sheaves-definition-composition-f-maps} 之后的讨论完全相同。

\medskip\noindent
利用上述讨论，我们可以定义“局部拓扑赋环空间”的范畴 $LTRS$。
其对象是一个二元组 $(X, \mathcal{O}_X)$，其中 $X$ 是拓扑空间，
$\mathcal{O}_X$ 是拓扑环层，并且各茎 $\mathcal{O}_{X, x}$
（忘掉拓扑后）都是局部环。$LTRS$ 中从
$(X, \mathcal{O}_X)$ 到 $(Y, \mathcal{O}_Y)$ 的态射是二元组
$(f, f^\sharp)$，其中 $f : X \to Y$ 是拓扑空间的连续映射，
$f^\sharp : \mathcal{O}_Y \to \mathcal{O}_X$ 是一个 $f$-映射，
并且对每个 $x \in X$，诱导映射
$$
f^\sharp_x : \mathcal{O}_{Y, f(x)} \longrightarrow \mathcal{O}_{X, x}
$$
（忘掉拓扑后）是局部环的局部同态。复合的定义方式与局部赋环空间态射的复合完全相同。

\medskip\noindent
现在假设拓扑空间 $X$ 有一个由拟紧开集组成的基。给定集合、群、环或某环上模的层
$\mathcal{F}$，可以赋予 $\mathcal{F}$ 以拓扑空间、拓扑群、拓扑环或拓扑模层的结构。
具体而言，若 $U \subset X$ 是拟紧开集，就给 $\mathcal{F}(U)$ 赋离散拓扑。
若 $U \subset X$ 任意，则选择由拟紧开集组成的开覆盖
$U = \bigcup_{i \in I} U_i$，并给 $\mathcal{F}(U)$ 赋予从
$\prod_{i \in I} \mathcal{F}(U_i)$ 诱导的拓扑
（按照上述讨论，本来就应如此）。读者可验证（此处略去）按这种方式确实得到拓扑空间、
拓扑群、拓扑环或拓扑模层。若对每个拟紧开集 $U \subset X$，
$\mathcal{F}(U)$ 上的拓扑都是离散的，就称一个拓扑空间、拓扑群、拓扑环或拓扑模层
为{\it 伪离散的}。上述构造是遗忘函子的一个伴随，并诱导集合层范畴与伪离散拓扑空间层范畴之间的等价
（群、环、模的情形同样如此）。

\medskip\noindent
Grothendieck 和 Dieudonn\'e 首先定义了仿射形式概形。
它们对应于可容许拓扑环 $A$，见《代数续篇》之定义
\ref{more-algebra-definition-topological-ring}。具体而言，给定 $A$，考虑环 $A$
的一族定义理想基本系 $I_\lambda$。（在任意可容许拓扑环中，所有定义理想组成的族都是基本系。）
对每个 $\lambda$，可以考虑概形 $\Spec(A/I_\lambda)$。若
$I_\lambda \subset I_\mu$，则诱导态射
$$
\Spec(A/I_\mu) \to \Spec(A/I_\lambda)
$$
是一个增厚，因为对某个 $n$ 有 $I_\mu^n \subset I_\lambda$。
另一种看法是注意到每个映射
$$
\Spec(A/I_\lambda) \to \Spec(A)
$$
的像都同胚于 $A$ 的开素理想集合。这促使我们定义
$$
\text{Spf}(A) = \{\text{开素理想 }\mathfrak p \subset A\}
$$
并赋予它来自 $\Spec(A)$ 的拓扑。对每个 $\lambda$，可以把结构层
$\mathcal{O}_{\Spec(A/I_\lambda)}$ 看作 $\text{Spf}(A)$ 上的层。
令 $\mathcal{O}_\lambda$ 为上文所述相应的伪离散拓扑环层。然后令
$$
\mathcal{O}_{\text{Spf}(A)} = \lim \mathcal{O}_\lambda
$$
其中极限在拓扑环层范畴中取。二元组
$(\text{Spf}(A), \mathcal{O}_{\text{Spf}(A)})$ 称为 $A$ 的{\it 形式谱}。

\medskip\noindent
此时应当核验几件事。第一，茎 $\mathcal{O}_{\text{Spf}(A), x}$
（忘掉拓扑后）是局部环。第二，给定 $f \in A$，对相应开集
$D(f) \cap \text{Spf}(A)$，有拓扑环的等式
$$
\Gamma(D(f) \cap \text{Spf}(A), \mathcal{O}_{\text{Spf}(A)})
= A_{\{f\}} = \lim (A/I_\lambda)_f
$$
其中 $I_\lambda$ 是上述定义理想基本系。此外，环 $A_{\{f\}}$ 也是可容许的，并且
% P10-E0290
$(\text{Spf}(A_f), \mathcal{O}_{\text{Spf}(A_{\{f\}})})$
同构于
$(D(f) \cap \text{Spf}(A),
\mathcal{O}_{\text{Spf}(A)}|_{D(f) \cap \text{Spf}(A)})$。
最后，给定一对可容许拓扑环 $A,B$，有
\begin{equation}
\label{formal-spaces-equation-morphisms-affine-formal-schemes}
\Mor_{LTRS}((\text{Spf}(B), \mathcal{O}_{\text{Spf}(B)}),
(\text{Spf}(A), \mathcal{O}_{\text{Spf}(A)}))
= \Hom_{cont}(A, B)
\end{equation}
其中 $LTRS$ 是上文定义的“局部拓扑赋环空间”范畴。

\medskip\noindent
在此基础上，\cite{EGA} 把{\it 形式概形}定义为二元组
$(\mathfrak X, \mathcal{O}_\mathfrak X)$，其中 $\mathfrak X$ 是拓扑空间，
$\mathcal{O}_\mathfrak X$ 是拓扑环层，并且每个点都有一个在 $LTRS$ 中
同构于某个仿射形式概形的开邻域。形式概形的{\it 态射}
$f : (\mathfrak X, \mathcal{O}_\mathfrak X) \to
(\mathfrak Y, \mathcal{O}_\mathfrak Y)$ 是范畴 $LTRS$ 中的态射。

\medskip\noindent
设 $A$ 是带离散拓扑的环。则 $A$ 是可容许的，并且形式概形
$\text{Spf}(A)$ 等于 $\Spec(A)$。结构层
$\mathcal{O}_{\text{Spf}(A)}$ 是与 $\mathcal{O}_{\Spec(A)}$ 相伴的伪离散拓扑环层；
换言之，其底层环层等于 $\mathcal{O}_{\Spec(A)}$，并且在拟紧开集 $U$ 上，
环 $\mathcal{O}_{\text{Spf}(A)}(U) = \mathcal{O}_{\Spec(A)}(U)$
带离散拓扑，但一般的开集上并非如此。因此，可以给每个仿射概形关联一个仿射形式概形。
完全类似地，可以从一般概形 $(X, \mathcal{O}_X)$ 出发，把它关联到
$(X, \mathcal{O}'_X)$，其中 $\mathcal{O}'_X$ 是底层环层为
$\mathcal{O}_X$ 的伪离散拓扑环层。此构造与态射相容，并定义一个函子
\begin{equation}
\label{formal-spaces-equation-compare-schemes-formal-schemes}
\textit{概形} \longrightarrow \textit{形式概形}
\end{equation}
由式 (\ref{formal-spaces-equation-morphisms-affine-formal-schemes}) 立即可知，这个函子全忠实。

\medskip\noindent
设 $\mathfrak X$ 是形式概形。我们用公式
$\text{size}(\mathfrak X) = \max(\aleph_0, \kappa_1, \kappa_2)$
定义形式概形的{\it 大小}，其中 $\kappa_1$ 是 $\mathfrak X$ 的仿射形式开集的基数，
$\kappa_2$ 是所有 $\mathcal{O}_\mathfrak X(\mathfrak U)$ 的基数的上确界，
这里 $\mathfrak U \subset \mathfrak X$ 取遍这类仿射形式开集。

\begin{lemma}
\label{formal-spaces-lemma-fully-faithful}
如《集合》之引理 \ref{sets-lemma-construct-category} 选取一个概形范畴
$\Sch_\alpha$。给定形式概形 $\mathfrak X$，令
$$
h_\mathfrak X : (\Sch_\alpha)^{opp} \longrightarrow \textit{Sets},\quad
h_\mathfrak X(S) = \Mor_{\textit{形式概形}}(S, \mathfrak X)
$$
为其点函子。若 $\mathfrak X$ 的大小不过大，则
$$
\Mor_{\textit{形式概形}}(\mathfrak X, \mathfrak Y) =
\Mor_{\textit{PSh}(\Sch_\alpha)}(h_\mathfrak X, h_\mathfrak Y)
$$
\end{lemma}

\begin{proof}
首先注意，对任意形式概形 $\mathfrak X$，$h_\mathfrak X$ 都满足 Zariski 拓扑的层条件
（见《概形》之定义 \ref{schemes-definition-representable-by-open-immersions}）。
这是形式概形范畴中态射的局部性所致。此外，对形式概形的开浸入
$\mathfrak V \to \mathfrak W$，相应的函子变换
$h_\mathfrak V \to h_\mathfrak W$ 是单射且可由开浸入表示
（见《概形》之定义 \ref{schemes-definition-representable-by-open-immersions}）。
选取形式概形由仿射形式概形组成的开覆盖
$\mathfrak X = \bigcup \mathfrak U_i$。则函子族 $h_{\mathfrak U_i}$ 覆盖
$h_\mathfrak X$（见《概形》之定义
\ref{schemes-definition-representable-by-open-immersions}）。最后注意
$$
h_{\mathfrak U_i} \times_{h_\mathfrak X} h_{\mathfrak U_j} =
h_{\mathfrak U_i \cap \mathfrak U_j}
$$
所以，给出一个映射 $h_\mathfrak X \to h_\mathfrak Y$，等价于给出一族在交叠处一致的映射
$h_{\mathfrak U_i} \to h_\mathfrak Y$。因此，可以把引理中映射的双射性（分别地，单射性）
归约为各对 $(\mathfrak U_i, \mathfrak Y)$ 的双射性（分别地，单射性），
以及各对 $(\mathfrak U_i \cap \mathfrak U_j, \mathfrak Y)$ 的单射性（分别地，无需条件）。
这样便归约到 $\mathfrak X$ 是仿射形式概形的情形。设
$\mathfrak X = \text{Spf}(A)$，其中 $A$ 是某个可容许拓扑环；
再选取 $A$ 的一族定义理想基本系 $I_\lambda \subset A$。

\medskip\noindent
我们也可以在 $\mathfrak Y$ 上局部化。事实上，设
$\mathfrak V \subset \mathfrak Y$ 是开形式子概形，并设
$\varphi : h_\mathfrak X \to h_\mathfrak Y$。则
$$
h_\mathfrak V \times_{h_\mathfrak Y, \varphi} h_\mathfrak X
\to h_\mathfrak X
$$
可由开浸入表示。对所有 $\lambda$ 拉回到 $\Spec(A/I_\lambda)$，
得到开子概形 $U_\lambda \subset \Spec(A/I_\lambda)$。然而，若
$I_\lambda \subset I_\mu$，则态射
% P10-E0291
$\Spec(A/I_\lambda) \to \Spec(A/I_\mu)$ 把 $U_\mu$ 拉回到 $U_\lambda$。
因而它们黏合成一个开形式子概形 $\mathfrak U \subset \mathfrak X$。
一个直接的论证（略去）表明
$$
h_\mathfrak U = h_\mathfrak V \times_{h_\mathfrak Y} h_\mathfrak X
$$
由此可见，给定开覆盖 $\mathfrak Y = \bigcup \mathfrak V_j$ 和函子变换
$\varphi : h_\mathfrak X \to h_\mathfrak Y$，可得到 $\mathfrak X$ 的相应开覆盖。
由于 $\mathfrak X$ 是仿射的，可以把该覆盖加细为由仿射形式子概形组成的有限开覆盖
$\mathfrak X = \mathfrak U_1 \cup \ldots \cup \mathfrak U_n$。
换言之，对每个 $i$，存在某个 $j$ 和映射
$\varphi_i : h_{\mathfrak U_i} \to h_{\mathfrak V_j}$，使图表
$$
\xymatrix{
h_{\mathfrak U_i} \ar[r]_{\varphi_i} \ar[d] & h_{\mathfrak V_j} \ar[d] \\
h_{\mathfrak X} \ar[r]^\varphi & h_\mathfrak Y
}
$$
交换。再加上几个论证（略去），这说明只须在
$\mathfrak X$ 和 $\mathfrak Y$ 都是仿射形式概形时证明引理中的双射性。

\medskip\noindent
假设 $\mathfrak X$ 和 $\mathfrak Y$ 是仿射形式概形。写成
$\mathfrak X = \text{Spf}(A)$ 和 $\mathfrak Y = \text{Spf}(B)$。
设 $\varphi : h_\mathfrak X \to h_\mathfrak Y$ 是函子变换，
$I_\lambda \subset A$ 是一族定义理想基本系。典范包含态射
$i_\lambda : \Spec(A/I_\lambda) \to \mathfrak X$ 被映到态射
$\varphi(i_\lambda) : \Spec(A/I_\lambda) \to \mathfrak Y$。
由式 (\ref{formal-spaces-equation-morphisms-affine-formal-schemes})，它对应于连续映射
$\chi_\lambda : B \to A/I_\lambda$。由于 $\varphi$ 是函子变换，
若 $I_\lambda \subset I_\mu$，则复合
$B \to A/I_\lambda \to A/I_\mu$ 等于 $\chi_\mu$。换言之，得到环映射
$$
\chi = \lim \chi_\lambda : B \longrightarrow \lim A/I_\lambda = A
$$
这是连续同态，因为对每个 $\lambda$，$I_\lambda$ 的逆像都是开集
（$A/I_\lambda$ 带离散拓扑，而 $\chi_\lambda$ 连续）。因此，由式
(\ref{formal-spaces-equation-morphisms-affine-formal-schemes}) 得到态射
$\text{Spf}(\chi) : \mathfrak X \to \mathfrak Y$。
我们略去核验：在此情形中，这个构造是引理中映射的逆。

\medskip\noindent
集合论说明。为了使上述论证在给定的概形范畴 $\Sch_\alpha$ 上成立，只须保证证明中使用的所有概形
都同构于 $\Sch_\alpha$ 的对象。事实上，仔细分析可知，只须上述出现的概形
$\Spec(A/I_\lambda)$ 同构于 $\Sch_\alpha$ 的对象。为此，假设
$\mathfrak X$ 的大小不超过 $\Sch_\alpha$ 中某个概形的大小当然就足够了。
\end{proof}


\begin{lemma}
\label{formal-spaces-lemma-formal-scheme-sheaf-fppf}
\begin{slogan}
形式概形是 fpqc 层
\end{slogan}
设 $\mathfrak X$ 是形式概形。其点函子 $h_\mathfrak X$
（见引理 \ref{formal-spaces-lemma-fully-faithful}）满足 fpqc 覆盖的层条件。
\end{lemma}

\begin{proof}
《拓扑》之引理 \ref{topologies-lemma-sheaf-property-fpqc}
把问题归约为 Zariski 覆盖以及 $R \to S$ 忠实平坦时的覆盖
$\{\Spec(S) \to \Spec(R)\}$。我们在引理 \ref{formal-spaces-lemma-fully-faithful}
的证明中已观察到，$h_\mathfrak X$ 满足 Zariski 覆盖的层条件。

\medskip\noindent
设 $R \to S$ 是忠实平坦环映射。用
$\pi : \Spec(S) \to \Spec(R)$ 表示相应的概形态射；它是满且平坦的。
设 $f : \Spec(S) \to \mathfrak X$ 是态射，并且作为从
$\Spec(S \otimes_R S)$ 到 $\mathfrak X$ 的映射，有
$f \circ \text{pr}_1 = f \circ \text{pr}_2$。
由《下降》之引理 \ref{descent-lemma-equiv-fibre-product}，作为底层集合间的映射，
$f$ 具有形式 $f = g \circ \pi$，其中
$g : \Spec(R) \to \mathfrak X$ 是某个（集合论意义下的）映射。
由《态射》之引理 \ref{morphisms-lemma-fpqc-quotient-topology}
以及 $f$ 连续可知 $g$ 连续。

\medskip\noindent
取 $y \in \Spec(R)$。选取包含 $g(y)$ 的仿射形式开子概形
$\mathfrak U \subset \mathfrak X$。设
$\mathfrak U = \text{Spf}(A)$，其中 $A$ 是某个可容许拓扑环。
由上述结论，可选取 $r \in R$，使得
$y \in D(r) \subset g^{-1}(\mathfrak U)$。
$f$ 在 $\pi^{-1}(D(r))$ 上到 $\mathfrak U$ 的限制，由式
(\ref{formal-spaces-equation-morphisms-affine-formal-schemes}) 对应于连续环映射
$A \to S_r$。按对 $f$ 的假设，两个诱导环映射
$A \to S_r \otimes_{R_r} S_r = (S \otimes_R S)_r$ 相等。
注意 $R_r \to S_r$ 忠实平坦。由《下降》之引理
\ref{descent-lemma-ff-exact}，两个箭头
$S_r \to S_r \otimes_{R_r} S_r$ 的等化子是 $R_r$。
于是 $A \to S_r$ 唯一地分解经过一个映射 $A \to R_r$；
这个映射也连续，因为它与映射 $A \to S_r$ 有相同的（开）核。
该映射再由式 (\ref{formal-spaces-equation-morphisms-affine-formal-schemes})
给出态射 $D(r) \to \mathfrak U$。

\medskip\noindent
至此证明了什么？我们证明了：对任意 $y \in \Spec(R)$，存在标准仿射开集
$y \in D(r) \subset \Spec(R)$，使得态射
$f|_{\pi^{-1}(D(r))} : \pi^{-1}(D(r)) \to \mathfrak X$
唯一地分解经过某个态射 $D(r) \to \mathfrak X$。
我们略去核验这些态射黏合成所需态射 $\Spec(R) \to \mathfrak X$。
\end{proof}

\begin{remark}[McQuillan 的变体]
\label{formal-spaces-remark-mcquillan}
McQuillan 给出了形式概形构造的一个变体，见 \cite{McQuillan}。
他建议稍微减弱可容许性条件。回忆，可容许拓扑环是一个完备的
（按我们的约定也分离的）线性拓扑环 $A$，且存在一个定义理想：
开理想 $I$，使得 $0$ 的任意邻域都包含某个 $I^n$，其中 $n \geq 1$。
McQuillan 使用我们将称为{\it 弱可容许}的拓扑环。
弱可容许拓扑环 $A$ 是一个完备的（按我们的约定也分离的）线性拓扑环，
且存在一个{\it 弱定义理想}：开理想 $I$，使得对所有 $f \in I$，
当 $n \to \infty$ 时有 $f^n \to 0$。与可容许情形类似，
若 $I$ 是弱定义理想而 $J \subset A$ 是开理想，则 $I \cap J$ 是弱定义理想。
因此，弱定义理想组成 $0$ 的开邻域基本系，可以沿着与上文大致相同的路线，
以此概念为基础定义一个更大的形式概形范畴。引理 \ref{formal-spaces-lemma-fully-faithful}
和 \ref{formal-spaces-lemma-formal-scheme-sheaf-fppf} 的对应结论在这个背景下仍然成立
（证明相同）。
\end{remark}

\begin{remark}[拓扑空间预层的层化]
\label{formal-spaces-remark-sheafification-of-presheaves-in-top}
\begin{reference}
\cite{Gray}
\end{reference}
本注简要讨论拓扑空间预层的层化。完全相同的论证适用于拓扑阿贝尔群、拓扑环以及
（给定拓扑环上的）拓扑模的预层。为了在正确的一般性下进行讨论，设我们工作在站点
$\mathcal{C}$ 上。只关心拓扑空间 $X$ 上（预）层情形的读者，应把
$\mathcal{C}$ 的对象理解为 $X$ 的开集，把 $\mathcal{C}$ 的态射理解为开集的包含，
并把 $\mathcal{C}$ 中的覆盖理解为 $X$ 中的覆盖，见《站点》之例
\ref{sites-example-site-topological}。以
$\Sh(\mathcal{C}, \textit{Top})$ 表示 $\mathcal{C}$ 上拓扑空间层的范畴，
以 $\textit{PSh}(\mathcal{C}, \textit{Top})$ 表示 $\mathcal{C}$ 上拓扑空间预层的范畴。
设 $\mathcal{F}$ 是 $\mathcal{C}$ 上的拓扑空间预层。其层化
$\mathcal{F}^\#$ 应函子性地对
$\Sh(\mathcal{C}, \textit{Top})$ 中的 $\mathcal{G}$ 满足公式
$$
\Mor_{\textit{PSh}(\mathcal{C}, \textit{Top})}(\mathcal{F}, \mathcal{G})
=
\Mor_{\Sh(\mathcal{C}, \textit{Top})}(\mathcal{F}^\#, \mathcal{G})
$$
换言之，我们要构造包含函子
$\Sh(\mathcal{C}, \textit{Top}) \to
\textit{PSh}(\mathcal{C}, \textit{Top})$ 的左伴随。

我们首先断言 $\Sh(\mathcal{C}, \textit{Top})$ 有极限，且包含函子与极限交换。
事实上，给定范畴 $\mathcal{I}$ 以及到
$\Sh(\mathcal{C}, \textit{Top})$ 的函子 $i \mapsto \mathcal{G}_i$，
我们只须定义
$$
(\lim \mathcal{G}_i)(U) = \lim \mathcal{G}_i(U)
$$
其中右侧极限在拓扑空间范畴中取
（《拓扑》之引理 \ref{topology-lemma-limits}）。这确实定义一个层，
因为极限与极限交换（《范畴》之引理
\ref{categories-lemma-colimits-commute}），特别是与层公理所使用的乘积和等化子交换。
最后，从 $\mathcal{F}$ 到 $\lim \mathcal{G}_i$ 的预层态射显然等同于一族相容的态射
$\mathcal{F} \to \mathcal{G}_i$。换言之，对象 $\lim \mathcal{G}_i$
是拓扑空间预层范畴中的极限，从而当然也是拓扑空间层范畴中的极限。

第二个断言是：任意预层态射 $\mathcal{F} \to \mathcal{G}$，其中
$\mathcal{G}$ 是 $\Sh(\mathcal{C}, \textit{Top})$ 的对象，都分解经过某个大小有界的子层
$\mathcal{G}' \subset \mathcal{G}$。这里，把拓扑空间层 $\mathcal{H}$ 的{\it 大小}
$|\mathcal{H}|$ 定义为基数
$\sup_{U \in \Ob(\mathcal{C})} |\mathcal{H}(U)|$。
为证明该断言，令
$$
\mathcal{G}'(U) =
\left\{
\quad
s \in \mathcal{G}(U)
\quad \middle| \quad
\begin{matrix}
\text{存在覆盖 }\{U_i \to U\}_{i \in I} \\
\text{使得 }s|_{U_i} \in \Im(\mathcal{F}(U_i) \to \mathcal{G}(U_i))
\end{matrix}
\quad
\right\}
$$
给 $\mathcal{G}'(U)$ 赋诱导拓扑。则 $\mathcal{G}'$ 是拓扑空间层（细节略），
$\mathcal{G}' \to \mathcal{G}$ 是一个态射，且给定映射
$\mathcal{F} \to \mathcal{G}$ 分解经过它。此外，$\mathcal{G}'$ 的大小由某个只依赖于
$\mathcal{C}$ 和预层 $\mathcal{F}$ 的基数 $\kappa$ 控制
（提示：按我们的约定，$\mathcal{C}$ 中的覆盖组成一个集合）。综合以上结果，
《范畴》之定理 \ref{categories-theorem-adjoint-functor} 的假设得到满足，
从而得到层化，它是从层到预层的包含函子的左伴随。

最后，设 $p$ 是站点 $\mathcal{C}$ 的一个点，由函子
$u : \mathcal{C} \to \textit{Sets}$ 给出，见《站点》之定义
\ref{sites-definition-point}。对拓扑空间 $M$，规则
$$
U \mapsto \text{Map}(u(U), M) = \prod\nolimits_{x \in u(U)} M
$$
定义的预层在带乘积拓扑后是拓扑空间层。因此，与《站点》之引理
\ref{sites-lemma-point-pushforward-sheaf} 的证明完全相同的论证表明
$\mathcal{F}_p = \mathcal{F}^\#_p$；换言之，层化与在一点取茎交换。
\end{remark}





\section{约定和记号}
\label{formal-spaces-section-conventions}

\noindent
从现在起采用的约定与《空间的性质》之第
\ref{spaces-properties-section-conventions} 节中的约定相似。
因此，从现在起始终假设所有概形都包含在一个大 fppf 站点 $\Sch_{fppf}$ 中。
所考虑的每个环 $A$ 也都满足：$\Spec(A)$（同构于）这个大站点的某个对象。
对拓扑环 $A$，只假设所有离散商都具有这个性质
（不过通常会作更强的假设，比较注 \ref{formal-spaces-remark-set-theoretic}）。

\medskip\noindent
设 $S$ 是概形，$X$ 是 $S$ 上的“空间”，即 $(\Sch/S)_{fppf}$ 上的层。
本章把 $X$ 在 $(\Sch/S)_{fppf}$ 上层范畴中与自身的乘积写作
$X \times_S X$，而不是 $X \times X$。此外，若 $X$ 和 $Y$ 是“空间”，
则我们说“设 $f : X \to Y$ 是态射”，意指 $f$ 是函子的自然变换，
也就是 $(\Sch/S)_{fppf}$ 上层的映射。类似地，若 $U$ 是 $S$ 上的概形，
$X$ 是 $S$ 上的“空间”，则我们说“设 $f : U \to X$ 是态射”或
“设 $g : X \to U$ 是态射”，意指 $f$ 或 $g$ 是层映射
$h_U \to X$ 或 $X \to h_U$，其中 $h_U$ 如《范畴》之例
\ref{categories-example-hom-functor} 所定义。






\section{拓扑环与拓扑模}
\label{formal-spaces-section-topological-module}

\noindent
本节是《代数续篇》之第 \ref{more-algebra-section-topological-ring} 节的继续。
设 $R$ 是拓扑环，$M$ 是线性拓扑化的 $R$-模。当我们说“{\it 设 $M_\lambda$
是一族开子模基本系}”时，意思是每个 $M_\lambda$ 都是开子模，且 $0$ 的任意邻域
都包含某个 $M_\lambda$。换言之，这表示 $M_\lambda$ 是 $M$ 中由子模组成的
$0$ 的邻域基本系。类似地，若 $R$ 是线性拓扑环，则说“{\it 设 $I_\lambda$
是一族开理想基本系}”，意思是 $I_\lambda$ 是 $R$ 中由理想组成的 $0$ 的邻域基本系。

\begin{example}
\label{formal-spaces-example-what-does-it-mean}
设 $R$ 是线性拓扑环，$M$ 是线性拓扑化的 $R$-模。设 $I_\lambda$
是 $R$ 中的一族开理想基本系，$M_\mu$ 是 $M$ 中的一族开子模基本系。
$+ : M \times M \to M$ 的连续性是自动的，而 $R \times M \to M$ 的连续性表示
$$
\forall f, x, \mu\ \exists \lambda, \nu,\ (f + I_\lambda)(x + M_\nu)
\subset fx + M_\mu
$$
当 $M_\nu \subset M_\mu$ 时，有
$fM_\nu + I_\lambda M_\nu \subset M_\mu$，所以该条件等价于
$$
\forall x, \mu\ \exists \lambda\ I_\lambda x \subset M_\mu
$$
然而，给定 $\mu$，未必存在 $\lambda$ 使
$I_\lambda M \subset M_\mu$。例如，考虑带 $t$-进拓扑的
$R = k[[t]]$，以及 $M = \bigoplus_{n \in \mathbf{N}} R$，其开子模基本系为
$$
M_m = \bigoplus\nolimits_{n \in \mathbf{N}} t^{nm}R
$$
由于每个 $x \in M$ 只有有限多个非零坐标，可见给定 $m$ 和 $x$，
存在 $k$ 使 $t^k x \in M_m$。因此，$M$ 是线性拓扑化的 $R$-模，
但给定 $m$ 时，未必存在 $k$ 使 $t^kM \subset M_m$。
另一方面，若 $R \to S$ 是线性拓扑环的连续映射，则相应陈述确实成立；
也就是说，对每个开理想 $J \subset S$，存在开理想 $I \subset R$，使得
$IS \subset J$（读者可容易地从映射 $R \to S$ 的连续性推出这一点）。
\end{example}

\begin{lemma}
\label{formal-spaces-lemma-closed}
设 $R$ 是拓扑环。设 $M$ 是线性拓扑化的 $R$-模，并设
$M_\lambda$（$\lambda \in \Lambda$）是一族开子模基本系。
设 $N \subset M$ 是子模。则 $N$ 的闭包为
$\bigcap_{\lambda \in \Lambda} (N + M_\lambda)$。
\end{lemma}

\begin{proof}
因为每个 $N + M_\lambda$ 都是开集，所以也都是闭集；故其交是闭集。
若 $x \in M$ 不在 $N$ 的闭包中，则对某个 $\lambda$ 有
% P10-E0292
$(x + M_\lambda) \cap N = 0$。所以 $x \not \in N + M_\lambda$。
这就证明了引理。
\end{proof}

\noindent
除非另有说明，我们总给子模和商模赋诱导拓扑。设 $M$ 是拓扑环 $R$ 上的线性拓扑化模，
$0 \to N \to M \to Q \to 0$ 是 $R$-模的短正合列。
若 $M_\lambda$ 是 $M$ 的一族开子模基本系，则
$N \cap M_\lambda$ 是 $N$ 的一族开子模基本系。若
$\pi : M \to Q$ 是商映射，则 $\pi(M_\lambda)$ 是 $Q$ 的一族开子模基本系。
特别地，这些诱导拓扑都是线性拓扑。

\begin{lemma}
\label{formal-spaces-lemma-closure}
设 $R$ 是拓扑环。设 $M$ 是线性拓扑化的 $R$-模，$N \subset M$ 是子模。则
\begin{enumerate}
\item $0 \to N^\wedge \to M^\wedge \to (M/N)^\wedge$ 正合；
\item $N^\wedge$ 是 $N \to M^\wedge$ 的像的闭包。
\end{enumerate}
\end{lemma}

\begin{proof}
设 $M_\lambda$（$\lambda \in \Lambda$）是一族开子模基本系。
则 $N \cap M_\lambda$ 是 $N$ 的一族开子模基本系，并且
% P10-E0293
$M_\lambda + N/N$ 是 $M/N$ 的一族开子模基本系。因此，(1) 来自正合列
$$
0 \to N/N \cap M_\lambda \to M/M_\lambda
\to M/(M_\lambda + N) \to 0
$$
的正合性以及取极限与极限交换这一事实。第二个陈述由此推出，因为
$N \to N^\wedge$ 的像稠密，而 $M^\wedge \to (M/N)^\wedge$ 的核是闭的。
\end{proof}

\begin{lemma}
\label{formal-spaces-lemma-quotient-by-closed}
设 $R$ 是拓扑环。设 $M$ 是完备的线性拓扑化 $R$-模，
$N \subset M$ 是闭子模。若 $M$ 有可数的 $0$ 邻域基本系，则
$M/N$ 完备，且映射 $M \to M/N$ 是开映射。
\end{lemma}

\begin{proof}
设 $M_n$（$n \in \mathbf{N}$）是 $M$ 的一族开子模基本系。
可以假设对所有 $n$ 都有 $M_{n + 1} \subset M_n$。
系统 $(M_n + N)/N$ 是 $M/N$ 中的一个基本系。因此，只须证明
$M/N = \lim M/(M_n + N)$。考虑短正合列
$$
0 \to N/N \cap M_n \to M/M_n \to M/(M_n + N) \to 0
$$
由于系统 $\{N/N\cap M_n\}$ 的转移映射都是满射，
《代数》之引理 \ref{algebra-lemma-ML-exact-sequence} 表明
$M = \lim M/M_n$（利用 $M$ 的完备性）满射到
$\lim M/(M_n + N)$。由于 $N$ 闭，映射
$M \to \lim M/(M_n + N)$ 的核为 $N$
（见引理 \ref{formal-spaces-lemma-closed}）。最后，按商拓扑的定义，$M \to M/N$ 是开映射。
\end{proof}

\begin{lemma}
\label{formal-spaces-lemma-ses}
\begin{reference}
\cite[Theorem 8.1]{Ma}
\end{reference}
设 $R$ 是拓扑环。设 $M$ 是线性拓扑化的 $R$-模，$N \subset M$ 是子模。
假设 $M$ 有可数的 $0$ 邻域基本系。则
\begin{enumerate}
\item $0 \to N^\wedge \to M^\wedge \to (M/N)^\wedge \to 0$ 正合；
\item $N^\wedge$ 是 $N \to M^\wedge$ 的像的闭包；
\item $M^\wedge \to (M/N)^\wedge$ 是开映射。
\end{enumerate}
\end{lemma}

\begin{proof}
由引理 \ref{formal-spaces-lemma-closure}，列
$0 \to N^\wedge \to M^\wedge \to (M/N)^\wedge$ 正合，且 (2) 成立。
这给出典范映射 $c : M^\wedge/N^\wedge \to (M/N)^\wedge$。
由引理 \ref{formal-spaces-lemma-quotient-by-closed}，模 $M^\wedge/N^\wedge$ 完备，
且 $M^\wedge \to M^\wedge/N^\wedge$ 是开映射。
由完备化的泛性质，得到典范映射
$b : (M/N)^\wedge \to M^\wedge/N^\wedge$。
由于 $b$ 和 $c$ 在一个稠密子集上互逆，它们彼此互逆。
\end{proof}

\begin{lemma}
\label{formal-spaces-lemma-completion-adic-star}
设 $R$ 是拓扑环，$M$ 是拓扑 $R$-模。设 $I \subset R$ 是有限生成理想。
假设 $M$ 有一个拓扑为 $I$-进拓扑的开子模。则 $M^\wedge$
也有一个拓扑为 $I$-进拓扑的开子模，并且对所有 $n \geq 1$ 有
$M^\wedge/I^n M^\wedge = M/I^nM$。
\end{lemma}

\begin{proof}
设 $M' \subset M$ 是一个拓扑为 $I$-进拓扑的开子模。
则 $\{I^nM'\}_{n \geq 1}$ 是 $M$ 的一族开子模基本系。因此
$M^\wedge = \lim M/I^nM'$ 包含
$(M')^\wedge = \lim M'/I^nM'$ 作为开子模，并且由《代数》之引理
\ref{algebra-lemma-hathat-finitely-generated}，$(M')^\wedge$ 上的拓扑是 $I$-进拓扑。
因为 $I$ 有限生成，$I^n$ 也有限生成，设其生成元为 $f_1, \ldots, f_r$。
注意，满射
$(f_1, \ldots, f_r) : M^{\oplus r} \to I^n M$
按上述对 $M$ 上拓扑的描述是连续且开的。把引理 \ref{formal-spaces-lemma-ses}
应用于这个满射以及短正合列
$0 \to I^nM \to M \to M/I^nM \to 0$，可知
$$
(f_1, \ldots, f_r) :
(M^\wedge)^{\oplus r} \longrightarrow M^\wedge
$$
满射到满射 $M^\wedge \to M/I^nM$ 的核。由于
$f_1, \ldots, f_r$ 生成 $I^n$，结论成立。
\end{proof}

\begin{definition}
\label{formal-spaces-definition-toplogy-tensor-product}
\begin{reference}
比较 \cite[0, Section 7.7]{EGA}
\end{reference}
设 $R$ 是拓扑环，$M$ 和 $N$ 是线性拓扑化的 $R$-模。
$M$ 与 $N$ 的{\it 张量积}是（通常的）张量积 $M \otimes_R N$，
并赋予如下线性拓扑：规定
$$
\Im(M_\mu \otimes_R N + M \otimes_R N_\nu \longrightarrow M \otimes_R N)
$$
组成一族开子模基本系，其中 $M_\mu \subset M$ 和 $N_\nu \subset N$
分别遍历 $M$ 和 $N$ 的开子模基本系。{\it 完备张量积}
$$
M \widehat{\otimes}_R N =
\lim M \otimes_R N/(M_\mu \otimes_R N + M \otimes_R N_\nu) =
\lim M/M_\mu \otimes_R N/N_\nu
$$
是该张量积的完备化。
\end{definition}

\noindent
注意，$R$ 上的拓扑对张量积或完备张量积的构造并无影响。
若 $R \to A$ 和 $R \to B$ 是线性拓扑环的连续映射，
则上述构造给出张量积 $A \otimes_R B$ 和完备张量积
$A \widehat{\otimes}_R B$。

\medskip\noindent
这里记录注 \ref{formal-spaces-remark-mcquillan} 中引入的概念。

\begin{definition}
\label{formal-spaces-definition-weakly-admissible}
设 $A$ 是线性拓扑环。
\begin{enumerate}
\item 若当 $n \to \infty$ 时 $f^n \to 0$，则称元素 $f \in A$
为{\it 拓扑幂零的}。
\item $A$ 的一个{\it 弱定义理想}是完全由拓扑幂零元素组成的开理想
$I \subset A$。
\item 若 $A$ 有弱定义理想，则称 $A$ {\it 弱预可容许}。
\item 若 $A$ 弱预可容许且完备\footnote{按我们的约定，这包括分离性。}，
则称 $A$ {\it 弱可容许}。
\end{enumerate}
\end{definition}

\noindent
给定线性拓扑环 $A$ 中的弱定义理想 $I$ 和开理想 $J$，
交 $I \cap J$ 是弱定义理想。因此，一旦存在一个弱定义理想，
就存在一族由弱定义理想组成的开理想基本系。特别地，给定弱可容许拓扑环 $A$，有
$A = \lim A/I_\lambda$，其中 $\{I_\lambda\}$ 是一族弱定义理想基本系。

\begin{lemma}
\label{formal-spaces-lemma-weakly-admissible-henselian}
设 $A$ 是弱可容许拓扑环，$I \subset A$ 是弱定义理想。
则 $(A, I)$ 是 Hensel 对。
\end{lemma}

\begin{proof}
设 $A \to A'$ 是 étale 环映射，$\sigma : A' \to A/I$ 是 $A$-代数映射。
由《代数续篇》之引理 \ref{more-algebra-lemma-characterize-henselian-pair}，
只须把 $\sigma$ 提升为 $A$-代数映射 $A' \to A$。
为此，由于 $A$ 完备，只须对每个开理想 $J \subset I$，找到唯一的
$A$-代数映射 $A' \to A/J$ 来提升 $\sigma$。由于 $I$ 是弱定义理想，
理想 $I/J$ 局部幂零。结论遂由《代数续篇》之引理
\ref{more-algebra-lemma-locally-nilpotent-henselian} 得出。
\end{proof}

\begin{lemma}
\label{formal-spaces-lemma-topologically-nilpotent}
设 $B$ 是线性拓扑环。$B$ 中拓扑幂零元素的集合是 $B$ 的闭根式理想。
设 $\varphi : A \to B$ 是线性拓扑环的连续映射。
\begin{enumerate}
\item 若 $f \in A$ 拓扑幂零，则 $\varphi(f)$ 拓扑幂零。
\item 若 $I \subset A$ 完全由拓扑幂零元素组成，则
$\varphi(I)B$ 的闭包完全由拓扑幂零元素组成。
\end{enumerate}
\end{lemma}

\begin{proof}
令 $\mathfrak b \subset B$ 为拓扑幂零元素的集合。
我们略去 $\mathfrak b$ 是根式理想的证明（这是练习定义的好题目）。
设 $g$ 是 $\mathfrak b$ 闭包中的元素。目标是证明 $g$ 拓扑幂零。
设 $J \subset B$ 是开理想。须证明对某个 $e \geq 1$ 有 $g^e \in J$。
由引理 \ref{formal-spaces-lemma-closed}，有 $g \in \mathfrak b + J$。
因此，对某个 $f \in \mathfrak b$ 和 $h \in J$ 有 $g = f + h$。
选取 $m \geq 1$ 使 $f^m \in J$，则如所需有 $g^{m + 1} \in J$。

\medskip\noindent
设 $\varphi : A \to B$ 如引理陈述。(1) 显然，(2) 由此以及
$\mathfrak b$ 是闭理想这一事实推出。
\end{proof}

\begin{lemma}
\label{formal-spaces-lemma-closure-image-ideal}
设 $A \to B$ 是线性拓扑环的连续映射，$I \subset A$ 是理想。
则 $IB$ 的闭包是 $B \to B \widehat{\otimes}_A A/I$ 的核。
\end{lemma}

\begin{proof}
设 $J_\mu$ 是 $B$ 的一族开理想基本系。由引理 \ref{formal-spaces-lemma-closed}，
% P10-E0294
$IB$ 的闭包是 $\bigcap (IB + J_\lambda)$。设 $I_\mu$ 是 $A$
中的一族开理想基本系。则
$$
B \widehat{\otimes}_A A/I =
\lim (B/J_\lambda \otimes_A A/(I_\mu + I)) =
\lim B/(J_\lambda + I_\mu B + I B)
$$
由于 $A \to B$ 连续，对每个 $\lambda$ 都存在 $\mu$ 使
$I_\mu B \subset J_\lambda$，见例 \ref{formal-spaces-example-what-does-it-mean} 中的讨论。
因此，该极限可写成 $\lim B/(J_\lambda + IB)$，结论显然。
\end{proof}

\begin{lemma}
\label{formal-spaces-lemma-completed-tensor-product}
设 $B \to A$ 和 $B \to C$ 是线性拓扑环的连续同态。
\begin{enumerate}
\item 若 $A$ 和 $C$ 弱预可容许，则 $A \widehat{\otimes}_B C$ 弱可容许。
\item 若 $A$ 和 $C$ 预可容许，则 $A \widehat{\otimes}_B C$ 可容许。
\item 若 $A$ 和 $C$ 有可数的开理想基本系，则
$A \widehat{\otimes}_B C$ 有可数的开理想基本系。
\item 若 $A$ 和 $C$ 是预-adic 的，并有有限生成的定义理想，
则 $A \widehat{\otimes}_B C$ 是 adic 的，并有有限生成的定义理想。
\item 若 $A$ 和 $C$ 是预-adic Noether 环，并且
$B/\mathfrak b \to A/\mathfrak a$ 是有限型的，其中
$\mathfrak a \subset A$ 和 $\mathfrak b \subset B$ 是拓扑幂零元素的理想，
则 $A \widehat{\otimes}_B C$ 是 adic Noether 环。
\end{enumerate}
\end{lemma}

\begin{proof}
设 $I_\lambda \subset A$（$\lambda \in \Lambda$）和
$J_\mu \subset C$（$\mu \in M$）分别是开理想基本系。按定义
$$
A \widehat{\otimes}_B C =
\lim_{\lambda, \mu} A/I_\lambda \otimes_B C/J_\mu
$$
并带极限拓扑。因此，一族开理想基本系由映射
$A \widehat{\otimes}_B C \to A/I_\lambda \otimes_B C/J_\mu$
的核 $K_{\lambda, \mu}$ 给出。注意，$K_{\lambda, \mu}$ 是理想
$I_\lambda(A \widehat{\otimes}_B C) +
J_\mu(A \widehat{\otimes}_B C)$ 的闭包。最后，有一个像稠密的环同态
$\tau : A \otimes_B C \to A \widehat{\otimes}_B C$。

\medskip\noindent
(1) 的证明。若 $I_\lambda$ 和 $J_\mu$ 完全由拓扑幂零元素组成，
则由引理 \ref{formal-spaces-lemma-topologically-nilpotent}，$K_{\lambda, \mu}$ 也如此。
所以按定义，$A \widehat{\otimes}_B C$ 弱可容许。

\medskip\noindent
(2) 的证明。假设对某个 $\lambda_0$ 和 $\mu_0$，理想
$I = I_{\lambda_0} \subset A$ 和 $J_{\mu_0} \subset C$ 是定义理想。
所以，对每个 $\lambda$，存在 $n$ 使 $I^n \subset I_\lambda$；
对每个 $\mu$，存在 $m$ 使 $J^m \subset J_\mu$。于是
$$
\left(I(A \widehat{\otimes}_B C) + J(A \widehat{\otimes}_B C)\right)^{n + m}
\subset
I_\lambda(A \widehat{\otimes}_B C) + J_\mu(A \widehat{\otimes}_B C)
$$
由此，开理想 $K = K_{\lambda_0, \mu_0}$ 满足
$K^{n + m} \subset K_{\lambda, \mu}$。所以 $K$ 是
$A \widehat{\otimes}_B C$ 的定义理想，而 $A \widehat{\otimes}_B C$ 按定义可容许。

\medskip\noindent
(3) 的证明。若 $\Lambda$ 和 $M$ 可数，则 $\Lambda \times M$ 也可数。

\medskip\noindent
(4) 的证明。假设 $\Lambda = \mathbf{N}$、$M = \mathbf{N}$，
并且有有限生成理想 $I \subset A$ 和 $J \subset C$，使
$I_n = I^n$ 且 $J_n = J^n$。则
$$
I(A \widehat{\otimes}_B C) + J(A \widehat{\otimes}_B C)
$$
是有限生成理想，并且容易看出 $A \widehat{\otimes}_B C$ 是
$A \otimes_B C$ 关于该理想的完备化。因此，(4) 由《代数》之引理
\ref{algebra-lemma-hathat-finitely-generated} 得出。

\medskip\noindent
(5) 的证明。令 $\mathfrak c \subset C$ 为拓扑幂零元素的理想。
由于 $A$ 和 $C$ 是 adic Noether 环，可知 $\mathfrak a$ 和
$\mathfrak c$ 是定义理想（细节略）。由 (4) 已知
$A \widehat{\otimes}_B C$ 是 adic 的，并且
$\mathfrak a(A \widehat{\otimes}_B C) +
\mathfrak c(A \widehat{\otimes}_B C)$ 是有限生成的定义理想。由于
$$
A \widehat{\otimes}_B C /
\left(\mathfrak a(A \widehat{\otimes}_B C) +
\mathfrak c(A \widehat{\otimes}_B C)\right)
=
A/\mathfrak a \otimes_{B/\mathfrak b} C/\mathfrak c
$$
是 Noether 环 $C/\mathfrak c$ 上的有限型代数，结论由《代数》之引理
\ref{algebra-lemma-completion-Noetherian} 得出。
\end{proof}









\section{taut 环映射}
\label{formal-spaces-section-taut-ring-maps}

\noindent
为线性拓扑环之间的连续映射所具有的下述性质取一个名称是方便的。

\begin{definition}
\label{formal-spaces-definition-taut}
设 $\varphi : A \to B$ 是线性拓扑环的连续映射。若对每个开理想
$I \subset A$，理想 $\varphi(I)B$ 的闭包都是开的，并且这些闭包组成一族开理想基本系，
则称 $\varphi$ 是 {\it taut} 的\footnote{这是非标准记号。该定义可推广到模：
若对每个开理想 $I \subset A$，$IM$ 在 $M$ 中的闭包都是开的，
且这些闭包组成 $M$ 中 $0$ 的一族邻域基本系，则称线性拓扑化的 $A$-模
$M$ 是 $A$-taut 的。}。
\end{definition}

\noindent
若 $\varphi : A \to B$ 是线性拓扑环的连续映射，而 $I_\lambda$
是 $A$ 的一族开理想基本系，则 $\varphi$ 是 taut 的，当且仅当
% P10-E0295
$I_\lambda B$ 的闭包都是开的，并组成 $A$ 中的一族开理想基本系。

\begin{lemma}
\label{formal-spaces-lemma-taut-weakly-admissible}
设 $\varphi : A \to B$ 是弱可容许拓扑环的连续映射。下列条件等价：
\begin{enumerate}
\item $\varphi$ 是 taut 的；
\item 对 $A$ 的每个弱定义理想 $I \subset A$，$\varphi(I)B$ 的闭包
都是 $B$ 的弱定义理想，并且这些闭包组成 $B$ 的一族弱定义理想基本系。
\end{enumerate}
\end{lemma}

\begin{proof}
定义 \ref{formal-spaces-definition-taut} 后的说明表明 (2) 推出 (1)。反过来，
假设 $\varphi$ 是 taut 的。若 $I \subset A$ 是弱定义理想，则按 taut 性的定义，
$\varphi(I)B$ 的闭包是开的；由引理 \ref{formal-spaces-lemma-topologically-nilpotent}，
它完全由拓扑幂零元素组成。因此 $\varphi(I)B$ 的闭包是弱定义理想。
此外，按 taut 性的定义，这些理想组成一族开理想基本系，所以 (2) 成立。
\end{proof}

\begin{lemma}
\label{formal-spaces-lemma-completion-taut}
设 $A$ 是线性拓扑环。从 $A$ 到其完备化的映射
$A \to A^\wedge$ 是 taut 的。
\end{lemma}

\begin{proof}
设 $I_\lambda$ 是 $A$ 的一族开理想基本系。回忆
$A^\wedge = \lim A/I_\lambda$ 带极限拓扑；这表示各核
$J_\lambda = \Ker(A^\wedge \to A/I_\lambda)$
组成 $A^\wedge$ 的一族开理想基本系。由于 $J_\lambda$ 是
$I_\lambda A^\wedge$ 的闭包（比较引理 \ref{formal-spaces-lemma-closure-image-ideal}），结论成立。
\end{proof}

\begin{lemma}
\label{formal-spaces-lemma-composition-taut}
设 $A \to B$ 和 $B \to C$ 是线性拓扑环的连续同态。
若 $A \to B$ 和 $B \to C$ 都是 taut 的，则 $A \to C$ 是 taut 的。
\end{lemma}

\begin{proof}
略。提示：若 $I \subset A$ 是理想，$J$ 是 $IB$ 的闭包，
则 $JC$ 的闭包等于 $IC$ 的闭包。
\end{proof}

\begin{lemma}
\label{formal-spaces-lemma-permanence-taut}
设 $A \to B$ 和 $B \to C$ 是线性拓扑环的连续同态。
若 $A \to C$ 是 taut 的，则 $B \to C$ 是 taut 的。
\end{lemma}

\begin{proof}
设 $J \subset B$ 是开理想，其逆像为 $I \subset A$。
则 $JC$ 的闭包包含 $IC$ 的闭包。由于 $A \to C$ 是 taut 的，这个闭包是开的。
设 $I_\lambda$ 是 $A$ 的一族开理想基本系，$K_\lambda$ 是
$I_\lambda C$ 的闭包。由于 $A \to C$ 是 taut 的，它们组成
$C$ 的一族开理想基本系。以 $J_\lambda \subset B$ 表示 $K_\lambda$ 的逆像。
则 $J_\lambda C$ 的闭包是 $K_\lambda$。因此可见，当 $J$ 遍历 $B$
的开理想时，理想 $JC$ 的闭包组成 $C$ 的一族开理想基本系。
\end{proof}

\begin{lemma}
\label{formal-spaces-lemma-base-change-taut}
设 $A \to B$ 和 $A \to C$ 是线性拓扑环的连续同态。
若 $A \to B$ 是 taut 的，则
$C \to B \widehat{\otimes}_A C$ 是 taut 的。
\end{lemma}

\begin{proof}
设 $K \subset C$ 是开理想。选取任意开理想 $I \subset A$，使其在 $C$ 中的像包含在
% P10-E0296
$J$ 中。按假设，$IB$ 的闭包 $J$ 是开的。由于 $A \to B$ 是 taut 的，
可见 $B \widehat{\otimes}_A C$ 是环
$B/J \otimes_{A/I} C/K$ 对所有 $K$ 和 $I$ 的选择所取的极限；也就是说，理想
$J(B \widehat{\otimes}_A C) + K(B \widehat{\otimes}_A C)$
组成一族开理想基本系。又因为
$B \to B \widehat{\otimes}_A C$ 连续，并且 $I$ 映入 $K$，
可知 $J$ 映入 $K(B \widehat{\otimes}_A C)$ 的闭包。
因此该闭包等于
$J(B \widehat{\otimes}_A C) + K(B \widehat{\otimes}_A C)$，证明完成。
\end{proof}

\begin{lemma}
\label{formal-spaces-lemma-taut-ascent-countable}
设 $\varphi : A \to B$ 是线性拓扑环的连续同态。
若 $\varphi$ 是 taut 的，且 $A$ 有可数的开理想基本系，
则 $B$ 有可数的开理想基本系。
\end{lemma}

\begin{proof}
由定义立即得到。
\end{proof}

\begin{lemma}
\label{formal-spaces-lemma-taut-ascent-weakly-admissible}
设 $\varphi : A \to B$ 是线性拓扑环的连续同态。
若 $\varphi$ 是 taut 的，且 $A$ 弱预可容许，则 $B$ 弱预可容许。
\end{lemma}

\begin{proof}
设 $I \subset A$ 是弱定义理想。则 $IB$ 的闭包 $J$ 是开的，
并且由引理 \ref{formal-spaces-lemma-topologically-nilpotent}，它完全由拓扑幂零元素组成。
因此 $J$ 是 $B$ 的弱定义理想。
\end{proof}

\begin{lemma}
\label{formal-spaces-lemma-taut-ascent-admissible}
设 $\varphi : A \to B$ 是线性拓扑环的连续同态。
若 $\varphi$ 是 taut 的，且 $A$ 预可容许，则 $B$ 预可容许。
\end{lemma}

\begin{proof}
设 $I \subset A$ 是定义理想，$I_\lambda \subset A$ 是一族开理想基本系。
则 $IB$ 的闭包 $J$ 是开的，而 $I_\lambda B$ 的闭包 $J_\lambda$
也是开的，并组成 $B$ 的一族开理想基本系。对每个 $\lambda$，
都存在 $n$ 使 $I^n \subset I_\lambda$。注意，$J^n$ 包含在 $I^nB$ 的闭包中。
所以 $J^n \subset J_\lambda$，从而 $J$ 是定义理想。
\end{proof}

\begin{lemma}
\label{formal-spaces-lemma-dense-image-surjective}
设 $\varphi : A \to B$ 是线性拓扑环的连续同态。假设
\begin{enumerate}
\item $\varphi$ 是 taut 的，且像稠密；
\item $A$ 完备，并有可数的开理想基本系；
\item $B$ 分离。
\end{enumerate}
则 $\varphi$ 是满的开映射，$B$ 完备，并且对某个闭理想
$K \subset A$ 有 $B = A/K$。
\end{lemma}

\begin{proof}
由开映射引理（《代数续篇》之引理
\ref{more-algebra-lemma-open-mapping}）和 $\varphi$ 的 taut 性，
可知映射 $\varphi$ 是开的。由于 $\varphi$ 的像稠密，$\varphi$ 是满射。
$\varphi$ 的核 $K$ 是闭的，因为 $\varphi$ 连续。因此
$B = A/K$ 完备，例如可见引理 \ref{formal-spaces-lemma-quotient-by-closed}。
\end{proof}







\section{Adic 环映射}
\label{formal-spaces-section-adic-ring-maps}

\noindent
作如下定义。

\begin{definition}
\label{formal-spaces-definition-adic-homomorphism}
设 $A$ 和 $B$ 是预-adic 拓扑环。若存在 $A$ 的定义理想
$I \subset A$，使得 $B$ 上的拓扑为 $I$-进拓扑，则称环同态
$\varphi : A \to B$ 是 {\it adic} 的\footnote{这可能是非标准术语。}。
\end{definition}

\noindent
若 $\varphi : A \to B$ 是预-adic 环的 adic 同态，则 $\varphi$ 连续，
并且对 $A$ 的每个定义理想 $I$，$B$ 上的拓扑都是 $I$-进拓扑。

\begin{lemma}
\label{formal-spaces-lemma-composition-adic}
设 $A \to B$ 和 $B \to C$ 是预-adic 环的连续同态。
若 $A \to B$ 和 $B \to C$ 都是 adic 的，则 $A \to C$ 是 adic 的。
\end{lemma}

\begin{proof}
略。
\end{proof}

\begin{lemma}
\label{formal-spaces-lemma-permanence-adic}
设 $A \to B$ 和 $B \to C$ 是预-adic 环的连续同态。
若 $A \to C$ 是 adic 的，则 $B \to C$ 是 adic 的。
\end{lemma}

\begin{proof}
选取 $A$ 的一个定义理想 $I$。因为 $A \to C$ 是 adic 的，
$IC$ 是 $C$ 的定义理想。因为 $B \to C$ 连续，可以找到映入 $IC$
的定义理想 $J \subset B$。因为 $A \to B$ 连续，$J$ 在 $I$ 中的逆像
$I' \subset I$ 也是 $A$ 的定义理想。所以
$I'C \subset JC \subset IC$ 夹在两个定义理想之间，本身也是定义理想。
\end{proof}

\begin{lemma}
\label{formal-spaces-lemma-adic-taut}
设 $\varphi : A \to B$ 是预-adic 拓扑环之间的连续同态。
若 $\varphi$ 是 adic 的，则 $\varphi$ 是 taut 的。
\end{lemma}

\begin{proof}
由定义立即得到。
\end{proof}

\noindent
下一个引理说明两件事：
\begin{enumerate}
\item adic 性沿完备线性拓扑环的 taut 映射上升；
\item 对 adic 环之间的连续映射 $\varphi : A \to B$，
若 $A$ 有有限生成的定义理想，则“$\varphi$ 是 taut 的”与
“$\varphi$ 是 adic 的”等价。
\end{enumerate}
由于 (2)，可以说“taut 性”把“adic 性”推广到了任意线性拓扑环之间的连续环映射。
另见第 \ref{formal-spaces-section-adic} 节。

\begin{lemma}
\label{formal-spaces-lemma-taut-is-adic}
设 $\varphi : A \to B$ 是线性拓扑环的连续映射。
若 $\varphi$ 是 taut 的，$A$ 是预-adic 的且有有限生成的定义理想，
而 $B$ 完备，则 $B$ 是 adic 的且有有限生成的定义理想，并且环映射
$\varphi$ 是 adic 的。
\end{lemma}

\begin{proof}
选取 $A$ 的一个有限生成定义理想 $I$。令 $J_n$ 为
$\varphi(I^n)B$ 在 $B$ 中的闭包。由于 $B$ 完备，有
$B = \lim B/J_n$。令 $B' = \lim B/I^nB$ 为 $B$ 的 $I$-进完备化。
由《代数》之引理 \ref{algebra-lemma-hathat-finitely-generated}，
$B'$ 上的 $I$-进拓扑完备，并且 $B'/I^nB' = B/I^nB$。
所以环映射 $B' \to B$ 连续，并且像稠密，因为对所有 $n$，
$B' \to B/I^nB \to B/J_n$ 都是满射。最后，映射 $B' \to B$ 是 taut 的，
因为 $(I^nB')B = I^nB$，且 $A \to B$ 是 taut 的。
由引理 \ref{formal-spaces-lemma-dense-image-surjective}，$B' \to B$ 是满的开映射。
因此 $B$ 上的拓扑是 $I$-进拓扑，证明完成。
\end{proof}








\section{弱 adic 环}
\label{formal-spaces-section-weakly-adic}

\noindent
建议读者跳过本节。下面是 adic 环的一个自然推广。

\begin{definition}
\label{formal-spaces-definition-weakly-adic}
\begin{reference}
\cite[Definition 8.3.8]{Gabber-Ramero}
\end{reference}
设 $A$ 是线性拓扑环。
\begin{enumerate}
\item 若存在理想 $I \subset A$，使得对所有 $n \geq 0$，$I^n$ 的闭包都是开的，
且这些闭包组成一族开理想基本系，则称 $A$ {\it 弱预-adic}
\footnote{在 \cite{Gabber-Ramero} 中，作者称 $A$ 是 {\it $c$-adic} 的。}。
\item 若 $A$ 弱预-adic 且完备\footnote{按我们的约定，这包括分离性。}，
则称 $A$ {\it 弱 adic}。
\end{enumerate}
\end{definition}

\noindent
对完备线性拓扑环，有下列蕴含关系
$$
\xymatrix{
\text{adic 且 Noether} \ar@{=>}[d] \\
\text{adic 且有有限生成的定义理想} \ar@{=>}[d] \\
\text{adic} \ar@{=>}[d] \\
\text{弱 adic} \ar@{=>}[d] \\
\text{可容许且第一可数} \ar@{=>}[d] \ar@{=>}[r] &
\text{可容许} \ar@{=>}[d] \\
\text{弱可容许且第一可数} \ar@{=>}[r] &
\text{弱可容许}
}
$$
这里“第一可数”表示拓扑环有可数的开理想基本系。
对不完备的线性拓扑环也有类似的蕴含图表
（即使用预-adic、弱预-adic、预可容许和弱预可容许这些概念）。
与预-adic 环的情形相反，弱预-adic 环的完备化是弱 adic 的；
下面刻画弱预-adic 环的引理说明了这一点。

\begin{lemma}
\label{formal-spaces-lemma-weakly-adic}
设 $A$ 是线性拓扑环。下列条件等价：
\begin{enumerate}
\item $A$ 弱预-adic；
\item 存在 taut 连续环映射 $A' \to A$，其中 $A'$ 是预-adic 拓扑环；
\item $A$ 预可容许，且存在一个定义理想 $I$，使得对所有 $n \geq 1$，
$I^n$ 的闭包都是开的；
\item $A$ 预可容许，并且对每个定义理想 $I$ 以及所有 $n \geq 1$，
$I^n$ 的闭包都是开的。
\end{enumerate}
弱预-adic 环的完备化是弱 adic 的。若 $A$ 弱 adic，则 $A$ 可容许，
并有可数的开理想基本系。
\end{lemma}

\begin{proof}
假设 (1)。选取理想 $I$，使得对所有 $n$，$I^n$ 的闭包都是开的，
且这些闭包组成一族开理想基本系。以 $A'$ 表示赋予 $I$-进拓扑的环 $A$。
则 $A' \to A$ 按定义是 taut 的，故 (2) 成立。

\medskip\noindent
假设 (2)。设 $I' \subset A'$ 是定义理想，以 $I$ 表示 $I'A$ 的闭包。
$A' \to A$ 的 taut 性表示，$(I')^nA$ 的闭包 $I_n$ 都是开的，
并组成一族开理想基本系。因此 $I = I_1$ 是开的，而 $I^n$ 的闭包等于
$I_n$，从而是开的并组成一族开理想基本系。特别地，$I$ 是一个定义理想，
且所有 $I^n$ 的闭包都是开的。所以 (3) 成立。

\medskip\noindent
若 $I \subset A$ 如 (3) 所述，则 $I$ 是定义 \ref{formal-spaces-definition-weakly-adic}
中的理想，从而 (1) 成立。又若 $I' \subset A$ 是任意其他定义理想，
则 $I'$ 是开的（见《代数续篇》之定义
\ref{more-algebra-definition-topological-ring}），故它包含某个 $I^n$，其中
$n \geq 1$。因此对所有 $m \geq 1$，$(I')^m$ 包含 $I^{nm}$，
从而 $(I')^m$ 的闭包对所有 $m$ 都是开的。这样便看出 (3) 推出 (4)。
(4) $\Rightarrow$ (3) 是平凡的。

\medskip\noindent
设 $A$ 弱预-adic。选取如 (2) 中的 $A' \to A$。
由引理 \ref{formal-spaces-lemma-completion-taut} 和 \ref{formal-spaces-lemma-composition-taut}，
复合 $A' \to A^\wedge$ 是 taut 的。因此，由 (2) 与 (1) 的等价性，
$A^\wedge$ 弱预-adic。因为线性拓扑环 $A$ 的完备化是完备的
（《代数续篇》之第 \ref{more-algebra-section-topological-ring} 节），
可见 $A^\wedge$ 弱 adic。

\medskip\noindent
设 $A$ 弱 adic。则 $A$ 完备，并且
% P10-E0297
由 (1) $\Rightarrow$ (3)，$A$ 且预可容许，所以 $A$ 可容许。
当然，按定义 $A$ 有可数的开理想基本系。
\end{proof}

\noindent
给出两个判据，保证弱 adic 环是 adic 的，并有有限生成的定义理想。

\begin{lemma}
\label{formal-spaces-lemma-weakly-adic-finite-generation-bis}
设 $A$ 是完备线性拓扑环。设 $I \subset A$ 是有限生成理想，
并且对所有 $n \geq 0$，$I^n$ 的闭包都是开的，且这些闭包组成一族开理想基本系。
则 $A$ 是 adic 的，并有有限生成的定义理想。
\end{lemma}

\begin{proof}
以 $A'$ 表示赋予 $I$-进拓扑的环 $A$。这些假设说明 $A' \to A$ 是 taut 的。
结论由引理 \ref{formal-spaces-lemma-taut-is-adic} 得出（确切地说，该引理还告诉我们 $I$ 是定义理想）。
\end{proof}

\begin{lemma}
\label{formal-spaces-lemma-weakly-adic-finite-generation}
设 $A$ 是弱 adic 拓扑环。设 $I$ 是定义理想，并且 $I/I_2$ 是有限生成模，
其中 $I_2$ 是 $I^2$ 的闭包。则 $A$ 是 adic 的，并有有限生成的定义理想。
\end{lemma}

\begin{proof}
我们使用引理 \ref{formal-spaces-lemma-weakly-adic} 的刻画而不再说明。
选取 $f_1, \ldots, f_r \in I$，其像生成 $I/I_2$。令
$I' = (f_1, \ldots, f_r)$，则 $I' + I_2 = I$。
此时 $I_2$ 是
$I^2 = (I' + I_2)^2 \subset I' + I_3$ 的闭包，其中 $I_3$ 是 $I^3$ 的闭包。
所以 $I' + I_3 = I$。如此继续，可见对所有 $n \geq 2$，
$I' + I_n = I$，其中 $I_n$ 是 $I^n$ 的闭包。换言之，$I'$ 在 $A$ 中的闭包为 $I$。
所以 $(I')^n$ 的闭包为 $I_n$。因此，$(I')^n$ 的闭包组成 $A$ 的一族开理想基本系。
结论由引理 \ref{formal-spaces-lemma-weakly-adic-finite-generation-bis} 得出。
\end{proof}

\noindent
“弱预-adic”性质的一个关键特征是：它沿线性拓扑环的 taut 环同态上升。

\begin{lemma}
\label{formal-spaces-lemma-taut-ascent-weakly-adic}
设 $\varphi : A \to B$ 是线性拓扑环的连续同态。
若 $\varphi$ 是 taut 的，且 $A$ 弱预-adic，则 $B$ 弱预-adic。
\end{lemma}

\begin{proof}
设 $I \subset A$ 是理想，使 $I^n$ 的闭包 $I_n$ 都是开的，
且这些闭包定义一族开理想基本系。则 $I^nB$ 的闭包等于 $I_nB$ 的闭包。
由于 $\varphi$ 是 taut 的，这些闭包都是开的，并组成 $B$ 的一族开理想基本系。
故 $B$ 弱预-adic。
\end{proof}

\begin{lemma}
\label{formal-spaces-lemma-completed-tensor-product-weakly-adic}
设 $B \to A$ 和 $B \to C$ 是线性拓扑环的连续同态。
若 $A$ 和 $C$ 弱预-adic，则 $A \widehat{\otimes}_B C$ 弱 adic。
\end{lemma}

\begin{proof}
我们使用引理 \ref{formal-spaces-lemma-weakly-adic} 的刻画而不再说明。
由引理 \ref{formal-spaces-lemma-completed-tensor-product}，已知
$A \widehat{\otimes}_B C$ 可容许。此外，该引理的证明表明：
% P10-E0298
当 $I \subset A$ 和 $J \subset C$ 是定义理想时，
$I(A \widehat{\otimes}_B C) + J(A \widehat{\otimes}_B C)$ 的闭包
$K \subset A \widehat{\otimes}_B C$ 是定义理想。
于是只须证明，对所有 $n \geq 1$，$K^n$ 的闭包都是开的。
由于理想 $K^n$ 包含
$I^n(A \widehat{\otimes}_B C) + J^n(A \widehat{\otimes}_B C)$，
而 $I^n$ 在 $A$ 中的闭包和 $J^n$ 在 $C$ 中的闭包都是开的，
可见 $K^n$ 在 $A \widehat{\otimes}_B C$ 中的闭包是开的。
\end{proof}






\section{性质的下降}
\label{formal-spaces-section-taut-descent}

\noindent
本节考虑如下情形：
\begin{enumerate}
\item $\varphi : A \to B$ 是线性拓扑化拓扑环的连续映射；
\item $\varphi$ 是 taut 的；
\item 对每个开理想 $I \subset A$，若 $J \subset B$ 表示 $IB$ 的闭包，
则映射 $A/I \to B/J$ 忠实平坦。
\end{enumerate}
我们将证明，在此情形下，$B$ 的性质由 $A$ 继承。

\begin{lemma}
\label{formal-spaces-lemma-taut-descent-countable}
在上述情形中，若 $B$ 有可数的开理想基本系，则 $A$ 有可数的开理想基本系。
\end{lemma}

\begin{proof}
选取一族开理想基本系
$B \supset J_1 \supset J_2 \supset \ldots$。由 $\varphi$ 的 taut 性，
对每个 $n$，可以找到开理想 $I_n$，使 $J_n \supset I_nB$。
我们断言 $I_n$ 是 $A$ 的一族开理想基本系。事实上，设 $I \subset A$ 开。
由于 $\varphi$ 是 taut 的，$IB$ 的闭包是开的，因此对充分大的某个 $n$ 包含 $J_n$。
所以 $I_nB \subset IB$。令 $J$ 为 $IB$ 在 $B$ 中的闭包。
因为 $A/I \to B/J$ 忠实平坦，所以它是单射。又由于
$I_nB \subset IB \subset J$，映射 $I_n \to A/I \to B/J$ 为零，
故 $I_n \to A/I$ 为零。因此 $I_n \subset I$，结论成立。
\end{proof}

\begin{lemma}
\label{formal-spaces-lemma-taut-descent-weakly-admissible}
在上述情形中，若 $B$ 弱预可容许，则 $A$ 弱预可容许。
\end{lemma}

\begin{proof}
设 $J \subset B$ 是弱定义理想。取开理想 $I \subset A$ 使 $IB \subset J$。
为证明 $I$ 是弱定义理想，须证明每个 $f \in I$ 都拓扑幂零。
设 $I' \subset A$ 是开理想，以 $J' \subset B$ 表示 $I'B$ 的闭包。
则 $A/I' \to B/J'$ 忠实平坦，从而是单射。因此，为证明 $f^n \in I'$，
只须证明 $\varphi(f)^n \in J'$。当 $n \gg 0$ 时这成立，因为
$\varphi(f) \in J$，理想 $J$ 是 $B$ 的弱定义理想，而 $J'$ 在 $B$ 中开。
\end{proof}

\begin{lemma}
\label{formal-spaces-lemma-taut-descent-admissible}
在上述情形中，若 $B$ 预可容许，则 $A$ 预可容许。
\end{lemma}

\begin{proof}
% P10-E0299
设 $J \subset B$ 是弱定义理想。取开理想 $I \subset A$ 使 $IB \subset J$。
设 $I' \subset A$ 是开理想。为证明 $I$ 是定义理想，须证明当 $n \gg 0$ 时
$I^n \subset I'$。以 $J' \subset B$ 表示 $I'B$ 的闭包。
则 $A/I' \to B/J'$ 忠实平坦，从而是单射。因此，为证明
$I^n \subset I'$，只须证明 $\varphi(I)^n \subset J'$。
当 $n \gg 0$ 时这成立，因为 $\varphi(I) \subset J$，理想 $J$ 是 $B$ 的定义理想，
而 $J'$ 在 $B$ 中开。
\end{proof}

\begin{lemma}
\label{formal-spaces-lemma-taut-descent-weakly-adic}
在上述情形中，若 $B$ 弱预-adic，则 $A$ 弱预-adic。
\end{lemma}

\begin{proof}
我们使用引理 \ref{formal-spaces-lemma-weakly-adic} 给出的弱预-adic 环刻画而不再说明。
由引理 \ref{formal-spaces-lemma-taut-descent-admissible}，拓扑环 $A$ 预可容许。
设 $I \subset A$ 是定义理想。固定 $n \geq 1$。为证明引理，须证明
$I^n$ 的闭包是开的。设 $I_\lambda \subset A$ 是一族开理想基本系。
以 $J \subset B$（分别地，$J_\lambda \subset B$）表示
$IB$（分别地，$I_\lambda B$）的闭包。由于 $B$ 弱预-adic，$J^n$ 的闭包是开的。
所以存在 $\lambda$，使得
$$
J_\lambda \subset \bigcap\nolimits_\mu (J^n + J_\mu)
$$
因为由引理 \ref{formal-spaces-lemma-closed}，右侧是 $J^n$ 的闭包。
这表示 $J_\lambda$ 在 $B/J_\mu$ 中的像包含在 $J^n$ 在 $B/J_\mu$ 中的像内。
注意，$J^n$ 在 $B/J_\mu$ 中的像等于 $I^nB$ 在 $B/J_\mu$ 中的像
（因为 $J$ 的每个元素模 $J_\mu$ 都同余于 $IB$ 的一个元素）。
由于 $A/I_\mu \to B/J_\mu$ 忠实平坦，并且
$I_\lambda B \subset J_\lambda$，可知 $I_\lambda$ 在 $A/I_\mu$ 中的像
包含在 $I^n$ 的像内。因此 $I_\lambda$ 包含在 $I^n$ 的闭包中，证明完成。
\end{proof}

\begin{lemma}
\label{formal-spaces-lemma-taut-descent-adic-star}
在上述情形中，若 $B$ 是 adic 的且有有限生成的定义理想，而 $A$ 完备，
则 $A$ 是 adic 的且有有限生成的定义理想。
\end{lemma}

\begin{proof}
由引理 \ref{formal-spaces-lemma-taut-descent-weakly-adic}，已知 $A$ 弱 adic，
从而当然可容许（并用引理 \ref{formal-spaces-lemma-weakly-adic} 看出 adic 环是弱 adic 的）。
设 $I \subset A$ 是定义理想，$J \subset B$ 是有限生成的定义理想。
由于 $IB$ 的闭包是开的，可以找到 $n > 0$，使 $J^n$ 包含在 $IB$ 的闭包中。
所以用 $J^n$ 替换 $J$ 后，可以假设 $J$ 是包含在 $IB$ 闭包中的有限生成定义理想。
由引理 \ref{formal-spaces-lemma-closed}，这当然蕴含
$$
J \subset IB + J^2
$$
考虑有限生成 $A$-模 $M = (J + IB)/IB$。所显示的方程表明 $JM = M$。
由引理 \ref{formal-spaces-lemma-weakly-admissible-henselian}（例如）可见，$J$ 包含在
$B$ 的 Jacobson 根中。所以由 Nakayama 引理，更确切地说由《代数》之引理
\ref{algebra-lemma-NAK} 的 (2)，得到 $M = 0$。因此 $J \subset IB$。

\medskip\noindent
因为 $J$ 有限生成，可以找到有限生成理想 $I' \subset I$，使 $J \subset I'B$。
因为 $A \to B$ 连续、$J \subset B$ 开，并且 $I$ 是定义理想，
可以找到 $n > 0$ 使 $I^nB \subset J$。令 $J_{n + 1} \subset B$
为 $I^{n + 1}B$ 的闭包。我们有
$$
I^n \cdot (B/J_{n + 1}) \subset J \cdot (B/J_{n + 1}) \subset
I' \cdot (B/J_{n + 1})
$$
由于 $A/I^{n + 1} \to B/J_{n + 1}$ 忠实平坦，这蕴含
$I^n \cdot (A/I^{n + 1}) \subset I' \cdot (A/I^{n + 1})$，
进而意味着
$$
I^n \subset I' + I^{n + 1}
$$
这蕴含对所有 $k \geq 1$ 有 $I^n \subset I' + I^{n + k}$，
进而蕴含对所有 $k,m \geq 1$ 有
$I^{nm} \subset (I')^m + I^{nm + k}$。这说明 $(I')^m$ 的闭包包含 $I^{nm}$。
由于 $A$ 弱 adic，$I^{nm}$ 的闭包是开的，故
% P10-E0300
对所有 $m$，$(I')^m$ 的闭包都是开的。由于这些闭包组成 $A$ 的一族开理想基本系
（因为 $I^n$ 的闭包也具有此性质），结论由引理
\ref{formal-spaces-lemma-weakly-adic-finite-generation-bis} 得出。
\end{proof}








\section{仿射形式代数空间}
\label{formal-spaces-section-affine-formal-algebraic-spaces}

\noindent
本节引入仿射形式代数空间。它们事实上与 \cite{BVGD} 中所谓的仿射形式概形相同。
不过，为避免与 \cite{EGA} 所定义的仿射形式概形概念混淆，
我们将把它们称为仿射形式代数空间。

\medskip\noindent
回忆，概形的增厚是一个在底层拓扑空间上诱导满射的闭浸入，
见《态射续篇》之定义 \ref{more-morphisms-definition-thickening}。

\begin{definition}
\label{formal-spaces-definition-affine-formal-algebraic-space}
设 $S$ 是概形。若存在
\begin{enumerate}
\item 一个有向集 $\Lambda$；
\item $(\Sch/S)_{fppf}$ 中位于 $\Lambda$ 上的系统
$(X_\lambda, f_{\lambda \mu})$，其中
\begin{enumerate}
\item 每个 $X_\lambda$ 都仿射；
\item 每个 $f_{\lambda \mu} : X_\lambda \to X_\mu$ 都是增厚，
\end{enumerate}
\end{enumerate}
使得作为 fppf 层有
$$
X \cong \colim_{\lambda \in \Lambda} X_\lambda
$$
并且 $X$ 满足一个集合论条件（见注 \ref{formal-spaces-remark-set-theoretic}），
则称 $(\Sch/S)_{fppf}$ 上的层 $X$ 是{\it 仿射形式代数空间}。
$S$ 上仿射形式代数空间的{\it 态射}是层的映射。
\end{definition}

\noindent
注意，系统 $(X_\lambda, f_{\lambda \mu})$ 不是数据的一部分。
设 $U$ 是 $S$ 上的拟紧概形。由于转移映射都是单态射，由《站点》之引理
\ref{sites-lemma-directed-colimits-sections} 可见
$$
X(U) = \colim X_\lambda(U)
$$
因此，定义中的余极限所内含的 fppf 层化其实是 Zariski 层化，
而它对拟紧概形不起任何作用。

\begin{lemma}
\label{formal-spaces-lemma-diagonal-affine-formal-algebraic-space}
设 $S$ 是概形。若 $X$ 是 $S$ 上的仿射形式代数空间，则对角态射
$\Delta : X \to X \times_S X$ 可表示且为闭浸入。
\end{lemma}

\begin{proof}
设给定 $U \to X$ 和 $V \to X$，其中 $U,V$ 是 $S$ 上的概形。
来证明 $U \times_X V$ 可表示。按定义
\ref{formal-spaces-definition-affine-formal-algebraic-space} 写
$X = \colim X_\lambda$。上述讨论表明，在 $U$ 和 $V$ 上 Zariski 局部地，
% P10-E0301
这些态射分解经过某个 $X_\lambda$。此时
$U \times_X V = U \times_{X_\lambda} V$，它是概形。
所以由《空间》之引理 \ref{spaces-lemma-representable-diagonal}，对角可表示。
给定 $(a,b) : W \to X \times_S X$，其中 $W$ 是 $S$ 上的概形，考虑映射
$X \times_{\Delta, X \times_S X, (a, b)} W \to W$。
与前面一样，在 $W$ 上局部地，态射 $a$ 和 $b$ 映入某个 $\lambda$ 对应的仿射概形
$X_\lambda$，于是得到态射
$X_\lambda
\times_{\Delta_\lambda, X_\lambda \times_S X_\lambda, (a, b)} W \to W$。
这是 $\Delta_\lambda : X_\lambda \to X_\lambda \times_S X_\lambda$
的基变换，而后者是闭浸入，因为 $X_\lambda \to S$ 分离
（$X_\lambda$ 仿射）。所以 $X \to X \times_S X$ 是闭浸入。
\end{proof}

\noindent
概形态射 $X \to X'$ 是增厚，当且仅当它是闭浸入并在底层点集上诱导满射，
见（《态射续篇》之定义 \ref{more-morphisms-definition-thickening}）。
因此，增厚这一性质在任意基变换下保持，并且在目标上是 fpqc 局部的，
见《空间》之第 \ref{spaces-section-lists} 节。所以《空间》之定义
\ref{spaces-definition-relative-representable-property} 适用于“增厚”，
而我们知道 $(\Sch/S)_{fppf}$ 上预层的可表示变换 $F \to G$ 是增厚是什么意思。
我们指出，在 $F$ 和 $G$ 是代数空间时，这与《空间态射续篇》之定义
\ref{spaces-more-morphisms-definition-thickening} 中的增厚定义并不冲突。

\begin{lemma}
\label{formal-spaces-lemma-covering-by-thickenings}
设 $X_\lambda$（$\lambda \in \Lambda$）和 $X = \colim X_\lambda$
如定义 \ref{formal-spaces-definition-affine-formal-algebraic-space}。则
$X_\lambda \to X$ 可表示且为增厚。
\end{lemma}

\begin{proof}
由《空间》之第 \ref{spaces-section-representable} 节和第
\ref{spaces-section-representable-properties} 节的讨论，该陈述有意义。
由引理 \ref{formal-spaces-lemma-diagonal-affine-formal-algebraic-space}，态射
$X_\lambda \to X$ 可表示。给定 $U \to X$，其中 $U$ 是概形，
定义 \ref{formal-spaces-definition-affine-formal-algebraic-space} 后的讨论表明，
在 $U$ 上 Zariski 局部地，该态射分解经过某个满足
$\lambda \leq \mu$ 的 $X_\mu$。此时
$U \times_X X_\lambda = U \times_{X_\mu} X_\lambda$，
所以 $U \times_X X_\lambda \to U$ 是增厚
$X_\lambda \to X_\mu$ 的基变换。
\end{proof}

\begin{lemma}
\label{formal-spaces-lemma-factor-through-thickening}
设 $X_\lambda$（$\lambda \in \Lambda$）和 $X = \colim X_\lambda$
如定义 \ref{formal-spaces-definition-affine-formal-algebraic-space}。
若 $Y$ 是 $S$ 上的拟紧代数空间，则任意态射 $Y \to X$ 都分解经过某个
$X_\lambda$。
\end{lemma}

\begin{proof}
选取一个仿射概形 $V$ 和满的 étale 态射 $V \to Y$。
由定义 \ref{formal-spaces-definition-affine-formal-algebraic-space} 后的讨论，
复合 $V \to Y \to X$ 分解经过某个 $X_\lambda$。
由于 $V \to Y$ 是层的满射，结论成立。
\end{proof}

\begin{lemma}
\label{formal-spaces-lemma-characterize-affine-formal-algebraic-space}
设 $S$ 是概形，$X$ 是 $(\Sch/S)_{fppf}$ 上的层。
则 $X$ 是仿射形式代数空间，当且仅当下列条件成立：
\begin{enumerate}
\item 任意态射 $U \to X$，其中 $U$ 是 $S$ 上的仿射概形，
都分解经过一个可表示且为增厚的态射 $T \to X$，其中 $T$ 是 $S$ 上的仿射概形；
\item 满足如注 \ref{formal-spaces-remark-set-theoretic} 中的集合论条件。
\end{enumerate}
\end{lemma}

\begin{proof}
由引理 \ref{formal-spaces-lemma-covering-by-thickenings} 和
\ref{formal-spaces-lemma-factor-through-thickening}，仿射形式代数空间满足 (1) 和 (2)。
为证明逆命题，可以假设 $X$ 非空。令 $\Lambda$ 为如下范畴：
对象是可表示且为增厚的态射 $T \to X$，其中 $T$ 是 $S$ 上的仿射概形。
该范畴有向。因为 $X$ 非空，$\Lambda$ 至少含一个对象。
若 $T \to X$ 和 $T' \to X$ 属于 $\Lambda$，则可以把
$T \amalg T' \to X$ 分解经过 $\Lambda$ 中的 $T'' \to X$。
$\Lambda$ 的任意两个对象之间至多有唯一箭头。因此 $\Lambda$ 是有向集，且按假设
$X = \colim_{T \to X\text{ 属于 }\Lambda} T$。
为完成证明，须说明 $\Lambda$ 中任意箭头 $T \to T'$ 都是增厚。
这是因为 $T' \to X$ 是层的单态射，所以
$T = T \times_{T'} T' = T \times_X T'$；因而态射 $T \to T'$
等于投影 $T \times_X T' \to T'$，它是增厚，因为 $T \to X$ 是增厚。
\end{proof}

\noindent
一般的仿射形式代数空间 $X$ 不一定有足够多的函数来分离点（例如）。
见《例》之第 \ref{examples-section-affine-formal-algebraic-space} 节。
为刻画确实具有这种性质的空间，给出下列引理。

\begin{lemma}
\label{formal-spaces-lemma-mcquillan-affine-formal-algebraic-space}
设 $S$ 是概形。设 $X$ 是 $(\Sch/S)_{fppf}$ 上的 fppf 层，
满足注 \ref{formal-spaces-remark-set-theoretic} 的集合论条件。下列条件等价：
\begin{enumerate}
\item 存在 $S$ 上的弱可容许拓扑环 $A$（见注 \ref{formal-spaces-remark-mcquillan}），使得
$X = \colim_{I \subset A\text{ 为弱定义理想}} \Spec(A/I)$；
\item $X$ 是仿射形式代数空间，并且存在 $S$-代数 $A$ 和映射
$X \to \Spec(A)$，使得对任意闭浸入 $T \to X$，其中 $T$ 是仿射概形，
复合 $T \to \Spec(A)$ 是闭浸入；
\item $X$ 是仿射形式代数空间，并且存在 $S$-代数 $A$ 和映射
$X \to \Spec(A)$，使得对任意闭浸入 $T \to X$，其中 $T$ 是概形，
复合 $T \to \Spec(A)$ 是闭浸入；
\item $X$ 是仿射形式代数空间，并且对定义
\ref{formal-spaces-definition-affine-formal-algebraic-space} 中某个表示
$X = \colim X_\lambda$，各投影
$\lim \Gamma(X_\lambda, \mathcal{O}_{X_\lambda})
\to \Gamma(X_\lambda, \mathcal{O}_{X_\lambda})$ 都是满射；
\item $X$ 是仿射形式代数空间，并且对定义
\ref{formal-spaces-definition-affine-formal-algebraic-space} 中任意表示
$X = \colim X_\lambda$，各投影
$\lim \Gamma(X_\lambda, \mathcal{O}_{X_\lambda})
\to \Gamma(X_\lambda, \mathcal{O}_{X_\lambda})$ 都是满射。
\end{enumerate}
此外，该弱可容许拓扑环是
$A = \lim \Gamma(X_\lambda, \mathcal{O}_{X_\lambda})$，
赋予其极限拓扑；弱定义理想恰好分类那些可表示且为增厚的态射 $T \to X$。
\end{lemma}

\begin{proof}
(5) 推出 (4) 是显然的。

\medskip\noindent
对定义 \ref{formal-spaces-definition-affine-formal-algebraic-space} 中的
$X = \colim X_\lambda$ 假设 (4)。令
$A = \lim \Gamma(X_\lambda, \mathcal{O}_{X_\lambda})$。
设 $T \to X$ 是闭浸入，其中 $T$ 是概形
（注意，由引理 \ref{formal-spaces-lemma-diagonal-affine-formal-algebraic-space}，
$T \to X$ 可表示）。由于 $X_\lambda \to X$ 是增厚，
$X_\lambda \times_X T \to T$ 也是增厚。另一方面，
$X_\lambda \times_X T \to X_\lambda$ 是闭浸入，所以
$X_\lambda \times_X T$ 仿射。于是由《极限》之命题
\ref{limits-proposition-affine}，$T$ 仿射。再由引理
\ref{formal-spaces-lemma-factor-through-thickening}，$T \to X$ 分解经过某个
$X_\lambda$。因此
$A \to \Gamma(X_\lambda, \mathcal{O}) \to \Gamma(T, \mathcal{O})$
是满射。由此可见 (3) 成立。

\medskip\noindent
(3) 推出 (2) 是显然的。

\medskip\noindent
对 $A$ 和 $X \to \Spec(A)$ 假设 (2)。按定义
\ref{formal-spaces-definition-affine-formal-algebraic-space} 写
$X = \colim X_\lambda$。则按 (2)，
$A_\lambda = \Gamma(X_\lambda, \mathcal{O})$ 是 $A$ 的商。
所以 $A^\wedge = \lim A_\lambda$ 是完备拓扑环，见《代数续篇》之第
\ref{more-algebra-section-topological-ring} 节的讨论。因为
$A \to A_\lambda$ 是满射，各映射 $A^\wedge \to A_\lambda$ 也是满射。
我们断言，对任意 $\lambda$，$A^\wedge \to A_\lambda$ 的核
$I_\lambda \subset A^\wedge$ 都是弱定义理想。事实上，按极限拓扑的定义，它是开的。
若 $f \in I_\lambda$，则对任意 $\mu \in \Lambda$，$f$ 在
$A_\mu$ 中的像在 $X_\mu$ 各点的所有剩余域中都为零。因此它是
$A_\mu$ 的幂零元素，故某个幂 $f^n \in I_\mu$。所以
% P10-E0302
当 $n \to 0$ 时 $f^n \to 0$。因此 $A^\wedge$ 弱可容许。
最后，设 $I \subset A^\wedge$ 是弱定义理想。则
$I \subset A^\wedge$ 开，所以存在某个 $\lambda$ 使
$I \supset I_\lambda$。于是得到态射
$\Spec(A^\wedge/I) \to \Spec(A_\lambda) \to X$。
因而 $X = \colim \Spec(A^\wedge/I)$，这里余极限遍历所有弱定义理想。
所以 (1) 成立。

\medskip\noindent
假设 (1)。此时 $X$ 显然是仿射形式代数空间。任取定义
\ref{formal-spaces-definition-affine-formal-algebraic-space} 中的表示
$X = \colim X_\lambda$。对每个 $\lambda$，由引理
\ref{formal-spaces-lemma-factor-through-thickening}，可找到弱定义理想
$I \subset A$，使 $X_\lambda \to X$ 分解经过
$\Spec(A/I) \to X$。于是
$X_\lambda = \Spec(A/I_\lambda)$，其中 $I \subset I_\lambda$。
反过来，对任意弱定义理想 $I \subset A$，态射
$\Spec(A/I) \to X$ 分解经过某个 $X_\lambda$，即
$I_\lambda \subset I$。所以，每个 $I_\lambda$ 都是弱定义理想，
且它们组成弱定义理想集合的一个共尾子集。因此
$A = \lim A/I = \lim A/I_\lambda$，可见 (5) 成立，并且还有
$A = \lim \Gamma(X_\lambda, \mathcal{O}_{X_\lambda})$。
\end{proof}

\noindent
有了这个引理，可以作如下定义。

\begin{definition}
\label{formal-spaces-definition-types-affine-formal-algebraic-space}
设 $S$ 是概形，$X$ 是 $S$ 上的仿射形式代数空间。
若 $X$ 满足引理 \ref{formal-spaces-lemma-mcquillan-affine-formal-algebraic-space}
中的等价条件，则称 $X$ 是 {\it McQuillan} 的。
设 $A$ 是与 $X$ 相伴的弱可容许拓扑环。我们称
\begin{enumerate}
\item 若 $X$ 是 McQuillan 的且 $A$ 可容许
（《代数续篇》之定义 \ref{more-algebra-definition-topological-ring}），
则 $X$ 是{\it 经典的}；
\item 若 $X$ 是 McQuillan 的且 $A$ 弱 adic
（定义 \ref{formal-spaces-definition-weakly-adic}），则 $X$ {\it 弱 adic}；
\item 若 $X$ 是 McQuillan 的且 $A$ adic
（《代数续篇》之定义 \ref{more-algebra-definition-topological-ring}），
则 $X$ {\it adic}；
\item 若 $X$ 是 McQuillan 的、$A$ adic 且 $A$ 有有限生成的定义理想，
则 $X$ {\it adic*}；
\item 若 $X$ 是 McQuillan 的，且 $A$ 既 Noether 又 adic，
则 $X$ 是 {\it Noether} 的。
\end{enumerate}
\end{definition}

\noindent
在 \cite{Fujiwara-Kato} 中，adic 拓扑环 $A$ 包含有限生成定义理想这一性质
被称为“有限理想型”。给定仿射形式代数空间 $X$，定义中引入的概念之间有如下蕴含：
$$
\xymatrix{
X\text{ 是 Noether 的} \ar@{=>}[r] &
X\text{ 是 adic* 的} \ar@{=>}[r] &
X\text{ 是 adic 的} \ar@{=>}[lld] \\
X\text{ 是弱 adic 的} \ar@{=>}[r] &
X\text{ 是经典的} \ar@{=>}[r] &
X\text{ 是 McQuillan 的}
}
$$
见第 \ref{formal-spaces-section-weakly-adic} 节的讨论；精确陈述见引理
\ref{formal-spaces-lemma-implications-between-types}。

\begin{remark}
\label{formal-spaces-remark-compare-with-affine-formal-schemes}
经典仿射形式代数空间对应于 EGA（\cite{EGA}）所考虑的仿射形式概形。
为解释这一点，假设基概形是 $\Spec(\mathbf{Z})$。设
$\mathfrak X = \text{Spf}(A)$ 是仿射形式概形，并令
$h_\mathfrak X$ 为引理 \ref{formal-spaces-lemma-fully-faithful} 中的点函子。
则 $h_\mathfrak X = \colim h_{\Spec(A/I)}$，其中余极限遍历可容许拓扑环
$A$ 的定义理想集合。在仿射概形上取值时，这由式
(\ref{formal-spaces-equation-morphisms-affine-formal-schemes}) 得出；只须在仿射概形上核验，
因为两边都是 fppf 层，见引理 \ref{formal-spaces-lemma-formal-scheme-sheaf-fppf}。
因此 $h_\mathfrak X$ 是仿射形式代数空间。事实上，由定义
\ref{formal-spaces-definition-types-affine-formal-algebraic-space}，它是经典仿射形式代数空间。
所以引理 \ref{formal-spaces-lemma-fully-faithful} 告诉我们，仿射形式概形范畴
等价于经典仿射形式代数空间范畴。
\end{remark}

\noindent
作出上述与仿射形式概形的联系后，作如下定义显得自然。

\begin{definition}
\label{formal-spaces-definition-affine-formal-spectrum}
设 $S$ 是概形。设 $A$ 是 $S$ 上的弱可容许拓扑环，见定义
\ref{formal-spaces-definition-weakly-admissible}\footnote{经典情形见《代数续篇》之定义
\ref{more-algebra-definition-topological-ring}，差异的讨论见注
\ref{formal-spaces-remark-mcquillan}。}。$A$ 的{\it 形式谱}是仿射形式代数空间
$$
\text{Spf}(A) = \colim \Spec(A/I)
$$
其中余极限遍历 $A$ 的弱定义理想集合，并在范畴
$\Sh((\Sch/S)_{fppf})$ 中取。
\end{definition}

\noindent
这样的形式谱按构造是 McQuillan 的；反过来，每个 McQuillan 仿射形式代数空间
都同构于某个形式谱。需要指出的是，在我们的理论中，确实存在不能写成任何弱可容许拓扑环
的形式谱的仿射形式代数空间。依照 \cite{Yasuda}，可以把 $S$-pro-环定义为
$S$-代数范畴中的 pro-对象，见《范畴》之注
\ref{categories-remark-pro-category}。这样，每个 $S$ 上的仿射形式代数空间
都会是某个 $S$-pro-环的形式谱。我们不这样做，而将直接处理相应的仿射形式代数空间。

\medskip\noindent
形式谱的构造是函子性的。为解释这一点，设
% P10-E0303
$\varphi : B \to A$ 是 $S$ 上弱可容许拓扑环的连续映射。则
$$
\text{Spf}(\varphi) : \text{Spf}(B) \to \text{Spf}(A)
$$
是唯一的仿射形式代数空间态射，使得对所有满足
$\varphi(I) \subset J$ 的弱定义理想 $I \subset A$ 和 $J \subset B$，图表
$$
\xymatrix{
\Spec(B/J) \ar[d] \ar[r] & \Spec(A/I) \ar[d] \\
\text{Spf}(B) \ar[r] & \text{Spf}(A)
}
$$
交换。由于 $\varphi$ 的连续性蕴含：对每个弱定义理想 $J \subset B$，
存在具有所需性质的弱定义理想 $I \subset A$，可见所要求的交换性唯一决定并定义
$\text{Spf}(\varphi)$。

\begin{lemma}
\label{formal-spaces-lemma-morphism-between-formal-spectra}
设 $S$ 是概形，$A,B$ 是 $S$ 上的弱可容许拓扑环。
$S$ 上仿射形式代数空间的任意态射
$f : \text{Spf}(B) \to \text{Spf}(A)$ 都等于
$\text{Spf}(f^\sharp)$，其中
$f^\sharp : A \to B$ 是唯一的连续 $S$-代数映射。
\end{lemma}

\begin{proof}
设 $f : \text{Spf}(B) \to \text{Spf}(A)$ 如引理所述。
设 $J \subset B$ 是弱定义理想。由引理
\ref{formal-spaces-lemma-factor-through-thickening}，存在弱定义理想 $I \subset A$，使复合
$\Spec(B/J) \to \text{Spf}(B) \to \text{Spf}(A)$
分解经过 $\Spec(A/I)$。由《概形》之引理
\ref{schemes-lemma-morphism-into-affine}，得到 $S$-代数映射
$A/I \to B/J$。这些映射随 $J$ 变化相容，并定义映射
$f^\sharp : A \to B$。该映射连续，因为对每个弱定义理想
$J \subset B$，都存在弱定义理想 $I \subset A$ 使
$f^\sharp(I) \subset J$。等式 $f = \text{Spf}(f^\sharp)$
由我们对组成 $f^\sharp$ 的环映射 $A/I \to B/J$ 的选择成立。
\end{proof}

\begin{lemma}
\label{formal-spaces-lemma-presentation-representable}
设 $S$ 是概形，$f : X \to Y$ 是 $(\Sch/S)_{fppf}$ 上预层的映射。
若 $X$ 是仿射形式代数空间，而 $f$ 可由代数空间表示且局部拟有限，
则 $f$ 可表示（由概形表示）。
\end{lemma}

\begin{proof}
设 $T$ 是 $S$ 上的概形，且给定映射 $T \to Y$。须证明代数空间
$X \times_Y T$ 是概形。按定义
\ref{formal-spaces-definition-affine-formal-algebraic-space} 写
$X = \colim X_\lambda$。设 $W \subset X \times_Y T$ 是拟紧开子空间。
投影 $X \times_Y T \to X$ 在 $W$ 上的限制分解经过某个 $X_\lambda$。于是
$$
W \to X_\lambda \times_S T
$$
是单态射（从而分离）且局部拟有限（因为按对 $X \to Y$ 的假设，
$W \to X \times_Y T \to T$ 局部拟有限；见《空间的态射》之引理
\ref{spaces-morphisms-lemma-permanence-quasi-finite}）。
所以由《空间的态射》之命题
\ref{spaces-morphisms-proposition-locally-quasi-finite-separated-over-scheme}，
$W$ 是概形。因此由《空间的性质》之引理
\ref{spaces-properties-lemma-subscheme}，$X \times_Y T$ 是概形。
\end{proof}






\section{可数指标的仿射形式代数空间}
\label{formal-spaces-section-countably-indexed}

\noindent
这类仿射形式代数空间由下列引理描述。

\begin{lemma}
\label{formal-spaces-lemma-countable-affine-formal-algebraic-space}
设 $S$ 是概形，$X$ 是 $S$ 上的仿射形式代数空间。下列条件等价：
\begin{enumerate}
\item 存在 $S$ 上仿射概形的一列增厚
$X_1 \to X_2 \to X_3 \to \ldots$，使得 $X = \colim X_n$；
\item 存在定义 \ref{formal-spaces-definition-affine-formal-algebraic-space}
中的一种选择 $X = \colim X_\lambda$，使得 $\Lambda$ 可数。
\end{enumerate}
\end{lemma}

\begin{proof}
这来自如下观察：可数有向集有一个同构于 $(\mathbf{N}, \geq)$ 的共尾子集。
见《代数》之引理 \ref{algebra-lemma-ML-limit-nonempty} 的证明。
\end{proof}

\begin{definition}
\label{formal-spaces-definition-countable}
设 $S$ 是概形，$X$ 是 $S$ 上的仿射形式代数空间。
若引理 \ref{formal-spaces-lemma-countable-affine-formal-algebraic-space} 的等价条件成立，
则称 $X$ 是{\it 可数指标的}。
\end{definition}

\noindent
用 \cite{BVGD} 的语言，这表述为 $X$ 是 $\aleph_0$-ind 概形。

\begin{lemma}
\label{formal-spaces-lemma-implications-between-types}
设 $X$ 是概形 $S$ 上的仿射形式代数空间。
\begin{enumerate}
\item 若 $X$ 是 Noether 的，则 $X$ adic*。
\item 若 $X$ adic*，则 $X$ adic。
\item 若 $X$ adic，则 $X$ 弱 adic。
\item 若 $X$ 弱 adic，则 $X$ 经典。
\item 若 $X$ 弱 adic，则 $X$ 是可数指标的。
\item 若 $X$ 是可数指标的，则 $X$ 是 McQuillan 的。
\end{enumerate}
\end{lemma}

\begin{proof}
% P10-E0304
把 $X$ 写成 $X = \text{Spf}(A)$，其中 $A$ 是弱可容许的（从而完备的）
线性拓扑环。利用第 \ref{formal-spaces-section-weakly-adic} 节讨论的这类环各种类型之间的蕴含，
即可得到 (1)、(2)、(3) 和 (4)。

\medskip\noindent
(5) 的证明。按定义，存在弱 adic 拓扑环 $A$，使
$X = \colim \Spec(A/I)$，其中余极限遍历 $A$ 的定义理想。
由于 $A$ 弱 adic，特别地，
% P10-E0305
存在一族可数的开理想基本系 $I_\lambda$，见定义
\ref{formal-spaces-definition-weakly-adic}。于是按 $\text{Spf}(A)$ 的定义，
% P10-E0306
$X = \colim \Spec(A/I_n)$。所以 $X$ 是可数指标的。

\medskip\noindent
(6) 的证明。把 $X = \colim X_n$ 写成 $S$ 上仿射概形的一列增厚
$X_1 \to X_2 \to X_3 \to \ldots$ 的余极限。则
$$
A = \lim \Gamma(X_n, \mathcal{O}_{X_n})
$$
满射到每个 $\Gamma(X_n, \mathcal{O}_{X_n})$，因为态射
$X_n \to X_{n + 1}$ 是闭浸入，所以转移映射是满射。因此 $X$ 是 McQuillan 的。
\end{proof}

\begin{lemma}
\label{formal-spaces-lemma-countably-indexed}
设 $S$ 是概形，$X$ 是 $(\Sch/S)_{fppf}$ 上的预层。下列条件等价：
\begin{enumerate}
\item $X$ 是可数指标的仿射形式代数空间；
\item $X = \text{Spf}(A)$，其中 $A$ 是有可数 $0$ 邻域基本系的弱可容许拓扑
$S$-代数；
\item $X = \text{Spf}(A)$，其中 $A$ 是弱可容许拓扑 $S$-代数，
且有一族弱定义理想基本系
$A \supset I_1 \supset I_2 \supset I_3 \supset \ldots$；
\item $X = \text{Spf}(A)$，其中 $A$ 是完备拓扑 $S$-代数，
它的一族 $0$ 的开邻域基本系由理想的可数序列
$A \supset I_1 \supset I_2 \supset I_3 \supset \ldots$ 给出，
并且 $I_n/I_{n + 1}$ 局部幂零；
\item $X = \text{Spf}(A)$，其中 $A = \lim B/J_n$ 带极限拓扑，
而 $B \supset J_1 \supset J_2 \supset J_3 \supset \ldots$
是 $S$-代数 $B$ 中的一列理想，且 $J_n/J_{n + 1}$ 局部幂零。
\end{enumerate}
\end{lemma}

\begin{proof}
假设 (1)。由引理 \ref{formal-spaces-lemma-implications-between-types}，
可写 $X = \text{Spf}(A)$，其中 $A$ 是弱可容许拓扑 $S$-代数。
对引理 \ref{formal-spaces-lemma-countable-affine-formal-algebraic-space} (1)
中的任意表示 $X = \colim X_n$，可见
$A = \lim A_n$，其中 $X_n = \Spec(A_n)$，并且对某个弱定义理想
$I_n \subset A$ 有 $A_n = A/I_n$。这来自引理
\ref{formal-spaces-lemma-mcquillan-affine-formal-algebraic-space} 的最后一个陈述，
该陈述还蕴含 $\{I_n\}$ 是一族 $0$ 的开邻域基本系。
因此有弱定义理想的序列
$$
A \supset I_1 \supset I_2 \supset I_3 \supset \ldots
$$
且 $A = \lim A/I_n$。由此可见，条件 (1) 推出条件 (2)--(5) 中的每一个。

\medskip\noindent
假设 (5)。首先注意，$A = \lim B/J_n$ 上的极限拓扑是完备的线性拓扑，
见《代数续篇》之第 \ref{more-algebra-section-topological-ring} 节。
若 $f \in A$ 在 $B/J_1$ 中映为零，则由于其在 $J_1/J_2$ 中的像幂零，
某个幂在 $B/J_2$ 中映为零；再取一个幂，它在 $J_2/J_3$ 中映为零，等等。
由此可见，开理想 $\Ker(A \to B/J_1)$ 是弱定义理想。所以 $A$ 弱可容许。
由此可见 (5) 推出 (2)。

\medskip\noindent
取 $B = A$，显然 (4) 是 (5) 的特殊情形。显然 (3) 是 (2) 的特殊情形。

\medskip\noindent
假设 $A$ 如 (2) 所述。设 $E_n$ 是 $A$ 中一族可数的 $0$ 邻域基本系。
由于 $A$ 是弱可容许拓扑环，可以找到开理想 $I_n \subset E_n$。
还可以选取弱定义理想 $J \subset A$。则 $J \cap I_n$ 是 $A$
的一族弱定义理想基本系，并得到
$X = \text{Spf}(A) = \colim \Spec(A/(J \cap I_n))$，
这表明 $X$ 是可数指标的仿射形式代数空间。
\end{proof}

\begin{lemma}
\label{formal-spaces-lemma-characterize-noetherian-affine}
设 $S$ 是概形，$X$ 是仿射形式代数空间。下列条件等价：
\begin{enumerate}
\item $X$ 是 Noether 的；
\item $X$ adic*，并且对每个闭浸入 $T \to X$，其中 $T$ 是概形，
$T$ 都是 Noether 的；
\item $X$ adic*，并且对定义
\ref{formal-spaces-definition-affine-formal-algebraic-space} 中某种选择
$X = \colim X_\lambda$，概形 $X_\lambda$ 都是 Noether 的；
\item $X$ 弱 adic，并且对定义
\ref{formal-spaces-definition-affine-formal-algebraic-space} 中某种选择
$X = \colim X_\lambda$，概形 $X_\lambda$ 都是 Noether 的。
\end{enumerate}
\end{lemma}

\begin{proof}
假设 $X$ 是 Noether 的。则 $X = \text{Spf}(A)$，其中 $A$ 是 Noether adic 环。
设 $T \to X$ 是闭浸入，其中 $T$ 是概形。由引理
\ref{formal-spaces-lemma-mcquillan-affine-formal-algebraic-space}，$T$ 仿射，且
$T \to \Spec(A)$ 是闭浸入。由于 $A$ Noether，可见 $T$ Noether。
由此可见 (1) $\Rightarrow$ (2)。

\medskip\noindent
蕴含 (2) $\Rightarrow$ (3) 和 (2) $\Rightarrow$ (4) 是立即的
（见引理 \ref{formal-spaces-lemma-implications-between-types}）。

\medskip\noindent
为证明 (3) $\Rightarrow$ (1)，把 $X = \text{Spf}(A)$ 写成某个
adic 环 $A$ 的形式谱，其中 $A$ 有有限生成的定义理想 $I$。
还给定：对某个开理想基本系 $I_\lambda$，环 $A/I_\lambda$ 都 Noether。
由于 $I$ 开，可以找到 $\lambda$ 使 $I_\lambda \subset I$。
则 $A/I$ Noether，由《代数》之引理
\ref{algebra-lemma-completion-Noetherian} 可知 $A$ Noether。

\medskip\noindent
为证明 (4) $\Rightarrow$ (3)，把 $X = \text{Spf}(A)$ 写成某个弱 adic 环
$A$ 的形式谱。则 $A$ 可容许，并有一个定义理想 $I$；而且 $I^2$ 的闭包
$I_2$ 是开的，见引理 \ref{formal-spaces-lemma-weakly-adic}。
还给定：对某个开理想基本系 $I_\lambda$，环 $A/I_\lambda$ 都 Noether。
选取 $\lambda$ 使 $I_\lambda \subset I_2$。则 $A/I_2$ 是
$A/I_\lambda$ 的商，所以 Noether。故 $I/I_2$ 是有限 $A$-模。
由引理 \ref{formal-spaces-lemma-weakly-adic-finite-generation}，$A$ 是有有限生成定义理想的
adic 环。因此 $X$ adic*，即 (3) 成立。
\end{proof}





\section{形式代数空间}
\label{formal-spaces-section-formal-algebraic-spaces}

\noindent
我们暂时打破先对环、再对概形、最后对代数空间引入新概念的习惯，
不先引入形式概形便直接引入形式代数空间。一般思路是：形式代数空间是 fppf 拓扑中的层，
并且 étale 局部地是 \cite{BVGD} 意义下的仿射形式概形。
相关材料可见 \cite{Yasuda}。

\medskip\noindent
在形式代数空间的定义中，我们将借用《自举》之第
\ref{bootstrap-section-morphism-representable-by-spaces} 节和第
\ref{bootstrap-section-representable-by-spaces-properties} 节中的若干术语。

\begin{definition}
\label{formal-spaces-definition-formal-algebraic-space}
设 $S$ 是概形。若存在一族层映射
$\{X_i \to X\}_{i \in I}$，使得
\begin{enumerate}
\item $X_i$ 是仿射形式代数空间；
\item $X_i \to X$ 可由代数空间表示且为 étale；
\item $\coprod X_i \to X$ 作为层映射是满射，
\end{enumerate}
并且 $X$ 满足一个集合论条件（见注 \ref{formal-spaces-remark-set-theoretic}），
则称 $(\Sch/S)_{fppf}$ 上的层 $X$ 是{\it 形式代数空间}。
$S$ 上形式代数空间的{\it 态射}是层的映射。
\end{definition}

\noindent
讨论。合理性检查：仿射形式代数空间是形式代数空间。
在定义的情形中，态射 $X_i \to X$ 可表示（由概形表示），
见引理 \ref{formal-spaces-lemma-presentation-representable}。由《自举》之引理
\ref{bootstrap-lemma-surjective-flat-locally-finite-presentation}，
也可以不要求 $\coprod X_i \to X$ 作为层映射是满射，而只要求它满
（这是有意义的，因为它可表示）。

\medskip\noindent
我们的形式代数空间概念{\bf 非常一般}。事实上，即使上面定义的仿射形式代数空间
也是非常难处理的对象。

\begin{lemma}
\label{formal-spaces-lemma-diagonal-formal-algebraic-space}
设 $S$ 是概形。若 $X$ 是 $S$ 上的形式代数空间，则对角态射
$\Delta : X \to X \times_S X$ 可表示、为单态射、局部拟有限、局部有限型且分离。
\end{lemma}

\begin{proof}
设给定 $U \to X$ 和 $V \to X$，其中 $U,V$ 是 $S$ 上的概形。
则 $U \times_X V$ 是层。选取定义 \ref{formal-spaces-definition-formal-algebraic-space}
中的 $\{X_i \to X\}$。对每个 $i$，态射
$$
(U \times_X X_i) \times_{X_i} (V \times_X X_i)
= (U \times_X V) \times_X X_i \to U \times_X V
$$
作为 $X_i \to X$ 的基变换可表示且为 étale，而其源是概形
（使用引理 \ref{formal-spaces-lemma-diagonal-affine-formal-algebraic-space} 和
\ref{formal-spaces-lemma-presentation-representable}）。这些映射联合满，所以由《自举》之定理
\ref{bootstrap-theorem-final-bootstrap}，$U \times_X V$ 是代数空间。
态射 $U \times_X V \to U \times_S V$ 是单态射。它还局部拟有限，
因为预复合上面显示的态射后，得到复合
$$
(U \times_X X_i) \times_{X_i} (V \times_X X_i)
\to (U \times_X X_i) \times_S (V \times_X X_i)
\to U \times_S V
$$
它是闭浸入（引理 \ref{formal-spaces-lemma-diagonal-affine-formal-algebraic-space}）
与 étale 态射的复合，故局部拟有限，见《空间上的下降》之引理
\ref{spaces-descent-lemma-locally-quasi-finite-etale-local-source}。
因此由《空间的态射》之命题
\ref{spaces-morphisms-proposition-locally-quasi-finite-separated-over-scheme}，
$U \times_X V$ 是概形。所以由《空间》之引理
\ref{spaces-lemma-representable-diagonal}，$\Delta$ 可表示。

\medskip\noindent
事实上，上面已经证明概形态射
$U \times_X V \to U \times_S V$ 总是单态射且局部拟有限，
所以由《空间》之引理 \ref{spaces-lemma-transformation-diagonal-properties}，
$\Delta : X \to X \times_S X$ 是单态射且局部拟有限。
然后可用《空间》之引理
\ref{spaces-lemma-representable-transformations-property-implication} 的原理，
看出 $\Delta$ 分离且局部有限型。事实上，概形的单态射分离
（《概形》之引理 \ref{schemes-lemma-monomorphism-separated}），
而概形的局部拟有限态射局部有限型
（由《态射》之第 \ref{morphisms-section-quasi-finite} 节中的定义得到）。
\end{proof}

\begin{lemma}
\label{formal-spaces-lemma-space-to-formal-space}
设 $S$ 是概形，$f : X \to Y$ 是从 $S$ 上代数空间到 $S$ 上形式代数空间的态射。
则 $f$ 可由代数空间表示。
\end{lemma}

\begin{proof}
设 $Z \to Y$ 是态射，其中 $Z$ 是 $S$ 上的概形。须证明
$X \times_Y Z$ 是代数空间。选取概形 $U$ 以及满的 étale 态射 $U \to X$。
则 $U \times_Y Z \to X \times_Y Z$ 是可表示、满且 étale 的
（《空间》之引理
\ref{spaces-lemma-base-change-representable-transformations-property}），
并且由引理 \ref{formal-spaces-lemma-diagonal-formal-algebraic-space}，
$U \times_Y Z$ 是概形。结论遂由《自举》之定理
\ref{bootstrap-theorem-final-bootstrap} 得出。
\end{proof}

\begin{remark}
\label{formal-spaces-remark-compare-with-formal-schemes}
略去集合论问题，EGA 意义下的形式概形范畴
（见第 \ref{formal-spaces-section-formal-schemes-EGA} 节）等价于形式代数空间范畴的一个全子范畴。
为解释这一点，假设基概形是 $\Spec(\mathbf{Z})$。
由引理 \ref{formal-spaces-lemma-formal-scheme-sheaf-fppf}，与形式概形
$\mathfrak X$ 相伴的点函子 $h_\mathfrak X$ 是 fppf 拓扑中的层。
由引理 \ref{formal-spaces-lemma-fully-faithful}，赋值
$\mathfrak X \mapsto h_\mathfrak X$ 是从形式概形范畴到 fppf 层范畴的全忠实嵌入。
给定形式概形 $\mathfrak X$，选取由仿射形式概形 $\mathfrak X_i$ 组成的开覆盖
$\mathfrak X = \bigcup \mathfrak X_i$。由注
\ref{formal-spaces-remark-compare-with-affine-formal-schemes}，$h_{\mathfrak X_i}$
是仿射形式代数空间。态射
$h_{\mathfrak X_i} \to h_\mathfrak X$ 可表示且为开浸入。
所以 $\{h_{\mathfrak X_i} \to h_\mathfrak X\}$ 是定义
\ref{formal-spaces-definition-formal-algebraic-space} 中的一族映射，
可见 $h_\mathfrak X$ 是形式代数空间。
\end{remark}

\begin{remark}
\label{formal-spaces-remark-set-theoretic}
设 $S$ 是概形，并设 $(\Sch/S)_{fppf}$ 是《拓扑》之定义
\ref{topologies-definition-big-small-fppf} 中的大 fppf 站点。
在定义 \ref{formal-spaces-definition-affine-formal-algebraic-space} 和
\ref{formal-spaces-definition-formal-algebraic-space} 中，采用如下 $X$ 的集合论条件：
存在 $(\Sch/S)_{fppf}$ 的对象 $U,R$、作为 fppf 层满射的态射 $U \to X$，
以及作为 fppf 层满射的态射 $R \to U \times_X U$。
换言之，要求我们的层是两个可表示层之间两映射的余等化子。
下面是说明这一概念表现合理的若干观察：
\begin{enumerate}
\item 设 $X = \colim_{\lambda \in \Lambda} X_\lambda$，
并且该系统满足定义 \ref{formal-spaces-definition-affine-formal-algebraic-space} 的条件 (1) 和 (2)。
则 $U = \coprod_{\lambda \in \Lambda} X_\lambda \to X$ 是 fppf 层的满射。
此外，由引理 \ref{formal-spaces-lemma-diagonal-affine-formal-algebraic-space}，
$U \times_X U$ 是 $U \times_S U$ 的闭子概形。
所以，若 $U$ 可由 $(\Sch/S)_{fppf}$ 的对象表示，则 $U \times_S U$ 也可
（见《集合》之引理 \ref{sets-lemma-what-is-in-it}），集合论条件得到满足。
当 $\Lambda$ 可数时总是如此，见《集合》之引理
\ref{sets-lemma-what-is-in-it}。
\item 合理性检查。设 $\{X_i \to X\}_{i \in I}$ 如定义
\ref{formal-spaces-definition-formal-algebraic-space}（集合论条件按上文表述），
并假设每个 $X_i$ 实际上都是仿射概形。则 $X$ 是代数空间。
事实上，若选取更大的大 fppf 站点 $(\Sch'/S)_{fppf}$，使
$U' = \coprod X_i$ 和 $R' = \coprod X_i \times_X X_j$
都可由其中的对象表示，则 $X' = U'/R'$ 是这一选择下代数空间范畴的对象。
然后应用《空间》之引理 \ref{spaces-lemma-fully-faithful}，可知
$X$ 是 $(\Sch/S)_{fppf}$ 上的代数空间。
\item 设 $\{X_i \to X\}_{i \in I}$ 是满足定义
\ref{formal-spaces-definition-formal-algebraic-space} 中条件 (1)、(2)、(3) 的一族层映射。
对每个 $i$，可以选取 $U_i \in \Ob((\Sch/S)_{fppf})$ 以及层的满射
$U_i \to X_i$。所以，若 $I$ 不过大（例如可数），则
$U = \coprod U_i \to X$ 是层的满射，并且 $U$ 可由
$(\Sch/S)_{fppf}$ 的对象表示。要得到满射到 $U \times_X U$ 的
$R \in \Ob((\Sch/S)_{fppf})$，只须假设对角
$\Delta : X \to X \times_S X$ 不太难处理；例如，若 $X$ 的对角拟紧，
即 $X$ 拟分离，就总能做到。
\end{enumerate}
\end{remark}







\section{约化}
\label{formal-spaces-section-reduction}

\noindent
如下引理所示，每个形式代数空间都有一个底层既约代数空间。

\begin{lemma}
\label{formal-spaces-lemma-reduction-formal-algebraic-space}
设 $S$ 是概形，$X$ 是 $S$ 上的形式代数空间。
存在既约代数空间 $X_{red}$ 以及可表示且为增厚的态射
$X_{red} \to X$。若 $U$ 是既约代数空间，则态射 $U \to X$
唯一地分解经过 $X_{red}$。
\end{lemma}

\begin{proof}
首先假设 $X$ 是仿射形式代数空间。按定义
\ref{formal-spaces-definition-affine-formal-algebraic-space} 写
$X = \colim X_\lambda$。由于转移态射都是增厚，仿射概形
$X_\lambda$ 都有彼此同构的约化 $X_{red}$。
由引理 \ref{formal-spaces-lemma-covering-by-thickenings} 以及增厚的复合仍为增厚这一事实，
态射 $X_{red} \to X$ 可表示且为增厚。我们略去泛性质的核验
（使用《概形》之定义 \ref{schemes-definition-reduced-induced-scheme}、
《概形》之引理 \ref{schemes-lemma-map-into-reduction}、
《空间的性质》之定义 \ref{spaces-properties-definition-reduced-induced-space}
以及《空间的性质》之引理
\ref{spaces-properties-lemma-map-into-reduction}）。

\medskip\noindent
设 $X$ 和 $\{X_i \to X\}_{i \in I}$ 如定义
\ref{formal-spaces-definition-formal-algebraic-space}。对每个 $i$，令
$X_{i, red} \to X_i$ 为上面构造的约化。对 $i,j \in I$，投影
$X_{i, red} \times_X X_j \to X_{i, red}$ 是概形的 étale 态射
（étale 由假设，源是概形由引理 \ref{formal-spaces-lemma-presentation-representable}）。
所以 $X_{i, red} \times_X X_j$ 既约
（见《下降》之引理 \ref{descent-lemma-reduced-local-smooth}）。
因此，由泛性质，投影
$X_{i, red} \times_X X_j \to X_j$ 分解经过 $X_{j, red}$。
由于态射 $X_{i, red} \to X_i$ 是层的单射，可得
$$
R_{ij} = X_{i, red} \times_X X_j = X_{i, red} \times_X X_{j, red} =
X_i \times_X X_{j, red}
$$
令 $U = \coprod X_{i, red}$、$R = \coprod R_{ij}$，并以
$s,t : R \to U$ 表示两个投影。作为层，$R = U \times_X U$，
并且 $s$ 和 $t$ 都是 étale 的。于是
% P10-E0307
$(t,s) : R \to U$ 由上述观察定义一个 étale 等价关系。
所以由《空间》之定理 \ref{spaces-theorem-presentation}，
$X_{red} = U/R$ 是代数空间。按构造，图表
$$
\xymatrix{
\coprod X_{i, red} \ar[r] \ar[d] & \coprod X_i \ar[d] \\
X_{red} \ar[r] & X
}
$$
是笛卡尔的。由于右侧竖直箭头 étale 且满，顶部水平箭头可表示且为增厚，
由《自举》之引理 \ref{bootstrap-lemma-after-fppf-sep-lqf}，
$X_{red} \to X$ 可表示
（核验该引理的假设时，使用满的 étale 态射是满、平坦且局部有限表示的，
并使用增厚分离且局部拟有限）。然后可以用《空间》之引理
\ref{spaces-lemma-descent-representable-transformations-property}
推出 $X_{red} \to X$ 是增厚
（使用增厚等价于满的闭浸入这一事实）。

\medskip\noindent
最后，设 $U \to X$ 是态射，其中 $U$ 是 $S$ 上的既约代数空间。
则每个 $X_i \times_X U$ 都在 $U$ 上 étale，因而既约
（按《空间的性质》之第
\ref{spaces-properties-section-types-properties} 节中既约代数空间的定义）。
所以 $X_i \times_X U \to X_i$ 分解经过 $X_{i,red}$。
由于 $\{X_i \times_X U \to U\}$ 是 étale 覆盖，$U \to X$
分解经过 $X_{red}$。
\end{proof}

\begin{example}
\label{formal-spaces-example-reduction-affine-formal-spectrum}
设 $A$ 是弱可容许拓扑环。此时
$$
\text{Spf}(A)_{red} = \Spec(A/\mathfrak a)
$$
其中 $\mathfrak a \subset A$ 是拓扑幂零元素的理想。
事实上，$\mathfrak a$ 是根式理想
（引理 \ref{formal-spaces-lemma-topologically-nilpotent}），并且因为 $A$ 弱可容许，它是开的。
\end{example}

\begin{lemma}
\label{formal-spaces-lemma-reduction-smooth}
设 $S$ 是概形，$f : X \to Y$ 是 $S$ 上形式代数空间的态射，
它可由代数空间表示且光滑（例如 étale）。则
$X_{red} = X \times_Y Y_{red}$。
\end{lemma}

\begin{proof}
（étale 情形直接来自引理 \ref{formal-spaces-lemma-reduction-formal-algebraic-space}
证明中对底层既约代数空间的构造。）假设 $f$ 光滑。
注意，$X \times_Y Y_{red} \to Y_{red}$ 是代数空间的光滑态射。
所以由《空间上的下降》之引理
\ref{spaces-descent-lemma-reduced-local-smooth}，
$X \times_Y Y_{red}$ 是既约代数空间。于是约化的泛性质表明典范态射
$X_{red} \to X \times_Y Y_{red}$ 是同构。
\end{proof}

\begin{lemma}
\label{formal-spaces-lemma-reduction-surjective}
设 $S$ 是概形，$f : X \to Y$ 是 $S$ 上形式代数空间的态射，
它可由代数空间表示。则 $f$ 在《自举》之定义
\ref{bootstrap-definition-property-transformation} 的意义下为满射，
当且仅当 $f_{red} : X_{red} \to Y_{red}$ 是代数空间的满态射。
\end{lemma}

\begin{proof}
略。
\end{proof}






\section{沿增厚的代数空间余极限}
\label{formal-spaces-section-global-colimits}

\noindent
一种特殊的形式代数空间可以像下面引理那样，在整体上写成沿增厚的代数空间余滤余极限。
稍后（在第 \ref{formal-spaces-section-quasi-compact-quasi-separated} 节）将看到，
任意拟紧且拟分离的形式代数空间都是这样的整体余极限。

\begin{lemma}
\label{formal-spaces-lemma-colimit-is-formal}
设 $S$ 是概形。给定有向集 $\Lambda$ 以及位于 $\Lambda$ 上的代数空间系统
$(X_\lambda, f_{\lambda \mu})$，其中每个
$f_{\lambda \mu} : X_\lambda \to X_\mu$ 都是增厚。
则 $X = \colim_{\lambda \in \Lambda} X_\lambda$ 是 $S$ 上的形式代数空间。
\end{lemma}

\begin{proof}
因为在 fppf 层范畴中取余极限，$X$ 是层。选定 $\lambda \in \Lambda$。
选取 étale 覆盖 $\{X_{i,\lambda} \to X_\lambda\}$，其中
% P10-E0308
$X_i$ 是 $S$ 上的仿射概形，见《空间的性质》之引理
\ref{spaces-properties-lemma-cover-by-union-affines}。
对每个 $\mu \geq \lambda$，存在笛卡尔图表
$$
\xymatrix{
X_{i, \lambda} \ar[r] \ar[d] & X_{i, \mu} \ar[d] \\
X_\lambda \ar[r] & X_\mu
}
$$
其竖直箭头为 étale，见《空间态射续篇》之定理
\ref{spaces-more-morphisms-theorem-topological-invariance}
（这里还用到增厚是满足该定理条件的满闭浸入）。此外，这些图表在唯一同构意义下唯一，
所以当 $\mu' \geq \mu$ 时
$X_{i,\mu} = X_\mu \times_{X_{\mu'}} X_{i,\mu'}$。
% P10-E0309
态射 $X_{i,\mu} \to X_{i,\mu'}$ 是增厚，因为它是增厚的基变换。
由《空间的极限》之命题 \ref{spaces-limits-proposition-affine}
以及 $X_{i,\lambda}$ 仿射这一事实，每个 $X_{i,\mu}$ 都是仿射概形。
令 $X_i = \colim_{\mu \geq \lambda} X_{i,\mu}$。
则 $X_i$ 是仿射形式代数空间。态射 $X_i \to X$ 是 étale 的，
因为给定仿射概形 $U$，任意 $U \to X$ 都分解经过某个
$\mu \geq \lambda$ 对应的 $X_\mu$（细节略）。由此可见 $X$ 是形式代数空间。
\end{proof}

\noindent
设 $S$ 是概形，$X$ 是 $S$ 上的形式代数空间。如何证明或检验
$X$ 是引理 \ref{formal-spaces-lemma-colimit-is-formal} 中那样的整体余极限？
为此，寻找映射 $i : Z \to X$，其中 $Z$ 是 $S$ 上的代数空间，而
$i$ 是满的闭浸入；换言之，$i$ 是增厚。这是有意义的，因为
$i$ 可由代数空间表示（引理 \ref{formal-spaces-lemma-space-to-formal-space}），
可以像前面一样使用《自举》之定义
\ref{bootstrap-definition-property-transformation}。

\begin{example}
\label{formal-spaces-example-david-hansen}
设 $(A, \mathfrak m, \kappa)$ 是赋值环，对某个非零
$\pi \in \mathfrak m$，它关于 $(\pi)$-进拓扑完备。
再假设 $\mathfrak m$ 不是有限生成的。一个例子是
$A = \mathcal{O}_{\mathbf{C}_p}$、$\pi = p$，其中
$\mathcal{O}_{\mathbf{C}_p}$ 是 $p$-进复数域 $\mathbf{C}_p$
的整数环（$\mathbf{C}_p$ 是 $\mathbf{Q}_p$ 的代数闭包的完备化）。
另一个例子是
$$
A =
\left\{
\sum\nolimits_{\alpha \in \mathbf{Q},\ \alpha \geq 0} a_\alpha t^\alpha
\middle|
\begin{matrix}
a_\alpha \in \kappa\text{，并且对所有 }n\text{，只有有限多个} \\
\alpha \leq n\text{ 对应的 }a_\alpha\text{ 非零}
\end{matrix}
\right\}
$$
以及 $\pi = t$。则 $X = \text{Spf}(A)$ 是仿射形式代数空间，
$\Spec(\kappa) \to X$ 是增厚，对应于弱定义理想
$\mathfrak m \subset A$，但后者不是定义理想。
\end{example}

\begin{remark}[弱定义理想]
\label{formal-spaces-remark-weak-ideals-of-definition}
设 $\mathfrak X$ 是 McQuillan 意义下的形式概形，见注
\ref{formal-spaces-remark-mcquillan}。$\mathfrak X$ 的一个{\it 弱定义理想}
是理想层 $\mathcal{I} \subset \mathcal{O}_\mathfrak X$，使得对每个仿射形式开子概形
$\mathfrak U \subset \mathfrak X$，理想
$\mathcal{I}(\mathfrak U) \subset \mathcal{O}_\mathfrak X(\mathfrak U)$
都是弱可容许拓扑环 $\mathcal{O}_\mathfrak X(\mathfrak U)$ 的弱定义理想。
只须在某个仿射开覆盖的成员上检验该条件。有一一对应
$$
\{\mathfrak X\text{ 的弱定义理想}\}
\leftrightarrow
\{\text{如上所述的增厚 }i : Z \to h_\mathfrak X\}
$$
该对应把 $\mathcal{I}$ 映到概形
$Z = (\mathfrak X, \mathcal{O}_\mathfrak X/\mathcal{I})$
及其到 $\mathfrak X$ 的显然态射。{\it 弱定义理想基本系}是弱定义理想族
$\mathcal{I}_\lambda$，使得对每个仿射形式开子概形
$\mathfrak U \subset \mathfrak X$，理想
$$
I_\lambda = \mathcal{I}_\lambda(\mathfrak U) \subset
A = \Gamma(\mathfrak U, \mathcal{O}_\mathfrak X)
$$
组成弱可容许拓扑环 $A$ 的一族弱定义理想基本系。
只须在某个仿射开覆盖的成员上检验。由此得到：与 McQuillan 形式概形
$\mathfrak X$ 相伴的形式代数空间 $h_\mathfrak X$ 是引理
\ref{formal-spaces-lemma-colimit-is-formal} 中那样的概形余极限，当且仅当
$\mathfrak X$ 存在一族弱定义理想基本系。
\end{remark}

\begin{remark}[定义理想]
\label{formal-spaces-remark-ideals-of-definition}
设 $\mathfrak X$ 是 EGA 意义下的形式概形。$\mathfrak X$ 的一个{\it 定义理想}
是理想层 $\mathcal{I} \subset \mathcal{O}_\mathfrak X$，使得对每个仿射形式开子概形
$\mathfrak U \subset \mathfrak X$，理想
$\mathcal{I}(\mathfrak U) \subset \mathcal{O}_\mathfrak X(\mathfrak U)$
都是可容许拓扑环 $\mathcal{O}_\mathfrak X(\mathfrak U)$ 的定义理想。
只须在某个仿射开覆盖的成员上检验该条件。这里{\bf 并不}像注
\ref{formal-spaces-remark-weak-ideals-of-definition} 那样得到定义理想与增厚
$Z \to h_\mathfrak X$ 之间的同一对应；例
\ref{formal-spaces-example-david-hansen} 给出了一个例子。
{\it 定义理想基本系}是定义理想族 $\mathcal{I}_\lambda$，
使得对每个仿射形式开子概形 $\mathfrak U \subset \mathfrak X$，理想
$$
I_\lambda = \mathcal{I}_\lambda(\mathfrak U) \subset
A = \Gamma(\mathfrak U, \mathcal{O}_\mathfrak X)
$$
组成可容许拓扑环 $A$ 的一族定义理想基本系。
只须在某个仿射开覆盖的成员上检验。假设 $\mathfrak X$ 拟紧，
而 $\{\mathcal{I}_\lambda\}_{\lambda \in \Lambda}$ 是一族弱定义理想基本系。
若 $A$ 是可容许拓扑环，则所有充分小的开理想都是定义理想
（即包含在某个定义理想中的任意开理想都是定义理想）。
所以，因为只需在一个仿射开覆盖的有限多个成员上检验，
对充分大的 $\lambda$，$\mathcal{I}_\lambda$ 是定义理想。
利用注 \ref{formal-spaces-remark-weak-ideals-of-definition} 中的讨论，得到：
与 EGA 意义下拟紧形式概形 $\mathfrak X$ 相伴的形式代数空间
$h_\mathfrak X$ 是引理 \ref{formal-spaces-lemma-colimit-is-formal} 中那样的概形余极限，
当且仅当 $\mathfrak X$ 存在一族定义理想基本系。
\end{remark}






\section{沿闭子集的完备化}
\label{formal-spaces-section-completion}

\noindent
我们的形式代数空间概念很适合沿闭子集取完备化。

\begin{lemma}
\label{formal-spaces-lemma-completion-affine-is-affine-formal-algebraic-space}
设 $S$ 是概形，$X$ 是 $S$ 上的仿射概形，$T \subset |X|$ 是闭子集。
则函子
$$
(\Sch/S)_{fppf} \longrightarrow \textit{Sets},\quad
U \longmapsto \{f : U \to X \mid f(|U|) \subset T\}
$$
是 McQuillan 仿射形式代数空间。
\end{lemma}

\begin{proof}
设 $X = \Spec(A)$，而 $T$ 对应于根式理想 $I \subset A$。
设 $U = \Spec(B)$ 是 $S$ 上的仿射概形，$f : U \to X$ 是 $F(U)$
的元素。则 $f$ 对应于环映射 $\varphi : A \to B$，使得 $B$
的每个素理想都包含 $\varphi(I)B$。所以 $\varphi(I)$ 的每个元素在 $B$
中幂零，见《代数》之引理 \ref{algebra-lemma-Zariski-topology}。
% P10-E0310
令 $J = \Ker(\varphi)$，可知 $I/J$ 是 $A/J$ 中的局部幂零理想。
等价地，$V(J) = V(I) = T$。换言之，引理中的函子等于
$\colim \Spec(A/J)$，其中余极限遍历满足 $V(J) = T$ 的理想 $J$。
因此该函子是仿射形式代数空间。它是 McQuillan 的
（定义 \ref{formal-spaces-definition-types-affine-formal-algebraic-space}），
因为映射 $A \to A/J$ 是满射，从而
$A^\wedge = \lim A/J \to A/J$ 是满射，见引理
\ref{formal-spaces-lemma-mcquillan-affine-formal-algebraic-space}。
\end{proof}

\begin{lemma}
\label{formal-spaces-lemma-completion-is-formal-algebraic-space}
设 $S$ 是概形，$X$ 是 $S$ 上的代数空间，$T \subset |X|$ 是闭子集。
则函子
$$
(\Sch/S)_{fppf} \longrightarrow \textit{Sets},\quad
U \longmapsto \{f : U \to X \mid f(|U|) \subset T\}
$$
是形式代数空间。
\end{lemma}

\begin{proof}
以 $F$ 表示该函子。设 $\{U_i \to U\}$ 是 fppf 覆盖。
则 $\coprod |U_i| \to |U|$ 是满射。由于 $X$ 是 fppf 层，
可知 $F$ 是 fppf 层。

\medskip\noindent
选取 étale 覆盖 $\{g_i : X_i \to X\}$，使每个 $X_i$ 都仿射，且对所有 $i$，
见《空间的性质》之引理
\ref{spaces-properties-lemma-cover-by-union-affines}。
态射 $F \times_X X_i \to F$ 是 étale 的
（见《空间》之引理
\ref{spaces-lemma-base-change-representable-transformations-property}），
而映射 $\coprod F \times_X X_i \to F$ 是层的满射。
因此，只须证明 $F \times_X X_i$ 是仿射形式代数空间。
$F \times_X X_i$ 的一个 $U$-值点是态射 $U \to X_i$，
其像包含在闭子集 $g_i^{-1}(T) \subset |X_i|$ 中。
所以结论由引理
\ref{formal-spaces-lemma-completion-affine-is-affine-formal-algebraic-space} 得出。
\end{proof}

\begin{definition}
\label{formal-spaces-definition-completion}
设 $S$ 是概形，$X$ 是 $S$ 上的代数空间，$T \subset |X|$ 是闭子集。
引理 \ref{formal-spaces-lemma-completion-is-formal-algebraic-space} 中的形式代数空间
称为 $X$ {\it 沿 $T$ 的完备化}。
\end{definition}

\noindent
\cite[Chapter I, Section 10.8]{EGA} 用记号 $X_{/T}$ 表示该完备化，
我们有时也使用这一记号。设 $f : X \to X'$ 是概形 $S$ 上代数空间的态射。
假设 $T \subset |X|$ 和 $T' \subset |X'|$ 是闭子集，且
$|f|(T) \subset T'$。则显然，$f$ 定义完备化之间的形式代数空间态射
$$
X_{/T} \longrightarrow X'_{/T'}
$$


\begin{lemma}
\label{formal-spaces-lemma-map-completions-representable}
设 $S$ 是概形，$f : X' \to X$ 是 $S$ 上代数空间的态射。
设 $T \subset |X|$ 是闭子集，并令
$T' = |f|^{-1}(T) \subset |X'|$。则
$$
\xymatrix{
X'_{/T'} \ar[r] \ar[d] & X' \ar[d]^f \\
X_{/T} \ar[r] & X
}
$$
是层的笛卡尔图表。特别地，态射
$X'_{/T'} \to X_{/T}$ 可由代数空间表示。
\end{lemma}

\begin{proof}
事实上，设 $Y \to X$ 是从概形到 $X$ 的态射，使 $|Y|$ 映入 $T$。
则
% P10-E0311
$Y \times_X X' \to X$ 是代数空间的态射，使
$|Y \times_X X'|$ 映入 $T'$。所以函子
$Y \times_{X_{/T}} X'_{/T'}$ 由 $Y \times_X X'$ 表示，
从而引理成立。
\end{proof}

\begin{lemma}
\label{formal-spaces-lemma-reduction-completion}
设 $S$ 是概形，$X$ 是 $S$ 上的代数空间，$T \subset |X|$ 是闭子集。
$X$ 沿 $T$ 的完备化 $X_{/T}$ 的约化 $(X_{/T})_{red}$，
是 $X$ 中与 $T$ 对应的既约诱导闭子空间 $Z$。
\end{lemma}

\begin{proof}
由引理 \ref{formal-spaces-lemma-reduction-formal-algebraic-space}、《空间的性质》之定义
\ref{spaces-properties-definition-reduced-induced-space}
（它使用《空间的性质》之引理
\ref{spaces-properties-lemma-reduced-closed-subspace} 构造 $Z$）
以及 $X_{/T}$ 的定义可知，$Z$ 和 $(X_{/T})_{red}$ 是由同一映射性质
% P10-E0312
刻画既约代数空间：源为既约代数空间的态射 $g : Y \to X$
分解经过它们，当且仅当 $|Y|$ 映入 $T \subset |X|$。
\end{proof}

\begin{lemma}
\label{formal-spaces-lemma-affine-formal-completion-types}
设 $S$ 是概形，$X = \Spec(A)$ 是 $S$ 上的仿射概形，
$T \subset X$ 是闭子集，$X_{/T}$ 是 $X$ 沿 $T$ 的形式完备化。
\begin{enumerate}
\item 若 $X \setminus T$ 拟紧，即 $T$ 可构造，则 $X_{/T}$ adic*。
\item 若对某个有限生成理想 $I \subset A$ 有 $T = V(I)$，
则 $X_{/T} = \text{Spf}(A^\wedge)$，其中 $A^\wedge$ 是 $A$
的 $I$-进完备化。
\item 若 $X$ Noether，则 $X_{/T}$ Noether。
\end{enumerate}
\end{lemma}

\begin{proof}
由《代数》之引理 \ref{algebra-lemma-qc-open}，若 (1) 成立，
则可以找到 (2) 中的理想 $I \subset A$。若 (3) 成立，
也可以找到 (2) 中的理想 $I \subset A$。此外，由《代数》之引理
\ref{algebra-lemma-completion-Noetherian-Noetherian}，Noether 环的完备化仍 Noether。
综上，只须证明 (2)。

\medskip\noindent
(2) 的证明。设 $I = (f_1, \ldots, f_r) \subset A$ 切出 $T$。
若 $Z = \Spec(B)$ 是仿射概形，$g : Z \to X$ 是满足
$g(Z) \subset T$（集合论意义下）的态射，则对每个 $i$，
$g^\sharp(f_i)$ 在 $B$ 中幂零。所以对某个 $n$，$I^n$ 在 $B$ 中映为零。
因此可见
$X_{/T} = \colim \Spec(A/I^n) = \text{Spf}(A^\wedge)$。
\end{proof}

\noindent
下列引理由 Ofer Gabber 给出。

\begin{lemma}
\label{formal-spaces-lemma-completion-countably-indexed}
\begin{reference}
Ofer Gabber 于 2014 年 9 月 11 日的电子邮件。
\end{reference}
设 $S$ 是概形，$X = \Spec(A)$ 是 $S$ 上的仿射概形，
$T \subset X$ 是闭子概形。
\begin{enumerate}
\item 若形式完备化 $X_{/T}$ 是可数指标的，并且存在可数多个
$f_1,f_2,f_3,\ldots \in A$ 使
$T = V(f_1,f_2,f_3,\ldots)$，则 $X_{/T}$ adic*。
\item 若去掉 $T$ 可由 $X$ 中可数多个函数切出这一假设，(1) 的结论不成立。
\end{enumerate}
\end{lemma}

\begin{proof}
$X_{/T}$ 是可数指标的这一假设表示，存在理想序列
$$
A \supset J_1 \supset J_2 \supset J_3 \supset \ldots
$$
满足 $V(J_n) = T$，并且
% P10-E0313
每个满足 $V(J) = T$ 的理想 $J \subset A$ 都有某个 $n$ 使 $J \supset J_n$。

\medskip\noindent
为构造 (2) 的例子，令 $\omega_1$ 为第一个不可数序数。
设 $k$ 是域，$A$ 是由 $x_\alpha$（$\alpha \in \omega_1$）和
$y_{\alpha\beta}$（$\alpha \in \beta \in \omega_1$）生成的 $k$-代数，
满足关系 $x_\alpha = y_{\alpha\beta}x_\beta$。
令 $T = V(x_\alpha)$，$J_n = (x_\alpha^n)$。
若理想 $J \subset A$ 满足 $V(J) = T$，则对某个
$n_\alpha \geq 1$ 有 $x_\alpha^{n_\alpha} \in J$。
集合 $\{\alpha \mid n_\alpha = n\}$ 中至少有一个必须在 $\omega_1$ 中无界。
于是这些关系蕴含 $J_n \subset J$。

\medskip\noindent
为看出 (2) 成立，现在只须证明
$A^\wedge = \lim A/J_n$ 不是关于某个有限生成理想完备的环。
对 $\gamma \in \omega_1$，令 $A_\gamma$ 为 $A$ 关于如下理想的商：
该理想由 $x_\alpha$（$\alpha \in \gamma$）以及
$y_{\alpha\beta}$（$\alpha \in \gamma$）生成。
因为 $A/J_1$ 既约，$\lim A/J_n$ 的每个拓扑幂零元素 $f$ 都属于
$J_1^\wedge = \lim J_1/J_n$。这表示 $f$ 是只涉及可数多个生成元的无穷级数。
所以对某个 $\gamma$，$f$ 在
$A_\gamma^\wedge = \lim A_\gamma/J_nA_\gamma$ 中消失。
注意，由引理 \ref{formal-spaces-lemma-ses}，
$A^\wedge \to A_\gamma^\wedge$ 连续且开。
% P10-E0314
若 $A^\wedge$ 上的拓扑是某个有限生成理想
$I \subset A^\wedge$ 的 $I$-进拓扑，则 $I$ 会在某个
$A_\gamma^\wedge$ 中映为零。这将意味着 $A_\gamma^\wedge$ 离散，
但事实并非如此，因为有一个满的连续开映射（开性由引理
\ref{formal-spaces-lemma-ses}）
$A_\gamma^\wedge \to k[[t]]$，它对满足
$\gamma = \alpha$ 或 $\gamma \in \alpha$ 的指标由
$x_\alpha \mapsto t$、$y_{\alpha\beta} \mapsto 1$ 给出。

\medskip\noindent
证明 (1) 之前，先证明如下事实：若 $I \subset A^\wedge$ 是有限生成理想，
其闭包 $\bar I$ 是开的，则 $I = \bar I$。
由于 $V(J_n^2) = T$，存在 $m$ 使 $J_n^2 \supset J_m$。
所以，取子序列后可以假设对所有 $n$ 都有
$J_n^2 \supset J_{n+1}$。令
$J_n^\wedge = \lim_{k \geq n} J_n/J_k \subset A^\wedge$。
由于闭包 $\bar I = \bigcap (I + J_n^\wedge)$
（引理 \ref{formal-spaces-lemma-closed}）是开的，可见存在 $m$，使得对所有
$n \geq m$ 有 $I + J_n^\wedge \supset J_m^\wedge$。固定这样的 $m$。
对所有 $n \geq m+1$，有
$$
J_{n - 1}^\wedge I + J_{n + 1}^\wedge \supset
J_{n - 1}^\wedge (I + J_{n + 1}^\wedge) \supset
J_{n - 1}^\wedge J_m^\wedge
$$
事实上，第一个包含显然，第二个已在上面证明。
因为 $J_{n-1}J_m \supset J_{n-1}^2 \supset J_n$，
这些包含表明 $J_n$ 在 $A^\wedge$ 中的像包含在理想
$J_{n-1}^\wedge I + J_{n+1}^\wedge$ 中。因为该理想是开的，得到
$$
J_{n - 1}^\wedge I + J_{n + 1}^\wedge \supset J_n^\wedge.
$$
设 $I = (g_1,\ldots,g_t)$。取 $f \in J_{m+1}^\wedge$。
利用最后显示的包含（它对所有 $n \geq m+1$ 都成立），
可以对 $c \geq 0$ 归纳写成
$$
f = \sum f_{i, c} g_i \mod J_{m + 1+ c}^\wedge
$$
其中 $f_{i,c} \in J_m^\wedge$，且
$f_{i,c} \equiv f_{i,c-1} \bmod J_{m+c}^\wedge$。
由此 $IJ_m^\wedge \supset J_{m+1}^\wedge$。
再结合 $I + J_{m+1}^\wedge \supset J_m^\wedge$，得到 $I$ 开。

\medskip\noindent
(1) 的证明。假设 $T = V(f_1,f_2,f_3,\ldots)$。
令 $I_m \subset A^\wedge$ 为由 $f_1,\ldots,f_m$ 生成的理想。
分两种情形。

\medskip\noindent
情形 I：对某个 $m$，$I_m$ 的闭包是开的。
由前一段的结果，$I_m$ 开。按构造，对任意 $n$ 有
$(J_n)^2 \supset J_{n+1}$，所以 $(J_n^\wedge)^2$ 的闭包包含
$J_{n+1}^\wedge$，因而是开的。取充分大的 $n$，可知
$A^\wedge$ 中任意两个开理想的乘积的闭包是开的。
用关于 $k$ 的归纳证明，对 $k \ge 1$，$I_m^k$ 开。
$k=1$ 的情形是情形 I 中对 $m$ 的假设。若 $k>1$，并假设
$I_m^{k-1}$ 开，则 $I_m^k = I_m^{k-1}\cdot I_m$ 是两个开理想的乘积，
因而其闭包是开的。但 $I_m^k$ 有限生成，所以由前一段
（应用于 $I=I_m^k$），$I_m^k$ 是开的，可以继续关于 $k$ 的归纳。
由于 $I_m$ 的每个元素都拓扑幂零，可知 $I_m$ 是定义理想。
这证明 $A^\wedge$ 是有有限生成定义理想的 adic 环；即 $X_{/T}$ adic*。

\medskip\noindent
情形 II：对所有 $m$，$I_m$ 的闭包 $\bar I_m$ 都不开。
则 $A^\wedge/\bar I_m$ 上的拓扑不离散。这表示可以选取
$\phi(m) \geq m$，使
$$
\Im(J_{\phi(m)} \to A/(f_1, \ldots, f_m)) \not =
\Im(J_{\phi(m) + 1} \to A/(f_1, \ldots, f_m))
$$
这里使用了
$A^\wedge/(\bar I_m + J_n^\wedge)
= A/((f_1,\ldots,f_m)+J_n)$。
选取指数 $e_i>0$，使得当 $0<m<i$ 时
$f_i^{e_i} \in J_{\phi(m)+1}$。令
$J = (f_1^{e_1}, f_2^{e_2}, f_3^{e_3},\ldots)$。
则 $V(J)=T$。我们断言对所有 $n$ 都有 $J \not\supset J_n$，
这给出矛盾，证明情形 II 不会发生。事实上，$J$ 在
$A/(f_1,\ldots,f_m)$ 中的像包含在 $J_{\phi(m)+1}$ 的像中，
而后者真包含于 $J_m$ 的像。
\end{proof}















\section{纤维积}
\label{formal-spaces-section-fibre-products}

\noindent
这是关于形式代数空间纤维积的必备一节。

\begin{lemma}
\label{formal-spaces-lemma-etale-covering-by-formal-algebraic-spaces}
设 $S$ 是概形，$\{X_i \to X\}_{i \in I}$ 是
$(\Sch/S)_{fppf}$ 上的一族层映射。假设：
(a) $X_i$ 是 $S$ 上的形式代数空间；
(b) $X_i \to X$ 可由代数空间表示且为 étale；
(c) $\coprod X_i \to X$ 是层的满射。
则 $X$ 是 $S$ 上的形式代数空间。
\end{lemma}

\begin{proof}
对每个 $i$，选取定义 \ref{formal-spaces-definition-formal-algebraic-space}
中的 $\{X_{ij} \to X_i\}_{j \in J_i}$。则
$\{X_{ij} \to X\}_{i \in I, j \in J_i}$ 是定义
\ref{formal-spaces-definition-formal-algebraic-space} 中 $X$ 所需的一族映射。
\end{proof}

\begin{lemma}
\label{formal-spaces-lemma-fibre-products-general}
设 $S$ 是概形，$X,Y$ 是 $S$ 上的形式代数空间，$Z$ 是对角可由代数空间表示的层。
设 $X \to Z$ 和 $Y \to Z$ 是层映射。则 $X \times_Z Y$ 是形式代数空间。
\end{lemma}

\begin{proof}
选取定义 \ref{formal-spaces-definition-formal-algebraic-space} 中的
$\{X_i \to X\}$ 和 $\{Y_j \to Y\}$。则
$\{X_i \times_Z Y_j \to X \times_Z Y\}$ 是一族可由代数空间表示且为 étale 的映射。
所以引理 \ref{formal-spaces-lemma-etale-covering-by-formal-algebraic-spaces} 告诉我们，
只须在 $X$ 和 $Y$ 是仿射形式代数空间时证明
$X \times_Z Y$ 是形式代数空间。

\medskip\noindent
假设 $X$ 和 $Y$ 是仿射形式代数空间。按定义
\ref{formal-spaces-definition-affine-formal-algebraic-space} 写
$X = \colim X_\lambda$ 和 $Y = \colim Y_\mu$。则
$X \times_Z Y = \colim X_\lambda \times_Z Y_\mu$。
每个 $X_\lambda \times_Z Y_\mu$ 都是代数空间。
对 $\lambda \leq \lambda'$ 和 $\mu \leq \mu'$，态射
$$
X_\lambda \times_Z Y_\mu \to
X_\lambda \times_Z Y_{\mu'} \to
X_{\lambda'} \times_Z Y_{\mu'}
$$
是增厚基变换的复合，因而是增厚。因此应用引理
\ref{formal-spaces-lemma-colimit-is-formal} 即得结论。
\end{proof}

\begin{lemma}
\label{formal-spaces-lemma-fibre-products}
设 $S$ 是概形。$S$ 上形式代数空间的范畴有纤维积。
\end{lemma}

\begin{proof}
这是引理 \ref{formal-spaces-lemma-fibre-products-general} 的特殊情形，
因为形式代数空间的对角可表示，见引理
\ref{formal-spaces-lemma-diagonal-formal-algebraic-space}。
\end{proof}

\begin{lemma}
\label{formal-spaces-lemma-reduction-fibre-products}
设 $S$ 是概形，$X \to Z$ 和 $Y \to Z$ 是 $S$ 上形式代数空间的态射。则
$(X \times_Z Y)_{red} = (X_{red} \times_{Z_{red}} Y_{red})_{red}$。
\end{lemma}

\begin{proof}
这来自引理 \ref{formal-spaces-lemma-reduction-formal-algebraic-space} 中约化的泛性质。
\end{proof}

\noindent
我们已经证明了下列引理（当时尚不知道纤维积存在）。

\begin{lemma}
\label{formal-spaces-lemma-diagonal-morphism-formal-algebraic-spaces}
设 $S$ 是概形，$f : X \to Y$ 是 $S$ 上形式代数空间的态射。
对角态射 $\Delta : X \to X \times_Y X$ 可表示（由概形表示）、
为单态射、局部拟有限、局部有限型且分离。
\end{lemma}

\begin{proof}
设 $T$ 是概形，并给定态射 $T \to X \times_Y X$。则
$$
T \times_{(X \times_Y X)} X = T \times_{(X \times_S X)} X
$$
所以结论立即来自引理 \ref{formal-spaces-lemma-diagonal-formal-algebraic-space}。
\end{proof}






\section{形式代数空间的分离公理}
\label{formal-spaces-section-separation}

\noindent
本节讨论形式代数空间的“绝对”分离条件。形式代数空间态射的分离条件将在后文讨论。

\begin{lemma}
\label{formal-spaces-lemma-characterize-quasi-separated}
设 $S$ 是概形，$X$ 是 $S$ 上的形式代数空间。下列条件等价：
\begin{enumerate}
\item $X$ 的约化（引理 \ref{formal-spaces-lemma-reduction-formal-algebraic-space}）
是拟分离代数空间；
\item 对 $U \to X$、$V \to X$，若 $U,V$ 是拟紧概形，
则纤维积 $U \times_X V$ 拟紧；
\item 对 $U \to X$、$V \to X$，若 $U,V$ 仿射，
则纤维积 $U \times_X V$ 拟紧。
\end{enumerate}
\end{lemma}

\begin{proof}
由引理 \ref{formal-spaces-lemma-diagonal-formal-algebraic-space}，
$U \times_X V$ 是概形。令 $U_{red},V_{red},X_{red}$ 分别为 $U,V,X$ 的约化。
则
$$
U_{red} \times_{X_{red}} V_{red} = U_{red} \times_X V_{red}
\to U \times_X V
$$
是概形的增厚。牢记代数空间的对应引理，(1) 与 (2) 的等价性由此显然，
见《空间的性质》之引理
\ref{spaces-properties-lemma-characterize-quasi-separated}。
我们略去 (2) 与 (3) 等价性的证明。
\end{proof}

\begin{lemma}
\label{formal-spaces-lemma-characterize-separated}
设 $S$ 是概形，$X$ 是 $S$ 上的形式代数空间。下列条件等价：
\begin{enumerate}
\item $X$ 的约化（引理 \ref{formal-spaces-lemma-reduction-formal-algebraic-space}）
是分离代数空间；
\item 对 $U \to X$、$V \to X$，若 $U,V$ 仿射，则纤维积
$U \times_X V$ 仿射，并且
$$
\mathcal{O}(U) \otimes_\mathbf{Z} \mathcal{O}(V)
\longrightarrow
\mathcal{O}(U \times_X V)
$$
是满射。
\end{enumerate}
\end{lemma}

\begin{proof}
若 (2) 成立，则把《空间的性质》之引理
% P10-E0315
\ref{spaces-properties-lemma-characterize-quasi-separated}
应用于态射 $U \to X_{red}$ 和 $V \to X_{red}$，其中 $U,V$ 仿射，
并使用 $U \times_{X_{red}}V = U \times_X V$，可知
$X_{red}$ 是分离代数空间。

\medskip\noindent
假设 (1)。设 $U \to X$ 和 $V \to X$ 如 (2) 所述。
由引理 \ref{formal-spaces-lemma-diagonal-formal-algebraic-space}，
$U \times_X V$ 是概形。令 $U_{red},V_{red},X_{red}$ 分别为 $U,V,X$ 的约化。
则
$$
U_{red} \times_{X_{red}} V_{red} = U_{red} \times_X V_{red}
\to U \times_X V
$$
是概形的增厚。于是
$(U \times_X V)_{red} = (U_{red} \times_{X_{red}} V_{red})_{red}$。
特别地，$(U \times_X V)_{red}$ 是仿射概形，并且
$$
\mathcal{O}(U) \otimes_\mathbf{Z} \mathcal{O}(V)
\longrightarrow
\mathcal{O}((U \times_X V)_{red})
$$
是满射，见《空间的性质》之引理
% P10-E0315
\ref{spaces-properties-lemma-characterize-quasi-separated}。
于是由《空间的极限》之命题
\ref{spaces-limits-proposition-affine}，$U \times_X V$ 仿射。
另一方面，仿射概形的态射 $U \times_X V \to U \times V$ 是复合
$$
U \times_X V = X \times_{(X \times_S X)} (U \times_S V)
\to U \times_S V \to U \times V
$$
第一个态射是单态射且局部有限型
（引理 \ref{formal-spaces-lemma-diagonal-formal-algebraic-space}）。
第二个态射是浸入（《概形》之引理
\ref{schemes-lemma-fibre-product-after-map}）。
所以该复合是局部有限型的单态射。另一方面，它是整的，因为底层既约仿射概形上的映射
由上述结论是闭浸入，从而普遍闭（使用《态射》之引理
\ref{morphisms-lemma-integral-universally-closed}）。因此环映射
$$
\mathcal{O}(U) \otimes_\mathbf{Z} \mathcal{O}(V)
\longrightarrow
\mathcal{O}(U \times_X V)
$$
是有限型的整满射态射，因而有限，进而满
（使用《态射》之引理 \ref{morphisms-lemma-finite-integral}
和《代数》之引理 \ref{algebra-lemma-finite-epimorphism-surjective}）。
\end{proof}

\begin{definition}
\label{formal-spaces-definition-separated}
设 $S$ 是概形，$X$ 是 $S$ 上的形式代数空间。我们称
\begin{enumerate}
\item 若引理 \ref{formal-spaces-lemma-characterize-quasi-separated} 的等价条件成立，
则 $X$ {\it 拟分离}；
\item 若引理 \ref{formal-spaces-lemma-characterize-separated} 的等价条件成立，
则 $X$ {\it 分离}。
\end{enumerate}
\end{definition}

\noindent
下列引理特别蕴含：弱可容许拓扑环的完备张量积仍是弱可容许拓扑环。

\begin{lemma}
\label{formal-spaces-lemma-fibre-product-affines-over-separated}
设 $S$ 是概形，$X \to Z$ 和 $Y \to Z$ 是 $S$ 上形式代数空间的态射。
假设 $Z$ 分离。
\begin{enumerate}
\item 若 $X$ 和 $Y$ 是仿射形式代数空间，则 $X \times_Z Y$ 也是。
\item 若 $X$ 和 $Y$ 是 McQuillan 仿射形式代数空间，则 $X \times_Z Y$ 也是。
\item 若 $X,Y,Z$ 是分别对应于弱可容许拓扑 $S$-代数
$A,B,C$ 的 McQuillan 仿射形式代数空间，则 $X \times_Z Y$
对应于 $A \widehat{\otimes}_C B$。
\end{enumerate}
\end{lemma}

\begin{proof}
按定义 \ref{formal-spaces-definition-affine-formal-algebraic-space} 写
$X = \colim X_\lambda$ 和 $Y = \colim Y_\mu$。则
$X \times_Z Y = \colim X_\lambda \times_Z Y_\mu$。
由于 $Z$ 分离，这些纤维积仿射，所以 (1) 成立。假设
% P10-E0316
$X$ 和 $Y$ 对应于弱可容许拓扑 $S$-代数 $A$ 和 $B$，并且
$X_\lambda = \Spec(A/I_\lambda)$、
$Y_\mu = \Spec(B/J_\mu)$。则
$$
X_\lambda \times_Z Y_\mu \to
X_\lambda \times Y_\mu \to \Spec(A \otimes B)
$$
是闭浸入。因此，引理
\ref{formal-spaces-lemma-mcquillan-affine-formal-algebraic-space} 的一个条件成立，
从而 $X \times_Z Y$ 是 McQuillan 的。
若 $Z$ 也为对应于 $C$ 的 McQuillan 空间，则
$$
X_\lambda \times_Z Y_\mu = \Spec(A/I_\lambda \otimes_C B/J_\mu)
$$
所以可见与 $X \times_Z Y$ 对应的弱可容许拓扑环是完备张量积
（见定义 \ref{formal-spaces-definition-toplogy-tensor-product}）。
\end{proof}

\begin{lemma}
\label{formal-spaces-lemma-separated-from-separated}
设 $S$ 是概形，$X$ 是 $S$ 上的形式代数空间。
设 $U \to X$ 是态射，其中 $U$ 是 $S$ 上的分离代数空间。
则 $U \to X$ 分离。
\end{lemma}

\begin{proof}
该陈述有意义，因为 $U \to X$ 可由代数空间表示
（引理 \ref{formal-spaces-lemma-space-to-formal-space}）。
设 $T$ 是概形并给定态射 $T \to X$。须证明
$U \times_X T \to T$ 分离。由于
$U \times_X T \to U \times_S T$ 是单态射，只须证明
$U \times_S T \to T$ 分离。后者是 $U \to S$ 的基变换，所以结论成立。
上述论证使用了《空间的态射》之引理
\ref{spaces-morphisms-lemma-base-change-separated}、
\ref{spaces-morphisms-lemma-composition-separated}、
\ref{spaces-morphisms-lemma-monomorphism-separated} 和
\ref{spaces-morphisms-lemma-separated-implies-morphism-separated}。
\end{proof}




\section{拟紧形式代数空间}
\label{formal-spaces-section-quasi-compact}

\noindent
下面刻画拟紧形式代数空间。

\begin{lemma}
\label{formal-spaces-lemma-characterize-quasi-compact}
设 $S$ 是概形，$X$ 是 $S$ 上的形式代数空间。下列条件等价：
\begin{enumerate}
\item $X$ 的约化（引理 \ref{formal-spaces-lemma-reduction-formal-algebraic-space}）
是拟紧代数空间；
\item 可以找到定义 \ref{formal-spaces-definition-formal-algebraic-space}
中的 $\{X_i \to X\}_{i \in I}$，其中 $I$ 有限；
\item 存在可由代数空间表示、étale 且满的态射 $Y \to X$，
其中 $Y$ 是仿射形式代数空间。
\end{enumerate}
\end{lemma}

\begin{proof}
略。
\end{proof}

\begin{definition}
\label{formal-spaces-definition-quasi-compact}
设 $S$ 是概形，$X$ 是 $S$ 上的形式代数空间。
若引理 \ref{formal-spaces-lemma-characterize-quasi-compact} 的等价条件成立，
则称 $X$ {\it 拟紧}。
\end{definition}

\begin{lemma}
\label{formal-spaces-lemma-characterize-quasi-compact-morphism}
设 $S$ 是概形，$f : X \to Y$ 是 $S$ 上形式代数空间的态射。
下列条件等价：
\begin{enumerate}
\item 约化之间的诱导映射
$f_{red} : X_{red} \to Y_{red}$
（引理 \ref{formal-spaces-lemma-reduction-formal-algebraic-space}）是代数空间的拟紧态射；
\item 对每个拟紧概形 $T$ 和态射 $T \to Y$，纤维积
$X \times_Y T$ 是拟紧形式代数空间；
\item 对每个仿射概形 $T$ 和态射 $T \to Y$，纤维积
$X \times_Y T$ 是拟紧形式代数空间；
\item 存在定义 \ref{formal-spaces-definition-formal-algebraic-space} 中的覆盖
$\{Y_j \to Y\}$，使每个 $X \times_Y Y_j$ 都是拟紧形式代数空间。
\end{enumerate}
\end{lemma}

\begin{proof}
略。
\end{proof}

\begin{definition}
\label{formal-spaces-definition-quasi-compact-morphism}
设 $S$ 是概形，$f : X \to Y$ 是 $S$ 上形式代数空间的态射。
若引理 \ref{formal-spaces-lemma-characterize-quasi-compact-morphism} 的等价条件成立，
则称 $f$ {\it 拟紧}。
\end{definition}

\noindent
当态射可由代数空间表示（特别地，当它可表示）时，
这与已有的概念一致。

\begin{lemma}
\label{formal-spaces-lemma-quasi-compact-representable}
设 $S$ 是概形，$f : X \to Y$ 是 $S$ 上形式代数空间的态射，
并且可由代数空间表示。则 $f$ 在定义
\ref{formal-spaces-definition-quasi-compact-morphism} 的意义下拟紧，当且仅当
$f$ 在《自举》之定义 \ref{bootstrap-definition-property-transformation}
的意义下拟紧。
\end{lemma}

\begin{proof}
这立即来自定义和引理 \ref{formal-spaces-lemma-characterize-quasi-compact-morphism}。
\end{proof}





\section{拟紧且拟分离的形式代数空间}
\label{formal-spaces-section-quasi-compact-quasi-separated}

\noindent
下列结果由 Yasuda 给出，见 \cite[Proposition 3.32]{Yasuda}。

\begin{lemma}
\label{formal-spaces-lemma-structure-quasi-compact-quasi-separated}
\begin{reference}
\cite[Proposition 3.32]{Yasuda}
\end{reference}
设 $S$ 是概形，$X$ 是 $S$ 上拟紧且拟分离的形式代数空间。
则对有向集 $\Lambda$ 上的某个代数空间系统
$(X_\lambda,f_{\lambda\mu})$，有 $X = \colim X_\lambda$，
其中每个 $f_{\lambda\mu} : X_\lambda \to X_\mu$ 都是增厚。
\end{lemma}

\begin{proof}
由引理 \ref{formal-spaces-lemma-characterize-quasi-compact}，可以选取仿射形式代数空间
$Y$ 以及可表示、满且 étale 的态射 $Y \to X$。
按定义 \ref{formal-spaces-definition-affine-formal-algebraic-space} 写
$Y = \colim Y_\lambda$。

\medskip\noindent
取 $\lambda \in \Lambda$。由引理 \ref{formal-spaces-lemma-presentation-representable}，
$Y_\lambda \times_X Y$ 是概形。$Y_\lambda \times_X Y$ 的约化
（引理 \ref{formal-spaces-lemma-reduction-formal-algebraic-space}）等于
$Y_{red} \times_{X_{red}}Y_{red}$ 的约化；后者拟紧，因为 $X$ 拟分离且
$Y_{red}$ 仿射。所以 $Y_\lambda \times_X Y$ 是拟紧概形。
因此存在 $\mu \geq \lambda$，使
$\text{pr}_2 : Y_\lambda \times_X Y \to Y$ 分解经过 $Y_\mu$，
见引理 \ref{formal-spaces-lemma-factor-through-thickening}。
令 $Z_\lambda$ 为态射
$\text{pr}_2 : Y_\lambda \times_X Y \to Y_\mu$ 的概形论像。
它与 $\mu$ 的选择无关，我们可以而且将把
$Z_\lambda \subset Y$ 看作态射
$\text{pr}_2 : Y_\lambda \times_X Y \to Y$ 的概形论像。
注意，$Z_\lambda$ 也等于态射
$\text{pr}_1 : Y \times_X Y_\lambda \to Y$ 的概形论像，
因为它同构于定义 $Z_\lambda$ 所用的态射。
我们断言，作为 $Y \times_X Y$ 的子函子，有
$Z_\lambda \times_X Y = Y \times_X Z_\lambda$。
事实上，因为 $Y \to X$ 是 étale 的，由《空间的态射》之引理
\ref{spaces-morphisms-lemma-quasi-compact-scheme-theoretic-image}，
$Z_\lambda \times_X Y$ 是态射
$$
\text{pr}_{13} = \text{pr}_1 \times \text{id}_Y :
Y \times_X Y_\lambda \times_X Y \longrightarrow Y \times_X Y
$$
的概形论像。同理，$Y \times_X Z_\lambda$ 是态射
$$
\text{pr}_{13} = \text{id}_Y \times \text{pr}_2 :
Y \times_X Y_\lambda \times_X Y \longrightarrow Y \times_X Y
$$
的概形论像。断言得证。于是
$R_\lambda = Z_\lambda \times_X Y = Y \times_X Z_\lambda$，
连同态射 $R_\lambda \to Z_\lambda \times_S Z_\lambda$，
定义一个 étale 等价关系。由此得到代数空间
$X_\lambda = Z_\lambda/R_\lambda$。按构造，图表
$$
\xymatrix{
Z_\lambda \ar[r] \ar[d] & Y \ar[d] \\
X_\lambda \ar[r] & X
}
$$
是笛卡尔的（因为 $X$ 是两个投影 $R=Y\times_XY\to Y$ 的余等化子，
$Z_\lambda\subset Y$ 在 $R$ 下不变，且 $R_\lambda$ 是 $R$ 到
$Z_\lambda$ 上的限制）。所以由《空间》之引理
\ref{spaces-lemma-morphism-sheaves-with-P-effective-descent-etale}，
$X_\lambda \to X$ 可表示且为闭浸入。另一方面，由于
$Y_\lambda \subset Z_\lambda$，有
$(X_\lambda)_{red} = X_{red}$；换言之，$X_\lambda \to X$ 是增厚。
最后，我们断言
$$
X = \colim X_\lambda
$$
我们有 $Y \times_X X_\lambda = Z_\lambda \supset Y_\lambda$。
每个态射 $T \to X$（其中 $T$ 是 $S$ 上的概形）都能 étale 局部地提升为到
$Y$ 的态射，而后者又能 étale 局部地提升为到某个 $Y_\lambda$ 的态射。
所以 $T \to X$ 在 $T$ 上 étale 局部地提升为到 $X_\lambda$ 的态射。
证明完成。
\end{proof}

\begin{remark}
\label{formal-spaces-remark-structure-quasi-compact-quasi-separated}
本注把引理 \ref{formal-spaces-lemma-structure-quasi-compact-quasi-separated}
的陈述和证明翻译为 EGA 意义下形式概形的语言。
考察注 \ref{formal-spaces-remark-ideals-of-definition}，可把引理翻译如下：
\begin{itemize}
\item[(*)] 每个拟紧且拟分离的形式概形都有一族定义理想基本系。
\end{itemize}
为证明这一点，先使用《概形的上同调》之引理
\ref{coherent-lemma-induction-principle} 的归纳原理
（为拟紧且拟分离的形式概形重新表述），把问题归约到如下情形：
$\mathfrak X = \mathfrak U \cup \mathfrak V$，其中
$\mathfrak U,\mathfrak V$ 是开形式子概形，$\mathfrak V$ 仿射，
且结果对 $\mathfrak U$、$\mathfrak V$ 和
$\mathfrak U\cap\mathfrak V$ 都成立。任取定义理想
$\mathcal{I}\subset\mathcal{O}_\mathfrak U$ 和
$\mathcal{J}\subset\mathcal{O}_\mathfrak V$。
根据 $\mathfrak U$ 和 $\mathfrak V$ 上都有定义理想基本系这一假设，
并利用 $\mathfrak U\cap\mathfrak V$ 拟紧，可以找到定义理想
$\mathcal{I}'\subset\mathcal{I}$ 和
$\mathcal{J}'\subset\mathcal{J}$，使得
$$
\mathcal{I}'|_{\mathfrak U \cap \mathfrak V} \subset
\mathcal{J}|_{\mathfrak U \cap \mathfrak V}
\quad\text{且}\quad
\mathcal{J}'|_{\mathfrak U \cap \mathfrak V} \subset
\mathcal{I}|_{\mathfrak U \cap \mathfrak V}
$$
设 $U \to U' \to \mathfrak U$ 和 $V \to V' \to \mathfrak V$
分别为定义理想
$\mathcal{I}'\subset\mathcal{I}\subset\mathcal{O}_\mathfrak U$
和 $\mathcal{J}'\subset\mathcal{J}\subset\mathcal{O}_\mathfrak V$
所确定的闭浸入。以 $\mathfrak U\cap V$ 表示 $V$ 的开子概形，
其底层拓扑空间为 $\mathfrak U\cap\mathfrak V$ 的底层拓扑空间。
由 $\mathcal{I}'$ 的选择，存在分解
$\mathfrak U\cap V\to U'$。类似地定义
$U\cap\mathfrak V$，它分解经过 $V'$。然后考虑
$$
Z_U = \text{下列态射的概形论像：}
U \amalg (\mathfrak U \cap V) \longrightarrow U'
$$
以及
$$
Z_V = \text{下列态射的概形论像：}
(U \cap \mathfrak V) \amalg V \longrightarrow V'
$$
由于取拟紧态射的概形论像与限制到开集交换
（《态射》之引理
\ref{morphisms-lemma-quasi-compact-scheme-theoretic-image}），可见
$Z_U\cap\mathfrak V = \mathfrak U\cap Z_V$。
所以 $Z_U$ 和 $Z_V$ 黏合成概形 $Z$，并带有态射 $Z\to\mathfrak X$。
与注 \ref{formal-spaces-remark-weak-ideals-of-definition} 中的讨论类似，
可见 $Z$ 对应于弱定义理想
$\mathcal{I}_Z\subset\mathcal{O}_\mathfrak X$。
注意 $Z_U\subset U'$ 且 $Z_V\subset V'$。因此，以这种方式构造的所有
$\mathcal{I}_Z$ 组成一族弱定义理想基本系。于是其中一个子族组成一族定义理想基本系，
见注 \ref{formal-spaces-remark-ideals-of-definition}。
\end{remark}

\begin{lemma}
\label{formal-spaces-lemma-characterize-affine}
设 $S$ 是概形，$X$ 是 $S$ 上的形式代数空间。
则 $X$ 是仿射形式代数空间，当且仅当其约化 $X_{red}$
（引理 \ref{formal-spaces-lemma-reduction-formal-algebraic-space}）仿射。
\end{lemma}

\begin{proof}
由引理 \ref{formal-spaces-lemma-characterize-quasi-separated}、
\ref{formal-spaces-lemma-characterize-quasi-compact} 以及定义
\ref{formal-spaces-definition-separated}、\ref{formal-spaces-definition-quasi-compact}，
可见 $X$ 拟紧且拟分离。由 Yasuda 引理
（引理 \ref{formal-spaces-lemma-structure-quasi-compact-quasi-separated}），
可以把 $X = \colim X_\lambda$ 写成沿代数空间增厚的滤过余极限。
然而，由《空间的极限》之引理
\ref{spaces-limits-lemma-reduction-scheme}，每个 $X_\lambda$ 都仿射，
因为 $(X_\lambda)_{red}=X_{red}$。所以按定义，$X$ 是仿射形式代数空间。
\end{proof}






\section{可由代数空间表示的态射}
\label{formal-spaces-section-representable}

\noindent
设 $f : X \to Y$ 是可由代数空间表示的形式代数空间态射。
借助代数空间各章中已经完成的工作，对于这类态射，我们有大量现成理论可用。

\begin{lemma}
\label{formal-spaces-lemma-composition-representable}
可由代数空间表示的态射之复合仍可由代数空间表示。
对可表示（由概形表示）的态射同样成立。
\end{lemma}

\begin{proof}
见《自举》之引理 \ref{bootstrap-lemma-composition-transformation}。
\end{proof}

\begin{lemma}
\label{formal-spaces-lemma-base-change-representable}
可由代数空间表示的态射的基变换仍可由代数空间表示。
对可表示（由概形表示）的态射同样成立。
\end{lemma}

\begin{proof}
见《自举》之引理 \ref{bootstrap-lemma-base-change-transformation}。
\end{proof}

\begin{lemma}
\label{formal-spaces-lemma-permanence-representable}
设 $S$ 是概形，$f : X \to Y$ 和 $g : Y \to Z$ 是 $S$ 上形式代数空间的态射。
若 $g \circ f : X \to Z$ 可由代数空间表示，则
$f : X \to Y$ 可由代数空间表示。
\end{lemma}

\begin{proof}
由引理 \ref{formal-spaces-lemma-diagonal-morphism-formal-algebraic-spaces}，
$Y \to Z$ 的对角态射可表示。因此，由《自举》之引理
\ref{bootstrap-lemma-representable-by-spaces-permanence}，
$X \to Y$ 可由代数空间表示。
\end{proof}

\noindent
可由代数空间表示这一性质在源和靶上都是局部的。

\begin{lemma}
\label{formal-spaces-lemma-representable-by-algebraic-spaces-local}
设 $S$ 是概形，$f : X \to Y$ 是 $S$ 上形式代数空间的态射。
下列条件等价：
\begin{enumerate}
\item 态射 $f$ 可由代数空间表示；
\item 存在交换图表
$$
\xymatrix{
U \ar[d] \ar[r] & V \ar[d] \\
X \ar[r] & Y
}
$$
其中 $U,V$ 是形式代数空间，竖直箭头可由代数空间表示，
$U \to X$ 是满 étale 态射，且 $U \to V$ 可由代数空间表示；
\item 对任意交换图表
$$
\xymatrix{
U \ar[d] \ar[r] & V \ar[d] \\
X \ar[r] & Y
}
$$
若 $U,V$ 是形式代数空间且竖直箭头可由代数空间表示，
则态射 $U \to V$ 可由代数空间表示；
\item 存在定义 \ref{formal-spaces-definition-formal-algebraic-space} 意义下的覆盖
$\{Y_j \to Y\}$，并且对每个 $j$，存在定义
\ref{formal-spaces-definition-formal-algebraic-space} 意义下的覆盖
$\{X_{ji} \to Y_j \times_Y X\}$，使得对每个 $j$ 和 $i$，
$X_{ji} \to Y_j$ 可由代数空间表示；
\item 存在定义 \ref{formal-spaces-definition-formal-algebraic-space} 意义下的覆盖
$\{X_i \to X\}$，并且对每个 $i$，存在分解
$X_i \to Y_i \to Y$，其中 $Y_i$ 是仿射形式代数空间，
$Y_i \to Y$ 可由代数空间表示，且 $X_i \to Y_i$ 可由代数空间表示；以及
% P10-E0318
\item 此处还可补充更多条件。
\end{enumerate}
\end{lemma}

\begin{proof}
显然 (1) 蕴含 (2)，因为可取 $U=X$ 和 $V=Y$。
反之，把《自举》之引理 \ref{bootstrap-lemma-representable-by-spaces-cover}
用于 $U \to X \to Y$，可知 (2) 蕴含 (1)。

\medskip\noindent
假设 (1) 成立，并考虑 (3) 中的图表。
则 $U \to Y$ 可由代数空间表示
（它是 $U \to X \to Y$ 的复合；见《自举》之引理
\ref{bootstrap-lemma-composition-transformation}），且分解经过 $V$。
所以由引理 \ref{formal-spaces-lemma-permanence-representable}，
$U \to V$ 可由代数空间表示。

\medskip\noindent
显然 (3) 蕴含 (2)。因此 (1)--(3) 现已等价。

\medskip\noindent
注意，(4) 中的条件是有意义的，因为由引理
\ref{formal-spaces-lemma-fibre-products}，纤维积 $Y_j \times_Y X$ 是形式代数空间。
显然 (4) 蕴含 (5)。

\medskip\noindent
假设有 (5) 中的 $X_i \to Y_i \to Y$。令
$V=\coprod Y_i$ 且 $U=\coprod X_i$，即可看出 (5) 蕴含 (2)。

\medskip\noindent
最后，假设 (1)--(3) 成立。于是可以任取定义
\ref{formal-spaces-definition-formal-algebraic-space} 意义下的覆盖
$\{Y_j \to Y\}$，并对每个 $j$ 任取定义
\ref{formal-spaces-definition-formal-algebraic-space} 意义下的覆盖
$\{X_{ji} \to Y_j \times_Y X\}$。
% P10-E0317
由 (3)，$X_{ij} \to Y_j$ 可由代数空间表示，故 (4) 成立。
证明完成。
\end{proof}

\begin{lemma}
\label{formal-spaces-lemma-algebraic-space-over-affine-formal}
设 $S$ 是概形，$Y$ 是 $S$ 上的仿射形式代数空间。
设 $f : X \to Y$ 是 $(\Sch/S)_{fppf}$ 上层的映射，并且可由代数空间表示。
则 $X$ 是形式代数空间。
\end{lemma}

\begin{proof}
按定义 \ref{formal-spaces-definition-affine-formal-algebraic-space} 写成
$Y = \colim Y_\lambda$。对每个 $\lambda$，纤维积
$X \times_Y Y_\lambda$ 都是代数空间。因此，由引理
\ref{formal-spaces-lemma-colimit-is-formal}，
$X = \colim X \times_Y Y_\lambda$ 是形式代数空间。
\end{proof}

\begin{lemma}
\label{formal-spaces-lemma-representable-by-algebraic-spaces}
设 $S$ 是概形，$Y$ 是 $S$ 上的形式代数空间。
设 $f : X \to Y$ 是 $(\Sch/S)_{fppf}$ 上层的映射，并且可由代数空间表示。
则 $X$ 是形式代数空间。
\end{lemma}

\begin{proof}
设 $\{Y_i \to Y\}$ 如定义
\ref{formal-spaces-definition-formal-algebraic-space} 所述。
则 $X \times_Y Y_i \to X$ 是一族可由代数空间表示、étale 且联合满的态射。
所以只需证明 $X \times_Y Y_i$ 是形式代数空间，见引理
\ref{formal-spaces-lemma-etale-covering-by-formal-algebraic-spaces}。
这由引理 \ref{formal-spaces-lemma-algebraic-space-over-affine-formal} 得出。
\end{proof}

\begin{lemma}
\label{formal-spaces-lemma-affine-representable-by-algebraic-spaces}
设 $S$ 是概形，$f : X \to Y$ 是仿射形式代数空间之间可由代数空间表示的态射。
则 $f$ 可表示（由概形表示）且为仿射态射。
\end{lemma}

\begin{proof}
我们将证明 $f$ 为仿射态射；于是由《空间的态射》之引理
\ref{spaces-morphisms-lemma-affine-local}，
$f$ 可表示且为仿射态射。按定义
\ref{formal-spaces-definition-affine-formal-algebraic-space} 写成
$Y = \colim Y_\mu$ 和 $X = \colim X_\lambda$。
设 $T \to Y$ 是态射，其中 $T$ 是 $S$ 上的概形。
由《自举》之定义 \ref{bootstrap-definition-property-transformation}，
必须证明 $X \times_Y T \to T$ 为仿射态射。
为此可以假设 $T$ 仿射，而需证明 $X \times_Y T$ 仿射。
此时，由引理 \ref{formal-spaces-lemma-factor-through-thickening}，
$T \to Y$ 对某个 $\mu$ 分解经过 $Y_\mu \to Y$。
由于 $f$ 拟紧，由引理 \ref{formal-spaces-lemma-characterize-quasi-compact-morphism}，
$X \times_Y T$ 拟紧。因此，$X \times_Y T \to X$ 对某个
$\lambda$ 分解经过 $X_\lambda$。类似地，增大 $\mu$ 后，
$X_\lambda \to Y$ 分解经过 $Y_\mu$。于是
$X \times_Y T = X_\lambda \times_{Y_\mu} T$。
因仿射概形的纤维积仿射，结论成立。
\end{proof}

\begin{lemma}
\label{formal-spaces-lemma-representable-affine}
设 $S$ 是概形，$\varphi : A \to B$ 是 $S$ 上弱可容许拓扑环的连续映射。
下列条件等价：
\begin{enumerate}
\item $\text{Spf}(\varphi) : \text{Spf}(B) \to \text{Spf}(A)$
可由代数空间表示；
\item $\text{Spf}(\varphi) : \text{Spf}(B) \to \text{Spf}(A)$
可表示（由概形表示）；
\item $\varphi$ 是 taut 的，见定义 \ref{formal-spaces-definition-taut}。
\end{enumerate}
\end{lemma}

\begin{proof}
由引理 \ref{formal-spaces-lemma-affine-representable-by-algebraic-spaces}，
(1) 和 (2) 等价。

\medskip\noindent
假设等价条件 (1) 和 (2) 成立。若 $I \subset A$ 是弱定义理想，则
$\Spec(A/I) \to \text{Spf}(A)$ 可表示且为增厚
（这从形式谱的构造显然可见，也可由引理
\ref{formal-spaces-lemma-mcquillan-affine-formal-algebraic-space} 得出）。
于是其基变换
$\Spec(A/I) \times_{\text{Spf}(A)} \text{Spf}(B) \to \text{Spf}(B)$
可表示且为增厚。所以由引理
\ref{formal-spaces-lemma-mcquillan-affine-formal-algebraic-space}，
存在弱定义理想 $J(I) \subset B$，使得作为
$\text{Spf}(B)$ 的子函子有
$\Spec(A/I) \times_{\text{Spf}(A)} \text{Spf}(B) = \Spec(B/J(I))$。
我们得到笛卡尔图表
$$
\xymatrix{
\Spec(B/J(I)) \ar[d] \ar[r] & \Spec(A/I) \ar[d] \\
\text{Spf}(B) \ar[r] & \text{Spf}(A)
}
$$
由引理 \ref{formal-spaces-lemma-fibre-product-affines-over-separated}，
$B/J(I) = B \widehat{\otimes}_A A/I$。
于是由引理 \ref{formal-spaces-lemma-closure-image-ideal}，
$J(I)$ 是理想 $\varphi(I)B$ 的闭包。
由于 $\text{Spf}(A) = \colim \Spec(A/I)$（其中 $I$ 如上），
可得 $\text{Spf}(B) = \colim \Spec(B/J(I))$。
因此，理想 $J(I)$ 组成一族弱定义理想基本系
（见引理 \ref{formal-spaces-lemma-mcquillan-affine-formal-algebraic-space}）。
所以 (3) 成立。

\medskip\noindent
假设 (3) 成立。实质上，我们只需反向进行上一段的论证。
设 $I \subset A$ 是弱定义理想。由引理
\ref{formal-spaces-lemma-fibre-product-affines-over-separated}，得到笛卡尔图表
$$
\xymatrix{
\text{Spf}(B \widehat{\otimes}_A A/I) \ar[d] \ar[r] & \Spec(A/I) \ar[d] \\
\text{Spf}(B) \ar[r] & \text{Spf}(A)
}
$$
若 $J(I)$ 是 $IB$ 的闭包，则由 $\varphi$ 的 taut 性，
$J(I)$ 在 $B$ 中是开的。因此，若 $J$ 在 $B$ 中是开的且
% P10-E0319
$J \subset J(B)$，则
$B/J \otimes_A A/I = B/(IB + J) = B/J(I)$，因为由引理
\ref{formal-spaces-lemma-closed}，
$J(I) = \bigcap_{J \subset B\text{ 开}} (IB + J)$。
所以定义完备张量积的极限稳定下来，给出
$B \widehat{\otimes}_A A/I = B/J(I)$。
因此
$\text{Spf}(B \widehat{\otimes}_A A/I) = \Spec(B/J(I))$。
这证明，对每个弱定义理想 $I \subset A$，
$\text{Spf}(B) \times_{\text{Spf}(A)} \Spec(A/I)$ 都可表示。
由于从拟紧概形 $T$ 出发的每个态射
$T \to \text{Spf}(A)$ 都对某个弱定义理想 $I$ 分解经过
$\Spec(A/I)$（引理 \ref{formal-spaces-lemma-factor-through-thickening}），
可知 $\text{Spf}(\varphi)$ 可表示，即 (2) 成立。证明完成。
\end{proof}

\begin{lemma}
\label{formal-spaces-lemma-property-goes-up-affine-morphism}
设 $S$ 是概形，$Y$ 是仿射形式代数空间。
设 $f : X \to Y$ 是 $(\Sch/S)_{fppf}$ 上层的映射，且可表示并为仿射态射。
则
\begin{enumerate}
\item $X$ 是仿射形式代数空间；
\item 若 $Y$ 可数指标化，则 $X$ 可数指标化；
\item 若 $Y$ 可数指标化且经典，则 $X$ 可数指标化且经典；
\item 若 $Y$ 弱 adic，则 $X$ 弱 adic；
\item 若 $Y$ 为 adic*，则 $X$ 为 adic*；以及
\item 若 $Y$ 是 Noether 的且 $f$（局部）有限型，则 $X$ 是 Noether 的。
\end{enumerate}
\end{lemma}

\begin{proof}
(1) 的证明。按定义 \ref{formal-spaces-definition-affine-formal-algebraic-space} 写成
$Y = \colim_{\lambda \in \Lambda} Y_\lambda$。
由于 $f$ 可表示且为仿射态射，纤维积
$X_\lambda = Y_\lambda \times_Y X$ 仿射，并且
% P10-E0320
$X = \colim Y_\lambda \times_Y X$。
所以 $X$ 是仿射形式代数空间。

\medskip\noindent
(2) 的证明。若 $Y$ 可数指标化，则在上述论证中可以假设
$\Lambda$ 可数。于是立即可见 $X$ 也可数指标化。

\medskip\noindent
(3)、(4) 和 (5) 的证明。在每种情形中，由假设和引理
\ref{formal-spaces-lemma-implications-between-types}，$Y$ 都是可数指标化的仿射形式代数空间；
再由 (2)，$X$ 亦然。所以由引理 \ref{formal-spaces-lemma-countably-indexed}，
可以对某些弱可容许拓扑 $S$-代数 $A,B$ 写成
$X = \text{Spf}(A)$ 和 $Y = \text{Spf}(B)$。
由引理 \ref{formal-spaces-lemma-morphism-between-formal-spectra}，态射 $f$ 对应于连续
$S$-代数同态 $\varphi : B \to A$。由引理
\ref{formal-spaces-lemma-representable-affine}，$\varphi$ 是 taut 的。
于是，(3) 由引理 \ref{formal-spaces-lemma-taut-ascent-admissible} 得出，
(4) 由引理 \ref{formal-spaces-lemma-taut-ascent-weakly-adic} 得出，
(5) 由引理 \ref{formal-spaces-lemma-taut-is-adic} 得出。

\medskip\noindent
(6) 的证明。结合 (3) 与引理 \ref{formal-spaces-lemma-implications-between-types}，
可见 $X$ 为 adic*。所以可以使用引理
\ref{formal-spaces-lemma-characterize-noetherian-affine} 的判据。
首先，它告诉我们仿射概形 $Y_\lambda$ 是 Noether 的。
于是 $X_\lambda \to Y_\lambda$ 有限型，故 $X_\lambda$ 也是 Noether 的
（《态射》之引理 \ref{morphisms-lemma-finite-type-noetherian}）。
再由该判据，$X$ 是 Noether 的，证明完成。
\end{proof}

\begin{lemma}
\label{formal-spaces-lemma-property-goes-up-affine}
设 $S$ 是概形，$f : X \to Y$ 是仿射形式代数空间之间可由代数空间表示的态射。
则
\begin{enumerate}
\item 若 $Y$ 可数指标化，则 $X$ 可数指标化；
\item 若 $Y$ 可数指标化且经典，则 $X$ 可数指标化且经典；
\item 若 $Y$ 弱 adic，则 $X$ 弱 adic；
\item 若 $Y$ 为 adic*，则 $X$ 为 adic*；以及
\item 若 $Y$ 是 Noether 的且 $f$（局部）有限型，则 $X$ 是 Noether 的。
\end{enumerate}
\end{lemma}

\begin{proof}
结合引理 \ref{formal-spaces-lemma-affine-representable-by-algebraic-spaces} 与
\ref{formal-spaces-lemma-property-goes-up-affine-morphism}。
\end{proof}

\begin{example}
\label{formal-spaces-example-representable-morphism-from-completion}
设 $B$ 是弱可容许拓扑环，$B \to A$ 是环映射（$A$ 上不带拓扑）。
于是可以考虑
$$
A^\wedge = \lim A/JA
$$
其中极限遍历 $B$ 的所有弱定义理想 $J$。
则 $A^\wedge$（赋予极限拓扑）是完备线性拓扑化环。
满射 $A^\wedge \to A/JA$ 的（开）核 $I$ 是 $JA^\wedge$ 的闭包，
见引理 \ref{formal-spaces-lemma-closed}。由引理
\ref{formal-spaces-lemma-topologically-nilpotent}，可见 $I$ 由拓扑幂零元组成。
所以 $I$ 是 $A^\wedge$ 的弱定义理想，进而 $A^\wedge$ 是弱可容许拓扑环。
% P10-E0321
因此，$\varphi : B \to A^\wedge$ 是弱可容许拓扑环的 taut 映射，且
$$
\text{Spf}(A^\wedge) \longrightarrow \text{Spf}(B)
$$
是引理 \ref{formal-spaces-lemma-representable-affine} 所研究现象的一个特例。
\end{example}

\begin{remark}[警告]
\label{formal-spaces-remark-warning}
引理 \ref{formal-spaces-lemma-representable-affine}、
\ref{formal-spaces-lemma-property-goes-up-affine-morphism} 和
\ref{formal-spaces-lemma-property-goes-up-affine} 中的讨论在以下两方面都是最精确的：
\begin{enumerate}
\item 若 $A,B$ 是弱可容许环且 $\varphi : A \to B$ 是连续映射，则一般而言，
$\text{Spf}(\varphi) : \text{Spf}(B) \to \text{Spf}(A)$ 不可表示。
\item 若 $f : Y \to X$ 是仿射形式代数空间之间的可表示态射，
且 $X = \text{Spf}(A)$ 是 McQuillan 的，则不能由此推出 $Y$ 是 McQuillan 的。
\end{enumerate}
(1) 的一个例子是取域 $A=k$（带离散拓扑）以及带 $t$-进拓扑的
$B=k[[t]]$。(2) 的一个例子见《例子》之第
\ref{examples-section-affine-formal-algebraic-space} 节。
\end{remark}

\noindent
尽管有上述警告，我们仍有如下结果。

\begin{lemma}
\label{formal-spaces-lemma-etale}
设 $S$ 是概形，$Y$ 是 $S$ 上的 McQuillan 仿射形式代数空间；
即对某个弱可容许拓扑 $S$-代数 $B$，有
$Y = \text{Spf}(B)$。则下列两个范畴等价：
\begin{enumerate}
\item 仿射形式代数空间之间可由代数空间表示且为 étale 的态射
$f : X \to Y$ 所成的范畴；以及
\item 形如 $A^\wedge$ 的拓扑 $B$-代数所成的范畴，其中
$A$ 是 étale $B$-代数，并且
$A^\wedge = \lim A/JA$，这里 $J \subset B$ 遍历 $B$ 的弱定义理想。
\end{enumerate}
该等价把 $A^\wedge$ 送到 $X = \text{Spf}(A^\wedge)$。
特别地，(1) 中的任意 $X$ 都是 McQuillan 的。
\end{lemma}

\begin{proof}
设 $A$ 是 étale $B$-代数。则对每个开理想 $J \subset B$，
$B/J \to A/JA$ 都是 étale 的。因此，态射
$\text{Spf}(A^\wedge) \to Y$ 可表示且为 étale。
由引理 \ref{formal-spaces-lemma-morphism-between-formal-spectra}，
函子 $\text{Spf}$ 全忠实。为完成证明，下一段将说明 (1) 中任意
$X \to Y$ 都属于本质像。

\medskip\noindent
选取弱定义理想 $J_0 \subset B$。令
$Y_0 = \Spec(B/J_0)$ 和 $X_0 = Y_0 \times_Y X$。
则 $X_0 \to Y_0$ 是仿射概形之间的 étale 态射
（见引理 \ref{formal-spaces-lemma-affine-representable-by-algebraic-spaces}）。
设 $X_0 = \Spec(A_0)$。由《代数》之引理
\ref{algebra-lemma-lift-etale}，可以找到 étale 代数映射
$B \to A$，使得 $A_0 \cong A/J_0A$。
% P10-E0322
考虑定义理想 $J \subset J_0$。如上，可以对某个 étale 环映射
$B/J \to \bar A$ 写成
$\Spec(B/J) \times_Y X = \Spec(\bar A)$。
于是 $B/J \to \bar A$ 和 $B/J \to A/JA$ 都是 étale 环映射，
并且都提升 étale 环映射 $B/J_0 \to A_0$。
由《更多代数》之引理
\ref{more-algebra-lemma-locally-nilpotent-henselian}，
存在唯一的 $B/J$-代数同构
$\varphi_J : A/JA \to \bar A$，它提升模 $J_0$ 的同一化。
由于映射 $\varphi_J$ 唯一，当 $J$ 变化时它们彼此相容。因此
% P10-E0323
$$
X = \colim \Spec(B/J) \times_Y X = \colim \Spec(A/JA) = \text{Spf}(A)
$$
故引理成立。
\end{proof}

\begin{lemma}
\label{formal-spaces-lemma-etale-surjective}
采用引理 \ref{formal-spaces-lemma-etale} 的记号和假设，设
$f : X \to Y$ 对应于 $B \to A^\wedge$。下列条件等价：
\begin{enumerate}
\item $f : X \to Y$ 满；
\item $B \to A$ 忠实平坦；
\item 对每个弱定义理想 $J \subset B$，环映射
$B/J \to A/JA$ 忠实平坦；以及
\item 对某个弱定义理想 $J \subset B$，环映射
$B/J \to A/JA$ 忠实平坦。
\end{enumerate}
\end{lemma}

\begin{proof}
设 $J \subset B$ 是弱定义理想。由于 $J$ 的每个元素都拓扑幂零，
$1+J$ 的每个元素都是单位。因此，$J$ 包含在 $B$ 的 Jacobson 根中
（《代数》之引理 \ref{algebra-lemma-contained-in-radical}）。
所以，平坦环映射 $B \to A$ 忠实平坦，当且仅当
$B/J \to A/JA$ 忠实平坦
（《代数》之引理 \ref{algebra-lemma-ff-rings}）。
由此可见 (2)--(4) 等价。
若 (1) 成立，则对每个弱定义理想 $J \subset B$，态射
$\Spec(A/JA) = \Spec(B/J) \times_Y X \to \Spec(B/J)$ 满，
这蕴含 (3)。反之，假设 (3)。从拟紧概形 $T$ 出发的态射
$T \to Y$ 对 $B$ 的某个定义理想 $J$ 分解经过 $\Spec(B/J)$
% P10-E0324
（引理 \ref{formal-spaces-lemma-factor-through-thickening}）。因此，
$X \times_Y T = \Spec(A/JA) \times_{\Spec(B/J)} T \to T$
是满态射 $\Spec(A/JA) \to \Spec(B/J)$ 的基变换，因而满。
所以 (1) 成立。
\end{proof}





\section{形式代数空间的类型}
\label{formal-spaces-section-types}

\noindent
本节定义“局部 Noether”、“局部 adic*”、“局部弱 adic”、
“局部可数指标化且经典”以及“局部可数指标化”的形式代数空间。
“局部 adic”、“局部经典”和“局部 McQuillan”这几种类型尚未定义，
因为我们不知道如何在这些情形中证明下列引理的类似结果
（只须对仿射形式代数空间之间的 étale 覆盖证明这些引理的类似结果就足够了）。

\begin{lemma}
\label{formal-spaces-lemma-iff-countably-indexed}
设 $S$ 是概形，$X \to Y$ 是仿射形式代数空间之间可由代数空间表示、
满且平坦的态射。则 $X$ 可数指标化，当且仅当 $Y$ 可数指标化。
\end{lemma}

\begin{proof}
假设 $X$ 可数指标化。按引理
\ref{formal-spaces-lemma-countable-affine-formal-algebraic-space} 写成
$X = \colim X_n$。按定义
\ref{formal-spaces-definition-affine-formal-algebraic-space} 写成
$Y = \colim Y_\lambda$。对每个 $n$，可以选取 $\lambda_n$，使得
$X_n \to Y$ 分解经过 $Y_{\lambda_n}$，见引理
\ref{formal-spaces-lemma-factor-through-thickening}。
另一方面，对每个 $\lambda$，概形 $Y_\lambda \times_Y X$ 仿射
（引理 \ref{formal-spaces-lemma-affine-representable-by-algebraic-spaces}），
所以 $Y_\lambda \times_Y X \to X$ 对某个 $n$ 分解经过 $X_n$
（引理 \ref{formal-spaces-lemma-factor-through-thickening}）。考虑图表
$$
\xymatrix{
Y_\lambda \times_Y X \ar[r] \ar[d] & X_n \ar[r] \ar[d] & X \ar[d] \\
Y_\lambda \ar@{..>}[r] \ar@/_1pc/[rr] & Y_{\lambda_n} \ar[r] & Y
}
$$
若能证明虚线箭头存在，则可推出
$Y = \colim Y_{\lambda_n}$，从而 $Y$ 可数指标化。
为此，选取 $\mu \geq \lambda$ 且 $\mu \geq \lambda_n$。
于是 $Y_\lambda \to Y$ 和 $Y_{\lambda_n} \to Y$ 都分解经过
$Y_\mu \to Y$。设 $Y_\mu = \Spec(B_\mu)$，闭子概形
$Y_\lambda$ 对应于 $J \subset B_\mu$，闭子概形
$Y_{\lambda_n}$ 对应于 $J' \subset B_\mu$。我们要证明 $J' \subset J$。
由上图可知 $J'A_\mu \subset JA_\mu$，其中
$Y_\mu \times_Y X = \Spec(A_\mu)$。
由于 $X \to Y$ 满且平坦，态射
$Y_\lambda \times_Y X \to Y_\lambda$ 是仿射概形之间的忠实平坦态射，
故 $B_\mu \to A_\mu$ 忠实平坦。因此如所需，$J' \subset J$。

\medskip\noindent
假设 $Y$ 可数指标化。则由引理
\ref{formal-spaces-lemma-property-goes-up-affine}，$X$ 可数指标化。
\end{proof}

\begin{lemma}
\label{formal-spaces-lemma-iff-cic}
设 $S$ 是概形，$X \to Y$ 是仿射形式代数空间之间可由代数空间表示、
满且平坦的态射。则 $X$ 可数指标化且经典，当且仅当
$Y$ 可数指标化且经典。
\end{lemma}

\begin{proof}
一个方向的蕴含已在引理 \ref{formal-spaces-lemma-property-goes-up-affine} 中证明。
对另一个方向，注意由引理 \ref{formal-spaces-lemma-iff-countably-indexed}，
可以假设 $X,Y$ 都可数指标化。因此，由引理
\ref{formal-spaces-lemma-countably-indexed}，对某些弱可容许拓扑 $S$-代数
$A,B$，有 $X = \text{Spf}(A)$ 和 $Y = \text{Spf}(B)$。
由引理 \ref{formal-spaces-lemma-morphism-between-formal-spectra}，
态射 $X \to Y$ 对应于连续 $S$-代数同态
$\varphi : B \to A$。由引理 \ref{formal-spaces-lemma-representable-affine}，
$\varphi$ 是 taut 的。设 $J \subset B$ 是开理想，
$I \subset A$ 是 $JA$ 的闭包。由引理
\ref{formal-spaces-lemma-fibre-product-affines-over-separated} 和
\ref{formal-spaces-lemma-closure-image-ideal}，可见
$\Spec(B/J) \times_Y X = \Spec(A/I)$。
因此，$B/J \to A/I$ 忠实平坦（因为 $X \to Y$ 满且平坦）。
这意味着 $\varphi : B \to A$ 满足第
\ref{formal-spaces-section-taut-descent} 节的条件（交换 $A$ 与 $B$ 的角色）。
所以由引理 \ref{formal-spaces-lemma-taut-descent-weakly-admissible}，结论成立。
\end{proof}

\begin{lemma}
\label{formal-spaces-lemma-iff-weakly-adic}
设 $S$ 是概形，$X \to Y$ 是仿射形式代数空间之间可由代数空间表示、
满且平坦的态射。则 $X$ 弱 adic，当且仅当 $Y$ 弱 adic。
\end{lemma}

\begin{proof}
证明与引理 \ref{formal-spaces-lemma-iff-cic} 的证明完全相同，
但最后使用引理 \ref{formal-spaces-lemma-taut-descent-weakly-adic}。
\end{proof}

\begin{lemma}
\label{formal-spaces-lemma-iff-adic-star}
设 $S$ 是概形，$X \to Y$ 是仿射形式代数空间之间可由代数空间表示、
满且平坦的态射。则 $X$ 为 adic*，当且仅当 $Y$ 为 adic*。
\end{lemma}

\begin{proof}
证明与引理 \ref{formal-spaces-lemma-iff-cic} 的证明完全相同，
但最后使用引理 \ref{formal-spaces-lemma-taut-descent-adic-star}。
\end{proof}

\begin{lemma}
\label{formal-spaces-lemma-iff-noetherian}
设 $S$ 是概形，$X \to Y$ 是仿射形式代数空间之间可由代数空间表示、
满、平坦且（局部）有限型的态射。则 $X$ 是 Noether 的，
当且仅当 $Y$ 是 Noether 的。
\end{lemma}

\begin{proof}
注意，Noether 仿射形式代数空间为 adic*，见引理
\ref{formal-spaces-lemma-implications-between-types}。所以由引理
\ref{formal-spaces-lemma-iff-adic-star}，可以假设 $X,Y$ 都为 adic*。
我们将使用引理 \ref{formal-spaces-lemma-characterize-noetherian-affine} 的判据证明结论。
具体地，按引理 \ref{formal-spaces-lemma-countable-affine-formal-algebraic-space} 写成
$Y = \colim Y_n$。对每个 $n$，令
$X_n = Y_n \times_Y X$。则 $X_n$ 是仿射概形
（引理 \ref{formal-spaces-lemma-affine-representable-by-algebraic-spaces}），且
$X = \colim X_n$。每个态射 $X_n \to Y_n$ 都忠实平坦且有限型。因此，引理由
$X_n$ 在这种情形中是 Noether 的当且仅当 $Y_n$ 是 Noether 的这一事实得出，见
《代数》之引理 \ref{algebra-lemma-descent-Noetherian}（下降）和
《代数》之引理 \ref{algebra-lemma-Noetherian-permanence}（上升）。
\end{proof}

\begin{lemma}
\label{formal-spaces-lemma-type-local}
设 $S$ 是概形。令
$$
P \in
\left\{
\begin{matrix}
\text{可数指标化},\\
\text{可数指标化且经典},\\
\text{弱 adic},\ adic*,\ \text{Noether}
\end{matrix}
\right\}
$$
设 $X$ 是 $S$ 上的形式代数空间。下列条件等价：
\begin{enumerate}
\item 若 $Y$ 是仿射形式代数空间，且
$f : Y \to X$ 可由代数空间表示并为 étale，则 $Y$ 具有性质 $P$；
\item 对定义 \ref{formal-spaces-definition-formal-algebraic-space} 中的某个
$\{X_i \to X\}_{i \in I}$，每个 $X_i$ 都具有性质 $P$。
\end{enumerate}
\end{lemma}

\begin{proof}
显然 (1) 蕴含 (2)。假设 (2)，并设 $Y \to X$ 如 (1) 所述。
由于纤维积 $X_i \times_X Y$ 是形式代数空间
（引理 \ref{formal-spaces-lemma-fibre-products-general}），可以选取定义
\ref{formal-spaces-definition-formal-algebraic-space} 意义下的覆盖
$\{X_{ij} \to X_i \times_X Y\}$。
由于 $Y$ 拟紧，存在 $(i_1,j_1),\ldots,(i_n,j_n)$，使得
$$
X_{i_1 j_1} \amalg \ldots \amalg X_{i_n j_n} \longrightarrow Y
$$
满且为 étale。于是 $X_{i_kj_k} \to X_{i_k}$ 可由代数空间表示且为 étale，
故由引理 \ref{formal-spaces-lemma-property-goes-up-affine}，
$X_{i_kj_k}$ 具有性质 $P$。因此
$X_{i_1 j_1} \amalg \ldots \amalg X_{i_n j_n}$ 是具有性质 $P$ 的
仿射形式代数空间（这里略去关于仿射形式代数空间有限不交并的一点细节）。
最后应用引理 \ref{formal-spaces-lemma-iff-countably-indexed}、
\ref{formal-spaces-lemma-iff-cic}、\ref{formal-spaces-lemma-iff-weakly-adic}、
\ref{formal-spaces-lemma-iff-adic-star} 和 \ref{formal-spaces-lemma-iff-noetherian} 之一即可得出结论。
\end{proof}

\noindent
上一引理为如下定义铺平了道路。

\begin{definition}
\label{formal-spaces-definition-types-formal-algebraic-spaces}
设 $S$ 是概形，$X$ 是 $S$ 上的形式代数空间。
如果引理 \ref{formal-spaces-lemma-type-local} 的等价条件对相应性质成立，
则称 $X$ \textit{局部可数指标化}、
\textit{局部可数指标化且经典}、
\textit{局部弱 adic}、\textit{局部 adic*} 或
\textit{局部 Noether}。
\end{definition}

\noindent
局部 Noether 代数空间沿闭子集的形式完备化是局部 Noether 形式代数空间。

\begin{lemma}
\label{formal-spaces-lemma-formal-completion-types}
设 $S$ 是概形，$X$ 是 $S$ 上的代数空间，$T \subset |X|$ 是闭子集。
设 $X_{/T}$ 是 $X$ 沿 $T$ 的形式完备化。
\begin{enumerate}
\item 若 $X \setminus T \to X$ 拟紧，则 $X_{/T}$ 局部 adic*。
\item 若 $X$ 局部 Noether，则 $X_{/T}$ 局部 Noether。
\end{enumerate}
\end{lemma}

\begin{proof}
选取满 étale 态射 $U \to X$，其中
$U=\coprod U_i$ 是仿射概形的不交并，见《空间的性质》之引理
\ref{spaces-properties-lemma-cover-by-union-affines}。
设 $T_i \subset U_i$ 是 $T$ 的逆像。由引理
\ref{formal-spaces-lemma-map-completions-representable}，有
$X_{/T} \times_X U_i = (U_i)_{/T_i}$。
所以 $\{(U_i)_{/T_i} \to X_{/T}\}$ 是定义
\ref{formal-spaces-definition-formal-algebraic-space} 意义下的覆盖。
此外，若 $X \setminus T \to X$ 拟紧，则
$U_i \setminus T_i \to U_i$ 亦然；若 $X$ 局部 Noether，
则 $U_i$ 亦然。因此，引理由仿射情形即引理
\ref{formal-spaces-lemma-affine-formal-completion-types} 得出。
\end{proof}

\begin{remark}[警告]
\label{formal-spaces-remark-warning-completion}
设 $X = \Spec(A)$，且 $T \subset X$ 是有限生成理想
$I \subset A$ 的零点集。令 $J=\sqrt I$ 为 $I$ 的根。
由定义可见 $X_{/T}=\text{Spf}(A^\wedge)$，其中
$A^\wedge=\lim A/I^n$ 是 $A$ 的 $I$-进完备化。
另一方面，从 $I$-进完备化到 $J$-进完备化的映射
$A^\wedge \to \lim A/J^n$ 可能不是环同构。例如，令
$$
A = \bigcup\nolimits_{n \geq 1} \mathbf{C}[t^{1/n}]
$$
且 $I=(t)$。则 $J=\mathfrak m$ 是赋值环 $A$ 的极大理想，
且 $J^2=J$。所以 $A$ 的 $J$-进完备化是 $\mathbf C$，
而 $I$-进完备化是例 \ref{formal-spaces-example-david-hansen} 所述的赋值环
（不过特别容易看出 $A \subset A^\wedge$）。
\end{remark}

\begin{lemma}
\label{formal-spaces-lemma-types-fibre-products}
% P10-E0325
设 $S$ 是概形，$X \to Y$ 和 $Z \to Y$ 是 $S$ 上形式代数空间的态射。
则
\begin{enumerate}
\item 若 $X,Z$ 局部可数指标化，则 $X \times_Y Z$ 局部可数指标化。
\item 若 $X,Z$ 局部可数指标化且经典，则
$X \times_Y Z$ 局部可数指标化且经典。
% P10-E0326
\item 若 $X,Z$ 弱 adic，则 $X \times_Y Z$ 弱 adic。
\item 若 $X,Z$ 局部 adic*，则 $X \times_Y Z$ 局部 adic*。
\item 若 $X,Z$ 局部 Noether，且 $X_{red} \to Y_{red}$ 局部有限型，
则 $X \times_Y Z$ 局部 Noether。
\end{enumerate}
\end{lemma}

\begin{proof}
选取定义 \ref{formal-spaces-definition-formal-algebraic-space} 意义下的覆盖
$\{Y_j \to Y\}$。对每个 $j$，选取定义
\ref{formal-spaces-definition-formal-algebraic-space} 意义下的覆盖
$\{X_{ji} \to Y_j \times_Y X\}$，并选取定义
\ref{formal-spaces-definition-formal-algebraic-space} 意义下的覆盖
$\{Z_{jk} \to Y_j \times_Y Z\}$。
由引理 \ref{formal-spaces-lemma-fibre-product-affines-over-separated}，
$X_{ji} \times_{Y_j} Z_{jk}$ 是仿射形式代数空间。因此
$$
\{X_{ji} \times_{Y_j} Z_{jk} \to X \times_Y Z\}
$$
是定义 \ref{formal-spaces-definition-formal-algebraic-space} 意义下的覆盖。
% P10-E0330
所以，只需在 $X,Y,Z$ 都是仿射形式代数空间的情形中证明 (1)、(2)、(3) 和 (4)。

\medskip\noindent
假设 $X,Z$ 可数指标化。按引理
\ref{formal-spaces-lemma-countable-affine-formal-algebraic-space} 写成
$X=\colim X_n$ 和 $Z=\colim Z_m$。
按定义 \ref{formal-spaces-definition-affine-formal-algebraic-space} 写成
$Y=\colim_{\lambda\in\Lambda}Y_\lambda$。
对每个 $n,m$，可以找到 $\lambda_{n,m}\in\Lambda$，使得
$X_n\to Y$ 和 $Z_m\to Y$ 都分解经过 $Y_{\lambda_{n,m}}$
（例如见引理 \ref{formal-spaces-lemma-factor-through-thickening}）。
选取 $\lambda_0\in\Lambda$。对 $t\geq1$ 归纳选取
$\lambda_t\in\Lambda$，使它对所有 $1\leq n,m\leq t$ 都满足
$\lambda_t\geq\lambda_{n,m}$，并且 $\lambda_t\geq\lambda_{t-1}$。
令 $Y'=\colim Y_{\lambda_t}$。则 $Y'\to Y$ 是单态射，且
$X\to Y$ 和 $Z\to Y$ 都分解经过 $Y'$。所以可以用 $Y'$ 代替 $Y$；
换言之，可以假设 $Y$ 可数指标化。

\medskip\noindent
假设 $X,Y,Z$ 可数指标化。由引理 \ref{formal-spaces-lemma-countably-indexed}，
可对某些弱可容许拓扑环 $A,B,C$ 写成
$X=\text{Spf}(A)$、$Y=\text{Spf}(B)$、$Z=\text{Spf}(C)$。
% P10-E0327
由引理 \ref{formal-spaces-lemma-morphism-between-formal-spectra}，
态射 $X\to Y$ 和 $Z\to Y$ 由连续环映射
$B\to A$ 和 $B\to C$ 给出。
由引理 \ref{formal-spaces-lemma-fibre-product-affines-over-separated}，可见
$X\times_Y Z=\text{Spf}(A\widehat{\otimes}_B C)$，并且
$A\widehat{\otimes}_B C$ 是弱可容许拓扑环。
特别地，由引理 \ref{formal-spaces-lemma-completed-tensor-product} 的第 (3) 部分，
$X\times_Y Z$ 可数指标化。这证明了 (1)。

\medskip\noindent
(2) 的证明。在这种情形中，$X,Z$ 可数指标化，故上述论证表明
$X\times_Y Z$ 是 $A\widehat{\otimes}_B C$ 的形式谱，其中 $A,C$ 可容许。
于是由引理 \ref{formal-spaces-lemma-completed-tensor-product} 的第 (2) 部分，
$A\widehat{\otimes}_B C$ 可容许。

\medskip\noindent
(3) 的证明。如前可知，$X\times_Y Z$ 是
$A\widehat{\otimes}_B C$ 的形式谱，其中 $A,C$ 弱 adic。
于是由引理 \ref{formal-spaces-lemma-completed-tensor-product-weakly-adic}，
$A\widehat{\otimes}_B C$ 弱 adic。

\medskip\noindent
(4) 的证明。与上面一样论证，结论由引理
\ref{formal-spaces-lemma-completed-tensor-product} 的第 (4) 部分得出。

\medskip\noindent
(5) 的证明。为了从引理 \ref{formal-spaces-lemma-completed-tensor-product} 的第 (5) 部分
推出情形 (5)，需要证明假设彼此吻合。
% P10-E0328
具体地，采用证明第一段的记号，若
$X_{red}\to Y_{red}$ 局部有限型，则
$(X_{ji})_{red}\to(Y_j)_{red}$ 局部有限型。
这由《空间的态射》之引理
\ref{spaces-morphisms-lemma-finite-type-local} 以及下列交换图表中竖直态射为
étale 这一事实得出：
$$
\xymatrix{
(X_{ji})_{red} \ar[d] \ar[r] & (Y_j)_{red} \ar[d] \\
X_{red} \ar[r] & Y_{red}
}
$$
事实上，由引理 \ref{formal-spaces-lemma-reduction-smooth}，
% P10-E0329
有 $(X_{ji})_{red}=X_{ij}\times_X X_{red}$ 和
$(Y_j)_{red}=Y_j\times_Y Y_{red}$。
所以如上可归约到如下情形：$X,Y,Z$ 是仿射形式代数空间，
$X,Z$ 是 Noether 的，且 $X_{red}\to Y_{red}$ 有限型。
其次，在证明第二段中我们用 $Y'$ 替换了 $Y$，但按构造
$Y_{red}=Y'_{red}$，故有限型假设在该替换下保持不变。
于是可见 $X,Y,Z$ 分别对应于 $A,B,C$，而
$X\times_Y Z$ 对应于 $A\widehat{\otimes}_B C$，其中 $A,C$ 是
Noether adic 环。最后，取约化对应于除去拓扑幂零元理想
（例 \ref{formal-spaces-example-reduction-affine-formal-spectrum}），
所以 $X_{red}\to Y_{red}$ 有限型确实意味着
$B/\mathfrak b\to A/\mathfrak a$ 有限型。证明完成。
\end{proof}

\begin{lemma}
\label{formal-spaces-lemma-structure-locally-noetherian}
设 $S$ 是概形，$X$ 是 $S$ 上的局部 Noether 形式代数空间。
则 $X=\colim X_n$，其中
$X_1\to X_2\to X_3\to\ldots$ 是 $S$ 上局部 Noether 代数空间的
有限阶增厚系统，$X_1=X_{red}$，并且对所有 $m\geq n$，
$X_n$ 是 $X_1$ 在 $X_m$ 中的第 $n$ 个无穷小邻域。
\end{lemma}

\begin{proof}
我们只概述证明，并略去若干细节。令 $X_1=X_{red}$。
按如下规则定义子函子 $X_n\subset X$：若 $T$ 是概形，
则态射 $f:T\to X$ 分解经过 $X_n$，当且仅当闭浸入
$X_1\times_XT\to T$ 的理想层的第 $n$ 次幂为零。
于是 $X_n\subset X$ 是子层，因为拟凝聚模的零化可以 fppf 局部地检验。
我们断言，$X_n\to X$ 可由概形表示且为闭浸入，并且作为 fppf 层有
$X=\colim X_n$。为检验这一点，可以在 $X$ 上 étale 局部地工作。
所以可假设 $X=\text{Spf}(A)$ 是 Noether 仿射形式代数空间。
于是 $X_1=\Spec(A/\mathfrak a)$，其中 $\mathfrak a\subset A$ 是
Noether adic 拓扑环 $A$ 的拓扑幂零元理想。此时
$X_n=\Spec(A/\mathfrak a^n)$，从而得到所需结论。
\end{proof}









\section{态射与连续环映射}
\label{formal-spaces-section-morphisms-rings}

\noindent
本节以 $\textit{WAdm}$ 表示弱可容许拓扑环与连续环同态所成的范畴。
我们定义全子范畴
$$
\textit{WAdm} \supset
\textit{WAdm}^{count} \supset
\textit{WAdm}^{cic} \supset
\textit{WAdm}^{weakly\ adic} \supset
\textit{WAdm}^{adic*} \supset
\textit{WAdm}^{Noeth}
$$
其对象分别为
\begin{enumerate}
\item $\textit{WAdm}^{count}$：具有可数开理想基本系的弱可容许拓扑环 $A$；
\item $\textit{WAdm}^{cic}$：具有可数开理想基本系的可容许拓扑环 $A$；
\item $\textit{WAdm}^{weakly\ adic}$：弱 adic 拓扑环
（第 \ref{formal-spaces-section-weakly-adic} 节）；
\item $\textit{WAdm}^{adic*}$：具有有限生成定义理想的 adic 拓扑环；以及
\item $\textit{WAdm}^{Noeth}$：Noether adic 拓扑环。
\end{enumerate}
显然，这些类型的环的形式谱分别是局部可数指标化、
局部可数指标化且经典、局部弱 adic、局部 adic* 以及局部 Noether
形式代数空间的基本构件。

\medskip\noindent
我们简要回顾可数指标化仿射形式代数空间的态射与
$\textit{WAdm}^{count}$ 中态射之间的关系。
设 $S$ 是概形，$X,Y$ 是可数指标化仿射形式代数空间。
% P10-E0331
写成 $X=\text{Spf}(A)$ 和 $Y=\text{Spf}(B)$，其中拓扑
$S$-代数 $A,B$ 属于 $\textit{WAdm}^{count}$，见引理
\ref{formal-spaces-lemma-countably-indexed}。由引理
\ref{formal-spaces-lemma-morphism-between-formal-spectra}，态射
$f:X\to Y$ 与拓扑 $S$-代数的连续映射
$$
\varphi : B \longrightarrow A
$$
之间存在一一对应。该对应由
$f\mapsto f^\sharp$ 和 $\varphi\mapsto\text{Spf}(\varphi)$ 给出。

\medskip\noindent
设 $S$ 是概形，$f:X\to Y$ 是局部可数指标化形式代数空间的态射。
考虑交换图表
$$
\xymatrix{
U \ar[d] \ar[r] & V \ar[d] \\
X \ar[r] & Y
}
$$
其中 $U,V$ 是仿射形式代数空间，且 $U\to X$ 和 $V\to Y$
可由代数空间表示并为 étale。由定义
\ref{formal-spaces-definition-types-formal-algebraic-spaces}
（因而也由引理 \ref{formal-spaces-lemma-type-local}），可知 $U,V$ 是可数指标化
仿射形式代数空间。由上一段的讨论，$U\to V$ 同构于
$\text{Spf}(\varphi)$，其中
$$
\varphi : B \longrightarrow A
$$
是 $\textit{WAdm}^{count}$ 中拓扑 $S$-代数的连续映射。

\begin{lemma}
\label{formal-spaces-lemma-completion-in-sub}
设 $A\in\Ob(\textit{WAdm})$，$A\to A'$ 是环映射（无拓扑）。
令
$(A')^\wedge=\lim_{I\subset A\text{ w.i.d}}A'/IA'$ 为例
\ref{formal-spaces-example-representable-morphism-from-completion} 中构造的
$\textit{WAdm}$ 对象。
\begin{enumerate}
\item 若 $A$ 属于 $\textit{WAdm}^{count}$，则 $(A')^\wedge$ 亦然。
\item 若 $A$ 属于 $\textit{WAdm}^{cic}$，则 $(A')^\wedge$ 亦然。
\item 若 $A$ 属于 $\textit{WAdm}^{weakly\ adic}$，则 $(A')^\wedge$ 亦然。
\item 若 $A$ 属于 $\textit{WAdm}^{adic*}$，则 $(A')^\wedge$ 亦然。
\item 若 $A$ 属于 $\textit{WAdm}^{Noeth}$ 且 $A'$ 是 Noether 的，
则 $(A')^\wedge$ 属于 $\textit{WAdm}^{Noeth}$。
\end{enumerate}
\end{lemma}

\begin{proof}
回忆 $A\to(A')^\wedge$ 是 taut 的，见例
\ref{formal-spaces-example-representable-morphism-from-completion} 中的讨论。
所以 (1)、(2)、(3) 和 (4) 分别由引理
\ref{formal-spaces-lemma-taut-ascent-countable}、
\ref{formal-spaces-lemma-taut-ascent-admissible}、
\ref{formal-spaces-lemma-taut-ascent-weakly-adic} 和
\ref{formal-spaces-lemma-taut-is-adic} 得出。
最后，假设 $A$ 是 Noether adic 环。由 (4)，
$(A')^\wedge$ 是 adic 的。由《代数》之引理
\ref{algebra-lemma-completion-Noetherian-Noetherian}，
$(A')^\wedge$ 是 Noether 的。所以 (5) 成立。
\end{proof}

\begin{situation}
\label{formal-spaces-situation-local-property}
设 $P$ 是 $\textit{WAdm}^{count}$ 态射的一个性质。
考虑交换图表
\begin{equation}
\label{formal-spaces-equation-localize}
\vcenter{
\xymatrix{
A \ar[r] & (A')^\wedge \\
B \ar[r] \ar[u]^\varphi & (B')^\wedge \ar[u]_{\varphi'}
}
}
\end{equation}
并假设满足下列条件：
\begin{enumerate}
\item $A,B$ 是 $\textit{WAdm}^{count}$ 的对象；
\item $A\to A'$ 和 $B\to B'$ 是 étale 环映射；
\item $(A')^\wedge=\lim A'/IA'$，而
$(B')^\wedge=\lim B'/JB'$，其中 $I\subset A$ 与 $J\subset B$
分别遍历 $A$ 与 $B$ 的弱可容许定义理想；
% P10-E0332
\item $\varphi:B\to A$ 和
$\varphi':(B')^\wedge\to(A')^\wedge$ 连续。
\end{enumerate}
由引理 \ref{formal-spaces-lemma-completion-in-sub}，拓扑环
$(A')^\wedge$ 和 $(B')^\wedge$ 是 $\textit{WAdm}^{count}$ 的对象。
若下列公理成立，则称 $P$ 是\textit{局部性质}：
\begin{enumerate}
\item
\label{formal-spaces-item-axiom-1}
对任意图表 (\ref{formal-spaces-equation-localize})，有
$P(\varphi)\Rightarrow P(\varphi')$；
\item
\label{formal-spaces-item-axiom-2}
对任意图表 (\ref{formal-spaces-equation-localize})，若 $A\to A'$ 忠实平坦，
则 $P(\varphi')\Rightarrow P(\varphi)$；
\item
\label{formal-spaces-item-axiom-3}
若对 $i=1,\ldots,n$ 都有 $P(B\to A_i)$，则
$P(B\to\prod_{i=1,\ldots,n}A_i)$。
\end{enumerate}
由于 $\textit{WAdm}^{count}$ 有有限积，公理
(\ref{formal-spaces-item-axiom-3}) 是有意义的。
\end{situation}

\begin{lemma}
\label{formal-spaces-lemma-property-defines-property-morphisms}
设 $S$ 是概形，$f:X\to Y$ 是 $S$ 上局部可数指标化形式代数空间的态射。
设 $P$ 是 $\textit{WAdm}^{count}$ 态射的局部性质。下列条件等价：
\begin{enumerate}
\item 对每个交换图表
$$
\xymatrix{
U \ar[d] \ar[r] & V \ar[d] \\
X \ar[r] & Y
}
$$
若 $U,V$ 是仿射形式代数空间，且 $U\to X$ 和 $V\to Y$
可由代数空间表示并为 étale，则态射 $U\to V$ 对应于
具有性质 $P$ 的 $\textit{WAdm}^{count}$ 态射；
\item 存在定义 \ref{formal-spaces-definition-formal-algebraic-space} 意义下的覆盖
$\{Y_j\to Y\}$，并且对每个 $j$ 存在定义
\ref{formal-spaces-definition-formal-algebraic-space} 意义下的覆盖
$\{X_{ji}\to Y_j\times_YX\}$，使每个 $X_{ji}\to Y_j$
都对应于具有性质 $P$ 的 $\textit{WAdm}^{count}$ 态射；以及
\item 存在定义 \ref{formal-spaces-definition-formal-algebraic-space} 意义下的覆盖
$\{X_i\to X\}$，并且对每个 $i$ 存在分解
$X_i\to Y_i\to Y$，其中 $Y_i$ 是仿射形式代数空间，
$Y_i\to Y$ 可由代数空间表示并为 étale，且 $X_i\to Y_i$
对应于具有性质 $P$ 的 $\textit{WAdm}^{count}$ 态射。
\end{enumerate}
\end{lemma}

\begin{proof}
显然 (1) 蕴含 (2)，且 (2) 蕴含 (3)。
假设 $\{X_i\to X\}$ 和 $X_i\to Y_i\to Y$ 如 (3) 所述，
并给定 (1) 中的图表。由于 $Y_i\times_YV$ 是形式代数空间
（引理 \ref{formal-spaces-lemma-fibre-products-general}），可以选取定义
\ref{formal-spaces-definition-formal-algebraic-space} 意义下的覆盖
$\{Y_{ij}\to Y_i\times_YV\}$。对每个 $(i,j)$，类似地选取定义
\ref{formal-spaces-definition-formal-algebraic-space} 意义下的覆盖
$\{X_{ijk}\to Y_{ij}\times_{Y_i}X_i\times_XU\}$。
由于 $U$ 拟紧，可以选择
$(i_1,j_1,k_1),\ldots,(i_n,j_n,k_n)$，使得
$$
X_{i_1 j_1 k_1} \amalg \ldots \amalg X_{i_n j_n k_n} \longrightarrow U
$$
满。对 $s=1,\ldots,n$，考虑交换图表
$$
\xymatrix{
& & & X_{i_s j_s k_s} \ar[ld] \ar[d] \ar[rd] \\
X \ar[d] & X_{i_s} \ar[l] \ar[d] &
X_{i_s} \times_X U \ar[l] \ar[d] & Y_{i_s j_s} \ar[ld] \ar[rd] &
X_{i_s} \times_X U \ar[d] \ar[r] &
U \ar[d] \ar[r] & X \ar[d] \\
Y & Y_{i_s} \ar[l] &
Y_{i_s} \times_Y V \ar[l] & &
Y_{i_s} \times_Y V \ar[r] &
V \ar[r] & Y
}
$$
若可数指标化仿射形式代数空间的一个态射所对应的
$\textit{WAdm}^{count}$ 态射具有 $P$，就称该态射具有 $P$。
注意，映射 $X_{i_sj_sk_s}\to X_{i_s}$ 和
$Y_{i_sj_s}\to Y_{i_s}$ 由 étale 环映射的完备化给出，
见引理 \ref{formal-spaces-lemma-etale}。所以由公理 (\ref{formal-spaces-item-axiom-1})，
$P(X_{i_s}\to Y_{i_s})$ 蕴含
$P(X_{i_sj_sk_s}\to Y_{i_sj_s})$。
注意，映射 $Y_{i_sj_s}\to V$ 由 étale 环映射的完备化给出
% P10-E0333
（仍用上一引理）。把公理 (\ref{formal-spaces-item-axiom-2}) 用于图表
$$
\xymatrix{
X_{i_s j_s k_s} \ar@{=}[r] \ar[d] & X_{i_s j_s k_s} \ar[d] \\
Y_{i_s j_s} \ar[r] & V
}
$$
（这是允许的，因为恒等映射是忠实平坦环映射），可得
$P(X_{i_sj_sk_s}\to V)$ 成立。由公理 (\ref{formal-spaces-item-axiom-3})，
$P(\coprod_{s=1,\ldots,n}X_{i_sj_sk_s}\to V)$ 成立。
按构造，态射 $\coprod X_{i_sj_sk_s}\to U$ 满，
所以由引理 \ref{formal-spaces-lemma-etale-surjective}，相应的
$\textit{WAdm}^{count}$ 态射是忠实平坦 étale 环映射的完备化。
再应用一次公理 (\ref{formal-spaces-item-axiom-2})（取 $B'=B$），
即可如所需得到 $P(U\to V)$ 成立。
\end{proof}

\begin{remark}[adic-star 变体]
\label{formal-spaces-remark-variant-adic-star}
设 $P$ 是 $\textit{WAdm}^{adic*}$ 态射的性质。
若公理 (\ref{formal-spaces-item-axiom-1})、(\ref{formal-spaces-item-axiom-2})、
(\ref{formal-spaces-item-axiom-3})（见情形 \ref{formal-spaces-situation-local-property}）
对 $\textit{WAdm}^{adic*}$ 的态射成立，则称 $P$ 是\textit{局部性质}。
以完全相同的方式，可得到引理
\ref{formal-spaces-lemma-property-defines-property-morphisms} 对 $S$ 上局部 adic*
形式代数空间之间态射的变体。
\end{remark}

\begin{remark}[Noether 变体]
\label{formal-spaces-remark-variant-Noetherian}
设 $P$ 是 $\textit{WAdm}^{Noeth}$ 态射的性质。
% P10-E0334
若公理 (\ref{formal-spaces-item-axiom-1})、(\ref{formal-spaces-item-axiom-2})、
(\ref{formal-spaces-item-axiom-3})（见情形 \ref{formal-spaces-situation-local-property}）
对 $\textit{WAdm}^{Noeth}$ 的态射成立，则称 $P$ 是\textit{局部性质}。
以完全相同的方式，可得到引理
\ref{formal-spaces-lemma-property-defines-property-morphisms} 对 $S$ 上局部 Noether
形式代数空间之间态射的变体。
\end{remark}

\begin{situation}
\label{formal-spaces-situation-base-change-local-property}
设 $P$ 是 $\textit{WAdm}^{count}$ 态射的局部性质，见情形
\ref{formal-spaces-situation-local-property}。若对 $\textit{WAdm}^{count}$ 中给定的
$B\to A$ 和 $B\to C$，有
$P(B\to A)\Rightarrow P(C\to A\widehat{\otimes}_BC)$，
则称 $P$ 在\textit{基变换下稳定}。
这是有意义的，因为由引理 \ref{formal-spaces-lemma-completed-tensor-product}，
$A\widehat{\otimes}_BC$ 是 $\textit{WAdm}^{count}$ 的对象。
\end{situation}

\begin{lemma}
\label{formal-spaces-lemma-base-change-property-morphisms}
设 $S$ 是概形，$P$ 是在基变换下稳定的
$\textit{WAdm}^{count}$ 态射的局部性质。
设 $f:X\to Y$ 和 $g:Z\to Y$ 是 $S$ 上局部可数指标化
形式代数空间的态射。若 $f$ 满足引理
\ref{formal-spaces-lemma-property-defines-property-morphisms} 的等价条件，
则 $\text{pr}_2:X\times_YZ\to Z$ 亦然。
\end{lemma}

\begin{proof}
选取定义 \ref{formal-spaces-definition-formal-algebraic-space} 意义下的覆盖
$\{Y_j\to Y\}$。对每个 $j$，选取定义
\ref{formal-spaces-definition-formal-algebraic-space} 意义下的覆盖
$\{X_{ji}\to Y_j\times_YX\}$，并选取定义
\ref{formal-spaces-definition-formal-algebraic-space} 意义下的覆盖
$\{Z_{jk}\to Y_j\times_YZ\}$。
由引理 \ref{formal-spaces-lemma-types-fibre-products}，
$X_{ji}\times_{Y_j}Z_{jk}$ 是可数指标化仿射形式代数空间。
于是
$$
\{X_{ji} \times_{Y_j} Z_{jk} \to X \times_Y Z\}
$$
是定义 \ref{formal-spaces-definition-formal-algebraic-space} 意义下的覆盖。
此外，态射 $X_{ji}\times_{Y_j}Z_{jk}\to Z$ 分解经过 $Z_{jk}$。
由假设，$X_{ji}\to Y_j$ 对应于具有性质 $P$ 的
$\textit{WAdm}^{count}$ 态射 $B_j\to A_{ji}$。
态射 $Z_{jk}\to Y_j$ 对应于 $\textit{WAdm}^{count}$ 中的态射
$B_j\to C_{jk}$。由引理
\ref{formal-spaces-lemma-fibre-product-affines-over-separated}，
$X_{ji} \times_{Y_j} Z_{jk} =
\text{Spf}(A_{ji} \widehat{\otimes}_{B_j} C_{jk})$，
所以只需证明
$C_{jk}\to A_{ji}\widehat{\otimes}_{B_j}C_{jk}$ 具有性质 $P$；
而这恰由 $P$ 在基变换下稳定这一条件保证。
\end{proof}

\begin{remark}[adic-star 变体]
\label{formal-spaces-remark-base-change-variant-adic-star}
设 $P$ 是 $\textit{WAdm}^{adic*}$ 态射的局部性质，见注
\ref{formal-spaces-remark-variant-adic-star}。若对 $\textit{WAdm}^{adic*}$ 中给定的
$B\to A$ 和 $B\to C$，有
$P(B\to A)\Rightarrow P(C\to A\widehat{\otimes}_BC)$，
则称 $P$ 在\textit{基变换下稳定}。
这是有意义的，因为由引理 \ref{formal-spaces-lemma-completed-tensor-product}，
$A\widehat{\otimes}_BC$ 是 $\textit{WAdm}^{adic*}$ 的对象。
以完全相同的方式，可得到引理
\ref{formal-spaces-lemma-base-change-property-morphisms} 对 $S$ 上局部 adic*
形式代数空间之间态射的变体。
\end{remark}

\begin{remark}[Noether 变体]
\label{formal-spaces-remark-base-change-variant-Noetherian}
设 $P$ 是 $\textit{WAdm}^{Noeth}$ 态射的局部性质，见注
\ref{formal-spaces-remark-variant-Noetherian}。若对 $\textit{WAdm}^{Noeth}$ 中给定的
$B\to A$ 和 $B\to C$，性质 $P(B\to A)$ 同时蕴含
$A\widehat{\otimes}_BC$ 是 Noether adic 的
\footnote{判据见引理 \ref{formal-spaces-lemma-completed-tensor-product}。}
以及 $P(C\to A\widehat{\otimes}_BC)$，则称 $P$ 在\textit{基变换下稳定}。
以完全相同的方式，可得到引理
\ref{formal-spaces-lemma-base-change-property-morphisms} 对 $S$ 上局部 Noether
形式代数空间之间态射的变体。
\end{remark}

\begin{remark}[另一个 Noether 变体]
\label{formal-spaces-remark-base-change-variant-variant-Noetherian}
设 $P,Q$ 是 $\textit{WAdm}^{Noeth}$ 态射的局部性质，见注
\ref{formal-spaces-remark-variant-Noetherian}。若对 $\textit{WAdm}^{Noeth}$ 中给定的
$B\to A$ 和 $B\to C$，在 $P(B\to A)$ 与 $Q(B\to C)$ 成立时，
$A\widehat{\otimes}_BC$ 是 Noether adic 的并且
$P(C\to A\widehat{\otimes}_BC)$ 成立，
则称\textit{$P$ 在按 $Q$ 的基变换下稳定}。
与引理 \ref{formal-spaces-lemma-base-change-property-morphisms} 的证明完全相同地论证，
得到如下陈述：给定 $S$ 上局部 Noether 形式代数空间的态射
% P10-E0335
$f:X\to Y$ 和 $g:Y\to Z$，使得
\begin{enumerate}
\item 引理 \ref{formal-spaces-lemma-property-defines-property-morphisms}
关于 $f$ 和 $P$ 的等价条件成立；
\item 引理 \ref{formal-spaces-lemma-property-defines-property-morphisms}
关于 $g$ 和 $Q$ 的等价条件成立，
\end{enumerate}
则引理 \ref{formal-spaces-lemma-property-defines-property-morphisms}
关于 $\text{pr}_2:X\times_YZ\to Z$ 和 $P$ 的等价条件成立。
\end{remark}

\begin{situation}
\label{formal-spaces-situation-composition-local-property}
设 $P$ 是 $\textit{WAdm}^{count}$ 态射的局部性质，见情形
\ref{formal-spaces-situation-local-property}。若对 $\textit{WAdm}^{count}$ 中给定的
$B\to A$ 和 $C\to B$，有
$P(B\to A)\wedge P(C\to B)\Rightarrow P(C\to A)$，
则称 $P$ 在\textit{复合下稳定}。
\end{situation}

\begin{lemma}
\label{formal-spaces-lemma-composition-property-morphisms}
设 $S$ 是概形，$P$ 是在复合下稳定的
$\textit{WAdm}^{count}$ 态射的局部性质。
设 $f:X\to Y$ 和 $g:Y\to Z$ 是 $S$ 上局部可数指标化形式代数空间的态射。
% P10-E0336
若 $f,g$ 满足引理 \ref{formal-spaces-lemma-property-defines-property-morphisms}
的等价条件，则 $g\circ f:X\to Z$ 亦然。
\end{lemma}

\begin{proof}
选取定义 \ref{formal-spaces-definition-formal-algebraic-space} 意义下的覆盖
$\{Z_k\to Z\}$。对每个 $k$，选取定义
\ref{formal-spaces-definition-formal-algebraic-space} 意义下的覆盖
$\{Y_{kj}\to Z_k\times_ZY\}$；再对每个 $k,j$ 选取定义
\ref{formal-spaces-definition-formal-algebraic-space} 意义下的覆盖
$\{X_{kji}\to Y_{kj}\times_YX\}$。
% P10-E0337
若 $f,g$ 满足引理 \ref{formal-spaces-lemma-property-defines-property-morphisms}
的等价条件，则
% P10-E0338
$X_{kji}\to Y_{jk}$ 和 $Y_{jk}\to Z_k$ 分别对应于
$\textit{WAdm}^{count}$ 中具有性质 $P$ 的箭头
$B_{kj}\to A_{kji}$ 和 $C_k\to B_{kj}$。
因此复合也具有该性质，结论成立。
\end{proof}

\begin{remark}[adic-star 变体]
\label{formal-spaces-remark-composition-variant-adic-star}
设 $P$ 是 $\textit{WAdm}^{adic*}$ 态射的局部性质，见注
\ref{formal-spaces-remark-variant-adic-star}。若对 $\textit{WAdm}^{adic*}$ 中给定的
$B\to A$ 和 $C\to B$，有
$P(B\to A)\wedge P(C\to B)\Rightarrow P(C\to A)$，
则称 $P$ 在\textit{复合下稳定}。
以完全相同的方式，可得到引理
\ref{formal-spaces-lemma-composition-property-morphisms} 对 $S$ 上局部 adic*
形式代数空间之间态射的变体。
\end{remark}

\begin{remark}[Noether 变体]
\label{formal-spaces-remark-composition-variant-Noetherian}
设 $P$ 是 $\textit{WAdm}^{Noeth}$ 态射的局部性质，见注
\ref{formal-spaces-remark-variant-Noetherian}。若对 $\textit{WAdm}^{Noeth}$ 中给定的
$B\to A$ 和 $C\to B$，有
$P(B\to A)\wedge P(C\to B)\Rightarrow P(C\to A)$，
则称 $P$ 在\textit{复合下稳定}。
以完全相同的方式，可得到引理
\ref{formal-spaces-lemma-composition-property-morphisms} 对 $S$ 上局部 Noether
形式代数空间之间态射的变体。
\end{remark}

\begin{situation}
\label{formal-spaces-situation-permanence-local-property}
设 $P$ 是 $\textit{WAdm}^{count}$ 态射的局部性质，见情形
\ref{formal-spaces-situation-local-property}。若对 $\textit{WAdm}^{count}$ 中给定的
$B\to A$ 和 $C\to B$，有
$P(C\to B)\wedge P(C\to A)\Rightarrow P(B\to A)$，
则称 $P$ \textit{具有消去性质}。
\end{situation}

\begin{lemma}
\label{formal-spaces-lemma-permanence-property-morphisms}
设 $S$ 是概形，$P$ 是具有消去性质的
$\textit{WAdm}^{count}$ 态射的局部性质。
设 $f:X\to Y$ 和 $g:Y\to Z$ 是 $S$ 上局部可数指标化形式代数空间的态射。
% P10-E0339
若 $g\circ f$ 和 $g$ 满足引理
\ref{formal-spaces-lemma-property-defines-property-morphisms} 的等价条件，
则 $f:X\to Y$ 亦然。
\end{lemma}

\begin{proof}
选取定义 \ref{formal-spaces-definition-formal-algebraic-space} 意义下的覆盖
$\{Z_k\to Z\}$。对每个 $k$，选取定义
\ref{formal-spaces-definition-formal-algebraic-space} 意义下的覆盖
$\{Y_{kj}\to Z_k\times_ZY\}$；再对每个 $k,j$ 选取定义
\ref{formal-spaces-definition-formal-algebraic-space} 意义下的覆盖
$\{X_{kji}\to Y_{kj}\times_YX\}$。
% P10-E0340
设 $X_{kji}\to Y_{jk}$ 和 $Y_{jk}\to Z_k$ 分别对应于
$\textit{WAdm}^{count}$ 中的箭头
$B_{kj}\to A_{kji}$ 和 $C_k\to B_{kj}$。
% P10-E0341
若 $g\circ f$ 和 $g$ 满足引理
\ref{formal-spaces-lemma-property-defines-property-morphisms} 的等价条件，
则 $C_k\to B_{kj}$ 和 $C_k\to A_{kji}$ 具有性质 $P$。
因此 $B_{kj}\to A_{kji}$ 也具有性质 $P$，结论成立。
\end{proof}

\begin{remark}[adic-star 变体]
\label{formal-spaces-remark-permanence-variant-adic-star}
设 $P$ 是 $\textit{WAdm}^{adic*}$ 态射的局部性质，见注
\ref{formal-spaces-remark-variant-adic-star}。若对 $\textit{WAdm}^{adic*}$ 中给定的
$B\to A$ 和 $C\to B$，有
$P(C\to A)\wedge P(C\to B)\Rightarrow P(B\to A)$，
则称 $P$ \textit{具有消去性质}。
以完全相同的方式，可得到引理
% P10-E0342
\ref{formal-spaces-lemma-composition-property-morphisms} 对 $S$ 上局部 adic*
形式代数空间之间态射的变体。
\end{remark}

\begin{remark}[Noether 变体]
\label{formal-spaces-remark-permanence-variant-Noetherian}
设 $P$ 是 $\textit{WAdm}^{Noeth}$ 态射的局部性质，见注
\ref{formal-spaces-remark-variant-Noetherian}。若对 $\textit{WAdm}^{Noeth}$ 中给定的
$B\to A$ 和 $C\to B$，有
% P10-E0343
$P(C\to B)\wedge P(C\to A)\Rightarrow P(C\to B)$，
则称 $P$ \textit{具有消去性质}。
以完全相同的方式，可得到引理
% P10-E0344
\ref{formal-spaces-lemma-composition-property-morphisms} 对 $S$ 上局部 Noether
形式代数空间之间态射的变体。
\end{remark}













\section{Taut 环映射与由代数空间表示}
\label{formal-spaces-section-taut}

\noindent
% P10-E0345
本节简要证明，局部可数指标化形式代数空间之间的态射，
étale 局部地对应于具有可数开理想基本系的弱可容许拓扑环之间的
taut 连续环同态。事实上，这从引理 \ref{formal-spaces-lemma-representable-affine}
已经相当清楚，我们鼓励读者跳过本节。

\begin{lemma}
\label{formal-spaces-lemma-representable-property-rings}
设 $B\to A$ 是 $\textit{WAdm}^{count}$ 中的箭头。下列条件等价：
\begin{enumerate}
\item[(a)] $B\to A$ 是 taut 的（定义 \ref{formal-spaces-definition-taut}）；
% P10-E0346
\item[(b)] 对弱定义理想基本系
$B\supset J_1\supset J_2\supset J_3\supset\ldots$，存在交换图表
$$
\xymatrix{
A \ar[r] & \ldots \ar[r] & A_3 \ar[r] & A_2 \ar[r] & A_1 \\
B \ar[r] \ar[u] & \ldots \ar[r] & B/J_3 \ar[r] \ar[u] &
B/J_2 \ar[r] \ar[u] & B/J_1 \ar[u]
}
$$
使得 $A_{n+1}/J_nA_{n+1}=A_n$，并且
% P10-E0347
$A=\lim A_n$ 作为拓扑环成立。
\end{enumerate}
此外，这些等价条件定义一个局部性质；即它们满足公理
(\ref{formal-spaces-item-axiom-1})、(\ref{formal-spaces-item-axiom-2})、(\ref{formal-spaces-item-axiom-3})。
\end{lemma}

\begin{proof}
(a) 与 (b) 的等价性是立即的。下面将给出这些公理的代数证明，
不过事实上我们已经证明过它们。具体地，利用引理
\ref{formal-spaces-lemma-representable-affine}，等价条件 (a)、(b) 转化为：
相应的仿射形式代数空间态射可由代数空间表示。
由引理 \ref{formal-spaces-lemma-representable-by-algebraic-spaces-local}，
这个条件“在源和靶上 étale 局部”，所以立即得到公理
(\ref{formal-spaces-item-axiom-1})、(\ref{formal-spaces-item-axiom-2}) 和
(\ref{formal-spaces-item-axiom-3})。

\medskip\noindent
公理 (\ref{formal-spaces-item-axiom-1})、(\ref{formal-spaces-item-axiom-2})、
(\ref{formal-spaces-item-axiom-3}) 的直接代数证明。
给定图表 (\ref{formal-spaces-equation-localize})，如情形
\ref{formal-spaces-situation-local-property} 所述。由例
\ref{formal-spaces-example-representable-morphism-from-completion}，
映射 $A\to(A')^\wedge$ 和 $B\to(B')^\wedge$ 满足 (a)、(b)。

\medskip\noindent
假设 $\varphi$ 满足 (a)、(b)。设 $J\subset B$ 是弱定义理想。
则 $JA$ 与 $J(B')^\wedge$ 的闭包分别是弱定义理想
$I\subset A$ 与 $J'\subset(B')^\wedge$。
于是 $I(A')^\wedge$ 的闭包是弱定义理想
$I'\subset(A')^\wedge$。拓扑论论证表明，$I'$ 也是
$J(A')^\wedge$ 与 $J'(A')^\wedge$ 的闭包。
最后，当 $J$ 遍历 $B$ 的弱定义理想基本系时，
$I$ 与 $I'$ 分别遍历 $A$ 与 $(A')^\wedge$ 中的弱定义理想基本系。
所以 (a) 对 $\varphi'$ 成立。这证明了 (\ref{formal-spaces-item-axiom-1})。

\medskip\noindent
假设 $A\to A'$ 忠实平坦，且 $\varphi'$ 满足 (a)、(b)。
设 $J\subset B$ 是弱定义理想。把 (a)、(b) 用于映射
$B\to(B')^\wedge\to(A')^\wedge$，可知
$J(A')^\wedge$ 的闭包 $I'$ 是弱定义理想。
特别地，$I'$ 是开的，故 $I'$ 在 $A$ 中的逆像是开的。
现在有（解释见下）
\begin{align*}
A \cap I'
& =
A \cap \bigcap (J(A')^\wedge + \Ker((A')^\wedge \to A'/I_0A')) \\
& =
A \cap \bigcap \Ker((A')^\wedge \to A'/JA' + I_0 A') \\
& = \bigcap (JA + I_0)
\end{align*}
由引理 \ref{formal-spaces-lemma-closed}，这就是 $JA$ 的闭包。
各交遍历弱定义理想 $I_0\subset A$。
% P10-E0348
第一个等式是因为 $(A')^\wedge$ 中 $0$ 的一族邻域基本系由映射
$(A')^\wedge\to A'/I_0A'$ 的核组成。第二个等式显然。
% P10-E0349
第三个等式是因为 $A\to A'$ 忠实平坦，见《代数》之引理
\ref{algebra-lemma-faithfully-flat-universally-injective}。
因此 $JA$ 的闭包是开的。由引理
\ref{formal-spaces-lemma-topologically-nilpotent}，$JA$ 的闭包是 $A$ 的弱定义理想。
最后，给定弱定义理想 $I\subset A$，由 $B\to(B')^\wedge$
和 $\varphi'$ 的性质 (a)，可找到 $J$，使
$J(A')^\wedge$ 包含在 $I(A')^\wedge$ 的闭包中。
所以 (a) 对 $\varphi$ 成立。这证明了 (\ref{formal-spaces-item-axiom-2})。

\medskip\noindent
我们略去 (\ref{formal-spaces-item-axiom-3}) 的证明。
\end{proof}

\begin{lemma}
\label{formal-spaces-lemma-representable-local-property}
设 $S$ 是概形，$f:X\to Y$ 是 $S$ 上局部可数指标化形式代数空间的态射。
下列条件等价：
\begin{enumerate}
\item 对每个交换图表
$$
\xymatrix{
U \ar[d] \ar[r] & V \ar[d] \\
X \ar[r] & Y
}
$$
若 $U,V$ 是仿射形式代数空间，且 $U\to X$ 和 $V\to Y$
可由代数空间表示并为 étale，则态射 $U\to V$ 对应于
$\textit{WAdm}^{count}$ 中的 taut 映射 $B\to A$；
\item 存在定义 \ref{formal-spaces-definition-formal-algebraic-space} 意义下的覆盖
$\{Y_j\to Y\}$，并且对每个 $j$，存在定义
\ref{formal-spaces-definition-formal-algebraic-space} 意义下的覆盖
$\{X_{ji}\to Y_j\times_YX\}$，使每个 $X_{ji}\to Y_j$
都对应于 $\textit{WAdm}^{count}$ 中的 taut 环映射；
\item 存在定义 \ref{formal-spaces-definition-formal-algebraic-space} 意义下的覆盖
$\{X_i\to X\}$，并且对每个 $i$ 存在分解
$X_i\to Y_i\to Y$，其中 $Y_i$ 是仿射形式代数空间，
$Y_i\to Y$ 可由代数空间表示并为 étale，且 $X_i\to Y_i$
对应于 $\textit{WAdm}^{count}$ 中的 taut 环映射；以及
\item $f$ 可由代数空间表示。
\end{enumerate}
\end{lemma}

\begin{proof}
由引理 \ref{formal-spaces-lemma-representable-property-rings}，
$\textit{WAdm}^{count}$ 中映射的“taut”性质是局部性质。
所以引理 \ref{formal-spaces-lemma-property-defines-property-morphisms}
恰好告诉我们 (1)、(2)、(3) 等价。
另一方面，由引理 \ref{formal-spaces-lemma-representable-affine}，
$\textit{WAdm}^{count}$ 中映射为“taut”恰好对应于相应的可数指标化
仿射形式代数空间态射“可由代数空间表示”。
% P10-E0350
因此，引理 \ref{formal-spaces-lemma-representable-by-algebraic-spaces-local} 的蕴含
(1) $\Rightarrow$ (2) 表明本引理的 (4) 蕴含 (1)。
类似地，引理 \ref{formal-spaces-lemma-representable-by-algebraic-spaces-local} 的蕴含
(4) $\Rightarrow$ (1) 表明本引理的 (2) 蕴含 (4)。
\end{proof}








\section{Adic 态射}
\label{formal-spaces-section-adic}

\noindent
本节把局部 adic* 形式代数空间 $X,Y$ 之间偶尔使用的“adic 态射”
$f:X\to Y$ 这一概念，一方面与 $f$ 可由代数空间表示相对应，
另一方面与我们的 taut 连续环同态概念相对应。
首先回忆：对具有有限生成定义理想的 adic 环，taut 性等价于 adic 性。

\begin{lemma}
\label{formal-spaces-lemma-adic-homomorphism}
设 $A,B$ 是 pre-adic 拓扑环，
$\varphi:A\to B$ 是连续环同态。
\begin{enumerate}
\item 若 $\varphi$ 是 adic 的，则 $\varphi$ 是 taut 的。
\item 若 $B$ 完备，$A$ 有有限生成定义理想，且 $\varphi$ 是 taut 的，
则 $\varphi$ 是 adic 的。
\end{enumerate}
特别地，在范畴 $\textit{WAdm}^{adic*}$ 上，“$\varphi$ 是 adic 的”
与“$\varphi$ 是 taut 的”这两个条件等价。
\end{lemma}

\begin{proof}
(1) 是引理 \ref{formal-spaces-lemma-adic-taut}，(2) 是引理
\ref{formal-spaces-lemma-taut-is-adic}。最后的陈述由 (1)、(2) 得出。
\end{proof}

\noindent
设 $S$ 是概形，$f:X\to Y$ 是 $S$ 上局部 adic* 形式代数空间的态射。
由引理 \ref{formal-spaces-lemma-representable-local-property}，下列条件等价：
\begin{enumerate}
\item $f$ 可由代数空间表示（换言之，引理
\ref{formal-spaces-lemma-representable-by-algebraic-spaces-local} 的等价条件成立）；
\item 对每个交换图表
$$
\xymatrix{
U \ar[d] \ar[r] & V \ar[d] \\
X \ar[r] & Y
}
$$
若 $U,V$ 是仿射形式代数空间，且 $U\to X$ 和 $V\to Y$
可由代数空间表示并为 étale，则态射 $U\to V$ 对应于
$\textit{WAdm}^{adic*}$ 中的 adic 映射
\footnote{由引理 \ref{formal-spaces-lemma-adic-homomorphism}，这等价于 taut。}。
\end{enumerate}
在这种情形中，我们将称\textit{$f$ 是 adic 态射}
（正式定义见下）。这一概念或术语只对局部 adic* 形式代数空间之间的态射定义和使用，
因为否则我们没有上述 (1)、(2) 的等价性。

\begin{definition}
\label{formal-spaces-definition-adic-morphism}
设 $S$ 是概形，$f:X\to Y$ 是 $S$ 上形式代数空间的态射。
假设 $X,Y$ 局部 adic*。若 $f$ 可由代数空间表示，
则称 $f$ 是\textit{adic 态射}。见上面的讨论。
\end{definition}









\section{有限型态射}
\label{formal-spaces-section-finite-type}

\noindent
% P10-E0351
鉴于 Stacks Project 中各项内容的设置方式，下面的做法才是真正正确的，
更强的概念应当另取名称。

\begin{definition}
\label{formal-spaces-definition-finite-type}
设 $S$ 是概形，$f:Y\to X$ 是 $S$ 上形式代数空间的态射。
\begin{enumerate}
\item 若 $f$ 可由代数空间表示，并且在《自举》之定义
\ref{bootstrap-definition-property-transformation} 的意义下局部有限型，
则称 $f$ \textit{局部有限型}。
\item 若 $f$ 局部有限型且拟紧
（定义 \ref{formal-spaces-definition-quasi-compact-morphism}），
则称 $f$ \textit{有限型}。
\end{enumerate}
\end{definition}

\noindent
我们将在第 \ref{formal-spaces-section-tft} 节讨论某些形式代数空间的有限型态射
与拓扑有限型连续环映射 $A\to B$ 之间的关系。

\begin{lemma}
\label{formal-spaces-lemma-characterize-finite-type}
设 $S$ 是概形，$f:X\to Y$ 是 $S$ 上形式代数空间的态射。
下列条件等价：
\begin{enumerate}
\item $f$ 有限型；
\item $f$ 可由代数空间表示，并且在《自举》之定义
\ref{bootstrap-definition-property-transformation} 的意义下有限型。
\end{enumerate}
\end{lemma}

\begin{proof}
这由《自举》之引理
\ref{bootstrap-lemma-transformations-property-implication}、
代数空间态射的蕴含“拟紧 $+$ 局部有限型 $\Rightarrow$ 有限型”，
以及引理 \ref{formal-spaces-lemma-quasi-compact-representable} 得出。
\end{proof}

\begin{lemma}
\label{formal-spaces-lemma-composition-finite-type}
有限型态射的复合有限型。对局部有限型同样成立。
\end{lemma}

\begin{proof}
见《自举》之引理
\ref{bootstrap-lemma-composition-transformation-property}，
并使用《空间的态射》之引理
\ref{spaces-morphisms-lemma-composition-finite-type}。
\end{proof}

\begin{lemma}
\label{formal-spaces-lemma-base-change-finite-type}
有限型态射的基变换有限型。对局部有限型同样成立。
\end{lemma}

\begin{proof}
见《自举》之引理
\ref{bootstrap-lemma-base-change-transformation-property}，
并使用《空间的态射》之引理
\ref{spaces-morphisms-lemma-base-change-finite-type}。
\end{proof}

\begin{lemma}
\label{formal-spaces-lemma-permanence-finite-type}
设 $S$ 是概形，$f:X\to Y$ 和 $g:Y\to Z$ 是 $S$ 上形式代数空间的态射。
若 $g\circ f:X\to Z$ 局部有限型，则 $f:X\to Y$ 局部有限型。
\end{lemma}

\begin{proof}
由引理 \ref{formal-spaces-lemma-permanence-representable}，$f$ 可由代数空间表示。
设 $T$ 是概形且 $T\to Z$ 是态射。把《空间的态射》之引理
\ref{spaces-morphisms-lemma-permanence-finite-type}
用于代数空间的态射
$T\times_ZX\to T\times_ZY\to T$，即可得出结论。
\end{proof}

\noindent
局部有限型这一性质在源和靶上都是局部的。

\begin{lemma}
\label{formal-spaces-lemma-finite-type-local}
设 $S$ 是概形，$f:X\to Y$ 是 $S$ 上形式代数空间的态射。
下列条件等价：
\begin{enumerate}
\item 态射 $f$ 局部有限型；
\item 存在交换图表
$$
\xymatrix{
U \ar[d] \ar[r] & V \ar[d] \\
X \ar[r] & Y
}
$$
其中 $U,V$ 是形式代数空间，竖直箭头可由代数空间表示并为 étale，
$U\to X$ 满，且 $U\to V$ 局部有限型；
\item 对任意交换图表
$$
\xymatrix{
U \ar[d] \ar[r] & V \ar[d] \\
X \ar[r] & Y
}
$$
若 $U,V$ 是形式代数空间，且竖直箭头可由代数空间表示并为 étale，
则态射 $U\to V$ 局部有限型；
\item 存在定义 \ref{formal-spaces-definition-formal-algebraic-space} 意义下的覆盖
$\{Y_j\to Y\}$，并且对每个 $j$，存在定义
\ref{formal-spaces-definition-formal-algebraic-space} 意义下的覆盖
$\{X_{ji}\to Y_j\times_YX\}$，使得对每个 $j,i$，
$X_{ji}\to Y_j$ 局部有限型；
\item 存在定义 \ref{formal-spaces-definition-formal-algebraic-space} 意义下的覆盖
$\{X_i\to X\}$，并且对每个 $i$ 存在分解
$X_i\to Y_i\to Y$，其中 $Y_i$ 是仿射形式代数空间，
$Y_i\to Y$ 可由代数空间表示并为 étale，且 $X_i\to Y_i$ 局部有限型；以及
% P10-E0352
\item 此处还可补充更多条件。
\end{enumerate}
\end{lemma}

\begin{proof}
% P10-E0353
在这 5 种情形中的每一种，态射 $f:X\to Y$ 都可由代数空间表示，
见引理 \ref{formal-spaces-lemma-representable-by-algebraic-spaces-local}。
以下将不再说明地使用这一点。

\medskip\noindent
显然 (1) 蕴含 (2)，因为可以取 $U=X$ 和 $V=Y$。
反之，假设给定 (2) 中的图表。设 $T$ 是概形且 $T\to Y$ 是态射。
则可以考虑
$$
\xymatrix{
U \times_Y T \ar[d] \ar[r] & V \times_Y T \ar[d] \\
X \times_Y T \ar[r] & T
}
$$
竖直箭头为 étale，顶部水平箭头局部有限型，因为它们是相应态射的基变换。
% P10-E0354
所以由《空间的态射》之引理
\ref{spaces-morphisms-lemma-finite-type-local}，
$X\times_YT\to T$ 局部有限型。换言之，(1) 成立。

\medskip\noindent
假设 (1) 成立，并考虑 (3) 中的图表。则 $U\to Y$ 局部有限型
（它是 $U\to X\to Y$ 的复合；见《自举》之引理
\ref{bootstrap-lemma-composition-transformation-property}）。
设 $T$ 是概形且 $T\to V$ 是态射。则投影
$T\times_VU\to T$ 分解为
$$
T \times_V U = (T \times_Y U) \times_{(V \times_Y V)} V
\to T \times_Y U \to T
$$
第二个箭头局部有限型（它是复合 $U\to X\to Y$ 的基变换），
第一个箭头是对角态射 $V\to V\times_YV$ 的基变换；
由引理 \ref{formal-spaces-lemma-diagonal-morphism-formal-algebraic-spaces}，
该对角态射局部有限型。

\medskip\noindent
显然 (3) 蕴含 (2)。因此 (1)--(3) 现已等价。

\medskip\noindent
注意，(4) 中的条件是有意义的，因为由引理
\ref{formal-spaces-lemma-fibre-products}，纤维积 $Y_j\times_YX$ 是形式代数空间。
显然 (4) 蕴含 (5)。

\medskip\noindent
假设有 (5) 中的 $X_i\to Y_i\to Y$。令
$V=\coprod Y_i$ 且 $U=\coprod X_i$，即可看出 (5) 蕴含 (2)。

\medskip\noindent
最后，假设 (1)--(3) 成立。于是可以任取定义
\ref{formal-spaces-definition-formal-algebraic-space} 意义下的覆盖
$\{Y_j\to Y\}$，并对每个 $j$ 任取定义
\ref{formal-spaces-definition-formal-algebraic-space} 意义下的覆盖
$\{X_{ji}\to Y_j\times_YX\}$。
% P10-E0355
由 (3)，$X_{ij}\to Y_j$ 局部有限型，故 (4) 成立。证明完成。
\end{proof}

\begin{example}
\label{formal-spaces-example-finite-type-from-finite-type-ring-map}
设 $S$ 是概形，$A$ 是 $S$ 上的弱可容许拓扑环，
$A\to A'$ 是有限型环映射。则
$$
(A')^\wedge = \lim_{I \subset A\ w.i.d.} A'/IA'
$$
是弱可容许环，且相应的态射
$\text{Spf}((A')^\wedge)\to\text{Spf}(A)$ 可表示，
见例 \ref{formal-spaces-example-representable-morphism-from-completion}。
若 $T\to\text{Spf}(A)$ 是态射且 $T$ 是拟紧概形，则由引理
\ref{formal-spaces-lemma-factor-through-thickening}，它对某个弱定义理想
$I\subset A$ 分解经过 $\Spec(A/I)$。于是
$T\times_{\text{Spf}(A)}\text{Spf}((A')^\wedge)$ 等于
$T\times_{\Spec(A/I)}\Spec(A'/IA')$，故
$\text{Spf}((A')^\wedge)\to\text{Spf}(A)$ 有限型。
\end{example}

\begin{lemma}
\label{formal-spaces-lemma-locally-finite-type-locally-noetherian}
设 $S$ 是概形，$f:X\to Y$ 是 $S$ 上形式代数空间的态射。
% P10-E0356
若 $Y$ 局部 Noether 且 $f$ 局部有限型，则 $X$ 局部 Noether。
\end{lemma}

\begin{proof}
按引理 \ref{formal-spaces-lemma-finite-type-local} 选取
$\{Y_j\to Y\}$ 和 $\{X_{ij}\to Y_j\times_YX\}$。
由引理 \ref{formal-spaces-lemma-property-goes-up-affine-morphism}，
每个 $X_{ij}$ 都是 Noether 的。引理得证。
\end{proof}

\begin{lemma}
\label{formal-spaces-lemma-fibre-product-Noetherian}
设 $S$ 是概形，$f:X\to Y$ 和 $Z\to Y$ 是 $S$ 上形式代数空间的态射。
若 $Z$ 局部 Noether 且 $f$ 局部有限型，则
$Z\times_YX$ 局部 Noether。
\end{lemma}

\begin{proof}
由引理 \ref{formal-spaces-lemma-base-change-finite-type}，
态射 $Z\times_YX\to Z$ 局部有限型。
所以结论由引理 \ref{formal-spaces-lemma-locally-finite-type-locally-noetherian} 得出。
\end{proof}








\section{满态射}
\label{formal-spaces-section-surjective}

\noindent
由引理 \ref{formal-spaces-lemma-reduction-surjective}，下述定义不会与已有的代数空间态射定义，
或可由代数空间表示的形式代数空间态射定义冲突。

\begin{definition}
\label{formal-spaces-definition-surjective}
设 $S$ 为概形。若 $S$ 上形式代数空间的态射 $f:X\to Y$
在底层约化代数空间上诱导满态射 $X_{red}\to Y_{red}$，
则称 $f$ 为\textit{满态射}。
\end{definition}

\begin{lemma}
\label{formal-spaces-lemma-composition-surjective}
两个满态射的复合是满态射。
\end{lemma}

\begin{proof}
省略。
\end{proof}

\begin{lemma}
\label{formal-spaces-lemma-base-change-surjective}
满态射的基变换是满态射。
\end{lemma}

\begin{proof}
省略。
\end{proof}

\begin{lemma}
\label{formal-spaces-lemma-characterize-surjective}
设 $S$ 为概形，$f:X\to Y$ 为 $S$ 上形式代数空间的态射。
下列条件等价：
\begin{enumerate}
\item $f$ 是满态射；
\item 对每个概形 $T$ 以及态射 $T\to Y$，投影
$X\times_YT\to T$ 是形式代数空间的满态射；
\item 对每个仿射概形 $T$ 以及态射 $T\to Y$，投影
$X\times_YT\to T$ 是形式代数空间的满态射；
\item 存在定义 \ref{formal-spaces-definition-formal-algebraic-space} 意义下的覆盖
$\{Y_j\to Y\}$，使得每个 $X\times_YY_j\to Y_j$ 都是形式代数空间的满态射；
\item 存在形式代数空间的满态射 $Z\to Y$，使得
$X\times_YZ\to Z$ 是满态射，并且
% P10-E0357
\item 此处还可补充更多条件。
\end{enumerate}
\end{lemma}

\begin{proof}
省略。
\end{proof}








\section{单态射}
\label{formal-spaces-section-monomorphisms}

\noindent
定义如下。

\begin{definition}
\label{formal-spaces-definition-monomorphism}
设 $S$ 为概形。若 $S$ 上形式代数空间的态射是层之间的单射映射，
则称其为\textit{单态射}。
\end{definition}

\noindent
下面是一个例子。设 $X$ 为代数空间，且 $T\subset |X|$ 为闭子集。
那么，从沿 $T$ 的 $X$ 的形式完备化到 $X$ 的态射
$X_{/T}\to X$ 是单态射。特别地，形式代数空间的单态射一般并不
可由代数空间表示。

\begin{lemma}
\label{formal-spaces-lemma-composition-monomorphism}
两个单态射的复合是单态射。
\end{lemma}

\begin{proof}
省略。
\end{proof}

\begin{lemma}
\label{formal-spaces-lemma-base-change-monomorphism}
单态射的基变换是单态射。
\end{lemma}

\begin{proof}
省略。
\end{proof}

\begin{lemma}
\label{formal-spaces-lemma-characterize-monomorphisms}
设 $S$ 为概形，$f:X\to Y$ 为 $S$ 上形式代数空间的态射。
下列条件等价：
\begin{enumerate}
\item $f$ 是单态射；
\item 对每个概形 $T$ 以及态射 $T\to Y$，投影
$X\times_YT\to T$ 是形式代数空间的单态射；
\item 对每个仿射概形 $T$ 以及态射 $T\to Y$，投影
$X\times_YT\to T$ 是形式代数空间的单态射；
\item 存在定义 \ref{formal-spaces-definition-formal-algebraic-space} 意义下的覆盖
$\{Y_j\to Y\}$，使得每个 $X\times_YY_j\to Y_j$ 都是形式代数空间的单态射，
并且
\item 存在一族态射 $\{Y_j\to Y\}$，使得
$\coprod Y_j\to Y$ 是 $(\Sch/S)_{fppf}$ 上层的满射，且每个
$X\times_YY_j\to Y_j$ 对所有 $j$ 都是单态射；
\item 存在形式代数空间的态射 $Z\to Y$，该态射可由代数空间表示、满、平坦且局部有限呈示，
使得
% P10-E0359
$X\times_YZ\to X$ 是单态射，并且
% P10-E0358
\item 此处还可补充更多条件。
\end{enumerate}
\end{lemma}

\begin{proof}
省略。
\end{proof}







\section{闭浸入}
\label{formal-spaces-section-closed-immersions}

\noindent
定义如下。

\begin{definition}
\label{formal-spaces-definition-closed-immersion}
设 $S$ 为概形，设 $f:Y\to X$ 为 $S$ 上形式代数空间的态射。
若 $f$ 可由代数空间表示，并且在《自举》之定义
\ref{bootstrap-definition-property-transformation} 的意义下是闭浸入，
则称 $f$ 为\textit{闭浸入}。
\end{definition}

\noindent
% P10-E0360
阅读本节时，请跳过开头那些必需的引理。

\begin{lemma}
\label{formal-spaces-lemma-composition-closed-immersion}
两个闭浸入的复合是闭浸入。
\end{lemma}

\begin{proof}
省略。
\end{proof}

\begin{lemma}
\label{formal-spaces-lemma-base-change-closed-immersion}
闭浸入的基变换是闭浸入。
\end{lemma}

\begin{proof}
省略。
\end{proof}

\begin{lemma}
\label{formal-spaces-lemma-characterize-closed-immersion}
设 $S$ 为概形，$f:X\to Y$ 为 $S$ 上形式代数空间的态射。
下列条件等价：
\begin{enumerate}
\item $f$ 是闭浸入；
\item 对每个概形 $T$ 以及态射 $T\to Y$，投影
$X\times_YT\to T$ 是闭浸入；
\item 对每个仿射概形 $T$ 以及态射 $T\to Y$，投影
$X\times_YT\to T$ 是闭浸入；
\item 存在定义 \ref{formal-spaces-definition-formal-algebraic-space} 意义下的覆盖
$\{Y_j\to Y\}$，使得每个 $X\times_YY_j\to Y_j$ 都是闭浸入，并且
\item 存在形式代数空间的态射 $Z\to Y$，该态射可由代数空间表示、满、平坦且局部有限呈示，
使得
% P10-E0361
$X\times_YZ\to X$ 是闭浸入，并且
% P10-E0362
\item 此处还可补充更多条件。
\end{enumerate}
\end{lemma}

\begin{proof}
省略。
\end{proof}

\begin{lemma}
\label{formal-spaces-lemma-closed-immersion-into-McQuillan}
设 $S$ 为概形，$X$ 为 $S$ 上的 McQuillan 仿射形式代数空间。
设 $f:Y\to X$ 为 $S$ 上形式代数空间的闭浸入。则 $Y$ 是 McQuillan
仿射形式代数空间，并且 $f$ 对应于弱可容许拓扑 $S$-代数的连续同态
$A\to B$；该同态是 taut 的，具有闭核，并且像稠密。
\end{lemma}

\begin{proof}
写作 $X=\text{Spf}(A)$，其中 $A$ 是弱可容许拓扑环。
令 $I_\lambda$ 为 $A$ 中弱可容许定义理想的基本系统。则
$Y\times_X\Spec(A/I_\lambda)$ 是 $\Spec(A/I_\lambda)$ 的闭子概形，
因而是仿射的（定义 \ref{formal-spaces-definition-closed-immersion}）。设
$Y\times_X\Spec(A/I_\lambda)=\Spec(B_\lambda)$。
$A/I_\lambda\to B_\lambda$ 为满射。因此投影
$$
B = \lim B_\lambda \longrightarrow B_\lambda
$$
都是满射，因为复合 $A\to B\to B_\lambda$ 都是满射。
由引理 \ref{formal-spaces-lemma-mcquillan-affine-formal-algebraic-space}，可知 $Y$ 是
McQuillan 的。由引理 \ref{formal-spaces-lemma-representable-affine}，环映射
$A\to B$ 是 taut 的。由于 $B$ 完备且 $A\to B$ 连续，核是闭的。
最后，由于对所有 $\lambda$，$A\to B_\lambda$ 都是满射，故 $A$ 在
$B$ 中的像稠密。
\end{proof}

\noindent
尽管我们已经得到上述结果，但一般而言，当靶为 McQuillan 仿射形式代数空间时，
我们并不知道闭浸入的行为；见备注 \ref{formal-spaces-remark-questions}。

\begin{example}
\label{formal-spaces-example-closed-immersion-from-quotient}
设 $S$ 为概形，设 $A$ 为 $S$ 上的弱可容许拓扑环。
设 $K\subset A$ 为闭理想。令
$$
B = (A/K)^\wedge = \lim_{I \subset A\ w.i.d.} A/(I + K)
$$
则态射 $\text{Spf}(B)\to\text{Spf}(A)$ 可由代数空间表示；见例
\ref{formal-spaces-example-representable-morphism-from-completion}。
若 $T\to\text{Spf}(A)$ 是态射且 $T$ 是拟紧概形，则由引理
\ref{formal-spaces-lemma-factor-through-thickening}，该态射经由某个弱定义理想
$I\subset A$ 因子分解到 $\Spec(A/I)$。于是
$T\times_{\text{Spf}(A)}\text{Spf}(B)$
等于 $T\times_{\Spec(A/I)}\Spec(A/(K + I))$，
由此可见 $\text{Spf}(B)\to\text{Spf}(A)$ 是闭浸入。
由于 $K$ 闭，$A\to B$ 的核是 $K$；但要注意，一般而言环映射
$A\to B=(A/K)^\wedge$ 未必是满射。
\end{example}

\begin{lemma}
\label{formal-spaces-lemma-monomorphism-iso-over-red}
设 $S$ 为概形，$f:X\to Y$ 为形式代数空间的态射。假设
\begin{enumerate}
\item $f$ 可由代数空间表示；
\item $f$ 是单态射；
\item 包含映射 $Y_{red}\to Y$ 经由 $f$ 因子分解；
\item $f$ 局部有限型，或 $Y$ 局部 Noether。
\end{enumerate}
则 $f$ 是闭浸入。
\end{lemma}

\begin{proof}
假设 (2) 和 (3)，则
$X_{red}=X\times_YY_{red}=Y_{red}$。
以下不再特别说明地使用这一点。

\medskip\noindent
若 $Y'\to Y$ 是 $S$ 上形式代数空间的 'etale 态射，
则基变换 $f':X\times_YY'\to Y'$ 满足条件 (1)--(4)。
因此由引理 \ref{formal-spaces-lemma-characterize-closed-immersion}，可以假设 $Y$ 是
仿射形式代数空间。

\medskip\noindent
按定义 \ref{formal-spaces-definition-affine-formal-algebraic-space}，写作
$Y=\colim_{\lambda\in\Lambda}Y_\lambda$。于是
$X_\lambda=X\times_YY_\lambda$ 是带有单态射
$f_\lambda:X_\lambda\to Y_\lambda$ 的代数空间，并且该态射诱导同构
$X_{\lambda, red}\to Y_{\lambda, red}$。
因此由《空间的极限》之命题
\ref{spaces-limits-proposition-affine}，$X_\lambda$ 是仿射概形
（因为 $X_{\lambda, red}\to X_\lambda$ 是满射且整）。
要完成证明，只需说明 $X_\lambda\to Y_\lambda$ 是闭浸入，
这将在下一段中完成。

\medskip\noindent
设 $X\to Y$ 是仿射概形的单态射，并且
$X_{red}=X\times_YY_{red}=Y_{red}$。一般而言，这并不意味着
$X\to Y$ 是闭浸入；见《例》之第
\ref{examples-section-epimorphism-not-surjective} 节。
然而，在我们的条件 (4) 下，我们知道上一段中
% P10-E0363
$X_\lambda\to Y_\lambda$ 要么是有限型，要么 $Y_\lambda$ 是 Noether。
这意味着 $X\to Y$ 对应于环映射 $R\to A$，其中
$I\subset R$ 是幂零根（即 $R$ 的最大局部幂零理想），且
$R/I\to A/IA$ 是同构；并且 $R\to A$ 是有限型，或 $R$ 是 Noether。
在第一种情形，由《代数》之引理
\ref{algebra-lemma-surjective-mod-locally-nilpotent}，$R\to A$ 是满射；
在第二种情形，$I$ 有限生成，因而是幂零的，所以由 Nakayama 引理，
$R\to A$ 是满射；见《代数》之引理 \ref{algebra-lemma-NAK} 第 (11) 部分。
\end{proof}








\section{受限幂级数}
\label{formal-spaces-section-restricted-power-series}

\noindent
设 $A$ 为关于线性拓扑完备的拓扑环（《代数更多内容》之定义
\ref{more-algebra-definition-topological-ring}）。设 $I_\lambda$ 为开理想的基本系统。
设 $r\geq 0$ 为整数。在此情形下，通常记
$$
A\{x_1, \ldots, x_r\} =
\lim_\lambda A/I_\lambda[x_1, \ldots, x_r] =
\lim_\lambda (A[x_1, \ldots, x_r]/I_\lambda A[x_1, \ldots, x_r])
$$
并赋予极限拓扑。换言之，这是关于理想 $I_\lambda$ 的多项式环完备化。
我们可以把 $A\{x_1, \ldots, x_r\}$ 的元素看作关于
$x_1, \ldots, x_r$ 的幂级数
$$
f = \sum\nolimits_{E = (e_1, \ldots, e_r)} a_E x_1^{e_1} \ldots x_r^{e_r}
$$
其中系数 $a_E\in A$ 在 $A$ 的拓扑中趋于零。换言之，对任意 $\lambda$，
除有限多个 $a_E$ 外，其余都属于 $I_\lambda$。
因此，$A\{x_1, \ldots, x_r\}$ 的元素有时称为\textit{受限幂级数}。
有时将此环记为 $A\langle x_1, \ldots, x_r\rangle$；我们将避免使用这一记号。

\begin{remark}[受限幂级数的泛性质]
\label{formal-spaces-remark-universal-property}
\begin{reference}
\cite[Chapter 0, 7.5.3]{EGA}
\end{reference}
设 $A\to C$ 为完备线性拓扑环之间的连续映射。
则任意 $A$-代数映射 $A[x_1, \ldots x_r]\to C$ 都唯一延拓为
受限幂级数上的连续映射 $A\{x_1, \ldots, x_r\}\to C$。
\end{remark}

\begin{remark}
\label{formal-spaces-remark-I-adic-completion-and-restricted-power-series}
设 $A$ 为环，$I\subset A$ 为理想。若 $A$ 是 $I$-进完备的，
则 $A[x_1, \ldots, x_r]$ 的 $I$-进完备化 $A[x_1, \ldots, x_r]^\wedge$
作为环是 $A$ 上的受限幂级数环。然而，目前尚不清楚
$A[x_1, \ldots, x_r]^\wedge$ 是否 $I$-进完备。
我们把 $A\{x_1, \ldots, x_r\}$ 上的拓扑看作极限拓扑（它总是完备的），
而常常把 $A[x_1, \ldots, x_r]^\wedge$ 上的拓扑看作 $I$-进拓扑
（不总是完备的）。若 $I$ 有限生成，则
$A\{x_1, \ldots, x_r\}=A[x_1, \ldots, x_r]^\wedge$ 作为拓扑环成立；
见《代数》之引理 \ref{algebra-lemma-hathat-finitely-generated}。
\end{remark}






\section{拓扑有限型代数}
\label{formal-spaces-section-tft}

\noindent
这是我们的定义。这个定义并未得到普遍认可。许多作者会附加进一步的条件，
通常是因为他们只关注特定类型的环，而不是最一般的情形。

\begin{definition}
\label{formal-spaces-definition-topologically-finite-type}
设 $A\to B$ 为拓扑环之间的连续映射（《代数更多内容》之定义
\ref{more-algebra-definition-topological-ring}）。若存在像在 $B$ 中稠密的
$A$-代数映射 $A[x_1,\ldots,x_n]\to B$，则称 $B$
\textit{在 $A$ 上拓扑有限型}。
\end{definition}

\noindent
若 $A$ 是完备线性拓扑环，则受限幂级数环
$A\{x_1,\ldots,x_r\}$ 在 $A$ 上拓扑有限型。若 $k$ 是域，则幂级数环
$k[[x_1,\ldots,x_r]]$ 在 $k$ 上拓扑有限型。

\medskip\noindent
对于弱可容许拓扑环之间的连续 taut 映射，拓扑有限型这一性质，恰好对应于
相应仿射形式代数空间之间的有限型态射。

\begin{lemma}
\label{formal-spaces-lemma-topologically-finite-type-finite-type}
设 $S$ 为概形，$\varphi:A\to B$ 为 $S$ 上弱可容许拓扑环之间的连续映射。
下列条件等价：
\begin{enumerate}
\item $\text{Spf}(\varphi):Y=\text{Spf}(B)\to\text{Spf}(A)=X$
是有限型；
\item $\varphi$ 是 taut 的，且 $B$ 在 $A$ 上拓扑有限型。
\end{enumerate}
\end{lemma}

\begin{proof}
可以使用引理 \ref{formal-spaces-lemma-representable-affine} 将 $\varphi$ 的 taut 性与
$\text{Spf}(\varphi)$ 的可表示性联系起来。以下不再特别说明地使用这一点。
于是 $X=\colim\Spec(A/I)$，而
$Y=\colim\Spec(B/J(I))$，其中 $I\subset A$ 遍历 $A$ 的弱定义理想，
$J(I)$ 是 $IB$ 在 $B$ 中的闭包。

\medskip\noindent
假设 (2)。选取像在 $B$ 中稠密的环映射
$A[x_1,\ldots,x_r]\to B$。于是
$A[x_1,\ldots,x_r]\to B\to B/J(I)$ 的像也稠密，
这意味着它是满射。因此 $B/J(I)$ 是 $A/I$ 上的有限型代数。
设 $T\to X$ 是态射且 $T$ 是拟紧概形。则由引理
\ref{formal-spaces-lemma-factor-through-thickening}，对某个 $I$，$T\to X$ 经由
$\Spec(A/I)$ 因子分解。于是
$T\times_XY=T\times_{\Spec(A/I)}\Spec(B/J(I))$；见引理
\ref{formal-spaces-lemma-representable-affine} 的证明。
% P10-E0364
因此 $T\times_YX\to T$ 是有限型，因为它是有限型态射
$\Spec(B/J(I))\to\Spec(A/I)$ 的基变换。所以 (1) 成立。

\medskip\noindent
假设 (1)。任取如上的 $I\subset A$。由于
$\Spec(A/I)\times_XY=\Spec(B/J(I))$，可知
$A/I\to B/J(I)$ 是有限型的。选取 $b_1,\ldots,b_r\in B$，
它们在 $A/I$ 上映到 $B/J(I)$ 的生成元。我们断言，将 $x_i$ 映到 $b_i$ 的
环映射 $A[x_1,\ldots,x_r]\to B$ 的像是稠密的。为证明这一点，
令 $I'\subset I$ 为另一个弱定义理想。则
$$
B/(J(I') + IB) = B/J(I)
$$
因为 $J(I)$ 是 $IB$ 的闭包，而 $J(I')$ 是开集。于是可以应用《代数》之引理
\ref{algebra-lemma-surjective-mod-locally-nilpotent}，得到
$(A/I')[x_1,\ldots,x_r]\to B/J(I')$ 是满射。
故 (2) 成立，证明完成。
\end{proof}

\noindent
设 $A$ 为关于线性拓扑完备的拓扑环，设 $(I_\lambda)$ 为开理想的基本系统。
令 $\mathcal{C}$ 为逆系统 $(B_\lambda)$ 所成的范畴，其中
\begin{enumerate}
\item $B_\lambda$ 是有限型 $A/I_\lambda$-代数，并且
\item $B_\mu\to B_\lambda$ 是 $A/I_\mu$-代数同态，且诱导同构
$B_\mu/I_\lambda B_\mu\to B_\lambda$。
\end{enumerate}
$\mathcal{C}$ 中的态射由相容的同态系统给出。

\begin{lemma}
\label{formal-spaces-lemma-category-affine-over}
设 $S$ 为概形，$X$ 为 $S$ 上的仿射形式代数空间。假设 $X$ 是 McQuillan，
并令 $A$ 为与 $X$ 关联的弱可容许拓扑环。则在以下两个范畴之间存在反等价：
\begin{enumerate}
\item 上面引入的范畴 $\mathcal{C}$；
\item $X$ 上有限型仿射形式代数空间的态射 $Y\to X$ 所成的范畴。
\end{enumerate}
\end{lemma}

\begin{proof}
令 $(I_\lambda)$ 为 $A$ 中弱可容许定义理想的基本系统。考虑 (2) 中的 $Y$。
则 $Y\times_X\Spec(A/I_\lambda)$ 是仿射的（定义
\ref{formal-spaces-definition-finite-type} 及引理 \ref{formal-spaces-lemma-affine-representable-by-algebraic-spaces}）。
设 $Y\times_X\Spec(A/I_\lambda)=\Spec(B_\lambda)$。
由于 $\Spec(B_\lambda)\to\Spec(A/I_\lambda)$ 是有限型（由定义
\ref{formal-spaces-definition-finite-type}），环映射 $A/I_\lambda\to B_\lambda$ 是有限型的。
于是 $(B_\lambda)$ 是 $\mathcal{C}$ 的对象。

\medskip\noindent
反之，给定 $\mathcal{C}$ 的对象 $(B_\lambda)$，令
$Y=\colim\Spec(B_\lambda)$。这是仿射形式代数空间。我们断言
$$
Y \times_X \Spec(A/I_\lambda) =
\left(\colim_\mu \Spec(B_\mu)\right) \times_X \Spec(A/I_\lambda) =
\Spec(B_\lambda)
$$
为证明这一点，只需在拟紧概形 $U$ 上取值时得到相同结果。
态射
$U\to\left(\colim_\mu\Spec(B_\mu)\right)\times_X\Spec(A/I_\lambda)$
来自某个 $\mu\geq\lambda$ 的态射
$U\to\Spec(B_\mu)\times_{\Spec(A/I_\mu)}\Spec(A/I_\lambda)$
（两次使用引理 \ref{formal-spaces-lemma-factor-through-thickening}）。由于
$\mathcal{C}$ 的对象所满足的第二个条件，
$\Spec(B_\mu)\times_{\Spec(A/I_\mu)}\Spec(A/I_\lambda)=\Spec(B_\lambda)$，
这就证明了所需结论。由此可以说明 $Y\to X$ 是有限型。确切地说，对于任意
$U\to X$ 且 $U$ 为拟紧概形，由上面引用的引理，对某个 $\lambda$ 有分解
$U\to\Spec(A/I_\lambda)\to X$。于是
$$
Y \times_X U =
Y \times_X \Spec(A/I_\lambda)) \times_{\Spec(A/I_\lambda)} U =
\Spec(B_\lambda) \times_{\Spec(A/I_\lambda)} U
$$
% P10-E0365
如所需，这是 $U$ 上的有限型概形。因此构造
$(B_\lambda)\mapsto\colim\Spec(B_\lambda)$ 确实给出从范畴 (1) 到范畴 (2) 的函子。

\medskip\noindent
为完成证明，我们说明上述构造在范畴 (1) 和 (2) 之间定义了准逆函子。
一个方向需要证明，对于范畴 $\mathcal{C}$ 中任意对象 $(B_\lambda)$，有
$$
\left(\colim_\mu \Spec(B_\mu)\right) \times_X \Spec(A/I_\lambda) =
\Spec(B_\lambda)
$$
上面已经证明了这一点。另一个方向需要证明，对范畴 (2) 中给定的 $Y$，有
$$
Y = \colim (Y \times_X \Spec(A/I_\lambda))
$$
这同样可通过在拟紧测试对象上取值，并利用
$X=\colim\Spec(A/I_\lambda)$ 得到。
\end{proof}

\begin{remark}
\label{formal-spaces-remark-questions}
设 $A$ 为弱可容许拓扑环，$(I_\lambda)$ 为弱定义理想的基本系统。
令 $X=\text{Spf}(A)$，换言之，$X$ 是 McQuillan 仿射形式代数空间。
设 $f:Y\to X$ 为仿射形式代数空间的态射。一般而言，$Y$ 未必是 McQuillan。
更具体地，我们可以提出下列问题：
\begin{enumerate}
\item 假设 $f:Y\to X$ 是闭浸入。则 $Y$ 是 McQuillan，且 $f$ 对应于弱可容许拓扑环的连续映射
$\varphi:A\to B$，它是 taut 的，核 $K\subset A$ 是闭理想，并且像
$\varphi(A)$ 在 $B$ 中稠密；见引理 \ref{formal-spaces-lemma-closed-immersion-into-McQuillan}。
什么条件能保证 $B=(A/K)^\wedge$，如例 \ref{formal-spaces-example-closed-immersion-from-quotient} 中那样？
\item 什么条件能保证闭浸入 $f:Y\to X$ 对应于 $A$ 按闭理想作商所得的商环 $A/K$；换言之，
相应的连续映射 $\varphi$ 是满射且开映射？
\item 假设 $f:Y\to X$ 是有限型。由引理 \ref{formal-spaces-lemma-category-affine-over}，有
$Y=\colim\Spec(B_\lambda)$，其中 $(B_\lambda)$ 是 $\mathcal{C}$ 的对象。
此时确实存在固定整数 $r$，使得对所有 $\lambda$，$B_\lambda$ 都由
$A/I_\lambda$ 上的 $r$ 个元素生成（论证实质上已在引理
\ref{formal-spaces-lemma-topologically-finite-type-finite-type} 的 (1) $\Rightarrow$ (2) 证明中给出）。
然而，目前尚不清楚投影 $\lim B_\lambda\to B_\lambda$ 是否为满射，即尚不清楚 $Y$ 是否为 McQuillan。
是否存在 $Y$ 不是 McQuillan 的例子？
\item 假设 $f:Y\to X$ 是有限型且 $Y$ 是 McQuillan。则 $f$ 对应于弱可容许拓扑环的连续映射
$\varphi:A\to B$。事实上，$\varphi$ 是 taut 的，且 $B$ 在 $A$ 上拓扑有限型；见引理
\ref{formal-spaces-lemma-topologically-finite-type-finite-type}。换言之，$f$ 可分解为
$$
Y \longrightarrow \mathbf{A}^r_X \longrightarrow X
$$
其中第一个箭头是 McQuillan 仿射形式代数空间的闭浸入。然而，此时对
$Y\to\mathbf{A}^r_X$，问题 (1) 和 (2) 仍然有效。
\end{enumerate}
下面我们将在 $X$ 可数指标化时回答这些问题，即 $A$ 具有可数的开理想基本系统时。
如果你在更一般的情形下有这些问题的答案，或者有反例，请发送电子邮件至
\href{mailto:stacks.project@gmail.com}{stacks.project@gmail.com}。
\end{remark}

\begin{lemma}
\label{formal-spaces-lemma-closed-immersion-into-countably-indexed}
设 $S$ 为概形，$X$ 为 $S$ 上可数指标化的仿射形式代数空间。
设 $f:Y\to X$ 为 $S$ 上形式代数空间的闭浸入。则 $Y$ 是可数指标化的仿射形式代数空间，
且 $f$ 对应于 $A\to A/K$，其中 $A$ 是
$\textit{WAdm}^{count}$ 中的对象（第 \ref{formal-spaces-section-morphisms-rings} 节），
而 $K\subset A$ 是闭理想。
\end{lemma}

\begin{proof}
由引理 \ref{formal-spaces-lemma-countably-indexed}，$X=\text{Spf}(A)$，其中 $A$ 是
$\textit{WAdm}^{count}$ 中的对象。由于闭浸入是可表示且仿射的，
由引理 \ref{formal-spaces-lemma-property-goes-up-affine-morphism} 可知 $Y$ 是仿射形式代数空间且可数
% P10-E0366
指标化。因此再次应用引理 \ref{formal-spaces-lemma-countably-indexed}，得到
$Y=\text{Spf}(B)$，其中 $B$ 是 $\textit{WAdm}^{count}$ 中的对象。
由引理 \ref{formal-spaces-lemma-closed-immersion-into-McQuillan}，可知 $f$ 由
$\textit{WAdm}^{count}$ 中一个 taut 且像稠密的态射 $A\to B$ 给出。
最后应用引理 \ref{formal-spaces-lemma-dense-image-surjective} 即可完成证明。
\end{proof}

\begin{lemma}
\label{formal-spaces-lemma-quotient-restricted-power-series}
设 $B\to A$ 为 $\textit{WAdm}^{count}$ 中的箭头，见第
\ref{formal-spaces-section-morphisms-rings} 节。下列条件等价：
\begin{enumerate}
\item[(a)] $B\to A$ 是 taut 的，并且对每个弱定义理想 $J\subset B$，
其中 $I\subset A$ 是 $JA$ 的闭包，$B/J\to A/I$ 是有限型；
\item[(b)] $B\to A$ 是 taut 的，并且对 $B$ 的弱定义理想的一个共尾系统
$(J_\lambda)$，$B/J_\lambda\to A/I_\lambda$ 是有限型，其中
$I_\lambda\subset A$ 是 $J_\lambda A$ 的闭包；
\item[(c)] $B\to A$ 是 taut 的，并且 $A$ 在 $B$ 上拓扑有限型；
\item[(d)] 作为拓扑 $B$-代数，$A$ 同构于
$B\{x_1, \ldots, x_n\}$ 对某个闭理想取商所得的代数。
\end{enumerate}
此外，这些等价条件定义了一个局部性质，即它们满足公理
\ref{formal-spaces-item-axiom-1}、\ref{formal-spaces-item-axiom-2}、\ref{formal-spaces-item-axiom-3}。
\end{lemma}

\begin{proof}
由定义，蕴含 (a) $\Rightarrow$ (b)、(c) $\Rightarrow$ (a) 以及
(d) $\Rightarrow$ (c) 都是直接的。假设 (b) 成立，令 $J\subset B$ 和
$I\subset A$ 如 (a) 中所述。取交换图表
$$
\xymatrix{
A \ar[r] & \ldots \ar[r] & A_3 \ar[r] & A_2 \ar[r] & A_1 \\
B \ar[r] \ar[u] & \ldots \ar[r] & B/J_3 \ar[r] \ar[u] &
B/J_2 \ar[r] \ar[u] & B/J_1 \ar[u]
}
$$
使 $A_{n + 1}/J_nA_{n + 1}=A_n$，并且如引理
\ref{formal-spaces-lemma-representable-property-rings} 中那样 $A=\lim A_n$。
对每个 $m$，存在 $\lambda$ 使 $J_\lambda\subset J_m$。
由于 $B/J_\lambda\to A/I_\lambda$ 是有限型，这说明
$B/J_m\to A/I_m$ 是有限型。令 $\alpha_1, \ldots, \alpha_n\in A_1$
为 $A_1$ 在 $B/J_1$ 上的生成元。由于 $A$ 是一个具有满过渡映射的系统的
可数极限，可以找到 $a_1, \ldots, a_n\in A$，它们在 $A_1$ 中分别映到
$\alpha_1, \ldots, \alpha_n$。由备注 \ref{formal-spaces-remark-universal-property}，
存在连续映射 $B\{x_1, \ldots, x_n\}\to A$，将 $x_i$ 映到 $a_i$。
由《代数》之引理 \ref{algebra-lemma-surjective-mod-locally-nilpotent}，
此映射诱导满射 $(B/J_m)[x_1, \ldots, x_n]\to A_m$。
对 $m\geq 1$，得到短正合序列
$$
0 \to K_m \to (B/J_m)[x_1, \ldots, x_n] \to A_m \to 0
$$
诱导的过渡映射 $K_{m + 1}\to K_m$ 是满射，因为
$A_{m + 1}/J_mA_{m + 1}=A_m$。因此这些短正合序列的逆极限是正合的；见《代数》之引理
\ref{algebra-lemma-ML-exact-sequence}。由于
$B\{x_1, \ldots, x_n\}=\lim(B/J_m)[x_1, \ldots, x_n]$
且 $A=\lim A_m$，可知 $B\{x_1, \ldots, x_n\}\to A$ 是满射且开映射。
由于 $A$ 完备，核是闭理想。由此得到 (a)、(b)、(c)、(d) 等价。

\medskip\noindent
给定图表 \ref{formal-spaces-equation-localize}，其情形如
\ref{formal-spaces-situation-local-property} 所述。由例 \ref{formal-spaces-example-finite-type-from-finite-type-ring-map}，
$A\to(A')^\wedge$ 和 $B\to(B')^\wedge$ 满足 (a)、(b)、(c)、(d)。此外，
由引理 \ref{formal-spaces-lemma-representable-property-rings}，为证明公理
\ref{formal-spaces-item-axiom-1} 和 \ref{formal-spaces-item-axiom-2}，可以假设 $B\to A$ 与
$(B')^\wedge\to(A')^\wedge$ 都是 taut 的。现在取 $B$ 中的弱定义理想
$J$。令 $J'\subset(B')^\wedge$、$I\subset A$、$I'\subset(A')^\wedge$
分别为 $J(B')^\wedge$、$JA$、$J(A')^\wedge$ 的闭包。由上面的结论，
只需考虑交换图表
$$
\xymatrix{
A/I \ar[r] & (A')^\wedge/I' \\
B/J \ar[r] \ar[u]^{\overline{\varphi}} &
(B')^\wedge/J' \ar[u]_{\overline{\varphi}'}
}
$$
并证明：(1) $\overline{\varphi}$ 有限型蕴含 $\overline{\varphi}'$
有限型；(2) 若 $A\to A'$ 忠实平坦，则 $\overline{\varphi}'$ 有限型蕴含
$\overline{\varphi}$ 有限型。注意，由 $(B')^\wedge$ 和 $(A')^\wedge$ 的拓扑构造，
$(B')^\wedge/J'=B'/JB'$ 且 $(A')^\wedge/I'=A'/IA'$。
特别地，图表中的水平映射是 \'etale 的。(1) 现在由《代数》之引理
\ref{algebra-lemma-compose-finite-type} 得出；(2) 则由《下降》之引理
\ref{descent-lemma-finite-type-local-source-fppf-algebra} 得出，因为环映射
$A/I\to(A')^\wedge/I'=A'/IA'$ 是忠实平坦且 \'etale 的。

\medskip\noindent
公理 \ref{formal-spaces-item-axiom-3} 的证明略去。
\end{proof}

\begin{lemma}
\label{formal-spaces-lemma-quotient-restricted-power-series-admissible}
在引理 \ref{formal-spaces-lemma-quotient-restricted-power-series} 中，若 $B$ 是可容许的
（例如 adic），则等价条件 (a)--(d) 还等价于
\begin{enumerate}
\item[(e)] $B\to A$ 是 taut 的，并且对某个定义理想 $J\subset B$，
其中 $I\subset A$ 是 $JA$ 的闭包，$B/J\to A/I$ 是有限型。
\end{enumerate}
\end{lemma}

\begin{proof}
只需证明 (e) 蕴含 (a)。设 $J'\subset B$ 为弱定义理想，令
$I'\subset A$ 为 $J'A$ 的闭包。我们需要证明 $B/J'\to A/I'$ 是有限型。
若相应结论对较小的弱定义理想 $J''=J'\cap J$ 成立，则对 $J'$ 也成立。
因此可假设 $J'\subset J$。由于 $J$ 是定义理想（而不仅是弱定义理想），
存在 $n\geq 1$ 使 $J^n\subset J'$。于是可以考虑图表
$$
\xymatrix{
0 \ar[r] & I/I' \ar[r] & A/I' \ar[r] & A/I \ar[r] & 0 \\
0 \ar[r] & J/J' \ar[r] \ar[u] & B/J' \ar[r] \ar[u] & B/J \ar[r] \ar[u] & 0
}
$$
其两行都是正合的。由于 $I'\subset A$ 是开集，并且 $I$ 是 $JA$ 的闭包，
可知 $I/I'=(J/J')\cdot A/I'$。由于 $J/J'$ 是幂零理想，且
$B/J\to A/I$ 是有限型，由《代数》之引理
\ref{algebra-lemma-finite-type-mod-nilpotent}，可知 $A/I'$ 在 $B/J'$ 上有限型，
这正是所需结论。
\end{proof}

\begin{lemma}
\label{formal-spaces-lemma-representable-affine-finite-type}
设 $S$ 为概形，$f:X\to Y$ 为仿射形式代数空间的态射。假设 $Y$ 可数指标化。
下列条件等价：
\begin{enumerate}
\item $f$ 局部有限型；
\item $f$ 有限型；
\item $f$ 对应于 $\textit{WAdm}^{count}$ 中的态射 $B\to A$
（第 \ref{formal-spaces-section-morphisms-rings} 节），并满足引理
\ref{formal-spaces-lemma-quotient-restricted-power-series} 的等价条件。
\end{enumerate}
\end{lemma}

\begin{proof}
由于 $X$ 和 $Y$ 是仿射的，条件 (1) 和 (2) 等价是显然的。在 (1) 或 (2) 的情形，
按定义 $f$ 可由代数空间表示，因而由引理
\ref{formal-spaces-lemma-affine-representable-by-algebraic-spaces} 可知它是仿射的。
因此若 (1) 或 (2) 成立，由引理 \ref{formal-spaces-lemma-property-goes-up-affine-morphism}，
可知 $X$ 是可数指标化的。按引理 \ref{formal-spaces-lemma-countably-indexed}，写作
$X=\text{Spf}(A)$、$Y=\text{Spf}(B)$，其中 $A$ 和 $B$ 是
$\textit{WAdm}^{count}$ 中的拓扑 $S$-代数。由引理
\ref{formal-spaces-lemma-morphism-between-formal-spectra}，可知 $f$ 对应于连续映射
$B\to A$。于是结论由引理
\ref{formal-spaces-lemma-topologically-finite-type-finite-type} 得出。
\end{proof}

\begin{lemma}
\label{formal-spaces-lemma-finite-type-local-property}
设 $S$ 为概形，$f:X\to Y$ 为 $S$ 上局部可数指标化形式代数空间的态射。
下列条件等价：
\begin{enumerate}
\item 对每个交换图表
$$
\xymatrix{
U \ar[d] \ar[r] & V \ar[d] \\
X \ar[r] & Y
}
$$
若 $U,V$ 是仿射形式代数空间，$U\to X$ 和 $V\to Y$ 可由代数空间表示且为
\'etale，则态射 $U\to V$ 对应于 $\textit{WAdm}^{count}$ 中的态射，
该态射是 taut 且拓扑有限型；
\item 存在定义 \ref{formal-spaces-definition-formal-algebraic-space} 意义下的覆盖
$\{Y_j\to Y\}$，且对每个 $j$，存在定义
\ref{formal-spaces-definition-formal-algebraic-space} 意义下的覆盖
$\{X_{ji}\to Y_j\times_Y X\}$，使每个 $X_{ji}\to Y_j$ 对应于
$\textit{WAdm}^{count}$ 中的态射，该态射是 taut 且拓扑有限型；
\item 存在定义 \ref{formal-spaces-definition-formal-algebraic-space} 意义下的覆盖
$\{X_i\to X\}$，并且对每个 $i$ 存在分解 $X_i\to Y_i\to Y$，其中 $Y_i$
是仿射形式代数空间，$Y_i\to Y$ 可由代数空间表示且为 \'etale，且
$X_i\to Y_i$ 对应于 $\textit{WAdm}^{count}$ 中的态射，
% P10-E0368
该态射是 taut 且拓扑有限型，并且
\item $f$ 局部有限型。
\end{enumerate}
\end{lemma}

\begin{proof}
由引理 \ref{formal-spaces-lemma-quotient-restricted-power-series}，性质
$P(\varphi)=$“$\varphi$ 是 taut 且拓扑有限型”在
$\text{WAdm}^{count}$ 上是局部的。因此由引理
\ref{formal-spaces-lemma-property-defines-property-morphisms}，条件 (1)、(2)、(3) 等价。
另一方面，由引理 \ref{formal-spaces-lemma-representable-affine-finite-type}，
$\textit{WAdm}^{count}$ 中态射的性质 $P$ 恰好对应于可数指标化仿射形式代数空间
态射的局部有限型性质。因此引理 \ref{formal-spaces-lemma-finite-type-local} 的蕴含
(1) $\Rightarrow$ (3) 表明，本引理的 (4) 蕴含 (1)。类似地，
引理 \ref{formal-spaces-lemma-finite-type-local} 的蕴含 (4) $\Rightarrow$ (1) 表明，
本引理的 (2) 蕴含 (4)。
\end{proof}






\section{态射的分离公理}
\label{formal-spaces-section-separation-axioms}

\noindent
本节是《空间的态射》第 \ref{spaces-morphisms-section-separation-axioms} 节
关于形式代数空间态射的对应版本。

\begin{definition}
\label{formal-spaces-definition-separated-morphism}
设 $S$ 为概形。设 $f:X\to Y$ 为 $S$ 上形式代数空间的态射。
设 $\Delta_{X/Y}:X\to X\times_YX$ 为对角态射。
\begin{enumerate}
\item 若 $\Delta_{X/Y}$ 是闭浸入，则称 $f$ 为\textit{分离的}；
\item 若 $\Delta_{X/Y}$ 是拟紧的，则称 $f$ 为\textit{拟分离的}。
\end{enumerate}
\end{definition}

\noindent
由于由引理 \ref{formal-spaces-lemma-diagonal-morphism-formal-algebraic-spaces}，
$\Delta_{X/Y}$ 可表示（由概形表示），所以可以考虑从仿射概形 $T$ 到
$X\times_YX$ 的态射，并检查
$$
E = T \times_{X \times_Y X} X \longrightarrow T
$$
是否拟紧或为闭浸入；见引理 \ref{formal-spaces-lemma-quasi-compact-representable} 或定义
\ref{formal-spaces-definition-closed-immersion}。
注意，概形 $E$ 是两个态射 $a,b:T\to X$ 的等化子，其中它们作为到
$Y$ 的态射相同；并且 $E\to T$ 是单态射且局部有限型。

\begin{lemma}
\label{formal-spaces-lemma-base-change-separated}
定义 \ref{formal-spaces-definition-separated-morphism} 中列出的所有分离公理在基变换下稳定。
\end{lemma}

\begin{proof}
设 $f:X\to Y$ 和 $Y'\to Y$ 为形式代数空间的态射。设
$f':X'\to Y'$ 是 $f$ 沿 $Y'\to Y$ 的基变换。则
$\Delta_{X'/Y'}$ 是 $\Delta_{X/Y}$ 沿态射
$X'\times_{Y'}X'\to X\times_YX$ 的基变换。
定义 \ref{formal-spaces-definition-separated-morphism} 中使用的对角态射的每个性质
都在基变换下稳定。因此引理成立。
\end{proof}

\begin{lemma}
\label{formal-spaces-lemma-fibre-product-after-map}
设 $S$ 为概形，$f:X\to Z$、$g:Y\to Z$ 以及 $Z\to T$ 为 $S$ 上形式代数空间的态射。
考虑诱导态射 $i:X\times_ZY\to X\times_TY$。则
\begin{enumerate}
\item $i$ 可表示（由概形表示），局部有限型、局部拟有限、分离且为单态射；
\item 若 $Z\to T$ 分离，则 $i$ 是闭浸入，并且
\item 若 $Z\to T$ 拟分离，则 $i$ 拟紧。
\end{enumerate}
\end{lemma}

\begin{proof}
由一般范畴论，下面的图表
% P10-E0369
$$
\xymatrix{
X \times_Z Y \ar[r]_i \ar[d] & X \times_T Y \ar[d] \\
Z \ar[r]^-{\Delta_{Z/T}} \ar[r] & Z \times_T Z
}
$$
是纤维积图表。因此 $i$ 是对角态射 $\Delta_{Z/T}$ 的基变换。
于是引理由引理 \ref{formal-spaces-lemma-diagonal-morphism-formal-algebraic-spaces} 得出。
\end{proof}

\begin{lemma}
\label{formal-spaces-lemma-composition-separated}
定义 \ref{formal-spaces-definition-separated-morphism} 中列出的所有分离公理在态射复合下稳定。
\end{lemma}

\begin{proof}
设 $f:X\to Y$ 和 $g:Y\to Z$ 为适用所讨论公理的形式代数空间态射。
对角态射 $\Delta_{X/Z}$ 是复合
$$
X \longrightarrow X \times_Y X \longrightarrow X \times_Z X
$$
我们的分离公理通过要求对角态射具有某个性质 $\mathcal{P}$ 来定义。
由上面的引理 \ref{formal-spaces-lemma-fibre-product-after-map}，第二个箭头也具有这一性质。
因此，由于具有性质 $\mathcal{P}$ 的（可表示）态射的复合仍是具有性质
$\mathcal{P}$ 的态射，引理得证。
\end{proof}

\begin{lemma}
\label{formal-spaces-lemma-separated-local}
设 $S$ 为概形，$f:X\to Y$ 为 $S$ 上形式代数空间的态射。令 $\mathcal{P}$ 为定义
\ref{formal-spaces-definition-separated-morphism} 中的任一分离公理。下列条件等价：
\begin{enumerate}
\item $f$ 具有性质 $\mathcal{P}$；
\item 对每个概形 $Z$ 以及态射 $Z\to Y$，$f$ 的基变换
$Z\times_YX\to Z$ 具有性质 $\mathcal{P}$；
\item 对每个仿射概形 $Z$ 以及态射 $Z\to Y$，$f$ 的基变换
$Z\times_YX\to Z$ 具有性质 $\mathcal{P}$；
\item 对每个仿射概形 $Z$ 以及态射 $Z\to Y$，形式代数空间
$Z\times_YX$ 具有性质 $\mathcal{P}$（见定义 \ref{formal-spaces-definition-separated}）；
\item 存在定义 \ref{formal-spaces-definition-formal-algebraic-space} 意义下的覆盖
$\{Y_j\to Y\}$，使得对所有 $j$，基变换 $Y_j\times_YX\to Y_j$ 具有性质 $\mathcal{P}$。
\end{enumerate}
\end{lemma}

\begin{proof}
以下将反复使用引理 \ref{formal-spaces-lemma-base-change-separated}，不再逐一说明。
特别地，(1) 蕴含 (2)，(2) 蕴含 (3) 是显然的。

\medskip\noindent
假设 (3)，设 $Z\to Y$ 为态射，其中 $Z$ 是仿射概形。令 $U,V$ 为仿射概形，
并设 $a:U\to Z\times_YX$ 和 $b:V\to Z\times_YX$ 为态射。则
$$
U \times_{Z \times_Y X} V =
(Z \times_Y X) \times_{\Delta, (Z \times_Y X) \times_Z (Z \times_Y X)}
(U \times_Z V)
$$
当 $\mathcal{P}=$“拟分离”时，这说明该空间是拟紧的；当
$\mathcal{P}=$“分离”时，这说明它是带有到 $U\times_ZV$ 的闭浸入的仿射概形。
因此 (4) 成立。

\medskip\noindent
假设 (4)，设 $Z\to Y$ 为态射，其中 $Z$ 是仿射概形。令 $U,V$ 为仿射概形，
并设 $a:U\to Z\times_YX$ 和 $b:V\to Z\times_YX$ 为态射。
将上面的论证反向进行，可知 $U\times_{Z\times_YX}V\to U\times_ZV$
在 $\mathcal{P}=$“拟分离”时拟紧，在 $\mathcal{P}=$“分离”时为闭浸入。
由于可以选择上述 $U,V$，使 $U$ 遍历 $Z\times_YX$ 的一个
\'etale 覆盖，我们得到相应态射
$$
U \times_Z V \to (Z \times_Y X) \times_Z (Z \times_Y X)
$$
构成由仿射概形组成的 \'etale 覆盖。因此
$\Delta:(Z\times_YX)\to(Z\times_YX)\times_Z(Z\times_YX)$
拟紧，或为闭浸入（分别对应两种 $\mathcal{P}$）。于是 (3) 成立。

\medskip\noindent
证明 (3) 蕴含 (5)。假设 (3)，令 $\{Y_j\to Y\}$ 如定义
\ref{formal-spaces-definition-formal-algebraic-space} 所述。我们需要证明态射
$$
\Delta_j :
Y_j \times_Y X
\longrightarrow
(Y_j \times_Y X) \times_{Y_j} (Y_j \times_Y X) =
Y_j \times_Y X \times_Y X
$$
具有相应性质（即拟紧或为闭浸入）。按定义
\ref{formal-spaces-definition-affine-formal-algebraic-space}，写作
$Y_j=\colim Y_{j,\lambda}$。在上式中以 $Y_{j,\lambda}$ 替代 $Y_j$，
由假设 (3) 可得所需性质。由于显示的箭头是箭头 $\Delta_{j,\lambda}$ 的余极限，
并且可以通过对映入 $Y_j\times_YX\times_YX$ 的仿射概形作基变换来检验
$\Delta_j$ 是否具有相应性质，故由引理
\ref{formal-spaces-lemma-factor-through-thickening} 得出结论。

\medskip\noindent
证明 (5) 蕴含 (1)。令 $\{Y_j\to Y\}$ 如 (5) 所述。则有纤维积图表
$$
\xymatrix{
\coprod Y_j \times_Y X \ar[r] \ar[d] &
X \ar[d] \\
\coprod Y_j \times_Y X \times_Y X \ar[r] &
X \times_Y X
}
$$
由假设，左侧竖直箭头是拟紧或闭浸入。由《空间》之引理
\ref{spaces-lemma-descent-representable-transformations-property}，
右侧竖直箭头也拟紧或为闭浸入。
\end{proof}






\section{固有态射}
\label{formal-spaces-section-proper}

\noindent
下面是我们将使用的定义。

\begin{definition}
\label{formal-spaces-definition-proper}
设 $S$ 为概形，设 $f:Y\to X$ 为 $S$ 上形式代数空间的态射。
若 $f$ 可由代数空间表示，并且在《自举》之定义
\ref{bootstrap-definition-property-transformation} 的意义下是固有的，
则称 $f$ 为\textit{固有态射}。
\end{definition}

\noindent
由定义可知，固有态射是有限型的。

\begin{lemma}
\label{formal-spaces-lemma-proper-local}
设 $S$ 为概形，$f:X\to Y$ 为 $S$ 上形式代数空间的态射。
下列条件等价：
\begin{enumerate}
\item $f$ 是固有态射；
\item 对每个概形 $Z$ 以及态射 $Z\to Y$，$f$ 的基变换
$Z\times_YX\to Z$ 是固有态射；
\item 对每个仿射概形 $Z$ 以及每个态射 $Z\to Y$，$f$ 的基变换
$Z\times_YX\to Z$ 是固有态射；
\item 对每个仿射概形 $Z$ 以及每个态射 $Z\to Y$，形式代数空间
$Z\times_YX$ 是在 $Z$ 上固有的代数空间；
\item 存在定义 \ref{formal-spaces-definition-formal-algebraic-space} 意义下的覆盖
$\{Y_j\to Y\}$，使得对所有 $j$，基变换 $Y_j\times_YX\to Y_j$ 是固有的。
\end{enumerate}
\end{lemma}

\begin{proof}
省略。
\end{proof}

\begin{lemma}
\label{formal-spaces-lemma-base-change-proper}
形式代数空间的固有态射在基变换下保持固有性。
\end{lemma}

\begin{proof}
这是引理 \ref{formal-spaces-lemma-proper-local} 及基变换的传递性的直接推论。
\end{proof}






\section{形式代数空间与 fpqc 覆盖}
\label{formal-spaces-section-fpqc}

\noindent
本节是《空间的性质》第 \ref{spaces-properties-section-fpqc} 节关于形式代数空间的对应版本。
请先阅读该节。

\begin{lemma}
\label{formal-spaces-lemma-sheaf-fpqc}
\begin{slogan}
形式代数空间是 fpqc 层
\end{slogan}
设 $S$ 为概形，设 $X$ 为 $S$ 上的形式代数空间。则 $X$ 满足 fpqc 拓扑的层性质。
\end{lemma}

\begin{proof}
证明与《空间的性质》之命题
\ref{spaces-properties-proposition-sheaf-fpqc} 的证明\textbf{完全相同}。
由于 $X$ 是 Zariski 拓扑下的层，只需证明如下事实。给定仿射概形之间的满平坦态射
$f:T'\to T$，有：$X(T)$ 是两个映射
$X(T')\to X(T'\times_TT')$ 的等化子。见《拓扑》之引理
\ref{topologies-lemma-sheaf-property-fpqc}。

\medskip\noindent
设 $a,b:T\to X$ 为两个满足 $a\circ f=b\circ f$ 的态射。我们需要证明 $a=b$。
考虑纤维积
$$
E = X \times_{\Delta_{X/S}, X \times_S X, (a, b)} T.
$$
由引理 \ref{formal-spaces-lemma-diagonal-formal-algebraic-space}，对角态射
$\Delta_{X/S}$ 是可表示的单态射。因此 $E\to T$ 是概形的单态射。
由假设 $a\circ f=b\circ f$，$T'\to T$ 唯一地经由 $E$ 因子分解。
考虑交换图表
$$
\xymatrix{
T' \times_T E \ar[r] \ar[d] & E \ar[d] \\
T' \ar[r] \ar@/^5ex/[u] \ar[ru] & T
}
$$
由于投影 $T'\times_TE\to T'$ 是带截面的单态射，可知它是同构。
于是由《下降》之引理
\ref{descent-lemma-descending-property-isomorphism}，$E\to T$ 是同构。
这就得到所需的 $a=b$。

\medskip\noindent
现在设 $c:T'\to X$ 为态射，且两个复合
$T'\times_TT'\to T'\to X$ 相同。我们需要找到态射 $a:T\to X$，
使其与 $T'\to T$ 的复合为 $c$。选取形式仿射概形 $U$ 及一个 \'etale 态射
$U\to X$，使 $|U|\to|X_{red}|$ 的像包含
$|c|:|T'|\to|X_{red}|$ 的像。这由定义
\ref{formal-spaces-definition-formal-algebraic-space}、《空间的性质》之引理
\ref{spaces-properties-lemma-topology-points}、有限个形式仿射代数空间的并仍为形式仿射代数空间
这一事实以及 $|T'|$ 拟紧这一事实得到（略去小论证）。
$U\to X$ 可由概形表示（引理 \ref{formal-spaces-lemma-presentation-representable}）且分离
（引理 \ref{formal-spaces-lemma-separated-from-separated}）。因此
$$
V = U \times_{X, c} T' \longrightarrow T'
$$
是概形的 \'etale 且分离的态射。由我们对 $U\to X$ 的选择，它也是满射
（若不想论证这一点，也可将 $U$ 换成形式仿射代数空间的互不相交并，使
$U\to X$ 满射；其余论证仍然成立）。条件
$c\circ\text{pr}_0=c\circ\text{pr}_1$ 意味着在 $V/T'/T$ 上得到下降数据
（《下降》之定义 \ref{descent-definition-descent-datum}），因为
\begin{align*}
V \times_{T'} (T' \times_T T')
& = U \times_{X, c \circ \text{pr}_0} (T' \times_T T') \\
& = (T' \times_T T') \times_{c \circ \text{pr}_1, X} U \\
& = (T' \times_T T') \times_{T'} V
\end{align*}
态射 $V\to T'$ 是 ind-拟仿射的，由《态射更多内容》之引理
\ref{more-morphisms-lemma-etale-separated-ind-quasi-affine} 得出
（因为 \'etale 态射局部拟有限；见《态射》之引理
\ref{morphisms-lemma-etale-locally-quasi-finite}）。由《群胚更多内容》之引理
\ref{more-groupoids-lemma-ind-quasi-affine}，下降数据是有效的。
设 $W\to T$ 为一个态射，使存在同构
$\alpha:T'\times_TW\to V$，该同构与 $V$ 上给定的下降数据以及
$T'\times_TW$ 上的典型下降数据相容。则 $W\to T$ 是满且 \'etale 的
（《下降》之引理 \ref{descent-lemma-descending-property-surjective} 和
\ref{descent-lemma-descending-property-etale}）。考虑复合
$$
b' : T' \times_T W \longrightarrow V = U \times_{X, c} T' \longrightarrow U
$$
由我们对 $\alpha$ 的选择以及 $c$ 的相应性质，两个复合
$b' \circ (\text{pr}_0, 1),
b' \circ (\text{pr}_1, 1) :
(T' \times_T T') \times_T W \to T' \times_T W \to U$
相同（计算略）。因此由《下降》之引理
\ref{descent-lemma-fpqc-universal-effective-epimorphisms}，$b'$ 下降为态射
$b:W\to U$。图表
$$
\xymatrix{
T' \times_T W \ar[r] \ar[d] & W \ar[r]_b & U \ar[d] \\
T' \ar[rr]^c &  & X
}
$$
交换。这意味着我们已经在 $T$ 的 \'etale 局部找到了 $a$，即得到
$a':W\to X$。然而，由于第一段已经证明唯一性，这个 \'etale 局部解满足粘合条件，
即 $\text{pr}_0^*a'=\text{pr}_1^*a'$ 作为 $X(W\times_TW)$ 的元素成立。
由于 $X$ 是 \'etale 层，得到唯一的 $a\in X(T)$，其在 $W$ 上限制为 $a'$。
\end{proof}





\section{从仿射形式概形出发的映射}
\label{formal-spaces-section-map-out-of}

\noindent
我们证明一些稍后会有用的结果。在论文 \cite{Bhatt-Algebraize} 中，读者可以找到
同类性质的非常一般的结果。

\begin{lemma}
\label{formal-spaces-lemma-map-into-affine}
设 $S$ 为概形，$A$ 为弱可容许拓扑 $S$-代数，$X$ 为 $S$ 上的仿射概形。
则自然映射
$$
\Mor_S(\Spec(A), X)
\longrightarrow
\Mor_S(\text{Spf}(A), X)
$$
是双射。
\end{lemma}

\begin{proof}
若 $X$ 是仿射的，写作 $X=\Spec(B)$，则由引理
\ref{formal-spaces-lemma-morphism-between-formal-spectra}，态射
$\text{Spf}(A)\to\Spec(B)$ 对应于连续的 $S$-代数映射 $B\to A$，
其中 $B$ 具有离散拓扑。这些正是 $S$-代数映射，对应于
$\Spec(A)\to\Spec(B)$ 的态射。
\end{proof}

\begin{lemma}
\label{formal-spaces-lemma-map-into-scheme}
设 $S$ 为概形，$A$ 为弱可容许拓扑 $S$-代数，并且对某个弱定义理想
$I\subset A$，$A/I$ 是局部环。设 $X$ 为 $S$ 上的概形。则自然映射
$$
\Mor_S(\Spec(A), X)
\longrightarrow
\Mor_S(\text{Spf}(A), X)
$$
是双射。
\end{lemma}

\begin{proof}
设 $\varphi:\text{Spf}(A)\to X$ 为态射。由于 $\Spec(A/I)$ 是局部的，
可知 $\varphi$ 将 $\Spec(A/I)$ 映入 $X$ 的某个仿射开子集 $U$。
然而，这又意味着对每个定义理想 $J$，$\Spec(A/J)$ 都映入 $U$。
因此可以应用引理 \ref{formal-spaces-lemma-map-into-affine}，可知 $\varphi$ 来自一个
$\Spec(A)\to X$ 的态射。这证明了该映射的满射性。
单射性的证明略去。
\end{proof}

\begin{lemma}
\label{formal-spaces-lemma-map-into-algebraic-space}
设 $S$ 为概形，$R$ 为完备局部 Noether $S$-代数，$X$ 为 $S$ 上的代数空间。
则自然映射
$$
\Mor_S(\Spec(R), X)
\longrightarrow
\Mor_S(\text{Spf}(R), X)
$$
是双射。
\end{lemma}

\begin{proof}
设 $\mathfrak m$ 为 $R$ 的极大理想。我们需要证明，对这样的 $R$，映射
$$
\Mor_S(\Spec(R), X) \longrightarrow
\lim \Mor_S(\Spec(R/\mathfrak m^n), X)
$$
是双射。

\medskip\noindent
单射性：设 $x,x':\Spec(R)\to X$ 为两个在右侧映到同一元素的态射。
考虑纤维积
$$
T = \Spec(R) \times_{(x, x'), X \times_S X, \Delta} X
$$
则 $T$ 是概形，且 $T\to\Spec(R)$ 局部有限型、为单态射、分离且局部拟有限；见《空间的态射》之引理
\ref{spaces-morphisms-lemma-properties-diagonal}。特别地，$T$ 局部 Noether；见《态射》之引理
\ref{morphisms-lemma-finite-type-noetherian}。令 $t\in T$ 为映到
$\Spec(R)$ 闭点的唯一点；由于 $x$ 与 $x'$ 在 $R/\mathfrak m$ 上相同，这一点存在。
则 $R\to\mathcal{O}_{T,t}$ 是 Noether 环的局部环映射，并且对所有 $n$，
$R/\mathfrak m^n\to\mathcal{O}_{T,t}/\mathfrak m^n\mathcal{O}_{T,t}$ 是同构
（因为 $x$ 和 $x'$ 在所有 $\Spec(R/\mathfrak m^n)$ 上相同）。
由于 $\mathcal{O}_{T,t}$ 单射地映入其完备化（见《代数》之引理
\ref{algebra-lemma-intersect-powers-ideal-module-zero}），可知
$R=\mathcal{O}_{T,t}$。因此 $x$ 和 $x'$ 在 $R$ 上相同。

\medskip\noindent
满射性：设 $(x_n)$ 是右侧的一个元素。选取概形 $U$ 及满的 \'etale 态射
$U\to X$。令 $x_0:\Spec(k)\to X$ 为在剩余域
$k=R/\mathfrak m$ 上诱导的态射。概形态射
$U\times_{X,x_0}\Spec(k)\to\Spec(k)$ 是满的 \'etale 态射。
因此 $U\times_{X,x_0}\Spec(k)$ 是 $k$ 的有限可分扩张的谱所成的非空互不相交并；见《态射》之引理
\ref{morphisms-lemma-etale-over-field}。于是可以找到有限可分扩张 $k'/k$ 和
$k'$-点 $u_0:\Spec(k')\to U$，使图表
$$
\xymatrix{
\Spec(k') \ar[d] \ar[r]_-{u_0} & U \ar[d] \\
\Spec(k) \ar[r]^-{x_0} & X
}
$$
交换。令 $R\subset R'$ 为 Noether 完备局部环的有限 \'etale 扩张，
其在剩余域上诱导 $k'/k$（见《代数》之引理
\ref{algebra-lemma-henselian-cat-finite-etale} 和
\ref{algebra-lemma-complete-henselian}）。令 $x'_n$ 为 $x_n$ 在
$\Spec(R'/\mathfrak m^nR')$ 上的限制。由《空间的态射更多内容》之引理
\ref{spaces-more-morphisms-lemma-etale-formally-etale}，可以找到
$(u'_n)\in\lim\Mor_S(\Spec(R'/\mathfrak m^nR'),U)$，其映到 $(x'_n)$。
由引理 \ref{formal-spaces-lemma-map-into-scheme}，族 $(u'_n)$ 来自唯一的态射
$u':\Spec(R')\to U$。令 $x':\Spec(R')\to X$ 为其复合。
注意，$R'\otimes_RR'$ 是完备局部 Noether 环的谱的有限积，仍适用当前讨论。
因此图表
$$
\xymatrix{
\Spec(R' \otimes_R R') \ar[r] \ar[d] & \Spec(R') \ar[d]^{x'} \\
\Spec(R') \ar[r]^{x'} & X
}
$$
由上面证明的单射性以及 $x'_n$ 是 $x_n$ 的限制（后者在
$R/\mathfrak m^n$ 上定义）而交换。由于 $\{\Spec(R')\to\Spec(R)\}$
是 fppf 覆盖，故 $x'$ 下降为态射 $x:\Spec(R)\to X$。
证明 $x_n$ 是 $x$ 在 $\Spec(R/\mathfrak m^n)$ 上的限制这一点略去。
\end{proof}

\begin{lemma}
\label{formal-spaces-lemma-adic-into-completion}
设 $S$ 为概形，$X$ 为 $S$ 上的代数空间。设 $T\subset |X|$ 为闭子集，且
$X\setminus T\to X$ 拟紧。设 $R$ 为完备局部 Noether $S$-代数。
则 adic 态射 $p:\text{Spf}(R)\to X_{/T}$ 对应于唯一的态射
$g:\Spec(R)\to X$，使得 $g^{-1}(T)=\{\mathfrak m_R\}$。
\end{lemma}

\begin{proof}
由于引理 \ref{formal-spaces-lemma-formal-completion-types}，$X_{/T}$ 是 adic* 的，
所以该陈述有意义（并且可以使用 adic 态射这一术语；见定义
\ref{formal-spaces-definition-adic-morphism}）。设 $p$ 已给定。由引理
\ref{formal-spaces-lemma-map-into-algebraic-space}，得到唯一态射 $g:\Spec(R)\to X$，
它对应于复合 $\text{Spf}(R)\to X_{/T}\to X$。
令 $Z\subset X$ 为 $T$ 上的约化诱导闭子空间结构。
% P10-E0371
包含态射 $Z\to X$ 对应于一个态射 $Z\to X_{/T}$。由于 $p$ 是 adic 的，
它可由代数空间表示，于是
$$
\text{Spf}(R) \times_{X_{/T}} Z = \text{Spf}(R) \times_X Z
$$
是带有到 $\text{Spf}(R)$ 的闭浸入的代数空间。（等式成立是因为
$X_{/T}\to X$ 是单态射。）因此这个纤维积等于 $\Spec(R/J)$，其中
理想 $J\subset R$ 包含某个 $n_0\geq1$ 的 $\mathfrak m_R^{n_0}$。
这意味着 $\Spec(R)\times_XZ$ 是 $\Spec(R)$ 的闭子概形，设
$\Spec(R)\times_XZ=\Spec(R/I)$，其与 $n\geq n_0$ 时的
$\Spec(R/\mathfrak m_R^n)$ 的交等于 $\Spec(R/J)$。
用代数语言说，这意味着对所有 $n\geq n_0$，
$I+\mathfrak m_R^n=J+\mathfrak m_R^n=J$。由 Krull 交定理，
这说明 $I=J$，结论得证。
\end{proof}





\section{形式代数空间的小 \'etale 位点}
\label{formal-spaces-section-etale-sites}

\noindent
下述定义的动机来自经典形式概形：形式概形
$(\mathfrak X,\mathcal{O}_\mathfrak X)$ 的底层拓扑空间是约化
$\mathfrak X_{red}$ 的底层拓扑空间。

\medskip\noindent
还有一个重要备注。假设 $X$ 是带有约化 $X_{red}$ 的代数空间
（《空间的性质》之定义 \ref{spaces-properties-definition-reduced-induced-space}）。则由《空间的态射更多内容》之定理
\ref{spaces-more-morphisms-theorem-topological-invariance} 和引理
\ref{spaces-more-morphisms-lemma-topological-invariance}，有
$$
X_{spaces, \etale} = X_{red, spaces, \etale},\quad
X_\etale = X_{red, \etale},\quad
X_{affine, \etale} = X_{red, affine, \etale}
$$
因此，在我们的形式代数空间恰好是代数空间的情形，下述定义不会与已有概念冲突。

\begin{definition}
\label{formal-spaces-definition-etale-sites}
设 $S$ 为概形，设 $X$ 为带有约化 $X_{red}$ 的形式代数空间
（引理 \ref{formal-spaces-lemma-reduction-formal-algebraic-space}）。
\begin{enumerate}
\item $X$ 的\textit{小 \'etale 位点} $X_\etale$ 是《空间的性质》之定义
\ref{spaces-properties-definition-etale-site} 中的位点 $X_{red,\etale}$；
\item 位点 $X_{spaces,\etale}$ 是《空间的性质》之定义
\ref{spaces-properties-definition-spaces-etale-site} 中的位点
$X_{red,spaces,\etale}$；
\item 位点 $X_{affine,\etale}$ 是《空间的性质》之引理
\ref{spaces-properties-lemma-alternative} 中的位点
$X_{red,affine,\etale}$。
\end{enumerate}
\end{definition}

\noindent
在引理 \ref{formal-spaces-lemma-identify-spaces-etale} 中我们将看到，
% P10-E0372
$X_{spaces,\etale}$ 可以用形式代数空间的态射来描述，这些态射可由代数空间表示且为 \'etale。
由《空间的性质》之引理
\ref{spaces-properties-lemma-compare-etale-sites} 和
\ref{spaces-properties-lemma-alternative}，有恒等
\begin{equation}
\label{formal-spaces-equation-etale-topos}
\Sh(X_\etale)=\Sh(X_{spaces,\etale})=\Sh(X_{affine,\etale})
\end{equation}
我们称之为 $X$ 的\textit{（小）\'etale 拓扑斯}。

\begin{lemma}
\label{formal-spaces-lemma-functoriality-etale-site}
设 $S$ 为概形，设 $f:X\to Y$ 为 $S$ 上形式代数空间的态射。
\begin{enumerate}
\item 存在连续函子
$Y_{spaces,\etale}\to X_{spaces,\etale}$，它诱导位点的态射
$$
f_{spaces, \etale}:X_{spaces,\etale}\to Y_{spaces,\etale}.
$$
\item 规则 $f\mapsto f_{spaces,\etale}$ 与复合相容，换言之
$(f\circ g)_{spaces,\etale}=f_{spaces,\etale}\circ g_{spaces,\etale}$；见《位点》之定义
\ref{sites-definition-composition-morphisms-sites}。
\item 与 $f_{spaces,\etale}$ 关联的拓扑斯态射，通过
\ref{formal-spaces-equation-etale-topos} 诱导拓扑斯态射
$f_{small}:\Sh(X_\etale)\to\Sh(Y_\etale)$，且其构造与复合相容。
\end{enumerate}
\end{lemma}

\begin{proof}
这里唯一要注意的是，由引理 \ref{formal-spaces-lemma-reduction-formal-algebraic-space}，
$f$ 诱导约化之间的态射 $X_{red}\to Y_{red}$。因此本引理立即由
代数空间态射的对应引理（《空间的性质》之引理
\ref{spaces-properties-lemma-functoriality-etale-site}）得到。
\end{proof}

\noindent
如果形式代数空间的态射 $X\to Y$ 是 \'etale 的，则拓扑斯态射
$\Sh(X_\etale)\to\Sh(Y_\etale)$ 是局部化。下面给出一个陈述。

\begin{lemma}
\label{formal-spaces-lemma-etale-morphism-topoi}
设 $S$ 为概形，$f:X\to Y$ 为 $S$ 上形式代数空间的态射。
假设 $f$ 可由代数空间表示且为 \'etale。此时存在余连续函子
$j:X_\etale\to Y_\etale$。拓扑斯态射 $f_{small}$ 是与 $j$ 关联的拓扑斯态射；见《位点》之引理
\ref{sites-lemma-cocontinuous-morphism-topoi}。此外，$j$ 也是连续的，
所以《位点》之引理 \ref{sites-lemma-when-shriek} 适用。
\end{lemma}

\begin{proof}
只要证明诱导的态射 $X_{red}\to Y_{red}$ 是 \'etale 的，结论就立即由代数空间的情形
（《空间的性质》之引理 \ref{spaces-properties-lemma-etale-morphism-topoi}）得到。
注意 $X\times_YY_{red}$ 是在约化代数空间 $Y_{red}$ 上 \'etale 的代数空间，
因而本身是约化的（依《空间的性质》第
\ref{spaces-properties-section-types-properties} 节中约化代数空间的定义）。
因此 $X_{red}=X\times_YY_{red}$，如所需。
\end{proof}

\begin{lemma}
\label{formal-spaces-lemma-affine-identify-affine-etale}
设 $S$ 为概形，$X$ 为 $S$ 上的仿射形式代数空间。则
$X_{affine,\etale}$ 等价于下列对象所成的范畴：对象是形式代数空间的态射
$\varphi:U\to X$，满足
\begin{enumerate}
\item $U$ 是仿射形式代数空间；
\item $\varphi$ 可由代数空间表示且为 \'etale。
\end{enumerate}
\end{lemma}

\begin{proof}
记 $\mathcal{C}$ 为引理中引入的范畴。注意，对 $\mathcal{C}$ 中的
$\varphi:U\to X$，$\varphi$ 可由概形表示且为仿射；见引理
\ref{formal-spaces-lemma-affine-representable-by-algebraic-spaces}。回忆
$X_{affine,\etale}=X_{red,affine,\etale}$。因此可以定义函子
$$
\mathcal{C}\longrightarrow X_{affine,\etale},\quad
(U\to X)\longmapsto U\times_XX_{red}
$$
因为 $U\times_XX_{red}$ 是仿射概形。

\medskip\noindent
为完成证明，我们构造准逆。按定义
\ref{formal-spaces-definition-affine-formal-algebraic-space}，写作 $X=\colim X_\lambda$。
对每个 $\lambda$，$X_{red}\subset X_\lambda$ 是一个增厚。因此对每个 $\lambda$ 有等价
$$
X_{red,affine,\etale}=X_{\lambda,affine,\etale}
$$
例如由《代数更多内容》之引理
\ref{more-algebra-lemma-locally-nilpotent-henselian} 得到。因此，若
$U_{red}\to X_{red}$ 是仿射的 \'etale 态射，则得到一族相容的仿射概形 \'etale 态射
$U_\lambda\to X_\lambda$，其与定义 $X$ 的系统中的过渡态射相容。于是令
$$
U=\colim U_\lambda
$$
作为 $X$ 上的仿射形式代数空间。构造给出
$U\times_XX_\lambda=U_\lambda$。这说明 $U\to X$ 可由代数空间表示且为 \'etale。
构造彼此为逆这一点的核验略去。
\end{proof}

\begin{lemma}
\label{formal-spaces-lemma-affine-etale-mcquillan}
设 $S$ 为概形，$X$ 为 $S$ 上的仿射形式代数空间。假设 $X$ 是 McQuillan，即
$X=\text{Spf}(A)$，其中 $A$ 是弱可容许拓扑 $S$-代数。
则 $(X_{affine,\etale})^{opp}$ 等价于下列范畴：
\begin{enumerate}
\item 对象是形如 $B^\wedge=\lim B/JB$ 的 $A$-代数，其中 $A\to B$ 是 \'etale 环映射，
且 $J$ 遍历 $A$ 的弱定义理想；
\item 态射是连续的 $A$-代数同态。
\end{enumerate}
\end{lemma}

\begin{proof}
结合引理 \ref{formal-spaces-lemma-affine-identify-affine-etale} 和引理
\ref{formal-spaces-lemma-etale} 即得。
\end{proof}

\begin{lemma}
\label{formal-spaces-lemma-identify-spaces-etale}
设 $S$ 为概形，$X$ 为 $S$ 上的形式代数空间。则 $X_{spaces,\etale}$ 等价于下列对象所成的范畴：
对象是形式代数空间的态射 $\varphi:U\to X$，且 $\varphi$ 可由代数空间表示并为 \'etale。
\end{lemma}

\begin{proof}
记 $\mathcal{C}$ 为引理中引入的范畴。
回忆 $X_{spaces, \etale} = X_{red, spaces, \etale}$。
因此可以定义函子
$$
\mathcal{C} \longrightarrow X_{spaces, \etale},\quad
(U \to X) \longmapsto U \times_X X_{red}
$$
因为 $U \times_X X_{red}$ 是在 $X_{red}$ 上 \'etale 的代数空间。

\medskip\noindent
为完成证明，我们构造一个准逆函子。
取 $X_{red, spaces, \etale}$ 中的对象
$\psi : V \to X_{red}$。
考虑函子
$U_{V, \psi} : (\Sch/S)_{fppf} \to \textit{Sets}$，其定义为
$$
U_{V, \psi}(T) = \{(a, b) \mid
a : T \to X,
\ b : T \times_{a, X} X_{red} \to V,
\ \psi \circ b = a|_{T \times_{a, X} X_{red}}\}
$$
我们断言，变换 $U_{V, \psi} \to X$，$(a, b) \mapsto a$，
定义了范畴 $\mathcal{C}$ 中的一个对象。
首先证明 $U_{V, \psi}$ 是形式代数空间。
注意，$U_{V, \psi}$ 是 fppf 位点上的层（略去一些细节）。
接下来，设 $X_i \to X$ 是由仿射形式代数空间给出的 \'etale 覆盖，
如定义 \ref{formal-spaces-definition-formal-algebraic-space} 中所述。
令 $V_i = V \times_{X_{red}} X_{i, red}$，并记
$\psi_i : V_i \to X_{i, red}$ 为投影。则
我们有
$$
U_{V, \psi} \times_X X_i = U_{V_i, \psi_i}
$$
这是形式上的论证，因为 $X_{i, red} = X_i \times_X X_{red}$
（这是由于 $X_i \to X$ 可由代数空间表示且为 \'etale）。
因此，只需证明 $U_{V_i, \psi_i}$ 是仿射形式代数空间，
因为这样就有一个覆盖 $U_{V_i, \psi_i} \to U_{V, \psi}$，
满足定义 \ref{formal-spaces-definition-formal-algebraic-space}。
另一方面，我们在引理 \ref{formal-spaces-lemma-etale-morphism-topoi} 的证明中已经看到，
$\psi_i : V_i \to X_i$ 是仿射形式代数空间之间的可表示且 \'etale
态射 $U_i \to X_i$ 的基变换。
于是，不难看出 $U_i = U_{V_i, \psi_i}$，如所需。

\medskip\noindent
我们略去验证 $U_{V, \psi} \to X$ 可由代数空间表示且为 \'etale。
这样便得到另一个方向的函子
$(V, \psi) \mapsto (U_{V, \psi} \to X)$。
构造彼此互逆的验证略去。
\end{proof}

\begin{lemma}
\label{formal-spaces-lemma-identify-affine-etale}
设 $S$ 为概形，$X$ 为 $S$ 上的形式代数空间。
则 $X_{affine, \etale}$ 等价于下列对象所成的范畴：对象是形式代数空间的态射
$\varphi : U \to X$，满足
\begin{enumerate}
\item $U$ 是仿射形式代数空间，
\item $\varphi$ 可由代数空间表示且为 \'etale。
\end{enumerate}
\end{lemma}

\begin{proof}
这由引理 \ref{formal-spaces-lemma-identify-spaces-etale} 和
\ref{formal-spaces-lemma-characterize-affine} 合并得到。
\end{proof}





\section{结构层}
\label{formal-spaces-section-structure-sheaf}


\noindent
如第 \ref{formal-spaces-section-etale-sites} 节所述，本节讨论的动机
来自经典形式概形。经典形式概形定义为二元组
$(\mathfrak X, \mathcal{O}_\mathfrak X)$，其中 $\mathcal{O}_\mathfrak X$
是拓扑环层。

\medskip\noindent
设 $X$ 为形式代数空间。$X$ 的一个{\it 结构层}
是 \'etale 位点 $X_\etale$ 上的拓扑环层 $\mathcal{O}_X$
（该位点在第 \ref{formal-spaces-section-etale-sites} 节中定义），并且对所有满足下列条件的情形，
都有
$$
\mathcal{O}_X(U_{red}) = \lim \Gamma(U_\lambda, \mathcal{O}_{U_\lambda})
$$
作为拓扑环成立：
\begin{enumerate}
\item $\varphi : U \to X$ 是形式代数空间的态射，
\item $U$ 是仿射形式代数空间，
\item $\varphi$ 可由代数空间表示且为 \'etale，
\item $U_{red} \to X_{red}$ 是 $X_\etale$ 中相应的仿射对象，见
引理 \ref{formal-spaces-lemma-identify-affine-etale}，
\item $U = \colim U_\lambda$ 是定义
\ref{formal-spaces-definition-affine-formal-algebraic-space} 中所述的 $U$ 的余极限表示。
\end{enumerate}
结构层确实存在，但其行为可能出人意料。

\begin{lemma}
\label{formal-spaces-lemma-structure-sheaf}
每个形式代数空间都有结构层。
\end{lemma}

\begin{proof}
设 $S$ 为概形，设 $X$ 为 $S$ 上的形式代数空间。
由（\ref{formal-spaces-equation-etale-topos}），只需在 $X_{affine, \etale}$ 上构造
$\mathcal{O}_X$ 为拓扑环层即可。
记 $\mathcal{C}$ 为如下范畴：对象是形式代数空间的态射
$\varphi : U \to X$，其中 $U$ 是仿射形式代数空间，
且 $\varphi$ 可由代数空间表示并为 \'etale。
由引理 \ref{formal-spaces-lemma-identify-affine-etale}，函子
$U \mapsto U_{red}$ 是等价
$\mathcal{C} \to X_{affine, \etale}$。
因此，按引理之前给出的规则，我们已经在
$X_{affine, \etale}$ 上得到拓扑环的预层 $\mathcal{O}_X$。
所以只需验证层条件。

\medskip\noindent
按 $X_{affine, \etale}$ 的定义，一个覆盖对应于有限族
$\{g_i : U_i \to U\}_{i = 1, \ldots, n}$，其中 $g_i$ 是 $\mathcal{C}$ 中的态射，
并且 $\{U_{i, red} \to U_{red}\}$ 是 \'etale 覆盖。
$g_i$ 可由代数空间表示（引理
\ref{formal-spaces-lemma-permanence-representable}），因而是仿射的（引理
\ref{formal-spaces-lemma-affine-representable-by-algebraic-spaces}）。
% P10-E0373
于是 $g_i$ 是 \'etale 的（这形式地由《空间的性质》之引理
\ref{spaces-properties-lemma-etale-permanence} 得到，因为 $U_i$ 与 $U$
在《Bootstrap》第 \ref{bootstrap-section-representable-by-spaces-properties} 节的意义下
均 \'etale 地态射到 $X$）。
最后，按定义 \ref{formal-spaces-definition-affine-formal-algebraic-space} 写
$U = \colim U_\lambda$。

\medskip\noindent
准备工作完成后，层性质如下证明。对每个 $\lambda$，令
$U_{i, \lambda} = U_i \times_U U_\lambda$，并令
$U_{ij, \lambda} = (U_i \times_U U_j) \times_U U_\lambda$。
由上述论证，这些是仿射概形，$\{U_{i, \lambda} \to U_\lambda\}$ 是 \'etale 覆盖，
且
$U_{ij, \lambda} = U_{i, \lambda} \times_{U_\lambda} U_{j, \lambda}$。
此外有 $U_i = \colim U_{i, \lambda}$ 以及
$U_i \times_U U_j = \colim U_{ij, \lambda}$。
对每个 $\lambda$，有正合序列
$$
0 \to
\Gamma(U_\lambda, \mathcal{O}_{U_\lambda}) \to
\prod\nolimits_i \Gamma(U_{i, \lambda}, \mathcal{O}_{U_{i, \lambda}}) \to
\prod\nolimits_{i, j} \Gamma(U_{ij, \lambda}, \mathcal{O}_{U_{ij, \lambda}})
$$
因为在 $U_\lambda$ 上、关于 \'etale 拓扑，结构层满足层条件
（见《\'Etale 上同调》之命题
\ref{etale-cohomology-proposition-quasi-coherent-sheaf-fpqc}）。
由于极限与极限交换，这些正合序列的逆极限仍为正合序列：
$$
0 \to
\lim \Gamma(U_\lambda, \mathcal{O}_{U_\lambda}) \to
\prod\nolimits_i \lim \Gamma(U_{i, \lambda}, \mathcal{O}_{U_{i, \lambda}}) \to
\prod\nolimits_{i, j} \lim
\Gamma(U_{ij, \lambda}, \mathcal{O}_{U_{ij, \lambda}})
$$
这意味着
$$
0 \to
\mathcal{O}_X(U_{red}) \to
\prod\nolimits_i \mathcal{O}_X(U_{i, red}) \to
\prod\nolimits_{i, j} \mathcal{O}_X((U_i \times_U U_j)_{red})
$$
作为阿贝尔群序列是正合的。我们仍需证明
$\mathcal{O}_X(U_{red})$ 上的极限拓扑由
$\prod_{i = 1, \ldots, n} \mathcal{O}_X(U_{i, red})$ 上的极限拓扑诱导。
映射
$\mathcal{O}_X(U_{red}) \to \mathcal{O}_{U_\lambda}(U_\lambda)$
的核构成开子模的基本系。由于这些核是映射
$\prod_i \mathcal{O}_X(U_{i, red}) \to
\prod_i \mathcal{O}_{U_{i, \lambda}}(U_{i, \lambda})$
的核的逆像，而这些核构成
% P10-E0374
$\mathcal{O}_X(U_{red}) \to \mathcal{O}_{U_\lambda}(U_\lambda)$
的开子模基本系，结论得证。
\end{proof}

\begin{remark}
\label{formal-spaces-remark-not-enough-sections}
结构层并不总是有“足够多的截面”。
在例子（第 \ref{examples-section-affine-formal-algebraic-space} 节）中，
我们已经看到存在不是 McQuillan 的仿射形式代数空间，
甚至还有一些例子，其点不能由正则函数分离。
\end{remark}


\noindent
在下一个引理中，我们证明可数指标的仿射形式代数空间上的结构层
具有消失的高阶上同调。
对于非可数指标的情形，推测这一结论一般不成立。

\begin{lemma}
\label{formal-spaces-lemma-higher-vanishing-structure-sheaf}
若 $X$ 是可数指标的仿射形式代数空间，则当 $n > 0$ 时有
$H^n(X_\etale, \mathcal{O}_X) = 0$。
\end{lemma}

\begin{proof}
我们可以在 $X_{affine, \etale}$ 上工作，因为它给出同一个拓扑斯。
我们将应用《位点上的上同调》之引理
\ref{sites-cohomology-lemma-cech-vanish-collection}
来证明消失性。由于 $X_{affine, \etale}$ 有有限不交并，
这将问题化为覆盖的 Čech 复形，其中覆盖由
$X_{affine, \etale} = X_{red, affine, \etale}$ 中的单个箭头
$\{U_{red} \to V_{red}\}$ 给出
（见《\'Etale 上同调》之引理
\ref{etale-cohomology-lemma-cech-complex}）。
因此我们必须证明
$$
0 \to \mathcal{O}_X(V_{red}) \to \mathcal{O}_X(U_{red}) \to
\mathcal{O}_X(U_{red} \times_{V_{red}} U_{red}) \to \ldots
$$
是正合的。下面在 $V_{red} = X_{red}$ 的情形证明这一点；
一般情形的证明完全相同。

\medskip\noindent
回忆 $X = \text{Spf}(A)$，其中 $A$ 是弱可容许拓扑环，
并且具有可数的弱定义理想基本系。
由引理 \ref{formal-spaces-lemma-affine-identify-affine-etale} 和
\ref{formal-spaces-lemma-affine-etale-mcquillan}，$X_{affine, \etale}$ 中的对象
$U_{red}$ 对应于仿射形式代数空间的态射 $U \to X$，
该态射可由代数空间表示且为 \'etale，并且
$U = \text{Spf}(B^\wedge)$，其中 $B$ 是 \'etale 的 $A$-代数。
根据结构层的定义规则，我们看到
$$
\mathcal{O}_X(U_{red}) = B^\wedge
$$
回忆 $B^\wedge = \lim B/JB$，其中极限遍历
$A$ 的弱定义理想 $J \subset A$。
逐一考察定义可得
$$
\mathcal{O}_X(U_{red} \times_{X_{red}} U_{red}) = (B \otimes_A B)^\wedge
$$
依此类推。由于 $U \to X$ 是覆盖，映射 $A \to B$
是忠实平坦的，见引理 \ref{formal-spaces-lemma-etale-surjective}。
因此复形
$$
0 \to A \to B \to B \otimes_A B \to B \otimes_A B \otimes_A B \to \ldots
$$
是{\bf 泛正合}的，见《下降》之引理
\ref{descent-lemma-ff-exact}。
我们的目标是证明
$$
H^n(0 \to A^\wedge \to B^\wedge \to (B \otimes_A B)^\wedge
\to (B \otimes_A B \otimes_A B)^\wedge \to \ldots)
$$
当 $n > 0$ 时为零。为看清其中的机制，将上述（完备化之前的）
正合复形拆成短正合序列
% P10-E0375
$$
0 \to A \to B \to M_1 \to 0,\quad
0 \to M_i \to B^{\otimes_A i + 1} \to M_{i + 1} \to 0
$$
由上面的论述，这些都是泛正合的短正合序列。
因此对 $A$ 的每个理想 $J$，都有
$JM_i = M_i \cap J(B^{\otimes_A i + 1})$。
特别地，$M_i$ 作为 $B^{\otimes_A i + 1}$ 的子模所承继的拓扑，
与 $M_i$ 作为 $B^{\otimes_A i}$ 的商模所承继的拓扑相同。
由于 $A$ 中存在可数的弱定义理想基本系，序列
$$
0 \to A^\wedge \to B^\wedge \to M_1^\wedge \to 0,\quad
0 \to M_i^\wedge \to (B^{\otimes_A i + 1})^\wedge \to M_{i + 1}^\wedge \to 0
$$
由引理 \ref{formal-spaces-lemma-ses} 保持正合。这就证明了引理。
\end{proof}

\begin{remark}
\label{formal-spaces-remark-bad-quasi-coherent}
即使结构层具有良好性质，也不意味着存在良好的拟凝聚模理论。
例如，在例子（第 \ref{examples-section-nonabelian-QCoh} 节）中，
我们看到，对于几乎所有诺特仿射形式代数空间，最自然的拟凝聚模概念
都会导出一个非阿贝尔的模范畴。
\end{remark}





\section{形式代数空间的余极限}
\label{formal-spaces-section-colimits-of-formal}

\noindent
本节将第 \ref{formal-spaces-section-global-colimits} 节的结果推广到
形式代数空间的态射系。
我们指出，在下面的引理中，条件
“$f_{\lambda \mu} : X_\lambda \to X_\mu$ 是闭浸入，
并诱导约化空间 $X_{\lambda, red} \to X_{\mu, red}$ 的同构”
可以改写为
“$f_{\lambda \mu}$ 可由代数空间表示且为增厚”。

\begin{lemma}
\label{formal-spaces-lemma-colimit-affine-formal}
设 $S$ 为概形。设给定有向集 $\Lambda$ 以及其上的仿射形式代数空间系统
$(X_\lambda, f_{\lambda \mu})$，其中每个
$f_{\lambda \mu} : X_\lambda \to X_\mu$ 都是闭浸入，
并诱导约化空间 $X_{\lambda, red} \to X_{\mu, red}$ 的同构。
则 $X = \colim_{\lambda \in \Lambda} X_\lambda$
是 $S$ 上的仿射形式代数空间。
\end{lemma}

\begin{proof}
我们可以写成
$X_\lambda = \colim_{\omega \in \Omega_\lambda} X_{\lambda, \omega}$，
其中这是有向集 $\Omega_\lambda$ 上仿射概形的余极限，
且过渡态射
$X_{\lambda, \omega} \to X_{\lambda, \omega'}$
是增厚。对每个 $\lambda, \mu \in \Lambda$ 以及
$\omega \in \Omega_\lambda$，当 $\mu \geq \lambda$ 时，
存在 $\omega' \in \Omega_\mu$，使得态射
$X_{\lambda, \omega} \to X_\mu$ 分解为经过 $X_{\mu, \omega'}$，
见引理 \ref{formal-spaces-lemma-factor-through-thickening}。于是
态射 $X_{\lambda, \omega} \to X_{\mu, \omega'}$
是闭浸入，并在约化上诱导同构，因而是增厚。令
$\Omega = \coprod_{\lambda \in \Lambda} \Omega_\lambda$，
并规定 $(\lambda, \omega) \leq (\mu, \omega')$ 当且仅当
$\lambda \leq \mu$ 且 $X_{\lambda, \omega} \to X_\mu$ 分解为经过
$X_{\mu, \omega'}$。由上述事实可知，$\Omega$ 是有向集，并且
$X = \colim_{\lambda \in \Lambda} X_\lambda =
\colim_{(\lambda, \omega) \in \Omega} X_{\lambda, \omega}$。
证明完成。
\end{proof}

\begin{lemma}
\label{formal-spaces-lemma-colimit-formal-spaces-is-formal-space}
设 $S$ 为概形。设给定有向集 $\Lambda$ 以及其上的形式代数空间系统
$(X_\lambda, f_{\lambda \mu})$，其中每个
$f_{\lambda \mu} : X_\lambda \to X_\mu$ 都是闭浸入，
并诱导约化空间 $X_{\lambda, red} \to X_{\mu, red}$ 的同构。
则 $X = \colim_{\lambda \in \Lambda} X_\lambda$
是 $S$ 上的形式代数空间。
\end{lemma}

\begin{proof}
由于我们在 fppf 层范畴中取余极限，可知 $X$ 是一个层。
选取并固定 $\lambda \in \Lambda$。按定义
\ref{formal-spaces-definition-formal-algebraic-space} 选取覆盖
$\{X_{i, \lambda} \to X_\lambda\}$。特别地，
$\{X_{i, \lambda, red} \to X_{\lambda, red}\}$ 是由仿射概形给出的
\'etale 覆盖。
对每个 $\mu \geq \lambda$，存在笛卡尔图
$$
\xymatrix{
X_{i, \lambda} \ar[r] \ar[d] & X_{i, \mu} \ar[d] \\
X_\lambda \ar[r] & X_\mu
}
$$
其竖直箭头为 \'etale。事实上，\'etale 态射
$X_{i, \lambda, red} \to X_{\lambda, red} = X_{\mu, red}$
对应于形式代数空间的 \'etale 态射 $X_{i, \mu} \to X_\mu$，
其中 $X_{i, \mu}$ 是仿射形式代数空间，见
引理 \ref{formal-spaces-lemma-affine-identify-affine-etale}。
同一个引理说明 $X_{i, \mu}$ 到 $X_\lambda$ 的基变换
与 $X_{i, \lambda}$ 一致。还可推出，当
$\mu' \geq \mu \geq \lambda$ 时，
$X_{i, \mu} = X_\mu \times_{X_{\mu'}} X_{i, \mu'}$。
令 $X_i = \colim X_{i, \mu}$。则
$X_{i, \mu} = X_i \times_X X_\mu$（作为函子）。
由于从仿射（或拟紧）概形 $T$ 到
$X = \colim X_\mu$ 的任意态射都落入某个 $X_\mu$，
% P10-E0376
因此可得 $\colim X_{i, \mu} \to \colim X_\mu$ 是 \'etale。
于是，只要证明 $\colim X_{i, \mu}$ 是仿射形式代数空间，
引理便得证。
注意，态射 $X_{i, \mu} \to X_{i, \mu'}$
是闭浸入 $X_\mu \to X_{\mu'}$ 的基变换，因而是闭浸入。
最后，由于 $X_{\mu, red} \to X_{\mu', red}$ 是同构，
态射 $X_{i, \mu, red} \to X_{i, \mu', red}$ 也是同构。
因此问题归结为引理 \ref{formal-spaces-lemma-colimit-affine-formal} 中讨论的情形。
\end{proof}





\section{再次完备化}
\label{formal-spaces-section-recompletion}

\noindent
本节定义沿约化空间的闭子集对形式代数空间进行完备化。
这是第 \ref{formal-spaces-section-completion} 节的自然推广。

\begin{lemma}
\label{formal-spaces-lemma-completion-affine-formal-is-affine-formal}
设 $S$ 为概形，设 $X$ 为 $S$ 上的仿射形式代数空间。
设 $T \subset |X_{red}|$ 为闭子集。则函子
$$
X_{/T} : (\Sch/S)_{fppf} \longrightarrow \textit{Sets},\quad
U \longmapsto \{f : U \to X : f(|U|) \subset T\}
$$
是仿射形式代数空间。
\end{lemma}

\begin{proof}
按定义 \ref{formal-spaces-definition-affine-formal-algebraic-space} 写
$X = \colim X_\lambda$。于是 $X_{\lambda, red} = X_{red}$，
我们可以并且将 $T$ 看作 $|X_\lambda| = |X_{\lambda, red}|$ 的闭子集。
由引理 \ref{formal-spaces-lemma-completion-affine-is-affine-formal-algebraic-space}，
对每个 $\lambda$，完备化 $(X_\lambda)_{/T}$ 是仿射形式代数空间。
过渡态射 $(X_\lambda)_{/T} \to (X_\mu)_{/T}$ 是闭浸入，
因为它们是过渡态射 $X_\lambda \to X_\mu$ 的基变换，
见引理 \ref{formal-spaces-lemma-map-completions-representable}。
此外，由引理 \ref{formal-spaces-lemma-reduction-completion}，
态射 $((X_\lambda)_{/T})_{red} \to ((X_\mu)_{/T})_{red}$
是同构。
由于 $X_{/T} = \colim (X_\lambda)_{/T}$，由引理
\ref{formal-spaces-lemma-colimit-affine-formal} 即得结论。
\end{proof}

\begin{lemma}
\label{formal-spaces-lemma-completion-fas-is-fas}
设 $S$ 为概形，设 $X$ 为 $S$ 上的形式代数空间。
设 $T \subset |X_{red}|$ 为闭子集。则函子
$$
X_{/T} : (\Sch/S)_{fppf} \longrightarrow \textit{Sets},\quad
U \longmapsto \{f : U \to X \mid f(|U|) \subset T\}
$$
是形式代数空间。
\end{lemma}

\begin{proof}
函子 $X_{/T}$ 是 fppf 层，因为若 $\{U_i \to U\}$ 是 fppf 覆盖，
则 $\coprod |U_i| \to |U|$ 是满射。

\medskip\noindent
按定义 \ref{formal-spaces-definition-formal-algebraic-space} 选取覆盖
$\{g_i : X_i \to X\}_{i \in I}$。
态射 $X_i \times_X X_{/T} \to X_{/T}$ 是 \'etale 的
（见《空间》之引理
\ref{spaces-lemma-base-change-representable-transformations-property}），
并且映射 $\coprod X_i \times_X X_{/T} \to X_{/T}$ 是层的满射。
因此只需证明 $X_{/T} \times_X X_i$ 是仿射形式代数空间。
$X_i \times_X X_{/T}$ 的一个 $U$-值点是态射 $U \to X_i$，
其像包含于闭子集
$|g_{i, red}|^{-1}(T) \subset |X_{i, red}|$。
于是结论由引理
\ref{formal-spaces-lemma-completion-affine-formal-is-affine-formal} 得到。
\end{proof}

\begin{definition}
\label{formal-spaces-definition-completion-formal-algebraic-space}
设 $S$ 为概形，设 $X$ 为 $S$ 上的形式代数空间。
设 $T \subset |X_{red}|$ 为闭子集。引理
\ref{formal-spaces-lemma-completion-is-formal-algebraic-space} 中的形式代数空间
$X_{/T}$ 称为沿 $T$ 的 $X$ 的{\it 完备化}。
\end{definition}

\noindent
设 $f : X \to X'$ 是概形 $S$ 上形式代数空间的态射。
设 $T \subset |X_{red}|$ 与 $T' \subset |X'_{red}|$ 为闭子集，
且 $|f_{red}|(T) \subset T'$。显然，$f$ 定义了完备化之间的形式代数空间态射
$$
X_{/T} \longrightarrow X'_{/T'}
$$

\begin{lemma}
\label{formal-spaces-lemma-map-recompletions}
设 $S$ 为概形，设 $f : X' \to X$ 是
$S$ 上形式代数空间的态射。设 $T \subset |X_{red}|$
为闭子集，并令 $T' = |f_{red}|^{-1}(T) \subset |X'_{red}|$。
则
$$
\xymatrix{
X'_{/T'} \ar[r] \ar[d] & X' \ar[d]^f \\
X_{/T} \ar[r] & X
}
$$
是 $S$ 上形式代数空间的笛卡尔图。
\end{lemma}

\begin{proof}
事实上，按构造水平箭头是单态射。因此只需证明：
概形 $U$ 到 $X'$ 的态射 $g : U \to X'$ 定义
$X'_{/T}$ 中的点，当且仅当 $f \circ g$ 定义 $X_{/T}$ 中的点。
换言之，我们需证明
% P10-E0377
$g(U)$ 包含于 $T' \subset |X'_{red}|$
当且仅当 $(f \circ g)(U)$ 包含于 $T \subset |X_{red}|$。
由我们对 $T'$ 取为 $T$ 的逆像，立即得到这一点。
\end{proof}

\begin{lemma}
\label{formal-spaces-lemma-reduction-recompletion}
设 $S$ 为概形，设 $X$ 为 $S$ 上的形式代数空间。
设 $T \subset |X_{red}|$ 为闭子集。完备化 $X_{/T}$ 沿 $T$ 的约化
$(X_{/T})_{red}$ 是 $X_{red}$ 的约化诱导闭子空间 $Z$，
它对应于 $T$。
\end{lemma}

\begin{proof}
这由引理 \ref{formal-spaces-lemma-reduction-formal-algebraic-space}、
《空间的性质》之定义
\ref{spaces-properties-definition-reduced-induced-space}
（该定义使用《空间的性质》之引理
\ref{spaces-properties-lemma-reduced-closed-subspace} 构造 $Z$）
以及 $X_{/T}$ 的定义得到：$Z$ 与 $(X_{/T})_{red}$ 都是约化代数空间，
并由同一映射性质刻画：
% P10-E0378
源为约化代数空间的态射 $g : Y \to X$ 可分解通过它们，
当且仅当 $|Y|$ 映入 $T \subset |X|$。
\end{proof}

\begin{lemma}
\label{formal-spaces-lemma-recompletion-affine-types}
设 $S$ 为概形，设 $X$ 为 $S$ 上的仿射形式代数空间。
设 $T \subset X_{red}$ 为闭子集，设 $X_{/T}$ 是沿 $T$ 的 $X$ 的形式完备化。
则
\begin{enumerate}
\item $X_{/T}$ 是仿射形式代数空间，
\item 若 $X$ 是 McQuillan，则 $X_{/T}$ 是 McQuillan，
\item 若 $|X_{red}| \setminus T$ 是拟紧的且 $X$ 是可数指标的，
则 $X_{/T}$ 是可数指标的，
\item 若 $|X_{red}| \setminus T$ 是拟紧的且 $X$ 是 adic* 的，
则 $X_{/T}$ 是 adic* 的，
\item 若 $X$ 是诺特的，则 $X_{/T}$ 是诺特的。
\end{enumerate}
\end{lemma}

\begin{proof}
(1) 是引理 \ref{formal-spaces-lemma-completion-affine-formal-is-affine-formal}。
若 $X$ 是 McQuillan，则对某个弱可容许拓扑环 $A$ 有
$X = \text{Spf}(A)$。于是
$X_{/T} \to X \to \Spec(A)$ 满足引理
\ref{formal-spaces-lemma-mcquillan-affine-formal-algebraic-space} 的性质 (2)，
因此 $X_{/T}$ 是 McQuillan，见定义
\ref{formal-spaces-definition-types-affine-formal-algebraic-space}。

\medskip\noindent
假设 $X$ 与 $T$ 如 (3) 中所述。
则 $X = \text{Spf}(A)$，其中 $A$ 有弱定义理想的基本系
$A \supset I_1 \supset I_2 \supset I_3 \supset \ldots$，
见引理 \ref{formal-spaces-lemma-countably-indexed}。
由《代数》之引理 \ref{algebra-lemma-qc-open}，
可找到有限生成理想
$\overline{J} = (\overline{f}_1, \ldots, \overline{f}_r) \subset A/I_1$，
使得 $T$ 在 $\Spec(A/I_1) = |X_{red}|$ 内由 $\overline{J}$ 定义。
取 $f_i \in A$ 为 $\overline{f}_i$ 的提升。
若 $Z = \Spec(B)$ 是仿射概形，且 $g : Z \to X$ 是
（集合论意义下）像包含于 $T$ 的态射，则
$g^\sharp : A \to B$ 对某个 $n$ 经由 $A/I_n$ 分解，
并且对每个 $i$，$g^\sharp(f_i)$ 在 $B$ 中幂零。
因此，对于某些 $n, m \geq 1$，理想
$J_{m, n} = (f_1, \ldots, f_r)^m + I_n$ 映到 $B$ 中为零。
由此 $X_{/T}$ 是
$\lim_{n, m} A/J_{m, n}$ 的形式谱，因而是可数指标的。
这证明了 (3)。

\medskip\noindent
(4) 的证明。这里论证与 (3) 相同。
不过此时可以取 $I_n = I^n$，其中 $I \subset A$ 是有限生成理想。
于是显然 $X_{/T}$ 是 $\lim A/J^n$ 的形式谱，
其中 $J = (f_1, \ldots, f_r) + I$。细节略去。

\medskip\noindent
(5) 的证明。在此情形，$X_{red}$ 是诺特环的谱，
因而满足 $|X_{red}| \setminus T$ 拟紧的假设。
因此如 (4) 的证明，我们看到 $X_{/T}$ 是 $\lim A/J^n$ 的谱，
这是诺特的 adic 拓扑环，见《代数》之引理
\ref{algebra-lemma-completion-Noetherian-Noetherian}。
\end{proof}

\begin{lemma}
\label{formal-spaces-lemma-recompletion-types}
设 $S$ 为概形，设 $X$ 为 $S$ 上的形式代数空间。
设 $T \subset X_{red}$ 为闭子集，设 $X_{/T}$ 是沿 $T$ 的 $X$ 的形式完备化。
则
\begin{enumerate}
\item 若 $X_{red} \setminus T \to X_{red}$ 是拟紧的且 $X$ 局部可数指标，
则 $X_{/T}$ 局部可数指标，
\item 若 $X_{red} \setminus T \to X_{red}$ 是拟紧的且 $X$ 局部 adic*，
则 $X_{/T}$ 局部 adic*，并且
\item 若 $X$ 局部诺特，则 $X_{/T}$ 局部诺特。
\end{enumerate}
\end{lemma}

\begin{proof}
按定义 \ref{formal-spaces-definition-formal-algebraic-space} 选取覆盖
$\{X_i \to X\}$。令 $T_i \subset X_{i, red}$ 为 $T$ 的逆像。
由引理 \ref{formal-spaces-lemma-map-recompletions}，有
$X_i \times_X X_{/T} = (X_i)_{/T_i}$。
因此 $\{(X_i)_{/T_i} \to X_{/T}\}$ 是
满足定义 \ref{formal-spaces-definition-formal-algebraic-space} 的覆盖。
此外，若 $X_{red} \setminus T \to X_{red}$ 是拟紧的，
则 $X_{i, red} \setminus T_i \to X_{i, red}$ 也是拟紧的；
若 $X$ 局部可数指标、局部 adic* 或局部诺特，则
% P10-E0379
$X_i$ 分别是可数指标、adic* 或诺特的。
因此引理由仿射情形（即引理
\ref{formal-spaces-lemma-recompletion-affine-types}）得到。
\end{proof}







\section{沿闭子空间的完备化}
\label{formal-spaces-section-completion-subspace}

\noindent
% P10-E0380
本节是关于相对于闭子空间的完备化的第
\ref{formal-spaces-section-completion} 节的类似物。

\begin{definition}
\label{formal-spaces-definition-completion-subspace}
设 $S$ 为概形，设 $X$ 为 $S$ 上的代数空间。
设 $Z \subset X$ 为闭子空间，并记 $Z_n \subset X$ 为 $n$ 阶无穷小邻域。
形式代数空间
$$
X^\wedge_Z = \colim Z_n
$$
（见引理 \ref{formal-spaces-lemma-colimit-formal-spaces-is-formal-space}）
称为沿 $Z$ 的 $X$ 的{\it 完备化}。
\end{definition}

\noindent
例如，若 $X = \Spec(A)$ 且 $Z$ 由理想 $I \subset A$ 定义，
则有
$$
X^\wedge_Z = \colim \Spec(A/I^n) = \text{Spf}(B)
$$
其中 $B = \lim A/I^n$，极限取极限拓扑。注意，$B$ 是
弱 adic 拓扑环，但一般不是 adic 的。

\medskip\noindent
回到一般情形，若 $T = |Z|$，则存在典范态射
$X^\wedge_Z \to X_{/T}$，比较沿 $Z$ 与沿 $T$ 的完备化
（第 \ref{formal-spaces-section-completion} 节），它不一定是同构
（见引理 \ref{formal-spaces-lemma-when-isomorphism}）。

\medskip\noindent
设 $f : X \to X'$ 是概形 $S$ 上代数空间的态射。
假设 $Z \subset X$ 与 $Z' \subset X'$ 是闭子空间，
且 $f|_Z$ 将 $Z$ 映入 $Z'$ 并诱导态射 $Z \to Z'$。
显然，$f$ 定义形式代数空间的态射
$$
X^\wedge_Z \longrightarrow (X')^\wedge_{Z'}
$$
介于这些完备化之间。

\begin{lemma}
\label{formal-spaces-lemma-map-completions-subspaces-representable}
设 $S$ 为概形，设 $f : X' \to X$ 是 $S$ 上代数空间的态射。
设 $Z \subset X$ 为闭子空间，并令 $Z' = f^{-1}(Z) = X' \times_X Z$。
则
$$
\xymatrix{
(X')^\wedge_{Z'} \ar[r] \ar[d] & X' \ar[d]^f \\
X^\wedge_Z \ar[r] & X
}
$$
是层的笛卡尔图。特别地，态射
$(X')^\wedge_{Z'} \to X^\wedge_Z$ 可由代数空间表示。
\end{lemma}

\begin{proof}
事实上，设 $Y \to X$ 是从概形到 $X$ 的态射，并且
$Y \to X$ 分解通过 $Z$。则 $Y \times_X X' \to X$
是代数空间的态射，且 $Y \times_X X' \to X'$
分解通过 $Z'$。由于对所有 $n \geq 1$ 都有
$Z'_n = X' \times_X Z_n$，无穷小邻域也具有同样性质。
因此函子的笛卡尔方图由公式
$X^\wedge_Z = \colim Z_n$ 以及
$(X')^\wedge_{Z'} = \colim Z'_n$ 得到。
\end{proof}

\begin{lemma}
\label{formal-spaces-lemma-reduction-completion-subspace}
设 $S$ 为概形，设 $X$ 为 $S$ 上的代数空间。
设 $Z \subset X$ 为闭子空间。完备化 $X^\wedge_Z$ 沿 $Z$ 的约化
$(X^\wedge_Z)_{red}$ 是 $Z_{red}$。
\end{lemma}

\begin{proof}
略。
\end{proof}

\begin{lemma}
\label{formal-spaces-lemma-affine-formal-completion-subspace-types}
设 $S$ 为概形，设 $X = \Spec(A)$ 是 $S$ 上的仿射概形。
设 $Z \subset X$ 是对应于理想 $I \subset A$ 的闭子概形。则
\begin{enumerate}
\item 仿射形式代数空间 $X^\wedge_Z$ 是弱 adic 的，
\item 若 $I$ 有限生成，则 $X^\wedge_Z = \text{Spf}(A^\wedge)$，
其中 $A^\wedge$ 是 $A$ 的 $I$-进完备化，
\item 若 $Z \to X$ 是有限呈示的，即 $I$ 有限生成，
则 $X^\wedge_Z$ 是 adic* 的，
\item 若 $X$ 是诺特的，则 $X^\wedge_Z$ 是诺特的。
\end{enumerate}
\end{lemma}

\begin{proof}
略。
\end{proof}

\begin{lemma}
\label{formal-spaces-lemma-formal-completion-subspace-types}
设 $S$ 为概形，设 $X$ 为 $S$ 上的代数空间。
设 $Z \subset X$ 是闭子空间，设 $X^\wedge_Z$ 是沿 $Z$ 的 $X$ 的形式完备化。
\begin{enumerate}
\item 形式代数空间 $X^\wedge_Z$ 是局部弱 adic 的，
\item 若 $Z \to X$ 是有限呈示的，则 $X^\wedge_Z$ 是局部 adic* 的，
\item 若 $X$ 局部诺特，则
% P10-E0381
$X_Z$ 是局部诺特的。
\end{enumerate}
\end{lemma}

\begin{proof}
略。
\end{proof}

\begin{lemma}
\label{formal-spaces-lemma-when-isomorphism}
设 $S$ 为概形，设 $X$ 为 $S$ 上的代数空间。
设 $Z \subset X$ 是闭子空间，令 $T = |Z|$。
典范态射 $c : X^\wedge_Z \to X_{/T}$ 是同构，当 $Z \to X$
有限呈示时成立，但一般不成立。
\end{lemma}

\begin{proof}
两种构造都与 \'etale 局部化交换，因此只需在 $X$ 仿射时证明。
设 $X = \Spec(A)$，且 $Z$ 对应于理想 $I \subset A$。
若 $I$ 有限生成，则 $X^\wedge_Z$ 与 $X_{/T}$ 都等于
$\text{Spf}(A^\wedge)$，其中 $A^\wedge$ 是 $A$ 的 $I$-进完备化，
见引理 \ref{formal-spaces-lemma-affine-formal-completion-types} 和
\ref{formal-spaces-lemma-affine-formal-completion-subspace-types}。
若 $A = \mathbf{Z}[x_1, x_2, \ldots]$ 且
$I = (x_1, x_2, \ldots)$，则
$\Spec(A/(x_n^n, n \geq 1)) \to X$ 经由 $X_{/T}$ 分解，
但不经由 $X^\wedge_Z$ 分解，因此 $c$ 不是同构。
\end{proof}
% Generated mechanically from the frozen restricted.tex target.
% Do not edit this generated file; edit the bound locale source instead.

% OK, start here.
%

\setcounter{chapter}{87}
\chapter{形式空间的代数化}


\phantomsection
\label{restricted-section-phantom}


\section{引言}
\label{restricted-section-introduction}

\noindent
本章的主要目标是证明 Artin 关于膨胀的定理，见定理
\ref{restricted-theorem-dilatations}；关于压缩的结果将在 Artin 公理的
\ref{artin-section-contractions} 节中讨论。这两个结果都用到关于形式代数空间的材料，
因此本章中部将延续上一章对形式代数空间的讨论，见《形式空间》之
\ref{formal-spaces-section-introduction} 节。本章第一部分专门讨论代数预备知识，
主要涉及 rig-\,'etale 代数的代数化。

\medskip\noindent
设 $A$ 为诺特环，$I\subset A$ 为理想。本章第一部分（第
\ref{restricted-section-two-categories}--\ref{restricted-section-approximation-principal} 节）讨论在诺特环
$A$ 上拓扑有限型的 $I$-进完备代数 $B$。我们将证明，对于某个受限幂级数环中的理想
$J$（闭理想），有 $B=A\{x_1,\ldots,x_n\}/J$（其中 $A$ 赋予 $I$-进拓扑）。
我们展示朴素余切复形 $\NL_{B/A}^\wedge$ 的良好定义。若 $I$ 的某个幂湮灭
$\NL_{B/A}^\wedge$，则称 $B$ 是 $(A,I)$ 上的 rig-\,'etale 代数；rig-光滑代数有类似概念。
若 $A$ 是 G-环，则利用 Popescu 定理可证明，$(A,I)$ 上的任意 rig-光滑代数 $B$ 都是某个有限型
$A$-代数的完备化；非正式地说，我们可以将 $B$“代数化”。不过，本部分的主要结果是：
任意 $(A,I)$ 上的 rig-\,'etale 代数 $B$ 都可以代数化（不要求 $A$ 是 G-环），见引理
\ref{restricted-lemma-approximate}。关于这类代数化的文献指引，见备注
\ref{restricted-remark-discussion}。本部分的一般参考文献为
\cite{EGA}、\cite{Abbes} 和 \cite{Fujiwara-Kato}。

\medskip\noindent
本章第二部分（第
\ref{restricted-section-finite-type-red}--\ref{restricted-section-formal-modifications} 节）在适当的一般性层面
讨论形式代数空间态射的类型（主要针对局部诺特形式代数空间）。其中最有趣的是最后一节
的“形式修改”概念。我们仔细核对我们的定义与 \cite{ArtinII} 中 Artin 的定义一致。

\medskip\noindent
最后，在本章第三也是最后一部分（第
\ref{restricted-section-completion-and-morphisms}--\ref{restricted-section-modifications} 节）中，我们证明主要定理
并给出若干应用。事实上，我们从更强的结果——定理
\ref{restricted-theorem-dilatations-general}——推出 Artin 定理。粗略地说，该定理断言：若
$f:\mathfrak X\to\mathfrak X'$ 是 rig-\,'etale 态射，且 $\mathfrak X'$ 是局部诺特代数空间
的形式完备化，则 $\mathfrak X$ 也是如此。在 Artin 的工作中，态射 $f$ 被假设为适当且
rig-满射。






\section{两个范畴}
\label{restricted-section-two-categories}

\noindent
设 $A$ 为环，$I\subset A$ 为理想。本节中 ${}^\wedge$ 表示 $I$-进完备化。
置 $A_n=A/I^n$，于是 $A$ 的 $I$-进完备化为 $A^\wedge=\lim A_n$。令
$\mathcal{C}$ 为范畴
\begin{equation}
\label{restricted-equation-C}
\mathcal{C} =
\left\{
\begin{matrix}
\text{逆系统 }\ldots\to B_3\to B_2\to B_1\\
\text{其中 }B_n\text{ 是有限型 }A_n\text{-代数，}\\
B_{n+1}\to B_n\text{ 是 }A_{n+1}\text{-代数映射，}\\
\text{并诱导 }B_{n+1}/I^nB_{n+1}\cong B_n
\end{matrix}
\right\}
\end{equation}
$\mathcal{C}$ 中的态射由同态系统给出。

令 $\mathcal{C}'$ 为范畴
\begin{equation}
\label{restricted-equation-C-prime}
\mathcal{C}' =
\left\{
\begin{matrix}
A\text{-代数 }B\text{，其 }I\text{-进完备}\\
\text{满足 }B/IB\text{ 为 }A/I\text{ 上的有限型代数}
\end{matrix}
\right\}
\end{equation}
$\mathcal{C}'$ 中的态射为 $A$-代数映射。存在函子
\begin{equation}
\label{restricted-equation-from-complete-to-systems}
\mathcal{C}' \longrightarrow \mathcal{C},\quad
B \longmapsto (B/I^nB)
\end{equation}
确实，由于 $B/IB$ 是 $A/I$ 上的有限型代数，环映射
$A_n=A/I^n\to B/I^nB$ 由代数引理
\ref{algebra-lemma-finite-type-mod-nilpotent} 是有限型的。

\begin{lemma}
\label{restricted-lemma-topologically-finite-type}
设 $A$ 为环，$I\subset A$ 为有限生成理想。函子
$$
\mathcal{C} \longrightarrow \mathcal{C}',\quad
(B_n) \longmapsto B = \lim B_n
$$
是 (\ref{restricted-equation-from-complete-to-systems}) 的准逆函子。
$A[x_1,\ldots,x_r]^\wedge$ 的完备化属于 $\mathcal{C}'$，且
$\mathcal{C}'$ 的任意对象都具有形式
$$
B=A[x_1,\ldots,x_r]^\wedge/J
$$
其中 $J\subset A[x_1,\ldots,x_r]^\wedge$ 为某个理想。
\end{lemma}

\begin{proof}
设 $(B_n)$ 是 $\mathcal{C}$ 的对象。由代数引理
\ref{algebra-lemma-limit-complete}，可知 $B=\lim B_n$ 是 $I$-进完备的，且
$B/I^nB=B_n$。因此 $B$ 是 $\mathcal{C}'$ 的对象，并且取商可恢复对象 $(B_n)$。
反之，若 $B$ 是 $\mathcal{C}'$ 的对象，则由假设
$B=\lim B/I^nB$。故 $B\mapsto(B/I^nB)$ 是本引理所述函子的准逆。

\medskip\noindent
由于 $A[x_1,\ldots,x_r]^\wedge=\lim A_n[x_1,\ldots,x_r]$，
由引理第一条它是 $\mathcal{C}'$ 的对象。最后，设 $B$ 是 $\mathcal{C}'$ 的对象。
选取 $b_1,\ldots,b_r\in B$，其在 $B/IB$ 中的像生成 $B/IB$ 作为 $A/I$-代数。
因为 $B$ 是 $I$-进完备的，$A$-代数映射
$A[x_1,\ldots,x_r]\to B$，$x_i\mapsto b_i$ 延拓为 $A$-代数映射
$A[x_1,\ldots,x_r]^\wedge\to B$。还需证明此映射满射；
这由代数引理 \ref{algebra-lemma-completion-generalities} 得出，因为映射
$A[x_1,\ldots,x_r]\to B$ 模 $I$ 满射，
且 $B=B^\wedge$。
\end{proof}

\noindent
提醒读者：当 $A$ 不是诺特环时，$\mathcal{C}'$ 对象的商未必仍是 $\mathcal{C}'$ 的对象。
见例子之引理 \ref{examples-lemma-noncomplete-quotient}。下面说明当 $A$ 为诺特环时不会发生这种情况。

\begin{lemma}
\label{restricted-lemma-topologically-finite-type-Noetherian}
\begin{reference}
\cite[Proposition 7.5.5]{EGA1}
\end{reference}
设 $A$ 为诺特环，$I\subset A$ 为理想。则
\begin{enumerate}
\item 范畴 $\mathcal{C}'$（\ref{restricted-equation-C-prime}）的每个对象都是诺特的；
\item 若 $B\in\Ob(\mathcal{C}')$ 且 $J\subset B$ 为理想，则 $B/J$ 是 $\mathcal{C}'$ 的对象；
\item 对有限型 $A$-代数 $C$，其 $I$-进完备化 $C^\wedge$ 属于 $\mathcal{C}'$；
\item 特别地，完备化 $A[x_1,\ldots,x_r]^\wedge$ 属于 $\mathcal{C}'$。
\end{enumerate}
\end{lemma}

\begin{proof}
(4) 由代数引理 \ref{algebra-lemma-completion-Noetherian-Noetherian} 得出，
因为 $A[x_1,\ldots,x_r]$ 是诺特环（代数引理
\ref{algebra-lemma-Noetherian-permanence}）。由引理
\ref{restricted-lemma-topologically-finite-type}，为证明 (1) 只需化归到刚刚证明的多项式环完备化情形。
(2) 由代数引理 \ref{algebra-lemma-completion-tensor} 得出，该引理说明每个有限
$B$-模都是 $IB$-进完备的。(3) 与 (4) 同理。
\end{proof}

\begin{remark}[基变换]
\label{restricted-remark-base-change}
设 $\varphi:A_1\to A_2$ 为环映射，$I_i\subset A_i$ 为理想，并且对某个 $c\geq1$ 有
$\varphi(I_1^c)\subset I_2$。这对所有 $n\geq1$ 诱导环映射
$A_{1,cn}=A_1/I_1^{cn}\to A_2/I_2^n=A_{2,n}$。令 $\mathcal{C}_i$ 为
$(A_i,I_i)$ 的范畴 (\ref{restricted-equation-C})。存在基变换函子
\begin{equation}
\label{restricted-equation-base-change-systems}
\mathcal{C}_1 \longrightarrow \mathcal{C}_2,\quad
(B_n) \longmapsto (B_{cn} \otimes_{A_{1,cn}} A_{2,n})
\end{equation}
令 $\mathcal{C}_i'$ 为 $(A_i,I_i)$ 的范畴 (\ref{restricted-equation-C-prime})。
若 $I_2$ 有限生成，则存在基变换函子
\begin{equation}
\label{restricted-equation-base-change-complete}
\mathcal{C}_1' \longrightarrow \mathcal{C}_2',\quad
B \longmapsto (B \otimes_{A_1} A_2)^\wedge
\end{equation}
因为此时完备化是完备的（代数引理
\ref{algebra-lemma-hathat-finitely-generated}）。若 $I_1,I_2$ 都有限生成，
则通过函子 (\ref{restricted-equation-from-complete-to-systems}) 两个基变换函子一致；
由引理 \ref{restricted-lemma-topologically-finite-type} 这些函子是等价。
\end{remark}

\begin{remark}[闭浸入下的基变换]
\label{restricted-remark-take-bar}
设 $A$ 为诺特环，$I\subset A$ 为理想。令 $\mathfrak a\subset A$ 为理想，
记 $\bar A=A/\mathfrak a$。令 $\bar I\subset\bar A$ 为理想，并对某些 $c,d\geq1$ 满足
$I^c\bar A\subset\bar I$ 且 $\bar I^d\subset I\bar A$。此时，$(A,I)$ 到
$(\bar A,\bar I)$ 的基变换函子 (\ref{restricted-equation-base-change-complete}) 由
$B\mapsto\bar B=B/\mathfrak aB$ 给出。事实上有
\begin{equation}
\label{restricted-equation-base-change-to-closed}
\bar B=(B\otimes_A\bar A)^\wedge=(B/\mathfrak aB)^\wedge
=B/\mathfrak aB
\end{equation}
最后一个等式是因为任意有限 $B$-模由代数引理
\ref{algebra-lemma-completion-tensor} 为 $I$-进完备；若再被 $\mathfrak a$ 湮灭，
则由代数引理 \ref{algebra-lemma-change-ideal-completion} 为 $\bar I$-进完备。
\end{remark}





\section{朴素余切复形}
\label{restricted-section-naive-cotangent-complex}

\noindent
设 $A$ 为诺特环，$I\subset A$ 为理想。设 $B$ 为 $A$-代数，
它是 $I$-进完备的，且 $A/I\to B/IB$ 为有限型映射，即为
(\ref{restricted-equation-C-prime}) 的对象。由引理
\ref{restricted-lemma-topologically-finite-type-Noetherian}，可写成
$$
B=A[x_1,\ldots,x_r]^\wedge/J
$$
其中 $J$ 为有限生成理想。对上述表示的一个选择，我们在此情形下用公式
\begin{equation}
\label{restricted-equation-NL}
\NL_{B/A}^\wedge=(J/J^2\longrightarrow\bigoplus B\text{d}x_i)
\end{equation}
定义\textit{朴素余切复形}；其项位于次数 $-1$ 和 $0$，其中该映射把
$g\in J$ 的剩余类送到微分
$\text{d}g=\sum(\partial g/\partial x_i)\text{d}x_i$。这里的偏导数把 $g$
看作幂级数来取。下面的引理说明 $\NL_{B/A}^\wedge$ 在同伦意义下定义良好。

\begin{lemma}
\label{restricted-lemma-NL-up-to-homotopy}
设 $A$ 为诺特环，$I\subset A$ 为理想。设 $B$ 是
(\ref{restricted-equation-C-prime}) 的对象。朴素余切复形 $\NL_{B/A}^\wedge$ 在 $K(B)$ 中定义良好。
\end{lemma}

\begin{proof}
该引理的含义是：给定第二个表示
$B=A[y_1,\ldots,y_s]^\wedge/K$，$B$-模复形
$$
(J/J^2\to B\text{d}x_i)
\quad\text{和}\quad
(K/K^2\to\bigoplus B\text{d}y_j)
$$
是同伦等价的。证明可完全仿照代数引理
\ref{algebra-lemma-NL-homotopy} 的证明。

\medskip\noindent
步骤 1。若选取 $g_i(y_1,\ldots,y_s)\in A[y_1,\ldots,y_s]^\wedge$，
其映到 $B$ 中 $x_i$ 的像，则得到唯一的连续 $A$-代数同态
$$
A[x_1,\ldots,x_r]^\wedge\to A[y_1,\ldots,y_s]^\wedge,\quad
x_i\mapsto g_i(y_1,\ldots,y_s)
$$
并与给定的到 $B$ 的满射相容。这样的映射称为表示之间的态射。
它把 $J$ 映入 $K$，从而诱导 $B$-模映射 $J/J^2\to K/K^2$。
把 $\text{d}x_i$ 送到 $\sum(\partial g_i/\partial y_j)\text{d}y_j$，
得到复形映射
$$
(J/J^2\to\bigoplus B\text{d}x_i)
\longrightarrow
(K/K^2\to\bigoplus B\text{d}y_j).
$$
当然，交换两个表示的角色也可作同样构造，从而得到反向的复形映射。

\medskip\noindent
步骤 2。上述构造与表示态射的复合相容。因此只需证明：给定
$g_i(x_1,\ldots,x_r)\in A[x_1,\ldots,x_n]^\wedge$，其映到 $B$ 中 $x_i$ 的像，
所诱导的复形映射
$$
(J/J^2\to\bigoplus B\text{d}x_i)
\longrightarrow
(J/J^2\to\bigoplus B\text{d}x_i)
$$
同伦于恒等映射。为此考虑映射
$h:\bigoplus B\text{d}x_i\to J/J^2$，其规则为
$\text{d}x_i\mapsto g_i(x_1,\ldots,x_n)-x_i$，直接计算即可。
\end{proof}

\begin{lemma}
\label{restricted-lemma-NL-is-completion}
设 $A$ 为诺特环，$I\subset A$ 为理想。设 $A\to B$ 为有限型环映射。
选取表示 $\alpha:A[x_1,\ldots,x_n]\to B$。则
$$
\NL_{B^\wedge/A}^\wedge=\lim\NL(\alpha)\otimes_B B^\wedge
$$
作为复形成立，并且在 $D(B^\wedge)$ 中有
$$
\NL_{B^\wedge/A}^\wedge=\NL_{B/A}\otimes_B^\mathbf{L}B^\wedge.
$$
\end{lemma}

\begin{proof}
所述命题有意义，因为 $B^\wedge$ 是
(\ref{restricted-equation-C-prime}) 的对象，这是引理
\ref{restricted-lemma-topologically-finite-type-Noetherian} 的结论。
令 $J=\Ker(\alpha)$。在 $A[x_1,\ldots,x_n]$ 上对有限模取 $I$-进完备化是正合的，
并且等于函子
$M\mapsto M\otimes_{A[x_1,\ldots,x_n]}A[x_1,\ldots,x_n]^\wedge$，
见代数引理 \ref{algebra-lemma-completion-tensor} 和
\ref{algebra-lemma-completion-flat}。此外，环映射
$A[x_1,\ldots,x_n]\to A[x_1,\ldots,x_n]^\wedge$ 及 $B\to B^\wedge$ 都是平坦的。
因此 $B^\wedge=A[x_1,\ldots,x_n]^\wedge/J^\wedge$，且
$$
(J/J^2)\otimes_BB^\wedge=(J/J^2)^\wedge=J^\wedge/(J^\wedge)^2.
$$
由于 $\NL(\alpha)=(J/J^2\to\bigoplus B\text{d}x_i)$，见代数之
\ref{algebra-section-netherlander} 节，故 $\NL_{B^\wedge/A}^\wedge$
等于 $\NL(\alpha)\otimes_BB^\wedge$。最后一个断言来自
$\NL_{B/A}$ 与 $\NL(\alpha)$ 同伦等价，且 $B\to B^\wedge$ 平坦（所以沿
$B\to B^\wedge$ 的导出基变换
就是普通基变换）。
\end{proof}

\begin{lemma}
\label{restricted-lemma-NL-is-limit}
设 $A$ 为诺特环，$I\subset A$ 为理想。设 $B$ 是
(\ref{restricted-equation-C-prime}) 的对象。则
\begin{enumerate}
\item 在 $D(B)$ 中，pro-对象
$\{\NL_{B/A}^\wedge\otimes_BB/I^nB\}$ 与 $\{\NL_{B_n/A_n}\}$
严格同构（详见证明）；
\item 在 $D(B)$ 中，$\NL_{B/A}^\wedge=R\lim\NL_{B_n/A_n}$。
\end{enumerate}
这里 $B_n,A_n$ 如第 \ref{restricted-section-two-categories} 节所定义。
\end{lemma}

\begin{proof}
命题的含义如下：对每个 $n$，有良定义的 $B_n$-模复形
$\NL_{B_n/A_n}$，并有过渡映射
$\NL_{B_{n+1}/A_{n+1}}\to\NL_{B_n/A_n}$。见代数之
\ref{algebra-section-netherlander} 节。因此可将
$$
\ldots\to\NL_{B_3/A_3}\to\NL_{B_2/A_2}\to\NL_{B_1/A_1}
$$
视为 $B$-模复形的逆系统，因而也是 $D(B)$ 中的逆系统。此外，
$R\lim\NL_{B_n/A_n}$ 是该逆系统的同伦极限，见导出范畴之
\ref{derived-section-derived-limit} 节。

\medskip\noindent
选取表示 $B=A[x_1,\ldots,x_r]^\wedge/J$。这给出表示
$$
B_n=B/I^nB=A_n[x_1,\ldots,x_r]/J_n
$$
其中
$$
J_n=JA_n[x_1,\ldots,x_r]
=J/(J\cap I^nA[x_1,\ldots,x_r]^\wedge).
$$
两项复形 $J_n/J_n^2\longrightarrow\bigoplus B_n\text{d}x_i$
表示 $\NL_{B_n/A_n}$，见代数之 \ref{algebra-section-netherlander} 节。
由诺特环 $A[x_1,\ldots,x_r]^\wedge$ 中的 Artin--Rees 定理
（代数引理 \ref{algebra-lemma-Artin-Rees}，以及引理
\ref{restricted-lemma-topologically-finite-type-Noetherian}），存在 $c\geq0$ 使得对所有 $n\geq c$
有典范满射
$$
J/I^nJ\to J_n\to J/I^{n-c}J\to J_{n-c},\quad n\geq c
$$
略作思考可知这些映射与微分相容，从而得到复形映射
$$
\NL_{B/A}^\wedge\otimes_BB/I^nB\to
\NL_{B_n/A_n}\to
\NL_{B/A}^\wedge\otimes_BB/I^{n-c}B\to
\NL_{B_{n-c}/A_{n-c}}
$$
并与逆系统 $\{\NL_{B/A}^\wedge\otimes_BB/I^nB\}$ 和
$\{\NL_{B_n/A_n}\}$ 的过渡映射相容。这就证明了引理的 (1)。

\medskip\noindent
由 (1)，且同构的 pro-系统具有相同的 $R\lim$，为证明 (2) 只需说明
$\NL_{B/A}^\wedge$ 等于 $R\lim\NL_{B/A}^\wedge\otimes_BB/I^nB$。
然而 $\NL_{B/A}^\wedge$ 是有限 $B$-模组成的两项复形 $M^\bullet$；
例如由代数引理 \ref{algebra-lemma-completion-tensor}，这些模是 $I$-进完备的。
因此
$M^\bullet=\lim M^\bullet/I^nM^\bullet=R\lim M^\bullet/I^nM^\bullet$，
见更多代数之引理 \ref{more-algebra-lemma-compute-Rlim-modules} 及备注
\ref{more-algebra-remark-how-unique}。
\end{proof}


\begin{lemma}
\label{restricted-lemma-NL-base-change}
设 $(A_1,I_1)\to(A_2,I_2)$ 如备注 \ref{restricted-remark-base-change}，
且 $A_1,A_2$ 为诺特环。设 $B_1$ 是 $(A_1,I_1)$ 的
(\ref{restricted-equation-C-prime}) 对象，令 $B_2$ 为 $B_1$ 的基变换。则存在典范映射
$$
\NL_{B_1/A_1}\otimes_{B_2}B_1\to\NL_{B_2/A_2}
$$
其在 $H^0$ 上诱导同构，在 $H^{-1}$ 上诱导满射。
\end{lemma}

\begin{proof}
选取表示 $B_1=A_1[x_1,\ldots,x_r]^\wedge/J_1$。由于
$$
A_2/I_2^n[x_1,\ldots,x_r]
=A_1/I_1^{cn}[x_1,\ldots,x_r]\otimes_{A_1/I_1^{cn}}A_2/I_2^n,
$$
有
$$
A_2[x_1,\ldots,x_r]^\wedge
=(A_1[x_1,\ldots,x_r]^\wedge\otimes_{A_1}A_2)^\wedge,
$$
两边都取 $I_2$-进完备化（当然，对 $A_1[x_1,\ldots,x_r]^\wedge$ 取的是
$I_1$-进完备化）。置 $J_2=J_1A_2[x_1,\ldots,x_r]^\wedge$。
同样论证得到 $B_2$ 在 $A_2$ 上的表示
\begin{align*}
B_2&=(B_1\otimes_{A_1}A_2)^\wedge\\
&=\lim\frac{A_1/I_1^{cn}[x_1,\ldots,x_r]}{J_1(A_1/I_1^{cn}[x_1,\ldots,x_r])}
\otimes_{A_1/I_1^{cn}}A_2/I_2^n\\
&=\lim\frac{A_2/I_2^n[x_1,\ldots,x_r]}{J_2(A_2/I_2^n[x_1,\ldots,x_r])}\\
&=A_2[x_1,\ldots,x_r]^\wedge/J_2.
\end{align*}
因此有交换图
$$
\xymatrix{
\NL^\wedge_{B_1/A_1} : \ar[d] &
J_1/J_1^2 \ar[r]_-{\text{d}} \ar[d] & \bigoplus B_1\text{d}x_i \ar[d] \\
\NL^\wedge_{B_2/A_2} : &
J_2/J_2^2 \ar[r]^-{\text{d}} & \bigoplus B_2\text{d}x_i
}
$$
诱导映射 $J_1/J_1^2\otimes_{B_1}B_2\to J_2/J_2^2$ 是满射，
因为 $J_2$ 由 $J_1$ 的像生成。这确定了引理中显示的箭头。
我们略去证明该箭头在同伦意义下良定义（即在同伦意义下与表示选择无关）。
关于诱导的上同调模上的映射的断言直接由上述讨论得到（细节略）。
\end{proof}

\begin{lemma}
\label{restricted-lemma-exact-sequence-NL}
设 $A$ 为诺特环，$I\subset A$ 为理想。设 $B\to C$ 为
(\ref{restricted-equation-C-prime}) 的态射。则存在正合序列
$$
\xymatrix{
C\otimes_BH^0(\NL_{B/A}^\wedge)\ar[r] &
H^0(\NL_{C/A}^\wedge)\ar[r] &
H^0(\NL_{C/B}^\wedge)\ar[r] & 0 \\
H^{-1}(\NL_{B/A}^\wedge\otimes_BC)\ar[r] &
H^{-1}(\NL_{C/A}^\wedge)\ar[r] &
H^{-1}(\NL_{C/B}^\wedge)\ar[llu]
}
$$
详见证明。
\end{lemma}

\begin{proof}
注意，取张量积 $\NL_{B/A}^\wedge\otimes_BC$ 是有意义的，因为由引理
\ref{restricted-lemma-NL-up-to-homotopy}，$\NL_{B/A}^\wedge$ 在同伦意义下定义良好。
此外，$(B,IB)$ 是一对，其中 $B$ 为诺特环（引理
\ref{restricted-lemma-topologically-finite-type-Noetherian}），而 $C$ 属于相应的范畴
(\ref{restricted-equation-C-prime})。因此 $6$ 项序列中的所有项都定义良好。

\medskip\noindent
选取表示 $B=A[x_1,\ldots,x_r]^\wedge/J$，再选取表示
$C=B[y_1,\ldots,y_s]^\wedge/J'$。合并这些表示得到
$$
C=A[x_1,\ldots,x_r,y_1,\ldots,y_s]^\wedge/K.
$$
于是读者可验证得到交换图
$$
\xymatrix{
0\ar[r] &
\bigoplus C\text{d}x_i\ar[r] &
\bigoplus C\text{d}x_i\oplus\bigoplus C\text{d}y_j\ar[r] &
\bigoplus C\text{d}y_j\ar[r] & 0\\
&
J/J^2\otimes_BC\ar[r]\ar[u] &
K/K^2\ar[r]\ar[u] &
J'/(J')^2\ar[r]\ar[u] &
0
}
$$
其两行均正合。注意，左侧竖直箭头是定义
$\NL_{B/A}^\wedge$ 的箭头与 $\text{id}_C$ 的张量积。
引理由蛇引理（代数引理 \ref{algebra-lemma-snake}）得到。
\end{proof}

\begin{lemma}
\label{restricted-lemma-transitive-lci-at-end}
在引理 \ref{restricted-lemma-exact-sequence-NL} 的假设下，假设对所有 $n$，
$B/I^nB\to C/I^nC$ 是局部完全交叉同态。则
$$
H^{-1}(\NL_{B/A}^\wedge\otimes_BC)\to H^{-1}(\NL_{C/A}^\wedge)
$$
是单射。
\end{lemma}

\begin{proof}
对每个 $n\geq1$，置 $A_n=A/I^n$、$B_n=B/I^nB$、$C_n=C/I^nC$。
有
\begin{align*}
H^{-1}(\NL_{B/A}^\wedge\otimes_BC)
&=\lim H^{-1}(\NL_{B/A}^\wedge\otimes_BC_n)\\
&=\lim H^{-1}(\NL_{B/A}^\wedge\otimes_BB_n\otimes_{B_n}C_n)\\
&=\lim H^{-1}(\NL_{B_n/A_n}\otimes_{B_n}C_n).
\end{align*}
第一个等式由更多代数之引理
\ref{more-algebra-lemma-consequence-Artin-Rees} 以及
$H^{-1}(\NL_{B/A}^\wedge\otimes_BC)$ 是有限 $C$-模这一事实得出，
故由代数引理 \ref{algebra-lemma-completion-tensor} 它是 $I$-进完备的。
第二个等式显然成立，第三个由引理 \ref{restricted-lemma-NL-is-limit} 得出。
映射
$H^{-1}(\NL_{B_n/A_n}\otimes_{B_n}C_n)\to H^{-1}(\NL_{C_n/A_n})$
由更多代数之引理 \ref{more-algebra-lemma-transitive-lci-at-end} 为单射。
证明完成，因为同样有
$H^{-1}(\NL_{C/A}^\wedge)=\lim H^{-1}(\NL_{C_n/A_n})$。
\end{proof}





\section{rig-光滑代数}
\label{restricted-section-rig-smooth}

\noindent
作为下述定义的动机，请参阅更多代数之备注
\ref{more-algebra-remark-smoothness-ext-1-zero}。

\begin{definition}
\label{restricted-definition-rig-smooth-homomorphism}
设 $A$ 为诺特环，$I\subset A$ 为理想。设 $B$ 是
(\ref{restricted-equation-C-prime}) 的对象。若存在整数 $c\geq0$，使得对每个
$B$-模 $N$，$I^c$ 湮灭 $\Ext^1_B(\NL_{B/A}^\wedge,N)$，
则称 $B$ 在 $(A,I)$ 上\textit{rig-光滑}。
\end{definition}

\noindent
下面说明这一条件的含义。

\begin{lemma}
\label{restricted-lemma-equivalent-with-artin-smooth}
设 $A$ 为诺特环，$I\subset A$ 为理想。设 $B$ 是
(\ref{restricted-equation-C-prime}) 的对象。写成
$B=A[x_1,\ldots,x_r]^\wedge/J$（引理
\ref{restricted-lemma-topologically-finite-type-Noetherian}），并令
$\NL_{B/A}^\wedge=(J/J^2\to\bigoplus B\text{d}x_i)$
为其朴素余切复形（\ref{restricted-equation-NL}）。以下条件等价：
\begin{enumerate}
\item $B$ 在 $(A,I)$ 上 rig-光滑；
\item $D(B)$ 中的对象 $\NL_{B/A}^\wedge$ 关于理想 $IB$ 满足更多代数之引理
\ref{more-algebra-lemma-ext-1-annihilated} 的等价条件 (1)--(6)；
\item 存在 $c\geq0$，使得对每个 $a\in I^c$，存在映射
$h:\bigoplus B\text{d}x_i\to J/J^2$，满足
$a:J/J^2\to J/J^2$ 等于 $h\circ\text{d}$；
\item 存在 $b_1,\ldots,b_s\in B$，使
$V(b_1,\ldots,b_s)\subset V(IB)$，并且对每个 $l=1,\ldots,s$，
存在 $m\geq0$、$f_1,\ldots,f_m\in J$ 及
$T\subset\{1,\ldots,n\}$ 且 $|T|=m$，使
\begin{enumerate}
\item $\det_{i\in T,j\leq m}(\partial f_j/\partial x_i)$ 在 $B$ 中整除 $b_l$；
\item $b_lJ\subset(f_1,\ldots,f_m)+J^2$。
\end{enumerate}
\end{enumerate}
\end{lemma}

\begin{proof}
(1)、(2)、(3) 的等价性立即由更多代数之引理
\ref{more-algebra-lemma-ext-1-annihilated} 得出。

\medskip\noindent
假设 $b_1,\ldots,b_s$ 如 (4) 所述。由于 $B$ 为诺特环，
$V(b_1,\ldots,b_s)\subset V(IB)$ 蕴含对某个 $c\geq0$ 有
$I^cB\subset(b_1,\ldots,b_s)$（例如由代数引理
\ref{algebra-lemma-Noetherian-power-ideal-kills-module}）。取
$1\leq l\leq s$、$m\geq0$、$f_1,\ldots,f_m\in J$ 及
$T\subset\{1,\ldots,n\}$ 且 $|T|=m$，满足 (4)(a)、(b)。在 $b_l$ 处取逆可得
$$
\NL_{B/A}^\wedge\otimes_BB_{b_l}=
\left(
\bigoplus\nolimits_{j\leq m}B_{b_l}f_j
\longrightarrow
\bigoplus\nolimits_{i=1,\ldots,n}B_{b_l}\text{d}x_i
\right)
$$
且该箭头同构于直和项
$\bigoplus_{i\in T}B_{b_l}\text{d}x_i$ 的嵌入。
因此 $H^{-1}(\NL_{B/A}^\wedge)$ 是 $b_l$-幂扭的，而
$H^0(\NL_{B/A}^\wedge)$ 在 $b_l$ 处取逆后成为有限自由模。
结合 $I^cB\subset(b_1,\ldots,b_s)$，得
$H^{-1}(\NL_{B/A}^\wedge)$ 是 $IB$-幂扭的。于是更多代数之引理
\ref{more-algebra-lemma-ext-1-annihilated} 的条件 (4) 成立，
从而 (4) 蕴含 (2)。

\medskip\noindent
现在假设等价条件 (1)、(2)、(3) 成立。我们将证明 (4)；但强烈建议读者自行核实。
复形 $\NL_{B/A}^\wedge$ 决定
$D^b_{\textit{Coh}}(\Spec(B))$ 的一个对象，其限制到 Zariski 开集
$U=\Spec(B)\setminus V(IB)$ 是置于次数 $0$ 的有限局部自由模
$\mathcal{E}$（例如由更多代数之引理
\ref{more-algebra-lemma-ext-1-annihilated} 的第四个等价条件可得）。
取 $J$ 的生成元 $f_1,\ldots,f_M$，得到正合序列
$$
\bigoplus\nolimits_{j=1,\ldots,M}\mathcal{O}_U\cdot f_j\to
\bigoplus\nolimits_{i=1,\ldots,n}\mathcal{O}_U\cdot\text{d}x_i\to
\mathcal{E}\to0.
$$
令 $U=\bigcup_{l=1,\ldots,s}U_l$ 为有限仿射开覆盖，使
$\mathcal{E}|_{U_l}$ 为秩 $r_l=n-m_l$ 的自由模，其中 $n\geq m_l\geq0$。
将各 $U_l$ 细化为仿射开覆盖后，可设存在
$T_l\subset\{1,\ldots,n\}$，使 $i\in\{1,\ldots,n\}\setminus T_l$ 的
$\text{d}x_i$ 映成 $\mathcal{E}|_{U_l}$ 的一组基。重复论证可设存在
$T'_l\subset\{1,\ldots,M\}$，其基数为 $m_l$，且 $j\in T'_l$ 的
$f_j$ 映成 $\mathcal{O}_{U_l}\cdot\text{d}x_i\to\mathcal{E}|_{U_l}$ 的核的一组基。
最后，由于开覆盖 $U=\bigcup U_l$ 可由标准开覆盖细化（代数引理
\ref{algebra-lemma-Zariski-topology}），可设 $U_l=D(g_l)$，其中 $g_l\in B$。
特别地 $V(g_1,\ldots,g_s)=V(IB)$。
由上述选择和线性代数，$\det_{i\in T_l,j\in T'_l}
(\partial f_j/\partial x_i)$ 在 $B_{b_l}$ 中为可逆元。
同样，$\NL_{B/A}^\wedge$ 在 $U_l$ 上次数 $-1$ 的上同调消失，
说明 $J/J^2+(f_j;j\in T')$ 被 $b_l$ 的某个幂湮灭。
将每个 $g_l$ 换成适当幂，即得引理的 (4)(a)、(4)(b)。
细节略。
\end{proof}

\begin{lemma}
\label{restricted-lemma-rig-smooth}
设 $A$ 为诺特环，$I$ 为理想，$B$ 为有限型 $A$-代数。
\begin{enumerate}
\item 若 $\Spec(B)\to\Spec(A)$ 在 $\Spec(A)\setminus V(I)$ 上光滑，
则 $B^\wedge$ 在 $(A,I)$ 上 rig-光滑；
\item 若 $B^\wedge$ 在 $(A,I)$ 上 rig-光滑，则存在
$g\in1+IB$，使 $\Spec(B_g)$ 在 $\Spec(A)\setminus V(I)$ 上光滑。
\end{enumerate}
\end{lemma}

\begin{proof}
以下不再重复提及引理 \ref{restricted-lemma-equivalent-with-artin-smooth}。

\medskip\noindent
假设 (1)。回忆 $\NL_{B/A}$ 与局部化交换，见代数引理
\ref{algebra-lemma-localize-NL}。因此由光滑环映射的定义（朴素余切复形拟同构于
置于次数 $0$ 的有限射影模），$\NL_{B/A}$ 关于理想 $IB$ 满足更多代数之引理
\ref{more-algebra-lemma-ext-1-annihilated} 的第四个等价条件（细节略）。
由引理 \ref{restricted-lemma-NL-is-completion}，
$\NL_{B^\wedge/A}^\wedge=\NL_{B/A}\otimes_BB^\wedge$，再由更多代数之引理
\ref{more-algebra-lemma-base-change-property-ext-1-annihilated} 得到 (2)。

\medskip\noindent
假设 (2)。选取表示 $B=A[x_1,\ldots,x_n]/J$，置 $N=J/J^2$，
并考虑由 $J/J^2$ 上恒等映射确定的元素
$\xi\in\Ext^1_B(\NL_{B/A},J/J^2)$。再次利用
$\NL_{B^\wedge/A}^\wedge=\NL_{B/A}\otimes_BB^\wedge$，可知假设蕴含
$$
\xi\otimes1\in
\Ext^1_{B^\wedge}(\NL_{B/A}\otimes_BB^\wedge,N\otimes_BB^\wedge)
=\Ext^1_{B^\wedge}(\NL_{B/A},N)\otimes_BB^\wedge
$$
被某个 $c\geq0$ 的 $I^c$ 湮灭。等式例如由更多代数之引理
\ref{more-algebra-lemma-base-change-RHom} 得出（也可由更简单的更多代数之引理
\ref{more-algebra-lemma-pseudo-coherence-and-base-change-ext} 推出）。
于是 $M=I^cB\xi\subset\Ext^1_B(\NL_{B/A},N)$ 是有限子模，并映到
$\Ext^1_B(\NL_{B/A},N)\otimes_BB^\wedge$ 的零元。由于 $B\to B^\wedge$ 平坦，
这意味着 $M\otimes_BB^\wedge$ 为零。由 Nakayama 引理（代数引理
\ref{algebra-lemma-NAK}），$M=I^cB\xi$ 被某个形如 $g=1+x$（$x\in IB$）的元素湮灭。
这说明对每个 $b\in I^cB$，存在使下图交换的 $B$-线性虚线箭头
$$
\xymatrix{
J/J^2 \ar[r] \ar[d]^b & \bigoplus B\text{d}x_i \ar@{..>}[d]^h \\
J/J^2 \ar[r] & (J/J^2)_g
}
$$
因此 $(\NL_{B/A})_{gb}$ 拟同构于有限射影模；细节略。
由于 $(\NL_{B/A})_{gb}=\NL_{B_{gb}/A}$ 在 $D(B_{gb})$ 中成立，
这说明 $B_{gb}$ 在 $\Spec(A)$ 上光滑。对所有 $b\in I^cB$ 均如此，
故 $\Spec(B_g)\to\Spec(A)$ 在 $\Spec(A)\setminus V(I)$ 上光滑。
\end{proof}

\begin{lemma}
\label{restricted-lemma-zero-ext-1-after-modding-out}
设 $(A_1,I_1)\to(A_2,I_2)$ 如备注 \ref{restricted-remark-base-change}，且
$A_1,A_2$ 为诺特环。设 $B_1$ 是 $(A_1,I_1)$ 的
(\ref{restricted-equation-C-prime}) 对象，$B_2$ 为 $B_1$ 的基变换。设
$f_1\in B_1$，其像为 $f_2\in B_2$。若对每个 $B_1$-模 $N_1$，
$\Ext^1_{B_1}(\NL_{B_1/A_1}^\wedge,N_1)$ 被 $f_1$ 湮灭，
则对每个 $B_2$-模 $N_2$，$\Ext^1_{B_2}(\NL_{B_2/A_2}^\wedge,N_2)$
被 $f_2$ 湮灭。
\end{lemma}

\begin{proof}
由引理 \ref{restricted-lemma-NL-base-change} 存在映射
$$
\NL_{B_1/A_1}\otimes_{B_2}B_1\to\NL_{B_2/A_2}
$$
其在 $H^0$ 上诱导同构、在 $H^{-1}$ 上诱导满射。
于是由更多代数之引理
\ref{more-algebra-lemma-two-term-base-change}、
\ref{more-algebra-lemma-base-change-property-ext-1-annihilated} 和
\ref{more-algebra-lemma-surjection-property-ext-1-annihilated} 得到结论；
后两个引理分别应用于主理想 $(f_1)\subset B_1$ 与 $(f_2)\subset B_2$。
\end{proof}

\begin{lemma}
\label{restricted-lemma-base-change-rig-smooth-homomorphism}
设 $A_1\to A_2$ 为诺特环映射。设 $I_i\subset A_i$ 为理想，且
$V(I_1A_2)=V(I_2)$。设 $B_1$ 是 $(A_1,I_1)$ 的
(\ref{restricted-equation-C-prime}) 对象，$B_2$ 是按备注 \ref{restricted-remark-base-change}
得到的 $B_1$ 的基变换。若 $B_1$ 在 $(A_1,I_1)$ 上 rig-光滑，
则 $B_2$ 在 $(A_2,I_2)$ 上 rig-光滑。
\end{lemma}

\begin{proof}
由引理 \ref{restricted-lemma-zero-ext-1-after-modding-out}、定义
\ref{restricted-definition-rig-smooth-homomorphism} 以及 $A_2$ 诺特，从而对某个
$c\geq0$ 有 $I_2^c$ 包含于 $I_1A_2$，立即得到。
\end{proof}









\section{环同态的变形}
\label{restricted-section-defos-ring-maps}

\noindent
关于从 rig-光滑代数提升环同态的一些工作。

\begin{remark}[线性近似]
\label{restricted-remark-linear-approximation}
设 $A$ 为环，且 $I \subset A$ 为有限生成理想。
设 $C$ 为 $I$-进完备的 $A$-代数。
设 $\psi : A[x_1, \ldots, x_r]^\wedge \to C$ 为连续的
$A$-代数映射。给定 $\delta_i \in C$（$i = 1, \ldots, r$）。
则可以考虑
$$
\psi' : A[x_1, \ldots, x_r]^\wedge \to C,\quad
x_i \longmapsto \psi(x_i) + \delta_i
$$
参见形式空间，备注 \ref{formal-spaces-remark-universal-property}。
于是有
$$
\psi'(g) = \psi(g) + \sum \psi(\partial g/\partial x_i)\delta_i + \xi
$$
其中误差项 $\xi \in (\delta_i\delta_j)$。这是将 $g$ 写成幂级数并逐项计算得到的。
由于 $g$ 的系数趋于零，收敛性是自动的。
细节略。
\end{remark}

\begin{remark}[提升映射]
\label{restricted-remark-improve-homomorphism}
设 $A$ 为诺特环，且 $I \subset A$ 为理想。
设 $B$ 是 \ref{restricted-equation-C-prime} 中的对象。
设 $C$ 为 $I$-进完备的 $A$-代数。
设 $\psi_n : B \to C/I^nC$ 为 $A$-代数同态。
将 $\psi_n$ 提升为到 $C/I^{2n}C$ 中的 $A$-代数同态的障碍是元素
$$
o(\psi_n) \in \Ext^1_B(\NL_{B/A}^\wedge, I^nC/I^{2n}C)
$$
正如我们将要说明的。具体地，取一个表示
$B = A[x_1, \ldots, x_r]^\wedge/J$。
取 $\psi : A[x_1, \ldots, x_r]^\wedge \to C$ 为 $\psi_n$ 的提升。
由于 $\psi(J) \subset I^nC$，有 $\psi(J^2) \subset I^{2n}C$，
从而得到 $B$-线性同态
$$
o(\psi) :
J/J^2 \longrightarrow I^nC/I^{2n}C, \quad g \longmapsto \psi(g)
$$
它当然延拓为 $C$-线性映射
$$
J/J^2 \otimes_B C \to I^nC/I^{2n}C
$$
由于 $\NL_{B/A}^\wedge = (J/J^2 \to \bigoplus B \text{d}x_i)$，
通过如下同一性，得到 $o(\psi_n)$，它是 $o(\psi)$ 的像：
\begin{align*}
& \Ext^1_B(\NL_{B/A}^\wedge, I^nC/I^{2n}C) \\
& =
\Coker\left(\Hom_B(\bigoplus B\text{d}x_i, I^nC/I^{2n}C) \to
\Hom_B(J/J^2, I^nC/I^{2n}C)\right)
\end{align*}
等式参见更多代数，引理
\ref{more-algebra-lemma-map-out-of-almost-free} 的 (1)。

\medskip\noindent
假设 $o(\psi_n)$ 在
$\Ext^1_B(\NL_{B/A}^\wedge, I^{n'}C/I^{2n'}C)$ 中映为零，
其中整数 $n'$ 满足 $n > n' > n/2$。我们声称，这意味着可以找到
$A$-代数同态 $\psi'_{2n'} : B \to C/I^{2n'}C$，
使其作为映射到 $C/I^{n'}C$ 的映射时与 $\psi_n$ 一致。
极端情形 $n'=n$ 说明了我们先前为何称 $o(\psi_n)$ 为将 $\psi_n$
提升到 $C/I^{2n}C$ 的障碍。
证明如下：假设 $o(\psi_n)$ 映为零，意味着可以找到一个 $B$-模映射
$$
h : \bigoplus B\text{d}x_i \longrightarrow I^{n'}C/I^{2n'}C
$$
使得 $o(\psi)$ 与 $h \circ \text{d}$ 作为映射到
$I^{n'}C/I^{2n'}C$ 时相同。设
$h(\text{d}x_i) = \delta_i \bmod I^{2n'}C$，其中 $\delta_i \in I^{n'}C$。
于是考虑映射
$$
\psi' : A[x_1, \ldots, x_r]^\wedge \to C,\quad
x_i \longmapsto \psi(x_i) - \delta_i
$$
幂级数计算表明 $\psi'(J) \subset I^{2n'}C$。事实上，对 $g \in J$ 有
$$
\psi'(g) \equiv
\psi(g) - \sum \psi(\partial g/\partial x_i)\delta_i \equiv
o(\psi)(g) - (h \circ \text{d})(g) \equiv
0 \bmod I^{2n'}C
$$
第一次等式参见备注 \ref{restricted-remark-linear-approximation}。
因此 $\psi'$ 诱导出所需的 $A$-代数同态
$$
\psi'_{2n'} : B \to C/I^{2n'}C
$$
\end{remark}

\begin{lemma}
\label{restricted-lemma-get-morphism-general-better}
假设给定以下数据：
\begin{enumerate}
\item 一个整数 $c \geq 0$；
\item 诺特环 $A$ 的一个理想 $I$；
\item 对于 $(A, I)$，\ref{restricted-equation-C-prime} 中的 $B$，满足对任意 $B$-模 $N$，
$I^c$ 湮灭 $\Ext^1_B(\NL_{B/A}^\wedge, N)$；
\item 一个诺特的 $I$-进完备 $A$-代数 $C$；记
$d = d(\text{Gr}_I(C))$ 和 $q_0 = q(\text{Gr}_I(C))$ 为局部上同调中
\ref{local-cohomology-section-uniform} 节所得的整数；
\item 一个整数 $n \geq \max(q_0 + (d + 1)c, 2(d + 1)c + 1)$；以及
\item 一个 $A$-代数同态 $\psi_n : B \to C/I^nC$。
\end{enumerate}
则存在 $A$-代数映射 $\varphi : B \to C$，使得
$$
\psi_n \bmod I^{n - (d + 1)c} = \varphi \bmod I^{n - (d + 1)c}
$$
\end{lemma}

\begin{proof}
考虑备注 \ref{restricted-remark-improve-homomorphism} 中的障碍类
$$
o(\psi_n) \in \Ext^1_B(\NL_{B/A}^\wedge, I^nC/I^{2n}C)
$$
对任意 $C/I^nC$-模 $N$，有
\begin{align*}
\Ext^1_B(\NL_{B/A}^\wedge, N)
& =
\Ext^1_{C/I^nC}(\NL_{B/A}^\wedge \otimes_B^\mathbf{L} C/I^nC, N) \\
& =
\Ext^1_{C/I^nC}(\NL_{B/A}^\wedge \otimes_B C/I^nC, N)
\end{align*}
第一个等式由更多代数，引理
\ref{more-algebra-lemma-upgrade-adjoint-tensor-RHom} 得到，第二个等式由
更多代数，引理 \ref{more-algebra-lemma-two-term-base-change} 得到。
特别地，看到对所有 $C/I^nC$-模 $N$，
$\Ext^1_{C/I^nC}(\NL_{B/A}^\wedge \otimes_B C/I^nC, N)$ 被 $I^cC$ 湮灭。
由此可应用
局部上同调，引理 \ref{local-cohomology-lemma-bound-two-term-complex}，
得到 $o(\psi_n)$ 映为零：
$$
\Ext^1_{C/I^nC}(\NL_{B/A}^\wedge \otimes_B C/I^nC, I^{n'}C/I^{2n'}C) =
\Ext^1_B(\NL_{B/A}^\wedge, I^{n'}C/I^{2n'}C) =
$$
其中 $n' = n - (d + 1)c$。根据备注
\ref{restricted-remark-improve-homomorphism} 的讨论，得到映射
$$
\psi'_{2n'} : B \to C/I^{2n'}C
$$
它模 $I^{n'}$ 与 $\psi_n$ 一致。
注意，由于 $n \geq 2(d + 1)c + 1$，有 $2n' > n$。

\medskip\noindent
可以重复此过程。从 $n_0 = n$ 和 $\psi^0 = \psi_n$ 开始，最终得到严格递增的整数序列
$$
n_0 < n_1 < n_2 < \ldots
$$
以及 $A$-代数同态 $\psi^i : B \to C/I^{n_i}C$，
使得 $\psi^{i + 1}$ 与 $\psi^i$ 模 $I^{n_i - tc}$ 一致。
由于 $C$ 是 $I$-进完备的，可以令 $\varphi$ 为映射
$$
\psi^i \bmod I^{n_i - (d + 1)c} : B \to C/I^{n_i - (d + 1)c}C
$$
的极限，引理得证。
\end{proof}

\noindent
建议读者跳到下一节。事实上，如果其中的代数 $C$ 假设为诺特的，
下面两个引理就是上述结果的推论。

\begin{lemma}
\label{restricted-lemma-get-morphism-nonzerodivisor}
设 $I = (a)$ 为诺特环 $A$ 的主理想。
设 $B$ 为 \ref{restricted-equation-C-prime} 中的对象。
假设给定整数 $c \geq 0$，使得对所有 $B$-模 $N$，
$\Ext^1_B(\NL_{B/A}^\wedge, N)$ 被 $a^c$ 湮灭。
设 $C$ 为 $I$-进完备的 $A$-代数，且 $a$ 在 $C$ 上为非零因子。
令 $n > 2c$。对任意 $A$-代数映射 $\psi_n : B \to C/a^nC$，
存在 $A$-代数映射 $\varphi : B \to C$，使得
$$
\psi_n \bmod a^{n - c}C = \varphi \bmod a^{n - c}C
$$
\end{lemma}

\begin{proof}
考虑备注 \ref{restricted-remark-improve-homomorphism} 中的障碍类
$$
o(\psi_n) \in \Ext^1_B(\NL_{B/A}^\wedge, a^nC/a^{2n}C)
$$
由于 $a$ 在 $C$ 上为非零因子，$C$-模范畴中的映射
$a^c : a^nC/a^{2n}C \to a^nC/a^{2n}C$ 同构于映射
$a^nC/a^{2n}C \to a^{n - c}C/a^{2n - c}C$。
因此根据关于 $\NL_{B/A}^\wedge$ 的假设，可知类 $o(\psi_n)$ 映为零于
$$
\Ext^1_B(\NL_{B/A}^\wedge, a^{n - c}C/a^{2n - c}C)
$$
从而当然也映为零于
$$
\Ext^1_B(\NL_{B/A}^\wedge, a^{n - c}C/a^{2n - 2c}C)
$$
根据备注 \ref{restricted-remark-improve-homomorphism} 的讨论，得到映射
$$
\psi_{2n - 2c} : B \to C/a^{2n - 2c}C
$$
它模 $a^{n - c}C$ 与 $\psi_n$ 一致。
注意，由于 $n > 2c$，有 $2n - 2c > n$。

\medskip\noindent
可以重复此过程。从 $n_0 = n$ 和 $\psi^0 = \psi_n$ 开始，最终得到严格递增的整数序列
$$
n_0 < n_1 < n_2 < \ldots
$$
以及 $A$-代数同态 $\psi^i : B \to C/a^{n_i}C$，
使得 $\psi^{i + 1}$ 与 $\psi^i$ 模 $a^{n_i - c}C$ 一致。
由于 $C$ 是 $I$-进完备的，可以令 $\varphi$ 为映射
$$
\psi^i \bmod a^{n_i - c}C : B \to C/a^{n_i - c}C
$$
的极限，结论成立。
\end{proof}

\begin{lemma}
\label{restricted-lemma-get-morphism-principal}
设 $I = (a)$ 为诺特环 $A$ 的主理想。
设 $B$ 为 \ref{restricted-equation-C-prime} 中的对象。
假设给定整数 $c \geq 0$，使得对所有 $B$-模 $N$，
$\Ext^1_B(\NL_{B/A}^\wedge, N)$ 被 $a^c$ 湮灭。
设 $C$ 为 $I$-进完备的 $A$-代数。
假设给定整数 $d \geq 0$，使得 $C[a^\infty] \cap a^dC = 0$。
令 $n > \max(2c, c + d)$。对任意 $A$-代数映射
$\psi_n : B \to C/a^nC$，存在 $A$-代数映射 $\varphi : B \to C$，使得
$$
\psi_n \bmod a^{n - c} = \varphi \bmod a^{n - c}
$$
\end{lemma}

\noindent
若 $C$ 为诺特环，则对某个 $e \geq 0$ 有 $C[a^\infty] = C[a^e]$。
由 Artin--Rees 定理（代数，引理 \ref{algebra-lemma-Artin-Rees}），
存在整数 $f$，使得对所有 $n \geq f$，
$a^nC \cap C[a^\infty] \subset a^{n - f}C[a^\infty]$。
于是 $d = e + f$ 是引理中的一个整数。特别地，当 $C$ 是
\ref{restricted-equation-C-prime} 中的对象时，由引理
\ref{restricted-lemma-topologically-finite-type-Noetherian}，此论证也适用。

\begin{proof}
令 $C \to C'$ 为 $C$ 对 $C[a^\infty]$ 的商。由代数，引理
\ref{algebra-lemma-quotient-complete} 以及
$\bigcap (C[a^\infty] + a^nC) = C[a^\infty]$ 的事实（因为当 $n \geq d$ 时，
和 $C[a^\infty] + a^nC$ 是直和），$A$-代数 $C'$ 是 $I$-进完备的。
对 $m \geq d$，图表
$$
\xymatrix{
0 \ar[r] &
C[a^\infty] \ar[r] \ar[d] &
C \ar[r] \ar[d] & C' \ar[r] \ar[d] & 0 \\
0 \ar[r] &
C[a^\infty] \ar[r] &
C/a^m C \ar[r] & C'/a^m C' \ar[r] & 0
}
$$
具有正合的行。因此对所有 $m \geq d$，$C$ 是 $C'$ 与 $C/a^mC$
沿 $C'/a^mC'$ 的纤维积。由引理
\ref{restricted-lemma-get-morphism-nonzerodivisor}，可以选择同态
$\varphi' : B \to C'$，使得 $\varphi'$ 与 $\psi_n$ 作为映射到
$C'/a^{n - c}C'$ 时一致。于是得到同态
$$
(\varphi', \psi_n \bmod a^{n - c}C) : B \to
C' \times_{C'/a^{n - c}C'} C/a^{n - c}C
$$
由于 $n - c \geq d$，这与同态 $\varphi : B \to C$ 是同一件事。
证明完毕。
\end{proof}


\section{G-环上 rig-光滑代数的代数化}
\label{restricted-section-over-G-ring}

\noindent
若基环 $A$ 是诺特 G-环，则对于任意理想 $I\subset A$（不必是主理想）
以及任意 rig-光滑代数，可以证明 \cite[III Theorem 7]{Elkik}。

\begin{lemma}
\label{restricted-lemma-close-enough}
设 $I$ 为诺特环 $A$ 的理想。令 $r\geq0$，并记
$P=A[x_1,\ldots,x_r]$ 为 $I$-进完备化。考虑 $P$ 的一个商的自由分解
$$
P^{\oplus t} \xrightarrow{K} P^{\oplus m}
\xrightarrow{g_1, \ldots, g_m} P \to B \to 0
$$
假设 $B$ 在 $(A,I)$ 上 rig-光滑。则存在整数 $n$，使得对任意复形
$$
P^{\oplus t} \xrightarrow{K'} P^{\oplus m}
\xrightarrow{g'_1, \ldots, g'_m} P
$$
若 $g_i-g'_i\in I^nP$ 且
$K-K'\in I^n\text{Mat}(m\times t,P)$，则存在 $A$-代数同构
$B\to B'$，其中 $B'=P/(g'_1,\ldots,g'_m)$。
\end{lemma}

\begin{proof}
(A) 由定义 \ref{restricted-definition-rig-smooth-homomorphism}，可以选择
$c\geq0$，使得对所有 $B$-模 $N$，$I^c$ 湮灭
$\Ext^1_B(\NL_{B/A}^\wedge,N)$。

\medskip\noindent
(B) 由更多代数之引理 \ref{more-algebra-lemma-approximate-complex} 和
\ref{more-algebra-lemma-approximate-complex-graded}，存在常数
$c_1=c(g_1,\ldots,g_m,K)$，使得当 $n\geq c_1+1$ 时，复形
$$
P^{\oplus t} \xrightarrow{K'} P^{\oplus m}
\xrightarrow{g'_1, \ldots, g'_m} P \to B' \to 0
$$
正合，并且 $\text{Gr}_I(B)\cong\text{Gr}_I(B')$。

\medskip\noindent
(C) 令 $d_0=d(\text{Gr}_I(B))$ 和 $q_0=q(\text{Gr}_I(B))$
为局部上同调之 \ref{local-cohomology-section-uniform} 节中所得的整数。

\medskip\noindent
我们声称
$$
n=\max(c_1+1,q_0+(d_0+1)c,2(d_0+1)c+1)
$$
可行，其中 $c$ 如 (A)，$c_1$ 如 (B)，而 $q_0,d_0$ 如 (C)。

\medskip\noindent
设 $g'_1,\ldots,g'_m$ 和 $K'$ 如引理所述。
由于 $g_i=g'_i\in I^nP$，得到典范的 $A$-代数同态
$$
\psi_n:B\longrightarrow B'/I^nB'
$$
它诱导同构 $B/I^nB\to B'/I^nB'$。因为
$\text{Gr}_I(B)\cong\text{Gr}_I(B')$，有
$d_0=d(\text{Gr}_I(B'))$ 及 $q_0=q(\text{Gr}_I(B'))$；
又因为 $n\geq\max(q_0+(1+d_0)c,2(d_0+1)c+1)$，
可应用引理 \ref{restricted-lemma-get-morphism-general-better}，找到 $A$-代数同态
$$
\varphi:B\longrightarrow B'
$$
满足
$$
\varphi\bmod I^{n-(d_0+1)c}B'
=\psi_n\bmod I^{n-(d_0+1)c}B'
$$
由于 $n-(d_0+1)c>0$，可见 $\varphi$ 是一个 $A$-代数同态，
且模 $I$ 诱导出上面所得的同构 $B/IB\to B'/IB'$。
证明的其余部分说明这些事实迫使 $\varphi$ 成为同构；
建议读者自行找出证明。

\medskip\noindent
具体地，例如应用代数引理
\ref{algebra-lemma-completion-generalities} 的 (1)，并利用 $B$ 与 $B'$
均完备，可知 $\varphi$ 为满射。
因此 $\varphi$ 诱导满射
$\text{Gr}_I(B)\to\text{Gr}_I(B')$；由于源和目标是彼此同构的诺特环，
此满射必为同构，见代数引理
\ref{algebra-lemma-surjective-endo-noetherian-ring-is-iso}
（当然，也可以证明 $\varphi$ 诱导的正是上面所得的同构，不过这还需要一个小论证）。
于是对所有 $e\geq0$，$\varphi$ 诱导单射
$I^eB/I^{e+1}B\to I^eB'/I^{e+1}B'$。
这意味着 $\varphi$ 是单射，因为由 Krull 交定理（代数引理
\ref{algebra-lemma-intersect-powers-ideal-module-zero}），
对任意 $b\in B$，存在 $e\geq0$ 使得 $b\in I^eB$ 且
$b\not\in I^{e+1}B$。证明完毕。
\end{proof}

\begin{lemma}
\label{restricted-lemma-algebraize-easy}
设 $I$ 为诺特环 $A$ 的理想。设 $C^h$ 为有限型 $A$-代数 $C$
关于理想 $IC$ 的 Hensel 化。设 $J\subset C^h$ 为理想。
则存在有限型 $A$-代数 $B$，使得
$B^\wedge\cong(C^h/J)^\wedge$。
\end{lemma}

\begin{proof}
由更多代数之引理
\ref{more-algebra-lemma-henselization-Noetherian-pair}，环 $C^h$ 是诺特的。
设 $J=(g_1,\ldots,g_m)$。环 $C^h$ 是平展 $C$-代数 $C'$ 的滤过余极限，
其中 $C/IC\to C'/IC'$ 是同构（见更多代数之引理
\ref{more-algebra-lemma-henselization} 的证明）。
选取一个 $C'$，使得 $g_1,\ldots,g_m$ 是
$g'_1,\ldots,g'_m\in C'$ 的像。置
$B=C'/(g'_1,\ldots,g'_m)$，得到有限型 $A$-代数。
当然，$(C,IC)$ 与 $C',IC')$ 具有相同的 Hensel 化和相同的完备化。
由此容易推出 $B^\wedge=(C^h/J)^\wedge$。
\end{proof}

\begin{proposition}
\label{restricted-proposition-approximate}
设 $I$ 为诺特 G-环 $A$ 的理想。设 $B$ 为
\ref{restricted-equation-C-prime} 中的对象。若 $B$ 在 $(A,I)$ 上 rig-光滑，
则存在有限型 $A$-代数 $C$ 以及 $A$-代数同构
$B\cong C^\wedge$。
\end{proposition}

\begin{proof}
选取表示 $B=A[x_1,\ldots,x_r]^\wedge/J$。记
$P=A[x_1,\ldots,x_r]^\wedge$。选取生成元
$g_1,\ldots,g_m\in J$。选取 $g_1,\ldots,g_m$ 之间关系模的生成元
$k_1,\ldots,k_t$，也就是说，使
$$
P^{\oplus t}\xrightarrow{k_1,\ldots,k_t}
P^{\oplus m}\xrightarrow{g_1,\ldots,g_m}
P\to B\to0
$$
成为一个自由分解。写 $k_i=(k_{i1},\ldots,k_{im})$，于是有
\begin{equation}
\label{restricted-equation-relations-straight-up}
\sum\nolimits_j k_{ij}g_j=0
\end{equation}
其中 $i=1,\ldots,t$。记 $K=(k_{ij})$ 为以 $k_{ij}$ 为元素的 $m\times t$ 矩阵。

\medskip\noindent
令 $A[x_1,\ldots,x_r]^h$ 为对
$(A[x_1,\ldots,x_r],IA[x_1,\ldots,x_r])$ 的 Hensel 化，见更多代数之引理
\ref{more-algebra-lemma-henselization}。
可以且将会把 $A[x_1,\ldots,x_r]^h$ 看作
$P=A[x_1,\ldots,x_r]^\wedge$ 的子环，见更多代数之引理
\ref{more-algebra-lemma-henselization-Noetherian-pair}。
由于 $A$ 是诺特 G-环，$A[x_1,\ldots,x_r]$ 也是诺特 G-环，见更多代数之命题
\ref{more-algebra-proposition-finite-type-over-G-ring}。
因此，对于理想 $I$ 所生成的理想，映射
$A[x_1,\ldots,x_r]^h\to A[x_1,\ldots,x_r]^\wedge=P$
具有逼近性质，见环映射的光滑化之引理
\ref{smoothing-lemma-henselian-pair}。
选择引理 \ref{restricted-lemma-close-enough} 中充分大的整数 $n$。
由逼近性质，可以选择
$g'_1,\ldots,g'_m\in A[x_1,\ldots,x_r]^h$ 以及矩阵
$$
K'=(k'_{ij})\in\text{Mat}(m\times t,A[x_1,\ldots,x_r]^h)
$$
使得在 $A[x_1,\ldots,x_r]^h$ 中
$\sum\nolimits_j k'_{ij}g'_j=0$，并且
$g_i-g'_i\in I^nP$ 及
$K-K'\in I^n\text{Mat}(m\times t,P)$。
由 $n$ 的选择，可知存在同构
$$
B\to P/(g'_1,\ldots,g'_m)=
\left(A[x_1,\ldots,x_r]^h/(g'_1,\ldots,g'_m)\right)^\wedge
$$
由引理 \ref{restricted-lemma-algebraize-easy}，证明完毕。
\end{proof}

\noindent
若 $A$ 只是诺特环而不是 G-环，下面的引理一般不成立。
事实上，若 $(A,\mathfrak m)$ 是局部环且 $I=\mathfrak m$，
则该引理等价于 $A^h$ 的 Artin 逼近性质（如环映射的光滑化之定理
\ref{smoothing-theorem-approximation-property}），
而这一性质并非对每个诺特局部环都成立。

\begin{lemma}
\label{restricted-lemma-fully-faithfulness}
设 $A$ 为诺特 G-环，且 $I\subset A$ 为理想。
设 $B,C$ 为有限型 $A$-代数。对 $I$-进完备化的任意 $A$-代数映射
$\varphi:B^\wedge\to C^\wedge$ 以及任意 $N\geq1$，存在
\begin{enumerate}
\item 一个平展环映射 $C\to C'$，它诱导同构
$C/IC\to C'/IC'$；
\item 一个 $A$-代数映射 $\varphi:B\to C'$；
\end{enumerate}
使得 $\varphi$ 与 $\psi$ 模 $I^N$ 作为到
$C^\wedge=(C')^\wedge$ 的映射时一致。
\end{lemma}

\begin{proof}
引理的陈述是有意义的：$C\to C'$ 是平坦的（代数，引理
\ref{algebra-lemma-etale}），因而对所有 $n$ 诱导同构
$C/I^nC\to C'/I^nC'$（更多代数，引理
\ref{more-algebra-lemma-neighbourhood-isomorphism}），进而诱导完备化之间的同构。
令 $C^h$ 为对 $(C,IC)$ 的 Hensel 化，见更多代数之引理
\ref{more-algebra-lemma-henselization}。
则 $C^h$ 是这些代数 $C'$ 的滤过余极限，并且映射
$C\to C'\to C^h$ 诱导完备化之间的同构（更多代数，引理
\ref{more-algebra-lemma-henselization-Noetherian-pair}）。
因此，只需证明存在一个模 $I^N$ 与 $\psi$ 同余的 $A$-代数映射
$B\to C^h$。
写
$$
B=A[x_1,\ldots,x_n]/(f_1,\ldots,f_m)
$$
环映射 $\psi$ 对应于元素
$\hat c_1,\ldots,\hat c_n\in C^\wedge$，满足对
$j=1,\ldots,m$ 有
$$
f_j(\hat c_1,\ldots,\hat c_n)=0
$$
由于 $A$ 是诺特 G-环，$C$ 也是诺特 G-环，见更多代数之命题
\ref{more-algebra-proposition-finite-type-over-G-ring}。
因此环映射的光滑化之引理
\ref{smoothing-lemma-henselian-pair} 适用，给出元素
$c_1,\ldots,c_n\in C^h$，使得对 $j=1,\ldots,m$ 有
$f_j(c_1,\ldots,c_n)=0$，并且
$\hat c_i-c_i\in I^NC^h$。
这就按要求确定了映射 $B\to C^h$。
\end{proof}

\section{rig-光滑代数的代数化}
\label{restricted-section-algebraization-rig-smooth}

\noindent
事实表明，若 rig-光滑代数具有一种特定的表示，则其代数化是直接的。
相关讨论还请参见备注 \ref{restricted-remark-discussion}。

\begin{lemma}
\label{restricted-lemma-presentation-rig-smooth}
设 $A$ 为环。设
$f_1,\ldots,f_m\in A[x_1,\ldots,x_n]$，并置
$B=A[x_1,\ldots,x_n]/(f_1,\ldots,f_m)$。
假设 $m\leq n$，并置
$g=\det_{1\leq i,j\leq m}(\partial f_j/\partial x_i)$。则
\begin{enumerate}
\item 对每个 $B$-模 $N$，$g$ 湮灭
$\Ext^1_B(\NL_{B/A},N)$；
\item 若 $n=m$，则在 $D(B)$ 中，$\NL_{B/A}$ 上乘以 $g$ 的映射为 $0$。
\end{enumerate}
\end{lemma}

\begin{proof}
令 $T$ 为 $m\times m$ 矩阵，其元素为
$\partial f_j/\partial x_i$（$1\leq i,j\leq n$）。
令 $K\in D(B)$ 由复形
$T:B^{\oplus m}\to B^{\oplus m}$ 表示，其项位于次数 $-1$ 和 $0$。
由更多代数之引理 \ref{more-algebra-lemma-silly}，
在 $D(B)$ 中 $g:K\to K$ 为零。置 $J=(f_1,\ldots,f_m)$。
回忆 $\NL_{B/A}$ 同伦等价于
$J/J^2\to\bigoplus_{i=1,\ldots,n}B\text{d}x_i$，见代数之
\ref{algebra-section-netherlander} 节。
记 $L$ 为复形
$J/J^2\to\bigoplus_{i=1,\ldots,m}B\text{d}x_i$，
从而有商映射 $\NL_{B/A}\to L$。
还有满的复形映射 $K\to L$：把次数 $-1$ 的项 $B^{\oplus m}$
中的第 $j$ 个基元素送到 $J/J^2$ 中 $f_j$ 的类。图示为
$$
\NL_{B/A}\to L\leftarrow K
$$
由更多代数之引理
\ref{more-algebra-lemma-two-term-surjection-map-zero}，
可知在 $D(B)$ 中 $L$ 上乘以 $g$ 的映射为 $0$。
另一方面，特异三角
$B^{\oplus n-m}[0]\to\NL_{B/A}\to L$
表明，对每个 $B$-模 $N$，映射
$\Ext^1_B(L,N)\to\Ext^1_B(\NL_{B/A},N)$ 满射，
故后者被 $g$ 湮灭。这证明了 (1)。
若 $n=m$，则 $\NL_{B/A}=L$，从而 (2) 成立。
\end{proof}

\begin{lemma}
\label{restricted-lemma-approximate-presentation-rig-smooth}
设 $I$ 为诺特环 $A$ 的理想。设 $B$ 为
\ref{restricted-equation-C-prime} 中的对象。取一个表示
$B=A[x_1,\ldots,x_r]^\wedge/J$。
假设存在元素 $b\in B$、整数 $0\leq m\leq r$ 以及
$f_1,\ldots,f_m\in J$，使得
\begin{enumerate}
\item 在 $\Spec(B)$ 中，$V(b)\subset V(IB)$；
\item
$\Delta=\det_{1\leq i,j\leq m}(\partial f_j/\partial x_i)$
在 $B$ 中的像整除 $b$；以及
\item $bJ\subset(f_1,\ldots,f_m)+J^2$。
\end{enumerate}
则存在有限型 $A$-代数 $C$ 以及 $A$-代数同构
$B\cong C^\wedge$。
\end{lemma}

\begin{proof}
这些条件蕴含 $B$ 在 $(A,I)$ 上 rig-光滑，见引理
\ref{restricted-lemma-equivalent-with-artin-smooth}。
对某个 $b'\in B$，在 $B$ 中写 $b'\Delta=b$。
设 $I=(a_1,\ldots,a_t)$。由于 $V(b)\subset V(IB)$，
存在整数 $c\geq0$，使得 $I^cB\subset bB$。
对某些 $b_i\in B$，在 $B$ 中写 $bb_i=a_i^c$。

\medskip\noindent
选择整数 $n\gg0$（稍后将看到需取多大）。
选择多项式 $f'_1,\ldots,f'_m\in A[x_1,\ldots,x_r]$，使得
$f_i-f'_i\in I^nA[x_1,\ldots,x_r]^\wedge$。
置
$\Delta'=\det_{1\leq i,j\leq m}(\partial f'_j/\partial x_i)$，
并考虑有限型 $A$-代数
$$
C=A[x_1,\ldots,x_r,z_1,\ldots,z_t]/
(f'_1,\ldots,f'_m,
z_1\Delta'-a_1^c,\ldots,z_t\Delta'-a_t^c)
$$
我们将对 $C$ 应用引理 \ref{restricted-lemma-presentation-rig-smooth}。
计算得
$$
\det\left(
\begin{matrix}
\text{由下列各式的偏导数组成的矩阵}\\
f'_1,\ldots,f'_m,z_1\Delta'-a_1^c,\ldots,z_t\Delta'-a_t^c\\
\text{关于变量}\\
x_1,\ldots,x_m,z_1,\ldots,z_t
\end{matrix}
\right)=(\Delta')^{t+1}
$$
因此，对所有 $C$-模 $N$，
$\Ext^1_C(\NL_{C/A},N)$ 被 $(\Delta')^{t+1}$ 湮灭。
由于在 $C$ 中 $a_i^c$ 被 $\Delta'$ 整除，
$a_i^{(t+1)c}$ 也湮灭这些 $\Ext^1$。
因此，对所有 $C$-模 $N$，$I^{c_1}$ 湮灭
$\Ext^1_C(\NL_{C/A},N)$，其中
$c_1=1+t((t+1)c-1)$。
$c_1$ 的确切值对其余论证并不重要；重要的是它与 $n$ 无关。

\medskip\noindent
由引理 \ref{restricted-lemma-NL-is-completion}，
$\NL_{C^\wedge/A}^\wedge=\NL_{C/A}\otimes_CC^\wedge$。
再由更多代数之引理
\ref{more-algebra-lemma-base-change-property-ext-1-annihilated} 和
\ref{more-algebra-lemma-two-term-base-change}，
可知对任意 $C^\wedge$-模 $N$，
$I^{c_1}$ 在
$\Ext^1_{C^\wedge}(\NL_{C^\wedge/A}^\wedge,N)$ 上的乘法也为零。
特别地，$C^\wedge$ 在 $(A,I)$ 上 rig-光滑。

\medskip\noindent
注意，有满的 $A$-代数同态
$$
\psi_n:C\longrightarrow B/I^nB
$$
它把 $x_i$ 的类送到 $x_i$ 的类，并把 $z_i$ 的类送到 $b_ib'$ 的类。
这由证明第一段中对 $b'$ 和 $b_i$ 的选择成立。

\medskip\noindent
令 $d=d(\text{Gr}_I(B))$ 和 $q_0=q(\text{Gr}_I(B))$
为局部上同调之 \ref{local-cohomology-section-uniform} 节中所得的整数。
由引理 \ref{restricted-lemma-get-morphism-general-better}，若取
$$
n\geq\max(q_0+(d+1)c_1,2(d+1)c_1+1),
$$
则可找到 $A$-代数同态
$\varphi:C^\wedge\to B$，它模
$I^{n-(d+1)c_1}B$ 与 $\psi_n$ 同余。

\medskip\noindent
由于 $\varphi:C^\wedge\to B$ 模 $I$ 满射，可知它本身满射
（例如使用代数引理 \ref{algebra-lemma-completion-generalities}）。
为完成证明，只需说明
$\Ker(\varphi)/\Ker(\varphi)^2$ 被 $I$ 的某个幂湮灭，见更多代数之引理
\ref{more-algebra-lemma-quotient-by-idempotent}。

\medskip\noindent
由于 $\varphi$ 满射，可知
$\NL_{B/C^{\wedge}}^\wedge$ 的上同调模为
$H^0(\NL_{B/C^{\wedge}}^\wedge)=0$ 以及
$H^{-1}(\NL_{B/C^{\wedge}}^\wedge)=\Ker(\varphi)/\Ker(\varphi)^2$。
由引理 \ref{restricted-lemma-exact-sequence-NL}，有正合序列
$$
H^{-1}(\NL_{C^\wedge/A}^\wedge\otimes_{C^\wedge}B)\to
H^{-1}(\NL_{B/A}^\wedge)\to
H^{-1}(\NL_{B/C^{\wedge}}^\wedge)\to
H^0(\NL_{C^\wedge/A}^\wedge\otimes_{C^\wedge}B)\to
H^0(\NL_{B/A}^\wedge)\to0
$$
前两个模被 $I$ 的某个幂湮灭，因为 $B$ 和 $C^\wedge$
在 $(A,I)$ 上 rig-光滑。因此，只需证明满射
$$
H^0(\NL_{C^\wedge/A}^\wedge\otimes_{C^\wedge}B)\to
H^0(\NL_{B/A}^\wedge)
$$
的核被 $I$ 的某个幂湮灭。
为此，只需证明它被 $b$ 的某个幂湮灭。换言之，只需证明
$$
H^0(\NL_{C^\wedge/A}^\wedge)\otimes_{C^\wedge}B[1/b]
\longrightarrow
H^0(\NL_{B/A}^\wedge)\otimes_BB[1/b]
$$
是同构。然而，两边都是秩为 $r-m$ 的自由 $B[1/b]$-模，
其基为 $\text{d}x_{m+1},\ldots,\text{d}x_r$。证明完毕。
\end{proof}

\begin{remark}
\label{restricted-remark-discussion}
设 $I$ 为诺特环 $A$ 的理想。设 $B$ 为
\ref{restricted-equation-C-prime} 中的对象，并且在 $(A,I)$ 上 rig-光滑。
\cite[Theorem 1.2]{gabber-zavyalov} 证明了 $B$ 同构于某个有限型
$A$-代数的 $I$-进完备化。该结果取代了下列部分结果：
\begin{enumerate}
\item 若 $A$ 是 G-环，则结论由命题
\ref{restricted-proposition-approximate} 得到；
\item 若 $B$ 在 $(A,I)$ 上 rig-\,'etale，则结论由引理
\ref{restricted-lemma-approximate} 得到；
\item 若 $I$ 为主理想，则结论由
\cite[III Theorem 7]{Elkik} 得到。
\end{enumerate}
\end{remark}

\section{rig-\,'etale 代数}
\label{restricted-section-rig-etale}

\noindent
鉴于 rig-光滑代数的定义（定义
\ref{restricted-definition-rig-smooth-homomorphism}），下面的定义并不令人意外。

\begin{definition}
\label{restricted-definition-rig-etale-homomorphism}
设 $A$ 为诺特环，且 $I\subset A$ 为理想。
设 $B$ 为 \ref{restricted-equation-C-prime} 中的对象。
若存在整数 $c\geq0$，使得对所有 $a\in I^c$，
在 $D(B)$ 中 $\NL_{B/A}^\wedge$ 上乘以 $a$ 的映射为零，
则称 $B$ {\it 在 $(A,I)$ 上 rig-\,'etale}。
\end{definition}

\noindent
下一引理中的条件 \ref{restricted-item-condition-artin} 是
\cite{ArtinII} 用来定义形式修改的条件之一。
我们将其加入条件列表，以便之后与我们的条件比较。

\begin{lemma}
\label{restricted-lemma-equivalent-with-artin}
设 $A$ 为诺特环，且 $I\subset A$ 为理想。
设 $B$ 为 \ref{restricted-equation-C-prime} 中的对象。写成
$B=A[x_1,\ldots,x_r]^\wedge/J$
（引理 \ref{restricted-lemma-topologically-finite-type-Noetherian}），
并令
$\NL_{B/A}^\wedge=(J/J^2\to\bigoplus B\text{d}x_i)$
为其朴素余切复形（\ref{restricted-equation-NL}）。
以下条件等价：
\begin{enumerate}
\item $B$ 在 $(A,I)$ 上 rig-\,'etale；
\item
\label{restricted-item-zero-on-NL}
存在 $c\geq0$，使得对所有 $a\in I^c$，在 $D(B)$ 中
$\NL_{B/A}^\wedge$ 上乘以 $a$ 的映射为零；
\item
\label{restricted-item-zero-on-cohomology-NL}
存在 $c\geq0$，使得 $H^i(\NL_{B/A}^\wedge)$（$i=-1,0$）
被 $I^c$ 湮灭；
\item
\label{restricted-item-zero-on-cohomology-NL-truncations}
存在 $c\geq0$，使得对所有 $n\geq1$，
$H^i(\NL_{B_n/A_n})$（$i=-1,0$）被 $I^c$ 湮灭，
其中 $A_n=A/I^n$ 且 $B_n=B/I^nB$；
\item
\label{restricted-item-condition-artin-pre-pre}
对每个 $a\in I$，存在 $c\geq0$，使得
\begin{enumerate}
\item $a^c$ 湮灭 $H^0(\NL_{B/A}^\wedge)$；并且
\item 存在 $f_1,\ldots,f_r\in J$，使得
$a^cJ\subset(f_1,\ldots,f_r)+J^2$。
\end{enumerate}
\item
\label{restricted-item-condition-artin-pre}
对每个 $a\in I$，存在 $f_1,\ldots,f_r\in J$ 及 $c\geq0$，使得
\begin{enumerate}
\item $\det_{1\leq i,j\leq r}(\partial f_j/\partial x_i)$
在 $B$ 中整除 $a^c$；并且
\item $a^cJ\subset(f_1,\ldots,f_r)+J^2$。
\end{enumerate}
\item
\label{restricted-item-condition-artin}
选取 $J$ 的生成元 $f_1,\ldots,f_t$，则
\begin{enumerate}
\item $B$ 在 $A$ 上的 Jacobi 理想，即 $B$ 中由矩阵
$(\partial f_j/\partial x_i)_{1\leq i\leq r,1\leq j\leq t}$
的 $r\times r$ 子式生成的理想，包含某个 $c$ 的理想 $I^cB$；并且
\item $B$ 在 $A$ 上的 Cramer 理想，即 $B$ 中由 $J$
作为 $A[x_1,\ldots,x_r]^\wedge$-模的第 $r$ 个 Fitting 理想在 $B$
中的像所生成的理想，包含某个 $c$ 的 $I^cB$。
\end{enumerate}
\end{enumerate}
\end{lemma}

\begin{proof}
(1) 与 \ref{restricted-item-zero-on-NL} 的等价性只是定义
\ref{restricted-definition-rig-etale-homomorphism} 的重述。

\medskip\noindent
\ref{restricted-item-zero-on-NL} 与
\ref{restricted-item-zero-on-cohomology-NL} 的等价性由更多代数之引理
\ref{more-algebra-lemma-zero-in-derived} 得到。

\medskip\noindent
\ref{restricted-item-zero-on-cohomology-NL} 与
\ref{restricted-item-zero-on-cohomology-NL-truncations} 的等价性来自系统
$\{\NL_{B_n/A_n}\}$ 与
$\NL_{B/A}^\wedge\otimes_BB_n$ 严格同构这一事实，见引理
\ref{restricted-lemma-NL-is-limit}。略去一些细节。

\medskip\noindent
假设 \ref{restricted-item-zero-on-NL}。令 $a\in I$。
选择 $c$，使得 $\NL_{B/A}^\wedge$ 上乘以 $a^c$ 的映射为零。
由更多代数之引理
\ref{more-algebra-lemma-map-out-of-almost-free} 的 (1)，
存在映射 $\alpha:\bigoplus B\text{d}x_i\to J/J^2$，
使得 $\text{d}\circ\alpha$ 与 $\alpha\circ\text{d}$
均为乘以 $a^c$ 的映射。
取 $f_i\in J$，使其模 $J^2$ 的类等于
$\alpha(\text{d}x_i)$。简单计算表明
\ref{restricted-item-condition-artin-pre} 的 (a)、(b) 成立。

\medskip\noindent
略去 \ref{restricted-item-condition-artin-pre} 蕴含
\ref{restricted-item-condition-artin-pre-pre} 的验证；这只是关于 $B$ 上形如
$M\to B^{\oplus r}$ 的两项复形的一个陈述。

\medskip\noindent
假设 \ref{restricted-item-condition-artin-pre-pre} 成立。
设 $I=(a_1,\ldots,a_t)$。令 $c_i\geq0$ 为使
\ref{restricted-item-condition-artin-pre-pre} 的 (a)、(b) 对
$a_i^{c_i}$ 成立的整数。于是
$I^{\sum c_i}$ 湮灭 $H^0(\NL_{B/A}^\wedge)$。
令 $f_{i,1},\ldots,f_{i,r}\in J$ 如
\ref{restricted-item-condition-artin-pre-pre} 的 (b) 中针对 $a_i$ 所述。
考虑复合
$$
B^{\oplus r}\to J/J^2\to\bigoplus B\text{d}x_i
$$
其中第 $j$ 个基向量映到 $J/J^2$ 中 $f_{i,j}$ 的类。
由 \ref{restricted-item-condition-artin-pre-pre} 的 (a)、(b)，
该复合的余核被 $a_i^{2c_i}$ 湮灭。
因此，该映射在对 $a_i^{c_i}$ 取逆后满射，进而为同构
（代数，引理 \ref{algebra-lemma-fun}）。
于是 $B^{\oplus r}\to\bigoplus B\text{d}x_i$ 的核是
$a_i$-幂挠模，从而
$$
H^{-1}(\NL_{B/A}^\wedge)
=\Ker(J/J^2\to\bigoplus B\text{d}x_i)
$$
也是 $a_i$-幂挠模。由于 $B$ 为诺特环（引理
\ref{restricted-lemma-topologically-finite-type-Noetherian}），
包括 $H^{-1}(\NL_{B/A}^\wedge)$ 在内的所有模都是有限的。
因此，对某个 $d_i\geq0$，$a_i^{d_i}$ 湮灭
$H^{-1}(\NL_{B/A}^\wedge)$。
由此，$I^{\sum d_i}$ 湮灭 $H^{-1}(\NL_{B/A}^\wedge)$，
故 \ref{restricted-item-zero-on-cohomology-NL} 成立。

\medskip\noindent
因此，条件
\ref{restricted-item-zero-on-NL}、
\ref{restricted-item-zero-on-cohomology-NL}、
\ref{restricted-item-zero-on-cohomology-NL-truncations}、
\ref{restricted-item-condition-artin-pre-pre} 和
\ref{restricted-item-condition-artin-pre} 彼此等价。
还需说明这些条件与 \ref{restricted-item-condition-artin} 等价。
注意，Cramer 理想
$\text{Fit}_r(J)B$ 等于
$\text{Fit}_r(J/J^2)$，因为
$J/J^2=J\otimes_{A[x_1,\ldots,x_r]^\wedge}B$，见更多代数之引理
\ref{more-algebra-lemma-fitting-ideal-basics} 的 (3)。
还要注意，Jacobi 理想正是
$\text{Fit}_0(H^0(\NL_{B/A}^\wedge))$。
因此，
\ref{restricted-item-zero-on-cohomology-NL} 与
\ref{restricted-item-condition-artin} 的等价性是一个纯代数问题，
下面一段讨论这一问题。

\medskip\noindent
设 $R$ 为诺特环，且 $I\subset R$ 为理想。
设 $M\xrightarrow{d}R^{\oplus r}$ 为两项复形。
需要证明以下条件等价：
\begin{enumerate}
\item[(A)] $M\to R^{\oplus r}$ 的上同调被 $I$ 的某个幂湮灭；以及
\item[(B)] 理想 $\text{Fit}_r(M)$ 和
$\text{Fit}_0(\text{Coker}(d))$ 包含 $I$ 的某个幂。
\end{enumerate}
由于 $R$ 是诺特环，可以把 (2) 改写为相应闭子概形之间的包含，见代数之引理
\ref{algebra-lemma-Zariski-topology} 和
\ref{algebra-lemma-Noetherian-power}。
另一方面，在 $V(\text{Fit}_0(\Coker(d)))$ 的补集上，
$d$ 的余核为零；而在 $V(\text{Fit}_r(M))$ 的补集上，
模 $M$ 局部由 $r$ 个元素生成，见更多代数之引理
\ref{more-algebra-lemma-fitting-ideal-generate-locally}。
因此，(B) 等价于
\begin{enumerate}
\item[(C)] 在 $V(I)$ 之外，$d$ 的余核为零，且模 $M$
局部由 $\leq r$ 个元素生成。
\end{enumerate}
当然，这等价于 $M\to R^{\oplus r}$ 在
$\Spec(R)\setminus V(I)$ 上具有零上同调，而后者又等价于 (A)。
证明完毕。
\end{proof}

\begin{lemma}
\label{restricted-lemma-rig-etale-rig-smooth}
设 $A$ 为诺特环，且 $I$ 为理想。
设 $B$ 为 \ref{restricted-equation-C-prime} 中的对象。
若 $B$ 在 $(A,I)$ 上 rig-\,'etale，则 $B$ 在 $(A,I)$ 上 rig-光滑。
\end{lemma}

\begin{proof}
立即由定义 \ref{restricted-definition-rig-smooth-homomorphism} 和
\ref{restricted-definition-rig-etale-homomorphism} 得到。
\end{proof}

\begin{lemma}
\label{restricted-lemma-rig-etale}
设 $A$ 为诺特环，且 $I$ 为理想。
设 $B$ 为有限型 $A$-代数。
\begin{enumerate}
\item 若 $\Spec(B)\to\Spec(A)$ 在
$\Spec(A)\setminus V(I)$ 上平展，则 $B^\wedge$ 满足引理
\ref{restricted-lemma-equivalent-with-artin} 的等价条件。
\item 若 $B^\wedge$ 满足引理
\ref{restricted-lemma-equivalent-with-artin} 的等价条件，则存在
$g\in1+IB$，使得 $\Spec(B_g)$ 在
$\Spec(A)\setminus V(I)$ 上平展。
\end{enumerate}
\end{lemma}

\begin{proof}
假设 $B^\wedge$ 满足引理
\ref{restricted-lemma-equivalent-with-artin} 的等价条件。
朴素余切复形 $\NL_{B/A}$ 是有限型 $B$-模组成的复形，
因而 $H^{-1}$ 和 $H^0$ 是有限 $B$-模。
完备化是有限 $B$-模上的正合函子（代数，引理
\ref{algebra-lemma-completion-flat}），并且
$\NL_{B^\wedge/A}^\wedge$ 是复形 $\NL_{B/A}$ 的完备化
（选取表示即可容易看出）。
因此，假设蕴含存在 $c\geq0$，使得对所有 $n$，
$H^{-1}/I^nH^{-1}$ 和 $H^0/I^nH^0$ 被 $I^c$ 湮灭。
由 Nakayama 引理（代数，引理 \ref{algebra-lemma-NAK}），
这意味着 $I^cH^{-1}$ 和 $I^cH^0$ 被某个形如
$g=1+x$（$x\in IB$）的元素湮灭。
对 $g$ 取逆后（这不改变商 $B/I^nB$），
可见 $\NL_{B/A}$ 的上同调被 $I^c$ 湮灭。
于是，由平展环映射的定义，
$A\to B$ 在 $B$ 中任何不位于 $V(I)$ 上方的素理想处平展，见代数之定义
\ref{algebra-definition-etale}。

\medskip\noindent
反之，假设 $\Spec(B)\to\Spec(A)$ 在
$\Spec(A)\setminus V(I)$ 上平展。
则对每个 $a\in I$，存在 $c\geq0$，使得
$\NL_{B/A}$ 上乘以 $a^c$ 的映射为零。
由于 $\NL_{B^\wedge/A}^\wedge$ 是 $\NL_{B/A}$ 的导出完备化
（见引理 \ref{restricted-lemma-NL-is-limit}），可知 $B^\wedge$
满足引理 \ref{restricted-lemma-equivalent-with-artin} 的等价条件。
\end{proof}

\begin{lemma}
\label{restricted-lemma-zero-after-modding-out}
设 $(A_1,I_1)\to(A_2,I_2)$ 如备注
\ref{restricted-remark-base-change}，且 $A_1,A_2$ 为诺特环。
设 $B_1$ 是 $(A_1,I_1)$ 的
\ref{restricted-equation-C-prime} 中的对象，并令 $B_2$ 为 $B_1$ 的基变换。
若在 $D(B_1)$ 中，$\NL^\wedge_{B_1/A_1}$ 上乘以
$f_1\in B_1$ 的映射为零，则在 $D(B_2)$ 中，
$\NL^\wedge_{B_2/A_2}$ 上乘以其像 $f_2\in B_2$ 的映射为零。
\end{lemma}

\begin{proof}
由引理 \ref{restricted-lemma-NL-base-change}，存在映射
$$
\NL_{B_1/A_1}\otimes_{B_2}B_1\to\NL_{B_2/A_2}
$$
它在 $H^0$ 上诱导同构，在 $H^{-1}$ 上诱导满射。
于是结论由更多代数之引理
\ref{more-algebra-lemma-two-term-surjection-map-zero} 得到。
\end{proof}

\begin{lemma}
\label{restricted-lemma-base-change-rig-etale-homomorphism}
设 $A_1\to A_2$ 为诺特环映射。设 $I_i\subset A_i$ 为理想，
使得 $V(I_1A_2)=V(I_2)$。设 $B_1$ 是 $(A_1,I_1)$ 的
\ref{restricted-equation-C-prime} 中的对象。
设 $B_2$ 为按备注 \ref{restricted-remark-base-change} 得到的 $B_1$ 的基变换。
若 $B_1$ 在 $(A_1,I_1)$ 上 rig-\,'etale，
则 $B_2$ 在 $(A_2,I_2)$ 上 rig-\,'etale。
\end{lemma}

\begin{proof}
由引理 \ref{restricted-lemma-zero-after-modding-out}、定义
\ref{restricted-definition-rig-etale-homomorphism} 以及以下事实得到：
由于 $A_2$ 诺特，对某个 $c\geq0$ 有 $I_2^c\subset I_1A_2$。
\end{proof}

\begin{lemma}
\label{restricted-lemma-fully-faithful-etale-over-complement}
设 $A$ 为诺特环，且 $I\subset A$ 为理想。
设 $B$ 为有限型 $A$-代数，使得
$\Spec(B)\to\Spec(A)$ 在 $\Spec(A)\setminus V(I)$ 上平展。
设 $C$ 为诺特 $A$-代数。则 $I$-进完备化的任意 $A$-代数映射
$B^\wedge\to C^\wedge$ 都来自唯一的 $A$-代数映射
$$
B\longrightarrow C^h
$$
其中 $C^h$ 是对 $(C,IC)$ 的 Hensel 化，如更多代数之引理
\ref{more-algebra-lemma-henselization} 所述。
此外，任意 $A$-代数同态 $B\to C^h$ 均经由某个平展
$C$-代数 $C'$ 分解，其中 $C/IC\to C'/IC'$ 是同构。
\end{lemma}

\begin{proof}
唯一性来自 $C^h$ 是 $C^\wedge$ 的子环这一事实，例如见更多代数之引理
\ref{more-algebra-lemma-henselization-Noetherian-pair}。
最后一个断言来自 $C^h$ 是这些 $C$-代数 $C'$ 的滤过余极限这一事实，
见更多代数之引理 \ref{more-algebra-lemma-henselization} 的证明。
说完这些，现在转向存在性的证明。

\medskip\noindent
设 $\varphi:B^\wedge\to C^\wedge$ 为给定的映射。
它定义了映射 $C\to B\otimes_AC$ 的完备化的截面
$$
\sigma:(B\otimes_AC)^\wedge\longrightarrow C^\wedge
$$
可以用 $(C,IC,B\otimes_AC,C,\sigma)$ 代替
$(A,I,B,C,\varphi)$。由此可见，可以假设 $A=C$。

\medskip\noindent
在 $A=C$ 情形下证明存在性。此时映射
$\varphi:B^\wedge\to A^\wedge$ 必为满射。
由引理 \ref{restricted-lemma-rig-etale} 和
\ref{restricted-lemma-exact-sequence-NL}，可见
$$
\NL_{A^\wedge/\!_\varphi B^\wedge}^\wedge
$$
的上同调群被 $I$ 的某个幂湮灭。
由于 $\varphi$ 满射，这蕴含
$\Ker(\varphi)/\Ker(\varphi)^2$ 被 $I$ 的某个幂湮灭。
因此，$\varphi:B^\wedge\to A^\wedge$ 是某个有限型
$B$-代数 $B\to D$ 的完备化，见更多代数之引理
\ref{more-algebra-lemma-quotient-by-idempotent}。
于是 $A\to D$ 是有限型代数映射，并诱导同构
$A^\wedge\to D^\wedge$。
由引理 \ref{restricted-lemma-rig-etale}，可以用一个局部化替换 $D$，
并假设 $A\to D$ 在 $V(I)$ 之外平展。
由于 $A^\wedge\to D^\wedge$ 为同构，
可见 $\Spec(D)\to\Spec(A)$ 在 $V(ID)$ 的某个邻域中也平展
（例如由更多态射之引理
\ref{more-morphisms-lemma-check-smoothness-on-infinitesimal-nbhds}）。
因此 $\Spec(D)\to\Spec(A)$ 平展。所以 $D$ 映到 $A^h$，引理得证。
\end{proof}

\section{一个推出论证}
\label{restricted-section-approximation}

\noindent
本节唯一的目标是证明下面的引理；它将在 rig-\,'etale 代数的代数化中起关键作用。
我们将用到一点代数空间理论来证明该引理。
本章较早的版本给出了一个使用代数与少量变形理论的（长得多的）证明；
有兴趣的读者可在 Stacks 项目的历史中找到它。

\begin{lemma}
\label{restricted-lemma-lift-approximation}
设 $A$ 为诺特环，且 $I\subset A$ 为理想。
设 $J\subset A$ 为幂零理想。考虑交换图表
$$
\xymatrix{
C \ar[r] & C_0 \ar@{=}[r] & C/JC \\
& B_0 \ar[u] \\
A \ar[r] \ar[uu] & A_0 \ar[u] \ar@{=}[r] & A/J
}
$$
其中竖直箭头均为有限型，并满足
\begin{enumerate}
\item $\Spec(C)\to\Spec(A)$ 在
$\Spec(A)\setminus V(I)$ 上平展；
\item $\Spec(B_0)\to\Spec(A_0)$ 在
$\Spec(A_0)\setminus V(IA_0)$ 上平展；并且
\item $B_0\to C_0$ 平展，且诱导同构
$B_0/IB_0=C_0/IC_0$。
\end{enumerate}
则可将上面的图表补成交换图表
$$
\xymatrix{
C \ar[r] & C/JC \\
B \ar[u] \ar[r] & B_0 \ar[u] \\
A \ar[r] \ar[u] & A/J \ar[u]
}
$$
其中 $A\to B$ 为有限型，$B/JB=B_0$，$B\to C$ 平展，并且
$\Spec(B)\to\Spec(A)$ 在 $\Spec(A)\setminus V(I)$ 上平展。
\end{lemma}

\begin{proof}
置 $X=\Spec(A)$、$X_0=\Spec(A_0)$、$Y_0=\Spec(B_0)$、
$Z=\Spec(C)$、$Z_0=\Spec(C_0)$。
此外，分别记
$U\subset X$、$U_0\subset X_0$、$V_0\subset Y_0$、
$W\subset Z$、$W_0\subset Z_0$ 为 $I$ 的零点集的补集。
下面的图有助于直观理解这一情形：
$$
\xymatrix{
Z \ar[dd] & Z_0 \ar[l] \ar[d] \\
& Y_0 \ar[d] \\
X & X_0 \ar[l]
}
\quad\quad\quad
\xymatrix{
W \ar[dd] & W_0 \ar[l] \ar[d] \\
& V_0 \ar[d] \\
U & U_0 \ar[l]
}
$$
引理中的条件保证
$$
\xymatrix{
W_0 \ar[r] \ar[d] & Z_0 \ar[d] \\
V_0 \ar[r] & Y_0
}
$$
是初等特异方形，见空间的导出范畴之定义
\ref{spaces-perfect-definition-elementary-distinguished-square}。
此外，已知 $W_0\to U_0$ 与 $V_0\to U_0$ 平展。
由于假设 $J$ 幂零，态射 $X_0\subset X$ 是有限阶加厚。
由平展位点的拓扑不变性，可以找到唯一的平展概形态射
$V\to X$，使得 $V_0=V\times_XX_0$；
并且可以把给定态射 $W_0\to V_0$ 提升为唯一的 $X$-态射
$W\to V$。见平展态射之定理
\ref{etale-theorem-remarkable-equivalence}。
由于 $W_0\to V_0$ 是分离的，态射 $W\to V$ 也分离，
例如见更多态射之引理
\ref{more-morphisms-lemma-deform-property}。
由空间的推出之引理
\ref{spaces-pushouts-lemma-construct-elementary-distinguished-square}，
可在 $X$ 上代数空间的范畴中构造初等特异方形
$$
\xymatrix{
W \ar[r] \ar[d] & Z \ar[d] \\
V \ar[r] & Y
}
$$
由于初等特异方形的基变换仍为初等特异方形（空间的导出范畴之引理
\ref{spaces-perfect-lemma-make-more-elementary-distinguished-squares}），可见
$$
\xymatrix{
W_0 \ar[r] \ar[d] & Z_0 \ar[d] \\
V_0 \ar[r] & Y\times_XX_0
}
$$
是初等特异方形。
由于初等特异方形是推出（空间的推出之引理
\ref{spaces-pushouts-lemma-elementary-distinguished-square-pushout}），
可知存在唯一的同构 $Y\times_XX_0=Y_0$，
且与涉及这些空间的两个方形相容。
由空间的极限之命题
\ref{spaces-limits-proposition-affine}，可知 $Y$ 是仿射的。
写 $Y=\Spec(B)$。显然，$B$ 可嵌入所需图表，并满足要求的所有性质。
\end{proof}

\section{rig-\,'etale 代数的代数化}
\label{restricted-section-approximation-principal}

\noindent
主要目标是证明 rig-\,'etale 代数的代数化，其中不假设底层诺特环 $A$
为 G-环，且理想 $I\subset A$ 是任意的——不必是主理想。
我们先证明主理想情形，再用第 \ref{restricted-section-approximation} 节的结果完成证明。

\begin{lemma}
\label{restricted-lemma-approximate-principal}
\begin{reference}
\cite[III Theorem 7]{Elkik} 的 rig-\,'etale 情形
\end{reference}
设 $A$ 为诺特环，且 $I=(a)$ 为主理想。
设 $B$ 为 \ref{restricted-equation-C-prime} 中的对象，并且
在 $(A,I)$ 上 rig-\,'etale。
则存在有限型 $A$-代数 $C$ 以及同构
$B\cong C^\wedge$。
\end{lemma}

\begin{proof}
选取表示 $B=A[x_1,\ldots,x_r]^\wedge/J$。
由引理 \ref{restricted-lemma-equivalent-with-artin} 的 (6)，可以找到
$c\geq0$ 和 $f_1,\ldots,f_r\in J$，使得
$\det_{1\leq i,j\leq r}(\partial f_j/\partial x_i)$
在 $B$ 中整除 $a^c$，并且
$a^cJ\subset(f_1,\ldots,f_r)+J^2$。
因此，引理 \ref{restricted-lemma-approximate-presentation-rig-smooth} 适用。
这已经完成证明，不过我们想指出，在此情形下，
引理 \ref{restricted-lemma-get-morphism-general-better} 可以由容易得多的引理
\ref{restricted-lemma-get-morphism-principal} 替代。
\end{proof}

\begin{lemma}
\label{restricted-lemma-approximate}
设 $A$ 为诺特环，且 $I\subset A$ 为理想。
设 $B$ 为 \ref{restricted-equation-C-prime} 中的对象，并且
在 $(A,I)$ 上 rig-\,'etale。
则存在有限型 $A$-代数 $C$ 以及同构
$B\cong C^\wedge$。
\end{lemma}

\begin{proof}
对 $I$ 的生成元个数作归纳来证明此引理。
设 $I=(a_1,\ldots,a_t)$。若 $t=0$，则 $I=0$，无事可证。
若 $t=1$，则引理由引理
\ref{restricted-lemma-approximate-principal} 得到。假设 $t>1$。

\medskip\noindent
对任意 $m\geq1$，置 $\bar A_m=A/(a_t^m)$。
考虑 $\bar A_m$ 中的理想
$\bar I_m=(\bar a_1,\ldots,\bar a_{t-1})$。
注意，$V(I\bar A_m)=V(\bar I_m)$。
令 $B_m=B/(a_t^m)$ 为映射
$(A,I)\to(\bar A_m,\bar I_m)$ 下 $B$ 的基变换，见备注
\ref{restricted-remark-take-bar}。
由引理 \ref{restricted-lemma-base-change-rig-etale-homomorphism}，
$B_m$ 在 $(\bar A_m,\bar I_m)$ 上 rig-\,'etale。

\medskip\noindent
由对 $t$ 的归纳假设，可以找到有限型
$\bar A_m$-代数 $C_m$ 以及映射 $C_m\to B_m$，
它诱导同构 $C_m^\wedge\cong B_m$，其中完备化是关于 $\bar I_m$ 的。
由引理 \ref{restricted-lemma-rig-etale}，可以假设
$\Spec(C_m)\to\Spec(\bar A_m)$ 在
$\Spec(\bar A_m)\setminus V(\bar I_m)$ 上平展。

\medskip\noindent
我们声称，可以如上一段选择 $A_m\to C_m\to B_m$，
并且还有同构
$C_m/(a_t^{m-1})\to C_{m-1}$，
它们与给定的 $A$-代数结构及到
$B_{m-1}=B_m/(a_t^{m-1})$ 的映射相容。
首先固定一个选择 $A_1\to C_1\to B_1$。
假设已经找到具有所需性质的
$C_{m-1}\to C_{m-2}\to\ldots\to C_1$。
注意，
$C_m/(a_t^{m-1})$ 在
$\Spec(\bar A_{m-1})\setminus V(\bar I_{m-1})$ 上平展。
因此由引理 \ref{restricted-lemma-fully-faithful-etale-over-complement}，
存在平展扩张 $C_{m-1}\to C'_{m-1}$，它模
$\bar I_{m-1}$ 诱导同构；并且存在 $\bar A_{m-1}$-代数映射
$$
C_m/(a_t^{m-1})\to C'_{m-1}
$$
它在完备化上诱导同构
$B_m/(a_t^{m-1})\to B_{m-1}$。
注意，由态射之引理 \ref{morphisms-lemma-etale-permanence}，
$C_m/(a_t^{m-1})\to C'_{m-1}$ 在 $V(\bar I_{m-1})$ 的补集上平展；
而在 $V(\bar I_{m-1})$ 上，它在完备化上诱导同构，因而在那里也平展
（例如由更多态射之引理
\ref{more-morphisms-lemma-check-smoothness-on-infinitesimal-nbhds}）。
因此 $C_m/(a_t^{m-1})\to C'_{m-1}$ 平展。
由平展态射的拓扑不变性（平展态射之定理
\ref{etale-theorem-remarkable-equivalence}），
存在平展环映射 $C_m\to C'_m$，使
$C_m/(a_t^{m-1})\to C'_{m-1}$ 同构于
$C_m/(a_t^{m-1})\to C'_m/(a_t^{m-1})$。
注意，$C'_m$ 的 $\bar I_m$-进完备化等于 $C_m$ 的
$\bar I_m$-进完备化，即 $B_m$（细节略）。
对图表
$$
\xymatrix{
& C'_m \ar[r] & C'_m/(a_t^{m-1}) \\
C''_m \ar@{..>}[ru] \ar@{..>}[rr] & & C_{m-1} \ar[u] \\
& \bar A_m \ar[r] \ar[uu] \ar@{..>}[lu] & \bar A_{m-1} \ar[u]
}
$$
应用引理 \ref{restricted-lemma-lift-approximation}，可知存在 $C''_m$ 的一个“提升”，
把 $C_{m-1}$ 提升为 $\bar A_m$ 上的代数，并具有所有所需性质。

\medskip\noindent
由构造，$(C_m)$ 是主理想 $(a_t)$ 的范畴
\ref{restricted-equation-C} 中的对象。
因此逆极限 $B'=\lim C_m$ 是 $(a_t)$-进完备的 $A$-代数，
且 $B'/a_tB'$ 是 $A/(a_t)$ 上的有限型代数，见引理
\ref{restricted-lemma-topologically-finite-type}。
由构造，有 $C_m=B'/(a_t^m)$，$C_m$ 的 $I$-进完备化为 $B_m$，
而 $B'$ 的 $I$-进完备化同构于 $B$。
对每个 $m$，复形 $\NL_{C_m/A_m}$ 的有限上同调模由构造支撑在
$V(IC_m)\subset\Spec(C_m)$ 上。
因此这些模是 $I$-进完备的（因为它们被 $I$ 的某个幂湮灭）。
由于 $\NL^\wedge_{B_m/A_m}$ 是 $\NL_{C_m/A_m}$ 的 $I$-进完备化，
见引理 \ref{restricted-lemma-NL-is-completion}，映射
$\NL_{C_m/A_m}\to\NL_{B_m/A_m}$ 在上同调上诱导同构。
于是，由于 $B$ 在 $A$ 上 rig-\,'etale，可知存在 $I$ 的一个固定幂，
它对所有 $m$ 都湮灭 $\NL_{C_m/A_m}$ 的上同调，见引理
\ref{restricted-lemma-equivalent-with-artin}。
由同一引理，$B'$ 在 $(A,(a_t))$ 上 rig-\,'etale。
最后，对 $(A,(a_t))$ 上的 $B'$ 应用引理
\ref{restricted-lemma-approximate-principal}，证明完成。
\end{proof}

\begin{lemma}
\label{restricted-lemma-approximate-by-etale-over-complement}
设 $A$ 为诺特环，且 $I\subset A$ 为理想。
设 $B$ 为 $I$-进完备的 $A$-代数，且
$A/I\to B/IB$ 为有限型。引理
\ref{restricted-lemma-equivalent-with-artin} 的等价条件也等价于
\begin{enumerate}
\item[(8)]
\label{restricted-item-algebraize}
存在有限型 $A$-代数 $C$，使得
$\Spec(C)\to\Spec(A)$ 在
$\Spec(A)\setminus V(I)$ 上平展，并且
$B\cong C^\wedge$。
\end{enumerate}
\end{lemma}

\begin{proof}
结合引理 \ref{restricted-lemma-equivalent-with-artin}、
\ref{restricted-lemma-approximate} 和
\ref{restricted-lemma-rig-etale}。略去一个小细节。
\end{proof}

\section{有限型态射}
\label{restricted-section-finite-type}

\noindent
在形式空间之 \ref{formal-spaces-section-finite-type} 节中，
我们定义了形式代数空间的有限型态射。
本节研究具有有限生成定义理想的 adic 拓扑环之间相应类型的连续环映射。
强烈建议读者跳过本节。

\begin{lemma}
\label{restricted-lemma-finite-type}
设 $A$ 和 $B$ 为具有有限生成定义理想的 adic 拓扑环。
设 $\varphi:A\to B$ 为连续环同态。以下条件等价：
\begin{enumerate}
\item $\varphi$ 是 adic 的，且 $B$ 在 $A$ 上拓扑有限型；
\item $\varphi$ 是 taut 的，且 $B$ 在 $A$ 上拓扑有限型；
\item 存在定义理想 $I\subset A$，使 $B$ 上的拓扑为 $I$-进拓扑，
并且存在定义理想 $I'\subset A$，使 $A/I'\to B/I'B$ 为有限型；
\item 对所有定义理想 $I\subset A$，$B$ 上的拓扑为 $I$-进拓扑，
且 $A/I\to B/IB$ 为有限型；
\item 存在定义理想 $I\subset A$，使 $B$ 上的拓扑为 $I$-进拓扑，
且 $B$ 属于范畴 \ref{restricted-equation-C-prime}；
\item 对所有定义理想 $I\subset A$，$B$ 上的拓扑为 $I$-进拓扑，
且 $B$ 属于范畴 \ref{restricted-equation-C-prime}；
\item 作为拓扑 $A$-代数，$B$ 是
$A\{x_1,\ldots,x_r\}$ 对某个闭理想的商；
\item 作为拓扑 $A$-代数，$B$ 是
$A[x_1,\ldots,x_r]^\wedge$ 对某个闭理想的商，其中
$A[x_1,\ldots,x_r]^\wedge$ 是 $A[x_1,\ldots,x_r]$
关于 $A$ 的某个定义理想的完备化；并且
\item 在此添加更多内容。
\end{enumerate}
此外，这些等价条件定义了
$\text{WAdm}^{adic*}$ 态射的一个局部性质；该范畴定义于形式空间之备注
\ref{formal-spaces-remark-variant-adic-star}。
\end{lemma}

\begin{proof}
taut 环同态定义于形式空间之定义
\ref{formal-spaces-definition-taut}。
adic 环同态定义于形式空间之定义
\ref{formal-spaces-definition-adic-homomorphism}。
引理由形式空间之引理
\ref{formal-spaces-lemma-quotient-restricted-power-series}、
\ref{formal-spaces-lemma-quotient-restricted-power-series-admissible} 和
\ref{formal-spaces-lemma-adic-homomorphism} 合并得到。细节略。
需要明确的是，拓扑环
$A[x_1,\ldots,x_n]^\wedge$ 与
$A\{x_1,\ldots,x_r\}$ 之间没有区别，见形式空间之备注
\ref{formal-spaces-remark-I-adic-completion-and-restricted-power-series}。
\end{proof}

\begin{remark}
\label{restricted-remark-NL-well-defined-topological}
设 $A\to B$ 为 $\text{WAdm}^{adic*}$ 中的箭头，
它是 adic 且拓扑有限型的（见引理 \ref{restricted-lemma-finite-type}）。
写 $B=A\{x_1,\ldots,x_r\}/J$。则可置\footnote{事实上，
这一构造对满足形式空间之引理
\ref{formal-spaces-lemma-quotient-restricted-power-series}
中等价条件的 $\text{WAdm}^{count}$ 箭头也适用。}
$$
\NL_{B/A}^\wedge=
\left(J/J^2\longrightarrow\bigoplus B\text{d}x_i\right)
$$
与引理 \ref{restricted-lemma-NL-up-to-homotopy} 的证明完全相同，
读者可以证明，该 $B$-模复形在同伦范畴 $K(B)$ 中
（相差唯一同构）定义良好。
现在，若 $A$ 为诺特环，而 $I\subset A$ 为定义理想，
则这一构造重新得到公式 \ref{restricted-equation-NL} 在第
\ref{restricted-section-naive-cotangent-complex} 节所定义的
$B$ 在 $(A,I)$ 上的朴素余切复形，
原因只是由形式空间之备注
\ref{formal-spaces-remark-I-adic-completion-and-restricted-power-series}，
$A[x_1,\ldots,x_n]^\wedge$ 与
$A\{x_1,\ldots,x_r\}$ 相同。
特别地，仍当 $A$ 是 adic 诺特拓扑环时，对象
$\NL_{B/A}^\wedge$ 与定义理想 $I\subset A$ 的选择无关。
\end{remark}

\begin{lemma}
\label{restricted-lemma-base-change-finite-type}
考虑引理 \ref{restricted-lemma-finite-type} 所定义的
$\textit{WAdm}^{adic*}$ 箭头的性质 $P$。
则 $P$ 在形式空间之备注
\ref{formal-spaces-remark-base-change-variant-adic-star}
所定义的基变换下稳定。
\end{lemma}

\begin{proof}
陈述由引理 \ref{restricted-lemma-finite-type} 是有意义的。
为证明它，假设在 $\textit{WAdm}^{adic*}$ 中有态射
$B\to A$ 和 $B\to C$，并且作为拓扑 $B$-代数，对某个闭理想 $J$ 有
$A=B\{x_1,\ldots,x_r\}/J$。
则 $A\widehat{\otimes}_BC$ 同构于
$C\{x_1,\ldots,x_r\}/J'$ 的商，其中
$J'$ 是 $JC\{x_1,\ldots,x_r\}$ 的闭包。
略去一些细节。
\end{proof}

\begin{lemma}
\label{restricted-lemma-composition-finite-type}
考虑引理 \ref{restricted-lemma-finite-type} 所定义的
$\textit{WAdm}^{adic*}$ 箭头的性质 $P$。
则 $P$ 在形式空间之备注
\ref{formal-spaces-remark-composition-variant-adic-star}
所定义的复合下稳定。
\end{lemma}

\begin{proof}
陈述由引理 \ref{restricted-lemma-finite-type} 是有意义的。
最容易的证明方式是说明：
(a) adic 拓扑环之间 adic 环映射的复合仍为 adic；
(b) 连续环映射的复合保持拓扑有限型这一性质。
细节略。
\end{proof}

\noindent
下面的引理说明，adic* 形式代数空间的态射局部有限型，
当且仅当它们在平展局部由引理
\ref{restricted-lemma-finite-type} 所述类型的拓扑环映射给出。

\begin{lemma}
\label{restricted-lemma-finite-type-morphisms}
设 $S$ 为概形。设 $f:X\to Y$ 为 $S$ 上局部 adic* 形式代数空间的态射。
以下条件等价：
\begin{enumerate}
\item 对每个交换图表
$$
\xymatrix{
U \ar[d] \ar[r] & V \ar[d] \\
X \ar[r] & Y
}
$$
其中 $U,V$ 是仿射形式代数空间，而 $U\to X$ 与 $V\to Y$
是由代数空间可表且平展的，态射 $U\to V$ 对应于
$\textit{WAdm}^{adic*}$ 中 adic 且拓扑有限型的箭头；
\item 存在形式空间之定义
\ref{formal-spaces-definition-formal-algebraic-space}
所述的覆盖 $\{Y_j\to Y\}$，并且对每个 $j$，
存在形式空间之定义
\ref{formal-spaces-definition-formal-algebraic-space}
所述的覆盖
$\{X_{ji}\to Y_j\times_YX\}$，使每个
$X_{ji}\to Y_j$ 对应于 $\textit{WAdm}^{adic*}$ 中 adic 且拓扑有限型的箭头；
\item 存在形式空间之定义
\ref{formal-spaces-definition-formal-algebraic-space}
所述的覆盖 $\{X_i\to X\}$，并且对每个 $i$，
存在分解 $X_i\to Y_i\to Y$，其中 $Y_i$ 是仿射形式代数空间，
$Y_i\to Y$ 由代数空间可表且平展，而
$X_i\to Y_i$ 对应于 $\textit{WAdm}^{adic*}$ 中 adic 且拓扑有限型的箭头；并且
\item $f$ 局部有限型。
\end{enumerate}
\end{lemma}

\begin{proof}
立即由引理 \ref{restricted-lemma-finite-type} 中 (1) 与 (2) 的等价性，以及形式空间之引理
\ref{formal-spaces-lemma-finite-type-local-property} 得到。
\end{proof}

\section{约化上的有限型}
\label{restricted-section-finite-type-red}

\noindent
本节略谈局部可数指标化形式代数空间的态射
$X\to Y$，其中 $X_{red}\to Y_{red}$ 局部有限型。
我们会把它转化为一个代数条件。为理解这一代数条件，记住以下事实会很有用：
\begin{itemize}
\item 若 $A$ 为弱容许拓扑环，则拓扑幂零元素的集合
$\mathfrak a\subset A$ 是开根理想，并且
$\text{Spf}(A)_{red}=\Spec(A/\mathfrak a)$。
\end{itemize}
见形式空间之定义
\ref{formal-spaces-definition-weakly-admissible}、引理
\ref{formal-spaces-lemma-topologically-nilpotent} 以及例
\ref{formal-spaces-example-reduction-affine-formal-spectrum}。

\begin{lemma}
\label{restricted-lemma-finite-type-red}
对 $\text{WAdm}^{count}$ 中的箭头 $\varphi:A\to B$，考虑性质
$P(\varphi)=$“诱导的环同态
$A/\mathfrak a\to B/\mathfrak b$ 为有限型”，
其中 $\mathfrak a\subset A$ 和 $\mathfrak b\subset B$
是拓扑幂零元素的理想。
则 $P$ 是形式空间之情形
\ref{formal-spaces-situation-local-property} 所定义的局部性质。
\end{lemma}

\begin{proof}
考虑形式空间之情形
\ref{formal-spaces-situation-local-property} 中的交换图表
$$
\xymatrix{
B \ar[r] & (B')^\wedge \\
A \ar[r] \ar[u]^\varphi & (A')^\wedge \ar[u]_{\varphi'}
}
$$
对该图表取 $\text{Spf}$，得到仿射形式代数空间的图表
$$
\xymatrix{
\text{Spf}(B) \ar[d] &
\text{Spf}((B')^\wedge) \ar[l] \ar[d] \\
\text{Spf}(A) &
\text{Spf}((A')^\wedge) \ar[l]
}
$$
由形式空间之引理 \ref{formal-spaces-lemma-etale}，
其中横向箭头由代数空间可表且平展。
因此得到仿射概形的交换图表
$$
\xymatrix{
\text{Spf}(B)_{red} \ar[d]^f &
\text{Spf}((B')^\wedge)_{red} \ar[l]^g \ar[d]^{f'} \\
\text{Spf}(A)_{red} &
\text{Spf}((A')^\wedge)_{red} \ar[l]
}
$$
由形式空间之引理
\ref{formal-spaces-lemma-reduction-smooth}，其中横向箭头平展。
由形式空间之例
\ref{formal-spaces-example-reduction-affine-formal-spectrum} 和引理
\ref{formal-spaces-lemma-etale-surjective}，
形式空间之情形
\ref{formal-spaces-situation-local-property} 中的条件 (1)、(2)、(3)
转化为以下陈述：
\begin{enumerate}
\item 若 $f$ 局部有限型，则 $f'$ 局部有限型；
\item 若 $f'$ 局部有限型且 $g$ 满，则 $f$ 局部有限型；并且
\item 若 $T_i\to S$（$i=1,\ldots,n$）局部有限型，则
$\coprod_{i=1,\ldots,n}T_i\to S$ 局部有限型。
\end{enumerate}
(1) 与 (2) 来自局部有限型这一性质在平展拓扑下对源和目标均为局部的事实，
见空间的态射之 \ref{spaces-morphisms-section-finite-type} 节中的讨论。
(3) 是定义的直接推论。
\end{proof}

\begin{lemma}
\label{restricted-lemma-base-change-finite-type-red}
考虑引理 \ref{restricted-lemma-finite-type-red} 所定义的
$\textit{WAdm}^{count}$ 箭头的性质 $P$。
则 $P$ 在基变换下稳定（形式空间之情形
\ref{formal-spaces-situation-base-change-local-property}）。
\end{lemma}

\begin{proof}
陈述由引理 \ref{restricted-lemma-finite-type-red} 是有意义的。
为证明它，假设在 $\textit{WAdm}^{count}$ 中有态射
$B\to A$ 与 $B\to C$，使得
$B/\mathfrak b\to A/\mathfrak a$ 为有限型，
其中 $\mathfrak b\subset B$ 与 $\mathfrak a\subset A$
是拓扑幂零元素的理想。
由于 $A$ 和 $B$ 弱容许，理想 $\mathfrak a$ 与 $\mathfrak b$ 是开的。
令 $\mathfrak c\subset C$ 为拓扑幂零元素的（开）理想。
则得到满射
$$
A\widehat{\otimes}_BC\to
A/\mathfrak a\otimes_{B/\mathfrak b}C/\mathfrak c
$$
其核是弱定义理想，因而由拓扑幂零元素组成
（请与形式空间之引理
\ref{formal-spaces-lemma-completed-tensor-product} 的证明比较）。
由于
$C/\mathfrak c\to
A/\mathfrak a\otimes_{B/\mathfrak b}C/\mathfrak c$
已经是 $B/\mathfrak b\to A/\mathfrak a$ 的基变换，所以为有限型。
结论随即得到。
\end{proof}

\begin{lemma}
\label{restricted-lemma-composition-finite-type-red}
考虑引理 \ref{restricted-lemma-finite-type-red} 所定义的
$\textit{WAdm}^{count}$ 箭头的性质 $P$。
则 $P$ 在复合下稳定（形式空间之情形
\ref{formal-spaces-situation-composition-local-property}）。
\end{lemma}

\begin{proof}
略。提示：有限型环映射的复合仍为有限型。
\end{proof}

\begin{lemma}
\label{restricted-lemma-finite-type-finite-type-red}
设 $\varphi:A\to B$ 为 $\textit{WAdm}^{count}$ 中的箭头。
若 $\varphi$ 是 taut 且拓扑有限型的，则 $\varphi$
满足引理 \ref{restricted-lemma-finite-type-red} 所定义的条件。
\end{lemma}

\begin{proof}
这是定义的一个容易推论。
\end{proof}

\begin{lemma}
\label{restricted-lemma-Noetherian-finite-type-red}
设 $\varphi:A\to B$ 为 $\textit{WAdm}^{Noeth}$ 中的箭头，
并满足引理 \ref{restricted-lemma-finite-type-red} 所定义的条件。
则 $A\to B$ 拓扑有限型。
\end{lemma}

\begin{proof}
令 $\mathfrak b\subset B$ 为拓扑幂零元素的理想。
选取 $b_1,\ldots,b_r\in B$，它们的像生成 $B/\mathfrak b$
作为 $A$-代数。再选取理想 $\mathfrak b$ 的生成元
$b_{r+1},\ldots,b_s$。
我们声称，映射
$$
\varphi:A[x_1,\ldots,x_s]\longrightarrow B,\quad
x_i\longmapsto b_i
$$
的像稠密。事实上，若对某个 $n\geq0$ 有 $b\in\mathfrak b^n$，
则可对多重指标
$E=(e_{r+1},\ldots,e_s)$ 写成
$$
b=\sum b_Eb_{r+1}^{e_{r+1}}\ldots b_s^{e_s},
|E|=\sum e_i=n, b_E\in B
$$
接着，可以写
$b_E=f_E(b_1,\ldots,b_r)+b'_E$，其中
$b'_E\in\mathfrak b$ 且
$f_E\in A[x_1,\ldots,x_r]$。
合起来得到 $b\in\Im(\varphi)+\mathfrak b^{n+1}$。
由归纳可见，对所有 $n\geq0$，
$B=\Im(\varphi)+\mathfrak b^n$；由于 $\mathfrak b$ 是 $B$ 的定义理想，
这就蕴含所需结论。
\end{proof}

\begin{lemma}
\label{restricted-lemma-Noetherian-adic-finite-type-red}
设 $\varphi:A\to B$ 为 $\textit{WAdm}^{Noeth}$ 中的箭头。
若 $\varphi$ 是 adic 的，则以下条件等价：
\begin{enumerate}
\item $\varphi$ 满足引理
\ref{restricted-lemma-finite-type-red} 所定义的条件；并且
\item $\varphi$ 满足引理
\ref{restricted-lemma-finite-type} 所定义的条件。
\end{enumerate}
\end{lemma}

\begin{proof}
略。提示：为证明 (1) $\Rightarrow$ (2)，使用引理
\ref{restricted-lemma-Noetherian-finite-type-red}。
\end{proof}

\begin{lemma}
\label{restricted-lemma-finite-type-red-morphisms}
设 $S$ 为概形。设 $f:X\to Y$ 为 $S$ 上局部可数指标化形式代数空间的态射。
以下条件等价：
\begin{enumerate}
\item 对每个交换图表
$$
\xymatrix{
U \ar[d] \ar[r] & V \ar[d] \\
X \ar[r] & Y
}
$$
其中 $U,V$ 是仿射形式代数空间，而 $U\to X$ 与 $V\to Y$
由代数空间可表且平展，态射 $U\to V$ 对应于
$\textit{WAdm}^{count}$ 中满足引理
\ref{restricted-lemma-finite-type-red} 所定义性质的箭头；
\item 存在形式空间之定义
\ref{formal-spaces-definition-formal-algebraic-space}
所述的覆盖 $\{Y_j\to Y\}$，并且对每个 $j$，
存在形式空间之定义
\ref{formal-spaces-definition-formal-algebraic-space}
所述的覆盖 $\{X_{ji}\to Y_j\times_YX\}$，
使每个 $X_{ji}\to Y_j$ 对应于
$\textit{WAdm}^{count}$ 中满足引理
\ref{restricted-lemma-finite-type-red} 所定义性质的箭头；
\item 存在形式空间之定义
\ref{formal-spaces-definition-formal-algebraic-space}
所述的覆盖 $\{X_i\to X\}$，并且对每个 $i$，
存在分解 $X_i\to Y_i\to Y$，其中 $Y_i$ 是仿射形式代数空间，
$Y_i\to Y$ 由代数空间可表且平展，而
$X_i\to Y_i$ 对应于
$\textit{WAdm}^{count}$ 中满足引理
\ref{restricted-lemma-finite-type-red} 所定义性质的箭头；并且
\item 态射 $f_{red}:X_{red}\to Y_{red}$ 局部有限型。
\end{enumerate}
\end{lemma}

\begin{proof}
(1)、(2)、(3) 的等价性由引理
\ref{restricted-lemma-finite-type-red} 以及形式空间之引理
\ref{formal-spaces-lemma-property-defines-property-morphisms} 得到。
令 $Y_j$ 与 $X_{ji}$ 如 (2)。则
\begin{itemize}
\item 族
$\{Y_{j,red}\to Y_{red}\}$ 与
$\{X_{ji,red}\to X_{red}\}$ 是仿射概形的平展覆盖。
这由形式空间之引理
\ref{formal-spaces-lemma-reduction-formal-algebraic-space}
的证明中的讨论得到，也可直接由形式空间之引理
\ref{formal-spaces-lemma-reduction-smooth} 得到。
\item 若 $X_{ji}\to Y_j$ 对应于
$\textit{WAdm}^{count}$ 中的态射 $B_j\to A_{ji}$，
则 $X_{ji,red}\to Y_{j,red}$ 对应于环映射
$B_j/\mathfrak b_j\to A_{ji}/\mathfrak a_{ji}$，
其中 $\mathfrak b_j$ 和 $\mathfrak a_{ji}$ 是拓扑幂零元素的理想。
这由形式空间之例
\ref{formal-spaces-example-reduction-affine-formal-spectrum} 得到。
因此，$X_{ji,red}\to Y_{j,red}$ 局部有限型，
当且仅当 $B_j\to A_{ji}$ 满足引理
\ref{restricted-lemma-finite-type-red} 所定义的性质。
\end{itemize}
由这些说明，(2) 与 (4) 等价，因为局部有限型是代数空间态射的一个性质，
且在平展拓扑下对源和目标均为局部的，见空间的态射之
\ref{spaces-morphisms-section-finite-type} 节中的讨论。
\end{proof}

\section{平坦态射}
\label{restricted-section-flat}

\noindent
本节定义局部诺特形式代数空间的平坦态射。

\begin{lemma}
\label{restricted-lemma-flat-axioms}
$\textit{WAdm}^{Noeth}$ 箭头上的性质
$P(\varphi)=$“$\varphi$ 平坦”是形式空间之备注
\ref{formal-spaces-remark-variant-Noetherian} 所定义的局部性质。
\end{lemma}

\begin{proof}
回忆该陈述的含义。首先，$\textit{WAdm}^{Noeth}$ 是以下范畴：
对象为 adic 诺特拓扑环，态射为连续环同态。
考虑满足以下条件的交换图表
$$
\xymatrix{
B \ar[r] & (B')^\wedge \\
A \ar[r] \ar[u]^\varphi & (A')^\wedge \ar[u]_{\varphi'}
}
$$
$A$ 与 $B$ 是 adic 诺特拓扑环，
$A\to A'$ 与 $B\to B'$ 是平展环映射，
对某个定义理想 $I\subset A$ 有
$(A')^\wedge=\lim A'/I^nA'$，
对某个定义理想 $J\subset B$ 有
$(B')^\wedge=\lim B'/J^nB'$，
并且 $\varphi:A\to B$ 与
$\varphi':(A')^\wedge\to(B')^\wedge$ 连续。
由形式空间之引理
\ref{formal-spaces-lemma-completion-in-sub}，
$(A')^\wedge$ 与 $(B')^\wedge$ 是 adic 诺特拓扑环。
需要证明
\begin{enumerate}
\item $\varphi$ 平坦 $\Rightarrow\varphi'$ 平坦；
\item 若 $B\to B'$ 忠实平坦，则
$\varphi'$ 平坦 $\Rightarrow\varphi$ 平坦；并且
\item 若对 $i=1,\ldots,n$，$A\to B_i$ 平坦，则
$A\to\prod_{i=1,\ldots,n}B_i$ 平坦。
\end{enumerate}
我们将不再说明地使用诺特环的完备化平坦这一事实
（代数，引理 \ref{algebra-lemma-completion-flat}）。
当然，$A\to A'$ 和 $B\to B'$ 平坦，
因此特别地，图表中的横向箭头平坦。

\medskip\noindent
(1) 的证明。若 $\varphi$ 平坦，则复合
$A\to(A')^\wedge\to(B')^\wedge$ 平坦。
因此由平坦性续篇之引理
\ref{flat-lemma-etale-flat-up-down}，
$A'\to(B')^\wedge$ 平坦。
再对 $R=A'$、理想 $I(A')$ 以及
$M=(B')^\wedge=M^\wedge$ 应用更多代数之引理
\ref{more-algebra-lemma-flat-after-completion}，
可见 $(A')^\wedge\to(B')^\wedge$ 平坦。

\medskip\noindent
(2) 的证明。假设 $\varphi'$ 平坦且 $B\to B'$ 忠实平坦。
则复合 $A\to(A')^\wedge\to(B')^\wedge$ 平坦。
又由形式空间之引理
\ref{formal-spaces-lemma-etale-surjective}，
$B\to(B')^\wedge$ 忠实平坦。
因此由代数之引理
\ref{algebra-lemma-flatness-descends-more-general}，
$\varphi:A\to B$ 平坦。

\medskip\noindent
(3) 的证明略。
\end{proof}

\begin{lemma}
\label{restricted-lemma-base-change-flat-continuous}
记 $P$ 为引理 \ref{restricted-lemma-flat-axioms} 所定义的
$\textit{WAdm}^{Noeth}$ 箭头的性质。
记 $Q$ 为引理 \ref{restricted-lemma-finite-type-red} 所定义的性质，
把它视为 $\textit{WAdm}^{Noeth}$ 箭头的性质。
记 $R$ 为引理 \ref{restricted-lemma-finite-type} 所定义的性质，
把它视为 $\textit{WAdm}^{Noeth}$ 箭头的性质。则
\begin{enumerate}
\item $P$ 在由 $Q$ 作基变换下稳定
（形式空间之备注
\ref{formal-spaces-remark-base-change-variant-variant-Noetherian}）；并且
\item $P+R$ 在基变换下稳定
（形式空间之备注
\ref{formal-spaces-remark-base-change-variant-Noetherian}）。
\end{enumerate}
\end{lemma}

\begin{proof}
陈述是有意义的，因为 $P,Q,R$ 都是
$\textit{WAdm}^{Noeth}$ 态射的局部性质。
设 $\varphi:B\to A$ 与 $\psi:B\to C$ 为
$\textit{WAdm}^{Noeth}$ 的态射。
若 $Q(\varphi)$ 或 $Q(\psi)$ 成立，则由形式空间之引理
\ref{formal-spaces-lemma-completed-tensor-product}，
$A\widehat{\otimes}_BC$ 是诺特环。
由于 $R$ 蕴含 $Q$（引理
\ref{restricted-lemma-finite-type-finite-type-red}），
可知在 (1)、(2) 两种情形下都如此。
这是首先必须核对的事项。还需证明
$C\to A\widehat{\otimes}_BC$ 平坦。

\medskip\noindent
(1) 的证明。固定定义理想 $I\subset A$ 与 $J\subset B$。
由引理 \ref{restricted-lemma-Noetherian-finite-type-red}，
环映射 $B\to C$ 拓扑有限型。
因此，对所有 $n\geq1$，$B\to C/J^n$ 为有限型。
于是环 $A\otimes_BC/J^n$ 是诺特的
（因为它在 $A$ 上有限型，且 $A$ 为诺特环）。
所以 $A\otimes_BC/J^n$ 的 $I$-进完备化
$A\widehat{\otimes}_BC/J^n$ 在 $C/J^n$ 上平坦：
$C/J^n\to A\otimes_BC/J^n$ 是 $B\to A$ 的基变换，因而平坦；
而由代数之引理 \ref{algebra-lemma-completion-flat}，
$A\otimes_BC/J^n\to A\widehat{\otimes}_BC/J^n$ 平坦。
注意
$$
A\widehat{\otimes}_BC/J^n=
(A\widehat{\otimes}_BC)/J^n(A\widehat{\otimes}_BC)
$$
细节略。由此，$M=A\widehat{\otimes}_BC$ 是 $C$-模，
关于 $J$-进拓扑完备，并且对所有 $n\geq1$，
$M/J^nM$ 在 $C/J^n$ 上平坦。
由更多代数之引理
\ref{more-algebra-lemma-limit-flat}，这蕴含 $M$ 在 $C$ 上平坦。

\medskip\noindent
(2) 的证明。此时 $B\to A$ 是 adic 的，因而只有
$$
A\widehat{\otimes}_BC=\lim A\otimes_BC/J^n
$$
对形式空间之引理
\ref{formal-spaces-lemma-completed-tensor-product}
应用 $C/J^n$ 代替 $C$，可知环
$A\otimes_BC/J^n$ 是诺特的。
与此前相同地得出结论。
\end{proof}

\begin{lemma}
\label{restricted-lemma-composition-flat-continuous}
记 $P$ 为引理 \ref{restricted-lemma-flat-axioms} 所定义的
$\textit{WAdm}^{Noeth}$ 箭头的性质。
则 $P$ 在复合下稳定（形式空间之备注
\ref{formal-spaces-remark-composition-variant-Noetherian}）。
\end{lemma}

\begin{proof}
这是因为平坦环映射的复合仍平坦。
\end{proof}

\begin{definition}
\label{restricted-definition-flat}
设 $S$ 为概形。设 $f:X\to Y$ 为 $S$ 上局部诺特形式代数空间的态射。
若对每个交换图表
$$
\xymatrix{
U \ar[d] \ar[r] & V \ar[d] \\
X \ar[r] & Y
}
$$
其中 $U,V$ 是仿射形式代数空间，而 $U\to X$ 与 $V\to Y$
由代数空间可表且平展，态射 $U\to V$
对应于 adic 诺特拓扑环之间的平坦映射，
则称 $f$ {\it 平坦}。
\end{definition}

\noindent
下面证明该条件可以在源和目标上平展局部地检验。

\begin{lemma}
\label{restricted-lemma-flat-morphisms}
设 $S$ 为概形。设 $f:X\to Y$ 为 $S$ 上局部诺特形式代数空间的态射。
以下条件等价：
\begin{enumerate}
\item $f$ 平坦；
\item 对每个交换图表
$$
\xymatrix{
U \ar[d] \ar[r] & V \ar[d] \\
X \ar[r] & Y
}
$$
其中 $U,V$ 是仿射形式代数空间，而 $U\to X$ 与 $V\to Y$
由代数空间可表且平展，态射 $U\to V$
对应于 $\textit{WAdm}^{Noeth}$ 中的平坦映射；
\item 存在形式空间之定义
\ref{formal-spaces-definition-formal-algebraic-space}
所述的覆盖 $\{Y_j\to Y\}$，并且对每个 $j$，
存在形式空间之定义
\ref{formal-spaces-definition-formal-algebraic-space}
所述的覆盖 $\{X_{ji}\to Y_j\times_YX\}$，
使每个 $X_{ji}\to Y_j$ 对应于
$\textit{WAdm}^{Noeth}$ 中的平坦映射；并且
\item 存在形式空间之定义
\ref{formal-spaces-definition-formal-algebraic-space}
所述的覆盖 $\{X_i\to X\}$，并且对每个 $i$，
存在分解 $X_i\to Y_i\to Y$，其中 $Y_i$ 是仿射形式代数空间，
$Y_i\to Y$ 由代数空间可表且平展，而
$X_i\to Y_i$ 对应于
$\textit{WAdm}^{Noeth}$ 中的平坦映射。
\end{enumerate}
\end{lemma}

\begin{proof}
(1) 与 (2) 的等价性就是定义 \ref{restricted-definition-flat}。
(2)、(3)、(4) 的等价性来自以下事实：
由引理 \ref{restricted-lemma-flat-axioms}，平坦是
$\text{WAdm}^{Noeth}$ 箭头的局部性质；
再应用形式空间之引理
\ref{formal-spaces-lemma-property-defines-property-morphisms}
对局部诺特代数空间之间态射的变体，
该变体在形式空间之备注
\ref{formal-spaces-remark-variant-Noetherian} 中提及。
\end{proof}

\begin{lemma}
\label{restricted-lemma-base-change-flat}
设 $S$ 为概形。设 $f:X\to Y$ 与 $g:Z\to Y$
为 $S$ 上局部诺特形式代数空间的态射。
\begin{enumerate}
\item 若 $f$ 平坦，且
$g_{red}:Z_{red}\to Y_{red}$ 局部有限型，
则基变换 $X\times_YZ\to Z$ 平坦。
\item 若 $f$ 平坦且局部有限型，则基变换
$X\times_YZ\to Z$ 平坦且局部有限型。
\end{enumerate}
\end{lemma}

\begin{proof}
(1) 由形式空间之备注
\ref{formal-spaces-remark-base-change-variant-variant-Noetherian}、
引理 \ref{restricted-lemma-base-change-flat-continuous} 的 (1)、
引理 \ref{restricted-lemma-flat-morphisms} 以及
引理 \ref{restricted-lemma-finite-type-red-morphisms} 合并得到。

\medskip\noindent
(2) 由形式空间之备注
\ref{formal-spaces-remark-base-change-variant-Noetherian}、
引理 \ref{restricted-lemma-base-change-flat-continuous} 的 (2)、
引理 \ref{restricted-lemma-flat-morphisms} 以及
引理 \ref{restricted-lemma-finite-type-morphisms} 合并得到。
\end{proof}

\begin{lemma}
\label{restricted-lemma-composition-flat}
设 $S$ 为概形。设 $f:X\to Y$ 与 $g:Y\to Z$
为 $S$ 上局部诺特形式代数空间的态射。
若 $f$ 与 $g$ 平坦，则 $g\circ f$ 也平坦。
\end{lemma}

\begin{proof}
结合形式空间之备注
\ref{formal-spaces-remark-composition-variant-Noetherian}
与引理 \ref{restricted-lemma-composition-flat-continuous}。
\end{proof}

\begin{lemma}
\label{restricted-lemma-representable-flat}
设 $S$ 为概形。设 $f:X\to Y$ 为 $S$ 上局部诺特形式代数空间的态射。
若 $f$ 由代数空间可表，并且在 Bootstrap 之定义
\ref{bootstrap-definition-property-transformation} 的意义下平坦，
则 $f$ 在定义 \ref{restricted-definition-flat} 的意义下平坦。
\end{lemma}

\begin{proof}
这是一个健全性检验，其证明本应平凡，但并不完全如此。
我们力劝读者跳过此证明。
假设 $f$ 由代数空间可表，并且在 Bootstrap 之定义
\ref{bootstrap-definition-property-transformation} 的意义下平坦。
考虑交换图表
$$
\xymatrix{
U \ar[d] \ar[r] & V \ar[d] \\
X \ar[r] & Y
}
$$
其中 $U,V$ 是仿射形式代数空间，而 $U\to X$ 与 $V\to Y$
由代数空间可表且平展。
则由形式空间之引理
\ref{formal-spaces-lemma-representable-local-property}，
态射 $U\to V$ 对应于
$\textit{WAdm}^{Noeth}$ 中的 taut 映射 $B\to A$。
注意，这意味着 $B\to A$ 是 adic 的（形式空间，引理
\ref{formal-spaces-lemma-adic-homomorphism}）。
特别地，对任意定义理想 $J\subset B$，
$A$ 上的拓扑为 $J$-进拓扑，并且图表
$$
\xymatrix{
\Spec(A/J^nA) \ar[r] \ar[d] & \Spec(B/J^n) \ar[d] \\
U \ar[r] & V
}
$$
是笛卡尔的。

\medskip\noindent
设 $T\to V$ 为态射，其中 $T$ 是概形。则
\begin{align*}
X\times_YT\to T\text{ 平坦}
&\Rightarrow
U\times_YT\to T\text{ 平坦}\\
&\Rightarrow
U\times_VV\times_YT\to T\text{ 平坦}\\
&\Rightarrow
U\times_VV\times_YT\to V\times_YT\text{ 平坦}\\
&\Rightarrow
U\times_VT\to T\text{ 平坦}
\end{align*}
第一个陈述是对 $f$ 的假设。
第一个蕴含成立，因为 $U\to X$ 平展，因而平坦，并且代数空间平坦态射的复合仍平坦。
第二个蕴含成立，因为
$U\times_YT=U\times_VV\times_YT$。
第三个蕴含由平坦性续篇之引理
\ref{flat-lemma-etale-flat-up-down} 得到。
第四个蕴含成立，因为可以由态射
$T\to V\times_YT$ 拉回。
由此可知，$U\to V$ 在 Bootstrap 之定义
\ref{bootstrap-definition-property-transformation} 的意义下平坦。
就连续环映射 $B\to A$ 而言，这意味着环映射
$B/J^n\to A/J^nA$ 平坦（见上图）。

\medskip\noindent
最后，例如由更多代数之引理
\ref{more-algebra-lemma-limit-flat}，可知 $B\to A$ 平坦。
\end{proof}

\section{rig-闭点}
\label{restricted-section-rig-points}

\noindent
我们只发展足够的理论，以便在后文中用它检验 rig-平坦性。
读者可在 \cite{BL-I} 中找到更多理论；该文除其他内容外，
还讨论了局部诺特形式概形的情形。

\begin{lemma}
\label{restricted-lemma-rig-point}
设 $A$ 为诺特 adic 拓扑环。设 $\mathfrak q\subset A$ 为素理想。
以下条件等价：
\begin{enumerate}
\item 对某个定义理想 $I\subset A$，
$I\not\subset\mathfrak q$，且 $\mathfrak q$ 对这一性质极大；
\item 对某个定义理想 $I\subset A$，
素理想 $\mathfrak q$ 定义了
$\Spec(A)\setminus V(I)$ 的一个闭点；
\item 对任意定义理想 $I\subset A$，
$I\not\subset\mathfrak q$，且 $\mathfrak q$ 对这一性质极大；
\item 对任意定义理想 $I\subset A$，
素理想 $\mathfrak q$ 定义了
$\Spec(A)\setminus V(I)$ 的一个闭点；
\item $\dim(A/\mathfrak q)=1$，且对某个定义理想
$I\subset A$，有 $I\not\subset\mathfrak q$；
\item $\dim(A/\mathfrak q)=1$，且对任意定义理想
$I\subset A$，有 $I\not\subset\mathfrak q$；
\item $\dim(A/\mathfrak q)=1$，且 $A/\mathfrak q$ 上诱导的拓扑非平凡；
\item $A/\mathfrak q$ 是 $1$ 维诺特完备局部整环，
其极大理想是 $A$ 的任意定义理想之像的根；并且
\item 在此添加更多内容。
\end{enumerate}
\end{lemma}

\begin{proof}
显然，(1) 与 (2) 等价；同理，(3) 与 (4) 等价。
由于 $V(I)$ 与 $A$ 的定义理想 $I$ 的选择无关，可知 (2) 与 (4) 等价。

\medskip\noindent
假设等价条件 (1)--(4) 成立。
若 $\dim(A/\mathfrak q)>1$，则可以选择极大理想
$\mathfrak q\subset\mathfrak m\subset A$，使得
$\dim((A/\mathfrak q)_\mathfrak m)>1$。
于是由代数之引理
\ref{algebra-lemma-Noetherian-local-domain-dim-2-infinite-opens}，
$$
\Spec((A/\mathfrak q)_\mathfrak m)
-V(I(A/\mathfrak q)_\mathfrak m)
$$
将是无限的。这与 $\mathfrak q$ 在
$\Spec(A)\setminus V(I)$ 中闭矛盾。
所以 (6) 成立。显然 (6) 蕴含 (5)。

\medskip\noindent
反之，假设 (5) 成立。设 $I\subset A$ 为定义理想。
由于 $A/\mathfrak q$ 关于 $I(A/\mathfrak q)$ 完备
（例如由代数之引理 \ref{algebra-lemma-completion-tensor}），
由代数之引理 \ref{algebra-lemma-radical-completion}，
$\Spec(A/\mathfrak q)$ 的所有闭点都包含在
$V(IA/\mathfrak q)$ 中。
由于 $\dim(A/\mathfrak q)=1$ 且
$I\not\subset\mathfrak q$，得到两点：
(a) $V(IA/\mathfrak q)$ 必须包含
$\Spec(A/\mathfrak q)$ 中除泛点以外的所有点；
(b) $V(IA/\mathfrak q)$ 必须是（有限）离散集。
由 (a)，$\mathfrak q$ 是
$\Spec(A)\setminus V(I)$ 的闭点，故 (2) 成立。

\medskip\noindent
继续假设 (5)。例如由更多代数之引理
\ref{more-algebra-lemma-irreducible-henselian-pair-connected}
（以及完备对是 Hensel 对这一事实，见更多代数之引理
\ref{more-algebra-lemma-complete-henselian}），
有限离散空间 $V(IA/\mathfrak q)$ 必为单点。
故 (8) 成立。反之，显然 (8) 蕴含 (5)。

\medskip\noindent
至此已知 (1)--(6) 与 (8) 等价。
略去它们也与 (7) 等价的验证。
\end{proof}

\noindent
为便于讨论这类素理想，引入如下非标准记号。

\begin{definition}
\label{restricted-definition-rig-closed}
设 $A$ 为诺特 adic 拓扑环。设 $\mathfrak q\subset A$ 为素理想。
若引理 \ref{restricted-lemma-rig-point} 的等价条件成立，
则称 $\mathfrak q$ {\it rig-闭}。
\end{definition}

\noindent
我们将需要几个引理；它们实质上说明，即使在相对情形下，
也存在很多 rig-闭素理想。

\begin{lemma}
\label{restricted-lemma-rig-closed-point-relative-residue-field}
设 $\varphi:A\to B$ 为 $\textit{WAdm}^{Noeth}$ 中的箭头。
记 $\mathfrak a\subset A$ 与 $\mathfrak b\subset B$
为拓扑幂零元素的理想。假设
$A/\mathfrak a\to B/\mathfrak b$ 为有限型。
设 $\mathfrak q\subset B$ 为 rig-闭素理想。
则局部环 $B/\mathfrak q$ 的剩余域 $\kappa$
是有限型 $A/\mathfrak a$-代数。
\end{lemma}

\begin{proof}
设 $\mathfrak q\subset\mathfrak m\subset B$
为包含 $\mathfrak q$ 的唯一极大理想。
则 $\mathfrak b\subset\mathfrak m$。因此
$A/\mathfrak a\to B/\mathfrak b\to B/\mathfrak m=\kappa$
为有限型。
\end{proof}

\begin{lemma}
\label{restricted-lemma-rig-closed-point-relative}
设 $\varphi:A\to B$ 为 $\textit{WAdm}^{Noeth}$ 中 adic 且拓扑有限型的箭头。
设 $\mathfrak q\subset B$ 为 rig-闭素理想。
令 $\mathfrak p=\varphi^{-1}(\mathfrak q)\subset A$。
令 $\mathfrak a\subset A$ 为拓扑幂零元素的理想。
以下条件等价：
\begin{enumerate}
\item $B/\mathfrak q$ 的剩余域 $\kappa$ 在
$A/\mathfrak a$ 上有限；
\item $\mathfrak p\subset A$ 是 rig-闭素理想；
\item $A/\mathfrak p\subset B/\mathfrak q$ 是有限环扩张。
\end{enumerate}
\end{lemma}

\begin{proof}
假设 (1)。回忆 $B/\mathfrak q$ 是 $1$ 维诺特局部环，
其拓扑是由极大理想给出的 adic 拓扑。
由于 $\varphi$ 是 adic 的，$A\to B/\mathfrak q$ 是 adic 的。
因此 $\varphi(\mathfrak a)$ 是 $B/\mathfrak q$ 中的非零理想。
所以 $B/\mathfrak q+\varphi(\mathfrak a)$ 具有有限长度。
再由假设，$B/\mathfrak q+\varphi(\mathfrak a)$
作为 $A/\mathfrak a$-模是有限的。
故由代数之引理
\ref{algebra-lemma-finite-over-complete-ring}，
$B/\mathfrak q$ 在 $A$ 上有限。因此 (3) 成立。

\medskip\noindent
假设 (3)。由代数之引理
\ref{algebra-lemma-integral-overring-surjective}，
$\Spec(B/\mathfrak q)\to\Spec(A/\mathfrak p)$ 满。
这蕴含 (2)。

\medskip\noindent
假设 (2)。记 $\kappa'$ 为 $A/\mathfrak p$ 的剩余域。
由引理 \ref{restricted-lemma-rig-closed-point-relative-residue-field}
（以及引理 \ref{restricted-lemma-finite-type-finite-type-red}），
扩张 $\kappa/\kappa'$ 作为代数是有限生成的。
由 Hilbert 零点定理（代数，引理
\ref{algebra-lemma-field-finite-type-over-domain}），
$\kappa/\kappa'$ 是有限扩张。因此 (1) 成立。
\end{proof}

\begin{lemma}
\label{restricted-lemma-rig-closed-jacobson}
设 $\varphi:A\to B$ 为 $\textit{WAdm}^{Noeth}$ 中 adic 且拓扑有限型的箭头。
设 $\mathfrak q\subset B$ 为 rig-闭素理想。
若对某个定义理想 $I\subset A$，$A/I$ 是 Jacobson 环，
则 $\mathfrak p=\varphi^{-1}(\mathfrak q)\subset A$ 是 rig-闭素理想。
\end{lemma}

\begin{proof}
由引理 \ref{restricted-lemma-rig-closed-point-relative-residue-field}
（与引理 \ref{restricted-lemma-finite-type-finite-type-red} 合用），
$B/\mathfrak q$ 的剩余域 $\kappa$ 在
$A/\mathfrak a$ 上有限型。
由于 $A/\mathfrak a$ 是 Jacobson 环，由代数之引理
\ref{algebra-lemma-silly-jacobson}，
$\kappa$ 在 $A/\mathfrak a$ 上有限。
结论由引理 \ref{restricted-lemma-rig-closed-point-relative} 得到。
\end{proof}

\begin{lemma}
\label{restricted-lemma-rig-closed-point-in-image}
设 $\varphi:A\to B$ 为 $\textit{WAdm}^{Noeth}$ 中 adic 且拓扑有限型的箭头。
设 $\mathfrak p\subset A$ 为 rig-闭素理想。
设 $\mathfrak a\subset A$ 与 $\mathfrak b\subset B$
为拓扑幂零元素的理想。若 $\varphi$ 平坦，则以下条件等价：
\begin{enumerate}
\item $A/\mathfrak p$ 的极大理想在映射
$\Spec(B/\mathfrak b)\to\Spec(A/\mathfrak a)$ 的像中；
\item 存在 rig-闭素理想 $\mathfrak q\subset B$，使得
$\mathfrak p=\varphi^{-1}(\mathfrak q)$。
\end{enumerate}
若这些条件成立，则 $\varphi$、$\mathfrak p$、$\mathfrak q$
满足引理 \ref{restricted-lemma-rig-closed-point-relative} 的结论。
\end{lemma}

\begin{proof}
(2) $\Rightarrow$ (1) 是立即的。假设 (1)。
为证明 $\mathfrak q$ 的存在性，可以用
$A/\mathfrak p$ 代替 $A$，并用
$B/\mathfrak pB$ 代替 $B$（略去一些细节）。
因此可假设 $(A,\mathfrak m,\kappa)$ 是 $1$ 维诺特完备局部环，
$\mathfrak m=\mathfrak a$，且 $\mathfrak p=(0)$。
条件 (1) 只是说
$$
B_0=B\otimes_A\kappa=B/\mathfrak mB=B/\mathfrak aB
$$
非零。注意，$B_0$ 在 $\kappa$ 上有限型。
因此可以对 $\dim(B_0)$ 作归纳。
若 $\dim(B_0)=0$，则任意极小素理想 $\mathfrak q\subset B$ 都可用
（$A\to B$ 的平坦性保证 $\mathfrak q$ 位于
$\mathfrak p=(0)$ 上方）。
若 $\dim(B_0)>0$，则可以找到 $b\in B$，其像
$b_0\in B_0$ 是非零因子且不是单位，见代数之引理
\ref{algebra-lemma-dim-not-zero-exists-nonzerodivisor-nonunit}。
由代数之引理 \ref{algebra-lemma-grothendieck}，
环 $B'=B/bB$ 在 $A$ 上平坦。
由于 $B'_0=B'\otimes_A\kappa=B_0/(b_0)$ 非零，
可对 $B'$ 应用归纳假设而得出结论。
引理的最后一个陈述显然由引理
\ref{restricted-lemma-rig-closed-point-relative} 得到。
\end{proof}

\noindent
引入一些记号。

\begin{definition}
\label{restricted-definition-completed-principal-localization}
设 $A$ 为具有有限生成定义理想的 adic 拓扑环。设 $f\in A$。
$A$ 的{\it 完备化主局部化} $A_{\{f\}}$ 定义为：
把 $A$ 在 $f$ 处的主局部化 $A_f=A[1/f]$
关于 $A$ 的任意定义理想作完备化。
\end{definition}

\noindent
需要说明的是，若 $f$ 拓扑幂零，则 $A_{\{f\}}$ 为零环。

\begin{lemma}
\label{restricted-lemma-rig-closed-point-in-localization}
设 $A$ 为 adic 诺特拓扑环。
设 $\mathfrak p\subset A$ 为素理想。
设 $f\in A$ 的像在 $A/\mathfrak p$ 中为单位。则
$$
\mathfrak pA_{\{f\}}=
\mathfrak p(A_f)^\wedge=
\mathfrak p\otimes_A(A_f)^\wedge=
(\mathfrak p_f)^\wedge
$$
是素理想，其商为
$$
A/\mathfrak p=(A/\mathfrak p)\otimes_A(A_f)^\wedge=
(A_f)^\wedge/\mathfrak p(A_f)^\wedge
=A_{\{f\}}/\mathfrak pA_{\{f\}}
$$
\end{lemma}

\begin{proof}
由于 $A_f$ 诺特，环映射
$A\to A_f\to(A_f)^\wedge$ 平坦。
对任意有限 $A$-模 $M$，
$M\otimes_A(A_f)^\wedge$ 是 $M_f$ 的完备化。
若 $f$ 在 $M$ 上为单位，则 $M_f=M$ 已经完备。
见代数之 \ref{algebra-section-completion-noetherian} 节中的讨论。
这些观察容易推出结论。
\end{proof}

\begin{lemma}
\label{restricted-lemma-rig-closed-point-after-localization}
设 $\varphi:A\to B$ 为 $\textit{WAdm}^{Noeth}$ 中 adic 且拓扑有限型的箭头。
设 $\mathfrak q\subset B$ 为 rig-闭素理想。
存在 $f\in A$，其像在 $B/\mathfrak q$ 中为单位，并得到图表
$$
\vcenter{
\xymatrix{
B \ar[r] &
B_{\{f\}} \\
A \ar[r] \ar[u]_\varphi &
A_{\{f\}} \ar[u]_{\varphi_{\{f\}}}
}
}
\quad\text{以及素理想}\quad
\vcenter{
\xymatrix{
\mathfrak q \ar@{-}[r] \ar@{-}[d] &
\mathfrak q' \ar@{-}[d] \ar@{=}[r] &
\mathfrak qB_{\{f\}} \\
\mathfrak p \ar@{-}[r] &
\mathfrak p'
}
}
$$
使得 $\mathfrak p'$ 是 rig-闭的；也就是说，映射
$A_{\{f\}}\to B_{\{f\}}$ 与素理想
$\mathfrak q'$、$\mathfrak p'$ 满足引理
\ref{restricted-lemma-rig-closed-point-relative} 的等价条件。
\end{lemma}

\begin{proof}
$\mathfrak q'$ 的描述见引理
\ref{restricted-lemma-rig-closed-point-in-localization}。
本引理唯一的断言是：可以适当选择 $f$，使素理想
$\mathfrak p'$ 具有性质
$\dim((A_f)^\wedge/\mathfrak p')=1$。
由引理 \ref{restricted-lemma-rig-closed-point-relative}，
这又只是说
$B/\mathfrak q=(B_f)^\wedge/\mathfrak q'$
的剩余域 $\kappa$ 在
$(A_f)^\wedge/\mathfrak a'=(A/\mathfrak a)_f$ 上有限。
由引理 \ref{restricted-lemma-rig-closed-point-relative-residue-field}，
已知 $A/\mathfrak a\to\kappa$ 是有限型代数同态。
由 Hilbert 零点定理的如下形式——代数之引理
\ref{algebra-lemma-field-finite-type-over-domain}，
可以找到 $f\in A$，其像在 $\kappa$ 中为单位，
使得 $\kappa$ 在 $A_f$ 上有限。证明完毕。
\end{proof}

\begin{lemma}
\label{restricted-lemma-rig-closed-point-variables}
设 $A$ 为诺特 adic 拓扑环。记
$A\{x_1,\ldots,x_n\}$ 为 $A$ 上的受限幂级数环。
设 $\mathfrak q\subset A\{x_1,\ldots,x_n\}$ 为素理想。
置
$\mathfrak q'=A[x_1,\ldots,x_n]\cap\mathfrak q$ 以及
$\mathfrak p=A\cap\mathfrak q$。
若 $\mathfrak q$ 与 $\mathfrak p$ 均为 rig-闭，则映射
$$
A[x_1,\ldots,x_n]_{\mathfrak q'}
\to
A\{x_1,\ldots,x_n\}_\mathfrak q
$$
在关于各自极大理想的完备化之间定义同构。
\end{lemma}

\begin{proof}
由引理 \ref{restricted-lemma-rig-closed-point-relative}，环映射
$A/\mathfrak p\to A\{x_1,\ldots,x_n\}/\mathfrak q$ 有限。
对每个 $m\geq1$，模
$\mathfrak q^m/\mathfrak q^{m+1}$ 在 $A$ 上有限，
因为它是有限
$A\{x_1,\ldots,x_n\}/\mathfrak q$-模。
所以 $A\{x_1,\ldots x_n\}/\mathfrak q^m$ 是有限 $A$-模。
因此
$A[x_1,\ldots,x_n]\to
A\{x_1,\ldots,x_n\}/\mathfrak q^m$
满射（因为其像稠密且是 $A$-子模）。
由此直接可知，对所有 $m$，
$$
A[x_1,\ldots,x_n]/(\mathfrak q')^m
\to A\{x_1,\ldots,x_n\}/\mathfrak q^m
$$
是同构，引理容易随之得到。
提示：选取拓扑幂零的 $g\in A$，且它不在 $\mathfrak p$ 中。
则完备化之间的映射是
$$
\lim_m\left(A[x_1,\ldots,x_n]/(\mathfrak q')^m\right)_g
\longrightarrow
\left(A\{x_1,\ldots,x_n\}/\mathfrak q^m\right)_g
$$
略去一些细节。
\end{proof}

\begin{lemma}
\label{restricted-lemma-rig-closed-point-etale}
设 $\varphi:A\to B$ 为 $\textit{WAdm}^{Noeth}$ 中的箭头。
假设 $\varphi$ 是 adic、拓扑有限型且平坦的，并且对
某个（相应地，任意）定义理想 $I\subset A$，
$A/I\to B/IB$ 平展。
设 $\mathfrak q\subset B$ 为 rig-闭素理想，且
$\mathfrak p=A\cap\mathfrak q$ 也是 rig-闭的。则
$\mathfrak pB_\mathfrak q=\mathfrak qB_\mathfrak q$。
\end{lemma}

\begin{proof}
令 $\kappa$ 为 $1$ 维完备局部环 $A/\mathfrak p$ 的剩余域。
由于 $A/I\to B/IB$ 平展，可知
$B\otimes_A\kappa$ 是 $\kappa$ 的有限可分扩张的有限积，
见代数之引理 \ref{algebra-lemma-etale-over-field}。
其中一个因子是 $B/\mathfrak q$ 的剩余域。
由代数之引理 \ref{algebra-lemma-finite-over-complete-ring}，
$B/\mathfrak pB$ 是有限 $A/\mathfrak p$-代数。
它也是平坦的。结合上述事实，由代数之引理
\ref{algebra-lemma-characterize-etale}，可知
$A/\mathfrak p\to B/\mathfrak pB$ 有限平展。
因此 $B/\mathfrak pB$ 既约，这蕴含引理的陈述（细节略）。
\end{proof}

\begin{lemma}
\label{restricted-lemma-fibre-regular}
设 $A$ 为 adic 诺特拓扑环。
设 $\mathfrak p\subset A$ 为 rig-闭素理想。
对任意 $n\geq1$，环映射
$$
A/\mathfrak p
\longrightarrow
A\{x_1,\ldots,x_n\}\otimes_AA/\mathfrak p=
A/\mathfrak p\{x_1,\ldots,x_n\}
$$
是正则的。特别地，代数
$A\{x_1,\ldots,x_n\}\otimes_A\kappa(\mathfrak p)$
在 $\kappa(\mathfrak p)$ 上几何正则。
\end{lemma}

\begin{proof}
我们将使用一些关于正则环映射的事实，读者可在更多代数之
\ref{more-algebra-section-regular} 节找到。
由于 $A/\mathfrak p$ 是完备诺特局部环，它是优秀的
（更多代数，命题
\ref{more-algebra-proposition-ubiquity-excellent}）。
因此 $A/\mathfrak p[x_1,\ldots,x_n]$ 也是优秀的
（由同一引用）。
所以由更多代数之引理
\ref{more-algebra-lemma-map-G-ring-to-completion-regular}，
$$
A/\mathfrak p[x_1,\ldots,x_n]
\to A/\mathfrak p\{x_1,\ldots,x_n\}
$$
是正则环同态。
当然，$A/\mathfrak p\to A/\mathfrak p[x_1,\ldots,x_n]$
光滑，因而正则。由于正则环映射的复合仍正则，证明完毕。
\end{proof}

\section{Rig-平坦同态}
\label{restricted-section-rig-flat-homomorphisms}

\noindent
本节定义 adic 诺特拓扑环的 rig-平坦同态。

\begin{lemma}
\label{restricted-lemma-naively-rig-flat-continuous}
设 $\varphi:A\to B$ 为 $\textit{WAdm}^{adic*}$ 中的态射
（形式空间，第 \ref{formal-spaces-section-morphisms-rings} 节）。
假设 $\varphi$ 是 adic 的。则下列条件等价：
\begin{enumerate}
\item 对每个拓扑幂零元 $f\in A$，$B_f$ 在 $A$ 上平坦；
\item 对每个拓扑幂零元 $g\in B$，$B_g$ 在 $A$ 上平坦；
\item 对每个不包含定义理想的素理想 $\mathfrak q\subset B$，
$B_\mathfrak q$ 在 $A$ 上平坦；
\item 对每个 rig-闭素理想 $\mathfrak q\subset B$，
$B_\mathfrak q$ 在 $A$ 上平坦；以及
\item 在此添加更多内容。
\end{enumerate}
\end{lemma}

\begin{proof}
这由定义和代数之引理
\ref{algebra-lemma-flat-localization} 得出。
\end{proof}

\begin{definition}
\label{restricted-definition-naively-rig-flat}
设 $\varphi:A\to B$ 为 adic 诺特拓扑环之间的连续环同态，
即 $\varphi$ 是 $\textit{WAdm}^{Noeth}$ 中的箭头。
如果 $\varphi$ 是 adic 的、拓扑有限型的，并满足引理
\ref{restricted-lemma-naively-rig-flat-continuous} 中的等价条件，
则称 $\varphi$ 是{\it 朴素 rig-平坦的}。
\end{definition}

\noindent
下例表明，这一概念并不能“局部化”。

\begin{example}
\label{restricted-example-recompletion-not-rig-flat}
由例子之引理 \ref{examples-lemma-nonreduced-recompletion}，
存在一个局部诺特 $2$ 维整环 $(A,\mathfrak m)$，它关于主理想
$I=(a)$ 完备；还存在 $f\in\mathfrak m$、$f\not\in I$，
使环 $A_{\{f\}}[1/a]$ 非约化。
这里，$A_{\{f\}}$ 是主局部化 $A_f$ 的 $I$-进完备化
$(A_f)^\wedge$。需要说明的是，环 $A_{\{f\}}[1/a]$ 非零。
令
$B=A_{\{f\}}/\text{nil}(A_{\{f\}})$
为模其幂零根的商。注意 $A\to B$ 是 adic 的且拓扑有限型。
事实上，在将 $x$ 映到 $1/f$ 在 $B$ 中的像的映射下，
$B$ 是 $A\{x\}=A[x]^\wedge$ 的一个商。
$B$ 中每个不包含 $a$ 的素理想 $\mathfrak q$ 都必定
位于 $(0)\subset A$ 之上\footnote{事实上，由于 $B$ 是 $a$-进完备的，
可找到 $\mathfrak q\subset\mathfrak q'\subset B$，使
$a\in\mathfrak q'$。于是
$\mathfrak p'=A\cap\mathfrak q'$ 包含 $a$ 而不包含 $f$，
故为高度 $1$ 的素理想。又因 $a\not\in\mathfrak p$，
$\mathfrak p=A\cap\mathfrak q$ 必真包含于 $\mathfrak p'$。
由 $\dim(A)=2$ 可知 $\mathfrak p=(0)$。}。
因此 $B_\mathfrak q$ 作为 $A$ 的分式域上的模，在 $A$ 上平坦。
所以 $A\to B$ 是朴素 rig-平坦的。另一方面，映射
$$
A_{\{f\}}
\longrightarrow
B_{\{f\}}=(B_f)^\wedge=B=A_{\{f\}}/\text{nil}(A_{\{f\}})
$$
在逆转 $a$ 后并不平坦，因为所得映射
$A_{\{f\}}[1/a]\to
A_{\{f\}}[1/a]/\text{nil}(A_{\{f\}}[1/a])$
是非平凡满射。因此
$A_{\{f\}}\to B_{\{f\}}^\wedge$ 并非朴素 rig-平坦！
\end{example}

\noindent
事实证明，采用下述定义便可轻易绕开这一问题。

\begin{definition}
\label{restricted-definition-rig-flat-continuous-homomorphism}
设 $\varphi:A\to B$ 为 adic 诺特拓扑环之间的连续环同态，
即 $\varphi$ 是 $\textit{WAdm}^{Noeth}$ 中的箭头。
如果 $\varphi$ 是 adic 的、拓扑有限型的，并且对每个 $f\in A$，
诱导映射
$$
A_{\{f\}}\longrightarrow B_{\{f\}}
$$
都是朴素 rig-平坦的（定义
\ref{restricted-definition-naively-rig-flat}），则称 $\varphi$ 是
{\it rig-平坦的}。
\end{definition}

\noindent
在上述定义中取 $f=1$，可见 rig-平坦性蕴含朴素 rig-平坦性。
上例表明其逆命题不成立。不过，在许多情形下，我们不必担心
rig-平坦性及其朴素版本之间的差别，如下一个引理所示。

\begin{lemma}
\label{restricted-lemma-rig-flat-naive}
设 $\varphi:A\to B$ 为 $\textit{WAdm}^{Noeth}$ 中的箭头。
若对某个（等价地，任意）定义理想 $I\subset A$，$A/I$ 是
Jacobson 环，并且 $\varphi$ 是朴素 rig-平坦的，
则 $\varphi$ 是 rig-平坦的。
\end{lemma}

\begin{proof}
假设 $\varphi$ 是朴素 rig-平坦的。先陈述这些假设的若干显然推论。
设 $f\in A$。则 $A,B,A_{\{f\}},B_{\{f\}}$ 都是诺特 adic
拓扑环。映射
$A\to A_{\{f\}}\to B_{\{f\}}$ 和
$A\to B\to B_{\{f\}}$ 都是 adic 的且拓扑有限型。
环映射 $A\to A_{\{f\}}$ 和 $B\to B_{\{f\}}$ 是平坦的，
因为它们分别是 $A\to A_f$ 和 $B\to B_f$ 与完备化映射的复合，
而完备化映射由代数之引理
\ref{algebra-lemma-completion-flat} 是平坦的。
环 $A,B,A_{\{f\}},B_{\{f\}}$ 各自模 $I$ 的商都是
$A/I$ 上的有限型代数，因而也都是 Jacobson 环
（代数，命题 \ref{algebra-proposition-Jacobson-permanence}）。

\medskip\noindent
设 $\mathfrak q'\subset B_{\{f\}}$ 是 rig-闭的。
只需证明 $(B_{\{f\}})_{\mathfrak q'}$ 在 $A_{\{f\}}$ 上平坦，
见引理 \ref{restricted-lemma-naively-rig-flat-continuous}。
由引理 \ref{restricted-lemma-rig-closed-jacobson}，位于 $\mathfrak q'$ 下方的
素理想 $\mathfrak q\subset B$、$\mathfrak p'\subset A_{\{f\}}$
和 $\mathfrak p\subset A$ 都是 rig-闭的。我们将把代数之引理
\ref{algebra-lemma-yet-another-variant-local-criterion-flatness}
应用于图表
$$
\xymatrix{
B_\mathfrak q \ar[r] &
(B_{\{f\}})_{\mathfrak q'} \\
A_\mathfrak p \ar[u] \ar[r] &
(A_{\{f\}})_{\mathfrak p'} \ar[u]
}
$$
并取 $M=B_\mathfrak q$。
唯一尚未核验的假设是：$\mathfrak p$ 生成
$(A_{\{f\}})_{\mathfrak p'}$ 的极大理想。
这由引理 \ref{restricted-lemma-rig-closed-point-in-localization} 得出；
这里利用 $\mathfrak p$ 和 $\mathfrak p'$ 都 rig-闭，以看出
$f$ 在 $A/\mathfrak p$ 中的像是单位
（这是证明中唯一会在没有 Jacobson 假设时失效的一步）。
具体而言，这说明
$A/\mathfrak p\to A_{\{f\}}/\mathfrak p'$
是局部环的有限包含（引理
\ref{restricted-lemma-rig-closed-point-relative}），并且 $f$ 在后一个环中的像是单位。
\end{proof}

\begin{lemma}
\label{restricted-lemma-rig-flat-base-change}
设 $\varphi:A\to B$ 和 $A\to C$ 为
$\textit{WAdm}^{Noeth}$ 中的箭头。
假设 $\varphi$ 是 rig-平坦的，且 $A\to C$ 是 adic 的并且拓扑有限型。
则 $C\to B\widehat{\otimes}_A C$ 是 rig-平坦的。
\end{lemma}

\begin{proof}
假设 $\varphi$ 是 rig-平坦的。设 $f\in C$。
我们须证明
$C_{\{f\}}\to B\widehat{\otimes}_A C_{\{f\}}$
是朴素 rig-平坦的。由于可以用 $C_{\{f\}}$ 代替 $C$，
只需证明 $C\to B\widehat{\otimes}_A C$ 是朴素 rig-平坦的。

\medskip\noindent
若 $A\to C$ 满射，或更一般地，若 $C$ 作为 $A$-模是有限的，
则
$B\otimes_A C=B\widehat{\otimes}_A C$，
因为完备诺特环上的有限模是完备的，见代数之引理
\ref{algebra-lemma-completion-tensor}。
由平坦性的通常基变换
（代数，引理 \ref{algebra-lemma-flat-base-change}），
可见在此情形下，$\varphi$ 的朴素 rig-平坦性蕴含
$C\to B\times_A C$ 的朴素 rig-平坦性。

\medskip\noindent
一般情形下，可将 $A\to C$ 分解为
$A\to A\{x_1,\ldots,x_n\}\to C$，
其中 $A\{x_1,\ldots,x_n\}$ 是受限幂级数环，且
$A\{x_1,\ldots,x_n\}\to C$ 满射。
因此，只需在 $C=A\{x_1,\ldots,x_n\}$ 时证明
$C\to B\widehat{\otimes}_A B$ 是朴素 rig-平坦的。
由于
$A\{x_1,\ldots,x_n\}=A\{x_1,\ldots,x_{n-1}\}\{x_n\}$，
对 $n$ 归纳便归结到下一段讨论的情形。

\medskip\noindent
这里 $C=A\{x\}$。注意
$B\widehat{\otimes}_A C=B\{x\}$。
我们须证明 $A\{x\}\to B\{x\}$ 是朴素 rig-平坦的。
设 $\mathfrak q\subset B\{x\}$ 为 rig-闭素理想。
我们须证明 $B\{x\}_{\mathfrak q}$ 在 $A\{x\}$ 上平坦。
置 $\mathfrak p=A\cap\mathfrak q$。
由引理 \ref{restricted-lemma-rig-closed-point-after-localization}，
可找到 $f\in A$，使 $f$ 在 $B\{x\}/\mathfrak q$ 中的像是单位，
且在 $A_{\{f\}}$ 中诱导出的素理想 $\mathfrak p'$ 是 rig-闭的。
下面将使用
$A_{\{f\}}\{x\}=A\{x\}_{\{f\}}$，并对 $B$ 作类似使用；
细节从略。考虑图表
$$
\xymatrix{
(B\{x\})_{\mathfrak q} \ar[r] &
(B_{\{f\}}\{x\})_{\mathfrak q'} \\
A\{x\} \ar[r] \ar[u] &
A_{\{f\}}\{x\} \ar[u]
}
$$
我们要证明左边的竖直箭头平坦。
上方水平箭头忠实平坦：它是局部环的局部同态，并且由于
$B_{\{f\}}\{x\}$ 是诺特环 $B\{x\}$ 的某个局部化的完备化，
它也是平坦的。类似地，下方水平箭头平坦。
所以只需证明右边的竖直箭头平坦。这把问题归结为下一段讨论的情形。

\medskip\noindent
这里 $C=A\{x\}$，并且有一个 rig-闭素理想
$\mathfrak q\subset B\{x\}$，使
$\mathfrak p=A\cap\mathfrak q$ 也 rig-闭。
由引理 \ref{restricted-lemma-rig-closed-point-relative}，这蕴含中间素理想
$B\cap\mathfrak q$ 和 $A\{x\}\cap\mathfrak q$ 也都是 rig-闭的。
考虑诺特局部环的局部同态图表
$$
\xymatrix{
(B[x])_{B[x]\cap\mathfrak q} \ar[r] &
(B\{x\})_{\mathfrak q} \\
(A[x])_{A[x]\cap\mathfrak q} \ar[r] \ar[u] &
(A\{x\})_{A\{x\}\cap\mathfrak q} \ar[u]
}
$$
由引理 \ref{restricted-lemma-rig-closed-point-variables}，
水平箭头在完备化上诱导同构。
我们已经知道左边的竖直箭头平坦
（因为 $A\to B$ 是朴素 rig-平坦的，故在由定义理想确定的
闭轨迹之外，$A[x]\to B[x]$ 是平坦的）。
最后由更多代数之引理
\ref{more-algebra-lemma-flat-completion} 得出结论。
\end{proof}

\begin{lemma}
\label{restricted-lemma-rig-flat-local-etale}
考虑 $\textit{WAdm}^{Noeth}$ 中的交换图表
$$
\xymatrix{
B \ar[r] & B' \\
A \ar[r] \ar[u]^\varphi & A' \ar[u]_{\varphi'}
}
$$
其中所有箭头都是 adic 的且拓扑有限型。
假设 $A\to A'$ 和 $B\to B'$ 平坦。
设 $I\subset A$ 为定义理想。
若 $\varphi$ 是 rig-平坦的且
$A/I\to A'/IA'$ 是 \'etale 的，则 $\varphi'$ 是 rig-平坦的。
\end{lemma}

\begin{proof}
给定 $f\in A'$，该引理的假设对图表
$$
\xymatrix{
B \ar[r] & (B')_{\{f\}} \\
A \ar[r] \ar[u]^\varphi & (A')_{\{f\}} \ar[u]
}
$$
仍成立。因此只需证明 $\varphi'$ 是朴素 rig-平坦的。

\medskip\noindent
取 rig-闭素理想 $\mathfrak q'\subset B'$。
我们须证明 $(B')_{\mathfrak q'}$ 在 $A'$ 上平坦。
可选择 $f\in A$，使其在 $B'/\mathfrak q'$ 中的像是单位，
且 $A_{\{f\}}$ 中诱导出的素理想 $\mathfrak p''$ 是 rig-闭的，
见引理 \ref{restricted-lemma-rig-closed-point-after-localization}。
准确地说，这里
$\mathfrak q''=\mathfrak q'B'_{\{f\}}$ 且
$\mathfrak p''=A_{\{f\}}\cap\mathfrak q''$。
考虑图表
$$
\xymatrix{
B'_{\mathfrak q'} \ar[r] &
(B'_{\{f\}})_{\mathfrak q''} \\
A \ar[r] \ar[u] &
A_{\{f\}} \ar[u]
}
$$
我们要证明左边的竖直箭头平坦。
上方水平箭头忠实平坦：它是局部环的局部同态，并且由于
$B'_{\{f\}}$ 是诺特环 $B'_f$ 的某个局部化的完备化，
它也是平坦的。类似地，下方水平箭头平坦。
所以只需证明右边的竖直箭头平坦。
最后，该引理的所有假设对图表
$$
\xymatrix{
B_{\{f\}} \ar[r] &
B'_{\{f\}} \\
A_{\{f\}} \ar[r] \ar[u] &
A'_{\{f\}} \ar[u]
}
$$
仍成立。这把问题归结为下一段讨论的情形。

\medskip\noindent
取 rig-闭素理想 $\mathfrak q'\subset B'$，并假设
$\mathfrak p=A\cap\mathfrak q'$ 也 rig-闭。
由引理 \ref{restricted-lemma-rig-closed-point-relative}，
这还蕴含素理想
$\mathfrak q=B\cap\mathfrak q'$ 和
$\mathfrak p'=A'\cap\mathfrak q'$ 都是 rig-闭的。
我们将把代数之引理
\ref{algebra-lemma-yet-another-variant-local-criterion-flatness}
应用于图表
$$
\xymatrix{
B_\mathfrak q \ar[r] &
B'_{\mathfrak q'} \\
A_\mathfrak p \ar[u] \ar[r] &
A'_{\mathfrak p'} \ar[u]
}
$$
并取 $M=B_\mathfrak q$。唯一尚未核验的假设是：
$\mathfrak p$ 生成 $A'_{\mathfrak p'}$ 的极大理想。
这由引理 \ref{restricted-lemma-rig-closed-point-etale} 得出。
\end{proof}

\begin{lemma}
\label{restricted-lemma-rig-flat-local-down}
考虑 $\textit{WAdm}^{Noeth}$ 中的交换图表
$$
\xymatrix{
B \ar[r] & B' \\
A \ar[r] \ar[u]^\varphi & A' \ar[u]_{\varphi'}
}
$$
其中所有箭头都是 adic 的且拓扑有限型。
假设 $A\to A'$ 平坦，且 $B\to B'$ 忠实平坦。
若 $\varphi'$ 是 rig-平坦的，则 $\varphi$ 是 rig-平坦的。
\end{lemma}

\begin{proof}
给定 $f\in A$，该引理的假设对图表
$$
\xymatrix{
B_{\{f\}} \ar[r] & (B')_{\{f\}} \\
A_{\{f\}} \ar[r] \ar[u]^\varphi & (A')_{\{f\}} \ar[u]
}
$$
仍成立（核验忠实平坦条件时，可用如下事实：
$B\to B'$ 忠实平坦，等价于 $B\to B'$ 平坦，且对于某个
定义理想 $I\subset A$，映射
$\Spec(B'/IB')\to\Spec(B/IB)$ 满射）。
因此只需证明 $\varphi$ 是朴素 rig-平坦的。
然而，已知 $\varphi'$ 是朴素 rig-平坦的，且
$\Spec(B')\to\Spec(B)$ 满射，结论立即得出。
\end{proof}

\noindent
最后，我们可以证明 rig-平坦性是一个局部性质。

\begin{lemma}
\label{restricted-lemma-rig-flat-axioms}
在 $\textit{WAdm}^{Noeth}$ 的箭头上，性质
$P(\varphi)=$“$\varphi$ 是 rig-平坦的”是形式空间之备注
\ref{formal-spaces-remark-variant-adic-star} 所定义的局部性质。
\end{lemma}

\begin{proof}
回顾一下该陈述的含义。首先，$\textit{WAdm}^{Noeth}$ 是以
adic 诺特拓扑环为对象、以连续环同态为态射的范畴。
考虑交换图表
$$
\xymatrix{
B \ar[r] & (B')^\wedge \\
A \ar[r] \ar[u]^\varphi & (A')^\wedge \ar[u]_{\varphi'}
}
$$
并假设它满足下列条件：$A$ 和 $B$ 是 adic 诺特拓扑环，
$A\to A'$ 和 $B\to B'$ 是 \'etale 环映射，
对于某个定义理想 $I\subset A$，
$(A')^\wedge=\lim A'/I^nA'$；
对于某个定义理想 $J\subset B$，
$(B')^\wedge=\lim B'/J^nB'$；并且
$\varphi:A\to B$ 与
$\varphi':(A')^\wedge\to(B')^\wedge$ 都连续。
注意，由形式空间之引理
\ref{formal-spaces-lemma-completion-in-sub}，
$(A')^\wedge$ 和 $(B')^\wedge$ 是 adic 诺特拓扑环。
我们须证明：
\begin{enumerate}
\item $\varphi$ 是 rig-平坦的 $\Rightarrow\varphi'$ 是 rig-平坦的；
\item 若 $B\to B'$ 忠实平坦，则 $\varphi'$ 是 rig-平坦的
$\Rightarrow\varphi$ 是 rig-平坦的；以及
\item 若对 $i=1,\ldots,n$，$A\to B_i$ 都是 rig-平坦的，
则 $A\to\prod_{i=1,\ldots,n}B_i$ 是 rig-平坦的。
\end{enumerate}
由引理 \ref{restricted-lemma-finite-type}，adic 且拓扑有限型这一性质满足
条件 (1)、(2) 和 (3)。因此，在对“rig-平坦”性质核验
(1)、(2) 和 (3) 时，可以已经假设所有环映射都是 adic 的且拓扑有限型。
此时，(1) 和 (2) 分别由引理
\ref{restricted-lemma-rig-flat-local-etale} 和
\ref{restricted-lemma-rig-flat-local-down} 得出。
平凡的 (3) 的证明从略。
\end{proof}

\begin{lemma}
\label{restricted-lemma-composition-rig-flat-continuous}
在 $\textit{WAdm}^{Noeth}$ 的箭头上，性质
$P(\varphi)=$“$\varphi$ 是 rig-平坦的”在复合下稳定，
其含义如形式空间之备注
\ref{formal-spaces-remark-composition-variant-Noetherian} 所定义。
\end{lemma}

\begin{proof}
由引理 \ref{restricted-lemma-rig-flat-axioms}，该陈述有意义。
为证明它，假设 $\textit{WAdm}^{Noeth}$ 中有 rig-平坦态射
$A\to B$ 和 $B\to C$。
由引理 \ref{restricted-lemma-composition-finite-type}，
$A\to C$ 是 adic 的且拓扑有限型。
为完成证明，须证明对每个 $f\in A$，映射
$A_{\{f\}}\to C_{\{f\}}$ 是朴素 rig-平坦的。
由于 $A_{\{f\}}\to B_{\{f\}}$ 和
$B_{\{f\}}\to C_{\{f\}}$ 都是朴素 rig-平坦的，
只需证明朴素 rig-平坦映射的复合仍朴素 rig-平坦。
这是代数之引理 \ref{algebra-lemma-composition-flat} 的推论。
\end{proof}

\section{Rig-平坦态射}
\label{restricted-section-rig-flat-morphisms}

\noindent
本节利用第 \ref{restricted-section-rig-flat-homomorphisms} 节中的工作，
定义局部诺特代数空间的 rig-平坦态射。

\begin{definition}
\label{restricted-definition-rig-flat}
设 $S$ 为概形。设 $f:X\to Y$ 为 $S$ 上局部诺特形式代数空间的态射。
若对每个交换图表
$$
\xymatrix{
U \ar[d] \ar[r] & V \ar[d] \\
X \ar[r] & Y
}
$$
其中 $U,V$ 是仿射形式代数空间，而 $U\to X$ 与 $V\to Y$
由代数空间可表且平展，态射 $U\to V$
对应于 adic 诺特拓扑环之间的 rig-平坦映射，
则称 $f$ 是{\it rig-平坦的}。
\end{definition}

\noindent
下面证明该条件可以在源和目标上平展局部地检验。

\begin{lemma}
\label{restricted-lemma-rig-flat-morphisms}
设 $S$ 为概形。设 $f:X\to Y$ 为 $S$ 上局部诺特形式代数空间的态射。
以下条件等价：
\begin{enumerate}
\item $f$ 是 rig-平坦的；
\item 对每个交换图表
$$
\xymatrix{
U \ar[d] \ar[r] & V \ar[d] \\
X \ar[r] & Y
}
$$
其中 $U,V$ 是仿射形式代数空间，而 $U\to X$ 与 $V\to Y$
由代数空间可表且平展，态射 $U\to V$
对应于 $\textit{WAdm}^{Noeth}$ 中的 rig-平坦映射；
\item 存在形式空间之定义
\ref{formal-spaces-definition-formal-algebraic-space}
所述的覆盖 $\{Y_j\to Y\}$，并且对每个 $j$，
存在形式空间之定义
\ref{formal-spaces-definition-formal-algebraic-space}
所述的覆盖 $\{X_{ji}\to Y_j\times_YX\}$，
使每个 $X_{ji}\to Y_j$ 对应于
$\textit{WAdm}^{Noeth}$ 中的 rig-平坦映射；并且
\item 存在形式空间之定义
\ref{formal-spaces-definition-formal-algebraic-space}
所述的覆盖 $\{X_i\to X\}$，并且对每个 $i$，
存在分解 $X_i\to Y_i\to Y$，其中 $Y_i$ 是仿射形式代数空间，
$Y_i\to Y$ 由代数空间可表且平展，而
$X_i\to Y_i$ 对应于
$\textit{WAdm}^{Noeth}$ 中的 rig-平坦映射。
\end{enumerate}
\end{lemma}

\begin{proof}
(1) 与 (2) 的等价性就是定义 \ref{restricted-definition-rig-flat}。
(2)、(3)、(4) 的等价性来自以下事实：
由引理 \ref{restricted-lemma-rig-flat-axioms}，rig-平坦是
$\text{WAdm}^{Noeth}$ 箭头的局部性质；
再应用形式空间之引理
\ref{formal-spaces-lemma-property-defines-property-morphisms}
对局部诺特代数空间之间态射的变体，
该变体在形式空间之备注
\ref{formal-spaces-remark-variant-Noetherian} 中提及。
\end{proof}

\begin{lemma}
\label{restricted-lemma-base-change-rig-flat}
设 $S$ 为概形。设 $f:X\to Y$ 与 $g:Z\to Y$
为 $S$ 上局部诺特形式代数空间的态射。
若 $f$ 是 rig-平坦的，且 $g$ 局部有限型，
则基变换 $X\times_YZ\to Z$ 是 rig-平坦的。
\end{lemma}

\begin{proof}
由形式空间之备注
\ref{formal-spaces-remark-base-change-variant-variant-Noetherian}
及形式空间之第 \ref{formal-spaces-section-adic} 节中的讨论，
这由引理 \ref{restricted-lemma-rig-flat-base-change} 得出。
\end{proof}

\begin{lemma}
\label{restricted-lemma-composition-rig-flat}
设 $S$ 为概形。设 $f:X\to Y$ 与 $g:Y\to Z$
为 $S$ 上局部诺特形式代数空间的态射。
若 $f$ 和 $g$ 都是 rig-平坦的，则 $g\circ f$ 也是。
\end{lemma}

\begin{proof}
由形式空间之备注
\ref{formal-spaces-remark-composition-variant-Noetherian}，
这由引理 \ref{restricted-lemma-composition-rig-flat-continuous} 得出。
\end{proof}

\section{Rig-光滑同态}
\label{restricted-section-rig-smooth-homomorphisms}

\noindent
本节证明 adic 诺特拓扑环的 rig-光滑同态的一些性质，
这些性质是引入局部诺特形式代数空间的 rig-光滑态射所需要的。

\begin{lemma}
\label{restricted-lemma-rig-smooth-continuous}
设 $A\to B$ 为 $\textit{WAdm}^{Noeth}$ 中的态射
（形式空间，第 \ref{formal-spaces-section-morphisms-rings} 节）。
下列条件等价：
\begin{enumerate}
\item[(a)] $A\to B$ 满足引理 \ref{restricted-lemma-finite-type} 中的等价条件，
并且存在定义理想 $I\subset B$，使 $B$ 在 $(A,I)$ 上 rig-光滑；以及
\item[(b)] $A\to B$ 满足引理 \ref{restricted-lemma-finite-type} 中的等价条件，
并且对每个定义理想 $I\subset A$，代数 $B$ 都在 $(A,I)$ 上 rig-光滑。
\end{enumerate}
\end{lemma}

\begin{proof}
设 $I$ 和 $I'$ 是 $A$ 的定义理想。则存在整数 $c\geq0$，使
$I^c\subset I'$ 且 $(I')^c\subset I$。因此，$B$ 在 $(A,I)$ 上
rig-光滑，当且仅当 $B$ 在 $(A,I')$ 上 rig-光滑。
这由定义 \ref{restricted-definition-rig-smooth-homomorphism}、包含关系
$I^c\subset I'$ 与 $(I')^c\subset I$，以及下述事实得出：
由备注 \ref{restricted-remark-NL-well-defined-topological}，
朴素余切复形 $\NL_{B/A}^\wedge$ 不依赖于 $A$ 的定义理想的选择。
\end{proof}

\begin{definition}
\label{restricted-definition-rig-smooth-continuous-homomorphism}
设 $\varphi:A\to B$ 为 adic 诺特拓扑环之间的连续环同态，
即 $\varphi$ 是 $\textit{WAdm}^{Noeth}$ 中的箭头。
若引理 \ref{restricted-lemma-rig-smooth-continuous} 中的等价条件成立，
则称 $\varphi$ 是{\it rig-光滑的}。
\end{definition}

\noindent
这定义了一个局部性质。

\begin{lemma}
\label{restricted-lemma-rig-smooth-axioms}
在 $\textit{WAdm}^{Noeth}$ 的箭头上，性质
$P(\varphi)=$“$\varphi$ 是 rig-光滑的”是形式空间之备注
\ref{formal-spaces-remark-variant-Noetherian} 所定义的局部性质。
\end{lemma}

\begin{proof}
回顾一下该陈述的含义。首先，$\textit{WAdm}^{Noeth}$ 是以
adic 诺特拓扑环为对象、以连续环同态为态射的范畴。
考虑交换图表
$$
\xymatrix{
B \ar[r] & (B')^\wedge \\
A \ar[r] \ar[u]^\varphi & (A')^\wedge \ar[u]_{\varphi'}
}
$$
并假设它满足下列条件：$A$ 和 $B$ 是 adic 诺特拓扑环，
$A\to A'$ 和 $B\to B'$ 是 \'etale 环映射，
对于某个定义理想 $I\subset A$，
$(A')^\wedge=\lim A'/I^nA'$；
对于某个定义理想 $J\subset B$，
$(B')^\wedge=\lim B'/J^nB'$；并且
$\varphi:A\to B$ 与
$\varphi':(A')^\wedge\to(B')^\wedge$ 都连续。
注意，由形式空间之引理
\ref{formal-spaces-lemma-completion-in-sub}，
$(A')^\wedge$ 和 $(B')^\wedge$ 是 adic 诺特拓扑环。
我们须证明：
\begin{enumerate}
\item $\varphi$ 是 rig-光滑的 $\Rightarrow\varphi'$ 是 rig-光滑的；
\item 若 $B\to B'$ 忠实平坦，则 $\varphi'$ 是 rig-光滑的
$\Rightarrow\varphi$ 是 rig-光滑的；以及
\item 若对 $i=1,\ldots,n$，$A\to B_i$ 都是 rig-光滑的，
则 $A\to\prod_{i=1,\ldots,n}B_i$ 是 rig-光滑的。
\end{enumerate}
引理 \ref{restricted-lemma-finite-type} 中的等价条件满足条件 (1)、(2) 和 (3)。
因此，在对“rig-光滑”性质核验 (1)、(2) 和 (3) 时，
可以已经假设每种情形下的环映射都满足引理
\ref{restricted-lemma-finite-type} 中的等价条件。

\medskip\noindent
选取定义理想 $I\subset A$。由以上所述，图表中每个环的拓扑都是
$I$-进拓扑，并且 $B$、$(A')^\wedge$、$(B')^\wedge$
属于 $(A,I)$ 的范畴 (\ref{restricted-equation-C-prime})。
由于 $A\to A'$ 和 $B\to B'$ 都是 \'etale 的，复形
$\NL_{A'/A}$ 和 $\NL_{B'/B}$ 为零，故由引理
\ref{restricted-lemma-NL-is-completion}，
$\NL_{(A')^\wedge/A}^\wedge$ 和
$\NL_{(B')^\wedge/B}^\wedge$ 为零。
将引理 \ref{restricted-lemma-exact-sequence-NL} 应用于
$A\to(A')^\wedge\to(B')^\wedge$，得到同构
$$
H^i(\NL_{(B')^\wedge/(A')^\wedge}^\wedge) \to H^i(\NL_{(B')^\wedge/A}^\wedge)
$$
因此
$\NL_{(B')^\wedge/A}^\wedge\to
\NL_{(B')^\wedge/(A')^\wedge}$ 是拟同构。
环映射 $B/I^nB\to B'/I^nB'$ 是 \'etale 的，
因而是局部完全交
（代数，引理 \ref{algebra-lemma-etale-standard-smooth}）。
所以可将引理 \ref{restricted-lemma-exact-sequence-NL} 和
\ref{restricted-lemma-transitive-lci-at-end} 应用于
$A\to B\to(B')^\wedge$，得到同构
$$
H^i(\NL_{B/A}^\wedge \otimes_B (B')^\wedge) \to
H^i(\NL_{(B')^\wedge/A}^\wedge)
$$
由此可知
$\NL_{B/A}^\wedge\otimes_B(B')^\wedge\to
\NL_{(B')^\wedge/A}^\wedge$ 是拟同构。
结合这两个观察，得到 $D((B')^\wedge)$ 中的同构
$$
\NL_{(B')^\wedge/(A')^\wedge}^\wedge \cong
\NL_{B/A}^\wedge \otimes_B (B')^\wedge
$$
准备工作至此完成，可以开始真正的证明。

\medskip\noindent
证明 (1)。假设 $\varphi$ 是 rig-光滑的。则存在 $c\geq0$，使对每个
$B$-模 $N$，$\Ext^1_B(\NL_{B/A}^\wedge,N)$ 都被 $I^c$ 零化。
由更多代数之引理
\ref{more-algebra-lemma-two-term-base-change} 和
\ref{more-algebra-lemma-base-change-property-ext-1-annihilated}，
这一性质在沿 $B\to(B')^\wedge$ 作基变换时保持。
因此，对每个 $(B')^\wedge$-模 $N$，
$\Ext^1_{(B')^\wedge}
(\NL_{(B')^\wedge/(A')^\wedge}^\wedge,N)$
都被 $I^c(A')^\wedge$ 零化，这说明 $\varphi'$ 是 rig-光滑的。
这证明了 (1)。

\medskip\noindent
为证明 (2)，假设 $B\to B'$ 忠实平坦且 $\varphi'$ 是 rig-光滑的。
则存在 $c\geq0$，使对每个 $(B')^\wedge$-模 $N'$，
$\Ext^1_{(B')^\wedge}
(\NL_{(B')^\wedge/(A')^\wedge}^\wedge,N')$
都被 $I^c(B')^\wedge$ 零化。
复合 $B\to B'\to(B')^\wedge$ 是平坦的
（代数，引理 \ref{algebra-lemma-completion-flat}），
故对任意 $B$-模 $N$，有
$$
\Ext^1_B(\NL_{B/A}^\wedge, N) \otimes_B (B')^\wedge =
\Ext^1_{(B')^\wedge}(\NL_{B/A}^\wedge \otimes_B (B')^\wedge,
N \otimes_B (B')^\wedge)
$$
其中使用了更多代数之引理
\ref{more-algebra-lemma-base-change-RHom} 的 (3)
（略去少量细节）。因此，该模被 $I^c$ 零化。
然而，由我们假设 $B\to B'$ 忠实平坦，
$B\to(B')^\wedge$ 实际上忠实平坦
（形式空间，引理 \ref{formal-spaces-lemma-etale-surjective}）。
由此可知 $\Ext^1_B(\NL_{B/A}^\wedge,N)$ 被 $I^c$ 零化。
所以 $\varphi$ 是 rig-光滑的。这证明了 (2)。

\medskip\noindent
为证明 (3)，置 $B=\prod_{i=1,\ldots,n}B_i$；
只需注意，$\NL_{B/A}^\wedge$ 是复形
$\NL_{B_i/A}^\wedge$ 的直和，其中各复形均视为 $B$-模复形。
\end{proof}

\begin{lemma}
\label{restricted-lemma-base-change-rig-smooth-continuous}
考虑 $\textit{WAdm}^{Noeth}$ 的箭头上的性质
$P(\varphi)=$“$\varphi$ 是 rig-光滑的”与
$Q(\varphi)$=“$\varphi$ 是 adic 的”。
则依照形式空间之备注
\ref{formal-spaces-remark-base-change-variant-variant-Noetherian}
中的定义，$P$ 在由 $Q$ 作基变换下稳定。
\end{lemma}

\begin{proof}
由引理 \ref{restricted-lemma-rig-smooth-continuous}，该陈述有意义。
为证明它，假设 $\textit{WAdm}^{Noeth}$ 中有态射
$B\to A$ 和 $B\to C$，且 $B\to A$ 是 rig-光滑的，
$B\to C$ 是 adic 的
（形式空间，定义 \ref{formal-spaces-definition-adic-homomorphism}）。
则可选取定义理想 $I\subset B$，使 $A$ 和 $C$ 上的拓扑都是
$I$-进拓扑。在这种情形下，由引理
\ref{restricted-lemma-base-change-rig-smooth-homomorphism} 立即可知，
$A\widehat{\otimes}_B C$ 在 $(C,IC)$ 上 rig-光滑。
\end{proof}

\begin{lemma}
\label{restricted-lemma-composition-rig-smooth-continuous}
在 $\textit{WAdm}^{Noeth}$ 的箭头上，性质
$P(\varphi)=$“$\varphi$ 是 rig-光滑的”在复合下稳定，
其含义如形式空间之备注
\ref{formal-spaces-remark-composition-variant-Noetherian} 所定义。
\end{lemma}

\begin{proof}
我们强烈建议读者自行寻找证明，不要阅读下面的证明。
由引理 \ref{restricted-lemma-rig-smooth-continuous}，该陈述有意义。
为证明它，假设 $\textit{WAdm}^{Noeth}$ 中有 rig-光滑态射
$A\to B$ 和 $B\to C$。
则可选取定义理想 $I\subset A$，使 $C$ 和 $B$ 上的拓扑都是
$I$-进拓扑。由引理 \ref{restricted-lemma-exact-sequence-NL}，得到正合列
$$
\xymatrix{
C \otimes_B H^0(\NL_{B/A}^\wedge) \ar[r] &
H^0(\NL_{C/A}^\wedge) \ar[r] &
H^0(\NL_{C/B}^\wedge) \ar[r] & 0 \\
H^{-1}(\NL_{B/A}^\wedge \otimes_B C) \ar[r] &
H^{-1}(\NL_{C/A}^\wedge) \ar[r] &
H^{-1}(\NL_{C/B}^\wedge) \ar[llu]
}
$$
注意，$H^{-1}(\NL_{B/A}^\wedge\otimes_BC)$ 和
$H^{-1}(\NL_{C/B}^\wedge)$ 都被 $I$ 的某个幂零化；
这由引理 \ref{restricted-lemma-equivalent-with-artin-smooth} 的 (2)，
结合更多代数之引理
\ref{more-algebra-lemma-two-term-base-change} 和
\ref{more-algebra-lemma-base-change-property-ext-1-annihilated}
得出（后两者用于处理沿 $B\to C$ 的基变换）。
因此 $H^{-1}(\NL_{C/A}^\wedge)$ 被 $I$ 的某个幂零化。
接着，根据引理 \ref{restricted-lemma-equivalent-with-artin-smooth} 的 (2)
所给出的 rig-光滑代数刻画，而该刻画又引用更多代数之引理
\ref{more-algebra-lemma-ext-1-annihilated} 的 (5)，
可以选取 $f_1,\ldots,f_s\in IB$ 和 $g_1,\ldots,g_t\in IC$，使
$V(f_1,\ldots,f_s)=V(IB)$ 且
$V(g_1,\ldots,g_t)=V(IC)$，并使
$H^0(\NL_{B/A}^\wedge)_{f_i}$ 是有限投射 $B_{f_i}$-模，
$H^0(\NL_{C/B}^\wedge)_{g_j}$ 是有限投射 $C_{g_j}$-模。
由于次数 $-1$ 的上同调在 $f_ig_j$ 处局部化后消失，得到短正合列
$$
0 \to
(C \otimes_B H^0(\NL_{B/A}^\wedge))_{f_ig_j} \to
H^0(\NL_{C/A}^\wedge)_{f_ig_j} \to
H^0(\NL_{C/B}^\wedge)_{f_ig_j} \to 0
$$
从而 $H^0(\NL_{C/A}^\wedge)_{f_ig_j}$ 作为同类模的扩张，
是有限投射 $C_{f_ig_j}$-模。
所以，由引理 \ref{restricted-lemma-equivalent-with-artin-smooth} 的 (2) 中的判据，
并进而由更多代数之引理
\ref{more-algebra-lemma-ext-1-annihilated} 的 (4) 中的判据，
可知 $C$ 在 $(A,I)$ 上 rig-光滑。
\end{proof}

\noindent
下述引理可解释为：rig-光滑同态是“rig-syntomic”的，
或说是“rig-平坦$+$rig-lci”的。

\begin{lemma}
\label{restricted-lemma-rig-smooth-rig-flat}
设 $\varphi:A\to B$ 为 $\textit{WAdm}^{Noeth}$ 中的箭头。
若 $\varphi$ 是 rig-光滑的，则 $\varphi$ 是 rig-平坦的；
并且对于任意表示 $B=A\{x_1,\ldots,x_n\}/J$ 以及不包含定义理想的
素理想 $J\subset\mathfrak q\subset A\{x_1,\ldots,x_n\}$，
理想
$J_\mathfrak q\subset A\{x_1,\ldots,x_n\}_\mathfrak q$
由正则序列生成。
\end{lemma}

\begin{proof}
设 $f\in A$。为证明 $\varphi$ 是 rig-平坦的，须证明
$\varphi_{\{f\}}:A_{\{f\}}\to B_{\{f\}}$ 是朴素 rig-平坦的。
现在，或者将 $\varphi_{\{f\}}$ 视为 $\varphi$ 的基变换并使用引理
\ref{restricted-lemma-base-change-rig-smooth-continuous}，
或者使用 rig-光滑是局部性质这一事实
（引理 \ref{restricted-lemma-rig-smooth-axioms}），
均可看出 $\varphi_{\{f\}}$ 是 rig-光滑的。
因此只需证明 $\varphi$ 是朴素 rig-平坦的。

\medskip\noindent
选取表示 $B=A\{x_1,\ldots,x_n\}/J$。
为核验引理的第二部分，只需对 $J\subset\mathfrak q$，
其中 $\mathfrak q$ 在不包含定义理想的素理想中极大，核验
$J_\mathfrak q\subset A\{x_1,\ldots,x_n\}_\mathfrak q$
由正则序列生成，见代数之引理
\ref{algebra-lemma-regular-sequence-in-neighbourhood}
（该引理表明，在 $V(J)$ 中，$J$ 可由正则序列生成的素理想构成开集）。
换言之，可以假设 $\mathfrak q$ 在
$A\{x_1,\ldots,x_n\}$ 中 rig-闭。
而为核验 $B$ 是朴素 rig-平坦的，也只需核验相应的局部化
$B_\mathfrak q$ 在 $A$ 上平坦。

\medskip\noindent
设 $\mathfrak q\subset A\{x_1,\ldots,x_n\}$ 是 rig-闭的，且
$J\subset\mathfrak q$。由引理
\ref{restricted-lemma-rig-closed-point-after-localization}，
可选取 $f\in A$，使其在
$A\{x_1,\ldots,x_n\}/\mathfrak q$ 中的像是单位，
且在 $A_{\{f\}}$ 中诱导出的素理想 $\mathfrak p'$ 是 rig-闭的。
下面将使用
$A_{\{f\}}\{x_1,\ldots,x_n\}=
A\{x_1,\ldots,x_n\}_{\{f\}}$；细节从略。
考虑图表
$$
\xymatrix{
A\{x_1, \ldots, x_n\}_{\mathfrak q} / J_\mathfrak q \ar[r] &
A_{\{f\}}\{x_1, \ldots, x_n\}_{\mathfrak q'}/
J A_{\{f\}}\{x_1, \ldots, x_n\}_{\mathfrak q'} \\
A\{x_1, \ldots, x_n\}_{\mathfrak q} \ar[r] \ar[u] &
A_{\{f\}}\{x_1, \ldots, x_n\}_{\mathfrak q'} \ar[u] \\
A \ar[r] \ar[u] &
A_{\{f\}} \ar[u]
}
$$
中间的水平箭头忠实平坦：它是局部环的局部同态，并且由于
$A_{\{f\}}\{x_1,\ldots,x_n\}$ 是诺特环
$A\{x_1,\ldots,x_n\}$ 的某个局部化的完备化，它也是平坦的。
类似地，底部水平箭头平坦。因此，为证明 $J_\mathfrak q$
由正则序列生成，且
$A\to A\{x_1,\ldots,x_n\}_{\mathfrak q}/J_\mathfrak q$
平坦，只需对
$J A_{\{f\}}\{x_1,\ldots,x_n\}_{\mathfrak q'}$ 和
$A_{\{f\}}\to A_{\{f\}}\{x_1,\ldots,x_n\}_{\mathfrak q'}/
J A_{\{f\}}\{x_1,\ldots,x_n\}_{\mathfrak q'}$
证明相同的陈述。关于正则序列的陈述，见代数之引理
\ref{algebra-lemma-flat-increases-depth} 或更多代数之引理
\ref{more-algebra-lemma-flat-descent-regular-ideal}。
最后，我们已经看到 $A_{\{f\}}\to B_{\{f\}}$ 是 rig-光滑的。
这把问题归结为下一段讨论的情形。

\medskip\noindent
设 $\mathfrak q\subset A\{x_1,\ldots,x_n\}$ 是 rig-闭的，
$J\subset\mathfrak q$，并且还假设
$\mathfrak p=A\cap\mathfrak q$ 也是 rig-闭的。
由引理 \ref{restricted-lemma-equivalent-with-artin-smooth} 给出的
rig-光滑代数刻画，重新排列变量 $x_1,\ldots,x_n$ 后，
可找到 $m\geq0$ 和 $f_1,\ldots,f_m\in J$，使：
\begin{enumerate}
\item $J_\mathfrak q$ 由 $f_1,\ldots,f_m$ 生成；并且
\item $\det_{1\leq i,j\leq m}(\partial f_j/\partial x_i)$
在 $A\{x_1,\ldots,x_n\}_\mathfrak q$ 中的像是单位。
\end{enumerate}
由引理 \ref{restricted-lemma-fibre-regular}，纤维环
$$
F = A\{x_1, \ldots, x_n\} \otimes_A \kappa(\mathfrak p)
$$
是正则的。注意，$A$-导子 $\partial/\partial x_i$
（唯一地）延拓为导子 $D_i:F\to F$。
由更多代数之引理
\ref{more-algebra-lemma-quotient-sequence-regular}，
$f_1,\ldots,f_m$ 在 $F_\mathfrak q$ 中的像构成正则序列。
由 $A\to A\{x_1,\ldots,x_n\}$ 的平坦性以及代数之引理
\ref{algebra-lemma-grothendieck-regular-sequence}，
这表明 $f_1,\ldots,f_m$ 在
$A\{x_1,\ldots,x_m\}_\mathfrak q$ 中的像构成正则序列，
并且模这些元素所得的商在 $A$ 上平坦。证明完毕。
\end{proof}

\begin{lemma}
\label{restricted-lemma-exact-sequence-NL-rig-smooth}
设 $A\to B\to C$ 为 $\textit{WAdm}^{Noeth}$ 中的箭头，
它们都是 adic 的且拓扑有限型。
若 $B\to C$ 是 rig-光滑的，则映射
$$
H^{-1}(\NL_{B/A}^\wedge \otimes_B C) \to H^{-1}(\NL_{C/A}^\wedge)
$$
的核（见引理 \ref{restricted-lemma-exact-sequence-NL}）被某个定义理想零化。
\end{lemma}

\begin{proof}
设 $\overline{\mathfrak q}\subset C$ 为不包含定义理想的素理想。
由于所论各模是有限的，只需证明
$$
H^{-1}(\NL_{B/A}^\wedge \otimes_B C)_{\overline{\mathfrak q}} \to
H^{-1}(\NL_{C/A}^\wedge)_{\overline{\mathfrak q}}
$$
是单射。如同引理 \ref{restricted-lemma-exact-sequence-NL} 的证明，
选取表示
$B=A\{x_1,\ldots,x_r\}/J$、
$C=B\{y_1,\ldots,y_s\}/J'$ 和
$C=A\{x_1,\ldots,x_r,y_1,\ldots,y_s\}/K$。
观察引理 \ref{restricted-lemma-exact-sequence-NL} 的证明中的图表，
可知只需证明 $J/J^2\otimes_BC\to K/K^2$
在与 $\overline{\mathfrak q}$ 对应的素理想
$\mathfrak q\subset A\{x_1,\ldots,x_r,y_1,\ldots,y_s\}$
处局部化后是单射。请与更多代数之引理
\ref{more-algebra-lemma-transitive-lci-at-end} 及其证明比较。
这等价于要求 $J/KJ\to K/K^2$ 在 $\mathfrak q$ 处局部化后是单射；
又等价于证明
$J_\mathfrak q\cap K^2_\mathfrak q=(KJ)_\mathfrak q$。
由引理 \ref{restricted-lemma-rig-smooth-rig-flat}，已知
$(K/J)_\mathfrak q=J'_\mathfrak q$ 由正则序列生成。
因此，所需的交性质由更多代数之引理
\ref{more-algebra-lemma-conormal-sequence-H1-regular-ideal} 得出
（并使用正则序列生成的理想是 $H_1$-正则的这一事实，见更多代数之
第 \ref{more-algebra-section-ideals} 节）。
\end{proof}

\section{Rig-光滑态射}
\label{restricted-section-rig-smooth-morphisms}

\noindent
本节利用第 \ref{restricted-section-rig-smooth-homomorphisms} 节中的工作，
定义局部诺特代数空间的 rig-光滑态射。

\begin{definition}
\label{restricted-definition-rig-smooth}
设 $S$ 为概形。设 $f:X\to Y$ 为 $S$ 上局部诺特形式代数空间的态射。
若对每个交换图表
$$
\xymatrix{
U \ar[d] \ar[r] & V \ar[d] \\
X \ar[r] & Y
}
$$
其中 $U,V$ 是仿射形式代数空间，而 $U\to X$ 与 $V\to Y$
由代数空间可表且平展，态射 $U\to V$
对应于 adic 诺特拓扑环之间的 rig-光滑映射，
则称 $f$ 是{\it rig-光滑的}。
\end{definition}

\noindent
下面证明该条件可以在源和目标上平展局部地检验。

\begin{lemma}
\label{restricted-lemma-rig-smooth-morphisms}
设 $S$ 为概形。设 $f:X\to Y$ 为 $S$ 上局部诺特形式代数空间的态射。
以下条件等价：
\begin{enumerate}
\item $f$ 是 rig-光滑的；
\item 对每个交换图表
$$
\xymatrix{
U \ar[d] \ar[r] & V \ar[d] \\
X \ar[r] & Y
}
$$
其中 $U,V$ 是仿射形式代数空间，而 $U\to X$ 与 $V\to Y$
由代数空间可表且平展，态射 $U\to V$
对应于 $\textit{WAdm}^{Noeth}$ 中的 rig-光滑映射；
\item 存在形式空间之定义
\ref{formal-spaces-definition-formal-algebraic-space}
所述的覆盖 $\{Y_j\to Y\}$，并且对每个 $j$，
存在形式空间之定义
\ref{formal-spaces-definition-formal-algebraic-space}
所述的覆盖 $\{X_{ji}\to Y_j\times_YX\}$，
使每个 $X_{ji}\to Y_j$ 对应于
$\textit{WAdm}^{Noeth}$ 中的 rig-光滑映射；并且
\item 存在形式空间之定义
\ref{formal-spaces-definition-formal-algebraic-space}
所述的覆盖 $\{X_i\to X\}$，并且对每个 $i$，
存在分解 $X_i\to Y_i\to Y$，其中 $Y_i$ 是仿射形式代数空间，
$Y_i\to Y$ 由代数空间可表且平展，而
$X_i\to Y_i$ 对应于
$\textit{WAdm}^{Noeth}$ 中的 rig-光滑映射。
\end{enumerate}
\end{lemma}

\begin{proof}
(1) 与 (2) 的等价性就是定义 \ref{restricted-definition-rig-smooth}。
(2)、(3)、(4) 的等价性来自以下事实：
由引理 \ref{restricted-lemma-rig-smooth-axioms}，rig-光滑是
$\text{WAdm}^{Noeth}$ 箭头的局部性质；
再应用形式空间之引理
\ref{formal-spaces-lemma-property-defines-property-morphisms}
对局部诺特代数空间之间态射的变体，
该变体在形式空间之备注
\ref{formal-spaces-remark-variant-Noetherian} 中提及。
\end{proof}

\begin{lemma}
\label{restricted-lemma-base-change-rig-smooth}
设 $S$ 为概形。设 $f:X\to Y$ 与 $g:Z\to Y$
为 $S$ 上局部诺特形式代数空间的态射。
若 $f$ 是 rig-光滑的且 $g$ 是 adic 的，
则基变换 $X\times_YZ\to Z$ 是 rig-光滑的。
\end{lemma}

\begin{proof}
由形式空间之备注
\ref{formal-spaces-remark-base-change-variant-variant-Noetherian}
及形式空间之第 \ref{formal-spaces-section-adic} 节中的讨论，
这由引理 \ref{restricted-lemma-base-change-rig-smooth-continuous} 得出。
\end{proof}

\begin{lemma}
\label{restricted-lemma-composition-rig-smooth}
设 $S$ 为概形。设 $f:X\to Y$ 与 $g:Y\to Z$
为 $S$ 上局部诺特形式代数空间的态射。
若 $f$ 和 $g$ 都是 rig-光滑的，则 $g\circ f$ 也是。
\end{lemma}

\begin{proof}
由形式空间之备注
\ref{formal-spaces-remark-composition-variant-Noetherian}，
这由引理 \ref{restricted-lemma-composition-rig-smooth-continuous} 得出。
\end{proof}

\begin{lemma}
\label{restricted-lemma-rig-smooth-rig-flat-morphism}
设 $S$ 为概形。设 $f:X\to Y$ 为 $S$ 上局部诺特形式代数空间的态射。
若 $f$ 是 rig-光滑的，则 $f$ 是 rig-平坦的。
\end{lemma}

\begin{proof}
这由引理 \ref{restricted-lemma-rig-smooth-rig-flat} 和定义立即得出。
\end{proof}

\section{Rig-\,'etale 同态}
\label{restricted-section-rig-etale-homomorphisms}

\noindent
本节证明 adic 诺特拓扑环的 rig-\,'etale 同态的一些性质，
这些性质是引入局部诺特代数空间的 rig-\,'etale 态射所需要的。

\begin{lemma}
\label{restricted-lemma-rig-etale-continuous}
设 $A\to B$ 为 $\textit{WAdm}^{Noeth}$ 中的态射
（形式空间，第 \ref{formal-spaces-section-morphisms-rings} 节）。
下列条件等价：
\begin{enumerate}
\item[(a)] $A\to B$ 满足引理 \ref{restricted-lemma-finite-type} 中的等价条件，
并且存在定义理想 $I\subset B$，使 $B$ 在 $(A,I)$ 上 rig-\,'etale；以及
\item[(b)] $A\to B$ 满足引理 \ref{restricted-lemma-finite-type} 中的等价条件，
并且对每个定义理想 $I\subset A$，代数 $B$ 都在 $(A,I)$ 上
rig-\,'etale。
\end{enumerate}
\end{lemma}

\begin{proof}
设 $I$ 和 $I'$ 是 $A$ 的定义理想。则存在整数 $c\geq0$，使
$I^c\subset I'$ 且 $(I')^c\subset I$。因此，$B$ 在 $(A,I)$ 上
rig-\,'etale，当且仅当 $B$ 在 $(A,I')$ 上 rig-\,'etale。
这由定义 \ref{restricted-definition-rig-etale-homomorphism}、包含关系
$I^c\subset I'$ 与 $(I')^c\subset I$，以及下述事实得出：
由备注 \ref{restricted-remark-NL-well-defined-topological}，
朴素余切复形 $\NL_{B/A}^\wedge$ 不依赖于 $A$ 的定义理想的选择。
\end{proof}

\begin{definition}
\label{restricted-definition-rig-etale-continuous-homomorphism}
设 $\varphi:A\to B$ 为 adic 诺特拓扑环之间的连续环同态，
即 $\varphi$ 是 $\textit{WAdm}^{Noeth}$ 中的箭头。
若引理 \ref{restricted-lemma-rig-etale-continuous} 中的等价条件成立，
则称 $\varphi$ 是{\it rig-etale 的}。
\end{definition}

\noindent
这定义了一个局部性质。

\begin{lemma}
\label{restricted-lemma-rig-etale-axioms}
在 $\textit{WAdm}^{Noeth}$ 的箭头上，性质
$P(\varphi)=$“$\varphi$ 是 rig-\,'etale 的”是形式空间之备注
\ref{formal-spaces-remark-variant-Noetherian} 所定义的局部性质。
\end{lemma}

\begin{proof}
本证明与引理 \ref{restricted-lemma-rig-smooth-axioms} 的证明完全相同。
回顾一下该陈述的含义。首先，$\textit{WAdm}^{Noeth}$ 是以
adic 诺特拓扑环为对象、以连续环同态为态射的范畴。
考虑交换图表
$$
\xymatrix{
B \ar[r] & (B')^\wedge \\
A \ar[r] \ar[u]^\varphi & (A')^\wedge \ar[u]_{\varphi'}
}
$$
并假设它满足下列条件：$A$ 和 $B$ 是 adic 诺特拓扑环，
$A\to A'$ 和 $B\to B'$ 是 \'etale 环映射，
对于某个定义理想 $I\subset A$，
$(A')^\wedge=\lim A'/I^nA'$；
对于某个定义理想 $J\subset B$，
$(B')^\wedge=\lim B'/J^nB'$；并且
$\varphi:A\to B$ 与
$\varphi':(A')^\wedge\to(B')^\wedge$ 都连续。
注意，由形式空间之引理
\ref{formal-spaces-lemma-completion-in-sub}，
$(A')^\wedge$ 和 $(B')^\wedge$ 是 adic 诺特拓扑环。
我们须证明：
\begin{enumerate}
\item $\varphi$ 是 rig-\,'etale 的
$\Rightarrow\varphi'$ 是 rig-\,'etale 的；
\item 若 $B\to B'$ 忠实平坦，则 $\varphi'$ 是 rig-\,'etale 的
$\Rightarrow\varphi$ 是 rig-\,'etale 的；以及
\item 若对 $i=1,\ldots,n$，$A\to B_i$ 都是 rig-\,'etale 的，
则 $A\to\prod_{i=1,\ldots,n}B_i$ 是 rig-\,'etale 的。
\end{enumerate}
引理 \ref{restricted-lemma-finite-type} 中的等价条件满足条件 (1)、(2) 和 (3)。
因此，在对“rig-\,'etale”性质核验 (1)、(2) 和 (3) 时，
可以已经假设每种情形下的环映射都满足引理
\ref{restricted-lemma-finite-type} 中的等价条件。

\medskip\noindent
选取定义理想 $I\subset A$。由以上所述，图表中每个环的拓扑都是
$I$-进拓扑，并且 $B$、$(A')^\wedge$、$(B')^\wedge$
属于 $(A,I)$ 的范畴 (\ref{restricted-equation-C-prime})。
由于 $A\to A'$ 和 $B\to B'$ 都是 \'etale 的，复形
$\NL_{A'/A}$ 和 $\NL_{B'/B}$ 为零，故由引理
\ref{restricted-lemma-NL-is-completion}，
$\NL_{(A')^\wedge/A}^\wedge$ 和
$\NL_{(B')^\wedge/B}^\wedge$ 为零。
将引理 \ref{restricted-lemma-exact-sequence-NL} 应用于
$A\to(A')^\wedge\to(B')^\wedge$，得到同构
$$
H^i(\NL_{(B')^\wedge/(A')^\wedge}^\wedge) \to H^i(\NL_{(B')^\wedge/A}^\wedge)
$$
因此
$\NL_{(B')^\wedge/A}^\wedge\to
\NL_{(B')^\wedge/(A')^\wedge}$ 是拟同构。
环映射 $B/I^nB\to B'/I^nB'$ 是 \'etale 的，
因而是局部完全交
（代数，引理 \ref{algebra-lemma-etale-standard-smooth}）。
所以可将引理 \ref{restricted-lemma-exact-sequence-NL} 和
\ref{restricted-lemma-transitive-lci-at-end} 应用于
$A\to B\to(B')^\wedge$，得到同构
$$
H^i(\NL_{B/A}^\wedge \otimes_B (B')^\wedge) \to
H^i(\NL_{(B')^\wedge/A}^\wedge)
$$
由此可知
$\NL_{B/A}^\wedge\otimes_B(B')^\wedge\to
\NL_{(B')^\wedge/A}^\wedge$ 是拟同构。
结合这两个观察，得到 $D((B')^\wedge)$ 中的同构
$$
\NL_{(B')^\wedge/(A')^\wedge}^\wedge \cong
\NL_{B/A}^\wedge \otimes_B (B')^\wedge
$$
准备工作至此完成，可以开始真正的证明。

\medskip\noindent
证明 (1)。假设 $\varphi$ 是 rig-\,'etale 的。
则存在 $c\geq0$，使对 $a\in I^c$ 的乘法在 $D(B)$ 中的
$\NL_{B/A}^\wedge$ 上为零。这一性质在沿
$B\to(B')^\wedge$ 作基变换时保持，见更多代数之引理
\ref{more-algebra-lemma-two-term-base-change}。
由上述同构，$\varphi'$ 是 rig-\,'etale 的。这证明了 (1)。

\medskip\noindent
为证明 (2)，假设 $B\to B'$ 忠实平坦且 $\varphi'$ 是
rig-\,'etale 的。则存在 $c\geq0$，使对 $a\in I^c$ 的乘法
在 $D((B')^\wedge)$ 中的
$\NL_{(B')^\wedge/(A')^\wedge}^\wedge$ 上为零。
由上述同构，$a^c$ 零化
$\NL_{B/A}^\wedge\otimes_B(B')^\wedge$ 的上同调模。
由我们假设 $B\to B'$ 忠实平坦，复合
$B\to(B')^\wedge$ 忠实平坦，见形式空间之引理
\ref{formal-spaces-lemma-etale-surjective}。
因此，$\NL_{B/A}^\wedge$ 的上同调模被 $I^c$ 零化。
由引理 \ref{restricted-lemma-equivalent-with-artin}，
$\varphi$ 是 rig-\,'etale 的。这证明了 (2)。

\medskip\noindent
为证明 (3)，置 $B=\prod_{i=1,\ldots,n}B_i$；
只需注意，$\NL_{B/A}^\wedge$ 是复形
$\NL_{B_i/A}^\wedge$ 的直和，其中各复形均视为 $B$-模复形。
\end{proof}

\begin{lemma}
\label{restricted-lemma-base-change-rig-etale-continuous}
考虑 $\textit{WAdm}^{Noeth}$ 的箭头上的性质
$P(\varphi)=$“$\varphi$ 是 rig-\,'etale 的”与
$Q(\varphi)$=“$\varphi$ 是 adic 的”。
则依照形式空间之备注
\ref{formal-spaces-remark-base-change-variant-variant-Noetherian}
中的定义，$P$ 在由 $Q$ 作基变换下稳定。
\end{lemma}

\begin{proof}
由引理 \ref{restricted-lemma-rig-etale-continuous}，该陈述有意义。
为证明它，假设 $\textit{WAdm}^{Noeth}$ 中有态射
$B\to A$ 和 $B\to C$，且 $B\to A$ 是 rig-\,'etale 的，
$B\to C$ 是 adic 的
（形式空间，定义 \ref{formal-spaces-definition-adic-homomorphism}）。
则可选取定义理想 $I\subset B$，使 $A$ 和 $C$ 上的拓扑都是
$I$-进拓扑。在这种情形下，由引理
\ref{restricted-lemma-base-change-rig-etale-homomorphism} 立即可知，
$A\widehat{\otimes}_B C$ 在 $(C,IC)$ 上 rig-\,'etale。
\end{proof}

\begin{lemma}
\label{restricted-lemma-composition-rig-etale-continuous}
在 $\textit{WAdm}^{Noeth}$ 的箭头上，性质
$P(\varphi)=$“$\varphi$ 是 rig-\,'etale 的”在复合下稳定，
其含义如形式空间之备注
\ref{formal-spaces-remark-composition-variant-Noetherian} 所定义。
\end{lemma}

\begin{proof}
由引理 \ref{restricted-lemma-rig-etale-continuous}，该陈述有意义。
为证明它，假设 $\textit{WAdm}^{Noeth}$ 中有 rig-\,'etale 态射
$A\to B$ 和 $B\to C$。
则可选取定义理想 $I\subset A$，使 $C$ 和 $B$ 上的拓扑都是
$I$-进拓扑。由引理 \ref{restricted-lemma-exact-sequence-NL}，得到正合列
$$
\xymatrix{
C \otimes_B H^0(\NL_{B/A}^\wedge) \ar[r] &
H^0(\NL_{C/A}^\wedge) \ar[r] &
H^0(\NL_{C/B}^\wedge) \ar[r] & 0 \\
H^{-1}(\NL_{B/A}^\wedge \otimes_B C) \ar[r] &
H^{-1}(\NL_{C/A}^\wedge) \ar[r] &
H^{-1}(\NL_{C/B}^\wedge) \ar[llu]
}
$$
存在 $c\geq0$，使对每个 $a\in I$，乘以 $a^c$ 在
$D(B)$ 中的 $\NL_{B/A}^\wedge$ 以及在
$D(C)$ 中的 $\NL_{C/B}^\wedge$ 上都为零。
当然，乘以 $a^c$ 在 $D(C)$ 中的
$\NL_{B/A}^\wedge\otimes_BC$ 上也为零。
所以
$H^0(\NL_{B/A}^\wedge)\otimes_A C$、
$H^0(\NL_{C/B}^\wedge)$、
$H^{-1}(\NL_{B/A}^\wedge\otimes_BC)$ 和
$H^{-1}(\NL_{C/B}^\wedge)$ 都被 $a^c$ 零化。
由正合列可知，乘以 $a^{2c}$ 在
$H^0(\NL_{C/A}^\wedge)$ 和
$H^{-1}(\NL_{C/A}^\wedge)$ 上为零。
由引理 \ref{restricted-lemma-equivalent-with-artin}，
$C$ 按所需在 $(A,I)$ 上 rig-\,'etale。
\end{proof}

\begin{lemma}
\label{restricted-lemma-permanence-rig-etale-continuous}
在 $\textit{WAdm}^{Noeth}$ 的箭头上，性质
$P(\varphi)=$“$\varphi$ 是 rig-\,'etale 的”具有形式空间之备注
\ref{formal-spaces-remark-permanence-variant-Noetherian}
所定义的消去性质。
\end{lemma}

\begin{proof}
由引理 \ref{restricted-lemma-rig-etale-continuous}，该陈述有意义。
为证明它，假设 $\textit{WAdm}^{Noeth}$ 中有映射
$A\to B$ 和 $B\to C$，且 $A\to C$ 与 $A\to B$
都是 rig-\,'etale 的。我们须证明 $B\to C$ 是 rig-\,'etale 的。
可选取定义理想 $I\subset A$，使 $C$ 和 $B$ 上的拓扑都是
$I$-进拓扑。由引理 \ref{restricted-lemma-exact-sequence-NL}，得到正合列
$$
\xymatrix{
C \otimes_B H^0(\NL_{B/A}^\wedge) \ar[r] &
H^0(\NL_{C/A}^\wedge) \ar[r] &
H^0(\NL_{C/B}^\wedge) \ar[r] & 0 \\
H^{-1}(\NL_{B/A}^\wedge \otimes_B C) \ar[r] &
H^{-1}(\NL_{C/A}^\wedge) \ar[r] &
H^{-1}(\NL_{C/B}^\wedge) \ar[llu]
}
$$
存在 $c\geq0$，使对每个 $a\in I$，乘以 $a^c$ 在
$D(B)$ 中的 $\NL_{B/A}^\wedge$ 以及在
$D(C)$ 中的 $\NL_{C/A}^\wedge$ 上都为零。
所以
$H^0(\NL_{B/A}^\wedge)\otimes_A C$、
$H^0(\NL_{C/A}^\wedge)$ 和
$H^{-1}(\NL_{C/A}^\wedge)$ 都被 $a^c$ 零化。
由正合列可知，乘以 $a^{2c}$ 在
$H^0(\NL_{C/B}^\wedge)$ 和
$H^{-1}(\NL_{C/B}^\wedge)$ 上为零。
由引理 \ref{restricted-lemma-equivalent-with-artin}，
$C$ 按所需在 $(B,IB)$ 上 rig-\,'etale。
\end{proof}

\section{Rig-\,'etale 态射}
\label{restricted-section-rig-etale-morphisms}

\noindent
本节利用第 \ref{restricted-section-rig-etale-homomorphisms} 节中的工作，
定义局部诺特代数空间的 rig-\,'etale 态射。

\begin{definition}
\label{restricted-definition-rig-etale}
设 $S$ 为概形。设 $f:X\to Y$ 为 $S$ 上局部诺特形式代数空间的态射。
若对每个交换图表
$$
\xymatrix{
U \ar[d] \ar[r] & V \ar[d] \\
X \ar[r] & Y
}
$$
其中 $U,V$ 是仿射形式代数空间，而 $U\to X$ 与 $V\to Y$
由代数空间可表且平展，态射 $U\to V$
对应于 adic 诺特拓扑环之间的 rig-\,'etale 映射，
则称 $f$ 是{\it rig-\,'etale 的}。
\end{definition}

\noindent
下面证明该条件可以在源和目标上平展局部地检验。

\begin{lemma}
\label{restricted-lemma-rig-etale-morphisms}
设 $S$ 为概形。设 $f:X\to Y$ 为 $S$ 上局部诺特形式代数空间的态射。
以下条件等价：
\begin{enumerate}
\item $f$ 是 rig-\,'etale 的；
\item 对每个交换图表
$$
\xymatrix{
U \ar[d] \ar[r] & V \ar[d] \\
X \ar[r] & Y
}
$$
其中 $U,V$ 是仿射形式代数空间，而 $U\to X$ 与 $V\to Y$
由代数空间可表且平展，态射 $U\to V$
对应于 $\textit{WAdm}^{Noeth}$ 中的 rig-\,'etale 映射；
\item 存在形式空间之定义
\ref{formal-spaces-definition-formal-algebraic-space}
所述的覆盖 $\{Y_j\to Y\}$，并且对每个 $j$，
存在形式空间之定义
\ref{formal-spaces-definition-formal-algebraic-space}
所述的覆盖 $\{X_{ji}\to Y_j\times_YX\}$，
使每个 $X_{ji}\to Y_j$ 对应于
$\textit{WAdm}^{Noeth}$ 中的 rig-\,'etale 映射；并且
\item 存在形式空间之定义
\ref{formal-spaces-definition-formal-algebraic-space}
所述的覆盖 $\{X_i\to X\}$，并且对每个 $i$，
存在分解 $X_i\to Y_i\to Y$，其中 $Y_i$ 是仿射形式代数空间，
$Y_i\to Y$ 由代数空间可表且平展，而
$X_i\to Y_i$ 对应于
$\textit{WAdm}^{Noeth}$ 中的 rig-\,'etale 映射。
\end{enumerate}
\end{lemma}

\begin{proof}
(1) 与 (2) 的等价性就是定义 \ref{restricted-definition-rig-etale}。
(2)、(3)、(4) 的等价性来自以下事实：
由引理 \ref{restricted-lemma-rig-etale-axioms}，rig-\,'etale 是
$\text{WAdm}^{Noeth}$ 箭头的局部性质；
再应用形式空间之引理
\ref{formal-spaces-lemma-property-defines-property-morphisms}
对局部诺特代数空间之间态射的变体，
该变体在形式空间之备注
\ref{formal-spaces-remark-variant-Noetherian} 中提及。
\end{proof}

\noindent
特别要说明的是，rig-\,'etale 态射局部有限型。

\begin{lemma}
\label{restricted-lemma-rig-etale-finite-type}
局部诺特形式代数空间的 rig-\,'etale 态射局部有限型。
\end{lemma}

\begin{proof}
引理 \ref{restricted-lemma-rig-etale-axioms} 中的性质 $P$
蕴含形式空间之引理
\ref{formal-spaces-lemma-quotient-restricted-power-series}
中的等价条件 (a)、(b)、(c)、(d)。
所以结论由形式空间之引理
\ref{formal-spaces-lemma-finite-type-local-property} 得出。
\end{proof}

\begin{lemma}
\label{restricted-lemma-rig-etale-rig-smooth-morphism}
局部诺特形式代数空间的 rig-\,'etale 态射是 rig-光滑的。
\end{lemma}

\begin{proof}
这由定义和引理 \ref{restricted-lemma-rig-etale-rig-smooth} 得出。
\end{proof}

\begin{lemma}
\label{restricted-lemma-base-change-rig-etale}
设 $S$ 为概形。设 $f:X\to Y$ 与 $g:Z\to Y$
为 $S$ 上局部诺特形式代数空间的态射。
若 $f$ 是 rig-\,'etale 的且 $g$ 是 adic 的，
则基变换 $X\times_YZ\to Z$ 是 rig-\,'etale 的。
\end{lemma}

\begin{proof}
由形式空间之备注
\ref{formal-spaces-remark-base-change-variant-variant-Noetherian}
及形式空间之第 \ref{formal-spaces-section-adic} 节中的讨论，
这由引理 \ref{restricted-lemma-base-change-rig-etale-continuous} 得出。
\end{proof}

\begin{lemma}
\label{restricted-lemma-composition-rig-etale}
设 $S$ 为概形。设 $f:X\to Y$ 与 $g:Y\to Z$
为 $S$ 上局部诺特形式代数空间的态射。
若 $f$ 和 $g$ 都是 rig-\,'etale 的，则 $g\circ f$ 也是。
\end{lemma}

\begin{proof}
由形式空间之备注
\ref{formal-spaces-remark-composition-variant-Noetherian}，
这由引理 \ref{restricted-lemma-composition-rig-etale-continuous} 得出。
\end{proof}

\begin{lemma}
\label{restricted-lemma-rig-etale-permanence}
设 $S$ 为概形。设 $f:X\to Y$ 与 $g:Y\to Z$
为 $S$ 上局部诺特形式代数空间的态射。
若 $g\circ f$ 和 $g$ 都是 rig-\,'etale 的，则 $f$ 也是。
\end{lemma}

\begin{proof}
由形式空间之备注
\ref{formal-spaces-remark-permanence-variant-Noetherian}，
这由引理 \ref{restricted-lemma-permanence-rig-etale-continuous} 得出。
\end{proof}

\begin{lemma}
\label{restricted-lemma-rig-etale-alternative-permanence}
设 $S$ 为概形。设 $f:X\to Y$ 与 $g:Y\to Z$
为 $S$ 上局部诺特形式代数空间的态射。
若 $g\circ f$ 是 rig-\,'etale 的且 $g$ 是 adic 单态射，
则 $f$ 是 rig-\,'etale 的。
\end{lemma}

\begin{proof}
使用引理 \ref{restricted-lemma-base-change-rig-etale} 以及如下事实：
$f$ 是 $g\circ f$ 沿 $g$ 的基变换。
\end{proof}

\begin{lemma}
\label{restricted-lemma-closed-immersion-rig-smooth}
设 $S$ 为概形。设 $f:X\to Y$ 为形式代数空间的态射。
假设 $X$ 和 $Y$ 局部诺特，且 $f$ 是闭浸入。下列条件等价：
\begin{enumerate}
\item $f$ 是 rig-光滑的；
\item $f$ 是 rig-\,'etale 的；
\item 对每个仿射形式代数空间 $V$ 以及每个由代数空间可表且平展的
态射 $V\to Y$，态射 $X\times_YV\to V$ 对应于
$\textit{WAdm}^{Noeth}$ 中的满态射 $B\to A$，其核 $J$
具有如下性质：对 $B$ 的某个定义理想 $I$，有 $I(J/J^2)=0$。
\end{enumerate}
\end{lemma}

\begin{proof}
先作如下观察：给定 (2) 中的 $V$ 和 $V\to Y$，而不对 $f$
施加任何进一步假设，由形式空间之引理
\ref{formal-spaces-lemma-closed-immersion-into-countably-indexed}，
可知态射 $X\times_YV\to V$ 对应于
$\textit{WAdm}^{Noeth}$ 中的满态射 $B\to A$。

\medskip\noindent
由引理 \ref{restricted-lemma-rig-etale-rig-smooth-morphism}，
(2) $\Rightarrow$ (1)。

\medskip\noindent
证明 (3) $\Rightarrow$ (2)。假设 (3)。由引理
\ref{restricted-lemma-rig-etale-morphisms}，只需证明 (3) 中出现的环映射
$B\to A$ 在定义
\ref{restricted-definition-rig-etale-continuous-homomorphism} 的意义下
是 rig-\,'etale 的。设 $I$ 如 (3) 中所述。
$A$ 在 $(B,I)$ 上的朴素余切复形
$\NL_{A/B}^\wedge$ 是把 $J/J^2$ 放在次数 $-1$ 所得的 $A$-模复形。
因此，由定义 \ref{restricted-definition-rig-etale-homomorphism}，
$A$ 在 $(B,I)$ 上 rig-\,'etale。

\medskip\noindent
假设 (1)，并设 $V$ 和 $B\to A$ 如 (3) 中所述。
由定义 \ref{restricted-definition-rig-smooth}，$B\to A$ 是 rig-光滑的。
任取定义理想 $I\subset B$。则 $A$ 在 $(B,I)$ 上 rig-光滑。
如上所述，复形 $\NL_{A/B}^\wedge$ 是把 $J/J^2$
放在次数 $-1$ 所得的复形。因此，由引理
\ref{restricted-lemma-equivalent-with-artin-smooth}，
$J/J^2$ 被某个 $n\geq1$ 的幂 $I^n$ 零化。
由于 $B$ 是 adic 的，$I^n$ 是 $B$ 的定义理想。证明完毕。
\end{proof}

\section{Rig-满射态射}
\label{restricted-section-rig-surjective}

\noindent
对于局部诺特形式代数空间之间的局部有限型态射，
可以采用一个源自 \cite{ArtinII} 的定义。
关于更一般情形中的处理方式，见备注
\ref{restricted-remark-rig-surjective-more-general} 中的讨论。

\begin{definition}
\label{restricted-definition-rig-surjective}
设 $S$ 为概形。设 $f:X\to Y$ 为 $S$ 上形式代数空间的态射。
假设 $X$ 和 $Y$ 局部诺特，且 $f$ 局部有限型。
如果对每个实线图表
$$
\xymatrix{
\text{Spf}(R') \ar@{..>}[r] \ar@{..>}[d] & X \ar[d]^f \\
\text{Spf}(R) \ar[r]^-p & Y
}
$$
其中 $R$ 是完备离散赋值环且 $p$ 是 adic 态射，
都存在完备离散赋值环的扩张 $R\subset R'$ 和态射
$\text{Spf}(R')\to X$，使所示图表交换，
则称 $f$ 是{\it rig-满射的}。
\end{definition}

\noindent
下面的引理将说明，在局部诺特形式代数空间和局部有限型态射的语境中，
这一概念表现得相当合理。下一个备注讨论如何修改该定义，
使之适用于更广泛的形式代数空间态射。

\begin{remark}
\label{restricted-remark-rig-surjective-more-general}
定义 \ref{restricted-definition-rig-surjective} 中表述的条件，
甚至对局部 adic* 形式代数空间的有限型态射也不合适。
例如，若
$A=(\bigcup_{n\geq1}k[t^{1/n}])^\wedge$，
其中完备化是 $t$-进完备化，则不存在 adic 态射
$\text{Spf}(R)\to\text{Spf}(A)$，其中 $R$ 是完备离散赋值环。
因而任意态射 $X\to\text{Spf}(A)$ 都会是 rig-满射的；
但是，由于 $A$ 是整环且 $t\in A$ 非零，我们希望把 $A$
看作至少具有一个“rig-点”，并不希望允许 $X=\emptyset$。
为了涵盖这一特定情形，可以考虑 adic 态射
$$
\text{Spf}(R) \longrightarrow Y
$$
其中 $R$ 是关于主理想 $J$ 完备的赋值环，并且
$\mathfrak m_R=\sqrt{J}$。在这种情形下，$R$ 的值群可嵌入
$(\mathbf{R},+)$，由此得到 Berkovich 在定义与 $Y$ 相关联的
解析空间时采用的观点，见 \cite{Berkovich}。
另一种方法由 Huber 倡导。在他的理论中，舍去
$\Spec(R/J)$ 是单点的假设，见 \cite{Huber-continuous-valuations}。
\end{remark}

\begin{lemma}
\label{restricted-lemma-composition-rig-surjective}
\begin{slogan}
局部有限型态射的 rig-满射性在复合下保持
\end{slogan}
设 $S$ 为概形。设 $f:X\to Y$ 与 $g:Y\to Z$
为 $S$ 上形式代数空间的态射。
假设 $X,Y,Z$ 局部诺特，且 $f,g$ 局部有限型。
若 $f,g$ 都是 rig-满射的，则 $g\circ f$ 也是。
\end{lemma}

\begin{proof}
这直接由定义得出
（并使用形式空间之引理
\ref{formal-spaces-lemma-composition-finite-type}）。
\end{proof}

\begin{lemma}
\label{restricted-lemma-base-change-rig-surjective}
设 $S$ 为概形。设 $f:X\to Y$ 与 $Z\to Y$
为 $S$ 上形式代数空间的态射。
假设 $X,Y,Z$ 局部诺特，且 $f,g$ 局部有限型。
若 $f$ 是 rig-满射的，则基变换 $Z\times_YX\to Z$ 也是。
\end{lemma}

\begin{proof}
这直接由定义得出（并使用形式空间之引理
\ref{formal-spaces-lemma-fibre-product-Noetherian} 和
\ref{formal-spaces-lemma-base-change-finite-type}）。
\end{proof}

\begin{lemma}
\label{restricted-lemma-rig-surjective-alternative-permanence}
设 $S$ 为概形。设 $f:X\to Y$ 与 $g:Y\to Z$
为 $S$ 上局部诺特形式代数空间的局部有限型态射。
若 $g\circ f$ 是 rig-满射的且 $g$ 是单态射，
则 $f$ 是 rig-满射的。
\end{lemma}

\begin{proof}
使用引理 \ref{restricted-lemma-base-change-rig-surjective} 以及如下事实：
$f$ 是 $g\circ f$ 沿 $g$ 的基变换。
\end{proof}

\begin{lemma}
\label{restricted-lemma-permanence-rig-surjective}
设 $S$ 为概形。设 $f:X\to Y$ 与 $g:Y\to Z$
为 $S$ 上形式代数空间的态射。
假设 $X,Y,Z$ 局部诺特，且 $f,g$ 局部有限型。
若 $g\circ f:X\to Z$ 是 rig-满射的，则
$g:Y\to Z$ 也是。
\end{lemma}

\begin{proof}
这由定义立即得出。
\end{proof}

\begin{lemma}
\label{restricted-lemma-etale-covering-rig-surjective}
设 $S$ 为概形。设 $f:X\to Y$ 为局部诺特形式代数空间的态射，
并假设它由代数空间可表、平展且满射。则 $f$ 是 rig-满射的。
\end{lemma}

\begin{proof}
设 $p:\text{Spf}(R)\to Y$ 为 adic 态射，其中 $R$ 是完备离散赋值环。
令 $Z=\text{Spf}(R)\times_YX$。则
$Z\to\text{Spf}(R)$ 由代数空间可表、平展且满射，
故 $Z$ 非空。选取非空仿射形式代数空间 $V$ 以及平展态射
$V\to Z$（按我们的定义，这是可以做到的）。
于是 $V\to\text{Spf}(R)$ 对应于 $R\to A^\wedge$，
其中 $R\to A$ 是平展环映射，见形式空间之引理
\ref{formal-spaces-lemma-etale}。由于 $A^\wedge\not=0$
（因为 $V\not=\emptyset$），可以找到位于 $\mathfrak m_R$ 上方的
$A$ 的极大理想 $\mathfrak m$。于是 $A_\mathfrak m$ 是离散赋值环
（更多代数，引理
\ref{more-algebra-lemma-Dedekind-etale-extension}）。
所以 $R'=A_\mathfrak m^\wedge$ 是完备离散赋值环
（更多代数，引理 \ref{more-algebra-lemma-completion-dvr}）。
应用形式空间之引理
\ref{formal-spaces-lemma-morphism-between-formal-spectra}，
得到所需态射 $\text{Spf}(R')\to V\to Z\to X$。
\end{proof}

\noindent
以上引理的结论是：可以在 $Y$ 上平展局部地检验
$f:X\to Y$ 是否 rig-满射。

\begin{lemma}
\label{restricted-lemma-upshot}
设 $S$ 为概形。设 $f:X\to Y$ 为局部诺特形式代数空间的
局部有限型态射。设 $\{g_i:Y_i\to Y\}$ 为形式代数空间态射族，
其中各态射由代数空间可表且平展，并且 $\coprod g_i$ 满射。
则 $f$ 是 rig-满射的，当且仅当每个
$f_i:X\times_YY_i\to Y_i$ 都是 rig-满射的。
\end{lemma}

\begin{proof}
事实上，若 $f$ 是 rig-满射的，则任意基变换也是
（引理 \ref{restricted-lemma-base-change-rig-surjective}）。
反之，若所有 $f_i$ 都是 rig-满射的，则
$\coprod f_i:\coprod X\times_YY_i\to\coprod Y_i$ 也是。
由引理 \ref{restricted-lemma-etale-covering-rig-surjective}，态射
$\coprod g_i:\coprod Y_i\to Y$ 是 rig-满射的。
所以 $\coprod X\times_YY_i\to Y$ 是 rig-满射的
（引理 \ref{restricted-lemma-composition-rig-surjective}）。
由于该态射经过 $X\to Y$ 分解，由引理
\ref{restricted-lemma-permanence-rig-surjective}，$X\to Y$ 是 rig-满射的。
\end{proof}

\begin{lemma}
\label{restricted-lemma-faithfully-flat-rig-surjective}
设 $A$ 为关于理想 $I$ 完备的诺特环。
设 $B$ 为 $I$-进完备的 $A$-代数。
若对每个 $n$，$A/I^n\to B/I^nB$ 都是有限型且平坦的，
并且在 $n=1$ 时忠实平坦，则
$\text{Spf}(B)\to\text{Spf}(A)$ 是 rig-满射的。
\end{lemma}

\begin{proof}
下文将不再另行说明地使用如下事实：形式谱之间的态射由相应拓扑环之间的
连续映射给出，见形式空间之引理
\ref{formal-spaces-lemma-morphism-between-formal-spectra}。
设 $\varphi:A\to R$ 为到完备离散赋值环 $A$ 的连续映射。
这蕴含 $\varphi(I)\subset\mathfrak m_R$。
另一方面，由于只须在 $\varphi$ 对应于 adic 态射时构造提升
$\varphi':B'\to R'$，可以假设 $\varphi(I)\not=0$。
因此可考虑基变换 $C=B\widehat{\otimes}_A R$，
例如见备注 \ref{restricted-remark-base-change}。
于是 $C$ 是 $\mathfrak m_R$-进完备的 $R$-代数，
$C/\mathfrak m_R^nC$ 在 $R/\mathfrak m_R^n$ 上有限型且平坦，
并且 $C/\mathfrak m_RC$ 非零。
任取位于 $\mathfrak m_R$ 上方的极大理想 $\mathfrak m\subset C$。
由平坦性（它蕴含下降）可知
$\Spec(C_\mathfrak m)\setminus V(\mathfrak m_RC_\mathfrak m)$
是非空开集。因此可以选取素理想
$\mathfrak q\subset\mathfrak m$，使 $\mathfrak q$ 给出
$\Spec(C_\mathfrak m)\setminus\{\mathfrak m\}$ 的闭点，
且 $\mathfrak q\not\in V(IC_\mathfrak m)$，见性质之引理
\ref{properties-lemma-complement-closed-point-Jacobson}。
于是 $C/\mathfrak q$ 是 $1$ 维局部整环，并且存在包含
$C/\mathfrak q\subset R'$，其中 $R'$ 是离散赋值环
（代数，引理 \ref{algebra-lemma-exists-dvr}）。
按构造，$\mathfrak m_RR'\subset\mathfrak m_{R'}$，
并且 $C\to R'$ 延拓为连续映射 $C\to(R')^\wedge$
（事实上，在当前情形下可选择 $R'$ 使
$R'=(R')^\wedge$，但我们不需要这一点）。
由于离散赋值环的完备化仍是离散赋值环，该假设给出环的交换图表
$$
\xymatrix{
(R')^\wedge & C \ar[l] & B \ar[l] \\
R \ar[u] & R \ar[l] \ar[u] & A \ar[l] \ar[u]
}
$$
它给出所需的提升。
\end{proof}

\begin{lemma}
\label{restricted-lemma-flat-rig-surjective}
设 $A$ 为关于理想 $I$ 完备的诺特环。
设 $B$ 为 $I$-进完备的 $A$-代数。假设：
\begin{enumerate}
\item $A$ 中的 $I$-挠为 $0$；
\item 对每个 $n$，$A/I^n\to B/I^nB$ 都平坦且有限型。
\end{enumerate}
则 $\text{Spf}(B)\to\text{Spf}(A)$ 是 rig-满射的，
当且仅当 $A/I\to B/IB$ 忠实平坦。
\end{lemma}

\begin{proof}
忠实平坦性由引理 \ref{restricted-lemma-faithfully-flat-rig-surjective}
蕴含 rig-满射性。为证明逆命题，下文将不再另行说明地使用如下事实：
$I$-挠消失等价于 $I$-幂挠消失
（更多代数，引理 \ref{more-algebra-lemma-torsion-free}）。
还将不再另行说明地使用如下事实：形式谱之间的态射由相应拓扑环之间的
连续映射给出，见形式空间之引理
\ref{formal-spaces-lemma-morphism-between-formal-spectra}。

\medskip\noindent
假设 $\text{Spf}(B)\to\text{Spf}(A)$ 是 rig-满射的。
选取极大理想 $I\subset\mathfrak m\subset A$。
$\Spec(A_\mathfrak m)$ 的开子集
$U=\Spec(A_\mathfrak m)\setminus V(I_\mathfrak m)$ 非空，
因为 $A_\mathfrak m$ 的 $I_\mathfrak m$-挠为零
（使用代数之引理
\ref{algebra-lemma-Noetherian-power-ideal-kills-module}）。
因此可以找到素理想 $\mathfrak q\subset A_\mathfrak m$，
它给出 $U$ 的一个点（即
$IA_\mathfrak m\not\subset\mathfrak q$），并对应于
$\Spec(A_\mathfrak m)\setminus\{\mathfrak m\}$ 的闭点，见性质之引理
\ref{properties-lemma-complement-closed-point-Jacobson}。
于是 $A_\mathfrak m/\mathfrak q$ 是 $1$ 维局部整环。
所以存在局部环的单局部同态
$A_\mathfrak m/\mathfrak q\subset R$，其中 $R$ 是离散赋值环
（代数，引理 \ref{algebra-lemma-exists-dvr}）。
按构造，$IR\subset\mathfrak m_R$，并且
$A\to R$ 延拓为连续映射 $A\to R^\wedge$。
由于离散赋值环的完备化仍是离散赋值环，该假设给出环的交换图表
$$
\xymatrix{
R' & B \ar[l] \\
R^\wedge \ar[u] & A \ar[l] \ar[u]
}
$$
因而找到 $B$ 中位于 $\mathfrak m$ 上方的素理想。
所以 $\Spec(B/IB)\to\Spec(A/I)$ 满射，进而
$A/I\to B/IB$ 忠实平坦
（代数，引理 \ref{algebra-lemma-ff-rings}）。
\end{proof}

\begin{lemma}
\label{restricted-lemma-monomorphism-rig-surjective}
设 $S$ 为概形。设 $f:X\to Y$ 为形式代数空间的态射。
假设 $X,Y$ 局部诺特，$f$ 局部有限型，且 $f$ 是单态射。
则 $f$ 是 rig-满射的，当且仅当每个 adic 态射
$\text{Spf}(R)\to Y$ 都经过 $X$ 分解，其中 $R$ 是完备离散赋值环。
\end{lemma}

\begin{proof}
一个方向是平凡的。为证明另一个方向，假设
$\text{Spf}(R)\to Y$ 是 adic 态射，并且存在完备离散赋值环的扩张
$R\subset R'$，使经过 $Y$ 分解的
$\text{Spf}(R')\to\text{Spf}(R)\to X$ 存在。
于是
$\Spec(R'/\mathfrak m_R^nR')\to\Spec(R/\mathfrak m_R^n)$
满且平坦；因 $X$ 满足 fpqc 覆盖的层条件，态射
$\Spec(R/\mathfrak m_R^n)\to X$ 经过 $X$ 分解，
见形式空间之引理 \ref{formal-spaces-lemma-sheaf-fpqc}。
换言之，$\text{Spf}(R)\to Y$ 经过 $X$ 分解。
\end{proof}

\begin{lemma}
\label{restricted-lemma-closed-immersion-rig-surjective}
设 $S$ 为概形。设 $f:X\to Y$ 为形式代数空间的态射。
假设 $X,Y$ 局部诺特，且 $f$ 是闭浸入。下列条件等价：
\begin{enumerate}
\item $f$ 是 rig-满射的；以及
\item 对每个仿射形式代数空间 $V$ 以及每个由代数空间可表且平展的
态射 $V\to Y$，态射 $X\times_YV\to V$ 对应于
$\textit{WAdm}^{Noeth}$ 中的满态射 $B\to A$，其核 $J$
具有如下性质：对 $B$ 的某个定义理想 $I$ 和某个 $n\geq1$，
有 $IJ^n=0$。
\end{enumerate}
\end{lemma}

\begin{proof}
先作如下观察：给定 (2) 中的 $V$ 和 $V\to Y$，而不对 $f$
施加任何进一步假设，由形式空间之引理
\ref{formal-spaces-lemma-closed-immersion-into-countably-indexed}，
可知态射 $X\times_YV\to V$ 对应于
$\textit{WAdm}^{Noeth}$ 中的满态射 $B\to A$。

\medskip\noindent
假设 (1)。由引理 \ref{restricted-lemma-base-change-rig-surjective}，
$\text{Spf}(A)\to\text{Spf}(B)$ 是 rig-满射的。
设 $I\subset B$ 为定义理想。由于 $B$ 是 adic 的，
对所有 $m\geq1$，$I^m\subset B$ 都是定义理想。
若对所有 $n,m\geq1$ 都有 $I^mJ^n\not=0$，
则 $IJ$ 不幂零，故 $V(IJ)\not=\Spec(B)$。
因而可找到素理想 $\mathfrak p\subset B$，使
$\mathfrak p\not\in V(I)\cup V(J)$。
注意 $I(B/\mathfrak p)\not=B/\mathfrak p$，故可找到极大理想
$\mathfrak p+I\subset\mathfrak m\subset B$。
由代数之引理 \ref{algebra-lemma-exists-dvr}，
可找到离散赋值环 $R$ 和单局部环同态
$(B/\mathfrak p)_\mathfrak m\to R$。
显然，环映射 $B\to R$ 不可能经过 $A=B/J$ 分解。
依照引理 \ref{restricted-lemma-monomorphism-rig-surjective}，
这与 $\text{Spf}(A)\to\text{Spf}(B)$ 是 rig-满射的相矛盾。
因此，对某些 $n,m$，确有 $I^nJ^m=0$，这证明 (2)。

\medskip\noindent
假设 (2)。由引理 \ref{restricted-lemma-upshot}，只需证明
$\text{Spf}(A)\to\text{Spf}(B)$ 是 rig-满射的。
选取定义理想 $I\subset B$ 和整数 $n$，使 $IJ^n=0$。
考虑环映射 $B\to R$，其中 $R$ 是离散赋值环，且 $I$ 的像非零。
由于 $R$ 是整环，$J$ 在 $R$ 中的像为零。
所以 $B\to R$ 经过满射 $B\to A$ 分解，结论由
rig-满射态射的定义得出。
\end{proof}

\begin{lemma}
\label{restricted-lemma-closed-immersion-rig-smooth-rig-surjective}
设 $S$ 为概形。设 $f:X\to Y$ 为形式代数空间的态射。
假设 $X,Y$ 局部诺特，且 $f$ 是闭浸入。下列条件等价：
\begin{enumerate}
\item $f$ 是 rig-光滑且 rig-满射的；
\item $f$ 是 rig-\,'etale 且 rig-满射的；以及
\item 对每个仿射形式代数空间 $V$ 以及每个由代数空间可表且平展的
态射 $V\to Y$，态射 $X\times_YV\to V$ 对应于
$\textit{WAdm}^{Noeth}$ 中的满态射 $B\to A$，其核 $J$
具有如下性质：对 $B$ 的某个定义理想 $I$，有 $IJ=0$。
\end{enumerate}
\end{lemma}

\begin{proof}
设 $I,J$ 为环 $B$ 的理想，使 $IJ^n=0$ 且 $I(J/J^2)=0$。
则 $I^nJ=0$（证明从略）。所以本引理由引理
\ref{restricted-lemma-closed-immersion-rig-smooth} 和
\ref{restricted-lemma-closed-immersion-rig-surjective} 的平凡组合得出。
\end{proof}

\begin{lemma}
\label{restricted-lemma-rig-etale-descent}
设 $S$ 为概形。设 $f:X\to Y$ 与 $g:Y\to Z$
为 $S$ 上局部诺特形式代数空间的态射。假设：
\begin{enumerate}
\item $g$ 局部有限型；
\item $f$ 是 rig-光滑（相应地，rig-\,'etale）且 rig-满射的；
\item $g\circ f$ 是 rig-光滑的（相应地，rig-\,'etale 的）。
\end{enumerate}
则 $g$ 是 rig-光滑的（相应地，rig-\,'etale 的）。
\end{lemma}

\begin{proof}
我们将证明 rig-光滑情形，并在证明末尾指出证明 rig-\,'etale 情形
所需的改动。考虑交换图表
$$
\xymatrix{
X \times_Y V \ar[r] \ar[d] &
V \ar[d] \ar[r] &
W \ar[d] \\
X \ar[r] &
Y \ar[r] &
Z
}
$$
其中 $V,W$ 是仿射形式代数空间，而 $V\to Y$ 与 $W\to Z$
由代数空间可表且平展。我们须证明 $V\to W$ 对应于
adic 诺特拓扑环之间的 rig-光滑映射，见定义
\ref{restricted-definition-rig-smooth}。
可以写成 $V=\text{Spf}(B)$ 和 $W=\text{Spf}(C)$，
并且 $V\to W$ 对应于 adic 环映射 $C\to B$；
该映射拓扑有限型，见引理 \ref{restricted-lemma-finite-type-morphisms}。

\medskip\noindent
下文将不再另行说明地使用如下事实：
$X\times_YV\to V$ 是 rig-光滑且 rig-满射的，见引理
\ref{restricted-lemma-base-change-rig-smooth} 和
\ref{restricted-lemma-base-change-rig-surjective}。
此外，复合 $X\times_YV\to V\to W$ 是 rig-光滑的，
因为 $g\circ f$ 是 rig-光滑的。

\medskip\noindent
设 $I\subset C$ 为定义理想。模
反设 $C\to B$ 不是 rig-光滑的，以求矛盾。
这意味着存在不包含 $IB$ 的素理想 $\mathfrak q\subset B$，
使得或者 $H^{-1}(\NL_{B/C}^\wedge)_\mathfrak p$ 非零，
或者 $H^0(\NL_{B/C}^\wedge)_\mathfrak p$ 不是有限自由
$B_\mathfrak q$-模。见引理
\ref{restricted-lemma-equivalent-with-artin-smooth}；略去若干细节。
可以选择极大理想 $IB+\mathfrak q\subset\mathfrak m$。
由代数之引理 \ref{algebra-lemma-exists-dvr}，
可以找到完备离散赋值环 $R$ 和单局部环同态
$(B/\mathfrak q)_\mathfrak m\to R$。

\medskip\noindent
用一个扩张代替 $R$ 后，可以假设 adic 态射
$\text{Spf}(R)\to V=\text{Spf}(B)$ 有提升
$\text{Spf}(R)\to X\times_YV$。
选取形式空间之定义
\ref{formal-spaces-definition-formal-algebraic-space}
所述的平展覆盖
$\{\text{Spf}(A_i)\to X\times_YV\}$。
由引理 \ref{restricted-lemma-etale-covering-rig-surjective}，
可以假设 $\text{Spf}(R)\to X\times_YV$ 提升为某个 $i$ 的态射
$\text{Spf}(R)\to\text{Spf}(A_i)$
（这可能需要再次用一个扩张代替 $R$）。
置 $A=A_i$。考虑环映射
$$
C \to B \to A \to R
$$
设 $\mathfrak p\subset A$ 为映射 $A\to R$ 的核；
注意 $\mathfrak p$ 位于 $\mathfrak q$ 上方。
已知 $C\to A$ 和 $B\to A$ 都是 rig-光滑的。
特别地，由引理 \ref{restricted-lemma-rig-smooth-rig-flat}，环映射
$B_\mathfrak q\to A_\mathfrak p$ 平坦。
考虑引理 \ref{restricted-lemma-exact-sequence-NL} 和
\ref{restricted-lemma-exact-sequence-NL-rig-smooth} 的相应正合列
$$
\xymatrix{
&
H^0(\NL_{B/C}^\wedge) \otimes_B A_\mathfrak p \ar[r] &
H^0(\NL_{A/C}^\wedge)_\mathfrak p \ar[r] &
H^0(\NL_{A/B}^\wedge)_\mathfrak p \ar[r] & 0 \\
0 \ar[r] &
H^{-1}(\NL_{B/C}^\wedge \otimes_B A)_\mathfrak p \ar[r] &
H^{-1}(\NL_{A/C}^\wedge)_\mathfrak p \ar[r] &
H^{-1}(\NL_{A/B}^\wedge)_\mathfrak p \ar[llu]
}
$$
由 $C\to A$ 和 $B\to A$ 的 rig-光滑性可知
$H^{-1}(\NL_{B/C}^\wedge\otimes_BA)_\mathfrak p=0$，
并且 $H^0(\NL_{B/C}^\wedge)\otimes_BA_\mathfrak p$
作为有限自由 $A_\mathfrak p$-模之间满射的核，是有限自由的。
由于 $B_\mathfrak q\to A_\mathfrak p$ 平坦，因而忠实平坦，
这蕴含
$H^{-1}(\NL_{B/C}^\wedge)_\mathfrak q=0$，
并且 $H^0(\NL_{B/C}^\wedge)_\mathfrak q$ 有限自由，
这正是所求的矛盾。

\medskip\noindent
在 rig-\,'etale 情形下，论证完全相同，
只是所得结论是
$H^{-1}(\NL_{B/C}^\wedge)_\mathfrak q$ 和
$H^0(\NL_{B/C}^\wedge)_\mathfrak q$ 都为零。
\end{proof}

\section{完备离散赋值环上的形式代数空间}
\label{restricted-section-over-cdrv}

\noindent
本节采用如下术语：若 $A$ 是弱可容许拓扑环，
则“{\it $X$ 是 $A$ 上的形式代数空间}”意指 $X$ 是一个形式代数空间，
并配备形式代数空间的态射
$p:X\to\text{Spf}(A)$。在这种情形下，称 $p$ 为{\it 结构态射}。

\begin{lemma}
\label{restricted-lemma-flat-locus}
设 $X$ 为完备离散赋值环 $A$ 上的局部诺特形式代数空间。
则存在形式代数空间的闭浸入 $X'\to X$，使 $X'$ 在 $A$ 上平坦，
并且任何局部诺特形式代数空间的态射 $Y\to X$，只要 $Y$
在 $A$ 上平坦，就经过 $X'$ 分解。
\end{lemma}

\begin{proof}
设 $\pi\in A$ 为一致化元。回顾：$A$-模平坦，
当且仅当其 $\pi$-幂挠为 $0$。

\medskip\noindent
先假设 $X$ 是仿射形式代数空间。
则 $X=\text{Spf}(B)$，其中 $B$ 是 adic 诺特 $A$-代数。
在这种情形下，令 $X'=\text{Spf}(B')$，其中
$B'=B/\pi\text{-幂挠子}$。
显然，$X'$ 在 $A$ 上平坦，且 $X'\to X$ 是闭浸入。
设 $g:Y\to X$ 为局部诺特形式代数空间的态射，
其中 $Y$ 在 $A$ 上平坦。选取形式空间之定义
\ref{formal-spaces-definition-formal-algebraic-space}
所述的覆盖 $\{Y_j\to Y\}$。
于是 $Y_j=\text{Spf}(C_j)$，其中 $C_j$ 在 $A$ 上平坦。
因而态射 $Y_j\to X$ 对应于连续 $R$-代数映射 $B\to C_j$；
由于 $B\to C_j$ 显然零化 $\pi$-幂挠，它经过 $X'$ 分解。
由于 $\{Y_j\to Y\}$ 是覆盖且 $X'\to X$ 是单态射，
可知 $g$ 经过 $X'$ 分解。

\medskip\noindent
设 $X$ 和 $\{X_i\to X\}_{i\in I}$ 如形式空间之定义
\ref{formal-spaces-definition-formal-algebraic-space} 中所述。
对每个 $i$，令 $X'_i\to X_i$ 为上述构造的平坦部分。
对 $i,j\in I$，投影 $X'_i\times_XX_j\to X'_i$
是概形的平展态射（前者由假设，后者由形式空间之引理
\ref{formal-spaces-lemma-presentation-representable}）。
因而 $X'_i\times_XX_j$ 在 $A$ 上平坦，因为由代数空间可表且平展的
态射是平坦的（引理 \ref{restricted-lemma-representable-flat}）。
所以由泛性质，投影 $X'_i\times_XX_j\to X_j$ 经过 $X'_j$ 分解。
由于态射 $X'_i\to X_i$ 是层的单射，得到
$$
R_{ij} = X'_i \times_X X_j = X'_i \times_X X'_j = X_i \times_X X'_j
$$
置 $U=\coprod X'_i$、$R=\coprod R_{ij}$，并以
$s,t:R\to U$ 表示两个投影。作为层，$R=U\times_XU$，
且 $s,t$ 都平展。由以上观察，$(t,s):R\to U$ 定义了平展等价关系。
所以由空间之定理 \ref{spaces-theorem-presentation}，
$X'=U/R$ 是代数空间。按构造，图表
$$
\xymatrix{
\coprod X'_i \ar[r] \ar[d] & \coprod X_i \ar[d] \\
X' \ar[r] & X
}
$$
是笛卡尔的。由于右侧竖直箭头平展满射，顶部水平箭头可表且为闭浸入，
由 Bootstrap 之引理
\ref{bootstrap-lemma-after-fppf-sep-lqf}，可知 $X'\to X$ 可表。
随后使用空间之引理
\ref{spaces-lemma-descent-representable-transformations-property}，
可知 $X'\to X$ 是闭浸入。

\medskip\noindent
最后，设 $Y\to X$ 为态射，其中 $Y$ 是在 $A$ 上平坦的
局部诺特形式代数空间。则每个 $X_i\times_XY$ 在 $Y$ 上平展，
因而在 $A$ 上平坦（见上文）。所以
$X_i\times_XY\to X_i$ 经过 $X'_i$ 分解。
由于 $\{X_i\times_XY\to Y\}$ 是平展覆盖，
$Y\to X$ 经过 $X'$ 分解。
\end{proof}

\begin{lemma}
\label{restricted-lemma-flat-and-diagonal-rig-surjective}
设 $X$ 为完备离散赋值环 $A$ 上的局部诺特形式代数空间，
且在该环上局部有限型。设 $X'\subset X$ 如引理
\ref{restricted-lemma-flat-locus} 中所述。
若 $X\to X\times_{\text{Spf}(A)}X$ 是 rig-\,'etale 且 rig-满射的，
则 $X'=\text{Spf}(A)$ 或 $X'=\emptyset$。
\end{lemma}

\begin{proof}
（顺便说明：由形式空间之引理
\ref{formal-spaces-lemma-diagonal-morphism-formal-algebraic-spaces}，
对角总是局部有限型；而由形式空间之引理
\ref{formal-spaces-lemma-base-change-finite-type} 和
\ref{formal-spaces-lemma-locally-finite-type-locally-noetherian}，
$X\times_{\text{Spf}(A)}X$ 局部诺特。
因此，对角态射上的这些条件是有意义的。）图表
$$
\xymatrix{
X' \ar[r] \ar[d] & X' \times_{\text{Spf}(A)} X' \ar[d] \\
X \ar[r] & X \times_{\text{Spf}(A)} X
}
$$
是笛卡尔的。因此，由引理 \ref{restricted-lemma-base-change-rig-surjective}，
$X'\to X'\times_{\text{Spf}(A)}X'$ 是 rig-\,'etale 且 rig-满射的。
选取仿射形式代数空间 $U$ 和由代数空间可表且平展的态射
$U\to X'$。则 $U=\text{Spf}(B)$，其中 $B$ 是 adic 诺特拓扑环，
是平坦 $A$-代数；其拓扑是 $\pi$-进拓扑，其中
$\pi\in A$ 是一致化元；并且对每个 $n$，
$A/\pi^nA\to B/\pi^nB$ 有限型。
为后用，注意这尤其蕴含：若 $B\not=0$，则映射
$\text{Spf}(B)\to\text{Spf}(A)$ 是层的满射
（请回顾，我们一贯使用 fppf 拓扑）。重复上述论证，可知
$$
W = U \times_{X'} U =
X' \times_{X' \times_{\text{Spf}(A)} X'} (U \times_{\text{Spf}(A)} U)
\longrightarrow
U \times_{\text{Spf}(A)} U
$$
是闭浸入，并且 rig-\,'etale 且 rig-满射。
由形式空间之引理
\ref{formal-spaces-lemma-fibre-product-affines-over-separated}，有
$U\times_{\text{Spf}(A)}U=
\text{Spf}(B\widehat{\otimes}_AB)$。
由于 $B\widehat{\otimes}_AB$ 是平坦 $A$-代数
$B\otimes_AB$ 的 $\pi$-进完备化，它是平坦 $A$-代数。
所以由引理 \ref{restricted-lemma-closed-immersion-rig-smooth-rig-surjective}，
$W=U\times_{\text{Spf}(A)}U$。
换言之，$U\times_{X'}U=U\times_{\text{Spf}(A)}U$；
这又意味着 $U\to X'$ 的像（视为层映射）单射地映到
$\text{Spf}(A)$。选取形式空间之定义
\ref{formal-spaces-definition-formal-algebraic-space}
所述的覆盖 $\{U_i\to X'\}$。
特别地，$\coprod U_i\to X'$ 是层的满射。
把以上论证应用于 $U_i\coprod U_j\to X'$
（使用 $U_i\amalg U_j$ 仍是仿射形式代数空间这一事实），
可知 $X'\to\text{Spf}(A)$ 是 fppf 层的单射。
由于 $X'$ 在 $A$ 上平坦，或者 $X'$ 为空
（若对所有 $i$，$U_i$ 都为空），或者该映射是同构
（若对某个 $i$，$U_i$ 非空；此时已看到
$U_i\to\text{Spf}(A)$ 是层的满射）。证明完毕。
\end{proof}

\begin{lemma}
\label{restricted-lemma-rig-monomorphism-rig-surjective}
设 $S$ 为概形。设 $f:X\to Y$ 为形式代数空间的态射。假设：
\begin{enumerate}
\item $X$ 和 $Y$ 局部诺特；
\item $f$ 局部有限型；
\item $\Delta_f:X\to X\times_YX$ 是 rig-\,'etale 且 rig-满射的。
\end{enumerate}
则 $f$ 是 rig-满射的，当且仅当每个 adic 态射
$\text{Spf}(R)\to Y$ 都提升为态射 $\text{Spf}(R)\to X$，
其中 $R$ 是完备离散赋值环。
\end{lemma}

\begin{proof}
一个方向是平凡的。为证明另一个方向，假设
$\text{Spf}(R)\to Y$ 是 adic 态射，并且存在完备离散赋值环的扩张
$R\subset R'$，使经过 $Y$ 分解的
$\text{Spf}(R')\to\text{Spf}(R)\to X$ 存在。
考虑纤维积图表
$$
\xymatrix{
\text{Spf}(R') \ar[r] \ar[rd] &
\text{Spf}(R) \times_Y X \ar[r] \ar[d]^p &
X \ar[d]^f \\
&
\text{Spf}(R) \ar[r] &
Y
}
$$
态射 $p$ 作为 $f$ 的基变换局部有限型，见形式空间之引理
\ref{formal-spaces-lemma-base-change-finite-type}。
对角态射 $\Delta_p$ 是 $\Delta_f$ 的基变换，
因而 rig-\,'etale 且 rig-满射。
由引理 \ref{restricted-lemma-flat-and-diagonal-rig-surjective}，
$\text{Spf}(R)\times_YX$ 在 $R$ 上的平坦轨迹或者是
$\emptyset$，或者等于 $\text{Spf}(R)$。
但是，由于 $\text{Spf}(R')$ 经过它分解，可知它非空，
因而按所需得到态射
$\text{Spf}(R)\to\text{Spf}(R)\times_YX\to X$。
\end{proof}

\section{完备化函子}
\label{restricted-section-completion-functor}

\noindent
本节考虑下列情形。首先固定基概形 $S$。
所有环、拓扑环、概形、代数空间、形式代数空间及其间态射都在 $S$ 上。
其次固定代数空间 $X$ 和闭子集 $T\subset|X|$。
以 $U\subset X$ 表示满足 $|U|=|X|\setminus T$ 的开子空间。图示为
$$
U \to X \quad |X| = |U| \amalg T
$$
在这种情形下，给定 $X$ 上的代数空间 $X'$，即配备态射
$f:X'\to X$ 的代数空间 $X'$，以 $T'\subset|X'|$
表示 $T$ 的逆像，并令 $U'\subset X'$ 为满足
$|U'|=|X'|\setminus T'$ 的开子空间。图示为
$$
U' = f^{-1}U
\quad\quad
\vcenter{
\xymatrix{
U' \ar[d] \ar[r] & X' \ar[d]_f \\
U \ar[r] & X
}
}
\quad\quad
\vcenter{
\xymatrix{
|U'| \ar[r] \ar[d] & |X'| \ar[d]^{|f|} & T' \ar[l] \ar[d] \\
|U| \ar[r] & |X| & T \ar[l]
}
}
\quad\quad
T' = |f|^{-1}T
$$
我们将把 $f$ 的性质同形式完备化的诱导态射
$$
f_{/T} : X'_{/T'} \longrightarrow X_{/T}
$$
的性质联系起来。如所示公式所表明的，以 $f_{/T}$ 表示该态射。
我们已在形式空间之引理
\ref{formal-spaces-lemma-map-completions-representable}
中看到 $f_{/T}$ 由代数空间可表。
事实上，如该引理的证明所示，图表
$$
\xymatrix{
X'_{/T'} \ar[d]_{f_{/T}} \ar[r] & X' \ar[d]^f \\
X_{/T} \ar[r] & X
}
$$
是笛卡尔的。阅读下面陈述并证明的各引理时，请记住这一事实。

\begin{lemma}
\label{restricted-lemma-map-completions-finite-type}
在上述情形中，若 $f$ 局部有限型，则 $f_{/T}$ 局部有限型。
\end{lemma}

\begin{proof}
（形式代数空间的有限型态射在形式空间之第
\ref{formal-spaces-section-finite-type} 节中讨论。）
事实上，设 $Z\to X$ 为从概形到 $X$ 的态射，
使 $|Z|$ 映入 $T$。由上述笛卡尔方形，
$Z\times_XX'$ 是表示 $Z\times_{X_{/T}}X'_{/T'}$ 的代数空间。
由于 $Z\times_XX'\to Z$ 由空间的态射之引理
\ref{spaces-morphisms-lemma-base-change-finite-type}
局部有限型，结论得证。
\end{proof}

\begin{lemma}
\label{restricted-lemma-map-completions-etale}
在上述情形中，若 $f$ 平展，则 $f_{/T}$ 平展。
\end{lemma}

\begin{proof}
采用引理 \ref{restricted-lemma-map-completions-finite-type} 证明中的相同论证，
结论由空间的态射之引理
\ref{spaces-morphisms-lemma-base-change-etale} 得出。
\end{proof}

\begin{lemma}
\label{restricted-lemma-closed-immersion-gives-closed-immersion}
在上述情形中，若 $f$ 是闭浸入，则 $f_{/T}$ 是闭浸入。
\end{lemma}

\begin{proof}
（形式代数空间的闭浸入在形式空间之第
\ref{formal-spaces-section-closed-immersions} 节中讨论。）
采用引理 \ref{restricted-lemma-map-completions-finite-type} 证明中的相同论证，
结论由空间之引理
\ref{spaces-lemma-base-change-immersions} 得出。
\end{proof}

\begin{lemma}
\label{restricted-lemma-proper-gives-proper}
在上述情形中，若 $f$ 是真的，则 $f_{/T}$ 是真的。
\end{lemma}

\begin{proof}
（形式代数空间的真态射在形式空间之第
\ref{formal-spaces-section-proper} 节中讨论。）
采用引理 \ref{restricted-lemma-map-completions-finite-type} 证明中的相同论证，
结论由空间的态射之引理
\ref{spaces-morphisms-lemma-base-change-proper} 得出。
\end{proof}

\begin{lemma}
\label{restricted-lemma-quasi-compact-gives-quasi-compact}
在上述情形中，若 $f$ 拟紧，则 $f_{/T}$ 拟紧。
\end{lemma}

\begin{proof}
（形式代数空间的拟紧态射在形式空间之第
\ref{formal-spaces-section-quasi-compact} 节中讨论。）
我们须证明
$(X'_{/T'})_{red}\to(X_{/T})_{red}$ 是代数空间的拟紧态射。
由形式空间之引理
\ref{formal-spaces-lemma-reduction-completion}，
这是态射 $Z'\to Z$，其中 $Z'\subset X'$（相应地，$Z\subset X$）
是 $T'$（相应地，$T$）上的约化诱导代数空间结构。
因此 $Z'\to f^{-1}Z=Z\times_XX'$ 是加厚
（即在底拓扑空间上诱导同构的闭浸入）。
由于 $Z\times_XX'\to Z$ 作为 $f$ 的基变换拟紧
（空间的态射，引理
\ref{spaces-morphisms-lemma-base-change-quasi-compact}），
由更多空间的态射之引理
\ref{spaces-more-morphisms-lemma-thicken-property-morphisms}，
$Z'\to Z$ 也拟紧。
\end{proof}

\begin{remark}
\label{restricted-remark-diagonal-gives-diagonal}
在上述情形中，考虑对角态射
$\Delta_f:X'\to X'\times_XX'$ 和
$\Delta_{f_{/T}}:X'_{/T'}\to
X'_{/T'}\times_{X_{/T}}X'_{/T'}$。
容易看出
$$
X'_{/T'} \times_{X_{/T}} X'_{/T'} = (X' \times_X X')_{/T''}
$$
作为 $X'\times_XX'$ 的子函子成立，其中
$T''\subset|X'\times_XX'|$ 是 $T$ 的逆像。
所以 $\Delta_{f_{/T}}=(\Delta_f)_{/T''}$。
下面将用此说明 $\Delta_f$ 的性质由 $\Delta_{f_{/T}}$ 继承。
\end{remark}

\begin{lemma}
\label{restricted-lemma-quasi-separated-gives-quasi-separated}
在上述情形中，若 $f$ 是（拟）分离的，则 $f_{/T}$ 也是。
\end{lemma}

\begin{proof}
（形式代数空间态射的分离性条件在形式空间之第
\ref{formal-spaces-section-separation-axioms} 节中讨论。）
我们须证明：若 $\Delta_f$ 拟紧（相应地，是闭浸入），
则 $\Delta_{f_{/T}}$ 也如此。
这由备注 \ref{restricted-remark-diagonal-gives-diagonal} 中的讨论以及引理
\ref{restricted-lemma-quasi-compact-gives-quasi-compact} 和
\ref{restricted-lemma-closed-immersion-gives-closed-immersion} 得出。
\end{proof}

\begin{lemma}
\label{restricted-lemma-smooth-gives-rig-smooth}
在上述情形中，若 $X$ 局部诺特，$f$ 局部有限型，
且 $U'\to U$ 光滑，则 $f_{/T}$ 是 rig-光滑的。
\end{lemma}

\begin{proof}
证明策略如下：归结到 $X,X'$ 都仿射的情形，
把仿射情形翻译为代数，最后应用引理 \ref{restricted-lemma-rig-smooth}。
我们建议读者略过细节。

\medskip\noindent
选取满平展态射 $W\to X$，其中 $W=\coprod W_i$
是仿射概形的不交并，见空间的性质之引理
\ref{spaces-properties-lemma-cover-by-union-affines}。
对每个 $i$，选取满平展态射
$W'_i\to W_i\times_XX'$，其中
$W'_i=\coprod W'_{ij}$ 是仿射概形的不交并。
特别地，$\coprod W'_{ij}\to X'$ 满且平展。
以 $f_{ij}:W_{ij}\to W_i$ 表示给定态射。
以 $T_i\subset W_i$ 和 $T'_{ij}\subset W_{ij}$ 表示 $T$ 的逆像。
由于沿 $T$ 的逆像取完备化产生笛卡尔图表（见上文），有
$(W_i)_{/T_i}=W_i\times_XX_{/T}$，并且类似地，
$(W'_{ij})_{/T'_{ij}}=W'_{ij}\times_{X'}X'_{/T'}$。
此外，回顾 $(W_i)_{/T_i}$ 和 $(W'_{ij})_{/T'_{ij}}$
是仿射形式代数空间。因此
$\{W'_{ij})_{/T'_{ij}}\to X'_{/T'}\}$ 是形式空间之定义
\ref{formal-spaces-definition-formal-algebraic-space} 所述的覆盖。
由引理 \ref{restricted-lemma-rig-smooth-morphisms}，只需证明
$$
(W'_{ij})_{/T'_{ij}} \longrightarrow (W_i)_{/T_i}
$$
是 rig-光滑的。注意，$W'_{ij}\to W_i$ 局部有限型，
并诱导光滑态射
$W'_{ij}\setminus T'_{ij}\to W_i\setminus T_i$
（因为 $f$ 具有这些性质，且这些态射性质在源和目标上平展局部）。
又注意 $W_i$ 局部诺特（因为 $X$ 局部诺特，
且该性质在代数空间上平展局部）。
所以只需证明 $X,X'$ 都为仿射概形时的引理。

\medskip\noindent
假设 $X=\Spec(A)$ 和 $X'=\Spec(A')$ 是仿射概形。
由于 $X$ 诺特，$A$ 是诺特环。
态射 $f$ 由有限型环映射 $A\to A'$ 给出。
设 $I\subset A$ 为截出 $T$ 的理想，则 $IA'$ 截出 $T'$。
此外，$\Spec(A')\to\Spec(A)$ 在
$\Spec(A)\setminus T$ 上光滑。
设 $A^\wedge$ 和 $(A')^\wedge$ 为 $I$-进完备化。
则 $X_{/T}=\text{Spf}(A^\wedge)$ 且
$X'_{/T'}=\text{Spf}((A')^\wedge)$，见形式空间之引理
\ref{formal-spaces-lemma-formal-completion-types} 的证明。
由引理 \ref{restricted-lemma-rig-smooth}，$(A')^\wedge$ 在 $(A.I)$ 上
rig-光滑；这又意味着 $A^\wedge\to(A')^\wedge$ 是 rig-光滑的；
最后，由引理 \ref{restricted-lemma-rig-smooth-morphisms}，
$X'_{/T'}\to X_{/T}$ 是 rig-光滑的。
\end{proof}

\begin{lemma}
\label{restricted-lemma-etale-gives-rig-etale}
在上述情形中，若 $X$ 局部诺特，$f$ 局部有限型，
且 $U'\to U$ 平展，则 $f_{/T}$ 是 rig-\,'etale 的。
\end{lemma}

\begin{proof}
证明与引理 \ref{restricted-lemma-smooth-gives-rig-smooth} 的证明完全相同，
只是把引理 \ref{restricted-lemma-rig-smooth} 和
\ref{restricted-lemma-rig-smooth-morphisms} 替换为引理
\ref{restricted-lemma-rig-etale} 和 \ref{restricted-lemma-rig-etale-morphisms}。
\end{proof}

\begin{lemma}
\label{restricted-lemma-completion-proper-surjective-rig-surjective}
在上述情形中，若 $X$ 局部诺特，$f$ 是真的，
且 $U'\to U$ 满射，则 $f_{/T}$ 是 rig-满射的。
\end{lemma}

\begin{proof}
（该陈述由引理 \ref{restricted-lemma-map-completions-finite-type} 和
形式空间之引理
\ref{formal-spaces-lemma-formal-completion-types} 是有意义的。）
设 $R$ 为分式域为 $K$ 的完备离散赋值环。
设 $p:\text{Spf}(R)\to X_{/T}$ 为形式代数空间的 adic 态射。
由形式空间之引理
\ref{formal-spaces-lemma-adic-into-completion}，
复合 $\text{Spf}(R)\to X_{/T}\to X$ 对应于态射
$q:\Spec(R)\to X$，它把 $\Spec(K)$ 映入 $U$。
由于 $U'\to U$ 真且满射，
$\Spec(K)\times_UU'$ 非空且在 $K$ 上真。
所以可以选取域扩张 $K'/K$ 和交换图表
$$
\xymatrix{
\Spec(K') \ar[r] \ar[d] & U' \ar[r] \ar[d] & X' \ar[d] \\
\Spec(K) \ar[r] & U \ar[r] & X
}
$$
设 $R'\subset K'$ 为支配 $R$ 且分式域为 $K'$ 的离散赋值环，
见代数之引理 \ref{algebra-lemma-exists-dvr}。
由于 $\Spec(K)\to X$ 延拓为 $\Spec(R)\to X$，
由真性的赋值判据
（空间的态射，引理
\ref{spaces-morphisms-lemma-characterize-proper}），
可以把 $U'$ 的 $K'$-点延拓为位于
$\Spec(R)\to X$ 上方的态射 $\Spec(R')\to X'$。
于是 $T'$ 在 $\Spec(R')$ 中的逆像是闭点，并得到提升 $p$ 的
adic 态射 $\text{Spf}((R')^\wedge)\to X'_{/T'}$，如所需
（注意，由更多代数之引理
\ref{more-algebra-lemma-completion-dvr}，
$(R')^\wedge$ 是完备离散赋值环）。
\end{proof}

\begin{lemma}
\label{restricted-lemma-separated-mono-open-diagonal-rig-surjective}
在上述情形中，若 $X$ 局部诺特，$f$ 分离且局部有限型，
并且 $U'\to U$ 是单态射，则
$\Delta_{f_{/T}}$ 是 rig-满射的。
\end{lemma}

\begin{proof}
对角 $\Delta_f:X'\to X'\times_XX'$ 是闭浸入，
而 $\Delta_f$ 的限制 $U'\to U'\times_UU'$ 满射。
所以结论由备注 \ref{restricted-remark-diagonal-gives-diagonal} 中的讨论和引理
\ref{restricted-lemma-completion-proper-surjective-rig-surjective} 得出。
\end{proof}

\section{形式修改}
\label{restricted-section-formal-modifications}

\noindent
本节定义并研究 Artin 的局部诺特形式代数空间之形式修改概念。
先给出定义。

\begin{definition}
\label{restricted-definition-formal-modification}
设 $S$ 为概形。设 $f:X\to Y$ 为 $S$ 上局部诺特形式代数空间的态射。
若满足下列条件，则称 $f$ 为{\it 形式修改}：
\begin{enumerate}
\item $f$ 是真态射（形式空间，定义
\ref{formal-spaces-definition-proper}）；
\item $f$ 是 rig-\,'etale 的；
\item $f$ 是 rig-满射的；
\item $\Delta_f:X\to X\times_YX$ 是 rig-满射的。
\end{enumerate}
\end{definition}

\noindent
引理 \ref{restricted-lemma-modification-gives-formal-modification}
给出一个典型例子；事实上，稍后将证明每个形式修改在“形式局部”
意义下都属于这一类型，见引理
\ref{restricted-lemma-formal-modifications-locally-algebraic}。
下面把这些条件与 Artin 论文中的条件作比较。

\begin{remark}
\label{restricted-remark-compare-formal-modification-artin}
在 \cite[Definition 1.7]{ArtinII} 中，形式修改被定义为局部诺特
形式代数空间的真态射 $f:X\to Y$，并满足下列三个条件\footnote{我们不会把
这些条件完全翻译为 Stacks project 中发展出的语言。
尽管如此，我们希望这里的讨论对读者有所帮助。}：
\begin{enumerate}
\item[(\romannumeral1)] $f$ 的 Cramer 理想和 Jacobian 理想
各自包含 $X$ 的一个定义理想；
\item[(\romannumeral2)] 定义对角映射
$\Delta:X\to X\times_YX$ 的理想被
$X\times_YX$ 的一个定义理想零化；以及
\item[(\romannumeral3)] 每当 $R$ 是完备离散赋值环时，
任意 adic 态射 $\text{Spf}(R)\to Y$ 都提升为
$\text{Spf}(R)\to X$。
\end{enumerate}
下面把这些条件与我们上列的条件作比较。

\medskip\noindent
关于 (\romannumeral1)。性质 (\romannumeral1) 与我们要求
$f$ 为 rig-\,'etale 态射的条件一致；这由引理
\ref{restricted-lemma-equivalent-with-artin} 的
(\ref{restricted-item-condition-artin}) 得出。

\medskip\noindent
关于 (\romannumeral2)。假设 $f$ 是 rig-\,'etale 的。
则 $\Delta_f:X\to X\times_YX$ 是局部诺特形式代数空间的
rig-\,'etale 态射，因为这两个空间都在 $X$ 上 rig-\,'etale
（第一个经由 $\text{id}_X$，第二个经由 $\text{pr}_1$）。
见引理 \ref{restricted-lemma-base-change-rig-etale} 和
\ref{restricted-lemma-rig-etale-permanence}。
所以由引理 \ref{restricted-lemma-closed-immersion-rig-smooth-rig-surjective}，
性质 (\romannumeral2) 与我们要求 $\Delta_f$ rig-满射的条件一致。

\medskip\noindent
关于 (\romannumeral3)。性质 (\romannumeral3) 与我们的
rig-满射态射概念并不完全一致，因为 Artin 要求所有 adic 态射
$\text{Spf}(R)\to Y$ 都提升为到 $X$ 的态射，而我们的 rig-满射概念
仅断言用一个扩张代替 $R$ 后存在提升。
不过，由 (\romannumeral1) 和 (\romannumeral2)，已经知道
$\Delta_f$ 是 rig-\,'etale 且 rig-满射的，因此由引理
\ref{restricted-lemma-rig-monomorphism-rig-surjective}，这些条件等价。
\end{remark}

\begin{lemma}
\label{restricted-lemma-modification-gives-formal-modification}
设 $S$、$f:X'\to X$、$T\subset|X|$、$U\subset X$、
$T'\subset|X'|$ 和 $U'\subset X'$ 如第
\ref{restricted-section-completion-functor} 节中所述。
若 $X$ 局部诺特，$f$ 是真的，且 $U'\to U$ 是同构，
则 $f_{/T}:X'_{/T'}\to X_{/T}$ 是形式修改。
\end{lemma}

\begin{proof}
由形式空间之引理
\ref{formal-spaces-lemma-formal-completion-types}，
该箭头的源和目标是局部诺特形式代数空间。
其余条件由引理 \ref{restricted-lemma-proper-gives-proper}、
\ref{restricted-lemma-etale-gives-rig-etale}、
\ref{restricted-lemma-completion-proper-surjective-rig-surjective} 和
\ref{restricted-lemma-separated-mono-open-diagonal-rig-surjective} 得出。
\end{proof}

\begin{lemma}
\label{restricted-lemma-base-change-formal-modification}
设 $S$ 为概形。设 $f:X\to Y$ 为 $S$ 上局部诺特形式代数空间的态射，
并且是形式修改。则对局部诺特形式代数空间的任意 adic 态射
$Y'\to Y$，基变换
$f':X\times_YY'\to Y'$ 是形式修改。
\end{lemma}

\begin{proof}
由形式空间之引理
\ref{formal-spaces-lemma-base-change-proper}，态射 $f'$ 是真的。
由引理 \ref{restricted-lemma-base-change-rig-etale}，态射 $f'$ 是 rig-etale 的。
于是由引理 \ref{restricted-lemma-base-change-rig-surjective}，
态射 $f'$ 是 rig-满射的。置 $X'=X\times_ Y'$。
态射 $\Delta_{f'}$ 是 $\Delta_f$ 沿 adic 态射
$X'\times_{Y'}X'\to X\times_YX$ 的基变换。
所以由引理 \ref{restricted-lemma-base-change-rig-surjective}，
$\Delta_{f'}$ 是 rig-满射的。
\end{proof}

\section{完备化与态射，I}
\label{restricted-section-completion-and-morphisms}

\noindent
本节给出一些关于完备化的预备结果，将用于定理
\ref{restricted-theorem-dilatations-general} 的证明。
尽管这里陈述并证明的引理并非平凡
（其中一些建立在我们关于 rig-\,'etale 代数之代数化的工作上），
我们仍建议读者初次阅读时跳过本节。

\begin{lemma}
\label{restricted-lemma-algebraize-rig-etale-affine}
设 $T\subset X$ 为诺特仿射概形 $X$ 的闭子集。
设 $W$ 为诺特仿射形式代数空间。
设 $g:W\to X_{/T}$ 为 rig-\,'etale 态射。
则存在仿射概形 $X'$ 和有限型态射 $f:X'\to X$，
该态射在 $X\setminus T$ 上平展，并且存在与 $f_{/T}$ 和 $g$
相容的同构 $X'_{/f^{-1}T}\cong W$。
此外，若 $W\to X_{/T}$ 平展，则 $X'\to X$ 平展。
\end{lemma}

\begin{proof}
$X'$ 的存在性是引理
\ref{restricted-lemma-approximate-by-etale-over-complement} 的重述。
最后一个陈述由更多态射之引理
\ref{more-morphisms-lemma-check-smoothness-on-infinitesimal-nbhds} 得出。
\end{proof}

\begin{lemma}
\label{restricted-lemma-algebraize-morphism-rig-etale}
假设给定：
\begin{enumerate}
\item 诺特仿射概形 $X,X',Y$；
\item 闭子集 $T\subset|X|$；
\item 局部有限型态射 $f:X'\to X$，它在 $X\setminus T$ 上平展；
\item 态射 $h:Y\to X$；
\item 位于 $X_{/T}$ 上的态射
$\alpha:Y_{/T}\to X'_{/T}$（记号见证明）。
\end{enumerate}
则存在仿射概形的平展态射 $b:Y'\to Y$，
它诱导同构 $b_{/T}:Y'_{/T}\to Y_{/T}$；
并且存在位于 $X$ 上的态射 $a:Y'\to X'$，使
$\alpha=a_{/T}\circ b_{/T}^{-1}$。
\end{lemma}

\begin{proof}
陈述中使用下标 ${}_{/T}$ 的记号指如下构造：
它把概形态射 $g:V\to X$ 关联到形式代数空间的态射
$g_{/T}:V_{/g^{-1}T}\to X_{/T}$；这是从 $X$ 上概形的范畴
到 $X_{/T}$ 上形式代数空间的范畴的函子，见第
\ref{restricted-section-completion-functor} 节。
说明这一点后，本引理只是引理
\ref{restricted-lemma-fully-faithful-etale-over-complement} 的改述。
\end{proof}

\begin{lemma}
\label{restricted-lemma-factor}
设 $S$ 为概形。设 $f:X\to Y$ 与 $g:Z\to Y$ 为代数空间的态射。
设 $T\subset|X|$ 闭。假设：
\begin{enumerate}
\item $X$ 局部诺特；
\item $g$ 是单态射且局部有限型；
\item $f|_{X\setminus T}:X\setminus T\to Y$ 经过 $g$ 分解；
\item $f_{/T}:X_{/T}\to Y$ 经过 $g$ 分解。
\end{enumerate}
则 $f$ 经过 $g$ 分解。
\end{lemma}

\begin{proof}
考虑纤维积 $E=X\times_YZ\to X$。
由假设，开浸入 $X\setminus T\to X$ 经过 $E$ 分解；
而任何满足 $|\varphi|(|X'|)\subset T$ 的态射
$\varphi:X'\to X$ 也经过 $E$ 分解，见形式空间之第
\ref{formal-spaces-section-completion} 节。
由更多空间的态射之引理
\ref{spaces-more-morphisms-lemma-check-smoothness-on-infinitesimal-nbhds}，
这蕴含 $E\to X$ 在 $E$ 中每个映到 $T$ 中某点的点处平展。
所以 $E\to X$ 是平展单态射，因而是开浸入
（空间的态射，引理
\ref{spaces-morphisms-lemma-etale-universally-injective-open}）。
于是由我们的假设蕴含 $|X|=|E|$，可知 $E=X$。
\end{proof}

\begin{lemma}
\label{restricted-lemma-faithful-general}
设 $S$ 为概形。设 $X,W$ 为 $S$ 上代数空间，其中 $X$ 局部诺特。
设 $T\subset|X|$ 为闭子集。
设 $a,b:X\to W$ 为 $S$ 上代数空间的态射，满足
$a|_{X\setminus T}=b|_{X\setminus T}$，并且作为态射
$X_{/T}\to W$ 有 $a_{/T}=b_{/T}$。则 $a=b$。
\end{lemma}

\begin{proof}
设 $E$ 为 $a,b$ 的等化子。则 $E$ 是代数空间，并且
$E\to X$ 局部有限型且为单态射，见空间的态射之引理
\ref{spaces-morphisms-lemma-properties-diagonal}。
我们的假设蕴含：可把引理 \ref{restricted-lemma-factor} 应用于两个态射
$f=\text{id}:X\to X$ 与 $g:E\to X$ 以及 $|X|$ 的闭子集 $T$。
\end{proof}

\begin{lemma}
\label{restricted-lemma-faithful}
设 $S$ 为概形。设 $X,Y$ 为 $S$ 上局部诺特代数空间。
设 $T\subset|X|$ 和 $T'\subset|Y|$ 为闭子集。
设 $a,b:X\to Y$ 为 $S$ 上代数空间的态射，满足
$a|_{X\setminus T}=b|_{X\setminus T}$、
$|a|(T)\subset T'$、$|b|(T)\subset T'$，并且作为态射
$X_{/T}\to Y_{/T'}$ 有 $a_{/T}=b_{/T}$。则 $a=b$。
\end{lemma}

\begin{proof}
这是更一般的引理 \ref{restricted-lemma-faithful-general} 的推论。
\end{proof}

\begin{lemma}
\label{restricted-lemma-equivalence-relation}
设 $S$ 为概形。设 $X$ 为 $S$ 上局部诺特代数空间，
并设 $T\subset|X|$ 为闭子集。
设 $s,t:R\to U$ 为 $X$ 上代数空间的两个态射。假设：
\begin{enumerate}
\item $R,U$ 在 $X$ 上局部有限型；
\item $s,t$ 到 $X\setminus T$ 的基变换是平展等价关系；
\item 形式完备化
$(t_{/T},s_{/T}):R_{/T}\to
U_{/T}\times_{X_{/T}}U_{/T}$ 也是等价关系（记号见证明）。
\end{enumerate}
则 $(t,s):R\to U\times_XU$ 是平展等价关系。
\end{lemma}

\begin{proof}
陈述中使用下标 ${}_{/T}$ 的记号指如下构造：
它把代数空间态射 $f:X'\to X$ 关联到形式代数空间的态射
$f_{/T}:X'_{/f^{-1}T}\to X_{/T}$，见第
\ref{restricted-section-completion-functor} 节。
由假设，态射 $s,t:R\to U$ 在 $X\setminus T$ 上平展。
由于映射 $s,t:R\to U$ 的形式完备化平展，例如由更多态射之引理
\ref{more-morphisms-lemma-check-smoothness-on-infinitesimal-nbhds}，
可知 $s,t$ 平展。把引理 \ref{restricted-lemma-factor} 应用于态射
$\text{id}:R\times_{U\times_XU}R\to R\times_{U\times_XU}R$
和 $\Delta:R\to R\times_{U\times_XU}R$，可知 $(t,s)$ 是单态射。
再次把它应用于
$(t\circ\text{pr}_0,s\circ\text{pr}_1):
R\times_{s,U,t}R\to U\times_XU$ 与
$(t,s):R\to U\times_XU$，得到“传递性”。
等价关系的另外两条公理的证明从略。
\end{proof}

\begin{lemma}
\label{restricted-lemma-smash-away-from-T}
设 $S$ 为概形。设 $X$ 为 $S$ 上局部诺特代数空间，
并设 $T\subset|X|$ 为闭子集。
设 $f:X'\to X$ 为局部有限型并在 $T$ 外平展的代数空间态射。
存在 $f$ 的分解
$$
X' \longrightarrow X'' \longrightarrow X
$$
并具有下列性质：$X''\to X$ 局部有限型，
$X''\to X$ 在 $X\setminus T$ 上为同构，并且
$X'_{/T}\to X''_{/T}$ 是同构（记号见证明）。
\end{lemma}

\begin{proof}
陈述中使用下标 ${}_{/T}$ 的记号指如下构造：
它把代数空间态射 $f:X'\to X$ 关联到形式代数空间的态射
$f_{/T}:X'_{/f^{-1}T}\to X_{/T}$，见第
\ref{restricted-section-completion-functor} 节。
还将使用第 \ref{restricted-section-completion-functor} 节中引入的记号
$U\subset X$ 和 $U'\subset X'$，表示满足
$|U|=|X|\setminus T$ 与
$U'=|X'|\setminus f^{-1}T$ 的开子空间。

\medskip\noindent
用 $X'\amalg U$ 代替 $X'$ 后，可以并且确实假设
$X'\to X$ 的像包含 $U$。令
$$
R = X' \amalg_{U'} (U' \times_U U')
$$
为 $U'\to X'$ 与对角态射
$U'\to U'\times_UU'=U'\times_XU'$ 的推出。
由于 $U'\to X$ 平展，该对角是开浸入，因而 $R$ 是代数空间
（例如由空间之引理
\ref{spaces-lemma-glueing-algebraic-spaces} 得出）。
两个投影 $U'\times_UU'\to U'$ 延拓到 $R$，
得到两个平展态射 $s,t:R\to X'$。
分别在每一片上核验，可知 $R$ 是 $X'$ 上的平展等价关系。
置 $X''=X'/R$；由 Bootstrap 之定理
\ref{bootstrap-theorem-final-bootstrap}，它是代数空间。
按构造得到引理中的分解，并且态射 $X''\to X$ 局部有限型
（因为这可以平展局部地核验，即在 $X'$ 上核验）。
由于 $U'\to U$ 是满平展态射，且
$s^{-1}(U')=t^{-1}(U')=U'\times_UU'$，
可知 $U''=U\times_XX''\to U$ 是同构。
最后，须证明态射 $X'\to X''$ 诱导同构
$X'_{/T}\to X''_{/T}$。为此，注意由我们对 $R$ 的选择，
$R$ 沿 $T$ 的逆像的形式完备化等于 $X'$ 沿 $T$ 的逆像的形式完备化！
由形式空间之第 \ref{formal-spaces-section-completion} 节中
对形式完备化的构造，作为层有
$X''_{/T}=(X'_{/T})/(R_{/T})$。
由于 $X'_{/T}=R_{/T}$，可知 $X'_{/T}=X''_{/T}$。证明完毕。
\end{proof}

\section{态射的 rig-粘合}
\label{restricted-section-glue-formal-to-actual}

\noindent
设 $X,W$ 为代数空间，其中 $X$ 诺特。
设 $Z\subset X$ 为闭子空间，其开补为 $U$。
粗略地说，下面的命题断言
$$
\{\text{态射 }X \to W\} =
\{\text{相容态射 }U \to W\text{ 且 }X_{/Z} \to W\}
$$
这里，$a:U\to W$ 与 $b:X_{/Z}\to W$ 的相容性意指
$a,b$ 定义同一个“rig-空间的态射”。
引入“rig-空间”的范畴需要大量工作，但为了准确陈述该条件在此情形下的
含义，我们无须这样做。

\begin{proposition}
\label{restricted-proposition-glue-modification}
设 $S$ 为概形。设 $X$ 为 $S$ 上局部诺特代数空间。
设 $T\subset|X|$ 为闭子集，其开补子空间为 $U\subset X$。
设 $f:X'\to X$ 为代数空间的真态射，且
$f^{-1}(U)\to U$ 是同构。对 $S$ 上任意代数空间 $W$，映射
$$
\Mor_S(X, W) \longrightarrow
\Mor_S(X', W) \times_{\Mor_S(X'_{/T}, W)} \Mor_S(X_{/T}, W)
$$
是双射。
\end{proposition}

\begin{proof}
设 $w':X'\to W$ 和 $\hat w:X_{/T}\to W$ 为态射，
它们分别与 $X'_{/T}\to X$ 和 $X'_{/T}\to X_{/T}$ 复合后，
给出同一个态射 $X'_{/T}\to W$。
我们须证明存在唯一态射 $w:X\to W$，它分别与
$X'\to X$ 和 $X_{/T}\to X$ 复合后恢复 $w'$ 和 $\hat w$。
唯一性立即由引理 \ref{restricted-lemma-faithful-general} 得出。

\medskip\noindent
关于 $T,f$ 的假设在沿代数空间的任意平展态射
$X_1\to X$ 作基变换时保持。
由于形式代数空间是平展拓扑的层，且已经有唯一性，
只需在用平展覆盖的成员代替 $X$ 后证明存在性。
所以可以假设 $X$ 是仿射诺特概形。

\medskip\noindent
假设 $X$ 是仿射诺特概形。我们将利用空间的推出之第
\ref{spaces-pushouts-section-coequalizer-glue} 节中的材料构造态射
$w:X\to W$。继续阅读本证明之前，先阅读该节的少量材料是有益的。

\medskip\noindent
置 $X''=X'\times_XX'$，并考虑两个态射
$a=w'\circ\text{pr}_1:X''\to W$ 和
$b=w'\circ\text{pr}_2:X''\to W$。
可见 $a,b$ 在开集 $U$ 上相等，并且
$a_{/T}$ 与 $b_{a/T}$ 相等
（因为二者都等于复合 $X''_{/T}\to X_{/T}\to W$，
其中第二个箭头是 $\hat w$）。
所以由引理 \ref{restricted-lemma-faithful-general}，有 $a=b$。

\medskip\noindent
以 $Z\subset X$ 表示 $T$ 上的约化诱导闭子概形结构。
对 $n\geq1$，以 $Z_n\subset X$ 表示 $Z$ 的第 $n$ 个无穷小邻域。
记 $w_n=\hat w|_{Z_n}:Z_n\to W$，于是
$X_{/T}=\colim Z_n$ 上有 $\hat w=\colim w_n$。
置 $Y_n=X'\amalg Z_n$。考虑两个投影
$$
s_n, t_n : R_n = Y_n \times_X Y_n \longrightarrow Y_n
$$
令 $Y_n\to X_n\to X$ 为空间的推出之第
\ref{spaces-pushouts-section-coequalizer-glue} 节中所述的
$s_n,t_n$ 的余等化子（特别地，该余等化子存在并具有良好性质等，
见空间的推出之引理
\ref{spaces-pushouts-lemma-coequalizer}）。
由上一段的结果 $a=b$ 以及 $w',\hat w$ 在 $X'_{/T}$ 上相等，
可知态射
$$
w' \amalg w_n : Y_n \longrightarrow W
$$
等化 $s_n,t_n$。因此，对所有 $n\geq1$，存在态射
$w^n:X_n\to W$，它与 $w',w_n$ 相容。
此外，对 $m\geq1$，复合
$$
X_n \to X_{n + m} \xrightarrow{w^{n + m}} W
$$
按构造等于 $w^n$（因为相应陈述对
$w'\amalg w_{n+m}$ 和 $w'\amalg w_n$ 成立）。
由空间的推出之引理
\ref{spaces-pushouts-lemma-essentially-constant} 和备注
\ref{spaces-pushouts-remark-essentially-constant}，
代数空间系统 $X_n$ 本质常值，其值为 $X$，结论得证。
\end{proof}

\section{Rig-\,'etale 态射的代数化}
\label{restricted-section-algebraization}

\noindent
本节证明 Artin 论文 \cite{ArtinII} 中关于膨胀之结果的一个推广。

\medskip\noindent
本节记号与第 \ref{restricted-section-completion-functor} 节一致，
但其中的代数空间和形式代数空间都局部诺特。

\medskip\noindent
因此，先固定基概形 $S$。所有环、拓扑环、概形、代数空间、
形式代数空间及其间态射都在 $S$ 上。
其次固定局部诺特代数空间 $X$ 和闭子集 $T\subset|X|$。
以 $U\subset X$ 表示满足 $|U|=|X|\setminus T$ 的开子空间。图示为
$$
U \to X \quad |X| = |U| \amalg T
$$
给定代数空间态射 $f:X'\to X$，将沿用第
\ref{restricted-section-completion-functor} 节中的记号
$U'=f^{-1}U$、$T'=|f|^{-1}(T)$ 和
$f_{/T}:X'_{/T'}\to X_{/T}$。
有时把 $X'_{/T'}$ 写成 $X'_{/T}$；更一般地，对 $X$ 上代数空间的态射
$a:X'\to X''$，以 $a_{/T}:X'_{/T}\to X''_{/T}$ 表示沿
$T$ 在 $X',X''$ 中的逆像完备化态射 $a$ 所得的形式代数空间态射。

\medskip\noindent
在此设定下，考虑函子
\begin{equation}
\label{restricted-equation-completion-functor}
\left\{
\begin{matrix}
\text{代数空间的态射}\\
f : X' \to X\text{，其局部}\\
\text{有限型，并且}\\
U' \to U\text{ 是同构}
\end{matrix}
\right\}
\longrightarrow
\left\{
\begin{matrix}
\text{态射 }g : W \to X_{/T}\\
\text{属于形式代数空间}\\
\text{其中 }W\text{ 局部诺特}\\
\text{并且 }g\text{ 是 rig-\,'etale 的}
\end{matrix}
\right\}
\end{equation}
它把 $f:X'\to X$ 映到 $f_{/T}:X'_{/T'}\to X_{/T}$。
这是有意义的，因为由引理 \ref{restricted-lemma-etale-gives-rig-etale}，
$f_{/T}$ 是 rig-\,'etale 的。

\begin{lemma}
\label{restricted-lemma-functoriality-completion-functor}
在上述情形中，设 $X_1\to X$ 为代数空间的态射，且 $X_1$ 局部诺特。
以 $T_1\subset|X_1|$ 表示 $T$ 的逆像，
以 $U_1\subset X_1$ 表示 $U$ 的逆像。记：
\begin{enumerate}
\item $\mathcal{C}_{X,T}$ 为如下范畴：其对象是局部有限型的
代数空间态射 $f:X'\to X$，并且
$U'=f^{-1}U\to U$ 是同构；
\item $\mathcal{C}_{X_1,T_1}$ 为如下范畴：其对象是局部有限型的
代数空间态射 $f_1:X_1'\to X_1$，并且
$f_1^{-1}U_1\to U_1$ 是同构；
\item $\mathcal{C}_{X_{/T}}$ 为如下范畴：其对象是形式代数空间的态射
$g:W\to X_{/T}$，其中 $W$ 局部诺特且 $g$ 是 rig-\,'etale 的；
\item $\mathcal{C}_{X_{1,/T_1}}$ 为如下范畴：其对象是形式代数空间的态射
$g_1:W_1\to X_{1,/T_1}$，其中 $W_1$ 局部诺特且
$g_1$ 是 rig-\,'etale 的。
\end{enumerate}
则图表
$$
\xymatrix{
\mathcal{C}_{X, T} \ar[d] \ar[r] &
\mathcal{C}_{X_{/T}} \ar[d] \\
\mathcal{C}_{X_1, T_1} \ar[r] &
\mathcal{C}_{X_{1, /T_1}}
}
$$
交换，其中水平箭头由 (\ref{restricted-equation-completion-functor}) 给出，
竖直箭头分别由沿 $X_1\to X$ 和沿
$X_{1,/T_1}\to X_{/T}$ 的基变换给出。
\end{lemma}

\begin{proof}
这是因为 $X$ 上代数空间范畴中的完备化函子
$(h:Y\to X)\mapsto Y_{/T}=Y_{/|h|^{-1}T}$ 与纤维积交换。
\end{proof}

\begin{lemma}
\label{restricted-lemma-completion-functor-fully-faithful}
在上述情形中，设 $f:X'\to X$ 为局部有限型且在 $U$ 上为同构的
代数空间态射。设 $g:Y\to X$ 为态射，其中 $Y$ 局部诺特。
则完备化给出双射
$$
\Mor_X(Y, X') \longrightarrow \Mor_{X_{/T}}(Y_{/T}, X'_{/T})
$$
特别地，函子 (\ref{restricted-equation-completion-functor}) 充分忠实。
\end{lemma}

\begin{proof}
设 $a,b:Y\to X'$ 为 $X$ 上态射，且 $a_{/T}=b_{/T}$。
则 $a,b$ 在开子空间 $g^{-1}U$ 上以及沿 $g^{-1}T$ 完备化后相等。
所以由引理 \ref{restricted-lemma-faithful}，$a=b$。
换言之，完备化映射总是单射。

\medskip\noindent
设 $\alpha:Y_{/T}\to X'_{/T}$ 为 $X_{/T}$ 上形式代数空间的态射。
我们须证明存在 $X$ 上态射 $a:Y\to X'$，使
$\alpha=a_{/T}$。证明先作一个标准但繁琐的归约到仿射情形，
然后应用引理 \ref{restricted-lemma-algebraize-morphism-rig-etale}。

\medskip\noindent
设 $\{h_i:Y_i\to Y\}$ 为代数空间的平展覆盖。
若对每个 $i$ 都能找到 $X$ 上态射 $a_i:Y_i\to X'$，
其完备化
$(a_i)_{/T}:(Y_i)_{/T}\to X'_{/T}$ 等于
$\alpha\circ(h_i)_{/T}$，则得到态射 $a:Y\to X'$，且
$\alpha=a_{/T}$。事实上，先注意，由于它们与 $\alpha$ 相等，
作为态射 $(Y_i\times_YY_j)_{/T}\to X'_{/T}$ 有
$(a_i)_{/T}\circ\text{pr}_1=
(a_j)_{/T}\circ\text{pr}_2$
（这里使用完备化 ${}_{/T}$ 与纤维积交换）。
由已经证明的单射性，这说明作为态射
$Y_i\times_YY_j\to X'$ 有
$a_i\circ\text{pr}_1=a_j\circ\text{pr}_2$。
由于 $X'$ 是 fppf 层，态射族 $a_i$ 下降为态射 $a:Y\to X'$。
又因为
$\{(a_i)_{/T}:(Y_i)_{/T}\to X'_{/T}\}$ 是平展覆盖，
有 $\alpha=a_{/T}$。

\medskip\noindent
由上一段的结果，为证明存在性，可以假设 $Y$ 仿射，
并且 $g:Y\to X$ 分解为 $g_1:Y\to X_1$ 和平展态射
$X_1\to X$，其中 $X_1$ 仿射。
令 $T_1\subset|X_1|$ 为 $T$ 的逆像，置
$X'_1=X'\times_XX_1$，投影为 $f_1:X'_1\to X_1$，并令
$$
\alpha_1 = (\alpha, (g_1)_{/T_1}) :
Y_{/T_1} = Y_{/T}
\longrightarrow
X'_{/T} \times_{X_{/T}} (X_1)_{/T_1} = (X'_1)_{/T_1}
$$
由此只需证明 $X_1$ 上 $\alpha_1$ 的存在性；换言之，
可用 $X_1,T_1,X'_1,Y,g_1,\alpha_1$ 分别代替
$X,T,X',Y,f,g,\alpha$。这把问题归结为下一段所述的情形。

\medskip\noindent
假设 $Y,X$ 仿射。回顾，$(Y_{/T})_{red}$ 是仿射概形
（同构于 $g^{-1}T\subset Y$ 上的约化诱导概形结构，见形式空间之引理
\ref{formal-spaces-lemma-reduction-completion}）。
因此
$\alpha_{red}:(Y_{/T})_{red}\to(X'_{/T})_{red}$
在 $f^{-1}T$ 中的像 $E$ 拟紧
（由同一引理，这是 $(X'_{/T})_{red}$ 的底拓扑空间）。
所以可以找到仿射概形 $V$ 和平展态射 $h:V\to X'$，
使 $h$ 的像包含 $E$。选取实线笛卡尔图表
$$
\xymatrix{
Y'_{/T} \ar@{..>}[rd] \ar@{..>}[r] &
W \ar[d] \ar[r] & V_{/T} \ar[d]^{h_{/T}} \\
& Y_{/T} \ar[r]^\alpha & X'_{/T}
}
$$
按构造，态射 $W\to Y_{/T}$ 由代数空间可表、平展且满射
（满射性可通过观察约化看出，见形式空间之引理
\ref{formal-spaces-lemma-reduction-surjective}）。
由引理 \ref{restricted-lemma-algebraize-rig-etale-affine}，
可写 $W=Y'_{/T}$，其中 $Y'\to Y$ 平展且 $Y'$ 仿射。
这给出图表中的虚线箭头。
由于 $W\to Y_{/T}$ 满射，$Y'\to Y$ 的像包含 $g^{-1}T$。
所以 $\{Y'\to Y,Y\setminus g^{-1}T\to Y\}$ 是平展覆盖。
由于 $f$ 在 $U$ 上为同构，存在（唯一的）位于 $X$ 上的态射
$Y\setminus g^{-1}T\to X'$，它在完备化上与 $\alpha$ 相等
（因为 $Y\setminus g^{-1}T$ 的完备化为空）。
所以只需证明 $Y'$ 的存在性；这把问题归结为下一段研究的情形。

\medskip\noindent
由上一段的结果，可以假设 $Y$ 仿射，并且 $\alpha$ 分解为
$Y_{/T}\to V_{/T}\to X'_{/T}$，其中 $V$ 是在 $X'$ 上平展的仿射概形。
仍可用仿射平展覆盖的成员代替 $Y$。
由引理 \ref{restricted-lemma-algebraize-morphism-rig-etale}，
可找到仿射概形的平展态射 $b:Y'\to Y$，
它诱导同构 $b_{/T}:Y'_{/T}\to Y_{/T}$；
还可找到态射 $c:Y'\to V$，使
$c_{/T}\circ b_{/T}^{-1}$ 是给定态射 $Y_{/T}\to V_{/T}$。
令 $a':Y'\to X'$ 为 $c$ 与 $V\to X'$ 的复合，则
$a'_{/T}=\alpha\circ b_{/T}$；换言之，
$Y'$ 和 $\alpha\circ b_{/T}$ 的存在性成立。
最后再次考虑平展覆盖
$\{Y'\to Y,Y\setminus g^{-1}T\to Y\}$，证明完成。
\end{proof}

\begin{lemma}
\label{restricted-lemma-dilatations-affine}
在上述情形中，假设 $X$ 仿射。则函子
(\ref{restricted-equation-completion-functor}) 是等价。
\end{lemma}

\noindent
证明本引理之前，先讨论一个例子。假设
$S=\Spec(k)$、$X=\mathbf{A}^1_k$、$T=\{0\}$。
则 $X_{/T}=\text{Spf}(k[[x]])$。
令 $W=\text{Spf}(k[[x]]\times k[[x]])$。
相应的 $f:X'\to X$ 是原点加倍的仿射直线到仿射直线的映射
（概形，例 \ref{schemes-example-affine-space-zero-doubled}）。
此外，这是从 $X\amalg X$ 到 $X$ 开始应用引理
\ref{restricted-lemma-smash-away-from-T} 中构造所得的输出。

\begin{proof}
已知该函子充分忠实，见引理
\ref{restricted-lemma-completion-functor-fully-faithful}。
下面证明本质满。设 $g:W\to X_{/T}$ 为形式代数空间的态射，
其中 $W$ 局部诺特且 $g$ 是 rig-\,'etale 的。
我们将分若干步证明 $W$ 属于本质像。

\medskip\noindent
第 1 步：$W$ 是仿射形式代数空间。
由引理 \ref{restricted-lemma-algebraize-rig-etale-affine}，
可找到有限型且在 $X\setminus T$ 上平展的 $U\to X$，
使 $U_{/T}$ 同构于 $W$。于是由引理
\ref{restricted-lemma-smash-away-from-T}，$W$ 属于本质像。

\medskip\noindent
第 2 步：$W$ 分离。选取形式空间之定义
\ref{formal-spaces-definition-formal-algebraic-space}
所述的 $\{W_i\to W\}$。
由第 1 步，形式代数空间 $W_i$ 和 $W_i\times_WW_j$ 都属于本质像。
设 $W_i=(X'_i)_{/T}$ 且
$W_i\times_WW_j=(X'_{ij})_{/T}$。
由充分忠实性，得到态射
$t_{ij}:X'_{ij}\to X'_i$ 和
$s_{ij}:X'_{ij}\to X'_j$，分别匹配投影
$W_i\times_WW_j\to W_i$ 和 $W_i\times_WW_j\to W_j$。
考虑结构
$$
R = \coprod X'_{ij},\quad
V = \coprod X'_i,\quad
s = \coprod s_{ij},\quad
t = \coprod t_{ij}
$$
（不能使用字母 $U$，因为它已经用过。）
应用引理 \ref{restricted-lemma-equivalence-relation}，可知
$(t,s):R\to V\times_XV$ 定义了 $X$ 上 $V$ 的平展等价关系。
所以可以取商 $X'=V/R$；由 Bootstrap 之定理
\ref{bootstrap-theorem-final-bootstrap}，它是代数空间。
由于完备化与纤维积和层的商交换，得到 $X'_{/T}\cong W$，
作为 $X_{/T}$ 上形式代数空间的同构。

\medskip\noindent
第 3 步：$W$ 是一般情形。选取形式空间之定义
\ref{formal-spaces-definition-formal-algebraic-space}
所述的 $\{W_i\to W\}$。
形式代数空间 $W_i$ 和 $W_i\times_WW_j$ 都分离。
所以由第 2 步，形式代数空间 $W_i$ 和 $W_i\times_WW_j$
都属于本质像。随后完全按照上一段论证，可知 $W$ 也属于本质像。
证明完毕。
\end{proof}

\begin{theorem}
\label{restricted-theorem-dilatations-general}
设 $S$ 为概形。设 $X$ 为 $S$ 上局部诺特代数空间。
设 $T\subset|X|$ 为闭子集。设 $U\subset X$ 为满足
$|U|=|X|\setminus T$ 的开子空间。完备化函子
(\ref{restricted-equation-completion-functor})
$$
\left\{
\begin{matrix}
\text{代数空间的态射}\\
f : X' \to X\text{，其局部}\\
\text{有限型，并且}\\
f^{-1}U \to U\text{ 是同构}
\end{matrix}
\right\}
\longrightarrow
\left\{
\begin{matrix}
\text{态射 }g : W \to X_{/T}\\
\text{属于形式代数空间}\\
\text{其中 }W\text{ 局部诺特}\\
\text{并且 }g\text{ 是 rig-\,'etale 的}
\end{matrix}
\right\}
$$
把 $f:X'\to X$ 映到
$f_{/T}:X'_{/T'}\to X_{/T}$；该函子是等价。
\end{theorem}

\begin{proof}
由引理 \ref{restricted-lemma-completion-functor-fully-faithful}，该函子充分忠实。
设 $g:W\to X_{/T}$ 为形式代数空间的态射，其中 $W$ 局部诺特，
且 $g$ 是 rig-\,'etale 的。为完成证明，将证明 $W$ 属于本质像。

\medskip\noindent
选取平展覆盖 $\{X_i\to X\}$，其中对所有 $i$，$X_i$ 都仿射。
以 $U_i\subset X_i$ 表示 $U$ 的逆像，以
$T_i\subset X_i$ 表示 $T$ 的逆像。回顾，
$(X_i)_{/T_i}=(X_i)_{/T}=(X_i\times_XX)_{/T}$，且
$W_i=X_i\times_XW=(X_i)_{/T}\times_{X_{/T}}W$，见引理
\ref{restricted-lemma-functoriality-completion-functor}。
注意，得到满足适当余圈条件的同构
$$
\alpha_{ij} :
W_i \times_{X_{/T}} (X_j)_{/T}
\longrightarrow 
(X_i)_{/T} \times_{X_{/T}} W_j
$$
把引理 \ref{restricted-lemma-dilatations-affine} 应用于
$X_i,T_i,U_i,W_i\to(X_i)_{/T}$，得到代数空间的态射
$X'_i\to X_i$，它局部有限型并在 $U_i$ 上为同构；
还得到 $(X_i)_{/T}$ 上的同构
$\beta_i:(X'_i)_{/T}\cong W_i$。
由充分忠实性，得到 $X_i\times_XX_j$ 上的同构
$$
a_{ij} : X'_i \times_X X_j \longrightarrow X_i \times_X X'_j
$$
使
$\alpha_{ij}=\beta_j|_{X_i\times_XX_j}
\circ(a_{ij})_{/T}\circ\beta_i^{-1}|_{X_i\times_XX_j}$。
再次使用充分忠实性（这次在
$X_i\times_XX_j\times_XX_k$ 上），可知态射 $a_{ij}$
满足与 $\alpha_{ij}$ 相同的余圈条件。
换言之，得到相对于族 $\{X_i\to X\}$ 的下降数据
（如空间上的下降之定义
\ref{spaces-descent-definition-descent-datum-for-family-of-morphisms}）
$(X'_i,a_{ij})$。
由 Bootstrap 之引理
\ref{bootstrap-lemma-descend-algebraic-space}，该下降数据有效。
所以得到代数空间的态射 $f:X'\to X$ 以及 $X_i$ 上的同构
$h_i:X'\times_XX_i\to X'_i$，使
$a_{ij}=h_j|_{X_i\times_XX_j}\circ
h_i^{-1}|_{X_i\times_XX_j}$。
读者可以核验，由此得到的 $X_i$ 上同构
$$
(X' \times_X X_i)_{/T}
\xrightarrow{\beta_i \circ (h_i)_{/T}}
W_i
$$
粘合为 $X_{/T}$ 上的同构 $X'_{/T}\to W$；略去若干细节。
\end{proof}

\section{完备化与态射，II}
\label{restricted-section-completion-and-morphisms-bis}

\noindent
为得到 Artin 的膨胀定理，需要在定理
\ref{restricted-theorem-dilatations-general} 给出的对应中，
把形式修改与实际修改相匹配。我们建议读者跳过本节。

\begin{lemma}
\label{restricted-lemma-output-quasi-compact}
在定理 \ref{restricted-theorem-dilatations-general} 的假设与记号下，
设 $f:X'\to X$ 对应于 $g:W\to X_{/T}$。
则 $f$ 拟紧，当且仅当 $g$ 拟紧。
\end{lemma}

\begin{proof}
若 $f$ 拟紧，则由引理
\ref{restricted-lemma-quasi-compact-gives-quasi-compact}，$g$ 拟紧。
反之，假设 $g$ 拟紧。选取平展覆盖 $\{X_i\to X\}$，其中 $X_i$ 仿射。
只需证明基变换 $X'\times_XX_i\to X_i$ 拟紧，见空间的态射之引理
\ref{spaces-morphisms-lemma-quasi-compact-local}。
由形式空间之引理
\ref{formal-spaces-lemma-characterize-quasi-compact-morphism}，
基变换
$W_i\times_{X_{/T}}(X_i)_{/T}\to(X_i)_{/T}$ 拟紧。
由引理 \ref{restricted-lemma-functoriality-completion-functor}，
归结为下一段所述的情形。

\medskip\noindent
假设 $X$ 仿射且 $g:W\to X_{/T}$ 拟紧。
我们须证明 $X'$ 拟紧。设 $V\to X'$ 为满平展态射，
其中 $V=\coprod_{j\in J}V_j$ 是仿射空间的不交并。
则 $V_{/T}\to X'_{/T}=W$ 是满平展态射。
由于 $W$ 拟紧，可以找到有限子集 $J'\subset J$，使
$\coprod_{j\in J'}(V_j)_{/T}\to W$ 满射。
于是
$$
U \amalg \coprod\nolimits_{j \in J'} V_j \longrightarrow X'
$$
满射（因而 $X'$ 拟紧）。事实上，由 $X'_{/T}=W$，有
$|X'|=|U|\amalg|W_{red}|$。
\end{proof}

\begin{lemma}
\label{restricted-lemma-output-quasi-separated}
在定理 \ref{restricted-theorem-dilatations-general} 的假设与记号下，
设 $f:X'\to X$ 对应于 $g:W\to X_{/T}$。
则 $f$ 拟分离，当且仅当 $g$ 拟分离。
\end{lemma}

\begin{proof}
若 $f$ 拟分离，则由引理
\ref{restricted-lemma-quasi-separated-gives-quasi-separated}，$g$ 拟分离。
反之，假设 $g$ 拟分离。我们须证明 $f$ 拟分离。
与引理 \ref{restricted-lemma-output-quasi-compact} 的证明完全相同，
使用空间的态射之引理
\ref{spaces-morphisms-lemma-separated-local} 和形式空间之引理
\ref{formal-spaces-lemma-separated-local}，可以在由仿射概形组成的
$X$ 的平展覆盖的成员上核验。因此可以并且确实假设 $X$ 仿射。

\medskip\noindent
设 $V\to X'$ 为满平展态射，其中
$V=\coprod_{j\in J}V_j$ 是仿射空间的不交并。
为证明 $X'$ 拟分离，只需证明对所有 $j,j'\in J$，
$V_j\times_{X'}V_{j'}$ 拟紧。
由于 $W$ 拟分离，纤维积
$(V_j\times_YV_{j'})_{/T}=
(V_j)_{/T}\times_{X'_{/T}}(V_{j'})_{/T}$
对所有 $j,j'\in J$ 都拟紧。
由于 $X$ 是诺特仿射概形，且 $U'\to U$ 是同构，
$$
(V_j \times_{X'} V_{j'}) \times_X U =
(V_j \times_X V_{j'}) \times_X U
$$
拟紧。因此，由等式
$$
|V_j \times_{X'} V_{j'}| =
|(V_j \times_{X'} V_{j'}) \times_X U| \amalg
|(V_j \times_{X'} V_{j'})_{/T, red}|
$$
以及形式代数空间拟紧当且仅当其相关联的约化代数空间拟紧这一事实，
得到结论。
\end{proof}

\begin{lemma}
\label{restricted-lemma-output-separated}
在定理 \ref{restricted-theorem-dilatations-general} 的假设与记号下，
设 $f:X'\to X$ 对应于 $g:W\to X_{/T}$。
则 $f$ 分离 $\Leftrightarrow$ $g$ 分离且
$\Delta_g:W\to W\times_{X_{/T}}W$ 是 rig-满射的。
\end{lemma}

\begin{proof}
若 $f$ 分离，则由引理
\ref{restricted-lemma-quasi-separated-gives-quasi-separated} 和
\ref{restricted-lemma-separated-mono-open-diagonal-rig-surjective}，
$g$ 分离且 $\Delta_g$ 是 rig-满射的。
假设 $g$ 分离且 $\Delta_g$ 是 rig-满射的。
与引理 \ref{restricted-lemma-output-quasi-compact} 的证明完全相同，
使用空间的态射之引理
\ref{spaces-morphisms-lemma-base-change-separated}
（代数空间态射的分离性在基上局部）、形式空间之引理
\ref{formal-spaces-lemma-base-change-separated}
（形式代数空间态射的分离性在基变换下保持），以及引理
\ref{restricted-lemma-base-change-rig-surjective}
（rig-满射性在基变换下保持），可以在由仿射概形组成的
$X$ 的平展覆盖的成员上核验。
因此可以并且确实假设 $X$ 仿射。此外，由引理
\ref{restricted-lemma-output-quasi-separated}，已经知道
$f:X'\to X$ 拟分离。

\medskip\noindent
由空间的上同调之引理
\ref{spaces-cohomology-lemma-check-separated-dvr} 和备注
\ref{spaces-cohomology-remark-variant}，只需证明：给定任意交换图表
$$
\xymatrix{
\Spec(K) \ar[r] \ar[d] & X' \ar[d] \\
\Spec(R) \ar[r]^p \ar@{-->}[ru] & X' \times_X X'
}
$$
其中 $R$ 是分式域为 $K$ 的完备离散赋值环，
存在使图表交换的虚线箭头
（因为这给出赋值判据的唯一性部分）。
令 $h:\Spec(R)\to X$ 为 $p$ 与态射
$Y\times_XY\to X$ 的复合。分三种情形：
情形 I：$h(\Spec(R))\subset U$。此情形平凡，
因为 $U'=X'\times_XU\to U$ 是同构。
情形 II：$h$ 把 $\Spec(R)$ 映入 $T$。
此情形由假设 $g:W\to X_{/T}$ 分离得出。
事实上，若 $Z$ 表示 $T$ 上的约化诱导闭子空间结构，
则 $h$ 经过 $Z$ 分解，并且
$$
W \times_{X_{/T}} Z = X' \times_X Z \longrightarrow Z
$$
由假设分离（并使用例如形式空间之引理
\ref{formal-spaces-lemma-separated-local}）。
把空间的上同调之引理
\ref{spaces-cohomology-lemma-check-separated-dvr}
应用于所示箭头，得到提升性质。
情形 III：$h(\Spec(K))$ 不在 $T$ 中，但 $h$ 把
$\Spec(R)$ 的闭点映入 $T$。在这种情形下，相应态射
$$
p_{/T} :
\text{Spf}(R)
\longrightarrow
(X' \times_X X')_{/T} =
W \times_{X_{/T}} W
$$
是 adic 态射（由形式空间之引理
\ref{formal-spaces-lemma-map-completions-representable} 和定义
\ref{formal-spaces-definition-adic-morphism}）。
因此，我们关于
$\Delta_g:W\to W\times_{X_{/T}}W$ rig-满射的假设蕴含：
可把 $p_{/T}$ 提升为态射
$\text{Spf}(R)\to W=X'_{/T}$，见引理
\ref{restricted-lemma-monomorphism-rig-surjective}。
使用形式空间之引理
\ref{formal-spaces-lemma-map-into-algebraic-space}
把复合 $\text{Spf}(R)\to X'$ 代数化，得到按所需提升 $p$ 的态射
$\Spec(R)\to X'$。
\end{proof}

\begin{lemma}
\label{restricted-lemma-output-proper}
在定理 \ref{restricted-theorem-dilatations-general} 的假设与记号下，
设 $f:X'\to X$ 对应于 $g:W\to X_{/T}$。
则 $f$ 为真，当且仅当 $g$ 是形式修改
（定义 \ref{restricted-definition-formal-modification}）。
\end{lemma}

\begin{proof}
若 $f$ 为真，则由引理
\ref{restricted-lemma-modification-gives-formal-modification}，$g$ 是形式修改。
假设 $g$ 是形式修改。由引理
\ref{restricted-lemma-output-quasi-compact} 和 \ref{restricted-lemma-output-separated}，
$f$ 拟紧且分离。

\medskip\noindent
由空间的上同调之引理
\ref{spaces-cohomology-lemma-check-proper-dvr} 和备注
\ref{spaces-cohomology-remark-variant}，只需证明：给定任意交换图表
$$
\xymatrix{
\Spec(K) \ar[r] \ar[d] & X' \ar[d]^f \\
\Spec(R) \ar[r]^p \ar@{-->}[ru] & X
}
$$
其中 $R$ 是分式域为 $K$ 的完备离散赋值环，
存在使图表交换的虚线箭头。分三种情形：
情形 I：$p(\Spec(R))\subset U$。此情形平凡，
因为 $U'\to U$ 是同构。
情形 II：$p$ 把 $\Spec(R)$ 映入 $T$。
此情形由假设 $g:W\to X_{/T}$ 为真得出。
事实上，若 $Z$ 表示 $T$ 上的约化诱导闭子空间结构，
则 $p$ 经过 $Z$ 分解，并且
$$
W \times_{X_{/T}} Z = X' \times_X Z \longrightarrow Z
$$
由假设为真。把空间的上同调之引理
\ref{spaces-cohomology-lemma-check-proper-dvr}
应用于所示箭头，得到提升性质。
情形 III：$p(\Spec(K))$ 不在 $T$ 中，但 $p$ 把
$\Spec(R)$ 的闭点映入 $T$。在这种情形下，相应态射
$$
p_{/T} : \text{Spf}(R) \longrightarrow X'_{/T} = W
$$
是 adic 态射（由形式空间之引理
\ref{formal-spaces-lemma-map-completions-representable} 和定义
\ref{formal-spaces-definition-adic-morphism}）。
因此，我们关于 $g:W\to X_{/T}$ rig-满射的假设蕴含：
可把 $g_{/T}$ 提升为态射
$\text{Spf}(R')\to W=X'_{/T}$，其中
$R\subset R'$ 是完备离散赋值环的某个扩张。
使用形式空间之引理
\ref{formal-spaces-lemma-map-into-algebraic-space}
把复合 $\text{Spf}(R')\to X'$ 代数化，得到按所需提升 $p$ 的态射
$\Spec(R')\to X'$。
\end{proof}

\begin{lemma}
\label{restricted-lemma-output-etale}
在定理 \ref{restricted-theorem-dilatations-general} 的假设与记号下，
设 $f:X'\to X$ 对应于 $g:W\to X_{/T}$。
则 $f$ 平展，当且仅当 $g$ 平展。
\end{lemma}

\begin{proof}
若 $f$ 平展，则由引理 \ref{restricted-lemma-map-completions-etale}，$g$ 平展。
反之，假设 $g$ 平展。由于 $f$ 在 $U$ 上为同构，
$f$ 在 $U$ 上平展。所以只需证明 $f$ 在 $X'$ 中每个
位于 $T$ 上方的点处平展。
以 $Z\subset X$ 表示底拓扑空间为
$|Z|=T\subset|X|$ 的约化闭子空间，见空间的性质之定义
\ref{spaces-properties-definition-reduced-induced-space}。
令 $Z_n\subset X$ 为第 $n$ 个无穷小邻域，则
$X_{/T}=\colim Z_n$。由于 $X'_{/T}=W\to X_{/T}$，
由所假设的 $g$ 的平展性可知
$f^{-1}(Z_n)=X'\times_XZ_n\to Z_n$ 平展。
由更多空间的态射之引理
\ref{spaces-more-morphisms-lemma-check-smoothness-on-infinitesimal-nbhds}，
$f$ 在位于 $T$ 上方的点处平展。
\end{proof}

\section{Artin 的膨胀定理}
\label{restricted-section-dilatations}

\noindent
本节对形式代数空间使用不同字体，以突出这些陈述与
\cite{ArtinII} 中相应陈述的相似性。下面是本章第一个主定理。

\begin{theorem}
\label{restricted-theorem-dilatations}
\begin{reference}
\cite[Theorem 3.2]{ArtinII}
\end{reference}
设 $S$ 为概形。设 $X$ 为 $S$ 上局部诺特代数空间。
设 $T\subset|X|$ 为闭子集。令
$\mathfrak X=X_{/T}$
为 $X$ 沿 $T$ 的形式完备化。设
$$
\mathfrak f : \mathfrak X' \to \mathfrak X
$$
为形式修改（定义 \ref{restricted-definition-formal-modification}）。
则存在唯一真态射 $f:X'\to X$，它在 $X$ 中 $T$ 的补上为同构，
并且其完备化 $f_{/T}$ 恢复 $\mathfrak f$。
\end{theorem}

\begin{proof}
这由定理 \ref{restricted-theorem-dilatations-general} 和引理
\ref{restricted-lemma-output-proper} 得出。
\end{proof}

\noindent
下面给出第 \ref{restricted-section-formal-modifications} 节中承诺的
形式修改刻画。

\begin{lemma}
\label{restricted-lemma-formal-modifications-locally-algebraic}
设 $S$ 为概形。设 $\mathfrak X'\to\mathfrak X$ 为 $S$ 上
局部诺特形式代数空间的形式修改
（定义 \ref{restricted-definition-formal-modification}）。给定：
\begin{enumerate}
\item 任意 adic 诺特拓扑环 $A$；
\item 任意 adic 态射
$\text{Spf}(A)\longrightarrow\mathfrak X$。
\end{enumerate}
则存在代数空间的真态射 $X\to\Spec(A)$，以及位于
$\text{Spf}(A)$ 上的同构
$$
\text{Spf}(A) \times_{\mathfrak X} \mathfrak X'
\longrightarrow
X_{/Z}
$$
它把 $\mathfrak X$ 的基变换与 $X$ 沿“闭纤维”
$Z=X\times_{\Spec(A)}\text{Spf}(A)_{red}$ 的形式完备化同构起来；
这里所说的闭纤维是 $X$ 在 $A$ 上的闭纤维。
\end{lemma}

\begin{proof}
由引理 \ref{restricted-lemma-base-change-formal-modification}，态射
$\text{Spf}(A)\times_{\mathfrak X}\mathfrak X'
\to\text{Spf}(A)$ 是形式修改。
所以结论由定理 \ref{restricted-theorem-dilatations} 得出。
\end{proof}

\section{对修改的应用}
\label{restricted-section-modifications}

\noindent
设 $A$ 为诺特环，设 $I\subset A$ 为理想。置
$X=\Spec(A)$ 和 $U=X\setminus V(I)$。本节考虑范畴
\begin{equation}
\label{restricted-equation-modification}
\left\{
f : X' \longrightarrow X
\quad \middle| \quad
\begin{matrix}
X'\text{ 是代数空间}\\
f\text{ 局部有限型}\\
f^{-1}(U) \to U\text{ 是同构}
\end{matrix}
\right\}
\end{equation}
从 $X'/X$ 到 $X''/X$ 的态射，是 $X$ 上代数空间的态射
$X'\to X''$。

\medskip\noindent
设 $A\to B$ 为诺特环同态，设 $J\subset B$ 为满足
$J=\sqrt{IB}$ 的理想。则沿态射
$\Spec(B)\to\Spec(A)$ 的基变换给出从 $A$ 的范畴
(\ref{restricted-equation-modification}) 到 $B$ 的范畴
(\ref{restricted-equation-modification}) 的函子。

\begin{lemma}
\label{restricted-lemma-Noetherian-local-ring}
设 $A\to B$ 为诺特环的环同态，并且对某个理想 $I\subset A$，
它在 $I$-进完备化上诱导同构
（例如，$B$ 是 $A$ 的 $I$-进完备化）。
则基变换定义了 $(A,I)$ 的范畴
(\ref{restricted-equation-modification}) 与 $(B,IB)$ 的范畴
(\ref{restricted-equation-modification}) 之间的范畴等价。
\end{lemma}

\begin{proof}
置 $X=\Spec(A)$、$T=V(I)$。
置 $X_1=\Spec(B)$、$T_1=V(IB)$。
由定理 \ref{restricted-theorem-dilatations-general}
（事实上只需引理 \ref{restricted-lemma-dilatations-affine} 中处理的仿射情形），
$X,T$ 的范畴 (\ref{restricted-equation-modification}) 等价于
局部诺特形式代数空间的 rig-\,'etale 态射
$W\to X_{/T}$ 的范畴。类似地，$X_1,T_1$ 的范畴
(\ref{restricted-equation-modification}) 等价于局部诺特形式代数空间的
rig-\,'etale 态射 $W_1\to X_{1,/T_1}$ 的范畴。
由于
$X_{/T}=\text{Spf}(A^\wedge)$ 且
$X_{1,/T_1}=\text{Spf}(B^\wedge)$
（形式空间，引理
\ref{formal-spaces-lemma-affine-formal-completion-types}），
由我们的假设 $A^\wedge\to B^\wedge$ 是同构，
可知这些范畴等价。略去该等价由基变换给出的核验。
\end{proof}

\begin{lemma}
\label{restricted-lemma-Noetherian-local-ring-properties}
记号与假设如引理 \ref{restricted-lemma-Noetherian-local-ring}。
设 $f:X'\to\Spec(A)$ 通过该等价对应于
$g:Y'\to\Spec(B)$。则 $f$ 拟紧、拟分离、分离、真、有限，
以及在此添加更多内容，当且仅当 $g$ 具有相应性质。
\end{lemma}

\begin{proof}
对于拟紧、拟分离、分离和真这些陈述，可以使用引理
\ref{restricted-lemma-output-quasi-compact}
\ref{restricted-lemma-output-quasi-separated}、
\ref{restricted-lemma-output-separated}、
\ref{restricted-lemma-output-quasi-separated} 和
\ref{restricted-lemma-output-proper}，把相应性质翻译成形式完备化的性质，
再使用引理 \ref{restricted-lemma-Noetherian-local-ring} 证明中的论证来推出。
不过，也有如下使用 fpqc 下降的直接论证。
首先，由于 $A^\wedge\to B^\wedge$ 是同构，
可以归结到对 $A\to A^\wedge$ 和 $B\to B^\wedge$ 证明本引理。
然后注意，
$\{U\to\Spec(A),\Spec(A^\wedge)\to\Spec(A)\}$ 是 fpqc 覆盖，
其中仍取 $U=\Spec(A)\setminus V(I)$。
按照范畴 (\ref{restricted-equation-modification}) 的定义，
$f$ 沿 $U\to\Spec(A)$ 的基变换是 $\text{id}_U$。
设 $P$ 为代数空间态射的一个在基上 fpqc 局部的性质
（空间上的下降，定义
\ref{spaces-descent-definition-property-morphisms-local}），
并且 $P$ 对恒等态射成立。
则 $P$ 对 $f$ 成立，当且仅当 $P$ 对 $g$ 成立。
由空间上的下降之引理
\ref{spaces-descent-lemma-descending-property-quasi-compact}、
\ref{spaces-descent-lemma-descending-property-quasi-separated}、
\ref{spaces-descent-lemma-descending-property-separated}、
\ref{spaces-descent-lemma-descending-property-proper} 和
\ref{spaces-descent-lemma-descending-property-finite}，
这适用于 $P$ 分别取拟紧、拟分离、分离、真和有限。
\end{proof}

\begin{lemma}
\label{restricted-lemma-equivalence-to-completion}
设 $A\to B$ 为局部诺特环的局部映射，满足：
\begin{enumerate}
\item $A\to B$ 平坦；
\item $\mathfrak m_B=\mathfrak m_AB$；
\item $\kappa(\mathfrak m_A)=\kappa(\mathfrak m_B)$。
\end{enumerate}
则从 $(A,\mathfrak m_A)$ 的范畴
(\ref{restricted-equation-modification}) 到 $(B,\mathfrak m_B)$ 的范畴
(\ref{restricted-equation-modification}) 的基变换函子是等价。
\end{lemma}

\begin{proof}
这些条件意指 $A\to B$ 在完备化上诱导同构，见更多代数之引理
\ref{more-algebra-lemma-flat-unramified}。
所以本引理是引理 \ref{restricted-lemma-Noetherian-local-ring} 的特例。
\end{proof}

\begin{lemma}
\label{restricted-lemma-dominate-by-admissible-blowup}
设 $(A,\mathfrak m,\kappa)$ 为诺特局部环。
设 $f:X\to S$ 为 (\ref{restricted-equation-modification}) 的对象，且 $f$ 为真。
则存在支配 $X$ 的 $U$-可容许爆破 $S'\to S$。
\end{lemma}

\begin{proof}
这是更多空间的态射之引理
\ref{spaces-more-morphisms-lemma-dominate-modification-by-blowup} 的特例。
\end{proof}
% Generated mechanically from the frozen spaces-resolve.tex target.
% Do not edit this generated file; edit the bound locale source instead.

% OK, start here.
%

\setcounter{chapter}{88}
\chapter{再论曲面解消}



\phantomsection
\label{spaces-resolve-section-phantom}


\section{引言}
\label{spaces-resolve-section-introduction}

\noindent
本章讨论维数为 $2$ 的诺特代数空间的奇点解消。
我们在前面的一章中已经依照 Lipman \cite{Lipman}
讨论过概形的曲面解消；见《曲面解消》第
\ref{resolve-section-introduction} 节。
本章大多数结果都是概形情形相应结果的直接推论。

\medskip\noindent
除非特别说明，本章中的所有几何对象都是代数空间。因此，若我们说
``设 $f : X \to Y$ 是一个修改''，则意指 $f$ 是《域上的空间》定义
\ref{spaces-over-fields-definition-modification} 中的态射。
固有态射等术语也作同样理解。









\section{修改}
\label{spaces-resolve-section-modifications}

\noindent
设 $(A, \mathfrak m, \kappa)$ 为诺特局部环。令
$S = \Spec(A)$ 且 $U = S \setminus \{\mathfrak m\}$。本节考虑范畴
\begin{equation}
\label{spaces-resolve-equation-modification}
\left\{
f : X \longrightarrow S
\quad \middle| \quad
\begin{matrix}
X\text{ 是代数空间}\\
f\text{ 是固有态射}\\
f^{-1}(U) \to U\text{ 是同构}
\end{matrix}
\right\}.
\end{equation}
从 $X/S$ 到 $X'/S$ 的态射，是与 $S$ 上结构态射相容的代数空间态射
$X \to X'$。我们在《形式空间的代数化》第
\ref{restricted-section-modifications} 节中已经看到，此范畴只依赖于
$A$ 的完备化，并证明了其中对象的一些初等性质。本节专门研究
$\dim(A) \leq 2$，或闭纤维维数至多为 $1$ 的情形。

\begin{lemma}
\label{spaces-resolve-lemma-modification}
设 $(A, \mathfrak m, \kappa)$ 为维数 $2$ 的诺特局部整环，并且
$U = \Spec(A) \setminus \{\mathfrak m\}$ 是正规概形。
则任意修改 $f : X \to \Spec(A)$ 都是
(\ref{spaces-resolve-equation-modification}) 中的态射。
\end{lemma}

\begin{proof}
设 $f : X \to S$ 为修改。我们须证明 $f^{-1}(U) \to U$ 是同构。
由于 $U$ 的每个闭点 $u$ 都具有余维 $1$，结论由《域上的空间》引理
\ref{spaces-over-fields-lemma-modification-normal-iso-over-codimension-1}
得到。
\end{proof}

\begin{lemma}
\label{spaces-resolve-lemma-closed-immersion-on-fibre}
设 $(A, \mathfrak m, \kappa)$ 为诺特局部环。
设 $g : X \to Y$ 为范畴 (\ref{spaces-resolve-equation-modification}) 中的态射。
若特殊纤维上的诱导态射 $X_\kappa \to Y_\kappa$ 是闭浸入，
则 $g$ 是闭浸入。
\end{lemma}

\begin{proof}
这是《空间态射进阶》引理
\ref{spaces-more-morphisms-lemma-where-closed-immersion}
的一个特殊情形。
\end{proof}

\begin{lemma}
\label{spaces-resolve-lemma-dimension-special-fibre}
设 $(A, \mathfrak m, \kappa)$ 为维数 $\geq 1$ 的诺特局部整环。
设 $f : X \to \Spec(A)$ 为代数空间态射。假设下列条件中至少一个成立：
\begin{enumerate}
\item $f$ 是修改（《域上的空间》定义
\ref{spaces-over-fields-definition-modification}）；
\item $f$ 是变换（《域上的空间》定义
\ref{spaces-over-fields-definition-alteration}）；
\item $f$ 局部有限型且拟分离，$X$ 整，并且 $|X|$ 中恰有一个点
映到 $\Spec(A)$ 的泛点；
\item $f$ 局部有限型，$X$ 合宜，并且 $|X|$ 中映到
$\Spec(A)$ 泛点的点，正是 $|X|$ 各不可约分支的泛点；
% P10-E0403
\item 此处还可增加其他条件。
\end{enumerate}
则 $\dim(X_\kappa) \leq \dim(A) - 1$。
\end{lemma}

\begin{proof}
情形 (1)、(2)、(3) 都是 (4) 的特殊情形。取仿射概形
$U = \Spec(B)$ 以及平展态射 $U \to X$。环映射 $A \to B$ 有限型。
我们须证明 $\dim(U_\kappa) \leq \dim(A) - 1$。
由于 $X$ 合宜，$U$ 各不可约分支的泛点就是位于 $|X|$
各不可约分支泛点之上的点；见《合宜空间》引理
\ref{decent-spaces-lemma-decent-generic-points}。
因此，$\Spec(B) \to \Spec(A)$ 在 $(0)$ 上的纤维是 $B$ 的极小素理想
$\mathfrak q_1, \ldots, \mathfrak q_r$ 所成的有限集合。
于是对某个非零 $f \in A$，$A_f \to B_f$ 有限
（《代数》引理 \ref{algebra-lemma-generically-finite}）。
由此 $\kappa(\mathfrak q_i)$ 是 $A$ 的分式域的有限扩张。
设 $\mathfrak q \subset B$ 是位于 $\mathfrak m$ 之上的素理想。则
$$
\dim(B_\mathfrak q) = \max \dim((B/\mathfrak q_i)_{\mathfrak q})
\leq \dim(A)
$$
其中不等式来自 $A \subset B/\mathfrak q_i$ 的维数公式；见《代数》引理
\ref{algebra-lemma-dimension-formula}。
然而，$B_\mathfrak q/\mathfrak m B_\mathfrak q$ 的维数
（这正是 $U_\kappa$ 在相应点处的局部环的维数）至少少一，
因为极小素理想 $\mathfrak q_i$ 不属于 $V(\mathfrak m)$。
最后应用《性质》引理 \ref{properties-lemma-dimension}。
\end{proof}

\begin{lemma}
\label{spaces-resolve-lemma-modification-of-dim-2-is-projective-over-complete}
若 $(A, \mathfrak m, \kappa)$ 为维数 $2$ 的完备诺特局部整环，
则 $\Spec(A)$ 的每个修改都在 $A$ 上射影。
\end{lemma}

\begin{proof}
由《空间态射进阶》引理
\ref{spaces-more-morphisms-lemma-projective-over-complete}，
只须证明 $\Spec(A)$ 的任意修改 $X$ 的特殊纤维维数 $\leq 1$。
这由引理 \ref{spaces-resolve-lemma-dimension-special-fibre} 得出。
\end{proof}






\section{策略}
\label{spaces-resolve-section-strategy}

\noindent
设 $S$ 为概形，$X$ 为 $S$ 上的合宜代数空间。
设 $x_1, \ldots, x_n \in |X|$ 为两两不同的闭点。
对每个 $i$，依照《合宜空间》引理
\ref{decent-spaces-lemma-decent-space-elementary-etale-neighbourhood}，
选取初等平展邻域 $(U_i, u_i) \to (X, x_i)$。
这意味着 $U_i$ 是仿射概形，$U_i \to X$ 平展，$u_i$ 是 $U_i$
中位于 $x_i$ 之上的唯一点，并且
$\Spec(\kappa(u_i)) \to X$ 是表示 $x_i$ 的单态射。
缩小 $U_i$ 后，我们可以并且确实假设：若 $j \not = i$，
则 $U_i$ 中不存在映到 $x_j$ 的点。注意 $u_i \in U_i$ 是闭点。

\medskip\noindent
以 $\mathcal{C}_{X, \{x_1, \ldots, x_n\}}$ 表示如下代数空间态射
$f : Y \to X$ 所成的范畴：它们诱导同构
$f^{-1}(X \setminus \{x_1, \ldots, x_n\}) \to
X \setminus \{x_1, \ldots, x_n\}$。
对每个 $i$，以 $\mathcal{C}_{U_i, u_i}$ 表示如下代数空间态射
$g_i : Y_i \to U_i$ 所成的范畴：它们诱导同构
$g_i^{-1}(U_i \setminus \{u_i\}) \to U_i \setminus \{u_i\}$。
基变换定义了一个函子
\begin{equation}
\label{spaces-resolve-equation-equivalence}
F :
\mathcal{C}_{X, \{x_1, \ldots, x_n\}}
\longrightarrow
\mathcal{C}_{U_1, u_1} \times \ldots \times \mathcal{C}_{U_n, u_n}.
\end{equation}
为了把本章至少一部分问题约化到概形情形，我们有下述引理。

\begin{lemma}
\label{spaces-resolve-lemma-equivalence}
函子 $F$（\ref{spaces-resolve-equation-equivalence}）是范畴等价。
\end{lemma}

\begin{proof}
当 $n = 1$ 时，这就是《空间的极限》引理
\ref{spaces-limits-lemma-excision-modifications}。
当 $n > 1$ 时，可以用完全相同的方法证明，也可以从前一情形推出。
例如，假设 $g_i : Y_i \to U_i$ 是
$\mathcal{C}_{U_i, u_i}$ 中的对象。由 $n = 1$ 的情形，
可以找到 $f'_i : Y'_i \to X$，它在 $X \setminus \{x_i\}$ 上为同构，
而其到 $U_i$ 的基变换为 $f_i$。% P10-E0404
于是可以令
$$
f : Y = Y'_1 \times_X \ldots \times_X Y'_n \to X
$$
这是 $\mathcal{C}_{X, \{x_1, \ldots, x_n\}}$ 中的对象，
沿 $U_i \to X$ 作基变换后恢复 $g_i$。因此该函子本质满。
完全忠实性的证明从略。% P10-E0405
\end{proof}

\begin{lemma}
\label{spaces-resolve-lemma-equivalence-properties}
设 $X, x_i, U_i \to X, u_i$ 如 (\ref{spaces-resolve-equation-equivalence}) 中所述。
若在 $F$ 下 $f : Y \to X$ 对应于 $g_i : Y_i \to U_i$，则
$f$ 拟紧、拟分离、分离、局部有限呈示、有限呈示、局部有限型、有限型、
固有、整或有限，当且仅当对 $i = 1, \ldots, n$，$g_i$ 具有相应性质。
\end{lemma}

\begin{proof}
这由《空间的极限》引理
\ref{spaces-limits-lemma-excision-modifications-properties} 得出。
\end{proof}

\begin{lemma}
\label{spaces-resolve-lemma-equivalence-fibre}
设 $X, x_i, U_i \to X, u_i$ 如 (\ref{spaces-resolve-equation-equivalence}) 中所述。
若在 $F$ 下 $f : Y \to X$ 对应于 $g_i : Y_i \to U_i$，
则代数空间 $Y_{x_i} \cong (Y_i)_{u_i}$。
\end{lemma}

\begin{proof}
这是显然的，因为 $u_i \to x_i$ 是同构。
\end{proof}







\section{由二次变换支配}
\label{spaces-resolve-section-quadratic-spaces}

\noindent
只有当 $X$ 合宜时，我们才定义空间在一点处的爆破。

\begin{definition}
\label{spaces-resolve-definition-blowup-at-point}
设 $S$ 为概形，$X$ 为 $S$ 上的合宜代数空间，
$x \in |X|$ 为闭点。由《合宜空间》引理
\ref{decent-spaces-lemma-decent-space-closed-point}，
可用闭浸入 $i : \Spec(k) \to X$ 表示 $x$。
所谓 $X$ 在 $x$ 处的{\it 爆破 $X' \to X$}，
是指 $X$ 沿闭子空间 $Z = i(\Spec(k)) \subset X$ 的爆破。
\end{definition}

\noindent
在这种一般性下，$X$ 在 $x$ 处的爆破未必固有。
然而，若 $X$ 局部诺特，则由《空间上的除子》引理
\ref{spaces-divisors-lemma-blowing-up-projective}，该爆破是固有的。
回忆，局部诺特代数空间为诺特的，当且仅当它拟紧且拟分离。
此外，对局部诺特代数空间而言，拟分离等价于合宜
（《合宜空间》引理
\ref{decent-spaces-lemma-locally-Noetherian-decent-quasi-separated}）。

\begin{lemma}
\label{spaces-resolve-lemma-equivalence-sequence-blowups}
设 $X, x_i, U_i \to X, u_i$ 如 (\ref{spaces-resolve-equation-equivalence}) 中所述，
并假设在 $F$ 下 $f : Y \to X$ 对应于 $g_i : Y_i \to U_i$。
则 $f$ 存在分解
$$
Y = Z_m \to Z_{m - 1} \to \ldots \to Z_1 \to Z_0 = X
$$
其中 $Z_{j + 1} \to Z_j$ 是 $Z_j$ 在位于
$\{x_1, \ldots, x_n\}$ 之上的闭点 $z_j$ 处的爆破，
当且仅当对每个 $i$，$g_i$ 存在分解
$$
Y_i = Z_{i, m_i} \to Z_{i, m_i - 1} \to \ldots \to
Z_{i, 1} \to Z_{i, 0} = U_i
$$
其中 $Z_{i, j + 1} \to Z_{i, j}$ 是 $Z_{i, j}$
在位于 $u_i$ 之上的闭点 $z_{i, j}$ 处的爆破。
\end{lemma}

\begin{proof}
爆破是可表态射。因此在任一情形中，归纳可见
$Z_j \to X$ 或 $Z_{i, j} \to U_i$ 可表。
于是由《合宜空间》引理
\ref{decent-spaces-lemma-representable-named-properties}，
每个 $Z_j$ 或 $Z_{i,j}$ 都是合宜代数空间。
这说明上述断言有意义（因为爆破只对合宜空间定义）。

为证明等价性，先从一列爆破
$Z_m \to Z_{m - 1} \to \ldots \to Z_1 \to Z_0 = X$ 出发。
第一个态射 $Z_1 \to X$ 是对某个 $x_i$ 的爆破，不妨设为 $x_1$。
把 $F$ 应用于 $Z_1 \to X$，便得到在 $u_1$ 处的爆破
$Z_{1, 1} \to U_1$ 是在 $u_1$ 处的爆破；% P10-E0406
而对 $i > 1$，$Z_{i,0}=U_i$。
下一步，我们或者在 $Z_1$ 上爆破某个 $x_i$（$i\geq 2$），
或者选取 $Z_1 \to X$ 在 $x_1$ 上纤维中的闭点 $z_1$。
第一种情形中做法显然；第二种情形中，利用
$(Z_1)_{x_1} \cong (Z_{1, 1})_{u_1}$
（引理 \ref{spaces-resolve-lemma-equivalence-fibre}），
得到与 $z_1$ 对应的闭点 $z_{1,1}\in Z_{1,1}$。
然后令 $Z_{1,2}\to Z_{1,1}$ 为在 $z_{1,1}$ 处的爆破。
依此继续，便构造出每个 $g_i$ 的分解。

\medskip\noindent
反过来，给定各列爆破
$Z_{i, m_i} \to Z_{i, m_i - 1} \to \ldots \to Z_{i, 1}
\to Z_{i, 0} = U_i$，以完全相同的方式构造 $X$ 的爆破序列。
\end{proof}

\begin{lemma}
\label{spaces-resolve-lemma-make-ideal-principal}
设 $S$ 为概形，$X$ 为 $S$ 上的诺特代数空间。
设 $T \subset |X|$ 为由闭点 $x$ 组成的有限集合，并且
(1) $X$ 在 $x$ 处正则；(2) $X$ 在 $x$ 处的局部环维数为 $2$。
设 $\mathcal{I} \subset \mathcal{O}_X$ 为拟凝聚理想层，
使 $\mathcal{O}_X/\mathcal{I}$ 支撑于 $T$。则存在序列
$$
X_m \to X_{m - 1} \to \ldots \to X_1 \to X_0 = X
$$
其中 $X_{j + 1} \to X_j$ 是 $X_j$ 在某个闭点 $x_j$ 处的爆破，
该点位于 $T$ 中一点之上，并且
$\mathcal{I}\mathcal{O}_{X_n}$ 是可逆理想层。% P10-E0407
\end{lemma}

\begin{proof}
设 $T = \{x_1, \ldots, x_r\}$。如第 \ref{spaces-resolve-section-strategy} 节，
选取初等平展邻域 $(U_i,u_i)\to(X,x_i)$。
对每个 $i$，限制
$\mathcal{I}_i=\mathcal{I}|_{U_i}\subset\mathcal{O}_{U_i}$
是支撑于 $u_i$ 的拟凝聚理想层。
$U_i$ 在 $u_i$ 处的局部环正则且维数为 $2$。
因此可以应用《曲面解消》引理
\ref{resolve-lemma-make-ideal-principal}，得到序列
$$
X_{i, m_i} \to X_{i, m_i - 1} \to \ldots \to X_1
\to X_{i, 0} = U_i
$$
% P10-E0408
其中各态射是在位于 $u_i$ 之上的闭点处作爆破，且
$\mathcal{I}_i \mathcal{O}_{X_{i, m_i}}$ 可逆。
由引理 \ref{spaces-resolve-lemma-equivalence-sequence-blowups}，得到一列爆破
$$
X_m \to X_{m - 1} \to \ldots \to X_1 \to X_0 = X
$$
如引理所述，它到各 $U_i$ 的基变换给出上述序列。
于是 $\mathcal{I}\mathcal{O}_{X_n}$ 是可逆理想层。
事实上，在 $X\setminus\{x_1,\ldots,x_n\}$ 上这是已知的，% P10-E0409
而在每个 $x_i$ 的纤维的一个平展邻域中，这由构造成立。
\end{proof}

\begin{lemma}
\label{spaces-resolve-lemma-dominate-by-blowing-up-in-points}
设 $S$ 为概形，$X$ 为 $S$ 上的诺特代数空间。
设 $T\subset|X|$ 为由闭点 $x$ 组成的有限集合，并且
(1) $X$ 在 $x$ 处正则；(2) $X$ 在 $x$ 处的局部环维数为 $2$。
设 $f:Y\to X$ 为代数空间的固有态射，并且在
$U=X\setminus T$ 上为同构。则存在序列
$$
X_n \to X_{n - 1} \to \ldots \to X_1 \to X_0 = X
$$
其中 $X_{i+1}\to X_i$ 是 $X_i$ 在闭点 $x_i$ 处的爆破，
该点位于 $T$ 中一点之上，并且复合态射存在分解
$X_n\to Y\to X$。
\end{lemma}

\begin{proof}
由《空间态射进阶》引理
\ref{spaces-more-morphisms-lemma-dominate-modification-by-blowup}，
存在支配 $Y\to X$ 的 $U$-容许爆破 $X'\to X$。
因此可以假设存在理想层 $\mathcal{I}\subset\mathcal{O}_X$，
使 $\mathcal{O}_X/\mathcal{I}$ 支撑于 $T$，并且 $Y$
是 $X$ 沿 $\mathcal{I}$ 的爆破。
由引理 \ref{spaces-resolve-lemma-make-ideal-principal}，存在序列
$$
X_n \to X_{n - 1} \to \ldots \to X_1 \to X_0 = X
$$
其中 $X_{i+1}\to X_i$ 是 $X_i$ 在位于 $T$ 中一点之上的闭点
$x_i$ 处的爆破，并且 $\mathcal{I}\mathcal{O}_{X_n}$ 是可逆理想层。
由爆破的泛性质（《空间上的除子》引理
\ref{spaces-divisors-lemma-universal-property-blowing-up}），
得到所需分解。
\end{proof}





\section{由正规化爆破支配}
\label{spaces-resolve-section-normalized-blowups}

\noindent
本节证明，曲面的一个修改可由一列在点处的正规化爆破支配。

\begin{definition}
\label{spaces-resolve-definition-normalized-blowup}
设 $S$ 为概形，$X$ 为 $S$ 上的合宜代数空间，并满足
《空间的态射》引理
\ref{spaces-morphisms-lemma-prepare-normalization} 中的等价条件。
设 $x\in|X|$ 为闭点。所谓 $X$ 在 $x$ 处的
{\it 正规化爆破}，是复合 $X''\to X'\to X$，其中
$X'\to X$ 是 $X$ 在 $x$ 处的爆破
（定义 \ref{spaces-resolve-definition-blowup-at-point}），
而 $X''\to X'$ 是 $X'$ 的正规化。
\end{definition}

\noindent
这里，正规化 $X''\to X'$ 作为代数空间有定义，% P10-E0410
因为 $X'$ 满足《空间的态射》引理
\ref{spaces-morphisms-lemma-prepare-normalization} 中的等价条件；
这由《空间上的除子》引理
\ref{spaces-divisors-lemma-blowup-finite-nr-irreducibles} 得到。
正规化的定义见《空间的态射》定义
\ref{spaces-morphisms-definition-normalization}。

\medskip\noindent
一般而言，即使 $X$ 诺特，正规化爆破也未必固有。
回忆，若一个代数空间有由 Nagata 环的谱所成的仿射平展覆盖，
则称它为 Nagata；见《空间的性质》定义
\ref{spaces-properties-definition-type-property}、
备注 \ref{spaces-properties-remark-list-properties-local-etale-topology}
以及《性质》定义 \ref{properties-definition-nagata}。

\begin{lemma}
\label{spaces-resolve-lemma-Nagata-normalized-blowup}
在定义 \ref{spaces-resolve-definition-normalized-blowup} 中，若 $X$ 为 Nagata，
则 $X$ 在 $x$ 处的正规化爆破是 $X$ 上固有的正规 Nagata 代数空间。
\end{lemma}

\begin{proof}
爆破态射 $X'\to X$ 固有（由于 $X$ 局部诺特，可以应用
《空间上的除子》引理
\ref{spaces-divisors-lemma-blowing-up-projective}）。
因此 $X'$ 为 Nagata（《空间的态射》引理
\ref{spaces-morphisms-lemma-finite-type-nagata}）。
故正规化 $X''\to X'$ 有限（《空间的态射》引理
\ref{spaces-morphisms-lemma-nagata-normalization}），
从而 $X''\to X$ 也固有（《空间的态射》引理
\ref{spaces-morphisms-lemma-finite-proper} 和
\ref{spaces-morphisms-lemma-composition-proper}）。
由此，正规化爆破是正规（《空间的态射》引理
\ref{spaces-morphisms-lemma-normalization-normal}）Nagata 代数空间。
\end{proof}

\noindent
下面是引理 \ref{spaces-resolve-lemma-equivalence-sequence-blowups}
对正规化爆破的类似结论。

\begin{lemma}
\label{spaces-resolve-lemma-equivalence-sequence-normalized-blowups}
设 $X,x_i,U_i\to X,u_i$ 如 (\ref{spaces-resolve-equation-equivalence}) 中所述，
并假设在 $F$ 下 $f:Y\to X$ 对应于 $g_i:Y_i\to U_i$。
假设 $X$ 满足《空间的态射》引理
\ref{spaces-morphisms-lemma-prepare-normalization} 中的等价条件。
则 $f$ 存在分解
$$
Y = Z_m \to Z_{m - 1} \to \ldots \to Z_1 \to Z_0 = X
$$
其中 $Z_{j+1}\to Z_j$ 是 $Z_j$ 在位于
$\{x_1,\ldots,x_n\}$ 之上的闭点 $z_j$ 处的正规化爆破，
当且仅当对每个 $i$，$g_i$ 存在分解
$$
Y_i = Z_{i,m_i} \to Z_{i,m_i-1} \to \ldots \to Z_{i,1}
\to Z_{i,0}=U_i
$$
其中 $Z_{i,j+1}\to Z_{i,j}$ 是 $Z_{i,j}$
在位于 $u_i$ 之上的闭点 $z_{i,j}$ 处的正规化爆破。
\end{lemma}

\begin{proof}
这由证明引理 \ref{spaces-resolve-lemma-equivalence-sequence-blowups}
时完全相同的论证得出。
\end{proof}

\noindent
Nagata 代数空间局部诺特。

\begin{lemma}
\label{spaces-resolve-lemma-dominate-by-normalized-blowing-up}
设 $S$ 为概形，$X$ 为 $S$ 上维数 $\dim(X)=2$ 的诺特
Nagata 代数空间，$f:Y\to X$ 为固有双有理态射。
则存在交换图
$$
\xymatrix{
X_n \ar[r] \ar[d] &
X_{n - 1} \ar[r] &
\ldots \ar[r] &
X_1 \ar[r] &
X_0 \ar[d] \\
Y \ar[rrrr]  & & & & X
}
$$
其中 $X_0\to X$ 是正规化，而
$X_{i+1}\to X_i$ 是 $X_i$ 在某个闭点处的正规化爆破。
\end{lemma}

\begin{proof}
虽然可以直接为代数空间证明本引理，
我们仍沿用上面的做法，将它约化到概形情形。

\medskip\noindent
我们将用到：诺特代数空间拟分离，因此其点有良定义的剩余域
（例如由《合宜空间》引理
\ref{decent-spaces-lemma-decent-space-elementary-etale-neighbourhood}）。
我们还将不再另行说明地使用《空间的态射》第
\ref{spaces-morphisms-section-nagata}、
\ref{spaces-morphisms-section-dimension-formula} 和
\ref{spaces-morphisms-section-normalization} 节的结果。
可用 $Y$ 的正规化替换 $Y$。令 $X_0\to X$ 为正规化。
态射 $Y\to X$ 分解经过 $X_0$，故可以假设 $X$ 与 $Y$ 都正规。

\medskip\noindent
假设 $X$ 与 $Y$ 正规。态射 $f:Y\to X$ 在一个开集上为同构，
该开集包含 $Y$ 中每个余维 $0$ 或 $1$ 的点，以及纤维为有限的每个点；
见《域上的空间》引理
\ref{spaces-over-fields-lemma-modification-normal-iso-over-codimension-1}。
由此存在闭点的有限集合 $T\subset|X|$，使 $f$
在 $X\setminus T$ 上为同构。
由《空间态射进阶》引理
\ref{spaces-more-morphisms-lemma-dominate-modification-by-blowup}，
存在支配 $Y$ 的 $X\setminus T$-容许爆破 $Y'\to X$。
用 $Y'$ 的正规化替换 $Y$ 后，可以假设 $Y\to X$ 可表。

\medskip\noindent
设 $T=\{x_1,\ldots,x_r\}$。如第 \ref{spaces-resolve-section-strategy} 节，
选取初等平展邻域 $(U_i,u_i)\to(X,x_i)$。
对每个 $i$，态射 $Y_i=Y\times_XU_i\to U_i$ 是固有双有理态射，
并且在 $U_i\setminus\{u_i\}$ 上为同构。
因此可以应用《曲面解消》引理
\ref{resolve-lemma-dominate-by-normalized-blowing-up}，得到序列
$$
X_{i,m_i}\to X_{i,m_i-1}\to\ldots\to X_1\to X_{i,0}=U_i
$$
% P10-E0411
其中各态射是在位于 $u_i$ 之上的闭点处作正规化爆破，
且 $X_{i,m_i}$ 支配 $Y_i$。
由引理 \ref{spaces-resolve-lemma-equivalence-sequence-normalized-blowups}，
得到一列正规化爆破
$$
X_m\to X_{m-1}\to\ldots\to X_1\to X_0=X
$$
如引理所述，其到各 $U_i$ 的基变换给出上述序列。
由引理 \ref{spaces-resolve-lemma-equivalence} 的范畴等价，$X_m$ 支配 $Y$。
\end{proof}















\section{到完备化的基变换}
\label{spaces-resolve-section-aux}

\noindent
下面这个简单引理将在后文中成为有用的工具。

\begin{lemma}
\label{spaces-resolve-lemma-iso-completions}
设 $(A,\mathfrak m,\kappa)$ 为局部环，极大理想 $\mathfrak m$
有限生成。设 $X$ 为 $A$ 上的合宜代数空间，并令
$Y=X\times_{\Spec(A)}\Spec(A^\wedge)$，其中 $A^\wedge$
为 $A$ 的 $\mathfrak m$-进完备化。
设 $q\in|Y|$ 的像 $p\in|X|$ 位于 $\Spec(A)$ 的闭点之上，
则 Hensel 局部环的映射
$\mathcal{O}_{X,p}^h\to\mathcal{O}_{Y,q}^h$
在完备化上诱导同构。
\end{lemma}

\begin{proof}
选取《合宜空间》引理
\ref{decent-spaces-lemma-decent-space-elementary-etale-neighbourhood}
中的初等平展邻域 $(U,u)\to(X,p)$，并令
$V=U\times_XY=U\times_{\Spec(A)}\Spec(A^\wedge)$、$v=(u,q)$，
便立刻从概形情形得到结论。概形情形见《曲面解消》引理
\ref{resolve-lemma-iso-completions}。
\end{proof}

\begin{lemma}
\label{spaces-resolve-lemma-port-regularity-to-completion}
设 $(A,\mathfrak m,\kappa)$ 为诺特局部环。
设 $X\to\Spec(A)$ 为局部有限型态射，且 $X$ 是合宜代数空间。
令 $Y=X\times_{\Spec(A)}\Spec(A^\wedge)$。
设 $y\in|Y|$ 的像为 $x\in|X|$。则：
\begin{enumerate}
\item 若 $\mathcal{O}_{Y,y}^h$ 正则，则 $\mathcal{O}_{X,x}^h$ 正则；
\item 若 $y$ 在闭纤维中，则 $\mathcal{O}_{Y,y}^h$ 正则
$\Leftrightarrow\mathcal{O}_{X,x}^h$ 正则；
\item 若 $X$ 在 $A$ 上固有，则 $X$ 正则当且仅当 $Y$ 正则。
\end{enumerate}
\end{lemma}

\begin{proof}
经过平展局部化，前两个断言立即归结为相应的概形结论；见《曲面解消》引理
\ref{resolve-lemma-port-regularity-to-completion}。
对 (3)，由于 $Y\to X$ 满（因为 $A\to A^\wedge$ 忠实平坦），
由 (1) 可知 $Y$ 正则蕴含 $X$ 正则。
反之，若 $X$ 正则，则 $Y$ 在特殊纤维所有点处的 Hensel 局部环都正则。
任取 $y\in|Y|$。在固有情形中，
$|Y|\to|\Spec(A^\wedge)|$ 是闭映射，故可找到特化
$y\leadsto y_0$，其中 $y_0$ 在闭纤维中。
依照《合宜空间》引理
\ref{decent-spaces-lemma-decent-space-elementary-etale-neighbourhood}，
选取初等平展邻域 $(V,v_0)\to(Y,y_0)$。
由于 $Y$ 合宜，可以把 $y\leadsto y_0$ 提升为 $V$ 中的特化
$v\leadsto v_0$（《合宜空间》引理
\ref{decent-spaces-lemma-decent-specialization}）。
于是 $\mathcal{O}_{V,v}$ 是 $\mathcal{O}_{V,v_0}$ 的局部化，故正则，
证明完成。
\end{proof}

\begin{lemma}
\label{spaces-resolve-lemma-formally-unramified}
设 $(A,\mathfrak m)$ 为局部诺特环，$X$ 为 $A$ 上的代数空间。
假设：
\begin{enumerate}
\item $A$ 解析无分歧
（《代数》定义 \ref{algebra-definition-analytically-unramified}）；
\item $X$ 在 $A$ 上局部有限型；
\item $X\to\Spec(A)$ 在 $X$ 中每个余维 $0$ 的点处平展。
\end{enumerate}
则 $X$ 的正规化在 $X$ 上有限。
\end{lemma}

\begin{proof}
选取概形 $U$ 及满平展态射 $U\to X$。
于是 $U\to\Spec(A)$ 满足《曲面解消》引理
\ref{resolve-lemma-formally-unramified} 的假设，因而也满足其结论。
\end{proof}











\section{推得的性质}
\label{spaces-resolve-section-existence-gives}

\noindent
本节证明：对诺特整代数空间，正则变换的存在会带来相当多的后果。
对“病态”诺特代数空间不感兴趣的读者可以跳过本节。

\begin{lemma}
\label{spaces-resolve-lemma-regular-alteration-implies}
设 $S$ 为概形，$Y$ 为 $S$ 上的诺特整代数空间。
假设存在变换 $f:X\to Y$，且 $X$ 正则。
则正规化 $Y^\nu\to Y$ 有限，并且 $Y$ 有一个正则的稠密开子空间。
\end{lemma}

\begin{proof}
通过平展局部化，只须在 $Y=\Spec(A)$ 时证明，其中 $A$ 是诺特整环。
设 $B$ 为 $A$ 在其分式域中的整闭包，并令
$C=\Gamma(X,\mathcal{O}_X)$。由《空间的上同调》引理
\ref{spaces-cohomology-lemma-proper-pushforward-coherent}，
$C$ 是有限 $A$-模。
由于 $X$ 正规（《空间的性质》引理
\ref{spaces-properties-lemma-regular-normal}），
$C$ 是正规整环（《域上的空间》引理
\ref{spaces-over-fields-lemma-normal-integral-sections}）。
因此 $B\subset C$；又因 $A$ 诺特，得出 $B$ 在 $A$ 上有限。

\medskip\noindent
存在非空开集 $V\subset Y$，使 $f^{-1}V\to V$ 有限；
见《域上的空间》定义 \ref{spaces-over-fields-definition-alteration}。
缩小 $V$ 后，可以假设 $f^{-1}V\to V$ 平坦
（《空间的态射》命题
\ref{spaces-morphisms-proposition-generic-flatness-reduced}）。
因此 $f^{-1}V\to V$ 忠实平坦。
由《代数》引理 \ref{algebra-lemma-descent-regular}，$V$ 正则。
\end{proof}

\begin{lemma}
\label{spaces-resolve-lemma-regular-alteration-implies-local}
设 $(A,\mathfrak m,\kappa)$ 为局部诺特整环。
假设存在变换 $f:X\to\Spec(A)$，且 $X$ 正则。则：
\begin{enumerate}
\item 存在非零 $f\in A$，使 $A_f$ 正则；
\item $A$ 在其分式域中的整闭包 $B$ 在 $A$ 上有限；
\item $B$ 的 $\mathfrak m$-进完备化是正规环；换言之，
$B$ 在各极大理想处的完备化都是正规整环；
\item $A$ 的泛形式纤维正则。
\end{enumerate}
\end{lemma}

\begin{proof}
(1) 与 (2) 由引理 \ref{spaces-resolve-lemma-regular-alteration-implies} 得出。
为给 (3) 的证明建立记号，需要重做该引理证明的一部分。
令 $C=\Gamma(X,\mathcal{O}_X)$。由《空间的上同调》引理
\ref{spaces-cohomology-lemma-proper-pushforward-coherent}，
$C$ 是有限 $A$-模。
由于 $X$ 正规（《空间的性质》引理
\ref{spaces-properties-lemma-regular-normal}），
$C$ 是正规整环（《域上的空间》引理
\ref{spaces-over-fields-lemma-normal-integral-sections}）。
因此 $B\subset C$；又因 $A$ 诺特，$B$ 在 $A$ 上有限。
由《曲面解消》引理 \ref{resolve-lemma-algebra-helper}，
为证明 (3)，只须证明 $C$ 的 $\mathfrak m$-进完备化 $C^\wedge$ 正规。

\medskip\noindent
由《代数》引理 \ref{algebra-lemma-completion-finite-extension}，
完备化 $C^\wedge$ 是 $C$ 在位于 $\mathfrak m$ 之上的各素理想处
完备化的乘积。这些素理想为有限多个，且就是 $C$ 的极大理想
$\mathfrak m_1,\ldots,\mathfrak m_r$。
（对 $B$ 的相应结果解释了引理中的最后一个说明。）
因此，用 $C_{\mathfrak m_i}$ 替换 $A$，并用
$X_i=X\times_{\Spec(C)}\Spec(C_{\mathfrak m_i})$ 替换 $X$，
便约化到下一段讨论的情形。
（注意，由《空间的上同调》引理
\ref{spaces-cohomology-lemma-flat-base-change-cohomology}，
$\Gamma(X_i,\mathcal{O})=C_{\mathfrak m_i}$。）

\medskip\noindent
现在 $A$ 是诺特局部正规整环，$f:X\to\Spec(A)$
是正则变换，并且 $\Gamma(X,\mathcal{O}_X)=A$。
我们须证明 $A$ 的完备化 $A^\wedge$ 是正规整环。
由引理 \ref{spaces-resolve-lemma-port-regularity-to-completion}，
$Y=X\times_{\Spec(A)}\Spec(A^\wedge)$ 正则。
由《空间的上同调》引理
\ref{spaces-cohomology-lemma-flat-base-change-cohomology}，
$\Gamma(Y,\mathcal{O}_Y)=A^\wedge$。% P10-E0412
与前面相同，可推出 $A^\wedge$ 正规。
具体而言，由《空间的性质》引理
\ref{spaces-properties-lemma-regular-normal}，$Y$ 正规。
它连通，因为 $\Gamma(Y,\mathcal{O}_Y)=A^\wedge$ 是局部环。
所以 $Y$ 正规且整（对分离代数空间，连通且正规蕴含整）。
于是由《域上的空间》引理
\ref{spaces-over-fields-lemma-normal-integral-sections}，
$\Gamma(Y,\mathcal{O}_Y)=A^\wedge$ 是正规整环。这证明了 (3)。

\medskip\noindent
证明 (4)。令 $\eta\in\Spec(A)$ 表示泛点，并以指标 $\eta$
表示到 $\eta$ 的基变换。由于 $f$ 是变换，概形 $X_\eta$
在 $\eta$ 上有限且忠实平坦。由引理
\ref{spaces-resolve-lemma-port-regularity-to-completion}，
$Y=X\times_{\Spec(A)}\Spec(A^\wedge)$ 正则，
故 $Y_\eta$ 正则（它是 $Y$ 中开集的极限）。
于是
$Y_\eta\to\Spec(A^\wedge\otimes_A\kappa(\eta))$
有限且忠实平坦地映到泛形式纤维。
结论由《代数》引理 \ref{algebra-lemma-descent-regular} 得出。
\end{proof}












\section{解消}
\label{spaces-resolve-section-resolution}

\noindent
先给出一个定义。

\begin{definition}
\label{spaces-resolve-definition-resolution}
设 $S$ 为概形，$Y$ 为 $S$ 上的诺特整代数空间。
所谓 $X$ 的一个{\it 奇点解消}，% P10-E0413
是修改 $f:X\to Y$，其中 $X$ 正则。
\end{definition}

\noindent
对曲面，我们有时还需要多一点信息。

\begin{definition}
\label{spaces-resolve-definition-resolution-surface}
设 $S$ 为概形，$Y$ 为 $S$ 上维数为 $2$ 的诺特整代数空间。
若存在序列
$$
Y_n \to X_{n - 1} \to \ldots \to Y_1 \to Y_0 \to Y
$$
% P10-E0414
满足
\begin{enumerate}
\item 对 $i=0,\ldots,n$，$Y_i$ 在 $Y$ 上固有；
\item $Y_0\to Y$ 是正规化；
\item 对 $i=1,\ldots,n$，$Y_i\to Y_{i-1}$ 是正规化爆破；
\item $Y_n$ 正则，
\end{enumerate}
则称 $Y$ 有一个{\it 由正规化爆破给出的奇点解消}。
\end{definition}

\noindent
注意，条件 (1) 蕴含 $Y$ 的正规化 $Y_0$ 在 $Y$ 上有限，
并且正规化爆破中所用的各正规化也都是有限的。
现在终于来到本章的主定理。

\begin{theorem}
\label{spaces-resolve-theorem-resolve}
设 $S$ 为概形，$Y$ 为 $S$ 上的二维整诺特代数空间。
下列条件等价：
\begin{enumerate}
\item 存在变换 $X\to Y$，其中 $X$ 正则；
\item $Y$ 存在奇点解消；
\item $Y$ 有由正规化爆破给出的奇点解消；
\item 正规化 $Y^\nu\to Y$ 有限，$Y^\nu$ 有有限多个奇点
$y_1,\ldots,y_m\in|Y|$，% P10-E0415
并且对每个 $i$，Hensel 局部环
$\mathcal{O}_{Y^\nu,y_i}^h$ 的完备化正规。
\end{enumerate}
\end{theorem}

\begin{proof}
蕴含关系 (3) $\Rightarrow$ (2) $\Rightarrow$ (1) 是立即的。

\medskip\noindent
设 $X\to Y$ 为变换且 $X$ 正则。
由引理 \ref{spaces-resolve-lemma-regular-alteration-implies}，$Y^\nu\to Y$ 有限。
考虑《空间的态射》引理
\ref{spaces-morphisms-lemma-normalization-normal}
给出的分解 $f:X\to Y^\nu$。
由《域上的空间》引理
\ref{spaces-over-fields-lemma-finite-in-codim-1}，
$f$ 在开集 $V\subset Y^\nu$ 上有限，而该开集包含 $Y^\nu$
中每个余维 $\leq1$ 的点。
由《代数》引理 \ref{algebra-lemma-CM-over-regular-flat}，
以及维数 $\leq 2$ 的正规局部环由 Serre 判别准则成为 Cohen--Macaulay 环这一事实
（《代数》引理 \ref{algebra-lemma-criterion-normal}），
$f$ 在 $V$ 上平坦。
再由《代数》引理 \ref{algebra-lemma-descent-regular}，$V$ 正则。
由于 $Y^\nu$ 诺特，得出
$Y^\nu\setminus V=\{y_1,\ldots,y_m\}$ 有限。
对每个 $i$，令 $\mathcal{O}_{Y^\nu,y_i}^h$ 为 Hensel 局部环。
则
$X\times_Y\Spec(\mathcal{O}_{Y^\nu,y_i}^h)$
是 $\Spec(\mathcal{O}_{Y^\nu,y_i}^h)$ 的正则变换
（省略一些细节）。由引理
\ref{spaces-resolve-lemma-regular-alteration-implies-local}，
$\mathcal{O}_{Y^\nu,y_i}^h$ 的完备化正规。
这就证明了 (1) $\Rightarrow$ (4)。

\medskip\noindent
假设 (4)，须证明 (3)。可以立刻用其正规化替换 $Y$。
设 $y_1,\ldots,y_m\in|Y|$ 为奇点。
如第 \ref{spaces-resolve-section-strategy} 节，选取一族初等平展邻域
$(V_i,v_i)\to(Y,y_i)$。
对每个 $i$，Hensel 局部环 $\mathcal{O}_{Y^\nu,y_i}^h$
是 $\mathcal{O}_{V_i,v_i}$ 的 Hensel 化，故这些环的完备化同构。
于是由概形情形的结果（《曲面解消》定理
\ref{resolve-theorem-resolve}），存在有限列正规化爆破
$$
X_{i,n_i}\to X_{i,n_i-1}\to\ldots\to V_i
$$
其中只爆破位于 $v_i$ 之上的点，并且 $X_{i,n_i}$ 正则。
由引理 \ref{spaces-resolve-lemma-equivalence-sequence-normalized-blowups}，
存在一列正规化爆破
$$
X_n\to X_{n-1}\to\ldots\to X_1\to Y
$$
当然 $X_n$ 也正则（考察局部环即可）。证明完成。
\end{proof}















\section{例子}
\label{spaces-resolve-section-examples}

\noindent
下面给出一些与本章前述结果有关的例子。

\begin{example}
\label{spaces-resolve-example-factorial}
\begin{reference}
\cite[4(c)]{Samuel-UFD}
\end{reference}
设 $k$ 为域。若 $r,s,t$ 两两互素，则环
$A=k[x,y,z]/(x^r+y^s+z^t)$ 是唯一分解整环（UFD）。
事实上，由于 $x^r+y^s+z^t$ 不可约，$A$ 是整环。
元素 $z$ 是素元，即它在 $A$ 中生成素理想。
另一方面，若对某个 $e$ 有 $t=1+ers$，则
$$
A[1/z]\cong k[x',y',1/z]
$$
其中 $x'=x/z^{es}$、$y'=y/z^{er}$，并且
$z=(x')^r+(y')^s$。
因此 $A[1/z]$ 是多项式环的一个局部化，故为 UFD。
由 Nagata 的一个论证，$A$ 也是 UFD；见《代数》引理
\ref{algebra-lemma-invert-prime-elements}。
当 $t$ 模 $rs$ 不同余于 $1$ 时，也可以给出类似论证。
\end{example}

\begin{example}
\label{spaces-resolve-example-completion-not-factorial}
\begin{reference}
一般情形中局部 Picard 群的非零性见
\cite{Brieskorn} 和 \cite{Lipman-rational}。
\end{reference}
若 $1<r<s<t$ 是两两互素且不等于 $2,3,5$ 的整数，则环
$A=\mathbf{C}[[x,y,z]]/(x^r+y^s+z^t)$ 不是 UFD。
例如考虑特殊情形
$A=\mathbf{C}[[x,y,z]]/(x^2+y^5+z^7)$。
考虑映射
$$
\psi_\zeta:\mathbf{C}[[x,y,z]]/(x^2+y^5+z^7)\to\mathbf{C}[[t]]
$$
其定义为
$$
x\mapsto t^7,\quad
y\mapsto t^3,\quad
z\mapsto-\zeta t^2(1+t)^{1/7}
$$
其中 $\zeta$ 是 $7$ 次单位根。
$\psi_\zeta$ 的核 $\mathfrak p_\zeta$ 是高度一素理想；
因而若 $A$ 是 UFD，它就是主理想，不妨设由
$f_\zeta\in\mathbf{C}[[x,y,z]]$ 生成。
注意 $V(x^3-y^7)=\bigcup V(\mathfrak p_\zeta)$，并且
$A/(x^3-y^7)$ 在闭点以外是既约的。
因此，仍假设 $A$ 为 UFD，就会有
$$
\prod\nolimits_\zeta f_\zeta
=u(x^3-y^7)+a(x^2+y^5+z^7)
\quad\text{于}\quad
\mathbf{C}[[x,y,z]]
$$
其中 $u\in\mathbf{C}[[x,y,z]]$ 是单位，
$a\in\mathbf{C}[[x,y,z]]$。
乘以一个常数后，可以假设 $u(0,0,0)=1$。
注意左端的消失阶为 $7$，故 $a=-x\bmod\mathfrak m^2$。
但这样一来，右端出现项 $xy^5$，而左端没有该项，矛盾。
\end{example}

\begin{example}
\label{spaces-resolve-example-not-blow-up}
存在一个优秀的维数为 $2$ 的诺特局部环，以及一个不是概形的修改
$X\to S=\Spec(A)$。我们概述一种构造。
设 $X$ 为 $\mathbf{C}$ 上的正规曲面，只有一个奇点 $x\in X$。
假设存在解消 $\pi:X'\to X$，使例外纤维
$C=\pi^{-1}(x)_{red}$ 是光滑射影曲线。
再假设存在 $c\in C$，使得只要 $\mathcal{O}_C(nc)$
在 $\Pic(X')\to\Pic(C)$ 的像中，就有 $n=0$。
令 $X''\to X'$ 为在非奇异点 $c$ 处的爆破。
设 $C'\subset X''$ 为 $C$ 的严格变换，$E\subset X''$ 为例外纤维。
由 Artin 的结果（\cite{ArtinII}；例如用 \cite{Mumford-topology}
可见 $C'$ 的法丛为负），可以在 $X''$ 中收缩曲线 $C'$，
得到代数空间 $X'''$。图示为
$$
\xymatrix{
& X'' \ar[ld] \ar[rd] \\
X' \ar[rd] &  & X''' \ar[ld] \\
& X
}
$$
我们断言 $X'''$ 不是概形。这给出了所需例子，
因为 $X'''$ 是概形，当且仅当 $X'''$ 到
$A=\mathcal{O}_{X,x}$ 的基变换是概形（省略细节）。
如果 $X'''$ 是概形，% P10-E0416
则 $C'$ 在 $X'''$ 中的像会有一个仿射邻域。
该邻域的补集会是 $X'''$ 上的有效 Cartier 除子
（因为 $X'''$ 除 $1$ 个点外非奇异）。
这个有效 Cartier 除子对应于 $X''$ 上一个与 $E$ 相交且避开 $C'$
的有效 Cartier 除子。取其在 $X'$ 中的像，就得到一个有效
Cartier 除子，它在集合论意义下与 $C$ 交于 $c$。
这不可能，因为按假设，$c$ 的任何倍数都不是某个 Cartier 除子的限制。

\medskip\noindent
最后还须找到这样的奇异曲面 $X$。可以直接取 $X$ 为
$\mathbf{A}^3_\mathbf{C}=\Spec(\mathbf{C}[x,y,z])$ 中由
$$
x^3+y^3+z^3+x^4+y^4+z^4=0
$$
给出的仿射曲面，奇点为 $(0,0,0)$。
此时 $(0,0,0)$ 是唯一奇点。
沿对应于 $(0,0,0)$ 的极大理想爆破 $X$，得到三个图册，
每个都同构于光滑仿射曲面
$$
1+s^3+t^3+x(1+s^4+t^4)=0
$$
它非奇异，例外除子 $C$ 由 $x=0$ 给出。
读者会认出 $C$ 是椭圆曲线。
此外，从 $(0,0,0)$ 作投影可见曲面 $X$ 是有理的；
也可在爆破方程中解出 $x$ 得到同一结论。
最后，非奇异有理曲面的 Picard 群可数，
而复数域上椭圆曲线的 Picard 群不可数。
故可以找到所要求的闭点 $c$。
\end{example}
% Generated mechanically from the frozen formal-defos.tex target.
% Do not edit this generated file; edit the bound locale source instead.


% OK, start here.
%

\setcounter{chapter}{89}
\chapter{形式形变理论}



\phantomsection
\label{formal-defos-section-phantom}




\section{引言}
\label{formal-defos-section-introduction}

\noindent
本章建立一种可在 Stacks 项目后续章节中应用的形式形变理论，
其论述紧随 Rim \cite[Exposee VI]{SGA7-I} 和 Schlessinger \cite{Sch}。
我们强烈建议初次接触这一主题的读者先阅读 Schlessinger 的论文；
该文的一般性已足以应付大多数应用，而且采用这种形式形变理论的论文
事实上大多都会使用 Schlessinger 的结果。

\medskip\noindent
设 $\Lambda$ 为完备诺特局部环，剩余域为 $k$；以
$\mathcal{C}_\Lambda$ 表示剩余域为 $k$ 的阿廷局部
$\Lambda$-代数所成的范畴。给定函子
$F : \mathcal{C}_\Lambda \to \textit{Sets}$，并假设 $F(k)$
是单元素集合。Schlessinger 的论文引入了条件 (H1)--(H4)，使得：
\begin{enumerate}
\item $F$ 有一个“包络”当且仅当 (H1)--(H3) 成立。
\item $F$ 是 pro-可表的当且仅当 (H1)--(H4) 成立。
\end{enumerate}
本章的目的是完全按照 Rim 的论文，从以下两方面推广这些结果：
\begin{enumerate}
\item[(A)] 以 $\mathcal{C}_\Lambda$ 上的群胚余纤维范畴
$\mathcal{F}$ 取代函子 $F$；见第 \ref{formal-defos-section-CLambda} 节。
\item[(B)] 我们允许 $\Lambda$ 为诺特环，并令 $\Lambda \to k$
为到一个域的有限环映射。范畴 $\mathcal{C}_\Lambda$ 的对象是阿廷局部
$\Lambda$-代数 $A$，并配备一个给定的等同 $A/\mathfrak m_A = k$。
\end{enumerate}
条件 $F(k)$ 为单元素集合的对应物，是 $\mathcal{F}(k)$ 为平凡群胚。
若 $\mathcal{F}$ 满足此条件，我们称它为一个{\it 预形变范畴}；
但一般而言我们并不作此假设。Rim 的论文
\cite[Exposee VI]{SGA7-I} 是本文各项结果的原始来源。
我们还要提及很有用的论文 \cite{Talpo-Vistoli}；该文也用群胚讨论形变理论，
但一般性不及本文。

\medskip\noindent
范畴 $\mathcal{C}_\Lambda$ 的“完备化”
$\widehat{\mathcal{C}}_\Lambda$ 起着重要作用。
$\widehat{\mathcal{C}}_\Lambda$ 的对象是诺特完备局部
$\Lambda$-代数 $R$，其剩余域配备了与 $k$ 的等同；见
第 \ref{formal-defos-section-category-completion-CLambda} 节。
一方面，$\mathcal{C}_\Lambda \subset \widehat{\mathcal{C}}_\Lambda$
是严格满子范畴；另一方面，$\widehat{\mathcal{C}}_\Lambda$
是 $\mathcal{C}_\Lambda$ 的 pro-对象范畴的满子范畴。若函子
$\mathcal{C}_\Lambda \to \textit{Sets}$ 同构于可表函子
$\underline{R} = \Mor_{\widehat{\mathcal{C}}_\Lambda}(R, -)$
在 $\mathcal{C}_\Lambda$ 上的限制，其中
$R \in \Ob(\widehat{\mathcal{C}}_\Lambda)$，则称该函子为
{\it pro-可表的}。

\medskip\noindent
{\it 群胚余纤维范畴}与群胚纤维范畴对偶；我们将在
第 \ref{formal-defos-section-preliminary} 节引入它们。
$\mathcal{C}_\Lambda$ 上群胚余纤维范畴之间的一个{\it 光滑}态射，
是指满足对象的无穷小提升判据的态射；见
第 \ref{formal-defos-section-smooth-morphisms} 节。
这与形式光滑环映射的定义类似；见《代数》定义
\ref{algebra-definition-formally-smooth}，并且恰好与
《可表示性判据》第 \ref{criteria-section-formally-smooth} 节中的概念对偶。
这是一个重要概念，因为我们最终要证明某些群胚余纤维范畴
具有光滑的 pro-可表表示，这与《代数叠》第
\ref{algebraic-section-stack-to-presentation} 和
\ref{algebraic-section-smooth-groupoid-gives-algebraic-stack} 节中
对代数叠的刻画十分相似。
$\mathcal{C}_\Lambda$ 上群胚余纤维范畴 $\mathcal{F}$ 的一个{\it 泛形式对象}，是其完备化中的对象
$\xi \in \widehat{\mathcal{F}}(R)$，使相伴态射
$\underline{\xi} : \underline{R}|_{\mathcal{C}_\Lambda} \to \mathcal{F}$
为光滑态射。

\medskip\noindent
在第 \ref{formal-defos-section-schlessinger-conditions} 节中，
我们在 $\mathcal{F}$ 上定义条件 (S1) 和 (S2)，以推广 Schlessinger 的
(H1) 和 (H2)。Schlessinger 的 (H3) 所对应的条件——即
$\mathcal{F}$ 有有限维切空间——不另命名。
发展这套理论的一个关键步骤，是为满足 (S1)、(S2) 和 (H3) 的
预形变范畴证明泛形式对象的存在性；见引理
\ref{formal-defos-lemma-versal-object-existence}。
Schlessinger 对函子 $F : \mathcal{C}_\Lambda \to \textit{Sets}$
所定义的{\it 包络}，在我们的术语中，是一个泛形式对象
$\xi \in \widehat{F}(R)$，且其诱导的切空间映射
$d\underline{\xi} : T\underline{R}|_{\mathcal{C}_\Lambda} \to TF$
为同构。
文献中常把包络称为“极小泛”对象。我们不采用这一叫法，理由如下。
一个函子可能有泛形式对象，却没有包络。此外，我们将在
第 \ref{formal-defos-section-minimal-versal} 节证明：若预形变范畴有泛形式对象，
则它总有一个{\it 极小}泛形式对象（定义见
定义 \ref{formal-defos-definition-minimal-versal}），并且该对象在同构意义下唯一；
见引理 \ref{formal-defos-lemma-minimal-versal}。
然而，极小泛形式对象未必在切空间上诱导同构！（见例
\ref{formal-defos-example-do-not-get-S2} 和
\ref{formal-defos-example-smooth-continued}。）

\medskip\noindent
在牢记上述差异的前提下，
定理 \ref{formal-defos-theorem-miniversal-object-existence}
正是上面 (1) 的直接推广：当 $\mathcal{F}$ 是函子时，它恢复
Schlessinger 的结果；在条件 (S1) 和 (S2) 成立时，它又通过切空间映射
$d\underline{\xi} : T\underline{R}|_{\mathcal{C}_\Lambda} \to TF$
刻画极小泛形式对象。

\medskip\noindent
在第 \ref{formal-defos-section-RS-condition} 节中，
我们在 $\mathcal{F}$ 上定义 Rim 的条件 (RS)，以推广 Schlessinger 的
(H4)。所谓{\it 形变范畴}，是指满足 (RS) 的预形变范畴。
pro-可表函子的对应物，是 $\mathcal{C}_\Lambda$ 上那些
{\it 可由 $\mathcal{C}_\Lambda$ 上函子中的光滑 pro-可表群胚表示}的群胚余纤维范畴；见定义
\ref{formal-defos-definition-groupoid-in-functors}、
\ref{formal-defos-definition-prorepresentable-groupoid-in-functors} 和
\ref{formal-defos-definition-smooth-groupoid-in-functors}。
这种表示概念把 $\mathcal{F}$ 各纤维的群胚结构纳入考虑。
在定理 \ref{formal-defos-theorem-presentation-deformation-groupoid} 中，
我们证明：$\mathcal{F}$ 可由函子中的光滑 pro-可表群胚表示，
当且仅当 $\mathcal{F}$ 的切空间和无穷小自同构空间均为有限维。
这是上面 (2) 的推广；当 $\mathcal{F}$ 是函子时，它归结为
Schlessinger 的结果。最后还有第 \ref{formal-defos-section-minimality} 节，
其中讨论如何利用极小泛形式对象构造函子中的光滑 pro-可表群胚的
极小表示；这种表示在同构意义下唯一。

\medskip\noindent
我们还得到对 Schlessinger 条件的如下概念性解释。
若预形变范畴 $\mathcal{F}$ 满足 (RS)，则相伴的同构类函子
$\overline{\mathcal{F}}: \mathcal{C}_\Lambda \to \textit{Sets}$
满足 (H1) 和 (H2)（见引理
\ref{formal-defos-lemma-RS-implies-S1-S2} 和
\ref{formal-defos-lemma-S1-S2-associated-functor}）。
反过来，若函子 $F : \mathcal{C}_\Lambda \to \textit{Sets}$
自然地产生于一个群胚余纤维范畴 $\mathcal{F}$ 的同构类函子，
那么实践中，用来证明 $F$ 满足 (H1) 和 (H2) 的论证，似乎也总能证明
$\mathcal{F}$ 满足 (RS)。相关例子见《形变问题》第
\ref{examples-defos-section-introduction} 节。
此外，若 $\mathcal{F}$ 满足 (RS)，则 $\overline{\mathcal{F}}$ 的条件
(H4) 可以用 $\mathcal{F}$ 中对象的自同构能否延拓来作简单解释
（引理 \ref{formal-defos-lemma-RS-associated-functor}）。
这些观察表明，应把 (RS) 看作形变理论中基本的黏合条件。




\section{记号与约定}
\label{formal-defos-section-notations-conventions}

\noindent
环均指带 $1$ 的交换环。局部环 $A$ 的极大理想记为
$\mathfrak{m}_A$。正整数集合记为
$\mathbf{N} = \{1, 2, 3, \ldots\}$。若 $U$ 是范畴
$\mathcal{C}$ 的一个对象，则以 $\underline{U}$ 表示函子
$\Mor_\mathcal{C}(U, -): \mathcal{C} \to \textit{Sets}$；见备注
\ref{formal-defos-remarks-cofibered-groupoids} 的
(\ref{formal-defos-item-definition-yoneda})。注意：这可能与其他章节的记号冲突；
在那里我们有时用 $\underline{U}$ 表示
$h_U(-) = \Mor_\mathcal{C}(-, U)$。

\medskip\noindent
本章始终设 $\Lambda$ 为诺特环，并设
$\Lambda \to k$ 是从 $\Lambda$ 到一个域的有限环映射。
该映射的核记为 $\mathfrak m_\Lambda$，像记为 $k' \subset k$。
可以证明 $\mathfrak m_\Lambda$ 是极大理想，
$k' = \Lambda/\mathfrak m_\Lambda$ 是域，且扩张 $k/k'$ 有限。
见公式 (\ref{formal-defos-equation-k-prime}) 附近的讨论。


\section{基范畴}
\label{formal-defos-section-CLambda}

\noindent
动机。形式形变理论的一个重要应用，是研究由代数空间表示的判据。
设给定局部诺特基底 $S$ 和函子
$F : (\Sch/S)_{fppf}^{opp} \to \textit{Sets}$。
设 $k$ 是 $S$ 上的有限型域，也就是说，给定有限型态射
$\Spec(k) \to S$。
Artin 的判据之一要求：对任意元素 $x \in F(\Spec(k))$，
与三元组 $(S, k, x)$ 相伴的预形变函子应为 pro-可表的。
由《态射》引理 \ref{morphisms-lemma-point-finite-type}，
$k$ 在 $S$ 上为有限型这一条件意味着存在仿射开集
$\Spec(\Lambda) \subset S$，使 $k$ 是有限 $\Lambda$-代数。
这说明了为何本章始终采用如下基范畴。

\begin{definition}
\label{formal-defos-definition-CLambda}
设 $\Lambda$ 为诺特环，$\Lambda \to k$ 为有限环映射，其中 $k$ 是域。
定义{\it $\mathcal{C}_\Lambda$} 为如下范畴：
\begin{enumerate}
\item 对象为二元组 $(A, \varphi)$，其中 $A$ 是阿廷局部
$\Lambda$-代数，而 $\varphi : A/\mathfrak m_A \to k$ 是
$\Lambda$-代数同构；
\item 态射 $f : (B, \psi) \to (A, \varphi)$ 为局部
$\Lambda$-代数同态，且满足
$\varphi \circ (f \bmod \mathfrak m) = \psi$。
\end{enumerate}
若 $\Lambda$ 是诺特完备局部环，而 $k$ 是其剩余域，
则称我们处于{\it 经典情形}。
\end{definition}

\noindent
注意，若 $\Lambda \to k$ 满射，且 $A$ 是阿廷局部
$\Lambda$-代数，则等同 $\varphi$ 一旦存在便是唯一的。
而且在这种情形下，任意 $\Lambda$-代数映射 $A \to B$
都与这些等同相容。因此，此时 $\mathcal{C}_\Lambda$
就是剩余域“为”$k$ 的阿廷局部 $\Lambda$-代数所成的范畴。
在一般情形下，我们也滥用记号，直接以 $A$ 表示
$\mathcal{C}_\Lambda$ 的对象。此外，我们常写
$A/\mathfrak m = k$，即假装 $\mathcal{C}_\Lambda$ 中所有环的
剩余域都是 $k$（由于 $\mathcal{C}_\Lambda$ 中所有环映射都与给定等同相容，
这不会造成任何问题）。
本章余下部分固定基环 $\Lambda$ 和域 $k$。
$\mathcal{C}_\Lambda$ 将作为下文所考虑余纤维范畴的基范畴。

\begin{definition}
\label{formal-defos-definition-small-extension}
设 $f: B \to A$ 是 $\mathcal{C}_\Lambda$ 中的环映射。
若 $f$ 为满射，并且 $\Ker(f)$ 是被 $\mathfrak{m}_B$ 零化的
非零主理想，则称 $f$ 为一个{\it 小扩张}。
\end{definition}

\noindent
由下面的引理，我们常可把 $\mathcal{C}_\Lambda$ 中涉及满环映射的论证
归约到小扩张的情形。

\begin{lemma}
\label{formal-defos-lemma-factor-small-extension}
设 $f: B \to A$ 是 $\mathcal{C}_\Lambda$ 中的满环映射。
则 $f$ 可分解为一列小扩张的复合。
\end{lemma}

\begin{proof}
设 $I$ 为 $f$ 的核。由于 $B$ 是阿廷环，极大理想
$\mathfrak{m}_B$ 幂零；设 $\mathfrak{m}_B^n = 0$。于是 $f$ 有分解
$$
B = B/I\mathfrak{m}_B^{n-1} \to B/I\mathfrak{m}_B^{n-2} \to
\ldots \to B/I \cong A
$$
其中每个映射均为满射，且其核被极大理想零化。
因此只需证明 $f$ 本身具有这一性质的情形，也就是说，
$I$ 被 $\mathfrak{m}_B$ 零化的情形。此时 $I$ 是有限维
$k$-向量空间；见《代数》引理
\ref{algebra-lemma-artinian-finite-length}。
取 $I$ 作为 $k$-向量空间的一组基 $x_1, \ldots, x_n$，便得到分解
$$
B \to B/(x_1) \to \ldots \to  B/(x_1, \ldots, x_n) \cong  A
$$
其中每个映射都是小扩张。
\end{proof}

\noindent
下一个引理说明：剩余域为 $k$ 的局部 $\Lambda$-代数上的模，
其长度可由它作为 $\Lambda$-模的长度计算。为解释陈述中的记号，
设 $k' \subset k$ 为我们固定的有限环映射 $\Lambda \to k$ 的像。
注意 $k' \subset k$ 是有限环扩张。因此 $k'$ 是域，而 $k/k'$
是有限域扩张；见《代数》引理
\ref{algebra-lemma-integral-under-field}。
又因 $\Lambda \to k'$ 为满射，其核是极大理想
$\mathfrak m_\Lambda$。因此
\begin{equation}
\label{formal-defos-equation-k-prime}
[k : k'] = [k : \Lambda/\mathfrak m_\Lambda] < \infty
\end{equation}
且在经典情形下 $k = k'$。本章始终固定记号
$k' = \Lambda/\mathfrak m_\Lambda$。

\begin{lemma}
\label{formal-defos-lemma-length}
设 $A$ 是剩余域为 $k$ 的局部 $\Lambda$-代数，$M$ 是 $A$-模。则
$[k : k'] \text{length}_A(M) = \text{length}_\Lambda(M)$。
在经典情形下，
$\text{length}_A(M) = \text{length}_\Lambda(M)$。
\end{lemma}

\begin{proof}
若 $M$ 是单 $A$-模，则 $M \cong k$（作为 $A$-模）；见《代数》引理
\ref{algebra-lemma-characterize-length-1}。
此时 $\text{length}_A(M) = 1$，且
$\text{length}_\Lambda(M) = [k' : k]$；% P10-E0420
见《代数》引理 \ref{algebra-lemma-dimension-is-length}。
若 $\text{length}_A(M)$ 有限，则取 $M$ 的一个 $A$-子模滤过，
使其逐次商均为单模，再由长度的可加性得到结论；见《代数》引理
\ref{algebra-lemma-length-additive}。
若 $\text{length}_A(M)$ 无限，则结论由显然不等式
$\text{length}_A(M) \leq \text{length}_\Lambda(M)$ 得出。
\end{proof}

\begin{lemma}
\label{formal-defos-lemma-surjective}
设 $A \to B$ 是 $\mathcal{C}_\Lambda$ 中的环映射。下列条件等价：
\begin{enumerate}
\item $f$ 为满射；% P10-E0421
\item $\mathfrak m_A/\mathfrak m_A^2 \to \mathfrak m_B/\mathfrak m_B^2$
为满射；
\item $\mathfrak m_A/(\mathfrak m_\Lambda A + \mathfrak m_A^2)
\to \mathfrak m_B/(\mathfrak m_\Lambda B + \mathfrak m_B^2)$ 为满射。
\end{enumerate}
\end{lemma}

\begin{proof}
对 $\mathcal{C}_\Lambda$ 中任意环映射 $f : A \to B$，都有
$f(\mathfrak m_A) \subset \mathfrak m_B$；例如，这是因为
$\mathfrak m_A$ 和 $\mathfrak m_B$ 分别是 $A$ 和 $B$ 中幂零元的集合。
设 $f$ 为满射。取 $y \in \mathfrak m_B$，再取 $x \in A$ 使
$f(x) = y$。由于 $f$ 诱导同构
$A/\mathfrak m_A \to B/\mathfrak m_B$，可知 $x \in \mathfrak m_A$。
故诱导映射
$\mathfrak m_A/\mathfrak m_A^2 \to \mathfrak m_B/\mathfrak m_B^2$
为满射。这证明了 (1) 蕴含 (2)。

\medskip\noindent
显然 (2) 蕴含 (3)。映射 $A \to B$ 给出典范交换图
$$
\xymatrix{
\mathfrak m_\Lambda/\mathfrak m_\Lambda^2 \otimes_{k'} k \ar[r] \ar[d] &
\mathfrak m_A/\mathfrak m_A^2 \ar[r] \ar[d] &
\mathfrak m_A/(\mathfrak m_\Lambda A + \mathfrak m_A^2) \ar[r] \ar[d] & 0 \\
\mathfrak m_\Lambda/\mathfrak m_\Lambda^2 \otimes_{k'} k \ar[r] &
\mathfrak m_B/\mathfrak m_B^2 \ar[r] &
\mathfrak m_B/(\mathfrak m_\Lambda B + \mathfrak m_B^2) \ar[r] & 0
}
$$
其两行正合。因此若 (3) 成立，则 (2) 也成立。

\medskip\noindent
设 (2) 成立。由 Nakayama 引理（《代数》引理
\ref{algebra-lemma-NAK}），要证明 $A \to B$ 为满射，
只需证明 $A/\mathfrak m_A \to B/\mathfrak m_AB$ 为满射。
（注意 $\mathfrak m_A$ 是幂零理想。）由于
$k = A/\mathfrak m_A = B/\mathfrak m_B$，只需证明
$\mathfrak m_AB \to \mathfrak m_B$ 为满射。再次应用 Nakayama 引理，
只需证明
$\mathfrak m_AB/\mathfrak m_A\mathfrak m_B \to
\mathfrak m_B/\mathfrak m_B^2$
为满射，而这正是我们的假设。
\end{proof}

\noindent
若 $A \to B$ 是 $\mathcal{C}_\Lambda$ 中的环映射，则映射
$\mathfrak m_A/(\mathfrak m_\Lambda A + \mathfrak m_A^2)
\to \mathfrak m_B/(\mathfrak m_\Lambda B + \mathfrak m_B^2)$
就是相对余切空间上的映射。正式定义如下。

\begin{definition}
\label{formal-defos-definition-tangent-space-ring}
设 $R \to S$ 是局部环之间的局部同态。所谓 $R$ 在 $S$ 上的
{\it 相对余切空间}% P10-E0422
\footnote{注意：我们稍后将看到，在本章的一般设定下，
$\Lambda$ 上对象 $A \in \mathcal{C}_\Lambda$ 的切空间，
不应简单地定义为相对余切空间的 $k$-线性对偶。
事实上，相对余切空间的正确定义是
$\Omega_{S/R} \otimes_S S/\mathfrak m_S$。}
是如下 $S/\mathfrak m_S$-向量空间：
$\mathfrak m_S/(\mathfrak m_R S + \mathfrak m_S^2)$。
\end{definition}

\noindent
若 $f_1: A_1 \to A$ 和 $f_2: A_2 \to A$ 是两个环映射，
则纤维积 $A_1 \times_A A_2$ 是 $A_1 \times A_2$ 的子环，
由两个投影到 $A$ 的像相等的元素组成。本章将反复考虑
$f_1$ 和 $f_2$ 属于 $\mathcal{C}_\Lambda$ 时涉及这种纤维积的条件。
但该纤维积未必是 $\mathcal{C}_\Lambda$ 的对象。

\begin{example}
\label{formal-defos-example-fibre-product}
设 $p$ 为素数，$n \in \mathbf{N}$。令
$\Lambda = \mathbf{F}_p(t_1, t_2, \ldots, t_n)$，并令
$k = \mathbf{F}_p(x_1, \ldots, x_n)$；映射 $\Lambda \to k$ 由
$t_i \mapsto x_i^p$ 给出。令 $A = k[\epsilon] = k[x]/(x^2)$。
则 $A$ 是 $\mathcal{C}_\Lambda$ 的对象。设
$D : k \to k$ 是 $k$ 在 $\Lambda$ 上的一个导子，例如
$D = \partial/\partial x_i$。则映射
$$
f_D : k \longrightarrow k[\epsilon], \quad
a \mapsto a + D(a)\epsilon
$$
是 $\mathcal{C}_\Lambda$ 中的态射。令 $A_1 = A_2 = k$，并令
$f_1 = f_{\partial/\partial x_1}$、$f_2(a) = a$。于是
$A_1 \times_A A_2 = \{a \in k \mid \partial/\partial x_1(a) = 0\}$，
它不满射到 $k$。因此该纤维积不是 $\mathcal{C}_\Lambda$ 的对象。
\end{example}

\noindent
可以证明，这一问题只会在剩余域扩张
$k/k'$（见 (\ref{formal-defos-equation-k-prime})）不可分，并且
$f_1$、$f_2$ 都不是满射时发生。

\begin{lemma}
\label{formal-defos-lemma-fiber-product-CLambda}
设 $f_1 : A_1 \to A$ 和 $f_2 : A_2 \to A$ 是
$\mathcal{C}_\Lambda$ 中的环映射。则：
\begin{enumerate}
\item 若 $f_1$ 或 $f_2$ 为满射，则
$A_1 \times_A A_2$ 属于 $\mathcal{C}_\Lambda$；
\item 若 $f_2$ 为小扩张，则
$A_1 \times_A A_2 \to A_1$ 也是小扩张；
\item 若域扩张 $k/k'$ 可分，则
$A_1 \times_A A_2$ 属于 $\mathcal{C}_\Lambda$。
\end{enumerate}
\end{lemma}

\begin{proof}
映射 $\Lambda \to A_1$ 和 $\Lambda \to A_2$ 诱导映射
$\Lambda \to A_1 \times_A A_2$，从而使环
$A_1 \times_A A_2$ 成为 $\Lambda$-代数。它是局部环，唯一极大理想为
$$
\mathfrak m_{A_1} \times_{\mathfrak m_A} \mathfrak m_{A_2} =
\Ker(A_1 \times_A A_2 \longrightarrow k)
$$
一个环是阿廷环，当且仅当它作为自身的模具有有限长度；见《代数》引理
\ref{algebra-lemma-artinian-finite-length}。
由于 $A_1$ 和 $A_2$ 是阿廷环，引理 \ref{formal-defos-lemma-length} 表明
$\text{length}_\Lambda(A_1)$ 和 $\text{length}_\Lambda(A_2)$，
从而也表明 $\text{length}_\Lambda(A_1 \times A_2)$，都是有限的。
又因 $A_1 \times_A A_2 \subset A_1 \times A_2$ 是
$\Lambda$-子模，便有
$\text{length}_{A_1 \times_A A_2}(A_1 \times_A A_2) \leq
\text{length}_\Lambda(A_1 \times_A A_2)$ 有限。因此
$A_1 \times_A A_2$ 是阿廷环。故阻碍
$A_1 \times_A A_2$ 成为 $\mathcal{C}_\Lambda$ 对象的唯一可能，
是它的剩余域经上述映射
$A_1 \times_A A_2 \to A \to A/\mathfrak m_A = k$
映到 $k$ 的一个真子域。

\medskip\noindent
(1) 的证明。若 $f_2$ 为满射，则投影
$A_1 \times_A A_2 \to A_1$ 为满射。因此复合
$A_1 \times_A A_2 \to A_1 \to A_1/\mathfrak m_{A_1} = k$
为满射，从而 $A_1 \times_A A_2$ 是
$\mathcal{C}_\Lambda$ 的对象。

\medskip\noindent
(2) 的证明。若 $f_2$ 为小扩张，则 $A_2 \to A$ 与
$A_1 \times_A A_2 \to A_1$ 都是满射，且二者具有相同的核。
因此 $A_1 \times_A A_2 \to A_1$ 的核是 $1$-维
$k$-向量空间，故 $A_1 \times_A A_2 \to A_1$ 为小扩张。

\medskip\noindent
(3) 的证明。选取 $\overline{x} \in k$，使
$k = k'(\overline{x})$（见《域》引理
\ref{fields-lemma-primitive-element}）。
设 $P'(T) \in k'[T]$ 是 $\overline{x}$ 在 $k'$ 上的极小多项式。
由于 $k/k'$ 可分，
$\text{d}P/\text{d}T(\overline{x}) \not = 0$。
在满射 $\Lambda[T] \to k'[T]$ 下，选取映到 $P'$ 的首一多项式
$P \in \Lambda[T]$。由于 $A, A_1, A_2$ 都是 Hensel 环；见《代数》引理
\ref{algebra-lemma-local-dimension-zero-henselian}，
可找到 $x, x_1, x_2 \in A, A_1, A_2$，使
$P(x) = 0, P(x_1) = 0, P(x_2) = 0$，且 $x,x_1,x_2$
在 $k$ 中的像均为 $\overline{x}$。
由唯一性，$x_1,x_2$ 都映到 $x \in A$；见《代数》引理
\ref{algebra-lemma-uniqueness}。所以
$(x_1, x_2) \in A_1 \times_A A_2$。
于是 $A_1 \times_A A_2$ 的剩余域包含 $k$ 在 $k'$ 上的一个生成元，
结论得证。
\end{proof}


\noindent
下面定义 $\mathcal{C}_\Lambda$ 中的本质满射。
引理 \ref{formal-defos-lemma-essential-surjection} 给出了
$\mathcal{C}_\Lambda$ 中一个满射为本质满射的充要条件。

\begin{definition}
\label{formal-defos-definition-essential-surjection}
设 $f: B \to A$ 是 $\mathcal{C}_\Lambda$ 中的环映射。
若其具有下列性质，则称 $f$ 为{\it 本质满射}：
\begin{enumerate}
\item $f$ 为满射；
\item 若 $g: C \to B$ 是 $\mathcal{C}_\Lambda$ 中的环映射，
且 $f \circ g$ 为满射，则 $g$ 为满射。
\end{enumerate}
\end{definition}

\noindent
利用引理 \ref{formal-defos-lemma-surjective}，可以如下刻画
$\mathcal{C}_\Lambda$ 中的本质满射。

\begin{lemma}
\label{formal-defos-lemma-essential-surjection-mod-squares}
设 $f: B \to A$ 是 $\mathcal{C}_\Lambda$ 中的环映射。
下列条件等价：
\begin{enumerate}
\item $f$ 是本质满射；
\item 映射 $B/\mathfrak m_B^2 \to A/\mathfrak m_A^2$ 是本质满射；
\item 映射
$B/(\mathfrak m_\Lambda B + \mathfrak m_B^2) \to
A/(\mathfrak m_\Lambda A + \mathfrak m_A^2)$ 是本质满射。
\end{enumerate}
\end{lemma}

\begin{proof}
假设 (3) 成立。设 $C \to B$ 是 $\mathcal{C}_\Lambda$ 中的环映射，
且 $C \to A$ 为满射。则
$C \to A/(\mathfrak m_\Lambda A + \mathfrak m_A^2)$ 也为满射。
由假设，$C \to B/(\mathfrak m_\Lambda B + \mathfrak m_B^2)$
为满射。两次应用引理 \ref{formal-defos-lemma-surjective}，可知
$C \to B$ 为满射。

\medskip\noindent
假设 (1) 成立。设
$C \to B/(\mathfrak m_\Lambda B + \mathfrak m_B^2)$
是 $\mathcal{C}_\Lambda$ 中的态射，且
$C \to A/(\mathfrak m_\Lambda A + \mathfrak m_A^2)$ 为满射。令
$C' = C \times_{B/(\mathfrak m_\Lambda B + \mathfrak m_B^2)} B$。
由引理 \ref{formal-defos-lemma-fiber-product-CLambda}，它是
$\mathcal{C}_\Lambda$ 的对象。
注意 $C' \to A/(\mathfrak m_\Lambda A + \mathfrak m_A^2)$
仍为满射，故由引理 \ref{formal-defos-lemma-surjective}，$C' \to A$ 为满射。
于是由假设，$C' \to B$ 为满射。因而
$C' \to B/(\mathfrak m_\Lambda B + \mathfrak m_B^2)$ 为满射；
由 $C'$ 的构造，这又蕴含
$C \to B/(\mathfrak m_\Lambda B + \mathfrak m_B^2)$ 为满射。

\medskip\noindent
第一段证明了 (3) $\Rightarrow$ (1)，第二段证明了
(1) $\Rightarrow$ (3)。(2) 与 (3) 的等价，是 (1) 与 (3)
等价的一个特殊情形，故证明完成。
\end{proof}

\noindent
为了进一步分析 $\mathcal{C}_\Lambda$ 中的本质满射，
引入一些记号。设 $A$ 是 $\mathcal{C}_\Lambda$ 的对象；
更一般地，也可取配备了 $\Lambda$-代数满射 $A \to k$
的任意 $\Lambda$-代数。有典范正合序列
\begin{equation}
\label{formal-defos-equation-sequence}
\mathfrak m_A/\mathfrak m_A^2 \xrightarrow{\text{d}_A}
\Omega_{A/\Lambda} \otimes_A k \to
\Omega_{k/\Lambda} \to 0
\end{equation}
见《代数》引理 \ref{algebra-lemma-differential-seq}。
注意 $\Omega_{k/\Lambda} = \Omega_{k/k'}$，其中 $k'$ 如
(\ref{formal-defos-equation-k-prime}) 所定义。以 $H_1(L_{k/\Lambda})$
表示 $k$ 在 $\Lambda$ 上的朴素余切复形的第一同调模；见《代数》定义
\ref{algebra-definition-naive-cotangent-complex}。
于是可以把 (\ref{formal-defos-equation-sequence}) 延长为正合序列
\begin{equation}
\label{formal-defos-equation-sequence-extended}
H_1(L_{k/\Lambda}) \to
\mathfrak m_A/\mathfrak m_A^2 \xrightarrow{\text{d}_A}
\Omega_{A/\Lambda} \otimes_A k \to
\Omega_{k/\Lambda} \to 0,
\end{equation}
见《代数》引理 \ref{algebra-lemma-exact-sequence-NL}。
若 $B \to A$ 是 $\mathcal{C}_\Lambda$ 中的环映射；
更一般地，若它是两个均配备到 $k$ 的 $\Lambda$-代数满射的
$\Lambda$-代数之间的映射，则得到交换图
\begin{equation}
\label{formal-defos-equation-sequence-functorial}
\vcenter{
\xymatrix{
H_1(L_{k/\Lambda}) \ar[r] \ar@{=}[d] &
\mathfrak m_B/\mathfrak m_B^2 \ar[r]_{\text{d}_B} \ar[d] &
\Omega_{B/\Lambda} \otimes_B k \ar[r] \ar[d] &
\Omega_{k/\Lambda} \ar[r] \ar@{=}[d] & 0 \\
H_1(L_{k/\Lambda}) \ar[r] &
\mathfrak m_A/\mathfrak m_A^2 \ar[r]^{\text{d}_A} &
\Omega_{A/\Lambda} \otimes_A k \ar[r] &
\Omega_{k/\Lambda} \ar[r] & 0
}
}
\end{equation}
其两行均正合。

\begin{lemma}
\label{formal-defos-lemma-H1-separable-case}
有典范映射
$$
\mathfrak m_\Lambda/\mathfrak m_\Lambda^2 \longrightarrow H_1(L_{k/\Lambda}).
$$
若 $k' \subset k$ 可分（例如 $k$ 的特征为零），则该映射诱导同构
$\mathfrak m_\Lambda/\mathfrak m_\Lambda^2 \otimes_{k'} k = H_1(L_{k/\Lambda})$。
若 $k = k'$（例如经典情形），则
$\mathfrak m_\Lambda/\mathfrak m_\Lambda^2 = H_1(L_{k/\Lambda})$。
复合
$$
\mathfrak m_\Lambda/\mathfrak m_\Lambda^2 \longrightarrow
H_1(L_{k/\Lambda}) \longrightarrow \mathfrak m_A/\mathfrak m_A^2
$$
来自典范映射 $\mathfrak m_\Lambda \to \mathfrak m_A$。
\end{lemma}

\begin{proof}
由于 $\Lambda \to k'$ 为满射且核为 $\mathfrak m_\Lambda$，
故 $H_1(L_{k'/\Lambda}) = \mathfrak m_\Lambda/\mathfrak m_\Lambda^2$。
上述映射来自朴素余切复形的函子性。若 $k' \subset k$ 可分，
则 $k' \to k$ 是平展环映射；见《代数》引理
\ref{algebra-lemma-etale-over-field}。
因此它的朴素余切复形具有平凡的同调群；见《代数》定义
\ref{algebra-definition-etale}。
把《代数》引理 \ref{algebra-lemma-exact-sequence-NL}
应用于环映射 $\Lambda \to k' \to k$，可得
$\mathfrak m_\Lambda/\mathfrak m_\Lambda^2 \otimes_{k'} k = H_1(L_{k/\Lambda})$。
最后一个断言的证明从略。
\end{proof}


\begin{lemma}
\label{formal-defos-lemma-essential-surjection}
设 $f: B \to A$ 是 $\mathcal{C}_\Lambda$ 中的环映射，
记号沿用 (\ref{formal-defos-equation-sequence-functorial})。
\begin{enumerate}
\item 引理 \ref{formal-defos-lemma-surjective} 中刻画 $f$ 为满射的等价条件，
还等价于下列条件：
\begin{enumerate}
\item $\Im(\text{d}_B) \to \Im(\text{d}_A)$ 为满射；
\item 映射 $\Omega_{B/\Lambda} \otimes_B k \to
\Omega_{A/\Lambda} \otimes_A k$ 为满射。
\end{enumerate}
\item 下列条件等价：
\begin{enumerate}
\item $f$ 是本质满射（见引理
\ref{formal-defos-lemma-essential-surjection-mod-squares}）；
\item 映射 $\Im(\text{d}_B) \to \Im(\text{d}_A)$ 为同构；
\item 映射 $\Omega_{B/\Lambda} \otimes_B k \to
\Omega_{A/\Lambda} \otimes_A k$ 为同构。
\end{enumerate}
\item 若 $k/k'$ 可分，则 $f$ 为本质满射，当且仅当映射
$\mathfrak m_B/(\mathfrak m_\Lambda B + \mathfrak m_B^2) \to
\mathfrak m_A/(\mathfrak m_\Lambda A + \mathfrak m_A^2)$
为同构。
\item 若 $f$ 为小扩张，则 $f$ 非本质，当且仅当 $f$ 在
$\mathcal{C}_\Lambda$ 中有截面 $s: A \to B$，且
$f \circ s = \text{id}_A$。
\end{enumerate}
\end{lemma}

\begin{proof}
(1) 的证明。由 (\ref{formal-defos-equation-sequence-functorial})，
(1)(a) 与 (1)(b) 等价。此外，若 $A \to B$ 为满射，% P10-E0423
则 (1)(a) 与 (1)(b) 成立。设 (1)(a) 成立。
由于 $\text{d}_A$ 的核是 $H_1(L_{k/\Lambda})$ 的像，
而后者也映到 $\mathfrak m_B/\mathfrak m_B^2$，可知
$\mathfrak m_B/\mathfrak m_B^2 \to \mathfrak m_A/\mathfrak m_A^2$
为满射。因此由引理 \ref{formal-defos-lemma-surjective}，$B \to A$ 为满射。
这就完成了 (1) 的证明。

\medskip\noindent
(2) 的证明。(2)(b) 与 (2)(c) 的等价性直接来自
(\ref{formal-defos-equation-sequence-functorial})。

\medskip\noindent
设 (2)(b) 成立。令 $g : C \to B$ 为
$\mathcal{C}_\Lambda$ 中的环映射，且 $f \circ g$ 为满射。
由引理 \ref{formal-defos-lemma-surjective}，
$\mathfrak m_C/\mathfrak m_C^2 \to \mathfrak m_A/\mathfrak m_A^2$
为满射。因此 $\Im(\text{d}_C) \to \Im(\text{d}_A)$ 为满射；
再由假设可知 $\Im(\text{d}_C) \to \Im(\text{d}_B)$ 为满射。
由 (1)，$C \to B$ 为满射。

\medskip\noindent
设 (2)(a) 成立。则 $f$ 为满射，且
$\Omega_{B/\Lambda} \otimes_B k \to \Omega_{A/\Lambda} \otimes_A k$
为满射。令 $K$ 为其核。由
(\ref{formal-defos-equation-sequence-functorial})，注意到
$K = \text{d}_B(\Ker(\mathfrak m_B/\mathfrak m_B^2 \to
\mathfrak m_A/\mathfrak m_A^2))$。
选取 $k$-向量空间的分裂
$$
\Omega_{B/\Lambda} \otimes_B k =
\Omega_{A/\Lambda} \otimes_A k \oplus K
$$
映射 $\text{d} : B \to \Omega_{B/\Lambda}$ 经投影到 $K$
诱导映射 $D : B \to K$。令
$C = \{b \in B \mid D(b) = 0\}$。Leibniz 法则表明，
这是 $B$ 的一个 $\Lambda$-子代数。取 $\overline{x} \in k$，
再取映到 $\overline{x}$ 的 $x \in B$。若 $D(x) \not = 0$，
则可找到 $y \in \mathfrak m_B$，使 $D(y) = D(x)$。
于是 $x-y \in C$ 映到 $\overline{x}$。因此 $C \to k$ 为满射，
且 $C$ 是 $\mathcal{C}_\Lambda$ 的对象。类似地，取
$\omega \in \Im(\text{d}_A)$。由 (1)，可找到
$x \in \mathfrak m_B$，使 $\text{d}_B(x)$ 映到 $\omega$。
若 $D(x) \not = 0$，则可找到 $y \in \mathfrak m_B$，
它在 $\mathfrak m_A/\mathfrak m_A^2$ 中映为零，且
$D(y) = D(x)$。于是 $z=x-y$ 属于 $\mathfrak m_C$，而其像
$\text{d}_C(z) \in \Omega_{C/k} \otimes_C k$% P10-E0424
映到 $\omega$。故
$\Im(\text{d}_C) \to \Im(\text{d}_A)$ 为满射。
由 (1)，$C \to A$ 为满射；再由假设，$C \to B$ 为满射。
所以 $D=0$，也就是 $K=0$，即 (2)(c) 成立。
这完成了 (2) 的证明。

\medskip\noindent
(3) 的证明。若 $k'/k$ 可分，% P10-E0425
则由引理 \ref{formal-defos-lemma-H1-separable-case}，
$H_1(L_{k/\Lambda}) =
\mathfrak m_\Lambda/\mathfrak m_\Lambda^2 \otimes_{k'} k$。
于是
$\Im(\text{d}_A) =
\mathfrak m_A/(\mathfrak m_\Lambda A + \mathfrak m_A^2)$，
对 $B$ 也同样成立。因此 (3) 由 (2) 得出。

\medskip\noindent
(4) 的证明。$f$ 的截面 $s$ 不是满射（按定义，小扩张有非平凡核），
故 $f$ 不是本质满射。反过来，设 $f$ 是小扩张但不是本质满射。
选取 $\mathcal{C}_\Lambda$ 中的非满环映射 $C \to B$，
使 $C \to A$ 为满射。令 $C' \subset B$ 为 $C \to B$ 的像。
则 $C' \not = B$，但 $C'$ 满射到 $A$。由于
$f : B \to A$ 是小扩张，
$\text{length}_C(B) = \text{length}_C(A) + 1$。
又因 $C'$ 是 $B$ 的真子环，
$\text{length}_C(C') \leq \text{length}_C(A)$。
但 $C' \to A$ 为满射，所以事实上
$\text{length}_C(C') = \text{length}_C(A)$，且
$C' \to A$ 为同构。这便给出所需截面。
\end{proof}

\begin{example}
\label{formal-defos-example-essential-surjection}
设 $\Lambda = k[[x]]$ 是 $k$ 上 $1$ 个变量的幂级数环。
令 $A=k$、$B=\Lambda/(x^2)$。则由引理
\ref{formal-defos-lemma-essential-surjection}，$B \to A$ 是本质满射；
这是因为它是小扩张，且映射 $B \to A$ 在范畴
$\mathcal{C}_\Lambda$ 中没有右逆。然而映射
$$
k \cong \mathfrak m_B/\mathfrak m_B^2
\longrightarrow
\mathfrak m_A/\mathfrak m_A^2 = 0
$$
不是同构。因此，在引理 \ref{formal-defos-lemma-essential-surjection} (3) 中，
必须考虑相对余切空间之间的映射
$\mathfrak m_B/(\mathfrak m_\Lambda B + \mathfrak m_B^2) \to
\mathfrak m_A/(\mathfrak m_\Lambda A + \mathfrak m_A^2)$。
\end{example}


\section{完备化的基范畴}
\label{formal-defos-section-category-completion-CLambda}

\noindent
下面对范畴 $\mathcal{C}_\Lambda$ 的“完备化”，将作为
$\mathcal{C}_\Lambda$ 上群胚余纤维范畴之完备化的基范畴
（第 \ref{formal-defos-section-formal-objects} 节）。

\begin{definition}
\label{formal-defos-definition-completion-CLambda}
设 $\Lambda$ 是诺特环，且 $\Lambda \to k$ 是到域 $k$ 的有限环映射。
定义{\it $\widehat{\mathcal{C}}_\Lambda$} 为如下范畴：
\begin{enumerate}
\item 对象为偶 $(R, \varphi)$，其中 $R$ 是诺特完备局部
$\Lambda$-代数，且 $\varphi : R/\mathfrak m_R \to k$ 是
$\Lambda$-代数同构；
\item 态射 $f : (S, \psi) \to (R, \varphi)$ 为局部
$\Lambda$-代数同态，且满足
$\varphi \circ (f \bmod \mathfrak m) = \psi$。
\end{enumerate}
\end{definition}

\noindent
与定义 \ref{formal-defos-definition-CLambda} 之后的讨论一样，
我们通常把 $\widehat{\mathcal{C}}_\Lambda$ 的对象简记为 $R$，
并约定识别 $R/\mathfrak m_R = k$。
本节讨论范畴 $\widehat{\mathcal{C}}_\Lambda$ 的对象与态射的一些基本性质，
这些性质与上一节关于范畴 $\mathcal{C}_\Lambda$ 的讨论相对应。

\medskip\noindent
首先注意，任意对象 $A \in \mathcal{C}_\Lambda$ 也是
$\widehat{\mathcal{C}}_\Lambda$ 的对象，因为阿廷局部环总是诺特的，
并且关于其极大理想完备（该理想毕竟是幂零理想）。此外，由定义显然可见
$\mathcal{C}_\Lambda \subset \widehat{\mathcal{C}}_\Lambda$
是由所有阿廷环构成的严格满子范畴。反过来，事实上
$\widehat{\mathcal{C}}_\Lambda$ 的每个对象都是
$\mathcal{C}_\Lambda$ 中对象的极限。

\medskip\noindent
设 $R$ 是 $\widehat{\mathcal{C}}_\Lambda$ 的对象。
对 $n \in \mathbf{N}$，考虑环 $R_n = R/\mathfrak m_R^n$。
它们是具有唯一幂零素理想的诺特局部环，因而是阿廷环；见
《代数》命题 \ref{algebra-proposition-dimension-zero-ring}。
环映射
$$
\ldots \to R_{n + 1} \to R_n \to \ldots \to R_2 \to R_1 = k
$$
全都满射。按定义，$R$ 的完备性意即 $R = \lim R_n$。
若 $f : R \to S$ 是 $\widehat{\mathcal{C}}_\Lambda$ 中的环映射，
则得到一族环映射 $f_n : R_n \to S_n$，其极限就是给定映射。

\begin{lemma}
\label{formal-defos-lemma-surjective-cotangent-space}
设 $f: R \to S$ 是 $\widehat{\mathcal{C}}_\Lambda$ 中的环映射。
下列条件等价：
\begin{enumerate}
\item $f$ 为满射；
\item 映射
$\mathfrak m_R/\mathfrak m_R^2 \to \mathfrak m_S/\mathfrak m_S^2$
为满射；
\item 映射
$\mathfrak m_R/(\mathfrak m_\Lambda R + \mathfrak m_R^2) \to
\mathfrak m_S/(\mathfrak m_\Lambda S + \mathfrak m_S^2)$
为满射。
\end{enumerate}
\end{lemma}

\begin{proof}
注意，对 $n \geq 2$，相对余切空间之间有等式
$$
\mathfrak m_R/(\mathfrak m_\Lambda R + \mathfrak m_R^2)
=
\mathfrak m_{R_n}/(\mathfrak m_\Lambda R_n + \mathfrak m_{R_n}^2)
$$
对 $S$ 也同样成立。因此由引理 \ref{formal-defos-lemma-surjective}，
$R_n \to S_n$ 对所有 $n$ 都是满射。令 $K_n$ 为
$R_n \to S_n$ 的核。则序列
$$
0 \to K_n \to R_n \to S_n \to 0
$$
组成逆系统的正合序列。由于每个 $K_n$ 都是阿廷模，系统 $(K_n)$
满足 Mittag-Leffler 条件。因此由《代数》引理
\ref{algebra-lemma-ML-exact-sequence}，取极限保持正合性。
所以 $\lim R_n \to \lim S_n$ 为满射，即 $f$ 为满射。
\end{proof}

\begin{lemma}
\label{formal-defos-lemma-CLambdahat-pushouts}
范畴 $\widehat{\mathcal{C}}_\Lambda$ 有推出。
\end{lemma}

\begin{proof}
设 $R \to S_1$ 和 $R \to S_2$ 是
$\widehat{\mathcal{C}}_\Lambda$ 中的态射。考虑环
$C = S_1 \otimes_R S_2$。这个环有有限生成的极大理想
$\mathfrak m = \mathfrak m_{S_1} \otimes S_2 +
S_1 \otimes \mathfrak m_{S_2}$，其剩余域为 $k$。
令 $C^\wedge$ 为 $C$ 关于 $\mathfrak m$ 的完备化。
则 $C^\wedge$ 是关于极大理想
$\mathfrak m^\wedge = \mathfrak mC^\wedge$ 完备的诺特环，
其剩余域与 $k$ 识别；见《代数》引理
\ref{algebra-lemma-completion-Noetherian}。
因此 $C^\wedge$ 是 $\widehat{\mathcal{C}}_\Lambda$ 的对象。
于是 $S_1 \to C^\wedge$ 和 $S_2 \to C^\wedge$ 使
$C^\wedge$ 成为 $\widehat{\mathcal{C}}_\Lambda$ 中
$R$ 上的推出（细节从略）。
\end{proof}

\noindent
我们不会用到下面的引理。

\begin{lemma}
\label{formal-defos-lemma-CLambdahat-coproducts}
范畴 $\widehat{\mathcal{C}}_\Lambda$ 有对象对的余积。
\end{lemma}

\begin{proof}
设 $R$ 和 $S$ 是 $\widehat{\mathcal{C}}_\Lambda$ 的对象。
考虑环 $C = R \otimes_\Lambda S$。有典范满射
$C \to R \otimes_\Lambda S \to k \otimes_\Lambda k \to k$，
其中最后一个映射为乘法映射。$C \to k$ 的核是极大理想
$\mathfrak m$。注意，$\mathfrak m$ 由 $\mathfrak m_R C$、
$\mathfrak m_S C$，以及 $C$ 中有限多个映到
$k \otimes_\Lambda k \to k$ 之核的生成元的元素所生成。
因此 $\mathfrak m$ 是有限生成理想。令 $C^\wedge$ 为
$C$ 关于 $\mathfrak m$ 的完备化。则 $C^\wedge$ 是关于极大理想
$\mathfrak m^\wedge = \mathfrak mC^\wedge$ 完备的诺特环，
剩余域为 $k$；见《代数》引理
\ref{algebra-lemma-completion-Noetherian}。
因此 $C^\wedge$ 是 $\widehat{\mathcal{C}}_\Lambda$ 的对象。
于是 $R \to C^\wedge$ 和 $S \to C^\wedge$ 使
$C^\wedge$ 成为 $\widehat{\mathcal{C}}_\Lambda$ 中的余积
（细节从略）。
\end{proof}

\noindent
范畴中的空余积就是该范畴的始对象。在经典情形下，
$\widehat{\mathcal{C}}_\Lambda$ 有始对象，即 $\Lambda$ 本身。
更一般地，若 $k' = k$，则 $\Lambda$ 关于
$\mathfrak m_\Lambda$ 的完备化 $\Lambda^\wedge$ 是始对象。
还可以更一般地说，若 $k' \subset k$ 可分，则
$\widehat{\mathcal{C}}_\Lambda$ 也有始对象。具体地，选取首一多项式
$P \in \Lambda[T]$，使 $k \cong k'[T]/(P')$，其中
$p' \in k'[T]$ 是 $P$ 的像。% P10-E0426
则 $R = \Lambda^\wedge[T]/(P)$ 是始对象；见引理
\ref{formal-defos-lemma-fiber-product-CLambda} 的证明。

\medskip\noindent
若 $R$ 是上述始对象，则有
$\mathcal{C}_\Lambda = \mathcal{C}_R$ 和
$\widehat{\mathcal{C}}_\Lambda = \widehat{\mathcal{C}}_R$，
这实际上把本章的全部讨论归结到经典情形。
但是，若 $k' \subset k$ 不可分，则始对象不存在。

\begin{lemma}
\label{formal-defos-lemma-derivations-finite}
设 $S$ 是 $\widehat{\mathcal{C}}_\Lambda$ 的对象。则
$\dim_k \text{Der}_\Lambda(S, k) < \infty$。
\end{lemma}

\begin{proof}
选取 $x_1, \ldots, x_n \in \mathfrak m_S$，使其像构成相对余切空间
$\mathfrak m_S/(\mathfrak m_\Lambda S + \mathfrak m_S^2)$
的一组 $k$-基。选取 $y_1, \ldots, y_m \in S$，使其在 $k$ 中的像
在 $k'$ 上生成 $k$。我们断言
$\dim_k \text{Der}_\Lambda(S, k) \leq n + m$。
为证明这一点，只需证明：若 $D(x_i) = 0$ 且 $D(y_j) = 0$，
则 $D = 0$。设 $a \in S$。可找到多项式
$P = \sum \lambda_J y^J$，其中 $\lambda_J \in \Lambda$，
使其在 $k$ 中的像与 $a$ 在 $k$ 中的像相同。
由对所有 $j$ 都有 $D(y_j) = 0$ 这一假设可知
$D(a - P) = D(a) - D(P) = D(a)$。因而可以假设
$a \in \mathfrak m_S$。写成 $a = \sum a_i x_i$，其中
$a_i \in S$。由 Leibniz 法则，
$$
D(a) = \sum x_iD(a_i) + \sum a_iD(x_i) = \sum x_iD(a_i)
$$
因为我们假设了 $D(x_i) = 0$。又因 $x_i$ 在 $k$ 上的乘法为零，
所以 $\sum x_iD(a_i) = 0$。
\end{proof}

\begin{lemma}
\label{formal-defos-lemma-derivations-surjective}
设 $f : R \to S$ 是 $\widehat{\mathcal{C}}_\Lambda$ 的态射。
若 $\text{Der}_\Lambda(S, k) \to \text{Der}_\Lambda(R, k)$ 为单射，
则 $f$ 为满射。
\end{lemma}

\begin{proof}
若 $f$ 不是满射，则由引理
\ref{formal-defos-lemma-surjective-cotangent-space}，
$\mathfrak m_S/(\mathfrak m_R S + \mathfrak m_S^2)$ 非零。
因而 $Q = S/(f(R) + \mathfrak m_R S + \mathfrak m_S^2)$ 也非零。
通过 $f$，$Q$ 是 $k = R/\mathfrak m_R$-向量空间。
再通过 $S \to k$ 把 $Q$ 视为 $S$-模。商映射
$D : S \to Q$ 是一个 $R$-导子：若 $a_1, a_2 \in S$，
则可写成
$a_1 = f(b_1) + a_1'$ 和 $a_2 = f(b_2) + a_2'$，
其中 $b_1, b_2 \in R$ 且
$a_1', a_2' \in \mathfrak m_S$。此时对 $i = 1, 2$，
$b_i$ 和 $a_i$ 在 $k$ 中有相同的像，并且在 $Q$ 中
\begin{align*}
a_1a_2 & = (f(b_1) + a_1')(f(b_2) + a_2') \\
& = f(b_1)a_2' + f(b_2)a_1' \\
& = f(b_1)(f(b_2) + a_2') + f(b_2)(f(b_1) + a_1') \\
& = f(b_1)a_2 + f(b_2)a_1
\end{align*}
这就证明了 Leibniz 法则。因此 $D : S \to Q$ 是一个
$\Lambda$-导子，并且与 $R \to S$ 复合为零。
由于 $Q \not = 0$，也存在与 $R \to S$ 复合为零的导子
$D : S \to k$；也就是说，
$\text{Der}_\Lambda(S, k) \to \text{Der}_\Lambda(R, k)$
不是单射。
\end{proof}

\begin{lemma}
\label{formal-defos-lemma-m-adic-topology}
设 $R$ 是 $\widehat{\mathcal{C}}_\Lambda$ 的对象，$(J_n)$ 是满足
$\mathfrak m_R^n \subset J_n$ 的理想递减列。令
$J = \bigcap J_n$。则序列 $(J_n/J)$ 在 $R/J$ 上定义
$\mathfrak m_{R/J}$-进拓扑。
\end{lemma}

\begin{proof}
显然 $\mathfrak m_{R/J}^n \subset J_n/J$。因此只需证明：
对每个 $n$，存在 $N$，使
$J_N/J \subset \mathfrak m_{R/J}^n$。这等价于
$J_N \subset \mathfrak m_R^n + J$。对每个 $n$，环
$R/\mathfrak m_R^n$ 都是阿廷环，故存在 $N_n$，使
$$
J_{N_n} + \mathfrak m_R^n = J_{N_n + 1} + \mathfrak m_R^n = \ldots
$$
令 $E_n = (J_{N_n} + \mathfrak m_R^n)/\mathfrak m_R^n$，
并令 $E = \lim E_n \subset \lim R/\mathfrak m_R^n = R$。
注意 $E \subset J$：因为对任意 $f \in E$ 和任意 $m$，
当所有 $n \gg 0$ 时均有
$f \in J_m + \mathfrak m_R^n$，故由 Krull 交定理
$f \in J_m$；见《代数》引理
\ref{algebra-lemma-intersect-powers-ideal-module-zero}。
由于过渡映射 $E_n \to E_{n - 1}$ 全都满射，
$J$ 满射到 $E_n$。所以取 $N = N_n$ 即可。% P10-E0427
\end{proof}

\begin{lemma}
\label{formal-defos-lemma-limit-artinian}
设 $\ldots \to A_3 \to A_2 \to A_1$ 是
$\mathcal{C}_\Lambda$ 中的一列满环映射。若
$\dim_k (\mathfrak m_{A_n}/\mathfrak m_{A_n}^2)$ 有界，则
$S = \lim A_n$ 是 $\widehat{\mathcal{C}}_\Lambda$ 的对象，
并且理想 $I_n = \Ker(S \to A_n)$ 在 $S$ 上定义
$\mathfrak m_S$-进拓扑。
\end{lemma}

\begin{proof}
下文直接使用映射 $S \to A_n$ 对所有 $n$ 均为满射这一事实。
注意，映射
$\mathfrak m_{A_{n + 1}}/\mathfrak m_{A_{n + 1}}^2 \to
\mathfrak m_{A_n}/\mathfrak m_{A_n}^2$ 为满射；见引理
\ref{formal-defos-lemma-surjective-cotangent-space}。
因此，当 $n$ 充分大时，维数
$\dim_k (\mathfrak m_{A_n}/\mathfrak m_{A_n}^2)$ 稳定为某个整数，
记为 $r$。于是可找到
$x_1, \ldots, x_r \in \mathfrak m_S$，使其在 $A_n$ 中的像生成
$\mathfrak m_{A_n}$。此外，选取 $y_1, \ldots, y_t \in S$，
使其在 $k$ 中的像在 $\Lambda$ 上生成 $k$。从而得到环映射
$P = \Lambda[z_1, \ldots, z_{r + t}] \to S$，
$z_i \mapsto x_i$ 且 $z_{r + j} \mapsto y_j$，
使复合 $P \to S \to A_n$ 对所有 $n$ 都为满射。
令 $\mathfrak m \subset P$ 为 $P \to k$ 的核。
令 $R = P^\wedge$ 为 $P$ 的 $\mathfrak m$-进完备化；
这是 $\widehat{\mathcal{C}}_\Lambda$ 的对象。
仍有相容的满射系 $R \to A_n$，故得到映射 $R \to S$。
令 $J_n = \Ker(R \to A_n)$，并令 $J = \bigcap J_n$。
由引理 \ref{formal-defos-lemma-m-adic-topology}，
$R/J = \lim R/J_n = \lim A_n = S$，而且理想
$J_n/J = I_n$ 定义 $\mathfrak m$-进拓扑。
（注意，对每个 $n$，都有
$\mathfrak m_R^{N_n} \subset J_n$，其中 $N_n$ 依指标而定，
未必有 $N_n = n$；所以在应用该引理之前，可能需要重编号理想
$J_n$。）
\end{proof}

\begin{lemma}
\label{formal-defos-lemma-power-series}
设 $R', R \in \Ob(\widehat{\mathcal{C}}_\Lambda)$。
假设对 $R$ 的某个理想 $I$ 有 $R = R' \oplus I$。
取 $x_1, \ldots, x_r \in I$，使其像构成
$I/\mathfrak m_R I$ 的一组基。令
$S = R'[[X_1, \ldots, X_r]]$，并考虑把 $X_i$ 映到 $x_i$ 的
$R'$-代数映射 $S \to R$。假设对每个 $n \gg 0$，映射
$S/\mathfrak m_S^n \to R/\mathfrak m_R^n$ 在
$\mathcal{C}_\Lambda$ 中都有左逆。则 $S \to R$ 为同构。
\end{lemma}

\begin{proof}
由 $R = R' \oplus I$，有
$$
\mathfrak m_R/\mathfrak m_R^2 =
\mathfrak m_{R'}/\mathfrak m_{R'}^2 \oplus I/\mathfrak m_RI
$$
类似地，
$$
\mathfrak m_S/\mathfrak m_S^2 =
\mathfrak m_{R'}/\mathfrak m_{R'}^2 \oplus \bigoplus kX_i
$$
因此，当 $n > 1$ 时，映射
$S/\mathfrak m_S^n \to R/\mathfrak m_R^n$
在余切空间上诱导同构。于是由引理
\ref{formal-defos-lemma-surjective-cotangent-space}，左逆
$h_n : R/\mathfrak m_R^n \to S/\mathfrak m_S^n$ 为满射。
又因 $h_n$ 作为左逆是单射，它是同构。
所以典范满射 $S/\mathfrak m_S^n \to R/\mathfrak m_R^n$
全都是同构，结论得证。
\end{proof}

\section{群胚余纤维范畴}
\label{formal-defos-section-preliminary}

\noindent
在发展这套理论时，我们使用群胚{\it 余纤维}范畴。
我们假定群胚{\it 纤维}范畴的定义和基本性质为已知；见
《范畴》第 \ref{categories-section-fibred-groupoids} 节。

\begin{definition}
\label{formal-defos-definition-category-cofibred-groupoids}
设 $\mathcal{C}$ 是范畴。$\mathcal{C}$ 上的一个
{\it 群胚余纤维范畴}，是配备函子
$p: \mathcal{F} \to \mathcal{C}$ 的范畴 $\mathcal{F}$，
使得经由
$p^{opp}: \mathcal{F}^{opp} \to \mathcal{C}^{opp}$，
$\mathcal{F}^{opp}$ 成为 $\mathcal{C}^{opp}$ 上的群胚纤维范畴。
\end{definition}

\noindent
明确地说，若下列两个条件成立，则
$p: \mathcal{F} \to \mathcal{C}$ 是群胚余纤维范畴：
\begin{enumerate}
\item 对 $\mathcal{C}$ 中每个态射 $f: U \to V$ 以及
$U$ 上的每个对象 $x$，存在 $\mathcal{F}$ 中位于 $f$ 上方的态射
$x \to y$。
\item 对 $\mathcal{F}$ 中每一对态射
$a: x \to y$ 和 $b: x \to z$，以及满足
$p(b) = f \circ p(a)$ 的任意态射
$f: p(y) \to p(z)$，存在唯一的 $\mathcal F$ 中态射
$c: y \to z$，它位于 $f$ 上方且满足 $b = c \circ a$。
\end{enumerate}

\begin{remarks}
\label{formal-defos-remarks-cofibered-groupoids}
关于群胚纤维范畴的一切内容都直接转化到余纤维情形。
下面这些说明用于固定记号。设 $\mathcal{C}$ 是范畴。
\begin{enumerate}
\item 我们常从记号中省略函子
$p: \mathcal{F} \to \mathcal{C}$。
\item $\mathcal{C}$ 中对象 $U$ 上方的纤维范畴记作
$\mathcal{F}(U)$。其对象是 $\mathcal{F}$ 中位于 $U$ 上方的对象，
其态射是 $\mathcal{F}$ 中位于 $\text{id}_U$ 上方的态射。
若 $x, y$ 是 $\mathcal{F}(U)$ 的对象，有时把
$\Mor_{\mathcal{F}(U)}(x, y)$ 写作 $\Mor_U(x, y)$。
\item 纤维范畴 $\mathcal{F}(U)$ 是群胚；见《范畴》引理
\ref{categories-lemma-fibred-groupoids}。
因此 $\mathcal{F}(U)$ 中的所有态射都是同构。
有时把 $\Mor_{\mathcal{F}(U)}(x, x)$ 写作
$\text{Aut}_U(x)$。
\item
\label{formal-defos-item-pushforward}
设 $\mathcal{F}$ 是 $\mathcal{C}$ 上的群胚余纤维范畴，
$f: U \to V$ 是 $\mathcal{C}$ 中的态射，且
$x \in \Ob(\mathcal{F}(U))$。
$x$ 沿 $f$ 的一个{\it 推前}，是 $\mathcal{F}$ 中位于
$f$ 上方的态射 $x \to y$。推前在唯一同构意义下唯一
（见《范畴》定义 \ref{categories-definition-cartesian-over-C}
之后的讨论）。有时把 $x$ 沿 $f$ 的“这个”推前写作
$x \to f_*x$。
\item $\mathcal{F}$ 的一个{\it 推前选取}，是对上述每一对
$(x, f)$ 选取 $x$ 沿 $f$ 的一个推前。
由选择公理，可以为 $\mathcal{F}$ 作出这样的推前选取。
\item 设 $\mathcal{F}$ 是 $\mathcal{C}$ 上的群胚余纤维范畴。
给定 $\mathcal{F}$ 的一个推前选取，就有相伴的伪函子
$\mathcal{C} \to \textit{Groupoids}$。
我们绝不会用到这一构造，故略去细节。
\item
\label{formal-defos-item-cofibered-morphism}
$\mathcal{C}$ 上群胚余纤维范畴之间的态射，是与到
$\mathcal{C}$ 的投影交换的函子。若 $\mathcal{F}$ 和
$\mathcal{F}'$ 是 $\mathcal{C}$ 上的群胚余纤维范畴，
则从 $\mathcal{F}$ 到 $\mathcal{F}'$ 的态射记为
$\Mor_\mathcal{C}(\mathcal{F}, \mathcal{F}')$。
\item
\label{formal-defos-item-definition-cofibered-groupoids-2-category}
群胚余纤维范畴组成一个 $(2, 1)$-范畴
$\text{Cof}(\mathcal{C})$。它的 1-态射是
(\ref{formal-defos-item-cofibered-morphism}) 中所述的态射。
若 $p : \mathcal{F} \to C$ 和
$p': \mathcal{F}' \to \mathcal{C}$ 是群胚余纤维范畴，% P10-E0428
而 $\varphi, \psi : \mathcal{F} \to \mathcal{F}'$ 是
$1$-态射，则 2-态射 $t : \varphi \to \psi$ 是满足下述条件的
函子态射：对所有 $x \in \Ob(\mathcal{F})$，
$p'(t_x) = \text{id}_{p(x)}$。
\item
\label{formal-defos-item-construction-associated-cofibered-groupoid}
设 $F : \mathcal{C} \to \textit{Groupoids}$ 是函子。
如下可构造与 $F$ 相伴的群胚余纤维范畴
$\mathcal{F} \to \mathcal{C}$。
$\mathcal{F}$ 的对象是偶 $(U, x)$，其中
$U \in \Ob(\mathcal{C})$ 且 $x \in \Ob(F(U))$。
一个态射 $(U, x) \to (V, y)$ 是偶 $(f, a)$，其中
$f \in \Mor_\mathcal{C}(U, V)$ 且
$a \in \Mor_{F(V)}(F(f)(x), y)$。
函子 $\mathcal{F} \to \mathcal{C}$ 把 $(U, x)$ 映到 $U$。
见《范畴》第 \ref{categories-section-presheaves-groupoids} 节。
\item
\label{formal-defos-item-associated-functor-isomorphism-classes}
设 $\mathcal{F}$ 是 $\mathcal{C}$ 上的群胚余纤维范畴。
对 $U \in \Ob(\mathcal{C})$，令
$\overline{\mathcal{F}}(U)$ 等于范畴 $\mathcal{F}(U)$
的同构类集合。% P10-E0429
若 $f : U \to V$ 是 $\mathcal{C}$ 的态射，
则把 $x$ 的同构类映到 $x$ 的一个推前 $f_*x$ 的同构类，
便得到集合映射
$\overline{\mathcal{F}}(U) \to \overline{\mathcal{F}}(V)$；
见 (\ref{formal-defos-item-pushforward})。% P10-E0430
于是 $\overline{\mathcal{F}} : \mathcal{C} \to \textit{Sets}$
是函子。类似地，若
$\varphi : \mathcal{F} \to \mathcal{G}$ 是余纤维范畴之间的态射，
则以
$\overline{\varphi}: \overline{\mathcal{F}} \to  \overline{\mathcal{G}}$
表示相伴的函子态射。
\item
\label{formal-defos-item-convention-cofibered-sets}
设 $F: \mathcal{C} \to \textit{Sets}$ 是函子。
可以把集合视为离散范畴，即只有恒等态射的群胚。
于是构造 (\ref{formal-defos-item-construction-associated-cofibered-groupoid})
给出与 $F$ 相伴的集合余纤维范畴。
这定义了从函子范畴
$\mathcal{C} \to \textit{Sets}$ 到 $\mathcal{C}$ 上
群胚余纤维范畴之范畴的一个充分忠实嵌入。
我们把函子范畴与它在这一嵌入下的像识别。
因此，若 $F : \mathcal{C} \to \textit{Sets}$ 是函子，
则仍以 $F$ 表示相伴的集合余纤维范畴；若
$\varphi : F \to G$ 是函子态射，则仍以 $\varphi$
表示对应的集合余纤维范畴态射，反之亦然。
见《范畴》第 \ref{categories-section-fibred-in-sets} 节。
\item
\label{formal-defos-item-definition-yoneda}
设 $U$ 是 $\mathcal{C}$ 的对象。以 $\underline{U}$ 表示函子
$\Mor_\mathcal{C}(U, -): \mathcal{C} \to
\textit{Sets}$。这定义了从 $\mathcal
C^{opp}$ 到函子范畴
$\mathcal{C} \to \textit{Sets}$ 的充分忠实嵌入。
因此，若 $f : U \to V$ 是态射，则仍以 $f$ 表示所诱导的态射
$\underline{V} \to \underline{U}$，反之亦然。
\item
\label{formal-defos-item-fibre-product}
群胚余纤维范畴的纤维积：若
$\mathcal{F} \to \mathcal{H}$ 和
$\mathcal{G} \to \mathcal{H}$ 是
$\mathcal{C}_\Lambda$ 上群胚余纤维范畴之间的态射，
则可把它们作为 $\mathcal{C}_\Lambda$ 上范畴的 2-纤维积构造，
用作其 2-纤维积；见《范畴》引理
\ref{categories-lemma-2-product-categories-over-C} 中的描述。
\item
\label{formal-defos-item-product}
群胚余纤维范畴的积：若 $\mathcal{F}$ 和 $\mathcal{G}$ 是
$\mathcal{C}_\Lambda$ 上的群胚余纤维范畴，则按定义，
其积为 $2$-纤维积
$\mathcal{F} \times_{\mathcal{C}_\Lambda} \mathcal{G}$；
见《范畴》引理
\ref{categories-lemma-2-product-categories-over-C} 中的描述。
\item
\label{formal-defos-item-definition-restricting-base-category}
限制基范畴：设 $p : \mathcal{F} \to \mathcal{C}$ 是
群胚余纤维范畴，且 $\mathcal{C}'$ 是 $\mathcal{C}$ 的满子范畴。
限制 $\mathcal{F}|_{\mathcal{C}'}$ 是 $\mathcal{F}$ 的满子范畴，
其对象是位于 $\mathcal{C}'$ 对象上方的那些对象。
经由函子
$p|_{\mathcal{C}'}: \mathcal{F}|_{\mathcal{C}'} \to \mathcal{C}'$，
它是群胚余纤维范畴。
\end{enumerate}
\end{remarks}

\section{pro-可表函子与预形变范畴}
\label{formal-defos-section-cofibered-groupoids}

\noindent
我们的基本目标，是理解 $\mathcal{C}_\Lambda$ 和
$\widehat{\mathcal{C}}_\Lambda$ 上的群胚余纤维范畴。
由于 $\mathcal{C}_\Lambda$ 是
$\widehat{\mathcal{C}}_\Lambda$ 的满子范畴，可以把
$\widehat{\mathcal{C}}_\Lambda$ 上的群胚余纤维范畴限制到
$\mathcal{C}_\Lambda$；见说明
\ref{formal-defos-remarks-cofibered-groupoids}
(\ref{formal-defos-item-definition-restricting-base-category})。
尤其可对函子这样做，因而也可对可表函子这样做。
由此在 $\mathcal{C}_\Lambda$ 上所得的函子称为 pro-可表函子。

\begin{definition}
\label{formal-defos-definition-prorepresentable}
设 $F : \mathcal{C}_\Lambda \to \textit{Sets}$ 是函子。
若对某个 $R \in \Ob(\widehat{\mathcal{C}}_\Lambda)$，
存在函子同构
$F \cong \underline{R}|_{\mathcal{C}_\Lambda}$，
则称 $F$ 是{\it pro-可表的}。
\end{definition}

\noindent
注意，若 $F : \mathcal{C}_\Lambda \to \textit{Sets}$
由 $R \in \Ob(\widehat{\mathcal{C}}_\Lambda)$ pro-表示，则
$$
F(k) = \Mor_{\widehat{\mathcal{C}}_\Lambda}(R, k) = \{*\}
$$
是单元集。形变理论中出现的
$\mathcal{C}_\Lambda$ 上群胚余纤维范畴也常满足类似条件。% P10-E0431

\begin{definition}
\label{formal-defos-definition-predeformation-category}
一个{\it 预形变范畴} $\mathcal{F}$，是
$\mathcal{C}_\Lambda$ 上的群胚余纤维范畴，并且
$\mathcal{F}(k)$ 等价于恰有一个对象和一个态射的范畴；
也就是说，$\mathcal{F}(k)$ 至少含有一个对象，且任意两个对象之间
都有唯一态射。预形变范畴之间的一个{\it 态射}，是
$\mathcal{C}_\Lambda$ 上群胚余纤维范畴之间的态射。
\end{definition}

\noindent
预形变范畴具有如下特征。设
$x_0 \in \Ob(\mathcal{F}(k))$。则 $\mathcal{F}$ 的每个对象
都配备了到 $x_0$ 的唯一态射。事实上，若 $x$ 是
$\mathcal{F}$ 中位于 $A$ 上方的对象，则可以选取推前
$x \to q_*x$，其中 $q : A \to k$ 为商映射。
存在唯一同构 $q_*x \to x_0$，而复合
$x \to q_*x \to x_0$ 就是所需态射。

\begin{remark}
\label{formal-defos-remark-predeformation-functor}
若函子 $F: \mathcal{C}_\Lambda \to \textit{Sets}$ 的相伴集合
余纤维范畴是预形变范畴，即 $F(k)$ 是单元集，
则称其为{\it 预形变函子}。
因此，若 $\mathcal{F}$ 是预形变范畴，则
$\overline{\mathcal{F}}$ 是预形变函子。
\end{remark}

\begin{remark}
\label{formal-defos-remark-localize-cofibered-groupoid}
设 $p : \mathcal{F} \to \mathcal{C}_\Lambda$ 是群胚余纤维范畴，
且 $x \in \Ob(\mathcal{F}(k))$。以 $\mathcal{F}_x$
表示位于 $x$ 上的对象所成范畴。
$\mathcal{F}_x$ 的对象是箭头 $y \to x$。
$\mathcal{F}_x$ 中的一个态射
$(y \to x) \to (z \to x)$ 是交换图
$$
\xymatrix{
y \ar[rr] \ar[dr] & & z \ar[dl] \\
& x &
}
$$
有遗忘函子 $\mathcal{F}_x \to \mathcal{F}$。
把函子 $p_x : \mathcal{F}_x \to \mathcal{C}_\Lambda$ 定义为复合
$\mathcal{F}_x \to \mathcal{F} \xrightarrow{p} \mathcal{C}_\Lambda$。
于是 $p_x : \mathcal{F}_x \to \mathcal{C}_\Lambda$
是预形变范畴（证明从略）。由此，对
$\mathcal{C}_\Lambda$ 上的任意群胚余纤维范畴以及任意
$x \in \Ob(\mathcal{F}(k))$，都可得到一个预形变范畴。
\end{remark}

\section{形式对象与完备化范畴}
\label{formal-defos-section-formal-objects}

\noindent
本节讨论如何在 $\mathcal{C}_\Lambda$ 上的群胚余纤维范畴与
$\widehat{\mathcal{C}}_\Lambda$ 上的群胚余纤维范畴之间
往返转换。

\begin{definition}
\label{formal-defos-definition-formal-objects}
设 $\mathcal{F}$ 是 $\mathcal{C}_\Lambda$ 上的群胚余纤维范畴。
$\mathcal{F}$ 的{\it 形式对象范畴 $\widehat{\mathcal{F}}$}
是具有如下对象和态射的范畴。
\begin{enumerate}
\item $\mathcal{F}$ 的一个
{\it 形式对象 $\xi = (R, \xi_n, f_n)$}，由
$\widehat{\mathcal{C}}_\Lambda$ 的一个对象 $R$、
一族以 $n \in \mathbf{N}$ 为指标的对象 $\xi_n$，它们属于
$\mathcal{F}(R/\mathfrak m_R^n)$，以及位于投影
$R/\mathfrak m_R^{n + 1} \to R/\mathfrak m_R^n$
上方的态射 $f_n : \xi_{n + 1} \to \xi_n$ 所组成。
\item 设 $\xi = (R, \xi_n, f_n)$ 和
$\eta = (S, \eta_n, g_n)$ 是 $\mathcal{F}$ 的形式对象。
形式对象的一个{\it 态射 $a : \xi \to \eta$}，由
$\widehat{\mathcal{C}}_\Lambda$ 中的映射
$a_0 : R \to S$，以及一族 $\mathcal{F}$ 中位于
$R/\mathfrak m_R^n \to S/\mathfrak m_S^n$ 上方的态射
$a_n : \xi_n \to \eta_n$ 所组成，并且对每个 $n$，图
$$
\xymatrix{
\xi_{n + 1} \ar[r]_{f_n} \ar[d]_{a_{n + 1}} & \xi_n \ar[d]^{a_n} \\
\eta_{n + 1} \ar[r]^{g_n} & \eta_n
}
$$
交换。
\end{enumerate}
\end{definition}

\noindent
形式对象范畴配备函子
$\widehat{p}: \widehat{\mathcal{F}} \to
\widehat{\mathcal{C}}_\Lambda$；它把对象
$(R, \xi_n, f_n)$ 映到 $R$，并把态射
$(R, \xi_n, f_n) \to (S, \eta_n, g_n)$ 映到 $R \to S$。

\begin{lemma}
\label{formal-defos-lemma-completion-cofibred}
设 $p : \mathcal{F} \to \mathcal{C}_\Lambda$ 是群胚余纤维范畴。
则
$\widehat{p} : \widehat{\mathcal{F}} \to
\widehat{\mathcal{C}}_\Lambda$
是群胚余纤维范畴。
\end{lemma}

\begin{proof}
设 $R \to S$ 是 $\widehat{\mathcal{C}}_\Lambda$ 中的环映射，
且 $(R, \xi_n, f_n)$ 是 $\widehat{\mathcal{F}}$ 的对象。
对每个 $n$，选取 $\xi_n$ 沿
$R/\mathfrak m_R^n \to S/\mathfrak m_S^n$ 的一个推前
$\xi_n \to \eta_n$。对每个 $n$，存在唯一的
$\mathcal{F}$ 中态射 $g_n : \eta_{n + 1} \to \eta_n$，
它位于 $S/\mathfrak m_S^{n + 1} \to S/\mathfrak m_S^n$
上方，且使图
$$
\xymatrix{
\xi_{n + 1} \ar[d] \ar[r]_{f_n} & \xi_n \ar[d] \\
\eta_{n + 1} \ar[r]^{g_n} & \eta_n
}
$$
交换（由群胚余纤维范畴的第一条公理）。% P10-E0432
因而得到位于 $R \to S$ 上方的态射
$(R, \xi_n, f_n) \to (S, \eta_n, g_n)$；
也就是说，群胚余纤维范畴的第一条公理对
$\widehat{\mathcal{F}}$ 成立。为验证第二条公理，设有
$\widehat{\mathcal{F}}$ 中的态射
$a : (R, \xi_n, f_n) \to (S, \eta_n, g_n)$ 和
$b : (R, \xi_n, f_n) \to (T, \theta_n, h_n)$，
以及 $\widehat{\mathcal{C}}_\Lambda$ 中满足
$c_0 \circ a_0 = b_0$ 的态射 $c_0 : S \to T$。
由 $\mathcal{F}$ 的群胚余纤维范畴第二条公理，
得到位于 $S/\mathfrak m_S^n \to T/\mathfrak m_T^n$
上方的唯一映射 $c_n : \eta_n \to \theta_n$，使
$c_n \circ a_n = b_n$。令 $c = (c_n)_{n \geq 0}$，
便得到所需的 $\widehat{\mathcal{F}}$ 中态射
$c : (S, \eta_n, g_n) \to (T, \theta_n, h_n)$
（验证 $h_n \circ c_{n + 1} = c_n \circ g_n$ 的步骤从略）。
\end{proof}

\begin{definition}
\label{formal-defos-definition-completion}
设 $p : \mathcal{F} \to \mathcal{C}_\Lambda$ 是群胚余纤维范畴。
群胚余纤维范畴
$\widehat{p} : \widehat{\mathcal  F} \to
\widehat{\mathcal{C}}_\Lambda$
称为 $\mathcal{F}$ 的{\it 完备化}。
\end{definition}

\noindent
若 $\mathcal{F}$ 是 $\mathcal C_\Lambda$ 上的群胚余纤维范畴，
我们已经用 $R$ 的极大理想幂滤过定义了
$R \in \Ob(\widehat{\mathcal{C}}_\Lambda)$ 时的
$\widehat{\mathcal{F}}(R)$。现在设
$\mathcal{I} = (I_n)$ 是由理想给出的 $R$ 的滤过，
它诱导 $\mathfrak{m}_R$-进拓扑。
定义 $\widehat{\mathcal{F}}_\mathcal{I}(R)$ 为具有如下对象和态射的范畴：
\begin{enumerate}
\item 对象是一族
$(\xi_n, f_n)_{n \in \mathbf{N}}$，其中
$\xi_n$ 是 $\mathcal{F}(R/I_n)$ 的对象，而
$f_n : \xi_{n + 1} \to \xi_n$ 是位于投影
$R/I_{n + 1} \to R/I_n$ 上方的态射。
\item 态射
$a : (\xi_n, f_n) \to (\eta_n, g_n)$ 由一族
$\mathcal{F}(R/I_n)$ 中的态射
$a_n : \xi_n \to \eta_n$ 组成，并且对每个 $n$，图
$$
\xymatrix{
\xi_{n + 1} \ar[r]^{f_n} \ar[d]_{a_{n + 1}} & \xi_n \ar[d]^{a_n} \\
\eta_{n + 1} \ar[r]^{g_n} & \eta_n
}
$$
交换。
\end{enumerate}

\begin{lemma}
\label{formal-defos-lemma-formal-objects-different-filtration}
在上述情形下，$\widehat{\mathcal{F}}_\mathcal{I}(R)$ 等价于范畴
$\widehat{\mathcal{F}}(R)$。
\end{lemma}

\begin{proof}
可如下定义一个等价
$\widehat{\mathcal{F}}_\mathcal{I}(R) \to
\widehat{\mathcal{F}}(R)$。对每个 $n$，令 $m(n)$ 为满足
$I_m \subset \mathfrak m_R^n$ 的最小 $m$。
给定 $\widehat{\mathcal{F}}_\mathcal{I}(R)$ 的对象
$(\xi_n, f_n)$，令 $\eta_n$ 为 $\xi_{m(n)}$ 沿
$R/I_{m(n)} \to R/\mathfrak m_R^n$ 的推前。
令 $g_n : \eta_{n + 1} \to \eta_n$ 为
$\mathcal{F}$ 中位于
$R/\mathfrak m_R^{n + 1} \to R/\mathfrak m_R^n$ 上方的唯一态射，
使图
$$
\xymatrix{
\xi_{m(n + 1)} \ar[rrr]_{f_{m(n)} \circ \ldots \circ f_{m(n + 1) - 1}} \ar[d]
& & & \xi_{m(n)} \ar[d] \\
\eta_{n + 1} \ar[rrr]^{g_n} & & & \eta_n
}
$$
交换（余纤维范畴的公理保证存在性和唯一性）。
函子
$\widehat{\mathcal{F}}_\mathcal{I}(R) \to
\widehat{\mathcal{F}}(R)$ 把 $(\xi_n, f_n)$ 映到
$(R, \eta_n, g_n)$。验证这确为范畴等价的步骤从略。
\end{proof}

\begin{remark}
\label{formal-defos-remark-different-sequence-ideals}
设 $p : \mathcal{F} \to \mathcal{C}_\Lambda$ 是群胚余纤维范畴。
假设对每个 $R \in \Ob(\widehat{\mathcal{C}}_\Lambda)$，
都给定了由理想构成的 $R$ 的滤过 $\mathcal{I}_R$。
若 $\mathcal{I}_R$ 对所有 $R$ 都在 $R$ 上诱导
$\mathfrak m_R$-进拓扑，则可模仿
$\widehat{\mathcal{F}}$ 的定义来定义范畴
$\widehat{\mathcal{F}}_\mathcal{I}$。
该范畴配备态射
$\widehat{p}_\mathcal{I} : \widehat{\mathcal{F}}_\mathcal{I} \to
\widehat{\mathcal{C}}_\Lambda$，从而成为群胚余纤维范畴，
并且 $\widehat{\mathcal{F}}_\mathcal{I}(R)$ 与上面定义的
$\widehat{\mathcal{F}}_{\mathcal{I}_R}(R)$ 同构。
在对象 $R \in \Ob(\widehat{\mathcal{C}}_\Lambda)$ 上使用引理
\ref{formal-defos-lemma-formal-objects-different-filtration} 的等价，
可知群胚余纤维范畴
$\widehat{\mathcal{F}}_\mathcal{I}$ 和
$\widehat{\mathcal{F}}$ 等价。
\end{remark}

\begin{remark}
\label{formal-defos-remark-completion-functor}
设 $F: \mathcal{C}_\Lambda \to \textit{Sets}$ 是函子。
把函子与集合余纤维范畴识别，则 $F$ 的完备化是函子
$\widehat{F} : \widehat{\mathcal{C}}_\Lambda \to \textit{Sets}$，
由 $\widehat{F}(S) = \lim F(S/\mathfrak{m}_S^{n})$ 给出。
这与 Schlessinger 论文 \cite{Sch} 中的定义一致。
\end{remark}

\begin{remark}
\label{formal-defos-remark-restrict-completion}
设 $\mathcal{F}$ 是 $\mathcal{C}_\Lambda$ 上的群胚余纤维范畴。
我们断言有典范等价
$$
can :
\widehat{\mathcal{F}}|_{\mathcal{C}_\Lambda}
\longrightarrow
\mathcal{F}.
$$
事实上，设 $A \in \Ob(\mathcal{C}_\Lambda)$，且
$(A, \xi_n, f_n)$ 是
$\widehat{\mathcal{F}}|_{\mathcal{C}_\Lambda}(A)$ 的对象。
由于 $A$ 是阿廷环，存在满足
$\mathfrak m_A^m = 0$ 的最小 $m \in \mathbf{N}$。
此时 $can$ 把 $(A, \xi_n, f_n)$ 映到 $\xi_m$。
由《范畴》引理
\ref{categories-lemma-equivalence-fibred-categories}，
这个函子是群胚余纤维范畴之间的等价；这是因为由引理
\ref{formal-defos-lemma-formal-objects-different-filtration}，它在每个纤维范畴上
都是等价，而且局部阿廷环 $A$ 上的
$\mathfrak m_A$-进拓扑来自零理想。
我们将经常通过函子 $can$ 的一个拟逆，把 $\mathcal{F}$
与 $\widehat{\mathcal{F}}$ 的一个满子范畴识别。
\end{remark}

\begin{remark}
\label{formal-defos-remark-completion-morphism}
设 $\varphi : \mathcal{F} \to \mathcal{G}$ 是
$\mathcal{C}_\Lambda$ 上群胚余纤维范畴之间的态射。
则诱导出 $\widehat{\mathcal{C}}_\Lambda$ 上群胚余纤维范畴之间的态射
$\widehat{\varphi}: \widehat{\mathcal{F}} \to
\widehat{\mathcal{G}}$。
它把 $\widehat{\mathcal{F}}$ 的对象
$\xi = (R, \xi_n, f_n)$ 映到
$(R, \varphi(\xi_n), \varphi(f_n))$，并把
$\widehat{\mathcal{F}}$ 的对象 $\xi$ 与 $\eta$ 之间的态射
$(a_0 : R \to S, a_n : \xi_n \to \eta_n)$ 映到
$(a_0 : R \to S, \varphi(a_n) : \varphi(\xi_n) \to \varphi(\eta_n))$。
最后，若 $t : \varphi \to \varphi'$ 是群胚余纤维范畴的
$1$-态射 $\varphi, \varphi': \mathcal{F} \to \mathcal{G}$
之间的 $2$-态射，则得到 $2$-态射
$\widehat{t} : \widehat{\varphi} \to \widehat{\varphi}'$。
具体地，对上述 $\xi = (R, \xi_n, f_n)$，令
$\widehat{t}_\xi = (t_{\varphi(\xi_n)})$。% P10-E0433
因此，完备化定义了 $2$-范畴之间的函子
$$
\widehat{~} :
\text{Cof}(\mathcal{C}_\Lambda)
\longrightarrow
\text{Cof}(\widehat{\mathcal{C}}_\Lambda)
$$
它从 $\mathcal{C}_\Lambda$ 上群胚余纤维范畴的 $2$-范畴
到 $\widehat{\mathcal{C}}_\Lambda$ 上群胚余纤维范畴的
$2$-范畴。
\end{remark}

\begin{remark}
\label{formal-defos-remark-completion-restriction-adjoint}
我们断言，说明 \ref{formal-defos-remark-completion-morphism} 的完备化函子与
说明 \ref{formal-defos-remarks-cofibered-groupoids}
(\ref{formal-defos-item-definition-restricting-base-category}) 的限制函子
$|_{\mathcal{C}_\Lambda} : \text{Cof}(\widehat{\mathcal{C}}_\Lambda)
\to \text{Cof}(\mathcal{C}_\Lambda)$
在如下精确意义下“2-伴随”。设
$\mathcal{F} \in \Ob(\text{Cof}(\mathcal{C}_\Lambda))$ 且
$\mathcal{G} \in \Ob(\text{Cof}(\widehat{\mathcal{C}}_\Lambda))$。
则有范畴等价
$$
\Phi :
\Mor_{\mathcal{C}_\Lambda}(
\mathcal{G}|_{\mathcal{C}_\Lambda}, \mathcal{F})
\longrightarrow
\Mor_{\widehat{\mathcal{C}}_\Lambda}(\mathcal{G}, \widehat{\mathcal{F}})
$$
为描述这一等价，我们如下定义典范态射
$\mathcal{G} \to
\widehat{\mathcal{G}|_{\mathcal{C}_\Lambda}}$ 和
$\widehat{\mathcal{F}}|_{\mathcal{C}_\Lambda} \to \mathcal{F}$：
\begin{enumerate}
\item 设 $R \in \Ob(\widehat{\mathcal{C}}_\Lambda))$，并设 $\xi$
为纤维范畴 $\mathcal{G}(R)$ 的对象。% P10-E0434
对每个 $n \in \mathbf{N}$，选取 $\xi$ 到
$R/\mathfrak m_R^n$ 的推前 $\xi \to \xi_n$，并令
$f_n : \xi_{n + 1} \to \xi_n$ 为所诱导的态射。
于是 $\mathcal{G} \to
\widehat{\mathcal{G}|_{\mathcal{C}_\Lambda}}$
把 $\xi$ 映到 $(R, \xi_n, f_n)$。
\item 这是说明 \ref{formal-defos-remark-restrict-completion} 的等价
$can : \widehat{\mathcal{F}}|_{\mathcal{C}_\Lambda} \to \mathcal{F}$。
\end{enumerate}
说明这些之后，等价
$\Phi : \Mor_{\mathcal{C}_\Lambda}(
\mathcal{G}|_{\mathcal{C}_\Lambda}, \mathcal{F})  \to
\Mor_{\widehat{\mathcal{C}}_\Lambda}(\mathcal{G},
\widehat{\mathcal{F}})$
把态射
$\varphi : \mathcal{G}|_{\mathcal{C}_\Lambda} \to \mathcal{F}$
映到
$$
\mathcal{G} \to \widehat{\mathcal{G}|_{\mathcal{C}_\Lambda}}
\xrightarrow{\widehat{\varphi}} \widehat{\mathcal{F}}
$$
$\Phi$ 有拟逆
$\Psi :
\Mor_{\widehat{\mathcal{C}}_\Lambda}(
\mathcal{G}, \widehat{\mathcal{F}}) \to
\Mor_{\mathcal{C}_\Lambda}(
\mathcal{G}|_{\mathcal{C}_\Lambda}, \mathcal{F})$，
它把 $\psi : \mathcal{G} \to \widehat{\mathcal{F}}$ 映到
$$
\mathcal{G}|_{\mathcal{C}_\Lambda}
\xrightarrow{\psi|_{\mathcal{C}_\Lambda}}
\widehat{\mathcal{F}}|_{\mathcal{C}_\Lambda} \to \mathcal{F}.
$$
验证 $\Phi$ 与 $\Psi$ 互为拟逆的步骤从略。
我们也不讨论 $\Phi$ 的函子性（因为这会进入我们无论如何都想避免的
3-范畴领域）。
\end{remark}

\begin{remark}
\label{formal-defos-remark-completion-restriction-cofset-adjoint}
对范畴 $\mathcal{C}$，以 $\text{CofSet}(\mathcal{C})$
表示 $\mathcal{C}$ 上集合余纤维范畴所成的范畴。
这是一个 $1$-范畴，与函子范畴
$\mathcal{C} \to \textit{Sets}$ 同构。% P10-E0435
见说明 \ref{formal-defos-remarks-cofibered-groupoids}
(\ref{formal-defos-item-convention-cofibered-sets})。
完备化函子和限制函子分别限制为函子
$\widehat{~} : \text{CofSet}(\mathcal{C}_\Lambda) \to
\text{CofSet}(\widehat{\mathcal{C}}_\Lambda)$ 和
$|_{\mathcal{C}_\Lambda} :
\text{CofSet}(\widehat{\mathcal{C}}_\Lambda) \to
\text{CofSet}(\mathcal{C}_\Lambda)$，
仍以相同符号表示。作为集合余纤维范畴之范畴上的函子，
完备化与限制在通常的 1-范畴意义下伴随：说明
\ref{formal-defos-remark-completion-restriction-adjoint} 中的同一构造，
对
$F \in \Ob(\text{CofSet}(\mathcal{C}_\Lambda))$ 和
$G \in \Ob(\text{CofSet}(\widehat{\mathcal{C}}_\Lambda))$
定义函子性双射
$$
\Mor_{\mathcal{C}_\Lambda}(G|_{\mathcal{C}_\Lambda}, F)
\longrightarrow
\Mor_{\widehat{\mathcal{C}}_\Lambda}(G, \widehat{F})
$$
同样，映射 $\widehat{F}|_{\mathcal{C}_\Lambda} \to F$
是同构。
\end{remark}

\begin{remark}
\label{formal-defos-remark-restrict-complete-continuous-functor}
设 $G : \widehat{\mathcal{C}}_\Lambda \to \textit{Sets}$
是与极限交换的函子。则说明
\ref{formal-defos-remark-completion-restriction-adjoint} 中所述的映射
$G \to \widehat{G|_{\mathcal{C}_\Lambda}}$ 是同构。
事实上，若 $S$ 是 $\widehat{\mathcal{C}}_\Lambda$ 的对象，
则有典范双射
$$
\widehat{G|_{\mathcal{C}_\Lambda}}(S) =
\lim_n G(S/\mathfrak{m}_S^n) =
G(\lim_n S/\mathfrak{m}_S^n) = G(S).
$$
特别地，若 $R$ 是 $\widehat{\mathcal{C}}_\Lambda$ 的对象，则
$\underline{R} =
\widehat{\underline{R}|_{\mathcal{C}_\Lambda}}$，
因为根据极限的定义，可表函子 $\underline{R}$ 与极限交换。
\end{remark}

\begin{remark}
\label{formal-defos-remark-formal-objects-yoneda}
设 $R$ 是 $\widehat{\mathcal{C}}_\Lambda$ 的对象。
它按说明 \ref{formal-defos-remarks-cofibered-groupoids}
(\ref{formal-defos-item-definition-yoneda}) 定义函子
$\underline{R}: \widehat{\mathcal{C}}_\Lambda \to \textit{Sets}$。
一如往常，我们把这个函子与相伴的集合余纤维范畴识别。
若 $\mathcal{F}$ 是 $\mathcal{C}_\Lambda$ 上的余纤维范畴，
则有范畴等价
\begin{equation}
\label{formal-defos-equation-formal-objects-maps}
\Mor_{\mathcal{C}_\Lambda}(
\underline{R}|_{\mathcal{C}_\Lambda}, \mathcal{F})
\longrightarrow
\widehat{\mathcal{F}}(R).
\end{equation}
它由复合
$$
\Mor_{\mathcal{C}_\Lambda}(
\underline{R}|_{\mathcal{C}_\Lambda}, \mathcal{F})
\xrightarrow{\Phi}
\Mor_{\widehat{\mathcal{C}}_\Lambda}(
\underline{R}, \widehat{\mathcal{F}})
\xrightarrow{\sim}
\widehat{\mathcal{F}}(R)
$$
给出，其中 $\Phi$ 如说明
\ref{formal-defos-remark-completion-restriction-adjoint} 所述，
而第二个等价来自 2-Yoneda 引理（即《范畴》引理
\ref{categories-lemma-yoneda-2category} 的余纤维版本）。
明确地说，这个等价把态射
$\varphi : \underline{R}|_{\mathcal{C}_\Lambda} \to \mathcal{F}$
映到 $\widehat{\mathcal{F}}(R)$ 中的形式对象
$(R, \varphi(R \to R/\mathfrak{m}_R^n), \varphi(f_n))$，
其中 $f_n : R/\mathfrak m_R^{n + 1} \to
R/\mathfrak m_R^n$ 是投影。

\medskip\noindent
假设已经为 $\mathcal{F}$ 选取了推前。
给定任意 $\xi \in \Ob(\widehat{\mathcal{F}}(R))$，
我们构造一个明确的态射
$\underline{\xi} :
\underline{R}|_{\mathcal{C}_\Lambda} \to \mathcal{F}$，
它在 (\ref{formal-defos-equation-formal-objects-maps}) 下映到 $\xi$。
具体地，设 $\xi = (R, \xi_n, f_n)$。
$\underline{R}|_{\mathcal{C}_\Lambda}$ 的对象 $\alpha$
等同于 $\widehat{\mathcal{C}}_\Lambda$ 中以阿廷环
$A$ 为靶的态射 $\alpha : R \to A$。
令 $m \in \mathbf{N}$ 为满足 $\mathfrak m_A^m = 0$ 的最小整数。
于是 $\alpha$ 唯一地经由
$\alpha_m : R/\mathfrak m_R^m \to A$ 分解，并可令
$\underline{\xi}(\alpha) = \alpha_{m, *}\xi_m$。
我们略去 $\underline{\xi}$ 在态射上的描述，也略去证明
$\underline{\xi}$ 经由
(\ref{formal-defos-equation-formal-objects-maps}) 映到 $\xi$ 的步骤。

\medskip\noindent
假设已经为 $\widehat{\mathcal{F}}$ 选取了推前。
在这种情况下，《范畴》引理
\ref{categories-lemma-yoneda-2category} 的证明给出 2-Yoneda
等价的一个明确拟逆
$$
\iota :
\widehat{\mathcal{F}}(R) \longrightarrow
\Mor_{\widehat{\mathcal{C}}_\Lambda}(
\underline{R}, \widehat{\mathcal{F}})
$$
它把 $\xi$ 映到态射
$\iota(\xi) : \underline{R} \to \widehat{\mathcal{F}}$；
该态射把
$f \in \underline{R}(S) =
\Mor_{\mathcal{C}_\Lambda}(R, S)$ 映到 $f_*\xi$。% P10-E0436
于是 (\ref{formal-defos-equation-formal-objects-maps}) 的一个拟逆为
$$
\widehat{\mathcal{F}}(R)
\xrightarrow{\iota}
\Mor_{\widehat{\mathcal{C}}_\Lambda}(
\underline{R}, \widehat{\mathcal{F}})
\xrightarrow{\Psi}
\Mor_{\mathcal{C}_\Lambda}(
\underline{R}|_{\mathcal{C}_\Lambda}, \mathcal{F})
$$
其中 $\Psi$ 如说明
\ref{formal-defos-remark-completion-restriction-adjoint} 所述。
给定 $\xi \in \Ob(\widehat{\mathcal{F}}(R))$，有
$\Psi(\iota(\xi)) \cong \underline{\xi}$，
其中 $\underline{\xi}$ 如上一段所述；这是因为在范畴等价
(\ref{formal-defos-equation-formal-objects-maps}) 下，两者都映到 $\xi$。
使用
$\underline{R} =
\widehat{\underline{R}|_{\mathcal{C}_\Lambda}}$
（见说明 \ref{formal-defos-remark-restrict-complete-continuous-functor}），
并展开 $\Phi$ 和 $\Psi$ 的定义，可得
$\iota(\xi)$ 与 $\underline{\xi}$ 的完备化同构。
\end{remark}

\begin{remark}
\label{formal-defos-remark-formal-objects-yoneda-map}
设 $\mathcal{F}$ 是 $\mathcal{C}_\Lambda$ 上的群胚余纤维范畴。
设 $\xi = (R, \xi_n, f_n)$ 和
$\eta = (S, \eta_n, g_n)$ 是 $\mathcal{F}$ 的形式对象。
令 $a = (a_n) : \xi \to \eta$ 为形式对象之间的态射，
即 $\widehat{\mathcal{F}}$ 的态射。
令 $f = \widehat{p}(a) = a_0 : R \to S$ 为 $a$ 在
$\widehat{\mathcal{C}}_\Lambda$ 中的投影。于是得到 $2$-交换图
$$
\xymatrix{
\underline{R}|_{\mathcal{C}_\Lambda} \ar[rd]_{\underline{\xi}} & &
\underline{S}|_{\mathcal{C}_\Lambda} \ar[ll]^f \ar[ld]^{\underline{\eta}} \\
& \mathcal{F}
}
$$
其中 $\underline{\xi}$ 和 $\underline{\eta}$ 是说明
\ref{formal-defos-remark-formal-objects-yoneda} 中构造的态射。
为说明这一点，设 $\alpha : S \to A$ 是
$\underline{S}|_{\mathcal{C}_\Lambda}$ 的对象（见上引文献）。
令 $m \in \mathbf{N}$ 为满足 $\mathfrak m_A^m = 0$ 的最小整数。
得到交换图
$$
\xymatrix{
R \ar[d]^f \ar[r] & R/\mathfrak m_R^m \ar[d]_{f_m} \ar[rd]^{\beta_m} \\
S \ar[r] & S/\mathfrak m_S^m \ar[r]^{\alpha_m} & A
}
$$
其底行箭头的复合为 $\alpha$。此时
$\underline{\eta}(\alpha) = \alpha_{m, *}\eta_m$ 且
$\underline{\xi}(\alpha \circ f) = \beta_{m, *}\xi_m$。
态射 $a_m : \xi_m \to \eta_m$ 位于 $f_m$ 上方，
故得到位于 $\text{id}_A$ 上方的典范态射
$$
\underline{\xi}(\alpha \circ f) = \beta_{m, *}\xi_m
\longrightarrow
\underline{\eta}(\alpha) = \alpha_{m, *}\eta_m
$$
使图
$$
\xymatrix{
\xi_m \ar[r] \ar[d]^{a_m} &  \beta_{m, *}\xi_m \ar[d] \\
\eta_m \ar[r] &  \alpha_{m, *}\eta_m
}
$$
由群胚余纤维范畴的公理而交换。这定义了函子变换
$\underline{\xi} \circ f \to \underline{\eta}$，
它见证了本说明第一个图的 2-交换性。
\end{remark}

\begin{remark}
\label{formal-defos-remark-spell-out-formal-object}
按照说明 \ref{formal-defos-remark-formal-objects-yoneda}，
给出 $\mathcal{F}$ 的形式对象 $\xi$，等价于给出一个
pro-可表函子 $U : \mathcal{C}_\Lambda \to \textit{Sets}$
以及一个态射 $U \to \mathcal{F}$。
\end{remark}

\section{光滑态射}
\label{formal-defos-section-smooth-morphisms}

\noindent
本节讨论 $\mathcal{C}_\Lambda$ 上群胚余纤维范畴之间的光滑态射。

\begin{definition}
\label{formal-defos-definition-smooth-morphism}
设 $\varphi : \mathcal{F} \to \mathcal{G}$ 是
$\mathcal{C}_\Lambda$ 上群胚余纤维范畴之间的态射。
若 $\varphi$ 满足下列条件，则称它{\it 光滑}：
设 $B \to A$ 是 $\mathcal{C}_\Lambda$ 中的满环映射，
$y \in \Ob(\mathcal{G}(B)), x \in
\Ob(\mathcal{F}(A))$，且
$y \to \varphi(x)$ 是位于 $B \to A$ 上方的态射。
则存在 $x' \in \Ob(\mathcal{F}(B))$、位于
$B \to A$ 上方的态射 $x' \to x$，以及位于
$\text{id}: B \to B$ 上方的态射 $\varphi(x') \to y$，
使图
$$
\xymatrix{
\varphi(x') \ar[r] \ar[dr] & y \ar[d] \\
& \varphi(x)
}
$$
交换。
\end{definition}

\begin{lemma}
\label{formal-defos-lemma-smoothness-small-extensions}
设 $\varphi : \mathcal{F} \to \mathcal{G}$ 是
$\mathcal{C}_\Lambda$ 上群胚余纤维范畴之间的态射。
若只对小扩张 $B \to A$ 假设定义
\ref{formal-defos-definition-smooth-morphism} 中的条件成立，
则 $\varphi$ 光滑。
\end{lemma}

\begin{proof}
设 $B \to A$ 是 $\mathcal{C}_\Lambda$ 中的满环映射。
设 $y \in \Ob(\mathcal{G}(B))$、
$x \in \Ob(\mathcal{F}(A))$，且
$y \to \varphi(x)$ 是位于 $B \to A$ 上方的态射。
由引理 \ref{formal-defos-lemma-factor-small-extension}，可把
$B \to A$ 分解为小扩张
$B = B_n \to B_{n-1} \to \ldots \to B_0 = A$。
对 $n$ 作归纳。若 $n = 1$，结论由假设成立。
若 $n > 1$，记
$f : B = B_n \to B_{n - 1}$ 及
$g : B_{n - 1} \to B_0 = A$。
选取 $y$ 沿 $f$ 的推前 $y \to f_* y$，
使态射 $y \to \varphi(x)$ 分解为
$y \to f_* y \to \varphi(x)$。
由归纳假设，可以找到位于
$g : B_{n - 1} \to A$ 上方的 $x_{n - 1} \to x$，
以及位于
$\text{id} : B_{n - 1} \to B_{n - 1}$ 上方的
$a : \varphi(x_{n - 1}) \to f_*y$，使图
$$
\xymatrix{
\varphi(x_{n - 1}) \ar[r]_-a \ar[dr] & f_*y \ar[d] \\
& \varphi(x)
}
$$
交换。把假设应用于 $y \to f_*y$ 与
$a^{-1} : f_*y \to \varphi(x_{n - 1})$ 的复合
$y \to \varphi(x_{n - 1})$，可得位于
$B_n \to B_{n - 1}$ 上方的
$x_n \to x_{n - 1}$，以及位于
$\text{id} : B_n \to B_n$ 上方的
$\varphi(x_n) \to y$，使图
$$
\xymatrix{
\varphi(x_n) \ar[r] \ar[d] & y \ar[d] \\
\varphi(x_{n - 1}) \ar[r]^-a \ar[dr] & f_*y \ar[d] \\
& \varphi(x)
}
$$
交换。于是复合 $x_n \to x_{n - 1} \to x$ 与
$\varphi(x_n) \to y$ 就是光滑性定义所需的态射。
\end{proof}

\begin{remark}
\label{formal-defos-remark-smoothness-2-categorical}
设 $\varphi : \mathcal{F} \to \mathcal{G}$ 是
$\mathcal{C}_\Lambda$ 上群胚余纤维范畴之间的态射，
且 $B \to A$ 是 $\mathcal{C}_\Lambda$ 中的环映射。
为纤维范畴 $\mathcal{F}(B)$ 和 $\mathcal{G}(B)$ 中的对象
选取沿 $B \to A$ 的推前，就确定了函子
$\mathcal{F}(B) \to \mathcal{F}(A)$ 和
$\mathcal{G}(B) \to \mathcal{G}(A)$，它们组成 $2$-交换图
$$
\xymatrix{
\mathcal{F}(B) \ar[r]^{\varphi} \ar[d] & \mathcal{G}(B) \ar[d] \\
\mathcal{F}(A) \ar[r]^{\varphi}        & \mathcal{G}(A) .
}
$$
因此诱导出函子
$\mathcal{F}(B) \to \mathcal{F}(A)
\times_{\mathcal{G}(A)} \mathcal{G}(B)$。
展开定义可知，$\varphi : \mathcal{F} \to \mathcal{G}$ 光滑，
当且仅当每当 $B \to A$ 为满射时，这个诱导函子本质满
（或者等价地，由引理
\ref{formal-defos-lemma-smoothness-small-extensions}，每当
$B \to A$ 为小扩张时它本质满）。
\end{remark}

\begin{remark}
\label{formal-defos-remark-compare-smooth-schlessinger}
说明 \ref{formal-defos-remark-smoothness-2-categorical} 中对光滑态射的刻画，
类似于 Schlessinger 的函子光滑态射概念；比较
\cite[Definition 2.2.]{Sch}。事实上，当
$\mathcal{F}$ 和 $\mathcal{G}$ 是集合余纤维范畴时，
我们的概念与 Schlessinger 的概念等价。
具体地，在这种情况下，设
$F, G : \mathcal{C}_\Lambda \to \textit{Sets}$ 为对应函子；
见说明 \ref{formal-defos-remarks-cofibered-groupoids}
(\ref{formal-defos-item-convention-cofibered-sets})。
则 $F \to G$ 光滑，当且仅当对
$\mathcal{C}_\Lambda$ 中的每个满环映射
$B \to A$，映射
$F(B) \to F(A) \times_{G(A)} G(B)$ 都满射。
\end{remark}

\begin{remark}
\label{formal-defos-remark-smooth-to-iso-classes}
设 $\mathcal{F}$ 是 $\mathcal{C}_\Lambda$ 上的群胚余纤维范畴。
则态射 $\mathcal{F} \to \overline{\mathcal{F}}$ 光滑。
事实上，设 $f : B \to A$ 是
$\mathcal{C}_\Lambda$ 中的环映射。
设 $x \in \Ob(\mathcal{F}(A))$，并设
$\overline{y} \in \overline{\mathcal{F}}(B)$
是某个 $y \in \Ob(\mathcal{F}(B))$ 的同构类，且满足
$\overline{f_*y} = \overline{x}$。
那么只需取 $x' = y$、隐含的位于 $B \to A$ 上方的态射
$x' = y \to x$，以及等式
$\overline{x'} = \overline{y}$，就解出了定义
\ref{formal-defos-definition-smooth-morphism} 中的问题。
\end{remark}

\noindent
若 $R \to S$ 是环映射
$\widehat{\mathcal{C}}_\Lambda$，则在函子
$\underline{S}, \underline{R}: \widehat{\mathcal{C}}_\Lambda
\to \textit{Sets}$ 之间诱导出态射
$\underline{S} \to \underline{R}$。% P10-E0437
在这种情况下，限制
$\underline{S}|_{\mathcal{C}_\Lambda} \to
\underline{R}|_{\mathcal{C}_\Lambda}$ 的光滑性是熟悉的概念：

\begin{lemma}
\label{formal-defos-lemma-smooth-morphism-power-series}
设 $R \to S$ 是 $\widehat{\mathcal{C}}_\Lambda$ 中的环映射。
则诱导态射
$\underline{S}|_{\mathcal{C}_\Lambda} \to
\underline{R}|_{\mathcal{C}_\Lambda}$
光滑，当且仅当 $S$ 是 $R$ 上的幂级数环。
\end{lemma}

\begin{proof}
先设 $S$ 是 $R$ 上的幂级数环，记
$S = R[[x_1, \ldots, x_n]]$。
$\underline{S}|_{\mathcal{C}_\Lambda} \to
\underline{R}|_{\mathcal{C}_\Lambda}$ 的光滑性意指如下性质
（见说明 \ref{formal-defos-remark-compare-smooth-schlessinger}）：
给定 $\mathcal{C}_\Lambda$ 中的满环映射 $B \to A$、
环映射 $R \to B$，以及使实线图
$$
\xymatrix{
S \ar[r] \ar@{..>}[rd] & A \\
R \ar[u] \ar[r] & B \ar[u]
}
$$
交换的环映射 $S \to A$，存在使该图交换的虚线箭头。
（注意这与《代数》定义
\ref{algebra-definition-formally-smooth} 的相似性。）
为构造虚线箭头，选取元素 $b_i \in B$，使其在 $A$ 中的像
等于 $x_i$ 在 $A$ 中的像。注意
$b_i \in \mathfrak m_B$，因为 $x_i$ 映到
$\mathfrak m_A$ 的元素。因此存在唯一的 $R$-代数映射
$R[[x_1, \ldots, x_n]] \to B$，它把 $x_i$ 映到 $b_i$，
并可作为所需虚线箭头。

\medskip\noindent
反过来，假设
$\underline{S}|_{\mathcal{C}_\Lambda} \to
\underline{R}|_{\mathcal{C}_\Lambda}$ 光滑。
选取 $x_1, \ldots, x_n \in S$，使其像构成 $S$ 在 $R$ 上的
相对余切空间
$\mathfrak m_S/(\mathfrak m_R S + \mathfrak m_S^2)$ 的一组基。
令 $T = R[[X_1, \ldots, X_n]]$。注意
$$
S/(\mathfrak m_R S + \mathfrak m_S^2) \cong
R/\mathfrak m_R[x_1, \ldots, x_n]/(x_ix_j)
$$
以及
$$
T/(\mathfrak m_R T + \mathfrak m_T^2) \cong
R/\mathfrak m_R[X_1, \ldots, X_n]/(X_iX_j).
$$
令
$S/(\mathfrak m_R S + \mathfrak m_S^2) \to
T/(\mathfrak m_R T + \mathfrak m_T^2)$
为把 $x_i$ 的类映到 $X_i$ 的类所定义的局部 $R$-代数同构。
令 $f_1 : S \to T/(\mathfrak m_R T + \mathfrak m_T^2)$
为复合
$S \to S/(\mathfrak m_R S + \mathfrak m_S^2)
\to T/(\mathfrak m_R T + \mathfrak m_T^2)$。
假设
$\underline{S}|_{\mathcal{C}_\Lambda} \to
\underline{R}|_{\mathcal{C}_\Lambda}$
光滑，意即可以把 $f_1$ 提升为映射
$f_2 : S \to T/\mathfrak{m}_T^2$，再提升为
$f_3 : S \to T/\mathfrak{m}_T^3$，并如此继续，对所有
$n \geq 1$ 都成立。于是得到局部 $R$-代数的诱导映射
$f : S \to T = \lim T/\mathfrak m_T^n$。
由 $f_1$ 的选取，映射 $f$ 在相对余切空间上诱导同构
$\mathfrak m_S/(\mathfrak m_R S + \mathfrak m_S^2) \to
\mathfrak m_T/(\mathfrak m_R T + \mathfrak m_T^2)$。
因此由引理 \ref{formal-defos-lemma-surjective-cotangent-space}，
$f$ 为满射（此处把 $f$ 视为
$\widehat{\mathcal{C}}_R$ 中的映射）。
选取 $X_i \in T$ 在 $f$ 下的原像 $y_i \in S$。
由于 $T$ 是 $R$ 上的幂级数环，存在局部 $R$-代数同态
$s : T \to S$，把 $X_i$ 映到 $y_i$。
由构造 $f \circ s = \text{id}$，故 $s$ 是单射。
但由于 $f$ 在相对余切空间上诱导同构，$s$ 也如此；
再次由引理 \ref{formal-defos-lemma-surjective-cotangent-space}，
它还是满射。因此 $s$ 和 $f$ 都是同构。
\end{proof}

\noindent
光滑态射具有如下函子性质。

\begin{lemma}
\label{formal-defos-lemma-smooth-properties}
设 $\varphi : \mathcal{F} \to \mathcal{G}$ 和
$\psi : \mathcal{G} \to \mathcal{H}$ 是
$\mathcal{C}_\Lambda$ 上群胚余纤维范畴之间的态射。
\begin{enumerate}
\item 若 $\varphi$ 和 $\psi$ 光滑，则
$\psi \circ \varphi$ 光滑。
\item 若 $\varphi$ 本质满且
$\psi \circ \varphi$ 光滑，则 $\psi$ 光滑。
\item 若 $\mathcal{G}' \to \mathcal{G}$ 是群胚余纤维范畴之间的态射，
且 $\varphi$ 光滑，则
$\mathcal{F} \times_\mathcal{G} \mathcal{G}' \to
\mathcal{G}'$ 光滑。
\end{enumerate}
\end{lemma}

\begin{proof}
(1) 和 (2) 直接由定义得出。(3) 的证明从略。提示：使用说明
\ref{formal-defos-remark-smoothness-2-categorical} 给出的光滑性表述，
并使用 $\mathcal{F} \times_\mathcal{G} \mathcal{G}'$
是 $2$-纤维积这一事实；见说明
\ref{formal-defos-remarks-cofibered-groupoids} (\ref{formal-defos-item-fibre-product})。
\end{proof}

\begin{lemma}
\label{formal-defos-lemma-smooth-morphism-essentially-surjective}
设 $\varphi : \mathcal{F} \to \mathcal{G}$ 是
$\mathcal{C}_\Lambda$ 上群胚余纤维范畴之间的光滑态射。
假设 $\varphi : \mathcal{F}(k) \to \mathcal{G}(k)$ 本质满。
则 $\varphi : \mathcal{F} \to \mathcal{G}$ 和
$\widehat{\varphi} : \widehat{\mathcal{F}} \to
\widehat{\mathcal{G}}$ 都本质满。
\end{lemma}

\begin{proof}
设 $y$ 是 $\mathcal{G}$ 中位于
$A \in \Ob(\mathcal{C}_\Lambda)$ 上方的对象。
设 $y \to y_0$ 是 $y$ 沿 $A \to k$ 的一个推前。
由 $\varphi : \mathcal{F}(k) \to \mathcal{G}(k)$
本质满这一假设，存在 $\mathcal{F}$ 中位于 $k$ 上方的对象
$x_0$ 以及同构 $y_0 \to \varphi(x_0)$。
$\varphi$ 的光滑性蕴含：存在 $\mathcal{F}$ 中位于
$A$ 上方的对象 $x$，使其像 $\varphi(x)$ 与 $y$ 同构。
因此 $\varphi : \mathcal{F} \to \mathcal{G}$ 本质满。

\medskip\noindent
设 $\eta = (R, \eta_n, g_n)$ 是
$\widehat{\mathcal{G}}$ 的对象。
我们构造 $\widehat{\mathcal{F}}$ 的对象 $\xi$ 以及同构
$\eta \to \varphi(\xi)$。由
$\varphi : \mathcal{F}(k) \to \mathcal{G}(k)$ 本质满这一假设，
对某个 $\xi_1 \in \Ob(\mathcal{F}(k))$，存在
$\mathcal{G}(k)$ 中的态射
$\eta_1 \to \varphi(\xi_1)$。态射
$\eta_2 \xrightarrow{g_1} \eta_1 \to \varphi(\xi_1)$
位于满环映射 $R/\mathfrak m_R^2 \to k$ 上方，
故由 $\varphi$ 的光滑性，存在
$\xi_2 \in \Ob(\mathcal{F}(R/\mathfrak m_R^2))$、
位于 $R/\mathfrak m_R^2 \to k$ 上方的态射
$f_1: \xi_2 \to \xi_1$，以及态射
$\eta_2 \to \varphi(\xi_2)$，使图
$$
\xymatrix{
\varphi(\xi_2)  \ar[r]^{\varphi(f_1)} &  \varphi(\xi_{1})   \\
\eta_2   \ar[u] \ar[r]^{g_1}  & \eta_1  \ar[u] \\
}
$$
交换。继续这一过程，就构造出
$\widehat{\mathcal{F}}$ 的对象
$\xi = (R, \xi_n, f_n)$，以及
$\widehat{\mathcal{G}}(R)$ 中的态射
$\eta \to \varphi(\xi) =
(R, \varphi(\xi_n), \varphi(f_n))$。
\end{proof}

\noindent
稍后我们要从 pro-可表函子构造到预形变范畴
$\mathcal{F}$ 的光滑态射。由说明
\ref{formal-defos-remark-formal-objects-yoneda} 中的讨论，
这些态射对应于 $\mathcal{F}$ 的某些形式对象。
更确切地说，它们就是所谓 $\mathcal{F}$ 的泛形式对象。

\begin{definition}
\label{formal-defos-definition-versal}
设 $\mathcal{F}$ 是群胚余纤维范畴，$\xi$ 是
$\mathcal{F}$ 的形式对象，位于
$R \in \Ob(\widehat{\mathcal{C}}_\Lambda)$ 上方。
若说明 \ref{formal-defos-remark-formal-objects-yoneda} 中的对应态射
$\underline{\xi}: \underline{R}|_{\mathcal{C}_\Lambda} \to
\mathcal{F}$ 光滑，则称 $\xi$ 是{\it 泛的}。
\end{definition}

\begin{remark}
\label{formal-defos-remark-versal-object}
设 $\mathcal{F}$ 是 $\mathcal C_\Lambda$ 上的群胚余纤维范畴，
且 $\xi$ 是 $\mathcal{F}$ 的形式对象。
由光滑性的定义，$\xi$ 的泛性等价于如下条件：
若
$$
\xymatrix{
& y \ar[d] \\
\xi \ar[r] & x
}
$$
是 $\widehat{\mathcal{F}}$ 中的图，且 $y \to x$
位于阿廷环之间的满射 $B \to A$ 上方
（可以假设它是小扩张），则存在态射 $\xi \to y$，使
$$
\xymatrix{
& y \ar[d] \\
\xi \ar[r] \ar[ur] & x
}
$$
交换。特别地，$\xi$ 为泛的这一条件不依赖于说明
\ref{formal-defos-remark-formal-objects-yoneda} 中构造
$\underline{\xi} :
\underline{R}|_{\mathcal{C}_\Lambda} \to \mathcal{F}$
时所作的推前选取。
\end{remark}

\begin{lemma}
\label{formal-defos-lemma-versal-object-quasi-initial}
设 $\mathcal{F}$ 是预形变范畴，且 $\xi$ 是
$\mathcal{F}$ 的泛形式对象。对
$\widehat{\mathcal{F}}$ 的任意形式对象 $\eta$，% P10-E0438
都存在态射 $\xi \to \eta$。
\end{lemma}

\begin{proof}
按假设，态射
$\underline{\xi} :
\underline{R}|_{\mathcal{C}_\Lambda} \to \mathcal{F}$
光滑。于是
$\iota(\xi) : \underline{R} \to \widehat{\mathcal{F}}$
是 $\underline{\xi}$ 的完备化；见说明
\ref{formal-defos-remark-formal-objects-yoneda}。
由引理 \ref{formal-defos-lemma-smooth-morphism-essentially-surjective}，
存在 $\underline{R}$ 的对象 $f$，使
$\iota(\xi)(f) = \eta$。此时 $f$ 是
$\widehat{\mathcal{C}}_\Lambda$ 中的环映射
$f : R \to S$。而 $\iota(\xi)(f) = \eta$ 意即
$f_*\xi \cong \eta$，这正是说存在位于 $f$ 上方的态射
$\xi \to \eta$。
\end{proof}

\section{光滑范畴或无阻碍范畴}
\label{formal-defos-section-smooth}

\noindent
设 $p : \mathcal{F} \to \mathcal{C}_\Lambda$ 是群胚余纤维范畴。
利用恒等函子，可以把 $\mathcal{C}_\Lambda$ 视为
$\mathcal{C}_\Lambda$ 上的群胚余纤维范畴。
这样，
$p : \mathcal{F} \longrightarrow \mathcal{C}_\Lambda$
就成为 $\mathcal{C}_\Lambda$ 上群胚余纤维范畴之间的态射。

\begin{definition}
\label{formal-defos-definition-cofibered-groupoid-projection-smooth}
设 $p : \mathcal{F} \to \mathcal{C}_\Lambda$ 是群胚余纤维范畴。
若其结构态射 $p$ 在定义 \ref{formal-defos-definition-smooth-morphism}
的意义下光滑，则称 $\mathcal{F}$ 是{\it 光滑的}或
{\it 无阻碍的}。
\end{definition}

\noindent
这是 $\mathcal{C}_\Lambda$ 上群胚余纤维范畴光滑性的“绝对”概念，
尽管更准确的说法是 $\mathcal{F}$ 在 $\Lambda$ 上光滑。
使用“{\it $\mathcal{F}$ 无阻碍}”这一说法时必须谨慎：
可能出现这种情况：$\mathcal{F}$ 有一个阻碍理论，
且该理论的阻碍空间非零，尽管 $\mathcal{F}$ 光滑。

\begin{remark}
\label{formal-defos-remark-smooth-on-top}
假设 $\mathcal{F}$ 是预形变范畴，并且有来自预形变范畴
$\mathcal{U}$ 的光滑态射
$\varphi : \mathcal U \to \mathcal{F}$。
则由引理 \ref{formal-defos-lemma-smooth-morphism-essentially-surjective}，
$\varphi$ 本质满；所以由引理 \ref{formal-defos-lemma-smooth-properties}，
$p: \mathcal{F} \to \mathcal{C}_\Lambda$ 光滑，
当且仅当复合
$\mathcal U \xrightarrow{\varphi} \mathcal{F} \xrightarrow{p}
\mathcal{C}_\Lambda$ 光滑；也就是说，$\mathcal{F}$ 光滑，
当且仅当 $\mathcal{U}$ 光滑。
\end{remark}

\begin{lemma}
\label{formal-defos-lemma-smooth}
设 $R \in \Ob(\widehat{\mathcal{C}}_\Lambda)$。下列条件等价：
\begin{enumerate}
\item $\underline{R}|_{\mathcal{C}_\Lambda}$ 光滑；
\item $\Lambda \to R$ 关于 $\mathfrak m_R$-进拓扑形式光滑；
\item $\Lambda \to R$ 平坦，并且
$R \otimes_\Lambda k'$ 在 $k'$ 上几何正则；
\item $\Lambda \to R$ 平坦，并且
$k' \to R \otimes_\Lambda k'$ 关于
$\mathfrak m_R$-进拓扑形式光滑。
\end{enumerate}
在经典情形下，这些条件还等价于
\begin{enumerate}
\item[(5)] 对某个 $n$，$R$ 同构于
$\Lambda[[x_1, \ldots, x_n]]$。
\end{enumerate}
\end{lemma}

\begin{proof}
态射
$p : \underline{R}|_{\mathcal{C}_\Lambda} \to
\mathcal{C}_\Lambda$
光滑，意即给定 $\mathcal{C}_\Lambda$ 中的满射
$B \to A$ 和映射 $R \to A$，可以在图
$$
\xymatrix{
R \ar[r] \ar@{-->}[rd] & A \\
\Lambda \ar[r] \ar[u] & B \ar[u]
}
$$
中找到虚线箭头。若 $\Lambda \to R$ 关于
$\mathfrak m_R$-进拓扑形式光滑，这当然成立；见
《代数进阶》定义
\ref{more-algebra-definition-formally-smooth-adic} 和
\ref{more-algebra-definition-formally-smooth}。
反过来，若上述性质成立，则由《代数进阶》引理
\ref{more-algebra-lemma-fs-local}，
$\Lambda \to R$ 关于 $\mathfrak m_R$-进拓扑形式光滑。
所以 (1) 与 (2) 等价。

\medskip\noindent
(2)、(3) 与 (4) 的等价性是《代数进阶》命题
\ref{more-algebra-proposition-fs-flat-fibre-fs}。
与 (5) 的等价性，例如可由引理
\ref{formal-defos-lemma-smooth-morphism-power-series} 以及经典情形下
$\mathcal{C}_\Lambda$ 与
$\underline{\Lambda}|_{\mathcal{C}_\Lambda}$ 相同这一事实得出。
\end{proof}

\begin{lemma}
\label{formal-defos-lemma-smooth-power-series-classical}
设 $\mathcal{F}$ 是预形变范畴，$\xi$ 是
$\mathcal{F}$ 的泛形式对象，位于
$R \in \Ob(\widehat{\mathcal{C}}_\Lambda)$ 上方。
下列条件等价：
\begin{enumerate}
\item $\mathcal{F}$ 无阻碍；
\item $\Lambda \to R$ 关于 $\mathfrak m_R$-进拓扑形式光滑。
\end{enumerate}
在经典情形下，这些条件还等价于
\begin{enumerate}
\item[(3)] 对某个 $n$，有
$R \cong \Lambda[[x_1, \ldots, x_n]]$。\ % P10-E0439
\end{enumerate}
\end{lemma}

\begin{proof}
若 (1) 成立，即 $\mathcal{F}$ 无阻碍，则复合
$$
\underline{R}|_{\mathcal{C}_\Lambda}
\xrightarrow{\underline{\xi}}
\mathcal{F}
\to
\mathcal{C}_\Lambda
$$
光滑；见引理 \ref{formal-defos-lemma-smooth-properties}。
所以由引理 \ref{formal-defos-lemma-smooth}，(2) 成立。
反过来，若 (2) 成立，则上述复合光滑，
而且由引理 \ref{formal-defos-lemma-versal-object-quasi-initial}，
第一个箭头本质满。因此由引理
\ref{formal-defos-lemma-smooth-properties}，第二个箭头光滑；
按定义，这意即 $\mathcal{F}$ 无阻碍。
经典情形下与 (3) 的等价性由引理 \ref{formal-defos-lemma-smooth} 得出。
\end{proof}

\begin{lemma}
\label{formal-defos-lemma-exists-smooth}
存在 $R \in \Ob(\widehat{\mathcal{C}}_\Lambda)$，
使引理 \ref{formal-defos-lemma-smooth} 的等价条件成立，并且还有
$H_1(L_{k/\Lambda}) = \mathfrak m_R/\mathfrak m_R^2$ 及
$\Omega_{R/\Lambda} \otimes_R k = \Omega_{k/\Lambda}$。
\end{lemma}

\begin{proof}
在经典情形下取 $R = \Lambda$。更一般地，若剩余域扩张
$k/k'$ 可分，则存在 $\Lambda$ 的完备化
$\Lambda^\wedge$ 的唯一有限平展扩张
$\Lambda^\wedge \to R$
（《代数》引理 \ref{algebra-lemma-complete-henselian} 和
\ref{algebra-lemma-henselian-cat-finite-etale}），
它在剩余域上诱导扩张 $k/k'$。

\medskip\noindent
一般情形下作如下处理。选取光滑 $\Lambda$-代数 $P$
和 $\Lambda$-代数满射 $P \to k$。
（例如，可令 $P$ 为多项式代数。）
以 $\mathfrak m_P$ 表示 $P \to k$ 的核。
Jacobi--Zariski 序列——见
(\ref{formal-defos-equation-sequence-extended}) 和《代数》引理
\ref{algebra-lemma-exact-sequence-NL}——是正合序列
$$
0 \to H_1(\NL_{k/\Lambda}) \to
\mathfrak m_P/\mathfrak m_P^2 \to
\Omega_{P/\Lambda} \otimes_P k \to
\Omega_{k/\Lambda} \to 0
$$
左端有 $0$，是因为 $P/k$ 光滑；% P10-E0440
所以 $\NL_{P/\Lambda}$ 拟同构于置于次数 $0$ 的有限投射模，
进而 $H_1(\NL_{P/\Lambda} \otimes_P k) = 0$。
假设 $f \in \mathfrak m_P$ 映到
$\Omega_{P/\Lambda} \otimes_P k$ 中的非零元素。
令 $P' = P/(f)$，便有 $\Lambda$-代数满射 $P' \to k$。
注意 $P'$ 在 $\mathfrak m_{P'}$ 处光滑：
这由《态射进阶》引理
\ref{more-morphisms-lemma-slice-smooth-given-element} 得出。
因此，用 $P'$ 的一个主局部化替换 $P$ 后，
$\dim(\mathfrak m_P/\mathfrak m_P^2)$ 减小。
有限次重复这一过程后，可以假设映射
$\mathfrak m_P/\mathfrak m_P^2 \to
\Omega_{P/\Lambda} \otimes_P k$ 为零，于是该正合序列分解为同构
$H_1(L_{k/\Lambda}) = \mathfrak m_P/\mathfrak m_P^2$ 和
$\Omega_{P/\Lambda} \otimes_P k = \Omega_{k/\Lambda}$。

\medskip\noindent
令 $R$ 为 $P$ 的 $\mathfrak m_P$-进完备化。
则 $R$ 是 $\widehat{\mathcal{C}}_\Lambda$ 的对象。
事实上，它是完备局部诺特环（见《代数》引理
\ref{algebra-lemma-completion-Noetherian-Noetherian}），
且其剩余域与 $k$ 识别。我们断言这个 $R$ 满足要求。

\medskip\noindent
首先注意，映射 $P \to R$ 诱导同构
$\mathfrak m_P/\mathfrak m_P^2 =
\mathfrak m_R/\mathfrak m_R^2$ 和
$\Omega_{P/\Lambda} \otimes_P k =
\Omega_{R/\Lambda} \otimes_R k$。
这是因为 $\mathfrak m_P/\mathfrak m_P^2$ 与
$\Omega_{P/\Lambda} \otimes_P k$ 都只依赖于
$\Lambda$-代数 $P/\mathfrak m_P^2$；见《代数》引理
\ref{algebra-lemma-differential-mod-power-ideal}。
对 $R$ 也同样如此，而且
$P/\mathfrak m_P^2 = R/\mathfrak m_R^2$。
利用 Jacobi--Zariski 序列
(\ref{formal-defos-equation-sequence-functorial}) 的函子性，
并利用同样结论对 $P$ 成立，可得
$H_1(L_{k/\Lambda}) = \mathfrak m_R/\mathfrak m_R^2$ 和
$\Omega_{R/\Lambda} \otimes_R k = \Omega_{k/\Lambda}$。

\medskip\noindent
最后，由于 $\Lambda \to P$ 光滑，故由《代数》命题
\ref{algebra-proposition-smooth-formally-smooth}，
$\Lambda \to P$ 形式光滑。
于是由《代数进阶》引理
\ref{more-algebra-lemma-formally-smooth}，
$\Lambda \to P$ 关于 $\mathfrak m_P$-进拓扑形式光滑。
由《代数进阶》引理
\ref{more-algebra-lemma-formally-smooth-completion}，
这一性质传递到完备化 $R$，证明完成。
事实上，只要
$\underline{R}|_{\mathcal{C}_\Lambda}$ 光滑，
$R$ 就同构于 $\Lambda$ 上某个光滑代数的完备化，
但我们不会用到这一点。
\end{proof}

\begin{example}
\label{formal-defos-example-smooth-explicit}
下面给出引理 \ref{formal-defos-lemma-exists-smooth} 中 $R$ 的一个更明确的例子。
设 $p$ 为素数，且 $n \in \mathbf{N}$。
令 $\Lambda = \mathbf{F}_p(t_1, t_2, \ldots, t_n)$，
$k = \mathbf{F}_p(x_1, \ldots, x_n)$，并令映射
$\Lambda \to k$ 由 $t_i \mapsto x_i^p$ 给出。此时可取
$$
R = \Lambda[x_1, \ldots, x_n]^\wedge_{(x_1^p - t_1, \ldots, x_n^p - t_n)}
$$
在这个例子中无法做得“更好”；也就是说，不能用
$\widehat{\mathcal{C}}_\Lambda$ 中更小的光滑对象来逼近
$\mathcal{C}_\Lambda$（可以论证 $R$ 的维数至少为 $n$，
因为映射
$\Omega_{R/\Lambda} \otimes_R k \to \Omega_{k/\Lambda}$
为满射）。稍后将更详细地讨论这一现象。
\end{example}

\section{Schlessinger 条件}
\label{formal-defos-section-schlessinger-conditions}

\noindent
下文经常考虑范畴 $\mathcal{C}_\Lambda$ 中环的纤维积
$A_1 \times_A A_2$。例 \ref{formal-defos-example-fibre-product}
已经表明，这样的纤维积未必总是 $\mathcal{C}_\Lambda$ 的对象。
不过，在下文几乎所有情形中，两个映射
$A_i \to A$ 之一都是满射；由引理
\ref{formal-defos-lemma-fiber-product-CLambda}，
$A_1 \times_A A_2$ 将是 $\mathcal{C}_\Lambda$ 的对象。
我们将不再说明而直接使用这一结果。

\medskip\noindent
以 $k[\epsilon]$ 表示 $k$ 上的对偶数环。
更一般地，对 $k$-向量空间 $V$，以 $k[V]$ 表示如下
$k$-代数：其底层向量空间为 $k \oplus V$，乘法由
$(a, v) \cdot (a', v') = (aa', av' + a'v)$ 给出。
当 $V = k$ 时，$k[V]$ 就是 $k$ 上的对偶数环。
对任意有限维 $k$-向量空间 $V$，环 $k[V]$ 都属于
$\mathcal{C}_\Lambda$。

\begin{definition}
\label{formal-defos-definition-S1-S2}
设 $\mathcal{F}$ 是 $\mathcal C_\Lambda$ 上的群胚余纤维范畴。
如下定义 $\mathcal{F}$ 上的{\it 条件 (S1) 和 (S2)}：
\begin{enumerate}
\item[(S1)] $\mathcal{F}$ 中每个图
$$
\vcenter{
\xymatrix{
           & x_2 \ar[d] \\
x_1 \ar[r] & x
}
}
\quad\text{位于下图上方}\quad
\vcenter{
\xymatrix{
           & A_2 \ar[d] \\
A_1 \ar[r] & A
}
}
$$
若其底图属于 $\mathcal{C}_\Lambda$ 且 $A_2 \to A$ 满射，
则可补成交换图
$$
\vcenter{
\xymatrix{
y \ar[r] \ar[d] & x_2 \ar[d] \\
x_1 \ar[r]      & x
}
}
\quad\text{位于下图上方}\quad
\vcenter{
\xymatrix{
A_1 \times_A A_2 \ar[r] \ar[d] & A_2 \ar[d] \\
A_1 \ar[r]      & A.
}
}
$$
\item[(S2)]
对 $\mathcal{F}$ 中位于 $\mathcal{C}_\Lambda$ 中如下形式图上方的图，
(S1) 条件成立：
$$
\xymatrix{
          & k[\epsilon] \ar[d] \\
A  \ar[r] & k.
}
$$
此外，若 $\mathcal{F}$ 中有两个交换图
$$
\vcenter{
\xymatrix{
y \ar[r]_c \ar[d]_a & x_\epsilon \ar[d]^e \\
x \ar[r]^d          & x_0
}
}
\quad\text{和}\quad
\vcenter{
\xymatrix{
y' \ar[r]_{c'} \ar[d]_{a'} & x_\epsilon \ar[d]^e \\
x \ar[r]^d                 & x_0
}
}
\quad\text{位于下图上方}\quad
\vcenter{
\xymatrix{
A \times_k k[\epsilon] \ar[r] \ar[d] & k[\epsilon] \ar[d] \\
A  \ar[r] & k
}
}
$$
则存在
$\mathcal{F}(A \times_k k[\epsilon])$ 中的态射
$b : y \to y'$，使 $a = a' \circ b$。
\end{enumerate}
\end{definition}

\noindent
可以用纤维范畴部分地解释条件 (S1) 与 (S2) 的含义。
设 $f_1 : A_1 \to A$ 和 $f_2 : A_2 \to A$ 是
$\mathcal{C}_\Lambda$ 中的环映射，且 $f_2$ 满射。
以 $p_i : A_1 \times_A A_2 \to A_i$ 表示投影映射。
假设已经为 $\mathcal{F}$ 作了推前选取。
于是环的交换图转化为纤维范畴的 $2$-交换图
$$
\xymatrix{
\mathcal{F}(A_1 \times_A A_2) \ar[r]_-{p_{2, *}} \ar[d]_{p_{1, *}} &
\mathcal{F}(A_2) \ar[d]^{f_{2, *}} \\
\mathcal{F}(A_1) \ar[r]^{f_{1, *}} & \mathcal{F}(A)
}
$$
从而得到到范畴 $2$-纤维积的函子
\begin{equation}
\label{formal-defos-equation-compare}
\mathcal{F}(A_1 \times_A A_2) \to
\mathcal{F}(A_1) \times_{\mathcal{F}(A)} \mathcal{F}(A_2)
\end{equation}
条件 (S1) 要求这个函子本质满。
条件 (S2) 的第一部分要求：若 $f_2$ 等于映射
$k[\epsilon] \to k$，这个函子是本质满的。% P10-E0441
此外，在这种情况下，(S2) 的第二部分蕴含：
在靶中变为同构的两个对象在源中同构
（但它与这一陈述{\it 并不}等价）。
按照定义中的方式陈述这些条件的优点，是无需作任何选取。

\begin{lemma}
\label{formal-defos-lemma-S1-small-extensions}
设 $\mathcal{F}$ 是 $\mathcal C_\Lambda$ 上的群胚余纤维范畴。
若只在 $A_2 \to A$ 为小扩张时假设 (S1) 的条件成立，
则 $\mathcal{F}$ 满足 (S1)。
\end{lemma}

\begin{proof}
证明从略。提示：应用引理
\ref{formal-defos-lemma-factor-small-extension}，并作类似于引理
\ref{formal-defos-lemma-smoothness-small-extensions} 证明中的归纳。
\end{proof}

\begin{remark}
\label{formal-defos-remark-compare-S1-S2-schlessinger}
当 $\mathcal{F}$ 是集合余纤维范畴时，条件 (S1) 与 (S2)
正是 Schlessinger 论文 \cite{Sch} 中的条件 (H1) 与 (H2)。
具体地，对函子
$F: \mathcal{C}_\Lambda \to \textit{Sets}$，
条件 (S1) 与 (S2) 表述为：
\begin{enumerate}
\item [(S1)] 若 $A_1 \to A$ 和 $A_2 \to A$ 是
$\mathcal{C}_\Lambda$ 中的映射，且 $A_2 \to A$ 满射，
则诱导映射
$F(A_1 \times_A A_2) \to F(A_1) \times_{F(A)} F(A_2)$
为满射。
\item [(S2)] 若 $A \to k$ 是 $\mathcal{C}_\Lambda$ 中的映射，
则诱导映射
$F(A \times_k k[\epsilon]) \to F(A) \times_{F(k)} F(k[\epsilon])$
为双射。
\end{enumerate}
映射
$F(A \times_k k[\epsilon]) \to F(A) \times_{F(k)} F(k[\epsilon])$
的单射性来自条件 (S2) 的第二部分以及所有态射都是恒等态射这一事实。
\end{remark}

\begin{lemma}
\label{formal-defos-lemma-S2-extensions}
设 $\mathcal{F}$ 是 $\mathcal{C}_\Lambda$ 上的群胚余纤维范畴。
若 $\mathcal{F}$ 满足 (S2)，则对任意有限维
$k$-向量空间 $V$，把 $k[\epsilon]$ 换成 $k[V]$ 后，
(S2) 的条件仍成立。
\end{lemma}

\begin{proof}
若 $\mathcal{F}$ 是集合余纤维范畴，即对应于函子
$F : \mathcal{C}_\Lambda \to \textit{Sets}$，
则结论来自说明 \ref{formal-defos-remark-compare-S1-S2-schlessinger}
中对 $F$ 的 (S2) 的描述，以及同构
$k[V] \cong k[\epsilon] \times_k \ldots \times_k k[\epsilon]$；
右边有 $\dim_k V$ 个因子。
本章余下部分只会用到函子情形。

\medskip\noindent
一般情形下，对 $\dim(V)$ 作归纳。
若 $\dim(V) = 1$，则 $k[V] \cong k[\epsilon]$，
结论由假设成立。若 $\dim(V) > 1$，写成
$V = V' \oplus k\epsilon$。选取图
$$
\vcenter{
\xymatrix{
& x_V \ar[d] \\
x \ar[r] & x_0
}
}
\quad\text{位于下图上方}\quad
\vcenter{
\xymatrix{
& k[V] \ar[d] \\
A \ar[r] & k
}
}
$$
选取位于 $k[V] \to k[V']$ 上方的态射
$x_V \to x_{V'}$，以及位于
$k[V] \to k[\epsilon]$ 上方的态射
$x_V \to x_\epsilon$。
注意，态射 $x_V \to x_0$ 分别分解为
$x_V \to x_{V'} \to x_0$ 和
$x_V \to x_\epsilon \to x_0$。
由归纳假设，可以找到图
$$
\vcenter{
\xymatrix{
y' \ar[d] \ar[r] & x_{V'} \ar[d] \\
x \ar[r] & x_0
}
}
\quad\text{位于下图上方}\quad
\vcenter{
\xymatrix{
A \times_k k[V'] \ar[d] \ar[r] & k[V'] \ar[d] \\
A \ar[r] & k
}
}
$$
这给出交换图
$$
\vcenter{
\xymatrix{
& x_\epsilon \ar[d] \\
y' \ar[r] & x_0
}
}
\quad\text{位于下图上方}\quad
\vcenter{
\xymatrix{
& k[\epsilon] \ar[d] \\
A \times_k k[V'] \ar[r] & k
}
}
$$
因此，由 (S2) 得到交换图
$$
\vcenter{
\xymatrix{
y \ar[d] \ar[r] & x_\epsilon \ar[d] \\
y' \ar[r] & x_0
}
}
\quad\text{位于下图上方}\quad
\vcenter{
\xymatrix{
(A \times_k k[V']) \times_k k[\epsilon] \ar[d] \ar[r] &
k[\epsilon] \ar[d] \\
A \times_k k[V'] \ar[r] & k
}
}
$$
注意
$(A \times_k k[V']) \times_k k[\epsilon] =
A \times_k k[V' \oplus k\epsilon]
= A \times_k k[V]$。
我们断言 $y$ 可置于正确的交换图中。
为说明这一点，取位于
$A \times_k k[V] \to k[V]$ 上方的态射
$y \to y_V$。可以把态射
$y \to y' \to x_{V'}$ 和 $y \to x_\epsilon$
经由态射 $y \to y_V$ 分解
（由群胚余纤维范畴的公理）。
因此，$y_V$ 与 $x_V$ 都可置于交换图
$$
\vcenter{
\xymatrix{
y_V \ar[r] \ar[d] & x_\epsilon \ar[d] \\
x_{V'} \ar[r]          & x_0
}
}
\quad\text{和}\quad
\vcenter{
\xymatrix{
x_V \ar[r] \ar[d] & x_\epsilon \ar[d] \\
x_{V'} \ar[r]                 & x_0
}
}
$$
中。因此，由 (S2) 的第二部分，存在同构
$y_V \to x_V$，它与 $y_V \to x_{V'}$ 和
$x_V \to x_{V'}$ 相容，特别地也与到 $x_0$ 的映射相容。
于是复合 $y \to y_V \to x_V$ 可置于所需交换图
$$
\vcenter{
\xymatrix{
y \ar[r] \ar[d] & x_V \ar[d] \\
x \ar[r] & x_0
}
}
\quad\text{位于下图上方}\quad
\vcenter{
\xymatrix{
A \times_k k[V] \ar[d] \ar[r] & k[V] \ar[d] \\
A \ar[r] & k
}
}
$$
中。由此可见，把 $k[\epsilon]$ 换成 $k[V]$ 后，
$(S2)$ 的第一部分成立。

\medskip\noindent
为证明第二部分，设给定两个交换图
$$
\vcenter{
\xymatrix{
y \ar[r] \ar[d] & x_V \ar[d] \\
x \ar[r] & x_0
}
}
\quad\text{和}\quad
\vcenter{
\xymatrix{
y' \ar[r] \ar[d] & x_V \ar[d] \\
x \ar[r] & x_0
}
}
\quad\text{位于下图上方}\quad
\vcenter{
\xymatrix{
A \times_k k[V] \ar[d] \ar[r] & k[V] \ar[d] \\
A \ar[r] & k
}
}
$$
我们将使用本证明第一段所引入的态射
$x_V \to x_{V'} \to x_0$ 和
$x_V \to x_\epsilon \to x_0$。
选取位于
$A \times_k k[V] \to A \times_k k[V']$ 上方的态射
$y \to y_{V'}$ 和 $y' \to y'_{V'}$。
余纤维范畴的公理蕴含可以找到交换图
$$
\vcenter{
\xymatrix{
y_{V'} \ar[r] \ar[d] & x_{V'} \ar[d] \\
x \ar[r] & x_0
}
}
\quad\text{和}\quad
\vcenter{
\xymatrix{
y'_{V'} \ar[r] \ar[d] & x_{V'} \ar[d] \\
x \ar[r] & x_0
}
}
\quad\text{位于下图上方}\quad
\vcenter{
\xymatrix{
A \times_k k[V'] \ar[d] \ar[r] & k[V'] \ar[d] \\
A \ar[r] & k
}
}
$$
由归纳假设，得到同构
$b : y_{V'} \to y'_{V'}$，
它与态射 $y_{V'} \to x$ 和
$y'_{V'} \to x$ 相容，特别地与到 $x_0$ 的态射相容。
于是有交换图
$$
\vcenter{
\xymatrix{
y \ar[r] \ar[d] & x_\epsilon \ar[d] \\
y'_{V'} \ar[r] & x_0
}
}
\quad\text{和}\quad
\vcenter{
\xymatrix{
y' \ar[r] \ar[d] & x_\epsilon \ar[d] \\
y'_{V'} \ar[r] & x_0
}
}
\quad\text{位于下图上方}\quad
\vcenter{
\xymatrix{
A \times_k k[\epsilon] \ar[d] \ar[r] & k[\epsilon] \ar[d] \\
A \ar[r] & k
}
}
$$
其中态射 $y \to y'_{V'}$ 是复合
$y \to y_{V'} \xrightarrow{b} y'_{V'}$，
而态射 $y \to x_\epsilon$ 和
$y' \to x_\epsilon$ 分别是映射
$y \to x_V$ 和 $y' \to x_V$ 与态射
$x_V \to x_\epsilon$ 的复合。
于是 (S2) 的第二部分保证存在同构
$y \to y'$，它与到 $y'_{V'}$ 的映射相容，
特别地与到 $x$ 的映射相容
（因为 $b$ 与到 $x$ 的映射相容）。
\end{proof}

\begin{lemma}
\label{formal-defos-lemma-S1-S2-associated-functor}
设 $\mathcal{F}$ 是 $\mathcal{C}_\Lambda$ 上的群胚余纤维范畴。
\begin{enumerate}
\item 若 $\mathcal{F}$ 满足 (S1)，则
$\overline{\mathcal{F}}$ 也满足 (S1)。
\item 若 $\mathcal{F}$ 满足 (S2)，并且下列条件至少有一个成立，
则 $\overline{\mathcal{F}}$ 也满足 (S2)：
\begin{enumerate}
\item $\mathcal{F}$ 是预形变范畴；
\item 范畴 $\mathcal{F}(k)$ 是集合或集合型群胚；
\item 对 $\mathcal{F}$ 中位于
$k[\epsilon] \to k$ 上方的任意态射
$x_\epsilon \to x_0$，推前映射
$\text{Aut}_{k[\epsilon]}(x_\epsilon) \to
\text{Aut}_k(x_0)$ 为满射。
\end{enumerate}
\end{enumerate}
\end{lemma}

\begin{proof}
假设 $\mathcal{F}$ 满足 (S1)。
设 $f_i : A_i \to A$ 是 $\mathcal{C}_\Lambda$ 中的环映射，
且 $f_2$ 满射。设 $x_i \in \mathcal{F}(A_i)$，
并且推前 $f_{1, *}(x_1)$ 与 $f_{2, *}(x_2)$ 同构。
于是可以用 $x$ 表示 $\mathcal{F}$ 中位于 $A$ 上方、
且与这两个推前都同构的对象，% P10-E0442
并得到 (S1) 中的图。因此，找到 $\mathcal{F}$ 中位于
$A_1 \times_A A_2$ 上方的对象 $y$，
其到 $A_1$、相应地到 $A_2$ 的推前分别与
$x_1$、相应地与 $x_2$ 同构。
由此可见，(S1) 对 $\overline{\mathcal{F}}$ 成立。

\medskip\noindent
假设 $\mathcal{F}$ 满足 (S2)。
$\overline{\mathcal{F}}$ 的 (S2) 第一部分由与上述相同的论证得出。
$\overline{\mathcal{F}}$ 的 (S2) 第二部分意指：对
$\mathcal{C}_\Lambda$ 中的任意环 $A$，映射
$$
\overline{\mathcal{F}}(A \times_k k[\epsilon]) \to
\overline{\mathcal{F}}(A)
\times_{\overline{\mathcal{F}}(k)}
\overline{\mathcal{F}}(k[\epsilon])
$$
为单射。设
$y, y' \in \mathcal{F}(A \times_k k[\epsilon])$。
利用余纤维范畴的公理，可以选取交换图
$$
\vcenter{
\xymatrix{
y \ar[r]_c \ar[d]_a & x_\epsilon \ar[d]^e \\
x \ar[r]^d          & x_0
}
}
\quad\text{和}\quad
\vcenter{
\xymatrix{
y' \ar[r]_{c'} \ar[d]_{a'} & x'_\epsilon \ar[d]^{e'} \\
x' \ar[r]^{d'}                 & x'_0
}
}
\quad\text{位于下图上方}\quad
\vcenter{
\xymatrix{
A \times_k k[\epsilon] \ar[d] \ar[r] & k[\epsilon] \ar[d] \\
A \ar[r] & k
}
}
$$
假设存在 $\mathcal{F}(A)$ 中的同构
$\alpha : x \to x'$，以及
$\mathcal{F}(k[\epsilon])$ 中的同构
$\beta : x_\epsilon \to x'_\epsilon$。
这也意味着存在与 $\alpha$ 相容的同构
$\gamma : x_0 \to x'_0$。
为证明 $\overline{\mathcal{F}}$ 的 (S2)，必须证明
$\mathcal{F}(A \times_k k[\epsilon])$ 中存在同构
$y \to y'$。
由 $\mathcal{F}$ 的 (S2)，如果能选取同构
$\alpha$、$\beta$ 和 $\gamma$，使
$$
\xymatrix{
x \ar[d]^\alpha \ar[r] & x_0 \ar[d]^\gamma &
x_\epsilon \ar[d]^\beta \ar[l]^e \\
x' \ar[r] & x'_0 & x'_\epsilon \ar[l]_{e'}
}
$$
交换，这样的态射就存在
（因为这时可在前一个显示图中用 $x'$ 替换 $x$，
并用 $x'_\epsilon$ 替换 $x_\epsilon$）。
左侧方块由 $\gamma$ 的选取而交换。
可把 $e' \circ \beta$ 分解成
$\gamma' \circ e$，其中
$\gamma' : x_0 \to x'_0$ 是另一个映射。
现在的问题是能否安排 $\gamma = \gamma'$？
若 $\mathcal{F}(k)$ 是集合或集合型群胚，这显然成立。
此外，若
$\text{Aut}_{k[\epsilon]}(x_\epsilon) \to
\text{Aut}_k(x_0)$ 为满射，
则可在 $\beta$ 前复合一个 $x_\epsilon$ 的自同构，
其像为 $\gamma^{-1} \circ \gamma'$，从而使要求成立。
\end{proof}

\begin{lemma}
\label{formal-defos-lemma-S1-S2-localize}
设 $\mathcal{F}$ 是 $\mathcal{C}_\Lambda$ 上的群胚余纤维范畴，
且 $x_0 \in \Ob(\mathcal{F}(k))$。
设 $\mathcal{F}_{x_0}$ 是说明
\ref{formal-defos-remark-localize-cofibered-groupoid} 中构造的
$\mathcal{C}_\Lambda$ 上群胚余纤维范畴。
\begin{enumerate}
\item 若 $\mathcal{F}$ 满足 (S1)，则
$\mathcal{F}_{x_0}$ 也满足 (S1)。
\item 若 $\mathcal{F}$ 满足 (S2)，则
$\mathcal{F}_{x_0}$ 也满足 (S2)。
\end{enumerate}
\end{lemma}

\begin{proof}
$\mathcal{F}_{x_0}$ 中每个如定义
\ref{formal-defos-definition-S1-S2} 所述的图，都产生 $\mathcal{F}$ 中的图；
而条件 (S1) 或 (S2) 对 $\mathcal{F}$ 中该图给出的输出，
也可视为对 $\mathcal{F}_{x_0}$ 的输出。
\end{proof}

\begin{lemma}
\label{formal-defos-lemma-lifting-section}
设 $p: \mathcal{F} \to \mathcal{C}_\Lambda$ 是群胚余纤维范畴。
考虑 $\mathcal{F}$ 中的图
$$
\vcenter{
\xymatrix{
y \ar[r] \ar[d]_a & x_\epsilon \ar[d]_e \\
x \ar[r]^d        & x_0
}
}
\quad\text{位于下图上方}\quad
\vcenter{
\xymatrix{
A \times_k k[\epsilon] \ar[r] \ar[d] & k[\epsilon] \ar[d] \\
A \ar[r] & k.
}
}
$$
其中底图属于 $\mathcal{C}_\Lambda$。
假设 $\mathcal{F}$ 满足 (S2)。
则存在满足 $a \circ s = \text{id}_x$ 的态射
$s : x \to y$，当且仅当存在满足
$e \circ s_\epsilon = d$ 的态射
$s_\epsilon : x \to x_\epsilon$。
\end{lemma}

\begin{proof}
“仅当”方向显然。反过来，假设存在态射
$s_\epsilon : x \to x_\epsilon$，使
$e \circ s_\epsilon = d$。
注意，$p(s_\epsilon) : A \to k[\epsilon]$ 是与映射
$A \to k$ 相容的环映射。因此得到
$$
\sigma = (\text{id}_A, p(s_\epsilon)) :
A \to A \times_k k[\epsilon].
$$
选取推前 $x \to \sigma_*x$。
由构造，可把 $s_\epsilon$ 分解为
$x \to \sigma_*x \to x_\epsilon$。
此外，由于 $\sigma$ 是
$A \times_k k[\epsilon] \to A$ 的截面，
得到态射 $\sigma_*x \to x$，使
$x \to \sigma_*x \to x$ 为 $\text{id}_x$。
因为 $e \circ s_\epsilon = d$，图
$$
\xymatrix{
\sigma_*x \ar[r] \ar[d] & x_\epsilon \ar[d]_e \\
x \ar[r]^d        & x_0
}
$$
交换。因此，由 (S2) 得到态射
$\sigma_*x \to y$，使复合
$\sigma_*x \to y \to x$ 等于给定映射
$\sigma_*x \to x$。现在问题的解是取
$a : x \to y$ 为复合
$x \to \sigma_*x \to y$。% P10-E0443
\end{proof}

\begin{lemma}
\label{formal-defos-lemma-lifting-along-small-extension}
考虑预形变范畴 $\mathcal{F}$ 中的交换图
$$
\vcenter{
\xymatrix{
y \ar[r] \ar[d] & x_2 \ar[d]^{a_2} \\
x_1 \ar[r]^{a_1}        & x
}
}
\quad\text{位于下图上方}
\vcenter{
\xymatrix{
A_1 \times_A A_2 \ar[r] \ar[d] & A_2 \ar[d]^{f_2} \\
A_1 \ar[r]^{f_1} & A
}
}
$$
上方，其中底图属于 $\mathcal{C}_\Lambda$，且
$f_2 : A_2 \to A$ 是小扩张。
假设存在映射 $h : A_1 \to A_2$，使
$f_2 = f_1 \circ h$。% P10-E0444
令 $I = \Ker(f_2)$。考虑环映射
$$
g : A_1 \times_A A_2 \longrightarrow k[I] = k \oplus I, \quad
(u, v) \longmapsto \overline{u} \oplus (v - h(u))
$$
选取推前 $y \to g_*y$。假设 $\mathcal{F}$ 满足 (S2)。
若存在态射 $x_1 \to g_*y$，则存在态射
$b: x_1 \to x_2$，使 $a_1 =  a_2 \circ b$。
\end{lemma}

\begin{proof}
注意，
$\text{id}_{A_1} \times g :
A_1 \times_A A_2 \to A_1 \times_k k[I]$
是同构，且 $k[I] \cong k[\epsilon]$。因此有图
$$
\vcenter{
\xymatrix{
y \ar[r] \ar[d] & g_*y \ar[d] \\
x_1 \ar[r]        & x_0
}
}
\quad\text{位于下图上方}\quad
\vcenter{
\xymatrix{
A_1 \times_k k[\epsilon] \ar[r] \ar[d] & k[\epsilon] \ar[d] \\
A_1 \ar[r] & k.
}
}
$$
其中 $x_0$ 是 $\mathcal{F}$ 中位于 $k$ 上方的对象
（$\mathcal{F}$ 的每个对象都有到 $x_0$ 的唯一态射；
见定义 \ref{formal-defos-definition-predeformation-category} 之后的讨论）。
若有态射 $x_1 \to g_*y$，则引理
\ref{formal-defos-lemma-lifting-section} 给出映射
$y \to x_1$ 的截面 $s : x_1 \to y$。
将它与映射 $y \to x_2$ 复合，得到
$b : x_1 \to x_2$；它满足
$a_1 =  a_2 \circ b$，因为引理中的图交换且 $s$ 是截面。
\end{proof}





\section{函子的切空间}
\label{formal-defos-section-tangent-spaces-functors}

\noindent
设 $R$ 是环。用 $\text{Mod}_R$ 表示 $R$-模范畴，
用 $\text{Mod}^{fg}_R$ 表示有限生成 $R$-模范畴。

\begin{definition}
\label{formal-defos-definition-linear}
设 $L: \text{Mod}^{fg}_R \to \text{Mod}_R$，
相应地，设 $L: \text{Mod}_R \to \text{Mod}_R$ 是函子。
若对 $\text{Mod}^{fg}_R$、相应地 $\text{Mod}_R$ 中的每一对对象
$M,N$，映射
$$
L : \Hom_R(M, N) \longrightarrow \Hom_R(L(M), L(N))
$$
都是 $R$-模映射，则称 $L$ 是 {\it $R$-线性的}。
\end{definition}

\begin{remark}
\label{formal-defos-remark-linear-enriched-over-modules}
对任意两个在 $\text{Mod}_R$ 上富集的范畴之间的函子，
都可以定义 $R$-线性的概念。这里专门对函子
$L: \text{Mod}^{fg}_R \to \text{Mod}_R$ 和
$L: \text{Mod}_R \to \text{Mod}_R$
作出定义，是因为迄今所需的正是这两种情形。
\end{remark}

\begin{remark}
\label{formal-defos-remark-linear-functor}
若 $L: \text{Mod}^{fg}_R \to \text{Mod}_R$ 是 $R$-线性函子，
则 $L$ 保持有限积，并把零模送到零模；见
《同调》引理 \ref{homology-lemma-additive-additive}。
另一方面，若函子 $\text{Mod}^{fg}_R \to \textit{Sets}$
保持有限积，并把零模送到单元素集合，
则它可唯一地提升为 $R$-线性函子；见引理
\ref{formal-defos-lemma-linear-functor}。
\end{remark}

\begin{lemma}
\label{formal-defos-lemma-linear-functor}
设 $L: \text{Mod}^{fg}_R \to \textit{Sets}$，
相应地，设 $L: \text{Mod}_R \to \textit{Sets}$ 是函子。
假设 $L(0)$ 是单元素集合，并且 $L$ 保持有限积。
则存在唯一的 $R$-线性函子
$\widetilde{L} : \text{Mod}^{fg}_R \to \text{Mod}_R$，
相应地，存在唯一的
$\widetilde{L} : \text{Mod}_R \to \text{Mod}_R$，使图
$$
\vcenter{
\xymatrix{
& \text{Mod}_R \ar[dr]^{\text{遗忘}} &   \\
\text{Mod}^{fg}_R  \ar[ur]^{\widetilde{L}} \ar[rr]^{L} &  &
\textit{Sets}
}
}
\quad\text{相应地}\quad
\vcenter{
\xymatrix{
& \text{Mod}_R \ar[dr]^{\text{遗忘}} &   \\
\text{Mod}_R  \ar[ur]^{\widetilde{L}} \ar[rr]^{L} &  &
\textit{Sets}
}
}
$$
交换。
\end{lemma}

\begin{proof}
这里只证明 $L: \text{Mod}^{fg}_R \to \textit{Sets}$ 的情形。
设 $M$ 是有限生成 $R$-模。我们把 $\widetilde{L}(M)$
定义为集合 $L(M)$ 配备以下 $R$-模结构。

\medskip\noindent
乘法：若 $r \in R$，则 $L(M)$ 上的乘以 $r$ 定义为由乘法映射
$r \cdot : M \to M$ 诱导的映射 $L(M) \to L(M)$。

\medskip\noindent
加法：求和映射 $M \times M \to M: (m_1, m_2) \mapsto m_1 + m_2$
诱导映射 $L(M \times M) \to L(M)$。由假设，$L(M \times M)$
典范同构于 $L(M) \times L(M)$。于是 $L(M)$ 上的加法定义为映射
$L(M) \times L(M) \cong L(M \times M) \to L(M)$。

\medskip\noindent
零元：存在唯一映射 $0 \to M$。$L(M)$ 的零元是映射
$L(0) \to L(M)$ 的像。

\medskip\noindent
我们略去验证：这定义了 $R$-模 $\widetilde{L}(M)$，
并且这是唯一使其关于 $M$ 具有 $R$-线性函子性的结构。% P10-E0445
\end{proof}

\begin{lemma}
\label{formal-defos-lemma-morphism-linear-functors}
设 $L_1, L_2: \text{Mod}^{fg}_R \to \textit{Sets}$ 是把 $0$
送到单元素集合并保持有限积的函子。
设 $t : L_1 \to L_2$ 是函子态射。
则 $t$ 在引理 \ref{formal-defos-lemma-linear-functor} 所保证的函子之间诱导态射
$\widetilde{t} : \widetilde{L}_1 \to \widetilde{L}_2$；
对每个 $M \in \Ob(\text{Mod}^{fg}_R)$，它就是
$\widetilde{t}_M = t_M: \widetilde{L}_1(M) \to \widetilde{L}_2(M)$。
换言之，$t_M: \widetilde{L}_1(M) \to \widetilde{L}_2(M)$
是 $R$-模映射。
\end{lemma}

\begin{proof}
证明从略。
\end{proof}

\noindent
当 $R = K$ 是域时，$K$-线性函子
$L : \text{Mod}^{fg}_K \to \text{Mod}_K$ 由它的值
$L(K)$ 决定。

\begin{lemma}
\label{formal-defos-lemma-linear-functor-over-field}
设 $K$ 是域，并设 $L: \text{Mod}^{fg}_K \to
\text{Mod}_K$ 是 $K$-线性函子。则 $L$ 同构于函子
$L(K) \otimes_K - : \text{Mod}^{fg}_K \to
\text{Mod}_K$。
\end{lemma}

\begin{proof}
对 $V \in \Ob(\text{Mod}^{fg}_K)$，同构
$L(K) \otimes_K V \to L(V)$ 在单纯张量上由
$x \otimes v \mapsto L(f_v)(x)$ 给出，其中
$f_v: K \to V$ 是把 $1 \mapsto v$ 的 $K$-线性映射。
当 $V = K$ 时，这是由 $K$ 上的乘法给出的同构
$L(K) \otimes_K K \to L(K)$。% P10-E0446
对一般的 $V$，由 $V=K$ 的情形以及 $L$ 与有限积交换这一事实
（说明 \ref{formal-defos-remark-linear-functor}），它是同构。
\end{proof}

\noindent
对环 $R$ 和 $R$-模 $M$，令 $R[M]$ 表示这样的 $R$-代数：
其底层 $R$-模是 $R \oplus M$，乘法由
$(r, m) \cdot (r', m') = (rr', rm' + r'm)$ 给出。
当 $M = R$ 时，这是 $R$ 上的对偶数环，记作 $R[\epsilon]$。

\medskip\noindent
现在设 $S$ 是环，并假设 $R$ 是 $S$-代数。
赋值 $M \mapsto R[M]$ 决定函子
$\text{Mod}_R \to S\text{-Alg}/R$，其中 $S\text{-Alg}/R$
表示以 $R$ 为靶的 $S$-代数范畴。注意，
$S\text{-Alg}/R$ 有有限积：若 $A_1 \to R$ 和
$A_2 \to R$ 是两个对象，则 $A_1 \times_R A_2$ 是它们的积。

\begin{lemma}
\label{formal-defos-lemma-preserves-products}
设 $R$ 是 $S$-代数。则上述函子
$\text{Mod}_R \to S\text{-Alg}/R$ 保持有限积。
\end{lemma}

\begin{proof}
这不过是说：若 $M$、$N$ 是 $R$-模，则映射
$R[M \times N] \to R[M] \times_R R[N]$ 是
$S\text{-Alg}/R$ 中的同构。
\end{proof}

\begin{lemma}
\label{formal-defos-lemma-tangent-space-functor}
设 $R$ 是 $S$-代数，并设 $\mathcal{C}$ 是 $S\text{-Alg}/R$
的严格满子范畴，它包含所有
$M \in \Ob(\text{Mod}^{fg}_R)$ 所对应的 $R[M]$。
设 $F: \mathcal{C} \to \textit{Sets}$ 是函子。
假设 $F(R)$ 是单元素集合，并且对任意
$M, N \in \Ob(\text{Mod}^{fg}_R)$，诱导映射
$$
F(R[M] \times_R R[N]) \to F(R[M]) \times F(R[N])
$$
都是双射。则对任意 $M \in \Ob(\text{Mod}^{fg}_R)$，
$F(R[M])$ 都有自然的 $R$-模结构。
\end{lemma}

\begin{proof}
注意，$R \cong R[0]$ 且
$R[M] \times_R R[N] \cong R[M \times N]$，
所以由对 $\mathcal{C}$ 的假设，$R$ 与
$R[M] \times_R R[N]$ 都是 $\mathcal{C}$ 的对象。
因此，对 $F$ 的条件是有意义的。
引理 \ref{formal-defos-lemma-preserves-products} 中的函子
$\text{Mod}_R \to S\text{-Alg}/R$，由对 $\mathcal{C}$ 的假设，
限制为函子 $\text{Mod}^{fg}_R \to \mathcal{C}$。
令 $L$ 为复合
$\text{Mod}^{fg}_R \to \mathcal{C} \to \textit{Sets}$，即
$L(M) = F(R[M])$。
由引理 \ref{formal-defos-lemma-preserves-products} 及对 $F$ 的假设，
$L$ 保持有限积。因此，引理 \ref{formal-defos-lemma-linear-functor} 表明：
对任意 $M \in \Ob(\text{Mod}^{fg}_R)$，
$L(M) = F(R[M])$ 都有自然的 $R$-模结构。
\end{proof}

\begin{definition}
\label{formal-defos-definition-tangent-space-over-R}
设 $\mathcal{C}$ 是引理 \ref{formal-defos-lemma-tangent-space-functor}
中那样的范畴。设 $F : \mathcal{C} \to \textit{Sets}$ 是函子，
且 $F(R)$ 是单元素集合。$F$ 的{\it 切空间 $TF$} 定义为
$F(R[\epsilon])$。
\end{definition}

\noindent
当 $F : \mathcal{C} \to \textit{Sets}$ 满足引理
\ref{formal-defos-lemma-tangent-space-functor} 的假设时，
切空间 $TF$ 具有自然的 $R$-模结构。

\begin{example}
\label{formal-defos-example-tangent-space-functor}
因为 $\mathcal{C}_\Lambda$ 包含有限维向量空间 $V$ 对应的所有
$k[V]$，所以定义 \ref{formal-defos-definition-tangent-space-over-R} 可应用于
$S = \Lambda$、$R = k$、$\mathcal{C} = \mathcal{C}_\Lambda$，
以及预形变函子 $F : \mathcal{C}_\Lambda \to \textit{Sets}$。
切空间为 $TF = F(k[\epsilon])$。
\end{example}

\begin{example}
\label{formal-defos-example-tangent-space-prorepresentable-functor}
在例 \ref{formal-defos-example-tangent-space-functor} 中，设
$F : \mathcal{C}_\Lambda \to \textit{Sets}$ 是 pro-可表函子；
更具体地，对某个
$S \in \Ob(\widehat{\mathcal{C}}_\Lambda)$，设
$F = \underline{S}|_{\mathcal{C}_\Lambda}$。
则 $F$ 与任意极限交换，因而满足引理
\ref{formal-defos-lemma-tangent-space-functor} 的假设。计算得
$$
TF = F(k[\epsilon]) = \Mor_{\mathcal{C}_\Lambda}(S, k[\epsilon])
= \text{Der}_\Lambda(S, k)
$$
更一般地，对有限维 $k$-向量空间 $V$，有
$$
F(k[V]) = \Mor_{\mathcal{C}_\Lambda}(S, k[V]) = \text{Der}_\Lambda(S, V).
$$
% P10-E0447
具体地，与增广 $q : S \to k$ 和 $k[V] \to k$ 相容的
$\Lambda$-代数映射 $f : S \to k[V]$，对应于由
$s \mapsto f(s) - q(s)$ 定义的导子 $D$。
反过来，$\Lambda$-导子 $D : S \to V$ 对应于
$\mathcal{C}_\Lambda$ 中按公式 $f(s) = q(s) + D(s)$ 定义的
$f : S \to k[V]$。由于这些识别具有函子性，
$TF$ 与 $\text{Der}_\Lambda(S, k)$ 上的 $k$-向量空间结构彼此对应
（见引理 \ref{formal-defos-lemma-morphism-linear-functors}）。% P10-E0448
由引理 \ref{formal-defos-lemma-derivations-finite}，可知 $\dim_k TF$ 有限。
\end{example}

\begin{example}
\label{formal-defos-example-tangent-space-classical-prorepresentable-functor}
在经典情形下，例
\ref{formal-defos-example-tangent-space-prorepresentable-functor} 的计算可以简化。
确切地说，此时函子
$F = \underline{S}|_{\mathcal{C}_\Lambda}$ 的切空间就是
$S$ 在 $\Lambda$ 上的相对余切空间，即
$TF = T_{S/\Lambda}$。事实上，当域扩张 $k/k'$ 可分时，
这一结论仍然成立。参见《习题》中的习题
\ref{exercises-exercise-tangent-space-Zariski}。
\end{example}

\begin{lemma}
\label{formal-defos-lemma-morphism-tangent-spaces}
设 $F, G: \mathcal{C} \to \textit{Sets}$ 是满足引理
\ref{formal-defos-lemma-tangent-space-functor} 假设的函子。
设 $t : F \to G$ 是函子态射。
对任意 $M \in \Ob(\text{Mod}^{fg}_R)$，映射
$t_{R[M]}: F(R[M]) \to G(R[M])$ 是 $R$-模映射，
其中 $F(R[M])$ 与 $G(R[M])$ 配备引理
\ref{formal-defos-lemma-tangent-space-functor} 给出的 $R$-模结构。
特别地，$t_{R[\epsilon]} : TF \to TG$ 是 $R$-模映射。
\end{lemma}

\begin{proof}
这由引理 \ref{formal-defos-lemma-morphism-linear-functors} 得出。
\end{proof}

\begin{example}
\label{formal-defos-example-tangent-space-map-prorepresentable-functor}
假设 $f : R \to S$ 是 $\widehat{\mathcal{C}}_\Lambda$ 中的环映射。
令 $F = \underline{R}|_{\mathcal{C}_\Lambda}$，
$G = \underline{S}|_{\mathcal{C}_\Lambda}$。
环映射 $f$ 诱导函子变换 $G \to F$。
由引理 \ref{formal-defos-lemma-morphism-tangent-spaces}，得到 $k$-线性映射
$TG \to TF$。这就是映射
$$
TG = \text{Der}_\Lambda(S, k) \longrightarrow \text{Der}_\Lambda(R, k) = TF
$$
这一点来自例 \ref{formal-defos-example-tangent-space-prorepresentable-functor}
中的典范识别
$F(k[V]) = \text{Der}_\Lambda(R, V)$ 和
$G(k[V]) = \text{Der}_\Lambda(S, V)$，
以及切空间上映射的计算规则。
\end{example}

\begin{lemma}
\label{formal-defos-lemma-tangent-space-tensor}
设 $F: \mathcal{C} \to \textit{Sets}$ 是满足引理
\ref{formal-defos-lemma-tangent-space-functor} 假设的函子。
假设 $R = K$ 是域。则对任意有限维 $K$-向量空间 $V$，有
$F(K[V]) \cong TF \otimes_K V$。
\end{lemma}

\begin{proof}
这由引理 \ref{formal-defos-lemma-linear-functor-over-field} 得出。
\end{proof}





\section{预形变范畴的切空间}
\label{formal-defos-section-tangent-spaces}

\noindent
我们将用一般的定义 \ref{formal-defos-definition-tangent-space-over-R}
定义预形变函子的切空间。例 \ref{formal-defos-example-tangent-space-functor}
已经具体说明了这一点。考察同构类的关联函子，
即可把该定义应用于预形变范畴。

\begin{definition}
\label{formal-defos-definition-tangent-space}
设 $\mathcal{F}$ 是预形变范畴。
$\mathcal{F}$ 的{\it 切空间 $T \mathcal{F}$} 是纤维范畴
$\mathcal F(k[\epsilon])$ 中对象的同构类集合
$\overline{\mathcal{F}}(k[\epsilon])$。
\end{definition}

\noindent
因此，$T \mathcal{F}$ 就是关联函子
$\overline{\mathcal{F}}: \mathcal{C}_\Lambda \to \textit{Sets}$
的切空间。当 $\mathcal{F}$ 满足 (S2) 时，
它具有自然的向量空间结构；事实上，只要
$\overline{\mathcal{F}}$ 满足 (S2) 即可。

\begin{lemma}
\label{formal-defos-lemma-tangent-space-vector-space}
设 $\mathcal{F}$ 是预形变范畴，且
$\overline{\mathcal{F}}$ 满足 (S2)\footnote{例如，若
$\mathcal{F}$ 满足 (S2)，则此条件成立；见引理
\ref{formal-defos-lemma-S1-S2-associated-functor}。}。
则 $T \mathcal{F}$ 具有自然的 $k$-向量空间结构。
对任意有限维向量空间 $V$，都有
$\overline{\mathcal{F}}(k[V]) = T\mathcal{F} \otimes_k V$，
且这一等式关于 $V$ 具有函子性。
\end{lemma}

\begin{proof}
记
$F = \overline{\mathcal{F}} : \mathcal{C}_\Lambda \to \textit{Sets}$。
这是预形变函子，并且 $F$ 满足 (S2)。由引理
\ref{formal-defos-lemma-S2-extensions}
（以及说明 \ref{formal-defos-remark-compare-S1-S2-schlessinger}
对集合值函子的对应表述）% P10-E0449
可知，对每个有限维向量空间 $V$ 和每个
$A \in \Ob(\mathcal{C}_\Lambda)$，映射
$$
F(A \times_k k[V]) \longrightarrow F(A) \times F(k[V])
$$
都是双射。特别地，取 $A = k[W]$，可得
$F(k[W] \times_k k[V]) = F(k[W]) \times F(k[V])$。
换言之，引理 \ref{formal-defos-lemma-tangent-space-functor} 的假设成立，
从而 $TF = T \mathcal{F}$ 具有自然的 $k$-向量空间结构。
最后一个断言由引理 \ref{formal-defos-lemma-tangent-space-tensor} 得出。
\end{proof}

\noindent
预形变范畴的态射在切空间上诱导映射。

\begin{definition}
\label{formal-defos-definition-differential}
设 $\varphi : \mathcal{F} \to \mathcal{G}$ 是预形变范畴的态射。% P10-E0450
$\varphi$ 的{\it 微分
$d \varphi : T \mathcal{F} \to T \mathcal{G}$}
是把函子态射
$\overline{\varphi}: \overline{\mathcal{F}} \to \overline{\mathcal{G}}$
在 $A = k[\epsilon]$ 处取值所得的映射。
\end{definition}

\begin{lemma}
\label{formal-defos-lemma-k-linear-differential}
设 $\varphi : \mathcal{F} \to \mathcal{G}$ 是预形变范畴的态射。
假设 $\overline{\mathcal{F}}$ 和 $\overline{\mathcal{G}}$
都满足 (S2)。则
$d \varphi : T \mathcal{F} \to T \mathcal{G}$ 是 $k$-线性的。
\end{lemma}

\begin{proof}
在引理 \ref{formal-defos-lemma-tangent-space-vector-space} 的证明中，
已经看到 $\overline{\mathcal{F}}$ 和 $\overline{\mathcal{G}}$
满足引理 \ref{formal-defos-lemma-tangent-space-functor} 的假设。
因此，本引理由引理 \ref{formal-defos-lemma-morphism-tangent-spaces} 得出。
\end{proof}

\begin{remark}
\label{formal-defos-remark-tangent-space-cofibered-groupoid}
可以把切空间和微分的概念如下全局化到任意群胚余纤维范畴。
设 $\mathcal{F}$ 是 $\mathcal{C}_\Lambda$ 上的群胚余纤维范畴，
且 $x \in \Ob(\mathcal{F}(k))$。
如说明 \ref{formal-defos-remark-localize-cofibered-groupoid}，
得到预形变范畴 $\mathcal{F}_x$。定义
$$
T_x\mathcal{F} = T\mathcal{F}_x
$$
为 $\mathcal{F}$ 在 $x$ 处的{\it 切空间}。
若 $\varphi : \mathcal{F} \to \mathcal{G}$ 是
$\mathcal{C}_\Lambda$ 上群胚余纤维范畴的态射，且
$x \in \Ob(\mathcal{F}(k))$，则有诱导态射
$\varphi_x: \mathcal{F}_x \to \mathcal{G}_{\varphi(x)}$。
把 $\varphi$ 在 $x$ 处的{\it 微分
$d_x \varphi : T_x \mathcal{F} \to T_{\varphi(x)} \mathcal{G}$}
定义为映射
$d \varphi_x: T \mathcal{F}_x \to T \mathcal{G}_{\varphi(x)}$。
若 $\mathcal{F}$ 和 $\mathcal{G}$ 都满足 (S2)，
则所有这些切空间都具有自然的 $k$-向量空间结构，
而所有微分
$d_x \varphi : T_x \mathcal{F} \to T_{\varphi(x)} \mathcal{G}$
都是 $k$-线性的（使用引理
\ref{formal-defos-lemma-S1-S2-localize} 和
\ref{formal-defos-lemma-k-linear-differential}）。
\end{remark}

\noindent
在经典情形下，或当 $k/k'$ 是可分域扩张时，
以下观察没有实质内容，因为这时
$\text{Der}_\Lambda(k, k)$ 和 $\text{Der}_\Lambda(k, V)$ 都为零。
有典范识别
$$
\Mor_{\mathcal{C}_\Lambda}(k, k[\epsilon]) =
\text{Der}_\Lambda(k, k).
$$
确切地，对 $D \in \text{Der}_\Lambda(k, k)$，
令 $f_D : k \to k[\epsilon]$ 为映射
$a \mapsto a + D(a)\epsilon$。更一般地，
给定 $k$ 上的有限维向量空间 $V$，有
$$
\Mor_{\mathcal{C}_\Lambda}(k, k[V]) =
\text{Der}_\Lambda(k, V)
$$
并且仍用 $f_D$ 表示与导子 $D$ 相伴的映射。还有
$$
\Mor_{\mathcal{C}_\Lambda}(k[W], k[V]) =
\Hom_k(V, W) \oplus \text{Der}_\Lambda(k, V)
$$
% P10-E0451
其中 $(\varphi, D)$ 对应于映射
$f_{\varphi, D} : a + w \mapsto a + \varphi(w) + D(a)$。
有时把由导子 $D : k \to V$ 决定的 $k[V]$ 自同构写成
$f_{1, D} : a + v \to a + v + D(a)$。% P10-E0452
注意，$f_{1, D} \circ f_{1, D'} = f_{1, D + D'}$。

\medskip\noindent
设 $\mathcal{F}$ 是 $\mathcal{C}_\Lambda$ 上的预形变范畴，
并设 $x_0 \in \Ob(\mathcal{F}(k))$。
由上可得典范映射
$$
\gamma_V :
\text{Der}_\Lambda(k, V)
\longrightarrow
\overline{\mathcal{F}}(k[V])
$$
定义为 $D \mapsto f_{D, *}(x_0)$。此外，还有作用
$$
a_V : \text{Der}_\Lambda(k, V) \times \overline{\mathcal{F}}(k[V])
\longrightarrow
\overline{\mathcal{F}}(k[V])
$$
定义为 $(D, x) \mapsto f_{1, D, *}(x)$。
这两个映射相容；也就是说，由这些映射复合的直接计算可得
$f_{1, D, *}f_{D', *}x_0 = f_{D + D', *}x_0$。
注意，$\gamma_V$ 与 $a_V$ 不依赖于 $x_0$ 的选取，
因为 $x_0$ 在同构意义下唯一。

\begin{lemma}
\label{formal-defos-lemma-action-linear}
设 $\mathcal{F}$ 是 $\mathcal{C}_\Lambda$ 上的预形变范畴。
若 $\overline{\mathcal{F}}$ 满足 (S2)，则映射 $\gamma_V$
是 $k$-线性的，并且 $a_V(D, x) = x + \gamma_V(D)$。
\end{lemma}

\begin{proof}
在引理 \ref{formal-defos-lemma-tangent-space-vector-space} 的证明中，
已经看到函子 $V \mapsto \overline{\mathcal{F}}(k[V])$
把 $0$ 送到单元素集合，并把积送到积。
函子 $V \mapsto \text{Der}_\Lambda(k, V)$ 也如此。
因此，由引理 \ref{formal-defos-lemma-morphism-linear-functors}，
$\gamma_V$ 是线性的。
设 $D : k \to V$ 是 $\Lambda$-导子。
令 $D_1 : k \to V^{\oplus 2}$ 等于 $a \mapsto (D(a), 0)$。
则图
$$
\xymatrix{
k[V \times V] \ar[r]_{+} \ar[d]^{f_{1, D_1}} & k[V] \ar[d]^{f_{1, D}} \\
k[V \times V] \ar[r]^{+} & k[V]
}
$$
交换。展开定义，并利用
$\overline{F}(V \times V) = \overline{F}(V) \times \overline{F}(V)$，% P10-E0453
这意味着对所有 $x_1, x_2 \in \overline{F}(V)$，有
$a_D(x_1) + x_2 = a_D(x_1 + x_2)$。
因此，只需证明 $a_V(D, 0) = 0 + \gamma_V(D)$，
其中 $0 \in \overline{F}(V)$ 是零向量。
按定义，这是元素 $f_{0, *}(x_0)$。
由于 $f_D = f_{1, D} \circ f_0$，所需结论随即得出。
\end{proof}

\noindent
上述构造的一个特殊情形是映射
\begin{equation}
\label{formal-defos-equation-map}
\gamma : \text{Der}_\Lambda(k, k) \longrightarrow T\mathcal{F}
\end{equation}
以及作用
\begin{equation}
\label{formal-defos-equation-action}
a : \text{Der}_\Lambda(k, k) \times T\mathcal{F} \longrightarrow T\mathcal{F}
\end{equation}
它们对任意预形变范畴 $\mathcal{F}$ 都有定义。
注意，若 $\varphi : \mathcal{F} \to \mathcal{G}$ 是预形变范畴的态射，
则得到交换图
$$
\vcenter{
\xymatrix{
\text{Der}_\Lambda(k, k) \ar[r]_-\gamma \ar[rd]_\gamma &
T\mathcal{F} \ar[d]_{d\varphi} \\
& T\mathcal{G}
}
}
\quad\text{和}\quad
\vcenter{
\xymatrix{
\text{Der}_\Lambda(k, k) \times T\mathcal{F} \ar[r]_-a
\ar[d]_{1 \times d\varphi} &
T\mathcal{F} \ar[d]^{d\varphi} \\
\text{Der}_\Lambda(k, k) \times T\mathcal{G} \ar[r]^-a &
T\mathcal{G}
}
}
$$





\section{万有形式对象}
\label{formal-defos-section-versal-objects}

\noindent
万有形式对象的存在迫使 $\mathcal{F}$ 具有性质 (S1)。

\begin{lemma}
\label{formal-defos-lemma-versal-object-S1}
设 $\mathcal{F}$ 是预形变范畴。
假设 $\mathcal{F}$ 有万有形式对象。
则 $\mathcal{F}$ 满足 (S1)。
\end{lemma}

\begin{proof}
设 $\xi$ 是 $\mathcal{F}$ 的万有形式对象。设
$$
\xymatrix{
           & x_2 \ar[d] \\
x_1 \ar[r] & x
}
$$
是 $\mathcal{F}$ 中的图，并且 $x_2 \to x$
位于一个满环映射上方。由于自然态射
$\widehat{\mathcal{F}}|_{\mathcal{C}_\Lambda} \xrightarrow{\sim} \mathcal{F}$
是等价（见说明 \ref{formal-defos-remark-restrict-completion}），
也可以把这个图视为 $\widehat{\mathcal{F}}$ 中的图。
由引理 \ref{formal-defos-lemma-versal-object-quasi-initial}，存在态射
$\xi \to x_1$；所以由说明 \ref{formal-defos-remark-versal-object}，
还得到态射 $\xi \to x_2$，使图
$$
\xymatrix{
\xi \ar[r] \ar[d]          & x_2 \ar[d] \\
x_1 \ar[r] & x
}
$$
交换。若 $x_1 \to x$ 与 $x_2 \to x$ 分别位于环映射
$A_1 \to A$ 与 $A_2 \to A$ 上方，
则把 $\xi$ 推前到 $A_1 \times_A A_2$，
便得到 (S1) 所要求的对象 $y$。
\end{proof}

\noindent
当余纤维范畴满足 (S1) 和 (S2) 时，
可以如下刻画万有形式对象。

\begin{lemma}
\label{formal-defos-lemma-versal-criterion}
设 $\mathcal{F}$ 是满足 (S1) 和 (S2) 的预形变范畴。
设 $\xi$ 是 $\mathcal{F}$ 的形式对象，对应于
$\underline{\xi} : \underline{R}|_{\mathcal{C}_\Lambda} \to \mathcal{F}$；
见说明 \ref{formal-defos-remark-formal-objects-yoneda}。
则 $\xi$ 是万有的，当且仅当以下两个条件成立：
\begin{enumerate}
\item 切空间上的映射
$d\underline{\xi} : T\underline{R}|_{\mathcal{C}_\Lambda} \to T\mathcal{F}$
是满射；并且
\item 给定 $\widehat{\mathcal{F}}$ 中的图
$$
\vcenter{
\xymatrix{
            &  y \ar[d] \\
\xi \ar[r]  &  x
}
}
\quad\text{位于下图上方}\quad
\vcenter{
\xymatrix{
         &   B  \ar[d]^{f} \\
R \ar[r] &   A
}
}
$$
其中底图属于 $\widehat{\mathcal{C}}_\Lambda$，
且 $B \to A$ 是阿廷环的小扩张，则存在环映射
$R \to B$，使图
$$
\xymatrix{
         &   B  \ar[d]^{f} \\
R \ar[ur] \ar[r] &   A
}
$$
交换。
\end{enumerate}
\end{lemma}

\begin{proof}
若 $\xi$ 是万有的，则 (1) 由引理
\ref{formal-defos-lemma-smooth-morphism-essentially-surjective} 成立，
(2) 由说明 \ref{formal-defos-remark-versal-object} 成立。
反过来，假设 (1) 和 (2) 成立。由说明
\ref{formal-defos-remark-versal-object}，必须证明：给定 (2) 中那样的
$\widehat{\mathcal{F}}$ 中的图，存在 $\xi \to y$，使图
$$
\xymatrix{
            &  y \ar[d] \\
\xi \ar[ur] \ar[r]  &  x
}
$$
交换。设 $b : R \to B$ 是 (2) 所保证的映射。
记 $y' = b_*\xi$，并选取给定态射 $\xi \to x$
在 $R \to B \to A$ 上方的分解
$\xi \to y' \to x$。由 (S1)，得到交换图
$$
\vcenter{
\xymatrix{
z  \ar[r] \ar[d]          &  y \ar[d] \\
y' \ar[r]  &  x
}
}
\quad\text{位于下图上方}\quad
\vcenter{
\xymatrix{
B \times_A B \ar[d] \ar[r] &   B  \ar[d]^{f} \\
B \ar[r]^{f} &   A .
}
}
$$
令 $I = \Ker(f)$。令
$\overline{g} : B \times_A B \to k[I]$
为环映射 $(u, v) \mapsto \overline{u} \oplus (v - u)$；
参见引理 \ref{formal-defos-lemma-lifting-along-small-extension}。
由 (1)，存在位于某个环映射
$i : R \to k[\epsilon]$ 上方的态射
$\xi \to \overline{g}_*z$。
选取阿廷商 $b_1 : R \to B_1$，使
$b : R \to B$ 和 $i : R \to k[\epsilon]$ 都经
$R \to B_1$ 分解；也就是说，得到
$h : B_1 \to B$ 和 $i' : B_1 \to k[\epsilon]$。
选取推前 $y_1 = b_{1, *}\xi$，选取 $\xi \to y'$
在 $R \to B_1 \to B$ 上方的分解
$\xi \to y_1 \to y'$，并选取
$\xi \to \overline{g}_*z$ 在
$R \to B_1 \to k[\epsilon]$ 上方的分解
$\xi \to y_1 \to \overline{g}_*z$。
再次应用 (S1)，得到
$$
\vcenter{
\xymatrix{
z_1  \ar[r] \ar[d] & z \ar[r] \ar[d] & y \ar[d] \\
y_1 \ar[r] & y' \ar[r] &  x
}
}
\quad\text{位于下图上方}\quad
\vcenter{
\xymatrix{
B_1 \times_A B \ar[d] \ar[r] & B \times_A B \ar[r] \ar[d] & B \ar[d]^{f} \\
B_1 \ar[r]  & B \ar[r] &  A .
}
}
$$
注意，引理 \ref{formal-defos-lemma-lifting-along-small-extension}
中（用 $h$ 定义的）映射
$g : B_1 \times_A B \to k[I]$，
是 $B_1 \times_A B \to B \times_A B$
与上述映射 $\overline{g}$ 的复合。
由构造，存在态射
$y_1 \to g_*z_1 \cong \overline{g}_*z$！
因此，引理 \ref{formal-defos-lemma-lifting-along-small-extension}
可应用于上述图中的外侧矩形，给出态射
$y_1 \to y$；再前复合 $\xi \to y_1$，
便得到所需态射 $\xi \to y$。
\end{proof}

\noindent
若 $\mathcal{F}$ 具有性质 (S1)，则“存在提升的最大商”存在。
精确表述如下。

\begin{lemma}
\label{formal-defos-lemma-largest-closed-where-lift}
设 $\mathcal{F}$ 是 $\mathcal{C}_\Lambda$ 上满足 (S1)
的群胚余纤维范畴。设 $B \to A$ 是
$\mathcal{C}_\Lambda$ 中的满射，其核 $I$ 被
$\mathfrak m_B$ 零化。设 $x \in \mathcal{F}(A)$。
理想集合
$$
\mathcal{J} = \{ J \subset I \mid
\text{存在位于 }B/J \to A\text{ 上方的 }y \to x\}
$$
有最小元。% P10-E0455
\end{lemma}

\begin{proof}
注意，$\mathcal{J}$ 非空，因为 $I \in \mathcal{J}$。
此外，若 $J \in \mathcal{J}$ 且 $J \subset J' \subset I$，
则 $J' \in \mathcal{J}$，因为可以把对象 $y$ 推前为
$B/J'$ 上方的对象 $y'$。设 $J$ 和 $K$ 是上述集合的元素。
我们断言 $J \cap K \in \mathcal{J}$，这将证明引理。
因为 $I$ 是 $k$-向量空间，可以找到理想
$J \subset J' \subset I$，使 $J \cap K = J' \cap K$ 且
$J' + K = I$。由上可用 $J'$ 代替 $J$，并假设
$J + K = I$。此时
$$
A/(J \cap K) = A/J \times_{A/I} A/K.
$$
% P10-E0454
因此，由 (S1) 以及假设存在于 $A/J$ 和 $A/K$ 上方的元素，
可知存在映到 $x$ 的元素
$z \in \mathcal{F}(A/(J \cap K))$。
\end{proof}

\noindent
稍后将改进下述结果。

\begin{lemma}
\label{formal-defos-lemma-versal-object-existence}
设 $\mathcal{F}$ 是 $\mathcal{C}_\Lambda$ 上的群胚余纤维范畴。
假设以下条件成立：
\begin{enumerate}
\item $\mathcal{F}$ 是预形变范畴。
\item $\mathcal{F}$ 满足 (S1)。
\item $\mathcal{F}$ 满足 (S2)。
\item $\dim_k T\mathcal{F}$ 有限。
\end{enumerate}
则 $\mathcal{F}$ 有万有形式对象。
\end{lemma}

\begin{proof}
假设 (1)、(2)、(3) 和 (4) 成立。
选取对象 $R \in \Ob(\widehat{\mathcal{C}}_\Lambda)$，
使 $\underline{R}|_{\mathcal{C}_\Lambda}$ 光滑；
见引理 \ref{formal-defos-lemma-exists-smooth}。
令 $r = \dim_k T\mathcal{F}$，并置
$S = R[[X_1, \ldots, X_r]]$。

\medskip\noindent
我们将对 $n \geq 2$ 归纳构造二元组
$(J_n, f_{n - 1} : \xi_n \to \xi_{n - 1})$，
其中 $J_n \subset S$ 构成递减理想序列，% P10-E0456
而 $f_{n - 1} : \xi_n \to \xi_{n - 1}$ 是
$\mathcal{F}$ 中位于投影
$S/J_n \to S/J_{n - 1}$ 上方的态射。

\medskip\noindent
第 1 步。令 $J_1 = \mathfrak m_S$。
令 $\xi_1$ 为 $k = S/J_1 = S/\mathfrak m_S$ 上方的
$\mathcal{F}$ 的对象；它在唯一同构意义下唯一。% P10-E0457

\medskip\noindent
第 2 步。令
$J_2 = \mathfrak m_S^2 + \mathfrak{m}_R S$。
则 $S/J_2 = k[V]$，其中
$V = kX_1 \oplus \ldots \oplus kX_r$。
由 $\overline{\mathcal{F}}$ 的 (S2)，得到双射
$$
\overline{\mathcal{F}}(S/J_2)
\longrightarrow
T\mathcal{F} \otimes_k V,
$$
见引理 \ref{formal-defos-lemma-S1-S2-associated-functor} 和
\ref{formal-defos-lemma-tangent-space-vector-space}。
选取 $T\mathcal{F}$ 的一组基
$\theta_1, \ldots, \theta_r$，并置
$\xi_2 = \sum \theta_i \otimes X_i \in \Ob(\mathcal{F}(S/J_2))$。% P10-E0458
这样选取的要点在于
$$
d\xi_2 :
\Mor_{\mathcal{C}_\Lambda}(S/J_2, k[\epsilon])
\longrightarrow
T\mathcal{F}
$$
是满射。令 $f_1 : \xi_2 \to \xi_1$ 为唯一态射。

\medskip\noindent
归纳步骤。假设对某个 $n \geq 2$，已经构造了
$(J_n, f_{n - 1} : \xi_n \to \xi_{n - 1})$。
满足以下条件的理想 $J \subset S$ 的集合有最小元
$J_{n + 1}$：
(a) $\mathfrak m_S J_n \subset J \subset J_n$；并且
(b) 存在位于 $S/J \to S/J_n$ 上方的态射
$\xi_{n + 1} \to \xi_n$；见引理
\ref{formal-defos-lemma-largest-closed-where-lift}。
令 $f_n : \xi_{n + 1} \to \xi_n$ 为
$\mathcal{F}$ 中位于 $S/J_{n + 1} \to S/J_n$ 上方的任意态射。

\medskip\noindent
令 $J = \bigcap J_n$，$\overline{S} = S/J$，并令
$\overline{J}_n = J_n/J$。由引理
\ref{formal-defos-lemma-m-adic-topology}，理想序列
$(\overline{J}_n)$ 在 $\overline{S}$ 上诱导
$\mathfrak m_{\overline{S}}$-进拓扑。
由于 $(\xi_n, f_n)$ 是
$\widehat{\mathcal{F}}_\mathcal{I}(\overline{S})$ 的对象，
其中 $\mathcal{I}$ 是 $\overline{S}$ 的滤过
$(\overline{J}_n)$，可见 $(\xi_n, f_n)$ 诱导
$\widehat{\mathcal{F}}(\overline{S})$ 的对象 $\xi$；
见引理 \ref{formal-defos-lemma-formal-objects-different-filtration}。% P10-E0459

\medskip\noindent
我们证明 $\xi$ 是万有的。为证明万有性，只需核验引理
\ref{formal-defos-lemma-versal-criterion} 的条件 (1) 与 (2)。
条件 (1) 来自上述第 2 步中对 $\xi_2$ 的选取。
假设给定 $\widehat{\mathcal{F}}$ 中的图
$$
\vcenter{
\xymatrix{
            &  y \ar[d] \\
\xi \ar[r]  &  x
}
}
\quad\text{位于下图上方}\quad
\vcenter{
\xymatrix{
         &   B  \ar[d]^{f} \\
\overline{S} \ar[r] &   A
}
}
$$
其中底图属于 $\widehat{\mathcal{C}}_\Lambda$，
且 $f: B \to A$ 是阿廷环的小扩张。
必须证明存在使右图补全的映射 $\overline{S} \to B$。
选取 $n$，使 $\overline{S} \to A$ 经
$\overline{S} \to S/J_n$ 分解。
这是可能的，因为如上所见，序列 $(\overline{J}_n)$
诱导 $\mathfrak m_{\overline{S}}$-进拓扑。
$\xi$ 沿 $\overline{S} \to S/J_n$ 的推前为 $\xi_n$。
可以把 $\xi \to x$ 分解为 $\xi \to \xi_n \to x$，
从而得到 $\mathcal{F}$ 中的图
$$
\vcenter{
\xymatrix{
            &  y \ar[d] \\
\xi_n \ar[r]  &  x
}
}
\quad\text{位于下图上方}\quad
\vcenter{
\xymatrix{
         &   B  \ar[d]^{f} \\
S/J_n \ar[r] &   A .
}
}
$$
为核验引理 \ref{formal-defos-lemma-versal-criterion} 的条件 (2)，
只需补全图
$$
\xymatrix{
S/J_{n + 1} \ar[d] \ar@{-->}[r] & B \ar[d]^{f} \\
S/J_n   \ar[r] & A
}
$$
或者等价地，补全图
$$
\xymatrix{
  &  S/J_n \times_A B \ar[d]^{p_1} \\
S/J_{n + 1} \ar@{-->}[ur] \ar[r] &  S/J_n.
}
$$
若 $p_1$ 有截面，则已经完成。否则，由引理
\ref{formal-defos-lemma-fiber-product-CLambda} (2)，
$p_1$ 是小扩张；所以由引理
\ref{formal-defos-lemma-essential-surjection} (4)，
$p_1$ 是本质满射。回忆
$S = R[[X_1, \ldots, X_r]]$，且选取 $R$ 时使
$\underline{R}|_{\mathcal{C}_\Lambda}$ 光滑。
因此，存在映射 $h : R \to B$，提升映射
$R \to S \to S/J_n \to A$。
由幂级数环的泛性质，存在 $R$-代数映射
$h : S = R[[X_1, \ldots, X_r]] \to B$，% P10-E0460
提升给定映射 $S \to S/J_n \to A$。
这诱导映射 $g: S \to S/J_n \times_A B$，
使图中的实线方块
$$
\xymatrix{
S \ar[d] \ar[r]_-g  &  S/J_n \times_A B \ar[d]^{p_1} \\
S/J_{n + 1} \ar@{-->}[ur] \ar[r] &  S/J_n
}
$$
交换。因为 $p_1$ 是本质满射，所以 $g$ 是满射。
我们断言 $S$ 的理想 $K = \Ker(g)$ 满足上述归纳步骤构造
$J_{n + 1}$ 时的条件 (a) 与 (b)。
确实，$K \subset J_n$ 显然，而因为 $p_1$ 是小扩张，
$\mathfrak m_SJ_n \subset K$；这证明了 (a)。
把 (S1) 应用于
$$
\xymatrix{
            &  y \ar[d] \\
\xi_n \ar[r]  &  x,
}
$$
可得 $\xi_n$ 到 $S/K \cong S/J_n \times_A B$ 的提升，
所以 (b) 成立。由于 $J_{n + 1}$ 是具有性质 (a) 与 (b)
的最小理想，这蕴含 $J_{n + 1} \subset K$。
因此，所需映射
$S/J_{n+1} \to S/K \cong S/J_n \times_A B$ 存在。
\end{proof}

\begin{remark}
\label{formal-defos-remark-compare-schlessinger-H3-pre}
设 $F : \mathcal{C}_\Lambda \to \textit{Sets}$ 是满足 (S1)
和 (S2) 的预形变函子。条件 $\dim_k TF < \infty$
正是 Schlessinger 论文中的条件 (H3)。
回忆，(S1) 与 (S2) 分别对应于条件 (H1) 与 (H2)；
见说明 \ref{formal-defos-remark-compare-S1-S2-schlessinger}。
因此，引理 \ref{formal-defos-lemma-versal-object-existence} 告诉我们：
$$
(H1) + (H2) + (H3)
\Rightarrow
\text{存在万有形式对象}
$$
对预形变函子成立。我们将在说明
\ref{formal-defos-remark-compare-schlessinger-H3} 中说明它与包络的联系。
\end{remark}





\section{极小的万有形式对象}
\label{formal-defos-section-minimal-versal}

\noindent
我们稍作准备，以理解万有形式对象的（非）唯一性。
结果表明，若预形变范畴有万有形式对象，
则它有极小的万有形式对象，并且任意两个这样的对象都同构。
此外，在用幂级数扩张替换基环的意义下，
所有万有形式对象都“差不多”相同。

\medskip\noindent
设 $\mathcal{F}$ 是 $\mathcal{C}_\Lambda$ 上的群胚余纤维范畴。
对 $\mathcal{F}$ 中位于
$A \in \Ob(\mathcal{C}_\Lambda)$ 上方的每个对象 $x$，
考虑范畴 $\mathcal{S}_x$，其对象为
$$
\Ob(\mathcal{S}_x) =
\{x' \to x \mid x' \to x\text{ 所在的子环为 }A' \subset A\}
$$
而态射是 $x$ 上方的态射。
对 $\mathcal{F}$ 中位于
$\mathcal{C}_\Lambda$ 的 $f : B \to A$ 上方的每个
$y \to x$，有如下定义的函子
$f_* : \mathcal{S}_y \to \mathcal{S}_x$：
给定 $B' \subset B$ 上方的 $y' \to y$，令
$A' = f(B')$，并令位于 $B' \to f(B')$ 上方的
$y' \to x'$ 是 $y'$ 的推前。% P10-E0461
由群胚余纤维范畴的公理，得到位于
$f(B') \to A$ 上方的唯一态射 $x' \to x$，使图
$$
\xymatrix{
y' \ar[d] \ar[r] & x' \ar[d] \\
y \ar[r] & x
}
$$
交换。于是 $x' \to x$ 是 $\mathcal{S}_x$ 的对象。
若 $\mathcal{S}_x$ 中的任意态射
$(x'_1 \to x) \to (x' \to x)$ 都是同构，
即 $x'$ 与 $x'_1$ 定义在 $A$ 的同一子环上，
则称 $\mathcal{S}_x$ 的对象 $x' \to x$ 是{\it 极小的}。
因为 $A$ 作为 $\Lambda$-模具有有限长度，
所以极小对象总是存在。

\begin{lemma}
\label{formal-defos-lemma-smallest-where-descends}
设 $\mathcal{F}$ 是 $\mathcal{C}_\Lambda$ 上满足 (S1)
的群胚余纤维范畴。
\begin{enumerate}
\item 对 $\mathcal{F}$ 中的 $y \to x$，
$\mathcal{S}_y$ 的极小对象映到 $\mathcal{S}_x$ 的极小对象。
\item 对 $\mathcal{F}$ 中位于
$\mathcal{C}_\Lambda$ 的满射 $f : B \to A$ 上方的
$y \to x$，$\mathcal{S}_x$ 的每个极小对象都是
$\mathcal{S}_y$ 的某个极小对象的像。
\end{enumerate}
\end{lemma}

\begin{proof}
(1) 的证明。设 $y \to x$ 位于 $f : B \to A$ 上方。
设位于 $B' \subset B$ 上方的 $y' \to y$ 是
$\mathcal{S}_y$ 的极小对象。令
$$
\vcenter{
\xymatrix{
y' \ar[d] \ar[r] & x' \ar[d] \\
y \ar[r] & x
}
}
\quad\text{位于下图上方}\quad
\vcenter{
\xymatrix{
B' \ar[d] \ar[r] & f(B') \ar[d] \\
B \ar[r] & A
}
}
$$
如上述 $f_*$ 的构造所给。假设
$(x'' \to x) \to (x' \to x)$ 是 $\mathcal{S}_x$ 的态射，
且 $x'' \to x'$ 位于 $A'' \subset f(B')$ 上方。
由 (S1)，存在位于
$B' \times_{f(B')} A'' \to B'$ 上方的 $y'' \to y'$。
由于 $y' \to y$ 极小，可知
$B' \times_{f(B')} A'' \to B'$ 是同构，
这蕴含 $A'' = f(B')$；即 $x' \to x$ 极小。

\medskip\noindent
(2) 的证明。假设 $f : B \to A$ 满射，且
$y \to x$ 位于 $f$ 上方。
设 $x' \to x$ 是 $\mathcal{S}_x$ 中位于
$A' \subset A$ 上方的极小对象。
由 (S1)，存在位于
$B' = f^{-1}(A') = B \times_A A' \to B$ 上方的
$y' \to y$，它在 $\mathcal{S}_x$ 中的像是 $x' \to x$。
所以 $f_*(y' \to y) = x' \to x$。
选取 $\mathcal{S}_y$ 中的态射
$(y'' \to y) \to (y' \to y)$，其中 $y'' \to y$ 是极小对象
（由引理之前关于长度的说明，这是可能的）。
于是 $f_*(y'' \to y)$ 是 $\mathcal{S}_x$ 中映到
$x' \to x$ 的对象（由 $f_*$ 的函子性）；
由于 $x' \to x$ 极小，它与 $x' \to x$ 同构。
\end{proof}

\begin{lemma}
\label{formal-defos-lemma-smallest-where-descends-versal}
设 $\mathcal{F}$ 是 $\mathcal{C}_\Lambda$ 上满足 (S1)
的群胚余纤维范畴。设 $\xi$ 是位于 $R$ 上方的
$\mathcal{F}$ 的万有形式对象。
存在位于 $R' \subset R$ 上方的态射 $\xi' \to \xi$，
具有以下极小性性质：
\begin{enumerate}
\item 对每个 $f : R \to A$，其中
$A \in \Ob(\mathcal{C}_\Lambda)$，推前图
$$
\vcenter{
\xymatrix{
\xi' \ar[d] \ar[r] & x' \ar[d] \\
\xi \ar[r] & x
}
}
\quad\text{位于下图上方}\quad
\vcenter{
\xymatrix{
R' \ar[d] \ar[r] & f(R') \ar[d] \\
R \ar[r] & A
}
}
$$
产生 $\mathcal{S}_x$ 的极小对象 $x' \to x$；并且
\item 对任意形式对象态射 $\xi'' \to \xi'$，
相应环映射 $R'' \to R'$ 都是满射。
\end{enumerate}
\end{lemma}

\begin{proof}
写成 $\xi = (R, \xi_n, f_n)$。令 $R'_1 = k$ 且
$\xi'_1 = \xi_1$。假设对 $m \leq n$，已经构造了
$\mathcal{S}_{\xi_m}$ 中位于
$R'_m \subset R/\mathfrak m_R^m$ 上方的极小对象
$\xi'_m \to \xi_m$；并且对 $m \leq n - 1$，
已经构造了与 $f_m$ 相容的态射
$f'_m : \xi'_{m + 1} \to \xi'_m$。
由引理 \ref{formal-defos-lemma-smallest-where-descends} (2)，
存在位于
$R'_{n + 1} \subset R/\mathfrak m_R^{n + 1}$ 上方的极小对象
$\xi'_{n + 1} \to \xi_{n + 1}$，其像是位于
$R'_n \subset R/\mathfrak m_R^n$ 上方的
$\xi'_n \to \xi_n$。这按构造给出交换图
$$
\xymatrix{
\xi'_{n + 1} \ar[r]_{f'_n} \ar[d] & \xi'_n \ar[d] \\
\xi_{n + 1} \ar[r]^{f_n} & \xi_n
}
$$
此外，环映射 $R'_{n + 1} \to R'_n$ 是满射。
令 $R' = \lim_n R'_n$。则 $R' \to R$ 是单射。

\medskip\noindent
不过，先验上并不清楚 $R'$ 是否诺特。
为证明这一点，利用 $\xi$ 的万有性。
确切地，万有性蕴含 $\widehat{\mathcal{F}}$ 中存在态射
$\xi \to \xi'_n$；见引理
\ref{formal-defos-lemma-versal-object-quasi-initial}。
相应映射 $R \to R'_n$ 必为满射
（因为 $\xi'_n \to \xi_n$ 在 $\mathcal{S}_{\xi_n}$ 中极小）。
因此余切空间的维数有界，而引理
\ref{formal-defos-lemma-limit-artinian} 蕴含 $R'$ 诺特，
即它是 $\widehat{\mathcal{C}}_\Lambda$ 的对象。
由引理 \ref{formal-defos-lemma-formal-objects-different-filtration}
（再加上引理 \ref{formal-defos-lemma-limit-artinian} 关于滤过的结果），
元素序列 $\xi'_n$ 定义 $R'$ 上方的形式对象 $\xi'$，
并且有映射 $\xi' \to \xi$。

\medskip\noindent
按构造，对每个 $n$，(1) 对
$R \to R/\mathfrak m_R^n$ 成立。
因为 (1) 中的每个 $R \to A$ 都经
$R \to R/\mathfrak m_R^n \to A$ 分解，
所以由引理 \ref{formal-defos-lemma-smallest-where-descends} (1)，
$\xi'_n \to \xi_n$ 在
$R'_n \to R/\mathfrak m_R^n$ 上方的极小性蕴含：
(1) 对位于 $f(R) \subset A$ 上方的 $x' \to x$ 成立。% P10-E0462

\medskip\noindent
若 (2) 中的 $R'' \to R'$ 不是满射，
则对某个 $n$，$R'' \to R' \to R'_n$ 不是满射，
于是 $\xi'_n \to \xi_n$ 不会是极小的，矛盾。
这个矛盾证明了 (2)。
\end{proof}

\begin{lemma}
\label{formal-defos-lemma-descends-versal}
设 $\mathcal{F}$ 是 $\mathcal{C}_\Lambda$ 上满足 (S1)
的群胚余纤维范畴。设 $\xi$ 是位于 $R$ 上方的
$\mathcal{F}$ 的万有形式对象。
设 $\xi' \to \xi$ 是引理
\ref{formal-defos-lemma-smallest-where-descends-versal} 所构造的、
位于 $R' \subset R$ 上方的形式对象态射。则
$$
R \cong R'[[x_1, \ldots, x_r]]
$$
是 $R'$ 上的幂级数环。
此外，$\xi'$ 也是万有形式对象。
\end{lemma}

\begin{proof}
由引理 \ref{formal-defos-lemma-versal-object-quasi-initial}，
存在态射 $\xi \to \xi'$。
由引理 \ref{formal-defos-lemma-smallest-where-descends-versal}，
相应映射 $f : R \to R'$ 诱导满射
$f|_{R'} : R' \to R'$。
由《代数》引理
\ref{algebra-lemma-surjective-endo-noetherian-ring-is-iso}，
这是同构。因此 $I = \Ker(f)$ 是 $R$ 的理想，
并且 $R = R' \oplus I$。
选取 $x_1, \ldots, x_n \in I$，使其构成
$I/\mathfrak m_RI$ 的一组基。考虑把 $X_i$ 映到 $x_i$ 的映射
$S = R'[[X_1, \ldots, X_r]] \to R$。% P10-E0463
对每个 $n \geq 1$，得到阿廷 $R'$-代数的满射
$B = S/\mathfrak m_S^n \to R/\mathfrak m_R^n = A$。
用 $y \in \Ob(\mathcal{F}(B)$、相应地
$x \in \Ob(\mathcal{F}(A))$ 表示 $\xi'$ 分别沿
$R' \to S \to B$、相应地沿 $R' \to S \to A$ 的推前。% P10-E0464
注意，$x$ 也是 $\xi$ 沿 $R \to A$ 的推前，
因为 $\xi$ 是 $\xi'$ 沿 $R' \to R$ 的推前。
因此有实线图
$$
\vcenter{
\xymatrix{
& y \ar[d] \\
\xi \ar[r] \ar@{..>}[ru] & x
}
}
\quad\text{位于下图上方}\quad
\vcenter{
\xymatrix{
& S/\mathfrak m_S^n \ar[d] \\
R \ar[r] \ar@{..>}[ru] & R/\mathfrak m_R^n
}
}
$$
因为 $\xi$ 万有，使用说明 \ref{formal-defos-remark-versal-object}，
得到补全这些图的虚线箭头。
特别地，映射
$S/\mathfrak m_S^n \to R/\mathfrak m_R^n$
有截面 $h_n : R/\mathfrak m_R^n \to S/\mathfrak m_S^n$。
由引理 \ref{formal-defos-lemma-power-series}，$S \to R$ 是同构。

\medskip\noindent
由于 $\xi$ 是 $\xi'$ 沿 $R' \to R$ 的推前，
说明 \ref{formal-defos-remark-formal-objects-yoneda-map} 给出交换图
$$
\xymatrix{
\underline{R}|_{\mathcal{C}_\Lambda} \ar[rr] \ar[rd]_{\underline{\xi}} & &
\underline{R'}|_{\mathcal{C}_\Lambda} \ar[ld]^{\underline{\xi'}} \\
& \mathcal{F}
}
$$
因为 $R' \to R$ 有左逆（即 $R \to R/I = R'$），
可知
$\underline{R}|_{\mathcal{C}_\Lambda} \to
\underline{R'}|_{\mathcal{C}_\Lambda}$ 本质满。
因此，由引理 \ref{formal-defos-lemma-smooth-properties}，
$\underline{\xi'}$ 光滑；即 $\xi'$ 是万有形式对象。
\end{proof}

\noindent
受前述引理启发，作如下定义。

\begin{definition}
\label{formal-defos-definition-minimal-versal}
设 $\mathcal{F}$ 是预形变范畴。
若对任意形式对象态射 $\xi' \to \xi$，
底层环映射都是满射，则称 $\mathcal{F}$ 的万有形式对象
$\xi$ 是{\it 极小的}\footnote{这可能不是标准术语。
许多作者把这一概念与切空间的性质联系起来。
我们将在第 \ref{formal-defos-section-miniversal-objects-existence} 节说明这种联系。}。
有时把极小的万有形式对象称为{\it 极小万有的}。
\end{definition}

\noindent
本节的工作说明这个定义是合理的。
首先，万有形式对象的存在蕴含 $\mathcal{F}$ 具有 (S1)。
其次，前述引理表明极小的万有形式对象存在。
最后，任意两个极小的万有形式对象都同构。
下面汇总这些结果（并给出详细证明）。

\begin{lemma}
\label{formal-defos-lemma-minimal-versal}
设 $\mathcal{F}$ 是具有万有形式对象的预形变范畴。则
\begin{enumerate}
\item $\mathcal{F}$ 有极小的万有形式对象；
\item 极小的万有形式对象在同构意义下唯一；并且
\item 任意万有形式对象都是某个极小的万有形式对象
沿幂级数环扩张的推前。
\end{enumerate}
\end{lemma}

\begin{proof}
假设 $\mathcal{F}$ 有位于 $R$ 上方的万有形式对象 $\xi$。
则它满足 (S1)；见引理 \ref{formal-defos-lemma-versal-object-S1}。
设位于 $R' \subset R$ 上方的 $\xi' \to \xi$
是引理 \ref{formal-defos-lemma-smallest-where-descends-versal}
所构造的任一态射。
由引理 \ref{formal-defos-lemma-descends-versal}，$\xi'$ 是万有的；
因此按构造，它是极小的万有形式对象。这证明了 (1)。
此外，$R \cong R'[[x_1, \ldots, x_n]]$，这证明了 (3)。

\medskip\noindent
假设 $\xi_i/R_i$ 是两个极小的万有形式对象。
由引理 \ref{formal-defos-lemma-versal-object-quasi-initial}，
存在态射 $\xi_1 \to \xi_2$ 和 $\xi_2 \to \xi_1$。
由极小性，相应环映射
$f : R_1 \to R_2$ 和 $g : R_2 \to R_1$ 都是满射。
因此，由《代数》引理
\ref{algebra-lemma-surjective-endo-noetherian-ring-is-iso}，
复合 $g \circ f : R_1 \to R_1$ 和
$f \circ g : R_2 \to R_2$ 都是同构。
从而 $f$ 与 $g$ 是同构，于是映射
$\xi_1 \to \xi_2$ 和 $\xi_2 \to \xi_1$ 也是同构
（因为由引理 \ref{formal-defos-lemma-completion-cofibred}，
$\widehat{\mathcal{F}}$ 余纤维于群胚）。
这证明了 (2)，并完成引理的证明。
\end{proof}








\section{极小万有形式对象与切空间}
\label{formal-defos-section-miniversal-objects-existence}

\noindent
定义 \ref{formal-defos-definition-minimal-versal} 所引入的一般极小性概念，
有时可以从切空间上的表现推导出来。
设 $\xi$ 是预形变范畴 $\mathcal{F}$ 的形式对象，并设
$\underline{\xi} : \underline{R}|_{\mathcal{C}_\Lambda} \to \mathcal{F}$
为相应态射。于是可以考虑条件% P10-E0465
\begin{equation}
\label{formal-defos-equation-bijective}
d\underline{\xi} : \text{Der}_\Lambda(R, k) \to T\mathcal{F}
\text{ 是双射}
\end{equation}
以及条件
\begin{equation}
\label{formal-defos-equation-bijective-orbits}
d\underline{\xi} : \text{Der}_\Lambda(R, k) \to T\mathcal{F}
\text{ 在 }\text{Der}_\Lambda(k, k)\text{-轨道上是双射。}
\end{equation}
这里使用了例 \ref{formal-defos-example-tangent-space-prorepresentable-functor}
中的识别
$T\underline{R}|_{\mathcal{C}_\Lambda} = \text{Der}_\Lambda(R, k)$，
以及导子在切空间上的作用（\ref{formal-defos-equation-action}）。
若 $k' \subset k$ 可分，则
$\text{Der}_\Lambda(k, k) = 0$，两个条件等价。
事实表明，在条件 (S2) 成立时，万有形式对象极小，
当且仅当 $\underline{\xi}$ 满足
（\ref{formal-defos-equation-bijective-orbits}）。
此外，若 $\underline{\xi}$ 满足（\ref{formal-defos-equation-bijective}），
则 $\mathcal{F}$ 满足 (S2)。

\begin{lemma}
\label{formal-defos-lemma-miniversal-object-unique}
设 $\mathcal{F}$ 是预形变范畴。
设 $\xi$ 是 $\mathcal{F}$ 的万有形式对象，
且满足（\ref{formal-defos-equation-bijective-orbits}）。
则 $\xi$ 是极小的万有形式对象。
特别地，这样的 $\xi$ 在同构意义下唯一。
\end{lemma}

\begin{proof}
若 $\xi$ 不极小，则存在位于 $R' \to R$ 上方的态射
$\xi' \to \xi$，使得
$R = R'[[x_1, \ldots, x_n]]$ 且 $n > 0$；
见引理 \ref{formal-defos-lemma-minimal-versal}。
因此 $d\underline{\xi}$ 分解为
$$
\text{Der}_\Lambda(R, k) \to
\text{Der}_\Lambda(R', k) \to T\mathcal{F}
$$
而（\ref{formal-defos-equation-bijective-orbits}）不可能成立，
因为
$D : f \mapsto \partial/\partial x_1(f) \bmod \mathfrak m_R$
属于第一个箭头的核，却不在
$\text{Der}_\Lambda(k, k) \to \text{Der}_\Lambda(R, k)$ 的像中。
\end{proof}

\begin{lemma}
\label{formal-defos-lemma-miniversal-object-existence-1}
设 $\mathcal{F}$ 是预形变范畴。
设 $\xi$ 是 $\mathcal{F}$ 的万有形式对象，
且满足（\ref{formal-defos-equation-bijective}）。则
\begin{enumerate}
\item $\mathcal{F}$ 满足 (S1)。
\item $\mathcal{F}$ 满足 (S2)。
\item $\dim_k T\mathcal{F}$ 有限。
\end{enumerate}
\end{lemma}

\begin{proof}
由引理 \ref{formal-defos-lemma-versal-object-S1}，条件 (S1) 成立。
由于 (S1) 成立，(S2) 的第一部分成立。设
$$
\vcenter{
\xymatrix{
y \ar[r]_c \ar[d]_a & x_\epsilon \ar[d]^e \\
x \ar[r]^d          & x_0
}
}
\quad\text{和}\quad
\vcenter{
\xymatrix{
y' \ar[r]_{c'} \ar[d]_{a'} & x_\epsilon \ar[d]^e \\
x \ar[r]^d                 & x_0
}
}
\quad\text{位于下图上方}\quad
\vcenter{
\xymatrix{
A \times_k k[\epsilon] \ar[r] \ar[d] & k[\epsilon] \ar[d] \\
A  \ar[r] & k
}
}
$$
为 (S2) 第二部分中的图。如上，可以找到态射
$b : \xi \to y$ 和 $b' : \xi \to y'$，使图
$$
\xymatrix{
\xi \ar[r]^{b'} \ar[d]_b          & y' \ar[d]^{a'} \\
y \ar[r]^{a} & x
}
$$
交换。用 $p : \mathcal{F} \to \mathcal{C}_\Lambda$
表示结构态射。设 $\widehat{p}(\xi) = R$，
即 $\xi$ 位于
$R \in \Ob(\widehat{\mathcal{C}}_\Lambda)$ 上方。
可见，$\xi$ 经 $p(c) \circ p(b)$ 的推前是 $x_\epsilon$，
而 $\xi$ 经 $p(c') \circ p(b')$ 的推前也是 $x_\epsilon$。
由于 $\xi$ 满足（\ref{formal-defos-equation-bijective}），可知
$p(c) \circ p(b) = p(c') \circ p(b')$
作为 $R \to k[\epsilon]$ 的映射相等。
因此 $p(b) = p(b')$ 作为
$R \to A \times_k k[\epsilon]$ 的映射相等。
于是 $y$ 与 $y'$ 都同构于 $\xi$ 沿此映射的推前，
从而如所需，得到与 $b$、$b'$ 相容的、
位于 $A \times_k k[\epsilon]$ 上方的唯一态射 $y \to y'$。

\medskip\noindent
最后，由例 \ref{formal-defos-example-tangent-space-prorepresentable-functor}，
可知
$\dim_k T\mathcal{F} = \dim_k T\underline{R}|_{\mathcal{C}_\Lambda}$
有限。
\end{proof}

\begin{example}
\label{formal-defos-example-do-not-get-S2}
存在这样的预形变范畴：它们有满足
（\ref{formal-defos-equation-bijective-orbits}）的万有形式对象，
但不满足 (S2)。一个简单例子是取
$F = \underline{k[\epsilon]}/G$，其中
$G \subset \text{Aut}_{\mathcal{C}_\Lambda}(k[\epsilon])$
是有限非平凡子群。确切地，映射
$\underline{k[\epsilon]} \to F$ 光滑，
但 $F$ 的切空间没有自然的 $k$-向量空间结构
（因为它是 $k$-向量空间关于有限群的商）。
\end{example}

\begin{lemma}
\label{formal-defos-lemma-construct-bijective-orbits}
设 $\mathcal{F}$ 是满足 (S2) 且具有万有形式对象的预形变范畴。
则它的极小万有形式对象满足
（\ref{formal-defos-equation-bijective-orbits}）。
\end{lemma}

\begin{proof}
设 $\xi$ 是 $\mathcal{F}$ 的极小万有形式对象；
见引理 \ref{formal-defos-lemma-minimal-versal}。
设 $\xi$ 位于
$R \in \Ob(\widehat{\mathcal{C}}_\Lambda)$ 上方。
为解释（\ref{formal-defos-equation-bijective-orbits}），指出如下事实：
$T\mathcal{F}$ 具有自然的 $k$-向量空间结构
（见引理 \ref{formal-defos-lemma-tangent-space-vector-space}）；
$d\underline{\xi} : \text{Der}_\Lambda(R, k) \to T\mathcal{F}$
是线性的（见引理 \ref{formal-defos-lemma-k-linear-differential}）；
而 $\text{Der}_\Lambda(k, k)$ 的作用由加法给出
（见引理 \ref{formal-defos-lemma-action-linear}）。
考虑图
$$
\xymatrix{
& \Hom_k(\mathfrak m_R/\mathfrak m_R^2, k) \\
K \ar[r] & \text{Der}_\Lambda(R, k) \ar[r]^{d\underline{\xi}} \ar[u] &
T\mathcal{F} \\
& \text{Der}_\Lambda(k, k) \ar[u] \ar[ru]
}
$$
向量空间 $K$ 是 $d\underline{\xi}$ 的核。
注意，中间一列在中项正合，因为它对偶于序列
（\ref{formal-defos-equation-sequence}）。
若（\ref{formal-defos-equation-bijective-orbits}）不成立，
则可以找到非零元素 $D \in K$，
它在 $\Hom_k(\mathfrak m_R/\mathfrak m_R^2, k)$ 中的像不为零。
这意味着存在 $t \in \mathfrak m_R$，使 $D(t) = 1$。% P10-E0466
令 $R' = \{a \in R \mid D(a) = 0\}$。
因为 $D$ 是导子，这是 $R$ 的子环。
由于 $D(t) = 1$，可知 $R' \to k$ 是满射
（比较引理 \ref{formal-defos-lemma-essential-surjection} 的证明）。
注意，$\mathfrak m_{R'} = \Ker(D : \mathfrak m_R \to k)$
是 $R$ 的理想，并且
$\mathfrak m_R^2 \subset \mathfrak m_{R'}$。因此
$$
\mathfrak m_R/\mathfrak m_R^2 =
\mathfrak m_{R'}/\mathfrak m_R^2 + k\overline{t}
$$
这蕴含把 $\epsilon$ 映到 $\overline{t}$ 的映射
$$
R'/\mathfrak m_R^2 \times_k k[\epsilon] \to R/\mathfrak m_R^2
$$
是同构。特别地，存在映射
$R/\mathfrak m_R^2 \to R'/\mathfrak m_R^2$。

\medskip\noindent
设 $\xi \to y$ 是位于
$R \to R/\mathfrak m_R^2$ 上方的态射。
设 $y \to x$ 是位于
$R/\mathfrak m_R^2 \to R'/\mathfrak m_R^2$ 上方的态射。
设 $y \to x_\epsilon$ 是位于
$R/\mathfrak m_R^2 \to k[\epsilon]$ 上方的态射。
设 $x_0$ 为 $\mathcal{F}$ 中位于 $k$ 上方的、
在唯一同构意义下唯一的对象。
由群胚余纤维范畴的公理，得到交换图
$$
\vcenter{
\xymatrix{
y \ar[r] \ar[d] & x_\epsilon \ar[d] \\
x \ar[r]   & x_0
}
}
\quad\text{位于下图上方}\quad
\vcenter{
\xymatrix{
R'/\mathfrak m_R^2 \times_k k[\epsilon] \ar[r] \ar[d] & k[\epsilon] \ar[d] \\
R'/\mathfrak m_R^2 \ar[r] & k.
}
}
$$
因为 $D \in K$，可知 $x_\epsilon$ 同构于
$0 \in \mathcal{F}(k[\epsilon])$；即 $x_\epsilon$ 是
$x_0$ 经 $k \to k[\epsilon], a \mapsto a$ 的推前。
因此，由引理 \ref{formal-defos-lemma-lifting-section}，
存在态射 $x \to y$。由于
$\text{length}_\Lambda(R'/\mathfrak m_R^2) <
\text{length}_\Lambda(R/\mathfrak m_R^2)$，
相应环映射
$R'/\mathfrak m_R^2 \to R/\mathfrak m_R^2$ 不是满射。
这与 $\xi/R$ 的极小性矛盾；见引理
\ref{formal-defos-lemma-smallest-where-descends-versal} 的 (1)。
此矛盾表明这样的 $D$ 不存在，结论得证。% P10-E0467
\end{proof}

\begin{theorem}
\label{formal-defos-theorem-miniversal-object-existence}
设 $\mathcal{F}$ 是预形变范畴。考虑以下条件：
\begin{enumerate}
\item $\mathcal{F}$ 有满足（\ref{formal-defos-equation-bijective}）
的极小万有形式对象；
\item $\mathcal{F}$ 有满足
（\ref{formal-defos-equation-bijective-orbits}）的极小万有形式对象；
\item 以下条件成立：
\begin{enumerate}
\item $\mathcal{F}$ 满足 (S1)。
\item $\mathcal{F}$ 满足 (S2)。
\item $\dim_k T\mathcal{F}$ 有限。
\end{enumerate}
\end{enumerate}
总有
$$
(1) \Rightarrow (3) \Rightarrow (2).
$$
若 $k' \subset k$ 可分，则三个条件等价。
\end{theorem}

\begin{proof}
引理 \ref{formal-defos-lemma-miniversal-object-existence-1}
表明 (1) $\Rightarrow$ (3)。
引理 \ref{formal-defos-lemma-versal-object-existence} 和
\ref{formal-defos-lemma-construct-bijective-orbits}
表明 (3) $\Rightarrow$ (2)。
若 $k' \subset k$ 可分，则
$\text{Der}_\Lambda(k, k) = 0$，从而
（\ref{formal-defos-equation-bijective}）$=$
（\ref{formal-defos-equation-bijective-orbits}）；即 (1) 与 (2) 相同。

\medskip\noindent
在经典情形下，(3) $\Rightarrow$ (1) 的另一种证明是：
在引理 \ref{formal-defos-lemma-versal-object-existence} 的证明中补充几句话，
即可看出此时可直接构造满足
（\ref{formal-defos-equation-bijective}）的万有对象。
这在经典情形下避免了使用引理
\ref{formal-defos-lemma-versal-object-existence}。细节从略。% P10-E0468
\end{proof}

\begin{remark}
\label{formal-defos-remark-compare-schlessinger-H3}
设 $F : \mathcal{C}_\Lambda \to \textit{Sets}$ 是满足
(S1)、(S2) 及 $\dim_k TF < \infty$ 的预形变函子。
回忆，这些条件分别对应于 Schlessinger 论文中的
(H1)、(H2) 与 (H3)；见说明
\ref{formal-defos-remark-compare-schlessinger-H3-pre}。
现在，在经典情形下（或当 $k' \subset k$ 可分时），
依照 Schlessinger 引入包络的概念：所谓{\it 包络}，
是满足如下条件的万有形式对象
$\xi \in \widehat{F}(R)$：
$d\underline{\xi} : T\underline{R}|_{\mathcal{C}_\Lambda} \to TF$
为同构；即（\ref{formal-defos-equation-bijective}）成立。
因此，定理 \ref{formal-defos-theorem-miniversal-object-existence} 表明，
在经典情形下
$$
(H1) + (H2) + (H3)
\Rightarrow
\text{存在包络}
$$
成立。换言之，我们的定理恢复了 Schlessinger
关于包络存在性的定理。
\end{remark}

\begin{remark}
\label{formal-defos-remark-compose-minimal-into-iso-classes}
设 $\mathcal{F}$ 是预形变范畴。
回忆，$\mathcal{F} \to \overline{\mathcal{F}}$ 光滑；
见说明 \ref{formal-defos-remark-smooth-to-iso-classes}。
因此，若
$\xi \in \widehat{\mathcal{F}}(R)$ 是万有形式对象，
则复合
$$
\underline{R}|_{\mathcal{C}_\Lambda} \longrightarrow
\mathcal{F} \longrightarrow \overline{\mathcal{F}}
$$
光滑（引理 \ref{formal-defos-lemma-smooth-properties}），
从而 $\xi$ 在 $\overline{\mathcal{F}}$ 中的像
$\overline{\xi}$ 是万有形式对象。
若（\ref{formal-defos-equation-bijective}）成立，
则 $\overline{\xi}$ 诱导同构
$T\underline{R}|_{\mathcal{C}_\Lambda} \to T\overline{\mathcal{F}}$，
因为 $\mathcal{F} \to \overline{\mathcal{F}}$ 识别切空间。
所以在这种情形下，$\overline{\xi}$ 是
$\overline{\mathcal{F}}$ 的包络；见说明
\ref{formal-defos-remark-compare-schlessinger-H3}。
由定理 \ref{formal-defos-theorem-miniversal-object-existence}，
若 $k' \subset k$ 可分，且 $\mathcal{F}$ 是满足
(S1)、(S2) 与 $\dim_k T\mathcal{F} < \infty$ 的预形变范畴，
则总能找到这样的 $\xi$。
\end{remark}

\begin{example}
\label{formal-defos-example-smooth-continued}
在引理 \ref{formal-defos-lemma-exists-smooth} 中，我们构造了对象
$R \in \widehat{\mathcal{C}}_\Lambda$，使
$\underline{R}|_{\mathcal{C}_\Lambda}$ 光滑，并且
$$
H_1(L_{k/\Lambda}) = \mathfrak m_R/\mathfrak m_R^2
\quad\text{和}\quad
\Omega_{R/\Lambda} \otimes_R k = \Omega_{k/\Lambda}
$$
下面用上述定理重新解释这一点。
确切地，把 $\mathcal{F} = \mathcal{C}_\Lambda$
视为其自身上方的群胚余纤维范畴（使用恒等函子）。
则 $\mathcal{F}$ 是预形变范畴，满足 (S1) 与 (S2)，
并且 $T\mathcal{F} = 0$。
因此，$\mathcal{F}$ 满足定理
\ref{formal-defos-theorem-miniversal-object-existence} 的条件 (3)。
该定理蕴含 (2) 成立；即可以找到某个
$S \in \widehat{\mathcal{C}}_\Lambda$ 上方、满足
（\ref{formal-defos-equation-bijective-orbits}）的极小万有形式对象
$\xi \in \widehat{\mathcal{F}}(S)$。
引理 \ref{formal-defos-lemma-smooth} 表明，
$\Lambda \to S$ 在 $\mathfrak m_S$-进拓扑中形式光滑
（因为
$\underline{\xi} : \underline{R}|_{\mathcal{C}_\Lambda} \to
\mathcal{F} = \mathcal{C}_\Lambda$ 光滑）。% P10-E0469
现在，条件（\ref{formal-defos-equation-bijective-orbits}）表明
$\text{Der}_\Lambda(S, k) \to 0$ 在
$\text{Der}_\Lambda(k, k)$-轨道上是双射。
这意味着单射
$\text{Der}_\Lambda(k, k) \to \text{Der}_\Lambda(S, k)$
也是满射。换言之，
$\Omega_{S/\Lambda} \otimes_S k = \Omega_{k/\Lambda}$。
由于 $\Lambda \to S$ 在 $\mathfrak m_S$-进拓扑中形式光滑，
可以应用《更多代数》引理
\ref{more-algebra-lemma-formally-smooth-JZ}，
得出正合序列（\ref{formal-defos-equation-sequence-extended}）
变成一对识别
$$
H_1(L_{k/\Lambda}) = \mathfrak m_S/\mathfrak m_S^2
\quad\text{和}\quad
\Omega_{S/\Lambda} \otimes_S k = \Omega_{k/\Lambda}
$$
逆向阅读该论证，可知引理 \ref{formal-defos-lemma-exists-smooth}
中构造的 $R$ 承载一个极小万有对象。
由极小万有对象的唯一性
（引理 \ref{formal-defos-lemma-minimal-versal}），
还可得 $R \cong S$；即两个构造给出相同答案。
\end{example}







\section{Rim--Schlessinger 条件与形变范畴}
\label{formal-defos-section-RS-condition}

\noindent
$\mathcal{C}_\Lambda$ 上的群胚余纤维范畴有一个非常自然的性质；
它在实践中容易核验，并且蕴含此前引入的 Schlessinger 性质
(S1) 与 (S2)。

\begin{definition}
\label{formal-defos-definition-RS}
设 $\mathcal{F}$ 是 $\mathcal C_\Lambda$ 上的群胚余纤维范畴。
若对 $\mathcal{F}$ 中的每个图
$$
\vcenter{
\xymatrix{
           & x_2 \ar[d] \\
x_1 \ar[r] & x
}
}
\quad\text{位于下图上方}\quad
\vcenter{
\xymatrix{
           & A_2 \ar[d] \\
A_1 \ar[r] & A
}
}
$$
其中底图属于 $\mathcal{C}_\Lambda$，且 $A_2 \to A$ 满射，
都存在 $\mathcal{F}$ 中的纤维积 $x_1 \times_x x_2$，
使图
$$
\vcenter{
\xymatrix{
x_1 \times_x x_2 \ar[r] \ar[d] & x_2 \ar[d] \\
x_1 \ar[r]      & x
}
}
\quad\text{位于下图上方}\quad
\vcenter{
\xymatrix{
A_1 \times_A A_2 \ar[r] \ar[d] & A_2 \ar[d] \\
A_1 \ar[r]      & A.
}
}
$$
则称 $\mathcal{F}$ 满足{\it 条件 (RS)}。
\end{definition}

\begin{lemma}
\label{formal-defos-lemma-RS-fiber-square}
设 $\mathcal{F}$ 是 $\mathcal{C}_\Lambda$ 上满足 (RS)
的群胚余纤维范畴。给定 $\mathcal{F}$ 中的交换图
$$
\vcenter{
\xymatrix{
y \ar[r] \ar[d] & x_2 \ar[d]   \\
x_1 \ar[r]      & x
}
}
\quad\text{位于下图上方}\quad
\vcenter{
\xymatrix{
A_1 \times_A A_2 \ar[r] \ar[d] & A_2 \ar[d] \\
A_1 \ar[r]      & A.
}
}
$$
若 $A_2 \to A$ 满射，则该图是纤维积方块。
\end{lemma}

\begin{proof}
因为 $\mathcal{F}$ 满足 (RS)，存在纤维积图
$$
\vcenter{
\xymatrix{
x_1 \times_x x_2 \ar[r] \ar[d] & x_2 \ar[d] \\
x_1 \ar[r]      & x
}
}
\quad\text{位于下图上方}\quad
\vcenter{
\xymatrix{
A_1 \times_A A_2 \ar[r] \ar[d] & A_2 \ar[d] \\
A_1 \ar[r]      & A.
}
}
$$
诱导映射 $y \to x_1 \times_x x_2$ 位于
$\text{id} : A_1 \times_A A_1 \to A_1 \times_A A_1$ 上方，% P10-E0470
因而是同构。
\end{proof}

\begin{lemma}
\label{formal-defos-lemma-RS-small-extension}
设 $\mathcal{F}$ 是 $\mathcal C_\Lambda$ 上的群胚余纤维范畴。
若只在 $A_2 \to A$ 是小扩张时假设定义
\ref{formal-defos-definition-RS} 中的条件成立，则 $\mathcal{F}$ 满足 (RS)。
\end{lemma}

\begin{proof}
应用引理 \ref{formal-defos-lemma-factor-small-extension}。
证明类似于引理 \ref{formal-defos-lemma-smoothness-small-extensions} 的证明。
\end{proof}

\begin{lemma}
\label{formal-defos-lemma-RS-2-categorical}
设 $\mathcal{F}$ 是 $\mathcal{C}_\Lambda$ 上的群胚余纤维范畴。
以下条件等价：
\begin{enumerate}
\item $\mathcal{F}$ 满足 (RS)；
\item 函子
$\mathcal{F}(A_1 \times_A A_2) \to
\mathcal{F}(A_1) \times_{\mathcal{F}(A)} \mathcal{F}(A_2)$
（见（\ref{formal-defos-equation-compare}））在 $A_2 \to A$ 满射时
都是范畴等价；% P10-E0471
\item 在 $A_2 \to A$ 是小扩张时，(2) 中的同一陈述成立。
\end{enumerate}
\end{lemma}

\begin{proof}
假设 (1)。由引理 \ref{formal-defos-lemma-RS-fiber-square}，
可知 $\mathcal{F}(A_1 \times_A A_2)$ 的每个对象都形如
$x_1 \times_x x_2$。此外，
$$
\Mor_{A_1 \times_A A_2}(x_1 \times_x x_2, y_1 \times_y y_2) =
\Mor_{A_1}(x_1, y_1) \times_{\Mor_A(x, y)}
\Mor_{A_2}(x_2, y_2).
$$
因此，由《范畴》说明
\ref{categories-remark-other-description-2-fibre-product}，
可知 $\mathcal{F}(A_1 \times_A A_2)$ 是
$\mathcal{F}(A_1)$ 与 $\mathcal{F}(A_2)$ 在
$\mathcal{F}(A)$ 上的 $2$-纤维积。
换言之，(2) 成立。

\medskip\noindent
蕴含 (2) $\Rightarrow$ (3) 是立即的。

\medskip\noindent
假设 (3)。给定 $q_1 : A_1 \to A$ 与 $q_2 : A_2 \to A$，
其中 $q_2$ 是小扩张。
将使用《范畴》说明
\ref{categories-remark-other-description-2-fibre-product}
中对 $2$-纤维积
$\mathcal{F}(A_1) \times_{\mathcal{F}(A)} \mathcal{F}(A_2)$ 的描述。
因此，设 $y \in \mathcal{F}(A_1 \times_A A_2)$ 对应于
$(x_1, x_2, x, a_1 : x_1 \to x, a_2 : x_2 \to x)$。
设 $z$ 是 $\mathcal{F}$ 中位于 $C$ 上方的对象。则
\begin{align*}
\Mor_\mathcal{F}(z, y) & =
\{(f, \alpha) \mid f : C \to A_1 \times_A A_2,
\alpha : f_*z \to y\} \\
& = \{(f_1, f_2, \alpha_1, \alpha_2) \mid
f_i : C \to A_i, \ \alpha_i : f_{i, *}z \to x_i, \\
& \quad\quad\quad\quad
q_1 \circ f_1 = q_2 \circ f_2, \ q_{1, *} \alpha_1 = q_{2, *}\alpha_2\} \\ % P10-E0472
& =
\Mor_\mathcal{F}(z, x_1) \times_{\Mor_\mathcal{F}(z, x)}
\Mor_\mathcal{F}(z, x_2)
\end{align*}
所以 $y$ 是 $x_1$ 与 $x_2$ 在 $x$ 上的纤维积。
由此可见，当 $A_2 \to A$ 是小扩张时，
$\mathcal{F}$ 满足 (RS)。于是由引理
\ref{formal-defos-lemma-RS-small-extension}，(RS) 成立。
\end{proof}

\begin{remark}
\label{formal-defos-remark-compare-schlessinger-H4}
当 $\mathcal{F}$ 余纤维于集合时，条件 (RS) 正是
Schlessinger 论文 \cite[Theorem 2.11]{Sch} 中的条件 (H4)。
确切地，对函子
$F: \mathcal{C}_\Lambda \to \textit{Sets}$，
条件 (RS) 表述为：若 $A_1 \to A$ 与 $A_2 \to A$
是 $\mathcal{C}_\Lambda$ 中的映射，且 $A_2 \to A$ 满射，
则诱导映射
$F(A_1 \times_A A_2) \to F(A_1) \times_{F(A)} F(A_2)$
是双射。
\end{remark}

\begin{lemma}
\label{formal-defos-lemma-RS-implies-S1-S2}
设 $\mathcal{F}$ 是 $\mathcal{C}_\Lambda$ 上的群胚余纤维范畴。
$\mathcal{F}$ 的条件 (RS) 蕴含 $\mathcal{F}$ 的 (S1) 与 (S2)。
\end{lemma}

\begin{proof}
使用引理 \ref{formal-defos-lemma-RS-2-categorical} 的重新表述，
以及定义 \ref{formal-defos-definition-S1-S2} 后对 (S1) 的解释，
立即可知 (RS) 蕴含 (S1)。这证明了 (S2) 的第一部分。
(S2) 的第二部分成立，是因为引理
\ref{formal-defos-lemma-RS-fiber-square} 告诉我们：若 $y,y'$ 如定义
\ref{formal-defos-definition-S1-S2} 中 (S2) 的第二部分所述，则
$y = x_1 \times_{d, x_0, e} x_2 = y'$。
（事实上，态射 $y \to y'$ 同时与 $a,a'$ 及 $c,c'$ 相容！）
\end{proof}

\noindent
下述引理类似于引理 \ref{formal-defos-lemma-S1-S2-associated-functor}。
回忆，若 $\mathcal{F}$ 是 $\mathcal{C}_\Lambda$
上的群胚余纤维范畴，且 $x$ 是 $\mathcal{F}$ 中位于 $A$ 上方的对象，
则记
$\text{Aut}_A(x) = \Mor_A(x, x) = \Mor_{\mathcal{F}(A)}(x, x)$。
若 $x' \to x$ 是 $\mathcal{F}$ 中位于
$A' \to A$ 上方的态射，则有良定义的群映射
$\text{Aut}_{A'}(x') \to \text{Aut}_A(x)$。

\begin{lemma}
\label{formal-defos-lemma-RS-associated-functor}
设 $\mathcal{F}$ 是 $\mathcal{C}_\Lambda$ 上满足 (RS)
的群胚余纤维范畴。以下条件等价：
\begin{enumerate}
\item $\overline{\mathcal{F}}$ 满足 (RS)。
\item 设 $f_1: A_1 \to A$ 和 $f_2: A_2 \to A$
是 $\mathcal{C}_\Lambda$ 中的环映射，且 $f_2$ 满射。
同构类集合的诱导映射
$$
\overline{\mathcal{F}(A_1) \times_{\mathcal{F}(A)} \mathcal{F}(A_2)}
\to \overline{\mathcal{F}}(A_1) \times_{\overline{\mathcal{F}}(A)}
\overline{\mathcal{F}}(A_2)
$$
是单射。
\item 对 $\mathcal{F}$ 中位于满环映射
$A' \to A$ 上方的每个态射 $x' \to x$，映射
$\text{Aut}_{A'}(x') \to \text{Aut}_A(x)$ 是满射。
\item 对 $\mathcal{F}$ 中位于小扩张
$A' \to A$ 上方的每个态射 $x' \to x$，映射
$\text{Aut}_{A'}(x') \to \text{Aut}_A(x)$ 是满射。
\end{enumerate}
\end{lemma}

\begin{proof}
我们证明 (1) 等价于 (2)，而 (2) 等价于 (3)。
(3) 与 (4) 的等价性由引理
\ref{formal-defos-lemma-factor-small-extension} 得出。

\medskip\noindent
设 $f_1: A_1 \to A$ 和 $f_2: A_2 \to A$
是 $\mathcal{C}_\Lambda$ 中的环映射，且 $f_2$ 满射。
由说明 \ref{formal-defos-remark-compare-schlessinger-H4}，
$\overline{\mathcal{F}}$ 满足 (RS)，当且仅当对任意这样的
$f_1,f_2$，映射
$$
\overline{\mathcal{F}}(A_1 \times_A A_2) \to \overline{\mathcal
F}(A_1) \times_{\overline{\mathcal{F}}(A)} \overline{\mathcal{F}}(A_2)
$$
都是双射。此映射至少是满射，因为这正是 (S1) 的条件，% P10-E0473
而由引理 \ref{formal-defos-lemma-RS-implies-S1-S2} 和
\ref{formal-defos-lemma-S1-S2-associated-functor}，
$\overline{\mathcal{F}}$ 满足 (S1)。
此外，该映射分解为
$$
\overline{\mathcal{F}}(A_1 \times_A A_2)
\longrightarrow
\overline{\mathcal{F}(A_1) \times_{\mathcal{F}(A)} \mathcal{F}(A_2)}
\longrightarrow
\overline{\mathcal{F}}(A_1) \times_{\overline{\mathcal{F}}(A)}
\overline{\mathcal{F}}(A_2),
$$
其中第一个映射是双射，因为
$$
\mathcal{F}(A_1 \times_A A_2)
\longrightarrow
\mathcal{F}(A_1) \times_{\mathcal{F}(A)} \mathcal{F}(A_2)
$$
由 $\mathcal{F}$ 的 (RS) 是等价。因此 (1) 等价于 (2)。

\medskip\noindent
假设 (2) 成立。设 $x' \to x$ 是 $\mathcal{F}$ 中
位于满环映射 $f: A' \to A$ 上方的态射。
设 $a \in \text{Aut}_A(x)$。对象
$$
(x', x', a : x \to x), \ (x', x', \text{id} : x \to x)
$$
属于
$\mathcal{F}(A') \times_{\mathcal{F}(A)} \mathcal{F}(A')$，
并在
$\overline{\mathcal{F}}(A') \times_{\overline{\mathcal{F}}(A)}
\overline{\mathcal{F}}(A')$ 中有相同的像。
由 (2)，存在映射 $b_1,b_2 : x' \to x'$，% P10-E0474
使图
$$
\xymatrix{
x \ar[r]_a \ar[d]_{f_*b_1} & x \ar[d]^{f_*b_2} \\
x \ar[r]^{\text{id}} & x
}
$$
交换。因此，
$b_2^{-1} \circ b_1 \in \text{Aut}_{A'}(x')$
的像是 $a \in \text{Aut}_A(x)$。所以 (3) 成立。

\medskip\noindent
假设 (3) 成立。设
$$
(x_1, x_2, a : (f_1)_*x_1 \to (f_2)_*x_2),
\ (x'_1, x'_2, a' : (f_1)_*x'_1 \to (f_2)_*x'_2)
$$
是
$\mathcal{F}(A_1) \times_{\mathcal{F}(A)} \mathcal{F}(A_2)$
中在
$\overline{\mathcal{F}}(A_1) \times_{\overline{\mathcal{F}}(A)}
\overline{\mathcal{F}}(A_2)$ 中具有相同像的对象。
则存在 $\mathcal{F}(A_1)$ 中的态射
$b_1: x_1 \to x'_1$，以及 $\mathcal F(A_2)$ 中的态射
$b_2: x_2 \to x'_2$。
由 (3)，可以用 $x_2$ 在 $A_2$ 上的自同构修改 $b_2$，
使图
$$
\xymatrix{
(f_1)_*x_1 \ar[r]_a \ar[d]_{(f_1)_*b_1} & (f_2)_*x_2 \ar[d]^{(f_2)_*b_2} \\
(f_1)_*x'_1 \ar[r]^{a'} & (f_2)_*x'_2.
}
$$
交换。这证明
$(x_1, x_2, a) \cong (x'_1, x'_2, a')$ 在
$\overline{\mathcal{F}(A_1) \times_{\mathcal{F}(A)} \mathcal{F}(A_2)}$
中成立。% P10-E0475
所以 (2) 成立。
\end{proof}

\noindent
最后，我们定义形变范畴的概念。

\begin{definition}
\label{formal-defos-definition-deformation-category}
满足 (RS) 的前形变范畴 $\mathcal{F}$ 称为一个{\it 形变范畴}。
形变范畴之间的态射，是它们作为 $\mathcal{C}_\Lambda$ 上范畴的态射。
\end{definition}

\begin{remark}
\label{formal-defos-remark-deformation-functor}
若函子 $F: \mathcal{C}_\Lambda \to \textit{Sets}$ 所对应的
余纤维集合是形变范畴，即 $F(k)$ 是单元素集合且 $F$ 满足 (RS)，
则称这样的函子为一个{\it 形变函子}。若 $\mathcal{F}$ 是形变范畴，
则 $\overline{\mathcal{F}}$ 是前形变函子，但未必是形变函子，
正如引理 \ref{formal-defos-lemma-RS-associated-functor} 所示。
\end{remark}

\begin{example}
\label{formal-defos-example-prorepresentable-deformation-functor}
pro-可表函子 $F$ 是形变函子。事实上，设
$R \in \Ob(\widehat{\mathcal{C}}_\Lambda)$ 且
$F(A) = \Mor_{\widehat{\mathcal{C}}_\Lambda}(R, A)$。
态射 $R \to k$ 唯一，所以 $F(k)$ 是单元素集合。由于
$$
\Hom_\Lambda(R, A_1 \times_A A_2) =
\Hom_\Lambda(R, A_1) \times_{\Hom_\Lambda(R, A)}
\Hom_\Lambda(R, A_2),
$$
对 $\widehat{\mathcal{C}}_\Lambda$ 中的映射也有同样的结论，
因而 $F$ 满足 (RS)。
\end{example}

\noindent
下面是我们关于以下过程的典型说明之一：从一个群胚余纤维范畴
转到其在 $k$ 上某一点处的前形变范畴。它说明该过程保持 (RS)。

\begin{lemma}
\label{formal-defos-lemma-localize-RS}
设 $\mathcal{F}$ 是 $\mathcal{C}_\Lambda$ 上的群胚余纤维范畴，
且 $x_0 \in \Ob(\mathcal{F}(k))$。设 $\mathcal{F}_{x_0}$ 是说明
\ref{formal-defos-remark-localize-cofibered-groupoid} 所构造的
$\mathcal{C}_\Lambda$ 上的群胚余纤维范畴。
若 $\mathcal{F}$ 满足 (RS)，则 $\mathcal{F}_{x_0}$ 也满足 (RS)。
特别地，$\mathcal{F}_{x_0}$ 是形变范畴。
\end{lemma}

\begin{proof}
$\mathcal{F}_{x_0}$ 中任何如定义 \ref{formal-defos-definition-RS} 所述的图表
都会给出 $\mathcal{F}$ 中的一个图表；在 $\mathcal{F}$ 中
对后一个图表应用 (RS) 所得的对象，
也可视为对 $\mathcal{F}_{x_0}$ 中原图表应用 (RS) 所得的对象。
\end{proof}

\noindent
下述引理类似于代数栈的 $2$-纤维积仍为代数栈这一事实。

\begin{lemma}
\label{formal-defos-lemma-RS-fiber-product-morphisms}
设
$$
\xymatrix{
\mathcal{H} \times_\mathcal{F} \mathcal{G} \ar[r] \ar[d] &
\mathcal{G} \ar[d]^g \\
\mathcal{H} \ar[r]^f & \mathcal{F}
}
$$
是 $\mathcal{C}_\Lambda$ 上群胚余纤维范畴的 $2$-纤维积。% P10-E0476
若 $\mathcal{F}, \mathcal{G}, \mathcal{H}$ 都满足 (RS)，
则 $\mathcal{H} \times_\mathcal{F} \mathcal{G}$ 满足 (RS)。
\end{lemma}

\begin{proof}
若 $A$ 是 $\mathcal{C}_\Lambda$ 的对象，则
$\mathcal{H} \times_\mathcal{F} \mathcal{G}$ 在 $A$ 上的纤维范畴中的对象
是三元组 $(u, v, a)$，其中 $u \in \mathcal{H}(A)$、
$v \in \mathcal{G}(A)$，且 $a : f(u) \to g(v)$ 是
$\mathcal{F}(A)$ 中的态射。考虑
$\mathcal{H} \times_\mathcal{F} \mathcal{G}$ 中的图表
$$
\vcenter{
\xymatrix{
           & (u_2, v_2, a_2) \ar[d] \\
(u_1, v_1, a_1) \ar[r] & (u, v, a)
}
}
\quad\text{位于}\quad
\vcenter{
\xymatrix{
           & A_2 \ar[d] \\
A_1 \ar[r] & A
}
}
$$
之上，其中后一个图表在 $\mathcal{C}_\Lambda$ 中且
$A_2 \to A$ 满射。由于 $\mathcal{H}$ 和 $\mathcal{G}$ 满足 (RS)，
存在位于 $A_1 \times_A A_2$ 上的纤维积
$u_1 \times_u u_2$ 与 $v_1 \times_v v_2$。由于
$\mathcal{F}$ 满足 (RS)，引理 \ref{formal-defos-lemma-RS-fiber-square} 表明
$$
\vcenter{
\xymatrix{
f(u_1 \times_u u_2) \ar[r] \ar[d] & f(u_2) \ar[d] \\
f(u_1) \ar[r] & f(u)
}
}
\quad\text{与}\quad
\vcenter{
\xymatrix{
g(v_1 \times_v v_2) \ar[r] \ar[d] & g(v_2) \ar[d] \\
g(v_1) \ar[r] & g(v)
}
}
$$
在 $\mathcal{F}$ 中都是纤维积方块。因此可将
$a_1 \times_a a_2$ 视为从 $f(u_1 \times_u u_2)$ 到
$g(v_1 \times_v v_2)$ 的、位于 $A_1 \times_A A_2$ 上的态射。
由此可知
$$
\xymatrix{
 (u_1 \times_u u_2, v_1 \times_v v_2, a_{1} \times_a a_2) \ar[d] \ar[r] &
(u_2, v_2, a_2) \ar[d] \\
(u_1, v_1, a_1) \ar[r] & (u, v, a)
}
$$
是 $\mathcal{H} \times_\mathcal{F} \mathcal{G}$ 中的纤维积方块，
这正是所需结论。
\end{proof}

\section{对象的提升}
\label{formal-defos-section-lifts}

\noindent
本节的内容是：在形变范畴的情形，切空间在沿小扩张的提升集合上
有一个单传递作用。

\begin{definition}
\label{formal-defos-definition-lifts}
设 $\mathcal{F}$ 是 $\mathcal{C}_\Lambda$ 上的群胚余纤维范畴。
设 $f: A' \to A$ 是 $\mathcal{C}_\Lambda$ 中的映射，并设
$x \in \mathcal{F}(A)$。$x$ 沿 $f$ 的提升范畴
$\textit{Lift}(x, f)$ 具有如下对象与态射。
\begin{enumerate}
\item 对象：$x$ 沿 $f$ 的一个{\it 提升}，是位于 $f$ 上的态射
$x' \to x$。
\item 态射：从提升 $a_1 : x'_1 \to x$ 到提升
$a_2 : x'_2 \to x$ 的一个{\it 提升态射}，是
$\mathcal{F}(A')$ 中满足 $a_2 = a_1 \circ b$ 的态射
$b : x'_1 \to x'_2$。% P10-E0477
\end{enumerate}
$x$ 沿 $f$ 的提升集合 $\text{Lift}(x, f)$，是
$\textit{Lift}(x, f)$ 的同构类集合。
\end{definition}

\begin{remark}
\label{formal-defos-remark-omit-arrow}
当映射 $f: A' \to A$ 由语境可知时，可以分别用
$\textit{Lift}(x, A')$ 和 $\text{Lift}(x, A')$ 代替
$\textit{Lift}(x, f)$ 和 $\text{Lift}(x, f)$。
\end{remark}

\begin{remark}
\label{formal-defos-remark-tangent-space-lifting}
设 $\mathcal{F}$ 是 $\mathcal{C}_\Lambda$ 上的群胚余纤维范畴，
且 $x_0 \in \Ob(\mathcal{F}(k))$。设 $V$ 是有限维向量空间。
则 $\text{Lift}(x_0, k[V])$ 是 $\mathcal{F}_{x_0}(k[V])$
的同构类集合，其中 $\mathcal{F}_{x_0}$ 是 $\mathcal{F}$ 中
位于 $x_0$ 上的对象所成的前形变范畴；见说明
\ref{formal-defos-remark-localize-cofibered-groupoid}。因此，若 $\mathcal{F}$
满足 (S2)，则 $\mathcal{F}_{x_0}$ 也满足 (S2)（见引理
\ref{formal-defos-lemma-S1-S2-localize}），并且由引理
\ref{formal-defos-lemma-tangent-space-vector-space} 可知
$$
\text{Lift}(x_0, k[V]) = T\mathcal{F}_{x_0} \otimes_k V
$$
是 $k$-向量空间之间的等式。
\end{remark}

\begin{remark}
\label{formal-defos-remark-lift-bijections}
设 $\mathcal{F}$ 是 $\mathcal
C_\Lambda$ 上满足 (RS) 的群胚余纤维范畴。设
$$
\xymatrix{
A_1 \times_A A_2 \ar[r] \ar[d] & A_2 \ar[d] \\
A_1 \ar[r] & A
}
$$
是 $\mathcal{C}_\Lambda$ 中的纤维积方块，并且
$A_1 \to A$ 或 $A_2 \to A$ 至少有一个满射。设
$x \in \Ob(\mathcal{F}(A))$。给定 $x$ 到 $A_1$ 和 $A_2$
的提升 $x_1 \to x$ 与 $x_2 \to x$，由 (RS) 得到
$x$ 到 $A_1 \times_A A_2$ 的提升
$x_1 \times_x x_2 \to x$。反过来，由引理
\ref{formal-defos-lemma-RS-fiber-square}，$x$ 到 $A_1 \times_A A_2$
的每个提升都具有这种形式。从而得到双射% P10-E0478
$$
\text{Lift}(x, A_1) \times \text{Lift}(x, A_2)
\longrightarrow
\text{Lift}(x, A_1 \times_A A_2).
$$
类似地，若 $x_1 \to x$ 是 $x$ 到 $A_1$ 的一个固定提升，
则有双射
$$
\text{Lift}(x_1, A_1 \times_A A_2)
\longrightarrow
\text{Lift}(x, A_2).
$$
现设
$$
\xymatrix{
A_1' \times_A A_2 \ar[r] \ar[d] & A_1 \times_A A_2 \ar[r] \ar[d] & A_2
\ar[d] \\
A_1' \ar[r] & A_1 \ar[r] & A
}
$$
是 $\mathcal{C}_\Lambda$ 中两个纤维积方块的复合，且
$A'_1 \to A_1$ 与 $A_1 \to A$ 都满射。设 $x_1 \to x$
是位于 $A_1 \to A$ 上的态射。由上可得双射
\begin{align*}
\text{Lift}(x_1, A_1' \times_A A_2)
& = \text{Lift}(x_1, A_1') \times \text{Lift}(x_1, A_1 \times_A A_2) \\
& = \text{Lift}(x_1, A_1') \times \text{Lift}(x, A_2).
\end{align*}
\end{remark}

\begin{lemma}
\label{formal-defos-lemma-free-transitive-action}
设 $\mathcal{F}$ 是形变范畴。设 $A' \to A$ 是
$\mathcal{C}_\Lambda$ 中的满环映射，其核 $I$ 被
$\mathfrak m_{A'}$ 零化。设 $x \in \Ob(\mathcal{F}(A))$。
若 $\text{Lift}(x, A')$ 非空，则
$T\mathcal{F} \otimes_k I$ 在 $\text{Lift}(x, A')$
上有一个自由且传递的作用。
\end{lemma}

\begin{proof}
考虑环映射 $g : A' \times_A A' \to k[I]$，其定义规则为
$g(a_1, a_2) = \overline{a_1} \oplus a_2 - a_1$（与引理
\ref{formal-defos-lemma-lifting-along-small-extension} 比较）。存在同构
$$
A' \times_A A' \xrightarrow{\sim} A' \times_k k[I]
$$
由
$(a_1, a_2) \mapsto (a_1, g(a_1, a_2))$
给出。
此同构与到第一因子 $A'$ 的投影相容，因而也与
$A' \times_A A'$ 和 $A' \times_k k[I]$ 到 $A$ 的投影相容。
所以有双射
\begin{equation}
\label{formal-defos-equation-one}
\text{Lift}(x, A' \times_A A')
\longrightarrow
\text{Lift}(x, A' \times_k k[I])
\end{equation}
由说明 \ref{formal-defos-remark-lift-bijections}，有双射
\begin{equation}
\label{formal-defos-equation-two}
\text{Lift}(x, A') \times \text{Lift}(x, A')
\longrightarrow
\text{Lift}(x, A' \times_A A')
\end{equation}
有交换图
$$
\xymatrix{
A' \times_k k[I] \ar[r] \ar[d] & A \times_k k[I] \ar[r] \ar[d] & k[I] \ar[d] \\
A' \ar[r] & A \ar[r] & k.
}
$$
因此，若选定 $x$ 沿 $A \to k$ 的一个推前 $x \to x_0$，
则由说明 \ref{formal-defos-remark-lift-bijections} 的最后部分得到双射
\begin{equation}
\label{formal-defos-equation-three}
\text{Lift}(x, A' \times_k k[I])
\longrightarrow
\text{Lift}(x, A') \times \text{Lift}(x_0, k[I])
\end{equation}
依次复合（\ref{formal-defos-equation-two}）、（\ref{formal-defos-equation-one}）和
（\ref{formal-defos-equation-three}），得到双射
$$
\Phi :
\text{Lift}(x, A') \times \text{Lift}(x, A')
\longrightarrow
\text{Lift}(x, A') \times \text{Lift}(x_0, k[I]).
$$
此双射与到第一因子的投影相容。由说明
\ref{formal-defos-remark-tangent-space-lifting} 可知
$\text{Lift}(x_0, k[I]) = T\mathcal{F} \otimes_k I$。
若 $\text{pr}_2$ 是
$\text{Lift}(x, A') \times \text{Lift}(x, A')$ 的第二投影，
则得到映射
$$
a = \text{pr}_2 \circ \Phi^{-1} :
\text{Lift}(x, A') \times (T\mathcal{F} \otimes_k I)
\longrightarrow
\text{Lift}(x, A').
$$
展开上述所有构造可知，$a(x' \to x, \theta)$ 是唯一的提升
$x'' \to x$，使得
$g_*(x', x'') = \theta$ 在
$\text{Lift}(x_0, k[I]) = T\mathcal{F} \otimes_k I$ 中成立。
为证明这是 $T\mathcal{F} \otimes_k I$ 在
$\text{Lift}(x, A')$ 上的作用，必须证明如下结论：
若 $x', x'', x'''$ 是 $x$ 的提升，并且
$g_*(x', x'') = \theta$、$g_*(x'', x''') = \theta'$，则
$g_*(x', x''') = \theta + \theta'$。这由交换图
$$
\xymatrix{
A' \times_A A' \times_A A'
\ar[rrrrr]_-{(a_1, a_2, a_3) \mapsto (g(a_1, a_2), g(a_2, a_3))}
\ar[rrrrrd]_{(a_1, a_2, a_3) \mapsto g(a_1, a_3)} & & & & &
k[I] \times_k k[I] = k[I \times I] \ar[d]^{+} \\
& & & & & k[I]
}
$$
得出。由于 $\Phi$ 是双射，此作用自由且传递。
\end{proof}

\begin{remark}
\label{formal-defos-remark-free-transitive-action-functorial}
引理 \ref{formal-defos-lemma-free-transitive-action} 的作用具有函子性。
设 $\varphi : \mathcal{F} \to \mathcal{G}$ 是形变范畴的态射。
设 $A' \to A$ 是满环映射，其核 $I$ 被
$\mathfrak m_{A'}$ 零化。设 $x \in \Ob(\mathcal{F}(A))$。
在此情形，$\varphi$ 诱导下列交换图中的竖直箭头：
$$
\xymatrix{
\text{Lift}(x, A') \times (T\mathcal{F} \otimes_k I)
\ar[d]_{(\varphi, d\varphi \otimes \text{id}_I)} \ar[r] &
\text{Lift}(x, A') \ar[d]^\varphi \\
\text{Lift}(\varphi(x), A') \times (T\mathcal{G} \otimes_k I) \ar[r] &
\text{Lift}(\varphi(x), A')
}
$$
该图交换，是因为映射（\ref{formal-defos-equation-two}）、
（\ref{formal-defos-equation-one}）和（\ref{formal-defos-equation-three}）中的每一个
（见引理 \ref{formal-defos-lemma-free-transitive-action} 的证明）
都给出一个类似的交换图。
\end{remark}

\section{Schlessinger 关于 pro-可表函子的定理}
\label{formal-defos-section-schlessingers-theorem}

\noindent
下面推出 Schlessinger 刻画 $\mathcal{C}_\Lambda$ 上
pro-可表函子的定理。

\begin{lemma}
\label{formal-defos-lemma-minimal-smooth-morphism-functors}
设 $F, G: \mathcal{C}_\Lambda \to \textit{Sets}$ 是形变函子。
设 $\varphi : F \to G$ 是光滑态射，并且诱导切空间的同构
$d\varphi : TF \to TG$。则 $\varphi$ 是同构。
\end{lemma}

\begin{proof}
对 $\text{length}_A(A)$ 作归纳，证明对所有
$A \in \Ob(\mathcal{C}_\Lambda)$，映射 $F(A) \to G(A)$
都是双射。当 $A = k$ 时，该陈述由 $F$ 和 $G$ 是形变函子的假设得出。
假设该陈述对长度小于 $n$ 的环成立，并设 $A'$ 是长度为 $n$ 的环。
选取小扩张 $f : A' \to A$。有交换图
$$
\xymatrix{
F(A') \ar[r] \ar[d]_{F(f)} & G(A') \ar[d]^{G(f)} \\
F(A) \ar[r]^{\sim} & G(A)
}
$$
其中映射 $F(A) \to G(A)$ 是双射。由 $F
\to G$ 的光滑性，$F(A') \to G(A')$ 是满射（引理
\ref{formal-defos-lemma-smooth-morphism-essentially-surjective}）。因此，只需对满足
$F(f)^{-1}(x)$ 非空的 $x \in F(A)$，在纤维
$F(f)^{-1}(x) \to G(f)^{-1}(\varphi(x))$ 上检验双射性。
这些纤维恰为 $\text{Lift}(x, A')$ 和
$\text{Lift}(\varphi(x), A')$，并且由假设有同构
$d\varphi \otimes \text{id} :
TF \otimes_k \Ker(f)  \to TG \otimes_k \Ker(f)$。
所以，由引理 \ref{formal-defos-lemma-free-transitive-action} 和说明
\ref{formal-defos-remark-free-transitive-action-functorial}，对满足
$F(f)^{-1}(x)$ 非空的 $x \in F(A)$，映射
$F(f)^{-1}(x) \to G(f)^{-1}(\varphi(x))$ 与
$TF \otimes_k \Ker(f)$ 的自由传递作用相容。因此它是双射。
\end{proof}

\noindent
注意，当 $k' \subset k$ 可分时，下述定理中的条件 (c) 是空条件。

\begin{theorem}
\label{formal-defos-theorem-Schlessinger-prorepresentability}
设 $F: \mathcal{C}_\Lambda \to \textit{Sets}$ 是函子。
则 $F$ pro-可表，当且仅当：(a) $F$ 是形变函子；
(b) $\dim_k TF$ 有限；以及
(c) $\gamma : \text{Der}_\Lambda(k, k) \to TF$ 是单射。
\end{theorem}

\begin{proof}
假设 $F$ 由 $R \in \widehat{\mathcal{C}}_\Lambda$ pro-表示。
由例 \ref{formal-defos-example-prorepresentable-deformation-functor} 可知
$F$ 是形变函子。由例
\ref{formal-defos-example-tangent-space-prorepresentable-functor} 可知
$\dim_k TF$ 有限。最后，由例
\ref{formal-defos-example-tangent-space-map-prorepresentable-functor}，
$\text{Der}_\Lambda(k, k) \to TF$ 等同于
$\text{Der}_\Lambda(k, k) \to \text{Der}_\Lambda(R, k)$；
后者是单射，因为 $R \to k$ 满射。

\medskip\noindent
反过来，假设 (a)、(b)、(c) 成立。由引理
\ref{formal-defos-lemma-RS-implies-S1-S2} 可知 (S1) 与 (S2) 成立。
所以由定理 \ref{formal-defos-theorem-miniversal-object-existence}，
存在 $F$ 的极小万有形式对象 $\xi$，使得
（\ref{formal-defos-equation-bijective-orbits}）成立。设 $\xi$ 位于 $R$ 上。
映射
$$
d\underline{\xi} : \text{Der}_\Lambda(R, k) \to T\mathcal{F}
$$
在 $\text{Der}_\Lambda(k, k)$-轨道上为双射。% P10-E0479
由 (c) 和引理 \ref{formal-defos-lemma-action-linear}，
$\text{Der}_\Lambda(k, k)$ 在左端的作用是自由的，
故该映射是双射。因此，由引理
\ref{formal-defos-lemma-minimal-smooth-morphism-functors}，
$\underline{\xi}$ 是同构。
\end{proof}

\section{无穷小自同构}
\label{formal-defos-section-infinitesimal-automorphisms}

\noindent
设 $\mathcal{F}$ 是 $\mathcal{C}_\Lambda$ 上的群胚余纤维范畴。
给定 $\mathcal{F}$ 中位于 $A' \to A$ 上的态射 $x' \to x$，
会诱导同态
$$
\text{Aut}_{A'}(x') \to \text{Aut}_A(x).
$$
引理 \ref{formal-defos-lemma-RS-associated-functor} 表明，此同态的余核决定
$\mathcal{F}$ 的条件 (RS) 是否能传递给
$\overline{\mathcal{F}}$。本节研究此同态的核。我们将看到，
它也衡量 $\mathcal{F}$ 与 $\overline{\mathcal{F}}$ 之间的差距。

\begin{definition}
\label{formal-defos-definition-relative-infinitesimal-auts}
设 $\mathcal{F}$ 是 $\mathcal
C_\Lambda$ 上的群胚余纤维范畴。设 $x' \to x$ 是
$\mathcal{F}$ 中位于 $A' \to A$ 上的态射。核
$$
\text{Inf}(x'/x) = \Ker(\text{Aut}_{A'}(x') \to \text{Aut}_A(x))
$$
称为 $x'$ 在 $x$ 上的{\it 无穷小自同构群}。
\end{definition}

\begin{definition}
\label{formal-defos-definition-infinitesimal-auts}
设 $\mathcal{F}$ 是 $\mathcal
C_\Lambda$ 上的群胚余纤维范畴。设
$x_0 \in \Ob(\mathcal{F}(k))$。假定已经选定 $x_0$ 沿映射
$k \to k[\epsilon], a \mapsto a$ 的一个推前 $x_0 \to x_0'$。
则存在唯一映射 $x'_0 \to x_0$，使得
$x_0 \to x_0' \to x_0$ 是 $x_0$ 上的恒等映射。于是
$$
\text{Inf}_{x_0}(\mathcal F) = \text{Inf}(x'_0/x_0)
$$
称为 $x_0$ 的{\it 无穷小自同构群}。% P10-E0480
\end{definition}

\begin{remark}
\label{formal-defos-remark-choice-pushforward-immaterial-infinitesimal-aut}
在典范同构的意义下，$\text{Inf}_{x_0}(\mathcal{F})$
与推前 $x_0 \to x_0'$ 的选择无关，因为任意两个推前典范同构。
此外，若 $y_0 \in \mathcal{F}(k)$ 且
$x_0 \cong y_0$ 在 $\mathcal{F}(k)$ 中成立，则
$\text{Inf}_{x_0}(\mathcal{F}) \cong \text{Inf}_{y_0}(\mathcal{F})$；
该同构（仅）依赖于同构 $x_0 \to y_0$ 的选择。
特别地，$\text{Aut}_k(x_0)$ 作用在
$\text{Inf}_{x_0}(\mathcal{F})$ 上。
\end{remark}

\begin{remark}
\label{formal-defos-remark-trivial-aut-point}
假设 $\mathcal{F}$ 是前形变范畴。则
\begin{enumerate}
\item 对 $x_0 \in \Ob(\mathcal{F}(k))$，自同构群
$\text{Aut}_k(x_0)$ 平凡，因而
$\text{Inf}_{x_0}(\mathcal{F}) = \text{Aut}_{k[\epsilon]}(x'_0)$；并且
\item 对 $x_0, y_0 \in \Ob(\mathcal{F}(k))$，存在唯一同构
$x_0 \to y_0$，因而有典范识别
$\text{Inf}_{x_0}(\mathcal{F}) = \text{Inf}_{y_0}(\mathcal{F})$。
\end{enumerate}
由于 $\mathcal{F}(k)$ 非空，选取
$x_0 \in \Ob(\mathcal{F}(k))$ 并置
$$
\text{Inf}(\mathcal{F}) = \text{Inf}_{x_0}(\mathcal{F}),
$$
便得到良定义的 $\mathcal{F}$ 的{\it 无穷小自同构群}。
用此记号，有
$\text{Inf}(\mathcal{F}_{x_0}) = \text{Inf}_{x_0}(\mathcal{F})$。
请与说明 \ref{formal-defos-remark-tangent-space-cofibered-groupoid} 中的等式
$T\mathcal{F}_{x_0} = T_{x_0}\mathcal{F}$ 比较。
\end{remark}

\noindent
我们将看到，当 $\mathcal{F}$ 满足 (RS) 时，
$\text{Inf}_{x_0}(\mathcal{F})$ 具有自然的 $k$-向量空间结构。
同时还将看到，若 $\mathcal{F}$ 满足 (RS)，则位于小扩张上方的态射
$x' \to x$ 的无穷小自同构 $\text{Inf}(x'/x)$ 由
$\text{Inf}_{x_0}(\mathcal{F})$ 控制，其中 $x_0$ 是 $x$
到 $\mathcal{F}(k)$ 的一个推前。为此，对任意对象
$x \in \Ob(\mathcal{F})$ 引入如下自同构函子。

\begin{definition}
\label{formal-defos-definition-automorphism-functor}
设 $p : \mathcal{F} \to \mathcal{C}$ 是任意基范畴
$\mathcal{C}$ 上的群胚余纤维范畴，并假定已经选定推前。
设 $x \in \Ob(\mathcal{F})$，并令 $U = p(x)$。
以 $U/\mathcal{C}$ 表示 $U$ 下的对象范畴。
$x$ 的{\it 自同构函子}是函子
$\mathit{Aut}(x) : U/\mathcal{C} \to \textit{Sets}$，
它把对象 $f : U \to V$ 送到 $\text{Aut}_V(f_*x)$，
并把态射
$$
\xymatrix{
V' \ar[rr] &                    & V\\
          & U \ar[ul]^{f'}  \ar[ur]_f &
}
$$
送到同态
$\text{Aut}_{V'}(f'_*x) \to \text{Aut}_V(f_*x)$；该同态来自
唯一的态射 $f'_*x \to f_*x$，它位于 $V' \to V$ 上，
并与 $x \to f'_*x$ 及 $x \to f_*x$ 相容。
\end{definition}

\noindent
我们将研究 $\mathcal{C}_\Lambda$ 上群胚余纤维范畴
$\mathcal{F}$ 中对象的自同构函子。若
$A \in \Ob(\mathcal{C}_\Lambda)$，则范畴
$A/\mathcal{C}_\Lambda$ 正是范畴 $\mathcal{C}_A$，
即第 \ref{formal-defos-section-CLambda} 节所定义的范畴，其中取
$\Lambda = A$ 且 $k = A/\mathfrak m_A$。因此，对象
$x \in \Ob(\mathcal{F}(A))$ 的自同构函子是函子
$\mathit{Aut}(x) : \mathcal{C}_A \to \textit{Sets}$。

\medskip\noindent
下述引理也可从引理 \ref{formal-defos-lemma-RS-fiber-product-morphisms} 推出：
只需考虑群胚余纤维范畴的“惯性”；例如见《栈》第
\ref{stacks-section-the-inertia-stack} 节和《范畴》第
\ref{categories-section-inertia} 节。不过，直接证明更容易。

\begin{lemma}
\label{formal-defos-lemma-Aut-functor-RS}
设 $\mathcal{F}$ 是 $\mathcal{C}_\Lambda$ 上满足 (RS)
的群胚余纤维范畴，并设 $x \in \Ob(\mathcal{F}(A))$。
则 $\mathit{Aut}(x): \mathcal{C}_A \to \textit{Sets}$ 满足 (RS)。
\end{lemma}

\begin{proof}
由引理 \ref{formal-defos-lemma-RS-2-categorical} 中函子
$\mathcal{F}(A_1 \times_A A_2) \to
\mathcal{F}(A_1) \times_{\mathcal{F}(A)} \mathcal{F}(A_2)$
的全忠实性可知，$\mathit{Aut}(x)$ 满足 (RS)。% P10-E0481
\end{proof}

\begin{lemma}
\label{formal-defos-lemma-Aut-functor-tangent-space}
设 $\mathcal{F}$ 是 $\mathcal{C}_\Lambda$ 上满足 (RS)
的群胚余纤维范畴。设 $x \in \Ob(\mathcal{F}(A))$，
并设 $x_0$ 是 $x$ 到 $\mathcal{F}(k)$ 的一个推前。
\begin{enumerate}
\item $T_{\text{id}_{x_0}} \mathit{Aut}(x)$ 具有自然的
$k$-向量空间结构，使得加法与
$T_{\text{id}_{x_0}} \mathit{Aut}(x)$ 中的复合一致。
特别地，$T_{\text{id}_{x_0}} \mathit{Aut}(x)$ 中的复合可交换。
\item 有 $k$-向量空间的典范同构
$T_{\text{id}_{x_0}} \mathit{Aut}(x) \to
T_{\text{id}_{x_0}} \mathit{Aut}(x_0)$。
\end{enumerate}
\end{lemma}

\begin{proof}
对函子 $\mathit{Aut}(x) : \mathcal{C}_A \to \textit{Sets}$
及元素 $\text{id}_{x_0} \in \mathit{Aut}(x)(k)$ 应用说明
\ref{formal-defos-remark-localize-cofibered-groupoid}，得到前形变函子
$F = \mathit{Aut}(x)_{\text{id}_{x_0}}$。由引理
\ref{formal-defos-lemma-Aut-functor-RS} 和 \ref{formal-defos-lemma-localize-RS}，
$F$ 是形变函子。根据定义，
$T_{\text{id}_{x_0}} \mathit{Aut}(x) = TF = F(k[\epsilon])$，
它具有引理 \ref{formal-defos-lemma-tangent-space-functor} 所给出的自然
$k$-向量空间结构。

\medskip\noindent
加法定义为复合
$$
F(k[\epsilon]) \times F(k[\epsilon]) \longrightarrow
F(k[\epsilon] \times_k k[\epsilon]) \longrightarrow
F(k[\epsilon]),
$$
其中第一个映射是由 (RS) 保证的双射的逆，而第二个映射由
$k$-代数映射
$k[\epsilon] \times_k k[\epsilon] \to k[\epsilon]$
诱导；该映射把 $(\epsilon, 0)$ 与 $(0, \epsilon)$ 都映到
$\epsilon$。若 $A \to B$ 是 $\mathcal{C}_\Lambda$ 中的环映射，
则 $F(A) \to F(B)$ 是同态，其中
$F(A) = \mathit{Aut}(x)_{\text{id}_{x_0}}(A)$ 与
$F(B) = \mathit{Aut}(x)_{\text{id}_{x_0}}(B)$
都在复合运算下成为群。因此，
$+ : F(k[\epsilon]) \times F(k[\epsilon])\to F(k[\epsilon])$
是同态，其中把 $F(k[\epsilon])$ 看作复合运算下的群。
以 $\text{id} \in F(k[\epsilon])$ 表示单位元。对任意
$v \in F(k[\epsilon])$，有
$+(v, \text{id}) = +(\text{id}, v) = v$，因为
$(\text{id}, v)$ 是 $v$ 沿环映射
$k[\epsilon] \to k[\epsilon] \times_k k[\epsilon]$ 的推前，
该映射满足 $\epsilon \mapsto (\epsilon, 0)$。
一般地，给定乘法为 $\circ$ 的群 $G$，若同态
$+ : G \times G \to G$ 满足
$+(g, 1) = +(1, g) = g$，其中 $1$ 是 $G$ 的单位元，
则 $+ = \circ$。% P10-E0482
这说明 $F(k[\epsilon])$ 的 $k$-向量空间结构中的加法与复合一致。

\medskip\noindent
最后，(2) 只需展开定义。具体地，
$T_{\text{id}_{x_0}} \mathit{Aut}(x)$ 是如下自同构
$\alpha$ 所成的集合：它们是 $x$ 沿
$A \to k \to k[\epsilon]$ 的推前的自同构，并且模
$\epsilon$ 为恒等。另一方面，
$T_{\text{id}_{x_0}} \mathit{Aut}(x_0)$ 是 $x_0$ 沿
$k \to k[\epsilon]$ 的推前的、模 $\epsilon$ 为恒等的自同构集合。
由于 $x_0$ 是 $x$ 沿 $A \to k$ 的推前，结论显然。
\end{proof}

\begin{remark}
\label{formal-defos-remark-infaut-lifting-equalities}
下面指出无穷小自同构群、提升以及自同构函子切空间之间的一些基本关系。
设 $\mathcal{F}$ 是 $\mathcal{C}_\Lambda$ 上的群胚余纤维范畴。
设 $x' \to x$ 是位于环映射 $A' \to A$ 上的态射。
由定义有等式
$$
\text{Inf}(x'/x) = \text{Lift}(\text{id}_x, A'),
$$
其中右端是把 $\text{id}_x$ 视为 $\mathit{Aut}(x')$ 的对象时
所得的提升。若 $x_0 \in \Ob(\mathcal{F}(k))$，且 $x'_0$
是到 $\mathcal{F}(k[\epsilon])$ 的推前，则对
$x'_0 \to x_0$ 应用上述等式，得到
$$
\text{Inf}_{x_0}(\mathcal{F}) =
\text{Lift}(\text{id}_{x_0}, k[\epsilon]) =
T_{\text{id}_{x_0}} \mathit{Aut}(x_0),
$$
最后一个等式直接由定义得出。
\end{remark}

\begin{lemma}
\label{formal-defos-lemma-infaut-vector-space}
设 $\mathcal{F}$ 是 $\mathcal{C}_\Lambda$ 上满足 (RS)
的群胚余纤维范畴，并设 $x_0 \in \Ob(\mathcal{F}(k))$。
则作为集合，$\text{Inf}_{x_0}(\mathcal{F})$ 等于
$T_{\text{id}_{x_0}} \mathit{Aut}(x_0)$，因而具有自然的
$k$-向量空间结构，使得加法与自同构的复合一致。
\end{lemma}

\begin{proof}
集合的等式如说明 \ref{formal-defos-remark-infaut-lifting-equalities} 末尾所述，
而关于向量空间结构的陈述由引理
\ref{formal-defos-lemma-Aut-functor-tangent-space} 得出。
\end{proof}

\begin{lemma}
\label{formal-defos-lemma-k-linear-infaut}
设 $\varphi : \mathcal{F} \to \mathcal{G}$ 是
$\mathcal{C}_\Lambda$ 上满足 (RS) 的群胚余纤维范畴之间的态射。
设 $x_0 \in \Ob(\mathcal{F}(k))$。则 $\varphi$ 诱导
$k$-线性映射
$\text{Inf}_{x_0}(\mathcal{F}) \to
\text{Inf}_{\varphi(x_0)}(\mathcal{G})$。
\end{lemma}

\begin{proof}
显然，$\varphi$ 诱导把恒等映到恒等的态射
$\mathit{Aut}(x_0) \to \mathit{Aut}(\varphi(x_0))$。
因此该结论由切空间的相应结论得出；见引理
\ref{formal-defos-lemma-k-linear-differential}。
\end{proof}

\begin{lemma}
\label{formal-defos-lemma-lifted-automorphisms-torsor}
设 $\mathcal{F}$ 是 $\mathcal{C}_\Lambda$ 上满足 (RS)
的群胚余纤维范畴。设 $x' \to x$ 是位于满环映射
$A' \to A$ 上的态射，该映射的核 $I$ 被
$\mathfrak m_{A'}$ 零化。设 $x_0$ 是 $x$ 到
$\mathcal{F}(k)$ 的一个推前。则
$T_{\text{id}_{x_0}} \mathit{Aut}(x') \otimes_k I
= \text{Inf}_{x_0}(\mathcal{F}) \otimes_k I$
在 $\text{Inf}(x'/x)$ 上有一个自由且传递的作用。
\end{lemma}

\begin{proof}
这就是引理 \ref{formal-defos-lemma-free-transitive-action} 在自同构层情形的类似结论。
确切地，对函子
$\mathit{Aut}(x') : \mathcal{C}_{A'} \to \textit{Sets}$
和元素 $\text{id}_{x_0} \in \mathit{Aut}(x)(k)$ 应用说明
\ref{formal-defos-remark-localize-cofibered-groupoid}，得到前形变函子
$F = \mathit{Aut}(x')_{\text{id}_{x_0}}$。% P10-E0483
由引理 \ref{formal-defos-lemma-Aut-functor-RS} 和 \ref{formal-defos-lemma-localize-RS}，
$F$ 是形变函子。因此，引理
\ref{formal-defos-lemma-free-transitive-action} 给出
$TF \otimes_k I$ 在 $\text{Lift}(\text{id}_x, A')$ 上的
自由且传递作用；由于 $\text{Lift}(\text{id}_x, A')$
是群，它总是非空。注意，由引理
\ref{formal-defos-lemma-Aut-functor-tangent-space}，有向量空间的等式
$$
TF = T_{\text{id}_{x_0}} \mathit{Aut}(x') \otimes_k I =
\text{Inf}_{x_0}(\mathcal{F}) \otimes_k I.
$$ % P10-E0484
最后，说明 \ref{formal-defos-remark-infaut-lifting-equalities} 中的等式
$\text{Inf}(x'/x) = \text{Lift}(\text{id}_x, A')$
完成证明。
\end{proof}

\begin{lemma}
\label{formal-defos-lemma-infaut-trivial}
设 $\mathcal{F}$ 是 $\mathcal{C}_\Lambda$ 上满足 (RS)
的群胚余纤维范畴。设 $x' \to x$ 是 $\mathcal{F}$ 中位于
一个满环映射上的态射。设 $x_0$ 是 $x$ 到
$\mathcal{F}(k)$ 的一个推前。若
$\text{Inf}_{x_0}(\mathcal{F}) = 0$，则
$\text{Inf}(x'/x) = 0$。
\end{lemma}

\begin{proof}
由引理 \ref{formal-defos-lemma-factor-small-extension} 和
\ref{formal-defos-lemma-lifted-automorphisms-torsor} 得出。
\end{proof}

\begin{lemma}
\label{formal-defos-lemma-infdef-trivial}
设 $\mathcal{F}$ 是 $\mathcal{C}_\Lambda$ 上满足 (RS)
的群胚余纤维范畴，并设
$x_0 \in \Ob(\mathcal{F}(k))$。则
$\text{Inf}_{x_0}(\mathcal{F}) = 0$，当且仅当群胚余纤维范畴的自然态射
$\mathcal{F}_{x_0} \to \overline{\mathcal{F}_{x_0}}$ 是等价。
\end{lemma}

\begin{proof}
态射 $\mathcal{F}_{x_0} \to \overline{\mathcal{F}_{x_0}}$ 是等价，
当且仅当 $\mathcal{F}_{x_0}$ 纤维化于集合胚；参见《范畴》第
\ref{categories-section-fibred-in-setoids} 节（根据定义，集合胚是
每个对象的自同构都只有恒等映射的群胚）。我们证明
$\text{Inf}_{x_0}(\mathcal{F}) = 0$ 当且仅当
$\mathcal{F}_{x_0}$ 满足此条件。显然，若
$\mathcal{F}_{x_0}$ 纤维化于集合胚，则
$\text{Inf}_{x_0}(\mathcal{F}) = 0$。反过来，假设
$\text{Inf}_{x_0}(\mathcal{F}) = 0$。设 $A$ 是
$\mathcal{C}_\Lambda$ 的对象。由引理
\ref{formal-defos-lemma-infaut-trivial}，对 $\mathcal{F}_{x_0}(A)$ 中的
任意对象 $x \to x_0$，都有 $\text{Inf}(x/x_0) = 0$。
根据定义，$\text{Inf}(x/x_0)$ 等于
$\mathcal{F}_{x_0}(A)$ 中 $x \to x_0$ 的自同构群，
所以 $\mathcal{F}_{x_0}(A)$ 是集合胚。
\end{proof}

\section{应用}
\label{formal-defos-section-applications}

\noindent
这里汇集一些稍后将用到的形变范畴结果。

\begin{lemma}
\label{formal-defos-lemma-deformation-categories-fiber-product-morphisms}
设 $f : \mathcal{H} \to \mathcal{F}$ 和
$g : \mathcal{G} \to \mathcal{F}$ 是形变范畴的 $1$-态射。则
\begin{enumerate}
\item $\mathcal{W} = \mathcal{H} \times_\mathcal{F} \mathcal{G}$
是形变范畴；并且
\item 有向量空间的 $6$ 项正合序列
$$
0 \to \text{Inf}(\mathcal{W})
\to \text{Inf}(\mathcal{H}) \oplus \text{Inf}(\mathcal{G})
\to \text{Inf}(\mathcal{F}) \to
T\mathcal{W} \to T\mathcal{H} \oplus T\mathcal{G} \to T\mathcal{F}.
$$
\end{enumerate}
\end{lemma}

\begin{proof}
(1) 由引理 \ref{formal-defos-lemma-RS-fiber-product-morphisms} 以及下述事实得出：
$\mathcal{W}(k)$ 是两个具有唯一同构类的集合胚在另一个具有唯一同构类的
集合胚上的纤维积。

\medskip\noindent
证明 (2)。设 $w_0 \in \Ob(\mathcal{W}(k))$，并设
$x_0, y_0, z_0$ 是 $w_0$ 在
$\mathcal{F}, \mathcal{H}, \mathcal{G}$ 中的像。
则 $\text{Inf}(\mathcal{W}) = \text{Inf}_{w_0}(\mathcal{W})$，
对 $\mathcal{H}$、$\mathcal{G}$ 和 $\mathcal{F}$ 也类似；见说明
\ref{formal-defos-remark-trivial-aut-point}。应用引理
\ref{formal-defos-lemma-k-linear-differential} 和
\ref{formal-defos-lemma-k-linear-infaut}，得到除“边界映射”
$\delta : \text{Inf}_{x_0}(\mathcal{F}) \to T\mathcal{W}$
之外的所有线性映射。稍后再插入适当的符号。

\medskip\noindent
构造 $\delta$。选取沿 $k \to k[\epsilon]$
的一个推前 $w_0 \to w'_0$。以 $x'_0,y'_0,z'_0$ 表示
$w'_0$ 在 $\mathcal{F},\mathcal{H},\mathcal{G}$ 中的像。
特别地，得到同构 $b' : f(y'_0) \to x'_0$ 和
$c' : x'_0 \to g(z'_0)$。以 $b : f(y_0) \to x_0$ 和
$c : x_0 \to g(z_0)$ 表示它们沿 $k[\epsilon] \to k$ 的推前。
注意，按照范畴纤维积的显式形式，这意味着
$w'_0 = (k[\epsilon], y'_0, z'_0, c' \circ b')$ 且
$w_0 = (k, y_0, z_0, c \circ b)$；见说明
\ref{formal-defos-remarks-cofibered-groupoids}（\ref{formal-defos-item-fibre-product}）。
给定 $\alpha : x'_0 \to x'_0$，置
$\delta(\alpha) = (k[\epsilon], y'_0, z'_0, c' \circ \alpha \circ b')$。
这确实是 $\mathcal{W}$ 中位于 $k[\epsilon]$ 上的对象，并且带有
位于 $k[\epsilon] \to k$ 上的态射
$(k[\epsilon], y'_0, z'_0, c' \circ \alpha \circ b') \to w_0$，
因为 $\alpha$ 推前到 $k$ 上成为恒等。更一般地，对任意
$k$-向量空间 $V$，可以用完全相同的公式定义映射
$$
\text{Lift}(\text{id}_{x_0}, k[V])
\longrightarrow
\text{Lift}(w_0, k[V]).
$$
此构造关于向量空间 $V$ 具有函子性（略去细节）。因此，
应用引理 \ref{formal-defos-lemma-morphism-linear-functors} 可知
$\delta$ 是 $k$-线性的。

\medskip\noindent
构造这些映射后，容易证明序列正合。第一个映射的单射性来自
$f \times g : \mathcal{W} \to \mathcal{H} \times \mathcal{G}$
是忠实的这一事实。若
$(\beta, \gamma) \in
\text{Inf}_{y_0}(\mathcal{H}) \oplus \text{Inf}_{z_0}(\mathcal{G})$
映到 $\text{Inf}_{x_0}(\mathcal{F})$ 中的同一元素，则
$(\beta, \gamma)$ 定义 $w'_0 =
(k[\epsilon], y'_0, z'_0, c' \circ b')$ 的自同构，
从而得到第二处的正合性。若上述 $\alpha$ 给出 $w_0$ 的平凡形变
$(k[\epsilon], y'_0, z'_0, c' \circ \alpha \circ b')$，
则同构
$w'_0 = (k[\epsilon], y'_0, z'_0, c' \circ b') \to
(k[\epsilon], y'_0, z'_0, c' \circ \alpha \circ b')$
产生一个作为 $\alpha$ 原像的对 $(\beta, \gamma)$。
若 $w = (k[\epsilon], y, z, \phi)$ 是 $w_0$ 的形变，且
$y'_0 \cong y$ 与 $z \cong z'_0$，则映射
$$
f(y'_0) \to f(y) \xrightarrow{\phi} g(z) \to g(z'_0)
$$
是一个在 $\delta$ 下映到 $w$ 的 $\alpha$。最后，若 $y$ 与 $z$
分别是 $y_0$ 与 $z_0$ 的形变，并且存在 $f(y_0)=x_0=g(z_0)$
的形变之间的同构 $\phi : f(y) \to g(z)$，则得到
$T\mathcal{W}$ 中 $(x, y)$ 的原像
$w = (k[\epsilon], y, z, \phi)$。% P10-E0485
证明完毕。
\end{proof}

\begin{lemma}
\label{formal-defos-lemma-map-fibre-products-smooth}
设 $\mathcal{H}_1 \to \mathcal{G}$、
$\mathcal{H}_2 \to \mathcal{G}$ 和 $\mathcal{G} \to \mathcal{F}$
是 $\mathcal{C}_\Lambda$ 上群胚余纤维范畴之间的映射。假设
\begin{enumerate}
\item $\mathcal{F}$ 和 $\mathcal{G}$ 是形变范畴；
\item $T\mathcal{G} \to T\mathcal{F}$ 是单射；并且
\item $\text{Inf}(\mathcal{G}) \to \text{Inf}(\mathcal{F})$ 是满射。
\end{enumerate}
则
$\mathcal{H}_1 \times_\mathcal{G} \mathcal{H}_2 \to
\mathcal{H}_1 \times_\mathcal{F} \mathcal{H}_2$ 是光滑的。
\end{lemma}

\begin{proof}
以 $p_i : \mathcal{H}_i \to \mathcal{G}$ 和
$q : \mathcal{G} \to \mathcal{F}$ 表示给定的映射。% P10-E0486
设 $A' \to A$ 是 $\mathcal{C}_\Lambda$ 中的小扩张。
$\mathcal{H}_1 \times_\mathcal{F} \mathcal{H}_2$ 中位于
$A'$ 上的对象，是三元组 $(x'_1, x'_2, a')$，其中
$x'_i$ 是 $\mathcal{H}_i$ 中位于 $A'$ 上的对象，而
$a' : q(p_1(x'_1)) \to q(p_2(x'_2))$ 是 $\mathcal{F}$
在 $A'$ 上的纤维范畴中的态射。沿 $A' \to A$ 推前，
得到 $(x_1,x_2,a)$。将它提升为
$\mathcal{H}_1 \times_\mathcal{G} \mathcal{H}_2$ 中位于 $A$
上的对象，意指找到位于 $A$ 上且满足 $q(b)=a$ 的态射
$b : p_1(x_1) \to p_2(x_2)$。因此必须证明，可以把 $b$
提升为态射 $b' : p_1(x'_1) \to p_2(x'_2)$，且其在 $q$ 下的像是
$a'$。

\medskip\noindent
注意，可以把
$$
p_1(x'_1) \to p_1(x_1) \xrightarrow{b} p_2(x_2)
\quad\text{和}\quad
p_2(x'_2) \to p_2(x_2)
$$
视为 $\textit{Lift}(p_2(x_2), A' \to A)$ 的两个对象。
函子 $q$ 把它们送到
$$
q(p_1(x'_1)) \to q(p_1(x_1)) \xrightarrow{b} q(p_2(x_2))
\quad\text{和}\quad % P10-E0487
q(p_2(x'_2)) \to q(p_2(x_2))
$$
这两个 $\textit{Lift}(q(p_2(x_2)), A' \to A)$ 的对象；
利用映射 $a' : q(p_1(x'_1)) \to q(p_2(x'_2))$，
它们彼此同构。另一方面，函子
$$
q : \textit{Lift}(p_2(x_2), A' \to A) \to
\textit{Lift}(q(p_2(x_2)), A' \to A)
$$
由引理 \ref{formal-defos-lemma-free-transitive-action} 和关于切空间的假设，
在同构类上诱导单射。因此存在态射
$b' : p_1(x_1') \to p_2(x'_2)$，其到 $A$ 的推前是 $b$。
不过，可能还需要调整 $b'$ 的选择，才能使 $q(b')=a'$。
为此，只需证明
$q : \text{Inf}(p_2(x'_2)/p_2(x_2)) \to
\text{Inf}(q(p_2(x'_2))/q(p_2(x_2)))$ 是满射。
这由关于无穷小自同构的假设和引理
\ref{formal-defos-lemma-lifted-automorphisms-torsor} 得出。
\end{proof}

\begin{lemma}
\label{formal-defos-lemma-easy-check-smooth}
设 $f : \mathcal{F} \to \mathcal{G}$ 是形变范畴之间的映射。
设 $x_0 \in \Ob(\mathcal{F}(k))$，其像为
$y_0 \in \Ob(\mathcal{G}(k))$。若
\begin{enumerate}
\item 映射 $T\mathcal{F} \to T\mathcal{G}$ 是满射；并且
\item 对 $\mathcal{C}_\Lambda$ 中的每个小扩张
$A' \to A$，以及像为 $y \in \mathcal{G}(A)$ 的每个
$x \in \mathcal{F}(A)$，若 $y$ 可以提升到 $A'$，
则 $x$ 也可以提升到 $A'$，
\end{enumerate}
则 $\mathcal{F} \to \mathcal{G}$ 光滑（反之亦然）。
\end{lemma}

\begin{proof}
设 $A' \to A$ 是小扩张，且 $x \in \mathcal{F}(A)$。
设 $y' \to f(x)$ 是 $\mathcal{G}$ 中位于 $A' \to A$ 上的态射。
考虑 $f$ 所诱导的函子
$\text{Lift}(A', x) \to \text{Lift}(A', f(x))$。% P10-E0488
必须证明存在 $\text{Lift}(A', x)$ 的对象 $x' \to x$，
其像为 $y' \to f(x)$；见引理
\ref{formal-defos-lemma-smoothness-small-extensions}。由条件 (2)，
$\text{Lift}(A', x)$ 不是空范畴。由条件 (2) 和引理
\ref{formal-defos-lemma-free-transitive-action}，可知同构类上的映射是满射，
正如所需。% P10-E0489 P10-E0490
\end{proof}

\begin{lemma}
\label{formal-defos-lemma-map-between-smooth}
设 $\mathcal{F} \to \mathcal{G} \to \mathcal{H}$ 是
$\mathcal{C}_\Lambda$ 上群胚余纤维范畴之间的映射。若
\begin{enumerate}
\item $\mathcal{F}$、$\mathcal{G}$ 是形变范畴；
\item 映射 $T\mathcal{F} \to T\mathcal{G}$ 是满射；并且
\item $\mathcal{F} \to \mathcal{H}$ 光滑，
\end{enumerate}
则 $\mathcal{F} \to \mathcal{G}$ 光滑。
\end{lemma}

\begin{proof}
设 $A' \to A$ 是 $\mathcal{C}_\Lambda$ 中的小扩张，
并设 $x \in \mathcal{F}(A)$，其像为
$y \in \mathcal{G}(A)$。假设存在提升
$y' \in \mathcal{G}(A')$。根据引理
\ref{formal-defos-lemma-easy-check-smooth}，只需检验 $x$ 也有提升。
取 $y'$ 的像 $z' \in \mathcal{H}(A')$。由于
$\mathcal{F} \to \mathcal{H}$ 光滑，存在
$x' \in \mathcal{F}(A')$，它同时映到
$x \in \mathcal{F}(A)$ 和 $z' \in \mathcal{H}(A')$；
见定义 \ref{formal-defos-definition-smooth-morphism}。证明完毕。
\end{proof}

\section{任意范畴上的函子中的群胚}
\label{formal-defos-section-groupoids-arbitrary}

\noindent
首先讨论任意范畴上的函子中的群胚的一般性质。下一节再转到范畴
$\mathcal{C}_\Lambda$。为清楚起见，有时把通常的群胚，
即所有态射都是同构的范畴，称为群胚范畴。

\begin{definition}
\label{formal-defos-definition-groupoid-in-functors}
设 $\mathcal{C}$ 是范畴。$\mathcal{C}$ 上的{\it 函子中的群胚范畴}
具有如下对象和态射。
\begin{enumerate}
\item 对象：$\mathcal{C}$ 上的一个{\it 函子中的群胚}是五元组
$(U, R, s, t, c)$，其中 $U, R : \mathcal{C} \to \textit{Sets}$
是函子，而 $s, t : R \to U$ 和
$c : R \times_{s, U, t} R \to R$ 是具有如下性质的态射：
对 $\mathcal{C}$ 的任意对象 $T$，五元组
$$
(U(T), R(T), s, t, c)
$$
是群胚范畴。
\item 态射：$\mathcal{C}$ 上函子中的群胚之间的一个
{\it 态射 $(U, R, s, t, c) \to (U', R', s', t', c')$}
由态射 $U \to U'$ 和 $R \to R'$ 组成，并具有如下性质：
对 $\mathcal{C}$ 的任意对象 $T$，诱导映射
$U(T) \to U'(T)$ 与 $R(T) \to R'(T)$ 定义群胚范畴之间的函子
$$
(U(T), R(T), s, t, c) \to (U'(T), R'(T), s', t', c').
$$
\end{enumerate}
\end{definition}

\begin{remark}
\label{formal-defos-remark-confusion-groupoids-in-functors}
$\mathcal{C}$ 上函子中的群胚等价于函子
$\mathcal{C} \to \textit{Groupoids}$ 的数据，而
$\mathcal{C}$ 上函子中的群胚态射等价于相应函子
$\mathcal{C} \to \textit{Groupoids}$ 之间的态射
（这里把 $\textit{Groupoids}$ 看作 1-范畴）。
不过，为了我们的目的，使用函子中的群胚这一术语更方便。
事实上，把函子中的群胚想成相应的函子
$\mathcal{C} \to \textit{Groupoids}$，或等价地想成与该函子关联的
群胚余纤维范畴，可能导致混淆（说明
\ref{formal-defos-remark-smooth-groupoid-in-functors-warning}）。
\end{remark}

\begin{remark}
\label{formal-defos-remark-identity-inverse}
设 $(U, R, s, t, c)$ 是范畴 $\mathcal{C}$ 上函子中的群胚。
存在唯一的态射 $e : U \to R$ 和 $i : R \to R$，使得对
$\mathcal{C}$ 的每个对象 $T$，$e: U(T) \to R(T)$ 把
$x \in U(T)$ 送到 $x$ 上的恒等态射，而
$i: R(T) \to R(T)$ 把 $a \in U(T)$ 送到 $a$ 在群胚范畴
$(U(T), R(T), s, t, c)$ 中的逆。% P10-E0491
有时分别把 $s$、$t$、$c$、$e$、$i$ 称为“源”、“靶”、
“复合”、“恒等”、“逆”。
\end{remark}

\begin{definition}
\label{formal-defos-definition-representable}
设 $\mathcal{C}$ 是范畴。若 $\mathcal{C}$ 上函子中的群胚
同构于形如 $(\underline{U}, \underline{R}, s, t, c)$ 的群胚，
其中 $U$ 和 $R$ 是 $\mathcal{C}$ 的对象，且推出
$R \amalg_{s, U, t} R$ 存在，则称它{\it 可表}。
\end{definition}

\begin{remark}
\label{formal-defos-remark-reason-existence-coproduct}
因此，$\mathcal{C}$ 上函子中的可表群胚由
$\mathcal{C}$ 的对象 $U$、$R$ 以及态射 $s,t:U\to R$、
$c:R\to R\amalg_{s,U,t}R$ 给出，并且
$(\underline{U}, \underline{R}, s, t, c)$ 满足定义
\ref{formal-defos-definition-groupoid-in-functors} 的条件。要求推出
$R \amalg_{s, U, t} R$ 存在，是为了使复合态射 $c$
能在 $\mathcal{C}$ 的态射层面定义。下文考虑
$\widehat{\mathcal{C}}_\Lambda$ 上函子中的可表群胚时，
此要求总会满足，因为由引理
\ref{formal-defos-lemma-CLambdahat-pushouts}，范畴
$\widehat{\mathcal{C}}_\Lambda$ 具有推出。
\end{remark}

\begin{remark}
\label{formal-defos-remark-simplify-terminology}
我们会用“{\it 设 $(\underline{U}, \underline{R}, s, t, c)$ 是
$\mathcal{C}$ 上函子中的群胚}”来表示所给的是函子中的可表群胚。
因此，这意味着 $U$、$R$ 是 $\mathcal{C}$ 的对象，存在态射
$s,t:U\to R$，推出 $R\amalg_{s,U,t}R$ 存在，
存在态射 $c:R\to R\amalg_{s,U,t}R$，并且
$(\underline{U}, \underline{R}, s, t, c)$ 是
$\mathcal{C}$ 上函子中的群胚。
\end{remark}

\noindent
下面引入函子中群胚的限制记号。稍后把
$\widehat{\mathcal
C}_\Lambda$ 限制到 $\mathcal{C}_\Lambda$ 时将用到它。

\begin{definition}
\label{formal-defos-definition-restricting-groupoids-in-functors}
设 $(U, R, s, t, c)$ 是范畴 $\mathcal{C}$ 上函子中的群胚，
并设 $\mathcal{C}'$ 是 $\mathcal{C}$ 的子范畴。
$(U, R, s, t, c)$ 到 $\mathcal{C}'$ 的{\it 限制
$(U, R, s, t, c)|_{\mathcal{C}'}$}，是 $\mathcal{C}'$
上函子中的群胚
$(U|_{\mathcal{C}'}, R|_{\mathcal
C'}, s|_{\mathcal{C}'}, t|_{\mathcal{C}'}, c|_{\mathcal{C}'})$。
\end{definition}

\begin{remark}
\label{formal-defos-remark-notation-restriction}
在定义 \ref{formal-defos-definition-restricting-groupoids-in-functors} 的情形，
常把 $s|_{\mathcal{C}'}, t|_{\mathcal{C}'}, c|_{\mathcal{C}'}$
简记为 $s,t,c$。
\end{remark}

\begin{definition}
\label{formal-defos-definition-quotient}
设 $(U, R, s, t, c)$ 是范畴 $\mathcal{C}$ 上函子中的群胚。
\begin{enumerate}
\item 对应 $T \mapsto (U(T), R(T), s, t, c)$ 确定函子
$\mathcal{C} \to \textit{Groupoids}$。{\it 商群胚余纤维范畴
$[U/R] \to \mathcal{C}$}，是与该函子关联的
$\mathcal{C}$ 上群胚余纤维范畴（如说明
\ref{formal-defos-remarks-cofibered-groupoids}（\ref{formal-defos-item-construction-associated-cofibered-groupoid}）
所述）。
\item {\it 商态射 $U \to [U/R]$} 是由如下规则定义的
$\mathcal{C}$ 上群胚余纤维范畴之间的态射：
\begin{enumerate}
\item $x \in U(T)$ 映到对象 $(T, x) \in \Ob([U/R](T))$；并且
\item $x \in U(T)$ 和 $f : T \to T'$ 给出位于
$f : T \to T'$ 上的态射
$(f, \text{id}_{U(f)(x)}): (T, x) \to (T, U(f)(x))$。
\end{enumerate}
\end{enumerate}
\end{definition}

\section{基范畴上的函子中的群胚}
\label{formal-defos-section-prorepresentable-groupoids-in-functors}

\noindent
本节讨论 $\mathcal{C}_\Lambda$ 上的函子中的群胚。
最终目标是说明，$\mathcal{C}_\Lambda$ 上函子中的 pro-可表群胚，
对性质良好的形变范畴起“呈现”作用，正如代数空间中的光滑群胚
对代数栈起呈现作用一样；参见《代数栈》第
\ref{algebraic-section-stack-to-presentation} 节。

\begin{definition}
\label{formal-defos-definition-prorepresentable-groupoid-in-functors}
若 $\mathcal{C}_\Lambda$ 上函子中的群胚同构于
$(\underline{R_0}, \underline{R_1}, s, t, c)|_{\mathcal{C}_\Lambda}$，
且 $(\underline{R_0}, \underline{R_1}, s, t, c)$ 是范畴
$\widehat{\mathcal{C}}_\Lambda$ 上函子中的某个可表群胚，
则称原群胚{\it pro-可表}。
\end{definition}

\noindent
设 $(U, R, s, t, c)$ 是 $\mathcal{C}_\Lambda$
上的函子中的群胚。取完备化，得到五元组
$(\widehat{U}, \widehat{R}, \widehat{s}, \widehat{t}, \widehat{c})$。
由说明 \ref{formal-defos-remark-completion-restriction-cofset-adjoint}，
$\text{CofSet}(\mathcal{C}_\Lambda)$ 上的完备化函子是右伴随，
所以与极限可交换。特别地，有典范同构
$$
\widehat{R \times_{s, U, t} R}
\longrightarrow
\widehat{R} \times_{\widehat{s}, \widehat{U}, \widehat{t}} \widehat{R},
$$
故可把 $\widehat{c}$ 看作函子
$\widehat{R} \times_{\widehat{s}, \widehat{U}, \widehat{t}} \widehat{R} \to
\widehat{R}$。于是
$(\widehat{U}, \widehat{R}, \widehat{s}, \widehat{t}, \widehat{c})$
是 $\widehat{\mathcal{C}}_\Lambda$ 上函子中的群胚，
其恒等态射和逆态射分别是 $(U,R,s,t,c)$ 相应态射的完备化。

\begin{definition}
\label{formal-defos-definition-completion-groupoid-in-functors}
设 $(U, R, s, t, c)$ 是 $\mathcal{C}_\Lambda$
上的函子中的群胚。$(U,R,s,t,c)$ 的{\it 完备化
$(U, R, s, t, c)^{\wedge}$}，是上面所述的
$\widehat{\mathcal{C}}_\Lambda$ 上函子中的群胚
$(\widehat{U}, \widehat{R}, \widehat{s}, \widehat{t}, \widehat{c})$。
\end{definition}

\begin{remark}
\label{formal-defos-remark-groupoid-in-functors-complete-restrict}
设 $(U, R, s, t, c)$ 是 $\mathcal{C}_\Lambda$
上的函子中的群胚。则有典范同构
$(U, R, s, t, c)^{\wedge}|_{\mathcal{C}_\Lambda}
\cong (U, R, s, t, c)$；见说明
\ref{formal-defos-remark-restrict-completion}。另一方面，设
$(U, R, s, t, c)$ 是 $\widehat{\mathcal{C}}_\Lambda$
上的函子中的群胚，并且
$U, R : \widehat{\mathcal{C}}_\Lambda \to \textit{Sets}$
都与极限可交换，例如 $U,R$ 可表。则有典范同构
$((U, R, s, t, c)|_{\mathcal{C}_\Lambda})^{\wedge}
\cong (U, R, s, t, c)$。这由说明
\ref{formal-defos-remark-restrict-complete-continuous-functor} 得出。
\end{remark}

\begin{lemma}
\label{formal-defos-lemma-groupoid-in-functors-prorep-equivalences}
设 $(U, R, s, t, c)$ 是 $\mathcal{C}_\Lambda$
上的函子中的群胚。
\begin{enumerate}
\item $(U,R,s,t,c)$ pro-可表，当且仅当其完备化作为
$\widehat{\mathcal{C}}_\Lambda$ 上函子中的群胚可表。
\item $(U,R,s,t,c)$ pro-可表，当且仅当 $U$ 和 $R$ pro-可表。
\end{enumerate}
\end{lemma}

\begin{proof}
(1) 由说明 \ref{formal-defos-remark-groupoid-in-functors-complete-restrict} 得出。
对 (2)，“仅当”方向由函子中 pro-可表群胚的定义立即得出。
反过来，假设 $U$、$R$ pro-可表，即对
$\widehat{\mathcal{C}}_\Lambda$ 的对象 $R_0$、$R_1$，有
$U \cong \underline{R_0}|_{\mathcal{C}_\Lambda}$ 且
$R \cong \underline{R_1}|_{\mathcal{C}_\Lambda}$。由说明
\ref{formal-defos-remark-restrict-complete-continuous-functor}，有
$\underline{R_0} \cong \widehat{\underline{R_0}|_{\mathcal{C}_\Lambda}}$
和
$\underline{R_1} \cong \widehat{\underline{R_1}|_{\mathcal{C}_\Lambda}}$，
所以完备化 $(U,R,s,t,c)^\wedge$ 是形如
$(\underline{R_0}, \underline{R_1}, \widehat{s}, \widehat{t}, \widehat{c})$
的函子中的群胚。由引理 \ref{formal-defos-lemma-CLambdahat-pushouts}，推出
$\underline{R_1} \times_{\widehat{s}, \underline{R_1}, \widehat{t}}
\underline{R_1}$ 存在。% P10-E0492
所以
$(\underline{R_0}, \underline{R_1}, \widehat{s}, \widehat{t}, \widehat{c})$
是 $\widehat{\mathcal{C}}_\Lambda$ 上函子中的可表群胚。
最后，由说明 \ref{formal-defos-remark-groupoid-in-functors-complete-restrict}，
限制
$(\underline{R_0}, \underline{R_1}, s, t, c)|_{\mathcal{C}_\Lambda}$
恢复 $(U,R,s,t,c)$，故根据定义 $(U,R,s,t,c)$ pro-可表。
\end{proof}

\section{基范畴上的函子中的光滑群胚}
\label{formal-defos-section-smooth-minimal-groupoids-in-functors}

\noindent
$\mathcal{C}_\Lambda$ 上函子中群胚的光滑性概念定义如下。

\begin{definition}
\label{formal-defos-definition-smooth-groupoid-in-functors}
设 $(U, R, s, t, c)$ 是 $\mathcal{C}_\Lambda$
上的函子中的群胚。若 $s,t:R\to U$ 光滑，
则称 $(U,R,s,t,c)$ {\it 光滑}。
\end{definition}

\begin{remark}
\label{formal-defos-remark-smooth-groupoid-in-functors-warning}
注意，此术语可能引起混淆：若 $(U,R,s,t,c)$ 是函子中的光滑群胚，
则商 $[U/R]$ 未必是定义
\ref{formal-defos-definition-cofibered-groupoid-projection-smooth} 所指的
光滑群胚余纤维范畴。不过，$(U,R,s,t,c)$ 的光滑性确实蕴含
商态射 $U\to[U/R]$ 光滑（事实上二者等价），正如引理
\ref{formal-defos-lemma-smooth-quotient-morphism} 将要说明的那样。
\end{remark}

\begin{remark}
\label{formal-defos-remark-smooth-power-series-prorepresentable-smooth-groupoid-in-functors}
设
$(\underline{R_0}, \underline{R_1}, s, t, c)|_{\mathcal{C}_\Lambda}$
是 $\mathcal{C}_\Lambda$ 上函子中的 pro-可表群胚。
则 $(\underline{R_0}, \underline{R_1}, s, t, c)|_{\mathcal{C}_\Lambda}$
光滑，当且仅当经 $s$ 和 $t$，$R_1$ 都是
$R_0$ 上的幂级数环。% P10-E0493
这由引理 \ref{formal-defos-lemma-smooth-morphism-power-series} 得出。
\end{remark}

\begin{lemma}
\label{formal-defos-lemma-smooth-quotient-morphism}
设 $(U,R,s,t,c)$ 是 $\mathcal{C}_\Lambda$ 上函子中的群胚。
以下条件等价：
\begin{enumerate}
\item 函子中的群胚 $(U,R,s,t,c)$ 光滑。
\item 态射 $s:R\to U$ 光滑。
\item 态射 $t:R\to U$ 光滑。
\item 商态射 $U\to[U/R]$ 光滑。
\end{enumerate}
\end{lemma}

\begin{proof}
(2) 等价于 (3)，因为 $(U,R,s,t,c)$ 的逆
$i:R\to R$ 是同构，且 $t=s\circ i$。根据定义，(1) 等价于
(2) 与 (3) 同时成立，因而等价于其中任意一个单独成立。

\medskip\noindent
最后证明 (2) 等价于 (4)。展开定义：
\begin{enumerate}
\item[(2)] $s:R\to U$ 光滑等价于下述条件：若
$f:B\to A$ 是 $\mathcal{C}_\Lambda$ 中的满环映射，
$a\in R(A)$、$y\in U(B)$ 且 $s(a)=U(f)(y)$，
则存在 $a'\in R(B)$，使 $R(f)(a')=a$ 且 $s(a')=y$。
\item[(4)] $U\to[U/R]$ 光滑等价于下述条件：若
$f:B\to A$ 是 $\mathcal{C}_\Lambda$ 中的满环映射，且
$(f,a):(B,y)\to(A,x)$ 是 $[U/R]$ 的态射，则存在
$x'\in U(B)$ 和 $b\in R(B)$，使 $s(b)=x'$、$t(b)=y$，
且 $c(a,R(f)(b))=e(x)$。这里 $e:U\to R$ 表示恒等，
而记号 $(f,a)$ 如说明
\ref{formal-defos-remarks-cofibered-groupoids}（\ref{formal-defos-item-construction-associated-cofibered-groupoid}）
所述；特别地，$a\in R(A)$，且
$s(a)=U(f)(y)$、$t(a)=x$。
\end{enumerate}
若 (4) 成立，并给定 (2) 中那样的 $f,a,y$，令 $x=t(a)$，
于是有态射 $(f,a):(B,y)\to(A,x)$。(4) 给出 $x'$ 和 $b$，
而 $a'=i(b)$ 满足 (2) 的要求。反过来，假设 (2) 成立，
并给定 (4) 中那样的 $(f,a):(B,y)\to(A,x)$。
则 (2) 给出 $a'\in R(B)$，而 $x'=t(a')$ 和 $b=i(a')$
满足 (4) 的要求。
\end{proof}

\section{作为函子中群胚之商的形变范畴}
\label{formal-defos-section-deformation-categories-as-quotients}

\noindent
我们讨论 $\mathcal{C}_\Lambda$ 上函子中的群胚应满足哪些条件，
才能保证其商是形变范畴，并计算这种商的切空间和无穷小自同构空间。

\begin{lemma}
\label{formal-defos-lemma-smooth-RS-groupoid-in-functors-quotient}
设 $(U,R,s,t,c)$ 是 $\mathcal{C}_\Lambda$ 上函子中的光滑群胚。
假设 $U$ 和 $R$ 满足 (RS)。则 $[U/R]$ 满足 (RS)。
\end{lemma}

\begin{proof}
设
$$
\xymatrix{
                           &     (A_2, x_2) \ar[d]^{(f_2, a_2)} \\
(A_1, x_1) \ar[r]^{(f_1, a_1)} &     (A, x)
}
$$
是 $[U/R]$ 中的图表，其中 $f_2:A_2\to A$ 为满射。
记号如说明 \ref{formal-defos-remarks-cofibered-groupoids}
（\ref{formal-defos-item-construction-associated-cofibered-groupoid}）所述。
因此 $f_1:A_1\to A, f_2:A_2\to A$ 是
$\mathcal{C}_\Lambda$ 中的映射，
$x\in U(A)$、$x_1\in U(A_1)$、$x_2\in U(A_2)$，
并且 $a_1,a_2\in R(A)$，其中
$s(a_1)=U(f_1)(x_1)$、$t(a_1)=x$，且
$s(a_2)=U(f_2)(x_2)$、$t(a_2)=x$。
我们如下构造该 $[U/R]$ 中图表的一个纤维积，
它位于 $A_1\times_A A_2$ 上方。

\medskip\noindent
令 $a=c(i(a_1),a_2)$，其中 $i:R\to R$ 是逆态射。
则 $a\in R(A)$、$x_2\in U(A_2)$，且
$s(a)=U(f_2)(x_2)$。因而有元素
$(a,x_2)\in R(A)\times_{s,U(A),U(f_2)}U(A_2)$。
由于 $s:R\to U$ 光滑，存在元素
$\widetilde{a}\in R(A_2)$，使
$R(f_2)(\widetilde{a})=a$ 且 $s(\widetilde{a})=x_2$。
特别地，
$U(f_2)(t(\widetilde{a}))=t(a)=U(f_1)(x_1)$。
所以 $x_1$ 和 $t(\widetilde{a})$ 确定元素
$$
(x_1, t(\widetilde{a})) \in U(A_1) \times_{U(A)} U(A_2).
$$
由 $U$ 满足 (RS) 的假设，有等同
$U(A_1)\times_{U(A)}U(A_2)=U(A_1\times_A A_2)$。
我们用
$x_1\times t(\widetilde{a})\in U(A_1\times_A A_2)$
表示与
$(x_1,t(\widetilde{a}))\in U(A_1)\times_{U(A)}U(A_2)$
对应的元素。% P10-E0494
令 $p_1,p_2$ 为 $A_1\times_A A_2$ 的投影。我们断言
$$
\xymatrix{
(A_1 \times_A A_2, x_1 \times t(\widetilde{a})) \ar[d]_{(p_1, e(x_1))}
\ar[rr]_-{(p_2, i(\widetilde{a}))} & & (A_2, x_2) \ar[d]^{(f_2, a_2)} \\
(A_1, x_1) \ar[rr]^{(f_1, a_1)} & & (A, x)
}
$$
是 $[U/R]$ 中的纤维方块。（注意，$e:U\to R$ 表示恒等。）

\medskip\noindent
该图表可交换，因为
$c(a_2,R(f_2)(i(\widetilde{a})))=c(a_2,i(a))=a_1$。
为了验证它是纤维方块，设
$$
\xymatrix{
(B, z) \ar[d]_{(g_1, b_1)} \ar[rr]_{(g_2, b_2)} & & (A_2, x_2)
\ar[d]^{(f_2, a_2)} \\
(A_1, x_1) \ar[rr]^{(f_1, a_1)} & & (A, x)
}
$$
是 $[U/R]$ 中的可交换图表。我们将证明存在唯一态射
$(g,b):(B,z)\to(A_1\times_A A_2,x_1\times t(\widetilde{a}))$，
与到 $(A_1,x_1)$ 和 $(A_2,x_2)$ 的态射相容。
必有 $g=(g_1,g_2):B\to A_1\times_A A_2$。
由于根据假设 $R$ 满足 (RS)，有等同
$R(A_1\times_A A_2)=R(A_1)\times_{R(A)}R(A_2)$。
因此可写成 $b=(b'_1,b'_2)$，其中
$b'_1\in R(A_1)$、$b'_2\in R(A_2)$，且二者在 $R(A)$ 中相等。
于是
$((g_1,g_2),(b'_1,b'_2)):(B,z)\to
(A_1\times_A A_2,x_1\times t(\widetilde{a}))$
与投影可交换，当且仅当
$b'_1=b_1$ 且 $b'_2=c(\widetilde{a},b_2)$；
这就证明了唯一性和存在性。
\end{proof}

\begin{lemma}
\label{formal-defos-lemma-deformation-groupoid-quotient}
设 $(U,R,s,t,c)$ 是 $\mathcal{C}_\Lambda$ 上函子中的光滑群胚。
假设 $U$ 和 $R$ 是形变函子。则：
\begin{enumerate}
\item 商 $[U/R]$ 是形变范畴。
\item $[U/R]$ 的切空间为
$$
T[U/R] = \Coker(ds-dt: TR \to TU).
$$
\item $[U/R]$ 的无穷小自同构空间为
$$
\text{Inf}([U/R]) =
\Ker(ds \oplus dt : TR \to TU \oplus TU).
$$
\end{enumerate}
\end{lemma}

\begin{proof}
由于 $U$ 和 $R$ 是形变函子，$[U/R]$ 是预形变范畴。
根据定义，形变函子满足 (RS)，而由引理
\ref{formal-defos-lemma-smooth-RS-groupoid-in-functors-quotient} 可见
[U/R] 满足 (RS)。% P10-E0495
所以 $[U/R]$ 是形变范畴。
(2) 和 (3) 直接由定义得出。
\end{proof}

\section{群胚余纤维范畴的呈现}
\label{formal-defos-section-presentation-categories-cofibred-in-groupoids}

\noindent
呈现定义如下。

\begin{definition}
\label{formal-defos-definition-presentation}
设 $\mathcal{F}$ 是范畴 $\mathcal{C}$ 上的群胚余纤维范畴。
设 $(U,R,s,t,c)$ 是 $\mathcal{C}$ 上函子中的群胚。
所谓{\it 由 $(U,R,s,t,c)$ 给出的 $\mathcal{F}$ 的呈现}，
是指 $\mathcal{C}$ 上群胚余纤维范畴的一个等价
$\varphi:[U/R]\to\mathcal{F}$。
\end{definition}

\noindent
下面两个一般性引理将用于得到呈现。

\begin{lemma}
\label{formal-defos-lemma-presentation-construction}
设 $\mathcal{F}$ 是范畴 $\mathcal{C}$ 上的群胚余纤维范畴。% P10-E0496
设 $U:\mathcal{C}\to\textit{Sets}$ 是函子。
设 $f:U\to\mathcal{F}$ 是 $\mathcal{C}$
上群胚余纤维范畴的态射。如下定义 $R,s,t,c$：
\begin{enumerate}
\item $R:\mathcal{C}\to\textit{Sets}$ 是函子
$U\times_{f,\mathcal{F},f}U$。
\item $t,s:R\to U$ 分别为第一投影和第二投影。
\item $c:R\times_{s,U,t}R\to R$ 是如下态射：
在典范同构
$R\times_{s,U,t}R\to
U\times_{f,\mathcal{F},f}U\times_{f,\mathcal{F},f}U$
下，它由
$U\times_{f,\mathcal{F},f}U\times_{f,\mathcal{F},f}U$
到第一因子和最后一个因子的投影给出。
\end{enumerate}
则 $(U,R,s,t,c)$ 是 $\mathcal{C}$ 上函子中的群胚。
\end{lemma}

\begin{proof}
略。
\end{proof}

\begin{lemma}
\label{formal-defos-lemma-presentation-morphism}
设 $\mathcal{F}$ 是范畴 $\mathcal{C}$ 上的群胚余纤维范畴。% P10-E0496
设 $U:\mathcal{C}\to\textit{Sets}$ 是函子。
设 $f:U\to\mathcal{F}$ 是 $\mathcal{C}$
上群胚余纤维范畴的态射。设 $(U,R,s,t,c)$ 是引理
\ref{formal-defos-lemma-presentation-construction} 中由
$f:U\to\mathcal{F}$ 构造的 $\mathcal{C}$ 上函子中的群胚。
则存在自然态射 $[f]:[U/R]\to\mathcal{F}$，满足：
\begin{enumerate}
\item $[f]:[U/R]\to\mathcal{F}$ 全忠实。
\item $[f]:[U/R]\to\mathcal{F}$ 是等价，当且仅当
$f:U\to\mathcal{F}$ 本质满。
\end{enumerate}
\end{lemma}

\begin{proof}
略。
\end{proof}

\section{形变范畴的呈现}
\label{formal-defos-section-presentation-deformation-categories}

\noindent
根据下一个引理，从预形变函子到预形变范畴 $\mathcal{F}$
的光滑态射，会给出一个由函子中光滑群胚构成的
$\mathcal{F}$ 的呈现。

\begin{lemma}
\label{formal-defos-lemma-smooth-groupoid-in-functors-construction}
设 $\mathcal{F}$ 是 $\mathcal{C}_\Lambda$ 上的群胚余纤维范畴。
设 $U:\mathcal{C}_\Lambda\to\textit{Sets}$ 是函子。
设 $f:U\to\mathcal{F}$ 是群胚余纤维范畴的光滑态射。则：
\begin{enumerate}
\item 若 $(U,R,s,t,c)$ 是引理
\ref{formal-defos-lemma-presentation-construction} 中由
$f:U\to\mathcal{F}$ 构造的 $\mathcal{C}_\Lambda$
上函子中的群胚，则 $(U,R,s,t,c)$ 光滑。
\item 若 $f:U(k)\to\mathcal{F}(k)$ 本质满，则引理
\ref{formal-defos-lemma-presentation-morphism} 中的态射
$[f]:[U/R]\to\mathcal{F}$ 是等价。
\end{enumerate}
\end{lemma}

\begin{proof}
根据引理 \ref{formal-defos-lemma-presentation-construction} 的构造，
有可交换图表
$$
\xymatrix{
R = U \times_{f, \mathcal{F}, f} U \ar[r]_-s \ar[d]_t & U
\ar[d]^f \\
U \ar[r]^f & \mathcal{F}
}
$$
其中 $t,s$ 分别为第一投影和第二投影。于是由引理
\ref{formal-defos-lemma-smooth-properties}，$t,s$ 光滑。因此 (1) 成立。

\medskip\noindent
若 (2) 的假设成立，则由引理
\ref{formal-defos-lemma-smooth-morphism-essentially-surjective}，
态射 $f:U\to\mathcal{F}$ 本质满。因此由引理
\ref{formal-defos-lemma-presentation-morphism}，态射
$[f]:[U/R]\to\mathcal{F}$ 是等价。
\end{proof}

\begin{lemma}
\label{formal-defos-lemma-deformation-functor-diagonal}
设 $\mathcal{F}$ 是形变范畴。
设 $U:\mathcal{C}_\Lambda\to\textit{Sets}$ 是形变函子。
设 $f:U\to\mathcal{F}$ 是群胚余纤维范畴的态射。
则 $U\times_{f,\mathcal{F},f}U$ 是形变函子，
其切空间嵌入如下 $k$-向量空间正合序列：
$$
0 \to \text{Inf}(\mathcal{F}) \to
T(U \times_{f, \mathcal{F}, f} U) \to TU \oplus TU
$$
\end{lemma}

\begin{proof}
这由引理 \ref{formal-defos-lemma-deformation-categories-fiber-product-morphisms}
以及 $\text{Inf}(U)=(0)$ 得出。
\end{proof}

\begin{lemma}
\label{formal-defos-lemma-prorepresentable-groupoid-in-functors-construction}
设 $\mathcal{F}$ 是形变范畴。
设 $U:\mathcal{C}_\Lambda\to\textit{Sets}$ 是 pro-可表函子。
设 $f:U\to\mathcal{F}$ 是群胚余纤维范畴的态射。
设 $(U,R,s,t,c)$ 是引理 \ref{formal-defos-lemma-presentation-construction}
中由 $f:U\to\mathcal{F}$ 构造的
$\mathcal{C}_\Lambda$ 上函子中的群胚。若
$\dim_k\text{Inf}(\mathcal{F})<\infty$，则
$(U,R,s,t,c)$ pro-可表。
\end{lemma}

\begin{proof}
注意，由例 \ref{formal-defos-example-prorepresentable-deformation-functor}，
$U$ 是形变函子。由引理
\ref{formal-defos-lemma-deformation-functor-diagonal} 可见，
$R=U\times_{f,\mathcal{F},f}U$ 是形变函子，其切空间
$TR=T(U\times_{f,\mathcal{F},f}U)$ 位于正合序列
$0\to\text{Inf}(\mathcal{F})\to TR\to TU\oplus TU$ 中。
由于我们已假设第一个空间维数有限，并且由例
\ref{formal-defos-example-tangent-space-prorepresentable-functor}，
$TU$ 的维数有限，所以 $\dim TR<\infty$。
映射 $\gamma:\text{Der}_\Lambda(k,k)\to TR$，见
（\ref{formal-defos-equation-map}），是单射；% P10-E0497
这是因为它与 $TR\to TU$ 的复合，由关于 pro-可表函子 $U$
的定理 \ref{formal-defos-theorem-Schlessinger-prorepresentability} 可知是单射。
于是由定理 \ref{formal-defos-theorem-Schlessinger-prorepresentability}，
$R$ pro-可表。再由引理
\ref{formal-defos-lemma-groupoid-in-functors-prorep-equivalences}，
$(U,R,s,t,c)$ pro-可表。
\end{proof}

\begin{theorem}
\label{formal-defos-theorem-presentation-deformation-groupoid}
设 $\mathcal{F}$ 是 $\mathcal{C}_\Lambda$ 上的群胚余纤维范畴。
则 $\mathcal{F}$ 存在一个由 $\mathcal{C}_\Lambda$
上函子中的光滑 pro-可表群胚给出的呈现，当且仅当以下条件成立：
\begin{enumerate}
\item $\mathcal{F}$ 是形变范畴。
\item $\dim_kT\mathcal{F}$ 有限。
\item $\dim_k\text{Inf}(\mathcal{F})$ 有限。
\end{enumerate}
\end{theorem}

\begin{proof}
回顾例 \ref{formal-defos-example-prorepresentable-deformation-functor}：
pro-可表函子是形变函子。因此，若 $\mathcal{F}$
等价于函子中的一个光滑 pro-可表群胚，% P10-E0498
则条件 (1)、(2)、(3) 分别由引理
\ref{formal-defos-lemma-deformation-groupoid-quotient} 的 (1)、(2)、(3) 得出。

\medskip\noindent
反过来，假设条件 (1)、(2)、(3) 成立。由条件 (1) 可知
(S1) 和 (S2) 成立；见引理 \ref{formal-defos-lemma-RS-implies-S1-S2}。
由引理 \ref{formal-defos-lemma-versal-object-existence}，存在万有形式对象
$\xi$。令 $U=\underline{R}|_{\mathcal{C}_\Lambda}$，% P10-E0499
则相应映射 $\underline{\xi}:U\to\mathcal{F}$ 光滑
（这正是万有形式对象的定义）。
设 $(U,R,s,t,c)$ 是引理
\ref{formal-defos-lemma-presentation-construction} 中由映射
$\underline{\xi}$ 构造的函子中的群胚。
由引理 \ref{formal-defos-lemma-smooth-groupoid-in-functors-construction} 可见，
$(U,R,s,t,c)$ 是函子中的光滑群胚，且
$[U/R]\to\mathcal{F}$ 是等价。由引理
\ref{formal-defos-lemma-prorepresentable-groupoid-in-functors-construction} 可见，
$(U,R,s,t,c)$ pro-可表。因此
$[U/R]\to\mathcal{F}$ 是所求的 $\mathcal{F}$ 的呈现。
\end{proof}

\section{关于极小性的说明}
\label{formal-defos-section-minimality}

\noindent
本章的主要定理是上面的定理
\ref{formal-defos-theorem-presentation-deformation-groupoid}。
它完全刻画了 $\mathcal{C}_\Lambda$ 上那些存在函子中
光滑 pro-可表群胚呈现的群胚余纤维范畴。
本节简要讨论如何利用第 \ref{formal-defos-section-minimal-versal} 节和
第 \ref{formal-defos-section-miniversal-objects-existence} 节所讨论的极小性，
得到一个“极小”光滑 pro-可表呈现。

\begin{definition}
\label{formal-defos-definition-minimal-groupoid-in-functors}
设 $(U,R,s,t,c)$ 是 $\mathcal{C}_\Lambda$
上函子中的光滑 pro-可表群胚。
\begin{enumerate}
\item 若群胚 $(U(k[\epsilon]),R(k[\epsilon]),s,t,c)$
全不连通，即任意两个不同对象之间都没有态射，
则称 $(U,R,s,t,c)$ {\it 规范化}。
\item 若 $U\to[U/R]$ 由 $[U/R]$
的一个极小万有形式对象给出，% P10-E0500
则称 $(U,R,s,t,c)$ {\it 极小}。
\end{enumerate}
\end{definition}

\noindent
这两个概念的差别与条件（\ref{formal-defos-equation-bijective}）和
（\ref{formal-defos-equation-bijective-orbits}）之间的差别有关，
并在 $k'\subset k$ 可分时消失。此外，函子中的规范化光滑
pro-可表群胚是极小的，正如下一个引理所示。准确陈述如下。

\begin{lemma}
\label{formal-defos-lemma-characterize-minimal-groupoid-in-functors}
设 $(U,R,s,t,c)$ 是 $\mathcal{C}_\Lambda$
上函子中的光滑 pro-可表群胚。
\begin{enumerate}
\item $(U,R,s,t,c)$ 规范化，当且仅当态射
$U\to[U/R]$ 在切空间上诱导同构；
\item $(U,R,s,t,c)$ 极小，当且仅当
$TU\to T[U/R]$ 的核包含在
$\text{Der}_\Lambda(k,k)\to TU$ 的像中。
\end{enumerate}
\end{lemma}

\begin{proof}
(1) 直接由定义得出。为证明 (2)，令
$\mathcal{F}=[U/R]$。由于 $\mathcal{F}$ 有一个呈现，
它是形变范畴；见定理
\ref{formal-defos-theorem-presentation-deformation-groupoid}。
特别地，它满足 (RS)、(S1)、(S2)；见引理
\ref{formal-defos-lemma-RS-implies-S1-S2}。
回顾引理 \ref{formal-defos-lemma-minimal-versal}：
极小万有形式对象在同构意义下唯一。
由定理 \ref{formal-defos-theorem-miniversal-object-existence}，
一个极小万有对象诱导满足
（\ref{formal-defos-equation-bijective-orbits}）的映射
$\underline{\xi}:\underline{R}|_{\mathcal{C}_\Lambda}
\to\mathcal{F}$。由于在 $\mathcal{F}$ 上有
$U\cong\underline{R}|_{\mathcal{C}_\Lambda}$，% P10-E0501
可见 $TU\to T\mathcal{F}=T[U/R]$
满足引理所述的性质。
\end{proof}

\noindent
$\mathcal C_\Lambda$ 上函子中的极小 pro-可表群胚之商，
不存在不是自同构的自等价。为证明这一点，先注意下列引理。

\begin{lemma}
\label{formal-defos-lemma-surjective-morphism-prorepresentable-functor}
设 $U:\mathcal{C}_\Lambda\to\textit{Sets}$ 是 pro-可表函子。
设 $\varphi:U\to U$ 是态射，并且
$d\varphi:TU\to TU$ 是同构。则 $\varphi$ 是同构。
\end{lemma}

\begin{proof}
若对某个 $R\in\Ob(\widehat{\mathcal{C}}_\Lambda)$，
$U\cong\underline{R}|_{\mathcal{C}_\Lambda}$，
则完备化 $\varphi$ 得到态射
$\underline{R}\to\underline{R}$。
若 $f:R\to R$ 是 $\widehat{\mathcal{C}}_\Lambda$
中相应的态射，则 $f$ 诱导同构
$\text{Der}_\Lambda(R,k)\to\text{Der}_\Lambda(R,k)$；
见例 \ref{formal-defos-example-tangent-space-map-prorepresentable-functor}。
特别地，由引理 \ref{formal-defos-lemma-derivations-surjective}，$f$ 是满射。
诺特环的满自同态是同构（见《代数》引理
\ref{algebra-lemma-surjective-endo-noetherian-ring-is-iso}），
所以 $f$ 是同构，从而 $\underline{R}\to\underline{R}$
是同构，进而 $\varphi:U\to U$ 是同构。
\end{proof}

\begin{lemma}
\label{formal-defos-lemma-minimal-prorepresentable-groupoid-autoequivalence}
设 $(U,R,s,t,c)$ 是 $\mathcal{C}_\Lambda$
上函子中的极小光滑 pro-可表群胚。
若 $\varphi:[U/R]\to[U/R]$ 是群胚余纤维范畴的等价，
则 $\varphi$ 是同构。
\end{lemma}

\begin{proof}
态射 $\varphi:[U/R]\to[U/R]$ 与 $\mathcal{C}_\Lambda$
上函子中群胚的态射
$\varphi:(U,R,s,t,c)\to(U,R,s,t,c)$ 是同一回事，% P10-E0502
这里后者如定义 \ref{formal-defos-definition-groupoid-in-functors} 所定义。
把相应态射记为 $\phi:U\to U$ 和 $\psi:R\to R$。
由于图表
$$
\xymatrix{
& \text{Der}_\Lambda(k, k) \ar[dr]_\gamma \ar[dl]^\gamma \\
TU \ar[rr]_{d\phi} \ar[d] & & TU \ar[d]  \\
T[U/R] \ar[rr]^{d\varphi} & & T[U/R]
}
$$
可交换，$d\varphi$ 双射，并且有引理
\ref{formal-defos-lemma-characterize-minimal-groupoid-in-functors}
中的极小性刻画，所以 $d\phi$ 是单射
（因维数理由，它也是双射）。因此由引理
\ref{formal-defos-lemma-surjective-morphism-prorepresentable-functor}，
$\phi:U\to U$ 是同构。
利用引理 \ref{formal-defos-lemma-deformation-functor-diagonal} 的正合序列
$$
0 \to \text{Inf}([U/R]) \to TR \to TU \oplus TU
$$
作类似论证，可证明 $\psi:R\to R$ 是同构。
但这也可由 $R=U\times_{[U/R]}U$ 以及
$\varphi$ 和 $\phi$ 都是同构直接得出。% P10-E0503
\end{proof}

\begin{lemma}
\label{formal-defos-lemma-minimal-prorepresentable-groupoid-equivalence}
设 $(U,R,s,t,c)$ 和 $(U',R',s',t',c')$ 是
$\mathcal{C}_\Lambda$ 上函子中的极小光滑 pro-可表群胚。
若 $\varphi:[U/R]\to[U'/R']$ 是群胚余纤维范畴的等价，
则 $\varphi$ 是同构。
\end{lemma}

\begin{proof}
设 $\psi:[U'/R']\to[U/R]$ 是 $\varphi$ 的拟逆。
由引理 \ref{formal-defos-lemma-minimal-prorepresentable-groupoid-autoequivalence}，
$\psi\circ\varphi$ 和 $\varphi\circ\psi$ 都是同构，
因此 $\varphi$ 和 $\psi$ 都是同构。
\end{proof}

\noindent
下一个引理概括了本章前面得到的一些结果。

\begin{lemma}
\label{formal-defos-lemma-minimal-groupoid-in-functors-construction}
设 $\mathcal{F}$ 是形变范畴，并且
$\dim_kT\mathcal{F}<\infty$ 且
$\dim_k\text{Inf}(\mathcal{F})<\infty$。
则 $\mathcal{F}$ 存在一个极小万有形式对象 $\xi$。
设 $\xi$ 位于 $R\in\Ob(\widehat{\mathcal{C}}_\Lambda)$ 上方。
令 $U=\underline{R}|_{\mathcal{C}_\Lambda}$。
令 $f=\underline{\xi}:U\to\mathcal{F}$ 为相应态射。
设 $(U,R,s,t,c)$ 是引理 \ref{formal-defos-lemma-presentation-construction}
中由 $f:U\to\mathcal{F}$ 构造的
$\mathcal{C}_\Lambda$ 上函子中的群胚。% P10-E0504
则 $(U,R,s,t,c)$ 是 $\mathcal{C}_\Lambda$
上函子中的极小光滑 pro-可表群胚，并且存在等价
$[U/R]\to\mathcal{F}$。
\end{lemma}

\begin{proof}
由于 $\mathcal{F}$ 是形变范畴，它满足 (S1) 和 (S2)；
见引理 \ref{formal-defos-lemma-RS-implies-S1-S2}。
由引理 \ref{formal-defos-lemma-versal-object-existence}，存在万有形式对象。
由引理 \ref{formal-defos-lemma-minimal-versal}，存在引理陈述中那样的
极小万有形式对象 $\xi/R$。令
$U=\underline{R}|_{\mathcal{C}_\Lambda}$，% P10-E0505
则相应映射 $\underline{\xi}:U\to\mathcal{F}$ 光滑
（这正是万有形式对象的定义）。
设 $(U,R,s,t,c)$ 是引理
\ref{formal-defos-lemma-presentation-construction} 中由映射
$\underline{\xi}$ 构造的函子中的群胚。由引理
\ref{formal-defos-lemma-smooth-groupoid-in-functors-construction} 可见，
$(U,R,s,t,c)$ 是函子中的光滑群胚，并且
$[U/R]\to\mathcal{F}$ 是等价。由引理
\ref{formal-defos-lemma-prorepresentable-groupoid-in-functors-construction} 可见，
$(U,R,s,t,c)$ pro-可表。
最后，$(U,R,s,t,c)$ 极小，因为
$U\to[U/R]=\mathcal{F}$ 对应于极小万有形式对象 $\xi$。% P10-E0506
\end{proof}

\noindent
由函子中的极小 pro-可表群胚给出的呈现满足下列唯一性。

\begin{lemma}
\label{formal-defos-lemma-minimal-presentations-equivalent}
设 $\mathcal{F}$ 是 $\mathcal{C}_\Lambda$
上的群胚余纤维范畴。% P10-E0507
假设 $\mathcal{F}$ 存在分别由函子中的极小光滑 pro-可表群胚
$(U,R,s,t,c)$ 和 $(U',R',s',t',c')$ 给出的呈现。
则 $(U,R,s,t,c)$ 与 $(U',R',s',t',c')$ 同构。
\end{lemma}

\begin{proof}
这由引理 \ref{formal-defos-lemma-minimal-prorepresentable-groupoid-equivalence}
以及下述观察得出：态射 $[U/R]\to[U'/R']$
与函子中群胚的态射是同一回事；% P10-E0502
这是因为定义 \ref{formal-defos-definition-quotient}
明确构造了 $[U/R]$。
\end{proof}

\noindent
综上，我们证明了下列定理。

\begin{theorem}
\label{formal-defos-theorem-minimal-smooth-prorepresentable-presentations}
设 $\mathcal{F}$ 是 $\mathcal{C}_\Lambda$ 上的群胚余纤维范畴。
考虑下列条件：
\begin{enumerate}
\item $\mathcal{F}$ 存在一个由 $\mathcal{C}_\Lambda$
上函子中的规范化光滑 pro-可表群胚给出的呈现；
\item $\mathcal{F}$ 存在一个由 $\mathcal{C}_\Lambda$
上函子中的光滑 pro-可表群胚给出的呈现；
\item $\mathcal{F}$ 存在一个由 $\mathcal{C}_\Lambda$
上函子中的极小光滑 pro-可表群胚给出的呈现；并且
\item $\mathcal{F}$ 满足以下条件：
\begin{enumerate}
\item $\mathcal{F}$ 是形变范畴。
\item $\dim_kT\mathcal{F}$ 有限。
\item $\dim_k\text{Inf}(\mathcal{F})$ 有限。
\end{enumerate}
\end{enumerate}
则 (2)、(3)、(4) 等价，并且都由 (1) 蕴含。
若 $k'\subset k$ 可分，则 (1)、(2)、(3)、(4) 全部等价。
此外，为 $\mathcal{F}$ 提供呈现的函子中极小光滑
pro-可表群胚在同构意义下唯一。
\end{theorem}

\begin{proof}
由引理 \ref{formal-defos-lemma-characterize-minimal-groupoid-in-functors}，
(1) 蕴含 (3)，并且在 $k'\subset k$ 可分时 (1) 等价于 (3)。
显然 (3) 蕴含 (2)。由定理
\ref{formal-defos-theorem-presentation-deformation-groupoid}，
(2) 蕴含 (4)。由引理
\ref{formal-defos-lemma-minimal-groupoid-in-functors-construction}，
(4) 蕴含 (3)。这就证明了所有蕴含关系。
最后的唯一性陈述由引理
\ref{formal-defos-lemma-minimal-presentations-equivalent} 得出。
\end{proof}

\section{万有环的唯一性}
\label{formal-defos-section-uniqueness}

\noindent
给定 $\widehat{\mathcal{C}}_\Lambda$ 中的 $R,S$，
若存在 $r\geq0$ 以及 $\widehat{\mathcal{C}}_\Lambda$ 中的映射
$h:R\to R[[t_1,\ldots,t_r]]$ 和
$k:R[[t_1,\ldots,t_r]]\to S$，使得对所有 $a\in R$，
\begin{enumerate}
\item $h(a)\bmod(t_1,\ldots,t_r)=a$；
\item $f(a)=k(a)$；
\item $g(a)=k(h(a))$，
\end{enumerate}
则称映射 $f,g:R\to S$ {\it 形式同伦}。
我们称 $(r,h,k)$ 是 $f$ 与 $g$ 之间的一个{\it 形式同伦}。

\begin{lemma}
\label{formal-defos-lemma-homotopy}
形式同伦在 $\widehat{\mathcal{C}}_\Lambda$
中的态射集合上构成等价关系。
\end{lemma}

\begin{proof}
假设任取 $r\geq1$，并有两个映射
$h_1,h_2:R\to R[[t_1,\ldots,t_r]]$，使得对所有 $a\in R$，
$h_1(a)\bmod(t_1,\ldots,t_r)
=h_2(a)\bmod(t_1,\ldots,t_r)=a$，
以及映射 $k:R[[t_1,\ldots,t_r]]\to S$。
我们断言 $k\circ h_1$ 与 $k\circ h_2$ 形式同伦。
该断言内在的对称性将说明形式同伦概念是对称的。% P10-E0508
事实上，映射
$$
\Psi :
R[[t_1, \ldots, t_r]]
\longrightarrow
R[[t_1, \ldots, t_r]],\quad
\sum a_I t^I \longmapsto \sum h_1(a_I)t^I
$$
是同构。对 $a\in R$，令
$h(a)=\Psi^{-1}(h_2(a))$，并令 $k'=k\circ\Psi$，
则可见 $(r,h,k')$ 是 $k\circ h_1$ 与
$k\circ h_2$ 之间的形式同伦，从而证明了该断言。% P10-E0509

\medskip\noindent
设有如上三个映射 $f_1,f_2,f_3:R\to S$，
以及 $f_1$ 与 $f_2$ 之间的形式同伦 $(r_1,h_1,k_1)$，
和 $f_3$ 与 $f_2$ 之间的形式同伦 $(r_2,h_2,k_2)$（！）。
重新标记坐标后，可假设
$h_2:R\to R[[t_{r_1+1},\ldots,t_{r_1+r_2}]]$，
且
$k_2:R[[t_{r_1+1},\ldots,t_{r_1+r_2}]]\to S$。
选择一个适当的同构
$$
R[[t_1, \ldots, t_{r_1 + r_2}]]
\longrightarrow
R[[t_{r_1 + 1}, \ldots, t_{r_1 + r_2}]]
\widehat{\otimes}_{h_2, R, h_1}
R[[t_1, \ldots, t_{r_1}]]
$$
后，可把这些映射置入可交换图表
$$
\xymatrix{
R \ar[r]_{h_1} \ar[d]^{h_2} &
R[[t_1, \ldots, t_{r_1}]] \ar[d]^{h_2'} \\
R[[t_{r_1 + 1}, \ldots, t_{r_1 + r_2}]] \ar[r]^{h_1'} &
R[[t_1, \ldots, t_{r_1 + r_2}]]
}
$$
中，其中当 $1\leq i\leq r_1$ 时 $h_2'(t_i)=t_i$，
当 $r_1+1\leq i\leq r_2$ 时 $h_1'(t_i)=t_i$。% P10-E0510
略去若干细节。
由于该图表是范畴 $\widehat{\mathcal{C}}_\Lambda$ 中的推出
（见引理 \ref{formal-defos-lemma-CLambdahat-pushouts} 的证明），
并且 $k_1\circ h_1=f_2=k_2\circ h_2$，
可知存在映射
$$
k : R[[t_1, \ldots, t_{r_1 + r_2}]] \to S
$$
使 $k_1=k\circ h_2'$ 且 $k_2=k\circ h_1'$。
记 $h=h_1'\circ h_2=h_2'\circ h_1$。于是
\begin{enumerate}
\item $k(h_1'(a))=k_2(a)=f_3(a)$；且
\item $k(h_2'(a))=k_1(a)=f_1(a)$。
\end{enumerate}
由证明第一段中的断言，这说明 $f_1$ 与 $f_3$ 形式同伦。
\end{proof}

\begin{lemma}
\label{formal-defos-lemma-composition-homotopic}
在范畴 $\widehat{\mathcal{C}}_\Lambda$ 中，若
$f_1,f_2:R\to S$ 形式同伦，且 $g:S\to S'$ 是态射，
则 $g\circ f_1$ 与 $g\circ f_2$ 形式同伦。
\end{lemma}

\begin{proof}
事实上，若 $(r,h,k)$ 是 $f_1$ 与 $f_2$ 之间的形式同伦，
则 $(r,h,g\circ k)$ 是 $g\circ f_1$ 与
$g\circ f_2$ 之间的形式同伦。
\end{proof}

\begin{lemma}
\label{formal-defos-lemma-versal-unique-up-to-homotopy}
设 $\mathcal{F}$ 是 $\mathcal{C}_\Lambda$ 上的形变范畴，
并且 $\dim_kT\mathcal{F}<\infty$、
$\dim_k\text{Inf}(\mathcal{F})<\infty$。
设 $\xi$ 是位于 $R$ 上方的万有形式对象。
设 $\eta$ 是位于 $S$ 上方的形式对象。
则任意两个满足
$$
f, g : R \to S
$$
且 $f_*\xi\cong\eta\cong g_*\xi$ 的映射都形式同伦。
\end{lemma}

\begin{proof}
根据定理 \ref{formal-defos-theorem-presentation-deformation-groupoid}
及其证明，$\mathcal{F}$ 有一个由函子中的光滑 pro-可表群胚
$$
(\underline{R}, \underline{R_1}, s, t, c, e, i)|_{\mathcal{C}_\Lambda}
$$
给出的呈现；该群胚位于 $\mathcal{C}_\lambda$ 上，
并且 $\mathcal{F}$。% P10-E0511 P10-E0512
于是映射 $s:R\to R_1$ 和 $t:R\to R_1$
是形式光滑环映射，且 $e:R_1\to R$ 是一个截面。
特别地，可对某个 $r\geq0$ 选择同构
$R_1=R[[t_1,\ldots,t_r]]$，
使 $s$ 是嵌入 $R\subset R[[t_1,\ldots,t_r]]$，
而 $t$ 对应于映射
$h:R\to R[[t_1,\ldots,t_r]]$，且对所有 $a\in R$，
$h(a)\bmod(t_1,\ldots,t_r)=a$。
同构 $\alpha:f_*\xi\to g_*\xi$ 的存在，恰好意味着
存在映射 $k:R_1\to S$，使
$f=k\circ s$ 且 $g=k\circ t$。
这恰好说明 $(r,h,k)$ 是 $f$ 与 $g$ 之间的形式同伦。
\end{proof}

\begin{lemma}
\label{formal-defos-lemma-homotopic-minimal-prime}
在范畴 $\widehat{\mathcal{C}}_\Lambda$ 中，若
$f_1,f_2:R\to S$ 形式同伦，且
$\mathfrak p\subset R$ 是极小素理想，则作为理想有
$f_1(\mathfrak p)S=f_2(\mathfrak p)S$。
\end{lemma}

\begin{proof}
设 $(r,h,k)$ 是 $f_1$ 与 $f_2$ 之间的形式同伦。
我们断言
$\mathfrak pR[[t_1,\ldots,t_r]]
=h(\mathfrak p)R[[t_1,\ldots,t_r]]$。
再与 $k$ 复合，该断言即蕴含本引理。
为证明断言，注意映射
$\mathfrak p\mapsto\mathfrak pR[[t_1,\ldots,t_r]]$
是 $R$ 的极小素理想与
$R[[t_1,\ldots,t_r]]$ 的极小素理想之间的双射。
最后，因为 $h$ 平坦，
$h(\mathfrak p)R[[t_1,\ldots,t_r]]$ 是极小素理想，
所以根据刚才所述，它形如
$\mathfrak qR[[t_1,\ldots,t_r]]$，
其中 $\mathfrak q\subset R$ 是某个极小素理想。
但根据形式同伦的定义，
$h\bmod(t_1,\ldots,t_r)=\text{id}_R$，
故 $\mathfrak q=\mathfrak p$，这正是所求。
\end{proof}

\section{更换剩余域}
\label{formal-defos-section-change-of-field}

\noindent
本节简要讨论用有限扩张替换剩余域 $k$ 后会发生什么。
设 $\Lambda$ 是诺特环，并设 $\Lambda\to k$ 是有限环映射，
其中 $k$ 是域。在本章中，我们一直用
$\mathcal{C}_\Lambda$ 表示剩余域与 $k$ 等同的
阿廷局部 $\Lambda$-代数范畴；见定义
\ref{formal-defos-definition-CLambda}。不过，本节将讨论更换 $k$ 后
会发生什么，% P10-E0513
所以改用记号 $\mathcal{C}_{\Lambda,k}$ 来指明对 $k$ 的依赖。

\begin{situation}
\label{formal-defos-situation-change-of-fields}
设 $\Lambda$ 是诺特环，并设
$\Lambda\to k\to l$ 是有限环映射，其中 $k$、$l$ 都是域。% P10-E0514
因此 $l/k$ 是有限域扩张。% P10-E0515
$\mathcal{C}_{\Lambda,l}$ 的典型对象记为 $B$，
$\mathcal{C}_{\Lambda,k}$ 的典型对象记为 $A$。定义
\begin{equation}
\label{formal-defos-equation-comparison}
\mathcal{C}_{\Lambda, l} \longrightarrow \mathcal{C}_{\Lambda, k},
\quad
B \longmapsto B \times_l k
\end{equation}
给定群胚余纤维范畴
$p:\mathcal{F}\to\mathcal{C}_{\Lambda,k}$，
令 $\mathcal{F}_{l/k}(B)=\mathcal{F}(B\times_l k)$，
便得到相应的群胚余纤维范畴
$$
p_{l/k} : \mathcal{F}_{l/k} \longrightarrow \mathcal{C}_{\Lambda, l}.
$$
\end{situation}

\noindent
函子（\ref{formal-defos-equation-comparison}）确有意义：
由于 $B\times_l k\subset B$，有
\begin{align*}
[k : k']\ \text{length}_{B \times_l k}(B \times_l k) & =
\text{length}_\Lambda(B \times_l k) \\
& \leq \text{length}_\Lambda(B) \\
& = [l : k']\ \text{length}_B(B) < \infty
\end{align*}
（见引理 \ref{formal-defos-lemma-length}），所以 $B\times_l k$ 是阿廷环
（见《代数》引理
\ref{algebra-lemma-artinian-finite-length}）。
因此 $B\times_l k$ 是剩余域为 $k$ 的阿廷局部环。
注意，（\ref{formal-defos-equation-comparison}）与纤维积可交换：
$$
(B_1 \times_B B_2) \times_l k =
(B_1 \times_l k) \times_{(B \times_l k)} (B_2 \times_l k),
$$
并且把满环映射变为满环映射。
我们使用较为“昂贵”的记号 $\mathcal{F}_{l/k}$，
以免与说明 \ref{formal-defos-remark-localize-cofibered-groupoid}
中的构造混淆。下面给出一些初等观察。

\begin{lemma}
\label{formal-defos-lemma-elementary-properties-change-of-field}
采用情形 \ref{formal-defos-situation-change-of-fields} 中的记号和假设。
\begin{enumerate}
\item 有
$\overline{\mathcal{F}_{l/k}}=(\overline{\mathcal{F}})_{l/k}$。
\item 若 $\mathcal{F}$ 是预形变范畴，则
$\mathcal{F}_{l/k}$ 是预形变范畴。
\item 若 $\mathcal{F}$ 满足 (S1)，则
$\mathcal{F}_{l/k}$ 满足 (S1)。
\item 若 $\mathcal{F}$ 满足 (S2)，则
$\mathcal{F}_{l/k}$ 满足 (S2)。
\item 若 $\mathcal{F}$ 满足 (RS)，则
$\mathcal{F}_{l/k}$ 满足 (RS)。
\end{enumerate}
\end{lemma}

\begin{proof}
(1) 直接由定义得出。

\medskip\noindent
由于 $\mathcal{F}_{l/k}(l)=\mathcal{F}(k)$，
(2) 由定义得出；见定义 \ref{formal-defos-definition-predeformation-category}。

\medskip\noindent
因为函子（\ref{formal-defos-equation-comparison}）与纤维积可交换，
且把满射变为满射，(3) 得出；见定义
\ref{formal-defos-definition-S1-S2}。

\medskip\noindent
证明 (4)。考虑 $\mathcal{C}_{\Lambda,l}$ 中如定义
\ref{formal-defos-definition-S1-S2} 所述的图表
$$
\xymatrix{
          & l[\epsilon] \ar[d] \\
B  \ar[r] & l
}.
$$
应用函子（\ref{formal-defos-equation-comparison}），得到
$$
\xymatrix{
          & k[l\epsilon] \ar[d] \\
B \times_l k  \ar[r] & k
},
$$
其中 $l\epsilon$ 表示有限维 $k$-向量空间
$l\epsilon\subset l[\epsilon]$。根据引理
\ref{formal-defos-lemma-S2-extensions}，$\mathcal{F}$ 的 (S2)
条件也对该图表成立。因此 (S2) 对
$\mathcal{F}_{l/k}$ 成立。

\medskip\noindent
(5) 由引理 \ref{formal-defos-lemma-RS-2-categorical} 的 (2)
对 (RS) 的刻画，以及
（\ref{formal-defos-equation-comparison}）与纤维积可交换这一事实得出。
\end{proof}

\noindent
下列引理尤其适用于 $\mathcal{F}$ 满足 (S2)
且为预形变范畴的情形；见引理
\ref{formal-defos-lemma-S1-S2-associated-functor}。

\begin{lemma}
\label{formal-defos-lemma-tangent-space-change-of-field}
采用情形 \ref{formal-defos-situation-change-of-fields} 中的记号和假设。
假设 $\mathcal{F}$ 是预形变范畴，且
$\overline{\mathcal{F}}$ 满足 (S2)。则切空间之间存在典范
$l$-向量空间同构
$$
T\mathcal{F} \otimes_k l \longrightarrow T\mathcal{F}_{l/k}.
$$
\end{lemma}

\begin{proof}
由引理 \ref{formal-defos-lemma-elementary-properties-change-of-field}，
可用 $\overline{\mathcal{F}}$ 替换 $\mathcal{F}$。
此外，由引理 \ref{formal-defos-lemma-tangent-space-vector-space}，
$T\mathcal{F}$、$T\mathcal{F}_{l/k}$ 分别具有典范
$k$-向量空间结构、$l$-向量空间结构。于是
$$
T\mathcal{F}_{l/k} = \mathcal{F}_{l/k}(l[\epsilon])
= \mathcal{F}(k[l\epsilon]) = T\mathcal{F} \otimes_k l,
$$
最后一个等号由引理
\ref{formal-defos-lemma-tangent-space-vector-space} 得出。
更一般地，给定有限维 $l$-向量空间 $V$，有
$$
\mathcal{F}_{l/k}(l[V]) = \mathcal{F}(k[V_k]) = T\mathcal{F} \otimes_k V_k,
$$
其中 $V_k$ 表示把 $V$ 看作 $k$-向量空间。
因此，作为取值于集合范畴的函子，函子
$V\mapsto\mathcal{F}_{l/k}(l[V])$ 和
$V\mapsto T\mathcal{F}\otimes_k V_k$ 典范等同。
由引理 \ref{formal-defos-lemma-linear-functor}，可见至多有一种方法
使任一函子成为 $l$-线性函子。因此这些同构与
$l$-向量空间结构相容，结论得证。
\end{proof}

\begin{lemma}
\label{formal-defos-lemma-inf-aut-change-of-field}
采用情形 \ref{formal-defos-situation-change-of-fields} 中的记号和假设。
假设 $\mathcal{F}$ 是形变范畴。
则无穷小自同构空间之间存在典范 $l$-向量空间同构
$$
\text{Inf}(\mathcal{F}) \otimes_k l
\longrightarrow
\text{Inf}(\mathcal{F}_{l/k}).
$$
\end{lemma}

\begin{proof}
设 $x_0\in\Ob(\mathcal{F}(k))$，并用 $x_{l,0}$
表示 $\mathcal{F}_{l/k}$ 在 $l$ 上方的相应对象。% P10-E0516
回顾
$\text{Inf}(\mathcal{F})=\text{Inf}_{x_0}(\mathcal{F})$ 且
$\text{Inf}(\mathcal{F}_{l/k})
=\text{Inf}_{x_{l,0}}(\mathcal{F}_{l/k})$；
见说明 \ref{formal-defos-remark-trivial-aut-point}。
回顾 $\text{Inf}_{x_0}(\mathcal{F})$ 的向量空间结构来自
将其等同为函子 $\mathit{Aut}(x_0)$ 的切空间；
该函子定义在剩余域为 $k$ 的阿廷局部 $k$-代数范畴
$\mathcal{C}_{k,k}$ 上。类似地，
$\text{Inf}_{x_{l,0}}(\mathcal{F}_{l/k})$
是 $\mathit{Aut}(x_{l,0})$ 的切空间；
后者定义在剩余域为 $l$ 的阿廷局部 $l$-代数范畴
$\mathcal{C}_{l,l}$ 上。展开定义可见，
$\mathit{Aut}(x_{l,0})$ 是
$\mathit{Aut}(x_0)_{l/k}$（它定义在
$\mathcal{C}_{k,l}$ 上）到 $\mathcal{C}_{l,l}$ 的限制。
把 $\mathit{Aut}(x_0)_{l/k}$ 看作
$\mathcal{C}_{k,l}$ 或 $\mathcal{C}_{l,l}$ 上的函子，
其切空间没有差别。因此本引理由引理
\ref{formal-defos-lemma-tangent-space-change-of-field} 以及下述事实得出：
由引理 \ref{formal-defos-lemma-Aut-functor-RS}，
$\mathit{Aut}(x_0)$ 满足 (RS)
（因而由引理 \ref{formal-defos-lemma-RS-implies-S1-S2} 满足 (S2)）。
\end{proof}

\begin{lemma}
\label{formal-defos-lemma-change-of-fields-smooth}
采用情形 \ref{formal-defos-situation-change-of-fields} 中的记号和假设。
若 $\mathcal{F}\to\mathcal{G}$ 是
$\mathcal{C}_{\Lambda,k}$ 上群胚余纤维范畴的光滑态射，
则 $\mathcal{F}_{l/k}\to\mathcal{G}_{l/k}$ 是
$\mathcal{C}_{\Lambda,l}$ 上群胚余纤维范畴的光滑态射。
\end{lemma}

\begin{proof}
这直接由定义以及
（\ref{formal-defos-equation-comparison}）保持满射这一事实得出。
\end{proof}

\noindent
关于 $\mathcal{F}$ 与 $\mathcal{F}_{l/k}$ 的关系，
还有更多可以说明的内容（特别是关于万有形变之间的关系）；
我们会在需要时把这些内容添加到这里。

\begin{lemma}
\label{formal-defos-lemma-change-of-field-versal-ring}
采用情形 \ref{formal-defos-situation-change-of-fields} 中的记号和假设。
设 $\xi$ 是 $\mathcal{F}$ 的万有形式对象，位于
$R\in\Ob(\widehat{\mathcal{C}}_{\Lambda,k})$ 上方。则存在：
\begin{enumerate}
\item 一个 $S\in\Ob(\widehat{\mathcal{C}}_{\Lambda,l})$
以及局部 $\Lambda$-代数同态 $R\to S$；
该同态在 $\mathfrak m_S$-进拓扑下形式光滑，
并在剩余域上诱导给定的域扩张 $l/k$；以及
\item $\mathcal{F}_{l/k}$ 的一个位于 $S$ 上方的万有形式对象。
\end{enumerate}
\end{lemma}

\begin{proof}
构造 $S$。选择 $R$-代数的满射
$R[x_1,\ldots,x_n]\to l$。其核为极大理想 $\mathfrak m$。
令 $S$ 为诺特环 $R[x_1,\ldots,x_n]$ 的
$\mathfrak m$-进完备化。于是由《代数》引理
\ref{algebra-lemma-completion-Noetherian-Noetherian}，
$S$ 属于 $\widehat{\mathcal{C}}_{\Lambda,l}$。
由《代数进阶》引理
\ref{more-algebra-lemma-formally-smooth} 和
\ref{more-algebra-lemma-formally-smooth-completion}，
以及 $R\to R[x_1,\ldots,x_n]$ 形式光滑这一事实，
映射 $R\to S$ 在 $\mathfrak m_S$-进拓扑下形式光滑。
（比较引理 \ref{formal-defos-lemma-exists-smooth} 的证明。）% P10-E0517

\medskip\noindent
由于 $\xi$ 万有，变换
$\underline{\xi}:\underline{R}|_{\mathcal{C}_{\Lambda,k}}
\to\mathcal{F}$ 光滑。由引理
\ref{formal-defos-lemma-change-of-fields-smooth}，诱导映射
$$
(\underline{R}|_{\mathcal{C}_{\Lambda, k}})_{l/k}
\longrightarrow
\mathcal{F}_{l/k}
$$
光滑。因此只需构造光滑态射
$\underline{S}|_{\mathcal{C}_{\Lambda,l}}\to
(\underline{R}|_{\mathcal{C}_{\Lambda,k}})_{l/k}$。
给出这样的映射，意即对
$\mathcal{C}_{\Lambda,l}$ 的每个对象 $B$，
给出关于 $B$ 函子性的集合映射
$$
\Mor_{\widehat{\mathcal{C}}_{\Lambda, l}}(S, B)
\longrightarrow
\Mor_{\widehat{\mathcal{C}}_{\Lambda, k}}(R, B \times_l k).
$$
给定左侧元素 $\varphi:S\to B$，
把它映到复合 $R\to S\to B$；
该复合的像包含在子 $\Lambda$-代数 $B\times_l k$ 中。
该映射的光滑性意味着：给定满射 $B'\to B$ 和可交换图表
$$
\xymatrix{
S \ar[r] & B \ar@{=}[r] & B \\
R \ar[u] \ar[r] & B' \times_l k \ar[u] \ar[r] & B' \ar[u]
},
$$
必须找到环映射 $S\to B'$，使其嵌入外侧矩形。
由于我们选择的 $R\to S$ 在 $\mathfrak m_S$-进拓扑下
形式光滑，这一映射的存在性由《代数进阶》引理
\ref{more-algebra-lemma-lift-continuous} 保证。
\end{proof}
% Generated mechanically from the frozen defos.tex target.
% Do not edit this generated file; edit the bound locale source instead.

% OK, start here.
%

\setcounter{chapter}{90}
\chapter{形变理论}



\phantomsection
\label{defos-section-phantom}


\section{引言}
\label{defos-section-introduction}

\noindent
本章旨在较为平缓地介绍模、态射等对象的形变理论。本章讨论能够利用
朴素余切复形证明的那些结果。在关于余切复形的一章中，我们还会略微推广
这些结果。希望深入了解的读者可以参阅 Illusie 关于这一主题的专著，见
\cite{cotangent}。





\section{环的形变与朴素余切复形}
\label{defos-section-deformations}

\noindent
本节利用朴素余切复形来处理一些形变理论问题。先取一个满环同态
$A' \to A$，其核为平方零理想 $I$。此外，设已给定环同态
$A \to B$、$B$-模 $N$ 以及 $A$-模同态 $c : I \to N$。
我们要问，能否找到一个对象来替代下图中的问号，使图表交换：
\begin{equation}
\label{defos-equation-to-solve}
\vcenter{
\xymatrix{
0 \ar[r] & N \ar[r] & {?} \ar[r] & B \ar[r] & 0 \\
0 \ar[r] & I \ar[u]^c \ar[r] & A' \ar[u] \ar[r] & A \ar[u] \ar[r] & 0
}
}
\end{equation}
并且，如果解存在，它在何种意义下是唯一的。更确切地说，我们寻找一个
满的 $A'$-代数同态 $B' \to B$，其核是与 $N$ 等同的平方零理想，
并且 $A' \to B'$ 诱导出给定同态 $c$。我们称 $B'$ 为
(\ref{defos-equation-to-solve}) 的一个{\it 解}。

\begin{lemma}
\label{defos-lemma-huge-diagram}
给定交换图
$$
\xymatrix{
& 0 \ar[r] & N_2 \ar[r] & B'_2 \ar[r] & B_2 \ar[r] & 0 \\
& 0 \ar[r]|\hole & I_2 \ar[u]_{c_2} \ar[r] &
A'_2 \ar[u] \ar[r]|\hole & A_2 \ar[u] \ar[r] & 0 \\
0 \ar[r] & N_1 \ar[ruu] \ar[r] & B'_1 \ar[r] & B_1 \ar[ruu] \ar[r] & 0 \\
0 \ar[r] & I_1 \ar[ruu]|\hole \ar[u]^{c_1} \ar[r] &
A'_1 \ar[ruu]|\hole \ar[u] \ar[r] & A_1 \ar[ruu]|\hole \ar[u] \ar[r] & 0
}
$$
其中前后两面均为 (\ref{defos-equation-to-solve}) 的解，则有：
\begin{enumerate}
\item 在
$\Ext^1_{B_1}(\NL_{B_1/A_1}, N_2)$
中存在一个典范元；该元为零，当且仅当存在环同态
$B'_1 \to B'_2$ 使上述图表交换。
\item 如果存在使图表交换的同态 $B'_1 \to B'_2$，那么所有这类
同态组成的集合是 $\Hom_{B_1}(\Omega_{B_1/A_1}, N_2)$ 作用下的
主齐性空间。
\end{enumerate}
\end{lemma}

\begin{proof}
令 $E$ 为 $B_1$ 的底层集合。考虑满同态 $A_1[E] \to B_1$，记其核为
$J$；朴素余切复形正是借助该同态按下式定义的：
$$
\NL_{B_1/A_1} = (J/J^2 \to \Omega_{A_1[E]/A_1} \otimes_{A_1[E]} B_1)
$$
参见《交换代数》第 \ref{algebra-section-netherlander} 节。由于
$\Omega_{A_1[E]/A_1} \otimes B_1$ 是自由 $B_1$-模，故
$$
\Ext^1_{B_1}(\NL_{B_1/A_1}, N_2) =
\frac{\Hom_{B_1}(J/J^2, N_2)}
{\Hom_{B_1}(\Omega_{A_1[E]/A_1} \otimes B_1, N_2)}
$$
我们将在右端的模中构造一个障碍。令
$J' = \Ker(A'_1[E] \to B_1)$。注意，存在满同态 $J' \to J$，
其核为 $I_1A'_1[E]$。对每个 $e \in E$，以 $x_e \in A_1[E]$
表示相应的变量。选取 $x_e$ 在 $B_1$ 中的像于 $B'_1$ 中的一个提升
$y_e$，并选取 $x_e$ 在 $B_2$ 中的像于 $B'_2$ 中的一个提升 $z_e$。
这些选择确定了 $A'_1$-代数同态
$$
A'_1[E] \to B'_1 \quad\text{且}\quad A'_1[E] \to B'_2
$$
其中第一个给出同态 $J' \to N_1,\ f' \mapsto f'(y_e)$，第二个给出
同态 $J' \to N_2,\ f' \mapsto f'(z_e)$。计算表明，这两个同态均在
$(J')^2$ 上为零。由于图表的左方块（涉及 $c_1$ 与 $c_2$）交换，
把它们视为取值于 $N_2$ 的同态时，它们在 $I_1A'_1[E]$ 上一致。
注意，$B'_1$ 是 $J' \to A'_1[E]$ 与 $J' \to N_1$ 的推出。因此，
若同态 $J' \to N_1 \to N_2$ 与 $J' \to N_2$ 一致，就得到一个
使图表交换的同态 $B'_1 \to B'_2$。于是，把下列同态的类定义为障碍：
$$
J/J^2 \to N_2,\quad f \mapsto f'(z_e) - \nu(f'(y_e))
$$
这里 $\nu : N_1 \to N_2$ 是给定同态，而 $f' \in J'$ 是 $f$ 的一个
提升。由上述讨论，该定义是良定的。注意，对某些
$\delta_{i,e} \in N_i$，我们可以把所选的 $z_e$ 改为
$z_e + \delta_{2,e}$，把 $y_e$ 改为 $y_e + \delta_{1,e}$。
这样一来，上述同态变为
% P11-E0001: frozen source formula retained; see control/ERRATA_CANDIDATES.jsonl.
$$
f \mapsto f'(z_e + \delta_{2, e}) - \nu(f'(y_e + \delta_{1, e})) =
f'(z_e) - \nu(f'(z_e)) +
\sum (\delta_{2, e} - \nu(\delta_{1, e}))\frac{\partial f}{\partial x_e}
$$
这恰好意味着，我们用复合同态
$J/J^2 \to \Omega_{A_1[E]/A_1} \otimes B_1 \to N_2$
来修正同态 $J/J^2 \to N_2$，其中第二个同态把 $\text{d}x_e$ 映为
$\delta_{2,e} - \nu(\delta_{1,e})$。因此，我们的障碍是良定义的，
并且它为零当且仅当存在提升。

\medskip\noindent
第 (2) 点来自如下观察：给定两个使图表交换的同态
$\varphi, \psi : B'_1 \to B'_2$，差 $\varphi - \psi$ 分解为一个同态
$D : B_1 \to N_2$，而它是一个 $A_1$-导子：
\begin{align*}
D(fg) & = \varphi(f'g') - \psi(f'g') \\
& =
\varphi(f')\varphi(g') - \psi(f')\psi(g') \\
& =
(\varphi(f') - \psi(f'))\varphi(g') + \psi(f')(\varphi(g') - \psi(g')) \\
& =
gD(f) + fD(g)
\end{align*}
因此，$D$ 对应于唯一的 $B_1$-线性同态
$\Omega_{B_1/A_1} \to N_2$。反过来，给定这样的线性同态便得到一个
导子 $D$；再给定一个使图表交换的环同态
$\psi : B'_1 \to B'_2$，则 $\psi + D$ 是另一个使图表交换的环同态。
\end{proof}

\begin{lemma}
\label{defos-lemma-choices}
如果 (\ref{defos-equation-to-solve}) 存在解，那么解的同构类集合是
$\Ext^1_B(\NL_{B/A}, N)$ 作用下的主齐性空间。
\end{lemma}

\begin{proof}
先注意，给定 (\ref{defos-equation-to-solve}) 的两个解 $B'_1$ 与 $B'_2$，
由引理 \ref{defos-lemma-huge-diagram} 可得一个障碍元
$o(B'_1, B'_2) \in \Ext^1_B(\NL_{B/A}, N)$；它阻碍同态
$B'_1 \to B'_2$ 的存在。显然，该元也是同构存在性的障碍，因而能够区分
不同的同构类。因此，为完成证明，只需说明：给定一个解 $B'$ 和一个元素
$\xi \in \Ext^1_B(\NL_{B/A}, N)$，可以找到另一个解 $B'_\xi$，使得
$o(B', B'_\xi) = \xi$。

\medskip\noindent
令 $E$ 为 $B$ 的底层集合。考虑满同态 $A[E] \to B$，记其核为 $J$；
朴素余切复形正是借助该同态按下式定义的：
$$
\NL_{B/A} = (J/J^2 \to \Omega_{A[E]/A} \otimes_{A[E]} B)
$$
参见《交换代数》第 \ref{algebra-section-netherlander} 节。由于
$\Omega_{A[E]/A} \otimes B$ 是自由 $B$-模，故
$$
\Ext^1_B(\NL_{B/A}, N) =
\frac{\Hom_B(J/J^2, N)}
{\Hom_B(\Omega_{A[E]/A} \otimes B, N)}
$$
因此，可以用一个同态 $\delta : J/J^2 \to N$ 的类来表示 $\xi$。

\medskip\noindent
对每个 $e \in E$，以 $x_e \in A[E]$ 表示相应的变量。选取 $x_e$ 在
$B$ 中的像于 $B'$ 中的一个提升 $y_e$。这些选择确定了一个
$A'$-代数同态 $\varphi : A'[E] \to B'$。令
$J' = \Ker(A'[E] \to B)$。注意，$\varphi$ 诱导同态
$\varphi|_{J'} : J' \to N$，且 $B'$ 是下图所示的推出：
$$
\xymatrix{
0 \ar[r] & N \ar[r] & B' \ar[r] & B \ar[r] & 0 \\
0 \ar[r] & J' \ar[u]^{\varphi|_{J'}} \ar[r] & A'[E] \ar[u] \ar[r] &
B \ar[u]_{=} \ar[r] & 0
}
$$
令 $\psi : J' \to N$ 为同态 $\varphi|_{J'}$ 与下列复合同态之和：
$$
J' \to J'/(J')^2 \to J/J^2 \xrightarrow{\delta} N.
$$
于是，沿 $\psi$ 所作的推出给出另一个环扩张 $B'_\xi$，它可置于上述
形式的图表中。计算表明 $o(B', B'_\xi) = \xi$，正如所需。
\end{proof}

\begin{lemma}
\label{defos-lemma-extensions-of-algebras}
设 $A$ 为环，$B$ 为 $A$-代数，$N$ 为 $B$-模。所有如下
$A$-代数扩张的同构类所成的集合
$$
0 \to N \to B' \to B \to 0
$$
（其中 $N$ 是平方零理想）与 $\Ext^1_B(\NL_{B/A}, N)$ 典范双射。
\end{lemma}

\begin{proof}
为证明这一点，把前述结果应用于 (\ref{defos-equation-to-solve}) 由下图给出的
情形：
$$
\xymatrix{
0 \ar[r] & N \ar[r] & {?} \ar[r] & B \ar[r] & 0 \\
0 \ar[r] & 0 \ar[u] \ar[r] & A \ar[u] \ar[r]^{\text{id}} & A \ar[u] \ar[r] & 0
}
$$
于是，本引理由引理 \ref{defos-lemma-choices} 以及解的存在性推出；一个解就是
$N \oplus B$。（该双射的直接构造见下述评注。）
\end{proof}

\begin{remark}
\label{defos-remark-extensions-of-algebras}
设 $A \to B$ 与 $N$ 如引理 \ref{defos-lemma-extensions-of-algebras} 所述。
令 $\alpha : P \to B$ 为 $B$ 在 $A$ 上的一个呈示，见《交换代数》
第 \ref{algebra-section-netherlander} 节。记 $J = \Ker(\alpha)$，
则与 $\alpha$ 相伴的朴素余切复形 $\NL(\alpha)$ 为复形
$J/J^2 \to \Omega_{P/A} \otimes_P B$。我们有
$$
\Ext^1_B(\NL(\alpha), N) =
\Coker\left(\Hom_B(\Omega_{P/A} \otimes_P B, N) \to \Hom_B(J/J^2, N)\right)
$$
因为 $\Omega_{P/A}$ 是自由模。考虑引理中所述的一个扩张
$0 \to N \to B' \to B \to 0$。由于 $P$ 是 $A$ 上的多项式代数，
可以把 $\alpha$ 提升为 $A$-代数同态 $\alpha' : P' \to B'$。
由于 $N$ 在 $B'$ 中平方为零，同态 $\alpha'|_J : J \to N$ 分解为
$J \to J/J^2 \to N$。该引理把这个扩张映到上述余核中此同态
$J/J^2 \to N$ 的类。
\end{remark}

\begin{lemma}
\label{defos-lemma-extensions-of-algebras-functorial}
给定环同态 $A \to B \to C$、$B$-模 $M$、$C$-模 $N$、$B$-线性同态
$c : M \to N$，以及如下具有平方零核的 $A$-代数扩张：
\begin{enumerate}
\item[(a)] 与 $\xi \in \Ext^1_B(\NL_{B/A}, M)$ 对应的
$0 \to M \to B' \to B \to 0$；
\item[(b)] 与 $\zeta \in \Ext^1_C(\NL_{C/A}, N)$ 对应的
$0 \to N \to C' \to C \to 0$。
\end{enumerate}
参见引理 \ref{defos-lemma-extensions-of-algebras}。那么，存在与 $B \to C$
及 $c$ 相容的 $A$-代数同态 $B' \to C'$，当且仅当 $\xi$ 与 $\zeta$
在 $\Ext^1_B(\NL_{B/A}, N)$ 中的像相同。
\end{lemma}

\begin{proof}
这个陈述是有意义的，因为我们有同态
$$
\Ext^1_B(\NL_{B/A}, M) \to \Ext^1_B(\NL_{B/A}, N)
$$
它由同态 $M \to N$ 诱导；还具有复合同态
$$
\Ext^1_C(\NL_{C/A}, N) \to \Ext^1_B(\NL_{C/A}, N) \to \Ext^1_B(\NL_{B/A}, N)
$$
其中第一个箭头使用限制同态 $D(C) \to D(B)$，第二个箭头使用复形的
典范同态 $\NL_{B/A} \to \NL_{C/A}$。把引理
\ref{defos-lemma-huge-diagram} 应用于下图，即可推出本引理的陈述：
$$
\xymatrix{
& 0 \ar[r] & N \ar[r] & C' \ar[r] & C \ar[r] & 0 \\
& 0 \ar[r]|\hole & 0 \ar[u] \ar[r] &
A \ar[u] \ar[r]|\hole & A \ar[u] \ar[r] & 0 \\
0 \ar[r] & M \ar[ruu] \ar[r] & B' \ar[r] & B \ar[ruu] \ar[r] & 0 \\
0 \ar[r] & 0 \ar[ruu]|\hole \ar[u] \ar[r] &
A \ar[ruu]|\hole \ar[u] \ar[r] & A \ar[ruu]|\hole \ar[u] \ar[r] & 0
}
$$
同时还要使用引理 \ref{defos-lemma-extensions-of-algebras} 与
\ref{defos-lemma-huge-diagram} 的证明中各项构造之间的一种相容性；其陈述和
证明从略。（直接论证见下述评注。）
\end{proof}

\begin{remark}
\label{defos-remark-extensions-of-algebras-functorial}
设 $A \to B \to C$、$M$、$N$、$c : M \to N$、
$0 \to M \to B' \to B \to 0$、$\xi \in \Ext^1_B(\NL_{B/A}, M)$、
$0 \to N \to C' \to C \to 0$ 以及
$\zeta \in \Ext^1_C(\NL_{C/A}, N)$ 如引理
\ref{defos-lemma-extensions-of-algebras-functorial} 所述。沿
$c : M \to N$ 作推出，可构造扩张
$$
\xymatrix{
0 \ar[r] & N \ar[r] & B'_1 \ar[r] & B \ar[r] & 0 \\
0 \ar[r] & M \ar[u]^c \ar[r] & B' \ar[u] \ar[r] & B \ar[u] \ar[r] & 0
}
$$
其中令 $B'_1 = (N \times B')/M$，而 $M$ 反对角嵌入。沿 $B \to C$
作拉回，可构造扩张
$$
\xymatrix{
0 \ar[r] & N \ar[r] & C' \ar[r] & C \ar[r] & 0 \\
0 \ar[r] & N \ar[u] \ar[r] & B'_2 \ar[u] \ar[r] & B \ar[u] \ar[r] & 0
}
$$
其中令 $B'_2 = C' \times_C B$（环的纤维积）。简单的追图表明，存在
与 $B \to C$ 及 $c$ 相容的 $A$-代数同态 $B' \to C'$，当且仅当把
$B'_1$ 与 $B'_2$ 都看作 $B$ 被 $N$ 扩张所得的 $A$-代数时二者同构。
因此，为证明引理 \ref{defos-lemma-extensions-of-algebras-functorial}，
只需说明：经引理 \ref{defos-lemma-extensions-of-algebras} 的双射，$B'_1$
对应于 $\xi$ 在同态
$\Ext^1_B(\NL_{B/A}, M) \to \Ext^1_B(\NL_{B/A}, N)$
下的像，并且 $B'_2$ 对应于 $\zeta$ 在同态
$\Ext^1_C(\NL_{C/A}, N) \to \Ext^1_B(\NL_{B/A}, N)$.
第一个陈述由评注 \ref{defos-remark-extensions-of-algebras} 中类的构造立即得到。
为证明第二个陈述，选取交换图
$$
\xymatrix{
Q \ar[r]_\beta & C \\
P \ar[u]^\varphi \ar[r]^\alpha & B \ar[u]
}
$$
其中各项都是 $A$-代数，$\alpha$ 是 $B$ 在 $A$ 上的一个呈示，而
$\beta$ 是 $C$ 在 $A$ 上的一个呈示。参见评注
\ref{defos-remark-extensions-of-algebras} 及其中的文献。令
$J = \Ker(\alpha)$、$K = \Ker(\beta)$。同态 $\varphi$ 诱导复形同态
$\NL(\alpha) \to \NL(\beta)$，特别地诱导
$\bar\varphi : J/J^2 \to K/K^2$。选取提升 $\beta$ 的 $A$-代数同态
$\beta' : Q \to C'$。于是
$\alpha' = (\beta' \circ \varphi, \alpha) : P \to B'_2 = C' \times_C B$
是 $\alpha$ 的一个提升。在这些选择下，由 $\beta'$ 诱导的同态
$K/K^2 \to N$ 与 $\bar\varphi : J/J^2 \to K/K^2$ 的复合，正是
$\alpha'$ 在 $J/J^2$ 上的限制。展开评注
\ref{defos-remark-extensions-of-algebras} 中各个类的构造，这确实表明
$B'_2$ 对应于 $\zeta$ 在同态
$\Ext^1_C(\NL_{C/A}, N) \to \Ext^1_B(\NL_{B/A}, N)$.
下的像。
\end{remark}

\begin{lemma}
\label{defos-lemma-parametrize-solutions}
设 $0 \to I \to A' \to A \to 0$、$A \to B$ 以及 $c : I \to N$
如 (\ref{defos-equation-to-solve}) 所述。记
$\xi \in \Ext^1_A(\NL_{A/A'}, I)$ 为经引理
\ref{defos-lemma-extensions-of-algebras} 与 $A$ 被 $I$ 扩张所得的 $A'$
相对应的元素。解的同构类集合与同态
$$
\Ext^1_B(\NL_{B/A'}, N) \to \Ext^1_A(\NL_{A/A'}, N)
$$
在 $\xi$ 的像上方的纤维典范双射。
\end{lemma}

\begin{proof}
把引理 \ref{defos-lemma-extensions-of-algebras} 应用于 $A' \to B$ 和
$B$-模 $N$，可见 $\Ext^1_B(\NL_{B/A'}, N)$ 中的元素 $\zeta$
参数化 $A'$-代数扩张 $0 \to N \to B' \to B \to 0$。把引理
\ref{defos-lemma-extensions-of-algebras-functorial} 应用于
$A' \to A \to B$ 和 $c : I \to N$，可见存在与 $c$ 及 $A \to B$
相容的 $A'$-代数同态 $A' \to B'$，当且仅当 $\zeta$ 映到 $\xi$。
这当然等价于说 $B'$ 是 (\ref{defos-equation-to-solve}) 的一个解。
\end{proof}

\begin{remark}
\label{defos-remark-parametrize-solutions}
注意，在引理 \ref{defos-lemma-parametrize-solutions} 的情形中有
$$
\Ext^1_A(\NL_{A/A'}, N) =
\Ext^1_B(\NL_{A/A'} \otimes_A^\mathbf{L} B, N) =
\Ext^1_B(\NL_{A/A'} \otimes_A B, N)
$$
第一个等号来自《更多交换代数》引理
\ref{more-algebra-lemma-tensor-hom-adjoint}，第二个等号来自
《更多交换代数》引理 \ref{more-algebra-lemma-tensor-NL}。我们有复形同态
$$
\NL_{A/A'} \otimes_A B \to \NL_{B/A'} \to \NL_{B/A}
$$
它们几乎构成一个可区别三角，见《交换代数》引理
\ref{algebra-lemma-exact-sequence-NL}。如果它确实是可区别三角，便可得出：
$\xi$ 在 $\Ext^2_B(\NL_{B/A}, N)$ 中的像，就是
(\ref{defos-equation-to-solve}) 存在解的障碍。
\end{remark}

\noindent
如果环同态 $A \to B$ 是局部完全交同态，那么解存在。这是一类提升结果；
注意，对于 syntomic 环同态，我们已在《环同态的光滑化》命题
\ref{smoothing-proposition-lift-smooth} 中证明了一个相当强的提升结果。

\begin{lemma}
\label{defos-lemma-existence-lci}
如果 $A \to B$ 是局部完全交环同态，那么
(\ref{defos-equation-to-solve}) 存在解。
\end{lemma}

\begin{proof}[第一种证明]
写成 $B = A[x_1, \ldots, x_n]/J$。由《更多交换代数》定义
\ref{more-algebra-definition-local-complete-intersection}，理想 $J$ 是
Koszul 正则的。这蕴含 $J$ 是 $H_1$-正则且拟正则的，
见《更多交换代数》第 \ref{more-algebra-section-ideals} 节。令
$J' \subset A'[x_1, \ldots, x_n]$ 为 $J$ 的原像。以
$I[x_1, \ldots, x_n]$ 表示同态
$A'[x_1, \ldots, x_n] \to A[x_1, \ldots, x_n]$ 的核。由
《更多交换代数》引理
\ref{more-algebra-lemma-conormal-sequence-H1-regular-ideal}，有
$I[x_1, \ldots, x_n] \cap (J')^2 = J'I[x_1, \ldots, x_n] =
JI[x_1, \ldots, x_n]$。因此得到短正合列
$$
0 \to I \otimes_A B \to J'/(J')^2 \to J/J^2 \to 0
$$
由于 $J/J^2$ 是投射模（《更多交换代数》引理
\ref{more-algebra-lemma-quasi-regular-ideal-finite-projective}），
可以选取该正合列的一个分裂
$$
J'/(J')^2 = I \otimes_A B \oplus J/J^2 
$$
令 $(J')^2 \subset J'' \subset J'$ 为映到上述分解中第二个直和项的那些
元素。于是
$$
0 \to I \otimes_A B \to A'[x_1, \ldots, x_n]/J'' \to B \to 0
$$
是 $N = I \otimes_A B$ 时 (\ref{defos-equation-to-solve}) 的一个解。
沿给定同态 $I \otimes_A B \to N$ 作推出，即得一般情形。
\end{proof}

\begin{proof}[第二种证明]
请先阅读评注 \ref{defos-remark-parametrize-solutions}。由《更多交换代数》引理
\ref{more-algebra-lemma-transitive-lci-at-end}，同态
$\NL_{A/A'} \otimes_A B \to \NL_{B/A'} \to \NL_{B/A}$
在 $D(B)$ 中确实构成一个可区别三角。因此，只需证明
$\Ext^2_B(\NL_{B/A}, N)$ 为零。由《更多交换代数》引理
\ref{more-algebra-lemma-lci-NL}，复形 $\NL_{B/A}$ 是 Tor 振幅位于
$[-1,0]$ 的完美复形。例如，由《更多交换代数》引理
\ref{more-algebra-lemma-splitting-unique} 第 (1) 点，这蕴含上述
$\Ext^2$ 为零。
\end{proof}







\section{环化空间的增厚}
\label{defos-section-thickenings-spaces}

\noindent
在下面几节中，我们将使用以下概念：
\begin{enumerate}
\item 环化空间 $(X', \mathcal{O}_{X'})$ 上的理想层
$\mathcal{I} \subset \mathcal{O}_{X'}$ 称为{\it 局部幂零的}，如果
$\mathcal{I}$ 的任意局部截面都是局部幂零的。与《交换代数》第
\ref{algebra-item-ideal-locally-nilpotent} 项比较。
\item 环化空间的一个{\it 增厚}是环化空间的态射
$i : (X, \mathcal{O}_X) \to (X', \mathcal{O}_{X'})$，满足：
\begin{enumerate}
\item $i$ 诱导同胚 $X \to X'$；
\item 同态 $i^\sharp : \mathcal{O}_{X'} \to i_*\mathcal{O}_X$ 是满的；
\item $i^\sharp$ 的核是局部幂零理想层。
\end{enumerate}
\item 环化空间的一个{\it 一阶增厚}是环化空间的增厚
$i : (X, \mathcal{O}_X) \to (X', \mathcal{O}_{X'})$，且
$\Ker(i^\sharp)$ 平方为零。
\item 显然可以定义{\it 增厚的态射}、{\it 基环化空间上增厚的态射}等概念。
\end{enumerate}
若 $i : (X, \mathcal{O}_X) \to (X', \mathcal{O}_{X'})$ 是环化空间的
一个增厚，我们便等同其底层拓扑空间，并把 $\mathcal{O}_X$、
$\mathcal{O}_{X'}$ 以及 $\mathcal{I} = \Ker(i^\sharp)$ 都视为
$X = X'$ 上的层。由此得到 $\mathcal{O}_{X'}$-模的短正合列
$$
0 \to \mathcal{I} \to \mathcal{O}_{X'} \to \mathcal{O}_X \to 0
$$
由《模层》引理 \ref{modules-lemma-i-star-equivalence}，
$\mathcal{O}_X$-模范畴等价于被 $\mathcal{I}$ 零化的
$\mathcal{O}_{X'}$-模范畴。特别地，如果 $i$ 是一阶增厚，那么
$\mathcal{I}$ 是一个 $\mathcal{O}_X$-模。

\begin{situation}
\label{defos-situation-morphism-thickenings}
增厚的一个态射 $(f,f')$ 由如下交换图给出：
\begin{equation}
\label{defos-equation-morphism-thickenings}
\vcenter{
\xymatrix{
(X, \mathcal{O}_X) \ar[r]_i \ar[d]_f & (X', \mathcal{O}_{X'}) \ar[d]^{f'} \\
(S, \mathcal{O}_S) \ar[r]^t & (S', \mathcal{O}_{S'})
}
}
\end{equation}
其中各项为环化空间，水平箭头为增厚。在这一情形中，令
$\mathcal{I} = \Ker(i^\sharp) \subset \mathcal{O}_{X'}$、
$\mathcal{J} = \Ker(t^\sharp) \subset \mathcal{O}_{S'}$。由于在底层
拓扑空间上 $f=f'$，我们等同（拓扑）逆像函子 $f^{-1}$ 与
$(f')^{-1}$。注意，
$(f')^\sharp : f^{-1}\mathcal{O}_{S'} \to \mathcal{O}_{X'}$
特别地诱导同态 $f^{-1}\mathcal{J} \to \mathcal{I}$，从而诱导
$\mathcal{O}_{X'}$-模同态
$$
(f')^*\mathcal{J} \longrightarrow \mathcal{I}
$$
如果 $i$ 与 $t$ 都是一阶增厚，那么
$(f')^*\mathcal{J}=f^*\mathcal{J}$，上述同态就成为
$f^*\mathcal{J} \to \mathcal{I}$。
\end{situation}

\begin{definition}
\label{defos-definition-strict-morphism-thickenings}
在情形 \ref{defos-situation-morphism-thickenings} 中，如果同态
$(f')^*\mathcal{J} \longrightarrow \mathcal{I}$ 是满的，则称
$(f,f')$ 是一个{\it 增厚的严格态射}。
\end{definition}

\noindent
下述引理特别说明：概形增厚的态射
$(f,f'):(X\subset X')\to(S\subset S')$ 是严格的，当且仅当
$X=S\times_{S'}X'$。

\begin{lemma}
\label{defos-lemma-strict-morphism-thickenings}
在情形 \ref{defos-situation-morphism-thickenings} 中，态射 $(f,f')$
是增厚的严格态射，当且仅当
(\ref{defos-equation-morphism-thickenings}) 在环化空间范畴中是笛卡儿的。
\end{lemma}

\begin{proof}
从略。
\end{proof}




\section{环化空间一阶增厚上的模}
\label{defos-section-modules-thickenings}

\noindent
本节讨论模的形变理论所需的一些预备知识。设
$i : (X,\mathcal{O}_X)\to(X',\mathcal{O}_{X'})$ 是环化空间的一阶增厚。
我们将随时使用第 \ref{defos-section-thickenings-spaces} 节引入的记号；
特别地，我们等同底层拓扑空间。本节考虑 $\mathcal{O}_{X'}$-模的短正合列
\begin{equation}
\label{defos-equation-extension}
0 \to \mathcal{K} \to \mathcal{F}' \to \mathcal{F} \to 0
\end{equation}
其中 $\mathcal{F}$、$\mathcal{K}$ 是 $\mathcal{O}_X$-模，而
$\mathcal{F}'$ 是 $\mathcal{O}_{X'}$-模。在这种情形中，有典范的
$\mathcal{O}_X$-模同态
$$
c_{\mathcal{F}'} :
\mathcal{I} \otimes_{\mathcal{O}_X} \mathcal{F}
\longrightarrow
\mathcal{K}
$$
其中 $\mathcal{I}=\Ker(i^\sharp)$。具体而言，给定 $\mathcal{I}$ 的
局部截面 $f$ 与 $\mathcal{F}$ 的局部截面 $s$，令
$c_{\mathcal{F}'}(f\otimes s)=fs'$，其中 $s'$ 是提升 $s$ 的
$\mathcal{F}'$ 的局部截面。

\begin{lemma}
\label{defos-lemma-inf-map}
设 $i : (X,\mathcal{O}_X)\to(X',\mathcal{O}_{X'})$ 是环化空间的一阶
增厚，并给定扩张
$$
0 \to \mathcal{K} \to \mathcal{F}' \to \mathcal{F} \to 0
\quad\text{且}\quad
0 \to \mathcal{L} \to \mathcal{G}' \to \mathcal{G} \to 0
$$
如 (\ref{defos-equation-extension}) 所述，以及同态
$\varphi:\mathcal{F}\to\mathcal{G}$ 与
$\psi:\mathcal{K}\to\mathcal{L}$。
\begin{enumerate}
\item 如果存在与 $\varphi$ 和 $\psi$ 相容的 $\mathcal{O}_{X'}$-模同态
$\varphi':\mathcal{F}'\to\mathcal{G}'$，那么图表
$$
\xymatrix{
\mathcal{I} \otimes_{\mathcal{O}_X} \mathcal{F}
\ar[r]_-{c_{\mathcal{F}'}} \ar[d]_{1 \otimes \varphi} &
\mathcal{K} \ar[d]^\psi \\
\mathcal{I} \otimes_{\mathcal{O}_X} \mathcal{G}
\ar[r]^-{c_{\mathcal{G}'}} &
\mathcal{L}
}
$$
交换。
\item 与 $\varphi$ 和 $\psi$ 相容的 $\mathcal{O}_{X'}$-模同态
$\varphi':\mathcal{F}'\to\mathcal{G}'$ 的集合，如果非空，就是
$\Hom_{\mathcal{O}_X}(\mathcal{F},\mathcal{L})$ 作用下的主齐性空间。
\end{enumerate}
\end{lemma}

\begin{proof}
第 (1) 点由各同态的描述立即得到。对于第 (2) 点，若 $\varphi'$ 与
$\varphi''$ 是两个与 $\varphi$ 和 $\psi$ 相容的同态
$\mathcal{F}'\to\mathcal{G}'$，那么 $\varphi'-\varphi''$ 分解为
$$
\mathcal{F}' \to \mathcal{F} \to \mathcal{L} \to \mathcal{G}'
$$
由《模层》引理 \ref{modules-lemma-i-star-equivalence}，中间的同态来自
$\Hom_{\mathcal{O}_X}(\mathcal{F},\mathcal{L})$ 中唯一的元素。反过来，
给定该群中的元素 $\alpha$，可以把上述复合（中间同态取 $\alpha$）
加到 $\varphi'$ 上。若干细节从略。
\end{proof}

\begin{lemma}
\label{defos-lemma-inf-obs-map}
设 $i:(X,\mathcal{O}_X)\to(X',\mathcal{O}_{X'})$ 是环化空间的一阶
增厚，并给定扩张
$$
0 \to \mathcal{K} \to \mathcal{F}' \to \mathcal{F} \to 0
\quad\text{且}\quad
0 \to \mathcal{L} \to \mathcal{G}' \to \mathcal{G} \to 0
$$
如 (\ref{defos-equation-extension}) 所述，以及同态
$\varphi:\mathcal{F}\to\mathcal{G}$ 与
$\psi:\mathcal{K}\to\mathcal{L}$。假设图表
$$
\xymatrix{
\mathcal{I} \otimes_{\mathcal{O}_X} \mathcal{F}
\ar[r]_-{c_{\mathcal{F}'}} \ar[d]_{1 \otimes \varphi} &
\mathcal{K} \ar[d]^\psi \\
\mathcal{I} \otimes_{\mathcal{O}_X} \mathcal{G}
\ar[r]^-{c_{\mathcal{G}'}} &
\mathcal{L}
}
$$
交换。那么存在元素
$$
o(\varphi, \psi) \in
\Ext^1_{\mathcal{O}_X}(\mathcal{F}, \mathcal{L})
$$
它为零，当且仅当存在与 $\varphi$ 和 $\psi$ 相容的同态
$\varphi':\mathcal{F}'\to\mathcal{G}'$。
\end{lemma}

\begin{proof}
可以显式构造扩张
$$
0 \to \mathcal{L} \to \mathcal{H} \to \mathcal{F} \to 0
$$
方法是令 $\mathcal{H}$ 为复形
$$
\mathcal{K}
\xrightarrow{1, - \psi}
\mathcal{F}' \oplus \mathcal{G}' \xrightarrow{\varphi, 1}
\mathcal{G}
$$
在中间项处的上同调层（记号含义显然）。利用引理中图表交换这一假设，
对局部截面作计算可知 $\mathcal{H}$ 被 $\mathcal{I}$ 零化。因此
$\mathcal{H}$ 在下述群中定义了一个类：
$$
\Ext^1_{\mathcal{O}_X}(\mathcal{F}, \mathcal{L})
\subset
\Ext^1_{\mathcal{O}_{X'}}(\mathcal{F}, \mathcal{L})
$$
最后，$\mathcal{H}$ 的类等于沿 $\psi$ 对扩张 $\mathcal{F}'$ 所作的
推出与沿 $\varphi$ 对扩张 $\mathcal{G}'$ 所作的拉回之差（计算从略）。
因此，$\mathcal{H}$ 的类为零，等价于存在交换图
$$
\xymatrix{
0 \ar[r] &
\mathcal{K} \ar[r] \ar[d]_{\psi} &
\mathcal{F}' \ar[r] \ar[d]_{\varphi'} &
\mathcal{F} \ar[r] \ar[d]_\varphi & 0\\
0 \ar[r] &
\mathcal{L} \ar[r] &
\mathcal{G}' \ar[r] &
\mathcal{G} \ar[r] & 0
}
$$
这正是所需。
\end{proof}

\begin{lemma}
\label{defos-lemma-inf-ext}
设 $i:(X,\mathcal{O}_X)\to(X',\mathcal{O}_{X'})$ 是环化空间的一阶
增厚。给定 $\mathcal{O}_X$-模 $\mathcal{F}$、$\mathcal{K}$ 以及
$\mathcal{O}_X$-线性同态
$c:\mathcal{I}\otimes_{\mathcal{O}_X}\mathcal{F}\to\mathcal{K}$。
如果存在满足 $c_{\mathcal{F}'}=c$ 的正合列
(\ref{defos-equation-extension})，那么这些扩张的同构类集合是
$\Ext^1_{\mathcal{O}_X}(\mathcal{F},\mathcal{K})$ 作用下的主齐性空间。
\end{lemma}

\begin{proof}
假设给定扩张
$$
0 \to \mathcal{K} \to \mathcal{F}'_1 \to \mathcal{F} \to 0
\quad\text{且}\quad
0 \to \mathcal{K} \to \mathcal{F}'_2 \to \mathcal{F} \to 0
$$
且 $c_{\mathcal{F}'_1}=c_{\mathcal{F}'_2}=c$。那么二者之差
（在扩张群中取差，见《同调代数》第 \ref{homology-section-extensions} 节）
是扩张
$$
0 \to \mathcal{K} \to \mathcal{E} \to \mathcal{F} \to 0
$$
其中 $\mathcal{E}$ 被 $\mathcal{I}$ 零化（局部计算从略）。所以该正合列
是 $\mathcal{O}_X$-模的扩张，见《模层》引理
\ref{modules-lemma-i-star-equivalence}。反过来，给定这样的扩张
$\mathcal{E}$，可将它加到 $\mathcal{O}_{X'}$-扩张 $\mathcal{F}'$ 上，
而不改变同态 $c_{\mathcal{F}'}$。若干细节从略。
\end{proof}

\begin{lemma}
\label{defos-lemma-inf-obs-ext}
设 $i:(X,\mathcal{O}_X)\to(X',\mathcal{O}_{X'})$ 是环化空间的一阶
增厚。给定 $\mathcal{O}_X$-模 $\mathcal{F}$、$\mathcal{K}$ 以及
$\mathcal{O}_X$-线性同态
$c:\mathcal{I}\otimes_{\mathcal{O}_X}\mathcal{F}\to\mathcal{K}$。
那么存在元素
$$
o(\mathcal{F}, \mathcal{K}, c) \in
\Ext^2_{\mathcal{O}_X}(\mathcal{F}, \mathcal{K})
$$
它为零，当且仅当存在满足 $c_{\mathcal{F}'}=c$ 的正合列
(\ref{defos-equation-extension})。
\end{lemma}

\begin{proof}
先证明：如果 $\mathcal{K}$ 是内射 $\mathcal{O}_X$-模，那么确实存在
满足 $c_{\mathcal{F}'}=c$ 的正合列 (\ref{defos-equation-extension})。
为此，选取平坦 $\mathcal{O}_{X'}$-模 $\mathcal{H}'$ 以及满同态
$\mathcal{H}'\to\mathcal{F}$（《模层》引理
\ref{modules-lemma-module-quotient-flat}）。令
$\mathcal{J}\subset\mathcal{H}'$ 为其核。由于 $\mathcal{H}'$ 平坦，故有
$$
\mathcal{I} \otimes_{\mathcal{O}_{X'}} \mathcal{H}' =
\mathcal{I}\mathcal{H}'
\subset \mathcal{J} \subset \mathcal{H}'
$$
注意，同态
$$
\mathcal{I}\mathcal{H}' =
\mathcal{I} \otimes_{\mathcal{O}_{X'}} \mathcal{H}'
\longrightarrow
\mathcal{I} \otimes_{\mathcal{O}_{X'}} \mathcal{F} =
\mathcal{I} \otimes_{\mathcal{O}_X} \mathcal{F}
$$
在 $\mathcal{I}\mathcal{J}$ 上为零。事实上，如果 $f$ 是
$\mathcal{I}$ 的局部截面、$s$ 是 $\mathcal{H}$ 的局部截面，那么
$fs$ 映到 $f\otimes\overline{s}$，其中 $\overline{s}$ 是 $s$ 在
$\mathcal{F}$ 中的像。因此得到如下 $\mathcal{O}_X$-模图表：
$$
\xymatrix{
\mathcal{I}\mathcal{H}'/\mathcal{I}\mathcal{J}
\ar@{^{(}->}[r] \ar[d] &
\mathcal{J}/\mathcal{I}\mathcal{J} \ar@{..>}[d]_\gamma \\
\mathcal{I} \otimes_{\mathcal{O}_X} \mathcal{F} \ar[r]^-c &
\mathcal{K}
}
$$
如果 $\mathcal{K}$ 是内射 $\mathcal{O}_X$-模，就得到虚线箭头。
以 $\gamma':\mathcal{J}\to\mathcal{K}$ 表示 $\gamma$ 与
$\mathcal{J}\to\mathcal{J}/\mathcal{I}\mathcal{J}$ 的复合。
局部计算表明，下列推出
$$
\xymatrix{
0 \ar[r] &
\mathcal{J} \ar[r] \ar[d]_{\gamma'} &
\mathcal{H}' \ar[r] \ar[d] &
\mathcal{F} \ar[r] \ar@{=}[d] &
0 \\
0 \ar[r] &
\mathcal{K} \ar[r] &
\mathcal{F}' \ar[r] &
\mathcal{F} \ar[r] &
0
}
$$
是本引理所提问题的一个解。

\medskip\noindent
一般情形。选取嵌入 $\mathcal{K}\subset\mathcal{K}'$，其中
$\mathcal{K}'$ 是内射 $\mathcal{O}_X$-模。令 $\mathcal{Q}$ 为商模，
于是有正合列
$$
0 \to \mathcal{K} \to \mathcal{K}' \to \mathcal{Q} \to 0
$$
以
$c':\mathcal{I}\otimes_{\mathcal{O}_X}\mathcal{F}\to\mathcal{K}'$
表示相应的复合。由上一段，存在正合列
$$
0 \to \mathcal{K}' \to \mathcal{E}' \to \mathcal{F} \to 0
$$
如 (\ref{defos-equation-extension}) 所述，且 $c_{\mathcal{E}'}=c'$。
注意，$c'$ 与同态 $\mathcal{K}'\to\mathcal{Q}$ 的复合为零；因此，
沿 $\mathcal{K}'\to\mathcal{Q}$ 对 $\mathcal{E}'$ 作推出所得的是扩张
$$
0 \to \mathcal{Q} \to \mathcal{D}' \to \mathcal{F} \to 0
$$
如 (\ref{defos-equation-extension}) 所述，且 $c_{\mathcal{D}'}=0$。
这恰好意味着 $\mathcal{D}'$ 被 $\mathcal{I}$ 零化；换言之，
$\mathcal{D}'$ 是 $\mathcal{O}_X$-模的扩张，从而定义元素
$$
o(\mathcal{F}, \mathcal{K}, c) \in
\Ext^1_{\mathcal{O}_X}(\mathcal{F}, \mathcal{Q}) =
\Ext^2_{\mathcal{O}_X}(\mathcal{F}, \mathcal{K})
$$
（等号来自与上述正合列相伴的长正合上同调列，以及取值于内射模
$\mathcal{K}'$ 的高次 Ext 群为零）。如果
$o(\mathcal{F},\mathcal{K},c)=0$，便可选取分裂
$s:\mathcal{F}\to\mathcal{D}'$，并令
$$
\mathcal{F}' = \Ker(\mathcal{E}' \to \mathcal{D}'/s(\mathcal{F}))
$$
于是得到下图：
$$
\xymatrix{
0 \ar[r] &
\mathcal{K} \ar[r] \ar[d] &
\mathcal{F}' \ar[r] \ar[d] &
\mathcal{F} \ar[r] \ar@{=}[d] &
0 \\
0 \ar[r] &
\mathcal{K}' \ar[r] &
\mathcal{E}' \ar[r] &
\mathcal{F} \ar[r] & 0
}
$$
两行均正合，这说明 $c_{\mathcal{F}'}=c$。反过来，如果
$\mathcal{F}'$ 存在，那么由引理 \ref{defos-lemma-inf-ext} 以及取值于内射模
$\mathcal{K}'$ 的高次 Ext 群为零，沿 $\mathcal{K}\to\mathcal{K}'$
对 $\mathcal{F}'$ 作推出所得的扩张同构于 $\mathcal{E}'$。由此得到
上述形式的图表，它蕴含 $\mathcal{D}'$ 作为扩张是分裂的；也就是说，
类 $o(\mathcal{F},\mathcal{K},c)$ 为零。
\end{proof}

\begin{remark}
\label{defos-remark-trivial-thickening}
设 $(X,\mathcal{O}_X)$ 为环化空间。如果存在环化空间的态射
$\pi:(X',\mathcal{O}_{X'})\to(X,\mathcal{O}_X)$，它是 $i$ 的左逆，
则称一阶增厚
$i:(X,\mathcal{O}_X)\to(X',\mathcal{O}_{X'})$ 是{\it 平凡的}。
这样的态射 $\pi$ 的一个选择称为该一阶增厚的一个{\it 平凡化}。
给定 $\pi$，利用 $\pi^\sharp$ 分裂满同态
$\mathcal{O}_{X'}\to\mathcal{O}_X$，便得到 $X$ 上代数层的分裂
\begin{equation}
\label{defos-equation-splitting}
\mathcal{O}_{X'} = \mathcal{O}_X \oplus \mathcal{I}
\end{equation}
反过来，这样的分裂确定一个态射 $\pi$。环化空间
$(X,\mathcal{O}_X)$ 的带平凡化一阶增厚所成的范畴，等价于
$\mathcal{O}_X$-模范畴。
\end{remark}

\begin{remark}
\label{defos-remark-trivial-extension}
设 $i:(X,\mathcal{O}_X)\to(X',\mathcal{O}_{X'})$ 是环化空间的平凡
一阶增厚，而 $\pi:(X',\mathcal{O}_{X'})\to(X,\mathcal{O}_X)$ 是一个
平凡化。给定任意三元组 $(\mathcal{F},\mathcal{K},c)$，其中
$\mathcal{F}$、$\mathcal{K}$ 是 $\mathcal{O}_X$-模，且
$c:\mathcal{I}\otimes_{\mathcal{O}_X}\mathcal{F}\to\mathcal{K}$，
可以令
$$
\mathcal{F}'_{c, triv} = \mathcal{F} \oplus \mathcal{K}
$$
再用与 $\pi$ 相伴的分裂 (\ref{defos-equation-splitting}) 和同态 $c$
定义 $\mathcal{O}_{X'}$-模结构，从而得到扩张
(\ref{defos-equation-extension})。我们称 $\mathcal{F}'_{c,triv}$ 为
由 $c$ 与平凡化 $\pi$ 所对应的 $\mathcal{F}$ 被 $\mathcal{K}$
扩张所得的{\it 平凡扩张}。给定 (\ref{defos-equation-extension}) 中的任意
扩张 $\mathcal{F}'$，可以利用
$\pi^\sharp:\mathcal{O}_X\to\mathcal{O}_{X'}$ 把 $\mathcal{F}'$
看成 $\mathcal{O}_X$-模的扩张，因而得到
$\Ext^1_{\mathcal{O}_X}(\mathcal{F},\mathcal{K})$ 中的类
$\xi_{\mathcal{F}'}$。引理 \ref{defos-lemma-inf-ext} 保证映射
$\mathcal{F}' \mapsto \xi_{\mathcal{F}'}$
诱导双射
$$
\left\{
\begin{matrix}
\text{扩张的同构类}\\
\mathcal{F}'\text{ 如 (\ref{defos-equation-extension}) 所述，且 }c=c_{\mathcal{F}'}
\end{matrix}
\right\}
\longrightarrow
\Ext^1_{\mathcal{O}_X}(\mathcal{F}, \mathcal{K})
$$
此外，平凡扩张 $\mathcal{F}'_{c,triv}$ 映到零类。
\end{remark}

\begin{remark}
\label{defos-remark-extension-functorial}
设 $(X,\mathcal{O}_X)$ 为环化空间，并设
$(X,\mathcal{O}_X)\to(X'_i,\mathcal{O}_{X'_i})$（$i=1,2$）为一阶
增厚，其理想层为 $\mathcal{I}_i$。设
$h:(X'_1,\mathcal{O}_{X'_1})\to(X'_2,\mathcal{O}_{X'_2})$
为 $(X,\mathcal{O}_X)$ 的一阶增厚之间的态射。图示如下：
$$
\xymatrix{
& (X, \mathcal{O}_X) \ar[ld] \ar[rd] & \\
(X'_1, \mathcal{O}_{X'_1}) \ar[rr]^h & & 
(X'_2, \mathcal{O}_{X'_2})
}
$$
注意，$h^\sharp:\mathcal{O}_{X'_2}\to\mathcal{O}_{X'_1}$ 特别地诱导
$\mathcal{O}_X$-模同态 $\mathcal{I}_2\to\mathcal{I}_1$。设
$\mathcal{F}$ 为 $\mathcal{O}_X$-模。对 $i=1,2$，设
$(\mathcal{K}_i,c_i)$ 由 $\mathcal{O}_X$-模 $\mathcal{K}_i$ 及同态
$c_i:\mathcal{I}_i\otimes_{\mathcal{O}_X}\mathcal{F}\to\mathcal{K}_i$
组成。再给定 $\mathcal{O}_X$-模同态
$\mathcal{K}_2\to\mathcal{K}_1$，使得
$$
\xymatrix{
\mathcal{I}_2 \otimes_{\mathcal{O}_X} \mathcal{F}
\ar[r]_-{c_2} \ar[d] &
\mathcal{K}_2 \ar[d] \\
\mathcal{I}_1 \otimes_{\mathcal{O}_X} \mathcal{F}
\ar[r]^-{c_1} &
\mathcal{K}_1
}
$$
交换。那么存在如下典范函子性：
$$
\left\{
\begin{matrix}
\mathcal{F}'_2\text{ 如 (\ref{defos-equation-extension}) 所述，且}\\
c_2=c_{\mathcal{F}'_2}\text{、}\mathcal{K}=\mathcal{K}_2
\end{matrix}
\right\}
\longrightarrow
\left\{
\begin{matrix}
\mathcal{F}'_1\text{ 如 (\ref{defos-equation-extension}) 所述，且}\\
c_1=c_{\mathcal{F}'_1}\text{、}\mathcal{K}=\mathcal{K}_1
\end{matrix}
\right\}
$$
具体而言，把 $\mathcal{O}_X$、$\mathcal{O}_{X'_i}$、$\mathcal{F}$、
$\mathcal{K}_i$ 等都视为 $X$ 上的层；给定 $\mathcal{F}'_2$ 后，
令层 $\mathcal{F}'_1$ 为推出，也就是说，它置于下列扩张图表中：
$$
\xymatrix{
0 \ar[r] &
\mathcal{K}_2 \ar[r] \ar[d] &
\mathcal{F}'_2 \ar[r] \ar[d] &
\mathcal{F} \ar@{=}[d] \ar[r] & 0 \\
0 \ar[r] &
\mathcal{K}_1 \ar[r] &
\mathcal{F}'_1 \ar[r] &
\mathcal{F} \ar[r] & 0
}
$$
推出上的 $\mathcal{O}_{X'_1}$-模结构之构造从略（该构造使用涉及
$c_1$ 与 $c_2$ 的图表的交换性）。
\end{remark}

\begin{remark}
\label{defos-remark-trivial-extension-functorial}
设 $(X,\mathcal{O}_X)$、
$(X,\mathcal{O}_X)\to(X'_i,\mathcal{O}_{X'_i})$、$\mathcal{I}_i$ 以及
$h:(X'_1,\mathcal{O}_{X'_1})\to(X'_2,\mathcal{O}_{X'_2})$ 如评注
\ref{defos-remark-extension-functorial} 所述。假设给定平凡化
$\pi_i:X'_i\to X$，满足 $\pi_1=h\circ\pi_2$。换言之，假设 $h$ 是
$(X,\mathcal{O}_X)$ 的带平凡化一阶增厚之间的态射。对 $i=1,2$，
设 $(\mathcal{K}_i,c_i)$ 由 $\mathcal{O}_X$-模 $\mathcal{K}_i$
及同态
$c_i:\mathcal{I}_i\otimes_{\mathcal{O}_X}\mathcal{F}\to\mathcal{K}_i$
组成。再给定 $\mathcal{O}_X$-模同态
$\mathcal{K}_2\to\mathcal{K}_1$，使得
$$
\xymatrix{
\mathcal{I}_2 \otimes_{\mathcal{O}_X} \mathcal{F}
\ar[r]_-{c_2} \ar[d] &
\mathcal{K}_2 \ar[d] \\
\mathcal{I}_1 \otimes_{\mathcal{O}_X} \mathcal{F}
\ar[r]^-{c_1} &
\mathcal{K}_1
}
$$
交换。在这种情形下，评注 \ref{defos-remark-trivial-extension} 的构造诱导
交换图
$$
\xymatrix{
\{\mathcal{F}'_2\text{ 如 (\ref{defos-equation-extension}) 所述，且 }
c_2=c_{\mathcal{F}'_2}\text{、}\mathcal{K}=\mathcal{K}_2\}
\ar[d] \ar[rr] & &
\Ext^1_{\mathcal{O}_X}(\mathcal{F}, \mathcal{K}_2) \ar[d] \\
\{\mathcal{F}'_1\text{ 如 (\ref{defos-equation-extension}) 所述，且 }
c_1=c_{\mathcal{F}'_1}\text{、}\mathcal{K}=\mathcal{K}_1\}
\ar[rr] & &
\Ext^1_{\mathcal{O}_X}(\mathcal{F}, \mathcal{K}_1)
}
$$
其中右侧竖直同态由 $\Ext$ 的函子性与同态
$\mathcal{K}_2\to\mathcal{K}_1$ 给出，左侧竖直映射则是评注
\ref{defos-remark-extension-functorial} 中的映射。
\end{remark}

\begin{remark}
\label{defos-remark-short-exact-sequence-thickenings}
设 $(X,\mathcal{O}_X)$ 为环化空间。对于 $(X,\mathcal{O}_X)$ 的一阶
增厚态射列
$$
(X'_1, \mathcal{O}_{X'_1}) \to
(X'_2, \mathcal{O}_{X'_2}) \to
(X'_3, \mathcal{O}_{X'_3})
$$
如果理想层 $\mathcal{I}_i$ 之间的相应同态给出
$\mathcal{O}_X$-模复形
$\mathcal{I}_3\to\mathcal{I}_2\to\mathcal{I}_1$
（即复合为零），则称该态射列为一个{\it 复形}。在这种情形下，复合
$(X'_1,\mathcal{O}_{X'_1})\to(X'_3,\mathcal{O}_{X'_3})$ 经
$(X,\mathcal{O}_X)\to(X'_3,\mathcal{O}_{X'_3})$ 分解；换言之，
$(X,\mathcal{O}_X)$ 的一阶增厚 $(X'_1,\mathcal{O}_{X'_1})$
是平凡的，并且带有典范平凡化
$\pi:(X'_1,\mathcal{O}_{X'_1})\to(X,\mathcal{O}_X)$。

\medskip\noindent
如果 $(X,\mathcal{O}_X)$ 的一阶增厚态射列
$$
(X'_1, \mathcal{O}_{X'_1}) \to
(X'_2, \mathcal{O}_{X'_2}) \to
(X'_3, \mathcal{O}_{X'_3})
$$
所对应的理想层同态构成 $\mathcal{O}_X$-模的短正合列
$$
0 \to \mathcal{I}_3 \to \mathcal{I}_2 \to \mathcal{I}_1 \to 0
$$
则称该态射列为一个{\it 短正合列}。
\end{remark}

\begin{remark}
\label{defos-remark-complex-thickenings-and-ses-modules}
设 $(X,\mathcal{O}_X)$ 为环化空间，$\mathcal{F}$ 为
$\mathcal{O}_X$-模。设
$$
(X'_1, \mathcal{O}_{X'_1}) \to
(X'_2, \mathcal{O}_{X'_2}) \to
(X'_3, \mathcal{O}_{X'_3})
$$
是 $(X,\mathcal{O}_X)$ 的一阶增厚所成的复形，见评注
\ref{defos-remark-short-exact-sequence-thickenings}。对 $i=1,2,3$，设
$(\mathcal{K}_i,c_i)$ 由 $\mathcal{O}_X$-模 $\mathcal{K}_i$ 与同态
$c_i:\mathcal{I}_i\otimes_{\mathcal{O}_X}\mathcal{F}\to\mathcal{K}_i$
组成。假设给定 $\mathcal{O}_X$-模的短正合列
$$
0 \to \mathcal{K}_3 \to \mathcal{K}_2 \to \mathcal{K}_1 \to 0
$$
使得图表
$$
\vcenter{
\xymatrix{
\mathcal{I}_2 \otimes_{\mathcal{O}_X} \mathcal{F}
\ar[r]_-{c_2} \ar[d] &
\mathcal{K}_2 \ar[d] \\
\mathcal{I}_1 \otimes_{\mathcal{O}_X} \mathcal{F}
\ar[r]^-{c_1} &
\mathcal{K}_1
}
}
\quad\text{且}\quad
\vcenter{
\xymatrix{
\mathcal{I}_3 \otimes_{\mathcal{O}_X} \mathcal{F}
\ar[r]_-{c_3} \ar[d] &
\mathcal{K}_3 \ar[d] \\
\mathcal{I}_2 \otimes_{\mathcal{O}_X} \mathcal{F}
\ar[r]^-{c_2} &
\mathcal{K}_2
}
}
$$
均交换。最后，假设给定 $\mathcal{O}_{X'_2}$-模的扩张
$$
0 \to \mathcal{K}_2 \to \mathcal{F}'_2 \to \mathcal{F} \to 0
$$
如 (\ref{defos-equation-extension}) 所述，其中
$\mathcal{K}=\mathcal{K}_2$ 且 $c_{\mathcal{F}'_2}=c_2$。在这种情形下，
应用评注 \ref{defos-remark-extension-functorial} 的函子性，可得 $X'_1$
上的扩张 $\mathcal{F}'_1$（稍后会描述这一特殊情形下的
$\mathcal{F}'_1$）。由评注 \ref{defos-remark-trivial-extension}，利用评注
\ref{defos-remark-short-exact-sequence-thickenings} 中的典范分裂
$\pi:(X'_1,\mathcal{O}_{X'_1})\to(X,\mathcal{O}_X)$，可得
$\xi_{\mathcal{F}'_1} \in
\Ext^1_{\mathcal{O}_X}(\mathcal{F}, \mathcal{K}_1)$.
最后，还有障碍
$$
o(\mathcal{F}, \mathcal{K}_3, c_3) \in
\Ext^2_{\mathcal{O}_X}(\mathcal{F}, \mathcal{K}_3)
$$
见引理 \ref{defos-lemma-inf-obs-ext}。在这种情形下，我们{\bf 断言}典范同态
$$
\partial :
\Ext^1_{\mathcal{O}_X}(\mathcal{F}, \mathcal{K}_1)
\longrightarrow
\Ext^2_{\mathcal{O}_X}(\mathcal{F}, \mathcal{K}_3)
$$
由短正合列
$0\to\mathcal{K}_3\to\mathcal{K}_2\to\mathcal{K}_1\to0$
诱导，并把 $\xi_{\mathcal{F}'_1}$ 映到障碍类
$o(\mathcal{F},\mathcal{K}_3,c_3)$。

\medskip\noindent
为证明该断言，选取嵌入 $j:\mathcal{K}_3\to\mathcal{K}$，其中
$\mathcal{K}$ 是内射 $\mathcal{O}_X$-模。可以把 $j$ 提升为同态
$j':\mathcal{K}_2\to\mathcal{K}$。令
$\mathcal{E}'_2=j'_*\mathcal{F}'_2$ 为沿 $j'$ 对
$\mathcal{F}'_2$ 所作的推出，于是
$c_{\mathcal{E}'_2}=j'\circ c_2$。图示如下：
$$
\xymatrix{
0 \ar[r] &
\mathcal{K}_2 \ar[r] \ar[d]_{j'} &
\mathcal{F}'_2 \ar[r] \ar[d] &
\mathcal{F} \ar[r] \ar[d] & 0 \\
0 \ar[r] &
\mathcal{K} \ar[r] &
\mathcal{E}'_2 \ar[r] &
\mathcal{F} \ar[r] & 0
}
$$
令 $\mathcal{E}'_3=\mathcal{E}'_2$，但经
$\mathcal{O}_{X'_3}\to\mathcal{O}_{X'_2}$ 把它视为
$\mathcal{O}_{X'_3}$-模。于是
$c_{\mathcal{E}'_3}=j\circ c_3$。引理 \ref{defos-lemma-inf-obs-ext} 的证明把
$o(\mathcal{F},\mathcal{K}_3,c_3)$ 构造为下列
$\mathcal{O}_X$-模扩张之类的边界：
$$
0 \to
\mathcal{K}/\mathcal{K}_3 \to
\mathcal{E}'_3/\mathcal{K}_3 \to
\mathcal{F} \to 0
$$
另一方面，注意 $\mathcal{F}'_1=\mathcal{F}'_2/\mathcal{K}_3$，故类
$\xi_{\mathcal{F}'_1}$ 是扩张
$$
0 \to \mathcal{K}_2/\mathcal{K}_3 \to \mathcal{F}'_2/\mathcal{K}_3
\to \mathcal{F} \to 0
$$
的类；这里利用 $\pi^\sharp$ 把该正合列视为 $\mathcal{O}_X$-模列，
而 $\pi:(X'_1,\mathcal{O}_{X'_1})\to(X,\mathcal{O}_X)$ 是典范分裂。
最后，该断言由下列交换图推出：
$$
\xymatrix{
0 \ar[r] &
\mathcal{K}_2/\mathcal{K}_3 \ar[r] \ar[d] &
\mathcal{F}'_2/\mathcal{K}_3 \ar[r] \ar[d] &
\mathcal{F} \ar[r] \ar[d] & 0 \\
0 \ar[r] &
\mathcal{K}/\mathcal{K}_3 \ar[r] &
\mathcal{E}'_3/\mathcal{K}_3 \ar[r] &
\mathcal{F} \ar[r] & 0
}
$$
该图对于上述 $\mathcal{O}_X$-模结构是 $\mathcal{O}_X$-线性的。
\end{remark}








\section{环化空间上模的无穷小形变}
\label{defos-section-deformation-modules}

\noindent
设 $i:(X,\mathcal{O}_X)\to(X',\mathcal{O}_{X'})$ 是环化空间的一阶
增厚。我们随时使用第 \ref{defos-section-thickenings-spaces} 节引入的记号。
设 $\mathcal{F}'$ 为 $\mathcal{O}_{X'}$-模，并令
$\mathcal{F}=i^*\mathcal{F}'$。在这种情形下，有
$\mathcal{O}_{X'}$-模的短正合列
$$
0 \to \mathcal{I}\mathcal{F}' \to \mathcal{F}' \to \mathcal{F} \to 0
$$
由于 $\mathcal{I}^2=0$，$\mathcal{I}\mathcal{F}'$ 上的
$\mathcal{O}_{X'}$-模结构来自唯一的 $\mathcal{O}_X$-模结构。因此，
上述正合列是 (\ref{defos-equation-extension}) 中的扩张。作为特例，如果
$\mathcal{F}'=\mathcal{O}_{X'}$，则
$i^*\mathcal{O}_{X'}=\mathcal{O}_X$ 且
$\mathcal{I}\mathcal{O}_{X'}=\mathcal{I}$，从而恢复结构层正合列
$$
0 \to \mathcal{I} \to \mathcal{O}_{X'} \to \mathcal{O}_X \to 0
$$

\begin{lemma}
\label{defos-lemma-inf-map-special}
设 $i:(X,\mathcal{O}_X)\to(X',\mathcal{O}_{X'})$ 是环化空间的一阶
增厚。设 $\mathcal{F}'$、$\mathcal{G}'$ 为
$\mathcal{O}_{X'}$-模。令
$\mathcal{F}=i^*\mathcal{F}'$、$\mathcal{G}=i^*\mathcal{G}'$，并设
$\varphi:\mathcal{F}\to\mathcal{G}$ 为 $\mathcal{O}_X$-线性同态。
$\varphi$ 到 $\mathcal{O}_{X'}$-线性同态
$\varphi':\mathcal{F}'\to\mathcal{G}'$ 的所有提升所成的集合，如果非空，
就是 $\Hom_{\mathcal{O}_X}(\mathcal{F},\mathcal{I}\mathcal{G}')$
作用下的主齐性空间。
\end{lemma}

\begin{proof}
这是引理 \ref{defos-lemma-inf-map} 的特例，不过我们也给出直接证明。我们有
模的短正合列
$$
0 \to \mathcal{I} \to \mathcal{O}_{X'} \to \mathcal{O}_X \to 0
\quad\text{且}\quad
0 \to \mathcal{I}\mathcal{G}' \to \mathcal{G}' \to \mathcal{G} \to 0
$$
$\mathcal{F}'$ 也有类似的正合列。由于 $\mathcal{I}$ 平方为零，
$\mathcal{I}$ 与 $\mathcal{I}\mathcal{G}'$ 上的
$\mathcal{O}_{X'}$-模结构均来自唯一的 $\mathcal{O}_X$-模结构。
因此
$$
\Hom_{\mathcal{O}_{X'}}(\mathcal{F}', \mathcal{I}\mathcal{G}') =
\Hom_{\mathcal{O}_X}(\mathcal{F}, \mathcal{I}\mathcal{G}')
\quad\text{且}\quad
\Hom_{\mathcal{O}_{X'}}(\mathcal{F}', \mathcal{G}) =
\Hom_{\mathcal{O}_X}(\mathcal{F}, \mathcal{G})
$$
本引理于是由正合列
$$
0 \to \Hom_{\mathcal{O}_{X'}}(\mathcal{F}', \mathcal{I}\mathcal{G}') \to
\Hom_{\mathcal{O}_{X'}}(\mathcal{F}', \mathcal{G}') \to
\Hom_{\mathcal{O}_{X'}}(\mathcal{F}', \mathcal{G})
$$
推出，见《同调代数》引理 \ref{homology-lemma-check-exactness}。
\end{proof}

\begin{lemma}
\label{defos-lemma-deform-module}
设 $(f,f')$ 是情形 \ref{defos-situation-morphism-thickenings} 中环化空间
一阶增厚的态射。设 $\mathcal{F}'$ 为 $\mathcal{O}_{X'}$-模，并令
$\mathcal{F}=i^*\mathcal{F}'$。假设 $\mathcal{F}$ 在 $S$ 上平坦，
且 $(f,f')$ 是增厚的严格态射（定义
\ref{defos-definition-strict-morphism-thickenings}）。那么下列条件等价：
\begin{enumerate}
\item $\mathcal{F}'$ 在 $S'$ 上平坦；
\item 典范同态
$f^*\mathcal{J} \otimes_{\mathcal{O}_X} \mathcal{F} \to
\mathcal{I}\mathcal{F}'$
是同构。
\end{enumerate}
此外，在这种情形下，同态
$$
f^*\mathcal{J} \otimes_{\mathcal{O}_X} \mathcal{F} \to
\mathcal{I} \otimes_{\mathcal{O}_X} \mathcal{F} \to
\mathcal{I}\mathcal{F}'
$$
均为同构。
\end{lemma}

\begin{proof}
由于 $(f,f')$ 是增厚的严格态射，同态
$f^*\mathcal{J}\to\mathcal{I}$ 是满的。因此，最后的陈述由 (2) 推出。

\medskip\noindent
证明 (1) 与 (2) 等价。可以在茎上检验这些条件。设
$x\in X\subset X'$，其像为 $s=f(x)\in S\subset S'$。令
$A'=\mathcal{O}_{S',s}$、$B'=\mathcal{O}_{X',x}$、
$A=\mathcal{O}_{S,s}$、$B=\mathcal{O}_{X,x}$。那么对某些平方零理想
$J$、$I$，有 $A=A'/J$、$B=B'/I$。由于 $(f,f')$ 是增厚的严格态射，
有 $I=JB'$。令 $M'=\mathcal{F}'_x$、$M=\mathcal{F}_x$。于是 $M'$
是 $B'$-模，$M$ 是 $B$-模。由 $\mathcal{F}=i^*\mathcal{F}'$ 可知，
满同态 $M'\to M$ 的核为 $IM'=JM'$。因而有短正合列
$$
0 \to JM' \to M' \to M \to 0
$$
利用《空间上的层》引理 \ref{sheaves-lemma-stalk-pullback-modules}
与《模层》引理 \ref{modules-lemma-stalk-tensor-product} 等同拉回与张量积
的茎，可见本引理中的典范同态在 $x$ 处的茎是同态
$$
(J \otimes_A B) \otimes_B M = J \otimes_A M = J \otimes_{A'} M'
\longrightarrow JM'
$$
假设 $\mathcal{F}$ 在 $S$ 上平坦，意味着 $M$ 是平坦 $A$-模。

\medskip\noindent
假设 (1) 成立。由《交换代数》引理
\ref{algebra-lemma-characterize-flat}，平坦性蕴含
$\text{Tor}_1^{A'}(M',A)=0$。由《交换代数》评注
\ref{algebra-remark-Tor-ring-mod-ideal}，这意味着
$J\otimes_{A'}M'\to M'$ 是单的。因此
$J\otimes_A M\to JM'$ 是同构。

\medskip\noindent
假设 (2) 成立。那么 $J\otimes_{A'}M'\to M'$ 是单的。由《交换代数》
评注 \ref{algebra-remark-Tor-ring-mod-ideal}，
$\text{Tor}_1^{A'}(M',A)=0$。再由《交换代数》引理
\ref{algebra-lemma-what-does-it-mean}，$M'$ 在 $A'$ 上平坦。
\end{proof}

\begin{lemma}
\label{defos-lemma-inf-map-rel}
设 $(f,f')$ 是情形 \ref{defos-situation-morphism-thickenings} 中一阶增厚的
态射。设 $\mathcal{F}'$、$\mathcal{G}'$ 为
$\mathcal{O}_{X'}$-模，并令
$\mathcal{F}=i^*\mathcal{F}'$、$\mathcal{G}=i^*\mathcal{G}'$。
设 $\varphi:\mathcal{F}\to\mathcal{G}$ 为 $\mathcal{O}_X$-线性同态。
假设 $\mathcal{G}'$ 在 $S'$ 上平坦，且 $(f,f')$ 是增厚的严格态射。
$\varphi$ 到 $\mathcal{O}_{X'}$-线性同态
$\varphi':\mathcal{F}'\to\mathcal{G}'$ 的所有提升所成的集合，如果非空，
就是下列群作用下的主齐性空间：
$$
\Hom_{\mathcal{O}_X}(\mathcal{F},
\mathcal{G} \otimes_{\mathcal{O}_X} f^*\mathcal{J})
$$
\end{lemma}

\begin{proof}
合并引理 \ref{defos-lemma-inf-map-special} 与 \ref{defos-lemma-deform-module} 即得。
\end{proof}

\begin{lemma}
\label{defos-lemma-inf-obs-map-special}
设 $i:(X,\mathcal{O}_X)\to(X',\mathcal{O}_{X'})$ 是环化空间的一阶
增厚。设 $\mathcal{F}'$、$\mathcal{G}'$ 为
$\mathcal{O}_{X'}$-模，并令
$\mathcal{F}=i^*\mathcal{F}'$、$\mathcal{G}=i^*\mathcal{G}'$。
设 $\varphi:\mathcal{F}\to\mathcal{G}$ 为 $\mathcal{O}_X$-线性同态。
存在元素
$$
o(\varphi) \in
\Ext^1_{\mathcal{O}_X}(Li^*\mathcal{F}',
\mathcal{I}\mathcal{G}')
$$
它为零，当且仅当 $\varphi$ 存在到 $\mathcal{O}_{X'}$-线性同态
$\varphi':\mathcal{F}'\to\mathcal{G}'$ 的提升。
\end{lemma}

\begin{proof}
由引理 \ref{defos-lemma-inf-map-special} 的证明可见，$\varphi$ 经同态
$$
\Hom_{\mathcal{O}_X}(\mathcal{F}, \mathcal{G}) =
\Hom_{\mathcal{O}_{X'}}(\mathcal{F}', \mathcal{G}) \longrightarrow
\Ext^1_{\mathcal{O}_{X'}}(\mathcal{F}', \mathcal{I}\mathcal{G}')
$$
所得的边界为零，当且仅当提升存在。又有
$$
\Ext^1_{\mathcal{O}_{X'}}(\mathcal{F}', \mathcal{I}\mathcal{G}') =
\Ext^1_{\mathcal{O}_X}(Li^*\mathcal{F}', \mathcal{I}\mathcal{G}')
$$
这是导出范畴上 $i_*=Ri_*$ 与 $Li^*$ 的伴随性（《层的上同调》引理
\ref{cohomology-lemma-adjoint}）。因此结论成立。
\end{proof}

\begin{lemma}
\label{defos-lemma-inf-obs-map-rel}
设 $(f,f')$ 是情形 \ref{defos-situation-morphism-thickenings} 中一阶增厚的
态射。设 $\mathcal{F}'$、$\mathcal{G}'$ 为
$\mathcal{O}_{X'}$-模，并令
$\mathcal{F}=i^*\mathcal{F}'$、$\mathcal{G}=i^*\mathcal{G}'$。
设 $\varphi:\mathcal{F}\to\mathcal{G}$ 为 $\mathcal{O}_X$-线性同态。
假设 $\mathcal{F}'$ 与 $\mathcal{G}'$ 均在 $S'$ 上平坦，且
$(f,f')$ 是增厚的严格态射。存在元素
$$
o(\varphi) \in  \Ext^1_{\mathcal{O}_X}(\mathcal{F},
\mathcal{G} \otimes_{\mathcal{O}_X} f^*\mathcal{J})
$$
它为零，当且仅当 $\varphi$ 存在到 $\mathcal{O}_{X'}$-线性同态
$\varphi':\mathcal{F}'\to\mathcal{G}'$ 的提升。
\end{lemma}

\begin{proof}[第一种证明]
这由引理 \ref{defos-lemma-inf-obs-map-special} 推出，因为我们断言在本引理的
假设下有
$$
\Ext^1_{\mathcal{O}_X}(Li^*\mathcal{F}',
\mathcal{I}\mathcal{G}') =
\Ext^1_{\mathcal{O}_X}(\mathcal{F},
\mathcal{G} \otimes_{\mathcal{O}_X} f^*\mathcal{J})
$$
事实上，由引理 \ref{defos-lemma-deform-module} 有
$\mathcal{I}\mathcal{G}' =
\mathcal{G} \otimes_{\mathcal{O}_X} f^*\mathcal{J}$
。另一方面，注意
$$
H^{-1}(Li^*\mathcal{F}') =
\text{Tor}_1^{\mathcal{O}_{X'}}(\mathcal{F}', \mathcal{O}_X)
$$
（局部计算从略）。利用短正合列
$$
0 \to \mathcal{I} \to \mathcal{O}_{X'} \to \mathcal{O}_X \to 0
$$
可见，该 $\text{Tor}_1$ 由同态
$\mathcal{I}\otimes_{\mathcal{O}_X}\mathcal{F}
\to\mathcal{I}\mathcal{F}'$ 的核计算；由引理
\ref{defos-lemma-deform-module} 的最后一个断言，该核为零。因此
$\tau_{\geq-1}Li^*\mathcal{F}'=\mathcal{F}$。另一方面，有
$$
\Ext^1_{\mathcal{O}_X}(Li^*\mathcal{F}',
\mathcal{I}\mathcal{G}') =
\Ext^1_{\mathcal{O}_X}(\tau_{\geq -1}Li^*\mathcal{F}',
\mathcal{I}\mathcal{G}')
$$
这是《导出范畴》引理 \ref{derived-lemma-negative-vanishing} 的对偶。
\end{proof}

\begin{proof}[第二种证明]
可以如下应用引理 \ref{defos-lemma-inf-obs-map}。注意，由引理
\ref{defos-lemma-deform-module} 有
$\mathcal{K}=\mathcal{I}\otimes_{\mathcal{O}_X}\mathcal{F}$、
$\mathcal{L}=\mathcal{I}\otimes_{\mathcal{O}_X}\mathcal{G}$
，并且 $c_{\mathcal{F}'}=1\otimes1$、
$c_{\mathcal{G}'}=1\otimes1$。取 $\psi=1\otimes\varphi$，引理中的
图表交换。因此可以取
$o(\varphi)=o(\varphi,1\otimes\varphi)$。
\end{proof}

\begin{lemma}
\label{defos-lemma-inf-ext-rel}
设 $(f,f')$ 是情形 \ref{defos-situation-morphism-thickenings} 中一阶增厚的
态射，$\mathcal{F}$ 为 $\mathcal{O}_X$-模。假设 $(f,f')$ 是增厚的
严格态射，且 $\mathcal{F}$ 在 $S$ 上平坦。如果存在二元组
$(\mathcal{F}',\alpha)$，其中 $\mathcal{F}'$ 是在 $S'$ 上平坦的
$\mathcal{O}_{X'}$-模，而
$\alpha:i^*\mathcal{F}'\to\mathcal{F}$ 是同构，那么这类二元组的
同构类集合是下列群作用下的主齐性空间：
$\Ext^1_{\mathcal{O}_X}(
\mathcal{F}, \mathcal{I} \otimes_{\mathcal{O}_X} \mathcal{F})$.
\end{lemma}

\begin{proof}
假设存在一个这样的模。那么由引理 \ref{defos-lemma-deform-module}，典范同态
$$
f^*\mathcal{J} \otimes_{\mathcal{O}_X} \mathcal{F} \to
\mathcal{I} \otimes_{\mathcal{O}_X} \mathcal{F}
$$
是同构。应用引理 \ref{defos-lemma-inf-ext}，取 $\mathcal{K}=
\mathcal{I} \otimes_{\mathcal{O}_X} \mathcal{F}$
、$c=1$。由引理 \ref{defos-lemma-deform-module}，相应的扩张
$\mathcal{F}'$ 均在 $S'$ 上平坦。
\end{proof}

\begin{lemma}
\label{defos-lemma-inf-obs-ext-rel}
设 $(f,f')$ 是情形 \ref{defos-situation-morphism-thickenings} 中一阶增厚的
态射，$\mathcal{F}$ 为 $\mathcal{O}_X$-模。假设 $(f,f')$ 是增厚的
严格态射，且 $\mathcal{F}$ 在 $S$ 上平坦。那么，存在在 $S'$ 上平坦的
$\mathcal{O}_{X'}$-模 $\mathcal{F}'$，满足
$i^*\mathcal{F}'\cong\mathcal{F}$，当且仅当：
\begin{enumerate}
\item 典范同态
$f^*\mathcal{J} \otimes_{\mathcal{O}_X} \mathcal{F} \to
\mathcal{I} \otimes_{\mathcal{O}_X} \mathcal{F}$
是同构；
\item 引理 \ref{defos-lemma-inf-obs-ext} 中的类
$o(\mathcal{F}, \mathcal{I} \otimes_{\mathcal{O}_X} \mathcal{F}, 1)
\in \Ext^2_{\mathcal{O}_X}(
\mathcal{F}, \mathcal{I} \otimes_{\mathcal{O}_X} \mathcal{F})$
为零。
\end{enumerate}
\end{lemma}

\begin{proof}
这立即由引理 \ref{defos-lemma-deform-module} 对在 $S'$ 上平坦的
$\mathcal{O}_{X'}$-模的刻画以及引理 \ref{defos-lemma-inf-obs-ext} 推出。
\end{proof}






\section{在环化空间平坦增厚上的平坦模中的应用}
\label{defos-section-flat}

\noindent
考虑交换图
$$
\xymatrix{
(X, \mathcal{O}_X) \ar[r]_i \ar[d]_f & (X', \mathcal{O}_{X'}) \ar[d]^{f'} \\
(S, \mathcal{O}_S) \ar[r]^t & (S', \mathcal{O}_{S'})
}
$$
其中各项为环化空间，水平箭头是情形
\ref{defos-situation-morphism-thickenings} 中的一阶增厚。令
$\mathcal{I}=\Ker(i^\sharp)\subset\mathcal{O}_{X'}$、
$\mathcal{J}=\Ker(t^\sharp)\subset\mathcal{O}_{S'}$。设
$\mathcal{F}$ 为 $\mathcal{O}_X$-模。假设：
\begin{enumerate}
\item $(f,f')$ 是增厚的严格态射；
\item $f'$ 平坦；
\item $\mathcal{F}$ 在 $S$ 上平坦。
\end{enumerate}
注意，(1) $+$ (2) 蕴含 $\mathcal{I}=f^*\mathcal{J}$（把引理
\ref{defos-lemma-deform-module} 应用于 $\mathcal{O}_{X'}$）。在这些假设下，
上一节的理论具有特别简洁的形式。下述引理汇总已经得到的结果。

\begin{lemma}
\label{defos-lemma-flat}
在上述情形中：
\begin{enumerate}
\item 存在在 $S'$ 上平坦的 $\mathcal{O}_{X'}$-模 $\mathcal{F}'$，
满足 $i^*\mathcal{F}'\cong\mathcal{F}$，当且仅当引理
\ref{defos-lemma-inf-obs-ext} 中的类
$o(\mathcal{F}, f^*\mathcal{J} \otimes_{\mathcal{O}_X} \mathcal{F}, 1)
\in \Ext^2_{\mathcal{O}_X}(
\mathcal{F}, f^*\mathcal{J} \otimes_{\mathcal{O}_X} \mathcal{F})$
为零。
\item 如果这样的模存在，那么提升的同构类集合是
$\Ext^1_{\mathcal{O}_X}(
\mathcal{F}, f^*\mathcal{J} \otimes_{\mathcal{O}_X} \mathcal{F})$
作用下的主齐性空间。
\item 给定提升 $\mathcal{F}'$，其拉回为
$\text{id}_{\mathcal{F}}$ 的 $\mathcal{F}'$ 自同构所成的集合，典范同构于
$\Ext^0_{\mathcal{O}_X}(
\mathcal{F}, f^*\mathcal{J} \otimes_{\mathcal{O}_X} \mathcal{F})$.
\end{enumerate}
\end{lemma}

\begin{proof}
由于上面已说明 $\mathcal{I}=f^*\mathcal{J}$，第 (1) 点由引理
\ref{defos-lemma-inf-obs-ext-rel} 推出。第 (2) 点由引理
\ref{defos-lemma-inf-ext-rel} 推出，第 (3) 点由引理
\ref{defos-lemma-inf-map-rel} 推出。
\end{proof}

\begin{situation}
\label{defos-situation-ses-flat-thickenings}
设 $f:(X,\mathcal{O}_X)\to(S,\mathcal{O}_S)$ 为环化空间的态射。
考虑交换图
$$
\xymatrix{
(X'_1, \mathcal{O}'_1) \ar[r]_h \ar[d]_{f'_1} &
(X'_2, \mathcal{O}'_2) \ar[r] \ar[d]_{f'_2} &
(X'_3, \mathcal{O}'_3) \ar[d]_{f'_3} \\
(S'_1, \mathcal{O}_{S'_1}) \ar[r] &
(S'_2, \mathcal{O}_{S'_2}) \ar[r] &
(S'_3, \mathcal{O}_{S'_3})
}
$$
其中：(a) 上行是 $X$ 的一阶增厚的短正合列；(b) 下行是 $S$ 的一阶
增厚的短正合列；(c) 每个 $f'_i$ 限制为 $f$；(d) 每个二元组
$(f,f'_i)$ 都是增厚的严格态射；(e) 每个 $f'_i$ 都平坦。最后，设
$\mathcal{F}'_2$ 为在 $S'_2$ 上平坦的 $\mathcal{O}'_2$-模，并令
$\mathcal{F}=\mathcal{F}'_2|_X$。设 $\pi:X'_1\to X$ 为典范分裂
（评注 \ref{defos-remark-short-exact-sequence-thickenings}）。
\end{situation}

\begin{lemma}
\label{defos-lemma-verify-iv}
在情形 \ref{defos-situation-ses-flat-thickenings} 中，模
$\pi^*\mathcal{F}$ 与 $h^*\mathcal{F}'_2$ 都是在 $S'_1$ 上平坦的
$\mathcal{O}'_1$-模，且在 $X$ 上限制为 $\mathcal{F}$。二者之差
（引理 \ref{defos-lemma-flat}）是
$\Ext^1_{\mathcal{O}_X}(
\mathcal{F}, f^*\mathcal{J}_1 \otimes_{\mathcal{O}_X} \mathcal{F})$
中的元素 $\theta$，它在
$\Ext^2_{\mathcal{O}_X}(
\mathcal{F}, f^*\mathcal{J}_3 \otimes_{\mathcal{O}_X} \mathcal{F})$
中的边界，等于把 $\mathcal{F}$ 提升为在 $S'_3$ 上平坦的
$\mathcal{O}'_3$-模的障碍（引理 \ref{defos-lemma-flat}）。
\end{lemma}

\begin{proof}
注意，$\pi^*\mathcal{F}$ 与 $h^*\mathcal{F}'_2$ 在 $X$ 上都限制为
$\mathcal{F}$，且同态 $\pi^*\mathcal{F}\to\mathcal{F}$ 与
$h^*\mathcal{F}'_2\to\mathcal{F}$ 的核均为
$f^*\mathcal{J}_1\otimes_{\mathcal{O}_X}\mathcal{F}$。因此，由引理
\ref{defos-lemma-deform-module} 得到平坦性。取边界是有意义的，因为模列
$$
0 \to f^*\mathcal{J}_3 \otimes_{\mathcal{O}_X} \mathcal{F} \to
f^*\mathcal{J}_2 \otimes_{\mathcal{O}_X} \mathcal{F} \to
f^*\mathcal{J}_1 \otimes_{\mathcal{O}_X} \mathcal{F} \to 0
$$
是短正合的；这是情形 \ref{defos-situation-ses-flat-thickenings} 中的假设以及
$\mathcal{F}$ 在 $S$ 上平坦的结果。关于障碍类的陈述，就是把评注
\ref{defos-remark-complex-thickenings-and-ses-modules} 的结果直接应用于
这一特殊情形。
\end{proof}






\section{环化空间的形变与朴素余切复形}
\label{defos-section-deformations-ringed-spaces}

\noindent
本节利用朴素余切复形来处理一些形变理论问题。先取环化空间的一阶增厚
$t:(S,\mathcal{O}_S)\to(S',\mathcal{O}_{S'})$。记
$\mathcal{J}=\Ker(t^\sharp)$，并等同 $S$ 与 $S'$ 的底层拓扑空间。
此外，设已给定环化空间的态射
$f:(X,\mathcal{O}_X)\to(S,\mathcal{O}_S)$、$\mathcal{O}_X$-模
$\mathcal{G}$ 以及模层的 $f$-映射 $c:\mathcal{J}\to\mathcal{G}$
（《空间上的层》定义 \ref{sheaves-definition-f-map} 及第
\ref{sheaves-section-ringed-spaces-functoriality-modules} 节）。
我们要问，能否找到一个对象来替代下图中的问号，使图表交换：
\begin{equation}
\label{defos-equation-to-solve-ringed-spaces}
\vcenter{
\xymatrix{
0 \ar[r] & \mathcal{G} \ar[r] & {?} \ar[r] & \mathcal{O}_X \ar[r] & 0 \\
0 \ar[r] & \mathcal{J} \ar[u]^c \ar[r] & \mathcal{O}_{S'} \ar[u] \ar[r] &
\mathcal{O}_S \ar[u] \ar[r] & 0
}
}
\end{equation}
（其中竖直箭头为 $f$-映射），并且，如果解存在，它在何种意义下是唯一的。
更确切地说，我们寻找一个一阶增厚
$i:(X,\mathcal{O}_X)\to(X',\mathcal{O}_{X'})$ 以及
(\ref{defos-equation-morphism-thickenings}) 中的增厚态射 $(f,f')$，其中
$\Ker(i^\sharp)$ 与 $\mathcal{G}$ 等同，且 $(f')^\sharp$ 诱导给定映射
$c$。我们称 $X'$ 为 (\ref{defos-equation-to-solve-ringed-spaces}) 的一个
{\it 解}。

\begin{lemma}
\label{defos-lemma-huge-diagram-ringed-spaces}
假设给定环化空间态射的交换图
\begin{equation}
\label{defos-equation-huge-1}
\vcenter{
\xymatrix{
& (X_2, \mathcal{O}_{X_2}) \ar[r]_{i_2} \ar[d]_{f_2} \ar[ddl]_g &
(X'_2, \mathcal{O}_{X'_2}) \ar[d]^{f'_2} \\
& (S_2, \mathcal{O}_{S_2}) \ar[r]^{t_2} \ar[ddl]|\hole &
(S'_2, \mathcal{O}_{S'_2}) \ar[ddl] \\
(X_1, \mathcal{O}_{X_1}) \ar[r]_{i_1} \ar[d]_{f_1} &
(X'_1, \mathcal{O}_{X'_1}) \ar[d]^{f'_1} \\
(S_1, \mathcal{O}_{S_1}) \ar[r]^{t_1} &
(S'_1, \mathcal{O}_{S'_1})
}
}
\end{equation}
其中水平箭头为一阶增厚。令
$\mathcal{G}_j=\Ker(i_j^\sharp)$，并假设给定模的 $g$-映射
$\nu:\mathcal{G}_1\to\mathcal{G}_2$，从而得到交换图
\begin{equation}
\label{defos-equation-huge-2}
\vcenter{
\xymatrix{
& 0 \ar[r] & \mathcal{G}_2 \ar[r] &
\mathcal{O}_{X'_2} \ar[r] &
\mathcal{O}_{X_2} \ar[r] & 0 \\
& 0 \ar[r]|\hole &
\mathcal{J}_2 \ar[u]_{c_2} \ar[r] &
\mathcal{O}_{S'_2} \ar[u] \ar[r]|\hole &
\mathcal{O}_{S_2} \ar[u] \ar[r] & 0 \\
0 \ar[r] & \mathcal{G}_1 \ar[ruu] \ar[r] &
\mathcal{O}_{X'_1} \ar[r] &
\mathcal{O}_{X_1} \ar[ruu] \ar[r] & 0 \\
0 \ar[r] & \mathcal{J}_1 \ar[ruu]|\hole \ar[u]^{c_1} \ar[r] &
\mathcal{O}_{S'_1} \ar[ruu]|\hole \ar[u] \ar[r] &
\mathcal{O}_{S_1} \ar[ruu]|\hole \ar[u] \ar[r] & 0
}
}
\end{equation}
其中前后两面均为 (\ref{defos-equation-to-solve-ringed-spaces}) 的解。
\begin{enumerate}
\item 在
$\Ext^1_{\mathcal{O}_{X_2}}(Lg^*\NL_{X_1/S_1}, \mathcal{G}_2)$
中存在一个典范元；该元为零，当且仅当存在与 $\nu$ 相容并可置入
(\ref{defos-equation-huge-1}) 的环化空间态射 $X'_2\to X'_1$。
\item 如果存在与 $\nu$ 相容并可置入 (\ref{defos-equation-huge-1}) 的态射
$X'_2\to X'_1$，那么所有这类态射组成的集合是下列群作用下的
主齐性空间：
$$
\Hom_{\mathcal{O}_{X_1}}(\Omega_{X_1/S_1}, g_*\mathcal{G}_2) =
\Hom_{\mathcal{O}_{X_2}}(g^*\Omega_{X_1/S_1}, \mathcal{G}_2) =
\Ext^0_{\mathcal{O}_{X_2}}(Lg^*\NL_{X_1/S_1}, \mathcal{G}_2).
$$
\end{enumerate}
\end{lemma}

\begin{proof}
朴素余切复形 $\NL_{X_1/S_1}$ 定义于《模层》定义
\ref{modules-definition-cotangent-complex-morphism-ringed-topoi}。
引理最后一个陈述中的等号来自以下事实：$g^*$ 左伴随于 $g_*$；
$H^0(\NL_{X_1/S_1})=\Omega_{X_1/S_1}$（由朴素余切复形的构造）；
$Lg^*$ 是 $g^*$ 的左导出函子。因此，在余下的证明中，我们使用群
$\Ext^k_{\mathcal{O}_{X_2}}(Lg^*\NL_{X_1/S_1}, \mathcal{G}_2)$,
其中 $k=0,1$。先说明，可以约化到本引理中所有环化空间的底层拓扑空间
都相同的情形。

\medskip\noindent
为此，注意 $g^{-1}\NL_{X_1/S_1}$ 等于环层同态
$g^{-1}f_1^{-1}\mathcal{O}_{S_1}\to g^{-1}\mathcal{O}_{X_1}$ 的
朴素余切复形，见《模层》引理 \ref{modules-lemma-pullback-NL}。
此外，$\NL_{X_1/S_1}$ 的 $0$ 次项是平坦
$\mathcal{O}_{X_1}$-模，因而典范同态
$$
Lg^*\NL_{X_1/S_1}
\longrightarrow
g^{-1}\NL_{X_1/S_1} \otimes_{g^{-1}\mathcal{O}_{X_1}} \mathcal{O}_{X_2}
$$
在 $0$ 次和 $-1$ 次上同调层上诱导同构。因此，可以把引理中的 Ext 群
替换为
$$
\Ext^k_{g^{-1}\mathcal{O}_{X_1}}(g^{-1}\NL_{X_1/S_1}, \mathcal{G}_2) =
\Ext^k_{g^{-1}\mathcal{O}_{X_1}}(
\NL_{g^{-1}\mathcal{O}_{X_1}/g^{-1}f_1^{-1}\mathcal{O}_{S_1}}, \mathcal{G}_2)
$$
与 $\nu$ 相容并可置入 (\ref{defos-equation-huge-1}) 的环化空间态射
$X'_2\to X'_1$ 所成的集合，与下列同态所成的集合一一对应：与
$f^\sharp$ 和 $\nu$ 相容的
$g^{-1}f_1^{-1}\mathcal{O}_{S'_1}$-代数同态
$g^{-1}\mathcal{O}_{X'_1}\to\mathcal{O}_{X'_2}$。由此可见，可以假设
我们有 $X$ 上层的图表 (\ref{defos-equation-huge-2})，并寻找可置入该图的
环层同态 $\mathcal{O}_{X'_1}\to\mathcal{O}_{X'_2}$。

\medskip\noindent
在本引理余下的证明中，假设所有底层拓扑空间都相同。也就是说，我们有
空间 $X$ 上层的图表 (\ref{defos-equation-huge-2})，并寻找可置入其中的环层
同态 $\mathcal{O}_{X'_1}\to\mathcal{O}_{X'_2}$。所用的 Ext 群为
$\Ext^k_{\mathcal{O}_{X_1}}(
\NL_{\mathcal{O}_{X_1}/\mathcal{O}_{S_1}}, \mathcal{G}_2)$, $k = 0, 1$.

\medskip\noindent
第 1 步：构造障碍类。考虑集合层
$$
\mathcal{E} = \mathcal{O}_{X'_1} \times_{\mathcal{O}_{X_2}} \mathcal{O}_{X'_2}
$$
它带有满映射 $\alpha:\mathcal{E}\to\mathcal{O}_{X_1}$，故可用
$\NL(\alpha)$ 代替
$\NL_{\mathcal{O}_{X_1}/\mathcal{O}_{S_1}}$，见《模层》引理
\ref{modules-lemma-NL-up-to-qis}。令
$$
\mathcal{I}' =
\Ker(\mathcal{O}_{S'_1}[\mathcal{E}] \to \mathcal{O}_{X_1})
\quad\text{且}\quad
\mathcal{I} =
\Ker(\mathcal{O}_{S_1}[\mathcal{E}] \to \mathcal{O}_{X_1})
$$
存在满同态 $\mathcal{I}'\to\mathcal{I}$，其核为
$\mathcal{J}_1\mathcal{O}_{S'_1}[\mathcal{E}]$。我们得到两个
$\mathcal{O}_{S'_1}$-代数同态
$$
a : \mathcal{O}_{S'_1}[\mathcal{E}] \to \mathcal{O}_{X'_1}
\quad\text{且}\quad
b : \mathcal{O}_{S'_1}[\mathcal{E}] \to \mathcal{O}_{X'_2}
$$
它们分别诱导同态
$a|_{\mathcal{I}'}:\mathcal{I}'\to\mathcal{G}_1$ 与
$b|_{\mathcal{I}'}:\mathcal{I}'\to\mathcal{G}_2$。$a$ 与 $b$ 均在
$(\mathcal{I}')^2$ 上为零。此外，由于 (\ref{defos-equation-huge-2}) 的左方块
交换，把 $a$ 与 $b$ 视为取值于 $\mathcal{G}_2$ 的同态时，它们在
$\mathcal{J}_1\mathcal{O}_{S'_1}[\mathcal{E}]$ 上一致。因此，差
$b|_{\mathcal{I}'} - \nu \circ a|_{\mathcal{I}'}$
诱导良定义的 $\mathcal{O}_{X_1}$-线性同态
$$
\xi : \mathcal{I}/\mathcal{I}^2 \longrightarrow \mathcal{G}_2
$$
它把 $\mathcal{I}$ 的局部截面 $f$ 的类映到
$\nu(a(f'))-b(f')$，其中 $f'$ 是 $f$ 到 $\mathcal{I}'$ 的局部截面的
一个提升。令
% P11-E0002: frozen source sign convention retained; see control/ERRATA_CANDIDATES.jsonl.
$[\xi] \in \Ext^1_{\mathcal{O}_{X_1}}(\NL(\alpha), \mathcal{G}_2)$
为其像（见下文）。

\medskip\noindent
第 2 步：$[\xi]$ 为零是必要条件。记
$\Omega = \Omega_{\mathcal{O}_{S_1}[\mathcal{E}]/\mathcal{O}_{S_1}}
\otimes_{\mathcal{O}_{S_1}[\mathcal{E}]} \mathcal{O}_{X_1}$.
注意，$\NL(\alpha)=(\mathcal{I}/\mathcal{I}^2\to\Omega)$ 可置于可区别三角
$$
\Omega[0] \to
\NL(\alpha) \to
\mathcal{I}/\mathcal{I}^2[1] \to
\Omega[1]
$$
因此可见，$[\xi]$ 为零，当且仅当对某个同态
$\Omega\to\mathcal{G}_2$，$\xi$ 是复合
$\mathcal{I}/\mathcal{I}^2\to\Omega\to\mathcal{G}_2$。假设存在可置入
(\ref{defos-equation-huge-2}) 的环层同态
$\varphi:\mathcal{O}_{X'_1}\to\mathcal{O}_{X'_2}$。此时考虑同态
$\mathcal{O}_{S'_1}[\mathcal{E}] \to \mathcal{G}_2$,
$f'\mapsto b(f')-\varphi(a(f'))$。计算表明，该同态在
$\mathcal{J}_1\mathcal{O}_{S'_1}[\mathcal{E}]$ 上为零，并诱导导子
$\mathcal{O}_{S_1}[\mathcal{E}]\to\mathcal{G}_2$。所得的线性同态
$\Omega\to\mathcal{G}_2$ 说明此时 $[\xi]=0$。

\medskip\noindent
第 3 步：$[\xi]$ 为零是充分条件。设
$\theta:\Omega\to\mathcal{G}_2$ 为 $\mathcal{O}_{X_1}$-线性同态，
使 $\xi=\theta\circ(\mathcal{I}/\mathcal{I}^2\to\Omega)$。
计算表明
$$
b + \theta \circ d : \mathcal{O}_{S'_1}[\mathcal{E}] \to \mathcal{O}_{X'_2}
$$
在 $\mathcal{I}'$ 上的限制与
$\nu\circ a:\mathcal{I}'\to\mathcal{G}_2$ 一致。由于
$\mathcal{O}_{X'_1}$ 是 $\mathcal{I}'\to
\mathcal{O}_{S'_1}[\mathcal{E}]$ 与
$\mathcal{I}'\to\mathcal{G}_1$ 的推出，同态
$b+\theta\circ d$ 与 $a$ 定义一个可置入
(\ref{defos-equation-huge-2}) 的同态
$\mathcal{O}_{X'_1}\to\mathcal{O}_{X'_2}$。

\medskip\noindent
上述特殊情形下第 (2) 点的证明从略。提示：这与引理
\ref{defos-lemma-huge-diagram} 第 (2) 点的证明完全相同。
\end{proof}

\begin{lemma}
\label{defos-lemma-NL-represent-ext-class}
设 $X$ 为拓扑空间，$\mathcal{A}\to\mathcal{B}$ 为环层同态，
$\mathcal{G}$ 为 $\mathcal{B}$-模。设
$\xi\in\Ext^1_\mathcal{B}(\NL_{\mathcal{B}/\mathcal{A}},\mathcal{G})$。
存在集合层的映射 $\alpha:\mathcal{E}\to\mathcal{B}$，使得把
$\xi\in\Ext^1_\mathcal{B}(\NL(\alpha),\mathcal{G})$ 表示为同态
$\mathcal{I}/\mathcal{I}^2\to\mathcal{G}$ 的类（记号见证明）。
\end{lemma}

\begin{proof}
回顾：给定 $\alpha:\mathcal{E}\to\mathcal{B}$，若
$\mathcal{A}[\mathcal{E}]\to\mathcal{B}$ 是满的且核为
$\mathcal{I}$，则复形
$\NL(\alpha) = (\mathcal{I}/\mathcal{I}^2 \to 
\Omega_{\mathcal{A}[\mathcal{E}]/\mathcal{A}}
\otimes_{\mathcal{A}[\mathcal{E}]} \mathcal{B})$
典范同构于 $\NL_{\mathcal{B}/\mathcal{A}}$，见《模层》引理
\ref{modules-lemma-NL-up-to-qis}。此外，注意
$\Omega = \Omega_{\mathcal{A}[\mathcal{E}]/\mathcal{A}}
\otimes_{\mathcal{A}[\mathcal{E}]} \mathcal{B}$
是与预层
$U \mapsto \bigoplus_{e \in \mathcal{E}(U)} \mathcal{B}(U)$.
相伴的层。换言之，$\Omega$ 是由集合层 $\mathcal{E}$ 生成的自由
$\mathcal{B}$-模；特别地，存在典范映射 $\mathcal{E}\to\Omega$。

\medskip\noindent
据此，选取某个 $\mathcal{E}$（例如像朴素余切复形定义中那样取
$\mathcal{E}=\mathcal{B}$）。把 $\xi$ 写成某个同态
$\mathcal{I}/\mathcal{I}^2\to\mathcal{G}$ 的类之障碍，是
$\Ext^1_\mathcal{B}(\Omega,\mathcal{G})$ 中的元素。设该元素由
$\mathcal{B}$-模扩张
$0\to\mathcal{G}\to\mathcal{H}\to\Omega\to0$ 表示。考虑集合层
$\mathcal{E}' = \mathcal{E} \times_\Omega \mathcal{H}$
它带有诱导映射 $\alpha':\mathcal{E}'\to\mathcal{B}$。令
$\mathcal{I}'=\Ker(\mathcal{A}[\mathcal{E}']\to\mathcal{B})$，
并令 $\Omega'=\Omega_{\mathcal{A}[\mathcal{E}']/\mathcal{A}}
\otimes_{\mathcal{A}[\mathcal{E}']} \mathcal{B}$.
在拟同构 $\NL(\alpha')\to\NL(\alpha)$ 下拉回 $\xi$，其在
$\Ext^1_\mathcal{B}(\Omega',\mathcal{G})$ 中的像为零。事实上，沿
$\Omega'\to\Omega$ 拉回扩张 $\mathcal{H}$ 所得的扩张是分裂的：
$\Omega'$ 是由集合层 $\mathcal{E}'$ 生成的自由 $\mathcal{B}$-模，
且按构造有交换图
$$
\xymatrix{
\mathcal{E}' \ar[r] \ar[d] & \mathcal{E} \ar[d] \\
\mathcal{H} \ar[r] & \Omega
}
$$
证明完成。
\end{proof}

\begin{lemma}
\label{defos-lemma-choices-ringed-spaces}
如果 (\ref{defos-equation-to-solve-ringed-spaces}) 存在解，那么解的同构类
集合是 $\Ext^1_{\mathcal{O}_X}(\NL_{X/S},\mathcal{G})$ 作用下的
主齐性空间。
\end{lemma}

\begin{proof}
先注意，给定 (\ref{defos-equation-to-solve-ringed-spaces}) 的两个解
$X'_1$ 与 $X'_2$，由引理 \ref{defos-lemma-huge-diagram-ringed-spaces}
可得障碍元
$o(X'_1,X'_2)\in\Ext^1_{\mathcal{O}_X}(\NL_{X/S},\mathcal{G})$；
它阻碍态射 $X'_1\to X'_2$ 的存在。显然，该元也是同构存在性的障碍，
因而能够区分不同的同构类。因此，为完成证明，只需说明：给定一个解
$X'$ 和元素
$\xi\in\Ext^1_{\mathcal{O}_X}(\NL_{X/S},\mathcal{G})$，可以找到
另一个解 $X'_\xi$，使得 $o(X',X'_\xi)=\xi$。

\medskip\noindent
按照引理 \ref{defos-lemma-NL-represent-ext-class}，为类 $\xi$ 选取
$\alpha:\mathcal{E}\to\mathcal{O}_X$。考虑满同态
$f^{-1}\mathcal{O}_S[\mathcal{E}]\to\mathcal{O}_X$，记其核为
$\mathcal{I}$；相应的朴素余切复形为
$\NL(\alpha) = (\mathcal{I}/\mathcal{I}^2 \to
\Omega_{f^{-1}\mathcal{O}_S[\mathcal{E}]/f^{-1}\mathcal{O}_S}
\otimes_{f^{-1}\mathcal{O}_S[\mathcal{E}]} \mathcal{O}_X)$.
由该引理，$\xi$ 是某个同态
$\delta:\mathcal{I}/\mathcal{I}^2\to\mathcal{G}$ 的类。把
$\mathcal{E}$ 替换为
$\mathcal{E}\times_{\mathcal{O}_X}\mathcal{O}_{X'}$ 后，还可假设
$\alpha$ 经映射 $\alpha':\mathcal{E}\to\mathcal{O}_{X'}$ 分解。

\medskip\noindent
这些选择确定 $f^{-1}\mathcal{O}_{S'}$-代数同态
$\varphi:f^{-1}\mathcal{O}_{S'}[\mathcal{E}]\to\mathcal{O}_{X'}$。
令
$\mathcal{I}'=\Ker(f^{-1}\mathcal{O}_{S'}[\mathcal{E}]
\to\mathcal{O}_X)$。注意，$\varphi$ 诱导同态
$\varphi|_{\mathcal{I}'}:\mathcal{I}'\to\mathcal{G}$，且
$\mathcal{O}_{X'}$ 是下图所示的推出：
$$
\xymatrix{
0 \ar[r] & \mathcal{G} \ar[r] & \mathcal{O}_{X'} \ar[r] &
\mathcal{O}_X \ar[r] & 0 \\
0 \ar[r] & \mathcal{I}' \ar[u]^{\varphi|_{\mathcal{I}'}} \ar[r] &
f^{-1}\mathcal{O}_{S'}[\mathcal{E}] \ar[u] \ar[r] &
\mathcal{O}_X \ar[u]_{=} \ar[r] & 0
}
$$
令 $\psi:\mathcal{I}'\to\mathcal{G}$ 为
$\varphi|_{\mathcal{I}'}$ 与下列复合同态之和：
$$
\mathcal{I}' \to \mathcal{I}'/(\mathcal{I}')^2 \to
\mathcal{I}/\mathcal{I}^2 \xrightarrow{\delta} \mathcal{G}.
$$
于是，沿 $\psi$ 所作的推出给出另一个环扩张
$\mathcal{O}_{X'_\xi}$，它可置于上述形式的图表中。计算（从略）表明
$o(X',X'_\xi)=\xi$，正如所需。
\end{proof}

\begin{lemma}
\label{defos-lemma-extensions-of-relative-ringed-spaces}
设 $f:(X,\mathcal{O}_X)\to(S,\mathcal{O}_S)$ 为环化空间的态射，
$\mathcal{G}$ 为 $\mathcal{O}_X$-模。所有如下
$f^{-1}\mathcal{O}_S$-代数扩张的同构类所成的集合
$$
0 \to \mathcal{G} \to \mathcal{O}_{X'} \to \mathcal{O}_X \to 0
$$
（其中 $\mathcal{G}$ 是平方零理想）\footnote{换言之，即 $S$ 上的一阶
增厚 $i:X\to X'$ 连同 $\mathcal{O}_X$-模同构
$\mathcal{G}\to\Ker(i^\sharp)$ 的同构类集合。}
与 $\Ext^1_{\mathcal{O}_X}(\NL_{X/S},\mathcal{G})$ 典范双射。
\end{lemma}

\begin{proof}
为证明这一点，把前述结果应用于
(\ref{defos-equation-to-solve-ringed-spaces}) 由下图给出的情形：
$$
\xymatrix{
0 \ar[r] &
\mathcal{G} \ar[r] &
{?} \ar[r] &
\mathcal{O}_X \ar[r] & 0 \\
0 \ar[r] &
0 \ar[u] \ar[r] &
\mathcal{O}_S \ar[u] \ar[r]^{\text{id}} &
\mathcal{O}_S \ar[u] \ar[r] & 0
}
$$
于是，本引理由引理 \ref{defos-lemma-choices-ringed-spaces} 以及解的存在性
推出；一个解就是 $\mathcal{G}\oplus\mathcal{O}_X$。（该双射的直接
构造见下述评注。）
\end{proof}

\begin{remark}
\label{defos-remark-extensions-of-relative-ringed-spaces}
设 $f:(X,\mathcal{O}_X)\to(S,\mathcal{O}_S)$ 与 $\mathcal{G}$
如引理 \ref{defos-lemma-extensions-of-relative-ringed-spaces} 所述。
考虑引理中的扩张
$0 \to \mathcal{G} \to \mathcal{O}_{X'} \to \mathcal{O}_X \to 0$
。可以选取集合层 $\mathcal{E}$ 以及交换图
$$
\xymatrix{
\mathcal{E} \ar[d]_{\alpha'} \ar[rd]^\alpha \\
\mathcal{O}_{X'} \ar[r] & \mathcal{O}_X
}
$$
使 $f^{-1}\mathcal{O}_S[\mathcal{E}]\to\mathcal{O}_X$ 是满的，核为
$\mathcal{J}$。（例如，可以取任意满射到 $\mathcal{O}_{X'}$ 的
集合层。）于是
$$
\NL_{X/S} \cong \NL(\alpha) = 
\left(
\mathcal{J}/\mathcal{J}^2
\longrightarrow
\Omega_{f^{-1}\mathcal{O}_S[\mathcal{E}]/f^{-1}\mathcal{O}_S}
\otimes_{f^{-1}\mathcal{O}_S[\mathcal{E}]} \mathcal{O}_X\right)
$$
见《模层》第 \ref{modules-section-netherlander} 节，特别是引理
\ref{modules-lemma-NL-up-to-qis}。当然，$\alpha'$ 确定同态
$f^{-1}\mathcal{O}_S[\mathcal{E}]\to\mathcal{O}_{X'}$，后者又确定同态
$$
\mathcal{J}/\mathcal{J}^2 \longrightarrow \mathcal{G}
$$
后者进一步确定
$\Ext^1_{\mathcal{O}_X}(\NL(\alpha), \mathcal{G}) =
\Ext^1_{\mathcal{O}_X}(\NL_{X/S}, \mathcal{G})$
中经本引理双射与 $\mathcal{O}_{X'}$ 对应的元素。
\end{remark}

\begin{lemma}
\label{defos-lemma-extensions-of-relative-ringed-spaces-functorial}
设 $f:(X,\mathcal{O}_X)\to(S,\mathcal{O}_S)$ 与
$g:(Y,\mathcal{O}_Y)\to(X,\mathcal{O}_X)$ 为环化空间的态射。
设 $\mathcal{F}$ 为 $\mathcal{O}_X$-模，$\mathcal{G}$ 为
$\mathcal{O}_Y$-模，$c:\mathcal{F}\to\mathcal{G}$ 为 $g$-映射。
最后，考虑
\begin{enumerate}
\item[(a)] 与 $\xi\in\Ext^1_{\mathcal{O}_X}(\NL_{X/S},\mathcal{F})$
对应的 $f^{-1}\mathcal{O}_S$-代数扩张
$0\to\mathcal{F}\to\mathcal{O}_{X'}\to\mathcal{O}_X\to0$；
\item[(b)] 与 $\zeta\in\Ext^1_{\mathcal{O}_Y}(\NL_{Y/S},\mathcal{G})$
对应的 $g^{-1}f^{-1}\mathcal{O}_S$-代数扩张
$0\to\mathcal{G}\to\mathcal{O}_{Y'}\to\mathcal{O}_Y\to0$。
\end{enumerate}
参见引理 \ref{defos-lemma-extensions-of-relative-ringed-spaces}。那么，存在与
$g$ 和 $c$ 相容的 $S$-态射 $g':Y'\to X'$，当且仅当 $\xi$ 与
$\zeta$ 在
$\Ext^1_{\mathcal{O}_Y}(Lg^*\NL_{X/S},\mathcal{G})$ 中的像相同。
\end{lemma}

\begin{proof}
这个陈述是有意义的，因为我们有同态
$$
\Ext^1_{\mathcal{O}_X}(\NL_{X/S}, \mathcal{F}) \to
\Ext^1_{\mathcal{O}_Y}(Lg^*\NL_{X/S}, Lg^*\mathcal{F}) \to
\Ext^1_{\mathcal{O}_Y}(Lg^*\NL_{X/S}, \mathcal{G})
$$
它使用同态
$Lg^*\mathcal{F}\to g^*\mathcal{F}\xrightarrow{c}\mathcal{G}$；还有
$$
\Ext^1_{\mathcal{O}_Y}(\NL_{Y/S}, \mathcal{G}) \to
\Ext^1_{\mathcal{O}_Y}(Lg^*\NL_{X/S}, \mathcal{G})
$$
它使用同态 $Lg^*\NL_{X/S}\to\NL_{Y/S}$。把引理
\ref{defos-lemma-huge-diagram-ringed-spaces} 应用于下图，即可推出本引理：
$$
\xymatrix{
& 0 \ar[r] &
\mathcal{G} \ar[r] &
\mathcal{O}_{Y'} \ar[r] &
\mathcal{O}_Y \ar[r] & 0 \\
& 0 \ar[r]|\hole & 0 \ar[u] \ar[r] &
\mathcal{O}_S \ar[u] \ar[r]|\hole &
\mathcal{O}_S \ar[u] \ar[r] & 0 \\
0 \ar[r] &
\mathcal{F} \ar[ruu] \ar[r] &
\mathcal{O}_{X'} \ar[r] &
\mathcal{O}_X \ar[ruu] \ar[r] & 0 \\
0 \ar[r] & 0 \ar[ruu]|\hole \ar[u] \ar[r] &
\mathcal{O}_S \ar[ruu]|\hole \ar[u] \ar[r] &
\mathcal{O}_S \ar[ruu]|\hole \ar[u] \ar[r] & 0
}
$$
同时还要使用引理 \ref{defos-lemma-extensions-of-relative-ringed-spaces} 与
\ref{defos-lemma-huge-diagram-ringed-spaces} 的证明中各项构造之间的一种
相容性；其陈述和证明从略。（直接论证见下述评注。）
\end{proof}

\begin{remark}
\label{defos-remark-extensions-of-relative-ringed-spaces-functorial}
设 $f:(X,\mathcal{O}_X)\to(S,\mathcal{O}_S)$、
$g:(Y,\mathcal{O}_Y)\to(X,\mathcal{O}_X)$、$\mathcal{F}$、
$\mathcal{G}$、$c:\mathcal{F}\to\mathcal{G}$、
$0\to\mathcal{F}\to\mathcal{O}_{X'}\to\mathcal{O}_X\to0$、
$\xi\in\Ext^1_{\mathcal{O}_X}(\NL_{X/S},\mathcal{F})$、
$0\to\mathcal{G}\to\mathcal{O}_{Y'}\to\mathcal{O}_Y\to0$ 以及
$\zeta\in\Ext^1_{\mathcal{O}_Y}(\NL_{Y/S},\mathcal{G})$
如引理 \ref{defos-lemma-extensions-of-relative-ringed-spaces-functorial}
所述。沿 $c:g^{-1}\mathcal{F}\to\mathcal{G}$ 作推出，可构造扩张
$$
\xymatrix{
0 \ar[r] &
\mathcal{G} \ar[r] &
\mathcal{O}'_1 \ar[r] &
g^{-1}\mathcal{O}_X \ar[r] & 0 \\
0 \ar[r] &
g^{-1}\mathcal{F} \ar[u]^c \ar[r] &
g^{-1}\mathcal{O}_{X'} \ar[u] \ar[r] &
g^{-1}\mathcal{O}_X \ar@{=}[u] \ar[r] & 0
}
$$
沿 $g^\sharp:g^{-1}\mathcal{O}_X\to\mathcal{O}_Y$ 作拉回，可构造扩张
$$
\xymatrix{
0 \ar[r] &
\mathcal{G} \ar[r] &
\mathcal{O}_{Y'} \ar[r] &
\mathcal{O}_Y \ar[r] & 0 \\
0 \ar[r] &
\mathcal{G} \ar@{=}[u] \ar[r] &
\mathcal{O}'_2 \ar[u] \ar[r] &
g^{-1}\mathcal{O}_X \ar[u] \ar[r] & 0
}
$$
追图表明，存在与 $g$ 和 $c$ 相容的 $S$-态射 $Y'\to X'$，当且仅当
把 $\mathcal{O}'_1$ 与 $\mathcal{O}'_2$ 都看作
$g^{-1}\mathcal{O}_X$ 被 $\mathcal{G}$ 扩张所得的
$g^{-1}f^{-1}\mathcal{O}_S$-代数时二者同构。由引理
\ref{defos-lemma-extensions-of-relative-ringed-spaces}，这些扩张由下式左端分类：
$$
\Ext^1_{g^{-1}\mathcal{O}_X}(
\NL_{g^{-1}\mathcal{O}_X/g^{-1}f^{-1}\mathcal{O}_S}, \mathcal{G}) =
\Ext^1_{\mathcal{O}_Y}(Lg^*\NL_{X/S}, \mathcal{G})
$$
这里的等号来自张量--Hom 伴随性以及等式
$$
\NL_{g^{-1}\mathcal{O}_X/g^{-1}f^{-1}\mathcal{O}_S} = g^{-1}\NL_{X/S}
\quad\text{且}\quad
Lg^*\NL_{X/S} =
g^{-1}\NL_{X/S} \otimes_{g^{-1}\mathcal{O}_X}^\mathbf{L} \mathcal{O}_Y
$$
第一个等式见《模层》引理 \ref{modules-lemma-pullback-NL}；第二个等式
由导出拉回的定义推出。因此，为证明引理
\ref{defos-lemma-extensions-of-relative-ringed-spaces-functorial}，
只需说明 $\mathcal{O}'_1$ 对应于 $\xi$ 的像，而
$\mathcal{O}'_2$ 对应于 $\zeta$ 的像。$\xi$ 与
$\mathcal{O}'_1$ 之间的对应由评注
\ref{defos-remark-extensions-of-relative-ringed-spaces} 中类 $\xi$ 的构造
立即得到。为说明 $\zeta$ 与 $\mathcal{O}'_2$ 之间的对应，先选取交换图
$$
\xymatrix{
\mathcal{E} \ar[d]_{\beta'} \ar[rd]^\beta \\
\mathcal{O}_{Y'} \ar[r] & \mathcal{O}_Y
}
$$
使 $g^{-1}f^{-1}\mathcal{O}_S[\mathcal{E}]\to\mathcal{O}_Y$
是满的，核为 $\mathcal{K}$。再选取交换图
$$
\xymatrix{
\mathcal{E} \ar[d]_{\beta'} &
\mathcal{E}' \ar[l]^\varphi \ar[d]_{\alpha'} \ar[rd]^\alpha \\
\mathcal{O}_{Y'} &
\mathcal{O}'_2 \ar[l] \ar[r] &
g^{-1}\mathcal{O}_X
}
$$
使 $g^{-1}f^{-1}\mathcal{O}_S[\mathcal{E}']\to
g^{-1}\mathcal{O}_X$ 是满的，核为 $\mathcal{J}$。（例如，只需把
$\mathcal{E}'=\mathcal{E}\amalg\mathcal{O}'_2$ 视为集合层。）
映射 $\varphi$ 诱导复形同态 $\NL(\alpha)\to\NL(\beta)$
（记号同《模层》第 \ref{modules-section-netherlander} 节），特别地诱导
$\bar\varphi:\mathcal{J}/\mathcal{J}^2\to
\mathcal{K}/\mathcal{K}^2$。于是 $\NL(\alpha)\cong\NL_{Y/S}$，且
$\NL(\beta) \cong \NL_{g^{-1}\mathcal{O}_X/g^{-1}f^{-1}\mathcal{O}_S}$
，而复形同态 $\NL(\alpha)\to\NL(\beta)$ 表示引理
\ref{defos-lemma-extensions-of-relative-ringed-spaces-functorial} 的陈述中
所用的同态 $Lg^*\NL_{X/S}\to\NL_{Y/S}$（见其证明的第一部分）。
现在，$\zeta$ 对应于由 $\beta'$ 诱导的同态
$\mathcal{K}/\mathcal{K}^2\to\mathcal{G}$ 的类，见评注
\ref{defos-remark-extensions-of-relative-ringed-spaces}。类似地，扩张
$\mathcal{O}'_2$ 对应于由 $\alpha'$ 诱导的同态
$\mathcal{J}/\mathcal{J}^2\to\mathcal{G}$。上述交换图表明，该同态是
由 $\beta'$ 诱导的同态
$\mathcal{K}/\mathcal{K}^2\to\mathcal{G}$ 与
$\bar\varphi:\mathcal{J}/\mathcal{J}^2\to
\mathcal{K}/\mathcal{K}^2$ 的复合。这证明了所需的相容性。
\end{remark}

\begin{lemma}
\label{defos-lemma-parametrize-solutions-ringed-spaces}
设 $t:(S,\mathcal{O}_S)\to(S',\mathcal{O}_{S'})$、
$\mathcal{J}=\Ker(t^\sharp)$、
$f:(X,\mathcal{O}_X)\to(S,\mathcal{O}_S)$、$\mathcal{G}$ 以及
$c:\mathcal{J}\to\mathcal{G}$ 如
(\ref{defos-equation-to-solve-ringed-spaces}) 所述。记
$\xi\in\Ext^1_{\mathcal{O}_S}(\NL_{S/S'},\mathcal{J})$ 为经引理
\ref{defos-lemma-extensions-of-relative-ringed-spaces} 与
$\mathcal{O}_S$ 被 $\mathcal{J}$ 扩张所得的 $\mathcal{O}_{S'}$
相对应的元素。解的同构类集合与同态
$$
\Ext^1_{\mathcal{O}_X}(\NL_{X/S'}, \mathcal{G}) \to
\Ext^1_{\mathcal{O}_X}(Lf^*\NL_{S/S'}, \mathcal{G})
$$
在 $\xi$ 的像上方的纤维典范双射。
\end{lemma}

\begin{proof}
把引理 \ref{defos-lemma-extensions-of-relative-ringed-spaces} 应用于
$X\to S'$ 和 $\mathcal{O}_X$-模 $\mathcal{G}$，可见
$\Ext^1_{\mathcal{O}_X}(\NL_{X/S'},\mathcal{G})$ 中的元素 $\zeta$
参数化 $f^{-1}\mathcal{O}_{S'}$-代数扩张
$0 \to \mathcal{G} \to \mathcal{O}_{X'} \to \mathcal{O}_X \to 0$
。把引理 \ref{defos-lemma-extensions-of-relative-ringed-spaces-functorial}
应用于 $X\to S\to S'$ 和 $c:\mathcal{J}\to\mathcal{G}$，可见存在与
$c$ 和 $f:X\to S$ 相容的 $S'$-态射 $X'\to S'$，当且仅当
$\zeta$ 映到 $\xi$。这当然等价于说 $\mathcal{O}_{X'}$ 是
(\ref{defos-equation-to-solve-ringed-spaces}) 的一个解。
\end{proof}

\begin{remark}
\label{defos-remark-parametrize-solutions-ringed-spaces}
在引理 \ref{defos-lemma-parametrize-solutions-ringed-spaces} 的情形中，
有复形同态
$$
Lf^*\NL_{S'/S} \to \NL_{X/S'} \to \NL_{X/S}
$$
% P11-E0003: frozen source complex retained; see control/ERRATA_CANDIDATES.jsonl.
这些同态几乎构成一个可区别三角，见《模层》引理
\ref{modules-lemma-exact-sequence-NL-ringed-topoi}。如果它确实是可区别
三角，便可得出：$\xi$ 在
$\Ext^2_{\mathcal{O}_X}(\NL_{X/S},\mathcal{G})$ 中的像，就是
(\ref{defos-equation-to-solve-ringed-spaces}) 存在解的障碍。
\end{remark}













\section{概形的形变}
\label{defos-section-deformations-schemes}

\noindent
本节具体说明第 \ref{defos-section-deformations-ringed-spaces} 节的结果对于
概形形变意味着什么。

\begin{lemma}
\label{defos-lemma-deform}
设 $S\subset S'$ 为概形的一阶增厚，$f:X\to S$ 为概形的平坦态射。
如果存在概形的平坦态射 $f':X'\to S'$ 以及 $S$ 上的同构
$a:X\to X'\times_{S'}S$，那么：
\begin{enumerate}
\item 二元组 $(f':X'\to S',a)$ 的同构类集合是
$\Ext^1_{\mathcal{O}_X}(\NL_{X/S},f^*\mathcal{C}_{S/S'})$
作用下的主齐性空间；
\item $S'$ 上在 $X'\times_{S'}S$ 上约化为恒等态射的自同构
$\varphi:X'\to X'$ 所成的集合为
$\Ext^0_{\mathcal{O}_X}(\NL_{X/S},f^*\mathcal{C}_{S/S'})$。
\end{enumerate}
\end{lemma}

\begin{proof}
先注意，《更多态射》第 \ref{more-morphisms-section-thickenings} 节所定义的
概形增厚，正是第 \ref{defos-section-thickenings-spaces} 节意义下为增厚的
概形态射。可以把 $X$ 视为 $X'$ 的闭子概形，使得
$(f,f'):(X\subset X')\to(S\subset S')$ 是一阶增厚的态射。于是，由
《更多态射》引理 \ref{more-morphisms-lemma-deform}（或更一般的引理
\ref{defos-lemma-deform-module}），$X$ 在 $X'$ 中的理想层等于
$f^*\mathcal{C}_{S/S'}$。因此有交换图
$$
\xymatrix{
0 \ar[r] & f^*\mathcal{C}_{S/S'} \ar[r] &
\mathcal{O}_{X'} \ar[r] &
\mathcal{O}_X \ar[r] & 0 \\
0 \ar[r] & \mathcal{C}_{S/S'} \ar[u] \ar[r] &
\mathcal{O}_{S'} \ar[u] \ar[r] &
\mathcal{O}_S \ar[u] \ar[r] & 0
}
$$
其中竖直箭头为 $f$-映射；请与
(\ref{defos-equation-to-solve-ringed-spaces}) 比较。因此，第 (1) 点由引理
\ref{defos-lemma-choices-ringed-spaces} 推出，第 (2) 点由引理
\ref{defos-lemma-huge-diagram-ringed-spaces} 的第 (2) 点推出。（注意，
《更多态射》第 \ref{more-morphisms-section-netherlander} 节为概形态射
所定义的 $\NL_{X/S}$，与第 \ref{defos-section-deformations-ringed-spaces} 节
所用的 $\NL_{X/S}$ 一致。）
\end{proof}










\section{环化拓扑斯的增厚}
\label{defos-section-thickenings-ringed-topoi}

\noindent
本节是第 \ref{defos-section-thickenings-spaces} 节对于环化拓扑斯的对应版本。
在下面几节中，我们将使用以下概念：
\begin{enumerate}
\item 环化拓扑斯 $(\Sh(\mathcal{D}),\mathcal{O}')$ 上的理想层
$\mathcal{I}\subset\mathcal{O}'$ 称为{\it 局部幂零的}，如果
$\mathcal{I}$ 的任意局部截面都是局部幂零的。
\item 环化拓扑斯的一个{\it 增厚}是环化拓扑斯的态射
$i : (\Sh(\mathcal{C}), \mathcal{O}) \to (\Sh(\mathcal{D}), \mathcal{O}')$
，满足：
\begin{enumerate}
\item $i_*$ 是等价 $\Sh(\mathcal{C})\to\Sh(\mathcal{D})$；
\item 同态 $i^\sharp:\mathcal{O}'\to i_*\mathcal{O}$ 是满的；
\item $i^\sharp$ 的核是局部幂零理想层。
\end{enumerate}
\item 环化拓扑斯的一个{\it 一阶增厚}是环化拓扑斯的增厚
$i : (\Sh(\mathcal{C}), \mathcal{O}) \to (\Sh(\mathcal{D}), \mathcal{O}')$
，且 $\Ker(i^\sharp)$ 平方为零。
\item 显然可以定义{\it 环化拓扑斯增厚的态射}、
{\it 基环化拓扑斯上环化拓扑斯增厚的态射}等概念。
\end{enumerate}
若
$i : (\Sh(\mathcal{C}), \mathcal{O}) \to (\Sh(\mathcal{D}), \mathcal{O}')$
是环化拓扑斯的增厚，我们便等同底层拓扑斯，并把 $\mathcal{O}$、
$\mathcal{O}'$ 以及 $\mathcal{I}=\Ker(i^\sharp)$ 都视为
$\mathcal{C}$ 上的层。由此得到 $\mathcal{O}'$-模的短正合列
$$
0 \to \mathcal{I} \to \mathcal{O}' \to \mathcal{O} \to 0
$$
由《位点上的模》引理 \ref{sites-modules-lemma-i-star-equivalence}，
$\mathcal{O}$-模范畴等价于被 $\mathcal{I}$ 零化的
$\mathcal{O}'$-模范畴。特别地，如果 $i$ 是一阶增厚，那么
$\mathcal{I}$ 是 $\mathcal{O}$-模。

\begin{situation}
\label{defos-situation-morphism-thickenings-ringed-topoi}
环化拓扑斯增厚的一个态射 $(f,f')$ 由如下交换图给出：
\begin{equation}
\label{defos-equation-morphism-thickenings-ringed-topoi}
\vcenter{
\xymatrix{
(\Sh(\mathcal{C}), \mathcal{O}) \ar[r]_i \ar[d]_f &
(\Sh(\mathcal{D}), \mathcal{O}') \ar[d]^{f'} \\
(\Sh(\mathcal{B}), \mathcal{O}_\mathcal{B}) \ar[r]^t &
(\Sh(\mathcal{B}'), \mathcal{O}_{\mathcal{B}'})
}
}
\end{equation}
其中各项为环化拓扑斯，水平箭头为增厚。在这一情形中，令
$\mathcal{I}=\Ker(i^\sharp)\subset\mathcal{O}'$、
$\mathcal{J}=\Ker(t^\sharp)\subset\mathcal{O}_{\mathcal{B}'}$。
由于在底层拓扑斯上 $f=f'$，我们等同逆像函子 $f^{-1}$ 与
$(f')^{-1}$。注意，
$(f')^\sharp : f^{-1}\mathcal{O}_{\mathcal{B}'} \to \mathcal{O}'$
特别地诱导同态 $f^{-1}\mathcal{J}\to\mathcal{I}$，从而诱导
$\mathcal{O}'$-模同态
$$
(f')^*\mathcal{J} \longrightarrow \mathcal{I}
$$
如果 $i$ 与 $t$ 都是一阶增厚，那么
$(f')^*\mathcal{J}=f^*\mathcal{J}$，上述同态就成为
$f^*\mathcal{J}\to\mathcal{I}$。
\end{situation}

\begin{definition}
\label{defos-definition-strict-morphism-thickenings-ringed-topoi}
在情形 \ref{defos-situation-morphism-thickenings-ringed-topoi} 中，如果同态
$(f')^*\mathcal{J}\longrightarrow\mathcal{I}$ 是满的，则称
$(f,f')$ 是一个{\it 增厚的严格态射}。
\end{definition}








\section{环化拓扑斯一阶增厚上的模}
\label{defos-section-modules-thickenings-ringed-topoi}

\noindent
本节讨论模的形变理论所需的一些预备知识。设
% P11-E0004: malformed frozen source object retained; see control/ERRATA_CANDIDATES.jsonl.
$i : (\Sh(\mathcal{C}, \mathcal{O}) \to (\Sh(\mathcal{D}), \mathcal{O}')$
为环化拓扑斯的一阶增厚。我们将随时使用第
\ref{defos-section-thickenings-ringed-topoi} 节引入的记号；特别地，我们等同
底层拓扑斯。本节考虑 $\mathcal{O}'$-模的短正合列
\begin{equation}
\label{defos-equation-extension-ringed-topoi}
0 \to \mathcal{K} \to \mathcal{F}' \to \mathcal{F} \to 0
\end{equation}
其中 $\mathcal{F}$、$\mathcal{K}$ 是 $\mathcal{O}$-模，而
$\mathcal{F}'$ 是 $\mathcal{O}'$-模。在这种情形中，有典范的
$\mathcal{O}$-模同态
$$
c_{\mathcal{F}'} :
\mathcal{I} \otimes_\mathcal{O} \mathcal{F}
\longrightarrow
\mathcal{K}
$$
其中 $\mathcal{I}=\Ker(i^\sharp)$。具体而言，给定 $\mathcal{I}$ 的
局部截面 $f$ 与 $\mathcal{F}$ 的局部截面 $s$，令
$c_{\mathcal{F}'}(f\otimes s)=fs'$，其中 $s'$ 是提升 $s$ 的
$\mathcal{F}'$ 的局部截面。

\begin{lemma}
\label{defos-lemma-inf-map-ringed-topoi}
设
$i:(\Sh(\mathcal{C}),\mathcal{O})\to
(\Sh(\mathcal{D}),\mathcal{O}')$ 是环化拓扑斯的一阶增厚，并给定扩张
$$
0 \to \mathcal{K} \to \mathcal{F}' \to \mathcal{F} \to 0
\quad\text{且}\quad
0 \to \mathcal{L} \to \mathcal{G}' \to \mathcal{G} \to 0
$$
如 (\ref{defos-equation-extension-ringed-topoi}) 所述，以及同态
$\varphi:\mathcal{F}\to\mathcal{G}$ 与
$\psi:\mathcal{K}\to\mathcal{L}$。
\begin{enumerate}
\item 如果存在与 $\varphi$ 和 $\psi$ 相容的 $\mathcal{O}'$-模同态
$\varphi':\mathcal{F}'\to\mathcal{G}'$，那么图表
$$
\xymatrix{
\mathcal{I} \otimes_\mathcal{O} \mathcal{F}
\ar[r]_-{c_{\mathcal{F}'}} \ar[d]_{1 \otimes \varphi} &
\mathcal{K} \ar[d]^\psi \\
\mathcal{I} \otimes_\mathcal{O} \mathcal{G}
\ar[r]^-{c_{\mathcal{G}'}} &
\mathcal{L}
}
$$
交换。
\item 与 $\varphi$ 和 $\psi$ 相容的 $\mathcal{O}'$-模同态
$\varphi':\mathcal{F}'\to\mathcal{G}'$ 的集合，如果非空，就是
$\Hom_\mathcal{O}(\mathcal{F},\mathcal{L})$ 作用下的主齐性空间。
\end{enumerate}
\end{lemma}

\begin{proof}
第 (1) 点由各同态的描述立即得到。对于第 (2) 点，若 $\varphi'$ 与
$\varphi''$ 是两个与 $\varphi$ 和 $\psi$ 相容的同态
$\mathcal{F}'\to\mathcal{G}'$，那么 $\varphi'-\varphi''$ 分解为
$$
\mathcal{F}' \to \mathcal{F} \to \mathcal{L} \to \mathcal{G}'
$$
由《位点上的模》引理 \ref{sites-modules-lemma-i-star-equivalence}，
中间的同态来自 $\Hom_\mathcal{O}(\mathcal{F},\mathcal{L})$ 中唯一的
元素。反过来，给定该群中的元素 $\alpha$，可以把上述复合（中间同态取
$\alpha$）加到 $\varphi'$ 上。若干细节从略。
\end{proof}

\begin{lemma}
\label{defos-lemma-inf-obs-map-ringed-topoi}
设
$i:(\Sh(\mathcal{C}),\mathcal{O})\to
(\Sh(\mathcal{D}),\mathcal{O}')$ 是环化拓扑斯的一阶增厚，并给定扩张
$$
0 \to \mathcal{K} \to \mathcal{F}' \to \mathcal{F} \to 0
\quad\text{且}\quad
0 \to \mathcal{L} \to \mathcal{G}' \to \mathcal{G} \to 0
$$
如 (\ref{defos-equation-extension-ringed-topoi}) 所述，以及同态
$\varphi:\mathcal{F}\to\mathcal{G}$ 与
$\psi:\mathcal{K}\to\mathcal{L}$。假设图表
$$
\xymatrix{
\mathcal{I} \otimes_\mathcal{O} \mathcal{F}
\ar[r]_-{c_{\mathcal{F}'}} \ar[d]_{1 \otimes \varphi} &
\mathcal{K} \ar[d]^\psi \\
\mathcal{I} \otimes_\mathcal{O} \mathcal{G}
\ar[r]^-{c_{\mathcal{G}'}} &
\mathcal{L}
}
$$
交换。那么存在元素
$$
o(\varphi, \psi) \in
\Ext^1_\mathcal{O}(\mathcal{F}, \mathcal{L})
$$
它为零，当且仅当存在与 $\varphi$ 和 $\psi$ 相容的同态
$\varphi':\mathcal{F}'\to\mathcal{G}'$。
\end{lemma}

\begin{proof}
可以显式构造扩张
$$
0 \to \mathcal{L} \to \mathcal{H} \to \mathcal{F} \to 0
$$
方法是令 $\mathcal{H}$ 为复形
$$
\mathcal{K}
\xrightarrow{1, - \psi}
\mathcal{F}' \oplus \mathcal{G}' \xrightarrow{\varphi, 1}
\mathcal{G}
$$
在中间项处的上同调层（记号含义显然）。利用引理中图表交换这一假设，
对局部截面作计算可知 $\mathcal{H}$ 被 $\mathcal{I}$ 零化。因此
$\mathcal{H}$ 在下述群中定义一个类：
$$
\Ext^1_\mathcal{O}(\mathcal{F}, \mathcal{L})
\subset
\Ext^1_{\mathcal{O}'}(\mathcal{F}, \mathcal{L})
$$
最后，$\mathcal{H}$ 的类等于沿 $\psi$ 对扩张 $\mathcal{F}'$ 所作的
推出与沿 $\varphi$ 对扩张 $\mathcal{G}'$ 所作的拉回之差（计算从略）。
因此，$\mathcal{H}$ 的类为零，等价于存在交换图
$$
\xymatrix{
0 \ar[r] &
\mathcal{K} \ar[r] \ar[d]_{\psi} &
\mathcal{F}' \ar[r] \ar[d]_{\varphi'} &
\mathcal{F} \ar[r] \ar[d]_\varphi & 0\\
0 \ar[r] &
\mathcal{L} \ar[r] &
\mathcal{G}' \ar[r] &
\mathcal{G} \ar[r] & 0
}
$$
这正是所需。
\end{proof}

\begin{lemma}
\label{defos-lemma-inf-ext-ringed-topoi}
设
$i:(\Sh(\mathcal{C}),\mathcal{O})\to
(\Sh(\mathcal{D}),\mathcal{O}')$ 是环化拓扑斯的一阶增厚。给定
$\mathcal{O}$-模 $\mathcal{F}$、$\mathcal{K}$ 以及
$\mathcal{O}$-线性同态
$c:\mathcal{I}\otimes_\mathcal{O}\mathcal{F}\to\mathcal{K}$。
如果存在满足 $c_{\mathcal{F}'}=c$ 的正合列
(\ref{defos-equation-extension-ringed-topoi})，那么这些扩张的同构类集合是
$\Ext^1_\mathcal{O}(\mathcal{F},\mathcal{K})$ 作用下的主齐性空间。
\end{lemma}

\begin{proof}
假设给定扩张
$$
0 \to \mathcal{K} \to \mathcal{F}'_1 \to \mathcal{F} \to 0
\quad\text{且}\quad
0 \to \mathcal{K} \to \mathcal{F}'_2 \to \mathcal{F} \to 0
$$
且 $c_{\mathcal{F}'_1}=c_{\mathcal{F}'_2}=c$。那么二者之差
（在扩张群中取差，见《同调代数》第 \ref{homology-section-extensions} 节）
是扩张
$$
0 \to \mathcal{K} \to \mathcal{E} \to \mathcal{F} \to 0
$$
其中 $\mathcal{E}$ 被 $\mathcal{I}$ 零化（局部计算从略）。所以该正合列
是 $\mathcal{O}$-模的扩张，见《位点上的模》引理
\ref{sites-modules-lemma-i-star-equivalence}。反过来，给定这样的扩张
$\mathcal{E}$，可将它加到 $\mathcal{O}'$-扩张 $\mathcal{F}'$ 上，
而不改变同态 $c_{\mathcal{F}'}$。若干细节从略。
\end{proof}

\begin{lemma}
\label{defos-lemma-inf-obs-ext-ringed-topoi}
设 $i : (\Sh(\mathcal{C}), \mathcal{O}) \to
(\Sh(\mathcal{D}), \mathcal{O}')$ 是环化拓扑斯的一阶增厚。给定
$\mathcal{O}$-模 $\mathcal{F}$、$\mathcal{K}$ 以及 $\mathcal{O}$-线性同态
$c : \mathcal{I} \otimes_\mathcal{O} \mathcal{F} \to \mathcal{K}$。
那么存在一个元素
$$
o(\mathcal{F}, \mathcal{K}, c) \in
\Ext^2_\mathcal{O}(\mathcal{F}, \mathcal{K})
$$
其消失是存在满足 $c_{\mathcal{F}'}=c$ 的正合列
(\ref{defos-equation-extension-ringed-topoi}) 的充要条件。
\end{lemma}

\begin{proof}
我们先证明：若 $\mathcal{K}$ 是内射 $\mathcal{O}$-模，
则确实存在满足 $c_{\mathcal{F}'}=c$ 的正合列
(\ref{defos-equation-extension-ringed-topoi})。为此，取一个平坦
$\mathcal{O}'$-模 $\mathcal{H}'$ 以及满射
$\mathcal{H}' \to \mathcal{F}$
（《位点上的模》引理 \ref{sites-modules-lemma-module-quotient-flat}）。
令 $\mathcal{J}\subset\mathcal{H}'$ 为其核。由于 $\mathcal{H}'$ 平坦，故有
$$
\mathcal{I} \otimes_{\mathcal{O}'} \mathcal{H}' =
\mathcal{I}\mathcal{H}'
\subset \mathcal{J} \subset \mathcal{H}'
$$
注意，同态
$$
\mathcal{I}\mathcal{H}' =
\mathcal{I} \otimes_{\mathcal{O}'} \mathcal{H}'
\longrightarrow
\mathcal{I} \otimes_{\mathcal{O}'} \mathcal{F} =
\mathcal{I} \otimes_\mathcal{O} \mathcal{F}
$$
零化 $\mathcal{I}\mathcal{J}$。事实上，若 $f$ 是 $\mathcal{I}$ 的局部截面，
$s$ 是 $\mathcal{H}$ 的局部截面，则 $fs$ 映为
$f\otimes\overline{s}$，其中 $\overline{s}$ 是 $s$ 在
$\mathcal{F}$ 中的像。因此得到 $\mathcal{O}$-模的图
$$
\xymatrix{
\mathcal{I}\mathcal{H}'/\mathcal{I}\mathcal{J}
\ar@{^{(}->}[r] \ar[d] &
\mathcal{J}/\mathcal{I}\mathcal{J} \ar@{..>}[d]_\gamma \\
\mathcal{I} \otimes_\mathcal{O} \mathcal{F} \ar[r]^-c &
\mathcal{K}
}
$$
若 $\mathcal{K}$ 作为 $\mathcal{O}$-模是内射的，就得到虚线箭头。
以 $\gamma':\mathcal{J}\to\mathcal{K}$ 记 $\gamma$ 与
$\mathcal{J}\to\mathcal{J}/\mathcal{I}\mathcal{J}$ 的复合。
局部计算表明，推出
$$
\xymatrix{
0 \ar[r] &
\mathcal{J} \ar[r] \ar[d]_{\gamma'} &
\mathcal{H}' \ar[r] \ar[d] &
\mathcal{F} \ar[r] \ar@{=}[d] &
0 \\
0 \ar[r] &
\mathcal{K} \ar[r] &
\mathcal{F}' \ar[r] &
\mathcal{F} \ar[r] &
0
}
$$
给出了本引理所述问题的一个解。

\medskip\noindent
一般情形。取嵌入 $\mathcal{K}\subset\mathcal{K}'$，其中
$\mathcal{K}'$ 是内射 $\mathcal{O}$-模。令 $\mathcal{Q}$ 为其商，
于是有正合列
$$
0 \to \mathcal{K} \to \mathcal{K}' \to \mathcal{Q} \to 0
$$
以 $c':\mathcal{I}\otimes_\mathcal{O}\mathcal{F}\to\mathcal{K}'$
记相应的复合同态。由上一段，存在正合列
$$
0 \to \mathcal{K}' \to \mathcal{E}' \to \mathcal{F} \to 0
$$
如 (\ref{defos-equation-extension-ringed-topoi}) 所示，且
$c_{\mathcal{E}'}=c'$。注意，$c'$ 与
$\mathcal{K}'\to\mathcal{Q}$ 的复合为零，故沿
$\mathcal{K}'\to\mathcal{Q}$ 对 $\mathcal{E}'$ 取推出所得扩张
$$
0 \to \mathcal{Q} \to \mathcal{D}' \to \mathcal{F} \to 0
$$
如 (\ref{defos-equation-extension-ringed-topoi}) 所示，且
$c_{\mathcal{D}'}=0$。这恰好意味着 $\mathcal{D}'$ 被
$\mathcal{I}$ 零化；换言之，$\mathcal{D}'$ 是 $\mathcal{O}$-模的扩张，
因而定义了元素
$$
o(\mathcal{F}, \mathcal{K}, c) \in
\Ext^1_\mathcal{O}(\mathcal{F}, \mathcal{Q}) =
\Ext^2_\mathcal{O}(\mathcal{F}, \mathcal{K})
$$
（该等式来自上述正合列所诱导的上同调长正合列，以及以内射模
$\mathcal{K}'$ 为第二变元的高阶 $\Ext$ 群消失）。若
$o(\mathcal{F},\mathcal{K},c)=0$，则可取分裂
$s:\mathcal{F}\to\mathcal{D}'$，并令
$$
\mathcal{F}' = \Ker(\mathcal{E}' \to \mathcal{D}'/s(\mathcal{F}))
$$
从而得到下图
$$
\xymatrix{
0 \ar[r] &
\mathcal{K} \ar[r] \ar[d] &
\mathcal{F}' \ar[r] \ar[d] &
\mathcal{F} \ar[r] \ar@{=}[d] &
0 \\
0 \ar[r] &
\mathcal{K}' \ar[r] &
\mathcal{E}' \ar[r] &
\mathcal{F} \ar[r] & 0
}
$$
其两行均正合，由此可见 $c_{\mathcal{F}'}=c$。反过来，若
$\mathcal{F}'$ 存在，则由引理 \ref{defos-lemma-inf-ext-ringed-topoi}
及以内射模 $\mathcal{K}'$ 为第二变元的高阶 $\Ext$ 群消失，
沿 $\mathcal{K}\to\mathcal{K}'$ 对 $\mathcal{F}'$ 取推出所得扩张
同构于 $\mathcal{E}'$。这给出如上图，从而说明
$\mathcal{D}'$ 作为扩张是分裂的，即类
$o(\mathcal{F},\mathcal{K},c)$ 为零。
\end{proof}

\begin{remark}
\label{defos-remark-trivial-thickening-ringed-topoi}
设 $(\Sh(\mathcal{C}),\mathcal{O})$ 是环化拓扑斯。若存在环化拓扑斯态射
$\pi : (\Sh(\mathcal{D}), \mathcal{O}') \to (\Sh(\mathcal{C}), \mathcal{O})$
作为 $i:(\Sh(\mathcal{C}),\mathcal{O})\to
(\Sh(\mathcal{D}),\mathcal{O}')$ 的左逆，就称这一阶增厚为
{\it 平凡的}。这样一个态射 $\pi$ 的选取称为该一阶增厚的一个
{\it 平凡化}。给定 $\pi$，利用 $\pi^\sharp$ 分裂满射
$\mathcal{O}'\to\mathcal{O}$，就得到 $\mathcal{C}$ 上代数层的分解
\begin{equation}
\label{defos-equation-splitting-ringed-topoi}
\mathcal{O}' = \mathcal{O} \oplus \mathcal{I}
\end{equation}
反过来，这样的分解决定一个态射 $\pi$。
$(\Sh(\mathcal{C}),\mathcal{O})$ 的带平凡化一阶增厚所成范畴
等价于 $\mathcal{O}$-模范畴。
\end{remark}

\begin{remark}
\label{defos-remark-trivial-extension-ringed-topoi}
设 $i:(\Sh(\mathcal{C}),\mathcal{O})\to
(\Sh(\mathcal{D}),\mathcal{O}')$ 是环化拓扑斯的平凡一阶增厚，
$\pi:(\Sh(\mathcal{D}),\mathcal{O}')\to
(\Sh(\mathcal{C}),\mathcal{O})$ 是一个平凡化。给定任意三元组
$(\mathcal{F},\mathcal{K},c)$，其中 $\mathcal{F}$、$\mathcal{K}$
是 $\mathcal{O}$-模，且有同态
$c : \mathcal{I} \otimes_\mathcal{O} \mathcal{F} \to \mathcal{K}$
可令
$$
\mathcal{F}'_{c, triv} = \mathcal{F} \oplus \mathcal{K}
$$
并利用与 $\pi$ 相伴的分解
(\ref{defos-equation-splitting-ringed-topoi}) 及同态 $c$ 定义
$\mathcal{O}'$-模结构，从而得到扩张
(\ref{defos-equation-extension-ringed-topoi})。我们称
$\mathcal{F}'_{c,triv}$ 为与 $c$ 和平凡化 $\pi$ 对应的、以
$\mathcal{K}$ 扩张 $\mathcal{F}$ 的{\it 平凡扩张}。
给定任意形如 (\ref{defos-equation-extension-ringed-topoi}) 的扩张
$\mathcal{F}'$，可利用 $\pi^\sharp:\mathcal{O}\to\mathcal{O}'$
把 $\mathcal{F}'$ 视为 $\mathcal{O}$-模扩张，因而得到
$\Ext^1_\mathcal{O}(\mathcal{F},\mathcal{K})$ 中的类
$\xi_{\mathcal{F}'}$。引理 \ref{defos-lemma-inf-ext-ringed-topoi}
保证 $\mathcal{F}'\mapsto\xi_{\mathcal{F}'}$ 诱导双射
$$
\left\{
\begin{matrix}
\text{满足下述条件的扩张 }\mathcal{F}'\text{ 的同构类}\\
\mathcal{F}'\text{ 如 (\ref{defos-equation-extension-ringed-topoi}) 所示，且 }
c=c_{\mathcal{F}'}
\end{matrix}
\right\}
\longrightarrow
\Ext^1_\mathcal{O}(\mathcal{F}, \mathcal{K})
$$
此外，平凡扩张 $\mathcal{F}'_{c,triv}$ 映到零类。
\end{remark}

\begin{remark}
\label{defos-remark-extension-functorial-ringed-topoi}
设 $(\Sh(\mathcal{C}),\mathcal{O})$ 是环化拓扑斯。设
$(\Sh(\mathcal{C}), \mathcal{O}) \to (\Sh(\mathcal{D}_i), \mathcal{O}'_i)$,
$i=1,2$ 是理想层为 $\mathcal{I}_i$ 的一阶增厚。设
$h:(\Sh(\mathcal{D}_1),\mathcal{O}'_1)\to
(\Sh(\mathcal{D}_2), \mathcal{O}'_2)$
是 $(\Sh(\mathcal{C}),\mathcal{O})$ 的一阶增厚之间的态射。图示为
$$
\xymatrix{
& (\Sh(\mathcal{C}), \mathcal{O}) \ar[ld] \ar[rd] & \\
(\Sh(\mathcal{D}_1), \mathcal{O}'_1) \ar[rr]^h & & 
(\Sh(\mathcal{D}_2), \mathcal{O}'_2)
}
$$
注意，$h^\sharp:\mathcal{O}'_2\to\mathcal{O}'_1$
特别诱导 $\mathcal{O}$-模同态 $\mathcal{I}_2\to\mathcal{I}_1$。
设 $\mathcal{F}$ 是 $\mathcal{O}$-模。设
$(\mathcal{K}_i,c_i)$，$i=1,2$，是两个二元组，
其中 $\mathcal{K}_i$ 是 $\mathcal{O}$-模，并有同态
$c_i : \mathcal{I}_i \otimes_\mathcal{O} \mathcal{F} \to
\mathcal{K}_i$。再给定 $\mathcal{O}$-模同态
$\mathcal{K}_2\to\mathcal{K}_1$，使图
$$
\xymatrix{
\mathcal{I}_2 \otimes_\mathcal{O} \mathcal{F}
\ar[r]_-{c_2} \ar[d] &
\mathcal{K}_2 \ar[d] \\
\mathcal{I}_1 \otimes_\mathcal{O} \mathcal{F}
\ar[r]^-{c_1} &
\mathcal{K}_1
}
$$
交换。那么存在典范的函子性映射
$$
\left\{
\begin{matrix}
\mathcal{F}'_2\text{ 如 (\ref{defos-equation-extension-ringed-topoi}) 所示，且 }\\
c_2=c_{\mathcal{F}'_2}\text{、}\mathcal{K}=\mathcal{K}_2
\end{matrix}
\right\}
\longrightarrow
\left\{
\begin{matrix}
\mathcal{F}'_1\text{ 如 (\ref{defos-equation-extension-ringed-topoi}) 所示，且 }\\
c_1=c_{\mathcal{F}'_1}\text{、}\mathcal{K}=\mathcal{K}_1
\end{matrix}
\right\}
$$
具体地，把 $\mathcal{O}$、$\mathcal{O}'_i$、$\mathcal{F}$、
$\mathcal{K}_i$ 等都视为 $\mathcal{C}$ 上的层；给定
$\mathcal{F}'_2$，令层 $\mathcal{F}'_1$ 等于推出，
即嵌入下列扩张图：
$$
\xymatrix{
0 \ar[r] &
\mathcal{K}_2 \ar[r] \ar[d] &
\mathcal{F}'_2 \ar[r] \ar[d] &
\mathcal{F} \ar@{=}[d] \ar[r] & 0 \\
0 \ar[r] &
\mathcal{K}_1 \ar[r] &
\mathcal{F}'_1 \ar[r] &
\mathcal{F} \ar[r] & 0
}
$$
推出上的 $\mathcal{O}'_1$-模结构之构造从略
（这里用到含 $c_1$、$c_2$ 的图的交换性）。
\end{remark}

\begin{remark}
\label{defos-remark-trivial-extension-functorial-ringed-topoi}
设 $(\Sh(\mathcal{C}),\mathcal{O})$、
$(\Sh(\mathcal{C}), \mathcal{O}) \to (\Sh(\mathcal{D}_i), \mathcal{O}'_i)$,
$\mathcal{I}_i$ 及 $h:(\Sh(\mathcal{D}_1),\mathcal{O}'_1)\to
(\Sh(\mathcal{D}_2),\mathcal{O}'_2)$ 如注
\ref{defos-remark-extension-functorial-ringed-topoi}。设给定平凡化
$\pi_i : (\Sh(\mathcal{D}_i), \mathcal{O}'_i) \to
(\Sh(\mathcal{C}), \mathcal{O})$，满足
$\pi_1=h\circ\pi_2$。换言之，设 $h$ 是
$(\Sh(\mathcal{C}),\mathcal{O})$ 的带平凡化一阶增厚之间的态射。
设 $(\mathcal{K}_i,c_i)$，$i=1,2$，是两个二元组，
其中 $\mathcal{K}_i$ 是 $\mathcal{O}$-模，并有同态
$c_i : \mathcal{I}_i \otimes_\mathcal{O} \mathcal{F} \to
\mathcal{K}_i$。再给定 $\mathcal{O}$-模同态
$\mathcal{K}_2\to\mathcal{K}_1$，使下图交换：
$$
\xymatrix{
\mathcal{I}_2 \otimes_\mathcal{O} \mathcal{F}
\ar[r]_-{c_2} \ar[d] &
\mathcal{K}_2 \ar[d] \\
\mathcal{I}_1 \otimes_\mathcal{O} \mathcal{F}
\ar[r]^-{c_1} &
\mathcal{K}_1
}
$$
在此情形，注 \ref{defos-remark-trivial-extension-ringed-topoi} 的构造
诱导交换图
$$
\xymatrix{
\{\mathcal{F}'_2\text{ 如 (\ref{defos-equation-extension-ringed-topoi}) 所示，且 }
c_2=c_{\mathcal{F}'_2}\text{、}\mathcal{K}=\mathcal{K}_2\}
\ar[d] \ar[rr] & &
\Ext^1_\mathcal{O}(\mathcal{F}, \mathcal{K}_2) \ar[d] \\
\{\mathcal{F}'_1\text{ 如 (\ref{defos-equation-extension-ringed-topoi}) 所示，且 }
c_1=c_{\mathcal{F}'_1}\text{、}\mathcal{K}=\mathcal{K}_1\}
\ar[rr] & &
\Ext^1_\mathcal{O}(\mathcal{F}, \mathcal{K}_1)
}
$$
其中右侧竖直箭头由 $\Ext$ 的函子性及同态
$\mathcal{K}_2\to\mathcal{K}_1$ 给出，左侧竖直箭头则是注
\ref{defos-remark-extension-functorial-ringed-topoi} 中的映射。
\end{remark}

\begin{remark}
\label{defos-remark-obstruction-extension-functorial-ringed-topoi}
设 $(\Sh(\mathcal{C}),\mathcal{O})$、
$(\Sh(\mathcal{C}), \mathcal{O}) \to (\Sh(\mathcal{D}_i), \mathcal{O}'_i)$,
$\mathcal{I}_i$ 及 $h:(\Sh(\mathcal{D}_1),\mathcal{O}'_1)\to
(\Sh(\mathcal{D}_2),\mathcal{O}'_2)$ 如注
\ref{defos-remark-extension-functorial-ringed-topoi}。注意，
$h^\sharp:\mathcal{O}'_2\to\mathcal{O}'_1$ 特别诱导
$\mathcal{O}$-模同态 $\mathcal{I}_2\to\mathcal{I}_1$。
设 $\mathcal{F}$ 是 $\mathcal{O}$-模。设
$(\mathcal{K}_i,c_i)$，$i=1,2$，是两个二元组，
其中 $\mathcal{K}_i$ 是 $\mathcal{O}$-模，并有同态
$c_i : \mathcal{I}_i \otimes_\mathcal{O} \mathcal{F} \to
\mathcal{K}_i$。再给定 $\mathcal{O}$-模同态
$\mathcal{K}_2\to\mathcal{K}_1$，使下图交换：
$$
\xymatrix{
\mathcal{I}_2 \otimes_\mathcal{O} \mathcal{F}
\ar[r]_-{c_2} \ar[d] &
\mathcal{K}_2 \ar[d] \\
\mathcal{I}_1 \otimes_\mathcal{O} \mathcal{F}
\ar[r]^-{c_1} &
\mathcal{K}_1
}
$$
我们{\bf 断言}，同态
$$
\Ext^2_\mathcal{O}(\mathcal{F}, \mathcal{K}_2)
\longrightarrow
\Ext^2_\mathcal{O}(\mathcal{F}, \mathcal{K}_1)
$$
把 $o(\mathcal{F},\mathcal{K}_2,c_2)$ 映到
$o(\mathcal{F},\mathcal{K}_1,c_1)$。

\medskip\noindent
为证明此断言，取嵌入 $j_2:\mathcal{K}_2\to\mathcal{K}_2'$，
其中 $\mathcal{K}_2'$ 是内射 $\mathcal{O}$-模。
如引理 \ref{defos-lemma-inf-obs-ext-ringed-topoi} 的证明所示，
可取一个 $\mathcal{O}_2$-模扩张
% P11-E0005: unprimed O_i symbols retained from frozen source; see control/ERRATA_CANDIDATES.jsonl.
$$
0 \to \mathcal{K}_2' \to \mathcal{E}_2 \to \mathcal{F} \to 0
$$
使 $c_{\mathcal{E}_2}=j_2\circ c_2$。引理
\ref{defos-lemma-inf-obs-ext-ringed-topoi} 的证明把
$o(\mathcal{F},\mathcal{K}_2,c_2)$ 构造为下列
$\mathcal{O}$-模正合列的 Yoneda 扩张类（其含义见《导出范畴》第
\ref{derived-section-ext} 节）：
$$
0 \to
\mathcal{K}_2 \to \mathcal{K}_2' \to
\mathcal{E}_2/\mathcal{K}_2 \to
\mathcal{F} \to 0
$$
令 $\mathcal{K}_1'$ 为
$\mathcal{K}_2\to\mathcal{K}_1\oplus\mathcal{K}_2'$ 的余核。
存在单射 $j_1:\mathcal{K}_1\to\mathcal{K}_1'$ 以及同态
$\mathcal{K}_2'\to\mathcal{K}_1'$，二者组成交换方块。取推出：
$$
\xymatrix{
0 \ar[r] &
\mathcal{K}_2' \ar[r] \ar[d] &
\mathcal{E}_2 \ar[r] \ar[d] &
\mathcal{F} \ar[r] \ar[d] & 0 \\
0 \ar[r] &
\mathcal{K}_1' \ar[r] &
\mathcal{E}_1 \ar[r] &
\mathcal{F} \ar[r] & 0
}
$$
$\mathcal{E}_1$ 上有典范的 $\mathcal{O}_1$-模结构，且在此结构下
$c_{\mathcal{E}_1}=j_1\circ c_1$（这里用到上面含
$c_1$、$c_2$ 的图的交换性）。引理
\ref{defos-lemma-inf-obs-ext-ringed-topoi} 的步骤告诉我们，
$o(\mathcal{F},\mathcal{K}_1,c_1)$ 是下列
$\mathcal{O}$-模正合列的 Yoneda 扩张类：
$$
0 \to
\mathcal{K}_1 \to
\mathcal{K}_1' \to
\mathcal{E}_1/\mathcal{K}_1 \to
\mathcal{F} \to 0
$$
由于存在正合列之间的同态
% P11-E0006: duplicated lower row retained from frozen source; see control/ERRATA_CANDIDATES.jsonl.
$$
\xymatrix{
0 \ar[r] &
\mathcal{K}_2 \ar[d] \ar[r] &
\mathcal{K}_2' \ar[d] \ar[r] &
\mathcal{E}_2/\mathcal{K}_2 \ar[r] \ar[d] &
\mathcal{F} \ar[r] \ar@{=}[d] &
0 \\
0 \ar[r] &
\mathcal{K}_2 \ar[r] &
\mathcal{K}_2' \ar[r] &
\mathcal{E}_2/\mathcal{K}_2 \ar[r] &
\mathcal{F} \ar[r] &
0
}
$$
故断言成立。
\end{remark}

\begin{remark}
\label{defos-remark-short-exact-sequence-thickenings-ringed-topoi}
设 $(\Sh(\mathcal{C}),\mathcal{O})$ 是环化拓扑斯。给定其一阶增厚的态射列
$$
(\Sh(\mathcal{D}_1), \mathcal{O}'_1) \to
(\Sh(\mathcal{D}_2), \mathcal{O}'_2) \to
(\Sh(\mathcal{D}_3), \mathcal{O}'_3)
$$
若理想层 $\mathcal{I}_i$ 之间的相应同态给出 $\mathcal{O}$-模复形
$\mathcal{I}_3 \to \mathcal{I}_2 \to \mathcal{I}_1$
（即复合为零），就称上述态射列为一个{\it 复形}。在此情形，复合
$(\Sh(\mathcal{D}_1), \mathcal{O}'_1) \to
(\Sh(\mathcal{D}_3), \mathcal{O}'_3)$ 分解通过
$(\Sh(\mathcal{C}), \mathcal{O}) \to
(\Sh(\mathcal{D}_3), \mathcal{O}'_3)$；换言之，
$(\Sh(\mathcal{D}_1),\mathcal{O}'_1)$ 作为
$(\Sh(\mathcal{C}),\mathcal{O})$ 的一阶增厚是平凡的，并带有典范平凡化
$\pi : (\Sh(\mathcal{D}_1), \mathcal{O}'_1) \to
(\Sh(\mathcal{C}), \mathcal{O})$。

\medskip\noindent
若 $(\Sh(\mathcal{C}),\mathcal{O})$ 的一阶增厚的态射列
$$
(\Sh(\mathcal{D}_1), \mathcal{O}'_1) \to
(\Sh(\mathcal{D}_2), \mathcal{O}'_2) \to
(\Sh(\mathcal{D}_3), \mathcal{O}'_3)
$$
所对应的理想层同态构成 $\mathcal{O}$-模短正合列
$$
0 \to \mathcal{I}_3 \to \mathcal{I}_2 \to \mathcal{I}_1 \to 0
$$
就称这一态射列为{\it 短正合列}。
\end{remark}

\begin{remark}
\label{defos-remark-complex-thickenings-and-ses-modules-ringed-topoi}
设 $(\Sh(\mathcal{C}),\mathcal{O})$ 是环化拓扑斯，
$\mathcal{F}$ 是 $\mathcal{O}$-模。设
$$
(\Sh(\mathcal{D}_1), \mathcal{O}'_1) \to
(\Sh(\mathcal{D}_2), \mathcal{O}'_2) \to
(\Sh(\mathcal{D}_3), \mathcal{O}'_3)
$$
是 $(\Sh(\mathcal{C}),\mathcal{O})$ 的一阶增厚复形，见注
\ref{defos-remark-short-exact-sequence-thickenings-ringed-topoi}。
设 $(\mathcal{K}_i,c_i)$，$i=1,2,3$，是三个二元组，
其中 $\mathcal{K}_i$ 是 $\mathcal{O}$-模，并有同态
$c_i : \mathcal{I}_i \otimes_\mathcal{O} \mathcal{F} \to
\mathcal{K}_i$。设给定 $\mathcal{O}$-模短正合列
$$
0 \to \mathcal{K}_3 \to \mathcal{K}_2 \to \mathcal{K}_1 \to 0
$$
使得
$$
\vcenter{
\xymatrix{
\mathcal{I}_2 \otimes_\mathcal{O} \mathcal{F}
\ar[r]_-{c_2} \ar[d] &
\mathcal{K}_2 \ar[d] \\
\mathcal{I}_1 \otimes_\mathcal{O} \mathcal{F}
\ar[r]^-{c_1} &
\mathcal{K}_1
}
}
\quad\text{且}\quad
\vcenter{
\xymatrix{
\mathcal{I}_3 \otimes_\mathcal{O} \mathcal{F}
\ar[r]_-{c_3} \ar[d] &
\mathcal{K}_3 \ar[d] \\
\mathcal{I}_2 \otimes_\mathcal{O} \mathcal{F}
\ar[r]^-{c_2} &
\mathcal{K}_2
}
}
$$
均为交换图。最后，设给定 $\mathcal{O}'_2$-模扩张
$$
0 \to \mathcal{K}_2 \to \mathcal{F}'_2 \to \mathcal{F} \to 0
$$
如 (\ref{defos-equation-extension-ringed-topoi}) 所示，其中
$\mathcal{K}=\mathcal{K}_2$，且 $c_{\mathcal{F}'_2}=c_2$。
在此情形，可应用注 \ref{defos-remark-extension-functorial-ringed-topoi}
的函子性，得到 $\mathcal{O}'_1$-模扩张 $\mathcal{F}'_1$
（下面会描述这一特殊情形中的 $\mathcal{F}'_1$）。由注
\ref{defos-remark-trivial-extension-ringed-topoi}，利用注
\ref{defos-remark-short-exact-sequence-thickenings-ringed-topoi} 中的典范分裂
$\pi : (\Sh(\mathcal{D}_1), \mathcal{O}'_1) \to
(\Sh(\mathcal{C}), \mathcal{O})$，得到
$\xi_{\mathcal{F}'_1} \in
\Ext^1_\mathcal{O}(\mathcal{F}, \mathcal{K}_1)$.
最后，还有障碍
$$
o(\mathcal{F}, \mathcal{K}_3, c_3) \in
\Ext^2_\mathcal{O}(\mathcal{F}, \mathcal{K}_3)
$$
见引理 \ref{defos-lemma-inf-obs-ext-ringed-topoi}。在此情形，
我们{\bf 断言}，由短正合列
$0\to\mathcal{K}_3\to\mathcal{K}_2\to\mathcal{K}_1\to0$
给出的典范同态
$$
\partial :
\Ext^1_\mathcal{O}(\mathcal{F}, \mathcal{K}_1)
\longrightarrow
\Ext^2_\mathcal{O}(\mathcal{F}, \mathcal{K}_3)
$$
把 $\xi_{\mathcal{F}'_1}$ 映到障碍类
$o(\mathcal{F},\mathcal{K}_3,c_3)$。

\medskip\noindent
为证明此断言，取嵌入 $j:\mathcal{K}_3\to\mathcal{K}$，
其中 $\mathcal{K}$ 是内射 $\mathcal{O}$-模。可把 $j$ 提升为同态
$j':\mathcal{K}_2\to\mathcal{K}$。令
$\mathcal{E}'_2=j'_*\mathcal{F}'_2$ 为沿 $j'$ 对
$\mathcal{F}'_2$ 取的推出，使
$c_{\mathcal{E}'_2}=j'\circ c_2$。图示如下：
$$
\xymatrix{
0 \ar[r] &
\mathcal{K}_2 \ar[r] \ar[d]_{j'} &
\mathcal{F}'_2 \ar[r] \ar[d] &
\mathcal{F} \ar[r] \ar[d] & 0 \\
0 \ar[r] &
\mathcal{K} \ar[r] &
\mathcal{E}'_2 \ar[r] &
\mathcal{F} \ar[r] & 0
}
$$
令 $\mathcal{E}'_3=\mathcal{E}'_2$，但通过
$\mathcal{O}'_3\to\mathcal{O}'_2$ 把它视为 $\mathcal{O}'_3$-模。
于是 $c_{\mathcal{E}'_3}=j\circ c_3$。引理
\ref{defos-lemma-inf-obs-ext-ringed-topoi} 的证明把
$o(\mathcal{F},\mathcal{K}_3,c_3)$ 构造为下列
$\mathcal{O}$-模扩张之类的边界：
$$
0 \to
\mathcal{K}/\mathcal{K}_3 \to
\mathcal{E}'_3/\mathcal{K}_3 \to
\mathcal{F} \to 0
$$
另一方面，注意 $\mathcal{F}'_1=\mathcal{F}'_2/\mathcal{K}_3$，
故类 $\xi_{\mathcal{F}'_1}$ 是扩张
$$
0 \to \mathcal{K}_2/\mathcal{K}_3 \to \mathcal{F}'_2/\mathcal{K}_3
\to \mathcal{F} \to 0
$$
之类；这里借助 $\pi^\sharp$ 把它视为 $\mathcal{O}$-模列，
其中 $\pi:(\Sh(\mathcal{D}_1),\mathcal{O}'_1)\to
(\Sh(\mathcal{C}),\mathcal{O})$ 是典范分裂。最后，存在交换图
$$
\xymatrix{
0 \ar[r] &
\mathcal{K}_2/\mathcal{K}_3 \ar[r] \ar[d] &
\mathcal{F}'_2/\mathcal{K}_3 \ar[r] \ar[d] &
\mathcal{F} \ar[r] \ar[d] & 0 \\
0 \ar[r] &
\mathcal{K}/\mathcal{K}_3 \ar[r] &
\mathcal{E}'_3/\mathcal{K}_3 \ar[r] &
\mathcal{F} \ar[r] & 0
}
$$
且该图关于上述 $\mathcal{O}$-模结构是 $\mathcal{O}$-线性的，
故断言成立。
\end{remark}







\section{环化拓扑斯上模的无穷小形变}
\label{defos-section-deformation-modules-ringed-topoi}

\noindent
设 $i:(\Sh(\mathcal{C}),\mathcal{O})\to
(\Sh(\mathcal{D}),\mathcal{O}')$ 是环化拓扑斯的一阶增厚。
我们沿用第 \ref{defos-section-thickenings-ringed-topoi} 节引入的记号。
设 $\mathcal{F}'$ 是 $\mathcal{O}'$-模，并令
$\mathcal{F}=i^*\mathcal{F}'$。此时有 $\mathcal{O}'$-模短正合列
$$
0 \to \mathcal{I}\mathcal{F}' \to \mathcal{F}' \to \mathcal{F} \to 0
$$
由于 $\mathcal{I}^2=0$，$\mathcal{I}\mathcal{F}'$ 上的
$\mathcal{O}'$-模结构来自唯一的 $\mathcal{O}$-模结构。
因此上述正合列是形如 (\ref{defos-equation-extension-ringed-topoi}) 的扩张。
特别地，若 $\mathcal{F}'=\mathcal{O}'$，则
$i^*\mathcal{O}'=\mathcal{O}$ 且
$\mathcal{I}\mathcal{O}'=\mathcal{I}$，于是恢复结构层正合列
$$
0 \to \mathcal{I} \to \mathcal{O}' \to \mathcal{O} \to 0
$$

\begin{lemma}
\label{defos-lemma-inf-map-special-ringed-topoi}
设 $i:(\Sh(\mathcal{C}),\mathcal{O})\to
(\Sh(\mathcal{D}),\mathcal{O}')$ 是环化拓扑斯的一阶增厚。
设 $\mathcal{F}'$、$\mathcal{G}'$ 是 $\mathcal{O}'$-模，令
$\mathcal{F}=i^*\mathcal{F}'$、$\mathcal{G}=i^*\mathcal{G}'$。
设 $\varphi:\mathcal{F}\to\mathcal{G}$ 是 $\mathcal{O}$-线性同态。
若把 $\varphi$ 提升为 $\mathcal{O}'$-线性同态
$\varphi':\mathcal{F}'\to\mathcal{G}'$ 的提升集合非空，
则它是 $\Hom_\mathcal{O}(\mathcal{F},\mathcal{I}\mathcal{G}')$
作用下的主齐性空间。
\end{lemma}

\begin{proof}
这是引理 \ref{defos-lemma-inf-map-ringed-topoi} 的特殊情形，
但我们也给出直接证明。有模的短正合列
$$
0 \to \mathcal{I} \to \mathcal{O}' \to \mathcal{O} \to 0
\quad\text{且}\quad
0 \to \mathcal{I}\mathcal{G}' \to \mathcal{G}' \to \mathcal{G} \to 0
$$
$\mathcal{F}'$ 也有类似的正合列。由于 $\mathcal{I}$ 平方为零，
$\mathcal{I}$ 与 $\mathcal{I}\mathcal{G}'$ 上的 $\mathcal{O}'$-模结构
都来自唯一的 $\mathcal{O}$-模结构。因此
$$
\Hom_{\mathcal{O}'}(\mathcal{F}', \mathcal{I}\mathcal{G}') =
\Hom_\mathcal{O}(\mathcal{F}, \mathcal{I}\mathcal{G}')
\quad\text{且}\quad
\Hom_{\mathcal{O}'}(\mathcal{F}', \mathcal{G}) =
\Hom_\mathcal{O}(\mathcal{F}, \mathcal{G})
$$
现在本引理由正合列
$$
0 \to \Hom_{\mathcal{O}'}(\mathcal{F}', \mathcal{I}\mathcal{G}') \to
\Hom_{\mathcal{O}'}(\mathcal{F}', \mathcal{G}') \to
\Hom_{\mathcal{O}'}(\mathcal{F}', \mathcal{G})
$$
得出，见《同调代数》引理 \ref{homology-lemma-check-exactness}。
\end{proof}

\begin{lemma}
\label{defos-lemma-deform-module-ringed-topoi}
设 $(f,f')$ 是情形 \ref{defos-situation-morphism-thickenings-ringed-topoi}
中的环化拓扑斯一阶增厚之间的态射。设 $\mathcal{F}'$ 是
$\mathcal{O}'$-模，并令 $\mathcal{F}=i^*\mathcal{F}'$。
假设 $\mathcal{F}$ 在 $\mathcal{O}_\mathcal{B}$ 上平坦，且
$(f,f')$ 是增厚的严格态射（定义
\ref{defos-definition-strict-morphism-thickenings-ringed-topoi}）。
那么下列条件等价：
\begin{enumerate}
\item $\mathcal{F}'$ 在 $\mathcal{O}_{\mathcal{B}'}$ 上平坦；
\item 典范同态
$f^*\mathcal{J} \otimes_\mathcal{O} \mathcal{F} \to
\mathcal{I}\mathcal{F}'$
是同构。
\end{enumerate}
此外，在此情形，同态
$$
f^*\mathcal{J} \otimes_\mathcal{O} \mathcal{F} \to
\mathcal{I} \otimes_\mathcal{O} \mathcal{F} \to
\mathcal{I}\mathcal{F}'
$$
均为同构。
\end{lemma}

\begin{proof}
由于 $(f,f')$ 是增厚的严格态射，同态
$f^*\mathcal{J}\to\mathcal{I}$ 是满射。因此最后一个断言由 (2) 得出。

\medskip\noindent
(1) 与 (2) 等价的证明。按定义，在 $\mathcal{O}_\mathcal{B}$ 上平坦
意指在 $f^{-1}\mathcal{O}_\mathcal{B}$ 上平坦；
$f^{-1}\mathcal{O}_{\mathcal{B}'}$ 上的平坦性同理。
注意，$(f,f')$ 的严格性及假设
$\mathcal{F}=i^*\mathcal{F}'$ 蕴含
$$
\mathcal{F} = \mathcal{F}'/(f^{-1}\mathcal{J})\mathcal{F}'
$$
（视为 $\mathcal{C}$ 上的层）。此外，注意
$f^*\mathcal{J} \otimes_\mathcal{O} \mathcal{F} =
f^{-1}\mathcal{J} \otimes_{f^{-1}\mathcal{O}_\mathcal{B}} \mathcal{F}$.
因此 (1) 与 (2) 的等价性由《位点上的模》引理
\ref{sites-modules-lemma-flat-over-thickening} 得出。
\end{proof}

\begin{lemma}
\label{defos-lemma-deform-fp-module-ringed-topoi}
设 $(f,f')$ 是情形 \ref{defos-situation-morphism-thickenings-ringed-topoi}
中的环化拓扑斯一阶增厚之间的态射。设 $\mathcal{F}'$ 是
$\mathcal{O}'$-模，并令 $\mathcal{F}=i^*\mathcal{F}'$。
假设 $\mathcal{F}'$ 在 $\mathcal{O}_{\mathcal{B}'}$ 上平坦，
且 $(f,f')$ 是增厚的严格态射。那么下列条件等价：
\begin{enumerate}
\item $\mathcal{F}'$ 是有限表示 $\mathcal{O}'$-模；
\item $\mathcal{F}$ 是有限表示 $\mathcal{O}$-模。
\end{enumerate}
\end{lemma}

\begin{proof}
(1) $\Rightarrow$ (2) 由《位点上的模》引理
\ref{sites-modules-lemma-local-pullback} 得出。
反过来，假设 $\mathcal{F}$ 有限表示。可以并且确实假设
$\mathcal{C}=\mathcal{C}'$。由引理
\ref{defos-lemma-deform-module-ringed-topoi}，有短正合列
% P11-E0007: undefined O_X subscript retained from frozen source; see control/ERRATA_CANDIDATES.jsonl.
$$
0 \to \mathcal{I} \otimes_{\mathcal{O}_X} \mathcal{F} \to
\mathcal{F}' \to \mathcal{F} \to 0
$$
取 $\mathcal{C}$ 的对象 $U$，使 $\mathcal{F}|_U$ 有表示
$$
\mathcal{O}_U^{\oplus m} \to \mathcal{O}_U^{\oplus n} \to \mathcal{F}|_U \to 0
$$
把 $U$ 换成某一覆盖的各成员后，可以假设同态
$\mathcal{O}_U^{\oplus n}\to\mathcal{F}|_U$ 提升为同态
$(\mathcal{O}'_U)^{\oplus n}\to\mathcal{F}'|_U$。
由 $\otimes$ 的右正合性，所诱导的同态
$\mathcal{I}^{\oplus n}\to\mathcal{I}\otimes\mathcal{F}$ 是满射。
因此再次把 $U$ 换成某一覆盖的各成员后，可找到给定同态
$\mathcal{O}_U^{\oplus m}\to\mathcal{O}_U^{\oplus n}$ 的提升
$(\mathcal{O}'|_U)^{\oplus m} \to (\mathcal{O}'|_U)^{\oplus n}$
使得
$$
(\mathcal{O}'_U)^{\oplus m} \to (\mathcal{O}'_U)^{\oplus n} \to
\mathcal{F}'|_U \to 0
$$
构成复形。再次利用 $\otimes$ 的右正合性可见此复形正合。
\end{proof}

\begin{lemma}
\label{defos-lemma-inf-map-rel-ringed-topoi}
设 $(f,f')$ 是情形 \ref{defos-situation-morphism-thickenings-ringed-topoi}
中的一阶增厚态射。设 $\mathcal{F}'$、$\mathcal{G}'$ 是
$\mathcal{O}'$-模，并令
$\mathcal{F}=i^*\mathcal{F}'$、$\mathcal{G}=i^*\mathcal{G}'$。
设 $\varphi:\mathcal{F}\to\mathcal{G}$ 是 $\mathcal{O}$-线性同态。
假设 $\mathcal{G}'$ 在 $\mathcal{O}_{\mathcal{B}'}$ 上平坦，
且 $(f,f')$ 是增厚的严格态射。若把 $\varphi$ 提升为
$\mathcal{O}'$-线性同态
$\varphi':\mathcal{F}'\to\mathcal{G}'$ 的提升集合非空，
则它是下列群作用下的主齐性空间：
$$
\Hom_\mathcal{O}(\mathcal{F},
\mathcal{G} \otimes_\mathcal{O} f^*\mathcal{J})
$$
\end{lemma}

\begin{proof}
合用引理 \ref{defos-lemma-inf-map-special-ringed-topoi} 与
\ref{defos-lemma-deform-module-ringed-topoi} 即得。
\end{proof}

\begin{lemma}
\label{defos-lemma-inf-obs-map-special-ringed-topoi}
设 $i:(\Sh(\mathcal{C}),\mathcal{O})\to
(\Sh(\mathcal{D}),\mathcal{O}')$ 是环化拓扑斯的一阶增厚。
设 $\mathcal{F}'$、$\mathcal{G}'$ 是 $\mathcal{O}'$-模，并令
$\mathcal{F}=i^*\mathcal{F}'$、$\mathcal{G}=i^*\mathcal{G}'$。
设 $\varphi:\mathcal{F}\to\mathcal{G}$ 是 $\mathcal{O}$-线性同态。
存在元素
$$
o(\varphi) \in
\Ext^1_\mathcal{O}(Li^*\mathcal{F}', \mathcal{I}\mathcal{G}')
$$
其消失是存在把 $\varphi$ 提升为 $\mathcal{O}'$-线性同态
$\varphi':\mathcal{F}'\to\mathcal{G}'$ 的提升的充要条件。
\end{lemma}

\begin{proof}
由引理 \ref{defos-lemma-inf-map-special-ringed-topoi} 的证明可知，
$\varphi$ 在同态
$$
\Hom_\mathcal{O}(\mathcal{F}, \mathcal{G}) =
\Hom_{\mathcal{O}'}(\mathcal{F}', \mathcal{G}) \longrightarrow
\Ext^1_{\mathcal{O}'}(\mathcal{F}', \mathcal{I}\mathcal{G}')
$$
下的边界消失，正是存在提升的充要条件。由导出范畴上
$i_*=Ri_*$ 与 $Li^*$ 的伴随性（《位点上的上同调》引理
\ref{sites-cohomology-lemma-adjoint}），有
$$
\Ext^1_{\mathcal{O}'}(\mathcal{F}', \mathcal{I}\mathcal{G}') =
\Ext^1_\mathcal{O}(Li^*\mathcal{F}', \mathcal{I}\mathcal{G}')
$$
结论随即得出。
\end{proof}

\begin{lemma}
\label{defos-lemma-inf-obs-map-rel-ringed-topoi}
设 $(f,f')$ 是情形 \ref{defos-situation-morphism-thickenings-ringed-topoi}
中的一阶增厚态射。设 $\mathcal{F}'$、$\mathcal{G}'$ 是
$\mathcal{O}'$-模，并令
$\mathcal{F}=i^*\mathcal{F}'$、$\mathcal{G}=i^*\mathcal{G}'$。
设 $\varphi:\mathcal{F}\to\mathcal{G}$ 是 $\mathcal{O}$-线性同态。
假设 $\mathcal{F}'$、$\mathcal{G}'$ 在
$\mathcal{O}_{\mathcal{B}'}$ 上平坦，且 $(f,f')$ 是增厚的严格态射。
存在元素
$$
o(\varphi) \in
\Ext^1_\mathcal{O}(\mathcal{F},
\mathcal{G} \otimes_\mathcal{O} f^*\mathcal{J})
$$
其消失是存在把 $\varphi$ 提升为 $\mathcal{O}'$-线性同态
$\varphi':\mathcal{F}'\to\mathcal{G}'$ 的提升的充要条件。
\end{lemma}

\begin{proof}[第一种证明]
这由引理 \ref{defos-lemma-inf-obs-map-special-ringed-topoi} 得出，
因为我们断言，在本引理的假设下有
$$
\Ext^1_\mathcal{O}(Li^*\mathcal{F}', \mathcal{I}\mathcal{G}') =
\Ext^1_\mathcal{O}(\mathcal{F},
\mathcal{G} \otimes_\mathcal{O} f^*\mathcal{J})
$$
事实上，由引理 \ref{defos-lemma-deform-module-ringed-topoi}，有
$\mathcal{I}\mathcal{G}' =
\mathcal{G}\otimes_\mathcal{O}f^*\mathcal{J}$。
另一方面，注意
$$
H^{-1}(Li^*\mathcal{F}') =
\text{Tor}_1^{\mathcal{O}'}(\mathcal{F}', \mathcal{O})
$$
（局部计算从略）。利用短正合列
$$
0 \to \mathcal{I} \to \mathcal{O}' \to \mathcal{O} \to 0
$$
可见，此 $\text{Tor}_1$ 由同态
$\mathcal{I} \otimes_\mathcal{O} \mathcal{F} \to \mathcal{I}\mathcal{F}'$
的核计算；由引理 \ref{defos-lemma-deform-module-ringed-topoi}
的最后一个断言，该核为零。因此
$\tau_{\geq-1}Li^*\mathcal{F}'=\mathcal{F}$。另一方面，
由《导出范畴》引理 \ref{derived-lemma-negative-vanishing} 的对偶形式，有
$$
\Ext^1_\mathcal{O}(Li^*\mathcal{F}',
\mathcal{I}\mathcal{G}') =
\Ext^1_\mathcal{O}(\tau_{\geq -1}Li^*\mathcal{F}',
\mathcal{I}\mathcal{G}')
$$
\end{proof}

\begin{proof}[第二种证明]
可如下应用引理 \ref{defos-lemma-inf-obs-map-ringed-topoi}。注意，由引理
\ref{defos-lemma-deform-module-ringed-topoi}，有
$\mathcal{K}=\mathcal{I}\otimes_\mathcal{O}\mathcal{F}$、
$\mathcal{L} = \mathcal{I} \otimes_\mathcal{O} \mathcal{G}$
且 $c_{\mathcal{F}'}=1\otimes1$、
$c_{\mathcal{G}'}=1\otimes1$。取 $\psi=1\otimes\varphi$，
则该引理中的图交换。因此取
$o(\varphi)=o(\varphi,1\otimes\varphi)$ 即可。
\end{proof}

\begin{lemma}
\label{defos-lemma-inf-ext-rel-ringed-topoi}
设 $(f,f')$ 是情形 \ref{defos-situation-morphism-thickenings-ringed-topoi}
中的一阶增厚态射。设 $\mathcal{F}$ 是 $\mathcal{O}$-模。
假设 $(f,f')$ 是增厚的严格态射，且 $\mathcal{F}$ 在
$\mathcal{O}_\mathcal{B}$ 上平坦。若存在二元组
$(\mathcal{F}',\alpha)$，其中 $\mathcal{F}'$ 是在
$\mathcal{O}_{\mathcal{B}'}$ 上平坦的 $\mathcal{O}'$-模，且
$\alpha:i^*\mathcal{F}'\to\mathcal{F}$ 是同构，
则这类二元组的同构类集合是下列群作用下的主齐性空间：
$\Ext^1_\mathcal{O}(
\mathcal{F}, \mathcal{I} \otimes_\mathcal{O} \mathcal{F})$.
\end{lemma}

\begin{proof}
若假设存在一个这样的模，则由引理
\ref{defos-lemma-deform-module-ringed-topoi}，典范同态
$$
f^*\mathcal{J} \otimes_\mathcal{O} \mathcal{F} \to
\mathcal{I} \otimes_\mathcal{O} \mathcal{F}
$$
是同构。应用引理 \ref{defos-lemma-inf-ext-ringed-topoi}，其中
$\mathcal{K}=\mathcal{I}\otimes_\mathcal{O}\mathcal{F}$ 且 $c=1$。
由引理 \ref{defos-lemma-deform-module-ringed-topoi}，相应扩张
$\mathcal{F}'$ 均在 $\mathcal{O}_{\mathcal{B}'}$ 上平坦。
\end{proof}

\begin{lemma}
\label{defos-lemma-inf-obs-ext-rel-ringed-topoi}
设 $(f,f')$ 是情形 \ref{defos-situation-morphism-thickenings-ringed-topoi}
中的一阶增厚态射。设 $\mathcal{F}$ 是 $\mathcal{O}$-模。
假设 $(f,f')$ 是增厚的严格态射，且 $\mathcal{F}$ 在
$\mathcal{O}_\mathcal{B}$ 上平坦。存在在
$\mathcal{O}_{\mathcal{B}'}$ 上平坦的 $\mathcal{O}'$-模
$\mathcal{F}'$ 满足 $i^*\mathcal{F}'\cong\mathcal{F}$，
当且仅当
\begin{enumerate}
\item 典范同态
$f^*\mathcal{J} \otimes_\mathcal{O} \mathcal{F} \to
\mathcal{I} \otimes_\mathcal{O} \mathcal{F}$
是同构；并且
\item 引理 \ref{defos-lemma-inf-obs-ext-ringed-topoi} 中的类
$o(\mathcal{F}, \mathcal{I} \otimes_\mathcal{O} \mathcal{F}, 1)
\in \Ext^2_\mathcal{O}(
\mathcal{F}, \mathcal{I} \otimes_\mathcal{O} \mathcal{F})$
为零。
\end{enumerate}
\end{lemma}

\begin{proof}
这立即由引理 \ref{defos-lemma-deform-module-ringed-topoi}
对在 $\mathcal{O}_{\mathcal{B}'}$ 上平坦的 $\mathcal{O}'$-模的刻画，
以及引理 \ref{defos-lemma-inf-obs-ext-ringed-topoi} 得出。
\end{proof}






\section{在环化拓扑斯平坦增厚上的平坦模中的应用}
\label{defos-section-flat-ringed-topoi}

\noindent
考虑环化拓扑斯的交换图
$$
\xymatrix{
(\Sh(\mathcal{C}), \mathcal{O}) \ar[r]_i \ar[d]_f &
(\Sh(\mathcal{D}), \mathcal{O}') \ar[d]^{f'} \\
(\Sh(\mathcal{B}), \mathcal{O}_\mathcal{B}) \ar[r]^t &
(\Sh(\mathcal{B}'), \mathcal{O}_{\mathcal{B}'})
}
$$
其水平箭头是情形 \ref{defos-situation-morphism-thickenings-ringed-topoi}
中的一阶增厚。令
$\mathcal{I}=\Ker(i^\sharp)\subset\mathcal{O}'$，并令
$\mathcal{J}=\Ker(t^\sharp)\subset\mathcal{O}_{\mathcal{B}'}$。
设 $\mathcal{F}$ 是 $\mathcal{O}$-模。假设
\begin{enumerate}
\item $(f,f')$ 是增厚的严格态射；
\item $f'$ 平坦；
\item $\mathcal{F}$ 在 $\mathcal{O}_\mathcal{B}$ 上平坦。
\end{enumerate}
注意，(1) $+$ (2) 蕴含 $\mathcal{I}=f^*\mathcal{J}$
（把引理 \ref{defos-lemma-deform-module-ringed-topoi} 应用于
$\mathcal{O}'$）。在这些假设下，上一节的理论格外简洁。
下列引理汇总已经得到的结果。

\begin{lemma}
\label{defos-lemma-flat-ringed-topoi}
在上述情形中：
\begin{enumerate}
\item 存在在 $\mathcal{O}_{\mathcal{B}'}$ 上平坦且满足
$i^*\mathcal{F}'\cong\mathcal{F}$ 的 $\mathcal{O}'$-模
$\mathcal{F}'$，当且仅当引理
\ref{defos-lemma-inf-obs-ext-ringed-topoi} 中的类
$o(\mathcal{F}, f^*\mathcal{J} \otimes_\mathcal{O} \mathcal{F}, 1)
\in \Ext^2_\mathcal{O}(
\mathcal{F}, f^*\mathcal{J} \otimes_\mathcal{O} \mathcal{F})$
为零。
\item 若这样的模存在，则提升的同构类集合是下列群作用下的主齐性空间：
$\Ext^1_\mathcal{O}(
\mathcal{F}, f^*\mathcal{J} \otimes_\mathcal{O} \mathcal{F})$.
\item 给定提升 $\mathcal{F}'$，$\mathcal{F}'$ 的所有拉回为
$\text{id}_\mathcal{F}$ 的自同构所成集合典范同构于
$\Ext^0_\mathcal{O}(
\mathcal{F}, f^*\mathcal{J} \otimes_\mathcal{O} \mathcal{F})$.
\end{enumerate}
\end{lemma}

\begin{proof}
由上面已见的 $\mathcal{I}=f^*\mathcal{J}$，(1) 由引理
\ref{defos-lemma-inf-obs-ext-rel-ringed-topoi} 得出。
(2) 由引理 \ref{defos-lemma-inf-ext-rel-ringed-topoi} 得出。
(3) 由引理 \ref{defos-lemma-inf-map-rel-ringed-topoi} 得出。
\end{proof}

\begin{situation}
\label{defos-situation-morphism-flat-thickenings-ringed-topoi}
设 $f:(\Sh(\mathcal{C}),\mathcal{O})\to
(\Sh(\mathcal{B}),\mathcal{O}_\mathcal{B})$ 是环化拓扑斯态射。
考虑交换图
$$
\xymatrix{
(\Sh(\mathcal{C}'_1), \mathcal{O}'_1) \ar[r]_h \ar[d]_{f'_1} &
(\Sh(\mathcal{C}'_2), \mathcal{O}'_2) \ar[d]_{f'_2} \\
(\Sh(\mathcal{B}'_1), \mathcal{O}_{\mathcal{B}'_1}) \ar[r] &
(\Sh(\mathcal{B}'_2), \mathcal{O}_{\mathcal{B}'_2})
}
$$
其中 $h$ 是 $(\Sh(\mathcal{C}),\mathcal{O})$ 的一阶增厚之间的态射，
下方水平箭头是 $(\Sh(\mathcal{B}),\mathcal{O}_\mathcal{B})$
的一阶增厚之间的态射，每个 $f'_i$ 都限制为 $f$，
两对 $(f,f_i')$ 都是增厚的严格态射，且两个 $f'_i$ 都平坦。
最后，设 $\mathcal{F}$ 是在 $\mathcal{O}_\mathcal{B}$ 上平坦的
$\mathcal{O}$-模。
\end{situation}

\begin{lemma}
\label{defos-lemma-functorial-ringed-topoi}
在情形 \ref{defos-situation-morphism-flat-thickenings-ringed-topoi} 中，
障碍类
$o(\mathcal{F}, f^*\mathcal{J}_2 \otimes_\mathcal{O} \mathcal{F}, 1)$
在典范同态
$$
\Ext^2_\mathcal{O}(
\mathcal{F}, f^*\mathcal{J}_2 \otimes_\mathcal{O} \mathcal{F})
\to \Ext^2_\mathcal{O}(
\mathcal{F}, f^*\mathcal{J}_1 \otimes_\mathcal{O} \mathcal{F})
$$
下映到障碍类
$o(\mathcal{F}, f^*\mathcal{J}_1 \otimes_\mathcal{O} \mathcal{F}, 1)$
。
\end{lemma}

\begin{proof}
由注 \ref{defos-remark-obstruction-extension-functorial-ringed-topoi} 得出。
\end{proof}

\begin{situation}
\label{defos-situation-ses-flat-thickenings-ringed-topoi}
设 $f:(\Sh(\mathcal{C}),\mathcal{O})\to
(\Sh(\mathcal{B}),\mathcal{O}_\mathcal{B})$ 是环化拓扑斯态射。
考虑交换图
$$
\xymatrix{
(\Sh(\mathcal{C}'_1), \mathcal{O}'_1) \ar[r]_h \ar[d]_{f'_1} &
(\Sh(\mathcal{C}'_2), \mathcal{O}'_2) \ar[r] \ar[d]_{f'_2} &
(\Sh(\mathcal{C}'_3), \mathcal{O}'_3) \ar[d]_{f'_3} \\
(\Sh(\mathcal{B}'_1), \mathcal{O}_{\mathcal{B}'_1}) \ar[r] &
(\Sh(\mathcal{B}'_2), \mathcal{O}_{\mathcal{B}'_2}) \ar[r] &
(\Sh(\mathcal{B}'_3), \mathcal{O}_{\mathcal{B}'_3})
}
$$
其中：(a) 上行是 $(\Sh(\mathcal{C}),\mathcal{O})$
的一阶增厚短正合列；(b) 下行是
$(\Sh(\mathcal{B}),\mathcal{O}_\mathcal{B})$ 的一阶增厚短正合列；
(c) 每个 $f'_i$ 都限制为 $f$；(d) 每一对 $(f,f_i')$
都是增厚的严格态射；(e) 每个 $f'_i$ 都平坦。最后，设
$\mathcal{F}'_2$ 是在 $\mathcal{O}_{\mathcal{B}'_2}$ 上平坦的
$\mathcal{O}'_2$-模，并令
$\mathcal{F}=\mathcal{F}'_2\otimes\mathcal{O}$。令
$\pi:(\Sh(\mathcal{C}'_1),\mathcal{O}'_1)\to
(\Sh(\mathcal{C}),\mathcal{O})$ 为典范分裂（注
\ref{defos-remark-short-exact-sequence-thickenings-ringed-topoi}）。
\end{situation}

\begin{lemma}
\label{defos-lemma-verify-iv-ringed-topoi}
在情形 \ref{defos-situation-ses-flat-thickenings-ringed-topoi} 中，
$\pi^*\mathcal{F}$ 与 $h^*\mathcal{F}'_2$ 是在
$\mathcal{O}_{\mathcal{B}'_1}$ 上平坦的 $\mathcal{O}'_1$-模，
且在 $(\Sh(\mathcal{C}),\mathcal{O})$ 上都限制为 $\mathcal{F}$。
二者之差（引理 \ref{defos-lemma-flat-ringed-topoi}）是
$\Ext^1_\mathcal{O}(\mathcal{F},
f^*\mathcal{J}_1 \otimes_\mathcal{O} \mathcal{F})$
中的元素 $\theta$，它在
$\Ext^2_\mathcal{O}(\mathcal{F},
f^*\mathcal{J}_3 \otimes_\mathcal{O} \mathcal{F})$
中的边界，等于把 $\mathcal{F}$ 提升为在
$\mathcal{O}_{\mathcal{B}'_3}$ 上平坦的 $\mathcal{O}'_3$-模的障碍
（引理 \ref{defos-lemma-flat-ringed-topoi}）。
\end{lemma}

\begin{proof}
注意，$\pi^*\mathcal{F}$ 与 $h^*\mathcal{F}'_2$ 在
$(\Sh(\mathcal{C}),\mathcal{O})$ 上都限制为 $\mathcal{F}$，
而 $\pi^*\mathcal{F}\to\mathcal{F}$ 与
$h^*\mathcal{F}'_2\to\mathcal{F}$ 的核均由
$f^*\mathcal{J}_1\otimes_\mathcal{O}\mathcal{F}$ 给出。
因此由引理 \ref{defos-lemma-deform-module-ringed-topoi} 可得平坦性。
取边界是有意义的，因为模列
$$
0 \to f^*\mathcal{J}_3 \otimes_\mathcal{O} \mathcal{F} \to
f^*\mathcal{J}_2 \otimes_\mathcal{O} \mathcal{F} \to
f^*\mathcal{J}_1 \otimes_\mathcal{O} \mathcal{F} \to 0
$$
由情形 \ref{defos-situation-ses-flat-thickenings-ringed-topoi} 的假设及
$\mathcal{F}$ 在 $\mathcal{O}_\mathcal{B}$ 上平坦这一事实而短正合。
关于障碍类的断言，正是把注
\ref{defos-remark-complex-thickenings-and-ses-modules-ringed-topoi}
的结果直接转译到这一特定情形。
\end{proof}







\section{环化拓扑斯的形变与朴素余切复形}
\label{defos-section-deformations-ringed-topoi}

\noindent
本节利用朴素余切复形研究一些形变理论。先取环化拓扑斯的一阶增厚
$t : (\Sh(\mathcal{B}), \mathcal{O}_\mathcal{B}) \to
(\Sh(\mathcal{B}'), \mathcal{O}_{\mathcal{B}'})$。
记 $\mathcal{J}=\Ker(t^\sharp)$，并等同 $\mathcal{B}$ 与
$\mathcal{B}'$ 的底层拓扑斯。再设给定环化拓扑斯态射
$f : (\Sh(\mathcal{C}), \mathcal{O}) \to
(\Sh(\mathcal{B}), \mathcal{O}_\mathcal{B})$、一个 $\mathcal{O}$-模
$\mathcal{G}$，以及 $f^{-1}\mathcal{O}_\mathcal{B}$-模层同态
$f^{-1}\mathcal{J}\to\mathcal{G}$。本节要问：能否找到下图问号处的对象
\begin{equation}
\label{defos-equation-to-solve-ringed-topoi}
\vcenter{
\xymatrix{
0 \ar[r] & \mathcal{G} \ar[r] & {?} \ar[r] & \mathcal{O} \ar[r] & 0 \\
0 \ar[r] & f^{-1}\mathcal{J} \ar[u]^c \ar[r] &
f^{-1}\mathcal{O}_{\mathcal{B}'} \ar[u] \ar[r] &
f^{-1}\mathcal{O}_\mathcal{B} \ar[u] \ar[r] & 0
}
}
\end{equation}
以及解（若存在）在何种意义下唯一。更确切地说，我们寻找一阶增厚
$i : (\Sh(\mathcal{C}), \mathcal{O}) \to (\Sh(\mathcal{C}'), \mathcal{O}')$
以及形如 (\ref{defos-equation-morphism-thickenings-ringed-topoi}) 的增厚态射
$(f,f')$，其中把 $\Ker(i^\sharp)$ 等同于 $\mathcal{G}$，
且 $(f')^\sharp$ 诱导给定同态 $c$。我们称
$(\Sh(\mathcal{C}'),\mathcal{O}')$ 是
(\ref{defos-equation-to-solve-ringed-topoi}) 的一个{\it 解}。


\begin{lemma}
\label{defos-lemma-huge-diagram-ringed-topoi}
设给定环化拓扑斯态射的交换图
\begin{equation}
\label{defos-equation-huge-1-ringed-topoi}
\vcenter{
\xymatrix{
& (\Sh(\mathcal{C}_2), \mathcal{O}_2) \ar[r]_{i_2} \ar[d]_{f_2} \ar[ddl]_g &
(\Sh(\mathcal{C}'_2), \mathcal{O}'_2) \ar[d]^{f'_2} \\
&
(\Sh(\mathcal{B}_2), \mathcal{O}_{\mathcal{B}_2}) \ar[r]^{t_2} \ar[ddl]|\hole &
(\Sh(\mathcal{B}'_2), \mathcal{O}_{\mathcal{B}'_2}) \ar[ddl] \\
(\Sh(\mathcal{C}_1), \mathcal{O}_1) \ar[r]_{i_1} \ar[d]_{f_1} &
(\Sh(\mathcal{C}'_1), \mathcal{O}'_1) \ar[d]^{f'_1} \\
(\Sh(\mathcal{B}_1), \mathcal{O}_{\mathcal{B}_1}) \ar[r]^{t_1} &
(\Sh(\mathcal{B}'_1), \mathcal{O}_{\mathcal{B}'_1})
}
}
\end{equation}
其水平箭头均为一阶增厚。令 $\mathcal{G}_j=\Ker(i_j^\sharp)$，
并设给定 $g^{-1}\mathcal{O}_1$-模同态
$\nu : g^{-1}\mathcal{G}_1 \to \mathcal{G}_2$
使得下图交换：
\begin{equation}
\label{defos-equation-huge-2-ringed-topoi}
\vcenter{
\xymatrix{
& 0 \ar[r] & \mathcal{G}_2 \ar[r] &
\mathcal{O}'_2 \ar[r] &
\mathcal{O}_2 \ar[r] & 0 \\
& 0 \ar[r]|\hole &
f_2^{-1}\mathcal{J}_2 \ar[u]_{c_2} \ar[r] &
f_2^{-1}\mathcal{O}_{\mathcal{B}'_2} \ar[u] \ar[r]|\hole &
f_2^{-1}\mathcal{O}_{\mathcal{B}_2} \ar[u] \ar[r] & 0 \\
0 \ar[r] &
\mathcal{G}_1 \ar[ruu] \ar[r] &
\mathcal{O}'_1 \ar[r] &
\mathcal{O}_1 \ar[ruu] \ar[r] & 0 \\
0 \ar[r] &
f_1^{-1}\mathcal{J}_1 \ar[ruu]|\hole \ar[u]^{c_1} \ar[r] &
f_1^{-1}\mathcal{O}_{\mathcal{B}'_1} \ar[ruu]|\hole \ar[u] \ar[r] &
f_1^{-1}\mathcal{O}_{\mathcal{B}_1} \ar[ruu]|\hole \ar[u] \ar[r] & 0
}
}
\end{equation}
其前后两个面都是 (\ref{defos-equation-to-solve-ringed-topoi}) 的解。
（西北偏北方向的箭头，是先对源施行 $g^{-1}$ 后得到的
$\mathcal{C}_2$ 上的同态。）
\begin{enumerate}
\item 存在典范元素
$\Ext^1_{\mathcal{O}_2}(
Lg^*\NL_{\mathcal{O}_1/\mathcal{O}_{\mathcal{B}_1}}, \mathcal{G}_2)$
其消失是存在环化拓扑斯态射
$(\Sh(\mathcal{C}'_2), \mathcal{O}'_2) \to
(\Sh(\mathcal{C}'_1), \mathcal{O}'_1)$ 使图
(\ref{defos-equation-huge-1-ringed-topoi})，并与 $\nu$ 相容的充要条件。
\item 若存在态射
$(\Sh(\mathcal{C}'_2), \mathcal{O}'_2) \to
(\Sh(\mathcal{C}'_1), \mathcal{O}'_1)$
使图 (\ref{defos-equation-huge-1-ringed-topoi}) 成立并与 $\nu$ 相容，
则所有这类态射所成集合是下列群作用下的主齐性空间：
$$
\Hom_{\mathcal{O}_1}(
\Omega_{\mathcal{O}_1/\mathcal{O}_{\mathcal{B}_1}}, g_*\mathcal{G}_2) =
\Hom_{\mathcal{O}_2}(
g^*\Omega_{\mathcal{O}_1/\mathcal{O}_{\mathcal{B}_1}}, \mathcal{G}_2) =
\Ext^0_{\mathcal{O}_2}(
Lg^*\NL_{\mathcal{O}_1/\mathcal{O}_{\mathcal{B}_1}}, \mathcal{G}_2).
$$
\end{enumerate}
\end{lemma}

\begin{proof}
本引理的证明与引理 \ref{defos-lemma-huge-diagram-ringed-spaces} 的证明相同。
我们建议读者阅读后者而非此处证明。我们将等同所涉每个增厚的底层拓扑斯
（陈述中已经采用这一约定）。引理最后一个断言中的等式由定义立即得出。
因此在余下的证明中，我们使用群
$\Ext^k_{\mathcal{O}_2}(
Lg^*\NL_{\mathcal{O}_1/\mathcal{O}_{\mathcal{B}_1}}, \mathcal{G}_2)$,
$k=0,1$。我们先说明，可归约到本引理所有环化拓扑斯的底层拓扑斯相同的情形。

\medskip\noindent
为此，注意
$g^{-1}\NL_{\mathcal{O}_1/\mathcal{O}_{\mathcal{B}_1}}$ 等于环层同态
$g^{-1}f_1^{-1}\mathcal{O}_{\mathcal{B}_1} \to g^{-1}\mathcal{O}_1$
的朴素余切复形，见《位点上的模》引理
\ref{sites-modules-lemma-pullback-differentials}。此外，
$\NL_{\mathcal{O}_1/\mathcal{O}_{\mathcal{B}_1}}$
的 $0$ 次项是平坦 $\mathcal{O}_1$-模，故典范同态
$$
Lg^*\NL_{\mathcal{O}_1/\mathcal{O}_{\mathcal{B}_1}}
\longrightarrow
g^{-1}\NL_{\mathcal{O}_1/\mathcal{O}_{\mathcal{B}_1}}
\otimes_{g^{-1}\mathcal{O}_1} \mathcal{O}_2
$$
在 $0$ 次与 $-1$ 次上同调层上诱导同构。因此可把引理中的
$\Ext$ 群替换为
$$
\Ext^k_{g^{-1}\mathcal{O}_1}(
g^{-1}\NL_{\mathcal{O}_1/\mathcal{O}_{\mathcal{B}_1}}, \mathcal{G}_2) =
\Ext^k_{g^{-1}\mathcal{O}_1}(
\NL_{g^{-1}\mathcal{O}_1/g^{-1}f_1^{-1}\mathcal{O}_{\mathcal{B}_1}},
\mathcal{G}_2)
$$
使图 (\ref{defos-equation-huge-1-ringed-topoi}) 成立并与 $\nu$ 相容的
环化拓扑斯态射
$(\Sh(\mathcal{C}'_2), \mathcal{O}'_2) \to
(\Sh(\mathcal{C}'_1), \mathcal{O}'_1)$ 所成集合，与所有同
$f^\sharp$、$\nu$ 相容的
$g^{-1}f_1^{-1}\mathcal{O}_{\mathcal{B}'_1}$-algebras
$g^{-1}\mathcal{O}'_1\to\mathcal{O}'_2$ 的代数同态集合一一对应。
由此可假设我们有位点 $\mathcal{C}$ 上层的图
(\ref{defos-equation-huge-2-ringed-topoi})（在底层拓扑斯上
$f_1=f_2=\text{id}$），并要寻找使该图成立的环层同态
$\mathcal{O}'_1\to\mathcal{O}'_2$。

\medskip\noindent
在本引理余下的证明中，假设所有底层拓扑空间相同；即有位点
$\mathcal{C}$ 上层的图 (\ref{defos-equation-huge-2-ringed-topoi})
（在底层拓扑斯上 $f_1=f_2=\text{id}$），并寻找使该图成立的环层同态
$\mathcal{O}'_1\to\mathcal{O}'_2$。我们使用的 $\Ext$ 群是
$\Ext^k_{\mathcal{O}_1}(
\NL_{\mathcal{O}_1/\mathcal{O}_{\mathcal{B}_1}}, \mathcal{G}_2)$, $k = 0, 1$.

\medskip\noindent
步骤 1：构造障碍类。考虑集合层
$$
\mathcal{E} = \mathcal{O}'_1 \times_{\mathcal{O}_2} \mathcal{O}'_2
$$
它带有满射 $\alpha:\mathcal{E}\to\mathcal{O}_1$，
故可用 $\NL(\alpha)$ 代替
$\NL_{\mathcal{O}_1/\mathcal{O}_{\mathcal{B}_1}}$，见
《位点上的模》引理 \ref{sites-modules-lemma-NL-up-to-qis}。令
$$
\mathcal{I}' =
\Ker(\mathcal{O}_{\mathcal{B}'_1}[\mathcal{E}] \to \mathcal{O}_1)
\quad\text{且}\quad
\mathcal{I} =
\Ker(\mathcal{O}_{\mathcal{B}_1}[\mathcal{E}] \to \mathcal{O}_1)
$$
存在满射 $\mathcal{I}'\to\mathcal{I}$，其核为
$\mathcal{J}_1\mathcal{O}_{\mathcal{B}'_1}[\mathcal{E}]$。
得到两个 $\mathcal{O}_{\mathcal{B}'_2}$-代数同态
$$
a : \mathcal{O}_{\mathcal{B}'_1}[\mathcal{E}] \to \mathcal{O}'_1
\quad\text{且}\quad
b : \mathcal{O}_{\mathcal{B}'_1}[\mathcal{E}] \to \mathcal{O}'_2
$$
它们诱导同态 $a|_{\mathcal{I}'}:\mathcal{I}'\to\mathcal{G}_1$ 与
$b|_{\mathcal{I}'}:\mathcal{I}'\to\mathcal{G}_2$。$a$、$b$ 均零化
$(\mathcal{I}')^2$。此外，由于
(\ref{defos-equation-huge-2-ringed-topoi}) 的左侧方块交换，
$a$、$b$ 作为取值于 $\mathcal{G}_2$ 的同态，在
$\mathcal{J}_1\mathcal{O}_{\mathcal{B}'_1}[\mathcal{E}]$
上一致。因此差
$b|_{\mathcal{I}'} - \nu \circ a|_{\mathcal{I}'}$
诱导良定义的 $\mathcal{O}_1$-线性同态
$$
\xi : \mathcal{I}/\mathcal{I}^2 \longrightarrow \mathcal{G}_2
$$
% P11-E0008: ill-typed reversed expression retained from frozen source; see control/ERRATA_CANDIDATES.jsonl.
它把 $\mathcal{I}$ 的局部截面 $f$ 之类映到
$a(f') - \nu(b(f'))$，
其中 $f'$ 是 $f$ 到 $\mathcal{I}'$ 的局部截面的一个提升。令
$[\xi] \in \Ext^1_{\mathcal{O}_1}(\NL(\alpha), \mathcal{G}_2)$
为其像（见下文）。

\medskip\noindent
步骤 2：$[\xi]$ 消失是必要的。记 $\Omega=
\Omega_{\mathcal{O}_{\mathcal{B}_1}[\mathcal{E}]/\mathcal{O}_{\mathcal{B}_1}}
\otimes_{\mathcal{O}_{\mathcal{B}_1}[\mathcal{E}]} \mathcal{O}_1$.
注意，$\NL(\alpha)=(\mathcal{I}/\mathcal{I}^2\to\Omega)$
嵌入可区别三角
$$
\Omega[0] \to
\NL(\alpha) \to
\mathcal{I}/\mathcal{I}^2[1] \to
\Omega[1]
$$
因此 $[\xi]$ 为零，当且仅当对于某个同态
$\Omega\to\mathcal{G}_2$，$\xi$ 是复合
$\mathcal{I}/\mathcal{I}^2\to\Omega\to\mathcal{G}_2$。
假设存在使图 (\ref{defos-equation-huge-2-ringed-topoi}) 成立的环层同态
$\varphi : \mathcal{O}'_1 \to \mathcal{O}'_2$
。在此情形，考虑同态
$\mathcal{O}'_1[\mathcal{E}] \to \mathcal{G}_2$,
$f'\mapsto b(f')-\varphi(a(f'))$。计算表明，它零化
$\mathcal{J}_1\mathcal{O}_{\mathcal{B}'_1}[\mathcal{E}]$，
并诱导导子 $\mathcal{O}_{\mathcal{B}_1}[\mathcal{E}]\to\mathcal{G}_2$。
所得线性同态 $\Omega\to\mathcal{G}_2$ 说明此时 $[\xi]=0$。

\medskip\noindent
步骤 3：$[\xi]$ 消失是充分的。设
$\theta:\Omega\to\mathcal{G}_2$ 是 $\mathcal{O}_1$-线性同态，
使 $\xi=\theta\circ(\mathcal{I}/\mathcal{I}^2\to\Omega)$。
计算表明
$$
b + \theta \circ d :
\mathcal{O}_{\mathcal{B}'_1}[\mathcal{E}]
\longrightarrow
\mathcal{O}'_2
$$
零化 $\mathcal{I}'$，因而定义使图
(\ref{defos-equation-huge-2-ringed-topoi}) 成立的同态
$\mathcal{O}'_1\to\mathcal{O}'_2$。

\medskip\noindent
上述特殊情形中 (2) 的证明从略。提示：这与引理
\ref{defos-lemma-huge-diagram} 中 (2) 的证明完全相同。
\end{proof}

\begin{lemma}
\label{defos-lemma-NL-represent-ext-class-ringed-topoi}
设 $\mathcal{C}$ 是位点，$\mathcal{A}\to\mathcal{B}$ 是
$\mathcal{C}$ 上的环层同态，$\mathcal{G}$ 是 $\mathcal{B}$-模。
设 $\xi\in\Ext^1_\mathcal{B}
(\NL_{\mathcal{B}/\mathcal{A}},\mathcal{G})$。
存在集合层同态 $\alpha:\mathcal{E}\to\mathcal{B}$，使
$\xi\in\Ext^1_\mathcal{B}(\NL(\alpha),\mathcal{G})$
是某个同态 $\mathcal{I}/\mathcal{I}^2\to\mathcal{G}$ 的类
（记号见证明）。
\end{lemma}

\begin{proof}
回忆：给定 $\alpha:\mathcal{E}\to\mathcal{B}$，使
$\mathcal{A}[\mathcal{E}]\to\mathcal{B}$ 是核为 $\mathcal{I}$ 的满射，
则复形
$\NL(\alpha) = (\mathcal{I}/\mathcal{I}^2 \to 
\Omega_{\mathcal{A}[\mathcal{E}]/\mathcal{A}}
\otimes_{\mathcal{A}[\mathcal{E}]} \mathcal{B})$
典范同构于 $\NL_{\mathcal{B}/\mathcal{A}}$，见《位点上的模》引理
\ref{sites-modules-lemma-NL-up-to-qis}。还要注意，
$\Omega = \Omega_{\mathcal{A}[\mathcal{E}]/\mathcal{A}}
\otimes_{\mathcal{A}[\mathcal{E}]} \mathcal{B}$ 是预层
$U \mapsto \bigoplus_{e \in \mathcal{E}(U)} \mathcal{B}(U)$.
的伴随层。换言之，$\Omega$ 是集合层 $\mathcal{E}$ 上的自由
$\mathcal{B}$-模；特别地，存在典范同态 $\mathcal{E}\to\Omega$。

\medskip\noindent
现在取某个 $\mathcal{E}$（例如像朴素余切复形的定义那样取
$\mathcal{E}=\mathcal{B}$）。把 $\xi$ 写成同态
$\mathcal{I}/\mathcal{I}^2\to\mathcal{G}$ 之类的障碍，是
$\Ext^1_\mathcal{B}(\Omega,\mathcal{G})$ 中的元素。
设它由 $\mathcal{B}$-模扩张
$0\to\mathcal{G}\to\mathcal{H}\to\Omega\to0$ 表示。
考虑集合层
$\mathcal{E}' = \mathcal{E} \times_\Omega \mathcal{H}$
它带有诱导同态 $\alpha':\mathcal{E}'\to\mathcal{B}$。令
$\mathcal{I}'=\Ker(\mathcal{A}[\mathcal{E}']\to\mathcal{B})$ 以及
$\Omega'=\Omega_{\mathcal{A}[\mathcal{E}']/\mathcal{A}}
\otimes_{\mathcal{A}[\mathcal{E}']}\mathcal{B}$。
$\xi$ 在拟同构 $\NL(\alpha')\to\NL(\alpha)$ 下的拉回，
在
$\Ext^1_\mathcal{B}(\Omega', \mathcal{G})$
中的像为零。这是因为，沿 $\Omega'\to\Omega$ 拉回扩张
$\mathcal{H}$ 所得扩张是分裂的：$\Omega'$ 是集合层
$\mathcal{E}'$ 上的自由 $\mathcal{B}$-模，而且按构造有交换图
$$
\xymatrix{
\mathcal{E}' \ar[r] \ar[d] & \mathcal{E} \ar[d] \\
\mathcal{H} \ar[r] & \Omega
}
$$
证明完毕。
\end{proof}

\begin{lemma}
\label{defos-lemma-choices-ringed-topoi}
若 (\ref{defos-equation-to-solve-ringed-topoi}) 存在解，则解的同构类集合是
$\Ext^1_\mathcal{O}
(\NL_{\mathcal{O}/\mathcal{O}_\mathcal{B}},\mathcal{G})$
作用下的主齐性空间。
\end{lemma}

\begin{proof}
立即注意到：给定 (\ref{defos-equation-to-solve-ringed-topoi}) 的两个解
$\mathcal{O}'_1$、$\mathcal{O}'_2$，由引理
\ref{defos-lemma-huge-diagram-ringed-topoi} 得到障碍元素
$o(\mathcal{O}'_1, \mathcal{O}'_2) \in \Ext^1_\mathcal{O}(
\NL_{\mathcal{O}/\mathcal{O}_\mathcal{B}}, \mathcal{G})$
它是存在同态 $\mathcal{O}'_1\to\mathcal{O}'_2$ 的障碍。
显然，此元素正是存在同构的障碍，因而区分同构类。
所以要完成证明，只需说明：给定一个解 $\mathcal{O}'$ 和元素
$\xi \in \Ext^1_\mathcal{O}(
\NL_{\mathcal{O}/\mathcal{O}_\mathcal{B}}, \mathcal{G})$
可找到第二个解 $\mathcal{O}'_\xi$，使
$o(\mathcal{O}',\mathcal{O}'_\xi)=\xi$。

\medskip\noindent
为类 $\xi$ 选取引理 \ref{defos-lemma-NL-represent-ext-class-ringed-topoi}
中的 $\alpha:\mathcal{E}\to\mathcal{O}$。考虑核为
$\mathcal{I}$ 的满射
$f^{-1}\mathcal{O}_\mathcal{B}[\mathcal{E}] \to \mathcal{O}$
及相应的朴素余切复形
$\NL(\alpha) = (\mathcal{I}/\mathcal{I}^2 \to
\Omega_{f^{-1}\mathcal{O}_\mathcal{B}[\mathcal{E}]/
f^{-1}\mathcal{O}_\mathcal{B}}
\otimes_{f^{-1}\mathcal{O}_\mathcal{B}[\mathcal{E}]} \mathcal{O})$.
由该引理，$\xi$ 是某个同态
$\delta:\mathcal{I}/\mathcal{I}^2\to\mathcal{G}$ 的类。
把 $\mathcal{E}$ 换成 $\mathcal{E}\times_\mathcal{O}\mathcal{O}'$
后，还可假设 $\alpha$ 分解通过同态
$\alpha':\mathcal{E}\to\mathcal{O}'$。

\medskip\noindent
这些选取决定 $f^{-1}\mathcal{O}_{\mathcal{B}'}$-代数同态
$\varphi:\mathcal{O}_{\mathcal{B}'}[\mathcal{E}]\to\mathcal{O}'$。
令 $\mathcal{I}'=\Ker(\varphi)$。注意，$\varphi$ 诱导同态
$\varphi|_{\mathcal{I}'}:\mathcal{I}'\to\mathcal{G}$，
且 $\mathcal{O}'$ 是下图中的推出：
$$
\xymatrix{
0 \ar[r] & \mathcal{G} \ar[r] & \mathcal{O}' \ar[r] &
\mathcal{O} \ar[r] & 0 \\
0 \ar[r] & \mathcal{I}' \ar[u]^{\varphi|_{\mathcal{I}'}} \ar[r] &
f^{-1}\mathcal{O}_{\mathcal{B}'}[\mathcal{E}] \ar[u] \ar[r] &
\mathcal{O} \ar[u]_{=} \ar[r] & 0
}
$$
令 $\psi:\mathcal{I}'\to\mathcal{G}$ 为同态
$\varphi|_{\mathcal{I}'}$ 与复合
$$
\mathcal{I}' \to \mathcal{I}'/(\mathcal{I}')^2 \to
\mathcal{I}/\mathcal{I}^2 \xrightarrow{\delta} \mathcal{G}.
$$
之和。沿 $\psi$ 取推出，得到另一个环扩张
$\mathcal{O}'_\xi$，它嵌入如上图。计算（从略）表明
$o(\mathcal{O}',\mathcal{O}'_\xi)=\xi$，正如所需。
\end{proof}

\begin{lemma}
\label{defos-lemma-extensions-of-relative-ringed-topoi}
设 $f:(\Sh(\mathcal{C}),\mathcal{O})\to
(\Sh(\mathcal{B}),\mathcal{O}_\mathcal{B})$ 是环化拓扑斯态射，
$\mathcal{G}$ 是 $\mathcal{O}$-模。所有
$f^{-1}\mathcal{O}_\mathcal{B}$-代数扩张
$$
0 \to \mathcal{G} \to \mathcal{O}' \to \mathcal{O} \to 0
$$
中 $\mathcal{G}$ 为平方零理想者的同构类集合\footnote{换言之，即
$(\Sh(\mathcal{B}),\mathcal{O}_\mathcal{B})$ 上一阶增厚
$i : (\Sh(\mathcal{C}), \mathcal{O}) \to (\Sh(\mathcal{C}), \mathcal{O}')$
连同 $\mathcal{O}$-模同构
$\mathcal{G}\to\Ker(i^\sharp)$ 的同构类集合。}，
与
$\Ext^1_\mathcal{O}(\NL_{\mathcal{O}/\mathcal{O}_\mathcal{B}}, \mathcal{G})$.
典范双射。
\end{lemma}

\begin{proof}
为证明此结论，把前述结果应用于由下图给出
(\ref{defos-equation-to-solve-ringed-topoi}) 的情形：
$$
\xymatrix{
0 \ar[r] &
\mathcal{G} \ar[r] &
{?} \ar[r] &
\mathcal{O} \ar[r] & 0 \\
0 \ar[r] &
0 \ar[u] \ar[r] &
f^{-1}\mathcal{O}_\mathcal{B} \ar[u] \ar[r]^{\text{id}} &
f^{-1}\mathcal{O}_\mathcal{B} \ar[u] \ar[r] & 0
}
$$
于是本引理由引理 \ref{defos-lemma-choices-ringed-topoi} 及解
$\mathcal{G}\oplus\mathcal{O}$ 的存在性得出。
（该双射的直接构造见下列注。）
\end{proof}

\begin{remark}
\label{defos-remark-extensions-of-relative-ringed-topoi}
% P11-E0009: missing Sh in frozen source codomain retained; see control/ERRATA_CANDIDATES.jsonl.
设 $f:(\Sh(\mathcal{C}),\mathcal{O})\to
(\mathcal{B},\mathcal{O}_\mathcal{B})$ 与 $\mathcal{G}$ 如引理
\ref{defos-lemma-extensions-of-relative-ringed-topoi}。考虑该引理中的扩张
$0 \to \mathcal{G} \to \mathcal{O}' \to \mathcal{O} \to 0$
。可以选取集合层 $\mathcal{E}$ 及交换图
$$
\xymatrix{
\mathcal{E} \ar[d]_{\alpha'} \ar[rd]^\alpha \\
\mathcal{O}' \ar[r] & \mathcal{O}
}
$$
使 $f^{-1}\mathcal{O}_\mathcal{B}[\mathcal{E}]\to\mathcal{O}$
是核为 $\mathcal{J}$ 的满射。（例如可取任意满射到
$\mathcal{O}'$ 的集合层。）于是
$$
\NL_{\mathcal{O}/\mathcal{O}_\mathcal{B}} \cong \NL(\alpha) = 
\left(
\mathcal{J}/\mathcal{J}^2
\longrightarrow
\Omega_{f^{-1}\mathcal{O}_\mathcal{B}[\mathcal{E}]/
f^{-1}\mathcal{O}_\mathcal{B}}
\otimes_{f^{-1}\mathcal{O}_\mathcal{B}[\mathcal{E}]} \mathcal{O}\right)
$$
见《位点上的模》第 \ref{sites-modules-section-netherlander} 节，
尤其是引理 \ref{sites-modules-lemma-NL-up-to-qis}。
当然，$\alpha'$ 决定同态
$f^{-1}\mathcal{O}_\mathcal{B}[\mathcal{E}]\to\mathcal{O}'$，
进而决定同态
$$
\mathcal{J}/\mathcal{J}^2 \longrightarrow \mathcal{G}
$$
后者又决定
$\Ext^1_\mathcal{O}(\NL(\alpha), \mathcal{G}) =
\Ext^1_\mathcal{O}(\NL_{\mathcal{O}/\mathcal{O}_\mathcal{B}}, \mathcal{G})$
中的元素；在本引理的双射下，该元素对应于 $\mathcal{O}'$。
\end{remark}

\begin{lemma}
\label{defos-lemma-extensions-of-relative-ringed-topoi-functorial}
设 $f:(\Sh(\mathcal{C}),\mathcal{O}_\mathcal{C})\to
(\Sh(\mathcal{B}),\mathcal{O}_\mathcal{B})$ 与
$g:(\Sh(\mathcal{D}),\mathcal{O}_\mathcal{D})\to
(\Sh(\mathcal{C}),\mathcal{O}_\mathcal{C})$ 是环化拓扑斯态射。
设 $\mathcal{F}$ 是 $\mathcal{O}_\mathcal{C}$-模，
$\mathcal{G}$ 是 $\mathcal{O}_\mathcal{D}$-模，且
$c:g^*\mathcal{F}\to\mathcal{G}$ 是 $\mathcal{O}_\mathcal{D}$-线性同态。
最后，考虑
\begin{enumerate}
\item[(a)]
$0 \to \mathcal{F} \to \mathcal{O}_{\mathcal{C}'} \to
\mathcal{O}_\mathcal{C} \to 0$
是对应于
$\xi\in\Ext^1_{\mathcal{O}_\mathcal{C}}
(\NL_{\mathcal{O}_\mathcal{C}/\mathcal{O}_\mathcal{B}},\mathcal{F})$
的 $f^{-1}\mathcal{O}_\mathcal{B}$-代数扩张；
\item[(b)]
$0 \to \mathcal{G} \to \mathcal{O}_{\mathcal{D}'} \to
\mathcal{O}_\mathcal{D} \to 0$
是对应于
$\zeta\in\Ext^1_{\mathcal{O}_\mathcal{D}}
(\NL_{\mathcal{O}_\mathcal{D}/\mathcal{O}_\mathcal{B}},\mathcal{G})$
的 $g^{-1}f^{-1}\mathcal{O}_\mathcal{B}$-代数扩张。
\end{enumerate}
见引理 \ref{defos-lemma-extensions-of-relative-ringed-topoi}。
存在环化拓扑斯态射
$$
g' :
(\Sh(\mathcal{D}), \mathcal{O}_{\mathcal{D}'})
\longrightarrow
(\Sh(\mathcal{C}), \mathcal{O}_{\mathcal{C}'})
$$
位于 $(\Sh(\mathcal{B}),\mathcal{O}_\mathcal{B})$ 上方且与
$g$、$c$ 相容，当且仅当 $\xi$、$\zeta$ 映到
$\Ext^1_{\mathcal{O}_\mathcal{D}}(
Lg^*\NL_{\mathcal{O}_\mathcal{C}/\mathcal{O}_\mathcal{B}}, \mathcal{G})$
中的同一元素。
\end{lemma}

\begin{proof}
该陈述有意义，因为利用同态
$Lg^*\mathcal{F}\to g^*\mathcal{F}\xrightarrow{c}\mathcal{G}$，
有同态
$$
\Ext^1_{\mathcal{O}_\mathcal{C}}(
\NL_{\mathcal{O}_\mathcal{C}/\mathcal{O}_\mathcal{B}}, \mathcal{F}) \to
\Ext^1_{\mathcal{O}_\mathcal{D}}(
Lg^*\NL_{\mathcal{O}_\mathcal{C}/\mathcal{O}_\mathcal{B}}, Lg^*\mathcal{F}) \to
\Ext^1_{\mathcal{O}_\mathcal{D}}
(Lg^*\NL_{\mathcal{O}_\mathcal{C}/\mathcal{O}_\mathcal{B}}, \mathcal{G})
$$
并且，利用同态
$Lg^*\NL_{\mathcal{O}_\mathcal{C}/\mathcal{O}_\mathcal{B}}\to
\NL_{\mathcal{O}_\mathcal{D}/\mathcal{O}_\mathcal{B}}$，有同态
% P11-E0010: undefined O_Y coefficient rings retained from frozen source; see control/ERRATA_CANDIDATES.jsonl.
$$
\Ext^1_{\mathcal{O}_Y}(
\NL_{\mathcal{O}_\mathcal{D}/\mathcal{O}_\mathcal{B}}, \mathcal{G}) \to
\Ext^1_{\mathcal{O}_Y}(
Lg^*\NL_{\mathcal{O}_\mathcal{C}/\mathcal{O}_\mathcal{B}}, \mathcal{G})
$$
把引理 \ref{defos-lemma-huge-diagram-ringed-topoi} 应用于图
$$
\xymatrix{
& 0 \ar[r] &
\mathcal{G} \ar[r] &
\mathcal{O}_{\mathcal{D}'} \ar[r] &
\mathcal{O}_\mathcal{D} \ar[r] & 0 \\
& 0 \ar[r]|\hole & 0 \ar[u] \ar[r] &
g^{-1}f^{-1}\mathcal{O}_\mathcal{B} \ar[u] \ar[r]|\hole &
g^{-1}f^{-1}\mathcal{O}_\mathcal{B} \ar[u] \ar[r] & 0 \\
0 \ar[r] &
\mathcal{F} \ar[ruu] \ar[r] &
\mathcal{O}_{\mathcal{C}'} \ar[r] &
\mathcal{O}_\mathcal{C} \ar[ruu] \ar[r] & 0 \\
0 \ar[r] & 0 \ar[ruu]|\hole \ar[u] \ar[r] &
f^{-1}\mathcal{O}_\mathcal{B} \ar[ruu]|\hole \ar[u] \ar[r] &
f^{-1}\mathcal{O}_\mathcal{B} \ar[ruu]|\hole \ar[u] \ar[r] & 0
}
$$
再结合引理 \ref{defos-lemma-extensions-of-relative-ringed-topoi} 与
\ref{defos-lemma-huge-diagram-ringed-topoi} 的证明中两个构造之间的相容性，
即可推出本引理。该相容性的陈述与证明从略。
（直接论证见下列注。）
\end{proof}

\begin{remark}
\label{defos-remark-extensions-of-relative-ringed-topoi-functorial}
设 $f : (\Sh(\mathcal{C}), \mathcal{O}_\mathcal{C}) \to
(\Sh(\mathcal{B}), \mathcal{O}_\mathcal{B})$,
$g : (\Sh(\mathcal{D}), \mathcal{O}_\mathcal{D}) \to
(\Sh(\mathcal{C}), \mathcal{O}_\mathcal{C})$,
$\mathcal{F}$,
$\mathcal{G}$,
$c : g^*\mathcal{F} \to \mathcal{G}$,
$0 \to \mathcal{F} \to \mathcal{O}_{\mathcal{C}'} \to
\mathcal{O}_\mathcal{C} \to 0$,
$\xi \in \Ext^1_{\mathcal{O}_\mathcal{C}}(
\NL_{\mathcal{O}_\mathcal{C}/\mathcal{O}_\mathcal{B}}, \mathcal{F})$,
$0 \to \mathcal{G} \to \mathcal{O}_{\mathcal{D}'} \to
\mathcal{O}_\mathcal{D} \to 0$，以及
$\zeta \in \Ext^1_{\mathcal{O}_\mathcal{D}}(
\NL_{\mathcal{O}_\mathcal{D}/\mathcal{O}_\mathcal{B}}, \mathcal{G})$
如引理 \ref{defos-lemma-extensions-of-relative-ringed-topoi-functorial}。
沿 $c:g^{-1}\mathcal{F}\to\mathcal{G}$ 取推出，可构造扩张
$$
\xymatrix{
0 \ar[r] &
\mathcal{G} \ar[r] &
\mathcal{O}'_1 \ar[r] &
g^{-1}\mathcal{O}_\mathcal{C} \ar[r] & 0 \\
0 \ar[r] &
g^{-1}\mathcal{F} \ar[u]^c \ar[r] &
g^{-1}\mathcal{O}_{\mathcal{C}'} \ar[u] \ar[r] &
g^{-1}\mathcal{O}_\mathcal{C} \ar@{=}[u] \ar[r] & 0
}
$$
沿 $g^\sharp:g^{-1}\mathcal{O}_\mathcal{C}\to
\mathcal{O}_\mathcal{D}$ 取拉回，可构造扩张
$$
\xymatrix{
0 \ar[r] &
\mathcal{G} \ar[r] &
\mathcal{O}_{\mathcal{D}'} \ar[r] &
\mathcal{O}_\mathcal{D} \ar[r] & 0 \\
0 \ar[r] &
\mathcal{G} \ar@{=}[u] \ar[r] &
\mathcal{O}'_2 \ar[u] \ar[r] &
g^{-1}\mathcal{O}_\mathcal{C} \ar[u] \ar[r] & 0
}
$$
追图可知，存在态射
$g' : (\Sh(\mathcal{D}), \mathcal{O}_{\mathcal{D}'}) \to
(\Sh(\mathcal{C}), \mathcal{O}_{\mathcal{C}'})$
位于 $(\Sh(\mathcal{B}),\mathcal{O}_\mathcal{B})$ 上方且与
$g$、$c$ 相容，当且仅当 $\mathcal{O}'_1$ 与 $\mathcal{O}'_2$
作为以 $\mathcal{G}$ 扩张 $g^{-1}\mathcal{O}_\mathcal{C}$ 的
$g^{-1}f^{-1}\mathcal{O}_\mathcal{B}$-代数扩张同构。
由引理 \ref{defos-lemma-extensions-of-relative-ringed-topoi}，这些扩张由下式左端分类：
$$
\Ext^1_{g^{-1}\mathcal{O}_\mathcal{C}}(
\NL_{g^{-1}\mathcal{O}_\mathcal{C}/g^{-1}f^{-1}\mathcal{O}_\mathcal{B}},
\mathcal{G}) =
\Ext^1_{\mathcal{O}_\mathcal{D}}(
Lg^*\NL_{\mathcal{O}_\mathcal{C}/\mathcal{O}_\mathcal{B}}, \mathcal{G})
$$
其中等式来自张量--Hom 伴随及等式
$$
\NL_{g^{-1}\mathcal{O}_\mathcal{C}/g^{-1}f^{-1}\mathcal{O}_\mathcal{B}} =
g^{-1}\NL_{\mathcal{O}_\mathcal{C}/\mathcal{O}_\mathcal{B}}
\quad\text{且}\quad
Lg^*\NL_{\mathcal{O}_\mathcal{C}/\mathcal{O}_\mathcal{B}} =
g^{-1}\NL_{\mathcal{O}_\mathcal{C}/\mathcal{O}_\mathcal{B}}
% P11-E0010: undefined O_X/O_Y coefficients retained from frozen source; see control/ERRATA_CANDIDATES.jsonl.
\otimes_{g^{-1}\mathcal{O}_X}^\mathbf{L} \mathcal{O}_Y
$$
第一个等式见《位点上的模》引理
\ref{sites-modules-lemma-pullback-NL}；第二个等式由导出拉回的定义得出。
因此，要说明引理
\ref{defos-lemma-extensions-of-relative-ringed-topoi-functorial} 成立，
只需证明 $\mathcal{O}'_1$ 对应于 $\xi$ 的像，且
$\mathcal{O}'_2$ 对应于 $\zeta$ 的像。$\xi$ 与
$\mathcal{O}'_1$ 的对应关系，立即来自注
\ref{defos-remark-extensions-of-relative-ringed-topoi} 中类 $\xi$ 的构造。
为说明 $\zeta$ 与 $\mathcal{O}'_2$ 的对应关系，先选取交换图
$$
\xymatrix{
\mathcal{E} \ar[d]_{\beta'} \ar[rd]^\beta \\
\mathcal{O}_{\mathcal{D}'} \ar[r] & \mathcal{O}_\mathcal{D}
}
$$
使 $g^{-1}f^{-1}\mathcal{O}_\mathcal{B}[\mathcal{E}]\to
\mathcal{O}_\mathcal{D}$ 是核为 $\mathcal{K}$ 的满射。再选取交换图
$$
\xymatrix{
\mathcal{E} \ar[d]_{\beta'} &
\mathcal{E}' \ar[l]^\varphi \ar[d]_{\alpha'} \ar[rd]^\alpha \\
\mathcal{O}_{\mathcal{D}'} &
\mathcal{O}'_2 \ar[l] \ar[r] &
g^{-1}\mathcal{O}_\mathcal{C}
}
$$
使 $g^{-1}f^{-1}\mathcal{O}_\mathcal{B}[\mathcal{E}']\to
g^{-1}\mathcal{O}_\mathcal{C}$ 是核为 $\mathcal{J}$ 的满射。
（例如，作为集合层直接取
$\mathcal{E}'=\mathcal{E}\amalg\mathcal{O}'_2$。）同态 $\varphi$
诱导复形同态 $\NL(\alpha)\to\NL(\beta)$
（记号见《模》第 \ref{modules-section-netherlander} 节），特别诱导
$\bar\varphi : \mathcal{J}/\mathcal{J}^2 \to \mathcal{K}/\mathcal{K}^2$.
% P11-E0011: reversed alpha/beta identifications retained from frozen source; see control/ERRATA_CANDIDATES.jsonl.
于是
$\NL(\alpha)\cong
\NL_{\mathcal{O}_\mathcal{D}/\mathcal{O}_\mathcal{B}}$
且
$\NL(\beta) \cong
\NL_{g^{-1}\mathcal{O}_\mathcal{C}/g^{-1}f^{-1}\mathcal{O}_\mathcal{B}}$
，而复形同态 $\NL(\alpha)\to\NL(\beta)$ 表示同态
$Lg^*\NL_{\mathcal{O}_\mathcal{C}/\mathcal{O}_\mathcal{B}} \to
\NL_{\mathcal{O}_\mathcal{D}/\mathcal{O}_\mathcal{B}}$
，即引理
\ref{defos-lemma-extensions-of-relative-ringed-topoi-functorial}
的陈述中所用的同态（见其证明第一部分）。现在，$\zeta$ 对应于
$\beta'$ 诱导的同态 $\mathcal{K}/\mathcal{K}^2\to\mathcal{G}$ 之类，
见注 \ref{defos-remark-extensions-of-relative-ringed-topoi}。类似地，扩张
$\mathcal{O}'_2$ 对应于 $\alpha'$ 诱导的同态
$\mathcal{J}/\mathcal{J}^2\to\mathcal{G}$。上面的交换图表明，
此同态是 $\beta'$ 诱导的同态
$\mathcal{K}/\mathcal{K}^2\to\mathcal{G}$ 与同态
$\bar\varphi:\mathcal{J}/\mathcal{J}^2\to\mathcal{K}/\mathcal{K}^2$
的复合。这证明了所需的相容性。
\end{remark}

\begin{lemma}
\label{defos-lemma-parametrize-solutions-ringed-topoi}
设 $t : (\Sh(\mathcal{B}), \mathcal{O}_\mathcal{B}) \to
(\Sh(\mathcal{B}'), \mathcal{O}_{\mathcal{B}'})$,
$\mathcal{J} = \Ker(t^\sharp)$,
$f : (\Sh(\mathcal{C}), \mathcal{O}) \to
(\Sh(\mathcal{B}), \mathcal{O}_\mathcal{B})$、$\mathcal{G}$ 以及
% P11-E0012: missing f^{-1} on J retained from frozen source; see control/ERRATA_CANDIDATES.jsonl.
$c : \mathcal{J} \to \mathcal{G}$ 如
(\ref{defos-equation-to-solve-ringed-topoi})。以
$\xi \in \Ext^1_{\mathcal{O}_\mathcal{B}}(
\NL_{\mathcal{O}_\mathcal{B}/\mathcal{O}_{\mathcal{B}'}}, \mathcal{J})$
记经引理 \ref{defos-lemma-extensions-of-relative-ringed-topoi}
对应于以 $\mathcal{J}$ 扩张 $\mathcal{O}_\mathcal{B}$ 所得扩张
$\mathcal{O}_{\mathcal{B}'}$ 的元素。解的同构类集合典范双射于同态
$$
\Ext^1_\mathcal{O}(\NL_{\mathcal{O}/\mathcal{O}_{\mathcal{B}'}},
\mathcal{G})\to
\Ext^1_\mathcal{O}(
Lf^*\NL_{\mathcal{O}_\mathcal{B}/\mathcal{O}_{\mathcal{B}'}}, \mathcal{G})
$$
在 $\xi$ 的像上方的纤维。
\end{lemma}

\begin{proof}
把引理 \ref{defos-lemma-extensions-of-relative-ringed-topoi}
应用于 $t\circ f:(\Sh(\mathcal{C}),\mathcal{O})\to
(\Sh(\mathcal{B}'), \mathcal{O}_{\mathcal{B}'})$
及 $\mathcal{O}$-模 $\mathcal{G}$，可见
$\Ext^1_\mathcal{O}(\NL_{\mathcal{O}/\mathcal{O}_{\mathcal{B}'}},
\mathcal{G})$
中的元素 $\zeta$ 参数化
$0 \to \mathcal{G} \to \mathcal{O}' \to \mathcal{O} \to 0$
这样的 $f^{-1}\mathcal{O}_{\mathcal{B}'}$-代数扩张。
把引理 \ref{defos-lemma-extensions-of-relative-ringed-topoi-functorial} 应用于
$$
(\Sh(\mathcal{C}), \mathcal{O}) \xrightarrow{f}
(\Sh(\mathcal{B}), \mathcal{O}_\mathcal{B}) \xrightarrow{t}
(\Sh(\mathcal{B}'), \mathcal{O}_{\mathcal{B}'})
$$
及 $c:\mathcal{J}\to\mathcal{G}$，可见存在态射
$$
f' : 
(\Sh(\mathcal{C}), \mathcal{O}')
\longrightarrow
(\Sh(\mathcal{B}'), \mathcal{O}_{\mathcal{B}'})
$$
位于 $(\Sh(\mathcal{B}'),\mathcal{O}_{\mathcal{B}'})$ 上方且与
$c$、$f$ 相容，当且仅当 $\zeta$ 映到 $\xi$。
这当然等价于说 $\mathcal{O}'$ 是
(\ref{defos-equation-to-solve-ringed-topoi}) 的一个解。
\end{proof}






















\section{代数空间的形变}
\label{defos-section-deformations-spaces}

\noindent
本节具体说明第 \ref{defos-section-deformations-ringed-topoi} 节的结果
对于代数空间的形变意味着什么。

\begin{lemma}
\label{defos-lemma-match-thickenings}
设 $S$ 是概形，$i:Z\to Z'$ 是 $S$ 上代数空间的态射。
下列条件等价：
\begin{enumerate}
\item $i$ 是《空间态射进阶》第
\ref{spaces-more-morphisms-section-thickenings} 节所定义的代数空间增厚；
\item 相伴的环化拓扑斯态射
$i_{small} : (\Sh(Z_\etale), \mathcal{O}_Z) \to
(\Sh(Z'_\etale), \mathcal{O}_{Z'})$
（《空间的性质》引理
\ref{spaces-properties-lemma-morphism-ringed-topoi}）
是第 \ref{defos-section-thickenings-ringed-topoi} 节意义下的增厚。
\end{enumerate}
\end{lemma}

\begin{proof}
我们强调，这并非显然之事。

\medskip\noindent
假设 (1)。由《空间态射进阶》引理
\ref{spaces-more-morphisms-lemma-thickening-equivalence}，
态射 $i$ 诱导小 étale 位点之间的等价，因而特别诱导拓扑斯之间的等价。
当然，按增厚的定义，$i^\sharp$ 是核局部幂零的满射。

\medskip\noindent
假设 (2)。（此方向不那么重要，更多是一则趣闻。）对任意 étale 态射
$Y'\to Z'$，可见 $Y=Z\times_{Z'}Y'$ 与 $Y'$ 有相同的 étale 拓扑斯。
特别地，$Y'$ 拟紧当且仅当 $Y$ 拟紧，因为拟紧性是拓扑斯论概念
（《位点》引理 \ref{sites-lemma-quasi-compact}）。由此还可见，
$Y'$ 拟紧且拟分离，当且仅当 $Y$ 拟紧且拟分离
（因为可用下述条件刻画 $Y'$ 拟分离：对于所有在 $Y'$ 上 étale 的
拟紧代数空间 $Y'_1,Y'_2$，纤维积
$Y'_1\times_{Y'}Y'_2$ 拟紧）。取 $Y'$ 仿射。
则代数空间 $Y$ 拟紧且拟分离。对于任意拟凝聚
$\mathcal{O}_Y$-模 $\mathcal{F}$，有
$H^q(Y, \mathcal{F}) = H^q(Y', (Y \to Y')_*\mathcal{F})$
，因为二者的 étale 拓扑斯相同。于是
$H^q(Y',(Y\to Y')_*\mathcal{F})=0$，因为推出是拟凝聚的
（《空间的态射》引理 \ref{spaces-morphisms-lemma-pushforward}），
且 $Y$ 仿射。
% P11-E0013: circular/reversed affine-space claims retained from frozen source; see control/ERRATA_CANDIDATES.jsonl.
由《空间的上同调》命题
\ref{spaces-cohomology-proposition-vanishing-affine}
可知 $Y'$ 仿射（本引理这一方向肯定有避开该引用的证明）。
因此 $i$ 是仿射态射。在仿射情形，由第
\ref{defos-section-thickenings-ringed-topoi} 节的条件容易得出
$i$ 是代数空间的增厚。
\end{proof}

\begin{lemma}
\label{defos-lemma-deform-spaces}
设 $S$ 是概形，$Y\subset Y'$ 是 $S$ 上代数空间的一阶增厚，
$f:X\to Y$ 是 $S$ 上代数空间的平坦态射。若存在 $S$ 上代数空间的
平坦态射 $f':X'\to Y'$ 以及 $Y$ 上的同构
$a:X\to X'\times_{Y'}Y$，则
\begin{enumerate}
\item 二元组 $(f':X'\to Y',a)$ 的同构类集合是
$\Ext^1_{\mathcal{O}_X}(\NL_{X/Y},f^*\mathcal{C}_{Y/Y'})$
作用下的主齐性空间；
\item 所有 $Y'$ 上自同构 $\varphi:X'\to X'$ 中，在
$X'\times_{Y'}Y$ 上约化为恒等者所成集合是
$\Ext^0_{\mathcal{O}_X}(\NL_{X/Y},f^*\mathcal{C}_{Y/Y'})$。
\end{enumerate}
\end{lemma}

\begin{proof}
我们把环化拓扑斯形变的材料应用于本引理各代数空间的小 étale 拓扑斯。
可把 $X$ 视为 $X'$ 的闭子空间，使
$(f,f'):(X\subset X')\to(Y\subset Y')$ 是一阶增厚之间的态射。
由引理 \ref{defos-lemma-match-thickenings}，这转化为环化拓扑斯增厚之间的态射。
由《空间态射进阶》引理
\ref{spaces-more-morphisms-lemma-deform}
（或更一般的引理 \ref{defos-lemma-deform-module-ringed-topoi}）可见，
% P11-E0014: reversed conormal indices retained from frozen source prose; see control/ERRATA_CANDIDATES.jsonl.
$X$ 在 $X'$ 中的理想层等于 $f^*\mathcal{C}_{Y'/Y}$，
而这实际上等价于 $X'$ 在 $Y'$ 上平坦。因此有交换图
$$
\xymatrix{
0 \ar[r] & f^*\mathcal{C}_{Y/Y'} \ar[r] &
\mathcal{O}_{X'} \ar[r] &
\mathcal{O}_X \ar[r] & 0 \\
0 \ar[r] &
f_{small}^{-1}\mathcal{C}_{Y/Y'} \ar[u] \ar[r] &
f_{small}^{-1}\mathcal{O}_{Y'} \ar[u] \ar[r] &
f_{small}^{-1}\mathcal{O}_Y \ar[u] \ar[r] & 0
}
$$
请与 (\ref{defos-equation-to-solve-ringed-topoi}) 比较。注意，(2) 中的自同构
$\varphi$ 给出嵌入上图的自同构
$\varphi^\sharp:\mathcal{O}_{X'}\to\mathcal{O}_{X'}$。
反过来，由《空间态射进阶》引理
\ref{spaces-more-morphisms-lemma-first-order-thickening-maps}.
，任意嵌入上方层图的自同构
$\alpha:\mathcal{O}_{X'}\to\mathcal{O}_{X'}$，
都等于 (2) 中某个自同构 $\varphi$ 的 $\varphi^\sharp$。
最后，由《空间态射进阶》引理
\ref{spaces-more-morphisms-lemma-first-order-thickening}
，若找到 $X_\etale$ 上另一个嵌入下图的环层 $\mathcal{A}$
$$
\xymatrix{
0 \ar[r] & f^*\mathcal{C}_{Y/Y'} \ar[r] &
\mathcal{A} \ar[r] &
\mathcal{O}_X \ar[r] & 0 \\
0 \ar[r] &
f_{small}^{-1}\mathcal{C}_{Y/Y'} \ar[u] \ar[r] &
f_{small}^{-1}\mathcal{O}_{Y'} \ar[u] \ar[r] &
f_{small}^{-1}\mathcal{O}_Y \ar[u] \ar[r] & 0
}
$$
，则存在一阶增厚 $X\subset X''$，满足
$\mathcal{O}_{X''}=\mathcal{A}$；再次应用《空间态射进阶》引理
\ref{spaces-more-morphisms-lemma-first-order-thickening-maps}，
就得到具备所有所需性质的态射
$(f,f''):(X\subset X'')\to(Y\subset Y')$。
因此 (1) 由引理 \ref{defos-lemma-choices-ringed-topoi} 得出，
(2) 由引理 \ref{defos-lemma-huge-diagram-ringed-topoi} 的 (2) 得出。
（注意，《空间态射进阶》第
\ref{spaces-more-morphisms-section-netherlander}
节为代数空间态射定义的 $\NL_{X/Y}$，与第
\ref{defos-section-deformations-ringed-topoi} 节所用的 $\NL_{X/Y}$ 一致。）
\end{proof}

\noindent
设 $S$ 是概形，$f:X\to B$ 是 $S$ 上代数空间的态射。
设 $\mathcal{F}\to\mathcal{G}$ 是 $\mathcal{O}_X$-模同态
（未必拟凝聚）。考虑由推出给出的函子
$$
F :
\left\{
\begin{matrix}
f^{-1}\mathcal{O}_B\text{-代数扩张}\\
0 \to \mathcal{F} \to \mathcal{O}' \to \mathcal{O}_X \to 0\\
\text{其中 }\mathcal{F}\text{ 是平方零理想}
\end{matrix}
\right\}
\longrightarrow
\left\{
\begin{matrix}
f^{-1}\mathcal{O}_B\text{-代数扩张}\\
0 \to \mathcal{G} \to \mathcal{O}' \to \mathcal{O}_X \to 0\\
\text{其中 }\mathcal{G}\text{ 是平方零理想}
\end{matrix}
\right\}
$$

\begin{lemma}
\label{defos-lemma-thickening-space-quasi-coherent}
在上述情形中，假设 $X$ 拟紧且拟分离，并且
$DQ_X(\mathcal{F})\to DQ_X(\mathcal{G})$
（《空间的导出范畴》第
\ref{spaces-perfect-section-better-coherator} 节）
是同构。那么函子 $F$ 是范畴等价。
\end{lemma}

\begin{proof}
回忆，$\NL_{X/B}$ 是 $D_\QCoh(\mathcal{O}_X)$ 的对象，见
《空间态射进阶》引理
\ref{spaces-more-morphisms-lemma-netherlander-quasi-coherent}.
因此我们的假设蕴含对每个 $i$，同态
$$
\Ext^i_X(\NL_{X/B}, \mathcal{F}) \longrightarrow
\Ext^i_X(\NL_{X/B}, \mathcal{G})
$$
都是同构。由引理 \ref{defos-lemma-huge-diagram-ringed-topoi}，
这蕴含该函子全忠实。另一方面，由引理
\ref{defos-lemma-choices-ringed-topoi}，该函子本质满射，
因为两个范畴中分别有解
$\mathcal{O}_X\oplus\mathcal{F}$ 与
$\mathcal{O}_X\oplus\mathcal{G}$。
\end{proof}

\noindent
设 $S$ 是概形，$B\subset B'$ 是 $S$ 上代数空间的一阶增厚，
其理想层为 $\mathcal{J}$。可把 $\mathcal{J}$ 视为拟凝聚
$\mathcal{O}_B$-模，也可视为 $B'$ 上的拟凝聚理想层；见
《空间态射进阶》第 \ref{spaces-more-morphisms-section-thickenings} 节。
设 $f:X\to B$ 是 $S$ 上代数空间的态射，
$\mathcal{F}\to\mathcal{G}$ 是 $\mathcal{O}_X$-模同态
（未必拟凝聚）。设
$c:f^{-1}\mathcal{J}\to\mathcal{F}$ 是
$f^{-1}\mathcal{O}_B$-模同态，并以
$c':f^{-1}\mathcal{J}\to\mathcal{G}$ 记相应复合。
考虑由推出给出的函子
$$
FT :
\{(\ref{defos-equation-to-solve-ringed-topoi})
\text{ 关于 }\mathcal{F}\text{ 与 }c\text{ 的解}\}
\longrightarrow
\{(\ref{defos-equation-to-solve-ringed-topoi})
\text{ 关于 }\mathcal{G}\text{ 与 }c'\text{ 的解}\}
$$

\begin{lemma}
\label{defos-lemma-thickening-over-thickening-space-quasi-coherent}
在上述情形中，假设 $X$ 拟紧且拟分离，并且
$DQ_X(\mathcal{F})\to DQ_X(\mathcal{G})$
（《空间的导出范畴》第
\ref{spaces-perfect-section-better-coherator} 节）
是同构。那么函子 $FT$ 是范畴等价。
\end{lemma}

\begin{proof}
(\ref{defos-equation-to-solve-ringed-topoi}) 关于 $\mathcal{F}$ 的一个解，
特别给出 $f^{-1}\mathcal{O}_{B'}$-代数扩张
$$
0 \to \mathcal{F} \to \mathcal{O}' \to \mathcal{O}_X \to 0
$$
其中 $\mathcal{F}$ 是平方零理想。$\mathcal{G}$ 的情形类似。
此外，给定这样的扩张，得到同态
$c_{\mathcal{O}'}:f^{-1}\mathcal{J}\to\mathcal{F}$。
因此所考察的是这类 $f^{-1}\mathcal{O}_{B'}$-代数扩张中
满足 $c=c_{\mathcal{O}'}$ 者所成的全子范畴。显然，若
$\mathcal{O}''=F(\mathcal{O}')$，其中 $F$ 是引理
\ref{defos-lemma-thickening-space-quasi-coherent} 的等价
（这次应用于 $X\to B'$），则 $c_{\mathcal{O}''}$ 是
$c_{\mathcal{O}'}$ 与同态 $\mathcal{F}\to\mathcal{G}$ 的复合。
本引理得证。
\end{proof}





\section{复形的形变}
\label{defos-section-deformations-complexes}

\noindent
本节为下一节作预备。我们会尽量利用交换代数各章的材料。

\begin{lemma}
\label{defos-lemma-canonical-class-algebra}
设 $R'\to R$ 是环满射，其核为平方零理想 $I$。
对于每个 $K\in D^-(R)$，$D(R)$ 中存在典范同态
$$
\omega(K) : K \longrightarrow K \otimes_R^\mathbf{L} I[2]
$$
具有下列性质：
\begin{enumerate}
\item $\omega(K)=0$ 当且仅当存在 $K'\in D(R')$，满足
$K'\otimes_{R'}^\mathbf{L}R=K$；
\item 给定 $D^-(R)$ 中的同态 $K\to L$，下图交换：
$$
\xymatrix{
K \ar[d] \ar[rr]_-{\omega(K)} & &
K \otimes^\mathbf{L}_R I[2] \ar[d] \\
L \ar[rr]^-{\omega(L)} & &
L \otimes^\mathbf{L}_R I[2]
}
$$
\item $\omega(K)$ 的形成与环同态 $R'\to S'$ 相容
（准确陈述见证明）。
\end{enumerate}
\end{lemma}

\begin{proof}
取表示 $K$ 的有上界自由 $R$-模复形 $K^\bullet$。
可选取提升 $K^n$ 的自由 $R'$-模 $(K')^n$，并选取提升
$K^\bullet$ 的微分 $d^n_K:K^n\to K^{n+1}$ 的 $R'$-模同态
$(d')^n_K:(K')^n\to(K')^{n+1}$。复合
$$
(d')^{n + 1}_K \circ (d')^n_K : (K')^n \to (K')^{n + 2}
$$
虽未必为零，却确实分解为
$$
(K')^n \to K^n \xrightarrow{\omega^n_K}
K^{n + 2} \otimes_R I = I(K')^{n + 2} \to (K')^{n + 2}
$$
，因为 $d^{n+1}\circ d^n=0$。计算表明，$\omega^n_K$
定义复形同态。此复形同态定义 $\omega(K)$。

\medskip\noindent
下面证明此构造与有上界自由 $R$-模复形同态
$\alpha^\bullet:K^\bullet\to L^\bullet$ 及给定的提升选取
$(K')^n,(L')^n,(d')^n_K,(d')^n_L$ 相容。具体地，取提升分量
$\alpha^n:K^n\to L^n$ 的同态
$(\alpha')^n:(K')^n\to(L')^n$。与前面一样，
$(d')^n_L\circ(\alpha')^n-(\alpha')^{n+1}\circ(d')_K^n$
分解为
% P11-E0015: wrong terminal cohomological degree retained from frozen source; see control/ERRATA_CANDIDATES.jsonl.
$$
(K')^n \to K^n \xrightarrow{h^n}
L^{n + 1} \otimes_R I = I(L')^{n + 1} \to (L')^{n + 2}
$$
。容易计算得到
$$
\omega^n_L \circ \alpha^n =
(d_L^{n + 1} \otimes \text{id}_I) \circ h^n + h^{n + 1} \circ d_K^n +
(\alpha^{n + 2} \otimes \text{id}_I) \circ \omega^n_K
$$
这证明了本引理 (2) 中的图在此特殊情形下交换。
把它应用于表示 $K$ 的两种不同有上界自由复形选取，
可知 $\omega(K)$ 良定义；当然，(2) 在一般情形下也成立。

\medskip\noindent
若 $K$ 提升为 $D^-(R')$ 中的 $K'$，则可用有上界自由
$R'$-模复形表示 $K'$，立即可见 $\omega(K)=0$。
反过来，回到选取 $K^\bullet,(K')^n,(d')^n_K$；若
$\omega(K)=0$，则可找到同态
$g^n:K^n\to K^{n+1}\otimes_R I$，满足
$$
\omega^n = (d_K^{n + 1} \otimes \text{id}_I) \circ g^n +
g^{n + 1} \circ d_K^n
$$
这意味着以 $(d')^n_K-g^n:(K')^n\to(K')^{n+1}$ 为微分，
得到提升 $K^\bullet$ 的自由 $R'$-模复形。这证明了 (1)。

\medskip\noindent
最后，(3) 的含义如下。设 $R'\to S'$ 是环同态，令
$S=S'\otimes_{R'}R$，并以 $J=IS'\subset S'$ 记
$S'\to S$ 的平方零核。给定 $K\in D^-(R)$，所得下图交换：
$$
\xymatrix{
K \otimes_R^\mathbf{L} S \ar[d] \ar[rr]_-{\omega(K) \otimes \text{id}} & &
(K \otimes^\mathbf{L}_R I[2]) \otimes_R^\mathbf{L} S \ar[d] \\
K \otimes_R^\mathbf{L} S \ar[rr]^-{\omega(K \otimes_R^\mathbf{L} S)} & &
(K \otimes_R^\mathbf{L} S) \otimes^\mathbf{L}_S J[2]
}
$$
其中右侧竖直箭头来自
$$
(K \otimes^\mathbf{L}_R I[2]) \otimes_R^\mathbf{L} S =
(K \otimes_R^\mathbf{L} S) \otimes_S^\mathbf{L}
(I \otimes_R^\mathbf{L} S)[2] \longrightarrow
(K \otimes_R^\mathbf{L} S) \otimes_S^\mathbf{L} J[2]
$$
如上选取 $K^\bullet,(K')^n,(d')^n_K$。于是可用
$K^\bullet\otimes_R S$、$(K')^n\otimes_{R'}S'$ 以及
$(d')^n_K\otimes\text{id}_{S'}$ 构造
$\omega(K\otimes_R^\mathbf{L}S)$。
在这些选取下，立即可在复形同态层次验证交换性。
\end{proof}







\section{环化拓扑斯上复形的形变}
\label{defos-section-thickenings-complexes}

\noindent
本节材料取自 \cite{lieblich-complexes}。

\medskip\noindent
本节材料适用于第 \ref{defos-section-thickenings-ringed-topoi} 节所定义的
环化拓扑斯一阶增厚情形。不过，为简化记号，我们假设底层位点
$\mathcal{C}$ 与 $\mathcal{D}$ 相同。此外，把环层的满同态
$\mathcal{O}' \to \mathcal{O}$ 记作
$\mathcal{O} \to \mathcal{O}_0$；这种记法在文献中或许更常见。

\begin{lemma}
\label{defos-lemma-lift-complex}
设 $\mathcal{C}$ 是位点，$\mathcal{O} \to \mathcal{O}_0$
是环层的满射。假设给定以下数据：
\begin{enumerate}
\item 平坦 $\mathcal{O}$-模 $\mathcal{G}^n$；
\item $\mathcal{O}$-模同态 $\mathcal{G}^n \to \mathcal{G}^{n + 1}$；
\item $\mathcal{O}_0$-模复形 $\mathcal{K}_0^\bullet$；
\item $\mathcal{O}$-模同态 $\mathcal{G}^n \to \mathcal{K}_0^n$。
\end{enumerate}
并且满足：
\begin{enumerate}
\item[(a)] 当 $n \gg 0$ 时，$H^n(\mathcal{K}_0^\bullet) = 0$；
\item[(b)] 当 $n \gg 0$ 时，$\mathcal{G}^n = 0$；
\item[(c)] 令
$\mathcal{G}^n_0 = \mathcal{G}^n \otimes_\mathcal{O} \mathcal{O}_0$
，则诱导同态给出复形 $\mathcal{G}_0^\bullet$ 及复形同态
$\mathcal{G}_0^\bullet \to \mathcal{K}_0^\bullet$。
\end{enumerate}
则存在：
\begin{enumerate}
\item[(\romannumeral1)]
平坦 $\mathcal{O}$-模 $\mathcal{F}^n$；
\item[(\romannumeral2)]
$\mathcal{O}$-模同态 $\mathcal{F}^n \to \mathcal{F}^{n + 1}$；
\item[(\romannumeral3)]
$\mathcal{O}$-模同态 $\mathcal{F}^n \to \mathcal{K}_0^n$；
\item[(\romannumeral4)]
$\mathcal{O}$-模同态 $\mathcal{G}^n \to \mathcal{F}^n$。
\end{enumerate}
它们满足：当 $n \gg 0$ 时 $\mathcal{F}^n = 0$；对所有 $n$，图
$$
\xymatrix{
\mathcal{G}^n \ar[r] \ar[d] & \mathcal{G}^{n + 1} \ar[d] \\
\mathcal{F}^n \ar[r] & \mathcal{F}^{n + 1}
}
$$
交换；复合
$\mathcal{G}^n \to \mathcal{F}^n \to \mathcal{K}_0^n$
等于给定同态 $\mathcal{G}^n \to \mathcal{K}_0^n$；并且令
$\mathcal{F}^n_0 = \mathcal{F}^n \otimes_\mathcal{O} \mathcal{O}_0$
时，得到复形 $\mathcal{F}_0^\bullet$ 和拟同构的复形同态
$\mathcal{F}_0^\bullet \to \mathcal{K}_0^\bullet$。
\end{lemma}

\begin{proof}
我们对 $e$ 作递降归纳，证明可对 $n \geq e$ 找到 $\mathcal{F}^n$、
$\mathcal{G}^n \to \mathcal{F}^n$ 及
$\mathcal{F}^n \to \mathcal{F}^{n + 1}$，使其嵌入交换图
$$
\xymatrix{
\ldots \ar[r] &
\mathcal{G}^{e - 1} \ar[r] \ar@/_2pc/[dd] &
\mathcal{G}^e \ar[d] \ar[r] \ar@/_2pc/[dd] &
\mathcal{G}^{e + 1} \ar[d] \ar[r] \ar@/_2pc/[dd]|\hole &
\ldots \\
& &
\mathcal{F}^e \ar[d] \ar[r] &
\mathcal{F}^{e + 1} \ar[d] \ar[r] & \ldots \\
\ldots  \ar[r] &
\mathcal{K}_0^{e - 1} \ar[r] &
\mathcal{K}_0^e \ar[r] &
\mathcal{K}_0^{e + 1} \ar[r] & \ldots
}
$$
，其中 $\mathcal{F}_0^\bullet$ 是复形，且诱导同态
$\mathcal{F}_0^\bullet \to \mathcal{K}_0^\bullet$
在 $n>e$ 时于 $H^n$ 上诱导同构，在 $n=e$ 时诱导满射。
当 $e\gg0$ 时，由假设 (a)、(b)，可对 $n\geq e$ 取
$\mathcal{F}^n=0$，故结论成立。

\medskip\noindent
归纳步骤。我们需要构造 $\mathcal{F}^{e - 1}$ 以及同态
$\mathcal{G}^{e - 1} \to \mathcal{F}^{e - 1}$、
$\mathcal{F}^{e - 1} \to \mathcal{F}^e$ 与
$\mathcal{F}^{e - 1} \to \mathcal{K}_0^{e - 1}$。
我们把 $\mathcal{F}^{e - 1} = A \oplus B \oplus C$
选为三个部分的直和。

\medskip\noindent
第一部分取 $A = \mathcal{G}^{e - 1}$，并把同态
$\mathcal{G}^{e - 1} \to \mathcal{F}^{e - 1}$ 选为第一直和项的嵌入。
同态 $A \to \mathcal{K}^{e - 1}_0$ 与
$A \to \mathcal{F}^e$ 取为显然的同态。

\medskip\noindent
为选取 $B$，考虑由归纳假设得到的满射
$$
\gamma :
\Ker(\mathcal{F}^e_0 \to \mathcal{F}^{e + 1}_0)
\longrightarrow
\Ker(\mathcal{K}^e_0 \to \mathcal{K}^{e + 1}_0)/
\Im(\mathcal{K}^{e - 1}_0 \to \mathcal{K}^e_0)
$$
可选取集合 $I$，并对每个 $i\in I$ 选取 $\mathcal{C}$ 的对象
$U_i$ 以及截面
$s_i \in \mathcal{F}^e(U_i)$、$t_i \in \mathcal{K}^{e - 1}_0(U_i)$
，使得：
\begin{enumerate}
\item $s_i$ 映为 $\Ker(\gamma) \subset
\Ker(\mathcal{F}^e_0 \to \mathcal{F}^{e + 1}_0)$ 的截面；
\item $s_i$ 与 $t_i$ 映为 $\mathcal{K}^e_0$ 的同一截面；
\item 各截面 $s_i$ 作为 $\mathcal{O}_0$-模生成 $\Ker(\gamma)$。
\end{enumerate}
我们略去完整论证；这里用到
$\mathcal{F}^e \to \mathcal{F}^e_0$ 是集合层的满射。令
$$
B = \bigoplus\nolimits_{i \in I} j_{U_i!}\mathcal{O}_{U_i}
$$
并用 $s_i$ 与 $t_i$ 确定各直和项
$j_{U_i!}\mathcal{O}_{U_i}$ 的像，从而定义同态
$B \to \mathcal{F}^e$ 与 $B \to \mathcal{K}_0^{e - 1}$。

\medskip\noindent
取 $\mathcal{F}^{e - 1} = A \oplus B$ 及上述同态，便得到对应于
$e-1$ 的上述图，并且
$\mathcal{F}_0^\bullet \to \mathcal{K}_0^\bullet$
在 $n\geq e$ 时于 $H^n$ 上诱导同构。为了使该同态在
$H^{e-1}$ 上为满射，按如下方式选取直和项 $C$。
选取集合 $J$，并对每个 $j\in J$ 选取 $\mathcal{C}$ 的对象
$U_j$ 及 $U_j$ 上
$\Ker(\mathcal{K}^{e - 1}_0 \to \mathcal{K}^e_0)$ 的截面 $t_j$，
使这些截面作为 $\mathcal{O}_0$-模生成该核。令
$$
C = \bigoplus\nolimits_{j \in J} j_{U_j!}\mathcal{O}_{U_j}
$$
% P11-E0016: source uses s_j although only t_j was introduced; frozen symbol retained; see control/ERRATA_CANDIDATES.jsonl.
再取零同态 $C \to \mathcal{F}^e$，并用 $s_j$ 确定直和项
$j_{U_j!}\mathcal{O}_{U_j}$ 的像，从而定义同态
$C \to \mathcal{K}_0^{e - 1}$。取
$\mathcal{F}^{e - 1} = A \oplus B \oplus C$ 及上述各同态，
便完成归纳步骤。
\end{proof}

\begin{lemma}
\label{defos-lemma-canonical-class}
设 $\mathcal{C}$ 是位点，$\mathcal{O} \to \mathcal{O}_0$
是环层的满射，其核是平方零理想层 $\mathcal{I}$。
对 $D^-(\mathcal{O}_0)$ 中每个对象 $K_0$，存在典范同态
$$
\omega(K_0) :
K_0 \longrightarrow
K_0 \otimes_{\mathcal{O}_0}^\mathbf{L} \mathcal{I}[2]
$$
属于 $D(\mathcal{O}_0)$，并且对 $D^-(\mathcal{O}_0)$ 中任意同态
$K_0 \to L_0$，图
$$
\xymatrix{
K_0 \ar[d] \ar[rr]_-{\omega(K_0)} & &
(K_0 \otimes^\mathbf{L}_{\mathcal{O}_0} \mathcal{I})[2] \ar[d] \\
L_0 \ar[rr]^-{\omega(L_0)} & &
(L_0 \otimes^\mathbf{L}_{\mathcal{O}_0} \mathcal{I})[2]
}
$$
交换。
\end{lemma}

\begin{proof}
用任意 $\mathcal{O}_0$-模复形 $\mathcal{K}_0^\bullet$ 表示 $K_0$。
在所有 $n$ 上取 $\mathcal{G}^n=0$，应用引理
\ref{defos-lemma-lift-complex}。把该引理给出的同态记为
$d : \mathcal{F}^n \to \mathcal{F}^{n + 1}$。于是
$d \circ d : \mathcal{F}^n \to \mathcal{F}^{n + 2}$
模 $\mathcal{I}$ 为零。由于 $\mathcal{F}^n$ 平坦，有
$\mathcal{I}\mathcal{F}^n =
\mathcal{F}^n \otimes_{\mathcal{O}} \mathcal{I} =
\mathcal{F}^n_0 \otimes_{\mathcal{O}_0} \mathcal{I}$。
因此得到典范复形同态
$$
d \circ d : \mathcal{F}_0^\bullet
\longrightarrow
(\mathcal{F}_0^\bullet \otimes_{\mathcal{O}_0} \mathcal{I})[2]
$$
由于 $\mathcal{F}_0^\bullet$ 是平坦 $\mathcal{O}_0$-模的有上界复形，
它是 K-平坦的，因而可用于计算导出张量积。此外，按构造，复形同态
$\mathcal{F}_0^\bullet \to \mathcal{K}_0^\bullet$ 是拟同构。
因此，刚构造的同态的源与靶分别表示 $K_0$ 和
$K_0 \otimes_{\mathcal{O}_0}^\mathbf{L} \mathcal{I}[2]$
，从而得到同态 $\omega(K_0)$。

\medskip\noindent
下面证明此过程与复形同态相容。设 $\mathcal{L}_0^\bullet$
表示 $D^-(\mathcal{O}_0)$ 中另一对象，并假设
$$
\mathcal{K}_0^\bullet \longrightarrow \mathcal{L}_0^\bullet
$$
是复形同态。对复形 $\mathcal{L}_0^\bullet$、平坦模
$\mathcal{F}^n$、同态 $\mathcal{F}^n \to \mathcal{F}^{n + 1}$
以及复合
$\mathcal{F}^n \to \mathcal{K}_0^n \to \mathcal{L}_0^n$
应用引理 \ref{defos-lemma-lift-complex}（这里字母的使用次序颠倒了，谨致歉意）。
得到平坦模 $\mathcal{G}^n$，以及同态
$\mathcal{F}^n \to \mathcal{G}^n$、
$\mathcal{G}^n \to \mathcal{G}^{n + 1}$ 和
$\mathcal{G}^n \to \mathcal{L}_0^n$，它们具备该引理中的全部性质。
于是显然
$$
\xymatrix{
\mathcal{F}_0^\bullet \ar[d] \ar[r] &
(\mathcal{F}_0^\bullet \otimes_{\mathcal{O}_0} \mathcal{I})[2] \ar[d] \\
\mathcal{G}_0^\bullet \ar[r] &
(\mathcal{G}_0^\bullet \otimes_{\mathcal{O}_0} \mathcal{I})[2]
}
$$
是复形的交换图。

\medskip\noindent
为证明 $\omega(K_0)$ 良定义，假设有两个表示 $K_0$ 的
$\mathcal{O}_0$-模复形
$\mathcal{K}_0^\bullet$ 与 $(\mathcal{K}'_0)^\bullet$，以及上述类型的两组数据
$(\mathcal{F}^n, d : \mathcal{F}^n \to \mathcal{F}^{n + 1},
\mathcal{F}^n \to \mathcal{K}_0^n)$
以及
$((\mathcal{F}')^n, d : (\mathcal{F}')^n \to (\mathcal{F}')^{n + 1},
(\mathcal{F}')^n \to \mathcal{K}_0^n)$
。可以选取复形 $(\mathcal{K}''_0)^\bullet$ 和拟同构
$\mathcal{K}_0^\bullet  \to (\mathcal{K}''_0)^\bullet$
以及
$(\mathcal{K}'_0)^\bullet  \to (\mathcal{K}''_0)^\bullet$
，以实现这两个复形在导出范畴中均表示 $K_0$ 这一事实。然后把上一段的结果应用于
$$
(\mathcal{K}_0)^\bullet \oplus (\mathcal{K}'_0)^\bullet
\longrightarrow
(\mathcal{K}''_0)^\bullet
$$
，得到交换图
$$
\xymatrix{
\mathcal{F}_0^\bullet \oplus (\mathcal{F}'_0)^\bullet
\ar[d] \ar[r] &
(\mathcal{F}_0^\bullet \otimes_{\mathcal{O}_0} \mathcal{I})[2] \oplus
((\mathcal{F}'_0)^\bullet \otimes_{\mathcal{O}_0} \mathcal{I})[2] \ar[d] \\
\mathcal{G}_0^\bullet \ar[r] &
(\mathcal{G}_0^\bullet \otimes_{\mathcal{O}_0} \mathcal{I})[2]
}
$$
由于竖直箭头在各直和项上给出拟同构，便得到
$D(\mathcal{O}_0)$ 中所需的交换性。

\medskip\noindent
既已证明良定义性，与同态相容的断言即由第二段的结果推出。
\end{proof}

\begin{lemma}
\label{defos-lemma-induced-map}
设 $(\mathcal{C}, \mathcal{O})$ 是环化位点，
$\alpha : K \to L$ 是 $D^-(\mathcal{O})$ 中的同态，
$\mathcal{F}$ 是 $\mathcal{O}$-模层，且 $n \in \mathbf{Z}$。
\begin{enumerate}
\item 若对 $i \geq n$，$H^i(\alpha)$ 是同构，则对 $i \geq n$，
$H^i(\alpha \otimes_\mathcal{O}^\mathbf{L} \text{id}_\mathcal{F})$
也是同构。
\item 若对 $i > n$，$H^i(\alpha)$ 是同构，且对 $i = n$ 它是满射，
则对 $i > n$，
$H^i(\alpha \otimes_\mathcal{O}^\mathbf{L} \text{id}_\mathcal{F})$
是同构，且对 $i = n$ 它是满射。
\end{enumerate}
\end{lemma}

\begin{proof}
选取可区别三角
$$
K \to L \to C \to K[1]
$$
在情形 (2) 中，对 $i\geq n$ 有 $H^i(C)=0$。因此由
《导出范畴》引理 \ref{derived-lemma-negative-vanishing}
的对偶，对 $i\geq n$ 有
$H^i(C \otimes_\mathcal{O}^\mathbf{L} \mathcal{F}) = 0$。
这又表明
$H^i(\alpha \otimes_\mathcal{O}^\mathbf{L} \text{id}_\mathcal{F})$
在 $i>n$ 时是同构，在 $i=n$ 时是满射。
在情形 (1) 中，还可知 $H^{n - 1}(L) \to H^{n - 1}(C)$
是满射。考虑图
$$
\xymatrix{
H^{n - 1}(L) \otimes_\mathcal{O} \mathcal{F} \ar[r] \ar[d] &
H^{n - 1}(C) \otimes_\mathcal{O} \mathcal{F} \ar@{=}[d] \\
H^{n - 1}(L \otimes_\mathcal{O}^\mathbf{L} \mathcal{F}) \ar[r] &
H^{n - 1}(C \otimes_\mathcal{O}^\mathbf{L} \mathcal{F})
}
$$
，可知下方水平箭头是满射。结合前述结论，这表明
$H^n(\alpha \otimes_\mathcal{O}^\mathbf{L} \text{id}_\mathcal{F})$
是同构。
\end{proof}

\begin{lemma}
\label{defos-lemma-canonical-class-obstruction}
设 $\mathcal{C}$ 是位点，$\mathcal{O} \to \mathcal{O}_0$
是环层的满射，其核是平方零理想层 $\mathcal{I}$。
对 $D^-(\mathcal{O}_0)$ 中每个对象 $K_0$，以下条件等价：
\begin{enumerate}
\item 引理 \ref{defos-lemma-canonical-class} 中构造的类
$\omega(K_0) \in
\Ext^2_{\mathcal{O}_0}(K_0, K_0 \otimes_{\mathcal{O}_0} \mathcal{I})$
为零；
\item 存在 $K \in D^-(\mathcal{O})$，使得
$K \otimes_\mathcal{O}^\mathbf{L} \mathcal{O}_0 = K_0$
在 $D(\mathcal{O}_0)$ 中成立。
\end{enumerate}
\end{lemma}

\begin{proof}
设 $K$ 如 (2) 所述。可用平坦 $\mathcal{O}$-模的有上界复形
$\mathcal{F}^\bullet$ 表示 $K$。于是
$\mathcal{F}_0^\bullet =
\mathcal{F}^\bullet \otimes_{\mathcal{O}} \mathcal{O}_0$
在 $D(\mathcal{O}_0)$ 中表示 $K_0$。由于 $\mathcal{F}^\bullet$ 是复形，
$d_{\mathcal{F}^\bullet} \circ d_{\mathcal{F}^\bullet} = 0$；
从 $\omega(K_0)$ 的构造即可看出它为零。

\medskip\noindent
假设 (1) 成立。取构造 $\omega(K_0)$ 时的
$\mathcal{F}^n$ 与 $d : \mathcal{F}^n \to \mathcal{F}^{n + 1}$。
$\omega(K_0)$ 为零意味着同态
$$
\omega = d \circ d : \mathcal{F}_0^\bullet
\longrightarrow
(\mathcal{F}_0^\bullet \otimes_{\mathcal{O}_0} \mathcal{I})[2]
$$
在 $D(\mathcal{O}_0)$ 中为零。按导出范畴作为
$\mathcal{O}_0$-模复形同伦范畴之局部化的定义，存在拟同构
$\alpha : \mathcal{G}_0^\bullet \to \mathcal{F}_0^\bullet$
以及 $\mathcal{O}_0$-模同态
$h^n : \mathcal{G}_0^n \to
\mathcal{F}_0^{n + 1} \otimes_\mathcal{O} \mathcal{I}$
，使得
$$
\omega \circ \alpha =
d_{\mathcal{F}_0^\bullet \otimes \mathcal{I}} \circ h +
h \circ d_{\mathcal{G}_0^\bullet}
$$
令
$$
\mathcal{H}^n = \mathcal{F}^n \times_{\mathcal{F}^n_0} \mathcal{G}_0^n
$$
并定义
$$
d' : \mathcal{H}^n \longrightarrow \mathcal{H}^{n + 1},\quad
(f^n, g_0^n) \longmapsto (d(f^n) - h^n(g_0^n), d(g_0^n))
$$
这里使用显然的记号，并用到
$\mathcal{F}_0^{n + 1} \otimes_{\mathcal{O}_0} \mathcal{I} =
\mathcal{F}^{n + 1} \otimes_\mathcal{O} \mathcal{I} =
\mathcal{I}\mathcal{F}^{n + 1} \subset \mathcal{F}^{n + 1}$。
由 $h^n$ 的选取和 $\omega$ 的定义可核验 $d'\circ d'=0$。
因此 $\mathcal{H}^\bullet$ 定义 $D(\mathcal{O})$ 中的对象。
另一方面，有 $\mathcal{O}$-模复形的短正合列
$$
0 \to \mathcal{F}_0^\bullet \otimes_{\mathcal{O}_0} \mathcal{I} \to
\mathcal{H}^\bullet \to \mathcal{G}_0^\bullet \to 0
$$
尚需证明
$\mathcal{H}^\bullet \otimes_\mathcal{O}^\mathbf{L} \mathcal{O}_0$
与 $K_0$ 同构。选取拟同构
$\mathcal{E}^\bullet \to \mathcal{H}^\bullet$
，其中 $\mathcal{E}^\bullet$ 是平坦 $\mathcal{O}$-模的有上界复形。
得到交换图
$$
\xymatrix{
0 \ar[r] &
\mathcal{E}^\bullet \otimes_\mathcal{O} \mathcal{I} \ar[d]^\beta \ar[r] &
\mathcal{E}^\bullet \ar[d]^\gamma \ar[r] &
\mathcal{E}_0^\bullet \ar[d]^\delta \ar[r] &
0 \\
0 \ar[r] &
\mathcal{F}_0^\bullet \otimes_{\mathcal{O}_0} \mathcal{I} \ar[r] &
\mathcal{H}^\bullet \ar[r] &
\mathcal{G}_0^\bullet \ar[r] &
0
}
$$
我们断言 $\delta$ 是拟同构。由于当 $i\gg0$ 时 $H^i(\delta)$
是同构，可以对满足“当 $i\geq n$ 时 $H^i(\delta)$ 为同构”的
$n$ 作递降归纳。注意，
$\mathcal{E}^\bullet \otimes_\mathcal{O} \mathcal{I}$
表示
$\mathcal{E}_0^\bullet \otimes_{\mathcal{O}_0}^\mathbf{L} \mathcal{I}$，
而
$\mathcal{F}_0^\bullet \otimes_{\mathcal{O}_0} \mathcal{I}$
表示
$\mathcal{G}_0^\bullet \otimes_{\mathcal{O}_0}^\mathbf{L} \mathcal{I}$，
并且作为 $D(\mathcal{O}_0)$ 中的同态，有
$\beta = \delta \otimes_{\mathcal{O}_0}^\mathbf{L} \text{id}_\mathcal{I}$
。这是因为
$\beta =
(\alpha \otimes \text{id}_\mathcal{I})
\circ
(\delta \otimes \text{id}_\mathcal{I})$。
假设在次数 $\geq n$ 上 $H^i(\delta)$ 是同构。由刚才所述及引理
\ref{defos-lemma-induced-map}，$\beta$ 也具有同一性质。考虑图
$$
\xymatrix{
H^{n - 1}(\mathcal{E}^\bullet \otimes_\mathcal{O} \mathcal{I})
\ar[r] \ar[d]^{H^{n - 1}(\beta)} &
H^{n - 1}(\mathcal{E}^\bullet) \ar[r] \ar[d] &
H^{n - 1}(\mathcal{E}_0^\bullet) \ar[r] \ar[d]^{H^{n - 1}(\delta)} &
H^n(\mathcal{E}^\bullet \otimes_\mathcal{O} \mathcal{I})
\ar[r] \ar[d]^{H^n(\beta)} &
H^n(\mathcal{E}^\bullet) \ar[d] \\
H^{n - 1}(\mathcal{F}_0^\bullet \otimes_\mathcal{O} \mathcal{I}) \ar[r] &
H^{n - 1}(\mathcal{H}^\bullet) \ar[r] &
H^{n - 1}(\mathcal{G}_0^\bullet) \ar[r] &
H^n(\mathcal{F}_0^\bullet \otimes_\mathcal{O} \mathcal{I}) \ar[r] &
H^n(\mathcal{H}^\bullet)
}
$$
由《同调代数》引理 \ref{homology-lemma-four-lemma}，
$H^{n - 1}(\delta)$ 是满射。再由引理 \ref{defos-lemma-induced-map}，
$H^{n - 1}(\beta)$ 是满射。再次应用《同调代数》引理
\ref{homology-lemma-four-lemma}，可知 $H^{n - 1}(\delta)$ 是同构。
归纳证明了该断言，故 $\delta$ 是拟同构，这正是所需结论。
\end{proof}

\begin{lemma}
\label{defos-lemma-lift-map-complexes}
设 $\mathcal{C}$ 是位点，$\mathcal{O} \to \mathcal{O}_0$
是环层的满射。假设给定以下数据：
\begin{enumerate}
\item $\mathcal{O}$-模复形 $\mathcal{F}^\bullet$；
\item $\mathcal{O}_0$-模复形 $\mathcal{K}_0^\bullet$；
\item 拟同构 $\mathcal{K}_0^\bullet \to
\mathcal{F}^\bullet \otimes_\mathcal{O} \mathcal{O}_0$。
\end{enumerate}
则存在拟同构 $\mathcal{G}^\bullet \to \mathcal{F}^\bullet$，使得复形同态
$\mathcal{G}^\bullet  \otimes_\mathcal{O} \mathcal{O}_0 \to
\mathcal{F}^\bullet \otimes_\mathcal{O} \mathcal{O}_0$
在 $\mathcal{O}_0$-模复形的同伦范畴中经由
$\mathcal{K}_0^\bullet$ 分解。
\end{lemma}

\begin{proof}
令 $\mathcal{F}_0^\bullet =
\mathcal{F}^\bullet \otimes_\mathcal{O} \mathcal{O}_0$。
由《导出范畴》引理 \ref{derived-lemma-make-surjective}，
给定同态有分解
$$
\mathcal{K}_0^\bullet \to \mathcal{L}_0^\bullet \to \mathcal{F}_0^\bullet
$$
，其中第一箭头有精确到同伦的逆，第二箭头逐项为分裂满射。
因此可以假设 $\mathcal{K}_0^\bullet \to \mathcal{F}_0^\bullet$
逐项为满射。此时取
$$
\mathcal{G}^n = \mathcal{F}^n \times_{\mathcal{F}^n_0} \mathcal{K}_0^n
$$
，一切即显然。
\end{proof}

\begin{lemma}
\label{defos-lemma-inf-obs-map-defo-complex}
设 $\mathcal{C}$ 是位点，$\mathcal{O} \to \mathcal{O}_0$
是环层的满射，其核是平方零理想层 $\mathcal{I}$。
设 $K,L\in D^-(\mathcal{O})$，并在 $D^-(\mathcal{O}_0)$ 中令
$K_0 = K \otimes_\mathcal{O}^\mathbf{L} \mathcal{O}_0$ 及
$L_0 = L \otimes_\mathcal{O}^\mathbf{L} \mathcal{O}_0$。
给定 $D(\mathcal{O}_0)$ 中的 $\alpha_0 : K_0 \to L_0$，存在典范元素
$$
o(\alpha_0) \in \Ext^1_{\mathcal{O}_0}(K_0,
L_0 \otimes_{\mathcal{O}_0}^\mathbf{L} \mathcal{I})
$$
；它为零当且仅当 $D(\mathcal{O})$ 中存在同态
$\alpha : K \to L$，满足
$\alpha_0 = \alpha \otimes_\mathcal{O}^\mathbf{L} \text{id}$。
\end{lemma}

\begin{proof}
寻找提升 $\alpha_0$ 的 $\alpha:K\to L$，等价于寻找
$\alpha:K\to L$，使复合 $K \xrightarrow{\alpha} L \to L_0$
等于复合 $K \to K_0 \xrightarrow{\alpha_0} L_0$。短正合列
$0 \to \mathcal{I} \to \mathcal{O} \to \mathcal{O}_0 \to 0$
给出典范可区别三角
$$
L \otimes_\mathcal{O}^\mathbf{L} \mathcal{I} \to
L \to
L_0 \to
(L \otimes_\mathcal{O}^\mathbf{L} \mathcal{I})[1]
$$
属于 $D(\mathcal{O})$。由《导出范畴》引理
\ref{derived-lemma-representable-homological}，复合
$$
K \to K_0 \xrightarrow{\alpha_0} L_0 \to
(L \otimes_\mathcal{O}^\mathbf{L} \mathcal{I})[1]
$$
为零当且仅当可找到提升 $\alpha_0$ 的 $\alpha:K\to L$。
此复合是下列群的元素：
$$
\Hom_{D(\mathcal{O})}(K, (L \otimes_\mathcal{O}^\mathbf{L} \mathcal{I})[1]) =
\Hom_{D(\mathcal{O}_0)}(K_0,
(L \otimes_\mathcal{O}^\mathbf{L} \mathcal{I})[1]) =
\Ext^1_{\mathcal{O}_0}(K_0,
L_0 \otimes_{\mathcal{O}_0}^\mathbf{L} \mathcal{I})
$$
，其中等号来自伴随。
\end{proof}

\begin{lemma}
\label{defos-lemma-first-order-defos-complex}
设 $\mathcal{C}$ 是位点，$\mathcal{O} \to \mathcal{O}_0$
是环层的满射，其核是平方零理想层 $\mathcal{I}$。
% P11-E0017: K_0 is typed over O although the lift is identified in D(O_0); frozen category retained; see control/ERRATA_CANDIDATES.jsonl.
设 $K_0 \in D^-(\mathcal{O})$。$K_0$ 的一个提升是二元组
$(K,\alpha_0)$，其中 $K$ 是 $D^-(\mathcal{O})$ 中的对象，而
$\alpha_0 : K \otimes_\mathcal{O}^\mathbf{L} \mathcal{O}_0 \to K_0$
是 $D(\mathcal{O}_0)$ 中的同构。
\begin{enumerate}
\item 给定提升 $(K,\alpha)$，该二元组的自同构群典范地等于同态
$$
\Ext^{-1}_{\mathcal{O}_0}(K_0, K_0)
\longrightarrow
\Hom_{\mathcal{O}_0}(K_0, K_0 \otimes_{\mathcal{O}_0}^\mathbf{L} \mathcal{I})
$$
的余核。
\item 若存在提升，则提升的同构类之集合在
$\Ext^1_{\mathcal{O}_0}(K_0,
K_0 \otimes_{\mathcal{O}_0}^\mathbf{L} \mathcal{I})$
的作用下为主齐性空间。
\end{enumerate}
\end{lemma}

\begin{proof}
$(K,\alpha)$ 的一个自同构是 $D(\mathcal{O})$ 中满足
$\varphi \otimes_\mathcal{O} \text{id}_{\mathcal{O}_0} = \text{id}$
的同态 $\varphi:K\to K$。这等价于说复合
$$
K \xrightarrow{\varphi - \text{id}} K \to
K \otimes_\mathcal{O}^\mathbf{L} \mathcal{O}_0
$$
为零。由此可知，自同构群是同态
$$
\Hom_\mathcal{O}(K, K_0[-1])
\longrightarrow
\Hom_\mathcal{O}(K, K_0 \otimes_{\mathcal{O}_0}^\mathbf{L} \mathcal{I})
$$
的余核；这是由 $D(\mathcal{O})$ 中的可区别三角
$$
K \otimes_\mathcal{O}^\mathbf{L} \mathcal{I} \to
K \to
K \otimes_\mathcal{O}^\mathbf{L} \mathcal{O}_0 \to
(K \otimes_\mathcal{O}^\mathbf{L} \mathcal{I})[1]
$$
及《导出范畴》引理 \ref{derived-lemma-representable-homological}
得出的。要把它改写为引理中的群，使用限制函子
$D(\mathcal{O}_0) \to D(\mathcal{O})$ 与
$- \otimes_\mathcal{O} \mathcal{O}_0 : D(\mathcal{O}) \to D(\mathcal{O}_0)$
之间的伴随。这证明了 (1)。

\medskip\noindent
证明 (2)。假设在 $D(\mathcal{O})$ 中
$K_0 = K \otimes_\mathcal{O}^\mathbf{L} \mathcal{O}_0$。
由引理 \ref{defos-lemma-inf-obs-map-defo-complex}，把提升
$(K',\alpha_0)$ 映到提升 $\alpha_0$ 的障碍 $o(\alpha_0)$ 的同态，
定义了从二元组的同构类集合到
$\Ext^1_{\mathcal{O}_0}(K_0,
K_0 \otimes_{\mathcal{O}_0}^\mathbf{L} \mathcal{I})$
的典范单射。为完成证明，只需证明它为满射。
取本引理的 $\Ext^1$ 中的
$\xi : K_0 \to (K_0 \otimes_{\mathcal{O}_0}^\mathbf{L} \mathcal{I})[1]$。
选取表示 $K$ 的平坦 $\mathcal{O}$-模的有上界复形
$\mathcal{F}^\bullet$。同态 $\xi$ 可表示为 $t\circ s^{-1}$，其中
$s : \mathcal{K}_0^\bullet \to \mathcal{F}_0^\bullet$ 是拟同构，而
$t : \mathcal{K}_0^\bullet \to
\mathcal{F}_0^\bullet \otimes_{\mathcal{O}_0} \mathcal{I}[1]$
是复形同态。由引理 \ref{defos-lemma-lift-map-complexes}，可以假设存在
$\mathcal{O}$-模复形间的拟同构
$\mathcal{G}^\bullet \to \mathcal{F}^\bullet$，使得
$\mathcal{G}_0^\bullet \to \mathcal{F}_0^\bullet$
在精确到同伦的意义下经由 $s$ 分解。我们可以并且确实把
$\mathcal{G}^\bullet$ 换成平坦 $\mathcal{O}$-模的有上界复形
（选取从这样一个复形到 $\mathcal{G}^\bullet$ 的拟同构并作替换）。
于是 $\xi$ 由复形同态
$t : \mathcal{G}_0^\bullet \to
\mathcal{F}_0^\bullet \otimes_{\mathcal{O}_0} \mathcal{I}[1]$
及拟同构 $\mathcal{G}_0^\bullet \to \mathcal{F}_0^\bullet$ 表示。令
$$
\mathcal{H}^n = \mathcal{F}^n \times_{\mathcal{F}_0^n} \mathcal{G}_0^n
$$
，其微分为
$$
\mathcal{H}^n \to \mathcal{H}^{n + 1},\quad
(f^n, g_0^n) \mapsto (d(f^n) + t(g_0^n), d(g_0^n))
$$
这是有意义的，因为
$\mathcal{F}_0^{n + 1} \otimes_{\mathcal{O}_0} \mathcal{I} =
\mathcal{F}^{n + 1} \otimes_\mathcal{O} \mathcal{I} =
\mathcal{I}\mathcal{F}^{n + 1} \subset \mathcal{F}^{n + 1}$。
我们略去证明 $\mathcal{H}^\bullet$ 是 $\mathcal{O}$-模复形的计算。
按构造有短正合列
$$
0 \to \mathcal{F}_0^\bullet \otimes_{\mathcal{O}_0} \mathcal{I} \to
\mathcal{H}^\bullet \to \mathcal{G}_0^\bullet \to 0
$$
，它由 $\mathcal{O}$-模复形构成。与引理
\ref{defos-lemma-canonical-class-obstruction} 的证明完全相同，
可证明此列诱导同构
$\alpha_0 :
\mathcal{H}^\bullet \otimes_\mathcal{O}^\mathbf{L} \mathcal{O}_0 \to
\mathcal{G}_0^\bullet$，属于 $D(\mathcal{O}_0)$。
换言之，我们构造出了二元组 $(\mathcal{H}^\bullet,\alpha_0)$。
我们略去 $o(\alpha_0)=\xi$ 的核验；提示：在此情形下，
由于存在提升 $\alpha_0$ 各分量的同态
$\mathcal{H}^n \to \mathcal{F}^n$（它们与微分不相容），
可以显式计算 $o(\alpha_0)$。证明完毕。
\end{proof}
% Generated mechanically from the frozen cotangent.tex target.
% Do not edit this generated file; edit the bound locale source instead.

% OK, start here.
%

\setcounter{chapter}{91}
\chapter{余切复形}



\phantomsection
\label{cotangent-section-phantom}


\section{引言}
\label{cotangent-section-introduction}

\noindent
本章的目标是构造环同态、概形态射以及代数空间态射的余切复形。
可供参考的文献包括讲义 \cite{quillenhomology}、论文
\cite{quillencohomology}，以及专著 \cite{Andre} 和 \cite{cotangent}。




\section{给读者的建议}
\label{cotangent-section-advice-reader}

\noindent
撰写本章时，我们尽量减少单纯形技巧的使用。对环 $A$ 上的环 $B$，
我们把选择一个预解 $P_\bullet$ 看作计算范畴 $\mathcal{C}_{B/A}$
上阿贝尔层同调的工具，见注 \ref{cotangent-remark-resolution}。这类似于用
{\v C}ech 复形计算上同调时“良覆盖”所起的作用。若要简要了解
范畴上的同调，请参阅《位点上的上同调》第
\ref{sites-cohomology-section-homology} 节。对于同调而言，导出下惊叹号
函子 $L\pi_!$ 所扮演的角色，正如
$R\Gamma(\mathcal{C}_{B/A}, -)$ 对上同调所扮演的角色。第
\ref{cotangent-section-compute-L-pi-shriek} 节研究的范畴
$\mathcal{C}_{B/A}$，是由分解 $A \to P \to B$（其中 $P$ 是 $A$
上的多项式代数）组成的范畴的反范畴。该范畴上有环层同态
$$
\underline{A} \longrightarrow \mathcal{O} \longrightarrow \underline{B}
$$
并且在对象 $U = (P \to B)$ 上有 $\mathcal{O}(U)=P$。最终可得 $B$
在 $A$ 上的余切复形为
$$
L_{B/A} =
L\pi_!(\Omega_{\mathcal{O}/\underline{A}} \otimes_\mathcal{O} \underline{B})
$$
见引理 \ref{cotangent-lemma-compute-cotangent-complex}。我们始终尽力采用这一观点
证明环同态余切复形的基本性质。尤其是，所有结果原则上都可在不依赖
标准预解存在性的情况下证明，尽管我们并未这样做。这一理论相当令人
满意；唯一可能的例外是基本三角（命题 \ref{cotangent-proposition-triangle}）
的证明需要稍多一点关于导出下惊叹号函子的理论。为向读者提供另一种
途径，我们在注 \ref{cotangent-remark-triangle} 和 \ref{cotangent-remark-explicit-map}
中给出一种仅依赖标准预解简单性质的、相当完整的证明纲要。

\medskip\noindent
对于环化拓扑斯态射、概形态射、代数空间态射等的余切复形，我们的做法是
尽可能从上面讨论的“普通环同态”情形推出结论。





\section{环同态的余切复形}
\label{cotangent-section-cotangent-ring-map}

\noindent
设 $A$ 为环，以 $\textit{Alg}_A$ 记 $A$-代数范畴。考虑一对伴随函子
$(U,V)$，其中 $V:\textit{Alg}_A\to\textit{Sets}$ 是遗忘函子，
$U:\textit{Sets}\to\textit{Alg}_A$ 则把集合 $E$ 送到 $A$ 上以
$E$ 为变量集的多项式代数 $A[E]$。令 $X_\bullet$ 为《单纯形》第
\ref{simplicial-section-standard} 节所构造的
$\text{Fun}(\textit{Alg}_A,\textit{Alg}_A)$ 中的单纯形对象。

\medskip\noindent
取一个 $A$-代数 $B$，以 $P_\bullet=X_\bullet(B)$ 记所得的单纯形
$A$-代数。回忆 $P_0=A[B]$、$P_1=A[A[B]]$，依此类推。特别地，
每一项 $P_n$ 都是多项式 $A$-代数。还记得存在增广
$$
\epsilon : P_\bullet \longrightarrow B
$$
这里把 $B$ 看作常值单纯形 $A$-代数。

\begin{definition}
\label{cotangent-definition-standard-resolution}
设 $A\to B$ 为环同态。$B$ 在 $A$ 上的{\it 标准预解}是增广
$\epsilon:P_\bullet\to B$，其各项为
$$
P_0 = A[B],\quad P_1 = A[A[B]],\quad \ldots
$$
而同态按《单纯形》例 \ref{simplicial-example-polynomial-algebra-maps}
中的方式构造。
\end{definition}

\noindent
我们将看到，在某些情形下可以利用标准预解计算左导出函子。

\begin{definition}
\label{cotangent-definition-cotangent-complex-ring-map}
环同态 $A\to B$ 的{\it 余切复形} $L_{B/A}$，是与单纯形
$B$-模
$$
\Omega_{P_\bullet/A} \otimes_{P_\bullet, \epsilon} B
$$
相伴的 $B$-模复形，其中 $\epsilon:P_\bullet\to B$ 是 $B$
在 $A$ 上的标准预解。
\end{definition}

\noindent
在《单纯形》第 \ref{simplicial-section-complexes} 节中，我们为单纯形模
配上链复形；但这里使用的是上链复形。因此，次数 $-n$ 上的项
$L_{B/A}^{-n}$ 是 $B$-模
$\Omega_{P_n/A}\otimes_{P_n,\epsilon_n}B$，且当 $m>0$ 时
$L_{B/A}^m=0$。

\begin{remark}
\label{cotangent-remark-variant-cotangent-complex}
设 $A\to B$ 为环同态。令 $\mathcal{A}$ 为 $A$-代数箭头
$\psi:C\to B$ 所成的范畴，令 $\mathcal{S}$ 为同态 $E\to B$
所成的范畴，其中 $E$ 是集合。存在伴随函子
$V:\mathcal{A}\to\mathcal{S}$（遗忘函子）与
$U:\mathcal{S}\to\mathcal{A}$；后者把 $E\to B$ 送到
$A[E]\to B$。令 $X_\bullet$ 为《单纯形》第
\ref{simplicial-section-standard} 节所构造的
$\text{Fun}(\mathcal{A},\mathcal{A})$ 中的单纯形对象。图
$$
\xymatrix{
\mathcal{A} \ar[d] \ar[r] & \mathcal{S} \ar@<1ex>[l] \ar[d] \\
\textit{Alg}_A \ar[r] & \textit{Sets} \ar@<1ex>[l]
}
$$
交换。因此，$X_\bullet(\text{id}_B:B\to B)$ 等于 $B$
在 $A$ 上的标准预解。
\end{remark}

\begin{lemma}
\label{cotangent-lemma-colimit-cotangent-complex}
设 $A_i\to B_i$ 是有向指标集 $I$ 上的一族环同态。那么
$\colim L_{B_i/A_i}=L_{\colim B_i/\colim A_i}$。
\end{lemma}

\begin{proof}
这是因为遗忘函子 $V:A\textit{-Alg}\to\textit{Sets}$ 及其伴随函子
$U:\textit{Sets}\to A\textit{-Alg}$ 都与滤余极限交换。此外，函子
$B/A\mapsto\Omega_{B/A}$ 也具有这一性质（《交换代数》引理
\ref{algebra-lemma-colimit-differentials}）。
\end{proof}





\section{单纯形预解与导出下惊叹号函子}
\label{cotangent-section-compute-L-pi-shriek}

\noindent
设 $A\to B$ 为环同态。考虑如下范畴：其对象是 $A$-代数同态
$\alpha:P\to B$，其中 $P$ 是 $A$ 上以某个集合\footnote{只需考虑
基数不超过 $B$ 的基数的集合。}为变量集的多项式代数；其态射
$s:(\alpha:P\to B)\to(\alpha':P'\to B)$ 是满足
$\alpha'\circ s=\alpha$ 的 $A$-代数同态 $s:P\to P'$。以
$\mathcal{C}=\mathcal{C}_{B/A}$ 记此范畴的{\bf 反范畴}。取反范畴的
原因是，我们希望把对象 $(P,\alpha)$ 看作对应于仿射概形图
$$
\xymatrix{
\Spec(B) \ar[d] \ar[r] & \Spec(P) \ar[ld] \\
\Spec(A)
}
$$
我们赋予 $\mathcal{C}$ 混沌拓扑（《位点》例
\ref{sites-example-indiscrete}）；换言之，把 $\mathcal{C}$ 视为一个
位点，其覆盖均由恒等态射给出，因而所有预层都是层。此外，在
$\mathcal{C}$ 上赋予两个环层。第一个是把对象 $(P,\alpha)$ 送到 $P$
的层 $\mathcal{O}$；第二个是常值层 $B$，记作 $\underline{B}$。
于是得到如下环化拓扑斯态射图
\begin{equation}
\label{cotangent-equation-pi}
\vcenter{
\xymatrix{
(\Sh(\mathcal{C}), \underline{B}) \ar[r]_i \ar[d]_\pi &
(\Sh(\mathcal{C}), \mathcal{O}) \\
(\Sh(*), B)
}
}
\end{equation}
态射 $i$ 在底层拓扑斯上是恒等态射，而
$i^\sharp:\mathcal{O}\to\underline{B}$ 是显然的同态。同态 $\pi$
如《位点上的上同调》例
\ref{sites-cohomology-example-category-to-point} 所述。下列导出函子将在
后文中起重要作用：
$
Li^* : D(\mathcal{O}) \longrightarrow D(\underline{B})
$
它左伴随于 $Ri_*=i_*:D(\underline{B})\to D(\mathcal{O})$；以及
$
L\pi_! : D(\underline{B}) \longrightarrow D(B)
$
它左伴随于 $\pi^*=\pi^{-1}:D(B)\to D(\underline{B})$。

\begin{lemma}
\label{cotangent-lemma-identify-pi-shriek}
沿用上述记号，设 $P_\bullet$ 是带有增广
$\epsilon:P_\bullet\to B$ 的单纯形 $A$-代数。假设每个 $P_n$
都是 $A$ 上的多项式代数，并且 $\epsilon$ 在底层单纯形集上是平凡
Kan 纤维化。那么
$$
L\pi_!(\mathcal{F}) = \mathcal{F}(P_\bullet, \epsilon)
$$
在 $D(\textit{Ab})$ 或 $D(B)$ 中成立，并且分别关于
$\textit{Ab}(\mathcal{C})$ 或 $\textit{Mod}(\underline{B})$ 中的
$\mathcal{F}$ 具有函子性。
\end{lemma}

\begin{proof}
我们将使用《位点上的上同调》引理
\ref{sites-cohomology-lemma-compute-by-cosimplicial-resolution} 的判据。
给定 $\mathcal{C}$ 的对象 $U=(Q,\beta)$，须证明
$$
S_\bullet = \Mor_\mathcal{C}((P_\bullet, \epsilon), (Q, \beta))
$$
与单点集同伦等价。把 $Q$ 写成 $A[E]$，其中 $E$ 为某个集合（由
$\mathcal{C}$ 的选取可如此表示）。于是
$$
S_\bullet = \Mor_{\textit{Sets}}((E, \beta|_E), (P_\bullet, \epsilon))
$$
令 $*$ 为单点集上的常值单纯形集。对 $b\in B$，以
$F_{b,\bullet}$ 记由笛卡儿图
$$
\xymatrix{
F_{b, \bullet} \ar[r] \ar[d] & P_\bullet \ar[d]_\epsilon \\
{*} \ar[r]^b & B
}
$$
定义的单纯形集。按此记号，
$S_\bullet=\prod_{e\in E}F_{\beta(e),\bullet}$。由于假设
$\epsilon$ 是平凡 Kan 纤维化，$F_{b,\bullet}\to *$ 也是平凡 Kan
纤维化（《单纯形》引理
\ref{simplicial-lemma-trivial-kan-base-change}）。因此
$S_\bullet\to *$ 是平凡 Kan 纤维化（《单纯形》引理
\ref{simplicial-lemma-product-trivial-kan}），从而 $S_\bullet$ 与 $*$
同伦等价（《单纯形》引理
\ref{simplicial-lemma-trivial-kan-homotopy}）。
\end{proof}

\noindent
特别地，可以用 $B$ 在 $A$ 上的标准预解计算导出下惊叹号函子。

\begin{lemma}
\label{cotangent-lemma-pi-shriek-standard}
设 $A\to B$ 为环同态，$\epsilon:P_\bullet\to B$ 为 $B$ 在 $A$
上的标准预解，并令 $\pi$ 如 (\ref{cotangent-equation-pi}) 所示。那么
$$
L\pi_!(\mathcal{F}) = \mathcal{F}(P_\bullet, \epsilon)
$$
在 $D(\textit{Ab})$ 或 $D(B)$ 中成立，并且分别关于
$\textit{Ab}(\mathcal{C})$ 或 $\textit{Mod}(\underline{B})$ 中的
$\mathcal{F}$ 具有函子性。
\end{lemma}

\begin{proof}[第一种证明]
应用引理 \ref{cotangent-lemma-identify-pi-shriek}。由于各项 $P_n$ 都是多项式
代数，该引理的第一个假设成立。第二个假设证明如下。由《单纯形》引理
\ref{simplicial-lemma-standard-simplicial-homotopy}，同态 $\epsilon$
在底层单纯形集上是同伦等价。由《单纯形》引理
\ref{simplicial-lemma-homotopy-equivalence}，这说明 $\epsilon$ 在相伴的
阿贝尔群复形之间诱导拟同构。再由《单纯形》引理
\ref{simplicial-lemma-qis-simplicial-abelian-groups}，$\epsilon$
在底层单纯形集上是平凡 Kan 纤维化。
\end{proof}

\begin{proof}[第二种证明]
我们将使用《位点上的上同调》引理
\ref{sites-cohomology-lemma-compute-by-cosimplicial-resolution} 的判据。
设 $U=(Q,\beta)$ 为 $\mathcal{C}$ 的对象。须证明
$$
S_\bullet = \Mor_\mathcal{C}((P_\bullet, \epsilon), (Q, \beta))
$$
与单点集同伦等价。把 $Q$ 写成 $A[E]$，其中 $E$ 是某个集合（由
$\mathcal{C}$ 的选取可如此表示）。使用注
\ref{cotangent-remark-variant-cotangent-complex} 的记号可得
$$
S_\bullet = \Mor_\mathcal{S}((E \to B), i(P_\bullet \to B))
$$
由《单纯形》引理
\ref{simplicial-lemma-standard-simplicial-homotopy}，同态
$i(P_\bullet\to B)\to i(B\to B)$ 在 $\mathcal{S}$ 中是同伦等价。
因此 $S_\bullet$ 与
$$
\Mor_\mathcal{S}((E \to B), (B \to B)) = \{*\}
$$
同伦等价，正合所需。
\end{proof}

\begin{lemma}
\label{cotangent-lemma-compute-cotangent-complex}
设 $A\to B$ 为环同态，并令 $\pi$ 与 $i$ 如 (\ref{cotangent-equation-pi})
所示。存在典范同构
$$
L_{B/A} = L\pi_!(Li^*\Omega_{\mathcal{O}/A}) =
L\pi_!(i^*\Omega_{\mathcal{O}/A}) =
L\pi_!(\Omega_{\mathcal{O}/A} \otimes_\mathcal{O} \underline{B})
$$
属于 $D(B)$。
\end{lemma}

\begin{proof}
对范畴 $\mathcal{C}$ 的对象 $\alpha:P\to B$，模 $\Omega_{P/A}$
是自由 $P$-模。因此 $\Omega_{\mathcal{O}/A}$ 是平坦
$\mathcal{O}$-模。于是
$Li^*\Omega_{\mathcal{O}/A}=i^*\Omega_{\mathcal{O}/A}$ 是这样的
$\underline{B}$-模层：它把 $\alpha:P\to A$ 送到 $B$-模
$\Omega_{P/A}\otimes_{P,\alpha}B$。
% P11-E0018: the frozen source says alpha:P->A although objects of C are alpha:P->B; retained here. See control/ERRATA_CANDIDATES.jsonl.
由引理 \ref{cotangent-lemma-pi-shriek-standard}，右端由该层在标准预解上的值
计算；而依定义 \ref{cotangent-definition-cotangent-complex-ring-map}，这正是左端
的定义。
\end{proof}

\begin{lemma}
\label{cotangent-lemma-pi-lower-shriek-constant-sheaf}
若 $A\to B$ 是环同态，且 $\pi$ 如 (\ref{cotangent-equation-pi}) 所示，则
$L\pi_!(\pi^{-1}M)=M$。
\end{lemma}

\begin{proof}
由引理 \ref{cotangent-lemma-identify-pi-shriek}，$L\pi_!(\pi^{-1}M)$ 可由
$(\pi^{-1}M)(P_\bullet,\epsilon)$ 计算，而后者是 $M$ 上的常值
单纯形对象。
\end{proof}

\begin{lemma}
\label{cotangent-lemma-identify-H0}
若 $A\to B$ 是环同态，则 $H^0(L_{B/A})=\Omega_{B/A}$。
\end{lemma}

\begin{proof}
我们直接计算。使用引理 \ref{cotangent-lemma-compute-cotangent-complex} 中的等同。
显然存在从 $\Omega_{\mathcal{O}/A}\otimes\underline{B}$ 到取值为
$\Omega_{B/A}$ 的常值层的同态。因此，它诱导同态
$$
H^0(L_{B/A}) = H^0(L\pi_!(\Omega_{\mathcal{O}/A} \otimes \underline{B}))
= \pi_!(\Omega_{\mathcal{O}/A} \otimes \underline{B}) \to \Omega_{B/A}
$$
在 $\mathcal{C}_{B/A}$ 中选取对象 $P\to B$，使 $P\to B$ 满射；
由《交换代数》引理 \ref{algebra-lemma-differential-surjective} 可知上述
同态满射。为证其单射，设 $P\to B$ 是 $\mathcal{C}_{B/A}$ 的对象，
且 $\xi\in\Omega_{P/A}\otimes_PB$ 在 $\Omega_{B/A}$ 中的像为零。
先选取分解 $P\to P'\to B$，使 $P'\to B$ 满射且 $P'$ 是 $A$
上的多项式代数。可以用 $P'$ 替换 $P$。若 $B=P/I$，则
$\Omega_{P/A}\otimes_PB\to\Omega_{B/A}$ 的核是 $I/I^2$ 的像
（《交换代数》引理 \ref{algebra-lemma-differential-seq}）。设 $\xi$
是 $f\in I$ 的像。考虑两个同态 $a,b:P'=P[x]\to P$：第一个把
$x$ 映到 $0$，第二个把 $x$ 映到 $f$（在两种情形中，
$P[x]\to B$ 都把 $x$ 映到零）。于是 $\xi$ 与 $0$ 都是
$\text{d}x\otimes1$ 在 $\Omega_{P'/A}\otimes_{P'}B$ 中的像。
故 $\xi$ 与 $0$ 在余极限
$\pi_!(\Omega_{\mathcal{O}/A}\otimes\underline{B})$ 中有相同的像
（见《位点上的上同调》例
\ref{sites-cohomology-example-category-to-point}），正合所需。
\end{proof}

\begin{lemma}
\label{cotangent-lemma-pi-lower-shriek-polynomial-algebra}
若 $B$ 是环 $A$ 上的多项式代数，且 $\pi$ 如
(\ref{cotangent-equation-pi}) 所示，则 $\pi_!$ 正合，且
$\pi_!\mathcal{F}=\mathcal{F}(B\to B)$。
\end{lemma}

\begin{proof}
由引理 \ref{cotangent-lemma-identify-pi-shriek}，可用 $B$ 上的常值单纯形代数
计算 $L\pi_!$。
\end{proof}

\begin{lemma}
\label{cotangent-lemma-cotangent-complex-polynomial-algebra}
若 $B$ 是环 $A$ 上的多项式代数，则 $L_{B/A}$ 与
$\Omega_{B/A}[0]$ 拟同构。
\end{lemma}

\begin{proof}
立即由引理 \ref{cotangent-lemma-compute-cotangent-complex} 和
\ref{cotangent-lemma-pi-lower-shriek-polynomial-algebra} 得出。
\end{proof}





\section{构造预解}
\label{cotangent-section-polynomial}

\noindent
在诺特有限型情形下，可以为有限型环同态构造一个“小”的单纯形预解。

\begin{lemma}
\label{cotangent-lemma-polynomial}
设 $A$ 是诺特环，$A\to B$ 是有限型环同态。令 $\mathcal{A}$ 为
$A$-代数同态 $C\to B$ 所成的范畴。设 $n\geq0$，并设
$P_\bullet$ 是 $\mathcal{A}$ 中满足下列条件的单纯形对象：
\begin{enumerate}
\item $P_\bullet\to B$ 是单纯形集的平凡 Kan 纤维化；
\item 当 $k\leq n$ 时，$P_k$ 在 $A$ 上有限型；
\item 作为 $\mathcal{A}$ 中的单纯形对象，
$P_\bullet=\text{cosk}_n\text{sk}_nP_\bullet$。
\end{enumerate}
那么 $P_{n+1}$ 是有限型 $A$-代数。
\end{lemma}

\begin{proof}
我们给出的证明虽然直接，却略显繁琐。为说明思路，先解释较小 $n$ 时的
情形，再给出一般证明。例如，若 $n=0$，则 (3) 意味着
$P_1=P_0\times_BP_0$。由于环同态 $P_0\to B$ 满射，由《更多交换代数》
引理 \ref{more-algebra-lemma-fibre-product-finite-type}，该代数在 $A$
上有限型。

\medskip\noindent
若 $n=1$，则 (3) 意味着
$$
P_2 = \{(f_0, f_1, f_2) \in P_1^3 \mid
d_0f_0 = d_0f_1,\ d_1f_0 = d_0f_2,\ d_1f_1 = d_1f_2 \}
$$
其中各等式均在 $P_0$ 中成立。注意，三元组
$$
(d_0f_0, d_1f_0, d_1f_1) = (d_0f_1, d_0f_2, d_1f_2)
$$
是 $B$ 上纤维积 $P_0\times_BP_0\times_BP_0$ 的元素，因为同态
$d_i:P_1\to P_0$ 都是 $B$ 上的态射。因此得到同态
$$
\psi : P_2 \longrightarrow P_0 \times_B P_0 \times_B P_0
$$
$\psi$ 在元素
$(g_0,g_1,g_2)\in P_0\times_BP_0\times_BP_0$ 上的纤维，是满足
$(d_0,d_1)(f_0)=(g_0,g_1)$、$(d_0,d_1)(f_1)=(g_0,g_2)$ 以及
$(d_0,d_1)(f_2)=(g_1,g_2)$ 的 $1$-单形三元组
$(f_0,f_1,f_2)$ 所成的集合。由于 $P_\bullet\to B$ 是平凡 Kan
纤维化，同态 $(d_0,d_1):P_1\to P_0\times_BP_0$ 满射。因此
$P_2$ 位于笛卡儿图
$$
\xymatrix{
P_2 \ar[d] \ar[r] & P_1^3 \ar[d] \\
P_0 \times_B P_0 \times_B P_0 \ar[r] & (P_0 \times_B P_0)^3
}
$$
中。由《更多交换代数》引理
\ref{more-algebra-lemma-formal-consequence} 即得结论。一般情形类似，
但需要稍多一些记号。

\medskip\noindent
现在处理 $n>1$ 的情形。由《单纯形》引理
\ref{simplicial-lemma-cosk-above-object}，条件
$P_\bullet=\text{cosk}_n\text{sk}_nP_\bullet$ 蕴含同一等式在单纯形
$A$-代数范畴中成立，因而也在集合范畴中成立（因为从 $A$-代数到集合
的遗忘函子与极限交换）。所以
$$
P_{n + 1} =
\Mor(\Delta[n + 1], P_\bullet) =
\Mor(\text{sk}_n \Delta[n + 1], \text{sk}_n P_\bullet)
$$
这里使用了《单纯形》引理 \ref{simplicial-lemma-simplex-map} 和等式
(\ref{simplicial-equation-cosk})。我们将对
$1\leq k<m\leq n+1$ 归纳证明环
$$
Q_{k, m} = \Mor(\text{sk}_k \Delta[m], \text{sk}_k P_\bullet)
$$
在 $A$ 上有限型。当 $k=1$ 且 $1<m\leq n+1$ 时，这与上面对
$n=1$ 的讨论完全类似。具体而言，存在笛卡儿图
$$
\xymatrix{
Q_{1, m} \ar[d] \ar[r] & P_1^N \ar[d] \\
P_0 \times_B \ldots \times_B P_0 \ar[r] & (P_0 \times_B P_0)^N
}
$$
其中 $N={m+1\choose2}$。同前可得结论。

\medskip\noindent
设 $1\leq k_0\leq n$，并假设对所有满足
$1\leq k\leq k_0$ 与 $k<m\leq n+1$ 的指标，$Q_{k,m}$ 在 $A$
上有限型。对 $k_0+1<m\leq n+1$，我们断言存在笛卡儿方块
$$
\xymatrix{
Q_{k_0 + 1, m} \ar[d] \ar[r] & P_{k_0 + 1}^N \ar[d] \\
Q_{k_0, m} \ar[r] & Q_{k_0, k_0 + 1}^N
}
$$
其中 $N$ 是 $\Delta[m]$ 的非退化 $(k_0+1)$-单形的个数。为看出这一点，
把 $Q_{k_0+1,m}$ 的元素看作从 $\Delta[m]$ 的 $(k_0+1)$-骨架到
$P_\bullet$ 的同态 $f$。把 $f$ 限制到 $k_0$-骨架，得到图中的左侧
竖直同态；把它限制到每个非退化 $(k_0+1)$-单形，则得到上方水平
同态。给定这样的 $f$，等价于给定它在 $k_0$-骨架和每个非退化
$(k_0+1)$-面上的限制，并要求这些限制在交集上一致；这正是该图所
表达的内容。此外，由 $P_\bullet\to B$ 是平凡 Kan 纤维化可知，同态
$$
P_{k_0} \to Q_{k_0, k_0 + 1} = \Mor(\partial \Delta[k_0 + 1], P_\bullet)
$$
满射，因为当 $k_0\geq1$ 时，每个同态
$\partial\Delta[k_0+1]\to B$ 都可延拓为
$\Delta[k_0+1]\to B$（这里略去关于常值单纯形集的一小段论证）。
% P11-E0020: the frozen source uses P_{k_0} for the boundary-filling map; retained here. See control/ERRATA_CANDIDATES.jsonl.
由归纳假设，环 $Q_{k_0,m}$ 和 $Q_{k_0,k_0+1}$ 都是有限型
$A$-代数。再应用《更多交换代数》引理
\ref{more-algebra-lemma-formal-consequence}，可知
$Q_{k_0+1,m}$ 也如此。
\end{proof}

\begin{proposition}
\label{cotangent-proposition-polynomial}
设 $A$ 是诺特环，$A\to B$ 是有限型环同态。存在带有增广
$\epsilon:P_\bullet\to B$ 的单纯形 $A$-代数 $P_\bullet$，使每个
$P_n$ 都是 $A$ 上的有限型多项式代数，并且 $\epsilon$ 是单纯形集的
平凡 Kan 纤维化。
\end{proposition}

\begin{proof}
令 $\mathcal{A}$ 为 $A$-代数同态 $C\to B$ 所成的范畴。本证明中的
单纯形对象、骨架函子和余骨架函子均取在此范畴中。

\medskip\noindent
选取 $A$ 上的有限型多项式代数 $P_0$ 以及满射 $P_0\to B$。作为
第一次近似，取 $P_\bullet=\text{cosk}_0(P_0)$。换言之，
$P_\bullet$ 是各项为
$P_n=P_0\times_A\ldots\times_AP_0$ 的单纯形 $A$-代数。
% P11-E0019: the frozen source uses fibre products over A in this coskeleton; retained here. See control/ERRATA_CANDIDATES.jsonl.
（在证明的最后一段，这个单纯形对象将记作 $P^0_\bullet$。）由
《单纯形》引理 \ref{simplicial-lemma-cosk-minus-one-equivalence}，同态
$P_\bullet\to B$ 是单纯形集的平凡 Kan 纤维化。还要注意
$P_\bullet=\text{cosk}_0\text{sk}_0P_\bullet$。

\medskip\noindent
假设对某个 $n\geq0$ 已构造出 $P_\bullet$（在证明的最后一段，它将
记作 $P^n_\bullet$），满足
\begin{enumerate}
\item[(a)] $P_\bullet\to B$ 是单纯形集的平凡 Kan 纤维化；
\item[(b)] 当 $0\leq k\leq n$ 时，$P_k$ 是有限生成多项式代数；
以及
\item[(c)] $P_\bullet = \text{cosk}_n \text{sk}_n P_\bullet$
\end{enumerate}
由引理 \ref{cotangent-lemma-polynomial}，可以找到 $A$ 上的有限生成多项式代数
$Q$ 以及满射 $Q\to P_{n+1}$。由于 $P_n$ 是多项式代数，$A$-代数
同态 $s_i:P_n\to P_{n+1}$ 可提升为同态 $s'_i:P_n\to Q$。令
$d'_j:Q\to P_n$ 等于 $Q\to P_{n+1}$ 与
$d_j:P_{n+1}\to P_n$ 的复合。规定当 $k\leq n$ 时
$P'_k=P_k$，并令 $P'_{n+1}=Q$；在次数 $k\leq n-1$ 上取态射
$d'_i=d_i$ 和 $s'_i=s_i$，在次数 $n$ 上使用态射 $d'_j$ 与
$s'_i$，由此得到 $\mathcal{A}$ 中的截断单纯形对象 $P'_\bullet$。
利用 $\text{cosk}_{n+1}$ 把它延拓为 $\mathcal{A}$ 中的完整单纯形
对象。由余骨架函子的函子性，存在单纯形对象态射
$P'_\bullet\to P_\bullet$，延拓给定的 $(n+1)$-截断单纯形对象态射。
（在证明的最后一段，此态射将记作
$P^{n+1}_\bullet\to P^n_\bullet$。）

\medskip\noindent
注意，把 $n$ 替换为 $n+1$ 后，$P'_\bullet$ 满足条件 (b) 与 (c)。
我们断言同态 $P'_\bullet\to P_\bullet$ 满足《单纯形》引理
\ref{simplicial-lemma-section} 中把 $n$ 替换为 $n+1$ 后的假设
(1)、(2)、(3) 与 (4)。条件 (1) 和 (2) 由构造成立。由《单纯形》
引理 \ref{simplicial-lemma-cosk-above-object}，有
$P_\bullet = \text{cosk}_{n + 1}\text{sk}_{n + 1}P_\bullet$
且
$P'_\bullet = \text{cosk}_{n + 1}\text{sk}_{n + 1}P'_\bullet$
这些等式不仅在 $\mathcal{A}$ 中成立，也在 $A$-代数范畴中成立，
从而在集合范畴中成立（因为从 $A$-代数到集合的遗忘函子与所有极限
交换）。这证明了 (3) 与 (4)。因此该引理可用，且
$P'_\bullet\to P_\bullet$ 是平凡 Kan 纤维化。由《单纯形》引理
\ref{simplicial-lemma-trivial-kan-composition}，可知
$P'_\bullet\to B$ 是平凡 Kan 纤维化，故 (a) 也成立。

\medskip\noindent
为完成证明，取单纯形代数序列
$$
\ldots \to P^2_\bullet \to P^1_\bullet \to P^0_\bullet
$$
的逆极限 $P_\bullet=\lim P^n_\bullet$。由《单纯形》引理
\ref{simplicial-lemma-limit-trivial-kan}，同态 $P_\bullet\to B$ 是
平凡 Kan 纤维化。另一方面，上述构造在每个次数上最终稳定为一个固定的
有限生成多项式代数，正合所需。
\end{proof}

\begin{lemma}
\label{cotangent-lemma-pi-shriek-finite}
设 $A$ 是诺特环，$A\to B$ 是有限型环同态，并令 $\pi$、
$\underline{B}$ 如 (\ref{cotangent-equation-pi}) 所示。若 $\mathcal{F}$ 是
$\underline{B}$-模，并且对每个
$\alpha:P=A[x_1,\ldots,x_n]\to B$，$\mathcal{F}(P,\alpha)$ 都是
有限 $B$-模，则 $L\pi_!(\mathcal{F})$ 的各上同调模都是有限 $B$-模。
\end{lemma}

\begin{proof}
由引理 \ref{cotangent-lemma-identify-pi-shriek} 与命题
\ref{cotangent-proposition-polynomial}，可以用一个由 $\mathcal{F}$ 在有限型
多项式代数上的取值构成的复形来计算 $L\pi_!(\mathcal{F})$。
\end{proof}

\begin{lemma}
\label{cotangent-lemma-cotangent-finite}
设 $A$ 是诺特环，$A\to B$ 是有限型环同态。那么对所有
$n\in\mathbf{Z}$，$H^n(L_{B/A})$ 都是有限 $B$-模。
\end{lemma}

\begin{proof}
应用引理 \ref{cotangent-lemma-compute-cotangent-complex} 和
\ref{cotangent-lemma-pi-shriek-finite}。
\end{proof}

\begin{remark}[预解]
\label{cotangent-remark-resolution}
设 $A\to B$ 为任意环同态。若增广单纯形 $A$-代数
$\epsilon:P_\bullet\to B$ 的每个 $P_n$ 都是多项式代数，并且
$\epsilon$ 是单纯形集的平凡 Kan 纤维化，则称它为 $B$ 在 $A$ 上的
{\it 预解}。若 $P_\bullet\to B$ 是单纯形 $A$-代数的增广，每个
$P_n$ 都是满射到 $B$ 的多项式代数，则下列条件等价：
\begin{enumerate}
\item $\epsilon:P_\bullet\to B$ 是 $B$ 在 $A$ 上的预解；
\item $\epsilon:P_\bullet\to B$ 在相伴复形上是拟同构；
\item $\epsilon:P_\bullet\to B$ 诱导单纯形集的同伦等价。
\end{enumerate}
为此应用《单纯形》中的引理
\ref{simplicial-lemma-trivial-kan-homotopy},
\ref{simplicial-lemma-homotopy-equivalence}，以及
\ref{simplicial-lemma-qis-simplicial-abelian-groups}.
$B$ 在 $A$ 上的预解 $P_\bullet$ 给出 $\mathcal{C}_{B/A}$ 中的余单纯形
对象 $U_\bullet$，如《位点上的上同调》引理
\ref{sites-cohomology-lemma-compute-by-cosimplicial-resolution} 所述，
并且
$$
L\pi_!\mathcal{F} = \mathcal{F}(P_\bullet)
$$
关于 $\mathcal{F}$ 具有函子性，见引理
\ref{cotangent-lemma-identify-pi-shriek}。命题 \ref{cotangent-proposition-polynomial} 证明的
形式部分表明预解存在。在引理 \ref{cotangent-lemma-pi-shriek-standard} 的第一种
证明中，我们还看到 $B$ 在 $A$ 上的标准预解确实是一个预解（故此术语
不会造成冲突）。不过，命题 \ref{cotangent-proposition-polynomial} 的证明无需
借助《单纯形》第 \ref{simplicial-section-standard} 节中的单纯形计算，
便已说明预解存在。而且，对{\it 任意}预解选择，均有典范同构
$$
L_{B/A} = \Omega_{P_\bullet/A} \otimes_{P_\bullet, \epsilon} B
$$
属于 $D(B)$；这是引理 \ref{cotangent-lemma-compute-cotangent-complex} 的结论。
任意选择预解的自由往往很有用。
\end{remark}

\begin{lemma}
\label{cotangent-lemma-O-homology-B-homology}
设 $A\to B$ 为环同态，并令 $\pi$、$\mathcal{O}$、$\underline{B}$
如 (\ref{cotangent-equation-pi}) 所示。对任意 $\mathcal{O}$-模 $\mathcal{F}$，有
$$
L\pi_!(\mathcal{F}) = L\pi_!(Li^*\mathcal{F}) =
L\pi_!(\mathcal{F} \otimes_\mathcal{O}^\mathbf{L} \underline{B})
$$
属于 $D(\textit{Ab})$。
\end{lemma}

\begin{proof}
只需验证《位点上的上同调》引理
\ref{sites-cohomology-lemma-O-homology-qis} 的假设对
$\mathcal{C}_{B/A}$ 上的 $\mathcal{O}\to\underline{B}$ 成立。以下不再
另行说明地使用注 \ref{cotangent-remark-resolution} 的结果。选取 $B$ 在 $A$
上的一个预解 $P_\bullet$，从而得到 $\mathcal{C}_{B/A}$ 中适当的余单纯形
对象 $U_\bullet$。由于 $P_\bullet\to B$ 在相伴的阿贝尔群复形上诱导
拟同构，故 $L\pi_!\mathcal{O}=B$。另一方面，
$L\pi_!\underline{B}$ 由 $\underline{B}(U_\bullet)=B$ 计算。这验证了
《位点上的上同调》引理
\ref{sites-cohomology-lemma-O-homology-qis} 的第二个假设，证明完成。
\end{proof}

\begin{lemma}
\label{cotangent-lemma-apply-O-B-comparison}
设 $A\to B$ 为环同态，并令 $\pi$、$\mathcal{O}$、$\underline{B}$
如 (\ref{cotangent-equation-pi}) 所示。则
$$
L\pi_!(\mathcal{O}) = L\pi_!(\underline{B}) = B
\quad\text{且}\quad
L_{B/A} = L\pi_!(\Omega_{\mathcal{O}/A} \otimes_\mathcal{O} \underline{B}) =
L\pi_!(\Omega_{\mathcal{O}/A})
$$
属于 $D(\textit{Ab})$。
\end{lemma}

\begin{proof}
这只是引理 \ref{cotangent-lemma-O-homology-B-homology} 的应用（而右侧第一个
等式是引理 \ref{cotangent-lemma-compute-cotangent-complex}）。
\end{proof}

\noindent
下面是基本三角的一个容易证明的特例。

\begin{lemma}
\label{cotangent-lemma-special-case-triangle}
设 $A\to B\to C$ 为环同态。若 $B$ 是 $A$ 上的多项式代数，则
$D(C)$ 中存在可区别三角
$L_{B/A} \otimes_B^\mathbf{L} C \to L_{C/A} \to L_{C/B} \to
L_{B/A} \otimes_B^\mathbf{L} C[1]$。
\end{lemma}

\begin{proof}
以下不再另行说明地使用注 \ref{cotangent-remark-resolution} 的观察。选取 $C$
在 $B$ 上的预解 $\epsilon:P_\bullet\to C$（例如标准预解）。因为
$B$ 是 $A$ 上的多项式代数，$P_\bullet$ 也是 $C$ 在 $A$ 上的预解。
因此 $L_{C/A}$ 由
$\Omega_{P_\bullet/A}\otimes_{P_\bullet,\epsilon}C$ 计算，而
$L_{C/B}$ 由 $\Omega_{P_\bullet/B}\otimes_{P_\bullet,\epsilon}C$
计算。对每个 $n$ 都有短正合列
$0 \to \Omega_{B/A} \otimes_B P_n \to \Omega_{P_n/A} \to \Omega_{P_n/B} \to 0$
（《交换代数》引理 \ref{algebra-lemma-ses-formally-smooth}），并且
$L_{B/A}=\Omega_{B/A}[0]$（引理
\ref{cotangent-lemma-cotangent-complex-polynomial-algebra}），故结论成立。
\end{proof}

\begin{example}
\label{cotangent-example-resolution-length-two}
设 $A\to B$ 为环同态。本例构造一个“显式”的、长度为 $2$ 的 $B$
在 $A$ 上的预解 $P_\bullet$。我们遵循命题
\ref{cotangent-proposition-polynomial} 证明中的步骤；另见注
\ref{cotangent-remark-resolution} 的讨论。

\medskip\noindent
选取满射 $P_0=A[u_i]\to B$，其中 $u_i$ 是一族变量。选取理想
$\Ker(P_0\to B)$ 的生成元 $f_t\in P_0$，$t\in T$。令
$P_1=A[u_i,x_t]$；面同态 $d_0$ 和 $d_1$ 是满足
$d_j(u_i)=u_i$、$d_0(x_t)=0$ 以及 $d_1(x_t)=f_t$ 的唯一
$A$-代数同态。同态 $s_0:P_0\to P_1$ 是满足 $s_0(u_i)=u_i$
的唯一 $A$-代数同态。显然
$$
P_1 \xrightarrow{d_0 - d_1} P_0 \to B \to 0
$$
正合；特别地，同态 $(d_0,d_1):P_1\to P_0\times_BP_0$ 满射。因此，
若以 $P_\bullet$ 记由 $P_0$、$P_1$、$d_0$、$d_1$ 与 $s_0$ 给出的
$1$-截断单纯形 $A$-代数，则增广
$\text{cosk}_1(P_\bullet)\to B$ 是平凡 Kan 纤维化。命题
\ref{cotangent-proposition-polynomial} 证明步骤的下一步，是选取多项式代数
$P_2$ 以及满射
$$
P_2 \longrightarrow \text{cosk}_1(P_\bullet)_2
$$
回忆
$$
\text{cosk}_1(P_\bullet)_2 =
\{(g_0, g_1, g_2) \in P_1^3 \mid d_0(g_0) = d_0(g_1),
d_1(g_0) = d_0(g_2), d_1(g_1) = d_1(g_2)\}
$$
把 $g_i\in P_1$ 看作关于 $x_t$ 的多项式，则这些条件为
$$
g_0(0) = g_1(0),\quad
g_0(f_t) = g_2(0),\quad
g_1(f_t) = g_2(f_t)
$$
因此 $\text{cosk}_1(P_\bullet)_2$ 包含元素
$y_t=(x_t,x_t,f_t)$ 和 $z_t=(0,x_t,x_t)$。
$\text{cosk}_1(P_\bullet)_2$ 中每个元素 $G$ 都可写成
$G=H+(0,0,g)$，其中 $H$ 属于同态
$A[u_i,y_t,z_t]\to\text{cosk}_1(P_\bullet)_2$ 的像，而
$g\in P_1$ 是常数项为零且在 $P_0$ 中满足 $g(f_t)=0$ 的多项式。
注意，
\begin{enumerate}
\item $g = x_t x_{t'} - f_t x_{t'}$；以及
\item 若 $r_t \in P_0$ 且 $P_0$ 中有 $\sum r_t f_t = 0$，则
$g = \sum r_t x_t$。
\end{enumerate}
都是 $P_1$ 中具有所需形式的元素。令
$$
Rel = \Ker(\bigoplus\nolimits_{t \in T} P_0 \longrightarrow P_0),\quad
(r_t) \longmapsto \sum r_tf_t
$$
令 $P_2=A[u_i,y_t,z_t,v_r,w_{t,t'}]$，其中 $r=(r_t)\in Rel$，
并取同态
$$
P_2 \longrightarrow \text{cosk}_1(P_\bullet)_2
$$
它由下列各式给出：$y_t\mapsto(x_t,x_t,f_t)$、
$z_t\mapsto(0,x_t,x_t)$、$v_r\mapsto(0,0,\sum r_tx_t)$，以及
$w_{t,t'}\mapsto(0,0,x_tx_{t'}-f_tx_{t'})$。一次计算（略去）表明
该同态满射。上面显示的同态选择确定了
$d_0,d_1,d_2:P_2\to P_1$。最后，命题
\ref{cotangent-proposition-polynomial} 证明中的步骤要求我们选取同态
$s_0,s_1:P_1\to P_2$，它们提升两个同态
$P_1\to\text{cosk}_1(P_\bullet)_2$。显然，可以取 $s_i$ 为由
$s_0(x_t)=y_t$ 与 $s_1(x_t)=z_t$ 确定的唯一 $A$-代数同态。
\end{example}








\section{函子性}
\label{cotangent-section-functoriality}

\noindent
本节考虑环同态的交换方块
\begin{equation}
\label{cotangent-equation-commutative-square}
\vcenter{
\xymatrix{
B \ar[r] & B' \\
A \ar[u] \ar[r] & A' \ar[u]
}
}
\end{equation}
我们断言，与此图相伴有典范的 $B$-线性复形同态
$$
L_{B/A} \longrightarrow L_{B'/A'}
$$
事实上，若 $P_\bullet\to B$ 是 $B$ 在 $A$ 上的标准预解，且
$P'_\bullet\to B'$ 是 $B'$ 在 $A'$ 上的标准预解，则存在与增广
$P_\bullet\to B$ 和 $P'_\bullet\to B'$ 相容的典范单纯形
$A$-代数同态 $P_\bullet\to P'_\bullet$。这可从《单纯形》第
\ref{simplicial-section-standard} 节对标准预解的构造看出；不过在当前
特例中，只需指出同态
$$
P_0 = A[B] \longrightarrow A'[B'] = P'_0,\quad
P_1 = A[A[B]] \longrightarrow A'[A'[B']] = P'_1,
$$
等等均由给定同态 $A\to A'$ 和 $B\to B'$ 给出。所需同态
$L_{B/A}\to L_{B'/A'}$ 随即来自相伴同态
$\Omega_{P_n/A}\to\Omega_{P'_n/A'}$。

\medskip\noindent
函子性同态还可描述如下。令
$\mathcal{C}=\mathcal{C}_{B/A}$ 与
$\mathcal{C}'=\mathcal{C}_{B'/A}'$ 为第
\ref{cotangent-section-compute-L-pi-shriek} 节所考虑的范畴。
% P11-E0021: the frozen source places the prime outside the B'/A subscript; retained here. See control/ERRATA_CANDIDATES.jsonl.
存在函子
$$
u : \mathcal{C} \longrightarrow \mathcal{C}',\quad
(P, \alpha) \longmapsto (P \otimes_A A', c \circ (\alpha \otimes 1))
$$
其中 $c:B\otimes_AA'\to B'$ 是显然的同态。按《位点上的上同调》例
\ref{sites-cohomology-example-morphism-categories} 的讨论，得到拓扑斯态射
$g:\Sh(\mathcal{C})\to\Sh(\mathcal{C}')$ 以及环化拓扑斯态射的交换图
\begin{equation}
\label{cotangent-equation-double-square}
\vcenter{
\xymatrix{
(\Sh(\mathcal{C}'), \underline{B}) \ar[d]_{\pi'} &
(\Sh(\mathcal{C}'), \underline{B'}) \ar[d]_{\pi'} \ar[l]^h &
(\Sh(\mathcal{C}), \underline{B'}) \ar[d]_\pi \ar[l]^g \\
(\Sh(*), B) &
(\Sh(*), B') \ar[l]_f &
(\Sh(*), B') \ar[l]
}
}
\end{equation}
这里 $h$ 在底层拓扑斯上是恒等态射，并在环层上由环同态 $B\to B'$
给出。由《位点上的上同调》注
\ref{sites-cohomology-remark-morphism-fibred-categories}，给定
$\mathcal{C}$ 上的 $\mathcal{F}$、$\mathcal{C}'$ 上的
$\mathcal{F}'$ 以及变换 $t:\mathcal{F}\to g^{-1}\mathcal{F}'$，
便得到典范同态
$L\pi_!(\mathcal{F})\to L\pi'_!(\mathcal{F}')$。把它应用于层
$$
\mathcal{F} : (P, \alpha) \mapsto \Omega_{P/A} \otimes_P B,\quad
\mathcal{F}' : (P', \alpha') \mapsto \Omega_{P'/A'} \otimes_{P'} B',
$$
以及由典范同态
$$
\Omega_{P/A} \otimes_P B \longrightarrow
\Omega_{P \otimes_A A'/A'} \otimes_{P \otimes_A A'} B'
$$
给出的变换 $t$，得到典范同态
$$
L\pi_!(\Omega_{\mathcal{O}/A} \otimes_\mathcal{O} \underline{B})
\longrightarrow
L\pi'_!(\Omega_{\mathcal{O}'/A'} \otimes_{\mathcal{O}'} \underline{B'})
$$
由引理 \ref{cotangent-lemma-compute-cotangent-complex}，这给出
$L_{B/A}\to L_{B'/A'}$。我们略去对此同态与上面通过单纯形预解定义的
同态一致性的验证。

\begin{lemma}
\label{cotangent-lemma-flat-base-change}
假设 (\ref{cotangent-equation-commutative-square}) 诱导拟同构
$B\otimes_A^\mathbf{L}A'=B'$。沿用
(\ref{cotangent-equation-double-square}) 的记号，并设
$\mathcal{F}'\in\textit{Ab}(\mathcal{C}')$，则
$L\pi_!(g^{-1}\mathcal{F}')=L\pi'_!(\mathcal{F}')$。
\end{lemma}

\begin{proof}
以下不再另行说明地使用注 \ref{cotangent-remark-resolution} 的结果。应用
《位点上的上同调》引理 \ref{sites-cohomology-lemma-get-it-now}。
设 $P_\bullet\to B$ 为预解。只要证明
$u(P_\bullet)=P_\bullet\otimes_AA'\to B'$ 是拟同构，结论即成立。
把 $P_\bullet$ 看作单纯形 $A$-模，则其相伴 $A$-模复形
$s(P_\bullet)$ 是 $B$ 的自由 $A$-模预解：确实，每个 $P_n$ 都是
自由 $A$-模，且 $s(P_\bullet)\to B$ 是拟同构。因此
$B\otimes_A^\mathbf{L}A'$ 由
$s(P_\bullet)\otimes_AA'=s(P_\bullet\otimes_AA')$ 计算。故本引理的
假设意味着 $\epsilon':P_\bullet\otimes_AA'\to B'$ 是拟同构。
\end{proof}

\noindent
特别地，当 $A\to A'$ 平坦且 $B'=B\otimes_AA'$ 时（平坦基变换），
下述引理适用。

\begin{lemma}
\label{cotangent-lemma-flat-base-change-cotangent-complex}
若 (\ref{cotangent-equation-commutative-square}) 诱导拟同构
$B\otimes_A^\mathbf{L}A'=B'$，则函子性同态诱导同构
$$
L_{B/A} \otimes_B^\mathbf{L} B' \longrightarrow L_{B'/A'}
$$
\end{lemma}

\begin{proof}
使用等式 (\ref{cotangent-equation-double-square}) 中引入的记号。我们有
$$
L_{B/A} \otimes_B^\mathbf{L} B' =
L\pi_!(\Omega_{\mathcal{O}/A} \otimes_\mathcal{O} \underline{B})
\otimes_B^\mathbf{L} B' =
L\pi_!(Lh^*(\Omega_{\mathcal{O}/A} \otimes_\mathcal{O} \underline{B}))
$$
第一个等式来自引理 \ref{cotangent-lemma-compute-cotangent-complex}，第二个来自
《位点上的上同调》引理 \ref{sites-cohomology-lemma-change-of-rings}。
由于 $\Omega_{\mathcal{O}/A}$ 是平坦 $\mathcal{O}$-模，
$\Omega_{\mathcal{O}/A}\otimes_\mathcal{O}\underline{B}$ 是平坦
$\underline{B}$-模。因此
$Lh^*(\Omega_{\mathcal{O}/A} \otimes_\mathcal{O} \underline{B}) =
\Omega_{\mathcal{O}/A} \otimes_\mathcal{O} \underline{B'}$
它等于
$g^{-1}(\Omega_{\mathcal{O}'/A'} \otimes_{\mathcal{O}'} \underline{B'})$
这可直接检验。再由引理 \ref{cotangent-lemma-flat-base-change}，并利用
$L_{B'/A'}$ 由
$L\pi'_!(\Omega_{\mathcal{O}'/A'}\otimes_{\mathcal{O}'}\underline{B'})$
计算，即得结论。
\end{proof}

\begin{remark}
\label{cotangent-remark-homotopy-triangle}
设给定方块 (\ref{cotangent-equation-commutative-square})，且存在使下图交换的
箭头 $\kappa:B\to A'$：
$$
\xymatrix{
B \ar[r]_\beta \ar[rd]_\kappa & B' \\
A \ar[u] \ar[r]^\alpha & A' \ar[u]
}
$$
我们断言，在此情形函子性同态 $P_\bullet\to P'_\bullet$ 与复合
$P_\bullet\to B\to A'\to P'_\bullet$ 同伦。事实上，利用 $\kappa$，
函子性同态分解为
$$
P_\bullet \to P_{A'/A', \bullet} \to P'_\bullet
$$
其中 $P_{A'/A',\bullet}$ 是 $A'$ 在 $A'$ 上的标准预解。由于 $A'$
是 $A'$ 上以空集为变量集的多项式代数，由《单纯形》引理
\ref{simplicial-lemma-standard-simplicial-homotopy} 可知，增广
$\epsilon_{A'/A'}:P_{A'/A',\bullet}\to A'$ 是单纯形环的同伦等价。
注意，该引理证明中构造的同伦逆同态
$c:A'\to P_{A'/A',\bullet}$ 正是结构同态。因此，下列两个复合
$$
\xymatrix{
P_\bullet \ar[r] &
P_{A'/A', \bullet} \ar@<1ex>[rr]^{\text{id}}
\ar@<-1ex>[rr]_{c \circ \epsilon_{A'/A'}} & &
P_{A'/A', \bullet} \ar[r] &
P'_\bullet
}
$$
就是上面讨论的两个同态，并且二者同伦（《单纯形》注
\ref{simplicial-remark-homotopy-pre-post-compose}），故得到所需结论。
由于第二个同态 $P_\bullet\to P'_\bullet$ 在
$\Omega_{P_\bullet/A}\to\Omega_{P'_\bullet/A'}$ 上诱导零同态，
可知在此情形函子性同态 $L_{B/A}\to L_{B'/A'}$ 零同伦。
\end{remark}

\begin{lemma}
\label{cotangent-lemma-cotangent-complex-product}
设 $A\to B$ 与 $A\to C$ 为环同态。那么同态
$L_{B\times C/A}\to L_{B/A}\oplus L_{C/A}$ 在
$D(B\times C)$ 中是同构。
\end{lemma}

\begin{proof}
虽然本引理可由基本三角推出，我们仍给出一个直接且初等的证明。把环同态
$A\to B\times C$ 分解为 $A\to A[x]\to B\times C$，其中
$x\mapsto(1,0)$。由引理 \ref{cotangent-lemma-special-case-triangle}，有可区别三角
$$
L_{A[x]/A} \otimes_{A[x]}^\mathbf{L} (B \times C) \to L_{B \times C/A} \to
L_{B \times C/A[x]} \to L_{A[x]/A} \otimes_{A[x]}^\mathbf{L} (B \times C)[1]
$$
属于 $D(B\times C)$。类似地，有可区别三角
$$
\begin{matrix}
L_{A[x]/A} \otimes_{A[x]}^\mathbf{L} B \to L_{B/A} \to L_{B/A[x]}
\to L_{A[x]/A} \otimes_{A[x]}^\mathbf{L} B[1] \\
L_{A[x]/A} \otimes_{A[x]}^\mathbf{L} C \to L_{C/A} \to L_{C/A[x]}
\to L_{A[x]/A} \otimes_{A[x]}^\mathbf{L} C[1]
\end{matrix}
$$
因此，只需证明 $A[x]$ 上的 $B\times C$ 的结论。注意，
$A[x]\to A[x,x^{-1}]$ 平坦，且
$(B\times C)\otimes_{A[x]}A[x,x^{-1}]
=B\otimes_{A[x]}A[x,x^{-1}]$，并且
$C\otimes_{A[x]}A[x,x^{-1}]=0$。所以由基变换（引理
\ref{cotangent-lemma-flat-base-change-cotangent-complex}），同态
$L_{B\times C/A[x]}\to L_{B/A[x]}\oplus L_{C/A[x]}$ 在逆转 $x$
后成为同构。同理可证，该同态在逆转 $x-1$ 后也成为同构。
本引理得证。
\end{proof}




\section{基本三角}
\label{cotangent-section-triangle}

\noindent
本节考虑一列环同态 $A\to B\to C$。我们的目标是证明，该三角诱导
可区别三角
\begin{equation}
\label{cotangent-equation-triangle}
L_{B/A} \otimes_B^\mathbf{L} C \to L_{C/A} \to L_{C/B} \to
L_{B/A} \otimes_B^\mathbf{L} C[1]
\end{equation}
属于 $D(C)$。这将在命题 \ref{cotangent-proposition-triangle} 中证明。另一种方法
见注 \ref{cotangent-remark-triangle}。

\medskip\noindent
考虑范畴 $\mathcal{C}_{C/B/A}$，它是如下范畴的{\bf 反范畴}：其对象
是 $(P\to B,Q\to C)$，其中
\begin{enumerate}
\item $P$ 是 $A$ 上的多项式代数；
\item $P\to B$ 是 $A$-代数同态；
\item $Q$ 是 $P$ 上的多项式代数；
\item $Q\to C$ 是 $P$-代数同态。
\end{enumerate}
取反范畴是因为我们希望把 $(P\to B,Q\to C)$ 看作对应于交换图
$$
\xymatrix{
\Spec(C) \ar[d] \ar[r] & \Spec(Q) \ar[d] \\
\Spec(B) \ar[d] \ar[r] & \Spec(P) \ar[dl] \\
\Spec(A)
}
$$
令 $\mathcal{C}_{B/A}$、$\mathcal{C}_{C/A}$、$\mathcal{C}_{C/B}$
为第 \ref{cotangent-section-compute-L-pi-shriek} 节所考虑的范畴。存在函子
$$
\begin{matrix}
u_1 : \mathcal{C}_{C/B/A} \to \mathcal{C}_{B/A}, &
(P \to B, Q \to C) \mapsto (P \to B) \\
u_2 : \mathcal{C}_{C/B/A} \to \mathcal{C}_{C/A}, &
(P \to B, Q \to C) \mapsto (Q \to C) \\
u_3 : \mathcal{C}_{C/B/A} \to \mathcal{C}_{C/B}, &
(P \to B, Q \to C) \mapsto (Q \otimes_P B \to C)
\end{matrix}
$$
这些函子诱导相应的拓扑斯态射 $g_i$。记
$\mathcal{O}_i=g_i^{-1}\mathcal{O}$，从而得到环化拓扑斯态射
\begin{equation}
\label{cotangent-equation-three-maps}
\begin{matrix}
g_1 : (\Sh(\mathcal{C}_{C/B/A}), \mathcal{O}_1)
\longrightarrow (\Sh(\mathcal{C}_{B/A}), \mathcal{O}) \\
g_2 : (\Sh(\mathcal{C}_{C/B/A}), \mathcal{O}_2)
\longrightarrow (\Sh(\mathcal{C}_{C/A}), \mathcal{O}) \\
g_3 : (\Sh(\mathcal{C}_{C/B/A}), \mathcal{O}_3)
\longrightarrow (\Sh(\mathcal{C}_{C/B}), \mathcal{O})
\end{matrix}
\end{equation}
记
$\pi : \Sh(\mathcal{C}_{C/B/A}) \to \Sh(*)$,
$\pi_1 : \Sh(\mathcal{C}_{B/A}) \to \Sh(*)$,
$\pi_2 : \Sh(\mathcal{C}_{C/A}) \to \Sh(*)$，以及
$\pi_3 : \Sh(\mathcal{C}_{C/B}) \to \Sh(*)$,
于是 $\pi=\pi_i\circ g_i$。利用引理
\ref{cotangent-lemma-compute-cotangent-complex} 中对余切复形的等同以及下列
引理，我们将得到所需的可区别三角。

\begin{lemma}
\label{cotangent-lemma-triangle-ses}
沿用 (\ref{cotangent-equation-three-maps}) 的记号，令
$$
\begin{matrix}
\Omega_1 = \Omega_{\mathcal{O}/A} \otimes_\mathcal{O} \underline{B}
\text{，定义在 }\mathcal{C}_{B/A}\text{ 上} \\
\Omega_2 = \Omega_{\mathcal{O}/A} \otimes_\mathcal{O} \underline{C}
\text{，定义在 }\mathcal{C}_{C/A}\text{ 上} \\
\Omega_3 = \Omega_{\mathcal{O}/B} \otimes_\mathcal{O} \underline{C}
\text{，定义在 }\mathcal{C}_{C/B}\text{ 上}
\end{matrix}
$$
那么在 $\mathcal{C}_{C/B/A}$ 上存在典范的 $\underline{C}$-模层短正合列
$$
0 \to g_1^{-1}\Omega_1 \otimes_{\underline{B}} \underline{C} \to
g_2^{-1}\Omega_2 \to
g_3^{-1}\Omega_3 \to 0
$$
它位于 $\mathcal{C}_{C/B/A}$ 上。
\end{lemma}

\begin{proof}
回忆 $g_i^{-1}$ 只需与 $u_i$ 预复合即可得到。给定对象
$U=(P\to B,Q\to C)$，有分裂短正合列
$$
0 \to \Omega_{P/A} \otimes Q \to \Omega_{Q/A} \to \Omega_{Q/P} \to 0
$$
例如这可由《交换代数》引理
\ref{algebra-lemma-ses-formally-smooth} 得出。在 $Q$ 上与 $C$
作张量积，得到短正合列
$$
0 \to \Omega_{P/A} \otimes C \to \Omega_{Q/A} \otimes C \to
\Omega_{Q/P} \otimes C \to 0
$$
有 $\Omega_{P/A}\otimes C=\Omega_{P/A}\otimes B\otimes C$，故这是
$g_1^{-1}\Omega_1\otimes_{\underline{B}}\underline{C}$ 在 $U$ 上的值。
模 $\Omega_{Q/A}\otimes C$ 是 $g_2^{-1}\Omega_2$ 在 $U$ 上的值。
由《交换代数》引理 \ref{algebra-lemma-differentials-base-change}，有
$\Omega_{Q/P}\otimes C=\Omega_{Q\otimes_PB/B}\otimes C$，故这是
$g_3^{-1}\Omega_3$ 在 $U$ 上的值。因此，本引理的短正合列来自把
最后显示的短正合列赋给每个 $U$。
\end{proof}

\begin{lemma}
\label{cotangent-lemma-polynomial-on-top}
沿用 (\ref{cotangent-equation-three-maps}) 的记号，假设 $C$ 是 $B$ 上的
多项式代数。那么对 $\mathcal{C}_{C/B}$ 上任意阿贝尔层
$\mathcal{F}$，有
$L\pi_!(g_3^{-1}\mathcal{F}) = L\pi_{3, !}\mathcal{F} = \pi_{3, !}\mathcal{F}$
\end{lemma}

\begin{proof}
把 $C$ 写成 $B[E]$，其中 $E$ 为某个集合。选取 $B$ 在 $A$ 上的
预解 $P_\bullet\to B$。对每个 $n$，考虑
$\mathcal{C}_{C/B/A}$ 的对象
$U_n=(P_n\to B,P_n[E]\to C)$。那么 $U_\bullet$ 是
$\mathcal{C}_{C/B/A}$ 中的余单纯形对象。注意，
$u_3(U_\bullet)$ 是 $\mathcal{C}_{C/B}$ 中取值为 $(C\to C)$ 的常值
余单纯形对象。我们将证明 $\mathcal{C}_{C/B/A}$ 的对象 $U_\bullet$
满足《位点上的上同调》引理
\ref{sites-cohomology-lemma-compute-by-cosimplicial-resolution} 的假设。
这便蕴含本引理，因为它说明 $L\pi_!(g_3^{-1}\mathcal{F})$ 由取值为
$\mathcal{F}(C\to C)$ 的常值单纯形阿贝尔群计算；而由引理
\ref{cotangent-lemma-pi-lower-shriek-polynomial-algebra}，这正是
$L\pi_{3,!}\mathcal{F}=\pi_{3,!}\mathcal{F}$ 的值。

\medskip\noindent
设 $U=(\beta:P\to B,\gamma:Q\to C)$ 为 $\mathcal{C}_{C/B/A}$ 的
对象。由此范畴的定义，可把 $P$ 写成 $A[S]$，把 $Q$ 写成
$A[S\amalg T]$。须证明
$$
\Mor_{\mathcal{C}_{C/B/A}}(U_\bullet, U)
$$
与单点单纯形集 $*$ 同伦等价。注意，该单纯形集是乘积
$$
\prod\nolimits_{s \in S} F_s \times \prod\nolimits_{t \in T} F'_t
$$
其中 $F_s$ 是对应于
$U_s = (A[\{s\}] \to B, A[\{s\}] \to C)$
的单纯形集，$F'_t$ 是对应于
$U_t=(A\to B,A[\{t\}]\to C)$ 的单纯形集。确实，在
$\mathcal{C}_{C/B/A}$ 中，对象 $U$ 是乘积
$\prod U_s\times\prod U_t$。只需每个 $F_s$ 与 $F'_t$ 都与 $*$
同伦等价，见《单纯形》引理
\ref{simplicial-lemma-products-homotopy}。对 $F_s$，结论来自
$P_\bullet\to B$ 作为预解是平凡 Kan 纤维化，而 $F_s$ 是该同态在
$\beta(s)$ 上的纤维。（使用《单纯形》引理
\ref{simplicial-lemma-trivial-kan-base-change} 和
\ref{simplicial-lemma-trivial-kan-homotopy}。）$F'_t$ 的情形更有意思。
这里须说明同态
$$
P_\bullet[E] \longrightarrow C = B[E]
$$
在 $\gamma(t)\in C$ 上的纤维与一点同伦等价。事实上，我们将证明该
同态是平凡 Kan 纤维化。确实，$P_\bullet\to B$ 是平凡 Kan 纤维化。
% P11-E0022: the frozen English source says “trivial can fibration”; Chinese renders the standard term Kan and the defect is logged. See control/ERRATA_CANDIDATES.jsonl.
对任意环 $R$，有
$$
R[E] =
\colim_{\Sigma \subset \text{Map}(E, \mathbf{Z}_{\geq 0})\text{ 有限}}
\prod\nolimits_{I \in \Sigma} R
$$
（滤余极限）。所以，上面显示的单纯形集同态是平凡 Kan 纤维化的滤余
极限；由《单纯形》引理
\ref{simplicial-lemma-filtered-colimit-trivial-kan}，它仍是平凡 Kan
纤维化。
\end{proof}

\begin{lemma}
\label{cotangent-lemma-triangle-compute-lower-shriek}
沿用 (\ref{cotangent-equation-three-maps}) 的记号，对 $i=1,2,3$ 有
$Lg_{i,!}\circ g_i^{-1}=\text{id}$，因而对 $i=1,2,3$ 也有
$L\pi_!\circ g_i^{-1}=L\pi_{i,!}$。
\end{lemma}

\begin{proof}
先证 $i=1$。我们断言函子 $\mathcal{C}_{C/B/A}$ 是
$\mathcal{C}_{B/A}$ 上的纤维化范畴。设给定
$(P\to B,Q\to C)$ 以及 $\mathcal{C}_{B/A}$ 中的态射
$(P'\to B)\to(P\to B)$。回忆，这意味着存在与到 $B$ 的同态相容的
$A$-代数同态 $P\to P'$。令 $Q'=Q\otimes_PP'$，并赋予其诱导的到
$C$ 的同态。于是态射
$$
(P' \to B, Q' \to C) \longrightarrow (P \to B, Q \to C)
$$
在 $\mathcal{C}_{C/B/A}$ 中（再次注意箭头反向），并且它是
$\mathcal{C}_{C/B/A}$ 中位于 $\mathcal{C}_{B/A}$ 上的强笛卡儿态射。
还要注意，$u_1$ 在 $P\to B$ 上的纤维范畴是
$\mathcal{C}_{C/P}$。令 $\mathcal{F}$ 为 $\mathcal{C}_{B/A}$ 上的
阿贝尔层。由于这是纤维化范畴，可以应用《位点上的上同调》引理
\ref{sites-cohomology-lemma-compute-left-derived-pi-shriek}。因此，
$L_ng_{1,!}g_1^{-1}\mathcal{F}$ 是把
$U\in\Ob(\mathcal{C}_{B/A})$ 送到 $g_1^{-1}\mathcal{F}$ 在 $U$
上纤维范畴的限制之第 $n$ 个同调的（预）层。由于这些限制均为常值，
结合上面对纤维范畴的等同，所需结论由引理
\ref{cotangent-lemma-pi-lower-shriek-constant-sheaf} 得出。

\medskip\noindent
现在处理 $i=2$。我们断言 $\mathcal{C}_{C/B/A}$ 是
$\mathcal{C}_{C/A}$ 上的纤维化范畴。设给定
$(P\to B,Q\to C)$ 以及 $\mathcal{C}_{C/A}$ 中的态射
$(Q'\to C)\to(Q\to C)$。回忆，这意味着存在与到 $C$ 的同态相容的
$B$-代数同态 $Q\to Q'$。于是
% P11-E0023: the frozen source calls Q->Q' a B-algebra map although C_{C/A} morphisms are A-algebra maps; retained here. See control/ERRATA_CANDIDATES.jsonl.
$$
(P \to B, Q' \to C) \longrightarrow (P \to B, Q \to C)
$$
是 $\mathcal{C}_{C/B/A}$ 中位于 $\mathcal{C}_{C/A}$ 上的强笛卡儿
态射。注意，$u_2$ 在 $Q\to C$ 上的纤维范畴有终对象（留意箭头反向），
即 $(A\to B,Q\to C)$。令 $\mathcal{F}$ 为
$\mathcal{C}_{C/A}$ 上的阿贝尔层。由于这是纤维化范畴，可以应用
《位点上的上同调》引理
\ref{sites-cohomology-lemma-compute-left-derived-pi-shriek}。因此，
$L_ng_{2,!}g_2^{-1}\mathcal{F}$ 是把
$U\in\Ob(\mathcal{C}_{C/A})$ 送到 $g_1^{-1}\mathcal{F}$ 在 $U$
上纤维范畴的限制之第 $n$ 个同调的（预）层。
% P11-E0024: frozen source uses g_1^{-1} in the i=2 case; retained here. See control/ERRATA_CANDIDATES.jsonl.
由于这些限制均为常值，且所有纤维范畴都有终对象，所需结论由
《位点上的上同调》引理 \ref{sites-cohomology-lemma-initial-final} 得出。

\medskip\noindent
最后处理 $i=3$。在此情形，把《位点上的上同调》引理
\ref{sites-cohomology-lemma-compute-left-derived-g-shriek} 应用于
$u=u_3:\mathcal{C}_{C/B/A}\to\mathcal{C}_{C/B}$ 以及
$\mathcal{F}'=g_3^{-1}\mathcal{F}$，其中 $\mathcal{F}$ 是
$\mathcal{C}_{C/B}$ 上的阿贝尔层。设
$U=(\overline{Q}\to C)$ 为 $\mathcal{C}_{C/B}$ 的对象。那么
$\mathcal{I}_U=\mathcal{C}_{\overline{Q}/B/A}$（再次留意箭头反向）。
层 $\mathcal{F}'_U$ 由规则
$(P\to B,Q\to\overline{Q})\mapsto
\mathcal{F}(Q\otimes_PB\to C)$ 给出。换言之，该层是
$\mathcal{C}_{\overline{Q}/C}$ 上一个层沿态射
$\Sh(\mathcal{C}_{\overline{Q}/B/A})
\to\Sh(\mathcal{C}_{\overline{Q}/B})$ 的拉回。
% P11-E0025: frozen source says the sheaf is on C_{Qbar/C} although the displayed morphism and application require C_{Qbar/B}; retained here. See control/ERRATA_CANDIDATES.jsonl.
因此，引理 \ref{cotangent-lemma-polynomial-on-top} 表明，当 $n>0$ 时
$H_n(\mathcal{I}_U,\mathcal{F}'_U)=0$，而当 $n=0$ 时它等于
$\mathcal{F}(\overline{Q}\to C)$。上述《位点上的上同调》引理
\ref{sites-cohomology-lemma-compute-left-derived-g-shriek} 蕴含
$Lg_{3,!}(g_3^{-1}\mathcal{F})=\mathcal{F}$，证明完成。
\end{proof}

\begin{proposition}
\label{cotangent-proposition-triangle}
设 $A\to B\to C$ 为环同态。存在典范可区别三角
$$
L_{B/A} \otimes_B^\mathbf{L} C \to L_{C/A} \to L_{C/B} \to
L_{B/A} \otimes_B^\mathbf{L} C[1]
$$
属于 $D(C)$。
\end{proposition}

\begin{proof}
考虑引理 \ref{cotangent-lemma-triangle-ses} 的层短正合列，并应用导出函子
$L\pi_!$，得到可区别三角
$$
L\pi_!(g_1^{-1}\Omega_1 \otimes_{\underline{B}} \underline{C}) \to
L\pi_!(g_2^{-1}\Omega_2) \to
L\pi_!(g_3^{-1}\Omega_3) \to
L\pi_!(g_1^{-1}\Omega_1 \otimes_{\underline{B}} \underline{C})[1]
$$
属于 $D(C)$。由引理 \ref{cotangent-lemma-triangle-compute-lower-shriek} 和
\ref{cotangent-lemma-compute-cotangent-complex}，第二项与第三项分别等于
$L_{C/A}$ 与 $L_{C/B}$，而第一项等于
$$
L\pi_{1, !}(\Omega_1 \otimes_{\underline{B}} \underline{C}) =
L\pi_{1, !}(\Omega_1) \otimes_B^\mathbf{L} C =
L_{B/A} \otimes_B^\mathbf{L} C
$$
第一个等式来自《位点上的上同调》引理
\ref{sites-cohomology-lemma-change-of-rings}（以及 $\Omega_1$ 作为
$\underline{B}$-模层的平坦性），第二个等式来自引理
\ref{cotangent-lemma-compute-cotangent-complex}。
\end{proof}

\begin{remark}
\label{cotangent-remark-triangle}
下面概述基本三角存在性的另一种、或许更简单的证明。设
$A\to B\to C$ 为环同态，并假设 $B\to C$ 单射。令
$P_\bullet\to B$ 为 $B$ 在 $A$ 上的标准预解，令
$Q_\bullet\to C$ 为 $C$ 在 $B$ 上的标准预解。图示如下：
% P11-E0026: the frozen Q-row below uses A-polynomial towers although Q is declared the standard resolution over B; retained here. See control/ERRATA_CANDIDATES.jsonl.
$$
\xymatrix{
P_\bullet : &
A[A[A[B]]] \ar[d]
\ar@<2ex>[r]
\ar@<0ex>[r]
\ar@<-2ex>[r]
&
A[A[B]] \ar[d]
\ar@<1ex>[r]
\ar@<-1ex>[r]
\ar@<1ex>[l]
\ar@<-1ex>[l]
&
A[B] \ar[d] \ar@<0ex>[l] \ar[r] &
B \\
Q_\bullet : &
A[A[A[C]]]
\ar@<2ex>[r]
\ar@<0ex>[r]
\ar@<-2ex>[r]
&
A[A[C]]
\ar@<1ex>[r]
\ar@<-1ex>[r]
\ar@<1ex>[l]
\ar@<-1ex>[l]
&
A[C] \ar@<0ex>[l] \ar[r] &
C
}
$$
注意，由于 $B\to C$ 单射，对每个 $n$，环 $Q_n$ 都是 $P_n$ 上的
多项式代数。因此得到 $\mathcal{C}_{C/B/A}$ 中的余单纯形对象（留意
箭头反向）。现在令
$\overline{Q}_\bullet=Q_\bullet\otimes_{P_\bullet}B$。命题
\ref{cotangent-proposition-triangle} 证明的关键，是说明
$\overline{Q}_\bullet$ 是 $C$ 在 $B$ 上的预解。把《位点上的上同调》
引理 \ref{sites-cohomology-lemma-O-homology-qis} 应用于
$\mathcal{C}=\Delta$、$\mathcal{O}=P_\bullet$、
$\mathcal{O}'=B$ 与 $\mathcal{F}=Q_\bullet$，即可得到此结论（这里
用到 $Q_n$ 在 $P_n$ 上平坦；关于单纯形模与层的关系，见《位点上的
上同调》注 \ref{sites-cohomology-remark-simplicial-modules}）。这一关键
事实表明，命题 \ref{cotangent-proposition-triangle} 的可区别三角，就是与下列
单纯形 $C$-模短正合列相伴的可区别三角：
$$
0 \to
\Omega_{P_\bullet/A} \otimes_{P_\bullet} C \to
\Omega_{Q_\bullet/A} \otimes_{Q_\bullet} C \to
\Omega_{\overline{Q}_\bullet/B} \otimes_{\overline{Q}_\bullet} C \to 0
$$
该短正合列由短正合列
$0 \to \Omega_{P_n/A} \otimes_{P_n} Q_n \to \Omega_{Q_n/A} \to
\Omega_{Q_n/P_n} \to 0$ 得到，见《交换代数》引理
\ref{algebra-lemma-ses-formally-smooth}。具体地，由注
\ref{cotangent-remark-resolution} 及上述关键事实，右端复形在 $D(C)$ 中表示
$L_{C/B}$。

\medskip\noindent
若 $B\to C$ 不单射，则可用上述方法为
$A\to B\to B\times C$ 得到基本三角。由于
$L_{B \times C/B} \to L_{B/B} \oplus L_{C/B}$ 以及
$L_{B \times C/A} \to L_{B/A} \oplus L_{C/A}$
在 $D(B\times C)$ 中都是拟同构（引理
\ref{cotangent-lemma-cotangent-complex-product}），再与平坦环同态
$B\times C\to C$ 作张量积，便在 $D(C)$ 中诱导所需的可区别三角。
\end{remark}

\begin{remark}
\label{cotangent-remark-explicit-map}
设 $A\to B\to C$ 为环同态，且 $B\to C$ 单射。回忆注
\ref{cotangent-remark-triangle} 中的记号
$P_\bullet$、$Q_\bullet$、$\overline{Q}_\bullet$。令
$R_\bullet$ 为 $C$ 在 $B$ 上的标准预解。本注说明如何得到
$\Omega_{\overline{Q}_\bullet/B}
\otimes_{\overline{Q}_\bullet}C$ 与
$L_{C/B}=\Omega_{R_\bullet/B}\otimes_{R_\bullet}C$ 的典范等同。
令 $S_\bullet\to B$ 为 $B$ 在 $B$ 上的标准预解。注意，由于
$B\to C$ 单射，函子性同态 $S_\bullet\to R_\bullet$ 把 $R_n$
表为 $S_n$ 上的多项式代数。例如，在次数 $0$ 上有同态
$B[B]\to B[C]$，在次数 $1$ 上有同态
$B[B[B]]\to B[B[C]]$，依此类推。因此
$\overline{R}_\bullet=R_\bullet\otimes_{S_\bullet}B$ 也是 $B$
上的单纯形多项式代数。与注 \ref{cotangent-remark-triangle} 中相同，由
《位点上的上同调》引理
\ref{sites-cohomology-lemma-O-homology-qis} 可知
$\overline{R}_\bullet\to C$ 是预解。由于存在交换图
$$
\xymatrix{
Q_\bullet \ar[r] & R_\bullet \\
P_\bullet \ar[u] \ar[r] & S_\bullet \ar[u] \ar[r] & B
}
$$
得到典范同态
$\overline{Q}_\bullet = Q_\bullet \otimes_{P_\bullet} B \to
\overline{R}_\bullet$。所以，同态
$$
L_{C/B} = \Omega_{R_\bullet/B} \otimes_{R_\bullet} C
\longrightarrow
\Omega_{\overline{R}_\bullet/B} \otimes_{\overline{R}_\bullet} C
\longleftarrow
\Omega_{\overline{Q}_\bullet/B} \otimes_{\overline{Q}_\bullet} C
$$
都是拟同构（注 \ref{cotangent-remark-resolution}）；把其中一个与另一个的逆
复合，便得到所需等同。
\end{remark}





\section{局部化与 \'etale 环同态}
\label{cotangent-section-localization}

\noindent
本节研究环局部化后的情形。设 $A\to A'\to B$ 为环同态，满足
$B=B\otimes_A^\mathbf{L}A'$。例如，当 $A'=S^{-1}A$ 是 $A$ 关于
乘法子集 $S\subset A$ 的局部化时，这一条件成立。在此情形，对
$\mathcal{C}_{B/A'}$ 上的阿贝尔层 $\mathcal{F}'$，
$g^{-1}\mathcal{F}'$ 在 $\mathcal{C}_{B/A}$ 上的同调，与
$\mathcal{F}'$ 在 $\mathcal{C}_{B/A'}$ 上的同调一致；精确陈述见
引理 \ref{cotangent-lemma-flat-base-change}。

\begin{lemma}
\label{cotangent-lemma-localize-at-bottom}
设 $A\to A'\to B$ 为环同态，满足
$B=B\otimes_A^\mathbf{L}A'$。那么在 $D(B)$ 中有
$L_{B/A}=L_{B/A'}$。
\end{lemma}

\begin{proof}
根据上面的讨论（即应用引理 \ref{cotangent-lemma-flat-base-change}）以及引理
\ref{cotangent-lemma-compute-cotangent-complex}，须证明
$\mathcal{C}_{B/A}$ 上由规则
$(P\to B)\mapsto\Omega_{P/A}\otimes_PB$ 给出的层，是由规则
$(P\to B)\mapsto\Omega_{P/A'}\otimes_PB$ 给出的层的拉回。拉回函子
$g^{-1}$ 由预复合函子
$u:\mathcal{C}_{B/A}\to\mathcal{C}_{B/A'}$ 给出，其中
$(P\to B)\mapsto(P\otimes_AA'\to B)$。因此只须证明
$$
\Omega_{P/A} \otimes_P B =
\Omega_{P \otimes_A A'/A'} \otimes_{(P \otimes_A A')} B
$$
由《交换代数》引理 \ref{algebra-lemma-differentials-base-change}，右端等于
$$
(\Omega_{P/A} \otimes_A A') \otimes_{(P \otimes_A A')} B
$$
由于 $P$ 是 $A$ 上的多项式代数，模 $\Omega_{P/A}$ 自由，故此等式
显然成立。
\end{proof}

\begin{lemma}
\label{cotangent-lemma-derived-diagonal}
设 $A\to B$ 为环同态，满足 $B=B\otimes_A^\mathbf{L}B$。那么在
$D(B)$ 中有 $L_{B/A}=0$。
\end{lemma}

\begin{proof}
由引理 \ref{cotangent-lemma-localize-at-bottom} 和
\ref{cotangent-lemma-cotangent-complex-polynomial-algebra}，有
$L_{B/A}=L_{B/B}=0$。
\end{proof}

\begin{lemma}
\label{cotangent-lemma-bootstrap}
设 $A\to B$ 为环同态，且当 $i>0$ 时
$\text{Tor}^A_i(B,B)=0$，并且 $L_{B/B\otimes_AB}=0$。那么在
$D(B)$ 中有 $L_{B/A}=0$。
\end{lemma}

\begin{proof}
由引理 \ref{cotangent-lemma-flat-base-change-cotangent-complex}，有
$L_{B/A}\otimes_B^\mathbf{L}(B\otimes_AB)=L_{B\otimes_AB/B}$。
现在对环同态 $B\to B\otimes_AB\to B$ 使用可区别三角
(\ref{cotangent-equation-triangle})：
$$
L_{B \otimes_A B/B} \otimes^\mathbf{L}_{(B \otimes_A B)} B \to
L_{B/B} \to L_{B/B \otimes_A B} \to
L_{B \otimes_A B/B} \otimes^\mathbf{L}_{(B \otimes_A B)} B[1]
$$
再利用 $L_{B/B}$ 的消失（引理
\ref{cotangent-lemma-cotangent-complex-polynomial-algebra}）以及假设中的
$L_{B/B\otimes_AB}$ 的消失，可得
$$
0 =
L_{B \otimes_A B/B} \otimes^\mathbf{L}_{(B \otimes_A B)} B =
L_{B/A} \otimes_B^\mathbf{L} (B \otimes_A B)
\otimes^\mathbf{L}_{(B \otimes_A B)} B = L_{B/A}
$$
正合所需。
\end{proof}

\begin{lemma}
\label{cotangent-lemma-when-zero}
在下列每种情形中，余切复形 $L_{B/A}$ 均为零：
\begin{enumerate}
\item $A\to B$ 与 $B\otimes_AB\to B$ 均平坦，即 $A\to B$ 弱
\'etale（《更多交换代数》定义
\ref{more-algebra-definition-weakly-etale}）；
\item $A\to B$ 是环的平坦满态射；
\item 对某个乘法子集 $S\subset A$，有 $B=S^{-1}A$；
\item $A\to B$ 非分歧且平坦；
\item $A\to B$ 是 \'etale 同态；
\item $A\to B$ 是一族余切复形消失的环同态的滤余极限；
\item $B$ 是 $A$ 的某个局部环的 Hensel 化；
\item $B$ 是 $A$ 的某个局部环的严格 Hensel 化；以及
\item 在此增加更多情形。
% P11-E0027: frozen source contains the editorial placeholder “add more here”; retained here. See control/ERRATA_CANDIDATES.jsonl.
\end{enumerate}
\end{lemma}

\begin{proof}
在情形 (1) 中，把引理 \ref{cotangent-lemma-derived-diagonal} 应用于满射平坦
环同态 $B\otimes_AB\to B$，得到 $L_{B/B\otimes_AB}=0$；再应用
引理 \ref{cotangent-lemma-bootstrap} 即得结论。情形 (2)--(5) 都是 (1) 的
特例。(6) 来自引理 \ref{cotangent-lemma-colimit-cotangent-complex}。(7) 与
(8) 来自如下事实：严格或非严格 Hensel 化是 $A$ 的 \'etale 环扩张
的滤余极限；见《交换代数》引理
\ref{algebra-lemma-henselization-different} 和
\ref{algebra-lemma-strict-henselization-different}。
\end{proof}

\begin{lemma}
\label{cotangent-lemma-localize-on-top}
设 $A\to B\to C$ 为环同态，且 $L_{C/B}=0$。那么
$L_{C/A}=L_{B/A}\otimes_B^\mathbf{L}C$。
\end{lemma}

\begin{proof}
这是可区别三角 (\ref{cotangent-equation-triangle}) 的直接推论。
\end{proof}

\begin{lemma}
\label{cotangent-lemma-localize}
设 $A\to B$ 为环同态，$S\subset A$ 与 $T\subset B$ 为乘法子集，
且 $S$ 的像包含在 $T$ 中。那么在 $D(T^{-1}B)$ 中有
$L_{T^{-1}B/S^{-1}A}=L_{B/A}\otimes_BT^{-1}B$。
\end{lemma}

\begin{proof}
引理 \ref{cotangent-lemma-localize-on-top} 给出
$L_{T^{-1}B/A}=L_{B/A}\otimes_BT^{-1}B$，引理
\ref{cotangent-lemma-localize-at-bottom} 给出
$L_{T^{-1}B/A}=L_{T^{-1}B/S^{-1}A}$。
\end{proof}

\begin{lemma}
\label{cotangent-lemma-cotangent-complex-henselization}
设 $A\to B$ 是局部环之间的局部同态。令 $A^h\to B^h$ 以及
$A^{sh}\to B^{sh}$ 分别为诱导的 Hensel 化同态与严格 Hensel 化同态。
那么
$$
L_{B^h/A^h} = L_{B^h/A} = L_{B/A} \otimes_B^\mathbf{L} B^h
\quad\text{分别}\quad
L_{B^{sh}/A^{sh}} = L_{B^{sh}/A} = L_{B/A} \otimes_B^\mathbf{L} B^{sh}
$$
分别在 $D(B^h)$ 与 $D(B^{sh})$ 中成立。
\end{lemma}

\begin{proof}
由引理 \ref{cotangent-lemma-when-zero}，复形 $L_{A^h/A}$、$L_{A^{sh}/A}$、
$L_{B^h/B}$ 与 $L_{B^{sh}/B}$ 均为零。对
$A\to B\to B^h$ 使用基本可区别三角
(\ref{cotangent-equation-triangle})，得到
$L_{B^h/A}=L_{B/A}\otimes_B^\mathbf{L}B^h$。再对
$A\to A^h\to B^h$ 使用基本三角，得到
$L_{B^h/A^h}=L_{B^h/A}$。严格 Hensel 化的情形同理。
\end{proof}




\section{光滑环同态}
\label{cotangent-section-smooth}

\noindent
设 $C\to B$ 为核是 $I$ 的环满射。若 $I/I^2$ 是平坦 $B$-模，且
$\text{Tor}_*^C(B,B)$ 是 $I/I^2$ 上的外代数，则称这样的环同态
“弱拟正则”。若要把引理 \ref{cotangent-lemma-when-zero} 对“\'etale 环同态”
所作的事情推广到“光滑环同态”，应考察满足乘法同态
$B\otimes_AB\to B$ 弱拟正则的平坦环同态 $A\to B$。目前我们只处理
光滑环同态。

\begin{lemma}
\label{cotangent-lemma-when-projective}
若 $A\to B$ 是光滑环同态，则 $L_{B/A}=\Omega_{B/A}[0]$。
\end{lemma}

\begin{proof}
由引理 \ref{cotangent-lemma-identify-H0}，次数 $0$ 的上同调已经一致。因此只需
证明其余上同调群均为零。由引理 \ref{cotangent-lemma-localize-on-top}，对
$g\in B$ 有 $L_{B_g/A}=(L_{B/A})_g$，故只需在 $\Spec(B)$ 上局部地
证明。于是可假设 $A\to B$ 标准光滑（《交换代数》引理
\ref{algebra-lemma-smooth-syntomic}），即可以把 $A\to B$ 分解为
$A\to A[x_1,\ldots,x_n]\to B$，其中
$A[x_1,\ldots,x_n]\to B$ 是 \'etale 同态。在此情形，引理
\ref{cotangent-lemma-when-zero} 与 \ref{cotangent-lemma-localize-on-top} 给出
$L_{B/A}=L_{A[x_1,\ldots,x_n]/A}\otimes B$，再由引理
\ref{cotangent-lemma-cotangent-complex-polynomial-algebra} 得出结论。
\end{proof}





\section{正特征}
\label{cotangent-section-positive-characteristic}

\noindent
本节固定一个素数 $p$。若环 $A$ 中 $p=0$，则以 $F_A:A\to A$
记 Frobenius 自同态 $a\mapsto a^p$。

\begin{lemma}
\label{cotangent-lemma-frobenius-homotopy}
设 $A\to B$ 为环同态，且在 $A$ 中 $p=0$。令 $P_\bullet$ 为 $B$
在 $A$ 上的标准预解。由图
$$
\xymatrix{
B \ar[r]_{F_B} & B \\
A \ar[u] \ar[r]^{F_A} & A \ar[u]
}
$$
按第 \ref{cotangent-section-functoriality} 节的方式诱导的同态
$P_\bullet\to P_\bullet$，与在每个 $P_n$ 上由 Frobenius 给出的自同态
$P_\bullet\to P_\bullet$ 同伦。
\end{lemma}

\begin{proof}
令 $\mathcal{A}$ 为 $\mathbf{F}_p$-代数同态 $A\to B$ 所成的范畴。
令 $\mathcal{S}$ 为二元组 $(A,E)$ 所成的范畴，其中 $A$ 是
$\mathbf{F}_p$-代数，$E$ 是集合。考虑伴随函子
$$
V : \mathcal{A} \to \mathcal{S}, \quad (A \to B) \mapsto (A, B)
$$
以及
$$
U : \mathcal{S} \to \mathcal{A}, \quad (A, E) \mapsto (A \to A[E])
$$
令 $X$ 为《单纯形》第 \ref{simplicial-section-standard} 节所构造的、
从 $\mathcal{A}$ 到 $\mathcal{A}$ 的函子范畴中的单纯形对象。
显然 $P_\bullet=X(A\to B)$，因为若固定 $A$，则……。
% P11-E0028: frozen source contains duplicated “in” and an incomplete sentence ending “because if we fix A then.”; retained as an explicit ellipsis. See control/ERRATA_CANDIDATES.jsonl.

\medskip\noindent
令 $Y=U\circ V$。回忆 $X$ 由 $Y$ 及若干同态构造，其各项
$X_n=Y\circ\ldots\circ Y$ 含 $n+1$ 个因子；该构造见《单纯形》例
\ref{simplicial-example-godement}，细节见《单纯形》引理
\ref{simplicial-lemma-standard-simplicial} 的证明。

\medskip\noindent
令 $f:\text{id}_\mathcal{A}\to\text{id}_\mathcal{A}$ 为恒等函子的
Frobenius 自同态。换言之，令
$f_{A\to B}=(F_A,F_B):(A\to B)\to(A\to B)$。于是
$X(A\to B)$ 上的两个同态分别由自然变换
$f\star1_X$ 与 $1_X\star f$ 给出。细节略去。结论遂由《单纯形》引理
\ref{simplicial-lemma-godement-before-after} 得出。
\end{proof}

\begin{lemma}
\label{cotangent-lemma-frobenius-acts-as-zero}
设 $p$ 为素数，$A\to B$ 为环同态，并假设在 $A$ 中 $p=0$。由
Frobenius 同态 $F_A$ 与 $F_B$ 按第 \ref{cotangent-section-functoriality} 节
诱导的同态 $L_{B/A}\to L_{B/A}$ 零同伦。
\end{lemma}

\begin{proof}
令 $P_\bullet$ 为 $B$ 在 $A$ 上的标准预解。由引理
\ref{cotangent-lemma-frobenius-homotopy}，$F_A$ 与 $F_B$ 诱导的同态
$P_\bullet\to P_\bullet$，与在每一项上由 Frobenius 给出的同态
$F_{P_\bullet}:P_\bullet\to P_\bullet$ 同伦。显然
$F_{P_\bullet}$ 在 $\Omega_{P_n/A}\to\Omega_{P_n/A}$ 上诱导零同态
（因为 $p$ 次幂的导数为零），故得到所需结论。
\end{proof}

\begin{lemma}
\label{cotangent-lemma-perfect-zero}
设 $p$ 为素数，$A\to B$ 为环同态，并假设在 $A$ 中 $p=0$。若
$A$ 与 $B$ 都是完美环，则 $L_{B/A}$ 在 $D(B)$ 中为零。
\end{lemma}

\begin{proof}
同态 $(F_A,F_B):(A\to B)\to(A\to B)$ 是同构，故在 $L_{B/A}$
上诱导同构；另一方面，由引理 \ref{cotangent-lemma-frobenius-acts-as-zero}，
它在 $L_{B/A}$ 上诱导零同态。
\end{proof}





\section{与朴素余切复形的比较}
\label{cotangent-section-surjections}

\noindent
朴素余切复形已在《代数》第 \ref{algebra-section-netherlander} 节中引入。

\begin{remark}
\label{cotangent-remark-make-map}
设 $A \to B$ 为环同态。依照第 \ref{cotangent-section-compute-L-pi-shriek} 节，
在 $\mathcal{C}_{B/A}$ 上工作，并令
$\mathcal{J} \subset \mathcal{O}$ 为
$\mathcal{O} \to \underline{B}$ 的核。由引理
\ref{cotangent-lemma-apply-O-B-comparison}，注意到 $L\pi_!(\mathcal{J}) = 0$。令
$\Omega =  \Omega_{\mathcal{O}/A} \otimes_\mathcal{O} \underline{B}$
，则由引理 \ref{cotangent-lemma-compute-cotangent-complex}，
$L_{B/A} = L\pi_!(\Omega)$。于是
$L\pi_!(\mathcal{J} \to \Omega) = L\pi_!(\Omega) = L_{B/A}$。
因此，对 $\mathcal{C}_{B/A}$ 的任一对象 $U = (P \to B)$，我们得到态射
\begin{equation}
\label{cotangent-equation-comparison-map-A}
(J \to \Omega_{P/A} \otimes_P B) \longrightarrow L_{B/A}
\end{equation}
其中 $J = \Ker(P \to B)$；这是 $D(A)$ 中的态射，参见
《站点上的上同调》评注
\ref{sites-cohomology-remark-map-evaluation-to-derived}。
继续以此方式进行。由引理
\ref{cotangent-lemma-O-homology-B-homology}，注意到
$L\pi_!(\mathcal{J} \otimes_\mathcal{O}^\mathbf{L} \underline{B}) =
L\pi_!(\mathcal{J}) = 0$。由于 $\text{Tor}_0^\mathcal{O}(\mathcal{J}, \underline{B}) =
\mathcal{J}/\mathcal{J}^2$
，谱序列
$$
H_p(\mathcal{C}_{B/A}, \text{Tor}_q^\mathcal{O}(\mathcal{J}, \underline{B}))
\Rightarrow 
H_{p + q}(\mathcal{C}_{B/A},
\mathcal{J} \otimes_\mathcal{O}^\mathbf{L} \underline{B}) = 0
$$
（《导出范畴》引理 \ref{derived-lemma-two-ss-complex-functor} 的对偶）
蕴含
$H_0(\mathcal{C}_{B/A}, \mathcal{J}/\mathcal{J}^2) = 0$
且 $H_1(\mathcal{C}_{B/A}, \mathcal{J}/\mathcal{J}^2) = 0$。
由此，$\underline{B}$-模复形
$\mathcal{J}/\mathcal{J}^2 \to \Omega$ 满足
$\tau_{\geq -1}L\pi_!(\mathcal{J}/\mathcal{J}^2 \to \Omega) =
\tau_{\geq -1}L_{B/A}$。
因此，对 $\mathcal{C}_{B/A}$ 的任一对象 $U = (P \to B)$，我们得到态射
\begin{equation}
\label{cotangent-equation-comparison-map}
(J/J^2 \to \Omega_{P/A} \otimes_P B) \longrightarrow \tau_{\geq -1}L_{B/A}
\end{equation}
；这是 $D(B)$ 中的态射，参见《站点上的上同调》评注
\ref{sites-cohomology-remark-map-evaluation-to-derived}。
\end{remark}

\noindent
第一种情形是环的满射。

\begin{lemma}
\label{cotangent-lemma-surjection}
\begin{slogan}
满环同态的余切复形在零次的上同调为零；其 $-1$ 次上同调是核模去其平方。
\end{slogan}
设 $A \to B$ 为核是 $I$ 的满环同态。则
$H^0(L_{B/A}) = 0$ 且 $H^{-1}(L_{B/A}) = I/I^2$。
此同构来自 $\mathcal{C}_{B/A}$ 的对象 $(A \to B)$ 所对应的态射
(\ref{cotangent-equation-comparison-map})。
\end{lemma}

\begin{proof}
我们将在下文利用 $A \to B$ 的满射性证明：在 $\mathcal{C}_{B/A}$ 上
存在短正合层列
$$
0 \to \pi^{-1}(I/I^2) \to \mathcal{J}/\mathcal{J}^2 \to \Omega \to 0
$$
。取 $L\pi_!$ 及其相应的同调长正合列，并使用评注
\ref{cotangent-remark-make-map} 所示的
$H_1(\mathcal{C}_{B/A}, \mathcal{J}/\mathcal{J}^2)$ 与
$H_0(\mathcal{C}_{B/A}, \mathcal{J}/\mathcal{J}^2)$ 的消失，再由引理
\ref{cotangent-lemma-pi-lower-shriek-constant-sheaf} 即得所需结论。

\medskip\noindent
剩下的是验证上面提到的局部陈述。对 $\mathcal{C}_{B/A}$ 的每个对象
$U = (P \to B)$，我们可以选取同构 $P = A[E]$，使得同态
$P \to B$ 将每个 $e \in E$ 映到零。于是
$J = \mathcal{J}(U) \subset P = \mathcal{O}(U)$
等于 $J = IP + (e; e \in E)$。上述短层列在 $U$ 上的值为
$$
0 \to I/I^2 \to J/J^2 \to \Omega_{P/A} \otimes_P B \to 0
$$
验证从略（提示：唯一棘手之处是 $IP \cap J^2 = IJ$；例如可由
《代数续篇》引理 \ref{more-algebra-lemma-conormal-sequence-H1-regular} 推出）。
\end{proof}

\begin{lemma}
\label{cotangent-lemma-relation-with-naive-cotangent-complex}
设 $A \to B$ 为环同态。则 $\tau_{\geq -1}L_{B/A}$
与朴素余切复形典范拟同构。
\end{lemma}

\begin{proof}
考虑核为 $I$ 的 $P = A[B] \to B$。$B$ 在 $A$ 上的朴素余切复形
$\NL_{B/A}$ 是复形 $I/I^2 \to \Omega_{P/A} \otimes_P B$，参见
《代数》定义 \ref{algebra-definition-naive-cotangent-complex}。
注意，我们已在 (\ref{cotangent-equation-comparison-map}) 中构造了典范态射
$$
c : \NL_{B/A} \longrightarrow \tau_{\geq -1}L_{B/A}
$$
考虑与环同态 $A \to A[B] \to B$ 相关的可区别三角
(\ref{cotangent-equation-triangle})
$$
L_{P/A} \otimes_P^\mathbf{L} B \to L_{B/A} \to L_{B/P} \to 
(L_{P/A} \otimes_P^\mathbf{L} B)[1]
$$
。我们知道
$L_{P/A} = \Omega_{P/A}[0] = \NL_{P/A}$ 在 $D(P)$ 中成立
（引理 \ref{cotangent-lemma-cotangent-complex-polynomial-algebra} 与
《代数》引理 \ref{algebra-lemma-NL-polynomial-algebra}），并且
$\tau_{\geq -1}L_{B/P} = I/I^2[1] = \NL_{B/P}$ 在 $D(B)$ 中成立
（引理 \ref{cotangent-lemma-surjection} 与
《代数》引理 \ref{algebra-lemma-NL-surjection}）。由
《代数》引理 \ref{algebra-lemma-exact-sequence-NL}
及与该可区别三角相关的上同调长正合列，要证明 $c$ 是拟同构，
只需证明态射 $L_{P/A} \to L_{B/A} \to L_{B/P}$ 在上同调群上
与朴素余切复形的相应态射
$\NL_{P/A} \to \NL_{B/A} \to \NL_{B/P}$
相容。验证从略。
\end{proof}

\begin{remark}
\label{cotangent-remark-explicit-comparison-map}
我们可以如下明确写出引理
\ref{cotangent-lemma-relation-with-naive-cotangent-complex} 的比较态射。
令 $P_\bullet$ 为 $B$ 在 $A$ 上的标准预解，并令
$I = \Ker(A[B] \to B)$。回忆 $P_0 = A[B]$。该引理中的态射
由交换图
$$
\xymatrix{
L_{B/A} \ar[d] & \ldots \ar[r] &
\Omega_{P_2/A} \otimes_{P_2} B
\ar[r] \ar[d] &
\Omega_{P_1/A} \otimes_{P_1} B
\ar[r] \ar[d] &
\Omega_{P_0/A} \otimes_{P_0} B
\ar[d] \\
\NL_{B/A} & \ldots \ar[r] &
0 \ar[r] & 
I/I^2 \ar[r] &
\Omega_{P_0/A} \otimes_{P_0} B
}
$$
给出。我们通过将 $\text{d}f \otimes b$ 映到
$(d_0(f) - d_1(f))b$ 在 $I/I^2$ 中的类，构造以 $I/I^2$ 为靶的向下箭头。
这里 $d_i : P_1 \to P_0$（$i = 0, 1$）是单纯形结构的两个面态射。
这是有意义的，因为 $d_0 - d_1$ 将 $P_1$ 映入
$I = \Ker(P_0 \to B)$。此规则良定义的验证从略。
我们的态射与微分
$\Omega_{P_1/A} \otimes_{P_1} B \to \Omega_{P_0/A} \otimes_{P_0} B$
相容，因为该微分将 $\text{d}f \otimes b$ 映到
$\text{d}(d_0(f) - d_1(f)) \otimes b$。此外，微分
$\Omega_{P_2/A} \otimes_{P_2} B \to \Omega_{P_1/A} \otimes_{P_1} B$
将 $\text{d}f \otimes b$ 映到
$\text{d}(d_0(f) - d_1(f) + d_2(f)) \otimes b$，后者被我们的向下箭头
消去。因此得到了复形态射。它与引理
\ref{cotangent-lemma-relation-with-naive-cotangent-complex} 中的态射相同这一验证从略。
\end{remark}

\begin{remark}
\label{cotangent-remark-surjection}
采用评注 \ref{cotangent-remark-make-map} 的记号。那里给出的论证表明，谱序列的微分
$$
H_2(\mathcal{C}_{B/A}, \mathcal{J}/\mathcal{J}^2)
\longrightarrow
H_0(\mathcal{C}_{B/A}, \text{Tor}_1^\mathcal{O}(\mathcal{J}, \underline{B}))
$$
是同构。令 $\mathcal{C}'_{B/A}$ 表示 $\mathcal{C}_{B/A}$ 中由满射
$P \to B$ 构成的全子范畴。余切复形与朴素余切复形相一致（引理
\ref{cotangent-lemma-relation-with-naive-cotangent-complex}），这表明我们在
$\mathcal{C}'_{B/A}$ 上有正合层列
$$
0 \to \underline{H_1(L_{B/A})} \to
\mathcal{J}/\mathcal{J}^2 \xrightarrow{\text{d}} \Omega \to
\underline{H_2(L_{B/A})} \to 0
$$
。由此，整个范畴 $\mathcal{C}_{B/A}$ 上的 $\Ker(d)$ 与 $\Coker(d)$
具有消失的高阶同调群，因为由引理 \ref{cotangent-lemma-identify-pi-shriek}，
这些群由常值单纯形阿贝尔群的同调群计算。因此我们推出
$$
H_n(\mathcal{C}_{B/A}, \mathcal{J}/\mathcal{J}^2) \to H_n(L_{B/A})
$$
对一切 $n \geq 2$ 都是同构。结合上面的评注，我们得到公式
$H_2(L_{B/A}) =
H_0(\mathcal{C}_{B/A}, \text{Tor}_1^\mathcal{O}(\mathcal{J}, \underline{B}))$。
\end{remark}





\section{Quillen 谱序列}
\label{cotangent-section-spectral-sequence}

\noindent
本节讨论一个联系导出张量积与余切复形的谱序列。

\begin{lemma}
\label{cotangent-lemma-vanishing-symmetric-powers}
记号和假设同《站点上的上同调》例
\ref{sites-cohomology-example-category-to-point}。假设 $\mathcal{C}$
具有《站点上的上同调》引理
\ref{sites-cohomology-lemma-compute-by-cosimplicial-resolution}
中的余单纯形对象。设 $\mathcal{F}$ 为满足
$H_0(\mathcal{C}, \mathcal{F}) = 0$ 的平坦 $\underline{B}$-模。
则当 $l < k$ 时，
$H_l(\mathcal{C}, \text{Sym}_{\underline{B}}^k(\mathcal{F})) = 0$。
\end{lemma}

\begin{proof}
在张量积、外幂和对称幂中省略下标 ${}_{\underline{B}}$。
我们对 $k$ 归纳证明该引理。$k = 0, 1$ 的情形由假设即得。
若 $k > 1$，考虑正合复形
$$
\ldots \to
\wedge^2\mathcal{F} \otimes \text{Sym}^{k - 2}\mathcal{F} \to
\mathcal{F} \otimes \text{Sym}^{k - 1}\mathcal{F} \to
\text{Sym}^k\mathcal{F} \to 0
$$
，其微分如 Koszul 复形中所定义。若将其视为
$\text{Sym}^k\mathcal{F}$ 的预解，便得到第一象限谱序列
$$
E_1^{p, q} =
H_p(\mathcal{C},
\wedge^{q + 1}\mathcal{F} \otimes \text{Sym}^{k - q - 1}\mathcal{F})
\Rightarrow
H_{p + q}(\mathcal{C}, \text{Sym}^k(\mathcal{F}))
$$
由《站点上的上同调》引理 \ref{sites-cohomology-lemma-eilenberg-zilber}，有
$$
L\pi_!(\wedge^{q + 1}\mathcal{F} \otimes \text{Sym}^{k - q - 1}\mathcal{F}) =
L\pi_!(\wedge^{q + 1}\mathcal{F}) \otimes_B^\mathbf{L}
L\pi_!(\text{Sym}^{k - q - 1}\mathcal{F}))
$$
% P11-E0029：冻结源在上一公式末尾保留一个疑似多余的右圆括号。
由导出张量积的构造可知，将归纳假设与
$H_0(\mathcal{C}, \wedge^{q + 1}(\mathcal{F})) = 0$ 的消失结合起来，
便可证明所需结论。该消失成立，是因为
$\wedge^{q + 1}(\mathcal{F})$ 是 $\mathcal{F}^{\otimes q + 1}$ 的商，且
$H_0(\mathcal{C}, \mathcal{F}^{\otimes q + 1})$ 是
$H_0(\mathcal{C}, \mathcal{F})^{\otimes q + 1}$ 的商，而后者为零。
\end{proof}

\begin{remark}
\label{cotangent-remark-first-homology-symmetric-power}
在引理 \ref{cotangent-lemma-vanishing-symmetric-powers} 的情形下，可以证明
$H_k(\mathcal{C}, \text{Sym}^k(\mathcal{F})) =
\wedge^k_B(H_1(\mathcal{C}, \mathcal{F}))$。事实上，从证明可以推出，
$H_k(\mathcal{C}, \text{Sym}^k(\mathcal{F}))$ 是下式的 $S_k$-余不变量：
$$
H^{-k}(L\pi_!(\mathcal{F}) \otimes_B^\mathbf{L}
L\pi_!(\mathcal{F}) \otimes_B^\mathbf{L}
\ldots \otimes_B^\mathbf{L} L\pi_!(\mathcal{F})) =
H_1(\mathcal{C}, \mathcal{F})^{\otimes k}
$$
因此，我们的断言是：这一作用由 $S_k$ 在张量积上的通常作用乘以符号特征给出。
要证明这一点，必须逐一核查与多重复形相关的全复形定义中的符号约定。
验证从略。
\end{remark}

\begin{lemma}
\label{cotangent-lemma-map-tors-zero}
设 $A$ 为环，$P = A[E]$ 为多项式环。令
$I = (e; e \in E) \subset P$。则态射
$\text{Tor}_i^P(A, I^{n + 1}) \to \text{Tor}_i^P(A, I^n)$
对一切 $i$ 和 $n$ 都为零。
\end{lemma}

\begin{proof}
以 $x_e \in P$ 表示对应于 $e \in E$ 的变量。$A$ 在 $P$ 上的一个自由预解
由 $x_e$ 上的 Koszul 复形 $K_\bullet$ 给出。这里 $K_i$ 的基由外积
$e_1 \wedge \ldots \wedge e_i$（$e_1, \ldots, e_i \in E$）给出，且
$d(e) = x_e$。因此 $K_\bullet \otimes_P I^n = I^nK_\bullet$ 计算
$\text{Tor}_i^P(A, I^n)$。注意一切均分次，其中
$\deg(x_e) = 1$、$\deg(e) = 1$，而对 $a \in A$ 有 $\deg(a) = 0$。
设 $\xi \in I^{n + 1}K_i$ 是次数为 $m$ 的齐次余圈。注意
$m \geq i + 1 + n$。由于 $K_\bullet$ 在次数 $>0$ 正合，存在
$\eta \in K_{i + 1}$ 使 $\xi = \text{d}\eta$（$i = 0$ 的情形留给读者）。
现在 $\deg(\eta) = m \geq i + 1 + n$。因而用上述基展开 $\eta$，
可见各坐标都属于 $I^n$。所以 $\xi$ 在 $I^nK_\bullet$ 的同调中映为零，
正如所需。
\end{proof}

\begin{theorem}[Quillen 谱序列]
\label{cotangent-theorem-quillen-spectral-sequence}
设 $A \to B$ 为满环同态。考虑 $\mathcal{C}_{B/A}$ 上的
$\underline{B}$-模层
$\Omega = \Omega_{\mathcal{O}/A} \otimes_\mathcal{O} \underline{B}$，
参见第 \ref{cotangent-section-compute-L-pi-shriek} 节。则存在 $E_1$ 页为
$$
E_1^{p, q} =
H_{- p - q}(\mathcal{C}_{B/A}, \text{Sym}^p_{\underline{B}}(\Omega))
\Rightarrow \text{Tor}^A_{- p - q}(B, B)
$$
的谱序列，其中 $d_r$ 的双次数为 $(r, -r + 1)$。此外，当 $i < k$ 时，
$H_i(\mathcal{C}_{B/A}, \text{Sym}^k_{\underline{B}}(\Omega)) = 0$。
\end{theorem}

\begin{proof}
令 $I \subset A$ 为 $A \to B$ 的核，并令
$\mathcal{J} \subset \mathcal{O}$ 为
$\mathcal{O} \to \underline{B}$ 的核。于是
$I\mathcal{O} \subset \mathcal{J}$。令
$\mathcal{K} = \mathcal{J}/I\mathcal{O}$ 且
$\overline{\mathcal{O}} = \mathcal{O}/I\mathcal{O}$。

\medskip\noindent
对 $\mathcal{C}_{B/A}$ 的每个对象 $U = (P \to B)$，我们可以选取同构
$P = A[E]$，使同态 $P \to B$ 将每个 $e \in E$ 映到零。于是
$J = \mathcal{J}(U) \subset P = \mathcal{O}(U)$
等于 $J = IP + (e; e \in E)$。此外，
$\overline{\mathcal{O}}(U) = B[E]$ 且 $K = \mathcal{K}(U) = (e; e \in E)$
，而且 $K$ 是多项式环 $B[E]$ 中由这些变量生成的理想。特别地，显然
$$
K/K^2 \xrightarrow{\text{d}} \Omega_{P/A} \otimes_P B
$$
是双射。换言之，
$\Omega = \mathcal{K}/\mathcal{K}^2$ 且
$\text{Sym}_B^k(\Omega) = \mathcal{K}^k/\mathcal{K}^{k + 1}$。
注意到 $\pi_!(\Omega) = \Omega_{B/A} = 0$（引理
\ref{cotangent-lemma-identify-H0}），因为 $A \to B$ 是满射（《代数》引理
\ref{algebra-lemma-trivial-differential-surjective}）。
由引理 \ref{cotangent-lemma-vanishing-symmetric-powers}，我们推出
$$
H_i(\mathcal{C}_{B/A}, \mathcal{K}^k/\mathcal{K}^{k + 1}) =
H_i(\mathcal{C}_{B/A}, \text{Sym}^k_{\underline{B}}(\Omega)) = 0
$$
对 $i < k$ 成立。这证明了定理的最后一个陈述。

\medskip\noindent
证明该定理的方法是注意到
$$
B \otimes_A^\mathbf{L} B = L\pi_!(\mathcal{O}) \otimes_A^\mathbf{L} B =
L\pi_!(\mathcal{O} \otimes_{\underline{A}}^\mathbf{L} \underline{B}) =
L\pi_!(\overline{\mathcal{O}})
$$
第一个等号由引理 \ref{cotangent-lemma-apply-O-B-comparison} 得到，第二个等号由
《站点上的上同调》引理 \ref{sites-cohomology-lemma-change-of-rings} 得到，
第三个等号则因为 $\mathcal{O}$ 在 $\underline{A}$ 上平坦。层
$\overline{\mathcal{O}}$ 有滤过
$$
\ldots \subset
\mathcal{K}^3 \subset
\mathcal{K}^2 \subset
\mathcal{K} \subset
\overline{\mathcal{O}}
$$
。这在表示 $L\pi_!(\overline{\mathcal{O}})$ 的复形 $C$ 上诱导滤过
$F$，其中 $F^pC$ 表示 $L\pi_!(\mathcal{K}^p)$（$C$ 与 $F$ 的构造从略）。
考虑《同调》第 \ref{homology-section-filtered-complex} 节中与
$(C, F)$ 相关的谱序列。其 $E_1$ 页为
$$
E_1^{p, q} = H_{- p - q}(\mathcal{C}_{B/A}, \mathcal{K}^p/\mathcal{K}^{p + 1})
\quad\Rightarrow\quad
H_{- p - q}(\mathcal{C}_{B/A}, \overline{\mathcal{O}}) = 
\text{Tor}_{- p - q}^A(B, B)
$$
，微分为 $E_r^{p, q} \to E_r^{p + r, q - r + 1}$。为了证明收敛性，
我们将证明：对每个 $k$，存在 $c$ 使得
$H_i(\mathcal{C}_{B/A}, \mathcal{K}^n) = 0$
对 $i < k$ 与 $n > c$ 成立\footnote{事后可以推出“正确的”消失：
当 $i < n$ 时，$H_i(\mathcal{C}_{B/A}, \mathcal{K}^n) = 0$。}。

\medskip\noindent
给定 $k \geq 0$，令 $c = k^2$。我们断言
$$
H_i(\mathcal{C}_{B/A}, \mathcal{K}^{n + c}) \to
H_i(\mathcal{C}_{B/A}, \mathcal{K}^n)
$$
对 $i < k$ 及一切 $n \geq 0$ 都是零。注意
$\mathcal{K}^n/\mathcal{K}^{n + c}$ 有有限滤过，其相继商为
$\mathcal{K}^m/\mathcal{K}^{m + 1}$（$n \leq m < n + c$），且当
$i < n$ 时，
$H_i(\mathcal{C}_{B/A}, \mathcal{K}^m/\mathcal{K}^{m + 1}) = 0$
（见上文）。因此，该断言蕴含：对 $i < k$ 及一切 $n \geq k$，
$H_i(\mathcal{C}_{B/A}, \mathcal{K}^{n + c}) = 0$，这正是我们需要证明的。

\medskip\noindent
证明该断言。回忆，对任一 $\mathcal{O}$-模 $\mathcal{F}$，态射
$\mathcal{F} \to \mathcal{F} \otimes_\mathcal{O}^\mathbf{L} B$
在施用 $L\pi_!$ 后诱导同构，参见引理
\ref{cotangent-lemma-O-homology-B-homology}。考虑态射
$$
\mathcal{K}^{n + k} \otimes_\mathcal{O}^\mathbf{L} B
\longrightarrow
\mathcal{K}^n \otimes_\mathcal{O}^\mathbf{L} B
$$
我们断言，此态射在次数 $0, -1, \ldots, - k + 1$ 的上同调层上诱导零态射。
若这一第二断言成立，则 $k$ 重复合
$$
\mathcal{K}^{n + c} \otimes_\mathcal{O}^\mathbf{L} B
\longrightarrow
\mathcal{K}^n \otimes_\mathcal{O}^\mathbf{L} B
$$
分解经过
$\tau_{\leq -k}\mathcal{K}^n \otimes_\mathcal{O}^\mathbf{L} B$，从而当
$i < k$ 时在 $H_i(\mathcal{C}_{B/A}, -) = L_i\pi_!( - )$ 上诱导零，
参见《导出范畴》引理
\ref{derived-lemma-trick-vanishing-composition}。由上面的评注，这意味着
$H_i(\mathcal{C}_{B/A}, \mathcal{K}^{n + c}) \to
H_i(\mathcal{C}_{B/A}, \mathcal{K}^n)$
同样为零，因而证明了第一个断言。

\medskip\noindent
证明第二个断言。该陈述是局部的，故可在上述对象
$U = (P \to B)$ 上工作。我们必须证明态射
$$
\text{Tor}_i^P(B, K^{n + k}) \to \text{Tor}_i^P(B, K^n)
$$
当 $i < k$ 时为零。存在谱序列
$$
\text{Tor}_a^P(P/IP, \text{Tor}_b^{P/IP}(B, K^n))
\Rightarrow
\text{Tor}_{a + b}^P(B, K^n),
$$
，参见《代数续篇》例 \ref{more-algebra-example-tor-change-rings}。
因此，只需证明态射
$$
\text{Tor}_i^{P/IP}(B, K^{n + 1}) \to \text{Tor}_i^{P/IP}(B, K^n)
$$
对一切 $i$ 都为零。这正是引理 \ref{cotangent-lemma-map-tors-zero}。
\end{proof}

\begin{remark}
\label{cotangent-remark-elucidate-ss}
在定理 \ref{cotangent-theorem-quillen-spectral-sequence} 的情形下，令
$I = \Ker(A \to B)$。则
$H^{-1}(L_{B/A}) = H_1(\mathcal{C}_{B/A}, \Omega) = I/I^2$，参见
引理 \ref{cotangent-lemma-surjection}。因此由评注
\ref{cotangent-remark-first-homology-symmetric-power}，
$H_k(\mathcal{C}_{B/A}, \text{Sym}^k(\Omega)) = \wedge^k_B(I/I^2)$。
于是 $E_1$ 页形如
$$
\begin{matrix}
B \\
0 \\
0 & I/I^2 \\
0 & H^{-2}(L_{B/A}) \\
0 & H^{-3}(L_{B/A}) & \wedge^2(I/I^2) \\
0 & H^{-4}(L_{B/A}) & H_3(\mathcal{C}_{B/A}, \text{Sym}^2(\Omega)) \\
0 & H^{-5}(L_{B/A}) & H_4(\mathcal{C}_{B/A}, \text{Sym}^2(\Omega)) &
\wedge^3(I/I^2)
\end{matrix}
$$
，其微分沿水平方向。因此得到边缘态射
$\text{Tor}_i^A(B, B) \to H^{-i}(L_{B/A})$（$i > 0$）及
$\wedge^i_B(I/I^2) \to \text{Tor}_i^A(B, B)$。最后，
$\text{Tor}_1^A(B, B) = I/I^2$，并且有五项正合列
$$
\text{Tor}_3^A(B, B) \to H^{-3}(L_{B/A}) \to \wedge^2_B(I/I^2) \to
\text{Tor}_2^A(B, B) \to H^{-2}(L_{B/A}) \to 0
$$
，它由低次数项组成。
\end{remark}

\begin{remark}
\label{cotangent-remark-elucidate-degree-two}
设 $A \to B$ 为环同态，$P_\bullet$ 为 $B$ 在 $A$ 上的预解
（评注 \ref{cotangent-remark-resolution}）。令
$J_n = \Ker(P_n \to B)$。注意到
$$
\text{Tor}_2^{P_n}(B, B) = 
\text{Tor}_1^{P_n}(J_n, B) =
\Ker(J_n \otimes_{P_n} J_n \to J_n^2).
$$
因此，由评注 \ref{cotangent-remark-surjection}，$H_2(L_{B/A})$ 典范地等于
$$
\Coker(\text{Tor}_2^{P_1}(B, B) \to \text{Tor}_2^{P_0}(B, B))
$$
。为使其更为明确，我们依照例 \ref{cotangent-example-resolution-length-two}
选取 $P_2$、$P_1$、$P_0$。我们断言
$$
\text{Tor}_2^{P_1}(B, B) =
\wedge^2(\bigoplus\nolimits_{t \in T} B)\ \oplus
\ \bigoplus\nolimits_{t \in T} J_0\ \oplus
\ \text{Tor}_2^{P_0}(B, B)
$$
确切地说，第一直和项的基元素 $x_t \wedge x_{t'}$ 对应于
$J_1 \otimes_{P_1} J_1$ 的元素
$x_t \otimes x_{t'} - x_{t'} \otimes x_t$。对 $f \in J_0$，
第二直和项的元素 $x_t \otimes f$ 对应于
$J_1 \otimes_{P_1} J_1$ 的元素
$x_t \otimes s_0(f) - s_0(f) \otimes x_t$。最后，态射
$\text{Tor}_2^{P_0}(B, B) \to \text{Tor}_2^{P_1}(B, B)$ 由 $s_0$ 给出。
态射
$d_0 - d_1 : \text{Tor}_2^{P_1}(B, B) \to \text{Tor}_2^{P_0}(B, B)$
在最后一个直和项上为零，将 $x_t \otimes f$ 映到
$f \otimes f_t - f_t \otimes f$，并将 $x_t \wedge x_{t'}$ 映到
$f_t \otimes f_{t'} - f_{t'} \otimes f_t$。综上，我们推出有正合列
$$
\wedge^2_B(J_0/J_0^2) \to \text{Tor}_2^{P_0}(B, B) \to H^{-2}(L_{B/A}) \to 0
$$
这样，我们便直接证明了评注 \ref{cotangent-remark-elucidate-ss}
中所讨论的 Quillen 谱序列的一个推论。
\end{remark}






\section{与 Lichtenbaum--Schlessinger 复形的比较}
\label{cotangent-section-compare-higher}

\noindent
设 $A \to B$ 为环同态。文献 \cite{Lichtenbaum-Schlessinger}
对 $\tau_{\geq -2}L_{B/A}$ 给出了相当明确的刻画；它常用于计算奇点的
半普遍形变空间。其构造如下。选取 $A$ 上的多项式代数 $P$ 以及核为
$I$ 的满射 $P \to B$。选取 $I$ 的生成元 $f_t$（$t \in T$），
从而诱导满射 $F = \bigoplus_{t \in T} P \to I$，其中 $F$ 是自由
$P$-模。令 $Rel \subset F$ 为 $F \to I$ 的核；换言之，$Rel$ 是
$f_t$ 之间的关系之集。令 $TrivRel \subset Rel$ 为平凡关系的子模，
即 $Rel$ 中由元素
$(\ldots, f_{t'}, 0, \ldots, 0, -f_t, 0, \ldots)$ 生成的子模。
考虑 $B$-模复形
\begin{equation}
\label{cotangent-equation-lichtenbaum-schlessinger}
Rel/TrivRel \longrightarrow
F \otimes_P B \longrightarrow
\Omega_{P/A} \otimes_P B
\end{equation}
，其中最后一项置于次数 $0$。第一个态射是显然的，第二个态射将对应于
$t \in T$ 的基元素映到 $\text{d}f_t \otimes 1$。

\begin{definition}
\label{cotangent-definition-biderivation}
设 $A \to B$ 为环同态，$M$ 为 $A$ 上的 $(B, B)$-双模。
所谓 {\it $A$-双导子}，是指满足
$\lambda(xy) = x\lambda(y) + \lambda(x)y$ 的 $A$-线性映射
$\lambda : B \to M$。
\end{definition}

\noindent
多项式代数上的双导子很容易描述。

\begin{lemma}
\label{cotangent-lemma-polynomial-ring-unique}
设 $P = A[S]$ 为 $A$ 上的多项式环，$M$ 为 $A$ 上的 $(P, P)$-双模。
给定每个 $s \in S$ 所对应的 $m_s \in M$，存在唯一的 $A$-双导子
$\lambda : P \to M$，它将 $s$ 映到 $m_s$。
\end{lemma}

\begin{proof}
我们在 $M$ 中令
$$
\lambda(s_1 \ldots s_t) =
\sum s_1 \ldots s_{i - 1} m_{s_i} s_{i + 1} \ldots s_t
$$
。按 $A$-线性延拓便得到一个双导子。
\end{proof}

\noindent
下面是比较陈述。读者也可参阅
\cite[page 206, Proposition 12]{Andre-Homologie}，或参阅文献
\cite{Doncel}；后者将复形 (\ref{cotangent-equation-lichtenbaum-schlessinger})
扩充一项，并将比较推广到 $\tau_{\geq -3}$。

\begin{lemma}
\label{cotangent-lemma-compare-higher}
在上述情形中，以 $L$ 表示复形
(\ref{cotangent-equation-lichtenbaum-schlessinger})。在 $D(B)$ 中存在典范态射
$L_{B/A} \to L$，它在 $D(B)$ 中诱导同构
$\tau_{\geq -2}L_{B/A} \to L$。
\end{lemma}

\begin{proof}
令 $P_\bullet \to B$ 为 $B$ 在 $A$ 上的预解（评注
\ref{cotangent-remark-resolution}）。我们将 $L_{B/A}$ 与
$\Omega_{P_\bullet/A} \otimes B$ 等同。为了构造该态射，需要作出若干选择。

\medskip\noindent
选取与给定态射 $P_0 \to B$ 和 $P \to B$ 相容的 $A$-代数同态
$\psi : P_0 \to P$。

\medskip\noindent
对某个集合 $S$ 写成 $P_1 = A[S]$。对 $s \in S$，可写成
$$
\psi(d_0(s) - d_1(s)) = \sum p_{s, t} f_t
$$
，其中 $p_{s, t} \in P$。通过态射
$(\psi \circ d_0, \psi \circ d_1)$，将
$F = \bigoplus_{t \in T} P$ 视为 $(P_1, P_1)$-双模。
由引理 \ref{cotangent-lemma-polynomial-ring-unique}，得到唯一的 $A$-双导子
$\lambda : P_1 \to F$，它将 $s$ 映到坐标为 $p_{s,t}$ 的向量。
按构造，复合
$$
P_1 \longrightarrow F \longrightarrow P
$$
将 $f \in P_1$ 映到 $\psi(d_0(f) - d_1(f))$，因为映射
$f \mapsto \psi(d_0(f) - d_1(f))$ 是一个 $A$-双导子，且在生成元上
与该复合一致。

\medskip\noindent
对 $g \in P_2$，我们断言
$\lambda(d_0(g) - d_1(g) + d_2(g))$ 是 $Rel$ 的元素。事实上，由上一段
最后的说明，它在 $P$ 中的像为
$$
\psi((d_0 - d_1)(d_0(g) - d_1(g) + d_2(g)))
$$
；由《单纯形》第 \ref{simplicial-section-complexes} 节，它为零）。
% P11-E0030：冻结源在上一句末尾保留一个疑似无匹配的右圆括号。

\medskip\noindent
$\psi$ 的选择确定态射
$$
\text{d}\psi \otimes 1 :
\Omega_{P_0/A} \otimes B
\longrightarrow
\Omega_{P/A} \otimes B
$$
将 $\lambda$ 与态射 $F \to F \otimes B$ 复合，得到通常的
$A$-导子，因为 $F \otimes B$ 上的两个 $P_1$-模结构相同。因此
$\lambda$ 确定态射
$$
\overline{\lambda} :
\Omega_{P_1/A} \otimes B
\longrightarrow
F \otimes B
$$
最后，我们得到 $B$-线性态射
$$
q :
\Omega_{P_2/A} \otimes B
\longrightarrow
Rel/TrivRel
$$
，它将 $\text{d}g$ 映到
$\lambda(d_0(g) - d_1(g) + d_2(g))$ 在该商中的类。

\medskip\noindent
图
$$
\xymatrix{
\Omega_{P_3/A} \otimes B \ar[r] \ar[d] &
\Omega_{P_2/A} \otimes B \ar[r] \ar[d]_q &
\Omega_{P_1/A} \otimes B \ar[r] \ar[d]_{\overline{\lambda}} &
\Omega_{P_0/A} \otimes B \ar[d]_{\text{d}\psi \otimes 1} \\
0 \ar[r] &
Rel/TrivRel \ar[r] &
F \otimes B \ar[r] &
\Omega_{P/A} \otimes B
}
$$
交换（计算从略），从而得到该引理的态射。由评注
\ref{cotangent-remark-explicit-comparison-map} 和引理
\ref{cotangent-lemma-relation-with-naive-cotangent-complex}，可见此态射诱导同构
$H_1(L_{B/A}) \to H_1(L)$ 与 $H_0(L_{B/A}) \to H_0(L)$。

\medskip\noindent
还须证明态射 $L_{B/A} \to L$ 诱导同构
$H_2(L_{B/A}) \to H_2(L)$。选取 $B$ 在 $A$ 上的预解，其中
$P_0 = P = A[u_i]$，再依照例 \ref{cotangent-example-resolution-length-two}
选取 $P_1$ 和 $P_2$。在评注 \ref{cotangent-remark-elucidate-degree-two}
中，我们构造了正合列
$$
\wedge^2_B(J_0/J_0^2) \to \text{Tor}_2^{P_0}(B, B) \to H^{-2}(L_{B/A}) \to 0
$$
，其中 $P_0 = P$ 且 $J_0 = \Ker(P \to B) = I$。利用短正合列
$0 \to I \to P \to B \to 0$ 与
$0 \to Rel \to F \to I \to 0$ 计算 Tor 群，得到
$\text{Tor}_2^P(B, B) = \Ker(Rel \otimes B \to F \otimes B)$。
在这一等同下，态射
$\wedge^2_B(I/I^2) \to \text{Tor}_2^P(B, B)$ 的像恰好是
$TrivRel \otimes B$ 的像。因此 $H_2(L_{B/A}) \cong H_2(L)$。

\medskip\noindent
最后还须核查，我们的态射 $L_{B/A} \to L$ 确实诱导上述同构。
以下将直接使用例 \ref{cotangent-example-resolution-length-two} 及评注
\ref{cotangent-remark-elucidate-degree-two}、\ref{cotangent-remark-surjection}
中讨论的记号和结果，不再另行说明。在
$\text{Tor}_2^{P_0}(B, B) = \Ker(I \otimes_P I \to I^2)$
中选取元素 $\xi$。对某些 $h_{t',t} \in P$，写成
$\xi = \sum h_{t', t}f_{t'} \otimes f_t$。沿上述正合列追踪可知，
$\xi$ 对应于元素
$r \in Rel \subset F = \bigoplus_{t \in T} P$ 在 $Rel \otimes B$
中的像，其中第 $t$ 个坐标为
$r_t = \sum_{t' \in T} h_{t', t}f_{t'}$。另一方面，$\xi$ 对应于
$H_2(L_{B/A}) = H_2(\Omega)$ 的如下元素：在二重扩张
$$
0 \to
\text{Tor}_2^\mathcal{O}(\underline{B}, \underline{B})
\to 
\mathcal{J} \otimes_\mathcal{O} \mathcal{J} \to \mathcal{J}
\to
\mathcal{J}/\mathcal{J}^2 \to 0
$$
下取 $\xi$ 的边界，再经
$\text{d} : H_2(\mathcal{J}/\mathcal{J}^2) \to H_2(\Omega)$
所得的像。我们计算该元素的相继超越映射。首先有
$$
\xi = (d_0 - d_1)(- \sum s_0(h_{t', t} f_{t'}) \otimes x_t)
$$
，接着有
$$
\sum s_0(h_{t', t} f_{t'}) x_t = d_0(v_r) - d_1(v_r) + d_2(v_r)
$$
，这是由于我们在例 \ref{cotangent-example-resolution-length-two} 中对变量 $v$
所作的选择。可以选取上面的态射 $\lambda$，使
$\lambda(u_i) = 0$ 且 $\lambda(x_t) = - e_t$，其中 $e_t \in F$
表示对应于 $t \in T$ 的基向量。因此，上述态射 $q$ 的构造将
$\text{d}v_r$ 映到
$$
\lambda(\sum s_0(h_{t', t} f_{t'}) x_t) =
\sum\nolimits_t \left(\sum\nolimits_{t'} h_{t', t}f_{t'}\right) e_t
$$
，这与 $\xi$ 在 $Rel \otimes B$ 中的像相符（上面出现的两个负号相消）。
这一相符性完成了证明。
\end{proof}

\begin{remark}[Lichtenbaum--Schlessinger 复形的函子性]
\label{cotangent-remark-functoriality-lichtenbaum-schlessinger}
考虑环同态的交换方块
$$
\xymatrix{
A' \ar[r] & B' \\
A \ar[u] \ar[r] & B \ar[u]
}
$$
。选取分解
$$
\xymatrix{
A' \ar[r] & P' \ar[r] & B' \\
A \ar[u] \ar[r] & P \ar[u] \ar[r] & B \ar[u]
}
$$
，其中 $P$ 是 $A$ 上的多项式代数，$P'$ 是 $A'$ 上的多项式代数。
选取 $\Ker(P \to B)$ 的生成元 $f_t$（$t \in T$）。
对 $t \in T$，以 $f'_t$ 表示 $f_t$ 在 $P'$ 中的像。选取
$f'_s \in P'$，使得当 $t \in T' = T \amalg S$ 时，元素 $f'_t$
生成 $P' \to B'$ 的核。令
$F = \bigoplus_{t \in T} P$ 且
$F' = \bigoplus_{t' \in T'} P'$。令
$Rel = \Ker(F \to P)$ 与 $Rel' = \Ker(F' \to P')$，其中各坐标上的
态射分别由乘以 $f_t$ 和 $f'_t$ 给出。最后，分别令 $TrivRel$ 与
$TrivRel'$ 等于 $Rel$ 与 $TrivRel$ 中由元素
% P11-E0031：冻结源把第二个承载子模写成 TrivRel，而不是 Rel'。
$(\ldots, f_{t'}, 0, \ldots, 0, -f_t, 0, \ldots)$
生成的子模，其中 $t,t'$ 分别取自 $T$ 与 $T'$。作出这些选择后，得到典范交换图
$$
\xymatrix{
L' : &
Rel'/TrivRel' \ar[r] &
F' \otimes_{P'} B' \ar[r] &
\Omega_{P'/A'} \otimes_{P'} B' \\
L : \ar[u] &
Rel/TrivRel \ar[r] \ar[u] &
F \otimes_P B \ar[r] \ar[u] &
\Omega_{P/A} \otimes_P B \ar[u]
}
$$
此外，追踪引理 \ref{cotangent-lemma-compare-higher} 的证明中所作的选择，
读者可见会得到交换图
$$
\xymatrix{
L_{B'/A'} \ar[r] & L' \\
L_{B/A} \ar[r] \ar[u] & L \ar[u]
}
$$
\end{remark}





\section{局部完全交的余切复形}
\label{cotangent-section-lci}

\noindent
若 $A \to B$ 是局部完全交同态，则 $L_{B/A}$ 是完美复形。
证明这一点的关键是下面的引理。

\begin{lemma}
\label{cotangent-lemma-special-case}
设 $A = \mathbf{Z}[x_1, \ldots, x_n] \to B = \mathbf{Z}$ 为将
$x_i$ 映到 $0$（$i = 1, \ldots, n$）的环同态。令
$I = (x_1, \ldots, x_n) \subset A$。则 $L_{B/A}$ 与 $I/I^2[1]$
拟同构。
\end{lemma}

\begin{proof}
有若干种证明方法。例如，可以明确构造 $B$ 在 $A$ 上的预解并作计算。
我们将使用 (\ref{cotangent-equation-triangle})。确切地说，考虑可区别三角
$$
L_{\mathbf{Z}[x_1, \ldots, x_n]/\mathbf{Z}}
\otimes_{\mathbf{Z}[x_1, \ldots, x_n]} \mathbf{Z} \to
L_{\mathbf{Z}/\mathbf{Z}} \to
L_{\mathbf{Z}/\mathbf{Z}[x_1, \ldots, x_n]}\to
L_{\mathbf{Z}[x_1, \ldots, x_n]/\mathbf{Z}}
\otimes_{\mathbf{Z}[x_1, \ldots, x_n]} \mathbf{Z}[1]
$$
由引理 \ref{cotangent-lemma-cotangent-complex-polynomial-algebra}，复形
$L_{\mathbf{Z}[x_1, \ldots, x_n]/\mathbf{Z}}$ 与
$\Omega_{\mathbf{Z}[x_1, \ldots, x_n]/\mathbf{Z}}$ 拟同构。
由引理 \ref{cotangent-lemma-when-zero}，复形 $L_{\mathbf{Z}/\mathbf{Z}}$
在 $D(\mathbf{Z})$ 中为零。因此，$L_{B/A}$ 只有一个非零上同调群；
由引理 \ref{cotangent-lemma-surjection}，它正如引理中所描述。
\end{proof}

\begin{lemma}
\label{cotangent-lemma-mod-regular-sequence}
设 $A \to B$ 为满环同态，其核 $I$ 由 Koszul 正则序列（例如正则序列）
生成。则 $L_{B/A}$ 与 $I/I^2[1]$ 拟同构。
\end{lemma}

\begin{proof}
设 $f_1, \ldots, f_r \in I$ 是生成 $I$ 的 Koszul 正则序列。
考虑将 $x_i$ 映到 $f_i$ 的环同态
$\mathbf{Z}[x_1, \ldots, x_r] \to A$。由于
$x_1, \ldots, x_r$ 是 $\mathbf{Z}[x_1, \ldots, x_r]$ 中的正则序列，
可见 $x_1, \ldots, x_r$ 上的 Koszul 复形是
$\mathbf{Z} = \mathbf{Z}[x_1, \ldots, x_r]/(x_1, \ldots, x_r)$
在 $\mathbf{Z}[x_1, \ldots, x_r]$ 上的自由预解（参见《代数续篇》引理
\ref{more-algebra-lemma-regular-koszul-regular}）。因此，
$f_1, \ldots, f_r$ 为 Koszul 正则序列这一假设恰好意味着
$B = A \otimes_{\mathbf{Z}[x_1, \ldots, x_r]}^\mathbf{L} \mathbf{Z}$。
所以，由引理 \ref{cotangent-lemma-flat-base-change-cotangent-complex} 与
\ref{cotangent-lemma-special-case}，
$L_{B/A} = L_{\mathbf{Z}/\mathbf{Z}[x_1, \ldots, x_r]}
\otimes_\mathbf{Z}^\mathbf{L} B$。
\end{proof}

\begin{lemma}
\label{cotangent-lemma-mod-Koszul-regular-ideal}
设 $A \to B$ 为满环同态，其核 $I$ 是 Koszul 理想。
则 $L_{B/A}$ 与 $I/I^2[1]$ 拟同构。
\end{lemma}

\begin{proof}
在 $\Spec(A)$ 上局部地，理想 $I$ 由 Koszul 正则序列生成，参见
《代数续篇》定义 \ref{more-algebra-definition-regular-ideal}。
因此结论由引理 \ref{cotangent-lemma-flat-base-change-cotangent-complex} 得到。
\end{proof}

\begin{proposition}
\label{cotangent-proposition-cotangent-complex-local-complete-intersection}
设 $A \to B$ 为局部完全交同态。则 $L_{B/A}$ 是 Tor 振幅位于
$[-1, 0]$ 的完美复形。
\end{proposition}

\begin{proof}
选取核为 $J$ 的满射 $P = A[x_1, \ldots, x_n] \to B$。
由引理 \ref{cotangent-lemma-relation-with-naive-cotangent-complex}，可见
$J/J^2 \to \bigoplus B\text{d}x_i$ 与
$\tau_{\geq -1}L_{B/A}$ 拟同构。注意 $J/J^2$ 是有限投射模
（《代数续篇》引理
\ref{more-algebra-lemma-quasi-regular-ideal-finite-projective}），
故 $\tau_{\geq -1}L_{B/A}$ 是 Tor 振幅位于 $[-1, 0]$ 的完美复形。
因此只须证明当 $i \not \in [-1, 0]$ 时 $H^i(L_{B/A}) = 0$。
这由 (\ref{cotangent-equation-triangle})
$$
L_{P/A} \otimes_P^\mathbf{L} B \to L_{B/A} \to L_{B/P} \to
L_{P/A} \otimes_P^\mathbf{L} B[1]
$$
及引理 \ref{cotangent-lemma-mod-Koszul-regular-ideal} 得到：除非
$i \in \{-1, 0\}$，否则 $H^i(L_{B/P})$ 为零。（对左端项还使用了
引理 \ref{cotangent-lemma-cotangent-complex-polynomial-algebra}。）
\end{proof}








\section{张量积与余切复形}
\label{cotangent-section-tensor-product}

\noindent
设 $R$ 为环，$A$、$B$ 为 $R$-代数。本节讨论
$L_{A \otimes_R B/R}$。我们所需的大部分信息都包含在下图中：
\begin{equation}
\label{cotangent-equation-tensor-product}
\vcenter{
\xymatrix{
L_{A/R} \otimes_A^\mathbf{L} (A \otimes_R B) \ar[r] &
L_{A \otimes_R B/B} \ar[r] &
E \\
L_{A/R} \otimes_A^\mathbf{L} (A \otimes_R B) \ar[r] \ar@{=}[u] &
L_{A \otimes_R B/R} \ar[r] \ar[u] &
L_{A \otimes_R B/A} \ar[u] \\
 &
L_{B/R} \otimes_B^\mathbf{L} (A \otimes_R B) \ar[u] \ar@{=}[r] &
L_{B/R} \otimes_B^\mathbf{L} (A \otimes_R B) \ar[u]
}
}
\end{equation}
说明：中间一行是环同态 $R \to A \to A \otimes_R B$ 的基本三角
(\ref{cotangent-equation-triangle})。中间一列是环同态
$R \to B \to A \otimes_R B$ 的基本三角
(\ref{cotangent-equation-triangle})。其次，$E$ 是 $D(A \otimes_R B)$ 中
“嵌入”右上角的对象，也就是说，它使最上面一行和最右边一列都成为
可区别三角。将《导出范畴》命题 \ref{derived-proposition-9}
应用于左下方块（在缺失位置放置 $0$），可知这样的 $E$ 存在。
更明确地说，例如可以把 $E$ 定义为复形态射
$$
L_{A/R} \otimes_A^\mathbf{L} (A \otimes_R B) \oplus
L_{B/R} \otimes_B^\mathbf{L} (A \otimes_R B)
\longrightarrow
L_{A \otimes_R B/R}
$$
的锥（《导出范畴》定义 \ref{derived-definition-cone}），再应用 TR3
得到以 $E$ 为靶的两个态射。在 Tor 独立的情形下，对象 $E$ 为零。

\begin{lemma}
\label{cotangent-lemma-tensor-product-tor-independent}
若 $A$ 与 $B$ 是 Tor 独立的 $R$-代数，则
(\ref{cotangent-equation-tensor-product}) 中的对象 $E$ 为零。在此情形下有
$$
L_{A \otimes_R B/R} =
L_{A/R} \otimes_A^\mathbf{L} (A \otimes_R B) \oplus
L_{B/R} \otimes_B^\mathbf{L} (A \otimes_R B)
$$
，它由 $A \otimes_R B$-模复形
$L_{A/R} \otimes_R B \oplus L_{B/R} \otimes_R A $
表示。
\end{lemma}

\begin{proof}
前两个陈述直接由引理 \ref{cotangent-lemma-flat-base-change-cotangent-complex} 得到。
最后一个陈述成立，是因为 $L_{A/R}$ 是自由 $A$-模的复形，故
$L_{A/R} \otimes_A^\mathbf{L} (A \otimes_R B)$ 由
$L_{A/R} \otimes_A (A \otimes_R B) = L_{A/R} \otimes_R B$
表示。
\end{proof}

\noindent
一般地，对象 $E$ 有如下性质。

\begin{lemma}
\label{cotangent-lemma-tensor-product}
设 $R$ 为环，$A$、$B$ 为 $R$-代数。对象 $E$ 在
(\ref{cotangent-equation-tensor-product}) 中满足
$$
H^i(E) =
\left\{
\begin{matrix}
0 & \text{若} & i \geq -1 \\
\text{Tor}_1^R(A, B) & \text{若} & i = -2
\end{matrix}
\right.
$$
\end{lemma}

\begin{proof}
我们采用 $E$ 是态射
$L_{B/R} \otimes_B^\mathbf{L} (A \otimes_R B) \to L_{A \otimes_R B/A}$
之锥的描述。由引理 \ref{cotangent-lemma-compare-higher}，典范截断
$\tau_{\geq -2}L_{B/R}$ 与 $\tau_{\geq -2}L_{A \otimes_R B/A}$
由 Lichtenbaum--Schlessinger 复形
(\ref{cotangent-equation-lichtenbaum-schlessinger}) 计算。这些同构与函子性相容
（评注 \ref{cotangent-remark-functoriality-lichtenbaum-schlessinger}）。
因此，本证明使用 Lichtenbaum--Schlessinger 复形。

\medskip\noindent
选取 $R$ 上的多项式代数 $P$ 和满射 $P \to B$。选取该满射之核的
生成元 $f_t \in P$（$t \in T$）。令
$Rel \subset F = \bigoplus_{t \in T} P$ 为态射 $F \to P$ 的核，
该态射将对应于 $t$ 的基向量映到 $f_t$。令
$P_A = A \otimes_R P$ 且
$F_A = A \otimes_R F = P_A \otimes_P F$。令 $Rel_A$ 为
$F_A \to P_A$ 的核。利用正合列
$$
0 \to Rel \to F \to P \to B \to 0
$$
及 Tor 的标准短正合列，得到正合列
$$
A \otimes_R Rel \to Rel_A \to \text{Tor}_1^R(A, B) \to 0
$$
注意 $P_A \to A \otimes_R B$ 是满射，其核由 $P_A$ 中的元素
$1 \otimes f_t$ 生成。以 $TrivRel_A \subset Rel_A$ 表示由元素
$(\ldots, 1 \otimes f_{t'}, 0, \ldots,
0, - 1 \otimes f_t \otimes 1, 0, \ldots)$
% P11-E0032：冻结源在上一坐标中保留疑似多余的 \otimes 1。
生成的 $P_A$-子模。由于
$TrivRel \otimes_R A \to TrivRel_A$ 是满射，得到典范正合列
$$
A \otimes_R (Rel/TrivRel) \to Rel_A/TrivRel_A \to \text{Tor}_1^R(A, B) \to 0
$$
Lichtenbaum--Schlessinger 复形之间的态射由下图给出：
$$
\xymatrix{
Rel_A/TrivRel_A \ar[r] &
F_A \otimes_{P_A} (A \otimes_R B) \ar[r] &
\Omega_{P_A/A \otimes_R B} \otimes_{P_A} (A \otimes_R B) \\
Rel/TrivRel \ar[r] \ar[u]_{-2} &
F \otimes_P B \ar[r] \ar[u]_{-1} &
\Omega_{P/A} \otimes_P B \ar[u]_0
}
$$
% P11-E0033：冻结源在上图右上项保留与该复形基环不符的微分模下标。
注意，对源施用函子 $A \otimes_R - = P_A \otimes_P -$ 后，纵向态射
$-1$ 与 $-0$ 诱导同构，而纵向态射 $-2$ 恰好给出以上所见的那个态射，
它的余核就是所需的 Tor 模。
% P11-E0034：冻结源把图中标为 0 的纵向态射称作 -0。
\end{proof}












\section{环同态的形变与余切复形}
\label{cotangent-section-deformations}

\noindent
本节续接《形变理论》第 \ref{defos-section-deformations} 节，
我们强烈建议读者先阅读该节。首先取满环同态 $A' \to A$，其核是平方为零的
理想 $I$。再假设给定环同态 $A \to B$、$B$-模 $N$ 以及 $A$-模态射
$c : I \to N$。本节要问的是，能否找到适合下图中问号位置的对象：
\begin{equation}
\label{cotangent-equation-to-solve}
\vcenter{
\xymatrix{
0 \ar[r] & N \ar[r] & {?} \ar[r] & B \ar[r] & 0 \\
0 \ar[r] & I \ar[u]^c \ar[r] & A' \ar[u] \ar[r] & A \ar[u] \ar[r] & 0
}
}
\end{equation}
并且还要问该解（若存在）在何种意义下唯一。更确切地说，我们寻找
$A'$-代数的满射 $B' \to B$，其核为平方零理想并与 $N$ 等同，且
$A' \to B'$ 诱导给定态射 $c$。我们称 $B'$ 是
(\ref{cotangent-equation-to-solve}) 的一个{\it 解}。

\begin{lemma}
\label{cotangent-lemma-find-obstruction}
在上述情形中，有：
\begin{enumerate}
\item 存在典范元素 $\xi \in \Ext^2_B(L_{B/A}, N)$；它为零是
(\ref{cotangent-equation-to-solve}) 存在解的充分必要条件。
\item 若存在解，则解的同构类之集在 $\Ext^1_B(L_{B/A}, N)$
的作用下为主齐性空间。
\item 给定一个解 $B'$，所有适合 (\ref{cotangent-equation-to-solve}) 的
$B'$ 自同构之集典范同构于 $\Ext^0_B(L_{B/A}, N)$。
\end{enumerate}
\end{lemma}

\begin{proof}
借助等同 $\NL_{B/A} = \tau_{\geq -1}L_{B/A}$（引理
\ref{cotangent-lemma-relation-with-naive-cotangent-complex}）与
$H^0(L_{B/A}) = \Omega_{B/A}$（引理 \ref{cotangent-lemma-identify-H0}），
我们已在《形变理论》引理 \ref{defos-lemma-huge-diagram} 与
\ref{defos-lemma-choices} 中见过第 (2)、(3) 项。

\medskip\noindent
证明 (1)。粗略地说，将朴素余切复形替换为完整余切复形后，结论便由
《形变理论》评注 \ref{defos-remark-parametrize-solutions} 中的讨论得到。
下面给出更详细的说明。由《形变理论》引理
\ref{defos-lemma-parametrize-solutions} 和评注
\ref{defos-remark-parametrize-solutions}，存在元素
$$
\xi' \in
\Ext^1_A(\NL_{A/A'}, N) =
\Ext^1_B(\NL_{A/A'} \otimes_A^\mathbf{L} B, N) =
\Ext^1_B(L_{A/A'} \otimes_A^\mathbf{L} B, N)
$$
（关于这些等号，参见《形变理论》评注
\ref{defos-remark-parametrize-solutions}，并使用
$\NL_{A'/A} = \tau_{\geq -1} L_{A'/A}$）
% P11-E0035：冻结源在上一等式中把 A/A' 反写为 A'/A。
，使得存在解当且仅当此元素属于态射
$$
\Ext^1_B(\NL_{B/A'}, N) = \Ext^1_B(L_{B/A'}, N)
\longrightarrow
\Ext^1_B(L_{A/A'} \otimes_A^\mathbf{L} B, N)
$$
的像。$A' \to A \to B$ 的可区别三角
(\ref{cotangent-equation-triangle}) 给出长正合列
$$
\ldots \to
\Ext^1_B(L_{B/A'}, N) \to
\Ext^1_B(L_{A/A'} \otimes_A^\mathbf{L} B, N) \to
\Ext^2_B(L_{B/A}, N) \to \ldots
$$
因此，取 $\xi$ 为 $\xi'$ 的像即可。
\end{proof}





\section{模的 Atiyah 类}
\label{cotangent-section-atiyah}

\noindent
设 $A \to B$ 为环同态，$M$ 为 $B$-模。设 $P \to B$ 为
$\mathcal{C}_{B/A}$ 的对象（第 \ref{cotangent-section-compute-L-pi-shriek} 节）。
考虑主部的扩张
$$
0 \to \Omega_{P/A} \otimes_P M \to P^1_{P/A}(M) \to M \to 0
$$
，参见《代数》引理 \ref{algebra-lemma-sequence-of-principal-parts}。
由《代数》评注 \ref{algebra-remark-functoriality-principal-parts}，
此列对 $P$ 具有函子性。因此在 $\mathcal{C}_{B/A}$ 上得到
$\mathcal{O}$-模层的短正合列
$$
0 \to \Omega_{\mathcal{O}/\underline{A}} \otimes_\mathcal{O} \underline{M} \to
P^1_{\mathcal{O}/\underline{A}}(M) \to \underline{M} \to 0
$$
。由引理 \ref{cotangent-lemma-pi-shriek-standard} 及 $L_{B/A}$ 各项的平坦性，有
$L\pi_!(\Omega_{\mathcal{O}/\underline{A}} \otimes_\mathcal{O} \underline{M})
= L_{B/A} \otimes_B M = L_{B/A} \otimes_B^\mathbf{L} M$
。由引理 \ref{cotangent-lemma-pi-lower-shriek-constant-sheaf}，有
$L\pi_!(\underline{M}) = M$。因此得到可区别三角
\begin{equation}
\label{cotangent-equation-atiyah}
L_{B/A} \otimes_B^\mathbf{L} M \to
L\pi_!\left(P^1_{\mathcal{O}/\underline{A}}(M)\right) \to M
\to L_{B/A} \otimes_B^\mathbf{L} M [1]
\end{equation}
；它位于 $D(B)$ 中。这里使用《站点上的上同调》评注
\ref{sites-cohomology-remark-O-homology-B-homology-general}，
以便在 $D(B)$ 中而不只是在 $D(A)$ 中得到可区别三角。

\begin{definition}
\label{cotangent-definition-atiyah-class}
设 $A \to B$ 为环同态，$M$ 为 $B$-模。
(\ref{cotangent-equation-atiyah}) 中的态射
$M \to L_{B/A} \otimes_B^\mathbf{L} M[1]$
称为 $M$ 的{\it Atiyah 类}。
\end{definition}




\section{余切复形}
\label{cotangent-section-cotangent-complex}

\noindent
本节讨论站点上环层同态的余切复形。在后续各节中，我们将其特殊化，
得到赋环拓扑斯态射、赋环空间态射、概形态射、代数空间态射等的余切复形。

\medskip\noindent
设 $\mathcal{C}$ 为站点，以 $\Sh(\mathcal{C})$ 表示相应的拓扑斯。
设 $\mathcal{A}$ 为 $\mathcal{C}$ 上的环层，以
$\mathcal{A}\textit{-Alg}$ 表示 $\mathcal{A}$-代数范畴。
考虑一对伴随函子 $(U, V)$，其中
$V : \mathcal{A}\textit{-Alg} \to \Sh(\mathcal{C})$ 是遗忘函子，而
$U : \Sh(\mathcal{C}) \to \mathcal{A}\textit{-Alg}$ 将集合层
$\mathcal{E}$ 送到 $\mathcal{A}$ 上以 $\mathcal{E}$ 为变量集的多项式代数
$\mathcal{A}[\mathcal{E}]$。令 $X_\bullet$ 为
$\text{Fun}(\mathcal{A}\textit{-Alg}, \mathcal{A}\textit{-Alg})$
中的单纯形对象，它构造于《单纯形》第 \ref{simplicial-section-standard} 节。

\medskip\noindent
现在假设 $\mathcal{A} \to \mathcal{B}$ 是环层同态。于是
$\mathcal{B}$ 是范畴 $\mathcal{A}\textit{-Alg}$ 的对象。以
$\mathcal{P}_\bullet = X_\bullet(\mathcal{B})$ 表示所得的单纯形
$\mathcal{A}$-代数。回忆
$\mathcal{P}_0 = \mathcal{A}[\mathcal{B}]$,
$\mathcal{P}_1 = \mathcal{A}[\mathcal{A}[\mathcal{B}]]$，依此类推。
还要回忆，存在增广
$$
\epsilon : \mathcal{P}_\bullet \longrightarrow \mathcal{B}
$$
，其中把 $\mathcal{B}$ 视为常值单纯形 $\mathcal{A}$-代数。

\begin{definition}
\label{cotangent-definition-standard-resolution-sheaves-rings}
设 $\mathcal{C}$ 为站点，$\mathcal{A} \to \mathcal{B}$ 为
$\mathcal{C}$ 上的环层同态。$\mathcal{B}$ 在 $\mathcal{A}$ 上的
{\it 标准预解}是增广
$\epsilon : \mathcal{P}_\bullet \to \mathcal{B}$
，其各项为
$$
\mathcal{P}_0 = \mathcal{A}[\mathcal{B}],\quad
\mathcal{P}_1 = \mathcal{A}[\mathcal{A}[\mathcal{B}]],\quad \ldots
$$
，态射如上所构造。
\end{definition}

\noindent
有了这一定义，环层同态的余切复形定义如下。我们将使用《站点上的模》
第 \ref{sites-modules-section-differentials} 节中定义的微分模。

\begin{definition}
\label{cotangent-definition-cotangent-complex-morphism-sheaves-rings}
设 $\mathcal{C}$ 为站点，$\mathcal{A} \to \mathcal{B}$ 为
$\mathcal{C}$ 上的环层同态。{\it 余切复形}
$L_{\mathcal{B}/\mathcal{A}}$ 是与单纯形模
$$
\Omega_{\mathcal{P}_\bullet/\mathcal{A}}
\otimes_{\mathcal{P}_\bullet, \epsilon} \mathcal{B}
$$
相伴的 $\mathcal{B}$-模复形，其中
$\epsilon : \mathcal{P}_\bullet \to \mathcal{B}$ 是
$\mathcal{B}$ 在 $\mathcal{A}$ 上的标准预解。我们通常把
$L_{\mathcal{B}/\mathcal{A}}$ 看作 $D(\mathcal{B})$ 的对象。
\end{definition}

\noindent
这些构造满足与第 \ref{cotangent-section-functoriality} 节所讨论者类似的函子性。
确切地说，给定 $\mathcal{C}$ 上环层的交换图
\begin{equation}
\label{cotangent-equation-commutative-square-sheaves}
\vcenter{
\xymatrix{
\mathcal{B} \ar[r] & \mathcal{B}' \\
\mathcal{A} \ar[u] \ar[r] & \mathcal{A}' \ar[u]
}
}
\end{equation}
，存在典范的 $\mathcal{B}$-线性复形态射
$$
L_{\mathcal{B}/\mathcal{A}} \longrightarrow L_{\mathcal{B}'/\mathcal{A}'}
$$
，其构造如下。若 $\mathcal{P}_\bullet \to \mathcal{B}$ 是
$\mathcal{B}$ 在 $\mathcal{A}$ 上的标准预解，而
$\mathcal{P}'_\bullet \to \mathcal{B}'$ 是
$\mathcal{B}'$ 在 $\mathcal{A}'$ 上的标准预解，则存在单纯形
$\mathcal{A}$-代数的典范态射
$\mathcal{P}_\bullet \to \mathcal{P}'_\bullet$，它与增广
$\mathcal{P}_\bullet \to \mathcal{B}$ 与
$\mathcal{P}'_\bullet \to \mathcal{B}'$ 相容。态射
$$
\mathcal{P}_0 = \mathcal{A}[\mathcal{B}]
\longrightarrow
\mathcal{A}'[\mathcal{B}'] = \mathcal{P}'_0,
\quad
\mathcal{P}_1 = \mathcal{A}[\mathcal{A}[\mathcal{B}]]
\longrightarrow
\mathcal{A}'[\mathcal{A}'[\mathcal{B}']] = \mathcal{P}'_1
$$
等由给定态射 $\mathcal{A} \to \mathcal{A}'$ 与
$\mathcal{B} \to \mathcal{B}'$ 给出。所需态射
$L_{\mathcal{B}/\mathcal{A}} \to L_{\mathcal{B}'/\mathcal{A}'}$
便来自微分层上的相伴态射。

\begin{lemma}
\label{cotangent-lemma-pullback-cotangent-morphism-topoi}
设 $f : \Sh(\mathcal{D}) \to \Sh(\mathcal{C})$ 为拓扑斯态射，
$\mathcal{A} \to \mathcal{B}$ 为 $\mathcal{C}$ 上的环层同态。则
$f^{-1}L_{\mathcal{B}/\mathcal{A}} = L_{f^{-1}\mathcal{B}/f^{-1}\mathcal{A}}$。
\end{lemma}

\begin{proof}
图
$$
\xymatrix{
\mathcal{A}\textit{-Alg} \ar[d]_{f^{-1}} \ar[r] &
\Sh(\mathcal{C}) \ar@<1ex>[l] \ar[d]^{f^{-1}} \\
f^{-1}\mathcal{A}\textit{-Alg} \ar[r] & \Sh(\mathcal{D}) \ar@<1ex>[l]
}
$$
交换。
\end{proof}

\begin{lemma}
\label{cotangent-lemma-compute-L-morphism-sheaves-rings}
设 $\mathcal{C}$ 为站点，$\mathcal{A} \to \mathcal{B}$ 为
$\mathcal{C}$ 上的环层同态。则
$H^i(L_{\mathcal{B}/\mathcal{A}})$ 是与预层
$U \mapsto H^i(L_{\mathcal{B}(U)/\mathcal{A}(U)})$ 相伴的层。
\end{lemma}

\begin{proof}
在 $\mathcal{C}$ 上赋予混沌拓扑（此时预层都是层），所得站点记为
$\mathcal{C}'$。存在拓扑斯态射
$f : \Sh(\mathcal{C}) \to \Sh(\mathcal{C}')$，其中 $f_*$
是层到预层的包含，而 $f^{-1}$ 是层化。由引理
\ref{cotangent-lemma-pullback-cotangent-morphism-topoi}，只须对
$\mathcal{C}'$ 证明结论，也就是只须处理 $\mathcal{C}$ 具有混沌拓扑的情形。

\medskip\noindent
若 $\mathcal{C}$ 带有混沌拓扑，则
$L_{\mathcal{B}/\mathcal{A}}(U)$ 等于
$L_{\mathcal{B}(U)/\mathcal{A}(U)}$，因为图
$$
\xymatrix{
\mathcal{A}\textit{-Alg} \ar[d]_{\text{$U$ 上的截面}} \ar[r] &
\Sh(\mathcal{C}) \ar@<1ex>[l] \ar[d]^{\text{$U$ 上的截面}} \\
\mathcal{A}(U)\textit{-Alg} \ar[r] & \textit{Sets} \ar@<1ex>[l]
}
$$
交换。
\end{proof}

\begin{remark}
\label{cotangent-remark-map-sections-over-U}
由引理 \ref{cotangent-lemma-compute-L-morphism-sheaves-rings} 的证明显然可知，
对任意 $U \in \Ob(\mathcal{C})$，存在 $\mathcal{B}(U)$-模复形的典范态射
$L_{\mathcal{B}(U)/\mathcal{A}(U)} \to L_{\mathcal{B}/\mathcal{A}}(U)$
。此外，这些态射与限制态射相容，而复形
$L_{\mathcal{B}/\mathcal{A}}$ 是规则
$U \mapsto L_{\mathcal{B}(U)/\mathcal{A}(U)}$ 的层化。
\end{remark}

\begin{lemma}
\label{cotangent-lemma-H0-L-morphism-sheaves-rings}
设 $\mathcal{C}$ 为站点，$\mathcal{A} \to \mathcal{B}$ 为
$\mathcal{C}$ 上的环层同态。则
$H^0(L_{\mathcal{B}/\mathcal{A}}) = \Omega_{\mathcal{B}/\mathcal{A}}$。
\end{lemma}

\begin{proof}
由引理 \ref{cotangent-lemma-compute-L-morphism-sheaves-rings}、
\ref{cotangent-lemma-identify-H0} 及《站点上的模》引理
\ref{sites-modules-lemma-differentials-sheafify} 得到。
\end{proof}

\begin{lemma}
\label{cotangent-lemma-compute-L-product-sheaves-rings}
设 $\mathcal{C}$ 为站点，$\mathcal{A} \to \mathcal{B}$ 与
$\mathcal{A} \to \mathcal{B}'$ 为 $\mathcal{C}$ 上的环层同态。则
$$
L_{\mathcal{B} \times \mathcal{B}'/\mathcal{A}}
\longrightarrow
L_{\mathcal{B}/\mathcal{A}} \oplus L_{\mathcal{B}'/\mathcal{A}}
$$
在 $D(\mathcal{B} \times \mathcal{B}')$ 中是同构。
\end{lemma}

\begin{proof}
由引理 \ref{cotangent-lemma-compute-L-morphism-sheaves-rings}，只须对环同态证明。
对环而言，这正是引理 \ref{cotangent-lemma-cotangent-complex-product}。
\end{proof}

\noindent
环层余切复形的基本三角是环同态情形之结果的直接推论。

\begin{lemma}
\label{cotangent-lemma-triangle-sheaves-rings}
设 $\mathcal{D}$ 为站点，$\mathcal{A} \to \mathcal{B} \to \mathcal{C}$
为 $\mathcal{D}$ 上的环层同态。存在典范可区别三角
$$
L_{\mathcal{B}/\mathcal{A}} \otimes_\mathcal{B}^\mathbf{L} \mathcal{C}
\to L_{\mathcal{C}/\mathcal{A}} \to L_{\mathcal{C}/\mathcal{B}} \to
L_{\mathcal{B}/\mathcal{A}} \otimes_\mathcal{B}^\mathbf{L} \mathcal{C}[1]
$$
；它位于 $D(\mathcal{C})$ 中。
\end{lemma}

\begin{proof}
我们将使用评注 \ref{cotangent-remark-triangle} 与 \ref{cotangent-remark-explicit-map}
中描述的方法构造该三角，并直接使用其中提到的结果。依照这些评注，
先在 $\mathcal{B} \to \mathcal{C}$ 是环层单射的情形下构造三角。
在此情形下，令
\begin{enumerate}
\item $\mathcal{P}_\bullet$ 为 $\mathcal{B}$ 在 $\mathcal{A}$ 上的标准预解；
\item $\mathcal{Q}_\bullet$ 为 $\mathcal{C}$ 在 $\mathcal{A}$ 上的标准预解；
\item $\mathcal{R}_\bullet$ 为 $\mathcal{C}$ 在 $\mathcal{B}$ 上的标准预解；
\item $\mathcal{S}_\bullet$ 为 $\mathcal{B}$ 在 $\mathcal{B}$ 上的标准预解；
\item $\overline{\mathcal{Q}}_\bullet =
\mathcal{Q}_\bullet \otimes_{\mathcal{P}_\bullet} \mathcal{B}$；并且
\item $\overline{\mathcal{R}}_\bullet =
\mathcal{R}_\bullet \otimes_{\mathcal{S}_\bullet} \mathcal{B}$。
\end{enumerate}
所求可区别三角是与单纯形 $\mathcal{C}$-模的短正合列
$$
0 \to
\Omega_{\mathcal{P}_\bullet/\mathcal{A}}
\otimes_{\mathcal{P}_\bullet} \mathcal{C} \to
\Omega_{\mathcal{Q}_\bullet/\mathcal{A}}
\otimes_{\mathcal{Q}_\bullet} \mathcal{C} \to
\Omega_{\overline{\mathcal{Q}}_\bullet/\mathcal{B}}
\otimes_{\overline{\mathcal{Q}}_\bullet} \mathcal{C} \to 0
$$
相伴的可区别三角。前两项等于引理陈述中三角的前两项。
将最后一项与 $L_{\mathcal{C}/\mathcal{B}}$ 等同，使用复形的拟同构
$$
L_{\mathcal{C}/\mathcal{B}} =
\Omega_{\mathcal{R}_\bullet/\mathcal{B}}
\otimes_{\mathcal{R}_\bullet} \mathcal{C}
\longrightarrow
\Omega_{\overline{\mathcal{R}}_\bullet/\mathcal{B}}
\otimes_{\overline{\mathcal{R}}_\bullet} \mathcal{C}
\longleftarrow
\Omega_{\overline{\mathcal{Q}}_\bullet/\mathcal{B}}
\otimes_{\overline{\mathcal{Q}}_\bullet} \mathcal{C}
$$
上述所有构造都可以先在预层层次完成，再进行层化。因此，要证明列正合，
或证明态射是拟同构，只须对环同态
% P11-E0036：冻结源英文把复数主语 maps 误写为单数 map。
$\mathcal{A}(U) \to \mathcal{B}(U) \to \mathcal{C}(U)$
证明相应陈述；这些结论已经知道。这样便完成了
$\mathcal{B} \to \mathcal{C}$ 为单射时的证明。

\medskip\noindent
一般地，必要时以 $\mathcal{B} \times \mathcal{C}$ 替换
$\mathcal{C}$，从而归约到 $\mathcal{B} \to \mathcal{C}$ 为单射的情形。
由评注 \ref{cotangent-remark-triangle} 中的论证和引理
\ref{cotangent-lemma-compute-L-product-sheaves-rings}，这是可行的。
\end{proof}

\begin{lemma}
\label{cotangent-lemma-stalk-cotangent-complex}
设 $\mathcal{C}$ 为站点，$\mathcal{A} \to \mathcal{B}$ 为
$\mathcal{C}$ 上的环层同态。若 $p$ 是 $\mathcal{C}$ 的点，则
$(L_{\mathcal{B}/\mathcal{A}})_p = L_{\mathcal{B}_p/\mathcal{A}_p}$。
\end{lemma}

\begin{proof}
这是引理 \ref{cotangent-lemma-pullback-cotangent-morphism-topoi} 的特殊情形。
\end{proof}

\noindent
关于朴素余切复形的构造及其性质，参见《站点上的模》
第 \ref{sites-modules-section-netherlander} 节。

\begin{lemma}
\label{cotangent-lemma-compare-cotangent-complex-with-naive}
设 $\mathcal{C}$ 为站点，$\mathcal{A} \to \mathcal{B}$ 为
$\mathcal{C}$ 上的环层同态。存在典范态射
$L_{\mathcal{B}/\mathcal{A}} \to \NL_{\mathcal{B}/\mathcal{A}}$
，它将朴素余切复形与截断
$\tau_{\geq -1}L_{\mathcal{B}/\mathcal{A}}$ 等同。
\end{lemma}

\begin{proof}
令 $\mathcal{P}_\bullet$ 为 $\mathcal{B}$ 在 $\mathcal{A}$ 上的标准预解，
并令 $\mathcal{I} = \Ker(\mathcal{A}[\mathcal{B}] \to \mathcal{B})$。
回忆 $\mathcal{P}_0 = \mathcal{A}[\mathcal{B}]$。该引理中的态射
由交换图
$$
\xymatrix{
L_{\mathcal{B}/\mathcal{A}} \ar[d] & \ldots \ar[r] &
\Omega_{\mathcal{P}_2/\mathcal{A}} \otimes_{\mathcal{P}_2} \mathcal{B}
\ar[r] \ar[d] &
\Omega_{\mathcal{P}_1/\mathcal{A}} \otimes_{\mathcal{P}_1} \mathcal{B}
\ar[r] \ar[d] &
\Omega_{\mathcal{P}_0/\mathcal{A}} \otimes_{\mathcal{P}_0} \mathcal{B}
\ar[d] \\
\NL_{\mathcal{B}/\mathcal{A}} & \ldots \ar[r] &
0 \ar[r] & 
\mathcal{I}/\mathcal{I}^2 \ar[r] &
\Omega_{\mathcal{P}_0/\mathcal{A}} \otimes_{\mathcal{P}_0} \mathcal{B}
}
$$
给出。我们将局部截面 $\text{d}f \otimes b$ 映到
$(d_0(f) - d_1(f))b$ 在 $\mathcal{I}/\mathcal{I}^2$ 中的类，
由此构造以 $\mathcal{I}/\mathcal{I}^2$ 为靶的向下箭头。
这里 $d_i : \mathcal{P}_1 \to \mathcal{P}_0$（$i = 0, 1$）
是单纯形结构的两个面态射。这是有意义的，因为 $d_0-d_1$ 将
$\mathcal{P}_1$ 映入
$\mathcal{I} = \Ker(\mathcal{P}_0 \to \mathcal{B})$。
此规则良定义的验证从略。我们的态射与微分
$\Omega_{\mathcal{P}_1/\mathcal{A}} \otimes_{\mathcal{P}_1} \mathcal{B}
\to \Omega_{\mathcal{P}_0/\mathcal{A}} \otimes_{\mathcal{P}_0} \mathcal{B}$
相容，因为该微分将局部截面 $\text{d}f \otimes b$ 映到
$\text{d}(d_0(f) - d_1(f)) \otimes b$。此外，微分
$\Omega_{\mathcal{P}_2/\mathcal{A}} \otimes_{\mathcal{P}_2} \mathcal{B}
\to \Omega_{\mathcal{P}_1/\mathcal{A}} \otimes_{\mathcal{P}_1} \mathcal{B}$
将局部截面 $\text{d}f \otimes b$ 映到
$\text{d}(d_0(f) - d_1(f) + d_2(f)) \otimes b$，后者被我们的向下箭头
消去。因此得到了复形态射。

\medskip\noindent
为了证明我们的态射在上同调层 $H^0$ 与 $H^{-1}$ 上诱导同构，
作如下论证。令 $\mathcal{C}'$ 为与 $\mathcal{C}$ 具有相同底范畴、
但赋予混沌拓扑的站点。令
$f : \Sh(\mathcal{C}) \to \Sh(\mathcal{C}')$ 为拉回函子是层化的
拓扑斯态射。把给定态射 $\mathcal{A}' \to \mathcal{B}'$
看作 $\mathcal{C}'$ 上的环层同态。上述构造在 $\mathcal{C}'$ 上给出态射
$L_{\mathcal{B}'/\mathcal{A}'} \to \NL_{\mathcal{B}'/\mathcal{A}'}$，
它在 $\mathcal{C}'$ 任一对象 $U$ 上的值正是态射
$$
L_{\mathcal{B}(U)/\mathcal{A}(U)} \to \NL_{\mathcal{B}(U)/\mathcal{A}(U)}
$$
；这是评注 \ref{cotangent-remark-explicit-comparison-map} 中的态射，并在
$H^0$ 与 $H^{-1}$ 上诱导同构。由于
$f^{-1}L_{\mathcal{B}'/\mathcal{A}'} = L_{\mathcal{B}/\mathcal{A}}$
（引理 \ref{cotangent-lemma-pullback-cotangent-morphism-topoi}），且
$f^{-1}\NL_{\mathcal{B}'/\mathcal{A}'} = \NL_{\mathcal{B}/\mathcal{A}}$
（《站点上的模》引理 \ref{sites-modules-lemma-pullback-NL}），
该引理得证。
\end{proof}











\section{模层的 Atiyah 类}
\label{cotangent-section-atiyah-general}

\noindent
设 $\mathcal{C}$ 为站点，$\mathcal{A} \to \mathcal{B}$ 为环层同态，
$\mathcal{F}$ 为 $\mathcal{B}$-模层。令
$\mathcal{P}_\bullet \to \mathcal{B}$ 为 $\mathcal{B}$ 在
$\mathcal{A}$ 上的标准预解（第 \ref{cotangent-section-cotangent-complex} 节）。
对每个 $n \geq 0$，考虑主部的扩张
\begin{equation}
\label{cotangent-equation-atiyah-extension}
0 \to
\Omega_{\mathcal{P}_n/\mathcal{A}} \otimes_{\mathcal{P}_n} \mathcal{F} \to
\mathcal{P}^1_{\mathcal{P}_n/\mathcal{A}}(\mathcal{F}) \to
\mathcal{F} \to 0
\end{equation}
，参见《站点上的模》引理
\ref{sites-modules-lemma-sequence-of-principal-parts}。此构造的函子性
（《站点上的模》评注
\ref{sites-modules-remark-functoriality-principal-parts}）告诉我们，
(\ref{cotangent-equation-atiyah-extension}) 是单纯形
$\mathcal{P}_\bullet$-模短正合列的第 $n$ 次部分（《站点上的上同调》
第 \ref{sites-cohomology-section-simplicial-modules} 节）。
使用《站点上的上同调》评注
\ref{sites-cohomology-remark-homology-augmentation} 中的函子
$L\pi_! : D(\mathcal{P}_\bullet) \to D(\mathcal{B})$
（这里使用 $\mathcal{P}_\bullet \to \mathcal{A}$ 是预解），
% P11-E0037：冻结源把标准预解的增广靶写成 A，而不是 B。
得到可区别三角
\begin{equation}
\label{cotangent-equation-atiyah-general}
L_{\mathcal{B}/\mathcal{A}} \otimes_\mathcal{B}^\mathbf{L} \mathcal{F} \to
L\pi_!\left(\mathcal{P}^1_{\mathcal{P}_\bullet/\mathcal{A}}(\mathcal{F})\right)
\to \mathcal{F} \to
L_{\mathcal{B}/\mathcal{A}} \otimes_\mathcal{B}^\mathbf{L} \mathcal{F} [1]
\end{equation}
；它位于 $D(\mathcal{B})$ 中。

\begin{definition}
\label{cotangent-definition-atiyah-class-general}
设 $\mathcal{C}$ 为站点，$\mathcal{A} \to \mathcal{B}$ 为环层同态，
$\mathcal{F}$ 为 $\mathcal{B}$-模层。
(\ref{cotangent-equation-atiyah-general}) 中的态射 $\mathcal{F} \to
L_{\mathcal{B}/\mathcal{A}} \otimes_\mathcal{B}^\mathbf{L} \mathcal{F}[1]$
称为 $\mathcal{F}$ 的{\it Atiyah 类}。
\end{definition}









\section{赋环空间态射的余切复形}
\label{cotangent-section-cotangent-morphism-ringed-spaces}

\noindent
赋环空间态射的余切复形用上面定义的余切复形来定义。

\begin{definition}
\label{cotangent-definition-cotangent-complex-morphism-ringed-spaces}
设 $f : (X, \mathcal{O}_X) \to (S, \mathcal{O}_S)$ 为赋环空间态射。
$f$ 的{\it 余切复形} $L_f$ 是
$L_f = L_{\mathcal{O}_X/f^{-1}\mathcal{O}_S}$。
我们也使用记号
$L_f = L_{X/S} = L_{\mathcal{O}_X/\mathcal{O}_S}$。
\end{definition}

\noindent
更确切地说，这意味着考虑环层同态
$f^\sharp : f^{-1}\mathcal{O}_S \to \mathcal{O}_X$ 的余切复形
（定义 \ref{cotangent-definition-cotangent-complex-morphism-sheaves-rings}）；
这些环层位于拓扑空间 $X$ 所对应的站点上（《站点》例
\ref{sites-example-site-topological}）。

\begin{lemma}
\label{cotangent-lemma-H0-L-morphism-ringed-spaces}
设 $f : (X, \mathcal{O}_X) \to (S, \mathcal{O}_S)$ 为赋环空间态射。
则 $H^0(L_{X/S}) = \Omega_{X/S}$。
\end{lemma}

\begin{proof}
这是引理 \ref{cotangent-lemma-H0-L-morphism-sheaves-rings} 的特殊情形。
\end{proof}

\begin{lemma}
\label{cotangent-lemma-triangle-ringed-spaces}
设 $f : X \to Y$ 与 $g : Y \to Z$ 为赋环空间态射。则存在典范可区别三角
$$
Lf^* L_{Y/Z} \to L_{X/Z} \to L_{X/Y} \to Lf^*L_{Y/Z}[1]
$$
；它位于 $D(\mathcal{O}_X)$ 中。
\end{lemma}

\begin{proof}
令 $h = g \circ f$，于是
$h^{-1}\mathcal{O}_Z = f^{-1}g^{-1}\mathcal{O}_Z$。由引理
\ref{cotangent-lemma-pullback-cotangent-morphism-topoi}，有
$f^{-1}L_{Y/Z} = L_{f^{-1}\mathcal{O}_Y/h^{-1}\mathcal{O}_Z}$
，且这是平坦 $f^{-1}\mathcal{O}_Y$-模的复形。因此，上述可区别三角
是引理 \ref{cotangent-lemma-triangle-sheaves-rings} 中可区别三角的一个实例，
其中 $\mathcal{A} = h^{-1}\mathcal{O}_Z$、
$\mathcal{B} = f^{-1}\mathcal{O}_Y$，且 $\mathcal{C} = \mathcal{O}_X$。
\end{proof}

\begin{lemma}
\label{cotangent-lemma-compare-cotangent-complex-with-naive-ringed-spaces}
设 $f : (X, \mathcal{O}_X) \to (Y, \mathcal{O}_Y)$ 为赋环空间态射。
存在典范态射 $L_{X/Y} \to \NL_{X/Y}$，它将朴素余切复形与截断
$\tau_{\geq -1}L_{X/Y}$ 等同。
\end{lemma}

\begin{proof}
这是引理 \ref{cotangent-lemma-compare-cotangent-complex-with-naive} 的特殊情形。
\end{proof}






\section{赋环空间的形变与余切复形}
\label{cotangent-section-deformations-ringed-spaces}

\noindent
本节续接《形变理论》第
\ref{defos-section-deformations-ringed-spaces} 节，我们强烈建议读者先阅读该节。
先简要回顾设置。给定赋环空间的一阶加厚
$t : (S, \mathcal{O}_S) \to (S', \mathcal{O}_{S'})$，其中
$\mathcal{J} = \Ker(t^\sharp)$；给定赋环空间态射
$f : (X, \mathcal{O}_X) \to (S, \mathcal{O}_S)$、$\mathcal{O}_X$-模
$\mathcal{G}$，以及模层的 $f$-态射
$c : \mathcal{J} \to \mathcal{G}$。我们要问的是，能否找到适合下图问号位置的对象：
\begin{equation}
\label{cotangent-equation-to-solve-ringed-spaces}
\vcenter{
\xymatrix{
0 \ar[r] & \mathcal{G} \ar[r] & {?} \ar[r] & \mathcal{O}_X \ar[r] & 0 \\
0 \ar[r] & \mathcal{J} \ar[u]^c \ar[r] & \mathcal{O}_{S'} \ar[u] \ar[r] &
\mathcal{O}_S \ar[u] \ar[r] & 0
}
}
\end{equation}
并且还要问该解（若存在）在何种意义下唯一。更确切地说，我们寻找一阶加厚
$i : (X, \mathcal{O}_X) \to (X', \mathcal{O}_{X'})$
和《形变理论》方程 (\ref{defos-equation-morphism-thickenings})
中的加厚态射 $(f, f')$，其中 $\Ker(i^\sharp)$ 与 $\mathcal{G}$ 等同，
且 $(f')^\sharp$ 诱导给定态射 $c$。我们称 $X'$ 是
(\ref{cotangent-equation-to-solve-ringed-spaces}) 的一个{\it 解}。

\begin{lemma}
\label{cotangent-lemma-find-obstruction-ringed-spaces}
在上述情形中，有：
\begin{enumerate}
\item 存在典范元素
$\xi \in \Ext^2_{\mathcal{O}_X}(L_{X/S}, \mathcal{G})$
；它为零是 (\ref{cotangent-equation-to-solve-ringed-spaces}) 存在解的充分必要条件。
\item 若存在解，则解的同构类之集在
$\Ext^1_{\mathcal{O}_X}(L_{X/S}, \mathcal{G})$ 的作用下为主齐性空间。
\item 给定一个解 $X'$，所有适合
(\ref{cotangent-equation-to-solve-ringed-spaces}) 的 $X'$ 自同构之集典范同构于
$\Ext^0_{\mathcal{O}_X}(L_{X/S}, \mathcal{G})$。
\end{enumerate}
\end{lemma}

\begin{proof}
借助等同 $\NL_{X/S} = \tau_{\geq -1}L_{X/S}$（引理
\ref{cotangent-lemma-compare-cotangent-complex-with-naive-ringed-spaces}）与
$H^0(L_{X/S}) = \Omega_{X/S}$（引理
\ref{cotangent-lemma-H0-L-morphism-ringed-spaces}），我们已在《形变理论》引理
\ref{defos-lemma-huge-diagram-ringed-spaces} 与
\ref{defos-lemma-choices-ringed-spaces} 中见过第 (2)、(3) 项。

\medskip\noindent
证明 (1)。粗略地说，将朴素余切复形替换为完整余切复形后，结论便由
《形变理论》评注 \ref{defos-remark-parametrize-solutions-ringed-spaces}
中的讨论得到。下面给出更详细的说明。由《形变理论》引理
\ref{defos-lemma-parametrize-solutions-ringed-spaces}，存在元素
$$
\xi' \in
\Ext^1_{\mathcal{O}_X}(Lf^*\NL_{S/S'}, \mathcal{G}) =
\Ext^1_{\mathcal{O}_X}(Lf^*L_{S/S'}, \mathcal{G})
$$
，使得存在解当且仅当此元素属于态射
$$
\Ext^1_{\mathcal{O}_X}(NL_{X/S'}, \mathcal{G}) =
% P11-E0038：冻结源在上一项用普通字母 NL，而不是本章宏 \NL。
\Ext^1_{\mathcal{O}_X}(L_{X/S'}, \mathcal{G})
\longrightarrow
\Ext^1_{\mathcal{O}_X}(Lf^*L_{S/S'}, \mathcal{G})
$$
的像。引理 \ref{cotangent-lemma-triangle-ringed-spaces} 中
$X \to S \to S'$ 的可区别三角给出长正合列
$$
\ldots \to
\Ext^1_{\mathcal{O}_X}(L_{X/S'}, \mathcal{G}) \to
\Ext^1_{\mathcal{O}_X}(Lf^*L_{S/S'}, \mathcal{G}) \to
\Ext^2_{\mathcal{O}_X}(L_{X/S}, \mathcal{G}) \to \ldots
$$
因此，取 $\xi$ 为 $\xi'$ 的像即可。
\end{proof}










\section{赋环拓扑斯态射的余切复形}
\label{cotangent-section-cotangent-morphism-ringed-topoi}

\noindent
赋环拓扑斯态射的余切复形用上面定义的余切复形来定义。

\begin{definition}
\label{cotangent-definition-cotangent-complex-morphism-ringed-topoi}
设 $(f, f^\sharp) : (\Sh(\mathcal{C}), \mathcal{O}_\mathcal{C}) \to
(\Sh(\mathcal{D}), \mathcal{O}_\mathcal{D})$ 为赋环拓扑斯态射。
$f$ 的{\it 余切复形} $L_f$ 是
$L_f = L_{\mathcal{O}_\mathcal{C}/f^{-1}\mathcal{O}_\mathcal{D}}$。
有时也写作 $L_f = L_{\mathcal{O}_\mathcal{C}/\mathcal{O}_\mathcal{D}}$。
\end{definition}

\noindent
此定义适用于许多情形，但并不总能给出预期的对象。例如，若
$f : X \to Y$ 是概形态射，则 $f$ 诱导大 \'etale 站点的态射
$f_{big} : (\Sch/X)_\etale \to (\Sch/Y)_\etale$
，后者是赋环拓扑斯态射（《下降》评注
\ref{descent-remark-change-topologies-ringed}）。然而，由于
$(f_{big})^\sharp$ 是同构，$L_{f_{big}} = 0$。另一方面，若在把 $f$
看作底层 Zariski 赋环拓扑斯之间态射的意义下取 $L_f$，那么 $L_f$
确实与余切复形 $L_{X/Y}$（定义见下文）一致，其零次上同调层为
$\Omega_{X/Y}$。

\begin{lemma}
\label{cotangent-lemma-H0-L-morphism-ringed-topoi}
设 $f : (\Sh(\mathcal{C}), \mathcal{O}) \to
(\Sh(\mathcal{B}), \mathcal{O}_\mathcal{B})$ 为赋环拓扑斯态射。
则 $H^0(L_f) = \Omega_f$。
\end{lemma}

\begin{proof}
这是引理 \ref{cotangent-lemma-H0-L-morphism-sheaves-rings} 的特殊情形。
\end{proof}

\begin{lemma}
\label{cotangent-lemma-triangle-ringed-topoi}
设 $f : (\Sh(\mathcal{C}_1), \mathcal{O}_1) \to
(\Sh(\mathcal{C}_2), \mathcal{O}_2)$ 与
$g : (\Sh(\mathcal{C}_2), \mathcal{O}_2) \to
(\Sh(\mathcal{C}_3), \mathcal{O}_3)$ 为赋环拓扑斯态射。
则存在典范可区别三角
$$
Lf^* L_g \to L_{g \circ f} \to L_f \to Lf^*L_g[1]
$$
；它位于 $D(\mathcal{O}_1)$ 中。
\end{lemma}

\begin{proof}
令 $h = g \circ f$，于是
$h^{-1}\mathcal{O}_3 = f^{-1}g^{-1}\mathcal{O}_3$。由引理
\ref{cotangent-lemma-pullback-cotangent-morphism-topoi}，有
$f^{-1}L_g = L_{f^{-1}\mathcal{O}_2/h^{-1}\mathcal{O}_3}$
，且这是平坦 $f^{-1}\mathcal{O}_2$-模的复形。因此，上述可区别三角
是引理 \ref{cotangent-lemma-triangle-sheaves-rings} 中可区别三角的一个实例，
其中 $\mathcal{A} = h^{-1}\mathcal{O}_3$、
$\mathcal{B} = f^{-1}\mathcal{O}_2$，且 $\mathcal{C} = \mathcal{O}_1$。
\end{proof}

\begin{lemma}
\label{cotangent-lemma-compare-cotangent-complex-with-naive-ringed-topoi}
设 $f : (\Sh(\mathcal{C}), \mathcal{O}) \to
(\Sh(\mathcal{B}), \mathcal{O}_\mathcal{B})$ 为赋环拓扑斯态射。
存在典范态射 $L_f \to \NL_f$，它将朴素余切复形与截断
$\tau_{\geq -1}L_f$ 等同。
\end{lemma}

\begin{proof}
这是引理 \ref{cotangent-lemma-compare-cotangent-complex-with-naive} 的特殊情形。
\end{proof}








\section{赋环拓扑斯的形变与余切复形}
\label{cotangent-section-deformations-ringed-topoi}

\noindent
本节续接《形变理论》第
\ref{defos-section-deformations-ringed-topoi} 节，我们强烈建议读者先阅读该节。
先简要回顾设置。给定赋环拓扑斯的一阶加厚
$t : (\Sh(\mathcal{B}), \mathcal{O}_\mathcal{B}) \to
(\Sh(\mathcal{B}'), \mathcal{O}_{\mathcal{B}'})$，其中
$\mathcal{J} = \Ker(t^\sharp)$；给定赋环拓扑斯态射
$f : (\Sh(\mathcal{C}), \mathcal{O}) \to
(\Sh(\mathcal{B}), \mathcal{O}_\mathcal{B})$、$\mathcal{O}$-模
$\mathcal{G}$，以及 $f^{-1}\mathcal{O}_\mathcal{B}$-模层的态射
$f^{-1}\mathcal{J} \to \mathcal{G}$。我们要问的是，能否找到适合下图问号位置的对象：
\begin{equation}
\label{cotangent-equation-to-solve-ringed-topoi}
\vcenter{
\xymatrix{
0 \ar[r] & \mathcal{G} \ar[r] & {?} \ar[r] & \mathcal{O} \ar[r] & 0 \\
0 \ar[r] & f^{-1}\mathcal{J} \ar[u]^c \ar[r] &
f^{-1}\mathcal{O}_{\mathcal{B}'} \ar[u] \ar[r] &
f^{-1}\mathcal{O}_\mathcal{B} \ar[u] \ar[r] & 0
}
}
\end{equation}
并且还要问该解（若存在）在何种意义下唯一。更确切地说，我们寻找一阶加厚
$i : (\Sh(\mathcal{C}), \mathcal{O}) \to (\Sh(\mathcal{C}'), \mathcal{O}')$
和《形变理论》方程
(\ref{defos-equation-morphism-thickenings-ringed-topoi}) 中的加厚态射
$(f,f')$，其中 $\Ker(i^\sharp)$ 与 $\mathcal{G}$ 等同，且
$(f')^\sharp$ 诱导给定态射 $c$。我们称
$(\Sh(\mathcal{C}'), \mathcal{O}')$ 是
(\ref{cotangent-equation-to-solve-ringed-topoi}) 的一个{\it 解}。

\begin{lemma}
\label{cotangent-lemma-find-obstruction-ringed-topoi}
在上述情形中，有：
\begin{enumerate}
\item 存在典范元素
$\xi \in \Ext^2_\mathcal{O}(L_f, \mathcal{G})$
；它为零是 (\ref{cotangent-equation-to-solve-ringed-topoi}) 存在解的充分必要条件。
\item 若存在解，则解的同构类之集在
$\Ext^1_\mathcal{O}(L_f, \mathcal{G})$ 的作用下为主齐性空间。
\item 给定一个解 $X'$，所有适合
(\ref{cotangent-equation-to-solve-ringed-topoi}) 的 $X'$ 自同构之集典范同构于
$\Ext^0_\mathcal{O}(L_f, \mathcal{G})$。
% P11-E0039：冻结源在拓扑斯情形中沿用了未定义的 X' 记号。
\end{enumerate}
\end{lemma}

\begin{proof}
借助等同 $\NL_f = \tau_{\geq -1}L_f$（引理
\ref{cotangent-lemma-compare-cotangent-complex-with-naive-ringed-topoi}）与
$H^0(L_f) = \Omega_f$（引理 \ref{cotangent-lemma-H0-L-morphism-ringed-topoi}），
我们已在《形变理论》引理 \ref{defos-lemma-huge-diagram-ringed-topoi}
与 \ref{defos-lemma-choices-ringed-topoi} 中见过第 (2)、(3) 项。

\medskip\noindent
证明 (1)。为使记号与《形变理论》第
\ref{defos-section-deformations-ringed-topoi} 节一致，我们写
$\NL_f = \NL_{\mathcal{O}/\mathcal{O}_\mathcal{B}}$ 及
$L_f = L_{\mathcal{O}/\mathcal{O}_\mathcal{B}}$，并对态射 $t$ 与
$t \circ f$ 使用类似记号。由《形变理论》引理
\ref{defos-lemma-parametrize-solutions-ringed-topoi}，存在元素
$$
\xi' \in
\Ext^1_\mathcal{O}(
Lf^*\NL_{\mathcal{O}_\mathcal{B}/\mathcal{O}_{\mathcal{B}'}}, \mathcal{G}) =
\Ext^1_\mathcal{O}(
Lf^*L_{\mathcal{O}_\mathcal{B}/\mathcal{O}_{\mathcal{B}'}}, \mathcal{G})
$$
，使得存在解当且仅当此元素属于态射
$$
\Ext^1_\mathcal{O}(
\NL_{\mathcal{O}/\mathcal{O}_{\mathcal{B}'}}, \mathcal{G}) =
\Ext^1_\mathcal{O}(
L_{\mathcal{O}/\mathcal{O}_{\mathcal{B}'}}, \mathcal{G})
\longrightarrow
\Ext^1_\mathcal{O}(
Lf^*L_{\mathcal{O}_\mathcal{B}/\mathcal{O}_{\mathcal{B}'}}, \mathcal{G})
$$
的像。引理 \ref{cotangent-lemma-triangle-ringed-topoi} 中 $f$ 与 $t$
的可区别三角给出长正合列
$$
\ldots \to
\Ext^1_\mathcal{O}(
L_{\mathcal{O}/\mathcal{O}_{\mathcal{B}'}}, \mathcal{G}) \to
\Ext^1_\mathcal{O}(
Lf^*L_{\mathcal{O}_\mathcal{B}/\mathcal{O}_{\mathcal{B}'}}, \mathcal{G})
\to
\Ext^1_\mathcal{O}(
L_{\mathcal{O}/\mathcal{O}_\mathcal{B}}, \mathcal{G})
% P11-E0040：冻结源把障碍项的次数写成 Ext^1，而不是 Ext^2。
$$
因此，取 $\xi$ 为 $\xi'$ 的像即可。
\end{proof}












\section{概形态射的余切复形}
\label{cotangent-section-cotangent-morphism-schemes}

\noindent
如上所述，我们如下定义概形态射的余切复形。

\begin{definition}
\label{cotangent-definition-cotangent-morphism-schemes}
设 $f : X \to Y$ 为概形态射。$X$ 在 $Y$ 上的{\it 余切复形
$L_{X/Y}$} 是把 $f$ 视为赋环空间态射时的余切复形
（定义 \ref{cotangent-definition-cotangent-complex-morphism-ringed-spaces}）。
\end{definition}

\noindent
特别地，第 \ref{cotangent-section-cotangent-morphism-ringed-spaces} 节的结果
适用于概形态射的余切复形。下一个引理表明，此定义与环同态的定义相容，
并且还蕴含 $L_{X/Y}$ 是 $D_\QCoh(\mathcal{O}_X)$ 的对象。

\begin{lemma}
\label{cotangent-lemma-morphism-affine-schemes}
设 $f : X \to Y$ 为概形态射。设 $U = \Spec(B) \subset X$ 与
$V = \Spec(A) \subset Y$ 为满足 $f(U) \subset V$ 的仿射开集。
存在复形的典范态射
$$
\widetilde{L_{B/A}} \longrightarrow L_{X/Y}|_U
$$
，它在 $D(\mathcal{O}_U)$ 中是同构。此态射与限制到
$X$ 和 $Y$ 中更小的仿射开集相容。
\end{lemma}

\begin{proof}
由评注 \ref{cotangent-remark-map-sections-over-U}，存在 $B=\mathcal{O}_X(U)$-模
复形的典范态射
$L_{\mathcal{O}_X(U)/f^{-1}\mathcal{O}_Y(U)} \to L_{X/Y}(U)$
，它与进一步的限制相容。利用典范态射
$A \to f^{-1}\mathcal{O}_Y(U)$，得到 $B$-模复形的典范态射
$L_{B/A} \to L_{\mathcal{O}_X(U)/f^{-1}\mathcal{O}_Y(U)}$
。利用 $\widetilde{\ }$ 函子的泛性质（参见《概形》引理
\ref{schemes-lemma-compare-constructions}），得到引理陈述中的态射。
要检验此态射在上同调层上为同构，只须检验它在茎上诱导同构。
这立即由引理 \ref{cotangent-lemma-stalk-cotangent-complex} 与
\ref{cotangent-lemma-localize} 得到（还使用了如下茎的描述：
$\mathcal{O}_X$ 与 $f^{-1}\mathcal{O}_Y$
在点 $\mathfrak p \in \Spec(B)$ 处分别为 $B_\mathfrak p$ 与
$A_\mathfrak q$，其中 $\mathfrak q = A \cap \mathfrak p$；
所用参考结果为《概形》引理 \ref{schemes-lemma-spec-sheaves} 和
《层》引理 \ref{sheaves-lemma-stalk-pullback}）。
\end{proof}

\begin{lemma}
\label{cotangent-lemma-scheme-over-ring}
设 $\Lambda$ 为环，$X$ 为 $\Lambda$ 上的概形。则
$$
L_{X/\Spec(\Lambda)} = L_{\mathcal{O}_X/\underline{\Lambda}}
$$
，其中 $\underline{\Lambda}$ 是 $X$ 上取值为 $\Lambda$ 的常值层。
\end{lemma}

\begin{proof}
令 $p : X \to \Spec(\Lambda)$ 为结构态射，并令
$q : \Spec(\Lambda) \to (*, \Lambda)$ 为显然态射。由引理
\ref{cotangent-lemma-triangle-ringed-spaces} 的可区别三角，只须证明 $L_q=0$。
为此，只须对 $\mathfrak p \in \Spec(\Lambda)$ 证明
$$
(L_q)_\mathfrak p =
L_{\mathcal{O}_{\Spec(\Lambda), \mathfrak p}/\Lambda} =
L_{\Lambda_\mathfrak p/\Lambda}
$$
为零（引理 \ref{cotangent-lemma-stalk-cotangent-complex}），而这由引理
\ref{cotangent-lemma-when-zero} 得到。
\end{proof}






\section{环上概形的余切复形}
\label{cotangent-section-cotangent-schemes-variant}

\noindent
设 $\Lambda$ 为环，$X$ 为 $\Lambda$ 上的概形。写作
$L_{X/\Spec(\Lambda)} = L_{X/\Lambda}$，其合理性由引理
\ref{cotangent-lemma-scheme-over-ring} 保证。本节给出一个类似于引理
\ref{cotangent-lemma-compute-cotangent-complex} 的 $L_{X/\Lambda}$ 描述。
确切地说，我们构造纤维化于 $X_{Zar}$ 上的范畴
$\mathcal{C}_{X/\Lambda}$，并在其上赋予一个（多项式）
$\Lambda$-代数层 $\mathcal{O}$，使得
$$
L_{X/\Lambda} =
L\pi_!(\Omega_{\mathcal{O}/\underline{\Lambda}} \otimes_\mathcal{O}
\underline{\mathcal{O}}_X).
$$
稍后将利用范畴 $\mathcal{C}_{X/\Lambda}$ 构造凝聚层叠的朴素障碍理论。

\medskip\noindent
设 $\Lambda$ 为环，$X$ 为 $\Lambda$ 上的概形。令
$\mathcal{C}_{X/\Lambda}$ 为以下交换图所构成的范畴：
\begin{equation}
\label{cotangent-equation-object}
\vcenter{
\xymatrix{
X \ar[d] & U \ar[l] \ar[d] \\
\Spec(\Lambda) & \mathbf{A} \ar[l]
}
}
\end{equation}
，其中
\begin{enumerate}
\item $U$ 是 $X$ 的开子概形；
\item 存在同构 $\mathbf{A} = \Spec(P)$，其中 $P$ 是 $\Lambda$
上以某个集合为变量集的多项式代数。
\end{enumerate}
换言之，$\mathbf{A}$ 是 $\Spec(\Lambda)$ 上的（无限维）仿射空间。
态射由交换图给出。回忆 $X_{Zar}$ 表示 $X$ 的小 Zariski 站点。
存在遗忘函子
$$
u : \mathcal{C}_{X/\Lambda} \to X_{Zar},\ (U \to \mathbf{A}) \mapsto U
$$
注意，$U$ 上的纤维范畴典范等价于第
\ref{cotangent-section-compute-L-pi-shriek} 节引入的范畴
$\mathcal{C}_{\mathcal{O}_X(U)/\Lambda}$。

\begin{lemma}
\label{cotangent-lemma-category-fibred}
在上述情形中，范畴 $\mathcal{C}_{X/\Lambda}$ 纤维化于 $X_{Zar}$ 上。
\end{lemma}

\begin{proof}
给定 $\mathcal{C}_{X/\Lambda}$ 的对象 $U \to \mathbf{A}$ 和
$X_{Zar}$ 的态射 $U' \to U$，考虑 $\mathcal{C}_{X/\Lambda}$ 的对象
$U' \to \mathbf{A}$，其中 $U' \to \mathbf{A}$ 是
$U' \to U$ 与 $U \to \mathbf{A}$ 的复合。$\mathcal{C}_{X/\Lambda}$ 中的态射
$(U' \to \mathbf{A}) \to (U \to \mathbf{A})$
在 $X_{Zar}$ 上是强笛卡儿的。
\end{proof}

\noindent
在 $\mathcal{C}_{X/\Lambda}$ 上赋予从 $X_{Zar}$ 继承的拓扑
（参见《叠》第 \ref{stacks-section-topology} 节）。函子 $u$ 定义拓扑斯态射
$\pi : \Sh(\mathcal{C}_{X/\Lambda}) \to \Sh(X_{Zar})$。
站点 $\mathcal{C}_{X/\Lambda}$ 带有若干环层：
\begin{enumerate}
\item 由规则
$(U \to \mathbf{A}) \mapsto \Gamma(\mathbf{A}, \mathcal{O}_\mathbf{A})$。
给出的层 $\mathcal{O}$；
\item 由规则 $(U \to \mathbf{A}) \mapsto \mathcal{O}_X(U)$ 给出的层
$\underline{\mathcal{O}}_X = \pi^{-1}\mathcal{O}_X$；
\item 常值层 $\underline{\Lambda}$。
\end{enumerate}
得到赋环拓扑斯态射
\begin{equation}
\label{cotangent-equation-pi-schemes}
\vcenter{
\xymatrix{
(\Sh(\mathcal{C}_{X/\Lambda}), \underline{\mathcal{O}}_X) \ar[r]_i \ar[d]_\pi &
(\Sh(\mathcal{C}_{X/\Lambda}), \mathcal{O}) \\
(\Sh(X_{Zar}), \mathcal{O}_X)
}
}
\end{equation}
态射 $i$ 在底层拓扑斯上是恒等态射，而
$i^\sharp : \mathcal{O} \to \underline{\mathcal{O}}_X$ 是显然同态。
态射 $\pi$ 是《站点上的上同调》情形
\ref{sites-cohomology-situation-fibred-category} 的特殊情形。
以下导出函子将在后文发挥重要作用：
$
Li^* : D(\mathcal{O}) \longrightarrow D(\underline{\mathcal{O}}_X)
$
左伴随于 $Ri_* = i_* : D(\underline{\mathcal{O}}_X) \to D(\mathcal{O})$；以及
$
L\pi_! : D(\underline{\mathcal{O}}_X) \longrightarrow D(\mathcal{O}_X)
$
左伴随于
$\pi^* = \pi^{-1} : D(\mathcal{O}_X) \to D(\underline{\mathcal{O}}_X)$。
借助前面的工作，可以计算 $L\pi_!$。

\begin{remark}
\label{cotangent-remark-compute-L-pi-shriek}
在上述情形中，对每个开集 $U \subset X$，令 $P_{\bullet,U}$ 为
$\mathcal{O}_X(U)$ 在 $\Lambda$ 上的标准预解，并令
$\mathbf{A}_{n,U}=\Spec(P_{n,U})$。于是
$\mathbf{A}_{\bullet, U}$
是 $\mathcal{C}_{X/\Lambda}$ 在 $U$ 上的纤维范畴
$\mathcal{C}_{\mathcal{O}_X(U)/\Lambda}$ 的余单纯形对象。此外，
如评注 \ref{cotangent-remark-resolution} 所讨论，
$\mathbf{A}_{\bullet,U}$ 是《站点上的上同调》引理
\ref{sites-cohomology-lemma-compute-by-cosimplicial-resolution}
中那样的 $\mathcal{C}_{\mathcal{O}_X(U)/\Lambda}$ 余单纯形对象。
由于构造 $U \mapsto \mathbf{A}_{\bullet,U}$ 对 $U$ 具有函子性，
给定 $\mathcal{C}_{X/\Lambda}$ 上的任一（阿贝尔）层 $\mathcal{F}$，
得到预层复形
$$
U \longmapsto \mathcal{F}(\mathbf{A}_{\bullet, U})
$$
，其上同调群计算 $\mathcal{F}$ 在该纤维范畴上的同调。由
《站点上的上同调》引理
\ref{sites-cohomology-lemma-compute-left-derived-pi-shriek}，可知其层化
计算 $L_n\pi_!(\mathcal{F})$。换言之，若一个层复形在次数 $-n$ 的项
是 $U \mapsto \mathcal{F}(\mathbf{A}_{n,U})$ 的层化，则该复形计算
$L\pi_!(\mathcal{F})$。
\end{remark}

\noindent
有了这一评注，现在可以陈述本节的主要结果。

\begin{lemma}
\label{cotangent-lemma-cotangent-morphism-schemes}
在上述情形中，$D(\mathcal{O}_X)$ 中存在典范同构
$$
L_{X/\Lambda} = 
L\pi_!(Li^*\Omega_{\mathcal{O}/\underline{\Lambda}}) =
L\pi_!(i^*\Omega_{\mathcal{O}/\underline{\Lambda}}) =
L\pi_!(\Omega_{\mathcal{O}/\underline{\Lambda}}
\otimes_\mathcal{O} \underline{\mathcal{O}}_X)
$$
\end{lemma}

\begin{proof}
首先注意，对 $\mathcal{C}_{X/\Lambda}$ 的任一对象
$(U \to \mathbf{A})$，层 $\mathcal{O}$ 的值都是 $\Lambda$ 上的多项式代数。
因此 $\Omega_{\mathcal{O}/\underline{\Lambda}}$ 是平坦
$\mathcal{O}$-模，从而引理陈述中的第二、第三个等号成立。

\medskip\noindent
由评注 \ref{cotangent-remark-compute-L-pi-shriek}，对象
$L\pi_!(\Omega_{\mathcal{O}/\underline{\Lambda}}
\otimes_\mathcal{O} \underline{\mathcal{O}}_X)$
由预层复形
$$
U \mapsto
\left(\Omega_{\mathcal{O}/\underline{\Lambda}}
\otimes_\mathcal{O} \underline{\mathcal{O}}_X\right)(\mathbf{A}_{\bullet, U})
=
\Omega_{P_{\bullet, U}/\Lambda} \otimes_{P_{\bullet, U}} \mathcal{O}_X(U) =
L_{\mathcal{O}_X(U)/\Lambda}
$$
的层化计算，其中使用评注 \ref{cotangent-remark-compute-L-pi-shriek} 的记号。
现在评注 \ref{cotangent-remark-map-sections-over-U} 表明
$L\pi_!(\Omega_{\mathcal{O}/\underline{\Lambda}}
\otimes_\mathcal{O} \underline{\mathcal{O}}_X)$
计算 $X$ 上环层同态
$\underline{\Lambda} \to \mathcal{O}_X$ 的余切复形。
由引理 \ref{cotangent-lemma-scheme-over-ring}，这正是所需结论。
\end{proof}








\section{代数空间态射的余切复形}
\label{cotangent-section-cotangent-morphism-spaces}

\noindent
我们利用相应的小 \'etale 站点之间的态射定义代数空间态射的余切复形。

\begin{definition}
\label{cotangent-definition-cotangent-morphism-spaces}
设 $S$ 为概形，$f : X \to Y$ 为 $S$ 上代数空间的态射。
$X$ 在 $Y$ 上的{\it 余切复形 $L_{X/Y}$}，是 $X$ 与 $Y$ 的小
\'etale 站点之间赋环拓扑斯态射 $f_{small}$ 的余切复形
（参见《空间的性质》引理
\ref{spaces-properties-lemma-morphism-ringed-topoi} 及定义
\ref{cotangent-definition-cotangent-complex-morphism-ringed-topoi}）。
\end{definition}

\noindent
特别地，第 \ref{cotangent-section-cotangent-morphism-ringed-topoi} 节的结果
适用于代数空间态射的余切复形。以下引理表明，此定义与环同态及概形的定义
相容，并且 $L_{X/Y}$ 是 $D_\QCoh(\mathcal{O}_X)$ 的对象。

\begin{lemma}
\label{cotangent-lemma-etale-localization}
设 $S$ 为概形。考虑 $S$ 上代数空间的交换图
$$
\xymatrix{
U \ar[d]_p \ar[r]_g & V \ar[d]^q \\
X \ar[r]^f & Y
}
$$
，其中 $p$ 与 $q$ 是 \'etale 的。则在 $D(\mathcal{O}_U)$ 中有典范等同
$L_{X/Y}|_{U_\etale} = L_{U/V}$。
\end{lemma}

\begin{proof}
余切复形的形成与拉回交换（引理
\ref{cotangent-lemma-pullback-cotangent-morphism-topoi}），并且
$p_{small}^{-1}\mathcal{O}_X = \mathcal{O}_U$，且
$g_{small}^{-1}\mathcal{O}_{V_\etale} =
p_{small}^{-1}f_{small}^{-1}\mathcal{O}_{Y_\etale}$
，因为 $q_{small}^{-1}\mathcal{O}_{Y_\etale} =
\mathcal{O}_{V_\etale}$
（《空间的性质》引理
\ref{spaces-properties-lemma-etale-exact-pullback}）。追踪各定义即得
$L_{X/Y}|_{U_\etale} = L_{U/V}$。
\end{proof}

\begin{lemma}
\label{cotangent-lemma-compare-spaces-schemes}
设 $S$ 为概形，$f : X \to Y$ 为 $S$ 上代数空间的态射。假设 $X$ 与
$Y$ 可分别由概形 $X_0$ 与 $Y_0$ 表示。则在 $D(\mathcal{O}_X)$ 中
有典范等同 $L_{X/Y} = \epsilon^*L_{X_0/Y_0}$，其中 $\epsilon$
如《空间的导出范畴》第
\ref{spaces-perfect-section-derived-quasi-coherent-etale} 节所定义，而
$L_{X_0/Y_0}$ 如定义 \ref{cotangent-definition-cotangent-morphism-schemes} 所定义。
% P11-E0041：冻结源英文在 representable 前遗漏系动词 are。
\end{lemma}

\begin{proof}
令 $f_0 : X_0 \to Y_0$ 为对应于 $f$ 的概形态射。存在典范态射
$\epsilon^{-1}f_0^{-1}\mathcal{O}_{Y_0} \to f_{small}^{-1}\mathcal{O}_Y$
，它与
$\epsilon^\sharp : \epsilon^{-1}\mathcal{O}_{X_0} \to \mathcal{O}_X$
相容，因为有交换图
$$
\xymatrix{
X_{0, Zar} \ar[d]_{f_0} & X_\etale \ar[l]^\epsilon \ar[d]^f \\
Y_{0, Zar} & Y_\etale \ar[l]_\epsilon
}
$$
，参见《空间的导出范畴》评注
\ref{spaces-perfect-remark-match-total-direct-images}。因此，由第
\ref{cotangent-section-cotangent-complex} 节所讨论的函子性及引理
\ref{cotangent-lemma-pullback-cotangent-morphism-topoi}，得到典范态射
$$
\epsilon^{-1}L_{X_0/Y_0} =
\epsilon^{-1}L_{\mathcal{O}_{X_0}/f_0^{-1}\mathcal{O}_{Y_0}} =
L_{\epsilon^{-1}\mathcal{O}_{X_0}/\epsilon^{-1}f_0^{-1}\mathcal{O}_{Y_0}}
\longrightarrow
L_{\mathcal{O}_X/f^{-1}_{small}\mathcal{O}_Y} = L_{X/Y}
$$
。为证明诱导态射 $\epsilon^*L_{X_0/Y_0} \to L_{X/Y}$ 是同构，
可以在几何点处的茎上检验（《空间的性质》定理
\ref{spaces-properties-theorem-exactness-stalks}）。我们将用引理
\ref{cotangent-lemma-stalk-cotangent-complex} 计算各茎。设
$\overline{x} : \Spec(k) \to X_0$ 为位于 $x \in X_0$ 上方的几何点，
并设 $\overline{y} = f \circ \overline{x}$ 位于 $y \in Y_0$ 上方。则
$$
L_{X/Y, \overline{x}} =
L_{\mathcal{O}_{X, \overline{x}}/\mathcal{O}_{Y, \overline{y}}}
$$
且
$$
(\epsilon^*L_{X_0/Y_0})_{\overline{x}} =
L_{X_0/Y_0, x} \otimes_{\mathcal{O}_{X_0, x}}
\mathcal{O}_{X, \overline{x}} =
L_{\mathcal{O}_{X_0, x}/\mathcal{O}_{Y_0, y}}
\otimes_{\mathcal{O}_{X_0, x}} \mathcal{O}_{X, \overline{x}}
$$
省略一些细节（提示：使用拉回的茎等于在像点处的茎，参见《站点》引理
\ref{sites-lemma-point-morphism-sites}；还要使用模的相应结果，参见
《站点上的模》引理 \ref{sites-modules-lemma-pullback-stalk}）。
注意，$\mathcal{O}_{X,\overline{x}}$ 是
$\mathcal{O}_{X_0,x}$ 的严格 Hensel 化；对
$\mathcal{O}_{Y,\overline{y}}$ 亦然（《空间的性质》引理
\ref{spaces-properties-lemma-describe-etale-local-ring}）。
因此结论由引理 \ref{cotangent-lemma-cotangent-complex-henselization} 得到。
\end{proof}

\begin{lemma}
\label{cotangent-lemma-space-over-ring}
设 $\Lambda$ 为环，$X$ 为 $\Lambda$ 上的代数空间。则
$$
L_{X/\Spec(\Lambda)} = L_{\mathcal{O}_X/\underline{\Lambda}}
$$
，其中 $\underline{\Lambda}$ 是 $X_\etale$ 上取值为 $\Lambda$ 的常值层。
\end{lemma}

\begin{proof}
令 $p : X \to \Spec(\Lambda)$ 为结构态射，并令
$q : \Spec(\Lambda)_\etale \to (*, \Lambda)$ 为显然态射。由引理
\ref{cotangent-lemma-triangle-ringed-topoi} 的可区别三角，只须证明 $L_q=0$。
为此，根据《空间的性质》定理
\ref{spaces-properties-theorem-exactness-stalks}，只须对几何点
$\overline{t} : \Spec(k) \to \Spec(\Lambda)$ 证明
$$
(L_q)_{\overline{t}} =
L_{\mathcal{O}_{\Spec(\Lambda)_\etale, \overline{t}}/\Lambda}
$$
为零（引理 \ref{cotangent-lemma-stalk-cotangent-complex}）。由于
$\mathcal{O}_{\Spec(\Lambda)_\etale, \overline{t}}$ 是
$\Lambda$ 的某个局部环的严格 Hensel 化（《空间的性质》引理
\ref{spaces-properties-lemma-describe-etale-local-ring}），
结论由引理 \ref{cotangent-lemma-when-zero} 得到。
\end{proof}










\section{环上代数空间的余切复形}
\label{cotangent-section-cotangent-spaces-variant}

\noindent
设 $\Lambda$ 为环，$X$ 为 $\Lambda$ 上的代数空间。写作
$L_{X/\Spec(\Lambda)} = L_{X/\Lambda}$，其合理性由引理
\ref{cotangent-lemma-space-over-ring} 保证。本节给出一个类似于引理
\ref{cotangent-lemma-compute-cotangent-complex} 的 $L_{X/\Lambda}$ 描述。
确切地说，我们构造纤维化于 $X_\etale$ 上的范畴
$\mathcal{C}_{X/\Lambda}$，并在其上赋予一个（多项式）
$\Lambda$-代数层 $\mathcal{O}$，使得
$$
L_{X/\Lambda} =
L\pi_!(\Omega_{\mathcal{O}/\underline{\Lambda}} \otimes_\mathcal{O}
\underline{\mathcal{O}}_X).
$$
稍后将利用范畴 $\mathcal{C}_{X/\Lambda}$ 构造凝聚层叠的朴素障碍理论。

\medskip\noindent
设 $\Lambda$ 为环，$X$ 为 $\Lambda$ 上的代数空间。令
$\mathcal{C}_{X/\Lambda}$ 为以下交换图所构成的范畴：
\begin{equation}
\label{cotangent-equation-object-space}
\vcenter{
\xymatrix{
X \ar[d] & U \ar[l] \ar[d] \\
\Spec(\Lambda) & \mathbf{A} \ar[l]
}
}
\end{equation}
，其中
\begin{enumerate}
\item $U$ 是概形；
\item $U \to X$ 是 \'etale 态射；
\item 存在同构 $\mathbf{A}=\Spec(P)$，其中 $P$ 是 $\Lambda$
上以某个集合为变量集的多项式代数。
\end{enumerate}
换言之，$\mathbf{A}$ 是 $\Spec(\Lambda)$ 上的（无限维）仿射空间。
态射由交换图给出。回忆 $X_\etale$ 表示 $X$ 的小 \'etale 站点，
其对象是 \'etale 于 $X$ 的概形。存在遗忘函子
$$
u : \mathcal{C}_{X/\Lambda} \to X_\etale,
\quad
(U \to \mathbf{A}) \mapsto U
$$
注意，$U$ 上的纤维范畴典范等价于第
\ref{cotangent-section-compute-L-pi-shriek} 节引入的范畴
$\mathcal{C}_{\mathcal{O}_X(U)/\Lambda}$。

\begin{lemma}
\label{cotangent-lemma-category-fibred-space}
在上述情形中，范畴 $\mathcal{C}_{X/\Lambda}$ 纤维化于 $X_\etale$ 上。
\end{lemma}

\begin{proof}
给定 $\mathcal{C}_{X/\Lambda}$ 的对象 $U \to \mathbf{A}$ 和
$X_\etale$ 的态射 $U' \to U$，考虑 $\mathcal{C}_{X/\Lambda}$ 的对象
$U' \to \mathbf{A}$，其中 $U' \to \mathbf{A}$ 是
$U' \to U$ 与 $U \to \mathbf{A}$ 的复合。态射
$(U' \to \mathbf{A}) \to (U \to \mathbf{A})$
在 $X_\etale$ 上是强笛卡儿的。
\end{proof}

\noindent
在 $\mathcal{C}_{X/\Lambda}$ 上赋予从 $X_\etale$ 继承的拓扑
（参见《叠》第 \ref{stacks-section-topology} 节）。函子 $u$ 定义拓扑斯态射
$\pi : \Sh(\mathcal{C}_{X/\Lambda}) \to \Sh(X_\etale)$。
站点 $\mathcal{C}_{X/\Lambda}$ 带有若干环层：
\begin{enumerate}
\item 由规则
$(U \to \mathbf{A}) \mapsto \Gamma(\mathbf{A}, \mathcal{O}_\mathbf{A})$。
给出的层 $\mathcal{O}$；
\item 由规则 $(U \to \mathbf{A}) \mapsto \mathcal{O}_X(U)$ 给出的层
$\underline{\mathcal{O}}_X = \pi^{-1}\mathcal{O}_X$；
\item 常值层 $\underline{\Lambda}$。
\end{enumerate}
得到赋环拓扑斯态射
\begin{equation}
\label{cotangent-equation-pi-spaces}
\vcenter{
\xymatrix{
(\Sh(\mathcal{C}_{X/\Lambda}), \underline{\mathcal{O}}_X) \ar[r]_i \ar[d]_\pi &
(\Sh(\mathcal{C}_{X/\Lambda}), \mathcal{O}) \\
(\Sh(X_\etale), \mathcal{O}_X)
}
}
\end{equation}
态射 $i$ 在底层拓扑斯上是恒等态射，而
$i^\sharp : \mathcal{O} \to \underline{\mathcal{O}}_X$ 是显然同态。
态射 $\pi$ 是《站点上的上同调》情形
\ref{sites-cohomology-situation-fibred-category} 的特殊情形。
以下导出函子将在后文发挥重要作用：
$
Li^* : D(\mathcal{O}) \longrightarrow D(\underline{\mathcal{O}}_X)
$
左伴随于 $Ri_* = i_* : D(\underline{\mathcal{O}}_X) \to D(\mathcal{O})$；以及
$
L\pi_! : D(\underline{\mathcal{O}}_X) \longrightarrow D(\mathcal{O}_X)
$
左伴随于
$\pi^* = \pi^{-1} : D(\mathcal{O}_X) \to D(\underline{\mathcal{O}}_X)$。
借助前面的工作，可以计算 $L\pi_!$。

\begin{remark}
\label{cotangent-remark-compute-L-pi-shriek-spaces}
在上述情形中，对 $X_\etale$ 的每个对象 $U \to X$，令
$P_{\bullet,U}$ 为 $\mathcal{O}_X(U)$ 在 $\Lambda$ 上的标准预解，并令
$\mathbf{A}_{n,U}=\Spec(P_{n,U})$。于是 $\mathbf{A}_{\bullet,U}$
是 $\mathcal{C}_{X/\Lambda}$ 在 $U$ 上的纤维范畴
$\mathcal{C}_{\mathcal{O}_X(U)/\Lambda}$ 的余单纯形对象。此外，
如评注 \ref{cotangent-remark-resolution} 所讨论，
$\mathbf{A}_{\bullet,U}$ 是《站点上的上同调》引理
\ref{sites-cohomology-lemma-compute-by-cosimplicial-resolution}
中那样的 $\mathcal{C}_{\mathcal{O}_X(U)/\Lambda}$ 余单纯形对象。
由于构造 $U \mapsto \mathbf{A}_{\bullet,U}$ 对 $U$ 具有函子性，
给定 $\mathcal{C}_{X/\Lambda}$ 上的任一（阿贝尔）层 $\mathcal{F}$，
得到预层复形
$$
U \longmapsto \mathcal{F}(\mathbf{A}_{\bullet, U})
$$
，其上同调群计算 $\mathcal{F}$ 在该纤维范畴上的同调。由
《站点上的上同调》引理
\ref{sites-cohomology-lemma-compute-left-derived-pi-shriek}，可知其层化
计算 $L_n\pi_!(\mathcal{F})$。换言之，若一个层复形在次数 $-n$ 的项
是 $U \mapsto \mathcal{F}(\mathbf{A}_{n,U})$ 的层化，则该复形计算
$L\pi_!(\mathcal{F})$。
\end{remark}

\noindent
有了这一评注，现在可以陈述本节的主要结果。

\begin{lemma}
\label{cotangent-lemma-cotangent-morphism-spaces}
在上述情形中，$D(\mathcal{O}_X)$ 中存在典范同构
$$
L_{X/\Lambda} = 
L\pi_!(Li^*\Omega_{\mathcal{O}/\underline{\Lambda}}) =
L\pi_!(i^*\Omega_{\mathcal{O}/\underline{\Lambda}}) =
L\pi_!(\Omega_{\mathcal{O}/\underline{\Lambda}}
\otimes_\mathcal{O} \underline{\mathcal{O}}_X)
$$
\end{lemma}

\begin{proof}
首先注意，对 $\mathcal{C}_{X/\Lambda}$ 的任一对象
$(U \to \mathbf{A})$，层 $\mathcal{O}$ 的值都是 $\Lambda$ 上的多项式代数。
因此 $\Omega_{\mathcal{O}/\underline{\Lambda}}$ 是平坦
$\mathcal{O}$-模，从而引理陈述中的第二、第三个等号成立。

\medskip\noindent
由评注 \ref{cotangent-remark-compute-L-pi-shriek-spaces}，对象
$L\pi_!(\Omega_{\mathcal{O}/\underline{\Lambda}}
\otimes_\mathcal{O} \underline{\mathcal{O}}_X)$
由预层复形
$$
U \mapsto
\left(\Omega_{\mathcal{O}/\underline{\Lambda}}
\otimes_\mathcal{O} \underline{\mathcal{O}}_X\right)(\mathbf{A}_{\bullet, U})
=
\Omega_{P_{\bullet, U}/\Lambda} \otimes_{P_{\bullet, U}} \mathcal{O}_X(U) =
L_{\mathcal{O}_X(U)/\Lambda}
$$
的层化计算，其中使用评注 \ref{cotangent-remark-compute-L-pi-shriek-spaces} 的记号。
现在评注 \ref{cotangent-remark-map-sections-over-U} 表明
$L\pi_!(\Omega_{\mathcal{O}/\underline{\Lambda}}
\otimes_\mathcal{O} \underline{\mathcal{O}}_X)$
计算 $X_\etale$ 上环层同态
$\underline{\Lambda} \to \mathcal{O}_X$ 的余切复形。
由引理 \ref{cotangent-lemma-space-over-ring}，这正是所需结论。
\end{proof}






\section{代数空间的纤维积与余切复形}
\label{cotangent-section-fibre-product}

\noindent
设 $S$ 为概形，$X \to B$ 与 $Y \to B$ 为 $S$ 上代数空间的态射。
考虑纤维积 $X \times_B Y$，其投影态射为
$p : X \times_B Y \to X$ 与 $q : X \times_B Y \to Y$。
本节讨论 $L_{X \times_B Y/B}$。我们所需的大部分信息都包含在下图中：
\begin{equation}
\label{cotangent-equation-fibre-product}
\vcenter{
\xymatrix{
Lp^*L_{X/B} \ar[r] &
L_{X \times_B Y/Y} \ar[r] &
E \\
Lp^*L_{X/B} \ar[r] \ar@{=}[u] &
L_{X \times_B Y/B} \ar[r] \ar[u] &
L_{X \times_B Y/X} \ar[u] \\
 &
Lq^*L_{Y/B} \ar[u] \ar@{=}[r] &
Lq^*L_{Y/B} \ar[u]
}
}
\end{equation}
说明：中间一行是引理 \ref{cotangent-lemma-triangle-ringed-topoi} 中态射
$X \times_B Y \to X \to B$ 的基本三角。中间一列是态射
$X \times_B Y \to Y \to B$ 的基本三角。其次，$E$ 是
$D(\mathcal{O}_{X \times_B Y})$ 中“嵌入”右上角的对象，也就是说，
它使最上面一行和最右边一列都成为可区别三角。将《导出范畴》命题
\ref{derived-proposition-9} 应用于左下方块（在缺失位置放置 $0$），
可知这样的 $E$ 存在。更明确地说，例如可以把 $E$ 定义为复形态射
$$
Lp^*L_{X/B} \oplus Lq^*L_{Y/B} \longrightarrow L_{X \times_B Y/B}
$$
的锥（《导出范畴》定义 \ref{derived-definition-cone}），再应用 TR3
得到以 $E$ 为靶的两个态射。在 Tor 独立的情形下，对象 $E$ 为零。

\begin{lemma}
\label{cotangent-lemma-fibre-product-tor-independent}
在上述情形中，若 $X$ 与 $Y$ 在 $B$ 上 Tor 独立，则
(\ref{cotangent-equation-fibre-product}) 中的对象 $E$ 为零。在此情形下有
$$
L_{X \times_B Y/B} = Lp^*L_{X/B} \oplus Lq^*L_{Y/B}
$$
\end{lemma}

\begin{proof}
选取概形 $W$ 和满 \'etale 态射 $W \to B$；选取概形 $U$ 和满
\'etale 态射 $U \to X \times_B W$；选取概形 $V$ 和满
\'etale 态射 $V \to Y \times_B W$。则
$U \times_W V \to X \times_B Y$ 也是满 \'etale 态射。
因此只须证明 $E$ 在 $U \times_W V$ 上的限制为零。由引理
\ref{cotangent-lemma-compare-spaces-schemes} 及《空间的导出范畴》引理
\ref{spaces-perfect-lemma-tor-independent}，归约到概形情形。
取适当的仿射开集后，进一步归约到仿射概形情形。再利用引理
\ref{cotangent-lemma-morphism-affine-schemes}，归约到环的张量积情形，
也就是引理 \ref{cotangent-lemma-tensor-product-tor-independent}。
\end{proof}

\noindent
一般地，对象 $E$ 有如下性质。

\begin{lemma}
\label{cotangent-lemma-fibre-product}
设 $S$ 为概形，$X \to B$ 与 $Y \to B$ 为 $S$ 上代数空间的态射。
(\ref{cotangent-equation-fibre-product}) 中的对象 $E$ 对 $i=0,-1$ 满足
$H^i(E)=0$，而对几何点
$(\overline{x}, \overline{y}) : \Spec(k) \to X \times_B Y$ 有
$$
H^{-2}(E)_{(\overline{x}, \overline{y})} =
\text{Tor}_1^R(A, B) \otimes_{A \otimes_R B} C
$$
，其中 $R = \mathcal{O}_{B, \overline{b}}$、
$A = \mathcal{O}_{X, \overline{x}}$、
$B = \mathcal{O}_{Y, \overline{y}}$，且
$C = \mathcal{O}_{X \times_B Y, (\overline{x}, \overline{y})}$。
\end{lemma}

\begin{proof}
余切复形的形成与取茎及拉回交换，参见引理
\ref{cotangent-lemma-stalk-cotangent-complex} 与
\ref{cotangent-lemma-pullback-cotangent-morphism-topoi}。注意，$C$ 是
$A \otimes_R B$ 的 Hensel 化。由第 \ref{cotangent-section-localization} 节的结果，
$L_{C/R} = L_{A \otimes_R B/R} \otimes_{A \otimes_R B} C$
。因此，$E$ 在该几何点处的茎是态射
$L_{A/R} \otimes C \to L_{A \otimes_R B/R} \otimes C$ 的锥。
所以，引理的结论由环的情形，即引理 \ref{cotangent-lemma-tensor-product} 得到。
\end{proof}
% Generated mechanically from the frozen examples-defos.tex target.
% Do not edit this generated file; edit the bound locale source instead.

% OK, start here.
%

\setcounter{chapter}{92}
\chapter{形变问题}


\phantomsection
\label{examples-defos-section-phantom}


\section{引言}
\label{examples-defos-section-introduction}

\noindent
本章旨在具体研究《形式形变理论》《形变理论》和《余切复形》三章中
发展的一般理论的若干例子。

\medskip\noindent
Schlessinger 的论文 \cite{Sch} 第 3 节也讨论了一些例子。






\section{形变问题的例子}
\label{examples-defos-section-examples}

\noindent
此处应加入的内容如下：
\begin{enumerate}
\item 概形的形变：
\begin{enumerate}
\item Rim--Schlessinger 条件；
\item 切空间的计算；
\item 无穷小形变的计算；
\item 仿射超曲面的形变范畴。
\end{enumerate}
\item 层的形变（例如固定 $X/S$、$S$ 的一个有限型点 $s$，以及
$X_s$ 上的拟凝聚层 $\mathcal{F}_s$）；
\item 代数空间的形变（与概形的形变非常相似，或许甚至更容易？）；
\item 映射的形变（例如概形之间的态射；可以固定靶与源中的一个或两个）；
\item 此处还需补充。
\end{enumerate}





\section{一般纲要}
\label{examples-defos-section-general}

\noindent
本节列出讨论下面几个例子时所遵循的步骤。

\medskip\noindent
步骤 I。在每一节中，我们固定一个诺特环 $\Lambda$，并固定一个有限环同态
$\Lambda \to k$，其中 $k$ 是域。照例，令
$\mathcal{C}_\Lambda = \mathcal{C}_{\Lambda, k}$ 为基范畴，参见
《形式形变理论》定义 \ref{formal-defos-definition-CLambda}。

\medskip\noindent
步骤 II。在每一节中，我们定义一个在 $\mathcal{C}_\Lambda$ 上
以群胚为纤维余纤维化的范畴 $\mathcal{F}$。有时，我们转而考虑函子
$F : \mathcal{C}_\Lambda \to \textit{Sets}$。

\medskip\noindent
步骤 III。我们说明 $\mathcal{F}$ 在何种程度上满足《形式形变理论》第
\ref{formal-defos-section-RS-condition} 节所讨论的
Rim--Schlessinger 条件 (RS)。类似地，我们可以讨论 $\mathcal{F}$
在何种程度上满足 (S1) 与 (S2)，或 $F$ 在何种程度上满足相应的
Schlessinger 条件 (H1) 与 (H2)。参见《形式形变理论》第
\ref{formal-defos-section-schlessinger-conditions} 节。

\medskip\noindent
步骤 IV。设 $x_0$ 为 $\mathcal{F}(k)$ 的对象，换言之，它是
$\mathcal{F}$ 在 $k$ 上的对象。本章使用记号
$$
\Deformationcategory_{x_0} = \mathcal{F}_{x_0}
$$
表示《形式形变理论》注
\ref{formal-defos-remark-localize-cofibered-groupoid} 中构造的预形变范畴。
若 $\mathcal{F}$ 满足 (RS)，则 $\Deformationcategory_{x_0}$ 是形变范畴
（《形式形变理论》引理 \ref{formal-defos-lemma-localize-RS}），
并满足 (S1) 与 (S2)（《形式形变理论》引理
\ref{formal-defos-lemma-RS-implies-S1-S2}）。若满足 (S1) 与 (S2)，
则一个重要问题是切空间
$$
T\Deformationcategory_{x_0} = T_{x_0}\mathcal{F} = T\mathcal{F}_{x_0}
$$
是否有限维（参见《形式形变理论》注
\ref{formal-defos-remark-tangent-space-cofibered-groupoid} 与定义
\ref{formal-defos-definition-tangent-space}）。有限维性保证
$\Deformationcategory_{x_0}$ 具有通用形式对象（《形式形变理论》引理
\ref{formal-defos-lemma-versal-object-existence}）。

\medskip\noindent
步骤 V。若 $\mathcal{F}$ 通过步骤 IV，则下一个问题是 $k$-向量空间
$$
\text{Inf}(\Deformationcategory_{x_0}) = \text{Inf}_{x_0}(\mathcal{F})
$$
（即 $x_0$ 的无穷小自同构空间）是否有限维。若是，则
$\Deformationcategory_{x_0}$ 可由 $\mathcal{C}_\Lambda$ 上函子范畴中的
光滑 pro-可表示群胚给出表示，参见《形式形变理论》定理
\ref{formal-defos-theorem-presentation-deformation-groupoid}。




\section{有限投射模}
\label{examples-defos-section-finite-projective-modules}

\noindent
本节只是热身。有限投射模当然不应有任何“模空间”。

\begin{example}[有限投射模]
\label{examples-defos-example-finite-projective-modules}
如下定义范畴 $\mathcal{F}$：
\begin{enumerate}
\item 对象是二元组 $(A,M)$，其中 $A$ 是 $\mathcal{C}_\Lambda$ 的对象，
$M$ 是有限投射 $A$-模；
\item 态射 $(f,g):(B,N)\to(A,M)$ 由 $\mathcal{C}_\Lambda$ 中的态射
$f:B\to A$ 与映射 $g:N\to M$ 组成，其中该映射是 $f$-线性的，并诱导同构
$N\otimes_{B,f}A\cong M$。
\end{enumerate}
函子 $p:\mathcal{F}\to\mathcal{C}_\Lambda$ 把 $(A,M)$ 送到 $A$，
把 $(f,g)$ 送到 $f$。显然，$p$ 以群胚为纤维余纤维化。
给定有限维 $k$-向量空间 $V$，令 $x_0=(k,V)$ 为
$\mathcal{F}(k)$ 的相应对象。记
$$
\Deformationcategory_V = \mathcal{F}_{x_0}
$$
\end{example}

\noindent
由于局部环上的每个有限投射模都是有限自由模（《代数》引理
\ref{algebra-lemma-finite-projective}），可见
$$
\begin{matrix}
\mathcal{F}(A)\text{ 的对象的} \\
\text{同构类}
\end{matrix}
= \coprod\nolimits_{n \geq 0} \{*\}
$$
尽管这意味着 $\mathcal{F}$ 的形变理论本质上是平凡的，我们仍依照
第 \ref{examples-defos-section-general} 节所列步骤完整推导，作为一个简单例子。

\begin{lemma}
\label{examples-defos-lemma-finite-projective-modules-RS}
例 \ref{examples-defos-example-finite-projective-modules} 满足 Rim--Schlessinger 条件
(RS)。特别地，对任意有限维 $k$-向量空间 $V$，
$\Deformationcategory_V$ 是形变范畴。
\end{lemma}

\begin{proof}
设 $A_1\to A$ 与 $A_2\to A$ 是 $\mathcal{C}_\Lambda$ 的态射，
并假设 $A_2\to A$ 满射。由《形式形变理论》引理
\ref{formal-defos-lemma-RS-2-categorical}，只须证明函子
$\mathcal{F}(A_1 \times_A A_2) \to
\mathcal{F}(A_1) \times_{\mathcal{F}(A)} \mathcal{F}(A_2)$
是范畴等价。

\medskip\noindent
因此，须证明 $A_1\times_AA_2$ 上的有限投射模范畴，等价于
$A_1$ 与 $A_2$ 上的有限投射模范畴在 $A$ 上有限投射模范畴之上的
纤维积。这是《更多代数》引理
\ref{more-algebra-lemma-finitely-presented-module-over-fibre-product}
的特殊情形。回忆，其逆函子把三元组 $(M_1,M_2,\varphi)$ 送到有限投射
$A_1\times_AA_2$-模 $M_1\times_\varphi M_2$；这里 $M_1$ 是有限投射
$A_1$-模，$M_2$ 是有限投射 $A_2$-模，而
$\varphi:M_1\otimes_{A_1}A\to M_2\otimes_{A_2}A$ 是 $A$-模同构。
\end{proof}

\begin{lemma}
\label{examples-defos-lemma-finite-projective-modules-TI}
在例 \ref{examples-defos-example-finite-projective-modules} 中，设 $V$ 为有限维
$k$-向量空间。则
$$
T\Deformationcategory_V = (0)
\quad\text{且}\quad
\text{Inf}(\Deformationcategory_V) = \text{End}_k(V)
$$
均为有限维。
\end{lemma}

\begin{proof}
令 $\mathcal{F}$ 如例 \ref{examples-defos-example-finite-projective-modules}，并置
$x_0=(k,V)\in\Ob(\mathcal{F}(k))$。回忆，
$T\Deformationcategory_V=T_{x_0}\mathcal{F}$ 是二元组 $(x,\alpha)$
的同构类之集；其中 $x$ 是 $\mathcal{F}$ 在对偶数 $k[\epsilon]$ 上的对象，
$\alpha:x\to x_0$ 是 $\mathcal{F}$ 中位于 $k[\epsilon]\to k$ 上方的态射。

\medskip\noindent
在同构意义下，存在唯一的二元组 $(M,\alpha)$，其中 $M$ 是
$k[\epsilon]$ 上的有限投射模，$\alpha:M\to V$ 是诱导同构
$M\otimes_{k[\epsilon]}k\to V$ 的 $k[\epsilon]$-线性映射。
例如，若 $V=k^{\oplus n}$，则取 $M=k[\epsilon]^{\oplus n}$，
并令 $\alpha$ 为显然的映射。

\medskip\noindent
类似地，$\text{Inf}(\Deformationcategory_V)=\text{Inf}_{x_0}(\mathcal{F})$
是 $x_0$ 在 $k[\epsilon]$ 上的平凡形变 $x'_0$ 的自同构之集。
详见《形式形变理论》定义
\ref{formal-defos-definition-infinitesimal-auts}。

\medskip\noindent
给定第二段中的 $(M,\alpha)$，可见 $\text{Inf}_{x_0}(\mathcal{F})$
的元素是满足 $\gamma\bmod\epsilon=\text{id}$ 的自同构
$\gamma:M\to M$。于是可写成
$\gamma=\text{id}_M+\epsilon\psi$，其中
$\psi:M/\epsilon M\to M/\epsilon M$ 是 $k$-线性的。利用 $\alpha$，
可将 $\psi$ 看作 $\text{End}_k(V)$ 的元素，证明完毕。
\end{proof}


\section{群表示}
\label{examples-defos-section-representations}

\noindent
表示的形变理论可能非常有趣。

\begin{example}[群表示]
\label{examples-defos-example-representations}
设 $\Gamma$ 为群。如下定义范畴 $\mathcal{F}$：
\begin{enumerate}
\item 对象是三元组 $(A,M,\rho)$，其中 $A$ 是 $\mathcal{C}_\Lambda$
的对象，$M$ 是有限投射 $A$-模，且
$\rho:\Gamma\to\text{GL}_A(M)$ 是同态；
\item 态射 $(f,g):(B,N,\tau)\to(A,M,\rho)$ 由
$\mathcal{C}_\Lambda$ 中的态射 $f:B\to A$ 与映射 $g:N\to M$
组成，其中该映射是 $f$-线性且 $\Gamma$-等变的，并诱导同构
$N\otimes_{B,f}A\cong M$。
\end{enumerate}
函子 $p:\mathcal{F}\to\mathcal{C}_\Lambda$ 把 $(A,M,\rho)$ 送到 $A$，
把 $(f,g)$ 送到 $f$。显然，$p$ 以群胚为纤维余纤维化。
给定有限维 $k$-向量空间 $V$ 与表示
$\rho_0:\Gamma\to\text{GL}_k(V)$，令 $x_0=(k,V,\rho_0)$ 为
$\mathcal{F}(k)$ 的相应对象。记
$$
\Deformationcategory_{V, \rho_0} = \mathcal{F}_{x_0}
$$
\end{example}

\noindent
由于局部环上的每个有限投射模都是有限自由模（《代数》引理
\ref{algebra-lemma-finite-projective}），可见
$$
\begin{matrix}
\mathcal{F}(A)\text{ 的对象的} \\
\text{同构类}
\end{matrix}
=
\coprod\nolimits_{n \geq 0}\quad
\begin{matrix}
\text{同态 }\rho:\Gamma\to\text{GL}_n(A)\text{ 的}\\
\text{GL}_n(A)\text{-共轭类}
\end{matrix}
$$
这已经比第 \ref{examples-defos-section-finite-projective-modules} 节的讨论更有趣。

\begin{lemma}
\label{examples-defos-lemma-representations-RS}
例 \ref{examples-defos-example-representations} 满足 Rim--Schlessinger 条件 (RS)。
特别地，对任意有限维表示 $\rho_0:\Gamma\to\text{GL}_k(V)$，
$\Deformationcategory_{V,\rho_0}$ 是形变范畴。
\end{lemma}

\begin{proof}
设 $A_1\to A$ 与 $A_2\to A$ 是 $\mathcal{C}_\Lambda$ 的态射，
并假设 $A_2\to A$ 满射。由《形式形变理论》引理
\ref{formal-defos-lemma-RS-2-categorical}，只须证明函子
$\mathcal{F}(A_1 \times_A A_2) \to
\mathcal{F}(A_1) \times_{\mathcal{F}(A)} \mathcal{F}(A_2)$
是范畴等价。

\medskip\noindent
考虑范畴 $\mathcal{F}(A_1)\times_{\mathcal{F}(A)}\mathcal{F}(A_2)$
的对象
$$
((A_1, M_1, \rho_1), (A_2, M_2, \rho_2), (\text{id}_A, \varphi))
$$
。如引理 \ref{examples-defos-lemma-finite-projective-modules-RS} 的证明所示，
可以考虑有限投射 $A_1\times_AA_2$-模 $M_1\times_\varphi M_2$。
由于 $\varphi$ 与给定作用相容，得到
$$
\rho_1 \times \rho_2 : \Gamma \longrightarrow
\text{GL}_{A_1 \times_A A_2}(M_1 \times_\varphi M_2)
$$
于是 $(M_1\times_\varphi M_2,\rho_1\times\rho_2)$ 是
$\mathcal{F}(A_1\times_AA_2)$ 的对象。此构造给出上述函子的拟逆。
\end{proof}

\begin{lemma}
\label{examples-defos-lemma-representations-TI}
在例 \ref{examples-defos-example-representations} 中，设
$\rho_0:\Gamma\to\text{GL}_k(V)$ 为有限维表示。则
$$
T\Deformationcategory_{V, \rho_0} = \Ext^1_{k[\Gamma]}(V, V) =
H^1(\Gamma, \text{End}_k(V))
\quad\text{且}\quad
\text{Inf}(\Deformationcategory_{V, \rho_0}) = H^0(\Gamma, \text{End}_k(V))
$$
因此，$\text{Inf}(\Deformationcategory_{V,\rho_0})$ 总是有限维；
若 $\Gamma$ 有限生成，则 $T\Deformationcategory_{V,\rho_0}$ 也是有限维。
\end{lemma}

\begin{proof}
先处理无穷小自同构。令 $M=V\otimes_k k[\epsilon]$，并赋予其诱导作用
$\rho_0':\Gamma\to\text{GL}_n(M)$。无穷小自同构，即
$\text{Inf}(\Deformationcategory_{V,\rho_0})$ 的元素，由自同构
$\gamma = \text{id} + \epsilon \psi : M \to M$
给出，正如引理 \ref{examples-defos-lemma-finite-projective-modules-TI} 的证明；
此外，$\psi$ 必须与 $\Gamma$ 的作用（由 $\rho_0$ 给出）交换。因此
$$
\text{Inf}(\Deformationcategory_{V, \rho_0}) = H^0(\Gamma, \text{End}_k(V))
$$
，正如引理所述。

\medskip\noindent
其次，设 $(k[\epsilon],M,\rho)$ 为 $\mathcal{F}$ 在 $k[\epsilon]$
上的对象，并设 $\alpha:M\to V$ 为诱导同构 $M/\epsilon M\to V$
的 $\Gamma$-等变映射。由于 $M$ 是自由 $k[\epsilon]$-模，得到
$\Gamma$-模的扩张
$$
0 \to V \to M \xrightarrow{\alpha} V \to 0
$$
左端映射的具体构造从略。反之，若给定如上的 $\Gamma$-模扩张，
便可据此在 $M$ 上赋予 $k[\epsilon]$-模结构，并得到
$\mathcal{F}(k[\epsilon])$ 的对象以及如上的映射 $\alpha$。因此
$$
T\Deformationcategory_{V, \rho_0} = \Ext^1_{k[\Gamma]}(V, V)
$$
，正如引理所述。由《Étale 上同调》引理
\ref{etale-cohomology-lemma-ext-modules-hom}，它等于
$H^1(\Gamma,\text{End}_k(V))$。

\medskip\noindent
关于维数的断言来自《Étale 上同调》引理
\ref{etale-cohomology-lemma-finite-dim-group-cohomology}。
\end{proof}

\noindent
在例 \ref{examples-defos-example-representations} 中，若 $\Gamma$ 有限生成，且
$(V,\rho_0)$ 是 $\Gamma$ 在 $k$ 上的有限维表示，则
$\Deformationcategory_{V,\rho_0}$ 可由 $\mathcal{C}_\Lambda$ 上
函子范畴中的光滑 pro-可表示群胚给出表示，因而具有一个（极小）通用
形式对象。这来自引理 \ref{examples-defos-lemma-representations-RS}、
\ref{examples-defos-lemma-representations-TI} 以及第 \ref{examples-defos-section-general} 节的一般讨论。

\begin{lemma}
\label{examples-defos-lemma-representations-hull}
在例 \ref{examples-defos-example-representations} 中，假设 $\Gamma$ 有限生成。
设 $\rho_0:\Gamma\to\text{GL}_k(V)$ 为有限维表示，并假设
$\Lambda$ 是剩余域为 $k$ 的完备局部环（经典情形）。则函子
$$
F : \mathcal{C}_\Lambda \longrightarrow \textit{Sets},\quad
A \longmapsto \Ob(\Deformationcategory_{V, \rho_0}(A))/\cong
$$
（它取对象的同构类）具有包络。若
$H^0(\Gamma,\text{End}_k(V))=k$，则 $F$ 是 pro-可表示的。
\end{lemma}

\begin{proof}
包络的存在性来自引理 \ref{examples-defos-lemma-representations-RS}、
\ref{examples-defos-lemma-representations-TI}、《形式形变理论》引理
\ref{formal-defos-lemma-RS-implies-S1-S2} 与注
\ref{formal-defos-remark-compose-minimal-into-iso-classes}。

\medskip\noindent
假设 $H^0(\Gamma,\text{End}_k(V))=k$。为证明 $F$ 是 pro-可表示的，
只须证明 $F$ 是形变函子，参见《形式形变理论》定理
\ref{formal-defos-theorem-Schlessinger-prorepresentability}。
换言之，须证明 $F$ 满足 (RS)。为此可使用《形式形变理论》引理
\ref{formal-defos-lemma-RS-associated-functor} 的判据。
若能证明
$$
A \cdot \text{id}_M =
\text{End}_{A[\Gamma]}(M)
$$
对 $\mathcal{F}$ 中每个满足 $M\otimes_Ak$ 作为 $\Gamma$-表示同构于
$V$ 的对象 $(A,M,\rho)$ 都成立，便可得到所需的自同构群满射性。
由于左端包含于右端，只须证明
$\text{length}_A \text{End}_{A[\Gamma]}(M) \leq \text{length}_A A$.
选取两两不同的理想
$(0) = I_n \subset \ldots \subset I_1 \subset A$
，其中 $n=\text{length}(A)$。相应地过滤 $M$，可见只须证明
$\Hom_{A[\Gamma]}(M,I_tM/I_{t+1}M)$ 的长度为 $1$。由于
$I_tM/I_{t+1}M\cong M\otimes_Ak$，且任意 $A[\Gamma]$-模映射
$M\to M\otimes_Ak$ 都唯一地通过商映射 $M\to M\otimes_Ak$ 分解，
从而给出
$$
\text{End}_{A[\Gamma]}(M \otimes_A k) = \text{End}_{k[\Gamma]}(V) = k
$$
，结论即得。
\end{proof}



\section{连续表示}
\label{examples-defos-section-continuous-representations}

\noindent
一个非常有趣的做法是选取无限 Galois 群，研究其表示的形变理论，
参见 \cite{Mazur-deforming}。

\begin{example}[拓扑群的表示]
\label{examples-defos-example-continuous-representations}
设 $\Gamma$ 为拓扑群。如下定义范畴 $\mathcal{F}$：
\begin{enumerate}
\item 对象是三元组 $(A,M,\rho)$，其中 $A$ 是 $\mathcal{C}_\Lambda$
的对象，$M$ 是有限投射 $A$-模，且
$\rho:\Gamma\to\text{GL}_A(M)$ 是连续同态；这里
$\text{GL}_A(M)$ 赋予离散拓扑\footnote{另一种做法是要求由 $\rho$
赋予 $G$-作用的 $A$-模 $M$，是《Étale 上同调》定义
\ref{etale-cohomology-definition-G-module-continuous} 所定义的
$A\text{-}G$-模。不过，由于 $M$ 是有限 $A$-模，两种要求等价。}；
\item 态射 $(f,g):(B,N,\tau)\to(A,M,\rho)$ 由
$\mathcal{C}_\Lambda$ 中的态射 $f:B\to A$ 与映射 $g:N\to M$
组成，其中该映射是 $f$-线性且 $\Gamma$-等变的，并诱导同构
$N\otimes_{B,f}A\cong M$。
\end{enumerate}
函子 $p:\mathcal{F}\to\mathcal{C}_\Lambda$ 把 $(A,M,\rho)$ 送到 $A$，
把 $(f,g)$ 送到 $f$。显然，$p$ 以群胚为纤维余纤维化。
给定有限维 $k$-向量空间 $V$ 与连续表示
$\rho_0:\Gamma\to\text{GL}_k(V)$，令 $x_0=(k,V,\rho_0)$ 为
$\mathcal{F}(k)$ 的相应对象。记
$$
\Deformationcategory_{V, \rho_0} = \mathcal{F}_{x_0}
$$
\end{example}

\noindent
由于局部环上的每个有限投射模都是有限自由模（《代数》引理
\ref{algebra-lemma-finite-projective}），可见
$$
\begin{matrix}
\mathcal{F}(A)\text{ 的对象的} \\
\text{同构类}
\end{matrix}
=
\coprod\nolimits_{n \geq 0}\quad
\begin{matrix}
\text{连续同态 }\rho:\Gamma\to\text{GL}_n(A)\text{ 的}\\
\text{GL}_n(A)\text{-共轭类}
\end{matrix}
$$

\begin{lemma}
\label{examples-defos-lemma-continuous-representations-RS}
例 \ref{examples-defos-example-continuous-representations} 满足 Rim--Schlessinger 条件
(RS)。特别地，对任意有限维连续表示
$\rho_0:\Gamma\to\text{GL}_k(V)$，
$\Deformationcategory_{V,\rho_0}$ 是形变范畴。
\end{lemma}

\begin{proof}
证明与引理 \ref{examples-defos-lemma-representations-RS} 的证明完全相同。
\end{proof}

\begin{lemma}
\label{examples-defos-lemma-continuous-representations-TI}
在例 \ref{examples-defos-example-continuous-representations} 中，设
$\rho_0:\Gamma\to\text{GL}_k(V)$ 为有限维连续表示。则
$$
T\Deformationcategory_{V, \rho_0} = H^1(\Gamma, \text{End}_k(V))
\quad\text{且}\quad
\text{Inf}(\Deformationcategory_{V, \rho_0}) = H^0(\Gamma, \text{End}_k(V))
$$
因此，$\text{Inf}(\Deformationcategory_{V,\rho_0})$ 总是有限维；
若 $\Gamma$ 拓扑有限生成，则 $T\Deformationcategory_{V,\rho_0}$
也是有限维。
\end{lemma}

\begin{proof}
证明与引理 \ref{examples-defos-lemma-representations-TI} 的证明完全相同。
\end{proof}

\noindent
在例 \ref{examples-defos-example-continuous-representations} 中，若 $\Gamma$ 拓扑有限生成，
且 $(V,\rho_0)$ 是 $\Gamma$ 在 $k$ 上的有限维连续表示，则
$\Deformationcategory_{V,\rho_0}$ 可由 $\mathcal{C}_\Lambda$ 上
函子范畴中的光滑 pro-可表示群胚给出表示，因而具有一个（极小）通用
形式对象。这来自引理 \ref{examples-defos-lemma-continuous-representations-RS}、
\ref{examples-defos-lemma-continuous-representations-TI} 以及第
\ref{examples-defos-section-general} 节的一般讨论。

\begin{lemma}
\label{examples-defos-lemma-continuous-representations-hull}
在例 \ref{examples-defos-example-continuous-representations} 中，假设 $\Gamma$
拓扑有限生成。设 $\rho_0:\Gamma\to\text{GL}_k(V)$ 为有限维表示，
并假设 $\Lambda$ 是剩余域为 $k$ 的完备局部环（经典情形）。则函子
$$
F : \mathcal{C}_\Lambda \longrightarrow \textit{Sets},\quad
A \longmapsto \Ob(\Deformationcategory_{V, \rho_0}(A))/\cong
$$
（它取对象的同构类）具有包络。若
$H^0(\Gamma,\text{End}_k(V))=k$，则 $F$ 是 pro-可表示的。
\end{lemma}

\begin{proof}
证明与引理 \ref{examples-defos-lemma-representations-hull} 的证明完全相同。
\end{proof}



\section{分次代数}
\label{examples-defos-section-graded-algebras}

\noindent
在证明极化固有概形之叠是代数叠时，我们将使用本节的例子。因此，
我们考虑齐次分量均为有限投射模的交换分次代数（有时称为“局部有限”）。

\begin{example}[分次代数]
\label{examples-defos-example-graded-algebras}
如下定义范畴 $\mathcal{F}$：
\begin{enumerate}
\item 对象是二元组 $(A,P)$，其中 $A$ 是 $\mathcal{C}_\Lambda$ 的对象，
$P$ 是分次 $A$-代数，并且对所有 $d\geq0$，$P_d$ 都是有限投射 $A$-模；
\item 态射 $(f,g):(B,Q)\to(A,P)$ 由 $\mathcal{C}_\Lambda$ 中的态射
$f:B\to A$ 与映射 $g:Q\to P$ 组成，其中该映射是 $f$-线性的，
并诱导同构 $Q\otimes_{B,f}A\cong P$。
\end{enumerate}
函子 $p:\mathcal{F}\to\mathcal{C}_\Lambda$ 把 $(A,P)$ 送到 $A$，
把 $(f,g)$ 送到 $f$。显然，$p$ 以群胚为纤维余纤维化。
给定满足对所有 $d\geq0$ 均有 $\dim_k(P_d)<\infty$ 的分次
$k$-代数 $P$，令 $x_0=(k,P)$ 为 $\mathcal{F}(k)$ 的相应对象。记
$$
\Deformationcategory_P = \mathcal{F}_{x_0}
$$
\end{example}

\begin{lemma}
\label{examples-defos-lemma-graded-algebras-RS}
例 \ref{examples-defos-example-graded-algebras} 满足 Rim--Schlessinger 条件 (RS)。
特别地，对任意分次 $k$-代数 $P$，$\Deformationcategory_P$ 是形变范畴。
\end{lemma}

\begin{proof}
设 $A_1\to A$ 与 $A_2\to A$ 是 $\mathcal{C}_\Lambda$ 的态射，
并假设 $A_2\to A$ 满射。由《形式形变理论》引理
\ref{formal-defos-lemma-RS-2-categorical}，只须证明函子
$\mathcal{F}(A_1 \times_A A_2) \to
\mathcal{F}(A_1) \times_{\mathcal{F}(A)} \mathcal{F}(A_2)$
是范畴等价。

\medskip\noindent
考虑范畴 $\mathcal{F}(A_1)\times_{\mathcal{F}(A)}\mathcal{F}(A_2)$
的对象
$$
((A_1, P_1), (A_2, P_2), (\text{id}_A, \varphi))
$$
。考虑 $P_1\times_\varphi P_2$。由于
$\varphi : P_1 \otimes_{A_1} A \to P_2 \otimes_{A_2} A$
是分次代数同构，可见 $P_1\times_\varphi P_2$ 的各分次分量都是有限投射
$A_1\times_AA_2$-模，参见引理
\ref{examples-defos-lemma-finite-projective-modules-RS} 的证明。因此，
$P_1\times_\varphi P_2$ 是 $\mathcal{F}(A_1\times_AA_2)$ 的对象。
此构造给出上述函子的拟逆，证明完毕。
\end{proof}

\begin{lemma}
\label{examples-defos-lemma-graded-algebras-TI}
在例 \ref{examples-defos-example-graded-algebras} 中，设 $P$ 为分次 $k$-代数。则
$$
T\Deformationcategory_P
\quad\text{且}\quad
\text{Inf}(\Deformationcategory_P) = \text{Der}_k(P, P)
$$
若 $P$ 在 $k$ 上有限生成，二者均为有限维。
\end{lemma}

\begin{proof}
先处理无穷小自同构。令 $Q=P\otimes_k k[\epsilon]$。
$\text{Inf}(\Deformationcategory_P)$ 的元素由自同构
$\gamma = \text{id} + \epsilon \delta : Q \to Q$
给出，其中 $\delta:P\to P$。$\gamma$ 是分次的，蕴含 $\delta$
是次数为 $0$ 的齐次映射；$\gamma$ 是 $k$-线性的，蕴含 $\delta$
是 $k$-线性的；$\gamma$ 保乘法，蕴含 $\delta$ 是 $k$-导子。
反之，给定次数为 $0$ 的齐次 $k$-导子 $\delta:P\to P$，便得到如上的
自同构 $\gamma=\text{id}+\epsilon\delta$。因此
$$
\text{Inf}(\Deformationcategory_P) = \text{Der}_k(P, P)
$$
，正如引理所述。显然，若 $P$ 由 $P_i$（$0\leq i\leq N$）中的元素生成，
则 $\delta$ 由线性映射 $\delta_i:P_i\to P_i$（$0\leq i\leq N$）确定，
因而
$$
\dim_k \text{Der}_k(P, P) < \infty
$$
，如所需。

\medskip\noindent
为完成引理的证明，我们证明形变空间是有限维的。为此选取分次
$k$-代数的呈示
$$
k[X_1, \ldots, X_n]/(F_1, \ldots, F_m) \longrightarrow P
$$
，其中 $\deg(X_i)=d_i$，且 $F_j$ 是次数为 $e_j$ 的齐次元。
设 $Q$ 为每一次数上有限自由的任意分次 $k[\epsilon]$-代数，并给定同构
$\alpha:Q/\epsilon Q\to P$，使 $(Q,\alpha)$ 定义
$T\Deformationcategory_P$ 的元素。选取次数为 $d_i$ 的齐次元
$q_i\in Q$，使其映到 $X_i$ 在 $P$ 中的像。于是得到
$$
k[\epsilon][X_1, \ldots, X_n] \longrightarrow Q,\quad
X_i \longmapsto q_i
$$
。由于 $P=Q/\epsilon Q$，Nakayama 引理表明此映射满射。
一个简短的追图表明，可以选取次数为 $e_j$ 的齐次元
$F_{\epsilon,j}\in k[\epsilon][X_1,\ldots,X_n]$，使其在 $Q$ 中映到零，
并在 $k[X_1,\ldots,X_n]$ 中映到 $F_j$。于是
$$
k[\epsilon][X_1, \ldots, X_n]/(F_{\epsilon, 1}, \ldots, F_{\epsilon, m})
\longrightarrow Q
$$
由 $Q$ 在 $k[\epsilon]$ 上的平坦性，这是 $Q$ 的一个呈示。写成
$$
F_{\epsilon, j} =  F_j + \epsilon G_j
$$
向量 $(G_1,\ldots,G_m)$ 存在某些不定性。首先，改变
$F_{\epsilon,j}$ 的选择，可以用核
映射 $k[X_1,\ldots,X_n]\to P$ 的核中次数为 $e_j$ 的任意元素修改 $G_j$。
因此，我们不记录 $(G_1,\ldots,G_m)$ 本身，而记录元素
$$
(g_1, \ldots, g_m) \in P_{e_1} \oplus \ldots \oplus P_{e_m}
$$
，其中 $g_j$ 是 $G_j$ 在 $P_{e_j}$ 中的像。此外，若把 $q_i$ 的选择改为
$q_i+\epsilon p_i$，其中 $p_i$ 的次数为 $d_i$，则一个从略的计算表明，
$g_j$ 变为
$$
g_j^{new} = g_j - \sum\nolimits_{i = 1}^n p_i \partial F_j / \partial X_i
$$
由此可见，$Q$ 的同构类由向量 $(G_1,\ldots,G_m)$ 在 $k$-向量空间
$$
W  = \Coker(P_{d_1} \oplus \ldots \oplus P_{d_n}
\xrightarrow{(\frac{\partial F_j}{\partial X_i})}
P_{e_1} \oplus \ldots \oplus P_{e_m})
$$
中的像确定。这样便得到单射
$$
T\Deformationcategory_P \longrightarrow W
$$
显然 $W$ 有限维，故引理成立。
\end{proof}

\noindent
在例 \ref{examples-defos-example-graded-algebras} 中，若 $P$ 是有限生成分次
$k$-代数，则 $\Deformationcategory_P$ 可由 $\mathcal{C}_\Lambda$ 上
函子范畴中的光滑 pro-可表示群胚给出表示，因而具有一个（极小）通用
形式对象。这来自引理 \ref{examples-defos-lemma-graded-algebras-RS}、
\ref{examples-defos-lemma-graded-algebras-TI} 以及第 \ref{examples-defos-section-general} 节的一般讨论。

\begin{lemma}
\label{examples-defos-lemma-graded-algebras-hull}
在例 \ref{examples-defos-example-graded-algebras} 中，假设 $P$ 是有限生成分次
$k$-代数，并假设 $\Lambda$ 是剩余域为 $k$ 的完备局部环（经典情形）。
则函子
$$
F : \mathcal{C}_\Lambda \longrightarrow \textit{Sets},\quad
A \longmapsto \Ob(\Deformationcategory_P(A))/\cong
$$
（它取对象的同构类）具有包络。
\end{lemma}

\begin{proof}
这立即来自引理 \ref{examples-defos-lemma-graded-algebras-RS}、
\ref{examples-defos-lemma-graded-algebras-TI}、《形式形变理论》引理
\ref{formal-defos-lemma-RS-implies-S1-S2} 与注
\ref{formal-defos-remark-compose-minimal-into-iso-classes}。
\end{proof}







\section{环}
\label{examples-defos-section-rings}

\noindent
环的形变理论与仿射概形的形变理论相同。对环和概形而言，我们所说的
形变都是指{\it 平坦}形变。

\begin{example}[环]
\label{examples-defos-example-rings}
如下定义范畴 $\mathcal{F}$：
\begin{enumerate}
\item 对象是二元组 $(A,P)$，其中 $A$ 是 $\mathcal{C}_\Lambda$ 的对象，
$P$ 是平坦 $A$-代数；
\item 态射 $(f,g):(B,Q)\to(A,P)$ 由 $\mathcal{C}_\Lambda$ 中的态射
$f:B\to A$ 与映射 $g:Q\to P$ 组成，其中该映射是 $f$-线性的，
并诱导同构 $Q\otimes_{B,f}A\cong P$。
\end{enumerate}
函子 $p:\mathcal{F}\to\mathcal{C}_\Lambda$ 把 $(A,P)$ 送到 $A$，
把 $(f,g)$ 送到 $f$。显然，$p$ 以群胚为纤维余纤维化。
给定 $k$-代数 $P$，令 $x_0=(k,P)$ 为 $\mathcal{F}(k)$ 的相应对象。记
$$
\Deformationcategory_P = \mathcal{F}_{x_0}
$$
\end{example}

\begin{lemma}
\label{examples-defos-lemma-rings-RS}
例 \ref{examples-defos-example-rings} 满足 Rim--Schlessinger 条件 (RS)。特别地，
对任意 $k$-代数 $P$，$\Deformationcategory_P$ 是形变范畴。
\end{lemma}

\begin{proof}
设 $A_1\to A$ 与 $A_2\to A$ 是 $\mathcal{C}_\Lambda$ 的态射，
并假设 $A_2\to A$ 满射。由《形式形变理论》引理
\ref{formal-defos-lemma-RS-2-categorical}，只须证明函子
$\mathcal{F}(A_1 \times_A A_2) \to
\mathcal{F}(A_1) \times_{\mathcal{F}(A)} \mathcal{F}(A_2)$
是范畴等价。这是《更多代数》引理
\ref{more-algebra-lemma-properties-algebras-over-fibre-product}
的特殊情形。
\end{proof}

\begin{lemma}
\label{examples-defos-lemma-rings-TI}
在例 \ref{examples-defos-example-rings} 中，设 $P$ 为 $k$-代数。则
$$
T\Deformationcategory_P = \text{Ext}^1_P(\NL_{P/k}, P)
\quad\text{且}\quad
\text{Inf}(\Deformationcategory_P) = \text{Der}_k(P, P)
$$
\end{lemma}

\begin{proof}
回忆，$\text{Inf}(\Deformationcategory_P)$ 是 $P$ 在 $k[\epsilon]$
上的平凡形变 $P[\epsilon]=P\otimes_k k[\epsilon]$ 中模 $\epsilon$
等于恒等的自同构之集。由《形变理论》引理
\ref{defos-lemma-huge-diagram}，它等于 $\Hom_P(\Omega_{P/k},P)$；
再由《代数》引理 \ref{algebra-lemma-universal-omega}，后者等于
$\text{Der}_k(P,P)$。

\medskip\noindent
回忆，$T\Deformationcategory_P$ 是 $P$ 在 $k[\epsilon]$ 上的平坦形变
$Q$ 的同构类之集；更确切地说，它是
$\Deformationcategory_P(k[\epsilon])$ 的同构类之集。再回忆，满足
$Q/\epsilon Q=P$ 的 $k[\epsilon]$-代数 $Q$ 在 $k[\epsilon]$ 上平坦，
当且仅当
$$
0 \to P \xrightarrow{\epsilon} Q \to P \to 0
$$
正合。这在《更多态射》引理 \ref{more-morphisms-lemma-deform} 中得到证明；
更一般的结果见《形变理论》引理 \ref{defos-lemma-deform-module}。
因此可应用《形变理论》引理 \ref{defos-lemma-choices}，得到这类形变的
同构类之集等于 $\text{Ext}^1_P(\NL_{P/k},P)$。
\end{proof}

\begin{lemma}
\label{examples-defos-lemma-smooth}
在例 \ref{examples-defos-example-rings} 中，设 $P$ 为光滑 $k$-代数。则
$T\Deformationcategory_P=(0)$。
\end{lemma}

\begin{proof}
由引理 \ref{examples-defos-lemma-rings-TI}，须证明
$\text{Ext}^1_P(\NL_{P/k},P)=(0)$。由于 $k\to P$ 光滑，
$\NL_{P/k}$ 拟同构于由置于次数 $0$ 的有限投射 $P$-模组成的复形。
\end{proof}

\begin{lemma}
\label{examples-defos-lemma-finite-type-rings-TI}
在引理 \ref{examples-defos-lemma-rings-TI} 中，若 $P$ 是有限型 $k$-代数，则
\begin{enumerate}
\item $\text{Inf}(\Deformationcategory_P)$ 有限维，当且仅当
$\dim(P)=0$；
\item 若 $\Spec(P)\to\Spec(k)$ 除有限多个点外均光滑，则
$T\Deformationcategory_P$ 有限维。
\end{enumerate}
\end{lemma}

\begin{proof}
证明 (1)。把 $\text{Der}_k(P,P)$ 看作 $P$-模。若它在 $k$ 上有限维，
则作为 $P$-模具有有限长度，因而其支集包含于 $\Spec(P)$ 的有限多个闭点
（《代数》引理 \ref{algebra-lemma-simple-pieces}）。由于
$\text{Der}_k(P,P)=\Hom_P(\Omega_{P/k},P)$，可见
$\text{Der}_k(P, P)_\mathfrak p = \text{Der}_k(P_\mathfrak p, P_\mathfrak p)$
对任意素理想 $\mathfrak p\subset P$ 都成立（这里使用《代数》引理
\ref{algebra-lemma-differentials-localize},
\ref{algebra-lemma-differentials-finitely-presented} 以及
\ref{algebra-lemma-hom-from-finitely-presented}）。设 $\mathfrak p$
为 $P$ 的极小素理想，对应于维数 $d>0$ 的不可约分支。那么
$P_\mathfrak p$ 是本质有限型于 $k$ 的 Artin 局部环，具有剩余域；并且，
例如由《代数》引理 \ref{algebra-lemma-characterize-smooth-over-field}，
$\Omega_{P_\mathfrak p/k}$ 非零。
% P11-E0047：冻结源句子在 “with residue field” 之后缺少应有的说明。
Artin 局部环上的任意非零有限模，都有同构于剩余域的子模与商模。因此
$\text{Der}_k(P_\mathfrak p, P_\mathfrak p) =
\Hom_{P_\mathfrak p}(\Omega_{P_\mathfrak p/k}, P_\mathfrak p)$
也非零。综合以上各点即得 (1)。

\medskip\noindent
证明 (2)。对 $P$ 的素理想 $\mathfrak p$，我们使用
$\NL_{P_\mathfrak p/k} = (\NL_{P/k})_\mathfrak p$
（《代数》引理 \ref{algebra-lemma-localize-NL}）以及
$\text{Ext}_P^1(\NL_{P/k}, P)_\mathfrak p =
\text{Ext}_{P_\mathfrak p}^1(\NL_{P_\mathfrak p/k}, P_\mathfrak p)$
（《更多代数》引理
\ref{more-algebra-lemma-pseudo-coherence-and-base-change-ext}）。
给定素理想 $\mathfrak p\subset P$，环同态 $k\to P$ 在
$\mathfrak p$ 处光滑，当且仅当 $(\NL_{P/k})_\mathfrak p$ 拟同构于
置于次数 $0$ 的有限投射模（这直接来自光滑环同态的定义，也可由更强的
《代数》引理 \ref{algebra-lemma-smooth-at-point} 得到）。

\medskip\noindent
假设 $P$ 在 $k$ 上除有限多个素理想外均光滑。由《代数》引理
\ref{algebra-lemma-finite-type-algebra-finite-nr-primes} 以及这些“坏”素理想
构成 $\Spec(P)$ 的闭子集这一事实，它们都是极大理想
$\mathfrak m_1,\ldots,\mathfrak m_n\subset P$。对
$\mathfrak p\not\in\{\mathfrak m_1,\ldots,\mathfrak m_n\}$，上述结果给出
$\text{Ext}^1_P(\NL_{P/k},P)_\mathfrak p=0$。因此
$\text{Ext}^1_P(\NL_{P/k},P)$ 是有限 $P$-模，其支集包含于
$\{\mathfrak m_1,\ldots,\mathfrak m_r\}$。
% P11-E0048：冻结源先编号至 m_n，此处却写成 m_r。
例如由《代数》命题
\ref{algebra-proposition-minimal-primes-associated-primes}，
$\text{Ext}^1_P(\NL_{P/k},P)$ 在 $k$ 上的维数是若干
$\dim_k\kappa(\mathfrak m_i)$ 的有限整数线性组合；由 Hilbert 零点定理
（《代数》定理 \ref{algebra-theorem-nullstellensatz}），它是有限的。
\end{proof}

\noindent
在例 \ref{examples-defos-example-rings} 中，设 $P$ 为有限型 $k$-代数。
$\Deformationcategory_P$ 可由 $\mathcal{C}_\Lambda$ 上函子范畴中的
光滑 pro-可表示群胚给出表示，当且仅当 $\dim(P)=0$。此外，若
$\Spec(P)\to\Spec(k)$ 只有有限多个奇点，则
$\Deformationcategory_P$ 具有通用形式对象。这来自引理
\ref{examples-defos-lemma-rings-RS}、\ref{examples-defos-lemma-finite-type-rings-TI} 以及第
\ref{examples-defos-section-general} 节的一般讨论。

\begin{lemma}
\label{examples-defos-lemma-rings-hull}
在例 \ref{examples-defos-example-rings} 中，假设 $P$ 是有限型 $k$-代数，且
$\Spec(P)\to\Spec(k)$ 除有限多个点外均光滑。再假设 $\Lambda$
是剩余域为 $k$ 的完备局部环（经典情形）。则函子
$$
F : \mathcal{C}_\Lambda \longrightarrow \textit{Sets},\quad
A \longmapsto \Ob(\Deformationcategory_P(A))/\cong
$$
（它取对象的同构类）具有包络。
\end{lemma}

\begin{proof}
这立即来自引理 \ref{examples-defos-lemma-rings-RS}、
\ref{examples-defos-lemma-finite-type-rings-TI}、《形式形变理论》引理
\ref{formal-defos-lemma-RS-implies-S1-S2} 与注
\ref{formal-defos-remark-compose-minimal-into-iso-classes}。
\end{proof}

\begin{lemma}
\label{examples-defos-lemma-localization}
在例 \ref{examples-defos-example-rings} 中，设 $P$ 为 $k$-代数，并设 $S\subset P$
为乘法子集。存在形变范畴之间的自然函子
$$
\Deformationcategory_P \longrightarrow \Deformationcategory_{S^{-1}P}
$$
\end{lemma}

\begin{proof}
给定 $P$ 的形变，可以将其局部化，从而得到该局部化的形变；这是显然的，
我们建议读者跳过此证明。更确切地说，设
$(A,Q)\to(k,P)$ 是 $\mathcal{F}$ 中的态射，即
$\Deformationcategory_P$ 的对象。令 $S_Q\subset Q$ 为 $S$ 的原像。
因此，$(A,S_Q^{-1}Q)\to(k,S^{-1}P)$ 是
$\Deformationcategory_{S^{-1}P}$ 的所需对象。
% P11-E0049：冻结源在此有悬空的 “Then” 紧接 “Hence”。
\end{proof}

\begin{lemma}
\label{examples-defos-lemma-henselization}
在例 \ref{examples-defos-example-rings} 中，设 $P$ 为 $k$-代数，$J\subset P$ 为理想。
以 $(P^h,J^h)$ 表示二元组 $(P,J)$ 的 Hensel 化。存在形变范畴之间的
自然函子
$$
\Deformationcategory_P \longrightarrow \Deformationcategory_{P^h}
$$
\end{lemma}

\begin{proof}
给定 $P$ 的形变，可以取其 Hensel 化，从而得到 Hensel 化的形变；
这是显然的，我们建议读者跳过此证明。更确切地说，设
$(A,Q)\to(k,P)$ 是 $\mathcal{F}$ 中的态射，即
$\Deformationcategory_P$ 的对象。以 $J_Q\subset Q$ 表示 $J$ 在 $Q$
中的原像，并设 $(Q^h,J_Q^h)$ 为二元组 $(Q,J_Q)$ 的 Hensel 化。
回忆，$Q\to Q^h$ 平坦（《更多代数》引理
\ref{more-algebra-lemma-henselization-flat}），故 $Q^h$ 在 $A$ 上平坦。
由《更多代数》引理 \ref{more-algebra-lemma-henselization-integral}，
态射 $Q^h\to P^h$ 诱导同构
$Q^h\otimes_Ak=Q^h\otimes_QP=P^h$。因此，
$(A,Q^h)\to(k,P^h)$ 是 $\Deformationcategory_{P^h}$ 的所需对象。
\end{proof}

\begin{lemma}
\label{examples-defos-lemma-strict-henselization}
在例 \ref{examples-defos-example-rings} 中，设 $P$ 为 $k$-代数。假设 $P$ 是局部环，
并设 $P^{sh}$ 为 $P$ 的一个严格 Hensel 化。存在形变范畴之间的自然函子
$$
\Deformationcategory_P \longrightarrow \Deformationcategory_{P^{sh}}
$$
\end{lemma}

\begin{proof}
给定 $P$ 的形变，可以取其严格 Hensel 化，从而得到严格 Hensel 化的
形变；这是显然的，我们建议读者跳过此证明。更确切地说，设
$(A,Q)\to(k,P)$ 是 $\mathcal{F}$ 中的态射，即
$\Deformationcategory_P$ 的对象。由于满射 $Q\to P$ 的核幂零，
$Q$ 是与 $P$ 具有相同剩余域的局部环。设 $Q^{sh}$ 为 $Q$ 的严格
Hensel 化。回忆，$Q\to Q^{sh}$ 平坦（《更多代数》引理
\ref{more-algebra-lemma-dumb-properties-henselization}），故 $Q^{sh}$
在 $A$ 上平坦。由《代数》引理
\ref{algebra-lemma-quotient-strict-henselization}，态射
$Q^{sh}\to P^{sh}$ 诱导同构
$Q^{sh}\otimes_Ak=Q^{sh}\otimes_QP=P^{sh}$。因此，
$(A,Q^{sh})\to(k,P^{sh})$ 是
$\Deformationcategory_{P^{sh}}$ 的所需对象。
\end{proof}

\begin{lemma}
\label{examples-defos-lemma-completion}
在例 \ref{examples-defos-example-rings} 中，设 $P$ 为 $k$-代数。假设 $P$ 是 Noether 环，
并设 $J\subset P$ 为理想，以 $P^\wedge$ 表示 $J$-进完备化。
存在形变范畴之间的自然函子
$$
\Deformationcategory_P \longrightarrow \Deformationcategory_{P^\wedge}
$$
\end{lemma}

\begin{proof}
给定 $P$ 的形变，可以将其完备化，从而得到完备化的形变；这是显然的，
我们建议读者跳过此证明。更确切地说，设
$(A,Q)\to(k,P)$ 是 $\mathcal{F}$ 中的态射，即
$\Deformationcategory_P$ 的对象。注意 $Q$ 是诺特环：满环同态
$Q\to P$ 的核幂零且有限生成，而 $P$ 是诺特环；应用《代数》引理
\ref{algebra-lemma-completion-Noetherian}。以 $J_Q\subset Q$ 表示 $J$
在 $Q$ 中的原像，并设 $Q^\wedge$ 为 $Q$ 的 $J_Q$-进完备化。
回忆，$Q\to Q^\wedge$ 平坦（《代数》引理
\ref{algebra-lemma-completion-flat}），故 $Q^\wedge$ 在 $A$ 上平坦。
例如由《代数》引理 \ref{algebra-lemma-completion-tensor}，诱导态射
$Q^\wedge\to P^\wedge$ 诱导同构
$Q^\wedge\otimes_Ak=Q^\wedge\otimes_QP=P^\wedge$。因此，
$(A,Q^\wedge)\to(k,P^\wedge)$ 是
$\Deformationcategory_{P^\wedge}$ 的所需对象。
\end{proof}

\begin{lemma}
\label{examples-defos-lemma-power-series-rings-TI}
在引理 \ref{examples-defos-lemma-rings-TI} 中，若
$P=k[[x_1,\ldots,x_n]]/(f)$，其中非零的
$f\in(x_1,\ldots,x_n)^2$，则
\begin{enumerate}
\item $\text{Inf}(\Deformationcategory_P)$ 有限维，当且仅当 $n=1$；
\item 若
$$
\sqrt{(f, \partial f/\partial x_1, \ldots,  \partial f/\partial x_n)} =
(x_1, \ldots, x_n)
$$
，则 $T\Deformationcategory_P$ 有限维。
\end{enumerate}
\end{lemma}

\begin{proof}
证明 (1)。考虑 $k[[x_1,\ldots,x_n]]$ 在 $k$ 上的导子
$\partial/\partial x_i$。记 $f_i=\partial f/\partial x_i$。导子
$$
\theta = \sum h_i \partial/\partial x_i
$$
是 $k[[x_1,\ldots,x_n]]$ 的导子；它在
$P=k[[x_1,\ldots,x_n]]/(f)$ 上诱导导子，当且仅当
$\sum h_if_i\in(f)$。此外，所诱导的 $P$ 上导子为零，当且仅当对
$i=1,\ldots,n$ 均有 $h_i\in(f)$。因此得到
$$
\Ker((f_1, \ldots, f_n) : P^{\oplus n} \longrightarrow P) \subset
\text{Der}_k(P, P)
$$
左端仅在 $n=1$ 时才是有限维 $k$-向量空间；证明从略。右端在 $n=1$
时有限维这一点也留给读者验证。由此得到 (1)。

\medskip\noindent
证明 (2)。令 $Q$ 为 $P$ 在 $k[\epsilon]$ 上的平坦形变，如引理
\ref{examples-defos-lemma-rings-TI} 的证明中那样。选取 $x_i$ 在 $P$ 中之像的提升
$q_i\in Q$。则 $Q$ 是完备局部环，其极大理想由
$q_1,\ldots,q_n$ 与 $\epsilon$ 生成（略去一个简短论证）。因此得到满射
$$
k[\epsilon][[x_1, \ldots, x_n]] \longrightarrow Q,\quad
x_i \longmapsto q_i
$$
选取映到 $Q$ 中零的元素
$f+\epsilon g\in k[\epsilon][[x_1,\ldots,x_n]]$。注意，$g$ 模 $(f)$
良定义。由于 $Q$ 在 $k[\epsilon]$ 上平坦，得到
$$
Q = k[\epsilon][[x_1, \ldots, x_n]]/(f + \epsilon g)
$$
最后，改变 $q_i$ 的选择，相当于对某些
$h_i\in k[[x_1,\ldots,x_n]]$，把坐标 $x_i$ 改为
$x_i+\epsilon h_i$。于是 $f+\epsilon g$ 变为
$f+\epsilon(g+\sum h_if_i)$，其中 $f_i=\partial f/\partial x_i$。
% P11-E0051：冻结源此句写作 “if we changing”。
因此，形变 $Q$ 的同构类由下列空间中的一个元素确定：
$$
k[[x_1, \ldots, x_n]]/
(f, \partial f/\partial x_1, \ldots,  \partial f/\partial x_n)
$$
它在 $k$ 上有限维，当且仅当其支集是
$k[[x_1,\ldots,x_n]]$ 的闭点；这又当且仅当
$\sqrt{(f, \partial f/\partial x_1, \ldots,  \partial f/\partial x_n)} =
(x_1, \ldots, x_n)$.
\end{proof}






\section{概形}
\label{examples-defos-section-schemes}

\noindent
概形的形变理论。

\begin{example}[概形]
\label{examples-defos-example-schemes}
如下定义范畴 $\mathcal{F}$：
\begin{enumerate}
\item 对象是二元组 $(A,X)$，其中 $A$ 是 $\mathcal{C}_\Lambda$ 的对象，
$X$ 是在 $A$ 上平坦的概形；
\item 态射 $(f,g):(B,Y)\to(A,X)$ 由 $\mathcal{C}_\Lambda$ 中的态射
$f:B\to A$ 与态射 $g:X\to Y$ 组成，使概形的交换图
$$
\xymatrix{
X \ar[r]_g \ar[d] & Y \ar[d] \\
\Spec(A) \ar[r]^f & \Spec(B)
}
$$
为笛卡儿图。
\end{enumerate}
函子 $p:\mathcal{F}\to\mathcal{C}_\Lambda$ 把 $(A,X)$ 送到 $A$，
把 $(f,g)$ 送到 $f$。显然，$p$ 以群胚为纤维余纤维化。
给定 $k$ 上概形 $X$，令 $x_0=(k,X)$ 为 $\mathcal{F}(k)$ 的相应对象。记
$$
\Deformationcategory_X = \mathcal{F}_{x_0}
$$
\end{example}

\begin{lemma}
\label{examples-defos-lemma-schemes-RS}
例 \ref{examples-defos-example-schemes} 满足 Rim--Schlessinger 条件 (RS)。特别地，
对任意 $k$ 上概形 $X$，$\Deformationcategory_X$ 是形变范畴。
\end{lemma}

\begin{proof}
设 $A_1\to A$ 与 $A_2\to A$ 是 $\mathcal{C}_\Lambda$ 的态射，
并假设 $A_2\to A$ 满射。由《形式形变理论》引理
\ref{formal-defos-lemma-RS-2-categorical}，只须证明函子
$\mathcal{F}(A_1 \times_A A_2) \to
\mathcal{F}(A_1) \times_{\mathcal{F}(A)} \mathcal{F}(A_2)$
是范畴等价。注意
$$
\xymatrix{
\Spec(A) \ar[r] \ar[d] & \Spec(A_2) \ar[d] \\
\Spec(A_1) \ar[r] &
\Spec(A_1 \times_A A_2)
}
$$
是《更多态射》引理
\ref{more-morphisms-lemma-pushout-along-thickening} 中的推出图。
因此，本引理是《更多态射》引理
\ref{more-morphisms-lemma-equivalence-categories-schemes-over-pushout-flat}
的特殊情形。
\end{proof}

\begin{lemma}
\label{examples-defos-lemma-schemes-TI}
在例 \ref{examples-defos-example-schemes} 中，设 $X$ 为 $k$ 上概形。则
$$
\text{Inf}(\Deformationcategory_X) =
\text{Ext}^0_{\mathcal{O}_X}(\NL_{X/k}, \mathcal{O}_X) =
\Hom_{\mathcal{O}_X}(\Omega_{X/k}, \mathcal{O}_X) =
\text{Der}_k(\mathcal{O}_X, \mathcal{O}_X)
$$
且
$$
T\Deformationcategory_X =
\text{Ext}^1_{\mathcal{O}_X}(\NL_{X/k}, \mathcal{O}_X)
$$
\end{lemma}

\begin{proof}
回忆，$\text{Inf}(\Deformationcategory_X)$ 是 $X$ 在 $k[\epsilon]$
上的平凡形变
$X'=X\times_{\Spec(k)}\Spec(k[\epsilon])$ 中模 $\epsilon$ 等于恒等的
自同构之集。由《形变理论》引理 \ref{defos-lemma-deform}，它等于
$\text{Ext}^0_{\mathcal{O}_X}(\NL_{X/k},\mathcal{O}_X)$。等式
$\text{Ext}^0_{\mathcal{O}_X}(\NL_{X/k},\mathcal{O}_X)=
\Hom_{\mathcal{O}_X}(\Omega_{X/k},\mathcal{O}_X)$ 来自《更多态射》引理
\ref{more-morphisms-lemma-netherlander-quasi-coherent}。等式
$\Hom_{\mathcal{O}_X}(\Omega_{X/k}, \mathcal{O}_X) =
\text{Der}_k(\mathcal{O}_X, \mathcal{O}_X)$
来自《态射》引理 \ref{morphisms-lemma-universal-derivation-universal}。

\medskip\noindent
回忆，$T_{x_0}\Deformationcategory_X$ 是 $X$ 在 $k[\epsilon]$ 上的
平坦形变 $X'$ 的同构类之集；更确切地说，它是
$\Deformationcategory_X(k[\epsilon])$ 的同构类之集。因此，引理的
第二个断言来自《形变理论》引理 \ref{defos-lemma-deform}。
\end{proof}

\begin{lemma}
\label{examples-defos-lemma-proper-schemes-TI}
在引理 \ref{examples-defos-lemma-schemes-TI} 中，若 $X$ 在 $k$ 上固有，则
$\text{Inf}(\Deformationcategory_X)$ 与 $T\Deformationcategory_X$
均为有限维。
\end{lemma}

\begin{proof}
由该引理，须证明
$\Ext^1_{\mathcal{O}_X}(\NL_{X/k},\mathcal{O}_X)$ 与
$\Ext^0_{\mathcal{O}_X}(\NL_{X/k},\mathcal{O}_X)$ 有限维。
由《更多态射》引理 \ref{more-morphisms-lemma-netherlander-fp}
以及 $X$ 为诺特概形这一事实，$\NL_{X/k}$ 的上同调层均凝聚，
且除次数 $0$ 与 $-1$ 外均为零。由《概形的导出范畴》引理
\ref{perfect-lemma-ext-finite}，所示 $\Ext$ 群都是有限维
$k$-向量空间，证明完毕。
\end{proof}

\noindent
在例 \ref{examples-defos-example-schemes} 中，若 $X$ 是 $k$ 上固有概形，则
$\Deformationcategory_X$ 可由 $\mathcal{C}_\Lambda$ 上函子范畴中的
光滑 pro-可表示群胚给出表示，因而具有一个（极小）通用形式对象。
这来自引理 \ref{examples-defos-lemma-schemes-RS}、\ref{examples-defos-lemma-proper-schemes-TI}
以及第 \ref{examples-defos-section-general} 节的一般讨论。

\begin{lemma}
\label{examples-defos-lemma-schemes-hull}
在例 \ref{examples-defos-example-schemes} 中，假设 $X$ 是固有 $k$-概形，并假设
$\Lambda$ 是剩余域为 $k$ 的完备局部环（经典情形）。则函子
$$
F : \mathcal{C}_\Lambda \longrightarrow \textit{Sets},\quad
A \longmapsto \Ob(\Deformationcategory_X(A))/\cong
$$
（它取对象的同构类）具有包络。若
$\text{Der}_k(\mathcal{O}_X,\mathcal{O}_X)=0$，则 $F$ 是
pro-可表示的。
\end{lemma}

\begin{proof}
包络的存在性立即来自引理 \ref{examples-defos-lemma-schemes-RS}、
\ref{examples-defos-lemma-proper-schemes-TI}、《形式形变理论》引理
\ref{formal-defos-lemma-RS-implies-S1-S2} 与注
\ref{formal-defos-remark-compose-minimal-into-iso-classes}。

\medskip\noindent
假设 $\text{Der}_k(\mathcal{O}_X,\mathcal{O}_X)=0$。由《形式形变理论》
引理 \ref{formal-defos-lemma-infdef-trivial}，
$\Deformationcategory_X$ 与 $F$ 等价。因此，$F$ 是具有有限维切空间的
形变函子（因为 $\Deformationcategory_X$ 是形变范畴），可以应用
《形式形变理论》定理
\ref{formal-defos-theorem-Schlessinger-prorepresentability}。
\end{proof}

\begin{lemma}
\label{examples-defos-lemma-open}
在例 \ref{examples-defos-example-schemes} 中，设 $X$ 为 $k$ 上概形，$U\subset X$
为开子概形。存在形变范畴之间的自然函子
$$
\Deformationcategory_X \longrightarrow \Deformationcategory_U
$$
\end{lemma}

\begin{proof}
给定 $X$ 的形变，可以取其中相应的开子概形，从而得到 $U$ 的形变。
细节从略。
\end{proof}

\begin{lemma}
\label{examples-defos-lemma-affine}
在例 \ref{examples-defos-example-schemes} 中，设 $X=\Spec(P)$ 为仿射 $k$-概形。
令 $\Deformationcategory_P$ 如例 \ref{examples-defos-example-rings}。存在形变范畴的
自然等价
$$
\Deformationcategory_X \longrightarrow \Deformationcategory_P
$$
\end{lemma}

\begin{proof}
该函子把 $(A,Y)$ 送到 $\Gamma(Y,\mathcal{O}_Y)$。这是可行的，因为由
《更多态射》引理 \ref{more-morphisms-lemma-thickening-affine-scheme}，
$X$ 的任意形变都是仿射的。
\end{proof}

\begin{lemma}
\label{examples-defos-lemma-local-ring}
在例 \ref{examples-defos-example-schemes} 中，设 $X$ 为 $k$ 上概形，并设 $p\in X$
为一点。令 $\Deformationcategory_{\mathcal{O}_{X,p}}$ 如例
\ref{examples-defos-example-rings}。存在形变范畴之间的自然函子
$$
\Deformationcategory_X
\longrightarrow
\Deformationcategory_{\mathcal{O}_{X, p}}
$$
\end{lemma}

\begin{proof}
选取包含 $p$ 的仿射开子概形 $U=\Spec(P)\subset X$。则
$\mathcal{O}_{X,p}$ 是 $P$ 的一个局部化。复合引理
\ref{examples-defos-lemma-open}、\ref{examples-defos-lemma-affine} 与 \ref{examples-defos-lemma-localization}
中的函子即可。
\end{proof}

\begin{situation}
\label{examples-defos-situation-glueing}
设 $\Lambda\to k$ 如第 \ref{examples-defos-section-general} 节。设 $X$ 为 $k$ 上概形，
具有仿射开覆盖 $X=U_1\cup U_2$，且
$U_{12}=U_1\cap U_2$ 也仿射。写成
$U_1=\Spec(P_1)$、$U_2=\Spec(P_2)$ 与 $U_{12}=\Spec(P_{12})$。
令 $\Deformationcategory_X$、$\Deformationcategory_{U_1}$、
$\Deformationcategory_{U_2}$ 与 $\Deformationcategory_{U_{12}}$
如例 \ref{examples-defos-example-schemes}，并令
$\Deformationcategory_{P_1}$、$\Deformationcategory_{P_2}$ 与
$\Deformationcategory_{P_{12}}$ 如例 \ref{examples-defos-example-rings}。
\end{situation}

\begin{lemma}
\label{examples-defos-lemma-glueing}
在情形 \ref{examples-defos-situation-glueing} 中，存在形变范畴的等价
$$
\Deformationcategory_X =
\Deformationcategory_{P_1}
\times_{\Deformationcategory_{P_{12}}}
\Deformationcategory_{P_2}
$$
，参见例 \ref{examples-defos-example-schemes} 与 \ref{examples-defos-example-rings}。
\end{lemma}

\begin{proof}
只须证明引理 \ref{examples-defos-lemma-open} 的函子定义等价
$$
\Deformationcategory_X \longrightarrow
\Deformationcategory_{U_1}
\times_{\Deformationcategory_{U_{12}}}
\Deformationcategory_{U_2}
$$
，因为这样便可应用引理 \ref{examples-defos-lemma-affine} 转化为环的情形。
为此构造一个拟逆。以
$F_i:\Deformationcategory_{U_i}\to\Deformationcategory_{U_{12}}$
表示引理 \ref{examples-defos-lemma-open} 的函子。右端的对象由
$\mathcal{C}_\Lambda$ 中的 $A$、对象
$(A,V_1)\to(k,U_1)$ 与 $(A,V_2)\to(k,U_2)$，以及态射
$$
g : F_1(A, V_1) \to F_2(A, V_2)
$$
给出。这里 $F_i(A,V_i)=(A,V_{i,3-i})$，其中
$V_{i,3-i}\subset V_i$ 是基变换到 $k$ 后成为
$U_{12}\subset U_i$ 的开子概形。态射 $g$ 定义 $A$ 上概形的同构
$V_{1,2}\to V_{2,1}$，且它与 $k$ 上的
$\text{id}:U_{12}\to U_{12}$ 相容。因此
$(\{1,2\},V_i,V_{i,3-i},g,g^{-1})$ 是《概形》第
\ref{schemes-section-glueing-schemes} 节中的粘合数据。设 $Y$ 为所得粘合，
参见《概形》引理 \ref{schemes-lemma-glue}。则 $Y$ 是 $A$ 上概形，
上述相容性给出典范同构
$Y\times_{\Spec(A)}\Spec(k)=X$。因此
$(A,Y)\to(k,X)$ 是 $\Deformationcategory_X$ 的对象。
关于此构造是函子且为给定函子的拟逆这一验证从略。
\end{proof}






\section{概形态射}
\label{examples-defos-section-schemes-morphisms}

\noindent
本节讨论概形态射的形变理论。
当然，这不过是概形图表形变的一个例子。

\begin{example}[概形态射]
\label{examples-defos-example-schemes-morphisms}
定义范畴 $\mathcal{F}$ 如下：
\begin{enumerate}
\item 对象是一个偶对 $(A, X \to Y)$，其中 $A$ 是
$\mathcal{C}_\Lambda$ 的对象，而 $X \to Y$ 是 $A$ 上的概形态射，
且 $X$ 与 $Y$ 都在 $A$ 上平坦；
\item 态射 $(f, g, h) : (A', X' \to Y') \to (A, X \to Y)$
由 $\mathcal{C}_\Lambda$ 中的态射 $f : A' \to A$，以及概形态射
$g : X \to X'$ 与 $h : Y \to Y'$ 组成，并使得
$$
\xymatrix{
X \ar[r]_g \ar[d] & X' \ar[d] \\
Y \ar[r]_h \ar[d] & Y' \ar[d] \\
\Spec(A) \ar[r]^f & \Spec(A')
}
$$
是概形的交换图，且其中两个方块都是笛卡儿的。
\end{enumerate}
函子 $p : \mathcal{F} \to \mathcal{C}_\Lambda$ 把 $(A, X \to Y)$
送到 $A$，把 $(f, g, h)$ 送到 $f$。显然，$p$ 以群胚为纤维余纤维化。
给定 $k$ 上的概形态射 $X \to Y$，令 $x_0 = (k, X \to Y)$
为 $\mathcal{F}(k)$ 中相应的对象。置
$$
\Deformationcategory_{X \to Y} = \mathcal{F}_{x_0}
$$
\end{example}

\begin{lemma}
\label{examples-defos-lemma-schemes-morphisms-RS}
例 \ref{examples-defos-example-schemes-morphisms}
满足 Rim--Schlessinger 条件 (RS)。
特别地，对 $k$ 上任意概形态射 $X \to Y$，
$\Deformationcategory_{X \to Y}$ 都是形变范畴。
\end{lemma}

\begin{proof}
设 $A_1 \to A$ 与 $A_2 \to A$ 是 $\mathcal{C}_\Lambda$ 中的态射，
并假设 $A_2 \to A$ 是满射。根据形式形变理论引理
\ref{formal-defos-lemma-RS-2-categorical}
，只需证明函子
$\mathcal{F}(A_1 \times_A A_2) \to
\mathcal{F}(A_1) \times_{\mathcal{F}(A)} \mathcal{F}(A_2)$
是范畴等价。注意到
$$
\xymatrix{
\Spec(A) \ar[r] \ar[d] & \Spec(A_2) \ar[d] \\
\Spec(A_1) \ar[r] &
\Spec(A_1 \times_A A_2)
}
$$
是态射进阶引理
\ref{more-morphisms-lemma-pushout-along-thickening} 中的推出图。
因而本引理由态射进阶引理
\ref{more-morphisms-lemma-equivalence-categories-schemes-over-pushout-flat}
立即得到：该引理把在 $A_1 \times_A A_2$ 上平坦的概形所成范畴描述为
在 $A_1$ 上平坦的概形所成范畴与在 $A_2$ 上平坦的概形所成范畴，
在 $A$ 上平坦的概形所成范畴之上的纤维积。
\end{proof}

\begin{lemma}
\label{examples-defos-lemma-schemes-morphisms-TI}
% P11-E0052：冻结源在本节两处引用 example-schemes，而所述对象属于 example-schemes-morphisms。
在例 \ref{examples-defos-example-schemes} 中，设 $f : X \to Y$ 是 $k$ 上的概形态射。
有如下典范的 $k$-向量空间正合列：
$$
\xymatrix{
0 \ar[r] &
\text{Inf}(\Deformationcategory_{X \to Y}) \ar[r] &
\text{Inf}(\Deformationcategory_X \times \Deformationcategory_Y) \ar[r] &
\text{Der}_k(\mathcal{O}_Y, f_*\mathcal{O}_X) \ar[lld] \\
& T\Deformationcategory_{X \to Y} \ar[r] &
T(\Deformationcategory_X \times \Deformationcategory_Y) \ar[r] &
\text{Ext}^1_{\mathcal{O}_X}(Lf^*\NL_{Y/k}, \mathcal{O}_X)
}
$$
\end{lemma}

\begin{proof}
形变范畴间的显然映射
$\Deformationcategory_{X \to Y} \to
\Deformationcategory_X \times \Deformationcategory_Y$
给出了引理中正合列的两个箭头。回忆，
$\text{Inf}(\Deformationcategory_{X \to Y})$
是平凡形变
$$
f' :  X' = X \times_{\Spec(k)} \Spec(k[\epsilon])
\xrightarrow{f \times \text{id}}
Y' = Y \times_{\Spec(k)} \Spec(k[\epsilon])
$$
在模 $\epsilon$ 意义下等于恒等的自同构所成的集合；这里该平凡形变是
$X \to Y$ 到 $k[\epsilon]$ 上的形变。这显然等同于偶对
$(\alpha, \beta) \in
\text{Inf}(\Deformationcategory_X \times \Deformationcategory_Y)$
，其中二者分别是 $X$ 与 $Y$ 的无穷小自同构，
且与 $f'$ 相容；也就是说 $f' \circ \alpha = \beta \circ f'$。
由形变理论引理 \ref{defos-lemma-huge-diagram-ringed-spaces}，
对任意偶对 $(\alpha, \beta)$，态射 $f' : X' \to Y'$ 与态射
$\beta^{-1} \circ f' \circ \alpha : X' \to Y'$ 之差定义了
$$
\text{Der}_k(\mathcal{O}_Y, f_*\mathcal{O}_X) =
\Hom_{\mathcal{O}_Y}(\Omega_{Y/k}, f_*\mathcal{O}_X)
$$
中的一个元素。这里的等号由态射进阶引理
\ref{more-morphisms-lemma-netherlander-quasi-coherent} 给出。
这定义了最上方最后一个水平箭头，并证明了前两处的正合性。对于映射
$$
\text{Der}_k(\mathcal{O}_Y, f_*\mathcal{O}_X)
\to
T\Deformationcategory_{X \to Y}
$$
，利用形变理论引理 \ref{defos-lemma-huge-diagram-ringed-spaces}，
我们把源中的元素解释为 $\Spec(k[\epsilon])$ 上的态射
$f_\epsilon : X' \to Y'$，并要求其模 $\epsilon$ 等于 $f$。
我们把 $f_\epsilon$ 送到
$T\Deformationcategory_{X \to Y}$ 中
$(f_\epsilon : X' \to Y')$ 的同构类。注意，
$(f_\epsilon : X' \to Y')$ 与平凡形变
$(f' : X' \to Y')$ 同构，当且仅当
% P11-E0053：冻结源下一公式中的 f 按类型应为 f'。
对某个偶对 $(\alpha, \beta)$ 有
$f_\epsilon  = \beta^{-1} \circ f \circ \alpha$；
这说明第三处正合。显然，若某个一阶形变
$(f_\epsilon : X_\epsilon \to Y_\epsilon)$
在 $T(\Deformationcategory_X \times \Deformationcategory_Y)$ 中映为零，
则可选取同构 $X' \to X_\epsilon$ 与 $Y' \to Y_\epsilon$，
从而得到它位于左下方向箭头的像中。因此第四处正合。
最后，给定 $X$、$Y$ 的两个一阶形变 $X_\epsilon$、$Y_\epsilon$，
存在障碍
$$
ob(X_\epsilon, Y_\epsilon) \in
\text{Ext}^1_{\mathcal{O}_X}(Lf^*\NL_{Y/k}, \mathcal{O}_X)
$$
；该障碍消失，当且仅当 $f : X \to Y$ 可提升为
$X_\epsilon \to Y_\epsilon$；参见形变理论引理
\ref{defos-lemma-huge-diagram-ringed-spaces}。证明完毕。
\end{proof}

\begin{lemma}
\label{examples-defos-lemma-proper-schemes-morphisms-TI}
在引理 \ref{examples-defos-lemma-schemes-morphisms-TI} 中，若 $X$ 与 $Y$
都在 $k$ 上固有，则
$\text{Inf}(\Deformationcategory_{X \to Y})$ 与
$T\Deformationcategory_{X \to Y}$ 都是有限维的。
\end{lemma}

\begin{proof}
从略。提示：仿照引理 \ref{examples-defos-lemma-proper-schemes-TI} 论证，
并使用本引理的正合列。
\end{proof}

\noindent
在例 \ref{examples-defos-example-schemes-morphisms} 中，若 $X \to Y$ 是
$k$ 上固有概形之间的态射，则
$\Deformationcategory_{X \to Y}$ 可由 $\mathcal{C}_\Lambda$
上函子中的一个光滑 pro-可表示群胚给出呈示，因而更有一个
（极小）通用形式对象。这由引理 \ref{examples-defos-lemma-schemes-morphisms-RS}、
\ref{examples-defos-lemma-proper-schemes-morphisms-TI} 以及
第 \ref{examples-defos-section-general} 节的一般讨论推出。

\begin{lemma}
\label{examples-defos-lemma-schemes-morphisms-hull}
在例 \ref{examples-defos-example-schemes-morphisms} 中，假设 $X \to Y$
是固有 $k$-概形之间的态射。再假设 $\Lambda$ 是剩余域为 $k$
的完备局部环（经典情形）。则函子
$$
F : \mathcal{C}_\Lambda \longrightarrow \textit{Sets},\quad
A \longmapsto \Ob(\Deformationcategory_{X \to Y}(A))/\cong
$$
把对象映到其同构类；它有包络。若
$\text{Der}_k(\mathcal{O}_X, \mathcal{O}_X) =
\text{Der}_k(\mathcal{O}_Y, \mathcal{O}_Y) = 0$，
则 $F$ 是 pro-可表示的。
\end{lemma}

\begin{proof}
包络的存在性立即由引理 \ref{examples-defos-lemma-schemes-morphisms-RS}、
\ref{examples-defos-lemma-proper-schemes-morphisms-TI}、形式形变理论引理
\ref{formal-defos-lemma-RS-implies-S1-S2} 以及注
\ref{formal-defos-remark-compose-minimal-into-iso-classes} 推出。

\medskip\noindent
假设 $\text{Der}_k(\mathcal{O}_X, \mathcal{O}_X) =
\text{Der}_k(\mathcal{O}_Y, \mathcal{O}_Y) = 0$。
引理 \ref{examples-defos-lemma-schemes-morphisms-TI} 的正合列结合引理
\ref{examples-defos-lemma-schemes-TI} 表明
$\text{Inf}(\Deformationcategory_{X \to Y}) = 0$。
于是由形式形变理论引理 \ref{formal-defos-lemma-infdef-trivial}，
$\Deformationcategory_{X \to Y}$ 与 $F$ 等价。因此 $F$ 是具有
有限维切空间的形变函子（因为 $\Deformationcategory_{X \to Y}$
是形变范畴），我们可以应用形式形变理论定理
\ref{formal-defos-theorem-Schlessinger-prorepresentability}。
\end{proof}

\begin{lemma}
\label{examples-defos-lemma-schemes-morphisms-smooth-to-base}
\begin{reference}
这在 \cite[Section 5.3]{Ravi-Murphys-Law} 与
\cite[Theorem 3.3]{Ran-deformations} 中有所讨论。
\end{reference}
在例 \ref{examples-defos-example-schemes} 中，设 $f : X \to Y$ 是 $k$ 上的概形态射。
若 $f_*\mathcal{O}_X = \mathcal{O}_Y$ 且 $R^1f_*\mathcal{O}_X = 0$，
则形变范畴的态射
$$
\Deformationcategory_{X \to Y} \to \Deformationcategory_X
$$
是等价。
\end{lemma}

\begin{proof}
我们构造引理中遗忘函子的一个拟逆。设 $(A,U)$ 是
$\Deformationcategory_X$ 的对象。给定映射 $X \to U$ 是有限阶加厚，
可据此把 $U$ 与 $X$ 的底拓扑空间等同起来；参见态射进阶第
\ref{more-morphisms-section-thickenings} 节。因此我们可以并确实把
$\mathcal{O}_U$ 看作 $X$ 上的 $A$-代数层；此外，
$U \to \Spec(A)$ 平坦意味着 $\mathcal{O}_U$ 作为 $A$-模层平坦。
特别地，有滤过
$$
0 = \mathfrak m_A^n\mathcal{O}_U \subset
\mathfrak m_A^{n - 1}\mathcal{O}_U \subset \ldots \subset
\mathfrak m_A^2\mathcal{O}_U \subset
\mathfrak m_A\mathcal{O}_U \subset \mathcal{O}_U
$$
；由平坦性，其相继商等于
$\mathcal{O}_X \otimes_k \mathfrak m_A^i/\mathfrak m_A^{i + 1}$
；参见态射进阶引理 \ref{more-morphisms-lemma-deform}，
或更一般的形变理论引理 \ref{defos-lemma-deform-module}。置
$$
\mathcal{O}_V = f_*\mathcal{O}_U
$$
，把它视作 $Y$ 上的 $A$-代数层。由于
$R^1f_*\mathcal{O}_X = 0$，由上述描述可知
$R^1f_*(\mathfrak m_A^i\mathcal{O}_U/\mathfrak m_A^{i + 1}\mathcal{O}_U) = 0$
对所有 $i$ 成立。这意味着序列
$$
0 \to
(f_*\mathcal{O}_X) \otimes_k \mathfrak m_A^i/\mathfrak m_A^{i + 1} \to
f_*(\mathcal{O}_U/\mathfrak m_A^{i + 1}\mathcal{O}_U) \to
f_*(\mathcal{O}_U/\mathfrak m_A^i\mathcal{O}_U) \to 0
$$
对所有 $i$ 都正合。反向应用上述文献（并使用归纳法），可知
$\mathcal{O}_V$ 是平坦的 $A$-代数层，且
$\mathcal{O}_V/\mathfrak m_A\mathcal{O}_V = \mathcal{O}_Y$。
应用态射进阶引理
\ref{more-morphisms-lemma-first-order-thickening}
可知 $(Y,\mathcal{O}_V)$ 是概形，记为 $V$。
等式 $\mathcal{O}_V = f_*\mathcal{O}_U$ 定义了带环空间的态射
$U \to V$，容易看出它是概形态射。结合已经证明的平坦性，
这就完成了证明。
\end{proof}







\section{代数空间}
\label{examples-defos-section-algebraic-spaces}

\noindent
本节讨论代数空间的形变理论。

\begin{example}[代数空间]
\label{examples-defos-example-spaces}
定义范畴 $\mathcal{F}$ 如下：
\begin{enumerate}
\item 对象是偶对 $(A,X)$，其中 $A$ 是 $\mathcal{C}_\Lambda$
的对象，$X$ 是在 $A$ 上平坦的代数空间；
\item 态射 $(f,g):(B,Y)\to(A,X)$ 由
$\mathcal{C}_\Lambda$ 中的态射 $f:B\to A$ 以及
$\Lambda$ 上代数空间的态射 $g:X\to Y$ 组成，并使得
$$
\xymatrix{
X \ar[r]_g \ar[d] & Y \ar[d] \\
\Spec(A) \ar[r]^f & \Spec(B)
}
$$
是代数空间的笛卡儿交换图。
\end{enumerate}
函子 $p:\mathcal{F}\to\mathcal{C}_\Lambda$ 把 $(A,X)$ 送到 $A$，
把 $(f,g)$ 送到 $f$。显然，$p$ 以群胚为纤维余纤维化。
给定 $k$ 上的代数空间 $X$，令 $x_0=(k,X)$ 为
$\mathcal{F}(k)$ 中相应的对象。置
$$
\Deformationcategory_X = \mathcal{F}_{x_0}
$$
\end{example}

\begin{lemma}
\label{examples-defos-lemma-spaces-RS}
例 \ref{examples-defos-example-spaces} 满足 Rim--Schlessinger 条件 (RS)。
特别地，对 $k$ 上任意代数空间 $X$，
$\Deformationcategory_X$ 都是形变范畴。
\end{lemma}

\begin{proof}
设 $A_1\to A$ 与 $A_2\to A$ 是 $\mathcal{C}_\Lambda$ 中的态射，
并假设 $A_2\to A$ 为满射。根据形式形变理论引理
\ref{formal-defos-lemma-RS-2-categorical}
，只需证明函子
$\mathcal{F}(A_1 \times_A A_2) \to
\mathcal{F}(A_1) \times_{\mathcal{F}(A)} \mathcal{F}(A_2)$
是范畴等价。注意到
$$
\xymatrix{
\Spec(A) \ar[r] \ar[d] & \Spec(A_2) \ar[d] \\
\Spec(A_1) \ar[r] &
\Spec(A_1 \times_A A_2)
}
$$
是空间的推出引理
\ref{spaces-pushouts-lemma-pushout-along-thickening} 中的推出图。
因而本引理是空间的推出引理
\ref{spaces-pushouts-lemma-equivalence-categories-spaces-pushout-flat}
的特殊情形。
\end{proof}

\begin{lemma}
\label{examples-defos-lemma-spaces-TI}
在例 \ref{examples-defos-example-spaces} 中，设 $X$ 是 $k$ 上的代数空间。则
$$
\text{Inf}(\Deformationcategory_X) =
\text{Ext}^0_{\mathcal{O}_X}(\NL_{X/k}, \mathcal{O}_X) =
\Hom_{\mathcal{O}_X}(\Omega_{X/k}, \mathcal{O}_X) =
\text{Der}_k(\mathcal{O}_X, \mathcal{O}_X)
$$
以及
$$
T\Deformationcategory_X =
\text{Ext}^1_{\mathcal{O}_X}(\NL_{X/k}, \mathcal{O}_X)
$$
\end{lemma}

\begin{proof}
回忆，$\text{Inf}(\Deformationcategory_X)$ 是 $X$ 到 $k[\epsilon]$
上的平凡形变 $X'=X\times_{\Spec(k)}\Spec(k[\epsilon])$
在模 $\epsilon$ 意义下等于恒等的自同构所成的集合。
由形变理论引理 \ref{defos-lemma-deform-spaces}，它等于
$\text{Ext}^0_{\mathcal{O}_X}(\NL_{X/k},\mathcal{O}_X)$。
等式 $\text{Ext}^0_{\mathcal{O}_X}(\NL_{X/k}, \mathcal{O}_X) =
\Hom_{\mathcal{O}_X}(\Omega_{X/k}, \mathcal{O}_X)$
由空间态射进阶引理
\ref{spaces-more-morphisms-lemma-netherlander-quasi-coherent}
得到。等式
$\Hom_{\mathcal{O}_X}(\Omega_{X/k}, \mathcal{O}_X) =
\text{Der}_k(\mathcal{O}_X, \mathcal{O}_X)$
由空间态射进阶定义
\ref{spaces-more-morphisms-definition-sheaf-differentials}
与位点上的模定义 \ref{sites-modules-definition-module-differentials} 得到。

\medskip\noindent
回忆，$T_{x_0}\Deformationcategory_X$ 是 $X$ 到 $k[\epsilon]$
上的平坦形变 $X'$ 的同构类所成的集合；更准确地说，
它是 $\Deformationcategory_X(k[\epsilon])$ 中同构类所成的集合。
因此引理的第二个断言由形变理论引理
\ref{defos-lemma-deform-spaces} 得到。
\end{proof}

\begin{lemma}
\label{examples-defos-lemma-proper-spaces-TI}
在引理 \ref{examples-defos-lemma-spaces-TI} 中，若 $X$ 在 $k$ 上固有，则
$\text{Inf}(\Deformationcategory_X)$ 与 $T\Deformationcategory_X$
都是有限维的。
\end{lemma}

\begin{proof}
由该引理，我们须证明
$\Ext^1_{\mathcal{O}_X}(\NL_{X/k}, \mathcal{O}_X)$ 与
$\Ext^0_{\mathcal{O}_X}(\NL_{X/k}, \mathcal{O}_X)$ 是有限维的。
由空间态射进阶引理
\ref{spaces-more-morphisms-lemma-netherlander-fp}
以及 $X$ 是 Noether 的这一事实可知，$\NL_{X/k}$ 的凝聚上同调层
除次数 $0$ 与 $-1$ 外均为零。由空间的导出范畴引理
\ref{spaces-perfect-lemma-ext-finite}，上述 $\Ext$ 群都是有限维
$k$-向量空间，证明完毕。
\end{proof}

\noindent
在例 \ref{examples-defos-example-spaces} 中，若 $X$ 是 $k$ 上的固有代数空间，
则 $\Deformationcategory_X$ 可由 $\mathcal{C}_\Lambda$
上函子中的一个光滑 pro-可表示群胚给出呈示，因而更有一个
（极小）通用形式对象。这由引理 \ref{examples-defos-lemma-spaces-RS}、
\ref{examples-defos-lemma-proper-spaces-TI} 以及第 \ref{examples-defos-section-general} 节的一般讨论推出。

\begin{lemma}
\label{examples-defos-lemma-spaces-hull}
在例 \ref{examples-defos-example-spaces} 中，假设 $X$ 是 $k$ 上的固有代数空间。
再假设 $\Lambda$ 是剩余域为 $k$ 的完备局部环（经典情形）。则函子
$$
F : \mathcal{C}_\Lambda \longrightarrow \textit{Sets},\quad
A \longmapsto \Ob(\Deformationcategory_X(A))/\cong
$$
把对象映到其同构类；它有包络。若
$\text{Der}_k(\mathcal{O}_X, \mathcal{O}_X)=0$，
则 $F$ 是 pro-可表示的。
\end{lemma}

\begin{proof}
包络的存在性立即由引理 \ref{examples-defos-lemma-spaces-RS}、
\ref{examples-defos-lemma-proper-spaces-TI}、形式形变理论引理
\ref{formal-defos-lemma-RS-implies-S1-S2} 以及注
\ref{formal-defos-remark-compose-minimal-into-iso-classes} 推出。

\medskip\noindent
假设 $\text{Der}_k(\mathcal{O}_X,\mathcal{O}_X)=0$。
由形式形变理论引理 \ref{formal-defos-lemma-infdef-trivial}，
$\Deformationcategory_X$ 与 $F$ 等价。因此 $F$ 是具有有限维切空间的
形变函子（因为 $\Deformationcategory_X$ 是形变范畴），
我们可以应用形式形变理论定理
\ref{formal-defos-theorem-Schlessinger-prorepresentability}。
\end{proof}





\section{完备化的形变}
\label{examples-defos-section-compare}

\noindent
本节比较由一个代数及其完备化所给出的形变问题。
我们先讨论“可提升性”。

\begin{lemma}
\label{examples-defos-lemma-lift-equivalence-module-derived}
设 $A'\to A$ 是核为幂零理想的环满射，
$A'\to P'$ 是平坦环同态。置 $P=P'\otimes_{A'}A$。
设 $M$ 是在 $A$ 上平坦的 $P$-模。则下列条件等价：
\begin{enumerate}
\item 存在在 $A'$ 上平坦的 $P'$-模 $M'$，使得
$M'\otimes_{P'}P=M$；
\item 存在对象 $K'\in D^-(P')$，使得
$K'\otimes_{P'}^\mathbf{L}P=M$。
\end{enumerate}
\end{lemma}

\begin{proof}
假设 $M'$ 如 (1) 所述。则
% P11-E0054：冻结源下列张量等式链中前、后三个底环的方向按类型不相容。
$$
M = M' \otimes_P P' = M' \otimes_{A'} A =
M' \otimes_A^\mathbf{L} A' = M' \otimes_{P'}^\mathbf{L} P
$$
前两个等式是显然的；第三个等式由 $M'$ 在 $A'$ 上平坦得到；
第四个等式由代数进阶引理
\ref{more-algebra-lemma-base-change-comparison} 得到。因此 (2) 成立。
反之，假设 $K'$ 如 (2) 所述。我们可以并确实假设 $M$ 非零。
令 $t$ 为使 $H^t(K')$ 非零的最大整数（它因 $M$ 非零而存在）。
若 $t>0$，则
$H^t(K')\otimes_{P'}P=H^t(K'\otimes_{P'}^\mathbf{L}P)$ 为零。
由于 $P'\to P$ 的核幂零，由 Nakayama 引理可得 $H^t(K')=0$，
矛盾。因此 $t=0$（$t<0$ 的情形同样荒谬）。于是
$M'=H^0(K')$ 是满足 $M=M'\otimes_{P'}P$ 的 $P'$-模，
而 Tor 谱序列给出单射
$$
\text{Tor}_1^{P'}(M', P) \to H^{-1}(M' \otimes_{P'}^\mathbf{L} P) = 0
$$
。由上述导出基变换文献，
$0 = \text{Tor}_1^{P'}(M', P) = \text{Tor}_1^{A'}(M', A)$.
再由代数引理 \ref{algebra-lemma-what-does-it-mean} 可知，
$M'$ 在 $A'$ 上平坦。
\end{proof}

\begin{lemma}
\label{examples-defos-lemma-lift-equivalence-module}
考虑 Noether 环的交换图
$$
\xymatrix{
A' \ar[d] \ar[r] &
P' \ar[d] \ar[r] &
Q' \ar[d] \\
A \ar[r] &
P \ar[r] &
Q
}
$$
，其中两个方块均为笛卡儿方块，水平箭头均平坦，
竖直箭头均为核幂零的满射。设 $J'\subset P'$ 是满足
$P'/J'=Q'/J'Q'$ 的理想，且 $M$ 是在 $A$ 上平坦的 $P$-模。
假设对每个 $g\in J'$，都存在在 $A'$ 上平坦的
$(P')_g$-模提升 $M_g$。则下列条件等价：
\begin{enumerate}
\item $M$ 有一个在 $A'$ 上平坦的 $P'$-模提升；
\item $M\otimes_PQ$ 有一个在 $A'$ 上平坦的 $Q'$-模提升。
\end{enumerate}
\end{lemma}

\begin{proof}
令 $I=\Ker(A'\to A)$。对满足 $I^n=0$ 的整数 $n>1$ 作归纳，
可约化到 $I$ 是平方零理想的情形；细节从略。
我们把 $M$ 的可提升条件转化为寻找引理
\ref{examples-defos-lemma-lift-equivalence-module-derived} 中那样的
$D^-(P')$ 对象的问题。相应的障碍是元素
$$
\omega(M) \in \text{Ext}^2_P(M, M \otimes_P^\mathbf{L} IP) =
\text{Ext}^2_P(M, M \otimes_P IP)
$$
，它由形变理论引理 \ref{defos-lemma-canonical-class-algebra} 构造。
显示公式中的等式成立，是因为 $M$ 与 $P$ 都在 $A$ 上平坦，故
$M\otimes_P^\mathbf{L}IP=M\otimes_PIP$\footnote{取由自由
$A$-模组成的解消 $F_\bullet\to I$。由于 $A\to P$ 平坦，
$P\otimes_AF_\bullet$ 是 $IP$ 的自由解消。因此
$M\otimes_P^\mathbf{L}IP$ 由
$M\otimes_PP\otimes_AF_\bullet=M\otimes_AF_\bullet$ 表示。
由于 $M$ 在 $A$ 上平坦，该复形只在次数 $0$ 有上同调。}。
同样，提升 $M\otimes_PQ$ 的障碍是元素
$$
\omega(M \otimes_P Q) \in
\text{Ext}^2_Q(M \otimes_P Q, (M \otimes_P Q) \otimes_Q IQ)
$$
；由形变理论引理 \ref{defos-lemma-canonical-class-algebra}
中构造 $\omega(-)$ 的函子性，它是 $\omega(M)$ 的像。
由代数进阶引理 \ref{more-algebra-lemma-base-change-RHom}，有
$$
\text{Ext}^2_Q(M \otimes_P Q, (M \otimes_P Q) \otimes_Q IQ) =
\text{Ext}^2_P(M, M \otimes_P IP) \otimes_P Q
$$
。这里使用了 $P$ 是 Noether 环且 $M$ 有限。
% P11-E0055：冻结源引理的假设并未要求 M 有限，但证明在此明确使用该条件。
关于 $P'\to Q'$ 的假设保证：对任意 $P$-模 $E$，映射
$E\to E\otimes_PQ$ 在 $J'$-幂挠子模上为双射；参见代数进阶引理
\ref{more-algebra-lemma-neighbourhood-equivalence}。
因此只需证明 $\omega(M)$ 是 $J'$-幂挠的。换言之，只需证明
$\omega(M)$ 在
$$
\text{Ext}^2_P(M, M \otimes_P IP)_g =
\text{Ext}^2_{P_g}(M_g, M_g \otimes_{P_g} IP_g)
$$
中对所有 $g\in J'$ 都变为零。由 $\omega(M)$ 的构造与基变换的相容性，
这确实成立，因为假设 $M_g$ 有提升（当然还须再次使用整串等价）。
% P11-E0056：冻结源将 However 误拼为 Howeover。
\end{proof}

\begin{lemma}
\label{examples-defos-lemma-lift-equivalence}
设 $A'\to A$ 是核幂零的 Noether 环满射，
$A\to B$ 是有限型平坦环同态。设 $\mathfrak b\subset B$ 是理想，
使得 $\Spec(B)\to\Spec(A)$ 在 $V(\mathfrak b)$ 的补集上为 syntomic。
则 $B$ 有到 $A'$ 的平坦提升，当且仅当其 $\mathfrak b$-进完备化
$B^\wedge$ 有到 $A'$ 的平坦提升。
\end{lemma}

\begin{proof}
选取 $A$-代数满射 $P=A[x_1,\ldots,x_n]\to B$。
令 $\mathfrak p\subset P$ 为 $\mathfrak b$ 的逆像。置
$P'=A'[x_1,\ldots,x_n]$，并以 $\mathfrak p'\subset P'$
表示 $\mathfrak p$ 的逆像（当然，这里的 $\mathfrak p$ 与
$\mathfrak p'$ 并不表示素理想）。以 $P^\wedge$ 与 $(P')^\wedge$
分别表示相应的完备化。

\medskip\noindent
假设 $A'\to B'$ 是 $A\to B$ 的平坦提升；换言之，
$A'\to B'$ 平坦，且有 $A$-代数同构 $B=B'\otimes_{A'}A$。
于是可选取提升给定满射 $P\to B$ 的 $A'$-代数同态 $P'\to B'$。
由 Nakayama 引理（代数引理 \ref{algebra-lemma-NAK}），
$B'$ 是 $P'$ 的商。因此可以赋予 $B'$ 一个在 $A'$ 上平坦的
$P'$-模结构，它提升作为在 $A$ 上平坦的 $P$-模的 $B$。
反之，若可把 $B$ 提升为在 $A'$ 上平坦的 $P'$-模 $M'$，
则再次应用 Nakayama 可知 $M'$ 是循环模 $M'\cong P'/J'$；
置 $B'=P'/J'$，便得到 $B$ 作为代数的平坦提升。

\medskip\noindent
置 $C=B^\wedge$ 与 $\mathfrak c=\mathfrak bC$。
假设 $A'\to C'$ 是 $A\to C$ 的平坦提升。则 $C'$ 关于
$\mathfrak c$ 的逆像 $\mathfrak c'$ 完备（代数引理
\ref{algebra-lemma-complete-modulo-nilpotent}）。选取提升
$A$-代数同态 $P\to C$ 的 $A'$-代数同态 $P'\to C'$。
这些同态经由完备化给出满射 $P^\wedge\to C$ 与
$(P')^\wedge\to C'$（对后一个再次使用 Nakayama 引理）。
特别地，可以赋予 $C'$ 一个在 $A'$ 上平坦的 $(P')^\wedge$-模结构，
它提升作为在 $A$ 上平坦的 $P^\wedge$-模的 $C$。
反之，若可把 $C$ 提升为在 $A'$ 上平坦的 $(P')^\wedge$-模 $N'$，
则再次应用 Nakayama 可知 $N'$ 是循环模
$N'\cong(P')^\wedge/\tilde J$；置
$C'=(P')^\wedge/\tilde J$，便得到 $C$ 作为代数的平坦提升。

\medskip\noindent
注意，$P'\to(P')^\wedge$ 是平坦环同态，并诱导同构
$P'/\mathfrak p'=(P')^\wedge/\mathfrak p'(P')^\wedge$。
因此，只要能证明对 $g\in\mathfrak p'$，$B_g$ 可提升为
在 $A'$ 上平坦的 $P'_g$-模，本引理便由引理
\ref{examples-defos-lemma-lift-equivalence-module} 推出。然而环同态
$A\to B_g$ 是 syntomic，故由环同态的光滑化命题
\ref{smoothing-proposition-lift-smooth}，它可提升为在 $A'$ 上平坦的代数
$B'$。由于 $A'\to P'_g$ 光滑，可以像此前一样把
$P_g\to B_g$ 提升为满射 $P'_g\to B'$，这正是所需结论。
\end{proof}

\noindent
记号。设 $A\to B$ 是环同态，$N$ 是 $B$-模。以
$\text{Exal}_A(B,N)$ 表示扩张
$$
0 \to N \to C \to B \to 0
$$
的同构类所成的集合；这里扩张属于 $A$-代数，且 $N$ 是 $C$
中的平方零理想。给定另一个这样的扩张
$0\to N\to C'\to B\to0$，所谓同构，是使得下图交换的
$A$-代数同构 $C\to C'$：
$$
\xymatrix{
0 \ar[r] &
N \ar[r] \ar[d]_{\text{id}} &
C \ar[r] \ar[d] &
B \ar[r] \ar[d]_{\text{id}} & 0 \\
0 \ar[r] &
N \ar[r] &
C' \ar[r] &
B \ar[r] & 0
}
$$
指派 $N\mapsto\text{Exal}_A(B,N)$ 是把积变成积的函子。
因此它是加性函子，而 $\text{Exal}_A(B,N)$ 具有自然的
$B$-模结构。事实上，由形变理论引理 \ref{defos-lemma-choices}，
有 $\text{Exal}_A(B,N)=\text{Ext}^1_B(\NL_{B/A},N)$。

\begin{lemma}
\label{examples-defos-lemma-first-order-completion}
设 $k$ 是域，$B$ 是有限型 $k$-代数。设 $J\subset B$ 是理想，
使得 $\Spec(B)\to\Spec(k)$ 在 $V(J)$ 的补集上光滑。
设 $N$ 是有限 $B$-模。则有典范双射
$$
\text{Exal}_k(B, N) \to \text{Exal}_k(B^\wedge, N^\wedge)
$$
。这里 $B^\wedge$ 与 $N^\wedge$ 是 $J$-进完备化。
\end{lemma}

\begin{proof}
该映射由完备化给出：给定 $\text{Exal}_k(B,N)$ 中的
$0\to N\to C\to B\to0$，把它送到 $C$ 关于 $J$ 的逆像的
完备化 $C^\wedge$。可与引理 \ref{examples-defos-lemma-completion} 的证明比较。

\medskip\noindent
由于 $k\to B$ 是有限呈示，复形 $\NL_{B/k}$ 可由复形
$N^{-1}\to N^0$ 表示，其中每个 $N^i$ 都是有限 $B$-模；
参见代数第 \ref{algebra-section-netherlander} 节，特别是代数引理
\ref{algebra-lemma-NL-homotopy}。由于 $B$ 是 Noether 环，
这意味着 $\NL_{B/k}$ 是伪凝聚的。对 $g\in J$，$k$-代数
$B_g$ 光滑，因而 $(\NL_{B/k})_g=\NL_{B_g/k}$ 拟同构于置于次数
$0$ 的有限投射 $B$-模。
% P11-E0058：冻结源此处的局部化复形自然应是有限投射 B_g-模，而非 B-模。
所以对任意 $B$-模 $N$ 及 $i\geq1$，有
$\text{Ext}^i_B(\NL_{B/k},N)_g=0$。由代数进阶引理
\ref{more-algebra-lemma-ext-annihilated-into} 可得
$$
\text{Ext}^1_B(\NL_{B/k}, N) \longrightarrow
\lim_n \text{Ext}^1_B(\NL_{B/k}, N/J^n N)
$$
对任意有限 $B$-模 $N$ 都是同构。

\medskip\noindent
映射的单射性。假设
$0\to N\to C\to B\to0$ 属于 $\text{Exal}_k(B,N)$，
并在 $\text{Exal}_k(B^\wedge,N^\wedge)$ 中映为零。
选取分裂 $C^\wedge=B^\wedge\oplus N^\wedge$。于是诱导映射
$C\to C^\wedge\to N^\wedge$ 对每个 $n$ 都给出映射
$C\to N/J^nN$。因此可见，我们的元素属于各映射
$$
\text{Ext}^1_B(\NL_{B/k}, N) \to
\text{Ext}^1_B(\NL_{B/k}, N/J^n N)
$$
的核，且这对所有 $n$ 均成立。由上一段可知，该元素为零。

\medskip\noindent
映射的满射性。设
$0\to N^\wedge\to C'\to B^\wedge\to0$ 是
$\text{Exal}_k(B^\wedge,N^\wedge)$ 中的元素。
沿 $B\to B^\wedge$ 拉回，得到 $\text{Exal}_k(B,N^\wedge)$
中的元素 $0\to N^\wedge\to C''\to B\to0$。有
$$
\text{Ext}^1_B(\NL_{B/k}, N^\wedge) =
\text{Ext}^1_B(\NL_{B/k}, N) \otimes_B B^\wedge =
\text{Ext}^1_B(\NL_{B/k}, N)
$$
。第一个等式由 $N^\wedge=N\otimes_BB^\wedge$（代数引理
\ref{algebra-lemma-completion-tensor}）以及代数进阶引理
\ref{more-algebra-lemma-pseudo-coherence-and-ext} 得到。
第二个等式成立，是因为 $\text{Ext}^1_B(\NL_{B/k},N)$
是 $J$-幂挠的（见上文），$B\to B^\wedge$ 平坦并诱导同构
$B/J\to B^\wedge/JB^\wedge$，以及代数进阶引理
\ref{more-algebra-lemma-neighbourhood-equivalence}。因此可找到
$C\in\text{Exal}_k(B,N)$，它在
$\text{Exal}_k(B,N^\wedge)$ 中映到 $C''$。于是
$$
0 \to N^\wedge \to C' \to B^\wedge \to 0
\quad\text{且}\quad
0 \to N^\wedge \to C^\wedge \to B^\wedge \to 0
$$
是 $\text{Exal}_k(B^\wedge,N^\wedge)$ 中的两个元素，
它们映到 $\text{Exal}_k(B,N^\wedge)$ 中的同一个元素。
取差，得到 $\text{Exal}_k(B^\wedge,N^\wedge)$ 中的元素
$0\to N^\wedge\to C'\to B^\wedge\to0$，其在
$\text{Exal}_k(B,N^\wedge)$ 中的像为零。这意味着存在
$$
\xymatrix{
0 \ar[r] &
N^\wedge \ar[r] &
C' \ar[r] &
B^\wedge \ar[r] & 0 \\
& & B \ar[u]^\sigma \ar[ru]
}
$$
。令 $J'\subset C'$ 为 $JB^\wedge\subset B^\wedge$ 的逆像。
为完成证明，只需注意：$\sigma$ 关于 $B$ 上的 $J$-进拓扑和
$C'$ 上的 $J'$-进拓扑连续，而且由代数引理
\ref{algebra-lemma-complete-modulo-nilpotent}，$C'$ 是 $J'$-进完备的
（这里还使用了 $C'$ 是 Noether 环；略去一个小细节）。
这就是说，$\sigma$ 经由完备化 $B^\wedge$ 分解，且
$C'=0$ 在 $\text{Exal}_k(B^\wedge,N^\wedge)$ 中成立。
\end{proof}

\begin{lemma}
\label{examples-defos-lemma-smooth-completion}
在例 \ref{examples-defos-example-rings} 中，设 $P$ 是 $k$-代数，
$J\subset P$ 是理想，并以 $P^\wedge$ 表示 $J$-进完备化。若
\begin{enumerate}
\item $k\to P$ 是有限型的；
\item $\Spec(P)\to\Spec(k)$ 在 $V(J)$ 的补集上光滑，
\end{enumerate}
则引理 \ref{examples-defos-lemma-completion} 中形变范畴之间的函子
$$
\Deformationcategory_P \longrightarrow \Deformationcategory_{P^\wedge}
$$
光滑，并在切空间上诱导同构。
\end{lemma}

\begin{proof}
由引理 \ref{examples-defos-lemma-rings-RS}，我们知道
$\Deformationcategory_P$ 与 $\Deformationcategory_{P^\wedge}$
都是形变范畴。因此只须检验该函子等同切空间并对应可提升性；
参见形式形变理论引理 \ref{formal-defos-lemma-easy-check-smooth}。
可提升性由引理 \ref{examples-defos-lemma-lift-equivalence} 证明，
切空间上的同构则是引理 \ref{examples-defos-lemma-first-order-completion}
中 $N=B$ 的特殊情形。
\end{proof}



\section{局部化的形变}
\label{examples-defos-section-compare-localization}

\noindent
本节比较由一个代数及其关于乘法子集的局部化所给出的形变问题。
我们先讨论“可提升性”。

\begin{lemma}
\label{examples-defos-lemma-lift-equivalence-localization}
设 $A'\to A$ 是核幂零的 Noether 环满射，
$A\to B$ 是有限型平坦环同态。设 $S\subset B$ 是乘法子集，
且只要 $\Spec(B)\to\Spec(A)$ 在 $\mathfrak q$ 处不是 syntomic，
就有 $S\cap\mathfrak q=\emptyset$。则 $B$ 有到 $A'$ 的平坦提升，
当且仅当 $S^{-1}B$ 有到 $A'$ 的平坦提升。
\end{lemma}

\begin{proof}
本证明与引理 \ref{examples-defos-lemma-lift-equivalence} 的证明相同，但更容易。
建议读者略过。选取 $A$-代数满射
$P=A[x_1,\ldots,x_n]\to B$。令 $S_P\subset P$ 为 $S$ 的逆像。
置 $P'=A'[x_1,\ldots,x_n]$，并以 $S_{P'}\subset P'$
表示 $S_P$ 的逆像。

\medskip\noindent
假设 $A'\to B'$ 是 $A\to B$ 的平坦提升；换言之，
$A'\to B'$ 平坦，且有 $A$-代数同构 $B=B'\otimes_{A'}A$。
于是可选取提升给定满射 $P\to B$ 的 $A'$-代数同态 $P'\to B'$。
由 Nakayama 引理（代数引理 \ref{algebra-lemma-NAK}），
$B'$ 是 $P'$ 的商。特别地，可以赋予 $B'$ 一个在 $A'$ 上平坦的
$P'$-模结构，它提升作为在 $A$ 上平坦的 $P$-模的 $B$。
反之，若可把 $B$ 提升为在 $A'$ 上平坦的 $P'$-模 $M'$，
则再次应用 Nakayama 可知 $M'$ 是循环模 $M'\cong P'/J'$；
置 $B'=P'/J'$，便得到 $B$ 作为代数的平坦提升。

\medskip\noindent
置 $C=S^{-1}B$。假设 $A'\to C'$ 是 $A\to C$ 的平坦提升。
$C'$ 中映到 $C$ 的可逆元的元素本身也可逆。选取提升
$A$-代数同态 $P\to C$ 的 $A'$-代数同态 $P'\to C'$。
由上述说明，这些同态经由局部化给出满射
$S_P^{-1}P\to C$ 与 $S_{P'}^{-1}P'\to C'$（对后一个使用
Nakayama 引理）。特别地，可以赋予 $C'$ 一个在 $A'$ 上平坦的
$S_{P'}^{-1}P'$-模结构，它提升作为在 $A$ 上平坦的
$S_P^{-1}P$-模的 $C$。反之，若可把 $C$ 提升为在 $A'$ 上平坦的
$S_{P'}^{-1}P'$-模 $N'$，则再次应用 Nakayama 可知 $N'$ 是循环模
$N'\cong S_{P'}^{-1}P'/\tilde J$；置
$C'=S_{P'}^{-1}P'/\tilde J$，便得到 $C$ 作为代数的平坦提升。

\medskip\noindent
概形态射的 syntomic 轨迹按定义是开集。令 $J_B\subset B$ 为
定义 $\Spec(B)$ 中 $\Spec(B)\to\Spec(A)$ 不是 syntomic 的点集的理想，
并以 $J_P\subset P$ 与 $J_{P'}\subset P'$ 表示相应理想。
注意，由引理中关于 $S$ 的假设，平坦环同态
$P'\to S_{P'}^{-1}P'$ 诱导同构
$P'/J_{P'}=S_{P'}^{-1}P'/J_{P'}S_{P'}^{-1}P'$；具体地说，
引理中的假设恰好是 $B/J_B=S^{-1}(B/J_B)$。
因此，只要能证明对 $g\in J_B$，$B_g$ 可提升为在 $A'$ 上平坦的
$P'_g$-模，本引理便由引理 \ref{examples-defos-lemma-lift-equivalence-module} 推出。
% P11-E0063：冻结源以 g∈J_B 开始，却随后写 P'_g；g 并非 P' 中已选定的元素。
然而环同态 $A\to B_g$ 是 syntomic，故由环同态的光滑化命题
\ref{smoothing-proposition-lift-smooth}，它可提升为在 $A'$ 上平坦的
代数 $B'$。由于 $A'\to P'_g$ 光滑，可以像此前一样把
$P_g\to B_g$ 提升为满射 $P'_g\to B'$，这正是所需结论。
\end{proof}

\begin{lemma}
\label{examples-defos-lemma-first-order-localization}
设 $k$ 是域，$B$ 是有限型 $k$-代数。设 $S\subset B$ 是乘法子集，
且只要 $\Spec(B)\to\Spec(k)$ 在 $\mathfrak q$ 处不光滑，
就有 $S\cap\mathfrak q=\emptyset$。设 $N$ 是有限 $B$-模。
则有典范双射
% P11-E0060：冻结源将 S 同时称为 multiplicative subset ideal，多出 ideal。
$$
\text{Exal}_k(B, N) \to \text{Exal}_k(S^{-1}B, S^{-1}N)
$$
\end{lemma}

\begin{proof}
本证明与引理 \ref{examples-defos-lemma-first-order-completion} 的证明相同，但更容易。
建议读者略过。该映射由局部化给出：给定
$\text{Exal}_k(B,N)$ 中的 $0\to N\to C\to B\to0$，
把它送到 $C$ 关于 $S$ 的逆像 $S_C\subset C$ 的局部化
$S_C^{-1}C$。可与引理 \ref{examples-defos-lemma-localization} 的证明比较。

\medskip\noindent
概形态射的光滑轨迹按定义是开集。令 $J\subset B$ 为定义
$\Spec(B)$ 中 $\Spec(B)\to\Spec(A)$ 不光滑的点集的理想。
% P11-E0061：本引理只定义了 k，此处冻结源却写 Spec(A)。
由于 $k\to B$ 是有限呈示，复形 $\NL_{B/k}$ 可由复形
$N^{-1}\to N^0$ 表示，其中每个 $N^i$ 都是有限 $B$-模；
参见代数第 \ref{algebra-section-netherlander} 节，特别是代数引理
\ref{algebra-lemma-NL-homotopy}。由于 $B$ 是 Noether 环，
这意味着 $\NL_{B/k}$ 是伪凝聚的。对 $g\in J$，$k$-代数
$B_g$ 光滑，因而 $(\NL_{B/k})_g=\NL_{B_g/k}$ 拟同构于置于次数
$0$ 的有限投射 $B$-模。
% P11-E0062：冻结源此处的局部化复形自然应是有限投射 B_g-模，而非 B-模。
所以对任意 $B$-模 $N$ 及 $i\geq1$，有
$\text{Ext}^i_B(\NL_{B/k},N)_g=0$。最后，有
$$
\text{Ext}^1_{S^{-1}B}(\NL_{S^{-1}B/k}, S^{-1}N) =
\text{Ext}^1_B(\NL_{B/k}, N) \otimes_B S^{-1}B =
\text{Ext}^1_B(\NL_{B/k}, N)
$$
。第一个等式由代数进阶引理
\ref{more-algebra-lemma-base-change-RHom} 与代数引理
\ref{algebra-lemma-localize-NL} 得到。第二个等式成立，是因为
$\text{Ext}^1_B(\NL_{B/k},N)$ 是 $J$-幂挠的，而 $S$ 中元素在
$J$-幂挠模上可逆地作用。结合引理 \ref{examples-defos-lemma-first-order-completion}
紧前方把 $\text{Exal}_A(B,N)$ 描述为
$\text{Ext}^1_B(\NL_{B/A},N)$ 的结论，证明完毕。
\end{proof}

\begin{lemma}
\label{examples-defos-lemma-smooth-localization}
在例 \ref{examples-defos-example-rings} 中，设 $P$ 是 $k$-代数，
$S\subset P$ 是乘法子集。若
\begin{enumerate}
\item $k\to P$ 是有限型的；
\item 对每个 $g\in S$，$\Spec(P)\to\Spec(k)$
在 $V(g)$ 的所有点处都光滑，
\end{enumerate}
则引理 \ref{examples-defos-lemma-localization} 中形变范畴之间的函子
$$
\Deformationcategory_P \longrightarrow \Deformationcategory_{S^{-1}P}
$$
光滑，并在切空间上诱导同构。
\end{lemma}

\begin{proof}
由引理 \ref{examples-defos-lemma-rings-RS}，我们知道
$\Deformationcategory_P$ 与 $\Deformationcategory_{S^{-1}P}$
都是形变范畴。因此只须检验该函子等同切空间并对应可提升性；
参见形式形变理论引理 \ref{formal-defos-lemma-easy-check-smooth}。
可提升性由引理 \ref{examples-defos-lemma-lift-equivalence-localization} 证明，
切空间上的同构则是引理 \ref{examples-defos-lemma-first-order-localization}
中 $N=B$ 的特殊情形。
\end{proof}



\section{Hensel 化的形变}
\label{examples-defos-section-compare-henselization}

\noindent
本节比较由一个代数及其完备化所给出的形变问题。
% P11-E0064：本节讨论 Hensel 化，冻结源却在此误写 completion。
我们先讨论“可提升性”。

\begin{lemma}
\label{examples-defos-lemma-lift-equivalence-henselization}
设 $A'\to A$ 是核幂零的 Noether 环满射，
$A\to B$ 是有限型平坦环同态。设 $\mathfrak b\subset B$ 是理想，
使得 $\Spec(B)\to\Spec(A)$ 在 $V(\mathfrak b)$ 的补集上为 syntomic。
设 $(B^h,\mathfrak b^h)$ 是对 $(B,\mathfrak b)$ 的 Hensel 化。
则 $B$ 有到 $A'$ 的平坦提升，当且仅当 $B^h$ 有到 $A'$ 的平坦提升。
\end{lemma}

\begin{proof}[第一种证明]
这个证明有所取巧。事实上，若 $B$ 有平坦提升 $B'$，
则取 Hensel 化 $(B')^h$，便得到 $B^h$ 的平坦提升
（可与引理 \ref{examples-defos-lemma-henselization} 的证明比较）。
反之，假设 $C'$ 是 $(B')^h$ 的一个在 $A'$ 上平坦的提升。
% P11-E0065：反向论证中冻结源写 (B')^h，但此处应为待提升的 B^h。
令 $\mathfrak c'\subset C'$ 为理想 $\mathfrak b^h$ 的逆像。
则 $C'$ 关于 $\mathfrak c'$ 的完备化 $(C')^\wedge$
是 $B^\wedge$ 的提升（细节从略）。因此由引理
\ref{examples-defos-lemma-lift-equivalence} 可知，$B$ 有平坦提升。
\end{proof}

\begin{proof}[第二种证明]
选取 $A$-代数满射 $P=A[x_1,\ldots,x_n]\to B$。
令 $\mathfrak p\subset P$ 为 $\mathfrak b$ 的逆像。置
$P'=A'[x_1,\ldots,x_n]$，并以 $\mathfrak p'\subset P'$
表示 $\mathfrak p$ 的逆像（当然，这里的 $\mathfrak p$ 与
$\mathfrak p'$ 并不表示素理想）。以 $P^h$ 与 $(P')^h$
分别表示相应的 Hensel 化。我们将使用 Hensel 化的函子性，
以及商的 Hensel 化就是 Hensel 化的相应商；参见代数进阶引理
\ref{more-algebra-lemma-irreducible-henselian-pair-connected} 与
\ref{more-algebra-lemma-henselization-integral}。

\medskip\noindent
假设 $A'\to B'$ 是 $A\to B$ 的平坦提升；换言之，
$A'\to B'$ 平坦，且有 $A$-代数同构 $B=B'\otimes_{A'}A$。
于是可选取提升给定满射 $P\to B$ 的 $A'$-代数同态 $P'\to B'$。
由 Nakayama 引理（代数引理 \ref{algebra-lemma-NAK}），
$B'$ 是 $P'$ 的商。特别地，可以赋予 $B'$ 一个在 $A'$ 上平坦的
$P'$-模结构，它提升作为在 $A$ 上平坦的 $P$-模的 $B$。
反之，若可把 $B$ 提升为在 $A'$ 上平坦的 $P'$-模 $M'$，
则再次应用 Nakayama 可知 $M'$ 是循环模 $M'\cong P'/J'$；
置 $B'=P'/J'$，便得到 $B$ 作为代数的平坦提升。

\medskip\noindent
置 $C=B^h$ 与 $\mathfrak c=\mathfrak bC$。假设 $A'\to C'$
是 $A\to C$ 的平坦提升。则 $C'$ 关于 $\mathfrak c$ 的逆像
$\mathfrak c'$ 是 Hensel 的（由代数进阶引理
\ref{more-algebra-lemma-henselian-henselian-pair} 以及
$C'\to C$ 的核幂零这一事实）。选取提升 $A$-代数同态
$P\to C$ 的 $A'$-代数同态 $P'\to C'$。这些同态经由 Hensel 化
给出满射 $P^h\to C$ 与 $(P')^h\to C'$（对后一个再次使用
Nakayama 引理）。特别地，可以赋予 $C'$ 一个在 $A'$ 上平坦的
$(P')^h$-模结构，它提升作为在 $A$ 上平坦的 $P^h$-模的 $C$。
反之，若可把 $C$ 提升为在 $A'$ 上平坦的 $(P')^h$-模 $N'$，
则再次应用 Nakayama 可知 $N'$ 是循环模
$N'\cong(P')^h/\tilde J$；置 $C'=(P')^h/\tilde J$，
便得到 $C$ 作为代数的平坦提升。

\medskip\noindent
注意，$P'\to(P')^h$ 是平坦环同态，并诱导同构
$P'/\mathfrak p'=(P')^h/\mathfrak p'(P')^h$（代数进阶引理
\ref{more-algebra-lemma-henselization-flat}）。因此，只要能证明对
$g\in\mathfrak p'$，$B_g$ 可提升为在 $A'$ 上平坦的 $P'_g$-模，
本引理便由引理 \ref{examples-defos-lemma-lift-equivalence-module} 推出。然而环同态
$A\to B_g$ 是 syntomic，故由环同态的光滑化命题
\ref{smoothing-proposition-lift-smooth}，它可提升为在 $A'$ 上平坦的
代数 $B'$。由于 $A'\to P'_g$ 光滑，可以像此前一样把
$P_g\to B_g$ 提升为满射 $P'_g\to B'$，这正是所需结论。
\end{proof}

\begin{lemma}
\label{examples-defos-lemma-first-order-henselization}
设 $k$ 是域，$B$ 是有限型 $k$-代数。设 $J\subset B$ 是理想，
使得 $\Spec(B)\to\Spec(k)$ 在 $V(J)$ 的补集上光滑。
设 $N$ 是有限 $B$-模。则有典范双射
$$
\text{Exal}_k(B, N) \to \text{Exal}_k(B^h, N^h)
$$
。这里 $(B^h,J^h)$ 是对 $(B,J)$ 的 Hensel 化，
且 $N^h=N\otimes_BB^h$。
\end{lemma}

\begin{proof}
本证明与引理 \ref{examples-defos-lemma-first-order-completion} 的证明相同，但更容易。
建议读者略过。该映射由 Hensel 化给出：给定
$\text{Exal}_k(B,N)$ 中的 $0\to N\to C\to B\to0$，
把它送到 $C$ 关于 $J$ 的逆像 $J_C\subset C$ 的 Hensel 化
$C^h$。可与引理 \ref{examples-defos-lemma-henselization} 的证明比较。

\medskip\noindent
由于 $k\to B$ 是有限呈示，复形 $\NL_{B/k}$ 可由复形
$N^{-1}\to N^0$ 表示，其中每个 $N^i$ 都是有限 $B$-模；
参见代数第 \ref{algebra-section-netherlander} 节，特别是代数引理
\ref{algebra-lemma-NL-homotopy}。由于 $B$ 是 Noether 环，
这意味着 $\NL_{B/k}$ 是伪凝聚的。对 $g\in J$，$k$-代数
$B_g$ 光滑，因而 $(\NL_{B/k})_g=\NL_{B_g/k}$ 拟同构于置于次数
$0$ 的有限投射 $B$-模。
% P11-E0066：冻结源此处的局部化复形自然应是有限投射 B_g-模，而非 B-模。
所以对任意 $B$-模 $N$ 及 $i\geq1$，有
$\text{Ext}^i_B(\NL_{B/k},N)_g=0$。最后，有
\begin{align*}
\text{Ext}^1_{B^h}(\NL_{B^h/k}, N^h)
& =
\text{Ext}^1_{B^h}(\NL_{B/k} \otimes_B B^h, N \otimes_B B^h) \\
& =
\text{Ext}^1_B(\NL_{B/k}, N) \otimes_B B^h \\
& =
\text{Ext}^1_B(\NL_{B/k}, N)
\end{align*}
第一个等式由代数进阶引理
\ref{more-algebra-lemma-henselization-NL}（更准确地说，是其关于
Hensel 对之 Hensel 化的类似结论）得到。第二个等式由代数进阶引理
\ref{more-algebra-lemma-base-change-RHom} 得到。第三个等式成立，
是因为 $\text{Ext}^1_B(\NL_{B/k},N)$ 是 $J$-幂挠的，映射
$B\to B^h$ 平坦并诱导同构 $B/J\to B^h/JB^h$（代数进阶引理
\ref{more-algebra-lemma-henselization-flat}），以及代数进阶引理
\ref{more-algebra-lemma-neighbourhood-equivalence}。结合引理
\ref{examples-defos-lemma-first-order-completion} 紧前方把
$\text{Exal}_A(B,N)$ 描述为 $\text{Ext}^1_B(\NL_{B/A},N)$ 的结论，
证明完毕。
\end{proof}

\begin{lemma}
\label{examples-defos-lemma-smooth-henselization}
在例 \ref{examples-defos-example-rings} 中，设 $P$ 是 $k$-代数，
$J\subset P$ 是理想。以 $(P^h,J^h)$ 表示对 $(P,J)$ 的 Hensel 化。若
\begin{enumerate}
\item $k\to P$ 是有限型的；
\item $\Spec(P)\to\Spec(k)$ 在 $V(J)$ 的补集上光滑，
\end{enumerate}
则引理 \ref{examples-defos-lemma-henselization} 中形变范畴之间的函子
$$
\Deformationcategory_P \longrightarrow \Deformationcategory_{P^h}
$$
光滑，并在切空间上诱导同构。
\end{lemma}

\begin{proof}
由引理 \ref{examples-defos-lemma-rings-RS}，我们知道
$\Deformationcategory_P$ 与 $\Deformationcategory_{P^h}$ 都是形变范畴。
因此只须检验该函子等同切空间并对应可提升性；参见形式形变理论引理
\ref{formal-defos-lemma-easy-check-smooth}。可提升性由引理
\ref{examples-defos-lemma-lift-equivalence-henselization} 证明，切空间上的同构则是
引理 \ref{examples-defos-lemma-first-order-henselization} 中 $N=B$ 的特殊情形。
\end{proof}



\section{孤立奇点的应用}
\label{examples-defos-section-isolated}

\noindent
我们应用上述讨论，研究只有有限多个奇点的有限型代数的形变理论。

\begin{lemma}
\label{examples-defos-lemma-isolated}
在例 \ref{examples-defos-example-rings} 中，设 $P$ 是 $k$-代数。
假设 $k\to P$ 是有限型的，且 $\Spec(P)\to\Spec(k)$
除在 $P$ 的极大理想 $\mathfrak m_1,\ldots,\mathfrak m_n$ 处外均光滑。
分别以 $P_{\mathfrak m_i}$、$P_{\mathfrak m_i}^h$、
$P_{\mathfrak m_i}^\wedge$ 表示局部环、Hensel 化和完备化。
则形变范畴的映射
$$
\Deformationcategory_P \to
\prod \Deformationcategory_{P_{\mathfrak m_i}} \to
\prod \Deformationcategory_{P_{\mathfrak m_i}^h} \to
\prod \Deformationcategory_{P_{\mathfrak m_i}^\wedge}
$$
光滑，并在其有限维切空间上诱导同构。
\end{lemma}

\begin{proof}
由引理 \ref{examples-defos-lemma-finite-type-rings-TI}，切空间是有限维的。
这些范畴之间的函子在引理 \ref{examples-defos-lemma-localization}、
\ref{examples-defos-lemma-henselization} 与 \ref{examples-defos-lemma-completion} 中构造
（我们略去一些形如“Hensel 化的完备化就是完备化”的验证）。

\medskip\noindent
置 $J=\mathfrak m_1\cap\ldots\cap\mathfrak m_n$，应用引理
\ref{examples-defos-lemma-smooth-completion} 可知
$\Deformationcategory_P \to \Deformationcategory_{P^\wedge}$
光滑，并在切空间上诱导同构；这里 $P^\wedge$ 是 $P$ 的
$J$-进完备化。然而，由于 $P^\wedge=\prod P_{\mathfrak m_i}^\wedge$，
可见映射 $\Deformationcategory_P\to
\prod \Deformationcategory_{P_{\mathfrak m_i}^\wedge}$
光滑，并在切空间上诱导同构。

\medskip\noindent
设 $(P^h,J^h)$ 是对 $(P,J)$ 的 Hensel 化。则
$P^h=\prod P_{\mathfrak m_i}^h$（考察幂等元并应用代数进阶引理
\ref{more-algebra-lemma-characterize-henselian-pair}）。
因此可应用引理 \ref{examples-defos-lemma-smooth-henselization}，
像完备化情形一样得到结论。

\medskip\noindent
为得到最后一种情形，只需分别对每个 $i$ 证明
$\Deformationcategory_{P_{\mathfrak m_i}} \to
\Deformationcategory_{P_{\mathfrak m_i}^\wedge}$
光滑，并在切空间上诱导同构。
% P11-E0067：冻结源的单数主语 map 后写 induce，缺少第三人称单数 s。
为此，可把 $P$ 替换为一个主局部化，使其唯一奇点是极大理想
$\mathfrak m$（对应于原来 $P$ 中的 $\mathfrak m_i$）。
然后对乘法子集 $S=P\setminus\mathfrak m$ 应用引理
\ref{examples-defos-lemma-smooth-localization} 即得结论。略去次要细节。
\end{proof}






\section{无阻碍形变问题}
\label{examples-defos-section-unobstructed}

\noindent
设 $p:\mathcal{F}\to\mathcal{C}_\Lambda$ 是以群胚为纤维的余纤维化范畴。
回忆，若 $p$ 光滑，就称 $\mathcal{F}$ {\it 光滑}或{\it 无阻碍}。
这意味着：给定 $\mathcal{C}_\Lambda$ 中的满射
$\varphi:A'\to A$ 以及 $x\in\Ob(\mathcal{F}(A))$，
存在 $\mathcal{F}$ 中的态射 $f:x'\to x$，满足
$p(f)=\varphi$。参见形式形变理论第
\ref{formal-defos-section-smooth} 节。本节给出若干具有几何意义的例子。

\begin{lemma}
\label{examples-defos-lemma-lci-unobstructed}
在例 \ref{examples-defos-example-rings} 中，设 $P$ 是 $k$ 上的局部完全交
（代数定义 \ref{algebra-definition-lci-field}）。
则 $\Deformationcategory_P$ 无阻碍。
\end{lemma}

\begin{proof}
设 $(A,Q)\to(k,P)$ 是 $\Deformationcategory_P$ 的对象。
由代数定义 \ref{algebra-definition-lci}，$A\to Q$ 是 syntomic 环同态。
因此，对 $\mathcal{C}_\Lambda$ 中任意满射 $A'\to A$，
由环同态的光滑化命题 \ref{smoothing-proposition-lift-smooth}，
存在提升 $A'\to A$ 的态射 $(A',Q')\to(A,Q)$。这就证明了引理。
\end{proof}

\begin{lemma}
\label{examples-defos-lemma-glueing-smooth}
在情形 \ref{examples-defos-situation-glueing} 中，若 $U_{12}\to\Spec(k)$ 光滑，
则态射
$$
\Deformationcategory_X
\longrightarrow
\Deformationcategory_{U_1} \times \Deformationcategory_{U_2} =
\Deformationcategory_{P_1} \times \Deformationcategory_{P_2}
$$
光滑。若此外 $U_1$ 是 $k$ 上的局部完全交，则
$$
\Deformationcategory_X
\longrightarrow
\Deformationcategory_{U_2} = \Deformationcategory_{P_2}
$$
光滑。
\end{lemma}

\begin{proof}
这些等号由引理 \ref{examples-defos-lemma-affine} 成立。按照形式形变理论第
\ref{formal-defos-section-smooth} 节，把 $\mathcal{C}_\Lambda$
看作 $\mathcal{C}_\Lambda$ 上的形变范畴。于是
$$
\Deformationcategory_{P_1} \times \Deformationcategory_{P_2} =
\Deformationcategory_{P_1}
\times_{\mathcal{C}_\Lambda}
\Deformationcategory_{P_2},
$$
参见形式形变理论注
\ref{formal-defos-remarks-cofibered-groupoids}
（\ref{formal-defos-item-product}）。利用引理 \ref{examples-defos-lemma-glueing}，
第一个断言是函子
$$
\Deformationcategory_{P_1}
\times_{\Deformationcategory_{P_{12}}}
\Deformationcategory_{P_2}
\longrightarrow
\Deformationcategory_{P_1}
\times_{\mathcal{C}_\Lambda}
\Deformationcategory_{P_2}
$$
光滑。只要能证明 $T\Deformationcategory_{P_{12}}=(0)$，
这就由形式形变理论引理
\ref{formal-defos-lemma-map-fibre-products-smooth} 得到。
由于 $P_{12}$ 在 $k$ 上光滑，该消失性由引理
\ref{examples-defos-lemma-smooth} 得到。对第二个断言，只需证明
$\Deformationcategory_{P_1} \to \mathcal{C}_\Lambda$
光滑；参见形式形变理论引理
\ref{formal-defos-lemma-smooth-properties}。换言之，须证明
$\Deformationcategory_{P_1}$ 无阻碍，这正是引理
\ref{examples-defos-lemma-lci-unobstructed}。
\end{proof}

\begin{lemma}
\label{examples-defos-lemma-curve-isolated}
在例 \ref{examples-defos-example-schemes} 中，设 $X$ 是 $k$ 上的概形。假设
\begin{enumerate}
\item $X$ 是分离的、在 $k$ 上有限型，且 $\dim(X)\leq1$；
\item $X\to\Spec(k)$ 除在闭点 $p_1,\ldots,p_n\in X$ 处外均光滑。
\end{enumerate}
分别以 $\mathcal{O}_{X,p_1}$、$\mathcal{O}_{X,p_1}^h$、
$\mathcal{O}_{X,p_1}^\wedge$ 表示局部环、Hensel 化和完备化。
% P11-E0068：冻结源在此三次写 p_1，而随后乘积及证明均按 p_i 取遍各奇点。
考虑形变范畴的映射
$$
\Deformationcategory_X
\longrightarrow
\prod \Deformationcategory_{\mathcal{O}_{X, p_i}}
\longrightarrow
\prod \Deformationcategory_{\mathcal{O}_{X, p_i}^h}
\longrightarrow
\prod \Deformationcategory_{\mathcal{O}_{X, p_i}^\wedge}
$$
第一个箭头光滑；第二、第三个箭头光滑，并在切空间上诱导同构。
\end{lemma}

\begin{proof}
选取仿射开集 $U_2\subset X$，使其包含 $p_1,\ldots,p_n$
以及 $X$ 每个不可约分支的泛点。这由代数簇引理
\ref{varieties-lemma-dim-1-quasi-projective} 与性质引理
\ref{properties-lemma-ample-finite-set-in-affine} 得以实现。
于是 $X\setminus U_2$ 是有限集，可以选取仿射开集
$U_1\subset X\setminus\{p_1,\ldots,p_n\}$，使得
$X=U_1\cup U_2$。置 $U_{12}=U_1\cap U_2$。
此时 $U_1$ 与 $U_{12}$ 都是 $k$ 上的光滑仿射概形。因此
$$
\Deformationcategory_X \longrightarrow \Deformationcategory_{U_2}
$$
由引理 \ref{examples-defos-lemma-glueing-smooth} 是光滑的。
应用引理 \ref{examples-defos-lemma-affine} 与 \ref{examples-defos-lemma-isolated} 即得结论。
\end{proof}

\begin{lemma}
\label{examples-defos-lemma-curve-isolated-lci}
在例 \ref{examples-defos-example-schemes} 中，设 $X$ 是 $k$ 上的概形。假设
\begin{enumerate}
\item $X$ 是分离的、在 $k$ 上有限型，且 $\dim(X)\leq1$；
\item $X$ 是 $k$ 上的局部完全交；
\item $X\to\Spec(k)$ 除在有限多个点处外均光滑。
\end{enumerate}
则 $\Deformationcategory_X$ 无阻碍。
\end{lemma}

\begin{proof}
设 $p_1,\ldots,p_n\in X$ 是 $X\to\Spec(k)$ 不光滑的点。
选取仿射开集 $U_2\subset X$，使其包含 $p_1,\ldots,p_n$
以及 $X$ 每个不可约分支的泛点。这由代数簇引理
\ref{varieties-lemma-dim-1-quasi-projective} 与性质引理
\ref{properties-lemma-ample-finite-set-in-affine} 得以实现。
于是 $X\setminus U_2$ 是有限集，可以选取仿射开集
$U_1\subset X\setminus\{p_1,\ldots,p_n\}$，使得
$X=U_1\cup U_2$。置 $U_{12}=U_1\cap U_2$。
此时 $U_1$ 与 $U_{12}$ 都是 $k$ 上的光滑仿射概形。因此
$$
\Deformationcategory_X \longrightarrow \Deformationcategory_{U_2}
$$
由引理 \ref{examples-defos-lemma-glueing-smooth} 是光滑的。
应用引理 \ref{examples-defos-lemma-affine} 与 \ref{examples-defos-lemma-lci-unobstructed} 即得结论。
\end{proof}




\section{光滑化}
\label{examples-defos-section-smoothing}

\noindent
设给定域 $k$ 上的有限型概形或代数空间 $X$。寻找有限型平坦态射
$Y\to\Spec(k[[t]])$ 往往很有用，其中一般纤维光滑，特殊纤维同构于
$X$。这样的对象称为 $X$ 的一个光滑化。本节将为具有孤立局部完全交
奇点的 $1$ 维分离 $X$ 找到光滑化。

\begin{lemma}
\label{examples-defos-lemma-criterion-smoothing}
设 $k$ 是域。置 $S=\Spec(k[[t]])$ 与
$S_n=\Spec(k[t]/(t^n))$。设 $Y\to S$ 是固有平坦概形态射，
其特殊纤维 $X$ 是 Cohen--Macaulay 且为 $d$ 维等维的。
记 $X_n=Y\times_SS_n$。若对某个 $n\geq1$，
$\Omega_{X_n/S_n}$ 的第 $d$ 个 Fitting 理想包含 $t^{n-1}$，
则 $Y\to S$ 的一般纤维光滑。
\end{lemma}

\begin{proof}
由态射进阶引理
\ref{more-morphisms-lemma-flat-finite-presentation-CM-open}
可知 $Y\to S$ 是 Cohen--Macaulay 态射。由态射引理
\ref{morphisms-lemma-flat-finite-presentation-CM-fibres-relative-dimension}
可知 $Y\to S$ 的相对维数为 $d$。由除子引理
\ref{divisors-lemma-d-fitting-ideal-omega-smooth}，
$\Omega_{Y/S}$ 的第 $d$ 个 Fitting 理想
$\mathcal{I}\subset\mathcal{O}_Y$ 定义了态射 $Y\to S$ 的奇异轨迹。
换言之，$V(\mathcal{I})\subset Y$ 是 $Y\to S$ 不光滑的点所成的闭子集。
由除子引理 \ref{divisors-lemma-base-change-and-fitting-ideal-omega}，
该 Fitting 理想的形成与基变换交换。根据假设，$t^{n-1}$ 是
$\mathcal{I}+t^n\mathcal{O}_Y$ 的截面。因此，对每个
$x\in X=V(t)\subset Y$，有 $t^{n-1}\in\mathcal{I}_x$；
这里 $\mathcal{I}_x$ 是在 $x$ 处的茎。这意味着在 $Y$ 中 $X$
的某个开邻域内有 $V(\mathcal{I})\subset V(t)$。由于 $Y\to S$
固有，可得所需的 $V(\mathcal{I})\subset V(t)$。
\end{proof}

\begin{lemma}
\label{examples-defos-lemma-jouanolou-type-thing}
设 $k$ 是域，$1\leq c\leq n$ 是整数。设
$f_1,\ldots,f_c\in k[x_1,\ldots x_n]$，并令
$a_{ij}$（$0\leq i\leq n$，$1\leq j\leq c$）为变量。考虑
% P11-E0069：冻结源的多项式环变量表在 \ldots 与 x_n 之间缺逗号。
$$
g_j = f_j + a_{0j} + a_{1j}x_1 + \ldots + a_{nj}x_n \in
k[a_{ij}][x_1, \ldots, x_n]
$$
以 $Y\subset\mathbf{A}^{n+c(n+1)}_k$ 表示由
$g_1,\ldots,g_c$ 定义的闭子概形。以
$\pi:Y\to\mathbf{A}^{c(n+1)}_k$ 表示到以 $a_{ij}$ 为变量的
仿射空间的投影。则 $\mathbf{A}^{c(n+1)}_k$ 中存在非空 Zariski
开子集，使得 $\pi$ 在其上光滑。
% P11-E0070：冻结源写 “a nonempty Zariski open of”，缺少 subset 等中心词。
\end{lemma}

\begin{proof}
回忆，$\pi$ 光滑的点所成的集合是开集，故其补集，即奇异轨迹，
是闭集。由 Chevalley 定理（采用态射引理
\ref{morphisms-lemma-chevalley} 的形式），奇异轨迹的像是可构造集。
因此，只要 $\mathbf{A}^{c(n+1)}_k$ 的泛点不在奇异轨迹的像中，
本引理便成立（例如可用拓扑引理
\ref{topology-lemma-generic-point-in-constructible}）。所以须证明：
不存在映到 $\mathbf{A}^{c(n+1)}_k$ 泛点的点 $y\in Y$，
使得 $\pi$ 在该点不光滑。考虑偏导数矩阵
$$
(\frac{\partial g_j}{\partial x_i}) =
(\frac{\partial f_j}{\partial x_i} + a_{ij})
$$
该矩阵在 $\kappa(y)$ 中的像必有秩 $<c$，否则 $\pi$ 将在 $y$
处光滑；参见环同态的光滑化第
\ref{smoothing-section-singular-ideal} 节的讨论。因此可找到不全为零的
$\lambda_1,\ldots,\lambda_c\in\kappa(y)$，使向量
$(\lambda_1,\ldots,\lambda_c)$ 位于该矩阵的核中。重新编号后，
可假设 $\lambda_1\not = 0$。除以 $\lambda_1$ 后，可假设该向量形如
$(1,\lambda_2,\ldots,\lambda_c)$。于是得到
% P11-E0071：冻结源下一公式首项含自由指标 j 且微分变量误为 x_1。
$$
a_{i1} = -
\frac{\partial f_j}{\partial x_1} -
\sum\nolimits_{j = 2, \ldots, c} \lambda_j(\frac{\partial f_j}{\partial x_i} + a_{ij})
$$
在 $\kappa(y)$ 中对 $i=1,\ldots,n$ 成立。此外，由于 $y\in Y$，还有
$$
a_{0j} = -f_j - a_{1j}x_1 - \ldots - a_{nj}x_n
$$
在 $\kappa(y)$ 中成立。这意味着 $\kappa(y)$ 中由 $a_{ij}$ 生成的
子域，包含于 $\kappa(y)$ 中由 $x_1,\ldots,x_n,\lambda_2,\ldots,\lambda_c$ 的像，
以及除 $a_{i1}$ 与 $a_{0j}$ 外的 $a_{ij}$ 所生成的子域。
计数可知，后者的超越次数至多为 $c(n+1)-1$。因此 $y$
不可能映到泛点，正如所需。
\end{proof}

\begin{lemma}
\label{examples-defos-lemma-smoothing-affine-lci}
设 $k$ 是域，$A$ 是 $k$ 上的全局完全交。存在有限型平坦环同态
$k[[t]]\to B$，满足 $B/tB\cong A$，并使 $B[1/t]$
在 $k((t))$ 上光滑。
\end{lemma}

\begin{proof}
按照代数定义 \ref{algebra-definition-lci-field}，写成
$A=k[x_1,\ldots,x_n]/(f_1,\ldots,f_c)$。我们将选取
$a_{ij}\in(t)\subset k[[t]]$，并置
$$
g_j = f_j + a_{0j} + a_{1j}x_1 + \ldots + a_{nj}x_n \in
k[[t]][x_1, \ldots, x_n]
$$
。随后取 $B=k[[t]][x_1,\ldots,x_n]/(g_1,\ldots,g_c)$。
我们断言 $k[[t]]\to B$ 在每个位于 $(t)$ 上方的素理想处都平坦。
事实上，元素 $f_1,\ldots,f_c$ 对 $k[x_1,\ldots,x_n]$ 中任何包含
$f_1,\ldots,f_c$ 的素理想 $\mathfrak p$，都在相应局部环中
构成正则序列（代数引理 \ref{algebra-lemma-lci}）。因此
$g_1,\ldots,g_c$ 局部地是正则序列的提升，可以应用代数引理
\ref{algebra-lemma-grothendieck-regular-sequence}。
在位于 $(0)\subset k[[t]]$ 上方的素理想处，平坦性自动成立，
因为 $k((t))=k[[t]]_{(0)}$ 是域。因此 $B$ 在 $k[[t]]$ 上平坦。

\medskip\noindent
只须再证明：适当选取 $a_{ij}$ 后，一般纤维 $B_{(0)}$
在 $k((t))$ 上光滑。为此，须证明可以选取 $a_{ij}$，使诱导态射
$$
(a_{ij}) : \Spec(k[[t]]) \longrightarrow \mathbf{A}^{c(n + 1)}_k
$$
映入引理 \ref{examples-defos-lemma-jouanolou-type-thing} 中的非空 Zariski 开集。
这是显然的，因为不存在关于 $a_{ij}$ 的非零多项式在
$(t)^{\oplus c(n+1)}$ 上恒为零（留作读者练习）。
\end{proof}

\begin{lemma}
\label{examples-defos-lemma-smoothing-artinian-lci}
设 $k$ 是域。设 $A$ 是有限维 $k$-代数，并且是 $k$ 上的局部完全交。
则存在有限平坦 $k[[t]]$-代数 $B$，满足 $B/tB\cong A$，
且 $B[1/t]$ 在 $k((t))$ 上为 \'etale。
\end{lemma}

\begin{proof}
由于 $A$ 是 Artin 环（代数引理
\ref{algebra-lemma-finite-dimensional-algebra}），可把 $A$ 写成局部
Artin 环的积（代数引理 \ref{algebra-lemma-artinian-finite-length}）。
因此只须证明 $A$ 为局部环的情形（这里使用了取主局部化保持
局部完全交性质；参见代数引理 \ref{algebra-lemma-localize-lci}）。
在这种情形下，$A$ 是全局完全交。考虑引理
\ref{examples-defos-lemma-smoothing-affine-lci} 中构造的代数 $B$。则
$k[[t]]\to B$ 在 $B$ 中唯一位于 $(t)$ 上方的素理想处拟有限
（代数定义 \ref{algebra-definition-quasi-finite}）。注意，$k[[t]]$
是 Hensel 局部环（代数引理 \ref{algebra-lemma-complete-henselian}）。
因此 $B=B'\times C$，其中 $B'$ 在 $k[[t]]$ 上有限，而 $C$
没有位于 $(t)$ 上方的素理想；参见代数引理
\ref{algebra-lemma-characterize-henselian}。于是 $B'$ 正是所求的环
（回忆，\'etale 等价于相对维数为 $0$ 的光滑）。
\end{proof}

\begin{lemma}
\label{examples-defos-lemma-smoothing-at-lci-point}
设 $k$ 是域，$A$ 是 $k$-代数。假设
\begin{enumerate}
\item $A$ 是本质有限型 $k$-局部环；
\item $A$ 是 $k$ 上的完全交
（代数定义 \ref{algebra-definition-lci-local-ring}）。
\end{enumerate}
置 $d=\dim(A)+\text{trdeg}_k(\kappa)$，其中 $\kappa$ 是 $A$
的剩余域。则存在整数 $n$ 以及本质有限型平坦环同态
$k[[t]]\to B$，满足 $B/tB\cong A$，并使 $t^n$ 属于
$\Omega_{B/k[[t]]}$ 的第 $d$ 个 Fitting 理想。
% P11-E0072：冻结源 “a flat, essentially of finite type ring map” 的修饰语次序不合英语语法。
\end{lemma}

\begin{proof}
由代数引理 \ref{algebra-lemma-lci-local}，可把 $A$ 写成 $k$
上的全局完全交 $P$ 在某个素理想 $\mathfrak p$ 处的局部化。
由代数引理 \ref{algebra-lemma-dimension-at-a-point-finite-type-field}，
注意到 $\dim(P)=d$。由引理 \ref{examples-defos-lemma-smoothing-affine-lci}，
可找到有限型平坦环同态 $k[[t]]\to Q$，使得 $P\cong Q/tQ$，
并使 $k((t))\to Q[1/t]$ 光滑。由该引理中 $Q$ 的构造可知，
$k[[t]]\to Q$ 是相对维数为 $d$ 的相对全局完全交；或者，
代数引理 \ref{algebra-lemma-syntomic} 表明 $Q$ 或 $Q$
的某个适当主局部化是这样的全局完全交。因此由除子引理
\ref{divisors-lemma-d-fitting-ideal-omega-smooth}，
$\Omega_{Q/k[[t]]}$ 的第 $d$ 个 Fitting 理想 $I\subset Q$
定义了 $\Spec(Q)\to\Spec(k[[t]])$ 的奇异轨迹。
于是对某个 $n$ 有 $t^n\in I$。令 $\mathfrak q\subset Q$
为 $\mathfrak p$ 的逆像，置 $B=Q_\mathfrak q$。引理得证。
\end{proof}

\begin{lemma}
\label{examples-defos-lemma-smoothing-proper-curve-isolated-lci}
设 $X$ 是域 $k$ 上的概形。假设
\begin{enumerate}
\item $X$ 在 $k$ 上固有；
\item $X$ 是 $k$ 上的局部完全交；
\item $X$ 的维数 $\leq1$；
\item $X\to\Spec(k)$ 除在有限多个点处外均光滑。
\end{enumerate}
则存在平坦射影态射 $Y\to\Spec(k[[t]])$，其一般纤维光滑，
特殊纤维同构于 $X$。
\end{lemma}

\begin{proof}
注意，$X$ 是 Cohen--Macaulay 的；参见代数引理
\ref{algebra-lemma-lci-CM}。因此 $X=X'\amalg X''$，其中
$\dim(X')=0$，而 $X''$ 是 $1$ 维等维的；参见态射引理
\ref{morphisms-lemma-flat-finite-presentation-CM-fibres-relative-dimension}。
由于 $X'$ 在 $k$ 上有限（代数簇引理
\ref{varieties-lemma-algebraic-scheme-dim-0}），由引理
\ref{examples-defos-lemma-smoothing-artinian-lci} 可找到
$Y'\to\Spec(k[[t]])$，其特殊纤维为 $X'$，一般纤维光滑。
因此只须对 $X''$ 证明本引理。以 $X''$ 替换 $X$ 后，
$X$ 是 Cohen--Macaulay 且为 $1$ 维等维的。

\medskip\noindent
我们将对情形 $\Lambda=k\to k$ 使用形变理论。令
$p_1,\ldots,p_r\in X$ 为 $X$ 的闭奇点，即 $X\to\Spec(k)$
不光滑的点。对每个 $i$，选取整数 $n_i$ 和本质有限型平坦环同态
$$
k[[t]] \longrightarrow B_i
$$
，满足 $B_i/tB_i\cong\mathcal{O}_{X,p_i}$，并使 $t^{n_i}$ 属于
$\Omega_{B_i/k[[t]]}$ 的第 $1$ 个 Fitting 理想。由引理
\ref{examples-defos-lemma-smoothing-at-lci-point}，这可以做到。注意，系统
$(B_i/t^nB_i)$ 定义了 $\Deformationcategory_{\mathcal{O}_{X,p_i}}$
在 $k[[t]]$ 上的形式对象。由引理 \ref{examples-defos-lemma-curve-isolated}，映射
$$
\Deformationcategory_X
\longrightarrow
\prod\nolimits_{i = 1, \ldots, r} \Deformationcategory_{\mathcal{O}_{X, p_i}}
$$
是形变范畴之间的光滑映射。因此由形式形变理论引理
\ref{formal-defos-lemma-smooth-morphism-essentially-surjective}
，$\Deformationcategory_X$ 中存在形式对象 $(X_n)$，
经上述箭头映到形式对象 $\prod_i(B_i/t^n)$。由空间态射进阶引理
\ref{spaces-more-morphisms-lemma-formal-algebraic-space-proper-reldim-1}
，存在 $k[[t]]$ 上的射影概形 $Y$ 以及相容同构
$Y\times_{\Spec(k[[t]])}\Spec(k[t]/(t^n))\cong X_n$。
由态射进阶引理
\ref{more-morphisms-lemma-check-flatness-on-infinitesimal-nbhds}
可知 $Y\to\Spec(k[[t]])$ 平坦。由于 $X$ 是 Cohen--Macaulay
且为 $1$ 维等维的，可应用引理 \ref{examples-defos-lemma-criterion-smoothing}
检验 $Y$ 的一般纤维光滑\footnote{警告：一般而言，$Y$ 在点
$p_i$ 处的局部环同构于 $B_i$ 这一说法{\bf 不}成立。
我们只知道两边都模去 $t^n$ 后该说法成立！}。
选取严格大于上述整数 $n_i$ 之最大值的 $n$。若能证明
$t^{n-1}$ 属于 $\Omega_{X_n/S_n}$ 的第一个 Fitting 理想，
其中 $S_n=\Spec(k[t]/(t^n))$，则证明完成。为此，只须在 $X_n$
每个闭点 $p$ 处的局部环中证明这一点。然而，若 $p$ 对应于
$X\to\Spec(k)$ 的光滑点，则 $\Omega_{X_n/S_n,p}$ 是秩为 $1$
的自由模，第一个 Fitting 理想等于该局部环。若对某个 $i$
有 $p=p_i$，则
$$
\Omega_{X_n/S_n, p_i} =
\Omega_{(B_i/t^nB_i)/(k[t]/(t^n))} =
\Omega_{B_i/k[[t]]}/t^n\Omega_{B_i/k[[t]]}
$$
由于取 Fitting 理想与基变换交换（我们已经使用过这一点；
在这里的代数情形中，它由代数进阶引理
\ref{more-algebra-lemma-fitting-ideal-basics} 得到），并且
$n-1\geq n_i$，可知 $t^{n-1}$ 属于该模作为
$B_i/t^nB_i$-模的 Fitting 理想，正如所需。
% P11-E0073：冻结源将 “we already used” 误写为 “with already used”。
\end{proof}

\begin{lemma}
\label{examples-defos-lemma-smoothing-curve-isolated-lci}
设 $k$ 是域，$X$ 是 $k$ 上的概形。假设
\begin{enumerate}
\item $X$ 是分离的、在 $k$ 上有限型，且 $\dim(X)\leq1$；
\item $X$ 是 $k$ 上的局部完全交；
\item $X\to\Spec(k)$ 除在有限多个点处外均光滑。
\end{enumerate}
则存在有限型平坦分离态射 $Y\to\Spec(k[[t]])$，其一般纤维光滑，
特殊纤维同构于 $X$。
\end{lemma}

\begin{proof}
若 $X$ 既约，则可像代数簇引理
\ref{varieties-lemma-reduced-dim-1-projective-completion} 中那样选取嵌入
$X\subset\overline{X}$。写成
$X=\overline{X}\setminus\{x_1,\ldots,x_n\}$，可见
$\mathcal{O}_{\overline{X},x_i}$ 是离散赋值环，因而特别是局部完全交
（代数定义 \ref{algebra-definition-lci-local-ring}）。因此
$\overline{X}$ 是 $k$ 上的局部完全交：在开集 $X$ 上这一点成立，
而在点 $x_i$ 处由代数引理 \ref{algebra-lemma-lci-local} 成立。
故可应用引理 \ref{examples-defos-lemma-smoothing-proper-curve-isolated-lci}，找到
射影平坦态射 $\overline{Y}\to\Spec(k[[t]])$，其一般纤维光滑，
特殊纤维为 $\overline{X}$。随后从 $\overline{Y}$ 中移去
$x_1,\ldots,x_n$，得到 $Y$。

\medskip\noindent
在一般情形下，写成 $X=X'\amalg X''$，其中
$\dim(X')=0$，而 $X''$ 是 $1$ 维等维的。
% P11-E0074：冻结源在 where 后多写了 with。
此时 $X''$ 既约，第一段适用于它。另一方面，可以像引理
\ref{examples-defos-lemma-smoothing-proper-curve-isolated-lci} 的证明中那样处理
$X'$。略去一些细节。
\end{proof}
% Generated mechanically from the frozen algebraic.tex target.
% Do not edit this generated file; edit the bound locale source instead.

% OK, start here.
%

\setcounter{chapter}{93}
\chapter{代数叠}



\phantomsection
\label{algebraic-section-phantom}


\section{引言}
\label{algebraic-section-introduction}

\noindent
我们在此定义代数叠，并给出一些非常初步的观察。一般原则是完全不预设
任何分离性条件，只在使引理、命题和定理成立或可证时加入必要条件。
因此，这里讨论的概念与文献中其他地方所用的概念略有不同，例如
\cite{LM-B}。

\medskip\noindent
本章并非代数叠的入门介绍。关于代数叠的非形式化讨论，请参阅
《代数叠导论》第
\ref{stacks-introduction-section-introduction} 节。


\section{约定}
\label{algebraic-section-conventions}

\noindent
本章采用的约定与《代数空间》一章相同。为方便起见，我们在此重述。

\medskip\noindent
我们在适当的大 fppf 位点 $\Sch_{fppf}$ 中工作，参见
《概形上的拓扑》定义 \ref{topologies-definition-big-fppf-site}。
因此，除非另有明确说明，所有概形都是 $\Sch_{fppf}$ 的对象。
第 \ref{algebraic-section-change-big-site} 节将讨论更换大 fppf 位点时
会发生哪些变化。

\medskip\noindent
我们始终相对于 $\Sch_{fppf}$ 中的一个基 $S$ 工作，并使用大 fppf
位点 $(\Sch/S)_{fppf}$，参见《概形上的拓扑》定义
\ref{topologies-definition-big-small-fppf}。取
$S = \Spec(\mathbf{Z})$ 即可恢复绝对情形。

\medskip\noindent
若 $U,T$ 是 $S$ 上的概形，则用 $U(T)$ 表示相应的
{\it 在} $S$ 上的 $T$-值点集。用公式写即
$U(T) = \Mor_S(T, U)$。

\medskip\noindent
注意，任意 fpqc 覆盖都是泛有效满态射，参见
《下降》引理 \ref{descent-lemma-fpqc-universal-effective-epimorphisms}。
因此，$\Sch_{fppf}$ 上的拓扑弱于典范拓扑，而所有可表预层都是层。








\section{记号}
\label{algebraic-section-notation}

\noindent
我们用字母 $S,T,U,V,X,Y$ 表示概形，用字母
$\mathcal{X}, \mathcal{Y}, \mathcal{Z}$ 表示
$(\Sch/S)_{fppf}$ 上的范畴（纤维化范畴、以群胚为纤维的
纤维化范畴、叠等等）。我们用小写字母 $f$、$g$ 表示
$(\Sch/S)_{fppf}$ 上的函子，例如
$f : \mathcal{X} \to \mathcal{Y}$。
我们用大写字母 $F$、$G$、$H$ 表示 $S$ 上的代数空间，更一般地，
表示 $(\Sch/S)_{fppf}$ 上的集合预层。（在后续各章中，
我们也将重新使用 $X$、$Y$ 等字母表示代数空间。）

\medskip\noindent
作出这些选择，是为了在奠定基础时清楚地区分本章中的各种对象。









\section{可表的群胚纤维化范畴}
\label{algebraic-section-representable}

\noindent
设 $S$ 是 $\Sch_{fppf}$ 中的概形。本章的基本研究对象是
群胚纤维化范畴
$p : \mathcal{X} \to (\Sch/S)_{fppf}$，参见《范畴》定义
\ref{categories-definition-fibred-groupoids}。我们常简称
“设 $\mathcal{X}$ 是 $(\Sch/S)_{fppf}$ 上的群胚纤维化范畴”
来表示这一情形。$(\Sch/S)_{fppf}$ 上群胚纤维化范畴之间的
$1$-态射 $\mathcal{X} \to \mathcal{Y}$，是
$(\Sch/S)_{fppf}$ 上群胚纤维化范畴所成 $2$-范畴中的
$1$-态射，参见《范畴》定义
\ref{categories-definition-categories-fibred-in-groupoids-over-C}。
它就是 $(\Sch/S)_{fppf}$ 上的函子
$\mathcal{X} \to \mathcal{Y}$。回忆，这实际上是一个
$(2, 1)$-范畴，并且所有 $2$-纤维积都存在。

\medskip\noindent
设 $\mathcal{X}$ 是 $(\Sch/S)_{fppf}$ 上的群胚纤维化范畴。
回忆，若存在概形 $U \in \Ob((\Sch/S)_{fppf})$ 以及范畴等价
$$
j : \mathcal{X} \longrightarrow (\Sch/U)_{fppf}
$$
（它们都是 $(\Sch/S)_{fppf}$ 上的范畴），则称 $\mathcal{X}$
{\it 可表}，参见《范畴》定义
\ref{categories-definition-representable-fibred-category}。
为区别于 $\mathcal{X}$ 可由代数空间表示的情形（见下文），
我们有时称 $\mathcal{X}$ {\it 可由概形表示}。

\medskip\noindent
若 $\mathcal{X},\mathcal{Y}$ 是群胚纤维化范畴，并分别由
$U,V$ 表示，则有
\begin{equation}
\label{algebraic-equation-morphisms-schemes}
\Mor_{\textit{Cat}/(\Sch/S)_{fppf}}(\mathcal{X}, \mathcal{Y})
\Big/
2\text{-同构}
=
\Mor_{\Sch/S}(U, V)
\end{equation}
参见《范畴》引理
\ref{categories-lemma-morphisms-representable-fibred-categories}。
更确切地，任意 $1$-态射 $\mathcal{X} \to \mathcal{Y}$
都给出态射 $U \to V$。反过来，给定 $S$ 上的概形态射
$U \to V$，存在一个给出 $U \to V$ 的 $1$-态射
$\phi : \mathcal{X} \to \mathcal{Y}$，且它在唯一
$2$-同构意义下唯一。






\section{2-Yoneda 引理}
\label{algebraic-section-2-yoneda}

\noindent
设 $U \in \Ob((\Sch/S)_{fppf})$，并设 $\mathcal{X}$ 是
$(\Sch/S)_{fppf}$ 上的群胚纤维化范畴。我们将频繁使用
$2$-Yoneda 引理，参见《范畴》引理
\ref{categories-lemma-yoneda-2category}。严格来说，它断言存在范畴等价
$$
\Mor_{\textit{Cat}/(\Sch/S)_{fppf}}(
(\Sch/U)_{fppf}, \mathcal{X})
\longrightarrow
\mathcal{X}_U, \quad
f \longmapsto f(U/U).
$$
也就是说，$1$-态射 $(\Sch/U)_{fppf} \to \mathcal{X}$
对应于纤维范畴 $\mathcal{X}_U$ 的对象 $x$。具体地，给定
$1$-态射 $f : (\Sch/U)_{fppf} \to \mathcal{X}$，可得到对象
$x = f(U/U) \in \Ob(\mathcal{X}_U)$。反过来，按《范畴》定义
\ref{categories-definition-pullback-functor-fibred-category}
为 $\mathcal{X}$ 选定拉回，再给定 $\mathcal{X}_U$ 的对象 $x$，
便得到函子 $(\Sch/U)_{fppf} \to \mathcal{X}$；它在对象上的规则为
$$
(\varphi : V \to U) \longmapsto \varphi^*x
$$
。沿用记号，我们仍用
$x : (\Sch/U)_{fppf} \to \mathcal{X}$ 表示该函子。
它确实满足 $x(U/U) = x$；此外，对任意满足 $f(U/U) = x$
的其他函子 $f$，存在唯一的 $2$-同构 $x \to f$。换言之，
对象 $x$ 在唯一 $2$-同构意义下唯一确定函子 $x$。

\medskip\noindent
下文将直接使用这一事实，不再另行说明。





\section{群胚纤维化范畴的可表态射}
\label{algebraic-section-representable-morphism}


\noindent
设 $\mathcal{X}$、$\mathcal{Y}$ 是 $(\Sch/S)_{fppf}$ 上的
群胚纤维化范畴。设 $f : \mathcal{X} \to \mathcal{Y}$ 是
{\it 可表 $1$-态射}，参见《范畴》定义
\ref{categories-definition-representable-map-categories-fibred-in-groupoids}。
这意味着，对每个 $U \in \Ob((\Sch/S)_{fppf})$ 以及任意
$y \in \Ob(\mathcal{Y}_U)$，$2$-纤维积
$(\Sch/U)_{fppf} \times_{y, \mathcal{Y}} \mathcal{X}$
作为 $(\Sch/U)_{fppf}$ 上的群胚纤维化范畴是可表的。
选取表示对象 $f_y : V_y \to U$ 以及等价
$$
(\Sch/V_y)_{fppf}
\longrightarrow
(\Sch/U)_{fppf} \times_{y, \mathcal{Y}} \mathcal{X}.
$$
态射 $f_y$ 对应于投影
$(\Sch/V_y)_{fppf} \to
(\Sch/U)_{fppf} \times_\mathcal{Y} \mathcal{Y}
\to (\Sch/U)_{fppf}$
；参见第 \ref{algebraic-section-representable} 节公式
(\ref{algebraic-equation-morphisms-schemes})。我们用下图表示这一情形：
\begin{equation}
\label{algebraic-equation-representable}
\vcenter{
\xymatrix{
V_y \ar@{~>}[r] \ar[d]_{f_y} &
(\Sch/V_y)_{fppf} \ar[d] \ar[r] &
\mathcal{X} \ar[d]^f \\
U \ar@{~>}[r] &
(\Sch/U)_{fppf} \ar[r]^-y &
\mathcal{Y}
}
}
\end{equation}
其中波浪箭头表示 $2$-Yoneda 嵌入。下面给出关于这一概念的若干引理，
它们在极其一般的情形下成立（即对任意具有纤维积的基范畴上的
群胚纤维化范畴都成立）。

\begin{lemma}
\label{algebraic-lemma-morphism-schemes-gives-representable-transformation}
设 $f : X \to Y$ 是 $(\Sch/S)_{fppf}$ 中的态射。则 $f$ 诱导的
$1$-态射
$$
(\Sch/X)_{fppf} \longrightarrow (\Sch/Y)_{fppf}
$$
是可表 $1$-态射。
\end{lemma}

\begin{proof}
这是形式的，只用到范畴 $(\Sch/S)_{fppf}$ 具有纤维积这一事实。
\end{proof}

\begin{lemma}
\label{algebraic-lemma-representable-morphism-equivalent}
设 $S$ 是 $\Sch_{fppf}$ 的对象。考虑
$(\Sch/S)_{fppf}$ 上群胚纤维化范畴的 $1$-态射所成的
$2$-交换图
$$
\xymatrix{
\mathcal{X}' \ar[r] \ar[d]_{f'} & \mathcal{X} \ar[d]^f \\
\mathcal{Y}' \ar[r] & \mathcal{Y}
}
$$
。假设水平箭头都是等价。则 $f$ 可表当且仅当 $f'$ 可表。
\end{lemma}

\begin{proof}
略。
\end{proof}

\begin{lemma}
\label{algebraic-lemma-composition-representable-transformations}
设 $S$ 是 $\Sch_{fppf}$ 中的概形。设
$\mathcal{X},\mathcal{Y},\mathcal{Z}$ 是
$(\Sch/S)_{fppf}$ 上的群胚纤维化范畴，并设
$f : \mathcal{X} \to \mathcal{Y}$、
$g : \mathcal{Y} \to \mathcal{Z}$ 是可表 $1$-态射。则
$$
g \circ f : \mathcal{X} \longrightarrow \mathcal{Z}
$$
是可表 $1$-态射。
\end{lemma}

\begin{proof}
这完全是形式的，并在任意范畴中成立。
\end{proof}

\begin{lemma}
\label{algebraic-lemma-base-change-representable-transformations}
设 $S$ 是 $\Sch_{fppf}$ 中的概形。设
$\mathcal{X},\mathcal{Y},\mathcal{Z}$ 是
$(\Sch/S)_{fppf}$ 上的群胚纤维化范畴。设
$f : \mathcal{X} \to \mathcal{Y}$ 是可表 $1$-态射，并设
$g : \mathcal{Z} \to \mathcal{Y}$ 是任意 $1$-态射。
考虑纤维积图
$$
\xymatrix{
\mathcal{Z} \times_{g, \mathcal{Y}, f} \mathcal{X} \ar[r]_-{g'} \ar[d]_{f'} &
\mathcal{X} \ar[d]^f \\
\mathcal{Z} \ar[r]^g & \mathcal{Y}
}
$$
则基变换 $f'$ 是可表 $1$-态射。
\end{lemma}

\begin{proof}
这完全是形式的，并在任意范畴中成立。
\end{proof}

\begin{lemma}
\label{algebraic-lemma-product-representable-transformations}
设 $S$ 是 $\Sch_{fppf}$ 中的概形。设
$\mathcal{X}_i,\mathcal{Y}_i$ 是 $(\Sch/S)_{fppf}$ 上的
群胚纤维化范畴，$i = 1, 2$。设
$f_i : \mathcal{X}_i \to \mathcal{Y}_i$（$i = 1, 2$）
是可表 $1$-态射。则
$$
f_1 \times f_2 :
\mathcal{X}_1 \times \mathcal{X}_2
\longrightarrow
\mathcal{Y}_1 \times \mathcal{Y}_2
$$
是可表 $1$-态射。
\end{lemma}

\begin{proof}
把 $f_1 \times f_2$ 写成复合
$\mathcal{X}_1 \times \mathcal{X}_2 \to
\mathcal{Y}_1 \times \mathcal{X}_2 \to
\mathcal{Y}_1 \times \mathcal{Y}_2$。
第一个箭头是 $f_1$ 沿映射
$\mathcal{Y}_1 \times \mathcal{X}_2 \to \mathcal{Y}_1$ 的基变换，
第二个箭头是 $f_2$ 沿映射
$\mathcal{Y}_1 \times \mathcal{Y}_2 \to \mathcal{Y}_2$ 的基变换。
因此，本引理是引理
\ref{algebraic-lemma-composition-representable-transformations} 与
\ref{algebraic-lemma-base-change-representable-transformations}
的形式推论。
\end{proof}



\section{分裂群胚纤维化范畴}
\label{algebraic-section-split}

\noindent
设 $S$ 是 $\Sch_{fppf}$ 中的概形。回忆，给定一个“群胚预层”
$$
F : (\Sch/S)_{fppf}^{opp} \longrightarrow \textit{Groupoids}
$$
，可得到 $(\Sch/S)_{fppf}$ 上的群胚纤维化范畴
$\mathcal{S}_F$，参见《范畴》例
\ref{categories-example-functor-groupoids}。与其中某个范畴同构（！）的
群胚纤维化范畴称为{\it 分裂群胚纤维化范畴}。
任意群胚纤维化范畴都等价于一个分裂的群胚纤维化范畴。

\medskip\noindent
若 $F$ 是集合预层，则 $\mathcal{S}_F$ 以集合为纤维，参见
《范畴》定义 \ref{categories-definition-category-fibred-sets}
与《范畴》例 \ref{categories-example-presheaf}。
规则 $F \mapsto \mathcal{S}_F$ 在某种意义下于预层上全忠实，
参见《范畴》引理 \ref{categories-lemma-2-category-fibred-sets}。
若 $F,G$ 是预层，则
$$
\mathcal{S}_{F \times G}
=
\mathcal{S}_F \times_{(\Sch/S)_{fppf}} \mathcal{S}_G
$$
；若 $F \to H$ 与 $G \to H$ 是集合预层的映射，则
$$
\mathcal{S}_{F \times_H G} =
\mathcal{S}_F \times_{\mathcal{S}_H} \mathcal{S}_G
$$
，其中右端都是 $2$-纤维积。这由定义立即可得，因为
$\mathcal{S}_F, \mathcal{S}_G, \mathcal{S}_H$ 的各纤维范畴
只有恒等态射。

\medskip\noindent
更特殊的情形是 $F = h_X$ 为可表预层。此时
$\mathcal{S}_{h_X} = (\Sch/X)_{fppf}$，参见《范畴》例
\ref{categories-example-fibred-category-from-functor-of-points}。

\medskip\noindent
下文将直接使用记号 $\mathcal{S}_F$，不再另行说明。




\section{可由代数空间表示的群胚纤维化范畴}
\label{algebraic-section-representable-by-algebraic-spaces}

\noindent
比可表性稍弱的概念是可由代数空间表示，本节对此加以讨论。
若我们使用代数空间的某个范畴 $\textit{Spaces}_{fppf}$
而非范畴 $\Sch_{fppf}$，或许可以避免这番讨论。然而，在我们看来，
把概形范畴视为定义代数叠各纤维范畴所用的“测试对象”的自然集合，
是很自然的。

\medskip\noindent
仿照《范畴》定义
\ref{categories-definition-representable-fibred-category}，
作如下定义。

\begin{definition}
\label{algebraic-definition-representable-by-algebraic-space}
设 $S$ 是 $\Sch_{fppf}$ 中的概形。若存在 $S$ 上的代数空间
$F$ 以及 $(\Sch/S)_{fppf}$ 上范畴的等价
$j : \mathcal{X} \to \mathcal{S}_F$，则称群胚纤维化范畴
$p : \mathcal{X} \to (\Sch/S)_{fppf}$
{\it 可由 $S$ 上的代数空间表示}。
\end{definition}

\noindent
遇到这种情形时，我们继续沿用记号而省略等价 $j$。
由上文形式地可知，若 $\mathcal{X}$（由概形）可表，则它可由
代数空间表示。下面是《范畴》引理
\ref{categories-lemma-characterize-representable-fibred-category}
的对应结论。

\begin{lemma}
\label{algebraic-lemma-characterize-representable-by-space}
设 $S$ 是 $\Sch_{fppf}$ 中的概形，并设
$p : \mathcal{X} \to (\Sch/S)_{fppf}$ 是群胚纤维化范畴。
则 $\mathcal{X}$ 可由 $S$ 上的代数空间表示，当且仅当满足以下条件：
\begin{enumerate}
\item $\mathcal{X}$ 以集合胚为纤维\footnote{这意味着它以群胚为纤维，
且各纤维范畴中的对象没有非平凡自同构；参见《范畴》定义
\ref{categories-definition-category-fibred-setoids}。}；
\item 预层 $U \mapsto \Ob(\mathcal{X}_U)/\!\!\cong$ 是代数空间。
\end{enumerate}
\end{lemma}

\begin{proof}
略；但可参见《范畴》引理
\ref{categories-lemma-characterize-representable-fibred-category}。
\end{proof}

\noindent
若 $\mathcal{X}, \mathcal{Y}$ 以群胚为纤维，并分别可由
$S$ 上的代数空间 $F,G$ 表示，则有
\begin{equation}
\label{algebraic-equation-morphisms-spaces}
\Mor_{\textit{Cat}/(\Sch/S)_{fppf}}(\mathcal{X}, \mathcal{Y})
\Big/
2\text{-同构}
=
\Mor_{\Sch/S}(F, G)
\end{equation}
参见《范畴》引理
\ref{categories-lemma-2-category-fibred-setoids}。更确切地，任意
$1$-态射 $\mathcal{X} \to \mathcal{Y}$ 都给出态射 $F \to G$。
反过来，给定 $S$ 上的层态射 $F \to G$，存在一个给出
$F \to G$ 的 $1$-态射 $\phi : \mathcal{X} \to \mathcal{Y}$，
且它在唯一 $2$-同构意义下唯一。



\section{可由代数空间表示的态射}
\label{algebraic-section-morphisms-representable-by-algebraic-spaces}

\noindent
仿照《范畴》定义
\ref{categories-definition-representable-map-categories-fibred-in-groupoids}，
作如下定义。

\begin{definition}
\label{algebraic-definition-representable-by-algebraic-spaces}
设 $S$ 是 $\Sch_{fppf}$ 中的概形。若对任意
$U \in \Ob((\Sch/S)_{fppf})$ 及任意
$y : (\Sch/U)_{fppf} \to \mathcal{Y}$，群胚纤维化范畴
$$
(\Sch/U)_{fppf} \times_{y, \mathcal{Y}} \mathcal{X}
$$
作为 $(\Sch/U)_{fppf}$ 上的范畴可由 $U$ 上的代数空间表示，
则称 $(\Sch/S)_{fppf}$ 上群胚纤维化范畴的 $1$-态射
$f : \mathcal{X} \to \mathcal{Y}$ {\it 可由代数空间表示}。
\end{definition}

\noindent
选取表示 $(\Sch/U)_{fppf} \times_{y, \mathcal{Y}} \mathcal{X}$
的 $U$ 上代数空间 $F_y$。可把 $F_y$ 看成 $S$ 上的代数空间，
它配备一个典范的 $S$ 上态射 $f_y : F_y \to U$；参见
《代数空间》第 \ref{spaces-section-change-base-scheme} 节。
相应图表为
\begin{equation}
\label{algebraic-equation-representable-by-algebraic-spaces}
\vcenter{
\xymatrix{
F_y \ar[d]_{f_y} &
(\Sch/U)_{fppf} \times_{y, \mathcal{Y}} \mathcal{X}
\ar@{~>}[l] \ar[d]_{\text{pr}_0} \ar[r]_-{\text{pr}_1} &
\mathcal{X} \ar[d]^f \\
U &
(\Sch/U)_{fppf} \ar@{~>}[l] \ar[r]^-y &
\mathcal{Y}
}
}
\end{equation}
其中波浪箭头表示如下构造：把以集合胚为纤维的叠送到其对象同构类所成的
相应层。右侧方块是 $2$-交换的，并且是 $2$-纤维积方块。

\medskip\noindent
下面是《范畴》引理
\ref{categories-lemma-criterion-representable-map-stack-in-groupoids}
的对应结论。

\begin{lemma}
\label{algebraic-lemma-criterion-map-representable-spaces-fibred-in-groupoids}
设 $S$ 是 $\Sch_{fppf}$ 中的概形。设
$f : \mathcal{X} \to \mathcal{Y}$ 是 $(\Sch/S)_{fppf}$ 上
群胚纤维化范畴的 $1$-态射。则 $f$ 可由代数空间表示的充要条件如下：
\begin{enumerate}
\item 对每个概形 $U/S$，纤维范畴之间的函子
$f_U : \mathcal{X}_U \longrightarrow \mathcal{Y}_U$ 是忠实的；
\item 对每个 $U$ 以及每个 $y \in \Ob(\mathcal{Y}_U)$，预层
$$
(h : V \to U)
\longmapsto
\{(x, \phi) \mid x \in \Ob(\mathcal{X}_V), \phi : h^*y \to f(x)\}/\cong
$$
是 $U$ 上的代数空间。
\end{enumerate}
这里我们为 $\mathcal{Y}$ 选定了拉回。
\end{lemma}

\begin{proof}
把《范畴》引理 \ref{categories-lemma-identify-fibre-product}
中对 $2$-纤维积
$(\Sch/U)_{fppf} \times_{y, \mathcal{Y}} \mathcal{X}$
的纤维范畴的描述，与引理
\ref{algebraic-lemma-characterize-representable-by-space} 结合即可。
\end{proof}

\noindent
下面给出关于这一概念的若干引理；它们在非常一般的情形下成立。

\begin{lemma}
\label{algebraic-lemma-representable-by-spaces-morphism-equivalent}
设 $S$ 是 $\Sch_{fppf}$ 的对象。考虑
$(\Sch/S)_{fppf}$ 上群胚纤维化范畴的 $1$-态射所成的
$2$-交换图
$$
\xymatrix{
\mathcal{X}' \ar[r] \ar[d]_{f'} & \mathcal{X} \ar[d]^f \\
\mathcal{Y}' \ar[r] & \mathcal{Y}
}
$$
。假设水平箭头都是等价。则 $f$ 可由代数空间表示，当且仅当
$f'$ 可由代数空间表示。
\end{lemma}

\begin{proof}
略。
\end{proof}

\begin{lemma}
\label{algebraic-lemma-morphism-spaces-gives-representable-by-spaces}
设 $S$ 是 $\Sch_{fppf}$ 的对象。设
$f : \mathcal{X} \to \mathcal{Y}$ 是 $S$ 上群胚纤维化范畴的
$1$-态射。若 $\mathcal{X}$ 与 $\mathcal{Y}$ 都可由
$S$ 上的代数空间表示，则 $1$-态射 $f$ 可由代数空间表示。
\end{lemma}

\begin{proof}
略。这只用到 $S$ 上代数空间的范畴具有纤维积这一事实；
参见《代数空间》引理 \ref{spaces-lemma-fibre-product-spaces}。
\end{proof}

\begin{lemma}
\label{algebraic-lemma-map-presheaves-representable-by-algebraic-spaces}
设 $S$ 是 $\Sch_{fppf}$ 的对象。设 $a : F \to G$ 是
$(\Sch/S)_{fppf}$ 上集合预层的映射，并用
$a' : \mathcal{S}_F  \to \mathcal{S}_G$
表示相应的集合纤维化范畴映射。则 $a$ 可由代数空间表示
（参见《自举》定义
\ref{bootstrap-definition-morphism-representable-by-spaces}），
当且仅当 $a'$ 可由代数空间表示。
\end{lemma}

\begin{proof}
略。
\end{proof}

\begin{lemma}
\label{algebraic-lemma-map-fibred-setoids-representable-algebraic-spaces}
设 $S$ 是 $\Sch_{fppf}$ 的对象。设
$f : \mathcal{X} \to \mathcal{Y}$ 是 $(\Sch/S)_{fppf}$ 上
集合胚纤维化范畴的 $1$-态射。令 $F$、$G$ 分别为如下预层：
它把 $T$ 分别送到 $\mathcal{X}_T$、$\mathcal{Y}_T$
的对象同构类集合。令 $a : F \to G$ 为与 $f$ 对应的预层映射。
则 $a$ 可由代数空间表示（参见《自举》定义
\ref{bootstrap-definition-morphism-representable-by-spaces}），
当且仅当 $f$ 可由代数空间表示。
\end{lemma}

\begin{proof}
略。提示：结合引理
\ref{algebraic-lemma-representable-by-spaces-morphism-equivalent} 与
\ref{algebraic-lemma-map-presheaves-representable-by-algebraic-spaces}。
\end{proof}

\begin{lemma}
\label{algebraic-lemma-base-change-representable-by-spaces}
设 $S$ 是 $\Sch_{fppf}$ 中的概形。设
$\mathcal{X}, \mathcal{Y}, \mathcal{Z}$ 是
$(\Sch/S)_{fppf}$ 上的群胚纤维化范畴。设
$f : \mathcal{X} \to \mathcal{Y}$ 是可由代数空间表示的
$1$-态射，并设 $g : \mathcal{Z} \to \mathcal{Y}$
是任意 $1$-态射。考虑纤维积图
$$
\xymatrix{
\mathcal{Z} \times_{g, \mathcal{Y}, f} \mathcal{X} \ar[r]_-{g'} \ar[d]_{f'} &
\mathcal{X} \ar[d]^f \\
\mathcal{Z} \ar[r]^g & \mathcal{Y}
}
$$
则基变换 $f'$ 是可由代数空间表示的 $1$-态射。
\end{lemma}

\begin{proof}
这是形式的。
\end{proof}

\begin{lemma}
\label{algebraic-lemma-base-change-by-space-representable-by-space}
设 $S$ 是 $\Sch_{fppf}$ 中的概形。设
$\mathcal{X}, \mathcal{Y}, \mathcal{Z}$ 是
$(\Sch/S)_{fppf}$ 上的群胚纤维化范畴。设
$f : \mathcal{X} \to \mathcal{Y}$、
$g : \mathcal{Z} \to \mathcal{Y}$ 是 $1$-态射。假设
\begin{enumerate}
\item $f$ 可由代数空间表示；
\item $\mathcal{Z}$ 可由 $S$ 上的代数空间表示。
\end{enumerate}
则 $2$-纤维积
$\mathcal{Z} \times_{g, \mathcal{Y}, f} \mathcal{X}$
可由代数空间表示。
\end{lemma}

\begin{proof}
这是《自举》引理
\ref{bootstrap-lemma-representable-by-spaces-over-space}
的改述。首先注意，
$\mathcal{Z} \times_{g, \mathcal{Y}, f} \mathcal{X}$
在 $(\Sch/S)_{fppf}$ 上以集合胚为纤维。因此，它等价于
$\mathcal{S}_F$，其中 $F$ 是 $(\Sch/S)_{fppf}$ 上的某个预层；
参见《范畴》引理 \ref{categories-lemma-setoid-fibres}。
再设 $G$ 是表示 $\mathcal{Z}$ 的代数空间。$1$-态射
$\mathcal{Z} \times_{g, \mathcal{Y}, f} \mathcal{X} \to \mathcal{Z}$
由引理 \ref{algebraic-lemma-base-change-representable-by-spaces}
可由代数空间表示。又由《范畴》引理
\ref{categories-lemma-2-category-fibred-setoids}，态射
$\mathcal{Z} \times_{g, \mathcal{Y}, f} \mathcal{X} \to \mathcal{Z}$
对应于态射 $F \to G$。由引理
\ref{algebraic-lemma-map-fibred-setoids-representable-algebraic-spaces}，
$F \to G$ 可由代数空间表示。因此，《自举》引理
\ref{bootstrap-lemma-representable-by-spaces-over-space}
说明 $F$ 是代数空间，如所需。
\end{proof}

\noindent
设 $S$、$\mathcal{X}$、$\mathcal{Y}$、$\mathcal{Z}$、$f$、$g$
如引理 \ref{algebraic-lemma-base-change-by-space-representable-by-space} 中所述。
设 $F$ 与 $G$ 是 $S$ 上的代数空间，其中 $F$ 表示
$\mathcal{Z} \times_{g, \mathcal{Y}, f} \mathcal{X}$，而 $G$ 表示
$\mathcal{Z}$。$1$-态射
$f' : \mathcal{Z} \times_{g, \mathcal{Y}, f} \mathcal{X} \to \mathcal{Z}$
由 (\ref{algebraic-equation-morphisms-spaces}) 对应于代数空间的态射
$f' : F \to G$。因此有图表
\begin{equation}
\label{algebraic-equation-representable-by-algebraic-spaces-on-space}
\vcenter{
\xymatrix{
F \ar[d]_{f'} &
\mathcal{Z} \times_{g, \mathcal{Y}, f} \mathcal{X}
\ar@{~>}[l] \ar[d] \ar[r] &
\mathcal{X} \ar[d]^f \\
G &
\mathcal{Z} \ar@{~>}[l] \ar[r]^-g &
\mathcal{Y}
}
}
\end{equation}
其中波浪箭头表示如下构造：把以集合胚为纤维的叠送到其对象同构类所成的
相应层。

\begin{lemma}
\label{algebraic-lemma-composition-representable-by-spaces}
设 $S$ 是 $\Sch_{fppf}$ 中的概形。设
$\mathcal{X}, \mathcal{Y}, \mathcal{Z}$ 是
$(\Sch/S)_{fppf}$ 上的群胚纤维化范畴。若
$f : \mathcal{X} \to \mathcal{Y}$、
$g : \mathcal{Y} \to \mathcal{Z}$ 是可由代数空间表示的
$1$-态射，则
$$
g \circ f : \mathcal{X} \longrightarrow \mathcal{Z}
$$
是可由代数空间表示的 $1$-态射。
\end{lemma}

\begin{proof}
这由引理 \ref{algebraic-lemma-base-change-by-space-representable-by-space}
推出。略去细节。
\end{proof}

\begin{lemma}
\label{algebraic-lemma-product-representable-by-spaces}
设 $S$ 是 $\Sch_{fppf}$ 中的概形。设
$\mathcal{X}_i, \mathcal{Y}_i$ 是 $(\Sch/S)_{fppf}$ 上的
群胚纤维化范畴，$i = 1, 2$。设
$f_i : \mathcal{X}_i \to \mathcal{Y}_i$，$i = 1, 2$，
是可由代数空间表示的 $1$-态射。则
$$
f_1 \times f_2 :
\mathcal{X}_1 \times \mathcal{X}_2
\longrightarrow
\mathcal{Y}_1 \times \mathcal{Y}_2
$$
是可由代数空间表示的 $1$-态射。
\end{lemma}

\begin{proof}
把 $f_1 \times f_2$ 写成复合
$\mathcal{X}_1 \times \mathcal{X}_2 \to
\mathcal{Y}_1 \times \mathcal{X}_2 \to
\mathcal{Y}_1 \times \mathcal{Y}_2$。第一个箭头是 $f_1$ 沿映射
$\mathcal{Y}_1 \times \mathcal{X}_2 \to \mathcal{Y}_1$ 的基变换，
第二个箭头是 $f_2$ 沿映射
$\mathcal{Y}_1 \times \mathcal{Y}_2 \to \mathcal{Y}_2$ 的基变换。
因此，本引理是引理
\ref{algebraic-lemma-composition-representable-by-spaces} 与
\ref{algebraic-lemma-base-change-representable-by-spaces} 的形式推论。
\end{proof}

\begin{lemma}
\label{algebraic-lemma-get-a-stack}
\begin{reference}
Matthew Emerton 在 2016 年 6 月 15 日的一封电子邮件中的引理
\end{reference}
设 $S$ 是 $\Sch_{fppf}$ 中的概形。设
$\mathcal{X} \to \mathcal{Z}$ 与 $\mathcal{Y} \to \mathcal{Z}$
是 $(\Sch/S)_{fppf}$ 上群胚纤维化范畴的 $1$-态射。
若 $\mathcal{X} \to \mathcal{Z}$ 可由代数空间表示，且
$\mathcal{Y}$ 是群胚叠，则
$\mathcal{X} \times_\mathcal{Z} \mathcal{Y}$ 是群胚叠。
\end{lemma}

\begin{proof}
态射可由代数空间表示这一性质在基变换下保持
（引理 \ref{algebraic-lemma-base-change-by-space-representable-by-space}）。
因此，对 $\mathcal{Y}$ 作基变换
$\mathcal{X} \times_\mathcal{Z} \mathcal{Y}$ 后，可归结到如下情形：
群胚纤维化范畴的态射 $\mathcal{X} \to \mathcal{Y}$
可由代数空间表示，且其靶是群胚叠；此时目标是证明
$\mathcal{X}$ 也是群胚叠。这由《叠》引理
\ref{stacks-lemma-relative-sheaf-over-stack-is-stack} 得出；
引理 \ref{algebraic-lemma-criterion-map-representable-spaces-fibred-in-groupoids}
保证了其假设成立。
\end{proof}







\section{可由代数空间表示的态射的性质}
\label{algebraic-section-representable-properties}

\noindent
下面的定义使这一说法得以成立。

\begin{definition}
\label{algebraic-definition-relative-representable-property}
设 $S$ 是 $\Sch_{fppf}$ 中的概形。设
$f : \mathcal{X} \to \mathcal{Y}$ 是 $(\Sch/S)_{fppf}$ 上
群胚纤维化范畴的 $1$-态射。假设 $f$ 可由代数空间表示。
设 $\mathcal{P}$ 是代数空间态射的一个性质，并且
\begin{enumerate}
\item 它在任意基变换下保持；
\item 它在基上是 fppf 局部的；参见《下降与代数空间》定义
\ref{spaces-descent-definition-property-morphisms-local}。
\end{enumerate}
此时，若对每个 $U \in \Ob((\Sch/S)_{fppf})$ 及任意
$y \in \mathcal{Y}_U$，所得的代数空间态射
$f_y : F_y \to U$（见图
(\ref{algebraic-equation-representable-by-algebraic-spaces})）都具有性质
$\mathcal{P}$，则称 $f$ {\it 具有性质 $\mathcal{P}$}。
\end{definition}

\noindent
需要特别注意，我们只把此定义用于在基变换下稳定、并且在靶的 fppf
拓扑中为局部的态射性质。这并非因为在其他情况下该定义没有意义，
而是因为针对所考虑的性质，我们可能希望给出另一种更合适的定义。

\begin{lemma}
\label{algebraic-lemma-property-morphism-equivalent}
设 $S$ 是 $\Sch_{fppf}$ 的对象。设 $\mathcal{P}$ 如定义
\ref{algebraic-definition-relative-representable-property} 中所述。考虑
$(\Sch/S)_{fppf}$ 上群胚纤维化范畴的 $1$-态射所成的
$2$-交换图
$$
\xymatrix{
\mathcal{X}' \ar[r] \ar[d]_{f'} & \mathcal{X} \ar[d]^f \\
\mathcal{Y}' \ar[r] & \mathcal{Y}
}
$$
。假设水平箭头都是等价，并且 $f$（等价地，$f'$）可由代数空间表示。
则 $f$ 具有 $\mathcal{P}$ 当且仅当 $f'$ 具有 $\mathcal{P}$。
\end{lemma}

\begin{proof}
由引理 \ref{algebraic-lemma-representable-by-spaces-morphism-equivalent}
可知该陈述有意义。略去证明。
\end{proof}

\noindent
下面作一个基本检验。

\begin{lemma}
\label{algebraic-lemma-map-presheaves-representable-by-spaces-transformation-property}
设 $S$ 是 $\Sch_{fppf}$ 中的概形。设 $a : F \to G$ 是
$(\Sch/S)_{fppf}$ 上的预层映射。设 $\mathcal{P}$ 如定义
\ref{algebraic-definition-relative-representable-property} 中所述，并假设
$a$ 可由代数空间表示。则 $a : F \to G$ 具有性质
$\mathcal{P}$（参见《自举》定义
\ref{bootstrap-definition-property-transformation}），当且仅当
群胚纤维化范畴的对应态射 $\mathcal{S}_F \to \mathcal{S}_G$
具有性质 $\mathcal{P}$。
\end{lemma}

\begin{proof}
由引理 \ref{algebraic-lemma-map-presheaves-representable-by-algebraic-spaces}
可知本引理有意义。略去证明。
\end{proof}

\begin{lemma}
\label{algebraic-lemma-map-fibred-setoids-property}
设 $S$ 是 $\Sch_{fppf}$ 的对象。设 $\mathcal{P}$ 如定义
\ref{algebraic-definition-relative-representable-property} 中所述。
设 $f : \mathcal{X} \to \mathcal{Y}$ 是 $(\Sch/S)_{fppf}$ 上
集合胚纤维化范畴的 $1$-态射。令 $F$、$G$ 分别为如下预层：
它把 $T$ 分别送到 $\mathcal{X}_T$、$\mathcal{Y}_T$
的对象同构类集合。令 $a : F \to G$ 为与 $f$ 对应的预层映射。
则 $a$ 具有 $\mathcal{P}$ 当且仅当 $f$ 具有 $\mathcal{P}$。
\end{lemma}

\begin{proof}
由引理 \ref{algebraic-lemma-map-fibred-setoids-representable-algebraic-spaces}
可知本引理有意义。把引理 \ref{algebraic-lemma-property-morphism-equivalent}
与引理
\ref{algebraic-lemma-map-presheaves-representable-by-spaces-transformation-property}
结合即可得到结论。
\end{proof}

\begin{lemma}
\label{algebraic-lemma-composition-representable-transformations-property}
设 $S$ 是 $\Sch_{fppf}$ 中的概形。设
$\mathcal{X}$、$\mathcal{Y}$、$\mathcal{Z}$ 是
$(\Sch/S)_{fppf}$ 上的群胚纤维化范畴。设 $\mathcal{P}$
是定义 \ref{algebraic-definition-relative-representable-property}
所述的性质，并且对复合稳定。设
$f : \mathcal{X} \to \mathcal{Y}$、
$g : \mathcal{Y} \to \mathcal{Z}$ 是可由代数空间表示的
$1$-态射。若 $f$ 与 $g$ 具有性质 $\mathcal{P}$，则
$g \circ f : \mathcal{X} \to \mathcal{Z}$ 也具有该性质。
\end{lemma}

\begin{proof}
由引理 \ref{algebraic-lemma-composition-representable-by-spaces}
可知本引理有意义。略去证明。
\end{proof}

\begin{lemma}
\label{algebraic-lemma-base-change-representable-transformations-property}
设 $S$ 是 $\Sch_{fppf}$ 中的概形。设
$\mathcal{X}, \mathcal{Y}, \mathcal{Z}$ 是
$(\Sch/S)_{fppf}$ 上的群胚纤维化范畴。设 $\mathcal{P}$
为定义 \ref{algebraic-definition-relative-representable-property}
所述的性质。设 $f : \mathcal{X} \to \mathcal{Y}$
是可由代数空间表示的 $1$-态射，$g : \mathcal{Z} \to \mathcal{Y}$
是任意 $1$-态射。考虑 $2$-纤维积图
$$
\xymatrix{
\mathcal{Z} \times_{g, \mathcal{Y}, f} \mathcal{X} \ar[r]_-{g'} \ar[d]_{f'} &
\mathcal{X} \ar[d]^f \\
\mathcal{Z} \ar[r]^g & \mathcal{Y}
}
$$
若 $f$ 具有 $\mathcal{P}$，则基变换 $f'$ 也具有 $\mathcal{P}$。
\end{lemma}

\begin{proof}
由引理 \ref{algebraic-lemma-base-change-representable-by-spaces}
可知本引理有意义。略去证明。
\end{proof}

\begin{lemma}
\label{algebraic-lemma-descent-representable-transformations-property}
设 $S$ 是 $\Sch_{fppf}$ 中的概形。设
$\mathcal{X}, \mathcal{Y}, \mathcal{Z}$ 是
$(\Sch/S)_{fppf}$ 上的群胚纤维化范畴。设 $\mathcal{P}$
为定义 \ref{algebraic-definition-relative-representable-property}
所述的性质。设 $f : \mathcal{X} \to \mathcal{Y}$
是可由代数空间表示的 $1$-态射，$g : \mathcal{Z} \to \mathcal{Y}$
是任意 $1$-态射。考虑纤维积图
$$
\xymatrix{
\mathcal{Z} \times_{g, \mathcal{Y}, f} \mathcal{X} \ar[r]_-{g'} \ar[d]_{f'} &
\mathcal{X} \ar[d]^f \\
\mathcal{Z} \ar[r]^g & \mathcal{Y}
}
$$
假设对每个概形 $U$ 以及 $\mathcal{Y}_U$ 的对象 $x$，都存在
fppf 覆盖 $\{U_i \to U\}$，使得 $x|_{U_i}$ 属于函子
$g : \mathcal{Z}_{U_i} \to \mathcal{Y}_{U_i}$ 的本质像。
此时，若 $f'$ 具有 $\mathcal{P}$，则 $f$ 具有 $\mathcal{P}$。
\end{lemma}

\begin{proof}
略去证明。提示：比较《代数空间》引理
\ref{spaces-lemma-descent-representable-transformations-property}
的证明。
\end{proof}

\begin{lemma}
\label{algebraic-lemma-product-representable-transformations-property}
设 $S$ 是 $\Sch_{fppf}$ 中的概形。设 $\mathcal{P}$
为定义 \ref{algebraic-definition-relative-representable-property}
所述且对复合稳定的性质。设 $\mathcal{X}_i, \mathcal{Y}_i$
是 $(\Sch/S)_{fppf}$ 上的群胚纤维化范畴，$i = 1, 2$。
设 $f_i : \mathcal{X}_i \to \mathcal{Y}_i$，$i = 1, 2$，
是可由代数空间表示的 $1$-态射。若 $f_1$ 与 $f_2$ 具有性质
$\mathcal{P}$，则
$
f_1 \times f_2 :
\mathcal{X}_1 \times \mathcal{X}_2
\to
\mathcal{Y}_1 \times \mathcal{Y}_2
$
也具有该性质。
\end{lemma}

\begin{proof}
由引理 \ref{algebraic-lemma-product-representable-by-spaces}
可知本引理有意义。略去证明。
\end{proof}

\begin{lemma}
\label{algebraic-lemma-representable-transformations-property-implication}
设 $S$ 是 $\Sch_{fppf}$ 中的概形。设 $\mathcal{X}$、$\mathcal{Y}$
是 $(\Sch/S)_{fppf}$ 上的群胚纤维化范畴。设
$f : \mathcal{X} \to \mathcal{Y}$ 是可由代数空间表示的
$1$-态射。设 $\mathcal{P}$、$\mathcal{P}'$ 是定义
\ref{algebraic-definition-relative-representable-property} 所述的性质。
假设对代数空间的任意态射 $a : F \to G$，都有
$\mathcal{P}(a) \Rightarrow \mathcal{P}'(a)$。若 $f$ 具有性质
$\mathcal{P}$，则 $f$ 具有性质 $\mathcal{P}'$。
\end{lemma}

\begin{proof}
形式成立。
\end{proof}

\begin{lemma}
\label{algebraic-lemma-open-fibred-category-is-full}
设 $S$ 是 $\Sch_{fppf}$ 中的概形。设
$j : \mathcal X \to \mathcal Y$ 是 $(\Sch/S)_{fppf}$ 上
群胚纤维化范畴的 $1$-态射。假设 $j$ 可由代数空间表示且为单态射
（参见定义 \ref{algebraic-definition-relative-representable-property}
以及《下降与代数空间》引理
\ref{spaces-descent-lemma-descending-property-monomorphism}）。
则 $j$ 在各纤维范畴上全忠实。
\end{lemma}

\begin{proof}
由引理 \ref{algebraic-lemma-criterion-map-representable-spaces-fibred-in-groupoids}
已知 $j$ 在各纤维范畴上忠实。考虑概形 $U$、
$\mathcal{X}_U$ 的两个对象 $u, v$，以及 $\mathcal{Y}_U$ 中的同构
$t : j(u) \to j(v)$。我们需要在 $\mathcal{X}_U$ 中构造
$u$ 与 $v$ 之间的同构。由 $2$-Yoneda 引理（见第
\ref{algebraic-section-2-yoneda} 节），把 $u$、$v$ 看成 $1$-态射
$u, v : (\Sch/U)_{fppf} \to \mathcal{X}$
，并考虑 $2$-纤维积
$$
(\Sch/U)_{fppf} \times_{j \circ v, \mathcal{Y}} \mathcal{X}.
$$
根据假设，它可由 $U$ 上的代数空间 $F_{j \circ v}$ 表示，
而态射 $F_{j \circ v} \to U$ 是单态射。但
$(1_U, v, 1_{j(v)})$ 给出从 $(\Sch/U)_{fppf}$ 到上述
$2$-纤维积的 $1$-态射，故 $F_{j \circ v} = U$
（这里用到：若代数空间的单态射 $V \to U$ 有截面，则 $V = U$）。
因此，投影到第一坐标的 $1$-态射
$$
(\Sch/U)_{fppf} \times_{j \circ v, \mathcal{Y}} \mathcal{X}
\to (\Sch/U)_{fppf}
$$
是纤维范畴的等价。由于 $(1_U, u, t)$ 与
$(1_U, v, 1_{j(v)})$ 给出
$((\Sch/U)_{fppf} \times_{j \circ v, \mathcal{Y}}
\mathcal{X})_U$ 中具有相同第一坐标的两个对象，故它们之间必有
$2$-纤维积中的一个 $2$-态射。按定义，这就是满足
$j(\tilde t) = t$ 的态射 $\tilde t : u \to v$。
\end{proof}

\noindent
下面刻画对角态射可由代数空间表示的群胚纤维化范畴。

\begin{lemma}
\label{algebraic-lemma-representable-diagonal}
设 $S$ 是 $\Sch_{fppf}$ 中的概形。设 $\mathcal{X}$ 是
$(\Sch/S)_{fppf}$ 上的群胚纤维化范畴。以下各条件等价：
\begin{enumerate}
\item 对角态射 $\mathcal{X} \to \mathcal{X} \times \mathcal{X}$
可由代数空间表示；
\item 对 $S$ 上的每个概形 $U$ 以及任意
$x, y \in \Ob(\mathcal{X}_U)$，同构层
$\mathit{Isom}(x, y)$ 是 $U$ 上的代数空间；
\item 对 $S$ 上的每个概形 $U$ 以及任意
$x \in \Ob(\mathcal{X}_U)$，相应的 $1$-态射
$x : (\Sch/U)_{fppf} \to \mathcal{X}$ 可由代数空间表示；
\item 对 $S$ 上的每对概形 $T_1, T_2$ 以及任意
$x_i \in \Ob(\mathcal{X}_{T_i})$，$i = 1, 2$，$2$-纤维积
$(\Sch/T_1)_{fppf} \times_{x_1, \mathcal{X}, x_2}
(\Sch/T_2)_{fppf}$
可由代数空间表示；
\item 对 $(\Sch/S)_{fppf}$ 上的每个可表群胚纤维化范畴
$\mathcal{U}$，每个 $1$-态射 $\mathcal{U} \to \mathcal{X}$
都可由代数空间表示；
\item 对 $(\Sch/S)_{fppf}$ 上的每对可表群胚纤维化范畴
$\mathcal{T}_1, \mathcal{T}_2$，以及任意 $1$-态射
$x_i : \mathcal{T}_i \to \mathcal{X}$，$i = 1, 2$，$2$-纤维积
$\mathcal{T}_1 \times_{x_1, \mathcal{X}, x_2} \mathcal{T}_2$
可由代数空间表示；
\item 对 $(\Sch/S)_{fppf}$ 上可由代数空间表示的每个
群胚纤维化范畴 $\mathcal{U}$，每个 $1$-态射
$\mathcal{U} \to \mathcal{X}$ 都可由代数空间表示；
\item 对 $(\Sch/S)_{fppf}$ 上可由代数空间表示的每对
群胚纤维化范畴 $\mathcal{T}_1, \mathcal{T}_2$，以及任意
$1$-态射 $x_i : \mathcal{T}_i \to \mathcal{X}$，$2$-纤维积
$\mathcal{T}_1 \times_{x_1, \mathcal{X}, x_2} \mathcal{T}_2$
可由代数空间表示。
\end{enumerate}
\end{lemma}

\begin{proof}
(1) 与 (2) 的等价性由《叠》引理
\ref{stacks-lemma-isom-as-2-fibre-product} 和定义得出。
下面证明 (1) 与 (3) 等价。记基范畴为
$\mathcal{C} = (\Sch/S)_{fppf}$。我们将使用类似的《范畴》引理
\ref{categories-lemma-representable-diagonal-groupoids}
的证明中的若干观察。我们用符号 $\cong$ 表示“在
$\mathcal{C} = (\Sch/S)_{fppf}$ 上群胚纤维化范畴的等价”。
假设 (1) 成立。给定 (3) 中的 $U$ 与 $x$。对任意概形 $V$ 以及
$y \in \Ob(\mathcal{X}_V)$，有（比较上述引用）
$$
\mathcal{C}/U
\times_{x, \mathcal{X}, y}
\mathcal{C}/V
\cong
(\mathcal{C}/U \times_S V)
\times_{(x, y), \mathcal{X} \times \mathcal{X}, \Delta}
\mathcal{X}
$$
，它根据假设可由代数空间表示。反过来，假设 (3) 成立。考虑
$S$ 上任意概形 $U$，以及 $\mathcal{X}$ 在 $U$ 上的一对对象
$(x, x')$。我们需要证明
$\mathcal{X} \times_{\Delta, \mathcal{X} \times \mathcal{X}, (x, x')} U$
可由代数空间表示。这很清楚，因为（比较上述引用）
$$
\mathcal{X}
\times_{\Delta, \mathcal{X} \times \mathcal{X}, (x, x')}
\mathcal{C}/U
\cong
(\mathcal{C}/U \times_{x, \mathcal{X}, x'} \mathcal{C}/U)
\times_{\mathcal{C}/U \times_S U, \Delta}
\mathcal{C}/U
$$
，而右端根据假设以及如下事实可由代数空间表示：$S$ 上代数空间的范畴
具有纤维积，并且包含 $U$ 与 $S$。

\medskip\noindent
等价关系 (3) $\Leftrightarrow$ (4)、(5) $\Leftrightarrow$ (6)
以及 (7) $\Leftrightarrow$ (8) 都是形式的。等价关系
(3) $\Leftrightarrow$ (5) 与 (4) $\Leftrightarrow$ (6)
由引理 \ref{algebraic-lemma-representable-by-spaces-morphism-equivalent} 得出。
假设 (3) 成立，并设 $\mathcal{U} \to \mathcal{X}$ 如 (7) 中所述。
为证明 (7)，必须说明：对每个概形 $V$ 以及每个 $1$-态射
$y : (\Sch/V)_{fppf} \to \mathcal{X}$，$2$-纤维积
$(\Sch/V)_{fppf} \times_{y, \mathcal{X}} \mathcal{U}$
可由代数空间表示。性质 (3) 告诉我们 $y$ 可由代数空间表示，
故引理 \ref{algebraic-lemma-base-change-by-space-representable-by-space}
给出所需结论。最后，(7) 直接推出 (3)。
\end{proof}

\noindent
在该引理的情形下，对其中任意 $1$-态射
$x : (\Sch/U)_{fppf} \to \mathcal{X}$，以及定义
\ref{algebraic-definition-relative-representable-property} 所述的任意性质，
说 $x$ 具有性质 $\mathcal{P}$ 都是有意义的。
特别地，这适用于 $\mathcal{P} = $“满射”、
$\mathcal{P} = $“光滑”以及 $\mathcal{P} = $“\'etale”；
参见《下降与代数空间》引理
\ref{spaces-descent-lemma-descending-property-surjective}、
\ref{spaces-descent-lemma-descending-property-smooth} 与
\ref{spaces-descent-lemma-descending-property-etale}。
在下文定义代数叠时，我们将使用这三种情形。








\section{群胚叠}
\label{algebraic-section-stacks}

\noindent
设 $S$ 是 $\Sch_{fppf}$ 中的概形。回忆，
$(\Sch/S)_{fppf}$ 上的范畴
$p : \mathcal{X} \to (\Sch/S)_{fppf}$
称为{\it 群胚叠}（参见《叠》定义
\ref{stacks-definition-stack-in-groupoids}），当且仅当
\begin{enumerate}
\item $p : \mathcal{X} \to (\Sch/S)_{fppf}$ 在
$(\Sch/S)_{fppf}$ 上以群胚为纤维；
\item 对所有 $U \in \Ob((\Sch/S)_{fppf})$ 以及所有
$x, y\in \Ob(\mathcal{X}_U)$，预层 $\mathit{Isom}(x, y)$
是位点 $(\Sch/U)_{fppf}$ 上的层；
\item 对 $(\Sch/S)_{fppf}$ 中的所有覆盖
$\mathcal{U} = \{U_i \to U\}$，关于 $\mathcal{U}$ 的所有
下降数据 $(x_i, \phi_{ij})$ 都是有效的。
\end{enumerate}
关于例子，参见《叠的例子》第
\ref{examples-stacks-section-examples-stacks-in-groupoids}
节及以下各节。










\section{代数叠}
\label{algebraic-section-algebraic-stacks}

\noindent
下面给出代数叠的定义。注意，条件 (2) 使得第 (3) 项中
$(\Sch/U)_{fppf} \to \mathcal{X}$ 光滑且满这一条件有意义；
参见引理 \ref{algebraic-lemma-representable-diagonal} 后的讨论。

\begin{definition}
\label{algebraic-definition-algebraic-stack}
设 $S$ 是 $\Sch_{fppf}$ 中的基概形。{\it $S$ 上的代数叠}
是范畴
$$
p : \mathcal{X} \to (\Sch/S)_{fppf}
$$
，它位于 $(\Sch/S)_{fppf}$ 上，并具有以下性质：
\begin{enumerate}
\item 范畴 $\mathcal{X}$ 是 $(\Sch/S)_{fppf}$ 上的群胚叠。
\item 对角态射
$\Delta : \mathcal{X} \to \mathcal{X} \times \mathcal{X}$
可由代数空间表示。
\item 存在概形 $U \in \Ob((\Sch/S)_{fppf})$ 以及满且光滑的
$1$-态射 $(\Sch/U)_{fppf} \to \mathcal{X}$
\footnote{在后续各章中，我们将遵循文献中的惯例，把它简记为
$U \to \mathcal{X}$。另一种同样合适的表述是：存在一个可表的
群胚纤维化范畴 $\mathcal{U}$ 以及满且光滑的 $1$-态射
$\mathcal{U} \to \mathcal{X}$。}。
\end{enumerate}
\end{definition}

\noindent
这一定义与文献中的其他定义有若干差别。

\medskip\noindent
第一，我们要求 $\mathcal{X}$ 在 fppf 拓扑中为群胚叠，而许多文献
使用 \'etale 拓扑。在我们看来，fppf 拓扑似乎是自然的工作拓扑。
最终所得的代数叠 $2$-范畴仍然相同。《可表性判据》第
\ref{criteria-section-stacks-etale} 节对此作了说明。

\medskip\noindent
第二，我们只要求 $\mathcal{X}$ 的对角态射可由代数空间表示，
而大多数文献还施加其他条件。我们的观点是：先尝试只在
$\mathcal{X}$ 的对角态射可由代数空间表示这一假设下证明随后的一些
结果，并仅在必要处加入额外假设。这样还有一个好处：任意代数空间
（定义见《代数空间》定义 \ref{spaces-definition-algebraic-space}）
都给出一个代数叠。

\medskip\noindent
第三，一些论文要求存在概形 $U$ 以及满且 \'etale 的态射
$U \to \mathcal{X}$。在首次引入代数叠的奠基性论文 \cite{DM} 中，
Deligne 与 Mumford 采用了这一定义，并证明亏格 $g > 1$ 的稳定曲线
模叠是具有概形 \'etale 覆盖的代数叠。Michael Artin（参见
\cite{ArtinVersal}）认识到，关于代数叠的许多自然结果都能推广到
只假设存在概形光滑覆盖的情形。因此我们作出上述选择。为区分这两种
情形，文献中使用“Deligne--Mumford 叠”和“Artin 叠”这两个术语。
我们把“Artin 叠”一词保留到以后使用（此处插入未来引用），并继续使用
“代数叠”；但用“Deligne--Mumford 叠”指具有概形 \'etale 覆盖的
代数叠。

\begin{definition}
\label{algebraic-definition-deligne-mumford}
设 $S$ 是 $\Sch_{fppf}$ 中的概形。设 $\mathcal{X}$ 是
$S$ 上的代数叠。若存在概形 $U$ 以及满 \'etale 态射
$(\Sch/U)_{fppf} \to \mathcal{X}$，则称 $\mathcal{X}$
为{\it Deligne--Mumford 叠}。
\end{definition}

\noindent
稍后我们将比较这里的 Deligne--Mumford 叠概念与 Deligne 和 Mumford
论文中定义的概念（见此处插入未来引用）。

\medskip\noindent
$S$ 上的代数叠所成范畴构成一个 $2$-范畴。下面给出精确定义。

\begin{definition}
\label{algebraic-definition-morphism-algebraic-stacks}
设 $S$ 是 $\Sch_{fppf}$ 中的概形。{\it $S$ 上代数叠的
$2$-范畴}是 $(\Sch/S)_{fppf}$ 上群胚纤维化范畴所成
$2$-范畴（参见《范畴》定义
\ref{categories-definition-categories-fibred-in-groupoids-over-C}）
的如下子 $2$-范畴：
\begin{enumerate}
\item 其对象是 $(\Sch/S)_{fppf}$ 上那些构成 $S$ 上代数叠的
群胚纤维化范畴。
\item 其 $1$-态射 $f : \mathcal{X} \to \mathcal{Y}$ 是
$(\Sch/S)_{fppf}$ 上范畴的任意函子，如《范畴》定义
\ref{categories-definition-categories-over-C}。
\item 其 $2$-态射是 $(\Sch/S)_{fppf}$ 上函子之间的变换，
如《范畴》定义 \ref{categories-definition-categories-over-C}。
\end{enumerate}
\end{definition}

\noindent
换言之，这个 $2$-范畴是 $\textit{Cat}/(\Sch/S)_{fppf}$
中以代数叠为对象的满子 $2$-范畴。注意，每个 $2$-态射自动为同构。
因此，这实际上是一个 $(2, 1)$-范畴，而不仅是 $2$-范畴。

\medskip\noindent
稍后将看到（此处插入未来引用），这个 $2$-范畴具有 $2$-纤维积。

\medskip\noindent
与上面的注记类似，$S$ 上代数叠的 $2$-范畴是
$(\Sch/S)_{fppf}$ 上群胚纤维化范畴所成 $2$-范畴的满子
$2$-范畴。事实证明，它在等价下封闭。精确陈述如下。

\begin{lemma}
\label{algebraic-lemma-equivalent}
设 $S$ 是 $\Sch_{fppf}$ 中的概形。设 $\mathcal{X}$、$\mathcal{Y}$
是 $(\Sch/S)_{fppf}$ 上的范畴。假设 $\mathcal{X}$、$\mathcal{Y}$
作为 $(\Sch/S)_{fppf}$ 上的范畴等价。则 $\mathcal{X}$ 是代数叠
当且仅当 $\mathcal{Y}$ 是代数叠。类似地，$\mathcal{X}$ 是
Deligne--Mumford 叠当且仅当 $\mathcal{Y}$ 是 Deligne--Mumford 叠。
\end{lemma}

\begin{proof}
假设 $\mathcal{X}$ 是代数叠（相应地，是 Deligne--Mumford 叠）。
由《叠》引理 \ref{stacks-lemma-stack-in-groupoids-equivalent}，
$\mathcal{Y}$ 是 $\Sch_{fppf}$ 上的群胚叠。选取
$\Sch_{fppf}$ 上的等价 $f : \mathcal{X} \to \mathcal{Y}$。
这给出 $2$-交换图
$$
\xymatrix{
\mathcal{X} \ar[r]_f \ar[d]_{\Delta_\mathcal{X}} &
\mathcal{Y} \ar[d]^{\Delta_\mathcal{Y}} \\
\mathcal{X} \times \mathcal{X} \ar[r]^{f \times f} &
\mathcal{Y} \times \mathcal{Y}
}
$$
，其中水平箭头都是等价。根据引理
\ref{algebraic-lemma-representable-by-spaces-morphism-equivalent}，
这说明 $\Delta_\mathcal{Y}$ 可由代数空间表示。最后，设 $U$ 是
$S$ 上的概形，并设 $x : (\Sch/U)_{fppf} \to \mathcal{X}$
是满且光滑（相应地，\'etale）的 $1$-态射。考虑图表
$$
\xymatrix{
(\Sch/U)_{fppf} \ar[r]_{\text{id}} \ar[d]_x &
(\Sch/U)_{fppf} \ar[d]^{f \circ x} \\
\mathcal{X} \ar[r]^f &
\mathcal{Y}
}
$$
并应用引理 \ref{algebraic-lemma-property-morphism-equivalent}，可知
$f \circ x$ 满且光滑（相应地，\'etale），如所需。
\end{proof}




\section{代数叠与代数空间}
\label{algebraic-section-stacks-spaces}

\noindent
本节讨论若干蕴含代数叠为代数空间的简单判据。主要结果是：这恰好发生在
各纤维范畴的对象没有非平凡自同构时。这并非显然！在证明之前，
先作一个基本检验。

\begin{lemma}
\label{algebraic-lemma-representable-algebraic}
设 $S$ 是 $\Sch_{fppf}$ 中的概形。
\begin{enumerate}
\item 若群胚纤维化范畴
$p : \mathcal{X} \to (\Sch/S)_{fppf}$
可由代数空间表示，则它是 Deligne--Mumford 叠。
\item 若 $F$ 是 $S$ 上的代数空间，则相应的群胚纤维化范畴
$p : \mathcal{S}_F \to (\Sch/S)_{fppf}$
是 Deligne--Mumford 叠。
\item 若 $X \in \Ob((\Sch/S)_{fppf})$，则
$(\Sch/X)_{fppf} \to (\Sch/S)_{fppf}$
是 Deligne--Mumford 叠。
\end{enumerate}
\end{lemma}

\begin{proof}
(2) 蕴含 (3) 是清楚的。由引理 \ref{algebraic-lemma-equivalent}，
(1) 与 (2) 等价。因此只需证明 (2)。首先，由于 $F$ 是层，
$\mathcal{S}_F$ 是集合叠（《叠》引理
\ref{stacks-lemma-stack-in-setoids-characterize}），因而更是群胚叠。
其次，对角态射
$\mathcal{S}_F \to \mathcal{S}_F  \times \mathcal{S}_F$
就是由 $F$ 的对角态射产生的态射
$\mathcal{S}_F \to \mathcal{S}_{F \times F}$。故由引理
\ref{algebraic-lemma-morphism-spaces-gives-representable-by-spaces}，
它可由代数空间表示。事实上，它甚至（由概形）可表，因为代数空间的
对角态射可表，但我们不需要这一点。设 $U$ 是概形，并设
$h_U \to F$ 是满 \'etale 态射。可把它看作代数空间的满
\'etale 态射。于是由引理
\ref{algebraic-lemma-map-presheaves-representable-by-spaces-transformation-property}，
相应的 $1$-态射 $(\Sch/U)_{fppf} \to \mathcal{S}_F$
满且 \'etale。
\end{proof}

\noindent
下面的结果说明，惯性平凡的 Deligne--Mumford 叠“就是”代数空间。
下方更强的命题
\ref{algebraic-proposition-algebraic-stack-no-automorphisms}
将取代本引理；该命题说明这一结论对一般代数叠也成立……

\begin{lemma}
\label{algebraic-lemma-algebraic-stack-no-automorphisms}
设 $S$ 是 $\Sch_{fppf}$ 中的概形。设 $\mathcal{X}$ 是
$S$ 上的代数叠。以下各条件等价：
\begin{enumerate}
\item $\mathcal{X}$ 是 Deligne--Mumford 叠且为集合胚叠；
\item $\mathcal{X}$ 是 Deligne--Mumford 叠，且典范 $1$-态射
$\mathcal{I}_\mathcal{X} \to \mathcal{X}$ 是等价；
\item $\mathcal{X}$ 可由代数空间表示。
\end{enumerate}
\end{lemma}

\begin{proof}
(1) 与 (2) 的等价性由《叠》引理
\ref{stacks-lemma-characterize-stack-in-setoids} 得出。
(3) $\Rightarrow$ (1) 由引理 \ref{algebraic-lemma-representable-algebraic}
得出。最后，假设 (1) 成立。由《叠》引理
\ref{stacks-lemma-stack-in-setoids-characterize}，存在
$(\Sch/S)_{fppf}$ 上的层 $F$ 以及等价
$j : \mathcal{X} \to \mathcal{S}_F$。由引理
\ref{algebraic-lemma-map-presheaves-representable-by-algebraic-spaces}，
$\Delta_\mathcal{X}$ 可由代数空间表示意味着
$\Delta_F : F \to F \times F$ 可由代数空间表示。
设 $U$ 是概形，并设 $x : (\Sch/U)_{fppf} \to \mathcal{X}$
是满 \'etale 态射。复合
$j \circ x : (\Sch/U)_{fppf} \to \mathcal{S}_F$
对应于层的态射 $h_U \to F$。由《自举》引理
\ref{bootstrap-lemma-representable-diagonal}，该态射可由代数空间
表示。于是由引理 \ref{algebraic-lemma-map-fibred-setoids-property}，
$h_U \to F$ 满且 \'etale。最后应用《自举》定理
\ref{bootstrap-theorem-bootstrap}，可知 $F$ 是代数空间。
\end{proof}

\begin{proposition}
\label{algebraic-proposition-algebraic-stack-no-automorphisms}
设 $S$ 是 $\Sch_{fppf}$ 中的概形。设 $\mathcal{X}$ 是
$S$ 上的代数叠。以下各条件等价：
\begin{enumerate}
\item $\mathcal{X}$ 是集合胚叠；
\item 典范 $1$-态射 $\mathcal{I}_\mathcal{X} \to \mathcal{X}$
是等价；
\item $\mathcal{X}$ 可由代数空间表示。
\end{enumerate}
\end{proposition}

\begin{proof}
(1) 与 (2) 的等价性由《叠》引理
\ref{stacks-lemma-characterize-stack-in-setoids} 得出。
(3) $\Rightarrow$ (1) 由引理
\ref{algebraic-lemma-algebraic-stack-no-automorphisms} 得出。最后，假设 (1)
成立。由《叠》引理 \ref{stacks-lemma-stack-in-setoids-characterize}，
存在等价 $j : \mathcal{X} \to \mathcal{S}_F$，其中 $F$ 是
$(\Sch/S)_{fppf}$ 上的层。由引理
\ref{algebraic-lemma-map-presheaves-representable-by-algebraic-spaces}，
$\Delta_\mathcal{X}$ 可由代数空间表示意味着
$\Delta_F : F \to F \times F$ 可由代数空间表示。
设 $U$ 是概形，并设 $x : (\Sch/U)_{fppf} \to \mathcal{X}$
是满光滑态射。复合
$j \circ x : (\Sch/U)_{fppf} \to \mathcal{S}_F$
对应于层的态射 $h_U \to F$。由《自举》引理
\ref{bootstrap-lemma-representable-diagonal}，该态射可由代数空间
表示。因此由引理 \ref{algebraic-lemma-map-fibred-setoids-property}，
$h_U \to F$ 满且光滑。特别地，它是满、平坦且局部有限表示的
（这由引理
\ref{algebraic-lemma-representable-transformations-property-implication}
以及代数空间的光滑态射平坦且局部有限表示这一事实得出；参见
《代数空间的态射》引理
\ref{spaces-morphisms-lemma-smooth-locally-finite-presentation} 与
\ref{spaces-morphisms-lemma-smooth-flat}）。最后应用《自举》定理
\ref{bootstrap-theorem-final-bootstrap}，可知 $F$ 是代数空间。
\end{proof}






\section{代数叠的 2-纤维积}
\label{algebraic-section-2-fibre-products}

\noindent
代数叠的 $2$-范畴具有积与 $2$-纤维积。第一个引理其实是引理
\ref{algebraic-lemma-2-fibre-product} 的特例，但其证明稍为简单。

\begin{lemma}
\label{algebraic-lemma-product-spaces}
设 $S$ 是 $\Sch_{fppf}$ 中的概形。设 $\mathcal{X}$、$\mathcal{Y}$
是 $S$ 上的代数叠。则
$\mathcal{X} \times_{(\Sch/S)_{fppf}} \mathcal{Y}$
是代数叠，并且是 $S$ 上代数叠的 $2$-范畴中的积。
\end{lemma}

\begin{proof}
$\mathcal{X} \times_{(\Sch/S)_{fppf}} \mathcal{Y}$ 在 $T$ 上的
对象就是一对 $(x, y)$，其中 $x$ 是 $\mathcal{X}_T$ 的对象，
$y$ 是 $\mathcal{Y}_T$ 的对象。因此由定义立即可知，
$\mathcal{X} \times_{(\Sch/S)_{fppf}} \mathcal{Y}$ 是群胚叠。
若 $(x, y)$ 与 $(x', y')$ 是
$\mathcal{X} \times_{(\Sch/S)_{fppf}} \mathcal{Y}$
在 $T$ 上的两个对象，则
$$
\mathit{Isom}((x, y), (x', y')) =
\mathit{Isom}(x, x') \times \mathit{Isom}(y, y').
$$
于是，由引理 \ref{algebraic-lemma-representable-diagonal} 中的等价条件，
以及代数空间的范畴具有积这一事实，可知
$\mathcal{X} \times_{(\Sch/S)_{fppf}} \mathcal{Y}$ 的对角态射
可由代数空间表示。最后，设
$U, V \in \Ob((\Sch/S)_{fppf})$，并设 $x,y$ 是满光滑态射
$x : (\Sch/U)_{fppf} \to \mathcal{X}$,
$y : (\Sch/V)_{fppf} \to \mathcal{Y}$.
注意，
$$
(\Sch/U \times_S V)_{fppf} =
(\Sch/U)_{fppf}
\times_{(\Sch/S)_{fppf}} (\Sch/V)_{fppf}.
$$
因此，
$\mathcal{X} \times_{(\Sch/S)_{fppf}} \mathcal{Y}$ 在
$(\Sch/U \times_S V)_{fppf}$ 上的对象
$(\text{pr}_U^*x, \text{pr}_V^*y)$ 定义一个 $1$-态射
$$
(\Sch/U \times_S V)_{fppf}
\longrightarrow
\mathcal{X} \times_{(\Sch/S)_{fppf}} \mathcal{Y}
$$
。它是 $x$ 与 $y$ 的基变换的复合，故满且光滑；参见引理
\ref{algebraic-lemma-base-change-representable-transformations-property} 与
\ref{algebraic-lemma-composition-representable-transformations-property}。
因此，$\mathcal{X} \times_{(\Sch/S)_{fppf}} \mathcal{Y}$
确实是代数叠。略去它确实为积的验证。
\end{proof}

\begin{lemma}
\label{algebraic-lemma-2-fibre-product-general}
设 $S$ 是 $\Sch_{fppf}$ 中的概形。设 $\mathcal{Z}$ 是
$(\Sch/S)_{fppf}$ 上的群胚叠，且其对角态射可由代数空间表示。
设 $\mathcal{X}$、$\mathcal{Y}$ 是 $S$ 上的代数叠。
设 $f : \mathcal{X} \to \mathcal{Z}$、
$g : \mathcal{Y} \to \mathcal{Z}$ 是群胚叠的 $1$-态射。
则 $2$-纤维积
$\mathcal{X} \times_{f, \mathcal{Z}, g} \mathcal{Y}$ 是代数叠。
\end{lemma}

\begin{proof}
需要验证定义 \ref{algebraic-definition-algebraic-stack} 的条件 (1)、(2) 与
(3)。第一个条件由《叠》引理
\ref{stacks-lemma-2-product-stacks-in-groupoids} 得出。

\medskip\noindent
需要验证的第二个条件是 $\mathit{Isom}$ 层可由代数空间表示。
为此，设 $T$ 是 $S$ 上的概形，而 $u,v$ 是
$(\mathcal{X} \times_{f, \mathcal{Z}, g} \mathcal{Y})_T$
的对象。根据我们对 $2$-纤维积的构造（其源头可追溯到《范畴》引理
\ref{categories-lemma-2-product-categories-over-C}），可写成
$u = (x, y, \alpha)$ 与 $v = (x', y', \alpha')$，其中
$\alpha : f(x) \to g(y)$，$\alpha'$ 也类似。显然，
$$
\xymatrix{
\mathit{Isom}(u, v) \ar[d] \ar[rr] & &
\mathit{Isom}(y, y') \ar[d]^{\phi \mapsto g(\phi) \circ \alpha} \\
\mathit{Isom}(x, x') \ar[rr]^-{\psi \mapsto \alpha' \circ f(\psi)} & &
\mathit{Isom}(f(x), g(y'))
}
$$
是 $(\Sch/T)_{fppf}$ 上层的笛卡儿图。根据假设，层
$\mathit{Isom}(y, y')$、$\mathit{Isom}(x, x')$、
$\mathit{Isom}(f(x), g(y'))$ 都是代数空间（见引理
\ref{algebraic-lemma-representable-diagonal}），故
$\mathit{Isom}(u, v)$ 是代数空间。

\medskip\noindent
设 $U, V \in \Ob((\Sch/S)_{fppf})$，并设 $x,y$ 是满光滑态射
$x : (\Sch/U)_{fppf} \to \mathcal{X}$,
$y : (\Sch/V)_{fppf} \to \mathcal{Y}$.
考虑态射
$$
(\Sch/U)_{fppf}
\times_{f \circ x, \mathcal{Z}, g \circ y}
(\Sch/V)_{fppf}
\longrightarrow
\mathcal{X} \times_{f, \mathcal{Z}, g} \mathcal{Y}.
$$
由于 $\mathcal{Z}$ 的对角态射可由代数空间表示，该箭头的源可由
代数空间 $F$ 表示；见引理 \ref{algebraic-lemma-representable-diagonal}。
此外，该态射是 $x$ 与 $y$ 的基变换的复合，故满且光滑；参见引理
\ref{algebraic-lemma-base-change-representable-transformations-property} 与
\ref{algebraic-lemma-composition-representable-transformations-property}。
选取概形 $W$ 以及满 \'etale 态射 $W \to F$，可知上述
$1$-态射与相应 $1$-态射
$$
(\Sch/W)_{fppf}
\longrightarrow
(\Sch/U)_{fppf}
\times_{f \circ x, \mathcal{Z}, g \circ y}
(\Sch/V)_{fppf}
$$
的复合满且光滑，从而证明最后一个条件。
\end{proof}

\begin{lemma}
\label{algebraic-lemma-2-fibre-product}
设 $S$ 是 $\Sch_{fppf}$ 中的概形。设
$\mathcal{X}, \mathcal{Y}, \mathcal{Z}$ 是 $S$ 上的代数叠。
设 $f : \mathcal{X} \to \mathcal{Z}$、
$g : \mathcal{Y} \to \mathcal{Z}$ 是代数叠的 $1$-态射。
则 $2$-纤维积
$\mathcal{X} \times_{f, \mathcal{Z}, g} \mathcal{Y}$ 是代数叠。
它也是 $(\Sch/S)_{fppf}$ 上代数叠的 $2$-范畴中的 $2$-纤维积。
\end{lemma}

\begin{proof}
$\mathcal{X} \times_{f, \mathcal{Z}, g} \mathcal{Y}$ 是代数叠，
由更强的引理 \ref{algebraic-lemma-2-fibre-product-general} 得出。
而 $\mathcal{X} \times_{f, \mathcal{Z}, g} \mathcal{Y}$
是 $S$ 上代数叠的 $2$-范畴中的 $2$-纤维积，则形式地得自如下事实：
$S$ 上代数叠的 $2$-范畴是 $(\Sch/S)_{fppf}$ 上群胚叠所成
$2$-范畴的满子 $2$-范畴。
\end{proof}







\section{代数叠：重新整理}
\label{algebraic-section-overhaul}

\noindent
关于代数叠的一些基本结果。

\begin{lemma}
\label{algebraic-lemma-lift-morphism-presentations}
设 $S$ 是 $\Sch_{fppf}$ 中的概形。设
$f : \mathcal{X} \to \mathcal{Y}$ 是 $S$ 上代数叠的 $1$-态射。
设 $V \in \Ob((\Sch/S)_{fppf})$，并设
$y : (\Sch/V)_{fppf} \to \mathcal{Y}$ 满且光滑。
则存在对象 $U \in \Ob((\Sch/S)_{fppf})$ 以及 $2$-交换图
$$
\xymatrix{
(\Sch/U)_{fppf} \ar[r]_a \ar[d]_x &
(\Sch/V)_{fppf} \ar[d]^y \\
\mathcal{X} \ar[r]^f & \mathcal{Y}
}
$$
，其中 $x$ 满且光滑。
\end{lemma}

\begin{proof}
首先选取 $W \in \Ob((\Sch/S)_{fppf})$ 以及满光滑 $1$-态射
$z : (\Sch/W)_{fppf} \to \mathcal{X}$。由于 $\mathcal{Y}$
是代数叠，可选取等价
$$
j :
\mathcal{S}_F
\longrightarrow
(\Sch/W)_{fppf}
\times_{f \circ z, \mathcal{Y}, y}
(\Sch/V)_{fppf}
$$
，其中 $F$ 是代数空间。由引理
\ref{algebraic-lemma-base-change-representable-transformations-property}，
态射 $\mathcal{S}_F \to (\Sch/W)_{fppf}$ 作为 $y$ 的基变换
满且光滑。故由引理
\ref{algebraic-lemma-composition-representable-transformations-property}，
$\mathcal{S}_F \to \mathcal{X}$ 满且光滑。选取对象
$U \in \Ob((\Sch/S)_{fppf})$ 以及满 \'etale 态射 $U \to F$。
再次应用引理
\ref{algebraic-lemma-composition-representable-transformations-property}，
即得所需性质。
\end{proof}

\noindent
本引理是命题 \ref{algebraic-proposition-algebraic-stack-no-automorphisms}
的推广。

\begin{lemma}
\label{algebraic-lemma-characterize-representable-by-algebraic-spaces}
设 $S$ 是 $\Sch_{fppf}$ 中的概形。设
$f : \mathcal{X} \to \mathcal{Y}$ 是 $S$ 上代数叠的 $1$-态射。
以下各条件等价：
\begin{enumerate}
\item 对 $U \in \Ob((\Sch/S)_{fppf})$，函子
$f : \mathcal{X}_U \to \mathcal{Y}_U$ 是忠实的；
\item 函子 $f$ 是忠实的；
\item $f$ 可由代数空间表示。
\end{enumerate}
\end{lemma}

\begin{proof}
由群胚纤维化范畴的 $1$-态射的一般性质，(1) 与 (2) 等价；
参见《范畴》引理 \ref{categories-lemma-equivalence-fibred-categories}。
由引理 \ref{algebraic-lemma-criterion-map-representable-spaces-fibred-in-groupoids}，
(3) 蕴含 (2)。最后，假设 (2) 成立。设 $U$ 是概形，并设
$y \in \Ob(\mathcal{Y}_U)$。需要证明
$$
\mathcal{W} = (\Sch/U)_{fppf} \times_{y, \mathcal{Y}} \mathcal{X}
$$
可由 $U$ 上的代数空间表示。由于 $(\Sch/U)_{fppf}$ 是代数叠，
引理 \ref{algebraic-lemma-2-fibre-product} 说明 $\mathcal{W}$ 是代数叠。
另一方面，把 $\mathcal{W}$ 的对象明确描述为三元组
$(V, x, \alpha : y(V) \to f(x))$，再利用 $f$ 的忠实性，
可知 $\mathcal{W}$ 的纤维范畴都是集合胚。因此，命题
\ref{algebraic-proposition-algebraic-stack-no-automorphisms}
保证 $\mathcal{W}$ 可由代数空间表示。
\end{proof}

\begin{lemma}
\label{algebraic-lemma-smooth-surjective-morphism-implies-algebraic}
设 $S$ 是 $\Sch_{fppf}$ 中的概形。设
$u : \mathcal{U} \to \mathcal{X}$ 是 $(\Sch/S)_{fppf}$ 上
群胚叠的 $1$-态射。若
\begin{enumerate}
\item $\mathcal{U}$ 可由代数空间表示；
\item $u$ 可由代数空间表示、满且光滑，
\end{enumerate}
则 $\mathcal X$ 是 $S$ 上的代数叠。
\end{lemma}

\begin{proof}
需要证明对角态射
$\Delta : \mathcal{X} \to \mathcal{X} \times \mathcal{X}$
可由代数空间表示；参见定义 \ref{algebraic-definition-algebraic-stack}。
给定 $S$ 上的两个概形 $T_1$、$T_2$，用
$\mathcal{T}_i = (\Sch/T_i)_{fppf}$ 表示相应的可表纤维范畴。
设给定 $1$-态射 $f_i : \mathcal{T}_i \to \mathcal{X}$。
根据引理 \ref{algebraic-lemma-representable-diagonal}，只需证明 $2$-纤维积
$\mathcal{T}_1 \times_\mathcal{X} \mathcal{T}_2$
可由代数空间表示。由《叠》引理
\ref{stacks-lemma-2-fibre-product-stacks-in-setoids-over-stack-in-groupoids}，
它无论如何都是集合胚叠。因此，
$\mathcal{T}_1 \times_\mathcal{X} \mathcal{T}_2$
对应于 $(\Sch/S)_{fppf}$ 上的某个层 $F$；参见《叠》引理
\ref{stacks-lemma-stack-in-setoids-characterize}。
设 $U$ 是表示 $\mathcal{U}$ 的代数空间。根据假设，
$$
\mathcal{T}_i' = \mathcal{U} \times_{u, \mathcal{X}, f_i} \mathcal{T}_i
$$
可由 $S$ 上的代数空间 $T'_i$ 表示。因此，
$\mathcal{T}_1' \times_\mathcal{U} \mathcal{T}_2'$
可由代数空间 $T'_1 \times_U T'_2$ 表示。考虑交换图
$$
\xymatrix{
&
\mathcal{T}_1 \times_{\mathcal X} \mathcal{T}_2 \ar[rr]\ar'[d][dd] & &
\mathcal{T}_1 \ar[dd] \\
\mathcal{T}_1' \times_\mathcal{U} \mathcal{T}_2' \ar[ur]\ar[rr]\ar[dd] & &
\mathcal{T}_1' \ar[ur]\ar[dd] \\
&
\mathcal{T}_2 \ar'[r][rr] & &
\mathcal X \\
\mathcal{T}_2' \ar[rr]\ar[ur] & &
\mathcal{U} \ar[ur] }
$$
在此图中，底面、右面、背面和正面方块都是 $2$-纤维积。于是形式论证
说明，态射
$\mathcal{T}_1' \times_\mathcal{U} \mathcal{T}_2' \to
\mathcal{T}_1 \times_{\mathcal X} \mathcal{T}_2$
是 $\mathcal{U} \to \mathcal{X}$ 的“基变换”；更确切地，图
$$
\xymatrix{
\mathcal{T}_1' \times_\mathcal{U} \mathcal{T}_2' \ar[d] \ar[r] &
\mathcal{U} \ar[d] \\
\mathcal{T}_1 \times_{\mathcal X} \mathcal{T}_2 \ar[r] &
\mathcal{X}
}
$$
是 $2$-纤维积方块。因此，$T'_1 \times_U T'_2 \to F$
可由代数空间表示、光滑且满；参见引理
\ref{algebraic-lemma-map-fibred-setoids-representable-algebraic-spaces}、
\ref{algebraic-lemma-base-change-representable-by-spaces}、
\ref{algebraic-lemma-map-fibred-setoids-property} 与
\ref{algebraic-lemma-base-change-representable-transformations-property}。
于是由《自举》定理 \ref{bootstrap-theorem-final-bootstrap}，
$F$ 是代数空间，证毕。
\end{proof}

\noindent
引理 \ref{algebraic-lemma-smooth-surjective-morphism-implies-algebraic}
的一个应用是：代数叠上的代数空间仍是代数叠。这对应于《自举》引理
\ref{bootstrap-lemma-representable-by-spaces-over-space}。
事实上，只需假设态射 $\mathcal{X} \to \mathcal{Y}$
是“代数的”；参见《可表性判据》引理
\ref{criteria-lemma-algebraic-morphism-to-algebraic}。

\begin{lemma}
\label{algebraic-lemma-representable-morphism-to-algebraic}
设 $S$ 是 $\Sch_{fppf}$ 中的概形。设
$\mathcal{X} \to \mathcal{Y}$ 是 $(\Sch/S)_{fppf}$ 上群胚叠的态射。
假设
\begin{enumerate}
\item $\mathcal{X} \to \mathcal{Y}$ 可由代数空间表示；
\item $\mathcal{Y}$ 是 $S$ 上的代数叠。
\end{enumerate}
则 $\mathcal{X}$ 是 $S$ 上的代数叠。
\end{lemma}

\begin{proof}
设 $\mathcal{V} \to \mathcal{Y}$ 是从可表群胚叠到
$\mathcal{Y}$ 的满光滑 $1$-态射。由定义
\ref{algebraic-definition-algebraic-stack}，这样的态射存在。于是 $2$-纤维积
$\mathcal{U} = \mathcal{V} \times_{\mathcal Y} \mathcal X$
由引理 \ref{algebraic-lemma-base-change-by-space-representable-by-space}
可由代数空间表示。$1$-态射 $\mathcal{U} \to \mathcal X$
可由代数空间表示、光滑且满；参见引理
\ref{algebraic-lemma-base-change-representable-by-spaces} 与
\ref{algebraic-lemma-base-change-representable-transformations-property}。
由引理 \ref{algebraic-lemma-smooth-surjective-morphism-implies-algebraic}，
$\mathcal{X}$ 是代数叠。
\end{proof}

\begin{lemma}
\label{algebraic-lemma-open-fibred-category-is-algebraic}
\begin{reference}
Matthew Emerton 在日期为 2016 年 6 月 15 日的一封电子邮件中指出，
可以去掉 $j$ 为单态射这一假设
\end{reference}
设 $S$ 是 $\Sch_{fppf}$ 中的概形。设
$j : \mathcal X \to \mathcal Y$ 是 $(\Sch/S)_{fppf}$ 上
群胚纤维化范畴的 $1$-态射。假设 $j$ 可由代数空间表示。
若 $\mathcal{Y}$ 是群胚叠（相应地，代数叠），则
$\mathcal{X}$ 也是。
\end{lemma}

\begin{proof}
由引理 \ref{algebraic-lemma-representable-morphism-to-algebraic}，
关于代数叠的陈述可由关于群胚叠的陈述推出。若 $j$ 可由代数空间表示，
则 $j$ 在各纤维范畴上忠实，并且对每个 $U$ 与每个
$y \in \Ob(\mathcal{Y}_U)$，预层
$$
(h : V \to U)
\longmapsto
% P11-E0075：冻结源在此使用了未定义的 f(x)，而本引理只定义了 j。
\{(x, \phi) \mid x \in \Ob(\mathcal{X}_V), \phi : h^*y \to f(x)\}/\cong
$$
是 $U$ 上的代数空间；见引理
\ref{algebraic-lemma-criterion-map-representable-spaces-fibred-in-groupoids}。
特别地，该预层是层，故结论由《叠》引理
\ref{stacks-lemma-relative-sheaf-over-stack-is-stack} 得出。
\end{proof}



\section{从代数叠到表现}
\label{algebraic-section-stack-to-presentation}

\noindent
给定 $S$ 上的代数叠，可以得到 $S$ 上代数空间中的一个群胚，
其相应的商叠就是该代数叠。

\medskip\noindent
回忆，若 $(U, R, s, t, c)$ 是 $S$ 上代数空间中的群胚，则
$[U/R]$ 表示与该数据相应的商叠；参见《代数空间中的群胚》定义
\ref{spaces-groupoids-definition-quotient-stack}。
一般来说，$[U/R]$ {\bf 不是}代数叠。特别地，下一个引理中出现的
叠 $[U/R]$ 一般不是代数的。

\begin{lemma}
\label{algebraic-lemma-map-space-into-stack}
设 $S$ 是 $\Sch_{fppf}$ 中的概形。设 $\mathcal{X}$ 是
$S$ 上的代数叠。设 $\mathcal{U}$ 是 $S$ 上可由代数空间表示的
代数叠。设 $f : \mathcal{U} \to \mathcal{X}$ 是 1-态射。则
\begin{enumerate}
\item $2$-纤维积
$\mathcal{R} = \mathcal{U} \times_{f, \mathcal{X}, f} \mathcal{U}$
可由代数空间表示；
\item 存在典范等价
$$
\mathcal{U} \times_{f, \mathcal{X}, f} \mathcal{U}
\times_{f, \mathcal{X}, f} \mathcal{U} =
\mathcal{R} \times_{\text{pr}_1, \mathcal{U}, \text{pr}_0} \mathcal{R},
$$
\item 投影 $\text{pr}_{02}$ 通过 (2) 诱导 $1$-态射
$$
\text{pr}_{02} :
\mathcal{R} \times_{\text{pr}_1, \mathcal{U}, \text{pr}_0} \mathcal{R}
\longrightarrow
\mathcal{R}
$$
\item 设 $U$、$R$ 是分别表示 $\mathcal{U}, \mathcal{R}$
的代数空间，并设 $t, s : R \to U$ 与
$c : R \times_{s, U, t} R \to R$ 是分别对应于上述 $1$-态射
$\text{pr}_0, \text{pr}_1 : \mathcal{R} \to \mathcal{U}$
以及
$\text{pr}_{02} :
\mathcal{R} \times_{\text{pr}_1, \mathcal{U}, \text{pr}_0} \mathcal{R} \to
\mathcal{R}$ 的态射，则五元组 $(U, R, s, t, c)$ 是
$S$ 上代数空间中的群胚；
\item 态射 $f$ 诱导 $(\Sch/S)_{fppf}$ 上群胚叠的典范
$1$-态射
$f_{can} : [U/R] \to \mathcal{X}$
；
\item $1$-态射 $f_{can} : [U/R] \to \mathcal{X}$ 全忠实。
\end{enumerate}
\end{lemma}

\begin{proof}
(1) 的证明。按定义，$\Delta_\mathcal{X}$ 可由代数空间表示，
故应用引理 \ref{algebraic-lemma-representable-diagonal} 可知
$\mathcal{U} \to \mathcal{X}$ 可由代数空间表示。因此结论由引理
\ref{algebraic-lemma-base-change-by-space-representable-by-space} 得出。

\medskip\noindent
设 $T$ 是 $S$ 上的概形。根据 $2$-纤维积的构造（见《范畴》引理
\ref{categories-lemma-2-product-categories-over-C}），纤维范畴
$\mathcal{R}_T$ 的对象是三元组 $(a, b, \alpha)$，其中
$a, b \in \Ob(\mathcal{U}_T)$，而
$\alpha : f(a) \to f(b)$ 是纤维范畴 $\mathcal{X}_T$ 中的态射。

\medskip\noindent
(2) 的证明。反复应用《范畴》引理
\ref{categories-lemma-associativity-2-fibre-product} 与
\ref{categories-lemma-2-fibre-product-erase-factor}
即可得到该等价。把
$\mathcal{U} \times_\mathcal{X} \mathcal{U} \times_\mathcal{X} \mathcal{U}$
与
$(\mathcal{U} \times_\mathcal{X} \mathcal{U})
\times_\mathcal{X} \mathcal{U}$.
等同。若 $T$ 是 $S$ 上的概形，则该等价在 $T$ 上的纤维范畴中，
把左端对象 $((a, b, \alpha), c, \beta)$ 送到右端对象
$((a, b, \alpha), (b, c, \beta))$。

\medskip\noindent
(3) 的证明。《范畴》引理
\ref{categories-lemma-triple-2-fibre-product-pr02}
的证明构造了 $1$-态射 $\text{pr}_{02}$。根据上面对纤维范畴对象
的描述，$((a, b, \alpha), (b, c, \beta))$ 被送到
$(a, c, \beta \circ \alpha)$。

\medskip\noindent
遗憾的是，这与我们的群胚约定{\it 不相容}：我们总令
$j = (t, s) : R \to U$，并把 $R$ 的 $T$-值点 $r$
“看成”态射 $r : s(r) \to t(r)$。不过，这不影响 (4) 的证明，
因为群胚的反范畴仍是群胚。但在 (5) 的证明中，正是这一点导致
下方显示公式中出现逆。

\medskip\noindent
(4) 的证明。回忆，层 $U$ 与层
$T \mapsto \Ob(\mathcal{U}_T)/\!\cong$ 同构，$R$ 也类似；
参见引理 \ref{algebraic-lemma-characterize-representable-by-space}。
由《范畴》引理
\ref{categories-lemma-category-fibred-setoids-presheaves-products}，
这一描述与 $2$-纤维积相容，故类似地有
$\mathcal{R} \times_{\text{pr}_1, \mathcal{U}, \text{pr}_0} \mathcal{R}$
与 $R \times_{s, U, t} R$ 之间的对应。态射
$t, s : R \to U$ 与 $c : R \times_{s, U, t} R \to R$
来自一般等式 (\ref{algebraic-equation-morphisms-spaces})。明确地说，
% P11-E0078：冻结源在此重复 \text{pr}_0；后续结构所用三个投影为 \text{pr}_0、\text{pr}_1、\text{pr}_{02}。
这些映射是由 $\text{pr}_0$、$\text{pr}_0$、$\text{pr}_{02}$
作用于纤维范畴的对象同构类而产生的函子变换。因此，为证明所得结构
是代数空间中的群胚，只需证明对 $S$ 上的每个概形 $T$，结构
$$
(\Ob(\mathcal{U}_T)/\!\cong,
\Ob(\mathcal{R}_T)/\!\cong,
\text{pr}_1, \text{pr}_0, \text{pr}_{02})
$$
是群胚；这由上面对 $\mathcal{R}_T$ 的对象的描述立即可见。

\medskip\noindent
(5) 的证明。我们最终将应用《代数空间中的群胚》引理
\ref{spaces-groupoids-lemma-quotient-stack-2-coequalizer}
得到函子 $[U/R] \to \mathcal{X}$。考虑 $1$-态射
$f : \mathcal{U} \to \mathcal{X}$。由 $\mathcal{R}$
作为 $2$-纤维积的定义，有 $2$-箭头
$\tau : f \circ \text{pr}_1 \to f \circ \text{pr}_0$。
具体地，在 $\mathcal{R}$ 于 $T$ 上的对象 $(a, b, \alpha)$ 上，
它是映射 $\alpha^{-1} : b \to a$。我们断言
$$
\tau \circ \text{id}_{\text{pr}_{02}} =
(\tau \star \text{id}_{\text{pr}_0})
\circ
(\tau \star \text{id}_{\text{pr}_1}).
$$
该恒等式表示：给定
$\mathcal{R} \times_{\text{pr}_1, \mathcal{U}, \text{pr}_0} \mathcal{R}$
在 $T$ 上的对象 $((a, b, \alpha), (b, c, \beta))$，则
$$
\xymatrix{
c \ar[r]^{\beta^{-1}} & b \ar[r]^{\alpha^{-1}} & a
}
$$
% P11-E0079：冻结源所写箭头的源、靶与前一复合相反；此处保留原式待 canon 核验。
的复合等于箭头 $(\beta \circ \alpha)^{-1} : a \to c$。这显然成立，
故断言得证。由此可见，对结构
$(\mathcal{U}, \mathcal{R}, \text{pr}_0, \text{pr}_1, \text{pr}_{02})$
以及 $1$-态射 $f$ 和 $2$-态射 $\tau$，《代数空间中的群胚》引理
\ref{spaces-groupoids-lemma-quotient-stack-2-coequalizer}
的全部假设都成立。不过，为应用该引理，我们需要证明它对结构
$(\mathcal{S}_U, \mathcal{S}_R, s, t, c)$ 及适当态射也成立。

\medskip\noindent
此时应当存在某种一般的抽象论证，在二者之间传递这些数据，
但它似乎相当冗长。我们改用如下技巧。选取典范等价
$j : \mathcal{U} \to \mathcal{S}_U$ 的一个拟逆
$j^{-1} : \mathcal{S}_U \to \mathcal{U}$；该等价来自
$U(T) = \Ob(\mathcal{U}_T)/\!\!\cong$。这仅意味着：
对每个概形 $T/S$ 以及每个对象 $a \in \mathcal{U}_T$，
我们从其同构类中选定一个特定元素，即 $j^{-1}(j(a))$。
因此，利用 $j^{-1}$ 可把 $\mathcal{S}_U$ 看成 $\mathcal{U}$
的子范畴。选定这个子范畴后，可以考虑 $\mathcal{R}_T$
中满足如下条件的对象 $(a, b, \alpha)$：$a,b$ 是
$(\mathcal{S}_U)_T$ 的对象，即
$j^{-1}(j(a)) = a$ 且 $j^{-1}(j(b)) = b$。显然，它们构成
$\mathcal{R}$ 的一个子范畴，并通过典范等价
$\mathcal{R} \to \mathcal{S}_R$ 同构地映到 $\mathcal{S}_R$。
此外，这显然与形成 $2$-纤维积
$\mathcal{R} \times_{\text{pr}_1, \mathcal{U}, \text{pr}_0} \mathcal{R}$.
相容。因此，只需把 $f$ 限制到 $\mathcal{S}_U$，并把 $\tau$
限制为函子 $\mathcal{S}_R \to \mathcal{X}$ 之间的变换。
于是《代数空间中的群胚》引理
\ref{spaces-groupoids-lemma-quotient-stack-2-coequalizer}
中的显示等式显然成立，因为在限制到子范畴
$\mathcal{S}_{R \times_{s, U, t} R}$ 之前，它甚至已经作为函子
$\mathcal{R} \times_{\text{pr}_1, \mathcal{U}, \text{pr}_0} \mathcal{R}
\to \mathcal{X}$ 的变换等式成立。

\medskip\noindent
这证明《代数空间中的群胚》引理
\ref{spaces-groupoids-lemma-quotient-stack-2-coequalizer}
适用，从而得到所需的叠态射
$f_{can} : [U/R] \to \mathcal{X}$。我们简要说明在这个特殊情形中
$f_{can}$ 如何定义。在 $\mathcal{S}_U$ 于 $T$ 上的对象 $a$ 上，
有 $f_{can}(a) = f(a)$；这里通过上面选定的嵌入，把
$\mathcal{S}_U \subset \mathcal{U}$。若 $a,b$ 是
$\mathcal{S}_U$ 在 $T$ 上的对象，则按定义，$[U/R]$ 中的态射
$\varphi : a \to b$ 是 $\mathcal{R}$ 在 $T$ 上形如
$\varphi = (b, a, \alpha)$ 的对象。（注意次序交换。）
《代数空间中的群胚》引理
\ref{spaces-groupoids-lemma-quotient-stack-2-coequalizer}
的证明中的规则为
\begin{equation}
\label{algebraic-equation-on-morphisms}
f_{can}(\varphi) = \Big(f(a) \xrightarrow{\alpha^{-1}} f(b)\Big).
\end{equation}
(6) 的证明。$[U/R]$ 与 $\mathcal{X}$ 都是叠。因此，给定概形
$T/S$ 以及 $[U/R]$ 在 $T$ 上的对象 $a,b$，可得到
fppf 层的变换
$$
\mathit{Isom}(a, b) \longrightarrow \mathit{Isom}(f_{can}(a), f_{can}(b))
$$
，它位于 $(\Sch/T)_{fppf}$ 上。需要证明这是同构。可以在 $T$ 上
fppf 局部地工作，故可假设 $a,b$ 来自态射 $a, b : T \to U$。
通过上述嵌入 $\mathcal{S}_U \subset \mathcal{U}$，还可把
$a,b$ 看成 $\mathcal{U}$ 在 $T$ 上的对象。我们已在
《代数空间中的群胚》引理
\ref{spaces-groupoids-lemma-quotient-stack-morphisms} 中看到，
左端的层由代数空间
$$
R \times_{(t, s), U \times_S U, (b, a)} T
$$
表示于 $T$ 上。另一方面，由《叠》引理
\ref{stacks-lemma-isom-as-2-fibre-product}，右端等于与如下集合胚叠
相应的层：
$$
\mathcal{X}
\times_{\mathcal{X} \times \mathcal{X}, (f \circ b, f \circ a)} T =
\mathcal{X}
\times_{\mathcal{X} \times \mathcal{X}, (f, f)}
(\mathcal{U} \times \mathcal{U})
\times_{\mathcal{U} \times \mathcal{U}, (b, a)} T =
\mathcal{R}
\times_{(\text{pr}_0, \text{pr}_1), \mathcal{U} \times \mathcal{U}, (b, a)} T
$$
，它由上面显示的纤维积表示。至此已经证明两个
$\mathit{Isom}$ 层同构。由公式 (\ref{algebraic-equation-on-morphisms})，
我们的 $1$-态射 $f_{can} : [U/R] \to \mathcal{X}$
在 $\mathit{Isom}$ 层上诱导这个同构。
\end{proof}

\noindent
可以使用上述高度抽象的引理来构造表现。

\begin{lemma}
\label{algebraic-lemma-stack-presentation}
设 $S$ 是 $\Sch_{fppf}$ 中的概形。设 $\mathcal{X}$ 是
$S$ 上的代数叠。设 $U$ 是 $S$ 上的代数空间，并设
$f : \mathcal{S}_U \to \mathcal{X}$ 是满光滑态射。
把引理 \ref{algebraic-lemma-map-space-into-stack} 应用于 $U$ 与 $f$，
得到代数空间中的群胚 $(U, R, s, t, c)$ 以及
$f_{can} : [U/R] \to \mathcal{X}$。则
\begin{enumerate}
\item 态射 $s$、$t$ 光滑；
\item $1$-态射 $f_{can} : [U/R] \to \mathcal{X}$ 是等价。
\end{enumerate}
\end{lemma}

\begin{proof}
由引理 \ref{algebraic-lemma-property-morphism-equivalent} 与
\ref{algebraic-lemma-map-presheaves-representable-by-spaces-transformation-property}，
态射 $s,t$ 光滑。由于 $1$-态射 $f$ 光滑且满，显然，对任意概形
$T$ 以及任意对象 $a \in \Ob(\mathcal{X}_T)$，存在光滑满态射
$T' \to T$，使得 $a|_T'$ 来自 $[U/R]_{T'}$ 的对象。
由于 $f_{can} : [U/R] \to \mathcal{X}$ 全忠实，而两端的对象下降
数据都有效，故 $[U/R] \to \mathcal{X}$ 本质满；参见《叠》引理
\ref{stacks-lemma-characterize-essentially-surjective-when-ff}。
\end{proof}

\begin{remark}
\label{algebraic-remark-flat-fp-presentation}
若对引理 \ref{algebraic-lemma-stack-presentation} 中的态射
$f : \mathcal{S}_U \to \mathcal{X}$ 只假设其满、平坦且局部有限表示，
则 $f_{can} : [U/R] \to \mathcal{X}$ 仍是等价。此时态射
$s$、$t$ 平坦且局部有限表示，但一般而言当然不光滑。
\end{remark}

\noindent
引理 \ref{algebraic-lemma-stack-presentation} 提示如下定义。

\begin{definition}
\label{algebraic-definition-smooth-groupoid}
设 $S$ 是概形。设 $B$ 是 $S$ 上的代数空间。设
$(U, R, s, t, c)$ 是 $B$ 上代数空间中的群胚。若
$s, t : R \to U$ 是代数空间的光滑态射，则称
$(U, R, s, t, c)$ 为{\it 光滑群胚}\footnote{这个术语可能稍显
混淆：它并不蕴含 $[U/R]$ 在任何对象上光滑。}。
\end{definition}

\begin{definition}
\label{algebraic-definition-presentation}
设 $\mathcal{X}$ 是 $S$ 上的代数叠。$\mathcal{X}$ 的一个
{\it 表现}由 $S$ 上代数空间中的光滑群胚
$(U, R, s, t, c)$ 以及等价 $f : [U/R] \to \mathcal{X}$ 给出。
\end{definition}

\noindent
上文已经看到，每个代数叠都有表现。下一步将证明，$S$ 上代数空间中的
每个光滑群胚都给出一个代数叠。


\section{与光滑群胚相应的代数叠}
\label{algebraic-section-smooth-groupoid-gives-algebraic-stack}

\noindent
本节从代数空间中的光滑群胚出发，证明相应的商叠是代数叠。

\begin{lemma}
\label{algebraic-lemma-diagonal-quotient-stack}
设 $S$ 是 $\Sch_{fppf}$ 中的概形。设 $(U, R, s, t, c)$ 是
$S$ 上代数空间中的群胚。则 $[U/R]$ 的对角态射可由代数空间表示。
\end{lemma}

\begin{proof}
只需证明 $\mathit{Isom}$ 层是代数空间；见引理
\ref{algebraic-lemma-representable-diagonal}。这由《自举》引理
\ref{bootstrap-lemma-quotient-stack-isom} 得出。
\end{proof}

\begin{lemma}
\label{algebraic-lemma-smooth-quotient-smooth-presentation}
设 $S$ 是 $\Sch_{fppf}$ 中的概形。设 $(U, R, s, t, c)$ 是
$S$ 上代数空间中的光滑群胚。则态射
$\mathcal{S}_U \to [U/R]$ 光滑且满。
\end{lemma}

\begin{proof}
设 $T$ 是概形，并设 $x : (\Sch/T)_{fppf} \to [U/R]$
是 $1$-态射。需要证明投影
$$
\mathcal{S}_U \times_{[U/R]} (\Sch/T)_{fppf}
\longrightarrow
(\Sch/T)_{fppf}
$$
满且光滑。由引理 \ref{algebraic-lemma-diagonal-quotient-stack} 与
\ref{algebraic-lemma-representable-diagonal} 已知，左端可由代数空间 $F$
表示。因此，需要证明相应的代数空间态射 $F \to T$ 满且光滑。
我们研究的是在靶的 fppf 拓扑中为局部的代数空间态射性质，
故可以在 $T$ 上 fppf 局部地验证。按构造，存在 $T$ 的 fppf 覆盖
$\{T_i \to T\}$，使得 $x|_{(\Sch/T_i)_{fppf}}$ 来自态射
$x_i : T_i \to U$。（注意，$F \times_T T_i$ 表示 $2$-纤维积
$\mathcal{S}_U \times_{[U/R]} (\Sch/T_i)_{fppf}$，故一切都与
沿 $T_i \to T$ 的基变换相容。）因此可假设 $x$ 来自
$x : T \to U$。此时有
$$
\mathcal{S}_U \times_{[U/R]} (\Sch/T)_{fppf}
=
(\mathcal{S}_U \times_{[U/R]} \mathcal{S}_U)
\times_{\mathcal{S}_U} (\Sch/T)_{fppf}
=
\mathcal{S}_R \times_{\mathcal{S}_U} (\Sch/T)_{fppf}
$$
第一个等式由《范畴》引理
\ref{categories-lemma-2-fibre-product-erase-factor} 得出，
第二个等式由《代数空间中的群胚》引理
\ref{spaces-groupoids-lemma-quotient-stack-2-cartesian} 得出。
显然，最后一个 $2$-纤维积由代数空间
$F = R \times_{s, U, x} T$ 表示，而投影
$R \times_{s, U, x} T \to T$ 是代数空间的光滑态射
$s : R \to U$ 的基变换，故光滑。它也满，因为 $s$ 有截面
（即群胚的恒等元 $e : U \to R$）。引理得证。
\end{proof}

\noindent
下面是本节的主要结果。

\begin{theorem}
\label{algebraic-theorem-smooth-groupoid-gives-algebraic-stack}
设 $S$ 是 $\Sch_{fppf}$ 中的概形。设 $(U, R, s, t, c)$ 是
$S$ 上代数空间中的光滑群胚。则商叠 $[U/R]$ 是 $S$ 上的代数叠。
\end{theorem}

\begin{proof}
验证定义 \ref{algebraic-definition-algebraic-stack} 的三个条件。
按构造，$[U/R]$ 是群胚叠，这就是第一个条件。

\medskip\noindent
第二个条件由更强的引理 \ref{algebraic-lemma-diagonal-quotient-stack} 得出。

\medskip\noindent
最后，需要证明存在 $S$ 上的概形 $W$ 以及满光滑 $1$-态射
% P11-E0080：冻结源在此使用未定义的 \mathcal{X}；本定理上下文的商叠为 [U/R]。
$(\Sch/W)_{fppf} \longrightarrow \mathcal{X}$.
先选取 $W \in \Ob((\Sch/S)_{fppf})$ 以及满 \'etale 态射
$W \to U$。注意，这给出集合纤维化范畴的满 \'etale 态射
$\mathcal{S}_W \to \mathcal{S}_U$；见引理
\ref{algebraic-lemma-map-presheaves-representable-by-spaces-transformation-property}。
当然，$\mathcal{S}_W \to \mathcal{S}_U$ 也满且光滑；
见引理 \ref{algebraic-lemma-representable-transformations-property-implication}。
因此，结合引理 \ref{algebraic-lemma-smooth-quotient-smooth-presentation} 与
\ref{algebraic-lemma-composition-representable-transformations-property}，
$\mathcal{S}_W \to \mathcal{S}_U \to [U/R]$ 满且光滑。
\end{proof}





\section{更换大位点}
\label{algebraic-section-change-big-site}

\noindent
本节简要讨论更换大位点时发生的情形。要点是，我们总能任意扩大大位点；
因此，可以假设想要考虑的任意概形集合都包含在我们研究代数空间所用的
大 fppf 位点中。我们建议读者跳过本节。

\medskip\noindent
叠的拉回定义见《叠》第 \ref{stacks-section-inverse-image} 节。

\begin{lemma}
\label{algebraic-lemma-change-big-site}
设给定大位点 $\Sch_{fppf}$ 与 $\Sch'_{fppf}$。假设
$\Sch_{fppf}$ 包含于 $\Sch'_{fppf}$；参见《概形上的拓扑》第
\ref{topologies-section-change-alpha} 节。设 $S$ 是
$\Sch_{fppf}$ 的对象。令
$f : (\Sch'/S)_{fppf} \to (\Sch/S)_{fppf}$ 为与包含函子
$u : (\Sch/S)_{fppf} \to (\Sch'/S)_{fppf}$ 相应的位点态射。
设 $\mathcal{X}$ 是 $(\Sch/S)_{fppf}$ 上的群胚叠。
\begin{enumerate}
\item 若 $\mathcal{X}$ 可由某个
$X \in \Ob((\Sch/S)_{fppf})$ 表示，则
$f^{-1}\mathcal{X}$ 也可表；事实上，它由同一个概形 $X$ 表示，
只是现在将该概形看成 $(\Sch'/S)_{fppf}$ 的对象；
\item 若 $\mathcal{X}$ 可由代数空间
$F \in \Sh((\Sch/S)_{fppf})$ 表示，则
$f^{-1}\mathcal{X}$ 可由代数空间 $f^{-1}F$ 表示；
\item 若 $\mathcal{X}$ 是代数叠，则 $f^{-1}\mathcal{X}$ 是代数叠；
\item 若 $\mathcal{X}$ 是 Deligne--Mumford 叠，则
$f^{-1}\mathcal{X}$ 也是 Deligne--Mumford 叠。
\end{enumerate}
\end{lemma}

\begin{proof}
证明 (3)。由引理 \ref{algebraic-lemma-stack-presentation}，可对某个
代数空间中的光滑群胚 $(U, R, s, t, c)$ 写成
$\mathcal{X} = [U/R]$。由《代数空间中的群胚》引理
\ref{spaces-groupoids-lemma-quotient-stack-change-big-site}，
有 $f^{-1}[U/R] = [f^{-1}U/f^{-1}R]$。当然，
$(f^{-1}U, f^{-1}R, f^{-1}s, f^{-1}t, f^{-1}c)$
也是代数空间中的光滑群胚。因此 (3) 得证。

\medskip\noindent
其他情形 (1)、(2)、(4) 各自都意味着 $\mathcal{X}$ 具有某种特定类型
的表现 $[U/R]$，因而转化为
$f^{-1}\mathcal{X} = [f^{-1}U/f^{-1}R]$ 的同类型表现。
引理得证。
\end{proof}

\noindent
一般而言，大位点上的代数空间限制到小位点后未必仍是小位点上的
代数空间（仅基数方面的原因就足以导致失败）。因此，这类简单引理
只能用来扩大基范畴，绝不能用来缩小它。

\begin{lemma}
\label{algebraic-lemma-fully-faithful}
假设 $\Sch_{fppf}$ 包含于 $\Sch'_{fppf}$。设 $S$ 是
$\Sch_{fppf}$ 的对象。用 $\textit{Algebraic-Stacks}/S$
表示用 $\Sch_{fppf}$ 定义的 $S$ 上代数叠的 $2$-范畴。
类似地，用 $\textit{Algebraic-Stacks}'/S$ 表示用
$\Sch'_{fppf}$ 定义的 $S$ 上代数叠的 $2$-范畴。引理
\ref{algebraic-lemma-change-big-site} 的规则
$\mathcal{X} \mapsto f^{-1}\mathcal{X}$
定义 $2$-范畴的函子
$$
\textit{Algebraic-Stacks}/S \longrightarrow \textit{Algebraic-Stacks}'/S
$$
，并对 $\textit{Algebraic-Stacks}/S$ 的每对对象
$\mathcal{X}, \mathcal{Y}$ 定义态射范畴的等价
$$
\Mor_{\textit{Algebraic-Stacks}/S}(\mathcal{X}, \mathcal{Y})
\longrightarrow
\Mor_{\textit{Algebraic-Stacks}'/S}(f^{-1}\mathcal{X}, f^{-1}\mathcal{Y})
$$
。$\textit{Algebraic-Stacks}'/S$ 的对象 $\mathcal{X}'$
等价于某个 $\textit{Algebraic-Stacks}/S$ 中的 $\mathcal{X}$
所对应的 $f^{-1}\mathcal{X}$，当且仅当它有一个表现
% P11-E0077：冻结源此处写成 \mathcal{X}=[U'/R']，上下文对象为 \mathcal{X}'。
$\mathcal{X} = [U'/R']$，其中 $U', R'$ 分别同构于某些
$U, R \in \textit{Spaces}/S$ 所对应的 $f^{-1}U$、$f^{-1}R$。
\end{lemma}

\begin{proof}
关于态射范畴的陈述是更一般的《叠》引理
\ref{stacks-lemma-bigger-site} 的推论。“本质像”的刻画则由引理
\ref{algebraic-lemma-change-big-site} 的证明中对 $f^{-1}$ 的描述得出。
\end{proof}


\section{更换基概形}
\label{algebraic-section-change-base-scheme}

\noindent
本节简要讨论更换基概形时会发生什么。结论是：给定基概形之间的一个态射
$S \to S'$，任意 $S$ 上的代数叠都可以视为 $S'$ 上的代数叠。

\begin{lemma}
\label{algebraic-lemma-category-of-spaces-over-smaller-base-scheme}
设 $\Sch_{fppf}$ 是一个大 fppf 位点，并设
$S \to S'$ 是该位点中的一个态射。上文
Stacks, Section \ref{stacks-section-localize}
中的构造 A 与构造 B 给出如下 $2$-范畴的同构：
$$
\left\{
\begin{matrix}
S\text{ 上代数叠 }\mathcal{X}\\
\text{所成的 }2\text{-范畴}
\end{matrix}
\right\}
\leftrightarrow
\left\{
\begin{matrix}
\text{二元组 }(\mathcal{X}', f)\text{ 所成的 }2\text{-范畴，其中}\\
\mathcal{X}'\text{ 是 }S'\text{ 上的代数叠，而}\\
f : \mathcal{X}' \to (\Sch/S)_{fppf}\text{ 是 }S'\text{ 上代数叠的态射}
\end{matrix}
\right\}
$$
\end{lemma}

\begin{proof}
该陈述是有意义的，因为函子
$j : (\Sch/S)_{fppf} \to (\Sch/S')_{fppf}$
是与 $(\Sch/S')_{fppf}$ 的对象 $S/S'$ 相联系的局部化函子。由
Stacks, Lemma \ref{stacks-lemma-localize-stacks}
只需证明构造 A 与构造 B 保持代数叠的子范畴。例如，若
$\mathcal{X} = [U/R]$，则把构造 A 应用于 $\mathcal{X}$ 所得到的正是
$\mathcal{X}' = \mathcal{X}$。反之，若 $\mathcal{X}' = [U'/R']$，
则态射 $p$ 诱导代数空间的态射 $U' \to S$ 与 $R' \to S$；于是
$\mathcal{X} = [U'/R']$，但现在将其视为 $S$ 上的叠。因此引理成立。
\end{proof}

\begin{definition}
\label{algebraic-definition-viewed-as}
设 $\Sch_{fppf}$ 是一个大 fppf 位点，并设
$S \to S'$ 是该位点中的一个态射。若
$p : \mathcal{X} \to (\Sch/S)_{fppf}$
是 $S$ 上的代数叠，则所谓把 $\mathcal{X}$
{\it 视为 $S'$ 上的代数叠}，是指代数叠
$$
\mathcal{X} \longrightarrow (\Sch/S')_{fppf}
$$
它由把引理 \ref{algebraic-lemma-category-of-spaces-over-smaller-base-scheme}
中的构造 A 应用于 $\mathcal{X}$ 而得到。
\end{definition}

\noindent
反过来，若从 $S'$ 上的一个代数叠 $\mathcal{X}'$ 出发，并希望得到
$S$ 上的代数叠，则考虑 $2$-纤维积
$$
\mathcal{X}'_S
=
(\Sch/S)_{fppf} \times_{(\Sch/S')_{fppf}} \mathcal{X}'
$$
由引理 \ref{algebraic-lemma-2-fibre-product}，它是 $S'$ 上的代数叠。此外，
它还配有自然的 $1$-态射
$p : \mathcal{X}'_S \to (\Sch/S)_{fppf}$，因而由
引理 \ref{algebraic-lemma-category-of-spaces-over-smaller-base-scheme}，
它以典范方式对应于 $S$ 上的一个代数叠。

\begin{definition}
\label{algebraic-definition-change-of-base}
设 $\Sch_{fppf}$ 是一个大 fppf 位点，设
$S \to S'$ 是该位点中的一个态射，并设 $\mathcal{X}'$ 是
$S'$ 上的一个代数叠。所谓 $\mathcal{X}'$ 的{\it 基变换}，
是指上面所述的 $S$ 上代数叠 $\mathcal{X}'_S$。
\end{definition}
% Generated mechanically from the frozen examples-stacks.tex target.
% Do not edit this generated file; edit the bound locale source instead.

% OK, start here.
%

\setcounter{chapter}{94}
\chapter{叠的例子}



\phantomsection
\label{examples-stacks-section-phantom}


\section{引言}
\label{examples-stacks-section-introduction}

\noindent
本章讨论代数几何中叠的一些例子。其中有些是代数叠，有些则不是。
我们将在后续章节讨论哪些例子是代数叠。因此，本章主要关注下降条件；
例如可参见 \cite{Vis2}。

\medskip\noindent
本章中的某些记号、约定和术语并不自然，对经验丰富的读者而言甚至可能
显得前后颠倒。这是有意为之。解释参见《Quot 空间与 Hilbert 空间》第
\ref{quot-section-conventions} 节。




\section{记号}
\label{examples-stacks-section-notation}

\noindent
本章固定一个合适的大 fppf 位点 $\Sch_{fppf}$，如《概形上的拓扑》定义
\ref{topologies-definition-big-fppf-site} 所述。因此，除非另有明示，
所有概形都是 $\Sch_{fppf}$ 的对象。我们始终相对于包含在
$\Sch_{fppf}$ 中的一个基 $S$ 工作，并使用大 fppf 位点
$(\Sch/S)_{fppf}$；参见《概形上的拓扑》定义
\ref{topologies-definition-big-small-fppf}。取
$S = \Spec(\mathbf{Z})$ 即可恢复绝对情形。




\section{叠的例子}
\label{examples-stacks-section-examples-stacks}

\noindent
我们先给出 $(\Sch/S)_{fppf}$ 上叠的若干重要例子。





\section{拟凝聚层}
\label{examples-stacks-section-stack-of-quasi-coherent-sheaves}

\noindent
如下定义范畴 $\QCohstack$：
\begin{enumerate}
\item $\QCohstack$ 的对象是二元组 $(X, \mathcal{F})$，其中
$X/S$ 是 $(\Sch/S)_{fppf}$ 的对象，而 $\mathcal{F}$ 是拟凝聚
$\mathcal{O}_X$-模；
\item 态射 $(f, \varphi) : (Y, \mathcal{G}) \to (X, \mathcal{F})$
是一个二元组，由 $S$ 上概形的态射 $f : Y \to X$ 以及 $f$-映射
（参见《空间上的层》第
\ref{sheaves-section-ringed-spaces-functoriality-modules} 节）
$\varphi : \mathcal{F} \to \mathcal{G}$ 组成；
\item 态射
$$
(Z, \mathcal{H}) \xrightarrow{(g, \psi)}
(Y, \mathcal{G}) \xrightarrow{(f, \phi)} (X, \mathcal{F})
$$
的复合是 $(f \circ g, \psi \circ \phi)$，其中 $\psi \circ \phi$
是 $f$-映射的复合。
\end{enumerate}
于是 $\QCohstack$ 是一个范畴，并且
$$
p : \QCohstack \to (\Sch/S)_{fppf},
\quad
(X, \mathcal{F}) \mapsto X
$$
是函子。注意，$\QCohstack$ 在概形 $X$ 上的纤维范畴，是拟凝聚
$\mathcal{O}_X$-模范畴 $\QCoh(\mathcal{O}_X)$ 的反范畴。
为供后文使用，我们还指出：给定
$(X, \mathcal{F}), (Y, \mathcal{G}) \in \Ob(\QCohstack)$
，有
\begin{equation}
\label{examples-stacks-equation-morphisms-qcoh}
\Mor_{\QCohstack}((Y, \mathcal{G}), (X, \mathcal{F}))
=
\coprod\nolimits_{f \in \Mor_S(Y, X)}
\Mor_{\QCoh(\mathcal{O}_Y)}(f^*\mathcal{F}, \mathcal{G})
\end{equation}
关于模的 $f$-映射的讨论，参见《空间上的层》第
\ref{sheaves-section-ringed-spaces-functoriality-modules} 节。

\medskip\noindent
范畴 $\QCohstack$ 不是 $(\Sch/S)_{fppf}$ 上的叠，因为它的对象全体
是一个真类。另一方面，我们将看到它确实满足叠的全部公理。
第 \ref{examples-stacks-section-stack-of-finitely-generated-quasi-coherent-sheaves} 节
将规避这一集合论问题。

\begin{lemma}
\label{examples-stacks-lemma-quasi-coherent-strongly-cartesian}
$\QCohstack$ 的态射
$(f, \varphi) : (Y, \mathcal{G}) \to (X, \mathcal{F})$
是强笛卡尔态射，当且仅当映射 $\varphi$ 诱导同构
$f^*\mathcal{F} \to \mathcal{G}$。
\end{lemma}

\begin{proof}
设 $(X, \mathcal{F}) \in \Ob(\QCohstack)$，并设
$f : Y \to X$ 是 $(\Sch/S)_{fppf}$ 的态射。注意，存在典范的
$f$-映射 $c : \mathcal{F} \to f^*\mathcal{F}$，因而得到态射
$(f, c) : (Y, f^*\mathcal{F}) \to (X, \mathcal{F})$.
我们断言 $(f, c)$ 是强笛卡尔态射。事实上，对 $\QCohstack$ 的任意对象
$(Z, \mathcal{H})$，有
\begin{align*}
\Mor_{\QCohstack}((Z, \mathcal{H}), (Y, f^*\mathcal{F}))
& =
\coprod\nolimits_{g \in \Mor_S(Z, Y)}
\Mor_{\QCoh(\mathcal{O}_Z)}(g^*f^*\mathcal{F}, \mathcal{H}) \\
& =
\coprod\nolimits_{g \in \Mor_S(Z, Y)}
\Mor_{\QCoh(\mathcal{O}_Z)}((f \circ g)^*\mathcal{F}, \mathcal{H}) \\
& =
\Mor_{\QCohstack}((Z, \mathcal{H}), (X, \mathcal{F}))
\times_{\Mor_S(Z, X)} \Mor_S(Z, Y)
\end{align*}
这里两次使用了式 (\ref{examples-stacks-equation-morphisms-qcoh})。这证明
$(f, c)$ 满足《范畴》定义 \ref{categories-definition-cartesian-over-C}
的条件，故断言成立。由《范畴》引理
\ref{categories-lemma-composition-cartesian}，同构都是强笛卡尔态射，
而强笛卡尔态射的复合仍是强笛卡尔态射；这证明了引理的充分性。
反过来，给定 $(X, \mathcal{F})$ 与 $f : Y \to X$，若存在以
$(X, \mathcal{F})$ 为目标、提升 $f$ 的强笛卡尔态射，则它必与
$(f, c)$ 同构（参见《范畴》定义
\ref{categories-definition-cartesian-over-C} 后的讨论）。因此引理的必要性也成立。
\end{proof}

\begin{lemma}
\label{examples-stacks-lemma-stack-of-quasi-coherent-sheaves}
函子 $p : \QCohstack \to (\Sch/S)_{fppf}$ 满足《叠》定义
\ref{stacks-definition-stack} 的条件 (1)、(2) 与 (3)。
\end{lemma}

\begin{proof}
由引理 \ref{examples-stacks-lemma-quasi-coherent-strongly-cartesian} 可知，
$\QCohstack$ 是 $(\Sch/S)_{fppf}$ 上的纤维化范畴。给定
$(\Sch/S)_{fppf}$ 的覆盖 $\mathcal{U} = \{X_i \to X\}_{i \in I}$，函子
$$
\QCoh(\mathcal{O}_X) \longrightarrow DD(\mathcal{U})
$$
是全忠实且本质满的；参见《下降》命题
\ref{descent-proposition-fpqc-descent-quasi-coherent}。因此可应用《叠》引理
\ref{stacks-lemma-stack-equivalences}，从而证明 $\QCohstack$
满足叠的全部公理。
\end{proof}





\section{有限型拟凝聚层的叠}
\label{examples-stacks-section-stack-of-finitely-generated-quasi-coherent-sheaves}

\noindent
事实证明，只考虑有限型拟凝聚模时，可以得到拟凝聚层的叠。用
$$
p_{fg} : \QCohstack_{fg} \to (\Sch/S)_{fppf}
$$
表示 $\QCohstack$ 在 $(\Sch/S)_{fppf}$ 上的全子范畴，其对象是满足
$\mathcal{F}$ 为有限型拟凝聚 $\mathcal{O}_T$-模的二元组
$(T, \mathcal{F})$。

\begin{lemma}
\label{examples-stacks-lemma-stack-of-finite-type-quasi-coherent-sheaves}
函子 $p_{fg} : \QCohstack_{fg} \to (\Sch/S)_{fppf}$ 满足《叠》定义
\ref{stacks-definition-stack} 的条件 (1)、(2) 与 (3)。
\end{lemma}

\begin{proof}
我们验证《叠》引理 \ref{stacks-lemma-substack} 的假设 (1)、(2)、(3)。
由引理 \ref{examples-stacks-lemma-quasi-coherent-strongly-cartesian}，态射
$(Y, \mathcal{G}) \to (X, \mathcal{F})$ 是强笛卡尔态射，当且仅当它诱导同构
$f^*\mathcal{F} \to \mathcal{G}$。由《模层》引理
\ref{modules-lemma-pullback-finite-type}，有限型 $\mathcal{O}_X$-模的拉回
仍为有限型。因此，《叠》引理 \ref{stacks-lemma-substack} 的假设 (1)
成立，而假设 (2) 显然成立。最后，为证明假设 (3)，需要证明：若
$\mathcal{F}$ 是拟凝聚 $\mathcal{O}_X$-模，且
$\{f_i : X_i \to X\}$ 是一个 fppf 覆盖，使每个
$f_i^*\mathcal{F}$ 都为有限型，则 $\mathcal{F}$ 也是有限型。
把 $\mathcal{F}$ 限制到 $X$ 的仿射开集后，问题归结为如下代数陈述：
设 $R \to S$ 是有限表现忠实平坦环映射，$M$ 是 $R$-模。若
$M \otimes_R S$ 是有限生成 $S$-模，则 $M$ 是有限生成 $R$-模。
这一代数事实的更强形式见《交换代数》引理
\ref{algebra-lemma-descend-properties-modules}。
\end{proof}

\begin{lemma}
\label{examples-stacks-lemma-finite-type}
设 $(X, \mathcal{O}_X)$ 是带环空间。
\begin{enumerate}
\item 有限型 $\mathcal{O}_X$-模范畴的同构类构成集合。
\item 有限型拟凝聚 $\mathcal{O}_X$-模范畴的同构类构成集合。
\end{enumerate}
\end{lemma}

\begin{proof}
(2) 中的范畴是 (1) 中范畴的全子范畴，故 (2) 由 (1) 得出。
考虑任意开覆盖 $\mathcal{U} : X = \bigcup_{i \in I} U_i$，并用
$j_i : U_i \to X$ 表示包含映射。再考虑任意映射
$r : I \to \mathbf{N}$。若 $\mathcal{F}$ 是一个 $\mathcal{O}_X$-模，
其在 $U_i$ 上的限制可由 $\mathcal{F}(U_i)$ 中至多 $r(i)$ 个截面生成，
则 $\mathcal{F}$ 是下列层的商：
$$
\mathcal{H}_{\mathcal{U}, r} =
\bigoplus\nolimits_{i \in I} j_{i, !}\mathcal{O}_{U_i}^{\oplus r(i)}
$$
按定义，若 $\mathcal{F}$ 为有限型，则存在一个指标集为 $I = X$ 的开覆盖
$\mathcal{U}$，使上述条件成立。因此，只需说明 $\mathcal{U}$ 的可能选择构成
集合（这是显然的），$r : I \to \mathbf{N}$ 的可能选择构成集合（也是显然的），
并且对每个 $\mathcal{U}$ 与 $r$，$\mathcal{H}_{\mathcal{U}, r}$ 的商模的可能
选择构成集合。换言之，只需证明：给定一个 $\mathcal{O}_X$-模
$\mathcal{H}$，其商的同构类至多构成集合。最后这一点很清楚，因为商映射
$\mathcal{H} \to \mathcal{F}$ 的核可由
$\prod_{U \subset X\text{ 开}} \mathcal{H}(U)$ 的幂集的一个子集参数化。
\end{proof}

\begin{lemma}
\label{examples-stacks-lemma-stack-fg-quasi-coherent}
存在子范畴
$\QCohstack_{fg, small} \subset \QCohstack_{fg}$
，它具有如下性质：
\begin{enumerate}
\item 包含函子
$\QCohstack_{fg, small} \to \QCohstack_{fg}$ 全忠实且本质满；
\item 函子
$p_{fg, small} : \QCohstack_{fg, small} \to (\Sch/S)_{fppf}$
使 $\QCohstack_{fg, small}$ 成为 $(\Sch/S)_{fppf}$ 上的叠。
\end{enumerate}
\end{lemma}

\begin{proof}
由引理 \ref{examples-stacks-lemma-stack-of-finite-type-quasi-coherent-sheaves} 与
\ref{examples-stacks-lemma-finite-type} 可知，
$p_{fg} : \QCohstack_{fg} \to (\Sch/S)_{fppf}$ 满足《叠》定义
\ref{stacks-definition-stack} 的 (1)、(2)、(3)，并满足《叠》注
\ref{stacks-remark-stack-make-small} 的附加条件 (4)。因此，该注中的讨论
给出 $\QCohstack_{fg, small}$。
\end{proof}

\noindent
我们以后常作替换
$$
\QCohstack_{fg} \leadsto \QCohstack_{fg, small}
$$
而不再说明，并且略用记号，仍以 $\QCohstack_{fg}$ 表示替换后的范畴。

\begin{remark}
\label{examples-stacks-remark-higher-rank}
注意，对某个无限基数 $\kappa$，若只考虑局部可由至多 $\kappa$ 个截面生成的
拟凝聚层，则本节的全部讨论仍然成立；例如可取 $\kappa = \aleph_0$。
\end{remark}





\section{有限 \'etale 覆盖}
\label{examples-stacks-section-finite-etale}

\noindent
如下定义范畴 $\textit{F\'Et}$：
\begin{enumerate}
\item $\textit{F\'Et}$ 的对象是概形的有限 \'etale 态射 $Y \to X$
（按我们的约定，这指 $(\Sch/S)_{fppf}$ 中的有限 \'etale 态射）；
\item $\textit{F\'Et}$ 的态射
$(b, a) : (Y \to X) \to (Y' \to X')$ 是交换图
$$
\xymatrix{
Y \ar[d] \ar[r]_b & Y' \ar[d] \\
X \ar[r]_a & X'
}
$$
，其中该图位于概形范畴中。
\end{enumerate}
于是 $\textit{F\'Et}$ 是一个范畴，并且
$$
p : \textit{F\'Et} \to (\Sch/S)_{fppf},
\quad
(Y \to X) \mapsto X
$$
是函子。注意，$\textit{F\'Et}$ 在概形 $X$ 上的纤维范畴，正是
《基本群》第 \ref{pione-section-finite-etale} 节所研究的范畴
$\textit{F\'Et}_X$。

\begin{lemma}
\label{examples-stacks-lemma-finite-etale-stack}
函子
$$
p : \textit{F\'Et} \longrightarrow (\Sch/S)_{fppf}
$$
定义了 $(\Sch/S)_{fppf}$ 上的叠。
\end{lemma}

\begin{proof}
有限 \'etale 态射的 fppf 下降由《下降》引理 \ref{descent-lemma-affine}、
\ref{descent-lemma-descending-property-finite} 以及
\ref{descent-lemma-descending-property-etale} 得出。细节从略。
\end{proof}






\section{代数空间}
\label{examples-stacks-section-stack-of-spaces}

\noindent
如下定义范畴 $\Spacesstack$：
\begin{enumerate}
\item $\Spacesstack$ 的对象是 $S$ 上代数空间的态射 $X \to U$，其中
$U$ 可由 $(\Sch/S)_{fppf}$ 的一个对象表示；
\item 态射 $(f, g) : (X \to U) \to (Y \to V)$ 是交换图
$$
\xymatrix{
X \ar[d] \ar[r]_f & Y \ar[d] \\
U \ar[r]^g & V
}
$$
，其中各箭头都是 $S$ 上代数空间的态射。
\end{enumerate}
于是 $\Spacesstack$ 是一个范畴，并且
$$
p : \Spacesstack \to (\Sch/S)_{fppf},
\quad
(X \to U) \mapsto U
$$
是函子。注意，$\Spacesstack$ 在概形 $U$ 上的纤维范畴，正是 $U$ 上
代数空间的范畴 $\textit{Spaces}/U$（参见《代数空间上的拓扑》第
\ref{spaces-topologies-section-procedure} 节）。因此，有时把
$\Spacesstack$ 的对象视为二元组 $X/U$，它由概形 $U$ 与 $U$ 上的
代数空间 $X$ 组成。为供后文使用，我们指出：给定
$(X/U), (Y/V) \in \Ob(\Spacesstack)$
，有
\begin{equation}
\label{examples-stacks-equation-morphisms-spaces}
\Mor_{\Spacesstack}(X/U, Y/V)
=
\coprod\nolimits_{g \in \Mor_S(U, V)}
\Mor_{\textit{Spaces}/U}(X, U \times_{g, V} Y)
\end{equation}
范畴 $\Spacesstack$ 几乎是 $(\Sch/S)_{fppf}$ 上的叠，但还不完全是。
问题出在集合论上，下面将予以说明。

\begin{lemma}
\label{examples-stacks-lemma-spaces-strongly-cartesian}
$\Spacesstack$ 的态射 $(f, g) : X/U \to Y/V$ 是强笛卡尔态射，
当且仅当映射 $f$ 诱导同构 $X \to U \times_{g, V} Y$。
\end{lemma}

\begin{proof}
设 $Y/V \in \Ob(\Spacesstack)$，并设 $g : U \to V$ 是
$(\Sch/S)_{fppf}$ 的态射。注意，投影 $p : U \times_{g, V} Y \to Y$
给出态射
% P11-E0082：冻结源此处的 “gives rise a morphism” 缺少介词 “to”。
$(p, g) : U \times_{g, V} Y/U \to Y/V$，它是 $\Spacesstack$ 的态射。
我们断言 $(p, g)$ 是强笛卡尔态射。事实上，对 $\Spacesstack$ 的任意对象
$Z/W$，有
\begin{align*}
\Mor_{\Spacesstack}(Z/W, U \times_{g, V} Y/U)
& =
\coprod\nolimits_{h \in \Mor_S(W, U)}
\Mor_{\textit{Spaces}/W}(Z, W \times_{h, U} U \times_{g, V} Y) \\
& =
\coprod\nolimits_{h \in \Mor_S(W, U)}
\Mor_{\textit{Spaces}/W}(Z, W \times_{g \circ h, V} Y) \\
& =
\Mor_{\Spacesstack}(Z/W, Y/V)
\times_{\Mor_S(W, V)} \Mor_S(W, U)
\end{align*}
这里两次使用了式 (\ref{examples-stacks-equation-morphisms-spaces})。这证明 $(p, g)$
满足《范畴》定义 \ref{categories-definition-cartesian-over-C} 的条件，
故断言成立。由《范畴》引理
\ref{categories-lemma-composition-cartesian}，同构都是强笛卡尔态射，
而强笛卡尔态射的复合仍是强笛卡尔态射；这证明了引理的充分性。
反过来，给定 $Y/V$ 与 $g : U \to V$，若存在以 $Y/V$ 为目标、
提升 $g$ 的强笛卡尔态射，则它必与 $(p, g)$ 同构（参见《范畴》定义
\ref{categories-definition-cartesian-over-C} 后的讨论）。因此引理的必要性也成立。
\end{proof}

\begin{lemma}
\label{examples-stacks-lemma-pre-stack-of-spaces}
函子 $p : \Spacesstack \to (\Sch/S)_{fppf}$ 满足《叠》定义
\ref{stacks-definition-stack} 的条件 (1) 与 (2)。
\end{lemma}

\begin{proof}
% P11-E0081：冻结源此处写作 “It is follows from”；应删去 “is”。
由引理 \ref{examples-stacks-lemma-spaces-strongly-cartesian} 可知，$\Spacesstack$ 是
$(\Sch/S)_{fppf}$ 上的纤维化范畴，这证明了 (1)。设
$\{U_i \to U\}_{i \in I}$ 是 $(\Sch/S)_{fppf}$ 的覆盖，设 $X, Y$
是 $U$ 上的代数空间。最后，设
$\varphi_i : X_{U_i} \to Y_{U_i}$ 是 $\textit{Spaces}/U_i$ 的态射，
并且 $\varphi_i$ 与 $\varphi_j$ 在 $U_i \times_U U_j$ 上都限制为
代数空间的同一态射 $X_{U_i \times_U U_j} \to Y_{U_i \times_U U_j}$。
为证明 (2)，需要说明存在唯一的 $U$ 上态射 $\varphi  : X \to Y$，
其到 $U_i$ 的基变换等于 $\varphi_i$。从 $X$ 到 $Y$ 的态射等同于层的映射，
故这直接由《位点与层》引理 \ref{sites-lemma-glue-maps} 得出。
\end{proof}

\begin{remark}
\label{examples-stacks-remark-stack-spaces}
暂且忽略集合论困难\footnote{困难并不在于 $\Spacesstack$ 是真类；
按我们对 $S$ 上代数空间的定义，$S$ 上代数空间的同构类只构成集合。
真正的问题是：代数空间的任意不交并最终可能过大，因而落在我们所选集合
“部分宇宙”之外。}，$\Spacesstack$ 也满足对象的下降，因而是叠。
具体而言，需要证明：给定
\begin{enumerate}
\item 一个 fppf 覆盖 $\{U_i \to U\}_{i \in I}$；
\item 对每个 $i \in I$，一个代数空间 $X_i/U_i$；
\item 对每个 $i, j \in I$，一个代数空间的同构
$\varphi_{ij} : X_i \times_U U_j \to U_i \times_U X_j$
，它位于 $U_i \times_U U_j$ 上，并在
$U_i \times_U U_j \times_U U_k$ 上满足上链条件。
\end{enumerate}
则存在代数空间 $X/U$ 以及 $U_i$ 上的同构 $X_{U_i} \cong X_i$，
它们恢复同构 $\varphi_{ij}$。首先，由《位点与层》引理
\ref{sites-lemma-glue-sheaves}，存在 $(\Sch/U)_{fppf}$ 上的层 $X$，
它恢复 $X_i$ 与 $\varphi_{ij}$。再由《自举》引理
\ref{bootstrap-lemma-locally-algebraic-space}，$X$ 是代数空间
（这里忽略该引理的集合论条件）。下一节将使用此论证证明：若只考虑
有限型代数空间，就会得到一个叠。
\end{remark}






\section{有限型代数空间的叠}
\label{examples-stacks-section-stack-of-finite-type-spaces}


\noindent
事实证明，只考虑有限型代数空间时，可以得到代数空间的叠。用
$$
p_{ft} : \Spacesstack_{ft} \to (\Sch/S)_{fppf}
$$
表示 $\Spacesstack$ 在 $(\Sch/S)_{fppf}$ 上的全子范畴，其对象是满足
$X \to U$ 为有限型态射的二元组 $X/U$。

\begin{lemma}
\label{examples-stacks-lemma-stack-of-finite-type-spaces}
函子
$p_{ft} : \Spacesstack_{ft} \to (\Sch/S)_{fppf}$
满足《叠》定义 \ref{stacks-definition-stack} 的条件 (1)、(2) 与 (3)。
\end{lemma}

\begin{proof}
下面将把论证写得极为琐细（这可能反而让人难以看清思路）。

\medskip\noindent
由引理 \ref{examples-stacks-lemma-spaces-strongly-cartesian} 已知：若诱导态射
$f : X \to U \times_V Y$ 是同构，则 $\Spacesstack$ 的态射
$(f, g) : X/U \to Y/V$ 是强笛卡尔态射。注意，若 $Y \to V$ 为有限型，
则 $U \times_V Y \to U$ 也为有限型；参见《代数空间的态射》引理
\ref{spaces-morphisms-lemma-base-change-finite-type}。因此，若 $\Spacesstack$ 的态射
$(f, g) : X/U \to Y/V$ 是 $\Spacesstack$ 中的强笛卡尔态射，且
$Y/V$ 是 $\Spacesstack_{ft}$ 的对象，则 $X/U$ 自动也是
$\Spacesstack_{ft}$ 的对象，当然 $(f, g)$ 在 $\Spacesstack_{ft}$ 中
也为强笛卡尔态射。由此可知，$\Spacesstack_{ft}$ 是
$(\Sch/S)_{fppf}$ 上的纤维化范畴。这证明了 (1)。

\medskip\noindent
上述论证还表明，包含函子 $\Spacesstack_{ft} \to \Spacesstack$
把强笛卡尔态射变为强笛卡尔态射。换言之，
$\Spacesstack_{ft} \to \Spacesstack$ 是 $(\Sch/S)_{fppf}$ 上
纤维化范畴的一个 $1$-态射。

\medskip\noindent
设 $U \in \Ob((\Sch/S)_{fppf})$，并设 $X, Y$ 是 $U$ 上的有限型
代数空间。由《叠》引理
\ref{stacks-lemma-presheaf-mor-map-fibred-categories}，得到预层的映射
$$
\mathit{Mor}_{\Spacesstack_{ft}}(X, Y)
\longrightarrow
\mathit{Mor}_{\Spacesstack}(X, Y)
$$
由于 $\Spacesstack_{ft}$ 是 $\Spacesstack$ 的全子范畴，该映射是同构。
引理 \ref{examples-stacks-lemma-pre-stack-of-spaces} 已证明右端是层，故左端也是层。
这证明了 (2)。

\medskip\noindent
为证明《叠》定义 \ref{stacks-definition-stack} 的条件 (3)，需要证明：给定
\begin{enumerate}
\item $(\Sch/S)_{fppf}$ 的一个覆盖 $\{U_i \to U\}_{i \in I}$；
\item 对每个 $i \in I$，一个 $U_i$ 上的有限型代数空间 $X_i$；
\item 对每个 $i, j \in I$，一个代数空间的同构
$\varphi_{ij} : X_i \times_U U_j \to U_i \times_U X_j$
，它位于 $U_i \times_U U_j$ 上，并在
$U_i \times_U U_j \times_U U_k$ 上满足上链条件。
\end{enumerate}
则存在 $U$ 上的有限型代数空间 $X$，以及 $U_i$ 上的同构
$X_{U_i} \cong X_i$，它们恢复同构 $\varphi_{ij}$。这由《自举》引理
\ref{bootstrap-lemma-descend-algebraic-space} 的 (2) 得出。再由
《下降与代数空间》引理
\ref{spaces-descent-lemma-descending-property-locally-finite-presentation}，
$X \to U$ 是有限型态射，证明完毕。
\end{proof}

\begin{lemma}
\label{examples-stacks-lemma-stack-ft-spaces}
存在子范畴
$\Spacesstack_{ft, small} \subset \Spacesstack_{ft}$
，它具有如下性质：
\begin{enumerate}
\item 包含函子
$\Spacesstack_{ft, small} \to \Spacesstack_{ft}$ 全忠实且本质满；
\item 函子
$p_{ft, small} : \Spacesstack_{ft, small} \to (\Sch/S)_{fppf}$
使 $\Spacesstack_{ft, small}$ 成为 $(\Sch/S)_{fppf}$ 上的叠。
\end{enumerate}
\end{lemma}

\begin{proof}
% P11-E0083：冻结源在单个引理引用前误用复数 “Lemmas”。
由引理 \ref{examples-stacks-lemma-stack-of-finite-type-spaces} 已知，
$p_{ft} : \Spacesstack_{ft} \to (\Sch/S)_{fppf}$ 满足《叠》定义
\ref{stacks-definition-stack} 的 (1)、(2)、(3)。《叠》注
\ref{stacks-remark-stack-make-small} 的附加条件 (4) 也成立，因为每个
$S$ 上的代数空间 $X$ 都可写成 $U/R$，其中
$U, R \in \Ob((\Sch/S)_{fppf})$；参见《代数空间》引理
\ref{spaces-lemma-space-presentation}。因此，对象的同构类只构成集合。
该注中的讨论于是给出 $\Spacesstack_{ft, small}$。
\end{proof}

\noindent
我们以后常作替换
$$
\Spacesstack_{ft} \leadsto \Spacesstack_{ft, small}
$$
而不再说明，并且略用记号，仍以 $\Spacesstack_{ft}$ 表示替换后的范畴。

\begin{remark}
\label{examples-stacks-remark-higher-cardinality-spaces}
注意，若只考虑如下代数空间 $X/U$，则本节的全部讨论仍然成立：它们局部
为有限型，并且 $U$ 的每个仿射开集在 $X$ 中的原像可由可数多个仿射开集覆盖。
如有需要，还可以对一个基数 $\kappa$ 引入 $\kappa$-型态射的概念
（即对环扩张生成元的数目，以及基中给定仿射开集上方仿射开集的基数，
分别施加某种界），于是可以完全依照上述方法构造叠
$$
\Spacesstack_\kappa \longrightarrow (\Sch/S)_{fppf}
$$
（但须确保 $\Sch$ 随 $\kappa$ 而取得足够大）。
\end{remark}





\section{群胚叠的例子}
\label{examples-stacks-section-examples-stacks-in-groupoids}

\noindent
上面的例子都是并非群胚叠的叠。本章余下部分将给出群胚叠在代数几何中的例子。



\section{与层相联系的叠}
\label{examples-stacks-section-stack-associated-to-sheaf}

\noindent
设 $F : (\Sch/S)_{fppf}^{opp} \to \textit{Sets}$ 是预层。由此得到集合纤维化范畴
$$
p_F : \mathcal{S}_F \to (\Sch/S)_{fppf},
$$
参见《范畴》例 \ref{categories-example-presheaf}。它是集合叠，当且仅当
$F$ 是层；参见《叠》引理 \ref{stacks-lemma-stack-in-setoids-characterize}。



\section{有限型拟凝聚层的群胚叠}
\label{examples-stacks-section-stack-in-groupoids-of-quasi-coherent-sheaves}

\noindent
设 $p : \QCohstack_{fg} \to (\Sch/S)_{fppf}$ 是第
\ref{examples-stacks-section-stack-of-finitely-generated-quasi-coherent-sheaves} 节引入的叠
（沿用该处的略用记号）。按照《范畴》引理
\ref{categories-lemma-fibred-gives-fibred-groupoids} 的步骤，可以把它变成群胚叠
$p' : \QCohstack_{fg}' \to (\Sch/S)_{fppf}$；参见《叠》引理
\ref{stacks-lemma-stack-gives-stack-groupoids}。在此特例中，这仅意味着
$\QCohstack_{fg}'$ 与 $\QCohstack_{fg}$ 具有相同对象，但它的态射是二元组
$(f, g) : (U, \mathcal{F}) \to (U', \mathcal{F}')$，其中 $g$ 是同构
$g : f^*\mathcal{F}' \to \mathcal{F}$。


\section{有限型代数空间的群胚叠}
\label{examples-stacks-section-stack-in-groupoids-of-finite-type-spaces}

\noindent
设 $p : \Spacesstack_{ft} \to (\Sch/S)_{fppf}$ 是第
\ref{examples-stacks-section-stack-of-finite-type-spaces} 节引入的叠（沿用该处的略用记号）。
按照《范畴》引理 \ref{categories-lemma-fibred-gives-fibred-groupoids} 的步骤，
可以把它变成群胚叠 $p' : \Spacesstack_{ft}' \to (\Sch/S)_{fppf}$；
参见《叠》引理 \ref{stacks-lemma-stack-gives-stack-groupoids}。
在此特例中，这仅意味着 $\Spacesstack_{ft}'$ 与 $\Spacesstack_{ft}$ 具有
相同对象，即有限型态射 $X \to U$，其中 $X$ 是 $S$ 上的代数空间，
$U$ 是 $S$ 上的概形。但态射 $(f, g) : X/U \to Y/V$ 现在是交换图
$$
\xymatrix{
X \ar[d] \ar[r]_f & Y \ar[d] \\
U \ar[r]^g & V
}
$$
，并且该图是笛卡尔的。



\section{商叠}
\label{examples-stacks-section-quotient-stacks}

\noindent
设 $(U, R, s, t, c)$ 是 $S$ 上代数空间中的群胚。此时商叠
$$
[U/R] \longrightarrow (\Sch/S)_{fppf}
$$
按构造是群胚叠；参见《代数空间中的群胚》定义
\ref{spaces-groupoids-definition-quotient-stack}。甚至其
$\mathit{Isom}$-层也可由代数空间表示；参见《自举》引理
\ref{bootstrap-lemma-quotient-stack-isom}。这些商叠在代数叠理论中具有
基础性的重要地位。

\medskip\noindent
上述构造的一个特例，是与数据 $(B, G/B, m, X/B, a)$ 相联系的商叠
$$
[X/G] \longrightarrow (\Sch/S)_{fppf}
$$
。这里
\begin{enumerate}
\item $B$ 是 $S$ 上的代数空间；
\item $(G, m)$ 是 $B$ 上的群代数空间；
\item $X$ 是 $B$ 上的代数空间；
\item $a : G \times_B X \to X$ 是 $G$ 在 $X$ 上的 $B$-作用。
\end{enumerate}
具体而言，由《代数空间中的群胚》定义
\ref{spaces-groupoids-definition-quotient-stack}，群胚叠 $[X/G]$
就是上面给出的商叠 $[X/G \times_B X]$。有必要具体说明范畴 $[X/G]$
究竟是什么样子；这将在第 \ref{examples-stacks-section-group-quotient-stacks} 节完成。



\section{挠子的分类}
\label{examples-stacks-section-torsors}

\noindent
我们将仔细说明：研究群代数空间 $G$ 或群层 $\mathcal{G}$ 的挠子叠，
可以具有哪些不同含义。



\subsection{群层的挠子}
\label{examples-stacks-subsection-torsors-sheaf}

\noindent
设 $\mathcal{G}$ 是 $(\Sch/S)_{fppf}$ 上的群层。对
$U \in \Ob((\Sch/S)_{fppf})$，用 $\mathcal{G}|_U$ 表示
$\mathcal{G}$ 在 $(\Sch/U)_{fppf}$ 上的限制。如下定义范畴
$\mathcal{G}\textit{-Torsors}$：
\begin{enumerate}
\item $\mathcal{G}\textit{-Torsors}$ 的对象是二元组
$(U, \mathcal{F})$，其中 $U$ 是 $(\Sch/S)_{fppf}$ 的对象，而
$\mathcal{F}$ 是一个 $\mathcal{G}|_U$-挠子；参见《位点上同调》定义
\ref{sites-cohomology-definition-torsor}。
\item 态射 $(U, \mathcal{F}) \to (V, \mathcal{H})$ 由二元组
$(f, \alpha)$ 给出，其中 $f : U \to V$ 是 $S$ 上概形的态射，而
$\alpha : f^{-1}\mathcal{H} \to \mathcal{F}$ 是
$\mathcal{G}|_U$-挠子的同构。
\end{enumerate}
于是 $\mathcal{G}\textit{-Torsors}$ 是一个范畴，并且
$$
p : \mathcal{G}\textit{-Torsors} \longrightarrow (\Sch/S)_{fppf},
\quad
(U, \mathcal{F}) \longmapsto U
$$
是函子。注意，$\mathcal{G}\textit{-Torsors}$ 在 $U$ 上的纤维范畴，
就是 $\mathcal{G}|_U$-挠子的范畴，而它是群胚。

\begin{lemma}
\label{examples-stacks-lemma-torsors-sheaf-stack-in-groupoids}
在作《叠》注 \ref{stacks-remark-stack-make-small} 所述的替换后，函子
$$
p : \mathcal{G}\textit{-Torsors} \longrightarrow (\Sch/S)_{fppf}
$$
定义了 $(\Sch/S)_{fppf}$ 上的群胚叠。
\end{lemma}

\begin{proof}
证明中最困难的部分是说明对象满足下降。设
$\{U_i \to U\}_{i \in I}$ 是 $(\Sch/S)_{fppf}$ 的覆盖。假设对每个
$i$ 给定一个 $\mathcal{G}|_{U_i}$-挠子 $\mathcal{F}_i$，并且对每个
$i, j \in I$ 给定一个同构
$\varphi_{ij} :
\mathcal{F}_i|_{U_i \times_U U_j} \to \mathcal{F}_j|_{U_i \times_U U_j}$
，它是 $\mathcal{G}|_{U_i \times_U U_j}$-挠子的同构，并在
$U_i \times_U U_j \times_U U_k$ 上满足适当的上链条件。于是，由
《位点与层》第 \ref{sites-section-glueing-sheaves} 节，得到
$(\Sch/U)_{fppf}$ 上的层 $\mathcal{F}$；它在每个 $U_i$ 上的限制恢复
$\mathcal{F}_i$，并恢复下降数据。由《位点与层》引理
\ref{sites-lemma-mapping-property-glue} 中的范畴等价，作用映射
$\mathcal{G}|_{U_i} \times \mathcal{F}_i \to \mathcal{F}_i$ 可黏合为映射
$a : \mathcal{G}|_U \times \mathcal{F} \to \mathcal{F}$。现在需要证明
$a$ 是作用，并且 $\mathcal{F}$ 成为 $\mathcal{G}|_U$-挠子。这两个性质
都可局部检验，因而由各作用
$\mathcal{G}|_{U_i} \times \mathcal{F}_i \to \mathcal{F}_i$ 的相应性质得出。
这证明 $\mathcal{G}\textit{-Torsors}$ 中的对象满足下降。若干细节从略。
\end{proof}



\subsection{群层挠子的变体}
\label{examples-stacks-subsection-variant-torsor-sheaf}

\noindent
小节 \ref{examples-stacks-subsection-torsors-sheaf} 的构造可以略作推广。具体而言，设
$\mathcal{G} \to \mathcal{B}$ 是 $(\Sch/S)_{fppf}$ 上层的映射，并设
$$
m :
\mathcal{G} \times_\mathcal{B} \mathcal{G}
\longrightarrow
\mathcal{G}
$$
是 $\mathcal{G}/\mathcal{B}$ 上的群律。换言之，二元组
$(\mathcal{G}, m)$ 是拓扑斯 $\Sh((\Sch/S)_{fppf})/\mathcal{B}$ 中的群对象。
关于拓扑斯的局部化，参见《位点与层》第
\ref{sites-section-localize-topoi} 节。在此情形下，可以如下定义范畴
$\mathcal{G}/\mathcal{B}\textit{-Torsors}$
（这里通过 Yoneda 嵌入把概形视为层）：
\begin{enumerate}
\item $\mathcal{G}/\mathcal{B}\textit{-Torsors}$ 的对象是三元组
$(U, b, \mathcal{F})$，其中
\begin{enumerate}
\item $U$ 是 $(\Sch/S)_{fppf}$ 的对象；
\item $b : U \to \mathcal{B}$ 是 $\mathcal{B}$ 在 $U$ 上的截面；
\item $\mathcal{F}$ 是 $U$ 上的
$U \times_{b, \mathcal{B}}\mathcal{G}$-挠子。
\end{enumerate}
\item 态射 $(U, b, \mathcal{F}) \to (U', b', \mathcal{F}')$ 由二元组
$(f, g)$ 给出，其中 $f : U \to U'$ 是 $S$ 上概形的态射，满足
$b = b' \circ f$，而 $g : f^{-1}\mathcal{F}' \to \mathcal{F}$ 是
$U \times_{b, \mathcal{B}} \mathcal{G}$-挠子的同构。
\end{enumerate}
于是 $\mathcal{G}/\mathcal{B}\textit{-Torsors}$ 是一个范畴，并且
$$
p :
\mathcal{G}/\mathcal{B}\textit{-Torsors}
\longrightarrow
(\Sch/S)_{fppf},
\quad
(U, b, \mathcal{F}) \longmapsto U
$$
是函子。注意，$\mathcal{G}/\mathcal{B}\textit{-Torsors}$ 在 $U$ 上的
纤维范畴，是遍历 $b : U \to \mathcal{B}$ 所得各
$U \times_{b, \mathcal{B}} \mathcal{G}$-挠子范畴的不交并，因而是群胚。

\medskip\noindent
在特例 $\mathcal{B} = S$ 中，便恢复小节
\ref{examples-stacks-subsection-torsors-sheaf} 引入的范畴 $\mathcal{G}\textit{-Torsors}$。

\begin{lemma}
\label{examples-stacks-lemma-variant-torsors-sheaf-stack-in-groupoids}
在作《叠》注 \ref{stacks-remark-stack-make-small} 所述的替换后，函子
$$
p :
\mathcal{G}/\mathcal{B}\textit{-Torsors}
\longrightarrow
(\Sch/S)_{fppf}
$$
定义了 $(\Sch/S)_{fppf}$ 上的群胚叠。
\end{lemma}

\begin{proof}
本证明重复引理 \ref{examples-stacks-lemma-torsors-sheaf-stack-in-groupoids} 的证明。
由于那里记号较为简洁，建议读者先阅读该证明。证明中最困难的部分是说明
对象满足下降。设 $\{U_i \to U\}_{i \in I}$ 是
$(\Sch/S)_{fppf}$ 的覆盖。假设对每个 $i$ 给定二元组
$(b_i, \mathcal{F}_i)$，它由态射 $b_i : U_i \to \mathcal{B}$ 与
$U_i \times_{b_i, \mathcal{B}} \mathcal{G}$-挠子 $\mathcal{F}_i$ 组成；
并且对每个 $i, j \in I$，有
$b_i|_{U_i \times_U U_j} = b_j|_{U_i \times_U U_j}$，且给定同构
$\varphi_{ij} :
\mathcal{F}_i|_{U_i \times_U U_j} \to \mathcal{F}_j|_{U_i \times_U U_j}$
，它是 $(U_i \times_U U_j) \times_\mathcal{B} \mathcal{G}$-挠子的同构，
并在 $U_i \times_U U_j \times_U U_k$ 上满足适当的上链条件。由
《位点与层》第 \ref{sites-section-glueing-sheaves} 节，得到
$(\Sch/U)_{fppf}$ 上的层 $\mathcal{F}$；它在每个 $U_i$ 上的限制恢复
$\mathcal{F}_i$，并恢复下降数据。由 $\mathcal{B}$ 的层公理，各态射
$b_i$ 来自唯一态射 $b : U \to \mathcal{B}$。由《位点与层》引理
\ref{sites-lemma-mapping-property-glue} 中的范畴等价，各作用映射
$(U_i \times_{b_i, \mathcal{B}} \mathcal{G}) \times_{U_i} \mathcal{F}_i
\to \mathcal{F}_i$
黏合为映射
$(U \times_{b, \mathcal{B}} \mathcal{G}) \times \mathcal{F} \to \mathcal{F}$.
现在需要证明这是一个作用，并且 $\mathcal{F}$ 成为
$U \times_{b, \mathcal{B}} \mathcal{G}$-挠子。这两个性质都可局部检验，
因而由各 $\mathcal{F}_i$ 上作用的相应性质得出。这证明
$\mathcal{G}/\mathcal{B}\textit{-Torsors}$ 中的对象满足下降。若干细节从略。
\end{proof}



\subsection{主齐性空间}
\label{examples-stacks-subsection-principal-homogeneous-spaces}

\noindent
设 $B$ 是 $S$ 上的代数空间，$G$ 是 $B$ 上的群代数空间。如下定义范畴
$G\textit{-Principal}$：
\begin{enumerate}
\item $G\textit{-Principal}$ 的对象是三元组 $(U, b, X)$，其中
\begin{enumerate}
\item $U$ 是 $(\Sch/S)_{fppf}$ 的对象；
\item $b : U \to B$ 是 $S$ 上的态射；
\item $X$ 是 $U$ 上的主齐性 $G_U$-空间，其中
$G_U = U \times_{b, B} G$。
\end{enumerate}
参见《代数空间中的群胚》定义
\ref{spaces-groupoids-definition-principal-homogeneous-space}。
\item 态射 $(U, b, X) \to (U', b', X')$ 由二元组 $(f, g)$ 给出，
其中 $f : U \to U'$ 是 $B$ 上概形的态射，而
$g : X \to U \times_{f, U'} X'$ 是主齐性 $G_U$-空间的同构。
\end{enumerate}
于是 $G\textit{-Principal}$ 是一个范畴，并且
$$
p : G\textit{-Principal} \longrightarrow (\Sch/S)_{fppf},
\quad
(U, b, X) \longmapsto U
$$
是函子。注意，$G\textit{-Principal}$ 在 $U$ 上的纤维范畴，是遍历
$b : U \to B$ 所得各主齐性 $U \times_{b, B} G$-空间范畴的不交并，
因而是群胚。

\medskip\noindent
在特例 $S = B$ 中，对象就是二元组 $(U, X)$，其中 $U$ 是 $S$ 上的概形，
$X$ 是 $U$ 上的主齐性 $G_U$-空间；态射就是笛卡尔图
$$
\xymatrix{
X \ar[d] \ar[r]_g & X' \ar[d] \\
U \ar[r]^f & U'
}
$$
，其中 $g$ 是 $G$-等变的。

\begin{remark}
\label{examples-stacks-remark-principal-stack-in-groupoids}
我们猜想，在作《叠》注 \ref{stacks-remark-stack-make-small} 所述的替换后，函子
$$
p : G\textit{-Principal} \longrightarrow (\Sch/S)_{fppf}
$$
定义了 $(\Sch/S)_{fppf}$ 上的群胚叠。只要能证明如下命题，该猜想便成立：给定
\begin{enumerate}
\item $(\Sch/S)_{fppf}$ 的一个覆盖 $\{U_i \to U\}_{i \in I}$；
% P11-E0084：冻结源此处写作 “an group algebraic space”；冠词应为 “a”。
\item $U$ 上的一个群代数空间 $H$；
\item 对每个 $i$，一个 $U_i$ 上的主齐性 $H_{U_i}$-空间 $X_i$；
\item $H$-等变同构
$\varphi_{ij} : X_{i, U_i \times_U U_j} \to X_{j, U_i \times_U U_j}$
，它们满足上链条件。
\end{enumerate}
则存在 $U$ 上的主齐性 $H$-空间 $X$，它恢复 $(X_i, \varphi_{ij})$。
《自举》引理 \ref{bootstrap-lemma-descent-torsor} 的证明方法把这个问题
归结为集合论问题，所以忽略集合论问题的读者将“知道”结论为真。
\url{https://math.columbia.edu/~dejong/wordpress/?p=591}
给出了处理这一问题的一项建议。
\end{remark}



\subsection{主齐性空间的变体}
\label{examples-stacks-subsection-variant-principal-homogeneous-spaces}

\noindent
设 $S$ 是概形，设 $B = S$，并设 $G$ 是 $B = S$ 上的群概形。
在此情形下，可以定义全子范畴
$G\textit{-Principal-Schemes} \subset G\textit{-Principal}$
，其对象是二元组 $(U, X)$，其中 $U$ 是 $(\Sch/S)_{fppf}$ 的对象，
而 $X \to U$ 是 $U$ 上可表的主齐性 $G$-空间，即一个概形。

\medskip\noindent
一般而言，$G\textit{-Principal-Schemes}$ 并不是 $(\Sch/S)_{fppf}$ 上的
群胚叠。原因是确实存在并非概形的主齐性空间，所以对象的下降一般不成立。



\subsection{fppf 拓扑中的挠子}
\label{examples-stacks-subsection-fppf-torsors}

\noindent
设 $B$ 是 $S$ 上的代数空间，$G$ 是 $B$ 上的群代数空间。如下定义范畴
$G\textit{-Torsors}$：
\begin{enumerate}
\item $G\textit{-Torsors}$ 的对象是三元组 $(U, b, X)$，其中
\begin{enumerate}
\item $U$ 是 $(\Sch/S)_{fppf}$ 的对象；
\item $b : U \to B$ 是一个态射；
\item $X$ 是 $U$ 上的 fppf $G_U$-挠子，其中
$G_U = U \times_{b, B} G$。
\end{enumerate}
参见《代数空间中的群胚》定义
\ref{spaces-groupoids-definition-principal-homogeneous-space}。
\item 态射 $(U, b, X) \to (U', b', X')$ 由二元组 $(f, g)$ 给出，
其中 $f : U \to U'$ 是 $B$ 上概形的态射，而
$g : X \to U \times_{f, U'} X'$ 是 $G_U$-挠子的同构。
\end{enumerate}
于是 $G\textit{-Torsors}$ 是一个范畴，并且
$$
p : G\textit{-Torsors} \longrightarrow (\Sch/S)_{fppf},
\quad
% P11-E0085：冻结源此处写成 (U,a,X)，但本小节对象三元组使用 b，而未定义 a。
(U, a, X) \longmapsto U
$$
是函子。注意，$G\textit{-Torsors}$ 在 $U$ 上的纤维范畴，是遍历
$b : U \to B$ 所得各 fppf $U \times_{b, B} G$-挠子范畴的不交并，
因而是群胚。

\medskip\noindent
在特例 $S = B$ 中，对象就是二元组 $(U, X)$，其中 $U$ 是 $S$ 上的概形，
$X$ 是 $U$ 上的 fppf $G_U$-挠子；态射就是笛卡尔图
$$
\xymatrix{
X \ar[d] \ar[r]_g & X' \ar[d] \\
U \ar[r]^f & U'
}
$$
，其中 $g$ 是 $G$-等变的。

\begin{lemma}
\label{examples-stacks-lemma-torsors-stack-in-groupoids}
在作《叠》注 \ref{stacks-remark-stack-make-small} 所述的替换后，函子
$$
p : G\textit{-Torsors} \longrightarrow (\Sch/S)_{fppf}
$$
定义了 $(\Sch/S)_{fppf}$ 上的群胚叠。
\end{lemma}

\begin{proof}
证明中最困难的部分是说明对象满足下降，这正是《自举》引理
\ref{bootstrap-lemma-descent-torsor}。我们省略《叠》定义
\ref{stacks-definition-stack-in-groupoids} 的公理 (1) 与 (2) 的证明。
\end{proof}

\begin{lemma}
\label{examples-stacks-lemma-compare-torsors}
设 $B$ 是 $S$ 上的代数空间，$G$ 是 $B$ 上的群代数空间。分别用
$\mathcal{G}$、$\mathcal{B}$ 表示把代数空间 $G$、$B$ 视为
$(\Sch/S)_{fppf}$ 上的层。函子
$$
G\textit{-Torsors} \longrightarrow \mathcal{G}/\mathcal{B}\textit{-Torsors}
$$
把三元组 $(U, b, X)$ 映到三元组 $(U, b, \mathcal{X})$，其中
$\mathcal{X}$ 是视为层的 $X$；该函子是 $(\Sch/S)_{fppf}$ 上群胚叠的等价。
\end{lemma}

\begin{proof}
使用《叠》引理
\ref{stacks-lemma-characterize-essentially-surjective-when-ff} 的结论。
$S$ 上代数空间的范畴是 $(\Sch/S)_{fppf}$ 上层范畴的全子范畴，
故该函子全忠实。此外，两边的所有对象都是局部平凡挠子，所以上述引理的
条件 (2) 成立。因此该函子是等价。
\end{proof}



\subsection{fppf 拓扑中挠子的变体}
\label{examples-stacks-subsection-variant-fppf-torsors}

\noindent
设 $S$ 是概形，设 $B = S$，并设 $G$ 是 $B = S$ 上的群概形。
在此情形下，可以定义全子范畴
$G\textit{-Torsors-Schemes} \subset G\textit{-Torsors}$
，其对象是二元组 $(U, X)$，其中 $U$ 是 $(\Sch/S)_{fppf}$ 的对象，
而 $X \to U$ 是 $U$ 上可表的 fppf $G$-挠子，即一个概形。

\medskip\noindent
一般而言，$G\textit{-Torsors-Schemes}$ 并不是 $(\Sch/S)_{fppf}$ 上的
群胚叠。原因是确实存在并非概形的 fppf $G$-挠子，所以对象的下降一般不成立。






\section{群作用的商}
\label{examples-stacks-section-group-quotient-stacks}

\noindent
至此已经引入了足够的记号，可以更具体地说明第
\ref{examples-stacks-section-quotient-stacks} 节的叠 $[X/G]$ 是什么样子。

\begin{situation}
\label{examples-stacks-situation-quotient-stack}
这里
\begin{enumerate}
\item $S$ 是包含在 $\Sch_{fppf}$ 中的概形；
\item $B$ 是 $S$ 上的代数空间；
\item $(G, m)$ 是 $B$ 上的群代数空间；
\item $\pi : X \to B$ 是 $B$ 上的代数空间；
\item $a : G \times_B X \to X$ 是 $G$ 在 $X$ 上的 $B$-作用。
\end{enumerate}
\end{situation}

\noindent
在此情形下，如下构造范畴 $[[X/G]]$\footnote{双括号记号 $[[X/G]]$
用于区分这一范畴与前面引入的叠 $[X/G]$。命题
\ref{examples-stacks-proposition-equal-quotient-stacks} 将证明两者典范等价。此后，
我们将用记号 $[X/G]$ 表示其中任意一个。}：
\begin{enumerate}
\item $[[X/G]]$ 的对象是四元组 $(U, b, P, \varphi : P \to X)$，其中
\begin{enumerate}
\item $U$ 是 $(\Sch/S)_{fppf}$ 的对象；
\item $b : U \to B$ 是 $S$ 上的态射；
\item $P$ 是 $U$ 上的 fppf $G_U$-挠子，其中
$G_U = U \times_{b, B} G$；
\item $\varphi : P \to X$ 是 $G$-等变态射，并嵌入交换图
$$
\xymatrix{
P \ar[d] \ar[r]_{\varphi} & X \ar[d] \\
U \ar[r]^b & B
}
$$
\end{enumerate}
\item $[[X/G]]$ 的态射是二元组
$(f, g) : (U, b, P, \varphi) \to (U', b', P', \varphi')$
，其中 $f : U \to U'$ 是 $B$ 上概形的态射，$g : P \to P'$ 是位于
$f$ 上方的 $G$-等变态射，它诱导同构 $P \cong U \times_{f, U'} P'$，
并满足 $\varphi = \varphi' \circ g$。换言之，$(f, g)$ 嵌入交换图
$$
\xymatrix{
P \ar[d] \ar[rrrd]_\varphi \ar[r]^g & P' \ar[d] \ar[rrd]^{\varphi'} \\
U \ar[rrrd]_b \ar[r]^f & U' \ar[rrd]^{b'} & & X \ar[d] \\
& & & B
}
$$
\end{enumerate}
于是 $[[X/G]]$ 是一个范畴，并且
$$
p : [[X/G]] \longrightarrow (\Sch/S)_{fppf},
\quad
(U, b, P, \varphi) \longmapsto U
$$
是函子。注意，$[[X/G]]$ 在 $U$ 上的纤维范畴，是遍历
$b \in \Mor_S(U, B)$ 所得如下对象范畴的不交并：配有到 $X$ 的
$G$-等变态射的 fppf $U \times_{b, B} G$-挠子 $P$。因此，
$[[X/G]]$ 的各纤维范畴都是群胚。

\medskip\noindent
注意，函子
$$
[[X/G]] \longrightarrow G\textit{-Torsors},
\quad
(U, b, P, \varphi) \longmapsto (U, b, P)
$$
是 $(\Sch/S)_{fppf}$ 上范畴的一个 $1$-态射。

\begin{lemma}
\label{examples-stacks-lemma-group-quotient-stack-in-groupoids}
在作《叠》注 \ref{stacks-remark-stack-make-small} 所述的替换后，函子
$$
p : [[X/G]] \longrightarrow (\Sch/S)_{fppf}
$$
定义了 $(\Sch/S)_{fppf}$ 上的群胚叠。
\end{lemma}

\begin{proof}
证明中最困难的部分是说明对象满足下降。设
$\{U_i \to U\}_{i \in I}$ 是 $(\Sch/S)_{fppf}$ 中的覆盖。设
$\xi_i = (U_i, b_i, P_i, \varphi_i)$ 是 $[[X/G]]$ 在 $U_i$ 上的对象，
并设 $\varphi_{ij} : \text{pr}_0^*\xi_i \to \text{pr}_1^*\xi_j$
是一项下降数据。把上面的函子 $[[X/G]] \to G\textit{-Torsors}$
应用于它，特别得到群胚叠 $G\textit{-Torsors}$ 中三元组
$(U_i, b_i, P_i)$ 的下降数据。已经知道 $G\textit{-Torsors}$ 是群胚叠
（引理 \ref{examples-stacks-lemma-torsors-stack-in-groupoids}）。因此，可以假设对某个态射
$b : U \to B$ 有 $b_i = b|_{U_i}$，并且对 $U$ 上的某个 fppf
$G_U = U \times_{b, B} G$-挠子 $P$ 有 $P_i = U_i \times_U P$。
态射 $\varphi_i$ 与各限制 $U_i \times_U P$ 上的典范下降数据相容，
因而定义态射 $\varphi : P \to X$。（例如，可用《位点与层》引理
\ref{sites-lemma-mapping-property-glue}，也可用《下降与代数空间》引理
\ref{spaces-descent-lemma-fpqc-universal-effective-epimorphisms} 得到
$\varphi$。）这证明了对象的下降。我们省略《叠》定义
\ref{stacks-definition-stack-in-groupoids} 的公理 (1) 与 (2) 的证明。
\end{proof}

\begin{proposition}
\label{examples-stacks-proposition-equal-quotient-stacks}
在情形 \ref{examples-stacks-situation-quotient-stack} 中，存在典范等价
$$
[X/G] \longrightarrow [[X/G]]
$$
，它是 $(\Sch/S)_{fppf}$ 上群胚叠的等价。
\end{proposition}

\begin{proof}
我们把细节完整写出，以确保所有定义恰好按正确方式配合。回忆，$[X/G]$
是与代数空间中的群胚 $(X, G \times_B X, s, t, c)$ 相联系的商叠；
参见《代数空间中的群胚》定义
\ref{spaces-groupoids-definition-quotient-stack}。这意味着 $[X/G]$ 是与函子
$$
(\Sch/S)_{fppf} \longrightarrow \textit{Groupoids},
\quad
U \longmapsto (X(U), G(U) \times_{B(U)} X(U), s, t, c)
$$
相联系的群胚纤维化范畴 $[X/_{\!p}G]$ 的叠化；这里
$s(g, x) = x$、$t(g, x) = a(g, x)$，并且
$c((g, x), (g', x')) = (m(g, g'), x')$。按《范畴》例
\ref{categories-example-functor-groupoids} 的构造，$[X/_{\!p}G]$ 的对象是
二元组 $(U, x)$，其中 $x \in X(U)$；而 $[X/_{\!p}G]$ 的态射
$(f, g) : (U, x) \to (U', x')$ 由概形态射 $f : U \to U'$ 与元素
$g \in G(U)$ 给出，并满足 $a(g, x) = x' \circ f$。因此，可以按如下规则
定义群胚叠的 $1$-态射
$$
F_p : [X/_{\!p}G] \longrightarrow [[X/G]]
$$
。在对象上令
$$
F_p(U, x) =
(U, \pi \circ x, G \times_{B, \pi \circ x} U, a \circ (\text{id}_G \times x))
$$
这是有意义的，因为图
$$
\xymatrix{
G \times_{B, \pi \circ x} U \ar[d] \ar[r]_{\text{id}_G \times x} &
G \times_{B, \pi} X \ar[r]_-a &
X \ar[d]^\pi \\
U \ar[rr]^{\pi \circ x} & & B
}
$$
交换；若把两个纤维积分别视为 $U$、$X$ 上的平凡 $G$-挠子，则两条水平箭头
都是 $G$-等变的。在态射 $(f, g) : (U, x) \to (U', x')$ 上，令
$F_p(f, g) = (f, R_{g^{-1}})$，其中 $R_{g^{-1}}$ 表示按 $g$ 的逆元作右平移。
更精确地，态射 $F_p(f, g) : F_p(U, x) \to F_p(U', x')$ 由笛卡尔图给出：
$$
\xymatrix{
G \times_{B, \pi \circ x} U \ar[d] \ar[r]_{R_{g^{-1}}} &
G \times_{B, \pi \circ x'} U' \ar[d] \\
U \ar[r]^f & U'
}
$$
其中 $R_{g^{-1}}$ 在 $T$-值点上由下式给出：
$$
R_{g^{-1}}(g', u) = (m(g', i(g(u))), f(u))
$$
为说明这一定义可行，需要验证
$$
a \circ (\text{id}_G \times x)
=
a \circ (\text{id}_G \times x') \circ R_{g^{-1}}
$$
。这确实成立，因为右端作用于 $T$-值点 $(g', u)$ 后，给出所需等式
\begin{align*}
a((\text{id}_G \times x')(m(g', i(g(u))), f(u)))
& =
a(m(g', i(g(u))), x'(f(u))) \\
& =
a(g', a(i(g(u)), x'(f(u)))) \\
& =
a(g', x(u))
\end{align*}
；这里使用了 $a(g, x) = x' \circ f$，从而
$a(i(g), x' \circ f) = x$。

\medskip\noindent
由《叠》引理 \ref{stacks-lemma-stackify-groupoids-universal-property}
所给叠化的泛性质，上述 $1$-态射 $F_p$ 具有典范延拓
$F : [X/G] \to [[X/G]]$。先证明 $F$ 全忠实。由于源与目标都是群胚叠，
只需证明 $F$ 识别各 $\mathit{Isom}$-层。选取概形 $U$ 以及 $[X/G]$
在 $U$ 上的对象 $\xi, \xi'$。我们要证明
$$
F :
\mathit{Isom}_{[X/G]}(\xi, \xi')
\longrightarrow
\mathit{Isom}_{[[X/G]]}(F(\xi), F(\xi'))
$$
是层的同构。只需在 $U$ 上局部地验证这一点，故可假设 $\xi, \xi'$ 来自
$[X/_{\!p}G]$ 在 $U$ 上的对象 $(U, x)$、$(U, x')$；这直接由叠化的构造
得出，《代数空间中的群胚》第
\ref{spaces-groupoids-section-explicit-quotient-stacks} 节也给出了详细说明。
直接使用上面对 $[X/_{\!p}G]$ 中态射的描述，或使用《代数空间中的群胚》引理
\ref{spaces-groupoids-lemma-quotient-stack-morphisms}，都可看出此时
$$
\mathit{Isom}_{[X/G]}(\xi, \xi') =
U \times_{(x, x'), X \times_S X, (s, t)} (G \times_B X)
$$
。这个纤维积的 $T$-值点对应于二元组 $(u, g)$，其中
$u \in U(T)$、$g \in G(T)$，并满足
$a(g, x \circ u) = x' \circ u$。（注意，这蕴含
$\pi \circ x \circ u = \pi \circ x' \circ u$。）另一方面，按定义，
$\mathit{Isom}_{[[X/G]]}(F(\xi), F(\xi'))$ 的 $T$-值点对应于满足
$\pi \circ x \circ u = \pi \circ x' \circ u : T \to B$ 的态射
$u : T \to U$，以及同构
$$
R :
G \times_{B, \pi \circ x \circ u} T
\longrightarrow
G \times_{B, \pi \circ x' \circ u} T
$$
；它是平凡 $G_T$-挠子的同构，并与给定的到 $X$ 的映射相容。
由于挠子平凡，存在某个 $g \in G(T)$，使
$R = R_{g^{-1}}$（右乘）。它与映射
$a \circ (1_G, x \circ u), a \circ (1_G, x' \circ u) : G \times_B T \to X$
相容，当且仅当 $a(g, x \circ u) = x' \circ u$。因此得到所需的
$\mathit{Isom}$-层相等性。

\medskip\noindent
既已知道 $F$ 全忠实，便可应用《叠》引理
\ref{stacks-lemma-characterize-essentially-surjective-when-ff}。因此，为证明
$F$ 是等价，只需说明 $[[X/G]]$ 的对象在 fppf 局部属于 $F$ 的本质像。
这一点很清楚，因为 fppf 挠子在 fppf 局部平凡。结论得证。
\end{proof}

\begin{lemma}
\label{examples-stacks-lemma-classifying-stacks}
\begin{slogan}
群概形或群代数空间的分类叠。
\end{slogan}
设 $S$ 是概形，$B$ 是 $S$ 上的代数空间，$G$ 是 $B$ 上的群代数空间。
则群胚叠
$$
[B/G],\quad
[[B/G]],\quad
G\textit{-Torsors},\quad
\mathcal{G}/\mathcal{B}\textit{-Torsors}
$$
彼此典范等价。若 $G \to B$ 平坦且局部有限表现，则它们还与
$G\textit{-Principal}$ 等价。
\end{lemma}

\begin{proof}
等价
$G\textit{-Torsors} \to \mathcal{G}/\mathcal{B}\textit{-Torsors}$
由引理 \ref{examples-stacks-lemma-compare-torsors} 给出。等价 $[B/G] \to [[B/G]]$
由命题 \ref{examples-stacks-proposition-equal-quotient-stacks} 给出。展开第
\ref{examples-stacks-section-group-quotient-stacks} 节中 $[[B/G]]$ 的定义，可见
% P11-E0086：冻结源此处写成 [[B//G]]；本段与定义所用记号均为 [[B/G]]。
$[[B//G]] = G\textit{-Torsors}$。

\medskip\noindent
最后，假设 $G \to B$ 平坦且局部有限表现。为证明自然函子
$G\textit{-Torsors} \to G\textit{-Principal}$ 是等价，只需说明：对
$B$ 上的概形 $U$，主齐性 $G_U$-空间 $X \to U$ 在 fppf 局部平凡。
按主齐性空间的定义（《代数空间中的群胚》定义
\ref{spaces-groupoids-definition-principal-homogeneous-space}），存在 fpqc 覆盖
$\{U_i \to U\}$，使得 $U_i \times_U X \cong G \times_B U_i$ 是
$U_i$ 上代数空间的同构。这蕴含 $X \to U$ 满、平坦且局部有限表现；
参见《下降与代数空间》引理
\ref{spaces-descent-lemma-descending-property-surjective},
\ref{spaces-descent-lemma-descending-property-flat} 以及
\ref{spaces-descent-lemma-descending-property-locally-finite-presentation}.
选取概形 $W$ 与满 \'etale 态射 $W \to X$。由上可知，
$\{W \to U\}$ 是 fppf 覆盖，并且 $X_W \to W$ 有截面。因此，
$X$ 是 fppf $G_U$-挠子。
\end{proof}

\begin{remark}
\label{examples-stacks-remark-X-mod-G-group}
设 $S$ 是概形，$G$ 是抽象群，$X$ 是 $S$ 上的代数空间，并设
$G \to \text{Aut}_S(X)$ 是群同态。在此情形下，可类似于上文如下定义
$[[X/G]]$：
\begin{enumerate}
\item $[[X/G]]$ 的对象是三元组 $(U, P, \varphi : P \to X)$，其中
\begin{enumerate}
\item $U$ 是 $(\Sch/S)_{fppf}$ 的对象；
\item $P$ 是 $(\Sch/U)_{fppf}$ 上的层，配有 $G$ 的作用，使它成为
取值为 $G$ 的常值层下的挠子；
\item $\varphi : P \to X$ 是 $G$-等变的层映射。
\end{enumerate}
\item 态射
$(f, g) : (U, P, \varphi) \to (U', P', \varphi')$
由概形态射
% P11-E0087：冻结源此处写 f:T\to T'，但对象的基为 U 与 U'。
$f : T \to T'$ 与 $G$-等变同构 $g : P \to f^{-1}P'$ 给出，
并满足 $\varphi = \varphi' \circ g$。
\end{enumerate}
完全依照上文的方法，得到函子
$$
[[X/G]] \longrightarrow (\Sch/S)_{fppf}
$$
，它使 $[[X/G]]$ 成为 $(\Sch/S)_{fppf}$ 上的群胚叠。只要 $G$ 的基数
不过大，常值层 $\underline{G}$ 就可由 $(\Sch/S)_{fppf}$ 上的 $G_S$ 表示；
而这个版本的 $[[X/G]]$ 与上面引入的叠 $[[X/G_S]]$ 等价。
\end{remark}






\section{Picard 叠}
\label{examples-stacks-section-picard-stack}

\noindent
本节在完全一般的情形下引入 Picard 叠。在《Quot 空间与 Hilbert 空间》一章中，
我们将证明它在适当假设下是代数叠；参见该章第
\ref{quot-section-picard-stack} 节。

\medskip\noindent
设 $S$ 是概形，并设 $\pi : X \to B$ 是 $S$ 上代数空间的态射。
如下定义范畴 $\Picardstack_{X/B}$：
\begin{enumerate}
\item 对象是三元组 $(U, b, \mathcal{L})$，其中
\begin{enumerate}
\item $U$ 是 $(\Sch/S)_{fppf}$ 的对象；
\item $b : U \to B$ 是 $S$ 上的态射；
% P11-E0088：冻结源此处写作 “is in invertible sheaf”；介词应为冠词 “an”。
\item $\mathcal{L}$ 是基变换 $X_U = U \times_{b, B} X$ 上的可逆层。
\end{enumerate}
\item 态射 $(f, g) : (U, b, \mathcal{L}) \to (U', b', \mathcal{L}')$
由 $B$ 上的概形态射 $f : U \to U'$ 与同构
$g : f^*\mathcal{L}' \to \mathcal{L}$ 给出。
\end{enumerate}
态射
$(f, g) : (U, b, \mathcal{L}) \to (U', b', \mathcal{L}')$
与
$(f', g') : (U', b', \mathcal{L}') \to (U'', b'', \mathcal{L}'')$
的复合由 $(f' \circ f, g \circ f^*(g'))$ 给出。这样得到范畴
$\Picardstack_{X/B}$，并且
$$
p : \Picardstack_{X/B} \longrightarrow (\Sch/S)_{fppf},
\quad
(U, b, \mathcal{L}) \longmapsto U
$$
是函子。注意，$\Picardstack_{X/B}$ 在 $U$ 上的纤维范畴，是遍历
$b \in \Mor_S(U, B)$ 所得各 $X_U = U \times_{b, B} X$ 上可逆层范畴的
不交并。因此，各纤维范畴都是群胚。

\begin{lemma}
\label{examples-stacks-lemma-picard-stack}
在作《叠》注 \ref{stacks-remark-stack-make-small} 所述的替换后，函子
$$
\Picardstack_{X/B} \longrightarrow (\Sch/S)_{fppf}
$$
定义了 $(\Sch/S)_{fppf}$ 上的群胚叠。
\end{lemma}

\begin{proof}
同往常一样，最困难的部分是说明对象满足下降。设
$\{U_i \to U\}$ 是 $(\Sch/S)_{fppf}$ 的覆盖。设
$\xi_i = (U_i, b_i, \mathcal{L}_i)$ 是 $\Picardstack_{X/B}$ 在 $U_i$
上方的对象，并设
$\varphi_{ij} : \text{pr}_0^*\xi_i \to \text{pr}_1^*\xi_j$
是一项下降数据。特别地，这蕴含各态射 $b_i$ 是某个态射
$b : U \to B$ 的限制。记 $X_U = U \times_{b, B} X$ 以及
$X_i = U_i \times_{b_i, B} X =
U_i \times_U U \times_{b, B} X = U_i \times_U X_U$.
注意，$\mathcal{L}_i$ 是可逆 $\mathcal{O}_{X_i}$-模，而
$\{X_i \to X_U\}$ 也构成 fppf 覆盖。此外，下降数据
$\varphi_{ij}$ 转化为可逆层 $\mathcal{L}_i$ 相对于 fppf 覆盖
$\{X_i \to X_U\}$ 的下降数据。因此，由《下降与代数空间》命题
\ref{spaces-descent-proposition-fpqc-descent-quasi-coherent}，得到 $X_U$
上的唯一可逆层 $\mathcal{L}$，它在 $X_i$ 上恢复 $\mathcal{L}_i$ 与下降数据。
于是，三元组 $(U, b, \mathcal{L})$ 正是所求的
$\Picardstack_{X/B}$ 在 $U$ 上的对象。细节从略。
\end{proof}





\section{惯性叠的例子}
\label{examples-stacks-section-examples-inertia}

\noindent
下面给出一些惯性叠的例子。

\begin{example}
\label{examples-stacks-example-inertia-stack-of-X-mod-G}
设 $S$ 是概形，$G$ 是交换群，$X \to S$ 是 $S$ 上的概形，并设
$a : G \times X \to X$ 是 $G$ 在 $X$ 上的作用。对 $g \in G$，
用 $g : X \to X$ 表示相应的自同构。此时，$[X/G]$ 的惯性叠
（参见注 \ref{examples-stacks-remark-X-mod-G-group}）由下式给出：
$$
I_{[X/G]} = \coprod\nolimits_{g\in G} [X^g/G],
$$
其中，对 $G$ 的元素 $g$，记号 $X^g$ 表示概形
$X^g = \{x \in X \mid g(x) = x\}$。用公式表达，$X^g$ 实际上是纤维积
$$
X^g =  X \times_{(1, 1), X \times_S X, (g, 1)} X.
$$
事实上，对任意 $S$-概形 $T$，$[X/G]$ 的惯性叠的一个 $T$-点由三元组
$(P/T, \phi, \alpha)$ 构成：这里 $P \to T$ 是 fppf $G$-挠子，
$\phi : P \to X$ 是 $G$-等变态射，而 $\alpha$ 是 $P$ 在 $T$ 上的自同构，
满足 $\phi \circ \alpha = \phi$。由于 $G$ 是\emph{交换}群的层，
在 $T$ 上的 fppf 拓扑中局部地，$\alpha$ 由乘以 $G$ 的某个元素 $g$ 给出。
条件 $\phi \circ \alpha = \phi$ 意味着 $\phi$ 通过 $X^g$ 到 $X$ 的包含
分解；也就是说，$\phi$ 由某个态射 $P \to X^g$ 与该包含复合而得。
上述讨论允许定义纤维化范畴的态射
$I_{[X/G]} \to \coprod_{g\in G} [X^g/G]$；它在 $T$-点上由上面的讨论给出。
我们省略证明它是等价。
\end{example}

\begin{example}
\label{examples-stacks-example-inertia-stack-of-picard}
设 $f : X \to S$ 是概形的态射。假设对任意 $T \to S$，基变换
$f_T : X_T \to T$ 都具有如下性质：映射
$\mathcal{O}_T \to f_{T, *}\mathcal{O}_{X_T}$ 是同构。（这强于
$f$ {\it 在维数 $0$ 上上同调平坦}（此处插入未来的参考文献），并蕴含该性质。）
考虑 Picard 叠 $\Picardstack_{X/S}$；参见第
\ref{examples-stacks-section-picard-stack} 节。它的惯性叠在 $S$-概形 $T$ 上的点由二元组
$(\mathcal{L}, \alpha)$ 构成，其中 $\mathcal{L}$ 是 $X_T$ 上的线丛，
$\alpha$ 是该线丛的自同构。也就是说，由我们的条件，可以把 $\alpha$
视为 $H^0(X_T, \mathcal{O}_{X_T})^\times = H^0(T, \mathcal{O}_T^*)$
的元素。注意，$H^0(T, \mathcal{O}_T^*) = \mathbf{G}_{m, S}(T)$；
参见《群胚概形》例 \ref{groupoids-example-multiplicative-group}。
因此，$\Picardstack_{X/S}$ 的惯性叠为
$$
I_{\Picardstack_{X/S}} = \mathbf{G}_{m, S} \times_S \Picardstack_{X/S}.
$$
% P11-E0089：冻结源在显示公式后写作句首 “. as”，标点与大小写均不成立。
，这是 $(\Sch/S)_{fppf}$ 上的叠。
\end{example}






\section{有限 Hilbert 叠}
\label{examples-stacks-section-hilbert-d-stack}

\noindent
我们在比严格所需稍为一般的情形下表述这一构造。固定
$(\Sch/S)_{fppf}$ 上群胚叠的一个 $1$-态射
$$
F : \mathcal{X} \longrightarrow \mathcal{Y}
$$
。对每个整数 $d \geq 1$，考虑如下定义的范畴
$\mathcal{H}_d(\mathcal{X}/\mathcal{Y})$：
\begin{enumerate}
% P11-E0090：冻结源本项缺少 “is”，含 “objects of in”，且把有限局部自由性质误写成缺少态射的 “Z is a finite locally free”。
\item 对象是五元组 $(U, Z, y, x, \alpha)$，其中 $U, Z$ 是
$(\Sch/S)_{fppf}$ 的对象，而 $Z$ 是 $U$ 上次数 $d$ 的有限局部自由对象，并且
$y \in \Ob(\mathcal{Y}_U)$, $x \in \Ob(\mathcal{X}_Z)$
，$\alpha : y|_Z \to F(x)$ 是同构\footnote{这意味着，经由 $2$-Yoneda 引理
（《范畴》引理 \ref{categories-lemma-yoneda-2category}），该数据给出
$2$-交换图
$$
\xymatrix{
(\Sch/Z)_{fppf} \ar[r]_-x \ar[d] & \mathcal{X} \ar[d]^F \\
(\Sch/U)_{fppf} \ar[r]^-y & \mathcal{Y}
}
$$
，其中各项是 $(\Sch/S)_{fppf}$ 上的群胚叠。或者，也可把 $\alpha$
画成 $2$-态射
$$
\xymatrix{
(\Sch/Z)_{fppf}
\rrtwocell^{y \circ (Z \to U)}_{F \circ x}{\alpha} & &
\mathcal{Y}.
}
$$
}.
\item 态射 $(U, Z, y, x, \alpha) \to (U', Z', y', x', \alpha')$ 由概形态射
$f : U \to U'$、概形态射 $g : Z \to Z'$ 以及同构
$b : y \to f^*y'$、$a : x \to g^*x'$ 给出；其中该态射诱导同构
% P11-E0091：冻结源的纤维积 Z'\times_U U' 无法由已给结构映射成型；沿 f 的基变换应位于 U 上。
$Z \to Z' \times_U U'$，而这些同构诱导交换图
$$
\xymatrix{
y|_Z \ar[rr]_\alpha \ar[d]_{b|_Z} & &
F(x) \ar[d]^{F(a)} \\
f^*y'|_Z \ar[rr]^{\alpha'} & &
F(g^*x') \\
}
$$
\end{enumerate}
由定义显然存在典范遗忘函子
$$
p :
\mathcal{H}_d(\mathcal{X}/\mathcal{Y})
\longrightarrow
(\Sch/S)_{fppf}
$$
，它把五元组 $(U, Z, y, x, \alpha)$ 映到概形 $U$，并把态射
$(f, g, b, a) : (U, Z, y, x, \alpha) \to (U', Z', y', x', \alpha')$
映到态射 $f : U \to U'$。

\begin{lemma}
\label{examples-stacks-lemma-hilbert-d-stack}
配有上述函子 $p$ 的范畴 $\mathcal{H}_d(\mathcal{X}/\mathcal{Y})$
定义了 $(\Sch/S)_{fppf}$ 上的群胚叠。
\end{lemma}

\begin{proof}
同往常一样，最困难的部分是说明对象满足下降。设
$\{U_i \to U\}$ 是 $(\Sch/S)_{fppf}$ 的覆盖。设
$\xi_i = (U_i, Z_i, y_i, x_i, \alpha_i)$ 是
$\mathcal{H}_d(\mathcal{X}/\mathcal{Y})$ 在 $U_i$ 上方的对象，并设
$\varphi_{ij} : \text{pr}_0^*\xi_i \to \text{pr}_1^*\xi_j$
是一项下降数据。首先，$\varphi_{ij}$ 诱导下降数据
$(Z_i/U_i, \varphi_{ij})$；由《下降》引理 \ref{descent-lemma-affine}，
它是有效的。再由《下降》引理
\ref{descent-lemma-descending-property-finite-locally-free}，这产生
次数 $d$ 的有限局部自由概形 $Z/U$。以下把 $Z_i$ 与
$Z \times_U U_i$ 识别。其次，由于 $\mathcal{Y}$ 是群胚叠，各纤维范畴
$\mathcal{Y}_{U_i}$ 中的对象 $y_i$ 下降为 $\mathcal{Y}_U$ 中的对象 $y$。
类似地，由于 $\mathcal{X}$ 是群胚叠，各纤维范畴
$\mathcal{X}_{Z_i}$ 中的对象 $x_i$ 下降为 $\mathcal{X}_Z$ 中的对象 $x$。
最后，给定的同构
$$
\alpha_i :
(y|_Z)_{Z_i} = y_i|_{Z_i}
\longrightarrow
F(x_i) = F(x|_{Z_i})
$$
可以黏合为态射 $\alpha : y|_Z \to F(x)$，因为
% P11-E0092：冻结源此处写作 “as the \mathcal{Y} is a stack”；定冠词多余。
$\mathcal{Y}$ 是叠，因而 $\mathit{Isom}_\mathcal{Y}(y|_Z, F(x))$ 是层。
细节从略。
\end{proof}

\begin{definition}
\label{examples-stacks-definition-hilbert-d-stack}
% P11-E0093：冻结源的 “denote X the ...” 缺少介词 “by” 或需调换宾语次序。
把上面构造的 $\mathcal{H}_d(\mathcal{X}/\mathcal{Y})$ 称为
{\it $\mathcal{X}$ 在 $\mathcal{Y}$ 上次数 $d$ 的有限 Hilbert 叠}。
若 $\mathcal{Y} = S$，记
$\mathcal{H}_d(\mathcal{X}) = \mathcal{H}_d(\mathcal{X}/\mathcal{Y})$.
若 $\mathcal{X} = \mathcal{Y} = S$，则记为 $\mathcal{H}_d$。
\end{definition}

\noindent
注意，给定上述 $F : \mathcal{X} \to \mathcal{Y}$，有如下
$(\Sch/S)_{fppf}$ 上群胚叠的自然 $1$-态射：
\begin{equation}
\label{examples-stacks-equation-diagram-hilbert-d-stack}
\vcenter{
\xymatrix{
\mathcal{H}_d(\mathcal{X}) \ar[rd] &
\mathcal{H}_d(\mathcal{X}/\mathcal{Y}) \ar[d] \ar[l] \ar[r] &
\mathcal{Y} \\
& \mathcal{H}_d
}
}
\end{equation}
每条箭头都由一个“遗忘函子”给出。

\begin{lemma}
\label{examples-stacks-lemma-faithful-hilbert}
$1$-态射
$\mathcal{H}_d(\mathcal{X}/\mathcal{Y}) \to \mathcal{H}_d(\mathcal{X})$
是忠实的。
\end{lemma}

\begin{proof}
为验证
$\mathcal{H}_d(\mathcal{X}/\mathcal{Y}) \to \mathcal{H}_d(\mathcal{X})$
忠实，只需证明它在各纤维范畴上忠实。设
$\xi = (U, Z, y, x, \alpha)$ 与 $\xi' = (U, Z', y', x', \alpha')$ 是
$\mathcal{H}_d(\mathcal{X}/\mathcal{Y})$ 在概形 $U$ 上的两个对象，并设
$(g, b, a), (g', b', a') : \xi \to \xi'$ 是
$\mathcal{H}_d(\mathcal{X}/\mathcal{Y})$ 在 $U$ 上纤维范畴中的
两个态射。这两个态射在 $\mathcal{H}_d(\mathcal{X})$ 中的像相等，
当且仅当 $g = g'$ 且 $a = a'$。此时交换图
$$
\xymatrix{
y|_Z \ar[rr]_\alpha \ar[d]_{b|_Z, \ b'|_Z} & &
F(x) \ar[d]^{F(a) = F(a')} \\
y'|_Z \ar[rr]^-{\alpha'} & &
F(g^*x') = F((g')^*x') \\
}
$$
蕴含 $b|_Z = b'|_Z$。由于 $Z \to U$ 是次数 $d$ 的有限局部自由态射，
$\{Z \to U\}$ 是 fppf 覆盖，故 $b = b'$。
\end{proof}
% Generated mechanically from the frozen stacks-sheaves.tex target.
% Do not edit this generated file; edit the bound locale source instead.

% OK, start here.
%

\setcounter{chapter}{95}
\chapter{代数叠上的层}



\phantomsection
\label{stacks-sheaves-section-phantom}


\section{引言}
\label{stacks-sheaves-section-introduction}

\noindent
关于如何理解代数叠上的层，有许多不同方式。本章讨论其中一种，
它尤其适合我们为代数叠所采用的基础。每当引入一类层时，
我们都会明确说明它与文献中类似概念的精确关系。
本章的目标，是陈述那些显然成立或不难证明的结果，
而把更复杂的构造留待后文处理。

\medskip\noindent
事实上，要建立完备的可构造 \'etale 层理论，和/或充分讨论
其上同调层拟凝聚的 $\mathcal{O}$-模复形的导出范畴，
需要投入大量工作；参见 \cite{olsson_sheaves}。我们将在
《叠的上同调》第 \ref{stacks-cohomology-section-introduction} 节回到这一问题。

\medskip\noindent
在文献以及研究代数叠上层的论文中，代数叠的光滑-\'etale 位点
往往扮演重要角色。然而，它颇为棘手，因为代数叠的态射并不诱导
光滑-\'etale 拓扑斯之间的态射。因此，我们作出如下设计选择：
尽可能久地避免提及光滑-\'etale 位点。传统上使用该位点的论证，
将由使用下文所定义位点 $\mathcal{X}_{smooth}$ 中一个 {\v C}ech 覆盖的论证取代。

\medskip\noindent
本章的一些记号、约定和术语显得别扭，在较有经验的读者看来甚至可能颠倒。
这是有意为之。解释见《Quot 与 Hilbert 空间》第
\ref{quot-section-conventions} 节。




\section{约定}
\label{stacks-sheaves-section-conventions}

\noindent
本章采用的约定与代数叠一章相同；参见《代数叠》第
\ref{algebraic-section-conventions} 节。为方便起见，在此重述。

\medskip\noindent
我们在《概形上的拓扑》定义
\ref{topologies-definition-big-fppf-site} 所述的一个合适的大 fppf 位点
$\Sch_{fppf}$ 中工作。因此，除非另有明示，所有概形都是
$\Sch_{fppf}$ 的对象。更换大 fppf 位点时需要作出的改动，将在别处记录
（此处插入未来的参考文献）。

\medskip\noindent
我们总是相对于 $\Sch_{fppf}$ 中的一个基 $S$ 工作，并采用大 fppf 位点
$(\Sch/S)_{fppf}$；参见《概形上的拓扑》定义
\ref{topologies-definition-big-small-fppf}。取
$S = \Spec(\mathbf{Z})$ 即可恢复绝对情形。





\section{预层}
\label{stacks-sheaves-section-presheaves}

\noindent
本节定义 $(\Sch/S)_{fppf}$ 上群胚纤维化范畴的预层，
不过大部分讨论对任意基范畴上的范畴都适用。
本节也用来引入后文所采用的记号。

\begin{definition}
\label{stacks-sheaves-definition-presheaves}
设 $p : \mathcal{X} \to (\Sch/S)_{fppf}$ 是群胚纤维化范畴。
\begin{enumerate}
\item {\it $\mathcal{X}$ 上的预层}，是 $\mathcal{X}$ 的底范畴上的预层。
\item {\it $\mathcal{X}$ 上预层的态射}，是 $\mathcal{X}$ 的底范畴上
预层的态射。
\end{enumerate}
% P11-E0094：冻结源的 “denote PSh(X) the category” 缺少介词 “by” 或需调换宾语次序。
以 $\textit{PSh}(\mathcal{X})$ 表示 $\mathcal{X}$ 上预层的范畴。
\end{definition}

\noindent
以上定义的是集合预层。当然，也可以讨论带基点集合、Abel 群、群、
幺半群、环、固定环上的模以及固定域上的 Lie 代数等的预层。
{\it Abel 预层}，即 Abel 群预层所成的范畴，记为
$\textit{PAb}(\mathcal{X})$。

\medskip\noindent
设 $f : \mathcal{X} \to \mathcal{Y}$ 是 $(\Sch/S)_{fppf}$
上群胚纤维化范畴的一个 $1$-态射。回忆，这仅仅意味着 $f$ 是
$(\Sch/S)_{fppf}$ 上的函子。《位点与层》第
\ref{sites-section-more-functoriality-PSh} 节给出一对伴随函子
\footnote{证明引理 \ref{stacks-sheaves-lemma-functoriality-sheaves} 后，
将以 $f^{-1}$ 和 $f_*$ 表示这两个函子。}
\begin{equation}
\label{stacks-sheaves-equation-pushforward-pullback}
f^p : \textit{PSh}(\mathcal{Y}) \longrightarrow \textit{PSh}(\mathcal{X})
\quad\text{且}\quad
{}_pf : \textit{PSh}(\mathcal{X}) \longrightarrow \textit{PSh}(\mathcal{Y}).
\end{equation}
伴随关系为
$$
\Mor_{\textit{PSh}(\mathcal{X})}(f^p\mathcal{G}, \mathcal{F})
=
\Mor_{\textit{PSh}(\mathcal{Y})}(\mathcal{G}, {}_pf\mathcal{F})
$$
其中 $\mathcal{F} \in \Ob(\textit{PSh}(\mathcal{X}))$，且
$\mathcal{G} \in \Ob(\textit{PSh}(\mathcal{Y}))$。称
$f^p\mathcal{G}$ 为 $\mathcal{G}$ 的{\it 拉回}。由定义可知
$$
f^p\mathcal{G}(x) = \mathcal{G}(f(x))
$$
对任意 $x \in \Ob(\mathcal{X})$ 都成立。预层 ${}_pf\mathcal{F}$
称为 $\mathcal{F}$ 的{\it 推前}，由公式
$$
({}_pf\mathcal{F})(y) = \lim_{f(x) \to y} \mathcal{F}(x).
$$
描述。本节余下内容或许应移入位点一章；无论如何，初读时应跳过。

\begin{lemma}
\label{stacks-sheaves-lemma-1-morphisms-presheaves}
设 $f : \mathcal{X} \to \mathcal{Y}$ 与 $g : \mathcal{Y} \to \mathcal{Z}$
是 $(\Sch/S)_{fppf}$ 上群胚纤维化范畴的 $1$-态射。则
$(g \circ f)^p = f^p \circ g^p$，并且存在典范同构
${}_p(g \circ f) \to {}_pg \circ {}_pf$
，它与 $(f^p, {}_pf)$、$(g^p, {}_pg)$ 以及
$((g \circ f)^p, {}_p(g \circ f))$ 的伴随关系相容。
\end{lemma}

\begin{proof}
设 $\mathcal{H}$ 是 $\mathcal{Z}$ 上的预层。等式
$(g \circ f)^p\mathcal{H} = f^p (g^p\mathcal{H})$ 由下列等式给出：
$$
(g \circ f)^p\mathcal{H}(x) = \mathcal{H}((g \circ f)(x))
= \mathcal{H}(g(f(x))) = f^p (g^p\mathcal{H})(x).
$$
省略其与限制映射相容的验证。

\medskip\noindent
接着定义变换 ${}_p(g \circ f) \to {}_pg \circ {}_pf$。
设 $\mathcal{F}$ 是 $\mathcal{X}$ 上的预层。若 $z$ 是
$\mathcal{Z}$ 的对象，则四元组
$(x, f(x) \to y, y, g(y) \to z)$ 构成范畴 $\mathcal{J}$，
二元组 $(x, g(f(x)) \to z)$ 构成范畴 $\mathcal{I}$。
存在典范函子 $\mathcal{J} \to \mathcal{I}$，它把对象
$(x, \alpha : f(x) \to y, y, \beta : g(y) \to z)$ 映到
$(x, \beta \circ f(\alpha) : g(f(x)) \to z)$。这给出下式中的箭头：
\begin{align*}
({}_p(g \circ f)\mathcal{F})(z) & =
\lim_{g(f(x)) \to z} \mathcal{F}(x) \\
& = \lim_\mathcal{I} \mathcal{F} \\
& \to \lim_\mathcal{J} \mathcal{F} \\
& = \lim_{g(y) \to z}
\Big(\lim_{f(x) \to y} \mathcal{F}(x)\Big) \\
& =
% P11-E0096：冻结源在以 z 为参数的推前公式末端写作 (x)；应为 (z)。
({}_pg \circ {}_pf\mathcal{F})(x)
\end{align*}
这里使用了《范畴》引理 \ref{categories-lemma-functorial-limit}。
省略其与限制映射相容的验证。也可以不作上述直接构造，而把
${}_p(g \circ f) \cong {}_pg \circ {}_pf$
定义为与伴随性质相容的唯一映射；这样还有一个优点：
不必另行证明相容性。

\medskip\noindent
与 $(f^p, {}_pf)$、$(g^p, {}_pg)$ 以及
$((g \circ f)^p, {}_p(g \circ f))$ 的伴随关系相容，是指对上述预层
$\mathcal{H}$ 与 $\mathcal{F}$，下图交换：
% P11-E0097：冻结源下图右上与右下的 Hom 范畴写成 PSh(Y)，但对象位于 Z 上；
% 同时下排把本段的 H 换成了未定义的 G。典范修正应为两个 PSh(Z) 及两处 H。
$$
\xymatrix{
\Mor_{\textit{PSh}(\mathcal{X})}(f^pg^p\mathcal{H}, \mathcal{F})
\ar@{=}[r] \ar@{=}[d] &
\Mor_{\textit{PSh}(\mathcal{Y})}(g^p\mathcal{H}, {}_pf\mathcal{F})
\ar@{=}[r] &
\Mor_{\textit{PSh}(\mathcal{Y})}(\mathcal{H}, {}_pg{}_pf\mathcal{F})
\\
\Mor_{\textit{PSh}(\mathcal{X})}((g \circ f)^p\mathcal{G}, \mathcal{F})
\ar@{=}[rr] & &
\Mor_{\textit{PSh}(\mathcal{Y})}(\mathcal{G}, {}_p(g \circ f)\mathcal{F})
\ar[u]
}
$$
证明从略。
\end{proof}

\begin{lemma}
\label{stacks-sheaves-lemma-2-morphisms-presheaves}
设 $f, g : \mathcal{X} \to \mathcal{Y}$ 是 $(\Sch/S)_{fppf}$
上群胚纤维化范畴的 $1$-态射，并设 $t : f \to g$ 是
$(\Sch/S)_{fppf}$ 上群胚纤维化范畴的 $2$-态射。与 $t$ 相联系，
有如下函子的典范同构：
$$
t^p : g^p \longrightarrow f^p
\quad\text{且}\quad
{}_pt : {}_pf \longrightarrow {}_pg
$$
% P11-E0098：冻结源的 “which compatible” 缺少有限动词 “are”。
它们与 $(f^p, {}_pf)$ 和 $(g^p, {}_pg)$ 的伴随关系相容，
也与 $2$-态射的竖直复合及水平复合相容。
\end{lemma}

\begin{proof}
设 $\mathcal{G}$ 是 $\mathcal{Y}$ 上的预层。则
$t^p : g^p\mathcal{G} \to f^p\mathcal{G}$ 由映射族
$$
g^p\mathcal{G}(x) = \mathcal{G}(g(x))
\xrightarrow{\mathcal{G}(t_x)}
\mathcal{G}(f(x)) = f^p\mathcal{G}(x)
$$
给出，其中参数 $x \in \Ob(\mathcal{X})$。这是有意义的，因为
$t_x : f(x) \to g(x)$，而 $\mathcal{G}$ 是反变函子。
省略其与限制映射相容的验证。

\medskip\noindent
为定义变换 ${}_pt$，对 $y \in \Ob(\mathcal{Y})$，令
${}_y^f\mathcal{I}$ 与 ${}_y^g\mathcal{I}$ 分别为二元组
$(x, \psi : f(x) \to y)$ 与 $(x, \psi : g(x) \to y)$ 所成的范畴；
参见《位点与层》第 \ref{sites-section-more-functoriality-PSh} 节。
注意，$t$ 定义函子
${}_yt : {}_y^g\mathcal{I} \to {}_y^f\mathcal{I}$
，其规则为
$$
(x, g(x) \to y) \longmapsto (x, f(x) \xrightarrow{t_x} g(x) \to y).
$$
若 $\mathcal{F}$ 是 $\mathcal{X}$ 上的预层，则 ${}_yt$ 与
映射 $\mathcal{F} : {}_y^f\mathcal{I}^{opp} \to \textit{Sets}$ 的复合，
其中 $(x, f(x) \to y) \mapsto \mathcal{F}(x)$，等于
映射 $\mathcal{F} : {}_y^g\mathcal{I}^{opp} \to \textit{Sets}$。
因此，由《范畴》引理 \ref{categories-lemma-functorial-limit}，
对每个 $y \in \Ob(\mathcal{Y})$ 都得到典范映射
$$
({}_pf\mathcal{F})(y) = \lim_{{}_y^f\mathcal{I}} \mathcal{F}
\longrightarrow
\lim_{{}_y^g\mathcal{I}} \mathcal{F} = ({}_pg\mathcal{F})(y)
$$
省略其与限制映射相容的验证。也可以不作上述直接构造，而把
变换 ${}_pt$ 定义为与伴随对 $(f^p, {}_pf)$ 和 $(g^p, {}_pg)$
的伴随性质相容的唯一映射（见下文）。这样还有一个优点：
不必另行证明相容性。

\medskip\noindent
与 $(f^p, {}_pf)$ 和 $(g^p, {}_pg)$ 的伴随关系相容，是指对上述预层
$\mathcal{G}$ 与 $\mathcal{F}$，下图交换：
$$
\xymatrix{
\Mor_{\textit{PSh}(\mathcal{X})}(f^p\mathcal{G}, \mathcal{F})
\ar@{=}[r] \ar[d]_{- \circ t^p} &
\Mor_{\textit{PSh}(\mathcal{Y})}(\mathcal{G}, {}_pf\mathcal{F})
\ar[d]^{{}_pt \circ -} \\
\Mor_{\textit{PSh}(\mathcal{X})}(g^p\mathcal{G}, \mathcal{F})
\ar@{=}[r] &
\Mor_{\textit{PSh}(\mathcal{Y})}(\mathcal{G}, {}_pg\mathcal{F})
}
$$
证明从略。提示：逐步检验《位点与层》引理
\ref{sites-lemma-adjoints-pu} 的证明，并利用图中水平映射和竖直映射的
显式描述看出相容性。

\medskip\noindent
省略其与竖直复合及水平复合相容的验证。提示：对 $t^p$ 而言，
证明是直接的；再利用与伴随关系的相容性，即可推出变换 ${}_pt$
也具有该性质。
\end{proof}







\section{层}
\label{stacks-sheaves-section-sheaves}

\noindent
先作一个重要而平凡的观察（对那些不担心集合论问题的读者尤其如此）。

\medskip\noindent
考虑《概形上的拓扑》定义
\ref{topologies-definition-big-fppf-site} 所述的大 fppf 位点
$\Sch_{fppf}$，并以 $\Sch_\alpha$ 表示它的底范畴。
除了作为 fppf 位点的底范畴外，$\Sch_\alpha$ 也可作为大 Zariski 位点、
大 \'etale 位点、大光滑位点和大 syntomic 位点的底范畴；参见
《概形上的拓扑》注 \ref{topologies-remark-choice-sites}。
% P11-E0099：冻结源的 “can also can serve” 重复情态动词 can。
分别以 $\Sch_{Zar}$、$\Sch_\etale$、$\Sch_{smooth}$ 和
$\Sch_{syntomic}$ 表示这些位点。在此情形下，我们把 $S$ 的大 Zariski 位点
$(\Sch/S)_{Zar}$、$S$ 的大 \'etale 位点 $(\Sch/S)_\etale$、
$S$ 的大光滑位点 $(\Sch/S)_{smooth}$、$S$ 的大 syntomic 位点
$(\Sch/S)_{syntomic}$ 以及 $S$ 的大 fppf 位点 $(\Sch/S)_{fppf}$
分别定义为这些（绝对）大位点的局部化
$\Sch_{Zar}/S$、$\Sch_\etale/S$、$\Sch_{smooth}/S$、
$\Sch_{syntomic}/S$ 和 $\Sch_{fppf}/S$（参见《位点与层》第
\ref{sites-section-localize} 节）。因此，它们都有同一个底范畴，
即 $\Sch_\alpha/S$。

\medskip\noindent
由此可知，若 $p : \mathcal{X} \to (\Sch/S)_{fppf}$ 是群胚纤维化范畴，
则 $\mathcal{X}$ 继承 Zariski、\'etale、光滑、syntomic 和 fppf 拓扑；
参见《叠》定义 \ref{stacks-definition-topology-inherited}。

\begin{definition}
\label{stacks-sheaves-definition-inherited-topologies}
设 $\mathcal{X}$ 是 $(\Sch/S)_{fppf}$ 上的群胚纤维化范畴。
\begin{enumerate}
\item {\it 相联系的 Zariski 位点}，记为 $\mathcal{X}_{Zar}$，
是 $\mathcal{X}$ 上从 $(\Sch/S)_{Zar}$ 继承的位点结构。
\item {\it 相联系的 \'etale 位点}，记为 $\mathcal{X}_\etale$，
是 $\mathcal{X}$ 上从 $(\Sch/S)_\etale$ 继承的位点结构。
\item {\it 相联系的光滑位点}，记为 $\mathcal{X}_{smooth}$，
是 $\mathcal{X}$ 上从 $(\Sch/S)_{smooth}$ 继承的位点结构。
\item {\it 相联系的 syntomic 位点}，记为 $\mathcal{X}_{syntomic}$，
是 $\mathcal{X}$ 上从 $(\Sch/S)_{syntomic}$ 继承的位点结构。
\item {\it 相联系的 fppf 位点}，记为 $\mathcal{X}_{fppf}$，
是 $\mathcal{X}$ 上从 $(\Sch/S)_{fppf}$ 继承的位点结构。
\end{enumerate}
\end{definition}

\noindent
由上述讨论，这一定义是有意义的。若 $\mathcal{X}$ 是代数叠，
文献把 $\mathcal{X}_{fppf}$（或与其等价的位点）称为
$\mathcal{X}$ 的{\it 大 fppf 位点}；其他拓扑也类似。
为区分这一构造与其他构造，我们偶尔会采用这种术语。

\begin{remark}
\label{stacks-sheaves-remark-ambiguity}
只有在记号 $\mathcal{X}$ 指群胚纤维化范畴，而不指概形、代数空间等时，
才采用这种记号。这样可以避免与概形或代数空间的小 \'etale 位点
$X_\etale$ 混淆（后一情形使用正体大写字母，而不用花体字母）。
\end{remark}

\noindent
定义这些拓扑后，便可以说明 $\mathcal{X}$ 上的层是什么意思，
也就是定义相应的拓扑斯。

\begin{definition}
\label{stacks-sheaves-definition-sheaves}
设 $\mathcal{X}$ 是 $(\Sch/S)_{fppf}$ 上的群胚纤维化范畴，
并设 $\mathcal{F}$ 是 $\mathcal{X}$ 上的预层。
\begin{enumerate}
\item 若 $\mathcal{F}$ 是相联系的 Zariski 位点
$\mathcal{X}_{Zar}$ 上的层，则称 $\mathcal{F}$ 是
{\it Zariski 层}，或{\it 关于 Zariski 拓扑的层}。
\item 若 $\mathcal{F}$ 是相联系的 \'etale 位点
$\mathcal{X}_\etale$ 上的层，则称 $\mathcal{F}$ 是
{\it \'etale 层}，或{\it 关于 \'etale 拓扑的层}。
\item 若 $\mathcal{F}$ 是相联系的光滑位点
$\mathcal{X}_{smooth}$ 上的层，则称 $\mathcal{F}$ 是
{\it 光滑层}，或{\it 关于光滑拓扑的层}。
\item 若 $\mathcal{F}$ 是相联系的 syntomic 位点
$\mathcal{X}_{syntomic}$ 上的层，则称 $\mathcal{F}$ 是
{\it syntomic 层}，或{\it 关于 syntomic 拓扑的层}。
\item 若 $\mathcal{F}$ 是相联系的 fppf 位点
$\mathcal{X}_{fppf}$ 上的层，则称 $\mathcal{F}$ 是
{\it fppf 层}，或简称{\it 层}，也称{\it 关于 fppf 拓扑的层}。
\end{enumerate}
层的态射就是预层的态射。相应的层范畴分别记为
$\Sh(\mathcal{X}_{Zar})$,
$\Sh(\mathcal{X}_\etale)$,
$\Sh(\mathcal{X}_{smooth})$,
$\Sh(\mathcal{X}_{syntomic})$，以及
$\Sh(\mathcal{X}_{fppf})$.
\end{definition}

\noindent
当然，也可以讨论带基点集合、Abel 群、群、幺半群、环、固定环上的模以及
固定域上的 Lie 代数等的层。{\it Abel 层}，即 Abel 群层所成的范畴，
记为 $\textit{Ab}(\mathcal{X}_{fppf})$；其他拓扑也类似。
若 $\mathcal{X}$ 是代数叠，则在不计集合论问题时，
$\Sh(\mathcal{X}_{fppf})$ 等价于文献所谓的
{\it $\mathcal{X}$ 的大 fppf 位点上的层范畴}。其他拓扑也类似。
为区分这一构造与其他构造，我们偶尔会采用这种术语。

\medskip\noindent
由于这些拓扑按强度递增排列，有如下严格全包含：
$$
\Sh(\mathcal{X}_{fppf}) \subset
\Sh(\mathcal{X}_{syntomic}) \subset
\Sh(\mathcal{X}_{smooth}) \subset
\Sh(\mathcal{X}_\etale) \subset
\Sh(\mathcal{X}_{Zar}) \subset \textit{PSh}(\mathcal{X})
$$
有时记
$\Sh(\mathcal{X}_{fppf}) = \Sh(\mathcal{X})$
以及
$\textit{Ab}(\mathcal{X}_{fppf}) = \textit{Ab}(\mathcal{X})$
；这与我们的术语一致，因为 $\mathcal{X}$ 上的层就是
$\mathcal{X}$ 上的 fppf 层。

\medskip\noindent
在这一设置下，这些拓扑斯的函子性是直接的，并且与上述包含函子相容。

\begin{lemma}
\label{stacks-sheaves-lemma-functoriality-sheaves}
设 $f : \mathcal{X} \to \mathcal{Y}$ 是 $(\Sch/S)_{fppf}$
上群胚纤维化范畴的一个 $1$-态射，并设
$\tau \in \{Zar,\allowbreak \etale,\allowbreak smooth,\allowbreak syntomic,\allowbreak fppf\}$。
式 (\ref{stacks-sheaves-equation-pushforward-pullback}) 中的函子 ${}_pf$ 与 $f^p$
把 $\tau$-层变为 $\tau$-层，并定义拓扑斯的态射
$f : \Sh(\mathcal{X}_\tau) \to \Sh(\mathcal{Y}_\tau)$。
\end{lemma}

\begin{proof}
这直接由《叠》引理 \ref{stacks-lemma-topology-inherited-functorial} 得出。
\end{proof}

\noindent
换言之，第 \ref{stacks-sheaves-section-presheaves} 节所定义的预层推前与拉回，
也产生 $\tau$-层的{\it 推前}与{\it 拉回}。
由以上讨论可知，可以毫无歧义地写成 $f^p = f^{-1}$ 以及
${}_pf = f_*$。

\begin{definition}
\label{stacks-sheaves-definition-morphism}
设 $f : \mathcal{X} \to \mathcal{Y}$ 是 $(\Sch/S)_{fppf}$
上群胚纤维化范畴的态射。把上面构造的
$$
f = (f^{-1}, f_*) :
\Sh(\mathcal{X}_{fppf})
\longrightarrow
\Sh(\mathcal{Y}_{fppf})
$$
称为{\it 相联系的 fppf 拓扑斯态射}。相联系的 Zariski、\'etale、
光滑和 syntomic 拓扑斯也类似。
\end{definition}

\noindent
如《位点与层》第 \ref{sites-section-sheaves-algebraic-structures} 节所述，
同一公式（作用于底层集合层）定义带基点集合、Abel 群、群、幺半群、环、
固定环上的模以及固定域上的 Lie 代数等的层（关于我们的一种拓扑）的推前与拉回。








\section{推前的计算}
\label{stacks-sheaves-section-pushforward}

\noindent
设 $f : \mathcal{X} \to \mathcal{Y}$ 是 $(\Sch/S)_{fppf}$
上群胚纤维化范畴的一个 $1$-态射，$\mathcal{F}$ 是
$\mathcal{X}$ 上的预层，并设 $y \in \Ob(\mathcal{Y})$。
可以如下计算 $f_*\mathcal{F}(y)$。假设 $y$ 位于概形 $V$ 上方，
并利用 $2$-Yoneda 引理把 $y$ 看作 $1$-态射。考虑投影
$$
\text{pr} :
(\Sch/V)_{fppf} \times_{y, \mathcal{Y}} \mathcal{X}
\longrightarrow
\mathcal{X}
$$
则有典范恒等式
\begin{equation}
\label{stacks-sheaves-equation-pushforward}
f_*\mathcal{F}(y) = \Gamma\Big(
(\Sch/V)_{fppf} \times_{y, \mathcal{Y}} \mathcal{X},
\ \text{pr}^{-1}\mathcal{F}\Big)
\end{equation}
事实上，该 $2$-纤维积的对象是三元组
$(h : U \to V, x, f(x) \to h^*y)$。从记号中省去 $h$ 后，
这等价于给定 $\mathcal{X}$ 的对象 $x$ 以及
$\mathcal{Y}$ 的态射 $\alpha : f(x) \to y$。
由定义，$f_*\mathcal{F}(y) = \lim_{f(x) \to y} \mathcal{F}(x)$，
故等式成立。

\medskip\noindent
由此得到下面关于推前的“基变换”结果。该结果是平凡的，
关键在于我们使用的是“大”位点。

\begin{lemma}
\label{stacks-sheaves-lemma-base-change}
设 $S$ 是概形。设
$$
\xymatrix{
\mathcal{Y}' \times_\mathcal{Y} \mathcal{X} \ar[r]_{g'} \ar[d]_{f'} &
\mathcal{X} \ar[d]^f \\
\mathcal{Y}' \ar[r]^g & \mathcal{Y}
}
$$
是 $S$ 上群胚纤维化范畴的 $2$-笛卡尔图。则有典范同构
$$
g^{-1}f_*\mathcal{F} \longrightarrow f'_*(g')^{-1}\mathcal{F}
$$
，它关于 $\mathcal{X}$ 上的预层 $\mathcal{F}$ 是函子性的。
\end{lemma}

\begin{proof}
给定 $\mathcal{Y}'$ 位于 $V$ 上方的对象 $y'$，有等价
$$
(\Sch/V)_{fppf} \times_{g(y'), \mathcal{Y}} \mathcal{X}
=
(\Sch/V)_{fppf} \times_{y', \mathcal{Y}'}
(\mathcal{Y}' \times_\mathcal{Y} \mathcal{X})
$$
因此，由式 (\ref{stacks-sheaves-equation-pushforward}) 得到双射
$g^{-1}f_*\mathcal{F}(y') \to f'_*(g')^{-1}\mathcal{F}(y')$。
省略其与限制映射相容的验证。
\end{proof}

\noindent
对于群胚纤维化范畴的可表态射，公式
(\ref{stacks-sheaves-equation-pushforward}) 可以简化。建议读者跳过本节余下内容。

\begin{lemma}
\label{stacks-sheaves-lemma-representable}
设 $f : \mathcal{X} \to \mathcal{Y}$ 是 $(\Sch/S)_{fppf}$
上群胚纤维化范畴的一个 $1$-态射。下列条件等价：
\begin{enumerate}
\item $f$ 是可表的；
\item 对每个 $y \in \Ob(\mathcal{Y})$，函子
$\mathcal{X}^{opp} \to \textit{Sets}$,
$x \mapsto \Mor_\mathcal{Y}(f(x), y)$
是可表的。
\end{enumerate}
\end{lemma}

\begin{proof}
由《代数叠》第 \ref{algebraic-section-representable-morphism} 节的讨论可知，
$f$ 可表，当且仅当对每个 $y \in \Ob(\mathcal{Y})$，若它位于 $U$ 上方，
则 $2$-纤维积
$(\Sch/U)_{fppf} \times_{y, \mathcal{Y}} \mathcal{X}$
可表，即对 $U$ 上某个概形 $V_y$ 具有形式 $(\Sch/V_y)_{fppf}$。
% P11-E0100：冻结源的 “Objects in this 2-fibre products” 单复数不一致。
该 $2$-纤维积的对象是三元组
$(h : V \to U, x, \alpha : f(x) \to h^*y)$，其中 $\alpha$ 位于
$\text{id}_V$ 上方。从记号中省去 $h$ 后，这等价于给定
$\mathcal{X}$ 的对象 $x$ 以及态射 $f(x) \to y$。
因此，该 $2$-纤维积由 $V_y$ 和 $f(x_y) \to y$ 表示
（其中 $x_y$ 是 $\mathcal{X}$ 位于 $V_y$ 上方的对象），
当且仅当 (2) 中的函子由 $x_y$ 表示，并以映射
$f(x_y) \to y$ 为泛对象。
\end{proof}

\noindent
设
$$
\xymatrix{
\mathcal{X} \ar[rr]_f \ar[rd]_p & &  \mathcal{Y} \ar[ld]^q \\
& (\Sch/S)_{fppf}
}
$$
是群胚纤维化范畴的 $1$-态射。假设 $f$ 可表。对每个
$y \in \Ob(\mathcal{Y})$，选择对象 $u(y) \in \Ob(\mathcal{X})$
来表示引理 \ref{stacks-sheaves-lemma-representable} 中的函子
$x \mapsto \Mor_\mathcal{Y}(f(x), y)$
（由选择公理，这是可能的）。
根据构造，这些对象带有典范态射 $f(u(y)) \to y$。
对 $\mathcal{Y}$ 的每个态射 $\beta : y' \to y$，得到
$\mathcal{X}$ 的唯一态射 $u(\beta) : u(y') \to u(y)$，使得下图
$$
\xymatrix{
f(u(y')) \ar[d] \ar[rr]_{f(u(\beta))} & & f(u(y)) \ar[d] \\
y' \ar[rr] & & y
}
$$
交换。换言之，$u : \mathcal{Y} \to \mathcal{X}$ 是函子。
事实上，还能得到更多信息。设
$V' = q(y')$、$V = q(y)$、$U' = p(u(y'))$ 且 $U = p(u(y))$。则
$$
\xymatrix{
U' \ar[rr]_{p(u(\beta))} \ar[d] & & U \ar[d] \\
V' \ar[rr]^{q(\beta)} & & V
}
$$
是纤维积方块。这是因为 $U' \to U$ 表示
$(\Sch/V')_{fppf} \times_{y', \mathcal{Y}} \mathcal{X} \to
(\Sch/V)_{fppf} \times_{y, \mathcal{Y}} \mathcal{X}$
沿 $V' \to V$ 的基变换。

\begin{lemma}
\label{stacks-sheaves-lemma-representable-pushforward}
设 $f : \mathcal{X} \to \mathcal{Y}$ 是 $(\Sch/S)_{fppf}$
上群胚纤维化范畴的可表 $1$-态射，并设
$\tau \in \{Zar,\allowbreak \etale,\allowbreak smooth,\allowbreak syntomic,\allowbreak fppf\}$。
则函子 $u : \mathcal{Y}_\tau \to \mathcal{X}_\tau$ 连续，
并定义位点的态射 $\mathcal{X}_\tau \to \mathcal{Y}_\tau$。
它诱导的拓扑斯态射
$\Sh(\mathcal{X}_\tau) \to \Sh(\mathcal{Y}_\tau)$
与引理 \ref{stacks-sheaves-lemma-functoriality-sheaves} 中构造的态射 $f$ 相同。
此外，对 $\mathcal{X}$ 上任意预层 $\mathcal{F}$，都有
$f_*\mathcal{F}(y) = \mathcal{F}(u(y))$。
\end{lemma}

\begin{proof}
设 $\{y_i \to y\}$ 是 $\mathcal{Y}$ 中的 $\tau$-覆盖。
按定义，这仅仅意味着 $\{q(y_i) \to q(y)\}$ 是概形的 $\tau$-覆盖。
由引理前的最后一个观察可知，$\{p(u(y_i)) \to p(u(y))\}$ 是
$\tau$-覆盖 $\{q(y_i) \to q(y)\}$ 沿
$p(u(y)) \to q(y)$ 的基变换；因此，由位点公理，它自身也是
$\tau$-覆盖。故 $\{u(y_i) \to u(y)\}$ 是 $\mathcal{X}$
中的 $\tau$-覆盖。这证明了 $u$ 连续。

\medskip\noindent
采用《位点与层》第 \ref{sites-section-functoriality-PSh} 节和
\ref{sites-section-continuous-functors} 节中的记号
$u_p, u_s, u^p, u^s$。若能证明引理的最后一个断言，
则由上面已证的 $u$ 的连续性，有 $f_* = u^p = u^s$；
再由伴随关系，有 $f^{-1} = u_s$。这将证明 $u_s$ 正合，
从而 $u$ 确定一个位点态射，并且相等性也随即清楚。
为说明 $f_*\mathcal{F}(y) = \mathcal{F}(u(y))$，注意按定义
$$
f_*\mathcal{F}(y) = ({}_pf\mathcal{F})(y) =
\lim_{f(x) \to y} \mathcal{F}(x).
$$
由于 $u(y)$ 是取此极限所遍历范畴的终对象，结论成立。
\end{proof}





\section{结构层}
\label{stacks-sheaves-section-structure-sheaf}

\noindent
设 $\tau \in \{Zar,\allowbreak \etale,\allowbreak smooth,\allowbreak syntomic,\allowbreak fppf\}$，并设
$p : \mathcal{X} \to (\Sch/S)_{fppf}$ 是群胚纤维化范畴。
$(\Sch/S)_{fppf}$ 上群胚纤维化范畴的 2-范畴有终对象，即
$\text{id} : (\Sch/S)_{fppf} \to (\Sch/S)_{fppf}$
；而 $p$ 是从 $\mathcal{X}$ 到该终对象的 $1$-态射。
因此，$(\Sch/S)_{fppf}$ 上任意预层 $\mathcal{G}$ 给出
$\mathcal{X}$ 上的预层 $p^{-1}\mathcal{G}$，其规则为
$p^{-1}\mathcal{G}(x) = \mathcal{G}(p(x))$。此外，第
\ref{stacks-sheaves-section-sheaves} 节的讨论表明，只要 $\mathcal{G}$ 是
$\tau$-层，$p^{-1}\mathcal{G}$ 也是 $\tau$-层。

\medskip\noindent
回忆，位点 $(\Sch/S)_{fppf}$ 是带环位点，其结构层
$\mathcal{O}$ 由下列规则定义：
$$
(\Sch/S)^{opp} \longrightarrow \textit{Rings},
\quad
U/S \longmapsto \Gamma(U, \mathcal{O}_U)
$$
参见《下降》定义 \ref{descent-definition-structure-sheaf}。

\begin{definition}
\label{stacks-sheaves-definition-structure-sheaf}
设 $p : \mathcal{X} \to (\Sch/S)_{fppf}$ 是群胚纤维化范畴。
{\it $\mathcal{X}$ 的结构层}是环层
$\mathcal{O}_\mathcal{X} = p^{-1}\mathcal{O}$。
\end{definition}

\noindent
对 $\mathcal{X}$ 位于 $U$ 上方的对象 $x$，有
$\mathcal{O}_\mathcal{X}(x) = \mathcal{O}(U) = \Gamma(U, \mathcal{O}_U)$。
不言而喻，$\mathcal{O}_\mathcal{X}$ 也是 Zariski、\'etale、
光滑和 syntomic 层；因此，位点
$\mathcal{X}_{Zar}$、$\mathcal{X}_\etale$、$\mathcal{X}_{smooth}$、
$\mathcal{X}_{syntomic}$ 与 $\mathcal{X}_{fppf}$ 都是带环位点。
该构造同样是函子性的。

\begin{lemma}
\label{stacks-sheaves-lemma-functoriality-structure-sheaf}
设 $f : \mathcal{X} \to \mathcal{Y}$ 是 $(\Sch/S)_{fppf}$
上群胚纤维化范畴的一个 $1$-态射，并设
$\tau \in \{Zar,\allowbreak \etale,\allowbreak smooth,\allowbreak syntomic,\allowbreak fppf\}$。
存在典范恒等式
$f^{-1}\mathcal{O}_\mathcal{Y} = \mathcal{O}_\mathcal{X}$，
它使 $f : \Sh(\mathcal{X}_\tau) \to \Sh(\mathcal{Y}_\tau)$
成为带环拓扑斯的态射。
\end{lemma}

\begin{proof}
% P11-E0102：冻结源的 “Denote p and q the structural functors” 缺少 by。
以 $p : \mathcal{X} \to (\Sch/S)_{fppf}$ 与
$q : \mathcal{Y} \to (\Sch/S)_{fppf}$ 表示结构函子。
于是 $p = q \circ f$；由引理 \ref{stacks-sheaves-lemma-1-morphisms-presheaves}，
$p^{-1} = f^{-1} \circ q^{-1}$。由于
$\mathcal{O}_\mathcal{X} = p^{-1}\mathcal{O}$ 以及
$\mathcal{O}_\mathcal{Y} = q^{-1}\mathcal{O}$
，结论成立。
\end{proof}

\begin{remark}
\label{stacks-sheaves-remark-flat}
在引理 \ref{stacks-sheaves-lemma-functoriality-structure-sheaf} 的情形下，
带环拓扑斯的态射
$f : \Sh(\mathcal{X}_\tau) \to \Sh(\mathcal{Y}_\tau)$
是平坦的；从等式
% P11-E0101：冻结源把结构层恒等式的 X/Y 方向写反；引理给出 f^{-1}O_Y=O_X。
$f^{-1}\mathcal{O}_\mathcal{X} = \mathcal{O}_\mathcal{Y}$.
可以清楚看出这一点。这多少有些反直觉，例如，代数叠的闭浸入
作为代数叠的态射通常并不平坦。然而，对概形的闭浸入
$i : X \to Y$，完全相同的现象也会发生：此时，相联系的大
$\tau$-位点态射
$i : (\Sch/X)_\tau \to (\Sch/Y)_\tau$
也是平坦的。
\end{remark}




\section{模层}
\label{stacks-sheaves-section-modules}

\noindent
既然有结构层，就有模。

\begin{definition}
\label{stacks-sheaves-definition-modules}
设 $\mathcal{X}$ 是 $(\Sch/S)_{fppf}$ 上的群胚纤维化范畴。
\begin{enumerate}
\item {\it $\mathcal{X}$ 上的模预层}，是
$\mathcal{O}_\mathcal{X}$-模预层。模预层的范畴记为
$\textit{PMod}(\mathcal{O}_\mathcal{X})$。
\item 若模预层 $\mathcal{F}$ 是 fppf 层，则称 $\mathcal{F}$ 为
{\it $\mathcal{O}_\mathcal{X}$-模}，更确切地说，是
{\it $\mathcal{O}_\mathcal{X}$-模层}。
$\mathcal{O}_\mathcal{X}$-模的范畴记为
$\textit{Mod}(\mathcal{O}_\mathcal{X})$。
\end{enumerate}
\end{definition}

\noindent
文献把这些模（预）层称为
{\it $\mathcal{X}$ 的大 fppf 位点上的 $\mathcal{O}_\mathcal{X}$-模（预）层}。
若要将这些范畴与其他范畴区分开，我们偶尔会采用这种术语。
还会遇到在 Zariski、\'etale、光滑或 syntomic 拓扑中是层、
却未必是 fppf 层的模预层。必要时，把相应范畴记为
$\textit{Mod}(\mathcal{X}_\etale, \mathcal{O}_\mathcal{X})$；
其他拓扑也类似。

\medskip\noindent
接着讨论函子性，先讨论模预层。设
$$
\xymatrix{
\mathcal{X} \ar[rr]_f \ar[rd]_p & &  \mathcal{Y} \ar[ld]^q \\
& (\Sch/S)_{fppf}
}
$$
是群胚纤维化范畴的 $1$-态射。Abel 预层上的函子
$f^{-1}$、$f_*$ 延拓为函子
\begin{equation}
\label{stacks-sheaves-equation-functoriality-presheaves-modules}
f^{-1} :
\textit{PMod}(\mathcal{O}_\mathcal{Y})
\longrightarrow
\textit{PMod}(\mathcal{O}_\mathcal{X})
\quad\text{且}\quad
f_* :
\textit{PMod}(\mathcal{O}_\mathcal{X})
\longrightarrow
\textit{PMod}(\mathcal{O}_\mathcal{Y})
\end{equation}
对 $f^{-1}$，这直接成立，因为
$f^{-1}\mathcal{G}(x) = \mathcal{G}(f(x))$ 是
$\mathcal{O}_\mathcal{Y}(f(x)) = \mathcal{O}(q(f(x))) = \mathcal{O}(p(x)) =
\mathcal{O}_\mathcal{X}(x)$ 上的模。也可由以下事实得到这一点：
$f^{-1}\mathcal{O}_\mathcal{Y} = \mathcal{O}_\mathcal{X}$
，并且 $f^{-1}$ 与（预层上的）极限交换。
由于 $f_*$ 是右伴随，它与（预层上的）一切极限交换，特别是与积交换。
因此，可以像《位点上的模》引理
\ref{sites-modules-lemma-pushforward-module} 的证明那样，
把 $f_*$ 延拓为模预层上的函子。我们断言，式
(\ref{stacks-sheaves-equation-functoriality-presheaves-modules}) 中的函子构成一对伴随函子：
$$
\Mor_{\textit{PMod}(\mathcal{O}_\mathcal{X})}(
f^{-1}\mathcal{G}, \mathcal{F})
=
\Mor_{\textit{PMod}(\mathcal{O}_\mathcal{Y})}(
\mathcal{G}, f_*\mathcal{F}).
$$
由于 $f^{-1}\mathcal{O}_\mathcal{Y} = \mathcal{O}_\mathcal{X}$，
给 $\mathcal{X}$ 与 $\mathcal{Y}$ 赋予混沌拓扑后，
这由《位点上的模》引理
\ref{sites-modules-lemma-adjoint-push-pull-modules} 得出。

\medskip\noindent
接着讨论模的函子性，即 fppf 拓扑中的模层。
仍以 $f$ 表示诱导的带环拓扑斯态射；参见引理
\ref{stacks-sheaves-lemma-functoriality-structure-sheaf}
（目前取 fppf 拓扑）。注意，式
(\ref{stacks-sheaves-equation-functoriality-presheaves-modules}) 中的函子
$f^{-1}$ 与 $f_*$ 保持模层子范畴；参见引理
\ref{stacks-sheaves-lemma-functoriality-sheaves}。因此直接可知
\begin{equation}
\label{stacks-sheaves-equation-functoriality-sheaves-modules}
f^{-1} :
\textit{Mod}(\mathcal{O}_\mathcal{Y})
\longrightarrow
\textit{Mod}(\mathcal{O}_\mathcal{X})
\quad\text{且}\quad
f_* :
\textit{Mod}(\mathcal{O}_\mathcal{X})
\longrightarrow
\textit{Mod}(\mathcal{O}_\mathcal{Y})
\end{equation}
构成一对伴随函子：
$$
\Mor_{\textit{Mod}(\mathcal{O}_\mathcal{X})}(
f^{-1}\mathcal{G}, \mathcal{F})
=
\Mor_{\textit{Mod}(\mathcal{O}_\mathcal{Y})}(
\mathcal{G}, f_*\mathcal{F}).
$$
由伴随的唯一性，得到 $f^* = f^{-1}$；这里 $f^*$ 是《位点上的模》第
\ref{sites-modules-section-functoriality-modules} 节对上述带环拓扑斯态射
$f$ 所定义的函子。当然，也可以由下式直接看出：
$f^*(-) = f^{-1}(-) \otimes_{f^{-1}\mathcal{O}_\mathcal{Y}}
\mathcal{O}_\mathcal{X}$，并且
$f^{-1}\mathcal{O}_\mathcal{Y} = \mathcal{O}_\mathcal{X}$.

\medskip\noindent
Zariski、\'etale、光滑和 syntomic 拓扑中的模层也类似。



\section{可表范畴}
\label{stacks-sheaves-section-representable}

\noindent
在这个短节中，我们把自己的定义与所论代数叠可表时的情形作比较。

\begin{lemma}
\label{stacks-sheaves-lemma-compare-with-scheme}
设 $S$ 是概形。设 $\mathcal{X}$ 是 $(\Sch/S)$ 上的群胚纤维化范畴。
假设 $\mathcal{X}$ 可由概形 $X$ 表示。对
$\tau \in \{Zar,\linebreak[0] \etale,\linebreak[0]
smooth,\linebreak[0] syntomic,\linebreak[0] fppf\}$，
有典范等价
$$
(\mathcal{X}_\tau, \mathcal{O}_\mathcal{X}) =
((\Sch/X)_\tau, \mathcal{O}_X)
$$
的带环位点。
\end{lemma}

\begin{proof}
选取 $(\Sch/S)_{fppf}$ 上群胚纤维化范畴的一个等价
$(\Sch/X)_\tau \to \mathcal{X}$，再利用构造
$\mathcal{X} \leadsto \mathcal{X}_\tau$ 的函子性即可。
\end{proof}

\begin{lemma}
\label{stacks-sheaves-lemma-compare-with-morphism-of-schemes}
设 $S$ 是概形。设 $f : \mathcal{X} \to \mathcal{Y}$ 是 $S$ 上
群胚纤维化范畴的态射。假设 $\mathcal{X}$、$\mathcal{Y}$ 分别可由概形
$X$、$Y$ 表示。设 $f : X \to Y$ 是与 $f$ 对应的概形态射。对
$\tau \in \{Zar,\linebreak[0] \etale,\linebreak[0]
smooth,\linebreak[0] syntomic,\linebreak[0] fppf\}$，
带环拓扑斯态射
$f : (\Sh(\mathcal{X}_\tau), \mathcal{O}_\mathcal{X}) \to
(\Sh(\mathcal{Y}_\tau), \mathcal{O}_\mathcal{Y})$
在引理 \ref{stacks-sheaves-lemma-compare-with-scheme} 的同一化之下，与带环拓扑斯态射
$f : (\Sh((\Sch/X)_\tau), \mathcal{O}_X) \to 
(\Sh((\Sch/Y)_\tau), \mathcal{O}_Y)$ 一致。
\end{lemma}

\begin{proof}
将定义展开即得。
\end{proof}




\section{限制}
\label{stacks-sheaves-section-restriction}


\noindent
一个平凡但有用的观察是：群胚纤维化范畴在一个对象处的局部化，
等价于该对象所位于的概形的大位点。

\begin{lemma}
\label{stacks-sheaves-lemma-localizing}
设 $p : \mathcal{X} \to (\Sch/S)_{fppf}$ 是群胚纤维化范畴。\newline
设 $\tau \in \{Zar,\allowbreak \etale,\allowbreak smooth,\allowbreak syntomic,\allowbreak fppf\}$。
设 $x \in \Ob(\mathcal{X})$ 位于 $U = p(x)$ 之上。
函子 $p$ 诱导位点的等价
$\mathcal{X}_\tau/x \to (\Sch/U)_\tau$。
\end{lemma}

\begin{proof}
这是 Stacks，引理 \ref{stacks-lemma-localizing} 的特例。
\end{proof}

\noindent
我们用上面的引理来讨论（预）层的拉回以及它到概形上的限制。

\begin{definition}
\label{stacks-sheaves-definition-pullback}
设 $p : \mathcal{X} \to (\Sch/S)_{fppf}$ 是群胚纤维化范畴。
设 $x \in \Ob(\mathcal{X})$ 位于 $U = p(x)$ 之上。
设 $\mathcal{F}$ 是 $\mathcal{X}$ 上的预层。
\begin{enumerate}
\item $\mathcal{F}$ 的{\it 拉回 $x^{-1}\mathcal{F}$}，是限制
$\mathcal{F}|_{(\mathcal{X}/x)}$；通过引理 \ref{stacks-sheaves-lemma-localizing} 的等价
$\mathcal{X}/x \to (\Sch/U)_{fppf}$，将它视为
$(\Sch/U)_{fppf}$ 上的预层。
\item $\mathcal{F}$ {\it 到 $U_\etale$ 的限制}是
$x^{-1}\mathcal{F}|_{U_\etale}$，并滥用记号写成
$\mathcal{F}|_{U_\etale}$。
\end{enumerate}
\end{definition}

\noindent
这个记号是合理的，因为 $2$-Yoneda 引理（参见
代数叠，第 \ref{algebraic-section-2-yoneda} 节）为对象 $x$
配给一个 $1$-态射 $x : (\Sch/U)_{fppf} \to \mathcal{X}/x$，
它是 $p : \mathcal{X}/x \to (\Sch/U)_{fppf}$ 的拟逆。
因此，$x^{-1}\mathcal{F}$ 确实是 $\mathcal{F}$ 经由这个
$1$-态射的拉回。特别地，由上文可知，若 $\mathcal{F}$
是层（或 Zariski、\'etale、光滑、syntomic 层），则
$x^{-1}\mathcal{F}$ 是 $(\Sch/U)_{fppf}$ 上的层（或分别是
$(\Sch/U)_{Zar}$、$(\Sch/U)_\etale$、
$(\Sch/U)_{smooth}$、$(\Sch/U)_{syntomic}$ 上的层）。

\medskip\noindent
设 $p : \mathcal{X} \to (\Sch/S)_{fppf}$ 是群胚纤维化范畴。
设 $\varphi : x \to y$ 是 $\mathcal{X}$ 的一个态射，它位于概形态射
$a : U \to V$ 之上。回忆 $a$ 诱导小 \'etale 位点的态射
$a_{small} : U_\etale \to V_\etale$；参见
\'Etale 上同调，第 \ref{etale-cohomology-section-functoriality} 节。
设 $\mathcal{F}$ 是 $\mathcal{X}$ 上的预层。设
$\mathcal{F}|_{U_\etale}$ 和 $\mathcal{F}|_{V_\etale}$
分别是 $\mathcal{F}$ 经由 $x$ 和 $y$ 的限制。有自然的{\it 比较}映射
\begin{equation}
\label{stacks-sheaves-equation-comparison-push}
c_\varphi :
\mathcal{F}|_{V_\etale}
\longrightarrow
a_{small, *}(\mathcal{F}|_{U_\etale})
\end{equation}
% P11-E0103
，它是 $U_\etale$ 上预层的映射。具体而言，若 $V' \to V$ 是
\'etale 的，置 $U' = V' \times_V U$，并通过下图在 $V'$ 上的截面定义
$c_\varphi$：
$$
\xymatrix{
a_{small, *}(\mathcal{F}|_{U_\etale})(V') &
\mathcal{F}|_{U_\etale}(U') \ar@{=}[l] &
\mathcal{F}(x') \ar@{=}[l] \\
\mathcal{F}|_{V_\etale}(V') \ar@{=}[rr] \ar[u]^{c_\varphi} &
&
\mathcal{F}(y') \ar[u]_{\mathcal{F}(\varphi')}
}
$$
这里，$\varphi' : x' \to y'$ 是 $\mathcal{X}$ 的一个态射，
它嵌入如下交换图：
$$
\vcenter{
\xymatrix{
x' \ar[r] \ar[d]_{\varphi'} & x \ar[d]^\varphi \\
y' \ar[r] & y
}
}
\quad\text{位于下列对象之上：}\quad
\vcenter{
\xymatrix{
U' \ar[r] \ar[d] & U \ar[d]^a \\
V' \ar[r] & V
}
}
$$
$\varphi'$ 的存在唯一性由群胚纤维化范畴的公理得到。
我们省略如下验证：这样定义的 $c_\varphi$ 确实是预层的映射
（即与限制映射相容），并且关于 $\mathcal{F}$ 是函子的。
当 $\mathcal{F}$ 是 \'etale 拓扑的层时，我们得到一个{\it 比较}映射
\begin{equation}
\label{stacks-sheaves-equation-comparison}
c_\varphi : a_{small}^{-1}(\mathcal{F}|_{V_\etale})
\longrightarrow
\mathcal{F}|_{U_\etale}
\end{equation}
，也如上记作 $c_\varphi$（不区分伴随映射是惯常的记号滥用）。

\begin{lemma}
\label{stacks-sheaves-lemma-comparison}
设 $\mathcal{F}$ 是 $\mathcal{X} \to (\Sch/S)_{fppf}$ 上的 \'etale 层。
\begin{enumerate}
\item 若 $\varphi : x \to y$ 和 $\psi : y \to z$ 是 $\mathcal{X}$ 的态射，
分别位于 $a : U \to V$ 和 $b : V \to W$ 之上，则复合
$$
a_{small}^{-1}(b_{small}^{-1} (\mathcal{F}|_{W_\etale}))
\xrightarrow{a_{small}^{-1}c_\psi}
a_{small}^{-1}(\mathcal{F}|_{V_\etale})
\xrightarrow{c_\varphi}
\mathcal{F}|_{U_\etale}
$$
在同一化
$$
(b \circ a)_{small}^{-1}(\mathcal{F}|_{W_\etale}) =
a_{small}^{-1}(b_{small}^{-1} (\mathcal{F}|_{W_\etale})).
$$
之下等于 $c_{\psi \circ \varphi}$。
\item 若 $\varphi : x \to y$ 位于概形的一个 \'etale 态射
$a : U \to V$ 之上，则 (\ref{stacks-sheaves-equation-comparison}) 是同构。
\item 设 $f : \mathcal{Y} \to \mathcal{X}$ 是 $(\Sch/S)_{fppf}$ 上
群胚纤维化范畴的一个 $1$-态射，且 $y$ 是 $\mathcal{Y}$ 的对象，
位于概形 $U$ 之上，其像为 $x = f(y)$。则有典范同一化
$f^{-1}\mathcal{F}|_{U_\etale} = \mathcal{F}|_{U_\etale}$。
\item 此外，给定 $\mathcal{Y}$ 中位于 $a : U' \to U$ 之上的
$\psi : y' \to y$，比较映射
$c_\psi : a_{small}^{-1}(f^{-1}\mathcal{F}|_{U_\etale}) \to
f^{-1}\mathcal{F}|_{U'_\etale}$，在 (3) 的同一化之下，等于比较映射
$c_{f(\psi)} : a_{small}^{-1}\mathcal{F}|_{U_\etale}
\to \mathcal{F}|_{U'_\etale}$。
\end{enumerate}
\end{lemma}

\begin{proof}
这些性质的验证从略。
\end{proof}

\noindent
下面转而讨论模（预）层的限制。

\begin{lemma}
\label{stacks-sheaves-lemma-localizing-structure-sheaf}
设 $p : \mathcal{X} \to (\Sch/S)_{fppf}$ 是群胚纤维化范畴。\newline
设 $\tau \in \{Zar,\allowbreak \etale,\allowbreak smooth,\allowbreak syntomic,\allowbreak fppf\}$。
设 $x \in \Ob(\mathcal{X})$ 位于 $U = p(x)$ 之上。
引理 \ref{stacks-sheaves-lemma-localizing} 的等价延拓为带环位点的等价
$(\mathcal{X}_\tau/x, \mathcal{O}_\mathcal{X}|_x) \to
((\Sch/U)_\tau, \mathcal{O})$.
\end{lemma}

\begin{proof}
这直接来自结构层的构造。
\end{proof}

\noindent
设 $\mathcal{X}$ 是 $(\Sch/S)_{fppf}$ 上的群胚纤维化范畴。
设 $\mathcal{F}$ 是定义 \ref{stacks-sheaves-definition-modules} 所述的
$\mathcal{X}$ 上模（预）层。设 $x$ 是 $\mathcal{X}$ 的一个对象，
位于 $U$ 之上。于是引理 \ref{stacks-sheaves-lemma-localizing-structure-sheaf} 保证限制
$x^{-1}\mathcal{F}$ 是 $(\Sch/U)_{fppf}$ 上的模（预）层。
在这种情形下，我们有时写作 $x^*\mathcal{F} = x^{-1}\mathcal{F}$。
类似地，若 $\mathcal{F}$ 是 Zariski、\'etale、光滑或 syntomic 拓扑的层，
则 $x^{-1}\mathcal{F}$ 也是相应拓扑的层。此外，到 $U$ 上的限制
$\mathcal{F}|_{U_\etale} = x^{-1}\mathcal{F}|_{U_\etale}$
是 $\mathcal{O}_{U_\etale}$-模预层。若 $\mathcal{F}$ 是 \'etale 拓扑的层，
则 $\mathcal{F}|_{U_\etale}$ 是模层。此外，若
$\varphi : x \to y$ 是 $\mathcal{X}$ 的一个态射，位于
$a : U \to V$ 之上，则比较映射 (\ref{stacks-sheaves-equation-comparison})
与 $a_{small}^\sharp$ 相容（参见下降，注
\ref{descent-remark-change-topologies-ringed}），并诱导一个{\it 比较}映射
\begin{equation}
\label{stacks-sheaves-equation-comparison-modules}
c_\varphi : a_{small}^*(\mathcal{F}|_{V_\etale})
\longrightarrow
\mathcal{F}|_{U_\etale}
\end{equation}
，它是 $\mathcal{O}_{U_\etale}$-模的映射。
注意，引理 \ref{stacks-sheaves-lemma-comparison} 的性质 (1)、(2)、(3) 和 (4)
对 \'etale 模层也成立。下文将直接使用这一点而不再说明。

\begin{lemma}
\label{stacks-sheaves-lemma-enough-points}
设 $p : \mathcal{X} \to (\Sch/S)_{fppf}$ 是群胚纤维化范畴。\newline
设 $\tau \in \{Zar,\allowbreak \etale,\allowbreak smooth,\allowbreak syntomic,\allowbreak fppf\}$。
位点 $\mathcal{X}_\tau$ 有足够多的点。
\end{lemma}

\begin{proof}
由位点，引理 \ref{sites-lemma-enough-points-local}，我们必须证明：
存在 $\mathcal{X}$ 的一族对象 $x$，使得 $\mathcal{X}_\tau/x$
有足够多的点，并且层 $h_x^\#$ 覆盖层范畴的终对象。
由引理 \ref{stacks-sheaves-lemma-localizing} 以及 \'Etale 上同调，引理
\ref{etale-cohomology-lemma-points-fppf} 可知，对每个对象 $x$，
$\mathcal{X}_\tau/x$ 都有足够多的点，结论得证。
\end{proof}







\section{限制到代数空间}
\label{stacks-sheaves-section-restriction-algebraic-spaces}

\noindent
本节考虑可由代数空间表示的范畴上的层。下一个引理是
拓扑，引理 \ref{topologies-lemma-at-the-bottom-etale}
对于代数空间的类似结果。

\begin{lemma}
\label{stacks-sheaves-lemma-compare}
设 $S$ 是概形。设 $\mathcal{X} \to (\Sch/S)_{fppf}$ 是群胚纤维化范畴。
假设 $\mathcal{X}$ 可由代数空间 $F$ 表示。则存在连续且余连续的函子
$
F_\etale \to \mathcal{X}_\etale
$
，它诱导带环位点的态射
$$
\pi_F :
(\mathcal{X}_\etale, \mathcal{O}_\mathcal{X})
\longrightarrow
(F_\etale, \mathcal{O}_F)
$$
以及带环拓扑斯的态射
$$
i_F :
(\Sh(F_\etale), \mathcal{O}_F)
\longrightarrow
(\Sh(\mathcal{X}_\etale), \mathcal{O}_\mathcal{X})
$$
，使得 $\pi_F \circ i_F = \text{id}$。此外，$\pi_{F, *} = i_F^{-1}$。
\end{lemma}

\begin{proof}
选取一个等价 $j : \mathcal{S}_F \to \mathcal{X}$；参见
代数叠，第 \ref{algebraic-section-split} 节和第
\ref{algebraic-section-representable-by-algebraic-spaces} 节。
$F_\etale$ 的一个对象是概形 $U$ 连同一个 \'etale 态射
$\varphi : U \to F$。于是 $\varphi$ 是 $\mathcal{S}_F$ 在 $U$ 上的对象，
故 $j(\varphi)$ 是 $\mathcal{X}$ 在 $U$ 上的对象。这样，$j$ 诱导函子
$u : F_\etale \to \mathcal{X}$。显然，对于 $\mathcal{X}$ 上的
\'etale 拓扑，$u$ 连续且余连续。由于 $j$ 是等价，函子 $u$ 是全忠实的。
此外，$F_\etale$ 中存在纤维积和等化子，而 $u$ 与它们交换，
因为它们是在 $F_\etale$ 中底层概形的层次上计算的。因此，位点，引理
\ref{sites-lemma-when-shriek}、\ref{sites-lemma-preserve-equalizers}
和 \ref{sites-lemma-back-and-forth} 适用。特别地，$u$ 定义拓扑斯态射
$i_F : \Sh(F_\etale) \to \Sh(\mathcal{X}_\etale)$，并且 $i_F^{-1}$
有一个左伴随 $i_{F, !}$，它与纤维积和等化子交换。

\medskip\noindent
我们断言 $i_{F, !}$ 是正合的。若此断言成立，就可以按
$\pi_F^{-1} = i_{F, !}$ 和 $\pi_{F, *} = i_F^{-1}$ 定义 $\pi_F$，
其余便都清楚了。为证明该断言，注意我们已经知道 $i_{F, !}$
右正合并保持纤维积。因此只须证明 $i_{F, !}* = *$，其中 $*$
表示集合层范畴的终对象。设 $U$ 是概形，并设
$\varphi : U \to F$ 是满的 \'etale 态射。置 $R = U \times_F U$。则
$$
\xymatrix{
h_R \ar@<1ex>[r] \ar@<-1ex>[r] & h_U \ar[r] & {*}
}
$$
是 $\Sh(F_\etale)$ 中的余等化子图。利用 $i_{F, !}$ 的右正合性、
$i_{F, !} = (u_p\ )^\#$ 以及位点，引理
\ref{sites-lemma-pullback-representable-presheaf}，可见
$$
\xymatrix{
h_{u(R)} \ar@<1ex>[r] \ar@<-1ex>[r] & h_{u(U)} \ar[r] & i_{F, !}{*}
}
$$
是 $\Sh(\mathcal{X}_\etale)$ 中的余等化子图。利用 $j$ 是等价以及
$F = U/R$，可知两个映射 $h_{u(R)} \to h_{u(U)}$ 在
$\Sh(\mathcal{X}_\etale)$ 中的余等化子是 $*$。
我们省略这些态射与结构层相容的证明。
\end{proof}

\begin{remark}
\label{stacks-sheaves-remark-compare}
引理 \ref{stacks-sheaves-lemma-compare} 中的构造与 \'etale 局部化相容。
准确表述如下。设 $S$ 是概形。设 $f : \mathcal{X} \to \mathcal{Y}$
是 $(\Sch/S)_{fppf}$ 上群胚纤维化范畴的态射。假设
$\mathcal{X}$、$\mathcal{Y}$ 可分别由代数空间 $F$、$G$ 表示，
并且所诱导的代数空间态射 $f : F \to G$ 是 \'etale 的。
% P11-E0104
将对应的带环拓扑斯态射记为 $f_{small} : F_\etale \to G_\etale$。则
$$
\xymatrix{
(\Sh(F_\etale), \mathcal{O}_F) \ar[rr]_{f_{small}} \ar[d]_{i_F} & &
(\Sh(G_\etale), \mathcal{O}_G) \ar[d]^{i_G} \\
(\Sh(\mathcal{X}_\etale), \mathcal{O}_\mathcal{X})
\ar[d]_{\pi_F} \ar[rr]_f & &
(\Sh(\mathcal{Y}_\etale), \mathcal{O}_\mathcal{Y}) \ar[d]^{\pi_G} \\
(\Sh(F_\etale), \mathcal{O}_F) \ar[rr]^{f_{small}} & &
(\Sh(G_\etale), \mathcal{O}_G)
}
$$
是带环拓扑斯的交换图。细节从略。
\end{remark}

\noindent
假设 $\mathcal{X}$ 是由代数空间 $F$ 表示的代数叠。
设 $j : \mathcal{S}_F \to \mathcal{X}$ 是等价，并将上面
引理 \ref{stacks-sheaves-lemma-compare} 的证明中的函子记为
$u : F_\etale \to \mathcal{X}_\etale$。
给定 $\mathcal{X}_\etale$ 上的层 $\mathcal{F}$，有
$$
\pi_{F, *}\mathcal{F}(U) = i_F^{-1}\mathcal{F}(U) = \mathcal{F}(u(U)).
$$
正因如此，我们常把 $i_F^{-1}$ 看作{\it 限制函子}，类似于
定义 \ref{stacks-sheaves-definition-pullback}，也类似于把概形的大 \'etale 位点上的层
限制到该概形的小 \'etale 位点。在这种情形下，我们常使用记号
\begin{equation}
\label{stacks-sheaves-equation-restrict}
\mathcal{F}|_{F_\etale} = i_F^{-1}\mathcal{F} = \pi_{F, *}\mathcal{F}
\end{equation}
。

\begin{lemma}
\label{stacks-sheaves-lemma-compare-morphism}
设 $S$ 是概形。设 $f : \mathcal{X} \to \mathcal{Y}$ 是
$(\Sch/S)_{fppf}$ 上群胚纤维化范畴的态射。假设
$\mathcal{X}$、$\mathcal{Y}$ 可分别由代数空间 $F$、$G$ 表示。
将所诱导的代数空间态射记为 $f : F \to G$，并将对应的带环拓扑斯态射记为
$f_{small} : F_\etale \to G_\etale$。则
$$
\xymatrix{
(\Sh(\mathcal{X}_\etale), \mathcal{O}_\mathcal{X})
\ar[d]_{\pi_F} \ar[rr]_f & &
(\Sh(\mathcal{Y}_\etale), \mathcal{O}_\mathcal{Y}) \ar[d]^{\pi_G} \\
(\Sh(F_\etale), \mathcal{O}_F) \ar[rr]^{f_{small}} & &
(\Sh(G_\etale), \mathcal{O}_G)
}
$$
是带环拓扑斯的交换图。
\end{lemma}

\begin{proof}
这与拓扑，引理 \ref{topologies-lemma-morphism-big-small-etale} (3) 类似，
不过由于 $F \to G$ 未必可由概形表示，存在一个小障碍。
特别地，我们得不到带环位点的交换图，而只能得到带环拓扑斯的交换图。

\medskip\noindent
在正式开始证明之前，先选取等价
$j : \mathcal{S}_F \to \mathcal{X}$ 和
$j' : \mathcal{S}_G \to \mathcal{Y}$；它们如引理
\ref{stacks-sheaves-lemma-compare} 的证明中那样诱导函子
$u : F_\etale \to \mathcal{X}$ 和 $u' : G_\etale \to \mathcal{Y}$。
由于 $\Sch_{fppf}$ 上群胚纤维化范畴之层具有 2-函子性
（参见第 \ref{stacks-sheaves-section-presheaves} 节的讨论），我们可以假设
$\mathcal{X} = \mathcal{S}_F$、$\mathcal{Y} = \mathcal{S}_G$，
并假设 $f : \mathcal{S}_F \to \mathcal{S}_G$ 是与态射
$f : F \to G$ 相联系的函子。相应地，记号中将省略 $u$ 和 $u'$；
也就是说，给定 $F_\etale$ 的一个对象 $U \to F$，我们把
$\mathcal{X}$ 的对应对象记为 $U/F$。对 $G$ 也同样处理。

\medskip\noindent
设 $\mathcal{G}$ 是 $\mathcal{X}_\etale$ 上的层。
% P11-E0105
为证明 (2)，我们计算 $\pi_{G, *}f_*\mathcal{G}$ 和
$f_{small, *}\pi_{F, *}\mathcal{G}$。为此，设 $V \to G$ 是
$G_\etale$ 的一个对象。则
$$
\pi_{G, *}f_*\mathcal{G}(V) = f_*\mathcal{G}(V/G) =
\Gamma\Big(
(\Sch/V)_{fppf} \times_{\mathcal{Y}} \mathcal{X},
\ \text{pr}^{-1}\mathcal{G}\Big)
$$
；参见 (\ref{stacks-sheaves-equation-pushforward})。公式中的纤维积是
$$
(\Sch/V)_{fppf} \times_{\mathcal{Y}} \mathcal{X} =
(\Sch/V)_{fppf} \times_{\mathcal{S}_G} \mathcal{S}_F =
\mathcal{S}_{V \times_G F}
$$
，也就是说，它是与代数空间 $V \times_G F$ 相联系的分裂群胚纤维化范畴。
而 $\text{pr}^{-1}\mathcal{G}$ 是 $\mathcal{S}_{V \times_G F}$ 上
关于 \'etale 拓扑的层。

\medskip\noindent
特别地，若 $V \times_G F$ 可表，即它是概形，则
$\pi_{G, *}f_*\mathcal{G}(V) = \mathcal{G}(V \times_G F/F)$，并且
$$
f_{small, *}\pi_{F, *}\mathcal{G}(V) =
\pi_{F, *}\mathcal{G}(V \times_G F) =
\mathcal{G}(V \times_G F/F)
$$
，这就证明了该特殊情形下所需的等式。

\medskip\noindent
一般地，选取概形 $U$ 和满的 \'etale 态射
$U \to V \times_G F$。置 $R = U \times_{V \times_G F} U$。于是
$U/V \times_G F$ 和 $R/V \times_G F$ 是上述纤维积范畴的对象。
由于 $\text{pr}^{-1}\mathcal{G}$ 是 $\mathcal{S}_{V \times_G F}$ 上
关于 \'etale 拓扑的层，图
$$
\xymatrix{
\Gamma\Big(
(\Sch/V)_{fppf} \times_{\mathcal{Y}} \mathcal{X},
\ \text{pr}^{-1}\mathcal{G}\Big)
\ar[r] &
\text{pr}^{-1}\mathcal{G}(U/V \times_G F) \ar@<1ex>[r] \ar@<-1ex>[r] &
\text{pr}^{-1}\mathcal{G}(R/V \times_G F)
}
$$
是等化子图。注意，由拉回的定义，
$\text{pr}^{-1}\mathcal{G}(U/V \times_G F) = \mathcal{G}(U/F)$ 且
$\text{pr}^{-1}\mathcal{G}(R/V \times_G F) = \mathcal{G}(R/F)$。
此外，由空间的性质，第 \ref{spaces-properties-section-etale-site} 节
中的材料（特别是空间的性质，注
\ref{spaces-properties-remark-explain-equivalence} 和引理
\ref{spaces-properties-lemma-functoriality-etale-site}），可见有等化子图
$$
\xymatrix{
f_{small, *}\pi_{F, *}\mathcal{G}(V)
\ar[r] &
\pi_{F, *}\mathcal{G}(U/F) \ar@<1ex>[r] \ar@<-1ex>[r] &
\pi_{F, *}\mathcal{G}(R/F)
}
$$
% P11-E0106
由于还有 $\pi_{F, *}\mathcal{G}(U/F) = \mathcal{G}(U/F)$ 以及
$\pi_{F, *}\mathcal{G}(U/F) = \mathcal{G}(U/F)$，我们得到典范同一化
$f_{small, *}\pi_{F, *}\mathcal{G}(V) = \pi_{G, *}f_*\mathcal{G}(V)$。
我们省略它与限制映射相容以及关于 $\mathcal{G}$ 为函子的证明。
\end{proof}

\noindent
% P11-E0107
设 $f : \mathcal{X} \to \mathcal{Y}$ 和 $f : F \to G$ 如上面引理的
第二部分所述。利用 (\ref{stacks-sheaves-equation-restrict})，该引理的一个推论是
\begin{equation}
\label{stacks-sheaves-equation-compare-big-small}
(f_*\mathcal{F})|_{G_\etale} =
f_{small, *}(\mathcal{F}|_{F_\etale})
\end{equation}
，对 $\mathcal{X}_\etale$ 上任意层 $\mathcal{F}$ 都成立。
此外，若 $\mathcal{F}$ 是 $\mathcal{O}$-模层，则
(\ref{stacks-sheaves-equation-compare-big-small}) 是 $G_\etale$ 上
$\mathcal{O}_G$-模的同构。

\medskip\noindent
最后，假设有一个 $2$-交换图
$$
\xymatrix{
\mathcal{U} \ar[r]^a \ar[dr]_f \drtwocell<\omit>{<-2>\varphi} &
\mathcal{V} \ar[d]^g \\
& \mathcal{X}
}
$$
，其中各箭头是 $(\Sch/S)_{fppf}$ 上群胚纤维化范畴的 $1$-态射；
假设 $\mathcal{F}$ 是 $\mathcal{X}_\etale$ 上的层，并且
$\mathcal{U}, \mathcal{V}$ 可分别由代数空间 $U, V$ 表示。
于是得到比较映射
\begin{equation}
\label{stacks-sheaves-equation-comparison-algebraic-spaces}
c_\varphi : a_{small}^{-1}(g^{-1}\mathcal{F}|_{V_\etale})
\longrightarrow
f^{-1}\mathcal{F}|_{U_\etale}
\end{equation}
，其中 $a : U \to V$ 表示与 $a$ 对应的代数空间态射。
这是 (\ref{stacks-sheaves-equation-comparison}) 的类似映射。我们把 $c_\varphi$
定义为如下映射的伴随：
$$
g^{-1}\mathcal{F}|_{V_\etale}
\longrightarrow
a_{small, *}(f^{-1}\mathcal{F}|_{U_\etale}) =
(a_*f^{-1}\mathcal{F})|_{V_\etale}
$$
（该等式由 (\ref{stacks-sheaves-equation-compare-big-small}) 得到）；此映射是如下映射
到 $V$ 上的限制 (\ref{stacks-sheaves-equation-restrict})：
$$
g^{-1}\mathcal{F} \to a_*a^{-1}g^{-1}\mathcal{F} = a_*f^{-1}\mathcal{F}
$$
，其中最后一个等式使用上图的 $2$-交换性。当 $\mathcal{F}$ 是
$\mathcal{O}_\mathcal{X}$-模层时，$c_\varphi$ 诱导一个{\it 比较}映射
\begin{equation}
\label{stacks-sheaves-equation-comparison-algebraic-spaces-modules}
c_\varphi : a_{small}^*(g^*\mathcal{F}|_{V_\etale})
\longrightarrow
f^*\mathcal{F}|_{U_\etale}
\end{equation}
，它是 $\mathcal{O}_{U_\etale}$-模的映射。这是
(\ref{stacks-sheaves-equation-comparison-modules}) 的类似映射。
注意，引理 \ref{stacks-sheaves-lemma-comparison} 的性质 (1)、(2)、(3) 和 (4)
在此情形下也成立。












\section{拟凝聚模}
\label{stacks-sheaves-section-quasi-coherent}

\noindent
现在可以把拟凝聚模的一般定义应用于本章所讨论的情形。

\begin{definition}
\label{stacks-sheaves-definition-quasi-coherent}
设 $p : \mathcal{X} \to (\Sch/S)_{fppf}$ 是群胚纤维化范畴。
{\it $\mathcal{X}$ 上的拟凝聚模}，或{\it 拟凝聚
$\mathcal{O}_\mathcal{X}$-模}，是位点上的模，定义
\ref{sites-modules-definition-site-local} 意义下带环位点
$(\mathcal{X}_{fppf}, \mathcal{O}_\mathcal{X})$ 上的拟凝聚模。
$\mathcal{X}$ 上拟凝聚层的范畴记为 $\QCoh(\mathcal{O}_\mathcal{X})$。
\end{definition}

\noindent
若 $\mathcal{X}$ 是代数叠，则这个定义与文献中的所有定义相符；
其含义是，$\QCoh(\mathcal{O}_\mathcal{X})$ 与文献中定义的这个范畴的
任一变体等价（不计集合论问题）。例如，在叠上的上同调，引理
\ref{stacks-sheaves-lemma-quasi-coherent} 中，我们会把自己的定义与
\cite[Definition 6.1]{olsson_sheaves} 中的定义相匹配。稍后还会看到
这个范畴的其他构造。

\medskip\noindent
% P11-E0108
一般而言（概形态射的情形也是如此），拟凝聚层沿 $1$-态射的推前
并非拟凝聚的。拉回则保持拟凝聚性。

\begin{lemma}
\label{stacks-sheaves-lemma-pullback-quasi-coherent}
设 $f : \mathcal{X} \to \mathcal{Y}$ 是 $(\Sch/S)_{fppf}$ 上
群胚纤维化范畴的 $1$-态射。拉回函子
$f^* = f^{-1} : \textit{Mod}(\mathcal{O}_\mathcal{Y}) \to
\textit{Mod}(\mathcal{O}_\mathcal{X})$
保持拟凝聚层。
\end{lemma}

\begin{proof}
这是一般事实；参见位点上的模，引理
\ref{sites-modules-lemma-local-pullback}。
\end{proof}

\noindent
拟凝聚层恰好可以借助其拉回来作非常简单的刻画。
关于借助限制所作的刻画，还可参见引理 \ref{stacks-sheaves-lemma-quasi-coherent}。

\begin{lemma}
\label{stacks-sheaves-lemma-characterize-quasi-coherent}
设 $p : \mathcal{X} \to (\Sch/S)_{fppf}$ 是群胚纤维化范畴。
设 $\mathcal{F}$ 是 $\mathcal{O}_\mathcal{X}$-模层。则 $\mathcal{F}$
拟凝聚，当且仅当对 $\mathcal{X}$ 的每个满足 $U = p(x)$ 的对象 $x$，
$x^*\mathcal{F}$ 都是 $(\Sch/U)_{fppf}$ 上的拟凝聚层。
\end{lemma}

\begin{proof}
由引理 \ref{stacks-sheaves-lemma-pullback-quasi-coherent}，该条件是必要的。
反过来，由于 $x^*\mathcal{F}$ 正是到 $\mathcal{X}_{fppf}/x$ 的限制，
从拟凝聚层的定义直接可知该条件也充分（这里还使用拟凝聚性是模层的
内蕴性质这一事实；参见位点上的模，第
\ref{sites-modules-section-intrinsic} 节）。
\end{proof}

\begin{lemma}
\label{stacks-sheaves-lemma-characterize-quasi-coherent-bis}
设 $p : \mathcal{X} \to (\Sch/S)_{fppf}$ 是群胚纤维化范畴。
设 $\mathcal{F}$ 是 $\mathcal{X}$ 上的模预层。下列条件等价：
\begin{enumerate}
\item $\mathcal{F}$ 是
$\textit{Mod}(\mathcal{X}_{Zar}, \mathcal{O}_\mathcal{X})$
的对象，并且 $\mathcal{F}$ 是位点上的模，定义
\ref{sites-modules-definition-site-local} 意义下
$(\mathcal{X}_{Zar}, \mathcal{O}_\mathcal{X})$ 上的拟凝聚模；
\item $\mathcal{F}$ 是
$\textit{Mod}(\mathcal{X}_\etale, \mathcal{O}_\mathcal{X})$
的对象，并且 $\mathcal{F}$ 是位点上的模，定义
\ref{sites-modules-definition-site-local} 意义下
$(\mathcal{X}_\etale, \mathcal{O}_\mathcal{X})$ 上的拟凝聚模；
\item $\mathcal{F}$ 是定义 \ref{stacks-sheaves-definition-quasi-coherent} 意义下
$\mathcal{X}$ 上的拟凝聚模。
\end{enumerate}
\end{lemma}

\begin{proof}
假设 (1)、(2) 或 (3) 中任一个成立。设 $x$ 是 $\mathcal{X}$ 的对象，
位于概形 $U$ 之上。回忆 $x^*\mathcal{F} = x^{-1}\mathcal{F}$
正是到 $\mathcal{X}/x = (\Sch/U)_\tau$ 的限制，其中
$\tau = fppf$、$\tau = \etale$ 或 $\tau = Zar$；参见
第 \ref{stacks-sheaves-section-restriction} 节。由带环位点上拟凝聚模的定义，
只要 $\mathcal{F}$ 拟凝聚，这个限制就拟凝聚。由下降，命题
\ref{descent-proposition-equivalence-quasi-coherent}，可见
$x^*\mathcal{F}$ 是与一个拟凝聚 $\mathcal{O}_U$-模相联系的层，
因而在 fppf、\'etale 和 Zariski 拓扑中都是拟凝聚模；
这里还使用下降，引理 \ref{descent-lemma-sheaf-condition-holds}
和定义 \ref{descent-definition-structure-sheaf}。
由于这对 $\mathcal{X}$ 的每个对象 $x$ 都成立，可见 $\mathcal{F}$
在这三种拓扑中的任一种下都是层。此外，由拟凝聚性的定义以及
$x$ 是 $\mathcal{X}$ 的任意对象这一事实，直接可知 $\mathcal{F}$
在这三种拓扑中的任一种下都是拟凝聚的。
\end{proof}







\section{局部拟凝聚模}
\label{stacks-sheaves-section-locally-quasi-coherent}

\noindent
尽管存在关于 Zariski 拓扑的变体，\'etale 拓扑看来才是下述定义中
自然应当使用的拓扑。

\begin{definition}
\label{stacks-sheaves-definition-locally-quasi-coherent}
设 $p : \mathcal{X} \to (\Sch/S)_{fppf}$ 是群胚纤维化范畴。
设 $\mathcal{F}$ 是 $\mathcal{O}_\mathcal{X}$-模预层。若
$\mathcal{F}$ 是 \'etale 拓扑的层，并且对 $\mathcal{X}$ 的每个对象 $x$，
限制 $x^*\mathcal{F}|_{U_\etale}$ 都是拟凝聚层，其中 $U = p(x)$，
则称 $\mathcal{F}$ {\it 局部拟凝聚}\footnote{这是非标准记号。}。
\end{definition}

\noindent
用 $\textit{LQCoh}(\mathcal{O}_\mathcal{X})$ 表示局部拟凝聚模的范畴。
现在有如下模范畴的图：
$$
\xymatrix{
\QCoh(\mathcal{O}_\mathcal{X}) \ar[r] \ar[d] &
\textit{Mod}(\mathcal{O}_\mathcal{X}) \ar[d] \\
\textit{LQCoh}(\mathcal{O}_\mathcal{X}) \ar[r] &
\textit{Mod}(\mathcal{X}_\etale, \mathcal{O}_\mathcal{X})
}
$$
，其中各箭头都是严格满嵌入。拟凝聚层的许多结果都有局部拟凝聚模的
% P11-E0109
对应结果。而且，从许多角度来看（稍后将会看到），这是一个自然的
研究范畴。例如，拟凝聚层恰好是那些“笛卡儿的”局部拟凝聚模，
即满足下一个引理的第二个条件的模。

\begin{lemma}
\label{stacks-sheaves-lemma-quasi-coherent}
设 $p : \mathcal{X} \to (\Sch/S)_{fppf}$ 是群胚纤维化范畴。
设 $\mathcal{F}$ 是 $\mathcal{O}_\mathcal{X}$-模预层。则 $\mathcal{F}$
拟凝聚，当且仅当下列两个条件成立：
\begin{enumerate}
\item $\mathcal{F}$ 局部拟凝聚；
\item 对 $\mathcal{X}$ 中任意位于 $f : U \to V$ 之上的态射
$\varphi : x \to y$，比较映射
$c_\varphi : f_{small}^*\mathcal{F}|_{V_\etale} \to
\mathcal{F}|_{U_\etale}$，即 (\ref{stacks-sheaves-equation-comparison-modules})，
是同构。
\end{enumerate}
\end{lemma}

\begin{proof}
假设 $\mathcal{F}$ 拟凝聚。则 $\mathcal{F}$ 是 fppf 拓扑的层，
因而也是 \'etale 拓扑的层。此外，$\mathcal{F}$ 到任意带环拓扑斯的
拉回都拟凝聚，故各限制 $x^*\mathcal{F}|_{U_\etale}$ 拟凝聚。
这证明了 $\mathcal{F}$ 局部拟凝聚。设 $y$ 是 $\mathcal{X}$ 的对象，
且 $V = p(y)$。我们已经看到 $\mathcal{X}/y = (\Sch/V)_{fppf}$。
由下降，命题 \ref{descent-proposition-equivalence-quasi-coherent}，
$y^*\mathcal{F}$ 是与概形 $V$ 上一个（通常意义下的）拟凝聚模
$\mathcal{F}_V$ 相联系的拟凝聚模。因此，比较映射
(\ref{stacks-sheaves-equation-comparison-modules}) 当然都是同构。

\medskip\noindent
反过来，假设 $\mathcal{F}$ 满足 (1) 和 (2)。设 $y$ 是
$\mathcal{X}$ 的对象，且 $V = p(y)$。把概形 $V$ 上与限制
$y^*\mathcal{F}|_{V_\etale}$ 对应的拟凝聚模记为 $\mathcal{F}_V$；
由假设 (1)，该限制是拟凝聚的；参见下降，命题
\ref{descent-proposition-equivalence-quasi-coherent}。
条件 (2) 现在表示：对 $y$ 上的每个 $x$，限制
$x^*\mathcal{F}|_{U_\etale}$ 都同构于 $\mathcal{F}_V$ 经相应概形态射
$U \to V$ 的拉回（所联系的 \'etale 层）。因此 $y^*\mathcal{F}$
是 $(\Sch/V)_{fppf}$ 上与 $\mathcal{F}_V$ 相联系的层。
所以它拟凝聚（再次使用下降，命题
\ref{descent-proposition-equivalence-quasi-coherent}），进而由引理
\ref{stacks-sheaves-lemma-characterize-quasi-coherent} 可知 $\mathcal{F}$ 在
$\mathcal{X}$ 上拟凝聚。
\end{proof}

\begin{lemma}
\label{stacks-sheaves-lemma-pullback-lqc}
设 $f : \mathcal{X} \to \mathcal{Y}$ 是 $(\Sch/S)_{fppf}$ 上
群胚纤维化范畴的 $1$-态射。拉回函子
$f^* = f^{-1} :
\textit{Mod}(\mathcal{Y}_\etale, \mathcal{O}_\mathcal{Y})
\to
\textit{Mod}(\mathcal{X}_\etale, \mathcal{O}_\mathcal{X})$
保持局部拟凝聚层。
\end{lemma}

\begin{proof}
设 $\mathcal{G}$ 在 $\mathcal{Y}$ 上局部拟凝聚。
选取 $\mathcal{X}$ 的一个对象 $x$，它位于概形 $U$ 之上。
限制 $x^*f^*\mathcal{G}|_{U_\etale}$ 等于
$(f \circ x)^*\mathcal{G}|_{U_\etale}$，因而由对 $\mathcal{G}$
的假设，它是拟凝聚层。
\end{proof}

\begin{lemma}
\label{stacks-sheaves-lemma-lqc-colimits}
设 $p : \mathcal{X} \to (\Sch/S)_{fppf}$ 是群胚纤维化范畴。
\begin{enumerate}
\item 范畴 $\textit{LQCoh}(\mathcal{O}_\mathcal{X})$ 有余极限，
并且它们与范畴
$\textit{Mod}(\mathcal{X}_\etale, \mathcal{O}_\mathcal{X})$
中的余极限一致。
\item 范畴 $\textit{LQCoh}(\mathcal{O}_\mathcal{X})$ 是 Abel 范畴，
其核与余核在
$\textit{Mod}(\mathcal{X}_\etale, \mathcal{O}_\mathcal{X})$
中计算；换言之，包含函子是正合的。
\item 给定
$0 \to \mathcal{F}_1 \to \mathcal{F}_2 \to \mathcal{F}_3 \to 0$，它是
$\textit{Mod}(\mathcal{X}_\etale, \mathcal{O}_\mathcal{X})$
中的短正合列；若三个模中的两个局部拟凝聚，则第三个也局部拟凝聚。
\item 给定 $\textit{LQCoh}(\mathcal{O}_\mathcal{X})$ 中的
$\mathcal{F}, \mathcal{G}$，在
$\textit{Mod}(\mathcal{X}_\etale, \mathcal{O}_\mathcal{X})$ 中的张量积
$\mathcal{F} \otimes_{\mathcal{O}_\mathcal{X}} \mathcal{G}$
是 $\textit{LQCoh}(\mathcal{O}_\mathcal{X})$ 的对象。
\item 给定 $\textit{LQCoh}(\mathcal{O}_\mathcal{X})$ 中的
$\mathcal{F}, \mathcal{G}$，其中 $\mathcal{F}$ 在
$\mathcal{X}_\etale$ 上有限表示，则层
$\SheafHom_{\mathcal{O}_\mathcal{X}}(\mathcal{F}, \mathcal{G})$
在 $\textit{Mod}(\mathcal{X}_\etale, \mathcal{O}_\mathcal{X})$
中是 $\textit{LQCoh}(\mathcal{O}_\mathcal{X})$ 的对象。
\end{enumerate}
\end{lemma}

\begin{proof}
在下面的论证中，$x$ 表示 $\mathcal{X}$ 的任意对象，它位于概形 $U$
之上。为了证明 $\textit{Mod}(\mathcal{X}_\etale,
\mathcal{O}_\mathcal{X})$ 的一个对象 $\mathcal{H}$ 属于
$\textit{LQCoh}(\mathcal{O}_\mathcal{X})$，我们将证明限制
$x^*\mathcal{H}|_{U_\etale} = \mathcal{H}|_{U_\etale}$ 是
$\textit{Mod}(U_\etale, \mathcal{O}_U)$ 的拟凝聚对象。

\medskip\noindent
证明 (1)。设 $\mathcal{I} \to \textit{LQCoh}(\mathcal{O}_\mathcal{X})$，
$i \mapsto \mathcal{F}_i$ 是一个图。考虑
$\textit{Mod}(\mathcal{X}_\etale, \mathcal{O}_\mathcal{X})$ 的对象
$\mathcal{F} = \colim_i \mathcal{F}_i$。拉回函子 $x^*$ 是左伴随，
故与所有余极限交换。因此
$x^*\mathcal{F} = \colim_i x^*\mathcal{F}_i$。类似地，有
$x^*\mathcal{F}|_{U_\etale} =
\colim_i x^*\mathcal{F}_i|_{U_\etale}$.
根据假设，每个 $x^*\mathcal{F}_i|_{U_\etale}$ 都拟凝聚。
因此，由下降，引理 \ref{descent-lemma-properties-quasi-coherent}，
$\colim_i x^*\mathcal{F}_i|_{U_\etale}$ 拟凝聚。
所以 $x^*\mathcal{F}|_{U_\etale}$ 如所需地拟凝聚。

\medskip\noindent
证明 (2)。由 (1)，$\textit{LQCoh}(\mathcal{O}_\mathcal{X})$ 中存在余核，
并且它们与在 $\textit{Mod}(\mathcal{X}_\etale,
\mathcal{O}_\mathcal{X})$ 中计算的余核一致。设
$\varphi : \mathcal{F} \to \mathcal{G}$ 是
$\textit{LQCoh}(\mathcal{O}_\mathcal{X})$ 的态射，并设
$\mathcal{K} = \Ker(\varphi)$ 是在
$\textit{Mod}(\mathcal{X}_\etale, \mathcal{O}_\mathcal{X})$ 中计算的。
若能证明 $\mathcal{K}$ 是局部拟凝聚模，(2) 的证明便完成了。
为此，注意核是在预层范畴中计算的（无需层化）。因此
$\mathcal{K}|_{U_\etale}$ 是映射
$\mathcal{F}|_{U_\etale} \to \mathcal{G}|_{U_\etale}$ 的核，
即 $U_\etale$ 上拟凝聚层之间一个映射的核；故由下降，引理
\ref{descent-lemma-properties-quasi-coherent}，它拟凝聚。这证明了 (2)。

\medskip\noindent
证明 (3)。设
$0 \to \mathcal{F}_1 \to \mathcal{F}_2 \to \mathcal{F}_3 \to 0$
是 $\textit{Mod}(\mathcal{X}_\etale, \mathcal{O}_\mathcal{X})$
中的短正合列。由于我们使用 \'etale 拓扑，限制
$0 \to \mathcal{F}_1|_{U_\etale} \to
\mathcal{F}_2|_{U_\etale} \to \mathcal{F}_3|_{U_\etale} \to 0$
也是短正合列。因此，(3) 来自下降，引理
\ref{descent-lemma-properties-quasi-coherent} 中的相应陈述。

\medskip\noindent
证明 (4)。设 $\mathcal{F}$ 和 $\mathcal{G}$ 属于
$\textit{LQCoh}(\mathcal{O}_\mathcal{X})$。由于到 $U_\etale$ 的限制
由沿带环拓扑斯态射的拉回给出
$U_\etale \to (\Sch/U)_\etale \to \mathcal{X}_\etale$
，可见张量积
$\mathcal{F} \otimes_{\mathcal{O}_\mathcal{X}} \mathcal{G}$
到 $U_\etale$ 的限制等于
$\mathcal{F}|_{U_\etale} \otimes_{\mathcal{O}_U} \mathcal{G}|_{U_\etale}$,
；参见位点上的模，引理 \ref{sites-modules-lemma-tensor-product-pullback}。
由于 $\mathcal{F}|_{U_\etale}$ 和 $\mathcal{G}|_{U_\etale}$
拟凝聚，它们的张量积也拟凝聚；参见下降，引理
\ref{descent-lemma-properties-quasi-coherent}。

\medskip\noindent
证明 (5)。设 $\mathcal{F}$ 和 $\mathcal{G}$ 属于
$\textit{LQCoh}(\mathcal{O}_\mathcal{X})$，且 $\mathcal{F}$ 有限表示。
由于
$(\Sch/U)_\etale = \mathcal{X}_\etale/x$
是 $\mathcal{X}_\etale$ 在一个对象处的局部化，可见
$\SheafHom_{\mathcal{O}_\mathcal{X}}(\mathcal{F}, \mathcal{G})$
到 $(\Sch/U)_\etale$ 的限制等于
$$
\mathcal{H} =
\SheafHom_{\mathcal{O}|_{(\Sch/U)_\etale}}(
\mathcal{F}|_{(\Sch/U)_\etale}, \mathcal{G}|_{(\Sch/U)_\etale})
$$
；这是由位点上的模，引理
\ref{sites-modules-lemma-internal-hom-restriction} 得到的。带环拓扑斯态射
$(U_\etale, \mathcal{O}_U) \to ((\Sch/U)_\etale, \mathcal{O})$
是平坦的，因为 $\mathcal{O}$ 的拉回是 $\mathcal{O}_U$。
因此，$\mathcal{H}$ 经此态射的拉回等于
$\SheafHom_{\mathcal{O}_U}(\mathcal{F}|_{U_\etale}, \mathcal{G}|_{U_\etale})$
；这是由位点上的模，引理
\ref{sites-modules-lemma-pullback-internal-hom} 得到的。换言之，
$\SheafHom_{\mathcal{O}_\mathcal{X}}(\mathcal{F}, \mathcal{G})$
到 $U_\etale$ 的限制是
$\SheafHom_{\mathcal{O}_U}(\mathcal{F}|_{U_\etale}, \mathcal{G}|_{U_\etale})$.
由于 $\mathcal{F}|_{U_\etale}$ 和 $\mathcal{G}|_{U_\etale}$ 拟凝聚，
$\SheafHom_{\mathcal{O}_U}(\mathcal{F}|_{U_\etale}, \mathcal{G}|_{U_\etale})$,
也拟凝聚；参见下降，引理 \ref{descent-lemma-properties-quasi-coherent}。
结论与前面相同。
\end{proof}

\noindent
在这里所讨论的一般性下，拟凝聚层的范畴不是 Abel 范畴。
参见例子，第 \ref{examples-section-nonabelian-QCoh} 节。
无需进一步工作，我们可以证明如下结果。

\begin{lemma}
\label{stacks-sheaves-lemma-qc-colimits}
设 $p : \mathcal{X} \to (\Sch/S)_{fppf}$ 是群胚纤维化范畴。
\begin{enumerate}
\item 范畴 $\QCoh(\mathcal{O}_\mathcal{X})$ 有余极限，
并且它们与下列范畴中的余极限一致：
$\textit{Mod}(\mathcal{X}_{Zar}, \mathcal{O}_\mathcal{X})$,
$\textit{Mod}(\mathcal{X}_\etale, \mathcal{O}_\mathcal{X})$,
$\textit{Mod}(\mathcal{O}_\mathcal{X})$ 以及
$\textit{LQCoh}(\mathcal{O}_\mathcal{X})$.
\item 给定 $\QCoh(\mathcal{O}_\mathcal{X})$ 中的
$\mathcal{F}, \mathcal{G}$，在
$\textit{Mod}(\mathcal{X}_{Zar}, \mathcal{O}_\mathcal{X})$、
$\textit{Mod}(\mathcal{X}_\etale, \mathcal{O}_\mathcal{X})$ 或
$\textit{Mod}(\mathcal{O}_\mathcal{X})$ 中计算的张量积
$\mathcal{F} \otimes_{\mathcal{O}_\mathcal{X}} \mathcal{G}$ 相同，
且其公共值是 $\QCoh(\mathcal{O}_\mathcal{X})$ 的对象。
\item 给定 $\QCoh(\mathcal{O}_\mathcal{X})$ 中的
$\mathcal{F}, \mathcal{G}$，其中 $\mathcal{F}$ 有限局部自由
（在 fppf 拓扑中，或等价地在 \'etale 拓扑中，或等价地在
Zariski 拓扑中），则在
$\textit{Mod}(\mathcal{X}_{Zar}, \mathcal{O}_\mathcal{X})$、
$\textit{Mod}(\mathcal{X}_\etale, \mathcal{O}_\mathcal{X})$ 或
$\textit{Mod}(\mathcal{O}_\mathcal{X})$ 中计算的内部 Hom
$\SheafHom_{\mathcal{O}_\mathcal{X}}(\mathcal{F}, \mathcal{G})$
相同，且其公共值是 $\QCoh(\mathcal{O}_\mathcal{X})$ 的对象。
\end{enumerate}
\end{lemma}

\begin{proof}
设 $x$ 是 $\mathcal{X}$ 的任意对象，它位于概形 $U$ 之上。
% P11-E0111
设 $\tau \in \{Zariski, \etale, fppf\}$。为了证明
$\textit{Mod}(\mathcal{X}_\tau, \mathcal{O}_\mathcal{X})$ 的对象
$\mathcal{H}$ 属于 $\QCoh(\mathcal{O}_\mathcal{X})$，只须证明限制
% P11-E0112
$x^*\mathcal{H}$（第 \ref{stacks-sheaves-section-restriction} 节）是
$\textit{Mod}((\Sch/U)_\tau, \mathcal{O})$ 的拟凝聚对象。
参见引理 \ref{stacks-sheaves-lemma-characterize-quasi-coherent} 和
\ref{stacks-sheaves-lemma-characterize-quasi-coherent-bis}。有限局部自由性也同样。
回忆 $(\Sch/U)_\tau = \mathcal{X}_\tau/x$ 是 $\mathcal{X}_\tau$
在一个对象处的局部化。因此，限制与余极限、张量积以及形成内部 Hom 交换
（参见位点上的模，引理
\ref{sites-modules-lemma-exactness-pushforward-pullback}、
\ref{sites-modules-lemma-tensor-product-pullback} 和
\ref{sites-modules-lemma-internal-hom-restriction}).
这把本引理约化为下降，引理 \ref{descent-lemma-qc-colimits}。
\end{proof}






\section{叠化与层}
\label{stacks-sheaves-section-stackification}

\noindent
群胚纤维化范畴上的层范畴，实际上只“知道”它的叠化。

\begin{lemma}
\label{stacks-sheaves-lemma-stackification}
设 $f : \mathcal{X} \to \mathcal{Y}$ 是 $(\Sch/S)_{fppf}$ 上
群胚纤维化范畴的 $1$-态射。若 $f$ 诱导叠化的等价，则拓扑斯态射
$f : \Sh(\mathcal{X}_{fppf}) \to \Sh(\mathcal{Y}_{fppf})$
是等价。
\end{lemma}

\begin{proof}
可以假设 $\mathcal{Y}$ 是 $\mathcal{X}$ 的叠化。我们断言
$f : \mathcal{X} \to \mathcal{Y}$ 是特殊余连续函子；参见位点，定义
\ref{sites-definition-special-cocontinuous-functor}；这将证明本引理。
由 Stacks，引理 \ref{stacks-lemma-topology-inherited-functorial}，
函子 $f$ 连续且余连续。由 Stacks，引理
\ref{stacks-lemma-stackify}，可见位点，引理
\ref{sites-lemma-equivalence} 的条件 (3)、(4) 和 (5) 成立。
\end{proof}

\begin{lemma}
\label{stacks-sheaves-lemma-stackification-quasi-coherent}
设 $f : \mathcal{X} \to \mathcal{Y}$ 是 $(\Sch/S)_{fppf}$ 上
群胚纤维化范畴的 $1$-态射。若 $f$ 诱导叠化的等价，则 $f^*$
诱导等价
% P11-E0116
$\textit{Mod}(\mathcal{O}_\mathcal{X}) \to
\textit{Mod}(\mathcal{O}_\mathcal{Y})$
以及
$\QCoh(\mathcal{O}_\mathcal{X}) \to
\QCoh(\mathcal{O}_\mathcal{Y})$.
\end{lemma}

\begin{proof}
可以假设 $\mathcal{Y}$ 是 $\mathcal{X}$ 的叠化。由引理
\ref{stacks-sheaves-lemma-stackification} 以及
$\mathcal{O}_\mathcal{X} = f^{-1}\mathcal{O}_\mathcal{Y}$，
第一个断言是清楚的。拟凝聚层的拉回拟凝聚；参见引理
% P11-E0114
\ref{stacks-sheaves-lemma-pullback-quasi-coherent}。因此，只须证明：若
$f^*\mathcal{G}$ 拟凝聚，则 $\mathcal{G}$ 拟凝聚。
为此，设 $y$ 是 $\mathcal{Y}$ 的一个对象。将 $\mathcal{Y}$ 是
$\mathcal{X}$ 的叠化这一条件展开，可见在 $\mathcal{Y}$ 中存在
一个 fppf 覆盖 $\{y_i \to y\}$，使得对 $\mathcal{X}$ 的某个对象
% P11-E0115
$x_i$，有 $y_i \cong f(x_i)$。设 $x_i$ 和 $y_i$ 位于概形 $U_i$ 之上。
于是，$f^*\mathcal{G}$ 拟凝聚意味着 $x_i^*f^*\mathcal{G}$ 拟凝聚。
由于 $x_i^*f^*\mathcal{G}$ 同构于 $y_i^*\mathcal{G}$（作为
% P11-E0113
$(\Sch/U_i)_{fppf}$ 上的层，可见 $y_i^*\mathcal{G}$ 拟凝聚。
由位点上的模，引理 \ref{sites-modules-lemma-local-final-object}，
$\mathcal{G}$ 到 $\mathcal{Y}/y$ 的限制拟凝聚。因此，由引理
\ref{stacks-sheaves-lemma-characterize-quasi-coherent}，$\mathcal{G}$ 拟凝聚。
\end{proof}





\section{拟凝聚层与表现}
\label{stacks-sheaves-section-quasi-coherent-presentation}

\noindent
首先把拟凝聚层与我们先前为概形和代数空间定义的概念相匹配。

\begin{lemma}
\label{stacks-sheaves-lemma-compare-locally-quasi-coherent}
设 $S$ 是概形。设 $\mathcal{X} \to (\Sch/S)_{fppf}$ 是可由代数空间
$F$ 表示的群胚纤维化范畴。若 $\mathcal{F}$ 属于
$\textit{LQCoh}(\mathcal{O}_\mathcal{X})$，则限制
$\mathcal{F}|_{F_\etale}$（\ref{stacks-sheaves-equation-restrict}）拟凝聚。
\end{lemma}

\begin{proof}
设 $U$ 是 $F$ 上的 \'etale 概形。则
$\mathcal{F}|_{U_\etale} = (\mathcal{F}|_{F_\etale})|_{U_\etale}$.
这是清楚的，不过也可参见注 \ref{stacks-sheaves-remark-compare}。
因此该断言由定义得到。
\end{proof}

\begin{lemma}
\label{stacks-sheaves-lemma-compare-quasi-coherent}
设 $S$ 是概形。设 $\mathcal{X} \to (\Sch/S)_{fppf}$ 是可由代数空间
$F$ 表示的群胚纤维化范畴。函子 (\ref{stacks-sheaves-equation-restrict}) 定义等价
$$
\QCoh(\mathcal{O}_\mathcal{X}) \to \QCoh(\mathcal{O}_F),\quad
\mathcal{F} \longmapsto \mathcal{F}|_{F_\etale}
$$
，其拟逆由 $\mathcal{G} \mapsto \pi_F^*\mathcal{G}$ 给出。
对于可由代数空间表示的群胚纤维化范畴之间的态射，这个等价与拉回相容。
\end{lemma}

\begin{proof}
由引理 \ref{stacks-sheaves-lemma-characterize-quasi-coherent-bis}，可以使用 \'etale 拓扑。
我们将直接使用引理 \ref{stacks-sheaves-lemma-compare} 的记号和结果而不再说明。
回忆限制函子
$\textit{Mod}(\mathcal{X}_\etale, \mathcal{O}_\mathcal{X}) \to
\textit{Mod}(F_\etale, \mathcal{O}_F)$,
$\mathcal{F} \mapsto \mathcal{F}|_{F_\etale}$ 由 $i_F^*$ 给出。
由引理 \ref{stacks-sheaves-lemma-compare-locally-quasi-coherent}，或由位点上的模，
引理 \ref{sites-modules-lemma-local-pullback}，可见若 $\mathcal{F}$
拟凝聚，则 $\mathcal{F}|_{F_\etale}$ 拟凝聚。因此，我们得到引理陈述中
所指的函子，并得到反方向的函子 $\pi_F^*$。由于
$\pi_F \circ i_F = \text{id}$，可见
$i_F^*\pi_F^*\mathcal{G} = \mathcal{G}$.

\medskip\noindent
对 $\textit{Mod}(\mathcal{X}_\etale, \mathcal{O}_\mathcal{X})$ 中的
$\mathcal{F}$，有典范映射
$\pi_F^*(\mathcal{F}|_{F_\etale}) \to \mathcal{F}$，即同一化
$\mathcal{F}|_{F_\etale} = \pi_{F, *}\mathcal{F}$ 的伴随映射。
我们将证明：若 $\mathcal{F}$ 是 $\mathcal{X}$ 上的拟凝聚模，
则此映射是同构。选取概形 $U$ 和满的 \'etale 态射 $U \to F$。
% P11-E0117
把 $\mathcal{X}$ 在 $U$ 上的对应对象记为 $x : U \to \mathcal{X}$。
只须证明 $\pi_F^*(\mathcal{F}|_{F_\etale}) \to \mathcal{F}$
限制到 $\mathcal{X}_\etale/x = (\Sch/U)_\etale$ 后是同构。
由于 $U \to F$ 是 \'etale 的，由注 \ref{stacks-sheaves-remark-compare} 可得
$$
\pi_F^*(\mathcal{F}|_{F_\etale})|_{\mathcal{X}_\etale/x} =
\pi_U^*(\mathcal{F}|_{U_\etale})
$$
，并且映射
$\pi_F^*(\mathcal{F}|_{F_\etale}) \to \mathcal{F}$
到 $\mathcal{X}_\etale/x = (\Sch/U)_\etale$ 的限制等于相应映射
$\pi_U^*(\mathcal{F}|_{U_\etale}) \to \mathcal{F}|_{(\Sch/U)_\etale}$.
我们已经在下降，第 \ref{descent-section-quasi-coherent-sheaves} 节
看到该结果对概形成立\footnote{具体而言，若 $U$ 是概形，且
$\mathcal{F}$ 在 $(\Sch/U)_\etale$ 上拟凝聚，则由下降，命题
\ref{descent-proposition-equivalence-quasi-coherent}，存在概形 $U$
上的某个拟凝聚模 $\mathcal{H}$，使得 $\mathcal{F} = \mathcal{H}^a$。
换言之，由下降，注 \ref{descent-remark-change-topologies-ringed-sites}，
并使用下降，引理 \ref{descent-lemma-compare-sites} 的记号，有
$\mathcal{F} = (\text{id}_{\etale,Zar})^*\mathcal{H}$。此外，
$\text{id}_{\etale,Zar} = \pi_U \circ \text{id}_{small,\etale,Zar}$
，因而 $\mathcal{F} = \pi_U^*\mathcal{G}$，其中
$\mathcal{G} = (\text{id}_{small,\etale,Zar})^*\mathcal{H}$ 拟凝聚。于是
$\pi_U^*i_U^*\mathcal{F} = \pi_U^*i_U^*\pi_U^*\mathcal{G} =
\pi_U^*\mathcal{G} = \mathcal{F}$，如所需。}，故结论成立。

\medskip\noindent
与拉回的相容性来自如下事实：拟逆由 $\pi_F^*$ 给出，
而且引理 \ref{stacks-sheaves-lemma-compare-morphism} 中的带环拓扑斯图交换。
\end{proof}

\noindent
在空间中的群胚，定义 \ref{spaces-groupoids-definition-groupoid-module} 中，
% P11-E0118
我们已经定义了任意群胚上的拟凝聚模。下列（形式）命题说明，
可以用表现上的拟凝聚模来研究商叠上的拟凝聚层。

\begin{proposition}
\label{stacks-sheaves-proposition-quasi-coherent}
设 $(U, R, s, t, c)$ 是 $S$ 上代数空间中的群胚。
设 $\mathcal{X} = [U/R]$ 是商叠。$\mathcal{X}$ 上的拟凝聚模范畴
等价于 $(U, R, s, t, c)$ 上的拟凝聚模范畴。
\end{proposition}

\begin{proof}
我们将构造拟逆函子
$$
\QCoh(\mathcal{O}_\mathcal{X})
\longleftrightarrow
\QCoh(U, R, s, t, c).
$$
，其中 $\QCoh(U, R, s, t, c)$ 表示群胚
$(U, R, s, t, c)$ 上的拟凝聚模范畴。

\medskip\noindent
设 $\mathcal{F}$ 是 $\QCoh(\mathcal{O}_\mathcal{X})$ 的对象。
% P11-E0119
把与 $U$ 和 $R$ 对应的群胚纤维化范畴记为 $\mathcal{U}$、$\mathcal{R}$，
并把 $\mathcal{X}$ 在 $U$ 上的（定义性）对象记为 $x$。
回忆有一个 $2$-交换图
$$
\xymatrix{
\mathcal{R} \ar[r]_s \ar[d]_t & \mathcal{U} \ar[d]^x \\
\mathcal{U} \ar[r]^x & \mathcal{X}
}
$$
；参见空间中的群胚，引理
\ref{spaces-groupoids-lemma-quotient-stack-2-arrow}。
由引理 \ref{stacks-sheaves-lemma-2-morphisms-presheaves}，图中内含的 $2$-箭头诱导同构
$\alpha : t^*x^*\mathcal{F} \to s^*x^*\mathcal{F}$
，它在 $\mathcal{R} \times_{s, \mathcal{U}, t} \mathcal{R}$ 上满足
余圈条件；这是空间中的群胚，引理
\ref{spaces-groupoids-lemma-quotient-stack-cocycle} 的推论。
因此，若置 $\mathcal{G} = x^*\mathcal{F}|_{U_\etale}$，则引理
\ref{stacks-sheaves-lemma-compare-quasi-coherent} 中的范畴等价
（以与拉回相容的方式使用数次）给出同构
$\alpha : t_{small}^*\mathcal{G} \to s_{small}^*\mathcal{G}$
，它在 $R \times_{s, U, t} R$ 上满足余圈条件；也就是说，
$(\mathcal{G}, \alpha)$ 是 $\QCoh(U, R, s, t, c)$ 的对象。
规则 $\mathcal{F} \mapsto (\mathcal{G}, \alpha)$ 就是从左到右的函子。

\medskip\noindent
构造反方向的函子。设 $(\mathcal{G}, \alpha)$ 是
$\QCoh(U, R, s, t, c)$ 的对象。由引理
\ref{stacks-sheaves-lemma-stackification-quasi-coherent}，叠化映射
$[U/_{\!p}R] \to [U/R]$（参见空间中的群胚，定义
\ref{spaces-groupoids-definition-quotient-stack}）诱导拟凝聚层范畴的等价。
因此，只须在 $[U/_{\!p}R]$ 上构造拟凝聚模 $\mathcal{F}$。

\medskip\noindent
回忆 $[U/_{\!p}R]$ 的对象 $x = (T, u)$ 由概形 $T$ 和态射
$u : T \to U$ 给出。态射 $(T, u) \to (T', u')$ 由一对 $(f, r)$ 给出，
其中 $f : T \to T'$、$r : T \to R$，且 $s \circ r = u$、
$t \circ r = u' \circ f$。把形如
$(f, e \circ u' \circ f) : (T, u' \circ f) \to (T', u')$.
的任意态射称为{\it 特殊态射}。以 $(T, u)$ 为对象、以特殊态射为态射的
范畴，正是 $U$ 上概形的范畴。

\medskip\noindent
采用上述记号，给定 $[U/_{\!p}R]$ 的对象 $(T, u)$，置
$$
\mathcal{F}(T, u) : = \Gamma(T, u_{small}^*\mathcal{G}).
$$
给定态射 $(f, r) : (T, u) \to (T', u')$，得到映射
\begin{align*}
\mathcal{F}(T', u') & = \Gamma(T', (u')_{small}^*\mathcal{G}) \\
& \to \Gamma(T, f_{small}^*(u')_{small}^*\mathcal{G}) =
\Gamma(T, (u' \circ f)_{small}^*\mathcal{G}) \\
& = \Gamma(T, (t \circ r)_{small}^*\mathcal{G}) =
\Gamma(T, r_{small}^*t_{small}^*\mathcal{G}) \\
& \to \Gamma(T, r_{small}^*s_{small}^*\mathcal{G}) =
\Gamma(T, (s \circ r)_{small}^*\mathcal{G}) \\
& = \Gamma(T, u_{small}^*\mathcal{G}) \\
& = \mathcal{F}(T, u)
\end{align*}
，其中第一个箭头是沿 $f$ 的拉回，第二个箭头是 $\alpha$。
注意，若 $(f, r)$ 是特殊态射，则由群胚上拟凝聚模层的公理，
$e_{small}^*\alpha = \text{id}$，因而这个映射正是沿 $f$ 的拉回。
余圈条件蕴含 $\mathcal{F}$ 是模预层（细节从略）。在特殊态射的情形下，
由 $\mathcal{F}$ 的限制映射的简单描述可知，$\mathcal{F}$ 到
$(\Sch/T)_{fppf}$ 的限制拟凝聚。因此，$\mathcal{F}$ 是
$[U/_{\!p}R]$ 上的层并且拟凝聚（引理
\ref{stacks-sheaves-lemma-characterize-quasi-coherent}）。

\medskip\noindent
我们省略上述两个函子互为拟逆的验证。
\end{proof}

\noindent
本节最后给出一个关于从拟凝聚层出发的映射的技术引理。
它是概形，引理 \ref{schemes-lemma-compare-constructions} 的类似结果。
稍后将看到（可表性判据，定理
\ref{criteria-theorem-flat-groupoid-gives-algebraic-stack}），
关于群胚的假设蕴含 $\mathcal{X}$ 是代数叠。

\begin{lemma}
\label{stacks-sheaves-lemma-map-from-quasi-coherent}
设 $(U, R, s, t, c)$ 是 $S$ 上代数空间中的群胚。
假设 $s, t$ 平坦且局部有限表示。设 $\mathcal{X} = [U/R]$ 是商叠。
% P11-E0120
把 $\mathcal{X}$ 在 $U$ 上的对象记为 $x$。设 $\mathcal{F}$ 是拟凝聚
$\mathcal{O}_\mathcal{X}$-模，并设 $\mathcal{H}$ 是
$\textit{Mod}(\mathcal{O}_\mathcal{X})$ 的任意对象。映射
$$
\Hom_{\mathcal{O}_\mathcal{X}}(\mathcal{F}, \mathcal{H})
\longrightarrow
\Hom_{\mathcal{O}_U}(x^*\mathcal{F}|_{U_\etale},
x^*\mathcal{H}|_{U_\etale}),
\quad
\phi \longmapsto x^*\phi|_{U_\etale}
$$
是单射，并且其像恰由如下映射组成：这些映射
$\varphi : x^*\mathcal{F}|_{U_\etale} \to
x^*\mathcal{H}|_{U_\etale}$ 使下图交换：
$$
\xymatrix{
s_{small}^*(x^*\mathcal{F}|_{U_\etale})
\ar[r] \ar[d]^{s_{small}^*\varphi} &
(x \circ s)^*\mathcal{F}|_{R_\etale} =
(x \circ t)^*\mathcal{F}|_{R_\etale} &
t_{small}^*(x^*\mathcal{F}|_{U_\etale})
\ar[l] \ar[d]_{t_{small}^*\varphi} \\
s_{small}^*(x^*\mathcal{H}|_{U_\etale})
\ar[r] &
(x \circ s)^*\mathcal{H}|_{R_\etale} =
(x \circ t)^*\mathcal{H}|_{R_\etale} &
t_{small}^*(x^*\mathcal{H}|_{U_\etale})
\ar[l]
}
$$
这是 $R_\etale$ 上模的图，其中水平箭头是比较映射
(\ref{stacks-sheaves-equation-comparison-algebraic-spaces-modules})。
\end{lemma}

\begin{proof}
由引理 \ref{stacks-sheaves-lemma-stackification-quasi-coherent}，叠化映射
$[U/_{\!p}R] \to [U/R]$（参见空间中的群胚，定义
\ref{spaces-groupoids-definition-quotient-stack}）诱导拟凝聚层范畴以及
fppf $\mathcal{O}$-模范畴的等价。因此，只须在
$\mathcal{X} = [U/_{\!p}R]$ 时证明本引理。由命题
\ref{stacks-sheaves-proposition-quasi-coherent} 及其证明，存在
$(U, R, s, t, c)$ 上的拟凝聚模 $(\mathcal{G}, \alpha)$，
使得 $\mathcal{F}$ 由规则
$\mathcal{F}(T, u) = \Gamma(T, u^*\mathcal{G})$.
给出。特别地，$x^*\mathcal{F}|_{U_\etale} = \mathcal{G}$，
而且引理陈述中的映射显然是单射。此外，给定映射
$\varphi : \mathcal{G} \to x^*\mathcal{H}|_{U_\etale}$
以及 $[U/_{\!p}R]$ 的任意对象 $y = (T, u)$，可以考虑映射
$$
\mathcal{F}(y) = \Gamma(T, u^*\mathcal{G})
\xrightarrow{u_{small}^*\varphi}
\Gamma(T, u_{small}^*x^*\mathcal{H}|_{U_\etale})
\rightarrow
\Gamma(T, y^*\mathcal{H}|_{T_\etale}) = \mathcal{H}(y)
$$
，其中第二个箭头是层 $\mathcal{H}$ 的比较映射
(\ref{stacks-sheaves-equation-comparison-modules})。若满足引理中的余圈条件，
则这个指派与 $[U/_{\!p}R]$ 的态射所给出的层
% P11-E0121
$\mathcal{F}$ 和 $\mathcal{G}$ 的限制映射相容。证明从略。
提示：在命题 \ref{stacks-sheaves-proposition-quasi-coherent} 的证明中，
$\mathcal{F}$ 的限制映射已用 $(\mathcal{G}, \alpha)$ 明确写出。
\end{proof}







\section{代数叠上的拟凝聚层}
\label{stacks-sheaves-section-quasi-coherent-algebraic-stacks}

\noindent
设 $\mathcal{X}$ 是 $S$ 上的代数叠。由代数叠，引理
\ref{algebraic-lemma-stack-presentation}，可以找到等价
$[U/R] \to \mathcal{X}$，其中 $(U, R, s, t, c)$
是代数空间中的光滑群胚。于是
$$
\QCoh(\mathcal{O}_\mathcal{X})
\cong
\QCoh(\mathcal{O}_{[U/R]})
\cong
\QCoh(U, R, s, t, c)
$$
，其中第二个等价是命题 \ref{stacks-sheaves-proposition-quasi-coherent}。
因此，代数叠上的拟凝聚层范畴等价于代数空间中光滑群胚上的
拟凝聚模范畴。特别地，由空间中的群胚，引理
\ref{spaces-groupoids-lemma-abelian}，可见
$\QCoh(\mathcal{O}_\mathcal{X})$ 是 Abel 范畴！

\medskip\noindent
当前的设置有一点令人不安：全忠实嵌入
$$
\QCoh(\mathcal{O}_\mathcal{X})
\longrightarrow
\textit{Mod}(\mathcal{O}_\mathcal{X})
$$
一般{\bf 并不}正合。不过，对概形恰好也会发生同样的事：
对大多数概形 $X$，嵌入
$$
\QCoh(\mathcal{O}_X) \cong
\QCoh((\Sch/X)_{fppf}, \mathcal{O}_X) \longrightarrow
\textit{Mod}((\Sch/X)_{fppf}, \mathcal{O}_X)
$$
并不正合；参见下降，引理
\ref{descent-lemma-equivalence-quasi-coherent-limits}。顺便指出，
下降，引理 \ref{descent-lemma-equivalence-quasi-coherent-limits}
的证明中的例子表明，一般而言，严格满嵌入
$\QCoh(\mathcal{O}_\mathcal{X}) \to
\textit{LQCoh}(\mathcal{O}_\mathcal{X})$ 也不正合。

\medskip\noindent
把迄今得到的所有结果汇集成一个陈述。

\begin{lemma}
\label{stacks-sheaves-lemma-quasi-coherent-algebraic-stack}
设 $\mathcal{X}$ 是 $S$ 上的代数叠。
\begin{enumerate}
\item 若 $[U/R] \to \mathcal{X}$ 是 $\mathcal{X}$ 的一个表现，
则有典范等价
$\QCoh(\mathcal{O}_\mathcal{X}) \cong
\QCoh(U, R, s, t, c)$.
\item 范畴 $\QCoh(\mathcal{O}_\mathcal{X})$ 是 Abel 范畴。
\item 包含函子 $\QCoh(\mathcal{O}_\mathcal{X}) \to
\textit{Mod}(\mathcal{O}_\mathcal{X})$ 右正合，但一般{\bf 不}正合。
\item 范畴 $\QCoh(\mathcal{O}_\mathcal{X})$ 有余极限，
并且它们与范畴 $\textit{Mod}(\mathcal{O}_\mathcal{X})$ 中的余极限一致。
\item 给定 $\QCoh(\mathcal{O}_\mathcal{X})$ 中的
$\mathcal{F}, \mathcal{G}$，在 $\textit{Mod}(\mathcal{O}_\mathcal{X})$
中的张量积 $\mathcal{F} \otimes_{\mathcal{O}_\mathcal{X}} \mathcal{G}$
是 $\QCoh(\mathcal{O}_\mathcal{X})$ 的对象。
\item 给定 $\QCoh(\mathcal{O}_\mathcal{X})$ 中的
$\mathcal{F}, \mathcal{G}$，其中 $\mathcal{F}$ 有限局部自由，
则 $\textit{Mod}(\mathcal{O}_\mathcal{X})$ 中的层
$\SheafHom_{\mathcal{O}_\mathcal{X}}(\mathcal{F}, \mathcal{G})$
是 $\QCoh(\mathcal{O}_\mathcal{X})$ 的对象。
\item 给定 $\textit{Mod}(\mathcal{O}_\mathcal{X})$ 中的短正合列
$0 \to \mathcal{F}_1 \to \mathcal{F}_2 \to \mathcal{F}_3 \to 0$
，若 $\mathcal{F}_1$ 和 $\mathcal{F}_3$ 拟凝聚，则
$\mathcal{F}_2$ 拟凝聚。
\end{enumerate}
\end{lemma}

\begin{proof}
性质 (4)、(5) 和 (6) 已在引理 \ref{stacks-sheaves-lemma-qc-colimits} 中证明。
(1) 就是命题 \ref{stacks-sheaves-proposition-quasi-coherent}。(2) 由 (1) 和
空间中的群胚，引理 \ref{spaces-groupoids-lemma-abelian} 得到，
如上所述。(3) 中包含函子的右正合性由 (4) 得到；请与同调，
引理 \ref{homology-lemma-exact-functor} 比较。关于 (3) 中包含函子
不正合，参见下降，引理
\ref{descent-lemma-equivalence-quasi-coherent-limits}。为证明 (7)，
注意只须检验 $\mathcal{F}_2$ 到一个概形的大位点的限制拟凝聚
（引理 \ref{stacks-sheaves-lemma-characterize-quasi-coherent}），
故结论来自下降，引理
\ref{descent-lemma-equivalence-quasi-coherent-limits} 的相应部分。
\end{proof}

\noindent
下面构造代数叠上模的凝聚化子。

\begin{proposition}
\label{stacks-sheaves-proposition-coherator}
设 $\mathcal{X}$ 是 $S$ 上的代数叠。
\begin{enumerate}
\item 范畴 $\QCoh(\mathcal{O}_\mathcal{X})$ 是 Grothendieck Abel 范畴。
因此，$\QCoh(\mathcal{O}_\mathcal{X})$ 有足够多的内射对象和所有极限。
\item 包含函子
$\QCoh(\mathcal{O}_\mathcal{X}) \to
\textit{Mod}(\mathcal{O}_\mathcal{X})$ 有右伴随\footnote{这个函子有时称为
{\it 凝聚化子}。}
$$
Q :
\textit{Mod}(\mathcal{O}_\mathcal{X})
\longrightarrow
\QCoh(\mathcal{O}_\mathcal{X})
$$
，并且对每个拟凝聚层 $\mathcal{F}$，伴随映射
$Q(\mathcal{F}) \to \mathcal{F}$ 是同构。
\end{enumerate}
\end{proposition}

\begin{proof}
这个证明重复概形情形的证明（参见性质，命题
\ref{properties-proposition-coherator}）以及代数空间情形的证明
（参见空间的性质，命题
\ref{spaces-properties-proposition-coherator}）。建议读者先阅读其中一个证明。

\medskip\noindent
(1) 表示 $\QCoh(\mathcal{O}_\mathcal{X})$ 满足：(a) 有所有余极限；
(b) 滤过余极限正合；(c) 有一个生成元；参见内射对象，第
\ref{injectives-section-grothendieck-conditions} 节。由引理
\ref{stacks-sheaves-lemma-quasi-coherent-algebraic-stack}，
% P11-E0122
$\QCoh(\mathcal{O}_X)$ 中存在余极限，并且它们与
$\textit{Mod}(\mathcal{O}_X)$ 中的余极限一致。由位点上的模，引理
\ref{sites-modules-lemma-limits-colimits}，滤过余极限正合。
所以 (a) 和 (b) 成立。

\medskip\noindent
选取表现 $\mathcal{X} = [U/R]$，使 $(U, R, s, t, c)$
是代数空间中的光滑群胚；特别地，$s$ 和 $t$ 是代数空间的平坦态射。
由上面的引理 \ref{stacks-sheaves-lemma-quasi-coherent-algebraic-stack}，有
$\QCoh(\mathcal{O}_\mathcal{X}) = \QCoh(U, R, s, t, c)$.
由空间中的群胚，引理 \ref{spaces-groupoids-lemma-set-generators}，
存在集合 $T$ 和 $\mathcal{X}$ 上一族拟凝聚层
$(\mathcal{F}_t)_{t \in T}$，使 $\mathcal{X}$ 上每个拟凝聚层都是
其子层的有向余极限，而这些子层各自同构于某个 $\mathcal{F}_t$。
因此 $\bigoplus_t \mathcal{F}_t$ 是 $\QCoh(\mathcal{O}_X)$ 的生成元，
从而 (c) 成立。关于极限和内射对象的断言在任意 Grothendieck Abel
范畴中都成立；参见内射对象，定理
\ref{injectives-theorem-injective-embedding-grothendieck} 和引理
\ref{injectives-lemma-grothendieck-products}。

\medskip\noindent
证明 (2)。为构造 $Q$，采用如下的一般步骤。给定
$\textit{Mod}(\mathcal{O}_\mathcal{X})$ 的对象 $\mathcal{F}$，
考虑函子
$$
\QCoh(\mathcal{O}_\mathcal{X})^{opp}
\longrightarrow
\textit{Sets},
\quad
\mathcal{G}
\longmapsto
\Hom_\mathcal{X}(\mathcal{G}, \mathcal{F})
$$
这个函子把余极限变为极限，因而可表；参见内射对象，引理
\ref{injectives-lemma-grothendieck-brown}。因此，存在拟凝聚层
$Q(\mathcal{F})$ 以及函子性同构
$\Hom_\mathcal{X}(\mathcal{G}, \mathcal{F}) =
\Hom_\mathcal{X}(\mathcal{G}, Q(\mathcal{F}))$
，其中 $\mathcal{G}$ 属于 $\QCoh(\mathcal{O}_\mathcal{X})$。
由 Yoneda 引理（范畴，引理 \ref{categories-lemma-yoneda}），
构造 $\mathcal{F} \leadsto Q(\mathcal{F})$ 关于 $\mathcal{F}$ 是函子的。
由构造，$Q$ 是包含函子的右伴随。当 $\mathcal{F}$ 拟凝聚时，
$Q(\mathcal{F}) \to \mathcal{F}$ 是同构；这是包含函子
$\QCoh(\mathcal{O}_\mathcal{X}) \to
\textit{Mod}(\mathcal{O}_\mathcal{X})$
全忠实这一事实的形式推论。
\end{proof}








\section{上同调}
\label{stacks-sheaves-section-cohomology-general}

\noindent
设 $S$ 是概形，并设 $\mathcal{X}$ 是 $(\Sch/S)_{fppf}$ 上的
群胚纤维化范畴。对任意
% P11-E0123
$\tau \in \{Zariski,\linebreak[0] \etale,\linebreak[0] smooth,\linebreak[0] syntomic,\linebreak[0] fppf\}$，范畴
$\textit{Ab}(\mathcal{X}_\tau)$ 和
$\textit{Mod}(\mathcal{X}_\tau, \mathcal{O}_\mathcal{X})$
都有足够多的内射对象；参见内射对象，定理
\ref{injectives-theorem-sheaves-injectives} 和
\ref{injectives-theorem-sheaves-modules-injectives}。
因此，可以使用位点上的上同调，第
\ref{sites-cohomology-section-cohomology-sheaves} 节的机制定义上同调群
$$
H^p(\mathcal{X}_\tau, \mathcal{F}) = H^p_\tau(\mathcal{X}, \mathcal{F})
\quad\text{且}\quad
H^p(x, \mathcal{F}) = H^p_\tau(x, \mathcal{F})
$$
，其中 $x \in \Ob(\mathcal{X})$ 任意，而 $\mathcal{F}$ 是
$\textit{Ab}(\mathcal{X}_\tau)$ 或
$\textit{Mod}(\mathcal{X}_\tau, \mathcal{O}_\mathcal{X})$ 的任意对象。
此外，若 $f : \mathcal{X} \to \mathcal{Y}$ 是 $(\Sch/S)_{fppf}$ 上
群胚纤维化范畴的 $1$-态射，则在 $\textit{Ab}(\mathcal{Y}_\tau)$ 或
$\textit{Mod}(\mathcal{Y}_\tau, \mathcal{O}_\mathcal{Y})$ 中得到
高阶正像 $R^if_*\mathcal{F}$。当然，如位点上的上同调，第
\ref{sites-cohomology-section-derived-functors} 节所述，
$H^p(-)$ 和 $R^if_*$ 也有导出版本。

\begin{lemma}
\label{stacks-sheaves-lemma-cohomology-restriction}
设 $S$ 是概形。设 $\mathcal{X}$ 是 $(\Sch/S)_{fppf}$ 上的
群胚纤维化范畴。设 $\tau \in \{Zariski,\linebreak[0] \etale,\linebreak[0] smooth,
\linebreak[0] syntomic,\linebreak[0] fppf\}$。设 $x \in \Ob(\mathcal{X})$ 是位于概形 $U$
之上的对象。设 $\mathcal{F}$ 是 $\textit{Ab}(\mathcal{X}_\tau)$ 或
$\textit{Mod}(\mathcal{X}_\tau, \mathcal{O}_\mathcal{X})$ 的对象。则
$$
H^p_\tau(x, \mathcal{F}) = H^p((\Sch/U)_\tau, x^{-1}\mathcal{F})
$$
；若 $\tau = \etale$，还有
$$
H^p_\etale(x, \mathcal{F}) =
H^p(U_\etale, \mathcal{F}|_{U_\etale}).
$$
\end{lemma}

\begin{proof}
第一个陈述来自位点上的上同调，引理
\ref{sites-cohomology-lemma-cohomology-of-open} 以及引理
\ref{stacks-sheaves-lemma-localizing-structure-sheaf} 的等价。
第二个陈述由第一个陈述和 \'Etale 上同调，引理
\ref{etale-cohomology-lemma-compare-cohomology-big-small} 合并得到。
\end{proof}














\section{内射层}
\label{stacks-sheaves-section-lower-shriek}

\noindent
内射 Abel 层或内射模的推前仍是内射的。

\begin{lemma}
\label{stacks-sheaves-lemma-pushforward-injective}
设 $f : \mathcal{X} \to \mathcal{Y}$ 是 $(\Sch/S)_{fppf}$ 上
群胚纤维化范畴的 $1$-态射。设
$\tau \in \{Zar,\allowbreak \etale,\allowbreak smooth,\allowbreak syntomic,\allowbreak fppf\}$。
\begin{enumerate}
\item 若 $\mathcal{I}$ 在 $\textit{Ab}(\mathcal{X}_\tau)$ 中内射，
则 $f_*\mathcal{I}$ 在 $\textit{Ab}(\mathcal{Y}_\tau)$ 中内射；
\item 若 $\mathcal{I}$ 在
$\textit{Mod}(\mathcal{X}_\tau, \mathcal{O}_\mathcal{X})$ 中内射，
则 $f_*\mathcal{I}$ 在
$\textit{Mod}(\mathcal{Y}_\tau, \mathcal{O}_\mathcal{Y})$
中内射。
\end{enumerate}
\end{lemma}

\begin{proof}
这形式地来自如下事实：$f^{-1}$ 是 $f_*$ 的正合左伴随；
参见同调，引理 \ref{homology-lemma-adjoint-preserve-injectives}。
\end{proof}

\noindent
本节余下部分证明：在 Abel 层和模上，拉回 $f^{-1}$ 有左伴随 $f_!$。
若 $f$ 可表（由概形或代数空间表示），那么 $f_!$ 是正合的，
而 $f^{-1}$ 保持内射对象。先证明几个预备引理，它们涉及群胚纤维化
范畴中的纤维积和等化子，以及这些构造关于态射的行为。

\begin{lemma}
\label{stacks-sheaves-lemma-fibre-products}
设 $p : \mathcal{X} \to (\Sch/S)_{fppf}$ 是群胚纤维化范畴。
\begin{enumerate}
\item 范畴 $\mathcal{X}$ 有纤维积。
\item 若 $\mathcal{X}$ 的 $\mathit{Isom}$-预层可由代数空间表示，
则 $\mathcal{X}$ 有等化子。
\item 若 $\mathcal{X}$ 是代数叠（或更一般地，是商叠），
则 $\mathcal{X}$ 有等化子。
\end{enumerate}
\end{lemma}

\begin{proof}
% P11-E0124
由于 $(\Sch/S)_{fppf}$ 有纤维积，(1) 来自范畴，引理
\ref{categories-lemma-fibred-groupoids-fibre-product-goes-up}。

\medskip\noindent
设 $a, b : x \to y$ 是 $\mathcal{X}$ 的态射。置 $U = p(x)$ 和
$V = p(y)$。概形范畴有等化子，故可令 $W \to U$ 是 $p(a)$ 与
$p(b)$ 的等化子。% P11-E0125
把 $\mathcal{X}$ 中位于 $W \to U$ 之上的一个态射记为
$c : z \to x$。若 $a$ 和 $b$ 的等化子存在，它就是 $a \circ c$
与 $b \circ c$ 的等化子。因此，可以假设
$p(a) = p(b) = f : U \to V$。由于 $\mathcal{X}$ 群胚纤维化，
在 $\mathcal{X}$ 位于 $U$ 上的纤维范畴中，存在唯一自同构
$i : x \to x$，使得 $a \circ i = b$。同样，$a$ 和 $b$ 的等化子
就是 $\text{id}_x$ 和 $i$ 的等化子。回忆
$\mathit{Isom}_\mathcal{X}(x)$ 是 $(\Sch/U)_{fppf}$ 上的预层，
它把 $T/U$ 映到 $\mathcal{X}$ 位于 $T$ 上的纤维范畴中
$x|_T$ 的自同构集合；参见 Stacks，定义
\ref{stacks-definition-mor-presheaf}。若
$\mathit{Isom}_\mathcal{X}(x)$ 可由代数空间 $G \to U$ 表示，
则 $\text{id}_x$ 和 $i$ 定义 $U$ 上的态射 $e, i : U \to G$。
置 $M = U \times_{e, G, i} U$；由空间的态射，引理
\ref{spaces-morphisms-lemma-section-immersion}，它是概形。
于是显然，$x|_M \to x$ 是 $\mathcal{X}$ 中映射
$\text{id}_x$ 和 $i$ 的等化子。这证明了 (2)。

\medskip\noindent
若对于 $S$ 上代数空间中的某个群胚 $(U, R, s, t, c)$，有
$\mathcal{X} = [U/R]$，则由自举，引理
\ref{bootstrap-lemma-quotient-stack-isom}，(2) 的假设成立。
若 $\mathcal{X}$ 是代数叠，则由代数叠，引理
\ref{algebraic-lemma-stack-presentation}，可以选取表现
$[U/R] \cong \mathcal{X}$。
\end{proof}

\begin{lemma}
\label{stacks-sheaves-lemma-fibre-products-morphism}
设 $f : \mathcal{X} \to \mathcal{Y}$ 是 $(\Sch/S)_{fppf}$ 上
群胚纤维化范畴的 $1$-态射。
\begin{enumerate}
\item 函子 $f$ 把纤维积变为纤维积。
\item 若 $f$ 忠实，则 $f$ 把等化子变为等化子。
\end{enumerate}
\end{lemma}

\begin{proof}
由范畴，引理 \ref{categories-lemma-fibred-groupoids-fibre-product-goes-up}，
可见 $\mathcal{X}$ 中的纤维积，是位于 $(\Sch/S)_{fppf}$ 中一个
纤维积图之上的任意交换方块。对 $\mathcal{Y}$ 也同样。因此 (1) 清楚。

\medskip\noindent
设 $x \to x'$ 是 $\mathcal{X}$ 中两个态射 $a, b : x' \to x''$
的等化子。我们将证明 $f(x) \to f(x')$ 是 $f(a)$ 和 $f(b)$ 的等化子。
设 $y \to f(x')$ 是 $\mathcal{Y}$ 的一个态射，它等化 $f(a)$ 和 $f(b)$。
设 $x, x', x''$ 位于概形 $U, U', U''$ 之上，而 $y$ 位于 $V$ 之上。
% P11-E0126
把 $y \to f(x')$ 在概形范畴中的像记为 $h : V \to U'$。
由纤维化范畴的公理，态射 $y \to f(x')$ 同构于
$f(h^*x') \to f(x')$。因此，由于 $f$ 忠实，可见
$h^*x' \to x'$ 等化 $a$ 和 $b$。于是得到唯一态射
$h^*x' \to x$，其像 $y = f(h^*x') \to f(x)$ 正是
$\mathcal{Y}$ 中所需的态射。
\end{proof}

\begin{lemma}
\label{stacks-sheaves-lemma-fibre-products-preserve-properties}
设 $f : \mathcal{X} \to \mathcal{Y}$、$g : \mathcal{Z} \to \mathcal{Y}$
是 $(\Sch/S)_{fppf}$ 上群胚纤维化范畴的忠实 $1$-态射。
\begin{enumerate}
\item 函子 $\mathcal{X} \times_\mathcal{Y} \mathcal{Z} \to \mathcal{Y}$
忠实；
\item 若 $\mathcal{X}, \mathcal{Z}$ 有等化子，则
$\mathcal{X} \times_\mathcal{Y} \mathcal{Z}$ 也有等化子。
\end{enumerate}
\end{lemma}

\begin{proof}
把 $\mathcal{X} \times_\mathcal{Y} \mathcal{Z}$ 的对象看作四元组
$(U, x, z, \alpha)$，其中 $\alpha : f(x) \to g(z)$ 是 $U$ 上的同构；
参见范畴，引理 \ref{categories-lemma-2-product-categories-over-C}。
态射 $(U, x, z, \alpha) \to (U', x', z', \alpha')$
是一对与 $\alpha$ 和 $\alpha'$ 相容的态射 $a : x \to x'$ 和
$b : z \to z'$。因此显然，若 $f$ 和 $g$ 忠实，则函子
$\mathcal{X} \times_\mathcal{Y} \mathcal{Z} \to \mathcal{Y}$ 也忠实。
现在，假设
$(a, b), (a', b') : (U, x, z, \alpha) \to (U', x', z', \alpha')$
是该 $2$-纤维积的两个态射。考虑 $a$ 与 $a'$ 的等化子
$x'' \to x$，以及 $b$ 与 $b'$ 的等化子 $z'' \to z$。
由于 $f$ 与等化子交换（引理 \ref{stacks-sheaves-lemma-fibre-products-morphism}），
$f(x'') \to f(x)$ 是 $f(a)$ 和 $f(a')$ 的等化子。类似地，
$g(z'') \to g(z)$ 是 $g(b)$ 和 $g(b')$ 的等化子。考虑下图：
$$
\xymatrix{
f(x'') \ar[r] \ar@{..>}[d]_{\alpha''}&
f(x) \ar[d]_\alpha
\ar@<0.5ex>[r]^{f(a)}
\ar@<-0.5ex>[r]_{f(a')}
 &
f(x') \ar[d]^{\alpha'} \\
g(z'') \ar[r] &
g(z)
\ar@<0.5ex>[r]^{g(b)}
\ar@<-0.5ex>[r]_{g(b')}
 &
g(z')
}
$$
虚线箭头显然存在并且是同构。然而，$\alpha''$ 在概形范畴中的像
先验地未必是其源的恒等态射。另一方面，$\alpha''$ 的存在意味着，
可以假设 $x''$ 和 $z''$ 定义在同一概形上，并且态射
$x'' \to x$ 和 $z'' \to z$ 在概形范畴中有相同的像。
重新作上图，可见虚线箭头现在确实投影到恒等态射，结论得证。
若干细节从略。
\end{proof}

\noindent
由于我们处理大位点，便有下面这个多少有悖直觉的结果
（它对概形的大位点之间的态射也成立）。警告：若去掉 $f$ 忠实的假设，
这个结果就不成立。

\begin{lemma}
\label{stacks-sheaves-lemma-pullback-injective}
设 $f : \mathcal{X} \to \mathcal{Y}$ 是 $(\Sch/S)_{fppf}$ 上
群胚纤维化范畴的 $1$-态射。设
$\tau \in \{Zar,\allowbreak \etale,\allowbreak smooth,\allowbreak syntomic,\allowbreak fppf\}$。函子
$f^{-1} : \textit{Ab}(\mathcal{Y}_\tau) \to \textit{Ab}(\mathcal{X}_\tau)$
有左伴随
$f_! : \textit{Ab}(\mathcal{X}_\tau) \to \textit{Ab}(\mathcal{Y}_\tau)$.
若 $f$ 忠实且 $\mathcal{X}$ 有等化子，则
\begin{enumerate}
\item $f_!$ 正合；
\item 若 $\mathcal{I}$ 在 $\textit{Ab}(\mathcal{Y}_\tau)$ 中内射，
则 $f^{-1}\mathcal{I}$ 在 $\textit{Ab}(\mathcal{X}_\tau)$ 中内射。
\end{enumerate}
\end{lemma}

\begin{proof}
由 Stacks，引理 \ref{stacks-lemma-topology-inherited-functorial}，
函子 $f$ 连续且余连续。因此，由位点上的模，引理
\ref{sites-modules-lemma-g-shriek-adjoint}，函子
$f^{-1} : \textit{Ab}(\mathcal{Y}_\tau) \to \textit{Ab}(\mathcal{X}_\tau)$
有左伴随
$f_! : \textit{Ab}(\mathcal{X}_\tau) \to \textit{Ab}(\mathcal{Y}_\tau)$.
为证明 (1)，应用位点上的模，引理
\ref{sites-modules-lemma-exactness-lower-shriek}；为检验该引理的假设，
使用上面的引理 \ref{stacks-sheaves-lemma-fibre-products} 和
\ref{stacks-sheaves-lemma-fibre-products-morphism}。(2) 形式地由此得到；参见同调，
引理 \ref{homology-lemma-adjoint-preserve-injectives}。
\end{proof}

\begin{lemma}
\label{stacks-sheaves-lemma-pullback-injective-modules}
设 $f : \mathcal{X} \to \mathcal{Y}$ 是 $(\Sch/S)_{fppf}$ 上
群胚纤维化范畴的 $1$-态射。设
$\tau \in \{Zar,\allowbreak \etale,\allowbreak smooth,\allowbreak syntomic,\allowbreak fppf\}$。函子
$f^* : \textit{Mod}(\mathcal{Y}_\tau, \mathcal{O}_\mathcal{Y}) \to
\textit{Mod}(\mathcal{X}_\tau, \mathcal{O}_\mathcal{X})$
有左伴随
$f_! : \textit{Mod}(\mathcal{X}_\tau, \mathcal{O}_\mathcal{X}) \to
\textit{Mod}(\mathcal{Y}_\tau, \mathcal{O}_\mathcal{Y})$；
在底层 Abel 层上，它与引理 \ref{stacks-sheaves-lemma-pullback-injective} 的函子
$f_!$ 一致。若 $f$ 忠实且 $\mathcal{X}$ 有等化子，则
\begin{enumerate}
\item $f_!$ 正合；
\item 若 $\mathcal{I}$ 在
% P11-E0127
$\textit{Mod}(\mathcal{Y}_\tau, \mathcal{O}_\mathcal{X})$ 中内射，
则 $f^{-1}\mathcal{I}$ 在
$\textit{Mod}(\mathcal{X}_\tau, \mathcal{O}_\mathcal{X})$
中内射。
\end{enumerate}
\end{lemma}

\begin{proof}
回忆 $f$ 是位点的连续且余连续函子，并且
$f^{-1}\mathcal{O}_\mathcal{Y} = \mathcal{O}_\mathcal{X}$。因此，
位点上的模，引理 \ref{sites-modules-lemma-lower-shriek-modules}
蕴含 $f^*$ 有左伴随 $f_!^{Mod}$。设 $x$ 是 $\mathcal{X}$ 的对象，
位于概形 $U$ 之上。则 $f$ 诱导带环位点的等价
$$
\mathcal{X}/x \longrightarrow \mathcal{Y}/f(x)
$$
，因为两边都等价于 $(\Sch/U)_\tau$；参见引理
\ref{stacks-sheaves-lemma-localizing-structure-sheaf}。位点上的模，注
\ref{sites-modules-remark-when-shriek-equal} 表明，
% P11-E0128
$f_!$ 与 Abel 层上的函子一致。

\medskip\noindent
现在假设 $\mathcal{X}$ 有等化子且 $f$ 忠实。引理
\ref{stacks-sheaves-lemma-pullback-injective} 告诉我们 $f_!$ 正合。最后，同调，引理
\ref{homology-lemma-adjoint-preserve-injectives} 蕴含关于内射模拉回的陈述。
\end{proof}




\section{{\v C}ech 复形}
\label{stacks-sheaves-section-cech}

\noindent
为了计算代数叠上一个层的上同调，我们把它与该层限制到这个代数叠的
覆盖后所得的上同调作比较。

\medskip\noindent
本节始终处于如下情形。给定群胚纤维化范畴的一个 $1$-态射
\begin{equation}
\label{stacks-sheaves-equation-covering}
\vcenter{
\xymatrix{
\mathcal{U} \ar[rr]_f \ar[rd]_q & &  \mathcal{X} \ar[ld]^p \\
& (\Sch/S)_{fppf}
}
}
\end{equation}
我们将把 $\mathcal{U}$ 看作 $\mathcal{X}$ 的一个“覆盖”。
因此，要考虑单纯对象
$$
\xymatrix{
\mathcal{U} \times_\mathcal{X} \mathcal{U} \times_\mathcal{X} \mathcal{U}
\ar@<1ex>[r]
\ar@<0ex>[r]
\ar@<-1ex>[r]
&
\mathcal{U} \times_\mathcal{X} \mathcal{U}
\ar@<0.5ex>[r]
\ar@<-0.5ex>[r]
&
\mathcal{U}
}
$$
，它位于 $(\Sch/S)_{fppf}$ 上群胚纤维化范畴的范畴中。
不过，由于这是 $(2, 1)$-范畴而不是范畴，必须明确说明其含义。
具体而言，令 $\mathcal{U}_n$ 为如下范畴，其对象为
$(u_0, \ldots, u_n, x, \alpha_0, \ldots, \alpha_n)$
，其中 $\alpha_i : f(u_i) \to x$ 是 $\mathcal{X}$ 中的同构。
把如下 $1$-态射记为 $f_n : \mathcal{U}_n \to \mathcal{X}$：
它把 $(u_0, \ldots, u_n, x, \alpha_0, \ldots, \alpha_n)$
映到对象 $x$。注意 $\mathcal{U}_0 = \mathcal{U}$ 且 $f_0 = f$。
给定映射 $\varphi : [m] \to [n]$，考虑 $1$-态射
% P11-E0129
$\mathcal{U}_\varphi : \mathcal{U}_n \longrightarrow \mathcal{U}_n$
，它在对象上由下式给出：
$$
(u_0, \ldots, u_n, x, \alpha_0, \ldots, \alpha_n)
\longmapsto
(u_{\varphi(0)}, \ldots, u_{\varphi(m)}, x,
\alpha_{\varphi(0)}, \ldots, \alpha_{\varphi(m)})
$$
所有这些 $1$-态射都严格正确地复合（无需 $2$-态射），而且所有这些
$1$-态射都是 $\mathcal{X}$ 上的 $1$-态射。把这个单纯对象记为
% P11-E0130
$\mathcal{U}_\bullet$。若 $\mathcal{F}$ 是 $\mathcal{X}$ 上的集合预层，
则得到余单纯集合
$$
\xymatrix{
\Gamma(\mathcal{U}_0, f_0^{-1}\mathcal{F})
\ar@<0.5ex>[r]
\ar@<-0.5ex>[r]
&
\Gamma(\mathcal{U}_1, f_1^{-1}\mathcal{F})
\ar@<1ex>[r]
\ar@<0ex>[r]
\ar@<-1ex>[r]
&
\Gamma(\mathcal{U}_2, f_2^{-1}\mathcal{F})
}
$$
这里的箭头是沿该单纯对象的相应态射所作的拉回映射。
若 $\mathcal{F}$ 是 Abel 群预层，这就是余单纯 Abel 群。

\medskip\noindent
设 $\mathcal{U} \to \mathcal{X}$ 如上，并设 $\mathcal{F}$ 是
$\mathcal{X}$ 上的 Abel 预层。与此情形相联系的{\it {\v C}ech 复形}
记为 $\check{\mathcal{C}}^\bullet(\mathcal{U} \to \mathcal{X},
\mathcal{F})$。它是与上面的余单纯 Abel 群相联系的上链复形；
参见单纯方法，第 \ref{simplicial-section-dold-kan-cosimplicial} 节。
其各项为
$$
\check{\mathcal{C}}^n(\mathcal{U} \to \mathcal{X}, \mathcal{F})
= \Gamma(\mathcal{U}_n, f_n^{-1}\mathcal{F}).
$$
边界映射为
$$
d^n = \sum\nolimits_{i = 0}^{n + 1} (-1)^i \delta^{n + 1}_i :
\Gamma(\mathcal{U}_n, f_n^{-1}\mathcal{F})
\longrightarrow
\Gamma(\mathcal{U}_{n + 1}, f_{n + 1}^{-1}\mathcal{F})
$$
，其中 $\delta^{n + 1}_i$ 对应于省略指标 $i$ 的映射
$[n] \to [n + 1]$。注意，映射
% P11-E0131
$\Gamma(\mathcal{X}, \mathcal{F}) \to
\Gamma(\mathcal{U}_0, f_0^{-1}\mathcal{F}_0)$
位于微分 $d^0$ 的核中。因此，把{\it 扩充 {\v C}ech 复形}定义为复形
$$
\ldots \to 0 \to
\Gamma(\mathcal{X}, \mathcal{F}) \to
\Gamma(\mathcal{U}_0, f_0^{-1}\mathcal{F}_0) \to
\Gamma(\mathcal{U}_1, f_1^{-1}\mathcal{F}_1) \to \ldots
$$
，其中 $\Gamma(\mathcal{X}, \mathcal{F})$ 置于次数 $-1$。
扩充 {\v C}ech 复形无环，当且仅当典范映射
$$
\Gamma(\mathcal{X}, \mathcal{F})[0]
\longrightarrow
\check{\mathcal{C}}^\bullet(\mathcal{U} \to \mathcal{X}, \mathcal{F})
$$
是复形的拟同构。

\begin{lemma}
\label{stacks-sheaves-lemma-generalities}
{\v C}ech 复形的一般性质。
\begin{enumerate}
\item 若
$$
\xymatrix{
\mathcal{V} \ar[d]_g \ar[r]_h & \mathcal{U} \ar[d]^f \\
\mathcal{Y} \ar[r]^e & \mathcal{X}
}
$$
% P11-E0132
是 $(\Sch/S)_{fppf}$ 上群胚纤维化范畴的 $2$-交换图，
则有 {\v C}ech 复形的态射
$$
\check{\mathcal{C}}^\bullet(\mathcal{U} \to \mathcal{X}, \mathcal{F})
\longrightarrow
\check{\mathcal{C}}^\bullet(\mathcal{V} \to \mathcal{Y}, e^{-1}\mathcal{F})
$$
\item 若 $h$ 和 $e$ 是等价，则 (1) 的映射是同构；
\item 若 $f, f' : \mathcal{U} \to \mathcal{X}$ 是 $2$-同构的，
则相联系的 {\v C}ech 复形同构。
\end{enumerate}
\end{lemma}

\begin{proof}
在 (1) 的情形下，设 $t : f \circ h \to e \circ g$ 是 $2$-态射。
复形之间的映射在次数 $n$ 处，由沿 $1$-态射
$\mathcal{V}_n \to \mathcal{U}_n$ 的拉回给出；该态射由规则
$$
(v_0, \ldots, v_n, y, \beta_0, \ldots, \beta_n)
\longmapsto
(h(v_0), \ldots, h(v_n), e(y),
e(\beta_0) \circ t_{v_0}, \ldots, e(\beta_n) \circ t_{v_n}).
$$
给出。对于 (2)，注意当拉回沿范畴的等价进行时，对任意集合预层，
整体截面上的拉回都是同构。把 (1) 和 (2) 合并即得 (3)。
\end{proof}

\begin{lemma}
\label{stacks-sheaves-lemma-homotopy}
若存在 $1$-态射 $s : \mathcal{X} \to \mathcal{U}$，使得
$f \circ s$ 与 $\text{id}_\mathcal{X}$ $2$-同构，
则扩充 {\v C}ech 复形零伦。
\end{lemma}

\begin{proof}
置 $\mathcal{U}' = \mathcal{U} \times_\mathcal{X} \mathcal{X}$，
它是范畴，引理 \ref{categories-lemma-2-product-categories-over-C}
所述的纤维积。令 $f' : \mathcal{U}' \to \mathcal{X}$ 为第二投影。
% P11-E0133
则 $\mathcal{U} \to \mathcal{U}'$，$u \mapsto (u, f(x), 1)$，
是 $\mathcal{X}$ 上的等价，故由引理 \ref{stacks-sheaves-lemma-generalities}，
可以用 $(\mathcal{U}', f')$ 代替 $(\mathcal{U}, f)$。
这样做的好处是，现在 $f'$ 有一个截面 $s'$，严格满足
$f' \circ s' = \text{id}_\mathcal{X}$。具体而言，若
% P11-E0134
$t : s \circ f \to \text{id}_\mathcal{X}$ 是 $2$-同构，
则可以置 $s'(x) = (s(x), x, t_x)$。因此，可以假设
$f \circ s = \text{id}_\mathcal{X}$。

\medskip\noindent
在 $f \circ s = \text{id}_\mathcal{X}$ 的情形下，结论来自一般原理。
下面明确给出这个伦。具体而言，对 $n \geq 0$，把
$s_n : \mathcal{U}_n \to \mathcal{U}_{n + 1}$ 定义为如下 $1$-态射，
它在对象上由规则
$$
(u_0, \ldots, u_n, x, \alpha_0, \ldots, \alpha_n)
\longmapsto
(u_0, \ldots, u_n, s(x), x,
\alpha_0, \ldots, \alpha_n, \text{id}_x).
$$
给出。定义
$$
h^{n + 1} :
\Gamma(\mathcal{U}_{n + 1}, f_{n + 1}^{-1}\mathcal{F})
\longrightarrow
\Gamma(\mathcal{U}_n, f_n^{-1}\mathcal{F})
$$
为沿 $s_n$ 的拉回。还置 $s_{-1} = s$，并令
$h^0 : \Gamma(\mathcal{U}_0, f_0^{-1}\mathcal{F}) \to
\Gamma(\mathcal{X}, \mathcal{F})$ 等于沿 $s_{-1}$ 的拉回。
于是，映射族 $\{h^n\}_{n \geq 0}$ 是扩充 {\v C}ech 复形上
$1$ 与 $0$ 之间的伦。
\end{proof}









\section{相对 {\v C}ech 复形}
\label{stacks-sheaves-section-sheaf-cech-complex}

\noindent
如 (\ref{stacks-sheaves-equation-covering}) 所示，设
$f : \mathcal{U} \to \mathcal{X}$ 是 $(\Sch/S)_{fppf}$ 上
群胚纤维化范畴的一个 $1$-态射。考虑相联系的单纯对象
$\mathcal{U}_\bullet$ 以及映射
$f_n : \mathcal{U}_n \to \mathcal{X}$。令
$\tau \in \{Zar,\allowbreak \etale,\allowbreak smooth,\allowbreak syntomic,\allowbreak fppf\}$.
最后，设 $\mathcal{F}$ 是 $\mathcal{X}_\tau$ 上的（集合）层。则
$$
\xymatrix{
f_{0, *}f_0^{-1}\mathcal{F}
\ar@<0.5ex>[r]
\ar@<-0.5ex>[r]
&
f_{1, *}f_1^{-1}\mathcal{F}
\ar@<1ex>[r]
\ar@<0ex>[r]
\ar@<-1ex>[r]
&
f_{2, *}f_2^{-1}\mathcal{F}
}
$$
是 $\mathcal{X}_\tau$ 上的余单纯层，其中使用位点，第
\ref{sites-section-pullback} 节引入的拉回映射。
若 $\mathcal{F}$ 是 Abel 层，则 $f_{n, *}f_n^{-1}\mathcal{F}$
构成 $\mathcal{X}_\tau$ 上的余单纯 Abel 层。
相联系的复形（参见单纯方法，第
\ref{simplicial-section-dold-kan-cosimplicial} 节）
$$
\ldots \to 0 \to
f_{0, *}f_0^{-1}\mathcal{F} \to
f_{1, *}f_1^{-1}\mathcal{F} \to
f_{2, *}f_2^{-1}\mathcal{F} \to \ldots
$$
称为与此情形相联系的{\it 相对 {\v C}ech 复形}。
把这个复形记为 $\mathcal{K}^\bullet(f, \mathcal{F})$。
{\it 扩充相对 {\v C}ech 复形}是复形
$$
\ldots \to 0 \to
\mathcal{F} \to
f_{0, *}f_0^{-1}\mathcal{F} \to
f_{1, *}f_1^{-1}\mathcal{F} \to
f_{2, *}f_2^{-1}\mathcal{F} \to \ldots
$$
，其中 $\mathcal{F}$ 位于次数 $-1$。扩充相对 {\v C}ech 复形无环，
当且仅当映射
$\mathcal{F}[0] \to \mathcal{K}^\bullet(f, \mathcal{F})$
是层复形的拟同构。

\begin{remark}
\label{stacks-sheaves-remark-cech-complex-presheaves}
当 $\mathcal{F}$ 是预层时，也可以定义复形
$\mathcal{K}^\bullet(f, \mathcal{F})$，只是不能引用位点，第
\ref{sites-section-pullback} 节来定义拉回映射。为说明这些拉回映射，
设给定交换图
$$
\xymatrix{
\mathcal{V} \ar[rd]_g \ar[rr]_h & &  \mathcal{U} \ar[ld]^f \\
& \mathcal{X}
}
$$
，其中各范畴都在 $(\Sch/S)_{fppf}$ 上纤维化为群胚；再设
$\mathcal{G}$ 是 $\mathcal{U}$ 上的预层。可以把拉回映射
$f_*\mathcal{G} \to g_*h^{-1}\mathcal{G}$ 定义为复合
$$
f_*\mathcal{G} \longrightarrow
f_*h_*h^{-1}\mathcal{G} = g_*h^{-1}\mathcal{G}
$$
，其中的映射来自伴随映射
$\mathcal{G} \to h_*h^{-1}\mathcal{G}$。这是可行的，因为在当前情形下，
函子 $h_*$ 与 $h^{-1}$ 在预层上互为伴随（并与它们在层上的对应函子一致）。
参见第 \ref{stacks-sheaves-section-presheaves} 节和第 \ref{stacks-sheaves-section-sheaves} 节。
\end{remark}

\begin{lemma}
\label{stacks-sheaves-lemma-generalities-sheafified}
相对 {\v C}ech 复形具有如下的一般性质。
\begin{enumerate}
\item 若
$$
\xymatrix{
\mathcal{V} \ar[d]_g \ar[r]_h & \mathcal{U} \ar[d]^f \\
\mathcal{Y} \ar[r]^e & \mathcal{X}
}
$$
% P11-E0132
是 $(\Sch/S)_{fppf}$ 上群胚纤维化范畴的 $2$-交换图，则存在态射
$e^{-1}\mathcal{K}^\bullet(f, \mathcal{F}) \to
\mathcal{K}^\bullet(g, e^{-1}\mathcal{F})$.
\item 若 $h$ 与 $e$ 是等价，则 (1) 中的映射是同构；
\item 若 $f, f' : \mathcal{U} \to \mathcal{X}$ 是 $2$-同构的，
则相联系的相对 {\v C}ech 复形同构。
\end{enumerate}
\end{lemma}

\begin{proof}
逐字采用引理 \ref{stacks-sheaves-lemma-generalities} 的证明，
并使用注 \ref{stacks-sheaves-remark-cech-complex-presheaves} 中的拉回映射。
\end{proof}

\begin{lemma}
\label{stacks-sheaves-lemma-homotopy-sheafified}
若存在 $1$-态射 $s : \mathcal{X} \to \mathcal{U}$，使得
$f \circ s$ 与 $\text{id}_\mathcal{X}$ 是 $2$-同构的，
则扩充相对 {\v C}ech 复形零伦。
\end{lemma}

\begin{proof}
逐字采用引理 \ref{stacks-sheaves-lemma-homotopy} 的证明。
\end{proof}

\begin{remark}
\label{stacks-sheaves-remark-cech-complex-sections}
来“计算”相对 {\v C}ech 复形在 $\mathcal{X}$ 的一个对象 $x$ 上的值。
设 $p(x) = U$。考虑 $2$-纤维积图（借此引入记号
$g : \mathcal{V} \to \mathcal{Y}$）
$$
\xymatrix{
\mathcal{V} \ar@{=}[r] \ar[d]_g &
(\Sch/U)_{fppf} \times_{x, \mathcal{X}} \mathcal{U} \ar[r] \ar[d] &
\mathcal{U} \ar[d]^f \\
\mathcal{Y} \ar@{=}[r] &
(\Sch/U)_{fppf} \ar[r]^-x & \mathcal{X}
}
$$
注意，引理 \ref{stacks-sheaves-lemma-generalities} 的证明中的态射
$\mathcal{V}_n \to \mathcal{U}_n$ 诱导等价
$\mathcal{V}_n =
(\Sch/U)_{fppf} \times_{x, \mathcal{X}} \mathcal{U}_n$.
因此，由 (\ref{stacks-sheaves-equation-pushforward}) 可见
$$
\Gamma(x, \mathcal{K}^\bullet(f, \mathcal{F})) =
\check{\mathcal{C}}^\bullet(\mathcal{V} \to \mathcal{Y}, x^{-1}\mathcal{F})
$$
换言之，相对 {\v C}ech 复形在 $\mathcal{X}$ 的对象 $x$ 上的值，
就是把 $f$ 基变换到 $\mathcal{X}/x \cong (\Sch/U)_{fppf}$ 后所得的
{\v C}ech 复形。例如，这说明引理 \ref{stacks-sheaves-lemma-homotopy} 蕴涵
引理 \ref{stacks-sheaves-lemma-homotopy-sheafified}；更一般地，关于（通常的）
{\v C}ech 复形的结果蕴涵相对 {\v C}ech 复形的相应结果。
\end{remark}

\begin{lemma}
\label{stacks-sheaves-lemma-base-change-cech-complex}
设
$$
\xymatrix{
\mathcal{V} \ar[d]_g \ar[r]_h & \mathcal{U} \ar[d]^f \\
\mathcal{Y} \ar[r]^e & \mathcal{X}
}
$$
是 $(\Sch/S)_{fppf}$ 上群胚纤维化范畴的 $2$-纤维积，且
$\mathcal{F}$ 是 $\mathcal{X}$ 上的 Abel 预层。则映射
$e^{-1}\mathcal{K}^\bullet(f, \mathcal{F}) \to
\mathcal{K}^\bullet(g, e^{-1}\mathcal{F})$
（见引理 \ref{stacks-sheaves-lemma-generalities-sheafified}）是 Abel 预层复形的同构。
\end{lemma}

\begin{proof}
设 $y$ 是 $\mathcal{Y}$ 中位于概形 $T$ 上的对象。置 $x = e(y)$。
下面证明该映射在 $y$ 上的截面上诱导同构。注意
$$
\Gamma(y, e^{-1}\mathcal{K}^\bullet(f, \mathcal{F})) =
\Gamma(x, \mathcal{K}^\bullet(f, \mathcal{F})) =
\check{\mathcal{C}}^\bullet(
(\Sch/T)_{fppf} \times_{x, \mathcal{X}} \mathcal{U} \to
(\Sch/T)_{fppf}, x^{-1}\mathcal{F})
$$
，这是由注 \ref{stacks-sheaves-remark-cech-complex-sections} 得到的。另一方面，
$$
\Gamma(y, \mathcal{K}^\bullet(g, e^{-1}\mathcal{F})) =
\check{\mathcal{C}}^\bullet(
(\Sch/T)_{fppf} \times_{y, \mathcal{Y}} \mathcal{V} \to
(\Sch/T)_{fppf}, y^{-1}e^{-1}\mathcal{F})
$$
仍由注 \ref{stacks-sheaves-remark-cech-complex-sections} 得到该式。
注意 $y^{-1}e^{-1}\mathcal{F} = x^{-1}\mathcal{F}$；由于该图是
$2$-Cartesian 的，$1$-态射
$$
(\Sch/T)_{fppf} \times_{y, \mathcal{Y}} \mathcal{V} \to
(\Sch/T)_{fppf} \times_{x, \mathcal{X}} \mathcal{U}
$$
是等价。因此，由引理 \ref{stacks-sheaves-lemma-generalities}，该映射在 $y$ 上的
截面上是同构。
\end{proof}

\noindent
正合性可以在一个“覆盖”上检验。

\begin{lemma}
\label{stacks-sheaves-lemma-check-exactness-covering}
设 $f : \mathcal{U} \to \mathcal{X}$ 是 $(\Sch/S)_{fppf}$ 上
群胚纤维化范畴的一个 $1$-态射。令
$\tau \in \{Zar,\allowbreak \etale,\allowbreak smooth,\allowbreak syntomic,\allowbreak fppf\}$.
设
$$
\mathcal{F} \to \mathcal{G} \to \mathcal{H}
$$
是 $\textit{Ab}(\mathcal{X}_\tau)$ 中的复形。假设
\begin{enumerate}
\item 对 $\mathcal{X}$ 的每个对象 $x$，都存在 $\mathcal{X}_\tau$ 中的覆盖
$\{x_i \to x\}$，使每个 $x_i$ 都同构于 $\mathcal{U}$ 中某个对象
$u_i$ 的像 $f(u_i)$；
\item $f^{-1}\mathcal{F} \to f^{-1}\mathcal{G} \to f^{-1}\mathcal{H}$ 正合。
\end{enumerate}
则序列 $\mathcal{F} \to \mathcal{G} \to \mathcal{H}$ 正合。
\end{lemma}

\begin{proof}
设 $x$ 是 $\mathcal{X}$ 中位于概形 $T$ 上的对象。考虑
$x^{-1}\mathcal{F} \to x^{-1}\mathcal{G} \to x^{-1}\mathcal{H}$
，这是 $(\Sch/T)_\tau$ 上的 Abel 层序列。只需证明这个序列正合。
由假设，存在一个 $\tau$-覆盖 $\{T_i \to T\}$，使得对位于 $T_i$ 上的
$\mathcal{U}$ 的某个对象 $u_i$，$x|_{T_i}$ 同构于 $f(u_i)$，而且
$u_i^{-1}f^{-1}\mathcal{F} \to u_i^{-1}f^{-1}\mathcal{G} \to
u_i^{-1}f^{-1}\mathcal{H}$ 这个 $(\Sch/T_i)_\tau$ 上的 Abel 层序列正合。
由于
$u_i^{-1}f^{-1}\mathcal{F} = x^{-1}\mathcal{F}|_{(\Sch/T_i)_\tau}$
，可知序列
$x^{-1}\mathcal{F} \to x^{-1}\mathcal{G} \to x^{-1}\mathcal{H}$
% P11-E0135
在覆盖的每个成员处局部化后变得正合，故该序列正合。
\end{proof}

\begin{proposition}
\label{stacks-sheaves-proposition-exactness-cech-complex}
设 $f : \mathcal{U} \to \mathcal{X}$ 是 $(\Sch/S)_{fppf}$ 上
群胚纤维化范畴的一个 $1$-态射。令
$\tau \in \{Zar,\allowbreak \etale,\allowbreak smooth,\allowbreak syntomic,\allowbreak fppf\}$.
若
\begin{enumerate}
\item $\mathcal{F}$ 是 $\mathcal{X}_\tau$ 上的 Abel 层；
\item 对 $\mathcal{X}$ 的每个对象 $x$，都存在 $\mathcal{X}_\tau$ 中的覆盖
$\{x_i \to x\}$，使每个 $x_i$ 都同构于 $\mathcal{U}$ 中某个对象
$u_i$ 的像 $f(u_i)$。
\end{enumerate}
则扩充相对 {\v C}ech 复形
$$
\ldots \to 0 \to
\mathcal{F} \to
f_{0, *}f_0^{-1}\mathcal{F} \to
f_{1, *}f_1^{-1}\mathcal{F} \to
f_{2, *}f_2^{-1}\mathcal{F} \to \ldots
$$
在 $\textit{Ab}(\mathcal{X}_\tau)$ 中正合。
\end{proposition}

\begin{proof}
由引理 \ref{stacks-sheaves-lemma-check-exactness-covering}，只需在拉回到 $\mathcal{U}$ 后
检验正合性。由引理 \ref{stacks-sheaves-lemma-base-change-cech-complex}，
扩充相对 {\v C}ech 复形的拉回同构于态射
$\mathcal{U} \times_\mathcal{X} \mathcal{U} \to \mathcal{U}$
% P11-E0136
及 $\mathcal{U}_\tau$ 上一个 Abel 层的扩充相对 {\v C}ech 复形。
由于存在截面
$\Delta_{\mathcal{U}/\mathcal{X}} : \mathcal{U} \to
\mathcal{U} \times_\mathcal{X} \mathcal{U}$，正合性由
引理 \ref{stacks-sheaves-lemma-homotopy-sheafified} 得出。
\end{proof}

\noindent
利用这一点，可以如下构造 {\v C}ech 到上同调谱序列。先给出一个技术性的
精确版本；下一节将给出一个仅适用于代数叠的版本。

\begin{lemma}
\label{stacks-sheaves-lemma-cech-to-cohomology}
设 $f : \mathcal{U} \to \mathcal{X}$ 是 $(\Sch/S)_{fppf}$ 上
群胚纤维化范畴的一个 $1$-态射。令
$\tau \in \{Zar,\allowbreak \etale,\allowbreak smooth,\allowbreak syntomic,\allowbreak fppf\}$.
假设
\begin{enumerate}
\item $\mathcal{F}$ 是 $\mathcal{X}_\tau$ 上的 Abel 层；
\item 对 $\mathcal{X}$ 的每个对象 $x$，都存在 $\mathcal{X}_\tau$ 中的覆盖
$\{x_i \to x\}$，使每个 $x_i$ 都同构于 $\mathcal{U}$ 中某个对象
$u_i$ 的像 $f(u_i)$；
\item 范畴 $\mathcal{U}$ 有等化子；
\item 函子 $f$ 忠实。
\end{enumerate}
则存在 Abel 群的第一象限谱序列
$$
E_1^{p, q} = H^q((\mathcal{U}_p)_\tau, f_p^{-1}\mathcal{F})
\Rightarrow
H^{p + q}(\mathcal{X}_\tau, \mathcal{F})
$$
，它收敛到 $\mathcal{F}$ 在 $\tau$-拓扑中的上同调。
\end{lemma}

\begin{proof}
证明开始前先作若干说明。由引理
\ref{stacks-sheaves-lemma-fibre-products-preserve-properties}（并作归纳），
所有群胚纤维化范畴 $\mathcal{U}_p$ 都有等化子，且所有态射
$f_p : \mathcal{U}_p \to \mathcal{X}$ 都忠实。设 $\mathcal{I}$ 是
$\textit{Ab}(\mathcal{X}_\tau)$ 的内射对象。由引理
\ref{stacks-sheaves-lemma-pullback-injective}，$f_p^{-1}\mathcal{I}$ 是
$\textit{Ab}((\mathcal{U}_p)_\tau)$ 的内射对象。因此，由引理
\ref{stacks-sheaves-lemma-pushforward-injective}，$f_{p, *}f_p^{-1}\mathcal{I}$ 是
$\textit{Ab}(\mathcal{X}_\tau)$ 的内射对象。于是命题
\ref{stacks-sheaves-proposition-exactness-cech-complex} 表明，扩充相对 {\v C}ech 复形
$$
\ldots \to 0 \to
\mathcal{I} \to
f_{0, *}f_0^{-1}\mathcal{I} \to
f_{1, *}f_1^{-1}\mathcal{I} \to
f_{2, *}f_2^{-1}\mathcal{I} \to \ldots
$$
是 $\textit{Ab}(\mathcal{X}_\tau)$ 中的正合复形，而且它的每一项都内射。
对这个复形取整体截面是正合的，因此可见 {\v C}ech 复形
$\check{\mathcal{C}}^\bullet(\mathcal{U} \to \mathcal{X}, \mathcal{I})$
拟同构于 $\Gamma(\mathcal{X}_\tau, \mathcal{I})[0]$。

\medskip\noindent
完成这些准备后，考虑与双复形（参见同调，第
\ref{homology-section-double-complex} 节）
$$
\check{\mathcal{C}}^\bullet(\mathcal{U} \to \mathcal{X}, \mathcal{I}^\bullet)
$$
相联系的两个谱序列，其中 $\mathcal{F} \to \mathcal{I}^\bullet$ 是
$\textit{Ab}(\mathcal{X}_\tau)$ 中的内射分解。以上讨论说明可以应用
同调，引理 \ref{homology-lemma-double-complex-gives-resolution}，从而
$\Gamma(\mathcal{X}_\tau, \mathcal{I}^\bullet)$
拟同构于与该双复形相联系的全复形。由以上说明，复形
$f_p^{-1}\mathcal{I}^\bullet$ 是 $f_p^{-1}\mathcal{F}$ 的内射分解。
因此，另一个谱序列正如引理中所述。
\end{proof}

\noindent
当然，也有模的版本。

\begin{lemma}
\label{stacks-sheaves-lemma-cech-to-cohomology-modules}
设 $f : \mathcal{U} \to \mathcal{X}$ 是 $(\Sch/S)_{fppf}$ 上
群胚纤维化范畴的一个 $1$-态射。令
$\tau \in \{Zar,\allowbreak \etale,\allowbreak smooth,\allowbreak syntomic,\allowbreak fppf\}$.
假设
\begin{enumerate}
\item $\mathcal{F}$ 是
$\textit{Mod}(\mathcal{X}_\tau, \mathcal{O}_\mathcal{X})$ 的对象；
\item 对 $\mathcal{X}$ 的每个对象 $x$，都存在 $\mathcal{X}_\tau$ 中的覆盖
$\{x_i \to x\}$，使每个 $x_i$ 都同构于 $\mathcal{U}$ 中某个对象
$u_i$ 的像 $f(u_i)$；
\item 范畴 $\mathcal{U}$ 有等化子；
\item 函子 $f$ 忠实。
\end{enumerate}
则存在 $\Gamma(\mathcal{O}_\mathcal{X})$-模的第一象限谱序列
$$
E_1^{p, q} = H^q((\mathcal{U}_p)_\tau, f_p^*\mathcal{F})
\Rightarrow
H^{p + q}(\mathcal{X}_\tau, \mathcal{F})
$$
，它收敛到 $\mathcal{F}$ 在 $\tau$-拓扑中的上同调。
\end{lemma}

\begin{proof}
本引理的证明与引理 \ref{stacks-sheaves-lemma-cech-to-cohomology} 的证明相同，
只是使用 $\textit{Mod}(\mathcal{X}_\tau, \mathcal{O}_\mathcal{X})$
中的内射分解，并以引理 \ref{stacks-sheaves-lemma-pullback-injective-modules}
代替引理 \ref{stacks-sheaves-lemma-pullback-injective}。
\end{proof}

\noindent
下面的引理把一种更常见的覆盖转化为以上遇到的那类覆盖。

\begin{lemma}
\label{stacks-sheaves-lemma-surjective-flat-locally-finite-presentation}
设 $f : \mathcal{X} \to \mathcal{Y}$ 是 $(\Sch/S)_{fppf}$ 上
群胚纤维化范畴的一个 $1$-态射。
\begin{enumerate}
\item 假设 $f$ 可由代数空间表示、满、平且局部有限表现。则对
$\mathcal{Y}$ 的任意对象 $y$，都存在一个 fppf 覆盖
$\{y_i \to y\}$ 以及 $\mathcal{X}$ 的对象 $x_i$，使得在
$\mathcal{Y}$ 中 $f(x_i) \cong y_i$。
\item 假设 $f$ 可由代数空间表示、满且光滑。则对 $\mathcal{Y}$ 的
任意对象 $y$，都存在一个 \'etale 覆盖 $\{y_i \to y\}$ 以及
$\mathcal{X}$ 的对象 $x_i$，使得在 $\mathcal{Y}$ 中
$f(x_i) \cong y_i$。
\end{enumerate}
\end{lemma}

\begin{proof}
证明 (1)。设 $y$ 位于概形 $V$ 上。可以把 $y$ 看成态射
$(\Sch/V)_{fppf} \to \mathcal{Y}$。依定义，$2$-纤维积
$\mathcal{X} \times_\mathcal{Y} (\Sch/V)_{fppf}$
可由代数空间 $W$ 表示，而且态射 $W \to V$ 满、平且局部有限表现。
选取概形 $U$ 以及满 \'etale 态射 $U \to W$。则 $U \to V$
也满、平且局部有限表现（参见空间的态射，引理
\ref{spaces-morphisms-lemma-etale-flat},
\ref{spaces-morphisms-lemma-etale-locally-finite-presentation},
\ref{spaces-morphisms-lemma-composition-surjective},
\ref{spaces-morphisms-lemma-composition-finite-presentation}，以及
\ref{spaces-morphisms-lemma-composition-flat}).
因此 $\{U \to V\}$ 是 fppf 覆盖。
% P11-E0137
把与 $1$-态射 $(\Sch/U)_{fppf} \to \mathcal{X}$ 对应的
$\mathcal{X}$ 中位于 $U$ 上的对象记为 $x$。则
$\{f(x) \to y\}$ 是所求的 $\mathcal{Y}$ 的 fppf 覆盖。

\medskip\noindent
证明 (2)。设 $y$ 位于概形 $V$ 上。可以把 $y$ 看成态射
$(\Sch/V)_{fppf} \to \mathcal{Y}$。依定义，$2$-纤维积
$\mathcal{X} \times_\mathcal{Y} (\Sch/V)_{fppf}$
可由代数空间 $W$ 表示，而且态射 $W \to V$ 满且光滑。
选取概形 $U$ 以及满 \'etale 态射 $U \to W$。则 $U \to V$
也满且光滑（参见空间的态射，引理
\ref{spaces-morphisms-lemma-etale-smooth},
\ref{spaces-morphisms-lemma-composition-surjective}，以及
\ref{spaces-morphisms-lemma-composition-smooth}).
因此 $\{U \to V\}$ 是光滑覆盖。由态射进阶，引理
\ref{more-morphisms-lemma-etale-dominates-smooth}，存在 \'etale 覆盖
$\{V_i \to V\}$，使每个 $V_i \to V$ 都经由 $U$ 分解。
% P11-E0138
把与 $1$-态射
$$
(\Sch/V_i)_{fppf} \to (\Sch/U)_{fppf} \to \mathcal{X}.
$$
对应的 $\mathcal{X}$ 中位于 $V_i$ 上的对象记为 $x_i$。
则 $\{f(x_i) \to y\}$ 是所求的 $\mathcal{Y}$ 的 \'etale 覆盖。
\end{proof}

\begin{lemma}
\label{stacks-sheaves-lemma-cech-to-cohomology-relative}
设 $f : \mathcal{U} \to \mathcal{X}$ 与
$g : \mathcal{X} \to \mathcal{Y}$
是 $(\Sch/S)_{fppf}$ 上群胚纤维化范畴的可复合 $1$-态射。令
$\tau \in \{Zar, \etale, smooth, syntomic, \linebreak[0] fppf\}$.
假设
\begin{enumerate}
\item $\mathcal{F}$ 是 $\mathcal{X}_\tau$ 上的 Abel 层；
\item 对 $\mathcal{X}$ 的每个对象 $x$，都存在 $\mathcal{X}_\tau$ 中的覆盖
$\{x_i \to x\}$，使每个 $x_i$ 都同构于 $\mathcal{U}$ 中某个对象
$u_i$ 的像 $f(u_i)$；
\item 范畴 $\mathcal{U}$ 有等化子；
\item 函子 $f$ 忠实。
\end{enumerate}
则存在 $\mathcal{Y}_\tau$ 上 Abel 层的第一象限谱序列
$$
E_1^{p, q} = R^q(g \circ f_p)_*f_p^{-1}\mathcal{F}
\Rightarrow
R^{p + q}g_*\mathcal{F}
$$
，其中所有高阶正像都在 $\tau$-拓扑中计算。
\end{lemma}

\begin{proof}
注意，对 $f : \mathcal{U} \to \mathcal{X}$ 与 $\mathcal{F}$ 的假设
同引理 \ref{stacks-sheaves-lemma-cech-to-cohomology} 中完全相同。因此，该引理证明中的
准备说明在这里同样成立。特别地，这些说明蕴涵
$$
0 \to g_*\mathcal{I} \to
(g \circ f_0)_*f_0^{-1}\mathcal{I} \to
(g \circ f_1)_*f_1^{-1}\mathcal{I} \to \ldots
$$
在 $\mathcal{I}$ 是 $\textit{Ab}(\mathcal{X}_\tau)$ 的内射对象时正合。
在此基础上，考虑同调，第 \ref{homology-section-double-complex} 节中
与双复形 $\mathcal{C}^{\bullet, \bullet}$ 相联系的两个谱序列；其各项为
% P11-E0139
$$
\mathcal{C}^{p, q} = (g \circ f_p)_*\mathcal{I}^q
$$
，其中 $\mathcal{F} \to \mathcal{I}^\bullet$ 是
$\textit{Ab}(\mathcal{X}_\tau)$ 中的内射分解。由第一个谱序列和同调，
引理 \ref{homology-lemma-double-complex-gives-resolution} 可知，
$g_*\mathcal{I}^\bullet$ 拟同构于与
$\mathcal{C}^{\bullet, \bullet}$ 相联系的全复形。由于
$f_p^{-1}\mathcal{I}^\bullet$ 是 $f_p^{-1}\mathcal{F}$ 的内射分解
（参见引理 \ref{stacks-sheaves-lemma-pullback-injective}），第二个谱序列的各项正是
引理陈述中的
$E_1^{p, q} = R^q(g \circ f_p)_*f_p^{-1}\mathcal{F}$。
\end{proof}

\begin{lemma}
\label{stacks-sheaves-lemma-cech-to-cohomology-relative-modules}
设 $f : \mathcal{U} \to \mathcal{X}$ 与
$g : \mathcal{X} \to \mathcal{Y}$
是 $(\Sch/S)_{fppf}$ 上群胚纤维化范畴的可复合 $1$-态射。令
$\tau \in \{Zar, \etale, smooth, syntomic, \linebreak[0] fppf\}$.
假设
\begin{enumerate}
\item $\mathcal{F}$ 是
$\textit{Mod}(\mathcal{X}_\tau, \mathcal{O}_\mathcal{X})$ 的对象；
\item 对 $\mathcal{X}$ 的每个对象 $x$，都存在 $\mathcal{X}_\tau$ 中的覆盖
$\{x_i \to x\}$，使每个 $x_i$ 都同构于 $\mathcal{U}$ 中某个对象
$u_i$ 的像 $f(u_i)$；
\item 范畴 $\mathcal{U}$ 有等化子；
\item 函子 $f$ 忠实。
\end{enumerate}
则在 $\textit{Mod}(\mathcal{Y}_\tau, \mathcal{O}_\mathcal{Y})$ 中存在
第一象限谱序列
% P11-E0140
$$
E_1^{p, q} = R^q(g \circ f_p)_*f_p^{-1}\mathcal{F}
\Rightarrow
R^{p + q}g_*\mathcal{F}
$$
，其中所有高阶正像都在 $\tau$-拓扑中计算。
\end{lemma}

\begin{proof}
证明与引理 \ref{stacks-sheaves-lemma-cech-to-cohomology-relative} 的证明相同，
只是使用 $\textit{Mod}(\mathcal{X}_\tau, \mathcal{O}_\mathcal{X})$
中的内射分解，并以引理 \ref{stacks-sheaves-lemma-pullback-injective-modules}
代替引理 \ref{stacks-sheaves-lemma-pullback-injective}。
\end{proof}






\section{代数叠上的上同调}
\label{stacks-sheaves-section-cohomology}

\noindent
设 $\mathcal{X}$ 是 $S$ 上的代数叠。以上各节已经说明如何定义
$\mathcal{X}$ 上 \'etale、……、fppf 拓扑的层。实际上，对每个
$\tau \in \{Zar,\allowbreak \etale,\allowbreak smooth,\allowbreak syntomic,\allowbreak fppf\}$，我们都构造了位点
$\mathcal{X}_\tau$。在这些位点上有 Abel 层 $\mathcal{F}$ 的概念。
在位点上同调一章中，我们已经说明如何定义上同调。综合这些内容，
把{\it 导出整体截面}或{\it 全上同调}
$$
R\Gamma_{Zar}(\mathcal{X}, \mathcal{F}),
R\Gamma_\etale(\mathcal{X}, \mathcal{F}), \ldots,
R\Gamma_{fppf}(\mathcal{X}, \mathcal{F})
$$
定义为 $\Gamma(\mathcal{X}_\tau, \mathcal{I}^\bullet)$，其中
$\mathcal{F} \to \mathcal{I}^\bullet$ 是
$\textit{Ab}(\mathcal{X}_\tau)$ 中的内射分解。$\mathcal{F}$ 的第 $i$
上同调群，就是全上同调的第 $i$ 个上同调，记为
$$
H^i_{Zar}(\mathcal{X}, \mathcal{F}),
H^i_\etale(\mathcal{X}, \mathcal{F}), \ldots,
H^i_{fppf}(\mathcal{X}, \mathcal{F}).
$$
由态射进阶，引理 \ref{more-morphisms-lemma-etale-dominates-smooth}，
稍后会得出 $H^i_\etale = H^i_{smooth}$。

\medskip\noindent
若 $\mathcal{F}$ 是 $\mathcal{O}_\mathcal{X}$-模预层，并且在
$\tau$-拓扑中是层，则使用
$\textit{Mod}(\mathcal{X}_\tau, \mathcal{O}_\mathcal{X})$ 中的内射分解
来计算其全上同调及上同调群。由非常一般的位点上同调，引理
\ref{sites-cohomology-lemma-cohomology-modules-abelian-agree}，所得结果分别
拟同构及同构于把 $\mathcal{F}$ 看作 Abel 群层时的上同调。

\medskip\noindent
到目前为止，计算上同调群的唯一工具是以上证明的关于 {\v C}ech 复形的
结果。这里用代数叠的语言，针对 \'etale 与 fppf 拓扑重新表述它。设
$f : \mathcal{U} \to \mathcal{X}$ 是代数叠的一个 $1$-态射。回顾
$$
f_p : \mathcal{U}_p =
\mathcal{U} \times_\mathcal{X} \ldots \times_\mathcal{X} \mathcal{U}
\longrightarrow
\mathcal{X}
$$
是结构态射，其中有 $(p + 1)$ 个因子。还要回顾，$\mathcal{X}$ 上的层
是关于 fppf 拓扑的层。注意，若 $\mathcal{U}$ 是代数空间，则
$f : \mathcal{U} \to \mathcal{X}$ 可由代数空间表示；参见代数叠，引理
\ref{algebraic-lemma-representable-diagonal}。因此，本命题尤其适用于
由概形给出的代数叠 $\mathcal{X}$ 的光滑覆盖。

\begin{proposition}
\label{stacks-sheaves-proposition-smooth-covering-compute-cohomology}
\begin{slogan}
叠的上同调可以在一个表现的 {\v C}ech 神经上计算
\end{slogan}
设 $f : \mathcal{U} \to \mathcal{X}$ 是代数叠的一个 $1$-态射。
\begin{enumerate}
\item 设 $\mathcal{F}$ 是 $\mathcal{X}$ 上的 Abel \'etale 层。
假设 $f$ 可由代数空间表示、满且光滑。则存在谱序列
$$
E_1^{p, q} = H^q_\etale(\mathcal{U}_p, f_p^{-1}\mathcal{F})
\Rightarrow
H^{p + q}_\etale(\mathcal{X}, \mathcal{F})
$$
\item 设 $\mathcal{F}$ 是 $\mathcal{X}$ 上的 Abel 层。假设 $f$
可由代数空间表示、满、平且局部有限表现。则存在谱序列
$$
E_1^{p, q} = H^q_{fppf}(\mathcal{U}_p, f_p^{-1}\mathcal{F})
\Rightarrow
H^{p + q}_{fppf}(\mathcal{X}, \mathcal{F})
$$
\end{enumerate}
\end{proposition}

\begin{proof}
为此，检验引理 \ref{stacks-sheaves-lemma-cech-to-cohomology} 的假设 (1) -- (4)。
由代数叠，引理
\ref{algebraic-lemma-characterize-representable-by-algebraic-spaces}，
$1$-态射 $f$ 忠实。这就证明了 (4)。假设 (3) 来自
$\mathcal{U}$ 是代数叠这一事实；参见引理 \ref{stacks-sheaves-lemma-fibre-products}。
为验证 (2)，应用引理
\ref{stacks-sheaves-lemma-surjective-flat-locally-finite-presentation}。条件 (1) 由约定成立。
\end{proof}










\section{高阶正像与代数叠}
\label{stacks-sheaves-section-higher-direct-images}

\noindent
设 $g : \mathcal{X} \to \mathcal{Y}$ 是 $S$ 上代数叠的一个 $1$-态射。
以上各节已经对每个
$\tau \in \{Zar,\allowbreak \etale,\allowbreak smooth,\allowbreak syntomic,\allowbreak fppf\}$ 构造了带环拓扑斯的态射
$g : \Sh(\mathcal{X}_\tau) \to \Sh(\mathcal{Y}_\tau)$。
在位点上同调一章中，我们已经说明如何定义高阶正像。因此，把
{\it 全正像} $Rg_*\mathcal{F}$ 定义为 $g_*\mathcal{I}^\bullet$，其中
$\mathcal{F} \to \mathcal{I}^\bullet$ 是
$\textit{Ab}(\mathcal{X}_\tau)$ 中的内射分解。第 $i$ 个高阶正像
$R^ig_*\mathcal{F}$ 是全正像的第 $i$ 个上同调。须特别注意：
这里使用哪一种拓扑 $\tau$ 至关重要！

\medskip\noindent
若 $\mathcal{F}$ 是 $\mathcal{O}_\mathcal{X}$-模预层，并且在
$\tau$-拓扑中是层，则使用
$\textit{Mod}(\mathcal{X}_\tau, \mathcal{O}_\mathcal{X})$ 中的内射分解
来计算全正像和高阶正像。

\medskip\noindent
到目前为止，计算 $g_*$ 的高阶正像的唯一工具，是以上证明的关于
{\v C}ech 复形的结果。这需要选取一个“覆盖”
$f : \mathcal{U} \to \mathcal{X}$。若 $\mathcal{U}$ 是代数空间，
则 $f : \mathcal{U} \to \mathcal{X}$ 可由代数空间表示；
参见代数叠，引理 \ref{algebraic-lemma-representable-diagonal}。
因此，本命题尤其适用于由概形给出的代数叠 $\mathcal{X}$ 的光滑覆盖。

\begin{proposition}
\label{stacks-sheaves-proposition-smooth-covering-compute-direct-image}
设 $f : \mathcal{U} \to \mathcal{X}$ 与 $g : \mathcal{X} \to \mathcal{Y}$
是代数叠的可复合 $1$-态射。
\begin{enumerate}
\item 假设 $f$ 可由代数空间表示、满且光滑。
\begin{enumerate}
\item 若 $\mathcal{F}$ 属于 $\textit{Ab}(\mathcal{X}_\etale)$，
则存在谱序列
$$
E_1^{p, q} = R^q(g \circ f_p)_*f_p^{-1}\mathcal{F}
\Rightarrow
R^{p + q}g_*\mathcal{F}
$$
，它位于 $\textit{Ab}(\mathcal{Y}_\etale)$ 中，其中的高阶正像
在 \'etale 拓扑中计算。
\item 若 $\mathcal{F}$ 属于
$\textit{Mod}(\mathcal{X}_\etale, \mathcal{O}_\mathcal{X})$，则存在谱序列
% P11-E0142
$$
E_1^{p, q} = R^q(g \circ f_p)_*f_p^{-1}\mathcal{F}
\Rightarrow
R^{p + q}g_*\mathcal{F}
$$
，它位于 $\textit{Mod}(\mathcal{Y}_\etale, \mathcal{O}_\mathcal{Y})$ 中。
\end{enumerate}
\item 假设 $f$ 可由代数空间表示、满、平且局部有限表现。
\begin{enumerate}
\item 若 $\mathcal{F}$ 属于 $\textit{Ab}(\mathcal{X})$，则存在谱序列
$$
E_1^{p, q} = R^q(g \circ f_p)_*f_p^{-1}\mathcal{F}
\Rightarrow
R^{p + q}g_*\mathcal{F}
$$
，它位于 $\textit{Ab}(\mathcal{Y})$ 中，其中的高阶正像在 fppf 拓扑中计算。
\item 若 $\mathcal{F}$ 属于 $\textit{Mod}(\mathcal{O}_\mathcal{X})$，
则存在谱序列
% P11-E0142
$$
E_1^{p, q} = R^q(g \circ f_p)_*f_p^{-1}\mathcal{F}
\Rightarrow
R^{p + q}g_*\mathcal{F}
$$
，它位于 $\textit{Mod}(\mathcal{O}_\mathcal{Y})$ 中。
\end{enumerate}
\end{enumerate}
\end{proposition}

\begin{proof}
为此，检验引理 \ref{stacks-sheaves-lemma-cech-to-cohomology-relative} 和引理
\ref{stacks-sheaves-lemma-cech-to-cohomology-relative-modules} 的假设 (1) -- (4)。
由代数叠，引理
\ref{algebraic-lemma-characterize-representable-by-algebraic-spaces}，
$1$-态射 $f$ 忠实。这就证明了 (4)。假设 (3) 来自
$\mathcal{U}$ 是代数叠这一事实；参见引理 \ref{stacks-sheaves-lemma-fibre-products}。
为验证 (2)，应用引理
\ref{stacks-sheaves-lemma-surjective-flat-locally-finite-presentation}。
在所有四种情形中，条件 (1) 都由约定成立。
\end{proof}

\noindent
下面描述代数叠态射的高阶正像。

\begin{lemma}
\label{stacks-sheaves-lemma-pushforward-restriction}
设 $S$ 是概形。设 $f : \mathcal{X} \to \mathcal{Y}$ 是 $S$ 上代数叠的
$1$-态射\footnote{这个结果应当对 $(\Sch/S)_{fppf}$ 上群胚纤维化范畴的
任意 $1$-态射都成立。}。
% P11-E0141
令 $\tau \in \{Zariski,\linebreak[0] \etale,\linebreak[0]
smooth,\linebreak[0] syntomic,\linebreak[0] fppf\}$.
设 $\mathcal{F}$ 是 $\textit{Ab}(\mathcal{X}_\tau)$ 或
$\textit{Mod}(\mathcal{X}_\tau, \mathcal{O}_\mathcal{X})$ 的对象。
则层 $R^if_*\mathcal{F}$ 是与预层
$$
y \longmapsto
H^i_\tau\Big((\Sch/V)_{fppf} \times_{y, \mathcal{Y}} \mathcal{X},
\ \text{pr}^{-1}\mathcal{F}\Big)
$$
相联系的层。这里 $y$ 是 $\mathcal{Y}$ 中位于概形 $V$ 上的对象。
\end{lemma}

\begin{proof}
选取内射分解 $\mathcal{F}[0] \to \mathcal{I}^\bullet$。
由推前公式 (\ref{stacks-sheaves-equation-pushforward}) 可知，$R^if_*\mathcal{F}$
是与如下预层相联系的层：它给 $y$ 配上复形
$$
\begin{matrix}
\Gamma\Big((\Sch/V)_{fppf} \times_{y, \mathcal{Y}} \mathcal{X},
\ \text{pr}^{-1}\mathcal{I}^{i - 1}\Big) \\
\downarrow \\
\Gamma\Big((\Sch/V)_{fppf} \times_{y, \mathcal{Y}} \mathcal{X},
\ \text{pr}^{-1}\mathcal{I}^i\Big) \\
\downarrow \\
\Gamma\Big((\Sch/V)_{fppf} \times_{y, \mathcal{Y}} \mathcal{X},
\ \text{pr}^{-1}\mathcal{I}^{i + 1}\Big)
\end{matrix}
$$
的上同调。由于 $\text{pr}^{-1}$ 正合，只需证明
$\text{pr}^{-1}$ 保持内射对象。这由引理
\ref{stacks-sheaves-lemma-pullback-injective}、引理
\ref{stacks-sheaves-lemma-pullback-injective-modules} 以及如下事实得出：
$\text{pr}$ 是代数叠的可表示态射（因此，由代数叠，引理
\ref{algebraic-lemma-characterize-representable-by-algebraic-spaces}，
$\text{pr}$ 忠实），而且
$(\Sch/V)_{fppf} \times_{y, \mathcal{Y}} \mathcal{X}$
由引理 \ref{stacks-sheaves-lemma-fibre-products} 有等化子。
\end{proof}

\noindent
下面是一个平凡的基变换结果。

\begin{lemma}
\label{stacks-sheaves-lemma-base-change-higher-direct-images}
设 $S$ 是概形。令
% P11-E0141
$\tau \in \{Zariski,\linebreak[0] \etale,\linebreak[0]
smooth,\linebreak[0] syntomic,\linebreak[0] fppf\}$。设
$$
\xymatrix{
\mathcal{Y}' \times_\mathcal{Y} \mathcal{X} \ar[r]_{g'} \ar[d]_{f'} &
\mathcal{X} \ar[d]^f \\
\mathcal{Y}' \ar[r]^g & \mathcal{Y}
}
$$
是 $S$ 上代数叠的 $2$-Cartesian 图。则基变换映射是同构
$$
g^{-1}Rf_*\mathcal{F} \longrightarrow Rf'_*(g')^{-1}\mathcal{F}
$$
，它对 $\textit{Ab}(\mathcal{X}_\tau)$ 中的 $\mathcal{F}$ 或
$\textit{Mod}(\mathcal{X}_\tau, \mathcal{O}_\mathcal{X})$ 中的
$\mathcal{F}$ 具有函子性。
\end{lemma}

\begin{proof}
同构 $g^{-1}f_*\mathcal{F} = f'_*(g')^{-1}\mathcal{F}$ 正是
引理 \ref{stacks-sheaves-lemma-base-change}（而且它对任意预层都成立）。
对于全正像，由于态射 $g$ 与 $g'$ 平，存在基变换映射；参见位点上同调，
第 \ref{sites-cohomology-section-base-change-map} 节。为证明该映射是拟同构，
可以使用如下事实：对 $\mathcal{Y}'$ 中位于概形 $V$ 上的对象 $y'$，
存在等价
$$
(\Sch/V)_{fppf} \times_{g(y'), \mathcal{Y}} \mathcal{X}
=
(\Sch/V)_{fppf} \times_{y', \mathcal{Y}'}
(\mathcal{Y}' \times_\mathcal{Y} \mathcal{X})
$$
由此及引理 \ref{stacks-sheaves-lemma-pushforward-restriction} 可知，诱导映射
$g^{-1}R^if_*\mathcal{F} \to R^if'_*(g')^{-1}\mathcal{F}$
是同构。
\end{proof}












\section{比较}
\label{stacks-sheaves-section-compare}

\noindent
本节汇集一些比较用叠定义的上同调和用代数空间定义的上同调的结果。

\begin{lemma}
\label{stacks-sheaves-lemma-compare-injectives}
设 $S$ 是概形。设 $\mathcal{X}$ 是 $S$ 上的代数叠，并可由代数空间
$F$ 表示。
\begin{enumerate}
% P11-E0143
\item 若 $\mathcal{I}$ 在 $\textit{Ab}(\mathcal{X}_\etale)$ 中内射，
则 $\mathcal{I}|_{F_\etale}$ 在 $\textit{Ab}(F_\etale)$ 中内射；
\item 若 $\mathcal{I}^\bullet$ 是 $\textit{Ab}(\mathcal{X}_\etale)$ 中的
K-内射复形，则 $\mathcal{I}^\bullet|_{F_\etale}$ 是
$\textit{Ab}(F_\etale)$ 中的 K-内射复形。
\end{enumerate}
对于模，同一结论不成立。
\end{lemma}

\begin{proof}
这形式地来自如下事实：限制函子 $\pi_{F, *} = i_F^{-1}$（参见引理
\ref{stacks-sheaves-lemma-compare}）是正合函子 $\pi_F^{-1}$ 的右伴随；参见同调，引理
\ref{homology-lemma-adjoint-preserve-injectives} 以及导出范畴，引理
\ref{derived-lemma-adjoint-preserve-K-injectives}。关于本引理对模不成立，
参见 \'Etale 上同调，引理
\ref{etale-cohomology-lemma-compare-injectives}。
\end{proof}

\begin{lemma}
\label{stacks-sheaves-lemma-compare-morphism-cohomology}
设 $S$ 是概形。设 $f : \mathcal{X} \to \mathcal{Y}$ 是 $S$ 上代数叠的
态射。假设 $\mathcal{X}$、$\mathcal{Y}$ 分别可由代数空间 $F$、$G$
表示。
% P11-E0144
把诱导的代数空间态射记为 $f : F \to G$。
\begin{enumerate}
\item 对任意 $\mathcal{F} \in \textit{Ab}(\mathcal{X}_\etale)$，有
$$
(Rf_*\mathcal{F})|_{G_\etale} =
Rf_{small, *}(\mathcal{F}|_{F_\etale})
$$
，这是 $D(G_\etale)$ 中的等式。
\item 对 $\textit{Mod}(\mathcal{X}_\etale, \mathcal{O}_\mathcal{X})$
的任意对象 $\mathcal{F}$，有
$$
(Rf_*\mathcal{F})|_{G_\etale} =
Rf_{small, *}(\mathcal{F}|_{F_\etale})
$$
，这是 $D(\mathcal{O}_G)$ 中的等式。
\end{enumerate}
\end{lemma}

\begin{proof}
选取 $\mathcal{F}$ 的一个内射分解后，(1) 立即由引理
\ref{stacks-sheaves-lemma-compare-injectives} 和 (\ref{stacks-sheaves-equation-compare-big-small}) 得出。

\medskip\noindent
可以如下证明 (2)。引理 \ref{stacks-sheaves-lemma-compare-morphism} 已经表明，作为带环位点
的态射有 $\pi_G \circ f = f_{small} \circ \pi_F$。因此，由位点上同调，
引理 \ref{sites-cohomology-lemma-derived-pushforward-composition}，得到
$R\pi_{G, *} \circ Rf_* = Rf_{small, *} \circ R\pi_{F, *}$
。由于限制函子 $\pi_{F, *}$ 与 $\pi_{G, *}$ 正合，结论成立。
\end{proof}

\begin{lemma}
\label{stacks-sheaves-lemma-compare-representable-morphism-cohomology}
设 $S$ 是概形。考虑 $2$-纤维积方块
$$
\xymatrix{
\mathcal{X}' \ar[r]_{g'} \ar[d]_{f'} & \mathcal{X} \ar[d]^f \\
\mathcal{Y}' \ar[r]^g & \mathcal{Y}
}
$$
，其中各项都是 $S$ 上的代数叠。假设 $f$ 可由代数空间表示，且
$\mathcal{Y}'$ 可由代数空间 $G'$ 表示。则 $\mathcal{X}'$ 可由代数空间
$F'$ 表示，并且
% P11-E0145
把诱导的代数空间态射记为 $f' : F' \to G'$ 后，有
$$
g^{-1}(Rf_*\mathcal{F})|_{G'_\etale} =
Rf'_{small, *}((g')^{-1}\mathcal{F}|_{F'_\etale})
$$
，其中 $\mathcal{F}$ 是 $\textit{Ab}(\mathcal{X}_\etale)$ 或
$\textit{Mod}(\mathcal{X}_\etale, \mathcal{O}_\mathcal{X})$ 中的任意对象。
% P11-E0146
\end{lemma}

\begin{proof}
形式地结合引理 \ref{stacks-sheaves-lemma-base-change-higher-direct-images} 与引理
\ref{stacks-sheaves-lemma-compare-morphism-cohomology} 即得。
\end{proof}











\section{拓扑的变换}
\label{stacks-sheaves-section-change-topology}

\noindent
下面的技术性引理说明：若对 $\mathcal{X}$ 中位于平态射之上的态射，
比较映射都是同构，则局部拟凝聚层的 fppf 上同调等于其 \'etale 上同调。

\begin{lemma}
\label{stacks-sheaves-lemma-lqc-flat-base-change-fppf-sheaf}
设 $S$ 是概形，$\mathcal{X}$ 是 $S$ 上的代数叠。设 $\mathcal{F}$ 是
$\mathcal{O}_\mathcal{X}$-模预层。假设
\begin{enumerate}
\item[(a)] $\mathcal{F}$ 局部拟凝聚；
\item[(b)] 对 $\mathcal{X}$ 的任意态射 $\varphi : x \to y$，若它位于
平且局部有限表现的概形态射 $f : U \to V$ 之上，则
比较映射
$c_\varphi : f_{small}^*\mathcal{F}|_{V_\etale} \to
\mathcal{F}|_{U_\etale}$（见 (\ref{stacks-sheaves-equation-comparison-modules})）是同构。
\end{enumerate}
则 $\mathcal{F}$ 是 fppf 拓扑的层。
\end{lemma}

\begin{proof}
设 $\{x_i \to x\}$ 是 $\mathcal{X}$ 的 fppf 覆盖，位于 $S$ 上概形的
fppf 覆盖 $\{f_i : U_i \to U\}$ 之上。由假设，限制
$\mathcal{G} = \mathcal{F}|_{U_\etale}$ 拟凝聚，且比较映射
$f_{i, small}^*\mathcal{G} \to \mathcal{F}|_{U_{i, \etale}}$
都是同构。因此，$\mathcal{F}$ 对覆盖 $\{x_i \to x\}$ 的层条件，
等价于 $\mathcal{G}^a$ 在 $(\Sch/U)_{fppf}$ 上对覆盖
$\{U_i \to U\}$ 的层条件；后者由下降，引理
\ref{descent-lemma-sheaf-condition-holds} 成立。
\end{proof}

\begin{lemma}
\label{stacks-sheaves-lemma-compare-fppf-etale}
设 $S$ 是概形，$\mathcal{X}$ 是 $S$ 上的代数叠。设 $\mathcal{F}$ 是
% P11-E0147
满足下列条件的预层 $\mathcal{O}_\mathcal{X}$-模：
\begin{enumerate}
\item[(a)] $\mathcal{F}$ 局部拟凝聚；
\item[(b)] 对 $\mathcal{X}$ 的任意态射 $\varphi : x \to y$，若它位于
平且局部有限表现的概形态射 $f : U \to V$ 之上，则比较映射
$c_\varphi : f_{small}^*\mathcal{F}|_{V_\etale} \to
\mathcal{F}|_{U_\etale}$（见 (\ref{stacks-sheaves-equation-comparison-modules})）是同构。
\end{enumerate}
则 $\mathcal{F}$ 是 $\mathcal{O}_\mathcal{X}$-模，并且有如下结论：
\begin{enumerate}
\item 若 $\epsilon : \mathcal{X}_{fppf} \to \mathcal{X}_\etale$
是比较态射，则
$R\epsilon_*\mathcal{F} = \epsilon_*\mathcal{F}$.
\item 上同调群 $H^p_{fppf}(\mathcal{X}, \mathcal{F})$ 等于在
$\mathcal{X}$ 的 \'etale 拓扑中计算的上同调群。上同调群
$H^p_{fppf}(x, \mathcal{F})$ 以及导出版本
$R\Gamma(\mathcal{X}, \mathcal{F})$ 和 $R\Gamma(x, \mathcal{F})$ 也类似。
\item 若 $f : \mathcal{X} \to \mathcal{Y}$ 是 $(\Sch/S)_{fppf}$ 上
群胚纤维化范畴的一个 $1$-态射，则 $R^if_*\mathcal{F}$ 等于在
\'etale 上同调中计算的高阶正像的 fppf 层化。
% P11-E0148
导出拉回也类似。
\end{enumerate}
\end{lemma}

\begin{proof}
关于 $\mathcal{F}$ 是 $\mathcal{O}_\mathcal{X}$-模的断言由引理
\ref{stacks-sheaves-lemma-lqc-flat-base-change-fppf-sheaf} 得出。注意，$\epsilon$ 是由
$\mathcal{X}$ 上的恒等函子给出的位点态射。因此，层
$R^p\epsilon_*\mathcal{F}$ 是与预层
$x \mapsto H^p_{fppf}(x, \mathcal{F})$ 相联系的层；参见位点上同调，
引理 \ref{sites-cohomology-lemma-higher-direct-images}。为证明 (1)，
只需证明
$H^p_{fppf}(x, \mathcal{F}) = 0$ 对 $p > 0$
在 $x$ 位于仿射概形 $U$ 上时成立。由引理
\ref{stacks-sheaves-lemma-cohomology-restriction}，有
$H^p_{fppf}(x, \mathcal{F}) = H^p((\Sch/U)_{fppf}, x^{-1}\mathcal{F})$.
结合下降，引理 \ref{descent-lemma-quasi-coherent-and-flat-base-change} 与
概形上同调，引理
\ref{coherent-lemma-quasi-coherent-affine-cohomology-zero}，可见这些
上同调群为零。

\medskip\noindent
以上已经看到，$\epsilon_*\mathcal{F}$ 与 $\mathcal{F}$ 分别是
$\mathcal{X}_\etale$ 和 $\mathcal{X}_{fppf}$ 上与 $\mathcal{X}$ 上同一个
预层对应的层（更一般地，这对 $\mathcal{X}$ 上 fppf 拓扑中的任意层
都成立）。我们常常略加滥用地认同 $\mathcal{F}$ 与
$\epsilon_*\mathcal{F}$；引理的 (2) 和 (3) 应按此意义理解。因此，
(2) 形式地由 (1) 和 Leray 谱序列得出；参见位点上同调，引理
\ref{sites-cohomology-lemma-apply-Leray}。

\medskip\noindent
最后证明 (3)。层 $R^if_*\mathcal{F}$（分别地，
$Rf_{\etale, *}\mathcal{F}$）是与预层
$$
y \longmapsto
H^i_\tau\Big((\Sch/V)_{fppf} \times_{y, \mathcal{Y}} \mathcal{X},
\ \text{pr}^{-1}\mathcal{F}\Big)
$$
相联系的层，其中 $\tau$ 分别取 $fppf$（分别地，$\etale$）；参见引理
\ref{stacks-sheaves-lemma-pushforward-restriction}。注意，由引理
\ref{stacks-sheaves-lemma-pullback-lqc} 与引理 \ref{stacks-sheaves-lemma-comparison}，
$\text{pr}^{-1}\mathcal{F}$ 也满足性质 (a) 和 (b)，故由 (2)，
这两个预层相等。这立即蕴涵 (3)。
\end{proof}

\noindent
我们将用下面的引理比较代数叠上层的 \'etale 上同调与
lisse-\'etale 拓扑斯上的上同调。

\begin{lemma}
\label{stacks-sheaves-lemma-cohomology-on-subcategory}
设 $S$ 是概形，$\mathcal{X}$ 是 $S$ 上的代数叠。令
$\tau = \etale$（分别地，$\tau = fppf$）。设
$\mathcal{X}' \subset \mathcal{X}$ 是满足下列性质的全子范畴：
\begin{enumerate}
\item 若 $x \to x'$ 是 $\mathcal{X}$ 的态射，位于概形的光滑（分别地，
平且局部有限表现）态射之上，并且 $x' \in \Ob(\mathcal{X}')$，
则 $x \in \Ob(\mathcal{X}')$；
\item 存在位于某个概形 $U$ 上的对象 $x \in \Ob(\mathcal{X}')$，
使得相联系的 $1$-态射
$x : (\Sch/U)_{fppf} \to \mathcal{X}$ 光滑且满。
\end{enumerate}
如下得到位点 $\mathcal{X}'_\tau$：规定 $\mathcal{X}'$ 的一个覆盖，是
$\mathcal{X}'$ 中的任意态射族 $\{x_i \to x\}$，且该族是
$\mathcal{X}_\tau$ 中的覆盖。
则嵌入函子 $\mathcal{X}' \to \mathcal{X}_\tau$ 全忠实、余连续且连续，
因而定义拓扑斯态射
$$
g : \Sh(\mathcal{X}'_\tau) \longrightarrow \Sh(\mathcal{X}_\tau)
$$
，并且对所有 $p \geq 0$ 及所有
$\mathcal{F} \in \textit{Ab}(\mathcal{X}_\tau)$，有
$H^p(\mathcal{X}'_\tau, g^{-1}\mathcal{F}) =
H^p(\mathcal{X}_\tau, \mathcal{F})$。
\end{lemma}

\begin{proof}
注意，假设 (1) 蕴涵：若 $\{x_i \to x\}$ 是 $\mathcal{X}_\tau$ 的覆盖，
且 $x \in \Ob(\mathcal{X}')$，则 $x_i \in \Ob(\mathcal{X}')$。
由于 $\mathcal{X}'_\tau$ 中对象的覆盖与把它们视为
$\mathcal{X}_\tau$ 中对象时的覆盖一致，可知
$\mathcal{X}' \to \mathcal{X}$ 连续且余连续。由此得到态射 $g$，而函子
$g^{-1}$ 可认同为限制函子；参见位点，引理 \ref{sites-lemma-when-shriek}。

\medskip\noindent
特别地，若 $\{x_i \to x\}$ 是 $\mathcal{X}'_\tau$ 中的覆盖，则对
$\mathcal{X}$ 上的任意 Abel 层 $\mathcal{F}$，
% P11-E0149
$$
\check H^p(\{x_i \to x\}, g^{-1}\mathcal{F}) =
\check H^p(\{x_i \to x\}, \mathcal{F})
$$
因此，若 $\mathcal{I}$ 是 $\mathcal{X}_\tau$ 上的内射 Abel 层，
则高阶 {\v C}ech 上同调群为零（位点上同调，引理
\ref{sites-cohomology-lemma-injective-trivial-cech}）。于是，对
$\mathcal{X}'$ 的所有对象 $x$，有 $H^p(x, g^{-1}\mathcal{I}) = 0$
（位点上同调，引理
\ref{sites-cohomology-lemma-cech-vanish-collection}）。换言之，
$\mathcal{X}_\tau$ 上的内射 Abel 层关于函子 $H^0(x, g^{-1}-)$ 右无环。
由此，对所有 $\mathcal{F} \in \textit{Ab}(\mathcal{X})$ 以及所有
$x \in \Ob(\mathcal{X}')$，有
$H^p(x, g^{-1}\mathcal{F}) = H^p(x, \mathcal{F})$。

\medskip\noindent
按假设 (2)，选取位于概形 $U$ 上的对象 $x \in \mathcal{X}'$。
特别地，$\mathcal{X}/x \to \mathcal{X}$ 是代数叠的态射，
% P11-E0150
它可由代数空间表示、满且光滑。（注意，$\mathcal{X}/x$ 等价于
$(\Sch/U)_{fppf}$；参见引理 \ref{stacks-sheaves-lemma-localizing}。）层映射
$$
h_x \longrightarrow *
$$
在 $\Sh(\mathcal{X}_\tau)$ 中是满的。事实上，对 $\mathcal{X}$ 的
任意对象 $x'$，都存在 $\tau$-覆盖 $\{x'_i \to x'\}$，使得存在态射
$x'_i \to x$；参见引理
\ref{stacks-sheaves-lemma-surjective-flat-locally-finite-presentation}。由于 $g$ 正合，
层映射
$$
g^{-1}h_x \longrightarrow * = g^{-1}*
$$
在 $\Sh(\mathcal{X}'_\tau)$ 中也是满的。令 $h_{x, n}$ 为
$(n + 1)$ 重乘积 $h_x \times \ldots \times h_x$。则有谱序列
\begin{equation}
\label{stacks-sheaves-equation-spectral-sequence-one}
E_1^{p, q} = H^q(h_{x, p}, \mathcal{F}) \Rightarrow
H^{p + q}(\mathcal{X}_\tau, \mathcal{F})
\end{equation}
以及
\begin{equation}
\label{stacks-sheaves-equation-spectral-sequence-two}
E_1^{p, q} = H^q(g^{-1}h_{x, p}, g^{-1}\mathcal{F}) \Rightarrow
H^{p + q}(\mathcal{X}'_\tau, g^{-1}\mathcal{F})
\end{equation}
；参见位点上同调，引理
\ref{sites-cohomology-lemma-cech-to-cohomology-sheaf-sets}。

\medskip\noindent
情形 I：$\mathcal{X}$ 有终对象 $x$，且它也是 $\mathcal{X}'$ 的对象。
此情形立即由以上第二段的讨论得出。

\medskip\noindent
情形 II：$\mathcal{X}$ 可由代数空间 $F$ 表示。此时，层 $h_{x, n}$
可由 $\mathcal{X}$ 中的对象 $x_n$ 表示。（具体而言，若
$\mathcal{S}_F = \mathcal{X}$，且 $x : U \to F$ 是给定对象，
则 $h_{x, n}$ 可由 $\mathcal{S}_F$ 的对象
$U \times_F \ldots \times_F U \to F$ 表示。）因此
$H^q(h_{x, p}, \mathcal{F}) = H^q(x_p, \mathcal{F})$。
态射 $x_n \to x$ 位于概形的光滑态射之上，故对所有 $n$，都有
$x_n \in \mathcal{X}'$。从而
$H^q(g^{-1}h_{x, p}, g^{-1}\mathcal{F}) = H^q(x_p, g^{-1}\mathcal{F})$。
于是，由第二段的讨论，以上两个谱序列
(\ref{stacks-sheaves-equation-spectral-sequence-one}) 和
(\ref{stacks-sheaves-equation-spectral-sequence-two}) 的 $E_1^{p, q}$ 项一致。
本引理在情形 II 中也成立。

\medskip\noindent
情形 III：$\mathcal{X}$ 是代数叠。我们断言，在此情形下，由以上情形 II，
上同调群 $H^q(h_{x, p}, \mathcal{F})$ 与
% P11-E0151
$H^q(g^{-1}h_{x, n}, g^{-1}\mathcal{F})$ 一致。一旦证明这一点，
结论就会像之前一样得出。

\medskip\noindent
具体而言，考虑范畴 $\mathcal{X}/h_{x, n}$；参见位点，引理
\ref{sites-lemma-localize-topos-site}。由于 $h_{x, n}$ 是 $h_x$ 的
$(n + 1)$ 重乘积，
% P11-E0152
这个范畴的一个对象是 $(n + 2)$-元组
$(y, s_0, \ldots, s_n)$，其中 $y$ 是 $\mathcal{X}$ 的对象，
而每个 $s_i : y \to x$ 都是 $\mathcal{X}$ 的态射。这是
$(\Sch/S)_{fppf}$ 上的范畴。存在等价
$$
\mathcal{X}/h_{x, n}
\longrightarrow
(\Sch/U)_{fppf} \times_\mathcal{X} \ldots \times_\mathcal{X} (\Sch/U)_{fppf}
=: \mathcal{U}_n
$$
，它位于 $(\Sch/S)_{fppf}$ 上。具体而言，若
$x : (\Sch/U)_{fppf} \to \mathcal{X}$ 也表示与 $x$ 相联系的 $1$-态射，
且 $p : \mathcal{X} \to (\Sch/S)_{fppf}$ 是结构函子，则可以把
$(y, s_0, \ldots, s_n)$ 看成
$(y, f_0, \ldots, f_n, \alpha_0, \ldots, \alpha_n)$
，其中 $y$ 是 $\mathcal{X}$ 的对象，$f_i : p(y) \to p(x)$ 是概形态射，
而
% P11-E0153
$\alpha_i : y \to x(f_i)$ 是同构。由这些 $2n+3$-元组
$(y, f_0, \ldots, f_n, \alpha_0, \ldots, \alpha_n)$
构成的范畴，是上面所示代数叠的 $(n + 1)$ 重纤维积
$\mathcal{U}_n$ 的一个实现，正如第 \ref{stacks-sheaves-section-cech} 节中所讨论的。
由位点上同调，引理
\ref{sites-cohomology-lemma-cohomology-on-sheaf-sets}，有
$$
H^p(\mathcal{U}_n, \mathcal{F}|_{\mathcal{U}_n}) =
H^p(\mathcal{X}/h_{x, n}, \mathcal{F}|_{\mathcal{X}/h_{x, n}}) =
H^p(h_{x, n}, \mathcal{F}).
$$
最后讨论它的“加撇”版本。具体而言，经由以上等价，
$\mathcal{X}'/h_{x, n}$ 对应于全子范畴
$\mathcal{U}'_n \subset \mathcal{U}_n$，它由如下元组组成：
$(y, f_0, \ldots, f_n, \alpha_0, \ldots, \alpha_n)$
，其中 $y \in \mathcal{X}'$。因此，对于嵌入
$\mathcal{U}'_n \subset \mathcal{U}_n$，引理陈述中的性质 (1) 显然成立。
为验证性质 (2)，选取位于概形 $W$ 上的对象
$\xi = (y, s_0, \ldots, s_n)$，使得
$(\Sch/W)_{fppf} \to \mathcal{U}_n$ 光滑且满（这是可能的，因为
$\mathcal{U}_n$ 是代数叠）。则
$(\Sch/W)_{fppf} \to \mathcal{U}_n \to (\Sch/U)_{fppf}$
是态射 $x : (\Sch/U)_{fppf} \to \mathcal{X}$ 的若干基变换的复合，
因而光滑；参见代数叠，引理
\ref{algebraic-lemma-base-change-representable-transformations-property} 与
引理 \ref{algebraic-lemma-composition-representable-transformations-property}。
因此，$\mathcal{X}$ 的公理 (1) 蕴涵 $y$ 是 $\mathcal{X}'$ 的对象，
从而 $\xi$ 是 $\mathcal{U}'_n$ 的对象。再次使用
$$
H^p(\mathcal{U}'_n, \mathcal{F}|_{\mathcal{U}'_n}) =
H^p(\mathcal{X}'/h_{x, n}, \mathcal{F}|_{\mathcal{X}'/h_{x, n}}) =
H^p(g^{-1}h_{x, n}, g^{-1}\mathcal{F}).
$$
，现在可以对 $\mathcal{U}'_n \subset \mathcal{U}_n$ 应用情形 II，
从而得出结论。
\end{proof}








\section{限制到仿射对象}
\label{stacks-sheaves-section-alternative}

\noindent
本节给定一个 $(\Sch/S)_{fppf}$ 上群胚纤维化范畴 $\mathcal{X}$，
考虑 $\mathcal{X}$ 的全子范畴 $\mathcal{X}_{affine}$，其对象是位于
仿射概形 $U$ 上的对象 $x$。我们将看到，对任意比 Zariski 拓扑更细的
拓扑 $\tau$，$\mathcal{X}$ 上与 $\mathcal{X}_{affine, \tau}$ 上的层范畴
% P11-E0155
彼此一致。

\begin{definition}
\label{stacks-sheaves-definition-alternative}
设 $p : \mathcal{X} \to (\Sch/S)_{fppf}$ 是群胚纤维化范畴。
{\it 相联系的仿射位点}是 $\mathcal{X}$ 的全子范畴
$\mathcal{X}_{affine}$，其对象是位于概形 $U$ 上且 $U$ 仿射的那些
$x \in \Ob(\mathcal{X})$。在 $\mathcal{X}_{affine}$ 上取混沌拓扑，
即 $\mathcal{X}_{affine}$ 上的层与预层相同。
\end{definition}

\noindent
因此，函子 $p : \mathcal{X} \to (\Sch/S)_{fppf}$ 限制为函子
$$
p : \mathcal{X}_{affine} \longrightarrow (\textit{Aff}/S)_{fppf}
$$
，其中右侧记号引自拓扑，定义
\ref{topologies-definition-big-small-fppf}。显然，
$\mathcal{X}_{affine}$ 在 $(\textit{Aff}/S)_{fppf}$ 上纤维化为群胚。
因此，$\mathcal{X}_{affine}$ 从 $(\textit{Aff}/S)_{Zar}$、
$(\textit{Aff}/S)_\etale$、$(\textit{Aff}/S)_{smooth}$、
$(\textit{Aff}/S)_{syntomic}$ 和 $(\textit{Aff}/S)_{fppf}$ 继承
Zariski、\'etale、光滑、syntomic 和 fppf 拓扑；参见叠，定义
\ref{stacks-definition-topology-inherited}。

\begin{definition}
\label{stacks-sheaves-definition-alternative-inherited-topologies}
设 $p : \mathcal{X} \to (\Sch/S)_{fppf}$ 是群胚纤维化范畴。
\begin{enumerate}
\item {\it 相联系的仿射 Zariski 位点} $\mathcal{X}_{affine, Zar}$，
是 $\mathcal{X}_{affine}$ 上从 $(\textit{Aff}/S)_{Zar}$ 继承的位点结构。
\item {\it 相联系的仿射 \'etale 位点} $\mathcal{X}_{affine, \etale}$，
是 $\mathcal{X}_{affine}$ 上从 $(\textit{Aff}/S)_\etale$ 继承的位点结构。
\item {\it 相联系的仿射光滑位点} $\mathcal{X}_{affine, smooth}$，
是 $\mathcal{X}_{affine}$ 上从 $(\textit{Aff}/S)_{smooth}$ 继承的位点结构。
\item {\it 相联系的仿射 syntomic 位点} $\mathcal{X}_{affine, syntomic}$，
是 $\mathcal{X}_{affine}$ 上从 $(\textit{Aff}/S)_{syntomic}$ 继承的位点结构。
\item {\it 相联系的仿射 fppf 位点} $\mathcal{X}_{affine, fppf}$，
是 $\mathcal{X}_{affine}$ 上从 $(\textit{Aff}/S)_{fppf}$ 继承的位点结构。
\end{enumerate}
\end{definition}

\noindent
由以上讨论，这个定义是有意义的。\newline
取
% P11-E0154
$\tau \in \{Zariski,\linebreak[0] \etale,\linebreak[0] smooth,\linebreak[0] syntomic,\linebreak[0] fppf\}$。\newline
则
$\mathcal{X}_{affine}$ 中具有固定目标的态射族
$\{x_i \to x\}_{i \in I}$ 是 $\mathcal{X}_{affine, \tau}$ 中的覆盖，
当且仅当仿射概形的态射族 $\{p(x_i) \to p(x)\}_{i \in I}$ 是下列
拓扑定义中所定义的标准 $\tau$-覆盖：
\ref{topologies-definition-standard-Zariski},
\ref{topologies-definition-standard-etale},
\ref{topologies-definition-standard-smooth},
\ref{topologies-definition-standard-syntomic}，以及
\ref{topologies-definition-standard-fppf}.

\begin{lemma}
\label{stacks-sheaves-lemma-alternative}
设 $p : \mathcal{X} \to (\Sch/S)_{fppf}$ 是群胚纤维化范畴。\newline
令
% P11-E0154
$\tau \in \{Zariski,\linebreak[0] \etale,\linebreak[0] smooth,\linebreak[0] syntomic,\linebreak[0] fppf\}$。\newline
函子
$\mathcal{X}_{affine, \tau} \to \mathcal{X}_\tau$ 是特殊余连续函子。
因此，它诱导从 $\Sh(\mathcal{X}_{affine, \tau})$ 到
$\Sh(\mathcal{X}_\tau)$ 的拓扑斯等价。
\end{lemma}

\begin{proof}
略。提示：证明与拓扑，引理
\ref{topologies-lemma-affine-big-site-Zariski},
\ref{topologies-lemma-affine-big-site-etale},
\ref{topologies-lemma-affine-big-site-smooth},
\ref{topologies-lemma-affine-big-site-syntomic} 以及
\ref{topologies-lemma-affine-big-site-fppf} 的证明完全相同。
\end{proof}

\noindent
设 $p : \mathcal{X} \to (\Sch/S)_{fppf}$ 是群胚纤维化范畴。
% P11-E0156
把 $\mathcal{O}_\mathcal{X}$ 到 $\mathcal{X}_{affine}$ 的限制记为
$\mathcal{O}$。则 $\mathcal{O}$ 是 $\mathcal{X}_{affine}$ 上
Zariski、\'etale、光滑、syntomic 和 fppf 拓扑中的层。此外，引理
\ref{stacks-sheaves-lemma-alternative} 的拓扑斯等价扩张为等价
\begin{equation}
\label{stacks-sheaves-equation-alternative-ringed}
(\Sh(\mathcal{X}_{affine, \tau}), \mathcal{O})
\longrightarrow
(\Sh(\mathcal{X}_\tau), \mathcal{O}_\mathcal{X})
\end{equation}
，这是对
% P11-E0154
$\tau \in \{Zariski,\linebreak[0] \etale,\linebreak[0] smooth,\linebreak[0] syntomic,\linebreak[0] fppf\}$ 的带环拓扑斯等价。









\section{拟凝聚模与仿射对象}
\label{stacks-sheaves-section-alternative-quasi-coherent}

\noindent
设 $p : \mathcal{X} \to (\Sch/S)_{fppf}$ 是群胚纤维化范畴。
第 \ref{stacks-sheaves-section-alternative} 节已经为它配上带环位点
$(\mathcal{X}_{affine}, \mathcal{O})$。

\begin{lemma}
\label{stacks-sheaves-lemma-quasi-coherent-alternative}
设 $p : \mathcal{X} \to (\Sch/S)_{fppf}$ 是群胚纤维化范畴。
设 $\mathcal{F}$ 是 $\mathcal{X}_{affine}$ 上的 $\mathcal{O}$-模。
下列条件等价：
\begin{enumerate}
\item 对 $\mathcal{X}_{affine}$ 的每个态射 $x \to x'$，映射
$\mathcal{F}(x') \otimes_{\mathcal{O}(x')} \mathcal{O}(x) \to \mathcal{F}(x)$
是同构；
\item 按位点上的模，定义 \ref{sites-modules-definition-site-local} 的意义，
$\mathcal{F}$ 是 $(\mathcal{X}_{affine}, \mathcal{O})$ 上的拟凝聚模；
\item $\mathcal{F}$ 是 $\mathcal{X}_{affine}$ 上 Zariski 拓扑的层，
并且按位点上的模，定义 \ref{sites-modules-definition-site-local} 的意义，
它是 $(\mathcal{X}_{affine, Zar}, \mathcal{O})$ 上的拟凝聚模；
\item 对 \'etale 拓扑，条件与 (3) 相同；
\item 对光滑拓扑，条件与 (3) 相同；
\item 对 syntomic 拓扑，条件与 (3) 相同；
\item 对 fppf 拓扑，条件与 (3) 相同；
\item 经由等价 (\ref{stacks-sheaves-equation-alternative-ringed})，$\mathcal{F}$ 对应于
$\mathcal{X}$ 上的拟凝聚模。
\end{enumerate}
\end{lemma}

\begin{proof}
为说明 (2) 的含义，回顾 $\mathcal{X}_{affine}$ 是在范畴
$\mathcal{X}_{affine}$ 上赋予混沌拓扑而得到的位点（定义
\ref{stacks-sheaves-definition-alternative}），所以 $\mathcal{O}$-模层 $\mathcal{F}$
与 $\mathcal{O}$-模预层是一回事。由位点上的模，引理
\ref{sites-modules-lemma-chaotic-quasi-coherent}，条件 (1) 与 (2) 等价。
注意，对
% P11-E0154
$\tau \in \{Zariski,\linebreak[0] \etale,\linebreak[0] smooth,\linebreak[0] syntomic,\linebreak[0] fppf\}$，预层
$\mathcal{F}$ 是 $\tau$-层，当且仅当对所有
$x \in \Ob(\mathcal{X}_{affine})$，它到 $\mathcal{X}_{affine}/x$
的限制是 $\tau$-层。置 $U = p(x)$。与第 \ref{stacks-sheaves-section-restriction} 节
的讨论类似，$\mathcal{X}_{affine}$ 的对象 $x$ 诱导位点的等价
$\mathcal{X}_{affine, \etale}/x \to (\textit{Aff}/U)_\etale$。
因此，对每个这样的位点应用下降，引理
\ref{descent-lemma-quasi-coherent-alternative}，便得到 (1) 与 (3) -- (7)
的等价。由“拟凝聚”是模层的内蕴性质这一事实，(8) 与 (7) 的等价
立即成立；参见位点上的模，第 \ref{sites-modules-section-intrinsic} 节。
% P11-E0157
\end{proof}

\begin{lemma}
\label{stacks-sheaves-lemma-locally-quasi-coherent-alternative}
设 $p : \mathcal{X} \to (\Sch/S)_{fppf}$ 是群胚纤维化范畴。
设 $\mathcal{F}$ 是 $\mathcal{X}_{affine}$ 上的 $\mathcal{O}$-模。
下列条件等价：
\begin{enumerate}
\item 对 $\mathcal{X}_{affine}$ 的每个态射 $x \to x'$，若
$p(x) \to p(x')$ 是（仿射概形的）\'etale 态射，则映射
$\mathcal{F}(x') \otimes_{\mathcal{O}(x')} \mathcal{O}(x) \to \mathcal{F}(x)$
是同构；
\item $\mathcal{F}$ 是 $\mathcal{X}_{affine}$ 上 \'etale 拓扑的层，
而且对 $\mathcal{X}_{affine}$ 的每个对象 $x$，限制
$x^*\mathcal{F}|_{U_{affine, \etale}}$ 拟凝聚，其中 $U = p(x)$；
\item 经由 \'etale 拓扑的等价 (\ref{stacks-sheaves-equation-alternative-ringed})，
$\mathcal{F}$ 对应于 $\mathcal{X}$ 上的局部拟凝聚模。
\end{enumerate}
\end{lemma}

\begin{proof}
为说明条件 (2) 的含义，回顾 $U_{affine, \etale}$ 是 $U_\etale$
中由仿射对象组成的全子范畴；参见拓扑，定义
\ref{topologies-definition-big-small-etale}。与第
\ref{stacks-sheaves-section-restriction} 节的讨论类似，$\mathcal{X}_{affine}$ 的对象
$x$ 诱导位点的等价
$\mathcal{X}_{affine, \etale}/x \to (\textit{Aff}/U)_\etale$。
于是，$x^*\mathcal{F}$ 是 $(\textit{Aff}/U)_\etale$ 上与限制
$\mathcal{F}|_{\mathcal{X}_{affine, \etale}/x}$ 对应的模层。
最后，利用连续且余连续的嵌入函子
$U_{affine, \etale} \to (\textit{Aff}/U)_\etale$，可以进一步限制，
得到 $x^*\mathcal{F}|_{U_{affine, \etale}}$。

\medskip\noindent
对 $\mathcal{X}$ 中位于仿射概形 $U$ 上的每个对象 $x$，把下降，引理
\ref{descent-lemma-quasi-coherent-alternative-small} 应用于
$\mathcal{F}$ 到 $U_{affine, \etale}$ 的限制，再结合以上说明，便得到
(1) 与 (2) 的等价。由定义以及 $U_{affine, \etale}$ 与 $U_\etale$ 上的
拟凝聚模相互对应这一事实（例如再次使用下降，引理
\ref{descent-lemma-quasi-coherent-alternative-small}），(2) 与 (3) 的等价
立即成立。
\end{proof}













\section{导出范畴中的拟凝聚对象}
\label{stacks-sheaves-section-QC}

\noindent
数十年来，代数几何学家一直在研究 $\Sch/S$ 上不可表示函子 $X$
（取值于集合或群胚）的不变量。例如，在叠的概念发明之前，Mumford
把椭圆曲线模函子 $X$ 的 Picard 群胚 $\Pic(X)$ 定义为带有到 $X$
的映射（即带有一族椭圆曲线）的所有概形 $U$ 所成范畴上的
$2$-极限 $Pic(U)$ \cite{mumford_picard}。类似地，Beilinson--Drinfeld
把 ind-概形 $X = \colim X_i$ 的范畴 $\QCoh(X)$ 定义为 $2$-极限
$\lim \QCoh(X_i)$ \cite{BVGD}。这种策略足以定义像 $\QCoh(-)$ 这样的
$1$-范畴不变量，却不足以定义导出范畴不变量（例如拟凝聚导出范畴），
因为三角范畴的 $2$-极限性质不佳。随着高阶范畴技术和导出代数几何的
出现，这个问题可以优雅地解决：可以把函子 $X$ 的拟凝聚导出
$\infty$-范畴 $\mathcal{D}_{qc}(X)$ 定义为极限
$\lim \mathcal{D}_{qc}(U)$，其中 $U$ 遍历 X 上的所有导出仿射对象
（参见 \cite{lurie-thesis}）。

\medskip\noindent
本节的目标是给上述函子 $X$（取值于集合或群胚）配上三角范畴
$\mathit{QC}(X)$。事实上，这个构造适用于任意群胚纤维化范畴
$p : \mathcal{X} \to (\Sch/S)_{fppf}$（不只适用于分裂的情形）。
在良好情形下，可以证明范畴 $\mathit{QC}(\mathcal{X})$ 与
$\mathcal{D}_{qc}(\mathcal{X})$ 的同伦范畴一致，不过说明这个比较超出
本文范围。这个构造的显著特点如下：
\begin{enumerate}
\item[(a)] 依构造，$\mathit{QC}(\mathcal{X})$ 是
$D(\mathcal{X}_{affine}, \mathcal{O})$ 的全子范畴；
\item[(b)] 当 $\mathcal{X}$ 可由代数空间 $X$ 表示时，
$\mathit{QC}(\mathcal{X})$ 与 $D_\QCoh(\mathcal{O}_X)$ 一致；
\item[(c)] 当 $\mathcal{X}$ 是代数叠时，$\mathit{QC}(\mathcal{X})$
与 $D_\QCoh(\mathcal{O}_\mathcal{X})$ 一致；
\item[(d)] 当 $X = \text{Spf}(A)$ 是与 Noether 环 $A$ 相联系的
仿射形式代数空间，并对理想 $I$ 带有 $I$-进拓扑时，三角范畴
$\mathit{QC}(X)$ 与导出完备对象所成的全子范畴
$D_{comp}(A, I) \subset D(A)$ 一致。
\end{enumerate}
这些结果分别在命题 \ref{stacks-sheaves-proposition-QC-compare}、叠的导出范畴，命题
\ref{stacks-perfect-proposition-QC-compare} 以及命题
\ref{stacks-sheaves-proposition-QC-derived-complete} 中证明。

\medskip\noindent
为说明 $\mathit{QC}(\mathcal{X})$ 精确定义的动机，请注意引理
\ref{stacks-sheaves-lemma-quasi-coherent-alternative} 中的刻画：$\mathcal{X}$ 上的
拟凝聚模可刻画为 $\mathcal{X}_{affine}$ 上满足某种基变换性质的
$\mathcal{O}$-模预层。

\begin{definition}
\label{stacks-sheaves-definition-QC}
设 $p : \mathcal{X} \to (\Sch/S)_{fppf}$ 是群胚纤维化范畴。
设 $\mathcal{O}$ 是第 \ref{stacks-sheaves-section-alternative} 节中引入的
$\mathcal{X}_{affine}$ 上的环层。用公式
$$
\mathit{QC}(\mathcal{X}) = \mathit{QC}(\mathcal{X}_{affine}, \mathcal{O})
$$
定义{\it 导出范畴中拟凝聚对象所成的三角范畴}，其中右侧如位点上同调，
定义 \ref{sites-cohomology-definition-cartesian} 所定义。
\end{definition}

\noindent
注意，这个定义是有意义的：$\mathcal{X}_{affine}$ 是一个范畴，
通过赋予混沌拓扑而被视为位点，而 $\mathcal{O}$ 是这个范畴上的环层；
这恰好符合位点上同调，定义
\ref{sites-cohomology-definition-cartesian} 的要求。

\medskip\noindent
一般而言，这个定义与 $\mathcal{X}$ 上拟凝聚模的范畴之间的关系并不
十分清楚！例如，设 $M$ 是 $\mathit{QC}(\mathcal{X})$ 的对象。
则 $M$ 的上同调层 $H^i(M)$ 是 $\mathcal{X}_{affine}$ 上的
$\mathcal{O}$-模（预）层，但一般并不拟凝聚。不过，最后一个非零
上同调层是拟凝聚的。

\begin{lemma}
\label{stacks-sheaves-lemma-QC-bounded-above}
在定义 \ref{stacks-sheaves-definition-QC} 的情形中，设 $M$ 是
$\mathit{QC}(\mathcal{X})$ 的对象，且 $b \in \mathbf{Z}$ 满足对所有
$i > b$ 都有 $H^i(M) = 0$。则 $H^b(M)$ 是
$(\mathcal{X}_{affine}, \mathcal{O})$ 上的拟凝聚模；参见引理
\ref{stacks-sheaves-lemma-quasi-coherent-alternative}。
\end{lemma}

\begin{proof}
这是位点上同调，引理
\ref{sites-cohomology-lemma-QC-bounded-above} 的特殊情形。
\end{proof}

\begin{lemma}
\label{stacks-sheaves-lemma-QC-compare}
设 $S$ 是概形。设 $\mathcal{X} \to (\Sch/S)_{fppf}$ 是群胚纤维化范畴。
比较态射
$\epsilon : \mathcal{X}_{affine, \etale} \to \mathcal{X}_{affine}$
满足位点上同调，引理
\ref{sites-cohomology-lemma-cartesion-plus-topology} 的假设和结论。
\end{lemma}

\begin{proof}
假设 (1) 由 $\mathcal{X}_{affine}$ 的定义成立。对于条件 (2)，
使用如下事实：对位于仿射概形 $U = p(x)$ 上的
$x \in \Ob(\mathcal{X})$，有与结构层相容的等价
$\mathcal{X}_{affine, \etale}/x = (\textit{Aff}/U)_\etale$；
参见第 \ref{stacks-sheaves-section-restriction} 节的讨论。因此，只需证明：
给定仿射概形 $U = \Spec(R)$ 与 $R$-模复形 $M^\bullet$，
$(\textit{Aff}/U)_\etale$ 上与 $M^\bullet$ 相联系的模复形的全上同调
拟同构于 $M^\bullet$。这由以下结果结合得出：概形的导出范畴，引理
\ref{perfect-lemma-affine-compare-bounded}（Zariski 拓扑中仿射对象上的
模复形的全上同调）；空间的导出范畴，注
\ref{spaces-perfect-remark-match-total-direct-images}（拟凝聚模复形在小
Zariski 拓扑和小 \'etale 拓扑中的全上同调一致）；以及 \'Etale 上同调，
引理 \ref{etale-cohomology-lemma-compare-cohomology}（概形的大 \'etale
位点上的模复形的 \'etale 上同调，可以在限制到小 \'etale 位点后计算）。
\end{proof}

\noindent
若把这个定义应用于群胚纤维化范畴 $\mathcal{X}$ 可由代数空间 $X$
表示的情形，就会重新得到 $D_\QCoh(\mathcal{O}_X)$。稍后将陈述并证明
代数叠的类似结果（此处插入未来的引用）。
% P11-E0158

\begin{proposition}
\label{stacks-sheaves-proposition-QC-compare}
设 $S$ 是概形。设 $\mathcal{X} \to (\Sch/S)_{fppf}$ 是群胚纤维化范畴。
假设 $\mathcal{X}$ 可由代数空间 $X$ 表示。则
$\mathit{QC}(\mathcal{X})$ 典范等价于 $D_\QCoh(\mathcal{O}_X)$。
\end{proposition}

\begin{proof}
% P11-E0159
把由 $X$ 上 \'etale 的仿射概形组成、赋予混沌拓扑及其结构层
$\mathcal{O}_X$ 的范畴记为 $X_{affine}$；参见空间的导出范畴，第
\ref{spaces-perfect-section-QC} 节。引理 \ref{stacks-sheaves-lemma-compare} 的函子
$u : X_\etale \to \mathcal{X}_\etale$ 给出函子
$X_{affine} \to \mathcal{X}_{affine}$。它与结构层相容，并产生函子
$$
G :
\mathit{QC}(\mathcal{X}) = \mathit{QC}(\mathcal{X}_{affine}, \mathcal{O})
\longrightarrow
\mathit{QC}(X_{affine}, \mathcal{O}_X)
$$
参见位点上同调，引理
\ref{sites-cohomology-lemma-cartesian-functorial}。由空间的导出范畴，引理
\ref{spaces-perfect-lemma-DQCoh-alternative-small}，三角范畴
$\mathit{QC}(X_{affine}, \mathcal{O}_X)$ 等价于
$D_\QCoh(\mathcal{O}_X)$。因此，只需证明 $G$ 是等价。

\medskip\noindent
考虑带环位点的平比较态射
$\epsilon_\mathcal{X} : \mathcal{X}_{affine, \etale} \to \mathcal{X}_{affine}$，\newline
以及 $\epsilon_X : X_{affine, \etale} \to X_{affine}$
。引理 \ref{stacks-sheaves-lemma-QC-compare} 与空间的导出范畴，引理
\ref{spaces-perfect-lemma-DQCoh-alternative-small}（的证明）说明，函子
$\epsilon_\mathcal{X}^*$ 与 $\epsilon_X^*$
把 $\mathit{QC}(\mathcal{X}_{affine}, \mathcal{O})$ 与
$\mathit{QC}(X_{affine}, \mathcal{O}_X)$ 分别认同为子范畴
$Q_\mathcal{X} \subset D(\mathcal{X}_{affine, \etale}, \mathcal{O})$
以及
$Q_X \subset D(X_{affine, \etale}, \mathcal{O}_X)$.
在这些认同下，第一段中的函子 $G$ 由函子
$$
Li_X^* = R\pi_{X, *}:
D(\mathcal{X}_{affine, \etale}, \mathcal{O})
\longrightarrow
D(X_{affine, \etale}, \mathcal{O}_X)
$$
诱导，其中 $i_X$ 与 $\pi_X$ 是引理 \ref{stacks-sheaves-lemma-compare} 中的态射，
但把 \'etale 位点换成了相应的仿射位点。读者可以直接对仿射位点
重新证明该引理，或者使用拓扑斯等价
$\Sh(\mathcal{X}_{affine, \etale}) = \Sh(\mathcal{X}_\etale)$ 以及
$\Sh(X_{affine, \etale}) = \Sh(X_\etale)$，来证明这种替换是允许的。
该引理还说明 $Li_X^*$ 有左伴随
$$
L\pi_X^*:
D(X_{affine, \etale}, \mathcal{O}_X)
\longrightarrow
D(\mathcal{X}_{affine, \etale}, \mathcal{O})
$$
；而且由于 $\pi_X \circ i_X$ 是恒等态射，有
$Li_X^* \circ L\pi_X^* = \text{id}$。因此，只需证明 (a)
$L\pi_X^*$ 把 $Q_X$ 映入 $Q_\mathcal{X}$，以及 (b) $Li_X^*$ 的核为
$0$。参见导出范畴，引理
\ref{derived-lemma-fully-faithful-adjoint-kernel-zero}。

\medskip\noindent
证明 (a)。由空间的导出范畴，引理
\ref{spaces-perfect-lemma-DQCoh-alternative-small}，有
$Q_X = D_\QCoh(X_{affine, \etale}, \mathcal{O}_X)$。设 $K$ 是 $Q_X$
的对象。设 $x$ 是 $\mathcal{X}_{affine, \etale}$ 中位于仿射概形
$U = p(x)$ 上的对象。
% P11-E0160
把与 $x$ 对应的态射记为 $f : U \to X$。则可见
$$
R\Gamma(x, L\pi_X^*K) = R\Gamma(U, Lf^*K)
$$
这来自拉回的传递性；参见第
\ref{stacks-sheaves-section-restriction-algebraic-spaces} 节的讨论。接下来，设
$x \to x'$ 是 $\mathcal{X}_{affine, \etale}$ 的态射，位于仿射概形的
态射 $h : U \to U'$ 之上。
% P11-E0161
与之前一样，把与 $x$、$x'$ 对应的态射分别记为
$f : U \to X$ 和 $f' : U' \to X$，于是 $f = f' \circ h$。则
\begin{align*}
R\Gamma(x, L\pi_X^*K)
& =
R\Gamma(U, Lf^*K) \\
& =
R\Gamma(U, Lh^*L(f')^*K) \\
& =
R\Gamma(U', L(f')^*K) \otimes_{\mathcal{O}(U')}^\mathbf{L} \mathcal{O}(U) \\
& =
R\Gamma(x', L\pi_X^*K) \otimes_{\mathcal{O}(x')}^\mathbf{L} \mathcal{O}(x)
\end{align*}
因此，由位点上同调，引理
\ref{sites-cohomology-lemma-cartesion-plus-topology} 陈述中的脚注，
得到 (a)。第三个等式是概形的导出范畴，引理
\ref{perfect-lemma-quasi-coherence-pullback}。

\medskip\noindent
证明 (b)。设 $M$ 是 $Q_\mathcal{X}$ 的对象，且 $Li_X^*M = 0$。
设 $x'$ 是 $\mathcal{X}_{affine, \etale}$ 中位于仿射概形
$U' = p(x')$ 上的对象，并假设相应的态射 $f' : U' \to X$ 是
\'etale 的。则 $f' : U' \to X$ 是 $X_{affine, \etale}$ 的对象，而条件
$Li_X^*M = 0$ 蕴涵 $M|_{U'_\etale} = 0$。特别地，
$R\Gamma(x', M) = 0$。另一方面，对位点
$\mathcal{X}_{affine, \etale}$ 的任意对象 $x$，存在覆盖
$\{x_i \to x\}$，使得对每个 $i$ 都存在态射 $x_i \to x'_i$，
其中 $x'_i$ 对应于 $X_{affine, \etale}$ 的对象。由于 $M$ 属于
$Q_\mathcal{X}$，有
$$
R\Gamma(x_i, M) =
R\Gamma(x_i', M) \otimes_{\mathcal{O}(x_i')}^\mathbf{L} \mathcal{O}(x_i) =
0
$$
由此得出 $M$ 为零，正如所需。
\end{proof}

\noindent
为说明这个构造在另一个情形中也产生有趣的范畴，下面对 Noether
adic 环 $A$ 的形式谱所对应的 $\mathit{QC}(\text{Spf}(A))$ 作出刻画
并加以证明。

\begin{proposition}
\label{stacks-sheaves-proposition-QC-derived-complete}
设 $S$ 是概形。设 $X = \text{Spf}(A)$，其中 $A$ 是以 $I$ 为定义理想的
% P11-E0162
adic Noether 拓扑 $S$-代数；参见代数进阶，定义
\ref{more-algebra-definition-topological-ring} 以及形式空间，定义
\ref{formal-spaces-definition-affine-formal-spectrum}。
% P11-E0163
设 $p : \mathcal{X} \to (\Sch/S)_{fppf}$ 是与函子 $X$ 相联系的
集合纤维化范畴；参见范畴，例 \ref{categories-example-presheaf}。
则 $\mathit{QC}(\mathcal{X})$ 典范等价于 $D(A)$ 中关于 $I$ 导出完备的
对象所成的范畴 $D_{comp}(A, I)$。
\end{proposition}

\begin{proof}
回顾，作为 fppf 层有 $X = \colim \Spec(A/I^n)$。
$\mathcal{X}_{affine}$ 的一个对象，等同于带有给定态射
$f : U \to X$ 的仿射概形 $U = \Spec(R)$。由形式空间，引理
\ref{formal-spaces-lemma-factor-through-thickening}，存在 $n \geq 1$，
使 $f$ 经由单态射 $\Spec(A/I^n) \to X$ 分解。考虑全子范畴
$\mathcal{C} \subset \mathcal{X}_{affine}$，它由对象
$\Spec(A/I^n) \to X$ 组成。由刚才的说明和微分分次层，引理
\ref{sdga-lemma-cartesian-cofinal-subcategory}，到 $\mathcal{C}$ 的限制是
正合等价
$\mathit{QC}(\mathcal{X}) \to
\mathit{QC}(\mathcal{C}, \mathcal{O}|_\mathcal{C})$.
为简单起见，假设对所有 $n \geq 1$ 都有
$I^n \not = I^{n + 1}$。则带环位点
$(\mathcal{C}, \mathcal{O}|_\mathcal{C})$ 同构于带环位点
$(\mathbf{N}, (A/I^n))$；参见微分分次层，第
\ref{sdga-section-qc-inverse-system} 节。因此，结论由微分分次层，命题
\ref{sdga-proposition-derived-complete} 得出。
\end{proof}

\noindent
当 $\mathcal{X}$ 是代数叠时，下面的引理将用于比较
$\mathit{QC}(\mathcal{X})$ 与 $D_\QCoh(\mathcal{O}_\mathcal{X})$。

\begin{lemma}
\label{stacks-sheaves-lemma-QC-compare-fppf}
设 $S$ 是概形。设 $\mathcal{X} \to (\Sch/S)_{fppf}$ 是群胚纤维化范畴。
比较态射
$\epsilon : \mathcal{X}_{affine, fppf} \to \mathcal{X}_{affine}$
满足位点上同调，引理
\ref{sites-cohomology-lemma-cartesion-plus-topology} 的假设和结论。
\end{lemma}

\begin{proof}
证明与引理 \ref{stacks-sheaves-lemma-QC-compare} 的证明完全相同。假设 (1) 由
$\mathcal{X}_{affine}$ 的定义成立。对于条件 (2)，使用如下事实：
对位于仿射概形 $U = p(x)$ 上的 $x \in \Ob(\mathcal{X})$，有等价
% P11-E0164
$\mathcal{X}_{affine, \etale}/x = (\textit{Aff}/U)_\etale$，且该等价
与结构层相容；参见第 \ref{stacks-sheaves-section-restriction} 节的讨论。
因此，只需证明：给定仿射概形 $U = \Spec(R)$ 与 $R$-模复形
$M^\bullet$，$(\textit{Aff}/U)_{fppf}$ 上与 $M^\bullet$ 相联系的
模复形的全上同调拟同构于 $M^\bullet$。这正是 \'Etale 上同调，引理
\ref{etale-cohomology-lemma-cohomological-descent-complex-modules}。
\end{proof}
% Generated mechanically from the frozen criteria.tex target.
% Do not edit this generated file; edit the bound locale source instead.

% OK, start here.
%

\setcounter{chapter}{96}
\chapter{可表示性判据}


\phantomsection
\label{criteria-section-phantom}





\section{引言}
\label{criteria-section-introduction}

\noindent
本章的目的是寻找一些判据，用以保证赋予 fppf 拓扑的概形范畴上的
群胚叠是代数叠。历史上，这往往涉及证明某些函子可表示；参见
Grothendieck 的讲义
\cite{Gr-I},
\cite{Gr-II},
\cite{Gr-III},
\cite{Gr-IV},
\cite{Gr-V}，以及
\cite{Gr-VI}.
本章的标题由此而来。本章素材的另一个重要来源是 Artin 的工作；参见
\cite{ArtinI},
\cite{ArtinII},
\cite{Artin-Theorem-Representability},
\cite{Artin-Construction-Techniques},
\cite{Artin-Algebraic-Spaces},
\cite{Artin-Algebraic-Approximation},
\cite{Artin-Implicit-Function}，以及
\cite{ArtinVersal}.

\medskip\noindent
% P11-E0165
本章中的某些记号、约定和术语颇为别扭，经验更丰富的读者甚至可能觉得
它们前后颠倒。这是有意为之。解释见《Quot 空间与 Hilbert 空间》第
\ref{quot-section-conventions} 节。



\section{约定}
\label{criteria-section-conventions}

\noindent
本章采用的约定与《代数叠》一章中的约定相同；参见《代数叠》第
\ref{algebraic-section-conventions} 节。




\section{已有结果}
\label{criteria-section-done-so-far}

\noindent
代数空间情形下与本章对应的是题为《自举》的一章；参见《自举》第
\ref{bootstrap-section-introduction} 节。该章已经包含一些可表示性结果。
此外，其中处理的一部分预备材料，我们已经在《代数叠》一章中建立。
具体如下：
\begin{enumerate}
\item 我们在《自举》第
\ref{bootstrap-section-morphism-representable-by-spaces} 节讨论可由代数空间
表示的预层态射。在《代数叠》第
\ref{algebraic-section-morphisms-representable-by-algebraic-spaces}
节，我们讨论群胚纤维化范畴的 $1$-态射可由代数空间表示这一概念。
\item 我们在《自举》第
\ref{bootstrap-section-representable-by-spaces-properties} 节讨论可由代数空间
表示的预层态射的性质。在《代数叠》第
\ref{algebraic-section-representable-properties}
节，我们讨论可由代数空间表示的群胚纤维化范畴的 $1$-态射的性质。
\item 我们已经证明：若 $F$ 是一个层，其对角线可由代数空间表示，
并且它有一个来自代数空间的 \'etale 覆盖，则 $F$ 是代数空间；参见
《自举》，定理 \ref{bootstrap-theorem-bootstrap}。
（这是下一项结果的一个较弱版本。）
\item
\label{criteria-item-bootstrap-final}
我们已经证明：若 $F$ 是一个层，并且存在代数空间 $U$ 以及态射
$U \to F$，此态射可由代数空间表示、满、平坦且局部有限表示，
则 $F$ 是代数空间；参见《自举》，定理
\ref{bootstrap-theorem-final-bootstrap}。
\item 我们还证明了 (\ref{criteria-item-bootstrap-final}) 对代数叠的“光滑”类似物：
% P11-E0166
若 $\mathcal{X}$ 是 $(\Sch/S)_{fppf}$ 上的群胚叠，并且存在
$(\Sch/S)_{fppf}$ 上可由代数空间表示的群胚叠 $\mathcal{U}$，以及
$1$-态射 $u : \mathcal{U} \to \mathcal{X}$，此态射可由代数空间表示、
满且光滑，则 $\mathcal{X}$ 是代数叠；参见《代数叠》，引理
\ref{algebraic-lemma-smooth-surjective-morphism-implies-algebraic}.
\end{enumerate}
我们现在的第一项任务，是证明 (\ref{criteria-item-bootstrap-final}) 对一般代数叠的
类似物，即定理 \ref{criteria-theorem-bootstrap}。



\section{群胚叠的态射}
\label{criteria-section-1-morphisms}

\noindent
本节为预备内容，初读时可略过。

\begin{lemma}
\label{criteria-lemma-etale-permanence}
设 $\mathcal{X} \to \mathcal{Y} \to \mathcal{Z}$ 是
$(\Sch/S)_{fppf}$ 上群胚纤维化范畴的 $1$-态射。
若 $\mathcal{X} \to \mathcal{Z}$ 与 $\mathcal{Y} \to \mathcal{Z}$
可由代数空间表示且为 \'etale，则
$\mathcal{X} \to \mathcal{Y}$ 也具有这些性质。
\end{lemma}

\begin{proof}
设 $\mathcal{U}$ 是 $S$ 上可表示的群胚纤维化范畴，并设
$f : \mathcal{U} \to \mathcal{Y}$ 是 $1$-态射。我们必须证明
$\mathcal{X} \times_\mathcal{Y} \mathcal{U}$ 可由代数空间表示，
且在 $\mathcal{U}$ 上为 \'etale。考虑复合
$h : \mathcal{U} \to \mathcal{Z}$。则
$$
\mathcal{X} \times_\mathcal{Z} \mathcal{U}
\longrightarrow
\mathcal{Y} \times_\mathcal{Z} \mathcal{U}
$$
% P11-E0167
是两个群胚纤维化范畴之间的 $1$-态射；这两个范畴都可由代数空间表示，
且在 $\mathcal{U}$ 上均为 \'etale。因此，由空间的性质，引理
\ref{spaces-properties-lemma-etale-permanence}，该态射由代数空间的
\'etale 态射表示。最后，态射 $f$ 诱导态射
$\mathcal{U} \to \mathcal{Y} \times_\mathcal{Z} \mathcal{U}$，而且有
$$
\mathcal{X} \times_\mathcal{Y} \mathcal{U} =
(\mathcal{X} \times_\mathcal{Z} \mathcal{U})
\times_{(\mathcal{Y} \times_\mathcal{Z} \mathcal{U})}
\mathcal{U}.
$$
\end{proof}

\begin{lemma}
\label{criteria-lemma-stack-in-setoids-descent}
设 $\mathcal{X}, \mathcal{Y}, \mathcal{Z}$ 是
$(\Sch/S)_{fppf}$ 上的群胚叠，并设
$\mathcal{X} \to \mathcal{Y}$ 与 $\mathcal{Z} \to \mathcal{Y}$
是 $1$-态射。若
\begin{enumerate}
\item $\mathcal{Y}$、$\mathcal{Z}$ 分别可由 $S$ 上的代数空间
$Y$、$Z$ 表示；
% P11-E0168
\item 相联系的代数空间态射 $Y \to Z$ 满、平坦且局部有限表示；
\item $\mathcal{Y} \times_\mathcal{Z} \mathcal{X}$ 是集合胚叠；
\end{enumerate}
则 $\mathcal{X}$ 是集合胚叠。
\end{lemma}

\begin{proof}
这是《叠》引理 \ref{stacks-lemma-stack-in-setoids-descent} 的特例。
\end{proof}

\noindent
下一个引理是《代数叠》引理
\ref{algebraic-lemma-smooth-surjective-morphism-implies-algebraic}
的类似物，并将被更强的定理 \ref{criteria-theorem-bootstrap} 所取代。

\begin{lemma}
\label{criteria-lemma-flat-finite-presentation-surjective-diagonal}
设 $S$ 是概形，且 $u : \mathcal{U} \to \mathcal{X}$ 是
$(\Sch/S)_{fppf}$ 上群胚叠的 $1$-态射。若
\begin{enumerate}
\item $\mathcal{U}$ 可由代数空间表示；
\item $u$ 可由代数空间表示、满、平坦且局部有限表示；
\end{enumerate}
则
% P11-E0169
$\Delta : \mathcal{X} \to \mathcal{X} \times \mathcal{X}$
可由代数空间表示。
\end{lemma}

\begin{proof}
给定 $S$ 上的两个概形 $T_1$、$T_2$，
% P11-E0170
把相联系的可表示纤维范畴记为
$\mathcal{T}_i = (\Sch/T_i)_{fppf}$。设给定 $1$-态射
$f_i : \mathcal{T}_i \to \mathcal{X}$。由《代数叠》引理
\ref{algebraic-lemma-representable-diagonal}，只需证明 $2$-纤维积
$\mathcal{T}_1 \times_\mathcal{X} \mathcal{T}_2$
可由代数空间表示。由《叠》引理
\ref{stacks-lemma-2-fibre-product-stacks-in-setoids-over-stack-in-groupoids}
，它无论如何都是集合胚叠。因此，
$\mathcal{T}_1 \times_\mathcal{X} \mathcal{T}_2$ 对应于
$(\Sch/S)_{fppf}$ 上的某个层 $F$；参见《叠》引理
\ref{stacks-lemma-stack-in-setoids-characterize}。
设 $U$ 是表示 $\mathcal{U}$ 的代数空间。由假设，
$$
\mathcal{T}_i' = \mathcal{U} \times_{u, \mathcal{X}, f_i} \mathcal{T}_i
$$
可由 $S$ 上的代数空间 $T'_i$ 表示。因此，
$\mathcal{T}_1' \times_\mathcal{U} \mathcal{T}_2'$ 可由代数空间
$T'_1 \times_U T'_2$ 表示。考虑交换图
$$
\xymatrix{
&
\mathcal{T}_1 \times_{\mathcal X} \mathcal{T}_2 \ar[rr]\ar'[d][dd] & &
\mathcal{T}_1 \ar[dd] \\
\mathcal{T}_1' \times_\mathcal{U} \mathcal{T}_2' \ar[ur]\ar[rr]\ar[dd] & &
\mathcal{T}_1' \ar[ur]\ar[dd] \\
&
\mathcal{T}_2 \ar'[r][rr] & &
\mathcal X \\
\mathcal{T}_2' \ar[rr]\ar[ur] & &
\mathcal{U} \ar[ur] }
$$
在此图中，下方、右侧、后侧和前侧的方形都是 $2$-纤维积。于是形式论证表明，
$\mathcal{T}_1' \times_\mathcal{U} \mathcal{T}_2' \to
\mathcal{T}_1 \times_{\mathcal X} \mathcal{T}_2$
是 $\mathcal{U} \to \mathcal{X}$ 的“基变换”；更确切地，图
$$
\xymatrix{
\mathcal{T}_1' \times_\mathcal{U} \mathcal{T}_2' \ar[d] \ar[r] &
\mathcal{U} \ar[d] \\
\mathcal{T}_1 \times_{\mathcal X} \mathcal{T}_2 \ar[r] &
\mathcal{X}
}
$$
是一个 $2$-纤维积方形。因此，$T'_1 \times_U T'_2 \to F$
可由代数空间表示、平坦、局部有限表示且为满；参见《代数叠》引理
\ref{algebraic-lemma-map-fibred-setoids-representable-algebraic-spaces},
\ref{algebraic-lemma-base-change-representable-by-spaces},
\ref{algebraic-lemma-map-fibred-setoids-property} 和
\ref{algebraic-lemma-base-change-representable-transformations-property}.
所以，由《自举》定理 \ref{bootstrap-theorem-final-bootstrap}，
$F$ 是代数空间，结论得证。
\end{proof}

\begin{lemma}
\label{criteria-lemma-second-diagonal}
设 $\mathcal{X}$ 是 $(\Sch/S)_{fppf}$ 上的群胚纤维化范畴。
下列条件等价：
\begin{enumerate}
\item $\Delta_\Delta : \mathcal{X} \to
\mathcal{X} \times_{\mathcal{X} \times \mathcal{X}} \mathcal{X}$
可由代数空间表示；
\item 对每个 $1$-态射 $\mathcal{V} \to \mathcal{X} \times \mathcal{X}$，
若 $\mathcal{V}$（由概形）可表示，则纤维积
$\mathcal{Y} =
\mathcal{X} \times_{\Delta, \mathcal{X} \times \mathcal{X}} \mathcal{V}$
的对角态射可由代数空间表示。
\end{enumerate}
\end{lemma}

\begin{proof}
尽管这略费思量，但完全是形式的。回顾
$\mathcal{X} \times_{\mathcal{X} \times \mathcal{X}} \mathcal{X} =
\mathcal{I}_\mathcal{X}$ 是 $\mathcal{X}$ 的惯性，并且
$\Delta_\Delta$ 是 $\mathcal{I}_\mathcal{X}$ 的恒等截面；参见
《范畴》第 \ref{categories-section-inertia} 节。因此，条件 (1) 的含义如下：
给定概形 $V$、$\mathcal{X}$ 在 $V$ 上的对象 $x$ 以及
$\mathcal{X}_V$ 的态射 $\alpha : x \to x$，条件
“$\alpha = \text{id}_x$”在 $V$ 上定义一个代数空间。
（换言之，存在代数空间的单态射 $W \to V$，使概形态射
$f : T \to V$ 通过 $W$ 分解，当且仅当
$f^*\alpha = \text{id}_{f^*x}$。）

\medskip\noindent
另一方面，设 $V$ 是概形，$x,y$ 是 $\mathcal{X}$ 在 $V$ 上的对象。
于是 $(x,y)$ 定义态射
$\mathcal{V} = (\Sch/V)_{fppf} \to \mathcal{X} \times \mathcal{X}$.
再设 $h : V' \to V$ 是概形态射，并设
$\alpha : h^*x \to h^*y$ 和 $\beta : h^*x \to h^*y$ 是
$\mathcal{X}_{V'}$ 的态射。于是 $(\alpha, \beta)$ 定义态射
% P11-E0171: frozen source appears to use V rather than V' below and again in the following sentence; preserved canonically.
$\mathcal{V}' = (\Sch/V)_{fppf} \to \mathcal{Y} \times \mathcal{Y}$.
条件 (2) 现在断言（对上述任意选择），条件“$\alpha = \beta$”
在 $V$ 上定义一个代数空间。

\medskip\noindent
为证明等价性，给定 (2) 中的 $(\alpha,\beta)$，由 (1) 可知条件
“$\alpha^{-1} \circ \beta = \text{id}_{h^*x}$”定义一个代数空间。
取 $h = \text{id}_V$ 且 $\beta = \text{id}_x$，即得
(2) $\Rightarrow$ (1)。
\end{proof}











\section{在对象上保极限}
\label{criteria-section-limit-preserving}

\noindent
设 $S$ 是概形，$p : \mathcal{X} \to \mathcal{Y}$ 是
$(\Sch/S)_{fppf}$ 上群胚纤维化范畴的 $1$-态射。如果下列条件成立，
则称 $p$ {\it 在对象上保极限}：给定由下列各项组成的任意数据
\begin{enumerate}
\item 仿射概形 $U = \lim_{i \in I} U_i$，它写成 $S$ 上仿射概形
$U_i$ 的有向极限；
\item 对某个 $i$，$\mathcal{Y}$ 在 $U_i$ 上的对象 $y_i$；
\item $\mathcal{X}$ 在 $U$ 上的对象 $x$；以及
\item 同构 $\gamma : p(x) \to y_i|_U$，
\end{enumerate}
则存在 $i' \geq i$、$\mathcal{X}$ 在 $U_{i'}$ 上的对象 $x_{i'}$、同构
$\beta : x_{i'}|_U \to x$ 以及同构
$\gamma_{i'} : p(x_{i'}) \to y_i|_{U_{i'}}$，使得
\begin{equation}
\label{criteria-equation-limit-preserving}
\vcenter{
\xymatrix{
p(x_{i'}|_U) \ar[d]_{p(\beta)} \ar[rr]_{\gamma_{i'}|_U} & &
(y_i|_{U_{i'}})|_U \ar@{=}[d] \\
p(x) \ar[rr]^\gamma & & y_i|_U
}
}
\end{equation}
交换。在这种情形下，我们称“$(i',x_{i'},\beta,\gamma_{i'})$
是由数据 (1)、(2)、(3)、(4) 所提出问题的一个{\it 解}”。
这一定义的动机来自《空间的极限》引理
\ref{spaces-limits-lemma-characterize-relative-limit-preserving}。

\begin{lemma}
\label{criteria-lemma-base-change-limit-preserving}
设 $p : \mathcal{X} \to \mathcal{Y}$ 和
$q : \mathcal{Z} \to \mathcal{Y}$ 是 $(\Sch/S)_{fppf}$ 上
群胚纤维化范畴的 $1$-态射。如果
$p : \mathcal{X} \to \mathcal{Y}$ 在对象上保极限，则 $p$ 沿 $q$ 的基变换
$p' : \mathcal{X} \times_\mathcal{Y} \mathcal{Z} \to \mathcal{Z}$
也在对象上保极限。
\end{lemma}

\begin{proof}
这是形式的。设 $U = \lim_{i \in I} U_i$ 是 $S$ 上仿射概形
$U_i$ 的有向极限；对某个 $i$，设 $z_i$ 是 $\mathcal{Z}$ 在
$U_i$ 上的对象，$w$ 是 $\mathcal{X} \times_\mathcal{Y} \mathcal{Z}$
在 $U$ 上的对象，而
$\delta : p'(w) \to z_i|_U$ 是同构。可以写成
$w = (U, x, z, \alpha)$，其中 $x$ 是 $\mathcal{X}$ 在 $U$ 上的某个对象，
$z$ 是 $\mathcal{Z}$ 在 $U$ 上的对象，而
$\alpha : p(x) \to q(z)$ 是同构。注意 $p'(w)=z$，故
$\delta : z \to z_i|_U$。令 $y_i=q(z_i)$ 且
$\gamma = q(\delta) \circ \alpha : p(x) \to y_i|_U$。
由于 $p$ 在对象上保极限，存在 $i'\geq i$ 和 $\mathcal{X}$ 在
$U_{i'}$ 上的对象 $x_{i'}$，以及同构
$\beta : x_{i'}|_U \to x$ 和
$\gamma_{i'} : p(x_{i'}) \to y_i|_{U_{i'}}$，使
(\ref{criteria-equation-limit-preserving}) 交换。于是考虑
$\mathcal{X} \times_\mathcal{Y} \mathcal{Z}$ 在 $U_{i'}$ 上的对象
$w_{i'} = (U_{i'}, x_{i'}, z_i|_{U_{i'}}, \gamma_{i'})$，并定义同构
$$
w_{i'}|_U = (U, x_{i'}|_U, z_i|_U, \gamma_{i'}|_U)
\xrightarrow{(\beta, \delta^{-1})}
(U, x, z, \alpha) = w
$$
以及
$$
p'(w_{i'}) = z_i|_{U_{i'}} \xrightarrow{\text{id}} z_i|_{U_{i'}}.
$$
二者合在一起给出该问题的一个解。
\end{proof}

\begin{lemma}
\label{criteria-lemma-composition-limit-preserving}
设 $p : \mathcal{X} \to \mathcal{Y}$ 和
$q : \mathcal{Y} \to \mathcal{Z}$ 是 $(\Sch/S)_{fppf}$ 上
群胚纤维化范畴的 $1$-态射。如果 $p$ 和 $q$ 均在对象上保极限，
则复合 $q \circ p$ 也在对象上保极限。
\end{lemma}

\begin{proof}
这是形式的。设 $U = \lim_{i \in I} U_i$ 是 $S$ 上仿射概形
$U_i$ 的有向极限；对某个 $i$，设 $z_i$ 是 $\mathcal{Z}$ 在
$U_i$ 上的对象，$x$ 是 $\mathcal{X}$ 在 $U$ 上的对象，而
$\gamma : q(p(x)) \to z_i|_U$ 是同构。由于 $q$ 在对象上保极限，
存在 $i' \geq i$、$\mathcal{Y}$ 在 $U_{i'}$ 上的对象 $y_{i'}$、同构
$\beta : y_{i'}|_U \to p(x)$ 以及同构
$\gamma_{i'} : q(y_{i'}) \to z_i|_{U_{i'}}$，使
(\ref{criteria-equation-limit-preserving}) 交换。由于 $p$ 在对象上保极限，
存在 $i'' \geq i'$、$\mathcal{X}$ 在 $U_{i''}$ 上的对象 $x_{i''}$、同构
$\beta' : x_{i''}|_U \to x$ 以及同构
$\gamma'_{i''} : p(x_{i''}) \to y_{i'}|_{U_{i''}}$，使
(\ref{criteria-equation-limit-preserving}) 交换。所求解取为 $U_{i''}$ 上的
$x_{i''}$，并配以同构
$$
q(p(x_{i''})) \xrightarrow{q(\gamma'_{i''})}
q(y_{i'})|_{U_{i''}} \xrightarrow{\gamma_{i'}|_{U_{i''}}}
z_i|_{U_{i''}}
$$
和同构 $\beta' : x_{i''}|_U \to x$。我们略去
(\ref{criteria-equation-limit-preserving}) 交换性的验证。
\end{proof}

\begin{lemma}
\label{criteria-lemma-representable-by-spaces-limit-preserving}
设 $p : \mathcal{X} \to \mathcal{Y}$ 是 $(\Sch/S)_{fppf}$ 上
群胚纤维化范畴的 $1$-态射。如果 $p$ 可由代数空间表示，
则下列条件等价：
\begin{enumerate}
\item $p$ 在对象上保极限；
\item $p$ 局部有限表示（参见《代数叠》定义
\ref{algebraic-definition-relative-representable-property}）。
\end{enumerate}
\end{lemma}

\begin{proof}
假设 (2)。设 $U = \lim_{i \in I} U_i$ 是 $S$ 上仿射概形
$U_i$ 的有向极限；对某个 $i$，设 $y_i$ 是 $\mathcal{Y}$ 在
$U_i$ 上的对象，$x$ 是 $\mathcal{X}$ 在 $U$ 上的对象，而
$\gamma : p(x) \to y_i|_U$ 是同构。记 $X_{y_i}$ 为表示 $2$-纤维积
$$
(\Sch/U_i)_{fppf} \times_{y_i, \mathcal{Y}, p} \mathcal{X}.
$$
的 $U_i$ 上的代数空间。注意，
$\xi = (U, U \to U_i, x, \gamma^{-1})$ 定义此 $2$-纤维积在
$U$ 上的一个对象。由 $2$-Yoneda 引理，$\xi$ 对应于 $U_i$ 上的态射
$f_\xi : U \to X_{y_i}$。由《空间的极限》命题
\ref{spaces-limits-proposition-characterize-locally-finite-presentation}
，存在 $i'\geq i$ 和态射 $f_{i'} : U_{i'} \to X_{y_i}$，使
$f_\xi$ 是 $f_{i'}$ 与投影态射 $U \to U_{i'}$ 的复合。
$2$-Yoneda 引理还告诉我们，$f_{i'}$ 对应于上述 $2$-纤维积在
$U_{i'}$ 上的一个对象
$\xi_{i'} = (U_{i'}, U_{i'} \to U_i, x_{i'}, \alpha)$，其到 $U$ 的限制
恢复 $\xi$。特别地，我们得到同构
% P11-E0172: frozen source uses |U rather than |_U in the following two displays; preserved canonically.
$\gamma : x_{i'}|U \to x$。注意，$\alpha : y_i|_{U_{i'}} \to p(x_{i'})$。
因此，取 $x_{i'}$、同构
$\gamma : x_{i'}|U \to x$ 以及同构
$\beta = \alpha^{-1} : p(x_{i'}) \to y_i|_{U_{i'}}$，即为该问题的一个解。

\medskip\noindent
假设 (1)。选取概形 $T$ 和 $1$-态射
$y : (\Sch/T)_{fppf} \to \mathcal{Y}$。设 $X_y$ 是表示 $2$-纤维积
$(\Sch/T)_{fppf} \times_{y, \mathcal{Y}, p} \mathcal{X}$.
的 $T$ 上的代数空间。我们必须证明 $X_y \to T$ 局部有限表示。
为此，将使用《空间的极限》评注
\ref{spaces-limits-remark-limit-preserving} 中的判据。考虑仿射概形
$U = \lim_{i \in I} U_i$，它写成 $T$ 上仿射概形的有向极限。
任取 $i \in I$ 并令 $y_i = y|_{U_i}$。另以 $i'$ 表示 $I$ 中
大于或等于 $i$ 的元素。由 $2$-Yoneda 引理，$T$ 上的态射
$U \to X_y$ 与二元组 $(x,\alpha)$ 的同构类一一对应，其中 $x$ 是
$\mathcal{X}$ 在 $U$ 上的对象，而
$\alpha : y|_U \to p(x)$ 是同构。当然，在取逆的意义下，给出
$\alpha$ 等同于给出同构 $\gamma : p(x) \to y_i|_U$。
$T$ 上的态射 $U_{i'} \to X_y$ 情形类似。因此，(1) 保证典范映射
$$
\colim_{i' \geq i} X_y(U_{i'}) \longrightarrow X_y(U)
$$
在此情形下为满。由《空间的极限》引理
\ref{spaces-limits-lemma-surjection-is-enough}，$X_y \to T$ 局部有限表示。
\end{proof}

\begin{lemma}
\label{criteria-lemma-open-immersion-limit-preserving}
设 $p : \mathcal{X} \to \mathcal{Y}$ 是 $(\Sch/S)_{fppf}$ 上
群胚纤维化范畴的 $1$-态射。假设 $p$ 可由代数空间表示，且为开浸入。
则 $p$ 在对象上保极限。
\end{lemma}

\begin{proof}
这由引理 \ref{criteria-lemma-representable-by-spaces-limit-preserving}，并经由
《代数叠》引理
\ref{algebraic-lemma-representable-transformations-property-implication}
中的一般原理，再结合下述事实推出：代数空间的开浸入局部有限表示；
参见《空间的态射》引理
\ref{spaces-morphisms-lemma-open-immersion-locally-finite-presentation}。
\end{proof}

\noindent
设 $S$ 是概形。在下一个引理中，我们需要 $S$ 上代数空间 $X$ 的
{\it 大小}这一概念。确切地，给定基数 $\kappa$，我们称 $X$ 满足
$\text{size}(X) \leq \kappa$，当且仅当存在概形 $U$，使
$\text{size}(U) \leq \kappa$（参见《集合》第
\ref{sets-section-categories-schemes} 节），并且存在满的平展态射 $U \to X$。

\begin{lemma}
\label{criteria-lemma-check-representable-limit-preserving}
设 $S$ 是概形。对某个 $T \in \Ob((\Sch/S)_{fppf})$，令
$\kappa = \text{size}(T)$。设
$f : \mathcal{X} \to \mathcal{Y}$ 是 $(\Sch/S)_{fppf}$ 上
群胚纤维化范畴的 $1$-态射，并满足
\begin{enumerate}
\item $\mathcal{Y} \to (\Sch/S)_{fppf}$ 在对象上保极限；
\item 对 $S$ 上局部有限表示的仿射概形 $V$ 和
$y \in \Ob(\mathcal{Y}_V)$，纤维积
$(\Sch/V)_{fppf} \times_{y, \mathcal{Y}} \mathcal{X}$ 可由大小
$\leq \kappa$ 的代数空间表示\footnote{忽略集合论问题的读者可以删去大小条件。}；
\item $\mathcal{X}$ 和 $\mathcal{Y}$ 都是 Zariski 拓扑的叠。
\end{enumerate}
则 $f$ 可由代数空间表示。
\end{lemma}

\begin{proof}
设 $V$ 是 $S$ 上的概形且 $y \in \mathcal{Y}_V$。我们必须证明
$(\Sch/V)_{fppf} \times_{y, \mathcal{Y}} \mathcal{X}$ 可由代数空间表示。

\medskip\noindent
情形 I：$V$ 为仿射概形，并映入仿射开集
$\Spec(\Lambda) \subset S$。于是可写 $V = \lim V_i$，其中每个 $V_i$
都是 $\Spec(\Lambda)$ 上有限表示的仿射概形；参见《代数》引理
\ref{algebra-lemma-ring-colimit-fp}。由假设 (1)，对某个 $i$，$y$
来自 $V_i$ 上的对象 $y_i$。由假设 (3)，纤维积
$(\Sch/V_i)_{fppf} \times_{y_i, \mathcal{Y}} \mathcal{X}$ 可由代数空间
$Z_i$ 表示。于是
$(\Sch/V)_{fppf} \times_{y, \mathcal{Y}} \mathcal{X}$ 可由
% P11-E0173: frozen source writes Z below after defining Z_i; likely Z_i is intended, but the canonical formula is preserved.
$Z \times_{V_i} V$ 表示。

\medskip\noindent
情形 II：$V$ 为一般概形。选取仿射开覆盖
$V = \bigcup_{i \in I} V_i$，使每个 $V_i$ 都映入 $S$ 的一个仿射开集。
先断言
$\mathcal{Z} = (\Sch/V)_{fppf} \times_{y, \mathcal{Y}} \mathcal{X}$
是 Zariski 拓扑的集合胚叠。事实上，由《叠》引理
\ref{stacks-lemma-2-product-stacks-in-groupoids}，它是 Zariski 拓扑的
群胚叠。再设 $z$ 是 $\mathcal{Z}$ 在概形 $T$ 上的对象。
记 $g : T \to V$ 为与 $z$ 到 $(\Sch/V)_{fppf}$ 的投影相对应的态射。
考虑 Zariski 层
$\mathit{I} = \mathit{Isom}_{\mathcal{Z}}(z, z)$。由情形 I，
$\mathit{I}|_{g^{-1}(V_i)} = *$（单点层）。故
% P11-E0174: frozen source changes mathit I to mathcal I here without defining a new object; preserved canonically.
$\mathcal{I} = *$。所以 $\mathcal{Z}$ 纤维化于集合胚。为完成证明，
须证明 Zariski 层 $Z : T \mapsto \Ob(\mathcal{Z}_T)/\cong$ 是代数空间；
参见《代数叠》引理
\ref{algebraic-lemma-characterize-representable-by-space}。
存在映射 $p : Z \to V$（函子之间的变换），而由情形 I，
$Z_i = p^{-1}(V_i)$ 是代数空间。态射 $Z_i \to Z$ 可由开浸入表示，
且 $\coprod Z_i \to Z$ 在 Zariski 拓扑下为满。因此，由《自举》引理
\ref{bootstrap-lemma-glueing-sheaves}，$Z$ 是 fppf 拓扑的层。
于是可应用《空间》引理 \ref{spaces-lemma-glueing-algebraic-spaces}，
从而 $Z$ 是代数空间\footnote{为验证该引理中的集合论条件，论证如下。
首先选取开覆盖，使 $|I| \leq \text{size}(V)$。其次，选取大小
$\leq \max(\kappa, \text{size}(V))$ 的概形 $U_i$ 和满平展态射
$U_i \to Z_i$；由假设 (2) 和《集合》引理
\ref{sets-lemma-bound-size-fibre-product} 可以做到这一点（略去细节）。
于是《集合》引理 \ref{sets-lemma-what-is-in-it} 表明
$\coprod U_i$ 是 $(\Sch/S)_{fppf}$ 的对象。因此，由《空间》引理
\ref{spaces-lemma-coproduct-algebraic-spaces}，$\coprod Z_i$ 是代数空间。}。
\end{proof}

\begin{lemma}
\label{criteria-lemma-check-property-limit-preserving}
设 $S$ 是概形，$f : \mathcal{X} \to \mathcal{Y}$ 是
$(\Sch/S)_{fppf}$ 上群胚纤维化范畴的 $1$-态射。设 $\mathcal{P}$ 是
《代数叠》定义 \ref{algebraic-definition-relative-representable-property}
中所述的代数空间态射的性质。如果
\begin{enumerate}
\item $f$ 可由代数空间表示；
\item $\mathcal{Y} \to (\Sch/S)_{fppf}$ 在对象上保极限；
\item 对 $S$ 上局部有限表示的仿射概形 $V$ 和
$y \in \mathcal{Y}_V$，所得代数空间态射 $f_y : F_y \to V$
具有性质 $\mathcal{P}$；参见《代数叠》公式
(\ref{algebraic-equation-representable-by-algebraic-spaces})。
\end{enumerate}
则 $f$ 具有性质 $\mathcal{P}$。
\end{lemma}

\begin{proof}
设 $V$ 是 $S$ 上的概形且 $y \in \mathcal{Y}_V$。我们必须证明
$F_y \to V$ 具有性质 $\mathcal{P}$。由于 $\mathcal{P}$ 在基底上为
fppf 局部的，可以假设 $V$ 是映入仿射开集
$\Spec(\Lambda) \subset S$ 的仿射概形。于是可写
$V = \lim V_i$，其中每个 $V_i$ 都是 $\Spec(\Lambda)$ 上有限表示的
仿射概形；参见《代数》引理 \ref{algebra-lemma-ring-colimit-fp}。
由假设 (2)，对某个 $i$，$y$ 来自 $V_i$ 上的对象 $y_i$。
由假设 (3)，态射 $F_{y_i} \to V_i$ 具有性质 $\mathcal{P}$。
由于 $\mathcal{P}$ 在任意基变换下稳定，并且
$F_y = F_{y_i} \times_{V_i} V$，故 $F_y \to V$ 如愿具有性质
$\mathcal{P}$。
\end{proof}



\section{在对象上形式光滑}
\label{criteria-section-formally-smooth}

\noindent
设 $S$ 是概形，$p : \mathcal{X} \to \mathcal{Y}$ 是
$(\Sch/S)_{fppf}$ 上群胚纤维化范畴的 $1$-态射。如果下列条件成立，
则称 $p$ {\it 在对象上形式光滑}：给定由下列各项组成的任意数据
\begin{enumerate}
\item $S$ 上仿射概形的一阶加厚 $U \subset U'$；
\item $\mathcal{Y}$ 在 $U'$ 上的对象 $y'$；
\item $\mathcal{X}$ 在 $U$ 上的对象 $x$；以及
\item 同构 $\gamma : p(x) \to y'|_U$，
\end{enumerate}
则存在 $\mathcal{X}$ 在 $U'$ 上的对象 $x'$，并有同构
$\beta : x'|_U \to x$ 和同构 $\gamma' : p(x') \to y'$，使得
\begin{equation}
\label{criteria-equation-formally-smooth}
\vcenter{
\xymatrix{
p(x'|_U) \ar[d]_{p(\beta)} \ar[rr]_{\gamma'|_U} & &
y'|_U \ar@{=}[d] \\
p(x) \ar[rr]^\gamma & & y'|_U
}
}
\end{equation}
交换。在这种情形下，我们称“$(x',\beta,\gamma')$ 是由数据
(1)、(2)、(3)、(4) 所提出问题的一个{\it 解}”。

\begin{lemma}
\label{criteria-lemma-base-change-formally-smooth}
设 $p : \mathcal{X} \to \mathcal{Y}$ 和
$q : \mathcal{Z} \to \mathcal{Y}$ 是 $(\Sch/S)_{fppf}$ 上
群胚纤维化范畴的 $1$-态射。如果
$p : \mathcal{X} \to \mathcal{Y}$ 在对象上形式光滑，则 $p$ 沿 $q$ 的基变换
$p' : \mathcal{X} \times_\mathcal{Y} \mathcal{Z} \to \mathcal{Z}$
也在对象上形式光滑。
\end{lemma}

\begin{proof}
这是形式的。设 $U \subset U'$ 是 $S$ 上仿射概形的一阶加厚，
$z'$ 是 $\mathcal{Z}$ 在 $U'$ 上的对象，$w$ 是
$\mathcal{X} \times_\mathcal{Y} \mathcal{Z}$ 在 $U$ 上的对象，而
$\delta : p'(w) \to z'|_U$ 是同构。可以写
$w = (U, x, z, \alpha)$，其中 $x$ 是 $\mathcal{X}$ 在 $U$ 上的某个对象，
$z$ 是 $\mathcal{Z}$ 在 $U$ 上的对象，而
$\alpha : p(x) \to q(z)$ 是同构。注意 $p'(w)=z$，故
% P11-E0175: frozen source writes z|_U below although delta was defined with target z'|_U; preserved canonically.
$\delta : z \to z|_U$。令 $y'=q(z')$ 且
$\gamma = q(\delta) \circ \alpha : p(x) \to y'|_U$。
由于 $p$ 在对象上形式光滑，存在 $\mathcal{X}$ 在 $U'$ 上的对象 $x'$，
以及同构 $\beta : x'|_U \to x$ 和 $\gamma' : p(x') \to y'$，使
(\ref{criteria-equation-formally-smooth}) 交换。于是考虑
% P11-E0176: frozen source names the following lifted object w, while subsequent formulas use w'; preserved canonically.
$\mathcal{X} \times_\mathcal{Y} \mathcal{Z}$ 在 $U'$ 上的对象
$w = (U', x', z', \gamma')$，并定义同构
$$
w'|_U = (U, x'|_U, z'|_U, \gamma'|_U)
\xrightarrow{(\beta, \delta^{-1})}
(U, x, z, \alpha) = w
$$
以及
$$
p'(w') = z' \xrightarrow{\text{id}} z'.
$$
二者合在一起给出该问题的一个解。
\end{proof}

\begin{lemma}
\label{criteria-lemma-composition-formally-smooth}
设 $p : \mathcal{X} \to \mathcal{Y}$ 和
$q : \mathcal{Y} \to \mathcal{Z}$ 是 $(\Sch/S)_{fppf}$ 上
群胚纤维化范畴的 $1$-态射。如果 $p$ 和 $q$ 均在对象上形式光滑，
则复合 $q \circ p$ 也在对象上形式光滑。
\end{lemma}

\begin{proof}
这是形式的。设 $U \subset U'$ 是 $S$ 上仿射概形的一阶加厚，
$z'$ 是 $\mathcal{Z}$ 在 $U'$ 上的对象，$x$ 是 $\mathcal{X}$ 在
$U$ 上的对象，而 $\gamma : q(p(x)) \to z'|_U$ 是同构。
由于 $q$ 在对象上形式光滑，存在 $\mathcal{Y}$ 在 $U'$ 上的对象
$y'$、同构 $\beta : y'|_U \to p(x)$ 以及同构
$\gamma' : q(y') \to z'$，使 (\ref{criteria-equation-formally-smooth}) 交换。
由于 $p$ 在对象上形式光滑，存在 $\mathcal{X}$ 在 $U'$ 上的对象
$x'$、同构 $\beta' : x'|_U \to x$ 以及同构
$\gamma'' : p(x') \to y'$，使 (\ref{criteria-equation-formally-smooth}) 交换。
所求解取为 $U'$ 上的 $x'$，并配以同构
$$
q(p(x')) \xrightarrow{q(\gamma'')} q(y') \xrightarrow{\gamma'} z'
$$
和同构 $\beta' : x'|_U \to x$。我们略去
(\ref{criteria-equation-formally-smooth}) 交换性的验证。
\end{proof}

\noindent
注意，代数空间的形式光滑态射类在任意基变换下稳定，并且在目标上相对于
fpqc 拓扑是局部的；参见《空间态射进阶》引理
\ref{spaces-more-morphisms-lemma-base-change-formally-smooth} 和
\ref{spaces-more-morphisms-lemma-descending-property-formally-smooth}。
因此，下一个引理中的条件 (2) 是有意义的。

\begin{lemma}
\label{criteria-lemma-representable-by-spaces-formally-smooth}
设 $p : \mathcal{X} \to \mathcal{Y}$ 是 $(\Sch/S)_{fppf}$ 上
群胚纤维化范畴的 $1$-态射。如果 $p$ 可由代数空间表示，
则下列条件等价：
\begin{enumerate}
\item $p$ 在对象上形式光滑；
\item $p$ 形式光滑（参见《代数叠》定义
\ref{algebraic-definition-relative-representable-property}）。
\end{enumerate}
\end{lemma}

\begin{proof}
假设 (2)。设 $U \subset U'$ 是 $S$ 上仿射概形的一阶加厚，
$y'$ 是 $\mathcal{Y}$ 在 $U'$ 上的对象，$x$ 是 $\mathcal{X}$ 在
$U$ 上的对象，而 $\gamma : p(x) \to y'|_U$ 是同构。
记 $X_{y'}$ 为表示 $2$-纤维积
$$
(\Sch/U')_{fppf} \times_{y', \mathcal{Y}, p} \mathcal{X}.
$$
的 $U'$ 上的代数空间。注意，
$\xi = (U, U \to U', x, \gamma^{-1})$ 定义此 $2$-纤维积在 $U$ 上的
一个对象。由 $2$-Yoneda 引理，$\xi$ 对应于 $U'$ 上的态射
$f_\xi : U \to X_{y'}$。按假设，$X_{y'} \to U'$ 形式光滑，故存在态射
$f' : U' \to X_{y'}$，使 $f_\xi$ 是 $f'$ 与态射 $U \to U'$ 的复合。
$2$-Yoneda 引理还告诉我们，$f'$ 对应于上述 $2$-纤维积在 $U'$ 上的对象
$\xi' = (U', U' \to U', x', \alpha)$，其到 $U$ 的限制恢复 $\xi$。
特别地，我们得到同构
% P11-E0178: frozen source uses |U rather than |_U in the following two displays; preserved canonically.
$\gamma : x'|U \to x$。注意，$\alpha : y' \to p(x')$。
因此，取 $x'$、同构 $\gamma : x'|U \to x$ 以及同构
$\beta = \alpha^{-1} : p(x') \to y'$，即为该问题的一个解。

\medskip\noindent
假设 (1)。选取概形 $T$ 和 $1$-态射
$y : (\Sch/T)_{fppf} \to \mathcal{Y}$。设 $X_y$ 是表示 $2$-纤维积
$(\Sch/T)_{fppf} \times_{y, \mathcal{Y}, p} \mathcal{X}$.
的 $T$ 上的代数空间。我们必须证明 $X_y \to T$ 形式光滑。
因而只需证明：给定 $T$ 上仿射概形的一阶加厚 $U \subset U'$，则
% P11-E0177: frozen source maps X_y(U') to itself here and again below; the formal smoothness criterion normally maps X_y(U') to X_y(U). Preserved canonically.
$X_y(U') \to X_y(U')$ 为满（这里的态射取自 $T$ 上代数空间的范畴）。
令 $y'=y|_{U'}$。由 $2$-Yoneda 引理，$T$ 上的态射
$U \to X_y$ 与二元组 $(x,\alpha)$ 的同构类一一对应，其中 $x$ 是
$\mathcal{X}$ 在 $U$ 上的对象，而 $\alpha : y|_U \to p(x)$ 是同构。
当然，在取逆的意义下，给出 $\alpha$ 等同于给出同构
$\gamma : p(x) \to y'|_U$。$T$ 上的态射 $U' \to X_y$ 情形类似。
因此，(1) 保证 $X_y(U') \to X_y(U')$ 在此情形下为满，结论得证。
\end{proof}







\section{在对象上为满}
\label{criteria-section-formally-surjective}

\noindent
设 $S$ 是概形，$p : \mathcal{X} \to \mathcal{Y}$ 是
$(\Sch/S)_{fppf}$ 上群胚纤维化范畴的 $1$-态射。如果下列条件成立，
则称 $p$ {\it 在对象上为满}：给定由下列各项组成的任意数据
\begin{enumerate}
\item $S$ 上的域 $k$；以及
\item $\mathcal{Y}$ 在 $\Spec(k)$ 上的对象 $y$，
\end{enumerate}
则存在 $S$ 上的域扩张 $K/k$ 以及 $\mathcal{X}$ 在 $\Spec(K)$ 上的对象
$x$，使 $p(x) \cong y|_{\Spec(K)}$。

\begin{lemma}
\label{criteria-lemma-base-change-surjective}
设 $p : \mathcal{X} \to \mathcal{Y}$ 和
$q : \mathcal{Z} \to \mathcal{Y}$ 是 $(\Sch/S)_{fppf}$ 上
群胚纤维化范畴的 $1$-态射。如果
$p : \mathcal{X} \to \mathcal{Y}$ 在对象上为满，则 $p$ 沿 $q$ 的基变换
$p' : \mathcal{X} \times_\mathcal{Y} \mathcal{Z} \to \mathcal{Z}$
也在对象上为满。
\end{lemma}

\begin{proof}
这是形式的。设 $z$ 是 $\mathcal{Z}$ 在域 $k$ 上的对象。
由于 $p$ 在对象上为满，存在扩张 $K/k$、$\mathcal{X}$ 在 $K$ 上的对象
$x$ 以及同构 $\alpha : p(x) \to q(z)|_{\Spec(K)}$。于是
$w = (\Spec(K), x, z|_{\Spec(K)}, \alpha)$ 是
$\mathcal{X} \times_\mathcal{Y} \mathcal{Z}$ 在 $K$ 上的对象，且
$p'(w) = z|_{\Spec(K)}$。
\end{proof}

\begin{lemma}
\label{criteria-lemma-composition-surjective}
设 $p : \mathcal{X} \to \mathcal{Y}$ 和
$q : \mathcal{Y} \to \mathcal{Z}$ 是 $(\Sch/S)_{fppf}$ 上
群胚纤维化范畴的 $1$-态射。如果 $p$ 和 $q$ 均在对象上为满，
则复合 $q \circ p$ 也在对象上为满。
\end{lemma}

\begin{proof}
这是形式的。设 $z$ 是 $\mathcal{Z}$ 在域 $k$ 上的对象。
由于 $q$ 在对象上为满，存在域扩张 $K/k$ 和 $\mathcal{Y}$ 在 $K$ 上的对象
$y$，使
% P11-E0179: frozen source uses x below before x is introduced; z is evidently the intended object, but the canonical formula is preserved.
$q(y) \cong x|_{\Spec(K)}$。由于 $p$ 在对象上为满，存在域扩张 $L/K$
和 $\mathcal{X}$ 在 $L$ 上的对象 $x$，使
$p(x) \cong y|_{\Spec(L)}$。于是域扩张 $L/k$ 和 $\mathcal{X}$ 在
$L$ 上的对象 $x$ 如愿满足 $q(p(x)) \cong z|_{\Spec(L)}$。
\end{proof}

\begin{lemma}
\label{criteria-lemma-representable-by-spaces-surjective}
设 $p : \mathcal{X} \to \mathcal{Y}$ 是 $(\Sch/S)_{fppf}$ 上
群胚纤维化范畴的 $1$-态射。如果 $p$ 可由代数空间表示，
则下列条件等价：
\begin{enumerate}
\item $p$ 在对象上为满；
\item $p$ 为满（参见《代数叠》定义
\ref{algebraic-definition-relative-representable-property}）。
\end{enumerate}
\end{lemma}

\begin{proof}
假设 (2)。设 $k$ 是域，$y$ 是 $\mathcal{Y}$ 在 $k$ 上的对象。
记 $X_y$ 为表示 $2$-纤维积
$$
(\Sch/\Spec(k))_{fppf} \times_{y, \mathcal{Y}, p} \mathcal{X}.
$$
的 $k$ 上的代数空间。由于已假设 $p$ 为满，可知 $X_y$ 非空。
因此可找到域扩张 $K/k$ 和 $X_y$ 的一个 $K$-值点 $x$。
由 $2$-Yoneda 引理，它对应于 $\mathcal{X}$ 在 $K$ 上的对象 $x$ 以及同构
$p(x) \cong y|_{\Spec(K)}$，故 (1) 成立。

\medskip\noindent
假设 (1)。选取概形 $T$ 和 $1$-态射
$y : (\Sch/T)_{fppf} \to \mathcal{Y}$。设 $X_y$ 是表示 $2$-纤维积
$(\Sch/T)_{fppf} \times_{y, \mathcal{Y}, p} \mathcal{X}$.
的 $T$ 上的代数空间。我们必须证明 $X_y \to T$ 为满。
由《空间的态射》定义 \ref{spaces-morphisms-definition-surjective}，
必须证明 $|X_y| \to |T|$ 为满。这恰好意味着：给定 $T$ 上的域 $k$
和态射 $t : \Spec(k) \to T$，存在域扩张 $K/k$ 和态射
$x : \Spec(K) \to X_y$，使得
$$
\xymatrix{
\Spec(K) \ar[d] \ar[r]_x & X_y \ar[d] \\
\Spec(k) \ar[r]^t & T
}
$$
交换。由 $2$-Yoneda 引理，这恰好意味着必须找到 $k \subset K$ 和
$\mathcal{X}$ 在 $K$ 上的对象 $x$，使
$p(x) \cong t^*y|_{\Spec(K)}$。因此 (1) 保证了所需结论。
\end{proof}












\section{代数态射}
\label{criteria-section-algebraic}

\noindent
下列概念有时很有用。

\begin{definition}
\label{criteria-definition-algebraic}
设 $S$ 是概形，$F : \mathcal{X} \to \mathcal{Y}$ 是
$(\Sch/S)_{fppf}$ 上群胚叠的 $1$-态射。如果对每个概形 $T$ 及
$\mathcal{Y}$ 在 $T$ 上的每个对象 $\xi$，$2$-纤维积
$$
(\Sch/T)_{fppf} \times_{\xi, \mathcal{Y}} \mathcal{X}
$$
都是 $S$ 上的代数叠，则称 $F$ 是{\it 代数的}。
\end{definition}

\noindent
采用这一术语后，得到下列推广《代数叠》引理
\ref{algebraic-lemma-representable-morphism-to-algebraic} 的结果。

\begin{lemma}
\label{criteria-lemma-algebraic-morphism-to-algebraic}
设 $S$ 是概形。设 $F : \mathcal{X} \to \mathcal{Y}$ 是
$(\Sch/S)_{fppf}$ 上群胚叠的 $1$-态射。如果
\begin{enumerate}
\item $\mathcal{Y}$ 是代数叠；
\item $F$ 是代数的（见上文），
\end{enumerate}
则 $\mathcal{X}$ 是代数叠。
\end{lemma}

\begin{proof}
由假设 (1)，存在概形 $T$ 和 $\mathcal{Y}$ 在 $T$ 上的对象 $\xi$，
使相应的 $1$-态射 $\xi : (\Sch/T)_{fppf} \to \mathcal{Y}$
% P11-E0180: frozen source says "smooth an surjective"; "and" is evidently intended. Meaning is translated without changing any formula.
光滑且为满。于是
$\mathcal{U} = (\Sch/T)_{fppf} \times_{\xi, \mathcal{Y}} \mathcal{X}$
由假设 (2) 是代数叠。选取概形 $U$ 和满光滑的 $1$-态射
$(\Sch/U)_{fppf} \to \mathcal{U}$。投影
$\mathcal{U} \longrightarrow \mathcal{X}$ 是态射
$\xi : (\Sch/T)_{fppf} \to \mathcal{Y}$ 的基变换，因而满且光滑；
参见《代数叠》引理
\ref{algebraic-lemma-base-change-representable-transformations-property}。
于是复合
$(\Sch/U)_{fppf} \to \mathcal{U} \to \mathcal{X}$
作为满光滑态射的复合，仍然满且光滑；参见《代数叠》引理
\ref{algebraic-lemma-composition-representable-transformations-property}。
因此，由《代数叠》引理
\ref{algebraic-lemma-smooth-surjective-morphism-implies-algebraic}，
$\mathcal{X}$ 是代数叠。
\end{proof}

\begin{lemma}
\label{criteria-lemma-map-from-algebraic}
设 $S$ 是概形，$F : \mathcal{X} \to \mathcal{Y}$ 是
$(\Sch/S)_{fppf}$ 上群胚叠的 $1$-态射。如果 $\mathcal{X}$ 是代数叠，
且 $\Delta : \mathcal{Y} \to \mathcal{Y} \times \mathcal{Y}$
可由代数空间表示，则 $F$ 是代数的。
\end{lemma}

\begin{proof}
选取可表示群胚叠 $\mathcal{U}$ 和满光滑的 $1$-态射
$\mathcal{U} \to \mathcal{X}$。设 $T$ 是概形，$\xi$ 是 $\mathcal{Y}$
在 $T$ 上的对象。$2$-纤维积之间的态射
$$
(\Sch/T)_{fppf} \times_{\xi, \mathcal{Y}} \mathcal{U}
\longrightarrow
(\Sch/T)_{fppf} \times_{\xi, \mathcal{Y}} \mathcal{X}
$$
是 $\mathcal{U} \to \mathcal{X}$ 的基变换，故可由代数空间表示、为满且光滑；
参见《代数叠》引理
\ref{algebraic-lemma-base-change-representable-by-spaces} 和
\ref{algebraic-lemma-base-change-representable-transformations-property}。
由关于 $\mathcal{Y}$ 的对角态射的条件，此态射的源可由代数空间表示；
参见《代数叠》引理 \ref{algebraic-lemma-representable-diagonal}。
因此，由《代数叠》引理
\ref{algebraic-lemma-smooth-surjective-morphism-implies-algebraic}，
其目标是代数叠。
\end{proof}

\begin{lemma}
\label{criteria-lemma-diagonals-and-algebraic-morphisms}
设 $S$ 是概形，$F : \mathcal{X} \to \mathcal{Y}$ 是
$(\Sch/S)_{fppf}$ 上群胚叠的 $1$-态射。如果 $F$ 是代数的，且
$\Delta : \mathcal{Y} \to \mathcal{Y} \times \mathcal{Y}$
可由代数空间表示，则
$\Delta : \mathcal{X} \to \mathcal{X} \times \mathcal{X}$
可由代数空间表示。
\end{lemma}

\begin{proof}
假设 $F$ 是代数的，且
$\Delta : \mathcal{Y} \to \mathcal{Y} \times \mathcal{Y}$
可由代数空间表示。取 $S$ 上的概形 $U$ 以及 $\mathcal{X}$ 在 $U$ 上的
两个对象 $x_1,x_2$。须证明 $\mathit{Isom}(x_1, x_2)$ 是 $U$ 上的代数空间；
参见《代数叠》引理 \ref{algebraic-lemma-representable-diagonal}。
令 $y_i=F(x_i)$。存在集合层的态射
$$
f : \mathit{Isom}(x_1, x_2) \to \mathit{Isom}(y_1, y_2)
$$
，而按假设，其目标是代数空间。因此只需证明 $f$ 可由代数空间表示；
参见《自举》引理
\ref{bootstrap-lemma-representable-by-spaces-over-space}。
所以可以选取 $U$ 上的概形 $V$ 和同构
$\beta : y_{1, V} \to y_{2, V}$，并且须证明函子
$$
(\Sch/V)_{fppf} \to \textit{Sets},\quad
T/V \mapsto \{\alpha : x_{1, T} \to x_{2, T}
\text{ 于 }\mathcal{X}_T \mid F(\alpha) = \beta|_T\}
$$
是代数空间。考虑
$$
\mathcal{Z} = (\Sch/V)_{fppf} \times_{y_{1, V}, \mathcal{Y}} \mathcal{X}
$$
的对象 $z_1 = (V, x_{1, V}, \text{id})$ 和
$z_2 = (V, x_{2, V}, \beta)$。直接验证可知，上述函子在
$(\Sch/V)_{fppf}$ 上等于 $\mathit{Isom}(z_1,z_2)$。
因此，由 $F$ 是代数的这一假设（以及代数叠的定义），它是代数空间。
\end{proof}
















\section{截面空间}
\label{criteria-section-spaces-sections}

\noindent
给定态射 $W \to Z \to U$，可以考虑如下函子：它把 $U$ 上的概形
$U'$ 映为态射 $W \to Z$ 的基变换 $W_{U'} \to Z_{U'}$ 的截面
$\sigma : Z_{U'} \to W_{U'}$ 所成的集合。本节证明关于此函子的若干预备引理。

\begin{lemma}
\label{criteria-lemma-surjection-space-of-sections}
设 $Z \to U$ 是概形的有限态射，$W$ 是代数空间，而
$W \to Z$ 是满平展态射。则存在满平展态射 $U' \to U$ 和截面
$$
\sigma : Z_{U'} \to W_{U'}
$$
，它是态射 $W_{U'} \to Z_{U'}$ 的截面。
\end{lemma}

\begin{proof}
可以选取分离概形 $W'$ 和满平展态射 $W' \to W$。因此，以 $W'$ 替换
$W$ 后，可以假设 $W$ 是分离概形。记 $f : W \to Z$ 且
$\pi : Z \to U$。由于 $W$ 分离，
% P11-E0181: frozen source writes f \circ \pi : W \to U although the typed composition is \pi \circ f; preserved canonically.
$f \circ \pi : W \to U$ 分离（参见《概形》引理
\ref{schemes-lemma-compose-after-separated}）。设 $u \in U$ 是一点。
显然，只需找到 $(U,u)$ 的平展邻域 $(U',u')$，使 $U'$ 上存在截面
$\sigma$。设 $z_1,\ldots,z_r$ 是 $Z$ 中位于 $u$ 上方的点。
对每个 $i$，选取映到 $z_i$ 的点 $w_i\in W$。可以选取平展邻域
$(U', u') \to (U, u)$，使《态射进阶》引理
\ref{more-morphisms-lemma-etale-splits-off-quasi-finite-part-technical-variant}
的结论同时适用于 $Z \to U$ 与诸点 $z_1,\ldots,z_r$，以及
$W \to U$ 与诸点 $w_1,\ldots,w_r$。因此，以 $(U',u')$ 替换
$(U,u)$ 并重新编号后，可以假设所有域扩张
$\kappa(z_i)/\kappa(u)$ 和 $\kappa(w_i)/\kappa(u)$ 都是纯不可分的，
而且存在由开闭子概形给出的不交并分解
$$
Z = V_1 \amalg \ldots \amalg V_r \amalg A, \quad
W = W_1 \amalg \ldots \amalg W_r \amalg B
$$
，其中 $z_i\in V_i$、$w_i\in W_i$，且 $V_i\to U$、$W_i\to U$ 有限。
以 $U \setminus \pi(A)$ 替换 $U$ 后，可以假设 $A=\emptyset$，即
$Z = V_1 \amalg \ldots \amalg V_r$。再以
$W_i \cap f^{-1}(V_i)$ 替换 $W_i$，并以
$B \cup \bigcup W_i \cap f^{-1}(Z \setminus V_i)$ 替换 $B$，
可以假设 $f$ 把 $W_i$ 映入 $V_i$。此时
$f_i = f|_{W_i} : W_i \to V_i$ 是 $U$ 上有限的概形态射，因而有限
（参见《态射》引理 \ref{morphisms-lemma-finite-permanence}）。
它还按假设为平展，满足 $f_i^{-1}(\{z_i\})=w_i$，并诱导剩余域同构
$\kappa(z_i) = \kappa(w_i)$（因为二者都是 $\kappa(u)$ 的纯不可分扩张，
而 $f$ 平展，所以 $\kappa(w_i)/\kappa(z_i)$ 可分）。因此，由
《平展态射》引理 \ref{etale-lemma-finite-etale-one-point}，$f_i$ 在
$z_i$ 的某个邻域 $V_i'$ 上是同构。由于 $\pi : Z \to U$ 是闭映射，
缩小 $U$ 后，可以假设 $W_i \to V_i$ 是同构。这就证明了引理。
\end{proof}

\begin{lemma}
\label{criteria-lemma-space-of-sections}
设 $Z \to U$ 是概形的有限局部自由态射，$W$ 是代数空间，而
$W \to Z$ 是平展态射。则函子
$$
F : (\Sch/U)_{fppf}^{opp} \longrightarrow \textit{Sets},
$$
由规则
$$
U' \longmapsto
F(U') =
\{\sigma : Z_{U'} \to W_{U'}\text{ 下列对象的截面： }W_{U'} \to Z_{U'}\}
$$
定义；它是代数空间，且态射 $F \to U$ 平展。
\end{lemma}

\begin{proof}
先假设 $W \to Z$ 还分离。设 $U'$ 是 $U$ 上的概形且
$\sigma \in F(U')$。由《空间的态射》引理
\ref{spaces-morphisms-lemma-section-immersion}，态射 $\sigma$ 是闭浸入。
此外，由《空间的性质》引理
\ref{spaces-properties-lemma-etale-permanence}，$\sigma$ 平展。
故 $\sigma$ 也是开浸入；参见《空间的态射》引理
\ref{spaces-morphisms-lemma-etale-universally-injective-open}。
换言之，$Z_\sigma = \sigma(Z_{U'}) \subset W_{U'}$ 是开子空间，
使态射 $Z_\sigma \to Z_{U'}$ 为同构。特别地，态射
$Z_\sigma \to U'$ 有限。因此得到函子之间的变换
$$
F \longrightarrow (W/U)_{fin}, \quad
\sigma \longmapsto (U' \to U, Z_\sigma)
$$
，其中 $(W/U)_{fin}$ 是《空间群胚进阶》第
\ref{spaces-more-groupoids-section-finite} 节所引入的态射 $W \to U$
的有限部分。显然，此函子变换为单（因为可从 $Z_\sigma$ 恢复
$\sigma$，即取同构 $Z_\sigma \to Z_{U'}$ 的逆）。由《空间群胚进阶》命题
\ref{spaces-more-groupoids-proposition-finite-algebraic-space}
，$(W/U)_{fin}$ 是 $U$ 上平展的代数空间。因此，为完成此情形的证明，
只需证明 $F \to (W/U)_{fin}$ 可表示且为开浸入。为此，假设给定概形态射
$U' \to U$ 和开子空间 $Z' \subset W_{U'}$，且 $Z' \to U'$ 有限。
只需证明存在开子概形 $U'' \subset U'$，使态射 $T \to U'$ 通过
$U''$ 分解，当且仅当 $Z' \times_{U'} T$ 同构地映到
$Z \times_{U'} T$。这由《空间态射进阶》引理
\ref{spaces-more-morphisms-lemma-where-isomorphism}
推出（这里用到
% P11-E0182: frozen source refers to an undefined B in Z \to B; preserved canonically.
$Z \to B$ 平坦、局部有限表示且有限）。由此已证明 $W \to Z$
既分离又平展的情形。

\medskip\noindent
在一般情形下，选取分离概形 $W'$ 和满平展态射 $W' \to W$。
注意，由于其源分离，态射 $W' \to W$ 和
% P11-E0183: frozen source also claims W \to Z is separated "as their source is separated" although W was not assumed separated in this general case; preserved canonically.
$W \to Z$ 都分离。记 $F'$ 为引理中与 $W' \to Z \to U$ 相伴的函子。
在证明的第一段中，我们已证明 $F'$ 可由 $U$ 上平展的代数空间表示。
由引理 \ref{criteria-lemma-surjection-space-of-sections}，函子映射
$F' \to F$ 对 $\Sch/U$ 上的平展拓扑为满。此外，如果 $U'$ 和
$\sigma : Z_{U'} \to W_{U'}$ 定义一点 $\xi \in F(U')$，则纤维积
$$
F'' = F' \times_{F, \xi} U'
$$
是 $\Sch/U'$ 上与态射
$$
W'_{U'} \times_{W_{U'}, \sigma} Z_{U'} \to Z_{U'} \to U'.
$$
相伴的函子。由于第一个态射是分离态射的基变换，因而分离；由第一段的
结果，$F''$ 是 $U'$ 上平展的代数空间。于是 $F' \to F$ 是满平展的
函子变换，并可由代数空间表示。因此，由《自举》定理
\ref{bootstrap-theorem-final-bootstrap}，$F$ 是代数空间。
由于 $F' \to F$ 是代数空间的满平展态射，而 $F' \to U$ 平展，
所以 $F \to U$ 平展。
\end{proof}









\section{相对态射}
\label{criteria-section-relative-morphisms}

\noindent
继续《态射进阶》第 \ref{more-morphisms-section-relative-morphisms} 节
开始的讨论。

\medskip\noindent
设 $S$ 是概形，$Z \to B$ 和 $X \to B$ 是 $S$ 上代数空间的态射。
给定概形 $T$，可以考虑二元组 $(a,b)$，其中 $a : T \to B$ 是态射，
而 $b : T \times_{a, B} Z \to T \times_{a, B} X$ 是 $T$ 上的态射。
图示为
\begin{equation}
\label{criteria-equation-hom}
\vcenter{
\xymatrix{
T \times_{a, B} Z \ar[rd] \ar[rr]_b & &
T \times_{a, B} X \ar[ld] & Z \ar[rd] & & X \ar[ld] \\
& T \ar[rrr]^a & & & B
}
}
\end{equation}
当然，也可以把 $b$ 看作满足下图交换的态射
$b : T \times_{a, B} Z \to X$：
$$
\xymatrix{
T \times_{a, B} Z \ar[r] \ar[d] \ar@/^1pc/[rrr]_-b &
Z \ar[rd] & & X \ar[ld] \\
T \ar[rr]^a & & B
}
$$
在这种情形下，可以定义函子
\begin{equation}
\label{criteria-equation-hom-functor}
\mathit{Mor}_B(Z, X) : (\Sch/S)^{opp} \longrightarrow \textit{Sets},
\quad
T \longmapsto \{(a, b)\text{ 如上}\}
\end{equation}
有时把它视为定义在 $B$ 上概形范畴上的函子，此时从记号中省去 $a$。

\begin{lemma}
\label{criteria-lemma-hom-functor-sheaf}
设 $S$ 是概形，$Z \to B$ 和 $X \to B$ 是 $S$ 上代数空间的态射。则
\begin{enumerate}
\item $\mathit{Mor}_B(Z, X)$ 是 $(\Sch/S)_{fppf}$ 上的层。
\item 如果 $T$ 是 $S$ 上的代数空间，则存在典范双射
$$
\Mor_{\Sh((\Sch/S)_{fppf})}(T, \mathit{Mor}_B(Z, X))
=
\{(a, b)\text{ 如 }(\ref{criteria-equation-hom})\}
$$
\end{enumerate}
\end{lemma}

\begin{proof}
设 $T$ 是 $S$ 上的代数空间，$\{T_i \to T\}$ 是 $T$ 的一个 fppf 覆盖
（如《空间上的拓扑》第 \ref{spaces-topologies-section-fppf} 节）。
假设 $(a_i, b_i) \in \mathit{Mor}_B(Z, X)(T_i)$，且对所有 $i,j$，
$(a_i, b_i)|_{T_i \times_T T_j} = (a_j, b_j)|_{T_i \times_T T_j}$。
由《空间上的下降》引理
\ref{spaces-descent-lemma-fpqc-universal-effective-epimorphisms}，
存在唯一态射 $a : T \to B$，使 $a_i$ 是 $T_i \to T$ 与 $a$ 的复合。
于是 $\{T_i \times_{a_i, B} Z \to T \times_{a, B} Z\}$ 也是 fppf 覆盖；
同一引理表明存在唯一态射
$b : T \times_{a, B} Z \to T \times_{a, B} X$，使 $b_i$ 是
$T_i \times_{a_i, B} Z \to T \times_{a, B} Z$ 与 $b$ 的复合。
因此，$(a, b) \in \mathit{Mor}_B(Z, X)(T)$ 在 $T_i$ 上限制为
$(a_i,b_i)$，对所有 $i$ 均如此。

\medskip\noindent
注意，上一段的结果特别蕴含 (1)。

\medskip\noindent
设 $T$ 是 $S$ 上的代数空间。为证明 (2)，将在所示集合之间构造互逆映射。
下文所说的“二元组”均指符合 (\ref{criteria-equation-hom}) 的二元组 $(a,b)$。

\medskip\noindent
设 $v : T \to \mathit{Mor}_B(Z, X)$ 是自然变换。选取概形 $U$ 和满平展态射
$p : U \to T$。则 $v(p) \in \mathit{Mor}_B(Z, X)(U)$ 对应于 $U$ 上的
二元组 $(a_U,b_U)$。令 $R = U \times_T U$，投影为 $t,s:R\to U$。
由于 $v$ 是函子变换，$(a_U,b_U)$ 沿 $s$ 和 $t$ 的拉回相同。因此，
由于 $\{U \to T\}$ 是 fppf 覆盖，可以应用第一段的结果，
% P11-E0184: frozen source says "apply ... that deduce"; a coordinating conjunction is missing. Meaning is translated without changing formulas.
从而推得 $T$ 上存在唯一的二元组 $(a,b)$。

\medskip\noindent
反过来，设 $(a,b)$ 是 $T$ 上的二元组。令 $U \to T$、
$R = U \times_T U$ 和 $t,s:R\to U$ 如上。由 Yoneda 引理
（《范畴》引理 \ref{categories-lemma-yoneda}），限制 $(a,b)|_U$
给出函子变换 $v : h_U \to \mathit{Mor}_B(Z, X)$。由于两个拉回
$s^*(a, b)|_U$ 和 $t^*(a, b)|_U$ 相等，可知 $v$ 使两个映射
$h_t, h_s : h_R \to h_U$ 余等。由《空间》引理
\ref{spaces-lemma-space-presentation}，$T=U/R$ 是 fppf 商层；又由 (1)，
$\mathit{Mor}_B(Z, X)$ 是 fppf 层，故 $v$ 通过映射
$T \to \mathit{Mor}_B(Z, X)$ 分解。

\medskip\noindent
我们略去上述两个构造互逆的验证。
\end{proof}

\begin{lemma}
\label{criteria-lemma-base-change-hom-functor}
设 $S$ 是概形，$Z \to B$、$X \to B$ 和 $B' \to B$ 是 $S$ 上
代数空间的态射。令 $Z' = B' \times_B Z$ 且 $X' = B' \times_B X$。则
$$
\mathit{Mor}_{B'}(Z', X')
=
B' \times_B \mathit{Mor}_B(Z, X)
$$
在 $\Sh((\Sch/S)_{fppf})$ 中成立。
\end{lemma}

\begin{proof}
函子层面的等式立即由定义得到。由引理
\ref{criteria-lemma-hom-functor-sheaf}，等式两边都是层；又因为层的纤维积
等于相应预层（即函子）的纤维积，故层层面的等式也随之成立。
\end{proof}

\begin{lemma}
\label{criteria-lemma-etale-covering-hom-functor}
设 $S$ 是概形，$Z \to B$ 和 $X' \to X \to B$ 是 $S$ 上代数空间的态射。
假设
\begin{enumerate}
\item $X' \to X$ 平展；
\item $Z \to B$ 有限局部自由。
\end{enumerate}
则 $\mathit{Mor}_B(Z, X') \to \mathit{Mor}_B(Z, X)$ 可由代数空间表示且平展。
如果 $X' \to X$ 还为满，则
$\mathit{Mor}_B(Z, X') \to \mathit{Mor}_B(Z, X)$ 为满。
\end{lemma}

\begin{proof}
设 $U$ 是概形，且 $\xi=(a,b)$ 是 $\mathit{Mor}_B(Z,X)(U)$ 的元素。
须证明函子
$$
h_U \times_{\xi, \mathit{Mor}_B(Z, X)} \mathit{Mor}_B(Z, X')
$$
可由 $U$ 上平展的代数空间表示。令
$Z_U = U \times_{a, B} Z$ 且 $W = Z_U \times_{b, X} X'$。
则 $W \to Z_U \to U$ 如引理 \ref{criteria-lemma-space-of-sections} 中所述，
而该处定义的层 $F$ 与上述纤维积等同。这证明引理的第一个断言。
第二个断言由此以及引理 \ref{criteria-lemma-surjection-space-of-sections} 推出；
后者保证上述情形中的 $F \to U$ 为满。
\end{proof}

\begin{proposition}
\label{criteria-proposition-hom-functor-algebraic-space}
设 $S$ 是概形，$Z \to B$ 和 $X \to B$ 是 $S$ 上代数空间的态射。
如果 $Z \to B$ 有限局部自由，则 $\mathit{Mor}_B(Z, X)$ 是代数空间。
\end{proposition}

\begin{proof}
选取概形 $B'=\coprod B'_i$，它是仿射概形 $B'_i$ 的不交并，并选取
满平展态射 $B'\to B$。还可以假设 $B'_i\times_B Z$ 是某个环的谱，
该环作为 $\Gamma(B'_i, \mathcal{O}_{B'_i})$-模有限自由。由引理
\ref{criteria-lemma-base-change-hom-functor} 和《空间》引理
\ref{spaces-lemma-base-change-representable-transformations-property}
，态射 $\mathit{Mor}_{B'}(Z', X') \to \mathit{Mor}_B(Z, X)$ 满且平展。
因此，由《自举》定理 \ref{bootstrap-theorem-final-bootstrap}，只需证明
$B=B'$ 是仿射概形 $B'_i$ 的不交并、且每个 $B'_i\times_B Z$ 在
$B'_i$ 上有限自由的情形。事实上，此时只需对每个 $B'_i$ 上的限制证明
该结果。因此，可以假设 $B$ 仿射，且 $\Gamma(Z, \mathcal{O}_Z)$ 是有限自由
$\Gamma(B, \mathcal{O}_B)$-模。

\medskip\noindent
选取由仿射概形不交并而成的概形 $X'$ 以及满平展态射 $X'\to X$。
由引理 \ref{criteria-lemma-etale-covering-hom-functor}，态射
$\mathit{Mor}_B(Z, X') \to \mathit{Mor}_B(Z, X)$ 可由代数空间表示、平展且为满。
因此，由《自举》定理 \ref{bootstrap-theorem-final-bootstrap}，只需证明
$X$ 是仿射概形不交并的情形。这把问题归结为下一段讨论的情形。

\medskip\noindent
假设 $X=\coprod_{i\in I}X_i$ 是仿射概形的不交并，$B$ 仿射，且
$\Gamma(Z, \mathcal{O}_Z)$ 是有限自由 $\Gamma(B, \mathcal{O}_B)$-模。
对任意有限子集 $E\subset I$，令
$$
F_E = \mathit{Mor}_B(Z, \coprod\nolimits_{i \in E} X_i).
$$
由《态射进阶》引理
\ref{more-morphisms-lemma-hom-from-finite-free-into-affine}，$F_E$ 是代数空间。
考虑态射
$$
\coprod\nolimits_{E \subset I\text{ 有限}} F_E
\longrightarrow
\mathit{Mor}_B(Z, X)
$$
每个态射 $F_E\to\mathit{Mor}_B(Z,X)$ 都是开浸入，因为它恰好是参数化
二元组 $(a,b)$ 的轨迹，其中 $b$ 映入 $X$ 的开子概形
$\coprod\nolimits_{i\in E}X_i$。此外，如果 $T$ 拟紧，则对任意二元组
$(a,b)$，$b$ 的像包含于某个有限 $E\subset I$ 所对应的
$\coprod\nolimits_{i\in E}X_i$ 中。因此，所示箭头实际上是开覆盖，
由此结论成立\footnote{这里略去一些集合论论证。确切地，必须证明 $\coprod F_E$
是代数空间。这是因为 $|I|\leq\text{size}(X)$，并且
$\text{size}(F_E)\leq\text{size}(X)$；后者由《态射进阶》引理
\ref{more-morphisms-lemma-hom-from-finite-free-into-affine} 的证明中
对 $F_E$ 的显式描述推出。略去若干细节。}；这里应用了《空间》引理
\ref{spaces-lemma-glueing-algebraic-spaces}。
\end{proof}










\section{标量限制}
\label{criteria-section-restriction-of-scalars}

\noindent
假设 $X \to Z \to B$ 是 $S$ 上代数空间的态射。给定概形 $T$，
可以考虑二元组 $(a,b)$，其中 $a:T\to B$ 是态射，而
$b:T\times_{a,B}Z\to X$ 是 $Z$ 上的态射。图示为
\begin{equation}
\label{criteria-equation-pairs}
\vcenter{
\xymatrix{
& X \ar[d] \\
T \times_{a, B} Z \ar[d] \ar[ru]^b \ar[r] & Z \ar[d] \\
T \ar[r]^a & B
}
}
\end{equation}
在这种情形下，可以定义函子
\begin{equation}
\label{criteria-equation-restriction-of-scalars}
\text{Res}_{Z/B}(X) : (\Sch/S)^{opp} \longrightarrow \textit{Sets},
\quad
T \longmapsto \{(a, b)\text{ 如上}\}
\end{equation}
有时把它视为定义在 $B$ 上概形范畴上的函子，此时从记号中省去 $a$。

\begin{lemma}
\label{criteria-lemma-restriction-of-scalars-sheaf}
设 $S$ 是概形，$X \to Z \to B$ 是 $S$ 上代数空间的态射。则
\begin{enumerate}
\item $\text{Res}_{Z/B}(X)$ 是 $(\Sch/S)_{fppf}$ 上的层。
\item 如果 $T$ 是 $S$ 上的代数空间，则存在典范双射
$$
\Mor_{\Sh((\Sch/S)_{fppf})}(T, \text{Res}_{Z/B}(X))
=
\{(a, b)\text{ 如 }(\ref{criteria-equation-pairs})\}
$$
\end{enumerate}
\end{lemma}

\begin{proof}
设 $T$ 是 $S$ 上的代数空间，$\{T_i \to T\}$ 是 $T$ 的一个 fppf 覆盖
（如《空间上的拓扑》第 \ref{spaces-topologies-section-fppf} 节）。
假设 $(a_i,b_i)\in\text{Res}_{Z/B}(X)(T_i)$，且对所有 $i,j$，
$(a_i, b_i)|_{T_i \times_T T_j} = (a_j, b_j)|_{T_i \times_T T_j}$。
由《空间上的下降》引理
\ref{spaces-descent-lemma-fpqc-universal-effective-epimorphisms}，
存在唯一态射 $a:T\to B$，使 $a_i$ 是 $T_i\to T$ 与 $a$ 的复合。
于是 $\{T_i \times_{a_i, B} Z \to T \times_{a, B} Z\}$ 也是 fppf 覆盖；
同一引理表明存在唯一态射 $b:T\times_{a,B}Z\to X$，使 $b_i$ 是
$T_i \times_{a_i, B} Z \to T \times_{a, B} Z$ 与 $b$ 的复合。
因此，$(a,b)\in\text{Res}_{Z/B}(X)(T)$ 在 $T_i$ 上限制为
$(a_i,b_i)$，对所有 $i$ 均如此。

\medskip\noindent
注意，上一段的结果特别蕴含 (1)。

\medskip\noindent
设 $T$ 是 $S$ 上的代数空间。为证明 (2)，将在所示集合之间构造互逆映射。
下文所说的“二元组”均指符合 (\ref{criteria-equation-pairs}) 的二元组 $(a,b)$。

\medskip\noindent
设 $v:T\to\text{Res}_{Z/B}(X)$ 是自然变换。选取概形 $U$ 和满平展态射
$p:U\to T$。则 $v(p)\in\text{Res}_{Z/B}(X)(U)$ 对应于 $U$ 上的二元组
$(a_U,b_U)$。令 $R=U\times_TU$，投影为 $t,s:R\to U$。
由于 $v$ 是函子变换，$(a_U,b_U)$ 沿 $s$ 和 $t$ 的拉回相同。因此，
由于 $\{U\to T\}$ 是 fppf 覆盖，可以应用第一段的结果，
% P11-E0185: frozen source repeats the "apply ... that deduce" grammatical defect here; translated to the intended grammar.
从而推得 $T$ 上存在唯一的二元组 $(a,b)$。

\medskip\noindent
反过来，设 $(a,b)$ 是 $T$ 上的二元组。令 $U\to T$、
$R=U\times_TU$ 和 $t,s:R\to U$ 如上。由 Yoneda 引理
（《范畴》引理 \ref{categories-lemma-yoneda}），限制 $(a,b)|_U$
给出函子变换 $v:h_U\to\text{Res}_{Z/B}(X)$。由于两个拉回
$s^*(a, b)|_U$ 和 $t^*(a, b)|_U$ 相等，可知 $v$ 使两个映射
$h_t,h_s:h_R\to h_U$ 余等。由《空间》引理
\ref{spaces-lemma-space-presentation}，$T=U/R$ 是 fppf 商层；又由 (1)，
$\text{Res}_{Z/B}(X)$ 是 fppf 层，故 $v$ 通过映射
$T\to\text{Res}_{Z/B}(X)$ 分解。

\medskip\noindent
我们略去上述两个构造互逆的验证。
\end{proof}

\noindent
当然，层 $\text{Res}_{Z/B}(X)$ 自带函子的自然变换
$\text{Res}_{Z/B}(X)\to B$。下文将不加说明地使用它。

\begin{lemma}
\label{criteria-lemma-etale-base-change-restriction-of-scalars}
设 $S$ 是概形，$X\to Z\to B$ 和 $B'\to B$ 是 $S$ 上代数空间的态射。
令 $Z'=B'\times_BZ$ 且 $X'=B'\times_BX$。则
$$
\text{Res}_{Z'/B'}(X')
=
B' \times_B \text{Res}_{Z/B}(X)
$$
在 $\Sh((\Sch/S)_{fppf})$ 中成立。
\end{lemma}

\begin{proof}
函子层面的等式立即由定义得到。由引理
\ref{criteria-lemma-restriction-of-scalars-sheaf}，等式两边都是层；又因为层的
纤维积等于相应预层（即函子）的纤维积，故层层面的等式也随之成立。
\end{proof}

\begin{lemma}
\label{criteria-lemma-etale-covering-restriction-of-scalars}
设 $S$ 是概形，$X'\to X\to Z\to B$ 是 $S$ 上代数空间的态射。假设
\begin{enumerate}
\item $X'\to X$ 平展；
\item $Z\to B$ 有限局部自由。
\end{enumerate}
则 $\text{Res}_{Z/B}(X')\to\text{Res}_{Z/B}(X)$ 可由代数空间表示且平展。
如果 $X'\to X$ 还为满，则
$\text{Res}_{Z/B}(X')\to\text{Res}_{Z/B}(X)$ 为满。
\end{lemma}

\begin{proof}
设 $U$ 是概形，且 $\xi=(a,b)$ 是 $\text{Res}_{Z/B}(X)(U)$ 的元素。
须证明函子
$$
h_U \times_{\xi, \text{Res}_{Z/B}(X)} \text{Res}_{Z/B}(X')
$$
可由 $U$ 上平展的代数空间表示。令
$Z_U=U\times_{a,B}Z$ 且 $W=Z_U\times_{b,X}X'$。
则 $W\to Z_U\to U$ 如引理 \ref{criteria-lemma-space-of-sections} 中所述，
而该处定义的层 $F$ 与上述纤维积等同。这证明引理的第一个断言。
第二个断言由此以及引理 \ref{criteria-lemma-surjection-space-of-sections} 推出；
后者保证上述情形中的 $F\to U$ 为满。
\end{proof}

\noindent
此时，可以利用上述引理证明：只要 $Z\to B$ 有限局部自由，
$\text{Res}_{Z/B}(X)$ 就是代数空间；其证明几乎完全类似于
$\mathit{Mor}_B(Z,X)$ 是代数空间的证明，参见命题
\ref{criteria-proposition-hom-functor-algebraic-space}。不过，我们将从下一个引理
以及 $\mathit{Mor}_B(Z,X)$ 是代数空间这一事实直接推出该结果。

\begin{lemma}
\label{criteria-lemma-fibre-diagram}
设 $S$ 是概形，$X\to Z\to B$ 是 $S$ 上代数空间的态射。下图
$$
\xymatrix{
\mathit{Mor}_B(Z, X) \ar[r] & \mathit{Mor}_B(Z, Z) \\
\text{Res}_{Z/B}(X) \ar[r] \ar[u] & B \ar[u]_{\text{id}_Z}
}
$$
是 $(\Sch/S)_{fppf}$ 上层的笛卡尔图。
\end{lemma}

\begin{proof}
略。提示：这是代数几何函子观点的一道练习。
\end{proof}

\begin{proposition}
\label{criteria-proposition-restriction-of-scalars-algebraic-space}
设 $S$ 是概形，$X\to Z\to B$ 是 $S$ 上代数空间的态射。
如果 $Z\to B$ 有限局部自由，则 $\text{Res}_{Z/B}(X)$ 是代数空间。
\end{proposition}

\begin{proof}
由命题 \ref{criteria-proposition-hom-functor-algebraic-space}，函子
$\mathit{Mor}_B(Z,X)$ 和 $\mathit{Mor}_B(Z,Z)$ 都是代数空间。
因此，结论由引理 \ref{criteria-lemma-fibre-diagram} 的笛卡尔图以及下述事实推出：
代数空间的纤维积存在，并由其底层集合层范畴中的纤维积给出
（参见《空间》引理
\ref{spaces-lemma-fibre-product-spaces-over-sheaf-with-representable-diagonal}）。
\end{proof}






\section{有限 Hilbert 叠}
\label{criteria-section-finite-hilbert-stacks}

\noindent
本节证明若干关于《叠的例子》第
\ref{examples-stacks-section-hilbert-d-stack} 节所引入的有限 Hilbert 叠
$\mathcal{H}_d(\mathcal{X}/\mathcal{Y})$ 的结果。

\begin{lemma}
\label{criteria-lemma-map-hilbert}
考虑 $2$-交换图
$$
\xymatrix{
\mathcal{X}' \ar[r]_G \ar[d]_{F'} & \mathcal{X} \ar[d]^F \\
\mathcal{Y}' \ar[r]^H & \mathcal{Y}
}
$$
，其中各项为 $(\Sch/S)_{fppf}$ 上的群胚叠，并给定 $2$-同构
$\gamma : H \circ F' \to F \circ G$。在这种情形下，得到典范 $1$-态射
$\mathcal{H}_d(\mathcal{X}'/\mathcal{Y}') \to
\mathcal{H}_d(\mathcal{X}/\mathcal{Y})$.
此态射与《叠的例子》公式
(\ref{examples-stacks-equation-diagram-hilbert-d-stack}) 中的遗忘 $1$-态射相容。
\end{lemma}

\begin{proof}
把对象 $(U, Z, y', x', \alpha')$ 映为对象
$(U, Z, H(y'), G(x'), \gamma \star \text{id}_H \star \alpha')$
，其中 $\star$ 表示 $2$-态射的水平复合；参见《范畴》定义
\ref{categories-definition-horizontal-composition}。对态射
$(f, g, b, a) :
(U_1, Z_1, y_1', x_1', \alpha_1') \to (U_2, Z_2, y_2', x_2', \alpha_2')$
，赋予态射
$(f, g, H(b), G(a))$.
略去它定义 $(\Sch/S)_{fppf}$ 上范畴之间函子的验证。
\end{proof}

\begin{lemma}
\label{criteria-lemma-cartesian-map-hilbert}
在引理 \ref{criteria-lemma-map-hilbert} 的情形下，假设给定方形是 $2$-笛卡尔的。
则图
$$
\xymatrix{
\mathcal{H}_d(\mathcal{X}'/\mathcal{Y}') \ar[r] \ar[d] &
\mathcal{H}_d(\mathcal{X}/\mathcal{Y}) \ar[d] \\
\mathcal{Y}' \ar[r] &
\mathcal{Y}
}
$$
是 $2$-笛卡尔的。
\end{lemma}

\begin{proof}
由引理 \ref{criteria-lemma-map-hilbert} 得到 $2$-交换图，因而得到 $1$-态射
（即函子）
$$
\mathcal{H}_d(\mathcal{X}'/\mathcal{Y}')
\longrightarrow
\mathcal{Y}' \times_\mathcal{Y} \mathcal{H}_d(\mathcal{X}/\mathcal{Y})
$$
说明此函子为何本质满。右侧范畴的对象由下列数据给出：$S$ 上的概形
$U$、$\mathcal{Y}'_U$ 的对象 $y'$、
$\mathcal{H}_d(\mathcal{X}/\mathcal{Y})$ 在 $U$ 上的对象
$(U,Z,y,x,\alpha)$，以及 $\mathcal{Y}_U$ 中的同构 $H(y')\to y$。
假设恰好意味着存在 $\mathcal{X}'_Z$ 的对象 $x'$，并存在与 $\alpha$
相容的同构 $G(x')\cong x$ 和 $\alpha':y'|_Z\to F'(x')$。
于是 $(U,Z,y',x',\alpha')$ 是
$\mathcal{H}_d(\mathcal{X}'/\mathcal{Y}')$ 在 $U$ 上的对象。略去细节。
\end{proof}

\begin{lemma}
\label{criteria-lemma-etale-covering-hilbert}
在引理 \ref{criteria-lemma-map-hilbert} 的情形下，假设
\begin{enumerate}
\item $\mathcal{Y}' = \mathcal{Y}$ 且 $H = \text{id}_\mathcal{Y}$；
\item $G$ 可由代数空间表示且平展。
\end{enumerate}
则 $\mathcal{H}_d(\mathcal{X}'/\mathcal{Y}) \to
\mathcal{H}_d(\mathcal{X}/\mathcal{Y})$ 可由代数空间表示且平展。
如果 $G$ 还为满，则
$\mathcal{H}_d(\mathcal{X}'/\mathcal{Y}) \to
\mathcal{H}_d(\mathcal{X}/\mathcal{Y})$ 为满。
\end{lemma}

\begin{proof}
设 $U$ 是概形，而 $\xi=(U,Z,y,x,\alpha)$ 是
$\mathcal{H}_d(\mathcal{X}/\mathcal{Y})$ 在 $U$ 上的对象。
须证明 $2$-纤维积
\begin{equation}
\label{criteria-equation-to-show}
(\Sch/U)_{fppf}
\times_{\xi, \mathcal{H}_d(\mathcal{X}/\mathcal{Y})}
\mathcal{H}_d(\mathcal{X}'/\mathcal{Y})
\end{equation}
可由 $U$ 上平展的代数空间表示。它在 $U'$ 上的对象对应于
$\mathcal{X}'$ 在 $Z_{U'}$ 上的纤维范畴中的对象 $x'$，并满足
$G(x')\cong x|_{Z_{U'}}$。按假设，$2$-纤维积
$$
(\Sch/Z)_{fppf} \times_{x, \mathcal{X}} \mathcal{X}'
$$
可由代数空间 $W$ 表示，且投影 $W\to Z$ 平展。于是
(\ref{criteria-equation-to-show}) 可由引理 \ref{criteria-lemma-space-of-sections} 中引入的
代数空间 $F$ 表示；该空间参数化 $U$ 上 $W\to Z$ 的截面。
由于 $F\to U$ 平展，可知
$\mathcal{H}_d(\mathcal{X}'/\mathcal{Y}) \to
\mathcal{H}_d(\mathcal{X}/\mathcal{Y})$ 可由代数空间表示且平展。
最后，如果 $\mathcal{X}'\to\mathcal{X}$ 还为满，则 $W\to Z$ 为满，
因而由引理 \ref{criteria-lemma-surjection-space-of-sections}，$F\to U$ 为满。
所以在此情形下
$\mathcal{H}_d(\mathcal{X}'/\mathcal{Y}) \to
\mathcal{H}_d(\mathcal{X}/\mathcal{Y})$ 也为满。
\end{proof}

\begin{lemma}
\label{criteria-lemma-etale-map-hilbert}
在引理 \ref{criteria-lemma-map-hilbert} 的情形下，假设 $G$、$H$ 均可由代数空间
表示且平展。则 $\mathcal{H}_d(\mathcal{X}'/\mathcal{Y}') \to
\mathcal{H}_d(\mathcal{X}/\mathcal{Y})$ 可由代数空间表示且平展。
如果 $H$ 还为满，且诱导函子
$\mathcal{X}' \to \mathcal{Y}' \times_\mathcal{Y} \mathcal{X}$
为满，则
$\mathcal{H}_d(\mathcal{X}'/\mathcal{Y}') \to
\mathcal{H}_d(\mathcal{X}/\mathcal{Y})$ 为满。
\end{lemma}

\begin{proof}
令 $\mathcal{X}'' = \mathcal{Y}' \times_\mathcal{Y} \mathcal{X}$。
由引理 \ref{criteria-lemma-etale-permanence}，$1$-态射
$\mathcal{X}' \to \mathcal{X}''$ 可由代数空间表示且平展（特别地，
引理第二个断言中要求 $\mathcal{X}'\to\mathcal{X}''$ 为满的条件是有意义的）。
得到 $2$-交换图
$$
\xymatrix{
\mathcal{X}' \ar[r] \ar[d] &
\mathcal{X}'' \ar[r] \ar[d] &
\mathcal{X} \ar[d] \\
\mathcal{Y}' \ar[r] &
\mathcal{Y}' \ar[r] &
\mathcal{Y}
}
$$
由引理 \ref{criteria-lemma-cartesian-map-hilbert}，
$\mathcal{H}_d(\mathcal{X}''/\mathcal{Y}')$ 是
$\mathcal{H}_d(\mathcal{X}/\mathcal{Y})$ 沿
$\mathcal{Y}'\to\mathcal{Y}$ 的基变换。特别地，
$\mathcal{H}_d(\mathcal{X}''/\mathcal{Y}') \to
\mathcal{H}_d(\mathcal{X}/\mathcal{Y})$ 可由代数空间表示且平展；参见
《代数叠》引理
\ref{algebraic-lemma-base-change-representable-transformations-property}.
此外，如果 $H$ 为满，它也为满。因此，只要证明该结果对图中的左方形成立，
证明即告完成。这样便归结到 $\mathcal{Y}'=\mathcal{Y}$ 的情形，
而这正是引理 \ref{criteria-lemma-etale-covering-hilbert} 的内容。
\end{proof}

\begin{lemma}
\label{criteria-lemma-relative-hilbert}
设 $F : \mathcal{X} \to \mathcal{Y}$ 是 $(\Sch/S)_{fppf}$ 上群胚叠的
$1$-态射。假设
$\Delta : \mathcal{Y} \to \mathcal{Y} \times \mathcal{Y}$
可由代数空间表示。则态射
$$
\mathcal{H}_d(\mathcal{X}/\mathcal{Y})
\longrightarrow
\mathcal{H}_d(\mathcal{X}) \times \mathcal{Y}
$$
（见《叠的例子》公式
(\ref{examples-stacks-equation-diagram-hilbert-d-stack})）可由代数空间表示。
\end{lemma}

\begin{proof}
设 $U$ 是概形，而 $\xi=(U,Z,p,x,1)$ 是
$\mathcal{H}_d(\mathcal{X})=\mathcal{H}_d(\mathcal{X}/S)$ 在 $U$ 上的对象。
这里 $p$ 只是 $U$ 的结构态射。由于一切都在 $S$ 上，第五个分量 $1$
存在且唯一。另设 $y$ 是 $\mathcal{Y}$ 在 $U$ 上的对象。须证明
$2$-纤维积
\begin{equation}
\label{criteria-equation-res-isom}
(\Sch/U)_{fppf}
\times_{\xi \times y, \mathcal{H}_d(\mathcal{X}) \times \mathcal{Y}}
\mathcal{H}_d(\mathcal{X}/\mathcal{Y})
\end{equation}
可由代数空间表示。为说明这一点，引入
$$
I = \mathit{Isom}_\mathcal{Y}(y|_Z, F(x))
$$
。按假设，它是 $Z$ 上的代数空间。设 $a:U'\to U$ 是 $U$ 上的概形。
给出 (\ref{criteria-equation-res-isom}) 的纤维范畴在 $U'$ 上的对象意味着什么？
这意味着给出 $\mathcal{H}_d(\mathcal{X}/\mathcal{Y})$ 在 $U'$ 上的对象
$\xi'=(U',Z',y',x',\alpha')$，以及同构
$(U', Z', p', x', 1) \cong (U, Z, p, x, 1)|_{U'}$ 和
$y' \cong y|_{U'}$。因此，$\xi'$ 同构于
$(U', U' \times_{a, U} Z, a^*y, x|_{U' \times_{a, U} Z}, \alpha)$
，其中有某个态射
$$
\alpha :
a^*y|_{U' \times_{a, U} Z}
\longrightarrow
F(x|_{U' \times_{a, U} Z})
$$
，它位于 $\mathcal{Y}$ 在 $U'\times_{a,U}Z$ 上的纤维范畴中。
因此，可以把 $\alpha$ 看作态射 $b:U'\times_{a,U}Z\to I$。
由此可知，(\ref{criteria-equation-res-isom}) 可由 $\text{Res}_{Z/U}(I)$ 表示；
由命题 \ref{criteria-proposition-restriction-of-scalars-algebraic-space}，后者是代数空间。
\end{proof}

\noindent
下一个引理是引理 \ref{criteria-lemma-etale-covering-hilbert} 的一个（部分）推广。

\begin{lemma}
\label{criteria-lemma-representable-on-top}
设 $F : \mathcal{X} \to \mathcal{Y}$ 和
$G : \mathcal{X}' \to \mathcal{X}$ 是 $(\Sch/S)_{fppf}$ 上群胚叠的
$1$-态射。如果 $G$ 可由代数空间表示，则 $1$-态射
$$
\mathcal{H}_d(\mathcal{X}'/\mathcal{Y})
\longrightarrow
\mathcal{H}_d(\mathcal{X}/\mathcal{Y})
$$
可由代数空间表示。
\end{lemma}

\begin{proof}
设 $U$ 是概形，而 $\xi=(U,Z,y,x,\alpha)$ 是
$\mathcal{H}_d(\mathcal{X}/\mathcal{Y})$ 在 $U$ 上的对象。
须证明 $2$-纤维积
\begin{equation}
\label{criteria-equation-to-show-again}
(\Sch/U)_{fppf}
\times_{\xi, \mathcal{H}_d(\mathcal{X}/\mathcal{Y})}
\mathcal{H}_d(\mathcal{X}'/\mathcal{Y})
\end{equation}
可由 $U$ 上的代数空间表示。它在 $a:U'\to U$ 上的对象对应于
$\mathcal{X}'$ 在 $U'\times_{a,U}Z$ 上的对象 $x'$，并满足
$G(x')\cong x|_{U'\times_{a,U}Z}$。按假设，$2$-纤维积
$$
(\Sch/Z)_{fppf} \times_{x, \mathcal{X}} \mathcal{X}'
$$
可由 $Z$ 上的代数空间 $W$ 表示。因此，(\ref{criteria-equation-to-show-again})
可由 $\text{Res}_{Z/U}(W)$ 表示；由命题
\ref{criteria-proposition-restriction-of-scalars-algebraic-space}，后者是代数空间。
\end{proof}

\begin{lemma}
\label{criteria-lemma-limit-preserving}
设 $F : \mathcal{X} \to \mathcal{Y}$ 是 $(\Sch/S)_{fppf}$ 上群胚叠的
$1$-态射。假设 $F$ 可由代数空间表示且局部有限表示。则
$$
p : \mathcal{H}_d(\mathcal{X}/\mathcal{Y}) \to \mathcal{Y}
$$
在对象上保极限。
\end{lemma}

\begin{proof}
这意味着必须证明如下命题：给定
\begin{enumerate}
\item 仿射概形 $U = \lim_i U_i$，它写成 $S$ 上仿射概形 $U_i$ 的有向极限；
\item 对某个 $i$，$\mathcal{Y}$ 在 $U_i$ 上的对象 $y_i$；以及
\item $\mathcal{H}_d(\mathcal{X}/\mathcal{Y})$ 在 $U$ 上的对象
$\Xi = (U, Z, y, x, \alpha)$，满足 $y = y_i|_U$，
\end{enumerate}
则存在 $i' \geq i$ 和
$\Xi_{i'} = (U_{i'}, Z_{i'}, y_{i'}, x_{i'}, \alpha_{i'})$，它是
$\mathcal{H}_d(\mathcal{X}/\mathcal{Y})$ 在 $U_{i'}$ 上的对象，并满足
$\Xi_{i'}|_U = \Xi$ 和 $y_{i'} = y_i|_{U_{i'}}$。最后这两个等式恰好保证
(\ref{criteria-equation-limit-preserving}) 交换。

\medskip\noindent
设 $X_{y_i} \to U_i$ 是表示 $2$-纤维积
$$
(\Sch/U_i)_{fppf} \times_{y_i, \mathcal{Y}, F} \mathcal{X}.
$$
的代数空间。由关于 $F$ 的假设，$X_{y_i} \to U_i$ 局部有限表示。
% P11-E0186: frozen source contains the isolated incomplete sentence "Write Xi." here; retained semantically without inventing missing content.
写下 $\Xi$。显然，$\xi = (Z, Z \to U_i, x, \alpha)$ 是上述
$2$-纤维积的对象，故 $\xi$ 给出 $U_i$ 上代数空间的态射
$f_\xi : Z \to X_{y_i}$（因为 $X_{y_i}$ 是
% P11-E0187: frozen source says "isomorphisms classes"; the intended singular attributive form is translated.
$(\Sch/U_i)_{fppf} \times_{y, \mathcal{Y}, F} \mathcal{X}$
的对象同构类函子；参见《代数叠》引理
\ref{algebraic-lemma-characterize-representable-by-space}）。
% P11-E0188: frozen source uses y rather than y_i in the displayed fibre product above; the canonical math token is preserved.
由《极限》引理 \ref{limits-lemma-descend-finite-presentation} 和
\ref{limits-lemma-descend-finite-locally-free}
，存在 $i' \geq i$ 和次数为 $d$ 的有限局部自由态射
$Z_{i'} \to U_{i'}$，其到 $U$ 的基变换为 $Z$。由《空间的极限》命题
\ref{spaces-limits-proposition-characterize-locally-finite-presentation}
，以更大的指标替换 $i'$ 后，可以假设存在态射
$f_{i'} : Z_{i'} \to X_{y_i}$，使下图交换：
$$
\xymatrix{
Z \ar[d] \ar[r] \ar@/^3ex/[rr]^{f_\xi} &
Z_{i'} \ar[d] \ar[r]_{f_{i'}} & X_{y_i} \ar[d] \\
U \ar[r] & U_{i'} \ar[r] & U_i
}
$$
令
$\Xi_{i'} = (U_{i'}, Z_{i'}, y_{i'}, x_{i'}, \alpha_{i'})$，其中
\begin{enumerate}
\item $y_{i'}$ 是 $\mathcal{Y}$ 在 $U_{i'}$ 上的对象，即 $y_i$ 到
$U_{i'}$ 的拉回；
\item $x_{i'}$ 是 $\mathcal{X}$ 在 $Z_{i'}$ 上的对象，它经由
$2$-Yoneda 引理对应于 $1$-态射
$$
(\Sch/Z_{i'})_{fppf} \to
\mathcal{S}_{X_{y_i}} \to
(\Sch/U_i)_{fppf} \times_{y_i, \mathcal{Y}, F} \mathcal{X} \to
\mathcal{X}
$$
，其中中间箭头是定义 $X_{y_i}$ 的等价（记号如《代数叠》第
\ref{algebraic-section-representable-by-algebraic-spaces} 节和第
\ref{algebraic-section-split} 节）。
\item $\alpha_{i'} : y_{i'}|_{Z_{i'}} \to F(x_{i'})$ 是由下图的
$2$-交换性得到的同构：
$$
\xymatrix{
(\Sch/Z_{i'})_{fppf} \ar[r] \ar[rd] &
(\Sch/U_i)_{fppf} \times_{y_i, \mathcal{Y}, F} \mathcal{X}
\ar[r] \ar[d] &
\mathcal{X} \ar[d]^F \\
& (\Sch/U_{i'})_{fppf} \ar[r] & \mathcal{Y}
}
$$
\end{enumerate}
回顾 $f_\xi : Z \to X_{y_i}$ 是与对象
$\xi = (Z, Z \to U_i, x, \alpha)$ 相对应的态射；该对象是
$(\Sch/U_i)_{fppf} \times_{y_i, \mathcal{Y}, F} \mathcal{X}$
在 $Z$ 上的对象。按构造，$f_{i'}$ 是与对象
$\xi_{i'} = (Z_{i'}, Z_{i'} \to U_i, x_{i'}, \alpha_{i'})$ 相对应的态射。
由于 $f_\xi = f_{i'} \circ (Z \to Z_{i'})$，可知对象
$\xi_{i'} = (Z_{i'}, Z_{i'} \to U_i, x_{i'}, \alpha_{i'})$ 到 $Z$ 的拉回为
$\xi$。因此，$x_{i'}$ 拉回为 $x$，而 $\alpha_{i'}$ 拉回为 $\alpha$。
这意味着 $\Xi_{i'}$ 到 $U$ 的拉回为 $\Xi$，结论得证。
\end{proof}










\section{一点的有限 Hilbert 叠}
\label{criteria-section-hilbert-point}

\noindent
设 $d\geq1$ 是整数。在《叠的例子》定义
\ref{examples-stacks-definition-hilbert-d-stack} 中，我们定义了群胚叠
$\mathcal{H}_d$。本节证明 $\mathcal{H}_d$ 是代数叠。全文假设
$S=\Spec(\mathbf{Z})$；一般情形由基变换随之得到。回顾
$\mathcal{H}_d$ 在概形 $T$ 上的纤维范畴，是次数为 $d$ 的有限局部自由态射
$\pi:Z\to T$ 所成的范畴。我们不直接分类这些态射，而先研究拟凝聚代数层
$\pi_*\mathcal{O}_Z$。

\medskip\noindent
设 $R$ 是环。暂作如下定义：一个{\it $R$ 上的自由 $d$ 维代数}
由 $R^{\oplus d}$ 上的交换 $R$-代数结构 $m$ 给出，并要求
$e_1=(1,0,\ldots,0)$ 为单位元\footnote{也许更适合把它看作由乘法映射
$m:R^{\oplus d}\otimes_RR^{\oplus d}\to R^{\oplus d}$ 和满足若干公理的
环映射 $\psi:R\to R^{\oplus d}$ 组成的二元组。}。把 $m$ 看作 $R$-线性映射
$$
m : R^{\oplus d} \otimes_R R^{\oplus d} \longrightarrow R^{\oplus d}
$$
，它满足 $m(e_1,x)=m(x,e_1)=x$，并且 $m$ 定义交换结合的环结构。
若写 $m(e_i,e_j)=\sum a_{ij}^ke_k$，则这归结为条件
$$
\left\{
\begin{matrix}
\sum_l a_{ij}^la_{lk}^m = \sum_l a_{il}^ma_{jk}^l & \forall i, j, k, m \\
a_{ij}^k = a_{ji}^k & \forall i, j, k \\
a_{i1}^j = \delta_{ij} & \forall i, j
\end{matrix}
\right.
$$
，其中 $\delta_{ij}$ 是 Kronecker $\delta$-函数。于是定义
$$
R_{univ} = \mathbf{Z}[a_{ij}^k]/J
$$
，其中理想 $J$ 由上述关系生成。记
$$
m_{univ} :
R_{univ}^{\oplus d} \otimes_{R_{univ}} R_{univ}^{\oplus d}
\longrightarrow
R_{univ}^{\oplus d}
$$
% P11-E0189: frozen source's naming construction "Denote [display] the ..." omits "by"; translated naturally.
为 $R_{univ}$ 上的自由 $d$ 维代数 $m$，其结构常数为 $a_{ij}^k$ 模 $J$
的类。显然，给定环 $R$ 上任意自由 $d$ 维代数 $m$，存在唯一的
$\mathbf{Z}$-代数同态 $\psi:R_{univ}\to R$，使
$\psi_*m_{univ}=m$（这意味着 $m$ 由基变换函子
$-\otimes_{R_{univ}}R$ 作用于 $m_{univ}$ 得到）。换言之，令
$X=\Spec(R_{univ})$，得到典范等同
$$
X(T) = \{\text{自由 }d\text{-维代数 }m\text{ 相对于 }R\}
$$
，其中 $T=\Spec(R)$ 变化。通过 Zariski 局部化，得到下列看似更一般的等同
\begin{equation}
\label{criteria-equation-objects}
X(T) = \{\text{自由 }d\text{-维代数 }
m\text{ 相对于 }\Gamma(T, \mathcal{O}_T)\}
\end{equation}
，对任意概形 $T$ 成立。

\medskip\noindent
接下来略述{\it 自由 $d$ 维 $R$-代数的同构}。设 $m$、$m'$ 是环 $R$
上的两个自由 $d$ 维代数。一个{\it 从 $m$ 到 $m'$ 的同构}由可逆
$R$-线性映射
$$
\varphi : R^{\oplus d} \longrightarrow R^{\oplus d}
$$
给出，并满足 $\varphi(e_1)=e_1$ 以及
$$
m \circ \varphi \otimes \varphi = \varphi \circ m'.
$$
注意，可以复合这些同构，从而使 $R$ 上自由 $d$ 维代数的全体成为范畴。
由此得到函子
\begin{equation}
\label{criteria-equation-FAd}
FA_d : \Sch_{fppf}^{opp} \longrightarrow \textit{Groupoids}
\end{equation}
，它从概形范畴取值于群胚：对概形 $T$，赋予
$\Gamma(T,\mathcal{O}_T)$ 上自由 $d$ 维代数的集合，并用刚才定义的
同构概念赋予其范畴结构。

\medskip\noindent
上述讨论提示我们考虑群值函子 $G$，它把任意概形 $T$ 映为群
$$
G(T) = \{g \in \text{GL}_d(\Gamma(T, \mathcal{O}_T)) \mid g(e_1) = e_1\}
$$
显然，$G\subset\text{GL}_d$（参见《群胚》例
\ref{groupoids-example-general-linear-group}）是由方程
$x_{11}=1$ 和 $x_{i1}=0$（$i>1$）截出的闭子群概形。因此，$G$ 是
$\Spec(\mathbf{Z})$ 上的光滑仿射群概形。考虑作用
$$
a : G \times_{\Spec(\mathbf{Z})} X \longrightarrow X
$$
，它把左侧的 $T$-值点 $(g,m)$（其中 $T=\Spec(R)$）映为如下给出的
$R$ 上自由 $d$ 维代数：
$$
a(g, m) = g^{-1} \circ m \circ g \otimes g.
$$
注意，这意味着 $g$ 定义自由 $d$ 维 $R$-代数的同构
$m\to a(g,m)$。略去 $a$ 确实定义群概形 $G$ 在概形 $X$ 上作用的验证。

\begin{lemma}
\label{criteria-lemma-represent-FAd}
在 (\ref{criteria-equation-FAd}) 中定义的群胚值函子 $FA_d$ 同构（！）于如下
群胚值函子：它把概形 $T$ 映为具有下列数据的范畴
\begin{enumerate}
\item 对象集为 $X(T)$；
\item 态射集为 $G(T) \times X(T)$；
\item $s : G(T) \times X(T) \to X(T)$ 是投影映射；
\item $t : G(T) \times X(T) \to X(T)$ 是 $a(T)$；
\item 复合\newline
$G(T) \times X(T) \times_{s, X(T), t} G(T) \times X(T)
\to G(T) \times X(T)$ 由
$((g, m), (g', m')) \mapsto (gg', m')$ 给出。
\end{enumerate}
\end{lemma}

\begin{proof}
在 (\ref{criteria-equation-objects}) 中已经看到对象上的规则。上文还看到，
对任意自由 $d$ 维代数 $m$，$g\in G(T)$ 可视为从 $m$ 到 $a(g,m)$
的态射。反过来，任意态射 $m\to m'$ 由可逆线性映射 $\varphi$ 给出；
它对应于元素 $g\in G(T)$，并满足 $m'=a(g,m)$。
\end{proof}

\noindent
事实上，上一个引理所述的群胚 $(X,G\times X,s,t,c)$ 正是《群胚》引理
\ref{groupoids-lemma-groupoid-from-action} 中与作用
$a:G\times X\to X$ 相伴的群胚。由于 $G$ 在 $\Spec(\mathbf{Z})$ 上光滑，
可知两个态射 $s,t:G\times X\to X$ 都光滑：由对称性，只需证明其中一个，
而 $s$ 是 $G\to\Spec(\mathbf{Z})$ 的基变换。因此
$(G\times X,X,s,t,c)$ 是光滑群胚概形，并且由《代数叠》定理
\ref{algebraic-theorem-smooth-groupoid-gives-algebraic-stack}，商叠
$[X/G]$ 是代数叠。

\begin{proposition}
\label{criteria-proposition-finite-hilbert-point}
叠 $\mathcal{H}_d$ 等价于上述商叠 $[X/G]$。特别地，$\mathcal{H}_d$
是代数叠。
\end{proposition}

\begin{proof}
注意，由《空间中的群胚》定义
\ref{spaces-groupoids-definition-quotient-stack}，商叠 $[X/G]$ 是与下述
“群胚预层”相伴的群胚纤维化范畴的叠化；该预层把概形 $T$ 映为群胚
$$
(X(T), G(T) \times X(T), s, t, c).
$$
由引理 \ref{criteria-lemma-represent-FAd}，此“群胚预层”同构于 $FA_d$。
因此，只需证明
% P11-E0190: frozen source has the superfluous article "the" before mathcal H_d; translated naturally.
$\mathcal{H}_d$ 是（与“群胚预层”相伴的群胚纤维化范畴）$FA_d$ 的叠化。
为此，先定义函子
$$
\Spec : FA_d \longrightarrow \mathcal{H}_d
$$
回顾 $\mathcal{H}_d$ 在概形 $T$ 上的纤维范畴，是次数为 $d$ 的有限局部自由
态射 $Z\to T$ 所成的范畴。因此，给定概形 $T$ 和自由 $d$ 维
$\Gamma(T,\mathcal{O}_T)$-代数 $m$，可以赋予它对象
$$
Z = \underline{\Spec}_T(\mathcal{A})
$$
，它属于 $\mathcal{H}_{d,T}$；其中
$\mathcal{A}=\mathcal{O}_T^{\oplus d}$ 经由 $m$ 赋有
$\mathcal{O}_T$-代数结构。此外，若 $m'$ 是另一个这样的自由 $d$ 维
$\Gamma(T,\mathcal{O}_T)$-代数，而 $\varphi:m\to m'$ 是二者的同构，
则诱导的 $\mathcal{O}_T$-线性映射
$\varphi : \mathcal{O}_T^{\oplus d} \to \mathcal{O}_T^{\oplus d}$
诱导拟凝聚 $\mathcal{O}_T$-代数的同构
$$
\varphi : \mathcal{A}' \longrightarrow \mathcal{A}
$$
。因此
$$
\underline{\Spec}_T(\varphi) :
\underline{\Spec}_T(\mathcal{A})
\longrightarrow
\underline{\Spec}_T(\mathcal{A}')
$$
是纤维范畴 $\mathcal{H}_{d,T}$ 中的态射。略去此构造与基变换相容的验证，
于是确实得到上述函子 $\Spec:FA_d\to\mathcal{H}_d$。

\medskip\noindent
为证明 $\Spec:FA_d\to\mathcal{H}_d$ 在 $FA_d$ 的叠化与
$\mathcal{H}_d$ 之间诱导等价，只需验证
\begin{enumerate}
\item $\mathit{Isom}(m, m') = \mathit{Isom}(\Spec(m), \Spec(m'))$
，对任意 $m,m'\in FA_d(T)$；
\item 对任意概形 $T$ 和 $\mathcal{H}_{d,T}$ 的任意对象 $Z\to T$，
存在覆盖 $\{T_i\to T\}$，使对某个 $m\in FA_d(T_i)$，$Z|_{T_i}$
同构于 $\Spec(m)$；
\end{enumerate}
参见《叠》引理 \ref{stacks-lemma-stackify-groupoids}。第一个断言由下述观察推出：
任意同构
$$
\underline{\Spec}_T(\mathcal{A})
\longrightarrow
\underline{\Spec}_T(\mathcal{A}')
$$
在 $\mathcal{A}=\mathcal{A}'=\mathcal{O}_T^{\oplus d}$ 作为模时，必由
整体可逆矩阵 $g$ 给出。为证明第二个断言，设 $\pi:Z\to T$ 是次数为
$d$ 的有限局部自由态射。于是，
% P11-E0191: frozen source says "a locally free sheaf O_T-modules"; "of" is missing. Translated naturally.
$\mathcal{A}$ 是秩为 $d$ 的局部自由 $\mathcal{O}_T$-模层。考虑元素
$1\in\Gamma(T,\mathcal{A})$。对每个 $t\in T$，此元素在
$\mathcal{A}\otimes_{\mathcal{O}_{T,t}}\kappa(t)$ 中非零，因为概形
$Z_t = \Spec(\mathcal{A} \otimes_{\mathcal{O}_{T, t}} \kappa(t))$
在 $\kappa(t)$ 上次数为 $d>0$，故非空。因此，
$1:\mathcal{O}_T\to\mathcal{A}$ 可局部用作第一个基元素（例如可用
《代数》引理 \ref{algebra-lemma-cokernel-flat} 的 (1)、(2) 看出这一点）。
所以在 $T$ 上局部化后，可以假设存在同构
$\varphi : \mathcal{A} \to \mathcal{O}_T^{\oplus d}$
，使 $1\in\Gamma(\mathcal{A})$ 对应于第一个基元素。在这种情形下，
乘法映射
$\mathcal{A} \otimes_{\mathcal{O}_T} \mathcal{A} \to \mathcal{A}$
经由 $\varphi$ 转化为 $\Gamma(T,\mathcal{O}_T)$ 上的自由 $d$ 维代数
$m$。证明完成。
\end{proof}




\section{空间的有限 Hilbert 叠}
\label{criteria-section-spaces-hilbert}

\noindent
代数空间的有限 Hilbert 叠是代数叠。

\begin{lemma}
\label{criteria-lemma-hilbert-stack-of-space}
设 $S$ 是概形，$X$ 是 $S$ 上的代数空间。则 $\mathcal{H}_d(X)$ 是代数叠。
\end{lemma}

\begin{proof}
$1$-态射
$$
\mathcal{H}_d(X) \longrightarrow \mathcal{H}_d
$$
由引理 \ref{criteria-lemma-representable-on-top} 可由代数空间表示。
由命题 \ref{criteria-proposition-finite-hilbert-point}，叠 $\mathcal{H}_d$ 是代数叠。
因此，由《代数叠》引理
\ref{algebraic-lemma-representable-morphism-to-algebraic}，
$\mathcal{H}_d(X)$ 是代数叠。
\end{proof}

\noindent
此引理使我们能够进行自举。

\begin{lemma}
\label{criteria-lemma-hilbert-stack-relative-space}
设 $S$ 是概形，$F:\mathcal{X}\to\mathcal{Y}$ 是
$(\Sch/S)_{fppf}$ 上群胚叠的 $1$-态射，并满足
\begin{enumerate}
\item $\mathcal{X}$ 可由代数空间表示；
\item $F$ 可由代数空间表示、为满、平坦且局部有限表示。
\end{enumerate}
则 $\mathcal{H}_d(\mathcal{X}/\mathcal{Y})$ 是代数叠。
\end{lemma}

\begin{proof}
选取 $S$ 上的可表示群胚叠 $\mathcal{U}$ 和 $1$-态射
$f:\mathcal{U}\to\mathcal{H}_d(\mathcal{X})$，它可由代数空间表示、
光滑且为满。由于引理 \ref{criteria-lemma-hilbert-stack-of-space} 表明
$\mathcal{H}_d(\mathcal{X})$ 是代数叠，这种选择是可能的。考虑 $2$-纤维积
$$
\mathcal{W} =
\mathcal{H}_d(\mathcal{X}/\mathcal{Y})
\times_{\mathcal{H}_d(\mathcal{X}), f}
\mathcal{U}.
$$
由于 $\mathcal{U}$ 可表示（特别地，它是集合胚叠），由《叠的例子》引理
\ref{examples-stacks-lemma-faithful-hilbert} 和《叠》引理
\ref{stacks-lemma-2-fibre-product-gives-stack-in-setoids}，
$\mathcal{W}$ 是集合胚叠。$1$-态射
$\mathcal{W}\to\mathcal{H}_d(\mathcal{X}/\mathcal{Y})$ 是态射 $f$ 的基变换，
故可由代数空间表示、光滑且为满（参见《代数叠》引理
\ref{algebraic-lemma-base-change-representable-by-spaces} 和
\ref{algebraic-lemma-base-change-representable-transformations-property}）。
因此，如果能证明 $\mathcal{W}$ 可由代数空间表示，则由《代数叠》引理
\ref{algebraic-lemma-smooth-surjective-morphism-implies-algebraic}，
引理即得证。

\medskip\noindent
由引理 \ref{criteria-lemma-flat-finite-presentation-surjective-diagonal}，
$\mathcal{Y}$ 的对角态射可由代数空间表示。应用引理
\ref{criteria-lemma-relative-hilbert} 可知 $1$-态射
$$
\mathcal{H}_d(\mathcal{X}/\mathcal{Y})
\longrightarrow
\mathcal{H}_d(\mathcal{X}) \times \mathcal{Y}
$$
可由代数空间表示。考虑 $2$-纤维积
$$
\mathcal{V} =
\mathcal{H}_d(\mathcal{X}/\mathcal{Y})
\times_{(\mathcal{H}_d(\mathcal{X}) \times \mathcal{Y}), f \times F}
(\mathcal{U} \times \mathcal{X}).
$$
投影态射 $\mathcal{V}\to\mathcal{U}\times\mathcal{X}$ 是上一个所示态射的
基变换，故可由代数空间表示。因此 $\mathcal{V}$ 是代数空间（参见《自举》引理
\ref{bootstrap-lemma-representable-by-spaces-over-space}，或《代数叠》引理
\ref{algebraic-lemma-base-change-by-space-representable-by-space}）。
% P11-E0192: frozen source says V -> U here, but the displayed vertical morphism and subsequent argument are V -> W; preserved canonically in the formula.
$1$-态射 $\mathcal{V}\to\mathcal{U}$ 嵌入下列 $2$-笛卡尔图
$$
\xymatrix{
\mathcal{V} \ar[d] \ar[r] & \mathcal{X} \ar[d]^F \\
\mathcal{W} \ar[r] & \mathcal{Y}
}
$$
，这是因为
$$
\mathcal{H}_d(\mathcal{X}/\mathcal{Y})
\times_{(\mathcal{H}_d(\mathcal{X}) \times \mathcal{Y}), f \times F}
(\mathcal{U} \times \mathcal{X})
=
(\mathcal{H}_d(\mathcal{X}/\mathcal{Y})
\times_{\mathcal{H}_d(\mathcal{X}), f}
\mathcal{U}) \times_{\mathcal{Y}, F} \mathcal{X}.
$$
因此，$\mathcal{V}\to\mathcal{W}$ 作为 $F$ 的基变换，可由代数空间表示、
为满、平坦且局部有限表示。于是，与 $\mathcal{V}$ 和 $\mathcal{W}$
相伴的相应集合层也具有相同性质；参见《代数叠》引理
\ref{algebraic-lemma-map-fibred-setoids-property}。故由《自举》定理
\ref{bootstrap-theorem-final-bootstrap}，与 $\mathcal{W}$ 相伴的层
是代数空间。
\end{proof}




\section{Hilbert 叠中的 LCI 轨迹}
\label{criteria-section-lci}

\noindent
记号参见《叠的例子》第 \ref{examples-stacks-section-hilbert-d-stack} 节。
固定 $(\Sch/S)_{fppf}$ 上群胚叠的 $1$-态射
$F:\mathcal{X}\longrightarrow\mathcal{Y}$。假设 $F$ 可由代数空间表示。
固定 $d\geq1$。考虑 $\mathcal{H}_d(\mathcal{X}/\mathcal{Y})$ 的对象
$(U,Z,y,x,\alpha)$。存在诱导的 $1$-态射
$$
(\Sch/Z)_{fppf}
\longrightarrow
(\Sch/U)_{fppf} \times_{y, \mathcal{Y}, F} \mathcal{X}
$$
（由 $2$-纤维积的泛性质），它可由 $U$ 上代数空间的态射表示。
确切地，由于 $F$ 可由代数空间表示，可以选取 $U$ 上的代数空间 $X_y$，
它表示 $2$-纤维积
$(\Sch/U)_{fppf} \times_{y, \mathcal{Y}, F} \mathcal{X}$.
由于 $\alpha:y|_Z\to F(x)$ 是同构，可知
$\xi=(Z,Z\to U,x,\alpha)$ 是 $2$-纤维积
$(\Sch/U)_{fppf}\times_{y,\mathcal{Y},F}\mathcal{X}$ 在 $Z$ 上的对象。
因此，$\xi$ 给出 $U$ 上代数空间的态射 $x_\alpha:Z\to X_y$，因为
% P11-E0193: frozen source says "isomorphisms classes"; translated as the intended "isomorphism classes".
$X_y$ 是
$(\Sch/U)_{fppf} \times_{y, \mathcal{Y}, F} \mathcal{X}$
的对象同构类函子；参见《代数叠》引理
\ref{algebraic-lemma-characterize-representable-by-space}。图示为
\begin{equation}
\label{criteria-equation-relative-map}
\vcenter{
\xymatrix{
Z \ar[r]_{x_\alpha} \ar[rd] & X_y \ar[d] \\
& U
}
}
\quad\quad
\vcenter{
\xymatrix{
(\Sch/Z)_{fppf} \ar[rd] \ar[r]_-{x, \alpha} &
(\Sch/U)_{fppf} \times_{y, \mathcal{Y}, F} \mathcal{X} \ar[r] \ar[d] &
\mathcal{X} \ar[d]^F \\
& (\Sch/U)_{fppf} \ar[r]^y & \mathcal{Y}
}
}
\end{equation}
注意，如果
$(f, g, b, a) : (U, Z, y, x, \alpha) \to (U', Z', y', x', \alpha')$
是 $\mathcal{H}_d$ 的对象之间的态射，则态射
$x'_{\alpha'}:Z'\to X'_{y'}$ 是态射 $x_\alpha$ 沿 $g:U'\to U$
的基变换（略去细节）。

\medskip\noindent
现在进一步假设 $F$ 平坦且局部有限表示。在这种情形下，定义全子范畴
$$
\mathcal{H}_{d, lci}(\mathcal{X}/\mathcal{Y}) \subset
\mathcal{H}_d(\mathcal{X}/\mathcal{Y})
$$
，它由 $\mathcal{H}_d(\mathcal{X}/\mathcal{Y})$ 中满足下列条件的对象
$(U,Z,y,x,\alpha)$ 组成：相应态射 $x_\alpha:Z\to X_y$ 非分歧，
且为局部完全交态射（定义参见《空间的态射》定义
\ref{spaces-morphisms-definition-unramified} 和《空间态射进阶》定义
\ref{spaces-more-morphisms-definition-lci}）。

\begin{lemma}
\label{criteria-lemma-lci-locus-stack-in-groupoids}
设 $S$ 是概形。固定 $(\Sch/S)_{fppf}$ 上群胚叠的 $1$-态射
$F:\mathcal{X}\longrightarrow\mathcal{Y}$。假设 $F$ 可由代数空间表示、
平坦且局部有限表示。则 $\mathcal{H}_{d,lci}(\mathcal{X}/\mathcal{Y})$
是群胚叠，且嵌入函子
$$
\mathcal{H}_{d, lci}(\mathcal{X}/\mathcal{Y})
\longrightarrow
\mathcal{H}_d(\mathcal{X}/\mathcal{Y})
$$
可表示且为开浸入。
\end{lemma}

\begin{proof}
设 $\Xi=(U,Z,y,x,\alpha)$ 是 $\mathcal{H}_d$ 的对象。由
(\ref{criteria-equation-relative-map}) 后的评注，$\Xi$ 沿 $U'\to U$ 的拉回属于
$\mathcal{H}_{d,lci}(\mathcal{X}/\mathcal{Y})$，当且仅当 $x_\alpha$
的基变换非分歧且为局部完全交态射。注意，$Z\to U$ 有限局部自由
（因而平坦、局部有限表示且万有闭），而由关于 $F$ 的假设，$X_y\to U$
平坦且局部有限表示。于是《空间态射进阶》引理
\ref{spaces-more-morphisms-lemma-where-unramified} 和
\ref{spaces-more-morphisms-lemma-where-lci}
% P11-E0194: frozen source says "imply exists"; "there" is missing. Translated naturally.
表明存在开子概形 $W\subset U$，使态射 $U'\to U$ 通过 $W$ 分解，
当且仅当 $x_\alpha$ 沿 $U'\to U$ 的基变换非分歧且为局部完全交态射。
这意味着
$$
(\Sch/U)_{fppf}
\times_{\Xi, \mathcal{H}_d(\mathcal{X}/\mathcal{Y})}
\mathcal{H}_{d, lci}(\mathcal{X}/\mathcal{Y})
$$
可由 $W$ 表示。因此，引理的最后一个断言成立。第一个断言（即
$\mathcal{H}_{d,lci}(\mathcal{X}/\mathcal{Y})$ 是群胚叠）由此和
《代数叠》引理 \ref{algebraic-lemma-open-fibred-category-is-algebraic}
推出。
\end{proof}

\noindent
局部完全交态射“局部无阻碍”。这一点在远比本章所需特殊情形更一般的
范围内成立。

\begin{lemma}
\label{criteria-lemma-lci-unobstructed}
设 $U\subset U'$ 是仿射概形的一阶加厚，$X'$ 是 $U'$ 上平坦的代数空间。
令 $X=U\times_{U'}X'$。设 $Z\to U$ 是次数为 $d$ 的有限局部自由态射。
最后，设 $f:Z\to X$ 非分歧且为局部完全交态射。则存在交换图
$$
\xymatrix{
(Z \subset Z') \ar[rd] \ar[rr]_{(f, f')} & & (X \subset X') \ar[ld] \\
& (U \subset U')
}
$$
，其中各项为 $U'$ 上的代数空间，且 $Z'\to U'$ 是次数为 $d$ 的有限局部自由
态射，并满足 $Z=U\times_{U'}Z'$。
\end{lemma}

\begin{proof}
由《空间态射进阶》引理 \ref{spaces-more-morphisms-lemma-unramified-lci}，
非分歧态射 $Z\to X$ 的余法层 $\mathcal{C}_{Z/X}$ 是有限局部自由
$\mathcal{O}_Z$-模；又由《空间态射进阶》引理
\ref{spaces-more-morphisms-lemma-transitivity-conormal-lci}，有余法层的正合列
$$
0 \to i^*\mathcal{C}_{X/X'} \to
\mathcal{C}_{Z/X'} \to
\mathcal{C}_{Z/X} \to 0
$$
。由于 $Z$ 仿射，此列分裂。选取分裂
$$
\mathcal{C}_{Z/X'} = i^*\mathcal{C}_{X/X'} \oplus \mathcal{C}_{Z/X}
$$
设 $Z\subset Z''$ 是 $Z$ 在 $X'$ 上的万有一阶加厚（参见《空间态射进阶》第
\ref{spaces-more-morphisms-section-universal-thickening} 节）。
% P11-E0195: frozen source's naming construction "Denote I ... the" omits "by"; translated naturally.
记 $\mathcal{I}\subset\mathcal{O}_{Z''}$ 为与 $Z\subset Z''$ 相应的拟凝聚理想层。
% P11-E0196: frozen source says "we have C is I"; the intended identification is translated naturally.
按定义，把 $\mathcal{I}$ 视为 $Z$ 上的层时，它就是 $\mathcal{C}_{Z/X'}$。
因此，上述分裂确定分裂
$$
\mathcal{I} = i^*\mathcal{C}_{X/X'} \oplus \mathcal{C}_{Z/X}
$$
把 $\mathcal{C}_{Z/X}\subset\mathcal{I}$ 视为 $Z''$ 上的拟凝聚理想层，
令 $Z'\subset Z''$ 为它截出的闭子概形。显然，$Z'$ 是 $Z$ 的一阶加厚，
并得到引理陈述中的一阶加厚交换图。

\medskip\noindent
由于 $X'\to U'$ 平坦且 $X=U\times_{U'}X'$，可知
$\mathcal{C}_{X/X'}$ 是 $\mathcal{C}_{U/U'}$ 到 $X$ 的拉回；参见
《空间态射进阶》引理 \ref{spaces-more-morphisms-lemma-deform}。
注意，按构造，$\mathcal{C}_{Z/Z'}=i^*\mathcal{C}_{X/X'}$，故
$\mathcal{C}_{Z/Z'}$ 同构于 $\mathcal{C}_{U/U'}$ 到 $Z$ 的拉回。
再次应用《空间态射进阶》引理 \ref{spaces-more-morphisms-lemma-deform}
（或其概形版本，参见《态射进阶》引理 \ref{more-morphisms-lemma-deform}），
可知 $Z'\to U'$ 平坦且 $Z=U\times_{U'}Z'$。最后，《态射进阶》引理
\ref{more-morphisms-lemma-deform-property} 表明 $Z'\to U'$ 是次数为 $d$
的有限局部自由态射。
\end{proof}

\begin{lemma}
\label{criteria-lemma-lci-formally-smooth}
设 $F:\mathcal{X}\to\mathcal{Y}$ 是 $(\Sch/S)_{fppf}$ 上群胚叠的
$1$-态射。假设 $F$ 可由代数空间表示、平坦且局部有限表示。则
$$
p : \mathcal{H}_{d, lci}(\mathcal{X}/\mathcal{Y}) \to \mathcal{Y}
$$
在对象上形式光滑。
\end{lemma}

\begin{proof}
须证明如下命题：给定
\begin{enumerate}
\item $\mathcal{H}_{d,lci}(\mathcal{X}/\mathcal{Y})$ 在仿射概形 $U$
上的对象 $(U,Z,y,x,\alpha)$；
\item 一阶加厚 $U\subset U'$；以及
\item $\mathcal{Y}$ 在 $U'$ 上的对象 $y'$，满足 $y'|_U=y$，
\end{enumerate}
则存在 $\mathcal{H}_{d,lci}(\mathcal{X}/\mathcal{Y})$ 在 $U'$ 上的对象
$(U',Z',y',x',\alpha')$，满足 $Z=U\times_{U'}Z'$、$x=x'|_Z$ 以及
$\alpha=\alpha'|_U$。最后两个等式恰好保证
(\ref{criteria-equation-formally-smooth}) 交换。

\medskip\noindent
考虑公式 (\ref{criteria-equation-relative-map}) 中构造的态射
$x_\alpha:Z\to X_y$。类似地，
% P11-E0197: frozen source's naming construction "Denote similarly X' the ..." omits "by"; translated naturally.
记 $X'_{y'}$ 为 $U'$ 上表示 $2$-纤维积
$(\Sch/U')_{fppf} \times_{y', \mathcal{Y}, F} \mathcal{X}$.
按假设，态射 $X'_{y'}\to U'$ 平坦（且局部有限表示）。由于
$y'|_U=y$，可知 $X_y=U\times_{U'}X'_{y'}$。因此，可以应用引理
\ref{criteria-lemma-lci-unobstructed}，找到次数为 $d$ 的有限局部自由态射
$Z'\to U'$，满足 $Z=U\times_{U'}Z'$，并使 $Z'\to X'_{y'}$ 延拓
$x_\alpha$。按构造，态射 $Z'\to X'_{y'}$ 对应于二元组 $(x',\alpha')$。
显然，$(U',Z',y',x',\alpha')$ 是
$\mathcal{H}_d(\mathcal{X}/\mathcal{Y})$ 在 $U'$ 上的对象，并满足
$Z=U\times_{U'}Z'$、$x=x'|_Z$ 和 $\alpha=\alpha'|_U$。
由引理 \ref{criteria-lemma-lci-locus-stack-in-groupoids} 已知
$\mathcal{H}_{d,lci}(\mathcal{X}/\mathcal{Y})\subset
\mathcal{H}_d(\mathcal{X}/\mathcal{Y})$ 是“开子叠”，故
$(U',Z',y',x',\alpha')$ 如愿是
$\mathcal{H}_{d,lci}(\mathcal{X}/\mathcal{Y})$ 的对象。
\end{proof}

\begin{lemma}
\label{criteria-lemma-lci-surjective}
设 $F:\mathcal{X}\to\mathcal{Y}$ 是 $(\Sch/S)_{fppf}$ 上群胚叠的
$1$-态射。假设 $F$ 可由代数空间表示、平坦、为满且局部有限表示。则
$$
\coprod\nolimits_{d \geq 1} \mathcal{H}_{d, lci}(\mathcal{X}/\mathcal{Y})
\longrightarrow
\mathcal{Y}
$$
在对象上为满。
\end{lemma}

\begin{proof}
只需证明如下命题：对任意域 $k$ 和 $\mathcal{Y}$ 在 $\Spec(k)$ 上的对象
$y$，存在整数 $d\geq1$ 和
$\mathcal{H}_{d,lci}(\mathcal{X}/\mathcal{Y})$ 的对象
$(U,Z,y,x,\alpha)$，其中 $U=\Spec(k)$。此时，$p$ 以无需域扩张的
强意义在对象上为满。

\medskip\noindent
% P11-E0198: frozen source's naming construction "Denote X_y the ..." omits "by"; translated naturally.
记 $X_y$ 为 $U=\Spec(k)$ 上表示 $2$-纤维积
% P11-E0199: frozen source uses U' and y' in the following fibre product although this proof has introduced U and y only; frozen math is preserved.
$(\Sch/U')_{fppf} \times_{y', \mathcal{Y}, F} \mathcal{X}$.
的代数空间。按假设，态射 $X_y\to\Spec(k)$ 为满且局部有限表示
（并且平坦）。特别地，$X_y$ 非空。选取非空仿射概形 $V$ 和平展态射
$V\to X_y$。注意，$V\to\Spec(k)$（平坦）、为满且局部有限表示
（由《空间的态射》定义
\ref{spaces-morphisms-definition-locally-finite-presentation}）。
选取闭点 $v\in V$，使 $V\to\Spec(k)$ 在该点为 Cohen--Macaulay
（即 $V$ 在 $v$ 处为 Cohen--Macaulay）；参见《态射进阶》引理
\ref{more-morphisms-lemma-flat-finite-presentation-CM-open}。应用
《态射进阶》引理 \ref{more-morphisms-lemma-slice-CM}，得到正则浸入
$Z\to V$，其中 $Z=\{v\}$。这蕴含 $Z\to V$ 是闭浸入。此外，
$Z\to\Spec(k)$ 有限（例如由《代数》引理
\ref{algebra-lemma-isolated-point}）。故 $Z\to\Spec(k)$ 是某个次数 $d$
的有限局部自由态射。现在，$Z\to X_y$ 是闭浸入后接平展态射的复合，
因而非分歧（参见《空间的态射》引理
\ref{spaces-morphisms-lemma-composition-unramified},
\ref{spaces-morphisms-lemma-etale-unramified} 和
\ref{spaces-morphisms-lemma-immersion-unramified}）。
最后，$Z \to X_y$ 是局部完全交态射，因为它是概形的正则浸入与
代数空间的平展态射的复合（参见《态射进阶》引理
\ref{more-morphisms-lemma-regular-immersion-lci}，
《空间的态射》引理 \ref{spaces-morphisms-lemma-etale-smooth} 和
\ref{spaces-morphisms-lemma-smooth-syntomic}，以及
《空间的态射进阶》引理
\ref{spaces-more-morphisms-lemma-flat-lci} 和
\ref{spaces-more-morphisms-lemma-composition-lci}）。
态射 $Z \to X_y$ 对应于 $\mathcal{X}$ 在 $Z$ 上的一个对象 $x$，
以及同构 $\alpha : y|_Z \to F(x)$。我们由此得到
$\mathcal{H}_d(\mathcal{X}/\mathcal{Y})$ 的一个对象
$(U, Z, y, x, \alpha)$。由上面对态射 $Z \to X_y$ 的论述可知，
它实际上是子范畴
$\mathcal{H}_{d, lci}(\mathcal{X}/\mathcal{Y})$ 的对象，结论得证。
\end{proof}




















\section{代数叠的自举}
\label{criteria-section-bootstrap}

\noindent
下述定理是本章的主要结果之一。

\begin{theorem}
\label{criteria-theorem-bootstrap}
\begin{slogan}
Artin 关于平坦群胚可表示性的定理
\end{slogan}
设 $S$ 是概形，$F : \mathcal{X} \to \mathcal{Y}$ 是
$(\Sch/S)_{fppf}$ 上群胚叠的一个 $1$-态射。若
\begin{enumerate}
\item $\mathcal{X}$ 可由代数空间表示，并且
\item $F$ 可由代数空间表示、为满、平坦且局部有限表示，
\end{enumerate}
则 $\mathcal{Y}$ 是代数叠。
\end{theorem}

\begin{proof}
由引理 \ref{criteria-lemma-flat-finite-presentation-surjective-diagonal}，
$\mathcal{Y}$ 的对角可由代数空间表示。因此，我们只需验证存在
$(\Sch/S)_{fppf}$ 上群胚叠的一个 $1$-态射
$f : \mathcal{V} \to \mathcal{Y}$，使得 $\mathcal{V}$ 可表示，
而 $f$ 为满且光滑。由引理 \ref{criteria-lemma-hilbert-stack-relative-space}，
我们知道
$$
\coprod\nolimits_{d \geq 1} \mathcal{H}_d(\mathcal{X}/\mathcal{Y})
$$
是代数叠。由引理 \ref{criteria-lemma-lci-locus-stack-in-groupoids} 和
《代数叠》引理 \ref{algebraic-lemma-open-fibred-category-is-algebraic}，
可知
$$
\coprod\nolimits_{d \geq 1} \mathcal{H}_{d, lci}(\mathcal{X}/\mathcal{Y})
$$
也是代数叠。选取 $(\Sch/S)_{fppf}$ 上一个可表示的群胚叠
$\mathcal{V}$，以及一个满且光滑的 $1$-态射
$$
\mathcal{V}
\longrightarrow
\coprod\nolimits_{d \geq 1} \mathcal{H}_{d, lci}(\mathcal{X}/\mathcal{Y}).
$$
我们断言复合态射
$$
\mathcal{V}
\longrightarrow
\coprod\nolimits_{d \geq 1} \mathcal{H}_{d, lci}(\mathcal{X}/\mathcal{Y})
\longrightarrow
\mathcal{Y}
$$
光滑且为满，这就完成定理的证明。事实上，光滑性将由引理
\ref{criteria-lemma-limit-preserving} 和 \ref{criteria-lemma-lci-formally-smooth}
推出，而满性将由引理 \ref{criteria-lemma-lci-surjective} 推出。
我们在下一段详述其细节。

\medskip\noindent
由构造，$\mathcal{V} \to
\coprod\nolimits_{d \geq 1} \mathcal{H}_{d, lci}(\mathcal{X}/\mathcal{Y})$
可由代数空间表示、为满且光滑（因而由一般原理
《代数叠》引理
\ref{algebraic-lemma-representable-transformations-property-implication}
和《空间的态射进阶》引理
\ref{spaces-more-morphisms-lemma-smooth-formally-smooth}，
它也局部有限表示且形式光滑）。应用引理
\ref{criteria-lemma-representable-by-spaces-limit-preserving}、
\ref{criteria-lemma-representable-by-spaces-formally-smooth} 和
\ref{criteria-lemma-representable-by-spaces-surjective}
可知 $\mathcal{V} \to
\coprod\nolimits_{d \geq 1} \mathcal{H}_{d, lci}(\mathcal{X}/\mathcal{Y})$
在对象上保极限、在对象上形式光滑且在对象上为满。$1$-态射
$\coprod\nolimits_{d \geq 1} \mathcal{H}_{d, lci}(\mathcal{X}/\mathcal{Y})
\to \mathcal{Y}$ 满足：
\begin{enumerate}
\item 在对象上保极限：对
$\mathcal{H}_d(\mathcal{X}/\mathcal{Y}) \to \mathcal{Y}$ 而言，
这是引理 \ref{criteria-lemma-limit-preserving}；将其与引理
\ref{criteria-lemma-lci-locus-stack-in-groupoids}、
\ref{criteria-lemma-open-immersion-limit-preserving} 和
\ref{criteria-lemma-composition-limit-preserving}
结合，便得到
$\mathcal{H}_{d, lci}(\mathcal{X}/\mathcal{Y}) \to \mathcal{Y}$
的相应结论；
\item 由引理 \ref{criteria-lemma-lci-formally-smooth}，在对象上形式光滑；并且
\item 由引理 \ref{criteria-lemma-lci-surjective}，在对象上为满。
\end{enumerate}
应用引理 \ref{criteria-lemma-composition-limit-preserving}、
\ref{criteria-lemma-composition-formally-smooth} 和
\ref{criteria-lemma-composition-surjective}
可知复合态射 $\mathcal{V} \to \mathcal{Y}$ 在对象上保极限、
在对象上形式光滑且在对象上为满。再应用引理
\ref{criteria-lemma-representable-by-spaces-limit-preserving}、
\ref{criteria-lemma-representable-by-spaces-formally-smooth} 和
\ref{criteria-lemma-representable-by-spaces-surjective}
可知 $\mathcal{V} \to \mathcal{Y}$ 局部有限表示、形式光滑且为满。
最后，通过一般原理《代数叠》引理
\ref{algebraic-lemma-representable-transformations-property-implication}
应用无穷小提升判据（《空间的态射进阶》引理
\ref{spaces-more-morphisms-lemma-smooth-formally-smooth}），可知
$\mathcal{V} \to \mathcal{Y}$ 光滑，结论得证。
\end{proof}








\section{应用}
\label{criteria-section-applications}

\noindent
我们的首要任务是证明，与“平坦且局部有限表示的群胚”相伴的商叠
$[U/R]$ 是代数叠。商叠的定义见《空间中的群胚》定义
\ref{spaces-groupoids-definition-quotient-stack}。下述预备引理是
《代数叠》引理
\ref{algebraic-lemma-smooth-quotient-smooth-presentation} 的对应版本。

\begin{lemma}
\label{criteria-lemma-flat-quotient-flat-presentation}
设 $S$ 是 $\Sch_{fppf}$ 中的概形，
$(U, R, s, t, c)$ 是 $S$ 上代数空间中的群胚。
假设 $s, t$ 平坦且局部有限表示。则态射
$\mathcal{S}_U \to [U/R]$ 平坦、局部有限表示且为满。
\end{lemma}

\begin{proof}
设 $T$ 是概形，$x : (\Sch/T)_{fppf} \to [U/R]$ 是一个
$1$-态射。我们必须证明投影
$$
\mathcal{S}_U \times_{[U/R]} (\Sch/T)_{fppf}
\longrightarrow
(\Sch/T)_{fppf}
$$
为满、平坦且局部有限表示。我们已经知道左边可由一个代数空间 $F$
表示，参见《代数叠》引理
\ref{algebraic-lemma-diagonal-quotient-stack} 和
\ref{algebraic-lemma-representable-diagonal}。因此，我们必须证明相应的
代数空间态射 $F \to T$ 为满、局部有限表示且平坦。由于此处所考察的
代数空间态射性质关于 fppf 拓扑在靶上是局部的，我们可以在 $T$ 上
fppf 局部地检验这些性质。由构造，存在 $T$ 的一个 fppf 覆盖
$\{T_i \to T\}$，使得 $x|_{(\Sch/T_i)_{fppf}}$ 来自某个态射
$x_i : T_i \to U$。（注意，$F \times_T T_i$ 表示 $2$-纤维积
$\mathcal{S}_U \times_{[U/R]} (\Sch/T_i)_{fppf}$，故一切都与经由
$T_i \to T$ 的基变换相容。）因此，可假设 $x$ 来自
$x : T \to U$。在这种情形下，我们看到
$$
\mathcal{S}_U \times_{[U/R]} (\Sch/T)_{fppf}
=
(\mathcal{S}_U \times_{[U/R]} \mathcal{S}_U)
\times_{\mathcal{S}_U} (\Sch/T)_{fppf}
=
\mathcal{S}_R \times_{\mathcal{S}_U} (\Sch/T)_{fppf}
$$
其中第一个等号由《范畴》引理
\ref{categories-lemma-2-fibre-product-erase-factor} 给出，第二个等号由
《空间中的群胚》引理
\ref{spaces-groupoids-lemma-quotient-stack-2-cartesian} 给出。显然，
最后一个 $2$-纤维积由代数空间 $F = R \times_{s, U, x} T$ 表示；
投影 $R \times_{s, U, x} T \to T$ 是平坦且局部有限表示的代数空间态射
$s : R \to U$ 的基变换，因而平坦且局部有限表示。它也为满，因为
$s$ 有截面（即群胚的恒等态射 $e : U \to R$）。引理得证。
\end{proof}

\noindent
下面是本节的第一个主要结果。

\begin{theorem}
\label{criteria-theorem-flat-groupoid-gives-algebraic-stack}
设 $S$ 是 $\Sch_{fppf}$ 中的概形，
$(U, R, s, t, c)$ 是 $S$ 上代数空间中的群胚。
假设 $s, t$ 平坦且局部有限表示。则商叠 $[U/R]$ 是 $S$ 上的代数叠。
\end{theorem}

\begin{proof}
对态射
$$
(\Sch/U)_{fppf} \longrightarrow [U/R].
$$
我们检验定理 \ref{criteria-theorem-bootstrap} 的两个条件。第一个条件是显然的
（因为 $U$ 是代数空间）。第二个条件即引理
\ref{criteria-lemma-flat-quotient-flat-presentation}。
\end{proof}









\section{商叠何时是代数的？}
\label{criteria-section-quotient-algebraic}

\noindent
在《空间中的群胚》第 \ref{spaces-groupoids-section-stacks} 节中，
我们定义了与代数空间中的群胚 $(U, R, s, t, c)$ 相伴的商叠 $[U/R]$。
注意，$[U/R]$ 是一个群胚叠，其对角可由代数空间表示（参见
《自举》引理 \ref{bootstrap-lemma-quotient-stack-isom} 和
《代数叠》引理 \ref{algebraic-lemma-representable-diagonal}），并且存在
代数空间 $U$ 和一个 $1$-态射 $(\Sch/U)_{fppf} \to [U/R]$，它是如下
意义下的“fppf 满射”：它在对象同构类的预层上诱导的映射经层化后为满。
然而，一般而言 $[U/R]$ 并不一定是代数叠。这与定理
\ref{criteria-theorem-bootstrap} 并不矛盾，因为 $1$-态射
$(\Sch/U)_{fppf} \to [U/R]$ 未必平坦且局部有限表示。

\medskip\noindent
构造非代数商叠例子的最简方法，是考察形如 $[S/G]$ 的商，其中 $S$
是概形，而 $G$ 是 $S$ 上平凡作用于 $S$ 的群概形。确切地说，下面将看到
（引理 \ref{criteria-lemma-BG-algebraic}），若 $[S/G]$ 是代数的，则
$G \to S$ 必须平坦且局部有限表示。一个显式例子见《例子》第
\ref{examples-section-not-algebraic-stack} 节。

\begin{lemma}
\label{criteria-lemma-quotient-algebraic}
设 $S$ 是概形，$B$ 是 $S$ 上的代数空间，
$(U, R, s, t, c)$ 是 $B$ 上代数空间中的群胚。商叠 $[U/R]$
是代数叠，当且仅当存在代数空间态射 $g : U' \to U$，使得
\begin{enumerate}
\item 复合态射
$U' \times_{g, U, t} R \to R \xrightarrow{s} U$ 是层的满射；并且
\item 态射 $s', t' : R' \to U'$ 平坦且局部有限表示，其中
$(U', R', s', t', c')$ 是 $(U, R, s, t, c)$ 沿 $g$ 的限制。
\end{enumerate}
\end{lemma}

\begin{proof}
首先，假设 $g : U' \to U$ 满足 (1) 和 (2)。性质 (1) 蕴含
$[U'/R'] \to [U/R]$ 是等价，参见《空间中的群胚》引理
\ref{spaces-groupoids-lemma-quotient-stack-restrict-equivalence}。
由定理 \ref{criteria-theorem-flat-groupoid-gives-algebraic-stack}，商叠
$[U'/R']$ 是代数叠。因此 $[U/R]$ 也是代数叠，参见《代数叠》引理
\ref{algebraic-lemma-equivalent}。

\medskip\noindent
反过来，假设 $[U/R]$ 是代数叠。我们可以选取概形 $W$ 和一个满光滑
$1$-态射
$$
f : (\Sch/W)_{fppf} \longrightarrow [U/R].
$$
由 $2$-Yoneda 引理（《代数叠》第 \ref{algebraic-section-2-yoneda} 节），
这对应于 $[U/R]$ 在 $W$ 上的一个对象 $\xi$。由《空间中的群胚》引理
\ref{spaces-groupoids-lemma-quotient-stack-objects} 对 $[U/R]$ 的描述，
我们可以找到概形的一个满、平坦且局部有限表示的态射
$b : U' \to W$，使得 $\xi' = b^*\xi$ 对应于某个态射
$g : U' \to U$。注意，与 $\xi'$ 对应的 $1$-态射
$$
f' : (\Sch/U')_{fppf} \longrightarrow [U/R].
$$
为满、平坦且局部有限表示，参见《代数叠》引理
\ref{algebraic-lemma-composition-representable-transformations-property}。
因此，
$(\Sch/U')_{fppf} \times_{[U/R]} (\Sch/U')_{fppf}$
由代数空间
$$
\mathit{Isom}_{[U/R]}(\text{pr}_0^*\xi', \text{pr}_1^*\xi') =
(U' \times_S U')
\times_{(g \circ \text{pr}_0, g \circ \text{pr}_1), U \times_S U} R = R'
$$
表示（第一个等号见《空间中的群胚》引理
\ref{spaces-groupoids-lemma-quotient-stack-morphisms}；第二个等号是限制的
定义），并且经由 $s'$ 和 $t'$ 在 $U'$ 上都平坦且局部有限表示
（由基变换，参见《代数叠》引理
\ref{algebraic-lemma-base-change-representable-transformations-property}）。
由这一对 $R'$ 的描述以及《代数叠》引理
\ref{algebraic-lemma-map-space-into-stack}，我们得到一个典范的全忠实
$1$-态射 $[U'/R'] \to [U/R]$。因为 $f'$ 平坦、局部有限表示且为满，
这个 $1$-态射本质满（参见《叠》引理
\ref{stacks-lemma-characterize-essentially-surjective-when-ff}）；另一种证明
方法是应用《代数叠》注记
\ref{algebraic-remark-flat-fp-presentation}。最后，应用《空间中的群胚》引理
\ref{spaces-groupoids-lemma-quotient-stack-restrict-equivalence}，可知复合态射
$U' \times_{g, U, t} R \to R \xrightarrow{s} U$ 是层的满射。
\end{proof}

\begin{lemma}
\label{criteria-lemma-group-quotient-algebraic}
设 $S$ 是概形，$B$ 是 $S$ 上的代数空间，$G$ 是 $B$ 上的群代数空间。
设 $X$ 是 $B$ 上的代数空间，而 $a : G \times_B X \to X$ 是 $G$
在 $B$ 上对 $X$ 的一个作用。商叠 $[X/G]$ 是代数叠，当且仅当存在
代数空间态射 $\varphi : X' \to X$，使得
\begin{enumerate}
\item $G \times_B X' \to X$，$(g, x') \mapsto a(g, \varphi(x'))$，
是层的满射；并且
\item 由规则
$$
T \longmapsto \{(x'_1, g, x'_2) \in (X' \times_B G \times_B X')(T)
\mid \varphi(x'_1) = a(g, \varphi(x'_2))\}
$$
给出的代数空间 $X''$ 的两个投影 $X'' \to X'$ 都平坦且局部有限表示。
\end{enumerate}
\end{lemma}

\begin{proof}
本引理是引理 \ref{criteria-lemma-quotient-algebraic} 的一个特例。确切地说，
由《空间中的群胚》定义
\ref{spaces-groupoids-definition-quotient-stack}，商叠 $[X/G]$ 等于
代数空间中的群胚 $(X, G \times_B X, s, t, c)$ 的商叠
$[X/G \times_B X]$；该群胚与《空间中的群胚》引理
\ref{spaces-groupoids-lemma-groupoid-from-action} 中的群作用相伴。
要得到条件 (1)，还需作一个小观察。态射 $s : G \times_B X \to X$
是第二投影，而态射 $t :  G \times_B X \to X$ 是作用态射 $a$。
因此，引理 \ref{criteria-lemma-quotient-algebraic} 中的态射
$h : U' \times_{g, U, t} R \to R \xrightarrow{s} U$ 对应于当前设置中的态射
$$
X' \times_{\varphi, X, a} (G \times_B X) \xrightarrow{\text{pr}_1} X
$$
不过，由 $G$ 的取逆所给出的对称性，此态射同构于引理陈述中的态射
$$
(G \times_B X) \times_{\text{pr}_1, X, \varphi} X' \xrightarrow{a} X
$$
。细节从略。
\end{proof}

\begin{lemma}
\label{criteria-lemma-BG-algebraic}
\begin{slogan}
Gerbe 是代数的，当且仅当相伴的群平坦且局部有限表示
\end{slogan}
设 $S$ 是概形，$B$ 是 $S$ 上的代数空间，$G$ 是 $B$ 上的群代数空间。
赋予 $B$ 以 $G$ 的平凡作用。则商叠 $[B/G]$ 是代数叠，当且仅当
$G$ 在 $B$ 上平坦且局部有限表示。
\end{lemma}

\begin{proof}
若 $G$ 在 $B$ 上平坦且局部有限表示，则由定理
\ref{criteria-theorem-flat-groupoid-gives-algebraic-stack}，$[B/G]$ 是代数叠。

\medskip\noindent
反过来，假设 $[B/G]$ 是代数叠。由引理
\ref{criteria-lemma-group-quotient-algebraic} 以及作用的平凡性，可知存在代数空间
$B'$ 和态射 $B' \to B$，使得 (1) $B' \to B$ 是层的满射，并且
(2) 投影
$$
B' \times_B G \times_B B' \to B'
$$
平坦且局部有限表示。注意，$B' \to B$ 的基变换
$B' \times_B G \times_B B' \to G \times_B B'$ 也是层的满射。因此，
由《空间上的下降》引理 \ref{spaces-descent-lemma-curiosity}，投影
$G \times_B B' \to B'$ 平坦且局部有限表示。由 (1)，可以找到一个
fppf 覆盖 $\{B_i \to B\}$，使得 $B_i \to B$ 经 $B' \to B$ 分解。
于是，由基变换，$G \times_B B_i \to B_i$ 平坦且局部有限表示。由
《空间上的下降》引理
\ref{spaces-descent-lemma-descending-property-flat} 和
\ref{spaces-descent-lemma-descending-property-locally-finite-presentation}，
可知 $G \to B$ 平坦且局部有限表示。
\end{proof}

\noindent
稍后将看到，光滑 $S$-空间关于群代数空间 $G$ 的商叠是光滑的，
即使 $G$ 并不光滑亦然（《叠的态射》引理
\ref{stacks-morphisms-lemma-smooth-quotient-stack}）。




\section{\'etale 拓扑中的代数叠}
\label{criteria-section-stacks-etale}

\noindent
设 $S$ 是概形。我们可以研究大 fppf 位点 $(\Sch/S)_{fppf}$ 上的群胚叠，
也可以改为研究大 \'etale 位点 $(\Sch/S)_\etale$ 上的群胚叠。
《代数叠》第
\ref{algebraic-section-representable},
\ref{algebraic-section-2-yoneda},
\ref{algebraic-section-representable-morphism},
\ref{algebraic-section-split},
\ref{algebraic-section-representable-by-algebraic-spaces},
\ref{algebraic-section-morphisms-representable-by-algebraic-spaces},
\ref{algebraic-section-representable-properties} 和
\ref{algebraic-section-stacks}
节的全部内容，对 $(\Sch/S)_\etale$ 上群胚纤维化范畴同样有意义。因此，
在 \'etale 拓扑中工作会给出代数叠的第二种概念。这个概念先验地弱于
《代数叠》定义 \ref{algebraic-definition-algebraic-stack} 引入的概念，
因为 fppf 拓扑中的叠当然也是 \'etale 拓扑中的叠。不过，下述引理表明
这两个概念等价。

\begin{lemma}
\label{criteria-lemma-stacks-etale}
用 $\Sch_\alpha$ 表示 $\Sch_{fppf}$ 与 $\Sch_\etale$ 共同的底层范畴
（参见《叠上的层》第 \ref{stacks-sheaves-section-sheaves} 节和
《拓扑》注记 \ref{topologies-remark-choice-sites}）。设 $S$ 是
$\Sch_\alpha$ 的对象，并设
$$
p : \mathcal{X} \to \Sch_\alpha/S
$$
是满足下列性质的群胚纤维化范畴：
\begin{enumerate}
\item $\mathcal{X}$ 是 $(\Sch/S)_\etale$ 上的群胚叠；
\item 对角 $\Delta : \mathcal{X} \to \mathcal{X} \times \mathcal{X}$
可由代数空间表示\footnote{这里，“代数空间”既可以指 \'etale 拓扑中
对角可表示且具有概形给出的满 \'etale 覆盖的层，也可以指《代数空间》
定义 \ref{spaces-definition-algebraic-space} 所定义的代数空间。事实上，
由《自举》引理 \ref{bootstrap-lemma-spaces-etale}，两者没有区别。}；并且
\item 存在 $U \in \Ob(\Sch_\alpha/S)$ 和一个满光滑 $1$-态射
$(\Sch/U)_\etale \to \mathcal{X}$。
\end{enumerate}
则 $\mathcal{X}$ 是《代数叠》定义
\ref{algebraic-definition-algebraic-stack} 意义下的代数叠。
\end{lemma}

\begin{proof}
注意，本引理的性质 (2)、(3) 以及《代数叠》定义
\ref{algebraic-definition-algebraic-stack} 中相应的性质 (2)、(3)
与拓扑无关。这是因为这些性质只涉及群胚纤维化范畴的 $2$-纤维积概念、
群胚纤维化范畴的 $1$-态射和 $2$-态射、可由代数空间表示的群胚纤维化
范畴 $1$-态射这一概念，以及这种 $1$-态射为满且光滑的含义。因此，
我们只需证明，满足性质 (2)、(3) 的 \'etale 群胚叠 $\mathcal{X}$
也是 fppf 群胚叠。

\medskip\noindent
应用 (2)，令 $R$ 为表示下式的代数空间：
$$
(\Sch_\alpha/U) \times_\mathcal{X} (\Sch_\alpha/U)
$$
由 (3)，投影 $s, t : R \to U$ 光滑。完全仿照《代数叠》引理
\ref{algebraic-lemma-map-space-into-stack} 的证明，存在空间中的群胚
$(U, R, s, t, c)$ 和典范全忠实 $1$-态射
$[U/R]_\etale \to \mathcal{X}$，其中 $[U/R]_\etale$ 是群胚预层
$$
T \longmapsto (U(T), R(T), s(T), t(T), c(T))
$$
的 \'etale 叠化。断言：若 $V \to T$ 是从代数空间 $V$ 到概形 $T$
的满光滑态射，则存在一个加细覆盖 $\{V \to T\}$ 的 \'etale 覆盖
$\{T_i \to T\}$。这由《态射进阶》引理
\ref{more-morphisms-lemma-etale-dominates-smooth} 或更一般的《叠上的层》
引理 \ref{stacks-sheaves-lemma-surjective-flat-locally-finite-presentation}
推出。应用此断言，并完全仿照《代数叠》引理
\ref{algebraic-lemma-stack-presentation} 的论证，可知
$[U/R]_\etale \to \mathcal{X}$ 是等价。

\medskip\noindent
接着，用 $[U/R]$ 表示 fppf 拓扑中的商叠；由《代数叠》定理
\ref{algebraic-theorem-smooth-groupoid-gives-algebraic-stack}，它是代数叠。
于是有 $1$-态射
$$
U \to [U/R]_\etale \to [U/R].
$$
$U \to [U/R]_\etale \cong \mathcal{X}$ 与 $U \to [U/R]$ 都满且光滑
（前者由假设，后者由该定理），而且在这两种情形下，纤维积
$U \times_\mathcal{X} U$ 和 $U \times_{[U/R]} U$ 都由 $R$ 表示。因此，
$1$-态射 $[U/R]_\etale \to [U/R]$ 全忠实（因为商叠中的态射由到
$R$ 的态射给出，参见《空间中的群胚》第
\ref{spaces-groupoids-section-explicit-quotient-stacks} 节）。

\medskip\noindent
最后，对任意概形 $T$ 和态射 $t : T \to [U/R]$，纤维积
$V = T \times_{U/R} U$ 是在 $T$ 上满且光滑的代数空间。由上述断言，
存在一个 \'etale 覆盖 $\{T_i \to T\}_{i \in I}$ 以及 $T$ 上的态射
$T_i \to V$。这证明 $[U/R]$ 在 $T$ 上的对象 $t$ 在 \'etale 局部来自
$U$。由《叠》引理
\ref{stacks-lemma-characterize-essentially-surjective-when-ff}，可知
$[U/R]_\etale \to [U/R]$ 是 $(\Sch/S)_\etale$ 上群胚叠的等价。
证明完毕。
\end{proof}
% Generated mechanically from the frozen artin.tex target.
% Do not edit this generated file; edit the bound locale source instead.

% OK, start here.
%

\setcounter{chapter}{97}
\chapter{Artin 公理}


\phantomsection
\label{artin-section-phantom}





\section{引言}
\label{artin-section-introduction}

\noindent
本章讨论函子可由代数空间表示的 Artin 公理。推荐参考论文
\cite{ArtinI}, \cite{ArtinII}, \cite{ArtinVersal}.

\medskip\noindent
% P11-E0200: frozen source has plural coordinated subjects with singular "is awkward"; translated naturally.
本章中的某些记号、约定和术语颇为别扭，经验更丰富的读者甚至可能觉得
它们前后颠倒。这是有意为之。解释见《Quot 空间与 Hilbert 空间》第
\ref{quot-section-conventions} 节。

\medskip\noindent
设 $S$ 是局部 Noether 基概形，并设
$$
p : \mathcal{X} \longrightarrow (\Sch/S)_{fppf}
$$
是群胚纤维化范畴。设 $x_0$ 是 $\mathcal{X}$ 在域 $k$ 上的一个对象，
其中该域在 $S$ 上有限型。贯穿本章，与 $k$ 上的 $x_0$ 相伴的预形变范畴
（见《形式形变理论》定义
\ref{formal-defos-definition-predeformation-category}）
$$
\mathcal{F}_{\mathcal{X}, k, x_0}
\longrightarrow
\{\text{Artinian local }S\text{-algebras with residue field }k\}
$$
起着重要作用。我们为 $\mathcal{X}$ 引入 Rim--Schlessinger 条件 (RS)，
并证明该条件保证 $\mathcal{F}_{\mathcal{X}, k, x_0}$ 是形变范畴；也就是说，
$\mathcal{F}_{\mathcal{X}, k, x_0}$ 本身满足 (RS)。我们讨论用有限扩张
替换 $k$ 时 $\mathcal{F}_{\mathcal{X}, k, x_0}$ 如何变化，并讨论切空间。

\medskip\noindent
接着，我们讨论 $\mathcal{X}$ 的形式对象 $\xi = (\xi_n)$；它们是位于商
$R/\mathfrak m^n$ 上的对象所成的逆系统，其中 $R$ 是 Noether 完备局部
$S$-代数，其剩余域在 $S$ 上有限型。这等价于对某个 $x_0$ 和 $k$，
在 $\mathcal{F}_{\mathcal{X}, k, x_0}$ 中有一个形式对象。若 $\mathcal{X}$
在 $R$ 上有一个对象可产生该逆系统，就称形式对象为有效的。若
$\mathcal{X}$ 的形式对象产生 $\mathcal{F}_{\mathcal{X}, k, x_0}$ 的通用
形式对象，就称它是通用的。最后，给定有限型 $S$-概形 $U$、
$\mathcal{X}$ 在 $U$ 上的对象 $x$ 和闭点 $u_0 \in U$，若完备局部环
$\mathcal{O}_{U, u_0}^\wedge$ 上诱导的形式对象是通用的，就称 $x$
在 $u_0$ 处通用。

\medskip\noindent
完成这些准备后，可以陈述 Artin 的著名定理：若下列条件成立，则
$\mathcal{X}$ 是代数叠：
\begin{enumerate}
\item 对所有 $s \in S$，$\mathcal{O}_{S, s}$ 都是 G-环\footnote{这意味着
$S$ 是 G-概形，参见《性质》第 \ref{properties-section-G-scheme} 节。}；
\item $\Delta : \mathcal{X} \to \mathcal{X} \times \mathcal{X}$
可由代数空间表示；
\item $\mathcal{X}$ 是 \'etale 拓扑的叠；
\item $\mathcal{X}$ 保极限；
\item $\mathcal{X}$ 满足 (RS)；
\item 形变范畴 $\mathcal{F}_{\mathcal{X}, k, x_0}$ 的切空间和无穷小
自同构空间是有限维的；
\item 形式对象有效；
\item $\mathcal{X}$ 满足通用性的开性。
\end{enumerate}
这就是引理 \ref{artin-lemma-diagonal-representable}；一个略有改进的版本另见
命题 \ref{artin-proposition-second-diagonal-representable}。还有一个类似命题刻画
哪些函子 $F : (\Sch/S)_{fppf}^{opp} \to \textit{Sets}$ 是代数空间，
参见第 \ref{artin-section-algebraic-spaces} 节。

\medskip\noindent
下面粗略概述 Artin 定理的证明。首先，利用 (RS) 以及切空间和自同构空间
的有限维性，证明存在充分多的通用形式对象；例如参见《形式形变理论》引理
\ref{formal-defos-lemma-minimal-groupoid-in-functors-construction}。按假设，
这些形式对象有效。有效形式对象可由有限型 $S$-概形 $U$ 上的对象 $x$
“逼近”；参见引理 \ref{artin-lemma-approximate}。
% P11-E0201: frozen source omits "that" before "the local rings of S are G-rings"; translated as the intended coordination.
这一逼近利用 $S$ 的局部环都是 G-环以及 $\mathcal{X}$ 保极限这两个事实；
这也许是证明中最困难的部分，因为它依赖一般 N\'eron 去奇异化，用
Hensel 化中的解去逼近 Noether 局部 G-环上代数方程的形式解。接着，
通用性的开性蕴含：缩小 $U$ 后，可以假设 $x$ 在 $U$ 的每个闭点处通用。
完成这一切后，我们证明 $U \to \mathcal{X}$ 是光滑态射。取充分多的
$U \to \mathcal{X}$，即可得到 $\mathcal{X}$ 的一个“光滑图册”，从而证明
$\mathcal{X}$ 是代数叠。

\medskip\noindent
对给定的群胚纤维化范畴 $\mathcal{X}$ 检验 Artin 公理时，最困难的一步
通常是验证通用性的开性。以下讨论假设 $\mathcal{X}/S$ 已满足上列其他条件。
本章给出两种方法，用以证明 $\mathcal{X}$ 满足通用性的开性：
\begin{enumerate}
\item 第一种方法是假设一个称为 (RS*) 的更强 Rim--Schlessinger 条件，
并假设更强版本的形式有效性；本质上，它要求加厚逆系统上的对象有效。
事实证明，在这些假设下，通用性的开性自动成立；参见引理
\ref{artin-lemma-SGE-implies-openness-versality}。
% P11-E0202: frozen source says "that category of all schemes" where the definite article "the" is intended.
请注意，这里至关重要地用到了 $\mathcal{X}$ 定义在 $S$ 上所有概形的范畴
之上，而不只定义在 Noether 概形的范畴之上！
\item 第二种方法依循 Artin，要求 $\mathcal{X}$ 配备障碍理论。若该障碍理论
在适当意义下“与积可交换”，则 $\mathcal{X}$ 满足通用性的开性；参见引理
\ref{artin-lemma-get-openness-obstruction-theory}。
\end{enumerate}
障碍理论可以用许多不同方式公理化；文献中确实可见许多变体，往往针对特定
模叠作了调整。引理 \ref{artin-lemma-dual-openness} 说明一个使用导出范畴的变体；
它往往自然产生于文献中的形变理论计算。

\medskip\noindent
在第 \ref{artin-section-algebraic-spaces-noetherian} 节中，我们讨论需要作何修改，
才能把结果用于定义在 $S$ 上局部 Noether 概形范畴
$(\textit{Noetherian}/S)_\etale$ 上的函子。

\medskip\noindent
在本章最后一节，作为 Artin 公理的应用，我们证明 Artin 关于收缩存在性的
定理；参见第 \ref{artin-section-contractions} 节。该定理粗略地说：给定
% P11-E0203: frozen source says "X' separated of finite type"; translated with the intended conjunction.
在 $S$ 上分离且有限型的代数空间 $X'$、闭子集 $T' \subset |X'|$，以及形式修改
$$
\mathfrak{f} : X'_{/T'} \longrightarrow \mathfrak{X}
$$
其中 $\mathfrak{X}$ 是 $S$ 上的 Noether 形式代数空间，则存在一个真态射
$f : X' \to X$，“实现该收缩”。这指的是：存在等同
$\mathfrak{X} = X_{/T}$，使得
$\mathfrak{f} = f_{/T'} : X'_{/T'} \to X_{/T}$，其中 $T = f(T')$；
而且 $f$ 在 $X \setminus T$ 上是同构。证明的做法是在 $S$ 上局部 Noether
概形的范畴上定义函子 $F$，并证明 $F$ 满足 Artin 公理。有趣的是，在
Artin 公理的这个应用中，通用性的开性并非最难证明之处；
% P11-E0204: frozen source joins two independent clauses with a comma before "instead"; translated as two coordinated clauses.
反而是证明 $F$ 保极限需要大量工作和预备结果。




\section{约定}
\label{artin-section-conventions}

\noindent
本章采用的约定与《代数叠》一章中的约定相同；参见《代数叠》第
\ref{algebraic-section-conventions} 节。本章中的基概形 $S$ 往往是局部
Noether 的（不过在陈述结果时，我们总会重申这一条件）。





\section{预形变范畴}
\label{artin-section-predeformation-categories}

\noindent
设 $S$ 是局部 Noether 基概形，并设
$$
p : \mathcal{X} \longrightarrow (\Sch/S)_{fppf}
$$
是群胚纤维化范畴。设 $k$ 是域，而 $\Spec(k) \to S$ 是有限型态射
（参见《态射》引理 \ref{morphisms-lemma-point-finite-type}）。有时，
我们简称 {\it $k$ 是 $S$ 上有限型的域}。设 $x_0$ 是 $\mathcal{X}$
中位于 $\Spec(k)$ 上的对象。给定 $S$、$\mathcal{X}$、$k$ 和 $x_0$，
我们将构造《形式形变理论》定义
\ref{formal-defos-definition-predeformation-category} 意义下的预形变范畴。
该构造类似于《形式形变理论》注记
\ref{formal-defos-remark-localize-cofibered-groupoid} 的构造。

\medskip\noindent
首先，由《态射》引理 \ref{morphisms-lemma-point-finite-type}，可以选取仿射开集
$\Spec(\Lambda) \subset S$，使 $\Spec(k) \to S$ 经 $\Spec(\Lambda)$ 分解，
且相伴的环同态 $\Lambda \to k$ 有限。这给出范畴 $\mathcal{C}_\Lambda$；
参见《形式形变理论》定义 \ref{formal-defos-definition-CLambda}。在典范等价
意义下，范畴 $\mathcal{C}_\Lambda$ 不依赖于 $S$ 的仿射开集
$\Spec(\Lambda)$ 的选择。确切地说，$\mathcal{C}_\Lambda$ 等价于下述分解
所成范畴的反范畴：
\begin{equation}
\label{artin-equation-factor}
\Spec(k) \to \Spec(A) \to S
\end{equation}
这里要求 $A$ 是 Artin 局部环，并且 $\Spec(k) \to \Spec(A)$ 对应于
环同态 $A \to k$，它把 $k$ 等同于 $A$ 的剩余域。

\medskip\noindent
令 $\mathcal{F} = \mathcal{F}_{\mathcal{X}, k, x_0}$ 为如下范畴：
\begin{enumerate}
\item 对象是 $\mathcal{X}$ 的态射 $x_0 \to x$，其中
$p(x) = \Spec(A)$，$A$ 是 Artin 局部环，而
$p(x_0) \to p(x) \to S$ 是形如 (\ref{artin-equation-factor}) 的分解；并且
\item 态射 $(x_0 \to x) \to (x_0 \to x')$ 是 $\mathcal{X}$ 中的交换图
$$
\xymatrix{
x & & x' \ar[ll] \\
& x_0 \ar[lu] \ar[ru]
}
$$
。（注意箭头方向的反转。）
\end{enumerate}
若 $x_0 \to x$ 是 $\mathcal{F}$ 的对象，写 $p(x) = \Spec(A)$，便得到
$\mathcal{C}_\Lambda$ 的对象 $A$。我们常说 $x_0 \to x$ 或 $x$ 位于
$A$ 上。$\mathcal{F}$ 中从位于 $A$ 上的对象 $x_0 \to x$ 到位于 $A'$
上的对象 $x_0 \to x'$ 的态射，对应于 $\mathcal{X}$ 的态射
$x' \to x$，从而对应于态射
$p(x' \to x) : \Spec(A') \to \Spec(A)$；后者又对应于环同态
$A \to A'$。由于 $\mathcal{X}$ 是 $S$ 上概形范畴之上的范畴，
% P11-E0205: frozen source omits the article "a" before Lambda-algebra homomorphism; translated naturally.
可知 $A \to A'$ 是 $\Lambda$-代数同态。于是得到函子
\begin{equation}
\label{artin-equation-predeformation-category}
p : \mathcal{F} = \mathcal{F}_{\mathcal{X}, k, x_0}
\longrightarrow
\mathcal{C}_\Lambda.
\end{equation}
我们用记号 $\mathcal{F}(A)$ 表示 $\mathcal{C}_\Lambda$ 的对象 $A$ 上的
纤维范畴。$\mathcal{F}(A)$ 的对象就是 $\mathcal{X}$ 的态射
$x_0 \to x$，使得 $x$ 位于 $\Spec(A)$ 上，而 $x_0 \to x$ 位于
$\Spec(k) \to \Spec(A)$ 上。

\begin{lemma}
\label{artin-lemma-predeformation-category}
上面定义的函子 $p : \mathcal{F} \to \mathcal{C}_\Lambda$ 是预形变范畴。
\end{lemma}

\begin{proof}
须证明 (a) $\mathcal{F}$ 在 $\mathcal{C}_\Lambda$ 上以群胚为纤维余纤维化；
并且 (b) $\mathcal{F}(k)$ 等价于仅有一个对象和一个态射的范畴。

\medskip\noindent
(a) 的证明。由于 $\mathcal{X}$ 的纤维范畴是群胚，$\mathcal{F}$ 在
$\mathcal{C}_\Lambda$ 上的纤维范畴也是群胚。设 $A \to A'$ 是
$\mathcal{C}_\Lambda$ 的态射，而 $x_0 \to x$ 是 $\mathcal{F}(A)$
的对象。由于 $\mathcal{X}$ 以群胚为纤维纤维化，可以找到位于
$\Spec(A') \to \Spec(A)$ 上的态射 $x' \to x$。
% P11-E0206: frozen source says "is equal the given map"; "to" is missing. Translated naturally.
由于复合 $A \to A' \to k$ 等于给定映射 $A \to k$，由拉回在同构意义下
的唯一性可知，$x'$ 经 $\Spec(k) \to \Spec(A')$ 的拉回是 $x_0$；
也就是说，存在位于 $\Spec(k) \to \Spec(A')$ 上的态射 $x_0 \to x'$，
它与 $x_0 \to x$ 和 $x' \to x$ 相容。这证明 $\mathcal{F}$ 具有推前。
结论由《范畴》引理 \ref{categories-lemma-fibred-groupoids} 的对偶形式给出。

\medskip\noindent
(b) 的证明。若 $A = k$，则 $\Spec(k) = \Spec(A)$。由于 $\mathcal{X}$
在 $(\Sch/S)_{fppf}$ 上以群胚为纤维纤维化，对 $\mathcal{F}(k)$
中的任意对象 $x_0 \to x$，态射 $x_0 \to x$ 都是同构。因此，
$\mathcal{F}(k)$ 的每个对象都同构于 $x_0 \to x_0$。显然，
$\mathcal{F}$ 中 $x_0 \to x_0$ 的唯一自态射是恒等态射。
\end{proof}

\noindent
设 $S$ 是局部 Noether 基概形，$F : \mathcal{X} \to \mathcal{Y}$ 是
$(\Sch/S)_{fppf}$ 上群胚纤维化范畴之间的 $1$-态射。
% P11-E0207: frozen source says "Let k is a field" instead of "Let k be a field"; translated naturally.
设 $k$ 是 $S$ 上有限型的域。设 $x_0$ 是 $\mathcal{X}$ 中位于
$\Spec(k)$ 上的对象。令 $y_0 = F(x_0)$；这是 $\mathcal{Y}$ 中位于
$\Spec(k)$ 上的对象。则 $F$ 诱导函子
\begin{equation}
\label{artin-equation-functoriality}
F :
\mathcal{F}_{\mathcal{X}, k, x_0}
\longrightarrow
\mathcal{F}_{\mathcal{Y}, k, y_0}
\end{equation}
，其源和靶都是 $\mathcal{C}_\Lambda$ 上的余纤维化范畴。确切地说，
它把 $\mathcal{F}_{\mathcal{X}, k, x_0}(A)$ 的对象 $x_0 \to x$
映为 $\mathcal{F}_{\mathcal{Y}, k, y_0}(A)$ 的对象 $F(x_0) \to F(x)$。

\begin{lemma}
\label{artin-lemma-formally-smooth-on-deformation-categories}
设 $S$ 是局部 Noether 概形，$F : \mathcal{X} \to \mathcal{Y}$ 是
$(\Sch/S)_{fppf}$ 上群胚纤维化范畴的一个 $1$-态射。假设下列条件之一成立：
\begin{enumerate}
\item $F$ 在对象上形式光滑（《可表示性判据》第
\ref{criteria-section-formally-smooth} 节）；
\item $F$ 可由代数空间表示且形式光滑；或
\item $F$ 可由代数空间表示且光滑。
\end{enumerate}
则对每个 $S$ 上有限型的域 $k$ 以及 $\mathcal{X}$ 在 $k$ 上的对象
$x_0$，函子 (\ref{artin-equation-functoriality}) 在《形式形变理论》定义
\ref{formal-defos-definition-smooth-morphism} 的意义下光滑。
\end{lemma}

\begin{proof}
情形 (1) 只需展开定义。由《可表示性判据》引理
\ref{criteria-lemma-representable-by-spaces-formally-smooth}，假设 (2)
蕴含 (1)。由《空间的态射进阶》引理
\ref{spaces-more-morphisms-lemma-smooth-formally-smooth} 以及《代数叠》引理
\ref{algebraic-lemma-representable-transformations-property-implication}
的一般原理，假设 (3) 蕴含 (2)。
\end{proof}

\begin{lemma}
\label{artin-lemma-fibre-product-deformation-categories}
设 $S$ 是局部 Noether 概形。设
$$
\xymatrix{
\mathcal{W} \ar[d] \ar[r] & \mathcal{Z} \ar[d] \\
\mathcal{X} \ar[r] & \mathcal{Y}
}
$$
是 $(\Sch/S)_{fppf}$ 上群胚纤维化范畴的 $2$-纤维积。设 $k$ 是
$S$ 上有限型的域，而 $w_0$ 是 $\mathcal{W}$ 在 $k$ 上的对象。
设 $x_0, z_0, y_0$ 是 $w_0$ 在图中各态射下的像。则
$$
\xymatrix{
\mathcal{F}_{\mathcal{W}, k, w_0} \ar[d] \ar[r] &
\mathcal{F}_{\mathcal{Z}, k, z_0} \ar[d] \\
\mathcal{F}_{\mathcal{X}, k, x_0} \ar[r] & \mathcal{F}_{\mathcal{Y}, k, y_0}
}
$$
是预形变范畴的纤维积。
\end{lemma}

\begin{proof}
只需展开定义。细节从略。
\end{proof}






\section{推出与叠}
\label{artin-section-pushouts}

\noindent
本节证明代数叠对于某些推出具有良好行为。本节结果对任意基概形都成立。

\medskip\noindent
% P11-E0208: the frozen source contains an unresolved editorial-reference placeholder.
当 $Y$、$X'$、$X$、$Y'$ 是代数空间时，下述引理也成立；
参见（此处补入将来的参考文献）。

\begin{lemma}
\label{artin-lemma-pushout}
\begin{slogan}
代数叠满足（强）Rim--Schlessinger 条件
\end{slogan}
设 $S$ 是概形。设
$$
\xymatrix{
X \ar[r] \ar[d] & X' \ar[d] \\
Y \ar[r] & Y'
}
$$
是 $S$ 上概形范畴中的推出，其中 $X \to X'$ 是加厚，且 $X \to Y$
是仿射态射；参见《更多态射》引理
\ref{more-morphisms-lemma-pushout-along-thickening}。
设 $\mathcal{Z}$ 是 $S$ 上的代数叠。则纤维范畴的函子
$$
\mathcal{Z}_{Y'}
\longrightarrow
\mathcal{Z}_Y \times_{\mathcal{Z}_X} \mathcal{Z}_{X'}
$$
是范畴等价。
\end{lemma}

\begin{proof}
% P11-E0209: the frozen full-faithfulness argument names one object and only its self-Isom sheaf.
设 $y'$ 是左端的一个对象。定义在 $Y'$ 上概形范畴上的层
$\mathit{Isom}(y', y')$ 可由 $Y'$ 上的代数空间 $I$ 表示；参见
《代数叠》引理 \ref{algebraic-lemma-representable-diagonal}。
由于 $Y'$ 不仅在概形范畴中，而且在代数空间范畴中也是该推出，
我们断定引理中的函子是全忠实的；参见《空间的推出》引理
\ref{spaces-pushouts-lemma-pushout-along-thickening-schemes}。

\medskip\noindent
设 $(y, x', f)$ 是右端的一个对象，其中
$f : y|_X \to x'|_X$ 是同构。为完成证明，我们必须构造
$\mathcal{Z}_{Y'}$ 的一个对象 $y'$，使其在 $Y$ 和 $X'$ 上的限制
分别与 $y$ 和 $x'$ 相符，并且这种相符方式与 $f$ 相容。事实上，
只需在 $Y'$ 上 fppf 局部地构造 $y'$；参见《叠》引理
\ref{stacks-lemma-characterize-essentially-surjective-when-ff}。
选取一个可表代数叠 $\mathcal{W}$ 和一个满光滑态射
$\mathcal{W} \to \mathcal{Z}$。于是
$$
(\Sch/Y)_{fppf} \times_{y, \mathcal{Z}} \mathcal{W}
\quad\text{且}\quad
(\Sch/X')_{fppf} \times_{x', \mathcal{Z}} \mathcal{W}
$$
是分别由代数空间 $V$ 和 $U'$ 表示的代数叠，且相应表示空间
分别在 $Y$ 和 $X'$ 上光滑。同构 $f$ 诱导 $X$ 上的同构
$\varphi : V \times_Y X \to U' \times_{X'} X$。由《空间的推出》引理
\ref{spaces-pushouts-lemma-pushout-along-thickening} 和
\ref{spaces-pushouts-lemma-equivalence-categories-spaces-pushout-flat}
可知，推出 $V' = V \amalg_{V \times_Y X} U'$ 是 $Y'$ 上光滑的
代数空间；它到 $Y$ 和 $X'$ 的基变换分别恢复 $V$ 和 $U'$，
并且恢复的方式与 $\varphi$ 相容。

\medskip\noindent
令 $W$ 为表示 $\mathcal{W}$ 的代数空间。投影 $V \to W$ 和
$U' \to W$ 在 $V \times_Y X \cong U' \times_{X'} X$ 上作为态射相同，
故推出的泛性质确定代数空间的态射 $V' \to W$。选取概形 $Y_1'$
和满 \'etale 态射 $Y_1' \to V'$。令
$Y_1 = Y \times_{Y'} Y_1'$、$X_1' = X' \times_{Y'} Y_1'$、
$X_1 = X \times_{Y'} Y_1'$。复合
$$
(\Sch/Y_1') \to (\Sch/V') \to (\Sch/W) = \mathcal{W} \to \mathcal{Z}
$$
由 $2$-Yoneda 引理对应于 $\mathcal{Z}$ 在 $Y_1'$ 上的一个对象 $y_1'$；
它在 $Y_1$ 和 $X_1'$ 上的限制分别与 $y|_{Y_1}$ 和 $x'|_{X_1'}$
相符，且相符方式与 $f|_{X_1}$ 相容。因此，我们已经在 $Y'$
上光滑局部地构造出所需对象，证明完成。
\end{proof}








\section{Rim--Schlessinger 条件}
\label{artin-section-RS}

\noindent
下述定义的动机来自引理 \ref{artin-lemma-pushout}，以及《形式形变理论》定义
\ref{formal-defos-definition-RS} 和引理
\ref{formal-defos-lemma-RS-2-categorical}。

\begin{definition}
\label{artin-definition-RS}
设 $S$ 是局部 Noether 概形，$\mathcal{Z}$ 是
$(\Sch/S)_{fppf}$ 上的群胚纤维化范畴。若对每个如下推出
$$
\xymatrix{
X \ar[r] \ar[d] & X' \ar[d] \\
Y \ar[r] & Y' = Y \amalg_X X'
}
$$
（它位于 $S$ 上概形的范畴中），其中
\begin{enumerate}
\item $X$、$X'$、$Y$、$Y'$ 都是 Artin 局部环的谱；
\item $X$、$X'$、$Y$、$Y'$ 在 $S$ 上都是有限型的；并且
\item $X \to X'$（因而 $Y \to Y'$）是闭浸入，
\end{enumerate}
纤维范畴的函子
$$
\mathcal{Z}_{Y'}
\longrightarrow
\mathcal{Z}_Y \times_{\mathcal{Z}_X} \mathcal{Z}_{X'}
$$
都是范畴等价，则称 $\mathcal{Z}$ 满足{\it 条件 (RS)}。
\end{definition}

\noindent
若 $A$ 是剩余域为 $k$ 的 Artin 局部环，则任意态射
$\Spec(A) \to S$ 都是仿射且有限型的，当且仅当诱导态射
$\Spec(k) \to S$ 是有限型的；参见《态射》引理
\ref{morphisms-lemma-Artinian-affine} 和
\ref{morphisms-lemma-artinian-finite-type}。

\begin{lemma}
\label{artin-lemma-algebraic-stack-RS}
设 $\mathcal{X}$ 是局部 Noether 基概形 $S$ 上的代数叠。
则 $\mathcal{X}$ 满足 (RS)。
\end{lemma}

\begin{proof}
由定义和引理 \ref{artin-lemma-pushout} 立即得到。
\end{proof}

\begin{lemma}
\label{artin-lemma-fibre-product-RS}
设 $S$ 是概形。设 $p : \mathcal{X} \to \mathcal{Y}$ 和
$q : \mathcal{Z} \to \mathcal{Y}$ 是 $(\Sch/S)_{fppf}$ 上
群胚纤维化范畴的 $1$-态射。若 $\mathcal{X}$、$\mathcal{Y}$、
$\mathcal{Z}$ 都满足 (RS)，则
$\mathcal{X} \times_\mathcal{Y} \mathcal{Z}$ 也满足 (RS)。
\end{lemma}

\begin{proof}
这是形式论证。设
$$
\xymatrix{
X \ar[r] \ar[d] & X' \ar[d] \\
Y \ar[r] & Y' = Y \amalg_X X'
}
$$
是定义 \ref{artin-definition-RS} 中的图。我们必须证明
$$
(\mathcal{X} \times_{\mathcal{Y}} \mathcal{Z})_{Y'}
\longrightarrow
(\mathcal{X} \times_{\mathcal{Y}} \mathcal{Z})_Y
\times_{(\mathcal{X} \times_{\mathcal{Y}} \mathcal{Z})_X}
(\mathcal{X} \times_{\mathcal{Y}} \mathcal{Z})_{X'}
$$
是等价。利用 $2$-纤维积的定义，这化为
\begin{equation}
\label{artin-equation-RS-fibre-product}
\mathcal{X}_{Y'} \times_{\mathcal{Y}_{Y'}} \mathcal{Z}_{Y'}
\longrightarrow
(\mathcal{X}_Y \times_{\mathcal{Y}_Y} \mathcal{Z}_Y)
\times_{(\mathcal{X}_X \times_{\mathcal{Y}_X} \mathcal{Z}_X)}
(\mathcal{X}_{X'} \times_{\mathcal{Y}_{X'}} \mathcal{Z}_{X'}).
\end{equation}
% P11-E0210: the three frozen-source functors below mix the X, Y, and Z fibres instead of stating each category's own RS equivalence.
已知下列各函子
$$
\mathcal{X}_{Y'} \to \mathcal{X}_Y \times_{\mathcal{Y}_Y} \mathcal{Z}_Y,
\quad
\mathcal{Y}_{Y'} \to \mathcal{X}_X \times_{\mathcal{Y}_X} \mathcal{Z}_X,
\quad
\mathcal{Z}_{Y'} \to
\mathcal{X}_{X'} \times_{\mathcal{Y}_{X'}} \mathcal{Z}_{X'}
$$
都是等价。式 (\ref{artin-equation-RS-fibre-product}) 右端的一个对象是系统
$$
((x_Y, z_Y, \phi_Y), (x_{X'}, z_{X'}, \phi_{X'}), (\alpha, \beta)).
$$
% P11-E0211: the frozen source next uses x_{Y'} and z_{Y'} inside the input triples where the displayed right-hand object supplies x_{X'} and z_{X'}.
于是 $(x_Y, x_{Y'}, \alpha)$ 同构于 $\mathcal{X}_{Y'}$ 中某个对象
$x_{Y'}$ 的像，而 $(z_Y, z_{Y'}, \beta)$ 同构于
$\mathcal{Z}_{Y'}$ 中某个对象 $z_{Y'}$ 的像。态射对
$(\phi_Y, \phi_{X'})$ 对应于 $x_{Y'}$ 和 $z_{Y'}$ 在
$\mathcal{Y}_{Y'}$ 中的像之间的一个态射 $\psi$。于是
$(x_{Y'}, z_{Y'}, \psi)$ 是式 (\ref{artin-equation-RS-fibre-product})
左端的一个对象，它映到给定的右端对象。这证明
(\ref{artin-equation-RS-fibre-product}) 是本质满的。其全忠实性的证明从略。
\end{proof}





\section{形变范畴}
\label{artin-section-deformation-categories}

\noindent
我们把上文引入的记号与《形式形变理论》一章中的记号对应起来。

\begin{lemma}
\label{artin-lemma-deformation-category}
设 $S$ 是局部 Noether 概形，$\mathcal{X}$ 是
$(\Sch/S)_{fppf}$ 上满足 (RS) 的群胚纤维化范畴。对任意在 $S$
上有限型的域 $k$，以及 $\mathcal{X}$ 在 $k$ 上的任意对象 $x_0$，
预形变范畴
$p : \mathcal{F}_{\mathcal{X}, k, x_0} \to \mathcal{C}_\Lambda$
（式 (\ref{artin-equation-predeformation-category})）是形变范畴；参见
《形式形变理论》定义
\ref{formal-defos-definition-deformation-category}。
\end{lemma}

\begin{proof}
令 $\mathcal{F} = \mathcal{F}_{\mathcal{X}, k, x_0}$。设
$f_1 : A_1 \to A$ 和 $f_2 : A_2 \to A$ 是
$\mathcal{C}_\Lambda$ 中的环同态，且 $f_2$ 为满射。须证明函子
$$
\mathcal{F}(A_1 \times_A A_2)
\longrightarrow
\mathcal{F}(A_1) \times_{\mathcal{F}(A)} \mathcal{F}(A_2)
$$
是等价；参见《形式形变理论》引理
\ref{formal-defos-lemma-RS-2-categorical}。令
$X = \Spec(A)$、$X' = \Spec(A_2)$、$Y = \Spec(A_1)$ 以及
$Y' = \Spec(A_1 \times_A A_2)$。注意，在概形范畴中
$Y' = Y \amalg_X X'$；参见《更多态射》引理
\ref{more-morphisms-lemma-pushout-along-thickening}。在纤维范畴函子的图
$$
\xymatrix{
\mathcal{X}_{Y'} \ar[r] \ar[d] &
\mathcal{X}_Y \times_{\mathcal{X}_X} \mathcal{X}_{X'} \ar[d] \\
\mathcal{X}_{\Spec(k)} \ar@{=}[r] & \mathcal{X}_{\Spec(k)}
}
$$
中，由定义 \ref{artin-definition-RS}，上方横箭头是等价。由于
$\mathcal{F}(B)$ 是 $\mathcal{X}_{\Spec(B)}$ 中带有在 $k$ 上与
$x_0$ 的一个认同的对象所成的范畴，结论成立。
\end{proof}

\begin{remark}
\label{artin-remark-deformation-category-implies}
设 $S$ 是局部 Noether 概形，$\mathcal{X}$ 是
$(\Sch/S)_{fppf}$ 上的群胚纤维化范畴。设 $k$ 是 $S$ 上有限型的域，
$x_0$ 是 $\mathcal{X}$ 在 $k$ 上的对象。设
$p : \mathcal{F} \to \mathcal{C}_\Lambda$ 如式
(\ref{artin-equation-predeformation-category})。若 $\mathcal{F}$ 是形变范畴，
即 $\mathcal{F}$ 满足 Rim--Schlessinger 条件 (RS)，则由
《形式形变理论》引理 \ref{formal-defos-lemma-RS-implies-S1-S2}，
$\mathcal{F}$ 满足 Schlessinger 条件 (S1) 和 (S2)。令
$\overline{\mathcal{F}}$ 为同构类函子；参见《形式形变理论》评注
\ref{formal-defos-remarks-cofibered-groupoids}
中的 (\ref{formal-defos-item-associated-functor-isomorphism-classes})。
于是 $\overline{\mathcal{F}}$ 也满足 (S1) 和 (S2)；参见
《形式形变理论》引理
\ref{formal-defos-lemma-S1-S2-associated-functor}。特别地，
这适用于引理 \ref{artin-lemma-deformation-category} 的情形。
\end{remark}




\section{域的变换}
\label{artin-section-change-of-field}

\noindent
本节对应于《形式形变理论》一节
\ref{formal-defos-section-change-of-field}。正如该处所指出的，为讨论
域变换下发生的情形，需要写成 $\mathcal{C}_{\Lambda, k}$，而不是
$\mathcal{C}_\Lambda$。在下述引理中，我们采用《形式形变理论》情形
\ref{formal-defos-situation-change-of-fields} 中引入的记号
$\mathcal{F}_{l/k}$。

\begin{lemma}
\label{artin-lemma-change-of-field}
设 $S$ 是局部 Noether 概形，$\mathcal{X}$ 是
$(\Sch/S)_{fppf}$ 上的群胚纤维化范畴。设 $k$ 是 $S$ 上有限型的域，
$l/k$ 是有限扩张。
% P11-E0212: the frozen source says x_0 is an object of undefined \mathcal{F}, while the surrounding construction requires \mathcal{X}; the mathematical token is preserved.
设 $x_0$ 是 $\mathcal{F}$ 位于 $\Spec(k)$ 上的一个对象。
% P11-E0213: the frozen source omits "by" in the naming construction "Denote x_{l,0} the restriction".
记 $x_{l, 0}$ 为 $x_0$ 在 $\Spec(l)$ 上的限制。则存在
在 $\mathcal{C}_{\Lambda, l}$ 上以群胚为纤维的余纤维化范畴的典范函子
$$
(\mathcal{F}_{\mathcal{X}, k , x_0})_{l/k}
\longrightarrow
\mathcal{F}_{\mathcal{X}, l, x_{l, 0}}
$$
。若 $\mathcal{X}$ 满足 (RS)，则此函子是等价。
\end{lemma}

\begin{proof}
考虑式 (\ref{artin-equation-factor}) 中的分解
$$
\Spec(l) \to \Spec(B) \to S
$$
。由定义，
$$
(\mathcal{F}_{\mathcal{X}, k , x_0})_{l/k}(B) =
\mathcal{F}_{\mathcal{X}, k, x_0}(B \times_l k)
$$
；参见《形式形变理论》情形
\ref{formal-defos-situation-change-of-fields}。因此，其中一个对象是
$\mathcal{X}$ 的态射 $x_0 \to x$，它位于态射
$\Spec(k) \to \Spec(B \times_l k)$ 之上。
% P11-E0214: the frozen source omits the article in "Choosing pullback functor".
为 $\mathcal{X}$ 选取拉回函子后，可将 $x_0 \to x$ 对应到态射
$x_{l, 0} \to x_B$，其中 $x_B$ 是 $x$ 在 $\Spec(B)$ 上的限制
（经由来自 $B \times_l k \subset B$ 的态射
$\Spec(B) \to \Spec(B \times_l k)$）。此构造关于 $B$ 函子性地变化，
并与态射相容。

\medskip\noindent
下面假设 $\mathcal{X}$ 满足 (RS)。考虑诸图
$$
\vcenter{
\xymatrix{
l & B \ar[l] \\
k \ar[u] & B \times_l k \ar[l] \ar[u]
}
}
\quad\text{且}\quad
\vcenter{
\xymatrix{
\Spec(l) \ar[d] \ar[r] & \Spec(B) \ar[d] \\
\Spec(k) \ar[r] & \Spec(B \times_l k)
}
}
$$
左图是环的纤维积，右图是概形范畴中的推出；参见《更多态射》引理
\ref{more-morphisms-lemma-pushout-along-thickening}。这些概形在 $S$
上都是有限型的（参见定义 \ref{artin-definition-RS} 后的评注）。因此，
(RS) 给出纤维范畴的等价
$$
\mathcal{X}_{\Spec(B \times_l k)}
\longrightarrow
\mathcal{X}_{\Spec(k)}
\times_{\mathcal{X}_{\Spec(l)}}
\mathcal{X}_{\Spec(B)}
$$
。这说明上面定义的函子给出纤维范畴的等价。因此，由《范畴》引理
\ref{categories-lemma-equivalence-fibred-categories} 的对偶，
该函子是以群胚为纤维的余纤维化范畴之间的等价。
\end{proof}








\section{切空间}
\label{artin-section-tangent-spaces}

\noindent
设 $S$ 是局部 Noether 概形，$\mathcal{X}$ 是
$(\Sch/S)_{fppf}$ 上的群胚纤维化范畴。设 $k$ 是 $S$ 上有限型的域，
$x_0$ 是 $\mathcal{X}$ 在 $k$ 上的对象。在《形式形变理论》一节
\ref{formal-defos-section-tangent-spaces} 中，我们定义了预形变范畴
$\mathcal{F}_{\mathcal{X}, k, x_0}$ 的{\it 切空间}
\begin{equation}
\label{artin-equation-tangent-space}
T\mathcal{F}_{\mathcal{X}, k, x_0} =
\left\{
\begin{matrix}
\text{isomorphism classes of morphisms}\\
x_0 \to x\text{ 相对于 }\Spec(k) \to \Spec(k[\epsilon])
\end{matrix}
\right\}
\end{equation}
。在《形式形变理论》一节
\ref{formal-defos-section-infinitesimal-automorphisms} 中，我们定义了
\begin{equation}
\label{artin-equation-infinitesimal-automorphisms}
\text{Inf}(\mathcal{F}_{\mathcal{X}, k, x_0}) =
\Ker\left(
\text{Aut}_{\Spec(k[\epsilon])}(x'_0) \to \text{Aut}_{\Spec(k)}(x_0)
\right)
\end{equation}
，其中 $x_0'$ 是 $x_0$ 到 $\Spec(k[\epsilon])$ 的拉回。
若 $\mathcal{X}$ 满足 Rim--Schlessinger 条件 (RS)，则由
《形式形变理论》引理
\ref{formal-defos-lemma-tangent-space-vector-space}
（由引理 \ref{artin-lemma-deformation-category} 和评注
\ref{artin-remark-deformation-category-implies}，其假设成立），
$T\mathcal{F}_{\mathcal{X}, k, x_0}$ 带有自然的 $k$-向量空间结构。
此外，《形式形变理论》引理
\ref{formal-defos-lemma-infaut-vector-space} 表明，
$\text{Inf}(\mathcal{F}_{\mathcal{X}, k, x_0})$ 带有自然的
$k$-向量空间结构，使加法与自同构的复合一致。一个自然条件是要求
这些向量空间为有限维。

\medskip\noindent
下述引理说明，若 $\mathcal{X}$ 在 $S$ 上局部有限型，则此条件成立
（参见《叠的态射》一节 \ref{stacks-morphisms-section-finite-type}）。

\begin{lemma}
\label{artin-lemma-finite-dimension}
设 $S$ 是局部 Noether 概形。假设
\begin{enumerate}
\item $\mathcal{X}$ 是代数叠；
\item $U$ 是 $S$ 上局部有限型的概形；并且
\item $(\Sch/U)_{fppf} \to \mathcal{X}$ 是满光滑态射。
\end{enumerate}
% P11-E0215: the frozen source omits punctuation between the section-reference modifier and the main clause.
则对第 \ref{artin-section-predeformation-categories} 节中的任意
$\mathcal{F} = \mathcal{F}_{\mathcal{X}, k, x_0}$，切空间
$T\mathcal{F}$ 和无穷小自同构空间 $\text{Inf}(\mathcal{F})$
在 $k$ 上都是有限维的。
\end{lemma}

\begin{proof}
记 $\mathcal{U} = (\Sch/U)_{fppf}$。由代数叠的定义，$1$-态射
$\mathcal{U} \to \mathcal{X}$ 可由代数空间表示。因此，特别地，
2-纤维积
$$
\mathcal{U}_{x_0} = (\Sch/\Spec(k))_{fppf} \times_\mathcal{X} \mathcal{U}
$$
可由 $\Spec(k)$ 上的代数空间 $U_{x_0}$ 表示。于是
$U_{x_0} \to \Spec(k)$ 光滑且满（特别地，$U_{x_0}$ 非空）。由
《域上的空间》引理
\ref{spaces-over-fields-lemma-smooth-separable-closed-points-dense}
，可找到有限扩张 $l/k$ 和 $k$ 上的点
$\Spec(l) \to U_{x_0}$。由引理 \ref{artin-lemma-change-of-field} 以及
$\mathcal{X}$ 满足 (RS) 这一事实，有
$$
(\mathcal{F}_{\mathcal{X}, k , x_0})_{l/k} =
\mathcal{F}_{\mathcal{X}, l, x_{l, 0}}
$$
。因而
$$
T\mathcal{F} \otimes_k l \cong T\mathcal{F}_{\mathcal{X}, l, x_{l, 0}}
\quad\text{且}\quad
\text{Inf}(\mathcal{F}) \otimes_k l \cong
\text{Inf}(\mathcal{F}_{\mathcal{X}, l, x_{l, 0}})
$$
；这是由《形式形变理论》引理
\ref{formal-defos-lemma-tangent-space-change-of-field} 和
\ref{formal-defos-lemma-inf-aut-change-of-field}
得到的（由引理 \ref{artin-lemma-algebraic-stack-RS}、
\ref{artin-lemma-deformation-category} 和评注
\ref{artin-remark-deformation-category-implies}，这些引理可以应用）。
因此，只需证明 $T\mathcal{F}_{\mathcal{X}, l, x_{l, 0}}$ 和
$\text{Inf}(\mathcal{F}_{\mathcal{X}, l, x_{l, 0}})$ 在 $l$ 上有限维。
注意，$x_{l, 0}$ 来自 $\mathcal{U}$ 在 $l$ 上的一个点 $u_0$。

\medskip\noindent
我们暂时中断论证，说明无穷小自同构的结论可由切空间的结论推出。
令 $\mathcal{R} = \mathcal{U} \times_\mathcal{X} \mathcal{U}$，
并令 $r_0$ 为 $\mathcal{R}$ 的 $l$-值点
$(u_0, u_0, \text{id}_{x_0})$。结合引理
\ref{artin-lemma-fibre-product-deformation-categories} 和
《形式形变理论》引理
\ref{formal-defos-lemma-deformation-functor-diagonal}
可知
$$
\text{Inf}(\mathcal{F}_{\mathcal{X}, l, x_{l, 0}})
\subset
T\mathcal{F}_{\mathcal{R}, l, r_0}
$$
。注意，$\mathcal{R}$ 是代数叠；参见《代数叠》引理
\ref{algebraic-lemma-2-fibre-product-general}。此外，$\mathcal{R}$
可由代数空间 $R$ 表示，且该空间经任一投影在 $U$ 上光滑；参见
《代数叠》引理 \ref{algebraic-lemma-stack-presentation}。
% P11-E0216: the frozen source has "choose an scheme" and joins the choice directly to "we see" without a grammatical connector.
因此，选取概形 $U'$ 和满 \'etale 态射 $U' \to R$，可知 $U'$
在 $U$ 上光滑，因而在 $S$ 上局部有限型。由于
$(\Sch/U')_{fppf} \to \mathcal{R}$ 满且光滑，问题已经归结为切空间情形。

\medskip\noindent
函子 (\ref{artin-equation-functoriality})
$$
\mathcal{F}_{\mathcal{U}, l, u_0}
\longrightarrow
\mathcal{F}_{\mathcal{X}, l, x_{l, 0}}
$$
由引理 \ref{artin-lemma-formally-smooth-on-deformation-categories} 是光滑的。
所诱导的切空间映射
$$
T\mathcal{F}_{\mathcal{U}, l, u_0}
\longrightarrow
T\mathcal{F}_{\mathcal{X}, l, x_{l, 0}}
$$
是 $l$-线性的（由《形式形变理论》引理
\ref{formal-defos-lemma-k-linear-differential}），也是满的（因为由
《形式形变理论》引理
\ref{formal-defos-lemma-smooth-morphism-essentially-surjective}，
预形变范畴的光滑映射诱导切空间的满映射）。因此，只需证明与可表代数叠
$\mathcal{U}$ 相联系的形变空间在点 $u_0$ 处的切空间为有限维。
% P11-E0217: the frozen source says "deformation space" although the surrounding construction is the predeformation/deformation category associated to \mathcal{U}.
取仿射开集 $\Spec(R) \subset U$，使 $u_0 : \Spec(l) \to U$ 分解通过
$\Spec(R)$，且 $\Spec(R) \to S$ 分解通过
$\Spec(\Lambda) \subset S$。令 $\mathfrak m_R \subset R$ 为与 $u_0$
对应的 $\Lambda$-代数同态 $\varphi_0 : R \to l$ 的核。注意，$R$
在 Noether 环 $\Lambda$ 上有限型，因而是 Noether 环。因此，
$\mathfrak m_R = (f_1, \ldots, f_n)$ 是有限生成理想。我们有
$$
T\mathcal{F}_{\mathcal{U}, l, u_0}
=
\{\varphi : R \to l[\epsilon] \mid
\varphi \text{ is a } \Lambda\text{-algebra map and }
\varphi \bmod \epsilon = \varphi_0\}
$$
% P11-E0218: the frozen source joins the independent clauses around "hence" without punctuation.
右端的一个元素由它在 $f_1, \ldots, f_n$ 上的值确定；因此其维数至多为
$n$，结论成立。这里略去若干细节。
\end{proof}

\begin{lemma}
\label{artin-lemma-fibre-product-tangent-spaces}
设 $S$ 是局部 Noether 概形。设 $p : \mathcal{X} \to \mathcal{Y}$ 和
$q : \mathcal{Z} \to \mathcal{Y}$ 是 $(\Sch/S)_{fppf}$ 上
群胚纤维化范畴的 $1$-态射。假设 $\mathcal{X}, \mathcal{Y}, \mathcal{Z}$
满足 (RS)。设 $k$ 是 $S$ 上有限型的域，$w_0$ 是
$\mathcal{W} = \mathcal{X} \times_\mathcal{Y} \mathcal{Z}$ 在 $k$
上的一个对象。
% P11-E0219: the frozen source omits "by" in "Denote x_0, y_0, z_0 the objects".
将由 $w_0$ 得到的 $\mathcal{X}$、$\mathcal{Y}$、$\mathcal{Z}$ 的对象
分别记为 $x_0, y_0, z_0$。则存在 $k$-向量空间的 $6$-项正合列
$$
\xymatrix{
0 \ar[r] &
\text{Inf}(\mathcal{F}_{\mathcal{W}, k, w_0}) \ar[r] &
\text{Inf}(\mathcal{F}_{\mathcal{X}, k, x_0}) \oplus
\text{Inf}(\mathcal{F}_{\mathcal{Z}, k, z_0}) \ar[r] &
\text{Inf}(\mathcal{F}_{\mathcal{Y}, k, y_0}) \ar[lld] \\
 &
T\mathcal{F}_{\mathcal{W}, k, w_0} \ar[r] &
T\mathcal{F}_{\mathcal{X}, k, x_0} \oplus
T\mathcal{F}_{\mathcal{Z}, k, z_0} \ar[r] &
T\mathcal{F}_{\mathcal{Y}, k, y_0}
}
$$
。
\end{lemma}

\begin{proof}
由引理 \ref{artin-lemma-fibre-product-RS}，$\mathcal{W}$ 满足 (RS)，
因而本引理的陈述有意义。为证明本引理，应用引理
\ref{artin-lemma-fibre-product-deformation-categories}、
\ref{artin-lemma-deformation-category} 以及《形式形变理论》引理
\ref{formal-defos-lemma-deformation-categories-fiber-product-morphisms}。
\end{proof}






\section{形式对象}
\label{artin-section-formal-objects}

\noindent
本节把《形式形变理论》一章中已经定义的若干概念移到当前情形。
下文说“$R$ 是 $S$-代数”，是指 $R$ 是一个环，并配备概形态射
$\Spec(R) \to S$。

\begin{definition}
\label{artin-definition-formal-objects}
设 $S$ 是局部 Noether 概形，并设
$p : \mathcal{X} \to (\Sch/S)_{fppf}$ 是群胚纤维化范畴。
\begin{enumerate}
\item $\mathcal{X}$ 的一个{\it 形式对象} $\xi = (R, \xi_n, f_n)$
由如下数据组成：Noether 完备局部 $S$-代数 $R$；$\mathcal{X}$ 中位于
$\Spec(R/\mathfrak m_R^n)$ 上的对象 $\xi_n$；以及 $\mathcal{X}$ 中位于
% P11-E0220: the frozen source switches from \mathfrak m_R to undefined \mathfrak m in the transition morphisms and residue-field condition.
$\Spec(R/\mathfrak m^n) \to \Spec(R/\mathfrak m^{n + 1})$ 上的态射
$f_n : \xi_n \to \xi_{n + 1}$。这里要求 $R/\mathfrak m$ 是 $S$
上有限型的域。
\item 一个{\it 形式对象的态射}
$a : \xi = (R, \xi_n, f_n) \to \eta = (T, \eta_n, g_n)$
由态射 $a_n : \xi_n \to \eta_n$ 给出，并且对每个 $n$，图
$$
\xymatrix{
\xi_n \ar[r]_{f_n} \ar[d]_{a_n} & \xi_{n + 1} \ar[d]^{a_{n + 1}} \\
\eta_n \ar[r]^{g_n} & \eta_{n + 1}
}
$$
交换。应用函子 $p$，得到相容的态射族
$\Spec(R/\mathfrak m_R^n) \to \Spec(T/\mathfrak m_T^n)$，从而得到
$S$ 上的态射 $a_0 : \Spec(R) \to \Spec(T)$。我们称 $a$
{\it 位于} $a_0$ 之上。
\end{enumerate}
\end{definition}

\noindent
这样就得到 $\mathcal{X}$ 的形式对象所成的范畴。

\begin{remark}
\label{artin-remark-formal-objects-match}
设 $S$ 是局部 Noether 概形，并设
$p : \mathcal{X} \to (\Sch/S)_{fppf}$ 是群胚纤维化范畴。
设 $\xi = (R, \xi_n, f_n)$ 是形式对象。令 $k = R/\mathfrak m$ 且
$x_0 = \xi_1$。形式对象 $\xi$ 定义了预形变范畴
$\mathcal{F}_{\mathcal{X}, k, x_0}$ 的一个形式对象 $\xi$。
这立即来自上面的定义 \ref{artin-definition-formal-objects}、
《形式形变理论》定义 \ref{formal-defos-definition-formal-objects}，
以及第 \ref{artin-section-predeformation-categories} 节中对预形变范畴
$\mathcal{F}_{\mathcal{X}, k, x_0}$ 的构造。
\end{remark}

\noindent
若 $F : \mathcal{X} \to \mathcal{Y}$ 是 $(\Sch/S)_{fppf}$ 上
群胚纤维化范畴的 $1$-态射，则 $F$ 也诱导形式对象范畴之间的函子。

\begin{lemma}
\label{artin-lemma-smooth-lift-formal}
设 $S$ 是局部 Noether 概形。设 $F : \mathcal{X} \to \mathcal{Y}$ 是
$(\Sch/S)_{fppf}$ 上群胚纤维化范畴的 $1$-态射。设
$\eta = (R, \eta_n, g_n)$ 是 $\mathcal{Y}$ 的形式对象，$\xi_1$ 是
$\mathcal{X}$ 的对象，且 $F(\xi_1) \cong \eta_1$。若 $F$ 在对象上
形式光滑（参见《可表示性判据》一节
\ref{criteria-section-formally-smooth}），则存在 $\mathcal{X}$ 的形式对象
$\xi = (R, \xi_n, f_n)$，使 $F(\xi) \cong \eta$。
\end{lemma}

\begin{proof}
注意，每个态射
$\Spec(R/\mathfrak m^n) \to \Spec(R/\mathfrak m^{n + 1})$ 都是 $S$
上仿射概形的一阶加厚。因此，对 $F$ 的假设意味着可逐次把 $\xi_1$
提升为 $\mathcal{X}$ 的对象 $\xi_2, \xi_3, \ldots$，并配备相容同构
% P11-E0221: the frozen source switches xi and eta in both displayed compatibility isomorphisms.
$\eta_n|_{\Spec(R/\mathfrak m^{n - 1})} \cong \eta_{n - 1}$
和 $F(\eta_n) \cong \xi_n$。
\end{proof}

\noindent
设 $S$ 是局部 Noether 概形，并设
$p : \mathcal{X} \to (\Sch/S)_{fppf}$ 是群胚纤维化范畴。假设 $x$
是 $\mathcal{X}$ 在 $R$ 上的对象，其中 $R$ 是 Noether 完备局部
$S$-代数，其剩余域在 $S$ 上有限型。于是，可以考虑限制系统
$\xi_n = x|_{\Spec(R/\mathfrak m^n)}$，并赋予它由限制的传递性所得的
自然态射 $\xi_1 \to \xi_2 \to \ldots$。因此，
$\xi = (R, \xi_n, \xi_n \to \xi_{n + 1})$ 是 $\mathcal{X}$ 的形式对象。
此构造关于对象 $x$ 函子性地变化。这样得到函子
\begin{equation}
\label{artin-equation-approximation}
\left\{
\begin{matrix}
\text{objects }x\text{ 属于 }\mathcal{X} \text{ 使得 }p(x) = \Spec(R) \\
\text{其中 }R\text{ is Noetherian complete local}\\
\text{其中 }R/\mathfrak m\text{ of finite type over }S
\end{matrix}
\right\}
\longrightarrow
\left\{
\begin{matrix}
\text{formal objects of }\mathcal{X}
\end{matrix}
\right\}
\end{equation}
确切地说，左端是 $\mathcal{X}$ 中由所述对象组成的全子范畴，
右端是定义 \ref{artin-definition-formal-objects} 中 $\mathcal{X}$ 的形式对象范畴。

\begin{definition}
\label{artin-definition-effective}
设 $S$ 是局部 Noether 概形，$\mathcal{X}$ 是
$(\Sch/S)_{fppf}$ 上的群胚纤维化范畴。若 $\mathcal{X}$ 的形式对象
$\xi = (R, \xi_n, f_n)$ 位于函子 (\ref{artin-equation-approximation})
的本质像中，则称它是{\it 有效的}。
\end{definition}

\noindent
若这个群胚纤维化范畴是代数叠，则由下述引理，每个形式对象都是有效的。

\begin{lemma}
\label{artin-lemma-effective}
设 $S$ 是局部 Noether 概形，$\mathcal{X}$ 是 $S$ 上的代数叠。
函子 (\ref{artin-equation-approximation}) 是等价。
\end{lemma}

\begin{proof}
情形 I：$\mathcal{X}$ 可由概形表示。设
$\mathcal{X} = (\Sch/X)_{fppf}$，其中 $X$ 是 $S$ 上的某个概形。
展开定义后，需要证明如下事实：给定 Noether 完备局部 $S$-代数 $R$，
且 $R/\mathfrak m$ 在 $S$ 上有限型，则
$$
\Mor_S(\Spec(R), X) \longrightarrow \lim \Mor_S(\Spec(R/\mathfrak m^n), X)
$$
为双射。这来自《形式空间》引理
\ref{formal-spaces-lemma-map-into-scheme}。

\medskip\noindent
情形 II：$\mathcal{X}$ 可由代数空间表示。设 $\mathcal{X}$ 由 $X$ 表示。
仍须证明
$$
\Mor_S(\Spec(R), X) \longrightarrow \lim \Mor_S(\Spec(R/\mathfrak m^n), X)
$$
对上述 $R$ 为双射。这就是《形式空间》引理
\ref{formal-spaces-lemma-map-into-algebraic-space}。

\medskip\noindent
情形 III：一般的代数叠情形。一般地，式
(\ref{artin-equation-approximation}) 的左右两端，都是下述范畴上的
群胚纤维化范畴：对象为 $S$ 上的仿射概形，它们是剩余域在 $S$
上有限型的 Noether 完备局部环之谱。在下面的证明中还会看到，
对于这个范畴上的某个拓扑，它们构成叠。

\medskip\noindent
% P11-E0222: the frozen source says "prove fully faithfulness" instead of "prove full faithfulness".
先证明全忠实性。设 $R$ 是 Noether 完备局部 $S$-代数，且
$k = R/\mathfrak m$ 在 $S$ 上有限型。设 $x, x'$ 是 $\mathcal{X}$
在 $R$ 上的对象。由于 $\mathcal{X}$ 是代数叠，
$\mathit{Isom}(x, x')$ 可由 $\Spec(R)$ 上的代数空间 $I$ 表示；
参见《代数叠》引理 \ref{algebraic-lemma-representable-diagonal}。
把情形 II 应用于 $\Spec(R)$ 上的 $I$，立即得到
(\ref{artin-equation-approximation}) 在 $\Spec(R)$ 上的纤维范畴之间全忠实。
因此，由《范畴》引理 \ref{categories-lemma-equivalence-fibred-categories}，
该函子全忠实。

\medskip\noindent
本质满性。设 $\xi = (R, \xi_n, f_n)$ 是 $\mathcal{X}$ 的形式对象。
选取 $S$ 上的概形 $U$ 和满光滑态射
$f : (\Sch/U)_{fppf} \to \mathcal{X}$。对每个 $n$，考虑纤维积
$$
(\Sch/\Spec(R/\mathfrak m^n))_{fppf}
\times_{\xi_n, \mathcal{X}, f}
(\Sch/U)_{fppf}
$$
。由假设，它可由代数空间 $V_n$ 表示，且该空间在
$\Spec(R/\mathfrak m^n)$ 上满且光滑。态射
$f_n : \xi_n \to \xi_{n + 1}$ 诱导代数空间的 Cartesian 方块
$$
\xymatrix{
V_{n + 1} \ar[d] & V_n \ar[d] \ar[l] \\
\Spec(R/\mathfrak m^{n + 1}) & \Spec(R/\mathfrak m^n) \ar[l]
}
$$
。由《域上的空间》引理
\ref{spaces-over-fields-lemma-smooth-separable-closed-points-dense}
，可找到有限可分扩张 $k'/k$ 和 $k$ 上的点
$v'_1 : \Spec(k') \to V_1$。令 $R \subset R'$ 为剩余域扩张是
$k'/k$ 的有限 \'etale 扩张；由《代数》引理
\ref{algebra-lemma-henselian-cat-finite-etale} 和
\ref{algebra-lemma-complete-henselian}，它存在且唯一。对
$n = 2, 3, 4, \ldots$，把光滑性的无穷小提升判据（参见
《代数空间态射进阶》引理
\ref{spaces-more-morphisms-lemma-smooth-formally-smooth})
应用于 $V_n \to \Spec(R/\mathfrak m^n)$，可逐次找到
$\Spec(R/\mathfrak m^n)$ 上的态射
$v'_n : \Spec(R'/(\mathfrak m')^n) \to V_n$，使下图交换：
$$
\xymatrix{
\Spec(R'/(\mathfrak m')^{n + 1}) \ar[d]_{v'_{n + 1}} &
\Spec(R'/(\mathfrak m')^n) \ar[d]^{v'_n} \ar[l] \\
V_{n + 1} & V_n \ar[l]
}
$$
与投影态射 $V_n \to U$ 复合，得到相容态射系
$u'_n : \Spec(R'/(\mathfrak m')^n) \to U$。由情形 I，族 $(u'_n)$
来自唯一态射 $u' : \Spec(R') \to U$。
% P11-E0223: the frozen source omits "by" in "Denote x' the object".
把 $1$-态射 $f$ 应用于 $u'$ 所得的 $\mathcal{X}$ 在
$\Spec(R')$ 上的对象记为 $x'$。由构造，存在形式对象的态射
$$
(\ref{artin-equation-approximation})(x') =
(R', x'|_{\Spec(R'/(\mathfrak m')^n)}, \ldots)
\longrightarrow
(R, \xi_n, f_n)
$$
，它位于 $\Spec(R') \to \Spec(R)$ 之上。
% P11-E0224: the frozen source says the ring R' tensor_R R' is a finite product of spectra, mixing ring and scheme types.
注意，$R' \otimes_R R'$ 是若干 Noether 完备局部环的谱的有限乘积，
而当前讨论适用于这些谱。
% P11-E0225: the frozen source omits "by" in the naming construction for p_0 and p_1.
把两个投影记为
$p_0, p_1 : \Spec(R' \otimes_R R') \to \Spec(R')$。
% P11-E0226: the frozen source says "By the fully faithfulness" instead of "By the full faithfulness".
由上面证明的全忠实性，存在典范同构
$\varphi : p_0^*x' \to p_1^*x'$，因为在
$\Spec((R' \otimes_R R')/\mathfrak m^n(R' \otimes_R R'))$ 上已有这样的同构。
同构 $\varphi$ 满足上链条件的证明从略（参见《叠》定义
\ref{stacks-definition-descent-data}）。由于
$\{\Spec(R') \to \Spec(R)\}$ 是 fppf 覆盖，断定 $x'$ 下降为
$\mathcal{X}$ 在 $\Spec(R)$ 上的对象 $x$。
% P11-E0227: the frozen source refers to x_n, whereas the formal object's components are xi_n.
$x_n$ 是 $x$ 在 $\Spec(R/\mathfrak m^n)$ 上的限制这一点的证明从略。
\end{proof}

\begin{lemma}
\label{artin-lemma-fibre-product-effective}
设 $S$ 是概形。设 $p : \mathcal{X} \to \mathcal{Y}$ 和
$q : \mathcal{Z} \to \mathcal{Y}$ 是 $(\Sch/S)_{fppf}$ 上
群胚纤维化范畴的 $1$-态射。若函子
(\ref{artin-equation-approximation}) 对 $\mathcal{X}$、$\mathcal{Y}$ 和
$\mathcal{Z}$ 都是等价，则它对
$\mathcal{X} \times_\mathcal{Y} \mathcal{Z}$ 也是等价。
\end{lemma}

\begin{proof}
对 $\mathcal{X} \times_\mathcal{Y} \mathcal{Z}$ 而言，式
(\ref{artin-equation-approximation}) 的左端和右端，分别就是关于
$\mathcal{X}$、$\mathcal{Z}$ 在 $\mathcal{Y}$ 上的式
(\ref{artin-equation-approximation}) 左端和右端的 $2$-纤维积。由于取
$2$-纤维积与范畴等价相容，结论随即得到；参见《范畴》引理
\ref{categories-lemma-equivalence-2-fibre-product}。
\end{proof}






\section{逼近}
\label{artin-section-approximation}

\noindent
Michael Artin 的一个根本洞见是：可以逼近保极限叠的对象。具体地，
给定这个叠在某个 Noether 完备局部环上的对象 $x$，可以找到代数环上的
对象 $x_A$，使其“接近”$x$。这里，代数环是指有限型 $S$-代数，
而接近是指在进拓扑意义下接近。本节以简单而一般的形式陈述这一结果。

\medskip\noindent
为表述这个结果，需要汇集 Stacks Project 不同部分中的若干定义。首先，
在《可表示性判据》一节 \ref{criteria-section-limit-preserving} 中，
我们对概形范畴上群胚纤维化范畴的 $1$-态射引入了
{\it 在对象上保极限}的概念。在《更多代数》定义
\ref{more-algebra-definition-G-ring} 中，我们定义了 {\it G-环}。
设 $S$ 是局部 Noether 概形，并设 $A$ 是 $S$-代数。若
$\Spec(A) \to S$ 为有限型态射，则称 $A$ {\it 在 $S$ 上有限型}，
或称它是{\it 有限型 $S$-代数}。在这种情形下，$A$ 是 Noether 环。
最后，给定环 $A$ 和理想 $I$，记
$\text{Gr}_I(A) = \bigoplus I^n/I^{n + 1}$。

\begin{lemma}
\label{artin-lemma-approximate}
设 $S$ 是局部 Noether 概形，并设
$p : \mathcal{X} \to (\Sch/S)_{fppf}$ 是群胚纤维化范畴。设 $x$
是 $\mathcal{X}$ 位于 $\Spec(R)$ 上的对象，其中 $R$ 是 Noether
完备局部环，其剩余域 $k$ 在 $S$ 上有限型。令 $s \in S$ 为
$\Spec(k) \to S$ 的像。假设 (a) $\mathcal{O}_{S, s}$ 是 G-环，
并且 (b) $p$ 在对象上保极限。则对每个整数 $N \geq 1$，存在
\begin{enumerate}
\item 有限型 $S$-代数 $A$；
\item 极大理想 $\mathfrak m_A \subset A$；
\item $\mathcal{X}$ 在 $\Spec(A)$ 上的对象 $x_A$；
\item $S$-同构 $R/\mathfrak m_R^N \cong A/\mathfrak m_A^N$；
\item 同构
$x|_{\Spec(R/\mathfrak m_R^N)} \cong x_A|_{\Spec(A/\mathfrak m_A^N)}$
，它与 (4) 相容；以及
\item 分次 $k$-代数的同构
$\text{Gr}_{\mathfrak m_R}(R) \cong \text{Gr}_{\mathfrak m_A}(A)$
。
\end{enumerate}
\end{lemma}

\begin{proof}
选取仿射开集 $\Spec(\Lambda) \subset S$，使 $k$ 是有限
$\Lambda$-代数；参见《态射》引理
\ref{morphisms-lemma-point-finite-type}。我们可以并且确实用
$\Spec(\Lambda)$ 代替 $S$。

\medskip\noindent
可将 $R$ 写成有向余极限 $R = \colim C_j$，其中每个 $C_j$ 都是
有限型 $\Lambda$-代数（参见《代数》引理
\ref{algebra-lemma-ring-colimit-fp}）。由假设 (b)，对象 $x$ 同构于
某个 $C_j$ 上对象的限制。因此，可以选取有限型 $\Lambda$-代数 $C$、
$\Lambda$-代数同态 $C \to R$，以及 $\mathcal{X}$ 在
$\Spec(C)$ 上的对象 $x_C$，使得 $x = x_C|_{\Spec(R)}$。对 $C$ 的选取
只是记账手段，可以避免。为后文使用，写成
$C = \Lambda[y_1, \ldots, y_u]/(f_1, \ldots, f_v)$，并把 $y_i$
在映射 $C \to R$ 下的像记为 $\overline{a}_i \in R$。令
$\mathfrak m_C = C \cap \mathfrak m_R$。

\medskip\noindent
选取 $\Lambda$-代数满射 $\Lambda[x_1, \ldots, x_s] \to k$。
% P11-E0228: the frozen source omits "by" in "denote m' the kernel".
把它的核记为 $\mathfrak m'$。由多项式环的泛性质，可将其提升为
$\Lambda$-代数同态 $\Lambda[x_1, \ldots, x_s] \to R$。再添加若干变量
（即略微增大 $s$），使它们映到 $\mathfrak m_R$ 的一组生成元。
完成这一步后，$\Lambda[x_1, \ldots, x_s] \to R/\mathfrak m_R^2$
是满射。因此，
\begin{equation}
\label{artin-equation-surjection}
P = \Lambda[x_1, \ldots, x_s]_{\mathfrak m'}^\wedge \longrightarrow R
\end{equation}
是 Noether 完备局部环的满同态；例如参见《形式形变理论》引理
\ref{formal-defos-lemma-surjective-cotangent-space}。

\medskip\noindent
选取上面得到的 $\overline{a}_i$ 的提升 $a_i \in P$。
选取式 (\ref{artin-equation-surjection}) 之核的一组生成元
$b_1, \ldots, b_r \in P$。选取 $c_{ji} \in P$，使
$$
f_j(a_1, \ldots, a_u) = \sum c_{ji} b_i
$$
在 $P$ 中成立；根据此前的选取，这是可能的。选取如下核的一组生成元
$$
k_1, \ldots, k_t \in
\Ker(P^{\oplus r} \xrightarrow{(b_1, \ldots, b_r)} P)
$$
，并写 $k_i = (k_{i1}, \ldots, k_{ir})$ 和 $K = (k_{ij})$，使
$$
P^{\oplus t} \xrightarrow{K}
P^{\oplus r} \xrightarrow{(b_1, \ldots, b_r)}
P \to R \to 0
$$
成为 $P$-模的正合列。特别地，$\sum k_{ij} b_j = 0$。必要时增大
$N$ 后，可假设 $N - 1$ 对此正合列的前两个映射适用 Artin--Rees
引理（术语参见《更多代数》一节
\ref{more-algebra-section-artin-rees}）。

\medskip\noindent
由假设，$\mathcal{O}_{S, s} = \Lambda_{\Lambda \cap \mathfrak m'}$ 是
G-环。因此，由《更多代数》命题
\ref{more-algebra-proposition-finite-type-over-G-ring}
，环 $\Lambda[x_1, \ldots, x_s]_{\mathfrak m'}$ 是 $G$-环。
因而，由《环同态的光滑化》定理
\ref{smoothing-theorem-approximation-property-variant}
，存在 \'etale 环同态
$$
\Lambda[x_1, \ldots, x_s]_{\mathfrak m'} \to B,
$$
、$B$ 中位于 $\mathfrak m'$ 之上的极大理想 $\mathfrak m_B$，以及
% P11-E0229: the frozen source places the approximating elements in undefined B' although only B has been introduced.
元素 $a'_i, b'_i, c'_{ij}, k'_{ij} \in B'$，使得
\begin{enumerate}
\item $\kappa(\mathfrak m') = \kappa(\mathfrak m_B)$；这蕴含
$\Lambda[x_1, \ldots, x_s]_{\mathfrak m'} \subset B_{\mathfrak m_B}
\subset P$，并且 $P$ 被认同为 $B$ 在 $\mathfrak m_B$ 处的完备化；
参见《环同态的光滑化》定理
\ref{smoothing-theorem-approximation-property-variant} 之前的评注；
\item $a_i - a'_i, b_i - b'_i, c_{ij} - c'_{ij}, k_{ij} - k'_{ij} \in
(\mathfrak m')^N P$；并且
\item $f_j(a'_1, \ldots, a'_u) = \sum c'_{ji} b'_i$ 且
$\sum k'_{ij}b'_j = 0$。
\end{enumerate}
令 $A = B/(b'_1, \ldots, b'_r)$。
% P11-E0230: the frozen source omits "by" in "denote m_A the image".
把 $\mathfrak m_B$ 在 $A$ 中的像记为 $\mathfrak m_A$。（注意，$A$
在 $\Lambda$ 上本质有限型；证明末尾将说明如何得到在 $\Lambda$
上有限型的 $A$。）存在环同态 $C \to A$，把 $y_i \mapsto a'_i$，
因为 $a'_i$ 模 $(b'_1, \ldots, b'_r)$ 满足所需方程。注意，由上面的
性质 (2)，作为 $P = B^\wedge$ 的商，有
$A/\mathfrak m_A^N = R/\mathfrak m_R^N$。令
$x_A = x_C|_{\Spec(A)}$。由于映射
$$
C \to A \to A/\mathfrak m_A^N \cong R/\mathfrak m_R^N
\quad\text{且}\quad
C \to R \to R/\mathfrak m_R^N
$$
相等，经同构 $A/\mathfrak m_A^N = R/\mathfrak m_R^N$ 可知 $x_A$
和 $x$ 模 $\mathfrak m_R^N$ 相符。至此，已经证明引理陈述中的性质
(1)--(5)。为证明 (6)，注意
$$
P^{\oplus t} \xrightarrow{K}
P^{\oplus r} \xrightarrow{(b_1, \ldots, b_r)}
P
\quad\text{且}\quad
P^{\oplus t} \xrightarrow{K'}
P^{\oplus r} \xrightarrow{(b'_1, \ldots, b'_r)}
P
$$
是两个 $P$-模复形，它们模 $(\mathfrak m')^N$ 同余，且第一个正合。
根据上面对 $N$ 的选取，由《更多代数》引理
\ref{more-algebra-lemma-approximate-complex-graded} 可知，
$R = P/(b_1, \ldots, b_r)$ 和
$P/(b'_1, \ldots, b'_r) = B^\wedge/(b'_1, \ldots, b'_r) = A^\wedge$
具有同构的相伴分次代数。这正是所要证明的。

\medskip\noindent
证明的最后一段用来处理 $A$ 在 $S$ 上本质有限型、但尚非有限型的问题。
上面的构造给出 $A = B/(b'_1, \ldots, b'_r)$ 和
$\mathfrak m_A \subset A$，其中 $B$ 在
$\Lambda[x_1, \ldots, x_s]_{\mathfrak m'}$ 上为 \'etale。因此，$A$
在 Noether 环 $\Lambda[x_1, \ldots, x_s]_{\mathfrak m'}$ 上有限型。
% P11-E0231: the frozen source changes the polynomial-variable upper index from s to n without introducing n.
于是，对某个有限型 $\Lambda[x_1, \ldots, x_n]$-代数 $A_0$，可写成
$A = (A_0)_{\mathfrak m'}$。然后
$A = \colim (A_0)_f$，其中
$f \in \Lambda[x_1, \ldots, x_n] \setminus \mathfrak m'$；参见
《代数》引理 \ref{algebra-lemma-localization-colimit}。由于
$p : \mathcal{X} \to (\Sch/S)_{fppf}$ 在对象上保极限，可知对某个
上述 $f$，$x_A$ 来自 $\Spec((A_0)_f)$ 上的某个对象
$x_{(A_0)_f}$。用 $(A_0)_f$ 代替 $A$、用 $x_{(A_0)_f}$ 代替
$x_A$，并用 $(A_0)_f \cap \mathfrak m_A$ 代替 $\mathfrak m_A$ 后，
证明完成。
\end{proof}





\section{保极限性}
\label{artin-section-limits}

\noindent
若下述定义中的函子本质满，则态射
$p : \mathcal{X} \to (\Sch/S)_{fppf}$ 在《可表示性判据》一节
\ref{criteria-section-limit-preserving} 所定义的意义下在对象上保极限。
不过，《例》一节 \ref{examples-section-limit-preserving} 中的例子表明，
这并不等价于保极限。

\begin{definition}
\label{artin-definition-limit-preserving}
设 $S$ 是概形，$\mathcal{X}$ 是 $(\Sch/S)_{fppf}$ 上的
群胚纤维化范畴。若对每个 $S$ 上的仿射概形 $T$，只要它是
$S$ 上仿射概形 $T_i$ 的某个有向逆系统的极限 $T = \lim T_i$，
就有纤维范畴的等价
$$
\colim \mathcal{X}_{T_i} \longrightarrow \mathcal{X}_T
$$
，则称 $\mathcal{X}$ {\it 保极限}。
\end{definition}

\noindent
详细说明其含义。首先，给定 $\mathcal{X}$ 在 $T_i$ 上的对象 $x, y$，
应有
$$
\Mor_{\mathcal{X}_T}(x|_T, y|_T) =
\colim_{i' \geq i} \Mor_{\mathcal{X}_{T_{i'}}}(x|_{T_{i'}}, y|_{T_{i'}})
$$
；其次，$\mathcal{X}_T$ 的每个对象都同构于某个 $T_i$ 上对象的限制，
其中 $i$ 适当选取。注意，第一个条件意味着预层
$\mathit{Isom}_\mathcal{X}(x, y)$ 保极限（参见《叠》定义
\ref{stacks-definition-mor-presheaf}）。

\begin{lemma}
\label{artin-lemma-fibre-product-limit-preserving}
设 $S$ 是概形。设 $p : \mathcal{X} \to \mathcal{Y}$ 和
$q : \mathcal{Z} \to \mathcal{Y}$ 是 $(\Sch/S)_{fppf}$ 上
群胚纤维化范畴的 $1$-态射。
\begin{enumerate}
\item 若 $\mathcal{X} \to (\Sch/S)_{fppf}$ 和
$\mathcal{Z} \to (\Sch/S)_{fppf}$ 在对象上保极限，且
$\mathcal{Y}$ 保极限，则
$\mathcal{X} \times_\mathcal{Y} \mathcal{Z} \to (\Sch/S)_{fppf}$
在对象上保极限。
\item 若 $\mathcal{X}$、$\mathcal{Y}$ 和 $\mathcal{Z}$ 保极限，则
$\mathcal{X} \times_\mathcal{Y} \mathcal{Z}$ 也保极限。
\end{enumerate}
\end{lemma}

\begin{proof}
这是形式论证。证明 (1)。设 $T = \lim_{i \in I} T_i$ 是 $S$ 上
仿射概形 $T_i$ 的有向极限。我们将证明函子
$\colim \mathcal{X}_{T_i} \to \mathcal{X}_T$ 本质满。回忆，纤维积在
$T$ 上的一个对象是四元组 $(T, x, z, \alpha)$，其中 $x$ 是
$\mathcal{X}$ 位于 $T$ 上的对象，$z$ 是 $\mathcal{Z}$ 位于 $T$
上的对象，而 $\alpha : p(x) \to q(z)$ 是 $\mathcal{Y}$ 在 $T$
上的纤维范畴中的态射。由对 $\mathcal{X}$ 和 $\mathcal{Z}$ 的假设，
可以找到 $i$ 以及 $T_i$ 上的对象 $x_i$ 和 $z_i$，使得
% P11-E0232: the frozen source writes x_i|_T congruent T where the intended object is x.
$x_i|_T \cong T$ 且 $z_i|_T \cong z$。于是，$\alpha$ 对应于同构
$p(x_i)|_T \to q(z_i)|_T$。由对 $\mathcal{Y}$ 的假设，此同构来自同构
$\alpha_{i'} : p(x_i)|_{T_{i'}} \to q(z_i)|_{T_{i'}}$。用 $i'$ 代替
$i$、用 $x_i|_{T_{i'}}$ 代替 $x_i$，并用 $z_i|_{T_{i'}}$ 代替
$z_i$ 后，可知 $(T_i, x_i, z_i, \alpha_i)$ 是纤维积在 $T_i$
上的一个对象，而它限制为 $T$ 上同构于 $(T, x, z, \alpha)$ 的对象，
如所要求。

\medskip\noindent
% P11-E0233: the frozen proof of part (2) names the X-fibre functor rather than the fibre-product functor whose full faithfulness is at issue.
在 (2) 中，证明 $\colim \mathcal{X}_{T_i} \to \mathcal{X}_T$
全忠实的论证从略。
\end{proof}

\begin{lemma}
\label{artin-lemma-limit-preserving-algebraic-space}
设 $S$ 是概形，$\mathcal{X}$ 是 $S$ 上的代数叠。则下列条件等价：
\begin{enumerate}
\item $\mathcal{X}$ 是广群叠，且
$\mathcal{X} \to (\Sch/S)_{fppf}$ 在对象上保极限；
\item $\mathcal{X}$ 是广群叠且保极限；
\item $\mathcal{X}$ 可由一个局部有限表示的代数空间表示。
\end{enumerate}
\end{lemma}

\begin{proof}
在三个条件中的每一个之下，$\mathcal{X}$ 都可由 $S$ 上的代数空间
$X$ 表示；参见《代数叠》命题
\ref{algebraic-proposition-algebraic-stack-no-automorphisms}。
(1) 与 (2) 显然等价，因为广群之间的函子是等价，当且仅当它在同构类上
为满射。最后，由《空间的极限》命题
\ref{spaces-limits-proposition-characterize-locally-finite-presentation}，
(1) 与 (3) 等价。
\end{proof}

\begin{lemma}
\label{artin-lemma-diagonal}
设 $S$ 是概形，$\mathcal{X}$ 是 $(\Sch/S)_{fppf}$ 上的
群胚纤维化范畴。假设
$\Delta : \mathcal{X} \to \mathcal{X} \times \mathcal{X}$ 可由代数空间
表示，且 $\mathcal{X}$ 保极限。则 $\Delta$ 局部有限型。
\end{lemma}

\begin{proof}
应用《可表示性判据》引理
\ref{criteria-lemma-check-property-limit-preserving}。
% P11-E0234: the frozen source redundantly says "Let V be an affine scheme V".
设 $V$ 是在 $S$ 上局部有限表示的仿射概形 $V$，并设 $\theta$ 是
$\mathcal{X} \times \mathcal{X}$ 在 $V$ 上的对象。设 $F_\theta$ 为表示
$\mathcal{X} \times_{\Delta, \mathcal{X} \times \mathcal{X}, \theta}
(\Sch/V)_{fppf}$ 的代数空间，并设 $f_\theta : F_\theta \to V$ 为
典范态射（参见《代数叠》一节
\ref{algebraic-section-morphisms-representable-by-algebraic-spaces}）。
只需证明 $F_\theta \to V$ 具有相应性质。由引理
\ref{artin-lemma-fibre-product-limit-preserving} 和
\ref{artin-lemma-limit-preserving-algebraic-space}，$F_\theta \to S$
局部有限表示。因此，由《空间的态射》引理
\ref{spaces-morphisms-lemma-permanence-finite-type}，
$F_\theta \to V$ 局部有限型。
\end{proof}






\section{通用性}
\label{artin-section-versality}

\noindent
上一节说明了如何用代数对象逼近完备局部环上的对象。不过，为证明叠
$\mathcal{X}$ 是代数叠，需要找到从概形到 $\mathcal{X}$ 的光滑
$1$-态射。由于不预先假设 $\mathcal{X}$ 具有可表对角态射，甚至无法谈论
到 $\mathcal{X}$ 的光滑态射。因此，我们借用形变理论的术语，引入通用对象。

\begin{definition}
\label{artin-definition-versal-formal-object}
设 $S$ 是局部 Noether 概形，并设
$p : \mathcal{X} \to (\Sch/S)_{fppf}$ 是群胚纤维化范畴。设
$\xi = (R, \xi_n, f_n)$ 是形式对象。令 $k = R/\mathfrak m$ 且
$x_0 = \xi_1$。若把 $\xi$ 视为 $\mathcal{F}_{\mathcal{X}, k, x_0}$
的形式对象（评注 \ref{artin-remark-formal-objects-match}）时，它在
《形式形变理论》定义 \ref{formal-defos-definition-versal} 的意义下通用，
则称 $\xi$ 是{\it 通用的}。
\end{definition}

\noindent
简要说明其含义。采用定义中的记号，假设给定 $\mathcal{X}$ 的态射
$\xi_1 = x_0 \to y \to z$，它们位于闭浸入
$\Spec(k) \to \Spec(A) \to \Spec(B)$
之上，其中 $A, B$ 是剩余域为 $k$ 的 Artin 局部环。再假设给定
$n \geq 1$ 和交换图
$$
\vcenter{
\xymatrix{
& y \ar[ld] \\
\xi_n & \xi_1 \ar[u] \ar[l]
}
}
\quad\text{位于下列对象之上：}\quad
\vcenter{
\xymatrix{
& \Spec(A) \ar[ld] \\
\Spec(R/\mathfrak m^n) & \Spec(k) \ar[u] \ar[l]
}
}
$$
通用性意味着，对任意上述数据，都存在 $m \geq n$ 和交换图
$$
\vcenter{
\xymatrix{
& & z \ar[lldd] \\
& & y \ar[ld] \ar[u] \\
\xi_m & \xi_n \ar[l] & \xi_1 \ar[u] \ar[l]
}
}
\quad\text{位于下列对象之上：}
\vcenter{
\xymatrix{
& & \Spec(B) \ar[lldd] \\
& & \Spec(A) \ar[ld] \ar[u] \\
\Spec(R/\mathfrak m^m) &
\Spec(R/\mathfrak m^n) \ar[l] &
\Spec(k) \ar[u] \ar[l]
}
}
$$
请与《形式形变理论》评注 \ref{formal-defos-remark-versal-object} 比较。

\medskip\noindent
设 $S$ 是局部 Noether 概形，$U$ 是 $S$ 上的概形，其结构态射
$U \to S$ 局部有限型。设 $u_0 \in U$ 是 $U$ 的有限型点；参见
《态射》定义 \ref{morphisms-definition-finite-type-point}。令
$k = \kappa(u_0)$。注意，复合 $\Spec(k) \to S$ 也是有限型的；
参见《态射》引理 \ref{morphisms-lemma-composition-finite-type}。设
$p : \mathcal{X} \to (\Sch/S)_{fppf}$ 是群胚纤维化范畴。设 $x$
是 $\mathcal{X}$ 位于 $U$ 上的对象。
% P11-E0235: the frozen source omits "by" in "Denote x_0 the pullback".
把 $x$ 由 $u_0$ 的拉回记为 $x_0$。由 $2$-Yoneda 引理，$x$ 对应于
$1$-态射
$$
x : (\Sch/U)_{fppf} \longrightarrow \mathcal{X},
$$
；参见《代数叠》一节 \ref{algebraic-section-2-yoneda}。由此得到
预形变范畴的态射
\begin{equation}
\label{artin-equation-hat-x}
\hat x :
\mathcal{F}_{(\Sch/U)_{fppf}, k, u_0}
\longrightarrow
\mathcal{F}_{\mathcal{X}, k, x_0},
\end{equation}
，它位于 $\mathcal{C}_\Lambda$ 之上；参见
(\ref{artin-equation-functoriality})。

\begin{definition}
\label{artin-definition-versal}
设 $S$ 是局部 Noether 概形，$\mathcal{X}$ 是
$(\Sch/S)_{fppf}$ 上的群胚纤维化范畴。设 $U$ 是 $S$ 上局部有限型的
概形，$x$ 是 $\mathcal{X}$ 位于 $U$ 上的对象，而 $u_0$ 是 $U$
的有限型点。若态射 $\hat x$（式 (\ref{artin-equation-hat-x})）光滑，
则称 $x$ 在 $u_0$ 处{\it 通用}；参见《形式形变理论》定义
\ref{formal-defos-definition-smooth-morphism}。
\end{definition}

\noindent
这个定义与 $\mathcal{X}$ 的形式对象的通用性概念相吻合。

\begin{lemma}
\label{artin-lemma-versality-matches}
% P11-E0236: the frozen source ends the introductory "With notation" phrase with a period, leaving a sentence fragment.
采用定义 \ref{artin-definition-versal} 中的记号。令
$R = \mathcal{O}_{U, u_0}^\wedge$。令 $\xi$ 为与
$x|_{\Spec(R)}$ 相联系的 $\mathcal{X}$ 在 $R$ 上的形式对象；
参见 (\ref{artin-equation-approximation})。则
$$
x\text{ is versal at }u_0
\Leftrightarrow
\xi\text{ is versal}
$$
\end{lemma}

\begin{proof}
注意，$\mathcal{O}_{U, u_0}$ 是剩余域为 $k$ 的 Noether 局部
$S$-代数。因此，$R = \mathcal{O}_{U, u_0}^\wedge$ 是
$\mathcal{C}_\Lambda^\wedge$ 的对象；参见《形式形变理论》定义
\ref{formal-defos-definition-completion-CLambda}。回忆，$\xi$ 通用是指
$\underline{\xi} : \underline{R}|_{\mathcal{C}_\Lambda} \to
\mathcal{F}_{\mathcal{X}, k, x_0}$
光滑，而 $x$ 在 $u_0$ 处通用是指
$\hat x : \mathcal{F}_{(\Sch/U)_{fppf}, k, u_0}
\to \mathcal{F}_{\mathcal{X}, k, x_0}$ 光滑。存在预形变范畴的认同
$$
\underline{R}|_{\mathcal{C}_\Lambda}
=
\mathcal{F}_{(\Sch/U)_{fppf}, k, u_0},
$$
；记号参见《形式形变理论》评注
\ref{formal-defos-remark-formal-objects-yoneda}。具体地，给定 Artin
局部 $S$-代数 $A$，并把它的剩余域与 $k$ 认同，就有
$$
\Mor_{\mathcal{C}_\Lambda^\wedge}(R, A) =
\{\varphi \in \Mor_S(\Spec(A), U) \mid \varphi|_{\Spec(k)} = u_0\}
$$
。展开定义即可验证，所得映射
$$
\underline{R}|_{\mathcal{C}_\Lambda} =
\mathcal{F}_{(\Sch/U)_{fppf}, k, u_0}
\xrightarrow{\hat x}
\mathcal{F}_{\mathcal{X}, k, x_0},
$$
等于 $\underline{\xi}$，故引理成立。
\end{proof}

\noindent
下面给出一个合理性检验。

\begin{lemma}
\label{artin-lemma-versal-implies-smooth}
设 $S$ 是局部 Noether 概形。设 $f : U \to V$ 是在 $S$ 上局部有限型的
概形态射，并设 $u_0 \in U$ 是有限型点。下列条件等价：
\begin{enumerate}
\item $f$ 在 $u_0$ 处光滑；
\item 把 $f$ 视为 $(\Sch/V)_{fppf}$ 位于 $U$ 上的对象时，
它在 $u_0$ 处通用。
\end{enumerate}
\end{lemma}

\begin{proof}
这只是《更多态射》引理
\ref{more-morphisms-lemma-lifting-along-artinian-at-point} 的重述。
\end{proof}

\noindent
这个概念对于域扩张具有良好行为。

\begin{lemma}
\label{artin-lemma-versal-change-of-field}
设 $S$、$\mathcal{X}$、$U$、$x$、$u_0$ 如定义
\ref{artin-definition-versal}。设 $l$ 是域，并设
$u_{l, 0} : \Spec(l) \to U$ 是像为 $u_0$ 的态射，且
% P11-E0237: the frozen source equates the extension l/k with the residue field kappa(u_0), rather than defining k as that residue field and requiring l/k finite.
$l/k = \kappa(u_0)$ 是有限的。令 $x_{l, 0} = x_0|_{\Spec(l)}$。
若 $\mathcal{X}$ 满足 (RS)，且 $x$ 在 $u_0$ 处通用，则
$$
\mathcal{F}_{(\Sch/U)_{fppf}, l, u_{l, 0}}
\longrightarrow
\mathcal{F}_{\mathcal{X}, l, x_{l, 0}}
$$
光滑。
\end{lemma}

\begin{proof}
注意，由引理 \ref{artin-lemma-algebraic-stack-RS}，
$(\Sch/U)_{fppf}$ 满足 (RS)。因此，引理中的函子就是与 $\hat x$
相联系的函子
$$
(\mathcal{F}_{(\Sch/U)_{fppf}, k , u_0})_{l/k}
\longrightarrow
(\mathcal{F}_{\mathcal{X}, k , x_0})_{l/k}
$$
；参见引理 \ref{artin-lemma-change-of-field}。因而，引理由《形式形变理论》引理
\ref{formal-defos-lemma-change-of-fields-smooth} 得到。
\end{proof}

\noindent
下述引理是另一个合理性检验。其含义大致是：若 $x$ 如定义
\ref{artin-definition-versal} 那样在 $u_0$ 处通用，则把 $x$ 视为从 $U$
到 $\mathcal{X}$ 的态射后，每当用概形作基变换时，它都是光滑的。

\begin{lemma}
\label{artin-lemma-base-change-versal}
设 $S$、$\mathcal{X}$、$U$、$x$、$u_0$ 如定义
\ref{artin-definition-versal}。假设
\begin{enumerate}
\item $\Delta : \mathcal{X} \to \mathcal{X} \times \mathcal{X}$
可由代数空间表示；
\item $\Delta$ 局部有限型（例如，当 $\mathcal{X}$ 保极限时）；并且
\item $\mathcal{X}$ 具有 (RS)。
\end{enumerate}
设 $V$ 是 $S$ 上局部有限型的概形，而 $y$ 是 $\mathcal{X}$ 在 $V$
上的对象。作 $2$-纤维积
$$
\xymatrix{
\mathcal{Z} \ar[r] \ar[d] & (\Sch/U)_{fppf} \ar[d]^x \\
(\Sch/V)_{fppf} \ar[r]^y & \mathcal{X}
}
$$
设 $Z$ 是表示 $\mathcal{Z}$ 的代数空间，并设 $z_0 \in |Z|$ 是位于
$u_0$ 之上的有限型点。若 $x$ 在 $u_0$ 处通用，则态射
$Z \to V$ 在 $z_0$ 处光滑。
\end{lemma}

\begin{proof}
（陈述中的括号评注由引理 \ref{artin-lemma-diagonal} 得到。）由假设 (1)
和《代数叠》引理 \ref{algebraic-lemma-representable-diagonal}，$Z$ 存在。
由假设 (2)，$Z \to V \times_S U$ 局部有限型。选取概形 $W$、闭点
$w_0 \in W$，以及把 $w_0$ 映到 $z_0$ 的 \'etale 态射 $W \to Z$；
参见《空间的态射》定义
\ref{spaces-morphisms-definition-finite-type-point}。于是 $W$ 在 $S$
上局部有限型，且 $w_0$ 是 $W$ 的有限型点。令 $l = \kappa(z_0)$。
% P11-E0238: the frozen source omits "by" in the four-object naming construction.
把通过拉回到 $\Spec(l) = w_0$ 所得的 $\mathcal{Z}$、
$(\Sch/V)_{fppf}$、$(\Sch/U)_{fppf}$ 和 $\mathcal{X}$ 在 $\Spec(l)$
上的对象分别记为 $z_{l, 0}$、$v_{l, 0}$、$u_{l, 0}$ 和
$x_{l, 0}$。考虑
$$
\xymatrix{
\mathcal{F}_{(\Sch/W)_{fppf}, l, w_0} \ar[r] &
\mathcal{F}_{\mathcal{Z}, l, z_{l, 0}} \ar[d] \ar[r] &
\mathcal{F}_{(\Sch/U)_{fppf}, l, u_{l, 0}} \ar[d] \\
& \mathcal{F}_{(\Sch/V)_{fppf}, l, v_{l, 0}} \ar[r] &
\mathcal{F}_{\mathcal{X}, l, x_{l, 0}}
}
$$
由引理 \ref{artin-lemma-fibre-product-deformation-categories}，此方块是
预形变范畴的纤维积。由引理 \ref{artin-lemma-versal-change-of-field}，
右边的竖箭头光滑。由《形式形变理论》引理
\ref{formal-defos-lemma-smooth-properties}，左边的竖箭头光滑。
由引理 \ref{artin-lemma-formally-smooth-on-deformation-categories}，
左边的横箭头光滑。由此断定映射
$$
\mathcal{F}_{(\Sch/W)_{fppf}, l, w_0} \to
\mathcal{F}_{(\Sch/V)_{fppf}, l, v_{l, 0}}
$$
由《形式形变理论》引理 \ref{formal-defos-lemma-smooth-properties}
是光滑的。因此，由《更多态射》引理
\ref{more-morphisms-lemma-lifting-along-artinian-at-point}，
$W \to V$ 在 $w_0$ 处光滑。这恰好意味着 $Z \to V$ 在 $z_0$
处光滑，证明完成。
\end{proof}

\noindent
用通用对象的语言重述逼近结果。

\begin{lemma}
\label{artin-lemma-approximate-versal}
设 $S$ 是局部 Noether 概形，并设
$p : \mathcal{X} \to (\Sch/S)_{fppf}$ 是群胚纤维化范畴。设
$\xi = (R, \xi_n, f_n)$ 是 $\mathcal{X}$ 的形式对象，其中 $\xi_1$
位于 $\Spec(k) \to S$ 之上，而该态射的像为 $s \in S$。假设
\begin{enumerate}
\item $\xi$ 通用；
\item $\xi$ 有效；
\item $\mathcal{O}_{S, s}$ 是 G-环；并且
\item $p : \mathcal{X} \to (\Sch/S)_{fppf}$ 在对象上保极限。
\end{enumerate}
则存在有限型态射 $U \to S$、剩余域为 $k$ 的有限型点
$u_0 \in U$，以及 $\mathcal{X}$ 在 $U$ 上的对象 $x$，使得 $x$
在 $u_0$ 处通用，并且
$x|_{\Spec(\mathcal{O}_{U, u_0}/\mathfrak m_{u_0}^n)} \cong \xi_n$。
\end{lemma}

\begin{proof}
选取 $\mathcal{X}$ 位于 $\Spec(R)$ 上的对象 $x_R$，使其相联系的形式对象
为 $\xi$。令 $N = 2$，并应用引理 \ref{artin-lemma-approximate}，得到
$A, \mathfrak m_A, x_A, \ldots$。令
$\eta = (A^\wedge, \eta_n, g_n)$ 为与
$x_A|_{\Spec(A^\wedge)}$ 相联系的形式对象。存在图
$$
\vcenter{
\xymatrix{
& \eta \ar[d] \\
\xi \ar[r] \ar@{..>}[ru] & \xi_2 = \eta_2
}
}
\quad\text{位于下列对象之上：}\quad
\vcenter{
\xymatrix{
& A^\wedge \ar[d] \\
R \ar[r] \ar@{..>}[ru] & R/\mathfrak m_R^2 = A/\mathfrak m_A^2
}
}
$$
由 $\xi$ 的通用性，恰好可以找到诸图中的虚线箭头，因为由
《形式形变理论》评注 \ref{formal-defos-remark-versal-object}，
可以逐次找到态射 $\xi \to \eta_3$、$\xi \to \eta_4$ 等等。
由《形式形变理论》引理
\ref{formal-defos-lemma-surjective-cotangent-space}，相应环同态
$R \to A^\wedge$ 是满射。另一方面，由构造，对所有 $n$ 都有
$\dim_k \mathfrak m_R^n/\mathfrak m_R^{n + 1} =
\dim_k \mathfrak m_A^n/\mathfrak m_A^{n + 1}$。因此，由长度的可加性和
《形式形变理论》引理 \ref{formal-defos-lemma-length}，
$R/\mathfrak m_R^n$ 和 $A/\mathfrak m_A^n$ 作为 $\Lambda$-模具有相同的
有限长度。于是，对所有 $n$，
$R/\mathfrak m_R^n \to A/\mathfrak m_A^n$ 都是同构，从而
$R \to A^\wedge$ 是同构。因此，$\eta$ 同构于通用对象，故其自身也通用。
由引理 \ref{artin-lemma-versality-matches}，$x_A$ 在
$U = \Spec(A)$ 中与 $\mathfrak m_A$ 对应的点 $u_0$ 处通用。
\end{proof}

\begin{example}
\label{artin-example-approximate-versal-implies}
本例说明，要使引理 \ref{artin-lemma-approximate-versal} 的结论成立，局部环
$\mathcal{O}_{S, s}$ 必须是 G-环。具体地，设 $\Lambda$ 是 Noether 环，
$\mathfrak m$ 是 $\Lambda$ 的极大理想。令
$R = \Lambda_\mathfrak m^\wedge$。设 $\Lambda \to C \to R$ 是一个分解，
其中 $C$ 在 $\Lambda$ 上有限型。令 $S = \Spec(\Lambda)$、
$U = S \setminus \{\mathfrak m\}$，以及 $S' = U \amalg \Spec(C)$。
考虑由下述规则定义的函子
$F : (\Sch/S)_{fppf}^{opp} \to \textit{Sets}$：
$$
F(T) = 
\left\{
\begin{matrix}
* & \text{若 }T \to S\text{ factors through }S' \\
\emptyset & \text{否则}
\end{matrix}
\right.
$$
% P11-E0239: the frozen source says "Let X = S_F is the category", combining incompatible let- and copular constructions.
令 $\mathcal{X} = \mathcal{S}_F$ 是与 $F$ 相联系的集合纤维化范畴；
参见《代数叠》一节 \ref{algebraic-section-split}。于是
$\mathcal{X} \to (\Sch/S)_{fppf}$ 在对象上保极限，并且在 $R$ 上存在
一个有效的通用形式对象 $\xi$。因此，若引理
\ref{artin-lemma-approximate-versal} 的结论对 $\mathcal{X}$ 成立，则存在
有限型环同态 $\Lambda \to A$ 和位于 $\mathfrak m$ 之上的极大理想
$\mathfrak m_A$，使得
\begin{enumerate}
\item $\kappa(\mathfrak m) = \kappa(\mathfrak m_A)$；
\item $\Lambda \to A$ 和 $\mathfrak m_A$ 满足《代数》引理
\ref{algebra-lemma-smooth-test-artinian} 的条件 (4)；并且
\item 存在 $\Lambda$-代数同态 $C \to A$。
\end{enumerate}
由所引引理，$\Lambda \to A$ 在 $\mathfrak m_A$ 处光滑。对 $A$ 作切片后，
可假设 $\Lambda \to A$ 在 $\mathfrak m_A$ 处为 \'etale；例如参见
《更多态射》引理 \ref{more-morphisms-lemma-slice-smooth}，也可以直接论证。
写 $C = \Lambda[y_1, \ldots, y_n]/(f_1, \ldots, f_m)$。于是，
$C \to R$ 对应于方程组 $f_1 = \ldots = f_m = 0$ 在 $R$ 中的一个解；
参见《环同态的光滑化》一节
\ref{smoothing-section-approximation-G-rings}。因此，若引理
\ref{artin-lemma-approximate-versal} 的结论对每个上述 $\mathcal{X}$ 成立，
则在 $R$ 中有解的方程组在 $\Lambda_{\mathfrak m}$ 的 Hensel 化中也有解。
换言之，逼近性质对 $\Lambda_{\mathfrak m}^h$ 成立。这蕴含
$\Lambda_{\mathfrak m}^h$ 是 G-环
% P11-E0240: the frozen source contains a second unresolved future-reference placeholder.
（此处补入将来的参考文献；另见《环同态的光滑化》一节
\ref{smoothing-section-introduction} 中的讨论），进而蕴含
$\Lambda_{\mathfrak m}$ 是 G-环。
\end{example}






\section{通用性的开性}
\label{artin-section-openness-versality}

\noindent
接下来讨论通用性的开性。

\begin{definition}
\label{artin-definition-openness-versality}
设 $S$ 是局部 Noether 概形。
\begin{enumerate}
\item 设 $\mathcal{X}$ 是 $(\Sch/S)_{fppf}$ 上的群胚纤维化范畴。
若给定 $S$ 上局部有限型的概形 $U$、$\mathcal{X}$ 位于 $U$ 上的对象
$x$ 和有限型点 $u_0 \in U$，且 $x$ 在 $u_0$ 处通用时，总存在开邻域
$u_0 \in U' \subset U$，使得 $x$ 在 $U'$ 的每个有限型点处都通用，
则称 $\mathcal{X}$ 满足{\it 通用性的开性}。
\item 设 $f : \mathcal{X} \to \mathcal{Y}$ 是 $(\Sch/S)_{fppf}$ 上
群胚纤维化范畴的 $1$-态射。若给定 $S$ 上局部有限型的概形 $U$ 和
$\mathcal{Y}$ 位于 $U$ 上的对象 $y$ 时，
$(\Sch/U)_{fppf} \times_\mathcal{Y} \mathcal{X}$ 满足通用性的开性，
则称 $f$ 满足{\it 通用性的开性}。
\end{enumerate}
\end{definition}

\noindent
通用性的开性往往是最难核验的条件。不过，下面的例子表明这项要求确属必要。

\begin{example}
\label{artin-example-versality}
设 $k$ 是域，并令 $\Lambda = k[s, t]$。考虑依下述规则定义的函子
$F : \Lambda\text{-algebras} \longrightarrow \textit{Sets}$
$$
F(A) =
\left\{
\begin{matrix}
* & \text{if there exist }f_1, \ldots, f_n \in A\text{ 使得 } \\
  & A = (s, t, f_1, \ldots, f_n)\text{ 且 } f_i s = 0\ \forall i \\
\emptyset & \text{否则}
\end{matrix}
\right.
$$
从几何上说，$F(A) = *$ 意味着 $V(s, t) \subset \Spec(A)$ 有一个
拟紧开邻域 $W$，使得 $s|_W = 0$。令
$\mathcal{X} \subset (\Sch/\Spec(\Lambda))_{fppf}$ 为如下概形 $T$
所成的全子范畴：它有仿射开覆盖 $T = \bigcup \Spec(A_j)$，且对每个
$j$ 都有 $F(A_j) = *$。于是 $\mathcal{X}$ 满足 [0]、[1]、[2]、[3]
和 [4]，但不满足 [5]。具体地，在
$U = \Spec(k[s, t]/(s))$ 上存在对象 $x$，它在 $u_0 = (s, t)$ 处通用，
但不在任何其他点处通用。细节从略。
\end{example}

\noindent
设 $S$ 是局部 Noether 概形。设
$f : \mathcal{X} \to \mathcal{Y}$ 是 $(\Sch/S)_{fppf}$ 上群胚纤维化
范畴的 $1$-态射。考虑如下性质
\begin{equation}
\label{artin-equation-smooth}
\begin{matrix}
\text{for all fields }k\text{ of finite type over }S
\text{ and all }x_0 \in \Ob(\mathcal{X}_{\Spec(k)})\text{ the}\\
\text{map }
\mathcal{F}_{\mathcal{X}, k, x_0} \to \mathcal{F}_{\mathcal{Y}, k, f(x_0)}
\text{ of predeformation categories is smooth}
\end{matrix}
\end{equation}
下面围绕这一概念表述若干引理。首先把它与通用性（及其开性）联系起来。

\begin{lemma}
\label{artin-lemma-versal-smooth}
设 $S$ 是局部 Noether 概形，$\mathcal{X}$ 是 $(\Sch/S)_{fppf}$ 上的
群胚纤维化范畴。设 $U$ 是 $S$ 上局部有限型的概形，$x$ 是
$\mathcal{X}$ 位于 $U$ 上的对象。假设 $x$ 在 $U$ 的每个有限型点处
都通用，且 $\mathcal{X}$ 满足 (RS)。则
$x : (\Sch/U)_{fppf} \to \mathcal{X}$ 满足 (\ref{artin-equation-smooth})。
\end{lemma}

\begin{proof}
设 $\Spec(l) \to U$ 是态射，且 $l$ 在 $S$ 上有限型。则其像
$u_0 \in U$ 是 $U$ 的有限型点，而 $l/\kappa(u_0)$ 是有限扩张；参见
《态射》一节 \ref{morphisms-section-points-finite-type} 中的讨论。因此，
$\mathcal{F}_{(\Sch/U)_{fppf}, l, u_{l, 0}} \to
\mathcal{F}_{\mathcal{X}, l, x_{l, 0}}$
由引理 \ref{artin-lemma-versal-change-of-field} 可知是光滑的。
\end{proof}

\begin{lemma}
\label{artin-lemma-composition-smooth}
设 $S$ 是局部 Noether 概形。设 $f : \mathcal{X} \to \mathcal{Y}$ 和
$g : \mathcal{Y} \to \mathcal{Z}$ 是 $(\Sch/S)_{fppf}$ 上群胚纤维化
范畴的可复合 $1$-态射。若 $f$ 和 $g$ 满足
(\ref{artin-equation-smooth})，则 $g \circ f$ 也满足它。
\end{lemma}

\begin{proof}
这由《形式形变理论》引理
\ref{formal-defos-lemma-smooth-properties} 形式地推出。
\end{proof}

\begin{lemma}
\label{artin-lemma-base-change-smooth}
设 $S$ 是局部 Noether 概形。设 $f : \mathcal{X} \to \mathcal{Y}$ 和
$\mathcal{Z} \to \mathcal{Y}$ 是 $(\Sch/S)_{fppf}$ 上群胚纤维化范畴的
$1$-态射。若 $f$ 满足 (\ref{artin-equation-smooth})，则投影
$\mathcal{X} \times_\mathcal{Y} \mathcal{Z} \to \mathcal{Z}$ 也满足它。
\end{lemma}

\begin{proof}
% P11-E0241: the frozen source proof begins with the sentence fragment "Follows immediately from".
这立即由引理 \ref{artin-lemma-fibre-product-deformation-categories} 和
《形式形变理论》引理
\ref{formal-defos-lemma-smooth-properties} 得到。
\end{proof}

\begin{lemma}
\label{artin-lemma-smooth-smooth}
% P11-E0242: the frozen source uses plural "morphisms" after the singular article and subject.
设 $S$ 是局部 Noether 概形。设 $f : \mathcal{X} \to \mathcal{Y}$ 是
$(\Sch/S)_{fppf}$ 上群胚纤维化范畴的 $1$-态射。若 $f$ 在对象上形式
光滑，则 $f$ 满足 (\ref{artin-equation-smooth})。若 $f$ 可由代数空间表示且
光滑，则 $f$ 满足 (\ref{artin-equation-smooth})。
\end{lemma}

\begin{proof}
% P11-E0243: the frozen source proof is a noun phrase rather than a complete sentence.
这就是引理 \ref{artin-lemma-formally-smooth-on-deformation-categories} 的改述。
\end{proof}

\begin{lemma}
\label{artin-lemma-implies-smooth}
设 $S$ 是局部 Noether 概形。设 $f : \mathcal{X} \to \mathcal{Y}$ 是
$(\Sch/S)_{fppf}$ 上群胚纤维化范畴的 $1$-态射。假设
\begin{enumerate}
\item $f$ 可由代数空间表示；
\item $f$ 满足 (\ref{artin-equation-smooth})；
\item $\mathcal{X} \to (\Sch/S)_{fppf}$ 在对象上保极限；并且
\item $\mathcal{Y}$ 保极限。
\end{enumerate}
则 $f$ 光滑。
\end{lemma}

\begin{proof}
证明的关键是《更多态射》引理
\ref{more-morphisms-lemma-lifting-along-artinian-at-point}；它（几乎）断言，
$S$ 上有限型概形之间满足 (\ref{artin-equation-smooth}) 的态射是光滑态射。
证明中的其他论证本质上只是组织记号。

\medskip\noindent
设 $V$ 是 $S$ 上的概形，而 $y$ 是 $\mathcal{Y}$ 位于 $V$ 上的对象。
% P11-E0244: the frozen 2-fibre-product subscript names X, although f and y both map to Y.
设 $Z$ 是表示 $2$-纤维积
$\mathcal{Z} = \mathcal{X} \times_{f, \mathcal{X}, y} (\Sch/V)_{fppf}$.
的代数空间。需要证明投影态射 $Z \to V$ 光滑；参见《代数叠》定义
\ref{algebraic-definition-relative-representable-property}。事实上，只须在
$V$ 是 $S$ 上局部有限表示的仿射概形时证明这一点；参见《可表性判据》引理
\ref{criteria-lemma-check-property-limit-preserving}。此时，由引理
\ref{artin-lemma-limit-preserving-algebraic-space}，$(\Sch/V)_{fppf}$ 保极限。
再由引理 \ref{artin-lemma-fibre-product-limit-preserving} 和
\ref{artin-lemma-limit-preserving-algebraic-space}，$Z \to S$ 局部有限表示。
选取概形 $W$ 和满 \'etale 态射 $W \to Z$。于是 $W$ 在 $S$ 上局部有限表示。

\medskip\noindent
由于 $f$ 满足 (\ref{artin-equation-smooth})，由引理
\ref{artin-lemma-base-change-smooth}，$\mathcal{Z} \to (\Sch/V)_{fppf}$ 也满足它。
其次，$(\Sch/W)_{fppf} \to \mathcal{Z}$ 满足
(\ref{artin-equation-smooth})；这是引理 \ref{artin-lemma-smooth-smooth} 的结论。因此，复合
$$
(\Sch/W)_{fppf} \to \mathcal{Z} \to (\Sch/V)_{fppf}
$$
满足 (\ref{artin-equation-smooth})；这由引理 \ref{artin-lemma-composition-smooth} 得到。
《更多态射》引理
\ref{more-morphisms-lemma-lifting-along-artinian-at-point} 表明，复合
$W \to Z \to V$ 在 $W$ 的每个有限型点 $w_0$ 处光滑。由于光滑点集是开的，
由《态射》引理 \ref{morphisms-lemma-enough-finite-type-points}，可知
$W \to V$ 是概形的光滑态射。因此按定义，$Z \to V$ 是代数空间的光滑态射。
\end{proof}

\noindent
下面的引理说明我们将如何使用通用性的开性。

\begin{lemma}
\label{artin-lemma-get-smooth}
设 $S$ 是局部 Noether 概形。设
$p : \mathcal{X} \to (\Sch/S)_{fppf}$ 是群胚纤维化范畴。设 $k$ 是
$S$ 上的有限型域，而 $x_0$ 是 $\mathcal{X}$ 位于 $\Spec(k)$ 上的对象，
其像为 $s \in S$。假设
\begin{enumerate}
\item $\Delta : \mathcal{X} \to \mathcal{X} \times \mathcal{X}$
可由代数空间表示；
\item $\mathcal{X}$ 满足公理 [1]、[2]、[3]（参见一节
\ref{artin-section-axioms}）；
\item $\mathcal{X}$ 的每个形式对象都有效；
\item $\mathcal{X}$ 满足通用性的开性；并且
\item $\mathcal{O}_{S, s}$ 是 G-环。
\end{enumerate}
则存在有限型态射 $U \to S$ 和 $\mathcal{X}$ 位于 $U$ 上的对象 $x$，使得
$$
x : (\Sch/U)_{fppf} \longrightarrow \mathcal{X}
$$
光滑，并且存在有限型点 $u_0 \in U$，其剩余域为 $k$，且
$x|_{u_0} \cong x_0$。
\end{lemma}

\begin{proof}
由公理 [2]、引理 \ref{artin-lemma-deformation-category} 和评注
\ref{artin-remark-deformation-category-implies}，可知
$\mathcal{F}_{\mathcal{X}, k, x_0}$ 满足 (S1) 和 (S2)。再由公理 [3]，
其切空间是有限维的，故《形式形变理论》引理
\ref{formal-defos-lemma-versal-object-existence} 表明
$\mathcal{F}_{\mathcal{X}, k, x_0}$ 有通用形式对象 $\xi$。假设 (3)
断言 $\xi$ 有效。由公理 [1] 和引理 \ref{artin-lemma-approximate-versal}，
存在有限型态射 $U \to S$、$\mathcal{X}$ 位于 $U$ 上的对象 $x$，以及
$U$ 中剩余域为 $k$ 的有限型点 $u_0$，使得 $x$ 在 $u_0$ 处通用，且
$x|_{\Spec(k)} \cong x_0$。由通用性的开性，可以缩小 $U$，并假设 $x$
在 $U$ 的每个有限型点处都通用。我们断言
$$
x : (\Sch/U)_{fppf} \longrightarrow \mathcal{X}
$$
光滑，这就证明了引理。事实上，由引理 \ref{artin-lemma-versal-smooth}，$x$
满足 (\ref{artin-equation-smooth})；再由引理 \ref{artin-lemma-implies-smooth} 即得结论。
\end{proof}






\section{公理}
\label{artin-section-axioms}

\noindent
设 $S$ 是局部 Noether 概形。设
$p : \mathcal{X} \to (\Sch/S)_{fppf}$ 是群胚纤维化范畴。下面列出我们将对
$\mathcal{X}$ 考察的公理。
\begin{enumerate}
% P11-E0245: the frozen source says "is <= than", redundantly combining the inequality with "than".
\item[{[-1]}] 一个集合论条件\footnote{该条件如下：在 $k$ 遍历 $S$ 上
有限型域时，所有基数
$|\Ob(\mathcal{X}_{\Spec(k)})/\cong|$ 和
$|\text{Arrows}(\mathcal{X}_{\Spec(k)})|$ 的上确界 $\leq$
$(\Sch/S)_{fppf}$ 某个对象的大小。}；不关心集合论问题的读者可以忽略它；
\item[{[0]}] $\mathcal{X}$ 是 \'etale 拓扑下的广群叠；
\item[{[1]}] $\mathcal{X}$ 保极限；
\item[{[2]}] $\mathcal{X}$ 满足 Rim-Schlessinger 条件 (RS)；
\item[{[3]}] 对每个 $k$ 和 $x_0$，空间
$T\mathcal{F}_{\mathcal{X}, k, x_0}$ 和
$\text{Inf}(\mathcal{F}_{\mathcal{X}, k, x_0})$
都是有限维的；参见 (\ref{artin-equation-tangent-space}) 和
(\ref{artin-equation-infinitesimal-automorphisms})；
\item[{[4]}] 函子 (\ref{artin-equation-approximation}) 是等价；
\item[{[5]}] $\mathcal{X}$ 和
$\Delta : \mathcal{X} \to \mathcal{X} \times \mathcal{X}$ 都满足
通用性的开性。
\end{enumerate}










\section{函子的公理}
\label{artin-section-axioms-functors}

\noindent
设 $S$ 是概形，并设 $F : (\Sch/S)_{fppf}^{opp} \to \textit{Sets}$ 是函子。
% P11-E0246: the frozen source omits "by" in the naming construction for X.
把与 $F$ 相联系的集合纤维化范畴记为
$\mathcal{X} = \mathcal{S}_F$；参见《代数叠》一节
\ref{algebraic-section-split}。本节说明前述内容在 $\mathcal{X}$ 上的表述
与关于 $F$ 的陈述如何互译。

\medskip\noindent
设 $S$ 是局部 Noether 概形。设
$F : (\Sch/S)_{fppf}^{opp} \to \textit{Sets}$ 是函子。设 $k$ 是 $S$
上的有限型域，并设 $x_0 \in F(\Spec(k))$。相联系的预形变范畴
(\ref{artin-equation-predeformation-category}) 对应于函子
$$
F_{k, x_0} : \mathcal{C}_\Lambda \longrightarrow \textit{Sets},
\quad
A \longmapsto \{x \in F(\Spec(A)) \mid x|_{\Spec(k)} = x_0 \}.
$$
% P11-E0247: the frozen source uses the bare singular "functor" after "between ... and".
回忆，我们不区分 $\mathcal{C}_\Lambda$ 上的集合余纤维化范畴与函子
$\mathcal{C}_\Lambda \to \textit{Sets}$；参见《形式形变理论》评注
\ref{formal-defos-remarks-cofibered-groupoids}
(\ref{formal-defos-item-convention-cofibered-sets}).
给定函子变换 $a : F \to G$，令 $y_0 = a(x_0)$，便得到态射
$$
F_{k, x_0} \longrightarrow G_{k, y_0}
$$
；参见 (\ref{artin-equation-functoriality})。引理
\ref{artin-lemma-formally-smooth-on-deformation-categories} 表明，若
$a : F \to G$ 形式光滑（在《空间的更多态射》定义
\ref{spaces-more-morphisms-definition-formally-smooth-etale-unramified}
的意义下），则 $F_{k, x_0} \longrightarrow G_{k, y_0}$ 在《形式形变理论》
评注 \ref{formal-defos-remark-compare-smooth-schlessinger} 的意义下光滑。

\medskip\noindent
引理 \ref{artin-lemma-pushout} 断言，若在 $S$ 上概形的范畴中
$Y' = Y \amalg_X X'$，其中 $X \to X'$ 是加厚而 $X \to Y$ 是仿射的，则映射
$$
F(Y \amalg_X X') \to F(Y) \times_{F(X)} F(X')
$$
在 $F$ 是代数空间时为双射。对一般函子 $F$，若每当给定推出
$Y' = Y \amalg_X X'$，其中 $Y, X, X'$ 都是 $S$ 上有限型 Artin 局部环
之谱时，映射
$$
F(Y \amalg_X X') \to F(Y) \times_{F(X)} F(X')
$$
都是双射，则称其满足{\it Rim-Schlessinger 条件}，或称 $F$
{\it 满足 (RS)}。因此，每个代数空间都满足 (RS)。

\medskip\noindent
引理 \ref{artin-lemma-deformation-category} 断言，若函子 $F$ 满足 (RS)，则所有
$F_{k, x_0}$ 都是《形式形变理论》定义
\ref{formal-defos-definition-deformation-category} 意义下的形变函子；换言之，
它们满足《形式形变理论》评注
\ref{formal-defos-remark-compare-schlessinger-H4} 意义下的 (RS)。特别地，切空间
$$
TF_{k, x_0} = \{x \in F(\Spec(k[\epsilon])) \mid x|_{\Spec(k)} = x_0\}
$$
由《形式形变理论》引理
\ref{formal-defos-lemma-tangent-space-vector-space} 具有 $k$-向量空间结构。

\medskip\noindent
引理 \ref{artin-lemma-finite-dimension} 断言，$S$ 上局部有限型的代数空间 $F$
给出形变函子 $F_{k, x_0}$，其切空间 $TF_{k, x_0}$ 是有限维的。

\medskip\noindent
$F$ 的一个{\it 形式对象}\footnote{Artin 称之为形式形变。}
$\xi = (R, \xi_n)$ 由如下数据组成：Noether 完备局部 $S$-代数 $R$，
其剩余域在 $S$ 上有限型；以及元素
$\xi_n \in F(\Spec(R/\mathfrak m^n))$，满足
$\xi_{n + 1}|_{\Spec(R/\mathfrak m^n)} = \xi_n$。形式对象 $\xi$ 定义了
$F_{R/\mathfrak m, \xi_1}$ 的形式对象 $\xi$。当且仅当 $\xi$ 在
《形式形变理论》定义 \ref{formal-defos-definition-versal} 的意义下通用时，
称它{\it 通用}。若存在 $x \in F(\Spec(R))$，使得对所有 $n \geq 1$ 都有
$\xi_n = x|_{\Spec(R/\mathfrak m^n)}$，则称形式对象
$\xi = (R, \xi_n)$ {\it 有效}。引理 \ref{artin-lemma-effective} 断言，若 $F$
是代数空间，则每个形式对象都有效。

\medskip\noindent
设 $U$ 是 $S$ 上局部有限型的概形，并设 $x \in F(U)$。设
$u_0 \in U$ 是有限型点。当且仅当
$\xi = (\mathcal{O}_{U, u_0}^\wedge,
x|_{\Spec(\mathcal{O}_{U, u_0}/\mathfrak m_{u_0}^n)})$
是上述意义下的通用形式对象时，称 $x$ 在 $u_0$ 处通用。

\medskip\noindent
设 $S$ 是局部 Noether 概形。设
$F : (\Sch/S)_{fppf}^{opp} \to \textit{Sets}$ 是函子。下面列出我们将对
$F$ 考察的公理。
\begin{enumerate}
% P11-E0248: the frozen source repeats the malformed comparison "is <= than".
\item[{[-1]}] 一个集合论条件\footnote{该条件如下：在 $k$ 遍历 $S$ 上
有限型域时，所有基数 $|F(\Spec(k))|$ 的上确界 $\leq$
$(\Sch/S)_{fppf}$ 某个对象的大小。}；不关心集合论问题的读者可以忽略它；
\item[{[0]}] $F$ 是 \'etale 拓扑下的层；
\item[{[1]}] $F$ 保极限；
\item[{[2]}] $F$ 满足 Rim-Schlessinger 条件 (RS)；
\item[{[3]}] 每个切空间 $TF_{k, x_0}$ 都是有限维的；
\item[{[4]}] 每个形式对象都有效；
\item[{[5]}] $F$ 满足通用性的开性。
\end{enumerate}
这里，{\it 保极限}是《空间的极限》定义
\ref{spaces-limits-definition-locally-finite-presentation} 中的概念，而
{\it 通用性的开性}是指：给定 $S$ 上局部有限型的概形 $U$、
$x \in F(U)$，以及有限型点 $u_0 \in U$，若 $x$ 在 $u_0$ 处通用，
则存在开邻域 $u_0 \in U' \subset U$，使得 $x$ 在 $U'$ 的每个有限型点处
都通用。






\section{代数空间}
\label{artin-section-algebraic-spaces}

\noindent
下面给出关于代数空间的第一个主要结果。

\begin{proposition}
\label{artin-proposition-spaces-diagonal-representable}
设 $S$ 是局部 Noether 概形。设
$F : (\Sch/S)_{fppf}^{opp} \to \textit{Sets}$ 是函子。假设
\begin{enumerate}
\item $\Delta : F \to F \times F$ 可由代数空间表示；
\item $F$ 满足公理 [-1]、[0]、[1]、[2]、[3]、[4]、[5]
（参见一节 \ref{artin-section-axioms-functors}）；并且
\item 对 $S$ 的每个有限型点 $s$，$\mathcal{O}_{S, s}$ 都是 G-环。
\end{enumerate}
则 $F$ 是代数空间。
\end{proposition}

\begin{proof}
引理 \ref{artin-lemma-get-smooth} 适用于 $F$。利用这一点，对 $S$ 上每个有限型域
$k$ 和每个 $x_0 \in F(\Spec(k))$，选取 $S$ 上有限型的仿射概形
$U_{k, x_0}$ 及光滑态射 $U_{k, x_0} \to F$，使得存在剩余域为 $k$
的有限型点 $u_{k, x_0} \in U_{k, x_0}$，而 $u_{k, x_0}$ 的像为 $x_0$。于是
$$
U = \coprod\nolimits_{k, x_0} U_{k, x_0} \longrightarrow F
$$
光滑\footnote{集合论评注：由公理 [-1]，指标集的大小有界，所以这个余积
（同构于）$(\Sch/S)_{fppf}$ 的一个对象；参见《集合》引理
\ref{sets-lemma-what-is-in-it}。}。为完成证明，只须证明此映射是满的；
参见《自举》引理 \ref{bootstrap-lemma-spaces-etale-smooth-cover}
（此处用到公理 [0]）。由《可表性判据》引理
\ref{criteria-lemma-check-property-limit-preserving}，只须对那些 $V \to F$
证明 $U \times_F V \to V$ 是满的，其中 $V$ 是 $S$ 上局部有限表示的仿射
概形。由于 $U \times_F V \to V$ 光滑，其像是开的。因此，只须证明
$U \times_F V \to V$ 的像包含 $V$ 的所有有限型点；参见《态射》引理
\ref{morphisms-lemma-enough-finite-type-points}。设 $v_0 \in V$ 是有限型点。
则 $k = \kappa(v_0)$ 是 $S$ 上的有限型域。
% P11-E0249: the frozen source omits "by" in "Denote x_0 the composition".
把复合 $\Spec(k) \xrightarrow{v_0} V \to F$ 记为 $x_0$。于是
$(u_{k, x_0}, v_0) : \Spec(k) \to U \times_F V$ 是映到 $v_0$ 的点，
即得结论。
\end{proof}

\begin{lemma}
\label{artin-lemma-monomorphism}
设 $S$ 是局部 Noether 概形。设 $a : F \to G$ 是函子
$(\Sch/S)_{fppf}^{opp} \to \textit{Sets}$ 的变换。假设
\begin{enumerate}
\item $a$ 是单射；
\item $F$ 满足公理 [0]、[1]、[2]、[4] 和 [5]；
\item 对 $S$ 的每个有限型点 $s$，$\mathcal{O}_{S, s}$ 都是 G-环；
\item $G$ 是 $S$ 上局部有限型的代数空间。
\end{enumerate}
则 $F$ 是代数空间。
\end{lemma}

\begin{proof}
由引理 \ref{artin-lemma-finite-dimension}，函子 $G$ 满足 [3]。由于
$F \to G$ 是单射，可知 $F$ 也满足 [3]。此外，由于 $F \to G$ 是单射，
给定概形 $U$、$V$ 以及态射 $U \to F$ 和 $V \to F$ 时，有
$U \times_F V = U \times_G V$。因此，$\Delta : F \to F \times F$
可（由概形）表示，因为按假设这对 $G$ 成立。于是命题
\ref{artin-proposition-spaces-diagonal-representable}
适用\footnote{由于集合论条件 [-1] 对 $G$ 成立，它对 $F$ 也成立。细节从略。}。
\end{proof}












\section{代数叠}
\label{artin-section-algebraic-stacks}

\noindent
命题 \ref{artin-proposition-second-diagonal-representable} 是关于代数叠的第一个
主要结果。

\begin{lemma}
\label{artin-lemma-diagonal-representable}
设 $S$ 是局部 Noether 概形。设
$p : \mathcal{X} \to (\Sch/S)_{fppf}$ 是群胚纤维化范畴。假设
\begin{enumerate}
\item $\Delta : \mathcal{X} \to \mathcal{X} \times \mathcal{X}$
可由代数空间表示；
\item $\mathcal{X}$ 满足公理 [-1]、[0]、[1]、[2]、[3]（参见一节
\ref{artin-section-axioms}）；
\item $\mathcal{X}$ 的每个形式对象都有效；
\item $\mathcal{X}$ 满足通用性的开性；并且
\item 对 $S$ 的每个有限型点 $s$，$\mathcal{O}_{S, s}$ 都是 G-环。
\end{enumerate}
则 $\mathcal{X}$ 是代数叠。
\end{lemma}

\begin{proof}
引理 \ref{artin-lemma-get-smooth} 适用于 $\mathcal{X}$。利用这一点，对 $S$ 上
每个有限型域 $k$ 和每个对象同构类
$x_0 \in \Ob(\mathcal{X}_{\Spec(k)})$，选取 $S$ 上有限型的仿射概形
$U_{k, x_0}$ 及光滑态射 $(\Sch/U_{k, x_0})_{fppf} \to \mathcal{X}$，
使得存在剩余域为 $k$ 的有限型点
$u_{k, x_0} \in U_{k, x_0}$，而 $x_0$ 是 $u_{k, x_0}$ 的像。于是
$$
(\Sch/U)_{fppf} \to \mathcal{X},
\quad\text{其中}\quad
U = \coprod\nolimits_{k, x_0} U_{k, x_0}
$$
光滑\footnote{集合论评注：由公理 [-1]，指标集的大小有界，所以这个余积
（同构于）$(\Sch/S)_{fppf}$ 的一个对象；参见《集合》引理
\ref{sets-lemma-what-is-in-it}。}。为完成证明，只须证明此映射是满的；
参见《可表性判据》引理 \ref{criteria-lemma-stacks-etale}
（此处用到公理 [0]）。由《可表性判据》引理
\ref{criteria-lemma-check-property-limit-preserving}，只须证明
$(\Sch/U)_{fppf} \times_\mathcal{X} (\Sch/V)_{fppf} \to (\Sch/V)_{fppf}$
对那些 $y : (\Sch/V)_{fppf} \to \mathcal{X}$ 是满的，其中 $V$ 是
$S$ 上局部有限表示的仿射概形。由假设 (1)，纤维积
$(\Sch/U)_{fppf} \times_\mathcal{X} (\Sch/V)_{fppf}$ 可由代数空间 $W$
表示。于是 $W \to V$ 光滑，故其像是开的。因此，只须
证明 $W \to V$ 的像包含 $V$ 的所有有限型点；参见《态射》引理
\ref{morphisms-lemma-enough-finite-type-points}。设 $v_0 \in V$ 是有限型点。
则 $k = \kappa(v_0)$ 是 $S$ 上的有限型域。
% P11-E0250: the frozen source uses an ungrammatical "Denote x_0 = ... the pullback" construction.
把 $y$ 沿 $v_0$ 的拉回记为 $x_0 = y|_{\Spec(k)}$。于是
$(u_{k, x_0}, v_0)$ 给出态射 $\Spec(k) \to W$，它与 $W \to V$ 的复合
为 $v_0$，即得结论。
\end{proof}

\begin{proposition}
\label{artin-proposition-second-diagonal-representable}
设 $S$ 是局部 Noether 概形。设
$p : \mathcal{X} \to (\Sch/S)_{fppf}$ 是群胚纤维化范畴。假设
\begin{enumerate}
\item $\Delta_\Delta : \mathcal{X} \to
\mathcal{X} \times_{\mathcal{X} \times \mathcal{X}} \mathcal{X}$
可由代数空间表示；
\item $\mathcal{X}$ 满足公理 [-1]、[0]、[1]、[2]、[3]、[4] 和 [5]
（参见一节 \ref{artin-section-axioms}）；
\item 对 $S$ 的每个有限型点 $s$，$\mathcal{O}_{S, s}$ 都是 G-环。
\end{enumerate}
则 $\mathcal{X}$ 是代数叠。
\end{proposition}

\begin{proof}
先证明 $\Delta : \mathcal{X} \to \mathcal{X} \times \mathcal{X}$ 可由
代数空间表示。为此，只须证明
$$
\mathcal{Y} =
\mathcal{X} \times_{\Delta, \mathcal{X} \times \mathcal{X}, y} (\Sch/V)_{fppf}
$$
对任意 $S$ 上局部有限表示的仿射概形 $V$ 和 $\mathcal{X} \times \mathcal{X}$
位于 $V$ 上的对象 $y$，都可由代数空间表示；参见《可表性判据》引理
\ref{criteria-lemma-check-representable-limit-preserving}\footnote{《可表性判据》
引理 \ref{criteria-lemma-check-representable-limit-preserving} 中的集合论条件
成立：表示 $\mathcal{Y}$ 的代数空间 $Y$ 的大小适当地有界。具体地，
$Y \to S$ 将局部有限型，并且 $Y$ 将满足公理 [-1]。细节从略。}。
注意，$\mathcal{Y}$ 以集合胚为纤维（《叠》引理
\ref{stacks-lemma-isom-as-2-fibre-product}）。令
$Y : (\Sch/S)_{fppf}^{opp} \to \textit{Sets}$，
$T \mapsto \Ob(\mathcal{Y}_T)/\cong$ 为同构类函子。我们将应用命题
\ref{artin-proposition-spaces-diagonal-representable}，证明 $Y$ 是代数空间。

\medskip\noindent
注意，
$\Delta_\mathcal{Y} : \mathcal{Y} \to \mathcal{Y} \times \mathcal{Y}$
（因而 $Y \to Y \times Y$ 也一样）由条件 (1) 和《可表性判据》引理
\ref{criteria-lemma-second-diagonal} 可由代数空间表示。由《叠》引理
\ref{stacks-lemma-stack-in-setoids-characterize} 和
\ref{stacks-lemma-2-fibre-product-gives-stack-in-setoids}，$Y$ 是 \'etale
拓扑下的层；换言之，公理 [0] 成立。又由引理
\ref{artin-lemma-fibre-product-limit-preserving}，$Y$ 保极限，即 [1] 成立。
由引理 \ref{artin-lemma-algebraic-stack-RS} 和
\ref{artin-lemma-fibre-product-RS}，$Y$ 具有 (RS)，即公理 [2] 成立。$Y$ 的
公理 [3] 由引理 \ref{artin-lemma-finite-dimension} 和
\ref{artin-lemma-fibre-product-tangent-spaces} 得到。公理 [4] 由引理
\ref{artin-lemma-effective} 和 \ref{artin-lemma-fibre-product-effective} 得到。
$Y$ 的公理 [5] 直接来自 $\Delta_\mathcal{X}$ 的通用性的开性；后者是
$\mathcal{X}$ 的公理 [5] 的一部分。因此，命题
\ref{artin-proposition-spaces-diagonal-representable} 的所有假设都满足，且 $Y$
是代数空间。

\medskip\noindent
此时，由引理 \ref{artin-lemma-diagonal-representable}，$\mathcal{X}$ 是代数叠。
\end{proof}





\section{强 Rim--Schlessinger 条件}
\label{artin-section-RS-star}

\noindent
在本章余下部分，下面这个严格更强的 Rim--Schlessinger 条件将起重要作用。

\begin{definition}
\label{artin-definition-RS-star}
设 $S$ 是概形，$\mathcal{X}$ 是 $(\Sch/S)_{fppf}$ 上的群胚纤维化范畴。
若给定 $S$-代数的纤维积图
$$
\xymatrix{
B' \ar[r] & B \\
A' = A \times_B B' \ar[u] \ar[r] & A \ar[u]
}
$$
，其中 $B' \to B$ 为具有平方零核的满射时，纤维范畴的函子
$$
\mathcal{X}_{\Spec(A')}
\longrightarrow
\mathcal{X}_{\Spec(A)} \times_{\mathcal{X}_{\Spec(B)}} \mathcal{X}_{\Spec(B')}
$$
是范畴等价，则称 $\mathcal{X}$ 满足{\it 条件 (RS*)}。
\end{definition}

\noindent
作如下观察。设 $A \to B \leftarrow B'$ 如定义 \ref{artin-definition-RS-star} 中所述。
\begin{enumerate}
\item 在概形范畴中有
$\Spec(A') = \Spec(A) \amalg_{\Spec(B)} \Spec(B')$；参见《更多态射》引理
\ref{more-morphisms-lemma-pushout-along-thickening}；并且
\item 若 $\mathcal{X}$ 是代数叠，则由引理
\ref{artin-lemma-algebraic-stack-RS-star}，$\mathcal{X}$ 满足 (RS*)。
\end{enumerate}
若 $S$ 局部 Noether，则
\begin{enumerate}
\item[(3)] 若 $A$、$B$、$B'$ 在 $S$ 上有限型，且 $B$ 在 $A$ 上有限，
则 $A'$ 在 $S$ 上有限型\footnote{若 $\Spec(A)$ 映入 $S$ 的某个仿射开集，
这由《更多代数》引理 \ref{more-algebra-lemma-fibre-product-finite-type} 得到。
一般情形使用《更多代数》引理
\ref{more-algebra-lemma-diagram-localize} 得到。}；并且
\item[(4)] 若 $\mathcal{X}$ 满足 (RS*)，则 $\mathcal{X}$ 满足 (RS)，
因为 (RS) 恰好涵盖 (RS*) 中 $A$、$B$、$B'$ 为 Artin 局部环的情形。
\end{enumerate}

\begin{lemma}
\label{artin-lemma-algebraic-stack-RS-star}
设 $\mathcal{X}$ 是基概形 $S$ 上的代数叠。则 $\mathcal{X}$ 满足 (RS*)。
\end{lemma}

\begin{proof}
这由引理 \ref{artin-lemma-pushout} 推出；参见定义
\ref{artin-definition-RS-star} 后的评注。
\end{proof}

\begin{lemma}
\label{artin-lemma-fibre-product-RS-star}
设 $S$ 是概形。设 $p : \mathcal{X} \to \mathcal{Y}$ 和
$q : \mathcal{Z} \to \mathcal{Y}$ 是 $(\Sch/S)_{fppf}$ 上群胚纤维化
范畴的 $1$-态射。若 $\mathcal{X}$、$\mathcal{Y}$ 和 $\mathcal{Z}$
满足 (RS*)，则 $\mathcal{X} \times_\mathcal{Y} \mathcal{Z}$ 也满足 (RS*)。
\end{lemma}

\begin{proof}
证明与引理 \ref{artin-lemma-fibre-product-RS} 的证明完全相同。
\end{proof}








\section{通用性与一般化}
\label{artin-section-generalize-versality}

\noindent
我们证明，对具有 (RS*) 且保极限的叠，通用性在一般化下保持。
初次阅读时，建议跳过本节。

\begin{lemma}
\label{artin-lemma-single-point}
设 $S$ 是局部 Noether 概形，$\mathcal{X}$ 是 $(\Sch/S)_{fppf}$ 上具有
(RS*) 的群胚纤维化范畴。设 $x$ 是 $\mathcal{X}$ 位于 $S$ 上有限型仿射
概形 $U$ 上的对象。设 $u \in U$ 是有限型点，且 $x$ 不在 $u$ 处通用。
则存在 $\mathcal{X}$ 的态射 $x \to y$，它位于 $U \to T$ 之上并满足
\begin{enumerate}
\item 态射 $U \to T$ 是一阶加厚；
\item 有短正合列
$$
0 \to \kappa(u) \to \mathcal{O}_T \to \mathcal{O}_U \to 0
$$
% P11-E0251: condition (3) names the pair with alpha but defines and later uses beta.
\item {\bf 不}存在二元组 $(W, \alpha)$，它由 $u$ 的开邻域
$W \subset T$ 和态射 $\beta : y|_W \to x$ 组成，并使得复合
$$
x|_{U \cap W} \xrightarrow{\text{下列对象的限制： }x \to y}
y|_W \xrightarrow{\beta} x
$$
是典范态射 $x|_{U \cap W} \to x$。
\end{enumerate}
\end{lemma}

\begin{proof}
令 $R = \mathcal{O}_{U, u}^\wedge$。令 $k = \kappa(u)$ 为 $R$ 的剩余域。
令 $\xi$ 为与 $x$ 相联系的 $\mathcal{X}$ 在 $R$ 上的形式对象。由于 $x$
不在 $u$ 处通用，由引理 \ref{artin-lemma-versality-matches}，$\xi$ 不通用。
根据定义 \ref{artin-definition-versal-formal-object} 后的讨论，这意味着可以找到
$\mathcal{X}$ 的态射 $\xi_1 \to x_A \to x_B$，它位于闭浸入
$\Spec(k) \to \Spec(A) \to \Spec(B)$ 之上，其中 $A,B$ 是剩余域为 $k$
的 Artin 局部环；还可以找到 $n \geq 1$ 和交换图
$$
\vcenter{
\xymatrix{
& x_A \ar[ld] \\
\xi_n & \xi_1 \ar[u] \ar[l]
}
}
\quad\text{位于下列对象之上：}\quad
\vcenter{
\xymatrix{
& \Spec(A) \ar[ld] \\
\Spec(R/\mathfrak m^n) & \Spec(k) \ar[u] \ar[l]
}
}
$$
，使得{\bf 不}存在 $m \geq n$ 和交换图
$$
\vcenter{
\xymatrix{
& & x_B \ar[lldd] \\
& & x_A \ar[ld] \ar[u] \\
\xi_m & \xi_n \ar[l] & \xi_1 \ar[u] \ar[l]
}
}
\quad\text{位于下列对象之上：}
\vcenter{
\xymatrix{
& & \Spec(B) \ar[lldd] \\
& & \Spec(A) \ar[ld] \ar[u] \\
\Spec(R/\mathfrak m^m) &
\Spec(R/\mathfrak m^n) \ar[l] &
\Spec(k) \ar[u] \ar[l]
}
}
$$
还可假设 $B \to A$ 是小扩张；换言之，满射 $B \to A$ 的核 $I$ 作为
$A$-模同构于 $k$。这由《形式形变理论》评注
\ref{formal-defos-remark-versal-object} 得到。然后直接定义
$$
T = U \amalg_{\Spec(A)} \Spec(B)
$$
由性质 (RS*)，可找到 $T$ 上的 $y$，其在 $\Spec(B)$ 上的限制为 $x_B$，
在 $U$ 上的限制为 $x$（这给出位于 $U \to T$ 之上的箭头 $x \to y$）。
为完成证明，核验条件 (1)、(2) 和 (3)。

\medskip\noindent
由推出的构造，有交换图
$$
\xymatrix{
0 \ar[r] &
I \ar[r] &
B \ar[r] &
A \ar[r] &
0 \\
0 \ar[r] &
I \ar[r] \ar[u] &
\Gamma(T, \mathcal{O}_T) \ar[r] \ar[u] &
\Gamma(U, \mathcal{O}_U) \ar[r] \ar[u] &
0
}
$$
，其两行正合。这立即证明了 (1) 和 (2)。为完成证明，采用反证法。
假设有 (3) 中的二元组 $(W, \beta)$。由于 $\Spec(B) \to T$ 分解经过
$W$，得到态射
$$
x_B \to y|_W \xrightarrow{\beta} x
$$
由于 $B$ 是剩余域为 $k = \kappa(u)$ 的 Artin 局部环，$x_B \to x$
位于某个态射 $\Spec(B) \to U$ 之上，而该态射对某个 $m \geq n$ 分解经过
$\Spec(\mathcal{O}_{U, u}/\mathfrak m_u^m)$。换言之，$x_B \to x$ 分解经过
$\xi_m$，从而给出映射 $x_B \to \xi_m$。
% P11-E0252: the proof invokes alpha although condition (3) defines the relevant morphism as beta.
条件 (3) 中态射 $\alpha$ 的相容条件转化为下图交换：
$$
\xymatrix{
x_B \ar[d] & x_A \ar[d] \ar[l] \\
\xi_m & \xi_n \ar[l]
}
$$
这就给出所需的矛盾。
\end{proof}

\begin{lemma}
\label{artin-lemma-generalization-versality}
设 $S$ 是局部 Noether 概形，$\mathcal{X}$ 是 $(\Sch/S)_{fppf}$ 上的
群胚纤维化范畴。假设
\begin{enumerate}
\item $\Delta : \mathcal{X} \to \mathcal{X} \times \mathcal{X}$
可由代数空间表示；
\item $\mathcal{X}$ 具有 (RS*)；
\item $\mathcal{X}$ 保极限。
\end{enumerate}
设 $x$ 是 $\mathcal{X}$ 位于 $S$ 上有限型概形 $U$ 上的对象。设
$u \leadsto u_0$ 是 $U$ 的有限型点之间的特殊化，且 $x$ 在 $u_0$ 处通用。
则 $x$ 在 $u$ 处通用。
\end{lemma}

\begin{proof}
缩小 $U$ 后，可假设 $U$ 仿射，且 $U$ 映入 $S$ 的一个仿射开集
$\Spec(\Lambda)$。若 $x$ 不在 $u$ 处通用，则可按引理
\ref{artin-lemma-single-point} 选取位于 $U \to T$ 之上的 $x \to y$。写成
$U = \Spec(R_0)$ 和 $T = \Spec(R)$。态射 $U \to T$ 对应于满环同态
$R \to R_0$，其核是平方零理想。由假设 (3)，存在某个有限型
$\Lambda$-子代数 $R' \subset R$，使得 $y$ 来自
$U' = \Spec(R')$ 上的对象 $x'$。增大 $R'$ 后，可以且确实假设
$R' \to R_0$ 是满射，因而 $U \subset U'$ 是一阶加厚。于是现在有
$$
x \to y \to x'
\text{ 位于下列对象之上： }
U \to T \to U'
$$
由假设 (1)，存在 $S$ 上的代数空间 $Z$ 表示
$$
(\Sch/U)_{fppf} \times_{x, \mathcal{X}, x'} (\Sch/U')_{fppf}
$$
；参见《代数叠》引理 \ref{algebraic-lemma-representable-diagonal}。由
$2$-纤维积的构造，$Z$ 的一个 $V$-值点对应于三元组
$(a, a', \alpha)$，它由态射 $a : V \to U$、$a' : V \to U'$ 和态射
$\alpha : a^*x \to (a')^*x'$ 组成。得到交换图
$$
\xymatrix{
U \ar[rd] \ar[rdd] \ar[rrd] \\
& Z \ar[r]_{p'} \ar[d]^p & U' \ar[d] \\
& U \ar[r] & S
}
$$
% P11-E0253: the frozen source says "comes the isomorphism", omitting "from".
态射 $i : U \to Z$ 来自同构 $x \to x'|_U$。令 $z_0 = i(u_0) \in Z$。
由引理 \ref{artin-lemma-base-change-versal}，$Z \to U'$ 在 $z_0$ 处光滑。
把 $U$ 替换为 $u_0$ 的一个仿射开邻域，把 $U'$ 替换为相应的开集，再把
$Z$ 替换为这些开集分别在 $p$ 和 $p'$ 下的逆像之交后，便达到
$Z \to U'$ 沿 $i(U)$ 光滑的情形。由于 $u \leadsto u_0$，点 $u$ 位于这个
开集中。缩小 $U$ 显然不破坏引理 \ref{artin-lemma-single-point} 的条件 (3)
（概形 $U$、$T$、$U'$ 具有相同的底拓扑空间）。由于 $U \to U'$ 是仿射
概形的一阶加厚，可以选取态射 $i' : U' \to Z$，使得
$p' \circ i' = \text{id}_{U'}$，且其在 $U$ 上的限制为 $i$（《空间的更多态射》
引理 \ref{spaces-more-morphisms-lemma-smooth-formally-smooth}）。把万有态射
$p^*x \to (p')^*x'$ 沿 $i'$ 拉回，得到态射
$$
x' \to x
$$
，它位于 $p \circ i' : U' \to U$ 之上，并使得复合
$$
x \to x' \to x
$$
是恒等态射。回忆，已有位于态射 $T \to U'$ 之上的 $y \to x'$。作复合便
得到态射 $y \to x$，它的存在与引理 \ref{artin-lemma-single-point} 的条件 (3)
矛盾。这个矛盾完成了证明。
\end{proof}






\section{强形式有效性}
\label{artin-section-strong-formal-effectiveness}

\noindent
本节说明，形式对象有效性的一个强版本如何蕴含通用性的开性。
\cite[Theorem 1.1]{Bhatt-Algebraize} 的证明表明，拟紧且拟分离的代数空间
满足评注 \ref{artin-remark-strong-effectiveness} 中讨论的强形式有效性。此外，
这里发展的理论并非空泛：后文将用它证明凝聚层叠和复形模空间的通用性的开性；
参见《Quot》定理
\ref{quot-theorem-coherent-algebraic-general} 和
\ref{quot-theorem-complexes-algebraic}.

\begin{lemma}
\label{artin-lemma-infinite-sequence-pre}
设 $S$ 是局部 Noether 概形，$\mathcal{X}$ 是 $(\Sch/S)_{fppf}$ 上具有
(RS*) 的群胚纤维化范畴。设 $x$ 是 $\mathcal{X}$ 位于 $S$ 上有限型
仿射概形 $U$ 上的对象。
% P11-E0254: the frozen source says "a finite type points".
设 $u_n \in U$、$n \geq 1$ 是有限型点，并满足：(a) 当 $n \not = m$ 时，
不存在特殊化 $u_n \leadsto u_m$；(b) 对所有 $n$，$x$ 都不在 $u_n$ 处通用。
则存在态射
$$
x \to x_1 \to x_2 \to \ldots
\quad\text{于 }\mathcal{X}\text{ 位于下列对象之上： }\quad
U \to U_1 \to U_2 \to \ldots
$$
，它们位于 $S$ 之上并满足
\begin{enumerate}
\item 对每个 $n$，态射 $U \to U_n$ 是一阶加厚；
\item 对每个 $n$，有短正合列
$$
0 \to \kappa(u_n) \to \mathcal{O}_{U_n} \to \mathcal{O}_{U_{n - 1}} \to 0
$$
，其中当 $n = 1$ 时 $U_0 = U$；
\item 对每个 $n$，{\bf 不}存在二元组 $(W, \alpha)$，它由 $u_n$ 的开邻域
$W \subset U_n$ 和态射 $\alpha : x_n|_W \to x$ 组成，并使得复合
$$
x|_{U \cap W} \xrightarrow{\text{下列对象的限制： }x \to x_n}
x_n|_W \xrightarrow{\alpha} x
$$
是典范态射 $x|_{U \cap W} \to x$。
\end{enumerate}
\end{lemma}

\begin{proof}
由于诸点 $u_n$ 之间没有特殊化（特别地，$u_n$ 两两不同），对每个 $n$
都可以找到开集 $U' \subset U$，使得 $u_n \in U'$，而当
$i = 1, \ldots, n - 1$ 时 $u_i \not \in U'$。由引理
\ref{artin-lemma-single-point}，对每个 $n \geq 1$ 都可以找到
$$
x \to y_n
\quad\text{于 }\mathcal{X}\text{ 位于下列对象之上：}\quad
U \to T_n
$$
，使得
\begin{enumerate}
\item 态射 $U \to T_n$ 是一阶加厚；
\item 有短正合列
$$
0 \to \kappa(u_n) \to \mathcal{O}_{T_n} \to \mathcal{O}_U \to 0
$$
% P11-E0255: the auxiliary condition names the pair with alpha but defines the morphism as beta.
\item {\bf 不}存在二元组 $(W, \alpha)$，它由 $u_n$ 的开邻域
$W \subset T_n$ 和态射 $\beta : y_n|_W \to x$ 组成，并使得复合
$$
x|_{U \cap W} \xrightarrow{\text{下列对象的限制： }x \to y_n}
y_n|_W \xrightarrow{\beta} x
$$
是典范态射 $x|_{U \cap W} \to x$。
\end{enumerate}
因此可以归纳定义
$$
U_1 = T_1, \quad
U_{n + 1} = U_n \amalg_U T_{n + 1}
$$
令 $x_1 = y_1$ 并使用 (RS*)，可归纳找到 $U_{n + 1}$ 上的 $x_{n + 1}$，
其在 $U_n$ 上的限制为 $x_n$，在 $T_{n + 1}$ 上的限制为 $y_{n + 1}$。
$U \to U_n$ 的性质 (1) 来自《更多态射》引理
\ref{more-morphisms-lemma-pushout-along-thickening} 中推出的构造。类似地，
$U_n$ 的性质 (2) 由推出的构造从 $T_n$ 的性质 (2) 得到。按上面的讨论，
缩小到 $u_n$ 的一个开邻域 $U'$ 后，$(U_n, x_n)$ 的性质 (3) 直接由
$(T_n, y_n)$ 的性质 (3) 得到，因为 $T_n$ 和 $U_n$ 的相应开子概形同构。
省略一些细节。
\end{proof}

\begin{remark}[强有效性]
\label{artin-remark-strong-effectiveness}
设 $S$ 是局部 Noether 概形，$\mathcal{X}$ 是 $(\Sch/S)_{fppf}$ 上的
群胚纤维化范畴。假设有
\begin{enumerate}
\item 仿射开集 $\Spec(\Lambda) \subset S$；
\item $\Lambda$-代数的逆系统 $(R_n)$，其转移映射为满射，且核局部幂零；
\item $\mathcal{X}$ 的对象系统 $(\xi_n)$，它位于系统
$(\Spec(R_n))$ 之上。
\end{enumerate}
在这种情形下，令 $R = \lim R_n$。若存在 $\mathcal{X}$ 位于
$\Spec(R)$ 上的对象 $\xi$，其在 $\Spec(R_n)$ 上的限制给出系统
$(\xi_n)$，则称 $(\xi_n)$ {\it 有效}。
\end{remark}

\noindent
并非 $S$ 上每个代数叠 $\mathcal{X}$ 都满足如下形式的强有效性公理：
评注 \ref{artin-remark-strong-effectiveness} 中的每个系统 $(\xi_n)$ 都有效。
《例子》一节 \ref{examples-section-non-formal-effectiveness} 给出了一个反例。

\begin{lemma}
\label{artin-lemma-SGE-implies-openness-versality}
设 $S$ 是局部 Noether 概形，$\mathcal{X}$ 是 $(\Sch/S)_{fppf}$ 上的
群胚纤维化范畴。假设
\begin{enumerate}
\item $\Delta : \mathcal{X} \to \mathcal{X} \times \mathcal{X}$
可由代数空间表示；
\item $\mathcal{X}$ 具有 (RS*)；
\item $\mathcal{X}$ 保极限；
\item 评注 \ref{artin-remark-strong-effectiveness} 中的系统 $(\xi_n)$ 都有效，
其中对所有 $m \geq n$，$\Ker(R_m \to R_n)$ 都是平方零理想。
\end{enumerate}
则 $\mathcal{X}$ 满足通用性的开性。
\end{lemma}

\begin{proof}
选取 $S$ 上局部有限型的概形 $U$、$U$ 的有限型点 $u_0$，以及
$\mathcal{X}$ 位于 $U$ 上的对象 $x$，使得 $x$ 在 $u_0$ 处通用。缩小
$U$ 后，可假设 $U$ 仿射，且 $U$ 映入 $S$ 的仿射开集
$\Spec(\Lambda)$。令 $E \subset U$ 为如下有限型点 $u$ 所成的集合：$x$
不在 $u$ 处通用。由引理 \ref{artin-lemma-generalization-versality}，若
$u \in E$，则 $u_0$ 不是 $u$ 的特殊化。若通用性的开性不成立，则 $u_0$
位于 $E$ 的闭包 $\overline{E}$ 中。由《性质》引理
\ref{properties-lemma-countable-dense-subset}，可选取与 $E$ 具有相同闭包的
可数子集 $E' \subset E$。由《性质》引理
\ref{properties-lemma-maximal-points}，可假设 $E'$ 的诸点之间没有特殊化。
按上面的观察，$u_0$ 不是 $E'$ 中任何点的特殊化，所以 $E'$ 必须为
（可数）无限集。因此可写成 $E' = \{u_1, u_2, u_3, \ldots\}$；诸 $u_i$
之间没有特殊化，且 $u_0$ 位于 $E'$ 的闭包中。

\medskip\noindent
按引理 \ref{artin-lemma-infinite-sequence-pre} 选取位于
$U \to U_1 \to U_2 \to \ldots$ 之上的
$x \to x_1 \to x_2 \to \ldots$。写成 $U_n = \Spec(R_n)$ 和
$U = \Spec(R_0)$。令 $R = \lim R_n$。注意，$R \to R_0$ 是满射，
其核是平方零理想。由假设 (4)，得到 $\Spec(R)$ 上的 $\xi$，其到 $R_n$
的基变换为 $x_n$。由假设 (3)，存在某个有限型 $\Lambda$-子代数
$R' \subset R$，使得 $\xi$ 来自 $U' = \Spec(R')$ 上的对象 $\xi'$。
增大 $R'$ 后，可以且确实假设 $R' \to R_0$ 是满射，因而
$U \subset U'$ 是一阶加厚。于是现在有
$$
x \to x_1 \to x_2 \to \ldots \to \xi'
\text{ 位于下列对象之上： }
U \to U_1 \to U_2 \to \ldots \to U'
$$
由假设 (1)，存在 $S$ 上的代数空间 $Z$ 表示
$$
(\Sch/U)_{fppf} \times_{x, \mathcal{X}, \xi'} (\Sch/U')_{fppf}
$$
；参见《代数叠》引理 \ref{algebraic-lemma-representable-diagonal}。由
$2$-纤维积的构造，$Z$ 的一个 $T$-值点对应于三元组
$(a, a', \alpha)$，它由态射 $a : T \to U$、$a' : T \to U'$ 和态射
$\alpha : a^*x \to (a')^*\xi'$ 组成。得到交换图
$$
\xymatrix{
U \ar[rd] \ar[rdd] \ar[rrd] \\
& Z \ar[r]_{p'} \ar[d]^p & U' \ar[d] \\
& U \ar[r] & S
}
$$
% P11-E0256: the frozen source again says "comes the isomorphism", omitting "from".
态射 $i : U \to Z$ 来自同构 $x \to \xi'|_U$。令 $z_0 = i(u_0) \in Z$。
由引理 \ref{artin-lemma-base-change-versal}，$Z \to U'$ 在 $z_0$ 处光滑。
把 $U$ 替换为 $u_0$ 的一个仿射开邻域，把 $U'$ 替换为相应的开集，再把
$Z$ 替换为这些开集分别在 $p$ 和 $p'$ 下的逆像之交后，便达到
$Z \to U'$ 沿 $i(U)$ 光滑的情形。注意，这还涉及把 $u_n$ 替换为一个
子列，即只保留 $u_n$ 位于该开集中的那些指标。此外，缩小 $U$ 显然不破坏
引理 \ref{artin-lemma-infinite-sequence-pre} 的条件 (3)（概形 $U$、$U_n$、$U'$
具有相同的底拓扑空间）。由于 $U \to U'$ 是仿射概形的一阶加厚，可以选取
态射 $i' : U' \to Z$，使得 $p' \circ i' = \text{id}_{U'}$，且其在 $U$
上的限制为 $i$（《空间的更多态射》引理
\ref{spaces-more-morphisms-lemma-smooth-formally-smooth}）。把万有态射
$p^*x \to (p')^*\xi'$ 沿 $i'$ 拉回，得到态射
$$
\xi' \to x
$$
，它位于 $p \circ i' : U' \to U$ 之上，并使得复合
$$
x \to \xi' \to x
$$
是恒等态射。回忆，已有位于态射 $U_1 \to U'$ 之上的
$x_1 \to \xi'$。作复合便得到态射 $x_1 \to x$，它的存在与引理
\ref{artin-lemma-infinite-sequence-pre} 的条件 (3) 矛盾。这个矛盾完成了证明。
\end{proof}

\begin{remark}
\label{artin-remark-trade-openness-versality-diagonal-with-strong-effectiveness}
% P11-E0257: the frozen source uses "an category" instead of "a category".
可以从一个强形式有效性公理推出群胚纤维化范畴的对角态射满足通用性的开性。
设 $S$ 是局部 Noether 概形，$\mathcal{X}$ 是 $(\Sch/S)_{fppf}$ 上的
群胚纤维化范畴。假设
\begin{enumerate}
\item $\Delta_\Delta : \mathcal{X} \to
\mathcal{X} \times_{\mathcal{X} \times \mathcal{X}} \mathcal{X}$
可由代数空间表示；
\item $\mathcal{X}$ 具有 (RS*)；
\item $\mathcal{X}$ 保极限；
\item 给定评注 \ref{artin-remark-strong-effectiveness} 中的 $S$-代数逆系统
$(R_n)$，其中对所有 $m \geq n$，$\Ker(R_m \to R_n)$ 都是平方零理想时，
函子
$$
\mathcal{X}_{\Spec(\lim R_n)} \longrightarrow
\lim_n \mathcal{X}_{\Spec(R_n)}
$$
全忠实。
\end{enumerate}
则 $\Delta : \mathcal{X} \to \mathcal{X} \times \mathcal{X}$ 满足通用性的
开性。这由把引理 \ref{artin-lemma-SGE-implies-openness-versality} 应用于如下形式的
纤维积得到：
$\mathcal{X} \times_{\Delta, \mathcal{X} \times \mathcal{X}, y}
 (\Sch/V)_{fppf}$，其中 $V$ 是任意 $S$ 上局部有限表示的仿射概形，
而 $y$ 是 $\mathcal{X} \times \mathcal{X}$ 位于 $V$ 上的对象。
若以后需要这个结果，我们会把本评注改成引理并给出详细证明。
\end{remark}











\section{无穷小形变}
\label{artin-section-inf}

\noindent
本节讨论一节 \ref{artin-section-tangent-spaces} 中引入的切空间概念的一种推广。
为妥善处理这一问题，借用《形式形变理论》各节
\ref{formal-defos-section-tangent-spaces-functors},
\ref{formal-defos-section-lifts} 和
\ref{formal-defos-section-infinitesimal-automorphisms} 中的一些记号。

\medskip\noindent
设 $S$ 是概形，$\mathcal{X}$ 是 $(\Sch/S)_{fppf}$ 上的群胚纤维化范畴。
给定 $S$-代数同态 $A' \to A$ 和 $\mathcal{X}$ 位于 $\Spec(A)$ 上的对象
$x$，用 $\textit{Lift}(x, A')$ 表示 $x$ 到 $\Spec(A')$ 的提升所成的范畴。
$\textit{Lift}(x, A')$ 的对象是 $\mathcal{X}$ 的态射 $x \to x'$，它位于
$\Spec(A) \to \Spec(A')$ 之上；$\textit{Lift}(x, A')$ 的态射定义为交换图。
$\textit{Lift}(x, A')$ 的同构类集合记为 $\text{Lift}(x, A')$。参见
《形式形变理论》定义 \ref{formal-defos-definition-lifts} 和评注
\ref{formal-defos-remark-omit-arrow}。若 $A' \to A$ 是具有局部幂零核的满射，
则称 $\text{Lift}(x, A')$ 的元素 $x'$ 为 $x$ 的{\it（无穷小）形变}。
在这种情形下，$x'$ 在 $x$ 之上的{\it 无穷小自同构群}是核
$$
\text{Inf}(x'/x) =
\Ker\left(
\text{Aut}_{\mathcal{X}_{\Spec(A')}}(x') \to
\text{Aut}_{\mathcal{X}_{\Spec(A)}}(x)\right)
$$
注意，$\text{Inf}(x'/x)$ 的元素等同于 $\mathit{Aut}_\mathcal{X}(x')$
（相联系的集合纤维化范畴）中 $\text{id}_x$ 在 $\Spec(A')$ 上的提升。
请与《形式形变理论》定义
\ref{formal-defos-definition-relative-infinitesimal-auts} 和《形式形变理论》评注
\ref{formal-defos-remark-infaut-lifting-equalities} 比较。

\medskip\noindent
% P11-E0258: the frozen source omits "by" in the naming construction for A[M].
若 $M$ 是 $A$-模，把 $A[M]$ 记为如下 $A$-代数：其底 $A$-模为
$A \oplus M$，乘法由
$(a, m) \cdot (a', m') = (aa', am' + a'm)$ 给出。当 $M = A$ 时，这是
$A$ 上的对偶数环，按惯例记为 $A[\epsilon]$。存在 $A$-代数同态
$A[M] \to A$。$x$ 到 $\Spec(A[M])$ 的拉回称为 $x$ 到
$\Spec(A[M])$ 的{\it 平凡形变}。

\begin{lemma}
\label{artin-lemma-functoriality}
设 $S$ 是概形。设 $f : \mathcal{X} \to \mathcal{Y}$ 是
$(\Sch/S)_{fppf}$ 上群胚纤维化范畴的 $1$-态射。设
$$
\xymatrix{
B' \ar[r] & B \\
A' \ar[u] \ar[r] & A \ar[u]
}
$$
是 $S$-代数的交换图。设 $x$ 是 $\mathcal{X}$ 位于 $\Spec(A)$ 上的对象，
$y$ 是 $\mathcal{Y}$ 位于 $\Spec(B)$ 上的对象，而
$\phi : f(x)|_{\Spec(B)} \to y$ 是 $\mathcal{Y}$ 位于 $\Spec(B)$ 上的态射。
则存在典范函子
$$
\textit{Lift}(x, A') \longrightarrow \textit{Lift}(y, B')
$$
，它由 $f$ 和 $\phi$ 在提升范畴之间诱导。此构造以显然的方式与群胚
纤维化范畴的 $1$-态射的复合相容。
\end{lemma}

\begin{proof}
这个引理不证自明。
\end{proof}

\noindent
设 $S$ 是基概形，$\mathcal{X}$ 是 $(\Sch/S)_{fppf}$ 上的群胚纤维化范畴。
定义一个范畴，其对象是二元组 $(x, A' \to A)$，其中
\begin{enumerate}
\item $A' \to A$ 是 $S$-代数的满射，其核为平方零理想；
\item $x$ 是 $\mathcal{X}$ 位于 $\Spec(A)$ 上的对象。
\end{enumerate}
态射 $(y, B' \to B) \to (x, A' \to A)$ 由如下数据给出：交换图
$$
\xymatrix{
B' \ar[r] & B \\
A' \ar[u] \ar[r] & A \ar[u]
}
$$
所示的 $S$-代数同态，以及位于 $\Spec(B)$ 之上的态射
$x|_{\Spec(B)} \to y$。称之为{\it 形变情形}范畴。

\begin{lemma}
\label{artin-lemma-properties-lift-RS-star}
设 $S$ 是概形，$\mathcal{X}$ 是 $(\Sch/S)_{fppf}$ 上的群胚纤维化范畴。
假设 $\mathcal{X}$ 满足条件 (RS*)。设 $A$ 是 $S$-代数，而 $x$ 是
$\mathcal{X}$ 位于 $\Spec(A)$ 上的对象。
\begin{enumerate}
\item 存在 $A$-线性函子
$\text{Inf}_x : \text{Mod}_A \to \text{Mod}_A$
，使得给定形变情形 $(x, A' \to A)$ 和提升 $x'$ 时，存在同构
$\text{Inf}_x(I) \to \text{Inf}(x'/x)$，其中 $I = \Ker(A' \to A)$。
\item 存在 $A$-线性函子
$T_x : \text{Mod}_A \to \text{Mod}_A$
，使得
\begin{enumerate}
\item 给定 $\text{Mod}_A$ 中的 $M$，存在双射
$T_x(M) \to \text{Lift}(x, A[M])$；
\item 给定形变情形 $(x, A' \to A)$，存在作用
$$
T_x(I) \times \text{Lift}(x, A') \to \text{Lift}(x, A')
$$
，其中 $I = \Ker(A' \to A)$。若
$\text{Lift}(x, A') \not = \emptyset$，则此作用单传递。
\end{enumerate}
\end{enumerate}
\end{lemma}

\begin{proof}
把 $\text{Inf}_x$ 定义为函子
$$
\text{Mod}_A \longrightarrow \textit{Sets},\quad
M \longrightarrow
\text{Inf}(x'_M/x) = \text{Lift}(\text{id}_x, A[M])
$$
；它把 $M$ 映到 $x$ 到 $\Spec(A[M])$ 的平凡形变 $x'_M$ 的无穷小
自同构群，等价地，映到 $\mathit{Aut}_\mathcal{X}(x'_M)$ 中
$\text{id}_x$ 的提升所成的群。把 $T_x$ 定义为函子
$$
\text{Mod}_A \longrightarrow \textit{Sets},\quad
M \longrightarrow \text{Lift}(x, A[M])
$$
；它取 $x$ 到 $\Spec(A[M])$ 的无穷小形变的同构类。把《形式形变理论》
引理 \ref{formal-defos-lemma-linear-functor} 应用于 $\text{Inf}_x$ 和
$T_x$。该引理可用，因为 (RS*) 表明
$$
\textit{Lift}(x, A[M \times N]) =
\textit{Lift}(x, A[M]) \times \textit{Lift}(x, A[N])
$$
是范畴等式（平凡形变也彼此对应）。

\medskip\noindent
设 $(x, A' \to A)$ 是形变情形。考虑由规则
$g(a_1, a_2) = \overline{a_1} \oplus a_2 - a_1$ 定义的环同态
$g : A' \times_A A' \to A[I]$。存在同构
$$
A' \times_A A' \longrightarrow A' \times_A A[I]
$$
，由 $(a_1, a_2) \mapsto (a_1, g(a_1, a_2))$ 给出。此同构与到第一因子
$A'$ 的投影交换，因而也与到 $A$ 的投影交换。因此，两次应用 (RS*)，
得到范畴等价
\begin{align*}
\textit{Lift}(x, A') \times \textit{Lift}(x, A')
& =
\textit{Lift}(x, A' \times_A A') \\
& =
\textit{Lift}(x, A' \times_A A[I]) \\
& =
\textit{Lift}(x, A') \times \textit{Lift}(x, A[I])
\end{align*}
使用这些映射以及到最后一个积的末因子的投影，得到“差映射”
$$
\text{Inf}(x'/x) \times  \text{Inf}(x'/x)
\longrightarrow
\text{Inf}_x(I)
\quad\text{且}\quad
\text{Lift}(x, A') \times \text{Lift}(x, A')
\longrightarrow
T_x(I)
$$
这些差映射满足传递律
“$(x'_1 - x'_2) + (x'_2 - x'_3) = x'_1 - x'_3$”，因为
$$
\xymatrix{
A' \times_A A' \times_A A'
\ar[rrrrr]_-{(a_1, a_2, a_3) \mapsto (g(a_1, a_2), g(a_2, a_3))}
\ar[rrrrrd]_{(a_1, a_2, a_3) \mapsto g(a_1, a_3)} & & & & &
A[I] \times_A A[I] = A[I \times I] \ar[d]^{+} \\
& & & & & A[I]
}
$$
交换。把上面的一串等价取逆，便得到一个作用；只要
$\text{Inf}(x'/x)$，相应地 $\text{Lift}(x, A')$ 非空，该作用就是自由且
传递的。注意，$\text{Inf}(x'/x)$ 是群，故总是非空。
\end{proof}

\begin{remark}[函子性]
\label{artin-remark-functoriality}
假设和记号如引理 \ref{artin-lemma-properties-lift-RS-star}。设 $A \to B$ 是
环同态，且 $y = x|_{\Spec(B)}$。设 $M \in \text{Mod}_A$、
$N \in \text{Mod}_B$。
% P11-E0259: the frozen source omits "be" in "let M -> N an A-linear map".
再设 $M \to N$ 是 $A$-线性映射。则存在典范映射
$\text{Inf}_x(M) \to \text{Inf}_y(N)$ 和
$T_x(M) \to T_y(N)$，因为存在拉回函子
$$
\textit{Lift}(x, A[M]) \to \textit{Lift}(y, B[N])
$$
，它来自环同态 $A[M] \to B[N]$。类似地，给定形变情形的态射
$(y, B' \to B) \to (x, A' \to A)$，得到拉回函子
$\textit{Lift}(x, A') \to \textit{Lift}(y, B')$。由于 $\text{Inf}_x$
和 $T_x$ 上的作用、加法和标量乘法的构造只使用提升范畴中的态射（参见
《形式形变理论》引理 \ref{formal-defos-lemma-linear-functor} 的证明），
上述构造具有函子性。换言之，得到 $A$-线性映射
$$
\text{Inf}_x(M) \to \text{Inf}_y(N)
\quad\text{且}\quad
T_x(M) \to T_y(N)
$$
，使得下列图
$$
\vcenter{
\xymatrix{
\text{Inf}_y(J) \ar[r] & \text{Inf}(y'/y) \\
\text{Inf}_x(I) \ar[r] \ar[u] & \text{Inf}(x'/x) \ar[u]
}
}
\quad\text{且}\quad
\vcenter{
\xymatrix{
T_y(J) \times \text{Lift}(y, B') \ar[r] & \text{Lift}(y, B') \\
T_x(I) \times \text{Lift}(x, A') \ar[r] \ar[u] & \text{Lift}(x, A') \ar[u]
}
}
$$
交换。这里 $I = \Ker(A' \to A)$、$J = \Ker(B' \to B)$；$x'$ 是 $x$
到 $A'$ 的提升（它不一定存在），而 $y' = x'|_{\Spec(B')}$。
\end{remark}

\begin{remark}[自同构]
\label{artin-remark-automorphisms}
假设和记号如引理 \ref{artin-lemma-properties-lift-RS-star}。设 $x', x''$ 是 $x$
到 $A'$ 的提升。则有复合映射
$$
\text{Inf}(x'/x) \times
\Mor_{\textit{Lift}(x, A')}(x', x'') \times \text{Inf}(x''/x)
\longrightarrow
\Mor_{\textit{Lift}(x, A')}(x', x'').
$$
由于 $\textit{Lift}(x, A')$ 是群胚，若
$\Mor_{\textit{Lift}(x, A')}(x', x'')$ 非空，则这定义了
$\text{Inf}(x'/x)$ 在 $\Mor_{\textit{Lift}(x, A')}(x', x'')$ 上的单传递
左作用，以及 $\text{Inf}(x''/x)$ 的单传递右作用。现在，引理断言
$\text{Inf}(x'/x) = \text{Inf}_x(I) = \text{Inf}(x''/x)$。我们断言，
经这些认同后，上述两个作用一致。事实上，或者 $x' \not \cong x''$，
此时断言显然；或者 $x' \cong x''$，此时可假设 $x'' = x'$，而结论来自
$\text{Inf}(x'/x)$ 的交换性。特别地，得到良定义的作用
$$
\text{Inf}_x(I) \times \Mor_{\textit{Lift}(x, A')}(x', x'')
\longrightarrow
\Mor_{\textit{Lift}(x, A')}(x', x'')
$$
；只要 $\Mor_{\textit{Lift}(x, A')}(x', x'')$ 非空，此作用就单传递。
\end{remark}

\begin{remark}
\label{artin-remark-short-exact-sequence-thickenings}
设 $S$ 是概形，$\mathcal{X}$ 是 $(\Sch/S)_{fppf}$ 上的群胚纤维化范畴，
$A$ 是 $S$-代数。可以如下定义形变情形的{\it 短正合列}
$$
(x, A_1' \to A) \to (x, A_2' \to A) \to (x, A_3' \to A)
$$
：要求核 $I_i = \Ker(A_i' \to A)$ 之间的相应映射
% P11-E0260: the frozen source omits "that" after "we ask".
给出 $A$-模的短正合列
$$
0 \to I_3 \to I_2 \to I_1 \to 0
$$
。注意，在这种情形下，映射 $A_3' \to A_1'$ 分解经过 $A$，故有典范同构
$A_1' = A[I_1]$。
\end{remark}

\begin{lemma}
\label{artin-lemma-ses-inf-and-T}
设 $S$ 是概形。设 $p : \mathcal{X} \to \mathcal{Y}$ 和
$q : \mathcal{Z} \to \mathcal{Y}$ 是 $(\Sch/S)_{fppf}$ 上群胚纤维化
范畴的 $1$-态射。假设 $\mathcal{X}$、$\mathcal{Y}$、$\mathcal{Z}$
满足 (RS*)。设 $A$ 是 $S$-代数，$w$ 是
$\mathcal{W} = \mathcal{X} \times_\mathcal{Y} \mathcal{Z}$ 位于 $A$
之上的对象。
% P11-E0261: the frozen source omits "by" in the three-object naming construction.
把从 $w$ 得到的 $\mathcal{X}, \mathcal{Y}, \mathcal{Z}$ 的对象分别记为
$x, y, z$。对任意 $A$-模 $M$，有 $A$-模的 $6$-项正合列
$$
\xymatrix{
0 \ar[r] &
\text{Inf}_w(M) \ar[r] &
\text{Inf}_x(M) \oplus \text{Inf}_z(M) \ar[r] &
\text{Inf}_y(M) \ar[lld] \\
 &
T_w(M) \ar[r] &
T_x(M) \oplus T_z(M) \ar[r] &
T_y(M)
}
$$
\end{lemma}

\begin{proof}
由引理 \ref{artin-lemma-fibre-product-RS-star}，$\mathcal{W}$ 满足 (RS*)，
因而 $T_w(M)$ 和 $\text{Inf}_w(M)$ 有定义。横向箭头利用引理
\ref{artin-lemma-functoriality} 的函子性定义。

\medskip\noindent
定义“边界”映射 $\delta : \text{Inf}_y(M) \to T_w(M)$。选取同构
$p(x) \to y$ 和 $y \to q(z)$，使得在《范畴》引理
\ref{categories-lemma-2-product-fibred-categories}，更确切地说《范畴》引理
\ref{categories-lemma-2-product-categories-over-C} 对 $2$-纤维积的描述中，
$w = (x, z, p(x) \to y \to q(z))$。令 $x', y', z', w'$ 分别表示
$x, y, z, w$ 在 $A[M]$ 上的平凡形变。通过拉回得到同构
$y' \to p(x')$ 和 $q(z') \to y'$。元素 $\alpha \in \text{Inf}_y(M)$
等同于 $A[M]$ 上的自同构 $\alpha : y' \to y'$，它在 $A$ 上限制为 $y$
的恒等态射。因此，令
$$
\delta(\alpha) =
(x', z', p(x') \to y' \xrightarrow{\alpha} y' \to q(z'))
$$
，便得到 $T_w(M)$ 的对象。由《形式形变理论》引理
\ref{formal-defos-lemma-morphism-linear-functors}，这是 $A$-模的映射。

\medskip\noindent
证明的其余部分与《形式形变理论》引理
\ref{formal-defos-lemma-deformation-categories-fiber-product-morphisms}
的证明完全相同。
\end{proof}

\begin{remark}[与前述切空间的相容性]
\label{artin-remark-compare-deformation-spaces}
设 $S$ 是局部 Noether 概形，$\mathcal{X}$ 是 $(\Sch/S)_{fppf}$ 上的
群胚纤维化范畴。假设 $\mathcal{X}$ 具有 (RS*)。设 $k$ 是 $S$ 上的
有限型域，$x_0$ 是 $\mathcal{X}$ 位于 $\Spec(k)$ 上的对象。则有
$k$-向量空间的等式
$$
T\mathcal{F}_{\mathcal{X}, k, x_0} = T_{x_0}(k)
\quad\text{且}\quad
\text{Inf}(\mathcal{F}_{\mathcal{X}, k, x_0}) =
\text{Inf}_{x_0}(k)
$$
其中等号左边的空间由 (\ref{artin-equation-tangent-space}) 和
(\ref{artin-equation-infinitesimal-automorphisms}) 给出，右边的空间由引理
\ref{artin-lemma-properties-lift-RS-star} 给出。
\end{remark}

\begin{remark}[典范元素]
\label{artin-remark-canonical-element}
假设和记号如引理 \ref{artin-lemma-properties-lift-RS-star}。选取仿射开集
$\Spec(\Lambda) \subset S$，使得 $\Spec(A) \to S$ 对应于环同态
$\Lambda \to A$。考虑环同态
$$
A \longrightarrow A[\Omega_{A/\Lambda}],
\quad
a \longmapsto (a, \text{d}_{A/\Lambda}(a))
$$
把 $x$ 沿相应态射 $\Spec(A[\Omega_{A/\Lambda}]) \to \Spec(A)$ 拉回，
得到 $x$ 在 $A[\Omega_{A/\Lambda}]$ 上的形变 $x_{can}$。称之为
{\it 典范元素}
$$
x_{can} \in T_x(\Omega_{A/\Lambda}) = \text{Lift}(x, A[\Omega_{A/\Lambda}]).
$$
其次，假设 $\Lambda$ 是 Noether 环，且 $\Lambda \to A$ 有限型。令
$k = \kappa(\mathfrak p)$ 为 $U = \Spec(A)$ 的有限型点 $u_0$ 处的剩余域。
令 $x_0 = x|_{u_0}$。由 (RS*) 以及 $A[k] = A \times_k k[k]$，空间
$T_x(k)$ 是形变函子 $\mathcal{F}_{\mathcal{X}, k, x_0}$ 的切空间。借助
$$
T\mathcal{F}_{U, k, u_0} =
\text{Der}_\Lambda(A, k) = \Hom_A(\Omega_{A/\Lambda}, k)
$$
（参见《形式形变理论》例
\ref{formal-defos-example-tangent-space-prorepresentable-functor}）和 $T_x$
的函子性，典范元素给出 $U$ 上的对象 $x$ 所诱导的切空间映射。具体地，
$\theta \in T\mathcal{F}_{U, k, u_0}$ 映到
$T_x(k) = T\mathcal{F}_{\mathcal{X}, k, x_0}$ 中的 $T_x(\theta)(x_{can})$。
\end{remark}

\begin{remark}[典范自同构]
\label{artin-remark-canonical-isomorphism}
设 $S$ 是局部 Noether 概形，$\mathcal{X}$ 是 $(\Sch/S)_{fppf}$ 上的
群胚纤维化范畴。假设 $\mathcal{X}$ 满足条件 (RS*)。设 $A$ 是 $S$-代数，
$\Spec(A) \to S$ 映入某个仿射开集，并设 $x, y$ 是 $\mathcal{X}$ 位于
$\Spec(A)$ 上的对象。再设 $A \to B$ 是环同态，而
$\alpha : x|_{\Spec(B)} \to y|_{\Spec(B)}$ 是 $\mathcal{X}$ 位于
$\Spec(B)$ 上的态射。考虑环同态
$$
B \longrightarrow B[\Omega_{B/A}],
\quad
b \longmapsto (b, \text{d}_{B/A}(b))
$$
把 $\alpha$ 沿相应态射 $\Spec(B[\Omega_{B/A}]) \to \Spec(B)$ 拉回，
得到 $B[\Omega_{B/A}]$ 上 $x$ 与 $y$ 的拉回之间的态射 $\alpha_{can}$。
% P11-E0262: the frozen source uses the noun "pullback" as a verb instead of "pull back".
另一方面，可以沿如下态射拉回 $\alpha$：
$\Spec(B[\Omega_{B/A}]) \to \Spec(B)$，它对应于 $B$ 到
$B[\Omega_{B/A}]$ 第一直和项的嵌入。由评注 \ref{artin-remark-automorphisms}
中的讨论，可以取差
$$
\varphi(x, y, \alpha) = \alpha_{can} - \alpha|_{\Spec(B[\Omega_{B/A}])} \in
\text{Inf}_{x|_{\Spec(B)}}(\Omega_{B/A}).
$$
称之为{\it 典范自同构}。它依赖于所有数据 $A$、$x$、$y$、$A \to B$
和 $\alpha$。
\end{remark}





\section{障碍理论}
\label{artin-section-obstruction-theory}

\noindent
本节说明何谓障碍理论。与无穷小形变空间和无穷小自同构空间不同，
障碍理论是一项附加数据。这一表述受引理
\ref{artin-lemma-properties-lift-RS-star} 和评注 \ref{artin-remark-functoriality}
的结果启发。

\begin{definition}
\label{artin-definition-obstruction-theory}
设 $S$ 是局部 Noether 基概形，$\mathcal{X}$ 是 $(\Sch/S)_{fppf}$ 上的
群胚纤维化范畴。一个{\it 障碍理论}由下列数据给出：
\begin{enumerate}
\item 对每个满足 $\Spec(A) \to S$ 映入某个仿射开集的 $S$-代数 $A$，
以及 $\mathcal{X}$ 位于 $\Spec(A)$ 上的每个对象 $x$，给定一个 $A$-线性函子
$$
\mathcal{O}_x : \text{Mod}_A \to \text{Mod}_A
$$
，称为{\it 障碍模}函子；
% P11-E0263: the frozen source omits the existential verb before the induced map.
\item 对如 (1) 中的 $(x, A)$、环同态 $A \to B$、
$M \in \text{Mod}_A$、$N \in \text{Mod}_B$ 和 $A$-线性映射
$M \to N$，给定诱导的 $A$-线性映射
$\mathcal{O}_x(M) \to \mathcal{O}_y(N)$，其中 $y = x|_{\Spec(B)}$；
\item 对每个形变情形 $(x, A' \to A)$，给定一个{\it 障碍}元素
$o_x(A') \in \mathcal{O}_x(I)$，其中 $I = \Ker(A' \to A)$。
\end{enumerate}
这些数据须满足下列条件：
\begin{enumerate}
\item[(i)] 函子性映射使障碍模成为从三元组 $(x, A, M)$ 的范畴到集合范畴的
函子；
\item[(ii)] 对每个形变情形的态射
$(y, B' \to B) \to (x, A' \to A)$，元素 $o_x(A')$ 映到
$o_y(B')$；
\item[(iii)] 对每个形变情形 $(x, A' \to A)$，有
$$
\text{Lift}(x, A') \not = \emptyset
\Leftrightarrow
o_x(A') = 0
$$
。
\end{enumerate}
\end{definition}

\noindent
最后一个条件解释了上述术语。
% P11-E0264: the frozen source calls O_x(A') the obstruction module although the datum is O_x(I).
模 $\mathcal{O}_x(A')$ 称为{\it 障碍模}，元素 $o_x(A')$ 称为{\it 障碍}。
大多数障碍理论还有附加性质；要使它们有用，还需施加附加条件。
此外，这只是一种示例性定义；例如，定义中可以只考虑 $S$ 上有限型的形变情形。

\medskip\noindent
引入障碍理论的主要理由之一，是检验通用性的开性。下面的引理
\ref{artin-lemma-get-openness-obstruction-theory} 就是这类结果的一个例子。
这样做的最初思想出自 Artin；参见引言中提到的 Artin 论文。例如，Flenner
\cite{Flenner}、Hall \cite{Hall-coherent}、Hall 和 Rydh
\cite{rydh_axioms}、Olsson \cite{olsson_deformation}、Olsson 和 Starr
\cite{olsson-starr} 以及 Lieblich \cite{lieblich-complexes} 的工作都采用了
这一思想（文献次序随意）。此外，对特定的群胚纤维化范畴，作者往往会发展少量
适合手头问题的理论。
% P11-E0265: the frozen source retains a literal future-reference placeholder.
我们稍后将发展这一理论（此处插入将来的文献引用）。

\begin{lemma}
\label{artin-lemma-get-openness-obstruction-theory}
\begin{reference}
这是 \cite[Theorem 4.4]{Hall-coherent}。
\end{reference}
设 $S$ 是局部 Noether 概形，$\mathcal{X}$ 是 $(\Sch/S)_{fppf}$ 上的
群胚纤维化范畴。假设
\begin{enumerate}
\item $\Delta : \mathcal{X} \to \mathcal{X} \times \mathcal{X}$
可由代数空间表示；
\item $\mathcal{X}$ 具有 (RS*)；
\item $\mathcal{X}$ 保极限；
\item 存在一个障碍理论\footnote{分析证明可见，实际上只需对满足
$B = A$ 和 $y = x$ 的映射
$(y, B' \to B) \to (x, A' \to A)$，检验定义
\ref{artin-definition-obstruction-theory} 中障碍类的函子性 (ii)。}；
\item 对 $\mathcal{X}$ 位于 $\Spec(A)$ 上的对象 $x$ 以及 $A$-模
$M_n$（$n \geq 1$），有
\begin{enumerate}
\item $T_x(\prod M_n) = \prod T_x(M_n)$,
\item $\mathcal{O}_x(\prod M_n) \to \prod \mathcal{O}_x(M_n)$
是单射。
\end{enumerate}
\end{enumerate}
则 $\mathcal{X}$ 满足通用性的开性。
\end{lemma}

\begin{proof}
我们通过验证引理 \ref{artin-lemma-SGE-implies-openness-versality} 的条件 (4)
来证明。设 $(\xi_n)$ 和 $(R_n)$ 如评注 \ref{artin-remark-strong-effectiveness}，
且对所有 $m \geq n$，$\Ker(R_m \to R_n)$ 都是平方零理想。令
$A = R_1$、$x = \xi_1$。
% P11-E0266: the frozen source uses the incomplete naming construction "Denote M_n = ...".
记 $M_n = \Ker(R_n \to R_1)$。则 $M_n$ 是 $A$-模。令
$R = \lim R_n$，并令
$$
% P11-E0270: the frozen source omits the comma between r_3 and the ellipsis.
\tilde R = \{(r_1, r_2, r_3 \ldots) \in \prod R_n
\text{ such that all have the same image in }A\}
$$
则 $\tilde R \to A$ 是满射，其核为 $M = \prod M_n$。存在映射
$R \to \tilde R$ 和映射 $\tilde R \to A[M]$；后者由
$(r_1, r_2, r_3, \ldots) \mapsto
(r_1, r_2 - r_1, r_3 - r_2, \ldots)$ 给出。二者合在一起给出形变情形的短正合列
$$
(x, R \to A) \to (x, \tilde R \to A) \to (x, A[M])
$$
；参见评注 \ref{artin-remark-short-exact-sequence-thickenings}。相应的核序列
$0 \to \lim M_n \to M \to M \to 0$，正是计算模系统 $(M_n)$ 的极限的
典范序列。

\medskip\noindent
设 $o_x(\tilde R) \in \mathcal{O}_x(M)$ 为障碍元素。由于已有提升
$\xi_n$，可知 $o_x(\tilde R)$ 在 $\mathcal{O}_x(M_n)$ 中映为零。由假设
(5)(b)，得到 $o_x(\tilde R) = 0$。选取 $x$ 到 $\Spec(\tilde R)$ 的提升
$\tilde \xi$，并令 $\tilde \xi_n$ 为 $\tilde \xi$ 到 $\Spec(R_n)$ 的限制。
% P11-E0271: the frozen source uses singular "There exists" with plural "elements".
由引理 \ref{artin-lemma-properties-lift-RS-star} 的 (2)(b)，存在元素
$t_n \in T_x(M_n)$，使得 $t_n \cdot \tilde \xi_n = \xi_n$。由假设
(5)(a)，可以找到 $t \in T_x(M)$，它在 $T_x(M_n)$ 中映到 $t_n$。
把 $\tilde \xi$ 替换为 $t \cdot \tilde \xi$ 后，对所有 $n$，
$\tilde \xi$ 在 $\Spec(R_n)$ 上限制为 $\xi_n$。特别地，由于
$\xi_{n + 1}$ 在 $\Spec(R_n)$ 上限制为 $\xi_n$，$\tilde \xi$ 到
$\Spec(A[M])$ 的限制 $\overline{\xi}$ 具有如下性质：对所有 $n$，它在
$\Spec(A[M_n])$ 上限制为平凡形变。因此由假设 (5)(a)，
$\overline{\xi}$ 是 $x$ 的平凡形变。把公理 (RS*) 应用于
$R = \tilde R \times_{A[M]} A$，便知 $\tilde \xi$ 是 $x$ 在 $R$ 上某个形变
$\xi$ 的拉回。证明完毕。
\end{proof}

\begin{example}
\label{artin-example-global-sections}
设 $S = \Spec(\Lambda)$，其中 $\Lambda$ 是 Noether 环。设 $W \to S$
为概形态射，$\mathcal{F}$ 为在 $S$ 上平坦的拟凝聚 $\mathcal{O}_W$-模。
考虑函子
$$
F : (\Sch/S)_{fppf}^{opp} \longrightarrow \textit{Sets},
\quad
T/S \longrightarrow H^0(W_T, \mathcal{F}_T)
$$
，其中 $W_T = T \times_S W$ 是基变换，$\mathcal{F}_T$ 是 $\mathcal{F}$
到 $W_T$ 的拉回。若 $T = \Spec(A)$，则记 $W_T = W_A$，其余类同。
设 $\mathcal{X} \to (\Sch/S)_{fppf}$ 为与 $F$ 相联系的群胚纤维化范畴。
则 $\mathcal{X}$ 具有一个障碍理论。具体地，
\begin{enumerate}
\item 给定 $\Lambda$-代数 $A$ 和
$x \in H^0(W_A, \mathcal{F}_A)$，令
$\mathcal{O}_x(M) = H^1(W_A, \mathcal{F}_A \otimes_A M)$,
% P11-E0267: the frozen source places o_x(A') in O_x(A), while the boundary target is O_x(I).
\item 给定形变情形 $(x, A' \to A)$，令
$o_x(A') \in \mathcal{O}_x(A)$ 为 $x$ 在边界映射
$$
H^0(W_A, \mathcal{F}_A) \longrightarrow H^1(W_A, \mathcal{F}_A \otimes_A I)
$$
下的像；该边界映射来自模的短正合列
$$
0 \to \mathcal{F}_A \otimes_A I \to
\mathcal{F}_{A'} \to \mathcal{F}_A \to 0.
$$
\end{enumerate}
我们略去了一些细节，尤其是上述短正合列的构造（它使用 $W_A$ 和
$W_{A'}$ 具有相同底拓扑空间这一事实），以及为何 $\mathcal{F}$ 在 $S$
上的平坦性蕴含上述序列短正合的说明。
\end{example}

\begin{example}[关键例]
\label{artin-example-key}
设 $S = \Spec(\Lambda)$，其中 $\Lambda$ 是 Noether 环。取
$\mathcal{X} = (\Sch/X)_{fppf}$，其中 $X = \Spec(R)$ 且
$R = \Lambda[x_1, \ldots, x_n]/J$。朴素余切复形 $\NL_{R/\Lambda}$
（典范地）同伦等价于
$$
J/J^2
\longrightarrow
\bigoplus\nolimits_{i = 1, \ldots, n} R\text{d}x_i,
$$
参见《代数》引理 \ref{algebra-lemma-NL-homotopy}。考虑形变情形
$(x, A' \to A)$。
% P11-E0268: the frozen source omits "by" in "Denote I the kernel".
把 $A' \to A$ 的核记为 $I$。对象 $x$ 对应于 $(a_1, \ldots, a_n)$，
其中 $a_i \in A$，且对所有 $f \in J$，$f(a_1, \ldots, a_n) = 0$ 在
$A$ 中成立。令
\begin{align*}
% P11-E0269: the frozen source writes O_x(A') although the displayed module depends on I.
\mathcal{O}_x(A')
& =
\Hom_R(J/J^2, I)/\Hom_R(R^{\oplus n}, I) \\
& =
\Ext^1_R(\NL_{R/\Lambda}, I) \\
& =
\Ext^1_A(\NL_{R/\Lambda} \otimes_R A, I).
\end{align*}
选取 $a_i' \in A'$，使其在 $A$ 中的像为 $a_i$。则 $o_x(A')$ 是映射
$J/J^2 \to I$ 的类；该映射把 $f \in J$ 送到
$f(a_1', \ldots, a'_n) \in I$。我们略去 $o_x(A')$ 与选取无关的验证。
显然，若 $o_x(A') = 0$，则该映射可以提升。最后，函子性是直接的。
于是得到一个障碍理论。注意，$o_x(A')$ 还可稍微更典范地描述为 $D(A)$
中的复合
$$
\NL_{R/\Lambda} \to \NL_{A/\Lambda} \to \NL_{A/A'} = I[1]
$$
；最后一个认同参见《代数》引理 \ref{algebra-lemma-NL-surjection}。
\end{example}









\section{朴素障碍理论}
\label{artin-section-naive-obstruction-theory}

\noindent
本节标题指的是：本节将使用朴素余切复形。设 $(x, A' \to A)$ 是给定的
局部 Noether 概形 $S$ 上群胚纤维化范畴的一个形变情形。关键例
\ref{artin-example-key} 表明，任何障碍理论都应当与 $D(A)$ 中以 $A$ 的朴素余切
复形为目标的映射密切相关。展开这一思路，引理
\ref{artin-lemma-characterize-versal} 给出通用性的判别准则，进而在引理
\ref{artin-lemma-openness} 中得到通用性的开性判据。为了进一步形式化这个概念，
我们在定义 \ref{artin-definition-naive-obstruction-theory} 中引入朴素障碍理论的概念。

\medskip\noindent
下文使用《代数》一节 \ref{algebra-section-netherlander} 中定义的朴素余切复形。
特别地，若 $A' \to A$ 是 $\Lambda$-代数的满射，且其平方零核为 $I$，则有映射
$$
\NL_{A'/\Lambda} \to \NL_{A/\Lambda} \to \NL_{A/A'}
$$
，其复合同伦等价于零（参见《代数》评注
\ref{algebra-remark-composition-homotopy-equivalent-to-zero}）。由于这里使用的是
朴素余切复形而非完整余切复形，一般而言这并不构成可区别三角。存在同伦等价
$\NL_{A/A'} \to I[1]$（即把 $I$ 置于次数 $-1$ 所得的复形；参见《代数》
引理 \ref{algebra-lemma-NL-surjection}）。最后，注意存在典范映射
$\NL_{A/\Lambda} \to \Omega_{A/\Lambda}$。

\begin{lemma}
\label{artin-lemma-compute-ext-into-field}
设 $A \to k$ 是环同态，其中 $k$ 是域。设 $E \in D^-(A)$。则
$\Ext^i_A(E, k) = \Hom_k(H^{-i}(E \otimes^\mathbf{L} k), k)$。
\end{lemma}

\begin{proof}
从略。提示：把 $E$ 替换为由自由 $A$-模组成的上有界复形，并计算等式两边。
\end{proof}

\begin{lemma}
\label{artin-lemma-construct-essential-surjection}
设 $\Lambda \to A \to k$ 是 Noether 环的有限型环同态，且对 $A$ 的某个素理想
$\mathfrak p$ 有 $k = \kappa(\mathfrak p)$。
% P11-E0272: the frozen source omits the article in "be morphism of D^-(A)".
设 $\xi : E \to \NL_{A/\Lambda}$ 是 $D^{-}(A)$ 中的态射，并且
$H^{-1}(\xi \otimes^{\mathbf{L}} k)$ 不是满射。则存在 $\Lambda$-代数的满射
$A' \to A$，使得
\begin{enumerate}
\item[(a)] $I = \Ker(A' \to A)$ 平方为零，并且作为 $A$-模同构于 $k$；
\item[(b)] $\Omega_{A'/\Lambda} \otimes k = \Omega_{A/\Lambda} \otimes k$；
\item[(c)] $E \to \NL_{A/A'}$ 为零。
\end{enumerate}
\end{lemma}

\begin{proof}
设 $f \in A$、$f \not \in \mathfrak p$。假设对诱导映射
$E \otimes_A A_f \to \NL_{A_f/\Lambda}$，$A'' \to A_f$ 满足 (a)、(b)、(c)；
参见《代数》引理 \ref{algebra-lemma-localize-NL}。则可令
$A' = A'' \times_{A_f} A$，从而得到一个解。确实，
% P11-E0273: the frozen source identifies the kernel with Ker(A'' -> A), although only A'' -> A_f is given.
$A' \to A$ 显然满足 (a)，因为
$\Ker(A' \to A) = \Ker(A'' \to A) = I$。选取提升 $f$ 的
$f'' \in A''$。则 $A'$ 在 $(f'', f)$ 处的局部化同构于 $A''$（例如由
《代数进阶》引理 \ref{more-algebra-lemma-diagram-localize}）。因此 (b) 和 (c)
对 $A'$ 也显然成立。由此可见，可以把 $A$ 替换为局部化 $A_f$（有限多次）。
特别地，在作此替换后，可以假设 $\mathfrak p$ 是 $A$ 的极大理想；参见
《态射》引理 \ref{morphisms-lemma-point-finite-type}。

\medskip\noindent
选取表示 $A = \Lambda[x_1, \ldots, x_n]/J$。则 $\NL_{A/\Lambda}$
（典范地）同伦等价于
$$
J/J^2
\longrightarrow
\bigoplus\nolimits_{i = 1, \ldots, n} A\text{d}x_i,
$$
参见《代数》引理 \ref{algebra-lemma-NL-homotopy}。必要时作局部化后（使用
Nakayama 引理），可以选取 $J$ 的生成元 $f_1, \ldots, f_m$，使得
$f_j \otimes 1$ 构成 $J/J^2 \otimes_A k$ 的一组基。进一步重新编号后，
可假设 $\text{d}f_1, \ldots, \text{d}f_r$ 的像构成映射
$J/J^2 \otimes k \to \bigoplus k\text{d}x_i$ 的像的一组基，并且
$\text{d}f_{r + 1}, \ldots, \text{d}f_m$ 在
$\bigoplus k\text{d}x_i$ 中映为零。在这些选取下，空间
$$
H^{-1}(\NL_{A/\Lambda} \otimes^{\mathbf{L}}_A k) =
H^{-1}(\NL_{A/\Lambda} \otimes_A k)
$$
以 $f_{r + 1} \otimes 1, \ldots, f_m \otimes 1$ 为一组基。再次换基后，
可以假设 $H^{-1}(\xi \otimes^{\mathbf{L}} k)$ 的像包含在
$f_{r + 1} \otimes 1, \ldots, f_{m - 1} \otimes 1$ 的 $k$-线性张成中。令
$$
A' = \Lambda[x_1, \ldots, x_n]/(f_1, \ldots, f_{m - 1}, \mathfrak pf_m)
$$
按构造，$A' \to A$ 满足 (a)。由于 $\text{d}f_m$ 在
$\bigoplus k\text{d}x_i$ 中映为零，可知 (b) 成立。最后，按构造，诱导映射
$E \to \NL_{A/A'} = I[1]$ 在
$H^{-1}(E \otimes_A^\mathbf{L} k) \to I \otimes_A k$ 上诱导零映射。
由引理 \ref{artin-lemma-compute-ext-into-field}，可知该复合为零。
\end{proof}

\noindent
下一个引理是我们的关键技术结果。

\begin{lemma}
\label{artin-lemma-characterize-versal}
设 $S$ 是局部 Noether 概形，$\mathcal{X}$ 是 $(\Sch/S)_{fppf}$ 上满足
(RS*) 的群胚纤维化范畴。设 $U = \Spec(A)$ 是 $S$ 上有限型仿射概形，
且它映入仿射开集 $\Spec(\Lambda)$。设 $x$ 是 $\mathcal{X}$ 位于 $U$ 上的对象，
$\xi : E \to \NL_{A/\Lambda}$ 是 $D^{-}(A)$ 中的态射。假设
\begin{enumerate}
\item[(i)] 对每个形变情形 $(x, A' \to A)$，$x$ 可提升到 $\Spec(A')$，
当且仅当 $E \to \NL_{A/\Lambda} \to \NL_{A/A'}$ 为零；
\item[(ii)] 存在函子同构
$T_x(-) \to \Ext^0_A(E, -)$
，使得 $E \to \NL_{A/\Lambda} \to \Omega^1_{A/\Lambda}$ 对应于典范元素
（参见评注 \ref{artin-remark-canonical-element}）。
\end{enumerate}
设 $u_0 \in U$ 是剩余域为 $k = \kappa(u_0)$ 的有限型点。考虑下列陈述：
\begin{enumerate}
\item $x$ 在 $u_0$ 处通用；
\item $\xi : E \to \NL_{A/\Lambda}$ 诱导满射
$H^{-1}(E \otimes_A^{\mathbf{L}} k) \to
H^{-1}(\NL_{A/\Lambda} \otimes_A^{\mathbf{L}} k)$
和单射
$H^0(E \otimes_A^{\mathbf{L}} k) \to
H^0(\NL_{A/\Lambda} \otimes_A^{\mathbf{L}} k)$.
\end{enumerate}
则总有 (2) $\Rightarrow$ (1)；若 $u_0$ 是闭点，则还有
(1) $\Rightarrow$ (2)。
\end{lemma}

\begin{proof}
令 $\mathfrak p = \Ker(A \to k)$ 为与 $u_0$ 对应的素理想。

\medskip\noindent
% P11-E0274: the frozen source omits "is" in "Assume that x versal".
假设 $x$ 在 $u_0$ 处通用，且 $u_0$ 是 $U$ 的闭点。若
$H^{-1}(\xi \otimes_A^{\mathbf{L}} k)$ 不是满射，则取引理
\ref{artin-lemma-construct-essential-surjection} 中核为 $I$ 的扩张 $A' \to A$。
因为 $u_0$ 是闭点，$I$ 是有限 $A$-模，因而 $A'$ 是有限型
$\Lambda$-代数（若 $u_0$ 非闭，这一点会失效）。特别地，$A'$ 是 Noether 环。
由 $A'$ 的性质 (c) 和 $\xi$ 的性质 (i)，$x$ 可提升为 $A'$ 上的对象 $x'$。
% P11-E0275: the frozen source omits the article before "kernel".
令 $\mathfrak p' \subset A'$ 为到 $k$ 的满同态之核。由 Artin--Rees 引理
（《代数》引理 \ref{algebra-lemma-Artin-Rees}），存在 $n > 1$，使得
$(\mathfrak p')^n \cap I = 0$。于是
$$
B' = A'/(\mathfrak p')^n \longrightarrow A/\mathfrak p^n = B
$$
是 Artin 局部环的小的本质扩张；参见《形式形变理论》引理
\ref{formal-defos-lemma-essential-surjection}。另一方面，由于 $x$ 在 $u_0$
处通用，而 $x'|_{\Spec(B')}$ 是 $x|_{\Spec(B)}$ 的提升，存在整数
$m \geq n$ 和映射 $q : A/\mathfrak p^m \to B'$，使得复合
$A/\mathfrak p^m \to B' \to B$ 是商映射。由于 $B'$ 的极大理想的
$n$ 次幂为零，$q$ 分解经过 $B$，这与 $B' \to B$ 是本质满射矛盾。
这一矛盾说明 $H^{-1}(\xi \otimes_A^{\mathbf{L}} k)$ 是满射。

\medskip\noindent
% P11-E0276: the frozen source again omits "is" in "Assume that x versal".
假设 $x$ 在 $u_0$ 处通用。由引理 \ref{artin-lemma-compute-ext-into-field}，映射
$H^0(\xi \otimes_A^{\mathbf{L}} k)$ 对偶于映射
$\Ext^0_A(\NL_{A/\Lambda}, k) \to \text{Ext}^0_A(E, k)$。注意
$$
\Ext^0_A(\NL_{A/\Lambda}, k) = \text{Der}_\Lambda(A, k)
\quad\text{且}\quad
T_x(k) = \Ext^0_A(E, k)
$$
条件 (ii) 保证映射
$\Ext^0_A(\NL_{A/\Lambda}, k) \to \text{Ext}^0_A(E, k)$ 把 $U$ 在 $u_0$
处的切向量 $\theta$ 送到 $x_0$ 的相应无穷小形变；参见评注
\ref{artin-remark-canonical-element}。因此，若 $x$ 通用，则此映射为满射；参见
《形式形变理论》引理 \ref{formal-defos-lemma-versal-criterion}。于是
$H^0(\xi \otimes_A^{\mathbf{L}} k)$ 是单射。这就证明了当 $u_0$ 是闭点时的
(1) $\Rightarrow$ (2)。

\medskip\noindent
在证明的余下部分，假设
$H^{-1}(E \otimes_A^\mathbf{L} k) \to
H^{-1}(\NL_{A/\Lambda} \otimes_A^\mathbf{L} k)$ 是满射，并且
% P11-E0277: the frozen source omits "is" before "injective".
$H^0(E \otimes_A^\mathbf{L} k) \to
H^0(\NL_{A/\Lambda} \otimes_A^\mathbf{L} k)$
是单射。令 $R = A_\mathfrak p^\wedge$，并令 $\eta$ 为与
$x|_{\Spec(R)}$ 相联系的 $R$ 上形式对象。切空间上的映射
$d\underline{\eta}$ 是满射，因为它被认同为单射
$H^0(E \otimes_A^{\mathbf{L}} k) \to
H^0(\NL_{A/\Lambda} \otimes_A^{\mathbf{L}} k)$
的对偶（参见上一段）。根据《形式形变理论》引理
\ref{formal-defos-lemma-versal-criterion}，只须证明如下断言：设
$C' \to C$ 是剩余域为 $k$ 的有限型 Artin 局部 $\Lambda$-代数的小扩张，
$R \to C$ 是与剩余域认同相容的 $\Lambda$-代数同态。令
$y = x|_{\Spec(C)}$，并设 $y'$ 是 $y$ 到 $C'$ 的提升。要证：可以把
$\Lambda$-代数同态 $R \to C$ 提升为 $R \to C'$。

\medskip\noindent
注意，只须提升 $\Lambda$-代数同态 $A \to C$。令
$I = \Ker(C' \to C)$。注意，$I$ 是 $1$ 维 $k$-向量空间。提升
$A \to C$ 的障碍 $ob$ 是 $\Ext^1_A(\NL_{A/\Lambda}, I)$ 的元素；参见例
\ref{artin-example-key}。由引理 \ref{artin-lemma-compute-ext-into-field} 和我们的假设，
映射 $\xi$ 诱导单射
$$
\Ext^1_A(\NL_{A/\Lambda}, I)
\longrightarrow
\Ext^1_A(E, I)
$$
按 $ob$ 的构造及 (i)，$ob$ 在 $\Ext^1_A(E, I)$ 中的像就是把 $x$ 提升到
$A \times_C C'$ 的障碍。由 (RS*)，$y/C$ 可提升到 $y'/C'$ 这一事实蕴含
$x$ 可提升到 $A \times_C C'$。因此 $ob = 0$，结论得证。
\end{proof}

\noindent
上面的关键引理使我们能在若干情形下推出通用性的开性。

\begin{lemma}
\label{artin-lemma-openness}
设 $S$ 是局部 Noether 概形，$\mathcal{X}$ 是 $(\Sch/S)_{fppf}$ 上满足
(RS*) 的群胚纤维化范畴。设 $U = \Spec(A)$ 是 $S$ 上有限型仿射概形，
且它映入仿射开集 $\Spec(\Lambda)$。设 $x$ 是 $\mathcal{X}$ 位于 $U$ 上的对象，
$\xi : E \to \NL_{A/\Lambda}$ 是 $D^{-}(A)$ 中的态射。假设
\begin{enumerate}
\item[(i)] 对每个形变情形 $(x, A' \to A)$，$x$ 可提升到 $\Spec(A')$，
当且仅当 $E \to \NL_{A/\Lambda} \to \NL_{A/A'}$ 为零；
\item[(ii)] 存在函子同构
$T_x(-) \to \Ext^0_A(E, -)$
，使得 $E \to \NL_{A/\Lambda} \to \Omega^1_{A/\Lambda}$ 对应于典范元素
（参见评注 \ref{artin-remark-canonical-element}）；
\item[(iii)] $E$ 的上同调群都是有限 $A$-模。
\end{enumerate}
若 $x$ 在闭点 $u_0 \in U$ 处通用，则存在开邻域
$u_0 \in U' \subset U$，使得 $x$ 在 $U'$ 的每个有限型点处通用。
\end{lemma}

\begin{proof}
令 $C$ 为 $\xi$ 的锥，于是有可区别三角
$$
E \to \NL_{A/\Lambda} \to C \to E[1]
$$
属于 $D^{-}(A)$。由引理 \ref{artin-lemma-characterize-versal}，$x$ 在 $u_0$
处通用这一假设蕴含 $H^{-1}(C \otimes^\mathbf{L} k) = 0$。由《代数进阶》
引理 \ref{more-algebra-lemma-cut-complex-in-two}，存在不属于与 $u_0$ 对应素理想的
$f \in A$，使得对任意 $A_f$-模 $M$，都有
$H^{-1}(C \otimes^\mathbf{L}_A M) = 0$。再次使用引理
\ref{artin-lemma-characterize-versal}，可知在开集 $D(f) \subset U$ 的所有有限型点处
都有通用性。
\end{proof}

\noindent
上述技术引理提示了下面的定义。

\begin{definition}
\label{artin-definition-naive-obstruction-theory}
设 $S$ 是局部 Noether 基概形，$\mathcal{X}$ 是 $(\Sch/S)_{fppf}$ 上的
群胚纤维化范畴。假设 $\mathcal{X}$ 满足 (RS*)。一个{\it 朴素障碍理论}
由下列数据给出：
\begin{enumerate}
\item
\label{artin-item-map}
对每个满足 $\Spec(A) \to S$ 映入仿射开集
$\Spec(\Lambda) \subset S$ 的 $S$-代数 $A$，以及 $\mathcal{X}$ 位于
$\Spec(A)$ 上的每个对象 $x$，给定对象 $E_x \in D^-(A)$ 和映射
% P11-E0278: the frozen source writes E rather than the just-defined E_x in xi_x.
$\xi_x : E \to \NL_{A/\Lambda}$；
\item
\label{artin-item-inf}
给定如 (\ref{artin-item-map}) 中的 $(x, A)$，存在函子变换
$$
\text{Inf}_x( - ) \to \Ext^{-1}_A(E_x, -)
\quad\text{且}\quad
T_x(-) \to \Ext^0_A(E_x, -)
$$
\item
\label{artin-item-functoriality}
对如 (\ref{artin-item-map}) 中的 $(x, A)$ 和环同态 $A \to B$，令
$y = x|_{\Spec(B)}$，则 $D(A)$ 中存在函子性映射 $E_x \to E_y$。
\end{enumerate}
这些数据须满足下列条件：
\begin{enumerate}
\item[(i)]
在 (\ref{artin-item-functoriality}) 的情形中，图
$$
\xymatrix{
E_y \ar[r]_{\xi_y} & \NL_{B/\Lambda} \\
E_x \ar[u] \ar[r]^{\xi_x} & \NL_{A/\Lambda} \ar[u]
}
$$
在 $D(A)$ 中交换；
\item[(ii)]
给定如 (\ref{artin-item-map}) 中的 $(x, A)$ 和 $A \to B \to C$，令
$y = x|_{\Spec(B)}$、$z = x|_{\Spec(C)}$，则函子性映射 $E_x \to E_y$
与 $E_y \to E_z$ 的复合就是函子性映射 $E_x \to E_z$；
\item[(iii)]
(\ref{artin-item-inf}) 中的映射是同构，并与函子性映射及评注
\ref{artin-remark-functoriality} 中的映射相容；
\item[(iv)]
复合 $E_x \to \NL_{A/\Lambda} \to \Omega_{A/\Lambda}$ 对应于
$T_x(\Omega_{A/\Lambda}) = \Ext^0(E_x, \Omega_{A/\Lambda})$ 的典范元素；
参见评注 \ref{artin-remark-canonical-element}；
\item[(v)]
给定形变情形 $(x, A' \to A)$，其中 $I = \Ker(A' \to A)$，复合
$E_x \to \NL_{A/\Lambda} \to \NL_{A/A'}$ 在
$$
% P11-E0279: the frozen source's first Hom target is NL_{A/Lambda}, contradicting the composition and following terms.
\Hom_A(E_x, \NL_{A/\Lambda}) = \Ext^0_A(E_x, \NL_{A/A'}) =
\Ext^1_A(E_x, I)
$$
中为零，当且仅当 $x$ 可提升到 $A'$。
\end{enumerate}
\end{definition}

\noindent
特别地可见，令 $\mathcal{O}_x( - ) = \Ext^1_A(E_x, -)$，便得到一节
\ref{artin-section-obstruction-theory} 意义下的障碍理论。

\begin{lemma}
\label{artin-lemma-naive-obstruction-theory-qis}
设 $S$ 和 $\mathcal{X}$ 如定义 \ref{artin-definition-naive-obstruction-theory}，
并为 $\mathcal{X}$ 配备一个朴素障碍理论。设 $A \to B$ 和 $y \to x$
如 (\ref{artin-item-functoriality})。设 $k$ 是作为域的 $B$-代数。则函子性映射
$E_x \to E_y$ 诱导双射
$$
H^i(E_x \otimes_A^{\mathbf{L}} k) \to H^i(E_y \otimes_B^{\mathbf{L}} k)
$$
，其中 $i = 0, 1$。
\end{lemma}

\begin{proof}
令 $z = x|_{\Spec(k)}$。由 (RS*)，
$$
\textit{Lift}(x, A[k]) = \textit{Lift}(z, k[k])
\quad\text{且}\quad
\textit{Lift}(y, B[k]) = \textit{Lift}(z, k[k])
$$
，因为 $A[k] = A \times_k k[k]$ 且 $B[k] = B \times_k k[k]$。因此，
朴素障碍理论的性质蕴含函子性映射 $E_x \to E_y$ 对 $i = -1, 0$ 诱导双射
% P11-E0280: the frozen source inconsistently typesets the target Ext operator with \text{Ext}.
$\Ext^i_A(E_x, k) \to \text{Ext}^i_B(E_y, k)$。由引理
\ref{artin-lemma-compute-ext-into-field}，映射
$H^i(E_x \otimes_A^{\mathbf{L}} k) \to H^i(E_y \otimes_B^{\mathbf{L}} k)$,
$i = 0, 1$ 在对偶向量空间上诱导同构，因而本身也是同构。
\end{proof}

\begin{lemma}
\label{artin-lemma-naive-obstruction-theory-gives-openness}
设 $S$ 是局部 Noether 概形。
% P11-E0281: the frozen source gives the base of the fibred category an opp, unlike the chapter's definitions and conventions.
设 $p : \mathcal{X} \to (\Sch/S)_{fppf}^{opp}$ 是群胚纤维化范畴。
假设 $\mathcal{X}$ 满足 (RS*)，并且 $\mathcal{X}$ 具有朴素障碍理论。若对
$A$ 在 $S$ 上有限型的二元组 $(A, x)$，定义
\ref{artin-definition-naive-obstruction-theory} 中的复形 $E_x$ 具有有限生成上同调群，
则 $\mathcal{X}$ 满足通用性的开性。
\end{lemma}

\begin{proof}
设 $U$ 是 $S$ 上局部有限型概形，$x$ 是 $\mathcal{X}$ 位于 $U$ 上的对象，
$u_0$ 是 $U$ 的有限型点，且 $x$ 在 $u_0$ 处通用。首先可把 $U$ 缩小为
仿射概形，使 $u_0$ 为闭点，并使 $U \to S$ 映入仿射开集
$\Spec(\Lambda)$。记 $U = \Spec(A)$。设
$\xi_x : E_x \to \NL_{A/\Lambda}$ 为障碍映射。此时应用引理
\ref{artin-lemma-openness} 即得结论。
\end{proof}









\section{一个对偶概念}
\label{artin-section-dual}

\noindent
设 $(x, A' \to A)$ 是给定的局部 Noether 概形 $S$ 上群胚纤维化范畴
$\mathcal{X}$ 的形变情形。假设 $\mathcal{X}$ 具有障碍理论；参见定义
\ref{artin-definition-obstruction-theory}。实践中，往往有 $A$-模复形
$K^\bullet$ 和函子同构
$$
\text{Inf}_x(-) \to H^0(K^\bullet \otimes_A^\mathbf{L} -),\quad
T_x(-) \to H^1(K^\bullet \otimes_A^\mathbf{L} -),\quad
\mathcal{O}_x(-) \to H^2(K^\bullet \otimes_A^\mathbf{L} -)
$$
本节稍加形式化这一情形，并说明它如何在若干情形下导出通用性的开性。

\begin{example}
\label{artin-example-global-sections-dual}
设 $\Lambda, S, W, \mathcal{F}$ 如例 \ref{artin-example-global-sections}。
假设 $W \to S$ 是真态射，且 $\mathcal{F}$ 凝聚。由《概形的上同调》评注
\ref{coherent-remark-explain-perfect-direct-image}，存在由有限射影
$\Lambda$-模组成的有限复形 $N^\bullet$，它普遍计算 $\mathcal{F}$ 的上同调。
特别地，例 \ref{artin-example-global-sections} 中的障碍空间为
$\mathcal{O}_x(M) = H^1(N^\bullet \otimes_\Lambda M)$。因此，取
$K^\bullet = N^\bullet \otimes_\Lambda A[-1]$，便有
$\mathcal{O}_x(M) = H^2(K^\bullet \otimes_A^\mathbf{L} M)$。
\end{example}

\begin{situation}
\label{artin-situation-dual}
设 $S$ 是局部 Noether 概形，$\mathcal{X}$ 是 $(\Sch/S)_{fppf}$ 上的
群胚纤维化范畴。假设 $\mathcal{X}$ 具有 (RS*)，因而可以谈论函子
$T_x(-)$；参见引理 \ref{artin-lemma-properties-lift-RS-star}。设
$U = \Spec(A)$ 是 $S$ 上有限型仿射概形，且它映入仿射开集
$\Spec(\Lambda)$。设 $x$ 是 $\mathcal{X}$ 位于 $U$ 上的对象。假设给定
\begin{enumerate}
\item $A$-模复形 $K^\bullet$；
\item 函子变换 $T_x(-) \to H^1(K^\bullet \otimes_A^\mathbf{L} -)$；
\item 对每个核为 $I = \Ker(A' \to A)$ 的形变情形 $(x, A' \to A)$，
给定元素
$o_x(A') \in H^2(K^\bullet \otimes_A^\mathbf{L} I)$
\end{enumerate}
，并满足下列（最小）条件：
\begin{enumerate}
\item[(i)] 变换
$T_x(-) \to H^1(K^\bullet \otimes_A^\mathbf{L} -)$
是同构；
\item[(ii)] 给定形变情形的态射
$(x, A'' \to A) \to (x, A' \to A)$，元素 $o_x(A')$ 经映射
$H^2(K^\bullet \otimes_A^\mathbf{L} I) \to
H^2(K^\bullet \otimes_A^\mathbf{L} I')$
映到元素 $o_x(A'')$，其中 $I' = \Ker(A'' \to A)$；
\item[(iii)] $x$ 可提升为 $\Spec(A')$ 上的对象，当且仅当
$o_x(A') = 0$。
\end{enumerate}
也可以把无穷小自同构纳入其中，但为了得到尽可能强的结果，我们不这样做。
\end{situation}

\noindent
在情形 \ref{artin-situation-dual} 中，
$K^\bullet \otimes_A^\mathbf{L} \NL_{A/\Lambda}$ 将起重要作用。假设给定元素
$\xi \in H^1(K^\bullet \otimes_A^\mathbf{L} \NL_{A/\Lambda})$。于是：
(1) 对任意核为平方零理想 $I$ 的 $\Lambda$-代数满射 $A' \to A$，典范映射
$\NL_{A/\Lambda} \to \NL_{A/A'} = I[1]$ 把 $\xi$ 送到元素
$\xi_{A'} \in H^2(K^\bullet \otimes_A^\mathbf{L} I)$；(2) 映射
$\NL_{A/\Lambda} \to \Omega_{A/\Lambda}$ 把 $\xi$ 送到
$H^1(K^\bullet \otimes_A^\mathbf{L} \Omega_{A/\Lambda})$ 的元素 $\xi_{can}$。

\begin{lemma}
\label{artin-lemma-dual-obstruction}
% P11-E0285: the frozen source leaves "In Situation ..." as a sentence fragment.
在情形 \ref{artin-situation-dual} 中，进一步假设
\begin{enumerate}
\item[(iv)] 给定评注 \ref{artin-remark-short-exact-sequence-thickenings} 中那样的
形变情形短正合列以及提升 $x'_2 \in \text{Lift}(x, A_2')$，则
$o_x(A_3') \in H^2(K^\bullet \otimes_A^\mathbf{L} I_3)$
等于 $\partial\theta$，其中
$\theta \in H^1(K^\bullet \otimes_A^\mathbf{L} I_1)$
是经 $A_1' = A[I_1]$ 及给定映射
$T_x(-) \to H^1(K^\bullet \otimes_A^\mathbf{L} -)$ 与
$x'_2|_{\Spec(A_1')}$ 对应的元素。
\end{enumerate}
在这种情形下，存在元素
$\xi \in H^1(K^\bullet \otimes_A^\mathbf{L} \NL_{A/\Lambda})$
，使得
\begin{enumerate}
\item 对每个形变情形 $(x, A' \to A)$，都有
$\xi_{A'} = o_x(A')$；
\item 经给定变换 $T_x(-) \to H^1(K^\bullet \otimes_A^\mathbf{L} -)$，
$\xi_{can}$ 与评注 \ref{artin-remark-canonical-element} 的典范元素相匹配。
\end{enumerate}
\end{lemma}

\begin{proof}
% P11-E0282: the frozen source uses "a" before alpha instead of "an".
选取核为 $J$ 的 $\alpha : \Lambda[x_1, \ldots, x_n] \to A$。记
$P = \Lambda[x_1, \ldots, x_n]$。在证明的余下部分，使用
$$
\NL(\alpha) = (J/J^2 \longrightarrow \bigoplus A \text{d}x_i)
$$
；《代数》引理 \ref{algebra-lemma-NL-homotopy} 和《代数进阶》引理
\ref{more-algebra-lemma-derived-tor-homotopy} 保证这样做是允许的。考虑元素
$o_x(P/J^2) \in H^2(K^\bullet \otimes_A^\mathbf{L} J/J^2)$ 以及商
$$
C = (P/J^2 \times \bigoplus A \text{d}x_i)/(J/J^2)
$$
，其中 $J/J^2$ 对角嵌入。注意，$C \to A$ 是核为
$\bigoplus A\text{d}x_i$ 的满射。此外，$C \to A$ 有截面 $A \to C$，
它把 $f \in P$ 的类映到推出中 $(f, \text{d}f)$ 的类。
% P11-E0283: the frozen source omits "by" in "denote x_C the pullback".
为后用，把 $x$ 沿相应态射 $\Spec(C) \to \Spec(A)$ 的拉回记为 $x_C$。
因此 $o_x(C) = 0$。由此，$o_x(P/J^2)$ 在
$H^2(K^\bullet \otimes_A^\mathbf{L} \bigoplus A\text{d}x_i)$ 中映为零。
于是存在元素 $\xi \in H^1(K^\bullet \otimes_A^\mathbf{L} \NL(\alpha))$，
它映到 $o_x(P/J^2)$。

\medskip\noindent
注意，对任意形变情形 $(x, A' \to A)$，存在与到 $A$ 的增广相容的
$\Lambda$-代数同态 $P/J^2 \to A'$。因此，按构造及情形
\ref{artin-situation-dual} 的性质 (ii)，元素 $\xi$ 满足引理的第一个性质。

\medskip\noindent
注意，$\xi$ 的选取只在相差
$H^1(K^\bullet \otimes_A^\mathbf{L} \bigoplus A\text{d}x_i)$ 的元素的意义下
良定义。我们将证明，用上述群的一个元素修正 $\xi$ 后，可使引理的第二个断言
成立。令 $C' \subset C$ 为 $P/J^2$ 在 $\Lambda$-代数同态
$P/J^2 \to C$（第一因子的包含）下的像。注意
$\Ker(C' \to A) = \Im(J/J^2 \to \bigoplus A\text{d}x_i)$.
令 $\overline{C} = A[\Omega_{A/\Lambda}]$。映射
$P/J^2 \times \bigoplus A \text{d}x_i \to \overline{C}$,
$(f, \sum f_i \text{d}x_i) \mapsto (f \bmod J, \sum f_i \text{d}x_i)$
分解经过满射 $C \to \overline{C}$。于是
$$
(x, \overline{C} \to A) \to (x, C \to A) \to (x, C' \to A)
$$
是形变情形的短正合列。与之相伴的分裂
$\overline{C} = A[\Omega_{A/\Lambda}]$（来自评注
\ref{artin-remark-short-exact-sequence-thickenings}）等于上面给定的分裂。此外，
截面 $A \to C$ 与映射 $C \to \overline{C}$ 的复合，就是评注
\ref{artin-remark-canonical-element} 中的映射
$(1, \text{d}) : A \to A[\Omega_{A/\Lambda}]$。因此，$x_C$ 限制为
$T_x(\Omega_{A/\Lambda}) = \text{Lift}(x, A[\Omega_{A/\Lambda}])$ 的典范元素
$x_{can}$。由条件 (iv)，$o_x(P/J^2)$ 在
$$
H^1(K^\bullet \otimes_A^\mathbf{L} \Im(J/J^2 \to \bigoplus A\text{d}x_i))
$$
中映到 $\partial x_{can}$。按构造，$\xi$ 映到 $o_x(P/J^2)$。因此，
$x_{can}$ 和 $\xi_{can}$ 在上面显示的群中映到同一元素；由上同调长正合列，
这意味着二者相差 $H^1(K^\bullet \otimes_A^\mathbf{L}
\bigoplus A\text{d}x_i)$ 的一个元素，正如所需。
\end{proof}

\begin{lemma}
\label{artin-lemma-dual-openness}
在情形 \ref{artin-situation-dual} 中，假设引理 \ref{artin-lemma-dual-obstruction} 的
(iv) 成立，并且 $K^\bullet$ 是 $D(A)$ 的完美对象。在这种情形下，若 $x$
在闭点 $u_0 \in U$ 处通用，则存在开邻域 $u_0 \in U' \subset U$，使得
$x$ 在 $U'$ 的每个有限型点处通用。
\end{lemma}

\begin{proof}
可以假设 $K^\bullet$ 是由有限射影 $A$-模组成的有限复形。于是，与
$K^\bullet$ 作导出张量积就等于直接与 $K^\bullet$ 作张量积。令
$E^\bullet$ 为 $K^\bullet$ 的对偶完美复形；参见《代数进阶》引理
\ref{more-algebra-lemma-dual-perfect-complex}。（因此
$E^n = \Hom_A(K^{-n}, A)$，其微分是 $K^\bullet$ 的微分的转置。）令
$E \in D^{-}(A)$ 表示由复形 $E^\bullet[-1]$ 表示的对象。令
$\xi \in H^1(\text{Tot}(K^\bullet \otimes_A \NL_{A/\Lambda}))$ 为引理
\ref{artin-lemma-dual-obstruction} 中构造的元素。
% P11-E0284: the frozen source omits "by" in the naming construction for the corresponding map xi.
把相应映射记为 $\xi : E = E^\bullet[-1] \to \NL_{A/\Lambda}$（同上）。
我们断言，二元组 $(E, \xi)$ 满足引理 \ref{artin-lemma-openness} 的全部假设，
从而证明完毕。

\medskip\noindent
确实，引理 \ref{artin-lemma-openness} 的假设 (i) 来自引理
\ref{artin-lemma-dual-obstruction} 的结论 (1) 以及同上所述的事实
$H^2(K^\bullet \otimes_A^\mathbf{L} -) = \Ext^1(E, -)$。引理
\ref{artin-lemma-openness} 的假设 (ii) 来自引理 \ref{artin-lemma-dual-obstruction} 的结论
(2) 以及同上所述的事实
$H^1(K^\bullet \otimes_A^\mathbf{L} -) = \Ext^0(E, -)$。引理
\ref{artin-lemma-openness} 的假设 (iii) 是显然的。
\end{proof}







\section{Noether 概形上的保极限函子}
\label{artin-section-noetherian}

\noindent
有时，只考虑定义在（局部）Noether 概形全子范畴上的函子或叠会很方便。
本节讨论代数空间的情形。

\medskip\noindent
设 $S$ 是局部 Noether 概形。为了正确对齐各个范畴，我们稍微严格一些；
忽略集合论问题的读者，可以在下文中把所选的概形集合直接替换为所有概形的类。
如《拓扑》评注 \ref{topologies-remark-choice-sites}，选取包含 $S$ 的概形范畴
$\Sch_\alpha$，使得得到大站点 $(\Sch/S)_{Zar}$、
$(\Sch/S)_\etale$, $(\Sch/S)_{smooth}$, $(\Sch/S)_{syntomic}$,
和 $(\Sch/S)_{fppf}$，它们具有相同的底范畴 $\Sch_\alpha/S$。
% P11-E0286: the frozen source omits "by" in the naming construction for Noetherian_alpha.
把
$$
\textit{Noetherian}_\alpha \subset \Sch_\alpha
$$
记为由局部 Noether 概形组成的全子范畴。这确定了全子范畴
$$
\textit{Noetherian}_\alpha/S \subset \Sch_\alpha/S
$$
对 $\tau \in \{Zariski, \etale, smooth, syntomic, fppf\}$，有
\begin{enumerate}
\item 若 $f : X \to Y$ 是 $\Sch_\alpha/S$ 的态射，
$Y$ 属于 $\textit{Noetherian}_\alpha/S$，且 $f$ 局部有限型，则
$X$ 属于 $\textit{Noetherian}_\alpha/S$；
\item 若 $f : X \to Y$ 和 $g : Z \to Y$ 是
$\textit{Noetherian}_\alpha/S$ 的态射，且 $f$ 局部有限型，则
$\textit{Noetherian}_\alpha/S$ 中的纤维积 $X \times_Y Z$ 存在，
并与 $\Sch_\alpha/S$ 中的纤维积一致；
\item 若 $\{X_i \to X\}_{i \in I}$ 是 $(\Sch/S)_\tau$ 的覆盖，且
$X$ 属于 $\textit{Noetherian}_\alpha/S$，则对象 $X_i$ 都属于
$\textit{Noetherian}_\alpha/S$；
\item 在范畴 $\textit{Noetherian}_\alpha/S$ 上赋予
$(\Sch/S)_\tau$ 中对象均属于 $\textit{Noetherian}_\alpha/S$ 的覆盖族，
便得到一个站点，记为 $(\textit{Noetherian}/S)_\tau$；
\item 包含函子
$(\textit{Noetherian}/S)_\tau \to (\Sch/S)_\tau$
是全忠实、连续且余连续的。
\end{enumerate}
由《站点》引理 \ref{sites-lemma-cocontinuous-morphism-topoi} 和
\ref{sites-lemma-when-shriek}，得到拓扑斯态射
$$
g_\tau : \Sh((\textit{Noetherian}/S)_\tau) \longrightarrow \Sh((\Sch/S)_\tau)
$$
；其拉回函子就是沿包含函子
$(\textit{Noetherian}/S)_\tau \to (\Sch/S)_\tau$ 限制层。

\begin{remark}[警告]
\label{artin-remark-no-fibre-products}
站点 $(\textit{Noetherian}/S)_\tau$ 不具有纤维积。因此，处理层时必须谨慎。
例如，连续包含函子
$(\textit{Noetherian}/S)_\tau \to (\Sch/S)_\tau$
并不定义站点态射。当 $\tau = fppf$ 时的一个例子，参见《例》一节
\ref{examples-section-sheaves-locally-Noetherian}。
\end{remark}

\noindent
设 $F : (\textit{Noetherian}/S)_\tau^{opp} \to \textit{Sets}$ 为函子。
若对 $(\textit{Noetherian}/S)_\tau$ 中仿射概形的任意有向极限
$X = \lim X_i$，都有 $F(X) = \colim F(X_i)$，则称 $F$ {\it 保极限}。

\begin{lemma}
\label{artin-lemma-canonical-extension}
设 $\tau \in \{Zariski, \etale, smooth, syntomic, fppf\}$。沿包含函子
$(\textit{Noetherian}/S)_\tau \to (\Sch/S)_\tau$ 作限制，定义下列两范畴之间
的等价：
\begin{enumerate}
\item $(\Sch/S)_\tau$ 上的保极限层所成的范畴；
\item $(\textit{Noetherian}/S)_\tau$ 上的保极限层所成的范畴。
\end{enumerate}
\end{lemma}

\begin{proof}
设 $F : (\textit{Noetherian}/S)_\tau^{opp} \to \textit{Sets}$ 是既保极限
又为层的函子。由《拓扑》引理 \ref{topologies-lemma-extend} 和
\ref{topologies-lemma-extend-sheaf-general}，存在唯一的函子
$F' : (\Sch/S)_\tau^{opp} \to \textit{Sets}$，它保极限、是层，且限制为
$F$。事实上，《拓扑》引理 \ref{topologies-lemma-extend} 中 $F'$ 的构造关于
$F$ 具有函子性，而且该构造是限制函子的拟逆。略去一些细节。
\end{proof}

\begin{lemma}
\label{artin-lemma-representable-limit-preserving}
设 $X$ 是 $(\textit{Noetherian}/S)_\tau$ 的对象。若点函子
$h_X : (\textit{Noetherian}/S)_\tau^{opp} \to \textit{Sets}$ 保极限，
则 $X$ 在 $S$ 上局部有限表示。
\end{lemma}

\begin{proof}
设 $V \subset X$ 是映入仿射开集 $U \subset S$ 的仿射开子概形。可以把
$V = \lim V_i$ 写成 $U$ 上有限表示仿射概形 $V_i$ 的有向极限；参见
《代数》引理 \ref{algebra-lemma-ring-colimit-fp}。按假设，对某个 $i$，箭头
$V \to X$ 分解为 $V \to V_i \to X$。增大 $i$ 后，可以假设
$V_i \to X$ 分解经过 $V$，因为由《极限》引理
\ref{limits-lemma-descend-opens}，$V \subset X$ 在 $V_i$ 中的逆像最终等于
$V_i$。于是，在 $U$ 上概形的范畴中，对某个 $i$，恒等态射 $V \to V$
分解经过 $V_i$。因此 $V \to U$ 有限表示；相应的代数事实是：若 $B$ 是
$A$-代数，且 $\text{id} : B \to B$ 分解经过某个有限表示 $A$-代数，
则 $B$ 在 $A$ 上有限表示（一个不错的练习）。所以 $X$ 在 $S$ 上局部有限表示。
\end{proof}

\noindent
下一个引理有一个适用于可由代数空间表示的变换的版本。

\begin{lemma}
\label{artin-lemma-representable}
设 $\tau \in \{Zariski, \etale, smooth, syntomic, fppf\}$。设
$F', G' : (\Sch/S)_\tau^{opp} \to \textit{Sets}$ 是保极限层，
$a' : F' \to G'$ 是函子变换。
% P11-E0287: the frozen source omits "by" in the naming construction for a : F -> G.
把 $a' : F' \to G'$ 到 $(\textit{Noetherian}/S)_\tau$ 的限制记为
$a : F \to G$。下列条件等价：
\begin{enumerate}
\item $a'$ 可表示（作为函子变换；参见《范畴》定义
\ref{categories-definition-representable-morphism}）；
\item 对 $(\textit{Noetherian}/S)_\tau$ 的每个对象 $V$ 和每个映射
$V \to G$，纤维积
$F \times_G V : (\textit{Noetherian}/S)_\tau^{opp} \to \textit{Sets}$
是可表示函子；
\item 与 (2) 相同，但只要求 $V$ 是映入 $S$ 的某个仿射开集的 $S$ 上有限型
仿射概形。
\end{enumerate}
\end{lemma}

\begin{proof}
假设 (1)。由《空间的极限》引理
\ref{spaces-limits-lemma-locally-finite-presentation-permanence}，变换 $a'$ 保极限
\footnote{即使 $\tau \not = fppf$，这也有意义，因为 $(\Sch/S)_\tau$ 的
底范畴等于 $(\Sch/S)_{fppf}$ 的底范畴，而该陈述与拓扑无关。}。取 (2) 中的
$\xi : V \to G$。
% P11-E0288: the frozen source uses the incomplete naming construction "Denote V' = V but viewed ...".
记 $V' = V$，但把它看作 $(\Sch/S)_\tau$ 的对象。由于 $G$ 是 $G'$ 到
$(\textit{Noetherian}/S)_\tau$ 的限制，$\xi \in G(V)$ 对应于
$\xi' \in G'(V')$。按假设，$V' \times_{\xi', G'} F'$ 可由概形 $U'$ 表示。
与投影 $V' \times_{\xi', G'} F' \to V'$ 对应的概形态射 $U' \to V'$，由
《空间的极限》引理
\ref{spaces-limits-lemma-base-change-locally-finite-presentation} 和《极限》命题
\ref{limits-proposition-characterize-locally-finite-presentation}，是局部有限表示的。
因此 $U'$ 是局部 Noether 概形，从而 $U'$ 同构于
$(\textit{Noetherian}/S)_\tau$ 的某个对象 $U$。于是 $U$ 表示
$F \times_G V$，正如所需。

\medskip\noindent
(2) $\Rightarrow$ (3) 是立即的。假设 (3)。我们来证明 (1)。设 $T$ 是
$(\Sch/S)_\tau$ 的对象，并设 $T \to G'$ 为态射。需要证明函子
$F' \times_{G'} T$ 可由 $T$ 上的某个概形 $X$ 表示。令 $\mathcal{B}$ 为
$T$ 中映入 $S$ 的某个仿射开集的仿射开集所成的集合。这是 $T$ 的拓扑的一组基。
下面证明：对 $W \in \mathcal{B}$，纤维积 $F' \times_{G'} W$ 可由 $W$ 上的
某个概形 $X_W$ 表示。若 $\mathcal{B}$ 中 $W_1 \subset W_2$，则得到同构
$X_{W_1} \to X_{W_2} \times_{W_2} W_1$，因为 $X_{W_1}$ 和
$X_{W_2} \times_{W_2} W_1$ 都表示函子 $F' \times_{G'} W_1$。这些同构是典范的，
并满足《构造》引理 \ref{constructions-lemma-relative-glueing} 中提到的余圈条件。
因此，可以把概形 $X_W$ 黏合成 $T$ 上的概形 $X$。
% P11-E0289: the frozen source has subject-verb disagreement in "Compatibility ... provide".
黏合映射与映射 $X_W \to F'$ 的相容性给出映射 $X \to F'$。所得映射
$X \to F' \times_{G'} T$ 是同构，因为这可在 $T$ 上局部检验（该箭头的源与
目标都是 Zariski 拓扑的层）。

\medskip\noindent
设 $W$ 是映入仿射开集 $U \subset S$ 的仿射概形，并设 $W \to G'$ 为映射。
仍假设 (3)，需要证明 $F' \times_{G'} W$ 可由概形表示。可以把
$W = \lim V'_i$ 写成 $U$ 上有限表示仿射概形 $V'_i$ 的有向极限；参见
《代数》引理 \ref{algebra-lemma-ring-colimit-fp}。
% P11-E0290: the frozen source uses "an Noetherian scheme" instead of "a Noetherian scheme".
由于 $V'_i$ 在某个 Noether 概形上有限型，$V'_i$ 是 Noether 概形。
% P11-E0291: the frozen source uses the incomplete naming construction "Denote V_i = V'_i but viewed ...".
记 $V_i = V'_i$，但把它看作 $(\textit{Noetherian}/S)_\tau$ 的对象。
% P11-E0292: the frozen source omits "we" in "As G' is limit preserving can choose".
由于 $G'$ 保极限，可以选取某个 $i$ 和映射 $V'_i \to G'$，使得
$W \to G'$ 是复合 $W \to V'_i \to G'$。由于 $G$ 是 $G'$ 到
$(\textit{Noetherian}/S)_\tau$ 的限制，态射 $V'_i \to G'$ 就等同于态射
$V_i \to G$（见上文）。由假设 (3)，函子 $F \times_G V_i$ 可由
$(\textit{Noetherian}/S)_\tau$ 的某个对象 $X_i$ 表示。函子
$F \times_G V_i$ 保极限，因为它是 $F' \times_{G'} V'_i$ 的限制，而后一个
函子由《空间的极限》引理
\ref{spaces-limits-lemma-fibre-product-locally-finite-presentation}、$F'$ 和 $G'$
保极限的假设，以及《极限》评注 \ref{limits-remark-limit-preserving}（它说明
$V'_i$ 的点函子保极限）而保极限。由引理
\ref{artin-lemma-representable-limit-preserving}，$X_i$ 在 $S$ 上局部有限表示。
% P11-E0293: the frozen source uses the incomplete naming construction "Denote X'_i = X_i but viewed ...".
记 $X'_i = X_i$，但把它看作 $(\Sch/S)_\tau$ 的对象。于是
$F' \times_{G'} V'_i$ 和点函子 $h_{X'_i}$ 都是
$h_{X_i} : (\textit{Noetherian}/S)_\tau^{opp} \to \textit{Sets}$ 到
$(\Sch/S)_\tau$ 上保极限层的扩张。由引理 \ref{artin-lemma-canonical-extension}
的范畴等价，推出 $X'_i$ 表示 $F' \times_{G'} V'_i$。最后，
$$
F' \times_{G'} W = F' \times_{G'} V'_i \times_{V'_i} W =
X'_i \times_{V'_i} W
$$
可表示，正如所需。
\end{proof}






\section{Noether 情形中的代数空间}
\label{artin-section-algebraic-spaces-noetherian}

\noindent
设 $S$ 是局部 Noether 概形。以
$(\textit{Noetherian}/S)_\etale \subset (\Sch/S)_\etale$ 表示一节
\ref{artin-section-noetherian} 中研究的站点。设
$F : (\textit{Noetherian}/S)_\etale^{opp} \to \textit{Sets}$ 为函子；
换言之，$F$ 是 $(\textit{Noetherian}/S)_\etale$ 上的预层。在此设置下，
一节 \ref{artin-section-axioms-functors} 的所有公理 [-1]、[0]、[1]、[2]、[3]、
[4]、[5] 都有意义。下面逐一回顾，以明确我们的确切含义。

\medskip\noindent
公理 [-1]。这是一个集合论条件；不关心集合论问题的读者可以忽略它。
它对 $F$ 有意义，因为它涉及在 $S$ 上有限型域的谱处计算 $F$，而这些谱都是
$(\textit{Noetherian}/S)_\etale$ 的对象。

\medskip\noindent
% P11-E0295: the frozen source says the sheaf is on the opposite category rather than on the site.
公理 [0]。这是说 $F$ 是 $(\textit{Noetherian}/S)_\etale^{opp}$ 上的层，
即满足 \'etale 覆盖的层条件。

\medskip\noindent
公理 [1]。这是说 $F$ 按一节 \ref{artin-section-noetherian} 的定义保极限：
对 $(\textit{Noetherian}/S)_\etale$ 中仿射概形的任意有向极限
$X = \lim X_i$，都有 $F(X) = \colim F(X_i)$。

\medskip\noindent
公理 [2]。这是说 $F$ 满足 Rim--Schlessinger 条件 (RS)。考察定义
\ref{artin-definition-RS} 中条件 (RS) 的定义及一节
\ref{artin-section-axioms-functors} 中的讨论，可知其含义如下：给定 $S$ 上有限型
概形的任意推出 $Y' = Y \amalg_X X'$，其中 $Y, X, X'$ 是 Artin 局部环的谱，
则
$$
F(Y \amalg_X X') \to F(Y) \times_{F(X)} F(X')
$$
为双射。此条件有意义，因为概形 $X$、$X'$、$Y$、$Y'$ 在 $S$ 上有限型，
故都属于 $(\text{Noetherian}/S)_\etale$。

\medskip\noindent
公理 [3]。这是说每个切空间 $TF_{k, x_0}$ 都是有限维的。此条件有意义，
因为切空间 $TF_{k, x_0}$ 由 $F$ 在 $\Spec(k)$ 和
$\Spec(k[\epsilon])$ 处的取值构造，其中 $k$ 是 $S$ 上有限型域；因此，
这些取值都在范畴 $(\textit{Noetherian}/S)_\etale$ 的对象处计算。

\medskip\noindent
% P11-E0296: the frozen source has both a missing article before axiom and a spurious "the" before every.
公理 [4]。这是说每个形式对象都有效。考察各节
\ref{artin-section-formal-objects} 和 \ref{artin-section-axioms-functors} 中的讨论，
可知这里只涉及在 Noether 概形处计算函子，因此此条件对 $F$ 有意义。

\medskip\noindent
公理 [5]。这是说 $F$ 满足通用性的开性。回忆，其含义如下：给定 $S$ 上局部
有限型概形 $U$、$x \in F(U)$，以及有限型点 $u_0 \in U$，且 $x$ 在
$u_0$ 处通用，则存在开邻域 $u_0 \in U' \subset U$，使得 $x$ 在 $U'$
的每个有限型点处通用。与前面一样，验证此条件只涉及在 Noether 概形处计算函子。

\begin{proposition}
\label{artin-proposition-spaces-diagonal-representable-noetherian}
设 $S$ 是局部 Noether 概形，
$F : (\textit{Noetherian}/S)_\etale^{opp} \to \textit{Sets}$ 是函子。
假设
\begin{enumerate}
\item $\Delta : F \to F \times F$ 可表示（作为函子变换；参见《范畴》定义
\ref{categories-definition-representable-morphism}）；
\item $F$ 满足公理 [-1]、[0]、[1]、[2]、[3]、[4]、[5]（见上文）；
\item 对 $S$ 的所有有限型点 $s$，$\mathcal{O}_{S, s}$ 都是 G-环。
\end{enumerate}
则存在唯一的代数空间
$F' : (\Sch/S)_{fppf}^{opp} \to \textit{Sets}$
，其到 $(\textit{Noetherian}/S)_\etale$ 的限制为 $F$（说明见证明）。
\end{proposition}

\begin{proof}
回忆，站点 $(\Sch/S)_{fppf}$ 和 $(\Sch/S)_\etale$ 具有相同的底范畴；
参见一节 \ref{artin-section-noetherian} 中的讨论。类似地，站点
$(\textit{Noetherian}/S)_\etale$ 和
$(\textit{Noetherian}/S)_{fppf}$ 具有相同的底范畴。
% P11-E0297: the frozen source uses plural "underlying categories" for one shared underlying category.
由公理 [0] 和 [1]，函子 $F$ 是层且保极限。设
$F' : (\Sch/S)_\etale^{opp} \to \textit{Sets}$ 为 $F$ 的唯一扩张，
它是层（关于 \'etale 拓扑）且保极限；参见引理
\ref{artin-lemma-canonical-extension}。于是 $F'$ 满足一节
\ref{artin-section-axioms-functors} 中给出的公理 [0] 和 [1]。由引理
\ref{artin-lemma-representable}，$\Delta' : F' \to F' \times F'$ 可表示（由概形表示）。
另一方面，$F'$ 显然满足一节 \ref{artin-section-axioms-functors} 的公理
[-1]、[2]、[3]、[4]、[5]，因为这些公理都只涉及在
$(\textit{Noetherian}/S)_\etale$ 的对象处计算 $F'$，而我们已假设 $F$
满足相应条件。因此，由命题 \ref{artin-proposition-spaces-diagonal-representable}，
$F'$ 是代数空间。
\end{proof}












\section{Artin 收缩定理}
\label{artin-section-contractions}

\noindent
本节将自由使用形式代数空间的语言；参见《形式空间》一节
\ref{formal-spaces-section-introduction}。Artin 收缩定理是 Artin 论文
\cite{ArtinII} 的两个主要定理之一；第一个定理是他的膨胀定理，我们已在
《形式空间的代数化》一节 \ref{restricted-section-dilatations} 中陈述并证明。

\begin{situation}
\label{artin-situation-contractions}
设 $S$ 是局部 Noether 概形，$X'$ 是 $S$ 上局部有限型的代数空间，
$T' \subset |X'|$ 是闭子集。设 $U' \subset X'$ 为满足
$|U'| = |X'| \setminus T'$ 的开子空间。设 $W$ 是 $S$ 上局部 Noether 的
形式代数空间，且 $W_{red}$ 在 $S$ 上局部有限型。最后，设
$$
g : X'_{/T'} \longrightarrow W
$$
为形式修正；参见《形式空间的代数化》定义
\ref{restricted-definition-formal-modification}。回忆，$X'_{/T'}$
表示 $X'$ 沿 $T'$ 的形式完备化；参见《形式空间》一节
\ref{formal-spaces-section-completion}。
\end{situation}

\noindent
在上述情形中，我们的目标是证明：存在 $S$ 上代数空间的真态射
$f : X' \to X$、闭子集 $T \subset |X|$，以及形式代数空间的同构
$a : X_{/T} \to W$，使得
\begin{enumerate}
\item $T'$ 是 $T$ 在 $|f| : |X'| \to |X|$ 下的逆像；
\item $f : X' \to X$ 将 $U'$ 同构地映到 $X$ 的某个开子空间 $U$；以及
\item $g = a \circ f_{/T}$，其中
$f_{/T} : X'_{/T'} \to X_{/T}$ 是诱导态射。
\end{enumerate}
称 $(f : X' \to X, T, a)$ 为一个{\it 解}。

\medskip\noindent
我们将依照 Artin 的策略：在 $S$ 上局部 Noether 概形的范畴上构造一个函子
$F$，利用命题 \ref{artin-proposition-spaces-diagonal-representable-noetherian}
证明 $F$ 是代数空间，并证明取 $X = F$ 即可。

\begin{remark}
\label{artin-remark-G-rings}
特别地，我们无法证明所求结果对情形 \ref{artin-situation-contractions} 的每个实例
都成立，因为我们将需要假设 $S$ 的局部环都是 G-环。若读者能在一般情形下证明
该结果，或有反例，请通过
\href{mailto:stacks.project@gmail.com}{stacks.project@gmail.com}
告知我们。
\end{remark}

\noindent
在情形 \ref{artin-situation-contractions} 中，设 $V$ 是 $S$ 上局部 Noether 概形。
我们的函子 $F$ 在 $V$ 上的取值由所有如下三元组组成：
$$
(Z, u' : V \setminus Z \to U', \hat x : V_{/Z} \to W)
$$
它们满足下列条件：
\begin{enumerate}
\item $Z \subset V$ 是闭子集；
\item $u' : V \setminus Z \to U'$ 是 $S$ 上的态射；
\item $\hat x : V_{/Z} \to W$ 是 $S$ 上形式代数空间的 adic 态射；
\item $u'$ 与 $\hat x$ 相容（见下文）。
\end{enumerate}
相容性条件如下：拉回形式修正 $g$，得到形式修正
$$
X'_{/T'} \times_{g, W, \hat x} V_{/Z} \longrightarrow V_{/Z}
$$
；参见《形式空间的代数化》引理
\ref{restricted-lemma-base-change-formal-modification}。由膨胀的主定理
（《形式空间的代数化》定理 \ref{restricted-theorem-dilatations}），存在唯一的
代数空间真态射 $V' \to V$，它在 $V \setminus Z$ 上为同构，并且
$V'_{/Z} \to V_{/Z}$ 与上面显示的箭头同构。换言之，对某个态射
$\hat x' : V'_{/Z} \to X'_{/T'}$，有形式代数空间的 Cartesian 图
$$
\xymatrix{
V'_{/Z} \ar[r] \ar[d]_{\hat x'} & V_{/Z} \ar[d]^{\hat x} \\
X'_{/T'} \ar[r]^g & W
}
$$
以后不再特别说明地把 $V \setminus Z$ 视为 $V'$ 的开子空间。相容性条件是：
应存在态射 $x' : V' \to X'$，它在 $V \setminus Z \subset V'$ 和
$V'_{/Z}$ 上分别限制为 $u'$ 和 $\hat x$。
% P11-E0298: the frozen prose says the restriction on V'_{/Z} is \hat x, while the displayed diagram and preceding definition use \hat x'.
换言之，使图
$$
\xymatrix{
V \setminus Z \ar[r] \ar[d]_{u'} &
V' \ar[d]^{x'} &
V'_{/Z} \ar[l] \ar[d]^{\hat x'} \ar[r] &
V_{/Z} \ar[d]^{\hat x} \\
U' \ar[r] &
X' &
X'_{/T'} \ar[r]^g \ar[l] &
W
}
$$
交换。注意，由《形式空间的代数化》引理
\ref{restricted-lemma-faithful}，若态射 $x'$ 存在，则它是唯一的。
我们将用
“{\it $V' \to V$、$\hat x'$ 和 $x'$ 见证 $u'$ 与
$\hat x$ 之间的相容性}”表示这种情形。

\begin{remark}
\label{artin-remark-how-to-think-compatibility}
在情形 \ref{artin-situation-contractions} 中，设 $V$ 是 $S$ 上局部 Noether 概形，
$(Z, u', \hat x)$ 是满足上述 (1)、(2)、(3) 的三元组。我们想说明一种理解
相容性条件 (4) 的方式。由于我们将使用 $S$ 上解析空间的虚构范畴
$\textit{An}_S$ 以及虚构的解析化函子，所述方式在数学上并不精确。
% P11-E0299: the frozen source omits "to" in "we are going use".
$$
\left\{
\begin{matrix}
\text{locally Noetherian formal} \\
\text{algebraic spaces over }S
\end{matrix}
\right\}
\longrightarrow
\textit{An}_S,
\quad\quad
Y \longmapsto Y^{an}
$$
例如，若在 $S = \Spec(k)$ 上有 $Y = \text{Spf}(k[[t]])$，则应把 $Y^{an}$
理解为开单位圆盘。若 $Y = \Spec(k)$，则 $Y^{an}$ 是单点。范畴
$\textit{An}_S$ 应有开浸入和闭浸入，并且我们应能取闭集的开补。给定 $Y$，
态射 $Y_{red} \to Y$ 应诱导闭浸入 $Y_{red}^{an} \to Y^{an}$。
% P11-E0300: the frozen source has the agreement error "should induces".
令 $Y^{rig} = Y^{an} \setminus Y_{red}^{an}$ 为其开补。若 $Y$ 是代数空间且
$Z \subset Y$ 为闭集，则态射 $Y_{/Z} \to Y$ 应诱导开浸入
$Y_{/Z}^{an} \to Y^{an}$，继而诱导开浸入
$$
can : (Y_{/Z})^{rig} \longrightarrow (Y \setminus Z)^{an}
$$
此外，给定局部 Noether 形式代数空间的形式修正 $g : Y' \to Y$，诱导态射
$g^{rig} : (Y')^{rig} \to Y^{rig}$ 应为同构。给定 $\text{An}_S$ 和解析化
函子，可以考虑如下要求：
$$
\xymatrix{
(V_{/Z})^{rig} \ar[rr]_{can} \ar[d]_{(g^{rig})^{-1} \circ \hat x^{an}} & &
(V \setminus Z)^{an} \ar[d]^{(u')^{an}} \\
(X'_{/T'})^{rig} \ar[rr]^{can} & & (X' \setminus T')^{an}
}
$$
交换。
% P11-E0301: the frozen source writes (X'_{T'})^{rig}; the formal completion used throughout is X'_{/T'}.
此要求有意义，因为 $g^{rig} : (X'_{T'})^{rig} \to W^{rig}$ 是同构，且
$U' = X' \setminus T'$。最后，在解析化函子满足某些忠实性假设时，此要求将
等价于上面表述的相容性条件。我们希望这能促使读者把 $u'$ 与 $\hat x$ 的
相容性理解为某些映射相等这一要求，而不是要求某个交换图存在。
\end{remark}

\begin{lemma}
\label{artin-lemma-functor}
在情形 \ref{artin-situation-contractions} 中，规则 $F$ 把 $S$ 上局部 Noether 概形
$V$ 映到满足相容性条件的三元组 $(Z, u', \hat x)$ 所成的集合，并把
$S$ 上局部 Noether 概形的态射 $\varphi : V_2 \to V_1$ 映到映射
% P11-E0302: the frozen source has the duplicated article "a a morphism".
$$
F(\varphi) : F(V_1) \longrightarrow F(V_2)
$$
；后者把 $F(V_1)$ 的元素 $(Z_1, u'_1, \hat x_1)$ 映到 $F(V_2)$ 中由下列
规则给出的 $(Z_2, u'_2, \hat x_2)$：
\begin{enumerate}
\item $Z_2 \subset V_2$ 是 $Z_1$ 在 $\varphi$ 下的逆像；
\item $u'_2$ 是 $u'_1$ 与
$\varphi|_{V_2 \setminus Z_2} : V_2 \setminus Z_2 \to V_1 \setminus Z_1$
的复合；
\item $\hat x_2$ 是 $\hat x_1$ 与
$\varphi_{/Z_2} : V_{2, /Z_2} \to V_{1, /Z_1}$
的复合。
\end{enumerate}
则该规则是反变函子。
\end{lemma}

\begin{proof}
为看出 $u'_2$ 与 $\hat x_2$ 之间的相容性，设
$V'_1 \to V_1$、$\hat x'_1$ 和 $x'_1$ 见证 $u'_1$ 与 $\hat x_1$ 之间的
相容性。令 $V'_2 = V_2 \times_{V_1} V'_1$；令 $\hat x'_2$ 为
$\hat x'_1$ 与 $V'_{2, /Z_2} \to V'_{1, /Z_1}$ 的复合；令 $x'_2$
为 $x'_1$ 与 $V'_2 \to V'_1$ 的复合。于是
$V'_2 \to V_2$、$\hat x'_2$ 和 $x'_2$ 见证 $u'_2$ 与 $\hat x_2$
之间的相容性。略去详细验证。
\end{proof}

\begin{lemma}
\label{artin-lemma-solution}
在情形 \ref{artin-situation-contractions} 中，若存在解
$(f : X' \to X, T, a)$，则在 $S$ 上局部 Noether 概形 $V$ 的范畴上有函子性
双射 $F(V) = \Mor_S(V, X)$。
\end{lemma}

\begin{proof}
设 $V$ 是 $S$ 上局部 Noether 概形，$x : V \to X$ 是 $S$ 上的态射。
于是按如下方式得到 $F(V)$ 中的元素 $(Z, u', \hat x)$：
\begin{enumerate}
\item $Z \subset V$ 是 $T$ 在 $x$ 下的逆像；
\item $u' : V \setminus Z \to U' = U$ 是 $x$ 到 $V \setminus Z$ 的限制；
\item $\hat x : V_{/Z} \to W$ 是 $x_{/Z} : V_{/Z} \to X_{/T}$
与同构 $a : X_{/T} \to W$ 的复合。
\end{enumerate}
此三元组满足相容性条件，因为可以取 $V' = V \times_{x, X} X'$，并取
$\hat x'$ 为投影 $x' : V' \to X'$ 的完备化。
% P11-E0303: the frozen source omits "to be" or "as" after "we can take \hat x'".

\medskip\noindent
反之，设给定 $F(V)$ 的元素 $(Z, u', \hat x)$。我们断言，存在唯一的态射
$x : V \to X$ 与 $u'$ 和 $\hat x$ 相容。事实上，设
$V' \to V$、$\hat x'$ 和 $x'$ 见证 $u'$ 与 $\hat x$ 之间的相容性。
《形式空间的代数化》命题 \ref{restricted-proposition-glue-modification}
恰好给出我们所需的唯一态射 $x : V \to X$：它在 $V_{/Z}$ 上与 $\hat x$
一致，在 $V'$ 上与 $x'$ 一致（从而在 $V \setminus Z$ 上与 $u'$ 一致）。

\medskip\noindent
略去验证上述两个构造在各自定义域之间给出互逆双射。
\end{proof}

\begin{lemma}
\label{artin-lemma-functor-is-solution}
在情形 \ref{artin-situation-contractions} 中，若存在 $S$ 上局部有限型的代数空间
$X$，并且在 $S$ 上局部 Noether 概形 $V$ 的范畴上有函子性双射
$F(V) = \Mor_S(V, X)$，则 $X$ 是一个解。
\end{lemma}

\begin{proof}
我们必须构造真态射 $f : X' \to X$、闭子集 $T \subset |X|$ 以及同构
$a : X_{/T} \to W$，使其具有情形 \ref{artin-situation-contractions} 之后列出的
性质 (1)、(2)、(3)。

\medskip\noindent
本证明中的讨论略显拘泥，因为我们想仔细匹配底范畴。本段说明愿意略去这些细节的
读者如何更大胆地推进。读者可以把函子 $F$ 的定义扩张到 $S$ 上所有局部
Noether 代数空间。这样便可断言 $F$ 与 $X$ 作为这些代数空间范畴上的函子一致，
即 $X$ 表示 $F$。然后考虑 $F(X)$ 中的通用对象 $(T, u', \hat x)$。
读者会发现：对于见证 $u'$ 与 $\hat x$ 之间相容性的三元组
$X'' \to X$、$\hat x'$、$x'$，态射 $x' : X'' \to X'$ 是同构；将 $x'$
取逆便得到 $f : X' \to X$。最后，我们已有 $T \subset |X|$，而读者可以证明
$\hat x$ 是同构，它可作为最后一个要素，即 $a$。
% P11-E0304: the frozen source has "which can served" instead of "which can serve".

\medskip\noindent
以 $h_X(-) = \Mor_S(-, X)$ 表示 $X$ 的点函子到一节
\ref{artin-section-noetherian} 的范畴 $(\textit{Noetherian}/S)_\etale$ 的限制。
% P11-E0305: the frozen source uses the incomplete naming construction "Denote h_X(...) the functor".
由《空间的极限》注 \ref{spaces-limits-remark-limit-preserving}，代数空间 $X$
和 $X'$ 保极限，故其限制 $h_X$ 和 $h_{X'}$ 也保极限。因此，为构造 $f$，
只需构造变换 $h_{X'} \to h_X = F$；参见引理
\ref{artin-lemma-canonical-extension}。于是，设 $V \to S$ 是
$(\textit{Noetherian}/S)_\etale$ 的对象，且
$\tilde x : V \to X'$ 属于 $h_{X'}(V)$。按如下方式得到 $F(V)$ 的元素
$(Z, u', \hat x)$：
\begin{enumerate}
\item $Z \subset V$ 是 $T'$ 在 $\tilde x$ 下的逆像；
\item $u' : V \setminus Z \to U'$ 是 $\tilde x$ 到 $V \setminus Z$ 的限制；
\item $\hat x : V_{/Z} \to W$ 是
$x_{/Z} : V_{/Z} \to X'_{/T'}$ 与 $g : X'_{/T'} \to W$ 的复合。
% P11-E0306: the frozen source uses x_{/Z}, although the morphism introduced here is \tilde x.
\end{enumerate}
此三元组满足相容性条件：首先，我们总能自然得到 $V' \to V$ 和
$\hat x' : V'_{/Z'} \to X'_{/T'}$（参见引理 \ref{artin-lemma-functor} 之前的讨论）。
% P11-E0307: the frozen source introduces Z' only in V'_{/Z'} here; the construction elsewhere uses the inverse image of Z without this notation.
然后只需把 $x' : V' \to X'$ 定义为 $V' \to V$ 与态射
$\tilde x : V \to X'$ 的复合。略去验证此构造确实可行。

\medskip\noindent
若 $\xi : V \to X$ 是 \'etale 态射且 $V$ 是概形，则得到
$\xi = (Z, u', \hat x) \in F(V) = h_X(V) = X(V)$。当然，若
$\varphi : V' \to V$ 还是概形的一个 \'etale 态射，则
$(Z, u', \hat x)$ 到 $F(V')$ 的拉回对应于 $\xi \circ \varphi$。
闭子集 $T \subset |X|$ 就定义为满足如下条件的闭子集：当
$\xi = (Z, u', \hat x)$ 时，$\xi : V \to X$ 把 $T$ 拉回到 $Z$。
% P11-E0308: the frozen source omits terminal punctuation after "pulls T back to Z".

\medskip\noindent
考虑 $S$ 上 Noether 概形 $V$，以及如上对应于 $(Z, u', \hat x)$ 的态射
$\xi : V \to X$。可见，$\xi(V)$ 在集合论意义下包含于 $T$，当且仅当
$V = Z$（作为拓扑空间）。因此，$X_{/T}$ 与 $W$ 作为函子一致。这给出同构
$a : X_{/T} \to W$。（这里略去了一个小细节：对局部 Noether 形式代数空间
$X_{/T}$ 和 $W$，只需检验它们在 $S$ 上局部 Noether 概形处取值后得到同构。）

\medskip\noindent
略去条件 (1)、(2)、(3) 的证明。
\end{proof}

\begin{remark}
\label{artin-remark-diagonal}
在情形 \ref{artin-situation-contractions} 中，设 $V$ 是 $S$ 上局部 Noether 概形。
% P11-E0309: the frozen source leaves "In Situation ..." as a sentence fragment.
对 $i = 1, 2$，设 $(Z_i, u'_i, \hat x_i) \in F(V)$，并设
$V'_i \to V$、$\hat x'_i$ 和 $x'_i$ 见证 $u'_i$ 与 $\hat x_i$ 之间的相容性
（$i = 1, 2$）。

\medskip\noindent
令 $V' = V'_1 \times_V V'_2$。以 $E' \to V'$ 表示下列两个态射的等化子：
$$
V' \to V'_1 \xrightarrow{x'_1} X'
\quad\text{且}\quad
V' \to V'_2 \xrightarrow{x'_2} X'
$$
令 $Z = Z_1 \cap Z_2$。以 $E_W \to V_{/Z}$ 表示下列两个态射的等化子：
$$
V_{/Z} \to V_{/Z_1} \xrightarrow{\hat x_1} W
\quad\text{且}\quad
V_{/Z} \to V_{/Z_2} \xrightarrow{\hat x_2} W
$$
注意，$E' \to V$ 是分离且局部有限型的，而 $E_W$ 是在 $V$ 上分离的局部
Noether 形式代数空间。所涉及各态射之间的相容性表明：
\begin{enumerate}
\item $\Im(E' \to V) \cap (Z_1 \cup Z_2)$ 包含于 $Z = Z_1 \cap Z_2$；
\item 态射 $E' \times_V (V \setminus Z) \to V \setminus Z$
是单态射，并且等于 $u'_1$ 与 $u'_2$ 到
$V \setminus (Z_1 \cup Z_2)$ 的限制的等化子；
\item 态射 $E'_{/Z} \to V_{/Z}$ 通过 $E_W$ 分解，并且图
$$
\xymatrix{
E'_{/Z} \ar[r] \ar[d] & X'_{/T'} \ar[d]^g \\
E_W \ar[r] & W
}
$$
为 Cartesian 图。特别地，态射 $E'_{/Z} \to E_W$ 作为 $g$ 的基变换，是形式修正；
\item $E'$、$(E' \to V)^{-1}Z$ 和 $E'_{/Z} \to E_W$ 构成情形
\ref{artin-situation-contractions} 中的三元组，其基概形为局部 Noether 概形 $V$；
% P11-E0310: the frozen source uses "E', ..., and ... is a triple"; "form a triple" is grammatically clearer.
\item 给定局部 Noether 概形的态射 $\varphi : A \to V$，下列条件等价：
\begin{enumerate}
\item $(Z_1, u'_1, \hat x_1)$ 和 $(Z_2, u'_2, \hat x_2)$
限制为 $F(A)$ 的同一个元素；
\item $A \setminus \varphi^{-1}(Z) \to V \setminus Z$
通过 $E' \times_V (V \setminus Z)$ 分解，且
$A_{/\varphi^{-1}(Z)} \to V_{/Z}$ 通过 $E_W$ 分解。
\end{enumerate}
\end{enumerate}
由引理 \ref{artin-lemma-solution} 和 \ref{artin-lemma-functor-is-solution} 可知：
若 (4) 中的三元组有一个解 $E \to V$，则在 $V$ 上局部 Noether 概形的范畴上，
$E$ 表示 $F \times_{\Delta, F \times F} V$。
\end{remark}

\begin{lemma}
\label{artin-lemma-closed-immersion}
在情形 \ref{artin-situation-contractions} 中，假设给定闭子集 $Z \subset S$，使得
\begin{enumerate}
\item $Z$ 在 $X'$ 中的逆像是 $T'$；
\item $U' \to S \setminus Z$ 是闭浸入；
\item $W \to S_{/Z}$ 是闭浸入。
\end{enumerate}
则存在解 $(f : X' \to X, T, a)$，而且 $X \to S$ 是闭浸入。
\end{lemma}

\begin{proof}
假设有闭子概形 $X \subset S$，使得
$X \cap (S \setminus Z) = U'$ 且 $X_{/Z} = W$。则 $X$ 表示函子 $F$
（略去一些细节），因而是一个解。寻找 $X$ 显然是 $S$ 上的局部问题。
这样便约化到下一段讨论的情形。

\medskip\noindent
假设 $S = \Spec(A)$ 为仿射概形。设 $I \subset A$ 是定义 $Z$ 的根理想，
写作 $I = (f_1, \ldots, f_r)$。由假设，给定
\begin{enumerate}
\item 闭浸入 $U' \to S \setminus Z$ 确定理想
$J_i \subset A[1/f_i]$，使得 $J_i$ 和 $J_j$ 在 $A[1/f_if_j]$ 中生成同一理想；
\item 闭浸入 $W \to S_{/Z}$ 是映射
$\text{Spf}(A^\wedge/J') \to \text{Spf}(A^\wedge)$，其中
$J' \subset A^\wedge$ 是理想，而 $A^\wedge$ 是 $A$ 的 $I$-adic 完备化。
\end{enumerate}
为完成证明，需要找到理想 $J \subset A$，使得 $J_i = J[1/f_i]$ 且
$J' = JA^\wedge$。由《代数续篇》命题
\ref{more-algebra-proposition-equivalence}，只需证明对所有 $i$，
$J_i$ 和 $J'$ 在 $A^\wedge[1/f_i]$ 中生成同一理想。

\medskip\noindent
回忆，$A' = H^0(X', \mathcal{O})$ 是有限 $A$-代数，且其构造与平坦基变换交换
（《空间的上同调》引理
\ref{spaces-cohomology-lemma-proper-over-affine-cohomology-finite} 和
\ref{spaces-cohomology-lemma-flat-base-change-cohomology}）。令
$J'' = \Ker(A \to A')$\footnote{与读者可能预期的不同，理想 $J$ 和 $J''$
一般并不相等。}。
% P11-E0311: the frozen source uses the incomplete naming construction "Denote J'' = ...".
对 $A[1/f_i]$ 的谱作基变换便得到 $J_i = J''A[1/f_i]$。注意，有交换图
$$
\xymatrix{
X' \ar[d] &
X'_{/T'} \times_{S_{/Z}} \text{Spf}(A^\wedge) \ar[l] \ar[d] &
X'_{/T'} \times_W \text{Spf}(A^\wedge/J') \ar@{=}[l] \ar[d] \\
\Spec(A) &
\text{Spf}(A^\wedge) \ar[l] &
\text{Spf}(A^\wedge/J') \ar[l]
}
$$
中间的竖直箭头是左侧竖直箭头沿显然闭子集的完备化。由形式函数定理，有
$$
(A')^\wedge = \Gamma(X' \times_S \Spec(A^\wedge), \mathcal{O}) =
\lim H^0(X' \times_S \Spec(A/I^n), \mathcal{O})
$$
；参见《空间的上同调》定理
\ref{spaces-cohomology-theorem-formal-functions}。由该图可知，$J'$ 在
$(A')^\wedge$ 中映为零。因此 $J' \subset J'' A^\wedge$。考虑箭头
$$
X'_{/T'} \to
\text{Spf}(A^\wedge/J''A^\wedge) \to
\text{Spf}(A^\wedge/J') = W
$$
已知复合 $g$ 是形式修正（特别地，它是 rig-\'etale 且 rig-满射的），而第二个
箭头是闭浸入（特别地，它是 adic 单态射）。因此
$X'_{/T'} \to \text{Spf}(A^\wedge/J''A^\wedge)$ 是 rig-满射且
rig-\'etale 的；参见《形式空间的代数化》引理
\ref{restricted-lemma-rig-surjective-alternative-permanence} 和
\ref{restricted-lemma-rig-etale-alternative-permanence}。
应用《形式空间的代数化》引理
\ref{restricted-lemma-rig-etale-descent} 和
\ref{restricted-lemma-permanence-rig-surjective}
可知，$\text{Spf}(A^\wedge/J''A^\wedge) \to W$ 是 rig-\'etale 且 rig-满射的。
由《形式空间的代数化》引理
\ref{restricted-lemma-closed-immersion-rig-smooth-rig-surjective}
可知，对某个 $n > 0$ 有 $I^n J'' A^\wedge \subset J'$。从而
$J'' A^\wedge[1/f_i] = J' A^\wedge[1/f_i]$，于是对所有 $i$ 都有
$J_i A^\wedge[1/f_i] = J' A^\wedge[1/f_i]$，正如所需。
\end{proof}

\begin{lemma}
\label{artin-lemma-diagonal-contractions}
在情形 \ref{artin-situation-contractions} 中，假设 $X' \to S$ 和 $W \to S$ 分离。
则对角态射 $\Delta : F \to F \times F$ 可由闭浸入表示。
\end{lemma}

\begin{proof}
结合引理 \ref{artin-lemma-closed-immersion} 与注 \ref{artin-remark-diagonal} 中的讨论。
\end{proof}

\begin{lemma}
\label{artin-lemma-sheaf}
在情形 \ref{artin-situation-contractions} 中，函子 $F$ 对 $S$ 上局部 Noether 概形的
所有 \'etale 覆盖都满足层条件。
\end{lemma}

\begin{proof}
略。提示：态射可以在 \'etale 局部定义。
\end{proof}

\begin{lemma}
\label{artin-lemma-limit-preserving}
在情形 \ref{artin-situation-contractions} 中，函子 $F$ 保极限：对 $S$ 上 Noether
仿射概形的任意有向极限 $V = \lim V_\lambda$，都有
$F(V) = \colim F(V_\lambda)$。
\end{lemma}

\begin{proof}
这是一个长得离谱的证明。其中很大一部分由关于极限和 \'etale 局部化的标准论证
组成。我们建议读者直接跳到证明的最后部分，那里才会出现有趣的内容。

\medskip\noindent
设 $V = \lim_{\lambda \in \Lambda} V_i$ 是 $S$ 上概形的有向极限，其中 $V$
和 $V_\lambda$ 都是 Noether 概形，转移态射为仿射态射。
% P11-E0312: the frozen source writes V_i in a limit indexed by lambda and then consistently uses V_lambda.
关于概形极限的材料，参见《极限》一节 \ref{limits-section-limits}。
我们将证明 $\colim F(V_\lambda) \to F(V)$ 为双射。

\medskip\noindent
单射性证明：记号。设 $\lambda \in \Lambda$，并设
$\xi_{\lambda, 1}, \xi_{\lambda, 2} \in F(V_\lambda)$ 是限制为 $F(V)$
中同一个元素的两个元素。写作
$\xi_{\lambda, 1} = (Z_{\lambda, 1}, u'_{\lambda, 1},
\hat x_{\lambda, 1})$ 以及
$\xi_{\lambda, 2} = (Z_{\lambda, 2}, u'_{\lambda, 2}, \hat x_{\lambda, 2})$。

\medskip\noindent
单射性证明：$Z_{\lambda, i}$ 的一致。由于 $Z_{\lambda, 1}$ 和
$Z_{\lambda, 2}$ 限制为 $V$ 的同一个闭子集，增大 $i$ 后可假设
$Z_{\lambda, 1} = Z_{\lambda, 2}$；参见《极限》引理
\ref{limits-lemma-inverse-limit-top} 和《拓扑》引理
\ref{topology-lemma-describe-limits}。
% P11-E0313: the frozen source says "after increasing i"; the directed-system index here is lambda.
以 $Z_\lambda \subset V_\lambda$ 表示这个共同值；对
$\mu \geq \lambda$，以 $Z_\mu \subset V_\mu$ 表示其在 $V_\mu$ 中的逆像，
并以 $Z$ 表示其在 $V$ 中的逆像。
% P11-E0314: the frozen source duplicates "and" before "denote Z".
下面将使用如下事实：若把 $Z$ 和 $Z_\mu$ 看作既约闭子概形，则作为概形有
$Z = \lim_{\mu \geq \lambda} Z_\mu$。

\medskip\noindent
单射性证明：$u'_{\lambda, i}$ 的一致。由于 $U'$ 在 $S$ 上局部有限型，
且 $u'_{\lambda, 1}$ 与 $u'_{\lambda, 2}$ 到 $V \setminus Z$ 的限制相同，
增大 $\lambda$ 后可假设 $u'_{\lambda, 1} = u'_{\lambda, 2}$；参见《极限》命题
\ref{limits-proposition-characterize-locally-finite-presentation}。
以 $u'_\lambda$ 表示共同值，并以 $u'$ 表示其到 $V \setminus Z$ 的限制。

\medskip\noindent
单射性证明：重述。此时有
$\xi_{\lambda, 1} = (Z_\lambda, u'_\lambda, \hat x_{\lambda, 1})$，以及
$\xi_{\lambda, 2} = (Z_\lambda, u'_\lambda, \hat x_{\lambda, 2})$。
本部分证明面临的主要问题是证明：增大 $\lambda$ 后，态射
$\hat x_{\lambda, 1}$ 与 $\hat x_{\lambda, 2}$ 变为相同。

\medskip\noindent
单射性证明：$\hat x_{\lambda, i}|_{Z_\lambda}$ 的一致。考虑态射
% P11-E0315: the frozen source misspells "morphisms" as "morpisms".
$\hat x_{\lambda, 1}|_{Z_\lambda}, \hat x_{\lambda, 2}|_{Z_\lambda} :
Z_\lambda \to W_{red}$.
这些态射限制为同一个态射 $Z \to W_{red}$。由于 $W_{red}$ 是 $S$ 上局部
有限型的概形，由《极限》命题
\ref{limits-proposition-characterize-locally-finite-presentation}
可知，把 $\lambda$ 换成更大的指标后，可假设
$\hat x_{\lambda, 1}|_{Z_\lambda} = \hat x_{\lambda, 2}|_{Z_\lambda} :
Z_\lambda \to W_{red}$.

\medskip\noindent
单射性证明：结束。接下来，把注 \ref{artin-remark-diagonal} 中的讨论应用于
$V_\lambda$ 以及两个元素
$\xi_{\lambda, 1}, \xi_{\lambda, 2} \in F(V_\lambda)$，得到
\begin{enumerate}
\item $e_\lambda : E_\lambda' \to V_\lambda$
分离且局部有限型；
\item $e_\lambda^{-1}(V_\lambda \setminus Z_\lambda) \to
V_\lambda \setminus Z_\lambda$ 是同构；
\item 单态射 $E_{W, \lambda} \to V_{\lambda, /Z_\lambda}$，它是
$\hat x_{\lambda, 1}$ 与 $\hat x_{\lambda, 2}$ 的等化子；
\item 形式修正 $E'_{\lambda, /Z_\lambda} \to E_{W, \lambda}$。
\end{enumerate}
由注 \ref{artin-remark-diagonal} 的断言 (2) 以及上面得到
$u'_{\lambda, 1} = u'_{\lambda, 2} = u'_\lambda$ 的准备工作，断言 (2) 成立。
由于
$\hat x_{\lambda, 1}|_{Z_\lambda} = \hat x_{\lambda, 2}|_{Z_\lambda}$
，$Z_\lambda = (V_{\lambda, /Z_\lambda})_{red}$ 通过
$E_{W, \lambda}$ 分解。因此，由《形式空间》引理
\ref{formal-spaces-lemma-monomorphism-iso-over-red} 可知，
$E_{W, \lambda} \to V_{\lambda, /Z_\lambda}$ 是闭浸入。再把注
\ref{artin-remark-diagonal} 的断言 (4) 和引理 \ref{artin-lemma-closed-immersion}
应用于 $V_\lambda$ 上的三元组
$E_\lambda'$、$e_\lambda^{-1}(Z_\lambda)$、
$E'_{\lambda, /Z_\lambda} \to E_{W, \lambda}$，可知存在闭浸入
$E_\lambda \to V_\lambda$，它是此三元组的解。
接下来，使用注 \ref{artin-remark-diagonal} 的断言 (5)，并结合引理
\ref{artin-lemma-solution}，可知 $E_\lambda$ 是 $\xi_{\lambda, 1}$ 与
$\xi_{\lambda, 2}$ 的“等化子”。特别地，$V \to V_\lambda$ 通过
$E_\lambda$ 分解。再次使用《极限》命题
\ref{limits-proposition-characterize-locally-finite-presentation}
，可找到 $\mu \geq \lambda$，使得 $V_\mu \to V_\lambda$ 通过
$E_\lambda$ 分解，因而 $\xi_{\lambda, 1}$ 和 $\xi_{\lambda, 2}$ 到
$V_\mu$ 的拉回相同，正如所需。

\medskip\noindent
满射性证明：陈述。设 $\xi = (Z, u', \hat x)$ 是 $F(V)$ 的元素。
必须找到 $\lambda \in \Lambda$ 和元素 $\xi_\lambda \in F(V_\lambda)$，
使其限制为 $\xi$。

\medskip\noindent
满射性证明：问题是 \'etale 局部的。由证明前一部分得到的唯一性以及引理
\ref{artin-lemma-sheaf} 中 $F$ 的层条件，此问题在 \'etale 拓扑下是 $V$ 上局部的。
更确切地，设 $v \in V$。我们断言，只需找到 \'etale 态射
$(\tilde V, \tilde v) \to (V, v)$、某个 $\lambda$、某个 \'etale 态射
$\tilde V_\lambda \to V_\lambda$，以及某个元素
$\tilde \xi_\lambda \in F(\tilde V_\lambda)$，使得
% P11-E0316: the frozen source has the malformed phrase "some an etale morphism".
$\tilde V = \tilde V_\lambda \times_{V_\lambda} V$
且 $\xi|_U = \tilde \xi_\lambda|_U$。
% P11-E0317: the frozen source uses U in both restrictions without defining U in this paragraph; the intended comparison appears to be after pullback to \tilde V.
略去此断言的详细证明\footnote{为证明这一点，把若干态射
$\tilde V \to V$ 组成一个有限 \'etale 覆盖，并证明相应态射
$\tilde V_\lambda \to V_\lambda$ 在增大 $\lambda$ 后也组成 \'etale 覆盖。
然后利用单射性看出诸元素 $\tilde \xi_\lambda$ 在增大 $\lambda$ 后可以黏合，
并利用 $F$ 的层条件把这些元素下降为 $F(V_\lambda)$ 的一个元素。}。

\medskip\noindent
满射性证明：重述问题。回忆，若 \'etale 态射
$(\tilde V, \tilde v) \to (V, v)$ 中的 $\tilde V$ 为仿射概形，则对某个
$\lambda$，它是某个 \'etale 态射 $\tilde V_\lambda \to V_\lambda$
的基变换，其中 $\tilde V_\lambda$ 为仿射概形；例如参见《拓扑》引理
\ref{topologies-lemma-limit-fppf-topology}。给定 $\tilde V_\lambda$，有
$\tilde V = \lim_{\mu \geq \lambda} \tilde V_\lambda \times_{V_\lambda} V_\mu$.
因此，给定 \'etale 态射 $(\tilde V, \tilde v) \to (V, v)$ 且
$\tilde V$ 仿射后，可以用 $(\tilde V, \tilde v)$ 替换 $(V, v)$，并用
$\xi$ 到 $\tilde V$ 的限制替换 $\xi$。

\medskip\noindent
满射性证明：约化到基为仿射的情形。特别地，设 $\tilde S \subset S$ 是仿射开
子概形，使得 $v \in V$ 映到 $\tilde S$ 的一个点。由上一段，可以用
$\tilde V = \tilde S \times_S V$ 替换 $V$。当然，此时只需对 $F$ 到
$\tilde S$ 上局部 Noether 概形范畴的限制解决问题。这个函子正是整个情形基变换到
$\tilde S$ 后所对应的函子。因此，在证明的余下部分中可以并且确实假设
$S = \Spec(R)$ 是 Noether 仿射概形。

\medskip\noindent
满射性证明：容易的情形。若 $v \in V \setminus Z$，则可取
$\tilde V = V \setminus Z$。由《极限》引理 \ref{limits-lemma-descend-opens}，
它下降为某个 $\lambda$ 下的开子概形
$\tilde V_\lambda \subset V_\lambda$。接着，增大 $\lambda$ 后可假设存在态射
$u'_\lambda : \tilde V_\lambda \to U'$，其限制为 $u'$。取
$\tilde \xi_\lambda = (\emptyset, u'_\lambda, \emptyset)$，便得到所需的
$F(\tilde V_\lambda)$ 元素。

\medskip\noindent
满射性证明：困难情形及到仿射 $W$ 的约化。最困难的情形来自考虑
$v \in Z \subset V$。我们断言可以约化到 $W$ 为仿射形式概形的情形；
建议读者跳过此论证\footnote{Artin 证明本引理的方法是绕过这一步，因而可以避免
先证明单射性。具体地，在证明过程中，Artin 始终使用所涉及所有空间的有限仿射
\'etale 覆盖，并追踪它们之间的映射。事后看来，这或许比我们这里的做法更可取。}。
% P11-E0318: the frozen source says "a finite affine etale coverings", mixing a singular article with a plural noun.
事实上，可以选取 \'etale 态射 $\tilde W \to W$，其中 $\tilde W$ 是仿射形式
代数空间，使得 $v$ 在 $\hat x : V_{/Z} \to W$ 下的像包含于
$\tilde W \to W$ 在既约空间上的像。于是态射
$$
p : \tilde W \times_{W, g} X'_{/T'} \longrightarrow X'_{/T'}
$$
以及
$$
q : \tilde W \times_{W, \hat x} V_{/Z} \to V_{/Z}
$$
都是局部 Noether 形式代数空间的 \'etale 态射。由《形式空间的代数化》定理
\ref{restricted-theorem-dilatations-general}
的一个容易情形，存在代数空间态射 $\tilde X' \to X'$；它局部有限型，在
$U'$ 上为同构，并且 $\tilde X'_{/T'} \to X'_{/T'}$ 与 $p$ 同构。
由《形式空间的代数化》引理 \ref{restricted-lemma-output-etale}，态射
$\tilde X' \to X'$ 是 \'etale 的。以
$\tilde T' \subset |\tilde X'|$ 表示 $T'$ 的逆像，以
$\tilde U' \subset \tilde X'$ 表示其补开子空间，并以
$\tilde g' : \tilde X'_{/\tilde T'} \to \tilde W$ 表示由
$\tilde W \to W$ 对 $g$ 作基变换所得的形式修正。
% P11-E0319: the frozen source uses three consecutive incomplete "Denote ... the ..." constructions.
于是
$$
\tilde X',\ \tilde T',\ \tilde U',\ \tilde W,
\ \tilde g : \tilde X'_{/\tilde T'} \to \tilde W
$$
是情形 \ref{artin-situation-contractions} 的另一个实例。
% P11-E0320: the frozen source defines \tilde g' in prose but uses \tilde g in the displayed tuple and subsequently.
以 $\tilde F$ 表示由此三元组构造的函子。
% P11-E0321: the frozen source uses the incomplete construction "Denote \tilde F the functor".
有函子变换
$$
\tilde F \longrightarrow F
$$
，它以显然的方式利用态射 $\tilde X' \to X'$ 和 $\tilde W \to W$ 构造；
略去细节。

\medskip\noindent
满射性证明：困难情形及到仿射 $W$ 的约化，第 2 部分。由上面使用的同一定理，
存在代数空间态射 $\tilde V \to V$；它局部有限型，在 $V \setminus Z$ 上为同构，
并且 $\tilde V_{/Z} \to V_{/Z}$ 与 $q$ 同构。以
$\tilde Z \subset \tilde V$ 表示 $Z$ 的逆像。
% P11-E0322: the frozen source uses the incomplete naming construction "Denote \tilde Z ... the inverse image".
由《形式空间的代数化》引理
\ref{restricted-lemma-output-etale} 和 \ref{restricted-lemma-output-separated}
，态射 $\tilde V \to V$ 是 \'etale 且分离的。特别地，$\tilde V$ 是（局部
Noether）概形；例如参见《空间的态射》命题
\ref{spaces-morphisms-proposition-locally-quasi-finite-separated-over-scheme}.
已有态射 $u'$，可将其视为态射
$$
\tilde u' : \tilde V \setminus \tilde Z \longrightarrow \tilde U'
$$
，其中 $\tilde U' \subset \tilde X'$ 是同构地映到 $U'$ 的开集。还有态射
$$
\tilde {\hat x} :
\tilde V_{/\tilde Z} = \tilde W \times_{W, \hat x} V_{/Z}
\longrightarrow
\tilde W
$$
，这里使用的就是投影。因此有三元组
$$
\tilde \xi = (\tilde Z, \tilde u', \tilde {\hat x}) \in \tilde F(\tilde V)
$$
略去相容性条件的证明。提示：若 $V' \to V$、$\hat x'$ 和 $x'$ 见证
$u'$ 与 $\hat x$ 之间的相容性，则令
$\tilde V' = V' \times_V \tilde V$；它带有态射
$\tilde{\hat x}'$ 和 $\tilde x'$，而这些数据确实可行。
% P11-E0323: the frozen source misspells "morphisms" as "morphsms" and uses "show" where the subject requires "shows".
$\tilde \xi$ 在变换 $\tilde F \to F$ 下的像是 $\xi$ 到 $\tilde V$ 的限制。

\medskip\noindent
满射性证明：困难情形及到仿射 $W$ 的约化，第 3 部分。由我们对
$\tilde W \to W$ 的选择，存在仿射开集
$\tilde V_{open} \subset \tilde V$（记号快不够用了），它在 $V$ 中的像包含
所选点 $v \in V$。由下一段研究的情形及前面的说明，可以把
$\tilde \xi|_{\tilde V_{open}}$ 下降为 $\tilde F$ 在
$\tilde V_{\lambda, open}$ 上的某个元素 $\tilde \xi_\lambda$；这里有某个
\'etale 态射 $\tilde V_{\lambda, open} \to V_\lambda$，其到 $V$ 的基变换为
$\tilde V_{open}$。应用函子变换 $\tilde F \to F$，便得到我们所求的
$F(\tilde V_{\lambda, open})$ 元素。这样便约化到下一段讨论的情形。

\medskip\noindent
满射性证明：$W$ 仿射的情形。此时 $v \in Z \subset V$，且 $W$ 是仿射形式
代数空间。回忆，
$$
\xi = (Z, u', \hat x) \in F(V)
$$
仍可用 $v$ 的 \'etale 邻域替换 $V$。特别地，可以并且确实假设 $V$ 和
$V_\lambda$ 都是仿射的。

\medskip\noindent
满射性证明：下降 $Z$。可以找到 $\lambda$ 以及闭子概形
$Z_\lambda \subset V_\lambda$，使得 $Z$ 是 $Z_\lambda$ 到 $V$ 的基变换；
参见《极限》引理 \ref{limits-lemma-descend-finite-presentation}。警告：我们不知道
$Z_\lambda$ 是否是 $V_\lambda$ 的既约闭子概形（一般而言也并非如此）。
对 $\mu \geq \lambda$，以 $Z_\mu \subset V_\mu$ 表示其在 $V_\mu$ 中的
概形论逆像。下面将使用作为概形的等式
$Z = \lim_{\mu \geq \lambda} Z_\mu$。

\medskip\noindent
满射性证明：下降 $u'$。由于 $U'$ 在 $S$ 上局部有限型，增大 $\lambda$ 后
可假设存在态射 $u'_\lambda : V_\lambda \setminus Z_\lambda \to U'$，
其到 $V \setminus Z$ 的限制为 $u'$；参见《极限》命题
\ref{limits-proposition-characterize-locally-finite-presentation}。
对 $\mu \geq \lambda$，以 $u'_\mu$ 表示 $u'_\lambda$ 到
$V_\mu \setminus Z_\mu$ 的限制。

\medskip\noindent
满射性证明：下降一个见证。设 $V' \to V$、$\hat x'$ 和 $x'$ 见证
$u'$ 与 $\hat x$ 之间的相容性。使用与上面相同的参考结果，增大
$\lambda$ 后可假设存在有限型态射 $V'_\lambda \to V_\lambda$，它到 $V$
的基变换为 $V' \to V$。再次增大 $\lambda$ 后，可假设
$V'_\lambda \to V_\lambda$ 为真态射（《极限》引理
\ref{limits-lemma-eventually-proper}）。接着，可假设
$V'_\lambda \to V_\lambda$ 在 $V_\lambda \setminus Z_\lambda$ 上为同构
（《极限》引理 \ref{limits-lemma-descend-isomorphism}）。再者，可假设存在态射
$x'_\lambda : V'_\lambda \to X'$，其到 $V'$ 的限制为 $x'$。再次增大
$\lambda$，可假设 $x'_\lambda$ 在 $V_\lambda \setminus Z_\lambda$
上与 $u'_\lambda$ 一致。对 $\mu \geq \lambda$，以 $V'_\mu$ 表示
$V'_\lambda$ 的基变换，并以 $x'_\mu$ 表示 $x'_\lambda$ 的限制。

\medskip\noindent
满射性证明：代数。写作 $W = \text{Spf}(B)$、$V = \Spec(A)$，并且对
$\mu \geq \lambda$ 写作 $V_\mu = \Spec(A_\mu)$。以
$I_\mu \subset A_\mu$ 和 $I \subset A$ 分别表示定义 $Z_\mu$ 和 $Z$ 的理想。
% P11-E0324: the frozen source uses the incomplete construction "Denote I_mu ... and I ... the ideals".
于是 $I_\lambda A_\mu = I_\mu$ 且 $I_\lambda A = I$。态射 $\hat x$
确定且由如下连续环同态确定：
$$
(\hat x)^\sharp : B \longrightarrow A^\wedge
$$
，其中 $A^\wedge$ 是 $A$ 的 $I$-adic 完备化。为完成证明，必须证明：
对某个充分大的 $\mu$，此映射下降为取值于 $A_\mu^\wedge$ 的映射，其中
$A_\mu^\wedge$ 是 $A_\mu$ 的 $I_\mu$-adic 完备化。这并非平凡事实；
Artin 在论文
\cite{ArtinII} 中写道：“由于数据 (3.5) 涉及 $I$-adic 完备化，而完备化不与
有向极限交换，验证颇为微妙。这是分析中收敛证明的一个代数类比。”

\medskip\noindent
满射性证明：代数，更多环。记
$$
C_\mu = \Gamma(V'_\mu, \mathcal{O})
\quad\text{且}\quad
C = \Gamma(V', \mathcal{O})
$$
注意，$A \to C$ 和 $A_\mu \to C_\mu$ 是有限环同态，因为
$V' \to V$ 和 $V'_\mu \to V_\mu$ 是真态射；参见《空间的上同调》引理
\ref{spaces-cohomology-lemma-proper-over-affine-cohomology-finite}.
由于 $V = \lim V_\mu$ 且 $V' = \lim V'_\mu$，有
$$
A = \colim A_\mu
\quad\text{且}\quad
C = \colim C_\mu
$$
；这是《极限》引理 \ref{limits-lemma-descend-section} 的结论\footnote{由于各个
态射并不平坦，我们不知道是否有
$C_\mu = C_\lambda \otimes_{A_\lambda} A_\mu$。}。对元素
$a \in I$（相应地，$a \in I_\mu$），由平坦基变换，映射 $A_a \to C_a$
（相应地，$(A_\mu)_a \to (C_\mu)_a$）是同构（《空间的上同调》引理
\ref{spaces-cohomology-lemma-flat-base-change-cohomology}）。
因此，$A \to C$ 的核与余核支撑在 $V(I)$ 上；$A_\mu \to C_\mu$ 也类似。
% P11-E0325: the frozen source uses singular "is supported" for the compound subject "kernel and cokernel".
由此可知，$A \to C$ 的核与余核被 $I$ 的某个幂零化，而
$A_\mu \to C_\mu$ 的核与余核被 $I_\mu$ 的某个幂零化；参见《代数》引理
\ref{algebra-lemma-Noetherian-power-ideal-kills-module}。

\medskip\noindent
满射性证明：代数，更多环同态。以 $Z_n \subset V$ 表示 $Z$ 的第 $n$ 阶
无穷小邻域，以 $Z_{\mu, n} \subset V_\mu$ 表示 $Z_\mu$ 的第 $n$ 阶
无穷小邻域。
% P11-E0326: the frozen source uses two incomplete "Denote ... the nth infinitesimal neighbourhood" constructions.
由形式函数定理（《空间的上同调》定理
\ref{spaces-cohomology-theorem-formal-functions})
，有
$$
C^\wedge = \lim_n H^0(V' \times_V Z_n, \mathcal{O})
\quad\text{且}\quad
C_\mu^\wedge =
\lim_n H^0(V'_\mu \times_{V_\mu} Z_{\mu, n}, \mathcal{O})
$$
，其中 $C^\wedge$ 和 $C_\mu^\wedge$ 分别是关于 $I$ 和 $I_\mu$ 的完备化。
% P11-E0327: the frozen source says two objects "are the completion" instead of "are the completions".
把态射 $x'_\mu : V'_\mu \to X'$ 的完备化与态射
$g : X'_{/T'} \to W$ 复合，得到
$$
g \circ x'_{\mu, /Z_\mu} :
V'_{\mu, /Z_\mu} = \colim V_\mu' \times_{V_\mu} Z_{\mu, n}
\longrightarrow
W
$$
，因而由上面对完备化 $C_\mu^\wedge$ 的描述，得到连续环同态
$$
(g \circ x'_{\mu, /Z_\mu})^\sharp : B \longrightarrow C_\mu^\wedge
$$
$V' \to V$、$\hat x'$、$x'$ 见证 $u'$ 与 $\hat x$ 之间的相容性这一事实
蕴含下图交换：
% P11-E0328: the frozen source uses singular "witnesses" for the compound subject V' -> V, \hat x', x'.
$$
\xymatrix{
C_\mu^\wedge \ar[r] &
C^\wedge \\
B \ar[u]^{(g \circ x'_{\mu, /Z_\mu})^\sharp} \ar[r]^{(\hat x)^\sharp} &
A^\wedge \ar[u]
}
$$

\medskip\noindent
满射性证明：更多代数论证。回忆，有限 $A$-模 $\Ker(A \to C)$ 和
$\Coker(A \to C)$ 被 $I$ 的某个幂零化；类似地，有限 $A_\mu$-模
$\Ker(A_\mu \to C_\mu)$ 和 $\Coker(A_\mu \to C_\mu)$ 被 $I_\mu$
的某个幂零化。这说明这些模等于其完备化。由于在有限 $A$-模范畴上作
$I$-adic 完备化是正合的（参见《代数》一节
\ref{algebra-section-completion-noetherian}），故有
$$
\Coker(A^\wedge \to C^\wedge) = \Coker(A \to C)
$$
；核以及映射 $A_\mu \to C_\mu$ 也类似。当然，还有
$$
\Ker(A \to C) = \colim \Ker(A_\mu \to C_\mu)
\quad\text{且}\quad
\Coker(A \to C) = \colim \Coker(A_\mu \to C_\mu)
$$
回忆，$S = \Spec(R)$ 为仿射概形。由于所有态射都是 $S$ 上的态射，上述所有
环同态都是 $R$-代数同态。由《形式空间的代数化》引理
\ref{restricted-lemma-Noetherian-finite-type-red}
可知，$B$ 在 $R$ 上拓扑有限型。设 $B$ 由 $b_1, \ldots, b_n$ 拓扑生成。
取某个 $\mu$（例如 $\lambda$），并考虑下列元素：
$$
\text{images of }
(g \circ x'_{\mu, /Z_\mu})^\sharp(b_1)
, \ldots,
(g \circ x'_{\mu, /Z_\mu})^\sharp(b_n)
\text{ 于 }\Coker(A_\mu \to C_\mu)
$$
由上面方块的交换性，这些元素在 $\Coker(\alpha)$ 中的像为零。
% P11-E0329: the frozen source refers to Coker(alpha), but no morphism alpha is defined in this proof.
由于 $\Coker(A \to C) = \colim \Coker(A_\mu \to C_\mu)$，且这些余核等于
其完备化，增大 $\mu$ 后可假设这些像全都为零。这意味着连续同态
$(g \circ x'_{\mu, /Z_\mu})^\sharp$ 的像包含于
$\Im(A_\mu \to C_\mu)$。选取元素 $a_{\mu, j} \in (A_\mu)^\wedge$，
使其在 $(C_\mu)^\wedge$ 中映到
$(g \circ x'_{\mu, /Z_\mu})^\sharp(b_1)$。
% P11-E0330: the frozen source chooses a_{\mu,j} but maps every one to the expression with b_1; the subsequent sentence uses b_j.
由上面方块的交换性，$a_{\mu, j} \in A_\mu^\wedge$ 和
$(\hat x)^\sharp(b_j) \in A^\wedge$ 映到 $C^\wedge$ 的同一个元素。
由于 $\Ker(A \to C) = \colim \Ker(A_\mu \to C_\mu)$，且这些核等于其
完备化，增大 $\mu$ 后可以调整 $a_{\mu, j}$ 的选择，使
$a_{\mu, j}$ 在 $A^\wedge$ 中的像等于 $(\hat x)^\sharp(b_j)$。

\medskip\noindent
满射性证明：最后的代数论证。设 $\mathfrak b \subset B$ 是拓扑幂零元素所成的
理想。设 $J \subset R[x_1, \ldots, x_n]$ 是由满足
$h(b_1, \ldots, b_n) \in \mathfrak b$ 的那些
$h(x_1, \ldots, x_n)$ 组成的理想。于是得到拓扑 $R$-代数的连续满射
$$
\Phi : R[x_1, \ldots, x_n]^\wedge \longrightarrow B,\quad
x_j \longmapsto b_j
$$
，其中左边关于 $J$ 完备化。由于 $R[x_1, \ldots, x_n]$ 是 Noether 环，
可为 $J$ 选取生成元 $h_1, \ldots, h_m$。由上面方块的交换性，
$h_j(a_{\mu, 1}, \ldots, a_{\mu, n})$ 是 $A_\mu^\wedge$ 的元素，且其在
$A^\wedge$ 中的像包含于 $IA^\wedge$。事实上，环同态
$(\hat x)^\sharp$ 连续，并且因为 $A^\wedge/IA^\wedge = A/I$ 既约，
$IA^\wedge$ 是 $A^\wedge$ 的拓扑幂零元素所成的理想。（关于 Noether 环中
完备化的结果，参见《代数》一节
\ref{algebra-section-completion-noetherian}。）由于
$A/I = \colim A_\mu/I_\mu$，增大 $\mu$ 后可假设
$h_j(a_{\mu, 1}, \ldots, a_{\mu, n})$ 属于 $I_\mu A_\mu^\wedge$。
特别地，$A_\mu^\wedge$ 的元素
$h_j(a_{\mu, 1}, \ldots, a_{\mu, n})$ 在 $A_\mu^\wedge$ 中拓扑幂零。
因此得到连续 $R$-代数同态
$$
\Psi : R[x_1, \ldots, x_n]^\wedge \longrightarrow A_\mu^\wedge,\quad
x_j \longmapsto a_{\mu, j}
$$
为得到所需结论，需要检验 $\Ker(\Phi)$ 是否被 $\Psi$ 零化。这未必立即成立，
但增大 $\mu$ 后可以做到。事实上，由于
$R[x_1, \ldots, x_n]^\wedge$ 是 Noether 环，可以为理想 $\Ker(\Phi)$
选取生成元 $g_1, \ldots, g_l$。于是
$$
\Psi(g_1), \ldots, \Psi(g_l) \in
\Ker(A_\mu^\wedge \to C_\mu^\wedge) = \Ker(A_\mu \to C_\mu)
$$
在 $\Ker(A \to C) = \colim \Ker(A_\mu \to C_\mu)$ 中映为零。
因此，像前面一样增大 $\mu$，便得到所需结果。

\medskip\noindent
满射性证明：收尾。上面构造的连续环同态 $B \to (A_\mu)^\wedge$ 确定态射
$\hat x_\mu : V_{\mu, /Z_\mu} \to W$。$\hat x_\mu$ 与 $u'_\mu$ 的相容性
来自如下事实：按构造，环同态 $B \to (A_\mu)^\wedge$ 与环同态
$A_\mu \to C_\mu$ 相容。事实上，我们在构造中使用的真态射
$V'_\mu \to V_\mu$ 以及态射 $x'_\mu$ 和
$\hat x'_\mu = x'_{\mu, /Z_\mu}$ 将见证这一相容性。证明完成。
\end{proof}

\begin{lemma}
\label{artin-lemma-rs}
在情形 \ref{artin-situation-contractions} 中，函子 $F$ 满足
Rim--Schlessinger 条件 (RS)。
\end{lemma}

\begin{proof}
回忆，此条件只涉及函子 $F$ 在 $S$ 上概形 $V$ 处的取值 $F(V)$，这些概形是
Artin 局部环的谱；以及由 $S$ 上概形态射 $V_1 \to V_2$ 诱导的限制映射
$F(V_2) \to F(V_1)$，这些概形同样是 Artin 局部环的谱。
因此，设 $V/S$ 是某个 Artin 局部环的谱。
% P11-E0331: the frozen source misspells "spectrum" as "spetruim".
若 $\xi = (Z, u', \hat x) \in F(V)$，则在集合论意义下要么
$Z = \emptyset$，要么 $Z = V$。在第一种情形中，$\hat x$ 是从空形式代数空间
到 $W$ 的态射。在第二种情形中，$u'$ 是从空概形到 $X'$ 的态射，而
$\hat x : V \to W$ 是到 $W$ 的态射。因此
$$
F(V) = U'(V) \amalg W(V)
$$
，而且对上述 $V_1 \to V_2$，诱导映射 $F(V_2) \to F(V_1)$ 与此分解相容。
因此，只需证明 $U'$ 和 $W$ 都满足 Rim--Schlessinger 条件。对 $U'$，
这由引理 \ref{artin-lemma-algebraic-stack-RS} 得到。为证明 $W$ 的情形，按
《形式空间》引理 \ref{formal-spaces-lemma-structure-locally-noetherian}
写作 $W = \colim W_n$。设 $V = \Spec(A)$，其中
$(A, \mathfrak m)$ 是 Artin 局部环。取 $n \geq 1$ 使
$\mathfrak m^n = 0$，则 $W(V) = W_n(V)$。故 $W$ 的 Rim--Schlessinger
条件由如下事实推出：对所有 $n$，$W_n$ 都满足 Rim--Schlessinger 条件；
而后者又由引理
\ref{artin-lemma-algebraic-stack-RS} 推出。
\end{proof}

\begin{lemma}
\label{artin-lemma-finite-dim}
在情形 \ref{artin-situation-contractions} 中，函子 $F$ 的切空间是有限维的。
\end{lemma}

\begin{proof}
在引理 \ref{artin-lemma-rs} 的证明中已经看到：若 $V$ 是 Artin 局部环的谱，则
$F(V) = U'(V) \amalg W(V)$。切空间完全由 $F$ 在 $S$ 上此类概形处的取值
计算。因此，只需证明函子 $U'$ 和 $W$ 的切空间有限维。对 $U'$，这由引理
\ref{artin-lemma-finite-dimension} 得到。像引理 \ref{artin-lemma-rs} 的证明中那样写作
$W = \colim W_n$。于是 $W$ 的切空间等于 $W_2$ 的切空间，因为为得到切空间，
只需在满足 $\mathfrak m^2 = 0$ 的 Artin 局部环
$(A, \mathfrak m)$ 的谱处计算 $W$。再次由引理
\ref{artin-lemma-finite-dimension} 可知，$W_2$ 的切空间有限维。
\end{proof}

\begin{lemma}
\label{artin-lemma-formal-object-effective}
在情形 \ref{artin-situation-contractions} 中，假设 $X' \to S$ 分离。则 $F$
的每个形式对象都有效。
\end{lemma}

\begin{proof}
$F$ 的形式对象 $\xi = (R, \xi_n)$ 由如下数据组成：Noether 完备局部
$S$-代数 $R$，其剩余域在 $S$ 上有限型；以及对所有 $n$ 的元素
$\xi_n \in F(\Spec(R/\mathfrak m^n))$，满足
$\xi_{n + 1}|_{\Spec(R/\mathfrak m^n)} = \xi_n$。由引理
\ref{artin-lemma-rs} 证明中的讨论，$\xi$ 要么是 $U'$ 的形式对象，要么是 $W$
的形式对象。在第一种情形中，由引理 \ref{artin-lemma-effective}，$\xi$ 有效。
第二种情形才有趣。令 $V = \Spec(R)$。我们将构造元素
$(Z, u', \hat x) \in F(V)$，使其在
$F(\Spec(R/\mathfrak m^n))$ 中的像对所有 $n \geq 1$ 都是 $\xi_n$。

\medskip\noindent
可以把元素族 $\xi_n$ 看作态射
$$
\xi : \text{Spf}(R) \longrightarrow W
$$
，这是 $S$ 上局部 Noether 形式代数空间的态射。注意，一般而言 $\xi$
{\it 不是} adic 态射。为修正这一点，设 $I \subset R$ 是对应于形式闭子空间
$$
\text{Spf}(R) \times_{\xi, W} W_{red} \subset \text{Spf}(R)
$$
的理想。注意 $I \subset \mathfrak m_R$。令
$Z = V(I) \subset V = \Spec(R)$。由于 $R$ 是 $\mathfrak m_R$-adic
完备的，它更是 $I$-adic 完备的（《代数》引理
\ref{algebra-lemma-complete-by-sub}）。此外，我们断言，对每个 $n \geq 1$，
态射
$$
\xi|_{\text{Spf}(R/I^n)} :
\text{Spf}(R/I^n)
\longrightarrow
W
$$
实际上来自态射
$$
\xi'_n :  \Spec(R/I^n) \longrightarrow W
$$
事实上，像引理 \ref{artin-lemma-rs} 的证明中那样写作 $W = \colim W_n$，
注意 $\xi|_{\text{Spf}(R/I^n)}$ 映入 $W_n$，再应用《形式空间》引理
\ref{formal-spaces-lemma-map-into-algebraic-space}
把它代数化为态射 $\Spec(R/I^n) \to W_n$，正如所需。以
$\text{Spf}'(R) = V_{/Z}$ 表示赋予 $I$-adic 拓扑的 $R$ 的形式谱——
等价地，它是 $V$ 沿 $Z$ 的形式完备化。利用态射 $\xi'_n$，得到 adic 态射
$$
\hat x = (\xi'_n) : \text{Spf}'(R) \longrightarrow W
$$
，它是 $S$ 上局部 Noether 形式代数空间的态射。考虑基变换
$$
\text{Spf}'(R) \times_{\hat x, W, g} X'_{/T'} \longrightarrow \text{Spf}'(R)
$$
由《形式空间的代数化》引理
\ref{restricted-lemma-base-change-formal-modification}，这是形式修正。
因此，由膨胀的主定理（《形式空间的代数化》定理
\ref{restricted-theorem-dilatations}），得到真态射
$$
V' \longrightarrow V = \Spec(R)
$$
，它在 $\Spec(R) \setminus V(I)$ 上为同构，且其完备化恢复上面的形式修正。
换言之，
$$
V' \times_{\Spec(R)} \Spec(R/I^n) =
\Spec(R/I^n) \times_{\xi'_n, W, g} X'_{/T'}
$$
特别地，这说明我们有相容的态射系
$$
V' \times_{\Spec(R)} \Spec(R/I^n) \longrightarrow X' \times_S \Spec(R/I^n)
$$
因此，由 Grothendieck 代数化定理（采用《空间的态射续篇》引理
\ref{spaces-more-morphisms-lemma-algebraize-morphism})
的形式），得到态射
$$
x' : V' \to X'
$$
，它在 $S$ 上并恢复上面显示的态射。最后，令
$u' : V \setminus Z \to X'$ 为 $x'$ 到 $V \setminus Z \subset V'$
的限制，便得到所求元素 $(Z, u', \hat x) \in F(V)$ 的第三个分量。
% P11-E0332: the frozen source uses the incomplete construction "setting u' ... the restriction".
\end{proof}

\begin{lemma}
\label{artin-lemma-openness-smoothness}
设 $S$ 是局部 Noether 概形，$V$ 是 $S$ 上局部有限型的概形，
$Z \subset V$ 为闭集。设 $W$ 是 $S$ 上局部 Noether 的形式代数空间，且
$W_{red}$ 在 $S$ 上局部有限型。设 $g : V_{/Z} \to W$ 是 $S$ 上形式代数
空间的 adic 态射。设 $v \in V$ 是闭点，使得 $g$ 在 $v$ 处通用（含义见一节
\ref{artin-section-axioms-functors}）。则用 $v$ 的一个开邻域替换 $V$ 后，态射
$g$ 光滑（见证明）。
\end{lemma}

\begin{proof}
由于 $g$ 是 adic 的，它可由代数空间表示（《形式空间》一节
\ref{formal-spaces-section-adic}）。因此，所谓 $g$ 光滑，是指 $g$ 按
《自举》定义 \ref{bootstrap-definition-property-transformation} 的意义光滑。

\medskip\noindent
像《形式空间》引理 \ref{formal-spaces-lemma-structure-locally-noetherian}
中那样写作 $W = \colim W_n$。令
$V_n = V_{/Z} \times_{\hat x, W} W_n$。
% P11-E0333: the frozen source uses \hat x in this fibre product, but the morphism introduced in the lemma is g.
则 $V_n$ 是底集合为 $Z$ 的闭子概形。$V \to W$ 光滑等价于所有态射
$V_n \to W_n$ 都光滑；这是因为当 $T$ 是拟紧概形时，任意态射
$T \to W$ 都通过某个 $n$ 所对应的 $W_n$ 分解。
% P11-E0334: the frozen source refers to smoothness of V -> W, although only g: V_{/Z} -> W has been defined.
由引理 \ref{artin-lemma-base-change-versal}，态射 $V_n \to W_n$ 在 $v$ 处光滑
\footnote{该引理适用，因为 $W$ 的对角态射可由代数空间表示且局部有限型；
参见《形式空间》引理
\ref{formal-spaces-lemma-diagonal-morphism-formal-algebraic-spaces}
，而且我们已在引理 \ref{artin-lemma-rs} 的证明中看到 $W$ 具有 (RS)。}。
当然，这意味着给定任意 $n$，可以缩小 $V$，使 $V_n \to W_n$ 光滑。
问题在于找到一个同时对所有 $n$ 都有效的开集。

\medskip\noindent
问题在 $V$ 上是局部的，故可假设 $S = \Spec(R)$ 和 $V = \Spec(A)$ 都是仿射的。

\medskip\noindent
本段约化到 $W$ 为仿射形式代数空间的情形。选取仿射形式概形 $W'$ 以及
\'etale 态射 $W' \to W$，使 $v$ 在 $W_{red}$ 中的像属于
$W'_{red} \to W_{red}$ 的像。则
$V_{/Z} \times_{g, W} W' \to V_{/Z}$ 是 $S$ 上形式代数空间的 adic
\'etale 态射，而 $V_{/Z} \times_{g, W} W'$ 是仿射形式代数空间。
由《形式空间的代数化》引理
\ref{restricted-lemma-algebraize-rig-etale-affine}，存在仿射概形的
\'etale 态射 $\varphi : V' \to V$，使得 $V'$ 沿
$Z' = \varphi^{-1}(Z)$ 的完备化在 $V_{/Z}$ 上同构于
$V_{/Z} \times_{g, W} W'$。注意，$v$ 是某个 $v' \in V'$ 的像。
由于光滑性在基变换下保持，对所有 $n$，$V'_n \to W'_n$ 都光滑。
下一段将证明：用 $v'$ 的一个开邻域替换 $V'$ 后，对所有 $n$，
态射 $V'_n \to W'_n$ 都光滑。然后，用 $V' \to V$ 的开像替换 $V$，
由光滑性的 \'etale 下降可知 $V_n \to W_n$ 光滑。略去一些细节。

\medskip\noindent
假设 $S = \Spec(R)$、$V = \Spec(A)$、$Z = V(I)$ 且
$W = \text{Spf}(B)$。设 $v$ 对应于极大理想
$I \subset \mathfrak m \subset A$。给定 adic 连续 $R$-代数同态
$$
B \longrightarrow A^\wedge
$$
设 $\mathfrak b \subset B$ 是拓扑幂零元素所成的理想（这是 Noether adic
拓扑环 $B$ 的最大定义理想）。注意，$\mathfrak b A^\wedge$ 和
$IA^\wedge$ 都是 Noether adic 环 $A^\wedge$ 的定义理想。此外，
$\mathfrak m A^\wedge$ 是 $A^\wedge$ 的极大理想，包含
$\mathfrak b A^\wedge$ 和 $IA^\wedge$。已知
$$
B_n = B/\mathfrak b^n \to A^\wedge/\mathfrak b^n A^\wedge = A_n
$$
对所有 $n$ 都在 $\mathfrak m$ 处光滑。由上面的讨论，可以并且确实假设
$B_1 \to A_1$ 是光滑环同态。以 $\mathfrak m_1 \subset A_1$ 表示对应于
$\mathfrak m$ 的极大理想。
% P11-E0335: the frozen source uses the incomplete construction "Denote m_1 ... the maximal ideal".
由于光滑性蕴含平坦性，对所有 $n \geq 1$，映射
$$
\mathfrak b^n/\mathfrak b^{n + 1} \otimes_{B_1} (A_1)_{\mathfrak m_1}
\longrightarrow
\left(\mathfrak b^nA^\wedge/\mathfrak b^{n + 1}A^\wedge\right)_{\mathfrak m_1}
$$
是同构；参见《代数》引理 \ref{algebra-lemma-what-does-it-mean-again}。
考虑 Rees 代数
$$
B' = \bigoplus\nolimits_{n \geq 0} \mathfrak b^n/\mathfrak b^{n + 1}
$$
，它是 Noether 环 $B_1$ 上有限型的分次代数；再考虑 Rees 代数
$$
A' = \bigoplus\nolimits_{n \geq 0}
\mathfrak b^nA^\wedge/\mathfrak b^{n + 1}A^\wedge
$$
，它是 Noether 环 $A_1$ 上有限型的分次代数。
% P11-E0336: the frozen source duplicates the article in "a a finite type graded algebra".
考虑分次 $A_1$-代数同态
$$
\Psi : B' \otimes_{B_1} A_1 \longrightarrow A'
$$
由上文，此映射在 $A_1$ 的极大理想 $\mathfrak m_1$ 处局部化后是同构。
因此，$\Ker(\Psi)$（相应地，$\Coker(\Psi)$）是
$B' \otimes_{B_1} A_1$（相应地，$A'$）上的有限模，且其在
$\mathfrak m_1$ 处的局部化为零。从而，用一个主局部化替换 $A_1$
（并相应替换 $A$）后，可假设 $\Psi$ 是同构。（这是证明的关键步骤。）
反向推导可知 $B_n \to A_n$ 平坦；参见《代数》引理
\ref{algebra-lemma-what-does-it-mean-again}。因此 $A_n \to B_n$ 光滑
（作为具有光滑纤维的平坦环同态；参见《代数》引理
\ref{algebra-lemma-flat-fibre-smooth}），证明完成。
% P11-E0337: the frozen source reverses the ring-map direction in the final smoothness conclusion; the map proved flat and used throughout is B_n -> A_n.
\end{proof}

\begin{lemma}
\label{artin-lemma-openness-versality}
在情形 \ref{artin-situation-contractions} 中，函子 $F$ 满足通用性的开性。
\end{lemma}

\begin{proof}
必须证明如下事实：给定 $S$ 上局部有限型的概形 $V$、$\xi \in F(V)$，
以及有限型点 $v_0 \in V$，且 $\xi$ 在 $v_0$ 处通用，则用 $v_0$ 的一个
开邻域替换 $V$ 后，$\xi$ 在 $V$ 的每个有限型点处都通用。写作
$\xi = (Z, u', \hat x)$。

\medskip\noindent
第一种情形：$v_0 \not \in Z$。首先可用 $V \setminus Z$ 替换 $V$。
于是 $\xi = (\emptyset, u', \emptyset)$，且态射
$u' : V \to X'$ 在 $v_0$ 处通用。由《空间的态射续篇》引理
\ref{spaces-more-morphisms-lemma-lifting-along-artinian-at-point}
，这意味着 $u' : V \to X'$ 在 $v_0$ 处光滑。由于一个态射的光滑点集合是开集，
缩小 $V$ 后可假设 $u'$ 光滑。
% P11-E0338: the frozen source has the spurious article in "the set of a points".
同一引理于是说明 $\xi$ 在每一点处通用，正如所需。

\medskip\noindent
第二种情形：$v_0 \in Z$。像《形式空间》引理
\ref{formal-spaces-lemma-structure-locally-noetherian} 中那样写作
$W = \colim W_n$。由引理 \ref{artin-lemma-openness-smoothness}，可假设
$\hat x : V_{/Z} \to W$ 是形式代数空间的光滑态射。立即可知，
$\xi = (Z, u', \hat x)$ 在 $Z$ 的所有有限型点处通用。设
$V' \to V$、$\hat x'$ 和 $x'$ 见证 $u'$ 与 $\hat x$ 之间的相容性。
作为 $\hat x$ 的基变换，$\hat x' : V'_{/Z} \to X'_{/T'}$ 光滑。
由于 $\hat x'$ 是 $x' : V' \to X'$ 的完备化，由已经使用过的
《空间的态射续篇》引理
\ref{spaces-more-morphisms-lemma-lifting-along-artinian-at-point}.
可知，$x' : V' \to X'$ 在
$(V' \to V)^{-1}(Z) = |x'|^{-1}(T') \subset |V'|$ 的所有点处光滑。
由于一个态射的光滑点集合是开集，$x'$ 不光滑的点所成的闭集
$B \subset |V'|$ 与 $(V' \to V)^{-1}(Z)$ 不相交。由于
$V' \to V$ 为真态射，因而是闭态射，故
$(V' \to V)(B) \subset V$ 是与 $Z$ 不相交的闭子集。因此，缩小 $V$ 后
可假设 $B = \emptyset$，即 $x'$ 光滑。由上一段的讨论，这恰好意味着
$\xi$ 在 $V$ 中不属于 $Z$ 的所有有限型点处通用，证明完成。
\end{proof}

\noindent
下面给出最终结果。

\begin{theorem}
\label{artin-theorem-contractions}
\begin{reference}
\cite[Theorem 3.1]{ArtinII}
\end{reference}
设 $S$ 是局部 Noether 概形，且对所有有限型点 $s \in S$，
$\mathcal{O}_{S, s}$ 都是 G-环。设 $X'$ 是 $S$ 上局部有限型的代数空间，
$T' \subset |X'|$ 是闭子集。设 $W$ 是 $S$ 上局部 Noether 的形式代数空间，
且 $W_{red}$ 在 $S$ 上局部有限型。最后，设
$$
g : X'_{/T'} \longrightarrow W
$$
为形式修正；参见《形式空间的代数化》定义
\ref{restricted-definition-formal-modification}。若 $X'$ 和 $W$ 在 $S$
上分离\footnote{参见注 \ref{artin-remark-separated-needed}。}，则存在 $S$ 上
代数空间的真态射 $f : X' \to X$、闭子集 $T \subset |X|$，以及形式代数空间的
同构 $a : X_{/T} \to W$，使得
\begin{enumerate}
\item $T'$ 是 $T$ 在 $|f| : |X'| \to |X|$ 下的逆像；
\item $f : X' \to X$ 将 $X' \setminus T'$ 同构地映到 $X \setminus T$；以及
\item $g = a \circ f_{/T}$，其中
$f_{/T} : X'_{/T'} \to X_{/T}$ 是诱导态射。
\end{enumerate}
换言之，$(f : X' \to X, T, a)$ 是本节前面定义的一个解。
\end{theorem}

\begin{proof}
设 $F$ 是本节中利用 $X'$、$T'$、$W$、$g$ 构造的函子。由引理
\ref{artin-lemma-functor-is-solution}，只需证明 $F$ 对应于 $S$ 上局部有限型的
代数空间 $X$。为此，将应用命题
\ref{artin-proposition-spaces-diagonal-representable-noetherian}。具体而言，
由引理 \ref{artin-lemma-diagonal-contractions}，$F$ 的对角态射可由闭浸入表示；
又由诸引理 \ref{artin-lemma-sheaf}、\ref{artin-lemma-limit-preserving}、
\ref{artin-lemma-rs}, \ref{artin-lemma-finite-dim},
\ref{artin-lemma-formal-object-effective} 和 \ref{artin-lemma-openness-versality}，
我们得到公理 [0]、[1]、[2]、[3]、[4]、[5]。
\end{proof}

\begin{remark}
\label{artin-remark-separated-needed}
定理 \ref{artin-theorem-contractions} 的证明在两处使用了 $X'$ 和 $W$ 在 $S$
上分离这一事实。第一处是用它证明
$\Delta : F \to F \times F$ 可由代数空间表示。若把分离情形下的定理应用于
注 \ref{artin-remark-diagonal} 中的三元组
$E'$、$(E' \to V)^{-1}Z$、$E'_{/Z} \to E_W$，便可证明 $\Delta$ 可表示，
从而完全避免这一次使用该假设（这是对角态射通常的自举过程）。
% P11-E0339: the frozen source says "the triples" although E', the inverse image, and the formal modification are the three components of one triple.
因此，引理 \ref{artin-lemma-formal-object-effective} 的证明，是我们对定理
\ref{artin-theorem-contractions} 的证明中唯一真正需要使用 $X' \to S$ 分离之处。
读者可以检验，我们只用该假设来得到态射
$x' : V' \to X'$。若 $X' \to S$ 拟分离，则利用文献中的结果可以证明
$x'$ 存在；参见《空间的态射续篇》注
\ref{spaces-more-morphisms-remark-weaken-separation-axioms-question}.
由此可知，该定理在把
“分离”替换为“拟分离”后仍按原陈述成立。若以后需要这一版本，我们会在这里
精确陈述并仔细证明。
\end{remark}
% Generated mechanically from the frozen quot.tex target.
% Do not edit this generated file; edit the bound locale source instead.

% OK, start here.
%

\setcounter{chapter}{98}
\chapter{Quot 空间与 Hilbert 空间}


\phantomsection
\label{quot-section-phantom}


\section{引言}
\label{quot-section-introduction}

\noindent
本章最初设想的目的，是讨论 Quot 与 Hilbert 函子，并证明在满足若干技术条件时，
它们是代数空间。对概形而言，这些材料见于 Grothendieck 在 Bourbaki 讨论班的讲义；
% P11-E0340: the frozen source misspells the French word "séminaire" as "séminair".
参见
\cite{Gr-I},
\cite{Gr-II},
\cite{Gr-III},
\cite{Gr-IV},
\cite{Gr-V}，以及
\cite{Gr-VI}。对射影概形，Quot 概形与 Hilbert 概形位于适当很充足可逆层的
截面空间的 Grassmann 簇中，这给出了构造这些概形的方法。我们采用不同的路径：
利用 Artin 公理证明 Quot 与 Hilb 函子是代数空间。

\medskip\noindent
进一步考虑后发现，为了发展 Stacks Project 中的理论，从
\cite{lieblich_remarks} 引入的凝聚层叠 $\Cohstack_{X/B}$
（其支撑在基底上为真）开始讨论更为便利。这里，$f : X \to B$
是满足适当技术条件的代数空间态射，不过这一点可以推广（见下文）。
给定 $X$ 上的模 $\mathcal{F}$ 和 $\mathcal{G}$，在适当假设下，函子
% P11-E0341: the frozen source says "under suitably hypotheses" instead of "under suitable hypotheses".
$T/B \mapsto \Hom_{X_T}(\mathcal{F}_T, \mathcal{G}_T)$
是 $B$ 上的代数空间 $\mathit{Hom}(\mathcal{F}, \mathcal{G})$；
见一节 \ref{quot-section-hom}。在一节 \ref{quot-section-isom} 中，我们证明由同构组成的
子函子 $\mathit{Isom}(\mathcal{F}, \mathcal{G})$ 是代数空间。随后几节利用这一点
证明叠 $\Cohstack_{X/B}$ 的对角态射可表示。当 $X \to B$ 平坦时，我们在一节
\ref{quot-section-stack-coherent-sheaves} 中证明 $\Cohstack_{X/B}$ 是代数叠；
一般情形见一节 \ref{quot-section-not-flat}。文献指引见该节的引言。

\medskip\noindent
证明这一点后，不难进一步证明
$\Quotfunctor_{\mathcal{F}/X/B}$、$\Hilbfunctor_{X/B}$ 和
$\Picardfunctor_{X/B}$ 是代数空间，而 $\Picardstack_{X/B}$ 是代数叠；
见诸节 \ref{quot-section-quot}、\ref{quot-section-hilb}、
\ref{quot-section-picard-functor} 和 \ref{quot-section-picard-stack}。

\medskip\noindent
用通常的方法，我们在一节 \ref{quot-section-relative-morphisms} 中推出：
在适当假设下，相对态射函子 $\mathit{Mor}_B(Z, X)$ 是代数空间。

\medskip\noindent
在一节 \ref{quot-section-stack-of-spaces} 中，我们证明参数化真代数空间平坦族的群胚叠
$$
\Spacesstack'_{fp, flat, proper}
$$
满足除形式有效性以外的所有 Artin 公理（包括通用性的开性）。我们有意为这个叠
选用非常笨拙的记号，因为读者使用其性质时应当谨慎。

\medskip\noindent
在一节 \ref{quot-section-polarized} 中，我们证明参数化极化真代数空间平坦族的叠
$\Polarizedstack$ 是代数叠。由于已经处理了真代数空间的平坦族，这归结为证明
极化概形的形式有效性；该结果通常称为 Grothendieck 代数化定理。

\medskip\noindent
在一节 \ref{quot-section-curves} 中，我们证明参数化曲线族的叠
$\Curvesstack$ 是代数的。

\medskip\noindent
在一节 \ref{quot-section-moduli-complexes} 中，我们研究真态射上复形的模，
并得到代数叠 $\Complexesstack_{X/B}$。陈述与证明的思路取自
\cite{lieblich-complexes}。

\medskip\noindent
本章不包含什么？对于所得模空间和模叠所具有的性质（除代数性以外），本章几乎
不作讨论；相关内容参见《模叠》一节 \ref{moduli-section-introduction}。
% P11-E0342: the frozen source omits "that" in "the properties the resulting moduli spaces ... possess".
在所讨论的大多数结果中，可以用代数叠的态射
$\mathcal{X} \to \mathcal{B}$ 代替代数空间的态射 $X \to B$，从而推广这些构造。
% P11-E0343: the frozen source has the singular phrase "a morphism ... of algebraic space".
我们将在别处讨论这一点（此处插入未来引用）。
% P11-E0344: the frozen source retains the unresolved placeholder "(insert future reference here)".
对 Hilbert 空间还有更一般的“Hilbert 叠”概念，我们将在单独的一章中讨论；参见
（此处插入未来引用）。
% P11-E0345: the frozen source retains a second unresolved future-reference placeholder.




\section{约定}
\label{quot-section-conventions}

\noindent
我们有意把本章以及《叠的例子》《代数叠上的层》《可表示性判据》和《Artin 公理》
诸章放在代数叠理论的一般发展之前。原因是，从下一章开始（见《代数叠的性质》一节
\ref{stacks-properties-section-conventions}），我们不再区分概形与它所产生的代数叠。
因此，我们的语言将变得更灵活、更便于人理解，但也不那么精确。开头这几章——包括
最初的《代数叠》一章——奠定了基础，使我们以后可以忽略理解代数叠的不同方式之间
某些非常技术性的区别。
% P11-E0346: the frozen source has "groundwork that later allow us" instead of "groundwork that later allows us".
不过，特别是在《Artin 公理》和《可表示性判据》两章中，我们必须非常精确地说明
自己究竟处理哪些对象，因为我们正试图证明某些构造会产生代数叠或代数空间。
% P11-E0347: the frozen source calls the chapter "Criteria of Representability", inconsistent with its title "Criteria for Representability".

\medskip\noindent
遗憾的是，这意味着有些记号、约定和术语颇为别扭，在经验较丰富的读者看来甚至可能
前后颠倒。希望读者谅解！

\medskip\noindent
始终假设所有概形都包含在一个大 fppf 站点 $\Sch_{fppf}$ 中，并且所考虑的每个环
$A$ 都满足：$\Spec(A)$ 同构于该大站点的某个对象。

\medskip\noindent
设 $S$ 是概形，$X$ 是 $S$ 上的代数空间。在本章及下一章中，我们用
$X \times_S X$ 表示 $X$ 与自身的积（在 $S$ 上代数空间的范畴中），
而不用 $X \times X$。
\section{Hom 函子}
\label{quot-section-hom}

\noindent
本节研究下面定义的同态函子。

\begin{situation}
\label{quot-situation-hom}
设 $S$ 是概形，$f : X \to B$ 是 $S$ 上代数空间的态射。设
$\mathcal{F}$、$\mathcal{G}$ 是拟凝聚 $\mathcal{O}_X$-模。
对任意 $B$ 上的概形 $T$，把 $\mathcal{F}$ 和 $\mathcal{G}$ 到 $T$ 的基变换
记为 $\mathcal{F}_T$ 和 $\mathcal{G}_T$；换言之，它们是经由投影态射
$X_T = X \times_B T \to X$ 的拉回。
% P11-E0348: the frozen source uses the incomplete construction "denote F_T and G_T the base changes".
考虑函子
\begin{equation}
\label{quot-equation-hom}
\mathit{Hom}(\mathcal{F}, \mathcal{G}) :
(\Sch/B)^{opp}
\longrightarrow
\textit{Sets},\quad
T
\longrightarrow
\Hom_{\mathcal{O}_{X_T}}(\mathcal{F}_T, \mathcal{G}_T)
\end{equation}
\end{situation}

\noindent
在情形 \ref{quot-situation-hom} 中，有时把
$\mathit{Hom}(\mathcal{F}, \mathcal{G})$ 看作配有态射
$\mathit{Hom}(\mathcal{F}, \mathcal{G}) \to B$ 的函子
$(\Sch/S)^{opp} \to \textit{Sets}$。
具体地，若 $T$ 是 $S$ 上的概形，则
$\mathit{Hom}(\mathcal{F}, \mathcal{G})(T)$ 的一个元素由一对
$(h, u)$ 组成，其中 $h$ 是态射 $h : T \to B$，而
$u : \mathcal{F}_T \to \mathcal{G}_T$ 是 $\mathcal{O}_{X_T}$-模同态；
这里 $X_T = T \times_{h, B} X$，且 $\mathcal{F}_T$ 和 $\mathcal{G}_T$
是到 $X_T$ 的拉回。特别地，当我们说
$\mathit{Hom}(\mathcal{F}, \mathcal{G})$ 是代数空间时，是指相应函子
$(\Sch/S)^{opp} \to \textit{Sets}$ 是代数空间。

\begin{lemma}
\label{quot-lemma-hom-sheaf}
在情形 \ref{quot-situation-hom} 中，函子
$\mathit{Hom}(\mathcal{F}, \mathcal{G})$
满足 fpqc 拓扑的层条件。
\end{lemma}

\begin{proof}
设 $\{T_i \to T\}_{i \in I}$ 是 $B$ 上概形的 fpqc 覆盖。令
$X_i = X_{T_i} = X \times_S T_i$、$\mathcal{F}_i = u_{T_i}$
和 $\mathcal{G}_i = \mathcal{G}_{T_i}$。
% P11-E0349: the frozen source writes X times_S T_i although Situation situation-hom defines X_{T_i} as X times_B T_i.
% P11-E0350: the frozen source sets F_i = u_{T_i}; no u is defined here and the intended module is F_{T_i}.
注意，$\{X_i \to X_T\}_{i \in I}$ 是 $X_T$ 的 fpqc 覆盖；参见
《空间上的拓扑》引理 \ref{spaces-topologies-lemma-fpqc}。
因此，若映射族 $u_i : \mathcal{F}_i \to \mathcal{G}_i$ 满足
$u_i$ 和 $u_j$ 在 $X_{T_i \times_T T_j}$ 上限制为同一个映射，则由下降
（《空间上的下降》命题
\ref{spaces-descent-proposition-fpqc-descent-quasi-coherent}），它唯一地来自映射
$u : \mathcal{F}_T \to \mathcal{G}_T$。
\end{proof}

\noindent
一致性检验：$\mathit{Hom}$ 层在 $S$ 上代数空间之间起着同样的作用。

\begin{lemma}
\label{quot-lemma-extend-hom-to-spaces}
在情形 \ref{quot-situation-hom} 中，设 $T$ 是 $S$ 上的代数空间。则
% P11-E0351: the frozen source leaves "In Situation ..." as a sentence fragment.
$$
\Mor_{\Sh((\Sch/S)_{fppf})}(T, \mathit{Hom}(\mathcal{F}, \mathcal{G})) =
\{(h, u) \mid h : T \to B, u : \mathcal{F}_T \to \mathcal{G}_T\}
$$
其中 $\mathcal{F}_T,\mathcal{G}_T$ 表示 $\mathcal{F}$ 和 $\mathcal{G}$
到代数空间 $X \times_{B, h} T$ 的拉回。
\end{lemma}

\begin{proof}
选取概形 $U$ 和满的 \'etale 态射 $p : U \to T$。令
$R = U \times_T U$，其投影为 $t, s : R \to U$。

\medskip\noindent
设 $v : T \to \mathit{Hom}(\mathcal{F}, \mathcal{G})$
是自然变换。则 $v(p)$ 对应于 $U$ 上的一对 $(h_U, u_U)$。
由于 $v$ 是函子变换，$(h_U, u_U)$ 沿 $s$ 与 $t$ 的拉回相同。
因为 $T = U/R$（《空间》引理 \ref{spaces-lemma-space-presentation}），
得到态射 $h : T \to B$，使得 $h_U = h \circ p$。于是
$\mathcal{F}_U$ 是 $\mathcal{F}_T$ 到 $X_U$ 的拉回，
$\mathcal{G}_U$ 亦然。因此，由《空间上的下降》命题
\ref{spaces-descent-proposition-fpqc-descent-quasi-coherent}，
$u_U$ 下降为 $\mathcal{O}_{X_T}$-模同态
$u : \mathcal{F}_T \to \mathcal{G}_T$。

\medskip\noindent
反之，设 $(h, u)$ 是 $T$ 上的一对。对概形 $T'$ 的每个态射
$a : T' \to T$，令其对应于 $(h \circ a, a^*u)$，便得到自然变换
$v : T \to \mathit{Hom}(\mathcal{F}, \mathcal{G})$。
略去验证本段与上一段的构造互为逆构造。
\end{proof}

\begin{remark}
\label{quot-remark-hom-base-change}
在情形 \ref{quot-situation-hom} 中，设 $B' \to B$ 是 $S$ 上代数空间的态射。
令 $X' = X \times_B B'$，并把 $\mathcal{F}$、$\mathcal{G}$ 到 $X'$
的拉回记为 $\mathcal{F}'$、$\mathcal{G}'$。
% P11-E0352: the frozen source uses the incomplete construction "denote F', G' the pullback".
于是，与基变换 $f' : X' \to B'$ 相联系，得到函子
$\mathit{Hom}(\mathcal{F}', \mathcal{G}') : (\Sch/B')^{opp} \to \textit{Sets}$。
对 $B'$ 上的概形 $T$，显然有
$$
\mathit{Hom}(\mathcal{F}', \mathcal{G}')(T) =
\mathit{Hom}(\mathcal{F}, \mathcal{G})(T)
$$
其中右边经由复合 $T \to B' \to B$ 把 $T$ 看作 $B$ 上的概形。
这个平凡的注记有时可用于改变基代数空间。
\end{remark}

\begin{lemma}
\label{quot-lemma-hom-sheaf-in-X}
在情形 \ref{quot-situation-hom} 中，设 $\{X_i \to X\}_{i \in I}$ 是 fppf 覆盖，
并对每个 $i, j \in I$ 设
$\{X_{ijk} \to X_i \times_X X_j\}$ 是 fppf 覆盖。
把 $\mathcal{F}$ 到 $X_i$、$X_{ijk}$ 的拉回分别记为
$\mathcal{F}_i$、$\mathcal{F}_{ijk}$；类似地定义
$\mathcal{G}_i$ 和 $\mathcal{G}_{ijk}$。
% P11-E0353: the frozen source uses the incomplete "Denote ..., resp. ... the pullback" construction.
对每个 $B$ 上的概形 $T$，图
$$
\xymatrix{
\mathit{Hom}(\mathcal{F}, \mathcal{G})(T) \ar[r] &
\prod\nolimits_i
\mathit{Hom}(\mathcal{F}_i, \mathcal{G}_i)(T)
\ar@<1ex>[r]^-{\text{pr}_0^*} \ar@<-1ex>[r]_-{\text{pr}_1^*}
&
\prod\nolimits_{i, j, k}
\mathit{Hom}(\mathcal{F}_{ijk}, \mathcal{G}_{ijk})(T)
}
$$
把第一个箭头表示为后两个箭头的等化子。
\end{lemma}

\begin{proof}
设 $u_i : \mathcal{F}_{i, T} \to \mathcal{G}_{i, T}$ 是
$\text{pr}_0^*$ 与 $\text{pr}_1^*$ 的等化子中的元素。由于 fppf 覆盖的基变换
仍是 fppf 覆盖（《空间上的拓扑》引理
\ref{spaces-topologies-lemma-fppf}），可知
$\{X_{i, T} \to X_T\}_{i \in I}$ 和
$\{X_{ijk, T} \to X_{i, T} \times_{X_T} X_{j, T}\}$ 都是 fppf 覆盖。
应用《空间上的下降》命题
\ref{spaces-descent-proposition-fpqc-descent-quasi-coherent}，
先得 $u_i$ 和 $u_j$ 在 $X_{i, T} \times_{X_T} X_{j, T}$ 上限制为同一个态射；
再应用一次该命题，得到唯一态射 $u : \mathcal{F}_T \to \mathcal{G}_T$，
其在每个 $i$ 上限制为 $u_i$。证毕。
\end{proof}

\begin{lemma}
\label{quot-lemma-hom-limits}
在情形 \ref{quot-situation-hom} 中，若 $\mathcal{F}$ 有限表示，且 $f$ 拟紧、拟分离，
则 $\mathit{Hom}(\mathcal{F}, \mathcal{G})$ 保极限。
\end{lemma}

\begin{proof}
设 $T = \lim_{i \in I} T_i$ 是仿射 $B$-概形的有向极限。需要证明
$$
\mathit{Hom}(\mathcal{F}, \mathcal{G})(T) =
\colim \mathit{Hom}(\mathcal{F}, \mathcal{G})(T_i).
$$
选取 $0 \in I$。可以把 $B$ 替换为 $T_0$，把 $X$ 替换为 $X_{T_0}$，
把 $\mathcal{F}$、$\mathcal{G}$ 分别替换为 $\mathcal{F}_{T_0}$、
$\mathcal{G}_{T_0}$，并把 $I$ 替换为 $\{i \in I \mid i \geq 0\}$；
参见注记 \ref{quot-remark-hom-base-change}。因此可假设
$B = \Spec(R)$ 为仿射的。

\medskip\noindent
当 $B$ 仿射时，$X$ 拟紧且拟分离。选取满的 \'etale 态射 $U \to X$，
其中 $U$ 是仿射概形（《空间的性质》引理
\ref{spaces-properties-lemma-quasi-compact-affine-cover}）。
因为 $X$ 拟分离，概形 $U \times_X U$ 拟紧，故可选取满的 \'etale 态射
$V \to U \times_X U$，其中 $V$ 是仿射概形。应用引理
\ref{quot-lemma-hom-sheaf-in-X}，可知
$\mathit{Hom}(\mathcal{F}, \mathcal{G})$ 是下列两个对象之间两个映射的等化子：
$$
\mathit{Hom}(\mathcal{F}|_U, \mathcal{G}|_U)
\quad\text{与}\quad
\mathit{Hom}(\mathcal{F}|_V, \mathcal{G}|_V).
$$
这把问题约化到 $X$ 仿射的情形。

\medskip\noindent
在仿射情形，引理的陈述归结为如下问题：给定环同态 $R \to A$、两个
$A$-模 $M$、$N$，以及 $R$-代数的有向系统 $C = \colim C_i$，何时映射
$$
\colim \Hom_{A \otimes_R C_i}(M \otimes_R C_i, N \otimes_R C_i)
\longrightarrow
\Hom_{A \otimes_R C}(M \otimes_R C, N \otimes_R C)
$$
为双射？由《代数》引理
\ref{algebra-lemma-module-map-property-in-colimit}，若
$M \otimes_R C$ 作为 $A \otimes_R C$-模有限表示，则确实如此；
这等价于 $M$ 作为 $A$-模有限表示。
\end{proof}

\begin{lemma}
\label{quot-lemma-hom-closed}
设 $S$ 是概形，$B$ 是 $S$ 上的代数空间。设
$i : X' \to X$ 是 $B$ 上代数空间的闭浸入。设 $\mathcal{F}$ 是拟凝聚
$\mathcal{O}_X$-模，$\mathcal{G}'$ 是拟凝聚 $\mathcal{O}_{X'}$-模。则
$$
\mathit{Hom}(\mathcal{F}, i_*\mathcal{G}') =
\mathit{Hom}(i^*\mathcal{F}, \mathcal{G}')
$$
作为 $(\Sch/B)$ 上的函子相等。
\end{lemma}

\begin{proof}
设 $g : T \to B$ 是态射，其中 $T$ 是概形。把 $i$ 的基变换记为
$i_T : X'_T \to X_T$，并把投影记为
$h : X_T \to X$ 和 $h' : X'_T \to X'$。
% P11-E0354: the frozen source uses the incomplete construction "Denote i_T ... the base change".
% P11-E0355: the frozen source uses the incomplete construction "Denote h ... and h' ... the projections".
注意，$(h')^*i^*\mathcal{F} = i_T^*h^*\mathcal{F}$。
由于闭浸入是仿射的（《空间的态射》引理
\ref{spaces-morphisms-lemma-closed-immersion-affine}），由
《空间的上同调》引理 \ref{spaces-cohomology-lemma-affine-base-change} 有
$h^*i_*\mathcal{G} = i_{T, *}(h')^*\mathcal{G}$。
% P11-E0356: the frozen proof switches from the stated module G' to an undefined unprimed G in the base-change identity and the middle three Hom terms.
因此
\begin{align*}
\mathit{Hom}(\mathcal{F}, i_*\mathcal{G}')(T)
& =
\Hom_{\mathcal{O}_{X_T}}(h^*\mathcal{F}, h^*i_*\mathcal{G}') \\
& =
\Hom_{\mathcal{O}_{X_T}}(h^*\mathcal{F}, i_{T, *}(h')^*\mathcal{G}) \\
& =
\Hom_{\mathcal{O}_{X'_T}}(i_T^*h^*\mathcal{F}, (h')^*\mathcal{G}) \\
& =
\Hom_{\mathcal{O}_{X'_T}}((h')^*i^*\mathcal{F}, (h')^*\mathcal{G}) \\
& =
\mathit{Hom}(i^*\mathcal{F}, \mathcal{G}')(T),
\end{align*}
正如所需。中间的等号来自函子 $i_{T, *}$ 与 $i_T^*$ 的伴随性。
\end{proof}

\begin{lemma}
\label{quot-lemma-cohomology-perfect-complex}
设 $S$ 是概形，$B$ 是 $S$ 上的代数空间，$K$ 是
$D(\mathcal{O}_B)$ 的伪凝聚对象。
\begin{enumerate}
\item 若对 $(\Sch/B)$ 中所有 $g : T \to B$，上同调层
$H^{-1}(Lg^*K)$ 都为零，则函子
$$
(\Sch/B)^{opp} \longrightarrow \textit{Sets},\quad
(g : T \to B) \longmapsto H^0(T, H^0(Lg^*K))
$$
是 $B$ 上仿射且有限表示的代数空间。
\item 若对 $(\Sch/B)$ 中所有 $g : T \to B$，上同调层
$H^i(Lg^*K)$ 对 $i < 0$ 都为零，则 $K$ 完美，$K$ 局部具有
$[0, b]$ 中的 Tor 振幅，而且函子
$$
(\Sch/B)^{opp} \longrightarrow \textit{Sets},\quad
(g : T \to B) \longmapsto H^0(T, Lg^*K)
$$
是 $B$ 上仿射且有限表示的代数空间。
\end{enumerate}
\end{lemma}

\begin{proof}
在 (2) 的假设下，有
$H^0(T, Lg^*K) = H^0(T, H^0(Lg^*K))$。
先证明规则 $T \mapsto H^0(T, H^0(Lg^*K))$ 满足 fppf 拓扑的层条件。
为此，设 $\{h_i : T_i \to T\}$ 是 $B$ 上概形
$g : T \to B$ 的 fppf 覆盖。令 $g_i = g \circ h_i$。
注意，因为 $h_i$ 平坦，故 $Lh_i^* = h_i^*$，且 $h_i^*$ 与取上同调可交换。因此
$$
H^0(T_i, H^0(Lg_i^*K)) =
H^0(T_i, H^0(h_i^*Lg^*K)) =
H^0(T, h_i^*H^0(Lg^*K)).
$$
% P11-E0357: the frozen third term uses H^0(T, ...) although h_i^*H^0(Lg^*K) is a sheaf on T_i and both preceding terms use T_i.
拉回到 $T_i \times_T T_j$ 时同样如此。
由于 $Lg^*K$ 是 $T$ 上的伪凝聚复形（《站点上的上同调》引理
\ref{sites-cohomology-lemma-pseudo-coherent-pullback}），上同调层
$\mathcal{F} = H^0(Lg^*K)$ 是拟凝聚的（《空间的导出范畴》引理
\ref{spaces-perfect-lemma-pseudo-coherent}）。因此，由《空间上的下降》命题
\ref{spaces-descent-proposition-fpqc-descent-quasi-coherent} 可知
$$
H^0(T, \mathcal{F}) = \Ker(
\prod H^0(T_i, h_i^*\mathcal{F}) \to
\prod H^0(T_i \times_T T_j, (T_i \times_T T_j \to T)^*\mathcal{F})).
$$
由此可见，(1) 与 (2) 中的规则满足 fppf 覆盖的层条件。这意味着可以应用
《构造代数空间》引理
\ref{bootstrap-lemma-locally-algebraic-space-finite-type}，
把可表示性的证明约化为 $B$ 上 \'etale 局部的证明。此外，结果是否仿射且有限表示
也可在 $B$ 上 \'etale 局部检验；参见《空间的态射》诸引理
\ref{spaces-morphisms-lemma-affine-local} 和
\ref{spaces-morphisms-lemma-finite-presentation-local}。
因此可假设 $B$ 是仿射概形。

\medskip\noindent
设 $B = \Spec(A)$ 是仿射概形。由《空间的导出范畴》诸引理
\ref{spaces-perfect-lemma-pseudo-coherent}、
\ref{spaces-perfect-lemma-derived-quasi-coherent-small-etale-site} 和
\ref{spaces-perfect-lemma-descend-pseudo-coherent}，可知在余下证明中，
可把 $K$ 看作 Zariski 拓扑下 $B$ 上模复形导出范畴的完美对象。
由《概形的导出范畴》诸引理
\ref{perfect-lemma-pseudo-coherent}、
\ref{perfect-lemma-affine-compare-bounded} 和
\ref{perfect-lemma-pseudo-coherent-affine}，可找到 $A$-模的伪凝聚复形
$M^\bullet$，使得 $K$ 是 $D(\mathcal{O}_B)$ 中的相应对象。
关于拉回的假设蕴含：对每个素理想 $\mathfrak p \subset A$，
$M^\bullet \otimes^\mathbf{L}_A \kappa(\mathfrak p)$ 的 $H^{-1}$ 为零。
由《代数续篇》引理 \ref{more-algebra-lemma-cut-complex-in-two}，可写成
$$
M^\bullet =
\tau_{\geq 0}M^\bullet \oplus \tau_{\leq - 1}M^\bullet,
$$
其中对某个 $b \geq 0$，$\tau_{\geq 0}M^\bullet$ 是 Tor 振幅位于
$[0, b]$ 的完美复形（这里还使用了《代数续篇》诸引理
\ref{more-algebra-lemma-glue-perfect} 和
\ref{more-algebra-lemma-glue-tor-amplitude}）。
注意，在情形 (2) 中还可看出
$\tau_{\leq - 1}M^\bullet = 0$ 于 $D(A)$ 中，因而
$M^\bullet$ 和 $K$ 完美且 Tor 振幅位于 $[0, b]$。
对任意 $B$-概形 $g : T \to B$，有
$$
H^0(T, H^0(Lg^*K)) = H^0(T, H^0(Lg^*\tau_{\geq 0}K))
$$
（由《导出范畴》引理 \ref{derived-lemma-negative-vanishing} 的对偶）。
因此可以把 $K$ 替换为 $\tau_{\geq 0}K$，相应地把
$M^\bullet$ 替换为 $\tau_{\geq 0}M^\bullet$。换言之，可假设
$M^\bullet$ 的 Tor 振幅位于 $[0, b]$。

\medskip\noindent
设 $M^\bullet$ 的 Tor 振幅位于 $[0, b]$。可以假设 $M^\bullet$
是有限自由 $A$-模的上有界复形（由伪凝聚复形的定义；参见《代数续篇》定义
\ref{more-algebra-definition-pseudo-coherent} 及该定义之后的讨论）。
由《代数续篇》引理 \ref{more-algebra-lemma-last-one-flat}，
$M = \Coker(M^{- 1} \to M^0)$ 平坦。再由《代数》引理
\ref{algebra-lemma-finite-projective}，$M$ 有限局部自由。因此，
$M^\bullet$ 拟同构于
$$
M \to M^1 \to M^2 \to \ldots \to M^d \to 0 \ldots.
$$
注意，这是 K-平坦复形（《上同调》引理
\ref{cohomology-lemma-bounded-flat-K-flat}），所以经态射 $T \to B$
对 $K$ 作导出拉回可由复形
$$
g^*\widetilde{M} \to g^*\widetilde{M^1} \to \ldots
$$
计算。因此，只需证明函子
$$
(g : T \to B) \longmapsto
\Ker(
\Gamma(T,g^*\widetilde{M})
\to
\Gamma(T, g^*(\widetilde{M^1})
)
$$
可由 $B$ 上有限表示的仿射概形表示。

\medskip\noindent
为了证明最后的陈述，仍可把 $B$ 替换为某个仿射开覆盖的各成员。因此可假设
$M$ 有限自由（回忆，$M^1$ 原本就是有限自由的）。写
$M = A^{\oplus n}$ 和 $M^1 = A^{\oplus m}$。设映射 $M \to M^1$
由系数在 $A$ 中的 $m \times n$ 矩阵 $(a_{ij})$ 给出。于是
$\widetilde{M} = \mathcal{O}_B^{\oplus n}$ 且
$\widetilde{M^1} = \mathcal{O}_B^{\oplus m}$。上面的函子等于
$$
(g : T \to B) \longmapsto
\{(f_1, \ldots, f_n) \in \Gamma(T, \mathcal{O}_T) \mid
\sum g^\sharp(a_{ij})f_i = 0,\ j = 1, \ldots, m\}.
$$
显然，它可由仿射概形
$$
\Spec\left(A[x_1, \ldots, x_n]/(\sum a_{ij}x_i; j = 1, \ldots, m)\right)
$$
表示。引理得证。
\end{proof}

\noindent
函子 $\mathit{Hom}(\mathcal{F}, \mathcal{G})$ 在许多情形中可表示。
我们的所有结果都以如下基本情形为基础。下面给出的引理证明，在某种意义下是
\cite[III, Cor 7.7.8]{EGA} 的证明的自然推广。
% P11-E0358: the frozen source says "the natural generalization to the proof" rather than "of the proof".

\begin{lemma}
\label{quot-lemma-noetherian-hom}
在情形 \ref{quot-situation-hom} 中，假设
\begin{enumerate}
\item $B$ 是 Noether 代数空间；
\item $f$ 局部有限型且拟分离；
\item $\mathcal{F}$ 是有限型 $\mathcal{O}_X$-模；以及
\item $\mathcal{G}$ 是有限型 $\mathcal{O}_X$-模，在 $B$ 上平坦，且其支撑在
$B$ 上真。
\end{enumerate}
则函子 $\mathit{Hom}(\mathcal{F}, \mathcal{G})$ 是
$B$ 上仿射且有限表示的代数空间。
\end{lemma}

\begin{proof}
可以把 $X$ 替换为 $\mathcal{G}$ 的支撑的一个拟紧开邻域，因而可假设
$X$ 为 Noether 的。在这种情形，$X$ 与 $f$ 均拟紧且拟分离。
选取三元组 $(X, \mathcal{F}, -1)$ 的一个由完美复形 $P$ 给出的逼近
$P \to \mathcal{F}$；参见《空间的导出范畴》定义
\ref{spaces-perfect-definition-approximation-holds} 和定理
\ref{spaces-perfect-theorem-approximation}。
% P11-E0359: the frozen source closes this reference sentence with an unmatched parenthesis.
于是诱导映射
$$
\Hom_{\mathcal{O}_X}(\mathcal{F}, \mathcal{G})
\longrightarrow
\Hom_{D(\mathcal{O}_X)}(P, \mathcal{G})
$$
为同构，因为 $P \to \mathcal{F}$ 诱导同构
$H^0(P) \to \mathcal{F}$，且对 $i > 0$ 有 $H^i(P) = 0$。
此外，对任意态射 $g : T \to B$，把投影
$h : X_T = T \times_B X \to X$ 记下，并令 $P_T = Lh^*P$。
% P11-E0360: the frozen source uses the incomplete naming construction "denote h ... the projection".
同样有同构
$$
\Hom_{\mathcal{O}_{X_T}}(\mathcal{F}_T, \mathcal{G}_T)
\longrightarrow
\Hom_{D(\mathcal{O}_{X_T})}(P_T, \mathcal{G}_T),
$$
因为 $P_T = Lh^*P \to Lh^*\mathcal{F} \to \mathcal{F}_T$
诱导同构 $H^0(P_T) \to \mathcal{F}_T$（这是因为 $h^*$ 右正合，且对
$i > 0$ 有 $H^i(P) = 0$）。因此，只需对函子
$$
T \longmapsto \Hom_{D(\mathcal{O}_{X_T})}(P_T, \mathcal{G}_T)
$$
证明结论。由 Leray 谱序列（见《站点上的上同调》注记
\ref{sites-cohomology-remark-before-Leray}），有
$$
\Hom_{D(\mathcal{O}_{X_T})}(P_T, \mathcal{G}_T) =
H^0(X_T, R\SheafHom(P_T, \mathcal{G}_T)) =
H^0(T, Rf_{T, *}R\SheafHom(P_T, \mathcal{G}_T)),
$$
其中 $f_T : X_T \to T$ 是 $f$ 的基变换。由《空间的导出范畴》引理
\ref{spaces-perfect-lemma-base-change-RHom}，有
$$
Rf_{T, *}R\SheafHom(P_T, \mathcal{G}_T) = Lg^*Rf_*R\SheafHom(P, \mathcal{G}).
$$
由《空间的导出范畴》引理
\ref{spaces-perfect-lemma-ext-perfect}，
$D(\mathcal{O}_B)$ 的对象
$K = Rf_*R\SheafHom(P, \mathcal{G})$ 是完美的。因此，只要能证明对上述每个
$g : T \to B$ 及所有 $i < 0$，上同调层 $H^i(Lg^*K)$ 都为 $0$，
便可应用引理 \ref{quot-lemma-cohomology-perfect-complex}。
由最后一个显示公式，这一点很清楚：$Rf_{T, *}R\SheafHom(P_T, \mathcal{G}_T)$
在负次数的上同调层为零，因为 $P_T$ 完美且正次数上同调层为零，所以
$R\SheafHom(P_T, \mathcal{G}_T)$ 的负次数上同调层为零。
\end{proof}

\noindent
下面是引理 \ref{quot-lemma-noetherian-hom} 的一个直接推论。

\begin{proposition}
\label{quot-proposition-hom}
在情形 \ref{quot-situation-hom} 中，假设
\begin{enumerate}
\item $f$ 有限表示；以及
\item $\mathcal{G}$ 是有限表示 $\mathcal{O}_X$-模，在 $B$ 上平坦，且其支撑在
$B$ 上真。
\end{enumerate}
则函子 $\mathit{Hom}(\mathcal{F}, \mathcal{G})$ 是
$B$ 上的仿射代数空间。若 $\mathcal{F}$ 有限表示，则
$\mathit{Hom}(\mathcal{F}, \mathcal{G})$ 在 $B$ 上有限表示。
\end{proposition}

\begin{proof}
由引理 \ref{quot-lemma-hom-sheaf}，函子
$\mathit{Hom}(\mathcal{F}, \mathcal{G})$ 满足 fppf 覆盖的层条件。
这意味着可以\footnote{略去验证所引用引理的集合论条件 (3)。}应用
% P11-E0361: the frozen source says "This mean we may" instead of "This means we may".
% P11-E0362: the frozen footnote uses the open compound "set theoretical" rather than "set-theoretical".
《构造代数空间》引理 \ref{bootstrap-lemma-locally-algebraic-space}，
在 $B$ 上 \'etale 局部检验可表示性。此外，结果是否仿射或有限表示
也可在 $B$ 上 \'etale 局部检验；参见《空间的态射》诸引理
\ref{spaces-morphisms-lemma-affine-local} 和
\ref{spaces-morphisms-lemma-finite-presentation-local}。
因此可假设 $B$ 是仿射概形。

\medskip\noindent
假设 $B$ 是仿射概形。由于 $f$ 有限表示，可知 $X$ 拟紧且拟分离。
% P11-E0363: the frozen source omits "that" in "it follows X is quasi-compact".
因此，可把 $\mathcal{F} = \colim \mathcal{F}_i$ 写成有限表示
$\mathcal{O}_X$-模的滤过余极限（《空间的极限》引理
\ref{spaces-limits-lemma-colimit-finitely-presented}）。显然
$$
\mathit{Hom}(\mathcal{F}, \mathcal{G}) =
\lim \mathit{Hom}(\mathcal{F}_i, \mathcal{G}).
$$
因此，若每个 $\mathit{Hom}(\mathcal{F}_i, \mathcal{G})$
都可由仿射概形表示，则 $\mathit{Hom}(\mathcal{F}, \mathcal{G})$ 亦然。
使用《极限》一节 \ref{limits-section-limits} 与《空间的极限》一节
\ref{spaces-limits-section-limits} 的材料。于是可假设
$\mathcal{F}$ 有限表示。

\medskip\noindent
写 $B = \Spec(R)$。把 $R = \colim R_i$ 写成极限，其中每个 $R_i$
是有限型 $\mathbf{Z}$-代数。令 $B_i = \Spec(R_i)$。
由《空间的极限》诸引理
\ref{spaces-limits-lemma-descend-finite-presentation} 和
\ref{spaces-limits-lemma-descend-modules-finite-presentation}，
可找到 $i$、代数空间的态射 $X_i \to B_i$，以及有限表示
$\mathcal{O}_{X_i}$-模 $\mathcal{F}_i$ 和 $\mathcal{G}_i$，使得
$(X_i, \mathcal{F}_i, \mathcal{G}_i)$ 到 $B$ 的基变换恢复
$(X, \mathcal{F}, \mathcal{G})$。由《空间的极限》引理
\ref{spaces-limits-lemma-descend-flat}，增大 $i$ 后可假设
$\mathcal{G}_i$ 在 $B_i$ 上平坦。类似地，由《空间的极限》引理
\ref{spaces-limits-lemma-eventually-proper-support}，可假设
$\mathcal{G}_i$ 的概形论支撑在 $B_i$ 上真。此时可应用引理
\ref{quot-lemma-noetherian-hom}，得到
$H_i = \mathit{Hom}(\mathcal{F}_i, \mathcal{G}_i)$ 是
$B_i$ 上仿射且有限表示的代数空间。把它拉回到 $B$（使用注记
\ref{quot-remark-hom-base-change}），得到
$H_i \times_{B_i} B = \mathit{Hom}(\mathcal{F}, \mathcal{G})$，证毕。
\end{proof}






\section{Isom 函子}
\label{quot-section-isom}

\noindent
在情形 \ref{quot-situation-hom} 中，可以考虑子函子
$$
\mathit{Isom}(\mathcal{F}, \mathcal{G}) \subset
\mathit{Hom}(\mathcal{F}, \mathcal{G}),
$$
它在 $B$ 上概形 $T$ 处的值，是所有{\it 可逆}
$\mathcal{O}_{X_T}$-同态 $u : \mathcal{F}_T \to \mathcal{G}_T$ 所成的集合。

\medskip\noindent
有时把 $\mathit{Isom}(\mathcal{F}, \mathcal{G})$ 看作配有态射
$\mathit{Isom}(\mathcal{F}, \mathcal{G}) \to B$ 的函子
$(\Sch/S)^{opp} \to \textit{Sets}$。
具体地，若 $T$ 是 $S$ 上的概形，则
$\mathit{Isom}(\mathcal{F}, \mathcal{G})(T)$ 的一个元素由一对
$(h, u)$ 组成，其中 $h$ 是态射 $h : T \to B$，而
$u : \mathcal{F}_T \to \mathcal{G}_T$ 是 $\mathcal{O}_{X_T}$-模同构；
这里 $X_T = T \times_{h, B} X$，且 $\mathcal{F}_T$ 和 $\mathcal{G}_T$
是到 $X_T$ 的拉回。特别地，当我们说
$\mathit{Isom}(\mathcal{F}, \mathcal{G})$ 是代数空间时，是指相应函子
$(\Sch/S)^{opp} \to \textit{Sets}$ 是代数空间。

\begin{lemma}
\label{quot-lemma-isom-sheaf}
在情形 \ref{quot-situation-hom} 中，函子
$\mathit{Isom}(\mathcal{F}, \mathcal{G})$
满足 fpqc 拓扑的层条件。
\end{lemma}

\begin{proof}
已经看到 $\mathit{Hom}(\mathcal{F}, \mathcal{G})$ 满足层条件。
因此只需证明：给定 $B$ 上概形的 fpqc 覆盖
$\{T_i \to T\}_{i \in I}$ 和 $\mathcal{O}_{X_T}$-线性映射
$u : \mathcal{F}_T \to \mathcal{G}_T$，若对每个 $i$，
$u_{T_i}$ 都是同构，则 $u$ 是同构。由于
$\{X_i \to X_T\}_{i \in I}$ 是 $X_T$ 的 fpqc 覆盖（见
《空间上的拓扑》引理 \ref{spaces-topologies-lemma-fpqc}），
结论由《空间上的下降》命题
\ref{spaces-descent-proposition-fpqc-descent-quasi-coherent} 得出。
\end{proof}

\noindent
一致性检验：$\mathit{Isom}$ 层在 $S$ 上代数空间之间起着同样的作用。

\begin{lemma}
\label{quot-lemma-extend-isom-to-spaces}
在情形 \ref{quot-situation-hom} 中，设 $T$ 是 $S$ 上的代数空间。则
% P11-E0364: the frozen source leaves "In Situation ..." as a sentence fragment.
$$
\Mor_{\Sh((\Sch/S)_{fppf})}(T, \mathit{Isom}(\mathcal{F}, \mathcal{G})) =
\{(h, u) \mid
h : T \to B, u : \mathcal{F}_T \to \mathcal{G}_T\text{ 为同构}\},
$$
其中 $\mathcal{F}_T,\mathcal{G}_T$ 表示 $\mathcal{F}$ 和 $\mathcal{G}$
到代数空间 $X \times_{B, h} T$ 的拉回。
\end{lemma}

\begin{proof}
注意，该等式的左右两边分别是引理
\ref{quot-lemma-extend-hom-to-spaces} 中等式左右两边的子集。
% P11-E0365: the frozen source twice uses singular "hand side" for the paired left and right sides.
略去验证在该引理证明给出的对应下，这两个子集彼此对应。
\end{proof}

\begin{proposition}
\label{quot-proposition-isom}
在情形 \ref{quot-situation-hom} 中，假设
\begin{enumerate}
\item $f$ 有限表示；以及
\item $\mathcal{F}$ 和 $\mathcal{G}$ 是有限表示
$\mathcal{O}_X$-模，在 $B$ 上平坦，且其支撑在 $B$ 上真。
\end{enumerate}
则函子 $\mathit{Isom}(\mathcal{F}, \mathcal{G})$ 是
$B$ 上仿射且有限表示的代数空间。
\end{proposition}

\begin{proof}
采用缩写
$H = \mathit{Hom}(\mathcal{F}, \mathcal{G})$、
$I = \mathit{Hom}(\mathcal{F}, \mathcal{F})$、
$H' = \mathit{Hom}(\mathcal{G}, \mathcal{F})$ 和
$I' = \mathit{Hom}(\mathcal{G}, \mathcal{G})$。
由命题 \ref{quot-proposition-hom}，函子 $H$、$I$、$H'$、$I'$ 是代数空间，
而态射 $H \to B$、$I \to B$、$H' \to B$ 和 $I' \to B$
都是仿射且有限表示的。映射的复合给出 $B$ 上代数空间的态射
$$
c : H' \times_B H \longrightarrow I \times_B I',\quad
(u', u) \longmapsto (u \circ u', u' \circ u).
$$
% P11-E0366: the frozen ordered pair reverses the two composites relative to I = Hom(F,F) and I' = Hom(G,G).
由于 $I \times_B I' \to B$ 分离，对应于
$(\text{id}_\mathcal{F}, \text{id}_\mathcal{G})$ 的截面
$\sigma : B \to I \times_B I'$ 是闭浸入
（《空间的态射》引理 \ref{spaces-morphisms-lemma-section-immersion}）。
此外，$\sigma$ 有限表示（《空间的态射》引理
\ref{spaces-morphisms-lemma-finite-presentation-permanence}）。
因此
$$
\mathit{Isom}(\mathcal{F}, \mathcal{G}) =
(H' \times_B H) \times_{c, I \times_B I', \sigma} B
$$
也是 $B$ 上仿射且有限表示的代数空间。略去若干细节。
\end{proof}






\section{凝聚层叠}
\label{quot-section-stack-coherent-sheaves}

\noindent
本节证明：在适当假设下，$X/B$ 上的凝聚层叠是代数的。
这是 \cite[定理 2.1.1]{lieblich_remarks} 的一个特例；
该定理处理基底上 Artin 叠的凝聚层叠。

\begin{situation}
\label{quot-situation-coherent}
设 $S$ 是概形，$f : X \to B$ 是 $S$ 上代数空间的态射。
假设 $f$ 有限表示。用 $\Cohstack_{X/B}$ 表示如下范畴，
其对象是三元组 $(T, g, \mathcal{F})$，其中
\begin{enumerate}
\item $T$ 是 $S$ 上的概形；
\item $g : T \to B$ 是 $S$ 上的态射，并记
$X_T = T \times_{g, B} X$
\item $\mathcal{F}$ 是有限表示的拟凝聚 $\mathcal{O}_{X_T}$-模，
在 $T$ 上平坦，且其支撑在 $T$ 上真。
\end{enumerate}
态射 $(T, g, \mathcal{F}) \to (T', g', \mathcal{F}')$
由一对 $(h, \varphi)$ 给出，其中
\begin{enumerate}
\item $h : T \to T'$ 是 $B$ 上概形的态射
（即 $g' \circ h = g$）；以及
\item $\varphi : (h')^*\mathcal{F}' \to \mathcal{F}$ 是
$\mathcal{O}_{X_T}$-模同构，其中 $h' : X_T \to X_{T'}$
是 $h$ 的基变换。
\end{enumerate}
\end{situation}

\noindent
于是 $\Cohstack_{X/B}$ 是范畴，而规则
$$
p : \Cohstack_{X/B} \longrightarrow (\Sch/S)_{fppf},
\quad
(T, g, \mathcal{F}) \longmapsto T
$$
是函子。对 $S$ 上概形 $T$，用 $\Cohstack_{X/B, T}$ 表示
$p$ 在 $T$ 上的纤维范畴。这些纤维范畴都是群胚。

\begin{lemma}
\label{quot-lemma-coherent-fibred-in-groupoids}
在情形 \ref{quot-situation-coherent} 中，函子
$p : \Cohstack_{X/B} \longrightarrow (\Sch/S)_{fppf}$
是群胚纤维化范畴。
\end{lemma}

\begin{proof}
验证《范畴论》定义 \ref{categories-definition-fibred-groupoids}
的条件 (1) 和 (2)，即可证明 $p$ 是群胚纤维化范畴。
给定 $\Cohstack_{X/B}$ 的对象 $(T', g', \mathcal{F}')$ 和
$S$ 上概形的态射 $h : T \to T'$，可以令
% P11-E0367: the frozen source writes the composition in the ill-typed order h \circ g'.
$g = h \circ g'$ 以及 $\mathcal{F} = (h')^*\mathcal{F}'$，
其中 $h' : X_T \to X_{T'}$ 是 $h$ 的基变换。显然，这给出
$\Cohstack_{X/B}$ 中位于 $h$ 上方的态射
$(T, g, \mathcal{F}) \to (T', g', \mathcal{F}')$。这证明了 (1)。
为证明 (2)，设给定态射
$$
(h_1, \varphi_1) : (T_1, g_1, \mathcal{F}_1) \to (T, g, \mathcal{F})
\quad\text{和}\quad
(h_2, \varphi_2) : (T_2, g_2, \mathcal{F}_2) \to (T, g, \mathcal{F})
$$
属于 $\Cohstack_{X/B}$，并设有态射 $h : T_1 \to T_2$ 满足
$h_2 \circ h = h_1$。令 $\varphi$ 为复合
$$
(h')^*\mathcal{F}_2
\xrightarrow{(h')^*\varphi_2^{-1}}
(h')^*(h_2)^*\mathcal{F} = (h_1)^*\mathcal{F}
\xrightarrow{\varphi_1}
\mathcal{F}_1
$$
便得到态射
$(h, \varphi) : (T_1, g_1, \mathcal{F}_1) \to (T_2, g_2, \mathcal{F}_2)$
，从而验证条件 (2)。
\end{proof}

\begin{lemma}
\label{quot-lemma-coherent-diagonal}
在情形 \ref{quot-situation-coherent} 中，记
% P11-E0368: the frozen source punctuates "In Situation ..." as a sentence fragment.
$\mathcal{X} = \Cohstack_{X/B}$。则
$\Delta : \mathcal{X} \to \mathcal{X} \times \mathcal{X}$
可由代数空间表示。
\end{lemma}

\begin{proof}
考虑概形 $T$ 上 $\mathcal{X}$ 的两个对象
$x = (T, g, \mathcal{F})$ 和 $y = (T, h, \mathcal{G})$。
必须证明 $\mathit{Isom}_\mathcal{X}(x, y)$ 是 $T$ 上的代数空间，
见《代数叠》引理 \ref{algebraic-lemma-representable-diagonal}。
若对 $a : T' \to T$，限制 $x|_{T'}$ 和 $y|_{T'}$ 在纤维范畴
$\mathcal{X}_{T'}$ 中同构，则 $g \circ a = h \circ a$。
因此有预层的自然变换
$$
\mathit{Isom}_\mathcal{X}(x, y) \longrightarrow \text{等化子}(g, h)
$$
由于 $B$ 的对角态射可表示（由概形表示），该等化子是概形。
于是可以用此等化子替换 $T$，并用相应拉回替换层
$\mathcal{F}$ 和 $\mathcal{G}$。因而可以假设 $g = h$。此时
$\mathit{Isom}_\mathcal{X}(x, y) = \mathit{Isom}(\mathcal{F}, \mathcal{G})$
，结论由命题 \ref{quot-proposition-isom} 得出。
\end{proof}

\begin{lemma}
\label{quot-lemma-coherent-stack}
在情形 \ref{quot-situation-coherent} 中，函子
$p : \Cohstack_{X/B} \longrightarrow (\Sch/S)_{fppf}$
是群胚叠。
\end{lemma}

\begin{proof}
为证明 $\Cohstack_{X/B}$ 是群胚叠，必须证明预层
$\mathit{Isom}$ 都是层，且下降数据有效。关于 $\mathit{Isom}$ 的断言
由引理 \ref{quot-lemma-coherent-diagonal} 得出；见《代数叠》引理
\ref{algebraic-lemma-representable-diagonal}。下面证明关于下降数据的断言。
设 $\{a_i : T_i \to T\}$ 是 $S$ 上概形的 fppf 覆盖，
而 $(\xi_i, \varphi_{ij})$ 是 $\{T_i \to T\}$ 关于
$\Cohstack_{X/B}$ 的下降数据。对每个 $i$，可写
$\xi_i = (T_i, g_i, \mathcal{F}_i)$。用
$\text{pr}_0 : T_i \times_T T_j \to T_i$ 和
$\text{pr}_1 : T_i \times_T T_j \to T_j$ 表示投影。
条件 $\xi_i|_{T_i \times_T T_j} = \xi_j|_{T_i \times_T T_j}$
特别蕴含 $g_i \circ \text{pr}_0 = g_j \circ \text{pr}_1$。
故存在唯一态射 $g : T \to B$ 使得 $g_i = g \circ a_i$；见
《空间上的下降》引理
\ref{spaces-descent-lemma-fpqc-universal-effective-epimorphisms}.
记 $X_T = T \times_{g, B} X$，并置
$X_i = X_{T_i} = T_i \times_{g_i, B} X = T_i \times_{a_i, T} X_T$
以及
$$
X_{ij} = X_{T_i} \times_{X_T} X_{T_j} = X_i \times_{X_T} X_j
$$
，其到 $X_i$ 和 $X_j$ 的投影分别记为 $\text{pr}_i$ 和
$\text{pr}_j$。注意，$(T_i, g_i, \mathcal{F}_i)$ 沿
$\text{pr}_0 : T_i \times_T T_j \to T_i$ 的拉回为
$(T_i \times_T T_j, g_i \circ \text{pr}_0, \text{pr}_i^*\mathcal{F}_i)$.
因此，$\Cohstack_{X/B}$ 中关于 $\{T_i \to T\}$ 的下降数据
由对象 $(T_i, g \circ a_i, \mathcal{F}_i)$ 以及对每一对 $i,j$
给出的 $\mathcal{O}_{X_{ij}}$-模同构
$$
\varphi_{ij} :
\text{pr}_i^*\mathcal{F}_i \longrightarrow \text{pr}_j^*\mathcal{F}_j
$$
组成；它们在 $T_i \times_T T_j \times_T T_k$（上面 $X$ 的拉回）
上满足上闭链条件。现在只需利用 $\{X_i \to X_T\}$ 是 fppf 覆盖，
从而可将 $(\mathcal{F}_i, \varphi_{ij})$ 视为该覆盖的下降数据。
由《空间上的下降》命题
\ref{spaces-descent-proposition-fpqc-descent-quasi-coherent}
，此下降数据有效，因而得到 $X_T$ 上的拟凝聚层 $\mathcal{F}$，
它在 $X_i$ 上的限制是 $\mathcal{F}_i$。
由《空间的态射》引理 \ref{spaces-morphisms-lemma-flat-permanence}，
$\mathcal{F}$ 在 $T$ 上平坦；《空间上的下降》引理
\ref{spaces-descent-lemma-finite-presentation-descends}
保证 $\mathcal{F}$ 作为 $\mathcal{O}_{X_T}$-模有限表示。最后，
由《空间上的下降》引理
\ref{spaces-descent-lemma-descending-property-proper}，
$\mathcal{F}$ 的概形论支撑在 $T$ 上真，因为已假设
$\mathcal{F}_i$ 的概形论支撑在 $T_i$ 上真（注意，由
《空间的态射》引理 \ref{spaces-morphisms-lemma-flat-pullback-support}，
取概形论支撑与平坦基变换可交换）。这样便得到所需的 $T$ 上对象。
\end{proof}

\begin{remark}
\label{quot-remark-coherent-base-change}
在情形 \ref{quot-situation-coherent} 中，规则
$(T, g, \mathcal{F}) \mapsto (T, g)$ 定义群胚叠的 $1$-态射
$$
\Cohstack_{X/B} \longrightarrow \mathcal{S}_B
$$
（见引理 \ref{quot-lemma-coherent-stack}、《代数叠》第
\ref{algebraic-section-split} 节和《叠的例子》第
\ref{examples-stacks-section-stack-associated-to-sheaf} 节）。
设 $B' \to B$ 是 $S$ 上代数空间的态射，并设
$\mathcal{S}_{B'} \to \mathcal{S}_B$ 是相应的集合纤维化叠的
$1$-态射。置 $X' = X \times_B B'$。与基变换
$f' : X' \to B'$ 相联系，得到群胚叠
$\Cohstack_{X'/B'} \to (\Sch/S)_{fppf}$。在此情形下，图
$$
\vcenter{
\xymatrix{
\Cohstack_{X'/B'} \ar[r] \ar[d] & \Cohstack_{X/B} \ar[d] \\
\mathcal{S}_{B'} \ar[r] & \mathcal{S}_B
}
}
\quad
\begin{matrix}
\text{或用} \\
\text{另一种} \\
\text{记号}
\end{matrix}
\quad
\vcenter{
\xymatrix{
\Cohstack_{X'/B'} \ar[r] \ar[d] & \Cohstack_{X/B} \ar[d] \\
\Sch/B' \ar[r] & \Sch/B
}
}
$$
% P11-E0369: the frozen source omits the article before "2-fibre product square".
是 $2$-纤维积方块。这个简单评注有时可用于更换基代数空间。
\end{remark}

\begin{lemma}
\label{quot-lemma-coherent-limits}
在情形 \ref{quot-situation-coherent} 中，假设 $B \to S$ 局部有限表示。
则 $p : \Cohstack_{X/B} \to (\Sch/S)_{fppf}$ 保极限
（《Artin 公理》定义 \ref{artin-definition-limit-preserving}）。
\end{lemma}

\begin{proof}
用 $B(T)$ 表示以 $S$-态射 $T \to B$ 为对象的离散范畴。
设 $T = \lim T_i$ 是 $S$ 上仿射概形的滤过极限。
把 $\Cohstack_{X/B, T}$ 的对象 $(T, h, \mathcal{F})$ 映到
$B(T)$ 的对象 $h$，得到纤维范畴的交换图
$$
\xymatrix{
\colim \Cohstack_{X/B, T_i} \ar[r] \ar[d] &
\Cohstack_{X/B, T} \ar[d] \\
\colim B(T_i) \ar[r] & B(T)
}
$$
必须证明上方水平箭头是等价。由于已假设 $B$ 在 $S$ 上局部有限表示，
由《空间的极限》评注 \ref{spaces-limits-remark-limit-preserving}，
下方水平箭头是等价。这意味着可以假设 $T = \lim T_i$ 是
$B$ 上仿射概形的滤过极限。用 $g_i : T_i \to B$ 和
$g : T \to B$ 表示相应态射。置
$X_i = T_i \times_{g_i, B} X$ 和 $X_T = T \times_{g, B} X$。
% P11-E0370: for inverse-limit base change, the frozen source writes X_T = \colim X_i.
注意，$X_T = \colim X_i$，并且代数空间 $X_i$ 和 $X_T$
拟分离且拟紧（因为它们分别在仿射概形 $T_i$ 和 $T$ 上有限表示）。
由《空间的极限》引理
\ref{spaces-limits-lemma-descend-modules-finite-presentation}
可知
$$
\colim \textit{FP}(X_i) = \textit{FP}(X_T).
$$
，其中 $\textit{FP}(W)$ 是有限表示 $\mathcal{O}_W$-模范畴的简写。
《空间的极限》引理 \ref{spaces-limits-lemma-descend-flat} 和
\ref{spaces-limits-lemma-eventually-proper-support}
表明：若将 $\textit{FP}(X_i)$ 和 $\textit{FP}(X_T)$ 分别换成
在 $T_i$ 和 $T$ 上平坦、且概形论支撑分别在 $T_i$ 和 $T$ 上真的
对象所成的全子范畴，同一结论仍成立。这证明了引理。
\end{proof}

\begin{lemma}
\label{quot-lemma-coherent-RS-star}
在情形 \ref{quot-situation-coherent} 中，设
$$
\xymatrix{
Z \ar[r] \ar[d] & Z' \ar[d] \\
Y \ar[r] & Y'
}
$$
是 $S$ 上概形范畴中的推出，其中 $Z \to Z'$ 是加厚，
$Z \to Y$ 是仿射态射；见《态射进阶》引理
\ref{more-morphisms-lemma-pushout-along-thickening}。则纤维范畴间的函子
$$
\Cohstack_{X/B, Y'}
\longrightarrow
\Cohstack_{X/B, Y} \times_{\Cohstack_{X/B, Z}} \Cohstack_{X/B, Z'}
$$
是等价。
\end{lemma}

\begin{proof}
注意，相应映射
$$
B(Y') \longrightarrow B(Y) \times_{B(Z)} B(Z')
$$
是双射；见《空间的推出》引理
\ref{spaces-pushouts-lemma-pushout-along-thickening-schemes}。
于是利用交换图
$$
\xymatrix{
\Cohstack_{X/B, Y'} \ar[r] \ar[d] &
\Cohstack_{X/B, Y} \times_{\Cohstack_{X/B, Z}} \Cohstack_{X/B, Z'}
\ar[d] \\
B(Y') \ar[r] & B(Y) \times_{B(Z)} B(Z')
}
$$
% P11-E0371: the frozen source introduces an undefined B' in "Y' is a scheme over B'".
可知可以假设 $Y'$ 是 $B'$ 上的概形。由评注
\ref{quot-remark-coherent-base-change}，可以用 $Y'$ 替换 $B$，
用 $X \times_B Y'$ 替换 $X$。故可假设 $B = Y'$。
此时断言由《空间的推出》引理
\ref{spaces-pushouts-lemma-space-over-pushout-flat-modules} 得出。
\end{proof}

\begin{lemma}
\label{quot-lemma-coherent-over-first-order-thickening}
设
$$
\xymatrix{
X \ar[d] \ar[r]_i & X' \ar[d] \\
T \ar[r] & T'
}
$$
是代数空间的笛卡儿方块，其中 $T \to T'$ 是一阶加厚。
设 $\mathcal{F}'$ 是在 $T'$ 上平坦的 $\mathcal{O}_{X'}$-模，
并置 $\mathcal{F} = i^*\mathcal{F}'$。下列条件等价：
\begin{enumerate}
\item $\mathcal{F}'$ 是有限表示的拟凝聚 $\mathcal{O}_{X'}$-模；
\item $\mathcal{F}'$ 是有限表示的 $\mathcal{O}_{X'}$-模；
\item $\mathcal{F}$ 是有限表示的拟凝聚 $\mathcal{O}_X$-模；
\item $\mathcal{F}$ 是有限表示的 $\mathcal{O}_X$-模。
\end{enumerate}
\end{lemma}

\begin{proof}
回忆有限表示模是拟凝聚的，故 (1) 与 (2) 等价，(3) 与 (4) 等价。
(2) 与 (4) 的等价性是《形变理论》引理
\ref{defos-lemma-deform-fp-module-ringed-topoi} 的特例。
\end{proof}

\begin{lemma}
\label{quot-lemma-coherent-tangent-space}
在情形 \ref{quot-situation-coherent} 中，假设 $S$ 是局部 Noether 概形，
且 $B \to S$ 局部有限表示。设 $k$ 是 $S$ 上有限型的域，并设
$x_0 = (\Spec(k), g_0, \mathcal{G}_0)$
是 $k$ 上 $\mathcal{X} = \Cohstack_{X/B}$ 的对象。则空间
$T\mathcal{F}_{\mathcal{X}, k, x_0}$ 和
$\text{Inf}(\mathcal{F}_{\mathcal{X}, k, x_0})$
（《Artin 公理》第 \ref{artin-section-tangent-spaces} 节）都是有限维的。
\end{lemma}

\begin{proof}
注意，由引理 \ref{quot-lemma-coherent-RS-star}，群胚叠 $\mathcal{X}$
满足《Artin 公理》第 \ref{artin-section-inf} 节定义的性质 (RS*)。
特别地，$\mathcal{X}$ 满足 (RS)。因此所有相联系的预形变范畴都是
形变范畴（《Artin 公理》引理
\ref{artin-lemma-deformation-category}），所述断言有意义。

\medskip\noindent
本段说明可归约到 $B = \Spec(k)$ 的情形。置
$X_0 = \Spec(k) \times_{g_0, B} X$，并记
$\mathcal{X}_0 = \Cohstack_{X_0/k}$。在评注
\ref{quot-remark-coherent-base-change} 中已经看到：作为
$(\Sch/S)_{fppf}$ 上的群胚纤维化范畴，$\mathcal{X}_0$ 是
$\mathcal{X}$ 与 $\Spec(k)$ 在 $B$ 上的 $2$-纤维积。因此由
《Artin 公理》引理 \ref{artin-lemma-fibre-product-tangent-spaces}，
问题归结为证明 $B$、$\Spec(k)$ 和 $\mathcal{X}_0$ 的切空间与
无穷小自同构空间有限维。由《Artin 公理》引理
\ref{artin-lemma-finite-dimension}，$B$ 和 $\Spec(k)$ 的切空间
有限维；当然，它们的 $\text{Inf}$ 为零。因此只需处理
$\mathcal{X}_0$。

\medskip\noindent
设 $k[\epsilon]$ 是 $k$ 上的对偶数。令
$\Spec(k[\epsilon]) \to B$ 为 $g_0 : \Spec(k) \to B$ 与由包含
$k \to k[\epsilon]$ 所诱导态射
$\Spec(k[\epsilon]) \to \Spec(k)$ 的复合。置
$X_0 = \Spec(k) \times_B X$ 和
$X_\epsilon = \Spec(k[\epsilon]) \times_B X$。注意，$X_\epsilon$
是 $X_0$ 的一阶加厚，并在一阶加厚
$\Spec(k) \to \Spec(k[\epsilon])$ 上平坦。展开定义并应用引理
\ref{quot-lemma-coherent-over-first-order-thickening}，可知
$T\mathcal{F}_{\mathcal{X}_0, k, x_0}$ 是将 $\mathcal{G}_0$
提升为 $X_\epsilon$ 上平坦模的所有方式所成的集合。由
《形变理论》引理 \ref{defos-lemma-flat-ringed-topoi}，得到
$$
T\mathcal{F}_{\mathcal{X}_0, k, x_0} =
\Ext^1_{\mathcal{O}_{X_0}}(\mathcal{G}_0, \mathcal{G}_0)
$$
这里使用了 $k[\epsilon]$-模的同一
$\epsilon k[\epsilon] \cong k$。再次应用《形变理论》引理
\ref{defos-lemma-flat-ringed-topoi}，可知
$$
\text{Inf}(\mathcal{F}_{\mathcal{X}, k, x_0}) =
\Ext^0_{\mathcal{O}_{X_0}}(\mathcal{G}_0, \mathcal{G}_0)
$$
% P11-E0372: after reducing to X_0, the frozen source changes the Inf subscript back from X_0 to X.
由于 $\mathcal{G}_0$ 的支撑在 $\Spec(k)$ 上真，这些空间在 $k$ 上
有限维。事实上，$X_0$ 在 $\Spec(k)$ 上有限表示，因而是 Noether 的。
由于 $\mathcal{G}_0$ 有限表示，它是凝聚 $\mathcal{O}_{X_0}$-模。
于是应用《空间的导出范畴》引理
\ref{spaces-perfect-lemma-ext-finite}，便得到所需的有限性。
\end{proof}

\begin{lemma}
\label{quot-lemma-coherent-existence}
在情形 \ref{quot-situation-coherent} 中，假设 $S$ 是局部 Noether 概形，
且 $f : X \to B$ 分离。令 $\mathcal{X} = \Cohstack_{X/B}$。
则《Artin 公理》公式 (\ref{artin-equation-approximation}) 中的函子
是等价。
\end{lemma}

\begin{proof}
设 $A$ 是 $S$-代数，同时也是极大理想为 $\mathfrak m$ 的完备局部
Noether 环，其剩余域 $k$ 在 $S$ 上有限型。必须证明 $A$ 上对象的
范畴等价于 $A$ 上形式对象的范畴。由《Artin 公理》引理
\ref{artin-lemma-effective}，这对与 $B$ 相联系的集合纤维化范畴
$\mathcal{S}_B$ 已经成立，故只需对位于给定态射
$\Spec(A) \to B$ 上方的对象证明此事。

\medskip\noindent
置 $X_A = \Spec(A) \times_B X$ 和
$X_n = \Spec(A/\mathfrak m^n) \times_B X$。由 Grothendieck 存在定理
（《空间的态射进阶》定理
\ref{spaces-more-morphisms-theorem-grothendieck-existence}），
$X_A$ 上支撑在 $\Spec(A)$ 上真的凝聚模 $\mathcal{F}$ 所成范畴，
等价于系统 $(\mathcal{F}_n)$ 所成的范畴；这里 $\mathcal{F}_n$
是 $X_n$ 上支撑在 $\Spec(A/\mathfrak m^n)$ 上真的凝聚模。
该等价把 $\mathcal{F}$ 映到系统
$(\mathcal{F} \otimes_A A/\mathfrak m^n)$。参见《空间的态射进阶》
评注 \ref{spaces-more-morphisms-remark-reformulate-existence-theorem}
中的讨论。为完成引理的证明，只需证明：$\mathcal{F}$ 在 $A$ 上
平坦，当且仅当所有 $\mathcal{F} \otimes_A A/\mathfrak m^n$
都在 $A/\mathfrak m^n$ 上平坦。这由《空间的态射进阶》引理
\ref{spaces-more-morphisms-lemma-flatness-over-Noetherian-ring} 得出。
\end{proof}

\begin{lemma}
\label{quot-lemma-coherent-defo-thy}
在情形 \ref{quot-situation-coherent} 中，假设 $S$ 是局部 Noether 概形，
$S = B$，且 $f : X \to B$ 平坦。令
$\mathcal{X} = \Cohstack_{X/B}$。则 $\mathcal{X}$ 满足通用性的开性
（见《Artin 公理》定义 \ref{artin-definition-openness-versality}）。
\end{lemma}

\begin{proof}[第一种证明]
本证明基于《Artin 公理》引理 \ref{artin-lemma-dual-openness} 的判据。
设 $U \to S$ 是概形的有限型态射，$x$ 是 $U$ 上 $\mathcal{X}$
的对象，$u_0 \in U$ 是有限型点，并且 $x$ 在 $u_0$ 处通用。
缩小 $U$ 后，可以假设 $u_0$ 是闭点（《态射》引理
\ref{morphisms-lemma-point-finite-type}），并且 $U = \Spec(A)$，
而 $U \to S$ 的像包含在 $S$ 的某个仿射开集
$\Spec(\Lambda)$ 中。设 $\mathcal{F}$ 是与给定对象 $x$ 相应的
$X_A = \Spec(A) \times_S X$ 上凝聚模，它在 $A$ 上平坦。

\medskip\noindent
由《形变理论》引理 \ref{defos-lemma-flat-ringed-topoi}，
有函子同构
$$
T_x(M) = \Ext^1_{X_A}(\mathcal{F}, \mathcal{F} \otimes_A M)
$$
并且，对任意核 $I$ 的平方为零的 $\Lambda$-代数满射
$A' \to A$，有障碍类
$$
\xi_{A'} \in \Ext^2_{X_A}(\mathcal{F}, \mathcal{F} \otimes_A I)
$$
这里用到了：对上述任意 $A' \to A$，基变换
$X_{A'} = \Spec(A') \times_B X$ 在 $A'$ 上平坦。此外，由
《形变理论》引理 \ref{defos-lemma-functorial-ringed-topoi}，
障碍类的构造关于满射 $A' \to A$（固定 $A$）具有函子性。
把《空间的导出范畴》引理
\ref{spaces-perfect-lemma-compute-ext} 应用于 Ext 群
$\Ext^i_{X_A}(\mathcal{F}, \mathcal{F} \otimes_A M)$
在 $i \leq m$、$m = 2$ 时的计算。得到完美对象 $K \in D(A)$
和函子同构
$$
H^i(K \otimes_A^\mathbf{L} M)
\longrightarrow
\Ext^i_{X_A}(\mathcal{F}, \mathcal{F} \otimes_A M)
$$
，它们对 $i \leq m$ 成立并与边缘映射相容。此对象 $K$ 与上述同一
一起给出《Artin 公理》情形 \ref{artin-situation-dual} 中的数据。
最后，《Artin 公理》引理
\ref{artin-lemma-dual-obstruction} 的条件 (iv) 由《形变理论》引理
\ref{defos-lemma-verify-iv-ringed-topoi} 得出。
所以确实可以应用《Artin 公理》引理
\ref{artin-lemma-dual-openness}，引理得证。
\end{proof}

\begin{proof}[第二种证明]
本证明基于《Artin 公理》引理
\ref{artin-lemma-get-openness-obstruction-theory}。该引理的条件
(1)、(2)、(3) 分别对应于引理 \ref{quot-lemma-coherent-diagonal}、
\ref{quot-lemma-coherent-RS-star} 和
\ref{quot-lemma-coherent-limits}。

\medskip\noindent
我们已在形变理论一章中构造了障碍理论。具体地，给定 $S$-代数
$A$，以及由 $X_A$ 上 $\mathcal{F}$ 给出的 $\Spec(A)$ 上
$\Cohstack_{X/B}$ 的对象 $x$，置
$\mathcal{O}_x(M) = \Ext^2_{X_A}(\mathcal{F}, \mathcal{F} \otimes_A M)$
；若 $A' \to A$ 是核为 $I$ 的满射，则取障碍元
$$
o_x(A') = o(\mathcal{F}, \mathcal{F} \otimes_A I, 1) \in
\mathcal{O}_x(I) = \Ext^2_{X_A}(\mathcal{F}, \mathcal{F} \otimes_A I)
$$
，即《形变理论》引理 \ref{defos-lemma-flat-ringed-topoi} 中的元素。
除该定义的条件 (ii) 所述障碍类函子性之外，《Artin 公理》定义
\ref{artin-definition-obstruction-theory} 中障碍理论的所有性质
均由该引理得出。不过，正如《Artin 公理》引理
\ref{artin-lemma-get-openness-obstruction-theory} 的假设 (4) 脚注所述，
只需对固定的 $A$ 验证障碍类的函子性；这由《形变理论》引理
\ref{defos-lemma-functorial-ringed-topoi} 得出。《形变理论》引理
\ref{defos-lemma-flat-ringed-topoi} 还表明，对任意 $A$-模 $M$，
$T_x(M) = \Ext^1_{X_A}(\mathcal{F}, \mathcal{F} \otimes_A M)$
。

\medskip\noindent
为完成证明，只需说明
$T_x(\prod M_n) = \prod T_x(M_n)$ 以及
$\mathcal{O}_x(\prod M_n) = \prod \mathcal{O}_x(M)$.
% P11-E0373: the frozen source drops the subscript n in the final product O_x(M).
把《空间的导出范畴》引理
\ref{spaces-perfect-lemma-compute-ext}
应用于 Ext 群
$\Ext^i_{X_A}(\mathcal{F}, \mathcal{F} \otimes_A M)$
在 $i \leq m$、$m = 2$ 时的计算。得到完美对象 $K \in D(A)$
和函子同构
$$
H^i(K \otimes_A^\mathbf{L} M)
\longrightarrow
\Ext^i_{X_A}(\mathcal{F}, \mathcal{F} \otimes_A M)
$$
，其中 $i = 1, 2$。直接论证表明
$$
H^i(K \otimes_A^\mathbf{L} \prod M_n) =
\prod H^i(K \otimes_A^\mathbf{L} M_n)
$$
，只要 $K$ 是 $D(A)$ 的伪凝聚对象。事实上，这一性质（对所有 $i$）
刻画伪凝聚复形；见《代数进阶》引理
\ref{more-algebra-lemma-pseudo-coherent-tensor}。
\end{proof}

\begin{theorem}[凝聚层叠的代数性：平坦情形]
\label{quot-theorem-coherent-algebraic}
设 $S$ 是概形，$f : X \to B$ 是 $S$ 上代数空间的态射。
假设 $f$ 有限表示、分离且平坦\footnote{这个假设并非必要；
见第 \ref{quot-section-not-flat} 节。}。则 $\Cohstack_{X/B}$ 是
$S$ 上的代数叠。
\end{theorem}

\begin{proof}
置 $\mathcal{X} = \Cohstack_{X/B}$。已经看到 $\mathcal{X}$ 是
$(\Sch/S)_{fppf}$ 上的群胚叠，其对角态射可由代数空间表示
（引理 \ref{quot-lemma-coherent-stack} 和
\ref{quot-lemma-coherent-diagonal}）。因此只需找到概形 $W$ 和满且光滑的
态射 $W \to \mathcal{X}$。

\medskip\noindent
设 $B'$ 是概形，且 $B' \to B$ 是满的 \'etale 态射。置
$X' = B' \times_B X$，并用 $f' : X' \to B'$ 表示投影。
则 $\mathcal{X}' = \Cohstack_{X'/B'}$ 等于 $\mathcal{X}$ 与
同 $B'$ 相联系的集合纤维化范畴，在同 $B$ 相联系的集合纤维化
范畴上的 $2$-纤维积（评注 \ref{quot-remark-coherent-base-change}）。
由《代数叠》第 \ref{algebraic-section-representable-properties} 节，
态射 $\mathcal{X}' \to \mathcal{X}$ 满且 \'etale。因此只需对
$\mathcal{X}'$ 证明结论。换言之，可以假设 $B$ 是概形。

\medskip\noindent
假设 $B$ 是概形。此时可以用 $B$ 替换 $S$；见《代数叠》第
\ref{algebraic-section-change-base-scheme} 节。故可假设 $S = B$。

\medskip\noindent
假设 $S = B$。选取仿射开覆盖 $S = \bigcup U_i$。用
$\mathcal{X}_i$ 表示 $\mathcal{X}$ 在 $(\Sch/U_i)_{fppf}$ 上的限制。
若能找到 $U_i$ 上的概形 $W_i$ 和满光滑态射
$W_i \to \mathcal{X}_i$，则置 $W = \coprod W_i$，便得到满光滑态射
$W \to \mathcal{X}$。因此可以假设 $S = B$ 是仿射的。

\medskip\noindent
假设 $S = B$ 是仿射的，记 $S = \Spec(\Lambda)$。把
$\Lambda = \colim \Lambda_i$ 写成滤过余极限，其中每个 $\Lambda_i$
都是 $\mathbf{Z}$ 上有限型的。对某个 $i$，可找到代数空间的态射
$X_i \to \Spec(\Lambda_i)$，它有限表示、分离且平坦，其到
$\Lambda$ 的基变换为 $X$。见《空间的极限》引理
\ref{spaces-limits-lemma-descend-finite-presentation}、
\ref{spaces-limits-lemma-descend-separated-morphism} 和
\ref{spaces-limits-lemma-descend-flat}。
若证明 $\Cohstack_{X_i/\Spec(\Lambda_i)}$ 是代数叠，则由基变换
（评注 \ref{quot-remark-coherent-base-change} 和《代数叠》第
\ref{algebraic-section-change-base-scheme} 节）可知
$\mathcal{X}$ 是代数叠。因此可以假设 $\Lambda$ 是有限型
$\mathbf{Z}$-代数。

\medskip\noindent
假设 $S = B = \Spec(\Lambda)$ 是 $\mathbf{Z}$ 上有限型的仿射概形。
此时验证《Artin 公理》引理
\ref{artin-lemma-diagonal-representable} 的条件 (1)、(2)、(3)、(4)、
(5)，即可断定 $\mathcal{X}$ 是代数叠。注意，$\Lambda$ 是 G-环；
见《代数进阶》命题 \ref{more-algebra-proposition-ubiquity-G-ring}。
因此 $S$ 的所有局部环都是 G-环，故 (5) 成立。由引理
\ref{quot-lemma-coherent-defo-thy}，$\mathcal{X}$ 满足通用性的开性，
所以 (4) 成立。为验证 (2)，必须验证《Artin 公理》第
\ref{artin-section-axioms} 节的公理 [-1]、[0]、[1]、[2]、[3]。
略去 [-1] 的验证；公理 [0]、[1]、[2]、[3] 分别对应于引理
\ref{quot-lemma-coherent-stack}、
\ref{quot-lemma-coherent-limits},
\ref{quot-lemma-coherent-RS-star},
\ref{quot-lemma-coherent-tangent-space}.
条件 (3) 由引理 \ref{quot-lemma-coherent-existence} 得出。
最后，条件 (1) 即引理 \ref{quot-lemma-coherent-diagonal}。
这就完成了定理的证明。
\end{proof}








\section{非平坦情形下的凝聚层叠}
\label{quot-section-not-flat}

\noindent
定理 \ref{quot-theorem-coherent-algebraic} 中关于 $f : X \to B$ 平坦的
假设并非必要。本节给出另一种证明：利用《空间上的平坦性》第
\ref{spaces-flat-section-existence} 节和《Artin 公理》第
\ref{artin-section-strong-formal-effectiveness} 节的结果，
既避开平坦性假设，也无需验证通用性的开性。

\medskip\noindent
关于此问题的另一种处理方法，读者可参阅 \cite{ArtinI}，并沿用论文
\cite{olsson-starr}、\cite{lieblich_remarks}、\cite{olsson_proper}、
\cite{Hall-Rydh}、\cite{Hall-Rydh-Hilbert}、\cite{rydh_representability}
中讨论的方法。其中一些论文处理代数叠上代数叠的凝聚层叠这一更一般的
情形，另一些则处理 Hilbert 叠或 Quot 函子情形下的类似问题。
我们的策略是从凝聚层叠的代数性推出某些 Hilbert 叠和 Quot 函子的
代数性。

\begin{theorem}[凝聚层叠的代数性：一般情形]
\label{quot-theorem-coherent-algebraic-general}
设 $S$ 是概形，$f : X \to B$ 是 $S$ 上代数空间的态射。
% P11-E0374: the frozen source omits the article in "be morphism".
假设 $f$ 有限表示且分离。则 $\Cohstack_{X/B}$ 是 $S$ 上的代数叠。
\end{theorem}

\begin{proof}
只有证明的最后一步与平坦情形不同，不过这里重述全部论证，
以确保每一步都成立。

\medskip\noindent
置 $\mathcal{X} = \Cohstack_{X/B}$。已经看到 $\mathcal{X}$ 是
$(\Sch/S)_{fppf}$ 上的群胚叠，其对角态射可由代数空间表示
（引理 \ref{quot-lemma-coherent-stack} 和
\ref{quot-lemma-coherent-diagonal}）。因此只需找到概形 $W$ 和满光滑态射
$W \to \mathcal{X}$。

\medskip\noindent
设 $B'$ 是概形，且 $B' \to B$ 是满的 \'etale 态射。置
$X' = B' \times_B X$，并用 $f' : X' \to B'$ 表示投影。
则 $\mathcal{X}' = \Cohstack_{X'/B'}$ 等于 $\mathcal{X}$ 与
同 $B'$ 相联系的集合纤维化范畴，在同 $B$ 相联系的集合纤维化
范畴上的 $2$-纤维积（评注 \ref{quot-remark-coherent-base-change}）。
由《代数叠》第 \ref{algebraic-section-representable-properties} 节，
态射 $\mathcal{X}' \to \mathcal{X}$ 满且 \'etale。因此只需对
$\mathcal{X}'$ 证明结论。换言之，可以假设 $B$ 是概形。

\medskip\noindent
假设 $B$ 是概形。此时可以用 $B$ 替换 $S$；见《代数叠》第
\ref{algebraic-section-change-base-scheme} 节。故可假设 $S = B$。

\medskip\noindent
假设 $S = B$。选取仿射开覆盖 $S = \bigcup U_i$。用
$\mathcal{X}_i$ 表示 $\mathcal{X}$ 在 $(\Sch/U_i)_{fppf}$ 上的限制。
若能找到 $U_i$ 上的概形 $W_i$ 和满光滑态射
$W_i \to \mathcal{X}_i$，则置 $W = \coprod W_i$，便得到满光滑态射
$W \to \mathcal{X}$。因此可以假设 $S = B$ 是仿射的。

\medskip\noindent
假设 $S = B$ 是仿射的，记 $S = \Spec(\Lambda)$。把
$\Lambda = \colim \Lambda_i$ 写成滤过余极限，其中每个 $\Lambda_i$
都是 $\mathbf{Z}$ 上有限型的。对某个 $i$，可找到分离且有限表示的
代数空间态射 $X_i \to \Spec(\Lambda_i)$，其到 $\Lambda$ 的基变换
为 $X$。见《空间的极限》引理
\ref{spaces-limits-lemma-descend-finite-presentation} 和
\ref{spaces-limits-lemma-descend-separated-morphism}。
若证明 $\Cohstack_{X_i/\Spec(\Lambda_i)}$ 是代数叠，则由基变换
（评注 \ref{quot-remark-coherent-base-change} 和《代数叠》第
\ref{algebraic-section-change-base-scheme} 节）可知
$\mathcal{X}$ 是代数叠。因此可以假设 $\Lambda$ 是有限型
$\mathbf{Z}$-代数。

\medskip\noindent
假设 $S = B = \Spec(\Lambda)$ 是 $\mathbf{Z}$ 上有限型的仿射概形。
此时验证《Artin 公理》引理
\ref{artin-lemma-diagonal-representable} 的条件 (1)、(2)、(3)、(4)、
(5)，即可断定 $\mathcal{X}$ 是代数叠。注意，$\Lambda$ 是 G-环；
见《代数进阶》命题 \ref{more-algebra-proposition-ubiquity-G-ring}。
因此 $S$ 的所有局部环都是 G-环，故 (5) 成立。为验证 (2)，
必须验证《Artin 公理》第 \ref{artin-section-axioms} 节的公理
[-1]、[0]、[1]、[2]、[3]。略去 [-1] 的验证；公理
[0]、[1]、[2]、[3] 分别对应于引理 \ref{quot-lemma-coherent-stack}、
\ref{quot-lemma-coherent-limits},
\ref{quot-lemma-coherent-RS-star},
\ref{quot-lemma-coherent-tangent-space}.
条件 (3) 即引理 \ref{quot-lemma-coherent-existence}。
条件 (1) 即引理 \ref{quot-lemma-coherent-diagonal}。

\medskip\noindent
还需证明条件 (4)，即通用性的开性。为此将使用《Artin 公理》引理
\ref{artin-lemma-SGE-implies-openness-versality}。已经看到
$\mathcal{X}$ 的对角态射可由代数空间表示，并且它满足
(RS*) 且保极限（见上面用到的引理）。因此只需证明
$\mathcal{X}$ 满足《Artin 公理》引理
\ref{artin-lemma-SGE-implies-openness-versality} 中所述的强形式有效性。
这正是《空间上的平坦性》定理 \ref{spaces-flat-theorem-existence}，
证明完毕。
\end{proof}









\section{商函子}
\label{quot-section-functor-quotients}

\noindent
本节讨论下面定义的函子 $Q_{\mathcal{F}/X/B}$ 的若干一般性质。
记号 $\Quotfunctor_{\mathcal{F}/X/B}$ 留给
$\text{Q}_{\mathcal{F}/X/B}$ 的一个子函子。建议初次阅读时跳过本节。

\begin{situation}
\label{quot-situation-q}
设 $S$ 是概形，$f : X \to B$ 是 $S$ 上代数空间的态射，
$\mathcal{F}$ 是拟凝聚 $\mathcal{O}_X$-模。对任意 $B$ 上概形
$T$，用 $X_T$ 表示 $X$ 到 $T$ 的基变换，用 $\mathcal{F}_T$
表示 $\mathcal{F}$ 沿投影态射 $X_T = X \times_B T \to X$ 的拉回。
对这样的 $T$，置
$$
\text{Q}_{\mathcal{F}/X/B}(T) =
\left\{
\begin{matrix}
\mathcal{F}_T \to \mathcal{Q}\text{ 的商，其中 }\mathcal{Q}
\text{ 是在 }T\text{ 上平坦的}\\
\text{拟凝聚 }\mathcal{O}_{X_T}\text{-模}
\end{matrix}
\right\}
$$
若两个商具有相同的核，便将它们视为同一。设 $T' \to T$ 是
$B$ 上概形的态射，且 $\mathcal{F}_T \to \mathcal{Q}$ 是
$\text{Q}_{\mathcal{F}/X/B}(T)$ 的元素。由《空间的态射》引理
\ref{spaces-morphisms-lemma-base-change-module-flat}，拉回
$\mathcal{Q}' = (X_{T'} \to X_T)^*\mathcal{Q}$ 是在 $T'$ 上平坦的
拟凝聚 $\mathcal{O}_{X_{T'}}$-模。因此得到函子
\begin{equation}
\label{quot-equation-q}
\text{Q}_{\mathcal{F}/X/B} : (\Sch/B)^{opp} \longrightarrow \textit{Sets}
\end{equation}
这称为{\it $\mathcal{F}/X/B$ 的商函子}。再定义子函子
\begin{equation}
\label{quot-equation-q-fp}
\text{Q}^{fp}_{\mathcal{F}/X/B} : (\Sch/B)^{opp} \longrightarrow \textit{Sets}
\end{equation}
：它把 $T$ 映到 $\text{Q}_{\mathcal{F}/X/B}(T)$ 中由如下商
$\mathcal{F}_T \to \mathcal{Q}$ 组成的子集，即 $\mathcal{Q}$
作为 $\mathcal{O}_{X_T}$-模有限表示。由《空间的性质》第
\ref{spaces-properties-section-properties-modules} 节，这是子函子。
\end{situation}

\noindent
在情形 \ref{quot-situation-q} 中，有时把
$\text{Q}_{\mathcal{F}/X/B}$ 看作函子
% P11-E0375: the frozen source changes the denominator from B to S in both Q_{F/X/S} occurrences.
$(\Sch/S)^{opp} \to \textit{Sets}$，它配有态射
$\text{Q}_{\mathcal{F}/X/S} \to B$。具体地，若 $T$ 是
$S$ 上的概形，则 $\text{Q}_{\mathcal{F}/X/B}(T)$ 的一个元素是
一对 $(h, \mathcal{Q})$，其中
% P11-E0377: the frozen source omits "is" in "where h a morphism".
$h$ 是态射 $h : T \to B$，而 $\mathcal{Q}$ 是 $X_T =
X \times_{B, h} T$ 上有限表示的、在 $T$ 上平坦的商
% P11-E0376: finite presentation belongs only to Q^fp, but the frozen prose imposes it on general Q here.
$\mathcal{F}_T \to \mathcal{Q}$。特别地，当我们说
$\text{Q}_{\mathcal{F}/X/S}$ 是代数空间时，是指相应函子
$(\Sch/S)^{opp} \to \textit{Sets}$ 是代数空间。类似评注也适用于
$\text{Q}^{fp}_{\mathcal{F}/X/B}$。

\begin{remark}
\label{quot-remark-q-base-change}
在情形 \ref{quot-situation-q} 中，设 $B' \to B$ 是 $S$ 上代数空间的态射。
置 $X' = X \times_B B'$，并用 $\mathcal{F}'$ 表示 $\mathcal{F}$
到 $X'$ 的拉回。于是有 $B'$ 上概形范畴上的函子
$Q_{\mathcal{F}'/X'/B'}$。对 $B'$ 上概形 $T$，显然有
$$
Q_{\mathcal{F}'/X'/B'}(T) = Q_{\mathcal{F}/X/B}(T)
$$
，其中右边把 $T$ 通过复合 $T \to B' \to B$ 看作 $B$ 上的概形。
类似评注也适用于 $\text{Q}^{fp}_{\mathcal{F}/X/B}$。
这些简单评注有时可用于更换基代数空间。
\end{remark}

\begin{remark}
\label{quot-remark-q-sheaf}
设 $S$ 是概形，$X$ 是 $S$ 上的代数空间，$\mathcal{F}$ 是拟凝聚
$\mathcal{O}_X$-模。假设
$\{f_i : X_i \to X\}_{i \in I}$
是 fpqc 覆盖，并且对每个 $i, j \in I$，给定 fpqc 覆盖
$\{X_{ijk} \to X_i \times_X X_j\}$。此时有双射
$$
\left\{
\begin{matrix}
\mathcal{F} \to \mathcal{Q}\text{ 的商，其中 }\\
\mathcal{Q}\text{ 是拟凝聚的}\\
\end{matrix}
\right\}
\longrightarrow
\left\{
\begin{matrix}
\text{商族 }f_i^*\mathcal{F} \to \mathcal{Q}_i
\text{，其中 }\\
\mathcal{Q}_i\text{ 是拟凝聚的，且 }
\mathcal{Q}_i\text{ 与 }\mathcal{Q}_j\\
\text{在 }X_{ijk}\text{ 上限制为同一个商}
\end{matrix}
\right\}
$$
具体地，设 $(f_i^*\mathcal{F} \to \mathcal{Q}_i)_{i \in I}$ 是右边的
元素。由于 $\{X_{ijk} \to X_i \times_X X_j\}$ 是 fpqc 覆盖，
可知 $\mathcal{Q}_i$ 和 $\mathcal{Q}_j$ 的拉回，限制为
$\mathcal{F}$ 到 $X_i \times_X X_j$ 的拉回的同一个商（由
% P11-E0378: the frozen source says "fully faithfulness" instead of "full faithfulness".
《空间上的下降》命题
\ref{spaces-descent-proposition-fpqc-descent-quasi-coherent}
中的全忠实性）。因此得到拟凝聚模关于
$\{X_i \to X\}_{i \in I}$ 的下降数据。由《空间上的下降》命题
\ref{spaces-descent-proposition-fpqc-descent-quasi-coherent}
，可找到拟凝聚 $\mathcal{O}_X$-模的映射
$\mathcal{F} \to \mathcal{Q}$，其在 $X_i$ 上的限制恢复给定映射
$f_i^*\mathcal{F} \to \mathcal{Q}_i$。由于态射族
$\{X_i \to X\}$ 联合满且平坦，对每个点 $x \in |X|$，存在某个
$i$ 和点 $x_i \in |X_i|$ 映到 $x$。注意，所诱导的局部环映射
$\mathcal{O}_{X, \overline{x}} \to \mathcal{O}_{X_i, \overline{x_i}}$
忠实平坦；见《空间的态射》第
\ref{spaces-morphisms-section-flat} 节。因此
$\mathcal{F} \to \mathcal{Q}$ 是满射。
\end{remark}

\begin{lemma}
\label{quot-lemma-q-sheaf}
在情形 \ref{quot-situation-q} 中，函子
% P11-E0379: the frozen source punctuates "In Situation ..." as a sentence fragment.
$\text{Q}_{\mathcal{F}/X/B}$ 和
$\text{Q}^{fp}_{\mathcal{F}/X/B}$
满足 fpqc 拓扑的层条件。
\end{lemma}

\begin{proof}
设 $\{T_i \to T\}_{i \in I}$ 是 $S$ 上概形的 fpqc 覆盖。
% P11-E0380: the frozen proof uses schemes over S and X ×_S T_i, although Q is defined on Sch/B with X_T = X ×_B T.
置 $X_i = X_{T_i} = X \times_S T_i$ 和
$\mathcal{F}_i = \mathcal{F}_{T_i}$。注意，
$\{X_i \to X_T\}_{i \in I}$ 是 $X_T$ 的 fpqc 覆盖
（《空间上的拓扑》引理 \ref{spaces-topologies-lemma-fpqc}），并且
$X_{T_i \times_T T_{i'}} = X_i \times_{X_T} X_{i'}$。
设 $\mathcal{F}_i \to \mathcal{Q}_i$ 是
$\text{Q}_{\mathcal{F}/X/B}(T_i)$ 的一族元素，并且
$\mathcal{Q}_i$ 和 $\mathcal{Q}_{i'}$ 在
$\text{Q}_{\mathcal{F}/X/B}(T_i \times_T T_{i'})$ 中限制为同一元素。
由评注 \ref{quot-remark-q-sheaf}，得到拟凝聚
$\mathcal{O}_{X_T}$-模的满射 $\mathcal{F}_T \to \mathcal{Q}$，
其在 $X_i$ 上的限制恢复给定的商。由《空间的态射》引理
\ref{spaces-morphisms-lemma-flat-permanence}，$\mathcal{Q}$ 在
$T$ 上平坦。最后，在 $\text{Q}^{fp}_{\mathcal{F}/X/B}$ 的情形，
即若 $\mathcal{Q}_i$ 有限表示，则《空间上的下降》引理
\ref{spaces-descent-lemma-finite-presentation-descends}
保证 $\mathcal{Q}$ 作为 $\mathcal{O}_{X_T}$-模有限表示。
\end{proof}

\noindent
一致性检验：$\text{Q}_{\mathcal{F}/X/B}$、
$\text{Q}^{fp}_{\mathcal{F}/X/B}$
在 $S$ 上代数空间之间起着同样的作用。

\begin{lemma}
\label{quot-lemma-extend-q-to-spaces}
在情形 \ref{quot-situation-q} 中，设 $T$ 是 $S$ 上的代数空间。
% P11-E0381: the frozen source punctuates "In Situation ..." as a sentence fragment.
则
$$
\Mor_{\Sh((\Sch/S)_{fppf})}(T,  \text{Q}_{\mathcal{F}/X/B}) =
\left\{
\begin{matrix}
(h, \mathcal{F}_T \to \mathcal{Q})\text{，其中 }
h : T \to B\text{，且}\\
\mathcal{Q}\text{ 是拟凝聚的并在 T 上平坦}
\end{matrix}
\right\}
$$
，其中 $\mathcal{F}_T$ 表示 $\mathcal{F}$ 到代数空间
$X \times_{B, h} T$ 的拉回。类似地，有
$$
\Mor_{\Sh((\Sch/S)_{fppf})}(T,  \text{Q}^{fp}_{\mathcal{F}/X/B}) =
\left\{
\begin{matrix}
(h, \mathcal{F}_T \to \mathcal{Q})\text{，其中 }
h : T \to B\text{，且}\\
\mathcal{Q}\text{ 有限表示并在 T 上平坦}
\end{matrix}
\right\}
$$
\end{lemma}

\begin{proof}
选取概形 $U$ 和满的 \'etale 态射 $p : U \to T$。令
$R = U \times_T U$，其投影为 $t, s : R \to U$。

\medskip\noindent
设 $v : T \to \text{Q}_{\mathcal{F}/X/B}$ 是自然变换。则
$v(p)$ 对应于 $U$ 上的一对
$(h_U, \mathcal{F}_U \to \mathcal{Q}_U)$。由于 $v$ 是函子间的变换，
$(h_U, \mathcal{F}_U \to \mathcal{Q}_U)$ 沿 $s$ 和 $t$ 的拉回相同。
由于 $T = U/R$（《空间》引理
\ref{spaces-lemma-space-presentation}），得到态射 $h : T \to B$
满足 $h_U = h \circ p$。由《空间上的下降》命题
\ref{spaces-descent-proposition-fpqc-descent-quasi-coherent}
，商 $\mathcal{Q}_U$ 下降为 $X_T$ 上的商
$\mathcal{F}_T \to \mathcal{Q}$。由于 $U \to T$ 满且平坦，
由《空间的态射》引理 \ref{spaces-morphisms-lemma-flat-permanence}，
$\mathcal{Q}$ 在 $T$ 上平坦。

\medskip\noindent
反之，设 $(h, \mathcal{F}_T \to \mathcal{Q})$ 是 $T$ 上的一对。
把态射 $a : T' \to T$（其中 $T'$ 是概形）映到
$(h \circ a, \mathcal{F}_{T'} \to a^*\mathcal{Q})$，便得到自然变换
$v : T \to \text{Q}_{\mathcal{F}/X/B}$。略去验证本段与前段的构造
互为逆构造。

\medskip\noindent
在 $\text{Q}^{fp}_{\mathcal{F}/X/B}$ 的情形还需补充：给定态射
$h : T \to B$，$X_T$ 上的拟凝聚层作为 $\mathcal{O}_{X_T}$-模
有限表示，当且仅当其到 $X_U$ 的拉回作为 $\mathcal{O}_{X_U}$-模
有限表示。这由 $X_U \to X_T$ 满且 \'etale 以及《空间上的下降》引理
\ref{spaces-descent-lemma-finite-presentation-descends} 得出。
\end{proof}

\begin{lemma}
\label{quot-lemma-q-sheaf-in-X}
在情形 \ref{quot-situation-q} 中，设 $\{X_i \to X\}_{i \in I}$ 是 fpqc
覆盖，并且对每个 $i, j \in I$，设
$\{X_{ijk} \to X_i \times_X X_j\}$ 是 fpqc 覆盖。用
$\mathcal{F}_i$、相应地 $\mathcal{F}_{ijk}$ 表示 $\mathcal{F}$
到 $X_i$、相应地到 $X_{ijk}$ 的拉回。对每个 $B$ 上概形 $T$，图
$$
\xymatrix{
Q_{\mathcal{F}/X/B}(T) \ar[r] &
\prod\nolimits_i
Q_{\mathcal{F}_i/X_i/B}(T)
\ar@<1ex>[r]^-{\text{pr}_0^*} \ar@<-1ex>[r]_-{\text{pr}_1^*}
&
\prod\nolimits_{i, j, k}
Q_{\mathcal{F}_{ijk}/X_{ijk}/B}(T)
}
$$
把第一条箭头表示为后两条箭头的等化子。对函子
$\text{Q}^{fp}_{\mathcal{F}/X/B}$，同一结论成立。
\end{lemma}

\begin{proof}
设 $\mathcal{F}_{i, T} \to \mathcal{Q}_i$ 是
$\text{pr}_0^*$ 与 $\text{pr}_1^*$ 的等化子中的元素。
由评注 \ref{quot-remark-q-sheaf}，得到拟凝聚 $\mathcal{O}_{X_T}$-模的
满射 $\mathcal{F}_T \to \mathcal{Q}$，其在 $X_{i, T}$ 上的限制
恢复 $\mathcal{F}_i \to \mathcal{Q}_i$。由《空间的态射》引理
\ref{spaces-morphisms-lemma-flat-permanence}，$\mathcal{Q}$ 如所需
在 $T$ 上平坦。在函子 $\text{Q}^{fp}_{\mathcal{F}/X/B}$ 的情形，
即若 $\mathcal{Q}_i$ 有限表示，则 $\mathcal{Q}$ 也有限表示；
这由《空间上的下降》引理
\ref{spaces-descent-lemma-finite-presentation-descends} 得出。
\end{proof}

\begin{lemma}
\label{quot-lemma-q-limit-preserving}
在情形 \ref{quot-situation-q} 中，进一步假设
(a) $f$ 拟紧且拟分离；并且
(b) $\mathcal{F}$ 有限表示。
则函子 $\text{Q}^{fp}_{\mathcal{F}/X/B}$ 在如下意义下保极限：
若 $T = \lim T_i$ 是 $B$ 上仿射概形的有向极限，则
$\text{Q}^{fp}_{\mathcal{F}/X/B}(T) =
\colim \text{Q}^{fp}_{\mathcal{F}/X/B}(T_i)$.
\end{lemma}

\begin{proof}
设 $T = \lim T_i$ 如引理所述。选取 $i_0 \in I$，并用
$\{i \in I \mid i \geq i_0\}$ 替换 $I$。可以置
$B = S = T_{i_0}$，用 $X_{T_0}$ 替换 $X$，并用到 $X_{T_0}$ 的
% P11-E0382: after choosing i_0, the frozen source uses undefined T_0/X_{T_0} instead of T_{i_0}/X_{T_{i_0}}.
拉回替换 $\mathcal{F}$。于是 $X_T = \lim X_{T_i}$；见
《空间的极限》引理
\ref{spaces-limits-lemma-directed-inverse-system-has-limit}.
设 $\mathcal{F}_T \to \mathcal{Q}$ 是
$\text{Q}^{fp}_{\mathcal{F}/X/B}(T)$ 的元素。由《空间的极限》引理
\ref{spaces-limits-lemma-descend-modules-finite-presentation}
，存在某个 $i$ 以及有限表示 $\mathcal{O}_{X_{T_i}}$-模的映射
$\mathcal{F}_{T_i} \to \mathcal{Q}_i$，其到 $X_T$ 的拉回就是给定
商映射。

\medskip\noindent
还需验证：适当增大 $i$ 后，映射
$\mathcal{F}_{T_i} \to \mathcal{Q}_i$ 满，并且 $\mathcal{Q}_i$
在 $T_i$ 上平坦。为此，选取仿射概形 $U$ 和满的 \'etale 态射
$U \to X$（见《空间的性质》引理
\ref{spaces-properties-lemma-quasi-compact-affine-cover}）。
按定义，可以在拉回到 \'etale 覆盖
$U_{T_i} \to X_{T_i}$ 后检验满性和在 $T_i$ 上的平坦性。
这把问题归约到如下情形：$X = \Spec(B_0)$ 是
$B = S = T_0 = \Spec(A_0)$ 上有限表示的仿射概形。写
$T_i = \Spec(A_i)$，则 $T = \Spec(A)$，其中 $A = \colim A_i$。
于是得到如下代数问题。设 $M_i \to N_i$ 是有限表示
$B_0 \otimes_{A_0} A_i$-模的映射，使得
$M_i \otimes_{A_i} A \to N_i \otimes_{A_i} A$ 满，并且
$N_i \otimes_{A_i} A$ 在 $A$ 上平坦。证明：对某个 $i' \geq i$，
$M_i \otimes_{A_i} A_{i'} \to N_i \otimes_{A_i} A_{i'}$ 满，且
% P11-E0383: the descended module is over A_{i'}, but the frozen conclusion says flat over A.
$N_i \otimes_{A_i} A_{i'}$ 在 $A$ 上平坦。前者由《代数》引理
\ref{algebra-lemma-module-map-property-in-colimit} 得出，后者由
《代数》引理 \ref{algebra-lemma-flat-finite-presentation-limit-flat}
得出。
\end{proof}

\begin{lemma}
\label{quot-lemma-q-RS-star}
在情形 \ref{quot-situation-q} 中，设
% P11-E0384: the frozen source punctuates "In Situation ..." as a sentence fragment.
$$
\xymatrix{
Z \ar[r] \ar[d] & Z' \ar[d] \\
Y \ar[r] & Y'
}
$$
是 $B$ 上概形范畴中的推出，其中 $Z \to Z'$ 是加厚，$Z \to Y$
是仿射态射；见《态射进阶》引理
\ref{more-morphisms-lemma-pushout-along-thickening}。则自然映射
$$
Q_{\mathcal{F}/X/B}(Y') \longrightarrow
Q_{\mathcal{F}/X/B}(Y) \times_{Q_{\mathcal{F}/X/B}(Z)} Q_{\mathcal{F}/X/B}(Z')
$$
是双射。若 $X \to B$ 局部有限表示，则对
$Q^{fp}_{\mathcal{F}/X/B}$ 同一结论成立。
\end{lemma}

\begin{proof}
下面构造逆映射。假设有 $\mathcal{F}_Y \to \mathcal{A}$、
$\mathcal{F}_{Z'} \to \mathcal{B}'$ 以及同构
$\mathcal{A}|_{X_Z} \to \mathcal{B}'|_{X_Z}$
，它与给定满射相容。应用《空间的推出》引理
\ref{spaces-pushouts-lemma-space-over-pushout-flat-modules}
，得到 $X_{Y'}$ 上在 $Y'$ 上平坦的拟凝聚模 $\mathcal{A}'$。
由于该层构造为纤维积（见所引引理的证明），有典范映射
$\mathcal{F}_{Y'} \to \mathcal{A}'$。该映射为满射，因为它分解为
$$
\begin{matrix}
\mathcal{F}_{Y'} \\
\downarrow \\
(X_Y \to X_{Y'})_*\mathcal{F}_Y
\times_{(X_Z \to X_{Y'})_*\mathcal{F}_Z}
(X_{Z'} \to X_{Y'})_*\mathcal{F}_{Z'} \\
\downarrow \\
\mathcal{A}' =
(X_Y \to X_{Y'})_*\mathcal{A}
\times_{(X_Z \to X_{Y'})_*\mathcal{A}|_{X_Z}}
(X_{Z'} \to X_{Y'})_*\mathcal{B}'
\end{matrix}
$$
，其中第一条箭头由《代数进阶》引理
\ref{more-algebra-lemma-module-over-fibre-product-bis} 为满射，
第二条箭头由《代数进阶》引理
\ref{more-algebra-lemma-surjection-module-over-fibre-product} 为满射。

\medskip\noindent
在 $Q^{fp}_{\mathcal{F}/X/B}$ 的情形，只需证明上述构造产生有限表示
模。这在交换代数语境下由《代数进阶》评注
\ref{more-algebra-remark-relative-modules-over-fibre-product}
说明。代数空间上模的当前情形由此经 \'etale 局部化得出。
\end{proof}

\begin{remark}[商的障碍]
\label{quot-remark-q-obs}
在情形 \ref{quot-situation-q} 中，{\bf 假设} $\mathcal{F}$ 在 $B$ 上平坦。
设 $T \subset T'$ 是 $B$ 上概形的一阶加厚，其理想层为
% P11-E0385: the frozen source says "an first order thickening".
$\mathcal{J}$。则 $X_T \subset X_{T'}$ 是代数空间的一阶加厚，
其理想层 $\mathcal{I}$ 是 $f_T^*\mathcal{J}$ 的商。利用《空间的态射
进阶》第 \ref{spaces-more-morphisms-section-thickenings} 节所述的
基本等价，把 $X_{T'}$ 上、相应地 $T'$ 上的层看作 $X_T$ 上、
相应地 $T$ 上的层。设
$$
0 \to \mathcal{K} \to \mathcal{F}_T \to \mathcal{Q} \to 0
$$
定义 $Q_{\mathcal{F}/X/B}(T)$ 的元素 $x$。由于 $\mathcal{F}_{T'}$
在 $T'$ 上平坦，有短正合列
$$
0 \to f_T^*\mathcal{J} \otimes_{\mathcal{O}_{X_T}} \mathcal{F}_T
\xrightarrow{i} \mathcal{F}_{T'} \xrightarrow{\pi} \mathcal{F}_T \to 0
$$
，并且有
$f_T^*\mathcal{J} \otimes_{\mathcal{O}_{X_T}} \mathcal{F}_T =
\mathcal{I} \otimes_{\mathcal{O}_{X_T}} \mathcal{F}_T$；见《形变理论》
引理 \ref{defos-lemma-deform-module-ringed-topoi}。对
$\mathcal{O}_{X_T}$-模 $\mathcal{G}$，采用简写
$
f_T^*\mathcal{J} \otimes_{\mathcal{O}_{X_T}} \mathcal{G} =
\mathcal{G} \otimes_{\mathcal{O}_T} \mathcal{J}
$
。由于 $\mathcal{Q}$ 在 $T$ 上平坦，得到短正合列
$$
0 \to
\mathcal{K} \otimes_{\mathcal{O}_T} \mathcal{J} \to
\mathcal{F}_T \otimes_{\mathcal{O}_T} \mathcal{J} \to
\mathcal{Q} \otimes_{\mathcal{O}_T} \mathcal{J} \to
\to 0
$$
% P11-E0386: the frozen short exact sequence contains a duplicated \to before 0.
结合以上结果，得到典范扩张
% P11-E0387: the frozen source says "an canonical extension".
$$
0 \to \mathcal{Q} \otimes_{\mathcal{O}_T} \mathcal{J} \to
\pi^{-1}(\mathcal{K})/i(\mathcal{K} \otimes_{\mathcal{O}_T} \mathcal{J}) \to
\mathcal{K} \to 0
$$
，它是 $\mathcal{O}_{X_T}$-模的扩张。这定义典范类
$$
o_x(T') \in
\Ext^1_{\mathcal{O}_{X_T}}(\mathcal{K},
\mathcal{Q} \otimes_{\mathcal{O}_T} \mathcal{J})
$$
若 $o_x(T')$ 为零，则定义它的短正合列分裂；换言之，得到
$\mathcal{O}_{X_{T'}}$-子模
$\mathcal{K}' \subset \pi^{-1}(\mathcal{K})$，它位于短正合列
$0 \to \mathcal{K} \otimes_{\mathcal{O}_T} \mathcal{J} \to
\mathcal{K}' \to \mathcal{K} \to 0$.
由上面所引的引理可知
% P11-E0388: the frozen source says "the lemma reference above" instead of "the lemma referenced above".
$\mathcal{Q}' = \mathcal{F}_{T'}/\mathcal{K}'$
是 $x$ 到 $Q_{\mathcal{F}/X/B}(T')$ 中元素的提升。反之，容易看出，
提升的存在性蕴含 $o_x(T')$ 为零。此外，若
$x \in Q_{\mathcal{F}/X/B}^{fp}(T)$，则由《形变理论》引理
\ref{defos-lemma-deform-fp-module-ringed-topoi}，自动有
$x' \in Q_{\mathcal{F}/X/B}^{fp}(T')$。若以后需要此评注，
我们会把它改写为引理，精确表述结果并给出详细证明（事实上，
以上全部论证在任意环化拓扑斯的语境下都成立）。
\end{remark}

\begin{remark}[商的形变]
\label{quot-remark-q-defos}
在情形 \ref{quot-situation-q} 中，{\bf 假设} $\mathcal{F}$ 在 $B$ 上平坦。
继续评注 \ref{quot-remark-q-obs} 的讨论。假设 $o_x(T') = 0$。
我们断言，提升 $x' \in Q_{\mathcal{F}/X/B}(T')$ 所成的集合是群
$$
\Hom_{\mathcal{O}_{X_T}}(\mathcal{K},
\mathcal{Q} \otimes_{\mathcal{O}_T} \mathcal{J})
$$
下的主齐性空间。具体地，给定在 $T'$ 上平坦且提升商
$\mathcal{Q}$ 的任意 $\mathcal{F}_{T'} \to \mathcal{Q}'$，
得到行列均正合的交换图
$$
\xymatrix{
& 0 \ar[d] & 0 \ar[d] & 0 \ar[d] \\
0 \ar[r] &
\mathcal{K} \otimes \mathcal{J} \ar[r] \ar[d] &
\mathcal{F}_T \otimes \mathcal{J} \ar[r] \ar[d] &
\mathcal{Q} \otimes \mathcal{J} \ar[r] \ar[d] &
0 \\
0 \ar[r] &
\mathcal{K}' \ar[r] \ar[d] &
\mathcal{F}_{T'} \ar[r] \ar[d] &
\mathcal{Q}' \ar[r] \ar[d] &
0 \\
0 \ar[r] &
\mathcal{K} \ar[d] \ar[r] &
\mathcal{F}_T \ar[d] \ar[r] &
\mathcal{Q} \ar[d] \ar[r] &
0 \\
& 0 & 0 & 0
}
$$
（为说明此点，使用前一评注中的观察）。给定映射
$\varphi : \mathcal{K} \to \mathcal{Q} \otimes \mathcal{J}$，
考虑子层 $\mathcal{K}'_\varphi \subset \mathcal{F}_{T'}$：
它由如下局部截面 $s$ 组成，其在 $\mathcal{F}_T$ 中的像是
$\mathcal{K}$ 的局部截面 $k$，而在 $\mathcal{Q}'$ 中的像是
$\mathcal{Q} \otimes \mathcal{J}$ 的局部截面 $\varphi(k)$。再置
$\mathcal{Q}'_\varphi = \mathcal{F}_{T'}/\mathcal{K}'_\varphi$。
反之，$x$ 的任意第二个提升都对应于以此方式构造的一个商。
% P11-E0389: the frozen source misspells "quotients" as "qotients".
若以后需要此评注，我们会把它改写为引理，精确表述结果并给出详细证明
（事实上，以上全部论证在任意环化拓扑斯的语境下都成立）。
\end{remark}










\section{Quot 函子}
\label{quot-section-quot}

\noindent
本节证明 Quot 函子是代数空间。

\begin{situation}
\label{quot-situation-quot}
设 $S$ 是概形，$f : X \to B$ 是 $S$ 上代数空间的态射。
假设 $f$ 有限表示。设 $\mathcal{F}$ 是拟凝聚
$\mathcal{O}_X$-模。对任意 $B$ 上概形 $T$，用 $X_T$ 表示
$X$ 到 $T$ 的基变换，用 $\mathcal{F}_T$ 表示 $\mathcal{F}$
沿投影态射 $X_T = X \times_S T \to X$ 的拉回。对这样的 $T$，置
$$
\Quotfunctor_{\mathcal{F}/X/B}(T) =
\left\{
\begin{matrix}
\text{商 }\mathcal{F}_T \to \mathcal{Q}\text{，其中 }\mathcal{Q}
\text{ 是拟凝聚的}\\
\mathcal{O}_{X_T}\text{-模，在 }T\text{ 上有限表示且平坦}\\
\text{其支撑在 }T\text{ 上真}
\end{matrix}
\right\}
$$
由《空间的导出范畴》引理
\ref{spaces-perfect-lemma-base-change-module-support-proper-over-base}
，这是我们在第 \ref{quot-section-functor-quotients} 节讨论的函子
$Q^{fp}_{\mathcal{F}/X/B}$ 的子函子。于是得到函子
\begin{equation}
\label{quot-equation-quot}
\Quotfunctor_{\mathcal{F}/X/B} : (\Sch/B)^{opp} \longrightarrow \textit{Sets}
\end{equation}
这称为与 $\mathcal{F}/X/B$ 相联系的{\it Quot 函子}。
\end{situation}

\noindent
在情形 \ref{quot-situation-quot} 中，有时把
$\Quotfunctor_{\mathcal{F}/X/B}$ 看作函子
$(\Sch/S)^{opp} \to \textit{Sets}$，它配有态射
$\Quotfunctor_{\mathcal{F}/X/B} \to B$。具体地，若 $T$ 是
$S$ 上的概形，则 $\Quotfunctor_{\mathcal{F}/X/B}(T)$ 的一个元素是
一对 $(h, \mathcal{Q})$，其中 $h$ 是态射 $h : T \to B$，
而 $Q$ 是 $X_T = X \times_{B, h} T$ 上有限表示、在 $T$ 上平坦、
且支撑在 $T$ 上真的商 $\mathcal{F}_T \to \mathcal{Q}$。
特别地，当我们说 $\Quotfunctor_{\mathcal{F}/X/B}$ 是代数空间时，
是指相应函子 $(\Sch/S)^{opp} \to \textit{Sets}$ 是代数空间。

\begin{lemma}
\label{quot-lemma-quot-sheaf}
在情形 \ref{quot-situation-quot} 中，函子
$\Quotfunctor_{\mathcal{F}/X/B}$ 满足 fpqc 拓扑的层条件。
\end{lemma}

\begin{proof}
在引理 \ref{quot-lemma-q-sheaf} 中已经看到，函子
$\text{Q}^{fp}_{\mathcal{F}/X/S}$ 是层。回忆，对 $S$ 上概形
$T$，子集
$\Quotfunctor_{\mathcal{F}/X/S}(T) \subset \text{Q}_{\mathcal{F}/X/S}(T)$
选出支撑在 $T$ 上真的商。由《空间上的下降》引理
\ref{spaces-descent-lemma-descending-property-proper} 与
《空间的态射》引理
\ref{spaces-morphisms-lemma-flat-pullback-support} 的结果，这定义了
一个子层；后一个引理表明取概形论支撑与平坦基变换可交换。
\end{proof}

\noindent
一致性检验：$\Quotfunctor_{\mathcal{F}/X/B}$ 在 $S$ 上代数空间
之间起着同样的作用。

\begin{lemma}
\label{quot-lemma-extend-quot-to-spaces}
在情形 \ref{quot-situation-quot} 中，设 $T$ 是 $S$ 上的代数空间。
$$
\Mor_{\Sh((\Sch/S)_{fppf})}(T,  \Quotfunctor_{\mathcal{F}/X/B}) =
\left\{
\begin{matrix}
(h, \mathcal{F}_T \to \mathcal{Q}) \text{，其中 }
h : T \to B \text{ 且}\\
\mathcal{Q}\text{ 有限表示且}\\
\text{在 }T\text{ 上平坦且支撑在 }T\text{ 上真}
\end{matrix}
\right\}
$$
其中 $\mathcal{F}_T$ 表示 $\mathcal{F}$ 到代数空间
$X \times_{B, h} T$ 的拉回。
\end{lemma}

\begin{proof}
注意，等式左右两边分别是引理
\ref{quot-lemma-extend-q-to-spaces} 中第二个等式左右两边的子集。
为看出这些子集在该引理证明给出的识别下相对应，只需证明：
给定 $h : T \to B$、满的 \'etale 态射 $U \to T$ 和有限型拟凝聚
$\mathcal{O}_{X_T}$-模 $\mathcal{Q}$，下列条件等价：
\begin{enumerate}
\item $\mathcal{Q}$ 的概形论支撑在 $T$ 上真；以及
\item $(X_U \to X_T)^*\mathcal{Q}$ 的概形论支撑在 $U$ 上真。
\end{enumerate}
这由《空间上的下降》引理
\ref{spaces-descent-lemma-descending-property-proper} 与
《空间的态射》引理
\ref{spaces-morphisms-lemma-flat-pullback-support} 得出；后者表明
取概形论支撑与平坦基变换可交换。
\end{proof}

\begin{proposition}
\label{quot-proposition-quot}
设 $S$ 是概形。设 $f : X \to B$ 是 $S$ 上代数空间的态射。
设 $\mathcal{F}$ 是 $X$ 上的拟凝聚层。若 $f$ 有限表示且分离，则
$\Quotfunctor_{\mathcal{F}/X/B}$ 是代数空间。若 $\mathcal{F}$
有限表示，则 $\Quotfunctor_{\mathcal{F}/X/B} \to B$ 在局部上有限表示。
\end{proposition}

\begin{proof}
由引理 \ref{quot-lemma-quot-sheaf}，我们知道
$\Quotfunctor_{\mathcal{F}/X/B}$ 是 fppf 拓扑中的层。令
$\textit{Quot}_{\mathcal{F}/X/B}$ 为与
$\Quotfunctor_{\mathcal{F}/X/S}$ 对应的群胚叠，见《代数叠》第
\ref{algebraic-section-split} 节。由《代数叠》命题
\ref{algebraic-proposition-algebraic-stack-no-automorphisms}，
只需证明 $\textit{Quot}_{\mathcal{F}/X/B}$ 是代数叠。
考虑群胚叠的 $1$-态射
$$
\textit{Quot}_{\mathcal{F}/X/S}
\longrightarrow
\Cohstack_{X/B}
$$
它在 $(\Sch/S)_{fppf}$ 上把商
$\mathcal{F}_T \to \mathcal{Q}$ 关联到模 $\mathcal{Q}$。
由定理 \ref{quot-theorem-coherent-algebraic-general}，我们知道
$\Cohstack_{X/B}$ 是代数叠。由《代数叠》引理
\ref{algebraic-lemma-representable-morphism-to-algebraic}，
只需证明此 $1$-态射可由代数空间表示。

\medskip\noindent
设 $T$ 是 $S$ 上的概形，且 $\Cohstack_{X/B}$ 在 $T$ 上的对象
$(h, \mathcal{G})$ 对应于 $1$-态射
$\xi : (\Sch/T)_{fppf} \to \Cohstack_{X/B}$。其 $2$-纤维积
$$
\mathcal{Z} =
(\Sch/T)_{fppf}
\times_{\xi, \Cohstack_{X/B}}
\textit{Quot}_{\mathcal{F}/X/S}
$$
是集合胚中的叠，见《叠》引理
\ref{stacks-lemma-2-fibre-product-gives-stack-in-setoids}。
相应的集合层（即函子，见
《叠》引理
\ref{stacks-lemma-2-fibre-product-gives-stack-in-setoids} 和
\ref{stacks-lemma-when-stack-in-sets}）把概形 $T'/T$ 关联到
$X_{T'}$ 上拟凝聚模的满射 $u : \mathcal{F}_{T'} \to \mathcal{G}_{T'}$
的集合。因此，$\mathcal{Z}$ 可由代数空间
$\mathit{Hom}(\mathcal{F}_T, \mathcal{G})$ 的一个开子空间表示
（由《空间上的平坦性》引理 \ref{spaces-flat-lemma-F-surj-open}），
该代数空间来自命题 \ref{quot-proposition-hom}。
\end{proof}

\begin{remark}[通过 Artin 公理得到 Quot 函子]
\label{quot-remark-quot-via-artins-axioms}
设 $S$ 是局部环全为 G-环的 Noether 概形。设 $X$ 是 $S$ 上的代数空间，
其结构态射 $f : X \to S$ 有限表示且分离。设 $\mathcal{F}$ 是 $X$ 上
有限表示的拟凝聚层，且在 $S$ 上平坦。本备注概述如何使用 Artin 公理证明
$\Quotfunctor_{\mathcal{F}/X/S}$ 是在 $S$ 上局部有限表示的代数空间，
从而无需像命题 \ref{quot-proposition-quot} 的证明那样使用凝聚层叠的代数性。

\medskip\noindent
我们检查《Artin 公理》命题
\ref{artin-proposition-spaces-diagonal-representable} 中所列的条件。
$\Quotfunctor_{\mathcal{F}/X/S}$ 的对角态射可表示性可如下看出：
设有两个商 $\mathcal{F}_T \to \mathcal{Q}_i$，$i = 1, 2$。记第一个商的核为
$\mathcal{K}_1$。那么需要证明使 $u : \mathcal{K}_1 \to \mathcal{Q}_2$
变为零的 $T$ 的轨迹是可表示的。例如，这可由《空间上的平坦性》引理
\ref{spaces-flat-lemma-F-zero-closed-proper}
或本章前面对 $\mathit{Hom}$ 层的讨论得到。公理 [0]（层）、[1]（极限）、
[2]（Rim--Schlessinger）由引理 \ref{quot-lemma-quot-sheaf}、
\ref{quot-lemma-q-limit-preserving} 和 \ref{quot-lemma-q-RS-star}
得到（再加上处理真性条件的一些额外工作）。
公理 [3]（切空间有限维）由备注 \ref{quot-remark-q-defos} 中对无穷小形变的描述，
以及域上真代数空间中凝聚层上同调的有限性（《空间的上同调》引理
\ref{spaces-cohomology-lemma-proper-pushforward-coherent}）得到。
公理 [4]（形式对象的有效性）由 Grothendieck 存在性定理
（《空间态射补充》定理 \ref{spaces-more-morphisms-theorem-grothendieck-existence}）得到。
通常最难验证的是公理 [5]（通用性的开性）。例如，可以使用备注
\ref{quot-remark-q-obs} 中描述的障碍理论和备注 \ref{quot-remark-q-defos} 中的形变描述，
再利用《Artin 公理》引理
\ref{artin-lemma-get-openness-obstruction-theory} 中的判据来完成。
请与引理 \ref{quot-lemma-coherent-defo-thy} 的第二个证明比较。
\end{remark}









\section{Hilbert 函子}
\label{quot-section-hilb}

\noindent
本节证明 Hilb 函子是代数空间。

\begin{situation}
\label{quot-situation-hilb}
设 $S$ 是概形。设 $f : X \to B$ 是 $S$ 上代数空间的态射。
假设 $f$ 有限表示。对 $B$ 上任意概形 $T$，用 $X_T$ 表示
$X$ 到 $T$ 的基变换。对这样的 $T$，置
$$
\Hilbfunctor_{X/B}(T) =
\left\{
\begin{matrix}
\text{闭子空间 }Z \subset X_T\text{，使得 }Z \to T\\
\text{有限表示、平坦且真}
\end{matrix}
\right\}
$$
由于基变换保持所需性质
（《空间》引理 \ref{spaces-lemma-base-change-immersions} 以及
《空间的态射》引理
\ref{spaces-morphisms-lemma-base-change-finite-presentation},
\ref{spaces-morphisms-lemma-base-change-flat} 和
\ref{spaces-morphisms-lemma-base-change-proper}），
得到函子
\begin{equation}
\label{quot-equation-hilb}
\Hilbfunctor_{X/B} : (\Sch/B)^{opp} \longrightarrow \textit{Sets}
\end{equation}
这称为与 $X/B$ 相联系的 {\it Hilbert 函子}。
\end{situation}

\noindent
在情形 \ref{quot-situation-hilb} 中，有时把 $\Hilbfunctor_{X/B}$
看作函子 $(\Sch/S)^{opp} \to \textit{Sets}$，它配有态射
$\Hilbfunctor_{X/S} \to B$。具体地，若 $T$ 是 $S$ 上的概形，则
$\Hilbfunctor_{X/B}(T)$ 的一个元素是一对 $(h, Z)$，其中 $h$ 是态射
$h : T \to B$，且 $Z \subset X_T = X \times_{B, h} T$
是闭子概形，在 $T$ 上平坦、真且有限表示。特别地，当我们说
$\Hilbfunctor_{X/B}$ 是代数空间时，是指相应函子
$(\Sch/S)^{opp} \to \textit{Sets}$ 是代数空间。

\medskip\noindent
当然，Hilbert 函子只是 Quot 函子的特例。

\begin{lemma}
\label{quot-lemma-hilb-is-quot}
在情形 \ref{quot-situation-hilb} 中，有
$\Hilbfunctor_{X/B} = \Quotfunctor_{\mathcal{O}_X/X/B}$.
\end{lemma}

\begin{proof}
设 $T$ 是 $B$ 上的概形。给定
$Z \in \Hilbfunctor_{X/B}(T)$，考虑商
$\mathcal{O}_{X_T} \to i_*\mathcal{O}_Z$，其中
$i : Z \to X_T$ 是包含态射。注意，$i_*\mathcal{O}_Z$ 是拟凝聚的。
由于 $Z \to T$ 和 $X_T \to T$ 都有限表示，我们看到 $i$ 有限表示
（《空间的态射》引理
\ref{spaces-morphisms-lemma-finite-presentation-permanence}），
$i_*\mathcal{O}_Z$ 是有限表示的 $\mathcal{O}_{X_T}$-模
（《空间上的下降》引理
\ref{spaces-descent-lemma-finite-finitely-presented-module}).
由于 $Z \to T$ 真，我们看到 $i_*\mathcal{O}_Z$ 的支撑在 $T$ 上真
（按《空间的导出范畴》第
\ref{spaces-perfect-section-proper-over-base}).
由于 $\mathcal{O}_Z$ 在 $T$ 上平坦且 $i$ 仿射，我们看到
$i_*\mathcal{O}_Z$ 在 $T$ 上平坦（略去一个小论证）。因此
$\mathcal{O}_{X_T} \to i_*\mathcal{O}_Z$
是 $\Quotfunctor_{\mathcal{O}_X/X/B}(T)$ 的元素。

\medskip\noindent
反之，给定 $\Quotfunctor_{\mathcal{O}_X/X/B}(T)$ 的元素
$\mathcal{O}_{X_T} \to \mathcal{Q}$，考虑与拟凝聚理想层
$\mathcal{I} = \Ker(\mathcal{O}_{X_T} \to \mathcal{Q})$
相对应的闭浸入 $i : Z \to X_T$
（《空间的态射》引理
\ref{spaces-morphisms-lemma-closed-immersion-ideals}）。
由 $Z$ 的构造，$\mathcal{Q} = i_*\mathcal{O}_Z$。于是反向阅读
上面的论证可知，$Z$ 给出 $\Hilbfunctor_{X/B}(T)$ 的一个元素。
例如，$\mathcal{I}$ 是有限型拟凝聚的
（《站点上的模》引理
\ref{sites-modules-lemma-kernel-surjection-finite-onto-finite-presentation})
因此 $i : Z \to X_T$ 有限表示
（《空间的态射》引理
\ref{spaces-morphisms-lemma-closed-immersion-finite-presentation})
从而 $Z \to T$ 有限表示
（《空间的态射》引理
\ref{spaces-morphisms-lemma-composition-finite-presentation}).
由《空间的导出范畴》第
\ref{spaces-perfect-section-proper-over-base} 节中所述，$Z \to T$ 为真。
$Z \to T$ 的平坦性由
$\mathcal{Q}$ 在 $T$ 上的平坦性推出。

\medskip\noindent
我们略去验证上述两个构造互为逆构造（这是直接的）。
\end{proof}

\noindent
一致性检验：$\Hilbfunctor_{X/B}$ 的层在 $S$ 上的代数空间之间起着同样的作用。

\begin{lemma}
\label{quot-lemma-extend-hilb-to-spaces}
在情形 \ref{quot-situation-hilb} 中，设 $T$ 是 $S$ 上的代数空间。我们有
$$
\Mor_{\Sh((\Sch/S)_{fppf})}(T, \Hilbfunctor_{X/B}) =
\left\{
\begin{matrix}
(h, Z)\text{，其中 }h : T \to B,\ Z \subset X_T \\
\text{在 }T\text{ 上有限表示、平坦且真}
\end{matrix}
\right\}
$$
其中 $X_T = X \times_{B, h} T$。
\end{lemma}

\begin{proof}
由引理 \ref{quot-lemma-hilb-is-quot}，有
$\Hilbfunctor_{X/B} = \Quotfunctor_{\mathcal{O}_X/X/B}$。
因此可应用引理 \ref{quot-lemma-extend-quot-to-spaces}，看出左边
与如下满射的集合一一对应：$\mathcal{O}_{X_T} \to \mathcal{Q}$，
它们有限表示、在 $T$ 上平坦且支撑在 $T$ 上真。完全按照
引理 \ref{quot-lemma-hilb-is-quot} 的证明，我们看到这些商恰好对应于
满足 $Z \to T$ 真、平坦且有限表示的闭浸入 $Z \to X_T$。
\end{proof}

\begin{proposition}
\label{quot-proposition-hilb}
设 $S$ 是概形。设 $f : X \to B$ 是 $S$ 上代数空间的态射。
若 $f$ 有限表示且分离，则 $\Hilbfunctor_{X/B}$ 是在 $B$ 上局部有限
表示的代数空间。
\end{proposition}

\begin{proof}
这立即由引理 \ref{quot-lemma-hilb-is-quot} 和命题
\ref{quot-proposition-quot} 推出。
\end{proof}






\section{Picard 叠}
\label{quot-section-picard-stack}

\noindent
代数空间态射的 Picard 叠见《叠的例子》第
\ref{examples-stacks-section-picard-stack} 节。下面的引理将推出，在良好情形下，
它是凝聚层叠的一个开子叠。

\begin{lemma}
\label{quot-lemma-picard-stack-open-in-coh}
设 $S$ 是概形。设 $f : X \to B$ 是 $S$ 上代数空间的态射，且平坦、
有限表示并且真。自然映射
$$
\Picardstack_{X/B} \longrightarrow \Cohstack_{X/B}
$$
可表示为开浸入。
\end{lemma}

\begin{proof}
注意，该映射只是把《叠的例子》第
\ref{examples-stacks-section-picard-stack} 节中的三元组 $(T, g, \mathcal{L})$
送到同一个三元组 $(T, g, \mathcal{L})$，但现在把它视为情形
\ref{quot-situation-coherent} 所述类型的三元组。之所以可行，是因为可逆的
$\mathcal{O}_{X_T}$-模 $\mathcal{L}$ 显然是有限表示的
$\mathcal{O}_{X_T}$-模；它在 $T$ 上平坦，因为 $X_T \to T$ 平坦；并且其支撑
在 $T$ 上为真，因为 $X_T \to T$ 为真
（《空间的态射》引理 \ref{spaces-morphisms-lemma-base-change-flat}
和 \ref{spaces-morphisms-lemma-base-change-proper}）。因此陈述是有意义的。

\medskip\noindent
说到这里，显然引理的内容如下：给定 $\Cohstack_{X/B}$ 的对象
$(T, g, \mathcal{F})$，存在开子概形 $U \subset T$，使得对于概形态射
$T' \to T$，下列条件等价：
\begin{enumerate}
\item[(a)] $T' \to T$ 经由 $U$ 分解；
\item[(b)] $\mathcal{F}$ 沿 $X_{T'} \to X_T$ 拉回所得的
$\mathcal{F}_{T'}$ 是可逆的。
\end{enumerate}
令 $W \subset |X_T|$ 为点 $x \in |X_T|$ 的集合，其中 $\mathcal{F}$
在 $x$ 的某个邻域内局部自由。由《空间的态射进阶》引理
\ref{spaces-more-morphisms-lemma-finite-free-open}.
，$W$ 是开集，且 $W$ 的构造与任意基变换相容。显然，若 $T' \to T$
满足 (b)，则 $|X_{T'}| \to |X_T|$ 的像落在 $W$ 中。因此，为构造 $U$，
可以把 $T$ 替换为开集 $T \setminus f_T(|X_T| \setminus W)$。这样便处于
$\mathcal{F}$ 有限局部自由的情形。在此情形有不交并分解
$X_T = X_0 \amalg X_1 \amalg X_2 \amalg \ldots$
为开且闭子空间，使得 $\mathcal{F}$ 在 $X_i$ 上的限制是秩为 $i$ 的局部自由模。
于是显然
$$
U = T \setminus f_T(|X_0| \cup |X_2| \cup |X_3| \cup \ldots )
$$
可行。（注意，若假设 $T$ 拟紧，则 $X_T$ 拟紧，因而只有有限多个 $X_i$
非空，所以 $U$ 确实是开集。）
\end{proof}

\begin{proposition}
\label{quot-proposition-pic}
设 $S$ 是概形。设 $f : X \to B$ 是 $S$ 上代数空间的态射。若 $f$ 平坦、
有限表示且真，则 $\Picardstack_{X/B}$ 是代数叠。
\end{proposition}

\begin{proof}
立即由引理 \ref{quot-lemma-picard-stack-open-in-coh}、《代数叠》引理
\ref{algebraic-lemma-representable-morphism-to-algebraic}
以及下列定理之一推出：
定理 \ref{quot-theorem-coherent-algebraic}
或
定理 \ref{quot-theorem-coherent-algebraic-general}
\end{proof}








\section{Picard 函子}
\label{quot-section-picard-functor}

\noindent
本节重新讨论《曲线的 Picard 概形》第 \ref{pic-section-picard-functor} 节中的
Picard 函子。这里的讨论更一般，因为我们要研究代数空间态射的 Picard 函子，
正如 Picard 叠一节中所做的，见第 \ref{quot-section-picard-stack} 节。

\medskip\noindent
设 $S$ 是概形，$X$ 是 $S$ 上的代数空间。$X$ 上的可逆层是
$X_\etale$ 上的可逆 $\mathcal{O}_X$-模，见《站点上的模》定义
\ref{sites-modules-definition-invertible-sheaf}。可逆模同构类所成的群记为
$\Pic(X)$，见《站点上的模》定义 \ref{sites-modules-definition-pic}。
给定 $S$ 上代数空间的态射 $f : X \to Y$，拉回给出群同态
$\Pic(Y) \to \Pic(X)$。赋值 $X \leadsto \Pic(X)$ 是从概形范畴到
阿贝尔群范畴的逆变函子。该函子不可表示，但这个构造的相对版本有时可表示。

\begin{situation}
\label{quot-situation-pic}
设 $S$ 是概形。设 $f : X \to B$ 是 $S$ 上代数空间的态射。定义
$$
\Picardfunctor_{X/B} : (\Sch/B)^{opp} \longrightarrow \textit{Sets}
$$
为如下函子的 fppf 层化：该函子把 $B$ 上概形 $T$ 关联到群 $\Pic(X_T)$。
\end{situation}

\noindent
在情形 \ref{quot-situation-pic} 中，有时把 $\Picardfunctor_{X/B}$ 看作函子
$(\Sch/S)^{opp} \to \textit{Sets}$，它配有态射
$\Picardfunctor_{X/B} \to B$。在这种观点下，定义 $\Picardfunctor_{X/B}$
为如下函子的 fppf 层化：
$$
T/S \longmapsto \{(h, \mathcal{L}) \mid
h : T \to B,\ \mathcal{L} \in \Pic(X \times_{B, h} T)\}
$$
特别地，当我们说 $\Picardfunctor_{X/B}$ 是代数空间时，是指相应函子
$(\Sch/S)^{opp} \to \textit{Sets}$ 是代数空间。

\medskip\noindent
常用的一个备注是：若 $T$ 是 $B$ 上的概形，则
$\Picardfunctor_{X_T/T}$ 是 $\Picardfunctor_{X/B}$ 到
$(\Sch/T)_{fppf}$ 的限制。

\begin{lemma}
\label{quot-lemma-pic-over-pic}
在情形 \ref{quot-situation-pic} 中，函子 $\Picardfunctor_{X/B}$ 是函子
$T \mapsto \Ob(\Picardstack_{X/B, T})/\cong$ 的层化。
\end{lemma}

\begin{proof}
由于 Picard 叠 $\Picardstack_{X/B}$ 在 $T$ 上的纤维范畴
$\Picardstack_{X/B, T}$ 是 $X_T$ 上可逆层所成的范畴
（见第 \ref{quot-section-picard-stack} 节和《叠的例子》第
\ref{examples-stacks-section-picard-stack} 节），这由定义立即得到。
\end{proof}

\noindent
要看出 $\Picardfunctor_{X/B}$ 在 $B$ 上概形 $T$ 处的取值并不容易。
下面的引理有助于完成这一任务。

\begin{lemma}
\label{quot-lemma-flat-geometrically-connected-fibres}
在情形 \ref{quot-situation-pic} 中，若对 $B$ 上所有概形 $T$，
$\mathcal{O}_T \to f_{T, *}\mathcal{O}_{X_T}$ 都是同构，则
$$
0 \to \Pic(T) \to \Pic(X_T) \to \Picardfunctor_{X/B}(T)
$$
对所有 $T$ 都是正合列。
\end{lemma}

\begin{proof}
可以把 $B$ 替换为 $T$、把 $X$ 替换为 $X_T$，并假设 $B = T$，
以简化记号。设 $\mathcal{N}$ 是可逆的 $\mathcal{O}_B$-模。若
$f^*\mathcal{N} \cong \mathcal{O}_X$，则由假设
$f_*f^*\mathcal{N} \cong f_*\mathcal{O}_X \cong \mathcal{O}_B$。
由于 $\mathcal{N}$ 局部平凡，典范映射
$\mathcal{N} \to f_*f^*\mathcal{N}$ 在局部上是同构（因为由假设
$\mathcal{O}_B \to f_*f^*\mathcal{O}_B$ 是同构）。因此
$\mathcal{N} \to f_*f^*\mathcal{N} \to \mathcal{O}_B$ 是同构，
从而 $\mathcal{N}$ 平凡。这证明第一支箭头是单射。

\medskip\noindent
设 $\mathcal{L}$ 是 $\Pic(X) \to \Picardfunctor_{X/B}(B)$ 的核中的
可逆 $\mathcal{O}_X$-模。存在 fppf 覆盖 $\{B_i \to B\}$，使得
$\mathcal{L}$ 拉回到 $X_{B_i}$ 上的平凡可逆层。取一个平凡化截面
$s_i$。则 $\text{pr}_0^*s_i$ 和 $\text{pr}_1^*s_j$ 都是
$X_{B_i \times_B B_j}$ 上 $\mathcal{L}$ 的平凡化截面，因而相差一个乘法单位元
$$
f_{ij} \in
\Gamma(X_{S_i \times_B B_j}, \mathcal{O}_{X_{B_i \times_B B_j}}^*) =
\Gamma(B_i \times_B B_j, \mathcal{O}_{B_i \times_N B_j}^*)
$$
（等式由我们关于结构层推前的假设得到）。显然，这些元素在
$B_i \times_B B_j \times_B B_k$ 上满足上闭链条件，因而它们为 fppf 覆盖
$\{B_i \to B\}$ 上的可逆层定义下降数据。由《下降》命题
\ref{descent-proposition-fpqc-descent-quasi-coherent}，存在可逆
$\mathcal{O}_B$-模 $\mathcal{N}$，它在各 $B_i$ 上带有平凡化，且相应下降数据
为 $\{f_{ij}\}$。（该命题适用，是因为在证明开始处的替换后 $B$ 是概形。）
由于从下降数据到模的函子是全忠实的，故 $f^*\mathcal{N} \cong \mathcal{L}$。
\end{proof}

\begin{lemma}
\label{quot-lemma-flat-geometrically-connected-fibres-with-section}
在情形 \ref{quot-situation-pic} 中，设 $\sigma : B \to X$ 是一个截面。
假设对 $B$ 上所有 $T$，$\mathcal{O}_T \to f_{T, *}\mathcal{O}_{X_T}$
都是同构。则
$$
0 \to \Pic(T) \to \Pic(X_T) \to \Picardfunctor_{X/B}(T) \to 0
$$
是分裂正合列，分裂由
$\sigma_T^* : \Pic(X_T) \to \Pic(T)$ 给出。
\end{lemma}

\begin{proof}
记 $K(T) = \Ker(\sigma_T^* : \Pic(X_T) \to \Pic(T))$。由于
$\sigma$ 是 $f$ 的截面，我们看到 $\Pic(X_T)$ 是 $\Pic(T)$ 与 $K(T)$
的直和。因此由引理 \ref{quot-lemma-flat-geometrically-connected-fibres}，
对所有 $T$ 有 $K(T) \subset \Picardfunctor_{X/B}(T)$。此外，由构造显然
$\Picardfunctor_{X/B}$ 是预层 $K$ 的层化。要完成证明，只需证明 $K$
满足 fppf 覆盖的层条件；这将在下一段完成。

\medskip\noindent
设 $\{T_i \to T\}$ 是 fppf 覆盖。设 $\mathcal{L}_i$ 是 $K(T_i)$ 的元素，
它们对所有 $i$ 和 $j$ 在 $K(T_i \times_T T_j)$ 中的像相同。取同构
$\alpha_i : \mathcal{O}_{T_i} \to \sigma_{T_i}^*\mathcal{L}_i$
对所有 $i$ 成立。再取同构
$$
\varphi_{ij} :
\mathcal{L}_i|_{X_{T_i \times_T T_j}}
\longrightarrow
\mathcal{L}_j|_{X_{T_i \times_T T_j}}
$$
若映射
$$
\alpha_j|_{T_i \times_T T_j} \circ
\sigma_{T_i \times_T T_j}^*\varphi_{ij} \circ
\alpha_i|_{T_i \times_T T_j} :
\mathcal{O}_{T_i \times_T T_j} \to \mathcal{O}_{T_i \times_T T_j}
$$
不等于乘以 $1$ 而是乘以某个 $u_{ij}$，则可将
$\varphi_{ij}$ 乘以 $u_{ij}^{-1}$ 来修正。完成后，考虑自映射
$$
\varphi_{ki}|_{X_{T_i \times_T T_j \times_T T_k}} \circ
\varphi_{jk}|_{X_{T_i \times_T T_j \times_T T_k}} \circ
\varphi_{ij}|_{X_{T_i \times_T T_j \times_T T_k}}
\quad\text{作用于}\quad
\mathcal{L}_i|_{X_{T_i \times_T T_j \times_T T_k}}
$$
它等于乘以 $X_{T_i \times_T T_j \times_T T_k}$ 的结构层的某个截面
$f_{ijk}$。由我们对 $\varphi_{ij}$ 的选择，该映射沿 $\sigma$ 的拉回
等于乘以 $1$。由关于 $X$ 上函数的假设，$f_{ijk} = 1$。因此得到 fppf 覆盖
$\{X_{T_i} \to X\}$ 的下降数据。由《空间上的下降》命题
\ref{spaces-descent-proposition-fpqc-descent-quasi-coherent}
存在可逆 $\mathcal{O}_{X_T}$-模 $\mathcal{L}$ 及同构
$\alpha : \mathcal{O}_T \to \sigma_T^*\mathcal{L}$，其到 $X_{T_i}$ 的拉回恢复
$(\mathcal{L}_i, \alpha_i)$（略去一个小细节）。因此 $\mathcal{L}$ 如愿定义了
$K(T)$ 的一个元素。
\end{proof}

\noindent
在情形 \ref{quot-situation-pic} 中，设 $\sigma : B \to X$ 是一个截面。
记 $\Picardstack_{X/B, \sigma}$ 为如下定义的范畴：
\begin{enumerate}
\item 对象是四元组 $(T, h, \mathcal{L}, \alpha)$，其中
$(T, h, \mathcal{L})$ 是 $T$ 上 $\Picardstack_{X/B}$ 的对象，且
$\alpha : \mathcal{O}_T \to \sigma_T^*\mathcal{L}$ 是同构。
\item 态射 $(g, \varphi) : (T, h, \mathcal{L}, \alpha)
\to (T', h', \mathcal{L}', \alpha')$ 由满足 $h = h' \circ g$ 的概形态射
$g : T \to T'$ 以及同构 $\varphi : (g')^*\mathcal{L}' \to \mathcal{L}$ 给出，
并满足 $\sigma_T^*\varphi \circ g^*\alpha' = \alpha$。
这里 $g' : X_{T'} \to X_T$ 是 $g$ 的基变换。
\end{enumerate}
有自然的忠实遗忘函子
$$
\Picardstack_{X/B, \sigma} \longrightarrow
\Picardstack_{X/B}
$$
这样，我们把 $\Picardstack_{X/B, \sigma}$ 视为 $(\Sch/S)_{fppf}$ 上的范畴。

\begin{lemma}
\label{quot-lemma-pic-with-section-stack}
在情形 \ref{quot-situation-pic} 中，设 $\sigma : B \to X$ 是一个截面。
则上述定义的 $\Picardstack_{X/B, \sigma}$ 是 $(\Sch/S)_{fppf}$ 上的群胚叠。
\end{lemma}

\begin{proof}
由《叠的例子》引理 \ref{examples-stacks-lemma-picard-stack}，我们已经知道
$\Picardstack_{X/B}$ 是 $(\Sch/S)_{fppf}$ 上的群胚叠。下面证明
$\Picardstack_{X/B, \sigma}$ 的对象下降。设 $\{T_i \to T\}$ 是 fppf 覆盖，
设 $\xi_i = (T_i, h_i, \mathcal{L}_i, \alpha_i)$ 是位于 $T_i$ 上的
$\Picardstack_{X/B, \sigma}$ 对象，并设
$\varphi_{ij} : \text{pr}_0^*\xi_i \to \text{pr}_1^*\xi_j$
是下降数据。应用 $\Picardstack_{X/B}$ 的结果，可假设存在 $T$ 上
$\Picardstack_{X/B}$ 的对象 $(T, h, \mathcal{L})$，其拉回对所有 $i$ 都是 $\xi_i$。
于是得到
$$
\alpha_i : \mathcal{O}_{T_i} \to \sigma_{T_i}^*\mathcal{L}_i =
(T_i \to T)^*\sigma_T^*\mathcal{L}
$$
由于映射 $\varphi_{ij}$ 与 $\alpha_i$ 相容，$\alpha_i$ 和 $\alpha_j$
在 $T_i \times_T T_j$ 上的拉回相同。由拟凝聚层的下降
（《下降》命题 \ref{descent-proposition-fpqc-descent-quasi-coherent}），
可知 $\alpha_i$ 是单一映射
$\alpha : \mathcal{O}_T \to \sigma_T^*\mathcal{L}$ 的限制，如所需。
我们略去态射下降的证明。
\end{proof}

\begin{lemma}
\label{quot-lemma-compare-pic-with-section}
在情形 \ref{quot-situation-pic} 中，设 $\sigma : B \to X$ 是一个截面。
态射 $\Picardstack_{X/B, \sigma} \to \Picardstack_{X/B}$
可表示、满射且光滑。
\end{lemma}

\begin{proof}
设 $T$ 是概形，且 $(\Sch/T)_{fppf} \to \Picardstack_{X/B}$
由 $T$ 上 $\Picardstack_{X/B}$ 的对象 $\xi = (T, h, \mathcal{L})$ 给出。
我们需要证明
$$
(\Sch/T)_{fppf} \times_{\xi, \Picardstack_{X/B}} \Picardstack_{X/B, \sigma}
$$
可由概形 $V$ 表示，且相应态射 $V \to T$ 满射且光滑。见《代数叠》各节
\ref{algebraic-section-representable-morphism}、
\ref{algebraic-section-morphisms-representable-by-algebraic-spaces} 和
\ref{algebraic-section-representable-properties}。遗忘函子
$\Picardstack_{X/B, \sigma} \to \Picardstack_{X/B}$ 在纤维范畴上忠实，且对
$T'/T$，同构类的集合就是同构的集合
$$
\alpha' : \mathcal{O}_{T'} \longrightarrow (T' \to T)^*\sigma_T^*\mathcal{L}
$$
见《代数叠》引理
\ref{algebraic-lemma-criterion-map-representable-spaces-fibred-in-groupoids}.
由命题 \ref{quot-proposition-isom}，我们知道该函子可由 $T$ 上有限表示的仿射概形
$U$ 表示（将命题应用于 $\text{id} : T \to T$、$\mathcal{O}_T$ 和
$\sigma^*\mathcal{L}$）。在 $T$ 上作 Zariski 局部化，可假设
$\sigma_T^*\mathcal{L}$ 同构于 $\mathcal{O}_T$，于是我们的函子由 $T$ 上的
$\mathbf{G}_m \times T$ 表示。因此，$U \to T$ 在 $T$ 上 Zariski 局部看起来
就是投影 $\mathbf{G}_m \times T \to T$，确实光滑且满射。
\end{proof}

\begin{lemma}
\label{quot-lemma-flat-geometrically-connected-fibres-with-section-functor-stack}
在情形 \ref{quot-situation-pic} 中，设 $\sigma : B \to X$ 是一个截面。
若对 $B$ 上所有 $T$，$\mathcal{O}_T \to f_{T, *}\mathcal{O}_{X_T}$
都是同构，则 $\Picardstack_{X/B, \sigma} \to (\Sch/S)_{fppf}$
是集合胚上的纤维化范畴，其在 $T$ 上的同构类集合为
$$
\coprod\nolimits_{h : T \to B}
\Ker(\sigma_T^* : \Pic(X \times_{B, h} T) \to \Pic(T))
$$
\end{lemma}

\begin{proof}
若 $\xi = (T, h, \mathcal{L}, \alpha)$ 是 $T$ 上
$\Picardstack_{X/B, \sigma}$ 的对象，则 $\xi$ 的自同构 $\varphi$
由 $X_T$ 结构层的可逆全局截面 $u$ 的乘法给出，并且满足
$\sigma_T^*u = 1$。由假设 $\mathcal{O}_T \to f_{T, *}\mathcal{O}_{X_T}$
是同构，故 $u = 1$。因此 $\Picardstack_{X/B, \sigma}$
是 $(\Sch/S)_{fppf}$ 上集合胚中的纤维化范畴。
给定 $T$ 和 $h : T \to B$，二元组 $(\mathcal{L}, \alpha)$ 的同构类集合
等于满足 $\sigma_T^*\mathcal{L} \cong \mathcal{O}_T$（同构未指定）的
$\mathcal{L}$ 的同构类集合。这是显然的，因为 $\alpha$ 的任意两个选择
相差一个 $T$ 上的全局单位元，而这与 $X_T$ 上的全局单位元相同。
\end{proof}

\begin{proposition}
\label{quot-proposition-pic-functor}
设 $S$ 是概形。设 $f : X \to B$ 是 $S$ 上代数空间的态射。假设
\begin{enumerate}
\item $f$ 平坦、有限表示且真；
\item 对 $B$ 上所有概形 $T$，$\mathcal{O}_T \to f_{T, *}\mathcal{O}_{X_T}$
都是同构。
\end{enumerate}
则 $\Picardfunctor_{X/B}$ 是代数空间。
\end{proposition}

\noindent
在该命题的情形下，代数叠 $\Picardstack_{X/B}$ 是代数空间
$\Picardfunctor_{X/B}$ 上的束叠。发展束叠的一般理论后，
这给出命题的一个较短证明（但使用了更一般的理论）。

\begin{proof}
存在代数空间的满射、平坦、有限表示态射 $B' \to B$，使得基变换
$X' = X \times_B B'$ 在 $B'$ 上有截面；具体地，可取 $B' = X$。
注意 $\Picardfunctor_{X'/B'} = B' \times_B \Picardfunctor_{X/B}$。
因此 $\Picardfunctor_{X'/B'} \to \Picardfunctor_{X/B}$ 可由代数空间表示，
且满射、平坦并有限表示。于是，若能证明
$\Picardfunctor_{X'/B'}$ 是代数空间，则由《自举》定理
\ref{bootstrap-theorem-final-bootstrap}.
可推出 $\Picardfunctor_{X/B}$ 是代数空间。这样就归约到下一段所述情形。

\medskip\noindent
除命题假设外，假设有截面 $\sigma : B \to X$。由命题
\ref{quot-proposition-pic}，$\Picardstack_{X/B}$ 是代数叠。由引理
\ref{quot-lemma-compare-pic-with-section} 和《代数叠》引理
\ref{algebraic-lemma-representable-morphism-to-algebraic}
可知 $\Picardstack_{X/B, \sigma}$ 是代数叠。由引理
\ref{quot-lemma-flat-geometrically-connected-fibres-with-section-functor-stack}
和《代数叠》引理
\ref{algebraic-lemma-characterize-representable-by-space}
可知 $T \mapsto \Ker(\sigma_T^* : \Pic(X_T) \to \Pic(T))$
是代数空间。由引理 \ref{quot-lemma-flat-geometrically-connected-fibres-with-section}，
该函子与 $\Picardfunctor_{X/B}$ 相同。
\end{proof}

\begin{lemma}
\label{quot-lemma-diagonal-pic}
采用命题 \ref{quot-proposition-pic-functor} 中的假设和记号。则对角态射
$\Picardfunctor_{X/B} \to \Picardfunctor_{X/B} \times_B \Picardfunctor_{X/B}$
可表示为浸入。换言之，$\Picardfunctor_{X/B} \to B$ 局部分离。
\end{lemma}

\begin{proof}
设 $T$ 是 $B$ 上的概形，且 $s, t \in \Picardfunctor_{X/B}(T)$。
我们要证明存在局部闭子概形 $Z \subset T$，使得 $s|_Z = t|_Z$，并且
态射 $T' \to T$ 经由 $Z$ 分解，当且仅当 $s|_{T'} = t|_{T'}$。

\medskip\noindent
先将一般问题归约到 $s$ 和 $t$ 来自 $X_T$ 上可逆模的情形。建议读者跳过此步。
取 fppf 覆盖 $\{T_i \to T\}_{i \in I}$，使得对所有 $i$，
$s|_{T_i}$ 和 $t|_{T_i}$ 来自 $\Pic(X_{T_i})$。假设对所有对
$s|_{T_i}, t|_{T_i}$ 都能证明结论，则得到具有所需泛性质的局部闭子概形
$Z_i \subset T_i$。于是 $Z_i$ 和 $Z_j$ 在 $T_i \times_T T_j$ 中的概形论逆像
相同。这给出 $Z_i/T_i$ 上的下降数据。由于 $Z_i \to T_i$ 局部拟有限，
由《空间态射补充》引理
\ref{more-morphisms-lemma-separated-locally-quasi-finite-morphisms-fppf-descend}
可得局部拟有限态射 $Z \to T$，其基变换恢复 $Z_i \to T_i$。再由《下降》引理
\ref{descent-lemma-descending-fppf-property-immersion}，$Z \to T$ 是浸入。
最后，由于 $\Picardfunctor_{X/B}$ 是 fppf 层，可知 $s|_Z = t|_Z$，且 $Z$
满足上面所述的泛性质。

\medskip\noindent
假设 $s$ 和 $t$ 来自 $X_T$ 上的可逆模 $\mathcal{V}$、$\mathcal{W}$。置
$\mathcal{L} = \mathcal{V} \otimes \mathcal{W}^{\otimes -1}$。
我们要找 $T$ 的局部闭子概形 $Z$，使得 $T' \to T$ 经由 $Z$ 分解，当且仅当
$\mathcal{L}_{X_{T'}}$ 是 $T'$ 上某个可逆层的拉回，见引理
\ref{quot-lemma-flat-geometrically-connected-fibres}。因此 $Z$ 的存在性由
《空间的态射进阶》引理
\ref{spaces-more-morphisms-lemma-diagonal-picard-flat-proper}.
\end{proof}











\section{相对态射}
\label{quot-section-relative-morphisms}

\noindent
我们继续讨论《可表示性判别》第
\ref{criteria-section-relative-morphisms} 节。
在该节中，从概形 $S$ 以及
在 $S$ 上的代数空间态射 $Z \to B$ 和 $X \to B$ 出发，
我们构造了函子
$$
\mathit{Mor}_B(Z, X) : (\Sch/B)^{opp} \longrightarrow \textit{Sets}, \quad
T \longmapsto \{f : Z_T \to X_T\}
$$
我们有时把 $\mathit{Mor}_B(Z, X)$
看作带有态射 $\mathit{Mor}_B(Z, X) \to B$
的函子 $(\Sch/S)^{opp} \to \textit{Sets}$。
也就是说，如果 $T$ 是一个在 $S$ 上的概形，那么
$\mathit{Mor}_B(Z, X)(T)$ 的一个元素是一个二元组 $(f, h)$，
其中 $h$ 是态射 $h : T \to B$，
而 $f : Z \times_{B, h} T \to X \times_{B, h} T$
是 $T$ 上的代数空间态射。特别地，当我们说
$\mathit{Mor}_B(Z, X)$ 是代数空间时，意思是相应的函子
$(\Sch/S)^{opp} \to \textit{Sets}$ 是代数空间。

\begin{lemma}
\label{quot-lemma-Mor-into-Hilb}
设 $S$ 为概形。考虑
在 $S$ 上的代数空间态射 $Z \to B$ 和 $X \to B$。
若 $X \to B$ 是分离的，且 $Z \to B$ 是
有限表示、平坦且真的，
则存在函子之间的自然
单射变换
$$
\mathit{Mor}_B(Z, X) \longrightarrow \Hilbfunctor_{Z \times_B X/B}
$$
它把态射 $f : Z_T \to X_T$ 映到其图。
\end{lemma}

\begin{proof}
给定一个在 $B$ 上的概形 $T$，以及一个
在 $T$ 上的态射 $f_T : Z_T \to X_T$，
$f$ 的图是态射
$\Gamma_f = (\text{id}, f) : Z_T \to Z_T \times_T X_T = (Z \times_B X)_T$。
回忆一下，分离、平坦、真或有限表示
都是代数空间态射的性质，并且在基变换下稳定（《空间的态射》引理
\ref{spaces-morphisms-lemma-base-change-separated},
\ref{spaces-morphisms-lemma-base-change-flat},
\ref{spaces-morphisms-lemma-base-change-proper}，以及
\ref{spaces-morphisms-lemma-base-change-finite-presentation}).
因此，由《空间的态射》引理
\ref{spaces-morphisms-lemma-semi-diagonal}，$\Gamma_f$ 是闭浸入。
此外，$\Gamma_f(Z_T)$ 在 $T$ 上平坦、真且有限表示。
所以 $\Gamma_f(Z_T)$ 给出
$\Hilbfunctor_{Z \times_B X/B}(T)$ 的一个元素。
要证明该变换是单射，只需证明
具有相同图的两个态射相同。这是因为：
若 $Y \subset (Z \times_B X)_T$ 是某个态射 $f$ 的图，
那么可以取 $\text{pr}_1|_Y : Y \to Z_T$ 的逆，
再与 $\text{pr}_2|_Y$ 复合，从而恢复 $f$。
\end{proof}

\begin{lemma}
\label{quot-lemma-Mor-into-Hilb-open}
假设和记号同引理 \ref{quot-lemma-Mor-into-Hilb}。
变换
$\mathit{Mor}_B(Z, X) \longrightarrow \Hilbfunctor_{Z \times_B X/B}$
可由开浸入表示。
\end{lemma}

\begin{proof}
设 $T$ 是在 $B$ 上的概形，并设 $Y \subset (Z \times_B X)_T$
是 $\Hilbfunctor_{Z \times_B X/B}(T)$ 的一个元素。则 $Y$
是 $T$ 上某个态射 $Z_T \to X_T$ 的图，当且仅当
$k = \text{pr}_1|_Y : Y \to Z_T$ 是同构。由《空间态射进阶》引理
\ref{spaces-more-morphisms-lemma-where-isomorphism}
存在开子概形 $V \subset T$，使得对于任意概形态射 $T' \to T$，
$k_{T'} : Y_{T'} \to Z_{T'}$ 是同构，当且仅当
$T' \to T$ 经由 $V$ 分解。
这就证明了引理。
\end{proof}

\begin{proposition}
\label{quot-proposition-Mor}
设 $S$ 为概形。设 $Z \to B$ 和 $X \to B$ 是在 $S$ 上的代数
空间态射。假设 $X \to B$ 有限表示且分离，而
$Z \to B$ 有限表示、平坦且真。则
$\mathit{Mor}_B(Z, X)$ 是在 $B$ 上局部有限
表示的代数空间。
\end{proposition}

\begin{proof}
这是引理 \ref{quot-lemma-Mor-into-Hilb-open}
和命题 \ref{quot-proposition-hilb} 的直接推论。
\end{proof}








\section{代数空间叠}
\label{quot-section-stack-of-spaces}

\noindent
本节延续《Stacks 示例》中以下各节开始的讨论：
\ref{examples-stacks-section-stack-of-spaces},
\ref{examples-stacks-section-stack-of-finite-type-spaces}，以及
\ref{examples-stacks-section-stack-in-groupoids-of-finite-type-spaces}.
在 $\mathbf{Z}$ 上工作时，其中的讨论表明我们有一个群胚叠
$$
p'_{ft} : \Spacesstack'_{ft} \longrightarrow \Sch_{fppf}
$$
它参数化有限型代数空间的（非平坦）族。
更具体地说，$\Spacesstack'_{ft}$ 的一个对象
\footnote{我们总是进行《Stacks 示例》引理
\ref{examples-stacks-lemma-stack-ft-spaces} 中所述的替换。}
是从代数空间 $X$ 到概形 $S$ 的有限型态射 $X \to S$；而态射
$(X' \to S') \to (X \to S)$ 由一对 $(f, g)$ 给出，
其中 $f : X' \to X$ 是代数空间态射，
$g : S' \to S$ 是概形态射，
并且它们嵌入交换图
$$
\xymatrix{
X' \ar[d] \ar[r]_f & X \ar[d] \\
S' \ar[r]^g & S
}
$$
并诱导同构 $X' \to S' \times_S X$；换言之，该图
在代数空间范畴中是笛卡尔的。
函子 $p'_{ft}$ 将 $(X \to S)$ 送到 $S$，将
$(f, g)$ 送到 $g$。我们定义一个充满子范畴
$$
\Spacesstack'_{fp, flat, proper} \subset
\Spacesstack'_{ft}
$$
它由 $\Spacesstack'_{ft}$ 中的对象 $X \to S$ 组成，其中
$X \to S$ 是有限表示、平坦且真的。
记
$$
p'_{fp, flat, proper} :
\Spacesstack'_{fp, flat, proper}
\longrightarrow
\Sch_{fppf}
$$
为函子 $p'_{ft}$ 在上述子范畴上的限制。
我们先回顾上述参考文献中已经得到的结果，
然后开始补充更多结果。

\begin{lemma}
\label{quot-lemma-spaces-fibred-in-groupoids}
范畴 $\Spacesstack'_{ft}$ 在 $\Sch_{fppf}$ 上是群胚纤维化的。
$\Spacesstack'_{fp, flat, proper}$ 亦然。
\end{lemma}

\begin{proof}
对于 $\Spacesstack'_{ft}$，我们已在《Stacks 示例》第
\ref{examples-stacks-section-stack-in-groupoids-of-finite-type-spaces}
节看到这一点，这很容易推出另一种情形的结论。
不过，我们也直接检验《范畴》定义
\ref{categories-definition-fibred-groupoids} 的条件 (1) 和 (2)
来证明它。

\medskip\noindent
条件 (1)。设 $X \to S$ 是 $\Spacesstack'_{ft}$ 的对象，
且 $S' \to S$ 是概形态射。令
$X' = S' \times_S X$。由《空间的态射》引理
\ref{spaces-morphisms-lemma-base-change-finite-type}.
注意 $X' \to S'$ 是有限型的，从而得到一个
位于 $S' \to S$ 之上的态射 $(X' \to S') \to (X \to S)$。
另一种情形同理，使用《空间的态射》引理
\ref{spaces-morphisms-lemma-base-change-finite-presentation},
\ref{spaces-morphisms-lemma-base-change-flat}，以及
\ref{spaces-morphisms-lemma-base-change-proper}.

\medskip\noindent
条件 (2)。考虑 $\Spacesstack'_{ft}$ 中的态射
$(f, g) : (X' \to S') \to (X \to S)$ 和
$(a, b) : (Y \to T) \to (X \to S)$。
给定满足 $g \circ h = b$ 的态射 $h : T \to S'$，
我们需证明存在唯一的
$\Spacesstack'_{ft}$ 中的态射
$(k, h) : (Y \to T) \to (X' \to S')$，使得
$(f, g) \circ (k, h) = (a, b)$。
这由 $X' = S' \times_S X$ 立即得到。
因此，满足条件 (1) 的
$\Spacesstack'_{ft}$ 的任意充满子范畴同样成立。
\end{proof}

\begin{lemma}
\label{quot-lemma-spaces-diagonal}
对角态射
$$
\Delta : \Spacesstack'_{fp, flat, proper} \longrightarrow
\Spacesstack'_{fp, flat, proper} \times \Spacesstack'_{fp, flat, proper}
$$
可由代数空间表示。
\end{lemma}

\begin{proof}
我们将使用《代数叠》引理
\ref{algebraic-lemma-representable-diagonal} 的判据 (2)。
设 $S$ 为概形，且 $X$、$Y$ 是在 $S$ 上有限表示、
在 $S$ 上平坦且在 $S$ 上真的代数空间。
我们需要证明函子
$$
\mathit{Isom}_S(X, Y) : (\Sch/S)_{fppf} \longrightarrow \textit{Sets}, \quad
T \longmapsto \{f : X_T \to Y_T \text{ isomorphism}\}
$$
是代数空间。一个初等论证表明
$\mathit{Isom}_S(X, Y)$ 位于纤维积中：
$$
\xymatrix{
\mathit{Isom}_S(X, Y) \ar[r] \ar[d] & S \ar[d]_{(\text{id}, \text{id})} \\
\mathit{Mor}_S(X, Y) \times \mathit{Mor}_S(Y, X) \ar[r] &
\mathit{Mor}_S(X, X) \times \mathit{Mor}_S(Y, Y)
}
$$
下方的态射把 $(\varphi, \psi)$ 送到
$(\psi \circ \varphi, \varphi \circ \psi)$。
由命题 \ref{quot-proposition-Mor}，下行的函子
是 $S$ 上的代数空间。
于是，由于 $S$ 上代数空间的范畴具有纤维积，
即可得到结论。
\end{proof}

\begin{lemma}
\label{quot-lemma-spaces-stack}
范畴 $\Spacesstack'_{ft}$ 是
$\Sch_{fppf}$ 上的群胚叠。对
$\Spacesstack'_{fp, flat, proper}$ 也同样如此。
\end{lemma}

\begin{proof}
该引理成立的原因可概括为：代数空间的任意 fppf 下降数据
都是有效的，见《Bootstrap》第
\ref{bootstrap-section-applications}.
更准确地说，$\Spacesstack'_{ft}$ 的情形由《Stacks 示例》引理
\ref{examples-stacks-lemma-stack-of-finite-type-spaces}
推出，正如我们在《Stacks 示例》第
\ref{examples-stacks-section-stack-in-groupoids-of-finite-type-spaces}.
不过我们还是回顾证明。我们需要检验《Stacks》定义
(1)、(2) 和 (3) 的条件
\ref{stacks-definition-stack-in-groupoids}.

\medskip\noindent
性质 (1) 已在引理 \ref{quot-lemma-spaces-fibred-in-groupoids} 中看到。

\medskip\noindent
在 $\Spacesstack'_{fp, flat, proper}$ 的情形，性质 (2) 由
引理 \ref{quot-lemma-spaces-diagonal} 得到。
在 $\Spacesstack'_{ft}$ 的情形，它由《Stacks 示例》引理
\ref{examples-stacks-lemma-pre-stack-of-spaces}
得到（这确实是“正确”的参考文献）。

\medskip\noindent
对 $\Spacesstack'_{ft}$，按如下方式检验条件 (3)。设给定
\begin{enumerate}
\item 在 $\Sch_{fppf}$ 中的 fppf 覆盖 $\{U_i \to U\}_{i \in I}$，
\item 对每个 $i \in I$，一个在 $U_i$ 上有限型的代数空间 $X_i$，以及
\item 对每个 $i, j \in I$，一个代数空间同构
$\varphi_{ij} : X_i \times_U U_j \to U_i \times_U X_j$，
它在 $U_i \times_U U_j$ 上，并且在
$U_i \times_U U_j \times_U U_k$ 上满足余循环条件。
\end{enumerate}
我们必须证明存在一个在 $U$ 上有限型的代数空间 $X$，以及
在 $U_i$ 上的同构 $X_{U_i} \cong X_i$，它们恢复
同构 $\varphi_{ij}$。这由
《Bootstrap》引理 \ref{bootstrap-lemma-descend-algebraic-space} 的第 (2) 部分得到。
由《空间上的下降》引理
\ref{spaces-descent-lemma-descending-property-finite-type}
可知 $X \to U$ 是有限型的。在
$\Spacesstack'_{fp, flat, proper}$ 的情形，还要使用
《空间上的下降》引理
\ref{spaces-descent-lemma-descending-property-finite-presentation},
\ref{spaces-descent-lemma-descending-property-flat}，以及
\ref{spaces-descent-lemma-descending-property-proper}
在最后一步得到结论。
\end{proof}

\noindent
核对：叠 $\Spacesstack'_{ft}$ 和
$\Spacesstack'_{fp, flat, proper}$
在代数空间中发挥相同作用。

\begin{lemma}
\label{quot-lemma-extend-spaces-to-spaces}
设 $T$ 是在 $\mathbf{Z}$ 上的代数空间。令 $\mathcal{S}_T$
表示相应的代数叠（《代数叠》第
\ref{algebraic-section-split},
\ref{algebraic-section-representable-by-algebraic-spaces}，以及
\ref{algebraic-section-stacks-spaces}).
我们有范畴等价
$$
\left\{
\begin{matrix}
\text{代数空间态射 }\\
X \to T\text{，且为有限型}
\end{matrix}
\right\}
\longrightarrow
\Mor_{\textit{Cat}/\Sch_{fppf}}(\mathcal{S}_T, \Spacesstack'_{ft})
$$
以及范畴等价
$$
\left\{
\begin{matrix}
\text{代数空间态射 }X \to T\\
\text{，且有限表示、平坦且真}
\end{matrix}
\right\}
\longrightarrow
\Mor_{\textit{Cat}/\Sch_{fppf}}(\mathcal{S}_T,
\Spacesstack'_{fp, flat, proper})
$$
\end{lemma}

\begin{proof}
我们将由该引理对概形成立这一事实（本质上来自叠的构造），
以及代数空间在代数空间上的 fppf 下降数据有效这一事实推出它。
我们强烈建议读者跳过本证明。

\medskip\noindent
任一箭头从左到右的构造都是直接的：给定有限型 $X \to T$，函子
$\mathcal{S}_T \to  \Spacesstack'_{ft}$ 将 $U/T$ 送到
基变换 $X_U \to U$。下面解释如何构造拟逆。

\medskip\noindent
若 $T$ 是概形，则由 $2$-Yoneda 引理存在拟逆，见
《范畴》引理 \ref{categories-lemma-yoneda-2category}。
设 $p : U \to T$ 是满射 \'etale 态射，其中 $U$ 是概形。
令 $R = U \times_T U$，投影为 $s, t : R \to U$。
注意我们得到态射
$$
\xymatrix{
\mathcal{S}_{U \times_T U \times_T U} \ar@<2ex>[r] \ar[r] \ar@<-2ex>[r] &
\mathcal{S}_R \ar@<1ex>[r] \ar@<-1ex>[r] &
\mathcal{S}_U \ar[r] &
\mathcal{S}_T
}
$$
并且它们（严格地）满足各种相容性。

\medskip\noindent
设 $G : \mathcal{S}_T \to \Spacesstack'_{ft}$ 是在 $\Sch_{fppf}$ 上的函子。
通过上面显示的映射，$G$ 在 $\mathcal{S}_U$ 上的限制
由 $2$-Yoneda 引理对应于代数空间的有限型态射 $X_U \to U$。
由于 $p \circ s = p \circ t$，可知
$R \times_{s, U} X_U$ 和 $R \times_{t, U} X_U$ 都对应于
$G$ 在 $\mathcal{S}_R$ 上的限制。因此得到典范同构
$\varphi : X_U \times_{U, t} R \to R \times_{s, U} X_U$（在 $R$ 上）。
由上图的各种相容性，该同构满足余循环条件。
于是得到下降数据；由《Bootstrap》引理
\ref{bootstrap-lemma-descend-algebraic-space} 的第 (2) 部分。
换言之，我们得到右侧范畴中的对象 $X \to T$。
我们略去验证构造 $G \leadsto X$
是函子的，且它是另一构造的拟逆。
在 $\Spacesstack'_{fp, flat, proper}$ 的情形，还要使用
《空间上的下降》引理
\ref{spaces-descent-lemma-descending-property-finite-presentation},
\ref{spaces-descent-lemma-descending-property-flat}，以及
\ref{spaces-descent-lemma-descending-property-proper}
在最后一步得到 $X \to T$ 有限表示、平坦且真。
\end{proof}

\begin{remark}
\label{quot-remark-spaces-base-change}
设 $B$ 是在 $\Spec(\mathbf{Z})$ 上的代数空间。
令 $B\textit{-Spaces}'_{ft}$ 为由以下二元组组成的范畴：
$(X \to S, h : S \to B)$，其中 $X \to S$ 是
$\Spacesstack'_{ft}$ 的对象，且 $h : S \to B$ 是态射。
$B\textit{-Spaces}'_{ft}$ 中的态射 $(X' \to S', h') \to (X \to S, h)$
是 $\Spacesstack'_{ft}$ 中满足 $h \circ g = h'$ 的态射 $(f, g)$。
在此情形，图
$$
\xymatrix{
B\textit{-Spaces}'_{ft} \ar[r] \ar[d] & \Spacesstack'_{ft} \ar[d] \\
(\Sch/B)_{fppf} \ar[r] & \Sch_{fppf}
}
$$
是 $2$-纤维积方图。这个平凡的备注有时可用于从
绝对情形 $\Spacesstack'_{ft}$ 推出
给定基代数空间上的族的情形。
当然，对 $B\textit{-Spaces}'_{fp, flat, proper}$
也有类似构造。
\end{remark}

\begin{lemma}
\label{quot-lemma-spaces-limits}
叠
$p'_{fp, flat, proper} :
\Spacesstack'_{fp, flat, proper} \to \Sch_{fppf}$ 是保极限的
（《Artin 公理》定义 \ref{artin-definition-limit-preserving}）。
\end{lemma}

\begin{proof}
设 $T = \lim T_i$ 是仿射概形的有向逆系统的极限。
由《空间的极限》引理
\ref{spaces-limits-lemma-descend-finite-presentation}
在 $T$ 上有限表示的代数空间范畴是
在 $T_i$ 上有限表示的代数空间范畴的余极限。
最后使用平坦性和真性沿极限下降这一事实，见
《空间的极限》引理
\ref{spaces-limits-lemma-descend-flat} 以及
\ref{spaces-limits-lemma-eventually-proper}.
\end{proof}

\begin{lemma}
\label{quot-lemma-spaces-RS-star}
设
$$
\xymatrix{
T \ar[r] \ar[d] & T' \ar[d] \\
S \ar[r] & S'
}
$$
是概形范畴中的推出，其中
$T \to T'$ 是加厚，且 $T \to S$ 是仿射，见
《空间态射进阶》引理 \ref{more-morphisms-lemma-pushout-along-thickening}。
则纤维范畴上的函子
$$
\begin{matrix}
\Spacesstack'_{fp, flat, proper, S'} \\
\downarrow \\
\Spacesstack'_{fp, flat, proper, S}
\times_{\Spacesstack'_{fp, flat, proper, T}}
\Spacesstack'_{fp, flat, proper, T'}
\end{matrix}
$$
是等价。
\end{lemma}

\begin{proof}
若从条件列表中删去“真”，并将“有限表示”替换为
“局部有限表示”，则该函子是等价；见《空间的推出》引理
\ref{spaces-pushouts-lemma-equivalence-categories-spaces-pushout-flat}.
因此只需证明：给定代数空间到 $S'$ 的态射
$X' \to S'$，它平坦且局部有限表示，则
$X' \to S'$ 为真，当且仅当 $S \times_{S'} X' \to S$
和 $T' \times_{S'} X' \to T'$ 都为真。
一个方向由真性在基变换下保持这一事实得到
（《空间的态射》引理 \ref{spaces-morphisms-lemma-base-change-proper}）；
另一个方向由 $S \times_{S'} X' \to S$ 为真推出
$X' \to S'$ 为真这一事实得到，依据《空间态射进阶》引理
\ref{spaces-more-morphisms-lemma-thicken-property-morphisms-cartesian}.
\end{proof}

\begin{lemma}
\label{quot-lemma-spaces-tangent-space}
设 $k$ 为域，且 $x = (X \to \Spec(k))$ 是
$\mathcal{X} = \Spacesstack'_{fp, flat, proper}$ 在 $\Spec(k)$ 上的对象。
\begin{enumerate}
\item 若 $k$ 在 $\mathbf{Z}$ 上有限型，则向量空间
$T\mathcal{F}_{\mathcal{X}, k, x}$ 和
$\text{Inf}(\mathcal{F}_{\mathcal{X}, k, x})$
（见《Artin 公理》第 \ref{artin-section-tangent-spaces} 节）
是有限维的，并且
\item 一般而言，向量空间 $T_x(k)$ 和 $\text{Inf}_x(k)$
（见《Artin 公理》第 \ref{artin-section-inf} 节）
是有限维的。
\end{enumerate}
\end{lemma}

\begin{proof}
《Artin 公理》第 \ref{artin-section-tangent-spaces} 节的讨论
只适用于在基概形 $\Spec(\mathbf{Z})$ 上有限型的域。
由引理 \ref{quot-lemma-spaces-RS-star}，我们的叠满足 (RS*)，
因此可以应用
《Artin 公理》引理 \ref{artin-lemma-properties-lift-RS-star}
得到 (2) 中提到的向量空间 $T_x(k)$ 和 $\text{Inf}_x(k)$。
此外，在有限型情形，这些空间与 (1) 中提到的空间相同，
依据《Artin 公理》注记 \ref{artin-remark-compare-deformation-spaces}。
处理完这些后，我们开始证明。
注意一阶加厚
$\Spec(k) \to \Spec(k[\epsilon]) = \Spec(k[k])$
的余法模为 $k$。因此《形变理论》引理
\ref{defos-lemma-deform-spaces} 中描述 $X$ 的无穷小形变
和 $X$ 的无穷小自同构的公式变为
$$
T_x(k) = \Ext^1_{\mathcal{O}_X}(\NL_{X/k}, \mathcal{O}_X)
\quad\text{且}\quad
\text{Inf}_x(k) = \Ext^0_{\mathcal{O}_X}(\NL_{X/k}, \mathcal{O}_X)
$$
由《空间态射进阶》引理
\ref{spaces-more-morphisms-lemma-netherlander-fp}
以及 $X$ 是 Noetherian 这一事实，可知
$\NL_{X/k}$ 的相干上同调层除
次数 $0$ 和 $-1$ 外均为零。
由《空间的导出范畴》引理 \ref{spaces-perfect-lemma-ext-finite}，
上述 $\Ext$-群是有限维的 $k$-向量空间，
证明完成。
\end{proof}

\noindent
请注意，正如下一个引理所证明的，我们的叠不满足形式有效性，
所以形式光滑性的开放性有些奇怪，见《Examples》一节
\ref{examples-section-proper-spaces-not-algebraic}。
稍后我们将把形式光滑性的开放性应用于
满足形式有效性的 $\Spacesstack'_{fp, flat, proper}$ 的适当子叠，
从而证明这些叠是代数的。

\begin{lemma}
\label{quot-lemma-spaces-defo-thy}
群胚叠 $\mathcal{X} = \Spacesstack'_{fp, flat, proper}$
在 $\Spec(\mathbf{Z})$ 上满足形式光滑性的开放性。
同样地，基变换后（注记 \ref{quot-remark-spaces-base-change}），
形式光滑性的开放性对任意 Noether 基概形 $S$ 成立。
\end{lemma}

\begin{proof}
关于这一事实的“通常”证明，请参见本证明后的注记。
这里我们使用《Artin 公理》引理
\ref{artin-lemma-SGE-implies-openness-versality} 来证明。
我们已经看到 $\mathcal{X}$ 的对角可由代数空间表示，
满足 (RS*) 且保极限，见引理
\ref{quot-lemma-spaces-diagonal}、\ref{quot-lemma-spaces-RS-star} 以及
\ref{quot-lemma-spaces-limits}。
因此只需证明 $\mathcal{X}$ 满足《Artin 公理》引理
\ref{artin-lemma-SGE-implies-openness-versality} 中表述的强
形式有效性。

\medskip\noindent
设 $(R_n)$ 是环的逆系统，满足对所有 $n \geq m$，
$R_n \to R_m$ 是核平方为零的满射。设
$X_n \to \Spec(R_n)$ 是有限表示、平坦、真的态射，
其中 $X_n$ 是代数空间；并设 $X_{n + 1} \to X_n$ 是在
$\Spec(R_{n + 1})$ 上的态射，诱导同构
$X_n = X_{n + 1} \times_{\Spec(R_{n + 1})} \Spec(R_n)$。
我们需要找到一个平坦、真且有限表示的态射
$X \to \Spec(\lim R_n)$，其源是代数空间，
并且对所有 $n$，$X_n$ 都是 $X$ 的基变换。

\medskip\noindent
令 $I_n = \Ker(R_n \to R_1)$。可将
$(X_1 \subset X_n) \to (\Spec(R_1) \subset \Spec(R_n))$
看作一阶加厚的态射。（继续之前，请阅读《空间态射进阶》
关于代数空间加厚的部分，即第
\ref{spaces-more-morphisms-section-thickenings} 节。）$X_n$ 的结构层是扩张
$$
0 \to \mathcal{O}_{X_1} \otimes_{R_1} I_n \to
\mathcal{O}_{X_n} \to \mathcal{O}_{X_1} \to 0
$$
它位于 $0 \to I_n \to R_n \to R_1$ 之上，见
《空间态射进阶》引理
\ref{spaces-more-morphisms-lemma-deform}.
考虑扩张
$$
0 \to \lim \mathcal{O}_{X_1} \otimes_{R_1} I_n \to
\lim \mathcal{O}_{X_n} \to \mathcal{O}_{X_1} \to 0
$$
它位于 $0 \to \lim I_n \to \lim R_n \to R_1 \to 0$ 之上。
上述序列是正合的，因为由《空间的导出范畴》引理
可知核系统的 $R^1\lim$ 为零
\ref{spaces-perfect-lemma-Rlim-quasi-coherent}.
注意映射
$$
\mathcal{O}_{X_1} \otimes_{R_1} \lim I_n \longrightarrow
\lim \mathcal{O}_{X_1} \otimes_{R_1} I_n
$$
在应用函子 $DQ_X$ 后诱导同构，见《空间的导出范畴》引理
\ref{spaces-perfect-lemma-pullback-and-limits}.
因此由《形变理论》引理
\ref{defos-lemma-thickening-over-thickening-space-quasi-coherent} 的范畴等价，得到唯一扩张
$$
0 \to \mathcal{O}_{X_1} \otimes_{R_1} \lim I_n \to
\mathcal{O}' \to \mathcal{O}_{X_1} \to 0
$$
在 $0 \to \lim I_n \to \lim R_n \to R_1 \to 0$ 之上。
层 $\mathcal{O}'$ 由
$\Spec(R_1) \subset \Spec(\lim R_n)$ 上代数空间
$X_1 \subset X$ 的一阶加厚决定，依据《空间态射进阶》引理
\ref{spaces-more-morphisms-lemma-first-order-thickening}。
注意，由此前使用的《空间态射进阶》引理
\ref{spaces-more-morphisms-lemma-deform}，
$X \to \Spec(\lim R_n)$ 是平坦的。
由《空间态射进阶》引理
\ref{spaces-more-morphisms-lemma-deform-property}，可知
$X \to \Spec(\lim R_n)$ 是真且有限表示的。
证明完成。
\end{proof}

\begin{remark}
\label{quot-remark-spaces-defo-thy}
引理 \ref{quot-lemma-spaces-defo-thy} 也可用以下任一方法证明：
《Artin 公理》引理 \ref{artin-lemma-dual-openness}
（如同引理 \ref{quot-lemma-coherent-defo-thy} 的第一个证明），
或使用《Artin 公理》引理 \ref{artin-lemma-get-openness-obstruction-theory}
中的障碍理论（如同引理 \ref{quot-lemma-coherent-defo-thy} 的第二个证明）。
两种方法都使用《余切复形》第 \ref{cotangent-section-deformations-ringed-topoi} 节中
发展的形变与障碍理论，将形变和障碍所需的性质
转化为可应用《空间的导出范畴》引理
\ref{spaces-perfect-lemma-compute-ext} 的 $\Ext$-群。
第二种方法（使用障碍理论，因而使用完整余切复形）
或许是大多数参考文献采用的“标准”方法。
\end{remark}





\section{极化真概形叠}
\label{quot-section-polarized}

\noindent
研究极化真概形的叠，只需在 $\mathbf{Z}$ 上工作，
因为随后可以将其拉回到任意所需的概形或代数空间
（见注记 \ref{quot-remark-polarized-base-change}）。

\begin{situation}
\label{quot-situation-polarized}
我们如下定义范畴 $\Polarizedstack$。对象是二元组
$(X \to S, \mathcal{L})$，其中
\begin{enumerate}
\item $X \to S$ 是真、平坦且
有限表示的概形态射，并且
\item $\mathcal{L}$ 是可逆的 $\mathcal{O}_X$-模，
在 $X/S$ 上相对充足，
（《态射》定义 \ref{morphisms-definition-relatively-ample}）。
\end{enumerate}
对象之间的态射 $(X' \to S', \mathcal{L}') \to (X \to S, \mathcal{L})$
由三元组 $(f, g, \varphi)$ 给出，其中 $f : X' \to X$ 和 $g : S' \to S$
是嵌入交换图的概形态射
$$
\xymatrix{
X' \ar[d] \ar[r]_f & X \ar[d] \\
S' \ar[r]^g & S
}
$$
并诱导同构 $X' \to S' \times_S X$；换言之，该图是笛卡尔的；
且 $\varphi : f^*\mathcal{L} \to \mathcal{L}'$ 是同构。
复合按显然方式定义（见
《叠的例子》第
\ref{examples-stacks-section-stack-of-spaces} 节以及
第 \ref{examples-stacks-section-stack-of-quasi-coherent-sheaves} 节）。
遗忘函子
$$
p : \Polarizedstack \longrightarrow \Sch_{fppf},\quad
(X \to S, \mathcal{L}) \longmapsto S
$$
就是我们把 $\Polarizedstack$ 看作 $\Sch_{fppf}$ 上范畴的方式
（记号见第 \ref{quot-section-conventions} 节）。
\end{situation}

\noindent
在上一节中，我们已经对叠
$\Spacesstack'_{fp, flat, proper}$（由有限表示、平坦、真的
代数空间组成）做了大量工作。为使用这些材料，考虑遗忘函子
\begin{equation}
\label{quot-equation-over-proper-spaces}
\Polarizedstack \longrightarrow
\Spacesstack'_{fp, flat, proper},\quad
(X \to S, \mathcal{L}) \longmapsto (X \to S)
\end{equation}
下面将使用这个函子作为有用工具。
注意，若 $(X \to S)$ 属于
（\ref{quot-equation-over-proper-spaces}）的本质像，则
$X$ 和 $S$ 是概形。

\begin{lemma}
\label{quot-lemma-polarized-fibred-in-groupoids}
范畴 $\Polarizedstack$ 在
$\Spacesstack'_{fp, flat, proper}$ 上是群胚纤维化的。
范畴 $\Polarizedstack$ 在 $\Sch_{fppf}$ 上是群胚纤维化的。
\end{lemma}

\begin{proof}
我们检验《范畴》定义 \ref{categories-definition-fibred-groupoids} 的条件 (1) 和 (2)。

\medskip\noindent
条件 (1)。设 $(X \to S, \mathcal{L})$ 是
$\Polarizedstack$ 的对象，且 $(X' \to S') \to (X \to S)$
是 $\Spacesstack'_{fp, flat, proper}$ 中的态射。令
$\mathcal{L}'$ 为 $\mathcal{L}$ 到 $X'$ 的拉回。
注意 $X, S, S'$ 是概形，因此 $X'$ 也是概形
（概形的纤维积）。由此
由《态射》引理 \ref{morphisms-lemma-ample-base-change}，
$\mathcal{L}'$ 在 $X'/S'$ 上充足。
这样得到态射
$(X' \to S', \mathcal{L}') \to (X \to S, \mathcal{L})$，
它位于 $(X' \to S') \to (X \to S)$ 之上。

\medskip\noindent
条件 (2)。考虑 $\Polarizedstack$ 中的态射
$(f, g, \varphi) : (X' \to S', \mathcal{L}') \to (X \to S, \mathcal{L})$ 和
$(a, b, \psi) : (Y \to T, \mathcal{N}) \to (X \to S, \mathcal{L})$。
给定 $\Spacesstack'_{fp, flat, proper}$ 中的态射
$(k, h) : (Y \to T) \to (X' \to S')$，且
$(f, g) \circ (k, h) = (a, b)$，我们需证明存在唯一态射
$(k, h, \chi) : (Y \to T, \mathcal{N}) \to (X' \to S', \mathcal{L}')$
属于 $\Polarizedstack$，并满足
$(f, g, \varphi) \circ (k, h, \chi) = (a, b, \psi)$。
只需取
$$
\chi = \psi \circ (k^*\varphi)^{-1}
$$
这证明了条件 (2)。定义纤维化范畴的函子的复合
定义纤维化范畴，见《范畴》引理
\ref{categories-lemma-fibred-over-fibred}。
因此 $\Polarizedstack$ 在 $\Sch_{fppf}$ 上是群胚纤维化的
（严格地说，还应检验纤维范畴是群胚，并应用
《范畴》引理 \ref{categories-lemma-fibred-groupoids}）。
\end{proof}

\begin{lemma}
\label{quot-lemma-polarized-stack}
范畴 $\Polarizedstack$ 是
$\Spacesstack'_{fp, flat, proper}$ 上的群胚叠（赋予继承拓扑，
见《Stacks》定义 \ref{stacks-definition-topology-inherited}）。
范畴 $\Polarizedstack$ 是 $\Sch_{fppf}$ 上的群胚叠。
\end{lemma}

\begin{proof}
我们通过检验《Stacks》定义
\ref{stacks-definition-stack-in-groupoids} 的条件 (1)、(2)、(3)，
证明 $\Polarizedstack$ 是
$\Spacesstack'_{fp, flat, proper}$ 上的群胚叠。
条件 (1) 已在引理
\ref{quot-lemma-polarized-fibred-in-groupoids} 中看到。

\medskip\noindent
$\Spacesstack'_{fp, flat, proper}$ 的覆盖如下产生：设 $X \to S$ 是
$\Spacesstack'_{fp, flat, proper}$ 的对象，并设
$\{S_i \to S\}_{i \in I}$ 是 $\Sch_{fppf}$ 的覆盖。
令 $X_i = S_i \times_S X$。则
$\{(X_i \to S_i) \to (X \to S)\}_{i \in I}$ 是
$\Spacesstack'_{fp, flat, proper}$ 的覆盖，并且
$\Spacesstack'_{fp, flat, proper}$ 的每个覆盖都同构于其中之一。
令 $S_{ij} = S_i \times_S S_j$，并令
$X_{ij} = S_{ij} \times_S X$，于是
$(X_{ij} \to S_{ij}) = (X_i \to S_i) \times_{(X \to S)} (X_j \to S_j)$。
接下来，设 $\mathcal{L}, \mathcal{N}$
是 $X/S$ 上的充足可逆层，使得
$(X \to S, \mathcal{L})$ 和 $(X \to S, \mathcal{N})$
是位于对象 $(X \to S)$ 之上的 $\Polarizedstack$ 的两个对象。
为检验态射的下降，假设给定态射
$(\text{id}, \text{id}, \varphi_i)$，从
$(X_i \to S_i, \mathcal{L}|_{X_i})$ 到
$(X_i \to S_i, \mathcal{N}|_{X_i})$，
其到态射
$(X_{ij} \to S_{ij}, \mathcal{L}|_{X_{ij}})$ 到
$(X_{ij} \to S_{ij}, \mathcal{N}|_{X_{ij}})$ 的基变换
彼此一致。则
$\varphi_i : \mathcal{L}|_{X_i} \to \mathcal{N}|_{X_i}$
是 $X_i$ 上可逆模的同构，并且
$\varphi_i$ 与 $\varphi_j$ 在 $X_{ij}$ 上限制为相同同构。
由拟凝聚层的下降（《空间上的下降》命题
\ref{spaces-descent-proposition-fpqc-descent-quasi-coherent})
得到唯一同构 $\varphi : \mathcal{L} \to \mathcal{N}$，
其在 $X_i$ 上的限制恢复 $\varphi_i$。

\medskip\noindent
对象的下降证明完全相同。
具体地，设
$\{(X_i \to S_i) \to (X \to S)\}_{i \in I}$
是如上的 $\Spacesstack'_{fp, flat, proper}$ 覆盖。
设有对象 $(X_i \to S_i, \mathcal{L}_i)$，
它们是位于 $(X_i \to S_i)$ 之上的 $\Polarizedstack$ 对象，
并有下降数据
$$
(\text{id}, \text{id}, \varphi_{ij}) :
(X_{ij} \to S_{ij}, \mathcal{L}_i|_{X_{ij}})
\to
(X_{ij} \to S_{ij}, \mathcal{L}_j|_{X_{ij}})
$$
对每个指标三元组，它在
$(X_{ijk} \to S_{ijk})$ 上满足显然的余循环条件。
于是由
拟凝聚层的下降（《空间上的下降》命题
\ref{spaces-descent-proposition-fpqc-descent-quasi-coherent})
得到唯一可逆 $\mathcal{O}_X$-模
$\mathcal{L}$ 以及同构 $\mathcal{L}|_{X_i} \to \mathcal{L}_i$，
恢复下降数据 $\varphi_{ij}$。
要证明 $(X \to S, \mathcal{L})$ 是
$\Polarizedstack$ 的对象，需证明
$\mathcal{L}$ 是充足的。这由
《空间上的下降》引理 \ref{spaces-descent-lemma-descending-property-ample} 得到。

\medskip\noindent
由于我们已经看到 $\Spacesstack'_{fp, flat, proper}$
是 $\Sch_{fppf}$ 上的群胚叠（引理 \ref{quot-lemma-spaces-stack}），
形式上即可推出 $\Polarizedstack$ 是 $\Sch_{fppf}$ 上的群胚叠。
见《Stacks》引理 \ref{stacks-lemma-stack-over-stack}。
\end{proof}

\noindent
核对：叠 $\Polarizedstack$ 在代数空间中发挥相同作用。

\begin{lemma}
\label{quot-lemma-extend-polarized-to-spaces}
设 $T$ 是在 $\mathbf{Z}$ 上的代数空间。令 $\mathcal{S}_T$
表示相应的代数叠（《代数叠》第
\ref{algebraic-section-split},
\ref{algebraic-section-representable-by-algebraic-spaces}，以及
\ref{algebraic-section-stacks-spaces}).
我们有范畴等价
$$
\left\{
\begin{matrix}
(X \to T, \mathcal{L})\text{，其中}\\
X \to T\text{ 是代数空间态射，}\\
\text{它真、平坦且有限表示，并且}\\
\mathcal{L}\text{ 在 }X/T\text{ 上充足}
\end{matrix}
\right\}
\longrightarrow
\Mor_{\textit{Cat}/\Sch_{fppf}}(\mathcal{S}_T, \Polarizedstack)
$$
\end{lemma}

\begin{proof}
略去。提示：完全仿照引理
\ref{quot-lemma-extend-spaces-to-spaces} 的证明，并使用
《空间上的下降》命题
\ref{spaces-descent-proposition-fpqc-descent-quasi-coherent}
在拟逆函子的构造中下降可逆层。相对充足性沿下降，依据
《空间上的下降》引理
\ref{spaces-descent-lemma-descending-property-ample}.
\end{proof}

\begin{remark}
\label{quot-remark-polarized-base-change}
设 $B$ 是在 $\Spec(\mathbf{Z})$ 上的代数空间。
令 $B\textit{-Polarized}$ 为由三元组
$(X \to S, \mathcal{L}, h : S \to B)$ 组成的范畴，
其中 $(X \to S, \mathcal{L})$ 是 $\Polarizedstack$ 的对象，
且 $h : S \to B$ 是态射。
$B\textit{-Polarized}$ 中的态射 $(X' \to S', \mathcal{L}', h') \to (X \to S, \mathcal{L}, h)$
是 $\Polarizedstack$ 中满足 $h \circ g = h'$ 的态射 $(f, g, \varphi)$。
在此情形，图
$$
\xymatrix{
B\textit{-Polarized} \ar[r] \ar[d] & \Polarizedstack \ar[d] \\
(\Sch/B)_{fppf} \ar[r] & \Sch_{fppf}
}
$$
是 $2$-纤维积方图。这个平凡备注有时可用于从
绝对情形 $\Polarizedstack$ 推出给定基代数空间上的
族的情形。
\end{remark}

\begin{lemma}
\label{quot-lemma-polarized-to-spaces-algebraic}
函子（\ref{quot-equation-over-proper-spaces}）定义一个 $1$-态射
$$
\Polarizedstack \to \Spacesstack'_{fp, flat, proper}
$$
它是 $\Sch_{fppf}$ 上群胚叠的态射，且依《可表示性判别》定义
\ref{criteria-definition-algebraic} 的意义是代数的。
\end{lemma}

\begin{proof}
由引理 \ref{quot-lemma-spaces-stack} 和 \ref{quot-lemma-polarized-stack}，
该陈述有意义。为证明它，取概形 $S$，并取对象
$\xi = (X \to S)$，它属于 $\Spacesstack'_{fp, flat, proper}$，
且位于 $S$ 上。需证明
$$
\mathcal{X} = (\Sch/S)_{fppf} \times_{\xi, \Spacesstack'_{fp, flat, proper}}
\Polarizedstack
$$
是 $S$ 上的代数叠。注意，$\mathcal{X}$ 的对象由二元组
$(T/S, \mathcal{L})$ 给出，其中 $T$ 是 $S$ 上的概形，
$\mathcal{L}$ 是可逆 $\mathcal{O}_{X_T}$-模，且在 $X_T/T$ 上充足。
态射按显然方式定义。特别地，立即得到嵌入
$$
\mathcal{X} \subset \Picardstack_{X/S}
$$
它是 $(\Sch/S)_{fppf}$ 上范畴的嵌入，并在态射集合上诱导等式。
由命题 \ref{quot-proposition-pic}，$\Picardstack_{X/S}$ 是代数叠，
所以只需证明上述嵌入可由开浸入表示。这正是《空间上的下降》引理
\ref{spaces-descent-lemma-ample-in-neighbourhood} 的内容。
\end{proof}

\begin{lemma}
\label{quot-lemma-polarized-diagonal}
对角态射
$$
\Delta : \Polarizedstack \longrightarrow
\Polarizedstack \times \Polarizedstack
$$
由代数空间表示。
\end{lemma}

\begin{proof}
这是引理 \ref{quot-lemma-polarized-to-spaces-algebraic} 和
\ref{quot-lemma-spaces-diagonal} 的形式推论。
见《可表示性判别》引理
\ref{criteria-lemma-diagonals-and-algebraic-morphisms}。
\end{proof}

\begin{lemma}
\label{quot-lemma-polarized-limits}
群胚叠 $\Polarizedstack$ 保持极限
（《Artin 公理》定义 \ref{artin-definition-limit-preserving}）。
\end{lemma}

\begin{proof}
设 $I$ 是有向集，且 $(A_i, \varphi_{ii'})$ 是 $I$ 上的环系统。
令 $S = \Spec(A)$ 以及 $S_i = \Spec(A_i)$。需证明纤维范畴满足
$$
\Polarizedstack_S = \colim \Polarizedstack_{S_i}
$$
我们知道，$S$ 上有限表示概形的范畴是
$S_i$ 上有限表示概形范畴的余极限，见《极限》引理
\ref{limits-lemma-descend-finite-presentation}。
此外，给定有限表示的 $X_i \to S_i$，其极限为 $X \to S$，则
可逆 $\mathcal{O}_X$-模 $\mathcal{L}$ 的范畴是可逆
$\mathcal{O}_{X_i}$-模 $\mathcal{L}_i$ 范畴的余极限，见《极限》引理
\ref{limits-lemma-descend-modules-finite-presentation} 和
\ref{limits-lemma-descend-invertible-modules}。
若 $X \to S$ 真且平坦，则当 $i$ 充分大时，态射
$X_i \to S_i$ 也真且平坦，见《极限》引理
\ref{limits-lemma-eventually-proper} 和
\ref{limits-lemma-descend-flat-finite-presentation}。
最后，若 $\mathcal{L}$ 在 $X$ 上充足，则当 $i$ 充分大时，
$\mathcal{L}_i$ 在 $X_i$ 上充足，见《极限》引理
\ref{limits-lemma-limit-ample}。合并这些事实即完成证明。
\end{proof}

\begin{lemma}
\label{quot-lemma-polarized-RS-star}
在情形 \ref{quot-situation-coherent} 中，设
$$
\xymatrix{
T \ar[r] \ar[d] & T' \ar[d] \\
S \ar[r] & S'
}
$$
是概形范畴中的推出，其中 $T \to T'$ 是加厚，且 $T \to S$ 仿射；见
《态射进阶》引理 \ref{more-morphisms-lemma-pushout-along-thickening}。
则纤维范畴上的函子
$$
\Polarizedstack_{S'}
\longrightarrow
\Polarizedstack_S \times_{\Polarizedstack_T} \Polarizedstack_{T'}
$$
是等价。
\end{lemma}

\begin{proof}
由《态射进阶》引理
\ref{more-morphisms-lemma-equivalence-categories-schemes-over-pushout-flat}
有等价
$$
\textit{flat-lfp}_{S'}
\longrightarrow
\textit{flat-lfp}_S \times_{\textit{flat-lfp}_T} \textit{flat-lfp}_{T'}
$$
$\textit{flat-lfp}_S$ 表示在 $S$ 上平坦且局部有限表示的概形范畴。
设左侧的 $X'/S'$ 对应于右侧的三元组
$(X/S, Y'/T', \varphi)$。
令 $Y = T \times_{T'} Y'$；通过 $\varphi$，它同构于
$T \times_S X$。于是《态射进阶》引理
\ref{more-morphisms-lemma-scheme-over-pushout-flat-modules}
给出等价
$$
\textit{QCoh-flat}_{X'/S'}
\longrightarrow
\textit{QCoh-flat}_{X/S}
\times_{\textit{QCoh-flat}_{Y/T}} \textit{QCoh-flat}_{Y'/T'}
$$
$\textit{QCoh-flat}_{X/S}$ 表示由在 $S$ 上平坦的
拟凝聚 $\mathcal{O}_X$-模组成的范畴。
由于 $X \to S$、$Y \to T$、$X' \to S'$、$Y' \to T'$ 均平坦，
这尤其适用于可逆模，从而给出范畴等价
$$
\textit{Pic}(X')
\longrightarrow
\textit{Pic}(X) \times_{\textit{Pic}(Y)} \textit{Pic}(Y')
$$
$\textit{Pic}(X)$ 表示可逆 $\mathcal{O}_X$-模的范畴。
这里有一点需要说明：若 $\textit{QCoh-flat}_{X'/S'}$ 中的对象
$\mathcal{F}'$ 在 $X$ 和 $Y'$ 上的拉回都是可逆模，
则必须证明 $\mathcal{F}'$ 是可逆 $\mathcal{O}_{X'}$-模。
由上述引理，$\mathcal{F}'$ 是有限表示的
$\mathcal{O}_{X'}$-模。根据《态射进阶》引理
\ref{more-morphisms-lemma-flat-and-free-at-point-fibre}
只需检验 $\mathcal{F}'$ 在 $X' \to S'$ 各纤维上的限制是可逆的。
但 $X' \to S'$ 的纤维与 $X \to S$ 的纤维相同，
故这些限制确实可逆。

\medskip\noindent
由以上论证可知，即使先去掉（在 $S$ 上对象的范畴中）
$X \to S$ 为真的条件以及 $\mathcal{L}$ 为充足的条件，
我们仍得到范畴等价。显然，若 $X' \to S'$ 为真，则
$X \to S$ 和 $Y' \to T'$ 都为真（《态射》引理
\ref{morphisms-lemma-base-change-proper}).
反之，若 $X \to S$ 与 $Y' \to T'$ 为真，则
$X' \to S'$ 也为真；见《态射进阶》引理
\ref{more-morphisms-lemma-thicken-property-morphisms-cartesian}。
类似地，若 $\mathcal{L}'$ 在 $X'/S'$ 上充足，
则 $\mathcal{L}'|_X$ 在 $X/S$ 上充足，且
$\mathcal{L}'|_{Y'}$ 在 $Y'/T'$ 上充足
（《态射》引理 \ref{morphisms-lemma-ample-base-change}）。
最后，若 $\mathcal{L}'|_X$ 在 $X/S$ 上充足，且
$\mathcal{L}'|_{Y'}$ 在 $Y'/T'$ 上充足，则
$\mathcal{L}'$ 在 $X'/S'$ 上充足；见《态射进阶》引理
\ref{more-morphisms-lemma-thicken-property-relatively-ample}。
\end{proof}

\begin{lemma}
\label{quot-lemma-polarized-tangent-space}
设 $k$ 为域，且 $x = (X \to \Spec(k), \mathcal{L})$
是 $\Spec(k)$ 上 $\mathcal{X} = \Polarizedstack$ 的对象。
\begin{enumerate}
\item 若 $k$ 在 $\mathbf{Z}$ 上有限型，则向量空间
$T\mathcal{F}_{\mathcal{X}, k, x}$ 和
$\text{Inf}(\mathcal{F}_{\mathcal{X}, k, x})$
（见《Artin 公理》第 \ref{artin-section-tangent-spaces} 节）
均为有限维；
\item 一般而言，向量空间 $T_x(k)$ 和 $\text{Inf}_x(k)$
（见《Artin 公理》第 \ref{artin-section-inf} 节）均为有限维。
\end{enumerate}
\end{lemma}

\begin{proof}
《Artin 公理》第 \ref{artin-section-tangent-spaces} 节中的讨论
只适用于在基概形 $\Spec(\mathbf{Z})$ 上有限型的域。
由引理 \ref{quot-lemma-polarized-RS-star}，我们的叠满足 (RS*)，
因此可应用《Artin 公理》引理
\ref{artin-lemma-properties-lift-RS-star}，得到 (2) 中所述的向量空间
$T_x(k)$ 和 $\text{Inf}_x(k)$。此外，在有限型情形，
由《Artin 公理》注记 \ref{artin-remark-compare-deformation-spaces}，
这些空间与 (1) 中所述的空间相同。现在开始证明。

\medskip\noindent
一种证明可以仿照引理 \ref{quot-lemma-spaces-tangent-space} 的证明；
这要求我们为由一个概形和一个拟凝聚模组成的配对发展形变理论。
另一种证明可以利用引理 \ref{quot-lemma-spaces-tangent-space} 的结果、
$\Polarizedstack \to \Spacesstack'_{fp, flat, proper}$,
的代数性，以及对可逆模形变空间的计算。
不过，我们将把问题转化为分次 $k$-代数上的形变问题，
并由此推出结论。

\medskip\noindent
令 $\mathcal{C}_k$ 为剩余域是 $k$ 的 Artin 局部 $k$-代数
$A$ 所成的范畴。由 $k$ 上 $\mathcal{X}$ 的对象 $x$，
得到预形变范畴 $p : \mathcal{F} \to \mathcal{C}_k$；见
《Artin 公理》第 \ref{artin-section-predeformation-categories} 节。
因此 $\mathcal{F}(A)$ 是三元组
$(X_A, \mathcal{L}_A, \alpha)$ 所成的范畴，其中
$(X_A, \mathcal{L}_A)$ 是 $A$ 上 $\Polarizedstack$ 的对象，
且 $\alpha$ 是同构
$(X_A, \mathcal{L}_A) \times_{\Spec(A)} \Spec(k) \cong (X, \mathcal{L})$.
另一方面，令 $q : \mathcal{G} \to \mathcal{C}_k$ 为
《形变问题》例 \ref{examples-defos-example-graded-algebras}
所定义的群胚上余纤维化范畴。
取 $d_0 \gg 0$（下文会说明需要多大）。令 $P$ 为分次 $k$-代数
$$
P = k \oplus \bigoplus\nolimits_{d \geq d_0} H^0(X, \mathcal{L}^{\otimes d})
$$
则 $y = (k, P)$ 是 $\mathcal{G}(k)$ 的对象。
令 $\mathcal{G}_y$ 为《形式形变理论》注记
\ref{formal-defos-remark-localize-cofibered-groupoid}.
中的预形变范畴。给定如上的 $(X_A, \mathcal{F}_A, \alpha)$，令
$$
Q = A \oplus \bigoplus\nolimits_{d \geq d_0} H^0(X_A, \mathcal{L}_A^{\otimes d})
$$
同构 $\alpha$ 诱导映射 $\beta : Q \to P$。
由射影概形的形变理论
（《态射进阶》引理 \ref{more-morphisms-lemma-deform-projective}），
得到群胚上余纤维化范畴的 $1$-态射
$$
\mathcal{F} \longrightarrow \mathcal{G}_y,\quad
(X_A, \mathcal{F}_A, \alpha) \longmapsto (Q, \beta : Q \to P)
$$
它们都位于 $\mathcal{C}_k$ 上。事实上，此函子是等价，
其拟逆由 $Q \mapsto \underline{\text{Proj}}_A(Q)$ 给出。
具体而言，由《除子》引理 \ref{divisors-lemma-relative-proj-flat}，
概形 $X_A = \underline{\text{Proj}}_A(Q)$ 在 $A$ 上平坦。
令 $\mathcal{L}_A = \mathcal{O}_{X_A}(1)$；由同一引理，
它在 $A$ 上平坦。由 $\beta$ 得到同构
$(X_A, \mathcal{L}_A) \times_{\Spec(A)} \Spec(k) = (X, \mathcal{L})$
。随后利用《态射进阶》第
\ref{more-morphisms-section-morphisms-thickenings} 节和第
\ref{more-morphisms-section-deform} 节中的方法，
可由 $(X, \mathcal{L})$ 的相应性质推出配对
$(X_A, \mathcal{L}_A)$ 所需的全部性质。略去若干细节。

\medskip\noindent
综上，$T\mathcal{F} = T\mathcal{G}_y = T_y\mathcal{G}$，且
$\text{Inf}(\mathcal{F}) = \text{Inf}_y(\mathcal{G})$。
由《形变问题》引理 \ref{examples-defos-lemma-graded-algebras-TI}，
这些向量空间都是有限维的，证明完毕。
\end{proof}

\begin{lemma}[极化概形的强形式有效性]
\label{quot-lemma-polarized-strong-effectiveness}
\begin{slogan}
若假设平坦且有限表示，则 Grothendieck 代数化定理
在非 Noether 情形仍然成立。
\end{slogan}
设 $(R_n)$ 是环的逆系统，其转移映射满射且核局部幂零。
令 $R = \lim R_n$。令 $S_n = \Spec(R_n)$ 及
$S = \Spec(R)$。考虑概形的交换图
$$
\xymatrix{
X_1 \ar[r]_{i_1} \ar[d] & X_2 \ar[r]_{i_2} \ar[d] & X_3 \ar[r] \ar[d] &
\ldots \\
S_1 \ar[r] & S_2 \ar[r] & S_3 \ar[r] & \ldots
}
$$
其中各方块均为笛卡儿方块。设给定 $(\mathcal{L}_n, \varphi_n)$，
其中每个 $\mathcal{L}_n$ 都是 $X_n$ 上的可逆层，且
$\varphi_n : i_n^*\mathcal{L}_{n + 1} \to \mathcal{L}_n$ 是同构。若
\begin{enumerate}
\item $X_n \to S_n$ 真、平坦且有限表示；
\item $\mathcal{L}_1$ 在 $X_1$ 上充足，
\end{enumerate}
则存在真、平坦且有限表示的概形态射 $X \to S$，
存在充足可逆 $\mathcal{O}_X$-模 $\mathcal{L}$，
并存在同构 $X_n \cong X \times_S S_n$ 及
$\mathcal{L}_n \cong \mathcal{L}|_{X_n}$；
这些同构与态射 $i_n$ 和 $\varphi_n$ 相容。
\end{lemma}

\begin{proof}
依照《态射进阶》引理 \ref{more-morphisms-lemma-deform-projective}，
为 $X_1 \to S_1$ 和 $\mathcal{L}_1$ 选取 $d_0$。
对任意 $n \geq 1$，令
$$
A_n = R_n \oplus
\bigoplus\nolimits_{d \geq d_0} H^0(X_n, \mathcal{L}_n^{\otimes d})
$$
由该引理，每个 $A_n$ 都是有限表示的分次 $R_n$-代数，
其齐次部分 $(A_n)_d$ 是有限射影 $R_n$-模，并且
$X_n = \text{Proj}(A_n)$ 且
$\mathcal{L}_n = \mathcal{O}_{\text{Proj}(A_n)}(1)$。
该引理还保证映射
$$
A_1 \leftarrow A_2 \leftarrow A_3 \leftarrow \ldots
$$
当 $n \leq m$ 时诱导同构 $A_n = A_m \otimes_{R_m} R_n$。令
$$
B = \bigoplus\nolimits_{d \geq 0} B_d
\quad\text{其中}\quad
B_d = \lim_n (A_n)_d
$$
由《代数进阶》引理
\ref{more-algebra-lemma-lim-finite-projective-gives-finite-projective}
可知 $B_d$ 是有限射影 $R$-模，且
$B \otimes_R R_n = A_n$。因此概形
$$
X = \text{Proj}(B)
\quad\text{且}\quad
\mathcal{L} = \mathcal{O}_X(1)
$$
在 $S$ 上平坦，且 $\mathcal{L}$ 是在 $S$ 上平坦的拟凝聚
$\mathcal{O}_X$-模；见《除子》引理
\ref{divisors-lemma-relative-proj-flat}。
由于 Proj 的形成与基变换可交换
（《构造》引理 \ref{constructions-lemma-base-change-map-proj}），
得到典范同构
$$
X \times_S S_n = X_n
\quad\text{且}\quad
\mathcal{L}|_{X_n} \cong \mathcal{L}_n
$$
它们与系统的转移映射相容。因此可把 $X_1 \subset X$
看作闭子概形。下文将证明 $B$ 在 $R$ 上有限表示。
由《除子》引理 \ref{divisors-lemma-relative-proj-proper} 和
\ref{divisors-lemma-relative-proj-finite-presentation}
可知 $X \to S$ 有限表示且为真，并且
$\mathcal{L} = \mathcal{O}_X(1)$ 作为 $\mathcal{O}_X$-模有限表示。
由于 $\mathcal{L}$ 在基变换 $X_1 \to S_1$ 上的限制可逆，
由《态射进阶》引理 \ref{more-morphisms-lemma-finite-free-open}，
$\mathcal{L}$ 在 $X$ 中 $X_1$ 的某个开邻域上可逆。
由于 $X \to S$ 为闭的，且 $\Ker(R \to R_1)$ 包含于 Jacobson 根中
（《代数进阶》引理 \ref{more-algebra-lemma-limit-henselian}），
可知 $X$ 中 $X_1$ 的任何开邻域都等于 $X$。
因此 $\mathcal{L}$ 可逆。最后，$S$ 中使得 $\mathcal{L}$
在纤维上充足的点集是 $S$ 的开集
（《态射进阶》引理 \ref{more-morphisms-lemma-ample-in-neighbourhood}），
且包含 $S_1$，故等于 $S$。于是 $X \to S$ 和 $\mathcal{L}$
具备引理陈述所要求的全部性质。

\medskip\noindent
现在证明上述断言。选取表示
$A_1 = R_1[X_1, \ldots, X_s]/(F_1, \ldots, F_t)$，
其中 $X_i$ 是次数为 $d_i$ 的变量，$F_j$ 是 $X_i$
中次数为 $e_j$ 的齐次多项式。可以选取映射
$$
\Psi : R[X_1, \ldots, X_s] \longrightarrow B
$$
提升映射 $R_1[X_1, \ldots, X_s] \to A_1$。
由于每个 $B_d$ 都是有限射影 $R$-模，再次利用
$\Ker(R \to R_1)$ 包含于 $R$ 的 Jacobson 根这一事实，
由 Nakayama 引理（《代数》引理 \ref{algebra-lemma-NAK}）可知
$\Psi$ 满射。由于 $- \otimes_R R_1$ 右正合，可以找到
$G_1, \ldots, G_t \in \Ker(\Psi)$，它们在
$R_1[X_1, \ldots, X_s]$ 中映到 $F_1, \ldots, F_t$。
注意，对所有 $d \geq 0$，$\Ker(\Psi)_d$ 是有限射影 $R$-模，
因为它是有限射影 $R$-模之间满射
$R[X_1, \ldots, X_s]_d \to B_d$ 的核。
再次由 Nakayama 引理可知，$\Ker(\Psi)$ 由
$G_1, \ldots, G_t$ 生成。
\end{proof}

\begin{lemma}
\label{quot-lemma-polarized-existence}
考虑基概形 $\Spec(\mathbf{Z})$ 上的叠
$\Polarizedstack$。则每个形式对象都是有效的。
\end{lemma}

\begin{proof}
引理中各概念的定义见《Artin 公理》第
\ref{artin-section-formal-objects} 节。
由定义可见，本引理立即推出自更一般的引理
\ref{quot-lemma-polarized-strong-effectiveness}。
\end{proof}

\begin{lemma}
\label{quot-lemma-polarized-defo-thy}
群胚叠 $\Polarizedstack$ 在 $\Spec(\mathbf{Z})$ 上
满足通用性的开放性。类似地，经基变换后
（注记 \ref{quot-remark-polarized-base-change}），
通用性的开放性在任意 Noether 基概形 $S$ 上成立。
\end{lemma}

\begin{proof}
这推出自《Artin 公理》引理
\ref{artin-lemma-SGE-implies-openness-versality} 以及引理
\ref{quot-lemma-polarized-diagonal}、
\ref{quot-lemma-polarized-RS-star},
\ref{quot-lemma-polarized-limits} 和
\ref{quot-lemma-polarized-strong-effectiveness}.
关于此事实的“通常”证明，见本证明之后注记中的讨论。
\end{proof}

\begin{remark}
\label{quot-remark-polarized-defo-thy}
引理 \ref{quot-lemma-polarized-defo-thy} 也可用《Artin 公理》引理
\ref{artin-lemma-get-openness-obstruction-theory} 中的障碍理论证明
（如引理 \ref{quot-lemma-coherent-defo-thy} 的第二个证明）。
为此，必须把《余切复形》第
\ref{cotangent-section-deformations-ringed-topoi} 节发展出的形变与障碍理论
推广到代数空间与拟凝聚模所成配对的情形。
另一种办法是利用 $1$-态射
$\Polarizedstack \to \Spacesstack'_{fp, flat, proper}$
为代数的（引理 \ref{quot-lemma-polarized-to-spaces-algebraic}），
以及我们已知其目标满足通用性的开放性
（引理 \ref{quot-lemma-spaces-defo-thy} 和
注记 \ref{quot-remark-spaces-defo-thy}）。
\end{remark}

\begin{theorem}[极化概形叠的代数性]
\label{quot-theorem-polarized-algebraic}
叠 $\Polarizedstack$（情形 \ref{quot-situation-polarized}）是代数的。
事实上，对任意代数空间 $B$，叠
$B\textit{-Polarized}$（注记 \ref{quot-remark-polarized-base-change}）
都是代数的。
\end{theorem}

\begin{proof}
绝对情形推出自《Artin 公理》引理
\ref{artin-lemma-diagonal-representable} 以及引理
\ref{quot-lemma-polarized-diagonal}、
\ref{quot-lemma-polarized-RS-star},
\ref{quot-lemma-polarized-limits},
\ref{quot-lemma-polarized-existence} 和
\ref{quot-lemma-polarized-defo-thy}.
关于 $B$ 的情形由此推出：$B\textit{-Polarized}$
在注记 \ref{quot-remark-polarized-base-change} 中被描述为
$2$-纤维积，而代数叠存在 $2$-纤维积；见
《代数叠》引理 \ref{algebraic-lemma-2-fibre-product}。
\end{proof}











\section{曲线叠}
\label{quot-section-curves}

\noindent
本节证明曲线叠是代数的。关于曲线模空间的进一步讨论，参见《曲线模空间》第
\ref{moduli-curves-section-introduction} 节。

\medskip\noindent
Stacks 项目中的曲线是维数为 $1$ 的簇。不过，当我们谈论曲线族时，通常允许
纤维可约和/或非约。在本节中，
曲线叠将“参数化维数 $\leq 1$ 的真概形”。然而，事实表明，为了得到曲线族的正确概念，
我们需要允许族的总空间是代数空间。这导出如下定义。

\begin{situation}
\label{quot-situation-curves}
我们定义范畴 $\Curvesstack$ 如下：
\begin{enumerate}
\item 对象是 {\it 曲线族}。更准确地说，对象是态射 $f : X \to S$，其中基底 $S$ 是概形，
{\it 总空间} $X$ 是代数空间，且 $f$ 平坦、真、有限表示，
相对维数 $\leq 1$（《空间的态射》定义
\ref{spaces-morphisms-definition-relative-dimension}）。
\item 对象之间的态射 $(X' \to S') \to (X \to S)$
由一对 $(f, g)$ 给出，其中 $f : X' \to X$ 是代数空间的态射，
$g : S' \to S$ 是概形的态射，并且它们适合交换图
$$
\xymatrix{
X' \ar[d] \ar[r]_f & X \ar[d] \\
S' \ar[r]^g & S
}
$$
诱导同构 $X' \to S' \times_S X$；换言之，该图是笛卡尔的。
\end{enumerate}
遗忘函子
$$
p : \Curvesstack \longrightarrow \Sch_{fppf},\quad
(X \to S) \longmapsto S
$$
就是我们把 $\Curvesstack$ 看作 $\Sch_{fppf}$ 上范畴的方式
（记号见第 \ref{quot-section-conventions} 节）。
\end{situation}

\noindent
由《域上的空间》引理
\ref{spaces-over-fields-lemma-codim-1-point-in-schematic-locus}
以及更一般地由《空间的更多态射》引理
\ref{spaces-more-morphisms-lemma-projective-over-complete}
可知，若 $S$ 是域的谱、阿廷局部环的谱，或诺特完备局部环的谱，
则对任意曲线族 $X \to S$，总空间 $X$ 是概形。
另一方面，在 $\mathbf{A}^1_k$ 上存在总空间不是概形的曲线族，见
《例子》第 \ref{examples-section-family-of-curves} 节。

\medskip\noindent
显然
\begin{equation}
\label{quot-equation-curves-over-proper-spaces}
\Curvesstack \subset \Spacesstack'_{fp, flat, proper}
\end{equation}
且 $\Spacesstack'_{fp, flat, proper}$ 的对象 $X \to S$
属于 $\Curvesstack$ 当且仅当 $X \to S$ 的相对维数 $\leq 1$。
我们将利用这一点验证 $\Curvesstack$ 的 Artin 公理。

\begin{lemma}
\label{quot-lemma-curves-fibred-in-groupoids}
范畴 $\Curvesstack$ 在 $\Sch_{fppf}$ 上是群胚纤维化的。
\end{lemma}

\begin{proof}
利用嵌入（\ref{quot-equation-curves-over-proper-spaces}）、像的描述以及
$\Spacesstack'_{fp, flat, proper}$ 的相应事实
（引理 \ref{quot-lemma-spaces-fibred-in-groupoids}），
问题归结为如下陈述：给定态射
$$
\xymatrix{
X' \ar[r] \ar[d] & X \ar[d] \\
S' \ar[r] & S
}
$$
在 $\Spacesstack'_{fp, flat, proper}$ 中（回忆这尤其意味着该图是笛卡尔的），
若 $X \to S$ 的相对维数 $\leq 1$，则 $X' \to S'$ 的相对维数
也 $\leq 1$。
这由《空间的态射》引理
\ref{spaces-morphisms-lemma-dimension-fibre-after-base-change}.
\end{proof}

\begin{lemma}
\label{quot-lemma-curves-stack}
范畴 $\Curvesstack$ 是 $\Sch_{fppf}$ 上的群胚叠。
\end{lemma}

\begin{proof}
利用嵌入（\ref{quot-equation-curves-over-proper-spaces}）、像的描述以及
$\Spacesstack'_{fp, flat, proper}$ 的相应事实
（引理 \ref{quot-lemma-spaces-stack}），
问题归结为如下陈述：设 $\Spacesstack'_{fp, flat, proper}$ 的对象
$X \to S$ 以及 fppf 覆盖 $\{S_i \to S\}_{i \in I}$，
则下列条件等价：
\begin{enumerate}
\item $X \to S$ 的相对维数 $\leq 1$；
\item 对每个 $i$，基变换 $X_i \to S_i$ 的相对维数
$\leq 1$。
\end{enumerate}
这由《空间的态射》引理
\ref{spaces-morphisms-lemma-dimension-fibre-after-base-change}.
\end{proof}

\begin{lemma}
\label{quot-lemma-curves-diagonal}
对角态射
$$
\Delta : \Curvesstack \longrightarrow \Curvesstack \times \Curvesstack
$$
由代数空间表示。
\end{lemma}

\begin{proof}
这立即由全忠实嵌入
（\ref{quot-equation-curves-over-proper-spaces}）以及
$\Spacesstack'_{fp, flat, proper}$ 的相应事实
（引理 \ref{quot-lemma-spaces-diagonal}）得到。
\end{proof}

\begin{remark}
\label{quot-remark-curves-base-change}
设 $B$ 是 $\Spec(\mathbf{Z})$ 上的代数空间。
令 $B\text{-}\Curvesstack$ 为由对 $(X \to S, h : S \to B)$
组成的范畴，其中 $X \to S$ 是 $\Curvesstack$ 的对象，
$h : S \to B$ 是态射。
在 $B\text{-}\Curvesstack$ 中，态射 $(X' \to S', h') \to (X \to S, h)$
是 $\Curvesstack$ 中满足 $h \circ g = h'$ 的态射 $(f, g)$。
在此情形，图
$$
\xymatrix{
B\text{-}\Curvesstack \ar[r] \ar[d] & \Curvesstack \ar[d] \\
(\Sch/B)_{fppf} \ar[r] & \Sch_{fppf}
}
$$
是 $2$-纤维积方形。这个平凡的备注有时很有用：
它可以把绝对情形 $\Curvesstack$ 的结果推导到
给定基代数空间上的曲线族情形。
\end{remark}

\begin{lemma}
\label{quot-lemma-curves-limits}
叠 $\Curvesstack \to \Sch_{fppf}$ 保持极限
（《Artin 公理》定义 \ref{artin-definition-limit-preserving}）。
\end{lemma}

\begin{proof}
利用嵌入（\ref{quot-equation-curves-over-proper-spaces}）、像的描述以及
$\Spacesstack'_{fp, flat, proper}$ 的相应事实
（引理 \ref{quot-lemma-spaces-limits}），
问题归结为如下陈述：
设 $T = \lim T_i$ 是仿射概形有向逆系统的极限。
设 $i \in I$，且 $X_i \to T_i$ 是
$T_i$ 上 $\Spacesstack'_{fp, flat, proper}$ 的对象。
假设 $T \times_{T_i} X_i \to T$ 的相对维数 $\leq 1$。
则存在某个 $i' \geq i$，使得态射
$T_{i'} \times_{T_i} X_i \to T_i$ 的相对维数 $\leq 1$。
这由《空间的极限》引理
\ref{spaces-limits-lemma-eventually-relative-dimension}.
\end{proof}

\begin{lemma}
\label{quot-lemma-curves-RS-star}
设
$$
\xymatrix{
T \ar[r] \ar[d] & T' \ar[d] \\
S \ar[r] & S'
}
$$
是概形范畴中的推出，其中 $T \to T'$ 是加厚且 $T \to S$ 是仿射；见
《更多态射》引理 \ref{more-morphisms-lemma-pushout-along-thickening}。
则纤维范畴上的函子
$$
\Curvesstack_{S'}
\longrightarrow
\Curvesstack_S
\times_{\Curvesstack_T}
\Curvesstack_{T'}
$$
是等价。
\end{lemma}

\begin{proof}
利用嵌入（\ref{quot-equation-curves-over-proper-spaces}）、像的描述以及
$\Spacesstack'_{fp, flat, proper}$ 的相应事实
（引理 \ref{quot-lemma-spaces-RS-star}），
问题归结为如下陈述：
设 $X' \to S'$ 是代数空间到 $S'$ 的态射，有限表示、平坦且真，
则 $X' \to S'$ 的相对维数 $\leq 1$
当且仅当 $S \times_{S'} X' \to S$
和 $T' \times_{S'} X' \to T'$ 的相对维数 $\leq 1$。
一个方向由相对维数 $\leq 1$ 在基变换下保持这一事实得到
（《空间的态射》引理
\ref{spaces-morphisms-lemma-dimension-fibre-after-base-change}).
另一个方向由以下事实得到：相对维数 $\leq 1$ 可在纤维上检验，
并且 $X' \to S'$ 的纤维（在概形 $S'$ 的点上）
与 $S \times_{S'} X' \to S$ 的纤维相同，
因为根据《更多态射》引理
\ref{more-morphisms-lemma-pushout-along-thickening}，
$S \to S'$ 是加厚。
\end{proof}

\begin{lemma}
\label{quot-lemma-curves-tangent-space}
设 $k$ 是域，且 $x = (X \to \Spec(k))$ 是
$\Spec(k)$ 上 $\mathcal{X} = \Curvesstack$ 的对象。
\begin{enumerate}
\item 若 $k$ 在 $\mathbf{Z}$ 上有限型，则向量空间
$T\mathcal{F}_{\mathcal{X}, k, x}$ 和
$\text{Inf}(\mathcal{F}_{\mathcal{X}, k, x})$
（见《Artin 公理》第 \ref{artin-section-tangent-spaces} 节）
是有限维的；并且
\item 一般地，向量空间 $T_x(k)$ 和 $\text{Inf}_x(k)$
（见《Artin 公理》第 \ref{artin-section-inf} 节）
是有限维的。
\end{enumerate}
\end{lemma}

\begin{proof}
这立即由全忠实嵌入
（\ref{quot-equation-curves-over-proper-spaces}）以及
$\Spacesstack'_{fp, flat, proper}$ 的相应事实
（引理 \ref{quot-lemma-spaces-tangent-space}）得到。
\end{proof}

\begin{lemma}
\label{quot-lemma-curves-existence}
考虑基概形 $\Spec(\mathbf{Z})$ 上的叠 $\Curvesstack$。
则每个形式对象都是有效的。
\end{lemma}

\begin{proof}
引理中各概念的定义见《Artin 公理》第
\ref{artin-section-formal-objects} 节。
设 $(A, \mathfrak m, \kappa)$ 是诺特完备局部环。
设 $(X_n \to \Spec(A/\mathfrak m^n))$
是 $A$ 上 $\Curvesstack$ 的形式对象。
由《空间的更多态射》引理
\ref{spaces-more-morphisms-lemma-formal-algebraic-space-proper-reldim-1}
存在射影态射 $X \to \Spec(A)$
以及同构的相容系统
$X \times_{\Spec(A)} \Spec(A/\mathfrak m^n) \cong X_n$。
由《更多态射》引理
\ref{more-morphisms-lemma-check-flatness-on-infinitesimal-nbhds}
可知 $X \to \Spec(A)$ 是平坦的。由《更多态射》引理
\ref{more-morphisms-lemma-dimension-fibres-proper-flat}
可知 $X \to \Spec(A)$ 的相对维数 $\leq 1$。
引理得证。
\end{proof}

\begin{lemma}
\label{quot-lemma-curves-defo-thy}
群胚叠 $\mathcal{X} = \Curvesstack$
在 $\Spec(\mathbf{Z})$ 上满足形式光滑性的开放性。
同样，在基变换之后（备注 \ref{quot-remark-curves-base-change}），
形式光滑性的开放性对任意诺特基概形 $S$ 成立。
\end{lemma}

\begin{proof}
这立即由全忠实嵌入
（\ref{quot-equation-curves-over-proper-spaces}）以及
$\Spacesstack'_{fp, flat, proper}$ 的相应事实
（引理 \ref{quot-lemma-spaces-defo-thy}）得到。
\end{proof}

\begin{theorem}[曲线叠的代数性]
\label{quot-theorem-curves-algebraic}
\begin{reference}
参见 \cite[命题 3.3，第 8 页]{dJHS} 以及
\cite[Jack Hall 的附录 B，定理 B.1]{Smyth}。
\end{reference}
叠 $\Curvesstack$（情形 \ref{quot-situation-curves}）
是代数的。事实上，对任意代数空间 $B$，叠
$B\text{-}\Curvesstack$（备注 \ref{quot-remark-curves-base-change}）
也是代数的。
\end{theorem}

\begin{proof}
绝对情形由《Artin 公理》引理
\ref{artin-lemma-diagonal-representable}
以及引理 \ref{quot-lemma-curves-diagonal}、
\ref{quot-lemma-curves-RS-star},
\ref{quot-lemma-curves-limits},
\ref{quot-lemma-curves-existence}，以及
\ref{quot-lemma-curves-defo-thy}.
$B$ 上的情形由此、备注 \ref{quot-remark-curves-base-change} 中
$B\text{-}\Curvesstack$ 作为 $2$-纤维积的描述，以及
代数叠具有 $2$-纤维积这一事实得到；见《代数叠》引理
\ref{algebraic-lemma-2-fibre-product}。
\end{proof}

\begin{lemma}
\label{quot-lemma-curves-open-and-closed-in-spaces}
这个 $1$-态射（\ref{quot-equation-curves-over-proper-spaces}）
$$
\Curvesstack \longrightarrow \Spacesstack'_{fp, flat, proper}
$$
由开浸入和闭浸入表示。
\end{lemma}

\begin{proof}
由于（\ref{quot-equation-curves-over-proper-spaces}）是范畴的全忠实嵌入，
只需证明如下事实：给定 $\Spacesstack'_{fp, flat, proper}$ 的对象
$X \to S$，存在开闭子概形 $U \subset S$，
使得态射 $S' \to S$ 经由 $U$ 分解，当且仅当
$X \to S$ 的基变换 $X' \to S'$ 的相对维数 $\leq 1$。
这立即由《空间的更多态射》引理
\ref{spaces-more-morphisms-lemma-dimension-fibres-proper-flat}.
\end{proof}

\begin{remark}
\label{quot-remark-alternative-approach-curves}
考虑 $2$-纤维积
$$
\xymatrix{
\Curvesstack \times_{\Spacesstack'_{fp, flat, proper}}
\Polarizedstack \ar[r] \ar[d] &
\Polarizedstack \ar[d] \\
\Curvesstack \ar[r] &
\Spacesstack'_{fp, flat, proper}
}
$$
这个纤维积参数化极化曲线，即带有相对充足可逆层的曲线族。
事实表明，左侧竖直箭头
$$
\textit{PolarizedCurves} \longrightarrow \Curvesstack
$$
是代数的、光滑且满射的。也就是说，这个 $1$-态射
是代数的（因为它是引理 \ref{quot-lemma-polarized-to-spaces-algebraic} 中箭头的基变换），
每个点都在像中，并且曲线上的可逆层不存在变形障碍
（见引理 \ref{quot-lemma-curves-existence} 的证明）。
这给出了证明 $\Curvesstack$ 代数性的另一种方法。
具体地，由引理 \ref{quot-lemma-curves-open-and-closed-in-spaces}，
可知 $\textit{PolarizedCurves}$ 是代数叠 $\Polarizedstack$ 的开闭子叠，
而任何从代数叠出发的光滑代数态射的目标群胚叠都是代数叠。
\end{remark}









\section{真态射上的复形模空间}
\label{quot-section-moduli-complexes}

\noindent
本节的标题和内容取自
\cite{lieblich-complexes}。设 $S$ 是概形，且
$f : X \to B$ 是代数空间的真、平坦、有限表示态射。
我们将证明存在一个代数叠
$$
\Complexesstack_{X/B}
$$
参数化纤维上负自扩张为零的 $D^b_{\textit{Coh}}$ 对象“族”。
更准确地说，族由总空间导出范畴中的相对完美对象给出；
这一略显技术性的概念见《空间的更多态射》第
\ref{spaces-more-morphisms-section-relatively-perfect}.

\medskip\noindent
即使 $X$ 是域 $k$ 上的真代数空间，我们也得到一个非常有趣的代数叠。
具体地，有嵌入
$$
\Cohstack_{X/k} \longrightarrow \Complexesstack_{X/k}
$$
因为对任意 $\mathcal{O}$-模 $\mathcal{F}$（在任意带环拓扑斯上），
都有 $\Ext^i_\mathcal{O}(\mathcal{F}, \mathcal{F}) = 0$（当 $i < 0$）。
虽然这当然说明我们的叠非空，但研究 $\Complexesstack_{X/k}$
的真正动机是，导出范畴中常常存在这样的对象：
$D^b_{\textit{Coh}}(\mathcal{O}_X)$ 中负自扩张为零的对象，
并且其上同调层在多个次数中非零。
例如，$X$ 可能与另一个真代数空间 $Y$（定义在 $k$ 上）导出等价，
即存在一个 $k$-线性等价
$$
F : D^b_{\textit{Coh}}(\mathcal{O}_Y)
\longrightarrow
D^b_{\textit{Coh}}(\mathcal{O}_X)
$$
有时确实如此，而 $F$ 并非由 $X$ 与 $Y$ 之间的同构给出；
例如阿贝尔簇及其对偶的情形。在此情形，$F$ 诱导代数叠之间的同构
$$
\Complexesstack_{Y/k}
\longrightarrow
\Complexesstack_{X/k}
$$
（此处待补未来的参考文献），特别地，$Y$ 上的凝聚层叠映入 $X$ 上的复形叠。
反过来，如果我们足够理解 $\Complexesstack_{X/k}$ 的几何，
就可以尝试利用它研究所有可能的导出等价空间 $Y$。

\begin{lemma}
\label{quot-lemma-complexes-open-neg-exts-vanishing}
设 $S$ 是概形。
设 $f : X \to Y$ 是 $S$ 上代数空间的态射。
假设 $f$ 真、平坦且有限表示。
设 $K, E \in D(\mathcal{O}_X)$。假设 $K$ 伪凝聚，
且 $E$ 是 $Y$-完美的（《空间的更多态射》定义
\ref{spaces-more-morphisms-definition-relatively-perfect}）。
对域 $k$ 及态射 $y : \Spec(k) \to Y$，记 $K_y$、$E_y$
为到纤维 $X_y$ 的拉回。
\begin{enumerate}
\item 存在开子集 $W \subset Y$，其特征性质为
$$
y \in |W|
\Leftrightarrow
\Ext^i_{\mathcal{O}_{X_y}}(K_y, E_y) = 0
\text{，其中 }i < 0.
$$
\item 对于任意态射 $V \to Y$，若其经由 $W$ 分解，则有
$$
\Ext^i_{\mathcal{O}_{X_V}}(K_V, E_V) = 0
\quad\text{当}\quad i < 0
$$
其中 $X_V$ 是 $X$ 的基变换，而 $K_V$、$E_V$
是 $K$、$E$ 到 $X_V$ 的导出拉回。
\item 函子 $V \mapsto \Hom_{\mathcal{O}_{X_V}}(K_V, E_V)$
是 $(\textit{Spaces}/W)_{fppf}$ 上的层，由一个在 $W$ 上仿射且有限表示的
代数空间表示。
\end{enumerate}
\end{lemma}

\begin{proof}
对任意态射 $V \to Y$，复形 $K_V$ 伪凝聚
（《站点上的上同调》引理
\ref{sites-cohomology-lemma-pseudo-coherent-pullback})
且 $E_V$ 是 $V$-完美的（《空间的更多态射》引理
\ref{spaces-more-morphisms-lemma-base-change-relatively-perfect}).
另一个观察是，给定 $y : \Spec(k) \to Y$
以及域扩张 $k'/k$，带有 $y' : \Spec(k') \to Y$，
所诱导的态射满足
$$
\Ext^i_{\mathcal{O}_{X_{y'}}}(K_{y'}, E_{y'}) =
\Ext^i_{\mathcal{O}_{X_y}}(K_y, E_y) \otimes_k k'
$$
由《概形的导出范畴》引理
\ref{perfect-lemma-affine-morphism-and-hom-out-of-perfect}.
因此，(1) 中的消失确实是所诱导点 $y \in |Y|$ 的性质。
证明中不再特别提及这两个观察。

\medskip\noindent
首先假设 $Y$ 是仿射概形。于是可应用
《空间的更多态射》引理
\ref{spaces-more-morphisms-lemma-compute-ext-rel-perfect}
并找到伪凝聚的 $L \in D(\mathcal{O}_Y)$，它按该引理所述的意义计算
“普遍地计算” $Rf_*R\SheafHom(K, E)$。
展开定义，对点 $y \in Y$ 有等式
$$
\Ext^i_{\kappa(y)}(L \otimes_{\mathcal{O}_Y}^\mathbf{L} \kappa(y),
\kappa(y)) = \Ext^i_{\mathcal{O}_{X_y}}(K_y, E_y)
$$
由此得到
$$
H^i(L \otimes_{\mathcal{O}_Y}^\mathbf{L} \kappa(y)) = 0
\text{ 对 } i > 0 \Leftrightarrow
\Ext^i_{\mathcal{O}_{X_y}}(K_y, E_y) = 0 \text{ 对 }i < 0.
$$
由《概形的导出范畴》引理 \ref{perfect-lemma-jump-loci}，
满足该条件的点集 $W$ 由 $y \in Y$ 构成，是 $Y$ 的开子集。
这个开集 $W$ 对所有来自域谱的态射，依据 $L$ 的“普遍性”，
满足 (1) 中的要求。

\medskip\noindent
回到一般的代数空间 $Y$。
取 $\{V_i \to Y\}$ 为由仿射概形 $V_i$ 给出的 \'etale 覆盖。
则子集 $W \subset |Y|$ 拉回为相应的子集
$W_i \subset |V_i|$（对应于 $X_{V_i}$、$K_{V_i}$、$E_{V_i}$）。
由上一段可知 $W_i$ 是开集，因而 $W$ 是开集。
这就一般地证明了 (1)。此外，(2)、(3) 完全用范畴
$\textit{Spaces}/W$ 及限制 $X_W$、$K_W$、$E_W$ 表述。
因此只需考虑 $W = Y$ 的情形。

\medskip\noindent
假设 $W = Y$。我们断言，对任意代数空间 $V$（在 $Y$ 上），
$Rf_{V, *}R\SheafHom(K_V, E_V)$ 在次数 $< 0$ 的上同调层均消失。
这将证明 (2)，因为
$$
\Ext^i_{\mathcal{O}_{X_V}}(K_V, E_V) =
H^i(X_V, R\SheafHom(K_V, E_V)) = 
H^i(V, Rf_{V, *}R\SheafHom(K_V, E_V))
$$
由《站点上的上同调》引理
\ref{sites-cohomology-lemma-section-RHom-over-U} 以及
\ref{sites-cohomology-lemma-Leray-unbounded}
并且由上同调层消失及《导出范畴》引理
\ref{derived-lemma-negative-vanishing}，上同调群 $H^i$ 在 $i < 0$ 时为零。

\medskip\noindent
为证明该断言，我们可以在 $V$ 上作 \'etale 局部工作。
特别地，可假设 $Y$ 仿射且 $W = Y$。
令 $L \in D(\mathcal{O}_Y)$ 如证明第二段所述。
对代数空间 $V$（在 $Y$ 上），记 $L_V$ 为 $L$ 到 $V$ 的导出拉回。
（我们将使用的一个重要特征是，$L$ 对所有代数空间 $V$（在 $Y$ 上）都适用，
而不只对仿射 $V$ 适用。）
由于 $W = Y$，有 $H^i(L) = 0$（当 $i > 0$）
（使用《更多代数》引理
\ref{more-algebra-lemma-lift-pseudo-coherent-from-residue-field}
从纤维转到茎）。因此 $H^i(L_V) = 0$（当 $i > 0$）。
定义 $L$ 的性质是
$$
Rf_{V, *}R\SheafHom(K_V, E_V) = R\SheafHom(L_V, \mathcal{O}_V)
$$
由于 $L_V$ 集中在次数 $\leq 0$，可知
$R\SheafHom(L_V, \mathcal{O}_V)$ 集中在次数 $\geq 0$，
从而证明断言。(2) 的证明完成。

\medskip\noindent
假设 $W = Y$，但对代数空间 $Y$ 不作任何假设。
由 (2) 以及《单纯空间》引理
\ref{spaces-simplicial-lemma-fppf-neg-ext-zero-hom}，
函子 $F$（由 $F(V) = \Hom_{\mathcal{O}_{X_V}}(K_V, E_V)$ 给出）
是一个层\footnote{要检验覆盖 $\{V_i \to V\}_{i \in I}$ 的层性质，先考虑
{\v C}ech fppf 超覆盖 $a : V_\bullet \to V$，其中
$V_n = \coprod_{i_0 \ldots i_n} V_{i_0} \times_V \ldots \times_V V_{i_n}$
然后令 $U_\bullet = V_\bullet \times_{a, V} X_V$。于是
$U_\bullet \to X_V$ 是一个 fppf 超覆盖，可应用《单纯空间》引理
\ref{spaces-simplicial-lemma-fppf-neg-ext-zero-hom}.}
在 $(\textit{Spaces}/Y)_{fppf}$ 上。要证明 $F$ 是代数空间且
$F \to Y$ 在 $Y$ 上仿射并有限表示，我们可以作 \'etale 局部工作；
见《引导》引理 \ref{bootstrap-lemma-locally-algebraic-space-finite-type}
以及《空间的态射》引理
\ref{spaces-morphisms-lemma-affine-local} 和
\ref{spaces-morphisms-lemma-finite-presentation-local}。
因此，只需证明 $F$ 在 $Y$ 为仿射概形时是 $Y$ 上仿射且有限表示的代数空间。
此时回到伪凝聚复形 $L \in D(\mathcal{O}_Y)$。
由于 $H^i(L) = 0$（当 $i > 0$），可将 $L$ 表示为如下形式的复形：
$$
\ldots \to \mathcal{O}_Y^{\oplus m_1} \to \mathcal{O}_Y^{\oplus m_0}
\to 0 \to \ldots
$$
末项位于次数 $0$；见《更多代数》引理
\ref{more-algebra-lemma-pseudo-coherent}。结合证明中前面两个显示公式可得
$$
F(V) =
\Ker(
\Hom_V(\mathcal{O}_V^{\oplus m_0}, \mathcal{O}_V)
\to 
\Hom_V(\mathcal{O}_V^{\oplus m_1}, \mathcal{O}_V)
)
$$
换言之，存在纤维积图
$$
\xymatrix{
F \ar[d] \ar[r] & Y \ar[d]^0 \\
\mathbf{A}_Y^{m_0} \ar[r] & \mathbf{A}_Y^{m_1}
}
$$
这就证明了所需结论。
\end{proof}

\begin{lemma}
\label{quot-lemma-complexes}
设 $S$ 是概形。设 $f : X \to Y$ 是 $S$ 上代数空间的态射。
假设 $f$ 真、平坦且有限表示。
设 $E \in D(\mathcal{O}_X)$。
假设
\begin{enumerate}
\item $E$ 是 $S$-完美的（《空间的更多态射》定义
\ref{spaces-more-morphisms-definition-relatively-perfect})，以及
\item 对每个点 $s \in S$，有
$$
\Ext^i_{\mathcal{O}_{X_s}}(E_s, E_s) = 0
\quad\text{当}\quad i < 0
$$
其中 $E_s$ 是到纤维 $X_s$ 的拉回。
\end{enumerate}
则
\begin{enumerate}
\item[(a)] (1)、(2) 在任意基变换 $V \to Y$ 下保持；
\item[(b)] 有
$\Ext^i_{\mathcal{O}_{X_V}}(E_V, E_V) = 0$（当 $i < 0$），
且对所有 $V$（在 $Y$ 上）均成立；
\item[(c)] $V \mapsto \Hom_{\mathcal{O}_{X_V}}(E_V, E_V)$
由一个在 $Y$ 上仿射且有限表示的代数空间表示。
\end{enumerate}
这里 $X_V$ 是 $X$ 的基变换，$E_V$ 是 $E$ 到 $X_V$ 的导出拉回。
\end{lemma}

\begin{proof}
这是引理 \ref{quot-lemma-complexes-open-neg-exts-vanishing} 的直接推论。
\end{proof}

\begin{situation}
\label{quot-situation-complexes}
设 $S$ 是概形。设 $f : X \to B$ 是 $S$ 上代数空间的态射。
假设 $f$ 真、平坦且有限表示。
记 $\Complexesstack_{X/B}$ 为如下范畴：其对象是三元组 $(T, g, E)$，其中
\begin{enumerate}
\item $T$ 是 $S$ 上的概形；
\item $g : T \to B$ 是 $S$ 上的态射，并令
$X_T = T \times_{g, B} X$
\item $E$ 是 $D(\mathcal{O}_{X_T})$ 的对象，满足引理
\ref{quot-lemma-complexes} 的条件 (1)、(2)。
\end{enumerate}
态射 $(T, g, E) \to (T', g', E')$
由一对 $(h, \varphi)$ 给出，其中
\begin{enumerate}
\item $h : T \to T'$ 是 $B$ 上概形的态射
（即 $g' \circ h = g$）；并且
\item $\varphi : L(h')^*E' \to E$ 是
$D(\mathcal{O}_{X_T})$ 中的同构，其中 $h' : X_T \to X_{T'}$
是 $h$ 的基变换。
\end{enumerate}
\end{situation}

\noindent
因此 $\Complexesstack_{X/B}$ 是一个范畴，且规则
$$
p : \Complexesstack_{X/B} \longrightarrow (\Sch/S)_{fppf},
\quad
(T, g, E) \longmapsto T
$$
是函子。对概形 $T$（在 $S$ 上），记 $\Complexesstack_{X/B, T}$
为 $p$ 在 $T$ 上的纤维范畴。这些纤维范畴是群胚。

\begin{lemma}
\label{quot-lemma-complexes-fibred-in-groupoids}
在情形 \ref{quot-situation-complexes} 中，函子
$p : \Complexesstack_{X/B} \longrightarrow (\Sch/S)_{fppf}$
是群胚纤维化的。
\end{lemma}

\begin{proof}
我们通过检验函子 $p$ 是否为群胚纤维化，使用《范畴》定义
\ref{categories-definition-fibred-groupoids}.
给定对象 $(T', g', E')$（属于 $\Complexesstack_{X/B}$）
以及态射 $h : T \to T'$（为 $S$ 上概形的态射），令 $g = h \circ g'$，
并令 $E = L(h')^*E'$，其中 $h' : X_T \to X_{T'}$
是 $h$ 的基变换。显然得到态射
$(T, g, E) \to (T', g', E')$，它属于 $\Complexesstack_{X/B}$
且位于 $h$ 之上。
这证明了 (1)。
对于 (2)，设给定态射
$$
(h_1, \varphi_1) : (T_1, g_1, E_1) \to (T, g, E)
\quad\text{以及}\quad
(h_2, \varphi_2) : (T_2, g_2, E_2) \to (T, g, E)
$$
它们属于 $\Complexesstack_{X/B}$，且有态射 $h : T_1 \to T_2$ 满足
$h_2 \circ h = h_1$。令 $\varphi$ 为复合
$$
L(h')^*E_2
\xrightarrow{L(h')^*\varphi_2^{-1}}
L(h')^*L(h_2)^*E = L(h_1)^*E
\xrightarrow{\varphi_1}
E_1
$$
即可得到态射
$(h, \varphi) : (T_1, g_1, E_1) \to (T_2, g_2, E_2)$，
从而证明条件 (2) 成立。
\end{proof}

\begin{lemma}
\label{quot-lemma-complexes-diagonal}
在情形 \ref{quot-situation-complexes} 中，记
$\mathcal{X} = \Complexesstack_{X/B}$。则
$\Delta : \mathcal{X} \to \mathcal{X} \times \mathcal{X}$ 是
由代数空间表示。
\end{lemma}

\begin{proof}
考虑两个对象 $x = (T, g, E)$ 和 $y = (T, g', E')$，它们属于
$\mathcal{X}$ 在概形 $T$ 上的纤维。需证明
$\mathit{Isom}_\mathcal{X}(x, y)$ 是 $T$ 上的代数空间，见
《代数叠》引理 \ref{algebraic-lemma-representable-diagonal}。
若对 $h : T' \to T$，限制对象 $x|_{T'}$ 与 $y|_{T'}$
在纤维范畴 $\mathcal{X}_{T'}$ 中同构，则 $g \circ h = g' \circ h$。
因此存在预层的变换
$$
\mathit{Isom}_\mathcal{X}(x, y) \longrightarrow \text{Equalizer}(g, g')
$$
由于 $B$ 的对角态射可表示（由概形表示），这个等化子是概形。
因此可用此等化子替换 $T$，并用其拉回替换 $E$、$E'$。
于是可假设 $g = g'$。

\medskip\noindent
假设 $g = g'$。以 $B$ 替换为 $T$，并以 $X$ 替换为 $X_T$ 后，
问题变为：给定 $E, E' \in D(\mathcal{O}_X)$，
满足引理 \ref{quot-lemma-complexes} 的条件 (1)、(2)，
证明 $\mathit{Isom}(E, E')$ 是代数空间。
这里 $\mathit{Isom}(E, E')$ 是函子
$$
(\Sch/B)^{opp} \to \textit{Sets},\quad
T \mapsto \{\varphi : E_T \to E'_T
\text{ isomorphism in }D(\mathcal{O}_{X_T})\}
$$
其中 $E_T$、$E'_T$ 是 $E$、$E'$ 到 $X_T$ 的导出拉回。
现在，令 $W \subset B$（分别令 $W' \subset B$）为 $B$ 中由引理
\ref{quot-lemma-complexes-open-neg-exts-vanishing} 与 $E, E'$（分别与 $E', E$）
关联的开子空间。
显然，若存在同构 $E_T \to E'_T$（如 $\mathit{Isom}(E, E')$ 的定义所述），
则 $T \to B$ 同时经由 $W$、$W'$ 分解（因为 $E$、$E'$ 满足条件 (1)，
且显然 $E_t \cong E'_t$，所以不存在非零态射
$E_t[i] \to E_t$ 或 $E'_t[i] \to E_t$；在纤维 $X_t$ 上当 $i > 0$ 时均如此。
因此可用 $B$ 替换为开集 $W \cap W'$。
在此情形，函子 $H = \SheafHom(E, E')$
$$
(\Sch/B)^{opp} \to \textit{Sets},\quad T \mapsto
\Hom_{\mathcal{O}_{X_T}}(E_T, E'_T)
$$
由引理 \ref{quot-lemma-complexes-open-neg-exts-vanishing}，
是 $B$ 上仿射且有限表示的代数空间。
对下列函子同样如此：
$H' = \SheafHom(E', E)$,
$I = \SheafHom(E, E)$，以及
$I' = \SheafHom(E', E')$.
因此可重复命题 \ref{quot-proposition-isom} 证明中的论证，得到
$$
\mathit{Isom}(E, E') = (H' \times_B H) \times_{c, I \times_B I', \sigma} B
$$
其中 $c$、$\sigma$ 是某些态射。因此
$\mathit{Isom}(E, E')$ 是代数空间。
\end{proof}

\begin{lemma}
\label{quot-lemma-complexes-stack}
在情形 \ref{quot-situation-complexes} 中，函子
$p : \Complexesstack_{X/B} \longrightarrow (\Sch/S)_{fppf}$
是群胚叠。
\end{lemma}

\begin{proof}
要证明 $\Complexesstack_{X/B}$ 是群胚叠，
需证明预层 $\mathit{Isom}$ 是层且下降数据有效。
关于 $\mathit{Isom}$ 的陈述由引理 \ref{quot-lemma-complexes-diagonal} 得到，见
《代数叠》引理 \ref{algebraic-lemma-representable-diagonal}。
下面证明下降数据的陈述。

\medskip\noindent
设 $\{a_i : T_i \to T\}$ 是 $S$ 上概形的 fppf 覆盖。
设 $(\xi_i, \varphi_{ij})$ 是 $\{T_i \to T\}$ 的下降数据，
取值于 $\Complexesstack_{X/B}$。
对每个 $i$，可写 $\xi_i = (T_i, g_i, E_i)$。
记投影 $\text{pr}_0 : T_i \times_T T_j \to T_i$ 和
$\text{pr}_1 : T_i \times_T T_j \to T_j$。
条件 $\xi_i|_{T_i \times_T T_j} \cong \xi_j|_{T_i \times_T T_j}$
尤其蕴含 $g_i \circ \text{pr}_0 = g_j \circ \text{pr}_1$。
因此存在唯一态射 $g : T \to B$ 使 $g_i = g \circ a_i$，见
《空间上的下降》引理
\ref{spaces-descent-lemma-fpqc-universal-effective-epimorphisms}.
记 $X_T = T \times_{g, B} X$。令
$X_i = X_{T_i} = T_i \times_{g_i, B} X = T_i \times_{a_i, T} X_T$
以及
$$
X_{ij} = X_{T_i} \times_{X_T} X_{T_j} = X_i \times_{X_T} X_j
$$
其投影 $\text{pr}_i$、$\text{pr}_j$ 分别到 $X_i$、$X_j$。
注意 $(T_i, g_i, E_i)$ 沿
$\text{pr}_0 : T_i \times_T T_j \to T_i$ 的拉回为
$(T_i \times_T T_j, g_i \circ \text{pr}_0, L\text{pr}_i^*E_i)$.
因此，$\{T_i \to T\}$ 在 $\Complexesstack_{X/B}$ 中的下降数据
由对象 $(T_i, g \circ a_i, E_i)$ 以及每一对 $i, j$ 的同构给出，该同构位于
$D\mathcal{O}_{X_{ij}})$
$$
\varphi_{ij} :
L\text{pr}_i^*E_i \longrightarrow L\text{pr}_j^*E_j
$$
并在 $X$ 到
$T_i \times_T T_j \times_T T_k$.
利用引理 \ref{quot-lemma-complexes} 的 (b) 所给出的负 Ext 消失，可应用
《单纯空间》引理 \ref{spaces-simplicial-lemma-fppf-glue-neg-ext-zero}
得到这些复形的下降\footnote{为检验这一点，先考虑
{\v C}ech fppf 超覆盖 $a : T_\bullet \to T$，其中
$T_n = \coprod_{i_0 \ldots i_n} T_{i_0} \times_T \ldots \times_T T_{i_n}$
然后令 $U_\bullet = T_\bullet \times_{a, T} X_T$。于是
$U_\bullet \to X_T$ 是 fppf 超覆盖，可应用《单纯空间》引理
\ref{spaces-simplicial-lemma-fppf-glue-neg-ext-zero}.}
对于这些复形。换言之，存在对象 $E$，属于
$D_\QCoh(\mathcal{O}_{X_T})$，将 $E_i$ 限制到 $X_{T_i}$ 上，
并与 $\varphi_{ij}$ 相容。回忆，$T$-完美意味着伪凝聚，
且相对于 $f^{-1}\mathcal{O}_T$ 具有局部有限 Tor 维。
因此 $E$ 是 $T$-完美的，这是由《空间的更多态射》引理
\ref{spaces-more-morphisms-lemma-pseudo-coherent-descends-fpqc} 以及
\ref{spaces-more-morphisms-lemma-tor-amplitude-descends-fppf}.
最后需检验 $E$ 满足引理 \ref{quot-lemma-complexes} 的条件 (2)。
这立即由引理 \ref{quot-lemma-complexes-open-neg-exts-vanishing} 中
开集 $W$ 的描述，以及 (2) 对 $E_i$ 在 $X_{T_i}/T_i$ 上成立这一事实得到。
\end{proof}

\begin{remark}
\label{quot-remark-complexes-base-change}
在情形 \ref{quot-situation-complexes} 中，规则
$(T, g, E) \mapsto (T, g)$ 定义 $1$-态射
$$
\Complexesstack_{X/B} \longrightarrow \mathcal{S}_B
$$
它是群胚叠的态射
（见引理 \ref{quot-lemma-complexes-stack}、《代数叠》第
\ref{algebraic-section-split} 节，以及《叠的例子》第
\ref{examples-stacks-section-stack-associated-to-sheaf}).
设 $B' \to B$ 是 $S$ 上代数空间的态射。设
$\mathcal{S}_{B'} \to \mathcal{S}_B$ 是相应的集合纤维化叠的
$1$-态射。令 $X' = X \times_B B'$。
由基变换得到群胚叠
$\Complexesstack_{X'/B'} \to (\Sch/S)_{fppf}$，其中
$f' : X' \to B'$。
在此情形，图
$$
\begin{matrix}
\vcenter{\xymatrix{
\Complexesstack_{X'/B'} \ar[r] \ar[d] &
\Complexesstack_{X/B} \ar[d] \\
\mathcal{S}_{B'} \ar[r] & \mathcal{S}_B
}} \\
\text{或者换一种记号：} \\
\vcenter{\xymatrix{
\Complexesstack_{X'/B'} \ar[r] \ar[d] &
\Complexesstack_{X/B} \ar[d] \\
\Sch/B' \ar[r] & \Sch/B
}}
\end{matrix}
$$
是 $2$-纤维积方形。这个平凡的备注有时可用于改变基代数空间。
\end{remark}

\begin{lemma}
\label{quot-lemma-complexes-limits}
在情形 \ref{quot-situation-complexes} 中，假设 $B \to S$
局部有限表示。则
$p : \Complexesstack_{X/B} \to (\Sch/S)_{fppf}$ 保持极限
（《Artin 公理》定义 \ref{artin-definition-limit-preserving}）。
\end{lemma}

\begin{proof}
记 $B(T)$ 为离散范畴，其对象是 $S$-态射 $T \to B$。
设 $T = \lim T_i$ 是 $S$ 上仿射概形的滤极限。
将对象 $(T, h, E)$（它是 $\Complexesstack_{X/B, T}$ 的对象）
所对应的 $h$ 视为 $B(T)$ 中的对象，得到纤维范畴的交换图
$$
\xymatrix{
\colim \Complexesstack_{X/B, T_i} \ar[r] \ar[d] &
\Complexesstack_{X/B, T} \ar[d] \\
\colim B(T_i) \ar[r] & B(T)
}
$$
需证明上方水平箭头是等价。由于假设 $B$ 在 $S$ 上局部有限表示，
由《空间的极限》备注 \ref{spaces-limits-remark-limit-preserving}
可知下方水平箭头是等价。这意味着可假设 $T = \lim T_i$
是 $B$ 上仿射概形的滤极限。记相应态射
$g_i : T_i \to B$、$g : T \to B$。令
$X_i = T_i \times_{g_i, B} X$，$X_T = T \times_{g, B} X$。
注意 $X_T = \colim X_i$。
由《空间的更多态射》引理
\ref{spaces-more-morphisms-lemma-descend-relatively-perfect} 表明，
$T$-完美对象构成的 $D(\mathcal{O}_{X_T})$ 范畴是
$T_i$-完美对象构成的 $D(\mathcal{O}_{X_{T_i}})$ 范畴的余极限。
因此只需证明：给定 $T_i$-完美对象 $E_i$，属于
$D(\mathcal{O}_{X_{T_i}})$，且导出拉回 $E$ 是 $E_i$ 到 $X_T$ 的导出拉回，
满足引理 \ref{quot-lemma-complexes} 的条件 (2)，
则增大 $i$ 后，$E_i$ 满足引理 \ref{quot-lemma-complexes} 的条件 (2)。
令 $W \subset |T_i|$ 为引理
\ref{quot-lemma-complexes-open-neg-exts-vanishing} 为 $E_i$、$E_i$ 构造的开集。
由 $E$ 的假设，$T \to T_i$ 经由 $T$ 分解。
于是存在 $i' \geq i$ 使 $T_{i'} \to T_i$ 经由 $W$ 分解，见
《极限》引理 \ref{limits-lemma-limit-contained-in-constructible}。
$i'$ 按 $W$ 的构造即可。
\end{proof}

\begin{lemma}
\label{quot-lemma-complexes-RS-star}
在情形 \ref{quot-situation-complexes} 中，设
$$
\xymatrix{
Z \ar[r] \ar[d] & Z' \ar[d] \\
Y \ar[r] & Y'
}
$$
是 $S$ 上概形范畴中的推出，其中 $Z \to Z'$ 是有限阶加厚，
$Z \to Y$ 是仿射；见《更多态射》引理
\ref{more-morphisms-lemma-pushout-along-thickening}。
则纤维范畴上的函子
$$
\Complexesstack_{X/B, Y'}
\longrightarrow
\Complexesstack_{X/B, Y}
\times_{\Complexesstack_{X/B, Z}}
\Complexesstack_{X/B, Z'}
$$
是等价。
\end{lemma}

\begin{proof}
注意相应的映射
$$
B(Y') \longrightarrow B(Y) \times_{B(Z)} B(Z')
$$
是双射，见《空间的推出》引理
\ref{spaces-pushouts-lemma-pushout-along-thickening-schemes}.
因此利用交换图
$$
\xymatrix{
\Complexesstack_{X/B, Y'} \ar[r] \ar[d] &
\Complexesstack_{X/B, Y}
\times_{\Complexesstack_{X/B, Z}}
\Complexesstack_{X/B, Z'}
\ar[d] \\
B(Y') \ar[r] & B(Y) \times_{B(Z)} B(Z')
}
$$
可知可假设 $Y'$ 是 $B'$ 上的概形。由备注
\ref{quot-remark-complexes-base-change}，可将 $B$ 替换为 $Y'$，
并将 $X$ 替换为 $X \times_B Y'$。因此可假设 $B = Y'$。

\medskip\noindent
假设 $B = Y'$。先证明函子的全忠实性。
为此，设 $\xi_1, \xi_2$ 是 $\Complexesstack_{X/B}$ 中位于 $Y'$ 上的两个对象。
需证明
$$
\mathit{Isom}(\xi_1, \xi_2)(Y') \longrightarrow
\mathit{Isom}(\xi_1, \xi_2)(Y)
\times_{\mathit{Isom}(\xi_1, \xi_2)(Z)}
\mathit{Isom}(\xi_1, \xi_2)(Z')
$$
是双射。然而我们已经知道 $\mathit{Isom}(\xi_1, \xi_2)$
是 $B = Y'$ 上的代数空间。因此由《Artin 公理》引理
\ref{artin-lemma-pushout}（或上述《空间的推出》引理
\ref{spaces-pushouts-lemma-pushout-along-thickening-schemes}).

\medskip\noindent
本质满射性。设 $(E_Y, E_{Z'}, \alpha)$ 是三元组，其中
$E_Y \in D(\mathcal{O}_Y)$、$E_{Z'} \in D(\mathcal{O}_{X_{Z'}})$
是满足下列条件的对象：$(Y, Y \to B, E_Y)$ 是
$Y$ 上属于 $\Complexesstack_{X/B}$ 的对象，并且
$(Z', Z' \to B, E_{Z'})$ 是 $Z'$ 上属于 $\Complexesstack_{X/B}$ 的对象，
且 $\alpha : L(X_Z \to X_Y)^*E_Y \to L(X_Z \to X_{Z'})^*E_{Z'}$
是 $D(\mathcal{O}_{Z'})$ 中的同构。也就是说
$$
((Y, Y \to B, E_Y), (Z', Z' \to B, E_{Z'}), \alpha)
$$
是本引理箭头目标中的对象。
注意图
$$
\xymatrix{
X_Z \ar[r] \ar[d] & X_{Z'} \ar[d] \\
X_Y \ar[r] & X_{Y'}
}
$$
是推出，其中 $X_Z \to X_Y$ 仿射且 $X_Z \to X_{Z'}$ 是加厚
（见《空间的推出》引理
\ref{spaces-pushouts-lemma-equivalence-categories-spaces-pushout-flat}).
因此由《空间的推出》引理
\ref{spaces-pushouts-lemma-pushout-along-thickening-derived}
可找到对象 $E_{Y'} \in D(\mathcal{O}_{X_{Y'}})$
以及同构
$L(X_Y \to X_{Y'})^*E_{Y'} \to E_Y$ 及
$L(X_{Z'} \to X_{Y'})^*E_{Y'} \to E_Z$
与 $\alpha$ 相容。显然，只要证明 $E_{Y'}$ 是 $Y'$-完美的，
就完成了，因为引理 \ref{quot-lemma-complexes} 的条件 (2)
是点上的性质（且 $Y$ 与 $Y'$ 具有相同的点）。
这由《空间的更多态射》引理
\ref{spaces-more-morphisms-lemma-thickening-relatively-perfect}.
\end{proof}

\begin{lemma}
\label{quot-lemma-complexes-tangent-space}
在情形 \ref{quot-situation-complexes} 中，假设 $S$ 是局部诺特概形，
且 $B \to S$ 局部有限表示。
设 $k$ 是 $S$ 上有限型的域，且
$x_0 = (\Spec(k), g_0, E_0)$
是 $\mathcal{X} = \Complexesstack_{X/B}$ 上位于 $k$ 的对象。
则空间 $T\mathcal{F}_{\mathcal{X}, k, x_0}$ 和
$\text{Inf}(\mathcal{F}_{\mathcal{X}, k, x_0})$
（《Artin 公理》第 \ref{artin-section-tangent-spaces} 节）
是有限维的。
\end{lemma}

\begin{proof}
注意由引理 \ref{quot-lemma-complexes-RS-star}，
我们的群胚叠 $\mathcal{X}$ 满足《Artin 公理》第
\ref{artin-section-RS-star} 节定义的性质 (RS*)。
特别地，$\mathcal{X}$ 满足 (RS)。
因此所有相关的预形变范畴都是形变范畴
（《Artin 公理》引理 \ref{artin-lemma-deformation-category}），
上述陈述有意义。

\medskip\noindent
本段将说明可归约到 $B = \Spec(k)$ 的情形。
令 $X_0 = \Spec(k) \times_{g_0, B} X$，并记
$\mathcal{X}_0 = \Complexesstack_{X_0/k}$。在备注
\ref{quot-remark-complexes-base-change} 中已看到，
$\mathcal{X}_0$ 是 $2$-纤维积：$\mathcal{X}$ 与 $\Spec(k)$ 在 $B$ 上的纤维积，
作为 $(\Sch/S)_{fppf}$ 上的群胚纤维化范畴。
因此由《Artin 公理》引理
\ref{artin-lemma-fibre-product-tangent-spaces}，
只需证明 $B$、$\Spec(k)$、$\mathcal{X}_0$ 的切空间和无穷小自同构空间有限维。
根据《Artin 公理》引理 \ref{artin-lemma-finite-dimension}，
$B$ 和 $\Spec(k)$ 的切空间有限维，且显然它们的 $\text{Inf}$ 消失。
因此只需处理 $\mathcal{X}_0$。

\medskip\noindent
令 $k[\epsilon]$ 为 $k$ 上的对偶数。
令 $\Spec(k[\epsilon]) \to B$ 为 $g_0 : \Spec(k) \to B$
与 $\Spec(k[\epsilon]) \to \Spec(k)$ 的复合，后一个态射来自包含
$k \to k[\epsilon]$。
令 $X_0 = \Spec(k) \times_B X$，以及
$X_\epsilon = \Spec(k[\epsilon]) \times_B X$.
注意 $X_\epsilon$ 是 $X_0$ 的一阶加厚，
并且在一阶加厚 $\Spec(k) \to \Spec(k[\epsilon])$ 上平坦。
注意 $X_0$、$X_\epsilon$ 导出典范等价的小 \'etale 拓扑斯，见
《空间的更多态射》第
\ref{spaces-more-morphisms-section-thickenings}.
由《空间的更多态射》引理
\ref{spaces-more-morphisms-lemma-thickening-relatively-perfect}
可知 $T\mathcal{F}_{\mathcal{X}_0, k, x_0}$ 是
$E_0$ 到 $X_\epsilon$ 的提升的同构类集合（意义见《形变理论》引理
\ref{defos-lemma-first-order-defos-complex}）。
因此
$$
T\mathcal{F}_{\mathcal{X}_0, k, x_0} =
\Ext^1_{\mathcal{O}_{X_0}}(E_0, E_0)
$$
这里使用了模同一 $\epsilon k[\epsilon] \cong k$，作为 $k[\epsilon]$-模。
再次使用《形变理论》引理
\ref{defos-lemma-first-order-defos-complex}，可知存在满射
$$
\text{Inf}(\mathcal{F}_{\mathcal{X}, k, x_0})
\leftarrow
\Ext^0_{\mathcal{O}_{X_0}}(E_0, E_0)
$$
这是 $k$-向量空间的满射。由于 $E_0$ 伪凝聚，它属于
$D^-_{\textit{Coh}}(\mathcal{O}_{X_0})$，由《空间的导出范畴》引理
\ref{spaces-perfect-lemma-identify-pseudo-coherent-noetherian}。
由于 $E_0$ 局部具有有限 Tor 维且 $X_0$ 拟紧，可知
$E_0 \in D^b_{\textit{Coh}}(\mathcal{O}_{X_0})$。
因此上述 $\Ext$s 是有限维 $k$-向量空间，依据《空间的导出范畴》引理
\ref{spaces-perfect-lemma-ext-finite}。
\end{proof}

\begin{lemma}
\label{quot-lemma-complexes-strong-effectiveness}
在情形 \ref{quot-situation-complexes} 中，假设 $B = S$ 局部诺特。
则依《Artin 公理》备注 \ref{artin-remark-strong-effectiveness}
的意义，$p : \Complexesstack_{X/S} \to (\Sch/S)_{fppf}$
满足强形式有效性。
\end{lemma}

\begin{proof}
设 $(R_n)$ 是 $S$-代数的逆系统，过渡映射满射且核局部幂零。
令 $R = \lim R_n$。设 $(\xi_n)$ 是
$\Complexesstack_{X/B}$ 中位于 $(\Spec(R_n))$ 之上的对象系统。需证明
$(\xi_n)$ 有效，即存在对象 $\xi$，属于 $\Complexesstack_{X/B}$，位于
$\Spec(R)$ 之上。

\medskip\noindent
记 $X_R = \Spec(R) \times_S X$，$X_n = \Spec(R_n) \times_S X$。
显然 $X_n$ 是 $X_R$ 沿 $R \to R_n$ 的基变换。由于 $S = B$，
$\xi_n$ 对应于 $R_n$-完美对象
$E_n \in D(\mathcal{O}_{X_n})$，满足引理 \ref{quot-lemma-complexes} 的条件 (2)。
特别地，$E_n$ 伪凝聚。
同构 $\xi_{n + 1}|_{\Spec(R_n)} \cong \xi_n$
对应同构 $L(X_n \to X_{n + 1})^*E_{n + 1} \to E_n$。
因此由《空间上的平坦性》定理
\ref{spaces-flat-theorem-existence-derived}，
可找到伪凝聚对象 $E$，属于 $D(\mathcal{O}_{X_R})$，
使 $E_n$ 都等于 $E$ 的导出拉回（对所有 $n$），
并与过渡同构相容。

\medskip\noindent
注意 $(R, \Ker(R \to R_1))$ 是 Hensel 对，见《更多代数》引理
\ref{more-algebra-lemma-limit-henselian}。
特别地，$\Ker(R \to R_1)$ 包含于 $R$ 的 Jacobson 根中。
于是可应用《空间的更多态射》引理
\ref{spaces-more-morphisms-lemma-henselian-relatively-perfect}
可知 $E$ 是 $R$-完美的。

\medskip\noindent
最后需检验引理 \ref{quot-lemma-complexes} 的条件 (2)。
由引理 \ref{quot-lemma-complexes-open-neg-exts-vanishing}，
点 $t$ 属于 $\Spec(R)$ 且 $E_t$ 的负自扩张消失的点集是开集。
由于该条件在 $V(\Ker(R \to R_1))$ 中成立，且
$\Ker(R \to R_1)$ 包含于 $R$ 的 Jacobson 根中，
可知它对所有点成立。
\end{proof}

\begin{theorem}[真态射上复形模空间的代数性]
\label{quot-theorem-complexes-algebraic}
\begin{reference}
\cite{lieblich-complexes}
\end{reference}
设 $S$ 是概形。设 $f : X \to B$ 是
$S$ 上代数空间的态射。假设 $f$ 真、平坦且有限表示。
则 $\Complexesstack_{X/B}$ 是 $S$ 上的代数叠。
\end{theorem}

\begin{proof}
令 $\mathcal{X} = \Complexesstack_{X/B}$。我们已经看到，$\mathcal{X}$
是 $(\Sch/S)_{fppf}$ 上的群胚叠，且其对角态射由代数空间表示
（引理 \ref{quot-lemma-complexes-stack} 和 \ref{quot-lemma-complexes-diagonal}）。
因此只需找到概形 $W$ 以及满射光滑态射 $W \to \mathcal{X}$。

\medskip\noindent
设 $B'$ 是概形，且 $B' \to B$ 是满射 \'etale 态射。
令 $X' = B' \times_B X$，并记投影 $f' : X' \to B'$。
则 $\mathcal{X}' = \Complexesstack_{X'/B'}$ 等于 $2$-纤维积：
$\mathcal{X}$ 与由 $B'$ 关联的集合纤维化范畴在由 $B$ 关联的集合纤维化范畴上的
纤维积（备注 \ref{quot-remark-complexes-base-change}）。由
《代数叠》第 \ref{algebraic-section-representable-properties} 节的材料，
态射 $\mathcal{X}' \to \mathcal{X}$ 是满射且 \'etale。
因此只需对 $\mathcal{X}'$ 证明结论。换言之，可假设 $B$ 是概形。

\medskip\noindent
假设 $B$ 是概形。在此情形可将 $S$ 替换为 $B$，见
《代数叠》第 \ref{algebraic-section-change-base-scheme} 节。
因此可假设 $S = B$。

\medskip\noindent
假设 $S = B$。取仿射开覆盖 $S = \bigcup U_i$。
记 $\mathcal{X}_i$ 为 $\mathcal{X}$ 在
$(\Sch/U_i)_{fppf}$ 上的限制。若能找到概形 $W_i$（定义在 $U_i$ 上）以及满射光滑态射
$W_i \to \mathcal{X}_i$，则令 $W = \coprod W_i$，得到满射光滑态射
$W \to \mathcal{X}$。因此可假设 $S = B$ 是仿射的。

\medskip\noindent
假设 $S = B$ 仿射，记 $S = \Spec(\Lambda)$。
将 $\Lambda = \colim \Lambda_i$ 写成滤余极限，其中每个 $\Lambda_i$
在 $\mathbf{Z}$ 上有限型。对某个 $i$，可找到代数空间态射
$X_i \to \Spec(\Lambda_i)$，它真、平坦且有限表示，并且其到
$\Lambda$ 的基变换是 $X$。见《空间的极限》引理
\ref{spaces-limits-lemma-descend-finite-presentation},
\ref{spaces-limits-lemma-descend-flat}，以及
\ref{spaces-limits-lemma-eventually-proper}。
若证明 $\Complexesstack_{X_i/\Spec(\Lambda_i)}$ 是代数叠，则由基变换
（备注 \ref{quot-remark-complexes-base-change} 以及《代数叠》第
\ref{algebraic-section-change-base-scheme} 节）可知
$\mathcal{X}$ 是代数叠。因此可假设 $\Lambda$ 是有限型
$\mathbf{Z}$-代数。

\medskip\noindent
假设 $S = B = \Spec(\Lambda)$ 是有限型的仿射 $\mathbf{Z}$-代数。
此时将验证《Artin 公理》引理
\ref{artin-lemma-diagonal-representable} 的条件 (1)、(2)、(3)、(4)、(5)，
从而推出 $\mathcal{X}$ 是代数叠。
注意 $\Lambda$ 是 G-环，见《更多代数》命题
\ref{more-algebra-proposition-ubiquity-G-ring}。
因此 $S$ 的所有局部环都是 G-环，故 (5) 成立。
要检验 (2)，需验证《Artin 公理》第
\ref{artin-section-axioms} 节的公理 [-1]、[0]、[1]、[2]、[3]。
我们省略 [-1] 的验证；公理
[0]、[1]、[2]、[3] 分别对应引理 \ref{quot-lemma-complexes-stack}、
\ref{quot-lemma-complexes-limits},
\ref{quot-lemma-complexes-RS-star},
\ref{quot-lemma-complexes-tangent-space}。
条件 (3) 由引理 \ref{quot-lemma-complexes-strong-effectiveness} 得出。
条件 (1) 就是引理 \ref{quot-lemma-complexes-diagonal}。

\medskip\noindent
剩下的是证明条件 (4)，即形式光滑性的开放性。
为此将使用《Artin 公理》引理
\ref{artin-lemma-SGE-implies-openness-versality}。
我们已经看到 $\mathcal{X}$ 的对角态射由代数空间表示，满足 (RS*)，
并保持极限（见上面使用的引理）。
因此只需证明 $\mathcal{X}$ 满足《Artin 公理》引理
\ref{artin-lemma-SGE-implies-openness-versality} 中表述的强形式有效性。
这由引理 \ref{quot-lemma-complexes-strong-effectiveness} 得出，证明完成。
\end{proof}
% Generated mechanically from the frozen stacks-properties.tex target.
% Do not edit this generated file; edit the bound locale source instead.

% OK, start here.
%

\setcounter{chapter}{99}
\chapter{代数叠的性质}



\phantomsection
\label{stacks-properties-section-phantom}


\section{引言}
\label{stacks-properties-section-introduction}

\noindent
关于代数叠的简要介绍，见《代数叠》第
\ref{algebraic-section-introduction} 节；也请阅读该章的部分内容，
作为我们所用代数叠理论的基础。该章有意仔细区分概形、代数空间和代数叠；
从本章开始，我们采用通常的术语滥用，把这些概念交替使用。

\medskip\noindent
本章旨在介绍代数叠的一些基本概念和性质。
对于对角态射可表示的拟分离代数叠，一个基础参考文献是 \cite{LM-B}。



\section{约定与术语滥用}
\label{stacks-properties-section-conventions}

\noindent
我们选定一个大 fppf 位点 $\Sch_{fppf}$。
所有概形都包含在 $\Sch_{fppf}$ 中。
所考虑的每个环 $A$ 都满足：$\Spec(A)$（同构于）此大位点的一个对象。

\medskip\noindent
我们还固定一个基概形 $S$；按上述约定，它是 $\Sch_{fppf}$ 的一个元素。
只关心绝对情形的读者可以取 $S = \Spec(\mathbf{Z})$。

\medskip\noindent
关于代数叠，我们采用如下约定：
\begin{enumerate}
\item 所谓{\it 代数叠}，是指 $S$ 上的代数叠，即群胚上的纤维化范畴
$p : \mathcal{X} \to (\Sch/S)_{fppf}$
，它满足《代数叠》定义 \ref{algebraic-definition-algebraic-stack}
中的条件。
\item 称 $f : \mathcal{X} \to \mathcal{Y}$ 为{\it 代数叠的态射}，
是指 $S$ 上代数叠的 $1$-态射，即 $(\Sch/S)_{fppf}$ 上
群胚纤维化范畴的 $1$-态射；见《代数叠》定义
\ref{algebraic-definition-morphism-algebraic-stacks}。
\item 所谓 {\it $2$-态射} $\alpha : f \to g$，是指
$S$ 上代数叠的 $2$-范畴中的 $2$-态射；见《代数叠》定义
\ref{algebraic-definition-morphism-algebraic-stacks}。
\item 给定代数叠的态射 $\mathcal{X} \to \mathcal{Z}$ 和
$\mathcal{Y} \to \mathcal{Z}$，我们滥用术语，把 $2$-纤维积
$\mathcal{X} \times_\mathcal{Z} \mathcal{Y}$ 称为{\it 纤维积}。
\item 用 $\mathcal{X} \times_S \mathcal{Y}$ 表示代数叠
$\mathcal{X}$ 与 $\mathcal{Y}$ 的积。
\item 我们常滥用记号：若两个代数叠 $\mathcal{X}$ 和 $\mathcal{Y}$
在此 $2$-范畴中等价，就称它们{\it 同构}。
\end{enumerate}
关于代数空间，我们采用如下约定：
\begin{enumerate}
\item 若称 $X$ 是{\it 代数空间}，是指 $X$ 是 $S$ 上的代数空间，
即 $X$ 是 $(\Sch/S)_{fppf}$ 上满足《空间》定义
\ref{spaces-definition-algebraic-space} 中条件的预层。
\item {\it 代数空间的态射} $f :X \to Y$，是指《空间》定义
\ref{spaces-definition-morphism-algebraic-spaces} 所定义的
$S$ 上代数空间的态射。
\item 我们{\bf 不}区分代数空间 $X$ 与它所给出的代数叠
$\mathcal{S}_X \to (\Sch/S)_{fppf}$；见《代数叠》引理
\ref{algebraic-lemma-representable-algebraic}。
\item 特别地，从 $X$ 到代数叠 $\mathcal{Y}$ 的{\it 态射}
$f : X \to \mathcal{Y}$，是指代数叠的态射
$f : \mathcal{S}_X \to \mathcal{Y}$。态射 $\mathcal{Y} \to X$ 亦然。
\item 此外，给定代数叠 $\mathcal{X}$，称
{\it $\mathcal{X}$ 是代数空间}，是指 $\mathcal{X}$ 可由代数空间表示；
见《代数叠》定义
\ref{algebraic-definition-representable-by-algebraic-space}。
\item 我们采用如下记号约定：若用罗马体大写字母
（如 $X, Y, Z, A, B, \ldots$）表示一个代数叠，
则其惰性叠平凡，因而它是代数空间；见《代数叠》命题
\ref{algebraic-proposition-algebraic-stack-no-automorphisms}。
\end{enumerate}
关于概形，我们采用如下约定：
\begin{enumerate}
\item 若称 $X$ 是{\it 概形}，是指 $X$ 是 $S$ 上的概形，
即 $X$ 是 $(\Sch/S)_{fppf}$ 的对象。
\item 所谓{\it 概形的态射}，是指 $S$ 上概形的态射。
\item 我们{\bf 不}区分概形 $X$ 与它所给出的代数叠
$\mathcal{S}_X \to (\Sch/S)_{fppf}$；见《代数叠》引理
\ref{algebraic-lemma-representable-algebraic}。
\item 特别地，从概形 $X$ 到代数叠 $\mathcal{Y}$ 的{\it 态射}
$f : X \to \mathcal{Y}$，是指代数叠的态射
$f : \mathcal{S}_X \to \mathcal{Y}$。态射 $\mathcal{Y} \to X$ 亦然。
\item 此外，给定代数叠 $\mathcal{X}$，称
{\it $\mathcal{X}$ 是概形}，是指 $\mathcal{X}$ 可表示；
见《代数叠》第 \ref{algebraic-section-representable} 节。
\end{enumerate}
关于代数叠的态射，我们采用如下约定：
\begin{enumerate}
\item 若对每个概形 $T$ 和态射 $T \to \mathcal{Y}$，纤维积
$T \times_\mathcal{Y} \mathcal{X}$ 都是概形，则称代数叠的态射
$f : \mathcal{X} \to \mathcal{Y}$ {\it 可表示}，或
{\it 可由概形表示}。见《代数叠》第
\ref{algebraic-section-representable-morphism} 节。
\item 若对每个概形 $T$ 和态射 $T \to \mathcal{Y}$，纤维积
$T \times_\mathcal{Y} \mathcal{X}$ 都是代数空间，则称代数叠的态射
$f : \mathcal{X} \to \mathcal{Y}$ {\it 可由代数空间表示}。
见《代数叠》定义
\ref{algebraic-definition-representable-by-algebraic-spaces}。
在此情形，只要 $Z \to \mathcal{Y}$ 是源为代数空间的态射，
$Z \times_\mathcal{Y} \mathcal{X}$ 就是代数空间；见《代数叠》引理
\ref{algebraic-lemma-base-change-by-space-representable-by-space}。
\item 我们可以滥用术语：若代数叠的一个图在代数叠的
$2$-范畴中 $2$-交换，就称该图交换。
\end{enumerate}
注意，从代数空间到代数叠的每个态射 $X \to \mathcal{Y}$
都可由代数空间表示；见《代数叠》引理
\ref{algebraic-lemma-representable-diagonal}。
以后将不加说明地使用这一基本结果。




\section{可由代数空间表示的态射的性质}
\label{stacks-properties-section-properties-morphisms}

\noindent
我们将在另一章研究代数叠（任意）态射的性质。
对于可由代数空间表示的态射，我们知道满射、光滑、\'etale 等性质的含义。
这尤其适用于从代数空间到代数叠的态射 $X \to \mathcal{Y}$。
本节回顾其机制，列出适用的性质，并证明几个简单引理。

\medskip\noindent
第一个引理说明：若一个态射经某个平坦、局部有限表示且满射的态射作基变换后
可由代数空间表示，则它本身也可由代数空间表示。

\begin{lemma}
\label{stacks-properties-lemma-check-representable-covering}
设 $f : \mathcal{X} \to \mathcal{Y}$ 是代数叠的态射。
设 $W$ 是代数空间，且 $W \to \mathcal{Y}$ 满射、局部有限表示并平坦。
下列条件等价：
\begin{enumerate}
\item $f$ 可由代数空间表示；
\item $W \times_\mathcal{Y} \mathcal{X}$ 是代数空间。
\end{enumerate}
\end{lemma}

\begin{proof}
(1) $\Rightarrow$ (2) 是《代数叠》引理
\ref{algebraic-lemma-base-change-by-space-representable-by-space}。
反之，设 $W \to \mathcal{Y}$ 如 (2) 所述。为证明 (1)，
只需证明 $f$ 在纤维范畴上忠实；见《代数叠》引理
\ref{algebraic-lemma-characterize-representable-by-algebraic-spaces}。
条件 (2) 尤其蕴含 $W \times_\mathcal{Y} \mathcal{X} \to W$ 忠实。
于是由《叠》引理 \ref{stacks-lemma-faithful-descent}，$f$ 也忠实。
\end{proof}

\noindent
设 $P$ 是代数空间态射的一个性质，它在靶上 fppf 局部，
且在任意基变换下保持。设 $f : \mathcal{X} \to \mathcal{Y}$
是可由代数空间表示的代数叠态射。我们称
{\it $f$ 具有性质 $P$}，当且仅当对每个概形 $T$ 及态射
$T \to \mathcal{Y}$，代数空间的态射
$T \times_\mathcal{Y} \mathcal{X} \to T$ 具有性质 $P$；
见《代数叠》定义 \ref{algebraic-definition-relative-representable-property}。

\medskip\noindent
事实上，若 $f : \mathcal{X} \to \mathcal{Y}$ 可由代数空间表示且具有性质
$P$，则对代数叠的任意态射 $\mathcal{Y}' \to \mathcal{Y}$，
基变换
$\mathcal{Y}' \times_\mathcal{Y} \mathcal{X} \to \mathcal{Y}'$
也具有性质 $P$；见《代数叠》引理
\ref{algebraic-lemma-base-change-representable-by-spaces} 和
\ref{algebraic-lemma-base-change-representable-transformations-property}。
若性质 $P$ 在复合下保持，则同一结论也适用于可由代数空间表示的
代数叠态射；见《代数叠》引理
\ref{algebraic-lemma-composition-representable-by-spaces} 和
\ref{algebraic-lemma-composition-representable-transformations-property}。
此外，在此情形，可由代数空间表示且具有性质 $\mathcal{P}$ 的态射的积
$\mathcal{X}_1 \times \mathcal{X}_2 \to \mathcal{Y}_1 \times \mathcal{Y}_2$
也具有性质 $\mathcal{P}$；见《代数叠》引理
\ref{algebraic-lemma-product-representable-transformations-property}.

\medskip\noindent
最后，设 $P, P'$ 是代数空间态射的两个性质，它们在靶上 fppf 局部，
且在任意基变换下保持。若对每个态射 $f$ 都有
$P(f) \Rightarrow P'(f)$，则对可由代数空间表示的代数叠态射，
相应性质之间也有同一蕴含；见《代数叠》引理
\ref{algebraic-lemma-representable-transformations-property-implication}.
下文及后续各章将{\it 不加说明地}使用这一点。

\medskip\noindent
上述讨论适用于代数空间态射的下列各个性质：
\begin{enumerate}
\item 拟紧；见《空间的态射》引理
\ref{spaces-morphisms-lemma-base-change-quasi-compact}
以及《空间上的下降》引理
\ref{spaces-descent-lemma-descending-property-quasi-compact}；
\item 拟分离；见《空间的态射》引理
\ref{spaces-morphisms-lemma-base-change-separated}
以及《空间上的下降》引理
\ref{spaces-descent-lemma-descending-property-quasi-separated}；
\item 泛闭；见《空间的态射》引理
\ref{spaces-morphisms-lemma-base-change-universally-closed}
以及《空间上的下降》引理
\ref{spaces-descent-lemma-descending-property-universally-closed}；
\item 泛开；见《空间的态射》引理
\ref{spaces-morphisms-lemma-base-change-universally-open}
以及《空间上的下降》引理
\ref{spaces-descent-lemma-descending-property-universally-open}；
\item 泛商映射；见《空间的态射》引理
\ref{spaces-morphisms-lemma-base-change-universally-submersive}
以及《空间上的下降》引理
\ref{spaces-descent-lemma-descending-property-universally-submersive}；
\item 泛同胚；见《空间的态射》引理
\ref{spaces-morphisms-lemma-base-change-universal-homeomorphism}
以及《空间上的下降》引理
\ref{spaces-descent-lemma-descending-property-universal-homeomorphism}；
\item 满射；见《空间的态射》引理
\ref{spaces-morphisms-lemma-base-change-surjective}
以及《空间上的下降》引理
\ref{spaces-descent-lemma-descending-property-surjective}；
\item 泛单射；见《空间的态射》引理
\ref{spaces-morphisms-lemma-base-change-universally-injective}
以及《空间上的下降》引理
\ref{spaces-descent-lemma-descending-property-universally-injective}；
\item 局部有限型；见《空间的态射》引理
\ref{spaces-morphisms-lemma-base-change-finite-type}
以及《空间上的下降》引理
\ref{spaces-descent-lemma-descending-property-locally-finite-type}；
\item 局部有限表示；见《空间的态射》引理
\ref{spaces-morphisms-lemma-base-change-finite-presentation}
以及《空间上的下降》引理
\ref{spaces-descent-lemma-descending-property-locally-finite-presentation}；
\item 有限型；见《空间的态射》引理
\ref{spaces-morphisms-lemma-base-change-finite-type}
以及《空间上的下降》引理
\ref{spaces-descent-lemma-descending-property-finite-type}；
\item 有限表示；见《空间的态射》引理
\ref{spaces-morphisms-lemma-base-change-finite-presentation}
以及《空间上的下降》引理
\ref{spaces-descent-lemma-descending-property-finite-presentation}；
\item 平坦；见《空间的态射》引理
\ref{spaces-morphisms-lemma-base-change-flat}
以及《空间上的下降》引理
\ref{spaces-descent-lemma-descending-property-flat}；
\item 开浸入；见《空间的态射》第
\ref{spaces-morphisms-section-immersions} 节
以及《空间上的下降》引理
\ref{spaces-descent-lemma-descending-property-open-immersion}；
\item 同构；见《空间上的下降》引理
\ref{spaces-descent-lemma-descending-property-isomorphism}；
\item 仿射；见《空间的态射》引理
\ref{spaces-morphisms-lemma-base-change-affine}
以及《空间上的下降》引理
\ref{spaces-descent-lemma-descending-property-affine}；
\item 闭浸入；见《空间的态射》第
\ref{spaces-morphisms-section-immersions} 节
以及《空间上的下降》引理
\ref{spaces-descent-lemma-descending-property-closed-immersion}；
\item 分离；见《空间的态射》引理
\ref{spaces-morphisms-lemma-base-change-separated}
以及《空间上的下降》引理
\ref{spaces-descent-lemma-descending-property-separated}；
\item 真；见《空间的态射》引理
\ref{spaces-morphisms-lemma-base-change-proper}
以及《空间上的下降》引理
\ref{spaces-descent-lemma-descending-property-proper}；
\item 拟仿射；见《空间的态射》引理
\ref{spaces-morphisms-lemma-base-change-quasi-affine}
以及《空间上的下降》引理
\ref{spaces-descent-lemma-descending-property-quasi-affine}；
\item 整；见《空间的态射》引理
\ref{spaces-morphisms-lemma-base-change-integral}
以及《空间上的下降》引理
\ref{spaces-descent-lemma-descending-property-integral}；
\item 有限；见《空间的态射》引理
\ref{spaces-morphisms-lemma-base-change-integral}
以及《空间上的下降》引理
\ref{spaces-descent-lemma-descending-property-finite}；
\item （局部）拟有限；见《空间的态射》引理
\ref{spaces-morphisms-lemma-base-change-quasi-finite}
以及《空间上的下降》引理
\ref{spaces-descent-lemma-descending-property-quasi-finite}；
\item syntomic；见《空间的态射》引理
\ref{spaces-morphisms-lemma-base-change-syntomic}
以及《空间上的下降》引理
\ref{spaces-descent-lemma-descending-property-syntomic}；
\item 光滑；见《空间的态射》引理
\ref{spaces-morphisms-lemma-base-change-smooth}
以及《空间上的下降》引理
\ref{spaces-descent-lemma-descending-property-smooth}；
\item 非分歧；见《空间的态射》引理
\ref{spaces-morphisms-lemma-base-change-unramified}
以及《空间上的下降》引理
\ref{spaces-descent-lemma-descending-property-unramified}；
\item \'etale；见《空间的态射》引理
\ref{spaces-morphisms-lemma-base-change-etale}
以及《空间上的下降》引理
\ref{spaces-descent-lemma-descending-property-etale}；
\item 有限局部自由；见《空间的态射》引理
\ref{spaces-morphisms-lemma-base-change-finite-locally-free}
以及《空间上的下降》引理
\ref{spaces-descent-lemma-descending-property-finite-locally-free}；
\item 单态射；见《空间的态射》引理
\ref{spaces-morphisms-lemma-base-change-monomorphism}
以及《空间上的下降》引理
\ref{spaces-descent-lemma-descending-property-monomorphism}；
\item 浸入；见《空间的态射》第
\ref{spaces-morphisms-section-immersions} 节
以及《空间上的下降》引理
\ref{spaces-descent-lemma-descending-fppf-property-immersion}；
\item 局部分离；见《空间的态射》引理
\ref{spaces-morphisms-lemma-base-change-separated}
以及《空间上的下降》引理
\ref{spaces-descent-lemma-descending-fppf-property-locally-separated}。
\end{enumerate}

\begin{lemma}
\label{stacks-properties-lemma-property-spaces-too}
设 $P$ 是如上的代数空间态射的性质。
设 $f : \mathcal{X} \to \mathcal{Y}$ 是可由代数空间表示的代数叠态射。
下列条件等价：
\begin{enumerate}
\item $f$ 具有性质 $P$；
\item 对每个代数空间 $Z$ 以及态射 $Z \to \mathcal{Y}$，
态射 $Z \times_\mathcal{Y} \mathcal{X} \to Z$ 具有性质 $P$。
\end{enumerate}
\end{lemma}

\begin{proof}
(2) $\Rightarrow$ (1) 是直接的。设 (1) 成立。
取 (2) 中的 $Z \to \mathcal{Y}$。选取概形 $U$ 以及满的 \'etale 态射
$U \to Z$。由假设，态射
$U \times_\mathcal{Y} \mathcal{X} \to U$ 具有性质 $P$。但图表
$$
\xymatrix{
U \times_\mathcal{Y} \mathcal{X} \ar[d] \ar[r] &
Z \times_\mathcal{Y} \mathcal{X} \ar[d] \\
U \ar[r] & Z
}
$$
是笛卡儿的；由于 $\{U \to Z\}$ 是 fppf 覆盖，右侧竖直箭头因而具有
性质 $P$。
\end{proof}

\noindent
下一个引理说明，只需在沿一个满、平坦且局部有限表示的态射作基变换后
检验性质 $P$。

\begin{lemma}
\label{stacks-properties-lemma-check-property-covering}
设 $P$ 是如上的代数空间态射的性质。
设 $f : \mathcal{X} \to \mathcal{Y}$ 是可由代数空间表示的代数叠态射。
设 $W$ 是代数空间，并设 $W \to \mathcal{Y}$ 满、局部有限表示且平坦。
令 $V = W \times_\mathcal{Y} \mathcal{X}$。则
$$
(f\text{ 具有性质 }P) \Leftrightarrow (\text{投影 }V \to W\text{ 具有性质 }P).
$$
\end{lemma}

\begin{proof}
由引理 \ref{stacks-properties-lemma-property-spaces-too} 可得从左到右的蕴含。
设 $V \to W$ 具有性质 $P$。令 $T$ 为概形，并取态射
$T \to \mathcal{Y}$。考虑代数空间的交换图表
$$
\xymatrix{
T \times_\mathcal{Y} \mathcal{X} \ar[d] &
T \times_\mathcal{Y} V \ar[d] \ar[l] \ar[r] &
V \ar[d] \\
T & T \times_\mathcal{Y} W \ar[l] \ar[r] & W
}
$$
其中两个方块都是笛卡儿的。
左下方态射满、平坦且局部有限表示，故
$\{T \times_\mathcal{Y} V \to T\}$ 是 fppf 覆盖。
因此，右侧竖直箭头具有性质 $P$ 蕴含左侧竖直箭头也具有性质 $P$。
\end{proof}

\begin{lemma}
\label{stacks-properties-lemma-check-property-weak-covering}
设 $P$ 是如上的代数空间态射的性质。
设 $f : \mathcal{X} \to \mathcal{Y}$ 是可由代数空间表示的代数叠态射。
设 $\mathcal{Z} \to \mathcal{Y}$ 是可由代数空间表示、满、平坦且
局部有限表示的代数叠态射。
令 $\mathcal{W} = \mathcal{Z} \times_\mathcal{Y} \mathcal{X}$。则
$$
(f\text{ 具有性质 }P) \Leftrightarrow
(\text{投影 }\mathcal{W} \to \mathcal{Z}\text{ 具有性质 }P).
$$
\end{lemma}

\begin{proof}
选取代数空间 $W$ 以及满、平坦且局部有限表示的态射
$W \to \mathcal{Z}$。由上述讨论，复合
$W \to \mathcal{Y}$ 也满、平坦且局部有限表示。记
$V = W \times_\mathcal{Z} \mathcal{W} = V \times_\mathcal{Y} \mathcal{X}$.
由引理 \ref{stacks-properties-lemma-check-property-covering} 可知，$f$ 具有性质
$\mathcal{P}$ 当且仅当 $V \to W$ 具有该性质，而
$\mathcal{W} \to \mathcal{Z}$ 具有性质 $\mathcal{P}$ 当且仅当
$V \to W$ 具有该性质。引理得证。
\end{proof}

\begin{lemma}
\label{stacks-properties-lemma-check-property-after-precomposing}
\leavevmode\par
设 $P$ 是如上的代数空间态射的性质。
设 $\tau \in \{\etale, smooth, syntomic, fppf\}$。
设 $\mathcal{X} \to \mathcal{Y}$ 和 $\mathcal{Y} \to \mathcal{Z}$
是可由代数空间表示的代数叠态射。假设
\begin{enumerate}
\item $\mathcal{X} \to \mathcal{Y}$ 满，并且是 \'etale、光滑、syntomic，
或者平坦且局部有限表示；
\item 该复合具有性质 $P$；
\item $P$ 在 $\tau$ 拓扑下对源是局部的。
\end{enumerate}
则 $\mathcal{Y} \to \mathcal{Z}$ 具有性质 $P$。
\end{lemma}

\begin{proof}
令 $Z$ 为概形，并取态射 $Z \to \mathcal{Z}$。令
$X = \mathcal{X} \times_\mathcal{Z} Z$、
$Y = \mathcal{Y} \times_\mathcal{Z} Z$。由 (1)，$\{X \to Y\}$
是代数空间的一个 $\tau$ 覆盖；由 (2)，$X \to Z$ 具有性质 $P$。
再由 (3)，这蕴含 $Y \to Z$ 具有性质 $P$，结论得证。
\end{proof}

\begin{lemma}
\label{stacks-properties-lemma-representable-in-terms-presentations}
设 $g : \mathcal{X}' \to \mathcal{X}$ 是可由代数空间表示的代数叠态射。
设 $[U/R] \to \mathcal{X}$ 是一个表示。令
$U' = U \times_\mathcal{X} \mathcal{X}'$、
$R' = R \times_\mathcal{X} \mathcal{X}'$。
则存在形如 $(U', R', s', t', c')$ 的代数空间群胚和一个表示
$[U'/R'] \to \mathcal{X}'$，并且图表
$$
\xymatrix{
[U'/R'] \ar[d]_{[\text{pr}]} \ar[r] & \mathcal{X}' \ar[d]^g \\
[U/R] \ar[r] & \mathcal{X}
}
$$
是 $2$-交换的，其中态射 $[\text{pr}]$ 来自群胚态射
$\text{pr} : (U', R', s', t', c') \to (U, R, s, t, c)$.
\end{lemma}

\begin{proof}
由于 $U \to \mathcal{Y}$ 满且光滑（见《代数叠》引理
\ref{algebraic-lemma-smooth-quotient-smooth-presentation}），
基变换 $U' \to \mathcal{X}'$ 也满且光滑。
因此，由《代数叠》引理 \ref{algebraic-lemma-stack-presentation}，
为得到光滑群胚 $(U', R', s', t', c')$ 和表示
$[U'/R'] \to \mathcal{X}'$，只需证明
$R' = U' \times_{\mathcal{X}'} U'$。
利用 $R = V \times_\mathcal{Y} V$（见《空间中的群胚》引理
\ref{spaces-groupoids-lemma-quotient-stack-2-cartesian}），
这由下式推出：
$$
R' =
U \times_\mathcal{X} U \times_\mathcal{X} \mathcal{X}' =
(U \times_\mathcal{X} \mathcal{X}')
\times_{\mathcal{X}'}
(U \times_\mathcal{X} \mathcal{X}')
$$
见《范畴》引理 \ref{categories-lemma-associativity-2-fibre-product} 和
\ref{categories-lemma-2-fibre-product-erase-factor}。
显然，投影态射 $U' \to U$ 和 $R' \to R$ 给出所需的群胚态射
$\text{pr} : (U', R', s', t', c') \to (U, R, s, t, c)$.
于是商叠的态射 $[\text{pr}]$，由《空间中的群胚》引理
\ref{spaces-groupoids-lemma-quotient-stack-functorial}。

\medskip\noindent
还需证明该图表 $2$-交换。显然，图表
$$
\xymatrix{
U' \ar[d]_{\text{pr}_U} \ar[r]_{f'} & \mathcal{X}' \ar[d]^g \\
U \ar[r]^f & \mathcal{X}
}
$$
$2$-交换，其中 $\text{pr}_U : U' \to U$ 是投影。
由 $R = U \times_\mathcal{X} U$、$t = \text{pr}_0$ 和
$s = \text{pr}_1$，在 $\Mor(R, \mathcal{X})$ 中有典范 $2$-箭头
$\tau : f \circ t \to f \circ s$。利用同构
$R' \to U' \times_{\mathcal{X}'} U'$，类似地得到同构
$\tau' : f' \circ t' \to f' \circ s'$。注意
$g \circ f' \circ t' = f \circ t \circ \text{pr}_R$ 以及
$g \circ f' \circ s' = f \circ s \circ \text{pr}_R$，其中
$\text{pr}_R : R' \to R$ 是投影。因此，可以考察是否有
\begin{equation}
\label{stacks-properties-equation-verify}
\tau \star \text{id}_{\text{pr}_R} = \text{id}_g \star \tau'.
\end{equation}
现作两个断言：(1) 若方程 (\ref{stacks-properties-equation-verify}) 成立，则该图表
$2$-交换；(2) 方程 (\ref{stacks-properties-equation-verify}) 成立。
我们略去两个断言的证明。提示：(1) 来自《代数叠》引理
\ref{algebraic-lemma-map-space-into-stack} 中
$f = f_{can}$ 和 $f' = f'_{can}$ 的构造；(2) 可由仔细展开定义得到。
\end{proof}

\begin{remark}
\label{stacks-properties-remark-representable-over}
设 $\mathcal{Y}$ 是代数叠。考虑如下 $2$-范畴：
\begin{enumerate}
\item 对象是可由代数空间表示的态射
$f : \mathcal{X} \to \mathcal{Y}$；
\item 一个 $1$-态射
$(g, \beta) :
(f_1 : \mathcal{X}_1 \to \mathcal{Y})
\to
(f_2 : \mathcal{X}_2 \to \mathcal{Y})$
由态射 $g : \mathcal{X}_1 \to \mathcal{X}_2$ 和 $2$-态射
$\beta : f_1 \to f_2 \circ g$ 组成；
\item 两个这样的态射
$(g, \beta), (g', \beta') :
(f_1 : \mathcal{X}_1 \to \mathcal{Y})
\to
(f_2 : \mathcal{X}_2 \to \mathcal{Y})$
之间的 $2$-态射，是满足
$(\text{id}_{f_2} \star \alpha) \circ \beta = \beta'$
的 $2$-态射 $\alpha : g \to g'$。
\end{enumerate}
类比《空间上的拓扑》第 \ref{spaces-topologies-section-procedure} 节的记号，
将此 $2$-范畴记作 $\textit{Spaces}/\mathcal{Y}$。
我们断言，在此 $2$-范畴中，态射范畴
$$
\Mor_{\textit{Spaces}/\mathcal{Y}}(
(f_1 : \mathcal{X}_1 \to \mathcal{Y}),
(f_2 : \mathcal{X}_2 \to \mathcal{Y}))
$$
全都是集合胚。事实上，一个 $2$-态射 $\alpha$ 是如下规则：它为
$\mathcal{X}_1$ 的每个对象 $x_1$ 指定相关纤维范畴
$\mathcal{X}_2$ 中的一个同构
$\alpha_{x_1} : g(x_1) \longrightarrow g'(x_1)$
使得图表
$$
\xymatrix{
& f_2(x_1) \ar[ld]_{\beta_{x_1}} \ar[rd]^{\beta'_{x_1}} \\
f_2(g(x_1)) \ar[rr]^{f_2(\alpha_{x_1})} & &
f_2(g'(x_1))
}
$$
交换。但由于 $f_2$ 忠实（见《代数叠》引理
\ref{algebraic-lemma-characterize-representable-by-algebraic-spaces}），
这意味着只要 $\alpha_{x_1}$ 存在，它便是唯一的！换言之，
$2$-范畴 $\textit{Spaces}/\mathcal{Y}$ 与一个范畴非常接近：
若用 $1$-态射的同构类替代 $1$-态射，就得到一个范畴。
我们常常不加说明地作此替换。
\end{remark}




\section{代数叠的点}
\label{stacks-properties-section-points}

\noindent
设 $\mathcal{X}$ 是代数叠，$K, L$ 是两个域，并设
$p : \Spec(K) \to \mathcal{X}$ 和
$q : \Spec(L) \to \mathcal{X}$ 是态射。
若存在域 $\Omega$ 和一个 $2$-交换图表
$$
\xymatrix{
\Spec(\Omega) \ar[r] \ar[d] &
\Spec(L) \ar[d]^q \\
\Spec(K) \ar[r]^p &
\mathcal{X}.
}
$$
则称 $p$ 与 $q$ {\it 等价}。

\begin{lemma}
\label{stacks-properties-lemma-equivalence}
上述概念确实在从域的谱到代数叠 $\mathcal{X}$ 的态射上定义了一个等价关系。
\end{lemma}

\begin{proof}
显然，该关系自反且对称。因此只需证明其传递性。
这归结为如下问题：给定图表
$$
\xymatrix{
\Spec(\Omega) \ar[r]_b \ar[d]_a &
\Spec(L) \ar[d]^q & \Spec(\Omega') \ar[l]^{b'} \ar[d]^{a'} \\
\Spec(K) \ar[r]^p &
\mathcal{X} &
\Spec(K') \ar[l]_{p'}
}
$$
且两个方块都 $2$-交换，需要证明 $p$ 与 $p'$ 等价。
由 $2$-Yoneda 引理（见《代数叠》第
\ref{algebraic-section-2-yoneda} 节），态射 $p$、$p'$ 和 $q$
分别由 $\mathcal{X}$ 在 $\Spec(K)$、$\Spec(K')$ 和 $\Spec(L)$
上纤维范畴中的对象 $x$、$x'$ 和 $y$ 给出。
两个方块的 $2$-交换性意味着存在同构
$\alpha : a^*x \to b^*y$ 和 $\alpha' : (a')^*x' \to (b')^*y$
分别位于 $\mathcal{X}$ 在 $\Spec(\Omega)$ 和 $\Spec(\Omega')$
上的纤维范畴中。任取域 $\Omega''$ 以及在 $L$ 上相同的嵌入
$\Omega \to \Omega''$ 和 $\Omega' \to \Omega''$。
于是可将上图扩充为
$$
\xymatrix{
& \Spec(\Omega'') \ar[ld]_c \ar[d]^{q'} \ar[rd]^{c'} \\
\Spec(\Omega) \ar[r]_b \ar[d]_a &
\Spec(L) \ar[d]^q & \Spec(\Omega') \ar[l]^{b'} \ar[d]^{a'} \\
\Spec(K) \ar[r]^p &
\mathcal{X} &
\Spec(K') \ar[l]_{p'}
}
$$
其中两个三角形交换，并且
$$
(q')^*(\alpha')^{-1} \circ (q')^*\alpha :
(a \circ c)^*x
\longrightarrow
(a' \circ c')^*x'
$$
是 $\Spec(\Omega'')$ 上纤维范畴中的同构。
故 $p$ 与 $p'$ 等价，正如所需。
\end{proof}

\begin{definition}
\label{stacks-properties-definition-points}
设 $\mathcal{X}$ 是代数叠。$\mathcal{X}$ 的一个{\it 点}，是从域的谱到
$\mathcal{X}$ 的态射所组成的一个等价类。
$\mathcal{X}$ 的点集记作 $|\mathcal{X}|$。
\end{definition}

\noindent
这与代数空间之点的定义一致；见《空间的性质》定义
\ref{spaces-properties-definition-points}。
此外，对概形而言，这恢复了通常的点的概念；见《空间的性质》引理
\ref{spaces-properties-lemma-scheme-points}。
若 $f : \mathcal{X} \to \mathcal{Y}$ 是代数叠态射，则它诱导映射
$|f| : |\mathcal{X}| \to |\mathcal{Y}|$，将代表元
$x : \Spec(K) \to \mathcal{X}$ 映为代表元
$f \circ x : \Spec(K) \to \mathcal{Y}$。该映射适定：
$2$-同构的 $1$-态射与一个 $1$-态射前复合或后复合后仍 $2$-同构，
因为可以与所给 $1$-态射的恒等态射作水平前复合或后复合。
这在任意（严格）$(2, 1)$-范畴中都成立。若
$$
\xymatrix{
\mathcal{X} \ar[d] \ar[r] & \mathcal{Y} \ar[d] \\
\mathcal{W} \ar[r] & \mathcal{Z}
}
$$
是代数叠的 $2$-交换图表，则集合图表
$$
\xymatrix{
|\mathcal{X}| \ar[d] \ar[r] & |\mathcal{Y}| \ar[d] \\
|\mathcal{W}| \ar[r] & |\mathcal{Z}|
}
$$
交换。特别地，若 $\mathcal{X} \to \mathcal{Y}$ 是等价，
则 $|\mathcal{X}| \to |\mathcal{Y}|$ 是双射。

\begin{lemma}
\label{stacks-properties-lemma-points-cartesian}
设
$$
\xymatrix{
\mathcal{Z} \times_\mathcal{Y} \mathcal{X} \ar[r] \ar[d] &
\mathcal{X} \ar[d] \\
\mathcal{Z} \ar[r] & \mathcal{Y}
}
$$
是代数叠的纤维积。则点集映射
$$
|\mathcal{Z} \times_\mathcal{Y} \mathcal{X}|
\longrightarrow
|\mathcal{Z}| \times_{|\mathcal{Y}|} |\mathcal{X}|
$$
是满射。
\end{lemma}

\begin{proof}
事实上，设给定域 $K$、$L$ 和态射
$\Spec(K) \to \mathcal{X}$、$\Spec(L) \to \mathcal{Z}$。
它们作为 $|\mathcal{Y}|$ 的元素相同这一假设，意味着存在公共扩张
$M/K$ 和 $M/L$，使得
$\Spec(M) \to \Spec(K) \to \mathcal{X} \to \mathcal{Y}$ 和
$\Spec(M) \to \Spec(L) \to \mathcal{Z} \to \mathcal{Y}$
是 $2$-同构的。这恰好就是得到态射
$\Spec(M) \to \mathcal{Z} \times_\mathcal{Y} \mathcal{X}$ 的条件。
\end{proof}

\begin{lemma}
\label{stacks-properties-lemma-characterize-surjective}
设 $f : \mathcal{X} \to \mathcal{Y}$ 是可由代数空间表示的代数叠态射。
下列条件等价：
\begin{enumerate}
\item $|f| : |\mathcal{X}| \to |\mathcal{Y}|$ 是满射；
\item $f$ 是满射（取第 \ref{stacks-properties-section-properties-morphisms} 节的含义）。
\end{enumerate}
\end{lemma}

\begin{proof}
设 (1) 成立。令 $T \to \mathcal{Y}$ 为源是概形的态射。
为证明 (2)，需证明代数空间态射
$T \times_\mathcal{Y} \mathcal{X} \to T$ 是满射。
由《空间的态射》定义 \ref{spaces-morphisms-definition-surjective}，
这意味着要证明 $|T \times_\mathcal{Y} \mathcal{X}| \to |T|$ 是满射。
应用引理 \ref{stacks-properties-lemma-points-cartesian}，可见这由 (1) 推出。

\medskip\noindent
反之，设 (2) 成立。令 $y : \Spec(K) \to \mathcal{Y}$
为从域的谱到 $\mathcal{Y}$ 的态射。由假设，代数空间态射
$\Spec(K) \times_{y, \mathcal{Y}} \mathcal{X} \to \Spec(K)$
是满射。由《空间的态射》定义
\ref{spaces-morphisms-definition-surjective}，这意味着存在域扩张
$K'/K$ 和态射
$\Spec(K') \to \Spec(K) \times_{y, \mathcal{Y}} \mathcal{X}$
使下图中的左方块交换：
$$
\xymatrix{
\Spec(K') \ar[r] \ar[d] &
\Spec(K) \times_{y, \mathcal{Y}} \mathcal{X} \ar[d] \ar[r] &
\mathcal{X} \ar[d]
\\
\Spec(K) \ar@{=}[r] &
\Spec(K) \ar[r]^-y &
\mathcal{Y}
}
$$
这说明 $|X| \to |\mathcal{Y}|$ 是满射。
\end{proof}

\noindent
下面的引理说明如何利用一个表示计算点集。

\begin{lemma}
\label{stacks-properties-lemma-points-presentation}
设 $\mathcal{X}$ 是代数叠。设 $\mathcal{X} = [U/R]$
是 $\mathcal{X}$ 的一个表示；见《代数叠》定义
\ref{algebraic-definition-presentation}。
则 $|R| \to |U| \times |U|$ 的像是一个等价关系，
且 $|\mathcal{X}|$ 是 $|U|$ 关于该等价关系的商。
\end{lemma}

\begin{proof}
该假设意味着我们有代数空间中的一个光滑群胚
$(U, R, s, t, c)$，以及等价
$f : [U/R] \to \mathcal{X}$。可以假设 $\mathcal{X} = [U/R]$。
诱导态射 $p : U \to \mathcal{X}$ 光滑且满；见《代数叠》引理
\ref{algebraic-lemma-smooth-quotient-smooth-presentation}。
因此，由引理 \ref{stacks-properties-lemma-characterize-surjective}，
$|U| \to |\mathcal{X}|$ 是满射。注意
$R = U \times_\mathcal{X} U$；见《空间中的群胚》引理
\ref{spaces-groupoids-lemma-quotient-stack-2-cartesian}。
故引理 \ref{stacks-properties-lemma-points-cartesian} 蕴含映射
$$
|R| \longrightarrow |U| \times_{|\mathcal{X}|} |U|
$$
是满射。因此，$|R| \to |U| \times |U|$ 的像恰好由所有满足
$u_1$ 与 $u_2$ 在 $|\mathcal{X}|$ 中像相同的偶
$(u_1, u_2) \in |U| \times |U|$ 构成。
合并这两个陈述即得引理。
\end{proof}

\begin{remark}
\label{stacks-properties-remark-more-general-presentation}
引理 \ref{stacks-properties-lemma-points-presentation} 的结果可作如下推广。
设 $\mathcal{X}$ 是代数叠。设 $U$ 是代数空间，并设
$f : U \to \mathcal{X}$ 是满射（由第
\ref{stacks-properties-section-properties-morphisms} 节，这一说法有意义）。
令 $R = U \times_\mathcal{X} U$，令 $(U, R, s, t, c)$
为代数空间中的群胚，并令 $f_{can} : [U/R] \to \mathcal{X}$
为《代数叠》引理 \ref{algebraic-lemma-map-space-into-stack}
中构造的典范态射。
则 $|R| \to |U| \times |U|$ 的像是一个等价关系，且
$|\mathcal{X}| = |U|/|R|$。引理 \ref{stacks-properties-lemma-points-presentation}
的证明无需改动便可适用。（当然，一般而言 $[U/R]$ 不是代数叠，
且一般而言 $f_{can}$ 不是同构。）
\end{remark}

\begin{lemma}
\label{stacks-properties-lemma-topology-points}
在代数叠的点集上，存在唯一的拓扑满足下列性质：
\begin{enumerate}
\item 对每个代数叠态射 $\mathcal{X} \to \mathcal{Y}$，
映射 $|\mathcal{X}| \to |\mathcal{Y}|$ 连续；
\item 对每个平坦且局部有限表示的态射 $U \to \mathcal{X}$，
其中 $U$ 是代数空间，拓扑空间映射
$|U| \to |\mathcal{X}|$ 连续且开。
\end{enumerate}
\end{lemma}

\begin{proof}
选取一个满、平坦且局部有限表示的态射
$p : U \to \mathcal{X}$，其中 $U$ 是代数空间。
由代数叠的定义，这样的态射存在，因为光滑态射平坦且局部有限表示
（见《空间的态射》引理
\ref{spaces-morphisms-lemma-smooth-locally-finite-presentation} 和
\ref{spaces-morphisms-lemma-smooth-flat}）。
按如下规则在 $|\mathcal{X}|$ 上定义拓扑：
$W \subset |\mathcal{X}|$ 是开集，当且仅当
$|p|^{-1}(W)$ 在 $|U|$ 中是开集。
为证明这与 $p$ 的选择无关，令 $p' : U' \to \mathcal{X}$
为另一个从代数空间到 $\mathcal{X}$ 的满、平坦且局部有限表示的态射。
令 $U'' = U \times_\mathcal{X} U'$，于是有 $2$-交换图表
$$
\xymatrix{
U'' \ar[r] \ar[d] & U' \ar[d] \\
U \ar[r] & \mathcal{X}
}
$$
由于 $U \to \mathcal{X}$ 和 $U' \to \mathcal{X}$ 满、平坦且
局部有限表示，由引理 \ref{stacks-properties-lemma-property-spaces-too} 可知
$U'' \to U'$ 和 $U'' \to U$ 也满、平坦且局部有限表示。
因此，映射 $|U''| \to |U'|$ 和 $|U''| \to |U|$
连续、开且满；见《空间的态射》定义
\ref{spaces-morphisms-definition-surjective} 和引理
\ref{spaces-morphisms-lemma-fppf-open}。
这显然说明我们的定义与 $p : U \to \mathcal{X}$ 的选择无关。

\medskip\noindent
设 $f : \mathcal{X} \to \mathcal{Y}$ 是代数叠态射。
由《代数叠》引理 \ref{algebraic-lemma-lift-morphism-presentations}，
可以找到一个 $2$-交换图表
$$
\xymatrix{
U \ar[d]_x \ar[r]_a & V \ar[d]^y \\
\mathcal{X} \ar[r]^f & \mathcal{Y}
}
$$
其中竖直箭头满且光滑。考虑相应的集合交换图表
$$
\xymatrix{
|U| \ar[d]_{|x|} \ar[r]_{|a|} & |V| \ar[d]^{|y|} \\
|\mathcal{X}| \ar[r]^{|f|} & |\mathcal{Y}|
}
$$
若 $W \subset |\mathcal{Y}|$ 是开集，则由上述定义，这恰好意味着
$|y|^{-1}(W)$ 在 $|V|$ 中是开集。由于 $|a|$ 连续，可得
$|a|^{-1}|y|^{-1}(W) = |x|^{-1}|f|^{-1}(W)$ 在 $|W|$ 中是开集；
按定义，这意味着 $|f|^{-1}(W)$ 在 $|\mathcal{X}|$ 中是开集。
故 $|f|$ 连续。

\medskip\noindent
最后，需要证明：若 $U$ 是代数空间，且
$U \to \mathcal{X}$ 平坦并局部有限表示，则
$|U| \to |\mathcal{X}|$ 是开映射。令 $V \to \mathcal{X}$
为满、平坦且局部有限表示的态射，其中 $V$ 是代数空间。
考虑交换图表
$$
\xymatrix{
|U \times_\mathcal{X} V| \ar[r]_e \ar[rd]_f &
|U| \times_{|\mathcal{X}|} |V| \ar[d]_c \ar[r]_d &
|V| \ar[d]^b \\
& |U| \ar[r]^a & |\mathcal{X}|
}
$$
态射 $U \times_\mathcal{X} V \to U$ 是满射，亦即
$f : |U \times_\mathcal{X} V| \to |U|$ 是满射。
左上方的水平箭头是满射；见引理 \ref{stacks-properties-lemma-points-cartesian}。
态射 $U \times_\mathcal{X} V \to V$ 平坦且局部有限表示，
故 $d \circ e : |U \times_\mathcal{X} V| \to |V|$ 是开映射；
见《空间的态射》引理 \ref{spaces-morphisms-lemma-fppf-open}。
任取开集 $W \subset |U|$。上述性质蕴含
$b^{-1}(a(W)) = (d \circ e)(f^{-1}(W))$ 是开集；
按构造，这意味着 $a(W)$ 是开集，正如所需。
\end{proof}

\begin{definition}
\label{stacks-properties-definition-topological-space}
设 $\mathcal{X}$ 是代数叠。$\mathcal{X}$ 的底层{\it 拓扑空间}，
是赋予引理 \ref{stacks-properties-lemma-topology-points} 中所构造拓扑的点集
$|\mathcal{X}|$。
\end{definition}

\noindent
若 $\mathcal{X}$ 是代数空间，则此定义与
$|\mathcal{X}|$ 上已有的拓扑一致。

\begin{lemma}
\label{stacks-properties-lemma-space-locally-quasi-compact}
设 $\mathcal{X}$ 是代数叠。$|\mathcal{X}|$ 的每个点都有一个
由拟紧开邻域组成的基本邻域系。
特别地，$|\mathcal{X}|$ 在《拓扑》定义
\ref{topology-definition-locally-quasi-compact} 的意义下局部拟紧。
\end{lemma}

\begin{proof}
这形式地来自如下事实：存在概形 $U$ 以及拓扑空间的满、开、连续映射
$U \to |\mathcal{X}|$。事实上，若 $U \to \mathcal{X}$ 满且光滑，
则引理 \ref{stacks-properties-lemma-topology-points} 保证
$|U| \to |\mathcal{X}|$ 连续、满且开。
\end{proof}











\section{满态射}
\label{stacks-properties-section-surjective}

\noindent
设 $f : \mathcal{X} \to \mathcal{Y}$ 是可由代数空间表示的代数叠态射。
我们已经在第 \ref{stacks-properties-section-properties-morphisms} 节定义了
$f$ 为满射的含义。引理 \ref{stacks-properties-lemma-characterize-surjective}
说明，这等价于要求
$|f| : |\mathcal{X}| \to |\mathcal{Y}|$ 是满射。
因此可以给出如下定义。

\begin{definition}
\label{stacks-properties-definition-surjective}
设 $f : \mathcal{X} \to \mathcal{Y}$ 是代数叠态射。
若相应拓扑空间的映射
$|f| : |\mathcal{X}| \to |\mathcal{Y}|$ 是满射，
则称 $f$ {\it 满}。
\end{definition}

\noindent
下面给出一些引理。

\begin{lemma}
\label{stacks-properties-lemma-composition-surjective}
满态射的复合是满态射。
\end{lemma}

\begin{proof}
略。
\end{proof}

\begin{lemma}
\label{stacks-properties-lemma-base-change-surjective}
满态射的基变换是满态射。
\end{lemma}

\begin{proof}
略。提示：使用引理 \ref{stacks-properties-lemma-points-cartesian}。
\end{proof}

\begin{lemma}
\label{stacks-properties-lemma-descent-surjective}
设 $f : \mathcal{X} \to \mathcal{Y}$ 是代数叠态射。
设 $\mathcal{Y}' \to \mathcal{Y}$ 是代数叠的满态射。
若 $f$ 的基变换
$f' : \mathcal{Y}' \times_\mathcal{Y} \mathcal{X}
\to \mathcal{Y}'$ 是满射，则 $f$ 是满射。
\end{lemma}

\begin{proof}
由引理 \ref{stacks-properties-lemma-points-cartesian} 立即可得。
\end{proof}

\begin{lemma}
\label{stacks-properties-lemma-surjective-permanence}
设 $\mathcal{X} \to \mathcal{Y} \to \mathcal{Z}$ 是代数叠态射。
若 $\mathcal{X} \to \mathcal{Z}$ 是满射，则
$\mathcal{Y} \to \mathcal{Z}$ 也是满射。
\end{lemma}

\begin{proof}
立即可得。
\end{proof}












\section{拟紧代数叠}
\label{stacks-properties-section-quasi-compact}

\noindent
由《空间的性质》引理
\ref{spaces-properties-lemma-quasi-compact-space}，
下述定义与代数空间的相应定义等价。

\begin{definition}
\label{stacks-properties-definition-quasi-compact}
设 $\mathcal{X}$ 是代数叠。当且仅当
$|\mathcal{X}|$ 拟紧时，称 $\mathcal{X}$ {\it 拟紧}。
\end{definition}

\begin{lemma}
\label{stacks-properties-lemma-quasi-compact-stack}
设 $\mathcal{X}$ 是代数叠。下列条件等价：
\begin{enumerate}
\item $\mathcal{X}$ 拟紧；
\item 存在满的光滑态射 $U \to \mathcal{X}$，其中 $U$ 是仿射概形；
\item 存在满的光滑态射 $U \to \mathcal{X}$，其中 $U$ 是拟紧概形；
\item 存在满的光滑态射 $U \to \mathcal{X}$，其中 $U$ 是拟紧代数空间；
\item 存在代数叠的满态射 $\mathcal{U} \to \mathcal{X}$，
其中 $\mathcal{U}$ 拟紧。
\end{enumerate}
\end{lemma}

\begin{proof}
我们将使用引理 \ref{stacks-properties-lemma-characterize-surjective}。
设 $\mathcal{U}$ 和 $\mathcal{U} \to \mathcal{X}$ 如 (5)。
由于 $|\mathcal{U}| \to |\mathcal{X}|$ 满且连续，
可知 $|\mathcal{X}|$ 拟紧。因此 (5) 蕴含 (1)。蕴含
(2) $\Rightarrow$ (3) $\Rightarrow$ (4) $\Rightarrow$ (5)
是直接的。设 (1) 成立，即 $\mathcal{X}$ 拟紧，也即
$|\mathcal{X}|$ 拟紧。选取概形 $U$ 以及满的光滑态射
$U \to \mathcal{X}$。由于 $|U| \to |\mathcal{X}|$ 是开映射，
存在拟紧开子集 $U' \subset U$，使得 $|U'| \to |X|$ 满
（并且仍然光滑）。选取有限仿射开覆盖
$U' = U_1 \cup \ldots \cup U_n$。则
$U_1 \amalg \ldots \amalg U_n \to \mathcal{X}$
是满的光滑态射，且其源是仿射概形（见《概形》引理
\ref{schemes-lemma-disjoint-union-affines}）。故 (2) 成立。
\end{proof}

\begin{lemma}
\label{stacks-properties-lemma-finite-disjoint-quasi-compact}
拟紧代数叠的有限不交并是拟紧代数叠。
\end{lemma}

\begin{proof}
这由相应的拓扑事实立即可得。
\end{proof}


\section{由概形性质定义的代数叠性质}
\label{stacks-properties-section-types-properties}

\noindent
概形的任一光滑局部性质都通过下述引理给出代数叠的相应性质。
注意，概形的光滑局部性质也是 \'etale 局部性质，
因为任一 \'etale 覆盖也是光滑覆盖。
因此，对概形的光滑局部性质 $P$，我们知道代数空间具有 $P$
的含义；见《空间的性质》第
\ref{spaces-properties-section-types-properties} 节。

\begin{lemma}
\label{stacks-properties-lemma-type-property}
设 $\mathcal{P}$ 是在光滑拓扑下局部的概形性质；见《下降》定义
\ref{descent-definition-property-local}。
设 $\mathcal{X}$ 是代数叠。下列条件等价：
\begin{enumerate}
\item 对某个概形 $U$ 和某个满的光滑态射 $U \to \mathcal{X}$，
概形 $U$ 具有性质 $\mathcal{P}$；
\item 对每个概形 $U$ 和每个光滑态射 $U \to \mathcal{X}$，
概形 $U$ 具有性质 $\mathcal{P}$；
\item 对某个代数空间 $U$ 和某个满的光滑态射 $U \to \mathcal{X}$，
代数空间 $U$ 具有性质 $\mathcal{P}$；
\item 对每个代数空间 $U$ 和每个光滑态射 $U \to \mathcal{X}$，
代数空间 $U$ 具有性质 $\mathcal{P}$。
\end{enumerate}
若 $\mathcal{X}$ 是概形 $U$，则这等价于 $\mathcal{P}(U)$。
若 $\mathcal{X}$ 是代数空间 $X$，则这等价于
$X$ 具有性质 $\mathcal{P}$。
\end{lemma}

\begin{proof}
设 $U \to \mathcal{X}$ 满且光滑，其中 $U$ 是代数空间。
设 $V \to \mathcal{X}$ 是光滑态射，其中 $V$ 是代数空间。
选取概形 $U'$、$V'$ 以及满的 \'etale 态射
$U' \to U$、$V' \to V$。最后，选取概形 $W$ 以及满的 \'etale 态射
$W \to V' \times_\mathcal{X} U'$。
由于 $W \to V'$ 和 $W \to U'$ 是代数空间的 \'etale 态射与
光滑态射的复合，它们是概形的光滑态射；见《空间的态射》引理
\ref{spaces-morphisms-lemma-etale-smooth} 和
\ref{spaces-morphisms-lemma-composition-smooth}。
此外，由于 $U' \to \mathcal{X}$ 满，$W \to V'$ 也满。因此
$$
\mathcal{P}(U) \Leftrightarrow
\mathcal{P}(U') \Rightarrow
\mathcal{P}(W) \Rightarrow
\mathcal{P}(V') \Leftrightarrow \mathcal{P}(V)
$$
其中两个等价来自代数空间的性质 $\mathcal{P}$ 之定义，
而两个蕴含来自《下降》定义
\ref{descent-definition-property-local}。
这证明了 (3) $\Rightarrow$ (4)。

\medskip\noindent
蕴含 (2) $\Rightarrow$ (1)、(1) $\Rightarrow$ (3)
以及 (4) $\Rightarrow$ (2) 都是直接的。
\end{proof}

\begin{definition}
\label{stacks-properties-definition-type-property}
设 $\mathcal{X}$ 是代数叠，$\mathcal{P}$ 是在光滑拓扑下局部的
概形性质。若引理 \ref{stacks-properties-lemma-type-property} 中任一等价条件成立，
则称 $\mathcal{X}$ {\it 具有性质 $\mathcal{P}$}。
\end{definition}

\begin{remark}
\label{stacks-properties-remark-list-properties-local-smooth-topology}
下面列出一些在光滑拓扑下局部的性质
（注意 fpqc、fppf 和 syntomic 拓扑都强于光滑拓扑）：
\begin{enumerate}
\item 局部 Noether；见《下降》引理
\ref{descent-lemma-Noetherian-local-fppf}；
\item Jacobson；见《下降》引理
\ref{descent-lemma-Jacobson-local-fppf}；
\item 局部 Noether 且满足 $(S_k)$；见《下降》引理
\ref{descent-lemma-Sk-local-syntomic}；
\item Cohen--Macaulay；见《下降》引理
\ref{descent-lemma-CM-local-syntomic}；
\item 既约；见《下降》引理
\ref{descent-lemma-reduced-local-smooth}；
\item 正规；见《下降》引理
\ref{descent-lemma-normal-local-smooth}；
\item 局部 Noether 且满足 $(R_k)$；见《下降》引理
\ref{descent-lemma-Rk-local-smooth}；
\item 正则；见《下降》引理
\ref{descent-lemma-regular-local-smooth}；
\item Nagata；见《下降》引理
\ref{descent-lemma-Nagata-local-smooth}。
\end{enumerate}
\end{remark}

\noindent
概形芽的任一光滑局部性质都给出代数叠的相应性质。
注意，芽的光滑局部性质也是 \'etale 局部性质。
因此，对概形芽的光滑局部性质 $P$，我们知道代数空间 $X$
在 $x \in |X|$ 处具有性质 $P$ 的含义；见《空间的性质》第
\ref{stacks-properties-section-types-properties} 节。

\begin{lemma}
\label{stacks-properties-lemma-local-source-target-at-point}
设 $\mathcal{X}$ 是代数叠，$x \in |\mathcal{X}|$ 是 $\mathcal{X}$ 的点。
设 $\mathcal{P}$ 是概形芽的光滑局部性质；见《下降》定义
\ref{descent-definition-local-at-point}。下列条件等价：
\begin{enumerate}
\item 对任一光滑态射 $U \to \mathcal{X}$，其中 $U$ 是概形，
以及任一满足 $a(u) = x$ 的 $u \in U$，都有 $\mathcal{P}(U, u)$；
\item 对某个光滑态射 $U \to \mathcal{X}$，其中 $U$ 是概形，
以及某个满足 $a(u) = x$ 的 $u \in U$，有 $\mathcal{P}(U, u)$；
\item 对任一光滑态射 $U \to \mathcal{X}$，其中 $U$ 是代数空间，
以及任一满足 $a(u) = x$ 的 $u \in |U|$，代数空间 $U$
在 $u$ 处具有性质 $\mathcal{P}$；
\item 对某个光滑态射 $U \to \mathcal{X}$，其中 $U$ 是代数空间，
以及某个满足 $a(u) = x$ 的 $u \in |U|$，代数空间 $U$
在 $u$ 处具有性质 $\mathcal{P}$。
\end{enumerate}
若 $\mathcal{X}$ 可表示，则这等价于
$\mathcal{P}(\mathcal{X}, x)$。若 $\mathcal{X}$ 是代数空间，
则这等价于 $\mathcal{X}$ 在 $x$ 处具有性质 $\mathcal{P}$。
\end{lemma}

\begin{proof}
取 (3) 中的 $a : U \to \mathcal{X}$ 和 $u \in |U|$。
再取另一个光滑态射 $b : V \to \mathcal{X}$，其中 $V$ 是代数空间，
并取满足 $b(v) = x$ 的 $v \in |V|$。
选取概形 $U'$、\'etale 态射 $U' \to U$ 以及映到 $u$ 的
$u' \in U'$。选取概形 $V'$、\'etale 态射 $V' \to V$
以及映到 $v$ 的 $v' \in V'$。
由引理 \ref{stacks-properties-lemma-points-cartesian}，存在点
$\overline{w} \in |V' \times_\mathcal{X} U'|$ 映到 $u'$ 和 $v'$。
选取概形 $W$ 以及满的 \'etale 态射
$W \to V' \times_\mathcal{X} U'$。可以选取映到 $\overline{w}$ 的
$w \in |W|$（见《空间的性质》引理
\ref{spaces-properties-lemma-characterize-surjective}）。
由于 $W \to V'$ 和 $W \to U'$ 是代数空间的 \'etale 态射与
光滑态射的复合，它们是概形的光滑态射；见《空间的态射》引理
\ref{spaces-morphisms-lemma-etale-smooth} 和
\ref{spaces-morphisms-lemma-composition-smooth}。因此
$$
\mathcal{P}(U, u)
\Leftrightarrow
\mathcal{P}(U', u')
\Leftrightarrow
\mathcal{P}(W, w)
\Leftrightarrow
\mathcal{P}(V', v')
\Leftrightarrow
\mathcal{P}(V, v)
$$
外侧两个等价来自《空间的性质》定义
\ref{spaces-properties-definition-property-at-point}，
另外两个等价来自概形芽的光滑局部性质之含义。
这证明了 (4) $\Rightarrow$ (3)。

\medskip\noindent
蕴含 (1) $\Rightarrow$ (2)、(2) $\Rightarrow$ (4)
以及 (3) $\Rightarrow$ (1) 都是直接的。
\end{proof}

\begin{definition}
\label{stacks-properties-definition-property-at-point}
设 $\mathcal{P}$ 是概形芽的光滑局部性质，$\mathcal{X}$ 是代数叠，
且 $x \in |\mathcal{X}|$。若引理
\ref{stacks-properties-lemma-local-source-target-at-point} 的任一等价条件成立，
则称 $\mathcal{X}$ {\it 在 $x$ 处具有性质 $\mathcal{P}$}。
\end{definition}










\section{代数叠的单态射}
\label{stacks-properties-section-monomorphisms}

\noindent
我们如下定义代数叠的单态射。引理 \ref{stacks-properties-lemma-monomorphism}
将说明，该定义与相应的 $2$-范畴论概念相容。

\begin{definition}
\label{stacks-properties-definition-monomorphism}
设 $f : \mathcal{X} \to \mathcal{Y}$ 是代数叠态射。
若它可由代数空间表示，并且是第
\ref{stacks-properties-section-properties-morphisms} 节意义下的单态射，
则称 $f$ 是{\it 单态射}。
\end{definition}

\noindent
先给出一些基本引理。

\begin{lemma}
\label{stacks-properties-lemma-base-change-monomorphism}
设 $\mathcal{X} \to \mathcal{Y}$ 是代数叠态射，
且 $\mathcal{Z} \to \mathcal{Y}$ 是单态射。
则 $\mathcal{Z} \times_\mathcal{Y} \mathcal{X} \to \mathcal{X}$
是单态射。
\end{lemma}

\begin{proof}
这由第 \ref{stacks-properties-section-properties-morphisms} 节的一般讨论推出。
\end{proof}

\begin{lemma}
\label{stacks-properties-lemma-composition-monomorphism}
代数叠单态射的复合是单态射。
\end{lemma}

\begin{proof}
这由第 \ref{stacks-properties-section-properties-morphisms} 节的一般讨论以及
《空间的态射》引理
\ref{spaces-morphisms-lemma-composition-monomorphism} 推出。
\end{proof}

\begin{lemma}
\label{stacks-properties-lemma-monomorphism}
设 $f : \mathcal{X} \to \mathcal{Y}$ 是代数叠态射。下列条件等价：
\begin{enumerate}
\item $f$ 是单态射；
\item $f$ 全忠实；
\item 对角态射
$\Delta_f : \mathcal{X} \to \mathcal{X} \times_\mathcal{Y} \mathcal{X}$
是等价；
\item 存在代数空间 $W$ 以及满、平坦且局部有限表示的态射
$W \to \mathcal{Y}$，使得
$V = \mathcal{X} \times_\mathcal{Y} W$ 是代数空间，并且
$V \to W$ 是代数空间的单态射。
\end{enumerate}
\end{lemma}

\begin{proof}
(1) 与 (4) 的等价性来自第 \ref{stacks-properties-section-properties-morphisms} 节的
一般讨论，特别是引理 \ref{stacks-properties-lemma-check-representable-covering}
和 \ref{stacks-properties-lemma-check-property-covering}。

\medskip\noindent
(2) 与 (3) 的等价性是《范畴》引理
\ref{categories-lemma-fully-faithful-diagonal-equivalence}。

\medskip\noindent
设等价条件 (2) 和 (3) 成立。由《代数叠》引理
\ref{algebraic-lemma-characterize-representable-by-algebraic-spaces}，
$f$ 可由代数空间表示。此外，$2$-Yoneda 引理与全忠实性共同蕴含：
对每个概形 $T$，函子
$$
\Mor(T, \mathcal{X})
\longrightarrow
\Mor(T, \mathcal{Y})
$$
全忠实。因此，给定态射 $y : T \to \mathcal{Y}$，
在唯一 $2$-同构的意义下，至多存在一个满足
$y \cong f \circ x$ 的态射 $x : T \to \mathcal{X}$。
特别地，给定概形态射 $h : T' \to T$，至多存在一个 $h$ 的提升
$\tilde h : T' \to T \times_\mathcal{Y} \mathcal{X}$。
故 $T \times_\mathcal{Y} \mathcal{X} \to T$ 是代数空间的单态射，
这证明了 (1)。

\medskip\noindent
最后，设 (1) 成立。则对任一概形 $T$ 和态射
$y : T \to \mathcal{Y}$，纤维积
$T \times_\mathcal{Y} \mathcal{X}$ 是代数空间，并且
$T \times_\mathcal{Y} \mathcal{X} \to T$ 是单态射。
因此，在唯一同构的意义下，恰好存在一个偶 $(x, \alpha)$，
其中 $x : T \to \mathcal{X}$ 是态射，而
$\alpha : f \circ x \to y$ 是 $2$-态射。
应用 $2$-Yoneda 引理，这恰好说明 $f$ 全忠实，即 (2) 成立。
\end{proof}

\begin{lemma}
\label{stacks-properties-lemma-monomorphism-injective-points}
\begin{slogan}
叠的单态射在点上是单射。
\end{slogan}
代数叠的单态射诱导点集之间的单射。
\end{lemma}

\begin{proof}
设 $f : \mathcal{X} \to \mathcal{Y}$ 是代数叠的单态射。
设 $x_i : \Spec(K_i) \to \mathcal{X}$ 是态射，并且
$f \circ x_1$ 与 $f \circ x_2$ 定义 $|\mathcal{Y}|$ 的同一元素。
应用定义，可找到公共扩张 $\Omega$、相应态射
$c_i : \Spec(\Omega) \to \Spec(K_i)$，以及 $2$-同构
$\beta : f \circ x_1 \circ c_1 \to f \circ x_1 \circ c_2$。
由于 $f$ 全忠实（见引理 \ref{stacks-properties-lemma-monomorphism}），
可以将 $\beta$ 提升为同构
$\alpha : x_1 \circ c_1 \to x_1 \circ c_2$。
故 $x_1$ 与 $x_2$ 定义 $|\mathcal{X}|$ 的同一点，正如所需。
\end{proof}

\begin{lemma}
\label{stacks-properties-lemma-monomorphism-diagonal}
设 $\mathcal{X} \to \mathcal{X}' \to \mathcal{Y}$ 是代数叠态射。
若 $\mathcal{X} \to \mathcal{X}'$ 是单态射，则典范图表
$$
\xymatrix{
\mathcal{X} \ar[r] \ar[d] &
\mathcal{X} \times_\mathcal{Y} \mathcal{X} \ar[d] \\
\mathcal{X}' \ar[r] &
\mathcal{X}' \times_\mathcal{Y} \mathcal{X}'
}
$$
是纤维积方块。
\end{lemma}

\begin{proof}
由引理 \ref{stacks-properties-lemma-monomorphism}，有
$\mathcal{X} = \mathcal{X} \times_{\mathcal{X}'} \mathcal{X}$。
因此，应用《范畴》引理
\ref{categories-lemma-fibre-product-after-map} 即得结论。
\end{proof}







\section{代数叠的浸入}
\label{stacks-properties-section-immersions}

\noindent
代数叠的浸入定义如下。

\begin{definition}
\label{stacks-properties-definition-immersion}
浸入。
\begin{enumerate}
\item 若一个代数叠态射可表示，并且是第
\ref{stacks-properties-section-properties-morphisms} 节意义下的开浸入，
则称它为{\it 开浸入}。
\item 若一个代数叠态射可表示，并且是第
\ref{stacks-properties-section-properties-morphisms} 节意义下的闭浸入，
则称它为{\it 闭浸入}。
\item 若一个代数叠态射可表示，并且是第
\ref{stacks-properties-section-properties-morphisms} 节意义下的浸入，
则称它为{\it 浸入}。
\end{enumerate}
\end{definition}

\noindent
对我们而言，这并不是理解浸入的最方便方式。更方便的看法是：
浸入是可由代数空间表示、并且是第
\ref{stacks-properties-section-properties-morphisms} 节意义下浸入的代数叠态射。
闭浸入和开浸入亦同。由于这显然等价于刚定义的概念，
以后将不加说明地使用这一刻画。
下面证明关于该概念的几个简单引理。

\begin{lemma}
\label{stacks-properties-lemma-base-change-immersion}
设 $\mathcal{X} \to \mathcal{Y}$ 是代数叠态射，
且 $\mathcal{Z} \to \mathcal{Y}$ 是（闭浸入，分别地，开浸入）。
则 $\mathcal{Z} \times_\mathcal{Y} \mathcal{X} \to \mathcal{X}$
是（闭浸入，分别地，开浸入）。
\end{lemma}

\begin{proof}
这由第 \ref{stacks-properties-section-properties-morphisms} 节的一般讨论推出。
\end{proof}

\begin{lemma}
\label{stacks-properties-lemma-composition-immersion}
代数叠浸入的复合是浸入。闭浸入和开浸入也分别如此。
\end{lemma}

\begin{proof}
这由第 \ref{stacks-properties-section-properties-morphisms} 节的一般讨论以及
《空间》引理 \ref{spaces-lemma-composition-immersions} 推出。
\end{proof}

\begin{lemma}
\label{stacks-properties-lemma-check-immersion-covering}
设 $f : \mathcal{X} \to \mathcal{Y}$ 是代数叠态射。
设 $W$ 是代数空间，且 $W \to \mathcal{Y}$
是满、平坦且局部有限表示的态射。下列条件等价：
\begin{enumerate}
\item $f$ 是（开浸入，分别地，闭浸入）；
\item $V = W \times_\mathcal{Y} \mathcal{X}$ 是代数空间，
且 $V \to W$ 是（开浸入，分别地，闭浸入）。
\end{enumerate}
\end{lemma}

\begin{proof}
这由第 \ref{stacks-properties-section-properties-morphisms} 节的一般讨论推出，
特别是引理 \ref{stacks-properties-lemma-check-representable-covering}
和 \ref{stacks-properties-lemma-check-property-covering}。
\end{proof}

\begin{lemma}
\label{stacks-properties-lemma-immersion-monomorphism}
浸入是单态射。
\end{lemma}

\begin{proof}
见《空间的态射》引理
\ref{spaces-morphisms-lemma-immersions-monomorphisms}。
\end{proof}

\begin{lemma}
\label{stacks-properties-lemma-subspace-induced-topology}
若 $f : \mathcal{X} \to \mathcal{Y}$ 是浸入，则
$|f| : |\mathcal{X}| \to |\mathcal{Y}|$ 是到某个局部闭子集上的同胚。
若 $f$ 是闭浸入，分别地，开浸入，则 $|f|$ 是闭映射，分别地，开映射。
\end{lemma}

\begin{proof}
略。
\end{proof}

\noindent
下面两个引理说明如何利用表示理解浸入。

\begin{lemma}
\label{stacks-properties-lemma-immersion-into-presentation}
设 $(U, R, s, t, c)$ 是代数空间中的光滑群胚，
且 $i : \mathcal{Z} \to [U/R]$ 是浸入。
则存在 $R$-不变的局部闭子空间 $Z \subset U$ 和表示
$[Z/R_Z] \to \mathcal{Z}$，其中 $R_Z$ 是 $R$ 在 $Z$ 上的限制，
使得
$$
\xymatrix{
[Z/R_Z] \ar[dr] \ar[rr] & & \mathcal{Z} \ar[ld]^i \\
& [U/R]
}
$$
为 $2$-交换图表。若 $i$ 是闭浸入（分别地，开浸入），
则 $Z$ 是 $U$ 的闭子空间（分别地，开子空间）。
\end{lemma}

\begin{proof}
由引理 \ref{stacks-properties-lemma-representable-in-terms-presentations}，
得到交换图表
$$
\xymatrix{
[U'/R'] \ar[dr] \ar[rr] & & \mathcal{Z} \ar[ld] \\
& [U/R]
}
$$
其中 $U' = \mathcal{Z} \times_{[U/R]} U$ 且
$R' = \mathcal{Z} \times_{[U/R]} R$。
由于 $\mathcal{Z} \to [U/R]$ 是浸入，$U' \to U$
是代数空间的浸入。令 $Z \subset U$ 为如下局部闭子空间：
$U' \to U$ 经 $Z$ 分解，并诱导同构 $U' \to Z$。
由 $R'$ 的构造显然有
$R' = U' \times_{U, t} R = R \times_{s, U} U'$.
这蕴含 $Z \cong U'$ 是 $R$-不变的，并且 $R' \to R$ 的像将
$R'$ 与 $R$ 在 $Z$ 上的限制
$R_Z = s^{-1}(Z) = t^{-1}(Z)$ 等同。故引理成立。
\end{proof}

\begin{lemma}
\label{stacks-properties-lemma-immersion-presentation}
设 $(U, R, s, t, c)$ 是代数空间中的光滑群胚。
令 $\mathcal{X} = [U/R]$ 为相应的代数叠；见《代数叠》定理
\ref{algebraic-theorem-smooth-groupoid-gives-algebraic-stack}。
设 $Z \subset U$ 是 $R$-不变的局部闭子空间。则
$$
[Z/R_Z] \longrightarrow [U/R]
$$
是代数叠的浸入，其中 $R_Z$ 是 $R$ 在 $Z$ 上的限制。
若 $Z \subset U$ 是开子空间（分别地，闭子空间），则该态射是
代数叠的开浸入（分别地，闭浸入）。
\end{lemma}

\begin{proof}
回忆，由《空间中的群胚》定义
\ref{spaces-groupoids-definition-invariant-open}
（另见该定义之后的讨论），作为 $R$ 的局部闭子空间，有
$R_Z = s^{-1}(Z) = t^{-1}(Z)$。因此，两个态射 $R_Z \to Z$
作为 $s$ 和 $t$ 的基变换都是光滑的。故
$(Z, R_Z, s|_{R_Z}, t|_{R_Z}, c|_{R_Z \times_{s, Z, t} R_Z})$ 是
代数空间中的光滑群胚，并且可见 $[Z/R_Z]$ 是代数叠；
见《代数叠》定理
\ref{algebraic-theorem-smooth-groupoid-gives-algebraic-stack}。
《空间中的群胚》引理
\ref{spaces-groupoids-lemma-criterion-fibre-product}
的假设全部满足，因而得到 $2$-纤维方块
$$
\xymatrix{
Z \ar[d] \ar[r] & [Z/R_Z] \ar[d] \\
U \ar[r] & [U/R]
}
$$
由此和引理 \ref{stacks-properties-lemma-check-representable-covering} 可知，
$[Z/R_Z] \to [U/R]$ 可由代数空间表示；进而由引理
\ref{stacks-properties-lemma-check-property-covering} 可知，右侧竖直箭头是浸入
（分别地，闭浸入；分别地，开浸入），当且仅当左侧竖直箭头如此。
\end{proof}

\noindent
可以如下定义开、闭及局部闭子叠。

\begin{definition}
\label{stacks-properties-definition-substacks}
设 $\mathcal{X}$ 是代数叠。
\begin{enumerate}
\item $\mathcal{X}$ 的一个{\it 开子叠}，是严格满子范畴
$\mathcal{X}' \subset \mathcal{X}$，使得 $\mathcal{X}'$ 是代数叠，
且 $\mathcal{X}' \to \mathcal{X}$ 是开浸入。
\item $\mathcal{X}$ 的一个{\it 闭子叠}，是严格满子范畴
$\mathcal{X}' \subset \mathcal{X}$，使得 $\mathcal{X}'$ 是代数叠，
且 $\mathcal{X}' \to \mathcal{X}$ 是闭浸入。
\item $\mathcal{X}$ 的一个{\it 局部闭子叠}，是严格满子范畴
$\mathcal{X}' \subset \mathcal{X}$，使得 $\mathcal{X}'$ 是代数叠，
且 $\mathcal{X}' \to \mathcal{X}$ 是浸入。
\end{enumerate}
\end{definition}

\noindent
使用该定义时应当谨慎。事实上，若
$f : \mathcal{X} \to \mathcal{Y}$ 是代数叠的等价，
且 $\mathcal{X}' \subset \mathcal{X}$ 是开子叠，
则子范畴 $f(\mathcal{X}')$ 未必是 $\mathcal{Y}$ 的开子叠。
问题在于它可能不是{\it 严格}满子范畴；但这也是唯一的问题。
下面给出形式化陈述。

\begin{lemma}
\label{stacks-properties-lemma-substack-image}
对任一浸入 $i : \mathcal{Z} \to \mathcal{X}$，存在唯一的局部闭子叠
$\mathcal{X}' \subset \mathcal{X}$，使得 $i$ 分解为等价
$i' : \mathcal{Z} \to \mathcal{X}'$ 与包含态射
$\mathcal{X}' \to \mathcal{X}$ 的复合。
若 $i$ 是闭浸入（分别地，开浸入），则 $\mathcal{X}'$
是 $\mathcal{X}$ 的闭子叠（分别地，开子叠）。
\end{lemma}

\begin{proof}
略。
\end{proof}

\begin{lemma}
\label{stacks-properties-lemma-substacks-presentation}
设 $[U/R] \to \mathcal{X}$ 是代数叠的一个表示。存在典范双射
$$
\mathcal{X}\text{ 的局部闭子叠 }\mathcal{Z}
\longrightarrow
U\text{ 的 }R\text{-不变局部闭子空间 }Z
$$
将 $\mathcal{Z}$ 映为 $U \times_\mathcal{X} \mathcal{Z}$。
此外，代数叠态射 $f : \mathcal{Y} \to \mathcal{X}$
经 $\mathcal{Z}$ 分解，当且仅当
$\mathcal{Y} \times_\mathcal{X} U \to U$ 经 $Z$ 分解。
闭子叠和开子叠也有相同结论。
\end{lemma}

\begin{proof}
由引理 \ref{stacks-properties-lemma-immersion-into-presentation} 和
\ref{stacks-properties-lemma-immersion-presentation}，该映射是双射。
若 $\mathcal{Y} \to \mathcal{X}$ 经 $\mathcal{Z}$ 分解，
则其基变换 $\mathcal{Y} \times_\mathcal{X} U \to U$
当然经 $Z$ 分解。反之，设态射 $\mathcal{Y} \to \mathcal{X}$
满足 $\mathcal{Y} \times_\mathcal{X} U \to U$ 经 $Z$ 分解。
我们将证明：对每个概形 $T$ 和每个态射 $T \to \mathcal{Y}$，
若该态射由 $\mathcal{Y}$ 在 $T$ 上纤维范畴中的对象 $y$ 给出，
则 $y$ 实际上属于 $\mathcal{Z}$ 在 $T$ 上的纤维范畴。
事实上，纤维积 $T \times_\mathcal{X} U$ 是代数空间，
且 $T \times_\mathcal{X} U \to T$ 是满的光滑态射。
故存在 fppf 覆盖 $\{T_i \to T\}$，使得对所有 $i$，
$T_i \to T$ 都经 $T \times_\mathcal{X} U \to T$ 分解。
于是 $T_i \to \mathcal{X}$ 经
$\mathcal{Y} \times_\mathcal{X} U$，进而经 $Z \subset U$ 分解。
因此 $y|_{T_i}$ 是 $\mathcal{Z}$ 的对象
（因为 $Z$ 是 $U$ 与 $\mathcal{Z}$ 在 $\mathcal{X}$ 上的纤维积）。
由于 $\mathcal{Z}$ 是严格满子叠，可知 $y$ 是
$\mathcal{Z}$ 的对象，正如所需。
\end{proof}

\begin{lemma}
\label{stacks-properties-lemma-open-substacks}
设 $\mathcal{X}$ 是代数叠。规则
$\mathcal{U} \mapsto |\mathcal{U}|$ 在 $\mathcal{X}$ 的开子叠与
$|\mathcal{X}|$ 的开子集之间定义一个保持包含关系的双射。
\end{lemma}

\begin{proof}
选取表示 $[U/R] \to \mathcal{X}$；见《代数叠》引理
\ref{algebraic-lemma-stack-presentation}。
由引理 \ref{stacks-properties-lemma-substacks-presentation}，
开子叠对应于 $U$ 的 $R$-不变开子概形。
另一方面，引理 \ref{stacks-properties-lemma-points-presentation} 和
\ref{stacks-properties-lemma-topology-points} 保证后者与
$|\mathcal{X}|$ 的开子集双射对应。
\end{proof}

\begin{lemma}
\label{stacks-properties-lemma-open-image-substack}
设 $\mathcal X$ 是代数叠，$U$ 是代数空间，并且
$U \to \mathcal X$ 是满的光滑态射。对任一开浸入
$V \hookrightarrow U$，存在代数叠 $\mathcal Y$、开浸入
$\mathcal Y \to \mathcal X$，以及满的光滑态射
$V \to \mathcal Y$。
\end{lemma}

\begin{proof}
定义群胚纤维化范畴 $\mathcal Y$：对
$(\Sch/S)_{fppf}$ 的对象 $T$，令其上纤维范畴
$\mathcal{Y}_T$ 为 $\mathcal{X}_T$ 的满子范畴，
其对象是所有满足投影态射
$V \times_{\mathcal X, y} T \to T$ 为满射的
$y \in \Ob(\mathcal{X}_T)$。
现对任一态射 $x : T \to \mathcal X$，$2$-纤维积
$T \times_{x, \mathcal X} \mathcal Y$ 在 $T'$ 上的纤维范畴
由三元组 $(f : T' \to T, y \in \mathcal{X}_{T'}, f^*x \simeq y)$
组成，并要求 $V \times_{\mathcal X, y} T' \to T'$ 为满射。
注意，$T \times_{x, \mathcal X} \mathcal Y$ 在集合胚中纤维化，
因为 $\mathcal Y \to \mathcal X$ 忠实
（见《叠》引理
\ref{stacks-lemma-2-fibre-product-gives-stack-in-setoids}）。
同构 $f^*x \simeq y$ 给出图表
$$
\xymatrix{
V \times_{\mathcal X, y} T' \ar[d] \ar[r] &
V \times_{\mathcal X, x} T \ar[r] \ar[d] &
V \ar[d] \\
T' \ar[r]^f &
T \ar[r]^x &
\mathcal X
}
$$
其中两个方块都是笛卡儿的。态射
$V \times_{\mathcal X, x} T \to T$ 作为基变换是光滑的，因而是开的。
令 $T_0 \subset T$ 为其像。由笛卡儿方块可得，
$V \times_{\mathcal X, y} T' \to T'$ 为满射，当且仅当
$f$ 落在 $T_0$ 中。因此，$T \times_{x, \mathcal X} \mathcal Y$
可由 $T_0$ 表示，所以包含 $\mathcal Y \to \mathcal X$ 是开浸入。
由《代数叠》引理
\ref{algebraic-lemma-open-fibred-category-is-algebraic}，
可知 $\mathcal{Y}$ 是代数叠。
最后，将态射 $V \to \mathcal X$ 记作 $g$，则
$V \times_{\mathcal X} V \to V$ 是满射（对角态射给出一个截面）。
故 $g$ 属于 $\mathcal{Y}_V \to \mathcal{X}_V$ 的像；亦即，
得到态射 $g' : V \to \mathcal{Y}$，它嵌入交换图表
$$
\xymatrix{
V \ar[r] \ar[d]^{g'} & U \ar[d] \\
\mathcal{Y} \ar[r] & \mathcal{X}
}
$$
由于 $V \times_{g, \mathcal X} \mathcal Y \to V$ 是单态射，
而 $(1, g')$ 给出一个截面，它实际上是同构。
因此，$g' : V \to \mathcal Y$ 是光滑态射，因为它是光滑态射
$g : V \to \mathcal{X}$ 的基变换。由 $\mathcal{Y}$ 的构造，
它还是满射；引理得证。
\end{proof}

\begin{lemma}
\label{stacks-properties-lemma-union-open-substacks}
设 $\mathcal X$ 是代数叠，且
$\mathcal{X}_i \subset \mathcal X$ 是由 $i \in I$ 标号的一族开子叠。
则存在一个开子叠，记作
$\bigcup_{i\in I} \mathcal{X}_i \subset \mathcal X$，
使得 $\mathcal{X}_i$ 是覆盖它的开子叠。
\end{lemma}

\begin{proof}
定义纤维化子范畴
$\mathcal{X}' = \bigcup_{i \in I} \mathcal{X}_i$：
对 $(\Sch/S)_{fppf}$ 的对象 $T$，令其上纤维范畴为
$\mathcal{X}_T$ 的满子范畴，其对象是所有满足态射
$\coprod_{i \in I} (\mathcal{X}_i \times_{\mathcal X} T) \to T$
为满射的 $x \in \Ob(\mathcal{X}_T)$。
令 $x_i \in \Ob((\mathcal{X}_i)_T)$。则 $(x_i, 1)$ 给出
$\mathcal{X}_i \times_{\mathcal X} T \to T$ 的截面，故有一个同构。
因此 $\mathcal{X}_i \subset \mathcal{X}'$ 是满子范畴。
现取 $x \in \Ob(\mathcal{X}_T)$。则
$\mathcal{X}_i \times_{\mathcal X} T$ 可由开子概形
$T_i \subset T$ 表示。$2$-纤维积
$\mathcal{X}' \times_{\mathcal X} T$ 在 $T'$ 上的纤维
由 $(y \in \mathcal{X}_{T'}, f : T' \to T, f^*x \simeq y)$ 组成，
并要求 $\coprod (\mathcal{X}_i \times_{\mathcal X, y} T') \to T'$
为满射。同构 $f^*x \simeq y$ 诱导同构
$\mathcal{X}_i \times_{\mathcal X, y} T' \simeq T_i \times_T T'$。
于是 $T_i \times_T T'$ 覆盖 $T'$，当且仅当
$f$ 落在 $\bigcup T_i$ 中。因此有图表
$$
\xymatrix{
T_i \ar[r] \ar[d] &
\bigcup T_i \ar[r] \ar[d] &
T \ar[d] \\
\mathcal{X}_i \ar[r] &
\mathcal{X}' \ar[r] &
\mathcal{X}
}
$$
其中两个方块均为笛卡儿的。由《代数叠》引理
\ref{algebraic-lemma-open-fibred-category-is-algebraic}，
可知 $\mathcal{X'} \subset \mathcal{X}$ 是代数的，且为开子叠。
由上述笛卡儿方块还显然可知，态射
$\coprod_{i \in I} \mathcal{X}_i \to \mathcal{X}'$，从而完成引理的证明。
\end{proof}

\begin{lemma}
\label{stacks-properties-lemma-quasi-compact-finite-subcover}
设 $\mathcal X$ 是代数叠，且 $\mathcal X' \subset \mathcal X$
是拟紧开子叠。设有由 $i \in I$ 标号的一族开子叠
$\mathcal{X}_i \subset \mathcal X$，满足
$\mathcal{X}' \subset \bigcup_{i \in I} \mathcal{X}_i$，
其中并按引理 \ref{stacks-properties-lemma-union-open-substacks} 定义。
则存在有限子集 $I' \subset I$，使得
$\mathcal{X}' \subset \bigcup_{i \in I'} \mathcal{X}_i$。
\end{lemma}

\begin{proof}
由于 $\mathcal X$ 是代数的，存在概形 $U$ 和满的光滑态射
$U \to \mathcal X$。令 $U_i \subset U$ 为表示
$\mathcal{X}_i \times_{\mathcal X} U$ 的开子概形，
并令 $U' \subset U$ 为表示
$\mathcal{X}' \times_{\mathcal X} U$ 的开子概形。
由假设，$U'\subset \bigcup_{i\in I} U_i$。
由引理 \ref{stacks-properties-lemma-quasi-compact-stack} 的证明，存在拟紧开子集
$V \subset U'$，使得 $V \to \mathcal{X}'$ 是满的光滑态射。
因此存在有限子集 $I' \subset I$，使得
$V \subset \bigcup_{i \in I'} U_i$。
我们断言 $\mathcal{X}' \subset \bigcup_{i \in I'} \mathcal{X}_i$。
对 $T \in \Ob((\Sch/S)_{fppf})$，取
$x \in \Ob(\mathcal{X}'_T)$。
由于 $\mathcal{X}' \to \mathcal{X}$ 是单态射，有笛卡儿方块
$$
\xymatrix{
V \times_\mathcal{X} T \ar[r] \ar[d] &
T \ar[d]^x \ar@{=}[r] &
T \ar[d]^x \\
V \ar[r] &
\mathcal{X}' \ar[r] &
\mathcal X
}
$$
由基变换，$V \times_{\mathcal X} T \to T$ 是满射。因此
$\bigcup_{i \in I'} U_i \times_{\mathcal X} T \to T$ 也是满射。
令 $T_i \subset T$ 为表示
$\mathcal{X}_i \times_{\mathcal X} T$ 的开子概形。
由形式论证，得到笛卡儿方块
$$
\xymatrix{
U_i \times_{\mathcal{X}_i} T_i \ar[r] \ar[d] &
U \times_{\mathcal X} T \ar[d] \\
T_i \ar[r] & T
}
$$
其中竖直箭头由基变换都是满射。由于
$U_i \times_{\mathcal{X}_i} T_i \simeq U_i \times_{\mathcal X} T$，
可得 $\bigcup_{i \in I'} T_i = T$。因此，由并的定义，
$x$ 是 $(\bigcup_{i\in I'} \mathcal{X}_i)_T$ 的对象。
注意，包含
$\mathcal{X}' \subset \bigcup_{i \in I'} \mathcal{X}_i$
自动是开子叠。
\end{proof}

\begin{lemma}
\label{stacks-properties-lemma-zariski-open-cover-stack-is-space}
设 $\mathcal X$ 是代数叠，且 $\mathcal{X}_i$（$i \in I$）
是一族 $\mathcal{X}$ 的开子叠。假设
\begin{enumerate}
\item $\mathcal{X} = \bigcup_{i \in I} \mathcal{X}_i$；
\item 每个 $\mathcal{X}_i$ 都是代数空间。
\end{enumerate}
则 $\mathcal{X}$ 是代数空间。
\end{lemma}

\begin{proof}
将《叠》引理 \ref{stacks-lemma-stack-in-setoids-descent}
应用于态射 $\coprod_{i \in I} \mathcal{X}_i \to \mathcal{X}$
和态射 $\text{id} : \mathcal{X} \to \mathcal{X}$，
可见 $\mathcal{X}$ 是集合胚中的叠。
因此 $\mathcal{X}$ 是代数空间；见《代数叠》命题
\ref{algebraic-proposition-algebraic-stack-no-automorphisms}。
\end{proof}

\begin{lemma}
\label{stacks-properties-lemma-zariski-open-cover-stack-is-scheme}
设 $\mathcal X$ 是代数叠，且 $\mathcal{X}_i$（$i \in I$）
是一族 $\mathcal{X}$ 的开子叠。假设
\begin{enumerate}
\item $\mathcal{X} = \bigcup_{i \in I} \mathcal{X}_i$；
\item 每个 $\mathcal{X}_i$ 都是概形。
\end{enumerate}
则 $\mathcal{X}$ 是概形。
\end{lemma}

\begin{proof}
由引理 \ref{stacks-properties-lemma-zariski-open-cover-stack-is-space}，
$\mathcal{X}$ 是代数空间。每个代数空间都有一个为概形的最大开子空间；
见《空间的性质》引理
\ref{spaces-properties-lemma-subscheme}。故 $\mathcal{X}$ 是概形。
\end{proof}


\noindent
下一个引理类似于《群胚进阶》引理
\ref{more-groupoids-lemma-local-source}。

\begin{lemma}
\label{stacks-properties-lemma-local-source}
设 $\mathcal{P}, \mathcal{Q}, \mathcal{R}$ 是代数空间态射的性质。假设
\begin{enumerate}
\item $\mathcal{P}, \mathcal{Q}, \mathcal{R}$ 在靶上 fppf 局部，
且在任意基变换下保持；
\item $\text{smooth} \Rightarrow \mathcal{R}$,
\item 对任一具有性质 $\mathcal{Q}$ 的态射 $f : X \to Y$，
存在最大开子空间 $W(\mathcal{P}, f) \subset X$，使得
$f|_{W(\mathcal{P}, f)}$ 具有性质 $\mathcal{P}$；
\item 对任一具有性质 $\mathcal{Q}$ 的态射 $f : X \to Y$，
以及任一具有性质 $\mathcal{R}$ 的态射 $Y' \to Y$，都有
$Y' \times_Y W(\mathcal{P}, f) = W(\mathcal{P}, f')$，
其中 $f' : X_{Y'} \to Y'$ 是 $f$ 的基变换。
\end{enumerate}
设 $f : \mathcal{X} \to \mathcal{Y}$ 是可由代数空间表示的代数叠态射，
并假设 $f$ 具有性质 $\mathcal{Q}$。则
\begin{enumerate}
\item[(A)] 存在最大开子叠 $\mathcal{X}' \subset \mathcal{X}$，
使得 $f|_{\mathcal{X}'}$ 具有性质 $\mathcal{P}$；
\item[(B)] 若 $\mathcal{Z} \to \mathcal{Y}$
是可由代数空间表示且具有性质 $\mathcal{R}$ 的代数叠态射，
则 $\mathcal{Z} \times_\mathcal{Y} \mathcal{X}'$
是 $\mathcal{Z} \times_\mathcal{Y} \mathcal{X}$ 的最大开子叠，
在它上面，基变换 $\text{id}_\mathcal{Z} \times f$
具有性质 $\mathcal{P}$。
\end{enumerate}
\end{lemma}

\begin{proof}
选取概形 $V$ 以及满的光滑态射 $V \to \mathcal{Y}$。
令 $U = V \times_\mathcal{Y} \mathcal{X}$，并令
$f' : U \to V$ 为 $f$ 的基变换。代数空间态射
$f' : U \to V$ 具有性质 $\mathcal{Q}$。
因此，由假设 (3) 得到开子空间 $W(\mathcal{P}, f') \subset U$。
注意
$U \times_\mathcal{X} U =
(V \times_\mathcal{Y} V) \times_\mathcal{Y} \mathcal{X}$
，故态射
$f'' : U \times_\mathcal{X} U \to V \times_\mathcal{Y} V$
是 $f$ 沿任一投影 $V \times_\mathcal{Y} V \to V$ 的基变换。
由 $V$ 的选择，这些投影光滑，故由 (2) 具有性质 $\mathcal{R}$。
因此，由 (4)，$W(\mathcal{P}, f')$ 在两个投影
$\text{pr}_i : U \times_\mathcal{X} U \to U$ 下的逆像相同。
换言之，$W(\mathcal{P}, f')$ 是 $U$ 的 $R$-不变子空间
（其中 $R = U \times_\mathcal{X} U$）。
令 $\mathcal{X}'$ 为通过引理
\ref{stacks-properties-lemma-immersion-into-presentation} 与
$W(\mathcal{P}, f)$ 对应的 $\mathcal{X}$ 的开子叠。
按构造，
$W(\mathcal{P}, f') = \mathcal{X}' \times_\mathcal{Y} V$，
故由引理 \ref{stacks-properties-lemma-check-property-covering}，
$f|_{\mathcal{X}'}$ 具有性质 $\mathcal{P}$。
此外，$\mathcal{X}'$ 是使得 $f|_{\mathcal{X}'}$ 具有
$\mathcal{P}$ 的最大开子叠，因为 $W(\mathcal{P}, f)$
也具有同样的极大性。这证明了 (A)。

\medskip\noindent
最后，若 $\mathcal{Z} \to \mathcal{Y}$ 是可由代数空间表示且
具有性质 $\mathcal{R}$ 的代数叠态射，则令
$T = V \times_\mathcal{Y} \mathcal{Z}$。
可见 $T \to V$ 是具有性质 $\mathcal{R}$ 的代数空间态射。
令 $f'_T : T \times_V U \to T$ 为 $f'$ 的基变换。
再次由 (4)，$W(\mathcal{P}, f'_T)$ 是
$W(\mathcal{P}, f)$ 在 $T \times_V U$ 中的逆像。
这蕴含 (B)；略去一些细节。
\end{proof}

\begin{remark}
\label{stacks-properties-remark-local-source-warning}
警告：使用引理 \ref{stacks-properties-lemma-local-source} 时应当谨慎。例如，它适用于
$\mathcal{P}=$``平坦''、$\mathcal{Q}=$``空性质''、
$\mathcal{R}=$``平坦且局部有限表示''的情形。
但给定代数空间态射 $f : X \to Y$，使得 $f|_W$ 平坦的最大开子空间
$W \subset X$ {\it 并不是} $f$ 平坦的点集！
\end{remark}

\begin{remark}
\label{stacks-properties-remark-local-source-apply}
尽管有注 \ref{stacks-properties-remark-local-source-warning} 中的警告，
在某些情形下使用引理 \ref{stacks-properties-lemma-local-source} 不会引起歧义。
下面列出这些情形。每种情形中，我们略去假设 (1) 和 (2) 的验证，
并给出蕴含 (3) 和 (4) 的参考文献。具体如下：
\begin{enumerate}
\item
\label{stacks-properties-item-rel-dim-leq-d}
$\mathcal{Q} = $``局部有限型''，$\mathcal{R} = \emptyset$，
$\mathcal{P} =$``相对维数 $\leq d$''。
见《空间的态射》定义
\ref{spaces-morphisms-definition-relative-dimension}
以及《空间的态射》引理
\ref{spaces-morphisms-lemma-openness-bounded-dimension-fibres} 和
\ref{spaces-morphisms-lemma-dimension-fibre-after-base-change}。
\item
\label{stacks-properties-item-loc-quasi-finite}
$\mathcal{Q} =$``局部有限型''，$\mathcal{R} = \emptyset$，
$\mathcal{P} =$``局部拟有限''。
这是上一项中 $d = 0$ 的情形；见《空间的态射》引理
\ref{spaces-morphisms-lemma-locally-quasi-finite-rel-dimension-0}.
另一方面，性质 (3) 和 (4) 在《空间的态射》引理
\ref{spaces-morphisms-lemma-locally-finite-type-quasi-finite-part}
中明确给出。
\item
\label{stacks-properties-item-unramified}
$\mathcal{Q} = $``局部有限型''，$\mathcal{R} = \emptyset$，
$\mathcal{P} =$``非分歧''。这就是《空间的态射》引理
\ref{spaces-morphisms-lemma-where-unramified}。
\item
\label{stacks-properties-item-flat}
$\mathcal{Q} =$``局部有限表示''，
$\mathcal{R} =$``平坦且局部有限表示''，
$\mathcal{P} =$``平坦''。见《空间态射进阶》定理
\ref{spaces-more-morphisms-theorem-openness-flatness} 和引理
\ref{spaces-more-morphisms-lemma-flat-locus-base-change}。
注意，此处 $W(\mathcal{P}, f)$ 总是恰好等于态射 $f$ 平坦的点集，
因为只在 $f$ 具有性质 $\mathcal{Q}$ 时才考虑这个开集
（见同处引文）。
\item
\label{stacks-properties-item-etale}
$\mathcal{Q} =$``局部有限表示''，
$\mathcal{R} =$``平坦且局部有限表示''，
$\mathcal{P}=$``\'etale''。这由合并
(\ref{stacks-properties-item-unramified}) 和 (\ref{stacks-properties-item-flat}) 得出，
因为平坦、局部有限表示的非分歧态射是 \'etale 态射；
见《空间的态射》引理
\ref{spaces-morphisms-lemma-unramified-flat-lfp-etale}。
\item 按需在此添加更多情形（与《群胚进阶》注
\ref{more-groupoids-remark-local-source-apply} 中更长的列表比较）。
\end{enumerate}
\end{remark}











\section{既约代数叠}
\label{stacks-properties-section-reduced}

\noindent
我们已经在第 \ref{stacks-properties-section-types-properties} 节定义了既约代数叠。

\begin{lemma}
\label{stacks-properties-lemma-reduced-closed-substack}
设 $\mathcal{X}$ 是代数叠，且 $T \subset |\mathcal{X}|$ 是闭子集。
存在唯一的闭子叠 $\mathcal{Z} \subset \mathcal{X}$，满足：
(a) $|\mathcal{Z}| = T$；(b) $\mathcal{Z}$ 既约。
\end{lemma}

\begin{proof}
设 $U \to \mathcal{X}$ 是满的光滑态射，其中 $U$ 是代数空间。
令 $R = U \times_\mathcal{X} U$，于是存在表示
$[U/R] \to \mathcal{X}$；见《代数叠》引理
\ref{algebraic-lemma-stack-presentation}。
照例将两个光滑投影态射记作 $s, t : R \to U$。
由引理 \ref{stacks-properties-lemma-points-presentation} 可知，$T$ 对应于闭子集
$T' \subset |U|$，并且
$|s|^{-1}(T') = |t|^{-1}(T')$。
令 $Z \subset U$ 为 $T'$ 上的既约诱导代数空间结构；
见《空间的性质》定义
\ref{spaces-properties-definition-reduced-induced-space}。
纤维积 $Z \times_{U, t} R$ 和 $R \times_{s, U} Z$
是 $R$ 的闭子空间
（见《空间》引理 \ref{spaces-lemma-base-change-immersions}）。
由《空间的态射》引理
\ref{spaces-morphisms-lemma-base-change-smooth}，投影
$Z \times_{U, t} R \to Z$ 和 $R \times_{s, U} Z \to Z$
都是光滑的。因此，由于 $Z$ 既约，注
\ref{stacks-properties-remark-list-properties-local-smooth-topology} 蕴含
$Z \times_{U, t} R$ 和 $R \times_{s, U} Z$ 都既约。
由于
$$
|Z \times_{U, t} R| = |t|^{-1}(T') = |s|^{-1}(T') = R \times_{s, U} Z
$$
由《空间的性质》引理
\ref{spaces-properties-lemma-reduced-closed-subspace} 中的唯一性，
可得 $Z \times_{U, t} R = R \times_{s, U} Z$。
故 $Z$ 是 $U$ 的 $R$-不变闭子空间。
通过引理 \ref{stacks-properties-lemma-substacks-presentation} 的对应，
得到闭子叠 $\mathcal{Z} \subset \mathcal{X}$，满足
$Z = \mathcal{Z} \times_\mathcal{X} U$。于是
$[Z/R_Z] \to \mathcal{Z}$ 是一个表示
（引理 \ref{stacks-properties-lemma-immersion-into-presentation}）。
因此 $|\mathcal{Z}| = |Z|/|R_Z| = |T'|/\sim$
就是给定闭子集 $T$。略去唯一性的证明。
\end{proof}

\begin{lemma}
\label{stacks-properties-lemma-reduced-stack-determined-by-points}
设 $\mathcal{X}$ 是代数叠。若
$\mathcal{X}' \subset \mathcal{X}$ 是闭子叠，
$\mathcal{X}$ 既约且
$|\mathcal{X}'| = |\mathcal{X}|$，则
$\mathcal{X}' = \mathcal{X}$。
\end{lemma}

\begin{proof}
选取表示 $[U/R] \to \mathcal{X}$，其中 $U$ 是概形。
由于 $\mathcal{X}$ 既约，由既约代数叠的定义，$U$ 既约。
由引理 \ref{stacks-properties-lemma-substacks-presentation}，
$\mathcal{X}'$ 对应于 $R$-不变闭子概形 $Z \subset U$。
现在 $|Z| \subset |U|$ 是 $|\mathcal{X}'|$ 的逆像，因而
$|Z| = |U|$。故 $Z$ 是 $U$ 的闭子概形，二者的底层点集相同。
由《概形》引理 \ref{schemes-lemma-map-into-reduction}，
映射 $\text{id}_U : U \to U$ 经 $Z \to U$ 分解，故
$Z = U$，亦即 $\mathcal{X}' = \mathcal{X}$。
\end{proof}

\begin{lemma}
\label{stacks-properties-lemma-map-into-reduction}
设 $\mathcal{X}$、$\mathcal{Y}$ 是代数叠。
设 $\mathcal{Z} \subset \mathcal{X}$ 是闭子叠。
假设 $\mathcal{Y}$ 既约。态射
$f : \mathcal{Y} \to \mathcal{X}$ 经 $\mathcal{Z}$ 分解，
当且仅当 $f(|\mathcal{Y}|) \subset |\mathcal{Z}|$。
\end{lemma}

\begin{proof}
假设 $f(|\mathcal{Y}|) \subset |\mathcal{Z}|$。考虑
$\mathcal{Y} \times_\mathcal{X} \mathcal{Z} \to \mathcal{Y}$。
存在等价
$\mathcal{Y} \times_\mathcal{X} \mathcal{Z} \to \mathcal{Y}'$，
其中 $\mathcal{Y}'$ 是 $\mathcal{Y}$ 的闭子叠；
见引理 \ref{stacks-properties-lemma-base-change-immersion} 和
\ref{stacks-properties-lemma-substack-image}。
利用引理 \ref{stacks-properties-lemma-points-cartesian}、
\ref{stacks-properties-lemma-monomorphism-injective-points} 和
\ref{stacks-properties-lemma-immersion-monomorphism}，可见
$|\mathcal{Y}'| = |\mathcal{Y}|$。因此，本引理归结为
引理 \ref{stacks-properties-lemma-reduced-stack-determined-by-points}。
\end{proof}

\begin{definition}
\label{stacks-properties-definition-reduced-induced-stack}
设 $\mathcal{X}$ 是代数叠，且
$Z \subset |\mathcal{X}|$ 是闭子集。
$Z$ 上的一个{\it 代数叠结构}，由 $\mathcal{X}$ 的一个闭子叠
$\mathcal{Z}$ 给出，并要求 $|\mathcal{Z}|$ 等于 $Z$。
$Z$ 上的{\it 既约诱导代数叠结构}，是引理
\ref{stacks-properties-lemma-reduced-closed-substack} 中构造的结构。
$\mathcal{X}$ 的{\it 既约化 $\mathcal{X}_{red}$}，
是 $|\mathcal{X}|$ 上的既约诱导代数叠结构。
\end{definition}

\noindent
事实上，可以用它定义局部闭子集上的既约诱导代数叠结构。

\begin{remark}
\label{stacks-properties-remark-stack-structure-locally-closed-subset}
设 $X$ 是代数叠，且
$T \subset |\mathcal{X}|$ 是局部闭子集。
令 $\partial T$ 为 $T$ 在拓扑空间 $|\mathcal{X}|$ 中的边界。用公式表示为
$$
\partial T = \overline{T} \setminus T.
$$
令 $\mathcal{U} \subset \mathcal{X}$ 为 $X$ 的开子叠，满足
$|\mathcal{U}| = |\mathcal{X}| \setminus \partial T$；
见引理 \ref{stacks-properties-lemma-open-substacks}。
令 $\mathcal{Z}$ 为 $\mathcal{U}$ 的既约闭子叠，满足
$|\mathcal{Z}| = T$；它通过取既约诱导闭子空间结构得到，
见定义 \ref{stacks-properties-definition-reduced-induced-stack}。
按构造，$\mathcal{Z} \to \mathcal{U}$ 是代数叠的闭浸入，
而 $\mathcal{U} \to \mathcal{X}$ 是开浸入。
因此，由引理 \ref{stacks-properties-lemma-composition-immersion}，
$\mathcal{Z} \to \mathcal{X}$ 是代数叠的浸入。
注意，$\mathcal{Z}$ 是既约代数叠，并且作为 $|X|$ 的子集，
$|\mathcal{Z}| = T$。有时称 $\mathcal{Z}$ 为 $T$ 上的
{\it 既约诱导子叠结构}。
\end{remark}













\section{剩余胚}
\label{stacks-properties-section-residual-gerbe}

\noindent
在 Stacks 项目中，我们希望将代数叠 $\mathcal{X}$ 在点
$x \in |\mathcal{X}|$ 处的{\it 剩余胚}定义为代数叠的单态射
$m_x : \mathcal{Z}_x \to \mathcal{X}$，其中
$\mathcal{Z}_x$ 是仅有一个点的既约代数叠，且该点在 $m_x$ 下映到 $x$。
事实证明，这一概念有许多问题：一般而言，其存在性不明确，唯一性也不明确。
我们通过对代数叠 $\mathcal{Z}_x$ 施加稍强的条件来解决唯一性问题。
下面通过几个关于仅有一个点的既约代数叠的简单引理作更详细的讨论。

\begin{lemma}
\label{stacks-properties-lemma-flat-cover-by-field}
设 $\mathcal{Z}$ 是代数叠，$k$ 是域，并设
$\Spec(k) \to \mathcal{Z}$ 满且平坦。
则对任一域 $k'$，任一态射
$\Spec(k') \to \mathcal{Z}$ 都满且平坦。
\end{lemma}

\begin{proof}
考虑纤维方块
$$
\xymatrix{
T \ar[d] \ar[r] & \Spec(k) \ar[d] \\
\Spec(k') \ar[r] & \mathcal{Z}
}
$$
注意，$T \to \Spec(k')$ 平坦且满，故 $T$ 非空。
另一方面，由于 $k$ 是域，$T \to \Spec(k)$ 平坦。
因此 $T \to \mathcal{Z}$ 平坦且满。
由《空间的态射》引理
\ref{spaces-morphisms-lemma-flat-permanence}
（通过第 \ref{stacks-properties-section-properties-morphisms} 节的讨论），
可知 $\Spec(k') \to \mathcal{Z}$ 平坦。
又因假设 $|\mathcal{Z}|$ 是单点集，它显然是满射。
\end{proof}

\begin{lemma}
\label{stacks-properties-lemma-unique-point}
设 $\mathcal{Z}$ 是代数叠。下列条件等价：
\begin{enumerate}
\item $\mathcal{Z}$ 既约且 $|\mathcal{Z}|$ 是单点集；
\item 存在满的平坦态射 $\Spec(k) \to \mathcal{Z}$，其中 $k$ 是域；
\item 存在局部有限型、满且平坦的态射
$\Spec(k) \to \mathcal{Z}$，其中 $k$ 是域。
\end{enumerate}
\end{lemma}

\begin{proof}
设 (1) 成立。令 $W$ 为概形，并设
$W \to \mathcal{Z}$ 是满的光滑态射。则 $W$ 是既约概形。
令 $\eta \in W$ 为 $W$ 某个不可约分支的泛点。
由于 $W$ 既约，有 $\mathcal{O}_{W, \eta} = \kappa(\eta)$。
因此典范态射 $\eta = \Spec(\kappa(\eta)) \to W$ 平坦。
于是复合 $\eta \to \mathcal{Z}$ 平坦
（见《空间的态射》引理
\ref{spaces-morphisms-lemma-composition-flat}）。
又因 $|\mathcal{Z}|$ 是单点集，该复合也是满射。换言之，(2) 成立。

\medskip\noindent
设 (2) 成立。令 $W$ 为概形，并设
$W \to \mathcal{Z}$ 是满的光滑态射。选取域 $k$
和满的平坦态射 $\Spec(k) \to \mathcal{Z}$。
则 $W \times_\mathcal{Z} \Spec(k)$ 是 $k$ 上光滑的代数空间，
因而正则（见《域上的空间》引理
\ref{spaces-over-fields-lemma-smooth-regular}），特别地既约。
由于 $W \times_\mathcal{Z} \Spec(k) \to W$ 满且平坦，
由《空间上的下降》引理
\ref{spaces-descent-lemma-descend-reduced} 可知 $W$ 既约。
换言之，(1) 成立。

\medskip\noindent
显然，(3) 蕴含 (2)。最后，设 (2) 成立。
选取非空仿射概形 $W$ 和光滑态射 $W \to \mathcal{Z}$。
选取闭点 $w \in W$，并令 $k = \kappa(w)$。复合
$$
\Spec(k) \xrightarrow{w} W \longrightarrow \mathcal{Z}
$$
由《空间的态射》引理
\ref{spaces-morphisms-lemma-composition-finite-type} 和
\ref{spaces-morphisms-lemma-smooth-locally-finite-type}.
为局部有限型。由引理 \ref{stacks-properties-lemma-flat-cover-by-field}，
它还平坦且满。因此 (3) 成立。
\end{proof}

\noindent
下一个引理从仅有一个点的代数叠中，挑出一个比上一引理略好的类。

\begin{lemma}
\label{stacks-properties-lemma-unique-point-better}
设 $\mathcal{Z}$ 是代数叠。下列条件等价：
\begin{enumerate}
\item $\mathcal{Z}$ 既约、局部 Noether，且
$|\mathcal{Z}|$ 是单点集；
\item 存在局部有限表示、满且平坦的态射
$\Spec(k) \to \mathcal{Z}$，其中 $k$ 是域。
\end{enumerate}
\end{lemma}

\begin{proof}
设 (2) 成立。由引理 \ref{stacks-properties-lemma-unique-point} 可知，
$\mathcal{Z}$ 既约且 $|\mathcal{Z}|$ 是单点集。
令 $W$ 为概形，并设 $W \to \mathcal{Z}$ 是满的光滑态射。
选取域 $k$ 以及局部有限表示、满且平坦的态射
$\Spec(k) \to \mathcal{Z}$。
则 $W \times_\mathcal{Z} \Spec(k)$ 是 $k$ 上光滑的代数空间，
因而局部 Noether（见《空间的态射》引理
\ref{spaces-morphisms-lemma-locally-finite-type-locally-noetherian}）。
由于 $W \times_\mathcal{Z} \Spec(k) \to W$
平坦、满且局部有限表示，
$\{W \times_\mathcal{Z} \Spec(k) \to W\}$ 是 fppf 覆盖。
由《空间上的下降》引理
\ref{spaces-descent-lemma-descend-locally-Noetherian}，
可知 $W$ 局部 Noether。换言之，(1) 成立。

\medskip\noindent
设 (1) 成立。选取非空仿射概形 $W$ 和光滑态射
$W \to \mathcal{Z}$。选取闭点 $w \in W$，并令
$k = \kappa(w)$。由于 $W$ 局部 Noether，态射
$w : \Spec(k) \to W$ 是有限表示的；见《态射》引理
\ref{morphisms-lemma-closed-immersion-finite-presentation}。
因此复合
$$
\Spec(k) \xrightarrow{w} W \longrightarrow \mathcal{Z}
$$
由《空间的态射》引理
\ref{spaces-morphisms-lemma-composition-finite-presentation} 和
\ref{spaces-morphisms-lemma-smooth-locally-finite-presentation}.
为局部有限表示。由引理 \ref{stacks-properties-lemma-flat-cover-by-field}，
它还平坦且满。因此 (2) 成立。
\end{proof}

\begin{lemma}
\label{stacks-properties-lemma-monomorphism-into-point}
设 $\mathcal{Z}' \to \mathcal{Z}$ 是代数叠的单态射。
假设存在域 $k$ 和局部有限表示、满且平坦的态射
$\Spec(k) \to \mathcal{Z}$。则 $\mathcal{Z}'$ 或为空，
或 $\mathcal{Z}' \to \mathcal{Z}$ 是等价。
\end{lemma}

\begin{proof}
可以假设 $\mathcal{Z}'$ 非空。此时纤维积
$T = \mathcal{Z}' \times_\mathcal{Z} \Spec(k)$ 非空；
见引理 \ref{stacks-properties-lemma-points-cartesian}。
现在 $T$ 是代数空间，且投影 $T \to \Spec(k)$ 是单态射。
故 $T = \Spec(k)$；见《空间的态射》引理
\ref{spaces-morphisms-lemma-monomorphism-toward-field}。
于是 $\Spec(k) \to \mathcal{Z}$ 经 $\mathcal{Z}'$ 分解。
设态射 $z : \Spec(k) \to \mathcal{Z}$
由 $\Spec(k)$ 上的对象 $\xi$ 给出。刚才已经看到，
$\xi$ 同构于 $\mathcal{Z}'$ 在 $\Spec(k)$ 上的一个对象 $\xi'$。
由于 $z$ 满、平坦且局部有限表示，
$\mathcal{Z}$ 在任一概形上的每个对象都在 fppf 局部同构于
$\xi$ 的拉回，因而也同构于 $\xi'$ 的拉回。
由群胚叠的对象下降，这蕴含
$\mathcal{Z}' \to \mathcal{Z}$ 本质满
（它也全忠实，见引理 \ref{stacks-properties-lemma-monomorphism}）。结论得证。
\end{proof}

\begin{lemma}
\label{stacks-properties-lemma-improve-unique-point}
设 $\mathcal{Z}$ 是代数叠。假设 $\mathcal{Z}$ 满足引理
\ref{stacks-properties-lemma-unique-point} 的等价条件。
则存在唯一的严格满子范畴
$\mathcal{Z}' \subset \mathcal{Z}$，使得
$\mathcal{Z}'$ 是满足引理
\ref{stacks-properties-lemma-unique-point-better} 之等价条件的代数叠。
包含态射 $\mathcal{Z}' \to \mathcal{Z}$ 是代数叠的单态射。
\end{lemma}

\begin{proof}
最后一个断言由第一个断言和引理 \ref{stacks-properties-lemma-monomorphism} 立即可得。
选取域 $k$ 和满、平坦且局部有限型的态射
$\Spec(k) \to \mathcal{Z}$。
令 $U = \Spec(k)$ 且 $R = U \times_\mathcal{Z} U$。
投影 $s, t : R \to U$ 局部有限型。
由于 $U$ 是域的谱，由《空间的态射》引理
\ref{spaces-morphisms-lemma-noetherian-finite-type-finite-presentation}，
可知 $s, t$ 平坦且局部有限表示。
由《可表示性判据》定理
\ref{criteria-theorem-flat-groupoid-gives-algebraic-stack}，
$\mathcal{Z}' = [U/R]$ 是代数叠。
由《代数叠》引理
\ref{algebraic-lemma-map-space-into-stack}，得到典范态射
$$
f : \mathcal{Z}' \longrightarrow \mathcal{Z}
$$
且它全忠实。因此，该态射可由代数空间表示
（见《代数叠》引理
\ref{algebraic-lemma-characterize-representable-by-algebraic-spaces}），
并且是单态射（见引理 \ref{stacks-properties-lemma-monomorphism}）。
由《可表示性判据》引理
\ref{criteria-lemma-flat-quotient-flat-presentation}，
态射 $U \to \mathcal{Z}'$ 满、平坦且局部有限表示。
故 $\mathcal{Z}'$ 是满足引理
\ref{stacks-properties-lemma-unique-point-better} 之等价条件的代数叠。
由《代数叠》引理 \ref{algebraic-lemma-equivalent}，
可以用 $\mathcal{Z}'$ 在 $\mathcal{Z}$ 中的本质像替代它。
因此，除了 $\mathcal{Z}' \subset \mathcal{Z}$ 的唯一性外，
引理的所有断言均已证明。设
$\mathcal{Z}'' \subset \mathcal{Z}$ 是另一个这样的代数叠。
则投影
$$
\mathcal{Z}'
\longleftarrow
\mathcal{Z}' \times_\mathcal{Z} \mathcal{Z}''
\longrightarrow
\mathcal{Z}''
$$
都是单态射。由引理 \ref{stacks-properties-lemma-points-cartesian}，
中间的代数叠非空。因此，由引理
\ref{stacks-properties-lemma-monomorphism-into-point}，两个投影均为同构，结论得证。
\end{proof}

\begin{example}
\label{stacks-properties-example-distinct}
下面给出一个引理 \ref{stacks-properties-lemma-improve-unique-point}
中所构造态射不是同构的例子。这个例子说明，在定义
\ref{stacks-properties-definition-residual-gerbe} 中要求剩余胚局部 Noether 是必要的。
事实上，该例甚至是代数空间！
令 $\text{Gal}(\overline{\mathbf{Q}}/\mathbf{Q})$ 为
$\mathbf{Q}$ 赋予有限群逆极限拓扑后的绝对 Galois 群。令
$$
U =
\Spec(\overline{\mathbf{Q}})
\times_{\Spec(\mathbf{Q})}
\Spec(\overline{\mathbf{Q}}) =
\text{Gal}(\overline{\mathbf{Q}}/\mathbf{Q}) \times
\Spec(\overline{\mathbf{Q}})
$$
（略去对最后一个等号含义的精确解释）。
令 $G$ 表示赋予离散拓扑的绝对 Galois 群
$\text{Gal}(\overline{\mathbf{Q}}/\mathbf{Q})$，
并将其视为 $\Spec(\overline{\mathbf{Q}})$ 上的常群概形；
见《群胚》例 \ref{groupoids-example-constant-group}。
于是 $G$ 自由且传递地作用在 $U$ 上。
令 $X = U/G$；见《空间》定义 \ref{spaces-definition-quotient}。
则 $X$ 是恰有一个点的非 Noether 既约代数空间。
此外，$X$ 有一个（局部）有限型的点：
$$
x :
\Spec(\overline{\mathbf{Q}}) \longrightarrow
U \longrightarrow X
$$
事实上，$U$ 的每个点都是闭点！
由于 $X$ 是 $\overline{\mathbf{Q}}$ 上的代数空间，
$x$ 是单态射。因此，$x$ 是引理
\ref{stacks-properties-lemma-improve-unique-point} 中构造的态射，但 $x$ 不是同构。
事实上，$\Spec(\overline{\mathbf{Q}}) \to X$
是 $X$ 在 $x$ 处的剩余胚。
\end{example}

\noindent
稍后将会看到，在对代数叠 $\mathcal{X}$ 的温和假设下，
下一个引理的等价条件对每个点 $x \in |\mathcal{X}|$ 都成立
（见《叠的态射》第
\ref{stacks-morphisms-section-existence-residual-gerbes} 节）。

\begin{lemma}
\label{stacks-properties-lemma-residual-gerbe}
设 $\mathcal{X}$ 是代数叠，且 $x \in |\mathcal{X}|$ 是一个点。
下列条件等价：
\begin{enumerate}
\item 存在代数叠 $\mathcal{Z}$ 和单态射
$\mathcal{Z} \to \mathcal{X}$，使得 $|\mathcal{Z}|$ 是单点集，
且 $|\mathcal{Z}|$ 在 $|\mathcal{X}|$ 中的像是 $x$；
\item 存在既约代数叠 $\mathcal{Z}$ 和单态射
$\mathcal{Z} \to \mathcal{X}$，使得 $|\mathcal{Z}|$ 是单点集，
且 $|\mathcal{Z}|$ 在 $|\mathcal{X}|$ 中的像是 $x$；
\item 存在代数叠 $\mathcal{Z}$、单态射
$f : \mathcal{Z} \to \mathcal{X}$，以及满的平坦态射
$z : \Spec(k) \to \mathcal{Z}$，其中 $k$ 是域，并且
$x = f(z)$。
\end{enumerate}
此外，若这些条件成立，则存在唯一的严格满子范畴
$\mathcal{Z}_x \subset \mathcal{X}$，使得
$\mathcal{Z}_x$ 是既约、局部 Noether 的代数叠，
且 $|\mathcal{Z}_x|$ 是单点集，并在映射
$|\mathcal{Z}_x| \to |\mathcal{X}|$ 下映到 $x$。
\end{lemma}

\begin{proof}
若 $\mathcal{Z} \to \mathcal{X}$ 如 (1)，则
$\mathcal{Z}_{red} \to \mathcal{X}$ 如 (2)。
（代数叠既约化的概念见第 \ref{stacks-properties-section-reduced} 节。）
故 (1) 蕴含 (2)。(2) 蕴含 (1) 是直接的。
(2) 与 (3) 的等价性由引理 \ref{stacks-properties-lemma-unique-point} 立即可得。

\medskip\noindent
至此已证明 (1) -- (3) 等价。选取 (2) 中的单态射
$f : \mathcal{Z} \to \mathcal{X}$。
注意，这蕴含 $f$ 全忠实；见引理 \ref{stacks-properties-lemma-monomorphism}。
将函子 $f$ 的本质像记作
$\mathcal{Z}' \subset \mathcal{X}$。则
$f : \mathcal{Z} \to \mathcal{Z}'$ 是等价，故
$\mathcal{Z}'$ 是代数叠；见《代数叠》引理
\ref{algebraic-lemma-equivalent}。
应用引理 \ref{stacks-properties-lemma-improve-unique-point}，
得到引理陈述中的严格满子范畴
$\mathcal{Z}_x \subset \mathcal{Z}'$。
这证明了除唯一性之外的所有断言。

\medskip\noindent
为证明唯一性，设
$\mathcal{Z}_x \subset \mathcal{X}$
和
$\mathcal{Z}'_x \subset \mathcal{X}$
是引理陈述中的两个严格满子范畴。则投影
$$
\mathcal{Z}'_x
\longleftarrow
\mathcal{Z}'_x \times_\mathcal{X} \mathcal{Z}_x
\longrightarrow
\mathcal{Z}_x
$$
都是单态射。由引理 \ref{stacks-properties-lemma-points-cartesian}，
中间的代数叠非空。因此，由引理
\ref{stacks-properties-lemma-monomorphism-into-point}，两个投影均为同构，结论得证。
\end{proof}

\noindent
经过上述说明，现在可以作如下定义。

\begin{definition}
\label{stacks-properties-definition-residual-gerbe}
设 $\mathcal{X}$ 是代数叠，且 $x \in |\mathcal{X}|$。
\begin{enumerate}
\item 若引理 \ref{stacks-properties-lemma-residual-gerbe} 的等价条件
(1)、(2)、(3) 成立，则称{\it $\mathcal{X}$ 在 $x$ 处的剩余胚存在}。
\item 若 $\mathcal{X}$ 在 $x$ 处的剩余胚存在，则
{\it $\mathcal{X}$ 在 $x$ 处的剩余胚}\footnote{这在精神上与
\cite{LM-B} 有冲突，但事实上并无冲突。具体地，他们在第 11 章
为任一拟分离代数叠上的任一点赋予一个胚（不必是代数的），
并称其为剩余胚。我们将在《叠的态射》引理
\ref{stacks-morphisms-lemma-every-point-residual-gerbe}
中看到，在拟分离代数叠上，每个点都有我们意义下的剩余胚，
且该剩余胚与他们的定义等价。关于这一主题的更多信息，见
\cite[Appendix B]{rydh_etale_devissage}。}
是引理 \ref{stacks-properties-lemma-residual-gerbe} 中构造的严格满子范畴
$\mathcal{Z}_x \subset \mathcal{X}$。
\end{enumerate}
\end{definition}

\noindent
特别地，我们知道 $\mathcal{Z}_x$（若存在）是局部 Noether 的
既约代数叠，并且存在一个域和一个满、平坦、局部有限表示的态射
$$
\Spec(k) \longrightarrow \mathcal{Z}_x.
$$
我们将在《叠的态射》引理
\ref{stacks-morphisms-lemma-noetherian-singleton-stack-gerbe}
中看到，$\mathcal{Z}_x$ 是胚。剩余胚的存在性在
《叠的态射》第
\ref{stacks-morphisms-section-existence-residual-gerbes} 节讨论。

\begin{example}
\label{stacks-properties-example-residual-gerbe}
设 $X$ 是概形，且 $x \in X$ 是一个点。则单态射
$x \to X$ 是 $X$ 在 $x$ 处的剩余胚；这里照例将
$x$ 与概形 $x = \Spec(\kappa(x))$ 等同。
若 $X$ 是代数空间且 $x \in |X|$，则 $x$ 处的剩余胚
（称为剩余空间）总是存在；见《适度空间》第
\ref{decent-spaces-section-singleton} 节。
\end{example}

\noindent
由下一个引理，剩余胚若存在，便是正则代数叠。

\begin{lemma}
\label{stacks-properties-lemma-residual-gerbe-regular}
若既约、局部 Noether 的代数叠 $\mathcal{Z}$
满足 $|\mathcal{Z}|$ 是单点集，则它正则。
\end{lemma}

\begin{proof}
设 $W \to \mathcal{Z}$ 是满的光滑态射，其中 $W$ 是概形。
设 $k$ 是域，并设 $\Spec(k) \to \mathcal{Z}$
满、平坦且局部有限表示（见引理 \ref{stacks-properties-lemma-unique-point-better}）。
代数空间 $T = W \times_\mathcal{Z} \Spec(k)$ 在 $k$ 上光滑，
特别地正则；见《域上的空间》引理
\ref{spaces-over-fields-lemma-smooth-regular}。
由于 $T \to W$ 局部有限表示、平坦且满，
由《空间上的下降》引理
\ref{spaces-descent-lemma-descend-regular}，可知 $W$ 正则。
按定义，这意味着 $\mathcal{Z}$ 正则。
\end{proof}

\begin{lemma}
\label{stacks-properties-lemma-residual-gerbe-points}
设 $\mathcal{X}$ 是代数叠，且 $x \in |\mathcal{X}|$。
假设 $\mathcal{X}$ 的剩余胚 $\mathcal{Z}_x$ 存在。
设 $f : \Spec(K) \to \mathcal{X}$ 是 $x$ 的等价类中的态射，
其中 $K$ 是域。则 $f$ 经包含态射
$\mathcal{Z}_x \to \mathcal{X}$ 分解。
\end{lemma}

\begin{proof}
选取域 $k$ 和满、平坦、局部有限表示的态射
$\Spec(k) \to \mathcal{Z}_x$。令
$T = \Spec(K) \times_\mathcal{X} \mathcal{Z}_x$。
由引理 \ref{stacks-properties-lemma-points-cartesian} 可知 $T$ 非空。
由于 $\mathcal{Z}_x \to \mathcal{X}$ 是单态射，
$T \to \Spec(K)$ 是单态射。因此，由《空间的态射》引理
\ref{spaces-morphisms-lemma-monomorphism-toward-field}，
可知 $T = \Spec(K)$，引理得证。
\end{proof}

\begin{lemma}
\label{stacks-properties-lemma-residual-gerbe-unique}
设 $\mathcal{X}$ 是代数叠，且 $x \in |\mathcal{X}|$。
设 $\mathcal{Z}$ 是满足引理
\ref{stacks-properties-lemma-unique-point-better} 之等价条件的代数叠，
并设 $\mathcal{Z} \to \mathcal{X}$ 是单态射，使得
$|\mathcal{Z}| \to |\mathcal{X}|$ 的像是 $x$。
则 $\mathcal{X}$ 在 $x$ 处的剩余胚
$\mathcal{Z}_x$ 存在，并且
$\mathcal{Z} \to \mathcal{X}$ 分解为
$\mathcal{Z} \to \mathcal{Z}_x \to \mathcal{X}$，
其中第一个箭头是等价。
\end{lemma}

\begin{proof}
令 $\mathcal{Z}_x \subset \mathcal{X}$ 为与函子
$\mathcal{Z} \to \mathcal{X}$ 的本质像对应的满子范畴。
则 $\mathcal{Z} \to \mathcal{Z}_x$ 是等价，故
$\mathcal{Z}_x$ 是代数叠；见《代数叠》引理
\ref{algebraic-lemma-equivalent}。
由于 $\mathcal{Z}_x$ 通过该等价继承 $\mathcal{Z}$ 的所有性质，
由引理 \ref{stacks-properties-lemma-residual-gerbe} 中的唯一性显然可知，
$\mathcal{Z}_x$ 是 $\mathcal{X}$ 在 $x$ 处的剩余胚。
\end{proof}

\begin{lemma}
\label{stacks-properties-lemma-residual-gerbe-functorial}
设 $f : \mathcal{X} \to \mathcal{Y}$ 是代数叠态射。
设 $x \in |\mathcal{X}|$ 的像为 $y \in |\mathcal{Y}|$。
若 $x$ 和 $y$ 的剩余胚
$\mathcal{Z}_x \subset \mathcal{X}$ 与
$\mathcal{Z}_y \subset \mathcal{Y}$ 均存在，
则 $f$ 诱导交换图表
$$
\xymatrix{
\mathcal{X} \ar[d]_f & \mathcal{Z}_x \ar[l] \ar[d] \\
\mathcal{Y} & \mathcal{Z}_y \ar[l]
}
$$
\end{lemma}

\begin{proof}
选取域 $k$ 和满、平坦、局部有限表示的态射
$\Spec(k) \to \mathcal{Z}_x$。由引理
\ref{stacks-properties-lemma-residual-gerbe-points}，态射
$\Spec(k) \to \mathcal{Y}$ 经 $\mathcal{Z}_y$ 分解。
因此 $\mathcal{Z}_x \times_\mathcal{Y} \mathcal{Z}_y$
是 $\mathcal{Z}_x$ 的非空子叠，故由引理
\ref{stacks-properties-lemma-monomorphism-into-point}，它等于 $\mathcal{Z}_x$。
\end{proof}

\begin{lemma}
\label{stacks-properties-lemma-residual-gerbe-isomorphic}
设 $f : \mathcal{X} \to \mathcal{Y}$ 是代数叠态射。
设 $x \in |\mathcal{X}|$ 的像为 $y \in |\mathcal{Y}|$。
假设 $x$ 和 $y$ 的剩余胚
$\mathcal{Z}_x \subset \mathcal{X}$ 与
$\mathcal{Z}_y \subset \mathcal{Y}$ 均存在，
并且存在 $x$ 的等价类中的态射
$\Spec(k) \to \mathcal{X}$，使得
$$
\Spec(k) \times_\mathcal{X} \Spec(k)
\longrightarrow
\Spec(k) \times_\mathcal{Y} \Spec(k)
$$
是同构。则 $\mathcal{Z}_x \to \mathcal{Z}_y$ 是同构。
\end{lemma}

\begin{proof}
设 $k'/k$ 是域扩张。则
$$
\Spec(k') \times_\mathcal{X} \Spec(k')
\longrightarrow
\Spec(k') \times_\mathcal{Y} \Spec(k')
$$
是引理中的态射沿忠实平坦态射
$\Spec(k' \otimes k') \to \Spec(k \otimes k)$ 的基变换。
因此，引理所述性质与 $x$ 的等价类中态射
$\Spec(k) \to \mathcal{X}$ 的选择无关。
故可以假设 $\Spec(k) \to \mathcal{Z}_x$
满、平坦且局部有限表示。此时有
$$
\mathcal{Z}_x = [\Spec(k)/R]
$$
，其中 $R = \Spec(k) \times_\mathcal{X} \Spec(k)$；
见引理 \ref{stacks-properties-lemma-improve-unique-point} 的证明。
由于还有 $R = \Spec(k) \times_\mathcal{Y} \Spec(k)$，
可知引理 \ref{stacks-properties-lemma-residual-gerbe-functorial} 中的态射
$\mathcal{Z}_x \to \mathcal{Z}_y$ 全忠实；这是由
《代数叠》引理
\ref{algebraic-lemma-map-space-into-stack} 得到的。
例如，再由引理 \ref{stacks-properties-lemma-residual-gerbe-unique} 即得结论。
\end{proof}








\section{叠的维数}
\label{stacks-properties-section-dimension}

\noindent
利用代数空间在一点处的维数概念（《空间的性质》定义
\ref{spaces-properties-definition-dimension-at-point}），
可以定义代数叠 $\mathcal{X}$ 在点 $x$ 处的维数。
下一个引理中的结果可能是 $\infty$：这或者因为
$\mathcal{X}$ 不拟紧，或者因为遇到《例》第
\ref{examples-section-Noetherian-infinite-dimension} 节所述现象。

\begin{lemma}
\label{stacks-properties-lemma-dimension-at-point-well-defined}
设 $\mathcal{X}$ 是概形 $S$ 上局部 Noether 的代数叠，
且 $x \in |\mathcal{X}|$ 是 $\mathcal{X}$ 的点。
设 $[U/R] \to \mathcal{X}$ 是一个表示
（《代数叠》定义 \ref{algebraic-definition-presentation}），
其中 $U$ 是概形。设 $u \in U$ 是映到 $x$ 的点。
令 $e : U \to R$ 为“恒等”映射，并令 $s : R \to U$
为“源”映射；后者是代数空间的光滑态射。
令 $R_u$ 为 $s : R \to U$ 在 $u$ 上的纤维。则元素
$$
\dim_x(\mathcal{X}) = \dim_u(U) - \dim_{e(u)}(R_u) \in \mathbf{Z} \cup \infty
$$
与表示和 $x$ 上方之点 $u$ 的选择无关。
\end{lemma}

\begin{proof}
由于 $R \to U$ 光滑，概形 $R_u$ 在 $\kappa(u)$ 上光滑，
因而具有有限维数。另一方面，概形 $U$ 局部 Noether，
但这并不保证 $\dim_u(U)$ 有限。因此，该差属于
$\mathbf{Z} \cup \{\infty\}$。

\medskip\noindent
设 $[U'/R'] \to \mathcal{X}$ 和 $u' \in U'$ 为第二组表示与点，
其中 $U'$ 是概形，且 $u'$ 映到 $x$。
考虑代数空间 $P = U \times_\mathcal{X} U'$。
由引理 \ref{stacks-properties-lemma-points-cartesian}，存在
$p \in |P|$ 映到 $u$ 和 $u'$。由于 $P \to U$ 和
$P \to U'$ 光滑，可得
$\dim_p(P) = \dim_u(U) + \dim_p(P_u)$ 以及
$\dim_p(P) = \dim_{u'}(U') + \dim_p(P_{u'})$；
见《空间的态射》引理
\ref{spaces-morphisms-lemma-smoothness-dimension-spaces}。
注意
$$
R'_{u'} = \Spec(\kappa(u')) \times_\mathcal{X} U'
\quad\text{且}\quad
P_u = \Spec(\kappa(u)) \times_\mathcal{X} U'
$$
用态射 $\Spec(\Omega) \to P$ 表示 $p \in |P|$。
由于 $p$ 同时映到 $u$ 和 $u'$，它诱导下列两个复合之间的 $2$-态射：
$\Spec(\Omega) \to \Spec(\kappa(u')) \to \mathcal{X}$ 和
$\Spec(\Omega) \to \Spec(\kappa(u)) \to \mathcal{X}$
，进而定义同构
$$
\Spec(\Omega) \times_{\Spec(\kappa(u'))} R'_{u'}
\cong
\Spec(\Omega) \times_{\Spec(\kappa(u))} P_u
$$
，它是 $\Spec(\Omega)$ 上代数空间的同构，并将
$\Omega$-有理点 $(1, e'(u'))$ 映到 $(1, p)$（略去一些细节）。
由此可得
$$
\dim_{e'(u')}(R'_{u'}) = \dim_p(P_u)
$$
，这是由《空间的态射》引理
\ref{spaces-morphisms-lemma-dimension-fibre-after-base-change}.
得到的。由对称性，有
$\dim_{e(u)}(R_u) = \dim_p(P_{u'})$。
综合以上各式，即得与选择无关。
\end{proof}

\noindent
可以利用上述引理作如下定义。

\begin{definition}
\label{stacks-properties-definition-dimension-at-point}
设 $\mathcal{X}$ 是概形 $S$ 上局部 Noether 的代数叠，
且 $x \in |\mathcal{X}|$ 是 $\mathcal{X}$ 的点。
设 $[U/R] \to \mathcal{X}$ 是一个表示
（《代数叠》定义 \ref{algebraic-definition-presentation}），
其中 $U$ 是概形；并设 $u \in U$ 是映到 $x$ 的点。
将{\it $\mathcal{X}$ 在 $x$ 处的维数}定义为元素
$\dim_x(\mathcal{X}) \in \mathbf{Z} \cup \infty$，满足
$$
\dim_x(\mathcal{X}) = \dim_u(U)-\dim_{e(u)}(R_u).
$$
记号同引理 \ref{stacks-properties-lemma-dimension-at-point-well-defined}。
\end{definition}

\noindent
当 $\mathcal{X}$ 是概形时（《拓扑》定义
\ref{topology-definition-Krull}），叠在一点处的维数与通常概念一致；
更一般地，当 $\mathcal{X}$ 是局部 Noether 代数空间时亦如此
（《空间的性质》定义
\ref{spaces-properties-definition-dimension-at-point}）。

\begin{definition}
\label{stacks-properties-definition-dimension}
设 $S$ 是概形，且 $\mathcal{X}$ 是 $S$ 上局部 Noether 的代数叠。
将 $\mathcal{X}$ 的{\it 维数} $\dim(\mathcal{X})$ 定义为
$$
\dim(\mathcal{X}) = \sup\nolimits_{x \in |\mathcal{X}|} \dim_x(\mathcal{X})
$$
\end{definition}

\noindent
若 $\mathcal{X}$ 是概形
（《性质》引理 \ref{properties-lemma-dimension}）
或代数空间（《空间的性质》定义
\ref{spaces-properties-definition-dimension}），
该维数定义与通常概念一致。

\begin{remark}
\label{stacks-properties-remark-dimension-empty-stack}
若 $\mathcal{X}$ 是域上非空的有限型叠，则
$\dim(\mathcal{X})$ 是整数。对任意局部 Noether 代数叠
$\mathcal{X}$，$\dim(\mathcal{X})$ 属于
$Z\cup \{\pm \infty\}$；
且 $\dim(\mathcal{X}) = -\infty$ 当且仅当 $\mathcal{X}$ 为空。
\end{remark}

\begin{example}
\label{stacks-properties-example-dimension-quotient-stack}
设 $X$ 是域 $k$ 上的有限型概形，且 $G$ 是作用在 $X$ 上的
$k$ 上有限型群概形。则商叠 $[X/G]$ 的维数等于
$\dim(X)-\dim(G)$。特别地，分类叠
$BG=[\Spec(k)/G]$ 的维数为 $-\dim(G)$。
因此，与概形或代数空间的情形不同，代数叠的维数可以是负整数。
\end{example}






\section{局部不可约性}
\label{stacks-properties-section-irreducible-local-ring}

\noindent
我们已经在《性质》第 \ref{properties-section-unibranch} 节
定义了概形在一点处的几何分支数，并在《空间的性质》第
\ref{spaces-properties-section-irreducible-local-ring} 节
定义了代数空间在一点处的几何分支数。
设 $n \in \mathbf{N}$。对局部环 $A$，令
$$
P_n(A) = \text{the number of geometric branches of }A\text{ is }n
$$
对光滑环同态 $A \to B$ 和 $B$ 中位于 $A$ 的素理想
$\mathfrak p$ 上方的素理想 $\mathfrak q$，有
$$
P_n(A_\mathfrak p) \Leftrightarrow P_n(B_\mathfrak q)
$$
，这是由《代数进阶》引理
\ref{more-algebra-lemma-invariance-number-branches-smooth} 得到的。
如同《空间的性质》注
\ref{spaces-properties-remark-list-properties-local-ring-local-etale-topology}，
可以利用 $P_n$ 定义概形芽 $(U, u)$ 的 \'etale 局部性质
$\mathcal{P}_n$，令
$\mathcal{P}_n(U, u) = P_n(\mathcal{O}_{U, u})$。
代数空间 $X$ 在点 $x$ 处的相应性质 $\mathcal{P}_n$
（见《空间的性质》定义
\ref{spaces-properties-definition-property-at-point}）
恰好是性质“$X$ 在 $x$ 处的几何分支数为 $n$”；
见《空间的性质》定义
\ref{spaces-properties-definition-number-of-geometric-branches}。
此外，性质 $\mathcal{P}_n$ 是光滑局部的；
见《下降》定义 \ref{descent-definition-local-at-point}。
这或者由上面显示的等价式推出，或者由《态射进阶》引理
\ref{more-morphisms-lemma-number-of-branches-and-smooth} 推出。
因此，定义 \ref{stacks-properties-definition-property-at-point} 适用，
从而得到代数叠在一点处的相应概念。

\begin{definition}
\label{stacks-properties-definition-number-of-geometric-branches}
设 $\mathcal{X}$ 是代数叠，且 $x \in |\mathcal{X}|$。
\begin{enumerate}
\item 若引理 \ref{stacks-properties-lemma-local-source-target-at-point} 的等价条件
对上述定义的 $\mathcal{P}_n$ 成立，则
{\it $\mathcal{X}$ 在 $x$ 处的几何分支数}为
$n \in \mathbf{N}$；否则为 $\infty$。
\item 若 $\mathcal{X}$ 在 $x$ 处的几何分支数为 $1$，
则称 $\mathcal{X}$ {\it 在 $x$ 处几何单分支}。
\end{enumerate}
\end{definition}






\section{有限性条件与点}
\label{stacks-properties-section-points-conditions}

\noindent
本节是《适度空间》第
\ref{decent-spaces-section-points-monomorphisms} 节
对代数叠之点的类似版本。

\begin{lemma}
\label{stacks-properties-lemma-UR-quasi-compact-above-x}
设 $\mathcal{X}$ 是代数叠，且 $x \in |\mathcal{X}|$ 是一个点。
下列条件等价：
\begin{enumerate}
\item $x$ 的等价类中的某个态射
$\Spec(k) \to \mathcal{X}$ 拟紧；
\item $x$ 的等价类中的任一态射
$\Spec(k) \to \mathcal{X}$ 拟紧。
\end{enumerate}
\end{lemma}

\begin{proof}
设 $\Spec(k) \to \mathcal{X}$ 属于 $x$ 的等价类，
并设 $k'/k$ 是域扩张。需要证明
$\Spec(k) \to \mathcal{X}$ 拟紧，当且仅当
$\Spec(k') \to \mathcal{X}$ 拟紧。
这由《空间的态射》引理
\ref{spaces-morphisms-lemma-surjection-from-quasi-compact}
以及《代数叠》引理
\ref{algebraic-lemma-representable-transformations-property-implication}
的原理推出。
\end{proof}

\noindent
有时将满足引理 \ref{stacks-properties-lemma-UR-quasi-compact-above-x}
之等价条件的点 $x \in |\mathcal{X}|$
称为“{\it 拟紧点}”。
% Generated mechanically from the frozen stacks-morphisms.tex target.
% Do not edit this generated file; edit the bound locale source instead.

% OK, start here.
%

\setcounter{chapter}{100}
\chapter{代数叠的态射}



\phantomsection
\label{stacks-morphisms-section-phantom}


\section{引言}
\label{stacks-morphisms-section-introduction}

\noindent
本章介绍若干类代数叠态射。对于对角态射可表示的拟分离代数叠，
可参考 \cite{LM-B}。

\medskip\noindent
我们的目标是把各类代数空间态射的定义推广到代数叠态射。
各情形略有不同，因此最好分别处理。

\medskip\noindent
对于可由代数空间表示的代数叠态射，我们已经定义了许多类态射；
见《叠的性质》第
\ref{stacks-properties-section-properties-morphisms} 节。
本章处理每个相应情形时，都必须确保一般情形下的定义与该处的定义相容。




\section{约定与术语滥用}
\label{stacks-morphisms-section-conventions}

\noindent
我们继续采用《叠的性质》第
\ref{stacks-properties-section-conventions} 节引入的约定与术语滥用。




\section{对角态射的性质}
\label{stacks-morphisms-section-diagonals}

\noindent
代数叠的对角态射与 $\mathit{Isom}$-层密切相关；见《代数叠》引理
\ref{algebraic-lemma-representable-diagonal}。
由代数叠定义中的第二个性质，这些 $\mathit{Isom}$-层总是代数空间。

\begin{lemma}
\label{stacks-morphisms-lemma-isom-locally-finite-type}
设 $\mathcal{X}$ 是代数叠。设 $T$ 是概形，且 $x,y$ 是
$\mathcal{X}$ 在 $T$ 上的纤维范畴中的对象。则态射
$\mathit{Isom}_\mathcal{X}(x, y) \to T$ 局部有限型。
\end{lemma}

\begin{proof}
由《代数叠》引理 \ref{algebraic-lemma-stack-presentation}，
可设对某个代数空间中的光滑群胚有
$\mathcal{X} = [U/R]$。由《空间的下降》引理
\ref{spaces-descent-lemma-descending-property-locally-finite-type}，
只需在 $T$ 上 fppf 局部地检验该性质。因此可设 $x,y$ 来自态射
$x', y' : T \to U$。由《空间中的群胚》引理
\ref{spaces-groupoids-lemma-quotient-stack-morphisms}，此时
$\mathit{Isom}_\mathcal{X}(x, y) = T \times_{(y', x'), U \times_S U} R$.
故只需证明 $R \to U \times_S U$ 局部有限型。这由复合
$s : R \to U \times_S U \to U$ 光滑而得（因而局部有限型；见
《空间的态射》引理
\ref{spaces-morphisms-lemma-smooth-locally-finite-presentation} 和
\ref{spaces-morphisms-lemma-finite-presentation-finite-type})
以及《空间的态射》引理
\ref{spaces-morphisms-lemma-permanence-finite-type}。
\end{proof}

\begin{lemma}
\label{stacks-morphisms-lemma-isom-pseudo-torsor-aut}
设 $\mathcal{X}$ 是代数叠。设 $T$ 是概形，且 $x,y$ 是
$\mathcal{X}$ 在 $T$ 上的纤维范畴中的对象。则
\begin{enumerate}
\item $\mathit{Isom}_\mathcal{X}(y, y)$ 是 $T$ 上的群代数空间，并且
\item $\mathit{Isom}_\mathcal{X}(x, y)$ 是 $T$ 上
$\mathit{Isom}_\mathcal{X}(y, y)$ 的伪挠子。
\end{enumerate}
\end{lemma}

\begin{proof}
见《空间中的群胚》定义
\ref{spaces-groupoids-definition-group-space} 和
\ref{spaces-groupoids-definition-pseudo-torsor}。
该引理由 $\mathcal{X}$ 是群胚中的叠立即得到。
\end{proof}

\noindent
设 $f : \mathcal{X} \to \mathcal{Y}$ 是代数叠的态射。
{\it $f$ 的对角态射}是态射
$$
\Delta_f :
\mathcal{X}
\longrightarrow
\mathcal{X} \times_\mathcal{Y} \mathcal{X}
$$
每个对角态射都具有以下两个性质。

\begin{lemma}
\label{stacks-morphisms-lemma-properties-diagonal}
\begin{slogan}
代数叠态射的对角态射可由代数空间表示，且局部有限型。
\end{slogan}
设 $f : \mathcal{X} \to \mathcal{Y}$ 是代数叠的态射。则
\begin{enumerate}
\item $\Delta_f$ 可由代数空间表示，并且
\item $\Delta_f$ 局部有限型。
\end{enumerate}
\end{lemma}

\begin{proof}
设 $T$ 是概形，且
$a : T \to \mathcal{X} \times_\mathcal{Y} \mathcal{X}$
是态射。由纤维积的定义和 $2$-Yoneda 引理，态射 $a$ 由三元组
$a = (x, x', \alpha)$ 给出，其中 $x,x'$ 是 $\mathcal{X}$
在 $T$ 上的对象，而 $\alpha : f(x) \to f(x')$ 是
$\mathcal{Y}$ 在 $T$ 上的纤维范畴中的态射。由代数叠的定义，层
$\mathit{Isom}_\mathcal{X}(x, x')$ 和
$\mathit{Isom}_\mathcal{Y}(f(x), f(x'))$ 是 $T$ 上的代数空间。
在此表述下，$\alpha$ 定义态射
$\mathit{Isom}_\mathcal{Y}(f(x), f(x')) \to T$ 的一个截面。对于
$T$ 上的概形 $T' \to T$，
$\mathcal{X} \times_{\mathcal{X} \times_\mathcal{Y} \mathcal{X}, a} T$
的一个 $T'$-值点等同于一个同构
$x|_{T'} \to x'|_{T'}$，其在 $f$ 下的像为 $\alpha|_{T'}$。
因此可见
\begin{equation}
\label{stacks-morphisms-equation-diagonal}
\vcenter{
\xymatrix{
\mathcal{X} \times_{\mathcal{X} \times_\mathcal{Y} \mathcal{X}, a} T
\ar[d] \ar[r] &
\mathit{Isom}_\mathcal{X}(x, x') \ar[d] \\
T\ar[r]^-\alpha &
\mathit{Isom}_\mathcal{Y}(f(x), f(x'))
}
}
\end{equation}
是 $T$ 上层的纤维方块。特别地，
$\mathcal{X} \times_{\mathcal{X} \times_\mathcal{Y} \mathcal{X}, a} T$
是代数空间，这证明了引理的 (1)。

\medskip\noindent
为证明第二个断言，需证明图 (\ref{stacks-morphisms-equation-diagonal}) 的左竖直箭头
局部有限型。由引理 \ref{stacks-morphisms-lemma-isom-locally-finite-type}，
代数空间 $\mathit{Isom}_\mathcal{X}(x, x')$ 在 $T$ 上局部有限型。
因此图 (\ref{stacks-morphisms-equation-diagonal}) 的右竖直箭头局部有限型；见
《空间的态射》引理 \ref{spaces-morphisms-lemma-permanence-finite-type}。
再由《空间的态射》引理
\ref{spaces-morphisms-lemma-base-change-finite-type} 即得结论。
\end{proof}

\begin{lemma}
\label{stacks-morphisms-lemma-properties-diagonal-representable}
设 $f : \mathcal{X} \to \mathcal{Y}$ 是可由代数空间表示的代数叠态射。则
\begin{enumerate}
\item $\Delta_f$ 可表示（由概形表示），
\item $\Delta_f$ 局部有限型，
\item $\Delta_f$ 是单态射，
\item $\Delta_f$ 分离，并且
\item $\Delta_f$ 局部拟有限。
\end{enumerate}
\end{lemma}

\begin{proof}
引理 \ref{stacks-morphisms-lemma-properties-diagonal} 已说明 $\Delta_f$
可由代数空间表示。因此断言 (2) -- (5) 有意义；见《叠的性质》第
\ref{stacks-properties-section-properties-morphisms} 节。
引理 \ref{stacks-morphisms-lemma-properties-diagonal} 也保证 (2) 成立。设
$T \to \mathcal{X} \times_\mathcal{Y} \mathcal{X}$ 是态射，并考察
图 (\ref{stacks-morphisms-equation-diagonal})。由《代数叠》引理
\ref{algebraic-lemma-criterion-map-representable-spaces-fibred-in-groupoids}
可知，右竖直箭头作为层态射是单射，即为代数空间的单态射。因此态射
$T \times_{\mathcal{X} \times_\mathcal{Y} \mathcal{X}} \mathcal{X} \to T$
也是单态射，故 (3) 成立。我们已经知道
$T \times_{\mathcal{X} \times_\mathcal{Y} \mathcal{X}} \mathcal{X} \to T$
局部有限型。因此《空间的态射》引理
\ref{spaces-morphisms-lemma-monomorphism-loc-finite-type-loc-quasi-finite}
说明
$T \times_{\mathcal{X} \times_\mathcal{Y} \mathcal{X}} \mathcal{X} \to T$
局部拟有限且分离。这证明了 (4) 和 (5)。最后，《空间的态射》命题
\ref{spaces-morphisms-proposition-locally-quasi-finite-separated-over-scheme}
说明
$T \times_{\mathcal{X} \times_\mathcal{Y} \mathcal{X}} \mathcal{X}$
是概形，从而证明 (1)。
\end{proof}

\begin{lemma}
\label{stacks-morphisms-lemma-representable-separated-diagonal-closed}
设 $f : \mathcal{X} \to \mathcal{Y}$ 是可由代数空间表示的代数叠态射。
则下列条件等价：
\begin{enumerate}
\item $f$ 分离；
\item $\Delta_f$ 是闭浸入；
\item $\Delta_f$ 是真态射；
\item $\Delta_f$ 泛闭。
\end{enumerate}
\end{lemma}

\begin{proof}
“$f$ 分离”“$\Delta_f$ 是闭浸入”“$\Delta_f$ 泛闭”以及
“$\Delta_f$ 是真态射”均指《叠的性质》第
\ref{stacks-properties-section-properties-morphisms} 节定义的概念。
选取概形 $V$ 以及满光滑态射 $V \to \mathcal{Y}$。
令 $U = \mathcal{X} \times_\mathcal{Y} V$；由假设它是代数空间，
而态射 $U \to \mathcal{X}$ 满且光滑。由《范畴》引理
\ref{categories-lemma-base-change-diagonal} 和《叠的性质》引理
\ref{stacks-properties-lemma-check-property-covering} 可知，对任意性质
$P$（如该引理中），$\Delta_f$ 具有 $P$ 当且仅当
$\Delta_{U/V} : U \to  U \times_V U$ 具有 $P$。
把《空间的态射》引理
\ref{spaces-morphisms-lemma-separated-diagonal-proper}
应用于 $U \to V$，即得 (2)、(3)、(4) 等价。
此外，若 (1) 成立，则 $U \to V$ 分离，故
$\Delta_{U/V}$ 是闭浸入，即 (2) 成立。
最后，设 (2) 成立。设 $T$ 是概形，且 $a : T \to \mathcal{Y}$
是态射。令 $T' = \mathcal{X} \times_\mathcal{Y} T$。
为证明 (1)，需证明代数空间的态射 $T' \to T$ 分离。
再次使用《范畴》引理 \ref{categories-lemma-base-change-diagonal}，
可见 $\Delta_{T'/T}$ 是 $\Delta_f$ 的基变换。
因此假设 (2) 蕴含 $\Delta_{T'/T}$ 是闭浸入，从而
$T' \to T$ 分离，正合所需。
\end{proof}

\begin{lemma}
\label{stacks-morphisms-lemma-representable-quasi-separated-diagonal-quasi-compact}
设 $f : \mathcal{X} \to \mathcal{Y}$ 是可由代数空间表示的代数叠态射。
则下列条件等价：
\begin{enumerate}
\item $f$ 拟分离；
\item $\Delta_f$ 拟紧；
\item $\Delta_f$ 有限型。
\end{enumerate}
\end{lemma}

\begin{proof}
“$f$ 拟分离”“$\Delta_f$ 拟紧”以及“$\Delta_f$ 有限型”
均指《叠的性质》第
\ref{stacks-properties-section-properties-morphisms} 节定义的概念。
注意，由引理 \ref{stacks-morphisms-lemma-properties-diagonal-representable}
（以及《代数叠》引理
\ref{algebraic-lemma-representable-transformations-property-implication}），
$\Delta_f$ 局部有限型，故 (2) 与 (3) 等价。
选取概形 $V$ 以及满光滑态射 $V \to \mathcal{Y}$。
令 $U = \mathcal{X} \times_\mathcal{Y} V$；由假设它是代数空间，
而态射 $U \to \mathcal{X}$ 满且光滑。由《范畴》引理
\ref{categories-lemma-base-change-diagonal} 和《叠的性质》引理
\ref{stacks-properties-lemma-check-property-covering} 可知：
$\Delta_f$ 拟紧当且仅当
$\Delta_{U/V} : U \to  U \times_V U$ 拟紧。
若 (1) 成立，则 $U \to V$ 拟分离，故
$\Delta_{U/V}$ 拟紧，即 (2) 成立。
设 (2) 成立。设 $T$ 是概形，且 $a : T \to \mathcal{Y}$
是态射。令 $T' = \mathcal{X} \times_\mathcal{Y} T$。
为证明 (1)，需证明代数空间的态射 $T' \to T$ 拟分离。
再次使用《范畴》引理 \ref{categories-lemma-base-change-diagonal}，
可见 $\Delta_{T'/T}$ 是 $\Delta_f$ 的基变换。
因此假设 (2) 蕴含 $\Delta_{T'/T}$ 拟紧，从而
$T' \to T$ 拟分离，正合所需。
\end{proof}

\begin{lemma}
\label{stacks-morphisms-lemma-representable-locally-separated-diagonal-immersion}
设 $f : \mathcal{X} \to \mathcal{Y}$ 是可由代数空间表示的代数叠态射。
则下列条件等价：
\begin{enumerate}
\item $f$ 局部分离；
\item $\Delta_f$ 是浸入。
\end{enumerate}
\end{lemma}

\begin{proof}
“$f$ 局部分离”和“$\Delta_f$ 是浸入”均指《叠的性质》第
\ref{stacks-properties-section-properties-morphisms} 节定义的概念。
证明略。提示：仿照引理
\ref{stacks-morphisms-lemma-representable-separated-diagonal-closed} 和
\ref{stacks-morphisms-lemma-representable-quasi-separated-diagonal-quasi-compact}
的证明论证。
\end{proof}






\section{分离公理}
\label{stacks-morphisms-section-separated}

\noindent
设 $\mathcal{X} = [U/R]$ 是代数叠的一个表现。则
$\mathcal{X}$ 在 $S$ 上的对角态射所具有的性质，就是态射
$j : R \to U \times_S U$ 所具有的性质。例如，若
$\mathcal{X} = [S/G]$，其中 $G$ 是 $S$ 上代数空间中的某个光滑群，
则 $j$ 是结构态射 $G \to S$。因此，对角态射本身并不自动分离
（这与概形和代数空间的情形不同）。说 $[S/G]$ 在 $S$ 上拟分离，
显然应蕴含 $G \to S$ 拟紧；但若不同时要求态射 $G \to S$ 拟分离，
我们不愿把 $[S/G]$ 称为在 $S$ 上拟分离。换言之，仅要求对角态射
拟紧并不真正符合我们对“拟分离代数叠”的直觉；还应要求对角态射
本身拟分离。

\medskip\noindent
那么“分离代数叠”应如何定义？《空间的态射》引理
\ref{spaces-morphisms-lemma-separated-diagonal-proper}
说明，代数空间分离当且仅当其对角态射为真态射。通常也用这一条件
定义分离代数叠。在上述例子 $[S/G] \to S$ 中，这意味着
$G \to S$ 是真群概形。因此，若 $E$ 是 $k$ 上的椭圆曲线，
则形如 $[\Spec(k)/E]$ 的代数叠在 $k$ 上为真
（此处以后补充参考文献）。在某些情形下，假设对角态射有限可能更自然。

\begin{definition}
\label{stacks-morphisms-definition-separated}
设 $f : \mathcal{X} \to \mathcal{Y}$ 是代数叠的态射。
\begin{enumerate}
\item 若 $\Delta_f$ 非分歧，则称 $f$ 为 {\it DM}\footnote{字母 DM
代表 Deligne--Mumford。若 $f$ 为 DM，则对任意概形 $T$ 和任意态射
$T \to \mathcal{Y}$，纤维积
$\mathcal{X}_T = \mathcal{X} \times_\mathcal{Y} T$
是 $T$ 上对角态射非分歧的代数叠，即 $\mathcal{X}_T$ 为 DM。
这蕴含 $\mathcal{X}_T$ 是 Deligne--Mumford 叠；见定理
\ref{stacks-morphisms-theorem-DM}。换言之，DM 态射的“纤维”是 Deligne--Mumford
叠。这至少为该术语提供了一些动机。}。
\item 若 $\Delta_f$ 局部拟有限，则称 $f$ 为 {\it 拟 DM}
\footnote{若 $f$ 为拟 DM，则 $\mathcal{X} \to \mathcal{Y}$ 的“纤维”
$\mathcal{X}_T$ 为拟 DM。代数叠 $\mathcal{X}$ 为拟 DM，当且仅当
存在概形 $U$ 以及有限表示、满、平坦且局部拟有限的态射
$U \to \mathcal{X}$；见定理 \ref{stacks-morphisms-theorem-quasi-DM}。
注意，这与 Deligne--Mumford 性质相似；后者以存在概形给出的
\'etale 覆盖来定义。}。
\item 若 $\Delta_f$ 为真态射，则称 $f$ {\it 分离}。
\item 若 $\Delta_f$ 拟紧且拟分离，则称 $f$ {\it 拟分离}。
\end{enumerate}
\end{definition}

\noindent
这里用到了 $\Delta_f$ 可由代数空间表示，并借助《叠的性质》第
\ref{stacks-properties-section-properties-morphisms} 节，
使对 $\Delta_f$ 施加这些条件有意义。注意，当 $f$ 可由代数空间表示时，
这些定义与已有概念并不冲突；见引理
\ref{stacks-morphisms-lemma-representable-quasi-separated-diagonal-quasi-compact} 和
\ref{stacks-morphisms-lemma-representable-separated-diagonal-closed}。
还可通过考察高阶对角态射来刻画这些条件；见引理
\ref{stacks-morphisms-lemma-definition-separated}。

\begin{definition}
\label{stacks-morphisms-definition-absolute-separated}
设 $\mathcal{X}$ 是基概形 $S$ 上的代数叠。记
$p : \mathcal{X} \to S$ 为结构态射。
\begin{enumerate}
\item 若 $p : \mathcal{X} \to S$ 为 DM，则称 $\mathcal{X}$
在 $S$ 上为 {\it DM}。
\item 若 $p : \mathcal{X} \to S$ 为拟 DM，则称 $\mathcal{X}$
在 $S$ 上{\it 拟 DM}。
\item 若 $p : \mathcal{X} \to S$ 分离，则称 $\mathcal{X}$
在 $S$ 上{\it 分离}。
\item 若 $p : \mathcal{X} \to S$ 拟分离，则称 $\mathcal{X}$
在 $S$ 上{\it 拟分离}。
\item 若 $\mathcal{X}$ 在 $\Spec(\mathbf{Z})$ 上为
DM\footnote{定理 \ref{stacks-morphisms-theorem-DM} 说明，这等价于
$\mathcal{X}$ 是 Deligne--Mumford 叠。}，则称 $\mathcal{X}$ 为 {\it DM}。
\item 若 $\mathcal{X}$ 在 $\Spec(\mathbf{Z})$ 上拟 DM，
则称 $\mathcal{X}$ {\it 拟 DM}。
\item 若 $\mathcal{X}$ 在 $\Spec(\mathbf{Z})$ 上分离，
则称 $\mathcal{X}$ {\it 分离}。
\item 若 $\mathcal{X}$ 在 $\Spec(\mathbf{Z})$ 上拟分离，
则称 $\mathcal{X}$ {\it 拟分离}。
\end{enumerate}
在最后四个定义中，我们通过《代数叠》定义
\ref{algebraic-definition-viewed-as}，
把 $\mathcal{X}$ 视为 $\Spec(\mathbf{Z})$ 上的代数叠。
\end{definition}

\noindent
因此，每个情形都有一个绝对概念和一个相对于给定基概形的概念
（按照《叠的性质》第 \ref{stacks-properties-section-conventions}
节引入的术语滥用，我们通常省略对基概形的说明）。
引理 \ref{stacks-morphisms-lemma-separated-implies-morphism-separated} 将说明
(1) $\Leftrightarrow$ (5) 且 (2) $\Leftrightarrow$ (6)。
下面花一些篇幅证明这些概念的若干标准结果。

\begin{lemma}
\label{stacks-morphisms-lemma-trivial-implications}
设 $f : \mathcal{X} \to \mathcal{Y}$ 是代数叠的态射。
\begin{enumerate}
\item 若 $f$ 分离，则 $f$ 拟分离。
\item 若 $f$ 为 DM，则 $f$ 拟 DM。
\item 若 $f$ 可由代数空间表示，则 $f$ 为 DM。
\end{enumerate}
\end{lemma}

\begin{proof}
为证 (1)，注意代数空间的真态射拟紧且拟分离；见《空间的态射》定义
\ref{spaces-morphisms-definition-proper}。为证 (2)，注意代数空间的
非分歧态射局部拟有限；见《空间的态射》引理
\ref{spaces-morphisms-lemma-unramified-quasi-finite}。
最后，(3) 由引理 \ref{stacks-morphisms-lemma-properties-diagonal-representable} 得到。
\end{proof}

\begin{lemma}
\label{stacks-morphisms-lemma-base-change-separated}
定义 \ref{stacks-morphisms-definition-separated} 中列出的所有分离公理都在基变换下稳定。
\end{lemma}

\begin{proof}
设 $f : \mathcal{X} \to \mathcal{Y}$ 和
$\mathcal{Y}' \to \mathcal{Y}$ 是代数叠的态射。设
$f' : \mathcal{Y}' \times_\mathcal{Y} \mathcal{X} \to \mathcal{Y}'$
是 $f$ 沿 $\mathcal{Y}' \to \mathcal{Y}$ 的基变换。
则 $\Delta_{f'}$ 是 $\Delta_f$ 沿态射
$\mathcal{X}' \times_{\mathcal{Y}'} \mathcal{X}' \to
\mathcal{X} \times_\mathcal{Y} \mathcal{X}$ 的基变换；见《范畴》引理
\ref{categories-lemma-base-change-diagonal}。
由《叠的性质》第
\ref{stacks-properties-section-properties-morphisms} 节的结果，
定义 \ref{stacks-morphisms-definition-separated} 中对角态射所用的每个性质都在基变换下稳定。
故引理成立。
\end{proof}

\begin{lemma}
\label{stacks-morphisms-lemma-check-separated-covering}
设 $f : \mathcal{X} \to \mathcal{Y}$ 是代数叠的态射。设 $W$ 是代数空间，
且 $W \to \mathcal{Y}$ 满、平坦并局部有限表示。若基变换
$W \times_\mathcal{Y} \mathcal{X} \to W$ 具有定义
\ref{stacks-morphisms-definition-separated} 中的某个分离性质，则 $f$ 也具有该性质。
\end{lemma}

\begin{proof}
记 $g : W \times_\mathcal{Y} \mathcal{X} \to W$ 为该基变换。
则 $\Delta_g$ 是 $\Delta_f$ 沿态射
$q : W \times_\mathcal{Y} (\mathcal{X} \times_\mathcal{Y} \mathcal{X})
\to \mathcal{X} \times_\mathcal{Y} \mathcal{X}$ 的基变换。由于 $q$
是 $W \to \mathcal{Y}$ 的基变换，$q$ 可由代数空间表示，并且满、
平坦且局部有限表示。因此结论由《叠的性质》引理
\ref{stacks-properties-lemma-check-property-weak-covering} 得到。
\end{proof}

\begin{lemma}
\label{stacks-morphisms-lemma-change-of-base-separated}
设 $S$ 是概形。在 $S$ 上拟 DM、在 $S$ 上拟分离或在 $S$ 上分离的性质
（见定义 \ref{stacks-morphisms-definition-absolute-separated}）在更换基概形下稳定；
见《代数叠》定义 \ref{algebraic-definition-change-of-base}。
\end{lemma}

\begin{proof}
由引理 \ref{stacks-morphisms-lemma-base-change-separated} 立即得到。
\end{proof}

\begin{lemma}
\label{stacks-morphisms-lemma-fibre-product-after-map}
设 $f : \mathcal{X} \to \mathcal{Z}$、$g : \mathcal{Y} \to \mathcal{Z}$
和 $\mathcal{Z} \to \mathcal{T}$ 是代数叠的态射。考察诱导态射
$i : \mathcal{X} \times_\mathcal{Z} \mathcal{Y} \to
\mathcal{X} \times_\mathcal{T} \mathcal{Y}$.
则
\begin{enumerate}
\item $i$ 可由代数空间表示，且局部有限型；
\item 若 $\Delta_{\mathcal{Z}/\mathcal{T}}$ 拟分离，则 $i$ 拟分离；
\item 若 $\Delta_{\mathcal{Z}/\mathcal{T}}$ 分离，则 $i$ 分离；
\item 若 $\mathcal{Z} \to \mathcal{T}$ 为 DM，则 $i$ 非分歧；
\item 若 $\mathcal{Z} \to \mathcal{T}$ 拟 DM，则 $i$ 局部拟有限；
\item 若 $\mathcal{Z} \to \mathcal{T}$ 分离，则 $i$ 为真态射；
\item 若 $\mathcal{Z} \to \mathcal{T}$ 拟分离，则 $i$ 拟紧且拟分离。
\end{enumerate}
\end{lemma}

\begin{proof}
下图
$$
\xymatrix{
\mathcal{X} \times_\mathcal{Z} \mathcal{Y} \ar[r]_i \ar[d] &
\mathcal{X} \times_\mathcal{T} \mathcal{Y} \ar[d] \\
\mathcal{Z} \ar[r]^-{\Delta_{\mathcal{Z}/\mathcal{T}}} \ar[r] &
\mathcal{Z} \times_\mathcal{T} \mathcal{Z}
}
$$
是 $2$-纤维积图；见《范畴》引理
\ref{categories-lemma-fibre-product-after-map}。
因此 $i$ 是对角态射 $\Delta_{\mathcal{Z}/\mathcal{T}}$ 的基变换。
故引理由引理 \ref{stacks-morphisms-lemma-properties-diagonal} 以及《叠的性质》第
\ref{stacks-properties-section-properties-morphisms} 节的内容得到。
\end{proof}

\begin{lemma}
\label{stacks-morphisms-lemma-semi-diagonal}
设 $\mathcal{T}$ 是代数叠。设 $g : \mathcal{X} \to \mathcal{Y}$
是 $\mathcal{T}$ 上代数叠的态射。考察 $g$ 的图
$i : \mathcal{X} \to \mathcal{X} \times_\mathcal{T} \mathcal{Y}$。则
\begin{enumerate}
\item $i$ 可由代数空间表示，且局部有限型；
\item 若 $\mathcal{Y} \to \mathcal{T}$ 为 DM，则 $i$ 非分歧；
\item 若 $\mathcal{Y} \to \mathcal{T}$ 拟 DM，则 $i$ 局部拟有限；
\item 若 $\mathcal{Y} \to \mathcal{T}$ 分离，则 $i$ 为真态射；
\item 若 $\mathcal{Y} \to \mathcal{T}$ 拟分离，则 $i$ 拟紧且拟分离。
\end{enumerate}
\end{lemma}

\begin{proof}
这是把引理 \ref{stacks-morphisms-lemma-fibre-product-after-map} 应用于态射
$\mathcal{X} = \mathcal{X} \times_\mathcal{Y} \mathcal{Y} \to
\mathcal{X} \times_\mathcal{T} \mathcal{Y}$ 所得的特殊情形。
\end{proof}

\begin{lemma}
\label{stacks-morphisms-lemma-section-immersion}
设 $f : \mathcal{X} \to \mathcal{T}$ 是代数叠的态射。
设 $s : \mathcal{T} \to \mathcal{X}$ 是态射，使得
$f \circ s$ 与 $\text{id}_\mathcal{T}$ $2$-同构。则
\begin{enumerate}
\item $s$ 可由代数空间表示，且局部有限型；
\item 若 $f$ 为 DM，则 $s$ 非分歧；
\item 若 $f$ 拟 DM，则 $s$ 局部拟有限；
\item 若 $f$ 分离，则 $s$ 为真态射；
\item 若 $f$ 拟分离，则 $s$ 拟紧且拟分离。
\end{enumerate}
\end{lemma}

\begin{proof}
这是把引理 \ref{stacks-morphisms-lemma-semi-diagonal} 应用于
$g = s$ 和 $\mathcal{Y} = \mathcal{T}$ 所得的特殊情形；此时
$i : \mathcal{T} \to \mathcal{T} \times_\mathcal{T} \mathcal{X}$
与 $s$ $2$-同构。
\end{proof}

\begin{lemma}
\label{stacks-morphisms-lemma-composition-separated}
定义 \ref{stacks-morphisms-definition-separated} 中列出的所有分离公理都在态射复合下稳定。
\end{lemma}

\begin{proof}
设 $f : \mathcal{X} \to \mathcal{Y}$ 和
$g : \mathcal{Y} \to \mathcal{Z}$ 是满足所论公理的代数叠态射。
对角态射 $\Delta_{\mathcal{X}/\mathcal{Z}}$ 是复合
$$
\mathcal{X} \longrightarrow
\mathcal{X} \times_\mathcal{Y} \mathcal{X} \longrightarrow
\mathcal{X} \times_\mathcal{Z} \mathcal{X}.
$$
所论分离公理通过要求对角态射具有某个性质 $\mathcal{P}$ 来定义。
由上面的引理 \ref{stacks-morphisms-lemma-fibre-product-after-map}，
第二个箭头也具有该性质。由于具有性质 $\mathcal{P}$ 且可由代数空间
表示的态射之复合仍具有性质 $\mathcal{P}$，引理即得；见
《叠的性质》第 \ref{stacks-properties-section-properties-morphisms} 节
的一般讨论以及《空间的态射》引理
\ref{spaces-morphisms-lemma-composition-unramified},
\ref{spaces-morphisms-lemma-composition-quasi-finite},
\ref{spaces-morphisms-lemma-composition-proper},
\ref{spaces-morphisms-lemma-composition-quasi-compact} 和
\ref{spaces-morphisms-lemma-composition-separated}.
\end{proof}

\begin{lemma}
\label{stacks-morphisms-lemma-separated-over-separated}
设 $f : \mathcal{X} \to \mathcal{Y}$ 是基概形 $S$ 上代数叠的态射。
\begin{enumerate}
\item 若 $\mathcal{Y}$ 在 $S$ 上为 DM 且 $f$ 为 DM，
则 $\mathcal{X}$ 在 $S$ 上为 DM。
\item 若 $\mathcal{Y}$ 在 $S$ 上拟 DM 且 $f$ 拟 DM，
则 $\mathcal{X}$ 在 $S$ 上拟 DM。
\item 若 $\mathcal{Y}$ 在 $S$ 上分离且 $f$ 分离，
则 $\mathcal{X}$ 在 $S$ 上分离。
\item 若 $\mathcal{Y}$ 在 $S$ 上拟分离且 $f$ 拟分离，
则 $\mathcal{X}$ 在 $S$ 上拟分离。
\item 若 $\mathcal{Y}$ 为 DM 且 $f$ 为 DM，则 $\mathcal{X}$ 为 DM。
\item 若 $\mathcal{Y}$ 拟 DM 且 $f$ 拟 DM，则 $\mathcal{X}$ 拟 DM。
\item 若 $\mathcal{Y}$ 分离且 $f$ 分离，则 $\mathcal{X}$ 分离。
\item 若 $\mathcal{Y}$ 拟分离且 $f$ 拟分离，则 $\mathcal{X}$ 拟分离。
\end{enumerate}
\end{lemma}

\begin{proof}
(1)、(2)、(3)、(4) 由引理 \ref{stacks-morphisms-lemma-composition-separated}
和定义 \ref{stacks-morphisms-definition-absolute-separated} 立即得到。
对于 (5)、(6)、(7)、(8)，把 $\mathcal{X}$ 和 $\mathcal{Y}$
视为 $\Spec(\mathbf{Z})$ 上的代数叠，并应用引理
\ref{stacks-morphisms-lemma-composition-separated}。细节略。
\end{proof}

\noindent
下一个引理与代数空间的相应结果略有不同。比较可见《空间的态射》引理
\ref{spaces-morphisms-lemma-compose-after-separated}。

\begin{lemma}
\label{stacks-morphisms-lemma-compose-after-separated}
设 $f : \mathcal{X} \to \mathcal{Y}$ 和
$g : \mathcal{Y} \to \mathcal{Z}$ 是代数叠的态射。
\begin{enumerate}
\item 若 $g \circ f$ 为 DM，则 $f$ 也为 DM。
\item 若 $g \circ f$ 拟 DM，则 $f$ 也拟 DM。
\item 若 $g \circ f$ 分离且 $\Delta_g$ 分离，则 $f$ 分离。
\item 若 $g \circ f$ 拟分离且 $\Delta_g$ 拟分离，则 $f$ 拟分离。
\end{enumerate}
\end{lemma}

\begin{proof}
考察 $g \circ f$ 的对角态射的分解
$$
\mathcal{X} \to
\mathcal{X} \times_\mathcal{Y} \mathcal{X} \to
\mathcal{X} \times_\mathcal{Z} \mathcal{X}
$$
两个态射都可由代数空间表示；见引理
\ref{stacks-morphisms-lemma-properties-diagonal} 和
\ref{stacks-morphisms-lemma-fibre-product-after-map}。
因此，对任意概形 $T$ 和态射
$T \to \mathcal{X} \times_\mathcal{Y} \mathcal{X}$
得到代数空间的态射
$$
A = \mathcal{X} \times_{(\mathcal{X} \times_\mathcal{Z} \mathcal{X})} T
\longrightarrow
B = (\mathcal{X} \times_\mathcal{Y} \mathcal{X})
\times_{(\mathcal{X} \times_\mathcal{Z} \mathcal{X})} T
\longrightarrow
T
$$
若 $g \circ f$ 为 DM（相应地，拟 DM），则复合 $A \to T$
非分歧（相应地，局部拟有限）。因此，由《空间的态射》引理
\ref{spaces-morphisms-lemma-permanence-unramified}
（相应地，《空间的态射》引理
\ref{spaces-morphisms-lemma-permanence-quasi-finite}），
$A \to B$ 非分歧（相应地，局部拟有限）。现考察图
$$
\xymatrix{
A \times_B T \ar[r] \ar[d] & T \ar[d] \\
A \ar[r] \ar[d] & B \ar[r] \ar[d] & T \ar[d] \\
\mathcal{X} \ar[r] &
\mathcal{X} \times_\mathcal{Y} \mathcal{X} \ar[r] &
\mathcal{X} \times_\mathcal{Z} \mathcal{X}
}
$$
其中所有方块均为 $2$-笛卡儿方块，而箭头 $T \to B$ 来自一开始的箭头
$T \to \mathcal{X} \times_\mathcal{Y} \mathcal{X}$。
由《空间的态射》引理
\ref{spaces-morphisms-lemma-base-change-unramified}
（相应地，《空间的态射》引理
\ref{spaces-morphisms-lemma-base-change-quasi-finite}），
箭头 $A \times_B T \to T$ 非分歧（相应地，局部拟有限）。
注意到 $A \times_B T$ 等于
$\mathcal{X} \times_{\mathcal{X} \times_\mathcal{Y} \mathcal{X}} T$
可知 $f$ 为 DM（相应地，拟 DM）。这证明了 (1) 和 (2)。

\medskip\noindent
(4) 的证明。设 $g \circ f$ 拟分离且 $\Delta_g$ 拟分离。考察
$g \circ f$ 的对角态射的分解
$$
\mathcal{X} \to
\mathcal{X} \times_\mathcal{Y} \mathcal{X} \to
\mathcal{X} \times_\mathcal{Z} \mathcal{X}
$$
两个态射都可由代数空间表示，且第二个态射拟分离；见引理
\ref{stacks-morphisms-lemma-properties-diagonal} 和
\ref{stacks-morphisms-lemma-fibre-product-after-map}。
因此，对任意概形 $T$ 和态射
$T \to \mathcal{X} \times_\mathcal{Y} \mathcal{X}$
得到代数空间的态射
$$
A = \mathcal{X} \times_{(\mathcal{X} \times_\mathcal{Z} \mathcal{X})} T
\longrightarrow
B = (\mathcal{X} \times_\mathcal{Y} \mathcal{X})
\times_{(\mathcal{X} \times_\mathcal{Z} \mathcal{X})} T
\longrightarrow
T
$$
其中 $B \to T$ 拟分离。由于假设 $g \circ f$ 拟分离，
复合 $A \to T$ 拟紧且拟分离。因此，由《空间的态射》引理
\ref{spaces-morphisms-lemma-compose-after-separated}，
$A \to B$ 拟分离；由《空间的态射》引理
\ref{spaces-morphisms-lemma-quasi-compact-permanence}，
$A \to B$ 拟紧。故 $f$ 拟分离。

\medskip\noindent
(3) 的证明。设 $g \circ f$ 分离且 $\Delta_g$ 分离。考察
$g \circ f$ 的对角态射的分解
$$
\mathcal{X} \to
\mathcal{X} \times_\mathcal{Y} \mathcal{X} \to
\mathcal{X} \times_\mathcal{Z} \mathcal{X}
$$
两个态射都可由代数空间表示，且第二个态射分离；见引理
\ref{stacks-morphisms-lemma-properties-diagonal} 和
\ref{stacks-morphisms-lemma-fibre-product-after-map}。
因此，对任意概形 $T$ 和态射
$T \to \mathcal{X} \times_\mathcal{Y} \mathcal{X}$
得到代数空间的态射
$$
A = \mathcal{X} \times_{(\mathcal{X} \times_\mathcal{Z} \mathcal{X})} T
\longrightarrow
B = (\mathcal{X} \times_\mathcal{Y} \mathcal{X})
\times_{(\mathcal{X} \times_\mathcal{Z} \mathcal{X})} T
\longrightarrow
T
$$
其中 $B \to T$ 分离。由于我们假设 $g \circ f$ 拟分离，
复合 $A \to T$ 为真态射。因此，由《空间的态射》引理
\ref{spaces-morphisms-lemma-universally-closed-permanence}，
$A \to B$ 为真态射；这意味着 $f$ 分离。
\end{proof}

\begin{lemma}
\label{stacks-morphisms-lemma-separated-implies-morphism-separated}
设 $\mathcal{X}$ 是基概形 $S$ 上的代数叠。
\begin{enumerate}
\item
$\mathcal{X}$ 为 DM $\Leftrightarrow$
$\mathcal{X}$ 在 $S$ 上为 DM。
\item
$\mathcal{X}$ 拟 DM $\Leftrightarrow$
$\mathcal{X}$ 在 $S$ 上拟 DM。
\item 若 $\mathcal{X}$ 分离，则 $\mathcal{X}$ 在 $S$ 上分离。
\item 若 $\mathcal{X}$ 拟分离，则 $\mathcal{X}$ 在 $S$ 上拟分离。
\end{enumerate}
设 $f : \mathcal{X} \to \mathcal{Y}$ 是基概形 $S$ 上代数叠的态射。
\begin{enumerate}
\item[(5)] 若 $\mathcal{X}$ 在 $S$ 上为 DM，则 $f$ 为 DM。
\item[(6)] 若 $\mathcal{X}$ 在 $S$ 上拟 DM，则 $f$ 拟 DM。
\item[(7)] 若 $\mathcal{X}$ 在 $S$ 上分离且
$\Delta_{\mathcal{Y}/S}$ 分离，则 $f$ 分离。
\item[(8)] 若 $\mathcal{X}$ 在 $S$ 上拟分离且
$\Delta_{\mathcal{Y}/S}$ 拟分离，则 $f$ 拟分离。
\end{enumerate}
\end{lemma}

\begin{proof}
(5)、(6)、(7)、(8) 由引理 \ref{stacks-morphisms-lemma-compose-after-separated}
和《空间》定义 \ref{spaces-definition-separated} 立即得到。
为证明 (3) 和 (4)，把 $X$ 和 $Y$ 视为
$\Spec(\mathbf{Z})$ 上的代数叠，并应用引理
\ref{stacks-morphisms-lemma-compose-after-separated}。
类似地，为证明 (1) 和 (2)，把 $\mathcal{X}$ 视为
$\Spec(\mathbf{Z})$ 上的代数叠，并考察态射
$$
\mathcal{X} \longrightarrow
\mathcal{X} \times_S \mathcal{X} \longrightarrow
\mathcal{X} \times_{\Spec(\mathbf{Z})} \mathcal{X}
$$
两个箭头都可由代数空间表示。第二个箭头是浸入
$\Delta_{S/\mathbf{Z}}$ 的基变换，因而非分歧且局部拟有限。
所以该复合非分歧（相应地，局部拟有限）当且仅当第一个箭头
非分歧（相应地，局部拟有限）；见《空间的态射》引理
\ref{spaces-morphisms-lemma-composition-unramified} 和
\ref{spaces-morphisms-lemma-permanence-unramified}
（相应地，《空间的态射》引理
\ref{spaces-morphisms-lemma-composition-quasi-finite} 和
\ref{spaces-morphisms-lemma-permanence-quasi-finite}).
\end{proof}

\begin{lemma}
\label{stacks-morphisms-lemma-properties-covering-imply-diagonal}
设 $\mathcal{X}$ 是代数叠。设 $W$ 是代数空间，且
$f : W \to \mathcal{X}$ 是满、平坦、局部有限表示的态射。
\begin{enumerate}
\item 若 $f$ 非分歧（即为 \'etale；也即 $\mathcal{X}$ 为
Deligne--Mumford），则 $\mathcal{X}$ 为 DM。
\item 若 $f$ 局部拟有限，则 $\mathcal{X}$ 拟 DM。
\end{enumerate}
\end{lemma}

\begin{proof}
注意，若 $f$ 非分歧，则由《空间的态射》引理
\ref{spaces-morphisms-lemma-unramified-flat-lfp-etale}，
该态射为 \'etale。这解释了 (1) 中的括注。
设 $f$ 非分歧（相应地，局部拟有限）。需证明
$\Delta_\mathcal{X} : \mathcal{X} \to \mathcal{X} \times \mathcal{X}$
非分歧（相应地，局部拟有限）。注意
$W \times W \to \mathcal{X} \times \mathcal{X}$ 也
满、平坦且局部有限表示。因此只需证明
$$
W \times_{\mathcal{X} \times \mathcal{X}, \Delta_\mathcal{X}} \mathcal{X}
=
W \times_\mathcal{X} W
\longrightarrow
W \times W
$$
非分歧（相应地，局部拟有限）；见《叠的性质》引理
\ref{stacks-properties-lemma-check-property-covering}。
由假设，态射 $\text{pr}_i : W \times_\mathcal{X} W \to W$
非分歧（相应地，局部拟有限）。因此，由《空间的态射》引理
\ref{spaces-morphisms-lemma-permanence-unramified}
（相应地，《空间的态射》引理
\ref{spaces-morphisms-lemma-permanence-quasi-finite}），
所显示的箭头非分歧（相应地，局部拟有限）。
\end{proof}

\begin{lemma}
\label{stacks-morphisms-lemma-monomorphism-separated}
代数叠的单态射分离且为 DM。代数叠的浸入亦如此。
\end{lemma}

\begin{proof}
若 $f : \mathcal{X} \to \mathcal{Y}$ 是代数叠的单态射，
则 $\Delta_f$ 是同构；见《叠的性质》引理
\ref{stacks-properties-lemma-monomorphism}。
代数空间的同构为真且非分歧，故 $f$ 分离且为 DM。
由于浸入是单态射，第二个断言由第一个断言得到；见《叠的性质》引理
\ref{stacks-properties-lemma-immersion-monomorphism}。
\end{proof}

\begin{lemma}
\label{stacks-morphisms-lemma-separation-properties-residual-gerbe}
设 $\mathcal{X}$ 是代数叠，且 $x \in |\mathcal{X}|$。
假设 $\mathcal{X}$ 在 $x$ 处的剩余胚 $\mathcal{Z}_x$ 存在。
若 $\mathcal{X}$ 为 DM、相应地拟 DM、相应地分离、相应地拟分离，
则 $\mathcal{Z}_x$ 也如此。
\end{lemma}

\begin{proof}
这是因为 $\mathcal{Z}_x \to \mathcal{X}$ 是单态射，故由引理
\ref{stacks-morphisms-lemma-monomorphism-separated}，它为 DM 且分离。
应用引理 \ref{stacks-morphisms-lemma-separated-over-separated} 即得结论。
\end{proof}


















\section{惯性叠}
\label{stacks-morphisms-section-inertia}

\noindent
群胚中的叠的（相对）惯性叠定义于《叠》第
\ref{stacks-section-the-inertia-stack} 节。
纤维范畴框架下的实际构造及其若干性质见《范畴》第
\ref{categories-section-inertia} 节。

\begin{lemma}
\label{stacks-morphisms-lemma-inertia}
设 $\mathcal{X}$ 是代数叠。则惯性叠
$\mathcal{I}_\mathcal{X}$ 也是代数叠。态射
$$
\mathcal{I}_\mathcal{X} \longrightarrow \mathcal{X}
$$
可由代数空间表示，且局部有限型。更一般地，设
$f : \mathcal{X} \to \mathcal{Y}$ 是代数叠的态射。
则相对惯性叠 $\mathcal{I}_{\mathcal{X}/\mathcal{Y}}$ 是代数叠，且态射
$$
\mathcal{I}_{\mathcal{X}/\mathcal{Y}} \longrightarrow \mathcal{X}
$$
可由代数空间表示，且局部有限型。
\end{lemma}

\begin{proof}
由《范畴》引理 \ref{categories-lemma-inertia-fibred-category}，有等价
$$
\mathcal{I}_\mathcal{X} \to
\mathcal{X} \times_{\Delta, \mathcal{X} \times_S \mathcal{X}, \Delta}
\mathcal{X}
\quad\text{且}\quad
\mathcal{I}_{\mathcal{X}/\mathcal{Y}} \to
\mathcal{X}
\times_{\Delta, \mathcal{X} \times_\mathcal{Y} \mathcal{X}, \Delta}
\mathcal{X}
$$
这说明这些惯性叠是代数叠。设态射 $T \to \mathcal{X}$
由 $\mathcal{X}$ 在 $T$ 上的纤维范畴中的对象 $x$ 给出。
于是得到 $2$-纤维积方块
$$
\xymatrix{
\mathit{Isom}_\mathcal{X}(x, x) \ar[d] \ar[r] &
\mathcal{I}_\mathcal{X} \ar[d] \\
T \ar[r]^x & \mathcal{X}
}
$$
这由 $\mathcal{I}_\mathcal{X}$ 的定义立即得到。
由于 $\mathit{Isom}_\mathcal{X}(x, x)$ 总是 $T$ 上局部有限型的
代数空间（见引理 \ref{stacks-morphisms-lemma-isom-locally-finite-type}），
可知 $\mathcal{I}_\mathcal{X} \to \mathcal{X}$
可由代数空间表示，且局部有限型。最后，对相对惯性叠有
$$
\vcenter{
\xymatrix{
\mathit{Isom}_\mathcal{X}(x, x) \ar[d] &
K \ar[l] \ar[d] \ar[r] &
\mathcal{I}_{\mathcal{X}/\mathcal{Y}} \ar[d] \\
\mathit{Isom}_\mathcal{Y}(f(x), f(x)) &
T \ar[l]_-e \ar[r]^x & \mathcal{X}
}
}
$$
其中两个方块都是 $2$-纤维积。这由《范畴》引理
\ref{categories-lemma-relative-inertia-as-fibre-product} 得到。
左竖直箭头是 $T$ 上局部有限型的代数空间之间的态射，因而局部有限型；
见《空间的态射》引理
\ref{spaces-morphisms-lemma-permanence-finite-type}。
故 $K$ 是代数空间，且 $K \to T$ 局部有限型。
这证明了关于相对惯性叠的断言。
\end{proof}

\begin{remark}
\label{stacks-morphisms-remark-inertia-is-group-in-spaces}
设 $f : \mathcal{X} \to \mathcal{Y}$ 是代数叠的态射。
《叠的性质》注
\ref{stacks-properties-remark-representable-over} 已说明，以
$\mathcal{X}$ 为靶、可由代数空间表示的态射
$\mathcal{Z} \to \mathcal{X}$ 所成的 $2$-范畴构成一个范畴。
在此范畴中，$\mathcal{X}/\mathcal{Y}$ 的惯性叠是一个{\it 群对象}。
回忆
$\mathcal{I}_{\mathcal{X}/\mathcal{Y}}$
的对象就是一对 $(x, \alpha)$，其中 $x$ 是 $\mathcal{X}$ 的对象，
而 $\alpha$ 是 $x$ 所在的 $\mathcal{X}$ 纤维范畴中 $x$ 的自同构，
并满足 $f(\alpha) = \text{id}$。复合
$$
c :
\mathcal{I}_{\mathcal{X}/\mathcal{Y}}
\times_\mathcal{X} \mathcal{I}_{\mathcal{X}/\mathcal{Y}}
\longrightarrow
\mathcal{I}_{\mathcal{X}/\mathcal{Y}}
$$
在对象上由规则
$$
((x, \alpha), (x', \alpha'), \beta) \mapsto
(x, \alpha \circ \beta^{-1} \circ \alpha' \circ \beta)
$$
给出。按照我们的纤维积定义，$\beta : x \to x'$ 是纤维范畴中的同构，
故该规则有意义。中性元
$e : \mathcal{X} \to \mathcal{I}_{\mathcal{X}/\mathcal{Y}}$
由函子 $x \mapsto (x, \text{id}_x)$ 给出。
群对象公理成立的证明从略。
\end{remark}

\noindent
设 $f : \mathcal{X} \to \mathcal{Y}$ 是代数叠的态射，并设
$\mathcal{I}_{\mathcal{X}/\mathcal{Y}}$ 为其惯性叠。
设 $T$ 是概形，且 $x$ 是 $\mathcal{X}$ 在 $T$ 上的对象。
令 $y = f(x)$。引理 \ref{stacks-morphisms-lemma-inertia} 的证明已说明：
对任意概形 $T$ 和 $\mathcal{X}$ 在 $T$ 上的对象 $x$，
有群层的正合列
\begin{equation}
\label{stacks-morphisms-equation-exact-sequence-isom}
0 \to
\mathit{Isom}_{\mathcal{X}/\mathcal{Y}}(x, x) \to
\mathit{Isom}_\mathcal{X}(x, x) \to
\mathit{Isom}_\mathcal{Y}(y, y)
\end{equation}
第二项和第三项上的群结构由引理
\ref{stacks-morphisms-lemma-isom-pseudo-torsor-aut} 定义，而该序列赋予第一项意义。
此外，有典范笛卡儿方块
$$
\xymatrix{
\mathit{Isom}_{\mathcal{X}/\mathcal{Y}}(x, x) \ar[d] \ar[r] &
\mathcal{I}_{\mathcal{X}/\mathcal{Y}} \ar[d] \\
T \ar[r]^x & \mathcal{X}
}
$$
事实上，注 \ref{stacks-morphisms-remark-inertia-is-group-in-spaces} 中讨论的
$\mathcal{I}_{\mathcal{X}/\mathcal{Y}}$ 上的群结构诱导出
$\mathit{Isom}_{\mathcal{X}/\mathcal{Y}}(x, x)$ 上的群结构。
这使我们也能对从代数空间到 $\mathcal{X}$ 的态射定义层
$\mathit{Isom}_{\mathcal{X}/\mathcal{Y}}(x, x)$
。下述定义将其形式化。

\begin{definition}
\label{stacks-morphisms-definition-isom}
设 $f : \mathcal{X} \to \mathcal{Y}$ 是代数叠的态射。
设 $Z$ 是代数空间。
\begin{enumerate}
\item 设 $x : Z \to \mathcal{X}$ 是态射。令
$$
\mathit{Isom}_{\mathcal{X}/\mathcal{Y}}(x, x) =
Z \times_{x, \mathcal{X}} \mathcal{I}_{\mathcal{X}/\mathcal{Y}}
$$
通过拉回注 \ref{stacks-morphisms-remark-inertia-is-group-in-spaces} 中讨论的复合法则，
赋予它 $Z$ 上群代数空间的结构。我们有时把
$\mathit{Isom}_{\mathcal{X}/\mathcal{Y}}(x, x)$
称为{\it $x$ 的相对自同构层}。
\item 设 $x_1, x_2 : Z \to \mathcal{X}$ 是态射。令
$y_i = f \circ x_i$。设 $\alpha : y_1 \to y_2$ 是 $2$-态射。
则 $\alpha$ 确定态射
$\Delta^\alpha : Z \to Z \times_{y_1, \mathcal{Y}, y_2} Z$，并令
$$
\mathit{Isom}_{\mathcal{X}/\mathcal{Y}}^\alpha(x_1, x_2) =
(Z \times_{x_1, \mathcal{X}, x_2} Z)
\times_{Z \times_{y_1, \mathcal{Y}, y_2} Z, \Delta^\alpha} Z.
$$
我们有时把
$\mathit{Isom}_{\mathcal{X}/\mathcal{Y}}^\alpha(x_1, x_2)$
称为{\it 从 $x_1$ 到 $x_2$ 的相对同构层}。
\end{enumerate}
若 $\mathcal{Y} = \Spec(\mathbf{Z})$，或更一般地，若 $\mathcal{Y}$
是代数空间，则使用记号
$\mathit{Isom}_\mathcal{X}(x, x)$ 和 $\mathit{Isom}_\mathcal{X}(x_1, x_2)$
，并使用术语{\it $x$ 的自同构层}和{\it 从 $x_1$ 到 $x_2$ 的同构层}。
\end{definition}

\begin{lemma}
\label{stacks-morphisms-lemma-isom-pseudo-torsor-aut-over-space}
设 $f : \mathcal{X} \to \mathcal{Y}$ 是代数叠的态射。
设 $Z$ 是代数空间，且 $x_i : Z \to \mathcal{X}$（$i = 1, 2$）
是态射。则
\begin{enumerate}
\item $\mathit{Isom}_{\mathcal{X}/\mathcal{Y}}(x_2, x_2)$
是 $Z$ 上的群代数空间；
\item 有群的正合列
$$
0 \to \mathit{Isom}_{\mathcal{X}/\mathcal{Y}}(x_2, x_2)
\to \mathit{Isom}_\mathcal{X}(x_2, x_2)
\to \mathit{Isom}_\mathcal{Y}(f \circ x_2, f \circ x_2)
$$
\item 有代数空间的态射
$
\mathit{Isom}_\mathcal{X}(x_1, x_2)
\to \mathit{Isom}_\mathcal{Y}(f \circ x_1, f \circ x_2)
$
使得对任意 $2$-态射 $\alpha : f \circ x_1 \to f \circ x_2$，
得到笛卡儿图
$$
\xymatrix{
\mathit{Isom}_{\mathcal{X}/\mathcal{Y}}^\alpha(x_1, x_2) \ar[d] \ar[r] &
Z \ar[d]^\alpha \\
\mathit{Isom}_\mathcal{X}(x_1, x_2) \ar[r] &
\mathit{Isom}_\mathcal{Y}(f \circ x_1, f \circ x_2)
}
$$
\item 对任意 $2$-态射 $\alpha : f \circ x_1 \to f \circ x_2$，
代数空间 $\mathit{Isom}_{\mathcal{X}/\mathcal{Y}}^\alpha(x_1, x_2)$
是 $Z$ 上 $\mathit{Isom}_{\mathcal{X}/\mathcal{Y}}(x_2, x_2)$ 的伪挠子。
\end{enumerate}
\end{lemma}

\begin{proof}
(1) 由定义 \ref{stacks-morphisms-definition-isom} 得到。
(2) 在 $Z$ 上 \'etale 局部地由正合列
(\ref{stacks-morphisms-equation-exact-sequence-isom}) 得到。展开定义即可看出 (3)。
在 $Z$ 的 \'etale 拓扑下局部地，(4) 归结为引理
\ref{stacks-morphisms-lemma-isom-pseudo-torsor-aut} 的 (2)。
\end{proof}

\begin{lemma}
\label{stacks-morphisms-lemma-cartesian-square-inertia}
设 $\pi : \mathcal{X} \to \mathcal{Y}$ 和
$f : \mathcal{Y}' \to \mathcal{Y}$ 是代数叠的态射。令
$\mathcal{X}' = \mathcal{X} \times_\mathcal{Y} \mathcal{Y}'$。
则下图中的两个方块
$$
\xymatrix{
\mathcal{I}_{\mathcal{X}'/\mathcal{Y}'} \ar[r]
\ar[d]_{
\text{《范畴》，公式}\ (\ref{categories-equation-functorial})
} &
\mathcal{X}' \ar[r]_{\pi'} \ar[d] & \mathcal{Y}' \ar[d]^f \\
\mathcal{I}_{\mathcal{X}/\mathcal{Y}} \ar[r] &
\mathcal{X} \ar[r]^\pi & \mathcal{Y}
}
$$
均为纤维积方块。
\end{lemma}

\begin{proof}
惯性叠 $\mathcal{I}_{\mathcal{X}'/\mathcal{Y}'}$ 定义为诸对
$(x', \alpha')$ 所成的范畴，其中 $x'$ 是 $\mathcal{X}'$ 的对象，
而 $\alpha'$ 是满足 $\pi'(\alpha') = \text{id}$ 的 $x'$ 的自同构；
见《范畴》第 \ref{categories-section-inertia} 节。
设 $x'$ 位于概形 $U$ 上，并映到 $\mathcal{X}$ 的对象 $x$。
由《范畴》引理 \ref{categories-lemma-2-product-categories-over-C}
中 $2$-纤维积的构造，
$x' = (U, x, y', \beta)$，其中 $y'$ 是 $\mathcal{Y}'$ 在 $U$ 上的对象，
而 $\beta$ 是 $\mathcal{Y}$ 在 $U$ 上的纤维范畴中的同构
$\beta : \pi(x) \to f(y')$。
由 $2$-纤维积的构造本身，自同构 $\alpha'$ 是一对
$(\alpha, \gamma)$，其中 $\alpha$ 是 $x$ 在 $U$ 上的自同构，
$\gamma$ 是 $y'$ 在 $U$ 上的自同构，并且
$\alpha$ 与 $\gamma$ 通过 $\beta$ 相容。条件
$\pi'(\alpha') = \text{id}$ 意味着 $\gamma = \text{id}$；
此时 $\alpha,\beta,\gamma$ 相容恰好等价于
$\pi(\alpha) = \text{id}$，即恰好意味着 $(x, \alpha)$ 是
$\mathcal{I}_{\mathcal{X}/\mathcal{Y}}$ 的对象。
由此可见左方块是纤维积方块（略去若干细节）。
\end{proof}

\begin{lemma}
\label{stacks-morphisms-lemma-monomorphism-cartesian-square-inertia}
设 $f : \mathcal{X} \to \mathcal{Y}$ 是代数叠的单态射。则图
$$
\xymatrix{
\mathcal{I}_\mathcal{X} \ar[r] \ar[d] &
\mathcal{X} \ar[d] \\
\mathcal{I}_\mathcal{Y} \ar[r] &
\mathcal{Y}
}
$$
是纤维积方块。
\end{lemma}

\begin{proof}
这由 $f$ 全忠实（见《叠的性质》引理
\ref{stacks-properties-lemma-monomorphism}）以及《范畴》第
\ref{categories-section-inertia} 节中惯性叠的定义立即得到。
具体地，$\mathcal{I}_\mathcal{X}$ 在概形 $T$ 上的对象等同于一对
$(x, \alpha)$，其中 $x$ 是 $\mathcal{X}$ 在 $T$ 上的对象，而
$\alpha : x \to x$ 是 $\mathcal{X}$ 在 $T$ 上的纤维范畴中的态射。
由于 $f$ 全忠实，$\alpha$ 等同于 $\mathcal{Y}$ 在 $T$ 上的
纤维范畴中的态射 $\beta : f(x) \to f(x)$。
因此可把 $\mathcal{I}_\mathcal{X}$ 在 $T$ 上的对象视为三元组
$((y, \beta), x, \gamma)$，其中 $y$ 是 $\mathcal{Y}$ 在 $T$ 上的对象，
$\beta : y \to y$ 属于 $\mathcal{Y}_T$，而
$\gamma : y \to f(x)$ 是 $T$ 上的同构；换言之，这是
$\mathcal{I}_\mathcal{Y} \times_\mathcal{Y} \mathcal{X}$
在 $T$ 上的对象。
\end{proof}

\begin{lemma}
\label{stacks-morphisms-lemma-presentation-inertia}
设 $\mathcal{X}$ 是代数叠。设 $[U/R] \to \mathcal{X}$ 是一个表现。
设 $G/U$ 是与群胚 $(U, R, s, t, c)$ 相关联的稳定子群代数空间。则
$$
\xymatrix{
G \ar[d] \ar[r] & U \ar[d] \\
\mathcal{I}_\mathcal{X} \ar[r] & \mathcal{X}
}
$$
是纤维积图。
\end{lemma}

\begin{proof}
由《空间中的群胚》引理
\ref{spaces-groupoids-lemma-2-cartesian-inertia} 立即得到。
\end{proof}









\section{高阶对角态射}
\label{stacks-morphisms-section-higher-diagonals}

\noindent
设 $f : \mathcal{X} \to \mathcal{Y}$ 是代数叠的态射。在这种情形下，
不仅可以考虑对角态射
$$
\Delta_f : \mathcal{X} \to \mathcal{X} \times_\mathcal{Y} \mathcal{X}
$$
还可以考虑对角态射的对角态射，即态射
$$
\Delta_{\Delta_f} :
\mathcal{X}
\longrightarrow
\mathcal{X} \times_{(\mathcal{X} \times_\mathcal{Y} \mathcal{X})} \mathcal{X}
$$
因此我们有时采用如下术语：称 $\Delta_{f, 0} = f$ 为{\it 第零对角态射}，
称 $\Delta_{f, 1} = \Delta_f$ 为{\it 第一对角态射}，并称
$\Delta_{f, 2} = \Delta_{\Delta_f}$ 为{\it 第二对角态射}。
注意，$\Delta_{f, 1}$ 可由代数空间表示且局部有限型；见引理
\ref{stacks-morphisms-lemma-properties-diagonal}。因此 $\Delta_{f, 2}$ 可表、是单态射、
局部有限型、分离且局部拟有限；见引理
\ref{stacks-morphisms-lemma-properties-diagonal-representable}。

\medskip\noindent
可以用相对惯性叠描述第二对角态射。具体而言，由《范畴》引理
\ref{categories-lemma-inertia-fibred-category}，纤维积
$\mathcal{X}
\times_{(\mathcal{X} \times_\mathcal{Y} \mathcal{X})} \mathcal{X}$
等价于相对惯性叠 $\mathcal{I}_{\mathcal{X}/\mathcal{Y}}$。
此外，在这一等同下，第二对角态射成为相对惯性叠的{\it 中性截面}
$$
\Delta_{f, 2} = e : \mathcal{X} \to \mathcal{I}_{\mathcal{X}/\mathcal{Y}}
$$
。类比代数空间中群胚的情形（《空间中的群胚》引理
\ref{spaces-groupoids-lemma-diagonal}），可得如下等价关系。

\begin{lemma}
\label{stacks-morphisms-lemma-diagonal-diagonal}
设 $f : \mathcal{X} \to \mathcal{Y}$ 是代数叠的态射。
\begin{enumerate}
\item
下列条件等价：
\begin{enumerate}
\item $\mathcal{I}_{\mathcal{X}/\mathcal{Y}} \to \mathcal{X}$
分离；
\item $\Delta_{f, 1} = \Delta_f :
\mathcal{X} \to \mathcal{X} \times_\mathcal{Y} \mathcal{X}$
分离；以及
\item $\Delta_{f, 2} = e :
\mathcal{X} \to \mathcal{I}_{\mathcal{X}/\mathcal{Y}}$
是闭浸入。
\end{enumerate}
\item
下列条件等价：
\begin{enumerate}
\item $\mathcal{I}_{\mathcal{X}/\mathcal{Y}} \to \mathcal{X}$
拟分离；
\item $\Delta_{f, 1} = \Delta_f :
\mathcal{X} \to \mathcal{X} \times_\mathcal{Y} \mathcal{X}$
拟分离；以及
\item $\Delta_{f, 2} = e :
\mathcal{X} \to \mathcal{I}_{\mathcal{X}/\mathcal{Y}}$
拟紧。
\end{enumerate}
\item
下列条件等价：
\begin{enumerate}
\item $\mathcal{I}_{\mathcal{X}/\mathcal{Y}} \to \mathcal{X}$
局部分离；
\item $\Delta_{f, 1} = \Delta_f :
\mathcal{X} \to \mathcal{X} \times_\mathcal{Y} \mathcal{X}$
局部分离；以及
\item $\Delta_{f, 2} = e :
\mathcal{X} \to \mathcal{I}_{\mathcal{X}/\mathcal{Y}}$
是浸入。
\end{enumerate}
\item
下列条件等价：
\begin{enumerate}
\item $\mathcal{I}_{\mathcal{X}/\mathcal{Y}} \to \mathcal{X}$
非分歧；
\item $f$ 是 DM。
\end{enumerate}
\item
下列条件等价：
\begin{enumerate}
\item $\mathcal{I}_{\mathcal{X}/\mathcal{Y}} \to \mathcal{X}$
局部拟有限；
\item $f$ 是拟 DM。
\end{enumerate}
\end{enumerate}
\end{lemma}

\begin{proof}
先证明 (1)、(2) 和 (3)。取代数空间 $U$ 以及满光滑态射
$U \to \mathcal{X}$。则
$G = U \times_\mathcal{X} \mathcal{I}_{\mathcal{X}/\mathcal{Y}}$
是 $U$ 上的代数空间（引理 \ref{stacks-morphisms-lemma-inertia}）。事实上，由注记
\ref{stacks-morphisms-remark-inertia-is-group-in-spaces} 中在相对惯性上构造的群律，
$G$ 是 $U$ 上的群代数空间。此外，作为 $U \to \mathcal{X}$ 的基变换，
$G \to \mathcal{I}_{\mathcal{X}/\mathcal{Y}}$ 满且光滑。最后，态射
$e : \mathcal{X} \to \mathcal{I}_{\mathcal{X}/\mathcal{Y}}$
沿 $G \to \mathcal{I}_{\mathcal{X}/\mathcal{Y}}$ 的基变换是 $G/U$ 的
单位截面 $U \to G$。因此，(a) 与 (c) 的等价性由《空间中的群胚》引理
\ref{spaces-groupoids-lemma-group-scheme-separated} 得到。
由于 $\Delta_{f, 2}$ 是 $\Delta_f$ 的对角态射，根据定义有
(b) $\Leftrightarrow$ (c)。

\medskip\noindent
再证明 (4) 和 (5)。回顾，(4)(b) 意味着 $\Delta_f$ 非分歧，
(5)(b) 意味着 $\Delta_f$ 局部拟有限。取概形 $Z$ 以及态射
$a : Z \to \mathcal{X} \times_\mathcal{Y} \mathcal{X}$。于是
$a = (x_1, x_2, \alpha)$，其中 $x_i : Z \to \mathcal{X}$，而
$\alpha : f \circ x_1  \to f \circ x_2$ 是 $2$-态射。回顾
$$
\vcenter{
\xymatrix{
\mathit{Isom}_{\mathcal{X}/\mathcal{Y}}^\alpha(x_1, x_2)
\ar[d] \ar[r] &
Z \ar[d] \\
\mathcal{X} \ar[r]^{\Delta_f} &
\mathcal{X} \times_\mathcal{Y} \mathcal{X}
}
}
\quad\text{且}\quad
\vcenter{
\xymatrix{
\mathit{Isom}_{\mathcal{X}/\mathcal{Y}}(x_2, x_2)
\ar[d] \ar[r] &
Z \ar[d]^{x_2} \\
\mathcal{I}_{\mathcal{X}/\mathcal{Y}} \ar[r] &
\mathcal{X}
}
}
$$
是笛卡尔方块。由引理 \ref{stacks-morphisms-lemma-isom-pseudo-torsor-aut-over-space}，
代数空间 $\mathit{Isom}_{\mathcal{X}/\mathcal{Y}}^\alpha(x_1, x_2)$
是 $\mathit{Isom}_{\mathcal{X}/\mathcal{Y}}(x_2, x_2)$ 在 $Z$ 上的伪挠子。
因此，(4) 和 (5) 中的等价性由《空间中的群胚》引理
\ref{spaces-groupoids-lemma-pseudo-torsor-implications} 得到。
\end{proof}

\begin{lemma}
\label{stacks-morphisms-lemma-second-diagonal}
设 $f : \mathcal{X} \to \mathcal{Y}$ 是代数叠的态射。下列条件等价：
\begin{enumerate}
\item 态射 $f$ 可由代数空间表示；
\item $f$ 的第二对角态射是同构；
\item 群叠 $ \mathcal{I}_{\mathcal{X}/\mathcal{Y}}$
在 $\mathcal X$ 上平凡；以及
\item 对概形 $T$ 和态射 $x : T \to \mathcal{X}$，
$\mathit{Isom}_\mathcal{X}(x, x) \to
\mathit{Isom}_\mathcal{Y}(f(x), f(x))$ 的核平凡。
\end{enumerate}
\end{lemma}

\begin{proof}
先证明 (1) 与 (2) 的等价性。事实上，$f$ 可由代数空间表示当且仅当
$f$ 忠实；见《代数叠》引理
\ref{algebraic-lemma-characterize-representable-by-algebraic-spaces}。
另一方面，$f$ 忠实当且仅当对概形 $T$ 上 $\mathcal{X}$ 的每个对象 $x$，
函子 $f$ 诱导单射
$\mathit{Isom}_\mathcal{X}(x, x) \to
\mathit{Isom}_\mathcal{Y}(f(x), f(x))$；这又当且仅当其核 $K$ 平凡，
亦即当且仅当对每个 $x : T \to \mathcal{X}$，$e : T \to K$ 是同构。
如上所述，$K = T \times_{x, \mathcal{X}} \mathcal{I}_{\mathcal{X}/\mathcal{Y}}$，
故这证明了 (1) 与 (2) 的等价性。为证明 (2) 与 (3) 的等价性，根据上述讨论，
只须注意群叠平凡当且仅当其单位截面是同构。最后，(3) 与 (4) 的等价性由定义
得到：在引理 \ref{stacks-morphisms-lemma-inertia} 的证明中，我们已经看到 (4) 中的核对应于
$T$ 上的纤维积
$T \times_{x, \mathcal{X}} \mathcal{I}_{\mathcal{X}/\mathcal{Y}}$。
\end{proof}

\noindent
该引理给出代数叠态射的如下层级关系。

\begin{lemma}
\label{stacks-morphisms-lemma-hierarchy}
代数叠的态射 $f : \mathcal{X} \to \mathcal{Y}$ 满足：
\begin{enumerate}
\item 该态射是单态射当且仅当 $\Delta_{f, 1}$ 是同构；并且
\item 该态射可由代数空间表示当且仅当 $\Delta_{f, 1}$ 是单态射。
\end{enumerate}
此外，第二对角态射 $\Delta_{f, 2}$ 总是单态射。
\end{lemma}

\begin{proof}
回顾《叠的性质》引理 \ref{stacks-properties-lemma-monomorphism}：
代数叠的态射是单态射，当且仅当其对角态射是叠的同构。因此，引理
\ref{stacks-morphisms-lemma-second-diagonal} 可以改述为：若一个态射的对角态射是单态射，
则该态射可由代数空间表示。特别地，这说明引理
\ref{stacks-morphisms-lemma-properties-diagonal-representable} 的条件 (3) 实际上是充要条件；
也就是说，代数叠的态射可由代数空间表示，当且仅当其对角态射是单态射。
\end{proof}

\begin{lemma}
\label{stacks-morphisms-lemma-first-diagonal-separated-second-diagonal-closed}
设 $f : \mathcal{X} \to \mathcal{Y}$ 是代数叠的态射。则
\begin{enumerate}
\item $\Delta_{f, 1}$ 分离 $\Leftrightarrow$
$\Delta_{f, 2}$ 是闭浸入 $\Leftrightarrow$
$\Delta_{f, 2}$ 为真 $\Leftrightarrow$
$\Delta_{f, 2}$ 泛闭；
\item $\Delta_{f, 1}$ 拟分离 $\Leftrightarrow$
$\Delta_{f, 2}$ 有限型 $\Leftrightarrow$ $\Delta_{f, 2}$ 拟紧；以及
\item $\Delta_{f, 1}$ 局部分离 $\Leftrightarrow$
$\Delta_{f, 2}$ 是浸入。
\end{enumerate}
\end{lemma}

\begin{proof}
将引理 \ref{stacks-morphisms-lemma-representable-separated-diagonal-closed}、
\ref{stacks-morphisms-lemma-representable-quasi-separated-diagonal-quasi-compact} 和
\ref{stacks-morphisms-lemma-representable-locally-separated-diagonal-immersion}
应用于 $\Delta_{f, 1}$ 即得。
\end{proof}

\noindent
下面这个引理颇为巧妙，并可能启发我们把这些条件推广到高阶代数叠。

\begin{lemma}
\label{stacks-morphisms-lemma-definition-separated}
设 $f : \mathcal{X} \to \mathcal{Y}$ 是代数叠的态射。则
\begin{enumerate}
\item $f$ 分离当且仅当 $\Delta_{f, 1}$ 和 $\Delta_{f, 2}$ 均泛闭；
\item $f$ 拟分离当且仅当 $\Delta_{f, 1}$ 和 $\Delta_{f, 2}$ 均拟紧；
\item $f$ 是拟 DM 当且仅当 $\Delta_{f, 1}$ 和 $\Delta_{f, 2}$
均局部拟有限；
\item $f$ 是 DM 当且仅当 $\Delta_{f, 1}$ 和 $\Delta_{f, 2}$ 均非分歧。
\end{enumerate}
\end{lemma}

\begin{proof}
证明 (1)。假设 $\Delta_{f, 2}$ 和 $\Delta_{f, 1}$ 均泛闭。由引理
\ref{stacks-morphisms-lemma-first-diagonal-separated-second-diagonal-closed}，
$\Delta_{f, 1}$ 分离且泛闭。再由《空间的态射》引理
\ref{spaces-morphisms-lemma-universally-closed-quasi-compact} 和
《代数叠》引理
\ref{algebraic-lemma-representable-transformations-property-implication}，
可知 $\Delta_{f, 1}$ 拟紧。因此它拟紧、分离、泛闭且局部有限型
（由引理 \ref{stacks-morphisms-lemma-properties-diagonal}），故为真。这证明了 (1) 的
``$\Leftarrow$''。反方向蕴涵的证明略去。

\medskip\noindent
证明 (2)。这立即由引理
\ref{stacks-morphisms-lemma-first-diagonal-separated-second-diagonal-closed} 得到。

\medskip\noindent
证明 (3)。把引理 \ref{stacks-morphisms-lemma-properties-diagonal-representable}
应用于 $\Delta_f = \Delta_{f, 1}$ 可知 $\Delta_{f, 2}$ 总是局部拟有限，
结论由此得到。

\medskip\noindent
证明 (4)。把引理 \ref{stacks-morphisms-lemma-properties-diagonal-representable}
应用于 $\Delta_f = \Delta_{f, 1}$，可知 $\Delta_{f, 2}$ 局部有限型且为
单态射，因而 $\Delta_{f, 2}$ 总是非分歧；结论由此得到。见《空间态射进阶》引理
\ref{spaces-more-morphisms-lemma-universally-injective-unramified}。
\end{proof}

\begin{lemma}
\label{stacks-morphisms-lemma-separated-implies-isom}
设 $f : \mathcal{X} \to \mathcal{Y}$ 是代数叠的分离
（相应地，拟分离、拟 DM、DM）态射。则
\begin{enumerate}
\item 给定代数空间 $T_i$（$i = 1, 2$）以及态射
$x_i : T_i \to \mathcal{X}$，令 $y_i = f \circ x_i$，则态射
$$
T_1 \times_{x_1, \mathcal{X}, x_2} T_2 \longrightarrow
T_1 \times_{y_1, \mathcal{Y}, y_2} T_2
$$
为真（相应地，拟紧且拟分离、局部拟有限、非分歧）；
\item 给定代数空间 $T$ 以及态射
$x_i : T \to \mathcal{X}$（$i = 1, 2$），令 $y_i = f \circ x_i$，则态射
$$
\mathit{Isom}_\mathcal{X}(x_1, x_2) \longrightarrow
\mathit{Isom}_\mathcal{Y}(y_1, y_2)
$$
为真（相应地，拟紧且拟分离、局部拟有限、非分歧）。
\end{enumerate}
\end{lemma}

\begin{proof}
证明 (1)。注意，图
$$
\xymatrix{
T_1 \times_{x_1, \mathcal{X}, x_2} T_2 \ar[d] \ar[r] &
T_1 \times_{y_1, \mathcal{Y}, y_2} T_2 \ar[d] \\
\mathcal{X} \ar[r] & \mathcal{X} \times_\mathcal{Y} \mathcal{X}
}
$$
是笛卡尔图。因此，结论由下述事实得到：$f$ 分离
（相应地，拟分离、拟 DM、DM）当且仅当其对角态射为真
（相应地，拟紧且拟分离、局部拟有限、非分歧）。

\medskip\noindent
证明 (2)。这是因为
$$
\mathit{Isom}_\mathcal{X}(x_1, x_2) =
(T \times_{x_1, \mathcal{X}, x_2} T) \times_{T \times T, \Delta_T} T
$$
，故 (2) 中的态射是 (1) 中态射的基变换。
\end{proof}








\section{拟紧态射}
\label{stacks-morphisms-section-quasi-compact}

\noindent
设 $f$ 是可由代数空间表示的代数叠态射。我们已在《叠的性质》第
\ref{stacks-properties-section-properties-morphisms} 节定义了 $f$ 拟紧的含义。
下面给出另一种刻画。

\begin{lemma}
\label{stacks-morphisms-lemma-characterize-representable-quasi-compact}
设 $f : \mathcal{X} \to \mathcal{Y}$ 是可由代数空间表示的代数叠态射。
下列条件等价：
\begin{enumerate}
\item $f$ 拟紧（定义见《叠的性质》第
\ref{stacks-properties-section-properties-morphisms} 节）；以及
\item 对每个拟紧代数叠 $\mathcal{Z}$ 及任意态射
$\mathcal{Z} \to \mathcal{Y}$，代数叠
$\mathcal{Z} \times_\mathcal{Y} \mathcal{X}$ 拟紧。
\end{enumerate}
\end{lemma}

\begin{proof}
假设 (1)，并设 $\mathcal{Z} \to \mathcal{Y}$ 是代数叠的态射，其中
$\mathcal{Z}$ 拟紧。由《叠的性质》引理
\ref{stacks-properties-lemma-quasi-compact-stack}，存在拟紧概形 $U$
及满光滑态射 $U \to \mathcal{Z}$。由于 $f$ 可由代数空间表示且拟紧，
根据定义，$U \times_\mathcal{Y} \mathcal{X}$ 是代数空间，并且
$U \times_\mathcal{Y} \mathcal{X} \to U$ 拟紧。因此
$U \times_\mathcal{Y} \mathcal{X}$ 是拟紧代数空间。态射
$U \times_\mathcal{Y} \mathcal{X} \to
\mathcal{Z} \times_\mathcal{Y} \mathcal{X}$
光滑且满（它是光滑满态射 $U \to \mathcal{Z}$ 的基变换）。
再次应用《叠的性质》引理
\ref{stacks-properties-lemma-quasi-compact-stack}，可知
$\mathcal{Z} \times_\mathcal{Y} \mathcal{X}$ 拟紧

\medskip\noindent
假设 (2)。设 $Z \to \mathcal{Y}$ 是态射，其中 $Z$ 是概形。
我们必须证明代数空间的态射
$p : Z \times_\mathcal{Y} \mathcal{X} \to Z$ 拟紧。
设 $U \subset Z$ 是仿射开子概形。则
$p^{-1}(U) = U \times_\mathcal{Y} Z$
，而代数空间 $U \times_\mathcal{Y} Z$ 由假设 (2) 拟紧。因此 $p$ 拟紧；
见《空间的态射》引理 \ref{spaces-morphisms-lemma-quasi-compact-local}。
\end{proof}

\noindent
这引出如下定义。

\begin{definition}
\label{stacks-morphisms-definition-quasi-compact}
设 $f : \mathcal{X} \to \mathcal{Y}$ 是代数叠的态射。若对每个拟紧代数叠
$\mathcal{Z}$ 及态射 $\mathcal{Z} \to \mathcal{Y}$，纤维积
$\mathcal{Z} \times_\mathcal{Y} \mathcal{X}$ 都拟紧，则称 $f$ {\it 拟紧}。
\end{definition}

\noindent
由上面的引理 \ref{stacks-morphisms-lemma-characterize-representable-quasi-compact}，
对可由代数空间表示的代数叠态射，这与已有概念一致。特别地，
这一概念与已经为代数叠和概形之间的态射定义的概念一致。

\begin{lemma}
\label{stacks-morphisms-lemma-base-change-quasi-compact}
代数叠的拟紧态射沿任意代数叠态射的基变换仍拟紧。
\end{lemma}

\begin{proof}
略。
\end{proof}

\begin{lemma}
\label{stacks-morphisms-lemma-composition-quasi-compact}
两个代数叠拟紧态射的复合拟紧。
\end{lemma}

\begin{proof}
略。
\end{proof}

\begin{lemma}
\label{stacks-morphisms-lemma-closed-immersion-quasi-compact}
代数叠的闭浸入拟紧。
\end{lemma}

\begin{proof}
这由浸入总是可表以及代数空间闭浸入的相应事实得到。
\end{proof}

\begin{lemma}
\label{stacks-morphisms-lemma-surjection-from-quasi-compact}
设
$$
\xymatrix{
\mathcal{X} \ar[rr]_f \ar[rd]_p & &
\mathcal{Y} \ar[dl]^q \\
& \mathcal{Z}
}
$$
是代数叠态射的 $2$-交换图。若 $f$ 满且 $p$ 拟紧，则 $q$ 拟紧。
\end{lemma}

\begin{proof}
设 $\mathcal{T}$ 是拟紧代数叠，并设 $\mathcal{T} \to \mathcal{Z}$ 是态射。
由《叠的性质》引理
\ref{stacks-properties-lemma-base-change-surjective}，态射
$\mathcal{T} \times_\mathcal{Z} \mathcal{X} \to
\mathcal{T} \times_\mathcal{Z} \mathcal{Y}$
满，而根据假设，$\mathcal{T} \times_\mathcal{Z} \mathcal{X}$ 拟紧。
因此，由《叠的性质》引理
\ref{stacks-properties-lemma-quasi-compact-stack}，
$\mathcal{T} \times_\mathcal{Z} \mathcal{Y}$ 拟紧。
\end{proof}

\begin{lemma}
\label{stacks-morphisms-lemma-quasi-compact-permanence}
设 $f : \mathcal{X} \to \mathcal{Y}$ 和
$g : \mathcal{Y} \to \mathcal{Z}$ 是代数叠的态射。若 $g \circ f$ 拟紧且
$g$ 拟分离，则 $f$ 拟紧。
\end{lemma}

\begin{proof}
这是因为 $f$ 等于如下复合：
$(1, f) : \mathcal{X} \to \mathcal{X} \times_\mathcal{Z} \mathcal{Y} \to
\mathcal{Y}$.
第一个态射是拟分离态射
$\mathcal{X} \times_\mathcal{Z} \mathcal{Y} \to \mathcal{X}$
的截面，故由引理 \ref{stacks-morphisms-lemma-section-immersion} 拟紧；该拟分离态射是 $g$
的基变换，见引理 \ref{stacks-morphisms-lemma-base-change-separated}。第二个态射是 $f$ 的
基变换，故拟紧；见引理 \ref{stacks-morphisms-lemma-base-change-quasi-compact}。
拟紧态射的复合拟紧；见引理 \ref{stacks-morphisms-lemma-composition-quasi-compact}。
\end{proof}

\begin{lemma}
\label{stacks-morphisms-lemma-quasi-compact-quasi-separated-permanence}
设 $f : \mathcal{X} \to \mathcal{Y}$ 是代数叠的态射。
\begin{enumerate}
\item 若 $\mathcal{X}$ 拟紧且 $\mathcal{Y}$ 拟分离，则 $f$ 拟紧。
\item 若 $\mathcal{X}$ 拟紧且拟分离，并且 $\mathcal{Y}$ 拟分离，
则 $f$ 拟紧且拟分离。
\item 拟紧且拟分离的代数叠的纤维积拟紧且拟分离。
\end{enumerate}
\end{lemma}

\begin{proof}
(1) 由引理 \ref{stacks-morphisms-lemma-quasi-compact-permanence} 得到。
(2) 由 (1) 和引理 \ref{stacks-morphisms-lemma-compose-after-separated} 得到。
对 (3)，设 $\mathcal{X} \to \mathcal{Y}$ 和 $\mathcal{Z} \to \mathcal{Y}$
是拟紧且拟分离代数叠的态射。利用 (2) 以及引理
\ref{stacks-morphisms-lemma-base-change-quasi-compact} 和
\ref{stacks-morphisms-lemma-base-change-separated}，态射
$\mathcal{X} \times_\mathcal{Y} \mathcal{Z} \to \mathcal{Z}$
作为 $\mathcal{X} \to \mathcal{Y}$ 的基变换拟紧且拟分离。因此，
$\mathcal{X} \times_\mathcal{Y} \mathcal{Z}$ 是 $\mathcal{Z}$ 上拟紧且
拟分离的代数叠，从而拟紧且拟分离；见引理
\ref{stacks-morphisms-lemma-separated-over-separated} 和
\ref{stacks-morphisms-lemma-composition-quasi-compact}。
\end{proof}

\begin{lemma}
\label{stacks-morphisms-lemma-reach-points-scheme-theoretic-image}
设 $f : \mathcal{X} \to \mathcal{Y}$ 是代数叠的拟紧态射。设
$y \in |\mathcal{Y}|$ 是 $|f|$ 的像的闭包中的一点。则存在以 $K$ 为分式域的
赋值环 $A$ 以及交换图
$$
\xymatrix{
\Spec(K) \ar[r] \ar[d] & \mathcal{X} \ar[d] \\
\Spec(A) \ar[r] & \mathcal{Y}
}
$$
使得 $\Spec(A)$ 的闭点映到 $y$。
\end{lemma}

\begin{proof}
取仿射概形 $V$、点 $v \in V$ 以及把 $v$ 映到 $y$ 的光滑态射
$V \to \mathcal{Y}$。考虑基变换图
$$
\xymatrix{
V \times_\mathcal{Y} \mathcal{X} \ar[r] \ar[d]_g & \mathcal{X} \ar[d]^f \\
V \ar[r] & \mathcal{Y}
}
$$
回顾，$|V \times_\mathcal{Y} \mathcal{X}| \to
|V| \times_{|\mathcal{Y}|} |\mathcal{X}|$ 满（《叠的性质》引理
\ref{stacks-properties-lemma-points-cartesian}）。由于
$|V| \to |\mathcal{Y}|$ 是开映射（《叠的性质》引理
\ref{stacks-properties-lemma-topology-points}），可知 $v$ 属于 $|g|$ 的像的闭包。
因此，只须对拟紧态射 $g$ 证明此引理（引理
\ref{stacks-morphisms-lemma-base-change-quasi-compact}）；下一段完成这一点。

\medskip\noindent
假设 $\mathcal{Y} = Y$ 是仿射概形。由于 $f$ 拟紧，$\mathcal{X}$ 拟紧
（定义 \ref{stacks-morphisms-definition-quasi-compact}）。取仿射概形 $W$ 以及满光滑态射
$W \to \mathcal{X}$。于是 $|f|$ 的像就是 $W \to Y$ 的像。
由《态射》引理 \ref{morphisms-lemma-reach-points-scheme-theoretic-image}，
可取图
$$
\xymatrix{
\Spec(K) \ar[r] \ar[d] & W \ar[d] \ar[r] & \mathcal{X} \ar[d] \\
\Spec(A) \ar[r] & Y \ar[r] & Y
}
$$
使得 $\Spec(A)$ 的闭点映到 $y$。再与 $W \to \mathcal{X}$ 复合即得所求。
\end{proof}

\begin{lemma}
\label{stacks-morphisms-lemma-check-quasi-compact-covering}
设 $f : \mathcal{X} \to \mathcal{Y}$ 是代数叠的态射。设
$W \to \mathcal{Y}$ 满、平坦且局部有限表示，其中 $W$ 是代数空间。
若基变换 $W \times_\mathcal{Y} \mathcal{X} \to W$ 拟紧，则 $f$ 拟紧。
\end{lemma}

\begin{proof}
假设 $W \times_\mathcal{Y} \mathcal{X} \to W$ 拟紧。设
$\mathcal{Z} \to \mathcal{Y}$ 是态射，其中 $\mathcal{Z}$ 是拟紧代数叠。
取概形 $U$ 以及满光滑态射
$U \to W \times_\mathcal{Y} \mathcal{Z}$。由于
$U \to \mathcal{Z}$ 平坦、满且局部有限表示，而 $\mathcal{Z}$ 拟紧，
可以找到拟紧开子概形 $U' \subset U$，使 $U' \to \mathcal{Z}$ 满。于是
$U' \times_\mathcal{Y} \mathcal{X} =
U' \times_W (W \times_\mathcal{Y} \mathcal{X})$
由假设拟紧，并满射到 $\mathcal{Z} \times_\mathcal{Y} \mathcal{X}$。
因此 $\mathcal{Z} \times_\mathcal{Y} \mathcal{X}$ 拟紧，正合所需。
\end{proof}











\section{诺特代数叠}
\label{stacks-morphisms-section-noetherian}

\noindent
我们已在《叠的性质》第 \ref{stacks-properties-section-types-properties} 节
定义了局部诺特代数叠。

\begin{definition}
\label{stacks-morphisms-definition-noetherian}
设 $\mathcal{X}$ 是代数叠。若 $\mathcal{X}$ 拟紧、拟分离且局部诺特，
则称 $\mathcal{X}$ {\it 诺特}。
\end{definition}

\noindent
注意，诺特代数叠 $\mathcal{X}$ 不仅拟紧且局部诺特，而且还拟分离。
用第 \ref{stacks-morphisms-section-higher-diagonals} 节的语言，若把“绝对”结构态射
$p : \mathcal{X} \to \Spec(\mathbf{Z})$（即把 $\mathcal{X}$ 看作
$\mathbf{Z}$ 上代数叠时的结构态射），则
$$
\mathcal{X}\text{ Noetherian}
\Leftrightarrow
\mathcal{X}\text{ locally Noetherian and }
\Delta_{p, 0}, \Delta_{p, 1}, \Delta_{p, 2}
\text{ quasi-compact}.
$$
这在后面将意味着：诺特代数叠上的有限型代数叠并不自动是诺特的。

\begin{lemma}
\label{stacks-morphisms-lemma-locally-closed-in-noetherian}
设 $j : \mathcal{X} \to \mathcal{Y}$ 是代数叠的浸入。
\begin{enumerate}
\item 若 $\mathcal{Y}$ 局部诺特，则 $\mathcal{X}$ 局部诺特且 $j$ 拟紧。
\item 若 $\mathcal{Y}$ 诺特，则 $\mathcal{X}$ 诺特。
\end{enumerate}
\end{lemma}

\begin{proof}
取概形 $V$ 以及满光滑态射 $V \to \mathcal{Y}$。则
$U = \mathcal{X} \times_\mathcal{Y} V$ 是概形，而 $V \to U$ 是浸入；
见《叠的性质》定义 \ref{stacks-properties-definition-immersion}。
回顾，$\mathcal{Y}$ 局部诺特当且仅当 $V$ 局部诺特。在这种情形下，
$U$ 也局部诺特（《态射》引理
\ref{morphisms-lemma-immersion-locally-finite-type} 和
\ref{morphisms-lemma-finite-type-noetherian}），且 $U \to V$ 拟紧
（《性质》引理 \ref{properties-lemma-immersion-into-noetherian}）。
这表明 $j$ 拟紧（引理 \ref{stacks-morphisms-lemma-check-quasi-compact-covering}），并且
$\mathcal{X}$ 局部诺特。最后，若 $\mathcal{Y}$ 诺特，则由上述论证可知
$\mathcal{X}$ 拟紧且局部诺特。为完成证明，注意 $j$ 分离；又因为
$\mathcal{Y}$ 拟分离，引理 \ref{stacks-morphisms-lemma-separated-over-separated} 表明
$\mathcal{X}$ 拟分离。
\end{proof}

\begin{lemma}
\label{stacks-morphisms-lemma-Noetherian-topology}
设 $\mathcal{X}$ 是代数叠。
\begin{enumerate}
\item 若 $\mathcal{X}$ 局部诺特，则 $|\mathcal{X}|$ 是局部诺特拓扑空间。
\item 若 $\mathcal{X}$ 拟紧且局部诺特，则 $|\mathcal{X}|$ 是诺特拓扑空间。
\end{enumerate}
\end{lemma}

\begin{proof}
假设 $\mathcal{X}$ 局部诺特。取概形 $U$ 以及满光滑态射
$U \to \mathcal{X}$。由于 $\mathcal{X}$ 局部诺特，$U$ 也局部诺特。
由《性质》引理 \ref{properties-lemma-Noetherian-topology}，这意味着
$|U|$ 是局部诺特拓扑空间。由于 $|U| \to |\mathcal{X}|$ 开且满，
由《拓扑》引理 \ref{topology-lemma-image-Noetherian}，
$|\mathcal{X}|$ 局部诺特。这证明了 (1)。若 $\mathcal{X}$ 拟紧且局部诺特，
则 $|\mathcal{X}|$ 拟紧且局部诺特。因此，由《拓扑》引理
\ref{topology-lemma-quasi-compact-locally-Noetherian-Noetherian}，
$|\mathcal{X}|$ 诺特。
\end{proof}

\begin{lemma}
\label{stacks-morphisms-lemma-locally-Noetherian-quasi-sober}
设 $\mathcal{X}$ 是局部诺特代数叠。则 $|\mathcal{X}|$ 拟清醒
（《拓扑》定义 \ref{topology-definition-generic-point}）。
\end{lemma}

\begin{proof}
我们必须证明每个不可约闭子集 $T \subset |\mathcal{X}|$ 都有一般点。
取仿射概形 $U$ 以及光滑态射 $f : U \to \mathcal{X}$，使
$f^{-1}(T) \subset |U|$ 非空。由于 $U$ 诺特，闭子集 $f^{-1}(T)$
只有有限多个不可约分支（《拓扑》引理 \ref{topology-lemma-Noetherian}）。
设 $f^{-1}(T) = Z_1 \cup \ldots \cup Z_n$ 是其不可约分支分解。
由于 $f$ 是开映射，$f|_{f^{-1}(T)} : f^{-1}(T) \to T$ 的像包含 $T$
的一个非空开子集。由于 $T$ 不可约，这意味着 $f(f^{-1}(T))$ 稠密。
仍由 $T$ 不可约，存在某个 $i$ 使 $f(Z_i)$ 稠密。若
$\xi_i \in Z_i$ 是一般点，则 $f(\xi_i)$ 是 $T$ 的一般点。
\end{proof}





\section{仿射态射}
\label{stacks-morphisms-section-affine}

\noindent
代数叠的仿射态射定义如下。

\begin{definition}
\label{stacks-morphisms-definition-affine}
若代数叠的态射可表，并且在《叠的性质》第
\ref{stacks-properties-section-properties-morphisms} 节的意义下仿射，
则称它{\it 仿射}。
\end{definition}

\noindent
对我们而言，把代数叠的仿射态射看作可由代数空间表示，并且在《叠的性质》第
\ref{stacks-properties-section-properties-morphisms} 节的意义下仿射的代数叠态射，
会稍方便一些。（回顾，在 Stacks 项目中，“可表”默认意为可由概形表示。）
这显然与刚才定义的概念等价，故下文将直接使用这一刻画，不再说明。
下面证明关于此概念的几个简单引理。

\begin{lemma}
\label{stacks-morphisms-lemma-base-change-affine}
设 $\mathcal{X} \to \mathcal{Y}$ 是代数叠的态射，
$\mathcal{Z} \to \mathcal{Y}$ 是代数叠的仿射态射。则
$\mathcal{Z} \times_\mathcal{Y} \mathcal{X} \to \mathcal{X}$
是代数叠的仿射态射。
\end{lemma}

\begin{proof}
这由《叠的性质》第
\ref{stacks-properties-section-properties-morphisms} 节的讨论得到。
\end{proof}

\begin{lemma}
\label{stacks-morphisms-lemma-composition-affine}
代数叠仿射态射的复合仿射。
\end{lemma}

\begin{proof}
这由《叠的性质》第
\ref{stacks-properties-section-properties-morphisms} 节的讨论以及
《空间的态射》引理 \ref{spaces-morphisms-lemma-composition-affine} 得到。
\end{proof}

\begin{lemma}
\label{stacks-morphisms-lemma-affine-permanence}
设
$$
\xymatrix{
\mathcal{X} \ar[rr]_f \ar[rd]_a & & \mathcal{Y} \ar[dl]^b \\
& \mathcal{Z}
}
$$
是代数叠态射的交换图。若 $a$ 仿射且 $\Delta_b$ 仿射，则 $f$ 仿射。
\end{lemma}

\begin{proof}
态射 $a$ 的基变换
$\text{pr}_2 : \mathcal{X} \times_\mathcal{Z} \mathcal{Y} \to \mathcal{Y}$
由引理 \ref{stacks-morphisms-lemma-base-change-affine} 仿射。态射
$(1, f) : \mathcal{X} \to \mathcal{X} \times_\mathcal{Z} \mathcal{Y}$
是
$\Delta_b : \mathcal{Y} \to \mathcal{Y} \times_\mathcal{Z} \mathcal{Y}$
沿态射 $\mathcal{X} \times_\mathcal{Z} \mathcal{Y} \to
\mathcal{Y} \times_\mathcal{Z} \mathcal{Y}$ 的基变换（见《范畴》第
\ref{categories-section-2-fibre-products} 节的材料）。因此，由引理
\ref{stacks-morphisms-lemma-base-change-affine}，它仿射。仿射态射的复合
$f = \text{pr}_2 \circ (1, f)$ 由引理 \ref{stacks-morphisms-lemma-composition-affine} 仿射，
证明完成。
\end{proof}





\section{整态射与有限态射}
\label{stacks-morphisms-section-integral}

\noindent
代数叠的整态射与有限态射定义如下。

\begin{definition}
\label{stacks-morphisms-definition-integral}
设 $f : \mathcal{X} \to \mathcal{Y}$ 是代数叠的态射。
\begin{enumerate}
\item 若 $f$ 可表，并且在《叠的性质》第
\ref{stacks-properties-section-properties-morphisms} 节的意义下为整，
则称 $f$ {\it 整}。
\item 若 $f$ 可表，并且在《叠的性质》第
\ref{stacks-properties-section-properties-morphisms} 节的意义下有限，
则称 $f$ {\it 有限}。
\end{enumerate}
\end{definition}

\noindent
对我们而言，把代数叠的整态射（相应地，有限态射）看作可由代数空间表示，
并且在《叠的性质》第
\ref{stacks-properties-section-properties-morphisms} 节的意义下为整
（相应地，有限）的代数叠态射，会稍方便一些。（回顾，在 Stacks 项目中，
“可表”默认意为可由概形表示。）这显然与刚才定义的概念等价，故下文将
直接使用这一刻画，不再说明。下面证明关于此概念的几个简单引理。

\begin{lemma}
\label{stacks-morphisms-lemma-base-change-integral}
设 $\mathcal{X} \to \mathcal{Y}$ 是代数叠的态射。设
$\mathcal{Z} \to \mathcal{Y}$ 是代数叠的整（或有限）态射。则
$\mathcal{Z} \times_\mathcal{Y} \mathcal{X} \to \mathcal{X}$
是代数叠的整（或有限）态射。
\end{lemma}

\begin{proof}
这由《叠的性质》第
\ref{stacks-properties-section-properties-morphisms} 节的讨论得到。
\end{proof}

\begin{lemma}
\label{stacks-morphisms-lemma-composition-integral}
代数叠整态射（相应地，有限态射）的复合为整（相应地，有限）。
\end{lemma}

\begin{proof}
这由《叠的性质》第
\ref{stacks-properties-section-properties-morphisms} 节的讨论以及
《空间的态射》引理 \ref{spaces-morphisms-lemma-composition-integral} 得到。
\end{proof}




\section{开态射}
\label{stacks-morphisms-section-open}

\noindent
设 $f$ 是可由代数空间表示的代数叠态射。我们已在《叠的性质》第
\ref{stacks-properties-section-properties-morphisms} 节定义了 $f$ 泛开的含义。
下面给出另一种刻画。

\begin{lemma}
\label{stacks-morphisms-lemma-characterize-representable-universally-open}
设 $f : \mathcal{X} \to \mathcal{Y}$ 是可由代数空间表示的代数叠态射。
下列条件等价：
\begin{enumerate}
\item $f$ 泛开（定义见《叠的性质》第
\ref{stacks-properties-section-properties-morphisms} 节）；以及
\item 对代数叠的每个态射 $\mathcal{Z} \to \mathcal{Y}$，拓扑空间的态射
$|\mathcal{Z} \times_\mathcal{Y} \mathcal{X}| \to |\mathcal{Z}|$ 是开映射。
\end{enumerate}
\end{lemma}

\begin{proof}
假设 (1)，并设 $\mathcal{Z} \to \mathcal{Y}$ 如 (2) 所述。取概形 $V$
及满光滑态射 $V \to \mathcal{Z}$。根据假设，代数空间的态射
$V \times_\mathcal{Y} \mathcal{X} \to V$ 泛开；特别地，映射
$|V \times_\mathcal{Y} \mathcal{X}| \to |V|$ 是开映射。由《叠的性质》第
\ref{stacks-properties-section-points} 节，在交换图
$$
\xymatrix{
|V \times_\mathcal{Y} \mathcal{X}| \ar[r] \ar[d] &
|\mathcal{Z} \times_\mathcal{Y} \mathcal{X}| \ar[d] \\
|V| \ar[r] & |\mathcal{Z}|
}
$$
中，水平箭头开且满，并且
$$
|V \times_\mathcal{Y} \mathcal{X}| \longrightarrow
|V| \times_{|\mathcal{Z}|} |\mathcal{Z} \times_\mathcal{Y} \mathcal{X}|
$$
满。因此，由于左侧竖直箭头是开映射，右侧竖直箭头也是开映射。
这证明了 (2)。蕴涵 (2) $\Rightarrow$ (1) 由定义得到。
\end{proof}

\noindent
因此可以采用如下自然定义。

\begin{definition}
\label{stacks-morphisms-definition-open}
设 $f : \mathcal{X} \to \mathcal{Y}$ 是代数叠的态射。
\begin{enumerate}
\item 若拓扑空间的映射 $|\mathcal{X}| \to |\mathcal{Y}|$ 是开映射，
则称 $f$ {\it 开}。
\item 若对代数叠的每个态射 $\mathcal{Z} \to \mathcal{Y}$，
拓扑空间的态射
$$
|\mathcal{Z} \times_\mathcal{Y} \mathcal{X}| \to |\mathcal{Z}|
$$
都是开映射，即基变换
$\mathcal{Z} \times_\mathcal{Y} \mathcal{X} \to \mathcal{Z}$ 开，
则称 $f$ {\it 泛开}。
\end{enumerate}
\end{definition}

\begin{lemma}
\label{stacks-morphisms-lemma-base-change-universally-open}
代数叠的泛开态射沿任意代数叠态射的基变换仍泛开。
\end{lemma}

\begin{proof}
这立即由定义得到。
\end{proof}

\begin{lemma}
\label{stacks-morphisms-lemma-composition-universally-open}
两个代数叠（泛）开态射的复合（泛）开。
\end{lemma}

\begin{proof}
略。
\end{proof}







\section{商映射态射}
\label{stacks-morphisms-section-submersive}

\noindent
设 $f$ 是可由代数空间表示的代数叠态射。我们已在《叠的性质》第
\ref{stacks-properties-section-properties-morphisms} 节定义了 $f$ 泛商映射的含义。
下面给出另一种刻画。

\begin{lemma}
\label{stacks-morphisms-lemma-characterize-representable-universally-submersive}
设 $f : \mathcal{X} \to \mathcal{Y}$ 是可由代数空间表示的代数叠态射。
下列条件等价：
\begin{enumerate}
\item $f$ 泛商映射（定义见《叠的性质》第
\ref{stacks-properties-section-properties-morphisms} 节）；以及
\item 对代数叠的每个态射 $\mathcal{Z} \to \mathcal{Y}$，拓扑空间的态射
$|\mathcal{Z} \times_\mathcal{Y} \mathcal{X}| \to |\mathcal{Z}|$ 是商映射。
\end{enumerate}
\end{lemma}

\begin{proof}
假设 (1)，并设 $\mathcal{Z} \to \mathcal{Y}$ 如 (2) 所述。取概形 $V$
及满光滑态射 $V \to \mathcal{Z}$。根据假设，代数空间的态射
$V \times_\mathcal{Y} \mathcal{X} \to V$ 泛商映射；特别地，映射
$|V \times_\mathcal{Y} \mathcal{X}| \to |V|$ 是商映射。由《叠的性质》第
\ref{stacks-properties-section-points} 节，在交换图
$$
\xymatrix{
|V \times_\mathcal{Y} \mathcal{X}| \ar[r] \ar[d] &
|\mathcal{Z} \times_\mathcal{Y} \mathcal{X}| \ar[d] \\
|V| \ar[r] & |\mathcal{Z}|
}
$$
中，水平箭头开且满，并且
$$
|V \times_\mathcal{Y} \mathcal{X}| \longrightarrow
|V| \times_{|\mathcal{Z}|} |\mathcal{Z} \times_\mathcal{Y} \mathcal{X}|
$$
满。因此，由于左侧竖直箭头是商映射，右侧竖直箭头也是商映射。
这证明了 (2)。蕴涵 (2) $\Rightarrow$ (1) 由定义得到。
\end{proof}

\noindent
因此可以采用如下自然定义。

\begin{definition}
\label{stacks-morphisms-definition-submersive}
设 $f : \mathcal{X} \to \mathcal{Y}$ 是代数叠的态射。
\begin{enumerate}
\item 若连续映射 $|\mathcal{X}| \to |\mathcal{Y}|$ 是商映射，
则称 $f$ {\it 是商映射}\footnote{这与微分流形的浸没概念截然不同。}；
见《拓扑》定义 \ref{topology-definition-submersive}。
\item 若对代数叠的每个态射 $\mathcal{Y}' \to \mathcal{Y}$，
基变换 $\mathcal{Y}' \times_\mathcal{Y} \mathcal{X} \to \mathcal{Y}'$
都是商映射，则称 $f$ {\it 泛商映射}。
\end{enumerate}
\end{definition}

\noindent
注意，商映射态射特别是满的。

\begin{lemma}
\label{stacks-morphisms-lemma-base-change-universally-submersive}
代数叠的泛商映射态射沿任意代数叠态射的基变换仍泛商映射。
\end{lemma}

\begin{proof}
这立即由定义得到。
\end{proof}

\begin{lemma}
\label{stacks-morphisms-lemma-composition-universally-submersive}
两个代数叠（泛）商映射态射的复合仍（泛）商映射。
\end{lemma}

\begin{proof}
略。
\end{proof}












\section{泛闭态射}
\label{stacks-morphisms-section-universally-closed}

\noindent
设 $f$ 是可由代数空间表示的代数叠态射。我们已在《叠的性质》第
\ref{stacks-properties-section-properties-morphisms} 节定义了 $f$ 泛闭的含义。
下面给出另一种刻画。

\begin{lemma}
\label{stacks-morphisms-lemma-characterize-representable-universally-closed}
设 $f : \mathcal{X} \to \mathcal{Y}$ 是可由代数空间表示的代数叠态射。
下列条件等价：
\begin{enumerate}
\item $f$ 泛闭（定义见《叠的性质》第
\ref{stacks-properties-section-properties-morphisms} 节）；以及
\item 对代数叠的每个态射 $\mathcal{Z} \to \mathcal{Y}$，拓扑空间的态射
$|\mathcal{Z} \times_\mathcal{Y} \mathcal{X}| \to |\mathcal{Z}|$ 是闭映射。
\end{enumerate}
\end{lemma}

\begin{proof}
假设 (1)，并设 $\mathcal{Z} \to \mathcal{Y}$ 如 (2) 所述。取概形 $V$
及满光滑态射 $V \to \mathcal{Z}$。根据假设，代数空间的态射
$V \times_\mathcal{Y} \mathcal{X} \to V$ 泛闭；特别地，映射
$|V \times_\mathcal{Y} \mathcal{X}| \to |V|$ 是闭映射。由《叠的性质》第
\ref{stacks-properties-section-points} 节，在交换图
$$
\xymatrix{
|V \times_\mathcal{Y} \mathcal{X}| \ar[r] \ar[d] &
|\mathcal{Z} \times_\mathcal{Y} \mathcal{X}| \ar[d] \\
|V| \ar[r] & |\mathcal{Z}|
}
$$
中，水平箭头开且满，并且
$$
|V \times_\mathcal{Y} \mathcal{X}| \longrightarrow
|V| \times_{|\mathcal{Z}|} |\mathcal{Z} \times_\mathcal{Y} \mathcal{X}|
$$
满。因此，由于左侧竖直箭头是闭映射，右侧竖直箭头也是闭映射。
这证明了 (2)。蕴涵 (2) $\Rightarrow$ (1) 由定义得到。
\end{proof}

\noindent
因此可以采用如下自然定义。

\begin{definition}
\label{stacks-morphisms-definition-closed}
设 $f : \mathcal{X} \to \mathcal{Y}$ 是代数叠的态射。
\begin{enumerate}
\item 若拓扑空间的映射 $|\mathcal{X}| \to |\mathcal{Y}|$ 是闭映射，
则称 $f$ {\it 闭}。
\item 若对代数叠的每个态射 $\mathcal{Z} \to \mathcal{Y}$，
拓扑空间的态射
$$
|\mathcal{Z} \times_\mathcal{Y} \mathcal{X}| \to |\mathcal{Z}|
$$
都是闭映射，即基变换
$\mathcal{Z} \times_\mathcal{Y} \mathcal{X} \to \mathcal{Z}$ 闭，
则称 $f$ {\it 泛闭}。
\end{enumerate}
\end{definition}

\begin{lemma}
\label{stacks-morphisms-lemma-base-change-universally-closed}
代数叠的泛闭态射沿任意代数叠态射的基变换仍泛闭。
\end{lemma}

\begin{proof}
这立即由定义得到。
\end{proof}

\begin{lemma}
\label{stacks-morphisms-lemma-composition-universally-closed}
两个代数叠（泛）闭态射的复合（泛）闭。
\end{lemma}

\begin{proof}
略。
\end{proof}

\begin{lemma}
\label{stacks-morphisms-lemma-universally-closed-local}
设 $f : \mathcal{X} \to \mathcal{Y}$ 是代数叠的态射。下列条件等价：
\begin{enumerate}
\item $f$ 泛闭；
\item 对每个概形 $Z$ 和每个态射 $Z \to \mathcal{Y}$，投影
$|Z \times_\mathcal{Y} \mathcal{X}| \to |Z|$ 是闭映射；
\item 对每个仿射概形 $Z$ 和每个态射 $Z \to \mathcal{Y}$，投影
$|Z \times_\mathcal{Y} \mathcal{X}| \to |Z|$ 是闭映射；以及
\item 存在代数空间 $V$ 及满光滑态射 $V \to \mathcal{Y}$，使得
$V \times_\mathcal{Y} \mathcal{X} \to V$ 是代数叠的泛闭态射。
\end{enumerate}
\end{lemma}

\begin{proof}
略去 (1) 蕴含 (2) 以及 (2) 蕴含 (3) 的证明。

\medskip\noindent
假设 (3)。取满光滑态射 $V \to \mathcal{Y}$。我们将证明
$V \times_\mathcal{Y} \mathcal{X} \to V$ 是代数叠的泛闭态射。
设 $\mathcal{Z} \to V$ 是从代数叠到 $V$ 的态射。设
$W \to \mathcal{Z}$ 是满光滑态射，其中 $W = \coprod W_i$ 是仿射概形的
不交并。于是有交换图
$$
\xymatrix{
\coprod_i |W_i \times_\mathcal{Y} \mathcal{X}| \ar@{=}[r] \ar[d] &
|W \times_\mathcal{Y} \mathcal{X}| \ar[r] \ar[d] &
|\mathcal{Z} \times_\mathcal{Y} \mathcal{X}| \ar[d] \ar@{=}[r] &
|\mathcal{Z} \times_V (V \times_\mathcal{Y} \mathcal{X})| \ar[ld] \\
\coprod |W_i| \ar@{=}[r] &
|W| \ar[r] &
|\mathcal{Z}|
}
$$
我们必须证明东南方向的箭头是闭映射。中间的水平箭头满且开
（《叠的性质》引理 \ref{stacks-properties-lemma-topology-points}）。
由假设 (3) 以及 $W_i$ 仿射这一事实，左侧竖直箭头都是闭映射，
从而右侧竖直箭头也是闭映射。

\medskip\noindent
假设 (4)。我们将证明 $f$ 泛闭。设 $\mathcal{Z} \to \mathcal{Y}$
是代数叠的态射。考虑图
$$
\xymatrix{
|(V \times_\mathcal{Y} \mathcal{Z})
\times_V (V \times_\mathcal{Y} \mathcal{X})| \ar@{=}[r] \ar[rd] &
|V \times_\mathcal{Y} \mathcal{X}| \ar[r] \ar[d] &
|Z \times_\mathcal{Y} \mathcal{X}| \ar[d] \\
 &
|V \times_\mathcal{Y} \mathcal{Z}| \ar[r] &
|\mathcal{Z}|
}
$$
根据假设，西南方向的箭头是闭映射。水平箭头满且开，因为相应的代数叠态射
满且光滑（见上面的文献）。因此，右侧竖直箭头是闭映射。
\end{proof}










\section{泛单射态射}
\label{stacks-morphisms-section-universally-injective}

\noindent
设 $f$ 是可由代数空间表示的代数叠态射。我们已在《叠的性质》第
\ref{stacks-properties-section-properties-morphisms} 节定义了 $f$ 泛单射的含义。
下面给出另一种刻画。

\begin{lemma}
\label{stacks-morphisms-lemma-characterize-representable-universally-injective}
设 $f : \mathcal{X} \to \mathcal{Y}$ 是可由代数空间表示的代数叠态射。
下列条件等价：
\begin{enumerate}
\item $f$ 泛单射（定义见《叠的性质》第
\ref{stacks-properties-section-properties-morphisms} 节）；以及
\item 对代数叠的每个态射 $\mathcal{Z} \to \mathcal{Y}$，映射
$|\mathcal{Z} \times_\mathcal{Y} \mathcal{X}| \to |\mathcal{Z}|$ 单射。
\end{enumerate}
\end{lemma}

\begin{proof}
假设 (1)，并设 $\mathcal{Z} \to \mathcal{Y}$ 如 (2) 所述。取概形 $V$
及满光滑态射 $V \to \mathcal{Z}$。根据假设，代数空间的态射
$V \times_\mathcal{Y} \mathcal{X} \to V$ 泛单射；特别地，映射
$|V \times_\mathcal{Y} \mathcal{X}| \to |V|$ 单射。由《叠的性质》第
\ref{stacks-properties-section-points} 节，在交换图
$$
\xymatrix{
|V \times_\mathcal{Y} \mathcal{X}| \ar[r] \ar[d] &
|\mathcal{Z} \times_\mathcal{Y} \mathcal{X}| \ar[d] \\
|V| \ar[r] & |\mathcal{Z}|
}
$$
中，水平箭头开且满，并且
$$
|V \times_\mathcal{Y} \mathcal{X}| \longrightarrow
|V| \times_{|\mathcal{Z}|} |\mathcal{Z} \times_\mathcal{Y} \mathcal{X}|
$$
满。因此，由于左侧竖直箭头单射，右侧竖直箭头也单射。
这证明了 (2)。蕴涵 (2) $\Rightarrow$ (1) 由定义得到。
\end{proof}

\noindent
因此可以采用如下自然定义。

\begin{definition}
\label{stacks-morphisms-definition-universally-injective}
设 $f : \mathcal{X} \to \mathcal{Y}$ 是代数叠的态射。若对代数叠的每个态射
$\mathcal{Z} \to \mathcal{Y}$，映射
$$
|\mathcal{Z} \times_\mathcal{Y} \mathcal{X}| \to |\mathcal{Z}|
$$
均单射，则称 $f$ {\it 泛单射}。
\end{definition}

\begin{lemma}
\label{stacks-morphisms-lemma-base-change-universally-injective}
代数叠的泛单射态射沿任意代数叠态射的基变换仍泛单射。
\end{lemma}

\begin{proof}
这立即由定义得到。
\end{proof}

\begin{lemma}
\label{stacks-morphisms-lemma-composition-universally-injective}
两个代数叠泛单射态射的复合泛单射。
\end{lemma}

\begin{proof}
略。
\end{proof}

\begin{lemma}
\label{stacks-morphisms-lemma-universally-injective}
设 $f : \mathcal{X} \to \mathcal{Y}$ 是代数叠的态射。下列条件等价：
\begin{enumerate}
\item $f$ 泛单射；
\item $\Delta : \mathcal{X} \to \mathcal{X} \times_\mathcal{Y} \mathcal{X}$
满；以及
\item 对代数闭域、态射 $x_1, x_2 : \Spec(k) \to \mathcal{X}$ 以及
$2$-箭头 $\beta : f \circ x_1 \to f \circ x_2$，存在 $2$-箭头
$\alpha : x_1 \to x_2$，使得 $\beta = \text{id}_f \star \alpha$。
\end{enumerate}
\end{lemma}

\begin{proof}
(1) $\Rightarrow$ (2)。若 $f$ 泛单射，则第一投影
$|\mathcal{X} \times_\mathcal{Y} \mathcal{X}| \to |\mathcal{X}|$
单射，这蕴含 $|\Delta|$ 满。

\medskip\noindent
(2) $\Rightarrow$ (1)。假设 $\Delta$ 满。则 $\Delta$ 的任意基变换都满
（见《叠的性质》第 \ref{stacks-properties-section-surjective} 节）。
由于 $f$ 的基变换的对角态射是 $\Delta$ 的基变换，只须证明
$|\mathcal{X}| \to |\mathcal{Y}|$ 单射。否则，由《叠的性质》引理
\ref{stacks-properties-lemma-points-cartesian}，第一投影
$|\mathcal{X} \times_\mathcal{Y} \mathcal{X}| \to |\mathcal{X}|$
不单射。这当然意味着 $|\Delta|$ 不满。

\medskip\noindent
(3) $\Rightarrow$ (2)。设
$t \in |\mathcal{X} \times_\mathcal{Y} \mathcal{X}|$。可以用态射
$t : \Spec(k) \to \mathcal{X} \times_\mathcal{Y} \mathcal{X}$
表示 $t$，其中 $k$ 是代数闭域。根据我们对 $2$-纤维积的构造，
可以用 $(x_1, x_2, \beta)$ 表示 $t$，其中
$x_1, x_2 : \Spec(k) \to \mathcal{X}$，而
$\beta : f \circ x_1 \to f \circ x_2$ 是 $2$-态射。于是 (3) 蕴含存在
映到 $\beta$ 的 $2$-态射 $\alpha : x_1 \to x_2$。这恰好意味着
$\Delta(x_1) = (x_1, x_1, \text{id})$ 与 $t$ 同构。因此 (2) 成立。

\medskip\noindent
(2) $\Rightarrow$ (3)。设 $x_1, x_2 : \Spec(k) \to \mathcal{X}$ 是态射，
其中 $k$ 是代数闭域。设
$\beta : f \circ x_1 \to f \circ x_2$ 是 $2$-态射。如上一段所述，得到态射
$t = (x_1, x_2, \beta) :
\Spec(k) \to \mathcal{X} \times_\mathcal{Y} \mathcal{X}$.
由引理 \ref{stacks-morphisms-lemma-properties-diagonal}，
$$
T = \mathcal{X}
\times_{\Delta, \mathcal{X} \times_\mathcal{Y} \mathcal{X}, t} \Spec(k)
$$
是 $\Spec(k)$ 上局部有限型的代数空间。条件 (2) 蕴含 $T$ 非空。
由于 $k$ 代数闭，$T$ 中有一个 $k$-点。展开定义，这意味着
$\Mor(\Spec(k), \mathcal{X})$ 中存在态射
$\alpha : x_1 \to x_2$，使得 $\beta = \text{id}_f \star \alpha$。
\end{proof}

\begin{lemma}
\label{stacks-morphisms-lemma-universally-injective-point}
设 $f : \mathcal{X} \to \mathcal{Y}$ 是代数叠的泛单射态射。设
$y : \Spec(k) \to \mathcal{Y}$ 是态射，其中 $k$ 是代数闭域。
若 $y$ 属于 $|\mathcal{X}| \to |\mathcal{Y}|$ 的像，则存在态射
$x : \Spec(k) \to \mathcal{X}$，使得 $y = f \circ x$。
\end{lemma}

\begin{proof}
先指出此引理并非显然，因为 $y$ 属于 $|f|$ 的像这一假设只意味着：
在可能用扩域替换 $k$ 之后，可以把 $y$ 提升为到 $\mathcal{X}$ 的态射。
为证明该引理，可以沿 $y$ 基变换 $f$；因而可以假设有非空代数叠
$\mathcal{X}$ 和泛单射态射 $\mathcal{X} \to \Spec(k)$，并希望找到
$\mathcal{X}$ 的一个 $k$-值点。可以用约化替换 $\mathcal{X}$。
可以选取域 $k'$ 以及满、平坦、局部有限型态射
$\Spec(k') \to \mathcal{X}$；见《叠的性质》引理
\ref{stacks-properties-lemma-unique-point}。由于
$\mathcal{X} \to \Spec(k)$ 泛单射，
$$
\Spec(k') \times_\mathcal{X} \Spec(k') \to \Spec(k' \otimes_k k')
$$
作为满态射
$\Delta : \mathcal{X} \to \mathcal{X} \times_{\Spec(k)} \mathcal{X}$
的基变换，它是满的（引理 \ref{stacks-morphisms-lemma-universally-injective}）。
由于 $k$ 代数闭，$k' \otimes_k k'$ 是整环（《代数》引理
\ref{algebra-lemma-geometrically-integral-any-integral-base-change}）。
设 $\xi \in \Spec(k') \times_\mathcal{X} \Spec(k')$ 是映到
$\Spec(k' \otimes_k k')$ 一般点的点。令 $U$ 为
$\Spec(k') \times_\mathcal{X} \Spec(k')$ 中含有 $\xi$ 的连通分支上
所赋的约化诱导闭子概形结构。于是两个投影 $U \to \Spec(k')$ 都局部有限型，
因为投影 $\Spec(k') \times_\mathcal{X} \Spec(k') \to \Spec(k')$
作为态射 $\Spec(k') \to \mathcal{X}$ 的基变换也具有这一性质。
应用《簇》命题 \ref{varieties-proposition-unique-base-field}，可知 $k'$ 在
$\Gamma(U, \mathcal{O}_U)$ 中的两个像的整闭包相等。在 $\kappa(\xi)$
中考察，这意味着每个形如 $\lambda \otimes 1$ 的元素都在子域
$$
1 \otimes k' \subset 
(\text{fraction field of }k' \otimes_k k') \subset
\kappa(\xi).
$$
上代数相关。由于 $k$ 代数闭，这只有在 $k' = k$ 时才可能，证明完成。
\end{proof}

\begin{lemma}
\label{stacks-morphisms-lemma-universally-injective-local}
设 $f : \mathcal{X} \to \mathcal{Y}$ 是代数叠的态射。下列条件等价：
\begin{enumerate}
\item $f$ 泛单射；
\item 对每个仿射概形 $Z$ 及任意态射 $Z \to \mathcal{Y}$，态射
$Z \times_\mathcal{Y} \mathcal{X} \to Z$ 泛单射；以及
\item 此处添加更多内容。
\end{enumerate}
\end{lemma}

\begin{proof}
蕴涵 (1) $\Rightarrow$ (2) 是立即的。假设 (2) 成立。我们将证明
$\Delta_f : \mathcal{X} \to \mathcal{X} \times_\mathcal{Y} \mathcal{X}$
满；由引理 \ref{stacks-morphisms-lemma-universally-injective}，这蕴含 (1)。考虑仿射概形 $V$
以及光滑态射 $V \to \mathcal{Y}$。由于
$g : V \times_\mathcal{Y} \mathcal{X} \to V$
由 (2) 泛单射，可知 $\Delta_g$ 满。然而，$\Delta_g$ 是 $\Delta_f$
沿光滑态射 $V \to \mathcal{Y}$ 的基变换。由于这些态射
$V \to \mathcal{Y}$ 合起来满，故 $\Delta_f$ 满。
\end{proof}

\begin{lemma}
\label{stacks-morphisms-lemma-check-universally-injective-covering}
设 $f : \mathcal{X} \to \mathcal{Y}$ 是代数叠的态射。设
$W \to \mathcal{Y}$ 满、平坦且局部有限表示，其中 $W$ 是代数空间。
若基变换 $W \times_\mathcal{Y} \mathcal{X} \to W$ 泛单射，
则 $f$ 泛单射。
\end{lemma}

\begin{proof}
注意，态射
$g : W \times_\mathcal{Y} \mathcal{X} \to W$
的对角态射 $\Delta_g$ 是 $\Delta_f$ 沿 $W \to \mathcal{Y}$ 的基变换。
因此，若 $\Delta_g$ 满，则由《叠的性质》引理
\ref{stacks-properties-lemma-check-property-covering}，$\Delta_f$ 也满。
于是该引理由引理 \ref{stacks-morphisms-lemma-universally-injective} 中的刻画 (2) 得到。
\end{proof}








\section{泛同胚}
\label{stacks-morphisms-section-universal-homeomorphisms}

\noindent
设 $f$ 是可由代数空间表示的代数叠态射。我们已在《叠的性质》第
\ref{stacks-properties-section-properties-morphisms} 节定义了 $f$ 为泛同胚的含义。
下面给出另一种刻画。

\begin{lemma}
\label{stacks-morphisms-lemma-characterize-representable-universal-homeomorphism}
设 $f : \mathcal{X} \to \mathcal{Y}$ 是可由代数空间表示的代数叠态射。
下列条件等价：
\begin{enumerate}
\item $f$ 是泛同胚（《叠的性质》第
\ref{stacks-properties-section-properties-morphisms} 节）；以及
\item 对代数叠的每个态射 $\mathcal{Z} \to \mathcal{Y}$，拓扑空间的映射
$|\mathcal{Z} \times_\mathcal{Y} \mathcal{X}| \to |\mathcal{Z}|$
是同胚。
\end{enumerate}
\end{lemma}

\begin{proof}
假设 (1)，并设 $\mathcal{Z} \to \mathcal{Y}$ 如 (2) 所述。取概形 $V$
及满光滑态射 $V \to \mathcal{Z}$。根据假设，代数空间的态射
$V \times_\mathcal{Y} \mathcal{X} \to V$ 是泛同胚；特别地，映射
$|V \times_\mathcal{Y} \mathcal{X}| \to |V|$ 是同胚。由《叠的性质》第
\ref{stacks-properties-section-points} 节，在交换图
$$
\xymatrix{
|V \times_\mathcal{Y} \mathcal{X}| \ar[r] \ar[d] &
|\mathcal{Z} \times_\mathcal{Y} \mathcal{X}| \ar[d] \\
|V| \ar[r] & |\mathcal{Z}|
}
$$
中，水平箭头开且满，并且
$$
|V \times_\mathcal{Y} \mathcal{X}| \longrightarrow
|V| \times_{|\mathcal{Z}|} |\mathcal{Z} \times_\mathcal{Y} \mathcal{X}|
$$
满。因此，由于左侧竖直箭头是同胚，右侧竖直箭头也是同胚。
这证明了 (2)。蕴涵 (2) $\Rightarrow$ (1) 由定义得到。
\end{proof}

\noindent
因此可以采用如下自然定义。

\begin{definition}
\label{stacks-morphisms-definition-universal-homeomorphism}
设 $f : \mathcal{X} \to \mathcal{Y}$ 是代数叠的态射。若对代数叠的每个态射
$\mathcal{Z} \to \mathcal{Y}$，拓扑空间的映射
$$
|\mathcal{Z} \times_\mathcal{Y} \mathcal{X}| \to |\mathcal{Z}|
$$
都是同胚，则称 $f$ 是{\it 泛同胚}。
\end{definition}

\begin{lemma}
\label{stacks-morphisms-lemma-base-change-universal-homeomorphism}
代数叠的泛同胚沿任意代数叠态射的基变换仍是泛同胚。
\end{lemma}

\begin{proof}
这立即由定义得到。
\end{proof}

\begin{lemma}
\label{stacks-morphisms-lemma-composition-universal-homeomorphism}
两个代数叠泛同胚的复合是泛同胚。
\end{lemma}

\begin{proof}
略。
\end{proof}

\begin{lemma}
\label{stacks-morphisms-lemma-check-universal-homeomorphism-covering}
设 $f : \mathcal{X} \to \mathcal{Y}$ 是代数叠的态射。设
$W \to \mathcal{Y}$ 满、平坦且局部有限表示，其中 $W$ 是代数空间。
若基变换 $W \times_\mathcal{Y} \mathcal{X} \to W$ 是泛同胚，
则 $f$ 是泛同胚。
\end{lemma}

\begin{proof}
假设 $g : W \times_\mathcal{Y} \mathcal{X} \to W$ 是泛同胚。
则 $g$ 泛单射，因而由引理
\ref{stacks-morphisms-lemma-check-universally-injective-covering}，$f$ 泛单射。
另一方面，设 $\mathcal{Z} \to \mathcal{Y}$ 是态射，其中 $\mathcal{Z}$
是代数叠。取概形 $U$ 及满光滑态射
$U \to W \times_\mathcal{Y} \mathcal{Z}$。考虑图
$$
\xymatrix{
W \times_\mathcal{Y} \mathcal{X} \ar[d]^g &
U \times_\mathcal{Y} \mathcal{X} \ar[d] \ar[l] \ar[r] &
\mathcal{Z} \times_\mathcal{Y} \mathcal{X} \ar[d] \\
W &
U \ar[l] \ar[r] &
\mathcal{Z}
}
$$
根据对 $g$ 的假设，中间的竖直箭头在拓扑空间上诱导同胚。态射
$U \to \mathcal{Z}$ 和
$U \times_\mathcal{Y} \mathcal{X} \to
\mathcal{Z} \times_\mathcal{Y} \mathcal{X}$
满、平坦且局部有限表示，故在拓扑空间上诱导开映射。由此可知
$|\mathcal{Z} \times_\mathcal{Y} \mathcal{X}| \to |\mathcal{Z}|$
是开映射。满性容易证明；我们略去其证明。
\end{proof}










\section{在源与靶上光滑局部的态射类型}
\label{stacks-morphisms-section-local-source-target}

\noindent
给定代数空间态射的一个{\it 在源与靶上光滑局部}的性质（见
《空间上的下降》定义 \ref{spaces-descent-definition-local-source-target}），
可以用它定义代数叠态射的相应性质，具体做法是要求下面引理中的任一等价条件。

\begin{lemma}
\label{stacks-morphisms-lemma-local-source-target}
设 $\mathcal{P}$ 是代数空间态射的一个在源与靶上光滑局部的性质。
设 $f : \mathcal{X} \to \mathcal{Y}$ 是代数叠的态射。考虑交换图
$$
\xymatrix{
U \ar[d]_a \ar[r]_h & V \ar[d]^b \\
\mathcal{X} \ar[r]^f & \mathcal{Y}
}
$$
，其中 $U$ 和 $V$ 是代数空间，竖直箭头光滑。下列条件等价：
\begin{enumerate}
\item 对任何如上并且 $U \to \mathcal{X} \times_\mathcal{Y} V$ 也光滑的图，
态射 $h$ 具有性质 $\mathcal{P}$；以及
\item 对某个如上的图，若 $a : U \to \mathcal{X}$ 满，则态射 $h$
具有性质 $\mathcal{P}$。
\end{enumerate}
若 $\mathcal{X}$ 和 $\mathcal{Y}$ 可由代数空间表示，则这些条件还等价于
$f$（作为代数空间的态射）具有性质 $\mathcal{P}$。若 $\mathcal{P}$
还在任意基变换下保持，并且在基上为 fppf 局部的，则对于可由代数空间表示的
态射 $f$，这些条件还等价于 $f$ 在《叠的性质》第
\ref{stacks-properties-section-properties-morphisms} 节的意义下具有性质
$\mathcal{P}$。
\end{lemma}

\begin{proof}
先证明蕴涵 (1) $\Rightarrow$ (2)。取代数空间 $V$ 以及满光滑态射
$V \to \mathcal{Y}$。取代数空间 $U$ 以及满光滑态射
$U \to \mathcal{X} \times_\mathcal{Y} V$。注意，$U \to \mathcal{X}$
作为基变换 $\mathcal{X} \times_\mathcal{Y} V \to \mathcal{X}$ 与所选态射
$U \to \mathcal{X} \times_\mathcal{Y} V$ 的复合，也满且光滑。因此得到
如 (1) 中的图。于是，若 (1) 成立，则 $h : U \to V$ 具有性质
$\mathcal{P}$；由于 $U \to \mathcal{X}$ 满，这意味着 (2) 成立。

\medskip\noindent
反之，假设 (2) 成立，并设 $U, V, a, b, h$ 如 (2) 所述。再设
$U', V', a', b', h'$ 是 (1) 中的任意图。图示为
$$
\xymatrix{
U \ar[d] \ar[r]_h & V \ar[d] \\
\mathcal{X} \ar[r]^f & \mathcal{Y}
}
\quad\quad
\xymatrix{
U' \ar[d] \ar[r]_{h'} & V' \ar[d] \\
\mathcal{X} \ar[r]^f & \mathcal{Y}
}
$$
为证明 (2) 蕴含 (1)，必须证明 $h'$ 具有 $\mathcal{P}$。为此考虑交换图
$$
\xymatrix{
U \ar[dd]^h &
U \times_\mathcal{X} U' \ar[d] \ar[l] \ar@/^6ex/[dd]^{(h, h')} \ar[r] &
U' \ar[dd]^{h'} \\
& U \times_\mathcal{Y} V' \ar[lu] \ar[d] & \\
V &
V \times_\mathcal{Y} V' \ar[l] \ar[r] &
V'
}
$$
，其中各对象都是代数空间。注意，水平箭头分别作为光滑态射
$V \to \mathcal{Y}$、$V' \to \mathcal{Y}$、$U \to \mathcal{X}$ 和
$U' \to \mathcal{X}$ 的基变换而光滑。还要注意，图
$$
\xymatrix{
U \times_\mathcal{X} U' \ar[d] \ar[r] & U' \ar[d] \\
U \times_\mathcal{Y} V' \ar[r] & \mathcal{X} \times_\mathcal{Y} V'
}
$$
是笛卡尔图；由于 $U', V', a', b', h'$ 如 (1) 所述，故左侧竖直箭头光滑。
由《空间上的下降》引理
\ref{spaces-descent-lemma-local-source-target-implies} 的 (2)，性质
$\mathcal{P}$ 在靶上光滑局部，因此基变换
$U \times_\mathcal{Y} V' \to V \times_\mathcal{Y} V'$ 具有
$\mathcal{P}$。由《空间上的下降》引理
\ref{spaces-descent-lemma-local-source-target-implies} 的 (1)，
$\mathcal{P}$ 在源上光滑局部，
所以可前复合光滑态射
$U \times_\mathcal{X} U' \to U \times_\mathcal{Y} V'$，并得出
$(h, h')$ 具有 $\mathcal{P}$。由于 $V \times_\mathcal{Y} V' \to V'$
光滑，由《空间上的下降》引理
\ref{spaces-descent-lemma-local-source-target-implies} 的 (3)，可知
$U \times_\mathcal{X} U' \to V'$ 具有 $\mathcal{P}$。最后，由于
$U \times_X U' \to U'$ 满且光滑，并且 $\mathcal{P}$ 在源上光滑局部
（同一引理），可知 $h'$ 具有 $\mathcal{P}$。这就完成了 (1) 与 (2)
等价性的证明。

\medskip\noindent
若 $\mathcal{X}$ 和 $\mathcal{Y}$ 可表，则应用《空间上的下降》引理
\ref{spaces-descent-lemma-local-source-target-characterize} 可知，
(1) 和 (2) 等价于 $f$ 具有 $\mathcal{P}$。

\medskip\noindent
最后，假设 $f$ 可表，$U, V, a, b, h$ 如引理的 (2) 所述，并且
$\mathcal{P}$ 在任意基变换下保持。必须证明：对任意概形 $Z$ 和态射
$Z \to \mathcal{Y}$，基变换
$Z \times_\mathcal{Y} \mathcal{X} \to Z$
具有性质 $\mathcal{P}$。考虑图
$$
\xymatrix{
Z \times_\mathcal{Y} U \ar[d] \ar[r] &
Z \times_\mathcal{Y} V \ar[d] \\
Z \times_\mathcal{Y} \mathcal{X} \ar[r] &
Z
}
$$
注意，上方水平箭头是 $h$ 的基变换，因而具有性质 $\mathcal{P}$。
左侧竖直箭头光滑且满，右侧竖直箭头光滑。因此，应用《空间上的下降》引理
\ref{spaces-descent-lemma-local-source-target-characterize} 可知
$Z \times_\mathcal{Y} \mathcal{X} \to Z$ 具有性质 $\mathcal{P}$。
\end{proof}

\begin{definition}
\label{stacks-morphisms-definition-P}
设 $\mathcal{P}$ 是代数空间态射的一个在源与靶上光滑局部的性质。
若引理 \ref{stacks-morphisms-lemma-local-source-target} 中的等价条件成立，则称代数叠态射
$f : \mathcal{X} \to \mathcal{Y}$ {\it 具有性质 $\mathcal{P}$}。
\end{definition}

\begin{remark}
\label{stacks-morphisms-remark-composition}
设 $\mathcal{P}$ 是代数空间态射的一个在源与靶上光滑局部且在复合下稳定的
性质。则定义 \ref{stacks-morphisms-definition-P} 所定义的代数叠态射性质在复合下稳定。
事实上，设 $f : \mathcal{X} \to \mathcal{Y}$ 和
$g : \mathcal{Y} \to \mathcal{Z}$ 是具有性质 $\mathcal{P}$ 的代数叠态射。
取代数空间 $W$ 以及满光滑态射 $W \to \mathcal{Z}$。取代数空间 $V$
以及满光滑态射 $V \to \mathcal{Y} \times_\mathcal{Z} W$。最后，
取代数空间 $U$ 以及满光滑态射
$U \to \mathcal{X} \times_\mathcal{Y} V$。根据定义，态射
$V \to W$ 和 $U \to V$ 具有性质 $\mathcal{P}$。由于假设
$\mathcal{P}$ 在复合下稳定，故 $U \to W$ 具有性质 $\mathcal{P}$。
再次由定义可知，$g \circ f : \mathcal{X} \to \mathcal{Z}$ 具有性质
$\mathcal{P}$。
\end{remark}

\begin{remark}
\label{stacks-morphisms-remark-base-change}
设 $\mathcal{P}$ 是代数空间态射的一个在源与靶上光滑局部且在基变换下稳定的
性质。则定义 \ref{stacks-morphisms-definition-P} 所定义的代数叠态射性质在基变换下稳定。
事实上，设 $f : \mathcal{X} \to \mathcal{Y}$ 和
$g : \mathcal{Y}' \to \mathcal{Y}$ 是代数叠的态射，并假设 $f$ 具有性质
$\mathcal{P}$。取代数空间 $V$ 以及满光滑态射 $V \to \mathcal{Y}$。
取代数空间 $U$ 以及满光滑态射
$U \to \mathcal{X} \times_\mathcal{Y} V$。最后，取代数空间 $V'$
以及满光滑态射 $V' \to \mathcal{Y}' \times_\mathcal{Y} V$。
根据定义，态射 $U \to V$ 具有性质 $\mathcal{P}$。由于假设
$\mathcal{P}$ 在基变换下稳定，故 $V' \times_V U \to V'$ 具有性质
$\mathcal{P}$。考虑图
$$
\xymatrix{
V' \times_V U \ar[r] \ar[d] &
\mathcal{Y}' \times_\mathcal{Y} \mathcal{X} \ar[r] \ar[d] &
\mathcal{X} \ar[d] \\
V' \ar[r] & \mathcal{Y}' \ar[r] & \mathcal{Y}
}
$$
可见左上方的水平箭头光滑且满，因此根据定义，投影
$\mathcal{Y}' \times_\mathcal{Y} \mathcal{X} \to \mathcal{Y}'$ 具有性质
$\mathcal{P}$。
\end{remark}

\begin{remark}
\label{stacks-morphisms-remark-implication}
设 $\mathcal{P}, \mathcal{P}'$ 是代数空间态射的两个在源与靶上光滑局部的
性质。假设对代数空间态射有
$\mathcal{P} \Rightarrow \mathcal{P}'$。则对于定义
\ref{stacks-morphisms-definition-P} 中用 $\mathcal{P}$ 和 $\mathcal{P}'$ 定义的代数叠态射
性质，同样有 $\mathcal{P} \Rightarrow \mathcal{P}'$。这由定义显然。
\end{remark}









\section{有限型态射}
\label{stacks-morphisms-section-finite-type}

\noindent
代数空间态射的“局部有限型”性质在源与靶上光滑局部；见《空间上的下降》
注记 \ref{spaces-descent-remark-list-local-source-target}。它还在基变换下稳定，
并且在靶上为 fpqc 局部的；见《空间的态射》引理
\ref{spaces-morphisms-lemma-base-change-finite-type} 和《空间上的下降》引理
\ref{spaces-descent-lemma-descending-property-locally-finite-type}。
因此，由上面的引理 \ref{stacks-morphisms-lemma-local-source-target}，可以如下定义代数空间的态射
局部有限型的含义；当态射可由代数空间表示时，这与《叠的性质》第
\ref{stacks-properties-section-properties-morphisms} 节定义的已有概念一致。

\begin{definition}
\label{stacks-morphisms-definition-locally-finite-type}
设 $f : \mathcal{X} \to \mathcal{Y}$ 是代数叠的态射。
\begin{enumerate}
\item 若引理 \ref{stacks-morphisms-lemma-local-source-target} 中的等价条件对
$\mathcal{P} = \text{locally of finite type}$ 成立，则称 $f$
{\it 局部有限型}。
\item 若 $f$ 局部有限型且拟紧，则称它{\it 有限型}。
\end{enumerate}
\end{definition}

\begin{lemma}
\label{stacks-morphisms-lemma-composition-finite-type}
有限型态射的复合为有限型。局部有限型也同样成立。
\end{lemma}

\begin{proof}
结合注记 \ref{stacks-morphisms-remark-composition} 与《空间的态射》引理
\ref{spaces-morphisms-lemma-composition-finite-type}。
\end{proof}

\begin{lemma}
\label{stacks-morphisms-lemma-base-change-finite-type}
有限型态射的基变换有限型。局部有限型也同样成立。
\end{lemma}

\begin{proof}
结合注记 \ref{stacks-morphisms-remark-base-change} 与《空间的态射》引理
\ref{spaces-morphisms-lemma-base-change-finite-type}。
\end{proof}

\begin{lemma}
\label{stacks-morphisms-lemma-immersion-locally-finite-type}
浸入局部有限型。
\end{lemma}

\begin{proof}
结合注记 \ref{stacks-morphisms-remark-implication} 与《空间的态射》引理
\ref{spaces-morphisms-lemma-immersion-locally-finite-type}。
\end{proof}

\begin{lemma}
\label{stacks-morphisms-lemma-locally-finite-type-locally-noetherian}
设 $f : \mathcal{X} \to \mathcal{Y}$ 是代数叠的态射。若 $f$ 局部有限型且
$\mathcal{Y}$ 局部诺特，则 $\mathcal{X}$ 局部诺特。
\end{lemma}

\begin{proof}
设
$$
\xymatrix{
U \ar[d] \ar[r] & V \ar[d] \\
\mathcal{X} \ar[r] & \mathcal{Y}
}
$$
是交换图，其中 $U$、$V$ 是概形，$V \to \mathcal{Y}$ 满且光滑，
$U \to V \times_\mathcal{Y} \mathcal{X}$ 满且光滑。则 $U \to V$
局部有限型。若 $\mathcal{Y}$ 局部诺特，则 $V$ 局部诺特。由《态射》引理
\ref{morphisms-lemma-finite-type-noetherian}，$U$ 局部诺特，
这意味着 $\mathcal{X}$ 局部诺特。
\end{proof}

\noindent
下面两个引理将在后文得到改进（在讨论局部有限表示的代数叠态射之后）。

\begin{lemma}
\label{stacks-morphisms-lemma-check-finite-type-covering}
设 $f : \mathcal{X} \to \mathcal{Y}$ 是代数叠的态射。设
$W \to \mathcal{Y}$ 满、平坦且局部有限表示，其中 $W$ 是代数空间。
若基变换 $W \times_\mathcal{Y} \mathcal{X} \to W$ 局部有限型，
则 $f$ 局部有限型。
\end{lemma}

\begin{proof}
取代数空间 $V$ 以及满光滑态射 $V \to \mathcal{Y}$。取代数空间 $U$
以及满光滑态射 $U \to V \times_\mathcal{Y} \mathcal{X}$。
必须证明 $U \to V$ 局部有限表示。现在把所有对象沿
$W \to \mathcal{Y}$ 基变换：令
$U' = W \times_\mathcal{Y} U$,
$V' = W \times_\mathcal{Y} V$,
$\mathcal{X}' = W \times_\mathcal{Y} \mathcal{X}$,
以及 $\mathcal{Y}' = W \times_\mathcal{Y} \mathcal{Y} = W$。
由基变换，$U' \to V' \times_{\mathcal{Y}'} \mathcal{X}'$ 仍光滑。
因此，根据代数叠局部有限型态射的定义以及
$\mathcal{X}' \to \mathcal{Y}'$ 局部有限型这一假设，
$U' \to V'$ 局部有限型。又因为 $V' \to V$ 作为
$W \to \mathcal{Y}$ 的基变换满、平坦且局部有限表示，由《空间上的下降》引理
\ref{spaces-descent-lemma-descending-property-locally-finite-type}，
$U \to V$ 局部有限型，证明完成。
\end{proof}

\begin{lemma}
\label{stacks-morphisms-lemma-check-finite-type-precompose}
设 $\mathcal{X} \to \mathcal{Y} \to \mathcal{Z}$ 是代数叠的态射。
假设 $\mathcal{X} \to \mathcal{Z}$ 局部有限型，并且
$\mathcal{X} \to \mathcal{Y}$ 可由代数空间表示、满、平坦且局部有限表示。
则 $\mathcal{Y} \to \mathcal{Z}$ 局部有限型。
\end{lemma}

\begin{proof}
取代数空间 $W$ 以及满光滑态射 $W \to \mathcal{Z}$。取代数空间 $V$
以及满光滑态射 $V \to W \times_\mathcal{Z} \mathcal{Y}$。令
$U = V \times_\mathcal{Y} \mathcal{X}$，它是代数空间。已知
$U \to V$ 满、平坦且局部有限表示，而 $U \to W$ 局部有限型。
因此，该引理归结为代数空间态射的情形，而该情形就是《空间上的下降》引理
\ref{spaces-descent-lemma-locally-finite-type-fppf-local-source}。
\end{proof}

\begin{lemma}
\label{stacks-morphisms-lemma-finite-type-permanence}
设 $f : \mathcal{X} \to \mathcal{Y}$、
$g : \mathcal{Y} \to \mathcal{Z}$ 是代数叠的态射。若
$g \circ f : \mathcal{X} \to \mathcal{Z}$ 局部有限型，
则 $f : \mathcal{X} \to \mathcal{Y}$ 局部有限型。
\end{lemma}

\begin{proof}
可以找到图
$$
\xymatrix{
U \ar[r] \ar[d] & V \ar[r] \ar[d] & W \ar[d] \\
\mathcal{X} \ar[r] & \mathcal{Y} \ar[r] & \mathcal{Z}
}
$$
，其中 $U$、$V$、$W$ 是概形，竖直箭头 $W \to \mathcal{Z}$ 满且光滑，
箭头 $V \to \mathcal{Y} \times_\mathcal{Z} W$ 满且光滑，并且箭头
$U \to \mathcal{X} \times_\mathcal{Y} V$ 满且光滑。于是
$U \to \mathcal{X} \times_\mathcal{Z} V$ 也满且光滑（它是满光滑态射与
此类态射的一个基变换的复合）。根据定义，$U \to W$ 局部有限型。
因此，由《态射》引理 \ref{morphisms-lemma-permanence-finite-type}，
$U \to V$ 局部有限型；根据定义，这又意味着
$\mathcal{X} \to \mathcal{Y}$ 局部有限型。
\end{proof}










\section{有限型点}
\label{stacks-morphisms-section-points-finite-type}

\noindent
设 $\mathcal{X}$ 是代数叠。有限型点 $x \in |\mathcal{X}|$ 是可以由局部有限型
态射 $\Spec(k) \to \mathcal{X}$ 表示的点。对于代数空间和代数叠，
有限型点适合作为闭点的替代。例如，它们总是“足够多”。

\begin{lemma}
\label{stacks-morphisms-lemma-point-finite-type}
设 $\mathcal{X}$ 是代数叠，$x \in |\mathcal{X}|$。下列条件等价：
\begin{enumerate}
\item 存在表示 $x$ 的局部有限型态射 $\Spec(k) \to \mathcal{X}$。
\item 存在概形 $U$、闭点 $u \in U$ 以及光滑态射
$\varphi : U \to \mathcal{X}$，使得 $\varphi(u) = x$。
\end{enumerate}
\end{lemma}

\begin{proof}
设 $u \in U$ 和 $U \to \mathcal{X}$ 如 (2) 所述。则
$\Spec(\kappa(u)) \to U$ 有限型，而 $U \to \mathcal{X}$ 可表且局部有限型
（由《空间的态射》引理
\ref{spaces-morphisms-lemma-etale-locally-finite-presentation} 和
\ref{spaces-morphisms-lemma-finite-presentation-finite-type}）。因此，
由引理 \ref{stacks-morphisms-lemma-composition-finite-type}，(1) 成立。

\medskip\noindent
反之，假设 $\Spec(k) \to \mathcal{X}$ 局部有限型并表示 $x$。设
$U \to \mathcal{X}$ 是满光滑态射，其中 $U$ 是概形。根据假设，
$U \times_\mathcal{X} \Spec(k) \to U$ 是局部有限型的代数空间态射。
取 $U \times_\mathcal{X} \Spec(k)$ 的一个有限型点 $v$（至少存在一个；
见《空间的态射》引理
\ref{spaces-morphisms-lemma-identify-finite-type-points}）。由《空间的态射》引理
\ref{spaces-morphisms-lemma-finite-type-points-morphism}，$v$ 的像
$u \in U$ 是 $U$ 的有限型点。因此，由《态射》引理
\ref{morphisms-lemma-identify-finite-type-points}，缩小 $U$ 后可以假设 $u$
是 $U$ 的闭点，即 (2) 成立。
\end{proof}

\begin{definition}
\label{stacks-morphisms-definition-finite-type-point}
设 $\mathcal{X}$ 是代数叠。若引理 \ref{stacks-morphisms-lemma-point-finite-type} 中的等价条件
成立，则称点 $x \in |\mathcal{X}|$ 是{\it 有限型点}\footnote{这略微滥用了
术语，因为称为“局部有限型点”或许更准确。}。以
$\mathcal{X}_{\text{ft-pts}}$ 表示 $\mathcal{X}$ 的有限型点集。
\end{definition}

\noindent
有限型点集可以描述如下。

\begin{lemma}
\label{stacks-morphisms-lemma-identify-finite-type-points}
设 $\mathcal{X}$ 是代数叠。则
$$
\mathcal{X}_{\text{ft-pts}} =
\bigcup\nolimits_{\varphi : U \to \mathcal{X}\text{ smooth}} |\varphi|(U_0)
$$
，其中 $U_0$ 是 $U$ 的闭点集。这里可以让 $U$ 遍历 $\mathcal{X}$ 上所有
光滑概形，或只遍历 $\mathcal{X}$ 上所有光滑仿射概形。
\end{lemma}

\begin{proof}
立即由引理 \ref{stacks-morphisms-lemma-point-finite-type} 得到。
\end{proof}

\begin{lemma}
\label{stacks-morphisms-lemma-finite-type-points-morphism}
设 $f : \mathcal{X} \to \mathcal{Y}$ 是代数叠的态射。若 $f$ 局部有限型，则
$f(\mathcal{X}_{\text{ft-pts}}) \subset \mathcal{Y}_{\text{ft-pts}}$.
\end{lemma}

\begin{proof}
取 $x \in \mathcal{X}_{\text{ft-pts}}$。用局部有限型态射
$x : \Spec(k) \to \mathcal{X}$ 表示 $x$。则由引理
\ref{stacks-morphisms-lemma-composition-finite-type}，$f \circ x$ 局部有限型。因此
$f(x) \in \mathcal{Y}_{\text{ft-pts}}$。
\end{proof}

\begin{lemma}
\label{stacks-morphisms-lemma-finite-type-points-surjective-morphism}
设 $f : \mathcal{X} \to \mathcal{Y}$ 是代数叠的态射。若 $f$ 局部有限型且满，则
$f(\mathcal{X}_{\text{ft-pts}}) = \mathcal{Y}_{\text{ft-pts}}$.
\end{lemma}

\begin{proof}
由引理 \ref{stacks-morphisms-lemma-finite-type-points-morphism}，有
$f(\mathcal{X}_{\text{ft-pts}}) \subset \mathcal{Y}_{\text{ft-pts}}$。
设 $y \in |\mathcal{Y}|$ 是有限型点。用局部有限型态射
$\Spec(k) \to \mathcal{Y}$ 表示 $y$。由于 $f$ 满，代数叠
$\mathcal{X}_k = \Spec(k) \times_\mathcal{Y} \mathcal{X}$ 非空，
故由引理 \ref{stacks-morphisms-lemma-identify-finite-type-points}，它有有限型点
$x \in |\mathcal{X}_k|$。态射 $\mathcal{X}_k \to \mathcal{X}$
是 $\Spec(k) \to \mathcal{Y}$ 的基变换，因而局部有限型
（引理 \ref{stacks-morphisms-lemma-base-change-finite-type}）。所以由引理
\ref{stacks-morphisms-lemma-finite-type-points-morphism}，$x$ 在 $\mathcal{X}$ 中的像是有限型点，
并且根据构造映到 $y$。
\end{proof}

\begin{lemma}
\label{stacks-morphisms-lemma-enough-finite-type-points}
设 $\mathcal{X}$ 是代数叠。对任意局部闭子集 $T \subset |\mathcal{X}|$，有
$$
T \not = \emptyset
\Rightarrow
T \cap \mathcal{X}_{\text{ft-pts}} \not = \emptyset.
$$
特别地，对任意闭子集 $T \subset |\mathcal{X}|$，
$T \cap \mathcal{X}_{\text{ft-pts}}$ 在 $T$ 中稠密。
\end{lemma}

\begin{proof}
令 $i : \mathcal{Z} \to \mathcal{X}$ 为 $T$ 上的约化诱导子叠结构；
见《叠的性质》注记
\ref{stacks-properties-remark-stack-structure-locally-closed-subset}。
浸入局部有限型；见引理 \ref{stacks-morphisms-lemma-immersion-locally-finite-type}。
因此，由引理 \ref{stacks-morphisms-lemma-finite-type-points-morphism}，有
$\mathcal{Z}_{\text{ft-pts}} \subset \mathcal{X}_{\text{ft-pts}} \cap T$.
最后，任何带有到 $\mathcal{Z}$ 的光滑态射的非空仿射概形 $U$ 至少有一个闭点，
故由引理 \ref{stacks-morphisms-lemma-identify-finite-type-points}，$\mathcal{Z}$ 至少有一个
有限型点。引理得证。
\end{proof}

\noindent
下面给出代数叠上有限型点的另一种、更具技术性的刻画。特别地，它告诉我们：
只要 $x$ 是有限型点，$\mathcal{X}$ 在 $x$ 处的剩余胚就存在！

\begin{lemma}
\label{stacks-morphisms-lemma-point-finite-type-monomorphism}
设 $\mathcal{X}$ 是代数叠，$x \in |\mathcal{X}|$。下列条件等价：
\begin{enumerate}
\item $x$ 是有限型点；
\item 存在代数叠 $\mathcal{Z}$，其底拓扑空间 $|\mathcal{Z}|$ 是单点集，
并且存在局部有限型态射 $f : \mathcal{Z} \to \mathcal{X}$，使得
$\{x\} = |f|(|\mathcal{Z}|)$；以及
\item $\mathcal{X}$ 在 $x$ 处的剩余胚 $\mathcal{Z}_x$ 存在，并且包含态射
$\mathcal{Z}_x \to \mathcal{X}$ 局部有限型。
\end{enumerate}
\end{lemma}

\begin{proof}
（本段出现的所有态射都可由代数空间表示，因此《叠的性质》第
\ref{stacks-properties-section-properties-morphisms} 节的约定和结果均适用。）
假设 $x$ 是有限型点。取仿射概形 $U$、闭点 $u \in U$ 以及光滑态射
$\varphi : U \to \mathcal{X}$，使得 $\varphi(u) = x$；见引理
\ref{stacks-morphisms-lemma-identify-finite-type-points}。照例令 $u = \Spec(\kappa(u))$。
令 $R = u \times_\mathcal{X} u$，于是得到代数空间中的群胚
$(u, R, s, t, c)$；见《代数叠》引理
\ref{algebraic-lemma-map-space-into-stack}。投影态射 $R \to u$ 是如下复合：
$$
R = u \times_\mathcal{X} u \to
u \times_\mathcal{X} U \to
u \times_\mathcal{X} X = u
$$
其中第一个箭头有限型（它是概形闭浸入 $u \to U$ 的基变换），
第二个箭头光滑（它是光滑态射 $U \to \mathcal{X}$ 的基变换）。因此，
$s, t : R \to u$ 作为复合局部有限型；见《空间的态射》引理
\ref{spaces-morphisms-lemma-composition-finite-type}。由于 $u$ 是域的谱，
$s, t$ 平坦且局部有限表示（由《空间的态射》引理
\ref{spaces-morphisms-lemma-noetherian-finite-type-finite-presentation}）。
由《可表性判据》定理
\ref{criteria-theorem-flat-groupoid-gives-algebraic-stack}，
$\mathcal{Z} = [u/R]$ 是代数叠。
由《代数叠》引理 \ref{algebraic-lemma-map-space-into-stack}，得到典范态射
$$
f : \mathcal{Z} \longrightarrow \mathcal{X}
$$
，它全忠实。因此，该态射可由代数空间表示；见《代数叠》引理
\ref{algebraic-lemma-characterize-representable-by-algebraic-spaces}；
并且是单态射；见《叠的性质》引理
\ref{stacks-properties-lemma-monomorphism}。由此，$\mathcal{X}$ 在 $x$ 处的
剩余胚 $\mathcal{Z}_x \subset \mathcal{X}$ 存在，并且 $f$ 通过一个等价
$\mathcal{Z} \to \mathcal{Z}_x$ 分解；见《叠的性质》引理
\ref{stacks-properties-lemma-residual-gerbe-unique}。根据构造，图
$$
\xymatrix{
u \ar[d] \ar[r] & U \ar[d] \\
\mathcal{Z} \ar[r]^f & \mathcal{X}
}
$$
交换。由《可表性判据》引理
\ref{criteria-lemma-flat-quotient-flat-presentation}，左侧竖直箭头满、
平坦且局部有限表示。考虑图
$$
\xymatrix{
u \times_\mathcal{X} U \ar[d] \ar[r] &
\mathcal{Z} \times_\mathcal{X} U \ar[r] \ar[d] & U \ar[d] \\
u \ar[r] & \mathcal{Z} \ar[r]^f & \mathcal{X}
}
$$
由于 $u \to \mathcal{X}$ 局部有限型，其基变换
$u \times_\mathcal{X} U \to U$ 局部有限型。此外，
$u \times_\mathcal{X} U \to \mathcal{Z} \times_\mathcal{X} U$
作为 $u \to \mathcal{Z}$ 的基变换满、平坦且局部有限表示。因此
$\{u \times_\mathcal{X} U \to \mathcal{Z} \times_\mathcal{X} U\}$
是代数空间的 fppf 覆盖，故由《空间上的下降》引理
\ref{spaces-descent-lemma-locally-finite-presentation-fppf-local-source}，
$\mathcal{Z} \times_\mathcal{X} U \to U$ 局部有限型。根据定义，
这意味着 $f$ 局部有限型（因为竖直箭头
$\mathcal{Z} \times_\mathcal{X} U \to \mathcal{Z}$ 作为
$U \to \mathcal{X}$ 的基变换光滑，并且由于 $\mathcal{Z}$ 只有一个点而满）。
由于 $\mathcal{Z} = \mathcal{Z}_x$，(3) 成立。

\medskip\noindent
(3) 蕴含 (2) 是显然的。若 (2) 成立，则由引理
\ref{stacks-morphisms-lemma-finite-type-points-morphism}，并用引理
\ref{stacks-morphisms-lemma-enough-finite-type-points} 看出
$\mathcal{Z}_{\text{ft-pts}}$ 非空，即 $\mathcal{Z}$ 的唯一点是
$\mathcal{Z}$ 的有限型点，可知 $x$ 是 $\mathcal{X}$ 的有限型点。
\end{proof}







\section{自同构群}
\label{stacks-morphisms-section-automorphism-groups}

\noindent
设 $\mathcal{X}$ 是代数叠。设 $x \in |\mathcal{X}|$ 对应于
$x : \Spec(k) \to \mathcal{X}$。在这种情形下，我们常说“设 $G_x/k$
是 $x$ 的自同构群代数空间”。这只是指
$$
G_x = \mathit{Isom}_\mathcal{X}(x, x) =
\Spec(k) \times_\mathcal{X} \mathcal{I}_\mathcal{X}
$$
是 $x$ 的自同构所成的群代数空间。它是 $\Spec(k)$ 上的群代数空间。
若 $k'/k$ 是域扩张，则所诱导态射 $x' : \Spec(k') \to \mathcal{X}$
的自同构群代数空间是 $G_x$ 到 $\Spec(k')$ 的基变换。

\begin{lemma}
\label{stacks-morphisms-lemma-automorphism-group-scheme}
在上述情形下，若下列条件之一成立，则 $G_x$ 是概形：
\begin{enumerate}
\item $\Delta : \mathcal{X} \to \mathcal{X} \times \mathcal{X}$
拟分离；
\item $\Delta : \mathcal{X} \to \mathcal{X} \times \mathcal{X}$
局部分离；
\item $\mathcal{X}$ 是拟 DM；
\item $\mathcal{I}_\mathcal{X} \to \mathcal{X}$
拟分离；
\item $\mathcal{I}_\mathcal{X} \to \mathcal{X}$
局部分离；或
\item $\mathcal{I}_\mathcal{X} \to \mathcal{X}$
局部拟有限。
\end{enumerate}
\end{lemma}

\begin{proof}
由引理 \ref{stacks-morphisms-lemma-diagonal-diagonal}，注意到 (1) $\Rightarrow$ (4)、
(2) $\Rightarrow$ (5) 且 (3) $\Rightarrow$ (6)。在情形 (4) 中，$G_x$
是拟分离代数空间；在情形 (5) 中，$G_x$ 是局部分离代数空间。
在两种情形下，$G_x$ 都是适宜代数空间（《适宜代数空间》第
\ref{decent-spaces-section-reasonable-decent} 节以及引理
\ref{decent-spaces-lemma-locally-separated-decent}）。于是，由《空间中群胚进阶》
引理 \ref{spaces-more-groupoids-lemma-group-scheme-over-field-separated}，
$G_x$ 分离；继而由《空间中群胚进阶》命题
\ref{spaces-more-groupoids-proposition-group-space-scheme-over-field}，
$G_x$ 是概形。在情形 (6) 中，$G_x \to \Spec(k)$ 局部拟有限，
故由《域上的空间》引理
\ref{spaces-over-fields-lemma-locally-quasi-finite-over-field}，$G_x$ 是概形。
\end{proof}

\begin{lemma}
\label{stacks-morphisms-lemma-property-automorphism-groups}
设 $\mathcal{X}$ 是代数叠，$x \in |\mathcal{X}|$ 是一点。设 $P$
是域上代数空间的一个在基域扩张下不变的性质；例如
$P(X/k) = X \to \Spec(k)\text{ is finite}$.
下列条件等价：
\begin{enumerate}
\item 对 $x$ 的类中的某个态射 $x : \Spec(k) \to \mathcal{X}$，
自同构群代数空间 $G_x/k$ 具有 $P$；以及
\item 对 $x$ 的类中的任意态射 $x : \Spec(k) \to \mathcal{X}$，
自同构群代数空间 $G_x/k$ 具有 $P$。
\end{enumerate}
\end{lemma}

\begin{proof}
略。
\end{proof}

\begin{remark}
\label{stacks-morphisms-remark-property-automorphism-groups}
设 $P$ 是域上代数空间的一个在基域扩张下不变的性质。给定代数叠
$\mathcal{X}$ 和 $x \in |\mathcal{X}|$，若引理
\ref{stacks-morphisms-lemma-property-automorphism-groups} 中的等价条件成立，则称
$\mathcal{X}$ 在 $x$ 处的自同构群具有 $P$。例如，若每当
$x : \Spec(k) \to \mathcal{X}$ 是 $x$ 的代表时，
$G_x \to \Spec(k)$ 都有限，则称{\it $\mathcal{X}$ 在 $x$ 处的自同构群有限}。
光滑、真等性质也同样处理。（这里显然有术语滥用，但我们相信它不会造成
混淆或不精确。）
\end{remark}

\begin{lemma}
\label{stacks-morphisms-lemma-iso-automorphism-groups}
设 $f : \mathcal{X} \to \mathcal{Y}$ 是代数叠的态射，
$x \in |\mathcal{X}|$ 是一点。下列条件等价：
\begin{enumerate}
\item 对 $x$ 的类中的某个态射 $x : \Spec(k) \to \mathcal{X}$，
令 $y = f \circ x$，则自同构群代数空间的映射
$G_x \to G_y$ 是同构；以及
\item 对 $x$ 的类中的任意态射 $x : \Spec(k) \to \mathcal{X}$，
令 $y = f \circ x$，则自同构群代数空间的映射
$G_x \to G_y$ 是同构。
\end{enumerate}
\end{lemma}

\begin{proof}
这归结为同构性在靶上为 fpqc 局部这一事实；见《空间上的下降》引理
\ref{spaces-descent-lemma-descending-property-isomorphism}。具体地，
假设 $k'/k$ 是域扩张，把相应的复合记作
$x' : \Spec(k') \to \mathcal{X}$，并令 $y' = f \circ x'$。则态射
$G_{x'} \to G_{y'}$ 是 $G_x \to G_y$ 沿
$\Spec(k') \to \Spec(k)$ 的基变换。因此，$G_x \to G_y$ 是同构当且仅当
$G_{x'} \to G_{y'}$ 是同构。由此可见，只要此性质对一个代表成立，
它就会传遍整个等价类。
\end{proof}

\begin{remark}
\label{stacks-morphisms-remark-identify-automorphism-groups}
设 $f : \mathcal{X} \to \mathcal{Y}$ 是代数叠的态射，
$x \in |\mathcal{X}|$ 是一点。文献中用“{\it $f$ 在 $x$ 处保持稳定子}”
或“{\it $f$ 在 $x$ 处反映不动点}”来表示引理
\ref{stacks-morphisms-lemma-iso-automorphism-groups} 的等价条件对 $f$ 和 $x$ 成立。
我们更愿意说“{\it $f$ 在 $x$ 与 $f(x)$ 处的自同构群之间诱导同构}”。
\end{remark}





\section{表现与代数叠的性质}
\label{stacks-morphisms-section-presentations-properties}

\noindent
设 $(U, R, s, t, c)$ 是代数空间中的群胚。若
$s, t : R \to U$ 平坦且局部有限表现，则商叠 $[U/R]$
是代数叠；参见《可表示性判据》定理
\ref{criteria-theorem-flat-groupoid-gives-algebraic-stack}.
本节研究 $(U, R, s, t, c)$ 的哪些性质会蕴含代数叠
$[U/R]$ 的相应性质。

\begin{lemma}
\label{stacks-morphisms-lemma-properties-diagonal-from-presentation}
设 $(U, R, s, t, c)$ 是代数空间中的群胚，且
$s, t : R \to U$ 平坦并局部有限表现。考虑代数叠
$\mathcal{X} = [U/R]$（见上文）。
\begin{enumerate}
\item 若 $R \to U \times U$ 分离，则 $\Delta_\mathcal{X}$ 分离。
\item 若 $U$、$R$ 分离，则 $\Delta_\mathcal{X}$ 分离。
\item 若 $R \to U \times U$ 局部拟有限，则 $\mathcal{X}$ 为拟 DM。
\item 若 $s, t : R \to U$ 局部拟有限，则 $\mathcal{X}$ 为拟 DM。
\item 若 $R \to U \times U$ 为真态射，则 $\mathcal{X}$ 分离。
\item 若 $s, t : R \to U$ 为真态射且 $U$ 分离，则
$\mathcal{X}$ 分离。
\item 此处添加更多内容。
\end{enumerate}
\end{lemma}

\begin{proof}
由《可表示性判据》引理
\ref{criteria-lemma-flat-quotient-flat-presentation}，
态射 $U \to \mathcal{X}$ 满、平坦且局部有限表现。因此
$U \times U \to \mathcal{X} \times \mathcal{X}$ 也有这些性质。
有笛卡儿图
$$
\xymatrix{
R  = U \times_\mathcal{X} U \ar[r] \ar[d] & U \times U \ar[d] \\
\mathcal{X} \ar[r] & \mathcal{X} \times \mathcal{X}
}
$$
（参见《空间中的群胚》引理
\ref{spaces-groupoids-lemma-quotient-stack-2-cartesian}).
于是，$\Delta_\mathcal{X}$ 具有《代数叠的性质》一节
\ref{stacks-properties-section-properties-morphisms}
所列性质之一，当且仅当态射 $R \to U \times U$ 具有该性质；参见
《代数叠的性质》引理
\ref{stacks-properties-lemma-check-property-covering}。
这说明了 (1)、(3) 和 (5)。条件 (2) 蕴含
$R \to U \times U$ 分离，故 (2) 由 (1) 得出。条件 (4)
蕴含条件 (3)，故 (4) 由 (3) 得出。由《空间的态射》引理
\ref{spaces-morphisms-lemma-universally-closed-permanence}
可知，条件 (6) 蕴含条件 (5)，故 (6) 由 (5) 得出。
\end{proof}

\begin{lemma}
\label{stacks-morphisms-lemma-points-presentation}
设 $(U, R, s, t, c)$ 是代数空间中的群胚，且
$s, t : R \to U$ 平坦并局部有限表现。考虑代数叠
$\mathcal{X} = [U/R]$（见上文）。则 $|R| \to |U| \times |U|$
的像是一个等价关系，而 $|\mathcal{X}|$ 是 $|U|$ 对此等价关系的商。
\end{lemma}

\begin{proof}
诱导态射 $p : U \to \mathcal{X}$ 满、平坦且局部有限表现；参见
《可表示性判据》引理
\ref{criteria-lemma-flat-quotient-flat-presentation}.
因此，由《代数叠的性质》引理
\ref{stacks-properties-lemma-characterize-surjective}.
可知 $|U| \to |\mathcal{X}|$ 满。注意到
$R = U \times_\mathcal{X} U$；参见《空间中的群胚》引理
\ref{spaces-groupoids-lemma-quotient-stack-2-cartesian}。于是，
《代数叠的性质》引理 \ref{stacks-properties-lemma-points-cartesian}
蕴含映射
$$
|R| \longrightarrow |U| \times_{|\mathcal{X}|} |U|
$$
满。因此，$|R| \to |U| \times |U|$ 的像恰为所有满足如下条件的
$(u_1, u_2) \in |U| \times |U|$ 所成的集合：$u_1$ 与 $u_2$
在 $|\mathcal{X}|$ 中有相同的像。合并这两个陈述即得引理。
\end{proof}






\section{代数叠的特殊表现}
\label{stacks-morphisms-section-presentations}

\noindent
本节证明两个重要定理。第一个把拟 DM 叠 $\mathcal{X}$ 刻画为如下形式的叠：
$\mathcal{X} = [U/R]$，其中 $s, t : R \to U$ 局部拟有限（并且平坦且
局部有限表现）。第二个定理断言，DM 代数叠是 Deligne--Mumford 叠。

\medskip\noindent
下述引理给出一个判据，用以判断表现的一个“截片”何时在该代数叠上仍然平坦。

\begin{lemma}
\label{stacks-morphisms-lemma-slice}
设 $\mathcal{X}$ 是代数叠。考虑笛卡儿图
$$
\xymatrix{
U \ar[d] & F \ar[l]^p \ar[d] \\
\mathcal{X} & \Spec(k) \ar[l]
}
$$
其中 $U$ 是代数空间，$k$ 是域，且 $U \to \mathcal{X}$ 平坦并局部有限表现。
设
$f_1, \ldots, f_r \in \Gamma(U, \mathcal{O}_U)$
且 $z \in |F|$，使 $f_1, \ldots, f_r$ 在局部环
$\mathcal{O}_{F, \overline{z}}$ 中的像构成正则序列。则以一个包含
$p(z)$ 的开子空间替换 $U$ 后，态射
$$
V(f_1, \ldots, f_r) \longrightarrow \mathcal{X}
$$
平坦且局部有限表现。
\end{lemma}

\begin{proof}
选取概形 $W$ 以及满光滑态射 $W \to \mathcal{X}$。选取域扩张
$k'/k$ 和态射 $w : \Spec(k') \to W$，使
$\Spec(k') \to W \to \mathcal{X}$ 与
$\Spec(k') \to \Spec(k) \to \mathcal{X}$ 是 $2$-同构的。
这是可能的，因为 $W \to \mathcal{X}$ 满。考虑交换图
$$
\xymatrix{
U \ar[d] &
U \times_\mathcal{X} W \ar[l]^-{\text{pr}_0} \ar[d] &
F' \ar[l]^-{p'} \ar[d] \\
\mathcal{X} &
W \ar[l] &
\Spec(k') \ar[l]
}
$$
其中两个方块都是笛卡儿的。由 $w$ 的选取可知
$F' = F \times_{\Spec(k)} \Spec(k')$。因此 $F' \to F$ 满，故可选取
映到 $z$ 的点 $z' \in |F'|$。由于 $F' \to F$ 平坦，
$\mathcal{O}_{F, \overline{z}} \to \mathcal{O}_{F', \overline{z}'}$
平坦；参见《空间的态射》引理
\ref{spaces-morphisms-lemma-flat-at-point-etale-local-rings}。
因此，$f_1, \ldots, f_r$ 在 $\mathcal{O}_{F', \overline{z}'}$ 中的像
构成正则序列；参见《代数》引理
\ref{algebra-lemma-flat-increases-depth}。注意，
$U \times_\mathcal{X} W \to W$ 是平坦且局部有限表现的代数空间态射。
因此，由《空间的态射进阶》引理
\ref{spaces-more-morphisms-lemma-slice} 可知，存在
$U \times_\mathcal{X} W$ 的开子空间 $U'$，它包含 $p'(z')$，并且交
$U' \cap (V(f_1, \ldots, f_r) \times_\mathcal{X} W)$ 在 $W$ 上平坦且
局部有限表现。注意，$\text{pr}_0(U')$ 是 $U$ 中包含 $p(z)$ 的开子空间，
因为 $\text{pr}_0$ 光滑，因而是开的。现在可知，态射
$U' \cap (V(f_1, \ldots, f_r) \times_\mathcal{X} W) \to \mathcal{X}$
作为复合
$$
U' \cap (V(f_1, \ldots, f_r) \times_\mathcal{X} W) \to W \to \mathcal{X}.
$$
平坦且局部有限表现。因此，《代数叠的性质》引理
\ref{stacks-properties-lemma-check-property-after-precomposing}
蕴含 $\text{pr}_0(U') \cap V(f_1, \ldots, f_r) \to \mathcal{X}$
平坦且局部有限表现，正合所需。
\end{proof}

\begin{lemma}
\label{stacks-morphisms-lemma-quasi-finite-at-point}
设 $\mathcal{X}$ 是代数叠。考虑笛卡儿图
$$
\xymatrix{
U \ar[d] & F \ar[l]^p \ar[d] \\
\mathcal{X} & \Spec(k) \ar[l]
}
$$
其中 $U$ 是代数空间，$k$ 是域，且 $U \to \mathcal{X}$ 局部有限型。
设 $z \in |F|$ 且 $\dim_z(F) = 0$。则以一个包含 $p(z)$ 的开子空间
替换 $U$ 后，态射
$$
U \longrightarrow \mathcal{X}
$$
局部拟有限。
\end{lemma}

\begin{proof}
由于 $f : U \to \mathcal{X}$ 局部有限型，存在极大开集
$W(f) \subset U$，使限制 $f|_{W(f)} : W(f) \to \mathcal{X}$
局部拟有限；参见《代数叠的性质》评注
\ref{stacks-properties-remark-local-source-apply}
(\ref{stacks-properties-item-loc-quasi-finite}).
因此，只需证明 $p(z)$ 是 $W(f)$ 的点。此外，上述评注还说明，
$W(f)$ 的构造与沿可由代数空间表示的态射所作的任意基变换交换。
故只需证明态射 $F \to \Spec(k)$ 在 $z$ 处局部拟有限。这立即由
《空间的态射》引理
\ref{spaces-morphisms-lemma-locally-quasi-finite-rel-dimension-0} 得出。
\end{proof}

\noindent
拟 DM 叠具有由概形给出的局部拟有限“覆盖”。

\begin{theorem}
\label{stacks-morphisms-theorem-quasi-DM}
设 $\mathcal{X}$ 是代数叠。下列条件等价：
\begin{enumerate}
\item $\mathcal{X}$ 是拟 DM 叠；
\item 存在概形 $W$ 以及满、平坦、局部有限表现且局部拟有限的态射
$W \to \mathcal{X}$。
\end{enumerate}
\end{theorem}

\begin{proof}
(2) $\Rightarrow$ (1) 即引理
\ref{stacks-morphisms-lemma-properties-covering-imply-diagonal}。假设 (1)。
设 $x \in |\mathcal{X}|$ 是有限型点。我们将构造一个在 $\mathcal{X}$
上、并在 $x$ 的某个邻域内“起作用”的概形。证明末尾取所有这些概形的
不交并即可得出结论。

\medskip\noindent
设 $U$ 是仿射概形，$U \to \mathcal{X}$ 是光滑态射，而
$u \in U$ 是映到 $x$ 的闭点；参见引理
\ref{stacks-morphisms-lemma-point-finite-type}。照例记 $u = \Spec(\kappa(u))$。
考虑交换图
$$
\xymatrix{
u \ar[d] & R \ar[l] \ar[d] \\
U \ar[d] & F \ar[d] \ar[l]^p \\
\mathcal{X} & u \ar[l]
}
$$
其中两个方块都是纤维积方块；特别地，$R = u \times_\mathcal{X} u$。
在引理 \ref{stacks-morphisms-lemma-point-finite-type-monomorphism} 的证明中，已经看到
$(u, R, s, t, c)$ 是代数空间中的群胚，且 $s, t$ 局部有限型。
令 $G \to u$ 为稳定子群代数空间（参见《空间中的群胚》定义
\ref{spaces-groupoids-definition-stabilizer-groupoid}).
注意
$$
G = R \times_{(u \times u)} u =
(u \times_\mathcal{X} u) \times_{(u \times u)} u =
\mathcal{X} \times_{\mathcal{X} \times \mathcal{X}} u.
$$
由于 $\mathcal{X}$ 为拟 DM，$G$ 在 $u$ 上局部拟有限。由
《空间中的群胚进阶》引理
\ref{spaces-more-groupoids-lemma-groupoid-on-field-dimension-equal-stabilizer}
可得 $\dim(R) = 0$。

\medskip\noindent
令 $e : u \to R$ 为该群胚的单位元。复合 $u \to R \to u$ 等于
$u$ 的恒等态射。由于 $u \subset U$ 是闭子概形，$R \subset F$
是闭子空间。因此，也可把 $e$ 看作 $F$ 的点，并将其记为 $z$。
考虑 \'etale 局部环之间的映射
$$
\mathcal{O}_{U, u}
\xrightarrow{p^\sharp}
\mathcal{O}_{F, \overline{z}}
\longrightarrow
\mathcal{O}_{R, \overline{e}}
$$
由第一段的结果，$\mathcal{O}_{R, \overline{e}}$ 的维数为 $0$。
另一方面，由于 $R$ 在 $F$ 中由 $\mathfrak m_u$ 截出，第二个箭头的核为
$p^\sharp(\mathfrak m_u)\mathcal{O}_{F, \overline{z}}$。因此
$$
\mathfrak m_{\overline{z}} =
\sqrt{p^\sharp(\mathfrak m_u)\mathcal{O}_{F, \overline{z}}}
$$
另一方面，由于态射 $U \to \mathcal{X}$ 光滑，$F \to u$ 是代数空间的
光滑态射。这意味着 $F$ 是正则代数空间（《域上的空间》引理
\ref{spaces-over-fields-lemma-smooth-regular}）。因此
$\mathcal{O}_{F, \overline{z}}$ 是正则局部环（《空间的性质》引理
\ref{spaces-properties-lemma-regular}）。注意，正则局部环是
Cohen--Macaulay 环（《代数》引理
\ref{algebra-lemma-regular-ring-CM}）。令
$d = \dim(\mathcal{O}_{F, \overline{z}})$。由《代数》引理
\ref{algebra-lemma-find-sequence-image-regular}，可找到
$f_1, \ldots, f_d \in \mathcal{O}_{U, u}$，使其像
$\varphi(f_1), \ldots, \varphi(f_d)$ 在
$\mathcal{O}_{F, \overline{z}}$ 中构成正则序列。由引理
\ref{stacks-morphisms-lemma-slice}，缩小 $U$ 后可假设
$Z = V(f_1, \ldots, f_d) \to \mathcal{X}$ 平坦且局部有限表现。
注意，按构造，$F_Z = Z \times_\mathcal{X} u$ 是
$F = U \times_\mathcal{X} u$ 的闭子空间，$z$ 是该闭子空间的点，并且
$$
\dim(\mathcal{O}_{F_Z, \overline{z}}) = 0.
$$
由于 $z$ 相对于 $u$ 的超越次数为零，由《空间的态射》引理
\ref{spaces-morphisms-lemma-dimension-fibre-at-a-point}
可得 $\dim_z(F_Z) = 0$。因而由引理
\ref{stacks-morphisms-lemma-quasi-finite-at-point} 可知，必要时再次缩小 $U$ 后，态射
$Z \to \mathcal{X}$ 局部拟有限。

\medskip\noindent
综上，对 $\mathcal{X}$ 的每个有限型点 $x$，存在局部拟有限、平坦且
局部有限表现的态射 $f_x : Z_x \to \mathcal{X}$，使 $x$ 属于
$|f_x|$ 的像。令 $W = \coprod_x Z_x$ 及 $f = \coprod f_x$。
则 $f$ 平坦、局部有限表现且局部拟有限。特别地，$|f|$ 的像是开集；
参见《代数叠的性质》引理
\ref{stacks-properties-lemma-topology-points}。按构造，该像包含
$\mathcal{X}$ 的所有有限型点，故由引理
\ref{stacks-morphisms-lemma-enough-finite-type-points}（以及《代数叠的性质》引理
\ref{stacks-properties-lemma-characterize-surjective}）可知 $f$ 满。
\end{proof}

\begin{lemma}
\label{stacks-morphisms-lemma-DM-residual-gerbe}
设 $\mathcal{Z}$ 是 DM、局部诺特且约化的代数叠，并且
$|\mathcal{Z}|$ 是单点集。则存在域 $k$ 以及满的 \'etale 态射
$\Spec(k) \to \mathcal{Z}$。
\end{lemma}

\begin{proof}
由《代数叠的性质》引理
\ref{stacks-properties-lemma-unique-point-better}，存在域 $k$ 以及满、平坦且
局部有限表现的态射 $\Spec(k) \to \mathcal{Z}$。令
$U = \Spec(k)$ 及 $R = U \times_\mathcal{Z} U$，于是得到代数空间中的群胚
$(U, R, s, t, c)$；参见《代数叠》引理
\ref{algebraic-lemma-criterion-map-representable-spaces-fibred-in-groupoids}.
注意，由《代数叠》评注
\ref{algebraic-remark-flat-fp-presentation}，有等价
$$
f_{can} : [U/R] \longrightarrow \mathcal{Z}
$$
投影 $s, t : R \to U$ 局部有限表现。由于 $\mathcal{Z}$ 为 DM，
稳定子群代数空间
$$
G = U \times_{U \times U} R = U \times_{U \times U} (U \times_\mathcal{Z} U) =
U \times_{\mathcal{Z} \times \mathcal{Z}, \Delta_\mathcal{Z}} \mathcal{Z}
$$
在 $U$ 上非分歧。特别地，$\dim(G) = 0$；由《空间中的群胚进阶》引理
\ref{spaces-more-groupoids-lemma-groupoid-on-field-dimension-equal-stabilizer}
可得 $\dim(R) = 0$。这蕴含 $R$ 是概形；参见《域上的空间》引理
\ref{spaces-over-fields-lemma-locally-finite-type-dim-zero}。由《簇》引理
\ref{varieties-lemma-algebraic-scheme-dim-0} 可知，$R$（以及 $G$）是若干
Artin 局部环之谱的不交并，而这些环经 $s$ 或 $t$ 在 $k$ 上有限。令
$P = \Spec(A) \subset R$ 为既开又闭的子概形，其底层点是群胚概形
$(U, R, s, t, c)$ 的单位元 $e$。由于
$s \circ e = t \circ e = \text{id}_{\Spec(k)}$，$A$ 是 Artin 局部环，
其剩余域经 $s^\sharp : k \to A$ 或 $t^\sharp : k \to A$ 与 $k$ 等同。
注意，$s, t : \Spec(A) \to \Spec(k)$ 都是有限态射（由上述引理）。
由于 $G \to \Spec(k)$ 非分歧，
$$
G \cap P = P \times_{U \times U} U = \Spec(A \otimes_{k \otimes k} k)
$$
在 $k$ 上非分歧。另一方面，$A \otimes_{k \otimes k} k$ 作为 $A$ 的商环
是局部环，并满射到 $k$。于是
$A \otimes_{k \otimes k} k = k$。由此，$P \to U \times U$ 泛单射
（因为 $P$ 只有一个点，且其剩余域为 $k$）、非分歧（由上面对唯一像点
之上的纤维的计算）且有限型（因为 $s, t$ 有限型），因而是单态射（参见
《\'Etale 态射》引理
\ref{etale-lemma-universally-injective-unramified}).
因此，$s|_P, t|_P : P \to U$ 定义有限平坦等价关系。应用《群胚》命题
\ref{groupoids-proposition-finite-flat-equivalence} 可知 $U/P$ 存在，且为概形
$\overline{U}$。此外，$U \to \overline{U}$ 有限局部自由，并且
$P = U \times_{\overline{U}} U$。事实上，
$\overline{U} = \Spec(k_0)$，其中 $k_0 \subset k$ 是 $R$-不变函数所成的环。
由于 $k$ 是域，由定义《群胚》等式
(\ref{groupoids-equation-invariants}) 可知 $k_0$ 是域。

\medskip\noindent
我们断言
\begin{equation}
\label{stacks-morphisms-equation-etale-covering}
\Spec(k_0) = \overline{U} = U/P \to [U/R] = \mathcal{Z}
\end{equation}
是所需的满 \'etale 态射。由《代数叠的性质》引理
\ref{stacks-properties-lemma-flat-cover-by-field}，该态射满。因此只需证明
(\ref{stacks-morphisms-equation-etale-covering}) 为 \'etale\footnote{我们诚请读者自行给出
这一事实的证明。事实上，该论证与《自举》定理
\ref{bootstrap-theorem-final-bootstrap} 的证明末尾的论证有很多共同之处，
因而或许应当在某处将其单独提炼为一个引理。}。我们不直接证明 \'etale 性，
而是先应用《自举》引理 \ref{bootstrap-lemma-divide-subgroupoid}，可知存在群胚概形
$(\overline{U}, \overline{R}, \overline{s}, \overline{t}, \overline{c})$
使 $(U, R, s, t, c)$ 是
$(\overline{U}, \overline{R}, \overline{s}, \overline{t}, \overline{c})$
经商态射 $U \to \overline{U}$ 所作的限制。（上面已验证该引理的全部假设，
只有 $j : R \to U \times U$ 分离且局部拟有限这一断言尚未验证；它由
$R$ 是在 $k$ 上局部拟有限的分离概形这一事实得出。）由于
$U \to \overline{U}$ 有限局部自由，$[U/R] \to [\overline{U}/\overline{R}]$
是等价；参见《空间中的群胚》引理
\ref{spaces-groupoids-lemma-quotient-stack-restrict-equivalence}。

\medskip\noindent
注意，$s, t$ 是态射 $\overline{s}, \overline{t}$ 沿
$U \to \overline{U}$ 的基变换。由于 $\{U \to \overline{U}\}$ 是 fppf 覆盖，
可知 $\overline{s}, \overline{t}$ 平坦、局部有限表现且局部拟有限；参见
《下降》引理 \ref{descent-lemma-descending-property-flat}、
\ref{descent-lemma-descending-property-locally-finite-presentation} 及
\ref{descent-lemma-descending-property-quasi-finite}。
考虑交换图
$$
\xymatrix{
U \times_{\overline{U}} U \ar@{=}[r] \ar[rd] & P \ar[r] \ar[d] & R \ar[d] \\
& \overline{U} \ar[r]^{\overline{e}} & \overline{R}
}
$$
关于限制的一般事实说明，外侧四个角构成笛卡儿图。由该等式可知内侧方块
也是笛卡儿的。由于 $P$ 在 $R$ 中为开集，由《下降》引理
\ref{descent-lemma-descending-property-open-immersion} 可知
$\overline{e}$ 是开浸入。

\medskip\noindent
当然，若 $\overline{e}$ 是开浸入，且 $\overline{s}, \overline{t}$ 平坦并
局部有限表现，则态射 $\overline{t}, \overline{s}$ 为 \'etale。例如，应用
《群胚进阶》引理 \ref{more-groupoids-lemma-sheaf-differentials} 可看出这一点：
该引理说明 $\Omega_{\overline{R}/\overline{U}} = 0$ 蕴含
$\overline{s}, \overline{t} : \overline{R} \to \overline{U}$ 非分歧（参见
《态射》引理 \ref{morphisms-lemma-unramified-omega-zero}），进而蕴含
$\overline{s}, \overline{t}$ 为 \'etale（参见《态射》引理
\ref{morphisms-lemma-flat-unramified-etale}）。因此，
$\mathcal{Z} = [\overline{U}/\overline{R}]$ 是代数叠 $\mathcal{Z}$ 的一个
\'etale 表现，并由《代数叠的性质》引理
\ref{stacks-properties-lemma-check-property-covering}
可知 $\overline{U} \to \mathcal{Z}$ 为 \'etale。
\end{proof}

\begin{lemma}
\label{stacks-morphisms-lemma-etale-at-point}
设 $\mathcal{X}$ 是代数叠。考虑笛卡儿图
$$
\xymatrix{
U \ar[d] & F \ar[l]^p \ar[d] \\
\mathcal{X} & \Spec(k) \ar[l]
}
$$
其中 $U$ 是代数空间，$k$ 是域，且 $U \to \mathcal{X}$ 平坦并局部有限表现。
设 $z \in |F|$，使 $F \to \Spec(k)$ 在 $z$ 处非分歧。则以一个包含
$p(z)$ 的开子空间替换 $U$ 后，态射
$$
U \longrightarrow \mathcal{X}
$$
为 \'etale。
\end{lemma}

\begin{proof}
由于 $f : U \to \mathcal{X}$ 平坦且局部有限表现，存在极大开集
$W(f) \subset U$，使限制 $f|_{W(f)} : W(f) \to \mathcal{X}$ 为
\'etale；参见《代数叠的性质》评注
\ref{stacks-properties-remark-local-source-apply}
(\ref{stacks-properties-item-etale}).
因此，只需证明 $p(z)$ 是 $W(f)$ 的点。此外，上述评注还说明，
$W(f)$ 的构造与沿可由代数空间表示的态射所作的任意基变换交换。
故只需证明态射 $F \to \Spec(k)$ 在 $z$ 处为 \'etale。它作为
$U \to \mathcal{X}$ 的基变换是平坦且局部有限表现的，并且按假设
$F \to \Spec(k)$ 在 $z$ 处非分歧，故结论由《空间的态射》引理
\ref{spaces-morphisms-lemma-unramified-flat-lfp-etale} 得出。
\end{proof}

\noindent
DM 叠是 Deligne--Mumford 叠。

\begin{theorem}
\label{stacks-morphisms-theorem-DM}
设 $\mathcal{X}$ 是代数叠。下列条件等价：
\begin{enumerate}
\item $\mathcal{X}$ 为 DM；
\item $\mathcal{X}$ 为 Deligne--Mumford；
\item 存在概形 $W$ 以及满 \'etale 态射 $W \to \mathcal{X}$。
\end{enumerate}
\end{theorem}

\begin{proof}
回忆，(3) 是 (2) 的定义；参见《代数叠》定义
\ref{algebraic-definition-deligne-mumford}。(3) $\Rightarrow$ (1) 即引理
\ref{stacks-morphisms-lemma-properties-covering-imply-diagonal}。假设 (1)。设
$x \in |\mathcal{X}|$ 是有限型点。我们将构造一个在 $\mathcal{X}$ 上、
并在 $x$ 的某个邻域内“起作用”的概形。证明末尾取所有这些概形的不交并
即可得出结论。

\medskip\noindent
由引理 \ref{stacks-morphisms-lemma-point-finite-type-monomorphism}，$\mathcal{X}$ 在 $x$
处的剩余胚 $\mathcal{Z}_x$ 存在，且 $\mathcal{Z}_x \to \mathcal{X}$
局部有限型。由引理 \ref{stacks-morphisms-lemma-separation-properties-residual-gerbe}，
代数叠 $\mathcal{Z}_x$ 为 DM。由引理
\ref{stacks-morphisms-lemma-DM-residual-gerbe}，存在域 $k$ 以及满 \'etale 态射
$z : \Spec(k) \to \mathcal{Z}_x$。特别地，复合
$x : \Spec(k) \to \mathcal{X}$ 局部有限型（由《空间的态射》引理
\ref{spaces-morphisms-lemma-composition-finite-type} 和
\ref{spaces-morphisms-lemma-etale-locally-finite-type}）。

\medskip\noindent
选取概形 $U$ 以及光滑态射 $U \to \mathcal{X}$，使 $x$ 属于
$|U| \to |\mathcal{X}|$ 的像。考虑纤维积方块
$$
\xymatrix{
U \ar[d] & F  \ar[l] \ar[d] \\
\mathcal{X} & \Spec(k) \ar[l]_-x
}
$$
换言之，$F = U \times_{\mathcal{X}, x} \Spec(k)$。由《代数叠的性质》引理
\ref{stacks-properties-lemma-points-cartesian} 可知 $F$ 非空。
由于 $\mathcal{Z}_x \to \mathcal{X}$ 是单态射，有
$$
\Spec(k) \times_{z, \mathcal{Z}_x, z} \Spec(k)
=
\Spec(k) \times_{x, \mathcal{X}, x} \Spec(k)
$$
按 $z$ 的构造，等式两端到 $\Spec(k)$ 的投影为 \'etale。由于
$$
F \times_U F =
(\Spec(k) \times_\mathcal{X} \Spec(k))
\times_{\Spec(k)} F
$$
可知投影映射 $F \times_U F \to F$ 也为 \'etale。因此
$\Delta_{F/U} : F \to F \times_U F$ 为 \'etale（参见《空间的态射》引理
\ref{spaces-morphisms-lemma-etale-permanence}）。由《空间的态射》引理
\ref{spaces-morphisms-lemma-etale-universally-injective-open}
可知 $\Delta_{F/U}$ 是开浸入，进而由《空间的态射》引理
\ref{spaces-morphisms-lemma-diagonal-unramified-morphism}
可知 $F \to U$ 非分歧。

\medskip\noindent
选取非空仿射概形 $V$ 以及 \'etale 态射 $V \to F$。（也可直接使用 $F$
而避免这一步，但这样做似乎更容易说明所发生的情形。）图示如下：
$$
\xymatrix{
U \ar[d] & F  \ar[l] \ar[d] & V \ar[l] \ar[ld] \\
\mathcal{X} & \Spec(k) \ar[l]_-x
}
$$
于是 $V \to \Spec(k)$ 是概形的光滑态射，而 $V \to U$ 是概形的非分歧态射
（参见《空间的态射》引理
\ref{spaces-morphisms-lemma-composition-smooth} 和
\ref{spaces-morphisms-lemma-composition-unramified}）。
选取闭点 $v \in V$，使 $k \subset \kappa(v)$ 有限可分；参见《簇》引理
\ref{varieties-lemma-smooth-separable-closed-points-dense}。令 $u \in U$
为其像点。局部环 $\mathcal{O}_{V, v}$ 正则（参见《簇》引理
\ref{varieties-lemma-smooth-regular}），且来自态射 $V \to U$ 的局部环同态
$$
\varphi : \mathcal{O}_{U, u} \longrightarrow \mathcal{O}_{V, v}
$$
满足 $\varphi(\mathfrak m_u)\mathcal{O}_{V, v} = \mathfrak m_v$；参见
《态射》引理 \ref{morphisms-lemma-unramified-at-point}。因此可找到
$f_1, \ldots, f_d \in \mathcal{O}_{U, u}$，使各像
$\varphi(f_1), \ldots, \varphi(f_d)$ 构成
$\mathfrak m_v/\mathfrak m_v^2$ 在 $\kappa(v)$ 上的一组基。由于
$\mathcal{O}_{V, v}$ 是正则局部环，这蕴含
$\varphi(f_1), \ldots, \varphi(f_d)$ 在 $\mathcal{O}_{V, v}$ 中构成
正则序列（参见《代数》引理 \ref{algebra-lemma-regular-ring-CM}）。
以 $u$ 的一个开邻域替换 $U$ 后，可假设
$f_1, \ldots, f_d \in \Gamma(U, \mathcal{O}_U)$。必要时再以 $u$ 的一个
更小开邻域替换 $U$，由引理 \ref{stacks-morphisms-lemma-slice} 可假设
$V(f_1, \ldots, f_d) \to \mathcal{X}$ 平坦且局部有限表现。按构造，
$$
V(f_1, \ldots, f_d) \times_\mathcal{X} \Spec(k)
\longleftarrow
V(f_1, \ldots, f_d) \times_U V
$$
为 \'etale，且 $V(f_1, \ldots, f_d) \times_U V$ 是由
$f_1|_V, \ldots, f_d|_V$ 截出的闭子概形 $T \subset V$。因此按构造
$v \in T$，并且
$$
\mathcal{O}_{T, v} =
\mathcal{O}_{V, v}/(\varphi(f_1), \ldots, \varphi(f_d)) = \kappa(v)
$$
是 $k$ 的有限可分扩张。由此 $T \to \Spec(k)$ 在 $v$ 处非分歧；参见
《态射》引理 \ref{morphisms-lemma-unramified-at-point}。根据代数空间的
非分歧态射的定义，这意味着
$V(f_1, \ldots, f_d) \times_\mathcal{X} \Spec(k) \to \Spec(k)$
在 $v$ 于 $V(f_1, \ldots, f_d) \times_\mathcal{X} \Spec(k)$ 中的像处
非分歧。应用引理 \ref{stacks-morphisms-lemma-etale-at-point} 可知，再将 $U$ 缩小为
$u$ 的另一个开邻域后，态射 $V(f_1, \ldots, f_d) \to \mathcal{X}$
为 \'etale。

\medskip\noindent
综上，对 $\mathcal{X}$ 的每个有限型点 $x$，存在 \'etale 态射
$f_x : W_x \to \mathcal{X}$，使 $x$ 属于 $|f_x|$ 的像。令
$W = \coprod_x W_x$ 及 $f = \coprod f_x$。则 $f$ 为 \'etale。特别地，
$|f|$ 的像是开集；参见《代数叠的性质》引理
\ref{stacks-properties-lemma-topology-points}。按构造，该像包含
$\mathcal{X}$ 的所有有限型点，故由引理
\ref{stacks-morphisms-lemma-enough-finite-type-points}（以及《代数叠的性质》引理
\ref{stacks-properties-lemma-characterize-surjective}）可知 $f$ 满。
\end{proof}

\noindent
下面这个有用推论说明，代数叠的 DM 态射的“纤维”是 Deligne--Mumford 的。

\begin{lemma}
\label{stacks-morphisms-lemma-DM}
设 $f : \mathcal{X} \to \mathcal{Y}$ 是代数叠的 DM 态射。则
\begin{enumerate}
\item 对每个 DM 代数叠 $\mathcal{Z}$ 和态射
$\mathcal{Z} \to \mathcal{Y}$，存在一个概形以及满 \'etale 态射
$U \to \mathcal{X} \times_\mathcal{Y} \mathcal{Z}$.
\item 对每个代数空间 $Z$ 和态射 $Z \to \mathcal{Y}$，存在一个概形以及
满 \'etale 态射
$U \to \mathcal{X} \times_\mathcal{Y} Z$.
\end{enumerate}
\end{lemma}

\begin{proof}
(1) 的证明。由于 $f$ 为 DM，由引理
\ref{stacks-morphisms-lemma-base-change-separated}，基变换
$\mathcal{X} \times_\mathcal{Y} \mathcal{Z} \to \mathcal{Z}$ 为 DM。
又因 $\mathcal{Z}$ 为 DM，引理 \ref{stacks-morphisms-lemma-separated-over-separated}
蕴含 $\mathcal{X} \times_\mathcal{Y} \mathcal{Z}$ 为 DM。因此存在概形
$U$ 以及满 \'etale 态射
$U \to \mathcal{X} \times_\mathcal{Y} \mathcal{Z}$；参见定理
\ref{stacks-morphisms-theorem-DM}。第 (2) 项是 (1) 的特例，因为代数空间在视为代数叠时
由引理 \ref{stacks-morphisms-lemma-trivial-implications} 可知是 DM。
\end{proof}






\section{Deligne--Mumford 轨迹}
\label{stacks-morphisms-section-DM}

\noindent
每个代数叠都有一个最大的开子叠，并且该子叠是 Deligne--Mumford 叠；
这基本显然，但我们仍在下面写出证明。当然，该子叠可能为空，例如当
$X = [\Spec(\mathbf{Z})/\mathbf{G}_{m, \mathbf{Z}}]$ 时如此。
下面将刻画 DM 轨迹中的点。

\begin{lemma}
\label{stacks-morphisms-lemma-open-DM-locus}
设 $\mathcal{X}$ 是代数叠。存在开子叠
$$
\mathcal{X}'' \subset \mathcal{X}' \subset \mathcal{X}
$$
使 $\mathcal{X}''$ 为 DM，$\mathcal{X}'$ 为拟 DM，并且它们分别是具有
这些性质的最大开子叠。
\end{lemma}

\begin{proof}
这里真正要说的只是：若 $\mathcal{U} \subset \mathcal{X}$ 和
$\mathcal{V} \subset \mathcal{X}$ 是 DM 开子叠，则满足
$|\mathcal{W}| = |\mathcal{U}| \cup |\mathcal{V}|$ 的开子叠
$\mathcal{W} \subset \mathcal{X}$ 也为 DM。（拟 DM 的情形类似。）
尽管这有些取巧，我们用定理 \ref{stacks-morphisms-theorem-DM} 来证明。由该定理，
可选取概形 $U$、$V$ 以及满 \'etale 态射
$U \to \mathcal{U}$、$V \to \mathcal{V}$。于是当然
$U \amalg V \to \mathcal{W}$ 满且为 \'etale。拟 DM 的情形使用定理
\ref{stacks-morphisms-theorem-quasi-DM} 以完全相同的方法证明。
\end{proof}

\begin{lemma}
\label{stacks-morphisms-lemma-points-DM-locus}
设 $\mathcal{X}$ 是代数叠。设 $x \in |\mathcal{X}|$ 对应于
$x : \Spec(k) \to \mathcal{X}$。令 $G_x/k$ 为 $x$ 的自同构群代数空间。
则
\begin{enumerate}
\item $x$ 属于 $\mathcal{X}$ 的 DM 轨迹，当且仅当
$G_x \to \Spec(k)$ 非分歧；
\item $x$ 属于 $\mathcal{X}$ 的拟 DM 轨迹，当且仅当
$G_x \to \Spec(k)$ 局部拟有限。
\end{enumerate}
\end{lemma}

\begin{proof}
(2) 的证明。选取概形 $U$ 以及满光滑态射 $U \to \mathcal{X}$。
考虑纤维积
$$
\xymatrix{
G \ar[r] \ar[d] & \mathcal{I}_\mathcal{X} \ar[d] \\
U \ar[r] & \mathcal{X}
}
$$
回忆，$G$ 是 $U \to \mathcal{X}$ 的自同构群代数空间。由
《空间中的群胚》引理
\ref{spaces-groupoids-lemma-open-over-which-unramified-or-lqf}
可知，存在最大开子概形 $U' \subset U$，使 $G_{U'} \to U'$ 局部拟有限。
此外，$U'$ 的构造与任意基变换交换。特别地，$U'$ 在
$R = U \times_\mathcal{X} U$ 中的两个逆像是 $R$ 的同一个开子空间
（毕竟，两个映射 $R \to \mathcal{X}$ 同构，因而有同构的自同构群空间）。
因此，由《代数叠的性质》引理
\ref{stacks-properties-lemma-substacks-presentation}
可知，$U'$ 是某个开子叠 $\mathcal{X}' \subset \mathcal{X}$ 的逆像，
并且有笛卡儿图
$$
\xymatrix{
G_{U'} \ar[r] \ar[d] & \mathcal{I}_{\mathcal{X}'} \ar[d] \\
U' \ar[r] & \mathcal{X}'
}
$$
因此，态射 $\mathcal{I}_{\mathcal{X}'} \to \mathcal{X}'$ 局部拟有限，
由引理 \ref{stacks-morphisms-lemma-diagonal-diagonal} 第 (5) 项可知 $\mathcal{X}'$ 为拟 DM。
另一方面，若 $\mathcal{W} \subset \mathcal{X}$ 是拟 DM 开子叠，则由
$U'$ 的构造，$\mathcal{W}$ 的逆像 $W \subset U$ 必须包含于
$U'$，因为
$\mathcal{I}_\mathcal{W} =
\mathcal{W} \times_\mathcal{X} \mathcal{I}_\mathcal{X}$
在 $\mathcal{W}$ 上局部拟有限。因此，$\mathcal{X}'$ 是拟 DM 轨迹。
最后，选取域扩张 $K/k$ 以及 $2$-交换图
$$
\xymatrix{
\Spec(K) \ar[r] \ar[d] & \Spec(k) \ar[d]^x \\
U \ar[r] & \mathcal{X}
}
$$
于是得到 $K$ 上群代数空间的同构
$G_x \times_{\Spec(k)} \Spec(K) \cong G \times_U \Spec(K)$
。因此，$G_x$ 在 $k$ 上局部拟有限，当且仅当
$\Spec(K) \to U$ 的像落在 $U'$ 中（为此使用 $U'$ 的构造与基变换交换，
以及对 $\Spec(K) \to \Spec(k)$ 和 $G_x$ 应用《空间中的群胚》引理
\ref{spaces-groupoids-lemma-open-over-which-unramified-or-lqf}
）。这就完成了 (2) 的证明。(1) 的证明完全相同。
\end{proof}









\section{局部拟有限态射}
\label{stacks-morphisms-section-locally-quasi-finite}

\noindent
代数空间态射的“局部拟有限”性质不是在源与靶上光滑局部的，故不能利用
第 \ref{stacks-morphisms-section-local-source-target} 节的材料来定义代数叠的局部拟有限态射。
对于可由代数空间表示的代数叠态射，我们已经知道局部拟有限的含义；参见
《代数叠的性质》一节
\ref{stacks-properties-section-properties-morphisms}.
为了找到适用于一般态射的条件，我们作如下观察。

\begin{lemma}
\label{stacks-morphisms-lemma-representable-by-spaces-quasi-finite}
设 $f : \mathcal{X} \to \mathcal{Y}$ 是代数叠的态射。假设 $f$ 可由
代数空间表示。下列条件等价：
\begin{enumerate}
\item $f$ 局部拟有限（其含义见《代数叠的性质》一节
\ref{stacks-properties-section-properties-morphisms}）；
\item $f$ 局部有限型，且对每个态射 $\Spec(k) \to \mathcal{Y}$，
其中 $k$ 为域，空间 $|\Spec(k) \times_\mathcal{Y} \mathcal{X}|$ 离散。
\end{enumerate}
\end{lemma}

\begin{proof}
假设 (1)。在此情形，代数空间的态射
$\mathcal{X}_k \to \Spec(k)$ 作为 $f$ 的基变换局部拟有限。因此由
《空间的态射》引理 \ref{spaces-morphisms-lemma-locally-quasi-finite}，
$|\mathcal{X}_k|$ 离散。反之，假设 (2)。选取满光滑态射
$V \to \mathcal{Y}$，其中 $V$ 是概形。只需证明代数空间态射
$V \times_\mathcal{Y} \mathcal{X} \to V$ 局部拟有限；参见
《代数叠的性质》引理
\ref{stacks-properties-lemma-check-property-covering}.
按假设，态射 $V \times_\mathcal{Y} \mathcal{X} \to V$ 局部有限型。
对任意态射 $\Spec(k) \to V$，其中 $k$ 为域，有
$$
\Spec(k) \times_V (V \times_\mathcal{Y} \mathcal{X}) =
\Spec(k) \times_\mathcal{Y} \mathcal{X}
$$
按假设，其点空间离散。因此，由《空间的态射》引理
\ref{spaces-morphisms-lemma-locally-quasi-finite} 可知
$V \times_\mathcal{Y} \mathcal{X} \to V$ 局部拟有限。
\end{proof}

\noindent
可由代数空间表示的代数叠态射为拟 DM；参见引理
\ref{stacks-morphisms-lemma-trivial-implications}。结合上述引理可知，在可由代数空间表示的
态射这一情形，下述定义与已有概念并不冲突。

\begin{definition}
\label{stacks-morphisms-definition-locally-quasi-finite}
设 $f : \mathcal{X} \to \mathcal{Y}$ 是代数叠的态射。若 $f$ 为拟 DM、
局部有限型，并且对每个态射 $\Spec(k) \to \mathcal{Y}$，其中 $k$ 为域，
空间 $|\mathcal{X}_k|$ 离散，则称 $f$ {\it 局部拟有限}。
\end{definition}

\noindent
要求 $f$ 为拟 DM 是自然的。例如，设 $k$ 为域，并考虑态射
$\pi : [\Spec(k)/\mathbf{G}_m] \to \Spec(k)$
。它具有单点纤维，并且局部有限型。稍后将看到，该态射光滑且相对维数为
$-1$，而我们希望局部拟有限态射的相对维数为 $0$。还要注意，截面
$\Spec(k) \to [\Spec(k)/\mathbf{G}_m]$ 不具有离散纤维，故不局部拟有限；
而我们希望局部拟有限态射具有如下恒久性质：若
$f : \mathcal{X} \to \mathcal{X}'$ 是代数叠的态射，并且在代数叠
$\mathcal{Y}$ 上局部拟有限，则 $f$ 局部拟有限（事实上有稍强的结论，
参见引理 \ref{stacks-morphisms-lemma-locally-quasi-finite-permanence}）。

\medskip\noindent
上述定义的另一个理由是下面的引理
\ref{stacks-morphisms-lemma-characterize-locally-quasi-finite}；它用 $\mathcal{X}$ 和
$\mathcal{Y}$ 的适当“表现”或“覆盖”的存在性刻画局部拟有限性。

\begin{lemma}
\label{stacks-morphisms-lemma-base-change-locally-quasi-finite}
局部拟有限态射的基变换局部拟有限。
\end{lemma}

\begin{proof}
对拟 DM 态射，这已见于引理 \ref{stacks-morphisms-lemma-base-change-separated}；
对局部有限型态射，这已见于引理
\ref{stacks-morphisms-lemma-base-change-finite-type}。关于纤维的条件显然由基变换继承。
\end{proof}

\begin{lemma}
\label{stacks-morphisms-lemma-locally-quasi-finite-over-field}
设 $\mathcal{X} \to \Spec(k)$ 是局部拟有限态射，其中 $\mathcal{X}$
为代数叠，$k$ 为域。设 $f : V \to \mathcal{X}$ 是局部拟有限态射，
其中 $V$ 为概形。则 $V \to \Spec(k)$ 局部拟有限。
\end{lemma}

\begin{proof}
由引理 \ref{stacks-morphisms-lemma-composition-finite-type} 可知
$V \to \Spec(k)$ 局部有限型。反设 $V \to \Spec(k)$ 不局部拟有限，
以求矛盾。则 $V$ 的点之间存在非平凡特化 $v \leadsto v'$；参见《态射》引理
\ref{morphisms-lemma-quasi-finite-at-point-characterize}。特别地，
$\text{trdeg}_k(\kappa(v)) > \text{trdeg}_k(\kappa(v'))$；参见《态射》引理
\ref{morphisms-lemma-dimension-fibre-specialization}。由于
$|\mathcal{X}|$ 离散，有 $|f|(v) = |f|(v')$。考虑
$R = V \times_\mathcal{X} V$。则 $R$ 是代数空间，而投影
$s, t : R \to V$ 作为 $V \to \mathcal{X}$ 的基变换局部拟有限
（后者可由代数空间表示，故这由《代数叠的性质》一节中的讨论得出
\ref{stacks-properties-section-properties-morphisms}）。由《代数叠的性质》引理
\ref{stacks-properties-lemma-points-cartesian} 可知，存在 $r \in |R|$，
使 $s(r) = v$ 且 $t(r) = v'$。由《空间的态射》引理
\ref{spaces-morphisms-lemma-compare-tr-deg} 可知，$v/k$、$r/k$ 与
$v'/k$ 的超越次数相等。此矛盾证明了引理。
\end{proof}

\begin{lemma}
\label{stacks-morphisms-lemma-composition-locally-quasi-finite}
局部拟有限态射的复合局部拟有限。
\end{lemma}

\begin{proof}
对拟 DM 态射，这已见于引理 \ref{stacks-morphisms-lemma-composition-separated}；
对局部有限型态射，这已见于引理
\ref{stacks-morphisms-lemma-composition-finite-type}。设 $\mathcal{X} \to \mathcal{Y}$ 和
$\mathcal{Y} \to \mathcal{Z}$ 局部拟有限。设 $k$ 为域，并设
$\Spec(k) \to \mathcal{Z}$ 为态射。只需证明
$|\mathcal{X}_k|$ 离散。由引理
\ref{stacks-morphisms-lemma-base-change-locally-quasi-finite}，态射
$\mathcal{X}_k \to \mathcal{Y}_k$ 和
$\mathcal{Y}_k \to \Spec(k)$ 局部拟有限。特别地，$\mathcal{Y}_k$
是拟 DM 代数叠；参见引理
\ref{stacks-morphisms-lemma-separated-implies-morphism-separated}。由定理
\ref{stacks-morphisms-theorem-quasi-DM}，可找到概形 $V$ 以及满、平坦、局部有限表现且
局部拟有限的态射 $V \to \mathcal{Y}_k$。由引理
\ref{stacks-morphisms-lemma-locally-quasi-finite-over-field} 可知，$V$ 在 $k$ 上局部拟有限；
特别地，$|V|$ 离散。态射
$V \times_{\mathcal{Y}_k} \mathcal{X}_k \to \mathcal{X}_k$ 满、平坦且
局部有限表现，故
$|V \times_{\mathcal{Y}_k} \mathcal{X}_k| \to |\mathcal{X}_k|$
满且开。因此只需证明 $|V \times_{\mathcal{Y}_k} \mathcal{X}_k|$ 离散。
注意，$V$ 是若干 Artin 局部 $k$-代数 $A_i$ 之谱的不交并，其剩余域为
$k_i$；参见《簇》引理
\ref{varieties-lemma-algebraic-scheme-dim-0}。因此，只需证明每个
$$
|\Spec(A_i) \times_{\mathcal{Y}_k} \mathcal{X}_k| =
|\Spec(k_i) \times_{\mathcal{Y}_k} \mathcal{X}_k| =
|\Spec(k_i) \times_\mathcal{Y} \mathcal{X}|
$$
离散；这由 $\mathcal{X} \to \mathcal{Y}$ 局部拟有限这一假设得出。
\end{proof}

\noindent
在用覆盖刻画局部拟有限态射之前，我们先对拟 DM 态射这样做。

\begin{lemma}
\label{stacks-morphisms-lemma-characterize-quasi-DM}
设 $f : \mathcal{X} \to \mathcal{Y}$ 是代数叠的态射。下列条件等价：
\begin{enumerate}
\item $f$ 为拟 DM；
\item 对任意态射 $V \to \mathcal{Y}$，其中 $V$ 为代数空间，存在满、平坦、
局部有限表现且局部拟有限的态射
$U \to \mathcal{X} \times_\mathcal{Y} V$，其中 $U$ 为代数空间；
\item 存在代数空间 $U$、$V$，满、平坦且局部有限表现的态射
$V \to \mathcal{Y}$，以及满、平坦、局部有限表现且局部拟有限的态射
$U \to \mathcal{X} \times_\mathcal{Y} V$。
\end{enumerate}
\end{lemma}

\begin{proof}
(2) $\Rightarrow$ (3) 是显然的。

\medskip\noindent
假设 (1)，并设 $V \to \mathcal{Y}$ 如 (2) 所述。则
$\mathcal{X} \times_\mathcal{Y} V \to V$ 为拟 DM；参见引理
\ref{stacks-morphisms-lemma-base-change-separated}。由引理
\ref{stacks-morphisms-lemma-trivial-implications}，代数空间 $V$ 为 DM，因而为拟 DM。
于是，由引理 \ref{stacks-morphisms-lemma-separated-over-separated}，
$\mathcal{X} \times_\mathcal{Y} V$ 为拟 DM。因此可以应用定理
\ref{stacks-morphisms-theorem-quasi-DM} 得到 (2) 所述的态射
$U \to \mathcal{X} \times_\mathcal{Y} V$。

\medskip\noindent
假设 (3)。设 $V \to \mathcal{Y}$ 和
$U \to \mathcal{X} \times_\mathcal{Y} V$ 如 (3) 所述。为证明 $f$ 为拟 DM，
只需证明 $\mathcal{X} \times_\mathcal{Y} V \to V$ 为拟 DM；参见引理
\ref{stacks-morphisms-lemma-check-separated-covering}。由引理
\ref{stacks-morphisms-lemma-properties-covering-imply-diagonal} 可知
$\mathcal{X} \times_\mathcal{Y} V$ 为拟 DM。于是由引理
\ref{stacks-morphisms-lemma-separated-implies-morphism-separated}，
$\mathcal{X} \times_\mathcal{Y} V \to V$ 为拟 DM，故 (1) 成立。
这就完成了引理的证明。
\end{proof}

\begin{lemma}
\label{stacks-morphisms-lemma-characterize-locally-quasi-finite}
设 $f : \mathcal{X} \to \mathcal{Y}$ 是代数叠的态射。下列条件等价：
\begin{enumerate}
\item $f$ 局部拟有限；
\item $f$ 为拟 DM，并且对任意代数空间 $V$ 给出的态射
$V \to \mathcal{Y}$，以及任意代数空间 $U$ 给出的局部拟有限态射
$U \to \mathcal{X} \times_\mathcal{Y} V$，态射 $U \to V$ 局部拟有限；
\item 对代数空间 $V$ 给出的任意态射 $V \to \mathcal{Y}$，存在满、平坦、
局部有限表现且局部拟有限的态射
$U \to \mathcal{X} \times_\mathcal{Y} V$，其中 $U$ 是代数空间，
并且 $U \to V$ 局部拟有限；
\item 存在代数空间 $U$、$V$，满、平坦且局部有限表现的态射
$V \to \mathcal{Y}$，以及满、平坦、局部有限表现且局部拟有限的态射
$U \to \mathcal{X} \times_\mathcal{Y} V$，并且 $U \to V$ 局部拟有限。
\end{enumerate}
\end{lemma}

\begin{proof}
假设 (1)。按假设，$f$ 为拟 DM。设 $V \to \mathcal{Y}$ 和
$U \to \mathcal{X} \times_\mathcal{Y} V$ 如 (2) 所述。由引理
\ref{stacks-morphisms-lemma-composition-locally-quasi-finite}，复合
$U \to \mathcal{X} \times_\mathcal{Y} V \to V$ 局部拟有限。
因此 (1) 蕴含 (2)。

\medskip\noindent
假设 (2)。设 $V \to \mathcal{Y}$ 如 (3) 所述。由引理
\ref{stacks-morphisms-lemma-characterize-quasi-DM}，可找到代数空间 $U$ 以及满、平坦、
局部有限表现且局部拟有限的态射
$U \to \mathcal{X} \times_\mathcal{Y} V$。由 (2)，复合
$U \to V$ 局部拟有限。因此 (2) 蕴含 (3)。

\medskip\noindent
(3) 蕴含 (4) 是显然的。

\medskip\noindent
假设 (4)。我们将证明 (1) 成立，从而完成证明。由引理
\ref{stacks-morphisms-lemma-characterize-quasi-DM} 可知 $f$ 为拟 DM。为证明 $f$ 局部有限型，
只需证明 $g : \mathcal{X} \times_\mathcal{Y} V \to V$ 局部有限型；参见引理
\ref{stacks-morphisms-lemma-check-finite-type-covering}。进而只需检验 $g$ 与
$h : U \to \mathcal{X} \times_\mathcal{Y} V$ 预复合后局部有限型；参见引理
\ref{stacks-morphisms-lemma-check-finite-type-precompose}。由于按假设
$g \circ h : U \to V$ 局部拟有限，这一点成立，故 $f$ 局部有限型。
最后，设 $k$ 为域，并设 $\Spec(k) \to \mathcal{Y}$ 为态射。则
$V \times_\mathcal{Y} \Spec(k)$ 是非空代数空间，且在 $k$ 上局部有限表现。
因此可找到有限扩张 $k'/k$ 以及态射 $\Spec(k') \to V$，使
$$
\xymatrix{
\Spec(k') \ar[r] \ar[d] & V \ar[d] \\
\Spec(k) \ar[r] & \mathcal{Y}
}
$$
交换（细节略）。于是 $\mathcal{X}_{k'} \to \mathcal{X}_k$ 可由概形表示、
满且有限局部自由。特别地，
$|\mathcal{X}_{k'}| \to |\mathcal{X}_k|$ 满且开。因此只需证明
$|\mathcal{X}_{k'}|$ 离散。由于
$$
U \times_V \Spec(k') =
U \times_{\mathcal{X} \times_\mathcal{Y} V} \mathcal{X}_{k'}
$$
可知 $U \times_V \Spec(k') \to \mathcal{X}_{k'}$ 满、平坦且局部有限表现
（作为 $U \to \mathcal{X} \times_\mathcal{Y} V$ 的基变换）。因此
$|U \times_V \Spec(k')| \to |\mathcal{X}_{k'}|$ 满且开。故只需证明
$|U \times_V \Spec(k')|$ 离散。这由 $U \to V$ 局部拟有限这一事实得出
（或者使用上面的定义，或者由代数空间态射的原定义并借助《空间的态射》引理
\ref{spaces-morphisms-lemma-locally-quasi-finite} 得出）。
\end{proof}

\begin{lemma}
\label{stacks-morphisms-lemma-locally-quasi-finite-permanence}
设 $\mathcal{X} \to \mathcal{Y} \to \mathcal{Z}$ 是代数叠的态射。
假设 $\mathcal{X} \to \mathcal{Z}$ 局部拟有限，且
$\mathcal{Y} \to \mathcal{Z}$ 为拟 DM。则
$\mathcal{X} \to \mathcal{Y}$ 局部拟有限。
\end{lemma}

\begin{proof}
把 $\mathcal{X} \to \mathcal{Y}$ 写成复合
$$
\mathcal{X} \longrightarrow
\mathcal{X} \times_\mathcal{Z} \mathcal{Y} \longrightarrow
\mathcal{Y}
$$
第二个箭头作为 $\mathcal{X} \to \mathcal{Z}$ 的基变换局部拟有限；参见引理
\ref{stacks-morphisms-lemma-base-change-locally-quasi-finite}。由于
$\mathcal{Y} \to \mathcal{Z}$ 为拟 DM，由引理
\ref{stacks-morphisms-lemma-semi-diagonal}，第一个箭头局部拟有限。因此由引理
\ref{stacks-morphisms-lemma-composition-locally-quasi-finite}，
$\mathcal{X} \to \mathcal{Y}$ 局部拟有限。
\end{proof}

















\section{拟有限态射}
\label{stacks-morphisms-section-quasi-finite}

\noindent
我们已在第 \ref{stacks-morphisms-section-locally-quasi-finite} 节定义了代数叠的“局部拟有限”
态射，并在第 \ref{stacks-morphisms-section-quasi-compact} 节定义了代数叠的“拟紧”态射。
根据定义，代数空间的态射拟有限，当且仅当它既局部拟有限又拟紧
（《空间的态射》定义
\ref{spaces-morphisms-definition-locally-quasi-finite}),
）。因此，可如下定义代数叠态射拟有限的含义；当该态射可由代数空间表示时，
它与《代数叠的性质》一节
\ref{stacks-properties-section-properties-morphisms} 中已经定义的概念一致。

\begin{definition}
\label{stacks-morphisms-definition-quasi-finite}
\begin{reference}
\cite{rydh_approx}
\end{reference}
设 $f : \mathcal{X} \to \mathcal{Y}$ 是代数叠的态射。若 $f$ 局部拟有限
（定义 \ref{stacks-morphisms-definition-locally-quasi-finite}）且拟紧
（定义 \ref{stacks-morphisms-definition-quasi-compact}），则称 $f$ {\it 拟有限}。
\end{definition}

\begin{lemma}
\label{stacks-morphisms-lemma-composition-quasi-finite}
拟有限态射的复合拟有限。
\end{lemma}

\begin{proof}
合并引理 \ref{stacks-morphisms-lemma-composition-locally-quasi-finite} 和
\ref{stacks-morphisms-lemma-composition-quasi-compact}。
\end{proof}

\begin{lemma}
\label{stacks-morphisms-lemma-base-change-quasi-finite}
拟有限态射的基变换拟有限。
\end{lemma}

\begin{proof}
合并引理 \ref{stacks-morphisms-lemma-base-change-locally-quasi-finite} 和
\ref{stacks-morphisms-lemma-base-change-quasi-compact}。
\end{proof}

\begin{lemma}
\label{stacks-morphisms-lemma-quasi-finite-permanence}
设 $f : \mathcal{X} \to \mathcal{Y}$ 和
$g : \mathcal{Y} \to \mathcal{Z}$ 是代数叠的态射。若 $g \circ f$ 拟有限，
且 $g$ 拟分离并为拟 DM，则 $f$ 拟有限。
\end{lemma}

\begin{proof}
合并引理 \ref{stacks-morphisms-lemma-locally-quasi-finite-permanence} 和
\ref{stacks-morphisms-lemma-quasi-compact-permanence}。
\end{proof}












\section{平坦态射}
\label{stacks-morphisms-section-flat}

\noindent
代数空间态射的“平坦”性质在源与靶上光滑局部；参见《空间上的下降》评注
\ref{spaces-descent-remark-list-local-source-target}。它还在基变换下稳定，
并且在靶上是 fpqc 局部的；参见《空间的态射》引理
\ref{spaces-morphisms-lemma-base-change-flat}
以及《空间上的下降》引理
\ref{spaces-descent-lemma-descending-property-flat}.
因此，由上述引理 \ref{stacks-morphisms-lemma-local-source-target}，可如下定义代数空间态射
何谓平坦；当该态射可由代数空间表示时，这与《代数叠的性质》一节
\ref{stacks-properties-section-properties-morphisms} 中已有的定义一致。

\begin{definition}
\label{stacks-morphisms-definition-flat}
设 $f : \mathcal{X} \to \mathcal{Y}$ 是代数叠的态射。若引理
\ref{stacks-morphisms-lemma-local-source-target} 的等价条件对
$\mathcal{P} = \text{flat}$ 成立，则称 $f$ {\it 平坦}。
\end{definition}

\begin{lemma}
\label{stacks-morphisms-lemma-composition-flat}
平坦态射的复合平坦。
\end{lemma}

\begin{proof}
合并评注 \ref{stacks-morphisms-remark-composition} 与《空间的态射》引理
\ref{spaces-morphisms-lemma-composition-flat}。
\end{proof}

\begin{lemma}
\label{stacks-morphisms-lemma-base-change-flat}
平坦态射的基变换平坦。
\end{lemma}

\begin{proof}
合并评注 \ref{stacks-morphisms-remark-base-change} 与《空间的态射》引理
\ref{spaces-morphisms-lemma-base-change-flat}。
\end{proof}

\begin{lemma}
\label{stacks-morphisms-lemma-descent-flat}
设 $f : \mathcal{X} \to \mathcal{Y}$ 是代数叠的态射。设
$\mathcal{Z} \to \mathcal{Y}$ 是代数叠的满平坦态射。若基变换
$\mathcal{Z} \times_\mathcal{Y} \mathcal{X} \to \mathcal{Z}$ 平坦，
则 $f$ 平坦。
\end{lemma}

\begin{proof}
选取代数空间 $W$ 以及满光滑态射 $W \to \mathcal{Z}$。则
$W \to \mathcal{Z}$ 满且平坦（《空间的态射》引理
\ref{spaces-morphisms-lemma-smooth-flat}），故 $W \to \mathcal{Y}$ 满且平坦
（由《代数叠的性质》引理
\ref{stacks-properties-lemma-composition-surjective}
和引理 \ref{stacks-morphisms-lemma-composition-flat}）。由于
$\mathcal{Z} \times_\mathcal{Y} \mathcal{X} \to \mathcal{Z}$
沿 $W \to \mathcal{Z}$ 的基变换是平坦态射（引理
\ref{stacks-morphisms-lemma-base-change-flat}），可用 $W$ 替换 $\mathcal{Z}$。

\medskip\noindent
选取代数空间 $V$ 以及满光滑态射 $V \to \mathcal{Y}$。选取代数空间 $U$
以及满光滑态射 $U \to V \times_\mathcal{Y} \mathcal{X}$。需要证明
$U \to V$ 平坦。现在沿 $W \to \mathcal{Y}$ 对所有对象作基变换：令
$U' = W \times_\mathcal{Y} U$，
$V' = W \times_\mathcal{Y} V$,
$\mathcal{X}' = W \times_\mathcal{Y} \mathcal{X}$,
以及 $\mathcal{Y}' = W \times_\mathcal{Y} \mathcal{Y} = W$。由基变换，
$U' \to V' \times_{\mathcal{Y}'} \mathcal{X}'$ 仍然光滑。因此，由我们对
代数叠的平坦态射的定义，以及 $\mathcal{X}' \to \mathcal{Y}'$ 平坦这一假设，
可知 $U' \to V'$ 平坦。又因 $V' \to V$ 作为
$W \to \mathcal{Y}$ 的基变换满，由《空间的态射》引理
\ref{spaces-morphisms-lemma-base-change-module-flat} (2)
可知 $U \to V$ 平坦，得证。
\end{proof}

\begin{lemma}
\label{stacks-morphisms-lemma-flat-permanence}
设 $\mathcal{X} \to \mathcal{Y} \to \mathcal{Z}$ 是代数叠的态射。
若 $\mathcal{X} \to \mathcal{Z}$ 平坦，且
$\mathcal{X} \to \mathcal{Y}$ 满且平坦，则
$\mathcal{Y} \to \mathcal{Z}$ 平坦。
\end{lemma}

\begin{proof}
选取代数空间 $W$ 以及满光滑态射 $W \to \mathcal{Z}$。选取代数空间 $V$
以及满光滑态射 $V \to W \times_\mathcal{Z} \mathcal{Y}$。选取代数空间
$U$ 以及满光滑态射 $U \to V \times_\mathcal{Y} \mathcal{X}$。
我们知道 $U \to V$ 平坦，且 $U \to W$ 平坦。此外，由于
$\mathcal{X} \to \mathcal{Y}$ 满，$U \to V$ 也满（作为满态射的复合）。
因此，引理归结为代数空间态射的情形。该情形即《空间的态射》引理
\ref{spaces-morphisms-lemma-flat-permanence}.
\end{proof}

\begin{lemma}
\label{stacks-morphisms-lemma-lift-valuation-ring-through-flat-morphism}
设 $f : \mathcal{X} \to \mathcal{Y}$ 是代数叠的平坦态射。设
$\Spec(A) \to \mathcal{Y}$ 为态射，其中 $A$ 是赋值环。若
$\Spec(A)$ 的闭点映到 $|\mathcal{Y}|$ 中属于
$|\mathcal{X|} \to |\mathcal{Y}|$ 之像的点，则存在交换图
$$
\xymatrix{
\Spec(A') \ar[r] \ar[d] & \mathcal{X} \ar[d] \\
\Spec(A) \ar[r] & \mathcal{Y}
}
$$
其中 $A \to A'$ 是赋值环扩张（《代数进阶》定义
\ref{more-algebra-definition-extension-valuation-rings}).
\end{lemma}

\begin{proof}
基变换 $\mathcal{X}_A \to \Spec(A)$ 平坦（引理
\ref{stacks-morphisms-lemma-base-change-flat}），且 $\Spec(A)$ 的闭点属于
$|\mathcal{X}_A| \to |\Spec(A)|$ 的像（《代数叠的性质》引理
\ref{stacks-properties-lemma-points-cartesian}）。因此可假设
$\mathcal{Y} = \Spec(A)$。设 $U \to \mathcal{X}$ 为满光滑态射，
其中 $U$ 为概形。于是对态射 $U \to \Spec(A)$ 应用《空间的态射》引理
\ref{spaces-morphisms-lemma-lift-valuation-ring-through-flat-morphism}
即得结论。
\end{proof}







\section{在一点处平坦}
\label{stacks-morphisms-section-flat-at-point}

\noindent
我们仍需发展一般机制，才能说明代数叠的态射在一点处具有某个给定性质的
含义。目前，下述引理已经足够。

\begin{lemma}
\label{stacks-morphisms-lemma-flat-at-point}
设 $f : \mathcal{X} \to \mathcal{Y}$ 是代数叠的态射。设
$x \in |\mathcal{X}|$。考虑交换图
$$
\vcenter{
\xymatrix{
U \ar[d]_a \ar[r]_h & V \ar[d]^b \\
\mathcal{X} \ar[r]^f & \mathcal{Y}
}
}
\quad\text{带标记点}
\vcenter{
\xymatrix{
u \in |U| \ar[d] \\
x \in |\mathcal{X}|
}
}
$$
其中 $U$、$V$ 是代数空间，$b$ 平坦，并且
$(a, h) : U \to \mathcal{X} \times_\mathcal{Y} V$
平坦。下列条件等价：
\begin{enumerate}
\item 对某一个上述图，$h$ 在 $u$ 处平坦；
\item 对每一个上述图，$h$ 在 $u$ 处平坦。
\end{enumerate}
\end{lemma}

\begin{proof}
假设给定引理中所述的第二个图 $U', V', u', a', b', h'$。于是可以考虑
$$
\xymatrix{
U \ar[d] & U \times_\mathcal{X} U' \ar[l] \ar[d] \ar[r] & U' \ar[d] \\
V & V \times_\mathcal{Y} V' \ar[l] \ar[r] & V'
}
$$
由《代数叠的性质》引理
\ref{stacks-properties-lemma-points-cartesian}，存在点
$u'' \in |U \times_\mathcal{X} U'|$ 映到 $u$ 和 $u'$。若 $h$ 在 $u$ 处平坦，
则基变换
$U \times_V (V \times_\mathcal{Y} V') \to V \times_\mathcal{Y} V'$
在 $u$ 上方的任一点处平坦；参见《空间的态射》引理
\ref{spaces-morphisms-lemma-base-change-module-flat}。另一方面，态射
$$
U \times_\mathcal{X} U' \to
U \times_\mathcal{X} (\mathcal{X} \times_\mathcal{Y} V') =
U \times_\mathcal{Y} V' =
U \times_V (V \times_\mathcal{Y} V')
$$
作为 $(a', h')$ 的基变换平坦；参见引理
\ref{stacks-morphisms-lemma-base-change-flat}。作复合并应用《空间的态射》引理
\ref{spaces-morphisms-lemma-composition-module-flat}，可知
$U \times_\mathcal{X} U' \to V \times_\mathcal{Y} V'$ 在 $u''$ 处平坦。
再与平坦映射 $V \times_\mathcal{Y} V' \to V'$ 复合，可知
$U \times_\mathcal{X} U' \to V'$ 在 $u''$ 处平坦。最后，由于
$U \times_\mathcal{X} U' \to U'$ 在 $u''$ 处平坦，且 $u''$ 映到 $u'$，
由《空间的态射》引理 \ref{spaces-morphisms-lemma-flat-permanence} 可知
$U' \to V'$ 在 $u'$ 处平坦。
\end{proof}

\begin{definition}
\label{stacks-morphisms-definition-flat-at-point}
设 $f : \mathcal{X} \to \mathcal{Y}$ 是代数叠的态射。设
$x \in |\mathcal{X}|$。若引理 \ref{stacks-morphisms-lemma-flat-at-point} 的等价条件成立，
则称 $f$ {\it 在 $x$ 处平坦}。
\end{definition}





\section{有限表现态射}
\label{stacks-morphisms-section-finite-presentation}

\noindent
代数空间态射的“局部有限表现”性质在源与靶上光滑局部；参见
《空间上的下降》评注
\ref{spaces-descent-remark-list-local-source-target}。它还在基变换下稳定，
并且在靶上是 fpqc 局部的；参见《空间的态射》引理
\ref{spaces-morphisms-lemma-base-change-finite-presentation} 以及
《空间上的下降》引理
\ref{spaces-descent-lemma-descending-property-locally-finite-presentation}.
因此，由上述引理 \ref{stacks-morphisms-lemma-local-source-target}，可如下定义代数叠态射
何谓局部有限表现；当该态射可由代数空间表示时，这与《代数叠的性质》一节
\ref{stacks-properties-section-properties-morphisms} 中已有的定义一致。

\begin{definition}
\label{stacks-morphisms-definition-locally-finite-presentation}
设 $f : \mathcal{X} \to \mathcal{Y}$ 是代数叠的态射。
\begin{enumerate}
\item 若引理 \ref{stacks-morphisms-lemma-local-source-target} 的等价条件对
$\mathcal{P} = \text{locally of finite presentation}$ 成立，则称 $f$
{\it 局部有限表现}。
\item 若该态射局部有限表现、拟紧且拟分离，则称 $f$ {\it 有限表现}。
\end{enumerate}
\end{definition}

\noindent
注意，有限表现态射{\bf 不只是}一个局部有限表现的拟紧态射。

\begin{lemma}
\label{stacks-morphisms-lemma-composition-finite-presentation}
有限表现态射的复合有限表现。对局部有限表现态射也有同样结论。
\end{lemma}

\begin{proof}
合并评注 \ref{stacks-morphisms-remark-composition} 与《空间的态射》引理
\ref{spaces-morphisms-lemma-composition-finite-presentation}。
\end{proof}

\begin{lemma}
\label{stacks-morphisms-lemma-base-change-finite-presentation}
有限表现态射的基变换有限表现。对局部有限表现态射也有同样结论。
\end{lemma}

\begin{proof}
合并评注 \ref{stacks-morphisms-remark-base-change} 与《空间的态射》引理
\ref{spaces-morphisms-lemma-base-change-finite-presentation}。
\end{proof}

\begin{lemma}
\label{stacks-morphisms-lemma-finite-presentation-finite-type}
局部有限表现态射局部有限型。有限表现态射有限型。
\end{lemma}

\begin{proof}
合并评注 \ref{stacks-morphisms-remark-implication} 与《空间的态射》引理
\ref{spaces-morphisms-lemma-finite-presentation-finite-type}。
\end{proof}

\begin{lemma}
\label{stacks-morphisms-lemma-noetherian-finite-type-finite-presentation}
设 $f : \mathcal{X} \to \mathcal{Y}$ 是代数叠的态射。
\begin{enumerate}
\item 若 $\mathcal{Y}$ 局部诺特且 $f$ 局部有限型，则 $f$ 局部有限表现。
\item 若 $\mathcal{Y}$ 局部诺特，且 $f$ 有限型并拟分离，则 $f$ 有限表现。
\end{enumerate}
\end{lemma}

\begin{proof}
假设 $f : \mathcal{X} \to \mathcal{Y}$ 局部有限型，且 $\mathcal{Y}$
局部诺特。这意味着存在引理 \ref{stacks-morphisms-lemma-local-source-target} 中所述的图，
其中 $h$ 局部有限型，竖直箭头 $a$ 满。由《空间的态射》引理
\ref{spaces-morphisms-lemma-noetherian-finite-type-finite-presentation}
可知 $h$ 局部有限表现。因此按定义，
$\mathcal{X} \to \mathcal{Y}$ 局部有限表现。这证明了 (1)。若 $f$
有限型且拟分离，则它也拟紧且拟分离，故 (2) 立即得出。
\end{proof}

\begin{lemma}
\label{stacks-morphisms-lemma-finite-presentation-permanence}
设 $f : \mathcal{X} \to \mathcal{Y}$ 和
$g : \mathcal{Y} \to \mathcal{Z}$ 是代数叠的态射。若 $g \circ f$
局部有限表现，且 $g$ 局部有限型，则 $f$ 局部有限表现。
\end{lemma}

\begin{proof}
选取代数空间 $W$ 以及满光滑态射 $W \to \mathcal{Z}$。选取代数空间 $V$
以及满光滑态射 $V \to \mathcal{Y} \times_\mathcal{Z} W$。选取代数空间
$U$ 以及满光滑态射 $U \to \mathcal{X} \times_\mathcal{Y} V$。
对态射 $U \to V \to W$ 应用《空间的态射》引理
\ref{spaces-morphisms-lemma-finite-presentation-permanence}
即得引理。
\end{proof}

\begin{lemma}
\label{stacks-morphisms-lemma-diagonal-morphism-finite-type}
设 $f : \mathcal{X} \to \mathcal{Y}$ 是代数叠的态射，其对角态射为
$\Delta : \mathcal{X} \to \mathcal{X} \times_\mathcal{Y} \mathcal{X}$。
若 $f$ 局部有限型，则 $\Delta$ 局部有限表现。若 $f$ 拟分离且局部有限型，
则 $\Delta$ 有限表现。
\end{lemma}

\begin{proof}
注意，$\Delta$ 是 $\mathcal{X}$ 上的态射（经第二投影）。若 $f$ 局部有限型，
则 $\mathcal{X}$ 在 $\mathcal{X}$ 上有限表现，且由引理
\ref{stacks-morphisms-lemma-base-change-finite-type}，
$\text{pr}_2 : \mathcal{X} \times_\mathcal{Y} \mathcal{X} \to \mathcal{X}$
局部有限型。因此，第一个陈述由引理
\ref{stacks-morphisms-lemma-finite-presentation-permanence} 得出。第二个陈述由第一个陈述
和定义得出（因为按定义，$f$ 拟分离意味着 $\Delta_f$ 拟紧且拟分离）。
\end{proof}

\begin{lemma}
\label{stacks-morphisms-lemma-open-immersion-locally-finite-presentation}
开浸入局部有限表现。
\end{lemma}

\begin{proof}
由《代数叠的性质》定义
\ref{stacks-properties-definition-immersion}
以及《空间的态射》引理
\ref{spaces-morphisms-lemma-open-immersion-locally-finite-presentation}
即得。
\end{proof}

\begin{lemma}
\label{stacks-morphisms-lemma-check-property-after-fppf-base-change}
设 $P$ 是代数空间态射的一个性质，它在靶上为 fppf 局部的，并在任意基变换
下保持。设 $f : \mathcal{X} \to \mathcal{Y}$ 是可由代数空间表示的代数叠态射。
设 $\mathcal{Z} \to \mathcal{Y}$ 是满、平坦且局部有限表现的代数叠态射。
令 $\mathcal{W} = \mathcal{Z} \times_\mathcal{Y} \mathcal{X}$。则
$$
(f\text{ 具有 }P) \Leftrightarrow
(\text{the projection }\mathcal{W} \to \mathcal{Z}\text{ 具有 }P).
$$
该陈述的含义见《代数叠的性质》一节
\ref{stacks-properties-section-properties-morphisms}.
\end{lemma}

\begin{proof}
选取代数空间 $W$ 以及满、平坦且局部有限表现的态射
$W \to \mathcal{Z}$。由《代数叠的性质》引理
\ref{stacks-properties-lemma-composition-surjective}
以及引理 \ref{stacks-morphisms-lemma-composition-flat} 和
\ref{stacks-morphisms-lemma-composition-finite-presentation}
可知，复合 $W \to \mathcal{Y}$ 也满、平坦且局部有限表现。记
$V = W \times_\mathcal{Z} \mathcal{W} = V \times_\mathcal{Y} \mathcal{X}$.
由《代数叠的性质》引理
\ref{stacks-properties-lemma-check-property-covering}
可知，$f$ 具有 $\mathcal{P}$ 当且仅当 $V \to W$ 具有该性质，
而 $\mathcal{W} \to \mathcal{Z}$ 具有 $\mathcal{P}$ 当且仅当
$V \to W$ 具有该性质。引理得证。
\end{proof}

\begin{lemma}
\label{stacks-morphisms-lemma-descent-property}
设 $\mathcal{P}$ 是代数空间态射的一个性质，它在源与靶上光滑局部，
并且在靶上为 fppf 局部。设 $f : \mathcal{X} \to \mathcal{Y}$ 是代数叠的态射。
设 $\mathcal{Z} \to \mathcal{Y}$ 是代数叠的满、平坦且局部有限表现的态射。
若基变换 $\mathcal{Z} \times_\mathcal{Y} \mathcal{X} \to \mathcal{Z}$
具有 $\mathcal{P}$，则 $f$ 具有 $\mathcal{P}$。
\end{lemma}

\begin{proof}
假设 $\mathcal{Z} \times_\mathcal{Y} \mathcal{X} \to \mathcal{Z}$
具有 $\mathcal{P}$。选取代数空间 $W$ 以及满光滑态射
$W \to \mathcal{Z}$。注意
$W \times_\mathcal{Z} \mathcal{Z} \times_\mathcal{Y} \mathcal{X} =
W \times_\mathcal{Y} \mathcal{X}$。根据
$\mathcal{Z} \times_\mathcal{Y} \mathcal{X} \to \mathcal{Z}$ 具有
$\mathcal{P}$ 的定义本身（参见定义 \ref{stacks-morphisms-definition-P} 和引理
\ref{stacks-morphisms-lemma-local-source-target}），可知
$W \times_\mathcal{Y} \mathcal{X} \to W$ 具有 $\mathcal{P}$。另一方面，
$W \to \mathcal{Z}$ 满、平坦且局部有限表现（《空间的态射》引理
\ref{spaces-morphisms-lemma-smooth-flat} 和
\ref{spaces-morphisms-lemma-smooth-locally-finite-presentation}）
），故 $W \to \mathcal{Y}$ 满、平坦且局部有限表现（由《代数叠的性质》引理
\ref{stacks-properties-lemma-composition-surjective}
以及引理 \ref{stacks-morphisms-lemma-composition-flat} 和
\ref{stacks-morphisms-lemma-composition-finite-presentation}）。
因此可用 $W$ 替换 $\mathcal{Z}$。

\medskip\noindent
选取代数空间 $V$ 以及满光滑态射 $V \to \mathcal{Y}$。选取代数空间 $U$
以及满光滑态射 $U \to V \times_\mathcal{Y} \mathcal{X}$。需要证明
$U \to V$ 具有 $\mathcal{P}$。现在沿 $W \to \mathcal{Y}$ 对所有对象
作基变换：令
$U' = W \times_\mathcal{Y} U$,
$V' = W \times_\mathcal{Y} V$,
$\mathcal{X}' = W \times_\mathcal{Y} \mathcal{X}$,
以及 $\mathcal{Y}' = W \times_\mathcal{Y} \mathcal{Y} = W$。由基变换，
$U' \to V' \times_{\mathcal{Y}'} \mathcal{X}'$ 仍然光滑。因此，在定义
$\mathcal{X}' \to \mathcal{Y}' = W$ 具有 $\mathcal{P}$ 时应用引理
\ref{stacks-morphisms-lemma-local-source-target}，可知 $U' \to V'$ 具有 $\mathcal{P}$。
又因 $V' \to V$ 作为 $W \to \mathcal{Y}$ 的基变换满、平坦且局部有限表现，
而 $\mathcal{P}$ 在靶上对 fppf 拓扑是局部的，故
$U \to V$ 具有 $\mathcal{P}$。
\end{proof}

\begin{lemma}
\label{stacks-morphisms-lemma-descent-finite-presentation}
设 $f : \mathcal{X} \to \mathcal{Y}$ 是代数叠的态射。设
$\mathcal{Z} \to \mathcal{Y}$ 是代数叠的满、平坦且局部有限表现的态射。
若基变换 $\mathcal{Z} \times_\mathcal{Y} \mathcal{X} \to \mathcal{Z}$
局部有限表现，则 $f$ 局部有限表现。
\end{lemma}

\begin{proof}
“局部有限表现”性质满足引理 \ref{stacks-morphisms-lemma-descent-property} 的条件。
在源与靶上光滑局部这一点已见于本节引言，而在靶上为 fppf 局部这一点即
《空间上的下降》引理
\ref{spaces-descent-lemma-descending-property-locally-finite-presentation}.
\end{proof}

\begin{lemma}
\label{stacks-morphisms-lemma-flat-finite-presentation-permanence}
设 $\mathcal{X} \to \mathcal{Y} \to \mathcal{Z}$ 是代数叠的态射。
若 $\mathcal{X} \to \mathcal{Z}$ 局部有限表现，且
$\mathcal{X} \to \mathcal{Y}$ 满、平坦且局部有限表现，则
$\mathcal{Y} \to \mathcal{Z}$ 局部有限表现。
\end{lemma}

\begin{proof}
选取代数空间 $W$ 以及满光滑态射 $W \to \mathcal{Z}$。选取代数空间 $V$
以及满光滑态射 $V \to W \times_\mathcal{Z} \mathcal{Y}$。选取代数空间
$U$ 以及满光滑态射 $U \to V \times_\mathcal{Y} \mathcal{X}$。我们知道
$U \to V$ 平坦且局部有限表现，而 $U \to W$ 局部有限表现。此外，
由于 $\mathcal{X} \to \mathcal{Y}$ 满，$U \to V$ 也满（作为满态射的复合）。
因此，引理归结为代数空间态射的情形。该情形即《空间上的下降》引理
\ref{spaces-descent-lemma-locally-finite-presentation-fppf-local-source}.
\end{proof}

\begin{lemma}
\label{stacks-morphisms-lemma-surjective-flat-locally-finite-presentation}
设 $f : \mathcal{X} \to \mathcal{Y}$ 是满、平坦且局部有限表现的代数叠态射。
则对每个概形 $U$ 以及 $\mathcal{Y}$ 在 $U$ 上的对象 $y$，存在 fppf 覆盖
$\{U_i \to U\}$ 以及 $\mathcal{X}$ 在 $U_i$ 上的对象 $x_i$，使在
$\mathcal{Y}_{U_i}$ 中有 $f(x_i) \cong y|_{U_i}$。
\end{lemma}

\begin{proof}
可把 $y$ 看作态射 $U \to \mathcal{Y}$。由《代数叠的性质》引理
\ref{stacks-properties-lemma-base-change-surjective}
以及引理 \ref{stacks-morphisms-lemma-base-change-finite-presentation} 和
\ref{stacks-morphisms-lemma-base-change-flat}
可知 $\mathcal{X} \times_\mathcal{Y} U \to U$ 满、平坦且局部有限表现。
设 $V$ 为概形，并设 $V \to \mathcal{X} \times_\mathcal{Y} U$ 光滑且满。
则 $V \to \mathcal{X} \times_\mathcal{Y} U$ 也满、平坦且局部有限表现
（参见《空间的态射》引理
\ref{spaces-morphisms-lemma-smooth-flat} 和
\ref{spaces-morphisms-lemma-smooth-locally-finite-presentation}）。
因此 $V \to U$ 也满、平坦且局部有限表现；参见《代数叠的性质》引理
\ref{stacks-properties-lemma-composition-surjective}
以及引理 \ref{stacks-morphisms-lemma-composition-finite-presentation} 和
\ref{stacks-morphisms-lemma-composition-flat}.
因此，$\{V \to U\}$ 是所需的 fppf 覆盖，而
$x : V \to \mathcal{X}$ 是所需对象。
\end{proof}

\begin{lemma}
\label{stacks-morphisms-lemma-surjective-family-flat-locally-finite-presentation}
设 $f_j : \mathcal{X}_j \to \mathcal{X}$（$j \in J$）是一族代数叠态射，
每个都平坦且局部有限表现，并且它们联合为满，即
$|\mathcal{X}| = \bigcup |f_j|(|\mathcal{X}_j|)$.
则对每个概形 $U$ 以及 $\mathcal{X}$ 在 $U$ 上的对象 $x$，存在 fppf 覆盖
$\{U_i \to U\}_{i \in I}$、映射 $a : I \to J$，以及
$\mathcal{X}_{a(i)}$ 在 $U_i$ 上的对象 $x_i$，使在
$\mathcal{X}_{U_i}$ 中有 $f_{a(i)}(x_i) \cong y|_{U_i}$。
\end{lemma}

\begin{proof}
对态射 $\coprod_{j \in J} \mathcal{X}_j \to \mathcal{X}$ 应用引理
\ref{stacks-morphisms-lemma-surjective-flat-locally-finite-presentation}。（由于我们的设置，
这里有一个轻微的集合论问题；我们将其忽略。）最后注意，态射
$x_i : U_i \to \coprod_{j \in J} \mathcal{X}_j$ 由不交并分解
$U_i = \coprod U_{i, j}$ 和态射 $U_{i, j} \to \mathcal{X}_j$ 给出。
于是 fppf 覆盖 $\{U_{i, j} \to U\}$ 以及态射
$U_{i, j} \to \mathcal{X}_j$ 满足要求。
\end{proof}

\begin{lemma}
\label{stacks-morphisms-lemma-fppf-open}
设 $f : \mathcal{X} \to \mathcal{Y}$ 平坦且局部有限表现。则
$|f| : |\mathcal{X}| \to |\mathcal{Y}|$ 是开映射。
\end{lemma}

\begin{proof}
选取概形 $V$ 以及满光滑态射 $V \to \mathcal{Y}$。选取概形 $U$ 以及
满光滑态射 $U \to V \times_\mathcal{Y} \mathcal{X}$。按假设，概形态射
$U \to V$ 平坦且局部有限表现。因此，由《态射》引理
\ref{morphisms-lemma-fppf-open}，$U \to V$ 是开映射。根据
$|\mathcal{Y}|$ 上拓扑的构造，映射 $|V| \to |\mathcal{Y}|$ 是开的。
映射 $|U| \to |\mathcal{X}|$ 满。由初等拓扑，这些事实蕴含结论。
\end{proof}

\begin{lemma}
\label{stacks-morphisms-lemma-descent-quasi-compact}
设 $f : \mathcal{X} \to \mathcal{Y}$ 是代数叠的态射。设
$\mathcal{Z} \to \mathcal{Y}$ 是代数叠的满、平坦且局部有限表现的态射。
若基变换 $\mathcal{Z} \times_\mathcal{Y} \mathcal{X} \to \mathcal{Z}$
拟紧，则 $f$ 拟紧。
\end{lemma}

\begin{proof}
需要证明：对任意态射 $\mathcal{Y}' \to \mathcal{Y}$，若
$\mathcal{Y}'$ 拟紧，则 $\mathcal{Y}' \times_\mathcal{Y} \mathcal{X}$ 拟紧。
记 $\mathcal{Z}' = \mathcal{Z} \times_\mathcal{Y} \mathcal{Y}'$。
由引理 \ref{stacks-morphisms-lemma-fppf-open}，
$|\mathcal{Z}'| \to |\mathcal{Y}'|$ 是开映射。因此可找到映满
$\mathcal{Y}'$ 的拟紧开子叠 $\mathcal{W} \subset \mathcal{Z}'$。由于
$\mathcal{Z} \times_\mathcal{Y} \mathcal{X} \to \mathcal{Z}$
拟紧，可知
$$
\mathcal{W} \times_\mathcal{Z} \mathcal{Z} \times_\mathcal{Y} \mathcal{X} =
\mathcal{W} \times_\mathcal{Y} \mathcal{X}
$$
拟紧。而映射
$\mathcal{W} \times_\mathcal{Y} \mathcal{X} \to
\mathcal{Y}' \times_\mathcal{Y} \mathcal{X}$
满，故得证。部分细节略。
\end{proof}

\begin{lemma}
\label{stacks-morphisms-lemma-check-separated-on-ui-cover}
设 $f : \mathcal{X} \to \mathcal{Y}$、$g : \mathcal{Y} \to \mathcal{Z}$
是可复合的代数叠态射，其复合为
$h = g \circ f : \mathcal{X} \to \mathcal{Z}$。若 $f$ 满、平坦、
局部有限表现且泛单射，并且 $h$ 分离，则 $g$ 分离。
\end{lemma}

\begin{proof}
考虑图
$$
\xymatrix{
\mathcal{X} \ar[r]_\Delta \ar[rd] &
\mathcal{X} \times_\mathcal{Y} \mathcal{X} \ar[r] \ar[d] &
\mathcal{X} \times_\mathcal{Z} \mathcal{X} \ar[d] \\
& \mathcal{Y} \ar[r] & \mathcal{Y} \times_\mathcal{Z} \mathcal{Y}
}
$$
该方块是笛卡儿的。需要证明底部水平箭头为真态射。我们已经知道它可由
代数空间表示且局部有限型（引理 \ref{stacks-morphisms-lemma-properties-diagonal}）。
由于右侧竖直箭头满、平坦且局部有限表现，只需证明右上方水平箭头为真态射
（引理 \ref{stacks-morphisms-lemma-check-property-after-fppf-base-change}）。由于 $h$ 分离，
顶部水平箭头的复合为真态射。

\medskip\noindent
由于 $f$ 泛单射，$\Delta$ 满（引理
\ref{stacks-morphisms-lemma-universally-injective}）。由于 $\Delta$ 与投影
$\mathcal{X} \times_\mathcal{Y} \mathcal{X} \to \mathcal{X}$ 的复合为恒等，
可知 $\Delta$ 泛闭。由《空间的态射》引理
\ref{spaces-morphisms-lemma-image-universally-closed-separated}
可知 $\mathcal{X} \times_\mathcal{Y} \mathcal{X} \to
\mathcal{X} \times_\mathcal{Z} \mathcal{X}$ 分离，因为
$\mathcal{X} \to \mathcal{X} \times_\mathcal{Z} \mathcal{X}$ 分离。
这里使用了如下事实：代数空间态射的性质之间的蕴含可转移为可由代数空间表示的
代数叠态射的相应性质之间的蕴含；这在《代数叠的性质》一节中讨论
\ref{stacks-properties-section-properties-morphisms}.
最后，使用同一原则，并由《空间的态射》引理
\ref{spaces-morphisms-lemma-image-proper-is-proper}
可知 $\mathcal{X} \times_\mathcal{Y} \mathcal{X} \to
\mathcal{X} \times_\mathcal{Z} \mathcal{X}$ 为真态射。
\end{proof}














\section{胚}
\label{stacks-morphisms-section-gerbes}

\noindent
一类重要的代数叠是形如 $[B/G]$ 的叠，其中 $B$ 是代数空间，而 $G$ 是
$B$ 上平坦且局部有限表现的群代数空间（在 $B$ 上平凡作用）；参见
《可表示性判据》引理 \ref{criteria-lemma-BG-algebraic}。事实证明，若一个
代数叠在 fppf 拓扑下局部地具有这种形式，则它是一个胚；参见引理
\ref{stacks-morphisms-lemma-local-structure-gerbe}。本节简要讨论这个概念及其相应的相对概念。

\begin{definition}
\label{stacks-morphisms-definition-gerbe}
设 $f : \mathcal{X} \to \mathcal{Y}$ 是代数叠的态射。若把源与靶
视为 $(\Sch/S)_{fppf}$ 上的群胚叠时，
$\mathcal{X}$ 是 $\mathcal{Y}$ 上的胚，则称 $\mathcal{X}$ 是
$\mathcal{Y}$ {\it 上的胚}；参见《叠》定义
\ref{stacks-definition-gerbe-over-stack-in-groupoids}。若存在态射
$\mathcal{X} \to X$，其中 $X$ 是代数空间，并且该态射使
$\mathcal{X}$ 成为 $X$ 上的胚，则称代数叠 $\mathcal{X}$ 是一个{\it 胚}。
\end{definition}

\noindent
$\mathcal{X}$ 是 $\mathcal{Y}$ 上的胚这一条件，纯粹用给定代数叠背后的
拓扑和范畴论来定义；但稍后将看到，该条件具有几何后果。例如，它蕴含
$\mathcal{X} \to \mathcal{Y}$ 满、平坦且局部有限表现；参见引理
\ref{stacks-morphisms-lemma-gerbe-fppf}。绝对概念更难解析，因为起初可能不清楚 $X$
是否被良好确定。事实上，它确实被良好确定。

\begin{lemma}
\label{stacks-morphisms-lemma-gerbe-over-iso-classes}
设 $\mathcal{X}$ 是代数叠。若 $\mathcal{X}$ 是胚，则预层
$$
(\Sch/S)_{fppf}^{opp} \to \textit{Sets}, \quad
U \mapsto \Ob(\mathcal{X}_U)/\!\!\cong
$$
的层化是代数空间，并且 $\mathcal{X}$ 是该空间上的胚。
\end{lemma}

\begin{proof}
（在这个证明中，第 \ref{stacks-morphisms-section-conventions} 节引入的术语滥用确实很有用。）
选取态射 $\pi : \mathcal{X} \to X$，其中 $X$ 是代数空间，并且该态射使
$\mathcal{X}$ 成为 $X$ 上的胚。只需证明 $X$ 是引理中所示预层
$\mathcal{F}$ 的层化。显然有映射 $c : \mathcal{F} \to X$。我们将使用
《叠》引理 \ref{stacks-lemma-when-gerbe} 的性质 (2)(a) 和 (2)(b)，
说明映射 $c^\# : \mathcal{F}^\# \to X$ 既满又单，因而是同构；参见
《站点》引理 \ref{sites-lemma-mono-epi-sheaves}。

满射性：设 $T$ 是概形，并设 $f : T \to X$。由性质 (2)(a)，存在 fppf 覆盖
$\{h_i : T_i \to T\}$ 和态射 $x_i : T_i \to \mathcal{X}$，使
$f \circ h_i$ 对应于 $\pi \circ x_i$。因此，$f|_{T_i}$ 属于 $c$ 的像。

单射性：设 $T$ 是概形，并设 $x, y : T \to \mathcal{X}$ 为态射，且
$c \circ x = c \circ y$。由 (2)(b)，可找到覆盖 $\{T_i \to T\}$ 以及
纤维范畴 $\mathcal{X}_{T_i}$ 中的态射
$x|_{T_i} \to y|_{T_i}$。因此，限制 $x|_{T_i}, y|_{T_i}$ 在
$\mathcal{F}(T_i)$ 中相等。这证明了 $x, y$ 按所需给出
$\mathcal{F}^\#$ 在 $T$ 上的同一截面。
\end{proof}

\begin{lemma}
\label{stacks-morphisms-lemma-base-change-gerbe}
设
$$
\xymatrix{
\mathcal{X}' \ar[r] \ar[d] & \mathcal{X} \ar[d] \\
\mathcal{Y}' \ar[r] & \mathcal{Y}
}
$$
为代数叠的纤维积。若 $\mathcal{X}$ 是 $\mathcal{Y}$ 上的胚，则
$\mathcal{X}'$ 是 $\mathcal{Y}'$ 上的胚。
\end{lemma}

\begin{proof}
立即由定义和《叠》引理 \ref{stacks-lemma-base-change-gerbe} 得出。
\end{proof}

\begin{lemma}
\label{stacks-morphisms-lemma-composition-gerbe}
设 $\mathcal{X} \to \mathcal{Y}$ 和 $\mathcal{Y} \to \mathcal{Z}$
是代数叠的态射。若 $\mathcal{X}$ 是 $\mathcal{Y}$ 上的胚，且
$\mathcal{Y}$ 是 $\mathcal{Z}$ 上的胚，则
$\mathcal{X}$ 是 $\mathcal{Z}$ 上的胚。
\end{lemma}

\begin{proof}
立即由《叠》引理 \ref{stacks-lemma-composition-gerbe} 得出。
\end{proof}

\begin{lemma}
\label{stacks-morphisms-lemma-gerbe-descent}
设
$$
\xymatrix{
\mathcal{X}' \ar[r] \ar[d] & \mathcal{X} \ar[d] \\
\mathcal{Y}' \ar[r] & \mathcal{Y}
}
$$
为代数叠的纤维积。若 $\mathcal{Y}' \to \mathcal{Y}$ 满、平坦且局部有限表现，
并且 $\mathcal{X}'$ 是 $\mathcal{Y}'$ 上的胚，则
$\mathcal{X}$ 是 $\mathcal{Y}$ 上的胚。
\end{lemma}

\begin{proof}
立即由引理 \ref{stacks-morphisms-lemma-surjective-flat-locally-finite-presentation} 和
《叠》引理 \ref{stacks-lemma-gerbe-descent} 得出。
\end{proof}

\begin{lemma}
\label{stacks-morphisms-lemma-gerbe-with-section}
设 $\pi : \mathcal{X} \to U$ 是从代数叠到代数空间的态射，并设
$x : U \to \mathcal{X}$ 是 $\pi$ 的截面。令
$G = \mathit{Isom}_\mathcal{X}(x, x)$；参见定义
\ref{stacks-morphisms-definition-isom}。若 $\mathcal{X}$ 是 $U$ 上的胚，则
\begin{enumerate}
\item 有群胚叠的典范等价
$$
x_{can} : [U/G] \longrightarrow \mathcal{X}.
$$
其中 $[U/G]$ 是 $G$ 在 $U$ 上平凡作用所对应的商叠；
\item $G \to U$ 平坦且局部有限表现；
\item $U \to \mathcal{X}$ 满、平坦且局部有限表现。
\end{enumerate}
\end{lemma}

\begin{proof}
令 $R = U \times_{x, \mathcal{X}, x} U$。态射 $R \to U \times U$
经过对角态射 $\Delta_U : U \to U \times U$ 分解，因为它经过
$U \times_U U = U$ 分解。因此 $R = G$，因为
\begin{align*}
G & = \mathit{Isom}_\mathcal{X}(x, x) \\
& = U \times_{x, \mathcal{X}} \mathcal{I}_\mathcal{X} \\
& = U \times_{x, \mathcal{X}}
(\mathcal{X}
\times_{\Delta, \mathcal{X} \times_S \mathcal{X}, \Delta}
\mathcal{X}) \\
& = (U \times_{x, \mathcal{X}, x} U) \times_{U \times U, \Delta_U} U \\
& = R \times_{U \times U, \Delta_U} U \\
& = R
\end{align*}
第四个等号使用《范畴》引理 \ref{categories-lemma-diagonal-2}。
令 $t, s : R \to U$ 为投影。《代数叠》引理
\ref{algebraic-lemma-map-space-into-stack} 在 $R$ 上构造的复合法则
$c : R \times_{s, U, t} R \to R$ 与 $G$ 上的群法则一致（证明略）。
因此，《代数叠》引理 \ref{algebraic-lemma-map-space-into-stack}
说明我们得到典范的全忠实 $1$-态射
$$
x_{can} : [U/G] \longrightarrow \mathcal{X}
$$
，它是 $(\Sch/S)_{fppf}$ 上的群胚叠之间的态射。为说明它是等价，只需证明
它本质满。为此，只需证明 $\mathcal{X}$ 在任一概形 $T$ 上的对象，在
fppf 局部下经某个态射 $T \to U$ 来自 $x$；参见《叠》引理
\ref{stacks-lemma-characterize-essentially-surjective-when-ff}。
而这由 $\pi$ 使 $\mathcal{X}$ 成为 $U$ 上的胚这一条件得出；参见
《叠》引理 \ref{stacks-lemma-when-gerbe} 的性质 (2)(a)。

\medskip\noindent
由《可表示性判据》引理 \ref{criteria-lemma-BG-algebraic} 可知
$G \to U$ 平坦且局部有限表现。最后，由《可表示性判据》引理
\ref{criteria-lemma-flat-quotient-flat-presentation}
可知 $U \to \mathcal{X}$ 满、平坦且局部有限表现。
\end{proof}

\begin{lemma}
\label{stacks-morphisms-lemma-local-structure-gerbe}
设 $\pi : \mathcal{X} \to \mathcal{Y}$ 是代数叠的态射。下列条件等价：
\begin{enumerate}
\item $\mathcal{X}$ 是 $\mathcal{Y}$ 上的胚；
\item 存在代数空间 $U$、在 $U$ 上平坦且局部有限表现的群代数空间 $G$，
以及满、平坦且局部有限表现的态射 $U \to \mathcal{Y}$，使在 $U$ 上有
$\mathcal{X} \times_\mathcal{Y} U \cong [U/G]$。
\end{enumerate}
\end{lemma}

\begin{proof}
假设 (2)。由引理 \ref{stacks-morphisms-lemma-gerbe-descent}，为证明 (1)，只需说明
$[U/G]$ 是 $U$ 上的胚。这立即由《空间中的群胚》引理
\ref{spaces-groupoids-lemma-group-quotient-gerbe} 得出。

\medskip\noindent
假设 (1)。$\pi$ 的任意基变换都是胚；参见引理
\ref{stacks-morphisms-lemma-base-change-gerbe}。第一步，选取概形 $V$ 以及满光滑态射
$V \to \mathcal{Y}$。于是可假设 $\pi : \mathcal{X} \to V$ 是概形上的胚。
这意味着存在 fppf 覆盖 $\{V_i \to V\}$，使纤维范畴
$\mathcal{X}_{V_i}$ 非空；参见《叠》引理
\ref{stacks-lemma-when-gerbe} (2)(a)。注意，
$U = \coprod V_i \to V$ 满、平坦且局部有限表现。因此，可用 $U$ 替换 $V$，
并假设 $\pi : \mathcal{X} \to U$ 是概形 $U$ 上的胚，且存在
$\mathcal{X}$ 在 $U$ 上的对象 $x$。由引理
\ref{stacks-morphisms-lemma-gerbe-with-section} 可知，在 $U$ 上有
$\mathcal{X} = [U/G]$，其中 $G$ 是 $U$ 上平坦且局部有限表现的群代数空间。
\end{proof}

\begin{lemma}
\label{stacks-morphisms-lemma-gerbe-fppf}
设 $\pi : \mathcal{X} \to \mathcal{Y}$ 是代数叠的态射。若
$\mathcal{X}$ 是 $\mathcal{Y}$ 上的胚，则 $\pi$ 满、平坦且局部有限表现。
\end{lemma}

\begin{proof}
由《代数叠的性质》引理
\ref{stacks-properties-lemma-descent-surjective}
以及引理 \ref{stacks-morphisms-lemma-descent-flat} 和
\ref{stacks-morphisms-lemma-descent-finite-presentation}
可知，只需在用满、平坦且局部有限表现的态射
$\mathcal{Y}' \to \mathcal{Y}$ 对 $\pi$ 作基变换后证明引理。由引理
\ref{stacks-morphisms-lemma-local-structure-gerbe}，可假设 $\mathcal{Y} = U$ 是代数空间，
并且在 $U$ 上有 $\mathcal{X} = [U/G]$。此时
$U \to [U/G]$ 满、平坦且局部有限表现；参见引理
\ref{stacks-morphisms-lemma-gerbe-with-section}。于是由《代数叠的性质》引理
\ref{stacks-properties-lemma-surjective-permanence} 以及引理
\ref{stacks-morphisms-lemma-flat-permanence} 和
\ref{stacks-morphisms-lemma-flat-finite-presentation-permanence}
可知 $\pi$ 满、平坦且局部有限表现。
\end{proof}

\begin{proposition}
\label{stacks-morphisms-proposition-when-gerbe}
设 $\mathcal{X}$ 是代数叠。下列条件等价：
\begin{enumerate}
\item $\mathcal{X}$ 是胚；
\item $\mathcal{I}_\mathcal{X} \to \mathcal{X}$ 平坦且局部有限表现。
\end{enumerate}
\end{proposition}

\begin{proof}
假设 (1)。选取到代数空间 $X$ 的态射 $\mathcal{X} \to X$，使
$\mathcal{X}$ 成为 $X$ 上的胚。设 $X' \to X$ 是满、平坦且局部有限表现的态射，
并令 $\mathcal{X}' = X' \times_X \mathcal{X}$。注意，由引理
\ref{stacks-morphisms-lemma-base-change-gerbe}，$\mathcal{X}'$ 是 $X'$ 上的胚。下图中的两个
方块都是纤维积方块：
$$
\xymatrix{
\mathcal{I}_{\mathcal{X}'} \ar[r] \ar[d] &
\mathcal{X}' \ar[r] \ar[d] & X' \ar[d] \\
\mathcal{I}_\mathcal{X} \ar[r] &
\mathcal{X} \ar[r] & X
}
$$
；参见引理 \ref{stacks-morphisms-lemma-cartesian-square-inertia}。因此，由引理
\ref{stacks-morphisms-lemma-descent-flat} 和
\ref{stacks-morphisms-lemma-descent-finite-presentation}，为证明
$\mathcal{I}_\mathcal{X} \to \mathcal{X}$ 平坦且局部有限表现，
只需在这样的基变换后证明。于是可应用引理
\ref{stacks-morphisms-lemma-local-structure-gerbe}，假设 $\mathcal{X} = [U/G]$。由引理
\ref{stacks-morphisms-lemma-gerbe-with-section} 可知，$G$ 在 $U$ 上平坦且局部有限表现，
而 $x : U \to [U/G]$ 满、平坦且局部有限表现。此外，
$\mathcal{I}_\mathcal{X}$ 沿 $x$ 的拉回为 $G$，再次由下降，即由引理
\ref{stacks-morphisms-lemma-descent-flat} 和
\ref{stacks-morphisms-lemma-descent-finite-presentation}，可知 (2) 成立。

\medskip\noindent
反之，假设 (2)。选取光滑表现 $\mathcal{X} = [U/R]$；参见《代数叠》一节
\ref{algebraic-section-stack-to-presentation}。记 $G \to U$ 为群胚
$(U, R, s, t, c, e, i)$ 的稳定子群代数空间；参见《空间中的群胚》定义
\ref{spaces-groupoids-definition-stabilizer-groupoid}.
由引理 \ref{stacks-morphisms-lemma-presentation-inertia} 可知，$G \to U$ 作为
$\mathcal{I}_\mathcal{X} \to \mathcal{X}$ 的基变换平坦且局部有限表现；
参见引理 \ref{stacks-morphisms-lemma-base-change-flat} 和
\ref{stacks-morphisms-lemma-base-change-finite-presentation}.
考虑如下 $G$ 在 $R$ 上的作用：
$$
a : G \times_{U, t} R \to R, \quad (g, r) \mapsto c(g, r)
$$
。由于 $R$ 是群胚，该作用在任意概形 $T$ 的 $T$-值点上自由。因此，
由《自举》引理 \ref{bootstrap-lemma-quotient-free-action}，$R' = R/G$
是代数空间，而商态射 $\pi : R \to R'$ 满、平坦且局部有限表现。
投影 $s, t : R \to U$ 是 $G$-不变的，故得到态射
$s' , t' : R' \to U$，使 $s = s' \circ \pi$ 且
$t = t' \circ \pi$。由于 $s, t : R \to U$ 平坦且局部有限表现，
由《空间的态射》引理
\ref{spaces-morphisms-lemma-flat-permanence} 和《空间上的下降》引理
\ref{spaces-descent-lemma-locally-finite-presentation-fppf-local-source}
可知 $s', t'$ 平坦且局部有限表现。考虑态射
$$
j' = (t', s') : R' \longrightarrow U \times U.
$$
我们断言它是单态射。事实上，设 $T$ 是概形，并设
$a, b : T \to R'$ 是在 $U \times U$ 中有相同像的态射。根据商
$R' = R/G$ 的定义，存在 fppf 覆盖 $\{h_j : T_j \to T\}$，使对某些态射
$a_j, b_j : T_j \to R$ 有 $a \circ h_j = \pi \circ a_j$ 且
$b \circ h_j = \pi \circ b_j$。由于 $a_j, b_j$ 在 $U \times U$ 中有相同像，
$g_j = c(a_j, i(b_j))$ 是 $G$ 的一个 $T_j$-值点，并且
$c(g_j, b_j) = a_j$。换言之，$a_j$ 和 $b_j$ 在 $R'$ 中有相同的像，
故断言得证。由于 $j : R \to U \times U$ 是预等价关系（参见
《空间中的群胚》引理
\ref{spaces-groupoids-lemma-groupoid-pre-equivalence}），且
$R \to R'$ 作为层的映射为满，可知
$j' : R' \to U \times U$ 是等价关系。因此，《自举》定理
\ref{bootstrap-theorem-final-bootstrap} 说明 $X = U/R'$ 是代数空间。
最后，我们断言态射
$$
\mathcal{X} = [U/R] \longrightarrow X = U/R'
$$
使 $\mathcal{X}$ 成为 $X$ 上的胚。这由《空间中的群胚》引理
\ref{spaces-groupoids-lemma-when-gerbe} 得出，因为 $R \to R'$ 满、平坦且
局部有限表现（必要时使用《自举》引理
\ref{bootstrap-lemma-surjective-flat-locally-finite-presentation}
以看出这蕴含所需假设）。
\end{proof}

\begin{lemma}
\label{stacks-morphisms-lemma-gerbe-isom-fppf}
设 $f : \mathcal{X} \to \mathcal{Y}$ 是使 $\mathcal{X}$ 成为
$\mathcal{Y}$ 上之胚的代数叠态射。则
\begin{enumerate}
\item $\mathcal{I}_{\mathcal{X}/\mathcal{Y}} \to \mathcal{X}$
平坦且局部有限表现；
\item $\mathcal{X} \to \mathcal{X} \times_\mathcal{Y} \mathcal{X}$
满、平坦且局部有限表现；
\item 给定代数空间 $T_i$（$i = 1, 2$）和态射
$x_i : T_i \to \mathcal{X}$，令 $y_i = f \circ x_i$，则态射
$$
T_1 \times_{x_1, \mathcal{X}, x_2} T_2 \longrightarrow
T_1 \times_{y_1, \mathcal{Y}, y_2} T_2
$$
满、平坦且局部有限表现；
\item 给定代数空间 $T$ 和态射 $x_i : T \to \mathcal{X}$（$i = 1, 2$），
令 $y_i = f \circ x_i$，则态射
$$
\mathit{Isom}_\mathcal{X}(x_1, x_2) \longrightarrow
\mathit{Isom}_\mathcal{Y}(y_1, y_2)
$$
满、平坦且局部有限表现。
\end{enumerate}
\end{lemma}

\begin{proof}
(1) 的证明。选取概形 $Y$ 以及满光滑态射 $Y \to \mathcal{Y}$。令
$\mathcal{X}' = \mathcal{X} \times_\mathcal{Y} Y$。由引理
\ref{stacks-morphisms-lemma-cartesian-square-inertia}，得到笛卡儿方块
$$
\xymatrix{
\mathcal{I}_{\mathcal{X}'} \ar[r] \ar[d] &
\mathcal{X}' \ar[r] \ar[d] & Y \ar[d] \\
\mathcal{I}_{\mathcal{X}/\mathcal{Y}} \ar[r] &
\mathcal{X} \ar[r] & \mathcal{Y}
}
$$
由引理 \ref{stacks-morphisms-lemma-descent-flat} 和
\ref{stacks-morphisms-lemma-descent-finite-presentation}
可知，只需证明 $\mathcal{I}_{\mathcal{X}'} \to \mathcal{X}'$
平坦且局部有限表现。这由命题 \ref{stacks-morphisms-proposition-when-gerbe} 得出
（因为由引理 \ref{stacks-morphisms-lemma-base-change-gerbe}，$\mathcal{X}'$ 是 $Y$ 上的胚）。

\medskip\noindent
(2) 的证明。沿用上述记号，注意可假设
$\mathcal{X}' = [Y/G]$，其中 $G$ 是 $Y$ 上平坦且局部有限表现的
群代数空间；参见引理 \ref{stacks-morphisms-lemma-local-structure-gerbe}。态射
$\Delta : \mathcal{X} \to \mathcal{X} \times_\mathcal{Y} \mathcal{X}$
在 $\mathcal{Y}$ 上沿态射 $Y \to \mathcal{Y}$ 的基变换是态射
$\Delta' : \mathcal{X}' \to \mathcal{X}' \times_Y \mathcal{X}'$.
因此，只需证明 $\Delta'$ 满、平坦且局部有限表现（参见引理
\ref{stacks-morphisms-lemma-descent-flat} 和
\ref{stacks-morphisms-lemma-descent-finite-presentation}).
换言之，需要证明
$$
[Y/G] \longrightarrow [Y/G \times_Y G]
$$
满、平坦且局部有限表现。这是因为沿满、平坦且局部有限表现的态射
$Y \to [Y/G \times_Y G]$ 作基变换后得到态射 $G \to Y$。

\medskip\noindent
(3) 的证明。注意图
$$
\xymatrix{
T_1 \times_{x_1, \mathcal{X}, x_2} T_2 \ar[d] \ar[r] &
T_1 \times_{y_1, \mathcal{Y}, y_2} T_2 \ar[d] \\
\mathcal{X} \ar[r] & \mathcal{X} \times_\mathcal{Y} \mathcal{X}
}
$$
是笛卡儿的。因此 (3) 由 (2) 得出。

\medskip\noindent
(4) 的证明。由于
$$
\mathit{Isom}_\mathcal{X}(x_1, x_2) =
(T \times_{x_1, \mathcal{X}, x_2} T) \times_{T \times T, \Delta_T} T
$$
，(4) 中的态射是 (3) 中态射的基变换，故结论成立。
\end{proof}

\begin{proposition}
\label{stacks-morphisms-proposition-when-gerbe-over}
设 $f : \mathcal{X} \to \mathcal{Y}$ 是代数叠的态射。下列条件等价：
\begin{enumerate}
\item $\mathcal{X}$ 是 $\mathcal{Y}$ 上的胚；
\item $f : \mathcal{X} \to \mathcal{Y}$ 与
$\Delta : \mathcal{X} \to \mathcal{X} \times_\mathcal{Y} \mathcal{X}$
均满、平坦且局部有限表现。
\end{enumerate}
\end{proposition}

\begin{proof}
(1) $\Rightarrow$ (2) 由引理 \ref{stacks-morphisms-lemma-gerbe-fppf} 和
\ref{stacks-morphisms-lemma-gerbe-isom-fppf} 得出。

\medskip\noindent
假设 (2)。只需对 $f$ 沿满、平坦且局部有限表现的态射
$\mathcal{Y}' \to \mathcal{Y}$ 所作的基变换证明 (1)；参见引理
\ref{stacks-morphisms-lemma-gerbe-descent}（注意，$f$ 的对角态射的基变换就是基变换的
对角态射）。因此可假设 $\mathcal{Y}$ 是概形 $Y$。在此情形，
$\mathcal{I}_\mathcal{X} \to \mathcal{X}$ 是 $\Delta$ 的基变换，故由命题
\ref{stacks-morphisms-proposition-when-gerbe} 可知 $\mathcal{X}$ 是胚。还需证明
$\mathcal{X}$ 是 $Y$ 上的胚。令 $\mathcal{X} \to X$ 为引理
\ref{stacks-morphisms-lemma-gerbe-over-iso-classes} 的态射，它使 $\mathcal{X}$ 成为代数空间
$X$ 上的胚，而后者分类 $\mathcal{X}$ 的对象的同构类。显然，
$f : \mathcal{X} \to Y$ 分解为 $\mathcal{X} \to X \to Y$。由于 $f$
满、平坦且局部有限表现，可知 $X \to Y$ 作为 fppf 层的映射为满
（例如使用引理 \ref{stacks-morphisms-lemma-surjective-flat-locally-finite-presentation}）。
另一方面，$X \to Y$ 也单：对任意概形 $T$ 以及 $X$ 的任意两个
$T$-值点 $x_1, x_2$，若它们映到 $Y$ 的同一点，则首先可在 $T$ 上
fppf 局部地把 $x_1, x_2$ 提升为 $\mathcal{X}$ 在 $T$ 上的对象
$\xi_1, \xi_2$；其次，由
$\Delta : \mathcal{X} \to \mathcal{X} \times_Y \mathcal{X}$
满、平坦且局部有限表现这一假设，可推出 $\xi_1$ 和 $\xi_2$ 在 fppf
局部同构。于是按 $X$ 的构造，$x_1 = x_2$。因此 $X = Y$，证明完成。
\end{proof}

\noindent
至此，我们已经发展出足够的机制，可以证明剩余胚在存在时确实是胚。

\begin{lemma}
\label{stacks-morphisms-lemma-noetherian-singleton-stack-gerbe}
设 $\mathcal{Z}$ 是约化、局部诺特的代数叠，并且 $|\mathcal{Z}|$ 是单点集。
则 $\mathcal{Z}$ 是某个约化、局部诺特的代数空间 $Z$ 上的胚，且
$|Z|$ 是单点集。
\end{lemma}

\begin{proof}
由《代数叠的性质》引理
\ref{stacks-properties-lemma-unique-point-better}，存在满、平坦且局部有限表现的
态射 $\Spec(k) \to \mathcal{Z}$，其中 $k$ 为域。于是
$\mathcal{I}_Z \times_\mathcal{Z} \Spec(k) \to \Spec(k)$ 可由代数空间表示，
并且局部有限型（作为 $\mathcal{I}_\mathcal{Z} \to \mathcal{Z}$ 的基变换；
参见引理 \ref{stacks-morphisms-lemma-inertia} 和 \ref{stacks-morphisms-lemma-base-change-finite-type}）。
因此，它局部有限表现；参见《空间的态射》引理
\ref{spaces-morphisms-lemma-noetherian-finite-type-finite-presentation}.
当然，因为 $k$ 是域，它也平坦。因此可以应用引理
\ref{stacks-morphisms-lemma-descent-flat} 和
\ref{stacks-morphisms-lemma-descent-finite-presentation}
，可知 $\mathcal{I}_\mathcal{Z} \to \mathcal{Z}$ 平坦且局部有限表现。
由命题 \ref{stacks-morphisms-proposition-when-gerbe}，$\mathcal{Z}$ 是胚。设
$\pi : \mathcal{Z} \to Z$ 为到代数空间的态射，使 $\mathcal{Z}$ 成为
$Z$ 上的胚。由引理 \ref{stacks-morphisms-lemma-gerbe-fppf}，$\pi$ 满、平坦且局部有限表现。
因此，作为复合，$\Spec(k) \to Z$ 满、平坦且局部有限表现；参见
《代数叠的性质》引理
\ref{stacks-properties-lemma-composition-surjective} 以及引理
\ref{stacks-morphisms-lemma-composition-flat} 和
\ref{stacks-morphisms-lemma-composition-finite-presentation}.
于是由《代数叠的性质》引理
\ref{stacks-properties-lemma-unique-point-better} 可知，$|Z|$ 是单点集，
且 $Z$ 局部诺特并约化。
\end{proof}

\begin{lemma}
\label{stacks-morphisms-lemma-gerbe-bijection-points}
设 $f : \mathcal{X} \to \mathcal{Y}$ 是代数叠的态射。若
$\mathcal{X}$ 是 $\mathcal{Y}$ 上的胚，则 $f$ 是泛同胚。
\end{lemma}

\begin{proof}
由引理 \ref{stacks-morphisms-lemma-base-change-gerbe}，关于 $f$ 的假设在基变换下保持。
因此，只需说明映射 $|\mathcal{X}| \to |\mathcal{Y}|$ 是拓扑空间的同胚。
设 $k$ 为域，并设 $y$ 是 $\mathcal{Y}$ 在 $\Spec(k)$ 上的对象。由
《叠》引理 \ref{stacks-lemma-when-gerbe} 的性质 (2)(a)，存在 fppf 覆盖
$\{T_i \to \Spec(k)\}$ 以及 $\mathcal{X}$ 在 $T_i$ 上的对象 $x_i$，
使 $f(x_i) \cong y|_{T_i}$。选取满足 $T_i \not = \emptyset$ 的 $i$。
对某个域 $K$，选取态射 $\Spec(K) \to T_i$。则 $k \subset K$，而
$x_i|_K$ 是位于 $y|_K$ 上方的 $\mathcal{X}$ 的对象。因此可知
$|\mathcal{Y}| \to |\mathcal{X}|$ 满。映射
$|\mathcal{Y}| \to |\mathcal{X}|$ 也单。事实上，若 $x, x'$ 是
$\mathcal{X}$ 在 $\Spec(k)$ 上的对象，并且它们的像 $f(x), f(x')$
在 $\mathcal{Y}$ 中经扩张后同构，则《叠》引理
\ref{stacks-lemma-when-gerbe} 的性质 (2)(b) 保证存在 $k$ 的一个扩张，
使 $x$ 和 $x'$ 同构（细节略）。因此
$|\mathcal{X}| \to |\mathcal{Y}|$ 连续且双射，只需再证明它是开的。
这由引理 \ref{stacks-morphisms-lemma-gerbe-fppf} 和 \ref{stacks-morphisms-lemma-fppf-open} 得出。
\end{proof}

\begin{lemma}
\label{stacks-morphisms-lemma-gerbe-diagonal-quasi-compact}
设 $f : \mathcal{X} \to \mathcal{Y}$ 是代数叠的态射，并且
$\mathcal{X}$ 是 $\mathcal{Y}$ 上的胚。若 $\Delta_\mathcal{X}$ 拟紧，
则 $\Delta_\mathcal{Y}$ 也拟紧。
\end{lemma}

\begin{proof}
考虑图
$$
\xymatrix{
\mathcal{X} \ar[r] &
\mathcal{X} \times_\mathcal{Y} \mathcal{X} \ar[r] \ar[d] &
\mathcal{X} \times \mathcal{X} \ar[d] \\
&
\mathcal{Y} \ar[r] &
\mathcal{Y} \times \mathcal{Y}
}
$$
由命题 \ref{stacks-morphisms-proposition-when-gerbe-over} 可知左上方箭头满。由于顶部水平箭头的
复合拟紧，由引理 \ref{stacks-morphisms-lemma-surjection-from-quasi-compact} 可知右上方箭头拟紧。
该方块是笛卡儿的，而右侧竖直箭头满、平坦且局部有限表现。因此由引理
\ref{stacks-morphisms-lemma-descent-quasi-compact} 得出结论。
\end{proof}

\noindent
下述引理说明，对任何本身为胚的代数叠，其所有点的剩余胚都存在。

\begin{lemma}
\label{stacks-morphisms-lemma-gerbe-residual-gerbe-exists}
设 $\mathcal{X}$ 是代数叠。若 $\mathcal{X}$ 是胚，则对每个
$x \in |\mathcal{X}|$，$\mathcal{X}$ 在 $x$ 处的剩余胚都存在。
\end{lemma}

\begin{proof}
设 $\pi : \mathcal{X} \to X$ 是从 $\mathcal{X}$ 到代数空间 $X$
的态射，并且它使 $\mathcal{X}$ 成为 $X$ 上的胚。设
$Z_x \to X$ 是 $X$ 在 $x$ 处的剩余空间，见《适宜代数空间》定义
\ref{decent-spaces-definition-residual-space}。令
$\mathcal{Z} = \mathcal{X} \times_X Z_x$。由引理
\ref{stacks-morphisms-lemma-base-change-gerbe}，代数叠 $\mathcal{Z}$ 是 $Z_x$ 上的胚。
因此 $|\mathcal{Z}| = |Z_x|$（引理
\ref{stacks-morphisms-lemma-gerbe-bijection-points}）是单点集。由于
$\mathcal{Z} \to Z_x$ 是 $\pi$ 的基变换，故其局部有限表现（见引理
\ref{stacks-morphisms-lemma-gerbe-fppf} 和 \ref{stacks-morphisms-lemma-base-change-finite-presentation}）。
于是 $\mathcal{Z}$ 局部诺特，见引理
\ref{stacks-morphisms-lemma-locally-finite-type-locally-noetherian}。
所以 $\mathcal{X}$ 在 $x$ 处的剩余胚 $\mathcal{Z}_x$ 存在，并且等于
$\mathcal{Z}_x = \mathcal{Z}_{red}$，即代数叠 $\mathcal{Z}$ 的约化。
事实上，上文已经说明 $|\mathcal{Z}_{red}|$ 是映到
$x \in |\mathcal{X}|$ 的单点集；它按构造是约化的，并且局部诺特
（因为局部诺特代数叠的约化仍局部诺特，细节略）。
\end{proof}









\section{按胚分层}
\label{stacks-morphisms-section-stratify}

\noindent
本节的目标是证明，许多代数叠 $\mathcal{X}$ 都有由局部闭子叠
$\mathcal{X}_i \subset \mathcal{X}$ 构成的“分层”，其中每个
$\mathcal{X}_i$ 都是胚。这说明在某种意义下，胚是构成任意代数叠的基本构件。
注意，在本节中，我们所谓的分层仅指
$$
|\mathcal{X}| = \bigcup\nolimits_i |\mathcal{X}_i|
$$
是与 $\mathcal{X}$ 关联的拓扑空间的分层，并无其他含义。因此，为研究这种分层，
可以无妨地用 $\mathcal{X}$ 的约化替换它（见《叠的性质》一节
\ref{stacks-properties-section-reduced}）。

\medskip\noindent
下述命题说明，$\mathcal{X}$ 的约化中（几乎总是）存在一个稠密开子叠。

\begin{proposition}
\label{stacks-morphisms-proposition-open-stratum}
设 $\mathcal{X}$ 是约化代数叠，且
$\mathcal{I}_\mathcal{X} \to \mathcal{X}$ 拟紧。则存在一个稠密开子叠
$\mathcal{U} \subset \mathcal{X}$，它是胚。
\end{proposition}

\begin{proof}
由命题 \ref{stacks-morphisms-proposition-when-gerbe}，只需找到一个稠密开子叠
$\mathcal{U}$，使得 $\mathcal{I}_\mathcal{U} \to \mathcal{U}$
平坦且局部有限表现。注意
$\mathcal{I}_\mathcal{U} =
\mathcal{I}_\mathcal{X} \times_\mathcal{X} \mathcal{U}$，见引理
\ref{stacks-morphisms-lemma-cartesian-square-inertia}。

\medskip\noindent
选取一个表现 $\mathcal{X} = [U/R]$。设 $G \to U$ 是群胚 $R$ 的稳定子群
代数空间。由引理 \ref{stacks-morphisms-lemma-presentation-inertia}，$G \to U$ 是
$\mathcal{I}_\mathcal{X} \to \mathcal{X}$ 的基变换，因而按假设拟紧，并由引理
\ref{stacks-morphisms-lemma-inertia} 局部有限型。设 $W \subset U$ 是最大的开子概形
（可能为空），使得限制 $G_W \to W$ 平坦且局部有限表现
（我们略去 $W$ 存在性的证明；提示：使用这些性质是局部的这一事实）。
由《代数空间的态射》命题
\ref{spaces-morphisms-proposition-generic-flatness-reduced}，
$W \subset U$ 稠密。注意，由《空间中群胚的进一步内容》引理
\ref{spaces-more-groupoids-lemma-property-G-invariant}，
$W \subset U$ 是 $R$-不变的。因此，由《叠的性质》引理
\ref{stacks-properties-lemma-substacks-presentation}，$W$ 对应于一个开子叠
$\mathcal{U} \subset \mathcal{X}$。由于 $|U| \to |\mathcal{X}|$ 是开映射，
且 $|W| \subset |U|$ 稠密，可知 $\mathcal{U}$ 在 $\mathcal{X}$ 中稠密。
最后，态射 $\mathcal{I}_\mathcal{U} \to \mathcal{U}$ 平坦且局部有限表现，
因为沿满光滑态射 $W \to \mathcal{U}$ 的基变换是态射 $G_W \to W$，
而后者按构造平坦且局部有限表现。见引理 \ref{stacks-morphisms-lemma-descent-flat} 和
\ref{stacks-morphisms-lemma-descent-finite-presentation}。
\end{proof}

\noindent
上述命题立即蕴含：在惯性拟紧的代数叠上，任意点都有剩余胚；我们将在引理
\ref{stacks-morphisms-lemma-every-point-residual-gerbe} 中证明这一点。另一方面，由胚构成的
有限分层并非总是存在。下面给出一个例子。

\begin{example}
\label{stacks-morphisms-example-infinite-stratification}
设 $k$ 是域。取 $U = \Spec(k[x_0, x_1, x_2, \ldots])$，并令
$\mathbf{G}_m$ 按下式作用：
$t(x_0, x_1, x_2, \ldots) = (tx_0, t^p x_1, t^{p^2} x_2, \ldots)$
，其中 $p$ 是素数。令 $\mathcal{X} = [U/\mathbf{G}_m]$。这是代数叠。
$\mathcal{X}$ 有如下分层：
\begin{enumerate}
\item $\mathcal{X}_0$ 是 $x_0$ 非零的部分；
\item $\mathcal{X}_1$ 是 $x_0$ 为零而 $x_1$ 非零的部分；
\item $\mathcal{X}_2$ 是 $x_0, x_1$ 为零而 $x_2$ 非零的部分；
\item 依此类推；并且
\item $\mathcal{X}_{\infty}$ 是所有 $x_i$ 都为零的部分。
\end{enumerate}
每一层都是某个概形上的胚：对 $\mathcal{X}_i$，其群为 $\mu_{p^i}$；
对 $\mathcal{X}_{\infty}$，其群为 $\mathbf{G}_m$。这些层是约化局部闭子叠。
不存在具有相同性质而更粗的分层。
\end{example}

\noindent
尽管如此，利用超限归纳与命题 \ref{stacks-morphisms-proposition-open-stratum}，我们可以找到
可能无限的胚分层……！

\begin{lemma}
\label{stacks-morphisms-lemma-every-point-in-a-stratum}
设 $\mathcal{X}$ 是代数叠，并且
$\mathcal{I}_\mathcal{X} \to \mathcal{X}$ 拟紧。则存在一个良序指标集 $I$，
并且对每个 $i \in I$，存在一个约化局部闭子叠
$\mathcal{U}_i \subset \mathcal{X}$，使得
\begin{enumerate}
\item 每个 $\mathcal{U}_i$ 都是胚；
\item 有 $|\mathcal{X}| = \bigcup_{i \in I} |\mathcal{U}_i|$；
\item $T_i = |\mathcal{X}| \setminus \bigcup_{i' < i} |\mathcal{U}_{i'}|$
对所有 $i \in I$ 都是 $|\mathcal{X}|$ 中的闭集；并且
\item $|\mathcal{U}_i|$ 是 $T_i$ 中的开集。
\end{enumerate}
此外，还可以作如下安排：或者 (a) $|\mathcal{U}_i| \subset T_i$ 稠密，
或者 (b) $\mathcal{U}_i$ 拟紧。在情形 (a) 中，若把 $\mathcal{U}_i$
取得尽可能大（细节见证明），则该分层是典范的。
\end{lemma}

\begin{proof}
设 $T \subset |\mathcal{X}|$ 是非空闭子集。对每个这样的 $T$，我们要找出
（分别地，选取）一个约化局部闭子叠
$\mathcal{U}(T) \subset \mathcal{X}$，使得
$|\mathcal{U}(T)| \subset T$ 开且稠密（分别地，非空且拟紧）。事实上，
由《叠的性质》引理 \ref{stacks-properties-lemma-reduced-closed-substack}，
存在唯一的约化闭子叠 $\mathcal{X}' \subset \mathcal{X}$，使得
$T = |\mathcal{X}'|$。注意，由引理
\ref{stacks-morphisms-lemma-monomorphism-cartesian-square-inertia}，
$\mathcal{I}_{\mathcal{X}'} =
\mathcal{I}_\mathcal{X} \times_\mathcal{X} \mathcal{X}'$。
因此，作为一个基变换，$\mathcal{I}_{\mathcal{X}'} \to \mathcal{X}'$ 拟紧；
见引理 \ref{stacks-morphisms-lemma-base-change-quasi-compact}。所以命题
\ref{stacks-morphisms-proposition-open-stratum} 蕴含：存在一个稠密且极大的开子叠
$\mathcal{U} \subset \mathcal{X}'$（见该命题的证明），它是胚。在情形 (a)
中令 $\mathcal{U}(T) = \mathcal{U}$（这是典范的）；在情形 (b) 中，
只需选取一个非空拟紧开子叠
$\mathcal{U}(T) \subset \mathcal{U}$，见《叠的性质》引理
\ref{stacks-properties-lemma-space-locally-quasi-compact}
（借助选择公理，可以对所有 $T$ 同时这样做）。

\medskip\noindent
利用超限递归，我们对每个序数 $\alpha$ 构造一个闭子集
$T_\alpha \subset |\mathcal{X}|$。对 $\alpha = 0$，令
$T_0 = |\mathcal{X}|$。给定 $T_\alpha$ 后，令
$$
T_{\alpha + 1} = T_\alpha \setminus |\mathcal{U}(T_\alpha)|.
$$
若 $\beta$ 是极限序数，则令
$$
T_\beta = \bigcap\nolimits_{\alpha < \beta} T_\alpha.
$$
我们断言，对所有充分大的 $\alpha$，都有 $T_\alpha = \emptyset$。事实上，
假设对所有 $\alpha$ 都有 $T_\alpha \not = \emptyset$。于是可以得到从
全体序数这个类到 $|\mathcal{X}|$ 的子集所成集合的单射，矛盾。

\medskip\noindent
该断言蕴含本引理。事实上，令
$$
I = \{\alpha \mid \mathcal{U}_\alpha \not = \emptyset \}.
$$
由该断言，这是良序集。对 $i = \alpha \in I$，令
$\mathcal{U}_i = \mathcal{U}_\alpha$。于是 $\mathcal{U}_i$
是约化局部闭子叠且是胚，即 (1) 成立。按构造，若
$i = \alpha \in I$，则 $T_i = T\alpha$，所以 (3) 成立。由我们对
$\mathcal{U}(T)$ 的选取，(4) 以及 (a) 或 (b) 也成立。最后，为证明 (2)，
取 $x \in |\mathcal{X}|$。存在使 $x \not \in T_\beta$ 的最小序数
$\beta$（因为序数是良序的）。由极限序数情形下 $T_\beta$ 的定义，
此时 $\beta$ 必为后继序数。因此 $\beta = \alpha + 1$，且
$x \in |\mathcal{U}(T_\alpha)|$，得证。
\end{proof}

\begin{remark}
\label{stacks-morphisms-remark-order-type}
我们可以问，引理 \ref{stacks-morphisms-lemma-every-point-in-a-stratum} 所构造的 (a) 型分层
给出的典范分层具有怎样的序型。一个自然猜想是，良序集 $I$ 的
{\it 基数}至多为 $\aleph_0$。我们不知道这是真是假。若你知道，请发送邮件至
\href{mailto:stacks.project@gmail.com}{stacks.project@gmail.com}.
\end{remark}







\section{代数叠的拓扑空间}
\label{stacks-morphisms-section-topology}

\noindent
本节把前面的结果应用于代数叠所关联的拓扑空间 $|\mathcal{X}|$。

\begin{lemma}
\label{stacks-morphisms-lemma-spectral-qc-diagonal-qc}
设 $\mathcal{X}$ 是拟紧代数叠，且其对角态射 $\Delta$ 拟紧。
则 $|\mathcal{X}|$ 是谱拓扑空间。
\end{lemma}

\begin{proof}
选取仿射概形 $U$ 以及满光滑态射 $U \to \mathcal{X}$，见《叠的性质》引理
\ref{stacks-properties-lemma-quasi-compact-stack}。于是
$|U| \to |\mathcal{X}|$ 连续、开且满，见《叠的性质》引理
\ref{stacks-properties-lemma-topology-points}。因此，$|\mathcal{X}|$
的拟紧开子集构成其拓扑的一组基。对拟紧开子集
$W_1, W_2 \subset |\mathcal{X}|$，可以选取 $U$ 的拟紧开子集
$V_1, V_2$，它们分别映到 $W_1$ 和 $W_2$。由于 $\Delta$ 拟紧，可知
$$
V_1 \times_\mathcal{X} V_2 =
(V_1 \times V_2) \times_{\mathcal{X} \times \mathcal{X}, \Delta} \mathcal{X}
$$
拟紧。由《叠的性质》引理
\ref{stacks-properties-lemma-points-cartesian}，
$|V_1 \times_\mathcal{X} V_2|$ 在 $|\mathcal{X}|$ 中的像是
$W_1 \cap W_2$。所以 $W_1 \cap W_2$ 拟紧。为完成证明，只需证明
$|\mathcal{X}|$ 是清醒的，见《拓扑学》定义
\ref{topology-definition-spectral-space}。

\medskip\noindent
设 $T \subset |\mathcal{X}|$ 是不可约闭子集。我们必须证明 $T$
有唯一的一般点。设 $\mathcal{Z} \subset \mathcal{X}$ 是与 $T$
对应的约化诱导闭子叠，见《叠的性质》定义
\ref{stacks-properties-definition-reduced-induced-stack}。
由于 $\mathcal{Z} \to \mathcal{X}$ 是闭浸入，可知
$\Delta_\mathcal{Z}$ 拟紧：首先，把
$\mathcal{Z} \to \mathcal{X} \times \mathcal{X}$ 写成
$\mathcal{Z} \to \mathcal{X}$ 与 $\Delta$ 的复合，从而证明它拟紧；
再把 $\mathcal{Z} \to \mathcal{X} \times \mathcal{X}$ 写成
$\Delta_\mathcal{Z}$ 与
$\mathcal{Z} \times \mathcal{Z} \to \mathcal{X} \times \mathcal{X}$
的复合，并使用引理 \ref{stacks-morphisms-lemma-quasi-compact-permanence} 以及
$\mathcal{Z} \times \mathcal{Z} \to \mathcal{X} \times \mathcal{X}$
分离这一事实。因此，问题归结为下一段讨论的情形。

\medskip\noindent
假设 $\mathcal{X}$ 拟紧，$\Delta$ 拟紧，$\mathcal{X}$ 约化，且
$|\mathcal{X}|$ 不可约。我们必须证明 $|\mathcal{X}|$ 有唯一的一般点。
由于 $\mathcal{I}_\mathcal{X} \to \mathcal{X}$ 是 $\Delta$ 的基变换，
可知 $\mathcal{I}_\mathcal{X} \to \mathcal{X}$ 拟紧
（引理 \ref{stacks-morphisms-lemma-base-change-quasi-compact}）。因此，由命题
\ref{stacks-morphisms-proposition-open-stratum}，存在一个稠密开子叠
$\mathcal{U} \subset \mathcal{X}$，它是胚。换言之，
$|\mathcal{U}| \subset |\mathcal{X}|$ 开且稠密。于是可以假设
$\mathcal{X}$ 是胚。设 $\mathcal{X} \to X$ 使 $\mathcal{X}$
成为代数空间 $X$ 上的胚。由引理 \ref{stacks-morphisms-lemma-gerbe-bijection-points}，
$|\mathcal{X}| \cong |X|$。特别地，$X$ 拟紧。由引理
\ref{stacks-morphisms-lemma-gerbe-diagonal-quasi-compact}，$X$ 的对角态射拟紧；也就是说，
$X$ 是拟分离代数空间。于是，由《代数空间的性质》引理
\ref{spaces-properties-lemma-quasi-compact-quasi-separated-spectral}，
$|X|$ 是谱空间，这就给出了所需结论。
\end{proof}

\begin{lemma}
\label{stacks-morphisms-lemma-spectral-qcqs}
设 $\mathcal{X}$ 是拟紧拟分离代数叠。则 $|\mathcal{X}|$ 是谱拓扑空间。
\end{lemma}

\begin{proof}
这是引理 \ref{stacks-morphisms-lemma-spectral-qc-diagonal-qc} 的特殊情形。
\end{proof}

\begin{lemma}
\label{stacks-morphisms-lemma-sober-qs}
设 $\mathcal{X}$ 是对角态射拟紧的代数叠（例如，若 $\mathcal{X}$ 拟分离）。
则存在开覆盖 $|\mathcal{X}| = \bigcup U_i$，其中每个 $U_i$ 都是谱空间。
特别地，$|\mathcal{X}|$ 是清醒拓扑空间。
\end{lemma}

\begin{proof}
这是引理 \ref{stacks-morphisms-lemma-spectral-qc-diagonal-qc} 的直接推论。
\end{proof}






\section{剩余胚的存在性}
\label{stacks-morphisms-section-existence-residual-gerbes}

\noindent
代数叠上一点的剩余胚定义见《叠的性质》定义
\ref{stacks-properties-definition-residual-gerbe}。我们已经证明，有限型点
（引理 \ref{stacks-morphisms-lemma-point-finite-type-monomorphism}）以及胚的任意一点
（引理 \ref{stacks-morphisms-lemma-gerbe-residual-gerbe-exists}）都有剩余胚。本节证明，
许多代数叠上都存在剩余胚。首先给出前面所承诺的命题
\ref{stacks-morphisms-proposition-open-stratum} 的应用。

\begin{lemma}
\label{stacks-morphisms-lemma-every-point-residual-gerbe}
设 $\mathcal{X}$ 是代数叠，并且
$\mathcal{I}_\mathcal{X} \to \mathcal{X}$ 拟紧。则对每个
$x \in |\mathcal{X}|$，$\mathcal{X}$ 在 $x$ 处的剩余胚都存在。
\end{lemma}

\begin{proof}
设 $T = \overline{\{x\}} \subset |\mathcal{X}|$ 是 $x$ 的闭包。
由《叠的性质》引理 \ref{stacks-properties-lemma-reduced-closed-substack}，
存在约化闭子叠 $\mathcal{X}' \subset \mathcal{X}$，使得
$T = |\mathcal{X}'|$。注意，由引理
\ref{stacks-morphisms-lemma-monomorphism-cartesian-square-inertia}，
$\mathcal{I}_{\mathcal{X}'} =
\mathcal{I}_\mathcal{X} \times_\mathcal{X} \mathcal{X}'$。因此，作为一个
基变换，$\mathcal{I}_{\mathcal{X}'} \to \mathcal{X}'$ 拟紧，见引理
\ref{stacks-morphisms-lemma-base-change-quasi-compact}。所以由命题
\ref{stacks-morphisms-proposition-open-stratum}，存在稠密开子叠
$\mathcal{U} \subset \mathcal{X}'$，它是胚。注意，
$x \in |\mathcal{U}|$，因为 $\{x\} \subset T$ 也是稠密子集。于是由引理
\ref{stacks-morphisms-lemma-gerbe-residual-gerbe-exists}，$\mathcal{U}$ 在 $x$ 处的剩余胚
$\mathcal{Z}_x \subset \mathcal{U}$ 存在。由定义立即可知，
$\mathcal{Z}_x \to \mathcal{X}$ 是 $\mathcal{X}$ 在 $x$ 处的剩余胚。
\end{proof}

\noindent
若叠是拟 DM 的，则剩余胚也存在。特别地，Deligne--Mumford 叠上总存在剩余胚。

\begin{lemma}
\label{stacks-morphisms-lemma-every-point-residual-gerbe-quasi-DM}
设 $\mathcal{X}$ 是拟 DM 代数叠。则对每个 $x \in |\mathcal{X}|$，
$\mathcal{X}$ 在 $x$ 处的剩余胚都存在。
\end{lemma}

\begin{proof}
选取概形 $U$ 以及满、平坦、局部有限表现且局部拟有限的态射
$U \to \mathcal{X}$，见定理 \ref{stacks-morphisms-theorem-quasi-DM}。令
$R = U \times_\mathcal{X} U$。作为态射 $U \to \mathcal{X}$ 的基变换，
投影 $s, t : R \to U$ 满、平坦、局部有限表现且局部拟有限。存在典范态射
$[U/R] \to \mathcal{X}$（见《代数叠》引理
\ref{algebraic-lemma-map-space-into-stack}）；由于
$U \to \mathcal{X}$ 满、平坦且局部有限表现，它是等价，见《代数叠》评注
\ref{algebraic-remark-flat-fp-presentation}。因此，可以假设
$\mathcal{X} = [U/R]$，其中 $(U, R, s, t, c)$ 是代数空间中的群胚，
且 $s, t : R \to U$ 满、平坦、局部有限表现且局部拟有限。令
$$
U' = \coprod\nolimits_{u \in U\text{ 位于下列对象之上： }x} \Spec(\kappa(u)).
$$
典范态射 $U' \to U$ 是单态射。令
$$
R' = U' \times_\mathcal{X} U' =
R \times_{(U \times U)} (U' \times U')
$$
由于 $U' \to U$ 是单态射，可知两个投影 $s', t' : R' \to U'$
都分解为一个单态射与一个局部拟有限态射的复合。因此，由于 $U'$
是域的谱的不交并，应用《域上的代数空间》引理
\ref{spaces-over-fields-lemma-mono-towards-locally-quasi-finite-over-field}，
可知态射 $s', t' : R' \to U'$ 局部拟有限。再次利用 $U'$ 是域的谱的不交并，
态射 $s', t'$ 也平坦。最后，$s', t'$ 局部拟有限蕴含 $s', t'$
局部有限型，进而 $s', t'$ 局部有限表现（因为 $U'$ 是域的谱的不交并，特别地局部诺特，
所以可以应用《代数空间的态射》引理
\ref{spaces-morphisms-lemma-noetherian-finite-type-finite-presentation}）。
于是，由《可表示性判据》定理
\ref{criteria-theorem-flat-groupoid-gives-algebraic-stack}，
$\mathcal{Z} = [U'/R']$ 是代数叠。由于 $R'$ 是 $R$ 沿 $U' \to U$
的限制，由《空间中的群胚》引理
\ref{spaces-groupoids-lemma-quotient-stack-restrict} 和《叠的性质》引理
\ref{stacks-properties-lemma-monomorphism}，可知
$\mathcal{Z} \to \mathcal{X}$ 是单态射。由于
$\mathcal{Z} \to \mathcal{X}$ 是单态射，因此
$|\mathcal{Z}| \to |\mathcal{X}|$ 单，见《叠的性质》引理
\ref{stacks-properties-lemma-monomorphism-injective-points}。由《叠的性质》引理
\ref{stacks-properties-lemma-points-cartesian}，可知
$$
|U'| = |\mathcal{Z} \times_\mathcal{X} U'|
\longrightarrow
|\mathcal{Z}| \times_{|\mathcal{X}|} |U'|
$$
满；结合我们对 $U'$ 的选取，这蕴含
$|\mathcal{Z}| \to |\mathcal{X}|$ 的像为 $\{x\}$。所以
$|\mathcal{Z}|$ 是单点集。最后，按构造，$U'$ 局部诺特且约化；也就是说，
$\mathcal{Z}$ 约化且局部诺特。这说明 $\mathcal{Z} \to \mathcal{X}$
的本质像是 $\mathcal{X}$ 在 $x$ 处的剩余胚，见《叠的性质》引理
\ref{stacks-properties-lemma-residual-gerbe-unique}。
\end{proof}

\begin{lemma}
\label{stacks-morphisms-lemma-every-point-residual-gerbe-locally-Noetherian}
设 $\mathcal{X}$ 是局部诺特代数叠。则对每个
$x \in |\mathcal{X}|$，$\mathcal{X}$ 在 $x$ 处的剩余胚都存在。
\end{lemma}

\begin{proof}
选取仿射概形 $U$ 和光滑态射 $U \to \mathcal{X}$，使 $x$ 位于开连续映射
$|U| \to |\mathcal{X}|$ 的像中。我们可以并且确实用与
$|U| \to |\mathcal{X}|$ 的像对应的开子叠替换 $\mathcal{X}$，见
《叠的性质》引理 \ref{stacks-properties-lemma-open-substacks}。因此，可以假设
$\mathcal{X} = [U/R]$，其中 $(U, R, s, t, c)$ 是代数空间中的光滑群胚，
且 $U$ 是诺特仿射概形；见《代数叠》一节
\ref{algebraic-section-stack-to-presentation}。

\medskip\noindent
设 $E \subset |U|$ 是 $\{x\} \subset |\mathcal{X}|$ 的逆像。当然，
$E \not = \emptyset$。由于 $|U|$ 是诺特拓扑空间，可以选取
$u \in E$，使得 $\overline{\{u\}} \cap E = \{u\}$。像往常一样，把
$u = \Spec(\kappa(u))$ 看作其剩余域的谱。记
$$
F = u \times_{U, t} R = u \times_\mathcal{X} U
\quad\text{且}\quad
R' = (u \times u) \times_{(U \times U), (t, s)} R = u \times_\mathcal{X} u
$$
。此外，记 $Z = \overline{\{u\}} \subset U$，赋予它约化诱导概形结构。
记 $p : F \to U$ 为第二投影诱导的态射（在 $F$ 的第一个纤维积描述中，
使用 $s : R \to U$）。则 $E$ 是 $p$ 的集合论像。态射 $R' \to F$
是单态射，并且经由 $Z$ 的逆像 $p^{-1}(Z)$ 分解。该逆像
$p^{-1}(Z) \subset F$ 是闭子概形；由于上面对 $u \in E$ 的谨慎选取，
限制 $p|_{p^{-1}(Z)} : p^{-1}(Z) \to Z$ 的集合论像包含于
$\{u\} \subset Z$。由于 $u = \lim W$，其中极限遍历不可约约化概形
$Z$ 的非空仿射开子概形，可知态射
$p|_{p^{-1}(Z)} : p^{-1}(Z) \to Z$ 经由态射 $u \to Z$ 分解。这显然蕴含
$R' = p^{-1}(Z)$。特别地，态射 $t' : R' \to u$ 局部有限表现，因为它是
局部诺特代数空间的闭浸入 $p^{-1}(Z) \to F$ 与光滑态射
$\text{pr}_1 : F \to u$ 的复合；使用《代数空间的态射》诸引理
\ref{spaces-morphisms-lemma-locally-finite-type-locally-noetherian}、
\ref{spaces-morphisms-lemma-closed-immersion-finite-presentation} 和
\ref{spaces-morphisms-lemma-composition-finite-presentation}。因此，
$(U, R, s, t, c)$ 沿 $u \to U$ 的限制 $(u, R', s', t', c')$
是代数空间中的群胚，其中 $s'$ 和 $t'$ 平坦且局部有限表现。所以由
《可表示性判据》定理 \ref{criteria-theorem-flat-groupoid-gives-algebraic-stack}，
$\mathcal{Z} = [u/R']$ 是代数叠。由于 $R'$ 是 $R$ 沿 $u \to U$
的限制，由《空间中的群胚》引理
\ref{spaces-groupoids-lemma-quotient-stack-restrict} 和《叠的性质》引理
\ref{stacks-properties-lemma-monomorphism}，可知
$\mathcal{Z} \to \mathcal{X}$ 是单态射。再由《叠的性质》一节
\ref{stacks-properties-section-residual-gerbe} 的内容，
$\mathcal{Z}$（同构于）该剩余胚。
\end{proof}













\section{\'Etale 局部结构}
\label{stacks-morphisms-section-etale-local}

\noindent
本节开始讨论关于代数叠的 \'etale 局部结构可以说些什么。

\begin{lemma}
\label{stacks-morphisms-lemma-quotient-etale}
设 $Y$ 是代数空间。设 $(U, R, s, t, c)$ 是 $Y$ 上代数空间中的群胚。
假设 $U \to Y$ 平坦且局部有限表现，并且
$R \to U \times_Y U$ 是开浸入。则 $X = [U/R] = U/R$ 是代数空间，
且 $X \to Y$ 是 \'etale 的。
\end{lemma}

\begin{proof}
由《可表示性判据》定理
\ref{criteria-theorem-flat-groupoid-gives-algebraic-stack}，
商叠 $[U/R]$ 是代数叠。另一方面，由于
$R \to U \times U$ 是单态射，它是代数空间（这里沿用我们的语言滥用，并见
《代数叠》命题
\ref{algebraic-proposition-algebraic-stack-no-automorphisms}），当然也等于
代数空间 $U/R$（《引导》定理 \ref{bootstrap-theorem-final-bootstrap}）。
由于 $U \to X$ 满、平坦且局部有限表现
（《引导》引理 \ref{bootstrap-lemma-covering-quotient}），由《代数空间的态射》
引理 \ref{spaces-morphisms-lemma-flat-permanence} 和《代数空间上的下降》引理
\ref{spaces-descent-lemma-flat-finitely-presented-permanence}，可知
$X \to Y$ 平坦且局部有限表现。最后，考虑笛卡儿图
$$
\xymatrix{
R \ar[d] \ar[r] & U \times_Y U \ar[d] \\
X \ar[r] & X \times_Y X
}
$$
由于右侧竖直箭头满、平坦且局部有限表现（略去一个小细节），而顶部水平箭头
按假设是开浸入，可知 $X \to X \times_Y X$ 是开浸入（使用《叠的性质》引理
\ref{stacks-properties-lemma-check-property-covering}）。所以由《代数空间的态射》
引理 \ref{spaces-morphisms-lemma-diagonal-unramified-morphism}，
$X \to Y$ 非分歧。最后应用《代数空间的态射》引理
\ref{spaces-morphisms-lemma-unramified-flat-lfp-etale} 即得结论。
\end{proof}

\begin{lemma}
\label{stacks-morphisms-lemma-quasi-splitting-etale}
设 $S$ 是概形。设 $(U, R, s, t, c)$ 是 $S$ 上代数空间中的群胚。
假设 $s, t$ 平坦且局部有限表现。设 $P \subset R$ 是开子空间，使得
$(U, P, s|_P, t|_P, c|_{P \times_{s, U, t} P})$
构成 $S$ 上代数空间中的群胚。则
$$
[U/P] \longrightarrow [U/R]
$$
是代数叠的态射，可由代数空间表示、满且 \'etale。
\end{lemma}

\begin{proof}
由于 $P \subset R$ 是开的，$s|_P$ 和 $t|_P$ 平坦且局部有限表现。因此，
由《可表示性判据》定理
\ref{criteria-theorem-flat-groupoid-gives-algebraic-stack}，
$[U/R]$ 和 $[U/P]$ 都是代数叠。为证明该态射可由代数空间表示，只需证明
$[U/P] \to [U/R]$ 在纤维范畴上忠实，见《代数叠》引理
\ref{algebraic-lemma-characterize-representable-by-algebraic-spaces}。
这立即由 $P \to R$ 是单态射以及《空间中的群胚》引理
\ref{spaces-groupoids-lemma-quotient-stack-objects} 给出的商叠对象的显式描述得出。
至此，由《代数叠》定义
\ref{algebraic-definition-relative-representable-property}，我们知道
$[U/P] \to [U/R]$ 满和 \'etale 分别意味着什么。例如，满性由
《可表示性判据》引理 \ref{criteria-lemma-representable-by-spaces-surjective}
以及域的谱上的商叠对象描述（见上引引理）得出。还需证明该态射是 \'etale 的。

\medskip\noindent
为此，只需证明 $U \times_{[U/R]} [U/P] \to U$ 是 \'etale 的，见
《叠的性质》引理 \ref{stacks-properties-lemma-check-property-covering}。
由《空间中的群胚》引理
\ref{spaces-groupoids-lemma-cartesian-square-of-morphism}，该纤维积等于
$[R/P \times_{s, U, t} R]$，而到 $U$ 的态射由
$s : R \to U$ 诱导。映射
$s', t' : P \times_{s, U, t} R \to R$ 分别由
$s' : (p, r) \mapsto r$ 和 $t' : (p, r) \mapsto c(p, r)$ 给出。由于
$P \subset R$ 是开的，可知
$(t', s') : P \times_{s, U, t} R \to R \times_{s, U, s} R$
是开浸入。因此应用引理 \ref{stacks-morphisms-lemma-quotient-etale} 即得结论。
\end{proof}

\begin{lemma}
\label{stacks-morphisms-lemma-etale-local-quasi-DM}
设 $\mathcal{X}$ 是代数叠。假设 $\mathcal{X}$ 拟 DM 且对角态射分离
（等价地，$\mathcal{I}_\mathcal{X} \to \mathcal{X}$
局部拟有限且分离）。设 $x \in |\mathcal{X}|$。则存在代数叠的态射
$$
\mathcal{U} \longrightarrow \mathcal{X}
$$
，具有下列性质：
\begin{enumerate}
\item 存在映到 $x$ 的点 $u \in |\mathcal{U}|$；
\item $\mathcal{U} \to \mathcal{X}$ 可由代数空间表示且 \'etale；
\item $\mathcal{U} = [U/R]$，其中 $(U, R, s, t, c)$ 是概形群胚，
$U$、$R$ 仿射，而 $s, t$ 有限、平坦且局部有限表现。
\end{enumerate}
\end{lemma}

\begin{proof}
（括号中的陈述由引理 \ref{stacks-morphisms-lemma-diagonal-diagonal} 中的等价条件得出。）
选取仿射概形 $U$ 以及平坦、局部有限表现、局部拟有限的态射
$U \to \mathcal{X}$，使得 $x$ 是某个点 $u \in U$ 的像。由定理
\ref{stacks-morphisms-theorem-quasi-DM} 以及 $\mathcal{X}$ 拟 DM 的假设，这可以做到。
令 $(U, R, s, t, c)$ 为取 $R = U \times_\mathcal{X} U$
所得的代数空间中的群胚，见《代数叠》引理
\ref{algebraic-lemma-map-space-into-stack}。设
$\mathcal{X}' \subset \mathcal{X}$ 是与 $|U| \to |\mathcal{X}|$
的开像对应的开子叠（《叠的性质》诸引理
\ref{stacks-properties-lemma-topology-points} 和
\ref{stacks-properties-lemma-open-substacks}）。显然，可以用开子叠
$\mathcal{X}'$ 替换 $\mathcal{X}$。于是可以假设
$U \to \mathcal{X}$ 满，此时《代数叠》评注
\ref{algebraic-remark-flat-fp-presentation} 给出
$\mathcal{X} = [U/R]$。注意，$s, t : R \to U$ 平坦、局部有限表现且
局部拟有限。由于
$R = U \times U \times_{(\mathcal{X} \times \mathcal{X})} \mathcal{X}$，
且 $\mathcal{X}$ 的对角态射分离，可知 $R$ 分离。因此
$s, t : R \to U$ 分离。将《代数空间的态射》命题
\ref{spaces-morphisms-proposition-locally-quasi-finite-separated-over-scheme}
应用于 $s : R \to U$，可知 $R$ 是概形。

\medskip\noindent
上文已经验证，$(U, R, s, t, c)$ 与 $u$ 满足
《空间中群胚的进一步内容》引理
\ref{spaces-more-groupoids-lemma-quasi-splitting-affine-scheme}
的所有假设。因此，可以找到基本 \'etale 邻域
$(U', u') \to (U, u)$，使得 $R$ 在 $U'$ 上的限制 $R'$ 在 $u$ 上拟分裂。
注意，$R' = U' \times_\mathcal{X} U'$（略去一个小细节；提示：使用纤维积的
传递性）。用 $(U', R', s', t', c')$ 替换 $(U, R, s, t, c)$，并像上文那样
缩小 $\mathcal{X}$，可以假设 $(U, R, s, t, c)$ 在 $u$ 上有一个拟分裂
（从《空间中群胚的进一步内容》定义
\ref{spaces-more-groupoids-definition-split-at-point} 的脚注可见，
从此以后点 $u$ 已无关紧要）。设 $P \subset R$ 是 $R$ 在 $u$ 上的一个拟分裂。
应用引理 \ref{stacks-morphisms-lemma-quasi-splitting-etale} 可知
$$
\mathcal{U} = [U/P] \longrightarrow [U/R] = \mathcal{X}
$$
具有所有所需性质。
\end{proof}

\begin{lemma}
\label{stacks-morphisms-lemma-etale-local-quasi-DM-at-x}
设 $\mathcal{X}$ 是代数叠。假设 $\mathcal{X}$ 拟 DM 且对角态射分离
（等价地，$\mathcal{I}_\mathcal{X} \to \mathcal{X}$ 局部拟有限且分离）。
设 $x \in |\mathcal{X}|$。假设 $\mathcal{X}$ 在 $x$ 处的自同构群有限
（评注 \ref{stacks-morphisms-remark-property-automorphism-groups}）。则存在代数叠的态射
$$
g : \mathcal{U} \longrightarrow \mathcal{X}
$$
，具有下列性质：
\begin{enumerate}
\item 存在映到 $x$ 的点 $u \in |\mathcal{U}|$，且 $g$ 诱导 $u$ 与 $x$
处自同构群之间的同构（评注 \ref{stacks-morphisms-remark-identify-automorphism-groups}）；
\item $\mathcal{U} \to \mathcal{X}$ 可由代数空间表示且 \'etale；
\item $\mathcal{U} = [U/R]$，其中 $(U, R, s, t, c)$ 是概形群胚，
$U$、$R$ 仿射，而 $s, t$ 有限、平坦且局部有限表现。
\end{enumerate}
\end{lemma}

\begin{proof}
由引理 \ref{stacks-morphisms-lemma-automorphism-group-scheme}，$G_x$ 是群概形。证明的第一部分
与引理 \ref{stacks-morphisms-lemma-etale-local-quasi-DM} 的证明第一部分{\bf 完全}相同。因此，
可以假设 $\mathcal{X} = [U/R]$，其中 $(U, R, s, t, c)$ 以及映到 $x$ 的
$u \in U$ 满足《空间中群胚的进一步内容》引理
\ref{spaces-more-groupoids-lemma-quasi-splitting-affine-scheme}
的全部假设。关于 $G_x$ 的假设蕴含 $G_u$ 在 $u$ 上有限。因此，
《空间中群胚的进一步内容》引理
\ref{spaces-more-groupoids-lemma-splitting-affine-scheme} 的全部假设都成立。
所以可以找到基本 \'etale 邻域 $(U', u') \to (U, u)$，使得 $R$ 在 $U'$
上的限制 $R'$ 在 $u$ 上分裂。注意，
$R' = U' \times_\mathcal{X} U'$（略去一个小细节；提示：使用纤维积的传递性）。
用 $(U', R', s', t', c')$ 替换 $(U, R, s, t, c)$，并像上文那样缩小
$\mathcal{X}$，可以假设 $(U, R, s, t, c)$ 在 $u$ 上有一个分裂。
设 $P \subset R$ 是 $R$ 在 $u$ 上的一个分裂。应用引理
\ref{stacks-morphisms-lemma-quasi-splitting-etale} 可知
$$
\mathcal{U} = [U/P] \longrightarrow [U/R] = \mathcal{X}
$$
可由代数空间表示且 \'etale。按构造，$G_u$ 包含于 $P$，所以该态射在
$u$ 处的自同构群上给出所需的同构。
\end{proof}

\begin{lemma}
\label{stacks-morphisms-lemma-etale-local-quasi-DM-at-x-inertia}
设 $\mathcal{X}$ 是代数叠。假设 $\mathcal{X}$ 拟 DM 且对角态射分离
（等价地，$\mathcal{I}_\mathcal{X} \to \mathcal{X}$ 局部拟有限且分离）。
设 $x \in |\mathcal{X}|$。假设 $x$ 可以由拟紧态射
$\Spec(k) \to \mathcal{X}$ 表示。则存在代数叠的态射
$$
g : \mathcal{U} \longrightarrow \mathcal{X}
$$
，具有下列性质：
\begin{enumerate}
\item 存在映到 $x$ 的点 $u \in |\mathcal{U}|$，且 $g$ 诱导 $u$ 与 $x$
处剩余胚之间的同构；
\item $\mathcal{U} \to \mathcal{X}$ 可由代数空间表示且 \'etale；
\item $\mathcal{U} = [U/R]$，其中 $(U, R, s, t, c)$ 是概形群胚，
$U$、$R$ 仿射，而 $s, t$ 有限、平坦且局部有限表现。
\end{enumerate}
\end{lemma}

\begin{proof}
证明的第一部分与引理 \ref{stacks-morphisms-lemma-etale-local-quasi-DM} 的证明第一部分
{\bf 完全}相同。因此，可以假设 $\mathcal{X} = [U/R]$，其中
$(U, R, s, t, c)$ 以及映到 $x$ 的 $u \in U$ 满足
《空间中群胚的进一步内容》引理
\ref{spaces-more-groupoids-lemma-quasi-splitting-affine-scheme}
的全部假设。注意，$u = \Spec(\kappa(u)) \to \mathcal{X}$ 拟紧，见
《叠的性质》引理
\ref{stacks-properties-lemma-UR-quasi-compact-above-x}。考虑笛卡儿图
$$
\xymatrix{
F \ar[d] \ar[r] & U \ar[d] \\
u \ar[r]^u & \mathcal{X}
}
$$
由于 $U$ 是仿射概形且 $F \to U$ 拟紧，可知 $F$ 拟紧。由于
$U \to \mathcal{X}$ 局部拟有限，可知 $F \to u$ 局部拟有限。因此
$F \to u$ 拟有限，并且 $F$ 是仿射概形，其底拓扑空间有限离散
（《域上的代数空间》引理
\ref{spaces-over-fields-lemma-locally-quasi-finite-over-field}）。
注意，我们有单态射 $u \times_\mathcal{X} u \to F$。特别地，集合
$\{r \in R : s(r) = u, t(r) = u\}$ 是
$|u \times_\mathcal{X} u| \to |R|$ 的像，故为有限集。于是，
《空间中群胚的进一步内容》引理
\ref{spaces-more-groupoids-lemma-strong-splitting-affine-scheme}
的全部假设成立。

\medskip\noindent
所以可以找到基本 \'etale 邻域 $(U', u') \to (U, u)$，使得 $R$ 在 $U'$
上的限制 $R'$ 在 $u'$ 上强分裂。注意，
$R' = U' \times_\mathcal{X} U'$（略去一个小细节；提示：使用纤维积的传递性）。
用 $(U', R', s', t', c')$ 替换 $(U, R, s, t, c)$，并像上文那样缩小
$\mathcal{X}$，可以假设 $(U, R, s, t, c)$ 在 $u$ 上有一个强分裂。
设 $P \subset R$ 是 $R$ 在 $u$ 上的一个强分裂。应用引理
\ref{stacks-morphisms-lemma-quasi-splitting-etale} 可知
$$
\mathcal{U} = [U/P] \longrightarrow [U/R] = \mathcal{X}
$$
可由代数空间表示且 \'etale。由于 $P \subset R$ 是开的，并且按构造包含
$\{r \in R : s(r) = u, t(r) = u\}$，可知
$u \times_\mathcal{U} u \to u \times_\mathcal{X} u$ 是同构。关于剩余胚的
陈述遂由《叠的性质》引理
\ref{stacks-properties-lemma-residual-gerbe-isomorphic} 得出
（注意，由引理 \ref{stacks-morphisms-lemma-every-point-residual-gerbe-quasi-DM}，
所涉及的剩余胚确实存在）。
\end{proof}







\section{光滑态射}
\label{stacks-morphisms-section-smooth}

\noindent
代数空间态射“为光滑”的性质在源与靶上光滑局部，见《代数空间上的下降》评注
\ref{spaces-descent-remark-list-local-source-target}。该性质在基变换下也稳定，
并且在靶上 fpqc 局部，见《代数空间的态射》引理
\ref{spaces-morphisms-lemma-base-change-smooth} 和《代数空间上的下降》引理
\ref{spaces-descent-lemma-descending-property-smooth}。因此，由上面的引理
\ref{stacks-morphisms-lemma-local-source-target}，可以如下定义代数空间的态射何时光滑；
当该态射可由代数空间表示时，这一定义与《叠的性质》一节
\ref{stacks-properties-section-properties-morphisms} 中已有的概念一致。

\begin{definition}
\label{stacks-morphisms-definition-smooth}
设 $f : \mathcal{X} \to \mathcal{Y}$ 是代数叠的态射。若引理
\ref{stacks-morphisms-lemma-local-source-target} 的等价条件对
$\mathcal{P} = \text{smooth}$ 成立，则称 $f$ {\it 光滑}。
\end{definition}

\begin{lemma}
\label{stacks-morphisms-lemma-composition-smooth}
光滑态射的复合仍光滑。
\end{lemma}

\begin{proof}
结合评注 \ref{stacks-morphisms-remark-composition} 与《代数空间的态射》引理
\ref{spaces-morphisms-lemma-composition-smooth}。
\end{proof}

\begin{lemma}
\label{stacks-morphisms-lemma-base-change-smooth}
光滑态射的基变换仍光滑。
\end{lemma}

\begin{proof}
结合评注 \ref{stacks-morphisms-remark-base-change} 与《代数空间的态射》引理
\ref{spaces-morphisms-lemma-base-change-smooth}。
\end{proof}

\begin{lemma}
\label{stacks-morphisms-lemma-descent-smooth}
设 $f : \mathcal{X} \to \mathcal{Y}$ 是代数叠的态射。设
$\mathcal{Z} \to \mathcal{Y}$ 是代数叠的满、平坦、局部有限表现态射。
若基变换
$\mathcal{Z} \times_\mathcal{Y} \mathcal{X} \to \mathcal{Z}$
光滑，则 $f$ 光滑。
\end{lemma}

\begin{proof}
“光滑”这一性质满足引理 \ref{stacks-morphisms-lemma-descent-property} 的条件。在源与靶上
光滑局部这一点已在本节导言中看到，而在靶上 fppf 局部是
《代数空间上的下降》引理
\ref{spaces-descent-lemma-descending-property-smooth}。
\end{proof}

\begin{lemma}
\label{stacks-morphisms-lemma-smooth-locally-finite-presentation}
代数叠的光滑态射局部有限表现。
\end{lemma}

\begin{proof}
略。
\end{proof}

\begin{lemma}
\label{stacks-morphisms-lemma-where-smooth}
设 $f : \mathcal{X} \to \mathcal{Y}$ 是代数叠的态射。存在最大的开子叠
$\mathcal{U} \subset \mathcal{X}$，使得
$f|_\mathcal{U} : \mathcal{U} \to \mathcal{Y}$ 光滑。此外，
该开子叠的构造与下列操作交换：
\begin{enumerate}
\item 前复合光滑态射；
\item 沿平坦且局部有限表现的态射作基变换；
\item 在 $f$ 局部有限表现时，沿平坦态射作基变换。
\end{enumerate}
\end{lemma}

\begin{proof}
选取交换图
$$
\xymatrix{
U \ar[d]_a \ar[r]_h & V \ar[d]^b \\
\mathcal{X} \ar[r]^f & \mathcal{Y}
}
$$
，其中 $U$ 和 $V$ 是代数空间，竖直箭头光滑，而
$a : U \to \mathcal{X}$ 满。存在最大的开子空间 $U' \subset U$，
使得 $h_{U'} : U' \to V$ 光滑，见《代数空间的态射》引理
\ref{spaces-morphisms-lemma-where-smooth}。设
$\mathcal{U} \subset \mathcal{X}$ 是与 $|U'| \to |\mathcal{X}|$
的像对应的开子叠（《叠的性质》诸引理
\ref{stacks-properties-lemma-topology-points} 和
\ref{stacks-properties-lemma-open-substacks}）。由引理
\ref{stacks-morphisms-lemma-local-source-target} 中的等价性，可知
$f|_\mathcal{U} : \mathcal{U} \to \mathcal{Y}$ 光滑，而且
$\mathcal{U}$ 是具有该性质的最大开子叠。

\medskip\noindent
断言 (1) 由光滑性在源上光滑局部这一事实得出（代数叠的光滑态射正是利用该性质
定义的）。断言 (2) 和 (3) 由代数空间的情形得出，见《代数空间的态射》引理
\ref{spaces-morphisms-lemma-where-smooth}。
\end{proof}

\begin{lemma}
\label{stacks-morphisms-lemma-smooth-quotient-stack}
设 $X \to Y$ 是代数空间的光滑态射。设 $G$ 是 $Y$ 上的群代数空间，
它在 $Y$ 上平坦且局部有限表现。设 $G$ 在 $Y$ 上作用于 $X$。
则商叠 $[X/G]$ 在 $Y$ 上光滑。
\end{lemma}

\noindent
即使 $G$ 在 $Y$ 上不光滑，这也成立！

\begin{proof}
由《可表示性判据》定理
\ref{criteria-theorem-flat-groupoid-gives-algebraic-stack}，
商 $[X/G]$ 是代数叠。$[X/G]$ 在 $Y$ 上的光滑性由光滑性可沿 fppf 覆盖
下降这一事实得出：选取满光滑态射 $U \to [X/G]$，其中 $U$ 是概形。
$[X/G]$ 在 $Y$ 上光滑，等价于 $U$ 在 $Y$ 上光滑。注意，
$U \times_{[X/G]} X$ 在 $X$ 上光滑，因而也在 $Y$ 上光滑
（因为光滑态射的复合仍光滑）。另一方面，
$U \times_{[X/G]} X \to U$ 局部有限表现、平坦且满
（因为它是具有这些性质的 $X \to [X/G]$ 的基变换；例如见
《可表示性判据》引理
\ref{criteria-lemma-flat-quotient-flat-presentation}）。因此，可以应用
《代数空间上的下降》引理 \ref{spaces-descent-lemma-smooth-permanence}。
\end{proof}

\begin{lemma}
\label{stacks-morphisms-lemma-gerbe-smooth}
设 $\pi : \mathcal{X} \to \mathcal{Y}$ 是代数叠的态射。若
$\mathcal{X}$ 是 $\mathcal{Y}$ 上的胚，则 $\pi$ 满且光滑。
\end{lemma}

\begin{proof}
满性已经在引理 \ref{stacks-morphisms-lemma-gerbe-fppf} 中得到。由引理
\ref{stacks-morphisms-lemma-descent-smooth}，只需沿满、平坦、局部有限表现的态射
$\mathcal{Y}' \to \mathcal{Y}$ 作基变换，用所得基变换替换 $\pi$ 后证明本引理。
由引理 \ref{stacks-morphisms-lemma-local-structure-gerbe}，可以假设
$\mathcal{Y} = U$ 是代数空间，而 $\mathcal{X} = [U/G]$ 是 $U$ 上的叠，
其中 $G \to U$ 平坦且局部有限表现。此时由引理
\ref{stacks-morphisms-lemma-smooth-quotient-stack} 得证。
\end{proof}






\section{在源与靶上 \'etale--光滑局部的态射类型}
\label{stacks-morphisms-section-etale-smooth-local-source-target}

\noindent
给定代数空间态射的一个{\it 在源与靶上 \'etale--光滑局部}的性质
（见《代数空间上的下降》定义
\ref{spaces-descent-definition-etale-smooth-local-source-target}），
可以用它定义代数叠的 DM 态射的相应性质，即要求下述引理的任一等价条件成立。

\begin{lemma}
\label{stacks-morphisms-lemma-etale-smooth-local-source-target}
设 $\mathcal{P}$ 是代数空间态射的一个在源与靶上 \'etale--光滑局部的性质。
设 $f : \mathcal{X} \to \mathcal{Y}$ 是代数叠的 DM 态射。考虑交换图
$$
\xymatrix{
U \ar[d]_a \ar[r]_h & V \ar[d]^b \\
\mathcal{X} \ar[r]^f & \mathcal{Y}
}
$$
，其中 $U$ 和 $V$ 是代数空间，$V \to \mathcal{Y}$ 光滑，而
$U \to \mathcal{X} \times_\mathcal{Y} V$ 是 \'etale 的。下列条件等价：
\begin{enumerate}
\item 对任意上述交换图，态射 $h$ 具有性质 $\mathcal{P}$；
\item 对某个满足 $a : U \to \mathcal{X}$ 为满的上述交换图，态射 $h$
具有性质 $\mathcal{P}$。
\end{enumerate}
若 $\mathcal{X}$ 和 $\mathcal{Y}$ 可由代数空间表示，则这些条件还等价于
$f$（作为代数空间的态射）具有性质 $\mathcal{P}$。若 $\mathcal{P}$
还在任意基变换下保持，并且在基上 fppf 局部，则对可由代数空间表示的态射
$f$，这些条件还等价于 $f$ 在《叠的性质》一节
\ref{stacks-properties-section-properties-morphisms} 的意义下具有性质
$\mathcal{P}$。
\end{lemma}

\begin{proof}
先证明 (1) $\Rightarrow$ (2)。选取代数空间 $V$ 以及满光滑态射
$V \to \mathcal{Y}$。由于 $f$ 是 DM 的，存在概形 $U$ 以及满 \'etale 态射
$U \to V \times_\mathcal{Y} \mathcal{X}$，见引理 \ref{stacks-morphisms-lemma-DM}。
于是可见 (2) 成立。注意，$U \to \mathcal{X}$ 也满且光滑，因为它是基变换
$\mathcal{X} \times_\mathcal{Y} V \to \mathcal{X}$ 与所选映射
$U \to \mathcal{X} \times_\mathcal{Y} V$ 的复合。因此得到了 (1)
中所述的交换图。所以若 (1) 成立，则 $h : U \to V$ 具有性质
$\mathcal{P}$；又因 $U \to \mathcal{X}$ 满，这正说明 (2) 成立。

\medskip\noindent
反之，假设 (2) 成立，并设 $U, V, a, b, h$ 如 (2) 所述。再任取 (1)
中所述的交换图 $U', V', a', b', h'$。图示为
$$
\xymatrix{
U \ar[d] \ar[r]_h & V \ar[d] \\
\mathcal{X} \ar[r]^f & \mathcal{Y}
}
\quad\quad
\xymatrix{
U' \ar[d] \ar[r]_{h'} & V' \ar[d] \\
\mathcal{X} \ar[r]^f & \mathcal{Y}
}
$$
为证明 (2) 蕴含 (1)，必须证明 $h'$ 具有 $\mathcal{P}$。为此，考虑代数空间的
交换图
$$
\xymatrix{
U \ar[d]^h &
U \times_\mathcal{X} U' \ar[l] \ar[d]^{(h, h')} \ar[r] &
U' \ar[d]^{h'} \\
V &
V \times_\mathcal{Y} V' \ar[l] \ar[r] &
V'
}
$$
。注意，水平箭头都是光滑态射
$V \to \mathcal{Y}$、$V' \to \mathcal{Y}$、$U \to \mathcal{X}$
和 $U' \to \mathcal{X}$ 的基变换，因而光滑。还要注意，下列方块
$$
\xymatrix{
U \ar[d] & U \times_\mathcal{X} U' \ar[l] \ar[d] &
U \times_\mathcal{X} U' \ar[d] \ar[r] & U' \ar[d] \\
V \times_\mathcal{Y} \mathcal{X} &
V \times_\mathcal{Y} U' \ar[l] &
U \times_\mathcal{Y} V' \ar[r] &
\mathcal{X} \times_\mathcal{Y} V'
}
$$
是笛卡儿的，故由对 $U', V', a', b', h'$ 和 $U, V, a, b, h$ 的假设，
竖直箭头都是 \'etale 的。由《代数空间上的下降》引理
\ref{spaces-descent-lemma-etale-smooth-local-source-target-implies} 的 (2)，
$\mathcal{P}$ 在靶上光滑局部，因此 $h$ 的基变换
$t : U \times_\mathcal{Y} V' \to V \times_\mathcal{Y} V'$
具有 $\mathcal{P}$。由《代数空间上的下降》引理
\ref{spaces-descent-lemma-etale-smooth-local-source-target-implies} 的 (1)，
$\mathcal{P}$ 在源上 \'etale 局部；
又因 $s : U \times_\mathcal{X} U' \to U \times_\mathcal{Y} V'$ 是
\'etale 的，可知态射 $(h, h') = t \circ s$ 具有 $\mathcal{P}$。考虑交换图
$$
\xymatrix{
U \times_\mathcal{X} U' \ar[r]_{(h, h')} \ar[d] &
V \times_\mathcal{Y} V' \ar[d] \\
U' \ar[r]^{h'} & V'
}
$$
左侧竖直箭头满，右侧竖直箭头光滑，而诱导态射
$$
U \times_\mathcal{X} U'
\longrightarrow
U' \times_{V'} (V \times_\mathcal{Y} V') = V \times_\mathcal{Y} U'
$$
如上所见是 \'etale 的。因此，由《代数空间上的下降》定义
\ref{spaces-descent-definition-etale-smooth-local-source-target} 的 (3)，
可知 $h'$ 具有 $\mathcal{P}$。这就证明了 (1) 与 (2) 的等价性。

\medskip\noindent
若 $\mathcal{X}$ 和 $\mathcal{Y}$ 可表，则《代数空间上的下降》引理
\ref{spaces-descent-lemma-etale-smooth-local-source-target-characterize}
适用，并说明 (1) 和 (2) 等价于 $f$ 具有 $\mathcal{P}$。

\medskip\noindent
最后，假设 $f$ 可表，$U, V, a, b, h$ 如本引理 (2) 所述，并且
$\mathcal{P}$ 在任意基变换下保持。必须证明，对任意概形 $Z$ 和态射
$Z \to \mathcal{X}$，基变换
$Z \times_\mathcal{Y} \mathcal{X} \to Z$
具有性质 $\mathcal{P}$。考虑交换图
$$
\xymatrix{
Z \times_\mathcal{Y} U \ar[d] \ar[r] &
Z \times_\mathcal{Y} V \ar[d] \\
Z \times_\mathcal{Y} \mathcal{X} \ar[r] &
Z
}
$$
注意，顶部水平箭头是 $h$ 的基变换，因而具有性质 $\mathcal{P}$。
左侧竖直箭头满，诱导态射
$$
Z \times_\mathcal{Y} U \longrightarrow
(Z \times_\mathcal{Y} \mathcal{X}) \times_Z (Z \times_\mathcal{Y} V)
$$
是 \'etale 的，而右侧竖直箭头光滑。因此，《代数空间上的下降》引理
\ref{spaces-descent-lemma-etale-smooth-local-source-target-characterize}
适用，并说明 $Z \times_\mathcal{Y} \mathcal{X} \to Z$ 具有性质
$\mathcal{P}$。
\end{proof}

\begin{definition}
\label{stacks-morphisms-definition-etale-smooth-P}
设 $\mathcal{P}$ 是代数空间态射的一个在源与靶上 \'etale--光滑局部的性质。
若引理 \ref{stacks-morphisms-lemma-local-source-target} 的等价条件成立，则称代数叠的 DM 态射
$f : \mathcal{X} \to \mathcal{Y}$ {\it 具有性质 $\mathcal{P}$}。
\end{definition}

\begin{remark}
\label{stacks-morphisms-remark-etale-smooth-composition}
设 $\mathcal{P}$ 是代数空间态射的一个在源与靶上 \'etale--光滑局部并且
在复合下稳定的性质。则定义 \ref{stacks-morphisms-definition-etale-smooth-P} 所定义的
代数叠 DM 态射的相应性质在复合下稳定。事实上，设
$f : \mathcal{X} \to \mathcal{Y}$ 和 $g : \mathcal{Y} \to \mathcal{Z}$
是具有性质 $\mathcal{P}$ 的代数叠 DM 态射。由引理
\ref{stacks-morphisms-lemma-composition-separated}，复合 $g \circ f$ 是 DM 的。选取代数空间
$W$ 以及满光滑态射 $W \to \mathcal{Z}$。选取代数空间 $V$ 以及满
\'etale 态射 $V \to \mathcal{Y} \times_\mathcal{Z} W$
（引理 \ref{stacks-morphisms-lemma-DM}）。再选取代数空间 $U$ 以及满 \'etale 态射
$U \to \mathcal{X} \times_\mathcal{Y} V$。按定义，态射
$V \to W$ 和 $U \to V$ 具有性质 $\mathcal{P}$。由于假设
$\mathcal{P}$ 在复合下稳定，故 $U \to W$ 具有性质 $\mathcal{P}$。
再次由定义可知，$g \circ f : \mathcal{X} \to \mathcal{Z}$ 具有性质
$\mathcal{P}$。
\end{remark}

\begin{remark}
\label{stacks-morphisms-remark-etale-smooth-base-change}
设 $\mathcal{P}$ 是代数空间态射的一个在源与靶上 \'etale--光滑局部并且
在基变换下稳定的性质。则定义 \ref{stacks-morphisms-definition-etale-smooth-P} 所定义的
代数叠 DM 态射的相应性质在任意基变换下稳定。事实上，设
$f : \mathcal{X} \to \mathcal{Y}$ 是代数叠的 DM 态射，而
$g : \mathcal{Y}' \to \mathcal{Y}$ 是代数叠的态射，并假设 $f$
具有性质 $\mathcal{P}$。则基变换
$\mathcal{Y}' \times_\mathcal{Y} \mathcal{X} \to \mathcal{Y}'$
由引理 \ref{stacks-morphisms-lemma-base-change-separated} 是 DM 态射。选取代数空间 $V$
以及满光滑态射 $V \to \mathcal{Y}$。选取代数空间 $U$ 以及满 \'etale 态射
$U \to \mathcal{X} \times_\mathcal{Y} V$（引理 \ref{stacks-morphisms-lemma-DM}）。
最后，选取代数空间 $V'$ 以及满光滑态射
$V' \to \mathcal{Y}' \times_\mathcal{Y} V$。按定义，态射
$U \to V$ 具有性质 $\mathcal{P}$。由于假设 $\mathcal{P}$ 在基变换下稳定，
$V' \times_V U \to V'$ 也具有性质 $\mathcal{P}$。考虑交换图
$$
\xymatrix{
V' \times_V U \ar[r] \ar[d] &
\mathcal{Y}' \times_\mathcal{Y} \mathcal{X} \ar[r] \ar[d] &
\mathcal{X} \ar[d] \\
V' \ar[r] & \mathcal{Y}' \ar[r] & \mathcal{Y}
}
$$
可知左上方水平箭头满，并且
$$
V' \times_V U \to V' \times_\mathcal{Y}
(\mathcal{Y}' \times_{\mathcal{Y}'} \mathcal{X}) =
V' \times_V (\mathcal{X} \times_\mathcal{Y} V)
$$
作为 $U \to \mathcal{X} \times_\mathcal{Y} V$ 的基变换是 \'etale 的。
所以由定义可知，投影
$\mathcal{Y}' \times_\mathcal{Y} \mathcal{X} \to \mathcal{Y}'$
具有性质 $\mathcal{P}$。
\end{remark}

\begin{remark}
\label{stacks-morphisms-remark-etale-smooth-implication}
设 $\mathcal{P}, \mathcal{P}'$ 是代数空间态射的两个在源与靶上
\'etale--光滑局部的性质。假设对代数空间的态射有
$\mathcal{P} \Rightarrow \mathcal{P}'$。则对利用 $\mathcal{P}$ 和
$\mathcal{P}'$ 在定义 \ref{stacks-morphisms-definition-etale-smooth-P} 中定义的代数叠态射的
性质，也有 $\mathcal{P} \Rightarrow \mathcal{P}'$。这由定义显然成立。
\end{remark}








\section{\'Etale 态射}
\label{stacks-morphisms-section-etale}

\noindent
不应把代数叠的 \'etale 态射定义为相对维数 $0$ 的光滑态射。事实上，态射
$$
[\mathbf{A}^1_k/\mathbf{G}_{m, k}] \longrightarrow \Spec(k)
$$
对群概形 $\mathbf{G}_{m, k}$ 在 $\mathbf{A}^1_k$ 上的任意作用都是相对维数
$0$ 的光滑态射。这与通常认为 \'etale 态射应当认同切空间的观念不符。
上例不是拟有限的，但态射
$$
\mathcal{X} = [\Spec(k)/\mu_{p, k}] \longrightarrow \Spec(k)
$$
光滑且拟有限（一节 \ref{stacks-morphisms-section-locally-quasi-finite}）。然而，若 $k$
的特征为 $p > 0$，则可表态射 $\Spec(k) \to \mathcal{X}$ 不是 \'etale 的，
因为其基变换
$\mu_{p, k} = \Spec(k) \times_\mathcal{X} \Spec(k) \to \Spec(k)$
是从非约化概形到域的谱的态射。因此，若把 \'etale 态射定义为光滑且局部
拟有限的态射，则《代数空间的态射》引理
\ref{spaces-morphisms-lemma-etale-permanence} 的类似结论将不成立。

\medskip\noindent
我们的做法则从下列要求出发：“\'etale 性”应当是在基变换下保持的性质；
并且若 $\mathcal{X} \to X$ 是从代数叠到概形的 \'etale 态射，则
$\mathcal{X}$ 应为 Deligne--Mumford 叠。换言之，我们要求 \'etale 态射为
DM 态射，并用一节 \ref{stacks-morphisms-section-etale-smooth-local-source-target} 的内容
定义代数叠的 \'etale 态射。

\medskip\noindent
在引理 \ref{stacks-morphisms-lemma-characterize-etale} 中，我们将把代数叠的 \'etale 态射刻画为
满足下列条件的态射：(a) 局部有限表现，(b) 平坦，且 (c) 对角态射为 \'etale。

\medskip\noindent
代数空间态射“为 \'etale”的性质在源与靶上 \'etale--光滑局部，见
《代数空间上的下降》评注
\ref{spaces-descent-remark-list-etale-smooth-local-source-target}。该性质在
基变换下也稳定，并且在靶上 fpqc 局部，见《代数空间的态射》引理
\ref{spaces-morphisms-lemma-base-change-etale} 和《代数空间上的下降》引理
\ref{spaces-descent-lemma-descending-property-etale}。因此，由上面的引理
\ref{stacks-morphisms-lemma-etale-smooth-local-source-target}，可以如下定义代数空间的态射何时为
\'etale；当该态射可由代数空间表示时，这一定义与《叠的性质》一节
\ref{stacks-properties-section-properties-morphisms} 中已有的概念一致，因为由引理
\ref{stacks-morphisms-lemma-trivial-implications}，这种态射自动是 DM 的。

\begin{definition}
\label{stacks-morphisms-definition-etale}
设 $f : \mathcal{X} \to \mathcal{Y}$ 是代数叠的态射。若 $f$ 是 DM 的，
并且引理 \ref{stacks-morphisms-lemma-etale-smooth-local-source-target} 的等价条件对
$\mathcal{P} = \etale$ 成立，则称 $f$ {\it 为 \'etale}。
\end{definition}

\noindent
下文将不再特别说明：当 $f$ 可由代数空间表示，或 $\mathcal{X}$ 和
$\mathcal{Y}$ 可由代数空间表示时，这一定义与已有的 \'etale 态射概念一致。

\begin{lemma}
\label{stacks-morphisms-lemma-composition-etale}
\'Etale 态射的复合仍为 \'etale。
\end{lemma}

\begin{proof}
结合评注 \ref{stacks-morphisms-remark-etale-smooth-composition} 与《代数空间的态射》引理
\ref{spaces-morphisms-lemma-composition-etale}。
\end{proof}

\begin{lemma}
\label{stacks-morphisms-lemma-base-change-etale}
\'Etale 态射的基变换仍为 \'etale。
\end{lemma}

\begin{proof}
结合评注 \ref{stacks-morphisms-remark-etale-smooth-base-change} 与《代数空间的态射》引理
\ref{spaces-morphisms-lemma-base-change-etale}。
\end{proof}

\begin{lemma}
\label{stacks-morphisms-lemma-open-immersion-etale}
开浸入是 \'etale 的。
\end{lemma}

\begin{proof}
设 $j : \mathcal{U} \to \mathcal{X}$ 是代数叠的开浸入。由于 $j$ 可表，
由引理 \ref{stacks-morphisms-lemma-trivial-implications}，它是 DM 的。另一方面，若
$X \to \mathcal{X}$ 是满光滑态射，其中 $X$ 是概形，则
$U = \mathcal{U} \times_\mathcal{X} X$ 是 $X$ 的开子概形。因此
$U \to X$ 是 \'etale 的（《态射》引理
\ref{morphisms-lemma-open-immersion-etale}），由定义可知 $j$ 是 \'etale 的。
\end{proof}

\begin{lemma}
\label{stacks-morphisms-lemma-etale}
设 $f : \mathcal{X} \to \mathcal{Y}$ 是代数叠的态射。下列条件等价：
\begin{enumerate}
\item $f$ 是 \'etale 的；
\item $f$ 是 DM 的，并且对任意态射 $V \to \mathcal{Y}$（其中 $V$
是代数空间）以及任意 \'etale 态射
$U \to V \times_\mathcal{Y} \mathcal{X}$（其中 $U$ 是代数空间），
态射 $U \to V$ 都是 \'etale 的；
\item 存在满、局部有限表现、平坦的态射 $W \to \mathcal{Y}$（其中 $W$
是代数空间）以及满 \'etale 态射
$T \to W \times_\mathcal{Y} \mathcal{X}$（其中 $T$ 是代数空间），
使得态射 $T \to W$ 是 \'etale 的。
\end{enumerate}
\end{lemma}

\begin{proof}
假设 (1)。则 $f$ 是 DM 的；由于 \'etale 性在基变换下保持，(2) 成立。

\medskip\noindent
假设 (2)。选取概形 $V$ 以及满 \'etale 态射 $V \to \mathcal{Y}$。
选取概形 $U$ 以及满 \'etale 态射
$U \to V \times_\mathcal{Y} \mathcal{X}$（引理 \ref{stacks-morphisms-lemma-DM}）。
于是 (3) 成立。

\medskip\noindent
假设 $W \to \mathcal{Y}$ 和
$T \to W \times_\mathcal{Y} \mathcal{X}$ 如 (3) 所述。先验证 $f$ 是 DM 的。
只需验证 $W \times_\mathcal{Y} \mathcal{X} \to W$ 是 DM 的，见引理
\ref{stacks-morphisms-lemma-check-separated-covering}。由引理
\ref{stacks-morphisms-lemma-compose-after-separated}，只需验证
$W \times_\mathcal{Y} \mathcal{X}$ 是 DM 叠。这由态射
$T \to W \times_\mathcal{Y} \mathcal{X}$ 的存在性以及定理
\ref{stacks-morphisms-theorem-DM} 的容易方向得出。

\medskip\noindent
假设 $f$ 是 DM 的，且 $W \to \mathcal{Y}$ 和
$T \to W \times_\mathcal{Y} \mathcal{X}$ 如 (3) 所述。设 $V$
是代数空间，$V \to \mathcal{Y}$ 满且光滑；设 $U$ 是代数空间，并令
$U \to V \times_\mathcal{Y} \mathcal{X}$ 满且 \'etale
（引理 \ref{stacks-morphisms-lemma-DM}）。必须验证 $U \to V$ 是 \'etale 的。由
《代数空间上的下降》引理
\ref{spaces-descent-lemma-descending-property-etale}，只需证明
$U \times_\mathcal{Y} W \to V \times_\mathcal{Y} W$ 是 \'etale 的。
可以依次用
$\mathcal{X} \times_\mathcal{Y} W, W, W, T, U \times_\mathcal{Y} W,
V \times_\mathcal{Y} W$ 替换
$\mathcal{X}, \mathcal{Y}, W, T, U, V$（略去一个小细节）。因此可以假设
$Y = \mathcal{Y}$ 是代数空间，存在代数空间 $T$ 以及满 \'etale 态射
$T \to \mathcal{X}$，使得 $T \to Y$ 是 \'etale 的，而 $U$ 和 $V$
仍如上所述。在此情形，有
$$
U \to V\text{ is \'etale}
\Leftrightarrow
\mathcal{X} \to Y\text{ is \'etale}
\Leftrightarrow
T \to Y\text{ is \'etale}
$$
，这是引理 \ref{stacks-morphisms-lemma-etale-smooth-local-source-target} 的性质 (1) 与 (2)
的等价性以及定义 \ref{stacks-morphisms-definition-etale} 所给出的。证明完成。
\end{proof}

\begin{lemma}
\label{stacks-morphisms-lemma-etale-permanence}
设 $\mathcal{X}, \mathcal{Y}$ 是在代数叠 $\mathcal{Z}$ 上 \'etale 的
代数叠。任意 $\mathcal{Z}$ 上的态射
$\mathcal{X} \to \mathcal{Y}$ 都是 \'etale 的。
\end{lemma}

\begin{proof}
由引理 \ref{stacks-morphisms-lemma-compose-after-separated}，态射
$\mathcal{X} \to \mathcal{Y}$ 是 DM 的。设 $W \to \mathcal{Z}$
是源为代数空间的满光滑态射。设
$V \to \mathcal{Y} \times_\mathcal{Z} W$ 是源为代数空间的满
\'etale 态射（引理 \ref{stacks-morphisms-lemma-DM}）。设
$U \to \mathcal{X} \times_\mathcal{Y} V$ 是源为代数空间的满
\'etale 态射（引理 \ref{stacks-morphisms-lemma-DM}）。则
$$
U \longrightarrow \mathcal{X} \times_\mathcal{Z} W
$$
是满 \'etale 态射，因为它是
$U \to \mathcal{X} \times_\mathcal{Y} V$ 与
$V \to \mathcal{Y} \times_\mathcal{Z} W$ 沿
$\mathcal{X} \times_\mathcal{Z} W \to \mathcal{Y} \times_\mathcal{Z} W$
所得基变换的复合。因此，只需证明 $U \to W$ 是 \'etale 的。由于
$\mathcal{X} \to \mathcal{Z}$ 和 $\mathcal{Y} \to \mathcal{Z}$
是 \'etale 的，$U \to W$ 和 $V \to W$ 也都是 \'etale 的；所需结论由
代数空间版本的引理，即《代数空间的态射》引理
\ref{spaces-morphisms-lemma-etale-permanence} 得出。
\end{proof}








\section{非分歧态射}
\label{stacks-morphisms-section-unramified}

\noindent
关于我们为何选择这个非分歧态射的定义，请读者参阅 \'etale 态射一节
\ref{stacks-morphisms-section-etale} 中的讨论。

\medskip\noindent
在引理 \ref{stacks-morphisms-lemma-characterize-unramified} 中，我们将把代数叠的非分歧态射
刻画为局部有限型且对角态射为 \'etale 的态射。

\medskip\noindent
代数空间态射“为非分歧”的性质在源与靶上 \'etale--光滑局部，见
《代数空间上的下降》评注
\ref{spaces-descent-remark-list-etale-smooth-local-source-target}。该性质在基变换下
也稳定，并且在靶上 fpqc 局部，见《代数空间的态射》引理
\ref{spaces-morphisms-lemma-base-change-unramified} 和《代数空间上的下降》引理
\ref{spaces-descent-lemma-descending-property-unramified}。因此，由上面的引理
\ref{stacks-morphisms-lemma-etale-smooth-local-source-target}，可以如下定义代数空间的态射何时
非分歧；当该态射可由代数空间表示时，这一定义与《叠的性质》一节
\ref{stacks-properties-section-properties-morphisms} 中已有的概念一致，因为由引理
\ref{stacks-morphisms-lemma-trivial-implications}，这种态射自动是 DM 的。

\begin{definition}
\label{stacks-morphisms-definition-unramified}
设 $f : \mathcal{X} \to \mathcal{Y}$ 是代数叠的态射。若 $f$ 是 DM 的，
并且引理 \ref{stacks-morphisms-lemma-etale-smooth-local-source-target} 的等价条件对
$\mathcal{P} =$“非分歧”成立，则称 $f$ {\it 非分歧}。
\end{definition}

\noindent
下文将不再特别说明：当 $f$ 可由代数空间表示，或 $\mathcal{X}$ 和
$\mathcal{Y}$ 可由代数空间表示时，这一定义与已有的非分歧态射概念一致。

\begin{lemma}
\label{stacks-morphisms-lemma-composition-unramified}
非分歧态射的复合仍非分歧。
\end{lemma}

\begin{proof}
结合评注 \ref{stacks-morphisms-remark-etale-smooth-composition} 与《代数空间的态射》引理
\ref{spaces-morphisms-lemma-composition-unramified}。
\end{proof}

\begin{lemma}
\label{stacks-morphisms-lemma-base-change-unramified}
非分歧态射的基变换仍非分歧。
\end{lemma}

\begin{proof}
结合评注 \ref{stacks-morphisms-remark-etale-smooth-base-change} 与《代数空间的态射》引理
\ref{spaces-morphisms-lemma-base-change-unramified}。
\end{proof}

\begin{lemma}
\label{stacks-morphisms-lemma-etale-unramified}
\'Etale 态射是非分歧的。
\end{lemma}

\begin{proof}
由评注 \ref{stacks-morphisms-remark-etale-smooth-implication} 和《代数空间的态射》引理
\ref{spaces-morphisms-lemma-etale-unramified} 得出。
\end{proof}

\begin{lemma}
\label{stacks-morphisms-lemma-immersion-unramified}
浸入是非分歧的。
\end{lemma}

\begin{proof}
设 $j : \mathcal{Z} \to \mathcal{X}$ 是代数叠的浸入。由于 $j$ 可表，
由引理 \ref{stacks-morphisms-lemma-trivial-implications}，它是 DM 的。另一方面，若
$X \to \mathcal{X}$ 是满光滑态射，其中 $X$ 是概形，则
$Z = \mathcal{Z} \times_\mathcal{X} X$ 是 $X$ 的局部闭子概形。因此
$Z \to X$ 非分歧（《态射》诸引理
\ref{morphisms-lemma-open-immersion-unramified} 和
\ref{morphisms-lemma-closed-immersion-unramified}），由定义可知 $j$ 非分歧。
\end{proof}

\begin{lemma}
\label{stacks-morphisms-lemma-unramified}
设 $f : \mathcal{X} \to \mathcal{Y}$ 是代数叠的态射。下列条件等价：
\begin{enumerate}
\item $f$ 非分歧；
\item $f$ 是 DM 的，并且对任意态射 $V \to \mathcal{Y}$（其中 $V$
是代数空间）以及任意 \'etale 态射
$U \to V \times_\mathcal{Y} \mathcal{X}$（其中 $U$ 是代数空间），
态射 $U \to V$ 都非分歧；
\item 存在满、局部有限表现、平坦的态射 $W \to \mathcal{Y}$（其中 $W$
是代数空间）以及满 \'etale 态射
$T \to W \times_\mathcal{Y} \mathcal{X}$（其中 $T$ 是代数空间），
使得态射 $T \to W$ 非分歧。
\end{enumerate}
\end{lemma}

\begin{proof}
假设 (1)。则 $f$ 是 DM 的；由于非分歧性在基变换下保持，(2) 成立。

\medskip\noindent
假设 (2)。选取概形 $V$ 以及满 \'etale 态射 $V \to \mathcal{Y}$。
选取概形 $U$ 以及满 \'etale 态射
$U \to V \times_\mathcal{Y} \mathcal{X}$（引理 \ref{stacks-morphisms-lemma-DM}）。
于是 (3) 成立。

\medskip\noindent
假设 $W \to \mathcal{Y}$ 和
$T \to W \times_\mathcal{Y} \mathcal{X}$ 如 (3) 所述。先验证 $f$ 是 DM 的。
只需验证 $W \times_\mathcal{Y} \mathcal{X} \to W$ 是 DM 的，见引理
\ref{stacks-morphisms-lemma-check-separated-covering}。由引理
\ref{stacks-morphisms-lemma-compose-after-separated}，只需验证
$W \times_\mathcal{Y} \mathcal{X}$ 是 DM 叠。这由态射
$T \to W \times_\mathcal{Y} \mathcal{X}$ 的存在性以及定理
\ref{stacks-morphisms-theorem-DM} 的容易方向得出。

\medskip\noindent
假设 $f$ 是 DM 的，且 $W \to \mathcal{Y}$ 和
$T \to W \times_\mathcal{Y} \mathcal{X}$ 如 (3) 所述。设 $V$
是代数空间，$V \to \mathcal{Y}$ 满且光滑；设 $U$ 是代数空间，并令
$U \to V \times_\mathcal{Y} \mathcal{X}$ 满且 \'etale
（引理 \ref{stacks-morphisms-lemma-DM}）。必须验证 $U \to V$ 非分歧。由
《代数空间上的下降》引理
\ref{spaces-descent-lemma-descending-property-unramified}，只需证明
$U \times_\mathcal{Y} W \to V \times_\mathcal{Y} W$ 非分歧。可以依次用
$\mathcal{X} \times_\mathcal{Y} W, W, W, T, U \times_\mathcal{Y} W,
V \times_\mathcal{Y} W$ 替换
$\mathcal{X}, \mathcal{Y}, W, T, U, V$（略去一个小细节）。因此可以假设
$Y = \mathcal{Y}$ 是代数空间，存在代数空间 $T$ 以及满 \'etale 态射
$T \to \mathcal{X}$，使得 $T \to Y$ 非分歧，而 $U$ 和 $V$
仍如上所述。在此情形，有
$$
U \to V\text{ is unramified}
\Leftrightarrow
\mathcal{X} \to Y\text{ is unramified}
\Leftrightarrow
T \to Y\text{ is unramified}
$$
，这是引理 \ref{stacks-morphisms-lemma-etale-smooth-local-source-target} 的性质 (1) 与 (2)
的等价性以及定义 \ref{stacks-morphisms-definition-unramified} 所给出的。证明完成。
\end{proof}

\begin{lemma}
\label{stacks-morphisms-lemma-unramified-quasi-finite}
代数叠的非分歧态射局部拟有限。
\end{lemma}

\begin{proof}
这由引理 \ref{stacks-morphisms-lemma-unramified}（非分歧态射的刻画）、引理
\ref{stacks-morphisms-lemma-characterize-locally-quasi-finite}（局部拟有限态射的刻画），
以及《代数空间的态射》引理
\ref{spaces-morphisms-lemma-unramified-quasi-finite}（代数空间的相应结果）得出。
\end{proof}

\begin{lemma}
\label{stacks-morphisms-lemma-permanence-unramified}
设 $\mathcal{X} \to \mathcal{Y} \to \mathcal{Z}$ 是代数叠的态射。
若 $\mathcal{X} \to \mathcal{Z}$ 非分歧，且
$\mathcal{Y} \to \mathcal{Z}$ 是 DM 的，则
$\mathcal{X} \to \mathcal{Y}$ 非分歧。
\end{lemma}

\begin{proof}
假设 $\mathcal{X} \to \mathcal{Z}$ 非分歧。由引理
\ref{stacks-morphisms-lemma-compose-after-separated}，态射
$\mathcal{X} \to \mathcal{Y}$ 是 DM 的。选取交换图
$$
\xymatrix{
U \ar[d] \ar[r] & V \ar[d] \ar[r] & W \ar[d] \\
\mathcal{X} \ar[r] & \mathcal{Y} \ar[r] & \mathcal{Z}
}
$$
，其中 $U, V, W$ 是代数空间，$W \to \mathcal{Z}$ 满且光滑，
$V \to \mathcal{Y} \times_\mathcal{Z} W$ 满且 \'etale，而
$U \to \mathcal{X} \times_\mathcal{Y} V$ 满且 \'etale
（见引理 \ref{stacks-morphisms-lemma-DM}）。于是
$U \to \mathcal{X} \times_\mathcal{Z} W$ 也满且 \'etale。因此已知
$U \to W$ 非分歧，而必须证明 $U \to V$ 非分歧。这由《代数空间的态射》引理
\ref{spaces-morphisms-lemma-permanence-unramified} 得出。
\end{proof}

\begin{lemma}
\label{stacks-morphisms-lemma-characterize-unramified}
设 $f : \mathcal{X} \to \mathcal{Y}$ 是代数叠的态射。下列条件等价：
\begin{enumerate}
\item $f$ 非分歧；
\item $f$ 局部有限型，且其对角态射为 \'etale。
\end{enumerate}
\end{lemma}

\begin{proof}
假设 $f$ 非分歧。则 $f$ 是 DM 的，因而可以选取代数空间 $U$、$V$，
满光滑态射 $V \to \mathcal{Y}$，以及满 \'etale 态射
$U \to \mathcal{X} \times_\mathcal{Y} V$（引理 \ref{stacks-morphisms-lemma-DM}）。
由于 $f$ 非分歧，诱导态射 $U \to V$ 非分歧。因此 $U \to V$ 局部有限型
（《代数空间的态射》引理
\ref{spaces-morphisms-lemma-unramified-locally-finite-type}），从而 $f$
局部有限型。对角态射
$\Delta : \mathcal{X} \to \mathcal{X} \times_\mathcal{Y} \mathcal{X}$
是 $\mathcal{Y}$ 上代数叠的态射。$\Delta$ 沿满光滑态射
$V \to \mathcal{Y}$ 的基变换，是 $f$ 的基变换，即
$\mathcal{X}_V = \mathcal{X} \times_\mathcal{Y} V \to V$ 的对角态射。
换言之，图
$$
\xymatrix{
\mathcal{X}_V \ar[r] \ar[d] &  \mathcal{X}_V \times_V \mathcal{X}_V \ar[d] \\
\mathcal{X} \ar[r] & \mathcal{X} \times_\mathcal{Y} \mathcal{X}
}
$$
是笛卡儿的。由于右侧竖直箭头满且光滑，由《叠的性质》引理
\ref{stacks-properties-lemma-check-property-weak-covering}，只需证明顶部
水平箭头为 \'etale。考虑交换图
$$
\xymatrix{
U \ar[d] \ar[r] & U \times_V U \ar[d] \\
\mathcal{X}_V \ar[r] & \mathcal{X}_V \times_V \mathcal{X}_V
}
$$
其中所有箭头都可由代数空间表示，竖直箭头为 \'etale，左侧竖直箭头满，
而由《代数空间的态射》引理
\ref{spaces-morphisms-lemma-diagonal-unramified-morphism}，顶部水平箭头是开浸入。
这蕴含所需结论：首先，作为 \'etale 态射的复合，
$U \to \mathcal{X}_V \times_V \mathcal{X}_V$ 是 \'etale 的；然后应用
《叠的性质》引理
\ref{stacks-properties-lemma-check-property-after-precomposing} 可知
$\mathcal{X}_V \to \mathcal{X}_V \times_V \mathcal{X}_V$ 是 \'etale 的，
因为对代数空间的态射而言，\'etale 性在源上对 \'etale 拓扑局部
（《代数空间上的下降》引理
\ref{spaces-descent-lemma-etale-etale-local-source}）。

\medskip\noindent
假设 $f$ 局部有限型且其对角态射为 \'etale。由于代数空间的 \'etale 态射
非分歧，按定义 $f$ 是 DM 的。因此，像上文一样，可以选取代数空间
$U$、$V$，满光滑态射 $V \to \mathcal{Y}$，以及满 \'etale 态射
$U \to \mathcal{X} \times_\mathcal{Y} V$（引理 \ref{stacks-morphisms-lemma-DM}）。
为完成证明，必须证明 $U \to V$ 非分歧。已经知道 $U \to V$ 局部有限型。
像上文一样，得到交换图
$$
\xymatrix{
U \ar[d] \ar[r] & U \times_V U \ar[d] \\
\mathcal{X}_V \ar[r] & \mathcal{X}_V \times_V \mathcal{X}_V
}
$$
，其中所有箭头都可由代数空间表示，竖直箭头为 \'etale，而底部水平箭头作为
$\Delta$ 的基变换也是 \'etale 的。例如由引理
\ref{stacks-morphisms-lemma-etale-permanence}\footnote{也很容易直接从《代数空间的态射》引理
\ref{spaces-morphisms-lemma-etale-permanence} 推出这一点。}，可知
$U \to U \times_V U$ 是 \'etale 的。因此，
$U \to U \times_V U$ 是 \'etale 单态射，故为开浸入
（《代数空间的态射》引理
\ref{spaces-morphisms-lemma-etale-universally-injective-open}）。于是由
《代数空间的态射》引理
\ref{spaces-morphisms-lemma-diagonal-unramified-morphism}，
$U \to V$ 非分歧。
\end{proof}

\begin{lemma}
\label{stacks-morphisms-lemma-characterize-etale}
设 $f : \mathcal{X} \to \mathcal{Y}$ 是代数叠的态射。下列条件等价：
\begin{enumerate}
\item $f$ 是 \'etale 的；
\item $f$ 局部有限表现、平坦且非分歧；
\item $f$ 局部有限表现、平坦，且其对角态射为 \'etale。
\end{enumerate}
\end{lemma}

\begin{proof}
(2) 与 (3) 的等价性立即由引理 \ref{stacks-morphisms-lemma-characterize-unramified} 得出。
所以在每种情形下，态射 $f$ 都是 DM 的。于是可以选取代数空间 $U$、$V$，
满光滑态射 $V \to \mathcal{Y}$，以及满 \'etale 态射
$U \to \mathcal{X} \times_\mathcal{Y} V$（引理 \ref{stacks-morphisms-lemma-DM}）。
为完成证明，必须证明 $U \to V$ 是 \'etale 的，当且仅当它局部有限表现、
平坦且非分歧。这由《代数空间的态射》引理
\ref{spaces-morphisms-lemma-unramified-flat-lfp-etale} 得出
（以及更为平凡的《代数空间的态射》诸引理
\ref{spaces-morphisms-lemma-etale-unramified}、
\ref{spaces-morphisms-lemma-etale-locally-finite-presentation} 和
\ref{spaces-morphisms-lemma-etale-flat}）。
\end{proof}








\section{固有态射}
\label{stacks-morphisms-section-proper}

\noindent
固有态射的概念在代数几何中起着重要作用。下面定义代数叠的固有态射。

\begin{definition}
\label{stacks-morphisms-definition-proper}
设 $f : \mathcal{X} \to \mathcal{Y}$ 是代数叠的态射。若 $f$ 分离、
有限型且泛闭，则称 $f$ {\it 固有}。
\end{definition}

\noindent
这与已有的代数空间固有态射概念并不冲突：代数空间的态射固有，当且仅当它
分离、有限型且泛闭（《代数空间的态射》定义
\ref{spaces-morphisms-definition-proper}）；而我们已经在引理
\ref{stacks-morphisms-lemma-representable-separated-diagonal-closed}、一节
\ref{stacks-morphisms-section-finite-type} 和引理
\ref{stacks-morphisms-lemma-characterize-representable-universally-closed} 中验证了这些概念的
相容性。类似地，若 $f : \mathcal{X} \to \mathcal{Y}$ 是可由代数空间表示的
代数叠态射，则《叠的性质》一节
\ref{stacks-properties-section-properties-morphisms} 已定义 $f$ 何时固有。
然而，该节的讨论说明，这等价于要求 $f$ 分离、有限型且泛闭，而上引同一批结果
给出了相容性。

\begin{lemma}
\label{stacks-morphisms-lemma-base-change-proper}
固有态射的基变换仍固有。
\end{lemma}

\begin{proof}
见引理 \ref{stacks-morphisms-lemma-base-change-separated}、
\ref{stacks-morphisms-lemma-base-change-finite-type} 和
\ref{stacks-morphisms-lemma-base-change-universally-closed}。
\end{proof}

\begin{lemma}
\label{stacks-morphisms-lemma-composition-proper}
固有态射的复合仍固有。
\end{lemma}

\begin{proof}
见引理 \ref{stacks-morphisms-lemma-composition-separated}、
\ref{stacks-morphisms-lemma-composition-finite-type} 和
\ref{stacks-morphisms-lemma-composition-universally-closed}。
\end{proof}

\begin{lemma}
\label{stacks-morphisms-lemma-closed-immersion-proper}
代数叠的闭浸入是代数叠的固有态射。
\end{lemma}

\begin{proof}
按定义，闭浸入可表（《叠的性质》定义
\ref{stacks-properties-definition-immersion}）。因此，结论由《叠的性质》一节
\ref{stacks-properties-section-properties-morphisms} 的讨论以及代数空间态射的
相应结果得出，见《代数空间的态射》引理
\ref{spaces-morphisms-lemma-closed-immersion-proper}。
\end{proof}

\begin{lemma}
\label{stacks-morphisms-lemma-universally-closed-permanence}
考虑代数叠的交换图
$$
\xymatrix{
\mathcal{X} \ar[rr] \ar[rd] & &
\mathcal{Y} \ar[ld] \\
& \mathcal{Z} &
}
$$
\begin{enumerate}
\item 若 $\mathcal{X} \to \mathcal{Z}$ 泛闭，而
$\mathcal{Y} \to \mathcal{Z}$ 分离，则态射
$\mathcal{X} \to \mathcal{Y}$ 泛闭。特别地，$|\mathcal{X}|$
在 $|\mathcal{Y}|$ 中的像是闭的。
\item 若 $\mathcal{X} \to \mathcal{Z}$ 固有，而
$\mathcal{Y} \to \mathcal{Z}$ 分离，则态射
$\mathcal{X} \to \mathcal{Y}$ 固有。
\end{enumerate}
\end{lemma}

\begin{proof}
假设 $\mathcal{X} \to \mathcal{Z}$ 泛闭，且
$\mathcal{Y} \to \mathcal{Z}$ 分离。把该态射分解为
$\mathcal{X} \to \mathcal{X} \times_\mathcal{Z} \mathcal{Y} \to \mathcal{Y}$.
第一个态射固有（引理 \ref{stacks-morphisms-lemma-semi-diagonal}），因而泛闭。投影
$\mathcal{X} \times_\mathcal{Z} \mathcal{Y} \to \mathcal{Y}$
是泛闭态射的基变换，故泛闭，见引理
\ref{stacks-morphisms-lemma-base-change-universally-closed}。所以，作为泛闭态射的复合，
$\mathcal{X} \to \mathcal{Y}$ 泛闭，见引理
\ref{stacks-morphisms-lemma-composition-universally-closed}。这证明了 (1)。为推出 (2)，
把 (1) 与引理 \ref{stacks-morphisms-lemma-compose-after-separated}、
\ref{stacks-morphisms-lemma-quasi-compact-permanence} 和
\ref{stacks-morphisms-lemma-finite-type-permanence} 结合即可。
\end{proof}

\begin{lemma}
\label{stacks-morphisms-lemma-image-proper-is-proper}
设 $\mathcal{Z}$ 是代数叠。设 $f : \mathcal{X} \to \mathcal{Y}$
是 $\mathcal{Z}$ 上代数叠的态射。若 $\mathcal{X}$ 在 $\mathcal{Z}$
上泛闭，且 $f$ 满，则 $\mathcal{Y}$ 在 $\mathcal{Z}$ 上泛闭。
特别地，若 $\mathcal{Y}$ 还在 $\mathcal{Z}$ 上分离且有限型，则
$\mathcal{Y}$ 在 $\mathcal{Z}$ 上固有。
\end{lemma}

\begin{proof}
假设 $\mathcal{X}$ 泛闭且 $f$ 满。记结构态射为
$p : \mathcal{X} \to \mathcal{Z}$ 和
$q : \mathcal{Y} \to \mathcal{Z}$。设
$\mathcal{Z}' \to \mathcal{Z}$ 是代数叠的态射。$f$ 沿
$\mathcal{Z}' \to \mathcal{Z}$ 的基变换
$f' : \mathcal{X}' \to \mathcal{Y}'$ 满
（《叠的性质》引理
\ref{stacks-properties-lemma-base-change-surjective}），而 $p$ 的基变换
$p' : \mathcal{X}' \to \mathcal{Z}'$ 是闭映射。若
$T \subset |\mathcal{Y}'|$ 闭，则
$(f')^{-1}(T) \subset |\mathcal{X}'|$ 闭，因而
$p'((f')^{-1}(T)) = q'(T)$ 闭。所以 $q'$ 是闭映射。
\end{proof}








\section{概形论像}
\label{stacks-morphisms-section-scheme-theoretic-image}

\noindent
定义如下。

\begin{definition}
\label{stacks-morphisms-definition-scheme-theoretic-image}
设 $f : \mathcal{X} \to \mathcal{Y}$ 是代数叠的态射。$f$ 的
{\it 概形论像}是使 $f$ 经由其分解的最小闭子叠
$\mathcal{Z} \subset \mathcal{Y}$\footnote{我们将在引理
\ref{stacks-morphisms-lemma-scheme-theoretic-image-existence} 中看到，概形论像总是存在。}。
\end{definition}

\noindent
我们常也用 $f : \mathcal{X} \to \mathcal{Z}$ 表示 $f$ 的这一分解。
若态射 $f$ 不拟紧，则概形论像的构造通常不与到 $\mathcal{Y}$ 的开子叠的
限制交换。然而，若 $f$ 拟紧，则概形论像与平坦基变换交换
（引理 \ref{stacks-morphisms-lemma-existence-plus-flat-base-change}）。

\begin{lemma}
\label{stacks-morphisms-lemma-cover-upstairs}
设 $f : \mathcal{X} \to \mathcal{Y}$ 是代数叠的态射。设
$g : \mathcal{W} \to \mathcal{X}$ 是代数叠的满、平坦、局部有限表现态射。
则 $f$ 的概形论像存在，当且仅当 $f \circ g$ 的概形论像存在；若存在，
二者相同。
\end{lemma}

\begin{proof}
假设 $\mathcal{Z} \subset \mathcal{Y}$ 是闭子叠，且 $f \circ g$
经由 $\mathcal{Z}$ 分解。为证明本引理，只需证明 $f$ 经由
$\mathcal{Z}$ 分解。考虑概形 $T$ 以及态射 $T \to \mathcal{X}$，
它由 $\mathcal{X}$ 在 $T$ 上的纤维范畴中的对象 $x$ 给出。我们要证明，
$f(x)$ 事实上属于 $\mathcal{Z}$ 在 $T$ 上的纤维范畴。投影
$T \times_\mathcal{X} \mathcal{W} \to T$ 满、平坦且局部有限表现。
因此存在 fppf 覆盖 $\{T_i \to T\}$，使得对所有 $i$，$T_i \to T$
都经由 $T \times_\mathcal{X} \mathcal{W} \to T$ 分解。于是
$T_i \to \mathcal{X}$ 经由 $\mathcal{W}$ 分解，进而
$T_i \to \mathcal{Y}$ 经由 $\mathcal{Z}$ 分解。所以
$f(x)|_{T_i}$ 是 $\mathcal{Z}$ 的对象。由于 $\mathcal{Z}$ 是严格满子叠，
可知 $f(x)$ 是 $\mathcal{Z}$ 的对象，正合所需。
\end{proof}

\begin{lemma}
\label{stacks-morphisms-lemma-scheme-theoretic-image-existence}
设 $f : \mathcal{Y} \to \mathcal{X}$ 是代数叠的态射。则 $f$
的概形论像存在。
\end{lemma}

\begin{proof}
选取概形 $V$ 以及满光滑态射 $V \to \mathcal{Y}$。由引理
\ref{stacks-morphisms-lemma-cover-upstairs}，可以用 $V$ 替换 $\mathcal{Y}$。因此，只需证明：
若 $X \to \mathcal{X}$ 是从概形到代数叠的态射，则其概形论像存在。
选取概形 $U$ 以及满光滑态射 $U \to \mathcal{X}$。令
$R = U \times_\mathcal{X} U$。由《代数叠》引理
\ref{algebraic-lemma-stack-presentation}，有 $\mathcal{X} = [U/R]$。
由《叠的性质》引理
\ref{stacks-properties-lemma-substacks-presentation}，$\mathcal{X}$
的闭子叠 $\mathcal{Z}$ 与 $R$-不变闭子概形 $Z \subset U$ 形成
$1$ 对 $1$ 的对应。
设 $Z_1 \subset U$ 是 $X \times_\mathcal{X} U \to U$ 的概形论像。
注意，$X \to \mathcal{X}$ 经由 $\mathcal{Z}$ 分解，当且仅当
$X \times_\mathcal{X} U \to U$ 经由相应的 $R$-不变闭子概形 $Z$ 分解
（细节略；提示：这是因为 $X \times_\mathcal{X} U \to X$ 满且光滑）。
因此，必须证明存在包含 $Z_1$ 的最小 $R$-不变闭子概形 $Z \subset U$。

\medskip\noindent
设 $\mathcal{I}_1 \subset \mathcal{O}_U$ 是与闭子概形
$Z_1 \subset U$ 对应的拟凝聚理想层。设 $Z_\alpha$（$\alpha \in A$）
为 $U$ 的所有包含 $Z_1$ 的 $R$-不变闭子概形所成的集合。对 $\alpha \in A$，
设 $\mathcal{I}_\alpha \subset \mathcal{O}_U$ 是与闭子概形
$Z_\alpha \subset U$ 对应的拟凝聚理想层。包含关系
$Z_1 \subset Z_\alpha$ 意味着
$\mathcal{I}_\alpha \subset \mathcal{I}_1$。$Z_\alpha$ 的 $R$-不变性意味着
$$
s^{-1}\mathcal{I}_\alpha \cdot \mathcal{O}_R =
t^{-1}\mathcal{I}_\alpha \cdot \mathcal{O}_R
$$
作为（代数空间）$R$ 上的（拟凝聚）理想层相等。考虑像
$$
\mathcal{I} =
\Im\left(
\bigoplus\nolimits_{\alpha \in A} \mathcal{I}_\alpha \to \mathcal{I}_1
\right) =
\Im\left(
\bigoplus\nolimits_{\alpha \in A} \mathcal{I}_\alpha \to \mathcal{O}_X
\right)
$$
由于拟凝聚层的直和仍拟凝聚，拟凝聚层之间映射的像也拟凝聚，可知
$\mathcal{I}$ 拟凝聚。由于拉回是正合的且与直和交换，可知
$$
s^{-1}\mathcal{I} \cdot \mathcal{O}_R =
t^{-1}\mathcal{I} \cdot \mathcal{O}_R
$$
因此，$\mathcal{I}$ 定义一个 $R$-不变闭子概形 $Z \subset U$，
它包含于每个 $Z_\alpha$ 且包含 $Z_1$，正合所需。
\end{proof}

\begin{lemma}
\label{stacks-morphisms-lemma-factor-factor}
设
$$
\xymatrix{
\mathcal{X}_1 \ar[d] \ar[r]_{f_1} & \mathcal{Y}_1 \ar[d] \\
\mathcal{X}_2 \ar[r]^{f_2} & \mathcal{Y}_2
}
$$
是代数叠的交换图。设 $\mathcal{Z}_i \subset \mathcal{Y}_i$（$i = 1, 2$）
是 $f_i$ 的概形论像。则态射
$\mathcal{Y}_1 \to \mathcal{Y}_2$ 诱导态射
$\mathcal{Z}_1 \to \mathcal{Z}_2$ 以及交换图
$$
\xymatrix{
\mathcal{X}_1 \ar[r] \ar[d] &
\mathcal{Z}_1 \ar[d] \ar[r] &
\mathcal{Y}_1 \ar[d] \\
\mathcal{X}_2 \ar[r] &
\mathcal{Z}_2 \ar[r] &
\mathcal{Y}_2
}
$$
\end{lemma}

\begin{proof}
$\mathcal{Z}_2$ 在 $\mathcal{Y}_1$ 中的概形论逆像是
$\mathcal{Y}_1$ 的闭子叠，$f_1$ 经由它分解。因此
$\mathcal{Z}_1$ 包含于其中。本引理得证。
\end{proof}

\begin{lemma}
\label{stacks-morphisms-lemma-existence-plus-flat-base-change}
设 $f : \mathcal{X} \to \mathcal{Y}$ 是代数叠的拟紧态射。
则概形论像的构造与平坦基变换交换。
\end{lemma}

\begin{proof}
设 $\mathcal{Y}' \to \mathcal{Y}$ 是代数叠的平坦态射。选取概形 $V$
以及满光滑态射 $V \to \mathcal{Y}$。选取概形 $V'$ 以及满光滑态射
$V' \to \mathcal{Y}' \times_\mathcal{Y} V$。可以并且确实假设
$V = \coprod_{i \in I} V_i$ 是仿射概形的不交并，并且
$V' = \coprod_{i \in I} \coprod_{j \in J_i} V_{i, j}$
是仿射概形的不交并，其中每个 $V_{i, j}$ 都映入 $V_i$。设
\begin{enumerate}
\item $\mathcal{Z} \subset \mathcal{Y}$ 是 $f$ 的概形论像；
\item $\mathcal{Z}' \subset \mathcal{Y}'$ 是 $f$ 沿
$\mathcal{Y}' \to \mathcal{Y}$ 所得基变换的概形论像；
\item $Z \subset V$ 是 $f$ 沿 $V \to \mathcal{Y}$ 所得基变换的概形论像；
\item $Z' \subset V'$ 是 $f$ 沿 $V' \to \mathcal{Y}$ 所得基变换的概形论像。
\end{enumerate}
若能证明
(a) $Z = V \times_\mathcal{Y} \mathcal{Z}$,
(b) $Z' = V' \times_{\mathcal{Y}'} \mathcal{Z}'$，以及
(c) $Z' = V' \times_V Z$
，则本引理成立：包含态射
$\mathcal{Z}' \to \mathcal{Z} \times_\mathcal{Y} \mathcal{Y}'$
（引理 \ref{stacks-morphisms-lemma-factor-factor}）必为同构，因为沿满光滑态射
$V' \to \mathcal{Y}'$ 作基变换后它是同构。

\medskip\noindent
证明 (a)。令 $R = V \times_\mathcal{Y} V$。由《叠的性质》引理
\ref{stacks-properties-lemma-substacks-presentation}，规则
$\mathcal{Z} \mapsto \mathcal{Z} \times_\mathcal{Y} V$
给出 $\mathcal{Y}$ 的闭子叠与 $V$ 的 $R$-不变闭子空间之间的
$1$ 对 $1$ 对应。
此外，$f : \mathcal{X} \to \mathcal{Y}$ 经由 $\mathcal{Z}$ 分解，
当且仅当其基变换 $g : \mathcal{X} \times_\mathcal{Y} V \to V$
经由 $\mathcal{Z} \times_\mathcal{Y} V$ 分解。我们断言：$g$
的概形论像 $Z \subset V$ 是 $R$-不变的。由刚才所述，该断言蕴含 (a)。

\medskip\noindent
对每个 $i$，态射 $\mathcal{X} \times_\mathcal{Y} V_i \to V_i$ 拟紧，
因而 $\mathcal{X} \times_\mathcal{Y} V_i$ 拟紧。所以可以选取仿射概形
$W_i$ 以及满光滑态射
$W_i \to \mathcal{X} \times_\mathcal{Y} V_i$。注意，
$W = \coprod W_i$ 是概形，带有满光滑态射
$W \to \mathcal{X} \times_\mathcal{Y} V$，并且与 $g$ 复合所得
$W \to V$ 拟紧。设 $Z \to V$ 是 $W \to V$ 的概形论像，见《态射》一节
\ref{morphisms-section-scheme-theoretic-image} 和《代数空间的态射》一节
\ref{spaces-morphisms-section-scheme-theoretic-image}。由引理
\ref{stacks-morphisms-lemma-cover-upstairs}，$Z \subset V$ 是 $g$ 的概形论像。为证明
$Z$ 是 $R$-不变的，我们断言二者
$$
\text{pr}_0^{-1}(Z), \text{pr}_1^{-1}(Z) \subset R = V \times_\mathcal{Y} V
$$
都是 $\mathcal{X} \times_\mathcal{Y} R \to R$ 的概形论像。事实上，
先用《代数空间的态射》引理
\ref{spaces-morphisms-lemma-flat-base-change-scheme-theoretic-image}
可知，$\text{pr}_0^{-1}(Z)$ 是复合
$$
W \times_{V, \text{pr}_0} R = W \times_\mathcal{Y} V \to
\mathcal{X} \times_\mathcal{Y} R \to R
$$
的概形论像。由于这里第一个箭头满且光滑，可知
$\text{pr}_0^{-1}(Z)$ 是
$\mathcal{X} \times_\mathcal{Y} R \to R$ 的概形论像。同一论证适用于
$\text{pr}_1^{-1}(Z)$。所以 $Z$ 是 $R$-不变的。

\medskip\noindent
陈述 (b) 的证明与 (a) 完全相同。

\medskip\noindent
证明 (c)。设 $Z_i \subset V_i$ 是
$\mathcal{X} \times_\mathcal{Y} V_i \to V_i$ 的概形论像，并设
$Z_{i, j} \subset V_{i, j}$ 是
$\mathcal{X} \times_\mathcal{Y} V_{i, j} \to V_{i, j}$ 的概形论像。
显然，只需证明 $Z_i$ 在 $V_{i, j}$ 中的逆像是 $Z_{i, j}$。上文已经看到，
$Z_i$ 是 $W_i \to V_i$ 的概形论像；同理，$Z_{i, j}$ 是
$W_i \times_{V_i} V_{i, j} \to V_{i, j}$ 的概形论像。因此，该等式由
概形的情形（《态射》引理
\ref{morphisms-lemma-flat-base-change-scheme-theoretic-image}）以及
$V_{i, j} \to V_i$ 平坦这一事实得出。
\end{proof}

\begin{lemma}
\label{stacks-morphisms-lemma-topology-scheme-theoretic-image}
设 $f : \mathcal{X} \to \mathcal{Y}$ 是代数叠的拟紧态射。设
$\mathcal{Z} \subset \mathcal{Y}$ 是 $f$ 的概形论像。则
$|\mathcal{Z}|$ 是 $|f|$ 的像的闭包。
\end{lemma}

\begin{proof}
设 $z \in |\mathcal{Z}|$ 是一点。选取仿射概形 $V$、点 $v \in V$
以及把 $v$ 映到 $z$ 的光滑态射 $V \to \mathcal{Y}$。则
$\mathcal{X} \times_\mathcal{Y} V$ 是拟紧代数叠。因此，可以找到仿射概形
$W$ 以及满光滑态射 $W \to \mathcal{X} \times_\mathcal{Y} V$。
由引理 \ref{stacks-morphisms-lemma-existence-plus-flat-base-change}，
$\mathcal{X} \times_\mathcal{Y} V \to V$ 的概形论像为
$Z = \mathcal{Z} \times_\mathcal{Y} V$。所以由《叠的性质》引理
\ref{stacks-properties-lemma-points-cartesian}，$|\mathcal{Z}|$
在 $|V|$ 中的逆像是 $|Z|$。由引理 \ref{stacks-morphisms-lemma-cover-upstairs}，
$Z$ 是 $W \to V$ 的概形论像。由《代数空间的态射》引理
\ref{spaces-morphisms-lemma-quasi-compact-scheme-theoretic-image}，
$|W| \to |Z|$ 的像稠密。因此，
$|\mathcal{X} \times_\mathcal{Y} V| \to |Z|$ 的像稠密。注意
$v \in Z$。由于 $|V| \to |\mathcal{Y}|$ 是开映射，一个拓扑论证说明
$z$ 位于 $|f|$ 的像的闭包中，正合所需。
\end{proof}

\begin{lemma}
\label{stacks-morphisms-lemma-scheme-theoretic-image-of-partial-section}
设 $f : \mathcal{X} \to \mathcal{Y}$ 是可由代数空间表示且分离的代数叠态射。
设 $\mathcal{V} \subset \mathcal{Y}$ 是开子叠，且
$\mathcal{V} \to \mathcal{Y}$ 拟紧。设
$s : \mathcal{V} \to \mathcal{X}$ 是态射，满足
$f \circ s = \text{id}_\mathcal{V}$。设 $\mathcal{Y}'$ 是 $s$
的概形论像。则 $\mathcal{Y}' \to \mathcal{Y}$ 在
$\mathcal{V}$ 上是同构。
\end{lemma}

\begin{proof}
由引理 \ref{stacks-morphisms-lemma-quasi-compact-permanence}，态射
$s : \mathcal{V} \to \mathcal{Y}$ 拟紧。因此，由引理
\ref{stacks-morphisms-lemma-existence-plus-flat-base-change}，$s$ 的概形论像
$\mathcal{Y}'$ 的构造与平坦基变换交换。为证明本引理，可以假设
$\mathcal{Y}$ 可由代数空间表示，从而归结到代数空间的情形，即
《代数空间的态射》引理
\ref{spaces-morphisms-lemma-scheme-theoretic-image-of-partial-section}。
\end{proof}









\section{赋值判据}
\label{stacks-morphisms-section-valuative}

\noindent
陈述赋值判据时必须谨慎。确切地说，陈述中需要谈论交换图，但我们是在
$2$-范畴中工作，还必须确保 $2$-态射也按正确方式复合！

\begin{definition}
\label{stacks-morphisms-definition-fill-in-diagram}
设 $f : \mathcal{X} \to \mathcal{Y}$ 是代数叠的态射。考虑实线部分
$2$-交换的图
\begin{equation}
\label{stacks-morphisms-equation-diagram}
\vcenter{
\xymatrix{
\Spec(K) \ar[r]_-x \ar[d]_j & \mathcal{X} \ar[d]^f \\
\Spec(A) \ar[r]^-y \ar@{..>}[ru] & \mathcal{Y}
}
}
\end{equation}
，其中 $A$ 是分式域为 $K$ 的赋值环。设
$$
\gamma : y \circ j \longrightarrow f \circ x
$$
是见证该图 $2$-交换的 $2$-态射
（记号同《范畴》诸节 \ref{categories-section-formal-cat-cat} 和
\ref{categories-section-2-categories}）。给定
(\ref{stacks-morphisms-equation-diagram}) 与 $\gamma$，一个{\it 虚线箭头}是三元组
$(a, \alpha, \beta)$，由态射 $a : \Spec(A) \to \mathcal{X}$
及 $2$-箭头 $\alpha : a \circ j \to x$、$\beta : y \to f \circ a$
组成，并且满足
$\gamma = (\text{id}_f \star \alpha) \circ (\beta \star \text{id}_j)$；
换言之，下图交换：
$$
\xymatrix{
& f \circ a \circ j \ar[rd]^{\text{id}_f \star \alpha} \\
y \circ j \ar[ru]^{\beta \star \text{id}_j} \ar[rr]^\gamma & &
f \circ x
}
$$
。一个{\it 虚线箭头的态射}
$(a, \alpha, \beta) \to (a', \alpha', \beta')$ 是满足
$\alpha = \alpha' \circ (\theta \star \text{id}_j)$ 以及
$\beta' = (\text{id}_f \star \theta) \circ \beta$ 的
$2$-箭头 $\theta : a \to a'$。
\end{definition}

\noindent
上述定义是《范畴》定义 \ref{categories-definition-dotted-arrows} 的特殊情形。
一般而言，虚线箭头的范畴依赖于 $\gamma$。若 $\mathcal{Y}$ 可由代数空间表示
（或者域上对象的自同构群平凡），则当然至多有一个 $\gamma$，无需检查三角形
的交换性。更一般地，有引理
\ref{stacks-morphisms-lemma-cat-dotted-arrows-independent}。三角形的交换性在证明与基变换的
相容性时很重要，见引理 \ref{stacks-morphisms-lemma-cat-dotted-arrows-base-change} 的证明。

\begin{lemma}
\label{stacks-morphisms-lemma-cat-dotted-arrows}
在定义 \ref{stacks-morphisms-definition-fill-in-diagram} 的情形下，虚线箭头的范畴是群胚。
若 $\Delta_f$ 分离，则它是集合胚。
\end{lemma}

\begin{proof}
由于 $2$-箭头可逆，虚线箭头的范畴显然是群胚。给定虚线箭头
$(a, \alpha, \beta)$，虚线箭头 $(a, \alpha, \beta)$ 的一个自同构
是满足两个条件的 $2$-态射
$\theta : a \to a$。第一个条件
$\beta = (\text{id}_f \star \theta) \circ \beta$ 表明 $\theta$ 定义态射
$(a, \theta) : \Spec(A) \to \mathcal{I}_{\mathcal{X}/\mathcal{Y}}$.
第二个条件
$\alpha = \alpha \circ (\theta \star \text{id}_j)$
蕴含 $(a, \theta)$ 到 $\Spec(K)$ 的限制为恒等态射。图示为
$$
\xymatrix{
\mathcal{I}_{\mathcal{X}/\mathcal{Y}} \ar[d] & &
\Spec(K) \ar[d]^j \ar[ll]_{(a \circ j, \text{id})} \\
\mathcal{X} & & \Spec(A) \ar[ll]_a \ar[llu]_{(a, \theta)}
}
$$
换言之，若 $G \to \Spec(A)$ 是沿 $a$ 拉回相对惯性所得的群代数空间，
则 $\theta$ 定义点 $\theta \in G(A)$，其在 $G(K)$ 中的像平凡。
当然，若单位元 $e : \Spec(A) \to G$ 是闭浸入，则这只有在
$\theta$ 为恒等态射时才可能发生。查看引理
\ref{stacks-morphisms-lemma-diagonal-diagonal} 即得所需结论。
\end{proof}

\begin{lemma}
\label{stacks-morphisms-lemma-cat-dotted-arrows-independent}
在定义 \ref{stacks-morphisms-definition-fill-in-diagram} 中，假设
$\mathcal{I}_\mathcal{Y} \to \mathcal{Y}$ 固有
（例如，若 $\mathcal{Y}$ 分离，或 $\mathcal{Y}$ 在某个代数空间上分离）。
则虚线箭头的范畴不依赖于 $\gamma$ 的选取（精确到非典范等价）；而虚线箭头
对某个 $\gamma$ 存在（从而等价地对所有选择存在），当且仅当存在图
$$
\xymatrix{
\Spec(K) \ar[r]_-x \ar[d]_j & \mathcal{X} \ar[d]^f \\
\Spec(A) \ar[r]^-y \ar[ru]_a & \mathcal{Y}
}
$$
，其两个三角形都 $2$-交换（无需检查 $2$-态射是否正确复合）。
\end{lemma}

\begin{proof}
设 $\gamma, \gamma' : y \circ j \longrightarrow f \circ x$ 是两个
$2$-态射。则 $\gamma^{-1} \circ \gamma'$ 是 $y$ 在 $\Spec(K)$ 上的
自同构。因此，若
$\mathit{Isom}_\mathcal{Y}(y, y) \to \Spec(A)$ 固有，则由固有性的赋值判据
（《代数空间的态射》引理
\ref{spaces-morphisms-lemma-characterize-proper}），可以找到
$\delta : y \to y$，其到 $\Spec(K)$ 的限制为
$\gamma^{-1} \circ \gamma'$。于是可以用 $\delta$ 定义 $\gamma$
的虚线箭头范畴与 $\gamma'$ 的虚线箭头范畴之间的等价：把
$(a, \alpha, \beta)$ 送到 $(a, \alpha, \beta \circ \delta)$。
最后一个陈述显然。
\end{proof}

\begin{lemma}
\label{stacks-morphisms-lemma-cat-dotted-arrows-base-change}
假设给定 $2$-交换图
$$
\xymatrix{
\Spec(K) \ar[r]_-{x'} \ar[d]_j &
\mathcal{X}' \ar[d]^p \ar[r]_q &
\mathcal{X} \ar[d]^f \\
\Spec(A) \ar[r]^-{y'} &
\mathcal{Y}' \ar[r]^g &
\mathcal{Y}
}
$$
，其中右侧方块为 $2$-笛卡儿的。选取 $2$-箭头
$\gamma' : y' \circ j \to p \circ x'$。令
$x = q \circ x'$、$y = g \circ y'$，并令
$\gamma : y \circ j \to f \circ x$ 为 $\gamma'$ 与右侧方块
$2$-交换性中隐含的 $2$-箭头的复合。则左侧方块及 $\gamma'$
的虚线箭头范畴，等价于外部矩形及 $\gamma$ 的虚线箭头范畴。
\end{lemma}

\begin{proof}
（若不用虚线箭头的范畴，而考察如引理
\ref{stacks-morphisms-lemma-cat-dotted-arrows-independent} 中产生两个 $2$-交换三角形的态射的
同构类集合，我们不知道如何证明本引理的类似结论；事实上，该类似结论很可能
是错的。）第一种证明：本引理是《范畴》引理
\ref{categories-lemma-cat-dotted-arrows-base-change} 的特殊情形。
第二种证明：按《范畴》引理
\ref{categories-lemma-2-product-categories-over-C} 所述，可以用 $2$-纤维积
$\mathcal{Y}' \times_\mathcal{Y} \mathcal{X}$ 替换
$\mathcal{X}'$。此时对象 $x'$ 变成三元组
$(y' \circ j, x, \gamma)$。然后，可以把外部矩形的虚线箭头
$(a, \alpha, \beta)$ 送到左侧方块的虚线箭头
$(a', \alpha', \beta')$，其中取 $a' = (y', a, \beta)$、
$\alpha' = (\text{id}_{y' \circ j}, \alpha)$ 以及
$\beta' = \text{id}_{y'}$。细节略。
\end{proof}

\begin{lemma}
\label{stacks-morphisms-lemma-cat-dotted-arrows-composition}
假设给定 $2$-交换图
$$
\xymatrix{
\Spec(K) \ar[r]_-x \ar[dd]_j & \mathcal{X} \ar[d]^f \\
& \mathcal{Y} \ar[d]^g \\
\Spec(A) \ar[r]^-z & \mathcal{Z}
}
$$
选取 $2$-箭头 $\gamma : z \circ j \to g \circ f \circ x$。设
$\mathcal{C}$ 是外部矩形及 $\gamma$ 的虚线箭头范畴。设
$\mathcal{C}'$ 是下列方块
$$
\xymatrix{
\Spec(K) \ar[r]_-{f \circ x} \ar[d]_j & \mathcal{Y} \ar[d]^g \\
\Spec(A) \ar[r]^-z & \mathcal{Z}
}
$$
及 $\gamma$ 的虚线箭头范畴。则 $\mathcal{C}$ 等价于一个范畴
$\mathcal{C}''$，它具有下述性质：存在函子
$\mathcal{C}'' \to \mathcal{C}'$，使 $\mathcal{C}''$ 成为
$\mathcal{C}'$ 上纤维化于群胚的范畴，并且其纤维范畴是某些如下形式方块
$$
\xymatrix{
\Spec(K) \ar[r]_-x \ar[d]_j & \mathcal{X} \ar[d]^f \\
\Spec(A) \ar[r]^-y & \mathcal{Y}
}
$$
及某些 $y \circ j \to f \circ x$ 的选取所对应的虚线箭头范畴。
\end{lemma}

\begin{proof}
本引理是《范畴》引理
\ref{categories-lemma-cat-dotted-arrows-composition} 的特殊情形。
\end{proof}

\begin{definition}
\label{stacks-morphisms-definition-uniqueness}
设 $f : \mathcal{X} \to \mathcal{Y}$ 是代数叠的态射。若对每个图
(\ref{stacks-morphisms-equation-diagram}) 及定义
\ref{stacks-morphisms-definition-fill-in-diagram} 中的 $\gamma$，虚线箭头的范畴或者为空，或者是
恰有一个同构类的集合胚，则称 $f$ 满足{\it 赋值判据的唯一性部分}。
\end{definition}

\begin{lemma}
\label{stacks-morphisms-lemma-base-change-uniqueness}
满足赋值判据唯一性部分的代数叠态射，沿任意代数叠态射作基变换后，
仍是满足赋值判据唯一性部分的代数叠态射。
\end{lemma}

\begin{proof}
由引理 \ref{stacks-morphisms-lemma-cat-dotted-arrows-base-change} 和定义得出。
\end{proof}

\begin{lemma}
\label{stacks-morphisms-lemma-composition-uniqueness}
满足赋值判据唯一性部分的代数叠态射的复合，仍是满足赋值判据唯一性部分的
代数叠态射。
\end{lemma}

\begin{proof}
由引理 \ref{stacks-morphisms-lemma-cat-dotted-arrows-composition} 和定义得出。
\end{proof}

\begin{lemma}
\label{stacks-morphisms-lemma-uniqueness-representable}
设 $f : \mathcal{X} \to \mathcal{Y}$ 是可由代数空间表示的代数叠态射。
下列条件等价：
\begin{enumerate}
\item $f$ 满足赋值判据的唯一性部分；
\item 对每个概形 $T$ 和态射 $T \to \mathcal{Y}$，态射
$\mathcal{X} \times_\mathcal{Y} T \to T$ 作为代数空间的态射，
满足赋值判据的唯一性部分。
\end{enumerate}
\end{lemma}

\begin{proof}
由引理 \ref{stacks-morphisms-lemma-cat-dotted-arrows-base-change} 和定义得出。
\end{proof}

\begin{definition}
\label{stacks-morphisms-definition-existence}
设 $f : \mathcal{X} \to \mathcal{Y}$ 是代数叠的态射。若对每个图
(\ref{stacks-morphisms-equation-diagram}) 及定义
\ref{stacks-morphisms-definition-fill-in-diagram} 中的 $\gamma$，都存在域扩张 $K'/K$
以及支配 $A$ 的赋值环 $A' \subset K'$，使得下图外部矩形
$$
\xymatrix{
\Spec(K') \ar[r] \ar@/^2em/[rr]_{x'} \ar[d]_{j'} &
\Spec(K) \ar[d]_j \ar[r]_-x &
\mathcal{X} \ar[d]^f \\
\Spec(A') \ar[r] \ar@/_2em/[rr]^{y'} &
\Spec(A) \ar[r]^-y &
\mathcal{Y}
}
$$
连同诱导的 $2$-箭头
$\gamma' : y' \circ j' \to f \circ x'$ 所对应的虚线箭头范畴非空，
则称 $f$ 满足{\it 赋值判据的存在性部分}。
\end{definition}

\noindent
我们已经在代数空间态射的情形中看到，为得到正确的赋值判据概念，必须允许
分式域的扩张。见《代数空间的态射》例
\ref{spaces-morphisms-example-finite-separable-needed}。不过，对分离代数空间之间的
态射，不需要这种扩张（《代数空间的态射》引理
\ref{spaces-morphisms-lemma-usual-enough}）。然而，对代数叠之间的态射，
即使 $\mathcal{X}$ 和 $\mathcal{Y}$ 都分离，也可能需要扩张。例如，考虑
代数叠态射
$$
[\Spec(\mathbf{C})/G] \to \Spec(\mathbf{C})
$$
，它在基概形 $\Spec(\mathbf{C})$ 上，其中 $G$ 是阶为 $2$ 的群。
源和靶都是分离代数叠，而且该态射固有。因此，它满足赋值判据的唯一性部分和
存在性部分（见引理 \ref{stacks-morphisms-lemma-criterion-proper}）。但另一方面，存在图
$$
\xymatrix{
\Spec(K) \ar[r] \ar[d] & [\Spec(\mathbf{C})/G] \ar[d] \\
\Spec(A) \ar[r] & \Spec(\mathbf{C})
}
$$
，当 $A = \mathbf{C}[[t]]$ 且 $K = \mathbf{C}((t))$ 时不存在虚线箭头。
事实上，顶部水平箭头由 $K = \mathbf{C}((t))$ 的谱上的一个 $G$-挠子给出。
由于严格 Hensel 局部环 $A = \mathbf{C}[[t]]$ 上的任意 $G$-挠子都平凡，
可见若虚线箭头总存在，则 $K$ 上的每个 $G$-挠子都平凡。这并不成立，
因为 $G = \{+1, -1\}$，且由 Kummer 理论，$K$ 上的 $G$-挠子由非平凡群
$K^*/(K^*)^2$ 分类。

\begin{lemma}
\label{stacks-morphisms-lemma-base-change-existence}
满足赋值判据存在性部分的代数叠态射，沿任意代数叠态射作基变换后，
仍是满足赋值判据存在性部分的代数叠态射。
\end{lemma}

\begin{proof}
由引理 \ref{stacks-morphisms-lemma-cat-dotted-arrows-base-change} 和定义得出。
\end{proof}

\begin{lemma}
\label{stacks-morphisms-lemma-composition-existence}
满足赋值判据存在性部分的代数叠态射的复合，仍是满足赋值判据存在性部分的
代数叠态射。
\end{lemma}

\begin{proof}
由引理 \ref{stacks-morphisms-lemma-cat-dotted-arrows-composition} 和定义得出。
\end{proof}

\begin{lemma}
\label{stacks-morphisms-lemma-existence-representable}
设 $f : \mathcal{X} \to \mathcal{Y}$ 是可由代数空间表示的代数叠态射。
下列条件等价：
\begin{enumerate}
\item $f$ 满足赋值判据的存在性部分；
\item 对每个概形 $T$ 和态射 $T \to \mathcal{Y}$，态射
$\mathcal{X} \times_\mathcal{Y} T \to T$ 作为代数空间的态射，
满足赋值判据的存在性部分。
\end{enumerate}
\end{lemma}

\begin{proof}
由引理 \ref{stacks-morphisms-lemma-cat-dotted-arrows-base-change} 和定义得出。
\end{proof}

\begin{lemma}
\label{stacks-morphisms-lemma-closed-immersion-valuative-criteria}
代数叠的闭浸入同时满足赋值判据的存在性部分和唯一性部分。
\end{lemma}

\begin{proof}
略。提示：利用引理 \ref{stacks-morphisms-lemma-uniqueness-representable} 和
\ref{stacks-morphisms-lemma-existence-representable}，归结到概形的闭浸入情形。
\end{proof}





\section{第二对角态射的赋值判据}
\label{stacks-morphisms-section-valuative-second}

\noindent
逆向陈述已经在引理 \ref{stacks-morphisms-lemma-cat-dotted-arrows} 中证明。判据本身如下。

\begin{lemma}
\label{stacks-morphisms-lemma-setoids-and-diagonal}
设 $f : \mathcal{X} \to \mathcal{Y}$ 是代数叠的态射。若
$\Delta_f$ 拟分离，并且对每个图 (\ref{stacks-morphisms-equation-diagram}) 及定义
\ref{stacks-morphisms-definition-fill-in-diagram} 中 $\gamma$ 的选取，虚线箭头范畴都是
集合胚，则 $\Delta_f$ 分离。
\end{lemma}

\begin{proof}
我们将写出详细证明，但强烈建议读者在阅读引理
\ref{stacks-morphisms-lemma-cat-dotted-arrows} 证明中的论证后，受其启发自行找出证明。

\medskip\noindent
假设 $\Delta_f$ 拟分离，并且对每个图
(\ref{stacks-morphisms-equation-diagram}) 及定义
\ref{stacks-morphisms-definition-fill-in-diagram} 中 $\gamma$ 的选取，虚线箭头范畴都是集合胚。
由引理 \ref{stacks-morphisms-lemma-diagonal-diagonal}，只需证明
$e : \mathcal{X} \to \mathcal{I}_{\mathcal{X}/\mathcal{Y}}$
是闭浸入。由引理
\ref{stacks-morphisms-lemma-first-diagonal-separated-second-diagonal-closed}，事实上只需证明
$e = \Delta_{f, 2}$ 泛闭。任一上述引理都说明，由 $\Delta_f$
拟分离的假设，$e = \Delta_{f, 2}$ 拟紧。

\medskip\noindent
本段证明 $e$ 满足赋值判据的存在性部分。考虑实线部分 $2$-交换的图
$$
\xymatrix{
\Spec(K) \ar[r]_x \ar[d]_j & \mathcal{X} \ar[d]^e \\
\Spec(A) \ar[r]^{(a, \theta)} & \mathcal{I}_{\mathcal{X}/\mathcal{Y}}
}
$$
，并任取见证该图 $2$-交换的 $2$-态射
$\alpha : (a, \theta) \circ j \to e \circ x$
（这里用 $\alpha$ 代替定义 \ref{stacks-morphisms-definition-fill-in-diagram} 中使用的字母
$\gamma$）。注意，$f \circ \theta = \text{id}$；下文将用到这一点。
又有 $e \circ x = (x, \text{id}_x)$ 以及
$(a, \theta) \circ j = (a \circ j, \theta \star \text{id}_j)$。
因此，$\alpha$ 是与 $\theta \star \text{id}_j$ 和 $\text{id}_x$
相容的 $2$-箭头 $\alpha : a \circ j \to x$。令
$y = f \circ x$ 且 $\beta = \text{id}_{f \circ a}$。反向阅读引理
\ref{stacks-morphisms-lemma-cat-dotted-arrows} 证明中的论证，可知 $\theta$
是虚线箭头 $(a, \alpha, \beta)$ 的自同构，其中
$$
\gamma : y \circ j \to f \circ x
\quad\text{equal to}\quad
\text{id}_f \star \alpha : f \circ a \circ j \to f \circ x
$$
。另一方面，$\text{id}_a$ 也是自同构，所以由关于 $f$ 的假设可知
$\theta = \text{id}_a$。于是可以取态射
$a : \Spec(A) \to \mathcal{X}$ 作为上面所示图的虚线箭头，其中
$2$-态射 $(a, \text{id}_a) \circ j \to (x, \text{id}_x)$ 由
$\alpha$ 给出，而 $(a, \theta) \to e \circ a$ 由
$\text{id}_a$ 给出。

\medskip\noindent
由引理 \ref{stacks-morphisms-lemma-base-change-existence}，$e$ 的任意基变换都满足赋值判据的
存在性部分。由于 $e$ 可由代数空间表示，只需证明沿满光滑态射
$I \to \mathcal{I}_{\mathcal{X}/\mathcal{Y}}$ 作基变换后 $e$ 泛闭，
其中 $I$ 是代数空间（见《叠的性质》一节
\ref{stacks-properties-section-properties-morphisms}）。该基变换
$e' : X' \to I'$ 是代数空间的拟紧态射，并满足赋值判据的存在性部分，
所以由《代数空间的态射》引理
\ref{spaces-morphisms-lemma-quasi-compact-existence-universally-closed} 泛闭。
\end{proof}






\section{对角态射的赋值判据}
\label{stacks-morphisms-section-valuative-diagonal}

\noindent
所需结果是引理 \ref{stacks-morphisms-lemma-uniqueness-and-diagonal}。先陈述并证明一个形式辅助引理。

\begin{lemma}
\label{stacks-morphisms-lemma-helper-diagonal}
设 $f : \mathcal{X} \to \mathcal{Y}$ 是代数叠的态射。考虑实线部分
$2$-交换的图
$$
\xymatrix{
\Spec(K) \ar[rr]_-x \ar[d]_j & &
\mathcal{X} \ar[d]^{\Delta_f} \\
\Spec(A) \ar[rr]^{(a_1, a_2, \varphi)} \ar@{..>}[rru] & &
\mathcal{X} \times_\mathcal{Y} \mathcal{X}
}
$$
，其中 $A$ 是分式域为 $K$ 的赋值环。设
$\gamma : (a_1, a_2, \varphi) \circ j \longrightarrow \Delta_f \circ x$
是见证该图 $2$-交换的 $2$-态射。则
\begin{enumerate}
\item 把 $\gamma = (\alpha_1, \alpha_2)$ 写成
$\alpha_i : a_i \circ j \to x$ 的形式，则在图
$$
\xymatrix{
\Spec(K) \ar[r]_-x \ar[d]_j & \mathcal{X} \ar[d]^f \\
\Spec(A) \ar[r]^-{f \circ a_1} \ar@{..>}[ru] & \mathcal{Y}
}
$$
中得到两个虚线箭头 $(a_1, \alpha_1, \text{id})$ 和
$(a_2, \alpha_2, \varphi)$。
\item 原图及 $\gamma$ 的虚线箭头范畴是集合胚，其对象的同构类集合等于
虚线箭头范畴中的态射集
$(a_1, \alpha_1, \text{id}) \to (a_2, \alpha_2, \varphi)$。
\end{enumerate}
\end{lemma}

\begin{proof}
由于 $\Delta_f$ 可由代数空间表示（因而 $\Delta_f$ 的对角态射分离），
由引理 \ref{stacks-morphisms-lemma-cat-dotted-arrows}，本引理第一个交换图的虚线箭头范畴
是集合胚。本引理的其余陈述都由 $2$-图式计算得出，略去这些计算。
\end{proof}

\begin{lemma}
\label{stacks-morphisms-lemma-uniqueness-and-diagonal}
设 $f : \mathcal{X} \to \mathcal{Y}$ 是代数叠的态射。假设 $f$ 拟分离。
若 $f$ 满足赋值判据的唯一性部分，则 $f$ 分离。
\end{lemma}

\begin{proof}
关于 $f$ 的假设意味着 $\Delta_f$ 拟紧且拟分离
（定义 \ref{stacks-morphisms-definition-separated}）。必须证明 $\Delta_f$ 固有。
引理 \ref{stacks-morphisms-lemma-setoids-and-diagonal} 说明 $\Delta_f$ 分离。由引理
\ref{stacks-morphisms-lemma-properties-diagonal}，$\Delta_f$ 局部有限型。为完成证明，
还需证明 $\Delta_f$ 泛闭。一个形式论证（见引理
\ref{stacks-morphisms-lemma-helper-diagonal}）说明，赋值判据的唯一性部分蕴含：任意如下
实线图中都有虚线箭头：
$$
\xymatrix{
\Spec(K) \ar[d] \ar[r] & \mathcal{X} \ar[d]^{\Delta_f} \\
\Spec(A) \ar[r] \ar@{..>}[ru] & \mathcal{X} \times_\mathcal{Y} \mathcal{X}
}
$$
利用该性质在任意基变换下保持这一点，可知 $\Delta_f$ 沿任意映入
$\mathcal{X} \times_\mathcal{Y} \mathcal{X}$ 的代数空间所作的基变换
都满足赋值判据的存在性部分；再由代数空间态射的赋值判据，可知它泛闭，
见《代数空间的态射》引理
\ref{spaces-morphisms-lemma-quasi-compact-existence-universally-closed}。
\end{proof}

\noindent
下面是逆命题。

\begin{lemma}
\label{stacks-morphisms-lemma-converse-uniqueness-and-diagonal}
设 $f : \mathcal{X} \to \mathcal{Y}$ 是代数叠的态射。若 $f$ 分离，
则 $f$ 满足赋值判据的唯一性部分。
\end{lemma}

\begin{proof}
由于 $f$ 分离，由引理 \ref{stacks-morphisms-lemma-cat-dotted-arrows}，所有虚线箭头范畴
都是集合胚。考虑图
$$
\xymatrix{
\Spec(K) \ar[r]_-x \ar[d]_j & \mathcal{X} \ar[d]^f \\
\Spec(A) \ar[r]^-y \ar@{..>}[ru] & \mathcal{Y}
}
$$
以及定义 \ref{stacks-morphisms-definition-fill-in-diagram} 中的 $2$-态射
$\gamma : y \circ j \to f \circ x$。考虑虚线箭头范畴的两个对象
$(a, \alpha, \beta)$ 和 $(a', \beta', \alpha')$。为完成证明，
必须证明这两个对象同构。同构
$$
f \circ a \xrightarrow{\beta^{-1}} y \xrightarrow{\beta'} f \circ a'
$$
意味着 $(a, a', \beta' \circ \beta^{-1})$ 是态射
$\Spec(A) \to \mathcal{X} \times_\mathcal{Y} \mathcal{X}$。
另一方面，$\alpha$ 和 $\alpha'$ 定义 $2$-箭头
$$
(a, a', \beta' \circ \beta^{-1}) \circ j =
(a \circ j, a' \circ j,
(\beta' \star \text{id}_j) \circ (\beta \star \text{id}_j)^{-1})
\xrightarrow{(\alpha, \alpha')} (x, x, \text{id}) = \Delta_f \circ x
$$
这里使用了 $(a, \alpha, \beta)$ 和 $(a', \alpha', \beta')$
都是关于 $\gamma$ 的虚线箭头这一事实。由此得到交换图
$$
\xymatrix{
\Spec(K) \ar[d]_j \ar[rr]_x & & \mathcal{X} \ar[d]^{\Delta_f} \\
\Spec(A) \ar[rr]^{(a, a', \beta' \circ \beta^{-1})} & &
\mathcal{X} \times_\mathcal{Y} \mathcal{X}
}
$$
，其 $2$-交换性由 $(\alpha, \alpha')$ 见证。现在，$\Delta_f$
可由代数空间表示（引理 \ref{stacks-morphisms-lemma-properties-diagonal}），并且因为 $f$
分离而固有。因此，由引理 \ref{stacks-morphisms-lemma-existence-representable}
和代数空间固有性的赋值判据（《代数空间的态射》引理
\ref{spaces-morphisms-lemma-characterize-proper}），可知存在虚线箭头。
展开该构造可见，这意味着 $(a, \alpha, \beta)$ 与
$(a', \alpha', \beta')$ 在虚线箭头范畴中同构，正合所需。
\end{proof}






\section{泛闭性的赋值判据}
\label{stacks-morphisms-section-valutive-criterion}

\noindent
先给出一个陈述。

\begin{lemma}
\label{stacks-morphisms-lemma-quasi-compact-existence-universally-closed}
设 $f : \mathcal{X} \to \mathcal{Y}$ 是代数叠的态射。假设
\begin{enumerate}
\item $f$ 拟紧；
\item $f$ 满足赋值判据的存在性部分。
\end{enumerate}
则 $f$ 泛闭。
\end{lemma}

\begin{proof}
由引理 \ref{stacks-morphisms-lemma-base-change-quasi-compact} 和
\ref{stacks-morphisms-lemma-base-change-existence}，性质 (1) 和 (2) 在任意基变换下保持。
由引理 \ref{stacks-morphisms-lemma-universally-closed-local}，只需证明：每当 $T$
是映入 $\mathcal{Y}$ 的仿射概形时，
$|T \times_\mathcal{Y} \mathcal{X}| \to |T|$ 是闭映射。因此，只需证明：
若 $f : \mathcal{X} \to Y$ 是从代数叠到仿射概形的拟紧态射，并满足
赋值判据的存在性部分，则 $|f|$ 是闭映射。设
$T \subset |\mathcal{X}|$ 是闭子集。为完成证明，必须证明 $f(T)$ 闭。

\medskip\noindent
设 $\mathcal{Z} \subset \mathcal{X}$ 是 $T$ 上的约化诱导代数叠结构
（《叠的性质》定义
\ref{stacks-properties-definition-reduced-induced-stack}）。则
$i : \mathcal{Z} \to \mathcal{X}$ 是闭浸入，而必须证明
$|\mathcal{Z}| \to |Y|$ 的像闭。闭浸入拟紧
（引理 \ref{stacks-morphisms-lemma-closed-immersion-quasi-compact}）且满足赋值判据的存在性部分
（引理 \ref{stacks-morphisms-lemma-closed-immersion-valuative-criteria}）；拟紧态射的复合拟紧
（引理 \ref{stacks-morphisms-lemma-composition-quasi-compact}）；复合还保持满足赋值判据存在性
部分这一性质（引理 \ref{stacks-morphisms-lemma-composition-existence}）。所以只需证明：
若 $f : \mathcal{X} \to Y$ 是从代数叠到仿射概形的拟紧态射，并满足
赋值判据的存在性部分，则 $|f|(|\mathcal{X}|)$ 闭。

\medskip\noindent
由于 $\mathcal{X}$ 拟紧（它在仿射概形 $Y$ 上拟紧），可以选取仿射概形
$U$ 以及满光滑态射 $U \to \mathcal{X}$
（《叠的性质》引理
\ref{stacks-properties-lemma-quasi-compact-stack}）。假设 $y \in Y$
位于 $U \to Y$ 的像的闭包中（换言之，位于 $|f|$ 的像的闭包中）。
由《态射》引理
\ref{morphisms-lemma-reach-points-scheme-theoretic-image}，可以找到分式域为
$K$ 的赋值环 $A$ 以及交换图
$$
\xymatrix{
\Spec(K) \ar[r] \ar[d] & U \ar[d] \\
\Spec(A) \ar[r] & Y
}
$$
，使得 $\Spec(A)$ 的闭点映到 $y$。由假设，得到扩张 $K'/K$、
支配 $A$ 的赋值环 $A' \subset K'$，以及下图中的虚线箭头：
$$
\xymatrix{
\Spec(K') \ar[r] \ar[d] &
\Spec(K) \ar[r] \ar[d] &
U \ar[d] \ar[r] &
\mathcal{X} \ar[d]^f \\
\Spec(A') \ar[r] \ar@{..>}[rrru] &
\Spec(A) \ar[r] &
Y  \ar@{=}[r] &
Y
}
$$
因此 $y$ 位于 $|f|$ 的像中，得证。
\end{proof}

\noindent
下面是逆命题。

\begin{lemma}
\label{stacks-morphisms-lemma-converse-existence-universally-closed}
设 $f : \mathcal{X} \to \mathcal{Y}$ 是代数叠的态射。假设
\begin{enumerate}
\item $f$ 拟分离；
\item $f$ 泛闭。
\end{enumerate}
则 $f$ 满足赋值判据的存在性部分。
\end{lemma}

\begin{proof}
考虑实线图
$$
\xymatrix{
\Spec(K) \ar[r]_-x \ar[d]_j & \mathcal{X} \ar[d]^f \\
\Spec(A) \ar[r]^-y \ar@{..>}[ru] & \mathcal{Y}
}
$$
，其中 $A$ 是分式域为 $K$ 的赋值环，并且
$\gamma : y \circ j \longrightarrow f \circ x$ 如定义
\ref{stacks-morphisms-definition-fill-in-diagram} 所述。由引理
\ref{stacks-morphisms-lemma-cat-dotted-arrows-base-change}，为找到虚线箭头
（可能先把 $K$ 换成一个扩张，并把 $A$ 换成支配它的赋值环），
可以用 $\Spec(A)$ 替换 $\mathcal{Y}$，并用
$\Spec(A) \times_\mathcal{Y} \mathcal{X}$ 替换 $\mathcal{X}$。
这里当然用到了拟分离和泛闭在基变换下保持。因此归结到下一段所述的情形。

\medskip\noindent
考虑实线图
$$
\xymatrix{
\Spec(K) \ar[r]_-x \ar[d]_j & \mathcal{X} \ar[d]^f \\
\Spec(A) \ar@{=}[r] \ar@{..>}[ru] & \Spec(A)
}
$$
，其中 $A$ 是分式域为 $K$ 的赋值环，如定义
\ref{stacks-morphisms-definition-fill-in-diagram} 所述。由引理
\ref{stacks-morphisms-lemma-quasi-compact-permanence} 以及 $f$ 拟分离这一事实，
态射 $x$ 拟紧。由于 $f$ 泛闭，特别地，
$|f|(\overline{\{x\}})$ 是 $\Spec(A)$ 中的闭集。
这个像包含 $\Spec(A)$ 的一般点，所以存在位于 $x$ 的闭包中的点
$x' \in |\mathcal{X}|$，它映到 $\Spec(A)$ 的闭点。
由引理 \ref{stacks-morphisms-lemma-reach-points-scheme-theoretic-image}，可以找到交换图
$$
\xymatrix{
\Spec(K') \ar[r] \ar[d] & \Spec(K) \ar[d] \\
\Spec(A') \ar[r] & \mathcal{X}
}
$$
，使得 $\Spec(A')$ 的闭点映到 $x' \in |\mathcal{X}|$。
于是 $\Spec(A') \to \Spec(A)$ 把闭点映到闭点，即 $A'$ 支配 $A$，
证明完成。
\end{proof}





\section{固有性的赋值判据}
\label{stacks-morphisms-section-valutive-criterion-properness}

\noindent
陈述如下。

\begin{lemma}
\label{stacks-morphisms-lemma-criterion-proper}
设 $f : \mathcal{X} \to \mathcal{Y}$ 是代数叠的态射。
假设 $f$ 有限型且拟分离。则下列条件等价：
\begin{enumerate}
\item $f$ 固有；
\item $f$ 同时满足赋值判据的唯一性部分和存在性部分。
\end{enumerate}
\end{lemma}

\begin{proof}
固有态射就是分离、有限型且泛闭的态射。因此，本引理由引理
\ref{stacks-morphisms-lemma-uniqueness-and-diagonal},
\ref{stacks-morphisms-lemma-converse-uniqueness-and-diagonal},
\ref{stacks-morphisms-lemma-quasi-compact-existence-universally-closed} 以及
\ref{stacks-morphisms-lemma-converse-existence-universally-closed} 得出。
\end{proof}







\section{局部完全交态射}
\label{stacks-morphisms-section-lci}

\noindent
代数空间态射“为局部完全交态射”的性质在源与靶上光滑局部；见
《代数空间上的下降》引理 \ref{spaces-descent-lemma-smooth-local-source-target}
以及《空间态射进阶》引理
\ref{spaces-more-morphisms-lemma-descending-property-lci} 和
\ref{spaces-more-morphisms-lemma-lci-syntomic-local-source}。
由上面的引理 \ref{stacks-morphisms-lemma-local-source-target}，可以如下定义一个
代数空间的态射为局部完全交态射的含义；当源和靶都是代数空间时，
它与《空间态射进阶》第 \ref{spaces-more-morphisms-section-lci} 节中
已经定义的概念一致。

\begin{definition}
\label{stacks-morphisms-definition-lci}
设 $f : \mathcal{X} \to \mathcal{Y}$ 是代数叠的态射。若引理
\ref{stacks-morphisms-lemma-local-source-target} 中的等价条件在
$\mathcal{P} = \text{local complete intersection}$ 时成立，
则称 $f$ 是{\it 局部完全交态射}或 {\it Koszul} 的。
\end{definition}

\begin{lemma}
\label{stacks-morphisms-lemma-composition-lci}
局部完全交态射的复合仍是局部完全交态射。
\end{lemma}

\begin{proof}
结合评注 \ref{stacks-morphisms-remark-composition} 与《空间态射进阶》引理
\ref{spaces-more-morphisms-lemma-composition-lci}。
\end{proof}

\begin{lemma}
\label{stacks-morphisms-lemma-flat-base-change-lci}
局部完全交态射的平坦基变换仍是局部完全交态射。
\end{lemma}

\begin{proof}
略。提示：完全按照评注 \ref{stacks-morphisms-remark-base-change} 论证（但只考虑平坦的
$\mathcal{Y}' \to \mathcal{Y}$），并使用《空间态射进阶》引理
\ref{spaces-more-morphisms-lemma-flat-base-change-lci}。
\end{proof}

\begin{lemma}
\label{stacks-morphisms-lemma-lci-permanence}
设
$$
\xymatrix{
\mathcal{X} \ar[rr]_f \ar[rd] & & \mathcal{Y} \ar[ld] \\
& \mathcal{Z}
}
$$
是代数叠态射的交换图。假设 $\mathcal{Y} \to \mathcal{Z}$ 光滑，且
$\mathcal{X} \to \mathcal{Z}$ 是局部完全交态射。则
$f : \mathcal{X} \to \mathcal{Y}$ 是局部完全交态射。
\end{lemma}

\begin{proof}
选择概形 $W$ 和满光滑态射 $W \to \mathcal{Z}$。选择概形 $V$
和满光滑态射 $V \to W \times_\mathcal{Z} \mathcal{Y}$。
再选择概形 $U$ 和满光滑态射
$U \to V \times_\mathcal{Y} \mathcal{X}$。于是 $U \to W$
是概形的局部完全交态射，而 $V \to W$ 是概形的光滑态射。
由概形情形的结论（《态射进阶》引理
\ref{more-morphisms-lemma-lci-permanence}），可知 $U \to V$
是局部完全交态射。按定义，这意味着 $f$ 是局部完全交态射。
\end{proof}





\section{稳定子保持态射}
\label{stacks-morphisms-section-stabilizer-preserving}

\noindent
文献中，若代数叠的态射 $f : \mathcal{X} \to \mathcal{Y}$ 所诱导的态射
$\mathcal{I}_\mathcal{X} \to
\mathcal{X} \times_\mathcal{Y} \mathcal{I}_\mathcal{Y}$
是同构，则称这个态射是{\it 稳定子保持的}或{\it 不动点反映的}。
这种态射在 $\mathcal{X}$ 的每一点处都诱导自同构群之间的同构
（评注 \ref{stacks-morphisms-remark-identify-automorphism-groups}）。
本节证明关于这一概念的一些简单引理。

\begin{lemma}
\label{stacks-morphisms-lemma-stabilizer-preserving}
设 $f : \mathcal{X} \to \mathcal{Y}$ 是代数叠的态射。若
$\mathcal{I}_\mathcal{X} \to
\mathcal{X} \times_\mathcal{Y} \mathcal{I}_\mathcal{Y}$ 是同构，
则 $f$ 可由代数空间表示。
\end{lemma}

\begin{proof}
由引理 \ref{stacks-morphisms-lemma-second-diagonal} 立即得到。
\end{proof}

\begin{remark}
\label{stacks-morphisms-remark-get-property-auts-from-diagonal}
设 $f : \mathcal{X} \to \mathcal{Y}$ 是代数叠的态射，并设
$U \to \mathcal{X}$ 是以代数空间为源的态射。令 $G \to H$ 为态射
$\mathcal{I}_\mathcal{X} \to
\mathcal{X} \times_\mathcal{Y} \mathcal{I}_\mathcal{Y}$
到 $U$ 上的拉回。若 $\Delta_f$ 非分歧、\'etale 等，则
$G \to H$ 也具有相应性质。这是因为
$$
\xymatrix{
U \times_\mathcal{X} U \ar[r] \ar[d] & \mathcal{X} \ar[d]^{\Delta_f} \\
U \times_\mathcal{Y} U \ar[r] & \mathcal{X} \times_\mathcal{Y} \mathcal{X}
}
$$
是笛卡儿图，并且态射 $G \to H$ 是左侧竖直箭头沿对角态射
$U \to U \times U$ 的基变换。比较引理
\ref{stacks-morphisms-lemma-separated-implies-isom} 的证明。
\end{remark}

\begin{lemma}
\label{stacks-morphisms-lemma-aut-iso-unramified}
设 $f : \mathcal{X} \to \mathcal{Y}$ 是代数叠的非分歧态射。
下列条件等价：
\begin{enumerate}
\item $\mathcal{I}_\mathcal{X} \to
\mathcal{X} \times_\mathcal{Y} \mathcal{I}_\mathcal{Y}$
是同构；
\item 对所有 $x \in |\mathcal{X}|$，$f$ 都诱导 $x$ 和 $f(x)$ 处的
自同构群之间的同构（评注 \ref{stacks-morphisms-remark-identify-automorphism-groups}）。
\end{enumerate}
\end{lemma}

\begin{proof}
选择概形 $U$ 和满光滑态射 $U \to \mathcal{X}$。记 $G \to H$ 为态射
$\mathcal{I}_\mathcal{X} \to
\mathcal{X} \times_\mathcal{Y} \mathcal{I}_\mathcal{Y}$
到 $U$ 上的拉回。由评注 \ref{stacks-morphisms-remark-get-property-auts-from-diagonal}
和引理 \ref{stacks-morphisms-lemma-characterize-unramified}，态射 $G \to H$ 是 \'etale 的。
条件 (1) 等价于 $G \to H$ 是同构（例如应用《叠的性质》引理
\ref{stacks-properties-lemma-check-property-covering} 即得）。
条件 (2) 等价于：对每个 $u \in U$，纤维的态射 $G_u \to H_u$
都是同构。因此 (1) $\Rightarrow$ (2) 是显然的。若 (2) 成立，则
$G \to H$ 是代数空间的满、泛单且 \'etale 的态射。由《代数空间的态射》引理
\ref{spaces-morphisms-lemma-etale-universally-injective-open}，
这种态射是同构。
\end{proof}

\begin{lemma}
\label{stacks-morphisms-lemma-stabilizer-preserving-unramified}
\begin{reference}
\cite[Proposition 3.5]{rydh_quotients} 和
\cite[Proposition 2.5]{alper_quotient}
\end{reference}
设 $f : \mathcal{X} \to \mathcal{Y}$ 是代数叠的态射。假设
\begin{enumerate}
\item $f$ 可由代数空间表示且非分歧；
\item $\mathcal{I}_\mathcal{Y} \to \mathcal{Y}$ 固有。
\end{enumerate}
则使 $f$ 诱导 $x$ 和 $f(x)$ 处自同构群之间同构的那些
$x \in |\mathcal{X}|$ 所成的集合是开的
（评注 \ref{stacks-morphisms-remark-identify-automorphism-groups}）。令
$\mathcal{U} \subset \mathcal{X}$ 为相应的开子叠，则态射
$\mathcal{I}_\mathcal{U} \to
\mathcal{U} \times_\mathcal{Y} \mathcal{I}_\mathcal{Y}$
是同构。
\end{lemma}

\begin{proof}
选择概形 $U$ 和满光滑态射 $U \to \mathcal{X}$。记 $G \to H$ 为态射
$\mathcal{I}_\mathcal{X} \to
\mathcal{X} \times_\mathcal{Y} \mathcal{I}_\mathcal{Y}$
到 $U$ 上的拉回。由评注 \ref{stacks-morphisms-remark-get-property-auts-from-diagonal}
和引理 \ref{stacks-morphisms-lemma-characterize-unramified}，态射 $G \to H$ 是 \'etale 的。
由于 $f$ 可由代数空间表示，$G \to H$ 是单态射。因此 $G \to H$
是开浸入；见《代数空间的态射》引理
\ref{spaces-morphisms-lemma-etale-universally-injective-open}。
依假设，$H \to U$ 固有。

\medskip\noindent
完成这些准备后，可如下证明本引理。引理中 $|\mathcal{X}|$ 的子集的
逆像，显然就是使 $G_u \to H_u$ 为同构的那些 $u \in U$
所成的集合（毕竟 $G_u$ 是 $H_u$ 的开子群代数空间）。这是开子集，
因为其补集是闭子集 $|H| \setminus |G|$ 的像，而
$|H| \to |U|$ 是闭映射。由《叠的性质》引理
\ref{stacks-properties-lemma-open-substacks}，可以考虑
$\mathcal{X}$ 的相应开子叠 $\mathcal{U}$。把引理
\ref{stacks-morphisms-lemma-aut-iso-unramified} 应用于
$\mathcal{U} \to \mathcal{Y}$，即得引理的最后一个陈述。
\end{proof}

\begin{lemma}
\label{stacks-morphisms-lemma-base-change-stabilizer-preserving}
设
$$
\xymatrix{
\mathcal{X}' \ar[r] \ar[d]_{f'} & \mathcal{X} \ar[d]^f \\
\mathcal{Y}' \ar[r] & \mathcal{Y}
}
$$
是代数叠的笛卡儿图。
\begin{enumerate}
\item 设 $x' \in |\mathcal{X}'|$ 的像为 $x \in |\mathcal{X}|$。
若 $f$ 诱导 $x$ 和 $f(x)$ 处自同构群之间的同构
（评注 \ref{stacks-morphisms-remark-identify-automorphism-groups}），则 $f'$ 诱导
$x'$ 和 $f(x')$ 处自同构群之间的同构。
\item 若 $\mathcal{I}_\mathcal{X} \to
\mathcal{X} \times_\mathcal{Y} \mathcal{I}_\mathcal{Y}$ 是同构，则
$\mathcal{I}_{\mathcal{X}'} \to
\mathcal{X}' \times_{\mathcal{Y}'} \mathcal{I}_{\mathcal{Y}'}$
是同构。
\end{enumerate}
\end{lemma}

\begin{proof}
略。
\end{proof}

\begin{lemma}
\label{stacks-morphisms-lemma-stabilizer-preserving-points-cartesian}
设
$$
\xymatrix{
\mathcal{X}' \ar[r] \ar[d]_{f'} & \mathcal{X} \ar[d]^f \\
\mathcal{Y}' \ar[r]^g & \mathcal{Y}
}
$$
是代数叠的笛卡儿图。若 $f$ 在各点处诱导自同构群之间的同构
（评注 \ref{stacks-morphisms-remark-identify-automorphism-groups}），则
$$
\Mor(\Spec(k), \mathcal{X}')
\longrightarrow
\Mor(\Spec(k), \mathcal{Y}') \times \Mor(\Spec(k), \mathcal{X})
$$
对任意域 $k$ 都在同构类上为单射。
\end{lemma}

\begin{proof}
必须证明：给定 $(y', x)$，至多有一个映到它的 $x'$。
按照我们对 $2$-纤维积的构造，态射 $x'$ 由三元组
$(x, y', \alpha)$ 给出，其中
$\alpha : g \circ y' \to f \circ x$ 是 $2$-态射。
现在假设还有第二个这样的三元组 $(x, y', \beta)$。则
$\alpha$ 和 $\beta$ 相差自同构群代数空间 $G_{f(x)}$ 的一个
$k$-值点 $\epsilon$。依假设，$f$ 诱导同构
$G_x \to G_{f(x)}$，所以可以把 $\epsilon$ 提升为 $G_x$ 的
$k$-值点 $\gamma$。于是
$(\gamma, \text{id}) : (x, y', \alpha) \to
(x, y', \beta)$ 是所需的同构。
\end{proof}

\begin{lemma}
\label{stacks-morphisms-lemma-etale-iso}
设 $f : \mathcal{X} \to \mathcal{Y}$ 是代数叠的态射。假设 $f$ 是
\'etale 的，$f$ 在各点处诱导自同构群之间的同构
（评注 \ref{stacks-morphisms-remark-identify-automorphism-groups}），并且对每个代数闭域
$k$，函子
$$
f : \Mor(\Spec(k), \mathcal{X}) \longrightarrow \Mor(\Spec(k), \mathcal{Y})
$$
都是等价。则 $f$ 是同构。
\end{lemma}

\begin{proof}
由引理 \ref{stacks-morphisms-lemma-universally-injective}，$f$ 泛单。结合引理
\ref{stacks-morphisms-lemma-stabilizer-preserving} 和
\ref{stacks-morphisms-lemma-aut-iso-unramified}
可知 $f$ 可由代数空间表示。因此，由《代数空间的态射》引理
\ref{spaces-morphisms-lemma-etale-universally-injective-open}，
$f$ 是开浸入。最后注意，本引理的条件还保证 $f$ 满。
\end{proof}







\section{正规化}
\label{stacks-morphisms-section-normalization}

\noindent
本节是《代数空间的态射》第
\ref{spaces-morphisms-section-normalization} 节的类比版本。

\begin{lemma}
\label{stacks-morphisms-lemma-prepare-normalization}
设 $\mathcal{X}$ 是代数叠。下列条件等价：
\begin{enumerate}
\item 存在满光滑态射 $U \to \mathcal{X}$，其中 $U$ 是概形，
且 $U$ 的每个拟紧开子集都有有限多个不可约分支；
\item 对每个概形 $U$ 和每个光滑态射 $U \to \mathcal{X}$，
$U$ 的每个拟紧开子集都有有限多个不可约分支；
\item 对每个代数空间 $Y$ 和光滑态射 $Y \to \mathcal{X}$，
空间 $Y$ 满足《代数空间的态射》引理
\ref{spaces-morphisms-lemma-prepare-normalization} 中的等价条件；
\item 对每个在 $\mathcal{X}$ 上光滑的拟紧代数叠 $\mathcal{Y}$，
空间 $|\mathcal{Y}|$ 有有限多个不可约分支。
\end{enumerate}
\end{lemma}

\begin{proof}
(1)、(2) 与 (3) 的等价性由《下降》引理
\ref{descent-lemma-locally-finite-nr-irred-local-fppf}、
《叠的性质》引理 \ref{stacks-properties-lemma-type-property} 以及
《代数空间的态射》引理
\ref{spaces-morphisms-lemma-prepare-normalization} 得出。
由这些引用也显然可知条件 (4) 蕴含条件 (1)。反过来，假设等价条件
(1)、(2) 与 (3) 成立，并设 $\mathcal{Y} \to \mathcal{X}$
是代数叠的光滑态射，且 $\mathcal{Y}$ 拟紧。由《叠的性质》引理
\ref{stacks-properties-lemma-quasi-compact-stack}，可以选择仿射概形
$V$ 和满光滑态射 $V \to \mathcal{Y}$。由 (2)，$V$ 有有限多个
不可约分支；又因 $|V| \to |\mathcal{Y}|$ 满且连续，故由《拓扑》引理
\ref{topology-lemma-surjective-continuous-irreducible-components}，
$|\mathcal{Y}|$ 有有限多个不可约分支。
\end{proof}

\begin{lemma}
\label{stacks-morphisms-lemma-normalization}
设 $\mathcal{X}$ 是满足引理 \ref{stacks-morphisms-lemma-prepare-normalization}
中等价条件的代数叠。则存在代数叠的整态射
$$
\mathcal{X}^\nu \longrightarrow \mathcal{X}
$$
，使得对每个概形 $U$ 和光滑态射 $U \to \mathcal{X}$，
纤维积 $\mathcal{X}^\nu \times_\mathcal{X} U$ 是 $U$ 的正规化。
\end{lemma}

\begin{proof}
设 $U \to \mathcal{X}$ 是满光滑态射，其中 $U$ 是概形。令
$R = U \times_\mathcal{X} U$。回顾一下，我们得到代数空间中的光滑群胚
$(U, R, s, t, c)$ 和 $\mathcal{X}$ 的表示
$\mathcal{X} = [U/R]$；见《代数叠》引理
\ref{algebraic-lemma-map-space-into-stack} 和
\ref{algebraic-lemma-stack-presentation} 以及定义
\ref{algebraic-definition-presentation}。关于 $\mathcal{X}$ 的假设意味着
$U$ 的正规化 $U^\nu$ 有定义；见《态射》定义
\ref{morphisms-definition-normalization}。由《代数空间的态射》引理
\ref{spaces-morphisms-lemma-normalization}
，取正规化与代数空间的光滑态射可交换。因此，$R$ 的正规化 $R^\nu$
同时同构于 $R \times_{s, U} U^\nu$ 和 $U^\nu \times_{U, t} R$。
于是得到代数空间的两个光滑态射
$s^\nu : R^\nu \to U^\nu$ 和 $t^\nu : R^\nu \to U^\nu$
。对纤维积作形式计算可知，
$R^\nu \times_{s^\nu, U^\nu, t^\nu} R^\nu$ 是
$R \times_{s, U, t} R$ 的正规化。因此，光滑态射
$c : R \times_{s, U, t} R \to R$ 也延拓为 $c^\nu$。
类似地，逆元 $i : R \to R$（一个同构）诱导同构
$i^\nu : R^\nu \to R^\nu$。最后，单位元 $e : U \to R$
提升为 $e^\nu : U^\nu \to R^\nu$；例如，这是因为 $e$ 是 $s$
的截面且 $R^\nu = R \times_{U, s} U^\nu$。我们断言
$(U^\nu, R^\nu, s^\nu, t^\nu, c^\nu)$ 是代数空间中的光滑群胚。
为此，需要对 $(U^\nu, R^\nu, s^\nu, t^\nu, c^\nu, e^\nu, i^\nu)$
验证《群胚》第 \ref{groupoids-section-groupoids} 节中的公理
(1)、(2)(a)、(2)(b)、(3)(a) 和 (3)(b)。例如，对 (1)，必须证明两个态射
$a, b : R^\nu \times_{s^\nu, U^\nu, t^\nu} R^\nu
\times_{s^\nu, U^\nu, t^\nu} R^\nu \to R^\nu$ 相同。
相同。这是因为已知相应的一对态射
$R \times_{s, U, t} R \times_{s, U, t} R \to R$
相同，而态射 $a$ 和 $b$ 都是这个态射到正规化上的唯一延拓。
其他公理同理。

\medskip\noindent
考虑代数叠 $\mathcal{X}^\nu = [U^\nu/R^\nu]$（《代数叠》定理
\ref{algebraic-theorem-smooth-groupoid-gives-algebraic-stack}).
由于有代数空间中群胚的态射
$(U^\nu, R^\nu, s^\nu, t^\nu, c^\nu) \to (U, R, s, t, c)$
，得到代数叠的态射 $\nu : \mathcal{X}^\nu \to \mathcal{X}$。
由于 $R^\nu = R \times_{s, U} U^\nu$，由《空间中的群胚》引理
\ref{spaces-groupoids-lemma-criterion-fibre-product}.
可知 $U^\nu = \mathcal{X}^\nu \times_\mathcal{X} U$。特别地，
由于 $U^\nu \to U$ 是整态射，$\nu$ 也是整态射。
略去如下验证：对每个从概形到 $\mathcal{X}$ 的光滑态射，
本引理中陈述的基变换性质都成立。
\end{proof}

\noindent
由此得到下面的定义。

\begin{definition}
\label{stacks-morphisms-definition-normalization}
设 $\mathcal{X}$ 是满足引理 \ref{stacks-morphisms-lemma-prepare-normalization}
中等价条件的代数叠。把态射
$$
\nu : \mathcal{X}^\nu \longrightarrow \mathcal{X}
$$
（由引理 \ref{stacks-morphisms-lemma-normalization} 构造）定义为 $\mathcal{X}$ 的
{\it 正规化}。
\end{definition}










\section{点与特化}
\label{stacks-morphisms-section-specializations}

\noindent
本节是《适宜代数空间》第
\ref{decent-spaces-section-specializations} 节的类比版本。

\begin{lemma}
\label{stacks-morphisms-lemma-specialization}
设 $\mathcal{X}$ 是代数叠。设 $f : U \to \mathcal{X}$ 是光滑态射，
其中 $U$ 是代数空间。设 $x' \leadsto x$ 是 $|\mathcal{X}|$
中点的一个特化。设 $u \in |U|$ 且 $f(u) = x$。若
$(\mathcal{X}, x')$ 满足《叠的性质》引理
\ref{stacks-properties-lemma-UR-quasi-compact-above-x} 中的等价条件，
则在 $|U|$ 中存在特化 $u' \leadsto u$，且 $f(u') = x'$。
\end{lemma}

\begin{proof}
选择 \'etale 态射 $(U_1, u_1) \to (U, u)$，其中 $U_1$ 是仿射概形。
于是可以并且确实用 $U_1$ 替换 $U$。因此可假设 $U$ 是仿射概形。
考虑代数空间 $R = U \times_\mathcal{X} U$ 及其光滑投影
$t, s : R \to U$。选择映到 $x'$ 的点 $w \in U$；这是可能的，
因为 $f : |U| \to |\mathcal{X}|$ 是开映射。由关于 $x'$ 的假设，
$t : R \to U$ 在 $w$ 上的纤维
$F' = t^{-1}(w) = R \times_{t, U} w$ 是拟紧代数空间。
选择仿射概形 $T$ 和满 \'etale 态射 $T \to F'$。
$x' \leadsto x$ 这一事实意味着 $u$ 位于态射
$$
T \to F' \to R \xrightarrow{s} U
$$
的像的闭包中。事实上，这个像是 $|U| \to |\mathcal{X}'|$ 在
$x'$ 上的纤维；若某个开子集 $u \in V \subset |U|$ 与该纤维不交，
则 $f(V)$ 是 $x$ 的一个不含 $x'$ 的开邻域，矛盾。因此，由《态射》引理
\ref{morphisms-lemma-reach-points-scheme-theoretic-image}
，在 $|U| \to |\mathcal{X}|$ 位于 $x'$ 上的纤维中存在
$u' \in |U|$，并且它特化到 $u$。
\end{proof}



\section{适宜代数叠}
\label{stacks-morphisms-section-reasonable-decent}

\noindent
本节是《适宜代数空间》第
\ref{decent-spaces-section-reasonable-decent} 节的类比版本。
特别地，下面的定义与该处定义的适宜代数空间概念相容。

\begin{definition}
\label{stacks-morphisms-definition-decent}
设 $\mathcal{X}$ 是代数叠。若对每个 $x \in |\mathcal{X}|$，
《叠的性质》引理
\ref{stacks-properties-lemma-UR-quasi-compact-above-x} 中的等价条件都成立，
则称 $\mathcal{X}$ 是{\it 适宜的}。
\end{definition}

\noindent
有人会把这个定义改述为：$\mathcal{X}$ 的每个点都是拟紧的。
稍强一些的条件是，要求 $x$ 的等价类中的任意态射
$\Spec(k) \to \mathcal{X}$ 既拟分离又拟紧。

\begin{lemma}
\label{stacks-morphisms-lemma-quasi-separated-decent}
拟分离代数叠 $\mathcal{X}$ 是适宜的。更一般地，若
$\Delta : \mathcal{X} \to \mathcal{X} \times \mathcal{X}$ 拟紧，
则 $\mathcal{X}$ 是适宜的。
\end{lemma}

\begin{proof}
具体地，若 $\mathcal{X}$ 拟分离，则任意源为拟紧概形 $T$ 的态射
$f : T \to \mathcal{X}$ 都拟紧；见引理
\ref{stacks-morphisms-lemma-quasi-compact-permanence}。若只知道 $\Delta$ 拟紧，
则使用描述
$$
T \times_{f, \mathcal{X}, f'} T' =
(T \times T')
\times_{(f, f'), \mathcal{X} \times \mathcal{X}, \Delta}
\mathcal{X}
$$
来证明这一点。细节略。
\end{proof}

\begin{lemma}
\label{stacks-morphisms-lemma-representable-decent}
设 $f : \mathcal{X} \to \mathcal{Y}$ 是代数叠的态射。假设 $Y$
是适宜的，并且 $f$ 可表（由概形表示），或者 $f$ 可由代数空间表示且拟分离。
则 $\mathcal{X}$ 是适宜的。
\end{lemma}

\begin{proof}
设 $x \in |\mathcal{X}|$ 的像为 $y \in |\mathcal{Y}|$。
在定义 $y$ 的等价类中选择态射 $y : \Spec(k) \to \mathcal{Y}$。
令 $\mathcal{X}_y = \Spec(k) \times_{y, \mathcal{Y}} \mathcal{X}$。
选择映到 $x$ 的点 $x' \in |\mathcal{X}_y|$；见《叠的性质》引理
\ref{stacks-properties-lemma-points-cartesian}。在 $x'$ 的等价类中选择态射
$x' : \Spec(k') \to \mathcal{X}_y$。有图
$$
\xymatrix{
\Spec(k') \ar[r]_{x'} &
\mathcal{X}_y \ar[r] \ar[d] &
\mathcal{X} \ar[d] \\
&
\Spec(k) \ar[r]^y &
\mathcal{Y}
}
$$
若 $\mathcal{Y}$ 适宜，则态射 $y$ 拟紧。因此，作为基变换，
$\mathcal{X}_y \to \mathcal{X}$ 拟紧
（引理 \ref{stacks-morphisms-lemma-base-change-quasi-compact}）。所以为得到结论，
只需证明 $x'$ 拟紧（引理 \ref{stacks-morphisms-lemma-composition-quasi-compact}）。
若 $f$ 可表，则 $\mathcal{X}_y$ 是概形，且 $x'$ 拟紧。
若 $f$ 可由代数空间表示且拟分离，则 $\mathcal{X}_y$ 是拟分离代数空间，
因而适宜（《适宜代数空间》引理
\ref{decent-spaces-lemma-properties-trivial-implications}）。
\end{proof}

\begin{lemma}
\label{stacks-morphisms-lemma-qc-decent}
设 $f : \mathcal{X} \to \mathcal{Y}$ 是代数叠的态射。
若 $f$ 拟紧且满，并且 $\mathcal{X}$ 适宜，则 $\mathcal{Y}$ 适宜。
\end{lemma}

\begin{proof}
设 $x : \Spec(k) \to \mathcal{X}$ 是态射，其中 $k$ 是域，并记
$y = f \circ x$。由于 $f$ 满，$|\mathcal{Y}|$ 的每个点都以这种方式得到；
见《叠的性质》引理
\ref{stacks-properties-lemma-characterize-surjective}。
考虑仿射概形 $T$ 和态射 $T \to \mathcal{Y}$。只需证明
$T \times_{\mathcal{Y}, y} \Spec(k)$ 拟紧；见引理
\ref{stacks-morphisms-lemma-check-quasi-compact-covering}。我们有
$$
(T \times_{\mathcal{Y}} \mathcal{X}) \times_{\mathcal{X}, x} \Spec(k) =
T \times_{\mathcal{Y}, y} \Spec(k)
$$
态射 $T \times_{\mathcal{Y}} \mathcal{X} \to T$ 拟紧，因而
$T \times_\mathcal{Y} \mathcal{X}$ 拟紧。又因 $\mathcal{X}$ 适宜，
$x$ 是拟紧态射，所以所示纤维积拟紧。
\end{proof}

\begin{lemma}
\label{stacks-morphisms-lemma-gerbe-decent}
设 $f : \mathcal{X} \to \mathcal{Y}$ 是代数叠的态射。
若 $\mathcal{X}$ 是 $\mathcal{Y}$ 上的胚，且 $\mathcal{X}$ 适宜，
则 $\mathcal{Y}$ 适宜。
\end{lemma}

\begin{proof}
假设 $\mathcal{X}$ 是 $\mathcal{Y}$ 上的胚，且 $\mathcal{X}$ 适宜。
注意，由引理 \ref{stacks-morphisms-lemma-gerbe-bijection-points}，$f$ 是泛同胚。
因此本引理由引理 \ref{stacks-morphisms-lemma-qc-decent} 得出。
\end{proof}



\section{适宜叠上的点}
\label{stacks-morphisms-section-points-decent}

\noindent
本节是《适宜代数空间》第 \ref{decent-spaces-section-points} 节的类比版本。
我们不知道适宜代数叠所关联的拓扑空间是否总是清醒的；稍弱的结果见命题
\ref{stacks-morphisms-proposition-decent-sober}。

\begin{lemma}
\label{stacks-morphisms-lemma-kolmogorov}
设 $\mathcal{X}$ 是适宜代数叠。则 $|\mathcal{X}|$ 是柯尔莫哥洛夫空间
（见《拓扑学》定义 \ref{topology-definition-generic-point}）。
\end{lemma}

\begin{proof}
设 $x_1, x_2 \in |\mathcal{X}|$，并且 $x_1 \leadsto x_2$、
$x_2 \leadsto x_1$。必须证明 $x_1 = x_2$。令
$\mathcal{Z} \subset \mathcal{X}$ 为约化闭子叠，其 $|\mathcal{Z}|$ 等于
$\overline{\{x_1\}} = \overline{\{x_2\}}$。由引理
\ref{stacks-morphisms-lemma-representable-decent}，
$\mathcal{Z}$ 适宜。用 $\mathcal{Z}$ 替换 $\mathcal{X}$ 后，
归结到下一段所述的情形。

\medskip\noindent
假设 $|\mathcal{X}|$ 不可约，并且有一般点 $x_1$ 和 $x_2$。
选取仿射概形 $U$、点 $u_1, u_2 \in U$ 以及光滑态射
$f : U \to \mathcal{X}$，使得 $f(u_i) = x_i$。于是可以找到第三个点
$u_3 \in U$，它是 $U$ 的某个不可约分支的一般点，而其像
$x_3 \in |\mathcal{X}|$ {\it 也是} $|\mathcal{X}|$ 的一般点。
具体地，只需把 $u_3$ 取为经过 $u_1$（或任选 $u_2$）的某个不可约分支的
任意一般点。下一段将证明 $x_1 = x_3$ 且 $x_2 = x_3$，从而得到所需结论。

\medskip\noindent
由对称性，只需证明 $x_1 = x_3$。由于 $x_1$ 是 $|\mathcal{X}|$
的一般点，有特化 $x_1 \leadsto x_3$。由引理
\ref{stacks-morphisms-lemma-specialization}，可以在 $U$ 中找到映到
$x_1 \leadsto x_3$ 的特化 $u'_1 \leadsto u_3$（！）。然而，
$u_3$ 是某个不可约分支的一般点，所以如所需，$u'_1 = u_3$。
\end{proof}

\begin{lemma}
\label{stacks-morphisms-lemma-decent-locally-noetherian-sober}
设 $\mathcal{X}$ 是适宜的局部诺特代数叠。则 $|\mathcal{X}|$
是清醒的局部诺特拓扑空间。
\end{lemma}

\begin{proof}
由引理 \ref{stacks-morphisms-lemma-Noetherian-topology}，拓扑空间 $|\mathcal{X}|$
局部诺特。由引理 \ref{stacks-morphisms-lemma-kolmogorov}，拓扑空间 $|\mathcal{X}|$
是柯尔莫哥洛夫空间。由引理 \ref{stacks-morphisms-lemma-locally-Noetherian-quasi-sober}，
拓扑空间 $|\mathcal{X}|$ 拟清醒。证明完成；见《拓扑学》定义
\ref{topology-definition-generic-point}。
\end{proof}

\begin{proposition}
\label{stacks-morphisms-proposition-decent-sober}
设 $\mathcal{X}$ 是适宜代数叠，并且
$\mathcal{I}_\mathcal{X} \to \mathcal{X}$ 拟紧。则
$|\mathcal{X}|$ 清醒。
\end{proposition}

\begin{proof}
由引理 \ref{stacks-morphisms-lemma-kolmogorov}，$|\mathcal{X}|$ 是柯尔莫哥洛夫空间
（实际上下面会重新证明这一点）。设
$T \subset |\mathcal{X}|$ 是不可约闭子集。必须证明 $T$ 有一般点。
令 $\mathcal{Z} \subset \mathcal{X}$ 为对应于 $T$ 的约化诱导闭子叠；
见《叠的性质》定义
\ref{stacks-properties-definition-reduced-induced-stack}。
由于 $\mathcal{Z} \to \mathcal{X}$ 是闭浸入，由引理
\ref{stacks-morphisms-lemma-representable-decent}，$\mathcal{Z}$ 是适宜代数叠。
此外，态射 $\mathcal{I}_\mathcal{Z} \to \mathcal{Z}$ 是
$\mathcal{I}_\mathcal{X} \to \mathcal{X}$ 的基变换
（引理 \ref{stacks-morphisms-lemma-monomorphism-cartesian-square-inertia}）。因此
$\mathcal{I}_\mathcal{Z} \to \mathcal{Z}$ 拟紧
（引理 \ref{stacks-morphisms-lemma-base-change-quasi-compact}）。由此归结到下一段所述的情形。

\medskip\noindent
假设 $\mathcal{X}$ 适宜，$\mathcal{I}_\mathcal{X} \to \mathcal{X}$
拟紧，$\mathcal{X}$ 约化，并且 $|\mathcal{X}|$ 不可约。必须证明
$|\mathcal{X}|$ 有一般点。由命题 \ref{stacks-morphisms-proposition-open-stratum}，
存在作为胚的稠密开子叠 $\mathcal{U} \subset \mathcal{X}$。
换言之，$|\mathcal{U}| \subset |\mathcal{X}|$ 开且稠密。
因此，除其他所有性质外，还可假设 $\mathcal{X}$ 是胚。
设 $\mathcal{X} \to X$ 使 $\mathcal{X}$ 成为代数空间 $X$ 上的胚。
由引理 \ref{stacks-morphisms-lemma-gerbe-bijection-points}，
$|\mathcal{X}| \cong |X|$。特别地，$X$ 拟紧且 $|X|$ 不可约。
又由引理 \ref{stacks-morphisms-lemma-gerbe-decent}，$X$ 是适宜代数空间。
于是，由《适宜代数空间》命题
\ref{decent-spaces-proposition-reasonable-sober}，
$|\mathcal{X}| = |X|$ 清醒，因而有一个（唯一的）一般点。
\end{proof}






\section{整代数叠}
\label{stacks-morphisms-section-integral-stacks}

\noindent
本节是《域上的空间》第
\ref{spaces-over-fields-section-integral-spaces} 节的类比版本。
受该节中的考察以及命题 \ref{stacks-morphisms-proposition-decent-sober} 的结果启发，
如下定义整代数叠（它与已有的整概形和整代数空间定义并不冲突）。

\begin{definition}
\label{stacks-morphisms-definition-integral-algebraic-stack}
若代数叠 $\mathcal{X}$ 约化、适宜，
$\mathcal{I}_\mathcal{X} \to \mathcal{X}$ 拟紧，并且
$|\mathcal{X}|$ 不可约，则称它是{\it 整的}。
\end{definition}

\noindent
注意，若 $\mathcal{X}$ 拟分离，则要使它为整的，只需
$\mathcal{X}$ 约化且 $|\mathcal{X}|$ 不可约；见引理
\ref{stacks-morphisms-lemma-totaro-case}。

\begin{lemma}
\label{stacks-morphisms-lemma-structure-integral-stack}
设 $\mathcal{X}$ 是整代数叠。则
\begin{enumerate}
\item $|\mathcal{X}|$ 清醒、不可约，并且有唯一的一般点；
\item 存在开子叠 $\mathcal{U} \subset \mathcal{X}$，它是整概形 $U$ 上的胚。
\end{enumerate}
\end{lemma}

\begin{proof}
命题 \ref{stacks-morphisms-proposition-decent-sober} 告诉我们 $|\mathcal{X}|$ 清醒。
它当然也不可约，因而有唯一的一般点 $x$（由清醒性的定义）。
命题 \ref{stacks-morphisms-proposition-open-stratum} 给出稠密开子叠
$\mathcal{U} \subset \mathcal{X}$，它是某个代数空间 $U$ 上的胚。
由引理 \ref{stacks-morphisms-lemma-gerbe-decent}，$U$ 是适宜代数空间（而
$\mathcal{U}$ 由引理 \ref{stacks-morphisms-lemma-representable-decent} 适宜）。
由于 $|U| = |\mathcal{U}|$，$|U|$ 不可约。最后，因为
$\mathcal{U}$ 约化，态射 $\mathcal{U} \to U$ 分解经过 $U_{red}$；
见《叠的性质》引理
\ref{stacks-properties-lemma-map-into-reduction}。现在，由于
$\mathcal{U} \to U$ 平坦、局部有限表现且满
（引理 \ref{stacks-morphisms-lemma-gerbe-fppf}），这蕴含 $U = U_{red}$，即 $U$
约化（略去一个小细节）。于是 $U$ 是整代数空间；见《域上的空间》定义
\ref{spaces-over-fields-definition-integral-algebraic-space}。
最后，可以用开子空间替换 $U$（并相应替换 $\mathcal{U}$），
并假设 $U$ 是整概形；见《域上的空间》第
\ref{spaces-over-fields-section-integral-spaces} 节中的讨论。
\end{proof}

\begin{lemma}
\label{stacks-morphisms-lemma-totaro-case}
设 $\mathcal{X}$ 是约化且拟分离的代数叠，并且其关联拓扑空间
$|\mathcal{X}|$ 不可约。则 $\mathcal{X}$ 是整的。
\end{lemma}

\begin{proof}
若 $\mathcal{X}$ 拟分离，则由引理
\ref{stacks-morphisms-lemma-quasi-separated-decent}，$\mathcal{X}$ 适宜。
若 $\mathcal{X}$ 拟分离，则
$\Delta : \mathcal{X} \to \mathcal{X} \times \mathcal{X}$ 拟紧；
因此，作为 $\Delta$ 沿 $\Delta$ 的基变换，
$\mathcal{I}_\mathcal{X} \to \mathcal{X}$ 拟紧；见引理
\ref{stacks-morphisms-lemma-base-change-quasi-compact}。由此可见，定义
\ref{stacks-morphisms-definition-integral-algebraic-stack} 的所有假设都成立
（同时也看出，可以把“拟分离”替换为“$\Delta_\mathcal{X}$ 拟紧”）。
\end{proof}

\begin{lemma}
\label{stacks-morphisms-lemma-decent-irreducible-closed}
设 $\mathcal{X}$ 是适宜代数叠，并且
$\mathcal{I}_\mathcal{X} \to \mathcal{X}$ 拟紧。
下列集合之间存在典范双射：
\begin{enumerate}
\item $\mathcal{X}$ 的点集，即 $|\mathcal{X}|$；
\item $|\mathcal{X}|$ 的不可约闭子集所成的集合；
\item $\mathcal{X}$ 的整闭子叠所成的集合。
\end{enumerate}
从 (1) 到 (2) 的双射把 $x$ 映到 $\overline{\{x\}}$。
从 (3) 到 (2) 的双射把 $\mathcal{Z}$ 映到 $|\mathcal{Z}|$。
\end{lemma}

\begin{proof}
由命题 \ref{stacks-morphisms-proposition-decent-sober}，$|\mathcal{X}|$ 清醒，
所以我们的映射定义了 (1) 与 (2) 之间的双射。给定闭且不可约的
$T \subset |\mathcal{X}|$，存在唯一的约化闭子叠
$\mathcal{Z} \subset \mathcal{X}$，使得 $|\mathcal{Z}| = T$；
即 $\mathcal{Z}$ 是 $T$ 上的约化诱导子空间结构；见《叠的性质》定义
\ref{stacks-properties-definition-reduced-induced-stack}。
于是 $\mathcal{Z}$ 是整代数叠，因为它适宜
（引理 \ref{stacks-morphisms-lemma-representable-decent}），态射
$\mathcal{I}_\mathcal{Z} \to \mathcal{Z}$ 拟紧
（它是 $\mathcal{I}_\mathcal{X} \to \mathcal{X}$ 的基变换；见引理
\ref{stacks-morphisms-lemma-monomorphism-cartesian-square-inertia}），
$\mathcal{Z}$ 约化，并且 $|\mathcal{Z}|$ 不可约。
\end{proof}










\section{剩余胚}
\label{stacks-morphisms-section-residual-gerbe}

\noindent
本节是《叠的性质》第
\ref{stacks-properties-section-residual-gerbe} 节的继续。

\begin{lemma}
\label{stacks-morphisms-lemma-gerbe-residual-gerbe}
设 $\pi : \mathcal{X} \to \mathcal{Y}$ 是代数叠的态射。设
$x \in |\mathcal{X}|$ 的像为 $y \in |\mathcal{Y}|$。假设
$\mathcal{Y}$ 在 $y$ 处的剩余胚
$\mathcal{Z}_y \subset \mathcal{Y}$ 存在，并且 $\mathcal{X}$ 是
$\mathcal{Y}$ 上的胚。则
$\mathcal{Z}_x = \mathcal{Z}_y \times_\mathcal{Y} \mathcal{X}$ 是
$\mathcal{X}$ 在 $x$ 处的剩余胚。
\end{lemma}

\begin{proof}
态射 $\mathcal{Z}_x \to \mathcal{X}$ 是单态射，因为它是单态射
$\mathcal{Z}_y \to \mathcal{Y}$ 的基变换。由引理
\ref{stacks-morphisms-lemma-gerbe-bijection-points}，态射 $\pi$ 是泛同胚，因而
$|\mathcal{Z}_x| = \{x\}$。最后，态射
$\mathcal{Z}_x \to \mathcal{Z}_y$ 是光滑态射 $\pi$ 的基变换，
所以它光滑；见引理 \ref{stacks-morphisms-lemma-gerbe-smooth}。因此，由于
$\mathcal{Z}_y$ 约化且局部诺特，$\mathcal{Z}_x$ 也如此
（细节略）。于是，由《叠的性质》定义
\ref{stacks-properties-definition-residual-gerbe}，$\mathcal{Z}_x$ 是
$\mathcal{X}$ 在 $x$ 处的剩余胚。
\end{proof}

\begin{lemma}
\label{stacks-morphisms-lemma-factor-through-residual-space}
设 $f : \mathcal{Y} \to \mathcal{X}$ 是代数叠的态射。
设 $x \in |\mathcal{X}|$ 是一点。假设
\begin{enumerate}
\item $\mathcal{X}$ 适宜或局部诺特（或两者皆是）；
\item $\mathcal{I}_\mathcal{X} \to \mathcal{X}$ 拟紧；
\item $|f|(|\mathcal{Y}|)$ 包含在 $\{x\} \subset |\mathcal{X}|$ 中；
\item $\mathcal{Y}$ 约化。
\end{enumerate}
则 $f$ 分解经过 $\mathcal{X}$ 在 $x$ 处的剩余胚 $\mathcal{Z}_x$
（其存在性由引理 \ref{stacks-morphisms-lemma-every-point-residual-gerbe} 或
\ref{stacks-morphisms-lemma-every-point-residual-gerbe-locally-Noetherian}).
保证）。
\end{lemma}

\begin{proof}
令 $T = \overline{\{x\}} \subset |\mathcal{X}|$ 为 $x$ 的闭包。
由《叠的性质》引理
\ref{stacks-properties-lemma-reduced-closed-substack}
，存在约化闭子叠 $\mathcal{X}' \subset \mathcal{X}$，使得
$T = |\mathcal{X}'|$。由《叠的性质》引理
\ref{stacks-properties-lemma-map-into-reduction}，态射 $f$ 分解经过
$\mathcal{X}'$。若 $\mathcal{X}$ 适宜，则由引理
\ref{stacks-morphisms-lemma-representable-decent}，叠 $\mathcal{X}'$ 适宜。
若 $\mathcal{X}$ 局部诺特，则 $\mathcal{X}'$ 局部诺特（细节略）。
注意，$\mathcal{I}_{\mathcal{X}'} \to \mathcal{X}'$ 是
$\mathcal{I}_\mathcal{X} \to \mathcal{X}$ 的基变换；由引理
\ref{stacks-morphisms-lemma-monomorphism-cartesian-square-inertia} 和
\ref{stacks-morphisms-lemma-base-change-quasi-compact}，
$\mathcal{I}_{\mathcal{X}'} \to \mathcal{X}'$ 拟紧。
由此归结到下一段所述的情形。

\medskip\noindent
假设 $\mathcal{X}$ 约化且 $x \in |\mathcal{X}|$ 是一般点。
由命题 \ref{stacks-morphisms-proposition-open-stratum}，存在作为胚的稠密开子叠
$\mathcal{U} \subset \mathcal{X}'$。注意 $x \in |\mathcal{U}|$。
重复上述论证，归结到下一段所述的情形。

\medskip\noindent
假设 $\mathcal{X} \to X$ 是代数空间 $X$ 上的胚。若
$\mathcal{X}$ 适宜，则由引理 \ref{stacks-morphisms-lemma-gerbe-bijection-points} 和
\ref{stacks-morphisms-lemma-qc-decent}，空间 $X$ 适宜。若 $\mathcal{X}$ 局部诺特，
则由 fppf 下降，$X$ 局部诺特（细节略）。因此相应结论对 $X$ 成立；
见《适宜代数空间》引理
\ref{decent-spaces-lemma-factor-through-residual-space-decent} 或
\ref{decent-spaces-lemma-factor-through-residual-space-Noetherian}
（略去一个小细节）。应用引理 \ref{stacks-morphisms-lemma-gerbe-residual-gerbe}，
可知结论对 $\mathcal{X}$ 也成立。
\end{proof}

\begin{remark}
\label{stacks-morphisms-lemma-factor-through-residual-space-Noetherian}
我们不知道仅假设 $\mathcal{X}$ 局部诺特时，引理
\ref{stacks-morphisms-lemma-factor-through-residual-space} 是否仍成立；换言之，
这里去掉关于惯性态射拟紧的假设。在这种情形下，若 $x$ 是闭点，
则下面这个简单得多的引理表明结论当然成立。
\end{remark}

\begin{lemma}
\label{stacks-morphisms-lemma-residual-space-closed}
设 $\mathcal{X}$ 是局部诺特代数叠。设 $x \in |\mathcal{X}|$ 的剩余胚为
$\mathcal{Z}_x \subset \mathcal{X}$
（引理 \ref{stacks-morphisms-lemma-every-point-residual-gerbe-locally-Noetherian}）。
则 $x$ 是 $|\mathcal{X}|$ 的闭点，当且仅当态射
$\mathcal{Z}_x \to \mathcal{X}$ 是闭浸入。
\end{lemma}

\begin{proof}
若 $\mathcal{Z}_x \to \mathcal{X}$ 是闭浸入，则 $x$ 是
$|\mathcal{X}|$ 的闭点；例如见引理 \ref{stacks-morphisms-lemma-closed-immersion-proper}。
反过来，假设 $x$ 是 $|\mathcal{X}|$ 的闭点。令
$\mathcal{Z} \subset \mathcal{X}$ 为满足 $|Z| = \{x\}$ 的约化闭子叠
（《叠的性质》引理
\ref{stacks-properties-lemma-reduced-closed-substack}）。由引理
\ref{stacks-morphisms-lemma-immersion-locally-finite-type} 和
\ref{stacks-morphisms-lemma-locally-finite-type-locally-noetherian}，$\mathcal{Z}$
是局部诺特代数叠。又因 $\mathcal{Z}$ 约化且
$|\mathcal{Z}| = \{x\}$，按定义，$\mathcal{Z} = \mathcal{Z}_x$
是剩余胚。
\end{proof}
% Generated mechanically from the frozen stacks-limits.tex target.
% Do not edit this generated file; edit the bound locale source instead.

% OK, start here.
%

\setcounter{chapter}{101}
\chapter{代数叠的极限}



\phantomsection
\label{stacks-limits-section-phantom}


\section{引言}
\label{stacks-limits-section-introduction}

\noindent
本章汇集与代数叠的极限有关的材料。
关于代数叠和代数空间极限的许多结果由 David Rydh
在 \cite{rydh_approx} 中得到。


\section{约定}
\label{stacks-limits-section-conventions}

\noindent
我们继续采用《叠的性质》第
\ref{stacks-properties-section-conventions} 节中引入的约定和语言滥用。







\section{有限表示态射}
\label{stacks-limits-section-finite-presentation}

\noindent
本节是《空间的极限》第
\ref{spaces-limits-section-finite-presentation} 节的对应版本。
在那里，我们定义了 $\Sch$ 上函子之间的变换保持极限的含义
（建议参看《空间的极限》引理
\ref{spaces-limits-lemma-characterize-relative-limit-preserving}
中的刻画）。在《可表示性判据》第
\ref{criteria-section-limit-preserving} 节中，我们定义了
“在对象上保持极限”的概念。回顾《Artin 公理》第
\ref{artin-section-limits} 节，我们已经定义了 $\Sch$
上的群胚纤维化范畴保持极限的含义。将这些定义结合起来，
得到如下概念。

\begin{definition}
\label{stacks-limits-definition-limit-preserving}
设 $S$ 为概形，并设 $f : \mathcal{X} \to \mathcal{Y}$ 是
$(\Sch/S)_{fppf}$ 上群胚纤维化范畴之间的一个 $1$-态射。
若对 $S$ 上仿射概形的每个有向极限 $U = \lim U_i$，纤维范畴图
$$
\xymatrix{
\colim \mathcal{X}_{U_i} \ar[r] \ar[d]_f & \mathcal{X}_U \ar[d]^f \\
\colim \mathcal{Y}_{U_i} \ar[r] & \mathcal{Y}_U
}
$$
都是 $2$-笛卡儿的，则称 $f$ {\it 保持极限}。
\end{definition}

\begin{lemma}
\label{stacks-limits-lemma-limit-preserving-objects}
设 $S$ 为概形，并设 $f : \mathcal{X} \to \mathcal{Y}$ 是
$(\Sch/S)_{fppf}$ 上群胚纤维化范畴之间的一个 $1$-态射。
若 $f$ 保持极限（定义 \ref{stacks-limits-definition-limit-preserving}），
则 $f$ 在对象上保持极限（《可表示性判据》第
\ref{criteria-section-limit-preserving} 节）。
\end{lemma}

\begin{proof}
% P12-ERR-0001: see ../control/ERRATA_CANDIDATES.jsonl
若对 $U$ 上仿射概形的每个有向极限 $U = \lim U_i$，函子
$$
\colim \mathcal{X}_{U_i} \longrightarrow
(\colim \mathcal{Y}_{U_i}) \times_{\mathcal{Y}_U} \mathcal{X}_U
$$
都是本质满的，则 $f$ 在对象上保持极限。
\end{proof}

\begin{lemma}
\label{stacks-limits-lemma-base-change-limit-preserving}
设 $p : \mathcal{X} \to \mathcal{Y}$ 和
$q : \mathcal{Z} \to \mathcal{Y}$ 是 $(\Sch/S)_{fppf}$
上群胚纤维化范畴之间的 $1$-态射。若
$p : \mathcal{X} \to \mathcal{Y}$ 保持极限，则 $p$ 沿 $q$
的基变换
$p' : \mathcal{X} \times_\mathcal{Y} \mathcal{Z} \to \mathcal{Z}$
也保持极限。
\end{lemma}

\begin{proof}
这是形式推论。设 $U = \lim_{i \in I} U_i$ 是 $S$ 上仿射概形
$U_i$ 的有向极限。对每个 $i$，有
$$
(\mathcal{X} \times_\mathcal{Y} \mathcal{Z})_{U_i} =
\mathcal{X}_{U_i} \times_{\mathcal{Y}_{U_i}} \mathcal{Z}_{U_i}
$$
滤过余极限与范畴的 $2$-纤维积可交换（略去细节），
因此若 $p$ 保持极限，就得到
\begin{align*}
\colim (\mathcal{X} \times_\mathcal{Y} \mathcal{Z})_{U_i}
& =
\colim \mathcal{X}_{U_i} \times_{\colim \mathcal{Y}_{U_i}}
\colim \mathcal{Z}_{U_i} \\
& =
\mathcal{X}_U \times_{\mathcal{Y}_U} \colim \mathcal{Y}_{U_i}
\times_{\colim \mathcal{Y}_{U_i}}
\colim \mathcal{Z}_{U_i} \\
& =
\mathcal{X}_U \times_{\mathcal{Y}_U} \colim \mathcal{Z}_{U_i} \\
& =
\mathcal{X}_U \times_{\mathcal{Y}_U} \mathcal{Z}_U \times_{\mathcal{Z}_U}
\colim \mathcal{Z}_{U_i} \\
& =
(\mathcal{X} \times_\mathcal{Y} \mathcal{Z})_U \times_{\mathcal{Z}_U}
\colim \mathcal{Z}_{U_i}
\end{align*}
这正是所需结论。
\end{proof}

\begin{lemma}
\label{stacks-limits-lemma-composition-limit-preserving}
设 $p : \mathcal{X} \to \mathcal{Y}$ 和
$q : \mathcal{Y} \to \mathcal{Z}$ 是 $(\Sch/S)_{fppf}$
上群胚纤维化范畴之间的 $1$-态射。若 $p$ 和 $q$ 都保持极限，
则复合 $q \circ p$ 也保持极限。
\end{lemma}

\begin{proof}
这是形式推论。设 $U = \lim_{i \in I} U_i$ 是 $S$ 上仿射概形
$U_i$ 的有向极限。若 $p$ 和 $q$ 保持极限，则
\begin{align*}
\colim \mathcal{X}_{U_i}
& =
\mathcal{X}_U \times_{\mathcal{Y}_U} \colim \mathcal{Y}_{U_i} \\
& =
\mathcal{X}_U \times_{\mathcal{Y}_U} \mathcal{Y}_U
\times_{\mathcal{Z}_U} \colim \mathcal{Z}_{U_i} \\
& =
\mathcal{X}_U \times_{\mathcal{Z}_U} \colim \mathcal{Z}_{U_i}
\end{align*}
这正是所需结论。
\end{proof}

\begin{lemma}
\label{stacks-limits-lemma-representable-by-spaces-limit-preserving}
设 $p : \mathcal{X} \to \mathcal{Y}$ 是 $(\Sch/S)_{fppf}$
上群胚纤维化范畴之间的一个 $1$-态射。若 $p$ 可由代数空间表示，
则下列条件等价：
\begin{enumerate}
\item $p$ 保持极限；
\item $p$ 在对象上保持极限；
\item $p$ 局部有限表示（参见《代数叠》定义
\ref{algebraic-definition-relative-representable-property}）。
\end{enumerate}
\end{lemma}

\begin{proof}
由《可表示性判据》引理
\ref{criteria-lemma-representable-by-spaces-limit-preserving}，
我们已经知道 (2) 与 (3) 等价。因此只需证明 (1) 与 (2) 等价。
其中一个方向已见引理 \ref{stacks-limits-lemma-limit-preserving-objects}。
为证另一方向，设 $U = \lim_{i \in I} U_i$ 是 $S$ 上仿射概形
$U_i$ 的有向极限。我们必须证明
$$
\colim \mathcal{X}_{U_i} \longrightarrow
\mathcal{X}_U \times_{\mathcal{Y}_U} \colim \mathcal{Y}_{U_i}
$$
是等价。由假设 (2)，它是本质满的，故只需证明其全忠实性。
由于 $p$ 在纤维范畴上忠实（《代数叠》引理
\ref{algebraic-lemma-criterion-map-representable-spaces-fibred-in-groupoids})
，可见该函子忠实。设 $x_i$ 和 $x'_i$ 是 $\mathcal{X}$ 在
$U_i$ 上的纤维范畴中的对象。上述函子将 $x_i$ 映为
$(x_i|_U, p(x_i), can)$，其中 $can$ 是典范同构
$p(x_i|_U) \to p(x_i)|_U$。于是，设给定态射
$$
(\alpha, \beta_i) : (x_i|_U, p(x_i), can) \longrightarrow
(x'_i|_U, p(x'_i), can)
$$
，它位于本证明第一个陈列箭头右端的范畴中。我们的任务是构造
$i' \geq i$ 以及一个态射
$x_i|_{U_{i'}} \to x'_i|_{U_{i'}}$，使其映到
$(\alpha, \beta_i|_{U_{i'}})$。

\medskip\noindent
令 $y_i = p(x_i)$、$y'_i = p(x'_i)$。由《代数叠》引理
\ref{algebraic-lemma-criterion-map-representable-spaces-fibred-in-groupoids})
，函子
$$ 
X_{y_i} : (\Sch/U_i)^{opp} \to \textit{Sets},\quad
V/U_i \mapsto
% P12-ERR-0002: see ../control/ERRATA_CANDIDATES.jsonl
\{(x, \phi) \mid x \in \Ob(\mathcal{X}_V), \phi : f(x) \to y_i|V\}/\cong
$$
是 $U_i$ 上的代数空间；类似定义的函子 $X_{y'_i}$ 亦然。
由于 (2) 与 (3) 等价，$X_{y'_i}$ 在 $U_i$ 上局部有限表示。
注意，$(x_i, \text{id})$ 和 $(x'_i, \text{id})$
分别定义了 $X_{y_i}$ 和 $X_{y'_i}$ 的 $U_i$-值点。
存在函子变换
$$
\beta_i : X_{y_i} \to X_{y'_i},\quad
(x/V, \phi) \mapsto (x/V, \beta_i|_V \circ \phi)
$$
；换言之，这是 $U_i$ 上代数空间之间的一个态射。我们断言
$$
\xymatrix{
U \ar[d] \ar[rr] & & U_i \ar[d]^{(x'_i, \text{id})} \\
U_i \ar[r]^{(x_i, \text{id})} & X_{y_i} \ar[r]^{\beta_i} & X_{y'_i}
}
$$
交换。事实上，这等价于函子 $X_{y'_i}$ 定义中的两对
$(x_i|_U, \beta_i|_U)$ 与 $(x'_i|_U, \text{id})$ 同构；
而态射 $\alpha : x_i|_U \to x'_i|_U$ 恰好给出这样的同构。
反向推理可见，若能找到 $i' \geq i$，使得图
$$
\xymatrix{
U_{i'} \ar[d] \ar[rr] & & U_i \ar[d]^{(x'_i, \text{id})} \\
U_i \ar[r]^{(x_i, \text{id})} & X_{y_i} \ar[r]^{\beta_i} & X_{y'_i}
}
$$
交换，那么就得到同构 $x_i|_{U_{i'}} \to x'_i|_{U_{i'}}$，
它解决了上一段提出的问题。另一方面，对角态射
$$
\Delta : X_{y'_i} \to X_{y'_i} \times_{U_i} X_{y'_i}
$$
局部有限表示（《空间的态射》引理
\ref{spaces-morphisms-lemma-diagonal-morphism-finite-type})
。因此，$U \to U_i$ 使到 $X_{y'_i}$ 的两个态射相等这一事实，
意味着对某个 $i' \geq i$，态射 $U_{i'} \to U_i$
已使这两个态射相等；参见《空间的极限》命题
\ref{spaces-limits-proposition-characterize-locally-finite-presentation}。
\end{proof}

\begin{lemma}
\label{stacks-limits-lemma-limit-preserving-diagonal}
设 $p : \mathcal{X} \to \mathcal{Y}$ 是 $(\Sch/S)_{fppf}$
上群胚纤维化范畴之间的一个 $1$-态射。下列条件等价：
\begin{enumerate}
\item 对角态射
$\Delta : \mathcal{X} \to \mathcal{X} \times_\mathcal{Y} \mathcal{X}$
保持极限；
\item 对 $S$ 上仿射概形的每个有向极限 $U = \lim U_i$，函子
$$
\colim \mathcal{X}_{U_i} \longrightarrow
\mathcal{X}_U \times_{\mathcal{Y}_U} \colim \mathcal{Y}_{U_i}
$$
是全忠实的。
\end{enumerate}
特别地，若 $p$ 保持极限，则 $\Delta$ 也保持极限。
\end{lemma}

\begin{proof}
设 $U = \lim U_i$ 是 $S$ 上仿射概形的有向极限。我们断言函子
$$
\colim \mathcal{X}_{U_i} \longrightarrow
\mathcal{X}_U \times_{\mathcal{Y}_U} \colim \mathcal{Y}_{U_i}
$$
全忠实，当且仅当函子
$$
\colim \mathcal{X}_{U_i} \longrightarrow
\mathcal{X}_U \times_{(\mathcal{X} \times_\mathcal{Y} \mathcal{X})_U}
\colim (\mathcal{X} \times_\mathcal{Y} \mathcal{X})_{U_i}
$$
是等价。这将证明本引理。由于
$(\mathcal{X} \times_\mathcal{Y} \mathcal{X})_U =
\mathcal{X}_U \times_{\mathcal{Y}_U} \mathcal{X}_U$
并且
$(\mathcal{X} \times_\mathcal{Y} \mathcal{X})_{U_i} =
\mathcal{X}_{U_i} \times_{\mathcal{Y}_{U_i}} \mathcal{X}_{U_i}$
，这纯粹是一个范畴论断言；下一段将予讨论。

\medskip\noindent
设 $\mathcal{I}$ 为滤过指标范畴，并设
$(\mathcal{C}_i)$ 和 $(\mathcal{D}_i)$ 是 $\mathcal{I}$
上的群胚系统。设 $p : (\mathcal{C}_i) \to (\mathcal{D}_i)$
为 $\mathcal{I}$ 上群胚系统之间的映射。再设有群胚函子
$p : \mathcal{C} \to \mathcal{D}$ 以及函子
$f : \colim \mathcal{C}_i \to \mathcal{C}$ 和
$g : \colim \mathcal{D}_i \to \mathcal{D}$，它们构成交换图
$$
\xymatrix{
\colim \mathcal{C}_i \ar[d]_p \ar[r]_f & \mathcal{C} \ar[d]^p \\
\colim \mathcal{D}_i \ar[r]^g & \mathcal{D}
}
$$
我们断言，
$$
A : \colim \mathcal{C}_i \longrightarrow
\mathcal{C} \times_\mathcal{D} \colim \mathcal{D}_i
$$
全忠实，当且仅当函子
$$
B : \colim \mathcal{C}_i \longrightarrow
\mathcal{C}
\times_{\Delta, \mathcal{C} \times_\mathcal{D} \mathcal{C}, f \times_g f}
\colim (\mathcal{C}_i \times_{\mathcal{D}_i} \mathcal{C}_i)
$$
是等价。令 $\mathcal{C}' = \colim \mathcal{C}_i$、
$\mathcal{D}' = \colim \mathcal{D}_i$。由于 $2$-纤维积与滤过余极限
可交换，$A$ 和 $B$ 分别成为函子
$$
A' : \mathcal{C}' \to \mathcal{C} \times_\mathcal{D} \mathcal{D}'
\quad\text{且}\quad
B' : \mathcal{C}' \longrightarrow
\mathcal{C}
\times_{\Delta, \mathcal{C} \times_\mathcal{D} \mathcal{C}, f \times_g f}
(\mathcal{C}' \times_{\mathcal{D}'} \mathcal{C}')
$$
因此只需证明：若
$$
\xymatrix{
\mathcal{C}' \ar[d]_p \ar[r]_f & \mathcal{C} \ar[d]^p \\
\mathcal{D}' \ar[r]^g & \mathcal{D}
}
$$
是群胚的交换图，则 $A'$ 全忠实，当且仅当 $B'$ 是等价。
这可由《范畴》引理
\ref{categories-lemma-fully-faithful-diagonal-equivalence}
（取平凡的，即单点的基范畴）推出，因为
$$
\mathcal{C}
\times_{\Delta, \mathcal{C} \times_\mathcal{D} \mathcal{C}, f \times_g f}
(\mathcal{C}' \times_{\mathcal{D}'} \mathcal{C}') =
\mathcal{C}'
\times_{A', \mathcal{C} \times_\mathcal{D} \mathcal{D}', A'}
\mathcal{C}'
$$
证明完毕。
\end{proof}

\begin{lemma}
\label{stacks-limits-lemma-locally-finite-presentation-limit-preserving}
设 $S$ 为概形，$\mathcal{X}$ 为 $S$ 上的代数叠。
若 $\mathcal{X} \to S$ 局部有限表示，则 $\mathcal{X}$
在《Artin 公理》定义 \ref{artin-definition-limit-preserving}
的意义下保持极限（等价地，态射 $\mathcal{X} \to S$ 保持极限）。
\end{lemma}

\begin{proof}
选取某个概形 $U$ 及一个光滑满射 $U \to \mathcal{X}$。
则 $U \to S$ 局部有限表示，参见《叠的态射》第
\ref{stacks-morphisms-section-finite-presentation} 节。
由《代数叠》引理 \ref{algebraic-lemma-stack-presentation}，
对代数空间中的某个光滑群胚 $(U, R, s, t, c)$，可写成
$\mathcal{X} = [U/R]$。由于 $U$ 在 $S$ 上局部有限表示，
代数空间 $R$ 在 $S$ 上也局部有限表示。
回顾，$[U/R]$ 是将如下群胚纤维化范畴叠化而得到的
$(\Sch/S)_{fppf}$ 上的群胚叠：它在 $T$ 上的纤维范畴是群胚
$(U(T), R(T), s, t, c)$。由于作为函子，$U$ 和 $R$ 都保持极限
（《空间的极限》命题
\ref{spaces-limits-proposition-characterize-locally-finite-presentation})
，该群胚纤维化范畴保持极限。因此，只需证明 fppf 叠化保留
保持极限这一性质。事实确实如此（提示：使用《拓扑》引理
\ref{topologies-lemma-limit-fppf-topology}).
不过，利用本情形中叠化的具体描述，我们在下面给出直接证明。

\medskip\noindent
设 $T = \lim T_\lambda$ 是 $S$ 上仿射概形的有向极限。
我们必须证明函子
$$
\colim [U/R]_{T_\lambda} \longrightarrow [U/R]_T
$$
是范畴等价。先证明该函子本质满。设 $x \in \Ob([U/R]_T)$。
《空间中的群胚》引理
\ref{spaces-groupoids-lemma-quotient-stack-objects}
给出了范畴 $[U/R]_T$ 的描述。特别地，$x$ 对应于一个 fppf 覆盖
$\{T_i \to T\}_{i \in I}$ 以及相对于该覆盖的一个
$[U/R]$-下降数据 $(u_i, r_{ij})$。细化此覆盖后，可设它是仿射概形
$T$ 的标准 fppf 覆盖。由《拓扑》引理
\ref{topologies-lemma-limit-fppf-topology}
，可选取某个 $\lambda$ 和标准 fppf 覆盖
$\{T_{\lambda, i} \to T_\lambda\}_{i \in I}$，使其到 $T$
的基变换等于 $\{T_i \to T\}_{i \in I}$。
对每个 $i$，增大 $\lambda$ 后，可找到
$u_{\lambda, i} : T_{\lambda, i} \to U$，使其与
$T_i \to T_{\lambda, i}$ 的复合为给定态射 $u_i$
（此处使用 $U$ 保持极限）。类似地，对每个 $i, j$，增大
$\lambda$ 后，可找到
$r_{\lambda, ij} : T_{\lambda, i} \times_{T_\lambda} T_{\lambda, j} \to R$，
使其与 $T_{ij} \to T_{\lambda, ij}$ 的复合为给定态射 $r_{ij}$
（此处使用 $R$ 保持极限）。进一步增大 $\lambda$ 后，可设
$$
s \circ r_{\lambda, ij} = u_{\lambda, i} \circ \text{pr}_0
\quad\text{且}\quad
t \circ r_{\lambda, ij} = u_{\lambda, j} \circ \text{pr}_1,
$$
并且
$$
c \circ (r_{\lambda, jk} \circ \text{pr}_{12}, r_{\lambda, ij}
\circ \text{pr}_{01}) = r_{\lambda, ik} \circ \text{pr}_{02}.
$$
换言之，可设 $(u_{\lambda, i}, r_{\lambda, ij})$ 是相对于覆盖
$\{T_{\lambda, i} \to T_\lambda\}_{i \in I}$ 的一个
$[U/R]$-下降数据。于是得到 $T_\lambda$ 上 $[U/R]$ 的相应对象，
其到 $T$ 的拉回与 $x$ 同构，正合所需。
全忠实性的证明完全相同，只需使用《空间中的群胚》引理
\ref{spaces-groupoids-lemma-quotient-stack-objects}
% P12-ERR-0003: see ../control/ERRATA_CANDIDATES.jsonl
给出的 $[U/T]$ 的纤维范畴中态射的描述。
\end{proof}

\begin{proposition}
\label{stacks-limits-proposition-characterize-locally-finite-presentation}
\begin{reference}
这是 \cite[Lemma 2.3.15]{Emerton-Gee} 的一个特例。
\end{reference}
设 $f : \mathcal{X} \to \mathcal{Y}$ 为代数叠之间的态射。
下列条件等价：
\begin{enumerate}
\item $f$ 保持极限；
\item $f$ 在对象上保持极限；
\item $f$ 局部有限表示。
\end{enumerate}
\end{proposition}

\begin{proof}
假设 (3) 成立。设 $T = \lim T_i$ 是仿射概形的有向极限。
考虑函子
$$
\colim \mathcal{X}_{T_i} \longrightarrow
\mathcal{X}_T \times_{\mathcal{Y}_T} \colim \mathcal{Y}_{T_i}
$$
设 $(x, y_i, \beta)$ 是右端的对象，即
$x \in \Ob(\mathcal{X}_T)$、$y_i \in \Ob(\mathcal{Y}_{T_i})$，
且 $\beta : f(x) \to y_i|_T$ 是 $\mathcal{Y}_T$ 中的态射。
于是可将 $(x, y_i, \beta)$ 视为 $T$ 上代数叠
$\mathcal{X}_{y_i} = \mathcal{X} \times_{\mathcal{Y}, y_i} T_i$
的一个对象。由于 $\mathcal{X}_{y_i} \to T_i$ 是 $f$ 的基变换，
故它局部有限表示；由引理
\ref{stacks-limits-lemma-locally-finite-presentation-limit-preserving}，
它保持极限。这意味着，对某个 $i' \geq i$，
$(x, y_i, \beta)$ 来自 $T_{i'}$ 上的对象。展开定义即知，
$(x, y_i, \beta)$ 属于所陈列函子的本质像。
换言之，该函子本质满；也就是说，$f$ 在对象上保持极限。
现在将此结论应用于 $f$ 的对角态射 $\Delta$。由《叠的态射》引理
\ref{stacks-morphisms-lemma-diagonal-morphism-finite-type}
，态射 $\Delta$ 局部有限表示。因此上述论证表明 $\Delta$
在对象上保持极限。由引理
\ref{stacks-limits-lemma-representable-by-spaces-limit-preserving}，
$\Delta$ 保持极限。再由引理 \ref{stacks-limits-lemma-limit-preserving-diagonal}，
上述陈列函子全忠实。它又已被证明本质满，故是等价，从而 (1) 成立。

\medskip\noindent
蕴含关系 (1) $\Rightarrow$ (2) 显然。现假设 (2)。
选取概形 $V$ 及光滑满射 $V \to \mathcal{Y}$。
由《可表示性判据》引理
\ref{criteria-lemma-base-change-limit-preserving}
，基变换 $\mathcal{X} \times_\mathcal{Y} V \to V$
在对象上保持极限。选取概形 $U$ 及光滑满射
$U \to \mathcal{X} \times_\mathcal{Y} V$。
由于光滑态射局部有限表示，证明的第一部分说明
$U \to \mathcal{X} \times_\mathcal{Y} V$ 保持极限。
由《可表示性判据》引理
\ref{criteria-lemma-composition-limit-preserving}
，复合 $U \to V$ 在对象上保持极限。因此 $U \to V$
局部有限表示，参见《可表示性判据》引理
\ref{criteria-lemma-representable-by-spaces-limit-preserving}。
这恰好是 $f$ 局部有限表示的条件，参见《叠的态射》定义
\ref{stacks-morphisms-definition-locally-finite-presentation}。
\end{proof}




\section{性质的下降}
\label{stacks-limits-section-descent}

\noindent
本节是《极限》第 \ref{limits-section-descent} 节的对应版本。

\begin{situation}
\label{stacks-limits-situation-descent}
设 $Y = \lim_{i \in I} Y_i$ 是一个代数空间有向系统的极限，
其转移态射均为仿射态射。假设对每个 $i \in I$，$X_i$
% P12-ERR-0004: see ../control/ERRATA_CANDIDATES.jsonl
都是拟紧且拟分离的。另选定一个元素 $0 \in I$。
\end{situation}

\begin{lemma}
\label{stacks-limits-lemma-eventually-separated}
在情形 \ref{stacks-limits-situation-descent} 中，设
$\mathcal{X}_0 \to Y_0$ 是从代数叠到 $Y_0$ 的态射，
并设 $\mathcal{X}_0$ 拟紧且拟分离。若
$Y \times_{Y_0} \mathcal{X}_0 \to Y$ 是分离态射，则对所有充分大的
$i \in I$，$Y_i \times_{Y_0} \mathcal{X}_0 \to Y_i$ 都是分离态射。
\end{lemma}

\begin{proof}
记 $\mathcal{X} = Y \times_{Y_0} \mathcal{X}_0$、
$\mathcal{X}_i = Y_i \times_{Y_0} \mathcal{X}_0$。
选取仿射概形 $U_0$ 及光滑满射 $U_0 \to \mathcal{X}_0$。
令 $U = Y \times_{Y_0} U_0$、$U_i = Y_i \times_{Y_0} U_0$。
则 $U$ 和 $U_i$ 均仿射，且 $U \to \mathcal{X}$ 和
$U_i \to \mathcal{X}_i$ 均为光滑满射。令
$R_0 = U_0 \times_{\mathcal{X}_0} U_0$、
$R = Y \times_{Y_0} R_0$、$R_i = Y_i \times_{Y_0} R_0$。
于是 $R = U \times_\mathcal{X} U$ 且
$R_i = U_i \times_{\mathcal{X}_i} U_i$。

\medskip\noindent
在此记号下，若 $\mathcal{X} \to Y$ 分离，则
$R \to U \times_Y U$ 是对角态射
$\mathcal{X} \to \mathcal{X} \times_Y \mathcal{X}$
沿 $U \times_Y U \to \mathcal{X} \times_Y \mathcal{X}$ 的基变换，
故为固有态射。反之，若 $R_i \to U_i \times_{Y_i} U_i$ 固有，
则 $\mathcal{X}_i \to Y_i$ 分离，因为
$U_i \times_{Y_i} U_i \to \mathcal{X}_i \times_{Y_i} \mathcal{X}_i$
是光滑满射；参见《叠的性质》引理
\ref{stacks-properties-lemma-check-property-covering}。
注意，$R_0 \to U_0 \times_{Y_0} U_0$ 局部有限型，且 $R_0$
拟紧、拟分离。由《空间的极限》引理
\ref{spaces-limits-lemma-eventually-proper}，对充分大的 $i$，
$R_i \to U_i \times_{Y_i} U_i$ 固有，证明完毕。
\end{proof}







\section{相对对象的下降}
\label{stacks-limits-section-descending-relative}

\noindent
本节是《空间的极限》第
\ref{spaces-limits-section-descending-relative} 节的对应版本。

\begin{lemma}
\label{stacks-limits-lemma-descend-a-stack-down}
设 $I$ 为有向集，$(X_i, f_{ii'})$ 为以 $I$ 为指标的代数空间逆系统。
假设
\begin{enumerate}
\item 态射 $f_{ii'} : X_i \to X_{i'}$ 均为仿射态射；
\item 空间 $X_i$ 均拟紧且拟分离。
\end{enumerate}
令 $X = \lim X_i$。若 $\mathcal{X}$ 是 $X$ 上有限表示的代数叠，
则存在 $i \in I$ 及 $X_i$ 上有限表示的代数叠 $\mathcal{X}_i$，
使得作为 $X$ 上的代数叠有
$\mathcal{X} \cong \mathcal{X}_i \times_{X_i} X$。
\end{lemma}

\begin{proof}
由《叠的态射》定义
\ref{stacks-morphisms-definition-locally-finite-presentation}
，态射 $\mathcal{X} \to X$ 拟紧、局部有限表示且拟分离。
由于 $X$ 拟紧且 $\mathcal{X} \to X$ 拟紧，$\mathcal{X}$ 也拟紧
（《叠的态射》定义
\ref{stacks-morphisms-definition-quasi-compact}）。
因此，可找到仿射概形 $U$ 及光滑满射
$U \to \mathcal{X}$（《叠的性质》引理
\ref{stacks-properties-lemma-quasi-compact-stack}）。
令 $R = U \times_\mathcal{X} U$。由《代数叠》引理
\ref{algebraic-lemma-stack-presentation}，得到 $X$ 上代数空间中的
光滑群胚 $(U, R, s, t, c)$，使得 $\mathcal{X} = [U/R]$。
由于 $\mathcal{X} \to X$ 拟分离且 $X$ 拟分离，
$\mathcal{X}$ 拟分离（《叠的态射》引理
\ref{stacks-morphisms-lemma-composition-separated}）。
故 $R \to U \times U$ 拟紧且拟分离（《叠的态射》引理
\ref{stacks-morphisms-lemma-fibre-product-after-map}），
从而 $R$ 是拟分离、拟紧的代数空间。另一方面，$U \to X$
局部有限表示，因此 $R \to X$ 也局部有限表示
（因为 $s : R \to U$ 光滑，故局部有限表示）。
于是 $(U, R, s, t, c)$ 是 $X$ 上有限表示代数空间范畴中的群胚对象。
由《空间的极限》引理
\ref{spaces-limits-lemma-descend-finite-presentation}
，存在某个 $i$ 及 $X_i$ 上代数空间中的群胚
$(U_i, R_i, s_i, t_i, c_i)$，其到 $X$ 的拉回与
$(U, R, s, t, c)$ 同构。增大 $i$ 后，可设 $s_i$ 和 $t_i$
均光滑，参见《空间的极限》引理
\ref{spaces-limits-lemma-descend-smooth}。商叠
$\mathcal{X}_i = [U_i/R_i]$ 是代数叠（《代数叠》定理
\ref{algebraic-theorem-smooth-groupoid-gives-algebraic-stack}）。

\medskip\noindent
由《空间中的群胚》引理
\ref{spaces-groupoids-lemma-quotient-stack-functorial}，
有态射 $[U/R] \to [U_i/R_i]$。我们断言，把它与态射
$[U/R] \to X$ 和 $[U_i/R_i] \to X_i$
（《空间中的群胚》引理
\ref{spaces-groupoids-lemma-quotient-stack-arrows}）结合起来，
得到同构（即等价）
$$
[U/R] \longrightarrow [U_i/R_i] \times_{X_i} X
$$
相应的映射
$$
[U/_{\!p}R] \longrightarrow [U_i/_{\!p}R_i] \times_{X_i} X
$$
在《空间中的群胚》公式
(\ref{spaces-groupoids-equation-quotient-stack}) 所述的“群胚预层”
层次上是同构。因此，该断言由叠化与纤维积可交换这一事实推出，
参见《叠》引理
\ref{stacks-lemma-stackification-fibre-product-fibred-categories}。
\end{proof}




\section{有限型对象嵌入有限表示对象}
\label{stacks-limits-section-finite-type-closed-in-finite-presentation}

\noindent
本节是《空间的极限》第
\ref{spaces-limits-section-finite-type-closed-in-finite-presentation}
节的对应版本。

\begin{lemma}
\label{stacks-limits-lemma-finite-type-closed-in-finite-presentation}
设 $f : \mathcal{X} \to Y$ 是从代数叠到代数空间的态射。假设：
\begin{enumerate}
\item $f$ 有限型且拟分离；
\item $Y$ 拟紧且拟分离。
\end{enumerate}
则存在有限表示态射 $f' : \mathcal{X}' \to Y$，以及 $Y$
上代数叠之间的闭浸入 $\mathcal{X} \to \mathcal{X}'$。
\end{lemma}

\begin{proof}
将 $Y = \lim_{i \in I} Y_i$ 写成以有向集 $I$ 为指标的代数空间极限，
其中转移态射为仿射态射，且 $Y_i$ 为诺特代数空间；参见
《空间的极限》命题 \ref{spaces-limits-proposition-approximate}。
我们将使用《空间的极限》第
\ref{spaces-limits-section-finite-type-quasi-separated} 节的材料。

\medskip\noindent
选取一个表示 $\mathcal{X} = [U/R]$，并以
$(U, R, s, t, c, e, i)$ 表示相应的 $Y$ 上代数空间群胚。
可以并且确实设 $U$ 仿射。于是 $U$、$R$ 和
$R \times_{s, U, t} R$ 都是 $Y$ 上有限型、拟分离的代数空间。
% P12-ERR-0005: see ../control/ERRATA_CANDIDATES.jsonl
有两个态射 $s, t : R \to U$，三个态射
$c : R \times_{s, U, t} R \to R$,
$\text{pr}_1 : R \times_{s, U, t} R \to R$,
$\text{pr}_2 : R \times_{s, U, t} R \to R$,
，一个态射 $e : U \to R$，以及最后一个态射 $i : R \to R$。
这些态射满足一组公理；详见《群胚》第
\ref{groupoids-section-groupoids} 节。

\medskip\noindent
根据《空间的极限》注
\ref{spaces-limits-remark-finite-type-gives-well-defined-system}
，可找到 $i_0 \in I$ 以及逆系统
\begin{enumerate}
\item $(U_i)_{i \geq i_0}$,
\item $(R_i)_{i \geq i_0}$,
\item $(T_i)_{i \geq i_0}$
\end{enumerate}
，它们位于 $(Y_i)_{i \geq i_0}$ 上，并满足
$U = \lim_{i \geq i_0} U_i$,
$R = \lim_{i \geq i_0} R_i$, and
$R \times_{s, U, t} R = \lim_{i \geq i_0} T_i$
，而且存在系统态射
\begin{enumerate}
\item $(s_i)_{i \geq i_0} : (R_i)_{i \geq i_0} \to (U_i)_{i \geq i_0}$,
\item $(t_i)_{i \geq i_0} : (R_i)_{i \geq i_0} \to (U_i)_{i \geq i_0}$,
\item $(c_i)_{i \geq i_0} : (T_i)_{i \geq i_0} \to (R_i)_{i \geq i_0}$,
\item $(p_i)_{i \geq i_0} : (T_i)_{i \geq i_0} \to (R_i)_{i \geq i_0}$,
\item $(q_i)_{i \geq i_0} : (T_i)_{i \geq i_0} \to (R_i)_{i \geq i_0}$,
\item $(e_i)_{i \geq i_0} : (U_i)_{i \geq i_0} \to (R_i)_{i \geq i_0}$,
\item $(i_i)_{i \geq i_0} : (R_i)_{i \geq i_0} \to (R_i)_{i \geq i_0}$
\end{enumerate}
使得
$s = \lim_{i \geq i_0} s_i$,
$t = \lim_{i \geq i_0} t_i$,
$c = \lim_{i \geq i_0} c_i$,
$\text{pr}_1 = \lim_{i \geq i_0} p_i$,
$\text{pr}_2 = \lim_{i \geq i_0} q_i$,
$e = \lim_{i \geq i_0} e_i$, and
$i = \lim_{i \geq i_0} i_i$.
由《空间的极限》引理
\ref{spaces-limits-lemma-morphism-good-diagram-smooth}
，可设 $s_i$ 和 $t_i$ 均光滑（这可能需要增大 $i_0$）。
由《空间的极限》引理
\ref{spaces-limits-lemma-morphism-good-diagram-flat}
，可设对所有 $i \geq i_0$，由 $s$ 和 $R \to R_i$ 给出的映射
$R \to U \times_{U_i, s_i} R_i$，以及由 $t$ 和 $R \to R_i$
给出的映射 $R \to U \times_{U_i, t_i} R_i$，都是同构。
由《空间的极限》引理
\ref{spaces-limits-lemma-good-diagram-fibre-product}，可设各图
$$
\xymatrix{
T_i \ar[r]_{q_i} \ar[d]_{p_i} & R_i \ar[d]^{t_i} \\
R_i \ar[r]^{s_i} & U_i
}
$$
均为笛卡儿图。《空间的极限》引理
\ref{spaces-limits-lemma-morphism-good-diagram} 的唯一性于是保证：
当 $i$ 充分大时，$s_i, t_i, c_i, e_i, i_i$
满足上述态射 $s, t, c, e, i$ 之间的关系。固定这样的一个 $i$。

\medskip\noindent
于是 $(U_i, R_i, s_i, t_i, c_i, e_i, i_i)$ 是 $Y_i$
上代数空间中的光滑群胚。因此 $\mathcal{X}_i = [U_i/R_i]$
是代数叠（《代数叠》定理
\ref{algebraic-theorem-smooth-groupoid-gives-algebraic-stack}).
群胚态射
$$
(U, R, s, t, c, e, i) \to (U_i, R_i, s_i, t_i, c_i, e_i, i_i)
$$
位于 $Y \to Y_i$ 上，并确定交换图
$$
\xymatrix{
\mathcal{X} \ar[d] \ar[r] & \mathcal{X}_i \ar[d] \\
Y \ar[r] & Y_i
}
$$
（《空间中的群胚》引理
\ref{spaces-groupoids-lemma-quotient-stack-functorial}).
我们断言态射 $\mathcal{X} \to Y \times_{Y_i} \mathcal{X}_i$
是闭浸入。该断言将完成证明，因为按构造，代数叠
$\mathcal{X}_i \to Y_i$ 有限表示。为证此断言，注意左图
$$
\xymatrix{
U \ar[d] \ar[r] & U_i \ar[d] \\
\mathcal{X} \ar[r] & \mathcal{X}_i
}
\quad\quad
\xymatrix{
U \ar[d] \ar[r] & Y \times_{Y_i} U_i \ar[d] \\
\mathcal{X} \ar[r] & Y \times_{Y_i} \mathcal{X}_i
}
$$
由《空间中的群胚》引理
\ref{spaces-groupoids-lemma-criterion-fibre-product}
及上述结果可知是笛卡儿图，故右侧交换图也是笛卡儿图。
根据《空间的极限》引理
\ref{spaces-limits-lemma-limit-from-good-diagram} 中逆系统
$(U_i)$ 的构造，$U \to Y \times_{Y_i} U_i$ 是闭浸入；
而 $Y \times_{Y_i} U_i \to Y \times_{Y_i} \mathcal{X}_i$
是光滑满射。结合《叠的性质》引理
\ref{stacks-properties-lemma-check-immersion-covering}，
即得所需结论。
\end{proof}

\noindent
对于分离代数叠，也有如下版本。

\begin{lemma}
\label{stacks-limits-lemma-separated-closed-in-finite-presentation}
设 $f : \mathcal{X} \to Y$ 是从代数叠到代数空间的态射。假设：
\begin{enumerate}
\item $f$ 有限型且分离；
\item $Y$ 拟紧且拟分离。
\end{enumerate}
则存在分离的有限表示态射 $f' : \mathcal{X}' \to Y$，
以及 $Y$ 上代数叠之间的闭浸入
$\mathcal{X} \to \mathcal{X}'$。
\end{lemma}

\begin{proof}
首先完全按照引理
\ref{stacks-limits-lemma-finite-type-closed-in-finite-presentation}
证明中的步骤（并沿用其中的记号），把嵌入
$\mathcal{X} \to \mathcal{X}'$ 构造成态射
$\mathcal{X} \to \mathcal{X}' = Y \times_{Y_i} \mathcal{X}_i$，
其中 $\mathcal{X}_i = [U_i/R_i]$。因此只需证明当 $i$
充分大时，$\mathcal{X}_i \to Y_i$ 分离。换言之，只需证明当 $i$
充分大时，$\mathcal{X}_i \to \mathcal{X}_i \times_{Y_i} \mathcal{X}_i$
固有。由于态射
$U_i \times_{Y_i} U_i \to \mathcal{X}_i \times_{Y_i} \mathcal{X}_i$
是光滑满射，且
$R_i = \mathcal{X}_i
\times_{\mathcal{X}_i \times_{Y_i} \mathcal{X}_i} U_i \times_{Y_i} U_i$
，只需证明态射
$(s_i, t_i) : R_i \to U_i \times_{Y_i} U_i$
在 $i$ 充分大时固有；参见《叠的性质》引理
\ref{stacks-properties-lemma-check-property-covering}。
下一段证明这一点。

\medskip\noindent
注意，$U \times_Y U \to Y$ 拟分离且有限型。因此可使用
《空间的极限》注
\ref{spaces-limits-remark-finite-type-gives-well-defined-system}
的构造，找到 $i_1 \in I$ 以及逆系统 $(V_i)_{i \geq i_1}$，
使 $U \times_Y U = \lim_{i \geq i_1} V_i$。
由《空间的极限》引理
\ref{spaces-limits-lemma-good-diagram-fibre-product}，
当 $i$ 充分大时，将该构造的函子性用于投影
$U \times_Y U \to U$，得到闭浸入
$$
V_i \to U_i \times_{Y_i} U_i
$$
（这里略有不符，因为严格说来应以 $Y \to Y_i$ 的概形论像替换
$Y_i$，但显然这不改变该纤维积。）另一方面，由《空间的极限》引理
\ref{spaces-limits-lemma-morphism-good-diagram-proper}
，把函子性用于固有态射
$(s, t) : R \to U \times_Y U$（此处使用 $\mathcal{X}$ 分离），
便得到当 $i$ 充分大时为固有态射的 $R_i \to V_i$。
复合这些态射可得：对所有充分大的 $i$，态射
$R_i \to U_i \times_{Y_i} U_i$ 固有。《空间的极限》注
\ref{spaces-limits-remark-finite-type-gives-well-defined-system}
% P12-ERR-0005: see ../control/ERRATA_CANDIDATES.jsonl
中构造的函子性表明，当 $i$ 充分大时，此态射就是
$(s_i, t_i)$，证明完毕。
\end{proof}






\section{泛闭态射}
\label{stacks-limits-section-universally-closed}

\noindent
本节是《空间的极限》第
\ref{spaces-limits-section-universally-closed} 节的对应版本。

\begin{lemma}
\label{stacks-limits-lemma-separate}
设 $g : Z \to Y$ 为仿射概形之间的态射，并设
$f : \mathcal{X} \to Y$ 为代数叠之间的拟紧态射。
设 $z \in Z$，且 $T \subset |\mathcal{X} \times_Y Z|$
是满足 $z \not \in \Im(T \to |Z|)$ 的闭子集。
若 $\mathcal{X}$ 拟紧，则存在开邻域 $V \subset Z$（含 $z$）、交换图
$$
\xymatrix{
V \ar[d] \ar[r]_a & Z' \ar[d]^b \\
Z \ar[r]^g & Y,
}
$$ 
% P12-ERR-0006: see ../control/ERRATA_CANDIDATES.jsonl
以及闭子集 $T' \subset |X \times_Y Z'|$，使得
\begin{enumerate}
\item $Z'$ 是 $Y$ 上有限表示的仿射概形；
\item 令 $z' = a(z)$，则 $z' \not \in \Im(T' \to |Z'|)$；
\item $T$ 在 $|\mathcal{X} \times_Y V|$ 中的逆像经由
$|\mathcal{X} \times_Y V| \to |\mathcal{X} \times_Y Z'|$
映入 $T'$。
\end{enumerate}
\end{lemma}

\begin{proof}
我们从概形态射的相应结果推出本结论。由于 $\mathcal{X}$ 拟紧，
可选取仿射概形 $W$ 及光滑满射 $W \to \mathcal{X}$。
令 $T_W \subset |W \times_Y Z|$ 为 $T$ 的逆像。
则 $z$ 不属于 $T_W$ 的像。由概形情形（《极限》引理
\ref{limits-lemma-separate}），可找到含 $z$ 的开邻域
$V \subset Z$、概形的交换图
$$
\xymatrix{
V \ar[d] \ar[r]_a & Z' \ar[d]^b \\
Z \ar[r]^g & Y,
}
$$
以及闭子集 $T' \subset |W \times_Y Z'|$，使得
\begin{enumerate}
\item $Z'$ 是 $Y$ 上有限表示的仿射概形；
\item 令 $z' = a(z)$，则 $z' \not \in \Im(T' \to |Z'|)$；
\item $T_1 = T_W \cap |W \times_Y V|$ 经由
$|W \times_Y V| \to |W \times_Y Z'|$ 映入 $T'$。
\end{enumerate}
交换图
$$
\xymatrix{
W \times_Y Z \ar[d] &
W \times_Y V \ar[l] \ar[rr]_{a_1} \ar[d]_c & &
W \times_Y Z' \ar[d]^q \\
\mathcal{X} \times_Y Z &
\mathcal{X} \times_Y V \ar[l] \ar[rr]^{a_2} & &
\mathcal{X} \times_Y Z'
}
$$
的两个方块均为笛卡儿方块，且竖直映射为光滑满射，因而尤其是开映射。
考察左侧方块可见，
$T_1 = T_W \cap |W \times_Y V|$ 是
$T_2 = T \cap |\mathcal{X} \times_Y V|$ 在 $c$ 下的逆像。
由《叠的性质》引理
\ref{stacks-properties-lemma-points-cartesian}，得到
$a_1(T_1) = q^{-1}(a_2(T_2))$。
由《拓扑》引理 \ref{topology-lemma-open-morphism-quotient-topology}，
得到
$$
q^{-1}\left(\overline{a_2(T_2)}\right) =
\overline{q^{-1}(a_2(T_2))} =
\overline{a_1(T_1)} \subset T'
$$
由于 $q$ 满射，$\overline{a_2(T_2)} \to |Z'|$ 的像不含 $z'$，
因为对 $T'$ 亦然。
因此，取上述含 $Z', V, a, b$ 的图以及闭子集
$\overline{a_2(T_2)} \subset |\mathcal{X} \times_Y Z'|$，
即得到本引理所述问题的解。
\end{proof}

\begin{lemma}
\label{stacks-limits-lemma-test-universally-closed}
设 $f : \mathcal{X} \to \mathcal{Y}$ 为代数叠之间的拟紧态射。
下列条件等价：
\begin{enumerate}
\item $f$ 泛闭；
% P12-ERR-0007: see ../control/ERRATA_CANDIDATES.jsonl
\item 对每个局部有限表示态射 $Z \to \mathcal{Y}$，其中 $Z$
为仿射概形，映射 $|\mathcal{X} \times_Y Z| \to |Z|$ 是闭映射；
\item 存在概形 $V$ 及光滑满射 $V \to \mathcal{Y}$，使得
$|\mathbf{A}^n \times (\mathcal{X} \times_\mathcal{Y} V)|
\to |\mathbf{A}^n \times V|$ 对所有 $n \geq 0$ 都是闭映射。
\end{enumerate}
\end{lemma}

\begin{proof}
(1) 蕴含 (2) 显然。

\medskip\noindent
假设 (2)。选取一个由仿射概形不交并而成的概形 $V$，以及光滑满射
$V \to \mathcal{Y}$。为证明 $f$ 泛闭，只需证明 $f$ 的基变换
$\mathcal{X} \times_\mathcal{Y} V \to V$ 泛闭；参见《叠的态射》引理
\ref{stacks-morphisms-lemma-universally-closed-local}。
注意，对该基变换，条件 (2) 仍成立。因此在证明 (2) 蕴含 (1) 时，
可设 $Y = \mathcal{Y}$ 是仿射概形。

\medskip\noindent
假设 (2)，并设 $\mathcal{Y} = Y$ 为仿射概形。若 $f$ 不泛闭，
则存在 $Y$ 上仿射概形 $Z$，使
$|\mathcal{X} \times_Y Z| \to |Z|$ 不是闭映射；参见《叠的态射》引理
\ref{stacks-morphisms-lemma-universally-closed-local}。
这意味着存在闭子集 $T \subset |\mathcal{X} \times_Y Z|$，
使 $\Im(T \to |Z|)$ 不闭。选取 $z \in |Z|$，使其属于 $T$
之像的闭包但不属于该像。应用引理 \ref{stacks-limits-lemma-separate}，
得到开邻域 $V \subset Z$、交换图
$$
\xymatrix{
V \ar[d] \ar[r]_a & Z' \ar[d]^b \\
Z \ar[r]^g & Y,
}
$$
以及闭子集 $T' \subset |\mathcal{X} \times_Y Z'|$，使得
\begin{enumerate}
\item $Z'$ 是 $Y$ 上有限表示的仿射概形；
\item 令 $z' = a(z)$，则 $z' \not \in \Im(T' \to |Z'|)$；
\item $T$ 在 $|\mathcal{X} \times_Y V|$ 中的逆像经由
$|\mathcal{X} \times_Y V| \to |\mathcal{X} \times_Y Z'|$
映入 $T'$。
\end{enumerate}
我们断言 $z'$ 属于 $\Im(T' \to |Z'|)$ 的闭包。这意味着
$|\mathcal{X} \times_Y Z'| \to |Z'|$ 不是闭映射，
与假设 (2) 矛盾。换言之，该断言证明了 (2) 蕴含 (1)。
为证明此断言，考察交换图
$$
\xymatrix{
\mathcal{X} \times_Y Z \ar[d] &
\mathcal{X} \times_Y V \ar[l] \ar[d] \ar[r] &
\mathcal{X} \times_Y Z' \ar[d] \\
Z & V \ar[l] \ar[r]^a & Z'
}
$$
令 $T_V \subset |\mathcal{X} \times_Y V|$ 为 $T$ 的逆像。
由《叠的性质》引理 \ref{stacks-properties-lemma-points-cartesian}，
$T_V$ 在 $|V|$ 中的像等于 $T$ 在 $|Z|$ 中之像的逆像。
由于 $z$ 属于 $T \to |Z|$ 之像的闭包，且 $|V| \to |Z|$ 是开映射，
$z$ 属于 $T_V \to |V|$ 之像的闭包。又因 $T_V$ 在
$|\mathcal{X} \times_Y Z'|$ 中的像包含于 $|T'|$，
立即得到 $z' = a(z)$ 属于 $T'$ 之像的闭包。

\medskip\noindent
(1) 蕴含 (3) 显然。设 $V \to \mathcal{Y}$ 如 (3) 所述。
% P12-ERR-0008: see ../control/ERRATA_CANDIDATES.jsonl
若能证明 $\mathcal{X} \times_Y V \to V$ 泛闭，则由《叠的态射》引理
\ref{stacks-morphisms-lemma-universally-closed-local}，$f$ 泛闭。
因此，只需证明：若 $f : \mathcal{X} \to \mathcal{Y}$
是代数叠之间的拟紧态射，$\mathcal{Y} = Y$ 是概形，且对所有 $n$，
$|\mathbf{A}^n \times \mathcal{X}| \to |\mathbf{A}^n \times Y|$
均为闭映射，则 $f$ 满足 (2)。设 $Z \to Y$ 局部有限表示，
其中 $Z$ 为仿射概形。我们必须证明映射
$|\mathcal{X} \times_Y Z| \to |Z|$ 是闭映射。
由于 $Y$ 是概形、$Z$ 仿射且 $Z \to Y$ 局部有限表示，
可找到浸入 $Z \to \mathbf{A}^n \times Y$；参见《态射》引理
\ref{morphisms-lemma-quasi-affine-finite-type-over-S}。考虑笛卡儿图
$$
\vcenter{
\xymatrix{
\mathcal{X} \times_Y Z \ar[d] \ar[r] &
\mathbf{A}^n \times \mathcal{X} \ar[d] \\
Z \ar[r] & \mathbf{A}^n \times Y
}
}
\quad
\begin{matrix}
\text{诱导出} \\
\text{Cartesian 方块}
\end{matrix}
\quad
\vcenter{
\xymatrix{
|\mathcal{X} \times_Y Z| \ar[d] \ar[r] &
|\mathbf{A}^n \times \mathcal{X}| \ar[d] \\
|Z| \ar[r] & |\mathbf{A}^n \times Y|
}
}
$$
，其中水平箭头都是到局部闭子集的同胚（《叠的性质》引理
\ref{stacks-properties-lemma-subspace-induced-topology}）。
% P12-ERR-0009: see ../control/ERRATA_CANDIDATES.jsonl
因此，$|X \times_Y Z|$ 的每个闭子集 $T$ 都是
$|\mathbf{A}^n \times Y|$ 中某个闭子集 $T'$ 的拉回。
假设说明 $T'$ 在 $|\mathbf{A}^n \times X|$ 中的像闭，
故 $T$ 在 $|Z|$ 中的像闭，正合所需。
\end{proof}
% Generated mechanically from the frozen stacks-cohomology.tex target.
% Do not edit this generated file; edit the bound locale source instead.

% OK, start here.
%

\setcounter{chapter}{102}
\chapter{代数叠的上同调}


\phantomsection
\label{stacks-cohomology-section-phantom}





\section{引言}
\label{stacks-cohomology-section-introduction}

\noindent
本章讨论代数叠的上同调，尤其包括拟凝聚层的上同调；
也就是说，我们将证明与《概形的上同调》和《代数空间的上同调》
两章结果相对应的结论。本章的结果与 \cite{LM-B} 中结果的主要区别是，
我们始终使用“大 site”。阅读本章之前，请先快速浏览《代数叠上的层》一章，
以熟悉其中引入的术语；参见《叠上的层》第
\ref{stacks-sheaves-section-introduction} 节。



\section{约定与语言滥用}
\label{stacks-cohomology-section-conventions}

\noindent
我们继续采用《叠的性质》第
\ref{stacks-properties-section-conventions} 节中引入的约定和语言滥用。








\section{记号}
\label{stacks-cohomology-section-notation}

\noindent
不同的拓扑。若用花体字母表示代数叠，例如
$\mathcal{X}, \mathcal{Y}, \mathcal{Z}$，则记号
$\mathcal{X}_{Zar}, \mathcal{X}_\etale, \mathcal{X}_{smooth},
\mathcal{X}_{syntomic}, \mathcal{X}_{fppf}$ 表示《叠上的层》定义
\ref{stacks-sheaves-definition-inherited-topologies} 中引入的 site
（可将其理解为“大 site”）。相应地，$\mathcal{X}$ 的结构层是
$\mathcal{X}_{fppf}$ 上的层。另一方面，代数空间和概形通常以正体大写字母
表示，例如 $X, Y, Z$；在这种情形下，$X_\etale$ 表示 $X$ 的小 \'etale site
（定义见《拓扑》定义
\ref{topologies-definition-big-small-etale} 或《空间的性质》定义
\ref{spaces-properties-definition-etale-site}）。二者的区别应当足够清楚。

\medskip\noindent
默认拓扑为 fppf 拓扑。因此，当我们所指的是 $\mathcal{X}_{fppf}$ 上的层
或 $\textit{Mod}(\mathcal{X}_{fppf}, \mathcal{O}_\mathcal{X})$ 的对象时，
有时会简称“$\mathcal{X}$ 上的层”或“$\mathcal{O}_\mathcal{X}$-模层”。

\medskip\noindent
若 $f : \mathcal{X} \to \mathcal{Y}$ 是代数叠之间的态射，
则对上述任一拓扑，预层上的函子 $f_*$ 和 $f^{-1}$ 都保持层。
因此，讨论层的直像或逆像时无需说明所采用的拓扑；但在计算上同调群和（或）
高阶直像时并非如此，此时我们总会明确所采用的拓扑。

\medskip\noindent
设 $f : X \to \mathcal{Y}$ 是从代数空间 $X$ 到代数叠
$\mathcal{Y}$ 的态射，并设 $\mathcal{G}$ 是某个拓扑 $\tau$ 下
$\mathcal{Y}_\tau$ 上的层。此时 $f^{-1}\mathcal{G}$ 是
$\mathcal{S}_X$（与 $X$ 相伴的代数叠）上 $\tau$ 拓扑的层，
因为按照我们的约定，$f$ 实际上是一个 $1$-态射
$f : \mathcal{S}_X \to \mathcal{Y}$。若 $\tau = \etale$ 或更强，
则以 $f^{-1}\mathcal{G}|_{X_\etale}$ 表示它在 $X$ 的 \'etale site 上的限制；
参见《叠上的层》第 \ref{stacks-sheaves-section-compare} 节。
若 $\mathcal{G}$ 是 $\mathcal{O}_\mathcal{X}$-模层，
我们有时改写为 $f^*\mathcal{G}$ 和 $f^*\mathcal{G}|_{X_\etale}$。




\section{拟凝聚模的拉回}
\label{stacks-cohomology-section-pullback}

\noindent
设 $f : \mathcal{X} \to \mathcal{Y}$ 是代数叠之间的态射。
赋环拓扑斯上的拟凝聚模与拉回相容，这是一个非常一般的事实。
特别地，拉回 $f^*$ 保持拟凝聚模，从而得到函子
$$
f^* :
\QCoh(\mathcal{O}_\mathcal{Y})
\longrightarrow
\QCoh(\mathcal{O}_\mathcal{X}),
$$
参见《叠上的层》引理
\ref{stacks-sheaves-lemma-pullback-quasi-coherent}。
一般而言，该函子并不正合；但当 $f$ 平坦时，它是正合的。

\begin{lemma}
\label{stacks-cohomology-lemma-flat-pullback-quasi-coherent}
若 $f : \mathcal{X} \to \mathcal{Y}$ 是代数叠之间的平坦态射，
则 $f^* : \QCoh(\mathcal{O}_\mathcal{Y}) \to
\QCoh(\mathcal{O}_\mathcal{X})$ 是正合函子。
\end{lemma}

\begin{proof}
选择概形 $V$ 以及满的光滑态射 $V \to \mathcal{Y}$。
再选择概形 $U$ 以及满的光滑态射
$U \to V \times_\mathcal{Y} \mathcal{X}$。
由于 $U \to \mathcal{X}$ 是两个此类态射的复合，它仍然光滑且满。
由交换图
$$
\xymatrix{
U \ar[d] \ar[r]_{f'} & V \ar[d] \\
\mathcal{X} \ar[r]^f & \mathcal{Y}
}
$$
得到阿贝尔范畴的交换图
$$
\xymatrix{
\QCoh(\mathcal{O}_U) & \QCoh(\mathcal{O}_V) \ar[l] \\
\QCoh(\mathcal{O}_\mathcal{X}) \ar[u] &
\QCoh(\mathcal{O}_\mathcal{Y}) \ar[l] \ar[u]
}
$$
我们关于该图下方两个范畴为阿贝尔范畴的证明表明，竖直函子都是忠实正合函子
（参见《叠上的层》引理
\ref{stacks-sheaves-lemma-quasi-coherent-algebraic-stack} 的证明）。
按照代数叠平坦态射的定义，$f'$ 是概形之间的平坦态射，
故 $(f')^*$ 在 $V$ 上的拟凝聚层范畴中是正合函子，结论随即成立。
\end{proof}

\begin{lemma}
\label{stacks-cohomology-lemma-quasi-coherent-check-exact}
设 $\mathcal{X}$ 为代数叠，$I$ 为集合；对每个 $i \in I$，
设 $x_i : U_i \to \mathcal{X}$ 是 $\mathcal{X}$ 的一个对象。
假设 $x_i$ 平坦，且
$\coprod x_i : \coprod U_i \to \mathcal{X}$ 为满射。
设 $\varphi : \mathcal{F} \to \mathcal{G}$ 是
$\QCoh(\mathcal{O}_\mathcal{X})$ 中的箭头，并以 $\varphi_i$
表示 $\varphi$ 在 $(U_i)_\etale$ 上的限制。
则 $\varphi$ 为单射、满射或同构，当且仅当每个 $\varphi_i$ 分别具有相应性质。
\end{lemma}

\begin{proof}
选择概形 $U$ 以及满的光滑态射 $x : U \to \mathcal{X}$，
并将 $x$ 视为 $\mathcal{X}$ 的一个对象。由此得到某个空间群胚
$(U, R, s, t, c)$ 给出的表示 $\mathcal{X} = [U/R]$，以及相应的等价
$$
\QCoh(\mathcal{O}_\mathcal{X}) = \QCoh(U, R, s, t, c)
$$
参见《叠上的层》第
\ref{stacks-sheaves-section-quasi-coherent-algebraic-stacks} 节中的讨论。
右端的阿贝尔范畴结构满足：$\varphi$ 为单射、满射或同构，当且仅当
$\varphi|_{U_\etale}$ 分别具有相应性质；参见《空间中的群胚》引理
\ref{spaces-groupoids-lemma-abelian}。

\medskip\noindent
% P12-ERR-0010: see ../control/ERRATA_CANDIDATES.jsonl
对每个 $i$，选择由概形组成的 \'etale 覆盖
$\{W_{i, j} \to V \times_\mathcal{X} U_i\}_{j \in J_i}$。
以 $g_{i, j} : W_{i, j} \to V$ 和
$h_{i, j} : W_{i, j} \to U_i$ 表示显然的箭头。
每个概形态射 $g_{i, j} : W_{i, j} \to U$ 都是平坦的，且它们联合满射。
类似地，对每个固定的 $i$，概形态射
$h_{i, j} : W_{i, j} \to U_i$ 都是平坦的，且联合满射。
由《叠上的层》引理 \ref{stacks-sheaves-lemma-quasi-coherent}，
$\varphi|_{U_\etale}$ 沿 $(g_{i, j})_{small}$ 的拉回是
$\varphi|_{(W_{i, j})_\etale}$，而
$\varphi|_{(U_i)_\etale}$ 沿 $(h_{i, j})_{small}$ 的拉回也是
$\varphi|_{(W_{i, j})_\etale}$。
概形之间的平坦态射对拟凝聚模的拉回是正合的；
沿一族联合满射的概形平坦态射拉回，会反映拟凝聚模的单射、满射和双射
（事实上，这对所有模都成立，因为可在茎上检验正合性）。因此
$$
\varphi|_{U_\etale} \text{ 单射}
\Leftrightarrow
\varphi|_{(W_{i, j})_\etale} \text{ 对所有 }i,j\text{ 均为单射}
\Leftrightarrow
\varphi|_{(U_i)_\etale} \text{ 对所有 }i\text{ 均为单射}
$$
证明完毕。
\end{proof}


\section{各类模的高阶直像}
\label{stacks-cohomology-section-key}

\noindent
下述引理是理解若干类模层之高阶直像的基础。
它有两个版本：一个用于 \'etale 拓扑，另一个用于 fppf 拓扑。

\begin{lemma}
\label{stacks-cohomology-lemma-general-pushforward}
设 $\mathcal{M}$ 是一条规则，它为每个代数叠 $\mathcal{X}$ 指定
$\textit{Mod}(\mathcal{X}_\etale, \mathcal{O}_\mathcal{X})$
的一个子范畴 $\mathcal{M}_\mathcal{X}$，并满足：
\begin{enumerate}
\item 对每个代数叠 $\mathcal{X}$，$\mathcal{M}_\mathcal{X}$ 是
$\textit{Mod}(\mathcal{X}_\etale, \mathcal{O}_\mathcal{X})$ 的弱 Serre 子范畴
（参见《同调》定义 \ref{homology-definition-serre-subcategory}）；
\item 对代数叠之间的光滑态射
$f : \mathcal{Y} \to \mathcal{X}$，函子 $f^*$ 将
$\mathcal{M}_\mathcal{X}$ 映入 $\mathcal{M}_\mathcal{Y}$；
\item 若 $f_i : \mathcal{X}_i \to \mathcal{X}$ 是一族代数叠之间的
光滑态射，且
$|\mathcal{X}| = \bigcup |f_i|(|\mathcal{X}_i|)$，则
$\textit{Mod}(\mathcal{X}_\etale, \mathcal{O}_\mathcal{X})$ 的对象
$\mathcal{F}$ 属于 $\mathcal{M}_\mathcal{X}$，当且仅当对每个 $i$，
$f_i^*\mathcal{F}$ 属于 $\mathcal{M}_{\mathcal{X}_i}$；
\item 若 $f : \mathcal{Y} \to \mathcal{X}$ 是代数叠之间的态射，
且 $\mathcal{X}$ 和 $\mathcal{Y}$ 均可由仿射概形表示，则 $R^if_*$
将 $\mathcal{M}_\mathcal{Y}$ 映入 $\mathcal{M}_\mathcal{X}$。
\end{enumerate}
于是，对代数叠之间任意拟紧且拟分离的态射
$f : \mathcal{Y} \to \mathcal{X}$，$R^if_*$ 将
$\mathcal{M}_\mathcal{Y}$ 映入 $\mathcal{M}_\mathcal{X}$。
（高阶直像在 \'etale 拓扑中计算。）
\end{lemma}

\begin{proof}
设 $f : \mathcal{Y} \to \mathcal{X}$ 是代数叠之间拟紧且拟分离的态射，
并设 $\mathcal{F}$ 是 $\mathcal{M}_\mathcal{Y}$ 的对象。
选择满的光滑态射 $\mathcal{U} \to \mathcal{X}$，其中
$\mathcal{U}$ 可由概形表示。由《叠上的层》引理
\ref{stacks-sheaves-lemma-base-change-higher-direct-images}，
取高阶直像与基变换可交换。假设 (2) 表明 $\mathcal{F}$ 在
$\mathcal{U} \times_\mathcal{X} \mathcal{Y}$ 上的拉回属于
$\mathcal{M}_{\mathcal{U} \times_\mathcal{X} \mathcal{Y}}$，
因为投影
$\mathcal{U} \times_\mathcal{X} \mathcal{Y} \to \mathcal{Y}$
是光滑态射的基变换，因而光滑。于是由 (3)，可将
$\mathcal{Y} \to \mathcal{X}$ 替换为投影
$\mathcal{U} \times_\mathcal{X} \mathcal{Y} \to \mathcal{U}$。
换言之，可假设 $\mathcal{X}$ 可由概形表示。
再用一次 (3)，可见该问题关于 $\mathcal{X}$ 是 Zariski 局部的，
故可假设 $\mathcal{X}$ 可由仿射概形表示。由于 $f$ 拟紧，
$\mathcal{Y}$ 也拟紧。因此可选择满的光滑态射
$g : \mathcal{V} \to \mathcal{Y}$，其中 $\mathcal{V}$ 可由仿射概形表示。

\medskip\noindent
在这种情形下，有《叠上的层》命题
\ref{stacks-sheaves-proposition-smooth-covering-compute-direct-image}
给出的谱序列
$$
E_2^{p, q} = R^q(f \circ g_p)_*g_p^*\mathcal{F}
\Rightarrow
R^{p + q}f_*\mathcal{F}
$$
回顾这是第一象限谱序列，因此可使用《同调》引理
\ref{homology-lemma-first-quadrant-ss} 的最后一部分。注意态射
$$
g_p : \mathcal{V}_p =
\mathcal{V} \times_\mathcal{Y} \ldots \times_\mathcal{Y} \mathcal{V}
\longrightarrow
\mathcal{Y}
$$
是光滑的，因为它们是光滑态射 $g$ 的若干基变换的复合。
因此由 (2)，层 $g_p^*\mathcal{F}$ 属于
$\mathcal{M}_{\mathcal{V}_p}$。所以只需证明：
$\mathcal{M}_{\mathcal{V}_p}$ 中对象沿态射
$$
\mathcal{V}_p =
\mathcal{V} \times_\mathcal{Y} \ldots \times_\mathcal{Y} \mathcal{V}
\longrightarrow
\mathcal{X}
$$
的高阶直像属于 $\mathcal{M}_\mathcal{X}$。
由《叠的态射》引理
\ref{stacks-morphisms-lemma-quasi-compact-quasi-separated-permanence}，
代数叠 $\mathcal{V}_p$ 拟紧且拟分离。当然，每个 $\mathcal{V}_p$
都可由代数空间表示（因为代数叠 $\mathcal{Y}$ 的对角可由代数空间表示）。
这把问题约化为如下情形：$\mathcal{Y}$ 可由代数空间表示，
而 $\mathcal{X}$ 可由仿射概形表示。

\medskip\noindent
在 $\mathcal{Y}$ 可由代数空间表示且 $\mathcal{X}$ 可由仿射概形表示的情形，
重新选择满的光滑态射 $\mathcal{V} \to \mathcal{Y}$，
其中 $\mathcal{V}$ 可由仿射概形表示。再次沿用上述论证，
又将问题约化到态射 $\mathcal{V}_p \to \mathcal{X}$。
但在当前情形下，代数叠 $\mathcal{V}_p$ 可由拟紧且拟分离的概形表示
（因为代数空间的对角可由概形表示）。

\medskip\noindent
因此，可假设 $\mathcal{Y}$ 可由概形表示，而 $\mathcal{X}$
可由仿射概形表示。再次选择满的光滑态射
$\mathcal{V} \to \mathcal{Y}$，其中 $\mathcal{V}$ 可由仿射概形表示。
此时，所有代数叠 $\mathcal{V}_p$ 都可由分离概形表示
（因为概形的对角是分离的）。

\medskip\noindent
因此，可假设 $\mathcal{Y}$ 可由分离概形表示，而 $\mathcal{X}$
可由仿射概形表示。又一次选择满的光滑态射
$\mathcal{V} \to \mathcal{Y}$，其中 $\mathcal{V}$ 可由仿射概形表示。
此时所有代数叠 $\mathcal{V}_p$ 都可由仿射概形表示
（因为分离概形的对角是闭浸入，因而是仿射的），
故该情形由假设 (4) 处理。证明完毕。
\end{proof}

\noindent
下面给出 fppf 拓扑的版本。

\begin{lemma}
\label{stacks-cohomology-lemma-general-pushforward-fppf}
设 $\mathcal{M}$ 是一条规则，它为每个代数叠 $\mathcal{X}$ 指定
$\textit{Mod}(\mathcal{O}_\mathcal{X})$ 的一个子范畴
$\mathcal{M}_\mathcal{X}$，并满足：
\begin{enumerate}
% P12-ERR-0011: see ../control/ERRATA_CANDIDATES.jsonl
\item $\mathcal{O}_\mathcal{X}$ 是
$\textit{Mod}(\mathcal{O}_\mathcal{X})$ 的弱 Serre 子范畴，
对所有代数叠 $\mathcal{X}$ 均如此；
\item 对代数叠之间的光滑态射
$f : \mathcal{Y} \to \mathcal{X}$，函子 $f^*$ 将
$\mathcal{M}_\mathcal{X}$ 映入 $\mathcal{M}_\mathcal{Y}$；
\item 若 $f_i : \mathcal{X}_i \to \mathcal{X}$ 是一族代数叠之间的
光滑态射，且
$|\mathcal{X}| = \bigcup |f_i|(|\mathcal{X}_i|)$，则
$\textit{Mod}(\mathcal{O}_\mathcal{X})$ 的对象 $\mathcal{F}$
属于 $\mathcal{M}_\mathcal{X}$，当且仅当对每个 $i$，
$f_i^*\mathcal{F}$ 属于 $\mathcal{M}_{\mathcal{X}_i}$；
\item 若 $f : \mathcal{Y} \to \mathcal{X}$ 是代数叠之间的态射，
且 $\mathcal{X}$ 和 $\mathcal{Y}$ 均可由仿射概形表示，
则 $R^if_*$ 将 $\mathcal{M}_\mathcal{Y}$ 映入
$\mathcal{M}_\mathcal{X}$。
\end{enumerate}
于是，对代数叠之间任意拟紧且拟分离的态射
$f : \mathcal{Y} \to \mathcal{X}$，$R^if_*$ 将
$\mathcal{M}_\mathcal{Y}$ 映入 $\mathcal{M}_\mathcal{X}$。
（高阶直像在 fppf 拓扑中计算。）
\end{lemma}

\begin{proof}
与引理 \ref{stacks-cohomology-lemma-general-pushforward} 的证明相同。
\end{proof}

\section{局部拟凝聚模}
\label{stacks-cohomology-section-locally-quasi-coherent}

\noindent
设 $\mathcal{X}$ 为代数叠，$\mathcal{F}$ 为
$\mathcal{O}_\mathcal{X}$-模预层。可以问 $\mathcal{F}$ 是否
{\it 局部拟凝聚}；参见《叠上的层》定义
\ref{stacks-sheaves-definition-locally-quasi-coherent}。
简言之，这表示 $\mathcal{F}$ 是 \'etale 拓扑下的
$\mathcal{O}_\mathcal{X}$-模层，并且对任意态射
$f : U \to \mathcal{X}$，限制 $f^*\mathcal{F}|_{U_\etale}$
在 $U_\etale$ 上拟凝聚。（实际定义略有不同，但与此等价。）
一个有用的事实是
$$
\textit{LQCoh}(\mathcal{O}_\mathcal{X}) \subset
\textit{Mod}(\mathcal{X}_\etale, \mathcal{O}_\mathcal{X})
$$
为弱 Serre 子范畴；参见《叠上的层》引理
\ref{stacks-sheaves-lemma-lqc-colimits}。

\begin{lemma}
\label{stacks-cohomology-lemma-check-lqc-on-etale-covering}
设 $\mathcal{X}$ 为代数叠，且
$f_j : \mathcal{X}_j \to \mathcal{X}$ 是一族代数叠之间的
光滑态射，满足
$|\mathcal{X}| =\bigcup |f_j|(|\mathcal{X}_j|)$。
设 $\mathcal{F}$ 是 $\mathcal{X}_\etale$ 上的
$\mathcal{O}_\mathcal{X}$-模层。若每个 $f_j^{-1}\mathcal{F}$
都局部拟凝聚，则 $\mathcal{F}$ 也局部拟凝聚。
\end{lemma}

\begin{proof}
可以把每个代数叠 $\mathcal{X}_j$ 换成概形 $U_j$
（这里用到任意代数叠都有概形给出的光滑覆盖，且光滑态射的复合仍光滑；
参见《叠的态射》引理
\ref{stacks-morphisms-lemma-composition-smooth}）。
$\mathcal{F}$ 在 $(\Sch/U_j)_\etale$ 上的拉回仍然局部拟凝聚；
参见《叠上的层》引理
\ref{stacks-sheaves-lemma-pullback-lqc}。
于是 $f = \coprod f_j : U = \coprod U_j \to \mathcal{X}$
是满的光滑态射。设 $x$ 为 $\mathcal{X}$ 的一个对象。
由《叠上的层》引理
\ref{stacks-sheaves-lemma-surjective-flat-locally-finite-presentation}，
存在 \'etale 覆盖 $\{x_i \to x\}_{i \in I}$，使每个 $x_i$
都可提升为 $(\Sch/U)_\etale$ 的一个对象 $u_i$。
这只是说，$x$ 和 $x_i$ 分别位于概形 $V$ 和 $V_i$ 上，
$\{V_i \to V\}$ 是 \'etale 覆盖，并且 $x_i$ 来自态射
$u_i : V_i \to U$。限制 $x_i^*\mathcal{F}|_{V_{i, \etale}}$
等于 $f^*\mathcal{F}$ 在 $V_{i, \etale}$ 上的限制；
参见《叠上的层》引理 \ref{stacks-sheaves-lemma-comparison}。
因此，$x^*\mathcal{F}|_{V_\etale}$ 是 $V$ 的小 \'etale site 上的层，
且对每个 $i$，其在 $V_{i, \etale}$ 上的限制都拟凝聚。
这蕴含它本身拟凝聚，正合所需；例如可用《空间的性质》引理
\ref{spaces-properties-lemma-characterize-quasi-coherent}。
\end{proof}

\begin{lemma}
\label{stacks-cohomology-lemma-pushforward-locally-quasi-coherent}
设 $f : \mathcal{X} \to \mathcal{Y}$ 是代数叠之间拟紧且拟分离的态射，
并设 $\mathcal{F}$ 是 $\mathcal{X}_\etale$ 上局部拟凝聚的
$\mathcal{O}_\mathcal{X}$-模层。则 $R^if_*\mathcal{F}$
（在 \'etale 拓扑中计算）在 $\mathcal{Y}_\etale$ 上局部拟凝聚。
\end{lemma}

\begin{proof}
我们用引理 \ref{stacks-cohomology-lemma-general-pushforward} 证明该结论，
并检验其中的假设 (1)--(4)。第 (1)、(2) 项由《叠上的层》引理
\ref{stacks-sheaves-lemma-lqc-colimits} 得出；第 (3) 项由引理
\ref{stacks-cohomology-lemma-check-lqc-on-etale-covering} 得出。因此只需证明 (4)。

\medskip\noindent
设 $f : \mathcal{X} \to \mathcal{Y}$ 是代数叠之间的态射，
且 $\mathcal{X}$ 和 $\mathcal{Y}$ 分别可由仿射概形 $X$ 和 $Y$ 表示。
任取位于概形 $V$ 上的 $\mathcal{Y}$ 的对象 $y$。
为清楚起见，以 $\mathcal{V} = (\Sch/V)_{fppf}$ 表示与 $V$ 相应的代数叠。
考察笛卡儿图
% P12-ERR-0012: see ../control/ERRATA_CANDIDATES.jsonl
$$
\xymatrix{
\mathcal{Z} \ar[d] \ar[r]_g \ar[d]_{f'} & \mathcal{X} \ar[d]^f \\
\mathcal{V} \ar[r]^y & \mathcal{Y}
}
$$
于是 $\mathcal{Z}$ 可由概形 $Z = V \times_Y X$ 表示，
且 $f'$ 拟紧并分离（甚至仿射）。由《叠上的层》引理
\ref{stacks-sheaves-lemma-compare-representable-morphism-cohomology}，有
$$
R^if_*\mathcal{F}|_{V_\etale} =
R^if'_{small, *}\big(g^*\mathcal{F}|_{Z_\etale}\big)
$$
由《空间的上同调》引理
\ref{spaces-cohomology-lemma-higher-direct-image}，右端是
$V_\etale$ 上的拟凝聚层。这就说明左端拟凝聚，正是所需结论。
\end{proof}

\begin{lemma}
\label{stacks-cohomology-lemma-check-lqc-on-flat-covering}
设 $\mathcal{X}$ 为代数叠，且
$f_j : \mathcal{X}_j \to \mathcal{X}$ 是一族代数叠之间平坦且
局部有限表示的态射，满足
$|\mathcal{X}| =\bigcup |f_j|(|\mathcal{X}_j|)$。
设 $\mathcal{F}$ 是 $\mathcal{X}_{fppf}$ 上的
$\mathcal{O}_\mathcal{X}$-模层。若每个 $f_j^{-1}\mathcal{F}$
都局部拟凝聚，则 $\mathcal{F}$ 也局部拟凝聚。
\end{lemma}

\begin{proof}
先假设存在态射 $a : \mathcal{U} \to \mathcal{X}$，它满、平坦、
局部有限表示、拟紧且拟分离，并且 $a^*\mathcal{F}$ 局部拟凝聚。
此时有正合列
$$
0 \to \mathcal{F} \to a_*a^*\mathcal{F} \to b_*b^*\mathcal{F}
$$
其中 $b$ 为态射
$b : \mathcal{U} \times_\mathcal{X} \mathcal{U} \to \mathcal{X}$；
参见《叠上的层》命题
\ref{stacks-sheaves-proposition-exactness-cech-complex} 和引理
\ref{stacks-sheaves-lemma-surjective-flat-locally-finite-presentation}。
此外，$b^*\mathcal{F}$ 是 $a^*\mathcal{F}$ 沿某个投影态射的拉回，
因而局部拟凝聚（《叠上的层》引理
\ref{stacks-sheaves-lemma-pullback-lqc}）。
由引理 \ref{stacks-cohomology-lemma-pushforward-locally-quasi-coherent}，
模 $a_*a^*\mathcal{F}$ 和 $b_*b^*\mathcal{F}$ 均局部拟凝聚。
（注意 $a_*$ 和 $b_*$ 与计算时所用的拓扑无关。）
于是由《叠上的层》引理
\ref{stacks-sheaves-lemma-lqc-colimits}，$\mathcal{F}$ 局部拟凝聚。

\medskip\noindent
下面把一般情形约化到第一段的情形。设 $x$ 是 $\mathcal{X}$
的一个对象，位于概形 $U$ 上。需证明 $\mathcal{F}|_{U_\etale}$
是拟凝聚的 $\mathcal{O}_U$-模层。在 $U$ 上 Zariski 局部地证明即可，
故可假设 $U$ 仿射。由《叠的态射》引理
\ref{stacks-morphisms-lemma-surjective-family-flat-locally-finite-presentation}，
存在 fppf 覆盖 $\{a_i : U_i \to U\}$，使每个 $x \circ a_i$
都经由某个 $f_j$ 分解。因此 $a_i^*\mathcal{F}$ 在
$(\Sch/U_i)_{fppf}$ 上局部拟凝聚。细化该覆盖后，
可假设 $\{U_i \to U\}_{i = 1, \ldots, n}$ 是标准 fppf 覆盖。
于是 $x^*\mathcal{F}$ 是 $(\Sch/U)_{fppf}$ 上的 fppf 模，
其沿态射 $a : U_1 \amalg \ldots \amalg U_n \to U$ 的拉回局部拟凝聚。
由第一段可知 $x^*\mathcal{F}$ 局部拟凝聚；
这当然蕴含 $\mathcal{F}|_{U_\etale}$ 拟凝聚。
\end{proof}


\section{平坦比较映射}
\label{stacks-cohomology-section-flat-comparison}

\noindent
设 $\mathcal{X}$ 为代数叠，且 $\mathcal{F}$ 是
$\textit{Mod}(\mathcal{X}_\etale, \mathcal{O}_\mathcal{X})$ 的对象。
给定 $\mathcal{X}$ 的一个对象 $x$，它位于概形 $U$ 上，则限制
$\mathcal{F}|_{U_\etale}$ 是 $x^{-1}\mathcal{F}$ 在 $U$ 的小
\'etale site 上的限制；参见《叠上的层》定义
\ref{stacks-sheaves-definition-pullback}。
再设 $\varphi : x \to x'$ 是 $\mathcal{X}$ 的一个态射，
位于概形态射 $f : U \to U'$ 上。于是有 $2$-交换图
$$
\xymatrix{
U \ar[rd]_x \ar[rr]_f & & U' \ar[ld]^{x'} \\
& \mathcal{X}
}
$$
与 $\varphi$ 相伴，我们得到限制之间的比较映射
\begin{equation}
\label{stacks-cohomology-equation-comparison-modules}
c_\varphi :
f_{small}^*(\mathcal{F}|_{U'_\etale})
\longrightarrow
\mathcal{F}|_{U_\etale}
\end{equation}
参见《叠上的层》公式
(\ref{stacks-sheaves-equation-comparison-modules})。
在这种情形下，可考察 $\mathcal{F}$ 的如下性质。

\begin{definition}
\label{stacks-cohomology-definition-flat-base-change}
设 $\mathcal{X}$ 为代数叠，并设 $\mathcal{F}$ 属于
$\textit{Mod}(\mathcal{X}_\etale, \mathcal{O}_\mathcal{X})$。
若且唯若每当 $f$ 平坦时 $c_\varphi$ 都是同构，
我们称 $\mathcal{F}$ 具有{\it 平坦基变换性质}\footnote{这可能不是标准记号。}
\end{definition}

\noindent
下面的引理给出这一概念的若干性质。

\begin{lemma}
\label{stacks-cohomology-lemma-check-flat-comparison-on-etale-covering}
设 $\mathcal{X}$ 为代数叠，且 $\mathcal{F}$ 是
$\mathcal{X}_\etale$ 上的 $\mathcal{O}_\mathcal{X}$-模层。
\begin{enumerate}
\item 若 $\mathcal{F}$ 具有平坦基变换性质，则对代数叠的任意态射
$g : \mathcal{Y} \to \mathcal{X}$，拉回 $g^*\mathcal{F}$ 也具有该性质。
\item $\textit{Mod}(\mathcal{X}_\etale, \mathcal{O}_\mathcal{X})$
中由具有平坦基变换性质的模组成的充满子范畴，是弱 Serre 子范畴。
\item 设 $f_i : \mathcal{X}_i \to \mathcal{X}$ 是一族代数叠之间的
光滑态射，满足
$|\mathcal{X}| = \bigcup_i |f_i|(|\mathcal{X}_i|)$。
若每个 $f_i^*\mathcal{F}$ 都具有平坦基变换性质，
则 $\mathcal{F}$ 也具有该性质。
\item $\mathcal{X}_\etale$ 上具有平坦基变换性质的
$\mathcal{O}_\mathcal{X}$-模范畴有余极限，且这些余极限与
$\textit{Mod}(\mathcal{X}_\etale, \mathcal{O}_\mathcal{X})$
中的余极限一致。
\item 给定 $\textit{Mod}(\mathcal{X}_\etale, \mathcal{O}_\mathcal{X})$
中具有平坦基变换性质的 $\mathcal{F}$ 和 $\mathcal{G}$，
则张量积 $\mathcal{F} \otimes_{\mathcal{O}_\mathcal{X}} \mathcal{G}$
具有平坦基变换性质。
\item 给定 $\textit{Mod}(\mathcal{X}_\etale, \mathcal{O}_\mathcal{X})$
中的 $\mathcal{F}$ 和 $\mathcal{G}$，若 $\mathcal{F}$ 有限表示，
且 $\mathcal{G}$ 具有平坦基变换性质，则层
$\SheafHom_{\mathcal{O}_\mathcal{X}}(\mathcal{F}, \mathcal{G})$
具有平坦基变换性质。
\end{enumerate}
\end{lemma}

\begin{proof}
设 $g : \mathcal{Y} \to \mathcal{X}$ 如 (1) 所述，
并设 $y$ 是 $\mathcal{Y}$ 的一个对象，位于概形 $V$ 上。
由《叠上的层》引理 \ref{stacks-sheaves-lemma-comparison}，有
$(g^*\mathcal{F})|_{V_\etale} = \mathcal{F}|_{V_\etale}$。
此外，$\mathcal{Y}$ 上层 $g^*\mathcal{F}$ 的比较映射，
是 $\mathcal{X}$ 上层 $\mathcal{F}$ 的比较映射的特例；
参见《叠上的层》引理 \ref{stacks-sheaves-lemma-comparison}。
由此 (1) 显然成立。

\medskip\noindent
证明 (2)。使用《同调》引理
\ref{homology-lemma-characterize-weak-serre-subcategory}
对弱 Serre 子范畴的刻画。具有平坦基变换性质的层之间的映射，
其核与余核也具有该性质。这是显然的，因为对概形的平坦态射，
$f_{small}^*$ 正合；而限制函子 $(-)|_{U_\etale}$ 也正合
（因为我们在 \'etale 拓扑中工作）。最后，若
$0 \to \mathcal{F}_1 \to \mathcal{F}_2 \to \mathcal{F}_3 \to 0$
是 $\textit{Mod}(\mathcal{X}_\etale, \mathcal{O}_\mathcal{X})$
中的短正合列，且外侧两个层具有平坦基变换性质，
则中间的层也具有该性质；这仍由 $f_{small}^*$ 与限制函子的正合性
以及五引理得出。

\medskip\noindent
证明 (3)。设 $f_i : \mathcal{X}_i \to \mathcal{X}$
是一族联合满射的代数叠光滑态射，并假设每个
$f_i^*\mathcal{F}$ 都具有平坦基变换性质。
由 (1)、代数叠的定义以及光滑态射的复合仍光滑这一事实
（参见《叠的态射》引理
\ref{stacks-morphisms-lemma-composition-smooth}），
可假设每个 $\mathcal{X}_i$ 都可由概形表示。
设 $\varphi : x \to x'$ 是 $\mathcal{X}$ 的一个态射，
位于概形的平坦态射 $a : U \to U'$ 上。
由《叠上的层》引理
\ref{stacks-sheaves-lemma-surjective-flat-locally-finite-presentation}，
存在一族联合满射的 \'etale 态射 $U'_i \to U'$，使
$U'_i \to U' \to \mathcal{X}$ 经由 $\mathcal{X}_i$ 分解。
于是得到交换图
$$
\xymatrix{
U_i = U \times_{U'} U_i' \ar[r]_-{a_i} \ar[d] &
U_i' \ar[r]_{x_i'} \ar[d] & \mathcal{X}_i \ar[d]^{f_i} \\
U \ar[r]^a & U' \ar[r]^{x'} & \mathcal{X}
}
$$
注意，每个 $a_i$ 都是 $a$ 的基变换，因而是概形的平坦态射。
以 $\psi_i : x_i \to x'_i$ 表示 $\mathcal{X}_i$ 的态射，
它位于 $a_i$ 上且目标为 $x_i'$。由假设，比较映射
$c_{\psi_i} :
(a_i)_{small}^*\big(f_i^*\mathcal{F}|_{(U'_i)_\etale}\big)
\to f_i^*\mathcal{F}|_{(U_i)_\etale}$ 是同构。
因为竖直箭头 $U_i' \to U'$ 和 $U_i \to U$ 是 \'etale 的，
层 $f_i^*\mathcal{F}|_{(U_i')_\etale}$ 和
$f_i^*\mathcal{F}|_{(U_i)_\etale}$ 分别是
$\mathcal{F}|_{U'_\etale}$ 和 $\mathcal{F}|_{U_\etale}$ 的限制，
而映射 $c_{\psi_i}$ 是 $c_\varphi$ 在 $(U_i)_\etale$ 上的限制；
参见《叠上的层》引理 \ref{stacks-sheaves-lemma-comparison}。
由于 $\{U_i \to U\}$ 是 \'etale 覆盖，这说明比较映射
$c_\varphi$ 是同构，正合所需。

\medskip\noindent
证明 (4)。设 $\mathcal{I} \to
\textit{Mod}(\mathcal{X}_\etale, \mathcal{O}_\mathcal{X})$，
$i \mapsto \mathcal{F}_i$ 是一个图，并假设每个 $\mathcal{F}_i$
都具有平坦基变换性质。设 $\varphi : x \to x'$ 是
$\mathcal{X}$ 的一个态射，位于概形的平坦态射 $f : U \to U'$ 上。
回顾 $\colim_i \mathcal{F}_i$ 是预层余极限的层化。
由于采用 \'etale 拓扑，显然有
$$
(\colim_i \mathcal{F}_i)|_{U_\etale} =
\colim_i {\mathcal{F}_i}|_{U_\etale}
$$
在 $U'_\etale$ 上的限制也同样如此。因而
\begin{align*}
f_{small}^*((\colim_i \mathcal{F}_i)|_{U'_\etale})
& =
f_{small}^*(\colim_i {\mathcal{F}_i}|_{U'_\etale}) \\
& =
\colim_i f_{small}^*({\mathcal{F}_i}|_{U'_\etale}) \\
& \xrightarrow{\colim c_\varphi}
\colim_i \mathcal{F}_i|_{U_\etale} \\
& =
(\colim_i \mathcal{F}_i)|_{U_\etale}
\end{align*}
第二个等号使用了 $f_{small}^*$ 与余极限可交换这一事实
（因为它是左伴随）。箭头为同构，因为每个 $\mathcal{F}_i$
都具有平坦基变换性质。因此余极限具有平坦基变换性质，(4) 成立。

\medskip\noindent
(5) 成立，因为张量积与拉回可交换；参见《site 上的模》引理
\ref{sites-modules-lemma-tensor-product-pullback}。略去细节。

\medskip\noindent
设 $\mathcal{F}$ 和 $\mathcal{G}$ 如 (6) 所述。
由于 $\mathcal{F}$ 拟凝聚，由《叠上的层》引理
\ref{stacks-sheaves-lemma-quasi-coherent}，它具有平坦基变换性质。
设 $\varphi : x \to x'$ 是 $\mathcal{X}$ 的一个态射，
位于概形的平坦态射 $f : U \to U'$ 上。
由于采用 \'etale 拓扑，有
$$
\SheafHom_{\mathcal{O}_\mathcal{X}}(\mathcal{F}, \mathcal{G})|_{U_\etale} =
\SheafHom_{\mathcal{O}_U}(\mathcal{F}|_{U_\etale}, \mathcal{G}|_{U_\etale})
$$
在 $U'_\etale$ 上的限制也同样如此（略去细节）。因此
% P12-ERR-0013: see ../control/ERRATA_CANDIDATES.jsonl
\begin{align*}
f_{small}^*(
\SheafHom_{\mathcal{O}_\mathcal{X}}(\mathcal{F}, \mathcal{G})|_{U'_\etale})
& =
f_{small}^*(
\SheafHom_{\mathcal{O}_{U'}}(
\mathcal{F}|_{U'_\etale}, \mathcal{G}|_{U'_\etale})) \\
& =
\SheafHom_{\mathcal{O}_{U'}}(
f_{small}^*(\mathcal{F}|_{U'_\etale}),
f_{small}^*(\mathcal{G}|_{U'_\etale})) \\
& \xrightarrow{c_\varphi}
\SheafHom_{\mathcal{O}_U}(\mathcal{F}|_{U_\etale}, \mathcal{G}|_{U_\etale}) \\
& =
\SheafHom_{\mathcal{O}_\mathcal{X}}(\mathcal{F}, \mathcal{G})|_{U_\etale}
\end{align*}
这里第二个等号是《site 上的模》引理
\ref{sites-modules-lemma-pullback-internal-hom}；该引理用到
$f : U \to U'$ 平坦，因而赋环 site 的态射 $f_{small}$ 也平坦。
箭头为同构，因为 $\mathcal{F}$ 和 $\mathcal{G}$ 都具有
平坦基变换性质。因此所述 $\SheafHom$ 也具有平坦基变换性质，正合所需。
\end{proof}

\begin{lemma}
\label{stacks-cohomology-lemma-flat-comparison}
设 $f : \mathcal{X} \to \mathcal{Y}$ 是代数叠之间拟紧且拟分离的态射。
设 $\mathcal{F}$ 是
$\textit{Mod}(\mathcal{X}_\etale, \mathcal{O}_\mathcal{X})$ 的对象，
它局部拟凝聚并具有平坦基变换性质。则每个 $R^if_*\mathcal{F}$
（在 \'etale 拓扑中计算）都具有平坦基变换性质。
\end{lemma}

\begin{proof}
我们用引理 \ref{stacks-cohomology-lemma-general-pushforward} 证明该结论。
对每个代数叠 $\mathcal{X}$，以
$\textit{LQCoh}^{fbc}(\mathcal{O}_\mathcal{X})$ 表示
$\textit{Mod}(\mathcal{X}_\etale, \mathcal{O}_\mathcal{X})$
中由局部拟凝聚且具有平坦基变换性质的层组成的充满子范畴。
一旦验证引理 \ref{stacks-cohomology-lemma-general-pushforward} 的条件 (1)--(4)，
结论即得。性质 (1)、(2)、(3) 由《叠上的层》引理
\ref{stacks-sheaves-lemma-pullback-lqc}、
\ref{stacks-sheaves-lemma-lqc-colimits}，以及引理
\ref{stacks-cohomology-lemma-check-lqc-on-etale-covering}、
\ref{stacks-cohomology-lemma-check-flat-comparison-on-etale-covering} 得出。
因此只需证明 (4)。

\medskip\noindent
设 $f : \mathcal{X} \to \mathcal{Y}$ 是代数叠之间的态射，
且 $\mathcal{X}$ 和 $\mathcal{Y}$ 分别可由仿射概形 $X$ 和 $Y$ 表示。
在这种情形下，设 $\psi : y \to y'$ 是 $\mathcal{Y}$ 的一个态射，
位于概形的平坦态射 $b : V \to V'$ 上。为清楚起见，记
$\mathcal{V} = (\Sch/V)_{fppf}$ 和
$\mathcal{V}' = (\Sch/V')_{fppf}$ 为相应的代数叠。
考察代数叠的图
$$
\xymatrix{
\mathcal{Z} \ar[d]_{f''} \ar[r]_a &
\mathcal{Z}' \ar[r]_{x'} \ar[d]_{f'} & \mathcal{X} \ar[d]^f \\
\mathcal{V} \ar[r]^b & \mathcal{V}' \ar[r]^{y'} & \mathcal{Y}
}
$$
其中两个方块都是笛卡儿的。由于 $f$ 可由概形表示
（且拟紧、分离，甚至仿射），$\mathcal{Z}$ 和 $\mathcal{Z}'$
分别可由概形 $Z$ 和 $Z'$ 表示，而且
$Z = V \times_{V'} Z'$。因为 $\mathcal{F}$ 具有平坦基变换性质，
映射
$$
a_{small}^*\big(\mathcal{F}|_{Z'_\etale}\big)
\longrightarrow
\mathcal{F}|_{Z_\etale}
$$
是同构。此外，由《叠上的层》引理
\ref{stacks-sheaves-lemma-compare-representable-morphism-cohomology}，有
$$
R^if_*\mathcal{F}|_{V'_\etale} =
R^i(f')_{small, *}\big(\mathcal{F}|_{Z'_\etale}\big)
$$
以及
$$
R^if_*\mathcal{F}|_{V_\etale} =
R^i(f'')_{small, *}\big(\mathcal{F}|_{Z_\etale}\big)
$$
因此，由《空间的上同调》引理
\ref{spaces-cohomology-lemma-flat-base-change-cohomology}，比较映射
$$
c_\psi :
b_{small}^*(R^if_*\mathcal{F}|_{V'_\etale})
\longrightarrow
R^if_*\mathcal{F}|_{V_\etale}
$$
是同构。故 $R^if_*\mathcal{F}$ 具有平坦基变换性质。
又由引理 \ref{stacks-cohomology-lemma-pushforward-locally-quasi-coherent}，
$R^if_*\mathcal{F}$ 局部拟凝聚，证明完毕。
\end{proof}


\section{具有平坦基变换性质的局部拟凝聚模}
\label{stacks-cohomology-section-loc-qcoh-flat-base-change}

\noindent
设 $\mathcal{X}$ 为代数叠。我们\footnote{抱歉使用了这样繁复的记号。}记
$$
\textit{LQCoh}^{fbc}(\mathcal{O}_\mathcal{X})
\subset
\textit{Mod}(\mathcal{X}_\etale, \mathcal{O}_\mathcal{X})
$$
为如下充满子范畴：其对象是 \'etale
$\mathcal{O}_\mathcal{X}$-模 $\mathcal{F}$，它既局部拟凝聚
（第 \ref{stacks-cohomology-section-locally-quasi-coherent} 节），
又具有平坦基变换性质（第 \ref{stacks-cohomology-section-flat-comparison} 节）。
由《叠上的层》引理 \ref{stacks-sheaves-lemma-quasi-coherent}，有
$$
\QCoh(\mathcal{O}_\mathcal{X}) \subset
\textit{LQCoh}^{fbc}(\mathcal{O}_\mathcal{X})
$$

\begin{proposition}
\label{stacks-cohomology-proposition-loc-qcoh-flat-base-change}
关于具有平坦基变换性质的局部拟凝聚模，有如下结果汇总。
\begin{enumerate}
\item 设 $\mathcal{X}$ 为代数叠。若 $\mathcal{F}$ 属于
$\textit{LQCoh}^{fbc}(\mathcal{O}_\mathcal{X})$，
则 $\mathcal{F}$ 是 fppf 拓扑的层；也就是说，它是
$\textit{Mod}(\mathcal{O}_\mathcal{X})$ 的对象。
\item 范畴 $\textit{LQCoh}^{fbc}(\mathcal{O}_\mathcal{X})$
既是 $\textit{Mod}(\mathcal{O}_\mathcal{X})$ 的弱 Serre 子范畴，
也是 $\textit{Mod}(\mathcal{X}_\etale, \mathcal{O}_\mathcal{X})$
的弱 Serre 子范畴。
\item 沿代数叠的任意态射 $f : \mathcal{X} \to \mathcal{Y}$ 的拉回 $f^*$，
会诱导函子
$f^* : \textit{LQCoh}^{fbc}(\mathcal{O}_\mathcal{Y}) \to
\textit{LQCoh}^{fbc}(\mathcal{O}_\mathcal{X})$。
\item 若 $f : \mathcal{X} \to \mathcal{Y}$ 是代数叠之间拟紧且
拟分离的态射，且 $\mathcal{F}$ 是
$\textit{LQCoh}^{fbc}(\mathcal{O}_\mathcal{X})$ 的对象，则
\begin{enumerate}
\item 总直像 $Rf_*\mathcal{F}$ 和高阶直像 $R^if_*\mathcal{F}$
既可在 \'etale 拓扑中计算，也可在 fppf 拓扑中计算，结果相同；
\item 每个 $R^if_*\mathcal{F}$ 都是
$\textit{LQCoh}^{fbc}(\mathcal{O}_\mathcal{Y})$ 的对象。
\end{enumerate}
\item 范畴 $\textit{LQCoh}^{fbc}(\mathcal{O}_\mathcal{X})$
有余极限，且这些余极限既与
$\textit{Mod}(\mathcal{X}_\etale, \mathcal{O}_\mathcal{X})$
中的余极限一致，也与
$\textit{Mod}(\mathcal{O}_\mathcal{X})$ 中的余极限一致。
\item 给定 $\textit{LQCoh}^{fbc}(\mathcal{O}_\mathcal{X})$ 中的
$\mathcal{F}$ 和 $\mathcal{G}$，则张量积
$\mathcal{F} \otimes_{\mathcal{O}_\mathcal{X}} \mathcal{G}$
属于 $\textit{LQCoh}^{fbc}(\mathcal{O}_\mathcal{X})$。
\item 给定有限表示的 $\mathcal{F}$ 以及
$\textit{LQCoh}^{fbc}(\mathcal{O}_\mathcal{X})$ 中的 $\mathcal{G}$，
则 $\SheafHom_{\mathcal{O}_\mathcal{X}}(\mathcal{F}, \mathcal{G})$
属于 $\textit{LQCoh}^{fbc}(\mathcal{O}_\mathcal{X})$。
\end{enumerate}
\end{proposition}

\begin{proof}
(1) 是《叠上的层》引理
\ref{stacks-sheaves-lemma-lqc-flat-base-change-fppf-sheaf}。

\medskip\noindent
对于嵌入
$\textit{LQCoh}^{fbc}(\mathcal{O}_\mathcal{X}) \subset
\textit{Mod}(\mathcal{X}_\etale, \mathcal{O}_\mathcal{X})$，
我们已经在引理 \ref{stacks-cohomology-lemma-flat-comparison} 的证明中看到 (2)。
现在证明嵌入
$\textit{LQCoh}^{fbc}(\mathcal{O}_\mathcal{X}) \subset
\textit{Mod}(\mathcal{O}_\mathcal{X})$
的 (2)。设 $\varphi : \mathcal{F} \to \mathcal{G}$ 是
$\textit{LQCoh}^{fbc}(\mathcal{O}_\mathcal{X})$ 中对象之间的态射。
由于无论在 \'etale 还是 fppf 拓扑中计算，$\Ker(\varphi)$ 都相同，
由 \'etale 情形可知 $\Ker(\varphi)$ 属于
$\textit{LQCoh}^{fbc}(\mathcal{O}_\mathcal{X})$。
另一方面，在 fppf 拓扑中计算的余核，是在 \'etale 拓扑中计算所得余核的
fppf 层化。但由 \'etale 情形，该 \'etale 余核属于
$\textit{LQCoh}^{fbc}(\mathcal{O}_\mathcal{X})$，
故由 (1) 它已经是 fppf 层。因此该余核属于
$\textit{LQCoh}^{fbc}(\mathcal{O}_\mathcal{X})$。
最后，假设
% P12-ERR-0014: see ../control/ERRATA_CANDIDATES.jsonl
$$
0 \to \mathcal{F}_1 \to \mathcal{F}_2 \to \mathcal{F}_3 \to 0
$$
是 $\textit{Mod}(\mathcal{O}_\mathcal{X})$ 中的正合列
（即使用 fppf 拓扑），其中 $\mathcal{F}_1$、$\mathcal{F}_2$
属于 $\textit{LQCoh}^{fbc}(\mathcal{O}_\mathcal{X})$。
为证明 $\mathcal{F}_2$ 是
$\textit{LQCoh}^{fbc}(\mathcal{O}_\mathcal{X})$ 的对象，
只需证明该序列在 \'etale 拓扑中也正合。为此，只需证明
$H^1_{fppf}(x, \mathcal{F}_1)$ 的任意元素，在 $x$ 的某个 \'etale
覆盖的各成员上都变为零（这里 $x$ 是 $\mathcal{X}$ 的任意对象）。
这是真的，因为由《叠上的层》引理
\ref{stacks-sheaves-lemma-compare-fppf-etale}，有
$H^1_{fppf}(x, \mathcal{F}_1) = H^1_\etale(x, \mathcal{F}_1)$；
并且上同调是局部的，参见《site 上的上同调》引理
\ref{sites-cohomology-lemma-kill-cohomology-class-on-covering}。
这证明了 (2)。

\medskip\noindent
(3) 由引理 \ref{stacks-cohomology-lemma-check-flat-comparison-on-etale-covering}
和《叠上的层》引理
\ref{stacks-sheaves-lemma-pullback-lqc} 得出。

\medskip\noindent
对于在 \'etale 上同调中计算的 $R^if_*\mathcal{F}$，
(4)(b) 由引理 \ref{stacks-cohomology-lemma-flat-comparison} 得出。
继而，将《叠上的层》引理
\ref{stacks-sheaves-lemma-compare-fppf-etale} 与上面的 (1) 结合，
便得到 (4)(a)。

\medskip\noindent
\'etale 拓扑的 (5) 由《叠上的层》引理
\ref{stacks-sheaves-lemma-lqc-colimits} 和引理
\ref{stacks-cohomology-lemma-check-flat-comparison-on-etale-covering} 得出。
于是 fppf 版本也成立，因为 \'etale 拓扑中的余极限由 (1)
已经是 fppf 层。

\medskip\noindent
(6)、(7) 由引理
\ref{stacks-cohomology-lemma-check-flat-comparison-on-etale-covering} 的相应部分和
《叠上的层》引理 \ref{stacks-sheaves-lemma-lqc-colimits} 得出。
\end{proof}

\begin{lemma}
\label{stacks-cohomology-lemma-check-lqc-fbc-on-covering}
设 $\mathcal{X}$ 为代数叠。
\begin{enumerate}
\item 设 $f_j : \mathcal{X}_j \to \mathcal{X}$ 是一族代数叠之间的
光滑态射，满足
$|\mathcal{X}| =\bigcup |f_j|(|\mathcal{X}_j|)$。
设 $\mathcal{F}$ 是 $\mathcal{X}_\etale$ 上的
$\mathcal{O}_\mathcal{X}$-模层。若每个 $f_j^{-1}\mathcal{F}$
% P12-ERR-0015: see ../control/ERRATA_CANDIDATES.jsonl
都属于 $\textit{LQCoh}^{fpc}(\mathcal{O}_{\mathcal{X}_i})$，
则 $\mathcal{F}$ 属于
$\textit{LQCoh}^{fbc}(\mathcal{O}_\mathcal{X})$。
\item 设 $f_j : \mathcal{X}_j \to \mathcal{X}$ 是一族代数叠之间
平坦且局部有限表示的态射，满足
$|\mathcal{X}| =\bigcup |f_j|(|\mathcal{X}_j|)$。
设 $\mathcal{F}$ 是 $\mathcal{X}_{fppf}$ 上的
$\mathcal{O}_\mathcal{X}$-模层。若每个 $f_j^{-1}\mathcal{F}$
都属于 $\textit{LQCoh}^{fbc}(\mathcal{O}_{\mathcal{X}_i})$，
则 $\mathcal{F}$ 属于
$\textit{LQCoh}^{fbc}(\mathcal{O}_\mathcal{X})$。
\end{enumerate}
\end{lemma}

\begin{proof}
(1) 由引理 \ref{stacks-cohomology-lemma-check-lqc-on-etale-covering} 和
\ref{stacks-cohomology-lemma-check-flat-comparison-on-etale-covering} 结合得出。
(2) 的证明与引理 \ref{stacks-cohomology-lemma-check-lqc-on-flat-covering} 的证明类似。
设 $\mathcal{F}$ 为 $\mathcal{X}_{fppf}$ 上的
$\mathcal{O}_\mathcal{X}$-模层。

\medskip\noindent
先假设存在态射 $a : \mathcal{U} \to \mathcal{X}$，它满、平坦、
局部有限表示、拟紧且拟分离，并且 $a^*\mathcal{F}$ 局部拟凝聚且
具有平坦基变换性质。此时有正合列
$$
0 \to \mathcal{F} \to a_*a^*\mathcal{F} \to b_*b^*\mathcal{F}
$$
其中 $b$ 为态射
$b : \mathcal{U} \times_\mathcal{X} \mathcal{U} \to \mathcal{X}$；
参见《叠上的层》命题
\ref{stacks-sheaves-proposition-exactness-cech-complex} 和引理
\ref{stacks-sheaves-lemma-surjective-flat-locally-finite-presentation}。
此外，$b^*\mathcal{F}$ 是 $a^*\mathcal{F}$ 沿某个投影态射的拉回，
因而局部拟凝聚并具有平坦基变换性质；参见命题
\ref{stacks-cohomology-proposition-loc-qcoh-flat-base-change}。
由命题 \ref{stacks-cohomology-proposition-loc-qcoh-flat-base-change}，
模 $a_*a^*\mathcal{F}$ 和 $b_*b^*\mathcal{F}$ 均局部拟凝聚且
具有平坦基变换性质。再由命题
\ref{stacks-cohomology-proposition-loc-qcoh-flat-base-change}，$\mathcal{F}$ 局部拟凝聚且
具有平坦基变换性质。

\medskip\noindent
选择概形 $U$ 和满的光滑态射 $x : U \to \mathcal{X}$。
由 (1)，只需证明 $x^*\mathcal{F}$ 局部拟凝聚并具有平坦基变换性质。
再由 (1)，在 $U$ 上 Zariski 局部地证明即可，故可假设 $U$ 仿射。
由《叠的态射》引理
\ref{stacks-morphisms-lemma-surjective-family-flat-locally-finite-presentation}，
存在 fppf 覆盖 $\{a_i : U_i \to U\}$，使每个 $x \circ a_i$
都经由某个 $f_j$ 分解。因此，$(\Sch/U_i)_{fppf}$ 上的模
$a_i^*\mathcal{F}$ 局部拟凝聚并具有平坦基变换性质。
细化该覆盖后，可假设
$\{U_i \to U\}_{i = 1, \ldots, n}$ 是标准 fppf 覆盖。
于是 $x^*\mathcal{F}$ 是 $(\Sch/U)_{fppf}$ 上的 fppf 模，
其沿态射 $a : U_1 \amalg \ldots \amalg U_n \to U$ 的拉回
局部拟凝聚并具有平坦基变换性质。因此由上一段，
$x^*\mathcal{F}$ 局部拟凝聚并具有平坦基变换性质，正合所需。
\end{proof}

\begin{lemma}
\label{stacks-cohomology-lemma-loc-qcoh-fbc-compute-higher-direct-image}
设 $f : \mathcal{X} \to \mathcal{Y}$ 是代数叠之间拟紧、拟分离且
可由代数空间表示的态射，并设 $\mathcal{F}$ 属于
$\textit{LQCoh}^{fbc}(\mathcal{O}_\mathcal{X})$。
则对 $\mathcal{Y}$ 的一个对象 $y : V \to \mathcal{Y}$，有
$$
(R^if_*\mathcal{F})|_{V_\etale} = R^if'_{small, *}(\mathcal{F}|_{U_\etale})
$$
其中 $f' : U = V \times_\mathcal{Y} \mathcal{X} \to V$
是 $f$ 的基变换。
\end{lemma}

\begin{proof}
由《叠上的层》引理
\ref{stacks-sheaves-lemma-base-change-higher-direct-images}，
可约化到 $\mathcal{X}$ 由 $U$ 表示、$\mathcal{Y}$ 由 $V$ 表示的情形。
这里当然还用到：由命题
\ref{stacks-cohomology-proposition-loc-qcoh-flat-base-change}，$\mathcal{F}$ 在 $U$
上的拉回属于 $\textit{LQCoh}^{fbc}(\mathcal{O}_U)$。
随后，结论由《叠上的层》引理
\ref{stacks-sheaves-lemma-compare-morphism-cohomology} 以及命题
\ref{stacks-cohomology-proposition-loc-qcoh-flat-base-change}（它说明可在 \'etale
拓扑中计算 $R^if_*$）得出。
\end{proof}

\begin{lemma}
\label{stacks-cohomology-lemma-loc-qcoh-fbc-affine-direct-image}
设 $f : \mathcal{X} \to \mathcal{Y}$ 是代数叠的仿射态射。
函子 $f_* : \textit{LQCoh}^{fbc}(\mathcal{O}_\mathcal{X}) \to
\textit{LQCoh}^{fbc}(\mathcal{O}_\mathcal{Y})$ 正合且与直和可交换。
对 $i > 0$，函子 $R^if_*$ 在
$\textit{LQCoh}^{fbc}(\mathcal{O}_\mathcal{X})$ 上为零。
\end{lemma}

\begin{proof}
这些函子由命题 \ref{stacks-cohomology-proposition-loc-qcoh-flat-base-change} 存在。
由引理 \ref{stacks-cohomology-lemma-loc-qcoh-fbc-compute-higher-direct-image}，
问题约化到代数空间的仿射态射，并在代数空间上的拟凝聚模语境中取高阶直像。
由《空间的上同调》第
\ref{spaces-cohomology-section-higher-direct-image} 节的讨论，
再约化到概形的仿射态射。对概形的仿射态射，
《概形的上同调》引理
\ref{coherent-lemma-relative-affine-vanishing} 给出拟凝聚模的高阶直像消失。
$R^1f_*$ 的消失蕴含 $f_*$ 正合。与直和可交换例如可由《态射》引理
\ref{morphisms-lemma-affine-equivalence-modules} 得出。
\end{proof}


\section{寄生模}
\label{stacks-cohomology-section-parasitic}

\noindent
下述定义与《下降》定义 \ref{descent-definition-parasitic} 相容。

\begin{definition}
\label{stacks-cohomology-definition-parasitic}
设 $\mathcal{X}$ 为代数叠。若对 $\mathcal{X}$ 的任意对象 $x$，
只要该对象位于概形 $U$ 上且相应态射
$x : U \to \mathcal{X}$ 平坦，就有 $\mathcal{F}(x) = 0$，
则称 $\mathcal{O}_\mathcal{X}$-模预层 $\mathcal{F}$ 是{\it 寄生的}。
\end{definition}

\noindent
下面的引理给出这一概念的若干性质。

\begin{lemma}
\label{stacks-cohomology-lemma-parasitic}
设 $\mathcal{X}$ 为代数叠，$\mathcal{F}$ 为
$\mathcal{O}_\mathcal{X}$-模预层。
\begin{enumerate}
\item 若 $\mathcal{F}$ 寄生，且
$g : \mathcal{Y} \to \mathcal{X}$ 是代数叠之间的平坦态射，
则 $g^*\mathcal{F}$ 寄生。
\item 对 $\tau \in \{Zariski, \etale, smooth, syntomic, fppf\}$，有：
\begin{enumerate}
\item 寄生模预层的 $\tau$-层化仍寄生；
\item $\textit{Mod}(\mathcal{X}_\tau, \mathcal{O}_\mathcal{X})$
中由寄生模组成的充满子范畴，是 Serre 子范畴。
\end{enumerate}
\item 假设 $\mathcal{F}$ 是 \'etale 拓扑的层。
设 $f_i : \mathcal{X}_i \to \mathcal{X}$ 是一族代数叠之间的
光滑态射，满足
$|\mathcal{X}| = \bigcup_i |f_i|(|\mathcal{X}_i|)$。
若每个 $f_i^*\mathcal{F}$ 都寄生，则 $\mathcal{F}$ 也寄生。
\item 假设 $\mathcal{F}$ 是 fppf 拓扑的层。
设 $f_i : \mathcal{X}_i \to \mathcal{X}$ 是一族代数叠之间
平坦且局部有限表示的态射，满足
$|\mathcal{X}| = \bigcup_i |f_i|(|\mathcal{X}_i|)$。
若每个 $f_i^*\mathcal{F}$ 都寄生，则 $\mathcal{F}$ 也寄生。
\end{enumerate}
\end{lemma}

\begin{proof}
为证明 (1)，设 $y$ 是 $\mathcal{Y}$ 的一个对象，位于概形 $V$ 上，
且相应态射 $y : V \to \mathcal{Y}$ 平坦。于是
$g(y) : V \to \mathcal{Y} \to \mathcal{X}$ 是平坦态射的复合，
因而平坦（参见《叠的态射》引理
\ref{stacks-morphisms-lemma-composition-flat}）。所以由假设，
$\mathcal{F}(g(y))$ 为零。由于
$g^*\mathcal{F} = g^{-1}\mathcal{F}(y) = \mathcal{F}(g(y))$，
可知 $g^*\mathcal{F}$ 寄生。

\medskip\noindent
为证明 (2)(a)，注意若 $\{x_i \to x\}$ 是 $\mathcal{X}$ 的
$\tau$-覆盖，则每个态射 $x_i \to x$ 都位于概形的平坦态射上。
因此，若 $x$ 位于概形 $U$ 上且 $x : U \to \mathcal{X}$ 平坦，
则所有对象 $x_i$ 也如此。故若预层 $\mathcal{F}$ 寄生，
则预层 $\mathcal{F}^+$（参见《site》第
\ref{sites-section-sheafification} 节）也寄生。
$\mathcal{F}$ 的层化是 $(\mathcal{F}^+)^+$，因而 (2)(a) 得证。

\medskip\noindent
设 $\mathcal{F}$ 为寄生 $\tau$-模。由定义立刻可知，
$\mathcal{F}$ 的任意子模都寄生。另一方面，若
$\mathcal{F}' \subset \mathcal{F}$ 是子模，则预层
$x \mapsto \mathcal{F}(x)/\mathcal{F}'(x)$
显然也寄生。因此由 (2)(a)，商 $\mathcal{F}/\mathcal{F}'$
是寄生模。最后，需证明：若有短正合列
$0 \to \mathcal{F}_1 \to \mathcal{F}_2 \to \mathcal{F}_3 \to 0$，
且 $\mathcal{F}_1$、$\mathcal{F}_3$ 寄生，则 $\mathcal{F}_2$ 寄生。
在位于对 $\mathcal{X}$ 平坦的概形上的 $x$ 处取值，结论立即得出。
这证明了 (2)(b)；参见《同调》引理
\ref{homology-lemma-characterize-serre-subcategory}。

\medskip\noindent
设 $f_i : \mathcal{X}_i \to \mathcal{X}$ 是一族联合满射的
代数叠光滑态射，并假设每个 $f_i^*\mathcal{F}$ 都寄生。
设 $x$ 是 $\mathcal{X}$ 的一个对象，位于概形 $U$ 上，且
$x : U \to \mathcal{X}$ 平坦。考察满的光滑覆盖
$W_i \to U \times_{x, \mathcal{X}} \mathcal{X}_i$，
并以 $y_i : W_i \to \mathcal{X}_i$ 表示投影。
于是 $\{f_i(y_i) \to x\}$ 是 $\mathcal{X}$ 的光滑拓扑覆盖。
由于平坦态射的复合仍平坦，有 $f_i^*\mathcal{F}(y_i) = 0$。
另一方面，如 (1) 的证明所见，
$f_i^*\mathcal{F}(y_i) = \mathcal{F}(f_i(y_i))$。
所以在 $\mathcal{X}$ 中存在某个光滑覆盖
$\{x_i \to x\}_{i \in I}$，使 $\mathcal{F}(x_i) = 0$。
由于光滑拓扑与 \'etale 拓扑相同，这蕴含
$\mathcal{F}(x) = 0$；参见《态射进阶》引理
\ref{more-morphisms-lemma-etale-dominates-smooth}。
具体地说，$\{x_i \to x\}_{i \in I}$ 位于概形的光滑覆盖
$\{U_i \to U\}_{i \in I}$ 上。由刚引用的引理，存在 \'etale 覆盖
$\{V_j \to U\}_{j \in J}$ 细化 $\{U_i \to U\}_{i \in I}$。
记 $x'_j = x|_{V_j}$。则 $\{x'_j \to x\}$ 是 $\mathcal{X}$
中的 \'etale 覆盖，并细化 $\{x_i \to x\}_{i \in I}$。
这说明映射
$\mathcal{F}(x) \to \prod_{j \in J} \mathcal{F}(x'_j)$
经由 $\mathcal{F}(x) \to \prod_{i \in I} \mathcal{F}(x_i)$ 分解；
前者因 $\mathcal{F}$ 是 \'etale 拓扑的层而单射，后者为零。
因此 $\mathcal{F}(x) = 0$，正合所需。

\medskip\noindent
(4) 的证明略去。提示：与 (3) 的证明类似，但更简单。
\end{proof}

\noindent
任意直像都保持寄生模。

\begin{lemma}
\label{stacks-cohomology-lemma-pushforward-parasitic}
设 $\tau \in \{\etale, fppf\}$，$\mathcal{X}$ 为代数叠，
且 $\mathcal{F}$ 是
$\textit{Mod}(\mathcal{X}_\tau, \mathcal{O}_\mathcal{X})$
的寄生对象。
\begin{enumerate}
\item 对所有 $i$，$H^i_\tau(\mathcal{X}, \mathcal{F}) = 0$。
\item 设 $f : \mathcal{X} \to \mathcal{Y}$ 是代数叠之间的态射。
则 $R^if_*\mathcal{F}$（在 $\tau$-拓扑中计算）是
$\textit{Mod}(\mathcal{Y}_\tau, \mathcal{O}_\mathcal{Y})$
的寄生对象。
\end{enumerate}
\end{lemma}

\begin{proof}
先将 (2) 约化到 (1)。由《叠上的层》引理
\ref{stacks-sheaves-lemma-pushforward-restriction} 可知，
$R^if_*\mathcal{F}$ 是与预层
$$
y \longmapsto
H^i_\tau\Big(V \times_{y, \mathcal{Y}} \mathcal{X},
\ \text{pr}^{-1}\mathcal{F}\Big)
$$
相伴的层。这里 $y$ 是 $\mathcal{Y}$ 的典型对象，位于概形 $V$ 上。
由引理 \ref{stacks-cohomology-lemma-parasitic}，只需证明当
$y : V \to \mathcal{Y}$ 平坦时，这些上同调群为零。
注意 $\text{pr} : V \times_{y, \mathcal{Y}} \mathcal{X} \to \mathcal{X}$
是 $y$ 的基变换，因而平坦。故由引理 \ref{stacks-cohomology-lemma-parasitic}，
$\text{pr}^{-1}\mathcal{F}$ 寄生。因此只需证明 (1)。

\medskip\noindent
为证明 (1)，可以使用《叠上的层》命题
\ref{stacks-sheaves-proposition-smooth-covering-compute-cohomology}
的谱序列，把问题约化到 $\mathcal{X}$ 可由代数空间表示的情形。
注意，在该谱序列中，由引理 \ref{stacks-cohomology-lemma-parasitic}，
每个 $f_p^{-1}\mathcal{F} = f_p^*\mathcal{F}$ 都是寄生模，
因为态射
$f_p : \mathcal{U}_p =
\mathcal{U} \times_\mathcal{X} \ldots
\times_\mathcal{X} \mathcal{U} \to \mathcal{X}$ 平坦。
再次使用这一谱序列（如引理
\ref{stacks-cohomology-lemma-general-pushforward} 的证明那样），
可约化到代数叠 $\mathcal{X}$ 可由概形 $X$ 表示的情形。
此时 $H^i_\tau(\mathcal{X}, \mathcal{F}) = H^i((\Sch/X)_\tau, \mathcal{F})$。
在这种情形下，消失性容易由 {\v C}ech 覆盖的论证得出；
参见《下降》引理 \ref{descent-lemma-cohomology-parasitic}。
\end{proof}

\noindent
下述引理是我们关注寄生模的主要原因之一。为理解其陈述，回顾函子
$\QCoh(\mathcal{O}_\mathcal{X}) \to
\textit{Mod}(\mathcal{X}_\etale, \mathcal{O}_\mathcal{X})$
以及
$\QCoh(\mathcal{O}_\mathcal{X}) \to
\textit{Mod}(\mathcal{O}_\mathcal{X})$
一般并不正合。

\begin{lemma}
\label{stacks-cohomology-lemma-exact-sequence-quasi-coherent-parasitic-cohomology}
设 $\mathcal{X}$ 为代数叠，并设
$\alpha : \mathcal{F} \to \mathcal{G}$ 和
$\beta : \mathcal{G} \to \mathcal{H}$
是 $\QCoh(\mathcal{O}_\mathcal{X})$ 中的映射，满足
$\beta \circ \alpha = 0$。下列条件等价：
\begin{enumerate}
\item 在阿贝尔范畴 $\QCoh(\mathcal{O}_\mathcal{X})$ 中，
复形 $\mathcal{F} \to \mathcal{G} \to \mathcal{H}$
在 $\mathcal{G}$ 处正合；
\item 无论在
$\textit{Mod}(\mathcal{X}_\etale, \mathcal{O}_\mathcal{X})$ 还是
$\textit{Mod}(\mathcal{X}_{fppf}, \mathcal{O}_\mathcal{X})$
中计算，$\Ker(\beta)/\Im(\alpha)$ 都寄生。
\end{enumerate}
\end{lemma}

\begin{proof}
有 $\QCoh(\mathcal{O}_\mathcal{X}) \subset
\textit{LQCoh}^{fbc}(\mathcal{O}_\mathcal{X})$；参见第
\ref{stacks-cohomology-section-loc-qcoh-flat-base-change} 节。
因此，在
$\textit{Mod}(\mathcal{X}_\etale, \mathcal{O}_\mathcal{X})$ 或
$\textit{Mod}(\mathcal{X}_{fppf}, \mathcal{O}_\mathcal{X})$
中计算的 $\Ker(\beta)/\Im(\alpha)$ 一致；参见命题
\ref{stacks-cohomology-proposition-loc-qcoh-flat-base-change}。
以下在 $\mathcal{X}$ 上采用 \'etale 拓扑。

\medskip\noindent
设 $\mathcal{E}$ 是复形
$\mathcal{F} \to \mathcal{G} \to \mathcal{H}$
在阿贝尔范畴 $\QCoh(\mathcal{O}_\mathcal{X})$ 中计算所得的上同调。
设 $x : U \to \mathcal{X}$ 是平坦态射，其中 $U$ 为概形。
由于采用 \'etale 拓扑，限制函子
$\textit{Mod}(\mathcal{X}_\etale, \mathcal{O}_\mathcal{X})
\to \textit{Mod}(U_\etale, \mathcal{O}_U)$ 正合。
另一方面，由引理 \ref{stacks-cohomology-lemma-flat-pullback-quasi-coherent} 和
《叠上的层》引理
\ref{stacks-sheaves-lemma-compare-quasi-coherent}，限制函子
$$
\QCoh(\mathcal{O}_\mathcal{X}) \xrightarrow{x^*}
\QCoh((\Sch/U)_\etale, \mathcal{O}) \xrightarrow{{-}|_{U_\etale}}
\QCoh(U_\etale, \mathcal{O}_U)
$$
也正合。因此
$\mathcal{E}|_{U_\etale} = (\Ker(\beta)/\Im(\alpha))|_{U_\etale}$。

\medskip\noindent
若 (1) 成立，则 $\mathcal{E} = 0$；所以对所有在
$\mathcal{X}$ 上平坦的 $U$，$\Ker(\beta)/\Im(\alpha)$
在 $U_\etale$ 上的限制为零，这正是寄生模的定义。
若 (2) 成立，则对所有在 $\mathcal{X}$ 上平坦的 $U$，
$\Ker(\beta)/\Im(\alpha)$ 在 $U_\etale$ 上的限制为零，
从而对所有在 $\mathcal{X}$ 上平坦的 $U$，
$\mathcal{E}$ 在 $U_\etale$ 上的限制为零。
这当然蕴含拟凝聚模 $\mathcal{E}$ 为零；例如，可对映射
$0 \to \mathcal{E}$ 应用引理 \ref{stacks-cohomology-lemma-quasi-coherent-check-exact}。
\end{proof}


\section{拟凝聚模}
\label{stacks-cohomology-section-quasi-coherent}

\noindent
我们已经看到，代数叠上的拟凝聚模范畴与其一个表示上的拟凝聚模范畴等价；
参见《叠上的层》第
\ref{stacks-sheaves-section-quasi-coherent-algebraic-stacks} 节。
这一事实是下述结果的基础。

\begin{lemma}
\label{stacks-cohomology-lemma-adjoint}
设 $\mathcal{X}$ 为代数叠，并设
$\textit{LQCoh}^{fbc}(\mathcal{O}_\mathcal{X})$
为具有平坦基变换性质的局部拟凝聚模范畴；参见第
\ref{stacks-cohomology-section-loc-qcoh-flat-base-change} 节。包含函子
$i : \QCoh(\mathcal{O}_\mathcal{X}) \to
\textit{LQCoh}^{fbc}(\mathcal{O}_\mathcal{X})$
有右伴随
$$
Q : \textit{LQCoh}^{fbc}(\mathcal{O}_\mathcal{X}) \to
\QCoh(\mathcal{O}_\mathcal{X})
$$
并且 $Q \circ i$ 是恒等函子。
\end{lemma}

\begin{proof}
选择概形 $U$ 和满的光滑态射 $f : U \to \mathcal{X}$。
令 $R = U \times_\mathcal{X} U$，从而得到代数空间中的光滑群胚
$(U, R, s, t, c)$，满足 $\mathcal{X} = [U/R]$；
参见《代数叠》引理 \ref{algebraic-lemma-stack-presentation}。
我们用 $[U/R]$ 替换 $\mathcal{X}$。由《叠上的层》命题
\ref{stacks-sheaves-proposition-quasi-coherent}，存在等价
$$
q_1 :
\QCoh(U, R, s, t, c)
\longrightarrow
\QCoh(\mathcal{O}_\mathcal{X})
$$
下面按如下规则构造函子
$$
q_2 :
\textit{LQCoh}^{fbc}(\mathcal{O}_\mathcal{X})
\longrightarrow
\QCoh(U, R, s, t, c)
$$
若 $\mathcal{F}$ 是
$\textit{LQCoh}^{fbc}(\mathcal{O}_\mathcal{X})$ 的对象，则令
$$
q_2(\mathcal{F}) = (f^*\mathcal{F}|_{U_\etale}, \alpha)
$$
其中 $\alpha$ 为同构
$$
t_{small}^*(f^*\mathcal{F}|_{U_\etale})
\to
t^*f^*\mathcal{F}|_{R_\etale} \to
s^*f^*\mathcal{F}|_{R_\etale} \to
s_{small}^*(f^*\mathcal{F}|_{U_\etale})
$$
外侧两个态射是比较映射。注意，$q_2(\mathcal{F})$ 拟凝聚，
恰是因为 $\mathcal{F}$ 局部拟凝聚；而构造下降数据 $\alpha$ 时，
我们使用了（并且需要）平坦基变换性质。略去余圈条件成立的验证
（参见《空间中的群胚》定义
\ref{spaces-groupoids-definition-groupoid-module}）。
查看《叠上的层》命题
\ref{stacks-sheaves-proposition-quasi-coherent} 的证明可见，
$q_2 \circ i$ 是 $q_1$ 的拟逆。定义 $Q = q_1 \circ q_2$。
设 $\mathcal{F}$ 是
$\textit{LQCoh}^{fbc}(\mathcal{O}_\mathcal{X})$ 的对象，
$\mathcal{G}$ 是 $\QCoh(\mathcal{O}_\mathcal{X})$ 的对象。则
\begin{align*}
\Mor_{\textit{LQCoh}^{fbc}(\mathcal{O}_\mathcal{X})}
(i(\mathcal{G}), \mathcal{F})
& =
\Mor_{\QCoh(U, R, s, t, c)}(q_2(i(\mathcal{G})), q_2(\mathcal{F})) \\
& =
\Mor_{\QCoh(\mathcal{O}_\mathcal{X})}(\mathcal{G}, Q(\mathcal{F}))
\end{align*}
第一个等号是《叠上的层》引理
\ref{stacks-sheaves-lemma-map-from-quasi-coherent}；
% P12-ERR-0016: see ../control/ERRATA_CANDIDATES.jsonl
第二个等号成立，是因为 $q_1 \circ i$ 与 $q_2$ 是范畴的互拟逆等价。
断言 $Q \circ i \cong \text{id}$ 是 $i$ 充满忠实这一事实的形式推论。
\end{proof}

\begin{lemma}
\label{stacks-cohomology-lemma-adjoint-kernel-parasitic}
设 $\mathcal{X}$ 为代数叠，并设
$Q : \textit{LQCoh}^{fbc}(\mathcal{O}_\mathcal{X}) \to
\QCoh(\mathcal{O}_\mathcal{X})$
为引理 \ref{stacks-cohomology-lemma-adjoint} 中构造的函子。
\begin{enumerate}
\item $Q$ 的核恰好是
$\textit{LQCoh}^{fbc}(\mathcal{O}_\mathcal{X})$ 的寄生对象全体。
\item 对 $\textit{LQCoh}^{fbc}(\mathcal{O}_\mathcal{X})$ 的任意对象
$\mathcal{F}$，伴随映射 $Q(\mathcal{F}) \to \mathcal{F}$
的核和余核都寄生。
\item 函子 $Q$ 正合，并与所有极限和余极限可交换。
\end{enumerate}
\end{lemma}

\begin{proof}
如引理 \ref{stacks-cohomology-lemma-adjoint} 的证明那样，写
$\mathcal{X} = [U/R]$。设 $\mathcal{F}$ 是
$\textit{LQCoh}^{fbc}(\mathcal{O}_\mathcal{X})$ 的对象。
由引理 \ref{stacks-cohomology-lemma-adjoint} 的证明显然可知，$\mathcal{F}$
属于 $Q$ 的核，当且仅当 $\mathcal{F}|_{U_\etale} = 0$。
特别地，若 $\mathcal{F}$ 寄生，则 $\mathcal{F}$ 属于该核。
接着，设 $x : V \to \mathcal{X}$ 是平坦态射，其中 $V$ 为概形。
令 $W = V \times_\mathcal{X} U$，并考察图
$$
\xymatrix{
W \ar[d]_p \ar[r]_q & V \ar[d] \\
U \ar[r] & \mathcal{X}
}
$$
注意，投影 $p : W \to U$ 平坦，而投影 $q : W \to V$ 光滑且满。
这蕴含 $q_{small}^*$ 在拟凝聚模上是忠实函子。
由假设，$\mathcal{F}$ 具有平坦基变换性质，故有
$p_{small}^*\mathcal{F}|_{U_\etale} \cong
q_{small}^*\mathcal{F}|_{V_\etale}$。
因此，若 $\mathcal{F}$ 属于 $Q$ 的核，则
$\mathcal{F}|_{V_\etale} = 0$，这就完成了 (1) 的证明。

\medskip\noindent
(2) 由上述讨论以及如下事实得出：映射
$Q(\mathcal{F}) \to \mathcal{F}$ 在限制到 $U_\etale$ 后成为同构。

\medskip\noindent
为证明 (3)，注意 $Q$ 作为右伴随是左正合的。设
$0 \to \mathcal{F} \to \mathcal{G} \to \mathcal{H} \to 0$
是 $\textit{LQCoh}^{fbc}(\mathcal{O}_\mathcal{X})$ 中的短正合列。
考察交换图
$$
\xymatrix{
0 \ar[r] &
Q(\mathcal{F}) \ar[r] \ar[d]_a &
Q(\mathcal{G}) \ar[r] \ar[d]_b &
Q(\mathcal{H}) \ar[r] \ar[d]_c & 0 \\
0 \ar[r] &
\mathcal{F} \ar[r] &
\mathcal{G} \ar[r] &
\mathcal{H} \ar[r] & 0
}
$$
由 (2)，$a$、$b$、$c$ 的核与余核都寄生；又因下行是短正合列，
可知上行作为 $\mathcal{O}_\mathcal{X}$-模复形，其上同调层寄生
（略去细节；这里用到寄生模范畴是全体模范畴的 Serre 子范畴）。
由 $Q$ 左正合，只需考察在 $Q(\mathcal{H})$ 处的正合性。
% P12-ERR-0017: see ../control/ERRATA_CANDIDATES.jsonl
然而，$Q(\mathcal{G}) \to Q(\mathcal{H}))$ 的余核
$\mathcal{Q}$，无论在 $\textit{Mod}(\mathcal{O}_\mathcal{X})$
还是在 $\QCoh(\mathcal{O}_\mathcal{X})$ 中计算，结果都相同，
因为包含函子
$\QCoh(\mathcal{O}_\mathcal{X}) \to
\textit{LQCoh}^{fbc}(\mathcal{O}_\mathcal{X})$
是左伴随，因而右正合。因此 $\mathcal{Q} = Q(\mathcal{Q})$
既拟凝聚又寄生，由 (1) 即为 $0$，正合所需。

\medskip\noindent
$Q$ 作为右伴随与所有极限可交换。由于 $Q$ 正合，
要证明 $Q$ 与所有余极限可交换，只需证明 $Q$ 与直和可交换；
参见《范畴》引理
\ref{categories-lemma-colimits-coproducts-coequalizers}。
设 $\mathcal{F}_i$（$i \in I$）是一族
$\textit{LQCoh}^{fbc}(\mathcal{O}_\mathcal{X})$ 的对象。
为说明 $Q(\bigoplus \mathcal{F}_i)$ 等于
$\bigoplus Q(\mathcal{F}_i)$，考察引理 \ref{stacks-cohomology-lemma-adjoint}
的证明中对 $Q$ 的构造。该构造使用表示
$\mathcal{X} = [U/R]$，其中 $U$ 是概形。此时，计算
$Q(\mathcal{F})$ 时先取 $\QCoh(U, R, s, t, c)$ 中的偶
$(\mathcal{F}|_{U_\etale}, \alpha)$，再使用等价
$\QCoh(U, R, s, t, c) \cong \QCoh(\mathcal{O}_\mathcal{X})$。
由于限制函子
$\textit{Mod}(\mathcal{O}_\mathcal{X}) \to
\textit{Mod}(\mathcal{O}_{U_\etale})$，
$\mathcal{F} \mapsto \mathcal{F}|_{U_\etale}$ 与直和可交换，
所需等式显然成立。
\end{proof}

\begin{lemma}
\label{stacks-cohomology-lemma-flat-morphism-and-Q}
设 $f : \mathcal{X} \to \mathcal{Y}$ 是代数叠之间的平坦态射。
则 $Q_\mathcal{X} \circ f^* = f^* \circ Q_\mathcal{Y}$，
其中 $Q_\mathcal{X}$ 和 $Q_\mathcal{Y}$ 如引理
\ref{stacks-cohomology-lemma-adjoint} 所述。
\end{lemma}

\begin{proof}
注意 $f^*$ 同时保持 $\QCoh$ 和 $\textit{LQCoh}^{fbc}$；
参见《叠上的层》引理
\ref{stacks-sheaves-lemma-pullback-quasi-coherent} 和命题
\ref{stacks-cohomology-proposition-loc-qcoh-flat-base-change}。
若 $\mathcal{F}$ 属于
$\textit{LQCoh}^{fbc}(\mathcal{O}_\mathcal{Y})$，
则由引理 \ref{stacks-cohomology-lemma-adjoint-kernel-parasitic}，
$Q_\mathcal{Y}(\mathcal{F}) \to \mathcal{F}$ 的核和余核都寄生。
由于 $f$ 平坦，由引理 \ref{stacks-cohomology-lemma-parasitic}，
$f^*Q_\mathcal{Y}(\mathcal{F}) \to f^*\mathcal{F}$ 的核和余核都寄生。
因而诱导映射
$f^*Q_\mathcal{Y}(\mathcal{F}) \to Q_\mathcal{X}(f^*\mathcal{F})$
的核和余核寄生，所以它是同构；例如可用引理
\ref{stacks-cohomology-lemma-exact-sequence-quasi-coherent-parasitic-cohomology}。
\end{proof}

\begin{lemma}
\label{stacks-cohomology-lemma-flat-object-and-Q}
设 $\mathcal{X}$ 为代数叠，并设 $x$ 是 $\mathcal{X}$ 的一个对象，
位于概形 $U$ 上，使得 $x : U \to \mathcal{X}$ 平坦。
% P12-ERR-0018: see ../control/ERRATA_CANDIDATES.jsonl
则对 $\QCoh^{fbc}(\mathcal{O}_\mathcal{X})$ 中的 $\mathcal{F}$，有
$Q(\mathcal{F})|_{U_\etale} = \mathcal{F}|_{U_\etale}$。
\end{lemma}

\begin{proof}
这是因为 $Q(\mathcal{F}) \to \mathcal{F}$ 的核和余核都寄生；
参见引理 \ref{stacks-cohomology-lemma-adjoint-kernel-parasitic}。
\end{proof}

\begin{remark}
\label{stacks-cohomology-remark-QCoh-abelian}
设 $\mathcal{X}$ 为代数叠。范畴 $\QCoh(\mathcal{O}_\mathcal{X})$
是阿贝尔范畴；包含函子
$\QCoh(\mathcal{O}_\mathcal{X}) \to \textit{Mod}(\mathcal{O}_\mathcal{X})$
右正合，但一般并不正合；参见《叠上的层》引理
\ref{stacks-sheaves-lemma-quasi-coherent-algebraic-stack}。
可以用引理 \ref{stacks-cohomology-lemma-adjoint} 和
\ref{stacks-cohomology-lemma-adjoint-kernel-parasitic} 中的函子 $Q$ 来理解这一点。
具体地，设 $\varphi : \mathcal{F} \to \mathcal{G}$
是拟凝聚 $\mathcal{O}_\mathcal{X}$-模之间的映射。则
\begin{enumerate}
\item 在 $\textit{Mod}(\mathcal{O}_\mathcal{X})$ 中计算的余核
$\Coker(\varphi)$ 拟凝聚，并且是 $\QCoh(\mathcal{O}_\mathcal{X})$
中 $\varphi$ 的余核；
\item 在 $\textit{Mod}(\mathcal{O}_\mathcal{X})$ 中计算的像
$\Im(\varphi)$ 拟凝聚，并且是 $\QCoh(\mathcal{O}_\mathcal{X})$
中 $\varphi$ 的像；
\item 在 $\textit{Mod}(\mathcal{O}_\mathcal{X})$ 中计算的核
$\Ker(\varphi)$，由命题
\ref{stacks-cohomology-proposition-loc-qcoh-flat-base-change} 属于
$\textit{LQCoh}^{fbc}(\mathcal{O}_\mathcal{X})$；而
$Q(\Ker(\varphi))$ 是 $\QCoh(\mathcal{O}_\mathcal{X})$ 中的核。
\end{enumerate}
这些结论由所给引文得出。
\end{remark}

\begin{remark}
\label{stacks-cohomology-remark-QCoh-tensor}
设 $\mathcal{X}$ 为代数叠。给定两个拟凝聚
$\mathcal{O}_\mathcal{X}$-模 $\mathcal{F}$ 和 $\mathcal{G}$，
张量积模 $\mathcal{F} \otimes_{\mathcal{O}_\mathcal{X}} \mathcal{G}$
拟凝聚；参见《叠上的层》引理
\ref{stacks-sheaves-lemma-quasi-coherent-algebraic-stack} 的 (5)。
类似地，给定两个具有平坦基变换性质的局部拟凝聚模，
其张量积也具有同样性质；参见命题
\ref{stacks-cohomology-proposition-loc-qcoh-flat-base-change}。因此包含函子
$$
\QCoh(\mathcal{O}_\mathcal{X}) \to
\textit{LQCoh}^{fbc}(\mathcal{O}_\mathcal{X}) \to
\textit{Mod}(\mathcal{O}_\mathcal{X})
$$
都是对称幺半范畴之间的函子。更值得注意的是，函子
$$
Q :
\textit{LQCoh}^{fbc}(\mathcal{O}_\mathcal{X})
\longrightarrow
\QCoh(\mathcal{O}_\mathcal{X})
$$
也是对称幺半范畴之间的函子。具体地，给定
$\textit{LQCoh}^{fbc}(\mathcal{O}_\mathcal{X})$ 中的
$\mathcal{F}$ 和 $\mathcal{G}$，得到
$$
\xymatrix{
Q(\mathcal{F})
\otimes_{\mathcal{O}_\mathcal{X}}
Q(\mathcal{G}) \ar[rr] \ar[rd] & &
\mathcal{F}
\otimes_{\mathcal{O}_\mathcal{X}}
\mathcal{G} \\
&
Q(\mathcal{F} \otimes_{\mathcal{O}_\mathcal{X}} \mathcal{G}) \ar[ru]
}
$$
其中西南方向的箭头来自西北方向箭头的泛性质
（以及左上角对象拟凝聚这一已经提及的事实）。
若把该图限制到 $U_\etale$，其中 $U \to \mathcal{X}$ 平坦，
则三个箭头都成为同构（参见引理 \ref{stacks-cohomology-lemma-adjoint}、
\ref{stacks-cohomology-lemma-adjoint-kernel-parasitic} 和定义
\ref{stacks-cohomology-definition-parasitic}）。因此
$Q(\mathcal{F}) \otimes_{\mathcal{O}_\mathcal{X}}
Q(\mathcal{G}) \to
Q(\mathcal{F} \otimes_{\mathcal{O}_\mathcal{X}} \mathcal{G})$
是同构；例如参见引理
\ref{stacks-cohomology-lemma-quasi-coherent-check-exact}。
\end{remark}

\begin{remark}
\label{stacks-cohomology-remark-bousfield-colocalization}
设 $\mathcal{X}$ 为代数叠，并以
$\textit{Parasitic}(\mathcal{O}_\mathcal{X}) \subset
\textit{Mod}(\mathcal{O}_\mathcal{X})$
表示由寄生模组成的充满子范畴。引理
\ref{stacks-cohomology-lemma-adjoint} 和 \ref{stacks-cohomology-lemma-adjoint-kernel-parasitic} 的结果蕴含
$$
\QCoh(\mathcal{O}_\mathcal{X}) =
\textit{LQCoh}^{fbc}(\mathcal{O}_\mathcal{X}) /
\textit{Parasitic}(\mathcal{O}_\mathcal{X})
\cap \textit{LQCoh}^{fbc}(\mathcal{O}_\mathcal{X})
$$
换言之，拟凝聚模范畴是具有平坦基变换性质的局部拟凝聚模范畴，
除以由寄生对象组成的 Serre 子范畴所得的商范畴。
参见《同调》引理
\ref{homology-lemma-serre-subcategory-is-kernel}。
该情形的一个关键特征是存在包含函子
$i : \QCoh(\mathcal{O}_\mathcal{X}) \to
\textit{LQCoh}^{fbc}(\mathcal{O}_\mathcal{X})$，
它是商函子的左伴随。在《叠的导出范畴》第
\ref{stacks-perfect-section-derived} 节，尤其是引理
\ref{stacks-perfect-lemma-bousfield-colocalization} 中，
我们证明在导出范畴层次上也有类似结果。
\end{remark}

\begin{lemma}
\label{stacks-cohomology-lemma-internal-hom-fp-into-qcoh}
设 $\mathcal{X}$ 为代数叠，$\mathcal{F}$ 是有限表示的
$\mathcal{O}_\mathcal{X}$-模，$\mathcal{G}$ 是拟凝聚
$\mathcal{O}_\mathcal{X}$-模。在
$\textit{Mod}(\mathcal{X}_\etale, \mathcal{O}_\mathcal{X})$ 或
$\textit{Mod}(\mathcal{O}_\mathcal{X})$ 中计算的内部 Hom
$\SheafHom_{\mathcal{O}_\mathcal{X}}(\mathcal{F}, \mathcal{G})$
彼此一致，其共同值是
$\textit{LQCoh}^{fbc}(\mathcal{O}_\mathcal{X})$ 的对象。
拟凝聚模
$
hom(\mathcal{F}, \mathcal{G}) =
Q(\SheafHom_{\mathcal{O}_\mathcal{X}}(\mathcal{F}, \mathcal{G}))
$
具有如下泛性质：
$$
\Hom_\mathcal{X}(\mathcal{H}, hom(\mathcal{F}, \mathcal{G})) =
\Hom_\mathcal{X}(\mathcal{H} \otimes_{\mathcal{O}_\mathcal{X}} \mathcal{F},
\mathcal{G})
$$
其中 $\mathcal{H}$ 属于 $\QCoh(\mathcal{O}_\mathcal{X})$。
\end{lemma}

\begin{proof}
《site 上的模》第 \ref{sites-modules-section-internal-hom} 节中
$\SheafHom_{\mathcal{O}_\mathcal{X}}(\mathcal{F}, \mathcal{G})$
的构造，只依赖于作为模预层的 $\mathcal{F}$ 和 $\mathcal{G}$；
输出的 $\SheafHom$ 是 fppf 拓扑的层，因为假设
$\mathcal{F}$ 和 $\mathcal{G}$ 是 fppf 拓扑的层；
参见《site 上的模》引理 \ref{sites-modules-lemma-internal-hom}。
由《叠上的层》引理 \ref{stacks-sheaves-lemma-lqc-colimits}，
$\SheafHom_{\mathcal{O}_\mathcal{X}}(\mathcal{F}, \mathcal{G})$
局部拟凝聚。由引理
\ref{stacks-cohomology-lemma-check-flat-comparison-on-etale-covering}，
$\SheafHom_{\mathcal{O}_\mathcal{X}}(\mathcal{F}, \mathcal{G})$
具有平坦基变换性质。因此
$\SheafHom_{\mathcal{O}_\mathcal{X}}(\mathcal{F}, \mathcal{G})$ 是
$\textit{LQCoh}^{fbc}(\mathcal{O}_\mathcal{X})$ 的对象，
可以对其应用引理 \ref{stacks-cohomology-lemma-adjoint} 的函子 $Q$。
由 $Q$ 的泛性质，对拟凝聚的 $\mathcal{H}$ 有
$$
\Hom_\mathcal{X}(\mathcal{H},
Q(\SheafHom_{\mathcal{O}_\mathcal{X}}(\mathcal{F}, \mathcal{G}))) =
\Hom_\mathcal{X}(\mathcal{H},
\SheafHom_{\mathcal{O}_\mathcal{X}}(\mathcal{F}, \mathcal{G}))
$$
所以引理中的显示公式由《site 上的模》引理
\ref{sites-modules-lemma-internal-hom-adjoint-tensor} 得出。
\end{proof}

\begin{lemma}
\label{stacks-cohomology-lemma-flat-morphism-and-hom}
设 $f : \mathcal{X} \to \mathcal{Y}$ 是代数叠之间的平坦态射。
设 $\mathcal{F}$ 是有限表示的 $\mathcal{O}_\mathcal{Y}$-模，
$\mathcal{G}$ 是拟凝聚 $\mathcal{O}_\mathcal{Y}$-模。
则 $f^*hom(\mathcal{F}, \mathcal{G}) = hom(f^*\mathcal{F}, f^*\mathcal{G})$，
记号如引理 \ref{stacks-cohomology-lemma-internal-hom-fp-into-qcoh}。
\end{lemma}

\begin{proof}
由《site 上的模》引理
\ref{sites-modules-lemma-pullback-internal-hom}，有
$f^*\SheafHom_{\mathcal{O}_\mathcal{Y}}(\mathcal{F}, \mathcal{G}) =
\SheafHom_{\mathcal{O}_\mathcal{X}}(f^*\mathcal{F}, f^*\mathcal{G})$。
（注意，这一步没有用到 $f$ 的平坦性，因为与 $f$ 相伴的赋环拓扑斯态射
总是平坦的；参见《叠上的层》注
\ref{stacks-sheaves-remark-flat}。）
然后应用引理 \ref{stacks-cohomology-lemma-flat-morphism-and-Q}
（而在这里确实用到了 $f$ 的平坦性）。
\end{proof}







\section{拟凝聚模的直像}
\label{stacks-cohomology-section-pushforward-quasi-coherent}

\noindent
设 $f : \mathcal{X} \to \mathcal{Y}$ 是代数叠之间的态射。
考虑直像
$$
f_* :
\textit{Mod}(\mathcal{O}_\mathcal{X})
\longrightarrow
\textit{Mod}(\mathcal{O}_\mathcal{Y})
$$
事实表明，该函子几乎从不保持拟凝聚层的子范畴。
例如，考虑概形态射
$$
j : X = \mathbf{A}^2_k \setminus \{0\} \longrightarrow \mathbf{A}^2_k = Y.
$$
与之相伴的代数叠态射为
$$
f = j_{big} : \mathcal{X} = (\Sch/X)_{fppf} \to
(\Sch/Y)_{fppf} = \mathcal{Y}
$$
结构层的直像 $f_*\mathcal{O}_\mathcal{X}$ 的全局截面为
$k[x, y]$。因此，若 $f_*\mathcal{O}_\mathcal{X}$ 在
$\mathcal{Y}$ 上拟凝聚，则必有
$f_*\mathcal{O}_\mathcal{X} = \mathcal{O}_\mathcal{Y}$。然而，
考虑映到 $0$ 的 $T = \Spec(k) \to \mathbf{A}^2_k = Y$。
因为 $X \times_Y T = \emptyset$，有
$\Gamma(T, f_*\mathcal{O}_\mathcal{X}) = 0$；而
$\Gamma(T, \mathcal{O}_\mathcal{Y}) = k$。
积极的一面是，对任意平坦态射 $T \to Y$，有等式
$\Gamma(T, f_*\mathcal{O}_\mathcal{X}) = \Gamma(T, \mathcal{O}_\mathcal{Y})$；
这由《概形上的上同调》引理
\ref{coherent-lemma-flat-base-change-cohomology} 得出，
其中用到 $j$ 是拟紧且拟分离的。

\medskip\noindent
设 $f : \mathcal{X} \to \mathcal{Y}$ 是代数叠之间拟紧且
拟分离的态射。我们利用以下三点来规避上述问题：
\begin{enumerate}
\item $f_*$ 确实保持局部拟凝聚模
（引理 \ref{stacks-cohomology-lemma-pushforward-locally-quasi-coherent}）；
\item $f_*$ 把拟凝聚层变为局部拟凝聚层，
且其平坦比较态射均为同构
（引理 \ref{stacks-cohomology-lemma-flat-comparison}）；
\item 具有平坦基变换性质的局部拟凝聚
$\mathcal{O}_\mathcal{Y}$-模，在 $\mathcal{Y}$ 的一个展示上给出
拟凝聚模，从而在 $\mathcal{Y}$ 上给出拟凝聚模；
参见《叠上的层》第
\ref{stacks-sheaves-section-quasi-coherent-algebraic-stacks} 节。
\end{enumerate}
因此得到函子
$$
f_{\QCoh, *} :
\QCoh(\mathcal{O}_\mathcal{X})
\longrightarrow
\QCoh(\mathcal{O}_\mathcal{Y})
$$
它是
$f^* : \QCoh(\mathcal{O}_\mathcal{Y}) \to
\QCoh(\mathcal{O}_\mathcal{X})$
的右伴随，并且还满足
$$
\Gamma(y, f_*\mathcal{F}) = \Gamma(y, f_{\QCoh, *}\mathcal{F})
$$
其中 $y \in \Ob(\mathcal{Y})$ 且相应的
$1$-态射 $y : V \to \mathcal{Y}$ 平坦；参见引理
\ref{stacks-cohomology-lemma-direct-image-quasi-coherent-over-flat}。
此外，类似的构造将给出函子
$R^if_{\QCoh, *}$。
但这些结果还不足以构造总直像函子
（作用于具有拟凝聚上同调层的复形）。

\begin{proposition}
\label{stacks-cohomology-proposition-direct-image-quasi-coherent}
设 $f : \mathcal{X} \to \mathcal{Y}$ 是代数叠之间拟紧且拟分离的态射。
函子
$f^* : \QCoh(\mathcal{O}_\mathcal{Y}) \to
\QCoh(\mathcal{O}_\mathcal{X})$
具有右伴随
$$
f_{\QCoh, *} :
\QCoh(\mathcal{O}_\mathcal{X})
\longrightarrow
\QCoh(\mathcal{O}_\mathcal{Y})
$$
它可定义为复合
$$
\QCoh(\mathcal{O}_\mathcal{X}) \to
\textit{LQCoh}^{fbc}(\mathcal{O}_\mathcal{X})
\xrightarrow{f_*} \textit{LQCoh}^{fbc}(\mathcal{O}_\mathcal{Y})
\xrightarrow{Q} \QCoh(\mathcal{O}_\mathcal{Y})
$$
其中函子 $f_*$ 和 $Q$ 分别如命题
\ref{stacks-cohomology-proposition-loc-qcoh-flat-base-change} 与引理
\ref{stacks-cohomology-lemma-adjoint} 所述。
此外，若把 $R^if_{\QCoh, *}$ 定义为复合
$$
\QCoh(\mathcal{O}_\mathcal{X}) \to
\textit{LQCoh}^{fbc}(\mathcal{O}_\mathcal{X})
\xrightarrow{R^if_*} \textit{LQCoh}^{fbc}(\mathcal{O}_\mathcal{Y})
\xrightarrow{Q} \QCoh(\mathcal{O}_\mathcal{Y})
$$
则函子序列 $\{R^if_{\QCoh, *}\}_{i \geq 0}$
构成一个上同调 $\delta$-函子。
\end{proposition}

\begin{proof}
这是命题中所述结果的组合。
伴随性可如下证明：设 $\mathcal{F}$ 是拟凝聚
$\mathcal{O}_\mathcal{X}$-模，$\mathcal{G}$ 是拟凝聚
$\mathcal{O}_\mathcal{Y}$-模。则有
\begin{align*}
\Mor_{\QCoh(\mathcal{O}_\mathcal{X})}(f^*\mathcal{G}, \mathcal{F})
& =
\Mor_{\textit{LQCoh}^{fbc}(\mathcal{O}_\mathcal{Y})}
(\mathcal{G}, f_*\mathcal{F}) \\
& =
\Mor_{\QCoh(\mathcal{O}_\mathcal{Y})}(\mathcal{G}, Q(f_*\mathcal{F}))
\\
& =
\Mor_{\QCoh(\mathcal{O}_\mathcal{Y})}(\mathcal{G},
f_{\QCoh, *}\mathcal{F})
\end{align*}
第一个等号由 $f_*$ 与 $f^*$ 的伴随性得到
（对任意模层而言）。由命题
\ref{stacks-cohomology-proposition-loc-qcoh-flat-base-change}，
$f_*\mathcal{F}$ 是
$\textit{LQCoh}^{fbc}(\mathcal{O}_\mathcal{Y})$
的对象（且可以在 fppf 拓扑或 \'etale 拓扑中计算）；
第二个等号由引理 \ref{stacks-cohomology-lemma-adjoint} 得到。
第三个等号是 $f_{\QCoh, *}$ 的定义。

\medskip\noindent
要说明 $\{R^if_{\QCoh, *}\}_{i \geq 0}$ 是《同调》定义
\ref{homology-definition-cohomological-delta-functor}
中的上同调 $\delta$-函子，设
$$
0 \to \mathcal{F}_1 \to \mathcal{F}_2 \to \mathcal{F}_3 \to 0
$$
是 $\QCoh(\mathcal{O}_\mathcal{X})$ 中的短正合列。
它在 $\textit{Mod}(\mathcal{O}_\mathcal{X})$ 中可能不是正合列，
但我们知道，它在忽略寄生模的意义下是正合的；
参见引理 \ref{stacks-cohomology-lemma-exact-sequence-quasi-coherent-parasitic-cohomology}。
因此可将该序列分解为
$$
\begin{matrix}
0 \to \mathcal{P}_1 \to \mathcal{F}_1 \to \mathcal{I}_2 \to 0 \\
0 \to \mathcal{I}_2 \to \mathcal{F}_2 \to \mathcal{Q}_2 \to 0 \\
0 \to \mathcal{P}_2 \to \mathcal{Q}_2 \to \mathcal{I}_3 \to 0 \\
0 \to \mathcal{I}_3 \to \mathcal{F}_3 \to \mathcal{P}_3 \to 0
\end{matrix}
$$
这些是 $\textit{Mod}(\mathcal{O}_\mathcal{X})$ 中的短正合列，
且 $\mathcal{P}_i$ 是寄生的。注意，每个层
$\mathcal{P}_j$、$\mathcal{I}_j$、$\mathcal{Q}_j$ 都是
$\textit{LQCoh}^{fbc}(\mathcal{O}_\mathcal{X})$ 的对象；
参见命题 \ref{stacks-cohomology-proposition-loc-qcoh-flat-base-change}。
应用 $R^if_*$ 得到长正合列
$$
\begin{matrix}
0 \to f_*\mathcal{P}_1 \to f_*\mathcal{F}_1 \to f_*\mathcal{I}_2 \to
R^1f_*\mathcal{P}_1 \to \ldots \\
0 \to f_*\mathcal{I}_2 \to f_*\mathcal{F}_2 \to f_*\mathcal{Q}_2 \to
R^1f_*\mathcal{I}_2 \to \ldots \\
0 \to f_*\mathcal{P}_2 \to f_*\mathcal{Q}_2 \to f_*\mathcal{I}_3 \to
R^1f_*\mathcal{P}_2 \to \ldots \\
0 \to f_*\mathcal{I}_3 \to f_*\mathcal{F}_3 \to f_*\mathcal{P}_3 \to
R^1f_*\mathcal{I}_3 \to \ldots
\end{matrix}
$$
% P12-ERR-0022: see ../control/ERRATA_CANDIDATES.jsonl
其中各项由命题
\ref{stacks-cohomology-proposition-loc-qcoh-flat-base-change}
都是 $\textit{LQCoh}^{fbc}(\mathcal{O}_\mathcal{Y})$ 的对象。
由引理 \ref{stacks-cohomology-lemma-pushforward-parasitic}，
层 $R^if_*\mathcal{P}_j$ 是寄生的，因而在应用函子
$Q$ 后消失；参见引理
\ref{stacks-cohomology-lemma-adjoint-kernel-parasitic}。
由于 $Q$ 正合，态射
$$
Q(R^if_*\mathcal{F}_3)
\cong
Q(R^if_*\mathcal{I}_3)
\cong
Q(R^if_*\mathcal{Q}_2)
\rightarrow
Q(R^{i + 1}f_*\mathcal{I}_2)
\cong
Q(R^{i + 1}f_*\mathcal{F}_1)
$$
可作为连接态射，从而使函子族
$\{R^if_{\QCoh, *}\}_{i \geq 0}$
成为上同调 $\delta$-函子。
\end{proof}

\begin{lemma}
\label{stacks-cohomology-lemma-direct-image-quasi-coherent-over-flat}
设 $f : \mathcal{X} \to \mathcal{Y}$ 是代数叠之间拟紧且拟分离的态射。
设 $y : V \to \mathcal{Y}$ 属于 $\Ob(\mathcal{Y})$，且 $y$ 为平坦态射。
设 $\mathcal{F}$ 属于 $\QCoh(\mathcal{O}_\mathcal{X})$。
则对所有 $i \in \mathbf{Z}$，有
$(f_*\mathcal{F})(y) = (f_{\QCoh, *}\mathcal{F})(y)$
以及 $(R^if_*\mathcal{F})(y) = (R^if_{\QCoh, *}\mathcal{F})(y)$。
\end{lemma}

\begin{proof}
这由命题 \ref{stacks-cohomology-proposition-direct-image-quasi-coherent}
中函子 $R^if_{\QCoh, *}$ 的构造、
定义 \ref{stacks-cohomology-definition-parasitic} 中寄生模的定义，
以及引理 \ref{stacks-cohomology-lemma-adjoint-kernel-parasitic} 的第 (2) 部分得出。
\end{proof}

\begin{remark}
\label{stacks-cohomology-remark-direct-image-quasi-coherent-tensor}
设 $f : \mathcal{X} \to \mathcal{Y}$ 是代数叠之间拟紧且拟分离的态射。
设 $\mathcal{F}$ 和 $\mathcal{G}$ 属于
$\QCoh(\mathcal{O}_\mathcal{X})$。则有典范交换图
$$
\xymatrix{
f_{\QCoh, *}\mathcal{F}
\otimes_{\mathcal{O}_\mathcal{Y}}
f_{\QCoh, *}\mathcal{G} \ar[r] \ar[d] &
f_*\mathcal{F}
\otimes_{\mathcal{O}_\mathcal{Y}}
f_*\mathcal{G} \ar[d]^c \\
f_{\QCoh, *}(\mathcal{F}
\otimes_{\mathcal{O}_\mathcal{X}}
\mathcal{G}) \ar[r] &
f_*(\mathcal{F}
\otimes_{\mathcal{O}_\mathcal{X}}
\mathcal{G})
}
$$
右边的竖直箭头 $c$ 是（次数 $0$ 的）朴素相对杯积；
参见《site 上的上同调》第
\ref{sites-cohomology-section-cup-product} 节。
% P12-ERR-0019: see ../control/ERRATA_CANDIDATES.jsonl
$c$ 的源与目标都属于
$\textit{LQCoh}^{fbc}(\mathcal{O}_\mathcal{X})$；参见命题
\ref{stacks-cohomology-proposition-loc-qcoh-flat-base-change}。
对 $c$ 应用 $Q$，便得到左边的竖直箭头，
因为 $Q$ 与张量积可交换；参见注
\ref{stacks-cohomology-remark-QCoh-tensor}。
该构造对 $\mathcal{F}$ 和 $\mathcal{G}$ 具有函子性。
\end{remark}

\begin{lemma}
\label{stacks-cohomology-lemma-leray}
设 $f : \mathcal{X} \to \mathcal{Y}$
是代数叠之间拟紧且拟分离的态射。
设 $\mathcal{F}$ 是 $\mathcal{X}$ 上的拟凝聚层。
则存在一个谱序列，其 $E_2$-页为
$$
E_2^{p, q} = H^p(\mathcal{Y}, R^qf_{\QCoh, *}\mathcal{F})
$$
并收敛到 $H^{p + q}(\mathcal{X}, \mathcal{F})$。
\end{lemma}

\begin{proof}
由《site 上的上同调》引理
\ref{sites-cohomology-lemma-Leray}，具有
$$
E_2^{p, q} = H^p(\mathcal{Y}, R^qf_*\mathcal{F})
$$
的 Leray 谱序列收敛到
$H^{p + q}(\mathcal{X}, \mathcal{F})$。伴随态射
$$
R^qf_{\QCoh, *}\mathcal{F} \longrightarrow R^qf_*\mathcal{F}
$$
的核与余核是 $\mathcal{Y}$ 上的寄生模
（引理 \ref{stacks-cohomology-lemma-adjoint-kernel-parasitic}），
因而其上同调为零
（引理 \ref{stacks-cohomology-lemma-pushforward-parasitic}）。
形式地可得
$H^p(\mathcal{Y}, R^qf_{\QCoh, *}\mathcal{F}) =
H^p(\mathcal{Y}, R^qf_*\mathcal{F})$，结论成立。
\end{proof}

\begin{lemma}
\label{stacks-cohomology-lemma-relative-leray}
设 $f : \mathcal{X} \to \mathcal{Y}$ 和 $g : \mathcal{Y} \to \mathcal{Z}$
是代数叠之间拟紧且拟分离的态射。
设 $\mathcal{F}$ 是 $\mathcal{X}$ 上的拟凝聚层。
则存在一个谱序列，其 $E_2$-页为
$$
E_2^{p, q} = R^pg_{\QCoh, *}(R^qf_{\QCoh, *}\mathcal{F})
$$
并收敛到 $R^{p + q}(g \circ f)_{\QCoh, *}\mathcal{F}$。
\end{lemma}

\begin{proof}
由《site 上的上同调》引理
\ref{sites-cohomology-lemma-relative-Leray}，具有
$$
E_2^{p, q} = R^pg_*(R^qf_*\mathcal{F})
$$
的 Leray 谱序列收敛到
$R^{p + q}(g \circ f)_*\mathcal{F}$。由命题
\ref{stacks-cohomology-proposition-loc-qcoh-flat-base-change} 的结果，
该谱序列的所有项都是
$\textit{LQCoh}^{fbc}(\mathcal{O}_\mathcal{Z})$ 的对象。
应用正合函子
$Q_\mathcal{Z} : \textit{LQCoh}^{fbc}(\mathcal{O}_\mathcal{Z}) \to
\QCoh(\mathcal{O}_\mathcal{Z})$，得到
% P12-ERR-0023: see ../control/ERRATA_CANDIDATES.jsonl
$\QCoh(\mathcal{O}_\mathcal{Z})$ 中一个收敛到
$R^{p + q}(g \circ f)_{\QCoh, *}\mathcal{F}$ 的谱序列。
因此，只要证明
% P12-ERR-0020: see ../control/ERRATA_CANDIDATES.jsonl
% P12-ERR-0021: see ../control/ERRATA_CANDIDATES.jsonl
$$
Q_\mathcal{Z}(R^pg_*(R^qf_*\mathcal{F})) =
Q_\mathcal{Z}(R^pg_*(Q_\mathcal{X}(R^qf_*\mathcal{F}))
$$
即可。这又由以下事实得出：态射
$$
Q_\mathcal{X}(R^qf_*\mathcal{F}) \longrightarrow R^qf_*\mathcal{F}
$$
的核与余核是寄生的
（引理 \ref{stacks-cohomology-lemma-adjoint-kernel-parasitic}），而
$R^pg_*$ 把寄生模变为寄生模
（引理 \ref{stacks-cohomology-lemma-pushforward-parasitic}）。
\end{proof}

\noindent
在本节最后，我们把与概形给出的光滑覆盖相联系的谱序列明确写出。
请与《叠上的层》第
\ref{stacks-sheaves-section-cohomology} 节和第
\ref{stacks-sheaves-section-higher-direct-images} 节比较。

\begin{proposition}
\label{stacks-cohomology-proposition-smooth-covering-compute-cohomology}
设 $f : \mathcal{U} \to \mathcal{X}$ 是代数叠之间的态射。
假设 $f$ 可由代数空间表示，满射、平坦且局部有限表示。
设 $\mathcal{F}$ 是拟凝聚 $\mathcal{O}_\mathcal{X}$-模。
则有谱序列
$$
E_2^{p, q} = H^q(\mathcal{U}_p, f_p^*\mathcal{F})
\Rightarrow
H^{p + q}(\mathcal{X}, \mathcal{F})
$$
其中 $f_p$ 是态射
$\mathcal{U} \times_\mathcal{X} \ldots \times_\mathcal{X} \mathcal{U} \to
\mathcal{X}$（含 $p + 1$ 个因子）。
\end{proposition}

\begin{proof}
这是《叠上的层》命题
\ref{stacks-sheaves-proposition-smooth-covering-compute-cohomology}
的一个特例。
\end{proof}

\begin{proposition}
\label{stacks-cohomology-proposition-smooth-covering-compute-direct-image}
设 $f : \mathcal{U} \to \mathcal{X}$ 和 $g : \mathcal{X} \to \mathcal{Y}$
是可复合的代数叠态射。假设
\begin{enumerate}
\item $f$ 可由代数空间表示，满射、平坦、局部有限表示、
拟紧且拟分离；
\item $g$ 拟紧且拟分离。
\end{enumerate}
若 $\mathcal{F}$ 属于 $\QCoh(\mathcal{O}_\mathcal{X})$，
则 $\QCoh(\mathcal{O}_\mathcal{Y})$ 中有谱序列
$$
E_2^{p, q} = R^q(g \circ f_p)_{\QCoh, *}f_p^*\mathcal{F}
\Rightarrow
R^{p + q}g_{\QCoh, *}\mathcal{F}
$$
。
\end{proposition}

\begin{proof}
注意，每个态射
$f_p : \mathcal{U} \times_\mathcal{X} \ldots \times_\mathcal{X} \mathcal{U} \to
\mathcal{X}$ 都拟紧且拟分离，因而 $g \circ f_p$
也拟紧且拟分离。故命题中的表述有意义
（即函子 $R^q(g \circ f_p)_{\QCoh, *}$ 已定义）。
由《叠上的层》命题
\ref{stacks-sheaves-proposition-smooth-covering-compute-direct-image}，
有谱序列
$$
E_2^{p, q} = R^q(g \circ f_p)_*f_p^{-1}\mathcal{F}
\Rightarrow
R^{p + q}g_*\mathcal{F}
$$
。应用正合函子
$Q_\mathcal{Y} : \textit{LQCoh}^{fbc}(\mathcal{O}_\mathcal{Y}) \to
\QCoh(\mathcal{O}_\mathcal{Y})$，便得到
$\QCoh(\mathcal{O}_\mathcal{Y})$ 中所需的谱序列。
\end{proof}

\section{关于拟凝聚模的进一步说明}
\label{stacks-cohomology-section-further-remarks}

\noindent
% P12-ERR-0024: see ../control/ERRATA_CANDIDATES.jsonl
本节汇集一些结果，以帮助理解如何使用代数叠上的拟凝聚模。

\medskip\noindent
设 $f : \mathcal{U} \to \mathcal{X}$ 是代数叠之间的态射。
假设 $\mathcal{U}$ 由代数空间 $U$ 表示。
考虑由拉回（《叠上的层》第
\ref{stacks-sheaves-section-modules} 节）之后再限制
（《叠上的层》第
\ref{stacks-sheaves-section-restriction-algebraic-spaces} 节）给出的函子
$$
a :
\textit{Mod}(\mathcal{X}_\etale, \mathcal{O}_\mathcal{X})
\longrightarrow
\textit{Mod}(U_\etale, \mathcal{O}_U),\quad
\mathcal{F}
\longmapsto
f^*\mathcal{F}|_{U_\etale}
$$
。把该函子应用于局部拟凝聚模，得到函子
$$
b : \textit{LQCoh}(\mathcal{O}_\mathcal{X})
\longrightarrow
\QCoh(U_\etale, \mathcal{O}_U)
$$
参见《叠上的层》引理
\ref{stacks-sheaves-lemma-pullback-lqc} 和
\ref{stacks-sheaves-lemma-compare-locally-quasi-coherent}。
进一步将函子限制到更小的子范畴，可得
$$
c :
\textit{LQCoh}^{fbc}(\mathcal{O}_\mathcal{X})
\longrightarrow
\QCoh(U_\etale, \mathcal{O}_U)
$$
以及
$$
d :
\QCoh(\mathcal{O}_\mathcal{X})
\longrightarrow
\QCoh(U_\etale, \mathcal{O}_U)
$$
关于这些函子，我们有以下结论：\footnote{我们建议在纸巾上
推导为何这些陈述成立，而不是追随所给的参照。}
\begin{enumerate}
\item 函子 $a$ 正合。事实上，拉回 $f^* = f^{-1}$ 正合
（《叠上的层》第 \ref{stacks-sheaves-section-modules} 节），
而到 $U_\etale$ 的限制也正合；参见《叠上的层》等式
(\ref{stacks-sheaves-equation-restrict})。
\item 函子 $b$ 正合。事实上，由《叠上的层》引理
\ref{stacks-sheaves-lemma-lqc-colimits}，包含
$\textit{LQCoh}(\mathcal{O}_\mathcal{X}) \to
\textit{Mod}(\mathcal{X}_\etale, \mathcal{O}_\mathcal{X})$ 正合。
\item 函子 $c$ 正合。事实上，由命题
\ref{stacks-cohomology-proposition-loc-qcoh-flat-base-change}，包含函子
$\textit{LQCoh}^{fbc}(\mathcal{O}_\mathcal{X}) \to
\textit{Mod}(\mathcal{X}_\etale, \mathcal{O}_\mathcal{X})$ 正合。
\item 函子 $d$ 右正合，但一般不正合。事实上，
由《叠上的层》引理
\ref{stacks-sheaves-lemma-qc-colimits}，包含函子
$\QCoh(\mathcal{O}_\mathcal{X}) \to
\textit{Mod}(\mathcal{X}_\etale, \mathcal{O}_\mathcal{X})$ 右正合。
我们略去说明非正合性的例子。
\item 若 $f$ 平坦，则 $d$ 正合。这由结合引理
\ref{stacks-cohomology-lemma-flat-pullback-quasi-coherent} 和《叠上的层》引理
\ref{stacks-sheaves-lemma-compare-quasi-coherent} 得出。
\item 若 $f$ 平坦，则 $c$ 消去寄生对象。
% P12-ERR-0025: see ../control/ERRATA_CANDIDATES.jsonl
事实上，由引理 \ref{stacks-cohomology-lemma-parasitic}，$f^*$ 保持寄生对象。
因此，对任意在 $U$ 上 \'etale、从而在 $\mathcal{X}$ 上平坦的概形
$V$，由限制与 \'etale 局部化的相容性，有
$0 = f^*\mathcal{F}|_{V_\etale} = c(\mathcal{F})|_{V_\etale}$；
参见《叠上的层》注
\ref{stacks-sheaves-remark-compare}。故显然 $c(\mathcal{F}) = 0$。
\item 若 $f$ 平坦，则 $c = d \circ Q$。事实上，
由引理 \ref{stacks-cohomology-lemma-adjoint-kernel-parasitic}，
$Q(\mathcal{F}) \to \mathcal{F}$ 的核与余核是寄生的。
因此，由于 $c$ 正合（(3)）并消去寄生对象（(6)），
对 $Q(\mathcal{F}) \to \mathcal{F}$ 应用 $c$ 得到同构。
\item 函子 $a, b, c, d$ 与余极限和任意直和可交换。
对 $f^*$ 和限制而言，这由它们是左伴随得出，
因而对 $a$ 成立。由上述参照，结论也对 $b$、$c$、$d$ 成立。
\item 函子 $a, b, c, d$ 与张量积可交换。
\item 若 $f$ 平坦且满射，$\mathcal{F}$ 属于
$\textit{LQCoh}^{fbc}(\mathcal{O}_\mathcal{X})$，且 $c(\mathcal{F}) = 0$，
则 $\mathcal{F}$ 是寄生的。事实上，由 (7) 得
$d(Q(\mathcal{F})) = 0$。由限制与 \'etale 局部化的相容性
（参见上述参照），可假设 $U$ 是概形。将引理
\ref{stacks-cohomology-lemma-quasi-coherent-check-exact} 应用于 $0 \to Q(\mathcal{F})$
和态射 $f : U \to \mathcal{X}$，可得 $Q(\mathcal{F}) = 0$。
因此，由引理 \ref{stacks-cohomology-lemma-adjoint-kernel-parasitic}，
$\mathcal{F}$ 是寄生的。
\item 若 $f$ 平坦且满射，则函子 $d$ 反映正合性。
更确切地，设 $\mathcal{F}^\bullet$ 是
$\QCoh(\mathcal{O}_\mathcal{X})$ 中的复形。
则 $\mathcal{F}^\bullet$ 在 $\QCoh(\mathcal{O}_\mathcal{X})$ 中正合，
当且仅当 $d(\mathcal{F}^\bullet)$ 正合。我们已在 (5) 中看到一个方向。
对另一方向，假设 $H^i(d(\mathcal{F}^\bullet)) = 0$。
则 $\mathcal{G} = H^i(\mathcal{F}^\bullet)$ 是
$\QCoh(\mathcal{O}_\mathcal{X})$ 中满足 $d(\mathcal{G}) = 0$ 的对象。
因此，由 (10)，$\mathcal{G}$ 同时拟凝聚且寄生，
从而为 $0$；例如参见注 \ref{stacks-cohomology-remark-bousfield-colocalization}。
\item
\label{stacks-cohomology-item-hom-restriction}
若 $f$ 平坦，
$\mathcal{F}, \mathcal{G} \in \Ob(\QCoh(\mathcal{O}_\mathcal{X}))$，
% P12-ERR-0026: see ../control/ERRATA_CANDIDATES.jsonl
且 $\mathcal{F}$ 有限表示，则有
$$
d(hom(\mathcal{F}, \mathcal{G})) =
\SheafHom_{\mathcal{O}_U}(d(\mathcal{F}), d(\mathcal{G}))
$$
记号如引理 \ref{stacks-cohomology-lemma-internal-hom-fp-into-qcoh}。
也许最容易的证法如下：
\begin{align*}
d(hom(\mathcal{F}, \mathcal{G}))
& =
d(Q(\SheafHom_{\mathcal{O}_\mathcal{X}}(\mathcal{F}, \mathcal{G}))) \\
& =
c(\SheafHom_{\mathcal{O}_\mathcal{X}}(\mathcal{F}, \mathcal{G})) \\
& =
f^*\SheafHom_{\mathcal{O}_\mathcal{X}}(\mathcal{F}, \mathcal{G})|_{U_\etale} \\
& =
\SheafHom_{\mathcal{O}_\mathcal{U}}(f^*\mathcal{F},
f^*\mathcal{G})|_{U_\etale} \\
& =
\SheafHom_{\mathcal{O}_U}(f^*\mathcal{F}|_{U_\etale},
f^*\mathcal{G}|_{U_\etale})
\end{align*}
第一个等号由 $hom$ 的构造得出，第二个等号由 (7) 得出，
第三个等号由 $c$ 的定义得出，第四个等号由《site 上的模》引理
\ref{sites-modules-lemma-pullback-internal-hom} 得出。
最后一个等号是对《叠上的层》引理
\ref{stacks-sheaves-lemma-compare} 中赋环拓扑斯的平坦态射
$i_U (U_\etale, \mathcal{O}_U) \to
(\mathcal{U}_\etale, \mathcal{O}_\mathcal{U})$
应用同一参照得出的。
\item 在此处添加更多内容。
\end{enumerate}

\section{余极限与上同调}
\label{stacks-cohomology-section-colimits}

\noindent
下述引理特别可用于拟凝聚层的图表。

\begin{lemma}
\label{stacks-cohomology-lemma-colimits}
设 $\mathcal{X}$ 是拟紧且拟分离的代数叠。
则对 $\mathcal{X}$ 上阿贝层的每个滤过图表，
$$
\colim_i H^p(\mathcal{X}, \mathcal{F}_i)
\longrightarrow
H^p(\mathcal{X}, \colim_i \mathcal{F}_i)
$$
都是同构。对 $\mathcal{X}_\etale$ 上的阿贝层，
在 \'etale 拓扑中取上同调，同样结论也成立。
\end{lemma}

\begin{proof}
令 $\tau = fppf$，或相应地令 $\tau = \etale$。
对 site $\mathcal{X}_\tau$ 应用《site 上的上同调》引理
\ref{sites-cohomology-lemma-colim-global}，便得到结论。
为检验其假设，我们使用《site 上的上同调》注
\ref{sites-cohomology-remark-colim-global}。
具体地，以 $\mathcal{B} \subset \Ob(\mathcal{X}_\tau)$ 表示
位于仿射概形上的对象集合。换言之，$\mathcal{B}$ 的一个元素是态射
$x : U \to \mathcal{X}$，其中 $U$ 仿射。
我们依次检验该注的条件 (1) -- (4)：
\begin{enumerate}
\item 由于 $\mathcal{X}$ 拟紧，存在满的光滑态射
$x : U \to \mathcal{X}$，且 $U$ 仿射
（《叠的性质》引理
\ref{stacks-properties-lemma-quasi-compact-stack}）。
因而 $h_x^\# \to *$ 是 $\mathcal{X}_\tau$ 上层的满态射。
\item 由于 $\mathcal{X}_\tau$ 中的覆盖分别是 fppf 覆盖或
\'etale 覆盖，$U \in \mathcal{B}$ 的每个覆盖都可由一个有限仿射
fppf 覆盖加细；分别参见《拓扑》引理
\ref{topologies-lemma-fppf-affine} 和引理
\ref{topologies-lemma-etale-affine}。
\item 设 $x : U \to \mathcal{X}$ 和 $x' : U' \to \mathcal{X}$
属于 $\mathcal{B}$。在 $\Sh(\mathcal{X}_\tau)$ 中的积
$h_x^\# \times h_{x'}^\#$ 等于 $\mathcal{X}_\tau$ 上由
$\mathcal{X}$ 上的代数空间
$W = U \times_{x, \mathcal{X}, x'} U'$ 确定的层：
对 $\mathcal{X}_\tau$ 的对象 $y : V \to \mathcal{X}$，有
$(h_x^\# \times h_{x'}^\#)(y) = \{f : V \to W \mid y =
x \circ \text{pr}_1 \circ f = x' \circ \text{pr}_2 \circ f\}$。
由于 $\mathcal{X}$ 拟分离，代数空间 $W$ 拟紧；
例如参见《叠的态射》引理
\ref{stacks-morphisms-lemma-quasi-compact-quasi-separated-permanence}。
因而可选取仿射概形 $U''$ 和满的 \'etale 态射
$U'' \to W$。以 $x'' : U'' \to \mathcal{X}$ 表示
$U'' \to W$ 与 $W \to \mathcal{X}$ 的复合。则
$h_{x''}^\# \to h_x^\# \times h_{x'}^\#$ 如需为满态射。
\item 设 $x : U \to \mathcal{X}$ 和 $x' : U' \to \mathcal{X}$
属于 $\mathcal{B}$。设 $a, b : U \to U'$ 是 $\mathcal{X}$ 上的态射，
即 $a, b  : x \to x'$ 是 $\mathcal{X}_\tau$ 中的态射。
则 $h_a$ 和 $h_b$ 的等化子由 $a, b : U \to U'$ 的等化子表示；
后者是 $\mathcal{X}$ 上的仿射概形，因而属于 $\mathcal{B}$。
\end{enumerate}
% P12-ERR-0027: see ../control/ERRATA_CANDIDATES.jsonl
证明完成。
\end{proof}

\begin{lemma}
\label{stacks-cohomology-lemma-colimit-cohomology}
设 $f : \mathcal{X} \to \mathcal{Y}$ 是代数叠之间拟紧且拟分离的态射。
设 $\mathcal{F} = \colim \mathcal{F}_i$ 是 $\mathcal{X}$ 上阿贝层的滤过余极限。
则对任意 $p \geq 0$，有
$$
R^pf_*\mathcal{F} = \colim R^pf_*\mathcal{F}_i.
$$
对 $\mathcal{X}_\etale$ 上的阿贝层，
在 \'etale 拓扑中取高阶直像，同样结论也成立。
\end{lemma}

\begin{proof}
我们对 fppf 拓扑证明此结论；\'etale 拓扑的证明相同。
回忆 $R^if_*\mathcal{F}$ 是 $\mathcal{Y}_{fppf}$ 上与预层
$$
(y : V \to \mathcal{Y})
\longmapsto
H^i(V \times_{y, \mathcal{Y}} \mathcal{X}, \text{pr}^{-1}\mathcal{F})
$$
相联系的层。参见《叠上的层》引理
\ref{stacks-sheaves-lemma-pushforward-restriction}。
回忆，余极限是与预层余极限相联系的层。
当 $V$ 仿射时，纤维积 $V \times_\mathcal{Y} \mathcal{X}$
拟紧且拟分离。因而对仿射的 $V$，可将引理
\ref{stacks-cohomology-lemma-colimits} 应用于 $H^p(V \times_\mathcal{Y} \mathcal{X}, -)$。
由于每个 $V$ 都有由仿射对象组成的 fppf 覆盖，引理得证。
略去一些细节。
\end{proof}

\begin{lemma}
\label{stacks-cohomology-lemma-quasi-coherent-pushforward-direct-sums}
设 $f : \mathcal{X} \to \mathcal{Y}$ 是代数叠之间拟紧且拟分离的态射。
函子 $f_{\QCoh, *}$ 与函子 $R^if_{\QCoh, *}$
均与直和和滤过余极限可交换。
\end{lemma}

\begin{proof}
由引理 \ref{stacks-cohomology-lemma-colimit-cohomology}，在全体模上，函子
$f_*$ 和 $R^if_*$ 与直和和滤过余极限可交换。
又因为 $f_{\QCoh, *} = Q \circ f_*$、
$R^if_{\QCoh, *} = Q \circ R^if_*$，且 $Q$ 与所有余极限可交换，
由引理 \ref{stacks-cohomology-lemma-adjoint-kernel-parasitic} 得到结论。
\end{proof}

\begin{lemma}
\label{stacks-cohomology-lemma-quasi-coherent-pushforward-affine}
设 $f : \mathcal{X} \to \mathcal{Y}$ 是代数叠之间的仿射态射。
对 $i > 0$，函子 $R^if_{\QCoh, *}$ 为零；函子
$f_{\QCoh, *}$ 正合，并与直和和所有余极限可交换。
\end{lemma}

\begin{proof}
由于 $R^if_{\QCoh, *} = Q \circ R^if_*$，引理
\ref{stacks-cohomology-lemma-loc-qcoh-fbc-affine-direct-image} 给出所述消失。
该消失与函子族 $\{R^if_{\QCoh, *}\}_{i \geq 0}$
构成 $\delta$-函子这一事实一起，表明 $f_{\QCoh, *}$ 正合；
参见命题 \ref{stacks-cohomology-proposition-direct-image-quasi-coherent}。
例如由引理 \ref{stacks-cohomology-lemma-quasi-coherent-pushforward-direct-sums}，
$f_{\QCoh, *}$ 与直和可交换。
与直和可交换的正合函子与所有余极限可交换。
\end{proof}

\noindent
下述引理告诉我们，在拟紧且拟分离的代数叠中，
有限表示模如预期般表现。

\begin{lemma}
\label{stacks-cohomology-lemma-finite-presentation-quasi-compact-colimit}
设 $\mathcal{X}$ 是拟紧且拟分离的代数叠。
设 $I$ 是有向集，$(\mathcal{F}_i, \varphi_{ii'})$ 是
$I$ 上的 $\mathcal{O}_\mathcal{X}$-模系。设 $\mathcal{G}$ 是有限表示的
$\mathcal{O}_\mathcal{X}$-模。则
$$
\colim_i \Hom_\mathcal{X}(\mathcal{G}, \mathcal{F}_i)
=
\Hom_\mathcal{X}(\mathcal{G}, \colim_i \mathcal{F}_i).
$$
特别地，$\Hom_\mathcal{X}(\mathcal{G}, -)$ 与
$\QCoh(\mathcal{O}_\mathcal{X})$ 中的滤过余极限可交换。
\end{lemma}

\begin{proof}
所显示的等式是《site 上的模》引理
\ref{sites-modules-lemma-finite-presentation-quasi-compact-colimit}
的一个特例。为应用它，我们需对 site $\mathcal{X}_{fppf}$
检验《Site》引理
\ref{sites-lemma-directed-colimits-global-sections} 第 (4) 部分的假设。
为此，我们将检验《Site》注
\ref{sites-remark-stronger-conditions} 的假设 (2)(a)、(2)(b)、(2)(c)。
具体地，以 $\mathcal{B} \subset \Ob(\mathcal{X}_{fppf})$ 表示
位于仿射概形上的对象集合。换言之，$\mathcal{B}$ 的一个元素是态射
$x : U \to \mathcal{X}$，其中 $U$ 仿射。
我们依次检验该注的条件 (2)(a)、(2)(b)、(2)(c)：
\begin{enumerate}
\item 由于 $\mathcal{X}$ 拟紧，存在满的光滑态射
$x : U \to \mathcal{X}$，且 $U$ 仿射
（《叠的性质》引理
\ref{stacks-properties-lemma-quasi-compact-stack}）。
因而 $h_x^\# \to *$ 是 $\mathcal{X}_{fppf}$ 上层的满态射。
\item 由于 $\mathcal{X}_{fppf}$ 中的覆盖是 fppf 覆盖，
$U \in \mathcal{B}$ 的每个覆盖都可由一个有限仿射 fppf 覆盖加细；
参见《拓扑》引理 \ref{topologies-lemma-fppf-affine}。
\item 设 $x : U \to \mathcal{X}$ 和 $x' : U' \to \mathcal{X}$
属于 $\mathcal{B}$。在 $\Sh(\mathcal{X}_{fppf})$ 中的积
$h_x^\# \times h_{x'}^\#$ 等于 $\mathcal{X}_{fppf}$ 上由
$\mathcal{X}$ 上的代数空间
$W = U \times_{x, \mathcal{X}, x'} U'$ 确定的层：
对 $\mathcal{X}_{fppf}$ 的对象 $y : V \to \mathcal{X}$，有
$(h_x^\# \times h_{x'}^\#)(y) = \{f : V \to W \mid y =
x \circ \text{pr}_1 \circ f = x' \circ \text{pr}_2 \circ f\}$。
由于 $\mathcal{X}$ 拟分离，代数空间 $W$ 拟紧；
例如参见《叠的态射》引理
\ref{stacks-morphisms-lemma-quasi-compact-quasi-separated-permanence}。
因而可选取仿射概形 $U''$ 和满的 \'etale 态射
$U'' \to W$。以 $x'' : U'' \to \mathcal{X}$ 表示
$U'' \to W$ 与 $W \to \mathcal{X}$ 的复合。则
$h_{x''}^\# \to h_x^\# \times h_{x'}^\#$ 如需为满态射。
\end{enumerate}
对最后一个陈述，注意包含函子
% P12-ERR-0028: see ../control/ERRATA_CANDIDATES.jsonl
$\QCoh(\mathcal{O}_X) \to \textit{Mod}(\mathcal{O}_X)$
与余极限可交换，且有限表示模拟凝聚。
参见《叠上的层》引理
\ref{stacks-sheaves-lemma-quasi-coherent-algebraic-stack}。
\end{proof}

\section{光滑-\'etale site 与平坦-fppf site}
\label{stacks-cohomology-section-lisse-etale}

\noindent
在著作 \cite{LM-B} 中，上述许多结果都是利用代数叠的
光滑-\'etale site 证明的。我们在此定义这个 site。
在《例》第 \ref{examples-section-lisse-etale-not-functorial} 节中，
我们说明光滑-\'etale site 不具函子性。
我们还定义它的类比物，即平坦-fppf site；
它更适合 Stacks Project 中对代数叠的发展
（因为我们以 fppf 拓扑为基础拓扑）。
当然，平坦-fppf site 也不具函子性。

\begin{definition}
\label{stacks-cohomology-definition-lisse-etale}
设 $\mathcal{X}$ 是代数叠。
\begin{enumerate}
\item $\mathcal{X}$ 的 {\it 光滑-\'etale site} 记为
$\mathcal{X}_{lisse,\etale}$\footnote{文献中，该 site 记为
$\text{Lis-\'et}(\mathcal{X})$ 或
$\text{Lis-Et}(\mathcal{X})$，相联系的拓扑斯记为
$\mathcal{X}_{\text{lis-\'e}t}$ 或 $\mathcal{X}_{\text{lis-et}}$。
在 Stacks Project 中，我们约定命名 site，并以
$\Sh(\mathcal{C})$ 表示相应的拓扑斯。}；它是 $\mathcal{X}$ 的充满子范畴，
其对象是那些位于概形 $U$ 上且态射
$x : U \to \mathcal{X}$ 光滑的 $x \in \Ob(\mathcal{X})$。
$\mathcal{X}_{lisse,\etale}$ 的一个覆盖是态射族
$\{x_i \to x\}_{i \in I}$，这是 $\mathcal{X}_{lisse,\etale}$ 中的态射族，它作为
$\mathcal{X}_\etale$ 中的态射族构成覆盖。
\item $\mathcal{X}$ 的 {\it 平坦-fppf site} 记为
$\mathcal{X}_{flat,fppf}$；它是 $\mathcal{X}$ 的充满子范畴，其对象是那些位于概形 $U$ 上且态射
$x : U \to \mathcal{X}$ 平坦的 $x \in \Ob(\mathcal{X})$。
$\mathcal{X}_{flat,fppf}$ 的一个覆盖是态射族
$\{x_i \to x\}_{i \in I}$，这是 $\mathcal{X}_{flat,fppf}$ 中的态射族，它作为
$\mathcal{X}_{fppf}$ 中的态射族构成覆盖。
\end{enumerate}
\end{definition}

\noindent
以 $\mathcal{O}_{\mathcal{X}_{lisse,\etale}}$ 表示
$\mathcal{O}_\mathcal{X}$ 到光滑-\'etale site 的限制；
$\mathcal{O}_{\mathcal{X}_{flat,fppf}}$ 类似。
光滑-\'etale site 与 \'etale site 之间的关系如下
（在下述引理中，我们主要关注“拓扑”性质）。

\begin{lemma}
\label{stacks-cohomology-lemma-lisse-etale}
设 $\mathcal{X}$ 是代数叠。
\begin{enumerate}
\item 包含函子
$\mathcal{X}_{lisse,\etale} \to \mathcal{X}_\etale$
充满忠实、连续且余连续。因此：
\begin{enumerate}
\item 有拓扑斯态射
$$
g :
\Sh(\mathcal{X}_{lisse,\etale})
\longrightarrow
\Sh(\mathcal{X}_\etale)
$$
且 $g^{-1}$ 由限制给出；
\item 在集合层上，函子 $g^{-1}$ 有左伴随 $g_!^{Sh}$；
\item 伴随态射 $g^{-1}g_* \to \text{id}$ 和
$\text{id} \to g^{-1}g_!^{Sh}$ 都是同构；
\item 在阿贝层上，函子 $g^{-1}$ 有左伴随 $g_!$；
\item 伴随态射 $\text{id} \to g^{-1}g_!$ 是同构；
\item 有 $g^{-1}\mathcal{O}_\mathcal{X} =
\mathcal{O}_{\mathcal{X}_{lisse,\etale}}$，因而 $g$ 诱导一个平坦赋环拓扑斯态射，
并且 $g^{-1} = g^*$。
\end{enumerate}
\item 包含函子
$\mathcal{X}_{flat,fppf} \to \mathcal{X}_{fppf}$
充满忠实、连续且余连续。因此：
\begin{enumerate}
\item 有拓扑斯态射
$$
g :
\Sh(\mathcal{X}_{flat,fppf})
\longrightarrow
\Sh(\mathcal{X}_{fppf})
$$
且 $g^{-1}$ 由限制给出；
\item 在集合层上，函子 $g^{-1}$ 有左伴随 $g_!^{Sh}$；
\item 伴随态射 $g^{-1}g_* \to \text{id}$ 和
$\text{id} \to g^{-1}g_!^{Sh}$ 都是同构；
\item 在阿贝层上，函子 $g^{-1}$ 有左伴随 $g_!$；
\item 伴随态射 $\text{id} \to g^{-1}g_!$ 是同构；
\item 有 $g^{-1}\mathcal{O}_\mathcal{X} =
\mathcal{O}_{\mathcal{X}_{flat,fppf}}$，因而 $g$ 诱导一个平坦赋环拓扑斯态射，
并且 $g^{-1} = g^*$。
\end{enumerate}
\end{enumerate}
\end{lemma}

\begin{proof}
在两种情形中，该函子显然充满忠实、连续且余连续
（参见《Site》定义 \ref{sites-definition-continuous} 和
\ref{sites-definition-cocontinuous}）。
因此，性质 (a)、(b)、(c) 由《Site》引理
\ref{sites-lemma-when-shriek} 和
\ref{sites-lemma-back-and-forth} 得出。
第 (d)、(e) 部分由《site 上的模》引理
\ref{sites-modules-lemma-g-shriek-adjoint} 和
\ref{sites-modules-lemma-back-and-forth} 得出。
第 (f) 部分显然。
\end{proof}

\begin{lemma}
\label{stacks-cohomology-lemma-lisse-etale-cohomology}
设 $\mathcal{X}$ 是代数叠。记号如引理
\ref{stacks-cohomology-lemma-lisse-etale}。
\begin{enumerate}
\item 对 $\mathcal{X}_\etale$ 上的阿贝层 $\mathcal{F}$，有
\begin{enumerate}
\item $H^p(\mathcal{X}_\etale, \mathcal{F}) =
H^p(\mathcal{X}_{lisse,\etale}, g^{-1}\mathcal{F})$；
\item 对 $\mathcal{X}_{lisse,\etale}$ 的任意对象 $x$，有
$H^p(x, \mathcal{F}) =
H^p(\mathcal{X}_{lisse,\etale}/x, g^{-1}\mathcal{F})$。
\end{enumerate}
对模层，同样结论也成立。
\item 对 $\mathcal{X}_{fppf}$ 上的阿贝层 $\mathcal{F}$，有
\begin{enumerate}
\item $H^p(\mathcal{X}_{fppf}, \mathcal{F}) =
H^p(\mathcal{X}_{flat,fppf}, g^{-1}\mathcal{F})$；
\item 对 $\mathcal{X}_{flat,fppf}$ 的任意对象 $x$，有
$H^p(x, \mathcal{F}) =
H^p(\mathcal{X}_{flat,fppf}/x, g^{-1}\mathcal{F})$。
\end{enumerate}
对模层，同样结论也成立。
\end{enumerate}
\end{lemma}

\begin{proof}
第 (1)(a) 部分由对包含函子
$\mathcal{X}_{lisse,\etale} \to \mathcal{X}_\etale$
应用《叠上的层》引理
\ref{stacks-sheaves-lemma-cohomology-on-subcategory} 得出。
第 (1)(b) 部分由第 (1)(a) 部分得出。事实上，若 $x$ 位于概形
$U$ 上，则 site $\mathcal{X}_\etale/x$ 等价于
$(\Sch/U)_\etale$，且
% P12-ERR-0029: see ../control/ERRATA_CANDIDATES.jsonl
$\mathcal{X}_{lisse,\etale}$ 等价于 $U_{lisse,\etale}$。
第 (2) 部分的证明相同。
\end{proof}

\begin{lemma}
\label{stacks-cohomology-lemma-lisse-etale-modules}
设 $\mathcal{X}$ 是代数叠。记号如引理
\ref{stacks-cohomology-lemma-lisse-etale}。
\begin{enumerate}
\item 存在函子
$$
g_! :
\textit{Mod}(\mathcal{X}_{lisse,\etale},
\mathcal{O}_{\mathcal{X}_{lisse,\etale}})
\longrightarrow
\textit{Mod}(\mathcal{X}_\etale, \mathcal{O}_{\mathcal{X}})
$$
它是 $g^*$ 的左伴随。此外，它与阿贝层上的函子 $g_!$ 一致，
且 $g^*g_! = \text{id}$。
\item 存在函子
$$
g_! :
\textit{Mod}(\mathcal{X}_{flat,fppf},
\mathcal{O}_{\mathcal{X}_{flat,fppf}})
\longrightarrow
\textit{Mod}(\mathcal{X}_{fppf}, \mathcal{O}_{\mathcal{X}})
$$
它是 $g^*$ 的左伴随。此外，它与阿贝层上的函子 $g_!$ 一致，
且 $g^*g_! = \text{id}$。
\end{enumerate}
\end{lemma}

\begin{proof}
在两种情形中，函子 $g_!$ 的存在性都由《site 上的模》引理
\ref{sites-modules-lemma-lower-shriek-modules} 得出。
要说明 $g_!$ 与阿贝层上的函子一致，我们将证明
《site 上的模》等式
(\ref{sites-modules-equation-compare-on-localizations})
中的态射是同构。

\medskip\noindent
光滑-\'etale 情形。设 $x \in \Ob(\mathcal{X}_{lisse,\etale})$
位于概形 $U$ 上，且 $x : U \to \mathcal{X}$ 光滑。
考虑诱导的充满忠实函子
$$
g' :
\mathcal{X}_{lisse,\etale}/x
\longrightarrow
\mathcal{X}_\etale/x
$$
右边与 $(\Sch/U)_\etale$ 识别，
左边与如下 $U'/U$ 概形组成的充满子范畴识别：
复合 $U' \to U \to \mathcal{X}$ 光滑。因此，《\'Etale 上同调》引理
\ref{etale-cohomology-lemma-compare-structure-sheaves}
可用。

\medskip\noindent
平坦-fppf 情形。设 $x \in \Ob(\mathcal{X}_{flat,fppf})$
位于概形 $U$ 上，且 $x : U \to \mathcal{X}$ 平坦。
考虑诱导的充满忠实函子
$$
g' :
\mathcal{X}_{flat,fppf}/x
\longrightarrow
\mathcal{X}_{fppf}/x
$$
右边与 $(\Sch/U)_{fppf}$ 识别，
左边与如下 $U'/U$ 概形组成的充满子范畴识别：
复合 $U' \to U \to \mathcal{X}$ 平坦。因此，《\'Etale 上同调》引理
\ref{etale-cohomology-lemma-compare-structure-sheaves}
可用。

\medskip\noindent
在两种情形中，等式 $g^*g_! = \text{id}$ 都由
$g^* = g^{-1}$ 以及引理 \ref{stacks-cohomology-lemma-lisse-etale}
中对阿贝层的等式得出。
\end{proof}

\begin{lemma}
\label{stacks-cohomology-lemma-lisse-etale-structure-sheaf}
设 $\mathcal{X}$ 是代数叠。记号如引理
\ref{stacks-cohomology-lemma-lisse-etale} 和 \ref{stacks-cohomology-lemma-lisse-etale-modules}。
\begin{enumerate}
\item 有 $g_!\mathcal{O}_{\mathcal{X}_{lisse,\etale}} =
\mathcal{O}_\mathcal{X}$。
\item 有 $g_!\mathcal{O}_{\mathcal{X}_{flat, fppf}} =
\mathcal{O}_\mathcal{X}$。
\end{enumerate}
\end{lemma}

\begin{proof}
在本证明中，我们记
$\mathcal{C} = \mathcal{X}_\etale$
（相应地，$\mathcal{C} = \mathcal{X}_{fppf}$），并记
$\mathcal{C}' = \mathcal{X}_{lisse,\etale}$
（相应地，$\mathcal{C}' = \mathcal{X}_{flat, fppf}$）。
则 $\mathcal{C}'$ 是 $\mathcal{C}$ 的充满子范畴。
我们把 $\mathcal{C}$ 的对象 $V$ 视为 $\mathcal{X}$ 上的概形，
把 $\mathcal{C}'$ 的对象 $U$ 视为在 $\mathcal{X}$ 上光滑
（相应地，平坦）的概形。
最后，记 $\mathcal{O} = \mathcal{O}_\mathcal{X}$，并记
$\mathcal{O}' = \mathcal{O}_{\mathcal{X}_{lisse,\etale}}$
（相应地，$\mathcal{O}' = \mathcal{O}_{\mathcal{X}_{flat,fppf}}$）。
按上述记号，有 $\mathcal{O}(V) = \Gamma(V, \mathcal{O}_V)$
和 $\mathcal{O}'(U) = \Gamma(U, \mathcal{O}_U)$。
考虑 $\mathcal{O}$-模同态
$g_!\mathcal{O}' \to \mathcal{O}$，它与识别
$\mathcal{O}' = g^{-1}\mathcal{O}$ 伴随。

\medskip\noindent
回忆 $g_!\mathcal{O}'$ 是与下列规则给出的预层
$g_{p!}\mathcal{O}'$ 相联系的层：
$$
V \longmapsto \colim_{V \to U} \mathcal{O}'(U)
$$
其中余极限在阿贝群范畴中取
（《site 上的模》定义 \ref{sites-modules-definition-g-shriek}）。
下文将频繁使用以下事实：若
$$
V \to U \to U'
$$
是态射，且 $f' \in \mathcal{O}'(U')$ 限制为
$f \in \mathcal{O}'(U)$，则 $(V \to U, f)$ 和
$(V \to U', f')$ 定义余极限中的同一元素。
此外，$g_!\mathcal{O}' \to \mathcal{O}$ 把元素
$(V \to U, f)$ 映到 $f$ 在 $V$ 上的拉回。

\medskip\noindent
我们证明 $g_!\mathcal{O}' \to \mathcal{O}$ 为满态射。
对 $\mathcal{C}$ 的某个对象 $V$，设 $h \in \mathcal{O}(V)$。
只需证明 $h$ 局部地属于像。选取 $\mathcal{C}'$ 的一个对象
$U$，它对应于满的光滑态射 $U \to \mathcal{X}$。
由于 $U \times_\mathcal{X} V \to V$ 满且光滑，用 $V$ 的一个
\'etale 覆盖的各成员取代 $V$ 后，可假设存在态射 $V \to U$；
参见《空间上的拓扑》引理
\ref{spaces-topologies-lemma-etale-dominates-smooth}。
利用 $h$，得到态射 $V \to U \times \mathbf{A}^1$，
使得写 $\mathbf{A}^1 = \Spec(\mathbf{Z}[t])$ 时，元素
$t \in \mathcal{O}(U \times \mathbf{A}^1)$ 拉回为 $h$。
由于 $U \times \mathbf{A}^1$ 是 $\mathcal{C}'$ 的对象，
$(V \to U \times \mathbf{A}^1, t)$ 是上述余极限的一个元素，
并如需映到 $h \in \mathcal{O}(V)$。

\medskip\noindent
假设 $s \in g_!\mathcal{O}'(V)$ 是一个在 $\mathcal{O}(V)$ 中映到零的截面。
要完成证明，需说明 $s$ 为零。用一个覆盖的各成员取代 $V$ 后，
可假设 $s$ 是余极限
$$
\colim_{V \to U} \mathcal{O}'(U)
$$
的元素。写 $s = \sum (\varphi_i, s_i)$ 为有限和，其中
$\varphi_i : V \to U_i$，$U_i$ 在 $\mathcal{X}$ 上光滑
（相应地，平坦），且
$s_i \in \Gamma(U_i, \mathcal{O}_{U_i})$。
选取一个概形 $W$，它到代数空间
$U = U_1 \times_\mathcal{X} \ldots \times_\mathcal{X} U_n$ 的态射满且 \'etale。
注意，$W$ 仍在 $\mathcal{X}$ 上光滑
（相应地，平坦），即它定义 $\mathcal{C}'$ 的一个对象。
纤维积
$$
V' = V \times_{(\varphi_1, \ldots, \varphi_n), U} W
$$
到 $V$ 满且 \'etale，因而只需证明 $s$ 在
$g_!\mathcal{O}'(V')$ 中映到零。注意，限制
$\sum (\varphi_i, s_i)|_{V'}$ 对应于函数 $s_i$ 到 $W$ 的拉回之和。
换言之，我们已归约到 $(\varphi, s)$ 的情形，其中
$\varphi : V \to U$ 是态射，$U$ 属于 $\mathcal{C}'$，且
$s \in \mathcal{O}'(U)$ 在 $\mathcal{O}(V)$ 中限制为零。
由交换图
$$
\xymatrix{
V \ar[rr]_-{(\varphi, 0)} \ar[rrd]_\varphi & & U \times \mathbf{A}^1 \\
& & U \ar[u]_{(\text{id}, 0)}
}
$$
可知
% P12-ERR-0030: see ../control/ERRATA_CANDIDATES.jsonl
$((\varphi, 0) : V \to U \times \mathbf{A}^1, \text{pr}_2^*x)$
在上述余极限中表示零。因而可以用
$U \times \mathbf{A}^1$ 取代 $U$，用 $(\varphi, 0)$ 取代 $\varphi$，
用 $\text{pr}_1^*s + \text{pr}_2^*x$ 取代 $s$。
因此可假设 $U$ 中 $s$ 的零点轨迹 $Z : s = 0$
在 $\mathcal{X}$ 上光滑（相应地，平坦）。
于是 $(V \to Z, 0)$ 与 $(\varphi, s)$ 在余极限中取同一值，
即元素 $s$ 如需为零。
\end{proof}

\noindent
可用光滑-\'etale site 和平坦-fppf site 如下刻画寄生模。

\begin{lemma}
\label{stacks-cohomology-lemma-parasitic-in-terms-flat-fppf}
设 $\mathcal{X}$ 是代数叠。
\begin{enumerate}
\item 设 $\mathcal{F}$ 是 $\mathcal{X}_\etale$ 上具有平坦基变换性质的
$\mathcal{O}_\mathcal{X}$-模。以下条件等价：
\begin{enumerate}
\item $\mathcal{F}$ 是寄生的；
\item $g^*\mathcal{F} = 0$，其中
$g : \Sh(\mathcal{X}_{lisse,\etale}) \to
\Sh(\mathcal{X}_\etale)$ 如引理 \ref{stacks-cohomology-lemma-lisse-etale} 所述。
\end{enumerate}
\item 设 $\mathcal{F}$ 是 $\mathcal{X}_{fppf}$ 上的
$\mathcal{O}_\mathcal{X}$-模。以下条件等价：
\begin{enumerate}
\item $\mathcal{F}$ 是寄生的；
\item $g^*\mathcal{F} = 0$，其中
$g :  \Sh(\mathcal{X}_{flat,fppf}) \to \Sh(\mathcal{X}_{fppf})$
如引理 \ref{stacks-cohomology-lemma-lisse-etale} 所述。
\end{enumerate}
\end{enumerate}
\end{lemma}

\begin{proof}
第 (2) 部分直接由定义得出
（这是平坦-fppf site 相比光滑-\'etale site 的一个优点）。
含意 (1)(a) $\Rightarrow$ (1)(b) 也很明显。
为证明 (1)(b) $\Rightarrow$ (1)(a)，设 $U$ 是概形，
$x : U \to \mathcal{X}$ 是满的光滑态射。
则 $x$ 是 $\mathcal{X}$ 的光滑-\'etale site 的对象。
因此 (1)(b) 意味着 $\mathcal{F}|_{U_\etale} = 0$。
% P12-ERR-0031: see ../control/ERRATA_CANDIDATES.jsonl
设 $V \to \mathcal{X}$ 是平坦态射，其中 $V$ 是概形。
置 $W = U \times_\mathcal{X} V$，并考虑图
$$
\xymatrix{
W \ar[d]_p \ar[r]_q & V \ar[d] \\
U \ar[r] & \mathcal{X}
}
$$
注意，投影 $p : W \to U$ 平坦，投影
$q : W \to V$ 光滑且满。这意味着 $q_{small}^*$
在拟凝聚模上是忠实函子。按假设，$\mathcal{F}$
具有平坦基变换性质，因而得到
$p_{small}^*\mathcal{F}|_{U_\etale} \cong
q_{small}^*\mathcal{F}|_{V_\etale}$。故若 $\mathcal{F}$
属于 $g^*$ 的核，则如需有 $\mathcal{F}|_{V_\etale} = 0$。
\end{proof}

\section{光滑-\'etale site 与平坦-fppf site 的函子性}
\label{stacks-cohomology-section-lisse-etale-functorial}

\noindent
光滑-\'etale site 对代数叠的光滑态射具有函子性，
平坦-fppf site 对代数叠的平坦态射具有函子性。
我们提醒读者，光滑-\'etale 拓扑斯和平坦-fppf 拓扑斯
对代数叠的全体态射并不具有函子性；参见《例》第
\ref{examples-section-lisse-etale-not-functorial} 节。

\begin{lemma}
\label{stacks-cohomology-lemma-lisse-etale-functorial}
设 $f : \mathcal{X} \to \mathcal{Y}$ 是代数叠之间的态射。
\begin{enumerate}
\item 若 $f$ 光滑，则 $f$ 限制为连续且余连续的函子
$\mathcal{X}_{lisse,\etale} \to \mathcal{Y}_{lisse,\etale}$，
它给出一个赋环拓扑斯态射，使下图交换：
$$
\xymatrix{
\Sh(\mathcal{X}_{lisse,\etale}) \ar[r]_{g'} \ar[d]_{f'} &
\Sh(\mathcal{X}_\etale) \ar[d]^f \\
\Sh(\mathcal{Y}_{lisse,\etale}) \ar[r]^g &
\Sh(\mathcal{Y}_\etale)
}
$$
有 $f'_*(g')^{-1} = g^{-1}f_*$ 和 $g'_!(f')^{-1} = f^{-1}g_!$。
\item 若 $f$ 平坦，则 $f$ 限制为连续且余连续的函子
$\mathcal{X}_{flat,fppf} \to \mathcal{Y}_{flat,fppf}$，
它给出一个赋环拓扑斯态射，使下图交换：
$$
\xymatrix{
\Sh(\mathcal{X}_{flat,fppf}) \ar[r]_{g'} \ar[d]_{f'} &
\Sh(\mathcal{X}_{fppf}) \ar[d]^f \\
\Sh(\mathcal{Y}_{flat,fppf}) \ar[r]^g &
\Sh(\mathcal{Y}_{fppf})
}
$$
有 $f'_*(g')^{-1} = g^{-1}f_*$ 和 $g'_!(f')^{-1} = f^{-1}g_!$。
\end{enumerate}
\end{lemma}

\begin{proof}
最初的陈述由以下事实得出：若 $x \in \Ob(\mathcal{X})$
位于概形 $U$ 上，且 $x : U \to \mathcal{X}$ 光滑
（相应地，平坦），而 $f$ 光滑
（相应地，平坦），则
$f(x) : U \to \mathcal{Y}$ 光滑（相应地，平坦）；
参见《叠的态射》引理
\ref{stacks-morphisms-lemma-composition-smooth} 和
\ref{stacks-morphisms-lemma-composition-flat}。
由我们对这些范畴中覆盖的定义，诱导函子
$\mathcal{X}_{lisse,\etale} \to \mathcal{Y}_{lisse,\etale}$
（相应地，$\mathcal{X}_{flat,fppf} \to \mathcal{Y}_{flat,fppf}$）
连续且余连续。最后，该图的交换性由水平态射由包含函子给出
（参见引理 \ref{stacks-cohomology-lemma-lisse-etale}）以及《Site》引理
\ref{sites-lemma-composition-cocontinuous} 得出。

\medskip\noindent
为证明 $f'_*(g')^{-1} = g^{-1}f_*$，设 $\mathcal{F}$ 是
$\mathcal{X}_\etale$ 上（相应地，$\mathcal{X}_{fppf}$ 上）的层。
有典范拉回态射
$$
g^{-1}f_*\mathcal{F} \longrightarrow f'_*(g')^{-1}\mathcal{F}
$$
参见《Site》第 \ref{sites-section-pullback} 节。
我们断言该态射是同构。为证明此点，选取
$\mathcal{Y}_{lisse,\etale}$ 的一个对象 $y$
（相应地，$\mathcal{Y}_{flat,fppf}$ 的一个对象）。
设 $y$ 位于概形 $V$ 上，且 $y : V \to \mathcal{Y}$
光滑（相应地，平坦）。由于 $g^{-1}$ 是限制，有
$$
\left(g^{-1}f_*\mathcal{F}\right)(y) =
\Gamma(V \times_{y, \mathcal{Y}} \mathcal{X},\ \text{pr}^{-1}\mathcal{F})
$$
这由《叠上的层》等式
(\ref{stacks-sheaves-equation-pushforward}) 得出。
以
$(V \times_{y, \mathcal{Y}} \mathcal{X})' \subset
V \times_{y, \mathcal{Y}} \mathcal{X}$
表示由下列对象组成的充满子范畴：
$z : W \to V \times_{y, \mathcal{Y}} \mathcal{X}$，且诱导态射
$W \to \mathcal{X}$ 光滑（相应地，平坦）。
以
$$
\text{pr}' :
(V \times_{y, \mathcal{Y}} \mathcal{X})'
\longrightarrow
\mathcal{X}_{lisse,\etale}
\ (\text{分别 }\mathcal{X}_{flat,fppf})
$$
表示上述公式中函子 $\text{pr}$ 的限制。
与《叠上的层》等式
(\ref{stacks-sheaves-equation-pushforward}) 的证明完全相同的论证表明，
对 $\mathcal{X}_{lisse,\etale}$ 上
（相应地，$\mathcal{X}_{flat,fppf}$ 上）的任意层 $\mathcal{H}$，有
\begin{equation}
\label{stacks-cohomology-equation-pushforward-lisse-etale}
f'_*\mathcal{H}(y) =
\Gamma((V \times_{y, \mathcal{Y}} \mathcal{X})',
\ (\text{pr}')^{-1}\mathcal{H})
\end{equation}
由于 $(g')^{-1}$ 是限制，可知
$$
\left(f'_*(g')^{-1}\mathcal{F}\right)(y) =
\Gamma((V \times_{y, \mathcal{Y}} \mathcal{X})',
\ \text{pr}^{-1}\mathcal{F}|_{(V \times_{y, \mathcal{Y}} \mathcal{X})'})
$$
由《叠上的层》引理
\ref{stacks-sheaves-lemma-cohomology-on-subcategory}，有
$$
\Gamma((V \times_{y, \mathcal{Y}} \mathcal{X})',
\ \text{pr}^{-1}\mathcal{F}|_{(V \times_{y, \mathcal{Y}} \mathcal{X})'})
=
\Gamma(V \times_{y, \mathcal{Y}} \mathcal{X},\ \text{pr}^{-1}\mathcal{F})
$$
如需相等。虽然我们略去了对该引理各假设的检验，
但指出，验证第二个条件时用到了 $V \to \mathcal{Y}$ 光滑
（相应地，平坦）这一事实。

\medskip\noindent
最后，等式 $g'_!(f')^{-1} = f^{-1}g_!$ 形式地由等式
$f'_*(g')^{-1} = g^{-1}f_*$、$f^{-1}$ 与 $f_*$ 的伴随性、
$g_!$ 与 $g^{-1}$ 的伴随性，以及它们的“带撇”版本得出。
\end{proof}

\begin{lemma}
\label{stacks-cohomology-lemma-lisse-etale-functorial-pushforward}
假设与记号如引理 \ref{stacks-cohomology-lemma-lisse-etale-functorial}。
设 $\mathcal{H}$ 是 $\mathcal{X}_{lisse,\etale}$ 上
（相应地，$\mathcal{X}_{flat,fppf}$ 上）的阿贝层。则
\begin{equation}
\label{stacks-cohomology-equation-higher-direct-image-lisse-etale}
R^pf'_*\mathcal{H} =
\text{下列对象的伴随层： }y \longmapsto
H^p((V \times_{y, \mathcal{Y}} \mathcal{X})', (\text{pr}')^{-1}\mathcal{H})
\end{equation}
这里 $y$ 是 $\mathcal{Y}_{lisse,\etale}$ 的对象
（相应地，$\mathcal{Y}_{flat,fppf}$ 的对象），它位于概形 $V$ 上；
记号 $(V \times_{y, \mathcal{Y}} \mathcal{X})'$ 和 $\text{pr}'$
在证明中解释。
\end{lemma}

\begin{proof}
与引理 \ref{stacks-cohomology-lemma-lisse-etale-functorial} 的证明一样，以
$(V \times_{y, \mathcal{Y}} \mathcal{X})' \subset
V \times_{y, \mathcal{Y}} \mathcal{X}$ 表示由对象
$(x, \varphi)$ 组成的充满子范畴，其中 $x$ 是
$\mathcal{X}_{lisse,\etale}$ 的对象
（相应地，$\mathcal{X}_{flat,fppf}$ 的对象），且
$\varphi : f(x) \to y$ 是 $\mathcal{Y}$ 中的态射。
由等式 (\ref{stacks-cohomology-equation-pushforward-lisse-etale})，有
$$
f'_*\mathcal{H}(y) =
\Gamma((V \times_{y, \mathcal{Y}} \mathcal{X})',
\ (\text{pr}')^{-1}\mathcal{H})
$$
其中 $\text{pr}'$ 是投影。对
$(V \times_{y, \mathcal{Y}} \mathcal{X})'$ 的对象 $(x, \varphi)$，
可将 $\varphi$ 视为 $x$ 上 $(f')^{-1}h_y$ 的一个截面。
因此 $(V \times_\mathcal{Y} \mathcal{X})'$ 是 site
$\mathcal{X}_{lisse,\etale}$
（相应地，$\mathcal{X}_{flat,fppf}$）
对集合层 $(f')^{-1}h_y$ 的局部化；参见《Site》引理
\ref{sites-lemma-localize-topos-site}。态射
$$
\text{pr}' : (V \times_{y, \mathcal{Y}} \mathcal{X})'
\to \mathcal{X}_{lisse,\etale}
\ (\text{分别 }
\text{pr}' : (V \times_{y, \mathcal{Y}} \mathcal{X})'
\to \mathcal{X}_{flat,fppf})
$$
是局部化态射。特别地，拉回 $(\text{pr}')^{-1}$ 保持内射阿贝层；
参见《site 上的上同调》引理
\ref{sites-cohomology-lemma-cohomology-on-sheaf-sets}。

\medskip\noindent
在 $\mathcal{X}_{lisse,\etale}$ 上
（相应地，$\mathcal{X}_{flat,fppf}$ 上）选取内射分解
$\mathcal{H} \to \mathcal{I}^\bullet$。
由直像公式，$R^if'_*\mathcal{H}$ 是与下列预层相联系的层：
该预层把 $y$ 映到复形
$$
\begin{matrix}
\Gamma\Big((V \times_{y, \mathcal{Y}} \mathcal{X})',
(\text{pr}')^{-1}\mathcal{I}^{i - 1}\Big) \\
\downarrow \\
\Gamma\Big((V \times_{y, \mathcal{Y}} \mathcal{X})',
(\text{pr}')^{-1}\mathcal{I}^i\Big) \\
\downarrow \\
\Gamma\Big((V \times_{y, \mathcal{Y}} \mathcal{X})',
(\text{pr}')^{-1}\mathcal{I}^{i + 1}\Big)
\end{matrix}
$$
的上同调。由于 $(\text{pr}')^{-1}$ 正合且保持内射对象，
复形 $(\text{pr}')^{-1}\mathcal{I}^\bullet$ 是
$(\text{pr}')^{-1}\mathcal{H}$ 的内射分解，证明完成。
\end{proof}

\begin{lemma}
\label{stacks-cohomology-lemma-lisse-etale-functorial-cohomology}
假设与记号如引理 \ref{stacks-cohomology-lemma-lisse-etale-functorial}。
对 $\mathcal{X}_\etale$ 上（相应地，$\mathcal{X}_{fppf}$ 上）的
任意阿贝层 $\mathcal{F}$，典范（基变换）态射
$$
g^{-1}Rf_*\mathcal{F} \longrightarrow Rf'_*(g')^{-1}\mathcal{F}
$$
是同构。
\end{lemma}

\begin{proof}
比较《叠上的层》引理
\ref{stacks-sheaves-lemma-pushforward-restriction} 和引理
\ref{stacks-cohomology-lemma-lisse-etale-functorial-pushforward}
中给出的 $g^{-1}R^pf_*\mathcal{F}$ 与
$R^pf'_*(g')^{-1}\mathcal{F}$ 的公式，可知只需证明
$$
H^p((V \times_{y, \mathcal{Y}} \mathcal{X})',
\ \text{pr}^{-1}\mathcal{F}|_{(V \times_{y, \mathcal{Y}} \mathcal{X})'})
=
H^p_\tau(V \times_{y, \mathcal{Y}} \mathcal{X},\ \text{pr}^{-1}\mathcal{F})
$$
其中 $\tau = \etale$（相应地，$\tau = fppf$）。
这里 $y$ 是 $\mathcal{Y}$ 的一个对象，位于概形 $V$ 上，
且态射 $y : V \to \mathcal{Y}$ 光滑（相应地，平坦）。
该等式由《叠上的层》引理
\ref{stacks-sheaves-lemma-cohomology-on-subcategory} 得出。
虽然我们略去了对该引理各假设的检验，
但指出，验证第二个条件时用到了
$V \to \mathcal{Y}$ 光滑（相应地，平坦）这一事实。
\end{proof}

\section{拟凝聚模与光滑-\'etale site、平坦-fppf site}
\label{stacks-cohomology-section-quasi-coherent-modules-II}

\noindent
本节说明如何利用代数叠的光滑-\'etale site 或平坦-fppf site
来理解其上的拟凝聚模。

\begin{lemma}
\label{stacks-cohomology-lemma-check-qc-on-etale-covering}
设 $\mathcal{X}$ 是代数叠。
\begin{enumerate}
\item 设 $f_j : \mathcal{X}_j \to \mathcal{X}$ 是代数叠的一族光滑态射，且
$|\mathcal{X}| =\bigcup |f_j|(|\mathcal{X}_j|)$。
设 $\mathcal{F}$ 是 $\mathcal{X}_\etale$ 上的
$\mathcal{O}_\mathcal{X}$-模层。若每个 $f_j^{-1}\mathcal{F}$ 都拟凝聚，
则 $\mathcal{F}$ 也拟凝聚。
\item 设 $f_j : \mathcal{X}_j \to \mathcal{X}$ 是代数叠的一族平坦且
局部有限表示的态射，且
$|\mathcal{X}| =\bigcup |f_j|(|\mathcal{X}_j|)$。
设 $\mathcal{F}$ 是 $\mathcal{X}_{fppf}$ 上的
$\mathcal{O}_\mathcal{X}$-模层。若每个 $f_j^{-1}\mathcal{F}$ 都拟凝聚，
则 $\mathcal{F}$ 也拟凝聚。
\end{enumerate}
\end{lemma}

\begin{proof}
证明 (1)。可以用概形 $U_j$ 取代每个代数叠 $\mathcal{X}_j$
（这里用到任意代数叠都有概形给出的光滑覆盖，
以及光滑态射的复合光滑；参见《叠的态射》引理
\ref{stacks-morphisms-lemma-composition-smooth}）。
$\mathcal{F}$ 到 $(\Sch/U_j)_\etale$ 的拉回仍拟凝聚；参见
《site 上的模》引理 \ref{sites-modules-lemma-local-pullback}。
因而 $f = \coprod f_j : U = \coprod U_j \to \mathcal{X}$
是光滑满态射。设 $x : V \to \mathcal{X}$ 是 $\mathcal{X}$ 的一个对象。
由《叠上的层》引理
\ref{stacks-sheaves-lemma-surjective-flat-locally-finite-presentation}，
存在 \'etale 覆盖 $\{x_i \to x\}_{i \in I}$，使得每个
$x_i$ 都提升为 $(\Sch/U)_\etale$ 的对象 $u_i$。
这只是说：$x_i$ 位于概形 $V_i$ 上，
$\{V_i \to V\}$ 是 \'etale 覆盖，且 $x_i$ 来自态射
$u_i : V_i \to U$。于是
$x_i^*\mathcal{F} = u_i^*f^*\mathcal{F}$ 拟凝聚。
这意味着 $(\Sch/V)_\etale$ 上的 $x^*\mathcal{F}$ 拟凝聚；
例如由《site 上的模》引理
\ref{sites-modules-lemma-local-final-object} 可得。
由《叠上的层》引理
\ref{stacks-sheaves-lemma-characterize-quasi-coherent-bis}，
$x^*\mathcal{F}$ 是 fppf 层。由于 $x$ 任意，可知
$\mathcal{F}$ 在 fppf 拓扑中是层。应用《叠上的层》引理
\ref{stacks-sheaves-lemma-characterize-quasi-coherent}，
可知 $\mathcal{F}$ 拟凝聚。

\medskip\noindent
证明 (2)。使用完全相同的论证；我们在此完整写出。
可以用概形 $U_j$ 取代每个代数叠 $\mathcal{X}_j$
（这里用到任意代数叠都有概形给出的光滑覆盖，
% P12-ERR-0032: see ../control/ERRATA_CANDIDATES.jsonl
以及平坦且局部有限表示的态射在复合下保持；参见《叠的态射》引理
\ref{stacks-morphisms-lemma-composition-flat} 和
\ref{stacks-morphisms-lemma-composition-finite-presentation}）。
$\mathcal{F}$ 到 $(\Sch/U_j)_\etale$ 的拉回仍局部拟凝聚；参见
《叠上的层》引理 \ref{stacks-sheaves-lemma-pullback-quasi-coherent}。
因而 $f = \coprod f_j : U = \coprod U_j \to \mathcal{X}$
是满的、平坦且局部有限表示的态射。
设 $x : V \to \mathcal{X}$ 是 $\mathcal{X}$ 的一个对象。
由《叠上的层》引理
\ref{stacks-sheaves-lemma-surjective-flat-locally-finite-presentation}，
存在 fppf 覆盖 $\{x_i \to x\}_{i \in I}$，使得每个
$x_i$ 都提升为 $(\Sch/U)_\etale$ 的对象 $u_i$。
这只是说：$x_i$ 位于概形 $V_i$ 上，
$\{V_i \to V\}$ 是 fppf 覆盖，且 $x_i$ 来自态射
$u_i : V_i \to U$。于是
$x_i^*\mathcal{F} = u_i^*f^*\mathcal{F}$ 拟凝聚。
这意味着 $(\Sch/V)_\etale$ 上的 $x^*\mathcal{F}$ 拟凝聚；
例如由《site 上的模》引理
\ref{sites-modules-lemma-local-final-object} 可得。
由《叠上的层》引理
\ref{stacks-sheaves-lemma-characterize-quasi-coherent}，
可知 $\mathcal{F}$ 拟凝聚。
\end{proof}

\noindent
我们回忆，在《site 上的模》第
\ref{sites-modules-section-local} 节中，已对任意赋环拓扑斯定义了拟凝聚模的概念。

\begin{lemma}
\label{stacks-cohomology-lemma-shriek-quasi-coherent}
设 $\mathcal{X}$ 是代数叠。记号如引理
\ref{stacks-cohomology-lemma-lisse-etale}。
\begin{enumerate}
\item 设 $\mathcal{H}$ 是 $\mathcal{X}$ 的光滑-\'etale site 上的拟凝聚
$\mathcal{O}_{\mathcal{X}_{lisse,\etale}}$-模。则 $g_!\mathcal{H}$
是 $\mathcal{X}$ 上的拟凝聚模。
\item 设 $\mathcal{H}$ 是 $\mathcal{X}$ 的平坦-fppf site 上的拟凝聚
$\mathcal{O}_{\mathcal{X}_{flat,fppf}}$-模。则 $g_!\mathcal{H}$
是 $\mathcal{X}$ 上的拟凝聚模。
\end{enumerate}
\end{lemma}

\begin{proof}
选取概形 $U$ 和满的光滑态射 $x : U \to \mathcal{X}$。
由《site 上的模》定义
\ref{sites-modules-definition-site-local}，存在 \'etale
（相应地，fppf）覆盖
$\{U_i \to U\}_{i \in I}$，使得每个拉回 $f_i^{-1}\mathcal{H}$
都有全局表示（参见《site 上的模》定义
\ref{sites-modules-definition-global}）。
这里 $f_i : U_i \to \mathcal{X}$ 是代数叠态射
$U_i \to U \to \mathcal{X}$ 的复合。
（回忆，拉回“就是”到 $\mathcal{X}/f_i$ 的限制；参见《叠上的层》定义
\ref{stacks-sheaves-definition-pullback} 及其后的讨论。）
由引理 \ref{stacks-cohomology-lemma-lisse-etale-functorial}，每个 $f_i$
光滑（相应地，平坦），因而
$f_i^{-1}g_!\mathcal{H} = g_{i, !}(f'_i)^{-1}\mathcal{H}$。
利用引理 \ref{stacks-cohomology-lemma-check-qc-on-etale-covering}，
可将引理的陈述归约到 $\mathcal{H}$ 有全局表示的情形。
设有 $\mathcal{O}$-模的正合列
$$
\bigoplus\nolimits_{j \in J} \mathcal{O} \longrightarrow
\bigoplus\nolimits_{i \in I} \mathcal{O} \longrightarrow
\mathcal{H} \longrightarrow 0
$$
其中 $\mathcal{O} = \mathcal{O}_{\mathcal{X}_{lisse,\etale}}$
（相应地，$\mathcal{O} = \mathcal{O}_{\mathcal{X}_{flat,fppf}}$）。
由于 $g_!$ 与任意余极限可交换
（因为它是左伴随函子；参见引理
\ref{stacks-cohomology-lemma-lisse-etale-modules} 和《范畴》引理
\ref{categories-lemma-adjoint-exact}），得到正合列
$$
\bigoplus\nolimits_{j \in J} g_!\mathcal{O} \longrightarrow
\bigoplus\nolimits_{i \in I} g_!\mathcal{O} \longrightarrow
g_!\mathcal{H} \longrightarrow 0
$$
引理 \ref{stacks-cohomology-lemma-lisse-etale-structure-sheaf} 表明
$g_!\mathcal{O} = \mathcal{O}_\mathcal{X}$。
在情形 (2) 中已得结论。在情形 (1) 中，应用《叠上的层》引理
\ref{stacks-sheaves-lemma-characterize-quasi-coherent-bis} 即得结论。
\end{proof}

\begin{lemma}
\label{stacks-cohomology-lemma-quasi-coherent}
设 $\mathcal{X}$ 是代数叠。
\begin{enumerate}
\item 对光滑-\'etale site，以 $g$ 表示引理
\ref{stacks-cohomology-lemma-lisse-etale} 中的态射。则：
\begin{enumerate}
\item 函子 $g^{-1}$ 和 $g_!$ 定义互逆函子
$$
\xymatrix{
\QCoh(\mathcal{O}_\mathcal{X}) \ar@<1ex>[r]^-{g^{-1}} &
\QCoh(\mathcal{X}_{lisse,\etale},
\mathcal{O}_{\mathcal{X}_{lisse,\etale}}) \ar@<1ex>[l]^-{g_!}
}
$$
\item 若 $\mathcal{F}$ 属于 $\textit{LQCoh}^{fbc}(\mathcal{O}_\mathcal{X})$，
则 $g^{-1}\mathcal{F}$ 属于
$\QCoh(\mathcal{O}_{\mathcal{X}_{lisse,\etale}})$；
\item $Q(\mathcal{F}) = g_!g^{-1}\mathcal{F}$，其中 $Q$ 如引理
\ref{stacks-cohomology-lemma-adjoint} 所述。
\end{enumerate}
\item 对平坦-fppf site，以 $g$ 表示引理
\ref{stacks-cohomology-lemma-lisse-etale} 中的态射。则：
\begin{enumerate}
\item 函子 $g^{-1}$ 和 $g_!$ 定义互逆函子
$$
\xymatrix{
\QCoh(\mathcal{O}_\mathcal{X}) \ar@<1ex>[r]^-{g^{-1}} &
\QCoh(\mathcal{X}_{flat,fppf},
\mathcal{O}_{\mathcal{X}_{flat,fppf}}) \ar@<1ex>[l]^-{g_!}
}
$$
\item 若 $\mathcal{F}$ 属于 $\textit{LQCoh}^{fbc}(\mathcal{O}_\mathcal{X})$，
则 $g^{-1}\mathcal{F}$ 属于
$\QCoh(\mathcal{O}_{\mathcal{X}_{flat,fppf}})$；
\item $Q(\mathcal{F}) = g_!g^{-1}\mathcal{F}$，其中 $Q$ 如引理
\ref{stacks-cohomology-lemma-adjoint} 所述。
\end{enumerate}
\end{enumerate}
\end{lemma}

\begin{proof}
沿任意赋环拓扑斯态射的拉回都保持拟凝聚模范畴；
参见《site 上的模》引理 \ref{sites-modules-lemma-local-pullback}。
因而 $g^{-1}$ 保持拟凝聚模范畴；这里我们用到
$
\QCoh(\mathcal{O}_\mathcal{X}) =
\QCoh(\mathcal{X}_\etale, \mathcal{O}_\mathcal{X})
$
这由《叠上的层》引理
\ref{stacks-sheaves-lemma-characterize-quasi-coherent-bis} 得出。
由引理 \ref{stacks-cohomology-lemma-shriek-quasi-coherent}，$g_!$ 也有同样的性质。
由引理 \ref{stacks-cohomology-lemma-lisse-etale}，我们知道
$\mathcal{H} \to g^{-1}g_!\mathcal{H}$ 是同构。
反之，若 $\mathcal{F}$ 属于 $\QCoh(\mathcal{O}_\mathcal{X})$，
则态射 $g_!g^{-1}\mathcal{F} \to \mathcal{F}$ 是 $\mathcal{X}$ 上拟凝聚模之间的态射，
其到任意在 $\mathcal{X}$ 上光滑的概形的限制都是同构。
因而《叠上的层》第
\ref{stacks-sheaves-section-quasi-coherent-presentation} 节和第
\ref{stacks-sheaves-section-quasi-coherent-algebraic-stacks} 节的讨论
（与展示上的拟凝聚模比较）表明它是同构。
这证明了 (1)(a) 和 (2)(a)。

\medskip\noindent
设 $\mathcal{F}$ 是 $\textit{LQCoh}^{fbc}(\mathcal{O}_\mathcal{X})$ 的对象。
由引理 \ref{stacks-cohomology-lemma-adjoint-kernel-parasitic}，态射
$Q(\mathcal{F}) \to \mathcal{F}$ 的核与余核是寄生的。
因此，由引理 \ref{stacks-cohomology-lemma-parasitic-in-terms-flat-fppf}，并利用
$g^* = g^{-1}$ 正合，可知
$g^*Q(\mathcal{F}) \to g^*\mathcal{F}$ 是同构。
于是 $g^*\mathcal{F}$ 拟凝聚。这证明了 (1)(b) 和 (2)(b)。
最后，(1)(c) 和 (2)(c) 成立，因为由上述论证，
$g_!g^*Q(\mathcal{F}) \to Q(\mathcal{F})$ 是同构。
\end{proof}

\begin{lemma}
\label{stacks-cohomology-lemma-quasi-coherent-weak-serre}
设 $\mathcal{X}$ 是代数叠。
\begin{enumerate}
\item $\QCoh(\mathcal{O}_{\mathcal{X}_{lisse,\etale}})$
是 $\textit{Mod}(\mathcal{O}_{\mathcal{X}_{lisse,\etale}})$
的弱 Serre 子范畴。
\item $\QCoh(\mathcal{O}_{\mathcal{X}_{flat,fppf}})$
是 $\textit{Mod}(\mathcal{O}_{\mathcal{X}_{flat,fppf}})$
的弱 Serre 子范畴。
\end{enumerate}
\end{lemma}

\begin{proof}
我们将验证《同调》引理
\ref{homology-lemma-characterize-weak-serre-subcategory} 的条件 (1)、(2)、(3)、(4)。

\medskip\noindent
由于 $0$ 是任意赋环 site 上的拟凝聚模，条件 (1) 成立。

\medskip\noindent
% P12-ERR-0033: see ../control/ERRATA_CANDIDATES.jsonl
按定义，$\QCoh(\mathcal{O})$ 是 $\textit{Mod}(\mathcal{O})$ 的严格充满子范畴，
故条件 (2) 成立。

\medskip\noindent
设 $\varphi : \mathcal{G} \to \mathcal{F}$ 是
$\mathcal{X}_{lisse,\etale}$ 或 $\mathcal{X}_{flat,fppf}$ 上拟凝聚模之间的态射。
有 $g^*g_!\mathcal{F} = \mathcal{F}$；对 $\mathcal{G}$ 和 $\varphi$ 也类似，
参见引理 \ref{stacks-cohomology-lemma-lisse-etale-modules}。
由引理 \ref{stacks-cohomology-lemma-shriek-quasi-coherent}，
$g_!\mathcal{F}$ 和 $g_!\mathcal{G}$ 是拟凝聚
$\mathcal{O}_\mathcal{X}$-模。由《叠上的层》引理
\ref{stacks-sheaves-lemma-quasi-coherent-algebraic-stack}，
$\Coker(g_!\varphi)$ 是 $\mathcal{X}$ 上的拟凝聚模
（并且是 $\mathcal{X}$ 上拟凝聚模范畴中的余核）。
由于 $g^*$ 正合（参见引理 \ref{stacks-cohomology-lemma-lisse-etale}），
$g^*\Coker(g_!\varphi) = \Coker(g^*g_!\varphi) = \Coker(\varphi)$
也拟凝聚（参见引理 \ref{stacks-cohomology-lemma-quasi-coherent}）。
由命题 \ref{stacks-cohomology-proposition-loc-qcoh-flat-base-change}，核
$\Ker(g_!\varphi)$ 属于
$\textit{LQCoh}^{fbc}(\mathcal{O}_\mathcal{X})$。
由于 $g^*$ 正合，有
$g^*\Ker(g_!\varphi) = \Ker(g^*g_!\varphi) = \Ker(\varphi)$。
由引理 \ref{stacks-cohomology-lemma-quasi-coherent}，$g^*$ 把
$\textit{LQCoh}^{fbc}(\mathcal{O}_\mathcal{X})$ 的对象映为拟凝聚模，
因而 $\Ker(\varphi)$ 也拟凝聚。这证明了 (3)。

\medskip\noindent
最后，假设
$$
0 \to \mathcal{F} \to \mathcal{E} \to \mathcal{G} \to 0
$$
是 $\mathcal{O}_{\mathcal{X}_{lisse,\etale}}$-模的扩张
（相应地，$\mathcal{O}_{\mathcal{X}_{flat,fppf}}$-模的扩张），
且 $\mathcal{F}$ 和 $\mathcal{G}$ 拟凝聚。为证明 (4) 并完成证明，
需说明 $\mathcal{E}$ 在 $\mathcal{X}_{lisse,\etale}$ 上
（相应地，在 $\mathcal{X}_{flat,fppf}$ 上）拟凝聚。
% P12-ERR-0034: see ../control/ERRATA_CANDIDATES.jsonl
设 $U$ 是 $\mathcal{X}_{lisse,\etale}$ 的对象
（相应地，$\mathcal{X}_{flat,fppf}$ 的对象）；我们把 $U$ 视为在
$\mathcal{X}$ 上光滑（相应地，平坦）的概形。
需证明 $\mathcal{E}$ 到 $U_{lisse,\etale}$ 的限制
% P12-ERR-0035: see ../control/ERRATA_CANDIDATES.jsonl
（相应地，$=U_{flat,fppf}$）拟凝聚。
因此可假设 $\mathcal{X} = U$ 是概形。
由于 $\mathcal{G}$ 在 $U_{lisse,\etale}$ 上
（相应地，在 $U_{flat,fppf}$ 上）拟凝聚，
用一个 \'etale（相应地，fppf）覆盖的各成员取代 $U$ 后，
可假设 $\mathcal{G}$ 有表示
$$
\bigoplus\nolimits_{j \in J} \mathcal{O} \longrightarrow
\bigoplus\nolimits_{i \in I} \mathcal{O} \longrightarrow
\mathcal{G} \longrightarrow 0
$$
该表示位于 $U_{lisse,\etale}$ 上
（相应地，$U_{flat,fppf}$ 上），其中 $\mathcal{O}$
是该 site 上的结构层。还可假设 $U$ 仿射。
由于 $\mathcal{F}$ 拟凝聚，有
$$
H^1(U_{lisse,\etale}, \mathcal{F}) = 0,
\quad\text{分别}\quad
H^1(U_{flat,fppf}, \mathcal{F}) = 0
$$
事实上，$\mathcal{F}$ 是 $U$ 的大 site 上某个拟凝聚模
$\mathcal{F}'$ 的拉回（由引理 \ref{stacks-cohomology-lemma-quasi-coherent}）；
$\mathcal{F}$ 与 $\mathcal{F}'$ 的上同调一致
（由引理 \ref{stacks-cohomology-lemma-lisse-etale-cohomology}）；
而我们知道，仿射概形 $U$ 的大 site 上 $\mathcal{F}'$ 的上同调为零
（要在当前情形中得到这一点，需将《下降》命题
\ref{descent-proposition-equivalence-quasi-coherent} 和
\ref{descent-proposition-same-cohomology-quasi-coherent}
与《概形上的上同调》引理
\ref{coherent-lemma-quasi-coherent-affine-cohomology-zero} 结合起来）。
因此，可将态射
$\bigoplus_{i \in I} \mathcal{O} \to \mathcal{G}$
提升到 $\mathcal{E}$。图追踪表明，得到正合列
$$
\bigoplus\nolimits_{j \in J} \mathcal{O} \to
\mathcal{F} \oplus \bigoplus\nolimits_{i \in I} \mathcal{O}
\to
\mathcal{E} \to 0
$$
由上面已证明的 (3)，可知 $\mathcal{E}$ 如需拟凝聚。
\end{proof}

\section{局部 Noether 叠上的凝聚层}
\label{stacks-cohomology-section-coherent-sheaves}

\noindent
本节对应于《空间的上同调》第
\ref{spaces-cohomology-section-coherent} 节。
在《site 上的模》第 \ref{sites-modules-section-local} 节中，
我们已对任意赋环拓扑斯定义凝聚模的概念。
然而，对任意代数叠 $\mathcal{X}$，凝聚
$\mathcal{O}_\mathcal{X}$-模范畴为零；本质上是因为 site $\mathcal{X}$
包含太多非 Noether 对象（即使 $\mathcal{X}$ 本身局部 Noether）。
我们改用下述引理定义凝聚模。

\begin{lemma}
\label{stacks-cohomology-lemma-coherent-Noetherian}
设 $\mathcal{X}$ 是局部 Noether 代数叠，
$\mathcal{F}$ 是 $\mathcal{O}_\mathcal{X}$-模。以下条件等价：
\begin{enumerate}
\item $\mathcal{F}$ 是拟凝聚的有限型
$\mathcal{O}_\mathcal{X}$-模；
\item $\mathcal{F}$ 是有限表示的
$\mathcal{O}_\mathcal{X}$-模；
\item $\mathcal{F}$ 拟凝聚，且对任意态射
$f : U \to \mathcal{X}$，若 $U$ 是局部 Noether 代数空间，
则拉回 $f^*\mathcal{F}|_{U_\etale}$ 凝聚；
\item $\mathcal{F}$ 拟凝聚，且存在代数空间 $U$ 和态射
$f : U \to \mathcal{X}$，其局部有限型、平坦且满，
并且拉回 $f^*\mathcal{F}|_{U_\etale}$ 凝聚。
\end{enumerate}
\end{lemma}

\begin{proof}
设 $f : U \to \mathcal{X}$ 如 (4) 所述。
则 $U$ 局部 Noether（《叠的态射》引理
\ref{stacks-morphisms-lemma-locally-finite-type-locally-noetherian}），
因而引理的陈述有意义。此外，由《叠的态射》引理
\ref{stacks-morphisms-lemma-noetherian-finite-type-finite-presentation}，
$f$ 局部有限表示。
设 $x$ 是 $\mathcal{X}$ 的一个对象，位于概形 $V$ 上。
为证明 (2)，需说明：用 $V$ 的某个 fppf 覆盖的各成员取代
$V$ 后，限制 $x^*\mathcal{F}$ 在
$\mathcal{X}/x \cong (\Sch/V)_{fppf}$ 上有全局有限表示。
投影 $W = U \times_\mathcal{X} V \to V$
局部有限表示、平坦且满。
因而，可以用 $W$ 的一个由概形组成的 \'etale 覆盖的各成员
取代 $V$，并假设有态射 $h : V \to U$ 满足 $f \circ h = x$。
由于 $\mathcal{F}$ 拟凝聚，限制 $x^*\mathcal{F}$ 是
$h_{small}^*(f^*\mathcal{F})|_{U_\etale}$ 沿 $\pi_V$ 的拉回；
参见《叠上的层》引理
\ref{stacks-sheaves-lemma-compare-quasi-coherent}。
% P12-ERR-0036: see ../control/ERRATA_CANDIDATES.jsonl
按假设，$f^*\mathcal{F}|_{U_\etale}$ 在 \'etale 拓扑中局部地有有限表示，
因而得到 (4) $\Rightarrow$ (2)。

\medskip\noindent
对任意赋环拓扑斯，(2) 都意味着 (1)（由定义立即得出）。
性质“有限型”和“拟凝聚”在任意赋环拓扑斯态射的拉回下保持；
参见《site 上的模》引理 \ref{sites-modules-lemma-local-pullback}。
因而 (1) 意味着 (3)；参见《空间的上同调》引理
\ref{spaces-cohomology-lemma-coherent-Noetherian}。
最后，(3) 显然意味着 (4)。
\end{proof}

\begin{definition}
\label{stacks-cohomology-definition-coherent}
设 $\mathcal{X}$ 是局部 Noether 代数叠。
设 $\mathcal{O}_\mathcal{X}$-模为 $\mathcal{F}$。若 $\mathcal{F}$ 满足引理
\ref{stacks-cohomology-lemma-coherent-Noetherian} 的一个（从而全部）等价条件，
则称其为 {\it 凝聚的}。
% P12-ERR-0037: see ../control/ERRATA_CANDIDATES.jsonl
凝聚 $\mathcal{O}_\mathcal{X}$-模范畴记为
$\textit{Coh}(\mathcal{O}_\mathcal{X})$。
\end{definition}

\begin{lemma}
\label{stacks-cohomology-lemma-elementary-coherent}
设 $\mathcal{X}$ 是局部 Noether 代数叠。
模 $\mathcal{O}_\mathcal{X}$ 是凝聚的；任意可逆
$\mathcal{O}_\mathcal{X}$-模是凝聚的；更一般地，任意有限局部自由
$\mathcal{O}_\mathcal{X}$-模是凝聚的。
\end{lemma}

\begin{proof}
由定义和《空间的上同调》引理
\ref{spaces-cohomology-lemma-coherent-Noetherian} 得出。
\end{proof}

\begin{lemma}
\label{stacks-cohomology-lemma-pullback-coherent}
设 $f : \mathcal{X} \to \mathcal{Y}$ 是局部 Noether 代数叠之间的态射。
则 $f^*$ 把 $\mathcal{Y}$ 上的凝聚模映为 $\mathcal{X}$ 上的凝聚模。
\end{lemma}

\begin{proof}
这直接由定义以及任意赋环拓扑斯态射的拉回保持有限表示模这一事实得出；
参见《site 上的模》引理 \ref{sites-modules-lemma-local-pullback}。
\end{proof}

\begin{lemma}
\label{stacks-cohomology-lemma-coherent-abelian-Noetherian}
设 $\mathcal{X}$ 是局部 Noether 代数叠。
凝聚 $\mathcal{O}_\mathcal{X}$-模范畴是阿贝范畴。
若 $\varphi : \mathcal{F} \to \mathcal{G}$ 是凝聚
$\mathcal{O}_\mathcal{X}$-模之间的态射，则：
\begin{enumerate}
\item 在 $\textit{Mod}(\mathcal{O}_\mathcal{X})$ 中计算的余核
$\Coker(\varphi)$ 是凝聚 $\mathcal{O}_\mathcal{X}$-模；
\item 在 $\textit{Mod}(\mathcal{O}_\mathcal{X})$ 中计算的像
$\Im(\varphi)$ 是凝聚 $\mathcal{O}_\mathcal{X}$-模；
\item 在 $\textit{Mod}(\mathcal{O}_\mathcal{X})$ 中计算的核
$\Ker(\varphi)$ 可能不凝聚，但它属于
$\textit{LQCoh}^{fbc}(\mathcal{O}_\mathcal{X})$，而 $Q(\Ker(\varphi))$
凝聚，并且是 $\textit{Coh}(\mathcal{O}_\mathcal{X})$ 中 $\varphi$ 的核。
\end{enumerate}
包含函子 $\textit{Coh}(\mathcal{O}_\mathcal{X}) \to
\QCoh(\mathcal{O}_\mathcal{X})$ 正合。
\end{lemma}

\begin{proof}
在 $\textit{Coh}(\mathcal{O}_\mathcal{X})$ 中取核、像和余核的规则，
与注 \ref{stacks-cohomology-remark-QCoh-abelian} 中对拟凝聚模的规定一致。
因而，只需证明拟凝聚模
$\Coker(\varphi)$、$\Im(\varphi)$ 和 $Q(\Ker(\varphi))$ 凝聚。
由引理 \ref{stacks-cohomology-lemma-coherent-Noetherian}，只需在到 $U_\etale$ 的限制上证明，
其中 $f : U \to \mathcal{X}$ 是某个满的光滑态射。
函子 $\mathcal{F} \mapsto f^*\mathcal{F}|_{U_\etale}$ 正合。
因此 $f^*\Coker(\varphi)$ 和 $f^*\Im(\varphi)$ 是凝聚
$\mathcal{O}_U$-模之间某个态射的余核和像，从而如需凝聚。
由引理 \ref{stacks-cohomology-lemma-parasitic}，函子
$\mathcal{F} \mapsto f^*\mathcal{F}|_{U_\etale}$ 消去寄生模。
因而，由引理 \ref{stacks-cohomology-lemma-adjoint-kernel-parasitic} 的第 (2) 部分，
$f^*Q(\Ker(\varphi))|_{U_\etale} = f^*\Ker(\varphi)|_{U_\etale}$。
以同样方式可得 $Q(\Ker(\varphi))$ 凝聚。
\end{proof}

\begin{lemma}
\label{stacks-cohomology-lemma-extension-coherent}
设 $\mathcal{X}$ 是局部 Noether 代数叠。
给定 $\textit{Mod}(\mathcal{O}_\mathcal{X})$ 中的短正合列
$0 \to \mathcal{F}_1 \to \mathcal{F}_2 \to \mathcal{F}_3 \to 0$，
若 $\mathcal{F}_1$ 和 $\mathcal{F}_3$ 凝聚，则 $\mathcal{F}_2$ 凝聚。
\end{lemma}

\begin{proof}
由《叠上的层》引理
\ref{stacks-sheaves-lemma-quasi-coherent-algebraic-stack} 的第 (7) 部分，
$\mathcal{F}_2$ 拟凝聚。于是，可以通过限制到 $U_\etale$
来检验 $\mathcal{F}_2$ 凝聚，其中 $U \to \mathcal{X}$ 满且光滑。
这由《空间的上同调》引理
\ref{spaces-cohomology-lemma-coherent-abelian-Noetherian} 得出。
略去一些细节。
\end{proof}

\noindent
凝聚模构成拟凝聚 $\mathcal{O}_\mathcal{X}$-模范畴的 Serre 子范畴。
对一般赋环拓扑斯上的模，这并不成立。

\begin{lemma}
\label{stacks-cohomology-lemma-coherent-Noetherian-quasi-coherent-sub-quotient}
设 $\mathcal{X}$ 是局部 Noether 代数叠。则
$\textit{Coh}(\mathcal{O}_\mathcal{X})$ 是
$\QCoh(\mathcal{O}_\mathcal{X})$ 的 Serre 子范畴。
设 $\varphi : \mathcal{F} \to \mathcal{G}$ 是拟凝聚
$\mathcal{O}_\mathcal{X}$-模之间的态射。则：
\begin{enumerate}
\item 若 $\mathcal{F}$ 凝聚且 $\varphi$ 满，则 $\mathcal{G}$ 凝聚；
\item 若 $\mathcal{F}$ 凝聚，则 $\Im(\varphi)$ 凝聚；
% P12-ERR-0038: see ../control/ERRATA_CANDIDATES.jsonl
\item 若 $\mathcal{G}$ 凝聚且 $\Ker(\varphi)$ 寄生，则
$\mathcal{F}$ 凝聚。
\end{enumerate}
\end{lemma}

\begin{proof}
选取概形 $U$ 和满的光滑态射 $f : U \to \mathcal{X}$。
则函子
$f^* : \QCoh(\mathcal{O}_\mathcal{X}) \to \QCoh(\mathcal{O}_U)$
正合（引理 \ref{stacks-cohomology-lemma-flat-pullback-quasi-coherent}）。
此外，按定义，$\textit{Coh}(\mathcal{O}_\mathcal{X})$
是 $\QCoh(\mathcal{O}_\mathcal{X})$ 的充满子范畴，
由这些满足 $f^*\mathcal{F}$ 属于
$\textit{Coh}(\mathcal{O}_U)$ 的对象 $\mathcal{F}$ 组成。
由此以及 $U$ 的相应事实，立即得到
$\textit{Coh}(\mathcal{O}_\mathcal{X})$ 是
$\QCoh(\mathcal{O}_\mathcal{X})$ 的 Serre 子范畴；
参见《空间的上同调》引理
\ref{spaces-cohomology-lemma-coherent-abelian-Noetherian} 和
\ref{spaces-cohomology-lemma-coherent-Noetherian-quasi-coherent-sub-quotient}。
我们略去 (1)、(2)、(3) 的证明。提示：与引理
\ref{stacks-cohomology-lemma-coherent-abelian-Noetherian} 的证明比较。
\end{proof}

\noindent
设 $\mathcal{X}$ 是局部 Noether 代数叠。
设 $U$ 是代数空间，$f : U \to \mathcal{X}$
满、局部有限表示且平坦。
注意 $U$ 局部 Noether
（《叠的态射》引理
\ref{stacks-morphisms-lemma-locally-finite-type-locally-noetherian}）。
设 $(U, R, s, t, c)$ 是代数空间群胚，
$f_{can} : [U/R] \to \mathcal{X}$ 是《代数叠》引理
\ref{algebraic-lemma-map-space-into-stack} 和注
\ref{algebraic-remark-flat-fp-presentation} 中构造的同构。
与《叠上的层》第
\ref{stacks-sheaves-section-quasi-coherent-algebraic-stacks} 节一样，得到等价
$$
\QCoh(\mathcal{O}_\mathcal{X})
\cong
\QCoh(\mathcal{O}_{[U/R]})
\cong
\QCoh(U, R, s, t, c)
$$
其中第二个等价是《叠上的层》命题
\ref{stacks-sheaves-proposition-quasi-coherent}。
回忆，在《空间中的群胚》第
\ref{spaces-groupoids-section-colimits} 节中，我们已定义充满子范畴
$$
\textit{Coh}(U, R, s, t, c) \subset \QCoh(U, R, s, t, c)
$$
其 {\it 凝聚模} 是那些满足下列条件的
$(\mathcal{G}, \alpha)$：$\mathcal{G}$ 是凝聚 $\mathcal{O}_U$-模。

\begin{lemma}
\label{stacks-cohomology-lemma-coherent-presentation}
在上述情形中，等价
$\QCoh(\mathcal{O}_\mathcal{X}) \cong \QCoh(U, R, s, t, c)$
把凝聚层映为凝聚层，反之亦然；即它诱导等价
$\textit{Coh}(\mathcal{O}_\mathcal{X}) \cong \textit{Coh}(U, R, s, t, c)$。
\end{lemma}

\begin{proof}
这直接由凝聚 $\mathcal{O}_\mathcal{X}$-模的定义得出。
为便于记录：上述材料用到《叠的态射》引理
\ref{stacks-morphisms-lemma-locally-finite-type-locally-noetherian}、
《代数叠》引理 \ref{algebraic-lemma-map-space-into-stack} 和注
\ref{algebraic-remark-flat-fp-presentation}、
《叠上的层》第
\ref{stacks-sheaves-section-quasi-coherent-algebraic-stacks} 节、
《叠上的层》命题 \ref{stacks-sheaves-proposition-quasi-coherent}，
以及《空间中的群胚》第
\ref{spaces-groupoids-section-colimits} 节。
\end{proof}

\begin{lemma}
\label{stacks-cohomology-lemma-coherent-hom}
设 $\mathcal{X}$ 是局部 Noether 代数叠。
% P12-ERR-0039: see ../control/ERRATA_CANDIDATES.jsonl
设 $\mathcal{F}$ 和 $\mathcal{G}$ 是凝聚 $\mathcal{O}_\mathcal{X}$-模。
则引理 \ref{stacks-cohomology-lemma-internal-hom-fp-into-qcoh}
中构造的内部 Hom $hom(\mathcal{F}, \mathcal{G})$
是凝聚 $\mathcal{O}_\mathcal{X}$-模。
\end{lemma}

\begin{proof}
设 $U \to \mathcal{X}$ 是由概形给出的光滑满态射。
由条目 (\ref{stacks-cohomology-item-hom-restriction})（位于第
\ref{stacks-cohomology-section-further-remarks} 节），$hom(\mathcal{F}, \mathcal{G})$
到 $U$ 的限制是各限制的 Hom 层。
因而引理由代数空间的情形得出；参见《空间的上同调》引理
\ref{spaces-cohomology-lemma-tensor-hom-coherent}。
\end{proof}

\section{Noether 叠上的凝聚层}
\label{stacks-cohomology-section-coherent-on-noetherian}

\noindent
本节对应于《空间的上同调》第
\ref{spaces-cohomology-section-coherent-quasi-compact} 节。

\begin{lemma}
\label{stacks-cohomology-lemma-directed-colimit-coherent}
设 $\mathcal{X}$ 是 Noether 代数叠。每个拟凝聚
$\mathcal{O}_\mathcal{X}$-模都是其凝聚子模的滤过余极限。
\end{lemma}

\begin{proof}
设 $\mathcal{F}$ 是拟凝聚 $\mathcal{O}_\mathcal{X}$-模。
若 $\mathcal{G}, \mathcal{H} \subset \mathcal{F}$ 是凝聚
$\mathcal{O}_\mathcal{X}$-子模，则态射
$\mathcal{G} \oplus \mathcal{H} \to \mathcal{F}$ 的像是另一个同时包含二者的凝聚
$\mathcal{O}_\mathcal{X}$-子模；参见引理
\ref{stacks-cohomology-lemma-coherent-Noetherian-quasi-coherent-sub-quotient}。
由此可知该系统有向。因而，现在只需证明
$\mathcal{F}$ 可写为凝聚模的滤过余极限，
因为随后可在 $\mathcal{F}$ 中取这些模的像，得到足够多的凝聚子模。

\medskip\noindent
设 $U$ 是仿射概形，$U \to \mathcal{X}$ 是满的光滑态射
（《叠的性质》引理
\ref{stacks-properties-lemma-quasi-compact-stack}）。
置 $R = U \times_\mathcal{X} U$，于是如《代数叠》引理
\ref{algebraic-lemma-stack-presentation} 中那样，有 $\mathcal{X} = [U/R]$。
由引理 \ref{stacks-cohomology-lemma-coherent-presentation}，
% P12-ERR-0040: see ../control/ERRATA_CANDIDATES.jsonl
有 $\QCoh(\mathcal{O}_X) = \QCoh(U, R, s, t, c)$ 和
$\textit{Coh}(\mathcal{O}_X) = \textit{Coh}(U, R, s, t, c)$。
这样便归约到对 $\QCoh(U, R, s, t, c)$ 证明相应陈述的问题。
这就是《空间中的群胚》引理
\ref{spaces-groupoids-lemma-colimit-coherent}；我们在下一段检验它的假设。

\medskip\noindent
我们敦促读者跳过证明的余下部分。
仿射概形 $U$ 是 Noether 的；这由我们对 $\mathcal{X}$
局部 Noether 的定义得出，参见《叠的性质》定义
\ref{stacks-properties-definition-type-property} 和注
\ref{stacks-properties-remark-list-properties-local-smooth-topology}。
投影态射 $s, t : R \to U$ 光滑（参见上述参照），
并且拟分离且拟紧（《叠的态射》引理
\ref{stacks-morphisms-lemma-quasi-compact-quasi-separated-permanence}）。
特别地，$R$ 是在 $U$ 上光滑的拟紧且拟分离代数空间，
因而是 Noether 的（《空间的态射》引理
\ref{spaces-morphisms-lemma-finite-presentation-noetherian}）。
\end{proof}
% Generated mechanically from the frozen stacks-perfect.tex target.
% Do not edit this generated file; edit the bound locale source instead.

% OK, start here.
%

\setcounter{chapter}{103}
\chapter{叠的导出范畴}


\phantomsection
\label{stacks-perfect-section-phantom}





\section{引言}
\label{stacks-perfect-section-introduction}

\noindent
本章讨论与代数叠相伴的导出范畴，尤其包括拟凝聚层的导出范畴；
也就是说，我们将证明概形情形（参见《概形的导出范畴》第
\ref{perfect-section-introduction} 节）和代数空间情形
（参见《空间的导出范畴》第
\ref{spaces-perfect-section-introduction} 节）相应结果的类似结论。
本章结果与 \cite{LM-B} 中结果的主要区别是，我们始终使用“大 site”。
阅读本章之前，请先快速浏览《代数叠上的层》和《代数叠的上同调》两章，
其中引入了本章所用的术语。



\section{约定、记号与语言滥用}
\label{stacks-perfect-section-conventions}

\noindent
我们继续采用《叠的性质》第
\ref{stacks-properties-section-conventions} 节中引入的约定和语言滥用。
记号依照《叠的上同调》第
\ref{stacks-cohomology-section-notation} 节。












\section{光滑-\'etale site 与平坦-fppf site}
\label{stacks-perfect-section-lisse-etale}

\noindent
本节是《叠的上同调》第
\ref{stacks-cohomology-section-lisse-etale} 节在导出范畴情形下的类似版本。

\begin{lemma}
\label{stacks-perfect-lemma-shriek-derived}
设 $\mathcal{X}$ 为代数叠。采用《叠的上同调》引理
\ref{stacks-cohomology-lemma-lisse-etale} 和
\ref{stacks-cohomology-lemma-lisse-etale-modules} 中的记号。
\begin{enumerate}
\item 函子
$g_! : \textit{Ab}(\mathcal{X}_{lisse,\etale}) \to
\textit{Ab}(\mathcal{X}_\etale)$
有左导出函子
$$
Lg_! :
D(\mathcal{X}_{lisse,\etale})
\longrightarrow
D(\mathcal{X}_\etale)
$$
它是 $g^{-1}$ 的左伴随，并且 $g^{-1}Lg_! = \text{id}$。
\item 函子 $g_! : 
\textit{Mod}(\mathcal{X}_{lisse,\etale},
\mathcal{O}_{\mathcal{X}_{lisse,\etale}}) \to
\textit{Mod}(\mathcal{X}_\etale, \mathcal{O}_{\mathcal{X}})$
有左导出函子
$$
Lg_! :
D(\mathcal{O}_{\mathcal{X}_{lisse,\etale}})
\longrightarrow
D(\mathcal{X}_\etale, \mathcal{O}_\mathcal{X})
$$
它是 $g^*$ 的左伴随，并且 $g^*Lg_! = \text{id}$。
\item 函子 $g_! : \textit{Ab}(\mathcal{X}_{flat,fppf}) \to
\textit{Ab}(\mathcal{X}_{fppf})$
有左导出函子
$$
Lg_! :
D(\mathcal{X}_{flat, fppf})
\longrightarrow
D(\mathcal{X}_{fppf})
$$
它是 $g^{-1}$ 的左伴随，并且 $g^{-1}Lg_! = \text{id}$。
\item 函子 $g_! :
\textit{Mod}(\mathcal{X}_{flat,fppf},
\mathcal{O}_{\mathcal{X}_{flat,fppf}}) \to
\textit{Mod}(\mathcal{X}_{fppf}, \mathcal{O}_{\mathcal{X}})$
有左导出函子
$$
Lg_! :
D(\mathcal{O}_{\mathcal{X}_{flat, fppf}})
\longrightarrow
D(\mathcal{O}_\mathcal{X})
$$
它是 $g^*$ 的左伴随，并且 $g^*Lg_! = \text{id}$。
\end{enumerate}
警告：先验地并不清楚模上的 $Lg_!$ 是否与阿贝尔层上的 $Lg_!$ 一致；
参见《site 上的上同调》评注
\ref{sites-cohomology-remark-when-derived-shriek-equal}。
\end{lemma}

\begin{proof}
函子 $Lg_!$ 的存在性及其与 $g^*$ 的伴随性由《site 上的上同调》引理
\ref{sites-cohomology-lemma-existence-derived-lower-shriek} 给出。
（对阿贝尔层的情形，取常值层 $\mathbf{Z}$ 为结构层。）
此外，它在复形 $\mathcal{H}^\bullet$ 上的计算方式是：
取一个适当的左分解
$\mathcal{K}^\bullet \to \mathcal{H}^\bullet$，
再将函子 $g_!$ 作用于 $\mathcal{K}^\bullet$。
由《叠的上同调》引理
\ref{stacks-cohomology-lemma-lisse-etale-modules} 和
\ref{stacks-cohomology-lemma-lisse-etale}，
$g^{-1}g_!\mathcal{K}^\bullet = \mathcal{K}^\bullet$，
故每种情形下的最后一个断言都成立。
\end{proof}










\begin{lemma}
\label{stacks-perfect-lemma-lisse-etale-functorial-derived}
采用《叠的上同调》引理
\ref{stacks-cohomology-lemma-lisse-etale-functorial} 中的假设和记号。
在无界导出范畴上有
$$
g^{-1} \circ Rf_* = Rf'_* \circ (g')^{-1}
\quad\text{且}\quad
L(g')_! \circ (f')^{-1} = f^{-1} \circ Lg_!
$$
（模的情形和阿贝尔层的情形均成立）。
\end{lemma}

\begin{proof}
令 $\tau = \etale$（相应地，令 $\tau = fppf$）。
设 $\mathcal{F}$ 为 $\mathcal{X}_\tau$ 上的阿贝尔层。
由《叠的上同调》引理
\ref{stacks-cohomology-lemma-lisse-etale-functorial-cohomology}，
典范的（基变换）映射
$$
g^{-1}Rf_*\mathcal{F} \longrightarrow Rf'_*(g')^{-1}\mathcal{F}
$$
是同构。证明的其余部分是形式的。由于阿贝尔群层与模层的上同调一致，
当 $\mathcal{F}$ 是 $\mathcal{X}_\tau$ 上的模层时，我们也得到
$g^{-1} Rf_*\mathcal{F} = Rf'_* (g')^{-1}\mathcal{F}$。

\medskip\noindent
接下来证明：对 $\mathcal{Y}_{lisse,\etale}$（相应地，
$\mathcal{Y}_{flat,fppf}$）上的 $\mathcal{G}$
（模层或阿贝尔群层），典范映射
$$
L(g')_!(f')^{-1}\mathcal{G} \to f^{-1}Lg_!\mathcal{G}
$$
是同构。为此，只需证明：对 $\mathcal{X}_\tau$ 上任意内射层
$\mathcal{I}$，诱导映射
$$
\Hom(L(g')_!(f')^{-1}\mathcal{G}, \mathcal{I}[n])
\leftarrow
\Hom(f^{-1}Lg_!\mathcal{G}, \mathcal{I}[n])
$$
对所有 $n \in \mathbf{Z}$ 都是同构。（这里的 Hom 取自适当的导出范畴。）
由 $f^{-1}$ 与 $Rf_*$ 的伴随性、$Lg_!$ 与 $g^{-1}$ 的伴随性，
以及相应的“加撇”版本，这归结为上面已经证明的同构
$g^{-1} Rf_*\mathcal{I} \to Rf'_* (g')^{-1}\mathcal{I}$。

\medskip\noindent
% P12-ERR-0042: see ../control/ERRATA_CANDIDATES.jsonl
若 $\mathcal{G}^\bullet$ 是
$\mathcal{Y}_{lisse,\etale}$（相应地，$\mathcal{Y}_{fppf}$）
上的有界复形（模复形或阿贝尔群复形），则典范映射
\begin{equation}
\label{stacks-perfect-equation-to-show}
L(g')_!(f')^{-1}\mathcal{G}^\bullet \to f^{-1}Lg_!\mathcal{G}^\bullet
\end{equation}
是同构。这由层的情形、涉及截断的通常论证，以及
$L(g')_!(f')^{-1}$ 和 $f^{-1}Lg_!$ 都是三角范畴的正合函子这一事实得到。

\medskip\noindent
设 $\mathcal{G}^\bullet$ 是
$\mathcal{Y}_{lisse,\etale}$（相应地，$\mathcal{Y}_{fppf}$）
上的上有界复形（模复形或阿贝尔群复形）。
典范映射（\ref{stacks-perfect-equation-to-show}）是同构，因为可用笨截断
$\sigma_{\geq -n}$（参见《同调》第
\ref{homology-section-truncations} 节）将
$\mathcal{G}^\bullet$ 写成有界复形的余极限
$\mathcal{G}^\bullet = \colim \mathcal{G}_n^\bullet$。
由此得到优三角
$$
\bigoplus\nolimits_{n \geq 1} \mathcal{G}_n^\bullet \to
\bigoplus\nolimits_{n \geq 1} \mathcal{G}_n^\bullet \to
\mathcal{G}^\bullet \to \ldots
$$
而函子 $L(g')_!$、$(f')^{-1}$、$f^{-1}$、$Lg_!$
中的每一个都与（复形的）直和交换。

\medskip\noindent
若 $\mathcal{G}^\bullet$ 是
$\mathcal{Y}_{lisse,\etale}$（相应地，$\mathcal{Y}_{fppf}$）
上的任意复形（模复形或阿贝尔群复形），则用典范截断
$\tau_{\leq n}$（参见《同调》第
\ref{homology-section-truncations} 节）将
$\mathcal{G}^\bullet$ 写成上有界复形的余极限，并重复上一段的论证。

\medskip\noindent
最后，由 $f^{-1}$ 与 $Rf_*$ 的伴随性、$Lg_!$ 与 $g^{-1}$ 的伴随性，
以及相应的“加撇”版本，我们得到：在完全一般的情形下，
引理的第一个恒等式由第二个恒等式推出。
\end{proof}
\begin{lemma}
\label{stacks-perfect-lemma-higher-shriek-quasi-coherent}
设 $\mathcal{X}$ 为代数叠。采用《叠的上同调》引理
\ref{stacks-cohomology-lemma-lisse-etale} 中的记号。
\begin{enumerate}
\item 设 $\mathcal{H}$ 为 $\mathcal{X}$ 的光滑-\'etale site 上的拟凝聚
$\mathcal{O}_{\mathcal{X}_{lisse,\etale}}$-模。
对所有 $p \in \mathbf{Z}$，层 $H^p(Lg_!\mathcal{H})$
都是 $\mathcal{X}$ 上具有平坦基变换性质的局部拟凝聚模。
\item 设 $\mathcal{H}$ 为 $\mathcal{X}$ 的平坦-fppf site 上的拟凝聚
$\mathcal{O}_{\mathcal{X}_{flat,fppf}}$-模。
对所有 $p \in \mathbf{Z}$，层 $H^p(Lg_!\mathcal{H})$
都是 $\mathcal{X}$ 上具有平坦基变换性质的局部拟凝聚模。
\end{enumerate}
\end{lemma}

\begin{proof}
选择概形 $U$ 以及满的光滑态射 $x : U \to \mathcal{X}$。
由《site 上的模》定义
\ref{sites-modules-definition-site-local}，存在一个 \'etale
（相应地，fppf）覆盖 $\{U_i \to U\}_{i \in I}$，
使得每个拉回 $f_i^{-1}\mathcal{H}$ 都有全局表示
（参见《site 上的模》定义
\ref{sites-modules-definition-global}）。
这里 $f_i : U_i \to \mathcal{X}$ 是复合
$U_i \to U \to \mathcal{X}$，因而是代数叠之间的态射。
（回忆拉回“就是”在 $\mathcal{X}/f_i$ 上的限制；参见《叠上的层》定义
\ref{stacks-sheaves-definition-pullback} 及其后的讨论。）
加细覆盖后，可设每个 $U_i$ 都是仿射概形。
由于每个 $f_i$ 都是光滑的（相应地，平坦的），由引理
\ref{stacks-perfect-lemma-lisse-etale-functorial-derived} 可知
$f_i^{-1}Lg_!\mathcal{H} = Lg_{i, !}(f'_i)^{-1}\mathcal{H}$。
利用《叠的上同调》引理
\ref{stacks-cohomology-lemma-check-lqc-fbc-on-covering}，
我们将引理的断言约化到如下情形：
$\mathcal{H}$ 有全局表示，并且对某个仿射概形
$X = \Spec(A)$ 有 $\mathcal{X} = (\Sch/X)_{fppf}$。

\medskip\noindent
设我们的表示形如
$$
\bigoplus\nolimits_{j \in J} \mathcal{O} \longrightarrow
\bigoplus\nolimits_{i \in I} \mathcal{O} \longrightarrow
\mathcal{H} \longrightarrow 0
$$
其中 $\mathcal{O} = \mathcal{O}_{\mathcal{X}_{lisse,\etale}}$
（相应地，$\mathcal{O} = \mathcal{O}_{\mathcal{X}_{flat,fppf}}$）。
注意 site $\mathcal{X}_{lisse,\etale}$
（相应地，$\mathcal{X}_{flat,fppf}$）有终对象，即
$X/X$，而且它是拟紧的（参见《site 上的上同调》第
\ref{sites-cohomology-section-limits} 节）。
因此由《Site》引理
\ref{sites-lemma-directed-colimits-sections} 有
$$
\Gamma(\bigoplus\nolimits_{i \in I} \mathcal{O}) =
\bigoplus\nolimits_{i \in I} A
$$
所以表示中的映射对应于 $A$-模 $M$ 的类似表示
$$
\bigoplus\nolimits_{j \in J} A \longrightarrow
\bigoplus\nolimits_{i \in I} A \longrightarrow
M \longrightarrow 0
$$
此外，$\mathcal{H}$ 等于与 $M$ 相伴的拟凝聚层 $M^a$
在光滑-\'etale（相应地，平坦-fppf）site 上的限制。
选择一个由自由 $A$-模组成的分解
$$
\ldots \to F_2 \to F_1 \to F_0 \to M \to 0
$$
复形
$$
\ldots \mathcal{O} \otimes_A F_2 \to \mathcal{O} \otimes_A F_1 \to
\mathcal{O} \otimes_A F_0 \to \mathcal{H} \to 0
$$
是 $\mathcal{H}$ 的自由 $\mathcal{O}$-模分解，因为对
$\mathcal{X}_{lisse,\etale}$（相应地，
$\mathcal{X}_{flat,fppf}$）的每个对象 $U/X$，
结构态射 $U \to X$ 都是平坦的。
因此按构造，$Lg_!\mathcal{H}$ 的值为
$$
\ldots \to
\mathcal{O}_\mathcal{X} \otimes_A F_2 \to
\mathcal{O}_\mathcal{X} \otimes_A F_1 \to
\mathcal{O}_\mathcal{X} \otimes_A F_0 \to 0 \to \ldots
$$
由于这是 $\mathcal{X}_\etale$（相应地，
$\mathcal{X}_{fppf}$）上的拟凝聚模复形，
由《叠的上同调》命题
\ref{stacks-cohomology-proposition-loc-qcoh-flat-base-change}
可知 $H^p(Lg_!\mathcal{H})$ 是拟凝聚的。
\end{proof}




\section{上同调与光滑-\'etale、平坦-fppf site}
\label{stacks-perfect-section-compare-unbounded-lisse-etale}

\noindent
我们已经看到，代数叠 $\mathcal{X}$ 上的层的上同调可在平坦-fppf site
% P12-ERR-0041: see ../control/ERRATA_CANDIDATES.jsonl
上计算。本节证明，对 $\mathcal{X}$ 的直接范畴中（可能）无界的对象，
同样的结论成立。

\begin{lemma}
\label{stacks-perfect-lemma-higher-shriek-Z}
设 $\mathcal{X}$ 为代数叠。无论 $Lg_!$ 取引理
\ref{stacks-perfect-lemma-shriek-derived} 第 (1) 部分中的函子，还是 $Lg_!$ 取引理
\ref{stacks-perfect-lemma-shriek-derived} 第 (3) 部分中的函子，
都有 $Lg_!\mathbf{Z} = \mathbf{Z}$。
\end{lemma}

\begin{proof}
我们先对平坦-fppf site 与 fppf site 的比较作证明；
光滑-\'etale site 的情形完全相同。
我们需要证明：当 $i \not = 0$ 时 $H^i(Lg_!\mathbf{Z})$ 为 $0$，
并且典范映射 $H^0(Lg_!\mathbf{Z}) \to \mathbf{Z}$ 是同构。
设 $f : \mathcal{U} \to \mathcal{X}$ 为满的平坦态射，
其中 $\mathcal{U}$ 是概形，并且 $f$ 还是局部有限表现的。
（例如，选择一个表示 $U \to \mathcal{X}$，并令
$\mathcal{U}$ 为与 $U$ 对应的代数叠。）
由《叠上的层》引理
\ref{stacks-sheaves-lemma-check-exactness-covering} 和
\ref{stacks-sheaves-lemma-surjective-flat-locally-finite-presentation}，
只需证明：当 $i \not = 0$ 时，拉回
$f^{-1}H^i(Lg_!\mathbf{Z})$ 为 $0$，并且
$H^0(Lg_!\mathbf{Z}) \to f^{-1}\mathbf{Z}$ 的拉回是同构。
由引理 \ref{stacks-perfect-lemma-lisse-etale-functorial-derived} 可得
$f^{-1}Lg_!\mathbf{Z} = L(g')_!\mathbf{Z}$，其中
$g' : \Sh(\mathcal{U}_{flat, fppf}) \to \Sh(\mathcal{U}_{fppf})$
是 $\mathcal{U}$ 对应的比较态射。
这将我们约化到下一段所研究的情形。

\medskip\noindent
设对某个概形 $X$ 有 $\mathcal{X} = (\Sch/X)_{fppf}$。
此时范畴 $\mathcal{X}_{flat, fppf}$ 有终对象 $e$，即 $X/X$；
而且函子
$u : \mathcal{X}_{flat, fppf} \to \mathcal{X}_{fppf}$
将 $e$ 映到终对象。由于 $\mathbf{Z}$ 是终对象上的自由阿贝尔层
（只要终对象存在），由《site 上的上同调》引理
\ref{sites-cohomology-lemma-existence-derived-lower-shriek}
中 $Lg_!$ 的构造本身可得 $Lg_!\mathbf{Z} = \mathbf{Z}$。
\end{proof}

\begin{lemma}
\label{stacks-perfect-lemma-lisse-etale-cohomology}
设 $\mathcal{X}$ 为代数叠。采用引理
\ref{stacks-perfect-lemma-shriek-derived} 中的记号。
\begin{enumerate}
\item 对 $D(\mathcal{X}_\etale)$ 中的 $K$，有
\begin{enumerate}
\item $R\Gamma(\mathcal{X}_\etale, K) =
R\Gamma(\mathcal{X}_{lisse,\etale}, g^{-1}K)$，以及
\item 对 $\mathcal{X}_{lisse,\etale}$ 的任意对象 $x$，有
$R\Gamma(x, K) =
R\Gamma(\mathcal{X}_{lisse,\etale}/x, g^{-1}K)$。
\end{enumerate}
\item 对 $D(\mathcal{X}_{fppf})$ 中的 $K$，有
\begin{enumerate}
\item $R\Gamma(\mathcal{X}_{fppf}, K) =
R\Gamma(\mathcal{X}_{flat,fppf}, g^{-1}K)$，以及
\item
% P12-ERR-0043: see ../control/ERRATA_CANDIDATES.jsonl
对 $\mathcal{X}_{flat,fppf}$ 的任意对象 $x$，有
$H^p(x, K) =
R\Gamma(\mathcal{X}_{flat,fppf}/x, g^{-1}K)$。
\end{enumerate}
\end{enumerate}
两种情形下，对模也有同样结论，因为 $g^{-1} = g^*$，
并且由《site 上的上同调》引理
\ref{sites-cohomology-lemma-modules-abelian-unbounded}，
上同调的计算没有区别。
\end{lemma}

\begin{proof}
我们先对平坦-fppf site 与 fppf site 的比较作证明；
光滑-\'etale site 的情形完全相同。
由引理 \ref{stacks-perfect-lemma-higher-shriek-Z}，有
$Lg_!\mathbf{Z} = \mathbf{Z}$。于是得到
\begin{align*}
R\Gamma(\mathcal{X}_{fppf}, K)
& =
R\Hom(\mathbf{Z}, K) \\
& =
R\Hom(Lg_!\mathbf{Z}, K) \\
& =
R\Hom(\mathbf{Z}, g^{-1}K) \\
& =
% P12-ERR-0044: see ../control/ERRATA_CANDIDATES.jsonl
R\Gamma(\mathcal{X}_{lisse,\etale}, g^{-1}K)
\end{align*}
% P12-ERR-0045: see ../control/ERRATA_CANDIDATES.jsonl
这证明了 (1)(a)。
(1)(b) 由 (1)(a) 推出。事实上，若 $x$ 位于概形 $U$ 之上，
则 site $\mathcal{X}_\etale/x$ 等价于
$(\Sch/U)_\etale$，而
% P12-ERR-0046: see ../control/ERRATA_CANDIDATES.jsonl
$\mathcal{X}_{lisse,\etale}$ 等价于 $U_{lisse, \etale}$。
\end{proof}
\section{拟凝聚模的导出范畴}
\label{stacks-perfect-section-derived}

\noindent
设 $\mathcal{X}$ 为代数叠。由于包含函子
$\QCoh(\mathcal{O}_\mathcal{X}) \to
\textit{Mod}(\mathcal{O}_\mathcal{X})$ 不是正合的，
我们不能把 $D_\QCoh(\mathcal{O}_\mathcal{X})$ 定义为
$D(\mathcal{O}_\mathcal{X})$ 中由上同调层拟凝聚的复形组成的充满子范畴。
作为替代，我们仿照《叠的上同调》评注
\ref{stacks-cohomology-remark-bousfield-colocalization}，
将拟凝聚模的导出范畴定义为一个商。

\medskip\noindent
回忆 $\textit{LQCoh}^{fbc}(\mathcal{O}_\mathcal{X}) \subset
\textit{Mod}(\mathcal{O}_\mathcal{X})$ 表示由具有平坦基变换性质的
局部拟凝聚 $\mathcal{O}_\mathcal{X}$-模组成的充满子范畴；
参见《叠的上同调》第
\ref{stacks-cohomology-section-loc-qcoh-flat-base-change} 节。
我们将缩写
$$
D_{\textit{LQCoh}^{fbc}}(\mathcal{O}_\mathcal{X}) =
D_{\textit{LQCoh}^{fbc}(\mathcal{O}_\mathcal{X})}(\mathcal{O}_\mathcal{X})
$$
由《导出范畴》引理
\ref{derived-lemma-cohomology-in-serre-subcategory} 和
《叠的上同调》命题
\ref{stacks-cohomology-proposition-loc-qcoh-flat-base-change} 第 (2) 部分，
可知 $D_{\textit{LQCoh}^{fbc}}(\mathcal{O}_\mathcal{X})$
是 $D(\mathcal{O}_\mathcal{X})$ 的严格充满饱和三角子范畴。

\medskip\noindent
令 $\textit{Parasitic}(\mathcal{O}_\mathcal{X}) \subset
\textit{Mod}(\mathcal{O}_\mathcal{X})$ 表示由寄生
$\mathcal{O}_\mathcal{X}$-模组成的充满子范畴；
参见《叠的上同调》第
\ref{stacks-cohomology-section-parasitic} 节。
我们将缩写
$$
D_{\textit{Parasitic}}(\mathcal{O}_\mathcal{X}) =
D_{\textit{Parasitic}(\mathcal{O}_\mathcal{X})}(\mathcal{O}_\mathcal{X})
$$
与前面一样，这是 $D(\mathcal{O}_\mathcal{X})$
的严格充满饱和三角子范畴，因为
$\textit{Parasitic}(\mathcal{O}_\mathcal{X})$
是 $\textit{Mod}(\mathcal{O}_\mathcal{X})$ 的 Serre 子范畴；
参见《叠的上同调》引理
\ref{stacks-cohomology-lemma-parasitic}。

\medskip\noindent
$\textit{Mod}(\mathcal{O}_\mathcal{X})$ 的弱 Serre 子范畴
$\textit{Parasitic}(\mathcal{O}_\mathcal{X}) \cap
\textit{LQCoh}^{fbc}(\mathcal{O}_\mathcal{X})$
的交仍是弱 Serre 子范畴。
类似地，我们缩写
\begin{align*}
D_{\textit{Parasitic} \cap \textit{LQCoh}^{fbc}}(\mathcal{O}_\mathcal{X})
& =
D_{\textit{Parasitic}(\mathcal{O}_\mathcal{X}) \cap
\textit{LQCoh}^{fbc}(\mathcal{O}_\mathcal{X})}(\mathcal{O}_\mathcal{X}) \\
& =
D_{\textit{Parasitic}}(\mathcal{O}_\mathcal{X})
\cap
D_{\textit{LQCoh}^{fbc}}(\mathcal{O}_\mathcal{X})
\end{align*}
与前面一样，这是 $D(\mathcal{O}_\mathcal{X})$
的严格充满饱和三角子范畴。因而更是
$D_{\textit{LQCoh}^{fbc}}(\mathcal{O}_\mathcal{X})$
的严格充满饱和三角子范畴。

\begin{definition}
\label{stacks-perfect-definition-derived}
设 $\mathcal{X}$ 为代数叠。采用上述记号，我们把
{\it 上同调层拟凝聚的 $\mathcal{O}_\mathcal{X}$-模导出范畴}
定义为 Verdier 商\footnote{此定义不同于文献中的定义；
参见 \cite[6.3]{olsson_sheaves}，但由引理
\ref{stacks-perfect-lemma-derived-quasi-coherent}，它与该定义一致。}
$$
D_\QCoh(\mathcal{O}_\mathcal{X}) =
D_{\textit{LQCoh}^{fbc}}(\mathcal{O}_\mathcal{X})/
D_{\textit{Parasitic} \cap \textit{LQCoh}^{fbc}}(\mathcal{O}_\mathcal{X})
$$
\end{definition}

\noindent
Verdier 商定义见《导出范畴》第
\ref{derived-section-quotients} 节。
$D_{\textit{LQCoh}^{fbc}}(\mathcal{O}_\mathcal{X})$
中的态射 $a : E \to E'$ 在
$D_\QCoh(\mathcal{O}_\mathcal{X})$ 中成为同构，当且仅当锥
$C(a)$ 的上同调层寄生；参见《导出范畴》引理
\ref{derived-lemma-operations}。

\medskip\noindent
考虑函子
$$
D_{\textit{LQCoh}^{fbc}}(\mathcal{O}_\mathcal{X})
\xrightarrow{H^i}
\textit{LQCoh}^{fbc}(\mathcal{O}_\mathcal{X})
\xrightarrow{Q}
\QCoh(\mathcal{O}_\mathcal{X})
$$
注意 $Q$ 零化子范畴
$\textit{Parasitic}(\mathcal{O}_\mathcal{X}) \cap
\textit{LQCoh}^{fbc}(\mathcal{O}_\mathcal{X})$；
参见《叠的上同调》引理
\ref{stacks-cohomology-lemma-adjoint-kernel-parasitic}。
由《导出范畴》引理
\ref{derived-lemma-universal-property-quotient}，
我们得到上同调函子
\begin{equation}
\label{stacks-perfect-equation-Hi-quasi-coherent}
H^i :
D_\QCoh(\mathcal{O}_\mathcal{X})
\longrightarrow
\QCoh(\mathcal{O}_\mathcal{X})
\end{equation}
此外，$E \in D_\QCoh(\mathcal{O}_\mathcal{X})$ 为零，
当且仅当对所有 $i \in \mathbf{Z}$ 都有 $H^i(E) = 0$，
因为由《叠的上同调》引理
\ref{stacks-cohomology-lemma-adjoint-kernel-parasitic}，
$Q$ 的核恰好等于
$\textit{Parasitic}(\mathcal{O}_\mathcal{X}) \cap
\textit{LQCoh}^{fbc}(\mathcal{O}_\mathcal{X})$。

\medskip\noindent
注意范畴
$\textit{Parasitic}(\mathcal{O}_\mathcal{X}) \cap
\textit{LQCoh}^{fbc}(\mathcal{O}_\mathcal{X})$ 和
$\textit{LQCoh}^{fbc}(\mathcal{O}_\mathcal{X})$
也是 \'etale 拓扑中模的阿贝尔范畴
$\textit{Mod}(\mathcal{X}_\etale, \mathcal{O}_\mathcal{X})$
的弱 Serre 子范畴；参见《叠的上同调》命题
\ref{stacks-cohomology-proposition-loc-qcoh-flat-base-change} 和引理
\ref{stacks-cohomology-lemma-parasitic}。
因此下述引理的陈述有意义。

\begin{lemma}
\label{stacks-perfect-lemma-compare-etale-fppf-QCoh}
设 $\mathcal{X}$ 为代数叠。缩写
$\mathcal{P}_\mathcal{X} = \textit{Parasitic}(\mathcal{O}_\mathcal{X}) \cap
\textit{LQCoh}^{fbc}(\mathcal{O}_\mathcal{X})$。
比较态射
$\epsilon : \mathcal{X}_{fppf} \to \mathcal{X}_\etale$
诱导交换图
$$
\xymatrix{
D_{\textit{Parasitic} \cap \textit{LQCoh}^{fbc}}(\mathcal{O}_\mathcal{X})
\ar[r] &
D_{\textit{LQCoh}^{fbc}}(\mathcal{O}_\mathcal{X}) \ar[r] &
D(\mathcal{O}_\mathcal{X}) \\
D_{\mathcal{P}_\mathcal{X}}(\mathcal{X}_\etale, \mathcal{O}_\mathcal{X})
\ar[r] \ar[u]^{\epsilon^*} &
D_{\textit{LQCoh}^{fbc}(\mathcal{O}_\mathcal{X})}(
\mathcal{X}_\etale, \mathcal{O}_\mathcal{X})
\ar[r] \ar[u]^{\epsilon^*} &
D(\mathcal{X}_\etale, \mathcal{O}_\mathcal{X})
\ar[u]^{\epsilon^*}
}
$$
此外，左边两个竖直箭头都是三角范畴的等价，因而还得到等价
$$
D_{\textit{LQCoh}^{fbc}(\mathcal{O}_\mathcal{X})}
(\mathcal{X}_\etale, \mathcal{O}_\mathcal{X})
/
D_{\mathcal{P}_\mathcal{X}}(\mathcal{X}_\etale, \mathcal{O}_\mathcal{X})
\longrightarrow
D_\QCoh(\mathcal{O}_\mathcal{X})
$$
\end{lemma}

\begin{proof}
由于 $\epsilon^*$ 正合，显然可得到引理陈述中的图。
我们将把《site 上的上同调》引理
\ref{sites-cohomology-lemma-compare-topologies-derived-adequate-modules}
应用于下述情形，以证明中间的竖直箭头是等价：
$\mathcal{C} = \mathcal{X}$，
$\tau = fppf$，
$\tau' = \etale$，
$\mathcal{O} = \mathcal{O}_\mathcal{X}$，
$\mathcal{A} = \textit{LQCoh}^{fbc}(\mathcal{O}_\mathcal{X})$，
而 $\mathcal{B}$ 是 $\mathcal{X}$ 中位于仿射概形之上的对象的集合。
要说明该引理适用，需要检验条件 (1)、(2)、(3)、(4)。
条件 (1) 和 (2) 由上面的讨论显然成立
（明确地说，这由《叠的上同调》命题
\ref{stacks-cohomology-proposition-loc-qcoh-flat-base-change} 得出）。
条件 (3) 成立是因为每个概形都有由仿射开集组成的 Zariski 开覆盖。
条件 (4) 由《下降》引理
\ref{descent-lemma-quasi-coherent-and-flat-base-change} 得出。

\medskip\noindent
我们略去下述验证：范畴的等价
$\epsilon^* : 
D_{\textit{LQCoh}^{fbc}(\mathcal{O}_\mathcal{X})}(
\mathcal{X}_\etale, \mathcal{O}_\mathcal{X})
\to
D_{\textit{LQCoh}^{fbc}}(\mathcal{O}_\mathcal{X})$
在上同调层寄生的复形子范畴之间诱导等价。
\end{proof}

\noindent
设 $\mathcal{X}$ 为代数叠。
由《叠的上同调》引理
\ref{stacks-cohomology-lemma-quasi-coherent-weak-serre}，
拟凝聚模范畴
$\QCoh(\mathcal{O}_{\mathcal{X}_{lisse,\etale}})$
构成
$\textit{Mod}(\mathcal{O}_{\mathcal{X}_{lisse,\etale}})$
的弱 Serre 子范畴，而
$\QCoh(\mathcal{O}_{\mathcal{X}_{flat,fppf}})$
构成
$\textit{Mod}(\mathcal{O}_{\mathcal{X}_{flat,fppf}})$
的弱 Serre 子范畴。因此可以考虑
$$
D_\QCoh(\mathcal{O}_{\mathcal{X}_{lisse,\etale}}) =
D_{\QCoh(\mathcal{O}_{\mathcal{X}_{lisse,\etale}})}(
\mathcal{O}_{\mathcal{X}_{lisse,\etale}})
\subset
D(\mathcal{O}_{\mathcal{X}_{lisse,\etale}})
$$
以及类似的
$$
D_\QCoh(\mathcal{O}_{\mathcal{X}_{flat,fppf}}) =
D_{\QCoh(\mathcal{O}_{\mathcal{X}_{flat,fppf}})}(
\mathcal{O}_{\mathcal{X}_{flat,fppf}})
\subset
D(\mathcal{O}_{\mathcal{X}_{flat,fppf}})
$$
与上面一样，它们都是严格充满饱和三角子范畴。
事实表明，$D_\QCoh(\mathcal{O}_\mathcal{X})$
与其中任意一个都等价。

\begin{lemma}
\label{stacks-perfect-lemma-derived-quasi-coherent}
设 $\mathcal{X}$ 为代数叠。令
$\mathcal{P}_\mathcal{X} = \textit{Parasitic}(\mathcal{O}_\mathcal{X}) \cap
\textit{LQCoh}^{fbc}(\mathcal{O}_\mathcal{X})$。
\begin{enumerate}
\item
设 $\mathcal{F}^\bullet$ 是
$D_{\textit{LQCoh}^{fbc}(\mathcal{O}_\mathcal{X})}
(\mathcal{X}_\etale, \mathcal{O}_\mathcal{X})$ 的对象。
对光滑-\'etale site，取《叠的上同调》引理
\ref{stacks-cohomology-lemma-lisse-etale} 中的 $g$，则
\begin{enumerate}
\item $g^*\mathcal{F}^\bullet$ 属于
$D_\QCoh(\mathcal{O}_{\mathcal{X}_{lisse,\etale}})$；
\item $g^*\mathcal{F}^\bullet = 0$，当且仅当
$\mathcal{F}^\bullet$ 属于
$D_{\mathcal{P}_\mathcal{X}}(\mathcal{X}_\etale, \mathcal{O}_\mathcal{X})$；
\item 对
$D_\QCoh(\mathcal{O}_{\mathcal{X}_{lisse,\etale}})$ 中的
$\mathcal{H}^\bullet$，
$Lg_!\mathcal{H}^\bullet$ 属于
$D_{\textit{LQCoh}^{fbc}(\mathcal{O}_\mathcal{X})}(
\mathcal{X}_\etale, \mathcal{O}_\mathcal{X})$；并且
\item 函子 $g^*$ 和 $Lg_!$ 给出互逆函子
$$
\xymatrix{
D_\QCoh(\mathcal{O}_\mathcal{X}) \ar@<1ex>[r]^-{g^*} &
D_\QCoh(\mathcal{O}_{\mathcal{X}_{lisse,\etale}})
\ar@<1ex>[l]^-{Lg_!}
}
$$
\end{enumerate}
\item
设 $\mathcal{F}^\bullet$ 是
$D_{\textit{LQCoh}^{fbc}}(\mathcal{O}_\mathcal{X})$ 的对象。
对平坦-fppf site，取《叠的上同调》引理
\ref{stacks-cohomology-lemma-lisse-etale} 中的 $g$，则
\begin{enumerate}
\item $g^*\mathcal{F}^\bullet$ 属于
$D_\QCoh(\mathcal{O}_{\mathcal{X}_{flat, fppf}})$；
\item $g^*\mathcal{F}^\bullet = 0$，当且仅当
$\mathcal{F}^\bullet$ 属于
$D_{\mathcal{P}_\mathcal{X}}(\mathcal{O}_\mathcal{X})$；
\item 对
$D_\QCoh(\mathcal{O}_{\mathcal{X}_{flat,fppf}})$ 中的
$\mathcal{H}^\bullet$，
$Lg_!\mathcal{H}^\bullet$ 属于
$D_{\textit{LQCoh}^{fbc}}(\mathcal{O}_\mathcal{X})$；并且
\item 函子 $g^*$ 和 $Lg_!$ 给出互逆函子
$$
\xymatrix{
D_\QCoh(\mathcal{O}_\mathcal{X}) \ar@<1ex>[r]^-{g^*} &
D_\QCoh(\mathcal{O}_{\mathcal{X}_{flat,fppf}}) \ar@<1ex>[l]^-{Lg_!}
}
$$
\end{enumerate}
\end{enumerate}
\end{lemma}

\begin{proof}
函子 $g^* = g^{-1}$ 正合，故 (1)(a)、(2)(a)、(1)(b) 和 (2)(b)
由《叠的上同调》引理
\ref{stacks-cohomology-lemma-quasi-coherent} 和
\ref{stacks-cohomology-lemma-parasitic-in-terms-flat-fppf} 得出。

\medskip\noindent
证明 (1)(c) 和 (2)(c)。
引理 \ref{stacks-perfect-lemma-shriek-derived} 中 $Lg_!$ 的构造
（通过《site 上的上同调》引理
\ref{sites-cohomology-lemma-existence-derived-lower-shriek}，
而该引理又使用《导出范畴》命题
\ref{derived-proposition-left-derived-exists}）
表明：对 $D(\mathcal{O}_{\mathcal{X}_{lisse,\etale}})$
中任意对象 $\mathcal{H}^\bullet$，$Lg_!$ 的计算式为
$$
Lg_!\mathcal{H}^\bullet = \colim g_!\mathcal{K}_n^\bullet =
g_! \colim \mathcal{K}_n^\bullet
$$
（逐项取余极限），其中拟同构
$\colim \mathcal{K}_n^\bullet \to \mathcal{H}^\bullet$
诱导拟同构
$\mathcal{K}_n^\bullet \to \tau_{\leq n} \mathcal{H}^\bullet$。
由于包含函子
$$
\textit{LQCoh}^{fbc}(\mathcal{O}_\mathcal{X}) \subset
\textit{Mod}(\mathcal{X}_\etale, \mathcal{O}_\mathcal{X})
\quad\text{且}\quad
\textit{LQCoh}^{fbc}(\mathcal{O}_\mathcal{X}) \subset
\textit{Mod}(\mathcal{O}_\mathcal{X})
$$
与滤过余极限相容，只需对
$D_\QCoh(\mathcal{O}_{\mathcal{X}_{lisse,\etale}})$ 和
$D_\QCoh(\mathcal{O}_{\mathcal{X}_{flat,fppf}})$ 中的上有界复形
$\mathcal{H}^\bullet$ 证明 (c)。
在此情形下，为证明 $H^n(Lg_!\mathcal{H}^\bullet)$ 属于
$\textit{LQCoh}^{fbc}(\mathcal{O}_\mathcal{X})$，
可对满足 $\mathcal{H}^i = 0$（当 $i > m$ 时）的整数 $m$ 作归纳。
若 $m < n$，则 $H^n(Lg_!\mathcal{H}^\bullet) = 0$，结论成立。
一般地，考虑优三角
$$
\tau_{\leq m - 1}\mathcal{H}^\bullet \to \mathcal{H}^\bullet \to
H^m(\mathcal{H}^\bullet)[-m] \to \ldots
$$
（《导出范畴》评注
\ref{derived-remark-truncation-distinguished-triangle}），
并作用函子 $Lg_!$。由于
$\textit{LQCoh}^{fbc}(\mathcal{O}_\mathcal{X})$
是模范畴的弱 Serre 子范畴，只需对三项中的两项证明 (c)。
对 $Lg_!\tau_{\leq m - 1}\mathcal{H}^\bullet$，
结论由归纳假设成立；对
$Lg_!H^m(\mathcal{H}^\bullet)[-m]$，
结论由引理 \ref{stacks-perfect-lemma-higher-shriek-quasi-coherent} 成立。
所以 (c) 成立。

\medskip\noindent
我们证明 (2)(d)。由 (2)(a) 和 (2)(b)，函子
$g^{-1} = g^*$ 诱导函子
$$
c :
D_\QCoh(\mathcal{O}_\mathcal{X})
\longrightarrow
D_\QCoh(\mathcal{O}_{\mathcal{X}_{flat, fppf}})
$$
参见《导出范畴》引理
\ref{derived-lemma-universal-property-quotient}。
因此有下述三角范畴的图
$$
\xymatrix{
D_{\textit{LQCoh}^{fbc}}(\mathcal{O}_\mathcal{X})
\ar[rd]^{g^{-1}} \ar[rr]_q & &
D_\QCoh(\mathcal{O}_\mathcal{X}) \ar[ld]^c \\
& D_\QCoh(\mathcal{O}_{\mathcal{X}_{flat, fppf}})
\ar@<1ex>[lu]^{Lg_!}
}
$$
其中 $q$ 为商函子，内三角交换，并且
$g^{-1}Lg_! = \text{id}$。
% P12-ERR-0047: see ../control/ERRATA_CANDIDATES.jsonl
对 $D_{\textit{LQCoh}^{fbc}}(\mathcal{O}_\mathcal{X})$
中 $E$ 的任意对象，映射 $a : Lg_!g^{-1}E \to E$
在 $D(\mathcal{O}_{\mathcal{X}_{flat, fppf}})$ 中映为拟同构。
因此 $a$ 的锥在 $g^{-1}$ 下映为零，由 (2)(b) 可知
$q(a)$ 是同构。所以 $q \circ Lg_!$ 是 $c$ 的拟逆。

\medskip\noindent
在光滑-\'etale site 的情形下，与上面完全相同的论证证明
$$
D_{\textit{LQCoh}^{fbc}(\mathcal{O}_\mathcal{X})}(
\mathcal{X}_\etale, \mathcal{O}_\mathcal{X})
/
D_{\mathcal{P}_\mathcal{X}}(
\mathcal{X}_\etale, \mathcal{O}_\mathcal{X})
$$
等价于
$D_\QCoh(\mathcal{O}_{\mathcal{X}_{lisse,\etale}})$。
应用引理 \ref{stacks-perfect-lemma-compare-etale-fppf-QCoh}
中的最后一个等价即可完成证明。
\end{proof}

\noindent
下述引理告诉我们，商函子
$D_{\textit{LQCoh}^{fbc}}(\mathcal{O}_\mathcal{X}) \to
D_\QCoh(\mathcal{O}_\mathcal{X})$
有左伴随。参见评注 \ref{stacks-perfect-remark-QCoh-admissible}。

\begin{lemma}
\label{stacks-perfect-lemma-bousfield-colocalization}
设 $\mathcal{X}$ 为代数叠。
设 $E$ 是 $D_{\textit{LQCoh}^{fbc}}(\mathcal{O}_\mathcal{X})$ 的对象。
在 $D_{\textit{LQCoh}^{fbc}}(\mathcal{O}_\mathcal{X})$ 中
存在典范优三角
$$
E' \to E \to P \to E'[1]
$$
使得 $P$ 属于
$D_{\textit{Parasitic} \cap \textit{LQCoh}^{fbc}}
(\mathcal{O}_\mathcal{X})$，并且对所有
$D_{\textit{Parasitic} \cap \textit{LQCoh}^{fbc}}(\mathcal{O}_\mathcal{X})$
中的 $P'$，都有
$$
\Hom_{D(\mathcal{O}_\mathcal{X})}(E', P') = 0
$$
\end{lemma}

\begin{proof}
考虑《叠的上同调》第
\ref{stacks-cohomology-section-lisse-etale} 节所研究的赋环拓扑斯态射
$g : \Sh(\mathcal{X}_{flat, fppf}) \longrightarrow \Sh(\mathcal{X}_{fppf})$。
令 $E' = Lg_!g^*E$，并令 $P$ 为伴随映射
$E' \to E$ 的锥；参见引理
\ref{stacks-perfect-lemma-shriek-derived} 第 (4) 部分。
由引理 \ref{stacks-perfect-lemma-derived-quasi-coherent} 的 (2)(a) 和 (2)(c)，
$E'$ 属于
$D_{\textit{LQCoh}^{fbc}}(\mathcal{O}_\mathcal{X})$。
所以 $P$ 也属于
$D_{\textit{LQCoh}^{fbc}}(\mathcal{O}_\mathcal{X})$。
由引理 \ref{stacks-perfect-lemma-shriek-derived} 第 (4) 部分，
$g^*Lg_! = \text{id}$，故映射 $g^*E' \to g^*E$ 是同构。
因此 $g^*P = 0$，进而由引理
\ref{stacks-perfect-lemma-derived-quasi-coherent} 的 (2)(b)，$P$ 是
$D_{\textit{Parasitic} \cap \textit{LQCoh}^{fbc}}(\mathcal{O}_\mathcal{X})$
的对象。
最后，对
$D_{\textit{Parasitic} \cap \textit{LQCoh}^{fbc}}(\mathcal{O}_\mathcal{X})$
中的 $P'$，有
$$
\Hom(E', P') = \Hom(Lg_!g^*E, P') = \Hom(g^*E, g^*P') = 0
$$
因为由引理 \ref{stacks-perfect-lemma-derived-quasi-coherent} 的 (2)(b)，
$g^*P' = 0$。
优三角 $E' \to E \to P \to E'[1]$ 是典范的
% P12-ERR-0048: see ../control/ERRATA_CANDIDATES.jsonl
（更精确地说，直到诱导 $E$ 上恒等映射的三角同构为止是唯一的）；
这由《导出范畴》第
\ref{derived-section-admissible} 节中的讨论得到。
\end{proof}

\begin{remark}
\label{stacks-perfect-remark-QCoh-admissible}
引理 \ref{stacks-perfect-lemma-bousfield-colocalization} 的结论告诉我们，
$$
D_{\textit{Parasitic} \cap \textit{LQCoh}^{fbc}}(\mathcal{O}_\mathcal{X})
\subset
D_{\textit{LQCoh}^{fbc}}(\mathcal{O}_\mathcal{X})
$$
是左可容许子范畴；参见《导出范畴》第
\ref{derived-section-admissible} 节。
特别地，若
$\mathcal{A} \subset D_{\textit{LQCoh}^{fbc}}(\mathcal{O}_\mathcal{X})$
表示它的左正交，则《导出范畴》命题
\ref{derived-proposition-summarize-admissible} 推出：
$\mathcal{A}$ 是
$D_{\textit{LQCoh}^{fbc}}(\mathcal{O}_\mathcal{X})$
中的右可容许子范畴，并且复合
$$
\mathcal{A} \longrightarrow
D_{\textit{LQCoh}^{fbc}}(\mathcal{O}_\mathcal{X}) \longrightarrow
D_\QCoh(\mathcal{O}_\mathcal{X})
$$
是等价。这意味着，可将
$D_\QCoh(\mathcal{O}_\mathcal{X})$
视为 $D_{\textit{LQCoh}^{fbc}}(\mathcal{O}_\mathcal{X})$
的严格充满饱和三角子范畴，也可视为
$D(\mathcal{X}_{fppf}, \mathcal{O}_\mathcal{X})$
的严格充满饱和三角子范畴。
\end{remark}





\section{拟凝聚模的导出直像}
\label{stacks-perfect-section-derived-pushforward}

\noindent
作为上面材料的第一个应用，我们构造导出直像。
读者可在《例》第
\ref{examples-section-derived-push-quasi-coherent} 节找到代数叠之间的
拟紧拟分离态射 $f : \mathcal{X} \to \mathcal{Y}$ 的一个例子，
其中直像函子 $Rf_*$ 并不诱导函子
$D_\QCoh(\mathcal{O}_\mathcal{X}) \to
D_\QCoh(\mathcal{O}_\mathcal{Y})$。
因此，有必要限制到下有界复形。

\begin{proposition}
\label{stacks-perfect-proposition-derived-direct-image-quasi-coherent}
设 $f : \mathcal{X} \to \mathcal{Y}$ 为代数叠之间的拟紧拟分离态射。
函子 $Rf_*$ 诱导交换图
$$
\xymatrix{
D^{+}_{\textit{Parasitic} \cap \textit{LQCoh}^{fbc}}(\mathcal{O}_\mathcal{X})
\ar[r] \ar[d]^{Rf_*} &
D^{+}_{\textit{LQCoh}^{fbc}}(\mathcal{O}_\mathcal{X})
\ar[r] \ar[d]^{Rf_*} &
D(\mathcal{O}_\mathcal{X})
\ar[d]^{Rf_*} \\
D^{+}_{\textit{Parasitic} \cap \textit{LQCoh}^{fbc}}(\mathcal{O}_\mathcal{Y})
\ar[r] &
D^{+}_{\textit{LQCoh}^{fbc}}(\mathcal{O}_\mathcal{Y}) \ar[r] &
D(\mathcal{O}_\mathcal{Y})
}
$$
从而在商范畴上诱导函子
$$
Rf_{\QCoh, *} :
D^{+}_\QCoh(\mathcal{O}_\mathcal{X})
\longrightarrow
D^{+}_\QCoh(\mathcal{O}_\mathcal{Y})
$$
此外，《叠的上同调》命题
\ref{stacks-cohomology-proposition-direct-image-quasi-coherent} 中的函子
% P12-ERR-0049: see ../control/ERRATA_CANDIDATES.jsonl
$R^if_\QCoh$ 等于 $H^i \circ Rf_{\QCoh, *}$，
其中 $H^i$ 如（\ref{stacks-perfect-equation-Hi-quasi-coherent}）所示。
\end{proposition}

\begin{proof}
需要证明：对
$D^{+}_{\textit{LQCoh}^{fbc}}(\mathcal{O}_\mathcal{X})$
中的 $E$，$Rf_*E$ 是
$D^{+}_{\textit{LQCoh}^{fbc}}(\mathcal{O}_\mathcal{Y})$ 的对象。
这由《叠的上同调》命题
\ref{stacks-cohomology-proposition-loc-qcoh-flat-base-change}
以及谱序列
$R^if_*H^j(E) \Rightarrow R^{i + j}f_*E$ 得出。
对寄生模的情形，用《叠的上同调》引理
\ref{stacks-cohomology-lemma-pushforward-parasitic}
作同样的论证。最后一个断言由
（\ref{stacks-perfect-equation-Hi-quasi-coherent}）中 $H^i$ 的定义显然成立。
\end{proof}




\section{拟凝聚模的导出拉回}
\label{stacks-perfect-section-derived-pullback}

\noindent
具有拟凝聚上同调层的复形的导出拉回总是存在。

\begin{proposition}
\label{stacks-perfect-proposition-derived-pullback-quasi-coherent}
设 $f : \mathcal{X} \to \mathcal{Y}$ 为代数叠之间的态射。
正合函子 $f^*$ 诱导交换图
$$
\xymatrix{
D_{\textit{LQCoh}^{fbc}}(\mathcal{O}_\mathcal{X}) \ar[r] &
D(\mathcal{O}_\mathcal{X}) \\
D_{\textit{LQCoh}^{fbc}}(\mathcal{O}_\mathcal{Y})
\ar[r] \ar[u]^{f^*} &
D(\mathcal{O}_\mathcal{Y}) \ar[u]^{f^*}
}
$$
复合
$$
D_{\textit{LQCoh}^{fbc}}(\mathcal{O}_\mathcal{Y})
\xrightarrow{f^*}
D_{\textit{LQCoh}^{fbc}}(\mathcal{O}_\mathcal{X})
\xrightarrow{q_\mathcal{X}}
D_\QCoh(\mathcal{O}_\mathcal{X})
$$
相对于局部化
$D_{\textit{LQCoh}^{fbc}}(\mathcal{O}_\mathcal{Y}) \to
D_\QCoh(\mathcal{O}_\mathcal{Y})$
是左可导的，并可把 $Lf^*_\QCoh$ 定义为它的左导出函子
$$
Lf_\QCoh^* :
D_\QCoh(\mathcal{O}_\mathcal{Y})
\longrightarrow
D_\QCoh(\mathcal{O}_\mathcal{X})
$$
（参见《导出范畴》定义
\ref{derived-definition-right-derived-functor-defined} 和
\ref{derived-definition-everywhere-defined}）。
若 $f$ 拟紧且拟分离，则 $Lf^*_\QCoh$ 和 $Rf_{\QCoh, *}$
满足下述伴随关系：
$$
\Hom_{D_\QCoh(\mathcal{O}_\mathcal{X})}(Lf^*_\QCoh A, B)
=
\Hom_{D_\QCoh(\mathcal{O}_\mathcal{Y})}(A, Rf_{\QCoh, *}B)
$$
其中 $A \in D_\QCoh(\mathcal{O}_\mathcal{Y})$ 且
$B \in D^{+}_\QCoh(\mathcal{O}_\mathcal{X})$。
\end{proposition}

\begin{proof}
为证明第一个断言，需要说明：对
$D_{\textit{LQCoh}^{fbc}}(\mathcal{O}_\mathcal{Y})$ 中的 $E$，
$f^*E$ 是
$D_{\textit{LQCoh}^{fbc}}(\mathcal{O}_\mathcal{X})$ 的对象。
由于 $f^* = f^{-1}$ 正合，而由《叠的上同调》命题
\ref{stacks-cohomology-proposition-loc-qcoh-flat-base-change}，
$f^*$ 将 $\textit{LQCoh}^{fbc}(\mathcal{O}_\mathcal{Y})$
映入 $\textit{LQCoh}^{fbc}(\mathcal{O}_\mathcal{X})$，
所以结论立即成立。

\medskip\noindent
令 $\mathcal{D} = D_{\textit{LQCoh}^{fbc}}(\mathcal{O}_\mathcal{Y})$。
令 $S$ 为 $\mathcal{D}$ 中由下述态射组成的类：
其锥是
$D_{\textit{Parasitic} \cap \textit{LQCoh}^{fbc}}(\mathcal{O}_\mathcal{Y})$
的对象。
令 $\mathcal{D}' = D_\QCoh(\mathcal{O}_\mathcal{X})$。
令 $F = q_\mathcal{X} \circ f^* : \mathcal{D} \to \mathcal{D}'$。
那么 $\mathcal{D}, S, \mathcal{D}', F$ 如《导出范畴》情形
\ref{derived-situation-derived-functor} 和定义
\ref{derived-definition-right-derived-functor-defined} 所述。
我们证明：对 $\mathcal{D}$ 的任意对象 $E$，$LF(E)$ 都有定义。
事实上，考虑引理 \ref{stacks-perfect-lemma-bousfield-colocalization} 中构造的三角
$$
E' \to E \to P \to E'[1]
$$
注意 $s : E' \to E$ 是 $S$ 的元素。我们断言 $E'$ 计算 $LF$。
事实上，设 $s' : E'' \to E$ 是 $S$ 的另一个元素，
也就是说，它嵌入三角
$E'' \to E \to P' \to E''[1]$，其中 $P'$ 属于
$D_{\textit{Parasitic} \cap \textit{LQCoh}^{fbc}}(\mathcal{O}_\mathcal{Y})$。
由引理 \ref{stacks-perfect-lemma-bousfield-colocalization}（及其证明），
$E' \to E$ 通过 $E'' \to E$ 分解。
因此 $E' \to E$ 在系统 $S/E$ 中是共终的。
于是显然 $E'$ 计算 $LF$。

\medskip\noindent
为证明最后一个断言，写
$B = q_\mathcal{X}(H)$ 和 $A = q_\mathcal{Y}(E)$。
如上选择 $E' \to E$。
一方面我们将使用
$Rf_{\QCoh, *}(B) = q_\mathcal{Y}(Rf_*H)$，
另一方面使用
$Lf^*_\QCoh(A) = q_\mathcal{X}(f^*E')$。
\begin{align*}
\Hom_{D_\QCoh(\mathcal{O}_\mathcal{X})}(Lf^*_\QCoh A, B)
& = 
\Hom_{D_\QCoh(\mathcal{O}_\mathcal{X})}(q_\mathcal{X}(f^*E'),
q_\mathcal{X}(H)) \\
& = 
\colim_{H \to H'} \Hom_{D(\mathcal{O}_\mathcal{X})}(f^*E', H') \\
& = \colim_{H \to H'} \Hom_{D(\mathcal{O}_\mathcal{Y})}(E', Rf_*H') \\
& = \Hom_{D(\mathcal{O}_\mathcal{Y})}(E', Rf_*H) \\
& =
\Hom_{D_\QCoh(\mathcal{O}_\mathcal{Y})}(A, Rf_{\QCoh, *}B)
\end{align*}
这里的余极限取遍
$D^+_{\textit{LQCoh}^{fbc}}(\mathcal{O}_\mathcal{X})$
中满足下述条件的态射 $s : H \to H'$：
其锥 $P(s)$ 是
$D^+_{\textit{Parasitic} \cap \textit{LQCoh}^{fbc}}(\mathcal{O}_\mathcal{X})$
的对象。
第一个等式已经在上面看到。
第二个等式由 Verdier 商的构造成立。
第三个等式由《site 上的上同调》引理
\ref{sites-cohomology-lemma-adjoint} 成立。
由命题 \ref{stacks-perfect-proposition-derived-direct-image-quasi-coherent}，
$Rf_*P(s)$ 是
$D^+_{\textit{Parasitic} \cap \textit{LQCoh}^{fbc}}(\mathcal{O}_\mathcal{Y})$
的对象，因而
$\Hom_{D(\mathcal{O}_\mathcal{Y})}(E', Rf_*P(s)) = 0$。
所以第四个等式成立。最后一个等式由 $E'$ 的构造成立。
\end{proof}








\section{导出范畴中的拟凝聚对象}
\label{stacks-perfect-section-QC}

\noindent
本节接续《叠上的层》第
\ref{stacks-sheaves-section-QC} 节。
设 $\mathcal{X}$ 为代数叠。在该节中，我们定义了三角范畴
$$
\mathit{QC}(\mathcal{X}) = \mathit{QC}(\mathcal{X}_{affine}, \mathcal{O})
$$
并证明：若 $\mathcal{X}$ 可由代数空间 $X$ 表示，则
$\mathit{QC}(\mathcal{X})$ 等价于
$D_\QCoh(\mathcal{O}_X)$。
事实表明，我们已经发展了恰好足够的理论来证明：
对任意代数叠，同一结论都成立。

\begin{lemma}
\label{stacks-perfect-lemma-cohomology-parasitic}
设 $\mathcal{X}$ 为代数叠。
设 $K$ 是 $D(\mathcal{X}_{fppf})$ 的对象，且其上同调层寄生。
那么，对 $\mathcal{X}$ 中位于概形 $U$ 之上的每个对象 $x$，
只要 $U \to \mathcal{X}$ 平坦，就有 $R\Gamma(x, K) = 0$。
\end{lemma}

\begin{proof}
以
$g : \Sh(\mathcal{X}_{flat, fppf}) \to \Sh(\mathcal{X}_{fppf})$
表示第 \ref{stacks-perfect-section-lisse-etale} 节讨论的拓扑斯态射。
设 $x$ 为 $\mathcal{X}$ 中位于概形 $U$ 之上的对象，
且 $U \to \mathcal{X}$ 平坦；也就是说，
$x$ 是 $\mathcal{X}_{flat, fppf}$ 的对象。
由引理 \ref{stacks-perfect-lemma-lisse-etale-cohomology} 的 (2)(b)，有
$R\Gamma(x, K) = R\Gamma(\mathcal{X}_{flat, fppf}/x, g^{-1}K)$。
然而，我们的假设意味着
$D(\mathcal{X}_{flat, fppf})$ 的对象 $g^{-1}K$ 的上同调层为零；
参见《叠的上同调》定义
\ref{stacks-cohomology-definition-parasitic}。
所以 $g^{-1}K = 0$，结论成立。
\end{proof}

\begin{lemma}
\label{stacks-perfect-lemma-cohomology-parasitic-complex}
设 $\mathcal{X}$ 为代数叠。
设 $K$ 是 $D(\mathcal{X}_{fppf})$ 的对象，并且对
$\mathcal{X}$ 中位于仿射概形 $U$ 之上的所有对象 $x$，
只要 $U \to \mathcal{X}$ 平坦，就有 $R\Gamma(x, K) = 0$。
那么对所有 $i$ 都有 $H^i(\mathcal{X}, K) = 0$。
\end{lemma}

\begin{proof}
以
$g : \Sh(\mathcal{X}_{flat, fppf}) \to \Sh(\mathcal{X}_{fppf})$
表示第 \ref{stacks-perfect-section-lisse-etale} 节讨论的拓扑斯态射。
由引理 \ref{stacks-perfect-lemma-lisse-etale-cohomology} 的 (2)(b)，
我们的假设意味着：在 $\mathcal{X}_{flat, fppf}$
中位于仿射概形之上的每个对象上，$g^{-1}K$ 的上同调都为零。
由于 $\mathcal{X}_{flat, fppf}$ 的每个对象 $x$ 都有由此类对象组成的覆盖，
可知 $g^{-1}K$ 的上同调层为零，即 $g^{-1}K = 0$。
于是当然有
$R\Gamma(\mathcal{X}_{flat, fppf}, g^{-1}K) = 0$，
再由引理 \ref{stacks-perfect-lemma-lisse-etale-cohomology} 的 (2)(a)，
这推出所需结论。
\end{proof}

\begin{lemma}
\label{stacks-perfect-lemma-higher-shriek-QC}
设 $\mathcal{X}$ 为代数叠。
设 $K$ 是
$D_\QCoh(\mathcal{O}_{\mathcal{X}_{flat, fppf}})$ 的对象。
那么 $Lg_!K$ 满足下述性质：对
$\mathcal{X}_{affine}$ 的任意态射 $x \to x'$，映射
$$
R\Gamma(x', Lg_!K) \otimes_{\mathcal{O}(x')}^\mathbf{L} \mathcal{O}(x)
\longrightarrow
R\Gamma(x, Lg_!K)
$$
是拟同构。
\end{lemma}

\begin{proof}
由引理 \ref{stacks-perfect-lemma-derived-quasi-coherent} 的 (2)(c)，对象
$Lg_!K$ 属于
$D_{\textit{LQCoh}^{fbc}}(\mathcal{O}_\mathcal{X})$。
由此很容易看出，若
$\mathcal{O}(x') \to \mathcal{O}(x)$ 是平坦环同态，
则引理中显示的映射是同构；细节从略。

\medskip\noindent
本段说明该问题对 \'etale 拓扑是局部的。
设 $x \to x'$ 是 $\mathcal{X}_{affine}$ 的一般态射。
设 $\{x'_i \to x'\}$ 是
$\mathcal{X}_{affine, \etale}$ 中的覆盖。
令 $x_i = x \times_{x'} x'_i$，则
$\{x_i \to x\}$ 也是
$\mathcal{X}_{affine, \etale}$ 的覆盖。
此时
$\mathcal{O}(x') \to \prod \mathcal{O}(x'_i)$
是忠实平坦 \'etale 环同态，并且
$$
\prod \mathcal{O}(x_i) =
\mathcal{O}(x) \otimes_{\mathcal{O}(x')} \left(\prod \mathcal{O}(x'_i)\right)
$$
因此，一个我们略去的简单代数论证表明：
只需对 $\mathcal{X}_{affine}$ 中的每个态射
$x_i \to x'_i$ 证明引理陈述中的结论。
换言之，该问题在 \'etale 拓扑中是局部的。

\medskip\noindent
选择概形 $X$ 以及满的光滑态射
$f : X \to \mathcal{X}$。
可以（按照我们的语言滥用）把 $f$ 视为 $\mathcal{X}$ 的对象；
于是 $(\Sch/X)_{fppf} = \mathcal{X}/f$，参见《叠上的层》第
\ref{stacks-sheaves-section-restriction} 节。
例如，由《叠上的层》引理
\ref{stacks-sheaves-lemma-surjective-flat-locally-finite-presentation}，
存在 \'etale 覆盖 $\{x'_i \to x'\}$，
使得 $x'_i : U'_i = p(x'_i) \to \mathcal{X}$ 通过 $f$ 分解。
由上一段的结果，可设 $x \to x'$ 是
$(\textit{Aff}/X)_{fppf}$ 中某个态射 $U \to U'$ 在函子
$(\Sch/X)_{fppf} \to \mathcal{X}$ 下的像。
% P12-ERR-0050: see ../control/ERRATA_CANDIDATES.jsonl
此时我们看出，$Lg_!K$ 在 $(\Sch/X)_{fppf}$ 上的限制等于
$f^*Lg_!K = L(g')_!(f')^*K$；这由引理
\ref{stacks-perfect-lemma-lisse-etale-functorial-derived} 得出。
这将我们约化到下一段讨论的情形。

\medskip\noindent
设 $\mathcal{X} = (\Sch/X)_{fppf}$，且 $x \to x'$
对应于仿射概形态射 $U \to U'$。
我们仍可在 $U'$ 上 \'etale（或 Zariski）局部地工作，
因而可设 $U' \to X$ 通过 $X$ 的某个仿射开集分解。
这将我们约化到下一段讨论的情形。

\medskip\noindent
设 $\mathcal{X} = (\Sch/X)_{fppf}$，其中
$X = \Spec(R)$ 是仿射概形，并且 $x \to x'$
对应于仿射概形态射 $U \to U'$。
令 $M^\bullet$ 为表示 $R\Gamma(X, K)$ 的 $R$-模复形。
由《代数进阶》引理
\ref{more-algebra-lemma-K-flat-resolution} 中的构造，
可设 $M^\bullet = \colim P_n^\bullet$，
其中每个 $P_n^\bullet$ 都是自由 $R$-模的上有界复形。
细节从略；另见《代数进阶》评注
\ref{more-algebra-remark-P-resolution}。
考虑 $X_{flat, fppf} = (\Sch/X)_{flat, fppf}$
上由规则
$$
U \longmapsto \Gamma(U, M^\bullet \otimes_R \mathcal{O}_U)
$$
给出的模复形 $M^\bullet_{flat, fppf}$。
由《下降》第
\ref{descent-section-quasi-coherent-sheaves} 节中的讨论，
这是层的复形。存在典范映射
$M^\bullet_{flat, fppf} \to K$；
由证明开头的评注，该映射在
$X_{flat, fppf}$ 的仿射对象上的截面上给出同构。
由于 $X_{flat, fppf}$ 的每个对象都有由仿射对象组成的覆盖，
可知 $M^\bullet_{flat, fppf}$ 与 $K$ 一致。

\medskip\noindent
令 $M^\bullet_{fppf}$ 为 $X_{fppf}$ 上由上述同一公式给出的模复形。
回忆 $Lg_!\mathcal{O} = g_!\mathcal{O} = \mathcal{O}$。
由于 $Lg_!$ 是 $g_!$ 的左导出函子，可得
$Lg_!P_{n, flat, fppf}^\bullet = P_{n, fppf}^\bullet$。
由于函子 $Lg_!$ 与同伦余极限交换
（也可使用它在《site 上的上同调》引理
\ref{sites-cohomology-lemma-existence-derived-lower-shriek} 中的构造），
并且 $M^\bullet = \colim P_n^\bullet$，可得
$Lg_!M^\bullet_{flat, fppf} = M^\bullet_{fppf}$。
设 $U = \Spec(A)$、$U' = \Spec(A')$，
并且 $U \to U'$ 对应于环同态 $A' \to A$。
由上可知
$$
R\Gamma(U, Lg_!K) = M^\bullet \otimes_R A
\quad\text{且}\quad
R\Gamma(U', Lg_!K) = M^\bullet \otimes_R A'
$$
由于 $M^\bullet$ 是 $R$-模的 K-平坦复形，
由张量积的传递性，
$$
R\Gamma(U', Lg_!K) \otimes_{A'}^\mathbf{L} A
\longrightarrow
R\Gamma(U, Lg_!K)
$$
是所需的拟同构。
\end{proof}

\begin{proposition}
\label{stacks-perfect-proposition-QC-compare}
设 $\mathcal{X}$ 为代数叠。那么
$\mathit{QC}(\mathcal{X})$ 典范等价于
$D_\QCoh(\mathcal{O}_\mathcal{X})$。
\end{proposition}

\begin{proof}
由《叠上的层》引理
\ref{stacks-sheaves-lemma-QC-compare-fppf}，
沿比较态射
$\epsilon : \mathcal{X}_{affine, fppf} \to \mathcal{X}_{affine}$
的拉回将 $\mathit{QC}(\mathcal{X})$ 与充满子范畴
$Q_\mathcal{X} \subset D(\mathcal{X}_{affine, fppf}, \mathcal{O})$
等同。利用《叠上的层》公式
（\ref{stacks-sheaves-equation-alternative-ringed}）中的赋环拓扑斯等价，
我们可以并将会把 $Q_\mathcal{X}$ 视为
$D(\mathcal{X}_{fppf}, \mathcal{O})$ 的充满子范畴。

\medskip\noindent
类似地，由引理 \ref{stacks-perfect-lemma-bousfield-colocalization} 和评注
\ref{stacks-perfect-remark-QCoh-admissible}，可将
$D_\QCoh(\mathcal{O}_\mathcal{X})$ 视为左可容许子范畴
$D_{\textit{Parasitic} \cap \textit{LQCoh}^{fbc}}(\mathcal{O}_\mathcal{X})$
在 $D_{\textit{LQCoh}^{fbc}}(\mathcal{O}_\mathcal{X})$
中的左正交 $\mathcal{A}$。

\medskip\noindent
为完成证明，我们将说明：作为
$D(\mathcal{X}_{fppf}, \mathcal{O})$ 的子范畴，
$Q_\mathcal{X}$ 等于 $\mathcal{A}$。

\medskip\noindent
步骤 1：$Q_\mathcal{X}$ 包含于
$D_{\textit{LQCoh}^{fbc}}(\mathcal{O}_\mathcal{X})$。
$Q_\mathcal{X}$ 的对象 $K$ 由下述性质刻画：
将 $K$ 视为 $D(\mathcal{X}_{affine, fppf}, \mathcal{O})$ 的对象时，
$R\epsilon_*K$ 是
$\mathit{QC}(\mathcal{X}_{affine}, \mathcal{O})$ 的对象。
这又恰好意味着：对 $\mathcal{X}_{affine}$ 的所有态射
$x \to x'$，映射
$$
R\Gamma(x', K) \otimes_{\mathcal{O}(x')}^\mathbf{L} \mathcal{O}(x)
\longrightarrow
R\Gamma(x, K)
$$
是同构；参见《site 上的上同调》引理
\ref{sites-cohomology-lemma-cartesion-plus-topology}
陈述中的脚注。
% P12-ERR-0051: see ../control/ERRATA_CANDIDATES.jsonl
现在，若 $x' \to x$ 位于仿射概形之间的某个平坦态射之上，
则这意味着
$$
H^i(x', K) \otimes_{\mathcal{O}(x')} \mathcal{O}(x)
\cong
H^i(x, K)
$$
这清楚地意味着 $H^i(K)$ 是 \'etale 拓扑的层
（《叠上的层》引理
\ref{stacks-sheaves-lemma-quasi-coherent-alternative}），
并且具有平坦基变换性质（略去一个小细节）。

\medskip\noindent
步骤 2：$Q_\mathcal{X}$ 包含于 $\mathcal{A}$。
为证明这一点，只需说明：对 $Q_\mathcal{X}$ 中的 $K$
以及
$D_{\textit{Parasitic} \cap \textit{LQCoh}^{fbc}}(\mathcal{O}_\mathcal{X})$
中的所有 $P$，都有 $\Hom(K, P) = 0$。
考虑对象
$$
H = R\SheafHom_{\mathcal{O}_\mathcal{X}}(K, P)
$$
设 $x$ 为 $\mathcal{X}$ 中位于仿射概形
$U = p(x)$ 之上的对象。
由《site 上的上同调》引理
\ref{sites-cohomology-lemma-section-RHom-over-U}，
下式中的第一个等式成立：
$$
R\Gamma(x, H) =
R\Hom_{\mathcal{O}_\mathcal{X}}(K|_{\mathcal{X}/x}, P|_{\mathcal{X}/x}) =
R\Hom_{\mathcal{O}}(K|_{\mathcal{X}_{affine}/x}, P|_{\mathcal{X}_{affine}/x})
$$
第二个等式来自下述事实：
site $\mathcal{X}/x$ 的拓扑斯等价于
site $\mathcal{X}_{affine}/x$ 的拓扑斯；
参见《叠上的层》公式
（\ref{stacks-sheaves-equation-alternative-ringed}）。
可对某个 $\mathit{QC}(\mathcal{O})$ 中的 $N$
写 $K = \epsilon^*N$。
于是由《site 上的上同调》引理
\ref{sites-cohomology-lemma-QC-hom-out-of-plus-topology}，有
$$
R\Gamma(x, H) =
R\Hom_{D(\mathcal{O}(x))}(R\Gamma(x, N), R\Gamma(x, P))
$$
由引理 \ref{stacks-perfect-lemma-cohomology-parasitic}，若
$U \to \mathcal{X}$ 平坦，则 $R\Gamma(x, P) = 0$，
因而在相同假设下 $R\Gamma(x, H) = 0$。
由引理 \ref{stacks-perfect-lemma-cohomology-parasitic-complex}，可得
$R\Gamma(\mathcal{X}, H) = 0$，从而 $\Hom(K, P) = 0$。

\medskip\noindent
步骤 3：$\mathcal{A}$ 包含于 $Q_\mathcal{X}$。
设 $K$ 是 $\mathcal{A}$ 的对象，并设
$x \to x'$ 为 $\mathcal{X}_{affine}$ 的态射。
需要证明
$$
R\Gamma(x', K) \otimes_{\mathcal{O}(x')}^\mathbf{L} \mathcal{O}(x)
\longrightarrow
R\Gamma(x, K)
$$
是拟同构；参见《site 上的上同调》引理
\ref{sites-cohomology-lemma-cartesion-plus-topology}
陈述中的脚注。
由引理 \ref{stacks-perfect-lemma-bousfield-colocalization} 的证明和评注
\ref{stacks-perfect-remark-QCoh-admissible} 中的讨论可知，
$\mathcal{A}$ 是 $Lg_!$ 在
$D_\QCoh(\mathcal{O}_{\mathcal{X}_{flat, fppf}})$
上的限制的像。因此可设对某个
$D_\QCoh(\mathcal{O}_{\mathcal{X}_{flat, fppf}})$ 中的 $M$，
有 $K = Lg_!M$。
% P12-ERR-0052: see ../control/ERRATA_CANDIDATES.jsonl
于是所需等式由引理 \ref{stacks-perfect-lemma-higher-shriek-QC} 得出。
\end{proof}
% Generated mechanically from the frozen stacks-introduction.tex target.
% Do not edit this generated file; edit the bound locale source instead.

% OK, start here.
%

\setcounter{chapter}{104}
\chapter{代数叠导论}


\phantomsection
\label{stacks-introduction-section-phantom}





\section{为什么阅读本章？}
\label{stacks-introduction-section-introduction}

\noindent
我们非正式地介绍代数叠。目标是迅速引入一套简单语言，
使你能够用它思考自己所关心的模问题的局部与整体性质。
做到这一点以后，在假定一般理论已经存在的前提下，你应当能够
对模问题提出适定的问题并开始求解。如果最终得到一个有趣的结果，
便可以回到 Stacks Project 其他部分的一般理论中，按需补齐细节。

\medskip\noindent
这里采用的观点接近 \cite{KatzMazur} 和 \cite{mumford_picard}
所采用的观点。






\section{预备知识}
\label{stacks-introduction-section-preliminary}

\noindent
设 $S$ 为概形。$S$ 上的一个{\it 椭圆曲线}是一个三元组
$(E, f, 0)$，其中 $E$ 是概形，且 $f : E \to S$ 和 $0 : S \to E$
是概形态射，并满足
\begin{enumerate}
\item $f : E \to S$ 是固有、光滑且相对维数为 $1$ 的态射，
\item 对每个 $s \in S$，纤维 $E_s$ 是亏格为 $1$ 的连通曲线，
即 $H^0(E_s, \mathcal{O})$ 和 $H^1(E_s, \mathcal{O})$
都是 $1$-维 $\kappa(s)$-向量空间，并且
\item $0$ 是 $f$ 的一个截面。
\end{enumerate}
给定椭圆曲线 $(E, f, 0)/S$ 和 $(E', f', 0')/S'$，
一个{\it 位于 $a : S \to S'$ 之上的椭圆曲线态射}
是态射 $\alpha : E \to E'$，使得图表
$$
\xymatrix{
E \ar[rr]_\alpha \ar[d]^f & & E' \ar[d]_{f'}  \\
S \ar@/^5ex/[u]^0 \ar[rr]^a & & S' \ar@/_5ex/[u]_{0'}
}
$$
交换且其中的方形为笛卡儿方形；换言之，态射 $\alpha$
诱导一个同构 $E \to S \times_{S'} E'$。
我们将定义椭圆曲线叠 $\mathcal{M}_{1, 1}$。
在 Stacks Project 的其余部分中，我们完整展开 Deligne 与 Mumford
在论文 \cite{DM} 中引入的方法：把 $\mathcal{M}_{1, 1}$ 表示为
一个带有函子的范畴
$$
p : \mathcal{M}_{1, 1} \longrightarrow \Sch, \quad
(E, f, 0)/S \longmapsto S
$$
% P12-ERR-0053: see ../control/ERRATA_CANDIDATES.jsonl
这意味着要使用概形范畴上的纤维化范畴、拓扑、群胚纤维化叠、
覆盖等等。在本章中，我们暂且把这些全部搁置，以稍微不同的方式
思考问题——这种方式或许更接近该理论创立者最初的思路。


\section{椭圆曲线的模叠}
\label{stacks-introduction-section-moduli-elliptic-curves}

\noindent
我们将做如下事情：
\begin{enumerate}
\item 从你所偏好的概形范畴 $\Sch$ 出发。
\item 加入一个新符号 $\mathcal{M}_{1, 1}$。
\item 一个态射 $S \to \mathcal{M}_{1, 1}$ {\bf 就是} $S$ 上的一条椭圆曲线
$(E, f, 0)$。
% P12-ERR-0054: see ../control/ERRATA_CANDIDATES.jsonl
\item 图表
$$
\xymatrix{
S \ar[rr]_a \ar[rd]_{(E, f, 0)} & & S' \ar[ld]^{(E', F', 0')} \\
& \mathcal{M}_{1, 1}
}
$$
{\bf 交换}，当且仅当存在一个位于 $a : S \to S'$ 之上的椭圆曲线态射
$\alpha : E \to E'$。我们称 $\alpha$ {\it 见证}了该图表的交换性。
\item 注意，交换图表按如下方式粘合：
$$
\xymatrix{
S \ar[rrr]_a \ar[rrrd]_{(E, f, 0)} & & &
S' \ar[d]_{(E', F', 0')} \ar[rrr]_{a'} & & &
S'' \ar[llld]^{(E'', F'', 0'')}
\\
& & & \mathcal{M}_{1, 1}
}
$$
因为若 $\alpha$ 和 $\alpha'$ 分别见证左、右三角形的交换性，
则 $\alpha' \circ \alpha$ 见证外侧三角形的交换性。
\item 复合
$$
S \xrightarrow{a} S' \xrightarrow{(E', f', 0')} \mathcal{M}_{1, 1}
$$
由 $(E' \times_{S'} S, f' \times_{S'} S, 0' \times_{S'} S)$ 给出。
\end{enumerate}
完成这一过程后，我们恰好用一个对象扩充了概形范畴 $\Sch$……

\medskip\noindent
不过，我们还没有定义从 $\mathcal{M}_{1, 1}$ 到概形 $T$ 的态射是什么。
答案是：采用使复合有意义的最弱概念。因此，态射
$F : \mathcal{M}_{1, 1} \to T$ 是一条规则，它为每条椭圆曲线
$(E, f, 0)/S$ 指定一个态射 $F(E, f, 0) : S \to T$，并且对任意交换图表
$$
\xymatrix{
S \ar[rr]_a \ar[rd]_{(E, f, 0)} & & S' \ar[ld]^{(E', F', 0')} \\
& \mathcal{M}_{1, 1}
}
$$
图表
$$
\xymatrix{
S \ar[rr]_a \ar[rd]_{F(E, f, 0)} & & S' \ar[ld]^{F(E', F', 0')} \\
& T
}
$$
也交换。你或许听说过的 $j$-不变量
$$
j : \mathcal{M}_{1, 1} \longrightarrow \mathbf{A}^1_{\mathbf{Z}}
$$
就是一个例子。这样看来，我们已经完成了……

\medskip\noindent
然而还没有！我们仍需定义态射
$\mathcal{M}_{1, 1} \to \mathcal{M}_{1, 1}$ 的概念。
做法与之前完全相同：态射
$F : \mathcal{M}_{1, 1} \to \mathcal{M}_{1, 1}$
是一条规则，它为每条椭圆曲线 $(E, f, 0)/S$ 指定另一条椭圆曲线
$F(E, f, 0)$，并保持上述图表的交换性。不过，由于我不知道
这种函子的非平凡例子，眼下就把从 $\mathcal{M}_{1, 1}$ 到自身的
态射集合定义为只含恒等态射。

\medskip\noindent
希望你已经看出怎样向这个扩充后的范畴加入其他对象。
直观上似乎很清楚：对任意“性态良好”的模问题，都可以执行上述构造，
向我们的范畴加入一个对象。事实上，现代代数几何的很大一部分
% P12-ERR-0055: see ../control/ERRATA_CANDIDATES.jsonl
正发生在这样一个世界中：$\Sch$ 被可数多个（明确构造的）模叠所扩充。

\medskip\noindent
你也许会反对说，由此得到的并不是范畴，因为我们必须给出交换性的见证，
所以对于图表何时交换、哪些图表的组合仍然交换，存在某种“含混性”。
然而事实证明，这种为交换性配备见证的想法正是处理 $2$-范畴的
一种有效方式！因此我们继续采用它。






\section{纤维积}
\label{stacks-introduction-section-fibre-products}

\noindent
这里要问的是，下面的纤维积应当是什么：
$$
\xymatrix{
& ? \ar@{..>}[rd] \ar@{..>}[ld] \\
S \ar[rd]_{(E, f, 0)} & & S' \ar[ld]^{(E', f', 0')} \\
& \mathcal{M}_{1, 1}
}
$$
答案是：从概形 $T$ 到 $? $ 的一个态射应当是三元组
$(a, a', \alpha)$，其中 $a : T \to S$、$a' : T \to S'$
是概形态射，而
$\alpha : E \times_{S, a} T \to E' \times_{S', a'} T$
是 $T$ 上椭圆曲线的一个同构。根据前面对复合与交换图表的定义，
这确实有意义。

\begin{lemma}[关键事实]
\label{stacks-introduction-lemma-key-fact}
函子 $\Sch^{opp} \to \textit{Sets}$，
$T \mapsto \{(a, a', \alpha)\text{ 如上所述}\}$
可由一个概形 $S \times_{\mathcal{M}_{1, 1}} S'$ 表示。
\end{lemma}

\begin{proof}
证明思路。把这个函子与
$$
\mathit{Isom}_{S \times S'}(E \times S', S \times E')
$$
联系起来，并使用 Grothendieck 的 Hilbert 概形理论。
\end{proof}

\begin{remark}
\label{stacks-introduction-remark-diagonal}
我们有公式
$S \times_{\mathcal{M}_{1, 1}} S' =
(S \times S')
\times_{\mathcal{M}_{1, 1} \times \mathcal{M}_{1, 1}}
\mathcal{M}_{1, 1}$。
因此，关键事实是 $\mathcal{M}_{1, 1}$ 的对角态射
$\Delta_{\mathcal{M}_{1, 1}}$ 的一个性质。
\end{remark}

\noindent
无论如何，关键事实使我们能够作如下定义。

\begin{definition}
\label{stacks-introduction-definition-smooth}
若对每个态射 $S' \to \mathcal{M}_{1, 1}$，投影态射
$$
S \times_{\mathcal{M}_{1, 1}} S' \longrightarrow S'
$$
都是光滑的，则称态射 $S \to \mathcal{M}_{1, 1}$ {\it 光滑}。
\end{definition}

\noindent
这与概形之间光滑态射的概念相容，因为光滑态射的基变换仍然光滑。
此外，很清楚怎样把这个定义推广到以 $\mathcal{M}_{1, 1}$
（或你所偏好的模叠）为靶的态射的其他性质。特别地，下面还会将它
用于{\it 满射}。





\section{定义}
\label{stacks-introduction-section-definition}

\noindent
我们把它表述为定义而不是结论，因为希望读者尝试其他情形
（不只是叠 $\mathcal{M}_{1, 1}$，也不只是全体概形的范畴 $\Sch$）。

\begin{definition}
\label{stacks-introduction-definition-algebraic-stack}
称 $\mathcal{M}_{1, 1}$ 为一个{\it 代数叠}，当且仅当
\begin{enumerate}
\item 对 $\Sch$ 上的 \'etale 拓扑，对象满足下降；
\item 关键事实成立；
% P12-ERR-0056: see ../control/ERRATA_CANDIDATES.jsonl
\item 存在满且光滑的态射 $S \to \mathcal{M}_{1, 1}$。
\end{enumerate}
\end{definition}

\noindent
第一个条件是一种“层性质”。这里把它详细写出，因为有一个技术点
需要说明。设给定概形 $S$、一个 \'etale 覆盖 $\{S_i \to S\}$，
以及态射 $e_i : S_i \to \mathcal{M}_{1, 1}$，使图表
$$
\xymatrix{
S_i \times_S S_j \ar[rd]_{e_i \circ \text{pr}_1} \ar[rr]_{\text{id}} & &
S_i \times_S S_j \ar[ld]^{e_j \circ \text{pr}_2} \\
& \mathcal{M}_{1, 1}
}
$$
交换。在这种情形下，层条件{\it 并不}保证存在态射
$e : S \to \mathcal{M}_{1, 1}$。事实上，我们还需为上述图表选取见证
$\alpha_{ij}$，并要求在 $S_i \times_S S_j \times_S S_k$ 上，作为见证有
$$
\text{pr}_{02}^*\alpha_{ik} =
\text{pr}_{12}^*\alpha_{jk} \circ \text{pr}_{01}^*\alpha_{ij}
$$
我想这是什么意思应当很清楚……若仍不清楚，恐怕就得阅读一些
关于群胚纤维化范畴等内容。无论如何，上式通常称为{\it 余圈条件}。
“层性质”的更精确表述是：给定 $\{S_i \to S\}$、
$e_i : S_i \to \mathcal{M}_{1, 1}$ 以及满足余圈条件的见证
$\alpha_{ij}$，存在一个在唯一同构意义下唯一的
$e : S \to \mathcal{M}_{1, 1}$，使得 $e_i \cong e|_{S_i}$，
并可恢复各 $\alpha_{ij}$。

\medskip\noindent
可以看到，即使精确表述这个命题也要花些工夫。
这种“层性质”的证明依赖于代数几何的一项基本技术，即下降理论。
我的建议是先直接接受“层性质”成立，再观察它在实际中蕴含什么。
事实上，要把“层性质”提炼成一个可以写在餐巾纸上的可操作命题，
需要相当的思维灵活性。也许最简单且已经颇有意思的变体如下：
设 $L/K$ 是域的有限伽罗瓦扩张，伽罗瓦群为
$G = \text{Gal}(L/K)$。令 $T = \Spec(L)$、$S = \Spec(K)$。
则 $\{T \to S\}$ 是一个 \'etale 覆盖。
令 $(E, f, 0)$ 为 $L$ 上的一条椭圆曲线。（是的，这只意味着
$E \subset \mathbf{P}^2_L$ 由一个 Weierstrass 方程给出，
而 $0$ 是通常的无穷远点。）记
$E_\sigma = E \times_{T, \Spec(\sigma)} T$ 为该基变换。
（是的，这对应于把 $\sigma$ 作用到 Weierstrass 方程的系数上，
还是应当用 $\sigma^{-1}$？）进一步设对每个 $\sigma \in G$，
给定 $T$ 上的同构
$$
\alpha_\sigma : E \longrightarrow E_\sigma
$$
此时，上述余圈条件意味着
$$
(\alpha_\tau)^\sigma \circ \alpha_\sigma = \alpha_{\tau\sigma}
$$
对所有 $\sigma, \tau \in G$ 成立。如果你学过群上同调，
这应当十分熟悉。无论如何，$\mathcal{M}_{1, 1}$ 上的“粘合”条件说，
若这一方程组有解，则存在 $S$ 上的一条椭圆曲线 $E'$，使得
$E \cong E' \times_S T$（实际上还说得更多，因为它也告诉你
怎样恢复各 $\alpha_\sigma$）。

\medskip\noindent
挑战：你能否完全只用 Weierstrass 方程所定义的椭圆曲线来证明这一点？




\section{一个光滑覆盖}
\label{stacks-introduction-section-smooth}

\noindent
最后要做的是找出 $\mathcal{M}_{1, 1}$ 的一个光滑覆盖。
事实上，在某种意义下，光滑覆盖的存在{\it 蕴含}
\footnote{这里稍有取巧，因为检验光滑性时必须证明与关键事实
非常接近的结论——毕竟光滑性是用纤维积定义的。好处是，
只需在纤维积一侧是你试图证明给出光滑覆盖的那个态射时，
证明这些纤维积存在。}关键事实！
对椭圆曲线，我们用 Weierstrass 方程来构造一个。

\medskip\noindent
令
$$
W = \Spec(\mathbf{Z}[a_1, a_2, a_3, a_4, a_6, 1/\Delta])
$$
其中 $\Delta \in \mathbf{Z}[a_1, a_2, a_3, a_4, a_6]$
是某个多项式（见下文）。令
$$
\mathbf{P}_W^2 \supset E_W :
zy^2 + a_1 xyz + a_3 yz^2 = x^3 + a_2x^2z + a_4xz^2 + a_6z^3.
$$
以 $f_W : E_W \to W$ 记投影。最后，以 $0_W : W \to E_W$
记由 $(0 : 1 : 0)$ 给出的 $f_W$ 的截面。事实证明，在赋权
$\deg(a_i) = i$ 下，存在一个关于 $a_i$ 的加权次数为 $12$ 的齐次多项式 $\Delta$，
使 $E_W \to W$ 光滑。可以通过计算 Weierstrass 方程的偏导数
显式求出它——当然也可以直接查表。还可以让 pari/gp 替你计算。
它就是
\begin{align*}
\Delta & = -a_6a_1^6 + a_4a_3a_1^5 + ((-a_3^2 - 12a_6)a_2 + a_4^2)a_1^4 + \\
& (8a_4a_3a_2 + (a_3^3 + 36a_6a_3))a_1^3 + \\
& ((-8a_3^2 - 48a_6)a_2^2 + 8a_4^2a_2 + (-30a_4a_3^2 + 72a_6a_4))a_1^2 + \\
& (16a_4a_3a_2^2 + (36a_3^3 + 144a_6a_3)a_2 - 96a_4^2a_3)a_1 + \\
& (-16a_3^2 - 64a_6)a_2^3 + 16a_4^2a_2^2 + (72a_4a_3^2 + 288a_6a_4)a_2 + \\
& -27a_3^4 - 216a_6a_3^2  -64a_4^3 - 432a_6^2
\end{align*}
你或许认得最后两项：在
$y^2 = x^3 + Ax + B$ 的情形，判别式为
$-64A^3 - 432B^2 = -16(4A^3 + 27B^2)$。

\begin{lemma}
\label{stacks-introduction-lemma-Weierstrass-smooth-cover}
态射 $W \xrightarrow{(E_W, f_W, 0_W)} \mathcal{M}_{1, 1}$
光滑且满。
\end{lemma}

\begin{proof}
满性来自每条域上的椭圆曲线都有 Weierstrass 方程这一事实。
% P12-ERR-0057: see ../control/ERRATA_CANDIDATES.jsonl
下面粗略说明一种证明光滑性的方法。考虑子群概形
$$
H =
\left\{
\left(
\begin{matrix}
u^2 & s & 0 \\
0 & u^3 & 0 \\
r & t & 1
\end{matrix}
\right)
\middle|
\begin{matrix}
u\text{ 为可逆元} \\
s, r, t\text{ 任意}
\end{matrix}
\right\}
\subset
\text{GL}_{3, \mathbf{Z}}
$$
$H$ 在 Weierstrass 概形 $W$ 上有作用 $H \times W \to W$。
为求出这一作用的方程，只需写出由 $H$ 中矩阵给出的坐标变换
如何作用于一般 Weierstrass 方程。随后可证明下列命题成立：
\begin{enumerate}
\item 任意椭圆曲线 $(E, f, 0)/S$ 在 $S$ 上 Zariski 局部地
都有一个 Weierstrass 方程；
\item $(E, f, 0)$ 的任意两个 Weierstrass 方程在 Zariski 局部
相差 $H$ 的一个元素。
\end{enumerate}
考察纤维积
$S \times_{\mathcal{M}_{1, 1}} W =
\mathit{Isom}_{S \times W}(E \times W, S \times E_W)$，
由此得出：这意味着态射 $W \to \mathcal{M}_{1, 1}$
是一个 $H$-挠子（即主齐性空间）。由于 $H \to \Spec(\mathbf{Z})$ 光滑，
而光滑群概形上的挠子也是光滑的，结论即得。
\end{proof}

\begin{remark}
\label{stacks-introduction-remark-quotient-stack}
上面概述的论证实际上表明
$\mathcal{M}_{1, 1} = [W/H]$ 是一个全局商叠。
证明一个模叠为代数叠的论证，约有 50\% 的时候还会表明
它是一个全局商叠。
\end{remark}



\section{代数叠的性质}
\label{stacks-introduction-section-properties}

\noindent
好，现在我们知道 $\mathcal{M}_{1, 1}$ 是代数叠了。
这能用来做什么？这里真正有帮助的，与其说是它为代数叠这一事实，
不如说是这样一种观点：$\mathcal{M}_{1, 1}$ 的性质应编码在态射
$S \to \mathcal{M}_{1, 1}$ 的性质中，也就是编码在椭圆曲线族中。
下面列出若干例子。

\medskip\noindent
{\bf 局部性质：}
$$
\mathcal{M}_{1, 1} \to \Spec(\mathbf{Z})\text{ 是光滑的}
\Leftrightarrow
W \to \Spec(\mathbf{Z})\text{ 是光滑的}
$$
{\bf 思路}。代数叠的局部性质编码在其光滑覆盖的局部性质中。

\medskip\noindent
{\bf 整体性质：}
$$
\begin{matrix}
\mathcal{M}_{1, 1}\text{ 是拟紧的} \Leftarrow W\text{ 是拟紧的}
\\
\mathcal{M}_{1, 1}\text{ 是不可约的} \Leftarrow W\text{ 是不可约的}
\end{matrix}
$$
{\bf 思路}。代数叠的某些整体性质可以从其一个{\it 合适的}
% P12-ERR-0058: see ../control/ERRATA_CANDIDATES.jsonl
\footnote{我猜想可能存在没有不可约光滑覆盖的不可约代数叠——
但如果确实存在，它一定相当糟糕！}光滑覆盖的相应性质读出。

\medskip\noindent
{\bf 拟凝聚层：}
$$
\QCoh(\mathcal{O}_{\mathcal{M}_{1, 1}}) =
H\text{-等变拟凝聚模，定义于 }W
$$
{\bf 思路}。一方面，$\mathcal{M}_{1, 1}$ 上的一个拟凝聚模
应当对应于：对每个态射 $e : S \to \mathcal{M}_{1, 1}$，
在 $S$ 上给定一个拟凝聚层 $\mathcal{F}_{S, e}$。
特别地，可取态射 $(E_W, f_W, 0_W) :  W \to \mathcal{M}_{1, 1}$。
由于这个态射是 $H$-等变的，所得拟凝聚模 $\mathcal{F}_W$
也是 $H$-等变的。反过来，给定一个 $H$-等变模，从
$S \times_{e, \mathcal{M}_{1, 1}} W$ 是 $H$-挠子这一观察出发，
% P12-ERR-0059: see ../control/ERRATA_CANDIDATES.jsonl
可以通过下降理论恢复各层 $\mathcal{F}_{S, e}$。

\medskip\noindent
{\bf Picard 群：}
$$
\Pic(\mathcal{M}_{1, 1}) = \Pic_H(W) = \mathbf{Z}/12\mathbf{Z}
$$
{\bf 思路}。第一个等式已在上面看到。注意 $\Pic(W) = 0$，
因为环 $\mathbf{Z}[a_1, a_2, a_3, a_4, a_6, 1/\Delta]$
的类群平凡。存在正合序列
$$
\mathbf{Z}\Delta \to
\Pic_H(\mathbf{A}^5_{\mathbf{Z}}) \to
\Pic_H(W) \to 0
$$
中间的群等于 $\Hom(H, \mathbf{G}_m) = \mathbf{Z}$。
% P12-ERR-0060: see ../control/ERRATA_CANDIDATES.jsonl
$\Delta$ 的像为 $12$，因为 $\Delta$ 的次数是 $12$。
这个论证大致正确，见 \cite{PicM11}。

\medskip\noindent
{\bf \'Etale 上同调：}设 $\Lambda$ 为环。
% P12-ERR-0061: see ../control/ERRATA_CANDIDATES.jsonl
存在一个第一象限谱序列，收敛到
$H^{p + q}_\etale(\mathcal{M}_{1, 1}, \Lambda)$，其 $E_2$-页为
$$
E_2^{p, q} = H_\etale^q(W \times H \times \ldots \times H, \Lambda)
\quad(p\text{ 个因子，每个均为 }H)
$$
{\bf 思路}。注意
$$
W \times_{\mathcal{M}_{1, 1}} W \times_{\mathcal{M}_{1, 1}}
\ldots \times_{\mathcal{M}_{1, 1}} W
= W \times H \times \ldots \times H
$$
因为 $W \to \mathcal{M}_{1, 1}$ 是一个 $H$-挠子。
这个谱序列是光滑覆盖 $\{W \to \mathcal{M}_{1, 1}\}$ 的
{\v C}ech-上同调谱序列。例如，由于 $W$ 连通，我们看到
$H^0_\etale(\mathcal{M}_{1, 1}, \Lambda) = \Lambda$；
又因为 $H^1_\etale(W, \Lambda) = 0$，所以
$H^1_\etale(\mathcal{M}_{1, 1}, \Lambda) = 0$
（当然，这需要证明）。当然，光滑覆盖
$W \to \mathcal{M}_{1, 1}$ 对计算 \'etale 上同调而言未必“最优”。
% Generated mechanically from the frozen stacks-more-morphisms.tex target.
% Do not edit this generated file; edit the bound locale source instead.

% OK, start here.
%

\setcounter{chapter}{105}
\chapter{叠态射的进一步讨论}



\phantomsection
\label{stacks-more-morphisms-section-phantom}


\section{引言}
\label{stacks-more-morphisms-section-introduction}

\noindent
本章继续研究代数叠态射的性质。对于具有可表示对角态射的
拟分离代数叠，可参见 \cite{LM-B}。




\section{约定与语言上的简化}
\label{stacks-more-morphisms-section-conventions}

\noindent
我们继续采用《叠的性质》第
\ref{stacks-properties-section-conventions} 节中引入的约定与
语言上的简化说法。






\section{增厚}
\label{stacks-more-morphisms-section-thickenings}

\noindent
下面的术语或许并非完全标准，但使用起来很方便。
若 $\mathcal{Y}$ 是代数叠 $\mathcal{X}$ 的闭子叠，
则态射 $\mathcal{Y} \to \mathcal{X}$ 是可表示的。

\begin{definition}
\label{stacks-more-morphisms-definition-thickening}
增厚。
\begin{enumerate}
\item 若 $\mathcal{X}$ 是代数叠 $\mathcal{X}'$ 的闭子叠，且相应的
拓扑空间相等，则称代数叠 $\mathcal{X}'$ 是代数叠 $\mathcal{X}$ 的一个
{\it 增厚}。
\item 给定两个增厚 $\mathcal{X} \subset \mathcal{X}'$ 与
$\mathcal{Y} \subset \mathcal{Y}'$，所谓{\it 增厚的态射}，是指代数叠
的态射 $f' : \mathcal{X}' \to \mathcal{Y}'$，使得
$f'|_\mathcal{X}$ 经过闭子叠 $\mathcal{Y}$ 分解。在这种情形下，令
$f = f'|_\mathcal{X} : \mathcal{X} \to \mathcal{Y}$，并称
$(f, f') : (\mathcal{X} \subset \mathcal{X}') \to
(\mathcal{Y} \subset \mathcal{Y}')$ 是一个增厚的态射。
\item 设 $\mathcal{Z}$ 为代数叠。类似地定义
{\it $\mathcal{Z}$ 上的增厚}以及
{\it $\mathcal{Z}$ 上增厚的态射}。
这意味着代数叠 $\mathcal{X}'$ 和 $\mathcal{Y}'$ 都配备到
$\mathcal{Z}$ 的结构态射，且 $f'$ 位于一个适当的代数叠
$2$-交换图中。
\end{enumerate}
\end{definition}

\noindent
设 $\mathcal{X} \subset \mathcal{X}'$ 为代数叠的增厚。
设 $U'$ 为概形，并设 $U' \to \mathcal{X}'$ 为满光滑态射。
令 $U = \mathcal{X} \times_{\mathcal{X}'} U'$，便得到增厚的态射
$$
(U \subset U') \longrightarrow (\mathcal{X} \subset \mathcal{X}')
$$
且 $U \to \mathcal{X}$ 是满光滑态射。我们经常可以从增厚
$U \subset U'$ 的相应性质推出增厚
$\mathcal{X} \subset \mathcal{X}'$ 的性质。
有时，作为语言上的简化，若态射 $\mathcal{X} \to \mathcal{X}'$
是闭浸入并诱导双射 $|\mathcal{X}| \to |\mathcal{X}'|$，
也称它为一个增厚。

\begin{lemma}
\label{stacks-more-morphisms-lemma-thickening}
设 $i : \mathcal{X} \to \mathcal{X}'$ 为代数叠的态射。
下列条件等价：
\begin{enumerate}
\item $i$ 是代数叠的增厚（采用上面的简化说法）；
\item $i$ 可由代数空间表示，并且是《叠的性质》第
\ref{stacks-properties-section-properties-morphisms} 节意义下的增厚。
\end{enumerate}
在这种情形下，$i$ 是闭浸入及万有同胚。
\end{lemma}

\begin{proof}
由《空间态射的进一步讨论》中的引理
\ref{spaces-more-morphisms-lemma-descending-property-thickening} 和
\ref{spaces-more-morphisms-lemma-base-change-thickening}，
代数空间的态射为（一阶）增厚这一性质 $P$ 在基底上是 fpqc 局部的，
且在基变换下稳定。因此，《叠的性质》第
\ref{stacks-properties-section-properties-morphisms} 节的讨论确实适用。
至此，(1) 与 (2) 的等价性来自 $P = P_1 + P_2$，其中 $P_1$
是成为闭浸入的性质，而 $P_2$ 是成为满态射的性质。
% P12-ERR-0062: see ../control/ERRATA_CANDIDATES.jsonl
（严格说来，读者还应参阅《空间态射的进一步讨论》中的定义
\ref{spaces-more-morphisms-definition-thickening}、
《叠的性质》中的定义 \ref{stacks-properties-definition-immersion}
及其后的讨论、《空间的态射》中的引理
\ref{spaces-morphisms-lemma-surjective-representable}，以及
《叠的性质》第 \ref{stacks-properties-section-surjective} 节，
以确认所有这些概念彼此吻合。）最后一个断言由前述内容立即可得。
\end{proof}











\noindent
下文将不再另加说明地使用这个引理。利用该引理所用的同一组参考文献，
即《空间态射的进一步讨论》中的引理
\ref{spaces-more-morphisms-lemma-descending-property-thickening} 和
\ref{spaces-more-morphisms-lemma-base-change-thickening}，
可以如下定义一阶增厚。

\begin{definition}
\label{stacks-more-morphisms-definition-first-order-thickening}
若 $\mathcal{X}$ 是代数叠 $\mathcal{X}'$ 的闭子叠，且
$\mathcal{X} \to \mathcal{X}'$ 是《叠的性质》第
\ref{stacks-properties-section-properties-morphisms} 节意义下的一阶增厚，
则称代数叠 $\mathcal{X}'$ 是代数叠 $\mathcal{X}$ 的一个
{\it 一阶增厚}。
\end{definition}

\noindent
若 $(U \subset U') \to (\mathcal{X} \subset \mathcal{X}')$
是如上的概形光滑覆盖，那么这恰好意味着 $U \subset U'$
是一阶增厚。下面陈述必需的几个引理。

\begin{lemma}
\label{stacks-more-morphisms-lemma-base-change-thickening}
设 $\mathcal{Y} \subset \mathcal{Y}'$ 为代数叠的增厚。
设 $\mathcal{X}' \to \mathcal{Y}'$ 为代数叠的态射，并令
$\mathcal{X} = \mathcal{Y} \times_{\mathcal{Y}'} \mathcal{X}'$。
则
$(\mathcal{X} \subset \mathcal{X}') \to (\mathcal{Y} \subset \mathcal{Y}')$
是增厚的态射。若 $\mathcal{Y} \subset \mathcal{Y}'$ 是一阶增厚，
则 $\mathcal{X} \subset \mathcal{X}'$ 也是一阶增厚。
\end{lemma}

\begin{proof}
参见上面的讨论、《叠的性质》第
\ref{stacks-properties-section-properties-morphisms} 节，以及
《空间态射的进一步讨论》中的引理
\ref{spaces-more-morphisms-lemma-base-change-thickening}。
\end{proof}

\begin{lemma}
\label{stacks-more-morphisms-lemma-composition-thickening}
若 $\mathcal{X} \subset \mathcal{X}'$ 与
$\mathcal{X}' \subset \mathcal{X}''$ 都是代数叠的增厚，
则 $\mathcal{X} \subset \mathcal{X}''$ 也是代数叠的增厚。
\end{lemma}

\begin{proof}
参见上面的讨论、《叠的性质》第
\ref{stacks-properties-section-properties-morphisms} 节，以及
《空间态射的进一步讨论》中的引理
\ref{spaces-more-morphisms-lemma-composition-thickening}
\end{proof}

\begin{example}
\label{stacks-more-morphisms-example-reduction-thickening}
设 $\mathcal{X}'$ 为代数叠。则 $\mathcal{X}'$ 是其约化
$\mathcal{X}'_{red}$ 的一个增厚，参见《叠的性质》中的定义
\ref{stacks-properties-definition-reduced-induced-stack}。
此外，若 $\mathcal{X} \subset \mathcal{X}'$ 是代数叠的增厚，则
$\mathcal{X}'_{red} = \mathcal{X}_{red} \subset \mathcal{X}$。
换言之，$\mathcal{X} = \mathcal{X}'_{red}$ 当且仅当
$\mathcal{X}$ 是约化代数叠。
\end{example}

\begin{lemma}
\label{stacks-more-morphisms-lemma-reduced-diagonal}
设 $(f, f') : (\mathcal{X} \subset \mathcal{X}') \to
(\mathcal{Y} \subset \mathcal{Y}')$ 为代数叠增厚的态射。则
$\mathcal{X} \times_\mathcal{Y} \mathcal{X} \to
\mathcal{X}' \times_{\mathcal{Y}'} \mathcal{X}'$
是增厚，并且典范图表
$$
\xymatrix{
\mathcal{X} \ar[r]_-\Delta \ar[d] &
\mathcal{X} \times_\mathcal{Y} \mathcal{X} \ar[d] \\
\mathcal{X}' \ar[r]^-{\Delta'} &
\mathcal{X}' \times_{\mathcal{Y}'} \mathcal{X}'
}
$$
是笛卡儿图表。
\end{lemma}

\begin{proof}
由于 $\mathcal{X} \to \mathcal{Y}'$ 经过闭子叠
$\mathcal{Y}$ 分解，我们有
$\mathcal{X} \times_\mathcal{Y} \mathcal{X} =
\mathcal{X} \times_{\mathcal{Y}'} \mathcal{X}$。
因而
$\mathcal{X} \times_\mathcal{Y} \mathcal{X} \to
\mathcal{X}' \times_{\mathcal{Y}'} \mathcal{X}'$
同构于复合
$$
\mathcal{X} \times_{\mathcal{Y}'} \mathcal{X} \to
\mathcal{X} \times_{\mathcal{Y}'} \mathcal{X}' \to
\mathcal{X}' \times_{\mathcal{Y}'} \mathcal{X}'
$$
其中两个态射都是增厚的基变换，故都是增厚
（引理 \ref{stacks-more-morphisms-lemma-base-change-thickening}）。因此其复合也是增厚
（引理 \ref{stacks-more-morphisms-lemma-composition-thickening}）。
由于 $\mathcal{X} \to \mathcal{X}'$ 是单态射，将《叠的性质》中的引理
\ref{stacks-properties-lemma-monomorphism-diagonal}
应用于 $\mathcal{X} \to \mathcal{X}' \to \mathcal{Y}'$，
便得到引理的最后一个断言。
\end{proof}

\begin{lemma}
\label{stacks-more-morphisms-lemma-thickening-diagonals}
设 $(f, f') : (\mathcal{X} \subset \mathcal{X}') \to
(\mathcal{Y} \subset \mathcal{Y}')$ 为代数叠增厚的态射。
令 $\Delta : \mathcal{X} \to \mathcal{X} \times_\mathcal{Y} \mathcal{X}$ 与
$\Delta' : \mathcal{X}' \to \mathcal{X}' \times_{\mathcal{Y}'} \mathcal{X}'$
为相应的对角态射。则下列各性质由 $\Delta$ 满足，当且仅当由
$\Delta'$ 满足：
(a) 可由概形表示，(b) 仿射，(c) 满，(d) 拟紧，
(e) 万有闭，(f) 整，(g) 拟分离，(h) 分离，
(i) 万有单射，(j) 万有开，(k) 局部拟有限，
(l) 有限，(m) 非分歧，(n) 单态射，(o) 浸入，
(p) 闭浸入，以及 (q) 固有。
\end{lemma}

\begin{proof}
注意到
$$
(\Delta, \Delta') :
(\mathcal{X} \subset \mathcal{X}')
\longrightarrow
(\mathcal{X} \times_\mathcal{Y} \mathcal{X} \subset
\mathcal{X}' \times_{\mathcal{Y}'} \mathcal{X}')
$$
是增厚的态射（引理 \ref{stacks-more-morphisms-lemma-reduced-diagonal}）。此外，由
《叠的态射》中的引理 \ref{stacks-morphisms-lemma-properties-diagonal}，
$\Delta$ 与 $\Delta'$ 均可由代数空间表示。因此，结合《叠的性质》第
\ref{stacks-properties-section-properties-morphisms} 节的讨论，
再应用《空间态射的进一步讨论》中的引理
\ref{spaces-more-morphisms-lemma-thicken-property-morphisms}，
便得到情形 (a)、(b)、(c)、(d)、(e)、(f)、(g)、(h)、(i) 与 (j)。

\medskip\noindent
引理 \ref{stacks-more-morphisms-lemma-reduced-diagonal} 告诉我们
$\mathcal{X} = (\mathcal{X} \times_\mathcal{Y} \mathcal{X})
\times_{(\mathcal{X}' \times_{\mathcal{Y}'} \mathcal{X}')} \mathcal{X}'$。
此外，由前述《叠的态射》中的引理
\ref{stacks-morphisms-lemma-properties-diagonal}，
$\Delta$ 与 $\Delta'$ 都是局部有限型的。
% P12-ERR-0063: see ../control/ERRATA_CANDIDATES.jsonl
因此，对情形 (k)、(l)、(m)、(n)、(o)、(p) 与 (q)，
应用《空间态射的进一步讨论》中的引理
\ref{spaces-more-morphisms-lemma-properties-that-extend-over-thickenings}
即可得到结论。
\end{proof}

\noindent
作为推论，得到下面这个令人满意的结果。

\begin{lemma}
\label{stacks-more-morphisms-lemma-thickening-properties}
\begin{reference}
\cite[Theorem 2.2.5]{Conrad-moduli}
\end{reference}
设 $\mathcal{X} \subset \mathcal{X}'$ 为代数叠的增厚。则
\begin{enumerate}
\item $\mathcal{X}$ 是代数空间，当且仅当 $\mathcal{X}'$ 是代数空间；
\item $\mathcal{X}$ 是概形，当且仅当 $\mathcal{X}'$ 是概形；
\item $\mathcal{X}$ 是 DM，当且仅当 $\mathcal{X}'$ 是 DM；
\item $\mathcal{X}$ 是拟-DM，当且仅当 $\mathcal{X}'$ 是拟-DM；
\item $\mathcal{X}$ 是分离的，当且仅当 $\mathcal{X}'$ 是分离的；
\item $\mathcal{X}$ 是拟分离的，当且仅当 $\mathcal{X}'$ 是拟分离的；并且
% P12-ERR-0064: see ../control/ERRATA_CANDIDATES.jsonl
\item 在此处补充更多内容。
\end{enumerate}
\end{lemma}

\begin{proof}
在每一种情形中，我们都把问题化归为关于对角态射的问题，
再把引理 \ref{stacks-more-morphisms-lemma-thickening-diagonals} 应用于增厚的态射
$$
(\mathcal{X} \subset \mathcal{X}') \to
\left(\Spec(\mathbf{Z}) \subset \Spec(\mathbf{Z})\right)
$$
这里先通过《代数叠》中的定义 \ref{algebraic-definition-viewed-as}，
把 $\mathcal{X} \subset \mathcal{X}'$ 看成
$\Spec(\mathbf{Z})$ 上代数叠的增厚。

\medskip\noindent
情形 (1)。代数叠是代数空间，当且仅当其对角态射为单态射，
参见《叠的态射》中的引理 \ref{stacks-morphisms-lemma-hierarchy}
（这也可直接由《代数叠》中的命题
\ref{algebraic-proposition-algebraic-stack-no-automorphisms} 得到）。

\medskip\noindent
情形 (2)。由 (1)，可设 $\mathcal{X}$ 与 $\mathcal{X}'$
均为代数空间，随后应用《空间态射的进一步讨论》中的引理
\ref{spaces-more-morphisms-lemma-thickening-scheme}。

\medskip\noindent
情形 (3)--(6)。每种情形都对应于对角态射的一个条件，
参见《叠的态射》中的定义
\ref{stacks-morphisms-definition-separated} 和
\ref{stacks-morphisms-definition-absolute-separated}。
\end{proof}









\section{增厚的态射}
\label{stacks-more-morphisms-section-morphisms-thickenings}

\noindent
若 $(f, f') : (\mathcal{X} \subset \mathcal{X}') \to
(\mathcal{Y} \subset \mathcal{Y}')$ 是代数叠增厚的态射，
那么态射 $f$ 的性质往往由 $f'$ 继承。这里有若干种变体。

\begin{lemma}
\label{stacks-more-morphisms-lemma-thicken-property-morphisms}
设 $(f, f') : (\mathcal{X} \subset \mathcal{X}') \to
(\mathcal{Y} \subset \mathcal{Y}')$
为代数叠增厚的态射。则
\begin{enumerate}
\item $f$ 是仿射态射，当且仅当 $f'$ 是仿射态射；
\item $f$ 是满态射，当且仅当 $f'$ 是满态射；
% P12-ERR-0065: see ../control/ERRATA_CANDIDATES.jsonl
\item $f$ 是拟紧的，当且仅当 $f'$ 是拟紧的；
\item $f$ 是万有闭的，当且仅当 $f'$ 是万有闭的；
\item $f$ 是整态射，当且仅当 $f'$ 是整态射；
\item $f$ 是万有单射的，当且仅当 $f'$ 是万有单射的；
\item $f$ 是万有开的，当且仅当 $f'$ 是万有开的；
\item $f$ 是拟-DM 的，当且仅当 $f'$ 是拟-DM 的；
\item $f$ 是 DM 的，当且仅当 $f'$ 是 DM 的；
\item $f$ 是（拟）分离的，当且仅当 $f'$ 是（拟）分离的；
\item $f$ 是可表示的，当且仅当 $f'$ 是可表示的；
\item $f$ 可由代数空间表示，当且仅当 $f'$ 可由代数空间表示；
% P12-ERR-0066: see ../control/ERRATA_CANDIDATES.jsonl
\item 在此处补充更多内容。
\end{enumerate}
\end{lemma}

\begin{proof}
由引理 \ref{stacks-more-morphisms-lemma-thickening}，态射
$\mathcal{X} \to \mathcal{X}'$ 与 $\mathcal{Y} \to \mathcal{Y}'$
都是万有同胚。因此，关于
$|f| : |\mathcal{X}| \to |\mathcal{Y}|$ 的任何条件，
都等价于 $|f'| : |\mathcal{X}'| \to |\mathcal{Y}'|$ 上的相应条件；
在用任意态射 $\mathcal{Z}' \to \mathcal{Y}'$ 作基变换后，
同样成立。这证明了 (2)、(3)、(4)、(6)、(7)。

\medskip\noindent
在情形 (8)、(9)、(10)、(12) 中，可以像引理
\ref{stacks-more-morphisms-lemma-thickening-diagonals} 那样，把 $f$ 与 $f'$ 上的条件
转化为对角态射 $\Delta$ 与 $\Delta'$ 上的条件。参见
《叠的态射》中的定义 \ref{stacks-morphisms-definition-separated}
和引理 \ref{stacks-morphisms-lemma-hierarchy}。
因此，这些情形由引理 \ref{stacks-more-morphisms-lemma-thickening-diagonals} 得出。

\medskip\noindent
(11) 的证明。若 $f'$ 可表示，则 $f$ 也可表示，因为对概形 $T$
及态射 $T \to \mathcal{Y}$，有
$\mathcal{X} \times_\mathcal{Y} T =
\mathcal{X} \times_{\mathcal{X}'} (\mathcal{X}' \times_{\mathcal{Y}'} T)$，
而 $\mathcal{X} \to \mathcal{X}'$ 是闭浸入（因而可表示）。
反过来，设 $f$ 可表示，并设 $T' \to \mathcal{Y}'$ 为态射，
其中 $T'$ 是概形。则
$$
\mathcal{X} \times_{\mathcal{Y}}
(\mathcal{Y} \times_{\mathcal{Y}'} T') =
\mathcal{X} \times_{\mathcal{X}'}
(\mathcal{X}' \times_{\mathcal{Y}'} T') \to
\mathcal{X}' \times_{\mathcal{Y}'} T'
$$
是一个增厚（由引理 \ref{stacks-more-morphisms-lemma-base-change-thickening}），
且其源是概形。故由引理 \ref{stacks-more-morphisms-lemma-thickening-properties}，
其靶也是概形。

\medskip\noindent
在情形 (1) 与 (5) 中，若 $f$ 或 $f'$ 中任一个具有所述性质，
则由 (11)，$f$ 与 $f'$ 都可表示。在这种情形下，选取代数空间 $V'$
及满光滑态射 $V' \to \mathcal{Y}'$。令
$V = \mathcal{Y} \times_{\mathcal{Y}'} V'$、
$U' = \mathcal{X}' \times_{\mathcal{Y}'} V'$，以及
$U = \mathcal{X} \times_{\mathcal{Y}'} V'$。
于是，由《叠的性质》中的引理
\ref{stacks-properties-lemma-check-property-covering} 所给原则，
所需结论可由代数空间增厚的态射
$(U \subset U') \to (V \subset V')$ 的相应结论得到。
关于代数空间情形的相应结果，参见《空间态射的进一步讨论》中的引理
\ref{spaces-more-morphisms-lemma-thicken-property-morphisms}。
\end{proof}






\section{代数叠的无穷小形变}
\label{stacks-more-morphisms-section-deform}

\noindent
本节对应于《空间态射的进一步讨论》第
\ref{spaces-more-morphisms-section-deform} 节。

\begin{lemma}
\label{stacks-more-morphisms-lemma-flatness-morphism-thickenings-fp-over-ft}
考虑代数叠增厚的交换图表
$$
\xymatrix{
(\mathcal{X} \subset \mathcal{X}') \ar[rr]_{(f, f')} \ar[rd] & &
(\mathcal{Y} \subset \mathcal{Y}') \ar[ld] \\
& (\mathcal{B} \subset \mathcal{B}')
}
$$
假设
\begin{enumerate}
\item $\mathcal{Y}' \to \mathcal{B}'$ 是局部有限型的；
\item $\mathcal{X}' \to \mathcal{B}'$ 是平坦且局部有限表示的；
\item $f$ 是平坦的；并且
\item $\mathcal{X} = \mathcal{B} \times_{\mathcal{B}'} \mathcal{X}'$ 且
$\mathcal{Y} = \mathcal{B} \times_{\mathcal{B}'} \mathcal{Y}'$。
\end{enumerate}
则 $f'$ 是平坦的，并且对 $|f'|$ 的像中的每个
$y' \in |\mathcal{Y}'|$，态射 $\mathcal{Y}' \to \mathcal{B}'$
在 $y'$ 处平坦。
\end{lemma}

\begin{proof}
选取代数空间 $U'$ 及满光滑态射 $U' \to \mathcal{B}'$。
选取代数空间 $V'$ 及满光滑态射
$V' \to U' \times_{\mathcal{B}'} \mathcal{Y}'$。
选取代数空间 $W'$ 及满光滑态射
$W' \to V' \times_{\mathcal{Y}'} \mathcal{X}'$。
令 $U, V, W$ 分别为 $U', V', W'$ 沿
$\mathcal{B} \to \mathcal{B}'$ 的基变换。
于是，$f'$ 的平坦性等价于 $W' \to V'$ 的平坦性，
而已知 $W \to V$ 平坦。因此，可以把代数空间情形的引理
应用于代数空间增厚的图表
$$
\xymatrix{
(W \subset W') \ar[rr] \ar[rd] & & (V \subset V') \ar[ld] \\
& (U \subset U')
}
$$
参见《空间态射的进一步讨论》中的引理
\ref{spaces-more-morphisms-lemma-flatness-morphism-thickenings-fp-over-ft}。
关于 $\mathcal{Y}'/\mathcal{B}'$ 在 $|f'|$ 像中各点处平坦的断言，
也以同样方式得出。
\end{proof}

\begin{lemma}
\label{stacks-more-morphisms-lemma-deform-property-fp-over-ft}
考虑代数叠增厚的交换图表
$$
\xymatrix{
(\mathcal{X} \subset \mathcal{X}') \ar[rr]_{(f, f')} \ar[rd] & &
(\mathcal{Y} \subset \mathcal{Y}') \ar[ld] \\
& (\mathcal{B} \subset \mathcal{B}')
}
$$
% P12-ERR-0067: see ../control/ERRATA_CANDIDATES.jsonl
假设 $\mathcal{Y}' \to \mathcal{B}'$ 是局部有限型的，
$\mathcal{X}' \to \mathcal{B}'$ 是平坦且局部有限表示的，
$\mathcal{X} = \mathcal{B} \times_{\mathcal{B}'} \mathcal{X}'$，并且
$\mathcal{Y} = \mathcal{B} \times_{\mathcal{B}'} \mathcal{Y}'$。则
\begin{enumerate}
\item $f$ 平坦，当且仅当 $f'$ 平坦；
\label{stacks-more-morphisms-item-flat-fp-over-ft}
\item $f$ 是同构，当且仅当 $f'$ 是同构；
\label{stacks-more-morphisms-item-isomorphism-fp-over-ft}
\item $f$ 是开浸入，当且仅当 $f'$ 是开浸入；
\label{stacks-more-morphisms-item-open-immersion-fp-over-ft}
\item $f$ 是单态射，当且仅当 $f'$ 是单态射；
\label{stacks-more-morphisms-item-monomorphism-fp-over-ft}
\item $f$ 是局部拟有限的，当且仅当 $f'$ 是局部拟有限的；
\label{stacks-more-morphisms-item-quasi-finite-fp-over-ft}
\item $f$ 是 syntomic 的，当且仅当 $f'$ 是 syntomic 的；
\label{stacks-more-morphisms-item-syntomic-fp-over-ft}
\item $f$ 是光滑的，当且仅当 $f'$ 是光滑的；
\label{stacks-more-morphisms-item-smooth-fp-over-ft}
\item $f$ 是非分歧的，当且仅当 $f'$ 是非分歧的；
\label{stacks-more-morphisms-item-unramified-fp-over-ft}
\item $f$ 是 \'etale 的，当且仅当 $f'$ 是 \'etale 的；
\label{stacks-more-morphisms-item-etale-fp-over-ft}
\item $f$ 是有限的，当且仅当 $f'$ 是有限的；并且
\label{stacks-more-morphisms-item-finite-fp-over-ft}
% P12-ERR-0068: see ../control/ERRATA_CANDIDATES.jsonl
\item 在此处补充更多内容。
\end{enumerate}
\end{lemma}

\begin{proof}
在情形 (\ref{stacks-more-morphisms-item-flat-fp-over-ft}) 中，结论由引理
\ref{stacks-more-morphisms-lemma-flatness-morphism-thickenings-fp-over-ft} 得出。

\medskip\noindent
在情形 (\ref{stacks-more-morphisms-item-syntomic-fp-over-ft})、
(\ref{stacks-more-morphisms-item-smooth-fp-over-ft}) 中，可以采用引理
\ref{stacks-more-morphisms-lemma-flatness-morphism-thickenings-fp-over-ft} 证明中所用的方法。
具体而言，选取代数空间 $U'$ 及满光滑态射
$U' \to \mathcal{B}'$。选取代数空间 $V'$ 及满光滑态射
$V' \to U' \times_{\mathcal{B}'} \mathcal{Y}'$。
选取代数空间 $W'$ 及满光滑态射
$W' \to V' \times_{\mathcal{Y}'} \mathcal{X}'$。
令 $U, V, W$ 分别为 $U', V', W'$ 沿
$\mathcal{B} \to \mathcal{B}'$ 的基变换。
% P12-ERR-0069: see ../control/ERRATA_CANDIDATES.jsonl
于是，$f$（相应地，\ $f'$）具有该性质，等价于
$W' \to V'$（相应地，\ $W \to V$）具有该性质。
因此，可以把代数空间情形的引理应用于代数空间增厚的图表
$$
\xymatrix{
(W \subset W') \ar[rr] \ar[rd] & & (V \subset V') \ar[ld] \\
& (U \subset U')
}
$$
参见《空间态射的进一步讨论》中的引理
\ref{spaces-more-morphisms-lemma-deform-property-fp-over-ft}。

\medskip\noindent
在情形 (\ref{stacks-more-morphisms-item-unramified-fp-over-ft}) 与
(\ref{stacks-more-morphisms-item-etale-fp-over-ft}) 中，首先由对 $f$ 或 $f'$ 的假设看出，
$f$ 与 $f'$ 都是代数叠的 DM 态射，参见引理
\ref{stacks-more-morphisms-lemma-thicken-property-morphisms}。然后选取代数空间 $U'$
及满光滑态射 $U' \to \mathcal{B}'$。
选取代数空间 $V'$ 及满光滑态射
$V' \to U' \times_{\mathcal{B}'} \mathcal{Y}'$。
选取代数空间 $W'$ 及满 \'etale(!) 态射
$W' \to V' \times_{\mathcal{Y}'} \mathcal{X}'$。
令 $U, V, W$ 分别为 $U', V', W'$ 沿
$\mathcal{B} \to \mathcal{B}'$ 的基变换。
于是 $W \to V \times_\mathcal{Y} \mathcal{X}$ 也满且 \'etale。
% P12-ERR-0069: see ../control/ERRATA_CANDIDATES.jsonl
因此，$f$（相应地，\ $f'$）具有该性质，等价于
$W' \to V'$（相应地，\ $W \to V$）具有该性质。
故可以把代数空间情形的引理应用于代数空间增厚的图表
$$
\xymatrix{
(W \subset W') \ar[rr] \ar[rd] & & (V \subset V') \ar[ld] \\
& (U \subset U')
}
$$
参见《空间态射的进一步讨论》中的引理
\ref{spaces-more-morphisms-lemma-deform-property-fp-over-ft}。

\medskip\noindent
在情形 (\ref{stacks-more-morphisms-item-isomorphism-fp-over-ft})、
(\ref{stacks-more-morphisms-item-open-immersion-fp-over-ft})、
(\ref{stacks-more-morphisms-item-monomorphism-fp-over-ft})、
(\ref{stacks-more-morphisms-item-finite-fp-over-ft}) 中，先由引理
\ref{stacks-more-morphisms-lemma-thicken-property-morphisms} 得到 $f$ 与 $f'$
都可由代数空间表示。因此，可以选取代数空间 $U'$ 及满光滑态射
$U' \to \mathcal{B}'$，再选取代数空间 $V'$ 及满光滑态射
$V' \to U' \times_{\mathcal{B}'} \mathcal{Y}'$；于是
$W' = V' \times_{\mathcal{Y}'} \mathcal{X}'$ 是代数空间。
令 $U, V, W$ 分别为 $U', V', W'$ 沿
$\mathcal{B} \to \mathcal{B}'$ 的基变换。
此时同样有 $W = V \times_\mathcal{Y} \mathcal{X}$。
接下来只需证明：$W' \to V'$ 是同构，\ 相应地为开浸入，\ 相应地为单态射，
\ 相应地为有限态射，
当且仅当 $W \to V$ 具有同一性质。参见《叠的性质》中的引理
\ref{stacks-properties-lemma-check-property-covering}。
由此应用上述代数空间的结果便得到结论。

\medskip\noindent
在情形 (\ref{stacks-more-morphisms-item-quasi-finite-fp-over-ft}) 中，首先由
《叠的态射》中的引理
\ref{stacks-morphisms-lemma-finite-type-permanence} 看到
$f$ 与 $f'$ 都是局部有限型的。另一方面，由引理
\ref{stacks-more-morphisms-lemma-thicken-property-morphisms}，$f$ 是拟-DM 的，
当且仅当 $f'$ 是拟-DM 的。要检验 $f$ 或 $f'$ 是否局部拟有限
（《叠的态射》中的定义
\ref{stacks-morphisms-definition-locally-quasi-finite}），
最后还需检验一个关于底层拓扑空间的条件；由证明第一段的讨论，
该条件对 $f$ 成立，当且仅当对 $f'$ 成立。
\end{proof}









\section{仿射概形的提升}
\label{stacks-more-morphisms-section-lifting-affines}

\noindent
考虑实线图表
$$
\xymatrix{
W \ar[d] \ar@{..>}[r] & W' \ar@{..>}[d] \\
\mathcal{X} \ar[r] & \mathcal{X}'
}
$$
其中 $\mathcal{X} \subset \mathcal{X}'$ 是代数叠的增厚，
$W$ 是仿射概形，且 $W \to \mathcal{X}$ 光滑。本节所研究的问题是：
能否找到 $W'$ 及虚线箭头，使得该方形为笛卡儿方形，
且 $W' \to \mathcal{X}'$ 光滑。一般情形下我们不知道答案，
但若 $\mathcal{X} \subset \mathcal{X}'$ 是一阶增厚，
将证明答案是肯定的。

\medskip\noindent
为研究这个问题，引入如下范畴。

\begin{remark}[提升范畴]
\label{stacks-more-morphisms-remark-gerbe-of-lifts}
考虑图表
$$
\xymatrix{
W \ar[d]_x \\
\mathcal{X} \ar[r] & \mathcal{X}'
}
$$
其中 $\mathcal{X} \subset \mathcal{X}'$ 是代数叠的增厚，
$W$ 是代数空间，且 $W \to \mathcal{X}$ 光滑。
如下构造范畴 $\mathcal{C}$ 及函子
$$
p : \mathcal{C} \longrightarrow W_{spaces, \etale}
$$
（记号见《空间的性质》中的定义
\ref{spaces-properties-definition-spaces-etale-site}）。
$\mathcal{C}$ 的一个对象是构成交换图表
\begin{equation}
\label{stacks-more-morphisms-equation-object}
\vcenter{
\xymatrix{
U \ar[d]_a \ar[r]_i & U' \ar[dd]^{x'} \\
W \ar[d]_x & \\
\mathcal{X} \ar[r] & \mathcal{X}'
}
}
\end{equation}
的系统 $(U, U', a, i, x', \alpha)$，其交换性由 $2$-态射
$\alpha : x \circ a \to x' \circ i$ 见证，并且
$U$ 与 $U'$ 是代数空间，$a : U \to W$ 是 \'etale 态射，
$x' : U' \to \mathcal{X}'$ 是光滑态射，且
$U = \mathcal{X} \times_{\mathcal{X}'} U'$。
特别地，$U \subset U'$ 是一个增厚。
一个态射
$$
(U, U', a, i, x', \alpha) \to (V, V', b, j, y', \beta)
$$
由 $(f, f', \gamma)$ 给出，其中 $f : U \to V$ 是 $W$ 上的态射，
$f' : U' \to V'$ 是一个态射，其在 $U$ 上的限制给出 $f$，并且
$\gamma : x' \circ f' \to y'$ 是见证下图右侧三角形交换性的
$2$-态射：
% P12-ERR-0070: see ../control/ERRATA_CANDIDATES.jsonl
% P12-ERR-0071: see ../control/ERRATA_CANDIDATES.jsonl
\begin{equation}
\label{stacks-more-morphisms-equation-morphism}
\vcenter{
\xymatrix{
& V \ar[ld]_f \ar[ldd]^b \ar[rr]_j & & V' \ar[ld]_{f'} \ar[lddd]^{y'} \\
U \ar[d]_a \ar[rr]_i & & U' \ar[dd]_{x'} \\
W \ar[d]_x & \\
\mathcal{X} \ar[rr] & & \mathcal{X}'
}
}
\end{equation}
最后，要求 $\gamma$ 与 $\alpha$、$\beta$ 相容：按照《范畴论》第
\ref{categories-section-formal-cat-cat} 节及第
\ref{categories-section-2-categories} 节中的 $2$-范畴演算，这写成
$$
\beta = (\gamma \star \text{id}_j) \circ (\alpha \star \text{id}_f)
$$
（更简洁地说：$\beta = j^*\gamma \circ f^*\alpha$）。
另一种表述是：对象是带有若干附加性质的交换图表
(\ref{stacks-more-morphisms-equation-object})，态射则是交换图表
(\ref{stacks-more-morphisms-equation-morphism})；这些图表位于《叠的性质》中的注
\ref{stacks-properties-remark-representable-over} 所引入的范畴
$\textit{Spaces}/\mathcal{X}'$ 中。由此清楚可见，$\mathcal{C}$ 是范畴，
而把 $(U, U', a, i, x', \alpha)$ 送到 $a : U \to W$ 的规则
$p : \mathcal{C} \to W_{spaces, \etale}$ 是函子。
\end{remark}

\begin{lemma}
\label{stacks-more-morphisms-lemma-morphisms-lifts-etale}
对任意态射 (\ref{stacks-more-morphisms-equation-morphism})，映射
$f' : V' \to U'$ 都是 \'etale 的。
\end{lemma}

\begin{proof}
事实上，$f : V \to U$ 作为 $W_{spaces, \etale}$ 中的态射是
\'etale 的；又因为 $U' \to \mathcal{X}'$ 与
$V' \to \mathcal{X}'$ 光滑，并且
$U = \mathcal{X} \times_{\mathcal{X}'} U'$、
$V = \mathcal{X} \times_{\mathcal{X}'} V'$，
所以可以应用引理 \ref{stacks-more-morphisms-lemma-deform-property-fp-over-ft}。
\end{proof}

\begin{lemma}
\label{stacks-more-morphisms-lemma-gerbe-of-lifts-fibred}
注 \ref{stacks-more-morphisms-remark-gerbe-of-lifts} 中构造的范畴
$p : \mathcal{C} \to W_{spaces, \etale}$ 是群胚纤维化的。
\end{lemma}

\begin{proof}
我们断言 $p$ 的各纤维范畴都是群胚。
若 (\ref{stacks-more-morphisms-equation-morphism}) 中的 $(f, f', \gamma')$ 是态射，
且 $f : U \to V$ 为同构，则由引理
\ref{stacks-more-morphisms-lemma-deform-property-fp-over-ft}，$f'$ 也是同构，
因而 $(f, f', \gamma')$ 是同构。

\medskip\noindent
考虑 $W_{spaces, \etale}$ 中的态射 $f : V \to U$，
以及 $\mathcal{C}$ 中位于 $U$ 上的对象
$\xi = (U, U', a, i, x', \alpha)$。我们构造位于 $V$ 上的“拉回”
$f^*\xi$。令 $b = a \circ f$。令 $f' : V' \to U'$ 为
其在 $V$ 上的限制是 $f$ 的 \'etale 态射（《空间态射的进一步讨论》中的
引理 \ref{spaces-more-morphisms-lemma-topological-invariance}）。
以 $j : V \to V'$ 记相应的增厚。令 $y' = x' \circ f'$，并令
$\gamma = \text{id} : x' \circ f' \to y'$。置
$$
\beta = \alpha \star \text{id}_f :
x \circ b = x \circ a \circ f \to
x' \circ i \circ f = x' \circ f' \circ j = y' \circ j
$$
显然，$(f, f', \gamma) : (V, V', b, j, y', \beta) \to
(U, U', a, i, x', \alpha)$ 是 (\ref{stacks-more-morphisms-equation-morphism}) 中的态射。
如此构造的态射 $(f, f', \gamma)$ 是强笛卡儿的
（《范畴论》中的定义 \ref{categories-definition-cartesian-over-C}）。
略去详细证明，但其根本原因如下。给定 $\mathcal{C}$ 中的态射
$(g, g', \epsilon) : (Y, Y', c, k, z', \delta) \to
(U, U', a, i, x', \alpha)$，若 $g$ 对某个 $h : Y \to V$ 有
$g = f \circ h$，则由《空间态射的进一步讨论》中的引理
\ref{spaces-more-morphisms-lemma-topological-invariance}，
得到唯一分解 $g' = f' \circ h'$；此后可构造所需的 $\zeta$，使得
$(h, h', \zeta) :  (Y, Y', c, k, z', \delta) \to
(V, V', b, j, y', \beta)$ 是 $\mathcal{C}$ 的态射，并满足
$(g, g', \epsilon) = (f, f', \gamma) \circ (h, h', \zeta)$。

\medskip\noindent
因此，$p : \mathcal{C} \to W_\etale$ 是纤维化范畴
（《范畴论》中的定义 \ref{categories-definition-fibred-category}）。
结合上面已经看到的各纤维范畴都是群胚这一事实，再由《范畴论》中的引理
\ref{categories-lemma-fibred-groupoids}，得到
$p : \mathcal{C} \to W_\etale$ 是群胚纤维化的。
\end{proof}

\begin{lemma}
\label{stacks-more-morphisms-lemma-gerbe-of-lifts-stack}
注 \ref{stacks-more-morphisms-remark-gerbe-of-lifts} 中构造的范畴
$p : \mathcal{C} \to W_{spaces, \etale}$ 是一个群胚叠。
\end{lemma}

\begin{proof}
由引理 \ref{stacks-more-morphisms-lemma-gerbe-of-lifts-fibred}，可见《叠》中的定义
\ref{stacks-definition-stack-in-groupoids} 的第一个条件成立。
按照惯例，我们检验对象的下降，而把态射的下降留给读者。
因此，设在 $W_{spaces, \etale}$ 中给定 $a : U \to W$、
$W_{spaces, \etale}$ 中的覆盖 $\{U_k \to U\}_{k \in K}$、
$\mathcal{C}$ 中位于 $U_k$ 上的对象
$\xi_k = (U_k, U'_k, a_k, i_k, x'_k, \alpha_k)$，以及满足余圈条件的
限制之间的态射
$$
\varphi_{kk'} = (f_{kk'}, f'_{kk'}, \gamma_{kk'}) :
\xi_k|_{U_k \times_U U_{k'}} \to
\xi_{k'}|_{U_k \times_U U_{k'}}
$$
为了证明有效性，可以先细化覆盖。因此，可以假设每个 $U_k$
都是概形（如果愿意，甚至可假设为仿射概形）。写成
% P12-ERR-0072: see ../control/ERRATA_CANDIDATES.jsonl
$$
\xi_k|_{U_k \times_U U_{k'}} =
(U_k \times_U U_{k'}, U'_{kk'}, a_{kk'}, x'_{kk'}, \alpha_{kk'})
$$
于是，把 $\mathcal{C}$ 中态射
$\xi_k|_{U_k \times_U U_{k'}} \to \xi_k$ 的第二个分量取出，
便得到 \'etale 态射
$s_{kk'} : U'_{kk'} \to U'_k$
（由引理 \ref{stacks-more-morphisms-lemma-morphisms-lifts-etale}）。
类似地，考察复合的第二个分量
$$
\xi_k|_{U_k \times_U U_{k'}} \xrightarrow{\varphi_{kk'}}
\xi_{k'}|_{U_k \times_U U_{k'}} \to \xi_{k'}
$$
便得到 \'etale 态射 $t_{kk'} : U'_{kk'} \to U'_{k'}$。
我们断言
$$
j :
\coprod\nolimits_{(k, k') \in K \times K} U'_{kk'}
\xrightarrow{(\coprod s_{kk'}, \coprod t_{kk'})}
(\coprod\nolimits_{k \in K} U'_k) \times (\coprod\nolimits_{k \in K} U'_k)
$$
是一个 \'etale 等价关系。首先，已经看到所示态射的两个分量
$s, t$ 都是 \'etale 的。用
$(\coprod U_k) \times (\coprod U_k) \to
(\coprod U'_k) \times (\coprod U'_k)$
对态射 $j$ 作基变换，得到单态射
$$
\coprod\nolimits_{(k, k') \in K \times K} U_k \times_U U_{k'}
\longrightarrow
(\coprod\nolimits_{k \in K} U_k) \times (\coprod\nolimits_{k \in K} U_k)
$$
因此，由《态射的进一步讨论》中的引理
\ref{more-morphisms-lemma-properties-that-extend-over-thickenings}，
$j$ 是单态射。最后，关系 $j$ 的对称性来自
$\varphi_{kk'}^{-1}$ 是 $\varphi_{k'k}$ 的“翻转”（参见《叠》中的注
\ref{stacks-remarks-definition-descent-datum}），而传递性来自余圈条件
（略去细节）。于是，$\coprod U'_k$ 关于 $j$ 的商是代数空间 $U'$
（《空间》中的定理 \ref{spaces-theorem-presentation}）。
上面已经说明存在增厚 $i : U \to U'$，因为 $j$ 在
$\coprod U_k$ 上的限制给出
$(\coprod U_k) \times_U (\coprod U_k)$。
最后，暂时把 $1$-态射 $x'_k : U'_k \to \mathcal{X}'$
视为叠 $\mathcal{X}'$ 在 $U'_k$ 上的对象，则可见这些对象配备了
相对于 \'etale 覆盖 $\{U'_k \to U'\}$ 的下降数据；该数据由
$\mathcal{C}$ 中态射 $\varphi_{kk'}$ 的第三个分量
$\gamma_{kk'}$ 给出。由于 $\mathcal{X}'$ 是叠，该下降数据有效；
翻译回原语言，得到 $1$-态射 $x' : U' \to \mathcal{X}'$，
使各复合 $U'_k \to U' \to \mathcal{X}'$ 配备到 $x'_k$ 的同构，
并与 $\gamma_{kk'}$ 相容。这意味着态射
$\alpha_k : x \circ a_k \to x'_k \circ i_k$ 粘合成态射
$\alpha : x \circ a \to x' \circ i$。于是
$\xi = (U, U', a, i, x', \alpha)$ 就是所求的位于 $U$ 上的对象。
\end{proof}

\begin{lemma}
\label{stacks-more-morphisms-lemma-etale-local-lifts}
设 $\mathcal{X} \subset \mathcal{X}'$ 是代数叠的一个增厚。
设 $W$ 是代数空间，且 $W \to \mathcal{X}$ 是光滑态射。
存在一个 \'etale 覆盖 $\{W_i \to W\}_{i \in I}$，并且对每个 $i$，
存在笛卡儿图表
$$
\xymatrix{
W_i \ar[r] \ar[d] & W_i' \ar[d] \\
\mathcal{X} \ar[r] & \mathcal{X}'
}
$$
其中 $W_i' \to \mathcal{X}'$ 是光滑的。
\end{lemma}

\begin{proof}
选取一个概形 $U'$ 以及一个满光滑态射 $U' \to \mathcal{X}'$。
像往常一样，置 $U = \mathcal{X} \times_{\mathcal{X}'} U'$。于是
$U \to \mathcal{X}$ 是满光滑态射。因此，基变换
$$
V = W \times_{\mathcal{X}} U \longrightarrow W
$$
是代数空间的满光滑态射。由《空间上的拓扑》中的引理
\ref{spaces-topologies-lemma-etale-dominates-smooth}，
可以找到一个 \'etale 覆盖 $\{W_i \to W\}$，使得
$W_i \to W$ 经由 $V \to W$ 分解。再用仿射概形覆盖 $W_i$
（《空间的性质》中的引理
\ref{spaces-properties-lemma-cover-by-union-affines}），
便可假设每个 $W_i$ 都是仿射的。我们可以并且确实以 $W_i$ 代替
$W$，从而归结到下一段讨论的情形。

\medskip\noindent
假设 $W$ 是仿射的，并且给定态射 $W \to \mathcal{X}$ 经由 $U$
分解。图示为
$$
W \xrightarrow{i} U \to \mathcal{X}
$$
由于 $W$ 和 $U$ 在 $\mathcal{X}$ 上都是光滑的，可见 $i$ 是局部有限型的
（《叠的态射》中的引理
\ref{stacks-morphisms-lemma-finite-type-permanence}）。
以 $\mathbf{A}^n_U$ 代替 $U$ 后，可以假设 $i$ 是浸入；参见
《态射》中的引理 \ref{morphisms-lemma-quasi-affine-finite-type-over-S}。
由《叠的态射》中的引理
\ref{stacks-morphisms-lemma-lci-permanence}，态射 $i$ 是局部完全交。
因此，$i$ 是 Koszul-正则浸入（定义见《除子》中的定义
\ref{divisors-definition-regular-immersion}）；这由《态射的进一步讨论》中的引理
\ref{more-morphisms-lemma-lci} 得出。

\medskip\noindent
我们仍可用仿射开覆盖代替 $W$。对每个点 $w \in W$，可以选取
仿射开集 $U'_w \subset U'$，使得若 $U_w \subset U$ 是对应的仿射开集，
则 $w \in i^{-1}(U_w)$，并且 $i^{-1}(U_w) \to U_w$ 是由
Koszul-正则序列
$f_1, \ldots, f_r \in \Gamma(U_w, \mathcal{O}_{U_w})$
截出的闭浸入。这由 Koszul-正则浸入的定义和《除子》中的引理
\ref{divisors-lemma-regular-ideal-sheaf-scheme} 得出。
置 $W_w = i^{-1}(U_w)$；这是 $w \in W$ 的一个仿射开邻域。
选取 $f_1, \ldots, f_r$ 的提升
$f'_1, \ldots, f'_r \in \Gamma(U'_w, \mathcal{O}_{U'_w})$。
这是可行的，因为 $U_w \to U'_w$ 是仿射概形的闭浸入。
令 $W'_w \subset U'_w$ 为由 $f'_1, \ldots, f'_r$ 截出的闭子概形。
我们断言 $W'_w \to \mathcal{X}'$ 是光滑的。该断言完成证明，因为依构造
$W_w = \mathcal{X} \times_{\mathcal{X}'} W'_w$。

\medskip\noindent
为检验该断言，只须对每个在 $\mathcal{X}'$ 上光滑的仿射概形 $X'$，
检验基变换 $W'_w \times_{\mathcal{X}'} X' \to X'$ 是光滑的。
选取一个 \'etale 态射
$$
Y' \to U'_w \times_{\mathcal{X}'} X'
$$
其中 $Y'$ 是仿射的。由于这些态射的像覆盖
$U'_w \times_{\mathcal{X}'} X'$，只须证明 $Y'$ 中由
$f'_1, \ldots, f'_r$ 截出的闭子概形 $Z'$ 在 $X'$ 上光滑。图示为
$$
\xymatrix{
Z' \ar[r] \ar[d] & Y' \ar[d] \\
W'_w \times_{\mathcal{X}'} X' \ar[d] \ar[r] &
U'_w \times_{\mathcal{X}'} X' \ar[d] \ar[r] & X' \\
W'_w = V(f'_1, \ldots, f'_r) \ar[r] & U'_w
}
$$
置 $X = \mathcal{X} \times_{\mathcal{X}'} X'$、
$Y = X \times_{X'} Y' = \mathcal{X} \times_{\mathcal{X}'} Y'$，以及
$Z = Y \times_{Y'} Z' = X \times_{X'} Z' =
\mathcal{X} \times_{\mathcal{X}'} Z'$。
于是 $(Z \subset Z') \to (Y \subset Y') \subset (X \subset X')$
是仿射概形增厚之间的（笛卡儿）态射，并且已知 $Z \to X$ 和
$Y' \to X'$ 都是光滑的。最后，由于 $Y' \to U'_w$ 光滑从而平坦，
函数序列 $f'_1, \ldots, f'_r$ 由《代数的进一步讨论》中的引理
\ref{more-algebra-lemma-koszul-regular-flat-base-change}
映成 $\Gamma(Y', \mathcal{O}_{Y'})$ 中的 Koszul-正则序列。
由《代数的进一步讨论》中的引理
\ref{more-algebra-lemma-cut-by-koszul}
（并利用 Koszul-正则序列是拟正则序列这一事实；参见《代数的进一步讨论》中的引理
\ref{more-algebra-lemma-regular-koszul-regular}、
\ref{more-algebra-lemma-koszul-regular-H1-regular} 和
\ref{more-algebra-lemma-H1-regular-quasi-regular}），
得出 $Z' \to X'$ 如所需是光滑的。
\end{proof}

\begin{lemma}
\label{stacks-more-morphisms-lemma-etale-local-lifts-isomorphic}
设 $\mathcal{X} \subset \mathcal{X}'$ 是代数叠的一个增厚。考虑交换图表
$$
\xymatrix{
W'' \ar[d]_{x''} & W \ar[l] \ar[r] \ar[d]_x & W' \ar[d]^{x'} \\
\mathcal{X}' & \mathcal{X} \ar[l] \ar[r] & \mathcal{X}'
}
$$
其中两个方形都是笛卡儿的，$W', W, W''$ 是代数空间，且竖直箭头
都是光滑的。那么存在
\begin{enumerate}
\item 一个 \'etale 覆盖 $\{f'_k : W'_k \to W'\}_{k \in K}$，
\item \'etale 态射 $f''_k : W'_k \to W''$，以及
\item $2$-态射 $\gamma_k : x'' \circ f''_k \to x' \circ f'_k$，
\end{enumerate}
使得 (a) $(f'_k)^{-1}(W) = (f''_k)^{-1}(W)$，(b)
$f'_k|_{(f'_k)^{-1}(W)} = f''_k|_{(f''_k)^{-1}(W)}$，并且
(c) 把 $\gamma_k$ 拉回到 (a) 中的闭子概形所得的态射，
与初始图表在 $W$ 上的交换性给出的 $2$-态射一致。
\end{lemma}

\begin{proof}
把给定的增厚记作 $i : W \to W'$ 和 $i'' : W \to W''$。
引理陈述中图表的交换性意味着存在一个 $2$-态射
% P12-ERR-0073: see ../control/ERRATA_CANDIDATES.jsonl
$\delta : x' \circ i' \to x'' \circ i''$。
这就是陈述的 (c) 中所指的 $2$-态射。考虑代数空间
$$
I' = W' \times_{x', \mathcal{X}', x''} W''
$$
以及投影 $p' : I' \to W'$ 和 $q' : I' \to W''$。
注意，存在一个“万有的”$2$-态射
$\gamma : x' \circ p' \to x'' \circ q'$（稍后将用到它）。
$\delta$ 的选择定义了态射
$$
\xymatrix{
W \ar[rr]_\delta & & I' \ar[ld]^{p'} \ar[rd]_{q'} \\
& W' & & W''
}
$$
使得复合 $W \to I' \to W'$ 和 $W \to I' \to W''$ 分别为
$i : W \to W'$ 和
% P12-ERR-0074: see ../control/ERRATA_CANDIDATES.jsonl
$i' : W \to W''$。由于 $x''$ 是光滑的，态射 $p' : I' \to W'$
作为 $x''$ 的基变换是光滑的。

\medskip\noindent
假设可以找到一个 \'etale 覆盖 $\{f'_k : W'_k \to W'\}$ 和态射
$\delta_k : W'_k \to I'$，使得 $\delta_k$ 在
% P12-ERR-0075: see ../control/ERRATA_CANDIDATES.jsonl
$W_k = (f'_k)^{-1}$ 上的限制等于 $\delta \circ f_k$，
其中 $f_k = f'_k|_{W_k}$。图示为
$$
\xymatrix{
W_k \ar[r]^{f_k} \ar[d] & W \ar[r]^\delta & I' \ar[d]^{p'} \\
W'_k \ar[rr]^{f'_k} \ar[rru]^{\delta_k} & & W'
}
$$
换言之，我们希望在可能以一个 \'etale 覆盖代替 $W'$ 后，
把 $p'$ 的给定截面 $\delta : W \to I'$ 延拓为 $W'$ 上的截面。

\medskip\noindent
若此事成立，则可置 $f''_k = q' \circ \delta_k$ 和
$\gamma_k = \gamma \star \text{id}_{\delta_k}$（更简洁地说，
$\gamma_k = \delta_k^*\gamma$）。也就是说，剩下只须证明
% P12-ERR-0076: see ../control/ERRATA_CANDIDATES.jsonl
态射 $f''_k$ 是 \'etale 的。依构造，态射 $x' \circ p'$ 与
$x'' \circ q'$ 是 $2$-同构的。因此，$x'' \circ f''_k$ 与
$x' \circ f'_k$ 是 $2$-同构的。由此，复合
$$
W'_k \xrightarrow{f''_k} W'' \xrightarrow{x''} \mathcal{X}'
$$
是光滑的，因为 $x' \circ f'_k$ 是光滑的。由于 $f_k$ 是 \'etale 的，
由引理 \ref{stacks-more-morphisms-lemma-deform-property-fp-over-ft} 得出 $f''_k$ 是 \'etale 的。

\medskip\noindent
如果该增厚是一阶增厚，那么可以选取任何使 $W_k'$ 为仿射的
% P12-ERR-0077: see ../control/ERRATA_CANDIDATES.jsonl
\'etale 覆盖 $\{W'_k \to W'\}$。事实上，由于 $p'$ 是光滑的，
无穷小提升判据表明 $p'$ 是形式光滑的（《空间态射的进一步讨论》中的引理
\ref{spaces-more-morphisms-lemma-smooth-formally-smooth}）。
由于 $W_k$ 是仿射的，并且 $W_k \to W'_k$ 是一阶增厚
（它是 $\mathcal{X} \to \mathcal{X}'$ 的基变换；参见引理
\ref{stacks-more-morphisms-lemma-base-change-thickening}），便得到所需的 $\delta_k$。

\medskip\noindent
在一般情形下，覆盖和态射 $\delta_k$ 的存在性由《空间态射的进一步讨论》中的引理
\ref{spaces-more-morphisms-lemma-smooth-strong-lift} 得出。
\end{proof}

\begin{lemma}
\label{stacks-more-morphisms-lemma-gerbe-of-lifts}
注 \ref{stacks-more-morphisms-remark-gerbe-of-lifts} 中构造的范畴
$p : \mathcal{C} \to W_{spaces, \etale}$ 是一个胚。
\end{lemma}

\begin{proof}
在引理 \ref{stacks-more-morphisms-lemma-gerbe-of-lifts-stack} 中已经看到，它是一个群胚叠。
因此只须检验《叠》中的定义 \ref{stacks-definition-gerbe}
的条件 (2) 和 (3)。条件 (2) 由引理
\ref{stacks-more-morphisms-lemma-etale-local-lifts} 得出。条件 (3) 由引理
\ref{stacks-more-morphisms-lemma-etale-local-lifts-isomorphic} 得出。
\end{proof}

\begin{lemma}
\label{stacks-more-morphisms-lemma-gerbe-of-lifts-first-order}
在注 \ref{stacks-more-morphisms-remark-gerbe-of-lifts} 中，假设
$\mathcal{X} \subset \mathcal{X}'$ 是一阶增厚。那么
\begin{enumerate}
\item 注 \ref{stacks-more-morphisms-remark-gerbe-of-lifts} 中构造的胚
$p : \mathcal{C} \to W_{spaces, \etale}$ 的对象的自同构层都是阿贝尔的；并且
\item 《叠》中的引理 \ref{stacks-lemma-gerbe-abelian-auts}
所构造的群层 $\mathcal{G}$ 是拟凝聚 $\mathcal{O}_W$-模。
\end{enumerate}
\end{lemma}

\begin{proof}
我们将同时证明两个断言。具体而言，给定对象
$\xi = (U, U', a, i, x', \alpha)$，我们将在
$U_{spaces, \etale}$ 上赋予 $\mathit{Aut}(\xi)$ 一个拟凝聚
$\mathcal{O}_U$-模结构，并证明该结构与拉回相容。
由层的粘合（《格址》第 \ref{sites-section-glueing-sheaves} 节）、
《叠》中的引理 \ref{stacks-lemma-gerbe-abelian-auts} 的证明中
把 $\mathcal{G}$ 构造为自同构层 $\mathit{Aut}(\xi)$ 的粘合，
以及只须在取一个 \'etale 覆盖之后检验模的拟凝聚性这一事实
（《空间的性质》中的引理
\ref{spaces-properties-lemma-characterize-quasi-coherent}），
这便足够了。

\medskip\noindent
我们将采用引理 \ref{stacks-more-morphisms-lemma-etale-local-lifts-isomorphic}
的证明中所用的相同方法来描述层 $\mathit{Aut}(\xi)$。
考虑代数空间
$$
I' = U' \times_{x', \mathcal{X}', x'} U'
$$
及其投影 $p' : I' \to U'$ 和 $q' : I' \to U'$。
在 $I'$ 上有一个万有 $2$-态射
$\gamma : x' \circ p' \to x' \circ q'$。
恒等态射 $x' \to x'$ 定义了对角态射
$$
\xymatrix{
U' \ar[rr]_{\Delta'} & & I' \ar[ld]^{p'} \ar[rd]_{q'} \\
& U' & & U'
}
$$
使复合 $U' \to I' \to U'$ 和 $U' \to I' \to U'$
都是恒等态射。我们用
$U, I, p, q, \Delta$ 表示 $U', I', p', q', \Delta'$ 到
$\mathcal{X}$ 的基变换。
% P12-ERR-0078: see ../control/ERRATA_CANDIDATES.jsonl
由于 $W' \to \mathcal{X}'$ 是光滑的，可见 $p' : I' \to U'$
作为基变换是光滑的。

\medskip\noindent
$\mathit{Aut}(\xi)$ 在 $U$ 上的一个截面是态射
$\delta' : U' \to I'$，它满足 $\delta'|_U = \Delta$ 且
$p' \circ \delta' = \text{id}_{U'}$。明确地说，
$(\text{id}_U, q' \circ \delta', (\delta')^*\gamma) : \xi \to \xi$
是相应自同构的公式。更一般地，若 $f : V \to U$ 是 \'etale 态射，
则存在增厚 $j : V \to V'$ 和 \'etale 态射 $f' : V' \to U'$，
后者在 $V$ 上的限制为 $f$；而 $f^*\xi$ 对应于
$(V, V', a \circ f, j, x' \circ f', f^*\alpha)$，参见引理
\ref{stacks-more-morphisms-lemma-gerbe-of-lifts-fibred} 的证明。
% P12-ERR-0079: see ../control/ERRATA_CANDIDATES.jsonl
$\mathit{Aut}(\xi)$ 在 $V$ 上的一个截面是态射
$\delta' : V' \to I'$，它满足 $\delta'|_V = \Delta \circ f$ 且
$p' \circ \delta' = f'$\footnote{相应自同构的一个公式是
$(\text{id}_V, h', (\delta')^*\gamma)$。这里
$h' : V' \to V'$ 是唯一的（同构）态射，满足
$h'|_V = \text{id}_V$，并使
$$
\xymatrix{
V' \ar[r]_{h'} \ar[rd]_{q' \circ \delta'} & V' \ar[d]^{f'} \\
& U'
}
$$
交换。$h'$ 的唯一性和存在性来自 \'etale 格址的拓扑不变性；
% P12-ERR-0080: see ../control/ERRATA_CANDIDATES.jsonl
参见《空间态射的进一步讨论》中的定理
\ref{spaces-more-morphisms-theorem-topological-invariance}。
读者可能觉得，我们反而应该考察态射
$\delta'' : V' \to V' \times_{\mathcal{X}'} V'$，使得
$\delta'' \circ j = \Delta_{V'/\mathcal{X}'}$ 且
$\text{pr}_1 \circ \delta'' = \text{id}_{V'}$。
这样做也可以：由于 $V' \times_{\mathcal{X}'} V' \to I'$ 是 \'etale 的，
同一拓扑不变性表明，把 $\delta''$ 送到
$\delta' = (V' \times_{\mathcal{X}'} V' \to I') \circ \delta''$，
给出了两组态射之间的双射。}。

\medskip\noindent
由此可知，$\mathit{Aut}(\xi)$ 作为集合层，与《空间态射的进一步讨论》中的注
\ref{spaces-more-morphisms-remark-another-special-case}
为增厚 $(U \subset U')$ 和 $(I \subset I')$ 所定义的层一致；
它们通过 $\text{id}_{U'}$ 和 $p'$ 位于 $(U \subset U')$ 上。
对角态射 $\Delta'$ 是该层的一个截面；用《空间态射的进一步讨论》中的引理
\ref{spaces-more-morphisms-lemma-action-sheaf} 作用于这个截面，得到同构
\begin{equation}
\label{stacks-more-morphisms-equation-isomorphism}
\SheafHom_{\mathcal{O}_U}(\Delta^*\Omega_{I/U}, \mathcal{C}_{U/U'})
\longrightarrow
\mathit{Aut}(\xi)
\end{equation}
在 $U_{spaces, \etale}$ 上成立。
% P12-ERR-0081: see ../control/ERRATA_CANDIDATES.jsonl
还有三件事需要检验：
\begin{enumerate}
\item (\ref{stacks-more-morphisms-equation-isomorphism}) 的构造与 \'etale 局部化交换；
\item $\SheafHom_{\mathcal{O}_U}(\Delta^*\Omega_{I/U}, \mathcal{C}_{U/U'})$
是 $U$ 上的拟凝聚模；
\item $\mathit{Aut}(\xi)$ 中的复合对应于这个拟凝聚模中截面的加法。
\end{enumerate}
我们将依次检验这些事项。

\medskip\noindent
为说明 (1)，必须证明：若 $f : V \to U$ 是 \'etale 的，则在 $U$ 上
用 $\xi$ 构造的 (\ref{stacks-more-morphisms-equation-isomorphism}) 限制为映射
(\ref{stacks-more-morphisms-equation-isomorphism})
$$
\SheafHom_{\mathcal{O}_V}(
\Delta_V^*\Omega_{V \times_\mathcal{X} V/V}, \mathcal{C}_{V/V'}) \to
\mathit{Aut}(\xi|_V)
$$
其中后一个映射是在 $V$ 上用 $\xi|_V$ 构造的，并定义于
$V_{spaces, \etale}$ 上。
这由上面脚注中的讨论和《空间态射的进一步讨论》中的引理
\ref{spaces-more-morphisms-lemma-action-by-derivations-etale-localization}
得出。

\medskip\noindent
(2) 的证明。由于 $p'$ 是光滑的，态射 $I \to U$ 是光滑的，
因而相对微分模 $\Omega_{I/U}$ 是有限局部自由的
（《空间态射的进一步讨论》中的引理
\ref{spaces-more-morphisms-lemma-smooth-omega-finite-locally-free}）。
另一方面，$\mathcal{C}_{U/U'}$ 是拟凝聚的
（《空间态射的进一步讨论》中的定义
\ref{spaces-more-morphisms-definition-conormal-sheaf}）。
由《空间的性质》中的引理
\ref{spaces-properties-lemma-properties-quasi-coherent} 即得结论。

\medskip\noindent
(3) 的证明。存在态射 $c' : I' \times_{p', U', q'} I' \to I'$，
使得 $(U', I', p', q', c')$ 是以 $\Delta'$ 为恒等元的代数空间群胚。
例如参见《代数叠》中的引理
\ref{algebraic-lemma-map-space-into-stack}。
$\mathit{Aut}(\xi)$ 中的复合由态射 $c'$ 如下诱导。假设有两个态射
$$
\delta'_1, \delta'_2 : U' \longrightarrow I'
$$
如上对应于 $\mathit{Aut}(\xi)$ 在 $U$ 上的截面；换言之，
% P12-ERR-0082: see ../control/ERRATA_CANDIDATES.jsonl
有 $\delta'_i|U = \Delta_U$ 且
$p' \circ \delta'_i = \text{id}_{U'}$。那么 $\mathit{Aut}(\xi)$ 中的复合为
$$
\delta'_1 \circ \delta'_2 = c'(\delta'_1 \circ q' \circ \delta'_2, \delta'_2)
$$
我们略去详细检验\footnote{读者可以立即看出，必须用
$q' \circ \delta'_2$ 预复合 $\delta'_1$，才能得到纤维积
$I' \times_{p', U', q'} I'$ 的一个良定义的 $U'$-值点。}。
因此，正处于《空间中群胚的进一步讨论》第
\ref{spaces-more-groupoids-section-groupoid-sections} 节所描述的情形，
而所需结论由《空间中群胚的进一步讨论》中的引理
\ref{spaces-more-groupoids-lemma-composition-is-addition} 得出。
\end{proof}

\begin{proposition}[Emerton]
\label{stacks-more-morphisms-proposition-affine-smooth-lift-to-first-order}
\begin{reference}
Matthew Emerton 于 2016 年 4 月 27 日的电子邮件。
\end{reference}
设 $\mathcal{X} \subset \mathcal{X}'$ 是代数叠的一阶增厚。
设 $W$ 是仿射概形，且 $W \to \mathcal{X}$ 是光滑态射。那么存在笛卡儿图表
$$
\xymatrix{
W \ar[d] \ar[r] & W' \ar[d] \\
\mathcal{X} \ar[r] & \mathcal{X}'
}
$$
其中 $W' \to \mathcal{X}'$ 是光滑的，并且 $W'$ 是仿射的。
\end{proposition}

\begin{proof}
考虑注 \ref{stacks-more-morphisms-remark-gerbe-of-lifts} 中引入的范畴
$p : \mathcal{C} \to W_{spaces, \etale}$。该命题断言存在
$\mathcal{C}$ 中位于 $W$ 上的对象。具体而言，如果有这样的对象
$(W, W', a, i, y', \alpha)$，则
$W = \mathcal{X} \times_{\mathcal{X}'} W'$。
因而 $W \to W'$ 是代数空间的增厚，所以由《空间态射的进一步讨论》中的引理
\ref{spaces-more-morphisms-lemma-thickening-scheme}
和《态射的进一步讨论》中的引理
\ref{more-morphisms-lemma-thickening-affine-scheme}，$W'$ 是仿射的。

\medskip\noindent
引理 \ref{stacks-more-morphisms-lemma-gerbe-of-lifts} 告诉我们，$\mathcal{C}$ 是
$W_{spaces, \etale}$ 上的胚。这意味着可以在 \'etale 局部找到解，
并且这些局部解在 \'etale 局部同构；这一部分不需要增厚是一阶的假设。
由引理 \ref{stacks-more-morphisms-lemma-gerbe-of-lifts-first-order}，该胚的对象的自同构层
是阿贝尔的，并且粘合成 $W_{spaces, \etale}$ 上的拟凝聚模
$\mathcal{G}$。我们将检验《格址上的上同调》中的引理
\ref{sites-cohomology-lemma-existence} 的条件 (1) 和 (2)，
以推出存在 $\mathcal{C}$ 中位于 $W$ 上的对象。
条件 (1) 成立：其中每个 $W_i$ 都是仿射的 \'etale 覆盖
$\{W_i \to W\}$ 在所有覆盖组成的集合中共终。
对这样的覆盖，$W_i$ 和 $W_i \times_W W_j$ 都是仿射的，而
$H^1(W_i, \mathcal{G})$ 和 $H^1(W_i \times_W W_j, \mathcal{G})$
都为零：例如，由《空间的上同调》中的命题
\ref{spaces-cohomology-proposition-vanishing}，
仿射代数空间上拟凝聚模的上同调为零。
最后，条件 (2) 是对于拟凝聚层 $\mathcal{G}$ 有
$H^2(W, \mathcal{G}) = 0$；这同样由《空间的上同调》中的命题
\ref{spaces-cohomology-proposition-vanishing} 得出。证明完毕。
\end{proof}

\section{无穷小形变}
\label{stacks-more-morphisms-section-inf}

\noindent
我们继续《Artin 公理》第 \ref{artin-section-inf} 节中的讨论。

\begin{lemma}
\label{stacks-more-morphisms-lemma-inf-quasi-coherent}
设 $\mathcal{X}$ 是概形 $S$ 上的代数叠。假设
$\mathcal{I}_\mathcal{X} \to \mathcal{X}$ 是局部有限表示的。
设 $A \to B$ 是平坦的 $S$-代数同态。设 $x$ 是
$\mathcal{X}$ 在 $A$ 上的对象，并置 $y = x|_B$。那么
$\text{Inf}_x(M) \otimes_A B = \text{Inf}_y(M \otimes_A B)$。
\end{lemma}

\begin{proof}
回忆，$\text{Inf}_x(M)$ 是 $x$ 到 $A[M]$ 的平凡形变的自同构组成的集合，
这些自同构在 $A$ 上诱导 $x$ 的恒等自同构。该平凡形变是经由
$\Spec(A[M]) \to \Spec(A)$ 把 $x$ 拉回到 $\Spec(A[M])$ 所得的。
令 $G \to \Spec(A)$ 为 $x$ 的自同构群代数空间
% P12-ERR-0083: see ../control/ERRATA_CANDIDATES.jsonl
（它存在是因为 $\mathcal{X}$ 是代数空间）。令
$e : \Spec(A) \to G$ 为单位元。《空间态射的进一步讨论》第
\ref{spaces-more-morphisms-section-action-by-derivations} 节的讨论给出
$$
\text{Inf}_x(M) = \Hom_A(e^*\Omega_{G/A}, M)
$$
同理，
$$
\text{Inf}_y(M \otimes_A B) = \Hom_B(e_B^*\Omega_{G_B/B}, M \otimes_A B)
$$
由于依假设 $G \to \Spec(A)$ 是局部有限表示的，可见
$\Omega_{G/A}$ 是局部有限表示的；参见《空间态射的进一步讨论》中的引理
\ref{spaces-more-morphisms-lemma-finite-presentation-differentials}。
因此，$e^*\Omega_{G/A}$ 是有限表示 $A$-模。此外，由《空间态射的进一步讨论》中的引理
\ref{spaces-more-morphisms-lemma-base-change-differentials}，
$\Omega_{G_B/B}$ 是 $\Omega_{G/A}$ 的拉回。因此
$e_B^*\Omega_{G_B/B} = e^*\Omega_{G/A} \otimes_A B$。
% P12-ERR-0084: see ../control/ERRATA_CANDIDATES.jsonl
由《代数的进一步讨论》中的引理
\ref{more-algebra-lemma-pseudo-coherence-and-base-change-ext} 即得结论。
\end{proof}

\begin{lemma}
\label{stacks-more-morphisms-lemma-sheaf-of-infinitesimal-lifts}
设 $\mathcal{X}$ 是基概形 $S$ 上的代数叠。假设
$\mathcal{I}_\mathcal{X} \to \mathcal{X}$ 是局部有限表示的。
设 $(A' \to A, x)$ 是一个形变情形。那么函子
$$
F : B' \longmapsto
\{\text{提升 }x|_{B' \otimes_{A'} A}\text{ 至 } B'\}/\text{同构}
$$
是《拓扑》中的定义 \ref{topologies-definition-big-small-fppf}
所定义的格址 $(\textit{Aff}/\Spec(A'))_{fppf}$ 上的层。
\end{lemma}

\begin{proof}
设 $\{T'_i \to T'\}_{i = 1, \ldots n}$ 是 $A'$ 上仿射概形的一个
标准 fppf 覆盖。写成 $T' = \Spec(B')$。像往常一样，记
$$
T'_{i_0 \ldots i_p} =
T'_{i_0} \times_{T'} \ldots \times_{T'} T'_{i_p} = \Spec(B'_{i_0 \ldots i_p})
$$
其中该环是适当的张量积。置 $B = B' \otimes_{A'} A$ 以及
$B_{i_0 \ldots i_p} = B'_{i_0 \ldots i_p} \otimes_{A'} A$。
记 $y = x|_B$ 和 $y_{i_0 \ldots i_p} = x|_{B_{i_0 \ldots i_p}}$。
设 $\gamma_i \in F(B'_i)$，并假设 $\gamma_{i_0}$ 与
$\gamma_{i_1}$ 在 $F(B'_{i_0i_1})$ 中的像相同。必须找到唯一的
$\gamma \in F(B')$，它在 $F(B'_i)$ 中映到 $\gamma_i$。

\medskip\noindent
在同构类 $\gamma_i$ 中选取 $\textit{Lift}(y_i, B'_i)$ 的一个具体对象
$y'_i$。选取范畴 $\textit{Lift}(y_{i_0i_1}, B'_{i_0i_1})$ 中的同构
$\varphi_{i_0i_1} : y'_{i_0}|_{B'_{i_0i_1}} \to
y'_{i_1}|_{B'_{i_0i_1}}$。如果这些映射
$\varphi_{i_0i_1}$ 满足余圈条件，那么由于 $\mathcal{X}$ 是
fppf 拓扑中的叠，便得到所需的对象 $\gamma$。余圈条件是复合
$$
y'_{i_0}|_{B'_{i_0i_1i_2}}
\xrightarrow{\varphi_{i_0i_1}|_{B'_{i_0i_1i_2}}}
y'_{i_1}|_{B'_{i_0i_1i_2}}
\xrightarrow{\varphi_{i_1i_2}|_{B'_{i_0i_1i_2}}}
y'_{i_2}|_{B'_{i_0i_1i_2}}
\xrightarrow{\varphi_{i_2i_0}|_{B'_{i_0i_1i_2}}}
y'_{i_0}|_{B'_{i_0i_1i_2}}
$$
是恒等态射。若非如此，则这些映射给出元素
$$
\delta_{i_0i_1i_2} \in
\text{Inf}_{y_{i_0i_1i_2}}(J_{i_0i_1i_2}) =
\text{Inf}_y(J) \otimes_B B_{i_0i_1i_2}
$$
这里 $J = \Ker(B' \to B)$，而
$J_{i_0 \ldots i_p} = \Ker(B'_{i_0 \ldots i_p} \to
B_{i_0 \ldots i_p})$。所示等式由引理
\ref{stacks-more-morphisms-lemma-inf-quasi-coherent} 得出；在应用该引理时取
% P12-ERR-0085: see ../control/ERRATA_CANDIDATES.jsonl
$B' \to B'_{i_0 \ldots i_p}$ 以及 $y$ 和 $y_{i_0 \ldots i_p}$。
映射 $B' \to B'_{i_0 \ldots i_p}$ 的平坦性还保证
$J_{i_0 \ldots i_p} = J \otimes_{B'} B'_{i_0 \ldots i_p}$。
一个略去的计算表明，$\delta_{i_0i_1i_2}$ 给出
{\v C}ech 复形中的一个 $2$-余圈
$$
\prod \text{Inf}_y(J) \otimes_B B_{i_0} \to
\prod \text{Inf}_y(J) \otimes_B B_{i_0i_1} \to
\prod \text{Inf}_y(J) \otimes_B B_{i_0i_1i_2} \to \ldots
$$
由《下降》中的引理 \ref{descent-lemma-standard-covering-Cech-quasi-coherent}，
该复形在正次数上无上同调，并且 $H^0 = \text{Inf}_y(J)$。
由于 $\text{Inf}_{y_{i_0i_1}}(J_{i_0i_1})$ 作用在态射上
（《Artin 公理》中的注 \ref{artin-remark-automorphisms}），
这意味着可以修改 $\varphi_{i_0i_1}$ 的选取，从而归结到
$\delta_{i_0i_1i_2} = 0$ 的情形。

\medskip\noindent
唯一性。还须证明，至多存在一个对所有 $i$ 都限制为 $\gamma_i$ 的
$\gamma$。假设有 $\textit{Lift}(y, B')$ 的对象 $y', z'$，以及
$\textit{Lift}(y_i, B'_i)$ 中的同构
$\psi_i : y'|_{B'_i} \to z'|_{B'_i}$。于是可以考察
$$
\psi_{i_1}^{-1} \circ \psi_{i_0} \in
\text{Inf}_{y_{i_0i_1}}(J_{i_0i_1}) =
\text{Inf}_y(J) \otimes_B B_{i_0i_1}
$$
与上面相同地论证可知，在 $B'$ 上寻找 $y'$ 与 $z'$ 之间同构的障碍，
是上述 {\v C}ech 复形的 $H^1$ 中的一个元素，而该群为零。
\end{proof}

\begin{lemma}
\label{stacks-more-morphisms-lemma-T-quasi-coherent}
设 $\mathcal{X}$ 是概形 $S$ 上的代数叠，并且其结构态射
$\mathcal{X} \to S$ 是局部有限表示的。设 $A \to B$ 是平坦的
$S$-代数同态。设 $x$ 是 $\mathcal{X}$ 在 $A$ 上的对象。那么
% P12-ERR-0086: see ../control/ERRATA_CANDIDATES.jsonl
$T_x(M) \otimes_A B = T_y(M \otimes_A B)$。
\end{lemma}

\begin{proof}
选取一个概形 $U$ 以及一个满光滑态射 $U \to \mathcal{X}$。
首先把该引理归结到 $x$ 可提升至 $U$ 的情形。回忆，$T_x(M)$ 是
$x$ 到 $A[M]$ 的提升的同构类组成的集合。因此，引理
\ref{stacks-more-morphisms-lemma-sheaf-of-infinitesimal-lifts}\footnote{该引理适用：
$\Delta : \mathcal{X} \to \mathcal{X} \times_S \mathcal{X}$
由《叠的态射》中的引理
\ref{stacks-morphisms-lemma-finite-presentation-permanence}
以及 $\mathcal{X} \to S$ 是局部有限表示的假设可知是局部有限表示的。
因此，$\mathcal{I}_\mathcal{X} \to \mathcal{X}$ 作为
$\Delta$ 的基变换是局部有限表示的。} 表明，规则
$$
A_1 \mapsto T_{x|_{A_1}}(M \otimes_A A_1)
$$
是 $\Spec(A)$ 的小 \'etale 格址上的层；这里需要张量积，才能使
$A[M] \to A_1[M \otimes_A A_1]$ 成为平坦环映射。
可以选取一个忠实平坦 \'etale 环映射 $A \to A_1$，使
$x|_{A_1}$ 可提升为态射 $u_1 : \Spec(A_1) \to U$；例如参见
《叠上的层》中的引理
\ref{stacks-sheaves-lemma-surjective-flat-locally-finite-presentation}。
写成 $A_2 = A_1 \otimes_A A_1$，并置
$B_1 = B \otimes_A A_1$ 和 $B_2 = B \otimes_A A_2$。考虑图表
$$
\xymatrix{
0 \ar[r] &
T_y(M \otimes_A B) \ar[r] &
T_{y|_{B_1}}(M \otimes_A B_1) \ar[r] &
T_{y|_{B_2}}(M \otimes_A B_2) \\
0 \ar[r] &
T_x(M) \ar[r] \ar[u] &
T_{x|_{A_1}}(M \otimes_A A_1) \ar[r] \ar[u] &
T_{x|_{A_2}}(M \otimes_A A_2) \ar[u]
}
$$
两行由层条件可知都是正合的。又有
$M \otimes_A B_i = (M \otimes_A A_i) \otimes_{A_i} B_i$。
因此，只要对中间和右边的竖直箭头证明结论，所需结果便成立。
这把问题归结到下一段讨论的情形。

\medskip\noindent
假设 $x$ 是态射 $u : \Spec(A) \to U$ 的像。注意，
$T_u(M) \to T_x(M)$ 是满射，因为 $U \to \mathcal{X}$ 光滑且
可由代数空间表示；参见《可表示性判据》中的引理
\ref{criteria-lemma-representable-by-spaces-formally-smooth}
（解释见其前的讨论）以及《空间态射的进一步讨论》中的引理
\ref{spaces-more-morphisms-lemma-smooth-formally-smooth}。
置 $R = U \times_\mathcal{X} U$。回忆，我们得到代数空间中的群胚
$(U, R, s, t, c, e, i)$，并且 $\mathcal{X} = [U/R]$。
由《Artin 公理》中的引理 \ref{artin-lemma-ses-inf-and-T}，有正合列
$$
T_{e \circ u}(M) \to T_u(M) \oplus T_u(M) \to T_x(M) \to 0
$$
其中右端的零由上面已经证明。对到 $B$ 的基变换也有类似的序列。
因此，只要对 $T_u(M)$ 和 $T_{e \circ u}(M)$ 证明该引理的结论，
便得到所需结果。这把问题归结到下一段讨论的情形。

\medskip\noindent
假设 $\mathcal{X} = X$ 是 $S$ 上局部有限表示的代数空间。那么
$$
T_x(M) = \Hom_A(x^*\Omega_{X/S}, M)
$$
这由《空间态射的进一步讨论》第
\ref{spaces-more-morphisms-section-action-by-derivations} 节的讨论得出。
同理，
$$
T_y(M \otimes_A B) = \Hom_B(y^*\Omega_{X/S}, M \otimes_A B)
$$
由于 $X \to S$ 是局部有限表示的，可见 $\Omega_{X/S}$ 是局部有限表示的；
参见《空间态射的进一步讨论》中的引理
\ref{spaces-more-morphisms-lemma-finite-presentation-differentials}。
因此，$x^*\Omega_{X/S}$ 是有限表示 $A$-模。显然，
$y^*\Omega_{X/S} = x^*\Omega_{X/S} \otimes_A B$。
% P12-ERR-0087: see ../control/ERRATA_CANDIDATES.jsonl
由《代数的进一步讨论》中的引理
\ref{more-algebra-lemma-pseudo-coherence-and-base-change-ext} 即得结论。
\end{proof}

\begin{lemma}
\label{stacks-more-morphisms-lemma-local-lift-enough}
设 $\mathcal{X}$ 是概形 $S$ 上的代数叠，并且其结构态射
$\mathcal{X} \to S$ 是局部有限表示的。设
$(A' \to A, x)$ 是一个形变情形。如果存在一个忠实平坦有限表示
$A'$-代数 $B'$，以及 $\mathcal{X}$ 在 $B'$ 上的对象 $y'$，
它提升 $x|_{B' \otimes_{A'} A}$，那么存在 $A'$ 上提升 $x$ 的对象 $x'$。
\end{lemma}

\begin{proof}
令 $I = \Ker(A' \to A)$。置
$B'_1 = B' \otimes_{A'} B'$ 和
$B'_2 = B' \otimes_{A'} B' \otimes_{A'} B'$。令
$J = IB'$、$J_1 = IB'_1$、$J_2 = IB'_2$，以及
$B = B'/J$、$B_1 = B'_1/J_1$、$B_2 = B'_2/J_2$。
置 $y = x|_B$、$y_1 = x|_{B_1}$、$y_2 = x|_{B_2}$。
令 $F$ 为引理 \ref{stacks-more-morphisms-lemma-sheaf-of-infinitesimal-lifts} 中的 fppf 层
（它适用；参见引理 \ref{stacks-more-morphisms-lemma-T-quasi-coherent} 的证明中的脚注）。
于是有等化子图表
$$
\xymatrix{
F(A') \ar[r] &
F(B') \ar@<1ex>[r] \ar@<-1ex>[r] &
F(B'_1)
}
$$
另一方面，采用《Artin 公理》第 \ref{artin-section-inf} 节的术语，有
$F(B') = \text{Lift}(y, B')$、
$F(B'_1) = \text{Lift}(y_1, B'_1)$ 和
$F(B'_2) = \text{Lift}(y_2, B'_2)$。这些集合非空，并且分别是
$T_y(J)$、$T_{y_1}(J_1)$、$T_{y_2}(J_2)$ 的（典范）主齐性空间；
参见《Artin 公理》中的引理
\ref{artin-lemma-properties-lift-RS-star}。
因此，$y'$ 在 $F(B'_1)$ 中的两个像之差是一个元素
$$
\delta_1 \in T_{y_1}(J_1) = T_x(I) \otimes_A B_1
$$
所示等式由引理 \ref{stacks-more-morphisms-lemma-T-quasi-coherent} 得出；在应用该引理时取
% P12-ERR-0088: see ../control/ERRATA_CANDIDATES.jsonl
$A' \to B'_1$ 以及 $x$ 和 $y_1$。映射 $A' \to B'_1$ 的平坦性
还保证 $J_1 = I \otimes_{A'} B'_1$。对于 $B'$ 和 $B'_2$
也有类似的等式。一个略去的计算表明，$\delta_1$ 给出
{\v C}ech 复形中的一个 $1$-余圈
$$
T_x(I) \otimes_A B \to
T_x(I) \otimes_A B_1 \to
T_x(I) \otimes_A B_2 \to \ldots
$$
由《下降》中的引理 \ref{descent-lemma-standard-covering-Cech-quasi-coherent}，
该复形在正次数上无上同调，并且 $H^0 = T_x(I)$。
因此，可以选取 $T_x(I) \otimes_A B = T_y(J)$ 中边界为
$\delta_1$ 的一个元素。以该元素作用所得结果代替 $y'$，
便得到满足 $\delta_1 = 0$ 的新选择 $y'$。于是 $y'$ 在两个映射
$F(B') \to F(B'_1)$ 下映到同一元素，并且由层条件得到
% P12-ERR-0089: see ../control/ERRATA_CANDIDATES.jsonl
$F(A')$ 的一个元素。
\end{proof}

\section{形式光滑态射}
\label{stacks-more-morphisms-section-formally-smooth}

\noindent
本节引入代数叠的形式光滑态射
$\mathcal{X} \to \mathcal{Y}$ 的概念。这种态射的特征是：
当 $T$ 仿射时，$\mathcal{X}$ 的 $T$-值点可以提升到
$T$ 的无穷小增厚。主要结果是，形式光滑且局部有限表示的态射是光滑的；
参见引理 \ref{stacks-more-morphisms-lemma-smooth-formally-smooth}。事实证明，
该判据通常比雅可比判据更容易使用。

\begin{definition}
\label{stacks-more-morphisms-definition-formally-smooth}
若代数叠的态射 $f : \mathcal{X} \to \mathcal{Y}$ 作为群胚纤维化范畴中的
$1$-态射在对象上形式光滑（见《可表示性判据》第
\ref{criteria-section-formally-smooth} 节的解释），则称它是
{\it 形式光滑的}。
\end{definition}

\noindent
我们把该定义的条件翻译成当前使用的语言（参见《叠的性质》第
\ref{stacks-properties-section-conventions} 节）。设
$f : \mathcal{X} \to \mathcal{Y}$ 是代数叠的态射。考虑一个
$2$-交换的实线图表
\begin{equation}
\label{stacks-more-morphisms-equation-diagram}
\vcenter{
\xymatrix{
T \ar[r]_-x \ar[d]_i & \mathcal{X} \ar[d]^f \\
T' \ar[r]^-y \ar@{..>}[ru] & \mathcal{Y}
}
}
\end{equation}
其中 $i : T \to T'$ 是仿射概形的一阶增厚。令
$$
\gamma : y \circ i \longrightarrow f \circ x
$$
为见证该图表 $2$-交换性的 $2$-态射。（记号同《范畴论》第
\ref{categories-section-formal-cat-cat} 节和第
\ref{categories-section-2-categories} 节。）
给定 (\ref{stacks-more-morphisms-equation-diagram}) 和 $\gamma$，一个
{\it 虚线箭头}是三元组 $(x', \alpha, \beta)$，其中包含态射
$x' : T' \to \mathcal{X}$ 和 $2$-箭头
$\alpha : x' \circ i \to x$、$\beta : y \to f \circ x'$，并满足
$\gamma = (\text{id}_f \star \alpha) \circ
(\beta \star \text{id}_i)$；换言之，图表
$$
\xymatrix{
& f \circ x' \circ i \ar[rd]^{\text{id}_f \star \alpha} \\
y \circ i \ar[ru]^{\beta \star \text{id}_i} \ar[rr]^\gamma & &
f \circ x
}
$$
交换。虚线箭头的一个{\it 态射}
$(x'_1, \alpha_1, \beta_1) \to (x'_2, \alpha_2, \beta_2)$
是 $2$-箭头 $\theta : x'_1 \to x'_2$，它满足
$\alpha_1 = \alpha_2 \circ (\theta \star \text{id}_i)$ 和
$\beta_2 = (\text{id}_f \star \theta) \circ \beta_1$。

\medskip\noindent
刚才描述的虚线箭头范畴是《范畴论》中的定义
\ref{categories-definition-dotted-arrows} 的一个特例。

\begin{lemma}
\label{stacks-more-morphisms-lemma-reformulate-formal-smoothness}
代数叠的态射 $f : \mathcal{X} \to \mathcal{Y}$ 是形式光滑的
（定义 \ref{stacks-more-morphisms-definition-formally-smooth}），当且仅当对于每个图表
(\ref{stacks-more-morphisms-equation-diagram}) 和 $\gamma$，虚线箭头范畴都非空。
\end{lemma}

\begin{proof}
略去不同语言之间的翻译。
\end{proof}

\begin{lemma}
\label{stacks-more-morphisms-lemma-base-change-formal-smoothness}
代数叠的形式光滑态射沿任意代数叠态射的基变换仍然形式光滑。
\end{lemma}

\begin{proof}
这由《范畴论》中的引理
\ref{categories-lemma-cat-dotted-arrows-base-change} 和定义得出。
\end{proof}

\begin{lemma}
\label{stacks-more-morphisms-lemma-composition-formal-smoothness}
代数叠的形式光滑态射的复合仍然形式光滑。
\end{lemma}

\begin{proof}
这由《范畴论》中的引理
\ref{categories-lemma-cat-dotted-arrows-composition} 和定义得出。
\end{proof}

\begin{lemma}
\label{stacks-more-morphisms-lemma-formal-smoothness-representable}
设 $f : \mathcal{X} \to \mathcal{Y}$ 是可由代数空间表示的代数叠态射。
那么下列条件等价：
\begin{enumerate}
\item $f$ 形式光滑；
\item 对每个概形 $T$ 和态射 $T \to \mathcal{Y}$，态射
$\mathcal{X} \times_\mathcal{Y} T \to T$
作为代数空间的态射是形式光滑的。
\end{enumerate}
\end{lemma}

\begin{proof}
这由《范畴论》中的引理
\ref{categories-lemma-cat-dotted-arrows-base-change} 和定义得出。
\end{proof}

\begin{lemma}
\label{stacks-more-morphisms-lemma-lift-to-smooth}
设 $T \to T'$ 是仿射概形的一阶增厚。设 $\mathcal{X}'$ 是
$T'$ 上的代数叠，并且其结构态射
$\mathcal{X}' \to T'$ 是光滑的。设 $x : T \to \mathcal{X}'$
是 $T'$ 上的态射。那么存在 $T'$ 上的
% P12-ERR-0090: see ../control/ERRATA_CANDIDATES.jsonl
态射 $x' : T' \to \mathcal{X}'$，满足 $x'|_T = x$。
\end{lemma}

\begin{proof}
可以应用引理 \ref{stacks-more-morphisms-lemma-local-lift-enough} 的结果。因此，只须构造
一个光滑满态射 $W' \to T'$，其中 $W'$ 仿射，并使
% P12-ERR-0091: see ../control/ERRATA_CANDIDATES.jsonl
$x|_{T \times_{W'} T'}$ 可提升到 $W'$。（我们促请读者使用已经建立的
代数空间的类似结果，自行找出这一事实的证明。）选取一个概形 $U'$
以及一个满光滑态射 $U' \to \mathcal{X}'$。注意，
$U' \to T'$ 是光滑的，并且投影
$T \times_{\mathcal{X}'} U' \to T$ 是满光滑的。
选取一个仿射概形 $W$ 和一个 \'etale 态射
$W \to T \times_{\mathcal{X}'} U'$，使 $W \to T$ 是满射。
于是 $W \to T$ 是仿射概形的光滑态射。以若干主仿射开集的不交并
代替 $W$ 后，可以假设存在仿射概形的光滑态射 $W' \to T'$，使得
$W = T \times_{T'} W'$；参见《代数》中的引理
\ref{algebra-lemma-lift-smooth}。由《空间态射的进一步讨论》中的引理
\ref{spaces-more-morphisms-lemma-smooth-formally-smooth}，
可以找到 $T'$ 上的态射 $W' \to U'$，它提升给定态射
$W \to U'$。证明完毕。
\end{proof}

\noindent
下面的引理是本节的主要结果。结合《叠的极限》中的命题
\ref{stacks-limits-proposition-characterize-locally-finite-presentation}，
它表明，可以用群胚叠的 $1$-态射
$\mathcal{X} \to \mathcal{Y}$ 的“简单”性质来判定代数叠的态射
$f : \mathcal{X} \to \mathcal{Y}$ 是否光滑。

\begin{lemma}[无穷小提升判据]
\label{stacks-more-morphisms-lemma-smooth-formally-smooth}
设 $f : \mathcal{X} \to \mathcal{Y}$ 是代数叠的态射。下列条件等价：
\begin{enumerate}
\item 态射 $f$ 是光滑的。
\item 态射 $f$ 是局部有限表示且形式光滑的。
\end{enumerate}
\end{lemma}

\begin{proof}
假设 $f$ 光滑。由《叠的态射》中的引理
\ref{stacks-morphisms-lemma-smooth-locally-finite-presentation}，
$f$ 是局部有限表示的。因此，给定一个图表
(\ref{stacks-more-morphisms-equation-diagram}) 和
$\gamma : y \circ i \to f \circ x$，只须找到一个虚线箭头
（参见引理 \ref{stacks-more-morphisms-lemma-reformulate-formal-smoothness}）。
构造纤维积，得到
$$
\xymatrix{
T \ar[d] \ar[r] &
T' \times_\mathcal{Y} \mathcal{X} \ar[d] \ar[r] &
\mathcal{X} \ar[d] \\
T' \ar[r] & T' \ar[r] & \mathcal{Y}
}
$$
由此可见，只须在左方形中找到虚线箭头。由于
$T' \times_\mathcal{Y} \mathcal{X} \to T'$ 是光滑的
（《叠的态射》中的引理
\ref{stacks-morphisms-lemma-base-change-smooth}），
引理 \ref{stacks-more-morphisms-lemma-lift-to-smooth} 保证左方形中的虚线箭头存在。

\medskip\noindent
反过来，假设 $f$ 是局部有限表示且形式光滑的。选取一个概形 $U$
以及一个满光滑态射 $U \to \mathcal{X}$。那么
$a : U \to \mathcal{X}$ 和 $b : U \to \mathcal{Y}$
均可由代数空间表示且局部有限表示（使用《叠的态射》中的引理
\ref{stacks-morphisms-lemma-composition-finite-presentation}
以及上面已见的光滑态射是局部有限表示的事实）。
我们将应用《代数叠》中的引理
\ref{algebraic-lemma-representable-transformations-property-implication}
所给出的一般原理，以《空间态射的进一步讨论》中的引理
\ref{spaces-more-morphisms-lemma-smooth-formally-smooth} 的等价性为输入，
同时使用《可表示性判据》中的引理
\ref{criteria-lemma-representable-by-spaces-formally-smooth} 的翻译。
首先把它应用于 $a$，可见 $a$ 在对象上形式光滑。其次，
利用依假设 $f$ 在对象上形式光滑
（参见引理 \ref{stacks-more-morphisms-lemma-reformulate-formal-smoothness}）以及
《可表示性判据》中的引理
\ref{criteria-lemma-composition-formally-smooth}，可见
$b = f \circ a$ 在对象上形式光滑。然后再次应用该原理，得出
$b$ 光滑。由代数叠态射之光滑性的定义，这意味着 $f$ 光滑，
证明完毕。
\end{proof}

\section{吹起与平坦性}
\label{stacks-more-morphisms-section-blowup-flat}

\noindent
本节简要讨论，对于代数空间上的代数叠，可以从
《空间态射的进一步讨论》第
\ref{spaces-more-morphisms-section-blowup-flat} 节和第
\ref{spaces-more-morphisms-section-applications-flattening-by-blowing-up}
节推出哪些结论。

\begin{lemma}
\label{stacks-more-morphisms-lemma-flatten-stack}
设 $f : \mathcal{X} \to Y$ 是从代数叠到代数空间的态射。
设 $V \subset Y$ 是开子空间。假设
\begin{enumerate}
\item $Y$ 是拟紧且拟分离的；
\item $f$ 是有限型且拟分离的；
\item $V$ 是拟紧的；并且
\item $\mathcal{X}_V$ 在 $V$ 上平坦且局部有限表示。
\end{enumerate}
那么存在一个 $V$-可容许吹起 $Y' \to Y$，以及闭子叠
$\mathcal{X}' \subset \mathcal{X}_{Y'}$，满足
$\mathcal{X}'_V = \mathcal{X}_V$，使得
$\mathcal{X}' \to Y'$ 平坦且有限表示。
\end{lemma}

\begin{proof}
注意，$\mathcal{X}$ 是拟紧的。选取一个仿射概形 $U$
以及一个满光滑态射 $U \to \mathcal{X}$。令
$R = U \times_\mathcal{X} U$，于是得到 $Y$ 上代数空间中的群胚
$(U, R, s, t, c)$，并且 $\mathcal{X} = [U/R]$
（《代数叠》中的引理 \ref{algebraic-lemma-stack-presentation}）。
可以把《空间态射的进一步讨论》中的引理
\ref{spaces-more-morphisms-lemma-flat-after-blowing-up}
应用于 $U \to Y$ 和开子空间 $V \subset Y$。于是得到一个
$V$-可容许吹起 $Y' \to Y$，使得严格变换
$U' \subset U_{Y'}$ 在 $Y'$ 上平坦且有限表示。
令 $R' \subset R_{Y'}$ 为 $R$ 的严格变换。由于 $s$ 和 $t$
都是光滑的（特别地也是平坦的），由《空间上的除子》中的引理
\ref{spaces-divisors-lemma-strict-transform-flat}，有笛卡儿图表
$$
\vcenter{
\xymatrix{
R' \ar[r] \ar[d] & R_{Y'} \ar[d]^{s_{Y'}} \\
U' \ar[r] & U_{Y'}
}
}
\quad\text{以及}\quad
\vcenter{
\xymatrix{
R' \ar[r] \ar[d] & R_{Y'} \ar[d]^{t_{Y'}} \\
U' \ar[r] & U_{Y'}
}
}
$$
换言之，$U'$ 是 $U_{Y'}$ 的一个 $R_{Y'}$-不变闭子空间。
因此，由《叠的性质》中的引理
\ref{stacks-properties-lemma-substacks-presentation}，
$U'$ 定义闭子叠 $\mathcal{X}' \subset \mathcal{X}_{Y'}$。
态射 $\mathcal{X}' \to Y'$ 平坦且局部有限表示，因为
$U' \to Y'$ 具有这些性质。另一方面，已经知道
$\mathcal{X}' \to Y'$ 是拟紧且拟分离的
（由对 $f$ 的假设以及闭浸入具有这些性质）。证明完毕。
\end{proof}

\section{代数叠的 Chow 引理}
\label{stacks-more-morphisms-section-chow}

\noindent
本节讨论代数叠的 Chow 引理。

\begin{lemma}
\label{stacks-more-morphisms-lemma-finite-cover-factor}
设 $Y$ 是拟紧拟分离代数空间。设 $V \subset Y$ 是拟紧开子空间。
设 $f : \mathcal{X} \to V$ 是满、平坦且局部有限表示的态射。
那么存在有限满态射 $g : Y' \to Y$，使得
% P12-ERR-0092: see ../control/ERRATA_CANDIDATES.jsonl
$V' = g^{-1}(V) \to Y$ 在 Zariski 局部经由 $f$ 分解。
\end{lemma}

\begin{proof}
先证明 $Y$ 是概形的情形。可以选取一个概形 $U$ 以及一个满光滑态射
$U \to \mathcal{X}$。于是 $\{U \to V\}$ 是概形的 fppf 覆盖。
由《态射的进一步讨论》中的引理
\ref{more-morphisms-lemma-dominate-fppf-finite}，存在有限满态射
$V' \to V$，使 $V' \to V$ 在 Zariski 局部经由 $U$ 分解。
由《态射的进一步讨论》中的引理
\ref{more-morphisms-lemma-extend-finite-surjective-morphisms}，
可以找到有限满态射 $Y' \to Y$，它在 $V$ 上的限制如所需为
$V' \to V$。

\medskip\noindent
若 $Y$ 是代数空间，则先沿一个有限满态射 $Y' \to Y$ 作有限基变换，
其中 $Y'$ 是概形，便可看出该引理成立。参见《空间的极限》中的命题
\ref{spaces-limits-proposition-there-is-a-scheme-finite-over}。
\end{proof}

\begin{lemma}
\label{stacks-more-morphisms-lemma-make-section}
设 $f : \mathcal{X} \to Y$ 是从代数叠到代数空间的态射。
设 $V \subset Y$ 是开子空间。假设
\begin{enumerate}
\item $f$ 是分离且有限型的；
\item $Y$ 是拟紧且拟分离的；
\item $V$ 是拟紧的；并且
\item $\mathcal{X}_V$ 是 $V$ 上的胚。
\end{enumerate}
那么存在交换图表
$$
\xymatrix{
\overline{Z} \ar[rd]_{\overline{g}} &
Z \ar[l]^j \ar[d]_g \ar[r]_h & \mathcal{X} \ar[ld]^f \\
& Y
}
$$
其中 $j$ 是开浸入，$\overline{g}$ 和 $h$ 是固有的，并且
$|V|$ 包含于 $|g|$ 的像中。
\end{lemma}

\begin{proof}
假设有交换图表
$$
\xymatrix{
\mathcal{X}' \ar[d]_{f'} \ar[r] & \mathcal{X} \ar[d]^f \\
Y' \ar[r] & Y
}
$$
以及拟紧开子空间 $V' \subset Y'$，使得 $Y' \to Y$
是代数空间的固有态射，$\mathcal{X}' \to \mathcal{X}$
是代数叠的固有态射，$V' \subset Y'$ 满射到 $V$，并且
$\mathcal{X}'_{V'}$ 是 $V'$ 上的胚。那么只须对二元组
$(f' : \mathcal{X}' \to Y', V')$ 证明该引理。略去一些细节。

\medskip\noindent
证明的总体策略。我们将反复应用上面的注记，把问题归结到
$f$ 的像是开集且 $f$ 在该开集上有截面的情形。每一步都很直接，
但步骤相当多，因而证明稍显繁复。

\medskip\noindent
使用《空间的极限》中的命题
\ref{spaces-limits-proposition-there-is-a-scheme-finite-over}，
把问题归结到 $Y$ 是概形的情形。（取一个有限满态射
$Y' \to Y$，其中 $Y'$ 是概形。置
$\mathcal{X}' = \mathcal{X}_{Y'}$，并应用证明开头的注记。）

\medskip\noindent
使用引理 \ref{stacks-more-morphisms-lemma-flatten-stack}
（并使用《叠的态射》中的引理
\ref{stacks-morphisms-lemma-gerbe-fppf} 说明胚是平坦且局部有限表示的），
把问题归结到 $f$ 平坦且有限表示的情形。

\medskip\noindent
由于 $f$ 平坦且局部有限表示，可见 $|f|$ 的像是开集 $W \subset Y$。
又因为 $\mathcal{X}$ 拟紧（这是因为 $f$ 有限型且 $Y$ 拟紧），
所以 $W$ 拟紧。由引理 \ref{stacks-more-morphisms-lemma-finite-cover-factor}，
可以找到有限满态射 $g : Y' \to Y$，使
$g^{-1}(W) \to Y$ 在 Zariski 局部经由
$\mathcal{X} \to Y$ 分解。以 $Y'$ 代替 $Y$，并以
$\mathcal{X} \times_Y Y'$ 代替 $\mathcal{X}$ 后，
归结到下一段所描述的情形。

\medskip\noindent
假设存在 $n \geq 0$、拟紧开集 $W_i \subset Y$
（$i = 1, \ldots, n$），以及态射 $x_i : W_i \to \mathcal{X}$，使得
(a) $f \circ x_i = \text{id}_{W_i}$，
(b) $W = \bigcup_{i = 1, \ldots, n} W_i$ 包含 $V$，并且
(c) $W$ 是 $|f|$ 的像。
对 $n$ 作归纳。起始情形为 $n = 0$：这蕴含 $V = \emptyset$，
此时可取 $\overline{Z} = \emptyset$。
若 $n > 0$，则对 $i = 1, \ldots, n$，考虑约化闭子概形
$Y_i$，其底拓扑空间为 $Y \setminus W_i$。考虑有限态射
$$
Y' = Y \amalg \coprod\nolimits_{i = 1, \ldots, n} Y_i \longrightarrow Y
$$
以及拟紧开集
$$
V' = (W_1 \cap \ldots \cap W_n \cap V) \amalg
\coprod_{i = 1, \ldots, n} (V \cap Y_i).
$$
由证明开头的注记，只要能对二元组
$$
(\mathcal{X} \to Y, W_1 \cap \ldots \cap W_n \cap V)
\quad\text{以及}\quad
(\mathcal{X} \times_Y Y_i \to Y_i, V \cap Y_i),\quad
i = 1, \ldots, n
$$
证明该引理，所需结果便成立。这里使用集合论等式
$V = (W_1 \cap \ldots \cap W_n \cap V) \cup
\bigcup\nolimits_{i = 1, \ldots n} (V \cap Y_i)$。
归纳假设适用于上述第二类二元组。因此归结到下一段所描述的情形。

\medskip\noindent
假设存在 $n \geq 0$、拟紧开集 $W_i \subset Y$
（$i = 1, \ldots, n$），以及态射 $x_i : W_i \to \mathcal{X}$，使得
(a) $f \circ x_i = \text{id}_{W_i}$，
(b) $W = \bigcup_{i = 1, \ldots, n} W_i$ 包含 $V$，
(c) $W$ 是 $|f|$ 的像，并且
(d) $V \subset W_1 \cap \ldots \cap W_n$。
态射
$$
T_{ij} = \mathit{Isom}_\mathcal{X}(x_i|_{W_i \cap W_j \cap V},
x_j|_{W_i \cap W_j \cap V}) \longrightarrow W_i \cap W_j \cap V
$$
是满、平坦且局部有限表示的（《叠的态射》中的引理
\ref{stacks-morphisms-lemma-gerbe-isom-fppf}）。
对每个拟紧开集 $W_i \cap W_j \cap V$ 和态射
$T_{ij} \to W_i \cap W_j \cap V$ 应用引理
\ref{stacks-more-morphisms-lemma-finite-cover-factor}，得到有限满态射
$Y'_{ij} \to Y$。以所有 $Y'_{ij}$ 在 $Y$ 上的纤维积代替 $Y$ 后，
归结到下一段所描述的情形。

\medskip\noindent
假设存在 $n \geq 0$、拟紧开集 $W_i \subset Y$
（$i = 1, \ldots, n$），以及态射 $x_i : W_i \to \mathcal{X}$，使得
(a) $f \circ x_i = \text{id}_{W_i}$，
(b) $W = \bigcup_{i = 1, \ldots, n} W_i$ 包含 $V$，
(c) $W$ 是 $|f|$ 的像，
(d) $V \subset W_1 \cap \ldots \cap W_n$，并且
(e) $x_i$ 与 $x_j$ 在 $W_i \cap W_j \cap V$ 上 Zariski 局部同构。
任取 $y \in V$。假设能找到拟紧开邻域
$y \in V_y \subset V$，使得该引理对二元组
$(\mathcal{X} \to Y, V_y)$ 成立，设其解为
$\overline{Z}_y, Z_y, \overline{g}_y, g_y, h_y$。
由于 $V$ 拟紧，可以找到有限多个点 $y_1, \ldots, y_m$，使
$V = V_{y_1} \cup \ldots \cup V_{y_m}$。于是置
$$
\overline{Z} = \coprod \overline{Z}_{y_j},\quad
Z = \coprod Z_{y_j},\quad
\overline{g} = \coprod \overline{g}_{y_j},\quad
g = \coprod g_{y_j},\quad
h = \coprod h_{y_j}
$$
便给出该引理的一个解。给定 $y$，由条件 (e)，可以选取拟紧开邻域
$y \in V_y \subset V$，以及同构
$\varphi_i : x_1|_{V_y} \to x_i|_{V_y}$（$i = 2, \ldots, n$）。
置 $\varphi_{ij} = \varphi_j \circ \varphi_i^{-1}$。
这把问题带到下一段所描述的情形。

\medskip\noindent
假设存在 $n \geq 0$、拟紧开集 $W_i \subset Y$
（$i = 1, \ldots, n$），以及态射 $x_i : W_i \to \mathcal{X}$，使得
(a) $f \circ x_i = \text{id}_{W_i}$，
(b) $W = \bigcup_{i = 1, \ldots, n} W_i$ 包含 $V$，
(c) $W$ 是 $|f|$ 的像，
(d) $V \subset W_1 \cap \ldots \cap W_n$，并且
(f) 存在同构 $\varphi_{ij} : x_i|_V \to x_j|_V$，满足
$\varphi_{jk} \circ \varphi_{ij} = \varphi_{ik}$。
态射
$$
I_{ij} = \mathit{Isom}_\mathcal{X}(x_i|_{W_i \cap W_j},
x_j|_{W_i \cap W_j}) \longrightarrow W_i \cap W_j
$$
是固有的，因为 $f$ 分离（《叠的态射》中的引理
\ref{stacks-morphisms-lemma-separated-implies-isom}）。
注意，$\varphi_{ij}$ 定义 $I_{ij} \to W_i \cap W_j$ 在 $V$ 上的截面
$V \to I_{ij}$。由《空间态射的进一步讨论》中的引理
\ref{spaces-more-morphisms-lemma-get-section-after-blowup}，
可以找到 $V$-可容许吹起 $p_{ij} : Y_{ij} \to Y$，使
% P12-ERR-0093: see ../control/ERRATA_CANDIDATES.jsonl
$s_{ij}$ 延拓到 $p_{ij}^{-1}(W_i \cap W_j)$。
以所有 $Y_{ij}$ 在 $Y$ 上的纤维积代替 $Y$ 后，
到达下一段所描述的情形。

\medskip\noindent
假设存在 $n \geq 0$、拟紧开集 $W_i \subset Y$
（$i = 1, \ldots, n$），以及态射 $x_i : W_i \to \mathcal{X}$，使得
(a) $f \circ x_i = \text{id}_{W_i}$，
(b) $W = \bigcup_{i = 1, \ldots, n} W_i$ 包含 $V$，
(c) $W$ 是 $|f|$ 的像，
(d) $V \subset W_1 \cap \ldots \cap W_n$，并且
(g) 存在同构
$\varphi_{ij} : x_i|_{W_i \cap W_j} \to x_j|_{W_i \cap W_j}$，满足
$$
\varphi_{jk}|_V \circ \varphi_{ij}|_V = \varphi_{ik}|_V.
$$
必要时再次以一个 $V$-可容许吹起代替 $Y$，可以假设 $V$ 在 $Y$
中稠密且概形论稠密，从而在 $Y$ 的任意包含 $V$ 的开子空间中也如此。
作此代替后，由《空间的态射》中的引理
\ref{spaces-morphisms-lemma-equality-of-morphisms} 以及
% P12-ERR-0094: see ../control/ERRATA_CANDIDATES.jsonl
$I_{ik} \to W_i \cap W_j$ 固有（因而分离）这一事实，得出
$$
\varphi_{jk}|_{W_i \cap W_j \cap W_k} \circ
\varphi_{ij}|_{W_i \cap W_j \cap W_k} =
\varphi_{ik}|_{W_i \cap W_j \cap W_k}
$$
当然，这意味着 $(x_i, \varphi_{ij})$ 是一个
% P12-ERR-0095: see ../control/ERRATA_CANDIDATES.jsonl
下降数据；由于 $\mathcal{X}$ 是叠，得到态射
$x : W \to \mathcal{X}$，它在 $W_i$ 上与 $x_i$ 一致。
由于 $x$ 是分离态射 $\mathcal{X} \to W$ 的一个截面，
可见 $x$ 是固有的（《叠的态射》中的引理
\ref{stacks-morphisms-lemma-section-immersion}）。
因此，该引理现在对
$\overline{Z} = Y$、$Z = W$、$\overline{g} = \text{id}_Y$、
$g = \text{id}_W$、$h = x$ 成立。
\end{proof}

\begin{theorem}[Chow 引理]
\label{stacks-more-morphisms-theorem-chow-finite-type}
\begin{reference}
这是 Ofer Gabber 的结果，参见
\cite[Theorem 1.1]{olsson_proper}
\end{reference}
设 $f : \mathcal{X} \to Y$ 是从代数叠到代数空间的态射。假设
\begin{enumerate}
\item $Y$ 是拟紧且拟分离的；
\item
% P12-ERR-0096: see ../control/ERRATA_CANDIDATES.jsonl
$f$ 是分离有限型的。
\end{enumerate}
那么存在交换图表
$$
\xymatrix{
\mathcal{X} \ar[rd] & X \ar[l] \ar[d] \ar[r] & \overline{X} \ar[ld] \\
& Y
}
$$
其中 $X \to \mathcal{X}$ 是固有满态射，
$X \to \overline{X}$ 是开浸入，并且
% P12-ERR-0097: see ../control/ERRATA_CANDIDATES.jsonl
$\overline{X} \to Y$ 是代数空间的固有态射。
\end{theorem}

\begin{proof}
粗略的想法是利用 $\mathcal{X}$ 有一个稠密开子叠为胚
（《叠的态射》中的命题
\ref{stacks-morphisms-proposition-open-stratum}），
并诉诸引理 \ref{stacks-more-morphisms-lemma-make-section}。
该方法不能直接奏效的原因是，这个开子叠可能并不拟紧，
从而会遇到技术困难。因此先作一个（标准的）归约，化到 Noether 情形。

\medskip\noindent
首先选取闭浸入
$\mathcal{X} \to \mathcal{X}'$，其中 $\mathcal{X}'$
是 $Y$ 上分离且有限型的代数叠。参见《叠的极限》中的引理
\ref{stacks-limits-lemma-separated-closed-in-finite-presentation}。
显然，只须对 $\mathcal{X}'$ 证明该定理；因此可以假设
$\mathcal{X} \to Y$ 是分离且有限表示的。

\medskip\noindent
假设 $\mathcal{X} \to Y$ 是分离且有限表示的。
由《空间的极限》中的命题
\ref{spaces-limits-proposition-approximate}，可以把
$Y = \lim Y_i$ 写成一个由 Noether 代数空间组成、转移态射仿射的系统的
有向极限。由《叠的极限》中的引理
\ref{stacks-limits-lemma-descend-a-stack-down}，存在某个 $i$
以及从代数叠到 $Y_i$ 的有限表示态射
$\mathcal{X}_i \to Y_i$，使得
$\mathcal{X} = Y \times_{Y_i} \mathcal{X}_i$。
增大 $i$ 后，可以假设 $\mathcal{X}_i \to Y_i$ 分离；
参见《叠的极限》中的引理
\ref{stacks-limits-lemma-eventually-separated}。
于是只须对 $\mathcal{X}_i \to Y_i$ 证明该定理。
这把问题归结到下一段讨论的情形。

\medskip\noindent
假设 $Y$ 是 Noether 的。可以用其约化代替
$\mathcal{X}$（《叠的性质》中的定义
\ref{stacks-properties-definition-reduced-induced-stack}）。
这把问题归结到下一段讨论的情形。

\medskip\noindent
假设 $Y$ 是 Noether 的且 $\mathcal{X}$ 是约化的。
由于 $\mathcal{X} \to Y$ 分离且 $Y$ 拟分离，可见
$\mathcal{X}$ 作为代数叠是拟分离的。因此惰性态射
$\mathcal{I}_\mathcal{X} \to \mathcal{X}$ 是拟紧的。
于是，由《叠的态射》中的命题
\ref{stacks-morphisms-proposition-open-stratum}，
存在一个作为胚的稠密开子叠
$\mathcal{V} \subset \mathcal{X}$。令
$\mathcal{V} \to V$ 为把 $\mathcal{V}$ 表示成代数空间 $V$
上的胚的态射。关于 $\mathcal{V} \to V$ 的构造，参见
《叠的态射》中的引理
\ref{stacks-morphisms-lemma-gerbe-over-iso-classes}。
特别地，该构造表明态射 $\mathcal{V} \to Y$ 分解为
$\mathcal{V} \to V \to Y$。图示为
$$
\xymatrix{
\mathcal{V} \ar[r] \ar[d] & \mathcal{X} \ar[d] \\
V \ar[r] & Y
}
$$
由于态射 $\mathcal{V} \to V$ 是满、平坦且有限表示的
（《叠的态射》中的引理
\ref{stacks-morphisms-lemma-gerbe-fppf}），并且
$\mathcal{V} \to Y$ 是局部有限表示的，可知
$V \to Y$ 是局部有限表示的（《叠的态射》中的引理
\ref{stacks-morphisms-lemma-flat-finite-presentation-permanence}）。
注意，$\mathcal{V} \to V$ 是万有同胚
（《叠的态射》中的引理
\ref{stacks-morphisms-lemma-gerbe-bijection-points}）。
由于 $\mathcal{V}$ 拟紧（参见《叠的态射》中的引理
\ref{stacks-morphisms-lemma-locally-closed-in-noetherian}），
可见 $V$ 拟紧。最后，由于 $\mathcal{V} \to Y$ 分离，
$V \to Y$ 也分离；这是把《叠的态射》中的引理
\ref{stacks-morphisms-lemma-check-separated-on-ui-cover}
应用于 $\mathcal{V} \to V \to Y$ 所得的
（其假设如上所见均满足）。

\medskip\noindent
以上说明，《空间的极限》中的引理
\ref{spaces-limits-lemma-embedding-into-affine-over-qs}
的假设适用于态射 $V \to Y$。因此可以找到稠密开子空间
$V' \subset V$ 以及 $Y$ 上的浸入
$V' \to \mathbf{P}^n_Y$。显然，可以用 $V'$ 代替 $V$，并用
$V'$ 在 $\mathcal{V}$ 中的逆像代替 $\mathcal{V}$
（回忆，如上所见 $|\mathcal{V}| = |V|$）。
于是可以假设有图表
$$
\xymatrix{
\mathcal{V} \ar[rr] \ar[d] & & \mathcal{X} \ar[d] \\
V \ar[r] & \mathbf{P}^n_Y \ar[r] & Y
}
$$
其中箭头 $V \to \mathbf{P}^n_Y$ 是浸入。
令 $\mathcal{X}'$ 为态射
$$
j : \mathcal{V} \longrightarrow \mathbf{P}^n_Y \times_Y \mathcal{X}
$$
的概形论像，并令 $Y'$ 为态射
$V \to \mathbf{P}^n_Y$ 的概形论像。得到交换图表
$$
\xymatrix{
\mathcal{V} \ar[r] \ar[d] &
\mathcal{X}' \ar[r] \ar[d] &
\mathbf{P}^n_Y \times_Y \mathcal{X} \ar[d] \ar[r] &
\mathcal{X} \ar[d] \\
V \ar[r] &
Y' \ar[r] &
\mathbf{P}^n_Y \ar[r] &
Y
}
$$
（参见《叠的态射》中的引理
\ref{stacks-morphisms-lemma-factor-factor}）。
我们断言 $\mathcal{V} = V \times_{Y'} \mathcal{X}'$，并且引理
\ref{stacks-more-morphisms-lemma-make-section} 适用于态射
$\mathcal{X}' \to Y'$ 和开子空间 $V \subset Y'$。
若该断言成立，则得到
$$
\xymatrix{
\overline{X} \ar[rd]_{\overline{g}} &
X \ar[l] \ar[d]_g \ar[r]_h & \mathcal{X}' \ar[ld]^f \\
& Y'
}
$$
其中 $X \to \overline{X}$ 是开浸入，
$\overline{g}$ 和 $h$ 是固有的，并且 $|V|$ 包含于 $|g|$ 的像中。
于是复合 $X \to \mathcal{X}' \to \mathcal{X}$ 是固有的
（因为它是固有态射的复合），且其像包含 $|\mathcal{V}|$；
因此该复合是满射。类似地，$\overline{X} \to Y' \to Y$
作为固有态射的复合是固有的。

\medskip\noindent
最后一步是证明该断言。注意，$\mathcal{X}' \to Y'$
是分离且有限型的，$Y'$ 是拟紧且拟分离的，而 $V$ 是拟紧的
（我们完全略去对所有细节的检验）。其次，注意
$b : \mathcal{X}' \to \mathcal{X}$ 在 $\mathcal{V}$ 上是同构；
这由《叠的态射》中的引理
\ref{stacks-morphisms-lemma-scheme-theoretic-image-of-partial-section}
得出。特别地，$\mathcal{V}$ 被识别为 $\mathcal{X}'$ 的开子叠。
态射 $j$ 是拟紧的（源拟紧且目标拟分离），所以由《叠的态射》中的引理
\ref{stacks-morphisms-lemma-existence-plus-flat-base-change}，
$j$ 的概形论像的形成与平坦基变换交换。特别地，
$V \times_{Y'} \mathcal{X}'$ 是态射
$\mathcal{V} \to V \times_{Y'} \mathcal{X}'$ 的概形论像。
然而，由《叠的态射》中的引理
\ref{stacks-morphisms-lemma-universally-closed-permanence}，
$|\mathcal{V}| \to |V \times_{Y'} \mathcal{X}'|$ 的像是闭的
（使用如上所见 $\mathcal{V} \to V$ 是万有同胚，从而是万有闭的）。
此外，该像是稠密的（把刚才所述与《叠的态射》中的引理
\ref{stacks-morphisms-lemma-topology-scheme-theoretic-image} 结合）。
% P12-ERR-0098: see ../control/ERRATA_CANDIDATES.jsonl
因此 $|\mathcal{V}| = |V \times_{Y'} \mathcal{X}'|$。
于是 $\mathcal{V} \to V \times_{Y'} \mathcal{X}'$ 是同构，
该断言证毕。
\end{proof}

\section{Noether 赋值判据}
\label{stacks-more-morphisms-section-Noetherian-valuative-criterion}

\noindent
本节讨论 Noether 情形下，仅使用离散赋值环来表述的代数叠态射的
（精细）赋值判据。判据有许多不同变体；今后将按需在此增补。

\medskip\noindent
设 $f : \mathcal{X} \to \mathcal{Y}$ 是代数叠（或代数空间或概形）
的态射。所谓{\it 精细}赋值判据，是指给定一个具有某些性质的态射
$\mathcal{U} \to \mathcal{X}$，并且只考察如下形式的实线图表中
虚线箭头的存在性或唯一性：
$$
\xymatrix{
\Spec(K) \ar[d] \ar[r] & \mathcal{U} \ar[r] & \mathcal{X} \ar[d] \\
\Spec(A) \ar[rr] \ar@{..>}[rru] & & \mathcal{Y}
}
$$
下文用这一术语描述
% P12-ERR-0099: see ../control/ERRATA_CANDIDATES.jsonl
迄今得到的结果。

\medskip\noindent
代数叠态射的非 Noether 赋值判据：
\begin{enumerate}
\item 《叠的态射》第
\ref{stacks-morphisms-section-valuative-second} 节
（用于对角态射的分离性）；
\item 《叠的态射》第
\ref{stacks-morphisms-section-valuative-diagonal} 节
（用于分离性）；
\item 《叠的态射》第
\ref{stacks-morphisms-section-valutive-criterion} 节
（用于万有闭性）；
\item 《叠的态射》第
\ref{stacks-morphisms-section-valutive-criterion-properness} 节
（用于固有性）。
\end{enumerate}
对于代数空间，有下列赋值判据：
\begin{enumerate}
\item 《空间的态射》第
\ref{spaces-morphisms-section-valuative-criterion-universally-closed} 节
（用于万有闭性）；
\item 《空间的态射》中的引理
\ref{spaces-morphisms-lemma-refined-valuative-criterion-universally-closed}
（万有闭性的精细判据）；
\item 《空间的态射》第
\ref{spaces-morphisms-section-valuative-separatedness} 节
（用于分离性）；
\item 《空间的态射》第
\ref{spaces-morphisms-section-valuative-criterion-properness} 节
（用于固有性）；
\item 《良好空间》第
\ref{decent-spaces-section-valuative-criterion-universally-closed} 节
（用于良好空间的万有闭性）；
\item 《良好空间》中的引理
\ref{decent-spaces-lemma-re-characterize-universally-closed}
（用于代数空间之间良好态射的万有闭性）；
\item 《空间的上同调》第
\ref{spaces-cohomology-section-Noetherian-valuative-criterion} 节包含
Noether 赋值判据：
\begin{enumerate}
\item 《空间的上同调》中的引理
\ref{spaces-cohomology-lemma-check-separated-dvr}
（使用离散赋值环判定分离性）；
\item 《空间的上同调》中的引理
\ref{spaces-cohomology-lemma-check-proper-dvr}
（使用离散赋值环判定固有性）；
\item 《空间的上同调》中的注
\ref{spaces-cohomology-remark-variant}
（讨论如何归约到完备离散赋值环）。
\end{enumerate}
\item 《空间的极限》第
\ref{spaces-limits-section-Noetherian-valuative-criterion} 节讨论
Noether 赋值判据：
\begin{enumerate}
\item 《空间的极限》中的引理
\ref{spaces-limits-lemma-Noetherian-dvr-valuative-separation}
（使用离散赋值环和泛点判定分离性）；
\item 《空间的极限》中的引理
\ref{spaces-limits-lemma-Noetherian-dvr-valuative-proper}
（使用离散赋值环和泛点判定固有性）；
\item 《空间的极限》中的引理
\ref{spaces-limits-lemma-check-universally-closed-Noetherian}
（使用离散赋值环判定万有闭性）。
\end{enumerate}
\item 《空间的极限》第
\ref{spaces-limits-section-refined-valuative-criteria} 节讨论
精细 Noether 赋值判据：
\begin{enumerate}
\item 《空间的极限》中的引理
\ref{spaces-limits-lemma-refined-valuative-criterion-proper} 和
\ref{spaces-limits-lemma-refined-valuative-criterion-universally-closed}
（使用离散赋值环给出的固有性精细判据）；
\item 《空间的极限》中的引理
\ref{spaces-limits-lemma-refined-valuative-criterion-separated}
（使用离散赋值环给出的分离性精细判据）。
\end{enumerate}
\end{enumerate}
对于概形，有下列赋值判据：
\begin{enumerate}
\item 《概形》第
\ref{schemes-section-valuative-criterion-universal-closedness} 节
（用于万有闭性）；
\item 《概形》第 \ref{schemes-section-valuative-separatedness} 节
（用于分离性）；
\item 《态射》第 \ref{morphisms-section-valuative-criteria} 节
（用于固有性）；
\item 《态射》中的引理
\ref{morphisms-lemma-refined-valuative-criterion-universally-closed}
（万有闭性的精细判据）；
\item 《极限》第 \ref{limits-section-Noetherian-valuative-criterion} 节
讨论 Noether 赋值判据：
\begin{enumerate}
\item 《极限》中的引理
\ref{limits-lemma-Noetherian-dvr-valuative-separation}
（使用离散赋值环和泛点判定分离性）；
\item 《极限》中的引理
\ref{limits-lemma-Noetherian-dvr-valuative-proper}
（使用离散赋值环和泛点判定固有性）；
\item 《极限》中的引理
\ref{limits-lemma-check-universally-closed-Noetherian}
（使用离散赋值环判定万有闭性）。
\end{enumerate}
\item 《极限》第 \ref{limits-section-refined-valuative-criteria} 节
讨论精细 Noether 赋值判据：
\begin{enumerate}
\item 《极限》中的引理
\ref{limits-lemma-refined-valuative-criterion-proper} 和
\ref{limits-lemma-refined-valuative-criterion-universally-closed}
（使用离散赋值环给出的固有性精细判据）；
\item 《极限》中的引理
\ref{limits-lemma-refined-valuative-criterion-separated}
（使用离散赋值环给出的分离性精细判据）。
\end{enumerate}
\item 《极限》第 \ref{limits-section-nagata-valuative} 节讨论
Noether 基底上的赋值判据；在该情形下，可以得到在基底上本质有限型的
离散赋值环。
\end{enumerate}
至此结束对先前结果的列举。

\medskip\noindent
本节的许多结果都可以（而且或许应该）诉诸下面的引理来证明，
尽管我们并非始终这样做。

\begin{lemma}
\label{stacks-more-morphisms-lemma-reach-point-closure-Noetherian}
设 $f : \mathcal{X} \to \mathcal{Y}$ 是代数叠的态射。假设
% P12-ERR-0100: see ../control/ERRATA_CANDIDATES.jsonl
$f$ 有限型且 $\mathcal{Y}$ 局部 Noether。设
$y \in |\mathcal{Y}|$ 是 $|f|$ 的像的闭包中的一点。
那么存在代数叠的交换图表
$$
\xymatrix{
\Spec(K) \ar[r] \ar[d] & \mathcal{X} \ar[d]^f \\
\Spec(A) \ar[r] & \mathcal{Y}
}
$$
其中 $A$ 是离散赋值环，$K$ 是其分式域，并且
$\Spec(A)$ 的闭点映到 $y$。
\end{lemma}

\begin{proof}
选取一个仿射概形 $V$、一点 $v \in V$，以及把 $v$ 映到 $y$ 的
光滑态射 $V \to \mathcal{Y}$。映射
$|V| \to |\mathcal{Y}|$ 是开的；并且由《叠的性质》中的引理
\ref{stacks-properties-lemma-points-cartesian}，
$|\mathcal{X} \times_\mathcal{Y} V| \to |V|$ 的像是
$|f|$ 的像的逆像。因此，点 $v$ 属于
$|\mathcal{X} \times_\mathcal{Y} V| \to |V|$ 的像的闭包。
若对 $\mathcal{X} \times_\mathcal{Y} V \to V$ 和点 $v$
证明该引理，则对 $f$ 和 $y$ 也得到该引理。这样便归结到下一段
所描述的情形。

\medskip\noindent
假设有引理中那样的 $f : \mathcal{X} \to Y$ 和 $y \in |Y|$，
其中 $Y$ 是 Noether 仿射概形。由于 $f$ 拟紧，可知
$\mathcal{X}$ 拟紧。因此可以选取一个仿射概形 $W$
以及一个满光滑态射 $W \to \mathcal{X}$。于是 $|f|$ 的像
与 $|W| \to |Y|$ 的像相同。这样便归结到概形的情形，
即《极限》中的引理
\ref{limits-lemma-reach-point-closure-Noetherian}。
\end{proof}

\begin{lemma}
\label{stacks-more-morphisms-lemma-check-separated-dvr}
设 $f : \mathcal{X} \to \mathcal{Y}$ 是代数叠的态射。假设
\begin{enumerate}
\item $\mathcal{Y}$ 局部 Noether；
\item $f$ 局部有限型且拟分离；
\item 对每个交换图表
$$
\xymatrix{
\Spec(K) \ar[r]_x \ar[d]_j & \mathcal{X} \ar[d]^f \\
\Spec(A) \ar[r]^y \ar@{-->}[ru] & \mathcal{Y}
}
$$
其中 $A$ 是离散赋值环、$K$ 是其分式域，并且对任意
$2$-箭头 $\gamma : y \circ j \to f \circ x$，虚线箭头范畴
（《叠的态射》中的定义
\ref{stacks-morphisms-definition-fill-in-diagram}）
或者为空，或者是恰有一个同构类的集合胚（setoid）。
\end{enumerate}
那么 $f$ 分离。
\end{lemma}

\begin{proof}
为证明 $f$ 分离，必须证明
$\Delta : \mathcal{X} \to \mathcal{X} \times_\mathcal{Y} \mathcal{X}$
固有。已经知道 $\Delta$ 可由代数空间表示且局部有限型
（《叠的态射》中的引理
\ref{stacks-morphisms-lemma-properties-diagonal}），并且拟紧、拟分离
（由 $f$ 拟分离的定义）。选取一个概形 $U$ 以及一个满光滑态射
$U \to \mathcal{X} \times_\mathcal{Y} \mathcal{X}$。置
$$
V = \mathcal{X} \times_{\Delta, \mathcal{X} \times_\mathcal{Y} \mathcal{X}} U
$$
只须证明代数空间的态射 $V \to U$ 固有
（《叠的性质》中的引理
\ref{stacks-properties-lemma-check-property-covering}）。
注意，$U$ 局部 Noether（使用《叠的态射》中的引理
\ref{stacks-morphisms-lemma-locally-finite-type-locally-noetherian}
以及 $U \to \mathcal{Y}$ 局部有限型这一事实），而
$V \to U$ 有限型且拟分离（它是 $\Delta$ 的基变换，并使用上面列出的
$\Delta$ 的性质）。应用《空间的上同调》中的引理
\ref{spaces-cohomology-lemma-check-proper-dvr}，只须证明：
给定交换图表
$$
\xymatrix{
\Spec(K) \ar[r]_v \ar[d]_j &
V \ar[d]^g \ar[r] &
\mathcal{X} \ar[d]^\Delta \\
\Spec(A) \ar[r]^u \ar@{-->}[ru] \ar@{..>}[rru] &
U \ar[r] &
\mathcal{X} \times_\mathcal{Y} \mathcal{X}
}
$$
其中 $A$ 是离散赋值环、$K$ 是其分式域，存在唯一使图表交换的
虚线箭头。由《叠的态射》中的引理
\ref{stacks-morphisms-lemma-cat-dotted-arrows-base-change}，
虚线箭头与点线箭头的范畴等价。假设 (3) 蕴含存在唯一的
点线箭头（在同构意义下）；参见《叠的态射》中的引理
\ref{stacks-morphisms-lemma-helper-diagonal}。因此如所需，
存在唯一的虚线箭头。
\end{proof}

\begin{lemma}
\label{stacks-more-morphisms-lemma-refined-valuative-criterion-proper}
设 $f : \mathcal{X} \to \mathcal{Y}$ 和
$h : \mathcal{U} \to \mathcal{X}$ 是代数叠的态射。
假设 $\mathcal{Y}$ 局部 Noether，$f$ 和 $h$ 都是有限型的，
$f$ 分离，并且 $|h| : |\mathcal{U}| \to |\mathcal{X}|$ 的像
在 $|\mathcal{X}|$ 中稠密。如果给定任意 $2$-交换图表
$$
\xymatrix{
\Spec(K) \ar[r]_-u \ar[d]_j & \mathcal{U} \ar[r]_h & \mathcal{X} \ar[d]^f \\
\Spec(A) \ar[rr]^-y & & \mathcal{Y}
}
$$
其中 $A$ 是分式域为 $K$ 的离散赋值环，并给定
% P12-ERR-0101: see ../control/ERRATA_CANDIDATES.jsonl
$\gamma : y \circ j \to f \circ h \circ u$，存在域扩张 $K'/K$
和支配 $A$ 的赋值环 $A' \subset K'$，使诱导图表
$$
\xymatrix{
\Spec(K') \ar[r]_-{x'} \ar[d]_{j'} & \mathcal{X} \ar[d]^f \\
\Spec(A') \ar[r]^-{y'} \ar@{..>}[ru] & \mathcal{Y}
}
$$
对于诱导的 $2$-箭头
$\gamma' : y' \circ j' \to f \circ x'$ 的虚线箭头范畴非空
（《叠的态射》中的定义
\ref{stacks-morphisms-definition-fill-in-diagram}），那么 $f$ 固有。
\end{lemma}

\begin{proof}
只须证明 $f$ 万有闭。设 $V \to \mathcal{Y}$ 是光滑态射，
其中 $V$ 是仿射概形。由《叠的性质》中的引理
\ref{stacks-properties-lemma-points-cartesian}，
$|\mathcal{U} \times_\mathcal{Y} V| \to
|\mathcal{X} \times_\mathcal{Y} V|$ 的像 $I$
是 $|h|$ 的像的逆像。由于
$|\mathcal{X} \times_\mathcal{Y} V| \to |\mathcal{X}|$ 是开映射
（《叠的态射》中的引理
\ref{stacks-morphisms-lemma-fppf-open}），可知 $I$ 在
$|\mathcal{X} \times_\mathcal{Y} V|$ 中稠密。
此外，由于虚线箭头范畴在基变换下表现良好
（《叠的态射》中的引理
\ref{stacks-morphisms-lemma-cat-dotted-arrows-base-change}），
关于（扩张后）虚线箭头存在性的假设被态射
$\mathcal{U} \times_\mathcal{Y} V \to
\mathcal{X} \times_\mathcal{Y} V \to V$ 继承。
因此，该引理的假设对态射
$\mathcal{U} \times_\mathcal{Y} V \to
\mathcal{X} \times_\mathcal{Y} V \to V$ 成立。于是可以假设
$\mathcal{Y}$ 是仿射概形。

\medskip\noindent
假设 $\mathcal{Y} = Y$ 是仿射概形。（从现在起，不必再操心
$2$-箭头 $\gamma$ 和 $\gamma'$；参见《叠的态射》中的引理
\ref{stacks-morphisms-lemma-cat-dotted-arrows-independent}。）
于是 $\mathcal{U}$ 拟紧。选取一个仿射概形 $U$ 以及一个满光滑态射
$U \to \mathcal{U}$。可以并且确实用 $U$ 代替 $\mathcal{U}$。
因此，可以假设 $\mathcal{U}$ 是仿射概形。

\medskip\noindent
假设 $\mathcal{Y} = Y$ 和 $\mathcal{U} = U$ 都是仿射概形。
由 Chow 引理（定理 \ref{stacks-more-morphisms-theorem-chow-finite-type}），
可以选取一个满固有态射 $X \to \mathcal{X}$，其中 $X$ 是代数空间。
下文将使用 $X \to Y$ 作为分离态射的复合是分离的这一事实。
考虑代数空间 $W = X \times_\mathcal{X} U$。投影态射
$W \to X$ 是有限型的。可以用 $W \to X$ 的概形论像代替 $X$，
从而假设 $|W|$ 在 $|X|$ 中的像在 $|X|$ 中稠密
（这里使用 $|h|$ 的像在 $|\mathcal{X}|$ 中稠密，所以作此代替后，
态射 $X \to \mathcal{X}$ 仍然是满射）。我们断言，对于每个
实线交换图表
$$
\xymatrix{
\Spec(K) \ar[r] \ar[d] & W \ar[r] & X \ar[d] \\
\Spec(A) \ar[rr] \ar@{..>}[rru] & & Y
}
$$
其中 $A$ 是分式域为 $K$ 的离散赋值环，存在使图表交换的点线箭头。
首先，只须在以一个扩张代替 $K$、并以该扩张中支配 $A$ 的一个赋值环
代替 $A$ 后证明点线箭头存在；参见《空间的态射》中的引理
\ref{spaces-morphisms-lemma-push-down-solution}。
由该引理的假设，得到扩张 $K'/K$、支配 $A$ 的赋值环
$A' \subset K'$，以及一个箭头
$\Spec(A') \to \mathcal{X}$；该箭头提升复合
$\Spec(A') \to \Spec(A) \to Y$，并与复合
$\Spec(K') \to \Spec(K) \to W \to X$ 相容。
由于 $X \to \mathcal{X}$ 固有，可以使用固有性的赋值判据
（《叠的态射》中的引理
\ref{stacks-morphisms-lemma-criterion-proper}），
找到扩张 $K''/K'$、支配 $A'$ 的赋值环 $A'' \subset K''$，
以及态射 $\Spec(A'') \to X$；该态射提升复合
$\Spec(A'') \to \Spec(A') \to \mathcal{X}$，并与复合
$\Spec(K'') \to \Spec(K') \to \Spec(K) \to X$ 相容。
于是
% P12-ERR-0102: see ../control/ERRATA_CANDIDATES.jsonl
$K''/K$、$A'' \subset K''$ 和态射 $\Spec(A'') \to X$
是上述问题的一个解，故断言成立。

\medskip\noindent
由代数空间情形的该引理（《空间的极限》中的引理
\ref{spaces-limits-lemma-refined-valuative-criterion-proper}），
该断言蕴含态射 $X \to Y$ 固有。由《叠的态射》中的引理
\ref{stacks-morphisms-lemma-image-proper-is-proper}，
得出 $\mathcal{X} \to Y$ 如所需固有。
\end{proof}

\begin{lemma}
\label{stacks-more-morphisms-lemma-refined-valuative-criterion-separated}
设 $f : \mathcal{X} \to \mathcal{Y}$ 和
$h : \mathcal{U} \to \mathcal{X}$ 是代数叠的态射。假设
$\mathcal{Y}$ 局部 Noether，$f$ 局部有限型且拟分离，
$h$ 有限型，并且
$|h| : |\mathcal{U}| \to |\mathcal{X}|$ 的像在
$|\mathcal{X}|$ 中稠密。如果对任意给定的 $2$-交换图表
$$
\xymatrix{
\Spec(K) \ar[r]_-u \ar[d]_j & \mathcal{U} \ar[r]_h & \mathcal{X} \ar[d]^f \\
\Spec(A) \ar[rr]^-y \ar@{..>}[rru] & & \mathcal{Y}
}
$$
其中 $A$ 是分式域为 $K$ 的离散赋值环，并且
$\gamma : y \circ j \to f \circ h \circ u$，虚线箭头范畴
或者为空，或者是恰有一个同构类的集合胚（setoid），那么 $f$ 分离。
\end{lemma}

\begin{proof}
我们必须证明 $\Delta$ 是固有态射。先假设 $\Delta$ 分离。
那么可以将引理 \ref{stacks-more-morphisms-lemma-refined-valuative-criterion-proper}
应用于态射 $\mathcal{U} \to \mathcal{X}$ 和
$\Delta : \mathcal{X} \to
\mathcal{X} \times_\mathcal{Y} \mathcal{X}$。
注意，由于 $f$ 拟分离，$\Delta$ 拟紧。当然，$\Delta$ 局部有限型
（这对任意对角态射都成立，参见《叠的态射》中的引理
\ref{stacks-morphisms-lemma-properties-diagonal}）。最后，假设给定一个
$2$-交换图表
$$
\xymatrix{
\Spec(K) \ar[r]_-u \ar[d]_j &
\mathcal{U} \ar[r]_h &
\mathcal{X} \ar[d]^\Delta \\
\Spec(A) \ar[rr]^-y \ar@{..>}[rru] & &
\mathcal{X} \times_\mathcal{Y} \mathcal{X}
}
$$
其中 $A$ 是分式域为 $K$ 的离散赋值环，并且
$\gamma : y \circ j \to \Delta \circ h \circ u$。
由《叠的态射》中的引理
\ref{stacks-morphisms-lemma-helper-diagonal} 以及本引理中的假设，
可知存在唯一的点线箭头。这证明引理
\ref{stacks-more-morphisms-lemma-refined-valuative-criterion-proper} 的最后一个假设成立，
从而得到结论。

\medskip\noindent
在一般情形下，只须证明 $\Delta$ 分离，因为这样便回到前一种情形。
事实上，我们断言本引理的假设对
$$
\mathcal{U} \to \mathcal{X}
\quad\text{和}\quad
\Delta :
\mathcal{X} \to
\mathcal{X} \times_\mathcal{Y} \mathcal{X}
$$
成立。确实，由于 $\Delta$ 可由代数空间表示，上一段那种图表的
虚线箭头范畴是一个集合胚（例如参见《叠的态射》中的引理
\ref{stacks-morphisms-lemma-cat-dotted-arrows}）。上一段的论证表明，
这些范畴或者为空，或者恰有一个同构类。因此 $\Delta$ 分离。
\end{proof}

\begin{lemma}
\label{stacks-more-morphisms-lemma-valuative-criterion-universally-closed}
设 $f : \mathcal{X} \to \mathcal{Y}$ 是代数叠的态射。
假设 $\mathcal{Y}$ 局部 Noether，并且 $f$ 有限型。如果对任意给定的
$2$-交换图表
$$
\xymatrix{
\Spec(K) \ar[r]_-x \ar[d]_j & \mathcal{X} \ar[d]^f \\
\Spec(A) \ar[r]^-y & \mathcal{Y}
}
$$
其中 $A$ 是分式域为 $K$ 的离散赋值环，
% P12-ERR-0103: see ../control/ERRATA_CANDIDATES.jsonl
并且 $\gamma : y \circ j \to f \circ x$，存在域扩张 $K'/K$
以及支配 $A$ 的赋值环 $A' \subset K'$，使得诱导图表
$$
\xymatrix{
\Spec(K') \ar[r]_-{x'} \ar[d]_{j'} & \mathcal{X} \ar[d]^f \\
\Spec(A') \ar[r]^-{y'} \ar@{..>}[ru] & \mathcal{Y}
}
$$
连同诱导 $2$-箭头 $\gamma' : y' \circ j' \to f \circ x'$ 的
虚线箭头范畴非空（《叠的态射》中的定义
\ref{stacks-morphisms-definition-fill-in-diagram}），那么 $f$ 万有闭。
\end{lemma}

\begin{proof}
设 $V \to \mathcal{Y}$ 是光滑态射，其中 $V$ 是仿射概形。
虚线箭头范畴在基变换下表现良好（《叠的态射》中的引理
\ref{stacks-morphisms-lemma-cat-dotted-arrows-base-change}）。
因此，关于（扩张后）虚线箭头存在性的假设被态射
$\mathcal{X} \times_\mathcal{Y} V \to V$ 继承。
所以本引理的假设对态射
$\mathcal{X} \times_\mathcal{Y} V \to V$ 成立。
于是可以假设 $\mathcal{Y}$ 是仿射概形。

\medskip\noindent
假设 $\mathcal{Y} = Y$ 是 Noether 仿射概形。（从现在起，
不必再操心 $2$-箭头 $\gamma$ 和 $\gamma'$；参见《叠的态射》中的引理
\ref{stacks-morphisms-lemma-cat-dotted-arrows-independent}。）
为证明 $f$ 万有闭，由《叠的极限》中的引理
\ref{stacks-limits-lemma-test-universally-closed}，只须证明对所有 $n$，
映射 $|\mathcal{X} \times \mathbf{A}^n| \to
|Y \times \mathbf{A}^n|$ 都是闭映射。由于本引理的假设被乘积态射
$\mathcal{X} \times \mathbf{A}^n \to Y \times \mathbf{A}^n$
继承（略去细节），问题约化为证明
$|\mathcal{X}| \to |Y|$ 是闭映射。

\medskip\noindent
假设 $Y$ 是 Noether 仿射概形。设
$T \subset |\mathcal{X}|$ 是闭子集。我们必须证明 $T$ 在
$|Y|$ 中的像是闭的。可以用 $T$ 上的约化诱导闭子空间结构
代替 $\mathcal{X}$；略去验证关于虚线箭头存在性的性质在该代替下
保持不变。因此，问题约化为证明
$|\mathcal{X}| \to |Y|$ 的像是闭的。

\medskip\noindent
设 $y \in |Y|$ 是 $|\mathcal{X}| \to |Y|$ 的像之闭包中的一点。
由引理 \ref{stacks-more-morphisms-lemma-reach-point-closure-Noetherian}，可以选取交换图表
$$
\xymatrix{
\Spec(K) \ar[r] \ar[d] & \mathcal{X} \ar[d]^f \\
\Spec(A) \ar[r] & Y
}
$$
其中 $A$ 是离散赋值环，$K$ 是其分式域，并且
$\Spec(A)$ 的闭点映到 $y$。由本引理的假设立即可知，
$y$ 位于 $|\mathcal{X}| \to |Y|$ 的像中，证明完成。
\end{proof}

\section{模空间}
\label{stacks-more-morphisms-section-moduli-spaces}

\noindent
本节讨论从代数叠到代数空间的态射
$f : \mathcal{X} \to Y$。在适当假设下，称 $Y$ 是
$\mathcal{X}$ 的一个{\it 模空间}。如果
$\mathcal{X} = [U/R]$ 是一个表现，那么可得到一个
$R$-不变态射 $U \to Y$，并且（在适当假设下）$Y$ 是群胚
$(U, R, s, t, c)$ 的一个{\it 商}。关于不同类型商的讨论，
可从《群胚的商》第
\ref{groupoids-quotients-section-introduction} 节开始查阅。

\begin{definition}
\label{stacks-more-morphisms-definition-categorical-quotient}
设 $\mathcal{X}$ 是代数叠。设
$f : \mathcal{X} \to Y$ 是到代数空间 $Y$ 的态射。
\begin{enumerate}
\item 如果任意到代数空间 $W$ 的态射
$\mathcal{X} \to W$ 都唯一地经由 $f$ 分解，则称 $f$ 是一个
{\it 范畴模空间}。
\item 如果对代数空间的任意平坦态射 $Y' \to Y$，基变换
$f' : Y' \times_Y \mathcal{X} \to Y'$ 都是范畴模空间，则称 $f$
是一个{\it 一致范畴模空间}。
\end{enumerate}
设 $\mathcal{C}$ 是代数空间范畴的一个全子范畴。
\begin{enumerate}
\item[(3)] 如果 $Y \in \Ob(\mathcal{C})$，并且对任意
$W \in \Ob(\mathcal{C})$，每个态射 $\mathcal{X} \to W$
都唯一地经由 $f$ 分解，则称 $f$ 是一个
{\it $\mathcal{C}$ 中的范畴模空间}。
% P12-ERR-0104: see ../control/ERRATA_CANDIDATES.jsonl
\item[(4)] 如果 $Y \in \Ob(\mathcal{C})$，并且对
$\mathcal{C}$ 中的每个平坦态射 $Y' \to Y$，基变换
$f' : Y' \times_Y \mathcal{X} \to Y'$ 都是
$\mathcal{C}$ 中的范畴模空间，则称其为一个
{\it $\mathcal{C}$ 中的一致范畴模空间}。
\end{enumerate}
\end{definition}

\noindent
由 Yoneda 引理，范畴模空间若存在便是唯一的。下面将它与此前为商
引入的语言相对应。

\begin{lemma}
\label{stacks-more-morphisms-lemma-quotient-compare}
设 $(U, R, s, t, c)$ 是代数空间中的群胚，其中
$s, t : R \to U$ 平坦且局部有限表现。考虑代数叠
$\mathcal{X} = [U/R]$。给定代数空间 $Y$，态射
$f : \mathcal{X} \to Y$ 与 $R$-不变态射
$\phi : U \to Y$ 之间存在一个 $1$-到-$1$ 对应。
\end{lemma}

\begin{proof}
《可表示性判据》中的定理
\ref{criteria-theorem-flat-groupoid-gives-algebraic-stack}
告诉我们 $\mathcal{X}$ 是代数叠。给定态射
$f : \mathcal{X} \to Y$，令 $\phi : U \to Y$ 为复合
$U \to \mathcal{X} \to Y$。由于
$R = U \times_\mathcal{X} U$（《空间中的群胚》中的引理
\ref{spaces-groupoids-lemma-quotient-stack-2-cartesian}），
立即可知 $\phi$ 是 $R$-不变的。反之，如果
$\phi : U \to Y$ 是到代数空间的 $R$-不变态射，那么由
《空间中的群胚》中的引理
\ref{spaces-groupoids-lemma-quotient-stack-2-coequalizer}，
可得到态射 $f : \mathcal{X} \to Y$。也可以利用
《空间中的群胚》中的引理
\ref{spaces-groupoids-lemma-quotient-stack-objects}
对商叠的显式描述，从 $\phi$ 构造 $f$。
\end{proof}

\begin{lemma}
\label{stacks-more-morphisms-lemma-categorical-quotient-compare}
% P12-ERR-0105: see ../control/ERRATA_CANDIDATES.jsonl
采用引理 \ref{stacks-more-morphisms-lemma-quotient-compare} 中的假设和记号。
那么 $f$ 是（一致）范畴模空间，当且仅当 $\phi$ 是（一致）范畴商。
对于全子范畴中的模空间也有类似结论。
\end{lemma}

\begin{proof}
由引理 \ref{stacks-more-morphisms-lemma-quotient-compare} 中建立的 $1$-到-$1$ 对应，立即可知
$f$ 是范畴模空间，当且仅当 $\phi$ 是范畴商
（《群胚的商》中的定义
\ref{groupoids-quotients-definition-categorical}）。
如果 $Y' \to Y$ 是态射，那么
$U' = Y' \times_Y U \to Y' \times_Y \mathcal{X} = \mathcal{X}'$
作为 $U \to \mathcal{X}$ 的基变换，是满、平坦且局部有限表现的态射
（《可表示性判据》中的引理
\ref{criteria-lemma-flat-quotient-flat-presentation}）。
又由纤维积的传递性，$R' = Y' \times_Y R$ 等于
$U' \times_{\mathcal{X}'} U'$。因此
$\mathcal{X}' = [U'/R']$；参见《代数叠》中的备注
\ref{algebraic-remark-flat-fp-presentation}。所以，把我们的情形基变换到
$Y'$ 后，仍得到本引理陈述中的情形。由此立即可知，$f$ 是一致范畴
模空间，当且仅当 $\phi$ 是一致范畴商。
\end{proof}

\begin{lemma}
\label{stacks-more-morphisms-lemma-check-uniform-categorical-quotient-on-affines}
设 $f : \mathcal{X} \to Y$ 是从代数叠到代数空间的态射。
如果对每个仿射概形 $Y'$ 和平坦态射 $Y' \to Y$，基变换
$f' : Y' \times_Y \mathcal{X} \to Y'$ 都是范畴模空间，
那么 $f$ 是一致范畴模空间。
\end{lemma}

\begin{proof}
选取一个 \'etale 覆盖 $\{Y_i \to Y\}$，其中 $Y_i$ 是仿射概形。
% P12-ERR-0106: see ../control/ERRATA_CANDIDATES.jsonl
对每个 $i$ 和 $j$，选取一个仿射开覆盖
$Y_i \times_Y Y_j = \bigcup Y_{ijk}$。置
$\mathcal{X}_i = Y_i \times_Y \mathcal{X}$ 和
$\mathcal{X}_{ijk} = Y_{ijk} \times_Y \mathcal{X}$。
设 $g : \mathcal{X} \to W$ 是到代数空间的态射。考虑图表
$$
\xymatrix{
\mathcal{X}_i \ar[r] \ar[d] & \mathcal{X} \ar[d] \ar[r]_g & W \\
Y_i \ar[r] \ar@{..>}[rru] & Y
}
$$
$\mathcal{X}_i \to Y_i$ 是范畴模空间这一假设给出唯一的
点线箭头 $h_i : Y_i \to W$。$\mathcal{X}_{ijk} \to Y_{ijk}$
是范畴模空间这一假设说明
% P12-ERR-0107: see ../control/ERRATA_CANDIDATES.jsonl
$h_i$ 和 $h_j$ 在 $Y_{ijk}$ 上的限制相等。因此，$h_i$ 和 $h_j$
在 $Y_i \times_Y Y_j$ 上相同。由于
$Y = \coprod Y_i / \coprod Y_i \times_Y Y_j$
（由《空间》第 \ref{spaces-section-presentations} 节），可知存在唯一的
态射 $Y \to W$，使得 $g$ 经由它分解。因此 $f$ 是范畴模空间。
相同论证可在平坦基变换后应用，故 $f$ 是一致范畴模空间。
\end{proof}

\section{Keel--Mori 定理}
\label{stacks-more-morphisms-section-Keel-Mori}

\noindent
本节开始在代数叠的框架下讨论 Keel 和 Mori 的定理。
关于相关文献的讨论，请参见《文献指南》中的第
\ref{guide-subsection-coarse-moduli-spaces} 小节。

\begin{definition}
\label{stacks-more-morphisms-definition-well-nigh-affine}
设 $\mathcal{X}$ 是代数叠。如果存在仿射概形 $U$ 和一个满、平坦、
有限且有限表现的态射 $U \to \mathcal{X}$，则称
$\mathcal{X}$ 是{\it 近乎仿射的}。
\end{definition}

\noindent
我们给这个性质取了一个多少有些荒唐的名字，因为并不打算过多使用它。

\begin{lemma}
\label{stacks-more-morphisms-lemma-well-nigh-affine}
设 $\mathcal{X}$ 是代数叠。下列条件等价：
\begin{enumerate}
\item $\mathcal{X}$ 近乎仿射；
\item 存在群胚概形 $(U, R, s, t, c)$，其中 $U$ 和 $R$ 仿射，
$s, t : R \to U$ 有限局部自由，并且 $\mathcal{X} = [U/R]$。
\end{enumerate}
若这些条件成立，则 $\mathcal{X}$ 拟紧、拟-DM 且分离。
\end{lemma}

\begin{proof}
假设 $\mathcal{X}$ 近乎仿射。选取仿射概形 $U$ 和一个满、平坦、
有限且有限表现的态射 $U \to \mathcal{X}$。置
$R = U \times_\mathcal{X} U$。于是得到代数空间中的群胚
$(U, R, s, t, c)$ 和同构 $[U/R] \to \mathcal{X}$；参见
《代数叠》中的引理 \ref{algebraic-lemma-map-space-into-stack}
及备注 \ref{algebraic-remark-flat-fp-presentation}。由于
$s, t : R \to U$ 是平坦、有限且有限表现的态射
% P12-ERR-0108: see ../control/ERRATA_CANDIDATES.jsonl
（作为 $U \to \mathcal{X})$ 的基变换），可知
$s, t$ 有限局部自由（《态射》中的引理
\ref{morphisms-lemma-finite-flat}）。这蕴含 $R$ 仿射
（因为有限态射是仿射的），故 (2) 成立。

\medskip\noindent
% P12-ERR-0109: see ../control/ERRATA_CANDIDATES.jsonl
假设有一个群胚概形 $(U, R, s, t, c)$，其中 $U$ 和 $R$ 仿射，
并且 $s, t : R \to U$ 有限局部自由。置
$\mathcal{X} = [U/R]$。由《可表示性判据》中的定理
\ref{criteria-theorem-flat-groupoid-gives-algebraic-stack}，
$\mathcal{X}$ 是代数叠（严格来说，这里并不需要这一点，但无论怎样
强调其正确性都不为过）。由《可表示性判据》中的引理
\ref{criteria-lemma-flat-quotient-flat-presentation}，态射
$U \to \mathcal{X}$ 满、平坦且局部有限表现。因此，根据《叠的性质》
中的引理 \ref{stacks-properties-lemma-check-property-covering}，
可以通过检验投影 $U \times_\mathcal{X} U \to U$ 是否有限，来检验
$U \to \mathcal{X}$ 是否有限。由《空间中的群胚》中的引理
\ref{spaces-groupoids-lemma-quotient-stack-2-cartesian}，
$U \times_\mathcal{X} U = R$，故确实如此。因此
$\mathcal{X}$ 近乎仿射。

\medskip\noindent
最后证明末一断言。由《叠的性质》中的引理
\ref{stacks-properties-lemma-quasi-compact-stack}，
$\mathcal{X}$ 拟紧。由《叠的态射》中的引理
\ref{stacks-morphisms-lemma-properties-diagonal-from-presentation}，
$\mathcal{X} = [U/R]$ 拟-DM 且分离。
\end{proof}

\begin{lemma}
\label{stacks-more-morphisms-lemma-affine-over-well-nigh-affine}
设代数叠 $\mathcal{X}$ 近乎仿射。
\begin{enumerate}
\item 如果 $\mathcal{X}$ 是代数空间，那么它是仿射的。
\item 如果 $\mathcal{X}' \to \mathcal{X}$ 是代数叠的仿射态射，
那么 $\mathcal{X}'$ 近乎仿射。
\end{enumerate}
\end{lemma}

\begin{proof}
% P12-ERR-0110: see ../control/ERRATA_CANDIDATES.jsonl
(1) 立即由《空间的极限》中的引理
\ref{spaces-limits-lemma-affine} 得出。不过，这未免大材小用，因为
把引理 \ref{stacks-more-morphisms-lemma-well-nigh-affine} 与《群胚》中的命题
\ref{groupoids-proposition-finite-flat-equivalence} 结合也可得到 (1)。

\medskip\noindent
为证明 (2)，选取仿射概形 $U$ 和一个满、平坦、有限且有限表现的态射
$U \to \mathcal{X}$。那么
$U' = \mathcal{X}' \times_\mathcal{X} U$ 到 $U$ 有仿射态射
（《叠的态射》中的引理
\ref{stacks-morphisms-lemma-base-change-affine}）。因此 $U'$ 是仿射概形。
当然，$U' \to \mathcal{X}'$ 作为 $U \to \mathcal{X}$ 的基变换，
是满、平坦、有限且有限表现的。
\end{proof}

\begin{lemma}
\label{stacks-more-morphisms-lemma-well-nigh-affine-moduli-space}
设代数叠 $\mathcal{X}$ 近乎仿射。在仿射概形范畴中存在一致范畴模空间
$$
f : \mathcal{X} \longrightarrow M
$$
而且 $f$ 分离、拟紧，并且是万有同胚。
\end{lemma}

\begin{proof}
写作 $\mathcal{X} = [U/R]$，其中 $(U, R, s, t, c)$
如引理 \ref{stacks-more-morphisms-lemma-well-nigh-affine} 所述。设 $C$ 是 $U$ 上
$R$-不变函数组成的环；参见《群胚》第
\ref{groupoids-section-finite-flat} 节。置 $M = \Spec(C)$。
由引理 \ref{stacks-more-morphisms-lemma-quotient-compare}，$R$-不变态射 $U \to M$
对应于态射 $f : \mathcal{X} \to M$。《概形》中的引理
\ref{schemes-lemma-morphism-into-affine} 给出的到仿射概形态射的刻画
立即保证
% P12-ERR-0111: see ../control/ERRATA_CANDIDATES.jsonl
$\phi : U \to M$ 是仿射概形范畴中的范畴商。因此 $f$ 是仿射概形
范畴中的范畴模空间（引理
\ref{stacks-more-morphisms-lemma-categorical-quotient-compare}）。

\medskip\noindent
由引理 \ref{stacks-more-morphisms-lemma-well-nigh-affine}，$\mathcal{X}$ 分离，故由
《叠的态射》中的引理
\ref{stacks-morphisms-lemma-compose-after-separated}，$f$ 分离。

\medskip\noindent
由于 $U \to \mathcal{X}$ 满，并且 $U \to M$ 拟紧，由《叠的态射》
中的引理 \ref{stacks-morphisms-lemma-surjection-from-quasi-compact}，
可知 $f$ 拟紧。

\medskip\noindent
由《群胚》中的引理 \ref{groupoids-lemma-integral-over-invariants}，复合
$$
U \to \mathcal{X} \to M
$$
是仿射概形的整态射。特别地，它万有闭
（《态射》中的引理
\ref{morphisms-lemma-integral-universally-closed}）。由于
$U \to \mathcal{X}$ 满，可知 $\mathcal{X} \to M$ 万有闭
（《叠的态射》中的引理
\ref{stacks-morphisms-lemma-image-proper-is-proper}）。为得出
$\mathcal{X} \to M$ 是万有同胚，只须证明它万有双射，即它满且万有单射。

\medskip\noindent
由《叠的态射》中的引理
\ref{stacks-morphisms-lemma-points-presentation}，有
$|\mathcal{X}| = |U|/|R|$。因此，由《群胚》中的引理
\ref{groupoids-lemma-points}，$|f|$ 满，甚至是双射。

\medskip\noindent
设 $C \to C'$ 是环映射。设 $(U', R', s', t', c')$ 是
$(U, R, s, t, c)$ 沿 $M' = \Spec(C') \to M$ 的基变换。
置 $\mathcal{X}' = [U'/R']$。由《群胚的商》中的引理
\ref{groupoids-quotients-lemma-base-change-quotient-stack}，有
$M' \times_M \mathcal{X} = \mathcal{X}'$。设 $C^1$ 是 $U'$ 上
$R'$-不变函数组成的环。置 $M^1 = \Spec(C^1)$，并考虑图表
$$
\xymatrix{
\mathcal{X}' \ar[d]^{f'} \ar[r] & \mathcal{X} \ar[dd]^f \\
M^1 \ar[d] \\
M' \ar[r] & M
}
$$
由《群胚》中的引理 \ref{groupoids-lemma-invariants-base-change}
和《代数》中的引理 \ref{algebra-lemma-universally-bijective}，态射
$M^1 \to M'$ 是同胚。另一方面，把上一段应用于
$(U', R', s', t', c')$，可知 $|f'|$ 是双射。我们得出，沿任意仿射
概形作基变换后，$f$ 都诱导点集之间的双射。因此，由《叠的态射》中的
引理 \ref{stacks-morphisms-lemma-universally-injective-local}，
$f$ 万有单射。

\medskip\noindent
最后，还须证明 $f$ 是仿射概形范畴中的
% P12-ERR-0112: see ../control/ERRATA_CANDIDATES.jsonl
一致模空间。这由以上讨论以及下述事实得出：如果环映射
$C \to C'$ 平坦，那么由《群胚》中的引理
\ref{groupoids-lemma-invariants-base-change}，$C' \to C^1$ 是同构。
\end{proof}

\begin{lemma}
\label{stacks-more-morphisms-lemma-well-nigh-affine-moduli-space-etale}
设 $h : \mathcal{X}' \to \mathcal{X}$ 是代数叠的态射。假设
$\mathcal{X}'$ 和 $\mathcal{X}$ 近乎仿射，$h$ 是 \'etale 的，
并且 $h$ 在自同构群上诱导同构（《叠的态射》中的备注
\ref{stacks-morphisms-remark-identify-automorphism-groups}）。
那么存在笛卡儿图表
$$
\xymatrix{
\mathcal{X}' \ar[d] \ar[r] & \mathcal{X} \ar[d] \\
M' \ar[r] & M
}
$$
其中 $M' \to M$ 是 \'etale 的，竖直箭头是引理
\ref{stacks-more-morphisms-lemma-well-nigh-affine-moduli-space} 中构造的模空间。
\end{lemma}

\begin{proof}
由《叠的态射》中的诸引理
\ref{stacks-morphisms-lemma-aut-iso-unramified} 和
\ref{stacks-morphisms-lemma-stabilizer-preserving}，注意到 $h$ 可由
代数空间表示。选取仿射概形 $U$ 和一个满、平坦、有限且有限表现的态射
$U \to \mathcal{X}$。于是
$U' = \mathcal{X}' \times_\mathcal{X} U$ 是代数空间，并且有一个有限
（特别地，仿射）态射 $U' \to \mathcal{X}'$。由引理
\ref{stacks-more-morphisms-lemma-affine-over-well-nigh-affine}，可知 $U'$ 仿射。
置 $R = U \times_\mathcal{X} U$ 和
$R' = U' \times_{\mathcal{X}'} U'$，便得到群胚
$(U, R, s, t, c)$ 和 $(U', R', s', t', c')$，使得
$\mathcal{X} = [U/R]$ 且 $\mathcal{X}' = [U'/R']$；参见引理
\ref{stacks-more-morphisms-lemma-well-nigh-affine} 的证明。
% P12-ERR-0113: see ../control/ERRATA_CANDIDATES.jsonl
可见下列图表
$$
\xymatrix{
R' \ar[d]_{s'} \ar[r]_f & R \ar[d]^s \\
U' \ar[r]^f & U
}
\quad
\quad
\xymatrix{
R' \ar[d]_{t'} \ar[r]_f & R \ar[d]^t \\
U' \ar[r]^f & U
}
\quad
\quad
\xymatrix{
G' \ar[d] \ar[r]_f & G \ar[d] \\
U' \ar[r]^f & U
}
$$
都是笛卡儿的，其中 $G$ 和 $G'$ 是稳定子群概形。前两个图表的断言
由纤维积的传递性得出；最后一个图表的断言则是因为它是同构
$\mathcal{I}_{\mathcal{X}'} \to
\mathcal{X}' \times_\mathcal{X} \mathcal{I}_\mathcal{X}$ 的拉回
（由已经使用过的《叠的态射》中的引理
\ref{stacks-morphisms-lemma-aut-iso-unramified}）。回顾引理
\ref{stacks-more-morphisms-lemma-well-nigh-affine-moduli-space} 中的构造：$M$，相应地，\ $M'$，
分别是 $U$ 上 $R$-不变函数环的谱和相应地，\ $U'$ 上
$R'$-不变函数环的谱。
因此，由《群胚》中的引理 \ref{groupoids-lemma-etale}，
$M' \to M$ 是 \'etale 的，并且 $U' = M' \times_M U$。
由此可得
% P12-ERR-0114: see ../control/ERRATA_CANDIDATES.jsonl
$R' = M' \times_M U$；换言之，群胚 $(U', R', s', t', c')$ 是
$(U, R, s, t, c)$ 沿 $M' \to M$ 的基变换。由《群胚的商》中的引理
\ref{groupoids-quotients-lemma-base-change-quotient-stack}，
这蕴含本引理中的图表为笛卡儿图表。
\end{proof}

\begin{lemma}
\label{stacks-more-morphisms-lemma-moduli-space-finite-affine}
设代数叠 $\mathcal{X}$ 近乎仿射。引理
\ref{stacks-more-morphisms-lemma-well-nigh-affine-moduli-space} 中的态射
$$
f : \mathcal{X} \longrightarrow M
$$
是一致范畴模空间。
\end{lemma}

\begin{proof}
我们已经知道
% P12-ERR-0115: see ../control/ERRATA_CANDIDATES.jsonl
$M$ 是仿射概形范畴中的一致范畴模空间。由引理
\ref{stacks-more-morphisms-lemma-check-uniform-categorical-quotient-on-affines}，只须证明：
对于仿射概形的任意平坦态射 $M' \to M$，基变换
$f' : M' \times_M \mathcal{X} \to M'$ 是范畴模空间。
由引理 \ref{stacks-more-morphisms-lemma-affine-over-well-nigh-affine}，注意到
$\mathcal{X}' = M' \times_M \mathcal{X}$ 近乎仿射。
% P12-ERR-0116: see ../control/ERRATA_CANDIDATES.jsonl
如此，以 $\mathcal{X}'$ 代替 $\mathcal{X}$ 并以 $M'$ 代替 $M$ 后，
问题约化为证明 $f$ 是范畴模空间。

\medskip\noindent
设 $g : \mathcal{X} \to Y$ 是态射，其中 $Y$ 是代数空间。必须证明
$g = h \circ f$，其中 $h : M \to Y$ 是唯一的态射。

\medskip\noindent
唯一性。假设有两个态射 $h_i : M \to Y$，满足
$g = h_1 \circ f = h_2 \circ f$。设 $M' \subset M$ 是 $h_1$ 和
$h_2$ 的等化子。那么 $M' \to M$ 是单态射，并且
$f : \mathcal{X} \to M$ 经由 $M'$ 分解。因此 $M' \to M$
是万有同胚。由此可知 $M'$ 仿射（《态射》中的引理
\ref{morphisms-lemma-universal-homeomorphism}）。但由于
$f : \mathcal{X} \to M$ 是仿射概形范畴中的范畴模空间，
于是 $M' = M$。

\medskip\noindent
存在性。下面将证明：对每个 $p \in M$，存在笛卡儿方形
$$
\xymatrix{
\mathcal{X}' \ar[r] \ar[d] & \mathcal{X} \ar[d] \\
M' \ar[r] & M
}
$$
其中 $M' \to M$ 是仿射概形之间的 \'etale 态射，$p$ 位于其像中，
并且复合 $\mathcal{X}' \to \mathcal{X} \to Y$ 经由 $M'$ 分解。
这意味着可以在 $M$ 上 \'etale 局部地构造映射 $h : M \to Y$。
由于 $Y$ 是 \'etale 拓扑的层，再由以上证明的唯一性，这便足够了
（略去一个小细节）。

\medskip\noindent
设 $y \in |Y|$ 是 $p$ 的像。设
$(V, v) \to (Y, y)$ 是 \'etale 态射，其中 $V$ 仿射。考虑
$\mathcal{X}' = V \times_Y \mathcal{X}$。注意，
$\mathcal{X}' \to \mathcal{X}$ 作为 $V \to Y$ 的基变换，是分离且
\'etale 的。此外，$\mathcal{X}' \to \mathcal{X}$ 在自同构群上诱导
同构（《叠的态射》中的备注
\ref{stacks-morphisms-remark-identify-automorphism-groups}），因为
$V \to Y$ 具有这一性质；参见《叠的态射》中的引理
\ref{stacks-morphisms-lemma-base-change-stabilizer-preserving}。
选取如引理 \ref{stacks-more-morphisms-lemma-well-nigh-affine} 中的表现
$\mathcal{X} = [U/R]$。置
$U' = \mathcal{X}' \times_\mathcal{X} U = V \times_Y U$，并选取
$u' \in U'$，使它映到 $p$ 和 $v$（这是可行的，由《空间的性质》中的
引理 \ref{spaces-properties-lemma-points-cartesian}）。由于
$U' \to U$ 分离且 \'etale，由《态射进阶》中的引理
\ref{more-morphisms-lemma-separated-locally-quasi-finite-over-affine}，
可知 $U'$ 的任意有限点集都包含在一个仿射开集中。另一方面，态射
$U' \to \mathcal{X}'$ 满、有限、平坦且局部有限表现。置
$R' = U' \times_{\mathcal{X}'} U'$，可知
$s', t' : R' \to U'$ 有限局部自由。由《群胚》中的引理
\ref{groupoids-lemma-find-invariant-affine}，存在包含 $u'$ 的
$R'$-不变仿射开子概形 $U'' \subset U'$。设
$\mathcal{X}'' \subset \mathcal{X}'$ 是相应的开子叠。那么
$\mathcal{X}''$ 近乎仿射。由引理
\ref{stacks-more-morphisms-lemma-well-nigh-affine-moduli-space-etale}，得到笛卡儿方形
$$
\xymatrix{
\mathcal{X}'' \ar[r] \ar[d] & \mathcal{X} \ar[d] \\
M'' \ar[r] & M
}
$$
其中 $M'' \to M$ 是 \'etale 的。由于
$\mathcal{X}'' \to M''$ 是仿射概形范畴中的范畴模空间，得到态射
$M'' \to V$，使复合 $\mathcal{X}'' \to \mathcal{X}' \to V$
等于复合 $\mathcal{X}'' \to M'' \to V$。这证明了断言并完成证明。
\end{proof}

\begin{lemma}
\label{stacks-more-morphisms-lemma-etale-separated-over-well-nigh-affine}
设 $h : \mathcal{X}' \to \mathcal{X}$ 是代数叠的态射。假设
$\mathcal{X}$ 近乎仿射，$h$ 是 \'etale 的，$h$ 分离，并且 $h$ 在自同构群
上诱导同构（《叠的态射》中的备注
\ref{stacks-morphisms-remark-identify-automorphism-groups}）。
那么存在笛卡儿图表
$$
\xymatrix{
\mathcal{X}' \ar[d] \ar[r] & \mathcal{X} \ar[d] \\
M' \ar[r] & M
}
$$
其中 $M' \to M$ 是概形之间分离的 \'etale 态射，并且
$\mathcal{X} \to M$ 是引理
\ref{stacks-more-morphisms-lemma-well-nigh-affine-moduli-space} 中构造的模空间。
\end{lemma}

\begin{proof}
选取仿射概形 $U$ 和一个满、平坦、有限且局部有限表现的态射
$U \to \mathcal{X}$。由于 $h$ 可由代数空间表示
（《叠的态射》中的诸引理
\ref{stacks-morphisms-lemma-aut-iso-unramified} 和
\ref{stacks-morphisms-lemma-stabilizer-preserving}），可知
$U' = \mathcal{X}' \times_\mathcal{X} U$ 是代数空间。由于
$U' \to U$ 分离且 \'etale，可知 $U'$ 是概形，并且 $U'$ 的任意有限
点集都包含在一个仿射开集中；参见《空间的态射》中的引理
\ref{spaces-morphisms-lemma-locally-quasi-finite-separated-representable}
和《态射进阶》中的引理
\ref{more-morphisms-lemma-separated-locally-quasi-finite-over-affine}。
置 $R' = U' \times_{\mathcal{X}'} U'$，可知
$s', t' : R' \to U'$ 有限局部自由。由《群胚》中的引理
\ref{groupoids-lemma-find-invariant-affine}，存在由 $R'$-不变仿射开
子概形 $U'_i \subset U'$ 组成的开覆盖
$U' = \bigcup U'_i$。设 $\mathcal{X}'_i \subset \mathcal{X}'$
是相应的开子叠。它们近乎仿射，因为
$U'_i \to \mathcal{X}'_i$ 满、平坦、有限且有限表现。由引理
\ref{stacks-more-morphisms-lemma-well-nigh-affine-moduli-space-etale}，得到笛卡儿图表
$$
\xymatrix{
\mathcal{X}'_i \ar[r] \ar[d] & \mathcal{X} \ar[d] \\
M'_i \ar[r] & M
}
$$
其中 $M'_i \to M$ 是仿射概形之间的 \'etale 态射，而竖直箭头如引理
\ref{stacks-more-morphisms-lemma-well-nigh-affine-moduli-space} 所述。注意
$\mathcal{X}'_{ij} = \mathcal{X}'_i \cap \mathcal{X}'_j$
% P12-ERR-0117: see ../control/ERRATA_CANDIDATES.jsonl
是 $\mathcal{X}'_i$ 和 $\mathcal{X}'_j$ 的开子空间。因此得到相应的
开子概形 $V_{ij} \subset M'_i$ 和 $V_{ji} \subset M'_j$。
由引理 \ref{stacks-more-morphisms-lemma-moduli-space-finite-affine} 的结论，可知
$\mathcal{X}'_{ij} \to V_{ij}$ 和
$\mathcal{X}'_{ji} \to V_{ji}$ 都是范畴模空间！因此得到唯一同构
$\varphi_{ij} : V_{ij} \to V_{ji}$，使得
$$
\xymatrix{
\mathcal{X}'_i \ar[d] & &
\mathcal{X}'_i \cap \mathcal{X}'_j \ar[rr] \ar[ll] \ar[ld] \ar[rd] & &
\mathcal{X}'_j \ar[d] \\
M'_i &
V_{ij} \ar[l] \ar[rr]^{\varphi_{ij}} & &
V_{ji} \ar[r] &
M'_j
}
$$
交换。这些同构满足《概形》第
% P12-ERR-0118: see ../control/ERRATA_CANDIDATES.jsonl
\ref{schemes-section-glueing-schemes} 节的余圈条件；这一点可由一个计算
（以及再次应用前一引理）得出，我们略去该计算。因此，可以把这些
% P12-ERR-0119: see ../control/ERRATA_CANDIDATES.jsonl
仿射概形黏合成概形 $M'$；参见《概形》中的引理
\ref{schemes-lemma-glue}。把 $M'_i$ 与它在 $M'$ 中的像等同起来。
我们断言存在态射 $\mathcal{X}' \to M'$，使它嵌入笛卡儿图表
$$
\xymatrix{
\mathcal{X}'_i \ar[r] \ar[d] & \mathcal{X}' \ar[d] \\
M'_i \ar[r] & M'
}
$$
这由《概形》中的引理 \ref{schemes-lemma-glue} 对到黏合概形 $M'$ 的
态射之描述以及下述事实显然可见：给出态射
% P12-ERR-0120: see ../control/ERRATA_CANDIDATES.jsonl
$\mathcal{X}' \to M'$ 等价于对任意态射 $T \to \mathcal{X}'$ 给定态射
$T \to M'$。类似地，存在态射 $M' \to M$，其在 $M'_i$ 上的限制是
给定态射 $M'_i \to M$。态射 $M' \to M$ 是 \'etale 的
（因为它在一个 \'etale 覆盖的成员上是 \'etale 的），而纤维积性质
成立，因为可以在（仿射）开覆盖 $M' = \bigcup M'_i$ 的成员上检验它。
最后，$M' \to M$ 分离，因为复合
$U' \to \mathcal{X}' \to M'$ 满且万有闭，所以可以应用《态射》中的
引理 \ref{morphisms-lemma-image-universally-closed-separated}。
\end{proof}

\begin{lemma}
\label{stacks-more-morphisms-lemma-etale-local-finite-inertia}
设 $\mathcal{X}$ 是代数叠。假设
$\mathcal{I}_\mathcal{X} \to \mathcal{X}$ 有限。那么存在集合 $I$，
并且对每个 $i \in I$，存在代数叠的态射
$$
g_i : \mathcal{X}_i \longrightarrow \mathcal{X}
$$
具有下列性质：
\begin{enumerate}
\item $|\mathcal{X}| = \bigcup |g_i|(|\mathcal{X}_i|)$；
\item $\mathcal{X}_i$ 近乎仿射；
\item $\mathcal{I}_{\mathcal{X}_i} \to
\mathcal{X}_i \times_\mathcal{X} \mathcal{I}_\mathcal{X}$
是同构；
\item $g_i : \mathcal{X}_i \to \mathcal{X}$ 可由代数空间表示、
分离且 \'etale。
\end{enumerate}
\end{lemma}

\begin{proof}
对任意 $x \in |\mathcal{X}|$，可以像《叠的态射》中的引理
\ref{stacks-morphisms-lemma-etale-local-quasi-DM-at-x} 那样选取
$g : \mathcal{U} \to \mathcal{X}$、$\mathcal{U} = [U/R]$ 和 $u$。
于是，由《叠的态射》中的引理
\ref{stacks-morphisms-lemma-stabilizer-preserving-unramified}，
存在包含 $u$ 的开子叠 $\mathcal{U}' \subset \mathcal{U}$，使得
$\mathcal{I}_{\mathcal{U}'} \to
\mathcal{U}' \times_\mathcal{X} \mathcal{I}_\mathcal{X}$
是同构。设 $U' \subset U$ 是与开子叠 $\mathcal{U}'$ 对应的
$R$-不变开集。设 $u' \in U'$ 是 $U'$ 中映到 $u$ 的一点。
注意，由于 $s : R \to U$ 有限，$t(s^{-1}(\{u'\}))$ 是有限集。
由《性质》中的引理 \ref{properties-lemma-ample-finite-set-in-affine}
和《群胚》中的引理 \ref{groupoids-lemma-find-invariant-affine}，
可以找到包含 $u'$ 的 $R$-不变仿射开集 $U'' \subset U'$。
设 $R''$ 是 $R$ 在 $U''$ 上的限制。那么
$\mathcal{U}'' = [U''/R'']$ 是 $\mathcal{U}'$ 中包含 $u$ 的开子叠，
它近乎仿射，并且
$\mathcal{I}_{\mathcal{U}''} \to
\mathcal{U}'' \times_\mathcal{X} \mathcal{I}_\mathcal{X}$
是同构，而 $\mathcal{U}'' \to \mathcal{X}$
% P12-ERR-0121: see ../control/ERRATA_CANDIDATES.jsonl
可由代数空间表示且 \'etale。最后，$\mathcal{U}'' \to \mathcal{X}$
分离，因为 $\mathcal{U}''$ 分离（引理
\ref{stacks-more-morphisms-lemma-well-nigh-affine}），
% P12-ERR-0122: see ../control/ERRATA_CANDIDATES.jsonl
$\mathcal{X}$ 的对角态射分离（《叠的态射》中的引理
\ref{stacks-morphisms-lemma-diagonal-diagonal}），并且分离性由
《叠的态射》中的引理
\ref{stacks-morphisms-lemma-compose-after-separated} 得出。
由于点 $x \in |\mathcal{X}|$ 是任意的，证明完成。
\end{proof}

\begin{theorem}[Keel-Mori]
\label{stacks-more-morphisms-theorem-keel-mori}
设 $\mathcal{X}$ 是代数叠。假设
$\mathcal{I}_\mathcal{X} \to \mathcal{X}$ 有限。那么存在一致范畴模空间
$$
f : \mathcal{X} \longrightarrow M
$$
并且 $f$ 分离、拟紧且为万有同胚。
\end{theorem}

\begin{proof}
选取集合 $I$\footnote{仍在留意集合论问题的读者应当确认 $I$
不会过大。}，并对每个 $i \in I$ 选取如引理
\ref{stacks-more-morphisms-lemma-etale-local-finite-inertia} 中的代数叠态射
$g_i : \mathcal{X}_i \to \mathcal{X}$；我们将不再特别说明地使用该引理
所列的全部性质。令
$$
f_i : \mathcal{X}_i \to M_i
$$
如引理 \ref{stacks-more-morphisms-lemma-well-nigh-affine-moduli-space} 所述。对
$i, j \in I$，考虑叠
$$
\mathcal{X}_{ij} = \mathcal{X}_i \times_{g_i, \mathcal{X}, g_j} \mathcal{X}_j
$$
由《叠的态射》中的引理
\ref{stacks-morphisms-lemma-base-change-separated}，投影
$\mathcal{X}_{ij} \to \mathcal{X}_i$ 和
$\mathcal{X}_{ij} \to \mathcal{X}_j$ 分离；由《叠的态射》中的引理
\ref{stacks-morphisms-lemma-base-change-etale}，它们是 \'etale 的；
它们还在自同构群上诱导同构（如《叠的态射》中的备注
\ref{stacks-morphisms-remark-identify-automorphism-groups}），这一点由
《叠的态射》中的引理
\ref{stacks-morphisms-lemma-base-change-stabilizer-preserving} 得出。因此，可以应用
引理 \ref{stacks-more-morphisms-lemma-etale-separated-over-well-nigh-affine} 得到交换图表
$$
\xymatrix{
\mathcal{X}_i \ar[d]_{f_i} &
\mathcal{X}_{ij} \ar[d]_{f_{ij}} \ar[l] \ar[r] &
\mathcal{X}_j \ar[d]_{f_j} \\
M_i &
M_{ij} \ar[l] \ar[r] &
M_j
}
$$
其中两个方形都是笛卡儿方形，而 $M_{ij} \to M_i$ 和
$M_{ij} \to M_j$ 是概形之间分离的 \'etale 态射；这里还由引理
\ref{stacks-more-morphisms-lemma-moduli-space-finite-affine} 使用了
% P12-ERR-0123: see ../control/ERRATA_CANDIDATES.jsonl
$f_i$ 是一致范畴商这一点。断言：
$$
\coprod M_{ij} \longrightarrow \coprod M_i \times \coprod M_i
$$
是 \'etale 等价关系。

\medskip\noindent
证明该断言。置 $R = \coprod M_{ij}$ 和 $U = \coprod M_i$。
已经看到 $t : R \to U$ 和 $s : R \to U$ 是 \'etale 的。
我们构造与
$\text{pr}_{13} : U \times U \times U \to U \times U$ 相容的态射
$c : R \times_{s, U, t} R \to R$。具体而言，对
$i, j, k \in I$，考虑
$$
\mathcal{X}_{ijk} =
\mathcal{X}_i \times_{g_i, \mathcal{X}, g_j} \mathcal{X}_j
\times_{g_j, \mathcal{X}, g_k} \mathcal{X}_k =
\mathcal{X}_{ij} \times_{\mathcal{X}_j} \mathcal{X}_{jk}
$$
完全沿用上一段的论证，可知
$M_{ijk} = M_{ij} \times_{M_j} M_{jk}$ 是
$\mathcal{X}_{ijk}$ 的范畴模空间。特别地，由投影
$\mathcal{X}_{ijk} \to \mathcal{X}_{ik}$ 得到典范态射
$M_{ijk} = M_{ij} \times_{M_j} M_{jk} \to M_{ik}$。
把这些态射合在一起，便得到态射 $c$。类似地，构造与
$\Delta : U \to U \times U$ 相容的态射 $e : U \to R$，以及与对换
$U \times U \to U \times U$ 相容的态射 $i : R \to R$。
设 $k$ 是代数闭域。那么
$$
\Mor(\Spec(k), \mathcal{X}_i) \to \Mor(\Spec(k), M_i) = M_i(k)
$$
在同构类上是双射，并且沿任意态射
% P12-ERR-0124: see ../control/ERRATA_CANDIDATES.jsonl
$M' \to M$ 作基变换后仍然如此。这由 $f_i$ 的选取以及《叠的态射》
中的诸引理 \ref{stacks-morphisms-lemma-universally-injective} 和
\ref{stacks-morphisms-lemma-universally-injective-point} 得出。
由 $2$-纤维积的构造，图表
$$
\xymatrix{
\Mor(\Spec(k), \mathcal{X}_{ij}) \ar[d] \ar[r] &
\Mor(\Spec(k), \mathcal{X}_j) \ar[d] \\
\Mor(\Spec(k), \mathcal{X}_i) \ar[r] &
\Mor(\Spec(k), \mathcal{X})
}
$$
是范畴的纤维积。由 $g_i$ 的选取，该图表中的函子在自同构群上诱导
双射。因此，该图表在同构类集合上诱导纤维积图表！于是
$$
R(k) = U(k) \times_{|\Mor(\Spec(k), \mathcal{X})|} U(k)
$$
其中 $|\Mor(\Spec(k), \mathcal{X})|$ 表示同构类集合。特别地，
对任意代数闭域 $k$，
% P12-ERR-0125: see ../control/ERRATA_CANDIDATES.jsonl
在 $k$-值点上的映射是等价关系。由《群胚》中的引理
\ref{groupoids-lemma-etale-equivalence-relation}，断言得证。

\medskip\noindent
设 $M = U/R$ 为上述 \'etale 等价关系的商代数空间；参见《空间》中的
定理 \ref{spaces-theorem-presentation}。存在典范态射
$f : \mathcal{X} \to M$，它嵌入交换图表
\begin{equation}
\label{stacks-more-morphisms-equation-fundamental-diagram}
\xymatrix{
\mathcal{X}_i \ar[r]_{g_i} \ar[d]_{f_i} & \mathcal{X} \ar[d]^f \\
M_i \ar[r] & M
}
\end{equation}
具体而言，这样的态射 $f$ 由对任意概形 $T$ 都与基变换相容的函子
$$
f : \Mor(T, \mathcal{X}) \longrightarrow \Mor(T, M)
$$
给出。设 $a : T \to \mathcal{X}$ 是左端的一个对象。得到 \'etale 覆盖
$\{T_i \to T\}$，其中
$T_i = \mathcal{X}_i \times_\mathcal{X} T$，以及态射
$a_i : T_i \to \mathcal{X}_i$。于是得到
$b_i = f_i \circ a_i : T_i \to M_i$。由于
$T_i \times_T T_j = \mathcal{X}_{ij} \times_\mathcal{X} T$，
还得到态射 $a_{ij} : T_i \times_T T_j \to \mathcal{X}_{ij}$。
置 $b_{ij} = f_{ij} \circ a_{ij}$，可知 $b_i \times b_j$
经由单态射 $M_{ij} \to M_i \times M_j$ 分解。因此，诸态射
$$
T_i \xrightarrow{b_i} M_i \to M
$$
在 $T_i \times_T T_j$ 上相同。由于 $M$ 是 \'etale 拓扑的层，
这些态射黏合成唯一态射 $b = f(a) : T \to M$。我们略去验证该构造
与基变换相容，也略去验证图表
(\ref{stacks-more-morphisms-equation-fundamental-diagram}) 交换。

\medskip\noindent
断言：图表 (\ref{stacks-more-morphisms-equation-fundamental-diagram}) 是笛卡儿的。
为看出这一点，考察诱导态射
$$
h_i : \mathcal{X}_i \longrightarrow M_i \times_M \mathcal{X}
$$
这是 $\mathcal{X}$ 上叠的 \'etale 态射，所以 $h_i$ 是 \'etale 的
（《叠的态射》中的引理
\ref{stacks-morphisms-lemma-etale-permanence}）。由于 $g_i$ 分离，
可知 $h_i$ 分离（使用《叠的态射》中的引理
\ref{stacks-morphisms-lemma-compose-after-separated} 以及以上看到的
$\mathcal{X}$ 的对角态射分离这一事实）。态射 $h_i$ 在自同构群上诱导
同构（《叠的态射》中的备注
\ref{stacks-morphisms-remark-identify-automorphism-groups}），
因为 $g_i$ 具有这一性质。对代数闭域 $k$，图表
$$
\xymatrix{
\Mor(\Spec(k), M_i \times_M \mathcal{X}) \ar[r] \ar[d] &
\Mor(\Spec(k), \mathcal{X}) \ar[d] \\
M_i(k) \ar[r] &
M(k)
}
$$
% P12-ERR-0126: see ../control/ERRATA_CANDIDATES.jsonl
是范畴的笛卡儿图表，并且上方箭头在自同构群上诱导双射。另一方面，
由以上所述，有
$$
M(k) = U(k)/R(k) = U(k)/
U(k) \times_{|\Mor(\Spec(k), \mathcal{X})|} U(k) =
|\Mor(\Spec(k), \mathcal{X})|
$$
因此，上方笛卡儿图表的右竖直箭头在同构类上是双射。可知
$|\Mor(\Spec(k), M_i \times_M \mathcal{X})| \to M_i(k)$ 是双射。
总结：
% P12-ERR-0127: see ../control/ERRATA_CANDIDATES.jsonl
$h_i$ 分离、\'etale，在自同构群上诱导同构（如《叠的态射》中的备注
\ref{stacks-morphisms-remark-identify-automorphism-groups}），并且在代数闭域
上的纤维范畴上诱导等价。因此，由《叠的态射》中的引理
\ref{stacks-morphisms-lemma-etale-iso}，它是同构。

\medskip\noindent
特别地，由该断言得到：存在满 \'etale 态射 $U \to M$，使得 $f$ 的
基变换分离、拟紧且为万有同胚。由此可知 $f$ 分离、拟紧且为万有同胚。
参见《叠的态射》中的诸引理
\ref{stacks-morphisms-lemma-check-separated-covering}、
\ref{stacks-morphisms-lemma-check-quasi-compact-covering} 和
% P12-ERR-0128: see ../control/ERRATA_CANDIDATES.jsonl
\ref{stacks-morphisms-lemma-check-universal-homeomorphism-covering}

\medskip\noindent
为完成证明，必须证明 $f : \mathcal{X} \to M$ 是一致范畴模空间。
为此，只须证明：给定代数空间的平坦态射 $M' \to M$，基变换
$$
M' \times_M \mathcal{X} \longrightarrow M'
$$
是范畴模空间。为此考虑态射
$$
\theta : M' \times_M \mathcal{X} \longrightarrow E
$$
其中 $E$ 是代数空间。对每个 $i$，已知 $f_i$ 是一致范畴模空间。
因此得到
$$
\xymatrix{
M' \times_M \mathcal{X}_i \ar[d] \ar[r] &
M' \times_M \mathcal{X} \ar[d]^\theta \\
M' \times_M M_i \ar[r]^{\psi_i} &
E
}
$$
由于 $\{M' \times_M M_i \to M'\}$ 是 \'etale 覆盖，为得到所需态射
$\psi : M' \to E$，只须证明 $\psi_i$ 和 $\psi_j$ 在
$M' \times_M M_i \times_M M_j = M' \times_M M_{ij}$ 上相同。
这容易由下述事实得出：
% P12-ERR-0129: see ../control/ERRATA_CANDIDATES.jsonl
$f_{ij} : \mathcal{X}_{ij} =
\mathcal{X}_i \times_\mathcal{X} \mathcal{X}_j \to M_{ij}$ 是一致范畴商；
略去细节。最后，再证明 $\psi$ 嵌入交换图表
$$
\xymatrix{
M' \times_M \mathcal{X} \ar[d] \ar[rd]^\theta \\
M' \ar[r]^\psi &
E
}
$$
因为“$\{M' \times_M \mathcal{X}_i \to M' \times_M \mathcal{X}\}$
是一个 \'etale 覆盖”，而诸态射 $\psi_i$ 按照构造嵌入相应的交换图表。
这完成了 Keel--Mori 定理的证明。
\end{proof}

\noindent
下一个引理强调该构造的 \'etale 局部性质。

\begin{lemma}
\label{stacks-more-morphisms-lemma-etale-separated-over-keel-mori}
设 $h : \mathcal{X}' \to \mathcal{X}$ 是代数叠的态射。假设
\begin{enumerate}
\item $\mathcal{I}_\mathcal{X} \to \mathcal{X}$ 有限；
\item $h$ 是 \'etale 且分离的，并且在自同构群上诱导同构
（《叠的态射》中的备注
\ref{stacks-morphisms-remark-identify-automorphism-groups}）。
\end{enumerate}
那么存在笛卡儿图表
$$
\xymatrix{
\mathcal{X}' \ar[d] \ar[r] &
\mathcal{X} \ar[d] \\
M' \ar[r] &
M
}
$$
其中 $M' \to M$ 是代数空间之间分离的 \'etale 态射，而竖直箭头是
定理 \ref{stacks-more-morphisms-theorem-keel-mori} 中构造的模空间。
\end{lemma}

\begin{proof}
由《叠的态射》中的引理
\ref{stacks-morphisms-lemma-aut-iso-unramified}，可知
$\mathcal{I}_{\mathcal{X}'} \to
\mathcal{X}' \times_\mathcal{X} \mathcal{I}_\mathcal{X}$
是同构。因此，$\mathcal{I}_{\mathcal{X}'} \to \mathcal{X}'$
作为 $\mathcal{I}_\mathcal{X} \to \mathcal{X}$ 的基变换而有限。
设 $f' : \mathcal{X}' \to M'$ 和 $f : \mathcal{X} \to M$ 如定理
\ref{stacks-more-morphisms-theorem-keel-mori} 所述。因为
% P12-ERR-0130: see ../control/ERRATA_CANDIDATES.jsonl
$f'$ 是范畴模空间，得到本引理中的交换图表。像引理
\ref{stacks-more-morphisms-lemma-etale-local-finite-inertia} 中那样选取 $I$ 和
$g'_i : \mathcal{X}'_i \to \mathcal{X}'$。注意，
$g_i = h \circ g'_i$ 是 \'etale 且分离的，并在自同构群上诱导同构
（《叠的态射》中的备注
\ref{stacks-morphisms-remark-identify-automorphism-groups}）。设
$f'_i : \mathcal{X}'_i \to M'_i$ 如引理
\ref{stacks-more-morphisms-lemma-well-nigh-affine-moduli-space} 所述。在定理
\ref{stacks-more-morphisms-theorem-keel-mori} 的证明中已经看到，图表
$$
\xymatrix{
\mathcal{X}'_i \ar[d]_{f'_i} \ar[r]_{g'_i} &
\mathcal{X}' \ar[d]^{f'} \\
M'_i \ar[r] &
M'
}
\quad\text{以及}\quad
\xymatrix{
\mathcal{X}'_i \ar[d]_{f'_i} \ar[r]_{g_i} &
\mathcal{X} \ar[d]^f \\
M'_i \ar[r] &
M
}
$$
都是笛卡儿的，并且 $M'_i \to M'$ 和 $M'_i \to M$ 是 \'etale 的
（完全以同样方式论证，这也直接由
% P12-ERR-0131: see ../control/ERRATA_CANDIDATES.jsonl
迄今所列态射 $g'_i, g_i, f', f'_i, f$ 的性质得出）。这首先蕴含
$M' \to M$ 是 \'etale 的，其次蕴含本引理中的图表是笛卡儿的。
仍须证明 $M' \to M$ 分离。为此考察图表
$$
\xymatrix{
\mathcal{X}' \ar[r] \ar[d] &
\mathcal{X}' \times_\mathcal{X} \mathcal{X}' \ar[d] \\
M' \ar[r] &
M' \times_M M'
}
$$
由于 $\mathcal{X}' \to \mathcal{X}$ 分离，上方水平箭头万有闭。
竖直箭头如定理 \ref{stacks-more-morphisms-theorem-keel-mori} 中所述
（是 $\mathcal{X} \to M$ 的平坦基变换），因而是万有同胚。
所以，下方水平箭头万有闭。这一点与它是代数空间的 \'etale 单态射
结合，证明它如所需为闭浸入。
\end{proof}

\section{模空间的性质}
\label{stacks-more-morphisms-section-properties-moduli-spaces}

\noindent
一旦证明模空间存在，就可以（而且通常很直接地）确定这些模空间的性质。

\begin{lemma}
\label{stacks-more-morphisms-lemma-keel-mori-finite-type}
设 $p : \mathcal{X} \to Y$ 是从代数叠到代数空间的态射。假设
\begin{enumerate}
\item $\mathcal{I}_\mathcal{X} \to \mathcal{X}$ 有限；
\item $Y$ 局部 Noether；
\item $p$ 局部有限型。
\end{enumerate}
设 $f : \mathcal{X} \to M$ 是定理
\ref{stacks-more-morphisms-theorem-keel-mori} 中构造的模空间。那么 $M \to Y$ 局部有限型。
\end{lemma}

\begin{proof}
由于 $f$ 是一致范畴模空间，得到态射 $M \to Y$。只须在 $M$ 和 $Y$
上 \'etale 局部地检验 $M \to Y$ 局部有限型。由于 $f$ 是一致范畴
模空间，可以先用一个在 $Y$ 上 \'etale 的仿射概形代替 $Y$。
接着，可以像引理 \ref{stacks-more-morphisms-lemma-etale-local-finite-inertia} 中那样选取
$I$ 和 $g_i : \mathcal{X}_i \to \mathcal{X}$。然后由引理
\ref{stacks-more-morphisms-lemma-etale-separated-over-keel-mori}，问题约化为
$\mathcal{X} = \mathcal{X}_i$ 的情形。换言之，可以假设
$\mathcal{X}$ 近乎仿射。在这种情形下，
$Y = \Spec(A_0)$，$\mathcal{X} = [U/R]$ 且
$U = \Spec(A)$，并且 $M = \Spec(C)$，其中 $C \subset A$ 是
$U$ 上所有 $R$-不变函数组成的集合。参见引理
\ref{stacks-more-morphisms-lemma-well-nigh-affine} 和
\ref{stacks-more-morphisms-lemma-well-nigh-affine-moduli-space}。于是 $A_0$ 是 Noether 环，
而 $A_0 \to A$ 有限型。此外，由《群胚》中的引理
\ref{groupoids-lemma-integral-over-invariants}，$A$ 在 $C$ 上整，
因而在 $C$ 上有限（因为它在 $A_0$ 上有限型）。最后应用《代数》中的
引理 \ref{algebra-lemma-Artin-Tate} 即得结论。
\end{proof}

\begin{lemma}
\label{stacks-more-morphisms-lemma-keel-mori-diagonal}
设 $\mathcal{X}$ 是代数叠。假设
$\mathcal{I}_\mathcal{X} \to \mathcal{X}$ 有限。设
$f : \mathcal{X} \to M$ 是定理
\ref{stacks-more-morphisms-theorem-keel-mori} 中构造的模空间。
\begin{enumerate}
\item 如果 $\mathcal{X}$ 拟分离，那么 $M$ 拟分离。
\item 如果 $\mathcal{X}$ 分离，那么 $M$ 分离。
% P12-ERR-0132: see ../control/ERRATA_CANDIDATES.jsonl
\item 此处还应添加更多内容，例如上述结论的相对版本。
\end{enumerate}
\end{lemma}

\begin{proof}
为证明这些断言，考虑图表
$$
\xymatrix{
\mathcal{X} \ar[d]_f \ar[r]_{\Delta_\mathcal{X}} &
\mathcal{X} \times \mathcal{X} \ar[d]^{f \times f} \\
M \ar[r]^{\Delta_M} &
M \times M
}
$$
由于 $f$ 是万有同胚，可知 $f \times f$ 是万有同胚。

\medskip\noindent
如果 $\mathcal{X}$ 分离，那么 $\Delta_\mathcal{X}$ 固有，故
$\Delta_\mathcal{X}$ 万有闭，进而 $\Delta_M$ 万有闭，所以由
《空间的态射》中的引理
\ref{spaces-morphisms-lemma-separated-diagonal-proper}，$M$ 分离。

\medskip\noindent
如果 $\mathcal{X}$ 拟分离，那么 $\Delta_\mathcal{X}$ 拟紧，
故 $\Delta_M$ 拟紧，所以 $M$ 拟分离。
\end{proof}

\begin{lemma}
\label{stacks-more-morphisms-lemma-keel-mori-proper}
设 $p : \mathcal{X} \to Y$ 是从代数叠到代数空间的态射。假设
\begin{enumerate}
\item $\mathcal{I}_\mathcal{X} \to \mathcal{X}$ 有限；
\item $p$ 固有；
\item $Y$ 局部 Noether。
\end{enumerate}
设 $f : \mathcal{X} \to M$ 是定理
\ref{stacks-more-morphisms-theorem-keel-mori} 中构造的模空间。那么 $M \to Y$ 固有。
\end{lemma}

\begin{proof}
由引理 \ref{stacks-more-morphisms-lemma-keel-mori-finite-type}，可知 $M \to Y$ 局部有限型。
由引理 \ref{stacks-more-morphisms-lemma-keel-mori-diagonal}，可知 $M \to Y$ 分离。
当然，$M \to Y$ 拟紧且万有闭，因为这些是拓扑性质，
$\mathcal{X} \to Y$ 具有这些性质，而
$\mathcal{X} \to M$ 是万有同胚。
\end{proof}

\section{叠与 fpqc 覆盖}
\label{stacks-more-morphisms-section-fpqc}

\noindent
某些代数叠满足 fpqc 下降。本节对代数空间的对应结论见《空间的性质》
第 \ref{spaces-properties-section-fpqc} 节。


\begin{proposition}
\label{stacks-more-morphisms-proposition-stack-fpqc}
\begin{reference}
Anatoly Preygel《代数叠导论》命题 3.3.6。
\end{reference}
设 $\mathcal{X}$ 是对角态射拟仿射\footnote{假设对角态射为
ind-拟仿射便已足够。}的代数叠。那么
$\mathcal{X}$ 对 fpqc 覆盖满足下降。
\end{proposition}

\begin{proof}
我们的约定是：$\mathcal{X}$ 是基概形 $S$ 上概形范畴（赋予 fppf 拓扑）
之上的群胚叠 $p : \mathcal{X} \to (\Sch/S)_{fppf}$。该陈述的含义如下：
给定 $S$ 上概形的 fpqc 覆盖
$\mathcal{U} = \{U_i \to U\}_{i \in I}$，函子
$$
\mathcal{X}_U \longrightarrow DD(\mathcal{U})
$$
是等价。左边是 $\mathcal{X}$ 在 $U$ 上的对象范畴，右边是
$\mathcal{X}$ 中相对于 $\mathcal{U}$ 的下降数据范畴。参见《叠》第
\ref{stacks-section-descent-data} 节的讨论。

\medskip\noindent
% P12-ERR-0133: see ../control/ERRATA_CANDIDATES.jsonl
全忠实性。假设有 $\mathcal{X}$ 在 $U$ 上的两个对象 $x, y$。
那么 $I = \mathit{Isom}(x, y)$ 是 $U$ 上的代数空间。因此，$I$ 在
$U_i$ 上的一族截面，若其限制在 $U_i \times_U U_j$ 上相同，
则由代数空间的对应命题来自 $U$ 上的唯一截面；参见《空间的性质》中的
命题 \ref{spaces-properties-proposition-sheaf-fpqc}。所以该函子全忠实。

\medskip\noindent
本质满射性。这里给定 $U_i$ 上的对象 $x_i$ 和
$U_i \times_U U_j$ 上的同构
$\varphi_{ij} : \text{pr}_0^*x_i \to \text{pr}_1^*x_j$，它们在
% P12-ERR-0134: see ../control/ERRATA_CANDIDATES.jsonl
$U_i \times_U U_j \times_U U_k$ 上满足余圈条件。

\medskip\noindent
设 $W$ 是仿射概形，并设 $W \to \mathcal{X}$ 是态射。对每个 $i$，
可以形成
$$
W_i = U_i \times_{x_i, \mathcal{X}} W
$$
投影 $W_i \to U_i$ 拟仿射，因为 $\mathcal{X}$ 的对角态射拟仿射。
对每对 $i, j \in I$，同构 $\varphi_{ij}$ 诱导同构
$$
W_i \times_U U_j =
(U_i \times_U U_j) \times_{x_i \circ \text{pr}_0, \mathcal{X}} W
\to
(U_i \times_U U_j) \times_{x_j \circ \text{pr}_1, \mathcal{X}} W =
U_i \times_U W_j
$$
而且，这些同构在 $U_i \times_U U_j \times_U U_k$ 上满足余圈条件。
换言之，这些同构定义了概形 $W_i/U_i$ 相对于 $\mathcal{U}$ 的下降数据。
由《下降》中的引理 \ref{descent-lemma-quasi-affine}，该下降数据有效
\footnote{在对角态射为 ind-拟仿射的情形，也可使用《群胚进阶》中的
引理 \ref{more-groupoids-lemma-ind-quasi-affine}。}。因此，存在拟仿射态射
$W' \to U$ 和交换图表
$$
\xymatrix{
W' \ar[d] &
W_i \ar[l] \ar[d] \ar[r] &
W \ar[d] \\
U &
U_i \ar[l] \ar[r]^{x_i} &
\mathcal{X}
}
$$
其中两个方形都是笛卡儿的。由于 $\{W_i \to W'\}_{i \in I}$
是 $\mathcal{U}$ 沿 $W' \to U$ 的基变换，它是 fpqc 覆盖。
由于 $W$ 对 fpqc 覆盖满足层条件，得到唯一态射 $W' \to W$，
使 $W_i \to W' \to W$ 为给定态射 $W_i \to W$。换言之，
有交换图表
$$
\xymatrix{
W_i \ar[d] \ar[r] &
W' \ar[d] \ar[r] &
W \ar[d] \\
U_i \ar[r] \ar@/_1pc/[rr]_{x_i} &
U &
\mathcal{X}
}
$$
它们与同构 $\varphi_{ij}$ 相容，并且其中的方形和矩形都是笛卡儿的。

\medskip\noindent
选取一族仿射概形 $W_\alpha$（$\alpha \in A$）以及光滑态射
$W_\alpha \to \mathcal{X}$，使得
$\coprod W_\alpha \to \mathcal{X}$ 满。通过上一段的步骤，
对每个 $\alpha$ 产生图表
$$
\xymatrix{
W_{\alpha, i} \ar[d] \ar[r] &
W_\alpha' \ar[d] \ar[r] &
W_\alpha \ar[d] \\
U_i \ar[r] \ar@/_1pc/[rr]_{x_i} &
U &
\mathcal{X}
}
$$
于是诸态射 $W_\alpha' \to U$ 光滑且联合满射。

\medskip\noindent
% P12-ERR-0135: see ../control/ERRATA_CANDIDATES.jsonl
用 $x_\alpha$ 表示 $\mathcal{X}$ 在 $W_\alpha'$ 上与
$W_\alpha' \to W_\alpha \to \mathcal{X}$ 对应的对象。由于
$\mathcal{X}$ 是 fppf 叠，并且
$\{W_\alpha' \to U\}$ 是 fppf 覆盖，只须证明在
$W_\alpha' \times_U W'_\beta$ 上存在满足余圈条件的同构
$\text{pr}_0^*x_\alpha \to \text{pr}_1^*x_\beta$。然而，拉回到
$W_{\alpha, i}$ 后，确实在
$W_{\alpha, i} \times_{U_i} W_{\beta, i} =
U_i \times_U (W_\alpha' \times_U W'_\beta)$ 上有这样的同构，
因为 $x_\alpha$ 到 $W_{\alpha, i}$ 的拉回同构于 $x_i$ 到
$W_{\alpha, i}$ 的拉回。由于
$\{U_i \times_U (W_\alpha' \times_U W'_\beta) \to
W_\alpha' \times_U W'_\beta\}_{i \in I}$ 是 fpqc 覆盖，又由于
上述图表与 $\varphi_{ij}$ 的相容性，这些同构下降到
$W_\alpha' \times_U W'_\beta$，证明完成。
\end{proof}

\section{张量函子}
\label{stacks-more-morphisms-section-tensor-functors}

\noindent
设 $f : \mathcal{Y} \to \mathcal{X}$ 是 Noether 代数叠的态射。拉回函子
$$
f^* :
\textit{Coh}(\mathcal{O}_\mathcal{X})
\longrightarrow
\textit{Coh}(\mathcal{O}_\mathcal{Y})
$$
是右正合张量函子：它可加、右正合，并与凝聚模的张量积交换。
可以问，任意右正合张量函子
$F : \textit{Coh}(\mathcal{O}_\mathcal{X}) \to
\textit{Coh}(\mathcal{O}_\mathcal{Y})$ 在多大程度上来自某个态射
$f : \mathcal{Y} \to \mathcal{X}$。关于这种类型的一个非常一般的结论，
读者可查阅 \cite{Hall-Rydh-coherent}。本节的目标是简短证明定理
\ref{stacks-more-morphisms-theorem-main}，作为这些思想的入门。

\medskip\noindent
先给出几个引理。

\begin{lemma}
\label{stacks-more-morphisms-lemma-extend}
设 $\mathcal{X}$ 和 $\mathcal{Y}$ 是 Noether 代数叠。任意右正合张量函子
$F : \textit{Coh}(\mathcal{O}_\mathcal{X}) \to
\textit{Coh}(\mathcal{O}_\mathcal{Y})$ 都唯一延拓为与所有余极限交换的
右正合张量函子
$F : \QCoh(\mathcal{O}_\mathcal{X}) \to \QCoh(\mathcal{O}_\mathcal{Y})$。
\end{lemma}

\begin{proof}
这种延拓的存在性和唯一性是一般事实；参见《范畴论》中的引理
\ref{categories-lemma-extend-functor-by-colim}。为看出该引理适用，
注意局部 Noether 代数叠上的凝聚模按定义就是有限表现模；参见
《叠的上同调》中的定义
\ref{stacks-cohomology-definition-coherent}。因此，由《叠的上同调》中的
引理 \ref{stacks-cohomology-lemma-finite-presentation-quasi-compact-colimit}，
$\mathcal{X}$ 上的凝聚模是
$\QCoh(\mathcal{O}_\mathcal{X})$ 的范畴紧对象。最后，由《叠的上同调》
中的引理 \ref{stacks-cohomology-lemma-directed-colimit-coherent}，
每个拟凝聚模都是其凝聚子模的滤过余极限。

\medskip\noindent
由于 $F$ 可加，$F$ 的延拓也可加（略去细节）。由于 $F$ 是张量函子，
而模的余极限与取张量积交换，$F$ 的延拓也是张量函子（略去细节）。

\medskip\noindent
本段证明该延拓与任意直和交换。如果
$\mathcal{F} = \bigoplus_{j \in J} \mathcal{H}_j$，其中
$\mathcal{H}_j$ 拟凝聚。

于是
$\mathcal{F} = \colim_{J' \subset J\text{，且为有限集}}
\bigoplus_{j \in J'} \mathcal{H}_j$。仍用 $F$ 表示 $F$ 的延拓，得到
\begin{align*}
F(\mathcal{F})
& =
\colim_{J' \subset J\text{，且为有限集}}
F(\bigoplus\nolimits_{j \in J'} \mathcal{H}_j) \\
& =
\colim_{J' \subset J\text{，且为有限集}}
\bigoplus\nolimits_{j \in J'} F(\mathcal{H}_j) \\
& =
\bigoplus\nolimits_{j \in J} F(\mathcal{H}_j)
\end{align*}
因此 $F$ 与任意直和交换。

\medskip\noindent
本段证明该延拓右正合。假设
$0 \to \mathcal{F} \to \mathcal{F}' \to \mathcal{F}'' \to 0$
是拟凝聚 $\mathcal{O}_\mathcal{X}$-模的短正合列。把
$\mathcal{F}' = \bigcup \mathcal{F}'_i$ 写成其凝聚子模之并
（见以上所给参考文献）。
% P12-ERR-0136: see ../control/ERRATA_CANDIDATES.jsonl
用 $\mathcal{F}''_i \subset \mathcal{F}''$ 表示 $\mathcal{F}'_i$ 的像，
并置 $\mathcal{F}_i = \mathcal{F} \cap \mathcal{F}'_i =
\Ker(\mathcal{F}'_i \to \mathcal{F}''_i)$。于是显然有
$\mathcal{F} = \bigcup \mathcal{F}_i$ 和
$\mathcal{F}'' = \bigcup \mathcal{F}''_i$，并有短正合列
$$
0 \to \mathcal{F}_i \to \mathcal{F}_i' \to \mathcal{F}_i'' \to 0
$$
由于该延拓与滤过余极限交换，有
$F(\mathcal{F}) = \colim_{i \in I} F(\mathcal{F}_i)$、
$F(\mathcal{F}') = \colim_{i \in I} F(\mathcal{F}'_i)$ 和
$F(\mathcal{F}'') = \colim_{i \in I} F(\mathcal{F}''_i)$。
% P12-ERR-0137: see ../control/ERRATA_CANDIDATES.jsonl
由于模层的滤过余极限正合，可知 $F$ 的延拓右正合。

\medskip\noindent
证明完成，因为与所有余积交换的右正合函子与所有余极限交换；参见
《范畴论》中的引理
\ref{categories-lemma-colimits-coproducts-coequalizers}。
\end{proof}

\begin{lemma}
\label{stacks-more-morphisms-lemma-affine}
设 $\mathcal{X}$ 是对角态射仿射的代数叠。设 $B$ 是环。设
$F : \QCoh(\mathcal{O}_\mathcal{X}) \to \text{Mod}_B$
是与直和交换的右正合张量函子。设
$g : U \to \mathcal{X}$ 是态射，其中 $U = \Spec(A)$ 仿射。那么
\begin{enumerate}
\item $C = F(g_{\QCoh, *}\mathcal{O}_U)$ 是交换 $B$-代数；
\item 存在环映射 $A \to C$；
\end{enumerate}
使得 $F \circ g_{\QCoh, *} : \text{Mod}_A \to \text{Mod}_B$
把 $M$ 映为视作 $B$-模的 $M \otimes_A C$。
\end{lemma}

\begin{proof}
注意 $g$ 拟紧且拟分离；参见《叠的态射》中的引理
\ref{stacks-morphisms-lemma-quasi-compact-quasi-separated-permanence}。
在《叠的上同调》中的命题
\ref{stacks-cohomology-proposition-direct-image-quasi-coherent} 中，
已经构造函子
$g_{\QCoh, *} : \QCoh(\mathcal{O}_U) \to
\QCoh(\mathcal{O}_\mathcal{X})$。由《叠的上同调》中的诸备注
\ref{stacks-cohomology-remark-direct-image-quasi-coherent-tensor} 和
\ref{stacks-cohomology-remark-QCoh-tensor}，得到乘法
$$
\mu :
g_{\QCoh, *}\mathcal{O}_U
\otimes_{\mathcal{O}_\mathcal{X}}
g_{\QCoh, *}\mathcal{O}_U
\longrightarrow
g_{\QCoh, *}\mathcal{O}_U
$$
它使 $g_{\QCoh, *}\mathcal{O}_U$ 成为交换
$\mathcal{O}_\mathcal{X}$-代数。因此
$C = F(g_{\QCoh, *}\mathcal{O}_U)$ 是 $\text{Mod}_B$ 中的交换代数对象；
换言之，$C$ 是交换 $B$-代数。注意，有映射
$\kappa : A \to
\text{End}_{\mathcal{O}_\mathcal{X}}(g_{\QCoh, *}\mathcal{O}_U)$，
使得对任意 $a \in A$，图表
% P12-ERR-0138: see ../control/ERRATA_CANDIDATES.jsonl
$$
\xymatrix{
g_{\QCoh, *}\mathcal{O}_U \otimes_{\mathcal{O}_\mathcal{X}}
g_{\QCoh, *}\mathcal{O}_U
\ar[d]_{\kappa(r) \otimes 1} \ar[rr]_-\mu & &
g_{\QCoh, *}\mathcal{O}_U \ar[d]^{\kappa(r)} \\
g_{\QCoh, *}\mathcal{O}_U \otimes_{\mathcal{O}_\mathcal{X}}
g_{\QCoh, *}\mathcal{O}_U
\ar[rr]^-\mu & &
g_{\QCoh, *}\mathcal{O}_U
}
$$
交换。由此得到映射
$\kappa' = F(\kappa) : A \to \text{End}_B(C)$，满足
$\kappa'(a)(c) c' = \kappa'(a)(cc')$；当然，这意味着
$a \mapsto \kappa'(a)(1)$ 是环映射 $A \to C$。

\medskip\noindent
态射 $g : U \to \mathcal{X}$ 仿射；参见《叠的态射》中的引理
\ref{stacks-morphisms-lemma-affine-permanence}。因此，由《叠的上同调》中的
引理 \ref{stacks-cohomology-lemma-quasi-coherent-pushforward-affine}，
$g_{\QCoh, *}$ 正合且与直和交换。所以
$F \circ g_{\QCoh, *} : \text{Mod}_A \to \text{Mod}_B$
是与直和交换的右正合函子，并且把 $A$ 映为 $C$。由《函子与态射》中的
引理 \ref{functors-lemma-functor}，函子 $F \circ g_{\QCoh, *}$
把 $A$-模 $M$ 映为视作 $B$-模的 $M \otimes_A C$。
\end{proof}

\begin{lemma}
\label{stacks-more-morphisms-lemma-universally-injective}
记号如引理 \ref{stacks-more-morphisms-lemma-affine}。假设 $\mathcal{X}$ 是 Noether 的，
并且 $g$ 满且平坦。那么 $B \to C$ 万有单射。
\end{lemma}

\begin{proof}
考虑 $\QCoh(\mathcal{O}_\mathcal{X})$ 中的自然映射
$1 : \mathcal{O}_\mathcal{X} \to g_{\QCoh, *}\mathcal{O}_U$。
拉回到 $U$ 并使用伴随，可知复合
$$
\mathcal{O}_U =
g^*\mathcal{O}_\mathcal{X} \xrightarrow{g^*1}
g^*g_{\QCoh, *}\mathcal{O}_U \to \mathcal{O}_U
$$
是 $\QCoh(\mathcal{O}_U)$ 中的恒等态射。把
$g_{\QCoh, *}\mathcal{O}_U = \colim \mathcal{F}_i$ 写成凝聚
$\mathcal{O}_\mathcal{X}$-模的滤过余极限；参见《叠的上同调》中的引理
\ref{stacks-cohomology-lemma-directed-colimit-coherent}。当 $i$ 足够大时，
映射 $1 : \mathcal{O}_\mathcal{X} \to g_{\QCoh, *}\mathcal{O}_U$
经由 $\mathcal{F}_i$ 分解；参见《叠的上同调》中的引理
\ref{stacks-cohomology-lemma-finite-presentation-quasi-compact-colimit}。
设 $s : \mathcal{O}_\mathcal{X} \to \mathcal{F}_i$ 为该分解。那么
$$
\mathcal{O}_U \xrightarrow{g^*s} g^*\mathcal{F}_i \to
g^*g_{\QCoh, *}\mathcal{O}_U \to \mathcal{O}_U
$$
是恒等态射。换言之，$s$ 拉回到 $U$ 后成为一个直和项的嵌入。
按《叠的上同调》中的引理
\ref{stacks-cohomology-lemma-internal-hom-fp-into-qcoh} 的记号，置
$\mathcal{F}_i^\vee = hom(\mathcal{F}_i, \mathcal{O}_\mathcal{X})$。
特别地，存在求值映射
$ev : \mathcal{F}_i \otimes_{\mathcal{O}_\mathcal{X}} \mathcal{F}_i^\vee
\to \mathcal{O}_\mathcal{X}$。在 $s$ 处求值定义映射
$s^\vee : \mathcal{F}_i^\vee \to \mathcal{O}_\mathcal{X}$。
与关于 $s$ 的陈述对偶，可知 $g^*(s^\vee)$ 满；关于 $hom$ 和
$\otimes$ 与限制到 $U$ 的相容性，参见
% P12-ERR-0139: see ../control/ERRATA_CANDIDATES.jsonl
《叠的上同调》第 \ref{stacks-cohomology-section-further-remarks} 节。
由于 $g$ 满且平坦，可知 $s^\vee$ 满（见同一出处）。由于 $F$ 右正合，
可知 $F(\mathcal{F}_i^\vee) \to F(\mathcal{O}_\mathcal{X}) = B$ 满。
选取 $\lambda \in F(\mathcal{F}_i^\vee)$，使它映到 $1 \in B$。
% P12-ERR-0140: see ../control/ERRATA_CANDIDATES.jsonl
用 $e = F(s)(1) \in F(\mathcal{F}_i)$ 表示 $1$ 在映射
$F(s) : B = F(\mathcal{O}_\mathcal{X}) \to F(\mathcal{F}_i)$ 下的像。
那么映射
$$
F(ev) :
F(\mathcal{F}_i) \otimes_B F(\mathcal{F}_i^\vee) =
F(\mathcal{F}_i \otimes_{\mathcal{O}_\mathcal{X}} \mathcal{F}_i^\vee)
\longrightarrow
F(\mathcal{O}_\mathcal{X}) = B
$$
按照构造把 $e \otimes \lambda$ 映到 $1$。因此映射
$B \to F(\mathcal{F}_i)$、$b \mapsto be$ 万有单射，因为它有单侧逆
$F(\mathcal{F}_i) \to B$、$\xi \mapsto F(ev)(\xi \otimes \lambda)$。
这对所有足够大的 $i$ 都成立，故得结论。
\end{proof}

\begin{lemma}
\label{stacks-more-morphisms-lemma-flat}
设 $B \to C$ 是环映射。如果
\begin{enumerate}
\item 余投影 $C \to C \otimes_B C$ 平坦；
\item $B \to C$ 万有单射，
\end{enumerate}
那么 $B \to C$ 忠实平坦。
\end{lemma}

\begin{proof}
由于 $B \to C$ 万有单射，映射
$\Spec(C) \to \Spec(B)$ 满。因此只须证明 $B \to C$ 平坦，
而这由《下降》中的定理
\ref{descent-theorem-descend-module-properties} 得出。
\end{proof}

\noindent
下面这个非常简单的 K\"unneth 版本在我们写出关于代数叠上拟凝聚模上同调
的 K\"unneth 定理一节后，
% P12-ERR-0141: see ../control/ERRATA_CANDIDATES.jsonl
应当变得过时。

\begin{lemma}
\label{stacks-more-morphisms-lemma-easy-kunneth}
设 $a : \mathcal{Y} \to \mathcal{X}$ 和
$b : \mathcal{Z} \to \mathcal{X}$ 可由概形表示、拟紧、拟分离且平坦。
那么
$a_{\QCoh, *}\mathcal{O}_\mathcal{Y}
\otimes_{\mathcal{O}_\mathcal{X}}
b_{\QCoh, *}\mathcal{O}_\mathcal{Z} =
f_{\QCoh, *}\mathcal{O}_{\mathcal{Y} \times_\mathcal{X} \mathcal{Z}}$，
其中 $f : \mathcal{Y} \times_\mathcal{X} \mathcal{Z} \to \mathcal{X}$
是显然的态射。
\end{lemma}

\begin{proof}
简记 $\mathcal{P} = \mathcal{Y} \times_\mathcal{X} \mathcal{Z}$。
由于 $a \circ \text{pr}_1 = f$ 且 $b \circ \text{pr}_2 = f$，
得到映射 $a_*\mathcal{O}_\mathcal{Y} \to f_*\mathcal{O}_\mathcal{P}$ 和
$b_*\mathcal{O}_\mathcal{Z} \to f_*\mathcal{O}_\mathcal{P}$
（使用相对拉回映射；参见《位点》第
\ref{sites-section-pullback} 节）。因此得到相对杯积
$$
\mu :
a_*\mathcal{O}_\mathcal{Y}
\otimes_{\mathcal{O}_\mathcal{X}}
b_*\mathcal{O}_\mathcal{Z} \longrightarrow
f_*\mathcal{O}_{\mathcal{Y} \times_\mathcal{X} \mathcal{Z}}
$$
应用 $Q$ 及其与张量积的相容性（《叠的上同调》中的备注
\ref{stacks-cohomology-remark-QCoh-tensor}），得到
$\QCoh(\mathcal{O}_\mathcal{X})$ 中的箭头
$Q(\mu) : a_{\QCoh, *}\mathcal{O}_\mathcal{Y}
\otimes_{\mathcal{O}_\mathcal{X}}
b_{\QCoh, *}\mathcal{O}_\mathcal{Z} \to
f_{\QCoh, *}\mathcal{O}_{\mathcal{Y} \times_\mathcal{X} \mathcal{Z}}$。
接着，选取概形 $U$ 和满光滑态射 $U \to \mathcal{X}$。
只须证明 $Q(\mu)$ 到 $U_\etale$ 的限制是同构；参见《叠的上同调》第
\ref{stacks-cohomology-section-further-remarks} 节。继而，由同一节的材料，
只须证明 $\mu$ 到 $U_\etale$ 的限制是同构
（这里使用 $\mu$ 的源和靶都是具有基变换性质的局部拟凝聚模）。
此外，可以在 \'etale 拓扑中计算直像；参见《叠的上同调》中的命题
\ref{stacks-cohomology-proposition-loc-qcoh-flat-base-change}。
最后可见
$a_*\mathcal{O}_\mathcal{Y}|_{U_\etale} =
(V \to U)_{small, *}\mathcal{O}_V$，其中
$V = U \times_\mathcal{X} \mathcal{Y}$。$b_*$ 和 $f_*$ 的情形类似。
因此，结论由概形之间平坦、拟紧、拟分离态射的 K\"unneth 公式得出；
参见《概形的导出范畴》中的引理
\ref{perfect-lemma-kunneth}。
\end{proof}

\begin{lemma}
\label{stacks-more-morphisms-lemma-fully-faithful}
设 $\mathcal{X}$ 是对角态射仿射的代数叠。设 $B$ 是环。设
$f_i : \Spec(B) \to \mathcal{X}$（$i = 1, 2$）是两个态射。设
$t : f_1^* \to f_2^*$ 是张量函子
$f_i^* : \QCoh(\mathcal{O}_\mathcal{X}) \to \text{Mod}_B$ 的同构。
那么存在诱导 $t$ 的 $2$-箭头 $f_1 \to f_2$。
\end{lemma}

\begin{proof}
选取仿射概形 $U = \Spec(A)$ 和满光滑态射
$g : U \to \mathcal{X}$；参见《叠的性质》中的引理
\ref{stacks-properties-lemma-quasi-compact-stack}。由于
$\mathcal{X}$ 的对角态射仿射，可知
$U_i = \Spec(B) \times_{f_i, \mathcal{X}, g} U$ 仿射。
设 $U_i = \Spec(C_i)$。那么 $C_i$ 是引理
\ref{stacks-more-morphisms-lemma-affine} 中利用函子 $F = f_i^*$ 构造的带有环映射
$A \to C_i$ 的 $B$-代数。因此，$t$ 诱导与环映射
$A \to C_1$ 和 $A \to C_2$ 相容的 $B$-代数同构 $C_1 \to C_2$。
换言之，我们有
% P12-ERR-0142: see ../control/ERRATA_CANDIDATES.jsonl
交换图表
$$
\xymatrix{
U_i \ar[r] \ar[d] & U \ar[d]^g \\
\Spec(B) \ar[r]^{f_i} & \mathcal{X}
}
\quad\quad
\xymatrix{
&
U_2 \ar[ld] \ar[d]^{\cong} \ar[rd] \\
\Spec(B) &
U_1 \ar[l] \ar[r] &
U
}
$$
这已经表明，$\mathcal{X}$ 在 $\Spec(B)$ 上的对象 $f_1$ 和 $f_2$
在光滑覆盖 $\{U_1 \to \Spec(B)\}$ 后同构。为证明该同构下降为
$\Spec(B)$ 上 $f_1$ 和 $f_2$ 的同构，必须证明我们的同构
（它来自以上交换图表）与 $U_1 \times_{\Spec(B)} U_1$ 上的下降数据
相容。为此，注意 $U \times_\mathcal{X} U$ 也仿射，存在态射
$g' : U \times_\mathcal{X} U \to \mathcal{X}$，并且
$$
U_i \times_{\Spec(B)} U_i =
\Spec(B) \times_{f_i, \mathcal{X}, g'}
(U \times_\mathcal{X} U)
$$
由此可知，来自同构 $C_1 \to C_2$ 的同构
$C_1 \otimes_B C_1 \to C_2 \otimes_B C_2$ 与态射
$U_i \times_{\Spec(B)} U_i \to U \times_\mathcal{X} U$ 相容。
略去一些细节。
\end{proof}

\begin{lemma}
\label{stacks-more-morphisms-lemma-main}
设 $\mathcal{X}$ 是对角态射仿射的 Noether 代数叠。设 $B$ 是环。
设 $F : \QCoh(\mathcal{O}_\mathcal{X}) \to \text{Mod}_B$
是与直和交换的右正合张量函子。那么 $F$ 来自唯一态射
$\Spec(B) \to \mathcal{X}$。
\end{lemma}

\begin{proof}
选取满光滑态射 $g : U \to \mathcal{X}$，其中
$U = \Spec(A)$ 仿射；参见《叠的性质》中的引理
\ref{stacks-properties-lemma-quasi-compact-stack}。应用引理
\ref{stacks-more-morphisms-lemma-affine}，得到有限型交换 $B$-代数
$C = F(g_{\QCoh, *}\mathcal{O}_U)$ 和环映射 $A \to C$。
由引理 \ref{stacks-more-morphisms-lemma-universally-injective}，环映射 $B \to C$ 万有单射。
考虑代数
$$
C \otimes_B C =
F(g_{\QCoh, *}\mathcal{O}_U
\otimes_{\mathcal{O}_\mathcal{X}}
g_{\QCoh, *}\mathcal{O}_U)
$$
由于 $g$ 平坦、拟紧且拟分离，由引理
\ref{stacks-more-morphisms-lemma-easy-kunneth}，下式第一个等号成立：
$$
g_{\QCoh, *}\mathcal{O}_U
\otimes_{\mathcal{O}_\mathcal{X}}
g_{\QCoh, *}\mathcal{O}_U
=
f_{\QCoh, *}\mathcal{O}_{U \times_\mathcal{X} U} =
g_{\QCoh, *}(\text{pr}_{2, *}\mathcal{O}_{U \times_\mathcal{X} U})
$$
其中 $f : U \times_\mathcal{X} U \to \mathcal{X}$ 是显然的态射，
而 $\text{pr}_2 : U \times_\mathcal{X} U \to U$ 是第二投影。
第二个等号由《叠的上同调》中的引理
\ref{stacks-cohomology-lemma-relative-leray} 及
$f = g \circ \text{pr}_2$ 得出。由于 $\mathcal{X}$ 的对角态射仿射，
可知 $U \times_\mathcal{X} U = \Spec(R)$ 仿射。用
$\text{pr}_2 : A \to R$ 把 $R$ 视为 $A$-代数。综上，得到
$$
C \otimes_B C =
F(g_{\QCoh, *}\mathcal{O}_U
\otimes_{\mathcal{O}_\mathcal{X}}
g_{\QCoh, *}\mathcal{O}_U) =
F(g_{\QCoh, *}(\text{pr}_{2, *}\mathcal{O}_{U \times_\mathcal{X} U})) =
R \otimes_A C
$$
其中最后一个等号由引理 \ref{stacks-more-morphisms-lemma-affine} 的末一断言得出。
由于 $A \to R$ 平坦（因为 $\text{pr}_2$ 作为
$U \to \mathcal{X}$ 的基变换而平坦），可知 $C \otimes_B C$ 在
$C$ 上平坦。由引理 \ref{stacks-more-morphisms-lemma-flat}，可知 $B \to C$ 忠实平坦。

\medskip\noindent
我们断言存在实线交换图表
$$
\xymatrix{
\Spec(C \otimes_B C) \ar@<1ex>[d] \ar@<-1ex>[d] \ar[r] &
U \times_\mathcal{X} U \ar@<1ex>[d] \ar@<-1ex>[d] \\
\Spec(C) \ar[d] \ar[r] &
U \ar[d] \\
\Spec(B) \ar@{..>}[r] &
\mathcal{X}
}
$$
箭头 $\Spec(C) \to U = \Spec(A)$ 来自引理
\ref{stacks-more-morphisms-lemma-affine} 陈述中的环映射 $A \to C$。类似地，箭头
$\Spec(C \otimes_B C) \to U \times_\mathcal{X} U$
来自环映射 $R \to C \otimes_B C$。为验证上方方形交换，使用引理
\ref{stacks-more-morphisms-lemma-fully-faithful}；略去细节。由命题
\ref{stacks-more-morphisms-proposition-stack-fpqc}，可知得到点线箭头
$\Spec(B) \to \mathcal{X}$。

\medskip\noindent
$F$ 是与该点线箭头的拉回相对应的函子这一陈述，也由此以及引理
\ref{stacks-more-morphisms-lemma-affine} 中的对应陈述显然可见。略去细节。
\end{proof}

\noindent
% P12-ERR-0143: see ../control/ERRATA_CANDIDATES.jsonl
对环 $B$，用 $\text{Mod}^{fg}_B$ 表示有限生成 $B$-模
（亦称有限 $B$-模）的范畴。

\begin{theorem}
\label{stacks-more-morphisms-theorem-main}
设 $\mathcal{X}$ 是对角态射仿射的 Noether 代数叠。设 $B$ 是
Noether 环。设
$F : \text{Coh}(\mathcal{O}_\mathcal{X}) \to \text{Mod}^{fg}_B$
是右正合张量函子。那么 $F$ 来自唯一态射
$\Spec(B) \to \mathcal{X}$。
\end{theorem}

\begin{proof}
由引理 \ref{stacks-more-morphisms-lemma-extend}，可以把 $F$ 唯一延拓为右正合张量函子
$F : \QCoh(\mathcal{O}_\mathcal{X}) \to \text{Mod}_B$，并且它与所有
% P12-ERR-0144: see ../control/ERRATA_CANDIDATES.jsonl
直和交换。于是可以应用引理 \ref{stacks-more-morphisms-lemma-main}。
\end{proof}
% Generated mechanically from the frozen stacks-geometry.tex target.
% Do not edit this generated file; edit the bound locale source instead.

% OK, start here.
%

\setcounter{chapter}{106}
\chapter{代数叠的几何}



\phantomsection
\label{stacks-geometry-section-phantom}


\section{引言}
\label{stacks-geometry-section-introduction}

\noindent
本章讨论代数叠的若干几何性质。第
\ref{stacks-geometry-section-multiplicities} 节和第
\ref{stacks-geometry-section-dimension-of-algebraic-stacks} 节的初始版本由
Matthew Emerton 与 Toby Gee 撰写，其原始形式见
\cite{Emerton-Gee-dim}。







\section{半通用环}
\label{stacks-geometry-section-versal}

\noindent
本节阐明形变环与局部 Noether 基上有限型代数叠的局部环之间的关系。

\begin{situation}
\label{stacks-geometry-situation-versal}
这里 $\mathcal{X}$ 是局部 Noether 概形 $S$ 上的局部有限型代数叠。
\end{situation}

\noindent
定义如下。

\begin{definition}
\label{stacks-geometry-definition-versal-ring-at-x}
在情形 \ref{stacks-geometry-situation-versal} 中，设
$x_0 : \Spec(k) \to \mathcal{X}$ 为一态射，其中 $k$ 是 $S$ 上的
有限型域。$\mathcal{X}$ 在 $x_0$ 处的一个
{\it 半通用环（versal ring）}，是一个剩余域为 $k$ 的完备
Noether 局部 $S$-代数 $A$，使得存在 Artin 公理中定义
\ref{artin-definition-versal-formal-object} 所述的半通用形式对象
$(A, \xi_n, f_n)$，且 $\xi_1 \cong x_0$（一个 $2$-同构）。
\end{definition}

\noindent
我们要证明半通用环存在，并且在光滑因子意义下唯一。为此将使用
《Artin 公理》第 \ref{artin-section-predeformation-categories} 节中的
预形变范畴。在我们的情形中，它们总是形变范畴。

\begin{lemma}
\label{stacks-geometry-lemma-deformation-category}
在情形 \ref{stacks-geometry-situation-versal} 中，设
$x_0 : \Spec(k) \to \mathcal{X}$ 为一态射，其中 $k$ 是 $S$ 上的
有限型域。那么 $\mathcal{F}_{\mathcal{X}, k, x_0}$
是形变范畴，并且
$T\mathcal{F}_{\mathcal{X}, k, x_0}$ 与
$\text{Inf}(\mathcal{F}_{\mathcal{X}, k, x_0})$
都是有限维 $k$-向量空间。
\end{lemma}

\begin{proof}
取仿射开子概形 $\Spec(\Lambda) \subset S$，使得
$\Spec(k) \to S$ 经它分解。由《Artin 公理》第
\ref{artin-section-predeformation-categories} 节，得到范畴
$\mathcal{C}_\Lambda$ 上的预形变范畴
$\mathcal{F}_{\mathcal{X}, k, x_0}$。
（正如所引之处指出的，该范畴只依赖于态射
$\Spec(k) \to S$，而不依赖于 $\Lambda$ 的选择。）
由《Artin 公理》中的引理
\ref{artin-lemma-deformation-category} 和
\ref{artin-lemma-algebraic-stack-RS}，
$\mathcal{F}_{\mathcal{X}, k, x_0}$ 实际上是形变范畴。
再由《Artin 公理》中的引理 \ref{artin-lemma-finite-dimension}，
$T\mathcal{F}_{\mathcal{X}, k, x_0}$ 与
$\text{Inf}(\mathcal{F}_{\mathcal{X}, k, x_0})$
是有限维 $k$-向量空间。
\end{proof}




\begin{lemma}
\label{stacks-geometry-lemma-versal-ring}
在情形 \ref{stacks-geometry-situation-versal} 中，设
$x_0 : \Spec(k) \to \mathcal{X}$ 为一态射，其中 $k$ 是 $S$ 上的
有限型域。那么 $\mathcal{X}$ 在 $x_0$ 处存在半通用环。
给定其中任意一对 $A$、$A'$，则对某个 $r$，作为 $S$-代数有
$A \cong A'[[t_1, \ldots, t_r]]$ 或
$A' \cong A[[t_1, \ldots, t_r]]$。
\end{lemma}

\begin{proof}
存在性由引理 \ref{stacks-geometry-lemma-deformation-category} 以及《形式形变论》中的
定义 \ref{formal-defos-definition-deformation-category} 和引理
\ref{formal-defos-lemma-RS-implies-S1-S2}、
\ref{formal-defos-lemma-versal-object-existence} 得出。
由《形式形变论》中的唯一性结果——引理
\ref{formal-defos-lemma-minimal-versal}，存在 $\mathcal{X}$ 在
$x_0$ 处的一个“极小”半通用环 $A$，使得 $\mathcal{X}$ 在 $x_0$
处的任何其他半通用环都同构于某个
$A[[t_1, \ldots, t_r]]$，其中 $r$ 适当。这显然推出第二个陈述。
\end{proof}

\begin{lemma}
\label{stacks-geometry-lemma-versal-ring-field-extension}
在情形 \ref{stacks-geometry-situation-versal} 中，设
$x_0 : \Spec(k) \to \mathcal{X}$ 为一态射，其中 $k$ 是 $S$ 上的
有限型域。设 $l/k$ 是有限域扩张，并用
% P12-ERR-0145: see ../control/ERRATA_CANDIDATES.jsonl
$x_{l, 0} : \Spec(l) \to \mathcal{X}$ 表示诱导的态射。
给定 $\mathcal{X}$ 在 $x_0$ 处的一个半通用环 $A$，存在
$\mathcal{X}$ 在 $x_{l, 0}$ 处的半通用环 $A'$，使得有一个
$S$-代数映射 $A \to A'$，它诱导给定的域扩张 $l/k$，并且在
$\mathfrak m_{A'}$-adic 拓扑下形式光滑。
\end{lemma}

\begin{proof}
这立即由《Artin 公理》中的引理
\ref{artin-lemma-change-of-field} 和《形式形变论》中的引理
\ref{formal-defos-lemma-change-of-field-versal-ring} 得出。
（还使用了 $\mathcal{X}$ 由《Artin 公理》中的引理
\ref{artin-lemma-algebraic-stack-RS} 满足 (RS)。）
\end{proof}

\begin{lemma}
\label{stacks-geometry-lemma-compare-versal-ring-completion}
在情形 \ref{stacks-geometry-situation-versal} 中，设 $x : U \to \mathcal{X}$ 为一态射，
其中 $U$ 是 $S$ 上的局部有限型概形。设 $u_0 \in U$ 是有限型点。
令 $k = \kappa(u_0)$，并用
% P12-ERR-0146: see ../control/ERRATA_CANDIDATES.jsonl
$x_0 : \Spec(k) \to \mathcal{X}$ 表示诱导的映射。下列条件等价：
\begin{enumerate}
\item $x$ 在 $u_0$ 处半通用
（《Artin 公理》中的定义 \ref{artin-definition-versal}）；
\item
$\hat x : \mathcal{F}_{U, k, u_0} \to \mathcal{F}_{\mathcal{X}, k, x_0}$
光滑；
\item 与
$x|_{\Spec(\mathcal{O}_{U, u_0}^\wedge)}$ 相关联的形式对象半通用；以及
% P12-ERR-0147: see ../control/ERRATA_CANDIDATES.jsonl
\item 存在 $U' \subset U$，它是 $x$ 的开邻域，并且
$x|_{U'} : U' \to \mathcal{X}$ 光滑。
\end{enumerate}
而且在此情形下，完备化 $\mathcal{O}_{U, u_0}^\wedge$ 是
$\mathcal{X}$ 在 $x_0$ 处的半通用环。
\end{lemma}

\begin{proof}
由于 $U \to S$ 局部有限型（它是这类态射的复合），
$\Spec(k) \to S$ 也是有限型的（同样是因为它是这类态射的复合）。
因此该陈述有意义。(1) 与 (2) 的等价性正是 $x$ 在 $u_0$ 处半通用
的定义。(1) 与 (3) 的等价性由《Artin 公理》中的引理
\ref{artin-lemma-versality-matches} 给出。因此 (1)、(2)、(3) 等价。

\medskip\noindent
若 $x|_{U'}$ 光滑，则由《Artin 公理》中的引理
\ref{artin-lemma-formally-smooth-on-deformation-categories}，函子
$\hat x : \mathcal{F}_{U, k, u_0} \to \mathcal{F}_{\mathcal{X}, k, x_0}$
光滑。因此 (4) 推出 (1)、(2)、(3)。
反过来，设 $x$ 在 $u_0$ 处半通用。取满的光滑态射
$y : V \to \mathcal{X}$，其中 $V$ 是概形。令
$Z = V \times_\mathcal{X} U$，并取位于 $u_0$ 上方的有限型点
$z_0 \in |Z|$（这是可能的，参见《空间态射》中的引理
\ref{spaces-morphisms-lemma-finite-type-points-surjective-morphism}）。
由《Artin 公理》中的引理 \ref{artin-lemma-base-change-versal}，
态射 $Z \to V$ 在 $z_0$ 处光滑。按定义可找到 $z_0$ 的开邻域
$W \subset Z$，使得 $W \to V$ 光滑。由于 $Z \to U$ 是开映射，
令 $U' \subset U$ 为 $W$ 的像。于是根据我们对叠的光滑态射的定义，
$U' \to \mathcal{X}$ 光滑。

\medskip\noindent
最后一个陈述由定义得出，因为 $\mathcal{O}_{U, u_0}^\wedge$
pro-表示 $\mathcal{F}_{U, k, u_0}$。
\end{proof}

\begin{lemma}
\label{stacks-geometry-lemma-characterize-smoothness}
% P12-ERR-0148: see ../control/ERRATA_CANDIDATES.jsonl
在情形 \ref{stacks-geometry-situation-versal} 中，设
$x_0 : \Spec(k) \to \mathcal{X}$ 为一态射，使得
$\Spec(k) \to S$ 是有限型的，其像为 $s$。设 $A$ 是
$\mathcal{X}$ 在 $x_0$ 处的半通用环。下列条件等价：
\begin{enumerate}
\item $x_0$ 位于 $\mathcal{X} \to S$ 的光滑轨迹中
（《叠的态射》中的引理 \ref{stacks-morphisms-lemma-where-smooth}）；
\item $\mathcal{O}_{S, s} \to A$ 在
$\mathfrak m_A$-adic 拓扑下形式光滑；以及
\item $\mathcal{F}_{\mathcal{X}, k, x_0}$ 无阻碍。
\end{enumerate}
\end{lemma}

\begin{proof}
(2) 与 (3) 的等价性立即由《形式形变论》中的引理
\ref{formal-defos-lemma-smooth-power-series-classical} 得出。

\medskip\noindent
注意，$\mathcal{O}_{S, s} \to A$ 在
$\mathfrak m_A$-adic 拓扑下形式光滑，当且仅当
$\mathcal{O}_{S, s} \to A' = A[[t_1, \ldots, t_r]]$
在 $\mathfrak m_{A'}$-adic 拓扑下形式光滑。
因此由引理 \ref{stacks-geometry-lemma-versal-ring}，(2) 不依赖于半通用环的选择。
接着，设 $l/k$ 是有限扩张，并按引理
\ref{stacks-geometry-lemma-versal-ring-field-extension} 选择 $A \to A'$。
若 $\mathcal{O}_{S, s} \to A$ 在 $\mathfrak m_A$-adic 拓扑下
形式光滑，则 $\mathcal{O}_{S, s} \to A'$ 在
$\mathfrak m_{A'}$-adic 拓扑下形式光滑；参见《代数的进一步讨论》
中的引理 \ref{more-algebra-lemma-compose-formally-smooth}。
反过来，若 $\mathcal{O}_{S, s} \to A'$ 在
$\mathfrak m_{A'}$-adic 拓扑下形式光滑，则
$\mathcal{O}_{S, s}^\wedge \to A'$ 与 $A \to A'$ 正则
（《代数的进一步讨论》中的命题
\ref{more-algebra-proposition-fs-regular}），从而
$\mathcal{O}_{S, s}^\wedge \to A$ 正则
（《代数的进一步讨论》中的引理
\ref{more-algebra-lemma-regular-permanence}），所以
$\mathcal{O}_{S, s} \to A$ 在 $\mathfrak m_A$-adic 拓扑下形式光滑
（仍用前一引理）。因此，对 $k$ 和 $x_0$，(2) 与 (1) 等价，
当且仅当对 $l$ 和 $x_{0, l}$ 也等价。

\medskip\noindent
取概形 $U$ 和光滑态射 $U \to \mathcal{X}$，使得
$\Spec(k) \times_\mathcal{X} U$ 非空。取有限扩张 $l/k$ 和一点
$w_0 : \Spec(l) \to \Spec(k) \times_\mathcal{X} U$。
设 $u_0 \in U$ 为 $w_0$ 的像。将上面的论证分别应用于
$l/k$ 与 $l/\kappa(u_0)$，可知能够约化到 $u_0$。因此可以设
$A = \mathcal{O}_{U, u_0}^\wedge$；参见引理
\ref{stacks-geometry-lemma-compare-versal-ring-completion}。
注意，$x_0$ 位于 $\mathcal{X} \to S$ 的光滑轨迹中，当且仅当
$u_0$ 位于 $U \to S$ 的光滑轨迹中；例如参见《叠的态射》中的引理
\ref{stacks-morphisms-lemma-where-smooth}。
于是 (1) 与 (2) 的等价性由《代数的进一步讨论》中的引理
\ref{more-algebra-lemma-formally-smooth-finite-type} 得出。
\end{proof}

\noindent
回顾 Artin 逼近的一个推论。

\begin{lemma}
\label{stacks-geometry-lemma-Artin-approximation-by-smooth-morphism}
% P12-ERR-0148: see ../control/ERRATA_CANDIDATES.jsonl
在情形 \ref{stacks-geometry-situation-versal} 中，设
$x_0 : \Spec(k) \to \mathcal{X}$ 为一态射，使得
$\Spec(k) \to S$ 是有限型的，其像为 $s$。设 $A$ 是
$\mathcal{X}$ 在 $x_0$ 处的半通用环。
若 $\mathcal{O}_{S, s}$ 是 G-环，则可找到一个光滑态射
$U \to \mathcal{X}$，其源是概形，并可找到一个剩余域为 $k$ 的点
$u_0 \in U$，使得
\begin{enumerate}
\item $\Spec(k) \to U \to \mathcal{X}$ 与给定态射 $x_0$ 相同；
\item 存在同构 $\mathcal{O}_{U, u_0}^\wedge \cong A$。
\end{enumerate}
\end{lemma}

\begin{proof}
设 $(\xi_n, f_n)$ 是 $A$ 上的半通用形式对象。
由《Artin 公理》中的引理 \ref{artin-lemma-effective}，
$\xi = (A, \xi_n, f_n)$ 是有效的。
根据假设，$\mathcal{X}$ 在 $S$ 上局部有限表示
（使用《叠的态射》中的引理
\ref{stacks-morphisms-lemma-noetherian-finite-type-finite-presentation}），
因而由《叠的极限》中的命题
\ref{stacks-limits-proposition-characterize-locally-finite-presentation}
保极限。因此，按《Artin 公理》中的引理
\ref{artin-lemma-approximate-versal} 所述的 Artin 逼近表明，
可以找到态射 $U \to \mathcal{X}$，其源为有限型 $S$-概形，
并含有剩余域为 $k$ 的一点 $u_0 \in U$，它满足 (1) 与 (2)，
且 $U \to \mathcal{X}$ 在 $u_0$ 处半通用。
由引理 \ref{stacks-geometry-lemma-compare-versal-ring-completion}，缩小 $U$ 后，
可以设 $U \to \mathcal{X}$ 光滑。
\end{proof}

\begin{remark}[提升半通用环]
\label{stacks-geometry-remark-upgrade}
在情形 \ref{stacks-geometry-situation-versal} 中，设
$x_0 : \Spec(k) \to \mathcal{X}$ 为一态射，其中 $k$ 是 $S$ 上的
有限型域。设 $A$ 是 $\mathcal{X}$ 在 $x_0$ 处的半通用环。
由《Artin 公理》中的引理 \ref{artin-lemma-effective}，
我们的半通用形式对象实际上来自 $S$ 上的一个态射
$$
\Spec(A) \longrightarrow \mathcal{X}
$$
此外，上述每个结果都可以提升为与该态射相容。列举如下：
% P12-ERR-0149: see ../control/ERRATA_CANDIDATES.jsonl
\begin{enumerate}
\item 在引理 \ref{stacks-geometry-lemma-versal-ring} 中，同构
$A \cong A'[[t_1, \ldots, t_r]]$ 或
$A' \cong A[[t_1, \ldots, t_r]]$
可以选成与这些态射相容；
\item 在引理 \ref{stacks-geometry-lemma-versal-ring-field-extension} 中，
同态 $A \to A'$ 可以选成与这些态射相容；
\item 在引理 \ref{stacks-geometry-lemma-compare-versal-ring-completion} 中，
态射 $\Spec(\mathcal{O}_{U, u_0}^\wedge) \to \mathcal{X}$
是典范映射 $\Spec(\mathcal{O}_{U, u_0}^\wedge) \to U$
与给定映射 $U \to \mathcal{X}$ 的复合；
\item 在引理 \ref{stacks-geometry-lemma-Artin-approximation-by-smooth-morphism} 中，
同构 $\mathcal{O}_{U, u_0}^\wedge \cong A$ 可以选成使
$\Spec(A) \to \mathcal{X}$ 对应于上一项中的典范映射。
\end{enumerate}
在每种情形中，该陈述都由我们的映射与半通用形式元素相容这一事实
得出；但要注意，所蕴含的图表仅在对一个 $2$-箭头作
（非典范）选择后才是 $2$-交换的。尽管如此，这意味着 (1) 中
所蕴含的映射 $A' \to A$ 或 $A \to A'$ 在形式同伦意义下是
良定义的；参见《形式形变论》中的引理
\ref{formal-defos-lemma-versal-unique-up-to-homotopy}。
\end{remark}

\begin{lemma}
\label{stacks-geometry-lemma-versal-ring-flat}
在情形 \ref{stacks-geometry-situation-versal} 中，设
$x_0 : \Spec(k) \to \mathcal{X}$ 为一态射，其中 $k$ 是 $S$ 上的
有限型域。设 $A$ 是 $\mathcal{X}$ 在 $x_0$ 处的半通用环。
那么注 \ref{stacks-geometry-remark-upgrade} 中的态射
$\Spec(A) \to \mathcal{X}$ 是平坦的。
\end{lemma}

\begin{proof}
如果 $S$ 在像点处的局部环是 G-环，则这立即由引理
\ref{stacks-geometry-lemma-Artin-approximation-by-smooth-morphism} 以及
Noether 局部环到其完备化的映射平坦这一事实得出。一般情形证明如下。

\medskip\noindent
步骤 I。若 $A$ 和 $A'$ 是 $\mathcal{X}$ 在 $x_0$ 处的两个
半通用环，则该结果对 $A$ 成立，当且仅当对 $A'$ 成立。
事实上，必要时交换 $A$ 与 $A'$ 后，可以设存在形式光滑映射
$\varphi : A \to A'$，使得复合
$$
\Spec(A') \to \Spec(A) \to \mathcal{X}
$$
就是态射 $\Spec(A') \to \mathcal{X}$；参见引理
\ref{stacks-geometry-lemma-versal-ring} 与注 \ref{stacks-geometry-remark-upgrade}。
由于 $A \to A'$ 忠实平坦，所需等价性由《叠的态射》中的引理
\ref{stacks-morphisms-lemma-composition-flat} 和
\ref{stacks-morphisms-lemma-flat-permanence} 得出。

\medskip\noindent
步骤 II。设 $l/k$ 是有限域扩张。设
$x_{l, 0} : \Spec(l) \to \mathcal{X}$ 为诱导的态射。
设 $A$ 是 $\mathcal{X}$ 在 $x_0$ 处的半通用环，并设
$A \to A'$ 如引理 \ref{stacks-geometry-lemma-versal-ring-field-extension} 中所述。
此时复合
$$
\Spec(A') \to \Spec(A) \to \mathcal{X}
$$
仍然是态射 $\Spec(A') \to \mathcal{X}$；参见注
\ref{stacks-geometry-remark-upgrade}。与之前相同地论证，并用步骤 I 说明半通用环的
选择无关紧要，就得到该引理对 $x_0$ 成立，当且仅当对
$x_{l, 0}$ 成立。

\medskip\noindent
步骤 III。取概形 $U$ 和满的光滑态射
$U \to \mathcal{X}$。于是可以在
$Z = U \times_\mathcal{X} x_0$（这是非空代数空间）上取有限型点
$z_0$。设 $u_0 \in U$ 为 $z_0$ 在 $U$ 中的像。
取一个概形和满的 \'etale 映射 $W \to Z$，使得 $z_0$ 是
闭点 $w_0 \in W$ 的像（参见《空间态射》第
\ref{spaces-morphisms-section-points-finite-type} 节）。
由于 $W \to \Spec(k)$ 和 $W \to U$ 都是有限型的，
$\kappa(w_0)/k$ 与 $\kappa(w_0)/\kappa(u_0)$ 都是有限域扩张
（参见《态射》第 \ref{morphisms-section-points-finite-type} 节）。
两次应用步骤 II，可以把 $x_0$ 替换为
$u_0 \to U \to \mathcal{X}$。于是我们的态射是复合
$$
\Spec(\mathcal{O}_{U, u_0}^\wedge) \to U \to \mathcal{X}
$$
% P12-ERR-0150: see ../control/ERRATA_CANDIDATES.jsonl
第一支箭头平坦，因为 Noether 局部环的完备化平坦
（《代数》中的引理 \ref{algebra-lemma-completion-flat}）；
第二支箭头平坦，因为光滑态射平坦。复合平坦，因为复合保持平坦性。
\end{proof}

\begin{remark}
\label{stacks-geometry-remark-groupoid-defo}
在情形 \ref{stacks-geometry-situation-versal} 中，设
$x_0 : \Spec(k) \to \mathcal{X}$ 为一态射，其中 $k$ 是 $S$ 上的
有限型域。由引理 \ref{stacks-geometry-lemma-deformation-category} 和《形式形变论》
中的定理 \ref{formal-defos-theorem-presentation-deformation-groupoid}，
可知 $\mathcal{F}_{\mathcal{X}, k, x_0}$ 在
$\mathcal{C}_\Lambda$ 上的函子中有一个由光滑 pro-可表群胚给出的
表示。展开定义，这意味着可以选择
\begin{enumerate}
\item 剩余域为 $k$ 的 Noether 完备局部 $\Lambda$-代数 $A$，
以及 $A$ 上的 $\mathcal{F}_{\mathcal{X}, k, x_0}$ 的半通用形式对象
$\xi$；
\item 剩余域为 $k$ 的 Noether 完备局部 $\Lambda$-代数 $B$，
以及同构
$$
\underline{B}|_{\mathcal{C}_\Lambda}
\longrightarrow
\underline{A}|_{\mathcal{C}_\Lambda}
\times_{\underline{\xi}, \mathcal{F}_{\mathcal{X}, k, x_0}, \underline{\xi}}
\underline{A}|_{\mathcal{C}_\Lambda}
$$
\end{enumerate}
两个投影对应于形式光滑映射
$t : A \to B$ 和 $s : A \to B$（因为 $\xi$ 半通用）。
存在映射 $c : B \to B \widehat{\otimes}_{s, A, t} B$，
它使 $(A, B, s, t, c)$ 成为剩余域为 $k$ 的 Noether 完备局部
$\Lambda$-代数范畴中的余群胚
（在 pro-可表函子上，该映射由《形式形变论》中的引理
\ref{formal-defos-lemma-presentation-construction} 构造）。
最后，所引定理说明 $\xi$ 诱导等价
$$
[\underline{A}|_{\mathcal{C}_\Lambda} / \underline{B}|_{\mathcal{C}_\Lambda}]
\longrightarrow
\mathcal{F}_{\mathcal{X}, k, x_0}
$$
它是 $\mathcal{C}_\Lambda$ 上余纤维化群胚之间的等价。
事实上，还得到等价
$$
[\underline{A}/\underline{B}]
\longrightarrow
\widehat{\mathcal{F}}_{\mathcal{X}, k, x_0}
$$
它是完备范畴 $\widehat{\mathcal{C}}_\Lambda$ 上余纤维化群胚之间的
等价（关于这为何成立，参见《形式形变论》第
\ref{formal-defos-section-prorepresentable-groupoids-in-functors} 节
中的讨论）。当然，$A$ 是 $\mathcal{X}$ 在 $x_0$ 处的半通用环。
\end{remark}








\section{代数叠各不可约分支的重数}
\label{stacks-geometry-section-multiplicities}

\noindent
若 $X$ 是局部 Noether 概形，则可将 $X$（仅视为拓扑空间）写成
不可约分支的并，记为 $X = \bigcup T_i$。每个不可约分支都是唯一
一般点 $\xi_i$ 的闭包，而局部环 $\mathcal O_{X,\xi_i}$ 是局部
Artin 环。可以把 {\it $X$ 沿 $T_i$ 的重数}，或称
{\it $T_i$ 在 $X$ 中的重数}，定义为
$$
m_{T_i, X} = \text{长度}_{\mathcal O_{X, \xi_i}} \mathcal O_{X, \xi_i}
$$
换言之，它就是该局部 Artin 环的长度。请与《Chow 同调》第
\ref{chow-section-cycle-of-closed-subscheme} 节比较。

\medskip\noindent
这里的目标是把该定义推广到局部 Noether 代数叠。若 $\mathcal{X}$
是叠，则其拓扑空间 $|\mathcal{X}|$（见《叠的性质》中的定义
\ref{stacks-properties-definition-topological-space}）是局部 Noether 的
（《叠的态射》中的引理
\ref{stacks-morphisms-lemma-Noetherian-topology}）。
$|\mathcal{X}|$ 的不可约分支有时也称为 $\mathcal{X}$ 的不可约分支。
若 $\mathcal{X}$ 拟分离，则 $|\mathcal{X}|$ 是清醒的（sober）
（《叠的态射》中的引理 \ref{stacks-morphisms-lemma-sober-qs}），
但在非拟分离情形下未必如此。例如可考虑非拟分离代数空间
$X = \mathbf{A}^1_\mathbf{C}/\mathbf{Z}$。此外，$|\mathcal{X}|$ 上
并不存在一个可以利用其茎来定义重数的结构层。

\begin{lemma}
\label{stacks-geometry-lemma-map-of-components}
设 $f : U \to \mathcal{X}$ 是从概形到局部 Noether 代数叠的光滑态射。
$|U|$ 的任一不可约分支之像的闭包，都是 $|\mathcal{X}|$ 的不可约分支。
若 $U \to \mathcal{X}$ 满射，则 $|\mathcal{X}|$ 的所有不可约分支都可
由此得到。
\end{lemma}

\begin{proof}
由《叠的性质》中的引理 \ref{stacks-properties-lemma-topology-points}，
映射 $|U| \to |\mathcal{X}|$ 连续且开。设 $T \subset |U|$ 是一个
不可约分支。由于 $U$ 局部 Noether，可以找到包含于 $T$ 的非空仿射开集
$W \subset U$。于是 $f(T) \subset |\mathcal{X}|$ 不可约，并且包含
非空开子集 $f(W)$。因此 $f(T)$ 的闭包不可约且包含一个非空开集，
从而该闭包是不可约分支。

\medskip\noindent
假设 $U \to \mathcal{X}$ 满射，并设 $Z \subset |\mathcal{X}|$ 是
不可约分支。取 $|\mathcal{X}|$ 的一个与 $Z$ 相交的 Noether 开子集
$V$。从 $V$ 中删去其他不可约分支后，可以设 $V \subset Z$。
取非空开集 $f^{-1}(V) \subset |U|$ 的一个不可约分支，并设
$T \subset |U|$ 为其闭包。它是 $|U|$ 的不可约分支；由 $T$ 的选择，
$f(T)$ 的闭包必与 $Z$ 相同。
\end{proof}

\noindent
上一引理尤其适用于局部 Noether 概形之间的光滑态射。下一个引理的
陈述中将隐含地使用这一特例。

\begin{lemma}
\label{stacks-geometry-lemma-multiplicities}
设 $U \to X$ 是局部 Noether 概形之间的光滑态射。
% P12-ERR-0151: see ../control/ERRATA_CANDIDATES.jsonl
设 $T'$ 是 $U$ 的一个不可约分支。设 $T$ 是 $T'$ 的像的闭包所给出的
$X$ 的不可约分支。那么 $m_{T', U} = m_{T, X}$。
\end{lemma}

\begin{proof}
用 $\xi'$ 表示 $T'$ 的一般点，用 $\xi$ 表示 $T$ 的一般点。
设 $A = \mathcal{O}_{X, \xi}$ 且 $B = \mathcal{O}_{U, \xi'}$。
需要证明 $\text{长度}_A A = \text{长度}_B B$。由于
$A \to B$ 是平坦的局部环同态（因为光滑态射平坦），由《代数》中的
引理 \ref{algebra-lemma-pullback-module} 有
$$
\text{长度}_A(A) \text{长度}_B(B/\mathfrak m_A B) =
\text{长度}_B(B)
$$
因此只需证明 $\mathfrak m_A B = \mathfrak m_B$，等价地，只需证明
$B/\mathfrak m_A B$ 约化。由于 $U \to X$ 光滑，其基变换
$U_{\xi} \to \Spec \kappa(\xi)$ 也光滑。由于 $U_{\xi}$ 是域上的
光滑概形，它是约化的，因而其任一点处的局部环也约化
% P12-ERR-0152: see ../control/ERRATA_CANDIDATES.jsonl
（《簇》中的引理 \ref{varieties-lemma-smooth-geometrically-normal}）。特别地，
$$
B/\mathfrak m_A B =
\mathcal{O}_{U, \xi'}/\mathfrak m_{X, \xi}\mathcal{O}_{U, \xi'} =
\mathcal{O}_{U_\xi, \xi'}
$$
约化，这正是所需。
\end{proof}

\noindent
利用这一结果，可以通过在光滑局部考察来说明重数有一个良定义的概念。

\begin{lemma}
\label{stacks-geometry-lemma-multiplicity}
设 $U_1 \to \mathcal{X}$ 和 $U_2 \to \mathcal{X}$ 是从概形到局部
Noether 代数叠 $\mathcal{X}$ 的两个光滑态射。分别设 $T_1'$ 和
$T_2'$ 是 $|U_1|$ 与 $|U_2|$ 的不可约分支。假设 $T_1'$ 与
$T_2'$ 的像的闭包，是 $|\mathcal{X}|$ 的同一个不可约分支 $T$。
那么 $m_{T_1', U_1} = m_{T_2', U_2}$。
\end{lemma}

\begin{proof}
分别设 $V_1$ 与 $V_2$ 是 $T_1'$ 与 $T'_2$ 的稠密子集，并分别在
$U_1$ 与 $U_2$ 中开（见引理 \ref{stacks-geometry-lemma-map-of-components} 的证明）。
$|V_1|$ 与 $|V_2|$ 在 $|\mathcal{X}|$ 中的像，都是不可约子集 $T$
的非空开子集，因而有非空交。由《叠的性质》中的引理
\ref{stacks-properties-lemma-points-cartesian}，映射
$|V_1 \times_\mathcal{X} V_2| \to |V_1| \times_{|\mathcal{X}|} |V_2|$
满射。因此 $V_1 \times_\mathcal{X} V_2$ 是非空代数空间；于是可取
一个源为（非空）概形的 \'etale 满射
$V \to V_1 \times_\mathcal{X} V_2$。
若以 $T'$ 表示 $V$ 的任一不可约分支，则引理
\ref{stacks-geometry-lemma-map-of-components} 表明，$T'$ 在 $U_1$ 中的像的闭包
（相应地，在 $U_2$ 中的像的闭包）等于 $T'_1$（相应地，等于 $T'_2$）。

\medskip\noindent
两次应用引理 \ref{stacks-geometry-lemma-multiplicities}，得到
$$
m_{T_1', U_1} = m_{T', V} = m_{T_2', U_2},
$$
这正是所需。
\end{proof}

\noindent
至此已经完成足够的工作，从而说明下述定义有意义。

\begin{definition}
\label{stacks-geometry-definition-multiplicity}
设 $\mathcal{X}$ 是局部 Noether 代数叠。设
$T \subset |\mathcal{X}|$ 是一个不可约分支。$T$ 在 $\mathcal{X}$ 中的
{\it 重数}定义为 $m_{T, \mathcal{X}} = m_{T', U}$，其中
$f : U \to \mathcal{X}$ 是从概形出发的光滑态射，而
$T' \subset |U|$ 是满足 $f(T') \subset T$ 的不可约分支。
\end{definition}

\noindent
由引理 \ref{stacks-geometry-lemma-map-of-components} 和 \ref{stacks-geometry-lemma-multiplicity}，
该定义与 $f : U \to \mathcal{X}$ 的选择以及映到 $T$ 的不可约分支
$T'$ 的选择无关。

\medskip\noindent
最后指出，把 $\mathcal{X}$ 的不可约分支视为闭子叠有时很方便。
为此，若 $T \subset |\mathcal{X}|$ 是不可约分支，则可考虑唯一的
约化闭子叠 $\mathcal{T} \subset \mathcal{X}$，使得
$|\mathcal{T}| = T$；见《叠的性质》中的定义
\ref{stacks-properties-definition-reduced-induced-stack}。
若 $\mathcal{X}$ 拟分离，则不可约分支是整叠；进一步讨论见
《叠的态射》第 \ref{stacks-morphisms-section-integral-stacks} 节。



\section{形式分支与重数}
\label{stacks-geometry-section-formal-branches}

\noindent
定义 \ref{stacks-geometry-definition-multiplicity} 给出的不可约分支重数概念，与
$\mathcal{X}$ 在有限型点处的半通用环（之谱）的不可约分支重数这一
相关概念之间，有一个比较将很方便。

\medskip\noindent
在情形 \ref{stacks-geometry-situation-versal} 中，设 $x_0 : \Spec(k) \to \mathcal{X}$
为一态射，其中 $k$ 是 $S$ 上的有限型域。设 $A$、$A'$ 是
$\mathcal{X}$ 在 $x_0$ 处的半通用环。必要时交换 $A$ 与 $A'$ 后，
可知存在一个形式光滑的\footnote{意指 $A'$ 同构于 $A$ 上的幂级数环。}
映射 $\varphi : A \to A'$，它与半通用形式对象相容；见引理
\ref{stacks-geometry-lemma-versal-ring} 和注 \ref{stacks-geometry-remark-upgrade}。
此外，$\varphi$ 在形式同伦意义下良定义；见《形式形变论》中的引理
\ref{formal-defos-lemma-versal-unique-up-to-homotopy}。特别地，由
《形式形变论》中的引理
\ref{formal-defos-lemma-homotopic-minimal-prime}，
$\varphi(\mathfrak p)A'$ 是 $A'$ 的良定义理想。由于 $A \to A'$
形式光滑，事实上 $\varphi(\mathfrak p)A'$ 是 $A'$ 的极小素理想，
而 $A'$ 的每个极小素理想都以此方式唯一地来自一个极小素理想
$\mathfrak p \subset A$（把 $A'$ 写成 $A$ 上的幂级数环即可容易证明
所有这些陈述）。因此，回顾极小素理想对应于不可约分支，下述定义有意义。

\begin{definition}
\label{stacks-geometry-definition-formal-branches}
设 $\mathcal{X}$ 是局部 Noether 概形 $S$ 上的局部有限型代数叠。
% P12-ERR-0153: see ../control/ERRATA_CANDIDATES.jsonl
设 $x_0 : \Spec(k) \to \mathcal{X}$ 是一态射，其中 $k$ 是 $S$ 上的
有限型域。{\it $\mathcal{X}$ 穿过 $x_0$ 的形式分支}，是任取
$\mathcal{X}$ 在 $x_0$ 处的一个半通用环时 $\Spec(A)$ 的不可约分支
所成的集合；对 $A$ 的不同选择，按上述步骤作认同。
% P12-ERR-0154: see ../control/ERRATA_CANDIDATES.jsonl
\end{definition}

\noindent
假设在定义 \ref{stacks-geometry-definition-formal-branches} 的情形中给定有限扩张
$l/k$。令 $x_{l, 0} : \Spec(l) \to \mathcal{X}$ 等于
$\Spec(l) \to \Spec(k)$ 与 $x_0$ 的复合。设
$A \to A'$ 如引理 \ref{stacks-geometry-lemma-versal-ring-field-extension} 中所述。
由于 $A \to A'$ 忠实平坦，态射
$$
\Spec(A') \to \Spec(A)
$$
把不可约分支（的一般点）映到不可约分支（的一般点）。
这将是一个满射，但一般并非双射。换言之，得到满射
$$
\text{形式分支：}\mathcal{X}\text{，穿过 }x_{l, 0}
\longrightarrow
\text{形式分支：}\mathcal{X}\text{，穿过 }x_0
$$
事实表明，若 $l/k$ 纯不可分，则该映射也单射
（如果将来需要，我们会在此补上精确陈述与证明）。

\begin{lemma}
\label{stacks-geometry-lemma-branches}
在定义 \ref{stacks-geometry-definition-formal-branches} 的情形中，存在一个典范满射：
从 $\mathcal{X}$ 穿过 $x_0$ 的形式分支集合，映到
$|\mathcal{X}|$ 中包含 $x_0$ 的 $|\mathcal{X}|$ 的不可约分支集合。
\end{lemma}

\begin{proof}
设 $A$ 如定义 \ref{stacks-geometry-definition-formal-branches} 中所述，并设
$\Spec(A) \to \mathcal{X}$ 如注 \ref{stacks-geometry-remark-upgrade} 中所述。
我们断言，$\Spec(A)$ 的一个不可约分支的一般点，映到
$|\mathcal{X}|$ 的一个不可约分支的一般点。取概形 $U$ 和满的光滑态射
$U \to \mathcal{X}$。考虑图表
$$
\xymatrix{
\Spec(A) \times_\mathcal{X} U \ar[d]_p \ar[r]_-q & U \ar[d]^f \\
\Spec(A) \ar[r]^j & \mathcal{X}
}
$$
由引理 \ref{stacks-geometry-lemma-versal-ring-flat}，$j$ 平坦，故 $q$ 平坦。
另一方面，$f$ 满且光滑，因而 $p$ 满且光滑。这说明不可约分支的任一
一般点 $\eta \in \Spec(A)$，都是代数空间
$\Spec(A) \times_\mathcal{X} U$ 的一个余维 $0$ 的点 $\eta'$ 的像
（记号见《空间的性质》第
\ref{spaces-properties-section-generic-points} 节，并在 \'etale
局部环上使用降链性）。由于 $q$ 平坦，$q(\eta')$ 是 $U$ 的余维 $0$
的点（论证相同）。由于 $U$ 是概形，$q(\eta')$ 是 $U$ 的一个
不可约分支的一般点。于是由引理 \ref{stacks-geometry-lemma-map-of-components}，
$q(\eta')$ 在 $|\mathcal{X}|$ 中的像的闭包是不可约分支，断言得证。

\medskip\noindent
显然，该断言给出了定义所需映射的机制。为证明其满射性，取
$u_0 \in U$，使其在 $|\mathcal{X}|$ 中映到 $x_0$。取 $u_0$ 的一个
仿射开邻域 $U' \subset U$。缩小 $U'$ 后，可以设 $U'$ 的每个
不可约分支都经过 $u_0$。于是可把 $\mathcal{X}$ 替换为
$|U'| \to |\mathcal{X}|$ 的像所对应的开子叠。因此可以设 $U$ 仿射，
% P12-ERR-0155: see ../control/ERRATA_CANDIDATES.jsonl
有一点 $u_0$ 映到 $x_0 \in |\mathcal{X}|$，并且 $U$ 的每个
不可约分支都经过 $u_0$。由《叠的性质》中的引理
\ref{stacks-properties-lemma-points-cartesian}，存在一点
$t \in |\Spec(A) \times_\mathcal{X} U|$，它映到 $\Spec(A)$ 的闭点以及
$u_0$。对平坦局部环同态使用降链性
$$
A \longrightarrow
\mathcal{O}_{\Spec(A) \times_\mathcal{X} U, \overline{t}}
\longleftarrow \mathcal{O}_{U, u_0}
$$
可知，$\mathcal{O}_{U, u_0}$ 的每个极小素理想，都是中间局部环的
一个极小素理想的像，而这样的极小素理想又映到 $A$ 的一个极小素理想。
这就证明了满射性。这里略去若干细节。
\end{proof}

\noindent
设 $A$ 是 Noether 完备局部环。于是 $\Spec(A)$ 的不可约分支具有重数；
见第 \ref{stacks-geometry-section-multiplicities} 节的引言。若
$A' = A[[t_1, \ldots, t_r]]$，则态射 $\Spec(A') \to \Spec(A)$
在不可约分支上诱导保持重数的双射（略去容易的证明）。
这一事实与定义 \ref{stacks-geometry-definition-formal-branches} 之前的讨论说明，
下述定义有意义。

\begin{definition}
\label{stacks-geometry-definition-multiplicity-formal-branches}
设 $\mathcal{X}$ 是局部 Noether 概形 $S$ 上的局部有限型代数叠。
% P12-ERR-0156: see ../control/ERRATA_CANDIDATES.jsonl
设 $x_0 : \Spec(k) \to \mathcal{X}$ 是一态射，其中 $k$ 是 $S$ 上的
有限型域。{\it $\mathcal{X}$ 穿过 $x_0$ 的一个形式分支的重数}，
是任取 $\mathcal{X}$ 在 $x_0$ 处的一个半通用环时，$\Spec(A)$ 的
相应不可约分支的重数（见上文讨论）。
\end{definition}

\begin{lemma}
\label{stacks-geometry-lemma-branches-multiplicity}
设 $\mathcal{X}$ 是局部 Noether 概形 $S$ 上的局部有限型代数叠。
% P12-ERR-0157: see ../control/ERRATA_CANDIDATES.jsonl
设 $x_0 : \Spec(k) \to \mathcal{X}$ 是一态射，其中 $k$ 是 $S$ 上的
有限型域，其像为 $s \in S$。若 $\mathcal{O}_{S, s}$ 是 G-环，
则引理 \ref{stacks-geometry-lemma-branches} 中的映射保持重数。
\end{lemma}

\begin{proof}
由引理 \ref{stacks-geometry-lemma-Artin-approximation-by-smooth-morphism}，可以设存在
一个光滑态射 $U \to \mathcal{X}$，其中 $U$ 是概形，并且 $U$ 有一个
$k$-值点 $u_0$，使得 $\mathcal{O}_{U, u_0}^\wedge$ 是
$\mathcal{X}$ 在 $x_0$ 处的半通用环。由引理
\ref{stacks-geometry-lemma-branches} 的证明中对映射的构造
（因为 $A = \mathcal{O}_{U, u_0}^\wedge$，此处大为简化），
只需证明：$U$ 的一个经过 $u_0$ 的不可约分支之重数，等于
$\Spec(\mathcal{O}_{U, u_0}^\wedge)$ 的映入它的任一不可约分支之重数。

\medskip\noindent
转化为交换代数，得到如下问题。设 $C = \mathcal{O}_{U, u_0}$。
它本质上是 $\mathcal{O}_{S, s}$ 上的有限型代数，因而是 G-环
（《代数的进一步讨论》中的命题
\ref{more-algebra-proposition-finite-type-over-G-ring}）。
于是 $A = C^\wedge$，所以 $C \to A$ 是正则环同态。
设 $\mathfrak q \subset C$ 是极小素理想，并设
$\mathfrak p \subset A$ 是位于 $\mathfrak q$ 上方的极小素理想。则
% P12-ERR-0158: see ../control/ERRATA_CANDIDATES.jsonl
$$
R = C_\mathfrak p \longrightarrow A_\mathfrak p = R'
$$
是 Artin 局部环之间的正则环同态。对这种环同态总有
$$
\text{长度}_R R = \text{长度}_{R'} R'
$$
这正是需要证明的，因为左边是 $U$ 上该分支的重数，右边是
$\Spec(A)$ 上该分支的重数。为看出这一等式，先由《代数》中的引理
\ref{algebra-lemma-pullback-module} 得
$$
\text{长度}_R(R) \text{长度}_{R'}(R'/\mathfrak m_R R') =
\text{长度}_{R'}(R')
$$
因此只需证明 $\mathfrak m_R R' = \mathfrak m_{R'}$；这是零维局部环
之间的正则同态的一个推论。
\end{proof}



\section{代数叠的维数理论}
\label{stacks-geometry-section-dimension-of-algebraic-stacks}

\noindent
就我们所知，文献中关于代数叠维数理论的主要结果出自
\cite{Osserman}，其中研究了余维和相对维数的概念。
% P12-ERR-0159: see ../control/ERRATA_CANDIDATES.jsonl
我们更细致地考察代数叠在一点处的维数，并证明若干结果，把一个态射在
源中某点处的纤维维数，与其源和目标的维数联系起来。我们还证明一项结果
（即下文引理 \ref{stacks-geometry-lemma-dimension-formula}），它使我们能够在适当假设下，
用半通用环计算代数叠在一点处的维数。
% P12-ERR-0160: see ../control/ERRATA_CANDIDATES.jsonl

\medskip\noindent
尽管我们并未总是试图把结果优化到最强形式，但大体上力求避免不必要的
假设。然而，在某些把代数叠的性质与它在一点处的半通用环之性质相比较的
结果中，我们只考虑如下代数叠：它在某个局部 Noether 概形基上局部有限
表示，并且该基的所有局部环都是 $G$-环。这样便可使用 Artin 逼近，
把半通用环的几何与叠本身的几何相比较。不过，我们证明的各项陈述未必都
真正需要这一限制性假设。由于我们所考虑的应用均满足它，在有帮助时采用
这一假设即可。

\medskip\noindent
若 $X$ 是概形，则把 $X$ 的维数 $\dim(X)$ 定义为 $X$ 的底拓扑空间的
Krull 维数；若 $x$ 是 $X$ 的一点，则把 $X$ 在 $x$ 处的维数
$\dim_x (X)$ 定义为所有包含 $x$ 的 $X$ 的开子集 $U$ 的维数之最小值；
见《性质》中的定义 \ref{properties-definition-dimension}。
有关系式 $\dim(X) = \sup_{x \in X} \dim_x(X)$；见《性质》中的引理
\ref{properties-lemma-dimension}。若 $X$ 局部 Noether，则
$\dim_x(X)$ 等于 $X$ 的所有经过 $x$ 的不可约分支在 $x$ 处维数的上确界。

\medskip\noindent
若 $X$ 是代数空间且 $x \in |X|$，则定义
$\dim_x X = \dim_u U,$，其中 $U$ 是任一具有 \'etale 满射
$U \to X$ 的概形，而 $u\in U$ 是位于 $x$ 上方的任一点；见
《空间的性质》中的定义
\ref{spaces-properties-definition-dimension-at-point}。令
$\dim(X) = \sup_{x \in |X|} \dim_x(X)$；见《空间的性质》中的定义
\ref{spaces-properties-definition-dimension}。

\begin{remark}
\label{stacks-geometry-remark-dimension-algebraic-space}
一般而言，代数空间 $X$ 在一点 $x$ 处的维数，未必等于其底拓扑空间
$|X|$ 在 $x$ 处的维数。例如，若 $k$ 是特征为零的域，且
$X =  \mathbf{A}^1_k / \mathbf{Z}$，则 $X$ 在其每一点处的维数都是
$1$（即 $\mathbf{A}^1_k$ 的维数）；而 $|X|$ 具有平凡拓扑，故其
Krull 维数为零。另一方面，《代数空间》中的例
\ref{spaces-example-infinite-product} 给出一个代数空间的例子：
它在每一点处的维数都是 $0$，但 $|X|$ 不可约、Krull 维数为 $1$，
并且具有一般点（所以 $|X|$ 在任一点处的维数都是 $1$）。
另见《空间的性质》第 \ref{spaces-properties-section-dimension} 节
对该例的讨论。

\medskip\noindent
另一方面，若 $X$ 是《合宜空间》中的定义
\ref{decent-spaces-definition-very-reasonable} 所称的
{\it 合宜的（decent）}代数空间（特别地，若 $X$ 拟分离；见
《合宜空间》第 \ref{decent-spaces-section-reasonable-decent} 节），
那么 $X$ 在 $x$ 处的维数实际上确实等于 $|X|$ 在 $x$ 处的维数；
见《合宜空间》中的引理
\ref{decent-spaces-lemma-dimension-decent-space}。
\end{remark}

\noindent
为了定义代数叠的维数，先引入如下相对维数概念会很有用：一个源为
代数空间、目标为代数叠的态射，在源中某点处的相对维数。这个定义稍显
繁复，仅仅是因为（与概形情形不同）代数叠或代数空间的点不能描述成
从一个域的谱出发的态射，而只能描述成这类态射的等价类。

\begin{definition}
\label{stacks-geometry-definition-relative-dimension}
若 $f : T \to \mathcal{X}$ 是从代数空间到代数叠的一个局部有限型态射，
% P12-ERR-0161: see ../control/ERRATA_CANDIDATES.jsonl
且 $t \in |T|$ 是一点，其像为 $x \in | \mathcal{X}|$，则如下定义
$f$ 在 $t$ 处的{\it 相对维数}，记作 $\dim_t(T_x),$：
选择一个代表 $x$ 的态射 $\Spec k \to \mathcal{X}$，其源为域的谱；
并选择一点 $t' \in |T \times_{\mathcal{X}} \Spec k|$，它在到
$|T|$ 的投影下映到 $t$（由《叠的性质》中的引理
\ref{stacks-properties-lemma-points-cartesian}，这样的点 $t'$ 存在）；
则
$$
\dim_t(T_x) = \dim_{t'}(T \times_{\mathcal{X}} \Spec k ).
$$
\end{definition}

\noindent
注意，由于 $T$ 是代数空间而 $\mathcal{X}$ 是代数叠，纤维积
$T \times_{\mathcal{X}} \Spec k$ 是代数空间，所以上述拟议定义右端的
量确实有定义（见上文讨论）。

\begin{remark}
\label{stacks-geometry-remark-relative-dimension}
(1)
容易验证 $\dim_t(T_x)$ 良定义，且与计算时所作选择无关。例如可使用
概形之间局部有限型态射的相对维数在基变换下不变这一事实；例如见
《态射》中的引理
\ref{morphisms-lemma-dimension-fibre-after-base-change}。

\medskip\noindent
(2)
当 $\mathcal{X}$ 也是代数空间时，直接检验即可确认这一定义与
《空间的态射》中的定义
\ref{spaces-morphisms-definition-dimension-fibre} 所给出的相对维数
定义一致。
\end{remark}

\noindent
接下来回顾下述引理；我们对局部 Noether 代数叠维数的研究以它为基础。

\begin{lemma}
\label{stacks-geometry-lemma-behaviour-of-dimensions-wrt-smooth-morphisms}
若 $f: U \to X$ 是局部 Noether 代数空间之间的光滑态射，且
$u \in |U|$ 的像为 $x \in |X|$，则
$$
\dim_u (U) = \dim_x(X) + \dim_{u} (U_x)
$$
其中 $\dim_u (U_x)$ 由定义 \ref{stacks-geometry-definition-relative-dimension} 给出。
\end{lemma}

\begin{proof}
见《空间的态射》中的引理
\ref{spaces-morphisms-lemma-smoothness-dimension-spaces}；注意，由注
\ref{stacks-geometry-remark-relative-dimension} (2)，此处使用的 $\dim_u (U_x)$
之定义与该处使用的定义一致。
\end{proof}

\begin{lemma}
\label{stacks-geometry-lemma-dimension-for-stacks}
% P12-ERR-0162: see ../control/ERRATA_CANDIDATES.jsonl
若 $\mathcal{X}$ 是局部 Noether 代数叠，且 $x \in |\mathcal{X}|$。
设 $U \to \mathcal{X}$ 是从代数空间到 $\mathcal{X}$ 的光滑态射，
并设 $u$ 是 $|U|$ 中映到 $x$ 的任一点。则
$$
\dim_x(\mathcal{X}) =  \dim_u(U) - \dim_{u}(U_x)
$$
其中相对维数 $\dim_u(U_x)$ 由定义
\ref{stacks-geometry-definition-relative-dimension} 给出，而 $\mathcal{X}$ 在 $x$
处的维数如《叠的性质》中的定义
\ref{stacks-properties-definition-dimension-at-point} 所述。
\end{lemma}

\begin{proof}
引理 \ref{stacks-geometry-lemma-behaviour-of-dimensions-wrt-smooth-morphisms} 可用来验证，
右端 $\dim_u(U) + \dim_u(U_x)$ 与光滑态射 $U \to \mathcal{X}$ 以及
$u \in |U|$ 的选择无关。
% P12-ERR-0163: see ../control/ERRATA_CANDIDATES.jsonl
略去细节。特别地，可以设 $U$ 是概形。在这种情形下，可以选取
$x$ 的代表为复合 $\Spec \kappa(u) \to U \to \mathcal{X}$ 来计算
$\dim_u(U_x)$，其中第一支态射是像为 $u \in U$ 的典范态射。
然后，若记 $R = U \times_{\mathcal{X}} U$，并用
$e : U \to R$ 表示对角态射，则相对维数在基变换下不变说明
$\dim_u(U_x) = \dim_{e(u)}(R_u)$。因此右端等于
$\dim_u (U) - \dim_{e(u)}(R_u) = \dim_x(\mathcal{X})$，正是所需。
\end{proof}

\begin{remark}
\label{stacks-geometry-remark-dimension-DM}
对于适当合宜的 Deligne--Mumford 叠（例如拟分离叠），
$\dim_x(\mathcal{X})$ 仍将等于拓扑定义的量
$\dim_x |\mathcal{X}|$。然而，对更一般的 Artin 叠，通常并非如此。
例如，若 $\mathcal{X} = [\mathbf{A}^1/\mathbf{G}_m]$
（在某个域上，并就 $\mathbf{G}_m$ 在 $\mathbf{A}^1$ 上的通常乘法
作用取商），则 $|\mathcal{X}|$ 有两个点，其中一点是另一点的特殊化
（对应于 $\mathbf{G}_m$ 在 $\mathbf{A}^1$ 上的两个轨道），故作为
拓扑空间其维数为 $1$；但对两个点 $x \in |\mathcal{X}|$，都有
$\dim_x (\mathcal{X}) = 0$。（一个更极端的例子是分类空间
$[\Spec k/\mathbf{G}_m]$，它在其唯一一点处的维数等于 $-1$。）
\end{remark}

\noindent
现在可以把定义 \ref{stacks-geometry-definition-relative-dimension} 推广到
（局部 Noether）代数叠之间的（局部有限型）态射。

\begin{definition}
\label{stacks-geometry-definition-relative-dimension-for-stacks}
若 $f : \mathcal{T} \to \mathcal{X}$ 是局部 Noether 代数叠之间的
局部有限型态射，且 $t \in |\mathcal{T}|$ 是一点，其像为
$x \in |\mathcal{X}|$，则如下定义 $f$ 在 $t$ 处的{\it 相对维数}，
记作 $\dim_t(\mathcal{T}_x),$：
选择一个代表 $x$ 的态射 $\Spec k \to \mathcal{X}$，其源为域的谱；
并选择一点 $t' \in |\mathcal{T} \times_{\mathcal{X}} \Spec k|$，
它在到 $|\mathcal{T}|$ 的投影下映到 $t$
（由《叠的性质》中的引理
\ref{stacks-properties-lemma-points-cartesian}，这样的点 $t'$ 存在；
% P12-ERR-0164: see ../control/ERRATA_CANDIDATES.jsonl
则
$$
\dim_t(\mathcal{T}_x) = \dim_{t'}(\mathcal{T} \times_{\mathcal{X}} \Spec k ).
$$
\end{definition}

\noindent
注意，由于 $\mathcal{T}$ 是代数叠而 $\mathcal{X}$ 是代数叠，纤维积
$\mathcal{T}\times_{\mathcal{X}} \Spec k$ 是代数叠，并由《叠的态射》
中的引理
\ref{stacks-morphisms-lemma-locally-finite-type-locally-noetherian}
知其局部 Noether。因此，上述拟议定义右端的量由《叠的性质》中的定义
\ref{stacks-properties-definition-dimension-at-point} 给出。

\begin{remark}
\label{stacks-geometry-remark-dimension-tangent-space-well-defined}
标准操作表明 $\dim_t(\mathcal{T}_x)$ 良定义，且与计算时所作选择无关。
\end{remark}

\noindent
现在建立相对维数的一些基本性质；它们是概形态射情形相应陈述的显然推广。

\begin{lemma}
\label{stacks-geometry-lemma-base-change-invariance-of-relative-dimension}
设有局部 Noether 叠的态射所成的笛卡儿方块
$$
\xymatrix{
\mathcal{T}' \ar[d]\ar[r] & \mathcal{T} \ar[d] \\
\mathcal{X}' \ar[r] & \mathcal{X}
}
$$
其中竖直态射局部有限型。若 $t' \in |\mathcal{T}'|$ 在
$|\mathcal{T}|$、$|\mathcal{X}'|$ 和 $|\mathcal{X}|$ 中的像分别为
$t$、$x'$ 和 $x$，则
$\dim_{t'}(\mathcal{T}'_{x'}) = \dim_{t}(\mathcal{T}_x).$
\end{lemma}

\begin{proof}
按定义，两边都可计算为同一个纤维积的维数。
\end{proof}

\begin{lemma}
\label{stacks-geometry-lemma-behaviour-of-dimensions-wrt-smooth-morphisms-stacky}
若 $f: \mathcal{U} \to \mathcal{X}$ 是局部 Noether 代数叠之间的
光滑态射，且 $u \in |\mathcal{U}|$ 的像为 $x \in |\mathcal{X}|$，则
$$
\dim_u (\mathcal{U}) = \dim_x(\mathcal{X}) + \dim_{u} (\mathcal{U}_x).
$$
\end{lemma}

\begin{proof}
取源为概形的光滑满射 $V \to \mathcal{U}$，并设 $v\in |V|$ 是映到
$u$ 的一点。复合 $V \to \mathcal{U} \to \mathcal{X}$ 也光滑，并由
引理 \ref{stacks-geometry-lemma-behaviour-of-dimensions-wrt-smooth-morphisms} 有
$\dim_x(\mathcal{X}) = \dim_v(V) - \dim_v(V_x),$
而 $\dim_u(\mathcal{U}) = \dim_v(V) - \dim_v(V_u).$ 因此
$$
\dim_u(\mathcal{U}) - \dim_x(\mathcal{X}) = \dim_v (V_x) - \dim_v (V_u).
$$

\medskip\noindent
选择 $x$ 的一个代表 $\Spec k \to \mathcal{X}$，并选择位于 $v$ 上方的
一点 $v' \in | V \times_{\mathcal{X}} \Spec k|$，其在
$|\mathcal{U}\times_{\mathcal{X}} \Spec k|$ 中的像为 $u'$；则按定义
$\dim_u(\mathcal{U}_x) = \dim_{u'}(\mathcal{U}\times_{\mathcal{X}} \Spec k),$
并且
$\dim_v(V_x) = \dim_{v'}(V\times_{\mathcal{X}} \Spec k).$

\medskip\noindent
现在
$V\times_{\mathcal{X}} \Spec k \to \mathcal{U}\times_{\mathcal{X}}\Spec k$
是光滑满射（因为它是这种态射的基变换），而其源是代数空间
（因为 $V$ 和 $\Spec k$ 是概形，$\mathcal{X}$ 是代数叠）。
因此再次按定义有
\begin{align*}
\dim_{u'}(\mathcal{U}\times_{\mathcal{X}} \Spec k)
& =
\dim_{v'}(V\times_{\mathcal{X}} \Spec k) -
\dim_{v'}(V \times_{\mathcal{X}} \Spec k)_{u'}) \\
& = \dim_v(V_x) -
\dim_{v'}( (V\times_{\mathcal{X}} \Spec k)_{u'}).
\end{align*}
又
$V\times_{\mathcal{X}} \Spec k \cong
V\times_{\mathcal{U}} (\mathcal{U}\times_{\mathcal{X}} \Spec k),$
所以引理 \ref{stacks-geometry-lemma-base-change-invariance-of-relative-dimension} 表明
$\dim_{v'}((V\times_{\mathcal{X}} \Spec k)_{u'})  = \dim_v(V_u).$
综合以上各式，得到
$$
\dim_u(\mathcal{U}) - \dim_x(\mathcal{X}) =
\dim_u(\mathcal{U}_x),
$$
正是所需。
\end{proof}

\begin{lemma}
\label{stacks-geometry-lemma-relative-dimension-is-semi-continuous}
设 $f: \mathcal{T} \to \mathcal{X}$ 是代数叠之间的局部有限型态射。
% P12-ERR-0165: see ../control/ERRATA_CANDIDATES.jsonl
\begin{enumerate}
\item
函数 $t \mapsto \dim_t(\mathcal{T}_{f(t)})$ 在 $|\mathcal{T}|$ 上
上半连续。
\item 若 $f$ 光滑，则函数
$t \mapsto \dim_t(\mathcal{T}_{f(t)})$ 在 $|\mathcal{T}|$ 上局部常值。
\end{enumerate}
\end{lemma}

\begin{proof}
先假设 $\mathcal{T}$ 是概形 $T$。设 $U \to \mathcal{X}$ 是源为概形的
光滑满射，并令 $T' = T \times_{\mathcal{X}} U$。设
$f': T' \to U$ 是 $f$ 在 $U$ 上的拉回，并设 $g: T' \to T$ 为投影。

\medskip\noindent
引理 \ref{stacks-geometry-lemma-base-change-invariance-of-relative-dimension} 表明，
对 $t' \in T'$ 有
$\dim_{t'}(T'_{f'(t')}) = \dim_{g(t')}(T_{f(g(t'))}),$。
由于 $g$ 光滑且满（它是光滑满射的基变换），$g$ 在底拓扑空间上
诱导的映射由《空间的性质》中的引理
\ref{spaces-properties-lemma-topology-points} 连续且开，并且满射。
因此，只需指出：态射 $f'$ 的第 (1) 项由《空间的态射》中的引理
\ref{spaces-morphisms-lemma-openness-bounded-dimension-fibres} 得出，
而第 (2) 项可由《态射》中的引理
\ref{morphisms-lemma-flat-finite-presentation-CM-fibres-relative-dimension}
或《态射》中的引理
\ref{morphisms-lemma-smooth-omega-finite-locally-free} 得出
（二者都给出概形情形的结果，由此可完全仿照《空间的态射》中的引理
\ref{spaces-morphisms-lemma-openness-bounded-dimension-fibres} 推出
代数空间的相应结果。
% P12-ERR-0166: see ../control/ERRATA_CANDIDATES.jsonl

\medskip\noindent
现在回到一般情形，并选择源为概形的光滑满射
$h:V \to \mathcal{T}$。若 $v \in V$，则基本上按定义有
$$
\dim_{h(v)}(\mathcal{T}_{f(h(v))}) =
\dim_{v}(V_{f(h(v))}) - \dim_{v}(V_{h(v)}).
$$
由于 $V$ 是概形，我们已经证明该等式右边第一项上半连续
（若 $f$ 光滑，它甚至局部常值），而第二项事实上局部常值。
所以它们的差上半连续（若 $f$ 光滑，则局部常值），因而函数
$\dim_{h(v)}(\mathcal{T}_{f(h(v))})$ 在 $|V|$ 上上半连续
（若 $f$ 光滑，则局部常值）。由于态射
$|V| \to |\mathcal{T}|$ 开且满，引理得证。
\end{proof}

\noindent
在继续展开理论之前，先证明两个与概形维数理论有关的引理。

\medskip\noindent
为说明第一个引理的背景，注意，若 $X$ 是有限维概形，则由于
$\dim X$ 被定义为各维数 $\dim_x X$ 的上确界，存在一点
$x \in X$，使得 $\dim_x X = \dim X$。下一个引理表明，还可以取
$x$ 为有限型点。

\begin{lemma}
\label{stacks-geometry-lemma-dimension-achieved-by-finite-type-point}
若 $X$ 是有限维概形，则存在闭点（从而是有限型点）$x \in X$，使得
$\dim_x X = \dim X$。
\end{lemma}

\begin{proof}
令 $d = \dim X$，并取 $X$ 的不可约闭子集的一条极大严格递降链，记为
\begin{equation}
\label{stacks-geometry-equation-maximal-chain}
Z_0 \supset Z_1 \supset \ldots \supset Z_d.
\end{equation}
子集 $Z_d$ 是 $X$ 的极小不可约闭子集，故 $Z_d$ 的每一点都是
$Z_d$ 的一般点。由于概形 $X$ 的底拓扑空间是清醒的，可知 $Z_d$
是单点集，仅由一个闭点 $x \in X$ 构成。若 $U$ 是 $x$ 的任一邻域，
则链
$$
U\cap Z_0 \supset U\cap Z_1 \supset \ldots \supset U\cap Z_d = Z_d =
\{x\}
$$
是 $U$ 的不可约闭子集的一条严格递降链，说明 $\dim U \geq d$。
因此 $\dim_x X \geq d$。反向不等式显然，引理得证。
\end{proof}

\noindent
下一个引理说明，在满足若干温和附加假设的不可约概形上，
$\dim_x X$ 是{\it 常值}函数。

\begin{lemma}
\label{stacks-geometry-lemma-constancy-of-dimension}
若 $X$ 是有限维的不可约、Jacobson、链状且局部 Noether 概形，
则对 $X$ 的每个非空开子集 $U$ 都有 $\dim U = \dim X$。
等价地，$\dim_x X$ 是 $X$ 上的常值函数。
\end{lemma}

\begin{proof}
两个断言的等价性直接由定义得出。于是，设 $U\subset X$ 是非空开子集。
显然 $\dim U \leq \dim X$，而需要证明 $\dim U \geq \dim X.$
令 $d = \dim X$，并取 $X$ 的不可约闭子集的一条极大严格递降链，记为
$$
X = Z_0 \supset Z_1 \supset \ldots \supset Z_d.
$$
由于 $X$ 是 Jacobson 的，极小不可约闭子集 $Z_d$ 等于某个闭点 $x$
所成的 $\{x\}$。

\medskip\noindent
若 $x \in U,$，则
$$
U = U \cap Z_0  \supset U\cap Z_1 \supset \ldots \supset
U\cap Z_d = \{x\}
$$
是 $U$ 的不可约闭子集的一条严格递降链，故
$\dim U \geq d$，正是所需。因此可设 $x \not\in U.$

\medskip\noindent
考虑平坦态射 $\Spec \mathcal{O}_{X,x} \to X$。$X$ 的非空
（从而稠密）开子集 $U$ 拉回为开子集
$V \subset \Spec \mathcal{O}_{X,x}$。把 $U$ 替换为一个非空拟紧、
从而 Noether 的开子集后，可以设嵌入 $U \to X$ 是拟紧态射。
由于拟紧态射的概形论像之形成与平坦基变换交换，
% P12-ERR-0167: see ../control/ERRATA_CANDIDATES.jsonl
由《态射》中的引理
\ref{morphisms-lemma-flat-base-change-scheme-theoretic-image} 可知
$V$ 在 $\Spec \mathcal{O}_{X,x}$ 中稠密，因而特别非空；当然
$x \not\in V.$（这里也用 $x$ 表示 $\Spec \mathcal{O}_{X,x}$ 的闭点，
因为其像等于给定点 $x \in X$。）现在
$\Spec \mathcal{O}_{X,x} \setminus \{x\}$ 是 Jacobson 的；
% P12-ERR-0168: see ../control/ERRATA_CANDIDATES.jsonl
见《性质》中的引理
\ref{properties-lemma-complement-closed-point-Jacobson}。所以 $V$ 包含
$\Spec \mathcal{O}_{X,x} \setminus \{x\}$ 的一个闭点 $z$。$z$ 的像
在 $X$ 中的闭包，是 $X$ 的一个包含 $x$ 的不可约闭子集 $Z$；
它与 $U$ 的交非空，并且不存在真包含于 $Z$ 而又真包含
$\{x\}$ 的不可约闭子集（因为拉回到 $\Spec \mathcal{O}_{X,x}$
在 $X$ 中包含 $x$ 的不可约闭子集，与
$\Spec \mathcal{O}_{X,x}$ 的不可约闭子集之间诱导双射）。
由于 $U \cap Z$ 是 $U$ 的非空闭子集，它包含一个在 $X$ 中闭的点
$u$（因为 $X$ 是 Jacobson 的）。又因为 $U\cap Z$ 是不可约集合 $Z$
的非空（从而稠密）开子集，而它包含一个不属于 $U$ 的点，即 $x$，
所以嵌入 $\{u\} \subset U\cap Z$ 是真包含。

\medskip\noindent
由于 $X$ 链状，链
$$
X = Z_0 \supset Z \supset \{x\} = Z_d
$$
可细化为长度 $d+1$ 的链，而该链必具有形式
$$
X = Z_0 \supset W_1 \supset \ldots \supset W_{d-1} = Z \supset \{x\} = Z_d.
$$
由于 $U\cap Z$ 非空，遂得到
$$
U = U \cap Z_0 \supset U \cap W_1\supset \ldots \supset U\cap W_{d-1}
= U\cap Z \supset \{u\}
$$
这是 $U$ 的不可约闭子集的一条长度为 $d+1$ 的严格递降链，说明
$\dim U \geq d$，正是所需。
\end{proof}

\noindent
引理 \ref{stacks-geometry-lemma-constancy-of-dimension} 的叠论类比将在下文引理
\ref{stacks-geometry-lemma-irreducible-implies-equidimensional} 中证明；但在此之前，
必须引入一个附加定义。之所以需要它，是因为概形为链状这一概念并非
\'etale 局部的（见《代数》中的注
\ref{algebra-remark-universally-catenary-does-not-descend} 的例子；
这使得很难定义代数空间或代数叠为链状是什么意思
（见 \cite[page 3]{Osserman} 中的讨论）。
% P12-ERR-0169: see ../control/ERRATA_CANDIDATES.jsonl
对维数理论的某些方面，下述定义似乎很好地替代了所缺少的链状代数叠概念。

\begin{definition}
\label{stacks-geometry-definition-pseudo-catenary}
称局部 Noether 代数叠 $\mathcal{X}$ 是{\it 拟链状的}，如果存在一个
光滑满射 $U \to \mathcal{X}$，其源是万有链状概形。
\end{definition}

\begin{example}
\label{stacks-geometry-example-pseudo-catenary}
若 $\mathcal{X}$ 在万有链状局部 Noether 概形 $S$ 上局部有限型，且
$U\to \mathcal{X}$ 是源为概形的光滑满射，则复合
$U \to \mathcal{X} \to S$ 局部有限型，所以 $U$ 万有链状；
% P12-ERR-0170: see ../control/ERRATA_CANDIDATES.jsonl
见《态射》中的引理
\ref{morphisms-lemma-universally-catenary-local}。因此
$\mathcal{X}$ 拟链状。
\end{example}

\noindent
下述引理表明，拟链状性质沿有限型态射传递。

\begin{lemma}
\label{stacks-geometry-lemma-catenary-covers}
若 $\mathcal{X}$ 是拟链状局部 Noether 代数叠，而
$\mathcal{Y} \to \mathcal{X}$ 是局部有限型态射，
% P12-ERR-0171: see ../control/ERRATA_CANDIDATES.jsonl
则存在一个光滑满射 $V \to \mathcal{Y}$，其源是万有链状概形；因此
$\mathcal{Y}$ 也拟链状。
\end{lemma}

\begin{proof}
根据假设，可以找到源为万有链状概形的光滑满射
$U \to \mathcal{X}$。基变换
$U\times_{\mathcal{X}} \mathcal{Y}$ 是代数叠；取源为概形的光滑满射
$V \to U\times_{\mathcal{X}} \mathcal{Y}$。于是复合
$V \to U\times_{\mathcal{X}} \mathcal{Y} \to \mathcal{Y}$ 光滑且满
（因为它是光滑满射的复合），而态射
$V \to U\times_{\mathcal{X}} \mathcal{Y} \to U$ 局部有限型
（因为它是局部有限型态射的复合）。由于 $U$ 万有链状，由《态射》中的
引理 \ref{morphisms-lemma-universally-catenary-local}，$V$ 万有链状，
正如所断言。
\end{proof}

\noindent
现在研究函数 $\dim_x(\mathcal{X})$ 在 $|\mathcal{X}|$ 上关于
$|\mathcal{X}|$ 的不可约分支的行为，以及若干相关问题；这里
$\mathcal{X}$ 是某个局部 Noether 叠。

\begin{lemma}
\label{stacks-geometry-lemma-irreducible-implies-equidimensional}
若 $\mathcal{X}$ 是 Jacobson、拟链状且局部 Noether 的代数叠，
并且 $|\mathcal{X}|$ 不可约，则 $\dim_x(\mathcal{X})$ 是
$|\mathcal{X}|$ 上的常值函数。
\end{lemma}

\begin{proof}
只需证明 $\dim_x(\mathcal{X})$ 在 $|\mathcal{X}|$ 上局部常值；
这样它必然常值（因为 $|\mathcal{X}|$ 不可约，故连通）。
由于 $\mathcal{X}$ 拟链状，可以找到光滑满射
$U \to \mathcal{X}$，其中 $U$ 是万有链状概形。若 $\{U_i\}$ 是
$U$ 的拟紧开子概形覆盖，则可把 $U$ 替换为 $\coprod U_i,$；
只需证明函数 $u \mapsto \dim_{f(u)}(\mathcal{X})$ 在 $U_i$ 上
局部常值。由于每次只对一个 $U_i$ 检验，现在略去下标，直接写作
$U$ 而非 $U_i$。由于 $U$ 拟紧，它是有限多个不可约分支的并，
记为 $T_1 \cup \ldots \cup T_n$。注意，每个 $T_i$ 都是 Jacobson、
链状且局部 Noether 的，因为它是 Jacobson、链状且局部 Noether
概形 $U$ 的闭子概形。

\medskip\noindent
由引理 \ref{stacks-geometry-lemma-behaviour-of-dimensions-wrt-smooth-morphisms} 有
$\dim_{f(u)}(\mathcal{X}) = \dim_{u}(U) - \dim_{u}(U_{f(u)}).$
引理 \ref{stacks-geometry-lemma-relative-dimension-is-semi-continuous} (2) 表明，
由于 $f$ 光滑，右端第二项在 $U$ 上局部常值；因此必须证明
$\dim_u(U)$ 在 $U$ 上局部常值。由于 $\dim_u(U)$ 是所有包含 $u$ 的
$U$ 的分支 $T_i$ 的维数 $\dim_u T_i$ 中的最大值，只需证明：
若一点 $u$ 位于两个不同分支 $T_i$ 与 $T_j$ 上（其中 $i \neq j$），
则 $\dim_u T_i = \dim_u T_j$；再注意，在不可约、Jacobson、链状且
局部 Noether 的概形 $T$ 上，$t\mapsto \dim_t T$ 是常值函数
（这由引理 \ref{stacks-geometry-lemma-constancy-of-dimension} 得出）。

\medskip\noindent
令 $V = T_i \setminus (\bigcup_{i' \neq i} T_{i'})$，并令
$W = T_j \setminus (\bigcup_{i' \neq j} T_{i'})$。于是 $V$ 和 $W$
各自都是 $U$ 的非空开子集，故各自在 $|\mathcal{X}|$ 中具有非空开像。
由于 $|\mathcal{X}|$ 不可约，这两个非空开子集在
$|\mathcal{X}|$ 中有非空交。取位于该交中的一点 $x$，并取映到 $x$
的点 $v \in V$ 和 $w\in W$。于是有
$$
\dim T_i = \dim V = \dim_v (U) = \dim_x (\mathcal{X}) + \dim_v (U_x)
$$
并且类似地有
$$
\dim T_j = \dim W = \dim_w (U) = \dim_x (\mathcal{X}) + \dim_w (U_x).
$$
由于 $u \mapsto \dim_u (U_{f(u)})$ 在 $U$ 上局部常值，而
$T_i \cup T_j$ 连通（它是两个不可约、从而连通，并具有非空交的集合之并），
可知 $\dim_v (U_x) = \dim_w(U_x)$。比较前两个等式，即得
$\dim T_i = \dim T_j$，正是所需。
\end{proof}

\begin{lemma}
\label{stacks-geometry-lemma-closed-immersions}
若 $\mathcal{Z} \hookrightarrow \mathcal{X}$ 是局部 Noether 代数叠的
闭浸入，而 $z \in |\mathcal{Z}|$ 的像为
$x \in |\mathcal{X}|$，则
$\dim_z (\mathcal{Z}) \leq \dim_x(\mathcal{X})$。
\end{lemma}

\begin{proof}
取源为概形的光滑满射 $U\to \mathcal{X}$；基变换后的态射
$V = U\times_{\mathcal{X}} \mathcal{Z} \to \mathcal{Z}$ 也光滑且满，
而投影 $V \to U$ 是闭浸入。若 $v \in |V|$ 映到
$z \in |\mathcal{Z}|$，并以 $u$ 表示 $v$ 在 $|U|$ 中的像，则显然
$\dim_v(V) \leq \dim_u(U)$；另一方面，由引理
\ref{stacks-geometry-lemma-base-change-invariance-of-relative-dimension} 有
$\dim_v (V_z) = \dim_u(U_x)$。因此
$$
\dim_z(\mathcal{Z})  = \dim_v(V) - \dim_v(V_z)
\leq \dim_u(U) - \dim_u(U_x) = \dim_x(\mathcal{X}),
$$
正如所断言。
\end{proof}

\begin{lemma}
\label{stacks-geometry-lemma-dimension-via-components}
若 $\mathcal{X}$ 是局部 Noether 代数叠，且
$x \in |\mathcal{X}|$，则
$\dim_x(\mathcal{X}) = \sup_{\mathcal{T}} \{ \dim_x(\mathcal{T}) \} $，
其中 $\mathcal{T}$ 遍历 $|\mathcal{X}|$ 的所有经过 $x$ 的不可约分支
（赋予其诱导约化结构）。
\end{lemma}

\begin{proof}
引理 \ref{stacks-geometry-lemma-closed-immersions} 表明，对每个经过点 $x$ 的不可约分支
$\mathcal{T}$，都有
$\dim_x (\mathcal{T}) \leq \dim_x(\mathcal{X})$。因此为证明引理，
只需证明
\begin{equation}
\label{stacks-geometry-equation-desired-inequality}
\dim_x(\mathcal{X}) \leq
\sup_{\mathcal{T}} \{\dim_x(\mathcal{T})\}.
\end{equation}
设 $U\to\mathcal{X}$ 是由概形给出的光滑覆盖。若 $T$ 是 $U$ 的一个
不可约分支，则以 $\mathcal{T}$ 表示其像在 $\mathcal{X}$ 中的闭包；
它是 $\mathcal{X}$ 的不可约分支。设 $u \in U$ 是映到 $x$ 的一点。
于是
$\dim_x(\mathcal{X})=\dim_uU-\dim_uU_x=\sup_T\dim_uT-\dim_uU_x$，
其中上确界取遍 $U$ 的所有经过 $u$ 的不可约分支。选择一个达到该
上确界的分支 $T$，并注意
$\dim_x(\mathcal{T})=\dim_uT-\dim_u T_x$。
所需不等式（\ref{stacks-geometry-equation-desired-inequality}）现在由显然的不等式
$\dim_u T_x \leq \dim_u U_x.$ 得出。
（注意，若 $\Spec k \to \mathcal{X}$ 是 $x$ 的一个代表，则
$T\times_{\mathcal{X}} \Spec k$ 是
$U\times_{\mathcal{X}} \Spec k$ 的闭子空间。）
\end{proof}

\begin{lemma}
\label{stacks-geometry-lemma-dimension-at-finite-type-point}
若 $\mathcal{X}$ 是局部 Noether 代数叠，且
$x \in |\mathcal{X}|$，则对 $\mathcal{X}$ 的任一包含 $x$ 的开子叠
$\mathcal{V}$，存在有限型点 $x_0 \in |\mathcal{V}|$，使得
$\dim_{x_0}(\mathcal{X}) = \dim_x(\mathcal{V})$。
\end{lemma}

\begin{proof}
取源为概形的光滑满射 $f:U \to \mathcal{X}$，并考虑函数
$u \mapsto \dim_{f(u)}(\mathcal{X});$。由于 $f$ 诱导的态射
$|U| \to |\mathcal{X}|$ 是开的（因为 $f$ 光滑）且满的
（根据假设），并且把有限型点映到有限型点（这正是
$|\mathcal{X}|$ 的有限型点之定义），只需证明：对任意 $u \in U$
以及 $u$ 的任一开邻域，都存在该邻域中的有限型点 $u_0$，使得
$\dim_{f(u_0)}(\mathcal{X}) = \dim_{f(u)}(\mathcal{X}).$
在对问题作此重述后，不再需要 $f$ 满，故可把 $U$ 替换为所考虑点
$u$ 的开邻域，从而把问题约化为证明：对每个 $u \in U$，存在有限型点
$u_0 \in U$，使得
$\dim_{f(u_0)}(\mathcal{X}) = \dim_{f(u)}(\mathcal{X}).$
由引理 \ref{stacks-geometry-lemma-behaviour-of-dimensions-wrt-smooth-morphisms}，
$\dim_{f(u)}(\mathcal{X}) = \dim_u(U) - \dim_u(U_{f(u)}),$
而
$\dim_{f(u_0)}(\mathcal{X}) = \dim_{u_0}(U) - \dim_{u_0}(U_{f(u_0)}).$
由于 $f$ 光滑，当 $u_0$ 在 $U$ 上变化时，表达式
$\dim_{u_0}(U_{f(u_0)})$ 局部常值
（由引理 \ref{stacks-geometry-lemma-relative-dimension-is-semi-continuous} (2)）。
必要时进一步在 $u$ 附近缩小 $U$，可以设它常值。于是问题变成证明：
可以找到有限型点 $u_0 \in U$，使得
$\dim_{u_0}(U) = \dim_u(U)$。按定义，$\dim_u U$ 是各维数
$\dim V$ 的最小值，其中 $V$ 取遍这样的开邻域 $V$：它们位于
原概形中并包含 $u$；所以可进一步
在 $u$ 附近缩小 $U$，使得 $\dim_u U = \dim U$。所需点 $u_0$
的存在性遂由引理 \ref{stacks-geometry-lemma-dimension-achieved-by-finite-type-point}
得出。
\end{proof}

\begin{lemma}
\label{stacks-geometry-lemma-monomorphing-a-component-in-of-the-right-dimension}
设 $\mathcal{T} \hookrightarrow \mathcal{X}$ 是代数叠的局部有限型
单态射，其中 $\mathcal{X}$（从而 $\mathcal{T}$ 也）是 Jacobson、
拟链状且局部 Noether 的。
% P12-ERR-0172: see ../control/ERRATA_CANDIDATES.jsonl
再假设 $\mathcal{T}$ 不可约，具有某个（有限）维数 $d$，而
$\mathcal{X}$ 约化且维数小于或等于 $d$。则存在 $\mathcal{T}$ 的
非空开子叠 $\mathcal{V}$，使得诱导的单态射
$\mathcal{V} \hookrightarrow \mathcal{X}$ 是开浸入，并把
$\mathcal{V}$ 认同为 $\mathcal{X}$ 的一个不可约分支的开子集。
\end{lemma}

\begin{proof}
取源为概形的光滑满射 $f:U \to \mathcal{X}$；由于 $\mathcal{X}$
约化，$U$ 必然也约化。记
$U' = \mathcal{T}\times_{\mathcal{X}} U$。基变换后的态射
$U' \to U$ 是代数空间的局部有限型单态射，因此可表；
% P12-ERR-0173: see ../control/ERRATA_CANDIDATES.jsonl
见《空间的态射》中的引理
\ref{spaces-morphisms-lemma-locally-quasi-finite-separated-representable}
和
\ref{spaces-morphisms-lemma-monomorphism-loc-finite-type-loc-quasi-finite}。
由于 $U$ 是概形，$U'$ 也是概形。投影
$f': U' \to \mathcal{T}$ 仍是光滑满射。设 $u' \in U'$，其像为
$u \in U$。引理
\ref{stacks-geometry-lemma-base-change-invariance-of-relative-dimension} 表明
$\dim_{u'}(U'_{f(u')}) = \dim_u(U_{f(u)}),$；另一方面，由引理
\ref{stacks-geometry-lemma-irreducible-implies-equidimensional} 以及对
$\mathcal{T}$ 和 $\mathcal{X}$ 的假设，有
$\dim_{f'(u')}(\mathcal{T}) =d
\geq \dim_{f(u)}(\mathcal{X})$。因此
\begin{equation}
\label{stacks-geometry-equation-dim-inequality}
\dim_{u'} (U') = \dim_{u'} (U'_{f(u')}) + \dim_{f'(u')}(\mathcal{T})
\\
\geq \dim_u (U_{f(u)}) + \dim_{f(u)}(\mathcal{X}) = \dim_u (U).
\end{equation}
由于 $U' \to U$ 是局部有限型单态射，特别地它非分歧；所以由非分歧
态射的 \'etale 局部结构，即《\'Etale 态射》中的引理
\ref{etale-lemma-finite-unramified-etale-local}，可找到交换图表
$$
\xymatrix{
V' \ar[r]\ar[d] & V \ar[d] \\
U' \ar[r] & U
}
$$
其中概形 $V'$ 非空，竖直箭头为 \'etale 态射，而上方水平箭头是闭浸入。
把 $V$ 替换为一个拟紧开子集，使其像与 $U'$ 的像有非空交，并把
$V'$ 替换为 $V$ 的原像，还可设 $V$（从而 $V'$）拟紧。
由于 $V$ 也局部 Noether，它遂为 Noether，因而是有限多个不可约分支的并。

\medskip\noindent
由于 \'etale 态射保持逐点维数，
% P12-ERR-0174: see ../control/ERRATA_CANDIDATES.jsonl
由《下降》中的引理 \ref{descent-lemma-dimension-at-point-local}
和（\ref{stacks-geometry-equation-dim-inequality}）可推出：对任一点
$v' \in V'$，若其像为 $v \in V$，则
$\dim_{v'}( V') \geq \dim_v(V)$。特别地，$V'$ 的像不可能包含于
$V$ 的两个不同不可约分支的交中；因此可以找到 $V$ 的至少一个
不可约开子集，它与 $V'$ 有非空交。把 $V$ 替换为此子集后，可以设
$V$ 是整概形（因为它既约化又不可约）。由上述维数不等式可知，
闭浸入 $V' \hookrightarrow V$ 实际上是同构。
若以 $W$ 表示 $V'$ 在 $U'$ 中的像，则 $W$ 是 $U'$ 的非空开子集
（因为 \'etale 态射是开映射），而诱导的单态射 $W \to U$ 为
\'etale 态射（因为它在源上 \'etale 局部，即拉回到 $V'$ 后如此），
故为开浸入（因为它是 \'etale 单态射）。因此，若以 $\mathcal{V}$
表示 $W$ 在 $\mathcal{T}$ 中的像，则 $\mathcal{V}$ 是
$\mathcal{T}$ 的稠密（等价地，非空）开子叠，其像在
$\mathcal{X}$ 的一个不可约分支中稠密。最后，注意态射
% P12-ERR-0175: see ../control/ERRATA_CANDIDATES.jsonl
$\mathcal{V} \to \mathcal{X}$ 光滑（因为它与光滑态射
$W\to \mathcal{V}$ 的复合光滑），并且也是单态射，故为开浸入。
\end{proof}

\begin{lemma}
\label{stacks-geometry-lemma-dims-of-images}
设 $f: \mathcal{T} \to \mathcal{X}$ 是 Jacobson、拟链状且局部
Noether 的代数叠之间的局部有限型态射，其源不可约，目标拟分离；
% P12-ERR-0176: see ../control/ERRATA_CANDIDATES.jsonl
并设 $\mathcal{Z} \hookrightarrow \mathcal{X}$ 表示
$\mathcal{T}$ 的概形论像。则对所有 $t \in |T|$，都有
% P12-ERR-0177: see ../control/ERRATA_CANDIDATES.jsonl
$\dim_t( \mathcal{T}_{f(t)}) \geq \dim \mathcal{T}  - \dim \mathcal{Z}$；
而且 $|\mathcal{T}|$ 有一个非空（等价地，稠密）开子集，使得其上等号成立。
\end{lemma}

\begin{proof}
用 $\mathcal{Z}$ 替换 $\mathcal{X}$ 后，可以并且确实设 $f$
在概形论意义下支配，并且 $\mathcal{X}$ 不可约。由纤维维数的
上半连续性（引理 \ref{stacks-geometry-lemma-relative-dimension-is-semi-continuous} (1)），
只需证明等式
$\dim_t( \mathcal{T}_{f(t)}) =\dim \mathcal{T}  - \dim \mathcal{Z}$
对位于 $\mathcal{T}$ 的某个非空开子叠中的 $t$ 成立。因此在论证中，
始终可以自由地把 $\mathcal{T}$ 替换为非空开子叠。

\medskip\noindent
设 $T' \to \mathcal{T}$ 是源为概形的光滑满射，并设 $T$ 是 $T'$ 的
非空拟紧开子集。由于 $\mathcal{Y}$ 拟分离，可知
% P12-ERR-0178: see ../control/ERRATA_CANDIDATES.jsonl
$T \to  \mathcal{Y}$ 拟紧（把《叠的态射》中的引理
\ref{stacks-morphisms-lemma-quasi-compact-permanence} 应用于态射
$T \to \mathcal{Y} \to \Spec \mathbf{Z}$）。因此，若把
$\mathcal{T}$ 替换为 $T$ 在 $\mathcal{T}$ 中的像，则可以设
（援引《叠的态射》中的引理
\ref{stacks-morphisms-lemma-surjection-from-quasi-compact}
% P12-ERR-0179: see ../control/ERRATA_CANDIDATES.jsonl
态射 $f:\mathcal{T} \to \mathcal{X}$ 拟紧。

\medskip\noindent
若取光滑满射 $U \to \mathcal{X}$，其中 $U$ 是概形，则引理
\ref{stacks-geometry-lemma-map-of-components} 保证可以找到 $U$ 的一个不可约开子集
$V$，使得 $V \to \mathcal{X}$ 光滑且在概形论意义下支配。
由于拟紧态射的概形论支配性在平坦基变换下保持，概形论支配态射 $f$
的基变换 $\mathcal{T} \times_{\mathcal{X}} V \to V$ 仍在概形论
意义下支配。设 $Z$ 是一个具有到该纤维积的光滑满射的概形；
于是 $Z \to \mathcal{T} \times_{\mathcal{X}} V \to V$
仍在概形论意义下支配。因此可找到 $Z$ 的一个不可约分支 $C$，
它在概形论意义下支配 $V$。由于复合
$Z \to \mathcal{T}\times_{\mathcal{X}} V \to \mathcal{T}$ 光滑，
而 $\mathcal{T}$ 不可约，引理 \ref{stacks-geometry-lemma-map-of-components} 表明，
源的任一不可约分支在 $|\mathcal{T}|$ 中具有稠密像。现在把 $C$
替换为一个与 $Z$ 的所有其他不可约分支不交的非空开子集 $W$，
再把 $\mathcal{T}$ 和 $\mathcal{X}$ 替换为 $W$ 和 $V$ 的像
（并应用引理 \ref{stacks-geometry-lemma-irreducible-implies-equidimensional}，
可知这不改变 $\mathcal{T}$ 或 $\mathcal{X}$ 的维数）。
若以 $\mathcal{W}$ 表示态射
$W \to \mathcal{T}\times_{\mathcal{X}} V$ 的像，则 $\mathcal{W}$
在 $\mathcal{T}\times_{\mathcal{X}} V$ 中开
（因为态射 $W \to \mathcal{T}\times_{\mathcal{X}} V$ 光滑），
并且不可约（因为它是不可约概形的像）。于是最终得到交换图表
$$
\xymatrix{
W \ar[dr] \ar[r]  & \mathcal{W} \ar[r] \ar[d] & V \ar[d] \\
& \mathcal{T} \ar[r] & \mathcal{X}
}
$$
其中 $W$ 和 $V$ 是概形，竖直箭头光滑且满，斜箭头以及左上水平箭头光滑，
% P12-ERR-0180: see ../control/ERRATA_CANDIDATES.jsonl
而诱导态射
$\mathcal{W} \to \mathcal{T}\times_{\mathcal{X}} V$ 是开浸入。
利用这一图表以及引理陈述中所涉各维数的定义，可把引理的验证约化到
已知的概形情形。

\medskip\noindent
固定 $w \in |W|$，并将其像依次记作
$w' \in |\mathcal{W}|$、$t \in |\mathcal{T}|$、
$v$（在 $|V|$ 中）和 $x$（在 $|\mathcal{X}|$ 中）。基本上按定义
（利用 $\mathcal{W}$ 在 $\mathcal{T}\times_{\mathcal{X}} V$ 中开，
以及基变换的纤维是纤维的基变换）得到等式
$$
\dim_v V_x = \dim_{w'} \mathcal{W}_t
$$
和
$$
\dim_t \mathcal{T}_x = \dim_{w'} \mathcal{W}_v.
$$
由引理 \ref{stacks-geometry-lemma-behaviour-of-dimensions-wrt-smooth-morphisms}
（图表中的斜箭头和右侧竖直箭头，把 $W$ 和 $V$ 实现为叠
$\mathcal{T}$ 和 $\mathcal{X}$ 的概形光滑覆盖），得到
$$
\dim_t \mathcal{T} = \dim_w W - \dim_w W_t
$$
以及
$$
\dim_x \mathcal{X} = \dim_v V - \dim_v V_x.
$$
合并这些等式，得到
$$
\dim_t \mathcal{T}_x - \dim_t \mathcal{T} + \dim_x \mathcal{X}
= \dim_{w'} \mathcal{W}_v - \dim_w W + \dim_w W_t + \dim_v V -
\dim_{w'} \mathcal{W}_t
$$
由于 $W \to \mathcal{W}$ 是光滑满射，在态射
$\Spec \kappa(v) \to V$ 上作基变换后仍然如此
（把 $W \to \mathcal{W}$ 视为 $V$ 上的态射）。
从这一光滑态射得到下列两个等式中的第一个：
$$
\dim_w W_v - \dim_{w'} \mathcal{W}_v = \dim_w (W_v)_{w'} = \dim_w W_{w'};
$$
第二个等式由直接比较所涉及的两个纤维得出。
类似地，若把 $W \to \mathcal{W}$ 视为 $\mathcal{T}$ 上的概形态射，
并在点 $t \in |\mathcal{T}|$ 的某个代表上作基变换，就得到等式
$$
\dim_w W_t - \dim_{w'} \mathcal{W}_t = \dim_w (W_t)_{w'} = \dim_w W_{w'}.
$$
综合以上各式，得到
$$
\dim_t \mathcal{T}_x - \dim_t \mathcal{T} + \dim_x \mathcal{X}
=  \dim_w W_v - \dim_w W + \dim_v V.
$$
我们的目标是证明，当 $t$ 取遍某个非空开子集时，该等式左边为零。
% P12-ERR-0181: see ../control/ERRATA_CANDIDATES.jsonl
当 $w$ 在 $W$ 的一个非空开子集上变化时，其像
$t \in |\mathcal{T}|$ 在 $|\mathcal{T}|$ 的一个非空开子集上变化
（因为 $W \to \mathcal{T}$ 光滑）。

\medskip\noindent
因此问题约化为证明：若 $W\to V$ 是不可约局部 Noether 概形之间的
局部有限型概形论支配态射，则存在其点的一个非空开子集，使得对
其中的 $w\in W$ 有
$\dim_w W_v =\dim_w W - \dim_v V$（其中 $v$ 表示 $w$ 在 $V$ 中的像）。
这是标准事实；为方便读者，回顾其证明。

\medskip\noindent
可把 $W$ 和 $V$ 替换为其底约化子概形，而不改变这一等式成立与否，
故可设它们事实上是整概形。由于 $\dim_w W_v$ 在 $W,$ 上局部常值，
% P12-ERR-0182: see ../control/ERRATA_CANDIDATES.jsonl
必要时把 $W$ 替换为非空开子集，可以设 $\dim_w W_v$ 常值，记其值为
$d$。选择该开子集为仿射的，还可设态射 $W\to V$ 实际上有限型。
必要时把 $V$ 替换为非空开子集
（随后在该开子集上拉回 $W$；所得拉回非空，因为拟紧概形论支配态射的
平坦基变换仍然概形论支配），还可设 $W$ 在 $V$ 上平坦。
于是态射 $W\to V$ 在《态射》中的定义
\ref{morphisms-definition-relative-dimension-d} 的意义下具有相对维数
$d$，并由《态射》中的引理
\ref{morphisms-lemma-rel-dimension-dimension} 得到
$\dim_w(W) = \dim_v(V) + d,$，正是所需。
\end{proof}

\begin{remark}
\label{stacks-geometry-remark-negative-dimension}
我们指出，在上一引理的情形中，未必有
$\dim \mathcal{T} \geq \dim \mathcal{Z}$；这并不与该引理陈述中的
不等式矛盾，因为态射 $f$ 的纤维仍是代数叠，因而维数可能为负。
例如，取 $k$ 为一个域，并把该引理应用于态射
$[\Spec k/\mathbf{G}_m] \to \Spec k$，即可看出这一点。

\medskip\noindent
若进一步假设该引理陈述中的态射 $f$ 是拟 DM 的
（其含义见《叠的态射》中的定义
\ref{stacks-morphisms-definition-separated}；例如，可由代数空间表示的
态射是拟 DM 的），则该态射在目标各点上的纤维都是拟 DM 代数叠，
因而具有非负维数。在这种情形下，该引理确实蕴含
$\dim \mathcal{T} \geq \dim \mathcal{Z}$。事实上，我们得到下面这个
更一般的结果。
\end{remark}

\begin{lemma}
\label{stacks-geometry-lemma-dims-of-images-two}
% P12-ERR-0183: see ../control/ERRATA_CANDIDATES.jsonl
设 $f: \mathcal{T} \to \mathcal{X}$ 是 Jacobson、拟链状且局部
Noether 的代数叠之间的局部有限型态射。假设它拟 DM，其源不可约，
其目标拟分离，并以
$\mathcal{Z} \hookrightarrow \mathcal{X}$ 表示 $\mathcal{T}$
的概形论像。则
$\dim \mathcal{Z} \leq \dim \mathcal{T}$，而且下列两个条件中
恰有一个成立：
\begin{enumerate}
\item 对每个有限型点 $t \in |T|,$ 都有
% P12-ERR-0184: see ../control/ERRATA_CANDIDATES.jsonl
$\dim_t(\mathcal{T}_{f(t)}) > 0,$ 在这种情形下，
$\dim \mathcal{Z} < \dim \mathcal{T}$；或者
\item $\mathcal{T}$ 和 $\mathcal{Z}$ 具有相同维数。
\end{enumerate}
\end{lemma}

\begin{proof}
正如上一备注所指出的，拟 DM 叠的维数总是非负的。因此，对所有
$t \in |\mathcal{T}|$ 都有
$\dim_t \mathcal{T}_{f(t)} \geq 0$，并且等式
$$
\dim_t \mathcal{T}_{f(t)} = \dim_t \mathcal{T} - \dim_{f(t)} \mathcal{Z}
$$
对一个稠密开子集中的点 $t\in |\mathcal{T}|$ 成立。
\end{proof}





\section{局部环的维数}
\label{stacks-geometry-section-dimension-local-ring}

\noindent
代数叠在通常意义下其实并没有局部环，但可以如下定义局部环的维数。

\begin{lemma}
\label{stacks-geometry-lemma-dimension-local-ring-pre}
设 $\mathcal{X}$ 是局部 Noether 代数叠。设 $U \to \mathcal{X}$
是光滑态射，并设 $u \in U$。则
$$
\dim(\mathcal{O}_{U, \overline{u}}) -
\dim(\mathcal{O}_{R_u, e(\overline{u})}) =
2\dim(\mathcal{O}_{U, \overline{u}}) -
\dim(\mathcal{O}_{R, e(\overline{u})})
$$
这里 $R = U \times_\mathcal{X} U$，投影为 $s, t : R \to U$，
对角态射为 $e : U \to R$，而 $R_u$ 是 $s : R \to U$ 在 $u$
上的纤维。
\end{lemma}

\begin{proof}
这是因为
$s : \mathcal{O}_{U, \overline{u}} \to \mathcal{O}_{R, e(\overline{u})}$
是 Noether 局部环之间的平坦局部同态，因而
$$
\dim(\mathcal{O}_{R, e(\overline{u})}) =
\dim(\mathcal{O}_{U, \overline{u}}) +
\dim(\mathcal{O}_{R_u, e(\overline{u})})
$$
由《代数》中的引理
\ref{algebra-lemma-dimension-base-fibre-equals-total} 成立。
\end{proof}

\begin{lemma}
\label{stacks-geometry-lemma-dimension-local-ring}
设 $\mathcal{X}$ 是局部 Noether 代数叠。
% P12-ERR-0185: see ../control/ERRATA_CANDIDATES.jsonl
设 $x \in |\mathcal{X}|$ 是有限型点
（见《叠的态射》中的定义
\ref{stacks-morphisms-definition-finite-type-point}）。
设 $d \in \mathbf{Z}$。下列条件等价：
\begin{enumerate}
\item 存在一个概形 $U$、一个光滑态射 $U \to \mathcal{X}$，以及
映到 $x$ 的有限型点 $u \in U$，使得
$2\dim(\mathcal{O}_{U, \overline{u}}) -
\dim(\mathcal{O}_{R, e(\overline{u})}) = d$；并且
\item 对任意概形 $U$、任意光滑态射 $U \to \mathcal{X}$，以及
任意映到 $x$ 的有限型点 $u \in U$，都有
$2\dim(\mathcal{O}_{U, \overline{u}}) -
\dim(\mathcal{O}_{R, e(\overline{u})}) = d$。
\end{enumerate}
这里 $R = U \times_\mathcal{X} U$，投影为 $s, t : R \to U$，
对角态射为 $e : U \to R$，而 $R_u$ 是 $s : R \to U$ 在 $u$
上的纤维。
\end{lemma}

\begin{proof}
假设有 $x$ 的两个光滑邻域 $(U, u)$ 和 $(U', u')$，其中 $u$ 和
$u'$ 都是有限型点。缩小 $U$ 和 $U'$ 后，可设 $u$ 和 $u'$ 都是
闭点（这由有限型点的定义得出）。然后选取满 \'etale 态射
$W \to U \times_\mathcal{X} U'$。设 $W_u$ 为 $W \to U$ 在 $u$
上的纤维，并设 $W_{u'}$ 为 $W \to U'$ 在 $u'$ 上的纤维。
由于 $u$ 和 $u'$ 映到 $|\mathcal{X}|$ 的同一点，
$W_u \cap W_{u'}$ 非空。因此可选取一个闭点 $w \in W$，它同时映到
$u$ 和 $u'$。这把问题约化为下一段中的讨论。

\medskip\noindent
设 $(U', u') \to (U, u)$ 是 $x$ 的两个光滑邻域之间的光滑态射，
且 $u$ 和 $u'$ 都是闭点。目标是证明：用 $(U, u)$ 定义的不变量
与用 $(U', u')$ 定义的不变量相同。为此注意，
$\mathcal{O}_{U, u} \to \mathcal{O}_{U', u'}$
是 Noether 局部环之间的平坦局部同态，因而
$$
\dim(\mathcal{O}_{U', \overline{u}'}) =
\dim(\mathcal{O}_{U, \overline{u}}) +
\dim(\mathcal{O}_{U'_u, \overline{u}'})
$$
由《代数》中的引理
\ref{algebra-lemma-dimension-base-fibre-equals-total} 成立。
（这里略去把局部环的性质同其严格 Hensel 化的性质联系起来所需的
全部步骤；见《代数进阶》中的章节
\ref{more-algebra-section-permanence-henselization}。）
另一方面，有
$$
R' = U' \times_{U, t} R \times_{s, U} U'
$$
因此可见
$$
\dim(\mathcal{O}_{R', e(\overline{u}')}) =
\dim(\mathcal{O}_{R, e(\overline{u})}) +
\dim(\mathcal{O}_{U'_u \times_u U'_u, (\overline{u}', \overline{u}')})
$$
要证明该引理，只需证明
$$
\dim(\mathcal{O}_{U'_u \times_u U'_u, (\overline{u}', \overline{u}')}) =
2\dim(\mathcal{O}_{U'_u, \overline{u}'})
$$
注意这并不总成立（例如，若 $U'_u$ 是一条曲线，而 $u'$ 是这条曲线的
一般点）。不过，我们知道 $u'$ 是在 $u$ 上局部有限型的代数空间
$U'_u$ 的闭点。在这种情形下，等式成立：首先，由
《簇》中的引理
\ref{varieties-lemma-dimension-product-locally-algebraic}，
有
$\dim_{(u', u')}(U'_u \times_u U'_u) = 2\dim_{u'}(U'_u)$；
其次，在局部代数概形的闭点处，维数与局部环的维数一致，见《簇》中的
引理 \ref{varieties-lemma-dimension-locally-algebraic}。
我们略去把这些概形结论转述为代数空间语言的过程。
\end{proof}

\begin{definition}
\label{stacks-geometry-definition-dimension-local-ring}
设 $\mathcal{X}$ 是局部 Noether 代数叠，并设
$x \in |\mathcal{X}|$ 是有限型点。若引理
\ref{stacks-geometry-lemma-dimension-local-ring} 中的等价条件成立，则称其中的
$d \in \mathbf{Z}$ 为
{\it $\mathcal{X}$ 在 $x$ 处的局部环的维数}。
\end{definition}

\noindent
准确地说，这一定义由引理 \ref{stacks-geometry-lemma-dimension-local-ring-pre}
和《叠的性质》中的定义
\ref{stacks-properties-definition-dimension-at-point} 所启发。
% P12-ERR-0186: see ../control/ERRATA_CANDIDATES.jsonl
本节最后建立一个公式，以便用 $\mathcal{X}$ 在 $x$ 处的半通用环的
性质计算 $\dim_x(\mathcal{X})$。

\begin{lemma}
\label{stacks-geometry-lemma-dimension-formula}
设 $\mathcal{X}$ 是局部 Noether 概形 $S$ 上的局部有限型代数叠。
设 $x_0 : \Spec(k) \to \mathcal{X}$ 是一个态射，其中 $k$ 是
$S$ 上的有限型域。按备注 \ref{stacks-geometry-remark-groupoid-defo}，用一个
余群胚 $(A, B, s, t, c)$ 表示
$\mathcal{F}_{\mathcal{X}, k, x_0}$；这里的两个环都是剩余域为
$k$ 的 Noether 完备局部 $S$-代数。则
$$
\text{局部环维数：}\mathcal{X}\text{ 在 }x_0 =
2\dim A - \dim B
$$
\end{lemma}

\begin{proof}
设 $s \in S$ 是 $x_0$ 的像。若 $\mathcal{O}_{S, s}$ 是 G-环
（这是实践中几乎总能满足的条件），则可如下证明该引理。
由引理 \ref{stacks-geometry-lemma-Artin-approximation-by-smooth-morphism}，可以找到
一个光滑态射 $U \to \mathcal{X}$，其源是概形，并且包含一个剩余域为
$k$ 的点 $u_0 \in U$，使得诱导态射
% P12-ERR-0187: see ../control/ERRATA_CANDIDATES.jsonl
$\Spec(k) \to U \to \mathcal{X}$ 与 $x_0$ 重合，且
$A = \mathcal{O}_{U, u_0}^\wedge$。记
$R = U \times_\mathcal{X} U$。于是可把
$\mathcal{O}_{R, e(u_0)}^\wedge$ 与 $B$ 认同。所需等式遂由定义得出。

\medskip\noindent
在证明余下部分中，我们说明如何在一般情形下证明该引理，但建议读者
跳过这部分。

\medskip\noindent
先证明右边与 $(A, B, s, t, c)$ 的选择无关。具体地，设
$(A', B', s', t', c')$ 是第二种选择。由于 $A$ 和 $A'$ 都是
$\mathcal{X}$ 在 $x_0$ 处的半通用环，必要时交换 $A$ 和 $A'$ 后，
可以选取一个与 $A$ 和 $A'$ 上给定的半通用形式对象
$\xi$ 和 $\xi'$ 相容的形式光滑映射 $A \to A'$。
回忆 $\widehat{\mathcal{C}}_\Lambda$ 有余积，而且这些余积由
$\Lambda$ 上的完备张量积给出；见《形式形变理论》中的引理
\ref{formal-defos-lemma-CLambdahat-coproducts}。
那么，$B$ pro-表示这样一个函子：它把 $\xi$ 推前到
$A \widehat{\otimes}_\Lambda A$ 后所得的两个对象之间的同构。
$B'$ 的情形类似。由此得到
$$
B' =
B \otimes_{(A \widehat{\otimes}_\Lambda A)}
(A' \widehat{\otimes}_\Lambda A')
$$
不难看出
$$
A \widehat{\otimes}_\Lambda A \longrightarrow
A \widehat{\otimes}_\Lambda A' \longrightarrow
A' \widehat{\otimes}_\Lambda A'
$$
是形式光滑的，其相对维数等于形式光滑映射 $A \to A'$ 的相对维数的
$2$ 倍。（这可由一般原理得出；也可直接看出，因为在此特定情形中，
$A'$ 是 $A$ 上关于 $r$ 个变量的幂级数环。）因此 $B \to B'$
形式光滑，其相对维数为 $2(\dim(A') - \dim(A))$，正是所需。

\medskip\noindent
接着，设 $l/k$ 是有限扩张。
% P12-ERR-0188: see ../control/ERRATA_CANDIDATES.jsonl
设 $x_{l, 0} : \Spec(l) \to \mathcal{X}$ 为诱导的点。我们断言，
该公式右边对 $x_0$ 和对 $x_{l, 0}$ 取相同的值。
这可以通过按引理 \ref{stacks-geometry-lemma-versal-ring-field-extension} 选取
$A \to A'$，并完全照搬上一段的论证来证明。细节从略。

\medskip\noindent
最后，照引理 \ref{stacks-geometry-lemma-versal-ring-flat} 的证明进行论证，可以利用
前两段中的相容性，把问题约化为第一段所讨论的情形：此时 $A$ 是某个
在 $\mathcal{X}$ 上光滑的概形 $U$ 在 $u_0$ 处的完备局部环，而
% P12-ERR-0189: see ../control/ERRATA_CANDIDATES.jsonl
$u_0$ 是有限型点。细节从略。
\end{proof}
% Generated mechanically from the frozen moduli.tex target.
% Do not edit this generated file; edit the bound locale source instead.

% OK, start here.
%

\setcounter{chapter}{107}
\chapter{模叠}


\phantomsection
\label{moduli-section-phantom}






\section{引言}
\label{moduli-section-introduction}

\noindent
本章验证若干模空间与模叠的基本性质，其中包括
$\mathit{Hom}$、$\mathit{Isom}$、$\Cohstack_{X/B}$、
$\Quotfunctor_{\mathcal{F}/X/B}$、$\Hilbfunctor_{X/B}$、
$\Picardstack_{X/B}$、$\Picardfunctor_{X/B}$、
$\mathit{Mor}_B(Z, X)$、$\Spacesstack'_{fp, flat, proper}$、
$\Polarizedstack$ 和 $\Complexesstack_{X/B}$。
% P12-ERR-0190: see ../control/ERRATA_CANDIDATES.jsonl
在适当假设下，我们已经证明这些对象是代数空间或代数叠；见
《Quot》中的章节
\ref{quot-section-hom}、
\ref{quot-section-isom}、
\ref{quot-section-stack-coherent-sheaves}、
\ref{quot-section-not-flat}、
\ref{quot-section-quot}、
\ref{quot-section-hilb}、
\ref{quot-section-picard-stack}、
\ref{quot-section-picard-functor}、
\ref{quot-section-relative-morphisms}、
\ref{quot-section-stack-of-spaces}、
\ref{quot-section-polarized} 和
\ref{quot-section-moduli-complexes}。
曲线叠记为 $\textit{Curves}$，它在《Quot》的章节
\ref{quot-section-curves} 中引入，并将在曲线模空间一章中讨论；见
《曲线模空间》中的章节
\ref{moduli-curves-section-stack-curves}。

\medskip\noindent
在某种意义上，本章沿着 Grothendieck 讲义
\cite{Gr-I}、
\cite{Gr-II}、
\cite{Gr-III}、
\cite{Gr-IV}、
\cite{Gr-V} 和
\cite{Gr-VI} 的思路展开。







\section{约定与语言滥用}
\label{moduli-section-conventions}

\noindent
我们继续采用《叠的性质》中的章节
\ref{stacks-properties-section-conventions} 所引入的约定和语言滥用。
除非另有说明，基概形均为 $\Spec(\mathbf{Z})$。






\section{Hom 与 Isom 的性质}
\label{moduli-section-hom-isom}

\noindent
设 $f : X \to B$ 是有限表示的代数空间态射。假设
$\mathcal{F}$ 和 $\mathcal{G}$ 是拟凝聚
$\mathcal{O}_X$-模。若 $\mathcal{G}$ 有限表示，在 $B$ 上平坦，
且其支撑在 $B$ 上真，则由
$$
T/B \longmapsto \Hom_{\mathcal{O}_{X_T}}(\mathcal{F}_T, \mathcal{G}_T)
$$
定义的函子 $\mathit{Hom}(\mathcal{F}, \mathcal{G})$
是 $B$ 上的仿射代数空间。若 $\mathcal{F}$ 有限表示，则
$\mathit{Hom}(\mathcal{F}, \mathcal{G}) \to B$
有限表示。见《Quot》中的命题 \ref{quot-proposition-hom}。

\medskip\noindent
若 $\mathcal{F}$ 和 $\mathcal{G}$ 都有限表示、在 $B$ 上平坦，
且其支撑在 $B$ 上真，则子函子
$$
\mathit{Isom}(\mathcal{F}, \mathcal{G}) \subset
\mathit{Hom}(\mathcal{F}, \mathcal{G})
$$
是 $B$ 上仿射且有限表示的代数空间。见《Quot》中的命题
\ref{quot-proposition-isom}。




\section{凝聚层的叠的性质}
\label{moduli-section-stack-coherent-sheaves}

\noindent
设 $f : X \to B$ 是分离且有限表示的代数空间态射。则参数化具有真支撑
的凝聚模平坦族的叠 $\Cohstack_{X/B}$ 是代数的。见《Quot》中的定理
\ref{quot-theorem-coherent-algebraic-general}。

\begin{lemma}
\label{moduli-lemma-coherent-diagonal-affine-fp}
$\Cohstack_{X/B}$ 在 $B$ 上的对角态射仿射且有限表示。
\end{lemma}

\begin{proof}
《Quot》中的引理 \ref{quot-lemma-coherent-diagonal}
已经证明该对角态射可由代数空间表示。从其证明可知，我们需要证明：
对一对有限表示的 $\mathcal{O}_{X_T}$-模
$\mathcal{F}$、$\mathcal{G}$，若它们在 $T$ 上平坦且支撑在 $T$
上真，则
$\mathit{Isom}(\mathcal{F}, \mathcal{G}) \to T$
仿射且有限表示。这已在章节 \ref{moduli-section-hom-isom} 中讨论。
\end{proof}


\begin{lemma}
\label{moduli-lemma-coherent-qs-lfp}
态射 $\Cohstack_{X/B} \to B$ 拟分离且局部有限表示。
\end{lemma}

\begin{proof}
要检验 $\Cohstack_{X/B} \to B$ 拟分离，需要证明其对角态射拟紧且
拟分离。这立即由引理 \ref{moduli-lemma-coherent-diagonal-affine-fp} 得出。
要证明 $\Cohstack_{X/B} \to B$ 局部有限表示，需要证明
$\Cohstack_{X/B} \to B$ 保持极限；见《叠的极限》中的命题
\ref{stacks-limits-proposition-characterize-locally-finite-presentation}。
这由《Quot》中的引理 \ref{quot-lemma-coherent-limits} 得出
（略去一个小细节）。
\end{proof}

\begin{lemma}
\label{moduli-lemma-coherent-existence-part}
再假设 $X \to B$ 为真且有限表示。则
$\Cohstack_{X/B} \to B$ 满足赋值判据的存在性部分
（《叠的态射》中的定义
\ref{stacks-morphisms-definition-existence}）。
\end{lemma}

\begin{proof}
作基变换后，问题立即约化为下述问题：给定分式域为 $K$ 的赋值环
$R$、一个在 $R$ 上真的代数空间 $X$，以及一个凝聚
$\mathcal{O}_{X_K}$-模 $\mathcal{F}_K$，证明存在一个有限表示的
$\mathcal{O}_X$-模 $\mathcal{F}$，它在 $R$ 上平坦，且其泛纤维为
$\mathcal{F}_K$。
注意，由《空间上的平坦性》中的定理
\ref{spaces-flat-theorem-finite-type-flat}，任意在 $R$ 上平坦的
有限型拟凝聚 $\mathcal{O}_X$-模 $\mathcal{F}$ 都有限表示。
% P12-ERR-0191: see ../control/ERRATA_CANDIDATES.jsonl
以 $j : X_K \to X$ 表示泛纤维的嵌入。由于它是仿射态射
$\Spec(K) \to \Spec(R)$ 的基变换，态射 $j$ 仿射。因此
$j_*\mathcal{F}_K$ 拟凝聚。把它写成
$$
j_*\mathcal{F}_K = \colim \mathcal{F}_i
$$
其中右边是它的有限型拟凝聚 $\mathcal{O}_X$-子模的滤过余极限；
见《空间的极限》中的引理
\ref{spaces-limits-lemma-directed-colimit-finite-type}。
由于 $j_*\mathcal{F}_K$ 是 $X$ 上的 $K$-向量空间层，它在
$\Spec(R)$ 上平坦。因此每个 $\mathcal{F}_i$ 都在 $R$ 上平坦；
这是因为在赋值环上平坦等价于无挠
（《代数进阶》中的引理
\ref{more-algebra-lemma-valuation-ring-torsion-free-flat}），
而子模继承无挠性。
最后，需要证明对某个 $i$，映射
$j^*\mathcal{F}_i \to \mathcal{F}_K$
是同构。由于
$j^*j_*\mathcal{F}_K = \mathcal{F}_K$（略去一个小细节），
且 $j^*$ 正合，可见对所有 $i$，映射
$j^*\mathcal{F}_i \to \mathcal{F}_K$ 都是单射。
由于 $j^*$ 与余极限交换，有
$\mathcal{F}_K = j^*j_*\mathcal{F}_K = \colim j^*\mathcal{F}_i$。
由于 $\mathcal{F}_K$ 凝聚（即有限表示），存在某个 $i$，使得
$j^*\mathcal{F}_i$ 包含 $X$ 的一个仿射 \'etale 覆盖上的全部
（有限多个）生成元。因此，当 $i$ 充分大时，
$j^*\mathcal{F}_i \to \mathcal{F}_K$ 满射。
\end{proof}
\begin{lemma}
\label{moduli-lemma-coherent-functorial}
设 $B$ 是代数空间。设 $\pi : X \to Y$ 是代数空间之间的拟有限态射，
其中两者均在 $B$ 上分离且有限表示。则 $\pi_*$ 诱导一个态射
$\Cohstack_{X/B} \to \Cohstack_{Y/B}$。
\end{lemma}

\begin{proof}
设 $(T \to B, \mathcal{F})$ 是 $\Cohstack_{X/B}$ 的一个对象。
我们断言：
\begin{enumerate}
\item[(a)] $(T \to B, \pi_{T, *}\mathcal{F})$ 是
$\Cohstack_{Y/B}$ 的对象；以及
\item[(b)] 对于 $T' \to T$，有
$\pi_{T', *}(X_{T'} \to X_T)^*\mathcal{F} =
(Y_{T'} \to Y_T)^*\pi_{T, *}\mathcal{F}$。
\end{enumerate}
(b) 保证这一构造按所需定义一个函子
$\Cohstack_{X/B} \to \Cohstack_{Y/B}$。

\medskip\noindent
设 $i : Z \to X_T$ 是由 $\mathcal{F}$ 的第零 Fitting 理想截出的
闭子空间
（《空间上的除子》中的章节
\ref{spaces-divisors-section-fitting-ideals}）。
% P12-ERR-0192: see ../control/ERRATA_CANDIDATES.jsonl
按假设，$Z \to B$ 为真（见《空间的导出范畴》中的章节
\ref{spaces-perfect-section-proper-over-base}）。
另一方面，$i$ 有限表示
（《空间上的除子》中的引理
\ref{spaces-divisors-lemma-fitting-ideal-of-finitely-presented}
以及《空间的态射》中的引理
\ref{spaces-morphisms-lemma-closed-immersion-finite-presentation}）。
存在有限型拟凝聚 $\mathcal{O}_Z$-模 $\mathcal{G}$，使得
$i_*\mathcal{G} = \mathcal{F}$
（《空间上的除子》中的引理
\ref{spaces-divisors-lemma-on-subscheme-cut-out-by-Fit-0}）。
事实上，由《空间上的下降》中的引理
\ref{spaces-descent-lemma-finite-finitely-presented-module}，
$\mathcal{G}$ 作为 $\mathcal{O}_Z$-模有限表示。
% P12-ERR-0193: see ../control/ERRATA_CANDIDATES.jsonl
注意 $\mathcal{G}$ 在 $B$ 上平坦；例如，这是因为
$\mathcal{G}$ 和 $\mathcal{F}$ 的茎相同
（《空间的态射》中的引理
\ref{spaces-morphisms-lemma-stalk-push-closed}）。
注意 $\pi_T \circ i : Z \to Y_T$ 作为拟有限态射的复合仍拟有限，并且
% P12-ERR-0194: see ../control/ERRATA_CANDIDATES.jsonl
$\pi_{T, *}\mathcal{F} = (\pi_T \circ i)_*\mathcal{G})$。
由于 $i$ 仿射，$i_*$ 的形成与基变换交换
（《空间的上同调》中的引理
\ref{spaces-cohomology-lemma-affine-base-change}）。
因此可把 $B$ 替换为 $T$，把 $X$ 替换为 $Z$，把
$\mathcal{F}$ 替换为 $\mathcal{G}$，并把 $Y$ 替换为 $Y_T$，
从而约化到下一段讨论的情形。

\medskip\noindent
假设 $X \to B$ 为真。则由《空间的态射》中的引理
\ref{spaces-morphisms-lemma-universally-closed-permanence}，
$\pi$ 为真，因而由《空间的态射进阶》中的引理
\ref{spaces-more-morphisms-lemma-characterize-finite}，它有限。
由于有限态射仿射，《空间的上同调》中的引理
\ref{spaces-cohomology-lemma-affine-base-change} 表明 (b) 成立。
另一方面，由《空间的态射》中的引理
\ref{spaces-morphisms-lemma-finite-presentation-permanence}，
$\pi$ 有限表示。因此，由《空间上的下降》中的引理
\ref{spaces-descent-lemma-finite-finitely-presented-module}，
$\pi_{T, *}\mathcal{F}$ 有限表示。
最后，例如利用《空间的上同调》中的引理
\ref{spaces-cohomology-lemma-stalk-push-finite} 考察茎，可知
$\pi_{T, *}\mathcal{F} $ 在 $B$ 上平坦。
\end{proof}

\begin{lemma}
\label{moduli-lemma-coherent-open}
设 $B$ 是代数空间。设 $\pi : X \to Y$ 是代数空间之间的开浸入，
其中两者均在 $B$ 上分离且有限表示。则引理
\ref{moduli-lemma-coherent-functorial} 中的态射
$\Cohstack_{X/B} \to \Cohstack_{Y/B}$ 是开浸入。
\end{lemma}

\begin{proof}
从略。提示：若 $\mathcal{F}$ 是 $\Cohstack_{Y/B}$ 在 $T$ 上的
一个对象，并且对 $t \in T$ 有
$\text{Supp}(\mathcal{F}_t) \subset |X_t|$，则对 $t$ 的某个邻域中的
$t' \in T$，同一结论也成立。
\end{proof}

\begin{lemma}
\label{moduli-lemma-coherent-closed}
设 $B$ 是代数空间。设 $\pi : X \to Y$ 是代数空间之间的闭浸入，
其中两者均在 $B$ 上分离且有限表示。则引理
\ref{moduli-lemma-coherent-functorial} 中的态射
$\Cohstack_{X/B} \to \Cohstack_{Y/B}$ 是闭浸入。
\end{lemma}

\begin{proof}
设 $\mathcal{I} \subset \mathcal{O}_Y$ 是把 $X$ 截为 $Y$ 的闭子空间的
理想层。回忆 $\pi_*$ 在拟凝聚 $\mathcal{O}_X$-模的范畴与被
$\mathcal{I}$ 零化的拟凝聚 $\mathcal{O}_Y$-模的范畴之间诱导等价；
见《空间的态射》中的引理
\ref{spaces-morphisms-lemma-i-star-equivalence}。
% P12-ERR-0195: see ../control/ERRATA_CANDIDATES.jsonl
对 $T \to B$ 作基变换并把 $\mathcal{I}$ 替换为理想层
$\mathcal{I}_T = \Im((Y_T \to Y)^*\mathcal{I} \to \mathcal{O}_{Y_T})$
后，作相应变更可得到同一结论。
分析引理 \ref{moduli-lemma-coherent-functorial} 的证明可知，态射
$\Cohstack_{X/B} \to \Cohstack_{Y/B}$ 的本质像恰好由对象
$\xi = (T \to B, \mathcal{F})$ 组成，其中
$\mathcal{F}$ 被 $\mathcal{I}_T$ 零化。
换言之，$\xi$ 位于本质像中，当且仅当乘法映射
$$
\mathcal{F} \otimes_{\mathcal{O}_{Y_T}} (Y_T \to Y)^*\mathcal{I}
\longrightarrow
\mathcal{F}
$$
为零，并且在任意进一步的基变换 $T' \to T$ 后也同样为零。
注意
$$
(Y_{T'} \to Y_T)^*(
\mathcal{F} \otimes_{\mathcal{O}_{Y_T}} (Y_T \to Y)^*\mathcal{I}) =
(Y_{T'} \to Y_T)^*\mathcal{F} \otimes_{\mathcal{O}_{Y_{T'}}}
(Y_{T'} \to Y)^*\mathcal{I})
$$
因此，由《空间上的平坦性》中的引理
\ref{spaces-flat-lemma-F-zero-closed-proper}，该乘法映射在 $T'$
上为零这一条件可由 $T$ 的某个闭子空间表示。
\end{proof}
\begin{situation}[数值不变量]
\label{moduli-situation-numerical}
设 $f : X \to B$ 如本节引言所述。设 $I$ 是一个集合，并对
$i \in I$ 设 $E_i \in D(\mathcal{O}_X)$ 为完美对象。给定
$\Cohstack_{X/B}$ 的一个对象 $(T \to B, \mathcal{F})$，
以 $E_{i, T}$ 表示 $E_i$ 到 $X_T$ 的导出拉回。$D(\mathcal{O}_T)$
中的对象
$$
K_i = Rf_{T, *}(E_{i, T} \otimes_{\mathcal{O}_{X_T}}^\mathbf{L} \mathcal{F})
$$
是完美的，且其形成与基变换交换；见《空间的导出范畴》中的引理
\ref{spaces-perfect-lemma-base-change-tensor-perfect}。
因此，由《空间的导出范畴》中的引理
\ref{spaces-perfect-lemma-chi-locally-constant}，函数
$$
\chi_i : |T| \longrightarrow \mathbf{Z},\quad
\chi_i(t) =
\chi(X_t, E_{i, t} \otimes_{\mathcal{O}_{X_t}}^\mathbf{L} \mathcal{F}_t) =
\chi(K_i \otimes_{\mathcal{O}_T}^\mathbf{L} \kappa(t))
$$
局部常值。设 $P : I \to \mathbf{Z}$ 是一个映射。考虑子叠
$$
\Cohstack^P_{X/B} \subset \Cohstack_{X/B}
$$
它由数值不变量与 $P$ 一致、具有真支撑的凝聚层平坦族组成。
更准确地，$\Cohstack_{X/B}$ 的一个对象
$(T \to B, \mathcal{F})$ 位于 $\Cohstack^P_{X/B}$ 中，当且仅当
对所有 $i \in I$ 和 $t \in T$ 都有 $\chi_i(t) = P(i)$。
\end{situation}

\begin{lemma}
\label{moduli-lemma-open-P}
在情形 \ref{moduli-situation-numerical} 中，叠
$\Cohstack^P_{X/B}$ 是代数的，而且
$$
\Cohstack^P_{X/B} \longrightarrow \Cohstack_{X/B}
$$
是平坦闭浸入。若 $I$ 有限，或 $B$ 局部 Noether，则
$\Cohstack^P_{X/B}$ 是 $\Cohstack_{X/B}$ 的既开又闭子叠。
\end{lemma}

\begin{proof}
若 $I$ 有限，这由函数 $t \mapsto \chi_i(t)$ 局部常值立即可知。
若 $I$ 无限，则写成
$$
I = \bigcup\nolimits_{I' \subset I\text{ 有限}} I'
$$
并记 $P' = P|_{I'}$。于是
$$
\Cohstack^P_{X/B} = \bigcap\nolimits_{I' \subset I\text{ 有限}}
\Cohstack^{P'}_{X/B}
$$
因此，$\Cohstack^P_{X/B}$ 总是代数叠，且态射
$\Cohstack^P_{X/B} \subset \Cohstack_{X/B}$ 总是平坦闭浸入，
但它可能不再是开子叠。（我们把例子的构造留给读者。）
不过，若 $B$ 局部 Noether，则由引理 \ref{moduli-lemma-coherent-qs-lfp}
和《叠的态射》中的引理
\ref{stacks-morphisms-lemma-locally-finite-type-locally-noetherian}，
$\Cohstack_{X/B}$ 也局部 Noether。因此，若
$U \to \Cohstack_{X/B}$ 是一个光滑满射，且 $U$ 是局部 Noether
概形，则既开又闭子叠 $\Cohstack^{P'}_{X/B}$ 的逆像在 $U$ 中的
交集是开集（因为局部 Noether 拓扑空间的连通分支是开的）。
这就证明了这种情形下的结论。
\end{proof}

\begin{lemma}
\label{moduli-lemma-finite-list-perfect-objects}
设 $f : X \to B$ 如本节引言所述。设
$E_1, \ldots, E_r \in D(\mathcal{O}_X)$ 为完美对象。
令 $I = \mathbf{Z}^{\oplus r}$，并考虑映射
$$
I \longrightarrow D(\mathcal{O}_X),\quad
(n_1, \ldots, n_r) \longmapsto
E_1^{\otimes n_1}
\otimes \ldots \otimes
E_r^{\otimes n_r}
$$
% P12-ERR-0196: see ../control/ERRATA_CANDIDATES.jsonl
设 $P : I \to \mathbf{Z}$ 是一个映射。则情形
\ref{moduli-situation-numerical} 中定义的
$\Cohstack^P_{X/B} \subset \Cohstack_{X/B}$
是既开又闭子叠。
\end{lemma}

\begin{proof}
可以在 $B$ 上 \'etale 局部地工作，故可设 $B$ 仿射。
在这种情形下，可以作绝对 Noether 约化；建议读者跳过此证明。
具体地，设 $B = \Spec(\Lambda)$。把
$\Lambda = \colim \Lambda_i$ 写成滤过余极限，其中每个
$\Lambda_i$ 都在 $\mathbf{Z}$ 上有限型。对某个 $i$，可以找到
一个分离且有限表示的代数空间态射
$X_i \to \Spec(\Lambda_i)$，其到 $\Lambda$ 的基变换为 $X$。
见《空间的极限》中的引理
\ref{spaces-limits-lemma-descend-finite-presentation} 和
\ref{spaces-limits-lemma-descend-separated-morphism}。
增大 $i$ 后，可进一步假设在 $D(\mathcal{O}_{X_i})$ 中存在完美对象
$E_{1, i}, \ldots, E_{r, i}$，它们到 $X$ 的导出拉回分别同构于
$E_1, \ldots, E_r$；见《空间的导出范畴》中的引理
\ref{spaces-perfect-lemma-perfect-on-limit}。
显然有 Cartesian 方块
$$
\xymatrix{
\Cohstack^P_{X/B} \ar[r] \ar[d] &
\Cohstack_{X/B} \ar[d] \\
\Cohstack^P_{X_i/\Spec(\Lambda_i)} \ar[r] &
\Cohstack_{X_i/\Spec(\Lambda_i)}
}
$$
因而可用引理 \ref{moduli-lemma-open-P} 完成证明。
\end{proof}

\begin{example}[具有固定 Hilbert 多项式的凝聚层]
\label{moduli-example-hilbert-polynomial}
设 $f : X \to B$ 如本节引言所述。设 $\mathcal{L}$ 是可逆
$\mathcal{O}_X$-模，并设
$P : \mathbf{Z} \to \mathbf{Z}$ 是数值多项式。于是可以考虑既开又闭的
代数子叠
$$
\Cohstack^P_{X/B} =
\Cohstack^{P, \mathcal{L}}_{X/B}
\subset \Cohstack_{X/B}
$$
它由数值不变量与 $P$ 一致、具有真支撑的凝聚层平坦族组成：
$\Cohstack_{X/B}$ 的一个对象
$(T \to B, \mathcal{F})$ 位于 $\Cohstack^P_{X/B}$ 中，当且仅当
$$
P(n) =
\chi(X_t, \mathcal{F}_t \otimes_{\mathcal{O}_{X_t}} \mathcal{L}_t^{\otimes n})
$$
对所有 $n \in \mathbf{Z}$ 和 $t \in T$ 成立。当然，这是情形
\ref{moduli-situation-numerical} 的特例，其中
% P12-ERR-0197: see ../control/ERRATA_CANDIDATES.jsonl
$I = \mathbf{Z} \to D(\mathcal{O}_X)$ 由
$n \mapsto \mathcal{L}^{\otimes n}$ 给出。
由引理 \ref{moduli-lemma-finite-list-perfect-objects}，这是既开又闭子叠。
由于函数
$n \mapsto
\chi(X_t, \mathcal{F}_t \otimes_{\mathcal{O}_{X_t}} \mathcal{L}_t^{\otimes n})$
总是数值多项式（《域上的空间》中的引理
\ref{spaces-over-fields-lemma-numerical-polynomial-from-euler}），
可得不交并分解
$$
\Cohstack_{X/B} = \coprod\nolimits_{P\text{ 为数值多项式}}
\Cohstack^P_{X/B}
$$
\end{example}






\section{Quot 的性质}
\label{moduli-section-quot}

\noindent
设 $f : X \to B$ 是分离且有限表示的代数空间态射，并设
$\mathcal{F}$ 是拟凝聚 $\mathcal{O}_X$-模。则
$\Quotfunctor_{\mathcal{F}/X/B}$ 是代数空间。若
$\mathcal{F}$ 有限表示，则
$\Quotfunctor_{\mathcal{F}/X/B} \to B$ 局部有限表示。
见《Quot》中的命题 \ref{quot-proposition-quot}。

\begin{lemma}
\label{moduli-lemma-quot-diagonal-closed}
$\Quotfunctor_{\mathcal{F}/X/B} \to B$ 的对角态射是闭浸入。
若 $\mathcal{F}$ 有限型，则该对角态射是有限表示的闭浸入。
\end{lemma}

\begin{proof}
设有一个概形 $T/B$，以及两个商
$\mathcal{F}_T \to \mathcal{Q}_i$（$i = 1, 2$），它们对应于
$\Quotfunctor_{\mathcal{F}/X/B}$ 在 $B$ 上的 $T$-值点。
以 $\mathcal{K}_1$ 表示第一个商的核，并以
$u : \mathcal{K}_1 \to \mathcal{Q}_2$ 表示复合。
由《空间上的平坦性》中的引理
\ref{spaces-flat-lemma-F-zero-closed-proper}，$T$ 中存在一个闭子空间，
使得 $T' \to T$ 通过它分解，当且仅当拉回 $u_{T'}$ 为零。
这就证明了对角态射是闭浸入。此外，若 $\mathcal{F}$ 有限型，则
$\mathcal{K}_1$ 有限型
（《位点上的模》中的引理
\ref{sites-modules-lemma-kernel-surjection-finite-onto-finite-presentation}），
并由同一引理可知该对角态射有限表示。
\end{proof}

\begin{lemma}
\label{moduli-lemma-quot-s-lfp}
态射 $\Quotfunctor_{\mathcal{F}/X/B} \to B$ 分离。
若 $\mathcal{F}$ 有限表示，则它还局部有限表示。
\end{lemma}

\begin{proof}
要检验 $\Quotfunctor_{\mathcal{F}/X/B} \to B$ 分离，需要证明其
对角态射为闭浸入。这由引理 \ref{moduli-lemma-quot-diagonal-closed} 得出。
第二个陈述是《Quot》中的命题 \ref{quot-proposition-quot} 的一部分。
\end{proof}

\begin{lemma}
\label{moduli-lemma-quot-existence-part}
假设 $X \to B$ 为真且有限表示，并且 $\mathcal{F}$ 是有限型拟凝聚模。
则 $\Quotfunctor_{\mathcal{F}/X/B} \to B$ 满足赋值判据的存在性部分
（《空间的态射》中的定义
\ref{spaces-morphisms-definition-valuative-criterion}）。
\end{lemma}

\begin{proof}
作基变换后，问题立即约化为下述问题：给定分式域为 $K$ 的赋值环
$R$、一个在 $R$ 上真的代数空间 $X$、一个有限型拟凝聚
$\mathcal{O}_X$-模 $\mathcal{F}$，以及一个凝聚商
$\mathcal{F}_K \to \mathcal{Q}_K$，证明存在一个商
$\mathcal{F} \to \mathcal{Q}$，其中 $\mathcal{Q}$ 是有限表示的
$\mathcal{O}_X$-模，在 $R$ 上平坦，且其泛纤维为 $\mathcal{Q}_K$。
注意，由《空间上的平坦性》中的定理
\ref{spaces-flat-theorem-finite-type-flat}，任意在 $R$ 上平坦的
有限型拟凝聚 $\mathcal{O}_X$-模 $\mathcal{F}$ 都有限表示。
先在仿射局部构造 $\mathcal{Q}$。

\medskip\noindent
仿射局部地，得到下述问题：设 $R \to A$ 是有限表示的环同态，
$M$ 是有限 $A$-模，且 $\varphi : M_K \to N_K$ 是一个 $A_K$-模商。
那么可以考虑
$$
L = \{x \in M \mid \varphi(x \otimes 1) = 0 \}
$$
% P12-ERR-0198: see ../control/ERRATA_CANDIDATES.jsonl
于是 $M \to M/L$ 是一个 $A$-模商，它作为 $R$-模无挠。因此，它作为
$R$-模平坦
（《代数进阶》中的引理
\ref{more-algebra-lemma-valuation-ring-torsion-free-flat}）。
由于 $M$ 作为 $A$-模有限，$L$ 也有限；由此得出
$L$ 作为 $A$-模有限表示（由上面的引用）。显然，$M/L$ 是满足
$(M/L)_K = N_K$ 的唯一这种商。

\medskip\noindent
上一段构造的唯一性保证这些商可以粘合，从而给出所需的
$\mathcal{Q}$。下面补充一些细节。选取一个满 \'etale 态射
$U \to X$，其中 $U$ 是仿射概形。用上述构造得到一个商
$\mathcal{F}|_U \to \mathcal{Q}_U$；它拟凝聚、在 $R$ 上平坦，
并在泛纤维上恢复 $\mathcal{Q}_K|U$。
由于 $X$ 分离，$U \times_X U$ 也是仿射概形，并且 \'etale 于
$X$。于是商
$\mathcal{F}|_{U \times_X U} \to \text{pr}_1^*\mathcal{Q}_U$ 和
$\mathcal{F}|_{U \times_X U} \to \text{pr}_2^*\mathcal{Q}_U$
由该构造的唯一性而一致。
% P12-ERR-0199: see ../control/ERRATA_CANDIDATES.jsonl
因此可把 $\mathcal{F}|_U \to \mathcal{Q}_U$ 下降为所需的满射
$\mathcal{F} \to \mathcal{Q}$
（《空间的性质》中的命题
\ref{spaces-properties-proposition-quasi-coherent}）。
\end{proof}

\begin{lemma}
\label{moduli-lemma-quot-functorial}
设 $B$ 是代数空间。设 $\pi : X \to Y$ 是代数空间之间的仿射拟有限态射，
且这些代数空间在 $B$ 上分离并有限表示。设 $\mathcal{F}$ 是拟凝聚
$\mathcal{O}_X$-模。则 $\pi_*$ 诱导态射
$\Quotfunctor_{\mathcal{F}/X/B} \to \Quotfunctor_{\pi_*\mathcal{F}/Y/B}$。
\end{lemma}

\begin{proof}
令 $\mathcal{G} = \pi_*\mathcal{F}$。由于 $\pi$ 仿射，对于任意
$B$ 上概形 $T$，由《空间的上同调》中的引理
\ref{spaces-cohomology-lemma-affine-base-change} 有
$\mathcal{G}_T = \pi_{T, *}\mathcal{F}_T$。此外，$\pi_T$ 仿射，
所以 $\pi_{T, *}$ 正合，并把商变为商。
% P12-ERR-0200: see ../control/ERRATA_CANDIDATES.jsonl
注意，拟凝聚商 $\mathcal{F}_T \to \mathcal{Q}$ 定义了
$\Quotfunctor_{X/B}$ 的一个点，当且仅当 $\mathcal{Q}$ 定义了
$\Cohstack_{X/B}$ 在 $T$ 上的一个对象（对 $\mathcal{G}$ 和 $Y$
也同样）。我们已在引理 \ref{moduli-lemma-coherent-functorial} 中看到，
$\pi_*$ 诱导态射 $\Cohstack_{X/B} \to \Cohstack_{Y/B}$。
因此，若 $\mathcal{F}_T \to \mathcal{Q}$ 属于
$\Quotfunctor_{\mathcal{F}/X/B}(T)$，则
$\mathcal{G}_T \to \pi_{T, *}\mathcal{Q}$ 属于
$\Quotfunctor_{\mathcal{G}/Y/B}(T)$。
\end{proof}

\begin{lemma}
\label{moduli-lemma-quot-open}
设 $B$ 是代数空间。设 $\pi : X \to Y$ 是代数空间之间的仿射开浸入，
且这些代数空间在 $B$ 上分离并有限表示。设 $\mathcal{F}$ 是拟凝聚
$\mathcal{O}_X$-模。则引理 \ref{moduli-lemma-quot-functorial} 的态射
$\Quotfunctor_{\mathcal{F}/X/B} \to
\Quotfunctor_{\pi_*\mathcal{F}/Y/B}$ 是开浸入。
\end{lemma}

\begin{proof}
略。提示：若 $(\pi_*\mathcal{F})_T \to \mathcal{Q}$ 是
$\Quotfunctor_{\pi_*\mathcal{F}/Y/B}(T)$ 的元素，并且对
$t \in T$ 有 $\text{Supp}(\mathcal{Q}_t) \subset |X_t|$，
那么对 $t$ 的某个邻域中的 $t' \in T$ 也有同一性质。
\end{proof}

\begin{lemma}
\label{moduli-lemma-quot-better-open}
设 $B$ 是代数空间。设 $j : X \to Y$ 是代数空间之间的开浸入，
且这些代数空间在 $B$ 上分离并有限表示。设 $\mathcal{G}$ 是拟凝聚
$\mathcal{O}_Y$-模，并令 $\mathcal{F} = j^*\mathcal{G}$。则存在
$B$ 上代数空间的开浸入
$$
\Quotfunctor_{\mathcal{F}/X/B}
\longrightarrow
\Quotfunctor_{\mathcal{G}/Y/B}
$$
\end{lemma}

\begin{proof}
若 $\mathcal{F}_T \to \mathcal{Q}$ 是
$\Quotfunctor_{\mathcal{F}/X/B}(T)$ 的元素，则可以考虑
$\mathcal{G}_T \to j_{T, *}\mathcal{F}_T \to j_{T, *}\mathcal{Q}$。
考察茎可知它是满射。由引理 \ref{moduli-lemma-coherent-functorial}，
% P12-ERR-0201: see ../control/ERRATA_CANDIDATES.jsonl
$j_{T, *}\mathcal{Q}$ 有限表示、在 $B$ 上平坦，且其支撑在 $B$ 上真。
因此得到 $\Quotfunctor_{\mathcal{G}/Y/B}$ 的一个 $T$-值点。
这就定义了本引理的态射。略去它是开浸入的证明。提示：若
$\mathcal{G}_T \to \mathcal{Q}$ 是
$\Quotfunctor_{\mathcal{G}/Y/B}(T)$ 的元素，并且对 $t \in T$ 有
$\text{Supp}(\mathcal{Q}_t) \subset |X_t|$，那么对 $t$ 的某个
邻域中的 $t' \in T$ 也有同一性质。
\end{proof}

\begin{lemma}
\label{moduli-lemma-quot-closed}
设 $B$ 是代数空间。设 $\pi : X \to Y$ 是代数空间之间的闭浸入，
且这些代数空间在 $B$ 上分离并有限表示。设 $\mathcal{F}$ 是拟凝聚
$\mathcal{O}_X$-模。则引理 \ref{moduli-lemma-quot-functorial} 的态射
$\Quotfunctor_{\mathcal{F}/X/B} \to
\Quotfunctor_{\pi_*\mathcal{F}/Y/B}$ 是同构。
\end{lemma}

\begin{proof}
对于每个 $B$ 上概形 $T$，态射 $\pi_T : X_T \to Y_T$ 是闭浸入。
于是，$\pi_{T, *}$ 给出 $\QCoh(\mathcal{O}_{X_T})$ 与
$\QCoh(\mathcal{O}_{Y_T})$ 的全子范畴之间的范畴等价；后者的对象
是被 $X_T$ 的理想层零化的拟凝聚模，见《空间的态射》中的引理
\ref{spaces-morphisms-lemma-i-star-equivalence}。
% P12-ERR-0202: see ../control/ERRATA_CANDIDATES.jsonl
由于 $(\pi_*\mathcal{F})_T$ 的商被该理想零化，便得到映射
$\Quotfunctor_{\mathcal{F}/X/B}(T) \to
\Quotfunctor_{\pi_*\mathcal{F}/Y/B}(T)$ 对所有 $T$ 都是双射，
这正是所需结论。
\end{proof}

\begin{lemma}
\label{moduli-lemma-quot-quotient}
设 $X \to B$ 如本节引言所述，并设
$\mathcal{F} \to \mathcal{G}$ 是拟凝聚 $\mathcal{O}_X$-模的满射。
则存在典范闭浸入
$\Quotfunctor_{\mathcal{G}/X/B} \to
\Quotfunctor_{\mathcal{F}/X/B}$。
\end{lemma}

\begin{proof}
令 $\mathcal{K} = \Ker(\mathcal{F} \to \mathcal{G})$。由拉回的
右正合性，对所有 $B$ 上概形 $T$，有
% P12-ERR-0203: see ../control/ERRATA_CANDIDATES.jsonl
$\mathcal{K}_T \to \mathcal{F}_T \to \mathcal{G}_T \to 0$
为正合列。特别地，$\mathcal{G}_T$ 的一个商确定
$\mathcal{F}_T$ 的一个商，从而得到函子间的变换
$\Quotfunctor_{\mathcal{G}/X/B} \to
\Quotfunctor_{\mathcal{F}/X/B}$。
由《空间上的平坦性》中的引理
\ref{spaces-flat-lemma-F-zero-closed-proper}，该变换是闭浸入。
% P12-ERR-0204: see ../control/ERRATA_CANDIDATES.jsonl
具体而言，给定 $\Quotfunctor_{\mathcal{F}/X/B}(T)$ 的元素
$\mathcal{F}_T \to \mathcal{Q}$，它向 $T'/T$ 的拉回位于该变换
的像中，当且仅当 $\mathcal{K}_{T'} \to \mathcal{Q}_{T'}$ 为零。
\end{proof}

\begin{remark}[数值不变量]
\label{moduli-remark-quot-numerical}
设 $f : X \to B$ 和 $\mathcal{F}$ 如本节引言所述。设 $I$ 是集合，
且对 $i \in I$，设 $E_i \in D(\mathcal{O}_X)$ 是完美对象。
设 $P : I \to \mathbf{Z}$ 是函数。回忆我们有态射
$$
\Quotfunctor_{\mathcal{F}/X/B} \longrightarrow \Cohstack_{X/B}
$$
它把 $\Quotfunctor_{\mathcal{F}/X/B}(T)$ 的元素
$\mathcal{F}_T \to \mathcal{Q}$ 送到 $\Cohstack_{X/B}$ 在 $T$
上的对象 $\mathcal{Q}$，见《Quot》中的命题
\ref{quot-proposition-quot} 的证明。因此可以作纤维积图
$$
\xymatrix{
\Quotfunctor^P_{\mathcal{F}/X/B} \ar[r] \ar[d] &
\Cohstack^P_{X/B} \ar[d] \\
\Quotfunctor_{\mathcal{F}/X/B} \ar[r] &
\Cohstack_{X/B}
}
$$
这就是左上角代数空间的定义图。左侧竖直箭头是平坦闭浸入；例如，
若 $I$ 有限，或 $B$ 局部 Noether，或 $I = \mathbf{Z}$ 且对某个
可逆 $\mathcal{O}_X$-模 $\mathcal{L}$ 有
$E_i = \mathcal{L}^{\otimes i}$，则它是既开又闭浸入（在最后一种
情形中，有时采用记号
$\Quotfunctor^{P, \mathcal{L}}_{\mathcal{F}/X/B}$）。
见情形 \ref{moduli-situation-numerical}、引理 \ref{moduli-lemma-open-P} 和
\ref{moduli-lemma-finite-list-perfect-objects}，以及例
\ref{moduli-example-hilbert-polynomial}。
\end{remark}

\begin{lemma}
\label{moduli-lemma-quot-tensor-invertible}
设 $f : X \to B$ 和 $\mathcal{F}$ 如本节引言所述。设
$\mathcal{L}$ 是可逆 $\mathcal{O}_X$-模。则与 $\mathcal{L}$
作张量积定义同构
$$
\Quotfunctor_{\mathcal{F}/X/B}
\longrightarrow
\Quotfunctor_{\mathcal{F} \otimes_{\mathcal{O}_X} \mathcal{L}/X/B}
$$
% P12-ERR-0205: see ../control/ERRATA_CANDIDATES.jsonl
给定数值多项式 $P(t)$，令 $P'(t) = P(t + 1)$，则此映射诱导既开又闭
子叠的同构
$\Quotfunctor^P_{\mathcal{F}/X/B}
\longrightarrow
\Quotfunctor^{P'}_{\mathcal{F} \otimes_{\mathcal{O}_X} \mathcal{L}/X/B}$。
\end{lemma}

\begin{proof}
令 $\mathcal{G} = \mathcal{F} \otimes_{\mathcal{O}_X} \mathcal{L}$。
注意
$\mathcal{G}_T = \mathcal{F}_T \otimes_{\mathcal{O}_{X_T}} \mathcal{L}_T$。
若 $\mathcal{F}_T \to \mathcal{Q}$ 是
$\Quotfunctor_{\mathcal{F}/X/B}(T)$ 的元素，就把它送到
$\Quotfunctor_{\mathcal{F} \otimes_{\mathcal{O}_X} \mathcal{L}/X/B}(T)$
的元素
$\mathcal{G}_T \to \mathcal{Q} \otimes_{\mathcal{O}_{X_T}} \mathcal{L}_T$。
这与拉回相容，因而按所需定义了函子间的变换。由于存在明显的逆变换，
它是同构。略去最后一个陈述的证明。
\end{proof}

\begin{lemma}
\label{moduli-lemma-quot-power-invertible}
设 $f : X \to B$ 和 $\mathcal{F}$ 如本节引言所述。设
$\mathcal{L}$ 是可逆 $\mathcal{O}_X$-模。则
$$
\Quotfunctor^{P, \mathcal{L}}_{\mathcal{F}/X/B} =
\Quotfunctor^{P', \mathcal{L}^{\otimes n}}_{\mathcal{F}/X/B}
$$
其中 $P'(t) = P(nt)$。
\end{lemma}

\begin{proof}
展开所有定义后立即得出。
\end{proof}






\section{Quot 的有界性}
\label{moduli-section-quot-bounded}

\noindent
与经典情形不同，我们已经知道 Quot 函子是代数空间，但尚不知道它是否
总能由有限型代数空间表示。

\begin{lemma}
\label{moduli-lemma-quot-Pn}
设 $n \geq 0$、$r \geq 1$、$P \in \mathbf{Q}[t]$。参数化
$\mathcal{O}_{\mathbf{P}^n_\mathbf{Z}}^{\oplus r}$ 的、Hilbert
多项式为 $P$ 的商的代数空间
$$
X = \Quotfunctor^P_{\mathcal{O}^{\oplus r}_{\mathbf{P}^n_\mathbf{Z}}/
\mathbf{P}^n_\mathbf{Z}/\mathbf{Z}}
$$
在 $\Spec(\mathbf{Z})$ 上真。
\end{lemma}

\begin{proof}
我们已经知道 $X \to \Spec(\mathbf{Z})$ 分离且局部有限表示
（引理 \ref{moduli-lemma-quot-s-lfp}）。我们还知道
$X \to \Spec(\mathbf{Z})$ 满足赋值判据的存在性部分，见引理
\ref{moduli-lemma-quot-existence-part}。由真性的赋值判据，只需证明我们的
Quot 空间拟紧，见《空间的态射》中的引理
\ref{spaces-morphisms-lemma-characterize-proper}。因此，只需找到一个
拟紧概形 $T$ 以及满态射 $T \to X$。设 $m$ 为《簇》中的引理
\ref{varieties-lemma-bound-quotients-free} 所给出的整数。令
$$
N = r{m + n \choose n} - P(m)
$$
以下把 $\mathbf{P}^n_\mathbf{Z} =
\text{Proj}(\mathbf{Z}[T_0, \ldots, T_n])$ 简记为
$\mathbf{P}^n$，未标底的积均指在 $\Spec(\mathbf{Z})$ 上的积。
证明的思路是在一个仿射概形 $T$ 上构造一个“万有”映射
$$
\Psi :
\mathcal{O}_{T \times \mathbf{P}^n}(-m)^{\oplus N}
\longrightarrow
\mathcal{O}_{T \times \mathbf{P}^n}^{\oplus r}
$$
并证明 $X$ 的每个点都对应于此映射在 $T$ 的某个点处的余核。

\medskip\noindent
$T$ 和 $\Psi$ 的定义。取 $T = \Spec(A)$，其中
$$
A = \mathbf{Z}[a_{i, j, E}]
$$
这里 $i \in \{1, \ldots, r\}$、$j \in \{1, \ldots, N\}$，
而 $E = (e_0, \ldots, e_n)$ 遍历总次数
% P12-ERR-0206: see ../control/ERRATA_CANDIDATES.jsonl
$|E| = \sum_{k = 0, \ldots n} e_k = m$ 的多重指标。然后把 $\Psi$
定义为一个矩阵，其 $(i, j)$ 元是映射
$$
\sum\nolimits_{E = (e_0, \ldots, e_n)}
a_{i, j, E} T_0^{e_0} \ldots T_n^{e_n} :
\mathcal{O}_{T \times \mathbf{P}^n}(-m)
\longrightarrow
\mathcal{O}_{T \times \mathbf{P}^n}
$$
其中求和遍历上述 $E$（当然，$i$ 和 $j$ 固定）。

\medskip\noindent
考虑 $T \times \mathbf{P}^n$ 上的商
$\mathcal{Q} = \Coker(\Psi)$。由《态射进阶》中的引理
\ref{more-morphisms-lemma-generic-flatness-stratification}，存在
$t \geq 0$ 以及闭子概形
$$
T = T_0 \supset T_1 \supset \ldots \supset T_t = \emptyset
$$
使得 $\mathcal{Q}$ 向
$(T_p \setminus T_{p + 1}) \times \mathbf{P}^n$ 的拉回
$\mathcal{Q}_p$ 在 $T_p \setminus T_{p + 1}$ 上平坦。注意，通过
拉回定义 $\mathcal{Q} = \Coker(\Psi)$ 的正合列，可得正合列
$$
\mathcal{O}_{(T_p \setminus T_{p + 1}) \times \mathbf{P}^n}(-m)^{\oplus N}
\to
\mathcal{O}_{(T_p \setminus T_{p + 1}) \times \mathbf{P}^n}^{\oplus r}
\to
\mathcal{Q}_p
\to
0
$$
因此得到态射
$$
% P12-ERR-0207: see ../control/ERRATA_CANDIDATES.jsonl
\coprod (T_p \setminus T_{p + 1})
\longrightarrow
% P12-ERR-0208: see ../control/ERRATA_CANDIDATES.jsonl
\Quotfunctor_{\mathcal{O}^{\oplus r}/\mathbf{P}/\mathbf{Z}}
\supset
\Quotfunctor^P_{\mathcal{O}^{\oplus r}/\mathbf{P}/\mathbf{Z}} = X
$$
由于左端是 Noether 概形，右端的包含是开的，只需证明 $X$ 的任意点
都在此态射的像中。

\medskip\noindent
设 $k$ 是域，并设 $x \in X(k)$。于是 $x$ 对应于凝聚
$\mathcal{O}_{\mathbf{P}^n_k}$-模的满射
$\mathcal{O}_{\mathbf{P}^n_k}^{\oplus r} \to \mathcal{F}$，
其中 $\mathcal{F}$ 的 Hilbert 多项式为 $P$。考虑短正合列
$$
0 \to \mathcal{K} \to
\mathcal{O}_{\mathbf{P}^n_k}^{\oplus r} \to
\mathcal{F} \to 0
$$
由《簇》中的引理 \ref{varieties-lemma-bound-quotients-free} 以及
$m$ 的选取，$\mathcal{K}$ 是 $m$-正则的。由《簇》中的引理
\ref{varieties-lemma-m-regular-globally-generated}，
$\mathcal{K}(m)$ 由整体截面生成。再由《簇》中的引理
\ref{varieties-lemma-m-regular-up} 和 $m$-正则性的定义，对
$i > 0$ 有
$H^i(\mathbf{P}^n_k, \mathcal{K}(m)) = 0$。因此
$$
\dim_k H^0(\mathbf{P}^n_k, \mathcal{K}(m)) =
\chi(\mathcal{K}(m)) =
\chi(\mathcal{O}_{\mathbf{P}^n_k}(m)^{\oplus r}) -
\chi(\mathcal{F}(m)) = N
$$
其中最后一个等式来自 $N$ 的选取。于是有满射
$$
\mathcal{O}_{\mathbf{P}^n_k}^{\oplus N}
\longrightarrow
\mathcal{K}(m)
$$
反向扭转并使用上述短正合列，可知 $\mathcal{F}$ 是映射
$$
\Psi_x :
\mathcal{O}_{\mathbf{P}^n_k}(-m)^{\oplus N}
\to
\mathcal{O}_{\mathbf{P}^n_k}^{\oplus r}
$$
的余核。存在唯一环同态 $\tau : A \to k$，使得 $\Psi$ 沿相应态射
$t = \Spec(\tau) : \Spec(k) \to T$ 的基变换等于 $\Psi_x$。
这是因为，定义 $\Psi_x$ 的
% P12-ERR-0209: see ../control/ERRATA_CANDIDATES.jsonl
$N \times r$ 矩阵的各项是 $T_0, \ldots, T_n$ 中次数为 $m$ 的
齐次多项式
$\sum \lambda_{i, j, E} T_0^{e_0} \ldots T_n^{e_n}$，其系数
$\lambda_{i, j, E} \in k$；我们可以令
$\tau(a_{i, j, E}) = \lambda_{i, j, E}$。于是对某个 $p$ 有
$t \in T_p \setminus T_{p + 1}$，且 $t$ 在上述态射下的像正是
$x$，如所需。
\end{proof}

\begin{lemma}
\label{moduli-lemma-quot-Pn-over-base}
设 $B$ 是代数空间，并令 $X = B \times \mathbf{P}^n_\mathbf{Z}$。
设 $\mathcal{L}$ 是 $\mathcal{O}_{\mathbf{P}^n}(1)$ 向 $X$ 的拉回。
设 $\mathcal{F}$ 是有限表示的 $\mathcal{O}_X$-模。参数化
$\mathcal{F}$ 的、关于 $\mathcal{L}$ 的 Hilbert 多项式为 $P$
的商的代数空间 $\Quotfunctor^P_{\mathcal{F}/X/B}$ 在 $B$ 上真。
\end{lemma}

\begin{proof}
此问题在 $B$ 上是 \'etale 局部的，见《空间的态射》中的引理
\ref{spaces-morphisms-lemma-proper-local}。因此可以假设 $B$ 是
仿射概形。此时 $\mathcal{L}$ 是 $X$ 上的丰富可逆模
（由《构造》中的引理 \ref{constructions-lemma-ample-on-proj}
以及《性质》中的定义 \ref{properties-definition-ample} 给出的
丰富可逆模之定义）。因此，由《性质》中的命题
\ref{properties-proposition-characterize-ample}，可以找到
$r' \geq 0$ 和 $r \geq 0$ 以及满射
$$
\mathcal{O}_X^{\oplus r} \longrightarrow
\mathcal{F} \otimes_{\mathcal{O}_X} \mathcal{L}^{\otimes r'}
$$
由引理 \ref{moduli-lemma-quot-tensor-invertible}，可以把 $\mathcal{F}$
替换为
$\mathcal{F} \otimes_{\mathcal{O}_X} \mathcal{L}^{\otimes r'}$，
并把 $P(t)$ 替换为 $P(t + r')$。由引理
\ref{moduli-lemma-quot-quotient}，得到闭浸入
$$
\Quotfunctor^P_{\mathcal{F}/X/B}
\longrightarrow
\Quotfunctor^P_{\mathcal{O}_X^{\oplus r}/X/B}
$$
由于我们已在引理 \ref{moduli-lemma-quot-Pn} 中证明
$\Quotfunctor^P_{\mathcal{O}_X^{\oplus r}/X/B} \to B$ 为真，
结论成立。
\end{proof}

\begin{lemma}
\label{moduli-lemma-quot-proper-over-base}
设 $f : X \to B$ 是代数空间之间的有限表示真态射。设
$\mathcal{F}$ 是有限表示的 $\mathcal{O}_X$-模。设
$\mathcal{L}$ 是在 $X/B$ 上丰富的可逆 $\mathcal{O}_X$-模，见
《空间上的除子》中的定义
\ref{spaces-divisors-definition-relatively-ample}。参数化
$\mathcal{F}$ 的、关于 $\mathcal{L}$ 的 Hilbert 多项式为 $P$
的商的代数空间 $\Quotfunctor^P_{\mathcal{F}/X/B}$ 在 $B$ 上真。
\end{lemma}

\begin{proof}
此问题在 $B$ 上是 \'etale 局部的，见《空间的态射》中的引理
\ref{spaces-morphisms-lemma-proper-local}。因此可以假设 $B$ 是
仿射概形。于是可以找到闭浸入 $i : X \to \mathbf{P}^n_B$，
使得对某个 $d \geq 1$ 有
$i^*\mathcal{O}_{\mathbf{P}^n_B}(1) \cong \mathcal{L}^{\otimes d}$。
见《态射》中的引理
\ref{morphisms-lemma-quasi-projective-finite-type-over-S}。
把 $\mathcal{L}$ 换成 $\mathcal{L}^{\otimes d}$，并把数值多项式
$P(t)$ 换成 $P(dt)$，不改变
$\Quotfunctor^P_{\mathcal{F}/X/B}$；略去一些细节。因此可以假设
$\mathcal{L} = i^*\mathcal{O}_{\mathbf{P}^n_B}(1)$。引理
\ref{moduli-lemma-quot-closed} 的同构
$\Quotfunctor_{\mathcal{F}/X/B} \to
\Quotfunctor_{i_*\mathcal{F}/\mathbf{P}^n_B/B}$ 诱导同构
$\Quotfunctor^P_{\mathcal{F}/X/B} \cong
\Quotfunctor^P_{i_*\mathcal{F}/\mathbf{P}^n_B/B}$。
由于由引理 \ref{moduli-lemma-quot-Pn-over-base}，
$\Quotfunctor^P_{i_*\mathcal{F}/\mathbf{P}^n_B/B}$ 在 $B$ 上真，
结论成立。
\end{proof}

\begin{lemma}
\label{moduli-lemma-quot-qc-over-base}
设 $f : X \to B$ 是代数空间之间的有限表示分离态射。设
$\mathcal{F}$ 是有限表示的 $\mathcal{O}_X$-模。设
$\mathcal{L}$ 是在 $X/B$ 上丰富的可逆 $\mathcal{O}_X$-模，见
《空间上的除子》中的定义
\ref{spaces-divisors-definition-relatively-ample}。参数化
$\mathcal{F}$ 的、关于 $\mathcal{L}$ 的 Hilbert 多项式为 $P$
的商的代数空间 $\Quotfunctor^P_{\mathcal{F}/X/B}$ 在 $B$
上分离且有限表示。
\end{lemma}

\begin{proof}
我们已经看到 $\Quotfunctor_{\mathcal{F}/X/B} \to B$ 分离且
局部有限表示，见引理 \ref{moduli-lemma-quot-s-lfp}。因此，只需证明注
\ref{moduli-remark-quot-numerical} 中的开子空间
$\Quotfunctor^P_{\mathcal{F}/X/B}$ 在 $B$ 上拟紧。

\medskip\noindent
此问题在 $B$ 上是 \'etale 局部的
（《空间的态射》中的引理
\ref{spaces-morphisms-lemma-quasi-compact-local}）。因此可以假设
$B$ 仿射。

\medskip\noindent
假设 $B = \Spec(\Lambda)$。把 $\Lambda = \colim \Lambda_i$
写成其有限型 $\mathbf{Z}$-子代数的余极限。于是可以找到一个 $i$，
以及引理中那样的、定义在 $B_i = \Spec(\Lambda_i)$ 上的系统
$X_i, \mathcal{F}_i, \mathcal{L}_i$，其向 $B$ 的基变换给出
$X, \mathcal{F}, \mathcal{L}$。这由《空间的极限》中的下列引理
得到：\ref{spaces-limits-lemma-descend-finite-presentation}
（用于找到 $X_i$）、
\ref{spaces-limits-lemma-descend-modules-finite-presentation}
（用于找到 $\mathcal{F}_i$）、
\ref{spaces-limits-lemma-descend-invertible-modules}
（用于找到 $\mathcal{L}_i$），以及
\ref{spaces-limits-lemma-descend-separated}
（用于使 $X_i$ 分离）。由于
$$
\Quotfunctor_{\mathcal{F}/X/B} = B \times_{B_i}
\Quotfunctor_{\mathcal{F}_i/X_i/B_i}
$$
并且 $\Quotfunctor^P_{\mathcal{F}/X/B}$ 也同样，问题约化为下一段
讨论的情形。

\medskip\noindent
假设 $B$ 仿射且 Noether。由引理
\ref{moduli-lemma-quot-power-invertible}，可以把 $\mathcal{L}$ 替换为
它的一个正幂。因此可以假设存在浸入
$i : X \to \mathbf{P}^n_B$，使得
$i^*\mathcal{O}_{\mathbf{P}^n}(1) = \mathcal{L}$。由《态射》中的
引理 \ref{morphisms-lemma-quasi-compact-immersion}，存在闭子概形
$X' \subset \mathbf{P}^n_B$，使得 $i$ 通过开浸入
$j : X \to X'$ 分解。由《性质》中的引理
\ref{properties-lemma-lift-finite-presentation}，存在有限表示的
$\mathcal{O}_{X'}$-模 $\mathcal{G}$，使得
$j^*\mathcal{G} = \mathcal{F}$。因此，由引理
\ref{moduli-lemma-quot-better-open} 得到开浸入
$$
\Quotfunctor_{\mathcal{F}/X/B}
\longrightarrow
\Quotfunctor_{\mathcal{G}/X'/B}
$$
显然，这个开浸入把 $\Quotfunctor^P_{\mathcal{F}/X/B}$ 送入
$\Quotfunctor^P_{\mathcal{G}/X'/B}$。现在，由引理
\ref{moduli-lemma-quot-proper-over-base}，
$\Quotfunctor^P_{\mathcal{G}/X'/B}$ 在 $B$ 上真。因此它是
Noether 的，而 Noether 代数空间的任意开子空间都拟紧，故结论成立。
\end{proof}




\section{Hilbert 函子的性质}
\label{moduli-section-hilb}

\noindent
设 $f : X \to B$ 是分离且有限表示的代数空间态射。则
$\Hilbfunctor_{X/B}$ 是在 $B$ 上局部有限表示的代数空间。
见《Quot》中的命题 \ref{quot-proposition-hilb}。

\begin{lemma}
\label{moduli-lemma-hilb-diagonal-closed}
$\Hilbfunctor_{X/B} \to B$ 的对角态射是有限表示的闭浸入。
\end{lemma}

\begin{proof}
我们已在《Quot》中的引理 \ref{quot-lemma-hilb-is-quot} 看到
$\Hilbfunctor_{X/B} = \Quotfunctor_{\mathcal{O}_X/X/B}$。
因此结论由引理 \ref{moduli-lemma-quot-diagonal-closed} 得出。
\end{proof}

\begin{lemma}
\label{moduli-lemma-hilb-s-lfp}
态射 $\Hilbfunctor_{X/B} \to B$ 分离且局部有限表示。
\end{lemma}

\begin{proof}
要检验 $\Hilbfunctor_{X/B} \to B$ 分离，需要证明其对角态射为
闭浸入。这由引理 \ref{moduli-lemma-hilb-diagonal-closed} 得出。
第二个陈述是《Quot》中的命题 \ref{quot-proposition-hilb} 的一部分。
\end{proof}

\begin{lemma}
\label{moduli-lemma-hilb-existence-part}
假设 $X \to B$ 为真且有限表示。则
$\Hilbfunctor_{X/B} \to B$ 满足赋值判据的存在性部分
（《空间的态射》中的定义
\ref{spaces-morphisms-definition-valuative-criterion}）。
\end{lemma}

\begin{proof}
我们已在《Quot》中的引理 \ref{quot-lemma-hilb-is-quot} 看到
$\Hilbfunctor_{X/B} = \Quotfunctor_{\mathcal{O}_X/X/B}$。
因此结论由引理 \ref{moduli-lemma-quot-existence-part} 得出。
\end{proof}

\begin{lemma}
\label{moduli-lemma-hilb-open}
设 $B$ 是代数空间。设 $\pi : X \to Y$ 是代数空间之间的开浸入，
且这些代数空间在 $B$ 上分离并有限表示。则 $\pi$ 诱导开浸入
$\Hilbfunctor_{X/B} \to \Hilbfunctor_{Y/B}$。
\end{lemma}

\begin{proof}
略。提示：若 $Z \subset X_T$ 是在 $T$ 上真的闭子概形，则 $Z$
也是 $Y_T$ 中的闭子概形。因此得到变换
$\Hilbfunctor_{X/B} \to \Hilbfunctor_{Y/B}$。若
$Z \subset Y_T$ 是 $\Hilbfunctor_{Y/B}(T)$ 的元素，并且对
$t \in T$ 有 $|Z_t| \subset |X_t|$，那么对 $t$ 的某个邻域中的
$t' \in T$ 也有同一性质。
\end{proof}

\begin{lemma}
\label{moduli-lemma-hilb-closed}
设 $B$ 是代数空间。设 $\pi : X \to Y$ 是代数空间之间的闭浸入，
且这些代数空间在 $B$ 上分离并有限表示。则 $\pi$ 诱导闭浸入
$\Hilbfunctor_{X/B} \to \Hilbfunctor_{Y/B}$。
\end{lemma}

\begin{proof}
由于 $\pi$ 是闭浸入，显然，给定闭子概形 $Z \subset X_T$，
% P12-ERR-0210: see ../control/ERRATA_CANDIDATES.jsonl
可把 $Z$ 看作 $X_T$ 的闭子概形。因此得到变换
$\Hilbfunctor_{X/B} \to \Hilbfunctor_{Y/B}$。立即可见该变换是
单态射。要证明它是闭浸入，可以对映射
$\mathcal{O}_Y \to \mathcal{O}_X$ 使用引理
\ref{moduli-lemma-quot-quotient}，并使用《Quot》中的引理
\ref{quot-lemma-hilb-is-quot} 的等同
$\Hilbfunctor_{X/B} = \Quotfunctor_{\mathcal{O}_X/X/B}$、
$\Hilbfunctor_{Y/B} = \Quotfunctor_{\mathcal{O}_Y/Y/B}$。
\end{proof}

\begin{remark}[数值不变量]
\label{moduli-remark-hilb-numerical}
设 $f : X \to B$ 如本节引言所述。设 $I$ 是集合，且对 $i \in I$，
设 $E_i \in D(\mathcal{O}_X)$ 是完美对象。设
$P : I \to \mathbf{Z}$ 是函数。回忆
$\Hilbfunctor_{X/B} = \Quotfunctor_{\mathcal{O}_X/X/B}$，见
《Quot》中的引理 \ref{quot-lemma-hilb-is-quot}。因此可以定义
$$
\Hilbfunctor^P_{X/B} = \Quotfunctor^P_{\mathcal{O}_X/X/B}
$$
其中 $\Quotfunctor^P_{\mathcal{O}_X/X/B}$ 如注
\ref{moduli-remark-quot-numerical} 所述。态射
$$
\Hilbfunctor^P_{X/B} \longrightarrow \Hilbfunctor_{X/B}
$$
是平坦闭浸入；例如，若 $I$ 有限，或 $B$ 局部 Noether，或
$I = \mathbf{Z}$ 且对某个可逆 $\mathcal{O}_X$-模 $\mathcal{L}$
有 $E_i = \mathcal{L}^{\otimes i}$，则它是既开又闭浸入。
在最后一种情形中，有时采用记号
$\Hilbfunctor^{P, \mathcal{L}}_{X/B}$。
\end{remark}

\begin{lemma}
\label{moduli-lemma-hilb-proper-over-base}
设 $f : X \to B$ 是代数空间之间的有限表示真态射。设
$\mathcal{L}$ 是在 $X/B$ 上丰富的可逆 $\mathcal{O}_X$-模，见
《空间上的除子》中的定义
\ref{spaces-divisors-definition-relatively-ample}。参数化关于
$\mathcal{L}$ 的 Hilbert 多项式为 $P$ 的闭子概形的代数空间
$\Hilbfunctor^P_{X/B}$ 在 $B$ 上真。
\end{lemma}

\begin{proof}
回忆 $\Hilbfunctor_{X/B} =
\Quotfunctor_{\mathcal{O}_X/X/B}$，见《Quot》中的引理
\ref{quot-lemma-hilb-is-quot}。因此本引理立即由引理
\ref{moduli-lemma-quot-proper-over-base} 得出。
\end{proof}

\begin{lemma}
\label{moduli-lemma-hilb-qc-over-base}
设 $f : X \to B$ 是代数空间之间的有限表示分离态射。设
$\mathcal{L}$ 是在 $X/B$ 上丰富的可逆 $\mathcal{O}_X$-模，见
《空间上的除子》中的定义
\ref{spaces-divisors-definition-relatively-ample}。参数化关于
$\mathcal{L}$ 的 Hilbert 多项式为 $P$ 的闭子概形的代数空间
$\Hilbfunctor^P_{X/B}$ 在 $B$ 上分离且有限表示。
\end{lemma}

\begin{proof}
回忆 $\Hilbfunctor_{X/B} =
\Quotfunctor_{\mathcal{O}_X/X/B}$，见《Quot》中的引理
\ref{quot-lemma-hilb-is-quot}。因此本引理立即由引理
\ref{moduli-lemma-quot-qc-over-base} 得出。
\end{proof}









\section{Picard 叠的性质}
\label{moduli-section-picard-stack}

\noindent
设 $f : X \to B$ 是平坦、真且有限表示的代数空间态射。则参数化
$X/B$ 上可逆层的叠 $\Picardstack_{X/B}$ 是代数的，见《Quot》中的
命题 \ref{quot-proposition-pic}。

\begin{lemma}
\label{moduli-lemma-pic-diagonal-affine-fp}
$\Picardstack_{X/B}$ 在 $B$ 上的对角态射仿射且有限表示。
\end{lemma}

\begin{proof}
我们已在《Quot》中的引理
\ref{quot-lemma-picard-stack-open-in-coh} 看到，
$\Picardstack_{X/B}$ 是 $\Cohstack_{X/B}$ 的开子叠。
因此结论由引理 \ref{moduli-lemma-coherent-diagonal-affine-fp} 得出。
\end{proof}

\begin{lemma}
\label{moduli-lemma-pic-qs-lfp}
态射 $\Picardstack_{X/B} \to B$ 拟分离且局部有限表示。
\end{lemma}

\begin{proof}
我们已在《Quot》中的引理
\ref{quot-lemma-picard-stack-open-in-coh} 看到，
$\Picardstack_{X/B}$ 是 $\Cohstack_{X/B}$ 的开子叠。
因此结论由引理 \ref{moduli-lemma-coherent-qs-lfp} 得出。
\end{proof}

\begin{lemma}
\label{moduli-lemma-pic-existence-part}
除为真之外，再假设 $X \to B$ 光滑。则
$\Picardstack_{X/B} \to B$ 满足赋值判据的存在性部分
（《叠的态射》中的定义
\ref{stacks-morphisms-definition-existence}）。
\end{lemma}

\begin{proof}
作基变换后，问题立即约化为下述问题：给定分式域为 $K$ 的赋值环
$R$、一个在 $R$ 上真且光滑的代数空间 $X$，以及一个可逆
$\mathcal{O}_{X_K}$-模 $\mathcal{L}_K$，证明存在一个可逆
$\mathcal{O}_X$-模 $\mathcal{L}$，其泛纤维为 $\mathcal{L}_K$。
注意，$X_K$ 是 Noether、分离且正则的
（使用《空间的态射》中的引理
\ref{spaces-morphisms-lemma-finite-presentation-noetherian}
和《域上的空间》中的引理
\ref{spaces-over-fields-lemma-smooth-regular}）。因此，对
$X_K$ 中的两个有效 Cartier 除子 $D_K, D'_K$，可以在 Picard 群中把
$\mathcal{L}_K$ 写成 $\mathcal{O}_{X_K}(D_K)$ 与
$\mathcal{O}_{X_K}(D'_K)$ 之差，见《空间上的除子》中的引理
\ref{spaces-divisors-lemma-Noetherian-regular-separated-pic-effective-Cartier}。
最后，我们知道 $D_K$ 和 $D'_K$ 分别是有效 Cartier 除子
$D, D' \subset X$ 的限制，见《空间上的除子》中的引理
\ref{spaces-divisors-lemma-smooth-over-valuation-ring-effective-Cartier}。
\end{proof}

\begin{lemma}
\label{moduli-lemma-pic-inertia}
假设对所有 $B$ 上概形 $T$ 都有
$f_{T, *}\mathcal{O}_{X_T} \cong \mathcal{O}_T$。则
$\Picardstack_{X/B}$ 的惯性叠等于
$\mathbf{G}_m \times \Picardstack_{X/B}$。
\end{lemma}

\begin{proof}
这在《叠的例子》中的例
\ref{examples-stacks-example-inertia-stack-of-picard} 中已有说明。
\end{proof}

\begin{lemma}
\label{moduli-lemma-pic-curves-smooth}
除本节的其他假设外，再假设 $f : X \to B$ 的相对维数
$\leq 1$。则 $\Picardstack_{X/B} \to B$ 光滑。
\end{lemma}

\begin{proof}
由引理 \ref{moduli-lemma-pic-qs-lfp}，我们已知
$\Picardstack_{X/B} \to B$ 局部有限表示。因此，由《叠的态射进阶》
中的引理 \ref{stacks-more-morphisms-lemma-smooth-formally-smooth}，
只需证明 $\Picardstack_{X/B} \to B$ 形式光滑。
作基变换后，问题立即约化为下述问题：给定仿射概形的一阶加厚
$T \subset T'$，给定真、平坦、有限表示且相对维数 $\leq 1$ 的
$X' \to T'$，并对 $X = T \times_{T'} X'$ 给定一个可逆
$\mathcal{O}_X$-模 $\mathcal{L}$，证明存在一个可逆
$\mathcal{O}_{X'}$-模 $\mathcal{L}'$，其在 $X$ 上的限制为
$\mathcal{L}$。由于 $T \subset T'$ 是一阶加厚，
$X \subset X'$ 也同样，见《空间的态射进阶》中的引理
\ref{spaces-more-morphisms-lemma-base-change-thickening}。
由《空间的态射进阶》中的引理
\ref{spaces-more-morphisms-lemma-picard-group-first-order-thickening}，
只需证明 $H^2(X, \mathcal{I}) = 0$，其中 $\mathcal{I}$ 是在
$X'$ 中截出 $X$ 的拟凝聚理想。
% P12-ERR-0211: see ../control/ERRATA_CANDIDATES.jsonl
以 $f : X \to T$ 表示结构态射。由《空间的上同调》中的引理
\ref{spaces-cohomology-lemma-higher-direct-images-zero-above-dimension-fibre}，
对 $p > 1$ 有 $R^pf_*\mathcal{I} = 0$。因此，由《空间的上同调》
中的引理
\ref{spaces-cohomology-lemma-quasi-coherence-higher-direct-images-application}
得到所需消失（这里最终用到了 $T$ 仿射）。
\end{proof}




\section{Picard 函子的性质}
\label{moduli-section-picard-functor}

\noindent
设 $f : X \to B$ 是平坦、真且有限表示的代数空间态射，并且还假设
对每个 $T/B$，典范映射
$$
\mathcal{O}_T \longrightarrow f_{T, *}\mathcal{O}_{X_T}
$$
都是同构。则 Picard 函子 $\Picardfunctor_{X/B}$ 是代数空间，
见《Quot》中的命题 \ref{quot-proposition-pic-functor}。
% P12-ERR-0212: see ../control/ERRATA_CANDIDATES.jsonl
它与 Picard 叠有紧密关系。

\begin{lemma}
\label{moduli-lemma-pic-gerbe-over-pic-functor}
态射 $\Picardstack_{X/B} \to \Picardfunctor_{X/B}$ 使 Picard 叠
成为 Picard 函子上的 gerbe。
\end{lemma}

\begin{proof}
$\Picardstack_{X/B} \to \Picardfunctor_{X/B}$ 为 gerbe 的定义
见《叠的态射》中的定义
\ref{stacks-morphisms-definition-gerbe}；该定义又引用《叠》中的定义
\ref{stacks-definition-gerbe-over-stack-in-groupoids}。为证明此事，
我们检验《叠》中的引理 \ref{stacks-lemma-when-gerbe} 的条件
(2)(a) 和 (2)(b)。这立即由《Quot》中的引理
\ref{quot-lemma-pic-over-pic} 得出；下面给出详细说明。

\medskip\noindent
条件 (2)(a)。设对某个 $B$ 上概形 $U$ 有
$\xi \in \Picardfunctor_{X/B}(U)$。由于
$\Picardfunctor_{X/B}$ 是 $B$ 上概形上的规则
$T \mapsto \Pic(X_T)$ 的 fppf 层化
（《Quot》中的情形 \ref{quot-situation-pic}），存在一个 fppf 覆盖
$\{U_i \to U\}$，使得 $\xi|_{U_i}$ 对应于 $X_{U_i}$ 上的某个
可逆模 $\mathcal{L}_i$。于是
$(U_i \to B, \mathcal{L}_i)$ 是 $\Picardstack_{X/B}$ 在 $U_i$
上的对象，并映到 $\xi|_{U_i}$。

\medskip\noindent
条件 (2)(b)。设 $U$ 是 $B$ 上概形，并设 $\mathcal{L}, \mathcal{N}$
是 $X_U$ 上的可逆模，它们映到 $\Picardfunctor_{X/B}(U)$ 的同一
元素。则存在 fppf 覆盖 $\{U_i \to U\}$，使得
$\mathcal{L}|_{X_{U_i}}$ 同构于
$\mathcal{N}|_{X_{U_i}}$。因此按所需找到同构
$(U \to B, \mathcal{L})|_{U_i} \to
(U \to B, \mathcal{N})|_{U_i}$。
\end{proof}

\begin{lemma}
\label{moduli-lemma-pic-functor-diagonal-qc-immersion}
$\Picardfunctor_{X/B}$ 在 $B$ 上的对角态射是拟紧浸入。
\end{lemma}

\begin{proof}
由《Quot》中的引理 \ref{quot-lemma-diagonal-pic}，对角态射是浸入。
为完成证明，只需说明它拟紧。由引理
\ref{moduli-lemma-pic-diagonal-affine-fp}，
$\Picardstack_{X/B}$ 的对角态射拟紧；由引理
\ref{moduli-lemma-pic-gerbe-over-pic-functor}，
$\Picardstack_{X/B}$ 是 $\Picardfunctor_{X/B}$ 上的 gerbe。
结论由《叠的态射》中的引理
\ref{stacks-morphisms-lemma-gerbe-diagonal-quasi-compact} 得出。
\end{proof}

\begin{lemma}
\label{moduli-lemma-pic-functor-qs-lfp}
态射 $\Picardfunctor_{X/B} \to B$ 拟分离且局部有限表示。
\end{lemma}

\begin{proof}
要检验 $\Picardfunctor_{X/B} \to B$ 拟分离，需要证明其对角态射
拟紧。这立即由引理
\ref{moduli-lemma-pic-functor-diagonal-qc-immersion} 得出。
态射 $\Picardstack_{X/B} \to \Picardfunctor_{X/B}$ 是满、平坦
且局部有限表示的
（由引理 \ref{moduli-lemma-pic-gerbe-over-pic-functor} 和《叠的态射》中的
引理 \ref{stacks-morphisms-lemma-gerbe-fppf}）。因此，由《叠的态射》
中的引理
\ref{stacks-morphisms-lemma-flat-finite-presentation-permanence}，
只需证明 $\Picardstack_{X/B} \to B$ 局部有限表示。
这由引理 \ref{moduli-lemma-pic-qs-lfp} 得出。
\end{proof}

\begin{lemma}
\label{moduli-lemma-pic-functor-uniqueness-part}
除本节的其他假设外，再假设 $X \to B$ 的几何纤维整。则
$\Picardfunctor_{X/B} \to B$ 分离。
\end{lemma}

\begin{proof}
由于 $\Picardfunctor_{X/B} \to B$ 拟分离，只需检验赋值判据的
唯一性部分，见《空间的态射》中的引理
\ref{spaces-morphisms-lemma-valuative-criterion-separatedness}。
这立即约化为下述问题：给定
\begin{enumerate}
\item 一个分式域为 $K$ 的赋值环 $R$；
% P12-ERR-0213: see ../control/ERRATA_CANDIDATES.jsonl
\item 一个在 $R$ 上真且平坦、具有整几何纤维的代数空间 $X$；
\item 一个元素 $a \in \Picardfunctor_{X/R}(R)$，满足
$a|_{\Spec(K)} = 0$，
\end{enumerate}
需要证明 $a = 0$。把《叠的态射》中的引理
\ref{stacks-morphisms-lemma-lift-valuation-ring-through-flat-morphism}
应用于满平坦态射
$\Picardstack_{X/R} \to \Picardfunctor_{X/R}$
（由引理 \ref{moduli-lemma-pic-gerbe-over-pic-functor} 和《叠的态射》中的
引理 \ref{stacks-morphisms-lemma-gerbe-fppf}，它是满且平坦的），
在把 $R$ 替换为一个扩张后，可以假设 $a$ 由一个可逆
$\mathcal{O}_X$-模 $\mathcal{L}$ 给出。由于
$a|_{\Spec(K)} = 0$，由《Quot》中的引理
\ref{quot-lemma-flat-geometrically-connected-fibres} 有
$\mathcal{L}_K \cong \mathcal{O}_{X_K}$。

\medskip\noindent
% P12-ERR-0214: see ../control/ERRATA_CANDIDATES.jsonl
以 $f : X \to \Spec(R)$ 表示结构态射。设
$\eta, 0 \in \Spec(R)$ 分别为泛点与闭点。考虑
$\Spec(R)$ 上的完美复形
$K = Rf_*\mathcal{L}$ 和
$M = Rf_*(\mathcal{L}^{\otimes -1})$，见《空间的导出范畴》中的
引理
\ref{spaces-perfect-lemma-flat-proper-perfect-direct-image-general}。
考虑《空间的导出范畴》中的引理
\ref{spaces-perfect-lemma-jump-loci} 与 $K$ 和 $M$ 关联的函数
$\beta_{K, i}, \beta_{M, i} : \Spec(R) \to \mathbf{Z}$。
% P12-ERR-0215: see ../control/ERRATA_CANDIDATES.jsonl
由于 $K$ 与 $M$ 的形成与基变换交换（见上述引理），由《域上的空间》
中的引理
\ref{spaces-over-fields-lemma-proper-geometrically-reduced-global-sections}
以及关于 $f$ 的纤维的假设，有
$\beta_{K, 0}(\eta) = \beta_{M, 0}(\eta) = 1$。
由上半连续性，
$\beta_{K, 0}(0) \geq 1$ 且 $\beta_{M, 0}(0) \geq 1$。
由《域上的空间》中的引理
\ref{spaces-over-fields-lemma-characterize-trivial-pic-integral}，
可知 $\mathcal{L}$ 在特殊纤维 $X_0$ 上的限制平凡。于是如上得到
% P12-ERR-0216: see ../control/ERRATA_CANDIDATES.jsonl
$\beta_{K, 0}(0) = \beta_{M, 0} = 1$。
然后由《代数进阶》中的引理
\ref{more-algebra-lemma-lift-pseudo-coherent-from-residue-field}，
可以用下列形式的复形表示 $K$：
$$
\ldots \to 0 \to R \to R^{\oplus \beta_{K, 1}(0)} \to
R^{\oplus \beta_{K, 2}(0)} \to \ldots
$$
现在 $R \to R^{\oplus \beta_{K, 1}(0)}$ 为零，因为
$\beta_{K, 0}(\eta) = 1$。换言之，在 $D(R)$ 中
$K = R \oplus \tau_{\geq 1}(K)$，其中对某个
$b \in \mathbf{Z}$，$\tau_{\geq 1}(K)$ 的 Tor 振幅位于
$[1, b]$。因此存在整体截面 $s \in H^0(X, \mathcal{L})$，
其在 $X_0$ 上的限制 $s_0$ 处处非零
（仍因为 $K$ 的形成与基变换交换）。于是
$s : \mathcal{O}_X \to \mathcal{L}$ 是可逆层之间的映射，其在
$X_0$ 上的限制是同构，因而它本身也是同构，如所需。
\end{proof}

\begin{lemma}
\label{moduli-lemma-pic-functor-curves-smooth}
除本节的其他假设外，再假设 $f : X \to B$ 的相对维数
$\leq 1$。则 $\Picardfunctor_{X/B} \to B$ 光滑。
\end{lemma}

\begin{proof}
由引理 \ref{moduli-lemma-pic-curves-smooth}，我们知道
$\Picardstack_{X/B} \to B$ 光滑。结合引理
\ref{moduli-lemma-pic-gerbe-over-pic-functor} 与《叠的态射》中的引理
\ref{stacks-morphisms-lemma-gerbe-smooth}，态射
$\Picardstack_{X/B} \to \Picardfunctor_{X/B}$ 满且光滑。
因此，若 $U$ 是概形且
$U \to \Picardstack_{X/B}$ 满且光滑，则
$U \to \Picardfunctor_{X/B}$ 满且光滑，而 $U \to B$ 也满且
光滑（因为这些性质在复合下保持）。所以，例如由《空间上的下降》
中的引理
\ref{spaces-descent-lemma-syntomic-smooth-etale-permanence}，
$\Picardfunctor_{X/B} \to B$ 光滑。
\end{proof}






\section{相对态射的性质}
\label{moduli-section-relative-morphisms}

\noindent
设 $B$ 是代数空间。设 $X$ 和 $Y$ 是 $B$ 上代数空间，其中
$Y \to B$ 平坦、真且有限表示，而 $X \to B$ 分离且有限表示。
则相对态射的函子 $\mathit{Mor}_B(Y, X)$ 是在 $B$ 上局部有限表示
的代数空间。见《Quot》中的命题 \ref{quot-proposition-Mor}。

\begin{lemma}
\label{moduli-lemma-Mor-diagonal-closed}
$\mathit{Mor}_B(Y, X) \to B$ 的对角态射是有限表示的闭浸入。
\end{lemma}

\begin{proof}
存在开浸入
$\mathit{Mor}_B(Y, X) \to \Hilbfunctor_{Y \times_B X/B}$，见
《Quot》中的引理 \ref{quot-lemma-Mor-into-Hilb-open}。
因此本引理由引理 \ref{moduli-lemma-hilb-diagonal-closed} 得出。
\end{proof}

\begin{lemma}
\label{moduli-lemma-Mor-s-lfp}
态射 $\mathit{Mor}_B(Y, X) \to B$ 分离且局部有限表示。
\end{lemma}

\begin{proof}
要检验 $\mathit{Mor}_B(Y, X) \to B$ 分离，需要证明其对角态射为
闭浸入。这由引理 \ref{moduli-lemma-Mor-diagonal-closed} 得出。
第二个陈述是《Quot》中的命题 \ref{quot-proposition-Mor} 的一部分。
\end{proof}

\begin{lemma}
\label{moduli-lemma-Isom-in-Mor}
设 $B, X, Y$ 如本节引言所述，并额外假设 $X \to B$ 为真。
则由同构组成的子函子
$\mathit{Isom}_B(Y, X) \subset \mathit{Mor}_B(Y, X)$
是开子空间。
\end{lemma}

\begin{proof}
这立即由《空间的态射进阶》中的引理
\ref{spaces-more-morphisms-lemma-where-isomorphism} 得出。
\end{proof}

\begin{remark}[数值不变量]
\label{moduli-remark-Mor-numerical}
设 $B, X, Y$ 如本节引言所述。设 $I$ 是集合，且对 $i \in I$，
设 $E_i \in D(\mathcal{O}_{Y \times_B X})$ 是完美对象。
设 $P : I \to \mathbf{Z}$ 是函数。回忆
$$
\mathit{Mor}_B(Y, X) \subset
\Hilbfunctor_{Y \times_B X/B}
$$
是开子空间，见《Quot》中的引理
\ref{quot-lemma-Mor-into-Hilb-open}。因此可以定义
$$
\mathit{Mor}^P_B(Y, X) =
\mathit{Mor}_B(Y, X) \cap \Hilbfunctor^P_{Y \times_B X/B}
$$
其中 $\Hilbfunctor^P_{Y \times_B X/B}$ 如注
\ref{moduli-remark-hilb-numerical} 所述。态射
$$
\mathit{Mor}^P_B(Y, X) \longrightarrow \mathit{Mor}_B(Y, X)
$$
是平坦闭浸入；例如，若 $I$ 有限，或 $B$ 局部 Noether，或
$I = \mathbf{Z}$ 且对某个可逆
$\mathcal{O}_{Y \times_B X}$-模 $\mathcal{L}$ 有
$E_i = \mathcal{L}^{\otimes i}$，则它是既开又闭浸入。
在最后一种情形中，有时采用记号
$\mathit{Mor}^{P, \mathcal{L}}_B(Y, X)$。
\end{remark}

\begin{lemma}
\label{moduli-lemma-Mor-qc-over-base}
设 $B, X, Y$ 如本节引言所述，并设 $\mathcal{L}$ 在 $X/B$ 上丰富，
$\mathcal{N}$ 在 $Y/B$ 上丰富。见《空间上的除子》中的定义
\ref{spaces-divisors-definition-relatively-ample}。
设 $P$ 是数值多项式。则
$$
\mathit{Mor}^{P, \mathcal{M}}_B(Y, X) \longrightarrow B
$$
分离且有限表示，其中
$\mathcal{M} = \text{pr}_1^*\mathcal{N}
\otimes_{\mathcal{O}_{Y \times_B X}} \text{pr}_2^*\mathcal{L}$。
\end{lemma}

\begin{proof}
由引理 \ref{moduli-lemma-Mor-s-lfp}，态射
$\mathit{Mor}_B(Y, X) \to B$ 分离且局部有限表示。因此只需证明
注 \ref{moduli-remark-Mor-numerical} 的既开又闭子空间
$\mathit{Mor}^{P, \mathcal{M}}_B(Y, X)$ 在 $B$ 上拟紧。

\medskip\noindent
此问题在 $B$ 上是 \'etale 局部的
（《空间的态射》中的引理
\ref{spaces-morphisms-lemma-quasi-compact-local}）。
因此可以假设 $B$ 仿射。

\medskip\noindent
假设 $B = \Spec(\Lambda)$。注意，$X$ 和 $Y$ 是概形，而
$\mathcal{L}$ 和 $\mathcal{N}$ 分别是 $X$ 和 $Y$ 上的丰富可逆层
（这立即由定义得出）。把 $\Lambda = \colim \Lambda_i$ 写成其
有限型 $\mathbf{Z}$-子代数的余极限。于是可以找到一个 $i$，以及
引理中那样的、定义在 $B_i = \Spec(\Lambda_i)$ 上的系统
$X_i, Y_i, \mathcal{L}_i, \mathcal{N}_i$，其向 $B$ 的基变换给出
$X, Y, \mathcal{L}, \mathcal{N}$。这由《极限》中的下列引理得到：
\ref{limits-lemma-descend-finite-presentation}
（用于找到 $X_i$、$Y_i$）、
\ref{limits-lemma-descend-invertible-modules}
（用于找到 $\mathcal{L}_i$、$\mathcal{N}_i$）、
\ref{limits-lemma-descend-separated-finite-presentation}
（用于使 $X_i \to B_i$ 分离）、
\ref{limits-lemma-eventually-proper}
（用于使 $Y_i \to B_i$ 为真），以及
\ref{limits-lemma-limit-ample}
（用于使 $\mathcal{L}_i$、$\mathcal{N}_i$ 丰富）。由于
$$
\mathit{Mor}_B(Y, X) = B \times_{B_i} \mathit{Mor}_{B_i}(Y_i, X_i)
$$
% P12-ERR-0217: see ../control/ERRATA_CANDIDATES.jsonl
并且 $\mathit{Mor}^P_B(Y, X)$ 也同样，问题约化为下一段讨论的情形。

\medskip\noindent
假设 $B$ 是 Noether 仿射概形。由《性质》中的引理
\ref{properties-lemma-ample-on-product}，$\mathcal{M}$ 丰富。
由引理 \ref{moduli-lemma-hilb-qc-over-base}，
$\Hilbfunctor^{P, \mathcal{M}}_{Y \times_B X/B}$ 在 $B$ 上
有限表示，因而是 Noether 的。按构造，
$$
\mathit{Mor}^{P, \mathcal{M}}_B(Y, X) =
\mathit{Mor}_B(Y, X) \cap
\Hilbfunctor^{P, \mathcal{M}}_{Y \times_B X/B}
$$
是 $\Hilbfunctor^{P, \mathcal{M}}_{Y \times_B X/B}$ 的开子空间，
因而拟紧（因为 Noether 代数空间的开子空间拟紧）。
\end{proof}






\section{极化真概形叠的性质}
\label{moduli-section-polarized}

\noindent
本节讨论模叠
$$
\Polarizedstack \longrightarrow \Spec(\mathbf{Z})
$$
的性质；它在概形 $S$ 上的截面范畴，是由配备相对丰富可逆层的、
% P12-ERR-0218: see ../control/ERRATA_CANDIDATES.jsonl
在 $S$ 上真、平坦且有限表示的概形组成的范畴。由《Quot》中的定理
\ref{quot-theorem-polarized-algebraic}，这是代数叠。

\begin{lemma}
\label{moduli-lemma-polarized-diagonal-separated-fp}
$\Polarizedstack$ 的对角态射分离且有限表示。
\end{lemma}

\begin{proof}
回忆 $\Polarizedstack$ 是保持极限的代数叠，见《Quot》中的引理
\ref{quot-lemma-polarized-limits}。由《叠的极限》中的引理
\ref{stacks-limits-lemma-limit-preserving-diagonal}，这蕴含
$\Delta : \Polarizedstack \to
\Polarizedstack \times \Polarizedstack$ 保持极限。因此，由
《叠的极限》中的命题
\ref{stacks-limits-proposition-characterize-locally-finite-presentation}，
$\Delta$ 局部有限表示。

\medskip\noindent
下面证明 $\Delta$ 分离。为此，只需证明：给定仿射概形 $U$ 以及
$\Polarizedstack$ 在 $U$ 上的两个对象
$\upsilon = (Y, \mathcal{N})$ 和
$\chi = (X, \mathcal{L})$，代数空间
$$
\mathit{Isom}_{\Polarizedstack}(\upsilon, \chi)
$$
分离。把同构 $\upsilon_T \to \chi_T$ 送到其底层同构
$Y_T \to X_T$ 的规则定义态射
$$
\mathit{Isom}_{\Polarizedstack}(\upsilon, \chi)
\longrightarrow
\mathit{Isom}_U(Y, X)
$$
由引理 \ref{moduli-lemma-Mor-s-lfp} 和 \ref{moduli-lemma-Isom-in-Mor}，
我们已经看到其目标是分离代数空间，所以只需证明此态射分离。
% P12-ERR-0219: see ../control/ERRATA_CANDIDATES.jsonl
给定某个概形 $T/U$ 上的同构 $f : Y_T \to X_T$，显然有
$$
\mathit{Isom}_{\Polarizedstack}(\upsilon, \chi)
\times_{\mathit{Isom}_U(Y, X), [f]} T
=
\mathit{Isom}(\mathcal{N}_T, f^*\mathcal{L}_T)
$$
这里 $[f] : T \to \mathit{Isom}_U(Y, X)$ 表示与 $f$ 对应的
$T$-值点，而
$\mathit{Isom}(\mathcal{N}_T, f^*\mathcal{L}_T)$ 是章节
\ref{moduli-section-hom-isom} 中讨论的代数空间。
% P12-ERR-0220: see ../control/ERRATA_CANDIDATES.jsonl
由于该代数空间在 $U$ 上仿射，所断言的结论蕴含 $\Delta$ 分离。

\medskip\noindent
最后证明 $\Delta$ 拟紧。由于 $\Delta$ 可由代数空间表示，
% P12-ERR-0221: see ../control/ERRATA_CANDIDATES.jsonl
只需检验 $\Delta$ 沿一个满光滑态射
$U \to \Polarizedstack \times \Polarizedstack$ 的基变换拟紧
（例如见《叠的性质》中的引理
\ref{stacks-properties-lemma-check-property-covering}）。
可以假设 $U = \coprod U_i$ 是仿射开集的不交并。
由于 $\Polarizedstack$ 保持极限（见上文），
$\Polarizedstack \to \Spec(\mathbf{Z})$ 局部有限表示，因而
$U_i \to \Spec(\mathbf{Z})$ 局部有限表示
（《叠的极限》中的命题
\ref{stacks-limits-proposition-characterize-locally-finite-presentation}
以及《叠的态射》中的引理
\ref{stacks-morphisms-lemma-composition-finite-presentation} 和
\ref{stacks-morphisms-lemma-smooth-locally-finite-presentation}）。
特别地，$U_i$ 是 Noether 仿射概形。于是问题约化为下一段讨论的情形。

\medskip\noindent
在本段中，给定 Noether 仿射概形 $U$ 以及
$\Polarizedstack$ 在 $U$ 上的两个对象
$\upsilon = (Y, \mathcal{N})$ 和
$\chi = (X, \mathcal{L})$，我们证明代数空间
$$
\mathit{Isom}_{\Polarizedstack}(\upsilon, \chi)
$$
拟紧。由于 $U$ 的连通分支既开又闭，可以用这些分支替换 $U$。
因此可以并且确实假设 $U$ 连通。设 $u \in U$ 是一个点。设 $P$
为 Hilbert 多项式
$n \mapsto \chi(Y_u, \mathcal{N}_u^{\otimes n})$，见《簇》中的
引理 \ref{varieties-lemma-numerical-polynomial-from-euler}。
由于 $U$ 连通，并且函数
$u \mapsto \chi(Y_u, \mathcal{N}_u^{\otimes n})$ 局部常值
（见《概形的导出范畴》中的引理
\ref{perfect-lemma-chi-locally-constant-geometric}），可知在 $U$ 的
每个点都得到同一 Hilbert 多项式。在 $Y \times_U X$ 上令
$\mathcal{M} = \text{pr}_1^*\mathcal{N}
\otimes_{\mathcal{O}_{Y \times_U X}} \text{pr}_2^*\mathcal{L}$。
% P12-ERR-0222: see ../control/ERRATA_CANDIDATES.jsonl
给定某个 $U$ 上概形 $T$ 的元素
$(f, \varphi) \in
\mathit{Isom}_{\Polarizedstack}(\upsilon, \chi)(T)$，则对每个
$t \in T$ 都有
$$
\chi(Y_t, (\text{id} \times f)^*\mathcal{M}^{\otimes n}) =
\chi(Y_t,
\mathcal{N}_t^{\otimes n} \otimes_{\mathcal{O}_{Y_t}}
f_t^*\mathcal{L}_t^{\otimes n}) =
\chi(Y_t, \mathcal{N}_t^{\otimes 2n}) = P(2n)
$$
其中在中间的等式中使用了同构
$\varphi : f^*\mathcal{L}_T \to \mathcal{N}_T$。
令 $P'(t) = P(2t)$，可知上述态射
$$
\mathit{Isom}_{\Polarizedstack}(\upsilon, \chi)
\longrightarrow
\mathit{Isom}_U(Y, X)
$$
的像包含在交
$$
\mathit{Isom}_U(Y, X) \cap \mathit{Mor}^{P', \mathcal{M}}_U(Y, X)
$$
中。该交是 $\mathit{Mor}_U(Y, X)$ 的两个开子空间之交
（见引理 \ref{moduli-lemma-Isom-in-Mor} 和注
\ref{moduli-remark-Mor-numerical}）。现在，
$\mathit{Mor}^{P', \mathcal{M}}_U(Y, X)$ 是 Noether 代数空间，
因为由引理 \ref{moduli-lemma-Mor-qc-over-base}，它在 $U$ 上有限表示。
所以该交也是 Noether 代数空间。由于态射
$$
\mathit{Isom}_{\Polarizedstack}(\upsilon, \chi)
\longrightarrow
\mathit{Isom}_U(Y, X) \cap \mathit{Mor}^{P', \mathcal{M}}_U(Y, X)
$$
仿射（见上文），结论成立。
\end{proof}

\begin{lemma}
\label{moduli-lemma-polarized-qs-lfp}
态射 $\Polarizedstack \to \Spec(\mathbf{Z})$ 拟分离且局部有限表示。
\end{lemma}

\begin{proof}
要检验 $\Polarizedstack \to \Spec(\mathbf{Z})$ 拟分离，需要证明其
对角态射拟紧且拟分离。这立即由引理
\ref{moduli-lemma-polarized-diagonal-separated-fp} 得出。
要证明 $\Polarizedstack \to \Spec(\mathbf{Z})$ 局部有限表示，
只需证明 $\Polarizedstack$ 保持极限，见《叠的极限》中的命题
\ref{stacks-limits-proposition-characterize-locally-finite-presentation}。
这就是《Quot》中的引理 \ref{quot-lemma-polarized-limits}。
\end{proof}

\begin{lemma}
\label{moduli-lemma-bounded-polarized}
设 $n \geq 1$ 是整数，$P$ 是数值多项式。设
$$
T \subset |\Polarizedstack|
$$
是具有下述性质的子集：对每个 $\xi \in T$，存在一个域 $k$ 和
$\Polarizedstack$ 在 $k$ 上的一个表示 $\xi$ 的对象
$(X, \mathcal{L})$，使得
\begin{enumerate}
\item $\mathcal{L}$ 在 $X$ 上的 Hilbert 多项式为 $P$；并且
\item 存在闭浸入 $i : X \to \mathbf{P}^n_k$，使得
$i^*\mathcal{O}_{\mathbf{P}^n}(1) \cong \mathcal{L}$。
\end{enumerate}
则 $T$ 是 Noether 拓扑空间，特别地，它拟紧。
\end{lemma}

\begin{proof}
注意，$|\Polarizedstack|$ 是局部 Noether 拓扑空间，见《叠的态射》
中的引理 \ref{stacks-morphisms-lemma-Noetherian-topology}
（这里还用到 $\Spec(\mathbf{Z})$ 是 Noether 的，因而由引理
\ref{moduli-lemma-polarized-qs-lfp} 和《叠的态射》中的引理
\ref{stacks-morphisms-lemma-locally-finite-type-locally-noetherian}，
$\Polarizedstack$ 是局部 Noether 代数叠）。因此，
$|\Polarizedstack|$ 的任意拟紧子集都是 Noether 拓扑空间，而这种
子集的任意子集也 Noether，见《拓扑》中的引理
\ref{topology-lemma-finite-union-Noetherian} 和
\ref{topology-lemma-Noetherian}。
% P12-ERR-0223: see ../control/ERRATA_CANDIDATES.jsonl
所以只需找到一个包含 $T$ 的拟紧子集。

\medskip\noindent
由引理 \ref{moduli-lemma-hilb-proper-over-base}，代数空间
$$
H =
\Hilbfunctor^{P, \mathcal{O}(1)}_{\mathbf{P}^n_\mathbf{Z}/\Spec(\mathbf{Z})}
$$
在 $\Spec(\mathbf{Z})$ 上真。由《Quot》中的引理
\ref{quot-lemma-extend-hilb-to-spaces}\footnote{
% P12-ERR-0224: see ../control/ERRATA_CANDIDATES.jsonl
我们将在后文看到（在此插入未来引用），$H$ 是概形，因而无需使用
本引理和《Quot》中的引理
\ref{quot-lemma-extend-polarized-to-spaces}。}，
$H$ 的恒等态射对应于闭子空间
$$
Z \subset \mathbf{P}^n_H
$$
它在 $H$ 上真、平坦且有限表示；限制
$\mathcal{N} = \mathcal{O}(1)|_Z$ 在 $Z/H$ 上相对丰富，并且在
$Z \to H$ 的纤维上具有 Hilbert 多项式 $P$。特别地，偶
$(Z \to H, \mathcal{N})$ 定义态射
$$
H \longrightarrow \Polarizedstack
$$
它把概形态射 $U \to H$ 送到族
$(Z_U \to U, \mathcal{N}_U)$ 的分类态射，见《Quot》中的引理
\ref{quot-lemma-extend-polarized-to-spaces}。
% P12-ERR-0225: see ../control/ERRATA_CANDIDATES.jsonl
由于 $H$ 是 Noether 代数空间（因为它在 $\mathbf{Z})$ 上真），
$|H|$ 是 Noether 的，因而拟紧。映射
$$
|H| \longrightarrow |\Polarizedstack|
$$
连续，所以其像拟紧。因此只需证明 $T$ 包含于
$|H| \to |\Polarizedstack|$ 的像中。然而，假设 (1) 和 (2)
恰好表达了这一事实：
% P12-ERR-0226: see ../control/ERRATA_CANDIDATES.jsonl
任取闭浸入 $i : X \to \mathbf{P}^n_k$，满足
$i^*\mathcal{O}_{\mathbf{P}^n}(1) \cong \mathcal{L}$，
由 $H$ 的模解释可得到 $H$ 的一个 $k$-值点。
这就完成了本引理的证明。
\end{proof}





\section{真态射上复形模空间的性质}
\label{moduli-section-complexes}

\noindent
设 $f : X \to B$ 是真、平坦且有限表示的代数空间态射。则参数化
负自 Ext 消失的相对完美复形的叠 $\Complexesstack_{X/B}$ 是代数的。
见《Quot》中的定理 \ref{quot-theorem-complexes-algebraic}。

\begin{lemma}
\label{moduli-lemma-complexes-diagonal-affine-fp}
$\Complexesstack_{X/B}$ 在 $B$ 上的对角态射仿射且有限表示。
\end{lemma}

\begin{proof}
《Quot》中的引理 \ref{quot-lemma-complexes-diagonal} 已证明该对角态射
可由代数空间表示。从其证明可知，需要证明下述事实：给定 $B$ 上概形
$T$ 以及对象 $E, E' \in D(\mathcal{O}_{X_T})$，使得
$(T, E)$ 和 $(T, E')$ 是 $\Complexesstack_{X/B}$ 在 $T$ 上的
纤维范畴中的对象，则
$\mathit{Isom}(E, E') \to T$ 仿射且有限表示。这里
$\mathit{Isom}(E, E')$ 是函子
$$
(\Sch/T)^{opp} \to \textit{Sets},\quad
T' \mapsto \{\varphi : E_{T'} \to E'_{T'}
\text{ 是下列范畴中的同构：}D(\mathcal{O}_{X_{T'}})\}
$$
其中 $E_{T'}$ 和 $E'_{T'}$ 是 $E$ 和 $E'$ 向 $X_{T'}$ 的导出拉回。
考虑由规则
$$
(\Sch/T)^{opp} \to \textit{Sets},\quad
% P12-ERR-0227: see ../control/ERRATA_CANDIDATES.jsonl
T' \mapsto \Hom_{\mathcal{O}_{X_{T'}}}(E_T, E'_T)
$$
定义的函子 $H = \SheafHom(E, E')$。由《Quot》中的引理
\ref{quot-lemma-complexes-open-neg-exts-vanishing}，这是在 $T$
上仿射且有限表示的代数空间。对
$H' = \SheafHom(E', E)$、$I = \SheafHom(E, E)$ 和
$I' = \SheafHom(E', E')$ 也同样。因此可见
$$
% P12-ERR-0228: see ../control/ERRATA_CANDIDATES.jsonl
\mathit{Isom}(E, E') = (H' \times_T H) \times_{c, I \times_T I', \sigma} T
$$
其中
$c(\varphi', \varphi) = (\varphi \circ \varphi', \varphi' \circ \varphi)$，
而 $\sigma = (\text{id}, \text{id})$
（与《Quot》中的命题 \ref{quot-proposition-isom} 的证明比较）。
因此，作为在 $T$ 上仿射的概形的纤维积，
$\mathit{Isom}(E, E')$ 在 $T$ 上仿射。类似地，
$\mathit{Isom}(E, E')$ 在 $T$ 上有限表示。
\end{proof}

\begin{lemma}
\label{moduli-lemma-complexes-qs-lfp}
态射 $\Complexesstack_{X/B} \to B$ 拟分离且局部有限表示。
\end{lemma}

\begin{proof}
要检验 $\Complexesstack_{X/B} \to B$ 拟分离，需要证明其对角态射
拟紧且拟分离。这立即由引理
\ref{moduli-lemma-complexes-diagonal-affine-fp} 得出。
要证明 $\Complexesstack_{X/B} \to B$ 局部有限表示，需要证明
$\Complexesstack_{X/B} \to B$ 保持极限，见《叠的极限》中的命题
\ref{stacks-limits-proposition-characterize-locally-finite-presentation}。
这由《Quot》中的引理 \ref{quot-lemma-complexes-limits} 得出
（略去一个小细节）。
\end{proof}
% Generated mechanically from the frozen moduli-curves.tex target.
% Do not edit this generated file; edit the bound locale source instead.

% OK, start here.
%

\setcounter{chapter}{108}
\chapter{曲线模空间}


\phantomsection
\label{moduli-curves-section-phantom}





\section{引言}
\label{moduli-curves-section-introduction}

\noindent
本章讨论若干熟知的曲线模叠。一个参考文献是
Deligne 与 Mumford 的著名论文，见 \cite{DM}。




\section{约定与语言滥用}
\label{moduli-curves-section-conventions}

\noindent
我们继续采用《叠的性质》中的章节
\ref{stacks-properties-section-conventions} 所引入的约定和语言滥用。
除非另有说明，基概形均为 $\Spec(\mathbf{Z})$。






\section{曲线叠}
\label{moduli-curves-section-stack-curves}

\noindent
本节是《Quot》中的章节 \ref{quot-section-curves} 的续篇。
设 $\Curvesstack$ 为这样一个叠：它在概形 $S$ 上的截面范畴是
$S$ 上曲线族的范畴。需要特别记住，所谓
{\it 曲线族}，是指态射 $f : X \to S$，其中 $X$
是代数空间（！），而 $f$ 平坦、真、有限表示且相对维数
$\leq 1$。我们已经知道 $\Curvesstack$ 是
$\mathbf{Z}$ 上的代数叠，见《Quot》中的定理
\ref{quot-theorem-curves-algebraic}。
如果在叠的定义中不允许代数空间，则该定理并不成立。

\medskip\noindent
基变换通常用下标表示，但对 $\Curvesstack$ 不能采用这种记号，
因为 $\Curvesstack_S$ 已用来表示 $S$ 上的纤维范畴。
因此，我们在《Quot》中的注
\ref{quot-remark-curves-base-change} 里用 $B\text{-}\Curvesstack$
表示到代数空间 $B$ 的基变换
$$
B\text{-}\Curvesstack = \Curvesstack \times B
$$
右侧的积取在终对象上，即取在 $\Spec(\mathbf{Z})$ 上。
左侧对象是这样一个叠：它分类 $B$ 上概形范畴中的曲线族。
特别地，若 $k$ 是域，则
$$
k\text{-}\Curvesstack = \Curvesstack \times \Spec(k)
$$
是分类 $k$ 上概形范畴中的曲线族的模叠。
在继续之前，先作如下相容性检验。

\begin{lemma}
\label{moduli-curves-lemma-extend-curves-to-spaces}
设 $T \to B$ 是代数空间的态射。范畴
$$
\Mor_B(T, B\text{-}\Curvesstack) = \Mor(T, \Curvesstack)
$$
就是 $T$ 上曲线族的范畴。
\end{lemma}

\begin{proof}
$T$ 上的曲线族是代数空间的态射 $f : X \to T$，它平坦、真、
有限表示且相对维数 $\leq 1$（《空间的态射》中的定义
\ref{spaces-morphisms-definition-relative-dimension}）。
这与《Quot》中的情形 \ref{quot-situation-curves} 的定义完全相同，
不同之处仅在于允许基 $T$ 为代数空间。
代数叠和代数空间的默认基范畴是概形范畴，因此本引理并不直接
由定义得出。不过，我们仍建议读者略过此证明。

\medskip\noindent
由上面给出的 $B\text{-}\Curvesstack$ 的积描述，只需证明绝对情形。
取概形 $U$ 以及满的 \'etale 态射 $p : U \to T$。
令 $R = U \times_T U$，其投影为 $s, t : R \to U$。

\medskip\noindent
设 $v : T \to \Curvesstack$ 是态射。则 $v \circ p$
对应于曲线族 $X_U \to U$。典范 $2$-态射
$v \circ p \circ t \to v \circ p \circ s$
是一个同构
$\varphi : X_U \times_{U, s} R \to X_U \times_{U, t} R$。
该同构在 $R \times_{s, t} R$ 上满足余圈条件。
由《自举》中的引理 \ref{bootstrap-lemma-descend-algebraic-space}，
得到代数空间的态射 $X \to T$；它在 $U$ 上的拉回等于 $X_U$，
并与 $\varphi$ 相容。由于 $\{U \to T\}$ 是 \'etale 覆盖，
《空间的下降》中的引理
\ref{spaces-descent-lemma-descending-property-flat}、
\ref{spaces-descent-lemma-descending-property-proper} 和
\ref{spaces-descent-lemma-descending-property-finite-presentation}
表明 $X \to T$ 平坦、真且有限表示。
此外，$X \to T$ 的相对维数 $\leq 1$，因为这是 \'etale 局部性质。
故 $X \to T$ 是 $T$ 上的曲线族。

\medskip\noindent
反之，设 $X \to T$ 是曲线族。则基变换 $X_U$ 确定态射
$w : U \to \Curvesstack$，而典范同构
$X_U \times_{U, s} R \to X_U \times_{U, t} R$
确定满足余圈条件的 $2$-箭头 $w \circ s \to w \circ t$。
% P12-ERR-0229: see ../control/ERRATA_CANDIDATES.jsonl
因此，由商 $[U/R]$ 的泛性质得到态射
$v : T = [U/R] \to \Curvesstack$；见《空间中的群胚》中的引理
\ref{spaces-groupoids-lemma-quotient-stack-2-coequalizer}。
% P12-ERR-0230: see ../control/ERRATA_CANDIDATES.jsonl
（事实上，在此情形下，更容易回到引入语言滥用之前，直接构造函子
$\Sch/T \to \Curvesstack$；该函子“就是”态射
$T \to \Curvesstack$。）

\medskip\noindent
我们略去如下验证：上述构造可扩张到对象之间的态射，且彼此拟逆。
\end{proof}













\section{极化曲线叠}
\label{moduli-curves-section-polarized-curves}

\noindent
本节展开《Quot》中的注
\ref{quot-remark-alternative-approach-curves} 所讨论的部分内容。
考虑 $2$-纤维积
$$
\xymatrix{
\Curvesstack \times_{\Spacesstack'_{fp, flat, proper}}
\Polarizedstack \ar[r] \ar[d] &
\Polarizedstack \ar[d] \\
\Curvesstack \ar[r] &
\Spacesstack'_{fp, flat, proper}
}
$$
我们将此 $2$-纤维积记为
$$
\textit{PolarizedCurves} =
\Curvesstack
\times_{\Spacesstack'_{fp, flat, proper}}
\Polarizedstack
$$
该纤维积参数化极化曲线，即赋有相对丰沛可逆层的曲线族。
更确切地，$\textit{PolarizedCurves}$ 的对象是二元组
$(X \to S, \mathcal{L})$，其中
\begin{enumerate}
\item $X \to S$ 是真、平坦、有限表示且相对维数
$\leq 1$ 的概形态射；以及
\item $\mathcal{L}$ 是在 $X/S$ 上相对丰沛的可逆
$\mathcal{O}_X$-模。
\end{enumerate}
$\textit{PolarizedCurves}$ 的对象之间的态射
$(X' \to S', \mathcal{L}') \to (X \to S, \mathcal{L})$
由三元组 $(f, g, \varphi)$ 给出，其中
$f : X' \to X$ 和 $g : S' \to S$ 是概形态射，并组成交换图
$$
\xymatrix{
X' \ar[d] \ar[r]_f & X \ar[d] \\
S' \ar[r]^g & S
}
$$
该图诱导同构 $X' \to S' \times_S X$；换言之，该图是 Cartesian 的，
且 $\varphi : f^*\mathcal{L} \to \mathcal{L}'$ 是同构。
复合按显然的方式定义。

\begin{lemma}
\label{moduli-curves-lemma-polarized-curves-in-polarized}
态射
$\textit{PolarizedCurves} \to
\Polarizedstack$ 是开闭浸入。
\end{lemma}

\begin{proof}
这是因为 $1$-态射
$\Curvesstack \to \Spacesstack'_{fp, flat, proper}$
可表为开闭浸入；见《Quot》中的引理
\ref{quot-lemma-curves-open-and-closed-in-spaces}。
\end{proof}

\begin{lemma}
\label{moduli-curves-lemma-polarized-curves-over-curves}
态射
$\textit{PolarizedCurves} \to \Curvesstack$
光滑且满。
\end{lemma}

\begin{proof}
满性。给定域 $k$ 以及 $k$ 上维数 $\leq 1$ 的真代数空间
$X$，即 $\Curvesstack$ 在 $k$ 上的对象。
由《域上的空间》中的引理
\ref{spaces-over-fields-lemma-codim-1-point-in-schematic-locus}，
代数空间 $X$ 是概形。因此 $X$ 是 $k$ 上维数 $\leq 1$ 的真概形。
由《簇》中的引理 \ref{varieties-lemma-dim-1-proper-projective}，
% P12-ERR-0231: see ../control/ERRATA_CANDIDATES.jsonl
可见 $X$ 在 $\kappa$ 上是 H-射影的。
特别地，$X$ 上存在丰沛可逆 $\mathcal{O}_X$-模
$\mathcal{L}$。于是 $(X, \mathcal{L})$ 是
$\textit{PolarizedCurves}$ 在 $k$ 上的对象，并映到 $X$。

\medskip\noindent
光滑性。设 $X \to S$ 是 $\Curvesstack$ 的对象，即态射
$S \to \Curvesstack$。显然
$$
\textit{PolarizedCurves}
\times_{\Curvesstack} S
\subset \Picardstack_{X/S}
$$
是由如下对象 $(T/S, \mathcal{L}/X_T)$ 构成的子叠：
$\mathcal{L}$ 在 $X_T/T$ 上丰沛。由《空间的下降》中的引理
\ref{spaces-descent-lemma-ample-in-neighbourhood}，这是开子叠。
又由《模叠》中的引理 \ref{moduli-lemma-pic-curves-smooth}，
$\Picardstack_{X/S} \to S$ 光滑，结论遂得。
\end{proof}

\begin{lemma}
\label{moduli-curves-lemma-etale-locally-scheme}
设 $X \to S$ 是曲线族。则存在 \'etale 覆盖
$\{S_i \to S\}$，使得 $X_i = X \times_S S_i$ 是概形。
甚至可以假设 $X_i$ 在 $S_i$ 上是 H-射影的。
\end{lemma}

\begin{proof}
这是引理 \ref{moduli-curves-lemma-polarized-curves-over-curves} 的直接推论。
% P12-ERR-0232: see ../control/ERRATA_CANDIDATES.jsonl
具体而言，展开定义，该引理给出一个满光滑态射 $S' \to S$，
使得 $X' = X \times_S S'$ 赋有在 $X'/S'$ 上丰沛的可逆
$\mathcal{O}_{X'}$-模 $\mathcal{L}'$。
于是可用一个 \'etale 覆盖 $\{S_i \to S\}$ 加细光滑覆盖
$\{S' \to S\}$；见《态射进阶》中的引理
\ref{more-morphisms-lemma-etale-dominates-smooth}。
用适当开覆盖替换 $S_i$ 后，可以假设 $X_i \to S_i$ 是
H-射影的；见《态射》中的引理
\ref{morphisms-lemma-proper-ample-locally-projective} 和
\ref{morphisms-lemma-characterize-locally-projective}
（《态射进阶》中的章节
\ref{more-morphisms-section-projective} 也对此作了详细讨论）。
\end{proof}




\section{曲线叠的性质}
\label{moduli-curves-section-properties}

\noindent
下面的引理对曲面的模空间并不成立；见注
\ref{moduli-curves-remark-boundedness-aut-does-not-work-surfaces}。

\begin{lemma}
\label{moduli-curves-lemma-curves-diagonal-separated-fp}
$\Curvesstack$ 的对角态射分离且有限表示。
\end{lemma}

\begin{proof}
回忆 $\Curvesstack$ 是保持极限的代数叠；见《Quot》中的引理
\ref{quot-lemma-curves-limits}。
由《叠的极限》中的引理
\ref{stacks-limits-lemma-limit-preserving-diagonal}，这蕴含
% P12-ERR-0233: see ../control/ERRATA_CANDIDATES.jsonl
$\Delta : \Polarizedstack \to \Polarizedstack \times \Polarizedstack$
保持极限。因此，由《叠的极限》中的命题
\ref{stacks-limits-proposition-characterize-locally-finite-presentation}，
$\Delta$ 局部有限表示。

\medskip\noindent
下面证明 $\Delta$ 分离。为此，只需说明：给定概形 $U$ 以及
$\Curvesstack$ 在 $U$ 上的两个对象 $Y \to U$ 和 $X \to U$，
代数空间
$$
\mathit{Isom}_U(Y, X)
$$
分离。
% P12-ERR-0234: see ../control/ERRATA_CANDIDATES.jsonl
这已经由《模叠》中的引理 \ref{moduli-lemma-Mor-s-lfp} 和
\ref{moduli-lemma-Isom-in-Mor} 得出：目标是分离代数空间。

\medskip\noindent
为完成证明，我们说明 $\Delta$ 拟紧。由于 $\Delta$
可由代数空间表示，只需检验 $\Delta$ 沿满光滑态射
$U \to \Curvesstack \times \Curvesstack$ 的基变换拟紧
（例如见《叠的性质》中的引理
\ref{stacks-properties-lemma-check-property-covering}）。
取 $U = \coprod U_i$ 为仿射开集的不交并，并取满光滑态射
$$
U \longrightarrow
\textit{PolarizedCurves} \times \textit{PolarizedCurves}
$$
由于引理 \ref{moduli-curves-lemma-polarized-curves-over-curves} 表明
$\textit{PolarizedCurves} \to \Curvesstack$ 满且光滑，
$U \to \Curvesstack \times \Curvesstack$ 也满且光滑。
又因为 $\textit{PolarizedCurves}$ 保持极限
（由《Artin 公理》中的引理
\ref{artin-lemma-fibre-product-limit-preserving} 以及《Quot》中的引理
\ref{quot-lemma-curves-limits}、
\ref{quot-lemma-polarized-limits} 和
\ref{quot-lemma-spaces-limits}），可见
$\textit{PolarizedCurves} \to \Spec(\mathbf{Z})$ 局部有限表示，
因而 $U_i \to \Spec(\mathbf{Z})$ 局部有限表示
（《叠的极限》中的命题
\ref{stacks-limits-proposition-characterize-locally-finite-presentation}
以及《叠的态射》中的引理
\ref{stacks-morphisms-lemma-composition-finite-presentation} 和
\ref{stacks-morphisms-lemma-smooth-locally-finite-presentation}）。
特别地，$U_i$ 是 Noether 仿射概形。这就化归到下一段所讨论的情形。

\medskip\noindent
在本段中，给定 Noether 仿射概形 $U$ 以及
$\textit{PolarizedCurves}$ 在 $U$ 上的两个对象
$(Y, \mathcal{N})$ 和 $(X, \mathcal{L})$，我们证明代数空间
$$
\mathit{Isom}_U(Y, X)
$$
拟紧。由于 $U$ 的连通分支既开又闭，可以用这些连通分支替换
$U$。因此可以并确实假设 $U$ 连通。取一点 $u \in U$。
令 $Q$、$P$ 为这些族的 Hilbert 多项式，即
$$
Q(n) = \chi(Y_u, \mathcal{N}_u^{\otimes n})
\quad\text{且}\quad
P(n) = \chi(X_u, \mathcal{L}_u^{\otimes n})
$$
见《簇》中的引理
\ref{varieties-lemma-numerical-polynomial-from-euler}。
由于 $U$ 连通，且函数
$u \mapsto \chi(Y_u, \mathcal{N}_u^{\otimes n})$ 和
$u \mapsto \chi(X_u, \mathcal{L}_u^{\otimes n})$
局部常值（见《概形的导出范畴》中的引理
\ref{perfect-lemma-chi-locally-constant-geometric}），
可见在 $U$ 的每一点都得到同一 Hilbert 多项式。令
$$
\mathcal{M} = \text{pr}_1^*\mathcal{N}
\otimes_{\mathcal{O}_{Y \times_U X}} \text{pr}_2^*\mathcal{L}
$$
为 $Y \times_U X$ 上的模。给定 $U$ 上某概形 $T$ 以及
% P12-ERR-0235: see ../control/ERRATA_CANDIDATES.jsonl
$(f, \varphi) \in \mathit{Isom}_U(Y, X)(T)$，则对每个 $t \in T$，
有
\begin{align*}
\chi(Y_t, (\text{id} \times f)^*\mathcal{M}^{\otimes n})
& =
\chi(Y_t,
\mathcal{N}_t^{\otimes n} \otimes_{\mathcal{O}_{Y_t}}
f_t^*\mathcal{L}_t^{\otimes n}) \\
& =
n\deg(\mathcal{N}_t) + n\deg(f_t^*\mathcal{L}_t) +
\chi(Y_t, \mathcal{O}_{Y_t}) \\
& =
Q(n) + n\deg(\mathcal{L}_t) \\
& =
Q(n) + P(n) - P(0)
\end{align*}
这里使用了真曲线的 Riemann--Roch；更确切地说，使用了
《簇》中的定义 \ref{varieties-definition-degree-invertible-sheaf}、
引理 \ref{varieties-lemma-degree-tensor-product}，以及
$f_t$ 是同构这一事实。令 $P'(t) = Q(t) + P(t) - P(0)$，得到
$$
\mathit{Isom}_U(Y, X) =
\mathit{Isom}_U(Y, X) \cap \mathit{Mor}^{P', \mathcal{M}}_U(Y, X)
$$
该交集是 $\mathit{Mor}_U(Y, X)$ 的两个开子空间之交；见
《模叠》中的引理 \ref{moduli-lemma-Isom-in-Mor} 和注
\ref{moduli-remark-Mor-numerical}。
现在，$\mathit{Mor}^{P', \mathcal{M}}_U(Y, X)$
在 $U$ 上有限表示，故由《模叠》中的引理
\ref{moduli-lemma-Mor-qc-over-base}，它是 Noether 代数空间。
因此该交集也是 Noether 代数空间，证明完成。
\end{proof}

\begin{remark}
\label{moduli-curves-remark-boundedness-aut-does-not-work-surfaces}
引理 \ref{moduli-curves-lemma-curves-diagonal-separated-fp} 的证明中的有界性论证
对曲面的模空间并不适用；事实上，其结论也是错误的，例如域上的
K3 曲面可以具有无限离散自同构群。该论证不适用的“原因”是：
对于域上的射影曲面 $S$，给定 Hilbert 多项式分别为 $Q$ 和 $P$
的丰沛可逆层 $\mathcal{N}$ 和 $\mathcal{L}$，
$\mathcal{N} \otimes_{\mathcal{O}_S} \mathcal{L}$ 的 Hilbert 多项式
并不存在先验界。用相交论的语言来说，若 $H_1$、$H_2$ 是 $S$
上的丰沛有效 Cartier 除子，则相交数 $H_1 \cdot H_2$
不能由 $H_1 \cdot H_1$ 和 $H_2 \cdot H_2$ 给出（上）界。
\end{remark}

\begin{lemma}
\label{moduli-curves-lemma-curves-qs-lfp}
态射 $\Curvesstack \to \Spec(\mathbf{Z})$ 拟分离且局部有限表示。
\end{lemma}

\begin{proof}
要检验 $\Curvesstack \to \Spec(\mathbf{Z})$ 拟分离，需要证明其
对角态射拟紧且拟分离。这立即由引理
\ref{moduli-curves-lemma-curves-diagonal-separated-fp} 得出。
要证明 $\Curvesstack \to \Spec(\mathbf{Z})$ 局部有限表示，只需证明
$\Curvesstack$ 保持极限；见《叠的极限》中的命题
\ref{stacks-limits-proposition-characterize-locally-finite-presentation}。
这就是《Quot》中的引理 \ref{quot-lemma-curves-limits}。
\end{proof}





\section{曲线叠的开子叠}
\label{moduli-curves-section-open}

\noindent
下文常用代数空间态射的一个性质 $P$ 来刻画
$\Curvesstack$ 的开子叠。要说明 $P$ 定义开子叠，只需检验
\begin{enumerate}
\item[(o)] 给定曲线族 $f : X \to S$，存在最大开子概形
$S' \subset S$，使得
$f|_{f^{-1}(S')} : f^{-1}(S') \to S'$ 具有性质 $P$，
并且 $S'$ 的构造与任意基变换交换。
\end{enumerate}
事实上，假设 (o) 成立。取概形 $U$ 以及满光滑态射
$m : U \to \Curvesstack$。令
$R = U \times_{\Curvesstack} U$，并以 $t, s : R \to U$
表示投影。回忆 $\Curvesstack = [U/R]$ 是一个展示；见
《代数叠》中的引理 \ref{algebraic-lemma-stack-presentation} 和定义
\ref{algebraic-definition-presentation}。
由 $\Curvesstack$ 作为曲线叠的构造，态射 $m$ 是曲线族
$C \to U$ 的分类态射。图
$$
\xymatrix{
R \ar[r]_s \ar[d]_t & U \ar[d] \\
U \ar[r] & \Curvesstack
}
$$
的 $2$-交换性蕴含
$C \times_{U, s} R  \cong C \times_{U, t} R$
（这是 $R$ 上曲线族的同构）。令 $W \subset U$ 为满足如下条件的
最大开子概形：
% P12-ERR-0236: see ../control/ERRATA_CANDIDATES.jsonl
$f|_{f^{-1}(W)} : f^{-1}(W) \to W$ 具有 (o) 中的性质 $P$。
由于按 (o)，$W$ 的构造与基变换交换，再结合上述同构，得到
$s^{-1}(W) = t^{-1}(W)$。因此，由《叠的性质》中的引理
\ref{stacks-properties-lemma-immersion-presentation}，
$W \subset U$ 对应于开子叠
$$
\Curvesstack^P \subset \Curvesstack
$$

\medskip\noindent
% P12-ERR-0237: see ../control/ERRATA_CANDIDATES.jsonl
继续采用上一段的设定。我们断言开子叠 $\Curvesstack^P$
具有以下两个泛性质：
\begin{enumerate}
\item 给定曲线族 $X \to S$，下列条件等价：
\begin{enumerate}
\item 分类态射 $S \to \Curvesstack$ 经 $\Curvesstack^P$ 分解；
\item 态射 $X \to S$ 具有性质 $P$；
\end{enumerate}
% P12-ERR-0238: see ../control/ERRATA_CANDIDATES.jsonl
\item 给定域 $k$ 上维数 $\leq 1$ 的真概形 $X$，下列条件等价：
\begin{enumerate}
\item 分类态射 $\Spec(k) \to \Curvesstack$ 经
$\Curvesstack^P$ 分解；
\item 态射 $X \to \Spec(k)$ 具有性质 $P$。
\end{enumerate}
\end{enumerate}
这可通过考察 $2$-纤维积
$$
\xymatrix{
T \ar[r]_p \ar[d]_q & U \ar[d] \\
S \ar[r] & \Curvesstack
}
$$
得到。注意，$T \to S$ 是 $U \to \Curvesstack$ 的基变换，
故满且光滑。因此，(o) 给出的开集 $S' \subset S$
由其在 $T$ 中的逆像确定。然而，由 (o) 中这些开集在基变换下
的不变性，并且由 $2$-交换性有
$X \times_S T \cong C \times_U T$，得到 $T$ 中开集的等式
$q^{-1}(S') = p^{-1}(W)$。这立即蕴含 (1)，而 (2) 是 (1) 的特例。

\medskip\noindent
给定代数空间态射的两个性质 $P$ 和 $Q$，假设已经证明
$\Curvesstack^Q$ 是 $\Curvesstack$ 的开子叠，则可以用完全相同的
方法证明
$\Curvesstack^{Q, P} \subset \Curvesstack^Q$
是开子叠。我们略去精确说明。

\section{具有有限约化自同构群的曲线}
\label{moduli-curves-section-finite-aut}

\noindent
设 $X$ 是域 $k$ 上维数 $\leq 1$ 的真概形，即
$\Curvesstack$ 在 $k$ 上的对象。由引理
\ref{moduli-curves-lemma-curves-diagonal-separated-fp}，自同构群代数空间
$\mathit{Aut}(X)$ 在 $k$ 上有限型且分离。
特别地，$\mathit{Aut}(X)$ 是群概形；见《空间中群胚进阶》中的引理
\ref{spaces-more-groupoids-lemma-group-space-scheme-locally-finite-type-over-k}。
若 $k$ 的特征为零，则 $\mathit{Aut}(X)$ 约化，甚至在 $k$
上光滑（《群胚》中的引理
\ref{groupoids-lemma-group-scheme-characteristic-zero-smooth}）。
然而，一般而言，即使 $X$ 几何约化，$\mathit{Aut}(X)$
也不约化。

\begin{example}[非约化自同构群]
\label{moduli-curves-example-non-reduced}
设 $k$ 是特征 $2$ 的代数闭域。令
$Y = Z = \mathbf{P}^1_k$。在 $\mathbf{A}^1_k$ 中取三个两两不同的
$k$-值点 $a, b, c$。把
% P12-ERR-0239: see ../control/ERRATA_CANDIDATES.jsonl
$\mathbf{A}^1_k \subset \mathbf{P}^1_k = Y = Z$ 看作开子概形，
得到闭浸入
$$
T =  \Spec(k[t]/(t - a)^2) \amalg \Spec(k[t]/(t - b)^2)
\amalg \Spec(k[t]/(t - c)^2)
\longrightarrow
\mathbf{P}^1_k
$$
令 $X$ 为图
$$
\xymatrix{
T \ar[r] \ar[d] & Y \ar[d] \\
Z \ar[r] & X
}
$$
中的推出。令 $U \subset X$ 为
$\mathbf{A}^1_k \amalg \mathbf{A}^1_k$ 的像所构成的仿射开部分。
则有等化子图
$$
\xymatrix{
\mathcal{O}_X(U) \ar[r] &
k[t] \times k[t] \ar@<1ex>[r] \ar@<-1ex>[r] &
k[t]/(t - a)^2 \times k[t]/(t - b)^2 \times k[t]/(t - c)^2
}
$$
在双数 $A = k[\epsilon]$ 上，该等化子图有一个非平凡自同构，
它把 $t$ 送到 $t + \epsilon$。留给读者验证：该自同构扩张为
$X$ 在 $A$ 上的自同构。另一方面，读者容易证明 $X$ 在 $k$
上的自同构群有限。因此 $\mathit{Aut}(X)$ 必不约化。
\end{example}

\noindent
设 $X$ 是域 $k$ 上维数 $\leq 1$ 的真概形，即
$\Curvesstack$ 在 $k$ 上的对象。即使 $X$ 光滑且几何连通，
若 $\mathit{Aut}(X)$ 几何约化，也未必有维数 $0$。

\begin{example}[正维光滑自同构群]
\label{moduli-curves-example-pos-dim}
设 $k$ 是代数闭域。若 $X$ 是光滑的亏格 $0$、相应地\ $1$ 的曲线，
则其自同构群的维数为 $3$、相应地\ $1$。
事实上，在亏格 $0$ 的情形，由《代数曲线》中的命题
\ref{curves-proposition-projective-line}，有
$X \cong \mathbf{P}^1_k$。由于函子同构
$$
\mathit{Aut}(\mathbf{P}^1_k) = \text{PGL}_{2, k}
$$
可见维数为 $3$。另一方面，若 $X$ 的亏格为 $1$，则映射
$X = \underline{\Hilbfunctor}^1_{X/k} \to
\underline{\Picardfunctor}^1_{X/k}$ 是同构；见
《曲线的 Picard 概形》中的引理 \ref{pic-lemma-picard-pieces} 以及
《代数曲线》中的定理
\ref{curves-theorem-curves-rational-maps}。
因此 $X$ 具有阿贝尔簇结构
（因为 $\underline{\Picardfunctor}^1_{X/k} \cong
\underline{\Picardfunctor}^0_{X/k}$）。
% P12-ERR-0240: see ../control/ERRATA_CANDIDATES.jsonl
特别地，$X$ 的（余）切丛平凡
（《群胚》中的引理
\ref{groupoids-lemma-group-scheme-module-differentials}）。
由此得到 $\dim_k H^0(X, T_X) = 1$，故
$\dim \mathit{Aut}(X) \leq 1$。另一方面，平移
（把 $X$ 看作群概形）给出 $\text{Aut}(X)$ 的一个 $1$-维部分，
% P12-ERR-0241: see ../control/ERRATA_CANDIDATES.jsonl
所以其维数确为 $1$。
\end{example}

\noindent
事实表明，$\Curvesstack$ 有一个开子叠，参数化自同构群
几何约化且有限的曲线。精确陈述如下。

\begin{lemma}
\label{moduli-curves-lemma-DM-curves}
% P12-ERR-0242: see ../control/ERRATA_CANDIDATES.jsonl
存在开子叠 $\Curvesstack^{DM} \subset \Curvesstack$，具有以下性质：
\begin{enumerate}
\item $\Curvesstack^{DM} \subset \Curvesstack$ 是最大的
DM 开子叠；
\item 给定曲线族 $X \to S$，下列条件等价：
\begin{enumerate}
\item 分类态射 $S \to \Curvesstack$ 经 $\Curvesstack^{DM}$ 分解；
\item 群代数空间 $\mathit{Aut}_S(X)$ 在 $S$ 上非分歧；
\end{enumerate}
% P12-ERR-0243: see ../control/ERRATA_CANDIDATES.jsonl
\item 给定域 $k$ 上维数 $\leq 1$ 的真概形 $X$，下列条件等价：
\begin{enumerate}
\item 分类态射 $\Spec(k) \to \Curvesstack$ 经
$\Curvesstack^{DM}$ 分解；
\item $\mathit{Aut}(X)$ 在 $k$ 上几何约化且维数为 $0$；
\item $\mathit{Aut}(X) \to \Spec(k)$ 非分歧。
\end{enumerate}
\end{enumerate}
\end{lemma}

\begin{proof}
具有性质 (1) 的开子叠之存在性由《叠的态射》中的引理
\ref{stacks-morphisms-lemma-open-DM-locus} 给出。
由《叠的态射》中的引理
\ref{stacks-morphisms-lemma-points-DM-locus}，该开子叠的点由
(3)(c) 刻画。(3)(b) 与 (3)(c) 的等价性是如下陈述：
若代数空间 $G$ 在域 $k$ 上局部有限型、几何约化且维数为 $0$，
则它在 $k$ 上非分歧。首先，由《域上的空间》中的引理
\ref{spaces-over-fields-lemma-locally-finite-type-dim-zero}，
$G$ 是概形。然后可在 $G$ 中取一个仿射开集，注意到它在 $k$
上为真，并应用《簇》中的引理
\ref{varieties-lemma-proper-geometrically-reduced-global-sections}。
略去次要细节。

\medskip\noindent
由 (3)，(2) 成立。事实上，态射
$\mathit{Aut}_S(X) \to S$ 局部有限型。因此，可以通过在概形 $S$
各点的纤维上检验，判断 $\mathit{Aut}_S(X) \to S$
是否在 $\mathit{Aut}_S(X)$ 的所有点处非分歧；见
《空间的态射》中的引理
\ref{spaces-morphisms-lemma-where-unramified}。
但基变换到 $S$ 的一点后，就回到 (3)(a) 与 (3)(c) 的等价性。
\end{proof}

\begin{lemma}
\label{moduli-curves-lemma-in-DM-locus-vector-fields}
设 $X$ 是域 $k$ 上维数 $\leq 1$ 的真概形。则性质
(3)(a)、(b)、(c) 还等价于
$\text{Der}_k(\mathcal{O}_X, \mathcal{O}_X) = 0$。
\end{lemma}

\begin{proof}
在上面的讨论中已经看到，$G = \mathit{Aut}(X)$ 是
$\Spec(k)$ 上有限型且分离的群概形；这里使用了引理
\ref{moduli-curves-lemma-curves-diagonal-separated-fp} 和《空间中群胚进阶》中的引理
\ref{spaces-more-groupoids-lemma-group-space-scheme-locally-finite-type-over-k}。
于是，$G$ 在 $k$ 上非分歧当且仅当 $\Omega_{G/k} = 0$
（《态射》中的引理 \ref{morphisms-lemma-unramified-omega-zero}）。
由《群胚》中的引理
\ref{groupoids-lemma-group-scheme-module-differentials}，
若 $T_{G/k, e} = 0$，则上述消失成立；这里 $T_{G/k, e}$ 是
$G$ 在单位元 $e \in G(k)$ 处的切空间，见《簇》中的定义
\ref{varieties-definition-tangent-space} 以及《簇》中的引理
\ref{varieties-lemma-tangent-space-cotangent-space} 里的公式。
由于 $\kappa(e) = k$，该切空间可用如下态射定义：
$\alpha : \Spec(k[\epsilon]) \to G = \mathit{Aut}(X)$，
其在 $\Spec(k)$ 上的限制为 $e$。于是，只需考察自同构
% P12-ERR-0244: see ../control/ERRATA_CANDIDATES.jsonl
$$
\alpha :
X \times_{\Spec(k)} \Spec(k[\epsilon])
\longrightarrow
X \times_{\Spec(k)} \Spec(k[\epsilon])
$$
它在 $\Spec(k[\epsilon])$ 上，且在 $\Spec(k)$ 上的限制为
$\text{id}_X$。这样的自同构称为无穷小自同构。

\medskip\noindent
$X$ 的无穷小自同构与 $\mathcal{O}_X$ 在 $k$ 上的导子按 $1$-对-$1$ 对应。
这由《态射进阶》中的引理
\ref{more-morphisms-lemma-difference-derivation} 和
\ref{more-morphisms-lemma-action-by-derivations} 得出
（我们只需要第一个引理，因为不关心反方向；另请参见
《态射进阶》中的注
\ref{more-morphisms-remark-another-special-case} 的说明）。
关于证明该等式的另一种论证，参见《形变问题》中的引理
\ref{examples-defos-lemma-schemes-TI}。
\end{proof}

\section{Cohen--Macaulay 曲线}
\label{moduli-curves-section-CM}

\noindent
$\Curvesstack$ 有一个开子叠，参数化 Cohen--Macaulay ``曲线''。

\begin{lemma}
\label{moduli-curves-lemma-CM-curves}
% P12-ERR-0245: see ../control/ERRATA_CANDIDATES.jsonl
存在开子叠 $\Curvesstack^{CM} \subset \Curvesstack$，使得：
\begin{enumerate}
\item 给定曲线族 $X \to S$，下列条件等价：
\begin{enumerate}
\item 分类态射 $S \to \Curvesstack$ 经 $\Curvesstack^{CM}$ 分解；
\item 态射 $X \to S$ 是 Cohen--Macaulay 态射；
\end{enumerate}
\item 给定域 $k$ 上的真概形 $X$，且 $\dim(X) \leq 1$，
下列条件等价：
\begin{enumerate}
\item 分类态射 $\Spec(k) \to \Curvesstack$ 经
$\Curvesstack^{CM}$ 分解；
\item $X$ 是 Cohen--Macaulay 概形。
\end{enumerate}
\end{enumerate}
\end{lemma}

\begin{proof}
设 $f : X \to S$ 是曲线族。由《空间的态射进阶》中的引理
\ref{spaces-more-morphisms-lemma-flat-finite-presentation-CM-open}，集合
$$
W = \{x \in |X| : f \text{ 在下述点为 Cohen--Macaulay：}x\}
$$
在 $|X|$ 中是开的，并且这个开集的构造与任意基变换交换。
由于 $f$ 为真，$S$ 的子集
$$
S' = S \setminus f(|X| \setminus W)
$$
是开的，且 $X \times_S S' \to S'$ 是 Cohen--Macaulay 态射。
此外，$S'$ 的构造与任意基变换交换，因为 $W$ 的构造具有此性质。
% P12-ERR-0246: see ../control/ERRATA_CANDIDATES.jsonl
因此，由第 \ref{moduli-curves-section-open} 节所讨论的方法，得到具有所需性质的开子叠。
\end{proof}

\begin{lemma}
\label{moduli-curves-lemma-CM-1-curves}
% P12-ERR-0247: see ../control/ERRATA_CANDIDATES.jsonl
存在开子叠 $\Curvesstack^{CM, 1} \subset \Curvesstack$，使得：
\begin{enumerate}
\item 给定曲线族 $X \to S$，下列条件等价：
\begin{enumerate}
\item 分类态射 $S \to \Curvesstack$ 经
$\Curvesstack^{CM, 1}$ 分解；
\item 态射 $X \to S$ 是 Cohen--Macaulay 态射，且具有相对维数 $1$
（《空间的态射》中的定义
\ref{spaces-morphisms-definition-relative-dimension}）；
\end{enumerate}
\item 给定域 $k$ 上的真概形 $X$，且 $\dim(X) \leq 1$，
下列条件等价：
\begin{enumerate}
\item 分类态射 $\Spec(k) \to \Curvesstack$ 经
$\Curvesstack^{CM, 1}$ 分解；
\item $X$ 是 Cohen--Macaulay 概形，且 $X$ 是维数为 $1$ 的等维概形。
\end{enumerate}
\end{enumerate}
\end{lemma}

\begin{proof}
由引理 \ref{moduli-curves-lemma-CM-curves}，若该开子叠存在，
则显然有 $\Curvesstack^{CM, 1} \subset \Curvesstack^{CM}$。
设 $f : X \to S$ 是曲线族，且 $f$ 是 Cohen--Macaulay 态射。
由《空间的态射进阶》中的引理
\ref{spaces-more-morphisms-lemma-lfp-CM-relative-dimension}，有按开闭子空间的分解
$$
X = X_0 \amalg X_1
$$
使得 $X_0 \to S$ 具有相对维数 $0$，而 $X_1 \to S$ 具有相对维数 $1$。
由于 $f$ 为真，$S$ 的子集
$$
S' = S \setminus f(|X_0|)
$$
是开的，且 $X \times_S S' \to S'$ 是相对维数为 $1$ 的
Cohen--Macaulay 态射。此外，$S'$ 的构造与任意基变换交换，
因为上述分解的构造具有此性质（相对维数在基变换下表现良好；见
《空间的态射》中的引理
\ref{spaces-morphisms-lemma-dimension-fibre-after-base-change}）。
因此，由第 \ref{moduli-curves-section-open} 节所讨论的方法，得到具有所需性质的开子叠。
\end{proof}

\section{给定亏格的曲线}
\label{moduli-curves-section-genus}

\noindent
Stacks Project 采用如下约定：域 $k$ 上的真 $1$-维概形 $X$
的亏格 $g$ 仅在 $H^0(X, \mathcal{O}_X) = k$ 时定义。
在此情况下，$g = \dim_k H^1(X, \mathcal{O}_X)$。
见《代数曲线》中的第 \ref{curves-section-genus} 节。
定义亏格所需的条件定义了一个开子叠；该开子叠又分解为开子叠的不交并，
每个亏格对应其中一个开子叠。

\begin{lemma}
\label{moduli-curves-lemma-pre-genus-curves}
% P12-ERR-0248: see ../control/ERRATA_CANDIDATES.jsonl
存在开子叠 $\Curvesstack^{h0, 1} \subset \Curvesstack$，使得：
\begin{enumerate}
\item 给定曲线族 $f : X \to S$，下列条件等价：
\begin{enumerate}
\item 分类态射 $S \to \Curvesstack$ 经
$\Curvesstack^{h0, 1}$ 分解；
\item $f_*\mathcal{O}_X = \mathcal{O}_S$，该等式在任意基变换后仍成立，
且 $f$ 的纤维维数为 $1$；
\end{enumerate}
\item 给定域 $k$ 上的真概形 $X$，且 $\dim(X) \leq 1$，
下列条件等价：
\begin{enumerate}
\item 分类态射 $\Spec(k) \to \Curvesstack$ 经
$\Curvesstack^{h0, 1}$ 分解；
\item $H^0(X, \mathcal{O}_X) = k$ 且 $\dim(X) = 1$。
\end{enumerate}
\end{enumerate}
\end{lemma}

\begin{proof}
给定曲线族 $X \to S$，由《空间的导出范畴》中的引理
\ref{spaces-perfect-lemma-jump-loci-geometric}，满足
$\kappa(s) = H^0(X_s, \mathcal{O}_{X_s})$ 的 $s \in S$
所成集合在 $S$ 中是开的。又由《空间的态射进阶》中的引理
\ref{spaces-more-morphisms-lemma-dimension-fibres-proper-flat}，
$S$ 中纤维维数为 $1$ 的点所成集合也是开的。
此外，若 $f : X \to S$ 是所有纤维维数均为 $1$ 的曲线族
（特别地，$f$ 是满射），则条件 (1)(b) 等价于对每个 $s \in S$ 都有
$\kappa(s) = H^0(X_s, \mathcal{O}_{X_s})$；见《空间的导出范畴》中的引理
\ref{spaces-perfect-lemma-proper-flat-h0}。
因此，该引理由第 \ref{moduli-curves-section-open} 节的一般讨论得出。
\end{proof}

\begin{lemma}
\label{moduli-curves-lemma-pre-genus-in-CM-1}
作为 $\Curvesstack$ 的开子叠，有
$\Curvesstack^{h0, 1} \subset \Curvesstack^{CM, 1}$。
\end{lemma}

\begin{proof}
见《代数曲线》中的引理 \ref{curves-lemma-automatic} 以及引理
\ref{moduli-curves-lemma-pre-genus-curves} 和 \ref{moduli-curves-lemma-CM-1-curves}。
\end{proof}

\begin{lemma}
\label{moduli-curves-lemma-genus}
设 $f : X \to S$ 是曲线族，并且对所有 $s \in S$ 都有
$\kappa(s) = H^0(X_s, \mathcal{O}_{X_s})$；换言之，分类态射
$S \to \Curvesstack$ 经 $\Curvesstack^{h0, 1}$ 分解
（引理 \ref{moduli-curves-lemma-pre-genus-curves}）。则：
\begin{enumerate}
\item $f_*\mathcal{O}_X = \mathcal{O}_S$，且该等式普遍成立；
\item $R^1f_*\mathcal{O}_X$ 是有限局部自由 $\mathcal{O}_S$-模；
\item 对任意态射 $h : S' \to S$，若 $f' : X' \to S'$ 是其基变换，
则 $h^*(R^1f_*\mathcal{O}_X) = R^1f'_*\mathcal{O}_{X'}$。
\end{enumerate}
\end{lemma}

\begin{proof}
应用《空间的导出范畴》中的引理
\ref{spaces-perfect-lemma-proper-flat-h0}。这证明了 (1)。
它还蕴含：在 $S$ 上局部可写成
$Rf_*\mathcal{O}_X = \mathcal{O}_S \oplus P$，其中 $P$
是 Tor 振幅位于 $[1, \infty)$ 的完美对象。回忆
$Rf_*\mathcal{O}_X$ 的构造与任意基变换交换
（《空间的导出范畴》中的引理
\ref{spaces-perfect-lemma-flat-proper-perfect-direct-image-general}）。
因此，对 $s \in S$，有
$$
H^i(P \otimes_{\mathcal{O}_S}^\mathbf{L} \kappa(s)) =
H^i(X_s, \mathcal{O}_{X_s})
\text{，其中 }i \geq 1
$$
由于 $X_s$ 是 $1$-维 Noether 概形，除非 $i = 1$，上式为零；见
《上同调》中的命题
\ref{cohomology-proposition-vanishing-Noetherian}。
于是 $P = H^1(P)[-1]$，并且 $H^1(P)$ 有限局部自由；例如可由
《代数进阶》中的引理
\ref{more-algebra-lemma-lift-perfect-from-residue-field} 得出。
由于所有构造都与基变换相容，(3) 也成立。
\end{proof}

\begin{lemma}
\label{moduli-curves-lemma-pre-genus-one-piece-per-genus}
存在开闭子叠的不交并分解
$$
\Curvesstack^{h0, 1} = \coprod\nolimits_{g \geq 0} \Curvesstack_g
$$
其中每个 $\Curvesstack_g$ 由下列条件刻画：
\begin{enumerate}
\item 给定曲线族 $f : X \to S$，下列条件等价：
\begin{enumerate}
\item 分类态射 $S \to \Curvesstack$ 经 $\Curvesstack_g$ 分解；
\item $f_*\mathcal{O}_X = \mathcal{O}_S$，该等式在任意基变换后仍成立，
$f$ 的纤维维数为 $1$，且 $R^1f_*\mathcal{O}_X$ 是秩为 $g$
的局部自由 $\mathcal{O}_S$-模；
\end{enumerate}
\item 给定域 $k$ 上的真概形 $X$，且 $\dim(X) \leq 1$，
下列条件等价：
\begin{enumerate}
\item 分类态射 $\Spec(k) \to \Curvesstack$ 经 $\Curvesstack_g$ 分解；
\item $\dim(X) = 1$、$k = H^0(X, \mathcal{O}_X)$，且 $X$
的亏格为 $g$。
\end{enumerate}
\end{enumerate}
\end{lemma}

\begin{proof}
由引理 \ref{moduli-curves-lemma-pre-genus-curves}，已经知道
$\Curvesstack^{h0, 1}$ 作为 $\Curvesstack$ 的开子叠存在，
并由本引理中不涉及 $R^1f_*$ 或 $H^1$ 的条件刻画。
由第 \ref{moduli-curves-section-open} 节的讨论和引理 \ref{moduli-curves-lemma-genus}，
立即得到开闭子叠分解的存在性。这证明了 (1) 中的刻画。
(2) 中的刻画由《代数曲线》中的定义
\ref{curves-definition-genus} 得出。
\end{proof}

\section{几何约化曲线}
\label{moduli-curves-section-geometrically-reduced}

\noindent
$\Curvesstack$ 有一个开子叠，参数化几何约化的 ``曲线''。

\begin{lemma}
\label{moduli-curves-lemma-geometrically-reduced-curves}
% P12-ERR-0249: see ../control/ERRATA_CANDIDATES.jsonl
存在开子叠 $\Curvesstack^{geomred} \subset \Curvesstack$，使得：
\begin{enumerate}
\item 给定曲线族 $X \to S$，下列条件等价：
\begin{enumerate}
\item 分类态射 $S \to \Curvesstack$ 经
$\Curvesstack^{geomred}$ 分解；
\item 态射 $X \to S$ 的纤维几何约化
（《空间的态射进阶》中的定义
\ref{spaces-more-morphisms-definition-geometrically-reduced-fibre}）；
\end{enumerate}
\item 给定域 $k$ 上的真概形 $X$，且 $\dim(X) \leq 1$，
下列条件等价：
\begin{enumerate}
\item 分类态射 $\Spec(k) \to \Curvesstack$ 经
$\Curvesstack^{geomred}$ 分解；
\item $X$ 在 $k$ 上几何约化。
\end{enumerate}
\end{enumerate}
\end{lemma}

\begin{proof}
设 $f : X \to S$ 是曲线族。由《空间的态射进阶》中的引理
\ref{spaces-more-morphisms-lemma-geometrically-reduced-open}，集合
$$
E = \{s \in S : \text{态射 }X \to S\text{ 在 }s
\text{ 处的纤维几何约化}\}
$$
在 $S$ 中是开的。由《空间的态射进阶》中的引理
\ref{spaces-more-morphisms-lemma-base-change-fibres-geometrically-reduced}，
这个开集的构造与任意基变换交换。因此，由第
\ref{moduli-curves-section-open} 节所讨论的方法，得到具有所需性质的开子叠。
\end{proof}

\begin{lemma}
\label{moduli-curves-lemma-geomred-in-CM}
作为 $\Curvesstack$ 的开子叠，有
$\Curvesstack^{geomred} \subset \Curvesstack^{CM}$。
\end{lemma}

\begin{proof}
这是因为维数 $\leq 1$ 的约化 Noether 概形是 Cohen--Macaulay 概形。
见《代数》中的引理 \ref{algebra-lemma-criterion-reduced}。
\end{proof}

\section{几何约化且连通的曲线}
\label{moduli-curves-section-geometrically-reduced-connected}

\noindent
$\Curvesstack$ 有一个开子叠，参数化几何约化且连通的 ``曲线''。
我们将立即排除 $0$-维对象。

\begin{lemma}
\label{moduli-curves-lemma-geometrically-reduced-connected-1-curves}
% P12-ERR-0250: see ../control/ERRATA_CANDIDATES.jsonl
存在开子叠 $\Curvesstack^{grc, 1} \subset \Curvesstack$，使得：
\begin{enumerate}
\item 给定曲线族 $X \to S$，下列条件等价：
\begin{enumerate}
\item 分类态射 $S \to \Curvesstack$ 经
$\Curvesstack^{grc, 1}$ 分解；
\item 态射 $X \to S$ 的几何纤维约化、连通且维数为 $1$；
\end{enumerate}
\item 给定域 $k$ 上的真概形 $X$，且 $\dim(X) \leq 1$，
下列条件等价：
\begin{enumerate}
\item 分类态射 $\Spec(k) \to \Curvesstack$ 经
$\Curvesstack^{grc, 1}$ 分解；
\item $X$ 几何约化、几何连通且维数为 $1$。
\end{enumerate}
\end{enumerate}
\end{lemma}

\begin{proof}
由引理 \ref{moduli-curves-lemma-geometrically-reduced-curves}、
\ref{moduli-curves-lemma-geomred-in-CM}、\ref{moduli-curves-lemma-CM-curves} 和
\ref{moduli-curves-lemma-CM-1-curves}，若所求开子叠存在，则显然有
$$
\Curvesstack^{grc, 1}
\subset
\Curvesstack^{geomred} \cap \Curvesstack^{CM, 1}
$$
设 $f : X \to S$ 是曲线族，其中 $f$ 是 Cohen--Macaulay 态射，
具有几何约化纤维且相对维数为 $1$。由《空间的态射进阶》中的引理
\ref{spaces-more-morphisms-lemma-stein-factorization-etale}，在 Stein 分解
$$
X \to T \to S
$$
中，态射 $T \to S$ 是 \'etale 态射。这蕴含存在开闭子概形
$S' \subset S$，使得 $X \times_S S' \to S'$ 具有几何连通纤维
（在《态射》中的引理 \ref{morphisms-lemma-finite-locally-free}
对有限局部自由态射 $T \to S$ 给出的分解中，这对应于 $S_1$）。
这个开集的构造与任意基变换交换，因为几何纤维的连通分支数在基变换下不变
（在当前特殊情形中，Stein 分解本身也与基变换交换，但得出结论并不需要这一点）。
因此，由第 \ref{moduli-curves-section-open} 节所讨论的方法，得到具有所需性质的开子叠。
\end{proof}

\begin{lemma}
\label{moduli-curves-lemma-geomredcon-in-h0-1}
作为 $\Curvesstack$ 的开子叠，有
$\Curvesstack^{grc, 1} \subset \Curvesstack^{h0, 1}$。
特别地，给定几何纤维约化、连通且维数为 $1$ 的曲线族
$f : X \to S$，则
% P12-ERR-0251: see ../control/ERRATA_CANDIDATES.jsonl
$R^1f_*\mathcal{O}_X$ 是有限局部自由 $\mathcal{O}_S$-模，
且其构造与任意基变换交换。
\end{lemma}

\begin{proof}
这由《簇》中的引理
\ref{varieties-lemma-proper-geometrically-reduced-global-sections} 以及引理
\ref{moduli-curves-lemma-pre-genus-curves} 和
\ref{moduli-curves-lemma-geometrically-reduced-connected-1-curves} 得出。
最后一个陈述由引理 \ref{moduli-curves-lemma-genus} 得出。
\end{proof}

\begin{lemma}
\label{moduli-curves-lemma-one-piece-per-genus}
存在开闭子叠的不交并分解
$$
\Curvesstack^{grc, 1} = \coprod\nolimits_{g \geq 0} \Curvesstack^{grc, 1}_g
$$
其中每个 $\Curvesstack^{grc, 1}_g$ 由下列条件刻画：
\begin{enumerate}
\item 给定曲线族 $f : X \to S$，下列条件等价：
\begin{enumerate}
\item 分类态射 $S \to \Curvesstack$ 经
$\Curvesstack^{grc, 1}_g$ 分解；
\item 态射 $f : X \to S$ 的几何纤维约化、连通且维数为 $1$，并且
$R^1f_*\mathcal{O}_X$ 是秩为 $g$ 的局部自由 $\mathcal{O}_S$-模；
\end{enumerate}
\item 给定域 $k$ 上的真概形 $X$，且 $\dim(X) \leq 1$，
下列条件等价：
\begin{enumerate}
\item 分类态射 $\Spec(k) \to \Curvesstack$ 经
$\Curvesstack^{grc, 1}_g$ 分解；
\item $X$ 几何约化、几何连通、维数为 $1$ 且亏格为 $g$。
\end{enumerate}
\end{enumerate}
\end{lemma}

\begin{proof}
第一种证明：令
$\Curvesstack^{grc, 1}_g = \Curvesstack^{grc, 1} \cap \Curvesstack_g$，
再结合引理 \ref{moduli-curves-lemma-geomredcon-in-h0-1} 和
\ref{moduli-curves-lemma-pre-genus-one-piece-per-genus}。
第二种证明：由第 \ref{moduli-curves-section-open} 节的讨论和引理
\ref{moduli-curves-lemma-geomredcon-in-h0-1}，立即得到开闭子叠分解的存在性。
这证明了 (1) 中的刻画。由于几何约化且连通的真 $1$-维概形 $X/k$
的亏格有定义（《代数曲线》中的定义 \ref{curves-definition-genus}
和《簇》中的引理
\ref{varieties-lemma-proper-geometrically-reduced-global-sections}），
且等于 $\dim_k H^1(X, \mathcal{O}_X)$，故 (2) 中的刻画也成立。
\end{proof}

\section{Gorenstein 曲线}
\label{moduli-curves-section-gorenstein}

\noindent
$\Curvesstack$ 有一个开子叠，参数化 Gorenstein ``曲线''。

\begin{lemma}
\label{moduli-curves-lemma-gorenstein-curves}
% P12-ERR-0252: see ../control/ERRATA_CANDIDATES.jsonl
存在开子叠 $\Curvesstack^{Gorenstein} \subset \Curvesstack$，使得：
\begin{enumerate}
\item 给定曲线族 $X \to S$，下列条件等价：
\begin{enumerate}
\item 分类态射 $S \to \Curvesstack$ 经
$\Curvesstack^{Gorenstein}$ 分解；
\item 态射 $X \to S$ 是 Gorenstein 态射；
\end{enumerate}
\item 给定域 $k$ 上的真概形 $X$，且 $\dim(X) \leq 1$，
下列条件等价：
\begin{enumerate}
\item 分类态射 $\Spec(k) \to \Curvesstack$ 经
$\Curvesstack^{Gorenstein}$ 分解；
\item $X$ 是 Gorenstein 概形。
\end{enumerate}
\end{enumerate}
\end{lemma}

\begin{proof}
设 $f : X \to S$ 是曲线族。由《空间的态射进阶》中的引理
\ref{spaces-more-morphisms-lemma-flat-finite-presentation-gorenstein-open}，集合
$$
W = \{x \in |X| : f \text{ 在下述点是 Gorenstein：}x\}
$$
在 $|X|$ 中是开的，并且这个开集的构造与任意基变换交换。
由于 $f$ 为真，$S$ 的子集
$$
S' = S \setminus f(|X| \setminus W)
$$
是开的，且 $X \times_S S' \to S'$ 是 Gorenstein 态射。
此外，$S'$ 的构造与任意基变换交换，因为 $W$ 的构造具有此性质。
% P12-ERR-0253: see ../control/ERRATA_CANDIDATES.jsonl
因此，由第 \ref{moduli-curves-section-open} 节所讨论的方法，得到具有所需性质的开子叠。
\end{proof}

\begin{lemma}
\label{moduli-curves-lemma-gorenstein-1-curves}
% P12-ERR-0254: see ../control/ERRATA_CANDIDATES.jsonl
存在开子叠
$\Curvesstack^{Gorenstein, 1} \subset \Curvesstack$，使得：
\begin{enumerate}
\item 给定曲线族 $X \to S$，下列条件等价：
\begin{enumerate}
\item 分类态射 $S \to \Curvesstack$ 经
$\Curvesstack^{Gorenstein, 1}$ 分解；
\item 态射 $X \to S$ 是 Gorenstein 态射，且具有相对维数 $1$
（《空间的态射》中的定义
\ref{spaces-morphisms-definition-relative-dimension}）；
\end{enumerate}
\item 给定域 $k$ 上的真概形 $X$，且 $\dim(X) \leq 1$，
下列条件等价：
\begin{enumerate}
\item 分类态射 $\Spec(k) \to \Curvesstack$ 经
$\Curvesstack^{Gorenstein, 1}$ 分解；
\item $X$ 是 Gorenstein 概形，且 $X$ 是维数为 $1$ 的等维概形。
\end{enumerate}
\end{enumerate}
\end{lemma}

\begin{proof}
回忆 Gorenstein 概形是 Cohen--Macaulay 概形
（《概形的对偶性》中的引理 \ref{duality-lemma-gorenstein-CM}），并且
% P12-ERR-0255: see ../control/ERRATA_CANDIDATES.jsonl
Gorenstein 态射是 Cohen--Macaulay 态射
（《概形的对偶性》中的引理 \ref{duality-lemma-gorenstein-CM-morphism}）。
因此，可以令 $\Curvesstack^{Gorenstein, 1}$ 等于
$\Curvesstack$ 内 $\Curvesstack^{Gorenstein}$ 与
$\Curvesstack^{CM, 1}$ 的交，并使用引理
\ref{moduli-curves-lemma-gorenstein-curves} 和 \ref{moduli-curves-lemma-CM-1-curves}。
\end{proof}

\section{局部完全交曲线}
\label{moduli-curves-section-lci}

\noindent
$\Curvesstack$ 有一个开子叠，参数化局部完全交 ``曲线''。

\begin{lemma}
\label{moduli-curves-lemma-lci-curves}
% P12-ERR-0256: see ../control/ERRATA_CANDIDATES.jsonl
存在开子叠 $\Curvesstack^{lci} \subset \Curvesstack$，使得：
\begin{enumerate}
\item 给定曲线族 $X \to S$，下列条件等价：
\begin{enumerate}
\item 分类态射 $S \to \Curvesstack$ 经 $\Curvesstack^{lci}$ 分解；
\item $X \to S$ 是局部完全交态射；
\item $X \to S$ 是 syntomic 态射。
\end{enumerate}
% P12-ERR-0257: see ../control/ERRATA_CANDIDATES.jsonl
\item 给定域 $k$ 上维数 $\leq 1$ 的真概形 $X$，下列条件等价：
\begin{enumerate}
\item 分类态射 $\Spec(k) \to \Curvesstack$ 经
$\Curvesstack^{lci}$ 分解；
\item $X$ 在 $k$ 上是局部完全交。
\end{enumerate}
\end{enumerate}
\end{lemma}

\begin{proof}
回忆 syntomic 态射等价于平坦的局部完全交态射；见
《空间的态射进阶》中的引理
\ref{spaces-more-morphisms-lemma-flat-lci}。因此 (1)(b) 等价于 (1)(c)。
在第 \ref{moduli-curves-section-open} 节中已经看到，只需证明：给定曲线族
$f : X \to S$，存在开子概形 $S' \subset S$，使得
$S' \times_S X \to S'$ 是局部完全交态射，并且 $S'$ 的构造
与任意基变换交换。这由更一般的《空间的态射进阶》中的引理
\ref{spaces-more-morphisms-lemma-where-lci} 得出。
\end{proof}

\section{具有孤立奇点的曲线}
\label{moduli-curves-section-curves-isolated}

\noindent
可以考察 $\Curvesstack$ 中参数化仅有有限多个奇点的 ``曲线''
所构成的开子叠（在我们的设定中，这些奇点可能对应于 $0$-维分支）。

\begin{lemma}
\label{moduli-curves-lemma-isolated-sings-curves}
% P12-ERR-0258: see ../control/ERRATA_CANDIDATES.jsonl
存在开子叠 $\Curvesstack^{+} \subset \Curvesstack$，使得：
\begin{enumerate}
\item 给定曲线族 $X \to S$，下列条件等价：
\begin{enumerate}
\item 分类态射 $S \to \Curvesstack$ 经 $\Curvesstack^{+}$ 分解；
\item 赋予 $X \to S$ 的奇异轨迹任意一个（或某一个）闭子空间结构后，
它在 $S$ 上有限。
\end{enumerate}
% P12-ERR-0259: see ../control/ERRATA_CANDIDATES.jsonl
\item 给定域 $k$ 上维数 $\leq 1$ 的真概形 $X$，下列条件等价：
\begin{enumerate}
\item 分类态射 $\Spec(k) \to \Curvesstack$ 经 $\Curvesstack^{+}$ 分解；
\item $X \to \Spec(k)$ 除有限多个点外处处光滑。
\end{enumerate}
\end{enumerate}
\end{lemma}

\begin{proof}
为证明该引理，只需证明：给定曲线族 $f : X \to S$，存在开子概形
$S' \subset S$，使得 $S' \times_S X \to S'$ 的纤维具有性质 (2)。
% P12-ERR-0260: see ../control/ERRATA_CANDIDATES.jsonl
（该开集的构造将自动与基变换交换。）按定义，$X \to S$
不光滑的点所成轨迹 $T \subset |X|$ 是闭的。令 $Z \subset X$
为在 $T$ 上赋予约化诱导代数空间结构所得到的闭子空间
（《空间的性质》中的定义
\ref{spaces-properties-definition-reduced-induced-space}）。
现在，若 $s \in S$ 是使 $Z_s$ 有限的点，则存在 $s$ 的开邻域
$U_s \subset S$，使得 $Z \cap f^{-1}(U_s) \to U_s$ 有限；见
《空间的态射进阶》中的引理
\ref{spaces-more-morphisms-lemma-proper-finite-fibre-finite-in-neighbourhood}。
这证明了该引理。
\end{proof}

\section{曲线叠的光滑轨迹}
\label{moduli-curves-section-smooth}

\noindent
态射
$$
\Curvesstack \longrightarrow \Spec(\mathbf{Z})
$$
在一个最大的开子叠上光滑；见《叠的态射》中的引理
\ref{stacks-morphisms-lemma-where-smooth}。我们希望给出曲线属于该轨迹的判据。
为此将使用少量形变理论。

\medskip\noindent
设 $k$ 是域，$X$ 是 $k$ 上维数 $\leq 1$ 的真概形。
为 $k$ 选取一个 Cohen 环 $\Lambda$；见《代数》中的引理
\ref{algebra-lemma-cohen-rings-exist}。于是处于《形变问题》中的例
\ref{examples-defos-example-schemes} 和引理
\ref{examples-defos-lemma-schemes-RS} 所描述的情形。因此，在剩余域为
$k$ 的 Artin 局部 $\Lambda$-代数所成范畴 $\mathcal{C}_\Lambda$ 上，
得到形变范畴 $\Deformationcategory_X$。

\begin{lemma}
\label{moduli-curves-lemma-in-smooth-locus}
在上述情形中，下列条件等价：
\begin{enumerate}
\item 分类态射 $\Spec(k) \to \Curvesstack$ 经
$\Curvesstack \to \Spec(\mathbf{Z})$ 光滑的开集分解；
\item 形变范畴 $\Deformationcategory_X$ 是无阻碍的。
\end{enumerate}
\end{lemma}

\begin{proof}
由于 $\Curvesstack \longrightarrow \Spec(\mathbf{Z})$ 局部有限表示
（引理 \ref{moduli-curves-lemma-curves-qs-lfp}），$\Curvesstack \longrightarrow
\Spec(\mathbf{Z})$ 光滑的开子叠之构造与平坦基变换交换
（《叠的态射》中的引理 \ref{stacks-morphisms-lemma-where-smooth}）。
由于 Cohen 环 $\Lambda$ 在 $\mathbf{Z}$ 上平坦，可以在 $\Lambda$ 上工作。
换言之，要证明
$$
\Lambda\text{-}\Curvesstack \longrightarrow \Spec(\Lambda)
$$
在由 $X/k$ 定义的点
$x_0 : \Spec(k) \to \Lambda\text{-}\Curvesstack$ 的一个开邻域上光滑，
当且仅当 $\Deformationcategory_X$ 无阻碍。

\medskip\noindent
该引理现在由《叠的几何》中的引理
\ref{stacks-geometry-lemma-characterize-smoothness} 和等式
$$
\Deformationcategory_X =
\mathcal{F}_{\Lambda\text{-}\Curvesstack, k, x_0}
$$
得出。建立这个等式并非完全平凡。事实上，左边是分类 $X$
作为 $A \in \Ob(\mathcal{C}_\Lambda)$ 上概形的所有平坦形变
$Y \to \Spec(A)$ 的形变范畴。右边是分类特殊纤维为 $X$
的所有平坦态射 $Y \to \Spec(A)$ 的形变范畴，其中 $Y$ 是代数空间，
且 $Y \to \Spec(A)$ 为真、有限表示并具有相对维数 $\leq 1$。
由于 $A$ 是 Artin 的，例如由《域上的空间》中的引理
\ref{spaces-over-fields-lemma-codim-1-point-in-schematic-locus} 可知
$Y$ 是概形。因此还需证明：$X$ 作为剩余域为 $k$ 的 Artin 局部环
$A$ 上概形的平坦形变 $Y \to \Spec(A)$ 为真、有限表示且相对维数
$\leq 1$。相对维数按纤维定义，故因为 $X/k$ 具有此性质，
$Y/A$ 也自动具有此性质。态射 $Y \to \Spec(A)$ 为真且局部有限表示，
因为 $X \to \Spec(k)$ 具有这些性质；见《态射进阶》中的引理
\ref{more-morphisms-lemma-deform-property}。
\end{proof}

\noindent
下面给出曲线叠中包含于光滑轨迹的一个 ``大'' 开子叠。

\begin{lemma}
\label{moduli-curves-lemma-big-smooth-part-curves}
开子叠
$$
\Curvesstack^{lci+} =
\Curvesstack^{lci} \cap \Curvesstack^{+}
\subset \Curvesstack
$$
具有以下性质：
\begin{enumerate}
\item $\Curvesstack^{lci+} \to \Spec(\mathbf{Z})$ 光滑；
\item 给定曲线族 $X \to S$，下列条件等价：
\begin{enumerate}
\item 分类态射 $S \to \Curvesstack$ 经 $\Curvesstack^{lci+}$ 分解；
\item $X \to S$ 是局部完全交态射，并且赋予 $X \to S$
的奇异轨迹任意一个（或某一个）闭子空间结构后，它在 $S$ 上有限；
\end{enumerate}
% P12-ERR-0261: see ../control/ERRATA_CANDIDATES.jsonl
\item 给定域 $k$ 上维数 $\leq 1$ 的真概形 $X$，下列条件等价：
\begin{enumerate}
\item 分类态射 $\Spec(k) \to \Curvesstack$ 经
$\Curvesstack^{lci+}$ 分解；
\item $X$ 在 $k$ 上是局部完全交，并且
$X \to \Spec(k)$ 除有限多个点外处处光滑。
\end{enumerate}
\end{enumerate}
\end{lemma}

\begin{proof}
若能证明存在开子叠 $\Curvesstack^{lci+}$，其点由 (2) 刻画，
则结合引理 \ref{moduli-curves-lemma-in-smooth-locus} 与《形变问题》中的引理
\ref{examples-defos-lemma-curve-isolated-lci}，可知 (1) 成立。由于
$$
\Curvesstack^{lci+} = \Curvesstack^{lci} \cap \Curvesstack^{+}
$$
在 $\Curvesstack$ 内成立，由引理 \ref{moduli-curves-lemma-lci-curves} 和
\ref{moduli-curves-lemma-isolated-sings-curves} 即得结论。
\end{proof}

\section{光滑曲线}
\label{moduli-curves-section-smooth-curves}

\noindent
本节研究 $\Curvesstack$ 中参数化光滑 ``曲线'' 的开子叠。

\begin{lemma}
\label{moduli-curves-lemma-smooth-curves}
% P12-ERR-0262: see ../control/ERRATA_CANDIDATES.jsonl
存在开子叠
$$
\Curvesstack^{smooth, 1} \subset \Curvesstack^{smooth} \subset \Curvesstack
$$
使得：
\begin{enumerate}
\item 给定曲线族 $f : X \to S$，下列条件等价：
\begin{enumerate}
\item 分类态射 $S \to \Curvesstack$ 经 $\Curvesstack^{smooth}$、
相应地，\ $\Curvesstack^{smooth, 1}$ 分解；
\item $f$ 光滑、相应地，\ 为相对维数 $1$ 的光滑态射；
\end{enumerate}
% P12-ERR-0263: see ../control/ERRATA_CANDIDATES.jsonl
\item 给定域 $k$ 上的真概形 $X$，且 $\dim(X) \leq 1$，
下列条件等价：
\begin{enumerate}
\item 分类态射 $\Spec(k) \to \Curvesstack$ 经 $\Curvesstack^{smooth}$、
相应地，\ $\Curvesstack^{smooth, 1}$ 分解；
\item $X$ 在 $k$ 上光滑；相应地，\ $X$ 在 $k$ 上光滑，且
$X$ 是维数为 $1$ 的等维概形。
\end{enumerate}
\end{enumerate}
\end{lemma}

\begin{proof}
为证明关于 $\Curvesstack^{smooth}$ 的陈述，只需证明：给定曲线族
$f : X \to S$，存在开子概形 $S' \subset S$，使得
$S' \times_S X \to S'$ 光滑，并且这个开集的构造与基变换交换。
已知存在最大的开集 $U \subset X$，使得 $U \to S$ 光滑，
而且 $U$ 的构造与任意基变换交换；见《空间的态射》中的引理
\ref{spaces-morphisms-lemma-where-smooth}。若
$T = |X| \setminus |U|$，则因为 $f$ 为真，$f(T)$ 在 $S$ 中闭。
令 $S' = S \setminus f(T)$，便得到所需开集。

\medskip\noindent
设 $f : X \to S$ 是曲线族，且 $f$ 光滑。则纤维 $X_s$
在 $\kappa(s)$ 上光滑，因而是 Cohen--Macaulay 概形
（例如可使用《代数》中的引理
\ref{algebra-lemma-smooth-over-field} 和
\ref{algebra-lemma-lci-CM} 看出）。因此，可以令
$$
\Curvesstack^{smooth, 1} = \Curvesstack^{smooth} \cap
\Curvesstack^{CM, 1}
$$
而所需等价性由已经对 $\Curvesstack^{smooth}$ 证明的结果以及引理
\ref{moduli-curves-lemma-CM-1-curves} 得出。
\end{proof}

\begin{lemma}
\label{moduli-curves-lemma-smooth-curves-smooth}
态射 $\Curvesstack^{smooth} \to \Spec(\mathbf{Z})$ 光滑。
\end{lemma}

\begin{proof}
由 $\Curvesstack^{smooth} \subset \Curvesstack^{lci+}$ 和引理
\ref{moduli-curves-lemma-big-smooth-part-curves} 立即得出。
\end{proof}

\begin{lemma}
\label{moduli-curves-lemma-smooth-curves-h0}
% P12-ERR-0264: see ../control/ERRATA_CANDIDATES.jsonl
存在开子叠 $\Curvesstack^{smooth, h0} \subset \Curvesstack$，使得：
\begin{enumerate}
\item 给定曲线族 $f : X \to S$，下列条件等价：
\begin{enumerate}
% P12-ERR-0265: see ../control/ERRATA_CANDIDATES.jsonl
\item 分类态射 $S \to \Curvesstack$ 经 $\Curvesstack^{smooth}$ 分解；
\item $f_*\mathcal{O}_X = \mathcal{O}_S$，该等式在任意基变换后仍成立，
且 $f$ 是相对维数为 $1$ 的光滑态射；
\end{enumerate}
% P12-ERR-0266: see ../control/ERRATA_CANDIDATES.jsonl
\item 给定域 $k$ 上的真概形 $X$，且 $\dim(X) \leq 1$，
下列条件等价：
\begin{enumerate}
\item 分类态射 $\Spec(k) \to \Curvesstack$ 经
$\Curvesstack^{smooth, h0}$ 分解；
\item $X$ 光滑，$\dim(X) = 1$，且 $k = H^0(X, \mathcal{O}_X)$；
\item $X$ 光滑，$\dim(X) = 1$，且 $X$ 几何连通；
\item $X$ 光滑，$\dim(X) = 1$，且 $X$ 几何整；
\item $X_{\overline{k}}$ 是光滑曲线。
\end{enumerate}
\end{enumerate}
\end{lemma}

\begin{proof}
若令
$$
\Curvesstack^{smooth, h0} = \Curvesstack^{smooth} \cap
\Curvesstack^{h0, 1}
$$
则由引理 \ref{moduli-curves-lemma-pre-genus-curves} 和
\ref{moduli-curves-lemma-smooth-curves} 可知 (1) 成立。事实上，这也给出
(2)(a) 与 (2)(b) 的等价性。为完成证明，需要说明 (2)(b)
分别等价于 (2)(c)、(2)(d) 和 (2)(e)。

\medskip\noindent
域上的光滑概形几何正规
（《簇》中的引理 \ref{varieties-lemma-smooth-geometrically-normal}），
光滑性在基变换下保持
（《态射》中的引理 \ref{morphisms-lemma-base-change-smooth}），
并且光滑性在目标上是 fpqc 局部的
（《下降》中的引理 \ref{descent-lemma-descending-property-smooth}）。
牢记这些事实，(2)(b)、(2)(c)、
% P12-ERR-0267: see ../control/ERRATA_CANDIDATES.jsonl
(2)(d) 和 (2)(e) 的等价性由《簇》中的引理
\ref{varieties-lemma-geometrically-normal-stein} 得出。
\end{proof}

\begin{definition}
\label{moduli-curves-definition-deligne-mumford-smooth}
\begin{reference}
\cite{DM}
\end{reference}
% P12-ERR-0269: see ../control/ERRATA_CANDIDATES.jsonl
我们用 $\mathcal{M}$ 表示引理 \ref{moduli-curves-lemma-smooth-curves-h0}
中引入的参数化曲线族的代数叠 $\Curvesstack^{smooth, h0}$，
并称其为{\it 光滑真曲线的模叠}。
对 $g \geq 0$，我们用 $\mathcal{M}_g$ 表示引理
\ref{moduli-curves-lemma-smooth-one-piece-per-genus} 中引入的代数叠，
并称其为{\it 亏格为 $g$ 的光滑真曲线的模叠}。
\end{definition}

\noindent
下面是必备引理。

\begin{lemma}
\label{moduli-curves-lemma-smooth-one-piece-per-genus}
存在开闭子叠的不交并分解
$$
\mathcal{M} = \coprod\nolimits_{g \geq 0} \mathcal{M}_g
$$
其中每个 $\mathcal{M}_g$ 由下列条件刻画：
\begin{enumerate}
\item 给定曲线族 $f : X \to S$，下列条件等价：
\begin{enumerate}
\item 分类态射 $S \to \Curvesstack$ 经 $\mathcal{M}_g$ 分解；
\item $X \to S$ 光滑，$f_*\mathcal{O}_X = \mathcal{O}_S$，
该等式在任意基变换后仍成立，且 $R^1f_*\mathcal{O}_X$
是秩为 $g$ 的局部自由 $\mathcal{O}_S$-模；
\end{enumerate}
% P12-ERR-0268: see ../control/ERRATA_CANDIDATES.jsonl
\item 给定域 $k$ 上的真概形 $X$，且 $\dim(X) \leq 1$，
下列条件等价：
\begin{enumerate}
\item 分类态射 $\Spec(k) \to \Curvesstack$ 经 $\mathcal{M}_g$ 分解；
\item $X$ 光滑，$\dim(X) = 1$，$k = H^0(X, \mathcal{O}_X)$，
且 $X$ 的亏格为 $g$；
\item $X$ 光滑，$\dim(X) = 1$，$X$ 几何连通，且 $X$ 的亏格为 $g$；
\item $X$ 光滑，$\dim(X) = 1$，$X$ 几何整，且 $X$ 的亏格为 $g$；
\item $X_{\overline{k}}$ 是亏格为 $g$ 的光滑曲线。
\end{enumerate}
\end{enumerate}
\end{lemma}

\begin{proof}
结合引理 \ref{moduli-curves-lemma-smooth-curves-h0} 和
\ref{moduli-curves-lemma-pre-genus-one-piece-per-genus}。也可以改用引理
\ref{moduli-curves-lemma-one-piece-per-genus}。
\end{proof}

\begin{lemma}
\label{moduli-curves-lemma-smooth-curves-h0-smooth}
态射 $\mathcal{M} \to \Spec(\mathbf{Z})$ 和
$\mathcal{M}_g \to \Spec(\mathbf{Z})$ 都光滑。
\end{lemma}

\begin{proof}
由于 $\mathcal{M}$ 是 $\Curvesstack^{lci+}$ 的开子叠，
该结论由引理 \ref{moduli-curves-lemma-big-smooth-part-curves} 得出。
\end{proof}

\section{光滑曲线的稠密性}
\label{moduli-curves-section-smooth-is-dense}

\noindent
本节标题有些误导，因为我们并不声称 $\Curvesstack^{smooth}$
在 $\Curvesstack$ 中稠密。事实上，这是错误的，Mumford 已在
\cite{PathologiesIV} 中证明这一点。不过，我们将看到光滑 ``曲线''
在一个大开集中稠密。

\begin{lemma}
\label{moduli-curves-lemma-smooth-dense}
包含关系
$$
|\Curvesstack^{smooth}| \subset |\Curvesstack^{lci+}|
$$
给出一个开稠密子集。
\end{lemma}

\begin{proof}
由《叠的性质》中的第 \ref{stacks-properties-section-points} 节所构造的
$|\Curvesstack^{lci+}|$ 上的拓扑，$|\Curvesstack^{smooth}|$
是开子集。设 $\xi \in |\Curvesstack^{lci+}|$ 是一点。
则存在域 $k$ 和 $k$ 上的概形 $X$，使得 $X$ 在 $k$ 上为真，
$\dim(X) \leq 1$，$X$ 在 $k$ 上是局部完全交，并且
% P12-ERR-0270: see ../control/ERRATA_CANDIDATES.jsonl
$X$ 在 $k$ 上除有限多个点外处处光滑，而且 $\xi$ 是由 $X$
确定的分类态射 $\Spec(k) \to \Curvesstack^{lci+}$ 的等价类。
见引理 \ref{moduli-curves-lemma-big-smooth-part-curves}。由《形变问题》中的引理
\ref{examples-defos-lemma-smoothing-proper-curve-isolated-lci}，
存在平坦射影态射 $Y \to \Spec(k[[t]])$，其泛纤维光滑，
其特殊纤维同构于 $X$。考虑由 $Y$ 确定的分类态射
$$
\Spec(k[[t]]) \longrightarrow \Curvesstack^{lci+}
$$
闭点的像为 $\xi$，而泛点的像属于 $|\Curvesstack^{smooth}|$。
由于泛点在 $|\Spec(k[[t]])|$ 中特化到闭点，故 $\xi$
属于 $|\Curvesstack^{smooth}|$ 的闭包，如所求。
\end{proof}

\section{结点曲线}
\label{moduli-curves-section-nodal-curves}

\noindent
结点曲线在代数几何中扮演特殊角色。建议读者简要浏览《代数曲线》中的
第 \ref{curves-section-nodal} 节和第
\ref{curves-section-families-nodal} 节，以及《空间的态射进阶》中的
第 \ref{spaces-more-morphisms-section-families-nodal} 节的部分讨论。

\begin{lemma}
\label{moduli-curves-lemma-nodal-curves}
% P12-ERR-0271: see ../control/ERRATA_CANDIDATES.jsonl
存在开子叠 $\Curvesstack^{nodal} \subset \Curvesstack$，使得：
\begin{enumerate}
\item 给定曲线族 $f : X \to S$，下列条件等价：
\begin{enumerate}
\item 分类态射 $S \to \Curvesstack$ 经 $\Curvesstack^{nodal}$ 分解；
\item $f$ 是相对维数为 $1$、至多有结点奇性的态射；
\end{enumerate}
% P12-ERR-0272: see ../control/ERRATA_CANDIDATES.jsonl
\item 给定域 $k$ 上的真概形 $X$，且 $\dim(X) \leq 1$，
下列条件等价：
\begin{enumerate}
\item 分类态射 $\Spec(k) \to \Curvesstack$ 经
$\Curvesstack^{nodal}$ 分解；
\item $X$ 的奇点至多为结点，且 $X$ 是维数为 $1$ 的等维概形。
\end{enumerate}
\end{enumerate}
\end{lemma}

\begin{proof}
事实上，只需证明：给定曲线族 $f : X \to S$，存在开子概形
$S' \subset S$，使得 $S' \times_S X \to S'$ 是相对维数为 $1$、
至多有结点奇性的态射，并且 $S'$ 的构造与任意基变换交换。
由《空间的态射进阶》中的引理
\ref{spaces-more-morphisms-lemma-locus-where-nodal}，存在最大的开子空间
$X' \subset X$，使得 $f|_{X'} : X' \to S$ 是相对维数为 $1$、
至多有结点奇性的态射。此外，$X'$ 的构造与基变换交换。
因此，可以取
$$
S' = S \setminus |f|(|X| \setminus |X'|)
$$
它是开的，因为真态射按定义普遍闭。
\end{proof}

\begin{lemma}
\label{moduli-curves-lemma-nodal-curves-smooth}
态射 $\Curvesstack^{nodal} \to \Spec(\mathbf{Z})$ 光滑。
\end{lemma}

\begin{proof}
由 $\Curvesstack^{nodal} \subset \Curvesstack^{lci+}$ 和引理
\ref{moduli-curves-lemma-big-smooth-part-curves} 立即得出。
\end{proof}

\section{相对对偶层}
\label{moduli-curves-section-relative-dualizing}

\noindent
本节主要用于在曲线族的情形引入记号。大部分工作已经在对偶性一章中完成。

\medskip\noindent
设 $f : X \to S$ 是曲线族。$D_\QCoh(\mathcal{O}_X)$ 中存在对象
$\omega_{X/S}^\bullet$，称为{\it 相对对偶复形}，具有如下性质：
对每个基变换图
$$
\xymatrix{
X_U \ar[d]_{f'} \ar[r]_{g'} & X \ar[d]^f \\
U \ar[r]^g & S
}
$$
若 $U = \Spec(A)$ 仿射，则复形
$\omega_{X_U/U}^\bullet = L(g')^*\omega_{X/S}^\bullet$ 表示函子
$$
D_\QCoh(\mathcal{O}_{X_U}) \longrightarrow \text{Mod}_A,\quad
K \longmapsto \Hom_U(Rf_*K, \mathcal{O}_U)
$$
更精确地，令 $(\omega_{X/S}^\bullet, \tau)$ 为《空间的对偶性》中的定义
\ref{spaces-duality-definition-relative-dualizing-proper-flat}
所定义的该曲线族的相对对偶复形。其存在性由《空间的对偶性》中的引理
\ref{spaces-duality-lemma-existence-relative-dualizing} 证明。
此外，$(\omega_{X/S}^\bullet, \tau)$ 的构造与任意基变换交换
（这本质上源于定义；精确参考是《空间的对偶性》中的引理
\ref{spaces-duality-lemma-base-change-relative-dualizing}）。
从现在起，我们将 $\omega_{X/S}^\bullet$ 的基变换与基变换后曲线族的
相对对偶复形等同，不再另行说明。

\medskip\noindent
设 $\{S_i \to S\}$ 是一个 \'etale 覆盖，其中 $S_i$ 仿射，且
$X_i = X \times_S S_i$ 是概形；见引理
\ref{moduli-curves-lemma-etale-locally-scheme}。由《空间的对偶性》中的引理
\ref{spaces-duality-lemma-compare}，$\omega_{X_i/S_i}^\bullet$
与《概形的对偶性》中的注
\ref{duality-remark-relative-dualizing-complex} 所讨论的概形间真、平坦、
有限表示态射 $f_i : X_i \to S_i$ 的相对对偶复形一致。
因此，为证明 $\omega_{X/S}^\bullet$ 的某个 \'etale 局部性质，
可以假设 $X \to S$ 是概形间态射，并使用概形对偶性一章中发展的理论。
更一般地，对 $X$ 的任意为概形的基变换，相对对偶复形都与
《概形的对偶性》中的注
\ref{duality-remark-relative-dualizing-complex} 的相对对偶复形一致。
从现在起，我们将使用这个等同，不再另行说明。

\medskip\noindent
特别地，设 $\Spec(k) \to S$ 是态射，其中 $k$ 是域。
% P12-ERR-0273: see ../control/ERRATA_CANDIDATES.jsonl
用 $X_k$ 表示该基变换（由《域上的空间》中的引理
\ref{spaces-over-fields-lemma-codim-1-point-in-schematic-locus}，
这是概形）。则 $\omega_{X_k/k}^\bullet$ 同构于《代数曲线》中的引理
\ref{curves-lemma-duality-dim-1} 的复形 $\omega_{X_k}^\bullet$
（两者表示同一个函子，故可使用 Yoneda 引理；而这实际上由上述评注得出）。
由此，余调层 $H^i(\omega_{X_k/k}^\bullet)$ 仅在 $i = 0, -1$
时可能非零。若 $X_k$ 是维数为 $1$ 的等维 Cohen--Macaulay 概形，
则仅有 $H^{-1}$；若 $X_k$ 还为 Gorenstein 概形，则
% P12-ERR-0274: see ../control/ERRATA_CANDIDATES.jsonl
$H^{-1}(\omega_{X_k/k})$ 可逆；见《代数曲线》中的引理
\ref{curves-lemma-duality-dim-1-CM} 和 \ref{curves-lemma-rr}。

\begin{lemma}
\label{moduli-curves-lemma-CM-dualizing}
设 $X \to S$ 是纤维为维数 $1$ 的等维 Cohen--Macaulay 概形的曲线族
（引理 \ref{moduli-curves-lemma-CM-1-curves}）。则
$\omega_{X/S}^\bullet = \omega_{X/S}[1]$，其中 $\omega_{X/S}$
是一个在 $S$ 上平坦的伪凝聚 $\mathcal{O}_X$-模，且其构造与任意基变换交换。
\end{lemma}

\begin{proof}
我们建议读者由上面对基变换到域后情形的讨论直接推出此结论。
我们的证明将使用一种稍显繁琐的归约，将问题化到 Noether 概形的情形。

\medskip\noindent
一旦证明 $\omega_{X/S}^\bullet = \omega_{X/S}[1]$ 且
$\omega_{X/S}$ 在 $S$ 上平坦，关于基变换的陈述就由
$\omega_{X/S}^\bullet$ 的构造与任意基变换交换这一已知事实得出。
此外，伪凝聚性将自动成立，因为按定义 $\omega_{X/S}^\bullet$
是伪凝聚的。其他余调层的消失与平坦性可在 \'etale 局部检验。
因此，可以假设 $f : X \to S$ 是概形间态射，且 $S$ 仿射
（见上面的讨论）。将 $S = \lim S_i$ 写成在 $\mathbf{Z}$ 上有限型的
仿射概形 $S_i$ 的余滤极限。由于 $\Curvesstack^{CM, 1}$ 在
$\mathbf{Z}$ 上局部有限表示（它是 $\Curvesstack$ 的开子叠；见引理
\ref{moduli-curves-lemma-CM-1-curves} 和 \ref{moduli-curves-lemma-curves-qs-lfp}），
可以找到一个 $i$ 以及曲线族 $X_i \to S_i$，其拉回为 $X \to S$
（《叠的极限》中的引理
\ref{stacks-limits-lemma-representable-by-spaces-limit-preserving}）。
必要时增大 $i$，可假设 $X_i$ 是概形；见《空间的极限》中的引理
\ref{spaces-limits-lemma-limit-is-scheme}。由于
$\omega_{X/S}^\bullet$ 的构造与任意基变换交换，可以用 $S_i$ 代替 $S$。
这样便可以并且确实假设 $S_i$ 是 Noether 概形。此时，由我们对纤维的假设，
$f$ 显然是 Cohen--Macaulay 态射
（《态射进阶》中的定义 \ref{more-morphisms-definition-CM}）。
此外，由《概形的对偶性》中的第
\ref{duality-section-upper-shriek} 节对 $f^!$ 的构造，有
$\omega_{X/S}^\bullet = f^!\mathcal{O}_S$。
% P12-ERR-0275: see ../control/ERRATA_CANDIDATES.jsonl
因此，该引理由《概形的对偶性》中的引理
\ref{duality-lemma-affine-flat-Noetherian-CM} 得出。
\end{proof}

\begin{definition}
\label{moduli-curves-definition-relative-dualizing-sheaf}
设 $f : X \to S$ 是纤维为维数 $1$ 的等维 Cohen--Macaulay 概形的曲线族
（引理 \ref{moduli-curves-lemma-CM-1-curves}）。则引理
\ref{moduli-curves-lemma-CM-dualizing} 中研究的 $\mathcal{O}_X$-模
$$
\omega_{X/S} = H^{-1}(\omega_{X/S}^\bullet)
$$
称为 $f$ 的{\it 相对对偶层}。
\end{definition}

\noindent
在定义 \ref{moduli-curves-definition-relative-dualizing-sheaf} 的情形，
相对对偶层 $\omega_{X/S}$ 具有以下性质（而且该性质在 $S$ 上局部刻画它）：
对每个基变换图
$$
\xymatrix{
X_U \ar[d]_{f'} \ar[r]_{g'} & X \ar[d]^f \\
U \ar[r]^g & S
}
$$
若 $U = \Spec(A)$ 仿射，则模
$\omega_{X_U/U} = (g')^*\omega_{X/S}$ 表示函子
$$
\QCoh(\mathcal{O}_{X_U}) \longrightarrow \text{Mod}_A,\quad
% P12-ERR-0276: see ../control/ERRATA_CANDIDATES.jsonl
\mathcal{F} \longmapsto \Hom_A(H^1(X, \mathcal{F}), A)
$$
这由上面给出的相对对偶复形的相应性质立即得出。特别地，若 $A = k$
是域，则恢复《代数曲线》中的引理
\ref{curves-lemma-duality-dim-1}、
\ref{curves-lemma-duality-dim-1-CM} 和 \ref{curves-lemma-rr}
所引入并研究的 $X_k$ 的对偶模。

\begin{lemma}
\label{moduli-curves-lemma-gorenstein-dualizing}
设 $X \to S$ 是纤维为维数 $1$ 的等维 Gorenstein 概形的曲线族
（引理 \ref{moduli-curves-lemma-gorenstein-1-curves}）。则相对对偶层
$\omega_{X/S}$ 是可逆 $\mathcal{O}_X$-模，且其构造与任意基变换交换。
\end{lemma}

\begin{proof}
这是因为由上面的讨论，相对对偶模到任一纤维的拉回均可逆。
或者，也可以完全仿照引理 \ref{moduli-curves-lemma-CM-dualizing} 的证明，
再由《概形的对偶性》中的引理
\ref{duality-lemma-affine-flat-Noetherian-gorenstein} 得到结论。
\end{proof}

\section{预稳定曲线}
\label{moduli-curves-section-prestable-curves}

\noindent
下述定义等价于通常公认的预稳定曲线族概念。

\begin{definition}
\label{moduli-curves-definition-prestable}
设 $f : X \to S$ 是曲线族。若满足下列条件，便称 $f$ 是一个
{\it 预稳定曲线族}：
\begin{enumerate}
\item $f$ 是相对维数为 $1$、至多有结点奇性的态射；并且
\item $f_*\mathcal{O}_X = \mathcal{O}_S$，而且该等式在任意基变换后仍成立
\footnote{事实上，只要求 $f_*\mathcal{O}_X = \mathcal{O}_S$ 即可，
因为此时 $f$ 的 Stein 分解是 \'etale 的；见《空间的态射进阶》中的引理
\ref{spaces-more-morphisms-lemma-stein-factorization-etale}。
也可以把该条件替换为要求几何纤维连通；见引理
\ref{moduli-curves-lemma-geomredcon-in-h0-1}。}。
\end{enumerate}
\end{definition}

\noindent
设 $X$ 是域 $k$ 上的真概形，且 $\dim(X) \leq 1$。
则 $X \to \Spec(k)$ 是曲线族，因此可以问它在上述定义的意义下是否预稳定
\footnote{此处不能使用“预稳定曲线”这一术语，因为“曲线”蕴含不可约性。
见《代数曲线》中的第 \ref{curves-section-families-nodal} 节的讨论。}。
展开定义可见，下列条件等价：
\begin{enumerate}
\item $X$ 预稳定；
\item $X$ 的奇点至多为结点，$\dim(X) = 1$，且
$k = H^0(X, \mathcal{O}_X)$；
\item $X_{\overline{k}}$ 连通，并且除有限多个结点外，
它在 $\overline{k}$ 上光滑
（《代数曲线》中的定义 \ref{curves-definition-multicross}）。
\end{enumerate}
这说明我们的定义与文献中的大多数定义一致。

\begin{lemma}
\label{moduli-curves-lemma-prestable-curves}
% P12-ERR-0277: see ../control/ERRATA_CANDIDATES.jsonl
存在开子叠 $\Curvesstack^{prestable} \subset \Curvesstack$，使得：
\begin{enumerate}
\item 给定曲线族 $f : X \to S$，下列条件等价：
\begin{enumerate}
\item 分类态射 $S \to \Curvesstack$ 经 $\Curvesstack^{prestable}$ 分解；
\item $X \to S$ 是预稳定曲线族；
\end{enumerate}
% P12-ERR-0278: see ../control/ERRATA_CANDIDATES.jsonl
\item 给定域 $k$ 上的真概形 $X$，且 $\dim(X) \leq 1$，
下列条件等价：
\begin{enumerate}
\item 分类态射 $\Spec(k) \to \Curvesstack$ 经
$\Curvesstack^{prestable}$ 分解；
\item $X$ 的奇点至多为结点，$\dim(X) = 1$，且
$k = H^0(X, \mathcal{O}_X)$。
\end{enumerate}
\end{enumerate}
\end{lemma}

\begin{proof}
给定曲线族 $X \to S$，它预稳定，当且仅当分类态射同时经
$\Curvesstack^{nodal}$ 和 $\Curvesstack^{h0, 1}$ 分解。
也可以使用 $\Curvesstack^{grc, 1}$（因为结点曲线几何约化，
所以它的 $H^0$ 等于基域，当且仅当它连通）。用公式表示即
$$
\Curvesstack^{prestable} =
\Curvesstack^{nodal} \cap \Curvesstack^{h0, 1} =
\Curvesstack^{nodal} \cap \Curvesstack^{grc, 1}
$$
因此，该引理由引理 \ref{moduli-curves-lemma-pre-genus-curves} 和
\ref{moduli-curves-lemma-nodal-curves} 得出。
\end{proof}

\noindent
对每个亏格 $g \geq 0$，都有分类亏格为 $g$ 的预稳定曲线的代数叠。
事实上，从现在起，我们称 $X \to S$ 是一个{\it 亏格为 $g$
的预稳定曲线族}，当且仅当分类态射 $S \to \Curvesstack$
经引理 \ref{moduli-curves-lemma-prestable-one-piece-per-genus} 的开子叠
$\Curvesstack^{prestable}_g$ 分解。

\begin{lemma}
\label{moduli-curves-lemma-prestable-one-piece-per-genus}
存在开闭子叠的分解
$$
\Curvesstack^{prestable} = \coprod\nolimits_{g \geq 0}
\Curvesstack^{prestable}_g
$$
其中每个 $\Curvesstack^{prestable}_g$ 由下列性质刻画：
\begin{enumerate}
\item 给定曲线族 $f : X \to S$，下列条件等价：
\begin{enumerate}
\item 分类态射 $S \to \Curvesstack$ 经
$\Curvesstack^{prestable}_g$ 分解；
\item $X \to S$ 是预稳定曲线族，且 $R^1f_*\mathcal{O}_X$
是秩为 $g$ 的局部自由 $\mathcal{O}_S$-模；
\end{enumerate}
\item 给定域 $k$ 上的真概形 $X$，且 $\dim(X) \leq 1$，
下列条件等价：
\begin{enumerate}
\item 分类态射 $\Spec(k) \to \Curvesstack$ 经
$\Curvesstack^{prestable}_g$ 分解；
\item $X$ 的奇点至多为结点，$\dim(X) = 1$，
$k = H^0(X, \mathcal{O}_X)$，且 $X$ 的亏格为 $g$。
\end{enumerate}
\end{enumerate}
\end{lemma}

\begin{proof}
我们已经看到 $\Curvesstack^{prestable}$ 包含于
$\Curvesstack^{h0, 1}$，故结论由引理
\ref{moduli-curves-lemma-prestable-curves} 和
\ref{moduli-curves-lemma-pre-genus-one-piece-per-genus} 得出。
\end{proof}

\begin{lemma}
\label{moduli-curves-lemma-prestable-curves-smooth}
态射 $\Curvesstack^{prestable} \to \Spec(\mathbf{Z})$ 和
$\Curvesstack^{prestable}_g \to \Spec(\mathbf{Z})$ 光滑。
\end{lemma}

\begin{proof}
由于 $\Curvesstack^{prestable}$ 是 $\Curvesstack^{nodal}$ 的开子叠，
结论由引理 \ref{moduli-curves-lemma-nodal-curves-smooth} 得出。
\end{proof}

\section{半稳定曲线}
\label{moduli-curves-section-semistable-curves}

\noindent
下述引理将帮助我们理解半稳定曲线族。

\begin{lemma}
\label{moduli-curves-lemma-semistable}
设 $f : X \to S$ 是亏格为 $g \geq 1$ 的预稳定曲线族。
设 $s \in S$ 是基概形上的一点，并设 $m \geq 2$。下列条件等价：
\begin{enumerate}
\item $X_s$ 没有有理尾
（《代数曲线》中的例 \ref{curves-example-rational-tail}）；
\item 对某个开子集 $s \in U \subset S$，态射
% P12-ERR-0279: see ../control/ERRATA_CANDIDATES.jsonl
$f^*f_*\omega_{X/S}^{\otimes m} \to \omega_{X/S}^{\otimes m}$
在 $f^{-1}(U)$ 上满射。
\end{enumerate}
\end{lemma}

\begin{proof}
假设 (2)。利用第 \ref{moduli-curves-section-relative-dualizing} 节的内容，
可知 $\omega_{X_s}^{\otimes m}$ 整体生成。然而，若
$C \subset X_s$ 是有理尾，则由《代数曲线》中的引理
\ref{curves-lemma-rational-tail-negative}，有
$\deg(\omega_{X_s}|_C) < 0$，故由《簇》中的引理
\ref{varieties-lemma-check-invertible-sheaf-trivial}，有
$H^0(C, \omega_{X_s}|_C) = 0$；这与它整体生成矛盾。
这证明了 (1)。

\medskip\noindent
假设 (1)。先假设 $g \geq 2$。由《代数曲线》中的引理
\ref{curves-lemma-contracting-rational-tails}，假设 (1) 蕴含
$\omega_{X_s}^{\otimes m}$ 整体生成。此外，由对偶性，有
$$
\Hom_{\kappa(s)}(H^1(X_s, \omega_{X_s}^{\otimes m}), \kappa(s)) =
H^0(X_s, \omega_{X_s}^{\otimes 1 - m})
$$
见《代数曲线》中的引理 \ref{curves-lemma-duality-dim-1-CM}。
由于 $\omega_{X_s}^{\otimes m}$ 整体生成，它在每个不可约分支上的限制
% P12-ERR-0280: see ../control/ERRATA_CANDIDATES.jsonl
次数非负。因此，$\omega_{X_s}^{\otimes 1 - m}$ 在每个不可约分支上的
限制次数非正。又由 Riemann--Roch（《代数曲线》中的引理
\ref{curves-lemma-rr}），有
$\deg(\omega_{X_s}^{\otimes 1 - m}) = (1 - m)(2g - 2) < 0$，
故由《簇》中的引理 \ref{varieties-lemma-no-sections-dual-nef}，
可知该 $H^0$ 为零。由上同调与基变换可知
$$
E = Rf_*\omega_{X/S}^{\otimes m}
$$
是完美复形，且其构造与任意基变换交换
（《空间的导出范畴》中的引理
\ref{spaces-perfect-lemma-flat-proper-perfect-direct-image-general}）。
上面证明的消失说明 $E \otimes^\mathbf{L} \kappa(s)$ 等于置于次数 $0$
的 $H^0(X_s, \omega_{X_s}^{\otimes m})$。缩小 $S$ 后可得，
$E = f_*\omega_{X/S}^{\otimes m}$ 是置于次数 $0$ 的局部自由
$\mathcal{O}_S$-模（而且如前所述，其构造与任意基变换交换）；见
《空间的导出范畴》中的引理
\ref{spaces-perfect-lemma-open-where-cohomology-in-degree-i-rank-r-geometric}。
态射 $f^*f_*\omega_{X/S}^{\otimes m} \to \omega_{X/S}^{\otimes m}$
限制到 $X_s$ 后满射。因此，它在 $X_s$ 的某个开邻域上满射。
由于 $f$ 是真态射，该开邻域包含 $f^{-1}(U)$，其中 $U$
是 $S$ 中 $s$ 的某个开邻域。

\medskip\noindent
假设 (1) 且 $g = 1$。由《代数曲线》中的引理
\ref{curves-lemma-contracting-rational-tails}，假设 (1) 意味着
$\omega_{X_s}$ 同构于 $\mathcal{O}_{X_s}$。若能证明缩小 $S$ 后
% P12-ERR-0281: see ../control/ERRATA_CANDIDATES.jsonl
可逆层 $\omega_{X/S}$ 变为平凡层，结论便成立。可以假设 $S$ 仿射。
进一步缩小 $S$ 后，可以写成
$$
Rf_*\mathcal{O}_X = (\mathcal{O}_S \xrightarrow{0} \mathcal{O}_S)
$$
它位于次数 $0$ 和 $1$，且与进一步的基变换相容；见引理
\ref{moduli-curves-lemma-genus}。由对偶性，这意味着
$$
Rf_*\omega_{X/S} = (\mathcal{O}_S \xrightarrow{0} \mathcal{O}_S)
$$
位于次数 $0$ 和 $1$\footnote{使用
$Rf_*\omega_{X/S}^\bullet =
% P12-ERR-0282: see ../control/ERRATA_CANDIDATES.jsonl
Rf_*R\SheafHom_{\mathcal{O}_X}(\mathcal{O}_X. \omega_{X/S}^\bullet) =
R\SheafHom_{\mathcal{O}_S}(Rf_*\mathcal{O}_X, \mathcal{O}_S)$；
这由《空间的对偶性》中的引理
\ref{spaces-duality-lemma-iso-on-RSheafHom} 和注
\ref{spaces-duality-remark-iso-on-RSheafHom} 得出；然后再使用由第
\ref{moduli-curves-section-relative-dualizing} 节中的定义得到的
$\omega_{X/S}^\bullet = \omega_{X/S}[1]$。}。
特别地，可以得到与基变换相容的同构
$\mathcal{O}_S \to f_*\omega_{X/S}$，因为
$Rf_*\omega_{X/S}$ 的构造与基变换相容（见上面给出的参考）。
由伴随性，得到整体截面 $\sigma \in \Gamma(X, \omega_{X/S})$。
该截面在纤维 $X_s$ 上的限制非零（事实上是一个基元素）；而且由于
$\omega_{X_s}$ 在该纤维上平凡，该截面在 $X_s$ 上处处非零。
% P12-ERR-0283: see ../control/ERRATA_CANDIDATES.jsonl
% P12-ERR-0284: see ../control/ERRATA_CANDIDATES.jsonl
因此，它在 $X_s$ 的某个开邻域上处处非零。由于 $f$ 是真态射，
该开邻域包含 $f^{-1}(U)$，其中 $U$ 是 $S$ 中 $s$ 的某个开邻域。
\end{proof}

\noindent
受引理 \ref{moduli-curves-lemma-semistable} 启发，我们作如下定义。

\begin{definition}
\label{moduli-curves-definition-semistable}
设 $f : X \to S$ 是曲线族。若
\begin{enumerate}
\item $X \to S$ 是预稳定曲线族；并且
\item 对所有 $s \in S$，$X_s$ 的亏格均为 $\geq 1$，且没有有理尾，
\end{enumerate}
则称 $f$ 是一个{\it 半稳定曲线族}。
\end{definition}

\noindent
特别地，亏格为 $0$ 的预稳定曲线族绝不半稳定。
设 $X$ 是域 $k$ 上的真概形，且 $\dim(X) \leq 1$。
则 $X \to \Spec(k)$ 是曲线族，因而可以问它是否半稳定。
展开定义可见，下列条件等价：
\begin{enumerate}
\item $X$ 半稳定；
\item $X$ 预稳定，亏格为 $\geq 1$，且没有有理尾；
\item $X_{\overline{k}}$ 连通，除有限多个结点外在 $\overline{k}$ 上光滑，
亏格为 $\geq 1$，并且没有同构于 $\mathbf{P}^1_{\overline{k}}$
且仅在一点与 $X_{\overline{k}}$ 的其余部分相交的不可约分支。
\end{enumerate}
要看出 (2) 与 (3) 等价，使用《代数曲线》中的引理
\ref{curves-lemma-contracting-rational-tails}：$X$ 没有有理尾，
当且仅当 $X_{\overline{k}}$ 没有有理尾。
这说明我们的定义与文献中的大多数定义一致。

\begin{lemma}
\label{moduli-curves-lemma-semistable-curves}
% P12-ERR-0285: see ../control/ERRATA_CANDIDATES.jsonl
存在开子叠 $\Curvesstack^{semistable} \subset \Curvesstack$，使得：
\begin{enumerate}
\item 给定曲线族 $f : X \to S$，下列条件等价：
\begin{enumerate}
\item 分类态射 $S \to \Curvesstack$ 经
$\Curvesstack^{semistable}$ 分解；
\item $X \to S$ 是半稳定曲线族；
\end{enumerate}
% P12-ERR-0286: see ../control/ERRATA_CANDIDATES.jsonl
\item 给定域 $k$ 上的真概形 $X$，且 $\dim(X) \leq 1$，
下列条件等价：
\begin{enumerate}
\item 分类态射 $\Spec(k) \to \Curvesstack$ 经
$\Curvesstack^{semistable}$ 分解；
\item $X$ 的奇点至多为结点，$\dim(X) = 1$，
$k = H^0(X, \mathcal{O}_X)$，$X$ 的亏格为 $\geq 1$，且
$X$ 没有有理尾；
\item $X$ 的奇点至多为结点，$\dim(X) = 1$，
$k = H^0(X, \mathcal{O}_X)$，而且对 $m \geq 2$，
% P12-ERR-0287: see ../control/ERRATA_CANDIDATES.jsonl
$\omega_{X_s}^{\otimes m}$ 整体生成。
\end{enumerate}
\end{enumerate}
\end{lemma}

\begin{proof}
(2)(b) 与 (2)(c) 的等价性是《代数曲线》中的引理
\ref{curves-lemma-contracting-rational-tails}。
在证明的其余部分，我们将按照定义 \ref{moduli-curves-definition-semistable} 使用 (2)(b)。

\medskip\noindent
由第 \ref{moduli-curves-section-open} 节的讨论，只需考虑预稳定曲线族
$f : X \to S$。由引理 \ref{moduli-curves-lemma-semistable}，得到所求轨迹的开性。
该开集的构造与任意基变换交换，因为由《代数曲线》中的引理
\ref{curves-lemma-contracting-rational-tails}，有理尾的存在或不存在
不受基域扩张影响。
\end{proof}

\begin{lemma}
\label{moduli-curves-lemma-semistable-one-piece-per-genus}
存在开闭子叠的分解
$$
\Curvesstack^{semistable} = \coprod\nolimits_{g \geq 1}
\Curvesstack^{semistable}_g
$$
其中每个 $\Curvesstack^{semistable}_g$ 由下列性质刻画：
\begin{enumerate}
\item 给定曲线族 $f : X \to S$，下列条件等价：
\begin{enumerate}
\item 分类态射 $S \to \Curvesstack$ 经
$\Curvesstack^{semistable}_g$ 分解；
\item $X \to S$ 是半稳定曲线族，且 $R^1f_*\mathcal{O}_X$
是秩为 $g$ 的局部自由 $\mathcal{O}_S$-模；
\end{enumerate}
% P12-ERR-0288: see ../control/ERRATA_CANDIDATES.jsonl
\item 给定域 $k$ 上的真概形 $X$，且 $\dim(X) \leq 1$，
下列条件等价：
\begin{enumerate}
\item 分类态射 $\Spec(k) \to \Curvesstack$ 经
$\Curvesstack^{semistable}_g$ 分解；
\item $X$ 的奇点至多为结点，$\dim(X) = 1$，
$k = H^0(X, \mathcal{O}_X)$，$X$ 的亏格为 $g$，且 $X$ 没有有理尾；
\item $X$ 的奇点至多为结点，$\dim(X) = 1$，
$k = H^0(X, \mathcal{O}_X)$，$X$ 的亏格为 $g$，而且对 $m \geq 2$，
% P12-ERR-0289: see ../control/ERRATA_CANDIDATES.jsonl
$\omega_{X_s}^{\otimes m}$ 整体生成。
\end{enumerate}
\end{enumerate}
\end{lemma}

\begin{proof}
合并引理 \ref{moduli-curves-lemma-semistable-curves} 和
\ref{moduli-curves-lemma-prestable-one-piece-per-genus} 即得。
\end{proof}

\begin{lemma}
\label{moduli-curves-lemma-semistable-curves-smooth}
态射 $\Curvesstack^{semistable} \to \Spec(\mathbf{Z})$ 和
$\Curvesstack^{semistable}_g \to \Spec(\mathbf{Z})$ 光滑。
\end{lemma}

\begin{proof}
由于 $\Curvesstack^{semistable}$ 是 $\Curvesstack^{nodal}$ 的开子叠，
结论由引理 \ref{moduli-curves-lemma-nodal-curves-smooth} 得出。
\end{proof}

\section{稳定曲线}
\label{moduli-curves-section-stable-curves}

\noindent
下述引理将帮助我们理解稳定曲线族。

\begin{lemma}
\label{moduli-curves-lemma-stable}
设 $f : X \to S$ 是亏格为 $g \geq 2$ 的预稳定曲线族，
并设 $s \in S$ 是基概形上的一点。下列条件等价：
\begin{enumerate}
\item $X_s$ 没有有理尾，也没有有理桥
（《代数曲线》中的例 \ref{curves-example-rational-tail} 和
\ref{curves-example-rational-bridge}）；
\item 对某个开子集 $s \in U \subset S$，$\omega_{X/S}$
在 $f^{-1}(U)$ 上丰沛。
\end{enumerate}
\end{lemma}

\begin{proof}
假设 (2)。则 $\omega_{X_s}$ 在 $X_s$ 上丰沛。
由《代数曲线》中的引理 \ref{curves-lemma-rational-tail-negative} 和
\ref{curves-lemma-rational-bridge-zero}，可知 (1) 成立
（还使用《簇》中的引理
\ref{varieties-lemma-ampleness-in-terms-of-degrees-components}
对丰沛可逆层的刻画）。

\medskip\noindent
假设 (1)。则由《代数曲线》中的引理
% P12-ERR-0290: see ../control/ERRATA_CANDIDATES.jsonl
\ref{curves-lemma-contracting-rational-bridges}，$\omega_{X_s}$
在 $X_s$ 上丰沛。结论由《空间的下降》中的引理
\ref{spaces-descent-lemma-ample-in-neighbourhood} 得出。
\end{proof}

\noindent
受引理 \ref{moduli-curves-lemma-stable} 启发，我们作如下定义。

\begin{definition}
\label{moduli-curves-definition-stable}
设 $f : X \to S$ 是曲线族。若
\begin{enumerate}
\item $X \to S$ 是预稳定曲线族；并且
% P12-ERR-0291: see ../control/ERRATA_CANDIDATES.jsonl
\item 对所有 $s \in S$，$X_s$ 的亏格均为 $\geq 2$，
且没有有理尾或有理桥，
\end{enumerate}
则称 $f$ 是一个{\it 稳定曲线族}。
\end{definition}

\noindent
特别地，亏格为 $0$ 或 $1$ 的预稳定曲线族绝不稳定。
设 $X$ 是域 $k$ 上的真概形，且 $\dim(X) \leq 1$。
则 $X \to \Spec(k)$ 是曲线族，因而可以问它是否稳定。
展开定义可见，下列条件等价：
\begin{enumerate}
\item $X$ 稳定；
\item $X$ 预稳定，亏格为 $\geq 2$，没有有理尾，也没有有理桥；
\item $X$ 几何连通，除有限多个结点外在 $k$ 上光滑，且
$\omega_X$ 丰沛。
\end{enumerate}
要看出 (2) 与 (3) 等价，使用上面的引理 \ref{moduli-curves-lemma-stable}。
这说明我们的定义与文献中的大多数定义一致。

\begin{lemma}
\label{moduli-curves-lemma-stable-curves}
% P12-ERR-0292: see ../control/ERRATA_CANDIDATES.jsonl
存在开子叠 $\Curvesstack^{stable} \subset \Curvesstack$，使得：
\begin{enumerate}
\item 给定曲线族 $f : X \to S$，下列条件等价：
\begin{enumerate}
\item 分类态射 $S \to \Curvesstack$ 经 $\Curvesstack^{stable}$ 分解；
\item $X \to S$ 是稳定曲线族；
\end{enumerate}
% P12-ERR-0293: see ../control/ERRATA_CANDIDATES.jsonl
\item 给定域 $k$ 上的真概形 $X$，且 $\dim(X) \leq 1$，
下列条件等价：
\begin{enumerate}
\item 分类态射 $\Spec(k) \to \Curvesstack$ 经
$\Curvesstack^{stable}$ 分解；
\item $X$ 的奇点至多为结点，$\dim(X) = 1$，
$k = H^0(X, \mathcal{O}_X)$，$X$ 的亏格为 $\geq 2$，
且 $X$ 没有有理尾或有理桥；
\item $X$ 的奇点至多为结点，$\dim(X) = 1$，
$k = H^0(X, \mathcal{O}_X)$，且
% P12-ERR-0294: see ../control/ERRATA_CANDIDATES.jsonl
$\omega_{X_s}$ 丰沛。
\end{enumerate}
\end{enumerate}
\end{lemma}

\begin{proof}
由第 \ref{moduli-curves-section-open} 节的讨论，只需考虑预稳定曲线族
$f : X \to S$。由引理 \ref{moduli-curves-lemma-stable}，得到所求轨迹的开性。
该开集的构造与任意基变换交换：一方面，由《代数曲线》中的引理
\ref{curves-lemma-contracting-rational-tails} 和
\ref{curves-lemma-contracting-rational-bridges}，有理尾或有理桥的
存在或不存在不受基域扩张影响；另一方面，由《下降》中的引理
\ref{descent-lemma-descending-property-ample}，丰沛性也不受基域扩张影响。
\end{proof}

\begin{definition}
\label{moduli-curves-definition-deligne-mumford}
\begin{reference}
\cite{DM}
\end{reference}
% P12-ERR-0295: see ../control/ERRATA_CANDIDATES.jsonl
我们用 $\overline{\mathcal{M}}$ 表示引理 \ref{moduli-curves-lemma-stable-curves}
中引入的、参数化稳定曲线族的代数叠 $\Curvesstack^{stable}$，
并称它为{\it 稳定曲线模叠}。对 $g \geq 2$，我们用
$\overline{\mathcal{M}}_g$ 表示引理
\ref{moduli-curves-lemma-stable-one-piece-per-genus} 中引入的代数叠，
并称它为{\it 亏格为 $g$ 的稳定曲线模叠}。
\end{definition}

\noindent
下面是必不可少的引理。

\begin{lemma}
\label{moduli-curves-lemma-stable-one-piece-per-genus}
存在开闭子叠的分解
$$
\overline{\mathcal{M}} = \coprod\nolimits_{g \geq 2} \overline{\mathcal{M}}_g
$$
其中每个 $\overline{\mathcal{M}}_g$ 由下列性质刻画：
\begin{enumerate}
\item 给定曲线族 $f : X \to S$，下列条件等价：
\begin{enumerate}
\item 分类态射 $S \to \Curvesstack$ 经
$\overline{\mathcal{M}}_g$ 分解；
\item $X \to S$ 是稳定曲线族，且 $R^1f_*\mathcal{O}_X$
是秩为 $g$ 的局部自由 $\mathcal{O}_S$-模；
\end{enumerate}
% P12-ERR-0296: see ../control/ERRATA_CANDIDATES.jsonl
\item 给定域 $k$ 上的真概形 $X$，且 $\dim(X) \leq 1$，
下列条件等价：
\begin{enumerate}
\item 分类态射 $\Spec(k) \to \Curvesstack$ 经
$\overline{\mathcal{M}}_g$ 分解；
\item $X$ 的奇点至多为结点，$\dim(X) = 1$，
$k = H^0(X, \mathcal{O}_X)$，$X$ 的亏格为 $g$，且
$X$ 没有有理尾或有理桥；
\item $X$ 的奇点至多为结点，$\dim(X) = 1$，
$k = H^0(X, \mathcal{O}_X)$，$X$ 的亏格为 $g$，且
% P12-ERR-0297: see ../control/ERRATA_CANDIDATES.jsonl
$\omega_{X_s}$ 丰沛。
\end{enumerate}
\end{enumerate}
\end{lemma}

\begin{proof}
合并引理 \ref{moduli-curves-lemma-stable-curves} 和
\ref{moduli-curves-lemma-prestable-one-piece-per-genus} 即得。
\end{proof}

\begin{lemma}
\label{moduli-curves-lemma-stable-curves-smooth}
态射 $\overline{\mathcal{M}} \to \Spec(\mathbf{Z})$ 和
$\overline{\mathcal{M}}_g \to \Spec(\mathbf{Z})$ 光滑。
\end{lemma}

\begin{proof}
由于 $\overline{\mathcal{M}}$ 是 $\Curvesstack^{nodal}$ 的开子叠，
结论由引理 \ref{moduli-curves-lemma-nodal-curves-smooth} 得出。
\end{proof}

\begin{lemma}
\label{moduli-curves-lemma-stable-curves-deligne-mumford}
叠 $\overline{\mathcal{M}}$ 和 $\overline{\mathcal{M}}_g$
是 $\Curvesstack^{DM}$ 的开子叠。特别地，
$\overline{\mathcal{M}}$ 和 $\overline{\mathcal{M}}_g$ 既是 DM 叠
（《叠的态射》中的定义
\ref{stacks-morphisms-definition-absolute-separated}），
也是 Deligne--Mumford 叠
（《代数叠》中的定义 \ref{algebraic-definition-deligne-mumford}）。
\end{lemma}

\begin{proof}
先证第一个断言。设 $X$ 是域 $k$ 上的真概形，其奇点至多为结点，
$\dim(X) = 1$，$k = H^0(X, \mathcal{O}_X)$，$X$ 的亏格为
$\geq 2$，而且 $X$ 没有有理尾或有理桥。需要证明分类态射
$\Spec(k) \to \overline{\mathcal{M}} \to \Curvesstack$
经 $\Curvesstack^{DM}$ 分解。可以先用代数闭包替换 $k$
（因为已经知道相关的叠是代数叠 $\Curvesstack$ 的开子叠）。
由引理 \ref{moduli-curves-lemma-stable-curves}、\ref{moduli-curves-lemma-DM-curves} 和
\ref{moduli-curves-lemma-in-DM-locus-vector-fields}，只需证明
$\text{Der}_k(\mathcal{O}_X, \mathcal{O}_X) = 0$。
这已在《代数曲线》中的引理
\ref{curves-lemma-stable-vector-fields} 中证明。

\medskip\noindent
由于 $\Curvesstack^{DM}$ 是 $\Curvesstack$ 中为 DM 的最大开子叠，
可知 $\Curvesstack^{DM}$ 的开子叠 $\overline{\mathcal{M}}$
也具有此性质。最后，由《叠的态射》中的定理
\ref{stacks-morphisms-theorem-DM}，DM 代数叠是 Deligne--Mumford 叠。
\end{proof}

\begin{lemma}
\label{moduli-curves-lemma-smooth-dense-in-stable}
设 $g \geq 2$。包含关系
$$
|\mathcal{M}_g| \subset |\overline{\mathcal{M}}_g|
$$
把左边实现为右边的开稠密子集。
\end{lemma}

\begin{proof}
由于 $\overline{\mathcal{M}}_g \subset \Curvesstack^{lci+}$ 是开的，
且 $\Curvesstack^{smooth} \cap \overline{\mathcal{M}}_g = \mathcal{M}_g$，
结论立即由引理 \ref{moduli-curves-lemma-smooth-dense} 得出。
\end{proof}

\section{收缩态射}
\label{moduli-curves-section-contracting}

\noindent
我们建议读者在继续本节之前，先熟悉《代数曲线》中的第
\ref{curves-section-contracting-rational-tails}、
\ref{curves-section-contracting-rational-bridges} 和
\ref{curves-section-contracting-to-stable} 节。本节的主要结果是“稳定化”态射
$$
\Curvesstack^{prestable}_g
\longrightarrow
\overline{\mathcal{M}}_g
$$
的存在性；见引理 \ref{moduli-curves-lemma-stabilization-morphism}。粗略地说，
该态射把亏格为 $g$ 的结点曲线的模点送到相伴稳定曲线的模点；
后者构造于《代数曲线》中的引理
\ref{curves-lemma-characterize-contraction-to-stable}。

\begin{lemma}
\label{moduli-curves-lemma-contract}
设 $S$ 是概形，$s \in S$ 是一点。设 $f : X \to S$ 和
$g : Y \to S$ 是曲线族，并设 $c : X \to Y$ 是 $S$ 上的态射。
若 $c_{s, *}\mathcal{O}_{X_s} = \mathcal{O}_{Y_s}$ 且
$R^1c_{s, *}\mathcal{O}_{X_s} = 0$，则以 $s$ 的一个开邻域替换 $S$ 后，
有 $\mathcal{O}_Y = c_*\mathcal{O}_X$ 和
$R^1c_*\mathcal{O}_X = 0$，而且经任意态射 $S' \to S$ 基变换后
这些等式仍成立。
\end{lemma}

\begin{proof}
设 $(U, u) \to (S, s)$ 是一个 \'etale 邻域，使得
$\mathcal{O}_{Y_U} = (X_U \to Y_U)_*\mathcal{O}_{X_U}$ 且
$R^1(X_U \to Y_U)_*\mathcal{O}_{X_U} = 0$，并且经任意
$U' \to U$ 基变换后也有相同结论。然后用 $U \to S$ 的开像替换 $S$。
给定 $S' \to S$，置 $U' = U \times_S S'$，便得到 \'etale 覆盖
$\{U' \to S'\}$ 和 $\{Y_{U'} \to Y_{S'}\}$。
因此，$c$ 经 $S' \to S$ 基变换后命题成立，可由
$X_U \to Y_U$ 经 $U' \to U$ 基变换后命题成立推出。
换言之，该问题在 $S$ 的 \'etale 拓扑下是局部的。
因此，由引理 \ref{moduli-curves-lemma-etale-locally-scheme}，可以假设 $X$ 和 $Y$
是概形。由《态射进阶》中的引理
\ref{more-morphisms-lemma-h1-fibre-zero-isom}，存在包含 $Y_s$
的开子概形 $V \subset Y$，使得
$c_*\mathcal{O}_X|_V = \mathcal{O}_V$ 且
$R^1c_*\mathcal{O}_X|_V = 0$，并且经任意 $S' \to S$ 基变换后
这些等式仍成立。由于 $g : Y \to S$ 是真态射，可以找到 $s$
的开邻域 $U \subset S$，使得 $g^{-1}(U) \subset V$。这个 $U$ 即可。
\end{proof}

\begin{lemma}
\label{moduli-curves-lemma-contract-basic-uniqueness}
设 $S$ 是概形，$s \in S$ 是一点。设 $f : X \to S$ 和
$g_i : Y_i \to S$（$i = 1, 2$）是曲线族，并设
$c_i : X \to Y_i$ 是 $S$ 上的态射。假设存在与 $c_{1, s}$ 和
$c_{2, s}$ 相容的纤维同构 $Y_{1, s} \cong Y_{2, s}$。
若 $c_{1, s, *}\mathcal{O}_{X_s} = \mathcal{O}_{Y_{1, s}}$ 且
$R^1c_{1, s, *}\mathcal{O}_{X_s} = 0$，则存在 $s$ 的开邻域 $U$
以及 $U$ 上曲线族的同构 $Y_{1, U} \cong Y_{2, U}$；该同构与给定的
纤维同构以及 $c_1$、$c_2$ 相容。
\end{lemma}

\begin{proof}
回顾 $\mathcal{O}_{S, s} = \colim \mathcal{O}_S(U)$，其中余极限取遍
$s$ 的仿射邻域 $U$ 所成的系统。因此，局部环上有限表示代数空间的范畴
是 $s$ 的仿射邻域上有限表示代数空间范畴的余极限。见《空间的极限》中的引理
\ref{spaces-limits-lemma-descend-finite-presentation}。
这样便归约到 $S$ 是某个局部环的谱且 $s$ 是闭点的情形。

\medskip\noindent
假设 $S = \Spec(A)$，其中 $A$ 是局部环，$s$ 是闭点。
将 $A = \colim A_j$ 写成局部 Noether 环 $A_j$ 的余极限
（例如可取在 $\mathbf{Z}$ 上本质有限型），其过渡同态均为局部同态。
置 $S_j = \Spec(A_j)$，其闭点为 $s_j$。可以找到一个 $j$ 以及曲线族
$X_j \to S_j$、$Y_{j, i} \to S_j$；见引理
\ref{moduli-curves-lemma-curves-qs-lfp} 和《叠的极限》中的引理
\ref{stacks-limits-lemma-representable-by-spaces-limit-preserving}。
必要时增大 $j$，可以找到态射 $c_{j, i} : X_j \to Y_{j, i}$，
其到 $s$ 的基变换为 $c_i$；见《空间的极限》中的引理
\ref{spaces-limits-lemma-descend-finite-presentation}。
由于 $\kappa(s) = \colim \kappa(s_j)$，同样可以假设存在与
$c_{j, 1, s_j}$ 和 $c_{j, 2, s_j}$ 相容的同构
$Y_{j, 1, s_j} \cong Y_{j, 2, s_j}$。最后，假设
$c_{1, s, *}\mathcal{O}_{X_s} = \mathcal{O}_{Y_{1, s}}$ 和
$R^1c_{1, s, *}\mathcal{O}_{X_s} = 0$ 可由 $c_{j, 1, s_j}$ 继承，因为
% P12-ERR-0298: see ../control/ERRATA_CANDIDATES.jsonl
$\{s_j \to s\}$ 是 fpqc 覆盖，而 $c_{1, s}$ 是
$c_{j, 1, s_j}$ 经该覆盖的基变换（细节略去）。
这样便把引理归约到下一段所讨论的情形。

\medskip\noindent
假设 $S$ 是 Noether 局部环 $\Lambda$ 的谱，$s$ 是闭点。
考虑态射
$$
(c_1, c_2) : X \longrightarrow Y_1 \times_S Y_2
$$
的概形论像 $Z$。引理的陈述等价于断言：$Z$ 经投影态射同构地映到
$Y_1$ 和 $Y_2$。由于取该态射的概形论像与平坦基变换交换
% P12-ERR-0299: see ../control/ERRATA_CANDIDATES.jsonl
（《空间的态射》中的引理
\ref{spaces-morphisms-lemma-flat-base-change-scheme-theoretic-image}，
可以用 $\Lambda$ 的完备化替换它
（《代数进阶》中的第 \ref{more-algebra-section-permanence-completion} 节）。

\medskip\noindent
假设 $S$ 是完备 Noether 局部环 $\Lambda$ 的谱。
注意，此时 $X$、$Y_1$、$Y_2$ 都是概形
（《空间的态射进阶》中的引理
\ref{spaces-more-morphisms-lemma-projective-over-complete}）。
% P12-ERR-0300: see ../control/ERRATA_CANDIDATES.jsonl
用 $X_n$、$Y_{1, n}$、$Y_{2, n}$ 表示 $X$、$Y_1$、$Y_2$
到 $\Spec(\Lambda/\mathfrak m^{n + 1})$ 的基变换。回顾箭头
$$
\Deformationcategory_{X_s \to Y_{2, s}} \cong
\Deformationcategory_{X_s \to Y_{1, s}} \longrightarrow
\Deformationcategory_{X_s}
$$
是等价；见《形变问题》中的引理
\ref{examples-defos-lemma-schemes-morphisms-smooth-to-base}。
因此，在 $\Deformationcategory_{X_s \to Y_{1, s}}$ 中存在形式对象的同构
$(X_n \to Y_{1, n}) \cong (X_n \to Y_{2, n})$。
最后，由 Grothendieck 代数化定理（《概形的上同调》中的引理
\ref{coherent-lemma-algebraize-morphism}），这给出与 $c_1$ 和 $c_2$
相容的同构 $Y_1 \to Y_2$。
\end{proof}

\begin{lemma}
\label{moduli-curves-lemma-contract-basic}
设 $f : X \to S$ 是曲线族，$s \in S$ 是一点。设
$h_0 : X_s \to Y_0$ 是到 $\kappa(s)$ 上真概形 $Y_0$ 的态射，
满足 $h_{0, *}\mathcal{O}_{X_s} = \mathcal{O}_{Y_0}$ 且
$R^1h_{0, *}\mathcal{O}_{X_s} = 0$。则存在基本 \'etale 邻域
$(U, u) \to (S, s)$、曲线族 $Y \to U$ 以及 $U$ 上的态射
$h : X_U \to Y$，使得它在 $u$ 上的纤维同构于 $h_0$。
\end{lemma}

\begin{proof}
先作一些归约；建议读者跳到后面。该问题在 $S$ 上是局部的，
故可以假设 $S$ 仿射。将 $S = \lim S_i$ 写成在 $\mathbf{Z}$ 上
有限型仿射概形 $S_i$ 的余滤极限。对某个 $i$，可以找到曲线族
$X_i \to S_i$，其基变换为 $X \to S$。这由引理
\ref{moduli-curves-lemma-curves-qs-lfp} 和《叠的极限》中的引理
\ref{stacks-limits-lemma-representable-by-spaces-limit-preserving} 得出。
设 $s_i \in S_i$ 是 $s$ 的像。注意
$\kappa(s) = \colim \kappa(s_i)$，而且 $X_s$ 是概形
（《域上的空间》中的引理
\ref{spaces-over-fields-lemma-codim-1-point-in-schematic-locus}）。
增大 $i$ 后，可以假设存在 $\kappa(s_i)$ 上有限型概形之间的态射
$h_{i, 0} : X_{i, s_i} \to Y_i$，其到 $\kappa(s)$ 的基变换为 $h_0$；
见《极限》中的引理 \ref{limits-lemma-descend-finite-presentation}。
再次增大 $i$ 后，可以假设 $Y_i$ 在 $\kappa(s_i)$ 上为真；
见《极限》中的引理 \ref{limits-lemma-eventually-proper}。
% P12-ERR-0301: see ../control/ERRATA_CANDIDATES.jsonl
设 $g_{i, 0} : Y_0 \to Y_{i, 0}$ 为投影。注意，它作为
$\Spec(\kappa(s)) \to \Spec(\kappa(s_i))$ 的基变换，是忠实平坦态射。
由平坦基变换，有
$$
h_{0, *}\mathcal{O}_{X_s} = g_{i, 0}^*h_{i, 0, *}\mathcal{O}_{X_{i, s_i}}
\quad\text{且}\quad
% P12-ERR-0302: see ../control/ERRATA_CANDIDATES.jsonl
R^1h_{0, *}\mathcal{O}_{X_s} = g_{i, 0}^*Rh_{i, 0, *}\mathcal{O}_{X_{i, s_i}}
$$
见《概形的上同调》中的引理
\ref{coherent-lemma-flat-base-change-cohomology}。
由忠实平坦性可知，$X_i \to S_i$、$s_i \in S_i$ 和
% P12-ERR-0303: see ../control/ERRATA_CANDIDATES.jsonl
$X_{i, s_i} \to Y_i$ 满足该引理的全部假设。
这便归约到下一段所讨论的情形。

\medskip\noindent
假设 $S$ 在 $\mathbf{Z}$ 上仿射且有限型。设
$\mathcal{O}_{S, s}^h$ 是 $S$ 在 $s$ 处局部环的 Hensel 化。
由《代数进阶》中的引理 \ref{more-algebra-lemma-henselization-G-ring}
和命题 \ref{more-algebra-proposition-ubiquity-G-ring}，
$\mathcal{O}_{S, s}^h$ 是 G-环。假设能够构造曲线族
$Y' \to \Spec(\mathcal{O}_{S, s}^h)$ 以及
$\Spec(\mathcal{O}_{S, s}^h)$ 上的态射
$$
h' : X \times_S \Spec(\mathcal{O}_{S, s}^h) \longrightarrow Y'
$$
其到闭点的基变换为 $h_0$。这已经足够。事实上，首先使用
$$
\mathcal{O}_{S, s}^h = \colim_{(U, u)} \mathcal{O}_U(U)
$$
其中余极限取遍基本 \'etale 邻域所成的滤过范畴
（《态射进阶》中的引理
\ref{more-morphisms-lemma-describe-henselization}）。
其次，再次使用如下事实：给定 $Y'$ 后，可以把它下降为某个 $U$ 上的
$Y \to U$（见上面给出的参考）。然后使用《极限》中的引理
\ref{limits-lemma-descend-finite-presentation}，把 $h'$ 下降为某个 $h$。
这便归约到下一段所讨论的情形。

\medskip\noindent
假设 $S = \Spec(\Lambda)$，其中
$(\Lambda, \mathfrak m, \kappa)$ 是 Hensel Noether 局部 G-环，
而 $s$ 是 $S$ 的闭点。回顾映射
$$
\Deformationcategory_{X_s \to Y_0} \to \Deformationcategory_{X_s}
$$
是等价；见《形变问题》中的引理
\ref{examples-defos-lemma-schemes-morphisms-smooth-to-base}。
（这是证明中唯一重要的一步；其余全是技巧。）
% P12-ERR-0304: see ../control/ERRATA_CANDIDATES.jsonl
用 $\Lambda^\wedge$ 表示 $\mathfrak m$-进完备化。
$X$ 到 $\Lambda/\mathfrak m^{n + 1}$ 的拉回 $X_n$
在 $\Lambda^\wedge$ 上定义了 $\Deformationcategory_{X_s}$ 的形式对象
$\xi$。由上述等价，得到 $\Lambda^\wedge$ 上
$\Deformationcategory_{X_s \to Y_0}$ 的形式对象 $\xi'$。
因此得到一个巨大的交换图
$$
\xymatrix{
\ldots \ar[r] &
X_n \ar[r] \ar[d] &
X_{n - 1} \ar[r] \ar[d] &
\ldots \ar[r] &
X_s \ar[d] \\
\ldots \ar[r] &
Y_n \ar[r] \ar[d] &
Y_{n - 1} \ar[r] \ar[d] &
\ldots \ar[r] &
Y_0 \ar[d] \\
\ldots \ar[r] &
\Spec(\Lambda/\mathfrak m^{n + 1}) \ar[r] &
\Spec(\Lambda/\mathfrak m^n) \ar[r] &
\ldots \ar[r] &
\Spec(\kappa)
}
$$
由《Quot》中的引理 \ref{quot-lemma-curves-existence}，
形式对象 $(Y_n)$ 来自曲线族 $Y' \to \Spec(\Lambda^\wedge)$。
由《空间的态射进阶》中的引理
\ref{spaces-more-morphisms-lemma-algebraize-morphism}，得到态射
$h' : X_{\Lambda^\wedge} \to Y'$，它对所有 $n$ 都诱导给定态射
$X_n \to Y_n$，特别地诱导给定态射 $X_s \to Y_0$。

\medskip\noindent
最后作一个标准的代数化/逼近论证。首先注意，可以找到有限生成
$\Lambda$-子代数 $\Lambda \subset A \subset \Lambda^\wedge$、
曲线族 $Y'' \to \Spec(A)$ 以及 $A$ 上的态射
$h'' : X_A \to Y''$，其到 $\Lambda^\wedge$ 的基变换为 $h'$。
这是因为 $\Lambda^\wedge$ 是这些环 $A$ 的滤过余极限；可以像前面那样，
利用 $\Curvesstack$ 局部有限表示这一事实（由《叠的极限》中的引理
\ref{stacks-limits-lemma-representable-by-spaces-limit-preserving}
得到 $A$ 上的 $Y''$），并使用《空间的极限》中的引理
\ref{spaces-limits-lemma-descend-finite-presentation}，
把 $h'$ 下降为某个 $h''$。然后，可以使用 G-环的逼近性质
（采用《环同态的光滑化》中的定理
\ref{smoothing-theorem-approximation-property} 的形式），
找到映射 $A \to \Lambda$，它所诱导的 $A \to \kappa$
与 $A \to \Lambda^\wedge$ 所诱导的映射相同。
把 $h''$ 基变换到 $\Lambda$，证明即告完成。
\end{proof}

\begin{lemma}
\label{moduli-curves-lemma-contract-prestable-to-stable}
设 $f : X \to S$ 是亏格为 $g \geq 2$ 的预稳定曲线族。
$f$ 存在分解 $X \to Y \to S$，其中
% P12-ERR-0305: see ../control/ERRATA_CANDIDATES.jsonl
$g : Y \to S$ 是稳定曲线族，而 $c : X \to Y$ 具有下列性质：
\begin{enumerate}
\item $\mathcal{O}_Y = c_*\mathcal{O}_X$ 且
$R^1c_*\mathcal{O}_X = 0$，并且经任意态射 $S' \to S$
基变换后这些等式仍成立；
\item 对任意 $s \in S$，态射 $c_s : X_s \to Y_s$ 是
《代数曲线》中的第 \ref{curves-section-contracting-to-stable} 节
所讨论的有理尾和有理桥的收缩。
\end{enumerate}
此外，$c : X \to Y$ 在唯一同构意义下唯一。
\end{lemma}

\begin{proof}
设 $s \in S$。令 $c_0 : X_s \to Y_0$ 为《代数曲线》中的第
\ref{curves-section-contracting-to-stable} 节的收缩
（更精确地说，见《代数曲线》中的引理
\ref{curves-lemma-characterize-contraction-to-stable}）。
由引理 \ref{moduli-curves-lemma-contract-basic}，存在基本 \'etale 邻域
$(U, u)$ 以及 $U$ 上曲线族之间的态射 $c : X_U \to Y$，
其在 $u$ 处的纤维恢复 $c_0$。由于 $\omega_{Y_0}$ 丰沛，
必要时缩小 $U$ 后，由引理 \ref{moduli-curves-lemma-stable-curves} 和
\ref{moduli-curves-lemma-stable-one-piece-per-genus} 所蕴含的开性可知，
$Y \to U$ 是亏格为 $g$ 的稳定曲线族。必要时再次缩小 $U$ 后，
引理 \ref{moduli-curves-lemma-contract} 给出 $c : X_U \to Y$ 的断言 (1)。
此外，(2) 由《代数曲线》中的引理
\ref{curves-lemma-characterize-contraction-to-stable} 的唯一性成立。
由此可知，引理中的态射 $c$ 在 $S$ 上 \'etale 局部存在。
更精确地，存在 \'etale 覆盖 $\{U_i \to S\}$ 以及 $U_i$ 上的态射
$c_i : X_{U_i} \to Y_i$，其中 $Y_i \to U_i$ 是具有引理中性质
(1) 和 (2) 的稳定曲线族。

\medskip\noindent
为完成证明，只需证明 $c : X \to Y$ 的唯一性（在唯一同构意义下）。
事实上，一旦完成这一步，便得到同构
% P12-ERR-0306: see ../control/ERRATA_CANDIDATES.jsonl
$$
\varphi_{ij} :
Y_i \times_{U_i} (U_i \times_S U_j)
\longrightarrow
Y_i \times_{U_j} (U_i \times_S U_j)
$$
% P12-ERR-0307: see ../control/ERRATA_CANDIDATES.jsonl
它们在 $U_i \times U_j \times U_k$ 上满足上链条件（由唯一性）。
% P12-ERR-0308: see ../control/ERRATA_CANDIDATES.jsonl
由于 $\overline{\mathcal{M}_g}$ 是代数叠，下降数据有效，
从而得到 $Y \to S$。态射 $c_i$ 下降为 $S$ 上的态射
$c : X \to Y$。最后，$c$ 的性质 (1) 和 (2) 立即由 $c_i$
的性质 (1) 和 (2) 得出。

\medskip\noindent
最后，若
% P12-ERR-0309: see ../control/ERRATA_CANDIDATES.jsonl
$c_1 : X \to Y_i$（$i = 1, 2$）是到 $S$ 上两个
% P12-ERR-0310: see ../control/ERRATA_CANDIDATES.jsonl
稳定曲线族的态射，并满足 (1) 和 (2)，则由引理
\ref{moduli-curves-lemma-contract-basic-uniqueness}，至少在 $S$ 上局部，
得到与 $c_1$ 和 $c_2$ 相容的态射 $Y_1 \to Y_2$。
略去验证这些态射唯一的过程（提示：这由 $c_1$ 的概形论像是
$Y_1$ 这一事实得出）。因此，这些局部给出的态射可以粘合，证明完成。
\end{proof}

\begin{lemma}
\label{moduli-curves-lemma-stabilization-morphism}
设 $g \geq 2$。存在 $\mathbf{Z}$ 上代数叠之间的态射
% P12-ERR-0311: see ../control/ERRATA_CANDIDATES.jsonl
$$
stabilization :
\Curvesstack^{prestable}_g
\longrightarrow
\overline{\mathcal{M}}_g
$$
它把亏格为 $g$ 的预稳定曲线族 $X \to S$ 送到引理
\ref{moduli-curves-lemma-contract-prestable-to-stable} 中与之相伴的稳定曲线族
$Y \to S$。
\end{lemma}

\begin{proof}
为说明这一点，只需检验引理
\ref{moduli-curves-lemma-contract-prestable-to-stable} 的构造与基变换相容
（也与同构相容，但这是立即的）；见《叠的性质》中的第
\ref{stacks-properties-section-conventions} 节所引入的关于代数叠的
（语言滥用）约定。为此，只需检验引理
\ref{moduli-curves-lemma-contract-prestable-to-stable} 的性质 (1) 和 (2)
在基变换下稳定。对 (1)，这立即清楚。对 (2)，一方面，
《代数曲线》中的引理 \ref{curves-lemma-contracting-rational-tails} 和
\ref{curves-lemma-contracting-rational-bridges} 的收缩在基域扩张下稳定；
另一方面，《代数曲线》中的引理
\ref{curves-lemma-characterize-contraction-to-stable}
对纤维上态射的刻画条件在基域扩张下保持。因此 (2) 成立。
\end{proof}

\section{稳定约化定理}
\label{moduli-curves-section-stable-reduction}

\noindent
在半稳定约化一章中，我们已经证明了著名的曲线半稳定约化定理。
设 $K$ 是离散赋值环 $R$ 的分式域。设 $C$ 是 $K$ 上的射影光滑曲线，
满足 $K = H^0(C, \mathcal{O}_C)$。按照《半稳定约化》中的定义
\ref{models-definition-semistable}，若 $R$ 上存在泛纤维为 $C$
的预稳定曲线族，或者 $C$ 在 $R$ 上的某个（等价地，任意）
极小正则模型预稳定，则称 $C$ 具有{\it 半稳定约化}。
本节将说明，对亏格为 $g \geq 2$ 的曲线，这也等价于稳定约化。

\begin{lemma}
\label{moduli-curves-lemma-stable-reduction}
设 $R$ 是分式域为 $K$ 的离散赋值环。设 $C$ 是 $K$ 上亏格为
$g \geq 2$ 的光滑射影曲线，满足 $K = H^0(C, \mathcal{O}_C)$。
下列条件等价：
\begin{enumerate}
\item $C$ 具有半稳定约化
（《半稳定约化》中的定义 \ref{models-definition-semistable}）；
\item $R$ 上存在泛纤维为 $C$ 的稳定曲线族。
\end{enumerate}
\end{lemma}

\begin{proof}
稳定曲线族也是预稳定曲线族，故 (2) 蕴含 (1) 是立即的。
反过来，给定 $R$ 上泛纤维为 $C$ 的预稳定曲线族，
由引理 \ref{moduli-curves-lemma-contract-prestable-to-stable}，可以把它收缩成稳定曲线族。
由于泛纤维已经稳定，该过程不改变泛纤维，证明完成。
\end{proof}

\noindent
下述引理说明，引理 \ref{moduli-curves-lemma-stable-reduction} 所保证的 $R$
上稳定曲线族在唯一同构意义下唯一。

\begin{lemma}
\label{moduli-curves-lemma-unique-stable-model}
设 $R$ 是分式域为 $K$ 的离散赋值环。设 $C$ 是 $K$ 上亏格为 $g$
的光滑真曲线，满足 $K = H^0(C, \mathcal{O}_C)$。
若 $X$ 和 $X'$ 是 $C$ 的模型
（《半稳定约化》中的第 \ref{models-section-models} 节），
而且 $X$ 和 $X'$ 是 $R$ 上亏格为 $g$ 的稳定曲线族，
则存在唯一的模型同构 $X \to X'$。
\end{lemma}

\begin{proof}
设 $Y$ 是 $C$ 的极小模型。回顾 $Y$ 存在且唯一，
并且在 $R$ 上是相对维数为 $1$、至多有结点奇性的态射；
见《半稳定约化》中的命题
\ref{models-proposition-exists-minimal-model} 以及引理
\ref{models-lemma-minimal-model-unique} 和
\ref{models-lemma-semistable}（后者适用是因为我们有 $X$）。
存在收缩态射
$$
Y \longrightarrow Z
$$
使得 $Z$ 是 $R$ 上亏格为 $g$ 的稳定曲线族
（引理 \ref{moduli-curves-lemma-contract-prestable-to-stable}）。我们断言，
存在唯一的模型同构 $X \to Z$。由对称性，对 $X'$ 也有相同结论，
从而完成证明。

\medskip\noindent
由《半稳定约化》中的引理
\ref{models-lemma-blowup-at-worst-nodal}，存在序列
$$
X_m \to \ldots \to X_1 \to X_0 = X
$$
使得 $X_{i + 1} \to X_i$ 是 $X_i$ 的奇异闭点 $x_i$ 的吹起，
$X_i \to \Spec(R)$ 是相对维数为 $1$、至多有结点奇性的态射，
而 $X_m$ 正则。由《半稳定约化》中的引理
\ref{models-lemma-pre-exists-minimal-model}，存在 $C$ 的真正则模型序列
$$
X_m = Y_n \to Y_{n - 1} \to \ldots \to Y_1 \to Y_0 = Y
$$
使得每个态射都是第一类例外曲线的收缩\footnote{事实上，
有 $X_m = Y$；即 $X_m$ 不含任何第一类例外曲线。
我们鼓励读者仔细思考这一点，因为它会稍微简化证明。}。
由《半稳定约化》中的引理
\ref{models-lemma-blowdown-at-worst-nodal}，每个 $Y_i$
在 $R$ 上都是相对维数为 $1$、至多有结点奇性的态射。
为证明该断言，只需证明存在同构 $X \to Z$，它与态射
$X_m \to X$ 和 $X_m = Y_n \to Y \to Z$ 相容。
设 $s \in \Spec(R)$ 是闭点。由引理
\ref{moduli-curves-lemma-contract-basic-uniqueness} 或引理
\ref{moduli-curves-lemma-contract-prestable-to-stable}，问题归约为证明态射
$X_{m, s} \to X_s$ 和 $X_{m, s} \to Z_s$
都等于《代数曲线》中的引理
\ref{curves-lemma-characterize-contraction-to-stable} 的典范态射。

\medskip\noindent
对 $\kappa(s)$ 上概形之间的态射 $c : U \to V$，若对
$v \in V$ 有 $\dim(U_v) \leq 1$，并且
$\mathcal{O}_V = c_*\mathcal{O}_U$、$R^1c_*\mathcal{O}_U = 0$，
则称 $c$ 具有性质 (*)。该性质在复合下稳定。
由于 $X_s$ 和 $Z_s$ 都是 $\kappa(s)$ 上亏格为 $g$ 的稳定曲线，
只需证明态射 $Y_s \to Z_s$、$X_{i + 1, s} \to X_{i, s}$ 和
% P12-ERR-0312: see ../control/ERRATA_CANDIDATES.jsonl
$Y_{i + 1, s} \to Y_{i, s}$ 中的每一个都满足性质 (*)；见
《代数曲线》中的引理
\ref{curves-lemma-characterize-contraction-to-stable}。

\medskip\noindent
性质 (*) 对 $Y_s \to Z_s$ 由构造成立。

\medskip\noindent
态射 $c : X_{i + 1, s} \to X_{i, s}$ 构造并研究于
《半稳定约化》中的引理
\ref{models-lemma-blowup-at-worst-nodal} 的证明中。
只需在 $X_{i, s}$ 上 \'etale 局部检验 (*)。
因而只需对《半稳定约化》中的例 \ref{models-example-blowup}
的态射“$X_1 \to X_0$”到 $R/\pi R$ 的基变换检验 (*)。
把具体计算留给读者。

\medskip\noindent
态射 $c : Y_{i + 1, s} \to Y_{i, s}$ 是第一类例外曲线
$E \subset Y_{i + 1}$ 的吹下的限制；也就是说，
$b : Y_{i + 1} \to Y_i$ 是 $E$ 的收缩，也就是说，$b$ 是曲面
$Y_i$ 上一个正则点的吹起
（《曲面的解消》中的第 \ref{resolve-section-minus-one} 节）。
于是 $\mathcal{O}_{Y_i} = b_*\mathcal{O}_{Y_{i + 1}}$ 且
$R^1b_*\mathcal{O}_{Y_{i + 1}} = 0$；例如见《曲面的解消》中的引理
\ref{resolve-lemma-cohomology-of-blowup}。
由《态射进阶》中的引理
\ref{more-morphisms-lemma-check-h1-fibre-zero}、
\ref{more-morphisms-lemma-h1-fibre-zero} 和
\ref{more-morphisms-lemma-h1-fibre-zero-check-h0-kappa}，可知
$\mathcal{O}_{Y_{i, s}} = c_*\mathcal{O}_{Y_{i + 1, s}}$ 且
$R^1c_*\mathcal{O}_{Y_{i + 1, s}} = 0$
（最后一个引理只给出
$\mathcal{O}_{Y_{i, s}} \to c_*\mathcal{O}_{Y_{i + 1, s}}$
的满射性，但单射性容易由如下事实得出：$Y_{i, s}$ 约化，
而 $c$ 只改变一个闭点上方的部分）。证明完成。
\end{proof}

\noindent
由引理 \ref{moduli-curves-lemma-stable-reduction} 和《半稳定约化》中的定理
\ref{models-theorem-semistable-reduction}，立即得到稳定约化定理。

\begin{theorem}
\label{moduli-curves-theorem-stable-reduction}
\begin{reference}
\cite[Corollary 2.7]{DM}
\end{reference}
设 $R$ 是分式域为 $K$ 的离散赋值环。设 $C$ 是 $K$ 上亏格为
$g \geq 2$ 的光滑射影曲线，满足 $H^0(C, \mathcal{O}_C) = K$。则：
\begin{enumerate}
\item 存在离散赋值环扩张 $R \subset R'$，它诱导分式域的有限可分扩张
$K'/K$；并且存在该环上亏格为 $g$ 的稳定曲线族
$Y \to \Spec(R')$，满足在 $K'$ 上
$Y_{K'} \cong C_{K'}$；
\item 存在有限可分扩张 $L/K$ 以及亏格为 $g$ 的稳定曲线族
$Y \to \Spec(A)$，其中 $A \subset L$ 是 $R$ 在 $L$ 中的整闭包，
并且在 $L$ 上有 $Y_L \cong C_L$。
\end{enumerate}
\end{theorem}

\begin{proof}
(1) 立即由引理 \ref{moduli-curves-lemma-stable-reduction} 和《半稳定约化》中的定理
\ref{models-theorem-semistable-reduction} 得出。

\medskip\noindent
证明 (2)。设 $L/K$ 是《半稳定约化》中的定理
\ref{models-theorem-semistable-reduction} 第 (3) 部分所得的有限可分扩张。
设 $A \subset L$ 是 $R$ 的整闭包。回顾 $A$ 是 $R$ 上有限的
Dedekind 整环，具有有限多个极大理想
$\mathfrak m_1, \ldots, \mathfrak m_n$；见《代数进阶》中的注
\ref{more-algebra-remark-finite-separable-extension}。
置 $S = \Spec(A)$、$S_i = \Spec(A_{\mathfrak m_i})$、
$U = \Spec(L)$ 和 $U_i = S_i \setminus \{\mathfrak m_i\}$。
注意，对 $i = 1, \ldots, n$，有 $U \cong U_i$。
置 $X = C_L$，把它视为 $S$ 的开子概形 $U$ 上的概形。
由 $L$、$A$ 的选取以及引理 \ref{moduli-curves-lemma-stable-reduction}，
存在稳定曲线族 $X_i \to S_i$ 以及同构
$X \times_U U_i \cong X_i \times_{S_i} U_i$。
由《空间的极限》中的引理
\ref{spaces-limits-lemma-glueing-near-multiple-closed-points}，
可以找到有限表示态射 $Y \to S$，它到每个 $S_i$ 的基变换
都同构于 $X_i$（$i = 1, \ldots, n$）。或者，可以使用如下事实：
$S = \bigcup_{i = 1, \ldots, n} S_i$ 是 $S$ 的开覆盖，且当
$i \not = j$ 时 $S_i \cap S_j = U$；再应用 $n - 1$ 次
《空间的极限》中的引理 \ref{spaces-limits-lemma-relative-glueing}，
得到 $Y \to S$，它到每个 $S_i$ 的基变换都同构于 $X_i$
（$i = 1, \ldots, n$）。显然，$Y \to S$ 正是所求的稳定曲线族。
\end{proof}

\section{稳定曲线叠的性质}
\label{moduli-curves-section-properties-stable}

\noindent
本节证明 $g \geq 2$ 时 $\overline{\mathcal{M}}_g$ 的基本结构结果。

\begin{lemma}
\label{moduli-curves-lemma-stable-separated}
设 $g \geq 2$。叠 $\overline{\mathcal{M}}_g$ 分离。
\end{lemma}

\begin{proof}
该陈述意味着态射
$\overline{\mathcal{M}}_g \to \Spec(\mathbf{Z})$ 分离。
我们将使用《叠的态射进阶》中的引理
\ref{stacks-more-morphisms-lemma-refined-valuative-criterion-separated}
所陈述的精细 Noether 赋值判据来证明这一点。

\medskip\noindent
由于 $\overline{\mathcal{M}}_g$ 是 $\Curvesstack$ 的开子叠，
由引理 \ref{moduli-curves-lemma-curves-qs-lfp} 可知，
$\overline{\mathcal{M}}_g \to \Spec(\mathbf{Z})$ 拟分离且局部有限表示。
特别地，叠 $\overline{\mathcal{M}}_g$ 局部 Noether
（《叠的态射》中的引理
\ref{stacks-morphisms-lemma-locally-finite-type-locally-noetherian}）。
由引理 \ref{moduli-curves-lemma-smooth-dense-in-stable}，开浸入
$\mathcal{M}_g \to \overline{\mathcal{M}}_g$ 的像稠密。
此外，$\mathcal{M}_g \to \overline{\mathcal{M}}_g$ 拟紧
（《叠的态射》中的引理
\ref{stacks-morphisms-lemma-locally-closed-in-noetherian}），
因而是有限型的。因此，《叠的态射进阶》中的引理
\ref{stacks-more-morphisms-lemma-refined-valuative-criterion-separated}
对态射
$$
\mathcal{M}_g \to \overline{\mathcal{M}}_g
\quad\text{且}\quad
\overline{\mathcal{M}}_g \to \Spec(\mathbf{Z})
$$
的所有预备假设均成立，只需检验以下条件：给定任意 $2$-交换图
$$
\xymatrix{
\Spec(K) \ar[r] \ar[d] &
\mathcal{M}_g \ar[r] &
\overline{\mathcal{M}}_g \ar[d] \\
\Spec(R) \ar[rr] \ar@{..>}[rru] & & \Spec(\mathbf{Z})
}
$$
其中 $R$ 是分式域为 $K$ 的离散赋值环，则虚线箭头范畴或者为空，
或者是恰有一个同构类的 setoid。（注意，无需过多担心 $2$-箭头；
见《叠的态射》中的引理
\ref{stacks-morphisms-lemma-cat-dotted-arrows-independent}。）
展开其含义，并使用如下事实：$\mathcal{M}_g$ 和
$\overline{\mathcal{M}}_g$ 分别是参数化亏格为 $g$ 的光滑曲线族和
稳定曲线族的代数叠，便发现所需证明的正是引理
\ref{moduli-curves-lemma-unique-stable-model} 中陈述并证明的唯一性结果。
\end{proof}

\begin{lemma}
\label{moduli-curves-lemma-stable-quasi-compact}
设 $g \geq 2$。叠 $\overline{\mathcal{M}}_g$ 拟紧。
\end{lemma}

\begin{proof}
我们使用第 \ref{moduli-curves-section-polarized-curves} 节的记号。考虑点集
$$
T \subset |\textit{PolarizedCurves}|
$$
其中的点 $\xi$ 满足：存在域 $k$ 以及 $k$ 上表示 $\xi$ 的偶
$(X, \mathcal{L})$，且具有下列两个性质：
\begin{enumerate}
\item $X$ 是亏格为 $g$ 的稳定曲线；
\item $\mathcal{L} = \omega_X^{\otimes 3}$。
\end{enumerate}
显然，在连续映射
$$
|\textit{PolarizedCurves}|
\longrightarrow
|\Curvesstack|
$$
下，集合 $T$ 的像恰为开子集
$$
|\overline{\mathcal{M}}_g| \subset |\Curvesstack|
$$
因此，只需证明 $T$ 拟紧。由引理
\ref{moduli-curves-lemma-polarized-curves-in-polarized}，可知
$$
|\textit{PolarizedCurves}| \subset |\Polarizedstack|
$$
是开闭浸入。因此，只需证明 $T$ 作为 $|\Polarizedstack|$
的子集拟紧。为此，使用《模叠》中的引理
\ref{moduli-lemma-bounded-polarized} 的判据。首先注意，对上述
$(X, \mathcal{L})$，由 Riemann--Roch，Hilbert 多项式 $P$ 是函数
$P(t) = (6g - 6)t + (1 - g)$；见《代数曲线》中的引理
\ref{curves-lemma-rr}。其次，由《代数曲线》中的引理
\ref{curves-lemma-tricanonical}，有 $H^1(X, \mathcal{L}) = 0$，
而且 $\mathcal{L}$ 极丰沛。这恰好意味着：取
% P12-ERR-0313: see ../control/ERRATA_CANDIDATES.jsonl
$n = P(3) - 1$，存在闭浸入
$$
i : X \longrightarrow \mathbf{P}^n_k
$$
使得
% P12-ERR-0314: see ../control/ERRATA_CANDIDATES.jsonl
$\mathcal{L} = i^*\mathcal{O}_{\mathbf{P}^1_k}(1)$，如所求。
\end{proof}

\noindent
下面是本节的主要定理。

\begin{theorem}
\label{moduli-curves-theorem-stable-smooth-proper}
设 $g \geq 2$。代数叠 $\overline{\mathcal{M}}_g$ 是
Deligne--Mumford 叠，在 $\Spec(\mathbf{Z})$ 上真且光滑。
此外，参数化光滑曲线的轨迹 $\mathcal{M}_g$ 是稠密开子叠。
\end{theorem}

\begin{proof}
陈述中提到的大多数性质已经证明。光滑性是引理
\ref{moduli-curves-lemma-stable-curves-smooth}；Deligne--Mumford 性是引理
\ref{moduli-curves-lemma-stable-curves-deligne-mumford}；$\mathcal{M}_g$
的开性是引理 \ref{moduli-curves-lemma-smooth-dense-in-stable}。
由引理 \ref{moduli-curves-lemma-stable-separated}，已知
$\overline{\mathcal{M}}_g \to \Spec(\mathbf{Z})$ 分离；
由引理 \ref{moduli-curves-lemma-stable-quasi-compact}，已知
$\overline{\mathcal{M}}_g$ 拟紧。因此，为证明
$\overline{\mathcal{M}}_g \to \Spec(\mathbf{Z})$ 为真并完成证明，
可以对态射 $\mathcal{M}_g \to \overline{\mathcal{M}}_g$ 和
$\overline{\mathcal{M}}_g \to \Spec(\mathbf{Z})$ 应用
《叠的态射进阶》中的引理
\ref{stacks-more-morphisms-lemma-refined-valuative-criterion-proper}。
因而只需检验以下条件：给定任意 $2$-交换图
$$
\xymatrix{
\Spec(K) \ar[r] \ar[d]_j &
\mathcal{M}_g \ar[r] &
\overline{\mathcal{M}}_g \ar[d] \\
\Spec(A) \ar[rr] & & \Spec(\mathbf{Z})
}
$$
其中 $A$ 是分式域为 $K$ 的离散赋值环，则存在域扩张 $K'/K$
以及支配 $A$ 的赋值环 $A' \subset K'$，使得所诱导的图
$$
\xymatrix{
\Spec(K') \ar[r] \ar[d]_{j'} & \overline{\mathcal{M}}_g \ar[d] \\
\Spec(A') \ar[r] \ar@{..>}[ru] & \Spec(\mathbf{Z})
}
$$
的虚线箭头范畴非空
（《叠的态射》中的定义
\ref{stacks-morphisms-definition-fill-in-diagram}）。
（注意，无需过多担心 $2$-箭头；见《叠的态射》中的引理
\ref{stacks-morphisms-lemma-cat-dotted-arrows-independent}。）
展开其含义，并使用如下事实：$\mathcal{M}_g$ 和
$\overline{\mathcal{M}}_g$ 分别是参数化亏格为 $g$ 的光滑曲线族和
稳定曲线族的代数叠，便发现所需证明的正是稳定约化定理，
即定理 \ref{moduli-curves-theorem-stable-reduction} 所含的结果。
\end{proof}
% Generated mechanically from the frozen examples.tex target.
% Do not edit this generated file; edit the bound locale source instead.

% OK, start here.
%

\setcounter{chapter}{109}
\chapter{例}



\phantomsection
\label{examples-section-phantom}


\section{引言}
\label{examples-section-introduction}

\noindent
本章将收录若干阐明相关理论的例子。





\section{空的极限}
\label{examples-section-empty-limit}

\noindent
本例由 Waterhouse 给出，见 \cite{Waterhouse}。
设 $S$ 是不可数集。对每个有限子集
$T \subset S$，考虑单射 $T \to \mathbf{N}$ 所成的集合 $M_T$。
若 $T \subset T' \subset S$ 均有限，则限制映射 $M_{T'} \to M_T$
是满射。于是，以 $S$ 的有限子集所成的有向偏序集为指标，
我们得到一个转移映射均为满射的逆系统。
但 $\lim M_T = \emptyset$，因为极限中的一个元素将给出单射
$S \to \mathbf{N}$。



\section{零极限}
\label{examples-section-zero-limit}

\noindent
设 $(S_i)_{i \in I}$ 是由非空集合构成的有向逆系统，
其转移映射均为满射，并且
$\lim S_i = \emptyset$；见第 \ref{examples-section-empty-limit} 节。
设 $K$ 是域，并令
$$
V_i = \bigoplus\nolimits_{s \in S_i} K
$$
则转移映射 $V_i \to V_j$ 在 $i \geq j$ 时是满射。
然而，$\lim V_i = 0$。事实上，若 $v = (v_i)$ 是极限中的元素，
则 $v_i$ 的支撑将是有限子集 $T_i \subset S_i$，并且
$\lim T_i \not = \emptyset$；见
《范畴》中的引理 \ref{categories-lemma-nonempty-limit}。

\medskip\noindent
对每个 $i$，考虑唯一的 $K$-线性映射 $V_i \to K$，
它把每个基向量 $s \in S_i$ 送到 $1$。令 $W_i \subset V_i$
为其核。则
$$
0 \to (W_i) \to (V_i) \to (K) \to 0
$$
是以有向集 $I$ 为指标的向量空间逆系统的不分裂短正合列。因此
% P12-ERR-0315
$W_i$ 是转移映射均为满射、极限为零而
$R^1\lim$ 非零的 $K$-向量空间有向系统。



\section{拟紧空间的非拟紧逆极限}
\label{examples-section-lim-not-quasi-compact}

\noindent
以 $\mathbf{N}$ 表示自然数集。
对每个整数 $n$，令 $I_n$ 表示所有自然数 $> n$ 所成的集合。
把 $T_n$ 定义为 $\mathbf{N}$ 上以
$\{1\}, \ldots , \{n\}, I_n$ 为基的唯一拓扑。以 $X_n$ 表示拓扑空间
$(\mathbf{N}, T_n)$。对每个 $m < n$，恒等映射
$$
f_{n, m} : X_n \longrightarrow X_m
$$
连续。显然，当 $m < n < p$ 时，复合
$f_{p, n} \circ f_{n, m}$ 等于 $f_{p, m}$。因此 $((X_n), (f_{n,m}))$
是由拟紧拓扑空间构成的有向逆系统。

\medskip\noindent
设 $T$ 是 $\mathbf{N}$ 上的离散拓扑，并令 $X$ 为
$(\mathbf{N}, T)$。则对每个整数 $n$，恒等映射
$$
f_n : X \longrightarrow X_n
$$
连续。我们断言，这是上述有向系统的逆极限。设 $(Y, S)$ 是任意拓扑空间。
对每个整数 $n$，设
$$
g_n : (Y, S) \longrightarrow (\mathbf{N}, T_n)
$$
为连续映射。假设对每个 $m < n$ 都有
$f_{n,m} \circ g_n = g_m$，即系统 $(g_n)$ 与上述有向系统相容。
特别地，所有集合映射 $g_n$ 都等于同一个集合映射
$$
g : Y \longrightarrow \mathbf{N}.
$$
此外，对每个整数 $n$，由于 $\{n\}$ 在 $X_n$ 中开，
$g^{-1}(\{n\}) = g_n^{-1}(\{n\})$ 在 $Y$ 中也开。
因此，集合映射 $g$ 对拓扑 $S$（定义在 $Y$ 上）与拓扑 $T$
（定义在 $\mathbf{N}$ 上）连续。故 $(X, (f_n))$ 是上述有向系统的逆极限。

\medskip\noindent
然而，$X$ 显然不是拟紧的，因为由单点集组成的无限开覆盖没有逆极限。
% P12-ERR-0316

\begin{lemma}
\label{examples-lemma-lim-not-quasi-compact}
存在一个以 $\mathbf{N}$ 为指标的拟紧拓扑空间逆系统，其极限并不拟紧。
\end{lemma}

\begin{proof}
见上面的讨论。
\end{proof}





\section{纤维积上的结构层}
\label{examples-section-silly}

\noindent
设 $X, Y, S, a, b, p, q, f$ 如《概形的导出范畴》引言
第 \ref{perfect-section-kunneth} 节中所述。图示为
$$
\xymatrix{
& X \times_S Y \ar[ld]^p \ar[rd]_q \ar[dd]^f \\
X \ar[rd]_a & & Y \ar[ld]^b \\
& S
}
$$
于是有典范映射
$$
can :
p^{-1}\mathcal{O}_X \otimes_{f^{-1}\mathcal{O}_S} q^{-1}\mathcal{O}_Y
\longrightarrow
\mathcal{O}_{X \times_S Y}
$$
它一般不是同构。

\medskip\noindent
例如，令 $S = \Spec(\mathbf{R})$、$X = \Spec(\mathbf{C})$，并令
$Y = \Spec(\mathbf{C})$。则
$X \times_S Y = \Spec(\mathbf{C}) \amalg \Spec(\mathbf{C})$
是由两个点构成的离散空间，而层 $p^{-1}\mathcal{O}_X$、
$q^{-1}\mathcal{O}_Y$ 与 $f^{-1}\mathcal{O}_S$ 分别是取值于
$\mathbf{C}$、$\mathbf{C}$ 与 $\mathbf{R}$ 的常值层。
因此，$can$ 的源是这个两点离散空间上取值于
$\mathbf{C} \otimes_\mathbf{R} \mathbf{C}$ 的常值层。
故其整体截面作为 $\mathbf{R}$-向量空间的维数为 $8$；
另一方面，对 $can$ 的目标取整体截面得到
$\mathbf{C} \times \mathbf{C}$，
它作为 $\mathbf{R}$-向量空间的维数为 $4$。

\medskip\noindent
另一个例子如下。设 $k$ 是代数闭域。考虑
$S = \Spec(k)$、$X = \mathbf{A}^1_k$ 与 $Y = \mathbf{A}^1_k$。
若 $U \subset X \times_S Y = \mathbf{A}^2_k$ 是非空开集，
则像 $p(U) \subset X = \mathbf{A}^1_k$ 与 $q(U) \subset \mathbf{A}^1_k$
均为开集；读者可以证明
$$
\left(
p^{-1}\mathcal{O}_X \otimes_{f^{-1}\mathcal{O}_S} q^{-1}\mathcal{O}_Y
\right)(U) = \mathcal{O}_X(p(U)) \otimes_k \mathcal{O}_Y(q(U))
$$
若 $U$ 是 $X \times_S Y = \mathbf{A}^2_k$ 中一条不可约曲线 $C$
的补集，并且 $p|_C$ 与 $q|_C$ 都非常值，
则这不等于 $\mathcal{O}_{X \times_S Y}(U)$。

\medskip\noindent
回到一般情形，设 $z = (x, y, s, \mathfrak p)$ 是
$X \times_S Y$ 的一点，如《概形》中的引理
\ref{schemes-lemma-points-fibre-product} 所述。
则映射 $can$ 在 $z$ 处的茎上给出映射
$$
can_z :
\mathcal{O}_{X, x} \otimes_{\mathcal{O}_{S, s}} \mathcal{O}_{Y, y}
\longrightarrow
\mathcal{O}_{X \times_S Y, z}
$$
这是平坦环同态，因为目标是源的一个局部化
（细节从略；提示：约化到 $X$、$Y$ 与 $S$ 均为仿射的情形）。
注意，源一般不是局部环；这也从另一个角度说明 $can$ 一般不是同构。

\medskip\noindent
更一般地，设给定一个 $\mathcal{O}_X$-模 $\mathcal{F}$ 与一个
$\mathcal{O}_Y$-模 $\mathcal{G}$。则有典范映射
% P12-ERR-0317
\begin{align*}
& p^{-1}\mathcal{F} \otimes_{f^{-1}\mathcal{O}_S} q^{-1}\mathcal{G} \\
& =
p^{-1}(\mathcal{F} \otimes_{\mathcal{O}_X} \mathcal{O}_X)
\otimes_{f^{-1}\mathcal{O}_S}
q^{-1}(\mathcal{O}_Y \otimes_{\mathcal{O}_Y} \mathcal{G}) \\
& =
p^{-1}\mathcal{F} \otimes_{p^{-1}\mathcal{O}_X} p^{-1}\mathcal{O}_X
\otimes_{f^{-1}\mathcal{O}_S}
q^{-1}\mathcal{O}_Y \otimes_{q^{-1}\mathcal{O}_Y} q^{-1}\mathcal{G} \\
& \xrightarrow{can}
p^{-1}\mathcal{F} \otimes_{q^{-1}\mathcal{O}_X}
\mathcal{O}_{X \times_S Y}
\otimes_{q^{-1}\mathcal{O}_Y} q^{-1}\mathcal{G}  \\
& =
p^{-1}\mathcal{F} \otimes_{q^{-1}\mathcal{O}_X}
\mathcal{O}_{X \times_S Y}
\otimes_{\mathcal{O}_{X \times_S Y}}
\mathcal{O}_{X \times_S Y}
\otimes_{q^{-1}\mathcal{O}_Y} q^{-1}\mathcal{G}  \\
& =
p^*\mathcal{F} \otimes_{\mathcal{O}_{X \times_S Y}} q^*\mathcal{G}
\end{align*}
它很少是同构。





\section{一个非整的连通概形，其局部环均为整环}
\label{examples-section-connected-locally-integral-not-integral}

\noindent
我们给出仿射概形 $X = \Spec(A)$ 的一个例子：它是连通的，
其所有局部环都是整环，但它不是整概形。
$X$ 的连通性意味着 $A$ 没有非平凡幂等元；见
《代数》中的引理 \ref{algebra-lemma-disjoint-decomposition}。
若每当 $A$ 中 $fg = 0$ 时，$X$ 的每一点都有一个邻域，
使得 $f$ 或 $g$ 之一在其上为零，则 $X$ 的局部环均为整环。
只要 $A$ 不是整环，$X$ 就不是整概形
（《性质》中的定义 \ref{properties-definition-integral}）。

\medskip\noindent
粗略地说，构造如下：令 $X_0$ 为仿射平面中的十字形
（即两条坐标轴的并）。然后，在平面的吹起上取 $X_0$ 的
赋予约化结构的全逆像 $X_1$（$X_1$ 由三条排成链的有理分支组成）。
再在所得曲面上沿 $X_1$ 的两个奇点吹起，并令 $X_2$ 为
$X_1$ 的约化逆像（它有五条有理分支），如此继续。
取 $X$ 为逆极限。这个构造唯一的问题是，吹起会粘入一条射影直线，
所以 $X_1$ 不是仿射的。我们改为粘入一条仿射直线以修正这一点
（因而我们的概形将是上述概形的一个开子集）。

\medskip\noindent
下面给出完全代数的构造：对每个 $k \ge 0$，令 $A_k$
为如下环：其元素是多项式族
$p_i \in \mathbf{C}[x]$，其中 $i = 0, \ldots, 2^k$，并且满足
$p_i(1) = p_{i + 1}(0)$。令 $X_k = \Spec(A_k)$。注意，$X_k$
是 $2^k + 1$ 条仿射直线的并，它们横截相交并排成一条链。
如下定义环同态 $A_k \to A_{k + 1}$：
$$
(p_0, \ldots, p_{2^k})
\longmapsto
(p_0, p_0(1), p_1, p_1(1), \ldots, p_{2^k}),
$$
换言之，每隔一个多项式便是常值的。这个映射把
$A_k$ 认作 $A_{k + 1}$ 的子环。令 $A$ 为 $A_k$ 的直极限
（本质上即它们的并）。令 $X = \Spec(A)$。对每个 $k$，都有
自然嵌入 $A_k \to A$，也就是映射 $X\to X_k$。
每个 $A_k$ 都连通但不整；这蕴含 $A$ 连通但不整。
尚需证明 $A$ 的局部环均为整环。

\medskip\noindent
取 $f, g \in A$，满足 $fg = 0$，并取 $x \in X$。我们来构造
$x$ 的一个邻域，使 $f$ 与 $g$ 之一在其上为零。选取 $k$，
使得 $f, g \in A_{k - 1}$（注意指标为 $k - 1$）。
令 $y$ 为 $x$ 在 $X_k$ 中的像。只需证明 $y$ 有一个邻域，
使得视为 $\mathcal{O}_{X_k}$ 截面的 $f$ 或 $g$ 之一在其上为零。
若 $y$ 是 $X_k$ 的光滑点，即它只位于 $2^k + 1$ 条直线中的一条上，
则这很显然。因此，可以假设 $y$ 是 $2^k$ 个奇点之一，
从而 $X_k$ 的两条分支经过 $y$。然而，在这两条分支中的一条
（指标为奇数的那条）上，$f$ 与 $g$ 都是常值的，因为它们是
$X_{k - 1}$ 上函数的拉回。由于处处都有 $fg = 0$，
$f$ 或 $g$ 之一（不妨设为 $f$）在另一条分支上为零。
这蕴含 $f$ 在两条分支上都为零，正合所需。

\section{非完备的完备化}
\label{examples-section-noncomplete-completion}

\noindent
设 $R$ 是环，$\mathfrak m$ 是极大理想。考虑完备化
$$
R^\wedge = \lim R/\mathfrak m^n.
$$
注意，$R^\wedge$ 是局部环，其极大理想为
$\mathfrak m' = \Ker(R^\wedge \to R/\mathfrak m)$。
事实上，若 $x = (x_n) \in R^\wedge$ 不属于 $\mathfrak m'$，则
$y = (x_n^{-1}) \in R^\wedge$ 满足 $xy = 1$，于是由
《代数》中的引理 \ref{algebra-lemma-characterize-local-ring}，
$R^\wedge$ 是局部环。现在，$R^\wedge$ 关于其极限拓扑总是完备的
% P12-ERR-0318
（见《代数进阶》第 \ref{more-algebra-section-topological-ring} 节中的讨论）。
但除此之外，还有如下问题：
\begin{enumerate}
\item 是否有 $\mathfrak m R^\wedge = \mathfrak m'$？
\item 把 $R^\wedge$ 视为 $R^\wedge$-模时，它是否 $\mathfrak m'$-进完备？
\item 把 $R^\wedge$ 视为 $R$-模时，它是否 $\mathfrak m$-进完备？
\end{enumerate}
事实证明，这些问题的答案都是否定的。
下面的例子取自 Bart de Smit 与 Hendrik Lenstra 的一份未发表笔记。
另见 \cite[Exercise III.2.12]{Bourbaki-CA} 与
\cite[Example 1.8]{Yekutieli}

\medskip\noindent
设 $k$ 是域，$R = k[x_1, x_2, x_3, \ldots]$，且
$\mathfrak m = (x_1, x_2, x_3, \ldots)$。
我们把 $R^\wedge$ 的元素 $f$ 看作一个（可能）无限和
$$
f = \sum a_I x^I
$$
（采用多重指标记号），并要求对每个 $d \geq 0$，
满足 $|I| = d$ 的非零 $a_I$ 只有有限多个。极大理想
$\mathfrak m' \subset R^\wedge$ 由常数项为零的 $f$ 组成。
特别地，元素
$$
f = x_1 + x_2^2 + x_3^3 + \ldots
$$
属于 $\mathfrak m'$，但不属于 $\mathfrak m R^\wedge$，
这说明本例中的 (1) 不成立。然而，若 (1) 不成立，
则 (3) 必然不成立，因为
$\mathfrak m' = \Ker(R^\wedge \to R/\mathfrak m)$，
并且可以对 $n = 1$ 应用
《代数》中的引理 \ref{algebra-lemma-hathat}。

\medskip\noindent
最后，我们证明 $R^\wedge$ 不是 $\mathfrak m'$-进完备的。
对 $n \geq 1$，令 $K_n = \Ker(R^\wedge \to R/\mathfrak m^n)$。于是有短正合列
$$
0 \to K_n/(\mathfrak m')^n \to R^\wedge/(\mathfrak m')^n \to
R/\mathfrak m^n \to 0
$$
投影映射 $R^\wedge \to R/\mathfrak m^{n + 1}$ 把
$(\mathfrak m')^n$ 满射到 $\mathfrak m^n/\mathfrak m^{n + 1}$。
由此，$K_{n + 1} \to K_n/(\mathfrak m')^n$ 是满射。
因此，逆系统 $\left(K_n/(\mathfrak m')^n\right)$
的转移映射均为满射；取逆极限，并应用
《代数》中的引理 \ref{algebra-lemma-Mittag-Leffler}，得到正合列
$$
0 \to \lim K_n/(\mathfrak m')^n \to
\lim R^\wedge/(\mathfrak m')^n \to
\lim R/\mathfrak m^n \to 0
$$
于是可见，$R^\wedge$ 关于 $\mathfrak m'$ 完备，当且仅当
对所有 $n \geq 1$ 都有 $K_n = (\mathfrak m')^n$。

\medskip\noindent
为说明本例中的 $R^\wedge$ 不是 $\mathfrak m'$-进完备的，
我们证明 $K_2 = \Ker(R^\wedge \to R/\mathfrak m^2)$
不等于 $(\mathfrak m')^2$。
注意，$(\mathfrak m')^2$ 的元素可以写成有限和
\begin{equation}
\label{examples-equation-sum}
\sum\nolimits_{i = 1, \ldots, t} f_i g_i
\end{equation}
其中 $f_i, g_i \in R^\wedge$ 的常数项为零。
为构造一个例子，我们将选取如下形式的 $z \in K_2$
% P12-ERR-0319
$$
z = z_1 + z_2 + z_3 + \ldots
$$
并使它满足下列性质：
\begin{enumerate}
\item 存在序列 $1 < d_1 < d_2 < d_3 < \ldots $ 与
$0 < n_1 < n_2 < n_3 < \ldots$。

其中
$z_i \in k[x_{n_i}, x_{n_i + 1}, \ldots, x_{n_{i + 1} - 1}]$
是次数为 $d_i$ 的齐次元；并且
\item 在环 $k[[x_{n_i}, x_{n_i + 1}, \ldots, x_{n_{i + 1} - 1}]]$
中，元素 $z_i$ 不能写成 $t \leq i$ 的和式 (\ref{examples-equation-sum})。
\end{enumerate}
这显然蕴含 $z$ 不属于 $(\mathfrak m')^2$，因为当 $i$ 足够大时，
关系式 (\ref{examples-equation-sum}) 在环
$k[[x_{n_i}, x_{n_i + 1}, \ldots, x_{n_{i + 1} - 1}]]$
中的像会导出矛盾。因此，只需证明：对每个 $t > 0$，
都存在 $d \gg 0$ 与整数 $n$，使我们能够找到
次数为 $d$ 的齐次元 $z \in k[x_1, \ldots, x_n]$，
它在 $k[[x_1, \ldots, x_n]]$ 中不能针对给定的 $t$
写成和式 (\ref{examples-equation-sum})。取 $n > 2t$，再取任意
与 $k$ 的特征互素的 $d > 1$，并令
$z = \sum_{i = 1, \ldots, n} x_i^d$。则理想
$$
(\frac{\partial z}{\partial x_1}, \ldots, \frac{\partial z}{\partial x_n})
=
(dx_1^{d - 1}, \ldots, dx_n^{d - 1})
$$
的零点集只含一个点。另一方面，由 Leibniz 法则，
$$
\frac{\partial ( \sum\nolimits_{i = 1, \ldots, t} f_i g_i ) }{\partial x_j}
\in (f_1, \ldots, f_t, g_1, \ldots, g_t)
$$
故这些偏导数的零点集至少包含
$$
V(f_1, \ldots, f_t, g_1, \ldots, g_t) \subset
\Spec(k[[x_1, \ldots, x_n]]).
$$
这便产生矛盾，因为
$V(f_1, \ldots, f_t, g_1, \ldots, g_t)$ 的维数至少为
$n - 2t \geq 1$。

\begin{lemma}
\label{examples-lemma-noncomplete-completion}
存在局部环 $R$ 与极大理想 $\mathfrak m$，使得 $R$ 关于
$\mathfrak m$ 的完备化 $R^\wedge$ 具有下列性质：
\begin{enumerate}
\item $R^\wedge$ 是局部环，但其极大理想不等于
$\mathfrak m R^\wedge$；
\item $R^\wedge$ 不是完备局部环；并且
\item 作为 $R$-模，$R^\wedge$ 不是 $\mathfrak m$-进完备的。
\end{enumerate}
\end{lemma}

\begin{proof}
由上面的讨论即得，因为（取 $R = k[x_1, x_2, x_3, \ldots]$）
局部化 $R_{\mathfrak m}$ 的完备化等于 $R$ 的完备化。
\end{proof}


\section{非完备商}
\label{examples-section-noncomplete-quotient}

\noindent
设 $k$ 是域。令
$$
R = k[t, z_1, z_2, z_3, \ldots, w_1, w_2, w_3, \ldots, x]/
(z_it - x^iw_i, z_i w_j)
$$
特别注意，在此环中 $z_iz_jt = 0$。$R$ 的任意元素 $f$
都能唯一地写成有限和
$$
f = \sum\nolimits_{i = 0, \ldots, d} f_i x^i
$$
其中每个 $f_i \in k[t, z_i, w_j]$ 都不含涉及乘积
$z_it$ 或 $z_iw_j$ 的项。此外，若 $f$ 以此方式写出，
则 $f \in (x^n)$ 当且仅当对 $i < n$ 有 $f_i = 0$。
所以 $x$ 是非零因子，且 $\bigcap (x^n) = 0$。
令 $R^\wedge$ 为 $R$ 关于理想 $(x)$ 的完备化。
注意，$R^\wedge$ 是 $(x)$-进完备的；见
《代数》中的引理 \ref{algebra-lemma-hathat-finitely-generated}。
由上可见，$R^\wedge$ 的元素能唯一地写成无限和
$$
f = \sum\nolimits_{i = 0}^\infty f_i x^i
$$
其中每个 $f_i \in k[t, z_i, w_j]$ 都不含涉及乘积
$z_it$ 或 $z_iw_j$ 的项。考虑元素
$$
f = \sum\nolimits_{i = 1}^\infty x^i w_i =
xw_1 + x^2w_2 + x^3w_3 + \ldots
$$
也就是说，$f_n = w_n$。注意，对每个 $n$ 都有
$f \in (t , x^n)$，因为对所有 $m$ 都有 $x^mw_m \in (t)$。
我们断言 $f \not \in (t)$。为证明这一点，假设
$tg = f$，其中 $g = \sum g_lx^l$ 采用上述典范形式。
由于 $tz_iz_j = 0$，不妨假设没有任何 $g_l$ 含涉及乘积
$z_iz_j$ 的项。考察把 $tg$ 化为典范形式的过程，可得如下规则：
给定 $g_l$ 的任意一项 $c m$，其中 $c \in k$，而 $m$ 是
$t, z_i, w_j$ 中的单项式，作如下替换：
\begin{enumerate}
\item 若单项式 $m$ 不含任何 $z_i$，则 $ctm$ 是 $f_l$ 的一项；
\item 若单项式 $m$ 含有某个 $z_i$，则它等于
$m = z_i$，并且 $cw_i$ 是 $f_{l + i}$ 中的一项。
% P12-ERR-0320
\end{enumerate}
由于 $g_0$ 是多项式，其中只出现有限多个变量 $z_i$。
选取 $n$，使 $z_n$ 不出现在 $g_0$ 中。
于是上述规则表明 $w_n$ 不出现在 $f_n$ 中，矛盾。
因此，$R^\wedge/(t)$ 并不完备；见
《代数》中的引理 \ref{algebra-lemma-quotient-complete}。

\begin{lemma}
\label{examples-lemma-noncomplete-quotient}
存在一个关于主理想 $I$ 完备的环 $R$ 以及一个主理想 $J$，
使得 $R/J$ 不是 $I$-进完备的。
\end{lemma}

\begin{proof}
见上面的讨论。
\end{proof}


\section{完备化不正合}
\label{examples-section-completion-not-exact}

\noindent
一个简明的例子如下。设 $R = k[t]$。以 $M^\wedge$
表示 $R$-模关于 $I = (t)$ 的完备化。令
$P = K = \bigoplus_{n \in \mathbf{N}} R$，并令
$M = \bigoplus_{n \in \mathbf{N}} R/(t^n)$。则有短正合列
$0 \to K \to P \to M \to 0$，其中第一个映射在第 $n$ 个直和项上
由乘以 $t^n$ 给出。我们断言
$0 \to K^\wedge \to P^\wedge \to M^\wedge \to 0$ 在中间不正合。
事实上，$\xi = (t^2, t^3, t^4, \ldots) \in P^\wedge$
在 $M^\wedge$ 中的像为零，但它不在 $K^\wedge \to P^\wedge$ 的像中，
因为若在像中，它将是 $(t, t, t, \ldots)$ 的像，而后者并非
$K^\wedge$ 的元素。

\medskip\noindent
另一个“更小”的例子如下。在引理
\ref{examples-lemma-noncomplete-quotient} 的情形下，短正合列
$0 \to J \to R \to R/J \to 0$ 在关于 $I$ 完备化后不再正合。
事实上，若 $f \in J$ 是生成元，则
$f : R \to J$ 是满射，故 $R \to J^\wedge$ 是满射，因而
$J^\wedge \to R$ 的像为 $(f) = J$；但 $R/J$ 不完备意味着满射
$R \to (R/J)^\wedge$ 的核严格大于 $J$，见
《代数》中的引理 \ref{algebra-lemma-completion-generalities} 与
\ref{algebra-lemma-quotient-complete}。
同理，序列
$R \to R \to R/(f) \to 0$ 在完备化后也不再正合。

\begin{lemma}
\label{examples-lemma-completion-not-exact}
\begin{slogan}
一般而言，完备化既非左正合，也非右正合。
\end{slogan}
一般而言，完备化不是正合函子；甚至一般也不是右正合函子。
即便 $I$ 有限生成，并把它限制在有限表示模的范畴上，这仍然成立。
\end{lemma}

\begin{proof}
见上面的讨论。
\end{proof}



\section{完备模范畴不是 Abel 范畴}
\label{examples-section-non-abelian}

\noindent
设 $R$ 是环，$I \subset R$ 是有限生成理想。考虑由 $I$-进完备
$R$-模组成的范畴 $\mathcal{A}$；见
《代数》中的定义 \ref{algebra-definition-complete}。
设 $\varphi : M \to N$ 是 $\mathcal{A}$ 中的态射。
$\varphi$ 在 $\mathcal{A}$ 中的余核是通常余核的 $I$-进完备化
$(\Coker(\varphi))^\wedge$
（由于 $I$ 有限生成，这个完备化是完备的；见
《代数》中的引理 \ref{algebra-lemma-hathat-finitely-generated}）。
令 $K = \Ker(\varphi)$；由于 $\varphi$ 连续，它是 $M$ 的闭子模。
由下面的引理 \ref{examples-lemma-closed-is-complete}，$K$ 是 $I$-进完备的，
因而就是 $\varphi$ 在 $\mathcal{A}$ 中的核。

\begin{lemma}
\label{examples-lemma-closed-is-complete}
设 $I \subset R$ 是环的理想，$M$ 是 $I$-进完备的 $R$-模，
$N \subset M$ 是子模。下列条件等价：
\begin{enumerate}
\item $N$ 在 $M$ 中闭；
\item $N = \bigcap_{n \geq 1} (N + I^n M)$；
\item $M/N$ 是 $I$-进完备的。
\end{enumerate}
若 $I$ 有限生成，则这些条件蕴含 $N$ 是 $I$-进完备的。
\end{lemma}

\begin{proof}
(1) 与 (2) 的等价性来自
《形式空间》中的引理 \ref{formal-spaces-lemma-closed}。
(2) 与 (3) 的等价性是
《代数》中的引理 \ref{algebra-lemma-quotient-complete}。
假设 $I$ 有限生成且 (1) 成立。令 $N^\wedge$ 为 $N$ 的 $I$-进完备化。
由于 $M$ 是 $I$-进完备的，包含映射具有分解
$$
N \to N^\wedge \to M
$$
由于 $N^\wedge \to M$ 连续，$N$ 在 $N^\wedge$ 中稠密，
而 $N$ 在 $M$ 中闭，可见 $N^\wedge \to M$ 经 $N$ 分解。
于是，对某个 $A$-模 $C$，有 $N^\wedge = N \oplus C$。
% P12-ERR-0321
因此，$N$ 是 $I$-进完备模 $N^\wedge$ 的直和项；见
《代数》中的引理 \ref{algebra-lemma-hathat-finitely-generated}
（此处用到了 $I$ 有限生成），故 $N$ 是 $I$-进完备的。
\end{proof}

\noindent
我们给出一个例子，说明一般在范畴 $\mathcal{A}$ 中
$\Im \not = \Coim$。取
$R = \mathbf{Z}_p = \lim_n \mathbf{Z}/p^n\mathbf{Z}$
为 $p$-进整数环，并取 $I = (p)$。考虑映射
$$
\text{diag}(1, p, p^2, \ldots) :
\left(\bigoplus\nolimits_{n \geq 1} \mathbf{Z}_p\right)^\wedge
\longrightarrow
\prod\nolimits_{n \geq 1} \mathbf{Z}_p
$$
其中左侧是直和的 $p$-进完备化。因此，左侧的元素是向量
$(x_1, x_2, x_3, \ldots)$，其中 $x_i \in \mathbf{Z}_p$，且当
$i \to \infty$ 时其 $p$-进赋值 $v_p(x_i) \to \infty$。
它被送到 $(x_1, px_2, p^2x_3, \ldots)$。由此可见，
$(1, p, p^2, \ldots)$ 属于像的闭包，但不属于像。
由上面对核与余核的描述，显然对该映射有
$\Im \not = \Coim$。

\begin{lemma}
\label{examples-lemma-complete-modules-not-abelian}
设 $R$ 是环，$I \subset R$ 是有限生成理想。
$I$-进完备 $R$-模的范畴具有核与余核，但一般不是 Abel 范畴，
即使 $R$ 是 Noether 环亦然。
\end{lemma}

\begin{proof}
见上文。
\end{proof}



\section{导出完备模范畴}
\label{examples-section-derived-complete-modules}

\noindent
阅读本节之前，请先阅读《代数进阶》第
\ref{more-algebra-section-derived-complete-modules} 节。

\medskip\noindent
设 $A$ 是环，$I$ 是 $A$ 的理想，并以 $\mathcal{C}$ 表示
《代数进阶》定义 \ref{more-algebra-definition-derived-complete}
所定义的导出完备模范畴。

\medskip\noindent
设 $T$ 是集合，$M_t$（$t \in T$）是一族导出完备模。
我们断言，一般而言，$\bigoplus M_t$ 不是导出完备模。
作为一个具体例子，令 $A = \mathbf{Z}_p$、$I = (p)$，并考虑
$\bigoplus_{n \in \mathbf{N}} \mathbf{Z}_p$。从
$\bigoplus_{n \in \mathbf{N}} \mathbf{Z}_p$ 到其 $p$-进完备化的映射
不是满射。这意味着 $\bigoplus_{n \in \mathbf{N}} \mathbf{Z}_p$
不可能导出完备，否则将蕴含该映射为满射；见
《代数进阶》中的引理 \ref{more-algebra-lemma-complete-derived-complete}。
因此，包含函子 $\mathcal{C} \to \text{Mod}_A$
既不与直和交换，也不与（滤过）余极限交换。

\medskip\noindent
假设 $I$ 有限生成。由《代数进阶》第
\ref{more-algebra-section-derived-complete-modules} 节中的讨论，
范畴 $\mathcal{C}$ 具有任意余极限。然而，我们断言，
滤过余极限在范畴 $\mathcal{C}$ 中并不正合。事实上，设
$A = \mathbf{Z}_p$ 且 $I = (p)$。存在 $p$-进完备 $A$-模的嵌入
$f_n : \mathbf{Z}_p/p\mathbf{Z}_p \to \mathbf{Z}_p/p^n\mathbf{Z}_p$，
由乘以 $p^{n - 1}$ 给出。存在交换图
$$
\xymatrix{
\mathbf{Z}_p/p\mathbf{Z}_p \ar[r]_{f_n} \ar[d]^1 &
\mathbf{Z}_p/p^n\mathbf{Z}_p \ar[d]_p \\
\mathbf{Z}_p/p\mathbf{Z}_p \ar[r]^{f_{n + 1}} &
\mathbf{Z}_p/p^{n + 1}\mathbf{Z}_p
}
$$
我们断言：这些嵌入在范畴 $\mathcal{C}$ 中的余极限给出映射
$\mathbf{Z}_p/p\mathbf{Z}_p \to 0$。事实上，右侧系统在
$\text{Mod}_A$ 中的余极限是 $\mathbf{Q}_p/\mathbf{Z}_p$。
因此，其在 $\mathcal{C}$ 中的余极限为
$$
H^0((\mathbf{Q}_p/\mathbf{Z}_p)^\wedge) =
H^0(\mathbf{Z}_p[1]) = 0
$$
这是《代数进阶》第
\ref{more-algebra-section-derived-complete-modules} 节的结论，
其中 ${}^\wedge$ 表示导出完备化。这证明了我们的断言。

\begin{lemma}
\label{examples-lemma-derived-complete-modules}
设 $A$ 是环，$I \subset A$ 是理想。
导出完备模的范畴 $\mathcal{C}$ 是 Abel 范畴，而包含函子
$F : \mathcal{C} \to \text{Mod}_A$ 正合且与任意极限交换。
若 $I$ 有限生成，则 $\mathcal{C}$ 具有任意直和与余极限，
但一般而言，$F$ 不与它们交换。最后，滤过余极限一般在
$\mathcal{C}$ 中并不正合，故 $\mathcal{C}$ 不是 Grothendieck Abel 范畴。
\end{lemma}

\begin{proof}
见《代数进阶》中的引理
\ref{more-algebra-lemma-derived-complete-modules} 以及上面的讨论。
\end{proof}



\section{非平坦完备化}
\label{examples-section-nonflat}

\noindent
与 Noether 情形不同，环关于理想的完备化并不总是平坦的。
我们已在《代数进阶》的例
\ref{more-algebra-example-not-glueing-pair} 中见过这种现象的两个例子。
本节再给出两个例子。

\begin{lemma}
\label{examples-lemma-countable-fg-tensor}
设 $R$ 是环，$M$ 是可数的 $R$-模。则 $M$ 是有限 $R$-模，
当且仅当 $M \otimes_R R^\mathbf{N} \to M^\mathbf{N}$ 是满射。
\end{lemma}

\begin{proof}
若 $M$ 是有限模，则由《代数》中的命题
\ref{algebra-proposition-fg-tensor}，该映射是满射。反之，假设该映射满射。
把 $M$ 的元素枚举为 $m_1, m_2, m_3, \ldots$。
设张量积中的元素 $\sum_{j = 1, \ldots, m} x_j \otimes a_j$
被送到 $(m_n) \in M^\mathbf{N}$。于是，正如《代数》命题
\ref{algebra-proposition-fg-tensor} 的证明中那样，可见
$x_1, \ldots, x_m$ 在 $R$ 上生成 $M$。
\end{proof}

\begin{lemma}
\label{examples-lemma-countable-fp-tensor}
设 $R$ 是可数环，$M$ 是可数的 $R$-模。则 $M$ 有限表示，
当且仅当典范映射
$M \otimes_R R^\mathbf{N} \to M^\mathbf{N}$ 是同构。
\end{lemma}

\begin{proof}
若 $M$ 是有限表示模，则由《代数》中的命题
\ref{algebra-proposition-fp-tensor}，该映射是同构。反之，假设该映射是同构。
由引理 \ref{examples-lemma-countable-fg-tensor}，模 $M$ 是有限的。
选取以 $K$ 为核的满射 $R^{\oplus m} \to M$。则 $K$ 作为
$R^{\oplus m}$ 的子模是可数的。仿照《代数》命题
\ref{algebra-proposition-fp-tensor} 的证明，可见
$K \otimes_R R^\mathbf{N} \to K^\mathbf{N}$ 是满射。
于是由引理 \ref{examples-lemma-countable-fg-tensor}，$K$ 是有限 $R$-模。
故 $M$ 有限表示。
\end{proof}

\begin{lemma}
\label{examples-lemma-countable-coherent}
设 $R$ 是可数环。则 $R$ 是凝聚环，当且仅当
$R^\mathbf{N}$ 是平坦 $R$-模。
\end{lemma}

\begin{proof}
若 $R$ 是凝聚环，则由《代数》中的命题
\ref{algebra-proposition-characterize-coherent}，$R^\mathbf{N}$ 是平坦模。
假设 $R^\mathbf{N}$ 平坦。设 $I \subset R$ 是有限生成理想。
为证明该引理，我们证明 $I$ 作为 $R$-模有限表示。事实上，
由于 $R^\mathbf{N}$ 平坦，映射
$I \otimes_R R^\mathbf{N} \to R^\mathbf{N}$ 是单射；
由引理 \ref{examples-lemma-countable-fg-tensor}，其像为 $I^\mathbf{N}$。
于是由引理 \ref{examples-lemma-countable-fp-tensor} 得出结论。
\end{proof}

\noindent
设 $R$ 是可数环。注意，作为 $R$-模，$R[[x]]$ 同构于
$R^\mathbf{N}$。由引理 \ref{examples-lemma-countable-coherent} 可见，
$R \to R[[x]]$ 平坦，当且仅当 $R$ 是凝聚环。
存在许多非凝聚的可数环，例如
$$
R = k[y, z, a_1, b_1, a_2, b_2, a_3, b_3, \ldots]/
(a_1 y + b_1 z, a_2 y + b_2 z, a_3 y + b_3 z, \ldots)
$$
其中 $k$ 是可数域。该环不是凝聚环，因为 $R$ 的理想
$(y, z)$ 不是有限表示 $R$-模。注意，$R[[x]]$ 是
$R[x]$ 关于主理想 $(x)$ 的完备化。

\begin{lemma}
\label{examples-lemma-completion-polynomial-ring-not-flat}
存在一个环，使得 $R[x]$ 在 $(x)$ 处的完备化 $R[[x]]$
在 $R$ 上不平坦，因而更不在 $R[x]$ 上平坦。
\end{lemma}

\begin{proof}
见上面的讨论。
\end{proof}

\noindent
事实证明，存在环 $R$，使得 $R[[x]]$ 在 $R$ 上平坦，
但 $R[[x]]$ 在 $R[x]$ 上不平坦。见 Badam Baplan 的
\href{https://math.stackexchange.com/users/164860/badam-baplan}{这篇帖子}。
具体而言，令 $R$ 为赋值环。则 $R$ 是凝聚环
（《代数》中的例 \ref{algebra-example-valuation-ring-coherent}），
因而由《代数》中的命题
\ref{algebra-proposition-characterize-coherent}，$R[[x]]$ 在 $R$ 上平坦。
另一方面，有如下引理。

\begin{lemma}
\label{examples-lemma-almost-integral-when-powerseries-flat}
设 $R$ 是分式域为 $K$ 的整环。若 $R[[x]]$ 在 $R[x]$ 上平坦，
则 $R$ 正规，当且仅当 $R$ 完全正规
（《代数》中的定义 \ref{algebra-definition-almost-integral}）。
\end{lemma}

\begin{proof}
假设存在 $\alpha \in K$ 与非零元 $r \in R$，使得对所有
$n \geq 1$ 都有 $r \alpha^n \in R$。在 $R[[x]]$ 中考虑
$f = \sum r \alpha^{n - 1} x^n$。写 $\alpha = a/b$，
其中 $a, b \in R$ 且 $b$ 非零。于是 $(a x - b)f = -rb$。
故 $rb$ 属于理想 $(ax - b)R[[x]]$。
令 $S = \{h \in R[x] : h(0) = 1\}$。这是乘法子集，
而 $R[x] \to R[[x]]$ 的平坦性蕴含
$S^{-1}R[x] \to R[[x]]$ 忠实平坦
（细节从略；提示：使用《代数》中的引理
\ref{algebra-lemma-ff-rings}）。因此
$$
S^{-1}R[x]/(ax - b)S^{-1}R[x] \to R[[x]]/(ax - b)R[[x]]
$$
是单射。由此存在 $h \in S$ 与 $g \in R[x]$，使得
$h rb = (ax - b) g$。写成
$h = 1 + h_1 x + \ldots + h_d x^d$，便得到
$$
1 + h_1 x + \ldots + h_d x^d = (1/r)(\alpha x - 1)g
$$
这个在 $K[x]$ 中的分解给出 $K[x^{-1}]$ 中相应的分解，
从而说明 $\alpha$ 是一个系数在 $R$ 中的首一多项式的根，正合所需。
\end{proof}

\begin{lemma}
\label{examples-lemma-completion-polynomial-ring-not-flat-bis}
若 $R$ 是维数 $> 1$ 的赋值环，则 $R[[x]]$ 在 $R$ 上平坦，
但在 $R[x]$ 上不平坦。
\end{lemma}

\begin{proof}
上面的论证表明，只要证明 $R$ 不是完全正规的，结论便成立
（赋值环是正规的；见《代数》中的引理
\ref{algebra-lemma-valuation-ring-normal}）。
取素理想链 $\mathfrak p \subset \mathfrak m \subset R$。
选取非零元 $x \in \mathfrak p$ 与
$y \in \mathfrak m \setminus \mathfrak p$。
则对所有 $n \geq 1$，都有 $x y^{-n} \in R$
（否则 $y^n/x \in R$，这与 $y \not \in \mathfrak p$ 矛盾）。
因此，$1/y$ 在 $R$ 上几乎整，但不属于 $R$。
\end{proof}

\noindent
下面构造一个局部化的完备化不平坦的例子。为此，考虑环
$$
R = k[y, z, a_1, a_2, a_3, \ldots]/(ya_i, a_i a_j)
$$
以 $f \in R$ 表示 $z$ 的剩余类。我们断言环映射
\begin{equation}
\label{examples-equation-nonflat}
R[[x]] \longrightarrow R_f[[x]]
\end{equation}
不平坦。令 $I$ 为 $y : R[[x]] \to R[[x]]$ 的核。
$I$ 的典型元素 $g$ 形如 $g = \sum g_{n, m} a_mx^n$，
其中 $g_{n, m} \in k[z]$，并且对给定的 $n$，
非零的 $g_{n, m}$ 只有有限多个。令 $J$ 为
$y : R_f[[x]] \to R_f[[x]]$ 的核。我们断言
$J \not = I R_f[[x]]$。事实上，若二者相等，则有
$$
\sum z^{-n} a_n x^n = \sum\nolimits_{i = 1, \ldots, m} h_i g_i
$$
其中某个 $m \geq 1$、$g_i \in I$，且 $h_i \in R_f[[x]]$。
设
$h_i = \bar h_i \bmod (y, a_1, a_2, a_3, \ldots)$，
其中 $\bar h_i \in k[z, 1/z][[x]]$。考察 $a_n$ 的系数，
并使用上面对元素 $g_i$ 的描述，将得到
$$
z^{-n} x^n = \sum \bar h_i \bar g_{i, n}
$$
其中某些 $\bar g_{i, n} \in k[z][[x]]$。这意味着所有
$z^{-n}x^n$ 都包含在由元素 $\bar h_i$ 生成的有限
$k[z][[x]]$-模中。由于 $k[z][[x]]$ 是 Noether 环，
这蕴含 $k[z, 1/z][[x]]$ 中由
$1, z^{-1}x, z^{-2}x^2, \ldots$ 生成的
$R[z][[x]]$-子模是有限的。
% P12-ERR-0322
由《代数》中的引理 \ref{algebra-lemma-characterize-integral-element}，
可得 $z^{-1}x$ 在 $k[z][[x]]$ 上整，矛盾。另一方面，
若 (\ref{examples-equation-nonflat}) 平坦，则把正合列
$0 \to I \to R[[x]] \xrightarrow{y} R[[x]]$ 与 $R_f[[x]]$ 张量，
将得到 $J = IR_f[[x]]$。

\begin{lemma}
\label{examples-lemma-nonflat-completion-localization}
存在关于主理想 $I$ 完备的环 $A$ 以及元素 $f \in A$，
使得 $A_f$ 的 $I$-进完备化 $A_f^\wedge$ 在 $A$ 上不平坦。
\end{lemma}

\begin{proof}
令 $A = R[[x]]$、$I = (x)$，并注意 $R_f[[x]]$
是 $R[[x]]_f$ 的完备化。
\end{proof}



\section{拟凝聚模的非 Abel 范畴}
\label{examples-section-nonabelian-QCoh}

\noindent
在《叠上的层》第 \ref{stacks-sheaves-section-quasi-coherent} 节中，
我们定义了 $\Sch$ 上群胚纤维化范畴上的拟凝聚模范畴。
虽然我们在《叠上的层》第
\ref{stacks-sheaves-section-quasi-coherent-algebraic-stacks} 节中证明，
对于代数叠，该范畴是 Abel 范畴；但本节将说明，
对于形式代数空间，情况并非如此。

\medskip\noindent
具体而言，把 $\mathbf{Z}_p$ 视为带有 $p$-进拓扑的拓扑环。
令 $X = \text{Spf}(\mathbf{Z}_p)$；见《形式空间》中的定义
\ref{formal-spaces-definition-affine-formal-spectrum}。
则 $X$ 是 $(\Sch/\mathbf{Z})_{fppf}$ 上的集合层，并给出一个
setoid 叠 $\mathcal{X}$；见《叠》中的引理
\ref{stacks-lemma-when-stack-in-sets}。因此，《叠上的层》第
\ref{stacks-sheaves-section-quasi-coherent-algebraic-stacks} 节的讨论适用。

\medskip\noindent
设 $\mathcal{F}$ 是 $\mathcal{X}$ 上的拟凝聚模。
由于 $X = \colim \Spec(\mathbf{Z}/p^n\mathbf{Z})$，由《叠上的层》中的引理
\ref{stacks-sheaves-lemma-quasi-coherent}，显然 $\mathcal{F}$
由序列 $(\mathcal{F}_n)$ 给出，其中
\begin{enumerate}
\item $\mathcal{F}_n$ 是 $\Spec(\mathbf{Z}/p^n\mathbf{Z})$ 上的拟凝聚模；并且
\item 转移映射给出同构
$\mathcal{F}_n = \mathcal{F}_{n + 1}/p^n\mathcal{F}_{n + 1}$。
\end{enumerate}
转化为模，可见 $\mathcal{F}$ 对应于系统 $(M_n)$，
其中每个 $M_n$ 都是被 $p^n$ 零化的 Abel 群，而转移映射诱导同构
$M_n = M_{n + 1}/p^n M_{n + 1}$。在此情形下，模
$M = \lim M_n$ 是 $p$-进完备模，且 $M_n = M/p^n M$；见
《代数》中的引理 \ref{algebra-lemma-limit-complete}。
由此可知，$X$ 上的拟凝聚模范畴等价于 $p$-进完备 Abel 群范畴。
该范畴不是 Abel 范畴；见第 \ref{examples-section-non-abelian} 节。

\begin{lemma}
\label{examples-lemma-quasi-coherent-not-abelian}
形式代数空间 $X$ 上的拟凝聚\footnote{此处采用上面所定义的拟凝聚模。
% P12-ERR-0323
鉴于 Stacks Project 中的设置方式，这确实是正确的定义；
如上所见，我们的定义与人们对 $\text{Spf}(A)$ 上拟凝聚模的朴素定义一致，
即完备 $A$-模。}模范畴一般不是 Abel 范畴，
即使 $X$ 是 Noether 仿射形式代数空间亦然。
\end{lemma}

\begin{proof}
见上面的讨论。
\end{proof}





\section{局部可分裂的不可分裂序列}
\label{examples-section-nonsplit-locally-split}

\noindent
考虑短正合序列
$$
0 \to M \to
\bigoplus\nolimits_{p\text{ 为素数}} \mathbf{Z}_{(p)} \to \mathbf{Q} \to 0
$$
其中 $M$ 是右侧映射的核，而该映射在每个直和项上均取嵌入
$\mathbf{Z}_{(p)} \to \mathbf{Q}$。这个 $\mathbf{Z}$-模序列不可分裂，
因为不存在非零同态 $\mathbf{Q} \to \mathbf{Z}_{(p)}$。另一方面，
若在任意素数 $p$ 处局部化，该序列便成为可分裂的。

\begin{lemma}
\label{examples-lemma-nonsplit-locally-split}
存在环 $R$ 以及一个不可分裂的模序列，使它在 Zariski 局部成为可分裂的。
\end{lemma}

\begin{proof}
见上述讨论。
\end{proof}












































\section{正则序列与基变换}
\label{examples-section-regular-base-change}

\noindent
我们将构造一个带有正则序列 $(x, y, z)$ 的环 $R$，使得存在非零元素
$\delta \in R/zR$ 满足 $x\delta = y\delta = 0$。

\medskip\noindent
为构造这个例，我们先在环 $k[x, y, z]$ 上构造一个特殊的模 $E$，
其中 $k$ 是任意域。具体地，$E$ 是下图中的推出：
$$
\xymatrix{
\frac{xk[x, y, z, y^{-1}]}{xyk[x, y, z]} \ar[r] \ar[d]^{z/x} &
\frac{k[x, y, z, x^{-1}, y^{-1}]}{yk[x, y, z, x^{-1}]} \ar[r] \ar[d] &
\frac{k[x, y, z, x^{-1}, y^{-1}]}{yk[x, y, z, x^{-1}] + xk[x, y, z, y^{-1}]}
\ar[d] \\
\frac{k[x, y, z, y^{-1}]}{yzk[x, y, z]} \ar[r] &
E \ar[r] &
\frac{k[x, y, z, x^{-1}, y^{-1}]}{yk[x, y, z, x^{-1}] + xk[x, y, z, y^{-1}]}
}
$$
其中各行都是短正合序列（由于排版上的困难，我们省略了两端的零）。
也可将 $E$ 描述为
$$
E = \{(f, g) \mid f \in k[x, y, z, x^{-1}, y^{-1}],
g \in k[x, y, z, y^{-1}] \}/\sim
$$
其中 $(f, g) \sim (f', g')$ 当且仅当存在
$h \in k[x, y, z, y^{-1}]$ 使得
$$
f = f' + xh \bmod yk[x, y, z, x^{-1}], \quad
g = g' - zh \bmod yzk[x, y, z]
$$
我们断言：(a) $x : E \to E$ 是单射；(b)
$y : E/xE \to E/xE$ 是单射；(c) $E/(x, y)E = 0$；(d) 存在
非零元素 $\delta \in E/zE$ 使得 $x\delta = y\delta = 0$。

\medskip\noindent
为证明 (a)，设 $(f, g)$ 是给出 $E$ 中一个元素的一对，并且
$(xf, xg) \sim 0$。则存在 $h \in k[x, y, z, y^{-1}]$，使得
$xf + xh \in yk[x, y, z, x^{-1}]$ 且
$xg - zh \in yzk[x, y, z]$。可以假设
$h = \sum a_{i, j, k}x^iy^jz^k$ 是仅出现 $j \leq 0$ 的单项式之和。
于是 $xg - zh \in yzk[x, y, z]$ 蕴含只出现 $i > 0$，即对某个
$h' \in k[x, y, z, y^{-1}]$ 有 $h = xh'$。
此时 $(f, g) \sim (f + xh', g - zh')$，故可设 $g = 0$ 且 $h = 0$。
在这种情形下，$xf \in yk[x, y, z, x^{-1}]$ 蕴含
$f \in yk[x, y, z, x^{-1}]$，从而 $(f, g) \sim 0$。
因此 $x : E \to E$ 是单射。

\medskip\noindent
由于乘以 $x$ 在
$\frac{k[x, y, z, x^{-1}, y^{-1}]}{yk[x, y, z, x^{-1}]}$ 上是同构，
可见 $E/xE$ 同构于
$$
\frac{k[x, y, z, y^{-1}]}{
yzk[x, y, z] + xk[x, y, z, y^{-1}] + zk[x, y, z, y^{-1}]}
=
\frac{k[x, y, z, y^{-1}]}{xk[x, y, z, y^{-1}] + zk[x, y, z, y^{-1}]}
$$
因而乘以 $y$ 在 $E/xE$ 上是同构。这显然蕴含 (b) 和 (c)。

\medskip\noindent
令 $e \in E$ 为 $(1, 0)$ 的等价类。假设 $e \in zE$。
则存在 $f \in k[x, y, z, x^{-1}, y^{-1}]$、
$g \in k[x, y, z, y^{-1}]$ 和 $h \in k[x, y, z, y^{-1}]$，使得
$$
1 + zf + xh \in yk[x, y, z, x^{-1}], \quad
0 + zg - zh \in yzk[x, y, z].
$$
这是不可能的：单项式 $1$ 既不可能出现在 $zf$ 中，也不可能出现在
$xh$ 中。另一方面，$ye = 0$，且
$xe = (x, 0) \sim (0, -z) = z(0, -1)$。因此，令 $\delta$ 等于
$e$ 在 $E/zE$ 中的同余类，即得到 (d)。

\begin{lemma}
\label{examples-lemma-strange-regular-sequence}
存在局部环 $R$ 和（位于极大理想中的）正则序列 $x, y, z$，
使得存在非零元素 $\delta \in R/zR$ 满足
$x\delta = y\delta = 0$。
\end{lemma}

\begin{proof}
令 $R = k[x, y, z] \oplus E$，其中把上述模 $E$ 视为平方零理想。
显然，$x, y, z$ 是 $R$ 中的正则序列，而元素
$\delta \in E/zE \subset R/zR$ 具有所需性质。
为得到局部的例，可将 $R$ 在极大理想 $\mathfrak m = (x, y, z, E)$ 处局部化。
序列 $x, y, z$ 仍为正则序列（因为局部化正合），而元素 $\delta$
仍非零，因为它的支撑位于 $\mathfrak m$。
\end{proof}

\begin{lemma}
\label{examples-lemma-base-change-regular-sequence}
存在局部环的局部同态 $A \to B$，以及 $B$ 的极大理想中的正则序列
$x, y$，使得 $B/(x, y)$ 在 $A$ 上平坦，但 $x, y$ 在
$B/\mathfrak m_AB$ 中的像 $\overline{x}, \overline{y}$ 不构成正则序列，
甚至不构成 Koszul 正则序列。
\end{lemma}

\begin{proof}
令 $A = k[z]_{(z)}$ 且
$B = (k[x, y, z] \oplus E)_{(x, y, z, E)}$。
由 $x, y, z$ 是 $B$ 中的正则序列（见引理
\ref{examples-lemma-strange-regular-sequence} 的证明），可知 $x, y$ 是 $B$ 中的
正则序列，并且 $B/(x, y)$ 是无挠 $A$-模，因而平坦。
另一方面，存在非零元素
$\delta \in B/\mathfrak m_AB = B/zB$ 被
$\overline{x}, \overline{y}$ 所零化。因此
$H_2(K_\bullet(B/\mathfrak m_AB, \overline{x}, \overline{y})) \not = 0$。
故 $\overline{x}, \overline{y}$ 不是 Koszul 正则的，特别地，它不是正则序列；
见 More on Algebra，引理
\ref{more-algebra-lemma-regular-koszul-regular}。
\end{proof}





\section{无限维 Noether 环}
\label{examples-section-Noetherian-infinite-dimension}

\noindent
正如我们在《代数》命题 \ref{algebra-proposition-dimension} 中所见，
Noether 局部环具有有限维数。不过，也存在无限维的 Noether 环。
见 \cite[附录，例 1]{Nagata}。

\medskip\noindent
具体地，设 $k$ 为域，并考虑环
$$
R = k[x_1, x_2, x_3, \ldots ].
$$
对 $i = 1, 2, \ldots$，令
$\mathfrak p_i = (x_{2^{i - 1}}, x_{2^{i - 1} + 1}, \ldots, x_{2^i - 1})$；
它们是 $R$ 的素理想。令 $S$ 为乘法子集
$$
S = \bigcap\nolimits_{i \geq 1} (R \setminus \mathfrak p_i).
$$
考虑环 $A = S^{-1}R$。我们断言：
\begin{enumerate}
\item 环 $A$ 的极大理想恰为理想 $\mathfrak m_i = \mathfrak p_iA$。
% P12-ERR-0324
\item 有 $A_{\mathfrak m_i} = R_{\mathfrak p_i}$，它是维数为 $2^i$ 的
Noether 局部环。
\item 环 $A$ 是 Noether 环。
\end{enumerate}
因此，这显然就是我们要找的例。细节从略。





\section{完备化非约化的局部环}
\label{examples-section-local-completion-nonreduced}

\noindent
在《代数》例 \ref{algebra-example-bad-dvr-char-p} 中，我们给出了一个
特征为 $p$、维数为 $1$ 的 Noether 局部整环 $R$，其完备化非约化。
本节给出 \cite[命题 3.1]{Ferrand-Raynaud} 中的例，它在特征零时给出类似的环。

\medskip\noindent
令 $\mathbf{C}\{x\}$ 为复数域 $\mathbf{C}$ 上的收敛幂级数环。
全体幂级数环 $\mathbf{C}[[x]]$ 是它的完备化。令
$K = \mathbf{C}\{x\}[1/x]$ 为收敛 Laurent 级数域。
$K$ 在 $\mathbf{C}$ 上的代数微分 $K$-模
$\Omega_{K/\mathbf{C}}$ 是无限维 $K$-向量空间（证明从略）。
我们可以选取 $f_n \in x\mathbf{C}\{x\}$（$n \geq 1$），使得
$
\text{d}x, \text{d}f_1, \text{d}f_2, \ldots
$
构成 $\Omega_{K/\mathbf{C}}$ 的某个基的一部分。因此，可以找到一个
$\mathbf{C}$-导子
$$
D : \mathbf{C}\{x\} \longrightarrow \mathbf{C}((x))
$$
% P12-ERR-0325
使得 $D(x) = 0$ 且 $D(f_i) = x^{-n}$。令
$$
A = \{f \in \mathbf{C}\{x\} \mid D(f) \in \mathbf{C}[[x]]\}
$$
我们断言：
\begin{enumerate}
\item $\mathbf{C}\{x\}$ 在 $A$ 上整；
\item $A$ 是局部整环；
\item $\dim(A) = 1$；
\item $A$ 的极大理想由 $x$ 和 $xf_1$ 生成；
\item $A$ 是 Noether 环；并且
\item $A$ 的完备化等于 $\mathbf{C}[[x]]$ 上的对偶数环。
\end{enumerate}
由于对偶数环非约化，环 $A$ 给出了所需的例。

\medskip\noindent
注意，若 $0 \not = f \in x\mathbf{C}\{x\}$，则对某个 $n \geq 0$
和 $h \in \mathbf{C}[[x]]$，可写成 $D(f) = h/f^n$。
于是 $D(f^{n + 1}/(n + 1)) \in \mathbf{C}[[x]]$ 且
$D(f^{n + 2}/(n + 2)) \in \mathbf{C}[[x]]$。因此可见
$f^{n + 1}, f^{n + 2} \in A$！特别地，(1) 成立。
我们还得出 $A$ 的分式域等于 $\mathbf{C}\{x\}$ 的分式域。
同时立即可见，$A \cap x\mathbf{C}\{x\}$ 是 $A$ 的非单位元集合，
故 $A$ 是维数为 $1$ 的局部整环。若能证明 (4)，便可推出 $A$ 是
Noether 环（证明从略）。设 $f \in A \cap x\mathbf{C}\{x\}$。
写成 $D(f) = h$，其中 $h \in \mathbf{C}[[x]]$。
再写成 $h = c + xh'$，其中 $c \in \mathbf{C}$、
$h' \in \mathbf{C}[[x]]$。则
$D(f - cxf_1) = c + xh' - c = xh'$。另一方面，
$f - cxf_1 = xg$，其中 $g \in \mathbf{C}\{x\}$；但由上述计算，
$D(g) = h' \in \mathbf{C}[[x]]$，因而 $g \in A$。
所以如所需，$f = cxf_1 + xg \in (x, xf_1)$。

\medskip\noindent
最后，为什么 $A$ 的完备化非约化？以 $\hat A$ 表示 $A$ 的完备化。
当然，它满射到 $\mathbf{C}\{x\}$ 的完备化 $\mathbf{C}[[x]]$，因为
$x \in A$。将此映射记为 $\psi : \hat A \to \mathbf{C}[[x]]$。
上面已经看到 $\mathfrak m_A = (x, xf_1)$，于是经简单计算有
$D(\mathfrak m_A^n) \subset (x^{n - 1})$。因此
$D : A \to \mathbf{C}[[x]]$ 连续，并在 $\psi$ 上诱导连续导子
$\hat D : \hat A \to \mathbf{C}[[x]]$。由此得到环映射
$$
\psi + \epsilon \hat D :
\hat A
\longrightarrow
\mathbf{C}[[x]][\epsilon].
$$
由于 $\hat A$ 是一维 Noether 完备局部环，只要证明该箭头满射，
便可推出 $\hat A$ 非约化。事实上该映射是同构，但我们省略对此的验证。
子环 $\mathbf{C}[x]_{(x)} \subset A$ 在完备化上诱导映射
$i : \mathbf{C}[[x]] \to \hat A$，使得
% P12-ERR-0326
$i \circ \psi = \text{id}$ 且 $D \circ i = 0$
（因为按构造 $D(x) = 0$）。考虑元素 $x^nf_n \in A$。有
$$
(\psi + \epsilon D)(x^nf_n) = x^n f_n + \epsilon
$$
对所有 $n \geq 1$ 均成立。由这些说明很容易推出满射性。





\section{另一个完备化非约化的局部环}
\label{examples-section-another-local-completion-nonreduced}

\noindent
本节构造一个维数为 $2$、关于某个主理想完备的 Noether 局部整环，
使得其一个局部化的再完备化非约化。

\medskip\noindent
设 $p$ 为素数。设 $k$ 为特征 $p$ 的域，并且 $k$ 在其 $p$ 次幂子域
$k^p$ 上的次数无限。例如 $k = \mathbf{F}_p(t_1, t_2, t_3, \ldots)$。
考虑环
$$
A =
\left\{
\begin{matrix}
\sum a_{i, j} x^iy^j \in k[[x, y]] \text{，满足对所有 }n \geq 0
\text{ 都有 } \\
[k^p(a_{n, n}, a_{n, n + 1}, a_{n + 1, n},
a_{n, n + 2}, a_{n + 2, n}, \ldots) : k^p] < \infty
\end{matrix}
\right\}
$$
作为集合，有
$$
k^p[[x, y]] \subset A  \subset k[[x, y]]
$$
$A$ 的每个元素 $f$ 都能唯一写成级数
$$
f = f_0 + f_1 xy + f_2 (xy)^2 + f_3 (xy)^3 + \ldots
$$
其中
$$
f_n = a_{n, n} + a_{n, n + 1} y + a_{n + 1, n} x +
a_{n, n + 2} y^2 + a_{n + 2, n} x^2 + \ldots
$$
而定义 $A$ 的公式中的条件意味着 $f_n$ 的系数生成 $k^p$ 的有限扩张。
从这一表示可清楚看出，$A$ 是 $k[[x, y]]$ 的一个
$k^p[[x, y]]$-子代数，并且关于理想 $xy$ 完备。此外，显然有
$$
A/xy A = C \times_k D
$$
其中 $k^p[[x]] \subset C \subset k[[x]]$ 和
$k^p[[y]] \subset D \subset k[[y]]$ 是《代数》例
\ref{algebra-example-bad-dvr-char-p} 中的幂级数子环。
因此，$C$ 和 $D$ 都是离散赋值环，从而 $A/ xy A$ 是 Noether 环。
由《代数》引理 \ref{algebra-lemma-completion-Noetherian}，可得 $A$ 是
Noether 环。由于 $\dim(k[[x, y]]) = 2$，再用《代数》引理
\ref{algebra-lemma-integral-sub-dim-equal}，可得 $\dim(A) = 2$。

\medskip\noindent
设 $f = \sum a_i x^i$ 为幂级数，使得
$k^p(a_0, a_1, a_2, \ldots)$ 在 $k^p$ 上的次数无限。
则 $f \not \in A$，但 $f^p \in A$。令
$$
B = A[f] \subset k[[x, y]]
$$
由于 $B$ 在 $A$ 上有限，$B$ 是 Noether 环。又因为 $B$ 关于由 $xy$
生成的理想完备；见《代数》引理 \ref{algebra-lemma-completion-tensor}。
事实上，$B$ 是以 $1, f, f^2, \ldots, f^{p - 1}$ 为基的自由 $A$-模；
证明从略。

\medskip\noindent
我们断言环
$$
(B_y)^\wedge = (B[1/y])^\wedge = \lim B[1/y]/(xy)^n B[1/y] =
\lim B[1/y] / x^n B[1/y]
$$
非约化。具体地，该环在
$$
(A_y)^\wedge = (A[1/y])^\wedge = \lim A[1/y]/(xy)^n A[1/y] =
\lim A[1/y] / x^n A[1/y]
$$
上自由，基为 $1, f, \ldots, f^{p - 1}$。然而，存在元素
$g \in (A_y)^\wedge$ 使得 $f^p = g^p$。事实上，只需取
$g = \sum a_i x^i$（与定义 $f$ 时使用相同的表达式），
该表达式在 $(A_y)^\wedge$ 中有意义。因此
$$
(B_y)^\wedge = (A_y)^\wedge[f]/(f^p - g^p) \cong
(A_y)^\wedge[\tau]/(\tau^p)
$$
非约化。事实上，这个例还说明稍强的结论。注意，$(A_y)^\wedge$
是以 $x$ 为一致化参数的离散赋值环，其剩余域是上述离散赋值环 $D$
的分式域。因此，甚至
$$
(B_y)^\wedge[1/(xy)] = ((B_y)^\wedge)_{xy}
$$
也非约化。由此得到如下类型的例。

\begin{lemma}
\label{examples-lemma-nonreduced-recompletion}
存在一个关于主理想 $I = (b)$ 完备的维数为 $2$ 的 Noether 局部整环
$(B, \mathfrak m)$，以及元素 $f \in \mathfrak m$、$f \not \in I$，
使得主元局部化 $B_f$ 的 $I$-进完备化
$C = (B_f)^\wedge$ 非约化，而且
$C_b = C[1/b] = (B_f)^\wedge[1/b]$ 也非约化。
\end{lemma}

\begin{proof}
见上述讨论。
\end{proof}





\section{一个非链状 Noether 局部环}
\label{examples-section-non-catenary-Noetherian-local}

\noindent
尽管 Noether 局部环有一套成功的维数理论，仍然存在非链状的 Noether
局部环。一个例见 \cite[附录，例 2]{Nagata}。事实上，我们将在最简单的
情形下给出这个例。具体地，我们将构造一个维数为 $2$、但并非泛链状的
Noether 局部整环 $A$。（注意，$A$ 自动是链状的；见《练习》中的练习
\ref{exercises-exercise-Noetherian-local-domain-dim-2-catenary}。）
一个并非泛链状的 Noether 局部环的存在性蕴含非链状 Noether 局部环的
存在性；在本节末尾，我们将对当前这个具体例说明这一点。

\medskip\noindent
设 $k$ 为域，并考虑 $k$ 上一个变量的形式幂级数环 $k[[x]]$。令
$$
z = \sum\nolimits_{i = 1}^\infty a_i x^i
$$
为形式幂级数。我们假设，作为 Laurent 级数域
$k((x)) = k[[x]][1/x]$ 的元素，$z$ 在 $k(x)$ 上超越。令
$$
z_j
=
x^{-j}(z - \sum\nolimits_{i = 1, \ldots, j - 1} a_i x^i)
=
\sum\nolimits_{i = j}^\infty a_i x^{i - j}
\in k[[x]].
$$
注意 $z = xz_1$。令 $R$ 为 $k[[x]]$ 中由 $x$、$z$ 以及所有 $z_j$
生成的子环，换言之，
$$
R = k[x, z_1, z_2, z_3, \ldots ] \subset k[[x]].
$$
考虑 $R$ 的理想 $\mathfrak m = (x)$ 和
$\mathfrak n = (x - 1, z_1, z_2 + a_1, z_3 + a_1 + a_2, \ldots)$。

\medskip\noindent
有 $xz_{j + 1} + a_j = z_j$。因此 $R/\mathfrak m = k$，且
$\mathfrak m$ 是极大理想。此外，$R$ 中不属于 $\mathfrak m$ 的任意元素
在 $k[[x]]$ 中映为单位元，故 $R_{\mathfrak m} \subset k[[x]]$。
% P12-ERR-0327
事实上，很容易推出 $R_{\mathfrak m}$ 是离散赋值环，且剩余域为 $k$。

\medskip\noindent
我们断言
$$
R/(x - 1) =
k[x, z_1, z_2, z_3, \ldots ]/(x - 1)
\cong
k[z].
$$
事实上，上述关系蕴含
$z_{j + 1} = z_j - a_j - (x - 1)z_{j + 1}$，因此在商
$R/(x - 1)$ 中，可以用 $z_j$ 表示 $z_{j + 1}$ 的类。
按构造，$R$ 的分式域在 $k$ 上的超越次数为 $2$，所以在
$R/(x - 1)$ 中 $z$ 在 $k$ 上超越，从而得到所需同构。
于是 $\mathfrak n = (x - 1, z)$，并且它是极大理想。事实上，映射
$$
k[x, x^{-1}, z]_{(x - 1, z)} \longrightarrow R_{\mathfrak n}
$$
是同构（因为 $x^{-1}$ 在 $R_{\mathfrak n}$ 中可逆，并且
$z_{j + 1} = x^{-1}z_j - a_j = \ldots = f_j(x, x^{-1}, z)$）。
这说明 $R_{\mathfrak n}$ 是维数为 $2$、剩余域为 $k$ 的正则局部环。

\medskip\noindent
令 $S$ 为乘法子集
$$
S =
(R \setminus \mathfrak m) \cap (R \setminus \mathfrak n) =
R \setminus (\mathfrak m \cup \mathfrak n)
$$
并令 $B = S^{-1}R$。我们断言：
\begin{enumerate}
\item 环 $B$ 是 $k$-代数。
\item 环 $B$ 的极大理想恰为两个理想 $\mathfrak mB$ 和 $\mathfrak nB$。
\item 在这两个极大理想处的剩余域均为 $k$。
\item 有 $B_{\mathfrak mB} = R_{\mathfrak m}$ 和
$B_{\mathfrak nB} = R_{\mathfrak n}$；它们分别是维数为 $1$ 和 $2$ 的
Noether 正则局部环。
\item 环 $B$ 是 Noether 环。
\end{enumerate}
验证细节从略。

\medskip\noindent
每当给定一个具有上述性质的 $k$-代数 $B$ 时，我们可以如下得到一个例。
取 $A = k + \text{rad}(B) \subset B$，其中
$\text{rad}(B) = \mathfrak mB \cap \mathfrak nB$ 是 Jacobson 根。
容易看出 $B$ 在 $A$ 上有限，因而由 Eakin 定理，$A$ 是 Noether 环
（见 \cite{Eakin}，或 \cite[附录 A1]{Nagata}，或
% P12-ERR-0328
在此插入将来的参考文献）。又因为 $A$ 是与 $B$ 具有相同分式域和剩余域
$k$ 的局部整环，而 $B$ 的维数为 $2$，由《代数》引理
\ref{algebra-lemma-integral-sub-dim-equal} 可知 $A$ 的维数也为 $2$。

\medskip\noindent
若 $A$ 是泛链状的，则维数公式《代数》引理
\ref{algebra-lemma-dimension-formula} 将给出
$\dim(B_{\mathfrak mB}) = 2$，矛盾。

\medskip\noindent
注意，$B$ 由一个元素在 $A$ 上生成。因此，对 $A[x]$ 的某个素理想
$\mathfrak p$，有 $B = A[x]/\mathfrak p$。令
$\mathfrak m' \subset A[x]$ 为对应于 $\mathfrak mB$ 的极大理想。
则一方面 $\dim(A[x]_{\mathfrak m'}) = 3$，另一方面
$$
(0)
\subset \mathfrak pA[x]_{\mathfrak m'}
\subset \mathfrak m'A[x]_{\mathfrak m'}
$$
是素理想的极大链。因此，$A[x]_{\mathfrak m'}$ 是一个非链状
Noether 局部环的例。





\section{不良 Noether 局部环的存在性}
\label{examples-section-bad}

\noindent
设 $(A, \mathfrak m, \kappa)$ 为 Noether 完备局部环。
文献 \cite{Lech} 证明：若 $\text{depth}(A) \geq 1$，并且 $A$ 包含
$\mathbf{Q}$ 或 $\mathbf{F}_p$ 作为子环，或者包含 $\mathbf{Z}$ 作为子环
且 $A$ 作为 $\mathbf{Z}$-模无挠，则 $A$ 是某个 Noether 局部整环的完备化。
这产生了许多具有“奇异”性质的 Noether 局部整环的例。

\medskip\noindent
例如，将此应用于 $A = \mathbf{C}[[x, y]]/(y^2)$，便得到一个完备化
非约化的 Noether 局部整环。请与第
\ref{examples-section-local-completion-nonreduced} 节比较。

\medskip\noindent
文献 \cite{LLPY} 给出了刻画 $A$ 何时为某个约化 Noether 局部环之
完备化的条件。

\medskip\noindent
文献 \cite{Heitmann-completion-UFD} 证明：若
$\text{depth}(A) \geq 2$，并且 $A$ 包含 $\mathbf{Q}$ 或
$\mathbf{F}_p$ 作为子环，或者包含 $\mathbf{Z}$ 作为子环且 $A$ 作为
$\mathbf{Z}$-模无挠，则 $A$ 是某个 Noether 局部唯一分解整环 $R$ 的
完备化。特别地，$R$ 正规（《代数》引理
\ref{algebra-lemma-UFD-normal}），因而 $R$ 的 Hensel 化也是正规整环
（More on Algebra，引理 \ref{more-algebra-lemma-henselization-normal}）。
所以，上述 $A$ 是某个 Hensel Noether 局部正规整环的完备化
（因为 $R$ 与其 Hensel 化的完备化相同；见 More on Algebra，引理
\ref{more-algebra-lemma-henselization-noetherian}）。

\medskip\noindent
应用这一结果，可找到一个 Noether 局部唯一分解整环 $R$，使得
$R^\wedge \cong \mathbf{C}[[x, y, z, w]]/(wx, wy)$。
注意，$\Spec(R^\wedge)$ 是一个维数为 $2$ 的正则分支和一个维数为 $3$ 的正则分支的并。
环 $R$ 不可能泛链状：令
$$
X \longrightarrow \Spec(R)
$$
为极大理想的爆破。于是 $X$ 是整概形。存在闭点 $x \in X$，使得
$\dim(\mathcal{O}_{X, x}) = 2$；具体地，在完备局部环的层面上，
我们选取 $x$ 位于维数为 $2$ 的分支的严格变换上，而不位于维数为 $3$ 的分支的严格变换上。
由《态射》引理 \ref{morphisms-lemma-dimension-formula}，可见 $R$ 并非
泛链状。请与第 \ref{examples-section-non-catenary-Noetherian-local} 节比较。

\medskip\noindent
上述环是链状的（因为它是维数为 $3$ 的 Noether 局部唯一分解整环）。
不过，在 \cite{Ogoma-example} 中，作者构造了一个 Noether 局部正规整环
$R$，满足 $R^\wedge \cong \mathbf{C}[[x, y, z, w]]/(wx, wy)$，
而 $R$ 不是链状的。另见 \cite{Heitmann-Ogoma} 和 \cite{Lech-YAPO}。

\medskip\noindent
文献 \cite{Heitmann-isolated} 证明：若 $A$ 包含 $\mathbf{Q}$ 或
$\mathbf{F}_p$ 作为子环，或者 $A$ 的剩余特征为 $p > 0$，且 $p$ 在
$A$ 的任意真局部化中都不可能映为非零零因子，则 $A$ 是某个具有
孤立奇点的 Noether 局部环 $R$ 的完备化。这里称 Noether 局部环 $R$
具有孤立奇点，是指对所有非极大素理想 $\mathfrak p \subset R$，
$R_\mathfrak p$ 都是正则局部环。

\medskip\noindent
论文 \cite{Nishimura-few} 和 \cite{Nishimura-few-II} 列出了许多具有
给定完备化的“不良” Noether 局部环。特别地，
% P12-ERR-0329
其中构造了一个维数为 $2$ 的 Nagata 局部正规整环，其完备化为
$\mathbf{C}[[x, y, z]]/(yz)$；还构造了一个完备化为
$\mathbf{C}[[x, y, z]]/(y^2 - z^3)$ 的例。

\medskip\noindent
顺便一提，文献 \cite{Loepp} 证明：若 $A$ 约化、等维，且 $A$ 中没有
整数是零因子，则 $A$ 是某个优秀 Noether 局部整环的完备化。
不过，这不会产生“不良” Noether 局部环，因为得到的是优秀环！





\section{Noether Jacobson 环中的维数}
\label{examples-section-noetherian-jacobson}

\noindent
设 $k$ 为有限域的代数闭包。令 $A = k[x, y]$ 且
$X = \Spec(A)$。令 $C = V(x)$ 为 $y$-轴（也可取 $X$ 中任意其他
维数为 $1$ 的整闭子概形）。把 $X$ 中其他维数为 $1$ 的整闭子概形
枚举为 $C_1, C_2, C_3, \ldots$，并把 $C$ 的闭点枚举为
$p_1, p_2, p_3, \ldots$。

\medskip\noindent
断言：对每个 $n$，存在维数为 $1$ 的不可约闭子集
$Z_n \subset X$，使得
$$
\{p_n\} = Z_n \cap (C \cup C_1 \cup \ldots \cup C_n)
$$
在集合论意义下成立。为此，令
$Y = C \cup C_1 \cup C_2 \cup \ldots \cup C_n$。
这是 $k$ 上维数为 $1$ 的约化仿射代数概形。只需找到
$f \in k[x, y]$，使得 $V(f) \cap  Y = \{p_n\}$ 在集合论意义下成立；
因为这时可取 $Z_n$ 为 $V(f)$ 的一个适当不可约分支。由于限制映射
$$
k[x, y] \longrightarrow \Gamma(Y, \mathcal{O}_Y)
$$
满射，只需在 $Y$ 上找到一个零点集在集合论意义下为 $\{p_n\}$ 的正则
函数 $g$。为说明这是可能的，选取支撑为 $p_n$ 的有效 Cartier 除子
$D \subset Y$（这可由《簇》引理
\ref{varieties-lemma-complement-codim-1-closed-points} 得到）。
% P12-ERR-0330
因此，只需证明对某个 $N > 0$ 有
$\mathcal{O}_X(ND) \cong \mathcal{O}_X$。不过，有限域代数闭包上的仿射
$1$-维代数概形的 Picard 群是挠群
% P12-ERR-0331
（在此插入将来的参考文献），故断言成立。

\medskip\noindent
对所有 $n$ 如上选取 $Z_n$。由于 $k[x, y]$ 是唯一分解整环，
可对某个不可约元素 $f_n \in A$ 写成 $Z_n = V(f_n)$。
令 $S \subset k[x, y]$ 为由 $f_1, f_2, f_3, \ldots$ 生成的乘法子集。
考虑 Noether 环 $B = S^{-1}A$。

\medskip\noindent
显然，环映射 $A \to B$ 识别局部环，并诱导单射
$\Spec(B) \to \Spec(A)$。此外，考察曲线 $C_1$ 可见，在从
$\Spec(A)$ 过渡到 $\Spec(B)$ 时，仅删去了 $C \cap C_1$ 中的点。
特别地，$\Spec(B)$ 有无穷多个对应于 $A$ 的极大理想的极大理想。
另一方面，$xB$ 是极大理想，因为 $B/xB$ 的谱仅含一个素理想：
我们删去了 $C = V(x)$ 的所有闭点（但没有删去泛点）。
最后，对 $i \geq 1$ 考虑曲线 $C_i$。取不可约元 $g_i \in A$，使得
$C_i = V(g_i)$。若对某个 $n$ 有 $C_i = Z_n$，则 $g_iB$ 是单位理想。
否则，在从 $A$ 过渡到 $B$ 时，$C_i$ 的闭点除有限多个外都保留下来：
具体地，从 $C_i$ 中仅删去了
$(Z_1 \cup \ldots \cup Z_{i - 1} \cup C) \cap C_i$ 中的点。

\medskip\noindent
上述 $B$ 的素谱结构及《代数》引理
\ref{algebra-lemma-noetherian-dim-1-Jacobson} 表明 $B$ 是 Jacobson 环。
极大理想包括位于 $\Spec(B)$ 中的 $A$ 的极大理想
% P12-ERR-0332
（而这样的理想有无穷多个），再加上极大理想 $xB$。
因此可见，出现了维数为 $1$ 和 $2$ 的局部环。

\begin{lemma}
\label{examples-lemma-Noetherian-Jacobson}
存在一个 Jacobson、泛链状的 Noether 整环 $B$，它有极大理想
$\mathfrak m_1, \mathfrak m_2$，并且
$\dim(B_{\mathfrak m_1}) = 1$ 且 $\dim(B_{\mathfrak m_2}) = 2$。
\end{lemma}

\begin{proof}
$B$ 的构造如上。只需指出，由《代数》引理
\ref{algebra-lemma-localization-catenary} 和《态射》引理
\ref{morphisms-lemma-ubiquity-uc}，$B$ 是泛链状的。
\end{proof}





\section{底空间为 Noether 而概形非 Noether}
\label{examples-section-noetherian-not-noetherian}

\noindent
我们给出两个例，说明底拓扑空间为 Noether 空间的概形未必是 Noether 概形。

\begin{example}
\label{examples-example-many-variables}
设 $k$ 为域，并令 $A = k[x_1, x_2, x_3, \dots] /
(x_1^2, x_2^2, x_3^2, \dots)$。$A$ 的任意素理想都包含幂零元
$x_1, x_2, x_3, \dots$，所以
$\mathfrak p = (x_1, x_2, x_3, \dots)$ 是 $A$ 唯一的素理想。
因此，$\operatorname{Spec} A$ 的底拓扑空间是单点，特别地是 Noether
空间。然而，$\mathfrak p$ 显然不是有限生成的。
\end{example}

\begin{example}
\label{examples-example-many-mononials}
设 $k$ 为域，并令 $A \subseteq k[x, y]$ 为由 $k$ 和单项式
$\{xy^i\}_{i \ge 0}$ 生成的子环。$A$ 中不包含 $x$ 的素理想与
$A_x \cong k[x, x^{-1}, y]$ 的素理想一一对应。若 $\mathfrak p$ 是
一个包含 $x$ 的素理想，则它包含每个 $xy^i$（$i \ge 0$），因为
$(xy^i)^2 = x(xy^{2i}) \in \mathfrak p$ 且 $\mathfrak p$ 是根理想。
因此，$\mathfrak p = (\{xy^i\}_{i \ge 0})$。所以，
$\operatorname{Spec} A$ 的底拓扑空间是 Noether 空间，因为它由
Noether 概形 $\Spec(A[x, x^{-1}, y])$ 的点以及素理想 $\mathfrak p$
组成。但是，环 $A$ 非 Noether，因为 $\mathfrak p$ 不是有限生成的。
注意，在这个例中，$A$ 还是整环。
\end{example}





\section{正规化拟仿射的非拟仿射簇}
\label{examples-section-nonquasi-affine}

\noindent
文献 \cite[II 注 6.6.13]{EGA} 提到了这类例的存在性。该文献为此类例
援引 EGA 第五卷，但第五卷并未出版。

\medskip\noindent
设 $k$ 为域。令 $Y = \mathbf{A}^2_k \setminus \{(0, 0)\}$。
我们将构造有限满双有理态射
$\pi : Y \longrightarrow X$，其中 $X$ 是 $k$ 上的簇，且 $X$ 非拟仿射。
具体地，考虑 $Y$ 中的下列曲线：
$$
\begin{matrix}
C_1 & : & x = 0 \\
C_2 & : & y = 0
\end{matrix}
$$
注意 $C_1 \cap C_2 = \emptyset$。选取同构
$\varphi : C_1 \to C_2$，$(0, y) \mapsto (y^{-1}, 0)$。
我们断言，存在唯一的上述态射 $\pi : Y \to X$，使得
$$
\xymatrix{
C_1
\ar@<1ex>[rr]^{\text{id}} \ar@<-1ex>[rr]_{\varphi}
& &
Y \ar[r]^\pi & X
}
$$
是簇范畴（甚至概形范畴）中的余等化子图。暂且接受这一点；我们来证明
这样的 $X$ 不可能拟仿射。事实上，显然有
$$
\Gamma(X, \mathcal{O}_X) =
\{ f \in k[x, y] \mid f(0, y) = f(y^{-1}, 0)\} =
k \oplus (xy) \subset k[x, y].
$$
特别地，这些函数不能分离点 $(1, 0)$ 和 $(-1, 0)$；它们在 $X$ 中的像
（如下所示）彼此不同（若 $k$ 的特征不是 $2$）。

\medskip\noindent
为证明 $X$ 存在，考虑 $Y$ 的 Zariski 开集 $D(x + y) \subset Y$。
它是环 $k[x, y, 1/(x + y)]$ 的谱，而曲线 $C_1$、$C_2$ 完全包含在
$D(x + y)$ 中。此外，态射
$$
C_1 \amalg C_2
\longrightarrow
D(x + y) \cap Y = \Spec(k[x, y, 1/(x + y)])
$$
是闭浸入。由 More on Algebra，引理
\ref{more-algebra-lemma-fibre-product-finite-type}，环
$$
A =
\{f \in k[x, y, 1/(x + y)] \mid f(0, y) = f(y^{-1}, 0)\}
$$
在 $k$ 上有限型。另一方面，$Y$ 的开集 $D(xy) \subset Y$ 与曲线
$C_1$、$C_2$ 不交。它是环
$$
B = k[x, y, 1/xy].
$$
% P12-ERR-0333
注意 $A_{xy} \cong B_{x + y}$（因为 $A$ 显然包含对任意多项式 $P$ 的元素
$xyP(x, y)$，以及元素 $xy/(x + y)$）。概形 $X$ 由仿射概形
$\Spec(A)$ 和 $\Spec(B)$ 通过同构 $A_{xy} \cong B_{x + y}$ 粘合而得，
因而显然在 $k$ 上有限型。为说明它是分离的，需要证明环映射
$A \otimes_k B \to B_{x + y}$ 满射。为此，使用
$A \otimes_k B$ 包含元素 $xy/(x + y) \otimes 1/xy$，它映到
$1/(x + y)$。态射 $Y \to X$ 由自然映射
$D(x + y) \to \Spec(A)$ 和 $D(xy) \to \Spec(B)$ 给出。
由于两者均有限，可得 $Y \to X$ 如所需是有限的。我们省略验证 $X$
确为上面所示图的余等化子；不过，见
% P12-ERR-0334
（在此插入关于概形范畴中推出的将来参考文献）。注意，态射
$\pi : Y \to X$ 确实把点 $(1, 0)$ 和 $(-1, 0)$ 映到 $X$ 中不同的点，
因为函数 $(x + y^3)/(x + y)^2 \in A$ 在这两点的值分别为
$1/1$ 与 $-1/(-1)^2 = -1$，它们总是不同
（除非特征为 $2$；在特征 $2$ 时请自行寻找两个点）。
我们把上述讨论总结为一个引理。

\begin{lemma}
\label{examples-lemma-quasi-affine-normalization-not-quasi-affine}
设 $k$ 为域。存在一个簇 $X$，它的正规化拟仿射，但它自身非拟仿射。
\end{lemma}

\begin{proof}
见上述讨论以及
% P12-ERR-0335
（在此插入关于正规化的将来参考文献）。
\end{proof}




\section{取概形论像}
\label{examples-section-scheme-theoretic-image}

\noindent
设 $k$ 为域，$t$ 为变量。令 $Y = \Spec(k[t])$ 且
$X = \coprod_{n \geq 1} \Spec(k[t]/(t^n))$。以 $f : X \to Y$ 表示如下态射：
对每个 $n \geq 1$ 使用闭浸入
$\Spec(k[t]/(t^n)) \to \Spec(k[t])$。在此情形下：
\begin{enumerate}
\item $f$ 的概形论像（《态射》定义
\ref{morphisms-definition-scheme-theoretic-image}）是 $Y$。
另一方面，$f$ 的像是 $Y$ 中由 $t = 0$ 给出的闭点。
因此，$f$ 的概形论像的底闭子集不等于 $f$ 的像的闭包。
\item 概形论像的构造不与到开子概形
$V = \Spec(k[t, 1/t]) \subset Y$ 上的限制可交换。
事实上，$V$ 在 $X$ 中的原像为空，故
$f|_{f^{-1}(V)} : f^{-1}(V) \to V$ 的概形论像是空概形。
它不等于 $Y \cap V$。
\end{enumerate}





\section{局部闭子集的像}
\label{examples-section-non-chevalley}

\noindent
Chevalley 定理断言，仿射概形之间有限表示态射下可构造集的像仍然可构造；
见《代数》定理 \ref{algebra-theorem-chevalley} 和《态射》第
\ref{morphisms-section-constructible} 节。我们将看到，对局部闭子集的像，
同样的结论并不成立。

\medskip\noindent
设 $k$ 为特征 $0$ 的域。考虑投影态射
$$
f :
X = \Spec(k[t, x_1, x_2, \ldots, y_1, y_2, \ldots])
\longrightarrow
\Spec(k[x_1, x_2, \ldots, y_1, y_2, \ldots]) = Y
$$
这是有限表示态射。令 $Z$ 为 $X$ 中由下列等式定义的闭子集：
$$
x_1(t - 1) = 0,\quad
x_2(t - 1)(t - 2) = 0,\quad
x_3(t - 1)(t - 2)(t - 3) = 0,\quad \ldots
$$
令 $U = \bigcup_{j \geq 1} U_j$ 为 $X$ 的开集，其中
$$
U_j = \text{使 }y_j(t - 1)(t - 2) ... (t - j)\text{ 非零的点}
$$
则有
$$
f(Z \cap U_j) =
\text{使 }x_1, \ldots, x_j\text{ 均为零而 }y_j\text{ 非零的点}
$$
我们断言，$B = f(Z \cap U) = \bigcup_{j \geq 1} f(Z \cap U_j)$
不是 $Y$ 中有限多个局部闭子集的并。

\medskip\noindent
证明该断言。假设 $B = A_1 \cup \cdots \cup A_m$ 是由 $Y$ 中局部闭
子集组成的 $B$ 的有限覆盖。我们将对 $n$ 归纳证明 $m \geq n$。
基始情形 $n = 1$ 成立，因为 $B$ 非空。假设 $n > 1$，且对 $n - 1$
归纳假设成立。由于 $B$ 的闭包是 $(x_1 = 0)$，某个 $A_i$ 必须包含
$(x_1 = 0)$ 的某个非空开子集。于是 $A_i$ 必须在 $(x_1 = 0)$ 中开。
但这样的开子集不可能包含满足 $y_1 = 0$ 的点；事实上，对 $B$ 中的点，
$y_1 = 0$ 迫使 $x_2 = 0$，这表明 $B$ 不包含 $(x_1 = 0)$ 内点
$(x, y)$ 的任何邻域。因此，余下 $m - 1$ 个集合在限制后构成
$B \cap (y_1 = 0)$ 的可构造覆盖。然而，注意右移映射
$x_i \mapsto x_{i + 1}$、$y_i \mapsto y_{i + 1}$ 将 $B$ 与
$B \cap (y_1 = 0)$ 识别！故由归纳假设，$m - 1 \geq n - 1$，
% P12-ERR-0336
从而 $m \geq n$。这完成归纳步骤的证明，也就证明了该断言。

\begin{lemma}
\label{examples-lemma-no-chevalley}
存在仿射概形之间的有限表示态射 $f : X \to Y$，以及 $X$ 的局部闭子集
$T$，使得 $f(T)$ 不是 $Y$ 中有限多个局部闭子集的并。
\end{lemma}

\begin{proof}
见上述讨论。
\end{proof}





\section{不在闭子概形中开的局部闭子概形}
\label{examples-section-strange-immersion}

\noindent
这是《态射》例 \ref{morphisms-example-thibaut} 的副本。
下面给出一个不能表示为先开浸入、后闭浸入之复合的浸入的例。
设 $k$ 为域。令 $X = \Spec(k[x_1, x_2, x_3, \ldots])$，并令
$U = \bigcup_{n = 1}^{\infty} D(x_n)$。于是 $U \to X$ 是开浸入。
考虑理想
$$
I_n =
(x_1^n, x_2^n, \ldots, x_{n - 1}^n, x_n - 1, x_{n + 1}, x_{n + 2}, \ldots)
\subset
k[x_1, x_2, x_3, \ldots][1/x_n].
$$
注意，对任意 $m \not = n$，有
$I_n k[x_1, x_2, x_3, \ldots][1/x_nx_m] = (1)$。
因此，$D(x_n)$ 上的拟凝聚理想 $\widetilde I_n$ 在 $D(x_nx_m)$ 上
彼此相同；即当 $n \not = m$ 时，
$\widetilde I_n|_{D(x_nx_m)} = \mathcal{O}_{D(x_n x_m)}$。
故这些理想粘合为拟凝聚理想层
$\mathcal{I} \subset \mathcal{O}_U$。令 $Z \subset U$ 为对应于
$\mathcal{I}$ 的闭子概形。于是 $Z \to X$ 是浸入。

\medskip\noindent
我们断言，不能把 $Z \to X$ 分解为
$Z \to \overline{Z} \to X$，其中 $\overline{Z} \to X$ 是闭浸入，
而 $Z \to \overline{Z}$ 是开浸入。事实上，$\overline{Z}$ 必须由理想
$I \subset k[x_1, x_2, x_3, \ldots]$ 定义，并满足
$I_n = I k[x_1, x_2, x_3, \ldots][1/x_n]$。
然而，在所有 $I_n$ 内都有像的元素
$f \in k[x_1, x_2, x_3, \ldots]$ 只有 $0$！所以这样的 $I$ 不存在。

\medskip\noindent
态射 $Z \to X$ 还给出了浸入的概形论像表现不良的例。
事实上，上述论证表明浸入 $Z \to X$ 的概形论像是 $X$。
另一方面，可见：
\begin{enumerate}
\item $Z$ 在 $X$ 中并非拓扑稠密；并且
\item $Z = Z \cap U \to U$ 的概形论像就是 $Z$。
这不等于 $U \cap X = U$，因此在此情形下，概形论像的构造不与到
开集上的限制可交换。
\end{enumerate}





\section{合适开集的不存在性}
\label{examples-section-nonexistence-opens}

\noindent
本节补充《性质》第 \ref{properties-section-finding-affine-opens} 节的结果。

\medskip\noindent
设 $k$ 为域，并令 $A = k[z_1, z_2, z_3, \ldots]/I$，其中 $I$ 是由
所有两两乘积 $z_iz_j$（$i \not = j$，$i, j \in \mathbf{N}$）生成的理想。
令 $S = \Spec(A)$。令 $s \in S$ 为对应于极大理想 $(z_i)$ 的闭点。
我们断言，不存在在 $S \setminus \{s\}$ 中稠密的拟紧开集
$V \subset S \setminus \{s\}$。注意
$S \setminus \{s\} = \bigcup D(z_i)$。每个 $D(z_i)$ 都是以
$\eta_i$ 为泛点的不可约开集。由此，对所有 $i$ 都有 $\eta_i \in V$。
然而，$S \setminus \{s\}$ 的主仿射开集形如 $D(f)$，其中
$f \in (z_1, z_2, \ldots)$。于是对某个 $n$ 有
$f \in (z_1, \ldots, z_n)$，从而可见 $D(f)$ 只包含有限多个点
$\eta_i$。所以 $V$ 不可能拟紧。

\medskip\noindent
设 $k$ 为域，并令 $B = k[x, z_1, z_2, z_3, \ldots]/J$，其中 $J$
由乘积 $xz_i$（$i \in \mathbf{N}$）以及所有两两乘积
$z_iz_j$（$i \not = j$，$i, j \in \mathbf{N}$）生成。
令 $T = \Spec(B)$。考虑主开集 $U = D(x)$。我们断言，不存在拟紧开集
$V \subset T$，使得 $V \cap U = \emptyset$ 且 $V \cup U$ 在 $T$ 中稠密。
令 $t \in T$ 为对应于极大理想 $(x, z_i)$ 的闭点。
$U$ 在 $T$ 中的闭包是 $\overline{U} = U \cup \{t\}$。
因此，$V \subset \bigcup_i D(z_i)$ 是拟紧开集。由上一段的论证，
可见 $V$ 不可能在 $\bigcup D(z_i)$ 中稠密。

\begin{lemma}
\label{examples-lemma-complement-of-affine-does-not-contain-qc-dense-open}
% P12-ERR-0337
仿射概形中拟紧开集的不存在性：
\begin{enumerate}
\item 存在仿射概形 $S$ 和仿射开集 $U \subset S$，使得不存在拟紧开集
$V \subset S$ 满足 $U \cap V = \emptyset$ 且 $U \cup V$ 在 $S$ 中稠密。
\item 存在仿射概形 $S$ 和闭点 $s \in S$，使得
$S \setminus \{s\}$ 不包含拟紧稠密开集。
\end{enumerate}
\end{lemma}

\begin{proof}
见上述讨论。
\end{proof}

\noindent
令 $X$ 为上述仿射概形 $T$ 的两个副本沿仿射开集 $U$ 的粘合。
于是存在态射 $\pi : X \to T$ 以及 $X = U_1 \cup U_2$，使得
$\pi$ 把 $U_i$ 同构地映到 $T$，并把 $U_1 \cap U_2$ 同构地映到 $U$。
注意，$X$ 拟分离（由《概形》引理
\ref{schemes-lemma-characterize-quasi-separated}）且拟紧。
我们断言，不存在分离、稠密的拟紧开集 $W \subset X$。
事实上，考虑映到上述闭点 $t \in T$ 的两个闭点
$x_1 \in U_1$、$x_2 \in U_2$。令
$\tilde \eta \in U_1 \cap U_2$ 为映到 $U$ 的（唯一）泛点 $\eta$ 的泛点。
注意，$\tilde\eta \leadsto x_1$ 且 $\tilde\eta \leadsto x_2$，
它们位于特化 $\eta \leadsto t$ 之上。由于
$\pi|_W : W \to T$ 是分离的，由分离性的赋值判据《概形》引理
\ref{schemes-lemma-valuative-criterion-separatedness}，不能同时有
$x_1$ 和 $x_2 \in W$。不妨设 $x_1 \not \in W$。
则 $W \cap U_1$ 是 $U_1$ 中不含 $x_1$ 的拟紧（因为 $X$ 拟分离）
稠密开集。现在注意，存在概形的同构
$(T, t) \cong (S, s)$（把 $x$ 映到 $z_1$，并把 $z_i$ 映到
$z_{i + 1}$）。因此，由本节第一段得到矛盾。

\begin{lemma}
\label{examples-lemma-no-dense-separated-quasi-compact-open-in-qcqs}
存在一个拟紧且拟分离的概形 $X$，它不包含分离的拟紧稠密开集。
\end{lemma}

\begin{proof}
见上述讨论。
\end{proof}





\section{拟紧稠密开子概形的不存在性}
\label{examples-section-nonexistence-qc-dense-open-subscheme}

\noindent
设 $X$ 为域 $k$ 上拟紧且拟分离的代数空间。我们知道概形轨迹
$X' \subset X$ 是稠密开子空间；见《空间的性质》命题
\ref{spaces-properties-proposition-locally-quasi-separated-open-dense-scheme}。
事实上，当 $X$ 合理时，这一结果也成立；见《良态空间》命题
\ref{decent-spaces-proposition-reasonable-open-dense-scheme}。
一个自然问题是，能否找到 $X$ 的拟紧稠密开子概形。一般而言，
这是不可能的。

\medskip\noindent
假设 $k$ 的特征不是 2。令
$B = k[x, z_1, z_2, z_3, \ldots]/J$，其中 $J$ 由乘积
$xz_i$（$i \in \mathbf{N}$）以及所有两两乘积
$z_iz_j$（$i \not = j$，$i, j \in \mathbf{N}$）生成。
令 $U = \Spec(B)$。以 $0 \in U$ 表示所有坐标均为零的闭点。令
$$
j : R = \Delta \amalg \Gamma \longrightarrow U \times_k U
$$
其中 $\Delta$ 是 $U$ 在 $k$ 上的对角态射的像，并且
$$
\Gamma = \{((x, 0, 0, 0, \ldots), (-x, 0, 0, 0, \ldots))
\mid x \in \mathbf{A}^1_k, x \not = 0\}.
$$
显然，$s, t : R \to U$ 是 étale 的，因而 $j$ 是 étale 等价关系。
商 $X = U/R$ 是代数空间（《空间》定理
\ref{spaces-theorem-presentation}）。注意，$j$ 不是浸入，因为
$(0, 0) \in \Delta$ 位于 $\Gamma$ 的闭包中。所以 $X$ 不是概形。
另一方面，由于 $R$ 拟紧，$X$ 拟分离。以 $0_X$ 表示点 $0 \in U$
的像。我们断言 $X \setminus \{0_X\}$ 是概形，具体为
$$
X \setminus \{0_X\} =
\Spec\left(k[x^2, x^{-2}]\right) \amalg
\Spec\left(k[z_1, z_2, z_3, \ldots]/(z_iz_j)\right) \setminus \{0\}
$$
（细节从略）。另一方面，我们已在第
\ref{examples-section-nonexistence-opens} 节看到，右侧的概形不包含拟紧稠密开集。

\begin{lemma}
\label{examples-lemma-nonexistence-qc-dense-open-subscheme}
存在一个拟紧且拟分离的代数空间，它不包含拟紧稠密开子概形。
\end{lemma}

\begin{proof}
见上述讨论。
\end{proof}

\noindent
使用《空间》例 \ref{spaces-example-non-representable-descent} 的构造，
其方式与我们上面使用《空间》例
\ref{spaces-example-affine-line-involution} 构造的方式相同，可以得到一个拟紧、
拟分离且局部分离的代数空间，它不包含拟紧稠密开子概形。





\section{代数空间上的仿射概形}
\label{examples-section-embedding-affines}

\medskip\noindent
假设 $f : Y \to X$ 是概形之间的态射，$f$ 局部有限型，且 $Y$ 仿射。
则存在从 $Y$ 到 $X$ 上仿射 $n$-空间的浸入
$Y \to \mathbf{A}^n_X$。见稍为一般的《态射》引理
\ref{morphisms-lemma-quasi-affine-finite-type-over-S}。

\medskip\noindent
现在假设 $f : Y \to X$ 是代数空间之间的态射，$f$ 局部有限型，
而 $Y$ 是仿射概形。此时一般不能找到从 $Y$ 到 $X$ 上仿射
$n$-空间的浸入。

\medskip\noindent
第一个（不太友好的）反例是 $Y = \Spec(k)$ 和
$X = [\mathbf{A}^1_k/\mathbf{Z}]$，其中 $k$ 是特征零的域，
$\mathbf{Z}$ 通过平移 $(n, t) \mapsto t + n$ 作用在
$\mathbf{A}^1_k$ 上。事实上，对任意 $X$ 上的态射
$Y \to \mathbf{A}^n_X$，可以拉回到 $X$ 的覆盖 $\mathbf{A}^1_k$，
于是得到映入 $\mathbf{A}^{n + 1}_k$ 的无穷多个
% P12-ERR-0338
$\mathbf{A}^1_k$ 的不交并，而该映射不是浸入。

\medskip\noindent
第二个反例是 $Y = \mathbf{A}^1_k \to X = \mathbf{A}^1_k/R$，其中
$R = \{(t, t)\} \amalg \{(t, -t), t \not = 0\}$。
事实上，在此情形下，态射 $Y \to \mathbf{A}^n_X$ 将由 $Y$ 上的某些
正则函数 $f_1, \ldots, f_n$ 给出，因而 $Y$ 与覆盖
$\mathbf{A}^{n + 1}_k \to \mathbf{A}^n_X$ 的纤维积是概形
$$
\{(f_1(t), \ldots, f_n(t), t)\} \amalg
\{(f_1(t), \ldots, f_n(t), -t), t \not = 0\}
$$
它到 $\mathbf{A}^{n + 1}_k$ 的显然态射不是浸入。
注意，这给出了 $X$ 拟分离的反例。

\begin{lemma}
\label{examples-lemma-cannot-embed-into-affine}
存在代数空间之间的有限型态射 $Y \to X$，其中 $Y$ 仿射、$X$ 拟分离，
并且不存在 $X$ 上的浸入 $Y \to \mathbf{A}^n_X$。
\end{lemma}

\begin{proof}
见上述讨论。
\end{proof}





\section{拟凝聚模的推前}
\label{examples-section-push-quasi-coherent}

\noindent
我们在《概形》引理 \ref{schemes-lemma-push-forward-quasi-coherent} 中证明：
当 $f$ 拟紧且拟分离时，$f_*$ 把拟凝聚模变为拟凝聚模。
下面给出一些例，说明这两个条件都是必要的。

\medskip\noindent
假设 $Y = \Spec(A)$ 是仿射概形，并且
$X = \coprod_{n \in \mathbf{N}} Y$。令 $f : X \to Y$ 为显然的态射。
我们断言，$f_*\mathcal{O}_X$ 不是拟凝聚的。事实上，对 $a \in A$ 有
$$
f_*\mathcal{O}_X(D(a)) = \prod\nolimits_{n \in \mathbf{N}} A_a
$$
因此，要使 $f_*\mathcal{O}_X$ 拟凝聚，就必须对所有 $a \in A$ 有
$$
\prod\nolimits_{n \in \mathbf{N}} A_a
=
\left(\prod\nolimits_{n \in \mathbf{N}} A\right)_a
$$
这一般不成立。例如，若 $A = \mathbf{Z}$ 且 $a = 2$，则
$(1, 1/2, 1/4, 1/8, \ldots)$ 是左侧的元素，但不属于右侧。
注意，$f$ 是非拟紧的分离态射。

\medskip\noindent
设 $k$ 为域。令
$$
A = k[t, z, x_1, x_2, x_3, \ldots]/(tx_1z, t^2x_2^2z, t^3x_3^3z, \ldots)
$$
令 $Y = \Spec(A)$。令 $V \subset Y$ 为开子概形
$V = D(x_1) \cup D(x_2) \cup \ldots$。令 $X$ 为 $Y$ 的两个副本沿
$V$ 的粘合，并令 $f : X \to Y$ 为显然的态射。于是有正合序列
$$
0 \to f_*\mathcal{O}_X \to
\mathcal{O}_Y \oplus \mathcal{O}_Y \xrightarrow{(1, -1)} j_*\mathcal{O}_V
$$
其中 $j : V \to Y$ 是包含态射。由于
$$
A \longrightarrow \prod A_{x_n}
$$
是单射（细节从略），可见 $\Gamma(Y, f_*\mathcal{O}_X) = A$。
另一方面，映射
$$
A_t \longrightarrow \prod A_{tx_n}
$$
的核非零，因为它包含元素 $z$。所以
$\Gamma(D(t), f_*\mathcal{O}_X)$ 严格大于 $A_t$，因为它包含
$(z, 0)$。因此，$f_*\mathcal{O}_X$ 不是拟凝聚的。
注意，$f$ 拟紧但非拟分离。

\begin{lemma}
\label{examples-lemma-pushforward-quasi-coherent}
《概形》引理 \ref{schemes-lemma-push-forward-quasi-coherent} 的条件是最优的：
拟紧性假设与拟分离性假设都不能去掉。
\end{lemma}

\begin{proof}
见上述讨论。
\end{proof}





\section{非有限生成而茎为秩 1 有限自由的模}
\label{examples-section-nonfree}

\noindent
令 $R = \mathbf{Q}[x]$。在 $R$ 的分式域中取子模
$M = \sum_{n \in \mathbf{N}} \frac{1}{x - n}R$。
则 $M$ 不是有限生成的，但对 $R$ 的每个素理想 $\mathfrak p$，
作为 $R_{\mathfrak p}$-模有
$M_{\mathfrak p} \cong R_{\mathfrak p}$。

\medskip\noindent
一个风格类似的例是 $R = \mathbf{Z}$ 和
$M = \sum_{p \text{ 为素数}} \frac{1}{p} \mathbf{Z} \subset \mathbf{Q}$；
它等于分母 $b$ 非零且无平方因子的分数 $\frac{a}{b}$ 所成的集合。








\section{在茎中可逆但本身不可逆的理想}
\label{examples-section-locally-invertible-not-invertible}

\noindent
设 $A$ 为整环，$I \subset A$ 为非零理想。回顾一下，我们说 $I$ 可逆，
是指 $I$ 作为 $A$-模可逆。我们将构造这种情形的一个例：$I$ 不可逆，
但对每个素理想 $\mathfrak q \subset A$，
$I_\mathfrak q = (f) \subset A_\mathfrak q$ 都是（非零）主理想。
文献中有时把对所有素理想 $\mathfrak q$，$I_\mathfrak q$ 都是主理想的
性质表述为“$I$ 是局部主理想”。我们不能采用这一术语，因为我们的
“局部”总是指“在 Zariski 拓扑中局部”（或在当前所用的任何拓扑中局部）。

\medskip\noindent
令 $R = \mathbf{Q}[x]$，并令 $M = \sum \frac{1}{x - n}R$ 为第
\ref{examples-section-nonfree} 节构造的模。考虑环\footnote{环 $A$ 是一个局部环
均为 Noether 环的非 Noether 整环的例。}
$$
A = \text{Sym}^*_R(M)
$$
以及理想 $I = M A = \bigoplus_{d \geq 1} \text{Sym}^d_R(M)$。
由于 $M$ 作为 $R$-模不是有限生成的，$I$ 不可能作为 $A$ 中的理想由
有限多个元素生成。由于可逆模是有限生成的，这意味着 $I$ 不可逆。
另一方面，令 $\mathfrak p \subset R$ 为素理想。按构造，
$M_\mathfrak p \cong R_\mathfrak p$。因此
$$
A_\mathfrak p = \text{Sym}^*_{R_\mathfrak p}(M_\mathfrak p) \cong
\text{Sym}^*_{R_\mathfrak p}(R_\mathfrak p) =
R_\mathfrak p[T]
$$
作为分次 $R_\mathfrak p$-代数成立。由此，
$I_\mathfrak p \subset A_\mathfrak p$ 由非零因子 $T$ 生成。
所以，对任意素理想 $\mathfrak q \subset A$，$I_\mathfrak q$ 确由单个元素生成。

\begin{lemma}
\label{examples-lemma-locally-principal-not-invertible}
存在整环 $A$ 和非零理想 $I \subset A$，使得对所有素理想
$\mathfrak q \subset A$，$I_\mathfrak q \subset A_\mathfrak q$ 都是主理想，
但 $I$ 不是可逆 $A$-模。
\end{lemma}

\begin{proof}
见上述讨论。
\end{proof}





\section{非投射的有限平坦模}
\label{examples-section-finite-flat-not-projective}

\noindent
这是《代数》注 \ref{algebra-remark-warning} 的副本。
有限 $R$-模若在 $R$ 上平坦，并不自动投射。一个反例是：
$R = \mathcal{C}^\infty(\mathbf{R})$ 为 $\mathbf{R}$ 上无穷次可微函数环，
并且 $M = R_{\mathfrak m} = R/I$，其中
$\mathfrak m = \{f \in R \mid f(0) = 0\}$ 且
$I = \{f \in R \mid \exists \epsilon, \epsilon > 0 :
f(x) = 0\ \forall x, |x| < \epsilon\}$。

\medskip\noindent
态射 $\Spec(R/I) \to \Spec(R)$ 也是平坦但非开的闭浸入的一个例。

\begin{lemma}
\label{examples-lemma-finite-flat-non-projective}
特殊的平坦模：
\begin{enumerate}
\item 存在环 $R$ 和有限平坦 $R$-模 $M$，而该模不投射。
\item 存在平坦但非开的闭浸入。
\end{enumerate}
\end{lemma}

\begin{proof}
见上述讨论。
\end{proof}





\section{非局部自由的投射模}
\label{examples-section-projective-not-locally-free}

\noindent
我们给出两个例。一个例的秩介于 $0$ 与 $1$ 之间，另一个例的秩为
$\aleph_0$。

\begin{lemma}
\label{examples-lemma-ideal-generated-by-idempotents-projective}
设 $R$ 为环。设 $I \subset R$ 为由可数个幂等元生成的理想。
则 $I$ 作为 $R$-模是投射的。
\end{lemma}

\begin{proof}
设 $I = (e_1, e_2, e_3, \ldots)$，其中 $e_n$ 是 $R$ 的幂等元。
归纳地把 $e_{n + 1}$ 替换为 $e_n + (1 - e_n)e_{n + 1}$ 后，
可设 $(e_1) \subset (e_2) \subset (e_3) \subset \ldots$，从而
$I = \bigcup_{n \geq 1} (e_n) = \colim_n e_nR$。在此情形下，
$$
\Hom_R(I, M) = \Hom_R(\colim_n e_nR, M)
= \lim_n \Hom_R(e_nR, M) = \lim_n e_nM
$$
注意，转移映射 $e_{n + 1}M \to e_nM$ 由乘以 $e_n$ 给出，并且满射。
所以，由《代数》引理 \ref{algebra-lemma-ML-exact-sequence}，函子
$\Hom_R(I, M)$ 正合；即 $I$ 是投射 $R$-模。
\end{proof}

\noindent
假设 $P \subset Q$ 是 $R$-模的包含，$Q$ 是有限 $R$-模，而 $P$
局部自由；见《代数》定义 \ref{algebra-definition-locally-free}。
假设 $Q$ 作为 $R$-模可由 $N$ 个元素生成。由《代数》引理
\ref{algebra-lemma-map-cannot-be-injective} 可知，$P$ 有限局部自由
（各自由部分的秩至多为 $N$）。在这种情形下，$P$ 是有限 $R$-模；
见《代数》引理 \ref{algebra-lemma-finite-projective}。

\medskip\noindent
结合上述结果，可见由可数个幂等元生成的非有限生成理想是投射的，
但不是局部自由的。一个显式例是
$R = \prod_{n \in \mathbf{N}} \mathbf{F}_2$，而 $I$ 是由幂等元
$$
e_n = (1, 1, \ldots, 1, 0, \ldots )
$$
生成的理想，其中 $1$ 所成序列的长度为 $n$。

\begin{lemma}
\label{examples-lemma-ideal-projective-not-locally-free}
存在环 $R$ 和理想 $I$，使得 $I$ 作为 $R$-模投射，但作为
$R$-模不是局部自由的。
\end{lemma}

\begin{proof}
见上文。
\end{proof}

\begin{remark}
\label{examples-remark-projective-with-dual-finite-not-finite}
上面 $R = \prod_{n \in \mathbf{N}} \mathbf{F}_2$、$I$ 是由幂等元
$e_n = (1, 1, \ldots, 1, 0, \ldots )$ 生成的理想这一例，还给出了
$M = I$ 的一个例：它是非有限生成投射模，而其对偶有限生成，甚至有限自由。
事实上，使用引理 \ref{examples-lemma-ideal-generated-by-idempotents-projective}
证明中的策略，读者可证明 $\Hom_R(I, R) \cong R$。
\end{remark}

\begin{lemma}
\label{examples-lemma-chow-group-product}
设 $K$ 为域。设 $C_i$（$i = 1, \ldots, n$）为 $K$ 上光滑、投射、
几何不可约的曲线。设 $P_i \in C_i(K)$ 为有理点，并设
$Q_i \in C_i$ 为满足 $[\kappa(Q_i) : K] = 2$ 的点。
则 $[P_1 \times \ldots \times P_n]$ 在
$\CH_0(U_1 \times_K \ldots \times_K U_n)$ 中非零，其中
$U_i = C_i \setminus \{Q_i\}$。
\end{lemma}

\begin{proof}
存在次数映射
$\deg : \CH_0(C_1 \times_K \ldots \times_K C_n) \to \mathbf{Z}$。
由于每个 $Q_i$ 在 $K$ 上的次数为 $2$，可见支撑在“边界”
$$
C_1 \times_K \ldots \times_K C_n
\setminus
U_1 \times_K \ldots \times_K U_n
$$
上的任意零维循环的次数都能被 $2$ 整除。
\end{proof}

\noindent
利用上述引理，还可如下构造另一个投射但非局部自由的模的例。
设 $C_n$（$n = 1, 2, 3, \ldots$）为 $\mathbf{Q}$ 上光滑、投射、
几何不可约的曲线，每条曲线上都有一对点 $P_n, Q_n \in C_n$，满足
$\kappa(P_n) = \mathbf{Q}$，且 $\kappa(Q_n)$ 是 $\mathbf{Q}$ 的二次扩张。
令 $U_n = C_n \setminus \{Q_n\}$；这是仿射曲线。
令 $\mathcal{L}_n$ 为 $U_n$ 上点 $P_n$ 的理想层之逆。
注意，在零维循环群 $\CH_0(U_n)$ 中，
$c_1(\mathcal{L}_n) = [P_n]$。令
$A_n = \Gamma(U_n, \mathcal{O}_{U_n})$，并令
$L_n = \Gamma(U_n, \mathcal{L}_n)$；这是 $A_n$ 上秩为 $1$ 的局部自由模。
令
$$
B_n = A_1 \otimes_{\mathbf{Q}} A_2 \otimes_{\mathbf{Q}} \ldots
\otimes_{\mathbf{Q}} A_n
$$
于是 $\Spec(B_n) = U_1 \times \ldots \times U_n$，所有积都在
$\Spec(\mathbf{Q})$ 上取。对 $i \leq n$，令
$$
L_{n, i} =
A_1 \otimes_{\mathbf{Q}} \ldots \otimes_{\mathbf{Q}} L_i
\otimes_{\mathbf{Q}} \ldots \otimes_{\mathbf{Q}} A_n
$$
它是秩为 $1$ 的局部自由 $B_n$-模。注意，
% P12-ERR-0339
这也是 $\text{pr}_i^*\mathcal{L}_n$ 的整体截面。令
$$
B_\infty = \colim_n B_n
\quad\text{且}\quad
L_{\infty, i} = \colim_n L_{n, i}
$$
最后，令
$$
M = \bigoplus\nolimits_{i \geq 1} L_{\infty, i}.
$$
这是有限局部自由模的直和，因而投射。我们断言 $M$ 不是局部自由的。
事实上，假设 $f \in B_\infty$ 是非零函数，使得 $M_f$ 在
$(B_\infty)_f$ 上自由。设 $e_1, e_2, \ldots$ 为一组基。
选取 $n \geq 1$，使得 $f \in B_n$。再选取 $m \geq n + 1$，
使得 $e_1, \ldots, e_{n + 1}$ 属于
$$
\bigoplus\nolimits_{1 \leq i \leq m} L_{m, i}.
$$
由于元素 $e_1, \ldots, e_{n + 1}$ 在忠实平坦基变换后成为一组基的一部分，
可知 Chern 类
$$
c_i(\text{pr}_1^*\mathcal{L}_1 \oplus \ldots \oplus
\text{pr}_m^*\mathcal{L}_m), \quad i = m, m - 1, \ldots, m - n
$$
在
$$
D(f) \subset U_1 \times \ldots \times U_m
$$
的 Chow 群中为零。由于 $f$ 是 $U_1 \times \ldots \times U_n$ 上一个
函数的拉回，这特别蕴含
$$
c_{m - n}(\mathcal{O}_W^{\oplus n} \oplus
\text{pr}_1^*\mathcal{L}_{n + 1} \oplus \ldots
\oplus \text{pr}_{m - n}^*\mathcal{L}_m) = 0.
$$
在域 $K = \mathbf{Q}(C_1 \times \ldots \times C_n)$ 上的簇
$$
W = (C_{n + 1} \times \ldots \times C_m)_K
$$
上成立。换言之，循环
$$
[(P_{n + 1} \times \ldots \times P_m)_K]
$$
在 $W$ 上的零维循环 Chow 群中为零。这与上面的引理
\ref{examples-lemma-chow-group-product} 矛盾，因为点 $Q_i$（
$n + 1 \leq i \leq m$）在
% P12-ERR-0340
$(C_n)_K$ 上诱导相应的点 $Q_i'$，并且由于
% P12-ERR-0341
$K/\mathbf{Q}$ 几何不可约，有 $[\kappa(Q_i') : K] = 2$。

\begin{lemma}
\label{examples-lemma-projective-not-locally-free}
存在可数环 $R$ 和投射模 $M$，它是可数多个秩为 $1$ 的局部自由模的直和，
但 $M$ 不是局部自由的。
\end{lemma}

\begin{proof}
见上文。
\end{proof}





\section{具有非零平坦理想的零维局部环}
\label{examples-section-zero-dimensional-flat-ideal}

\noindent
文献 \cite{Lazard} 和 \cite{Autour} 给出了具有非零平坦理想的零维
局部环的例。构造如下。设 $k$ 为域。设 $X_i, Y_i$（$i \geq 1$）为变量。
取 $R = k[X_i, Y_i]/(X_i - Y_i X_{i + 1}, Y_i^2)$。
分别以 $x_i$ 和 $y_i$ 表示 $X_i$ 和 $Y_i$ 在此环中的像。注意，在此环中
$$
x_i = y_i x_{i + 1} = y_i y_{i +1} x_{i + 2} =
y_i y_{i + 1} y_{i + 2} x_{i + 3} = \ldots
$$
环 $R$ 只有一个素理想，即 $\mathfrak m = (x_i, y_i)$。
我们断言，理想 $I = (x_i)$ 作为 $R$-模是平坦的。

\medskip\noindent
注意，$x_i$ 在 $R$ 中的零化理想为
$(x_1, x_2, x_3, \ldots, y_i, y_{i + 1}, y_{i + 2}, \ldots)$。
考虑由元素 $e_i$（$i \geq 1$）和关系 $e_i = y_i e_{i + 1}$ 生成的
$R$-模 $M$。$M$ 是平坦的，因为它是 $R$ 的诸副本经转移映射
$$
R \xrightarrow{y_1} R \xrightarrow{y_2} R \xrightarrow{y_3} \ldots
$$
所得的余极限 $\colim_i R$。注意，$e_i$ 在 $M$ 中的零化理想为
$(x_1, x_2, x_3, \ldots, y_i, y_{i + 1}, y_{i + 2}, \ldots)$。
由于 $M$ 与 $I$ 的每个元素都可分别写成 $f e_i$ 与 $h x_i$，
其中某些 $f, h \in R$，可见映射 $M \to I$、$e_i \to x_i$ 是同构，
且 $I$ 平坦。

\begin{lemma}
\label{examples-lemma-zero-dimensional-flat-ideal}
\begin{slogan}
具有平坦理想的零维环。
\end{slogan}
存在仅有一个素理想的局部环 $R$，以及非零理想 $I \subset R$，
它是平坦 $R$-模。
\end{lemma}

\begin{proof}
见上述讨论。
\end{proof}





\section{非满射的零维环上态射}
\label{examples-section-epimorphism-not-surjective}

\noindent
文献 \cite{Lazard-deux} 和 \cite{Autour} 中可找到如下例。
设 $k$ 为域。考虑环同态
$$
k[x_1, x_2, \ldots, z_1, z_2, \ldots]/
(x_i^{4^i}, z_i^{4^i})
\longrightarrow
k[x_1, x_2, \ldots, y_1, y_2, \ldots]/
(x_i^{4^i}, y_i - x_{i + 1}y_{i + 1}^2)
$$
它把 $x_i$ 映到 $x_i$，把 $z_i$ 映到 $x_iy_i$。注意，
$y_i^{4^{i + 1}}$ 在右侧为零，而 $y_1$ 非零（细节从略）。
该映射不满射：令 $\deg(x_i) = -1$、$\deg(z_i) = 0$、
$\deg(y_i) = 1$，便可把上述映射看成 $\mathbf{Z}$-分次代数的映射；
此时显然 $y_1$ 不在像中。最后，该映射是上态射，因为
$$
y_{i - 1} \otimes 1 = x_i y_i^2 \otimes 1 = y_i \otimes x_i y_i =
x_i y_i \otimes y_i = 1 \otimes x_i y_i^2.
$$
% P12-ERR-0342
所以，目标在源上的张量积同构于目标。

\begin{lemma}
\label{examples-lemma-epi-not-surjective}
存在维数为 $0$ 的局部环之间的上态射，而它不是满射。
\end{lemma}

\begin{proof}
见上述讨论。
\end{proof}





\section{有限型、非有限表示、在素理想处平坦}
\label{examples-section-ft-not-fp-flat-at-prime}

\noindent
设 $k$ 为域。考虑局部环 $A_0 = k[x, y]_{(x, y)}$。
记 $\mathfrak p_{0, n} = (y + x^n + x^{2n + 1})$。这是素理想。令
$$
A = A_0[z_1, z_2, z_3, \ldots]/(z_n z_m, z_n(y + x^n + x^{2n + 1}))
$$
注意，$A \to A_0$ 是满射，其核是平方零理想。因此，$A$ 也是局部环，
而 $A$ 的素理想与 $A_0$ 的素理想一一对应。以 $\mathfrak p_n$ 表示
$A$ 中对应于 $\mathfrak p_{0, n}$ 的素理想。注意，$\mathfrak p_n$ 是
$z_n$ 在 $A$ 中的零化理想。令
$$
C = A[z]/(xz^2 + z + y)[\frac{1}{2zx + 1}]
$$
注意，$A \to C$ 是 étale 环映射；见《代数》例
\ref{algebra-example-make-standard-smooth}。令 $\mathfrak q \subset C$
为由 $x$、$y$、$z$ 以及所有 $z_n$ 生成的极大理想。由于
$A \to C$ 平坦，$z_n$ 在 $C$ 中的零化理想为 $\mathfrak p_nC$。
计算得
\begin{align*}
C/\mathfrak p_n C
& =
A_0[z]/(xz^2 + z + y, y + x^n + x^{2n + 1})[1/(2zx + 1)] \\
& =
k[x]_{(x)}[z]/(xz^2 + z - x^n - x^{2n + 1})[1/(2zx + 1)] \\
& =
k[x]_{(x)}[z]/(z - x^n) \times
k[x]_{(x)}[z]/(xz + x^{n + 1} + 1)[1/(2zx + 1)] \\
& =
k[x]_{(x)} \times k(x)
\end{align*}
因为 $(z - x^n)(xz + x^{n + 1} + 1) = xz^2 + z - x^n - x^{2n + 1}$。
因此，$\mathfrak p_nC = \mathfrak r_n \cap \mathfrak q_n$，其中
$\mathfrak r_n = \mathfrak p_nC + (z - x^n)C$ 且
$\mathfrak q_n = \mathfrak p_nC + (xz + x^{n + 1} + 1)C$。
由于 $\mathfrak q_n + \mathfrak r_n = C$，还得到
$\mathfrak p_nC = \mathfrak r_n \mathfrak q_n$。
由此，$\mathfrak q_n$ 是 $\xi_n = (z - x^n)z_n$ 的零化理想。
注意，一方面 $\mathfrak r_n \subset \mathfrak q$，另一方面
$\mathfrak q_n + \mathfrak q = C$。例如，这可由 $\mathfrak q_n$ 是
$C$ 中不同于 $\mathfrak q$ 的极大理想推出。类似地，当
$n \not = m$ 时，有 $\mathfrak q_n + \mathfrak q_m = C$。
现在令
$$
B = \Im(C \longrightarrow C_{\mathfrak q})
$$
注意，由于 $xz + x^{n + 1} + 1$ 不属于 $\mathfrak q$，元素 $\xi_n$
在 $B$ 中映为零。以 $\mathfrak q' \subset B$ 表示 $\mathfrak q$ 的像。
按构造，$B$ 是有限型 $A$-代数，且
$B_{\mathfrak q'} \cong C_{\mathfrak q}$。特别地，
$B_{\mathfrak q'}$ 在 $A$ 上平坦。

\medskip\noindent
我们断言，不存在 $g' \in B$、$g' \not \in \mathfrak q'$，使得
$B_{g'}$ 在 $A$ 上有限表示。我们简述该断言的证明。
选取映到 $g' \in B$ 的元素 $g \in C$。考虑映射
$C_g \to B_{g'}$。由《代数》引理
\ref{algebra-lemma-finite-presentation-independent}，
% P12-ERR-0343
$B_g$ 在 $A$ 上有限表示，当且仅当 $C_g \to B_{g'}$ 的核有限生成。
但是，元素 $g \in C$ 不属于 $\mathfrak q$，故在
$A_0[z]/(xz^2 + z + y)$ 中映为非零元素。因此，$g$ 只能属于有限多个
素理想 $\mathfrak q_n$，因为素理想
$(y + x^n + x^{2n + 1}, xz + x^{n + 1} + 1)$ 构成二维不可约 Noether
空间 $\Spec(k[x, y, z]/(xz^2 + z + y))$ 的无穷多个余维 1 的点。
映射
$$
\bigoplus\nolimits_{g \not \in \mathfrak q_n} C/\mathfrak q_n
\longrightarrow
C_g, \quad
(c_n) \longrightarrow \sum c_n \xi_n
$$
是单射，其像是 $C_g \to B_{g'}$ 的核。我们省略这一陈述的证明。
（提示：把 $A$ 作为 $A_0$-模写成 $A = A_0 \oplus I$，其中 $I$ 是
$A \to A_0$ 的核。类似地，写成 $C = C_0 \oplus IC$。再写成
$IC = \bigoplus Cz_n \cong \bigoplus (C/\mathfrak r_n \oplus C/\mathfrak q_n)$，
并研究乘以 $g$ 对各直和项的作用。）
这完成了该断言证明的略述。这还证明，对上述任意 $g'$，$B_{g'}$
在 $A$ 上都不平坦。事实上，若它平坦，则 $z_n$ 在 $B_{g'}$ 中之像的
零化理想将是 $\mathfrak p_nB_{g'}$，并且不会包含 $z - x^n$。

\medskip\noindent
由此，我们可以对 \cite[第一部分，注 (3.4.7) (\romannumeral 5)]{GruRay}
中提出的问题给出否定回答。精确陈述如下。

\begin{lemma}
\label{examples-lemma-example-raynaud-gruson}
存在局部环 $A$、有限型环映射 $A \to B$，以及位于
$\mathfrak m_A$ 上方的素理想 $\mathfrak q$，使得
$B_{\mathfrak q}$ 在 $A$ 上平坦；并且对任意
$g \in B$、$g \not \in \mathfrak q$，环 $B_g$ 在 $A$ 上既非有限表示，
也不平坦。
\end{lemma}

\begin{proof}
见上述讨论。
\end{proof}





\section{有限型、平坦但非有限表示}
\label{examples-section-finite-type-flat-not-finite-presentation}

\noindent
本节给出一些有限型且平坦、但非有限表示的环映射和态射的例。

\medskip\noindent
设 $R$ 为具有理想 $I$ 的环，使得 $R/I$ 是有限平坦但非投射的模；
显式例见第 \ref{examples-section-finite-flat-not-projective} 节。
注意，这意味着 $I$ 不是有限生成的；见《代数》引理
\ref{algebra-lemma-finitely-generated-pure-ideal}。还要注意，
$I = I^2$；见《代数》引理 \ref{algebra-lemma-pure}。
各例的基环将取为 $R$，相应的基概形为 $S = \Spec(R)$。

\medskip\noindent
考虑环映射 $R \to R \oplus R/I\epsilon$，其中约定
$\epsilon^2 = 0$。这是有限平坦但非有限表示的环映射。
所有纤维环都是完全交且几何不可约。

\medskip\noindent
令 $A = R[x, y]/(xy, ay; a \in I)$。注意，作为 $R$-模，有
% P12-ERR-0344
$A = \bigoplus_{i \geq 0} Ry^i \oplus \bigoplus_{j > 0} R/Ix^j$。
所以 $R \to A$ 是平坦有限型但非有限表示的环映射。
每个纤维环同构于 $\kappa(\mathfrak p)[x, y]/(xy)$ 或
$\kappa(\mathfrak p)[x]$。

\medskip\noindent
可通过取 $B = R[X_0, X_1, X_2]/(X_1X_2, aX_2; a \in I)$，
把上一个例变成投射态射。在此情形下，
$X = \text{Proj}(B) \to S$ 是非有限表示的固有平坦态射，并且对每个
$s \in S$，纤维 $X_s$ 同构于 $\mathbf{P}^1_s$，或同构于
$\mathbf{P}^2_s$ 中由 $X_1X_2$ 的零化定义的闭子概形
（这是一条算术亏格为 $0$ 的投射结点曲线）。

\medskip\noindent
令 $M = R \oplus R \oplus R/I$。令 $B = \text{Sym}_R(M)$ 为 $M$
上的对称代数，并令 $X = \text{Proj}(B)$。于是
$X \to S$ 是非有限表示的固有平坦态射，并且对 $s \in S$，几何纤维
同构于 $\mathbf{P}^1_s$ 或 $\mathbf{P}^2_s$。特别地，这些纤维光滑且
几何不可约。

\begin{lemma}
\label{examples-lemma-finite-type-flat-not-finitely-presented}
存在下列各类例：
\begin{enumerate}
\item 非有限表示的平坦有限型环映射，其完全交纤维环几何不可约；
\item 非有限表示的平坦有限型环映射，其纤维环几何连通、几何约化、
维数为 1 且为完全交；
\item 非有限表示的概形固有平坦态射 $X \to S$，其每个纤维同构于
$\mathbf{P}^1_s$，或同构于 $\mathbf{P}^2_s$ 中 $X_1X_2$ 的零化轨迹；以及
\item 非有限表示的概形固有平坦态射 $X \to S$，其每个纤维同构于
$\mathbf{P}^1_s$ 或 $\mathbf{P}^2_s$。
\end{enumerate}
\end{lemma}

\begin{proof}
见上述讨论。
\end{proof}





\section{有限型环映射的拓扑}
\label{examples-section-topology-finite-type}

\noindent
设 $A \to B$ 为局部整环之间的局部映射。若 $A$ 为 Noether 环、
$A \to B$ 本质有限型，且 $A/\mathfrak m_A \subset B/\mathfrak m_B$
有限，则存在素理想 $\mathfrak q \subset B$、
$\mathfrak q \not = \mathfrak m_B$，使得 $A \to B/\mathfrak q$ 是某个
拟有限环映射的局部化。见 More on Morphisms，引理
\ref{more-morphisms-lemma-quasi-finite-quasi-section-meeting-nearby-open}。

\medskip\noindent
% P12-ERR-0345
本节给出一个例，说明当 $A$ 不再是 Noether 环时，上述结果不成立。
具体地，设 $k$ 为域，并令
$$
A = \{a_0 + a_1 x + a_2 x^2 + \ldots \mid a_0 \in k, a_i \in k((y))
\text{，对 }i\geq 1\}
$$
以及
$$
C = \{a_0 + a_1 x + a_2 x^2 + \ldots \mid a_0 \in k[y], a_i \in k((y))
\text{，对 }i\geq 1\}.
$$
包含 $A \to C$ 有限型，因为 $C$ 由 $y$ 在 $A$ 上生成。
我们断言，$A$ 是以
$\mathfrak m = \{a_1 x + a_2 x^2 + \ldots \in A\}$ 为极大理想的局部环，
且除 $(0)$ 和 $\mathfrak m$ 外没有其他素理想。事实上，$A$ 的元素
$f = a_0 + a_1 x + a_2 x^2 + \ldots$ 只要满足 $a_0 \not = 0$ 就可逆。
若 $\mathfrak q \subset A$ 是非零素理想，而
$f = a_i x^i + \ldots \in \mathfrak q$，则利用幂级数的性质可见，
对任意 $g \in k((y))$，元素
$g^{i + 1} x^{i + 1} \in \mathfrak q$；即 $gx \in \mathfrak q$。
这证明 $\mathfrak q = \mathfrak m$。

\medskip\noindent
至于环 $C$ 的谱，用与上面相同的论证可见，任意非零素理想都包含
位于 $\mathfrak m$ 上方的素理想
$\mathfrak p = \{a_1 x + a_2 x^2 + \ldots \in C\}$。
因此，$C$ 中位于 $(0)$ 上方的素理想只有 $(0)$。令
$\mathfrak m_C = yC + \mathfrak p$ 且 $B = C_{\mathfrak m_C}$。
则 $A \to B$ 就是所需的例。

\begin{lemma}
\label{examples-lemma-topology-finite-type}
存在局部整环之间的局部同态 $A \to B$，它本质有限型，且
$A/\mathfrak m_A \to B/\mathfrak m_B$ 有限，并且对 $B$ 的每个素理想
$\mathfrak q \not = \mathfrak m_B$，环映射 $A \to B/\mathfrak q$
都不是某个拟有限环映射的局部化。
\end{lemma}

\begin{proof}
见上述讨论。
\end{proof}





\section{纯但非泛纯}
\label{examples-section-pure-not-universally}

\noindent
设 $k$ 为域。令
$$
R = k[[x, xy, xy^2, \ldots]] \subset k[[x, y]].
$$
换言之，幂级数 $f \in k[[x, y]]$ 属于 $R$，当且仅当
$f(0, y)$ 是常数。特别地，$R[1/x] = k[[x, y]][1/x]$，而
$R/xR$ 是一个局部环，其极大理想的平方为零。
以 $R[y] \subset k[[x, y]]$ 表示满足 $f(0, y)$ 是 $y$ 的多项式的
幂级数 $f \in k[[x, y]]$ 所成的集合。于是 $R \to R[y]$ 是有限型但
非有限表示的环映射，并且在逆转 $x$ 后诱导同构。
还有一个满射 $R[y]/xR[y] \to k[y]$，其核的平方为零。
考虑有限表示环映射 $R \to S = R[t]/(xt - xy)$。
同样，$R[1/x] \to S[1/x]$ 是同构；在此情形下，
$S/xS \cong (R/xR)[t]/(xy)$ 映满 $k[t]$，且核幂零。
存在满射 $S \to R[y]$、$t \longmapsto y$，它在逆转 $x$ 后诱导同构，
并在模 $x$ 后诱导核幂零的满射。因此，$S \to R[y]$ 的核局部幂零。
特别地，$S \to R[y]$ 在谱上诱导泛同胚。

\medskip\noindent
首先，我们断言，$S$ 作为 $S$-模在 $R$ 上相对纯。
由于逆转 $x$ 后得到同构，只需在极大理想 $\mathfrak m \subset R$ 处
验证这一点。由于 $R$ 关于其极大理想完备，它是 Hensel 环；因此只需
验证：对每个素理想 $\mathfrak p \subset R$、
$\mathfrak p \not = \mathfrak m$，$S$ 中位于 $\mathfrak p$ 上方的唯一
素理想 $\mathfrak q$ 满足 $\mathfrak mS + \mathfrak q \not = S$。
由于 $\mathfrak p \not = \mathfrak m$，它对应于
$k[[x, y]][1/x]$ 的唯一素理想。因此，或者 $\mathfrak p = (0)$，
或者对某个不与 $x$ 相伴的不可约元素 $f \in k[[x, y]]$，有
$\mathfrak p = (f)$（这里使用 $k[[x, y]]$ 是唯一分解整环——
% P12-ERR-0346
在此插入将来的参考文献）。在第一种情形下，$\mathfrak q = (0)$，
结论显然。在第二种情形下，可将 $f$ 乘以一个单位元，使得
$f \in R[y]$（Weierstrass 预备定理；细节从略）。
于是容易看出 $R[y]/fR[y] \cong k[[x, y]]/(f)$，故 $f$ 在 $R[y]$
中定义素理想，且 $\mathfrak mR[y] + fR[y] \not = R[y]$。
由于 $S \to R[y]$ 在谱上诱导泛同胚，也就得到 $S$ 所需的结论。

\medskip\noindent
其次，我们断言，$S$ 在 $R$ 上并非泛相对纯。事实上，为说明这一点，
只需找到赋值环 $\mathcal{O}$ 和局部环映射
$R \to \mathcal{O}$，使得
$\Spec(R[y] \otimes_R \mathcal{O}) \to \Spec(\mathcal{O})$
不命中 $\Spec(\mathcal{O})$ 的闭点。等价地，需要找到
$\varphi : R \to \mathcal{O}$，使得 $\varphi(x) \not = 0$ 且
$v(\varphi(x)) > v(\varphi(xy))$，其中 $v$ 是 $\mathcal{O}$ 的赋值。
（因为这意味着 $y$ 的赋值为负。）为此，考虑按显然方式定义的环映射
$$
R
\longrightarrow
\{a_0 + a_1 x + a_2 x^2 + \ldots \mid a_0 \in k[y^{-1}], a_i \in k((y))\}
$$
可以找到一个赋值环 $\mathcal{O}$，它支配右侧在极大理想
$(y^{-1}, x)$ 处的局部化，由此即得结论。

\begin{lemma}
\label{examples-lemma-pure-not-universally-pure}
存在仿射概形之间的有限表示态射 $X \to S$，以及有限表示
$\mathcal{O}_X$-模 $\mathcal{F}$，使得 $\mathcal{F}$ 相对于 $S$ 纯，
但相对于 $S$ 并非泛纯。
\end{lemma}

\begin{proof}
见上述讨论。
\end{proof}





\section{非平坦的形式光滑环映射}
\label{examples-section-formally-smooth-nonflat}

\noindent
设 $k$ 为域。考虑 $k$-代数 $k[\mathbf{Q}]$。这是以
$x_\alpha, \alpha \in \mathbf{Q}$ 为基、乘法由
$x_\alpha x_\beta = x_{\alpha + \beta}$ 决定的 $k$-代数。
（特别地，$x_0 = 1$。）考虑 $k$-代数同态
$$
k[\mathbf{Q}] \longrightarrow k, \quad
x_\alpha \longmapsto 1.
$$
它满射，核 $J$ 由元素 $x_\alpha - 1$ 生成。我们来计算 $J/J^2$。
注意，乘以 $x_\alpha$ 在 $J/J^2$ 上是恒等映射。以 $z_\alpha$ 表示
$x_\alpha - 1$ 模 $J^2$ 的类。这些类生成 $J/J^2$。由于
$$
(x_\alpha - 1)(x_\beta - 1) = x_{\alpha + \beta} - x_\alpha - x_\beta + 1 =
(x_{\alpha + \beta} - 1) - (x_\alpha - 1) - (x_\beta - 1)
$$
可见在 $J/J^2$ 中，$z_{\alpha + \beta} = z_\alpha + z_\beta$。
$J/J^2$ 的一般元素形如 $\sum \lambda_\alpha z_\alpha$，其中
$\lambda_\alpha \in k$（只有有限多个非零）。注意，若 $k$ 的特征为
$p > 0$，则
$$
0 = pz_{\alpha/p} = z_{\alpha/p} + \ldots + z_{\alpha/p} = z_\alpha
$$
从而 $J/J^2 = 0$。若 $k$ 的特征为零，则
$$
J/J^2 = \mathbf{Q} \otimes_{\mathbf{Z}} k \cong k
$$
（细节从略），它非零。

\medskip\noindent
我们断言，当 $k$ 的特征为正时，$k[\mathbf{Q}] \to k$ 是形式光滑
环映射。事实上，假设给定实线交换图
$$
\xymatrix{
k \ar[r] \ar@{..>}[rd] & A \\
k[\mathbf{Q}] \ar[u] \ar[r]^\varphi & A' \ar[u]
}
$$
其中 $A' \to A$ 是满射，其核 $I$ 的平方为零。
为证明 $k[\mathbf{Q}] \to k$ 形式光滑，需要证明 $\varphi$ 经由 $k$
分解。由于 $\varphi(x_\alpha - 1)$ 在 $A$ 中映为零，$\varphi$ 诱导映射
$\overline{\varphi} : J/J^2 \to I$；该映射的零化恰是所需分解的障碍。
当特征为 $p > 0$ 时，$J/J^2 = 0$，所以得到所需结果，即此时
$k[\mathbf{Q}] \to k$ 形式光滑。最后，该环映射不平坦；例如，
非零因子 $x_2 - 1$ 被映为零。

\begin{lemma}
\label{examples-lemma-formally-smooth-nonflat}
存在形式光滑但非平坦的环映射。
\end{lemma}

\begin{proof}
见上述讨论。
\end{proof}





\section{非平坦的形式 étale 环映射}
\label{examples-section-formally-etale-nonflat}

\noindent
本节给出 \cite[0，例 19.10.3(i)]{EGA} 最后一句的反例
（这不在其后来的勘误表所捕获的条目之中）。考虑局部环 $A$ 及满足
$J^2 = J$ 的非零真理想 $J$ 所给出的 $A \to A/J$
（所以 $J$ 不是有限生成的）；一个代数闭非 Archimedean 域的赋值环，
取 $J$ 为其极大理想，就是这类 $(A, J)$ 的来源。
这些非平坦商映射形式 étale。事实上，假设给定交换图
$$
\xymatrix{
A/J \ar[r] & R/I \\
A \ar[u] \ar[r]^\varphi & R \ar[u]
}
$$
其中 $I$ 是环 $R$ 中满足 $I^2 = 0$ 的理想。则 $A \to R$ 唯一地经由
$A/J$ 分解，因为
$$
% P12-ERR-0347
\varphi(J) = \varphi(J^2) \subset (\varphi(J)A)^2 \subset I^2 = 0.
$$
因此，这也给出了 \cite[IV，定理 18.4.6(i)]{EGA} 中关于局部有限型、
形式 étale 态射“结构定理”的形式 étale 情形的反例
（但不是第 (ii) 部分的反例，而人们在实践中实际使用的正是第 (ii) 部分）。
后一证明中的错误是，证明的最后一步援引了错误的
% P12-ERR-0348
\cite[0，例 19.3.10(i)]{EGA}，上述反例正是由此混入的。

\begin{lemma}
\label{examples-lemma-formally-etale-not-presented}
存在形式 étale 的非平坦环映射。
\end{lemma}

\begin{proof}
见上述讨论。
\end{proof}





\section{余切复形非平凡的形式 étale 环映射}
\label{examples-section-formally-etale-nontrivial-cotangent-complex}

\noindent
设 $k$ 为域。考虑环
$$
R = k[\{x_n\}_{n \geq 1}, \{y_n\}_{n \geq 1}]/(
x_1y_1, x_{nm}^m - x_n, y_{nm}^m - y_n)
$$
令 $A$ 为在由所有 $x_n, y_n$ 生成的极大理想处的局部化，并以
$J \subset A$ 表示该极大理想。令 $B = A/J$。按构造，$J^2 = J$，
因而 $A \to B$ 形式 étale（见第
\ref{examples-section-formally-etale-nonflat} 节）。我们断言，元素
$x_1 \otimes y_1$ 是映射
$$
J \otimes_A J \longrightarrow J.
$$
之核中的非零元素。事实上，$(A, J)$ 是环
$$
R_n = k[x_n, y_n]/(x_n^n y_n^n)
$$
在相应极大理想处之局部化 $(A_n, J_n)$ 的余极限。
于是 $x_1 \otimes y_1$ 对应于元素
$x_n^n \otimes y_n^n \in J_n \otimes_{A_n} J_n$，而该元素非零
（由一个从略的显式计算可知）。由于 $\otimes$ 与余极限可交换，
即得结论。由 \cite[III 第 3.3 节]{cotangent}，$J$ 不是弱正则的。
因此，由 \cite[III 命题 3.3.3]{cotangent}，余切复形 $L_{B/A}$ 非零。
事实上，可以说得更精确。因为 $J/J^2 = 0$，有
$H_0(L_{B/A}) = \Omega_{B/A}$ 且 $H_1(L_{B/A}) = 0$。
但由 Quillen 基本谱序列的五项正合序列
（见《余切》注 \ref{cotangent-remark-elucidate-ss} 或
\cite[推论 8.2.6]{Reinhard}），以及
$\text{Tor}_2^A(B, B) = \Ker(J \otimes_A J \to J)$ 的非零性，
可得 $H_2(L_{B/A})$ 非零。

\begin{lemma}
\label{examples-lemma-formally-etale-nontrivial-cotangent-complex}
存在形式 étale 的满射环映射 $A \to B$，而 $L_{B/A}$ 不等于零。
\end{lemma}

\begin{proof}
见上述讨论。
\end{proof}





\section{平坦且形式非分歧并不蕴含形式 étale}
\label{examples-section-flat-formally-unramified-not-formally-etale}

\noindent
More on Morphisms，引理
\ref{more-morphisms-lemma-unramified-flat-formally-etale} 证明了概形的
非分歧平坦态射 $X \to S$ 形式 étale。本节的目标是给出两个例，
说明不能把“非分歧”替换为“形式非分歧”。第一个例利用完美环的特殊性质，
第二个例则说明，即使对 Noether 正则环之间的映射，该结果也会失效。

\begin{lemma}
\label{examples-lemma-perfect-closure-polynomial-ring}
设 $A = \mathbb{F}_p[T]$ 为 $\mathbb{F}_p$ 上单变量多项式环。
以 $A_{perf}$ 表示 $A$ 的完美闭包。则
$A \rightarrow A_{perf}$ 平坦且形式非分歧，但并非形式 étale。
\end{lemma}

\begin{proof}
注意，在 Frobenius 映射 $F_A : A \to A$ 下，目标副本 $A$ 是域上的
自由模，其基为 $\{1, T, \dots, T^{p - 1}\}$。因此，$F_A$ 忠实平坦；
所以，作为忠实平坦映射的余极限，$A \to A_{perf}$ 也忠实平坦。
由于 $A_{perf}$ 是完美环，相对 Frobenius
$F_{A_{perf}/A}$ 是满射。换言之，
$A_{perf} = A[A_{perf}^p]$，这立即蕴含
$\Omega_{A_{perf}/A} = 0$。于是，由 More on Morphisms，引理
\ref{more-morphisms-lemma-formally-unramified-differentials}，
$A \rightarrow A_{perf}$ 形式非分歧。

\medskip\noindent
只需证明 $A \rightarrow A_{perf}$ 不是形式光滑的。
注意，由于 $A$ 是光滑 $\mathbb{F}_p$-代数，余切复形
% P12-ERR-0349
$L_{A/\mathbb{F}_P} \simeq \Omega_{A/\mathbb{F}_p}[0]$
集中在次数 $0$；见《余切》引理 \ref{cotangent-lemma-when-projective}。
此外，由《余切》引理 \ref{cotangent-lemma-perfect-zero}，
$L_{A_{perf}/\mathbb{F}_p} = 0$ 在 $D(A_{perf})$ 中成立。
考虑 $D(A_{perf})$ 中余切复形的特异三角
$$
L_{A/\mathbb{F}_p} \otimes_A A_{perf} \to
L_{A_{perf}/\mathbb{F}_p} \to
L_{A_{perf}/A} \to
(L_{A/\mathbb{F}_p} \otimes_A A_{perf})[1]
$$
见《余切》第 \ref{cotangent-section-triangle} 节。得到
$L_{A_{perf}/A} = \Omega_{A/\mathbb{F}_p} \otimes_A A_{perf}[1]$；
即 $L_{A_{perf}/A}$ 等于置于次数 $-1$ 的秩为 $1$ 的自由
$A_{perf}$-模。因此，由 More on Morphisms，引理
\ref{more-morphisms-lemma-NL-formally-smooth} 和《余切》引理
\ref{cotangent-lemma-relation-with-naive-cotangent-complex}，
$A \rightarrow A_{perf}$ 不是形式光滑的。
\end{proof}

\noindent
下一个例也涉及素特征环，但可能更令人意外。缺点是，与前一个例相比，
它需要更多关于特征 $p$ 现象的知识。回顾一下，若素特征环 $A$ 的
Frobenius 映射有限，则称 $A$ 为 $F$-有限。

\begin{lemma}
\label{examples-lemma-completion-etale}
设 $(A, \mathfrak m, \kappa)$ 为素特征 $p > 0$ 的 Noether 局部环，
且 $[\kappa : \kappa^p] < \infty$。则从 $A$ 到其完备化的典范映射
$A \to A^\wedge$ 平坦且形式非分歧。不过，若 $A$ 正则但不优秀，
则该映射并非形式 étale。
\end{lemma}

\begin{proof}
完备化的平坦性是《代数》引理
\ref{algebra-lemma-completion-flat}。为证明该映射形式非分歧，
只需证明 $\Omega_{A^\wedge/A} = 0$；见《代数》引理
\ref{algebra-lemma-characterize-formally-unramified}。

\medskip\noindent
我们略述一个证明。选取 $x_1, \ldots, x_r \in A$，使它们映到
$\kappa$ 的一个 $p$-基 $\overline{x}_1, \ldots, \overline{x}_r$；
即 $\kappa$ 由 $\overline{x}_i$ 在 $\kappa^p$ 上极小生成。
选取 $\mathfrak m$ 的一组极小生成元 $y_1, \ldots, y_s$。
对每个 $n$，元素 $x_1, \ldots, x_r, y_1, \ldots, y_s$ 经 Frobenius
在 $(A/\mathfrak m^n)^p$ 上生成 $A/\mathfrak m^n$。略去若干细节。
由此，$F : A^\wedge \to A^\wedge$ 有限。因此，
$\Omega_{A^\wedge/A}$ 是有限 $A^\wedge$-模。另一方面，对任意
$a \in A^\wedge$ 和 $n$，可以找到 $a_0 \in A$，使得
$a - a_0 \in \mathfrak m^nA^\wedge$。于是
$\text{d}(a) \in \bigcap \mathfrak m^n \Omega_{A^\wedge/A}$；
由《代数》引理
\ref{algebra-lemma-intersect-powers-ideal-module-zero}，这蕴含
$\text{d}(a)$ 为零。因此，$\Omega_{A^\wedge/A} = 0$。

\medskip\noindent
假设 $A$ 正则。则由 Cohen 结构定理，
$x_1, \ldots, x_r, y_1, \ldots, y_s$ 是环 $A^\wedge$ 的一个
$p$-基；即有
$$
A^\wedge = \bigoplus\nolimits_{I, J} (A^\wedge)^p
x_1^{i_1} \ldots x_r^{i_r} y_1^{j_1} \ldots y_s^{j_s}
$$
其中 $I = (i_1, \ldots, i_r)$、$J = (j_1, \ldots, j_s)$，且
$0 \leq i_a, j_b \leq p - 1$。细节从略。特别地，
$\Omega_{A^\wedge}$ 是自由 $A^\wedge$-模，其基为
$\text{d}(x_1), \ldots, \text{d}(x_r), \text{d}(y_1), \ldots, \text{d}(y_s)$。

\medskip\noindent
现在，若 $A \to A^\wedge$ 形式 étale，甚至只需形式光滑，则由《代数》
命题 \ref{algebra-proposition-characterize-formally-smooth}，
$\NL_{A^\wedge/A}$ 在次数 $-1, 0$ 的上同调均为零。
对环映射 $\mathbf{F}_p \to A \to A^\wedge$ 使用 Jacobi--Zariski 序列
（《代数》引理 \ref{algebra-lemma-exact-sequence-NL}），得到同构
$\Omega_A \otimes_A A^\wedge \cong \Omega_{A^\wedge}$。
所以，$\Omega_A$ 以
$\text{d}(x_1), \ldots, \text{d}(x_r), \text{d}(y_1), \ldots, \text{d}(y_s)$
为自由基。考察分式域并利用 $A$ 正规，可得 $a \in A$ 是 $p$ 次幂，
当且仅当它在 $A^\wedge$ 中的像是 $p$ 次幂
（细节从略；使用《代数》引理
\ref{algebra-lemma-derivative-zero-pth-power}）。另一个推论是算子
$\partial/\partial x_a$ 和 $\partial/\partial y_b$ 在 $A$ 上有定义。

\medskip\noindent
% P12-ERR-0350
我们将证明，上述事实蕴含 $A$ 在 $A^p$ 上有限，并以
$x_1, \ldots, x_r, y_1, \ldots, y_s$ 为 $p$-基。
由 Kunz 的一个结果，这将与 $A$ 的非优秀性矛盾；见
\cite[推论 2.6]{Kun76}。具体地，设 $a \in A$，并写成
$$
a =  \sum\nolimits_{I, J} (a_{I, J})^p
x_1^{i_1} \ldots x_r^{i_r} y_1^{j_1} \ldots y_s^{j_s}
$$
其中 $a_{I, J} \in A^\wedge$。为完成证明，只需说明
$a_{I, J} \in A$。在等式两侧作用算子
$$
(\partial/\partial x_1)^{p - 1} \ldots
(\partial/\partial x_r)^{p - 1}
(\partial/\partial y_1)^{p - 1} \ldots
(\partial/\partial y_s)^{p - 1}
$$
可得 $a_{I, J}^p \in A$，其中
$I = (p - 1, \ldots, p - 1)$ 且 $J = (p - 1, \ldots, p - 1)$。
由上面的说明，这也蕴含 $a_{I, J} \in A$。把 $a$ 替换为
$a' = a - a_{I, J}^p x^I y^J$ 后，
% P12-ERR-0351
可以使用低 $1$ 阶的微分算子，得到另一个系数 $a_{I, J}$ 属于 $A$。
依此类推。
\end{proof}

\begin{remark}
\label{examples-remark-reference-existence-regular-nonexcellent-rings}
剩余域具有有限 $p$-基的非优秀正则环，甚至可以在特征 $p$ 的代数闭域
$k = \overline{k}$ 上的 $\mathbb{P}^2_k$ 的函数域中构造。
见 \cite[$\mathsection 4.1$]{DS18}。
\end{remark}

\noindent
引理 \ref{examples-lemma-completion-etale} 的证明实际上还说明了稍强的结论。

\begin{lemma}
\label{examples-lemma-excellent-regular-local-rings}
设 $(A, \mathfrak m, \kappa)$ 为特征 $p > 0$ 的正则局部环。
假设 $[\kappa : \kappa^p] < \infty$。则 $A$ 优秀，当且仅当
$A \to A^\wedge$ 形式 étale。
\end{lemma}

\begin{proof}
反向蕴含由引理 \ref{examples-lemma-completion-etale} 得到。对于正向蕴含，
注意，我们已经由引理 \ref{examples-lemma-completion-etale} 知道
$A \to A^\wedge$ 形式非分歧；等价地，
$\Omega_{A^\wedge/A}$ 为零。因此，当 $A$ 优秀时，只需证明完备化映射
形式光滑。由 Néron--Popescu 解奇异，$A \to A^\wedge$ 可写成光滑
$A$-代数的滤过余极限（《平滑化环映射》定理
\ref{smoothing-theorem-popescu}）。因此，$\NL_{A^\wedge/A}$ 在次数
$-1$ 的上同调为零。故由《代数》命题
\ref{algebra-proposition-characterize-formally-smooth}，
$A \to A^\wedge$ 形式光滑。
\end{proof}



\section{由幂等元集合生成的理想与局部化}
\label{examples-section-ideal-locally-idempotents}

\noindent
设 $R$ 为环。考虑环
$$
B(R) = R[x_n; n \in \mathbf{Z}]/(x_n(x_n - 1), x_nx_m; n \not = m)
$$
容易证明，每个素理想 $\mathfrak q \subset B(R)$ 或者形如
$$
\mathfrak q = \mathfrak pB(R) + (x_n; n \in \mathbf{Z})
$$
或者形如
$$
\mathfrak q =
\mathfrak pB(R) + (x_n - 1) + (x_m; n \not = m, m \in \mathbf{Z}).
$$
因此可见
$$
\Spec(B(R)) =
\Spec(R) \amalg \coprod\nolimits_{n \in \mathbf{Z}} \Spec(R)
$$
但其拓扑并不只是无交并拓扑。它具有如下性质：由
$n \in \mathbf{Z}$ 标号的每一份都是开子概形，确切地说就是标准开集
$D(x_n)$。“中央”的那一份 $\Spec(R)$ 位于其余任意无穷多份
$\Spec(R)$ 之并的闭包中。注意，最后这一份 $\Spec(R)$ 由理想
$(x_n, n \in \mathbf{Z})$ 截出，而该理想由幂等元 $x_n$ 生成。
因此，若 $\Spec(R)$ 连通，则上述分解恰好是 $\Spec(B(R))$
分解为连通分支的分解。

\medskip\noindent
其次，令
$A = \mathbf{C}[x, y]/((y - x^2 + 1)(y + x^2 - 1))$。
$A$ 的谱由两个不可约分支
$C_1 = \Spec(A_1)$、$C_2 = \Spec(A_2)$ 组成，其中
$A_1 = \mathbf{C}[x, y]/(y - x^2 + 1)$，且
$A_2 = \mathbf{C}[x, y]/(y + x^2 - 1)$。注意，它们分别由
$(x, y) = (t, t^2 - 1)$ 和 $(x, y) = (t, -t^2 + 1)$ 参数化，
并相交于 $P = (-1, 0)$ 和 $Q = (1, 0)$。我们可以构造 $B(A)$
的一个扭曲版本，按如下方式把 $B(A_1)$ 与 $B(A_2)$ 粘合起来：
在 $P$ 上方，使
$x_n \in B(A_1) \otimes \kappa(P)$
对应于 $x_n \in B(A_2) \otimes \kappa(P)$；但在 $Q$ 上方，使
$x_n \in B(A_1) \otimes \kappa(Q)$
对应于 $x_{n + 1} \in B(A_2) \otimes \kappa(Q)$。
把所得的 $A$-代数记作 $B^{twist}(A)$。细节从略。由构造，
$B^{twist}(A)$ 在 $A$ 上 Zariski 局部同构于未扭曲的版本。
具体而言，在主开集 $\Spec(A) \setminus \{P\}$ 和主开集
$\Spec(A) \setminus \{Q\}$ 上都是如此。然而，我们所选的粘合
产生了足够的“单值作用”，使得 $\Spec(B^{twist}(A))$ 连通
（细节从略）。最后，存在一个中央副本
% P12-ERR-0357
$\Spec(A) \to \Spec(B^{twist}(A))$，它给出一个闭子概形；其理想在
$B^{twist}(A)$ 上 Zariski 局部由
幂等元生成的理想截出，但整体上并非如此（因为 $B^{twist}(A)$
没有非平凡幂等元）。

\begin{lemma}
\label{examples-lemma-not-generated-idempotents}
存在一个仿射概形 $X = \Spec(A)$ 及其闭子概形 $T \subset X$，使得
$T$ 在 $X$ 上 Zariski 局部由幂等元生成的理想截出，但 $T$
并非由某个幂等元生成的理想截出。
\end{lemma}

\begin{proof}
见上述讨论。
\end{proof}



% P12-ERR-0352
\section{识别局部环但并非 ind-étale 的环映射}
\label{examples-section-not-ind-etale}

\noindent
注意，第 \ref{examples-section-ideal-locally-idempotents} 节构造的环映射
$R \to B(R)$ 是有限个 $R$ 的副本之积的余极限。因此
$R \to B(R)$ 是 ind-Zariski 的；见《Pro-étale 上同调》定义
\ref{proetale-definition-ind-zariski}。其次，考虑第
\ref{examples-section-ideal-locally-idempotents} 节构造的环映射
$A \to B^{twist}(A)$。由于该环映射在 $\Spec(A)$ 上 Zariski 局部
同构于某个 ind-Zariski 环映射 $R \to B(R)$，可知它识别局部环
（见《Pro-étale 上同调》引理
\ref{proetale-lemma-ind-zariski-implies}）。第
\ref{examples-section-ideal-locally-idempotents} 节的讨论表明，存在一个截面
$B^{twist}(A) \to A$，其核并非由幂等元生成。现在，若
$A \to B^{twist}(A)$ 是 ind-étale 的，即
$B^{twist}(A) = \colim A_i$，其中 $A \to A_i$ 为 étale，则
$A_i \to A$ 的核将由一个幂等元生成
（《代数》引理 \ref{algebra-lemma-map-between-etale} 和
\ref{algebra-lemma-surjective-flat-finitely-presented}）。
这与上述结果矛盾。

\begin{lemma}
\label{examples-lemma-not-ind-etale}
存在一个识别局部环但并非 ind-étale 的环映射 $A \to B$。
尤其，它并非 ind-Zariski。
\end{lemma}

\begin{proof}
见上述讨论。
\end{proof}




\section{与内射模伴随的非松弛拟凝聚层}
\label{examples-section-nonflasque}

\noindent
关于这类例子的更多内容，见 \cite[Exposé II，附录 I]{SGA6}，
其中 Illusie 解释了 Verdier 给出的一些例子。

\medskip\noindent
考虑仿射概形 $X = \Spec(A)$，其中
$$
A = k[x, y, z_1, z_2, \ldots]/(x^nz_n)
$$
是《性质》例 \ref{properties-example-does-not-work-in-general}
中的环。令 $I = (x) \subset A$。考虑 $X$ 的拟紧开集 $U = D(x)$。
我们已经在上述引文中看到，存在一个截面
$s \in \mathcal{O}_X(U)$，它对任意 $n \geq 0$ 都不来自某个
$A$-模映射 $I^n \to A$。

\medskip\noindent
令 $\alpha : A \to J$ 为 $A$ 到某个内射 $A$-模中的嵌入。
令 $Q = J/\alpha(A)$，并记商映射为 $\beta : J \to Q$。
我们断言，映射
$$
\Gamma(X, \widetilde{J})
\longrightarrow
\Gamma(U, \widetilde{J})
$$
不是满射。确切地说，我们断言 $\alpha(s)$ 不在其像中。
为证明这一点，使用反证法。假设 $x \in J$ 是一个在 $U$
上限制为 $\alpha(s)$ 的元素。于是 $\beta(x) \in Q$ 在 $U$
上的限制为 $0$。因此，对某个 $n$ 有 $I^n\beta(x) = 0$；见
《性质》引理
\ref{properties-lemma-sections-over-quasi-compact-open-in-affine}。
这蕴含我们得到一个态射
$\varphi : I^n \to A$，$h \mapsto \alpha^{-1}(hx)$。容易看出，
经《性质》引理
\ref{properties-lemma-sections-over-quasi-compact-open-in-affine}
中的映射，这个态射 $\varphi$ 给出截面 $s$，矛盾。

\begin{lemma}
\label{examples-lemma-nonflasque}
存在一个仿射概形 $X = \Spec(A)$ 和一个内射 $A$-模 $J$，使得
$\widetilde{J}$ 不是 $X$ 上的松弛层。即使 $U$ 是标准开集，限制映射
$\Gamma(X, \widetilde{J}) \to \Gamma(U, \widetilde{J})$
也未必是满射。
\end{lemma}

\begin{proof}
见上述讨论。
\end{proof}

\noindent
事实上，可以用类似构造得到一个内射模的例子，使其伴随拟凝聚层在某个
拟紧开集上具有非零上同调。具体地，从环
$$
A = k[x, y, w_1, u_1, w_2, u_2, \ldots]/(x^nw_n, y^nu_n, u_n^2, w_n^2)
$$
出发，其中 $k$ 为域。选取一个内射映射 $A \to I$，其中 $I$
为内射 $A$-模。我们断言，$A_{xy} \subset I_{xy}$ 中的元素
$1/xy$ 不在 $I_x \oplus I_y \to I_{xy}$ 的像中。反证之，假设
$$
\frac{1}{xy} = \frac{i}{x^n} + \frac{j}{y^n}
$$
对某个 $n \geq 1$ 及 $i, j \in I$ 成立。清除分母，得到
$$
(xy)^{n + m - 1} = x^my^{n + m}i + x^{n + m}y^mj
$$
对某个 $m \geq 0$ 成立。乘以 $u_{n + m}w_{n + m}$，可见
% P12-ERR-0353
$u_{n + m}w_{n + m}(xy)^{n + m - 1} = 0$ 在 $A$ 中成立，
这就是所需的矛盾。
令 $U = D(x) \cup D(y) \subset X = \Spec(A)$。对任意 $A$-模
$M$，由 Mayer--Vietoris 序列有正合列
$$
0 \to H^0(U, \widetilde{M}) \to M_x \oplus M_y \to M_{xy} \to
H^1(U, \widetilde{M}) \to 0
$$
由此可知 $H^1(U, \widetilde{I})$ 非零。

\begin{lemma}
\label{examples-lemma-nonvanishing}
存在一个仿射概形 $X = \Spec(A)$，其底层拓扑空间是 Noetherian 的；
还存在一个内射 $A$-模 $I$，使得 $\widetilde{I}$ 在 $X$ 的某个拟紧
开集 $U$ 上具有非零的 $H^1$。
\end{lemma}

\begin{proof}
见上述讨论。注意，作为拓扑空间，
$\Spec(A) = \Spec(k[x, y])$。
\end{proof}








\section{非分离平坦群概形}
\label{examples-section-non-separated-group-scheme}

\noindent
域上的每个群概形都是分离的；见《群胚》引理
\ref{groupoids-lemma-group-scheme-over-field-separated}。
对一般基上的群概形，这不成立。

\medskip\noindent
设 $k$ 为域。令 $S = \Spec(k[x]) = \mathbf{A}^1_k$。
令 $G$ 为原点 $0$ 加倍的仿射直线（见《概形》例
\ref{schemes-example-affine-space-zero-doubled}），并把它视为 $S$
上的概形。因此，$G \to S$ 的一个纤维或者是单点集，或者是一个二元集
（其中一个元素在 $U$ 中，另一个在 $V$ 中）。于是，可以在这些纤维上
赋予群结构（令 $U$ 中的元素为群结构的零元）。更精确地说，$G$
具有两个开集 $U, V$，它们同构地映到 $S$，而 $U \cap V$
同构地映到 $S \setminus \{0\}$。于是
$$
G \times_S G = U \times_S U \cup V \times_S U \cup
U \times_S V \cup V \times_S V
$$
其中每一片都同构于 $S$。因此，可以把乘法
$m : G \times_S G \to G$ 定义为唯一的 $S$-态射：它把第一片和
最后一片映入 $U$，把中间两片映入 $V$。这与上述逐点描述相符。
我们略去它确实定义群概形结构的验证。

\begin{lemma}
\label{examples-lemma-non-separated-group-scheme}
仿射直线上存在一个有限型平坦群概形，它不是分离的。
\end{lemma}

\begin{proof}
见上述讨论。
\end{proof}

\begin{lemma}
\label{examples-lemma-non-quasi-separated-group-scheme}
无穷维仿射空间上存在一个有限型平坦群概形，它不是拟分离的。
\end{lemma}

\begin{proof}
上述同一构造可以用无穷维仿射空间
$S = \mathbf{A}^\infty_k = \Spec k[x_1, x_2, \ldots]$ 作为基来实施，
并把与极大理想 $(x_1, x_2, \ldots)$ 对应的原点 $0 \in S$
作为在 $G$ 中加倍的闭点。所得群概形 $G \rightarrow S$
并非拟分离的，正如《概形》例
\ref{schemes-example-not-quasi-separated} 中所解释的那样。
\end{proof}



\section{单位分支平坦的非平坦群概形}
\label{examples-section-non-flat-group-scheme}

\noindent
设 $X \to S$ 为概形的单态射。令 $G = S \amalg X$。
令 $m : G \times_S G \to G$ 为 $S$-态射
$$
G \times_S G = X \times_S X \amalg X \amalg X \amalg S
\longrightarrow G = X \amalg S
$$
它把分量 $X \times_S X$ 和 $S$ 映入 $S$，并以恒等态射把两个
分量 $X$ 映入 $X$。这定义了一个群律。为看清这一点，需要证明
$m \circ (m \times \text{id}_G) = m \circ (\text{id}_G \times m)$
是 $G \times_S G \times_S G \to G$ 的两个相同映射。像上面一样把
$G \times_S G \times_S G$ 分解为各分量，可见只需在这 $8$ 片中的
每一片上验证。每片分别同构于 $S$、$X$、$X \times_S X$ 或
$X \times_S X \times_S X$ 之一。此外，两个映射分别把这些片映到
$S$、$X$、$S$、$X$。另一方面，$X \to S$ 是单态射这一事实蕴含
$X \times_S X \cong X$ 和 $X \times_S X \times_S X \cong X$，
并且在每种情形下恰有一个 $S$-态射 $S \to S$ 或 $X \to X$。
因此
$m \circ (m \times \text{id}_G) = m \circ (\text{id}_G \times m)$。
于是，取 $X \to S$ 为任意非平坦概形单态射（例如闭浸入），
便得到基 $S$ 上的一个群概形例子：其单位分支为 $S$（因而平坦），
但它本身不平坦。

\begin{lemma}
\label{examples-lemma-non-flat-group-scheme}
存在基 $S$ 上的群概形 $G$，其单位分支在 $S$ 上平坦，但它本身在
$S$ 上不平坦。
\end{lemma}

\begin{proof}
见上述讨论。
\end{proof}




\section{域上的非分离群代数空间}
\label{examples-section-non-separated-group-space}

\noindent
域上的每个群概形都是分离的；见《群胚》引理
\ref{groupoids-lemma-group-scheme-over-field-separated}。
但域上的群代数空间并非如此（不过，正面结果见本节末尾）。

\medskip\noindent
设 $k$ 为特征零的域。考虑《空间》例
\ref{spaces-example-affine-line-translation} 中的代数空间
$G = \mathbf{A}^1_k/\mathbf{Z}$。由构造，$G$ 是 $k$ 上的概形范畴中的预层
$$
T \longmapsto \Gamma(T, \mathcal{O}_T) / \mathbf{Z}
$$
在 fppf 拓扑下的伴随层。预层上显然的加法法则诱导出加法
$m : G \times G \to G$，它使 $G$ 成为 $\Spec(k)$ 上的群代数空间。
注意，$G$ 不是分离的（甚至不是拟分离或局部分离的）。另一方面，
$G \to \Spec(k)$ 却是有限型的！

\begin{lemma}
\label{examples-lemma-non-separated-group-space}
域上存在一个有限型群代数空间，它不是分离的（甚至不是拟分离或
局部分离的）。
\end{lemma}

\begin{proof}
见上述讨论。
\end{proof}

\noindent
正面结果如下：若群代数空间 $G$ 拟分离、局部分离，或者更一般地是
适度代数空间，则 $G$ 实际上是分离的；见《空间中的更多群胚》引理
\ref{spaces-more-groupoids-lemma-group-scheme-over-field-separated}。
此外，域上的有限型分离群代数空间实际上是概形；见《空间中的更多群胚》
引理
\ref{spaces-more-groupoids-lemma-group-space-scheme-locally-finite-type-over-k}。
证明的思路是，概形轨迹是稠密开集；见《空间的性质》命题
\ref{spaces-properties-proposition-locally-quasi-separated-open-dense-scheme}
或《适度空间》定理
\ref{decent-spaces-theorem-decent-open-dense-scheme}。
平移这个开集，便可看出 $G$ 的每一点都有一个为概形的开邻域。




% P12-ERR-0354
\section{étale 态射纤维中各点之间的专化}
\label{examples-section-specializations-fibre-etale}

\noindent
若 $f : X \to Y$ 是概形的 étale 态射，或更一般地是局部拟有限态射，
则纤维的各点之间不存在专化；见《态射》引理
\ref{morphisms-lemma-locally-quasi-finite-fibres}。
然而，对代数空间的态射，这一般不成立。

\medskip\noindent
为给出一个例子，设 $k$ 为域。令
$$
P = k[u, u^{-1}, y, \{x_n\}_{n \in \mathbf{Z}}].
$$
考虑由自同构 $\tau$ 生成的 $\mathbf{Z}$ 在 $P$ 上通过 $k$-代数映射
给出的作用，其中
$\tau(u) = u$、$\tau(y) = uy$ 且 $\tau(x_n) = x_{n + 1}$。
对 $d \geq 1$，令
$I_d = ((1 - u^d)y, x_n - x_{n + d}, n \in \mathbf{Z})$。
于是，
% P12-ERR-0355
$V(I_d) \subset \Spec(P)$ 是 $\tau^d$ 的不动点轨迹。
令 $S \subset P$ 为由 $y$ 及所有
$1 - u^d$（$d \in \mathbf{N}$）生成的乘法子集。于是可见，
$\mathbf{Z}$ 自由作用于 $U = \Spec(S^{-1}P)$。令
$X = U/\mathbf{Z}$ 为商代数空间；见《空间》定义
\ref{spaces-definition-quotient}。

\medskip\noindent
考虑 $S^{-1}P$ 中的素理想
$\mathfrak p_n = (x_n, x_{n + 1}, \ldots)$。注意
$\tau(\mathfrak p_n) = \mathfrak p_{n + 1}$。因此，它们中的每一个都
% P12-ERR-0356
定义一点 $\xi_n \in U$，其在 $X$ 中的像都是同一点 $x$，而该点属于
$X$。
此外，有专化
$$
\ldots \leadsto \xi_n \leadsto \xi_{n - 1} \leadsto \ldots
$$
因此，$U \to X$ 就是所求类型的例子。

\begin{lemma}
\label{examples-lemma-specializations-fibre-etale}
存在代数空间的 étale 态射 $f : X \to Y$，以及 $|X|$ 中点的一个
非平凡专化 $x \leadsto x'$，使得 $|Y|$ 中有
$f(x) = f(x')$。
\end{lemma}

\begin{proof}
见上述讨论。
\end{proof}







\section{不是 fppf 挠子的挠子}
\label{examples-section-torsor-not-fppf}

\noindent
在《群胚》注 \ref{groupoids-remark-fun-with-torsors} 中，我们提出了如下问题：
任意 $G$-挠子是否都是 fppf 拓扑下的 $G$-挠子？本节说明，情况并非
总是如此。

\medskip\noindent
设 $k$ 为域。以下所有概形和叠均在 $k$ 上。令
$G \to \Spec(k)$ 为群概形
$$
G = (\mu_{2, k})^\infty =
\mu_{2, k} \times_k \mu_{2, k} \times_k \mu_{2, k} \times_k \ldots =
\lim_n (\mu_{2, k})^n
$$
其中 $\mu_{2, k}$ 是 $\Spec(k)$ 上的二次单位根群概形；见《群胚》例
% P12-ERR-0358
\ref{groupoids-example-roots-of-unity}。由于 $G$ 是仿射概形的逆极限，
可见它是仿射群概形。事实上，它是环
$k[t_1, t_2, t_3, \ldots]/(t_i^2 - 1)$ 的谱。乘法映射
$m : G \times_k G \to G$ 在代数层面由
$t_i \mapsto t_i \otimes t_i$ 给出。

\medskip\noindent
我们断言，$k$ 上的任意 $G$-挠子都形如
$$
P = \Spec(k[x_1, x_2, x_3, \ldots]/(x_i^2 - a_i))
$$
其中 $a_i \in k^*$，而 $G$-作用 $G \times_k P \to P$
在代数层面由 $x_i \to t_i \otimes x_i$ 给出。证明从略。
事实上，对本例只需知道 $P$ 是 $G$-挠子；这是清楚的，因为在
$k' = k(\sqrt{a_1}, \sqrt{a_2}, \ldots)$ 上，概形 $P$ 以
$G$-等变方式同构于 $G$。注意，$P$ 平凡当且仅当 $k' = k$，
因为若 $P$ 有 $k$-有理点，则所有 $a_i$ 都是平方。

\medskip\noindent
我们断言，$P$ 是 fppf 挠子，当且仅当域扩张
$k' = k(\sqrt{a_1}, \sqrt{a_2}, \ldots)/k$ 有限。若 $k'$ 在 $k$
上有限，则
$\{\Spec(k') \to \Spec(k)\}$
是平凡化 $P$ 的 fppf 覆盖，因此 $P$ 确实是 fppf 挠子。反之，
假设 $P$ 是 fppf $G$-挠子。这意味着存在 fppf 覆盖
$\{S_i \to \Spec(k)\}$，使得每个 $P_{S_i}$ 都平凡。选取一个
$i$，使 $S_i$ 非空，并令 $s \in S_i$ 为闭点。由《簇》引理
\ref{varieties-lemma-locally-finite-type-Jacobson}，域扩张
$\kappa(s)/k$ 有限；由构造，$P_{\kappa(s)}$ 有一个
$\kappa(s)$-有理点。因此
$k \subset k' \subset \kappa(s)$，且 $k'$ 在 $k$ 上有限。

\medskip\noindent
为得到一个显式例子，例如可取 $k = \mathbf{Q}$ 和 $a_i = i$
（或者，如果愿意，可令 $a_i$ 为第 $i$ 个素数）。

\begin{lemma}
\label{examples-lemma-torsors-principal-spaces-not-equal}
设 $S$ 为概形，$G$ 为 $S$ 上的群概形。分类主齐 $G$-空间的叠
$G\textit{-Principal}$（见《叠的例子》小节
\ref{examples-stacks-subsection-principal-homogeneous-spaces}）与分类 fppf
$G$-挠子的叠 $G\textit{-Torsors}$（见《叠的例子》小节
\ref{examples-stacks-subsection-fppf-torsors}）一般并不等价。
\end{lemma}

\begin{proof}
上述讨论表明，函子
$G\textit{-Torsors} \to G\textit{-Principal}$ 一般不是本质满的。
\end{proof}








\section{具有拟紧平坦覆盖但非代数的叠}
\label{examples-section-not-algebraic-stack}

\noindent
本节简要描述 Brian Conrad 给出的一个例子。可以在
\href{https://mathoverflow.net/questions/15082/fpqc-covers-of-stacks/15269#15269}{这里}
在线找到该例。我们的例子略有不同。

\medskip\noindent
设 $k$ 为代数闭域。以下所有概形和叠均在 $k$ 上。
令 $G \to \Spec(k)$ 为仿射群概形。在《叠的例子》引理
\ref{examples-stacks-lemma-classifying-stacks} 中，我们给出了若干种彼此
等价的方式，把 $\mathcal{X} = [\Spec(k)/G]$ 看成
$(\Sch/\Spec(k))_{fppf}$ 上的群胚叠。特别地，$\mathcal{X}$ 分类
fppf $G$-挠子。更精确地说，一个 $1$-态射
$T \to \mathcal{X}$ 对应于 $T$ 上的一个 fppf $G_T$-挠子 $P$，而
$2$-箭头对应于挠子的同构。由此可知，对角 $1$-态射
$$
\Delta :
\mathcal{X}
\longrightarrow
\mathcal{X} \times_{\Spec(k)} \mathcal{X}
$$
可表且仿射。确切地说，给定概形 $T/k$ 上任意一对 fppf
$G_T$-挠子 $P_1, P_2$，概形 $\mathit{Isom}(P_1, P_2)$ 在 $T$
上仿射。$\Spec(k)$ 上平凡的 $G$-挠子定义一个 $1$-态射
$$
f : \Spec(k) \longrightarrow \mathcal{X}.
$$
我们断言，这是满的 $1$-态射。理由很简单：由定义，对任意 $1$-态射
$T \to \mathcal{X}$，存在一个 fppf 覆盖 $\{T_i \to T\}$，使得
$P_{T_i}$ 同构于平凡的 $G_{T_i}$-挠子。因此，复合
$T_i \to T \to \mathcal{X}$ 经过 $f$ 分解。于是显然，投影
$T \times_\mathcal{X} \Spec(k) \to T$ 是满射（这正是我们定义
$f$ 为满射的方式；见《代数叠》定义
\ref{algebraic-definition-relative-representable-property}）。
类似地可证明 $f$ 拟紧且平坦（细节从略）。还要在此记录如下观察：
$$
\Spec(k) \times_\mathcal{X} \Spec(k) \cong G
$$
是 $k$ 上概形的同构。

\medskip\noindent
假设存在满光滑态射 $p : U \to \mathcal{X}$，其中 $U$ 为概形。
考虑纤维积
$$
\xymatrix{
W \ar[d] \ar[r] & U \ar[d] \\
\Spec(k) \ar[r] & \mathcal{X}
}
$$
于是，$W$ 是 $k$ 上非空光滑概形，因而有一个 $k$-点。
这意味着 $f$ 可经过 $U$ 分解。因此得到
$$
G \cong
\Spec(k) \times_\mathcal{X} \Spec(k) \cong
(\Spec(k) \times_k \Spec(k))
\times_{(U \times_k U)}
(U \times_\mathcal{X} U)
$$
由于投影 $U \times_\mathcal{X} U \to U$ 按假设是光滑的，可知
$U \times_\mathcal{X} U \to U \times_k U$ 是局部有限型的；见
《态射》引理 \ref{morphisms-lemma-permanence-finite-type}。
因此，在这种情形下，$G$ 在 $k$ 上局部有限型。综上，我们证明了
下述引理（它可以显著推广）。

\begin{lemma}
\label{examples-lemma-BG-algebraic}
设 $k$ 为域，$G$ 为 $k$ 上的仿射群概形。若叠
$[\Spec(k)/G]$ 有一个由概形给出的光滑覆盖，则 $G$ 在 $k$ 上有限型。
\end{lemma}

\begin{proof}
见上述讨论。
\end{proof}

\noindent
为得到本节标题所述的显式例子，例如取第
\ref{examples-section-torsor-not-fppf} 节中的群概形
$G = (\mu_{2, k})^\infty$；它在 $k$ 上不是局部有限型。由上述讨论可见，
$\mathcal{X} = [\Spec(k)/G]$ 具有《代数叠》定义
\ref{algebraic-definition-algebraic-stack} 中的性质 (1) 和 (2)，
但不具有性质 (3)。因此 $\mathcal{X}$ 不是代数叠。另一方面，确实存在
一个概形 $U$ 和一个满、平坦、拟紧态射 $U \to \mathcal{X}$；它就是
上面研究过的态射 $f : \Spec(k) \to \mathcal{X}$。






\section{在对象上保持极限，但不保持极限}
\label{examples-section-limit-preserving}

\noindent
设 $S$ 为非空概形，$\mathcal{G}$ 为 $(\Sch/S)_{fppf}$ 上的内射 Abel 层。
我们得到 $S$ 上的群胚叠
$$
\mathcal{G}\textit{-Torsors} \longrightarrow (\Sch/S)_{fppf}
$$
见《叠的例子》引理
\ref{examples-stacks-lemma-torsors-sheaf-stack-in-groupoids}。
这个叠在 $(\Sch/S)_{fppf}$ 上的对象层面保持极限（见《可表性判据》第
\ref{criteria-section-limit-preserving} 节），因为每个
$\mathcal{G}$-挠子都是平凡的。另一方面，
$\mathcal{G}\textit{-Torsors}$ 一般不保持极限（见《Artin 公理》定义
\ref{artin-definition-limit-preserving}），因为作为层，$\mathcal{G}$
未必保持极限。例如，取任意非零内射层 $\mathcal{I}$，并令
$\mathcal{G} = \prod_{n \in \mathbf{Z}} \mathcal{I}$，便得到一个例子。

\begin{lemma}
\label{examples-lemma-limit-preserving-on-objects-not-limit-preserving}
设 $S$ 为非空概形。存在群胚叠
$p : \mathcal{X} \to (\Sch/S)_{fppf}$，使得 $p$ 在对象层面保持极限，
但 $\mathcal{X}$ 不保持极限。
\end{lemma}

\begin{proof}
见上述讨论。
\end{proof}



\section{一个非代数分类叠}
\label{examples-section-non-algebraic}

\noindent
令 $S = \Spec(\mathbf{F}_p)$，并以 $\mu_p$ 表示 $S$ 上的
$p$ 次单位根群概形。在《空间中的群胚》第
\ref{spaces-groupoids-section-stacks} 节中，我们引入了商叠
$[S/\mu_p]$；在《叠的例子》第
\ref{examples-stacks-section-group-quotient-stacks} 节中，我们证明了
$[S/\mu_p]$ 是 fppf $\mu_p$-挠子的分类叠：给定 $S$ 上概形 $T$，
范畴 $\Mor_S(T, [S/\mu_p])$ 典范等价于 $T$ 上 fppf
$\mu_p$-挠子的范畴。最后，我们已经在《可表性判据》定理
\ref{criteria-theorem-flat-groupoid-gives-algebraic-stack} 中看到，
$[S/\mu_p]$ 是代数叠。

\medskip\noindent
现在可以提出如下问题：“分类 étale $\mu_p$-挠子的群胚纤维范畴
$\mathcal{S}$ 又如何？”（换言之，$\mathcal{S}$ 是 $\Sch/S$ 上的范畴，
其在概形 $T$ 上方的纤维范畴是 $T$ 上 étale $\mu_p$-挠子的范畴。）

\medskip\noindent
第一个问题是，它并非 fppf 拓扑下的叠，因为对象的下降不会成立。
例如，$T = \Spec(\mathbf{F}_p(T))$ 上的 $\mu_p$-挠子
% P12-ERR-0359
$\Spec(\mathbf{F}_p(t)[x]/(x^p - t))$ 在 fppf 局部平凡，
但并非 étale 局部平凡。

\medskip\noindent
解决第一个问题的一种办法是改用 étale 拓扑；在这种情形下，对象的下降
确实成立。事实上，$\mathcal{S}$ 是 $(\Sch/S)_\etale$ 上的群胚叠。
此外，对角
$\Delta : \mathcal{S} \to \mathcal{S} \times \mathcal{S}$
也是（由概形）可表的。这是因为，给定两个 $\mu_p$-挠子（无论它们是否
étale 局部平凡），它们之间的同构层均由概形表示。

\medskip\noindent
因此，最后可以问：是否存在概形 $U$ 和光滑满 $1$-态射
$U \to \mathcal{S}$？我们将以两种方式证明这是不可能的：一种是直接
论证（建议读者跳过），另一种使用一般结果。

\medskip\noindent
直接论证（梗概）：注意，$1$-态射
$\mathcal{S} \to \Spec(\mathbf{F}_p)$ 满足形式光滑性的无穷小提升判据。
这是因为，给定概形的一阶无穷小加厚 $T \to T'$，映射
$\mu_p(T') \to \mu_p(T)$ 的核同构于 $T$ 在 $T'$ 中的理想层的截面，
因而
$H^1_\etale(T, \mu_p) = H^1_\etale(T', \mu_p)$。
此外，$\mathcal{S}$ 是保持极限的叠。因此，若
$U \to \mathcal{S}$ 光滑，则 $U \to \Spec(\mathbf{F}_p)$ 保持极限并满足
% P12-ERR-0360
形式光滑性的无穷小提升判据。这蕴含 $U$ 在 $\mathbf{F}_p$ 上光滑。
特别地，$U$ 是约化的，故
$H^1_\etale(U, \mu_p) = 0$。
因此，$U \to \mathcal{S}$ 分解为
$U \to \Spec(\mathbf{F}_p) \to \mathcal{S}$，且第一支箭头光滑。
由光滑性的下降可知，$U \to \mathcal{S}$ 光滑将蕴含
$\Spec(\mathbf{F}_p) \to \mathcal{S}$ 光滑。然而，事实并非如此，因为
$\Spec(\mathbf{F}_p) \times_\mathcal{S} \Spec(\mathbf{F}_p)$
是 $\mu_p$，它在 $\Spec(\mathbf{F}_p)$ 上不光滑。


\medskip\noindent
结构论证：在《可表性判据》第 \ref{criteria-section-stacks-etale} 节中，
我们已经看到，可把代数叠理解为 étale 拓扑下如下的群胚叠：其对角由
代数空间表示，并具有光滑覆盖。因此，若存在光滑满射
$U \to \mathcal{S}$，则 $\mathcal{S}$ 是代数叠，特别是满足 fppf
拓扑下的下降。但上面已经看到，$\mathcal{S}$
% P12-ERR-0361
不满足 fppf 拓扑下的下降。

\medskip\noindent
粗略地说，上述论证表明：对基 $B$ 上的群概形 $G$，étale 拓扑下
分类其 étale 局部平凡挠子的叠是代数的，当且仅当 $G$ 在 $B$
上光滑。采用 fppf 拓扑的一个优点是，只需假设
$G \to B$ 平坦且局部有限表示。事实上，fppf 拓扑下的商叠
$[B/G]$ 是代数的，当且仅当 $G \to B$ 平坦且局部有限表示；见
《可表性判据》引理 \ref{criteria-lemma-BG-algebraic}。







\section{具有拟紧平坦覆盖但非代数的层}
% P12-ERR-0362
\label{examples-section-not-algebraic}

\noindent
考虑函子 $F = (\mathbf{P}^1)^\infty$；也就是说，对概形 $T$，
$F(T)$ 是所有 $f = (f_1, f_2, f_3, \ldots)$ 的集合，其中每个
$f_i : T \to \mathbf{P}^1$ 都是概形态射。注意，$\mathbf{P}^1$
对 fpqc 覆盖满足层性质；见《下降》引理
\ref{descent-lemma-fpqc-universal-effective-epimorphisms}。
层的积仍是层，故 $F$ 也对 fpqc 拓扑满足层性质。$F$ 的对角可表：
若 $f : T \to F$ 和 $g : S \to F$ 为态射，则 $T \times_F S$
是概形 $T \times S$ 内各闭子概形
$T \times_{f_i, \mathbf{P}^1, g_i} S$ 的概形论交。考虑群概形
$\text{SL}_2$；它带有一个满光滑仿射态射
$\text{SL}_2 \to \mathbf{P}^1$。其次，考虑
$U = (\text{SL}_2)^\infty$ 及其典范（积）态射 $U \to F$。
注意，$U$ 是仿射概形。我们断言，态射 $U \to F$ 平坦、满且泛开。
确切地说，假设 $f : T \to F$ 为态射。于是
$Z = T \times_F U$ 是各概形
$Z_i = T \times_{f_i, \mathbf{P}^1} \text{SL}_2$ 在 $T$ 上的无穷纤维积。
每个态射 $Z_i \to T$ 都满、光滑且仿射，这蕴含
$$
Z = Z_1 \times_T Z_2 \times_T Z_3 \times_T \ldots
$$
% P12-ERR-0363
是一个在 $Z$ 上平坦且仿射的概形。一个简单的极限论证还表明
$Z \to T$ 是开的。

\medskip\noindent
另一方面，我们断言 $F$ 不是代数空间。事实上，若 $F$ 是代数空间，
则它将是拟紧且分离的代数空间（由我们对 $F$ 上纤维积的描述）。
因此，拟凝聚层的上同调将在某个截断次数以上消失（见《空间的上同调》
命题 \ref{spaces-cohomology-proposition-vanishing} 及其前面的评注）。
但显然，经投影
$$
(\mathbf{P}^1)^\infty \longrightarrow (\mathbf{P}^1)^n
$$
（它有一个截面）拉回 $\mathcal{O}(-2, -2, \ldots, -2)$，便可得到一个
拟凝聚层，其 $n$ 次上同调非零。综上，我们回答了 Anton Geraschenko
在 MathOverflow 上提出的一个问题。

\begin{lemma}
\label{examples-lemma-not-algebraic}
存在一个满足 fpqc 拓扑的层条件并具有可表对角
$\Delta : F \to F \times F$ 的函子
$F : \Sch^{opp} \to \textit{Sets}$；还存在概形 $U$ 及满、平坦、泛开、
拟紧态射 $U \to F$，但 $F$ 不是代数空间。
\end{lemma}

\begin{proof}
见上述讨论。
\end{proof}







\section{étale 拓扑与 Zariski 覆盖及有限 étale 覆盖的比较}
\label{examples-section-etale-zariski-finite-etale}

\noindent
本节给出 B.~Bhatt 为回答 C.~Deninger 的一个问题而找到的例子；
它说明，$\mathbf{C}$ 上代数簇的 étale 拓扑严格细于由开浸入和有限
étale 态射生成的拓扑。这类例子无疑为专家所知，也许可以一直追溯到
Artin 和 Grothendieck\footnote{据轶闻，Grothendieck 最初曾试图用
Zariski 开覆盖和有限 étale 覆盖定义 étale 拓扑，后来 Artin 说服他
放弃这种做法（大概正是依据本文所记录的这类例子）。}。
若读者知道已发表的例子，请发送电子邮件至
\href{mailto:stacks.project@gmail.com}{stacks.project@gmail.com}。

\medskip\noindent
设 $X$ 为 $\mathbf{C}$ 上的光滑连通仿射曲线。选取来自光滑连通曲线
$Y$ 的次数为 $3$ 的有限态射 $g: Y \to X$，使得存在一点
$x \in X$，其纤维 $g^{-1}(x) = \{u, y\}$ 由一个非分歧点 $u$
（次数 1）和一个分歧点 $y$（次数 2）组成。把 $X$ 替换为 $x$
的一个 Zariski 开邻域后，可以假设 $g$ 除了在 $y$ 处以外处处非分歧。
令 $U = Y \setminus \{y\}$。于是
$f=g|_U: U \to X$ 是 étale 满射，且 $f^{-1}(x) = \{u\}$。

\begin{lemma}
\label{examples-lemma-etale-not-dominated-zariski-finite-etale}
étale 覆盖 $U \to X$ 不能由开浸入和有限 étale 态射的任何有限复合
加细。
\end{lemma}

\begin{proof}
反证之，假设存在态射塔
$$
W_m \to W_{m-1} \to \dots \to W_1 \to W_0 = X
$$
其中每个映射 $W_i \to W_{i - 1}$ 或为开浸入，或为有限 étale 态射；
还存在 $W_m \to X$ 沿 $f : U \to X$ 的一个提升 $W_m \to U$，
以及一点 $w \in W_m$，它映到 $u \in U$（因而映到 $x \in X$）。
把每个 $W_k$ 替换为含 $w$ 的像 $w_k \in W_k$ 的连通分支后，
可以假设每个 $W_k$ 都光滑且连通。

\medskip\noindent
对每个 $0 \le i \le m$，考虑纤维积
$Z_i = Y \times_X W_i$。由于 $W_i \to X$ 为 étale 且 $Y$ 光滑，
曲线 $Z_i$ 是若干光滑连通曲线的无交并，而由基变换，映射
$Z_i \to W_i$ 均为次数 $3$ 的有限态射。显然，$Z_i \to W_i$
在 $w_i$ 上的纤维由两点 $z_i, v_i$ 组成，它们在 $Y$ 中分别映到
$y, u$。态射 $Z_i \to W_i$ 在 $v_i$ 处非分歧，而在 $z_i$ 处分歧
次数为 $2$。令 $n_i = \#\pi_0(Z_i)$。由于 $Z_i \to W_i$ 的纤维有
$2$ 个点，可知 $n_i \in \{1, 2\}$。此外，因为转移映射
$Z_i \to Z_{i - 1}$ 是占优的，$n_i$ 随 $i$ 只能增加。显然
$n_0 = 1$。而且必有 $n_m \geq 2$：存在 $X$-态射
$W_m \to U \subset Y$，故基变换 $Z_m \to W_m$ 有截面。因此，
存在最小指标 $i \ge 1$，使得 $n_{i - 1} = 1$ 且 $n_i = 2$。
% P12-ERR-0364
由于 $n_j$ 由 $Z_i$ 是否不可约决定，映射 $W_i \to W_{i - 1}$
不可能是开浸入，所以它必为有限 étale 态射。因为 $n_i = 2$，
光滑曲线 $Z_i$ 有两个连通分支；又因到 $W_i$ 的映射有限，其中一个
分支必同构地映到 $W_i$。因此可写
$Z_i = W_i \sqcup Z'_i$，其中 $Z'_i \to W_i$ 的次数为 $2$。
诱导映射 $W_i \subset Z_i \to Z_{i - 1}$ 是连通曲线之间的占优映射，
而这两条曲线在 $W_{i - 1}$ 上都是有限的。任何这样的映射都必为满射。
所以，可以选取一点 $t \in W_i$，它映到 $z_{i - 1}$。得到离散赋值环的
扩张
$$
% P12-ERR-0365
\mathcal{O}_{W_{i - 1}, w_{i - 1}} \to
\mathcal{O}_{Z_i, z_i} \to
\mathcal{O}_{W_i, z}
$$
其中第一个扩张是分歧的，而复合扩张是非分歧的，因为
$W_i \to W_{i - 1}$ 为 étale；这是不可能的，故得到矛盾。
\end{proof}










\section{层与专化}
\label{examples-section-sheaves}

\noindent
以下固定《拓扑》定义 \ref{topologies-definition-big-etale-site}
所构造的大 étale 位点 $\Sch_\etale$。此外，概形均视为该位点的对象。
回忆，若 $x, x'$ 是概形 $X$ 的点，而
$x \in \overline{\{x'\}}$，则称 $x$ 是 $x'$ 的一个{\it 专化}，
或写成 $x' \leadsto x$。特别地，当 $x = x'$ 时也成立。

\medskip\noindent
考虑由下述规则定义的函子 $F : \Sch_\etale \to \textit{Ab}$：
$$
F(X) = \prod\nolimits_{x \in X}
\prod\nolimits_{x' \in X, x' \leadsto x, x' \not = x}
\mathbf{Z}/2\mathbf{Z}
$$
给定概形 $X$，以 $|X|$ 表示其底层点集。把元素 $a \in F(X)$
视为集合映射
$|X| \times |X| \to \mathbf{Z}/2\mathbf{Z}$，
$(x, x') \mapsto a(x, x')$；当 $x = x'$ 或 $x$ 不是 $x'$ 的专化时，
它取零。给定概形态射 $f : X \to Y$，定义
$$
F(f) : F(Y) \longrightarrow F(X)
$$
如下：对 $b \in F(Y)$，令
$$
F(f)(b)(x, x') =
\left\{
\begin{matrix}
0 & \text{若 }x\text{ 并非专化自 }x' \\
b(f(x), f(x')) & \text{否则。}
\end{matrix}
\right.
$$
注意，这确实定义了 $F(X)$ 的一个元素。我们断言，若
$f : X \to Y$ 与 $g : Y \to Z$ 是可复合态射，则
$F(f) \circ F(g) = F(g \circ f)$。确切地说，令 $c \in F(Z)$，并设
$x' \leadsto x$ 是 $X$ 中各点的一个专化，则
$$
% P12-ERR-0366
F(g \circ f)(x, x') = c(g(f(x)), g(f(x'))) = F(g)(F(f)(c))(x, x')
$$
因为 $f(x') \leadsto f(x)$。（即使 $f(x) = f(x')$，论证也成立。）

\medskip\noindent
令 $G$ 为 $F$ 在 étale 拓扑下的层化。

\medskip\noindent
我们断言，若 $X$ 是概形，而 $x' \leadsto x$ 是满足
$x' \not = x$ 的专化，则 $G(X) \not = 0$。确切地说，取
$a \in F(X)$，使得在把 $a$ 视为函数
$|X| \times |X| \to \mathbf{Z}/2\mathbf{Z}$ 时，它在 $(x, x')$
处非零。令 $\{f_i : U_i \to X\}$ 为 $X$ 的 étale 覆盖。
可以选取一个 $i$ 和一点 $u_i \in U_i$，使 $f_i(u_i) = x$。
因为一般化可沿平坦态射提升（见《态射》引理
\ref{morphisms-lemma-generalizations-lift-flat}），可以找到专化
$u'_i \leadsto u_i$，使 $f_i(u'_i) = x'$。由上述构造可见
$F(f_i)(a) \not = 0$。因此，$a$ 确定 $G(X)$ 中的一个非零元素。

\medskip\noindent
注意，若 $X = \Spec(k)$，其中 $k$ 为域（或更一般地，是所有素理想均为
极大理想的环），则 $F(X) = 0$；而对每个 étale 态射 $U \to X$，
也有 $F(U) = 0$，因为 étale 态射的纤维中的不同点之间没有专化。
因此 $G(X) = 0$。

\medskip\noindent
假设 $X \subset X'$ 是加厚；见《态射的更多内容》定义
\ref{more-morphisms-definition-thickening}。由基变换函子
$U' \mapsto U = U' \times_{X'} X$，$X'$ 上 étale 概形的范畴与
$X$ 上 étale 概形的范畴等价；见《Étale 上同调》定理
\ref{etale-cohomology-theorem-topological-invariance}。
由于在这种情形下总有 $F(U) = F(U')$，所以也有 $G(X) = G(X')$。

\medskip\noindent
作为变体，可以考虑预层 $F_n$：它把概形 $X$ 映到所有映射
$a : |X|^{n + 1} \to \mathbf{Z}/2\mathbf{Z}$ 的集合，其中
$a(x_0, \ldots, x_n)$ 仅在
$x_n \leadsto \ldots \leadsto x_0$ 是专化序列且
$x_n \not = x_{n - 1} \not = \ldots \not = x_0$ 时才可能非零。
令 $G_n$ 为与 $F_n$ 伴随的层。与上面完全相同的方法表明，
$G_n$ 在 $\dim(X) \geq n$ 时非零，而在 $\dim(X) < n$ 时为零。

\begin{lemma}
\label{examples-lemma-sheaf-zero-on-low-dimension}
$\Sch_\etale$ 上存在一个 Abel 群层 $G$，具有如下性质：
\begin{enumerate}
\item 只要 $\dim(X) < n$，就有 $G(X) = 0$；
\item 若 $\dim(X) \geq n$，则 $G(X)$ 非零；且
\item 若 $X \subset X'$ 是加厚，则 $G(X) = G(X')$。
\end{enumerate}
\end{lemma}

\begin{proof}
见上述讨论。
\end{proof}

\begin{remark}
\label{examples-remark-specialization}
下面是若干评注：
\begin{enumerate}
\item 预层 $F$ 和 $F_n$ 是分离预层。
\item 事实上，$F$、$F_n$ 都不是层。
\item 可以证明，$G$、$G_n$
% P12-ERR-0367
实际上是 fppf 拓扑下的层。
\end{enumerate}
若以后需要，我们将证明这些结果。
\end{remark}


\section{层与可构造函数}
\label{examples-section-constructible-functions}

\noindent
以下固定《拓扑》定义 \ref{topologies-definition-big-etale-site}
所构造的大 étale 位点 $\Sch_\etale$。此外，概形均视为该位点的对象。
在本节中，概形 $X$ 的一个{\it 可构造分划}是局部有限的无交并分解
$X = \coprod_{i \in I} X_i$，其中每个 $X_i \subset X$ 都是
$X$ 的局部可构造子集。所谓局部有限，是指对任意拟紧开集
$U \subset X$，只有有限多个 $i \in I$ 满足
$X_i \cap U$ 非空。注意，若 $f : X \to Y$ 为概形态射，而
$Y = \coprod Y_j$ 是可构造分划，则
$X = \coprod f^{-1}(Y_j)$ 是 $X$ 的可构造分划。给定集合 $S$ 和概形
$X$，称映射 $f : |X| \to S$ 为{\it 可构造函数}，若
$X = \coprod_{s \in S} f^{-1}(s)$ 是 $X$ 的可构造分划。
若 $G$ 是（抽象）群，而 $a, b : |X| \to G$ 是可构造函数，则
$ab : |X| \to G, x \mapsto a(x)b(x)$ 也是可构造函数。
理由是，任意两个可构造分划都有一个同时加细二者的第三个分划。

\medskip\noindent
设 $A$ 为任意 Abel 群。对任意概形 $X$，定义
$$
F(X) =
\frac{\{a : |X| \to A \mid a \text{ 是可构造函数}\}}{\{\text{局部常值函数 }|X| \to A\}}
$$
把 $F(X)$ 的一个元素 $a$ 简单地视为一个模加局部常值函数后定义的函数。
给定概形态射 $f : X \to Y$ 和元素 $b \in F(Y)$，定义
$F(f)(b) = b \circ f$。于是 $F$ 是 $\Sch_\etale$ 上的预层。

\medskip\noindent
注意，若 $\{f_i : U_i \to X\}$ 是 fppf 覆盖，而 $a \in F(X)$
满足 $F(f_i)(a) = 0$ 于 $F(U_i)$，则对每个 $i$，$a \circ f_i$
都是局部常值函数。反过来，因为态射 $f_i$ 是开的，这意味着 $a$
是局部常值函数。因此 $a = 0$ 于 $F(X)$。由此可见，$F$ 是分离预层
（对 fppf 拓扑如此，因而更对 étale 拓扑如此）。

\medskip\noindent
令 $G$ 为 $F$ 在 étale 拓扑下的层化。由于 $F$ 分离，而且例如当
$X$ 是离散赋值环的谱时 $F(X) \not = 0$，可见 $G$ 非零。

\medskip\noindent
令 $X = \Spec(k)$，其中 $k$ 为域。于是，$X$ 的任意 étale 覆盖都可由
某个覆盖 $\{\Spec(k') \to \Spec(k)\}$ 控制，其中 $k'/k$
是有限可分域扩张。由于 $F(\Spec(k')) = 0$，可知 $G(X) = 0$。

\medskip\noindent
假设 $X \subset X'$ 是加厚；见《态射的更多内容》定义
\ref{more-morphisms-definition-thickening}。由基变换函子
$U' \mapsto U = U' \times_{X'} X$，$X'$ 上 étale 概形的范畴与
$X$ 上 étale 概形的范畴等价；见《Étale 上同调》定理
\ref{etale-cohomology-theorem-topological-invariance}。
因为在这种情形下 $F(U) = F(U')$，所以也有 $G(X) = G(X')$。

\medskip\noindent
层 $G$ 保持极限；见《空间的极限》定义
\ref{spaces-limits-definition-locally-finite-presentation}。
确切地说，设环 $R$ 写成环的有向余极限 $R = \colim_i R_i$。
令 $X = \Spec(R)$ 且 $X_i = \Spec(R_i)$，于是
$X = \lim_i X_i$。则 $G(X) = \colim_i G(X_i)$。为证明这一点，
先证明 $\Spec(R)$ 的一个可构造分划来自某个 $\Spec(R_i)$ 的一个
% P12-ERR-0368
可构造分划，由此得到 $F$ 的结论。为得到层化的结论，利用如下事实：
任意 étale 环映射 $R \to R'$ 都来自某个 $i$ 的 étale 环映射
$R_i \to R_i'$。细节从略。

\begin{lemma}
\label{examples-lemma-weird-sheaf}
$\Sch_\etale$ 上存在一个 Abel 群层 $G$，具有如下性质：
\begin{enumerate}
\item 对任意域 $k$，都有 $G(\Spec(k)) = 0$；
\item $G$ 保持极限；
\item 若 $X \subset X'$ 是加厚，则 $G(X) = G(X')$；且
\item $G$ 非零。
\end{enumerate}
\end{lemma}

\begin{proof}
见上述讨论。
\end{proof}




\section{光滑-étale 位点并非函子性的}
\label{examples-section-lisse-etale-not-functorial}

\noindent
$X$ 的{\it 光滑-étale} 位点 $X_{lisse,\etale}$ 是 $X$ 上光滑概形的
范畴，赋以（通常的）étale 覆盖；见《叠的上同调》第
\ref{stacks-cohomology-section-lisse-etale} 节。令
$f : X  \to Y$ 为概形态射。有函子
$$
u : Y_{lisse,\etale} \longrightarrow X_{lisse,\etale},\quad
V/Y \longmapsto V \times_Y X
$$
且它连续。因此，得到一对伴随函子
$$
u^s :
\Sh(X_{lisse,\etale})
\longrightarrow
\Sh(Y_{lisse,\etale}),
\quad
u_s :
\Sh(Y_{lisse,\etale})
\longrightarrow
\Sh(X_{lisse,\etale}),
$$
见《位点》第 \ref{sites-section-continuous-functors} 节。
我们断言，一般而言，$u$ 并{\bf 不}定义位点态射；见《位点》定义
\ref{sites-definition-morphism-sites}。换言之，我们断言 $u_s$
一般不是左正合的。注意，可表预层在光滑-étale 位点上是层。因此，
由《位点》引理 \ref{sites-lemma-pullback-representable-sheaf}，有
$u_sh_V = h_{V \times_Y X}$。现在考虑 $Y$ 上光滑概形 $V_1, V_2$
之间的两个态射
$$
\xymatrix{
V_1 \ar[rd] \ar@<1ex>[rr]^a \ar@<-1ex>[rr]_b & & V_2 \ar[ld] \\
& Y
}
$$
若 $u_s$ 左正合，则应有
$$
u_s \text{Equalizer}(h_a, h_b : h_{V_1} \to h_{V_2})
=
\text{Equalizer}(h_{a \times 1}, h_{b \times 1} :
h_{V_1 \times_Y X} \to h_{V_2 \times_Y X})
$$
我们将选取态射 $a, b : V_1 \to V_2$，使得不存在从 $Y$ 上光滑概形
到 $(a = b) \subset V_1$ 的态射，即等式左侧为空层；但基变换到
$X$ 后，等化子非空且在 $X$ 上光滑。一个简单例子是取
$X = \Spec(\mathbf{F}_p)$、$Y = \Spec(\mathbf{Z})$，并令
$V_1 = V_2 = \mathbf{A}^1_\mathbf{Z}$，态射为
$a(x) = x$ 和 $b(x) = x + p$。注意，$a$ 与 $b$ 的等化子是
$\mathbf{A}^1_\mathbf{Z}$ 在 $(p)$ 上的纤维。

\begin{lemma}
\label{examples-lemma-lisse-etale-not-functorial}
光滑-étale 位点并非函子性的，即使对概形态射也如此。
\end{lemma}

\begin{proof}
见上述讨论。
\end{proof}





\section{Noetherian 概形范畴上的层}
\label{examples-section-sheaves-locally-Noetherian}

\noindent
设 $S$ 为局部 Noetherian 概形。像《Artin 公理》第
\ref{artin-section-noetherian} 节一样，考虑包含函子
$$
u : (\textit{Noetherian}/S)_{fppf} \longrightarrow (\Sch/S)_{fppf}
$$
它把 $S$ 上局部 Noetherian 概形的 fppf 位点嵌入 $S$ 的一个大 fppf
位点。正如所引小节解释的，这个函子连续。因此，得到一对伴随函子
$$
u^s :
\Sh((\Sch/S)_{fppf})
\longrightarrow
\Sh((\textit{Noetherian}/S)_{fppf})
$$
和
$$
u_s :
\Sh((\textit{Noetherian}/S)_{fppf})
\longrightarrow
\Sh((\Sch/S)_{fppf})
$$
见《位点》第 \ref{sites-section-continuous-functors} 节。
然而，我们断言，一般而言，$u$ 不定义位点态射；见《位点》定义
\ref{sites-definition-morphism-sites}。换言之，函子 $u_s$
一般不是左正合的。

\medskip\noindent
令 $p$ 为素数，并置 $S = \Spec(\mathbf{F}_p)$。考虑
$(\textit{Noetherian}/S)_{fppf}$ 上如下定义的层的单射
$$
a : \mathcal{F} \longrightarrow \mathcal{G}
$$
对 $S$ 上局部 Noetherian 概形 $U$，定义
$$
\mathcal{G}(U) = \Gamma(U, \mathcal{O}_U)^* =
\Mor_S(U, \mathbf{G}_{m, S})
$$
并取
$$
\mathcal{F}(U) = \{f \in \mathcal{G}(U) \mid \text{fppf 局部地，}f
\text{ 有任意 }p\text{ 次幂根}\}
$$
Noetherian $\mathbf{F}_p$-代数 $A$ 的幂零根理想 $I \subset A$ 是
幂零的；$A$ 中 $1$ 的 $p$ 次幂根形如 $1 + a$，其中 $a \in I$；
% P12-ERR-0369
% P12-ERR-0370
因此，$1$ 的非平凡 $p$ 次幂根中没有一个具有任意 $p$ 次幂根。
由此可知，对任意局部 Noetherian 概形 $U$，$\mathcal{F}(U)$
都是无 $p$-挠的 Abel 群；一些细节从略。于是
$p : \mathcal{F} \to \mathcal{F}$ 是
$(\textit{Noetherian}/S)_{fppf}$ 上 Abel 层的单射。

\medskip\noindent
% P12-ERR-0371
为得到矛盾，假设 $u_s$ 正合。于是
$p : u_s\mathcal{F} \to u_s\mathcal{F}$ 也是单射，并可知对 $S$
上的任意 $V$，$(u_s\mathcal{F})(V)$ 都是无 $p$-挠的 Abel 群。
由于可表预层在 fppf 位点上是层，由《位点》引理
\ref{sites-lemma-pullback-representable-sheaf} 可见，
$u_s\mathcal{G}$ 由 $\mathbf{G}_{m, S}$ 表示。利用
$u_s\mathcal{F} \to u_s\mathcal{G}$ 是单射这一点，对 $S$ 上每个概形
$V$，得到一个无 $p$-挠子群
$$
(u_s\mathcal{F})(V) \subset \Gamma(V, \mathcal{O}_V)^*
$$
它具有如下性质：对 $S$ 上概形的每个态射 $V \to U$，若 $U$
局部 Noetherian，则子群
$$
\mathcal{F}(U) \subset \Gamma(U, \mathcal{O}_U)^*
$$
经限制映射
$\Gamma(U, \mathcal{O}_U)^* \to \Gamma(V, \mathcal{O}_V)^*$
映入子群 $(u_s\mathcal{F})(V)$。

\medskip\noindent
现在如下得到实际的矛盾：令
$k = \bigcup_{n \geq 0} \mathbf{F}_p(t^{1/{p^n}})$，并置
$$
B = k \otimes_{\mathbf{F}_p(t)} k
$$
以及 $V = \Spec(B)$。因为与两个余投影 $k \to B$ 对应，有两个投影态射
$V \to \Spec(k)$；又因为 $\Spec(k)$ 是 Noetherian 的，所以子群
$$
(u_s\mathcal{F})(V) \subset B^*
$$
含有 $k^* \otimes 1$ 和 $1 \otimes k^*$。这与如下事实矛盾：
$$
(t^{1/p} \otimes 1) \cdot (1 \otimes t^{-1/p}) =
t^{1/p} \otimes t^{-1/p}
$$
是 $B$ 的一个非平凡 $p$-挠单位。

\begin{lemma}
\label{examples-lemma-not-a-morphism-of-sites-noetherian-to-all}
取 $S = \Spec(\mathbf{F}_p)$ 时，包含函子
$(\textit{Noetherian}/S)_{fppf} \to (\Sch/S)_{fppf}$
不定义位点态射。
\end{lemma}

\begin{proof}
见上述讨论。
\end{proof}





\section{拟凝聚模的导出推前}
\label{examples-section-derived-push-quasi-coherent}

\noindent
设 $k$ 为特征 $p > 0$ 的域，并令 $S = \Spec(k[x])$。
令 $G = \mathbf{Z}/p\mathbf{Z}$；既可把它视为抽象群，也可视为 $S$
上的常群概形。考虑代数叠 $\mathcal{X} = [S/G]$，其中 $G$
在 $S$ 上平凡作用；见《叠的例子》注
\ref{examples-stacks-remark-X-mod-G-group} 和《可表性判据》引理
\ref{criteria-lemma-BG-algebraic}。考虑结构态射
$$
f : \mathcal{X} \longrightarrow S
$$
这个态射拟紧且拟分离。因此得到函子
$$
Rf_{\QCoh, *} :
D^{+}_\QCoh(\mathcal{O}_\mathcal{X})
\longrightarrow
D^{+}_\QCoh(\mathcal{O}_S),
$$
见《叠的导出范畴》命题
\ref{stacks-perfect-proposition-derived-direct-image-quasi-coherent}。
下面计算 $Rf_{\QCoh, *}\mathcal{O}_\mathcal{X}$。因为
$D_\QCoh(\mathcal{O}_S)$ 等价于 $k[x]$-模的导出范畴（见《概形的
导出范畴》引理 \ref{perfect-lemma-affine-compare-bounded}），
这等价于计算 $R\Gamma(\mathcal{X}, \mathcal{O}_\mathcal{X})$。
为此，可以使用覆盖 $S \to \mathcal{X}$ 及谱序列
$$
% P12-ERR-0372
H^q(S \times_\mathcal{X} \ldots \times_\mathcal{X} S, O)
\Rightarrow H^{p + q}(\mathcal{X}, \mathcal{O}_\mathcal{X})
$$
见《叠的上同调》命题
\ref{stacks-cohomology-proposition-smooth-covering-compute-cohomology}。
注意
$$
S \times_\mathcal{X} \ldots \times_\mathcal{X} S = S \times G^p
$$
是仿射的。因此，复形
$$
k[x] \to \text{Map}(G, k[x]) \to \text{Map}(G^2, k[x]) \to \ldots
$$
计算 $R\Gamma(\mathcal{X}, \mathcal{O}_\mathcal{X})$。这里，对
$\varphi \in \text{Map}(G^{p - 1}, k[x])$，其微分是把
$(g_1, \ldots, g_p)$ 映到下式的映射：
$$
\varphi(g_2, \ldots, g_p) +
\sum\nolimits_{i = 1}^{p - 1}
(-1)^i\varphi(g_1, \ldots, g_i + g_{i + 1}, \ldots, g_p)
+ (-1)^p\varphi(g_1, \ldots, g_{p - 1}).
$$
% P12-ERR-0373
这正是计算 $G$ 在 $k[x]$ 上平凡作用之群上同调的复形
（在此插入将来的参考文献）。循环群 $G$ 在 $k[x]$ 上的上同调，
在每个上同调次数 $\geq 0$ 中恰好都是一份 $k[x]$
（在此插入将来的参考文献）。由此得到
$$
Rf_*\mathcal{O}_\mathcal{X} = \bigoplus\nolimits_{n \geq 0} \mathcal{O}_S[-n]
$$
现在考虑复形
$$
E = \bigoplus\nolimits_{m \geq 0} \mathcal{O}_\mathcal{X}[m]
$$
它是 $D_\QCoh(\mathcal{O}_\mathcal{X})$ 的对象。我们先中断讨论，
给出一个一般结果。

\begin{lemma}
\label{examples-lemma-is-limit}
设 $\mathcal{X}$ 为代数叠。设 $K$ 是 $D(\mathcal{O}_\mathcal{X})$
的对象，其上同调层局部拟凝聚（《叠上的层》定义
\ref{stacks-sheaves-definition-locally-quasi-coherent}），并满足平坦基变换
性质（《叠的上同调》定义
\ref{stacks-cohomology-definition-flat-base-change}）。则在
$D(\mathcal{O}_\mathcal{X})$ 中存在特异三角
$$
K \to
\prod\nolimits_{n \geq 0} \tau_{\geq -n} K \to
\prod\nolimits_{n \geq 0} \tau_{\geq -n} K \to K[1]
$$
换言之，$K$ 是其典范截断的导出极限。
\end{lemma}

\begin{proof}
回忆，我们在 $\mathcal{X}$ 的“大 fppf 位点”
$\mathcal{X}_{fppf}$ 上工作（这是《叠上的层》和《叠的上同调》各章中
关于 $\mathcal{O}_\mathcal{X}$-模层的约定）。令 $\mathcal{B}$ 为
$\mathcal{X}_{fppf}$ 中位于某个仿射概形 $U$ 上方的对象 $x$ 的集合。
综合《叠上的层》引理
\ref{stacks-sheaves-lemma-compare-fppf-etale}、
\ref{stacks-sheaves-lemma-cohomology-restriction}，《下降》引理
\ref{descent-lemma-quasi-coherent-and-flat-base-change}，以及《概形的上同调》
引理 \ref{coherent-lemma-quasi-coherent-affine-cohomology-zero}，可见若
$\mathcal{F}$ 局部拟凝聚且 $x \in \mathcal{B}$，则
$H^p(x, \mathcal{F}) = 0$。现在，取 $d = 0$，结论即由《位点上的
上同调》引理 \ref{sites-cohomology-lemma-is-limit-dimension} 得到。
\end{proof}

\begin{lemma}
\label{examples-lemma-sum-is-product}
设 $\mathcal{X}$ 为代数叠。若 $\mathcal{F}_n$ 是 $\mathcal{X}$
上一族满足平坦基变换性质的局部拟凝聚层，则
$\oplus_n \mathcal{F}_n[n] \to \prod_n \mathcal{F}_n[n]$
是 $D(\mathcal{O}_\mathcal{X})$ 中的同构。
\end{lemma}

\begin{proof}
这是因为，由引理 \ref{examples-lemma-is-limit}，直和同构于直积。
\end{proof}

\noindent
继续上述讨论。因为拟凝聚模是局部拟凝聚的，并满足平坦基变换性质
（《叠上的层》引理 \ref{stacks-sheaves-lemma-quasi-coherent}），得到
$$
E = \prod\nolimits_{m \geq 0} \mathcal{O}_\mathcal{X}[m]
$$
由于上同调与极限交换，可见
$$
Rf_*E = \prod\nolimits_{m \geq 0}
\left(\bigoplus\nolimits_{n \geq 0} \mathcal{O}_S[m - n]\right)
$$
注意，这个复形不是 $D_\QCoh(\mathcal{O}_S)$ 的对象，因为其次数
$0$ 的上同调层是无穷多个 $\mathcal{O}_S$ 副本的直积，甚至不是局部
拟凝聚的 $\mathcal{O}_S$-模。

\begin{lemma}
\label{examples-lemma-push-not-OK}
代数叠的拟紧拟分离态射
$f : \mathcal{X} \to \mathcal{Y}$ 未必诱导函子
$Rf_* : D_\QCoh(\mathcal{O}_\mathcal{X}) \to
D_\QCoh(\mathcal{O}_\mathcal{Y})$。
\end{lemma}

\begin{proof}
见上述讨论。
\end{proof}



\section{大 Abel 范畴}
\label{examples-section-big}

\noindent
本节的目的是给出一个“大”Abel 范畴 $\mathcal{A}$ 及对象 $M, N$ 的例子，
使得扩张的同构类之汇集
% P12-ERR-0374
$\Ext_\mathcal{A}(M, N)$ 不是集合。该例由 Freyd 给出；见
\cite[第 131 页，练习 A]{Freyd}。

\medskip\noindent
如下定义 $\mathcal{A}$。$\mathcal{A}$ 的对象是三元组
$(M, \alpha, f)$，其中 $M$ 为 Abel 群，$\alpha$ 为序数，且
$f : \alpha \to \text{End}(M)$ 为映射。态射
$(M, \alpha, f) \to (M', \alpha', f')$ 由 Abel 群同态
$\varphi : M \to M'$ 给出，并要求对{\it 任意}序数 $\beta$ 都有
$$
\varphi \circ f(\beta) = f'(\beta) \circ \varphi
$$
这里的约定是：若 $\beta$ 不属于 $\alpha$，则置 $f(\beta) = 0$；类似地，
若 $\beta$ 不是 $\alpha'$ 的元素，则把 $f'(\beta)$ 置为零。
我们略去如此定义的范畴为 Abel 范畴的验证。

\medskip\noindent
考虑对象 $Z = (\mathbf{Z}, \emptyset, f)$，即所有算子均为零。
关键观察是，在 $\mathcal{A}$ 中计算的
$\Ext^1_\mathcal{A}(Z, Z)$ 是真类而非集合。确切地说，对每个序数
$\alpha$，可以找到一个 $Z$ 被 $Z$ 扩张得到的对象
$(M, \alpha + 1, f)$，其底层群为
$M = \mathbf{Z} \oplus \mathbf{Z}$，而 $f$ 的值除下式以外总为零：
$$
f(\alpha) =
\left(
\begin{matrix}
0 & 1 \\
0 & 0
\end{matrix}
\right).
$$
显然，这产生一个由扩张的同构类组成的真类。特别地，
$\mathcal{A}$ 的导出范畴的态射汇集是一些真类；见《导出范畴》引理
\ref{derived-lemma-ext-1}。这意味着，在定义三角范畴的 Verdier 商时，
必须格外谨慎。

\begin{lemma}
\label{examples-lemma-big-abelian-category}
存在一个“大”Abel 范畴 $\mathcal{A}$，其 $\Ext$-群为真类。
\end{lemma}

\begin{proof}
见上述讨论。
\end{proof}





\section{弱伴随点与概形论稠密性}
\label{examples-section-weak-ass-dense}

\noindent
设 $k$ 为域。令 $R = k[z, x_i, y_i]/(z^2, zx_iy_i)$，其中 $i$
遍历 $\mathbf{N}$ 的元素。注意，$R = R_0 \oplus M_0$，其中
$R_0 = k[x_i, y_i]$ 是子环，而 $M_0$ 是平方为零的理想，并且作为
$R_0$-模有 $M_0 \cong R_0/(x_iy_i)$。素理想
$\mathfrak p = (z, x_i)$ 弱伴随于作为 $R$-模的 $R$
（《代数》定义 \ref{algebra-definition-weakly-associated}）。
事实上，$R_\mathfrak p$ 中的元素 $z$ 非零，但被
$\mathfrak pR_\mathfrak p$ 零化。另一方面，考虑开子概形
$$
U = \bigcup D(x_i) \subset \Spec(R) = S
$$
我们断言 $U \subset S$ 是概形论稠密的
（《态射》定义 \ref{morphisms-definition-scheme-theoretically-dense}）。
为证明这一点，只需说明 $\mathcal{O}_S \to j_*\mathcal{O}_U$ 是单射，
其中 $j : U \to S$ 是包含态射；见《态射》引理
\ref{morphisms-lemma-characterize-scheme-theoretically-dense}。
译成代数语言，我们必须说明，对所有 $g \in R$，映射
$$
R_g \longrightarrow \prod R_{x_ig}
$$
都是单射。写成 $g = g_0 + m_0$，其中 $g_0 \in R_0$ 且
$m_0 \in M_0$。于是 $R_g = R_{g_0}$（略去细节）。故可假设
$g \in R_0$。还可假设 $g$ 非零。此时
$R_g = (R_0)_g \oplus (M_0)_g$。因为 $R_0$ 是整环，映射
$(R_0)_g \to \prod (R_0)_{x_ig}$ 是单射。若
$g \in (x_iy_i)$，则 $(M_0)_g = 0$，无须证明任何事情。若
$g \not \in (x_iy_i)$，则因 $(x_iy_i)$ 是 $R_0$ 的根理想，
只须说明 $M_0 \to \prod (M_0)_{x_ig}$ 是单射。映射
$R_0 \to M_0 \to (M_0)_{x_n}$ 的核为 $(x_iy_i, y_n)$。由于
$(x_iy_i, y_n)$ 是根理想，若 $g \not \in (x_iy_i, y_n)$，则映射
$R_0 \to M_0 \to (M_0)_{x_ng}$ 的核为 $(x_iy_i, y_n)$。因为对所有
$n \gg 0$ 都有 $g \not \in (x_iy_i, y_n)$，可知该核包含于
$\bigcap_{n \gg 0} (x_iy_i, y_n) = (x_iy_i)$，这正是所需结论。

\medskip\noindent
第二个例由 Ofer Gabber 给出。设 $k$ 为域，并分别令 $R$ 和 $R'$ 为
所有函数 $\mathbf{N} \to k$ 所成的环和所有最终为常值的函数
$\mathbf{N} \to k$ 所成的环。于是 $\Spec(R)$ 和 $\Spec(R')$ 分别是
Stone-{\v C}ech 紧化\footnote{每个 $f \in R$ 都可写成 $ue$，其中
$u$ 是单位而 $e$ 是幂等元。于是《代数》引理
\ref{algebra-lemma-ring-with-only-minimal-primes} 表明 $\Spec(R)$
是 Hausdorff 的。另一方面，可把带离散拓扑的 $\mathbf{N}$ 看成稠密开子集。
给定到 Hausdorff 拟紧拓扑空间 $X$ 的集合映射 $\mathbf{N} \to X$，得到环映射
$\mathcal{C}^0(X; k) \to R$，其中 $\mathcal{C}^0(X; k)$ 是局部常值映射
$X \to k$ 所成的 $k$-代数。这给出
$\Spec(R) \to \Spec(\mathcal{C}^0(X; k)) = X$，从而证明泛性质。}
$\beta\mathbf{N}$ 和单点紧化\footnote{此处可论证 $R'$ 中确实只有一个额外的
极大理想。}$\mathbf{N}^* = \mathbf{N} \cup \{\infty\}$。
所有点都是弱伴随点，因为 $R$ 和 $R'$ 中的所有素理想都是极小素理想。

\begin{lemma}
\label{examples-lemma-example-schematically-dense-missing-weakly-associated-point}
存在约化概形 $X$ 和概形论稠密开子概形 $U \subset X$，使得某个弱伴随点
$x \in X$ 不属于 $U$。
\end{lemma}

\begin{proof}
在第一个例中，由构造即有 $\mathfrak p \not \in U$。在 Gabber 的例中，
概形 $\Spec(R)$ 或 $\Spec(R')$ 是约化的。
\end{proof}



\section{迹不具可加性的例}
\label{examples-section-non-additive}

\noindent
设 $k$ 为域，并令 $R = k[\epsilon]$ 为 $k$ 上的对偶数环。换言之，
$R = k[x]/(x^2)$，且 $\epsilon$ 是 $x$ 在 $R$ 中的同余类。考虑复形的短正合列
$$
\xymatrix{
0 \ar[d] \ar[r] &
R \ar[d]^\epsilon \ar[r]_1 &
R \ar[d] \\
R \ar[r]^1 & R \ar[r] & 0
}
$$
这里各列为复形，第一行置于次数 $0$，第二行置于次数 $1$。分别以
$A^\bullet$、$B^\bullet$ 和 $C^\bullet$ 表示第一个复形（即左列）、
第二个复形和第三个复形。我们断言，图表
\begin{equation}
\label{examples-equation-commutes-up-to-homotopy}
\vcenter{
\xymatrix{
A^\bullet \ar[d]_{1 + \epsilon} \ar[r] &
B^\bullet \ar[r] \ar[d]_1 &
C^\bullet \ar[d]_1 \\
A^\bullet \ar[r] & B^\bullet \ar[r] & C^\bullet
}
}
\end{equation}
在 $K(R)$ 中交换，即它是一个在同伦意义下交换的复形图表。事实上，
右边的方块交换，而左边方块与交换相差同伦 $1 : A^1 \to B^0$。
另一方面，
$$
\text{Tr}_{A^\bullet}(1 + \epsilon) + \text{Tr}_{C^\bullet}(1)
\not = \text{Tr}_{B^\bullet}(1).
$$

\begin{lemma}
\label{examples-lemma-nonadditivity-of-trace}
存在环 $R$、同伦范畴 $K(R)$ 中的特异三角
$(K, L, M, \alpha, \beta, \gamma)$，以及该特异三角的自同态
$(a, b, c)$，使得 $K$、$L$、$M$ 都是完美复形，并且
$\text{Tr}_K(a) + \text{Tr}_M(c) \not = \text{Tr}_L(b)$。
\end{lemma}

\begin{proof}
考虑上述例子。映射 $\gamma : C^\bullet \to A^\bullet[1]$ 在次数 $0$
由乘以 $\epsilon$ 给出；见《导出范畴》定义
\ref{derived-definition-distinguished-triangle}。因此同样有
$$
\xymatrix{
C^\bullet \ar[d] \ar[r]_\gamma & A^\bullet[1] \ar[d] \\
C^\bullet \ar[r]^\gamma & A^\bullet[1]
}
$$
在 $K(R)$ 中交换，因为 $\epsilon(1 + \epsilon) = \epsilon$。
故我们确实得到特异三角的一个态射。
\end{proof}




\section{射影概形结构层的非有限上同调}
\label{examples-section-nonfinite-h0}

\noindent
设 $R$ 为任意非 Noether 环，$I$ 为 $R$ 的非有限生成理想。令
$S = R[x, y]/xyI$，并把它看成分次环，其中 $x$ 和 $y$ 的次数均为 $1$。
则 $X = \text{Proj}(S)$ 是 $R$ 上的射影概形。我们断言
$H^0(X, \mathcal{O}_X)$ 不是有限生成 $R$-模。

\medskip\noindent
按双次数分解环 $S_x$、$S_y$ 和 $S_{xy}$，可写成
$$
S_x = \bigoplus_{\substack{(a, b) \in \mathbf{Z}^2 \\ b \geq 0}} R'_{a, b},
\quad\quad
S_y = \bigoplus_{\substack{(a, b) \in \mathbf{Z}^2 \\ a \geq 0}} R''_{a, b},
\quad\quad
S_{xy} = \bigoplus_{\substack{(a, b) \in \mathbf{Z}^2}} R/I,
$$
其中 $(a, b)$-直和项位于次数 $a + b$，并且
$$
R'_{a, b} =
\left\{
\begin{matrix}
R & b = 0, \\
R/I & b \ne 0,
\end{matrix}
\right.
\quad\quad
R''_{a, b} =
\left\{
\begin{matrix}
R & a = 0, \\
R/I & a \ne 0.
\end{matrix}
\right.
$$
$R$-模 $H^0(X, \mathcal{O}_X)$ 可计算为映射
$$
(S_x)_0 \times (S_y)_0 \to (S_{xy})_0
$$
的核，该映射由 $(F, G) \mapsto F - G$ 定义。这里记号 $A_0$ 表示
$\mathbf{Z}$-分次环 $A$ 的次数 $0$ 分量。由上面对
$S_x, S_y, S_{xy}$ 的显式描述，该核同构于由满足
$F - G \in I$ 的偶对 $(F, G)$ 组成的子模 $M \subset R^2$。由于
$R$-模映射 $M \to I$，$(F, G) \mapsto F - G$ 是满射，
$M$ 不是有限生成 $R$-模。

\medskip\noindent
满射 $R[x, y] \to S$ 对应于闭浸入
$i : X \to \mathbf{P}^1_R$。于是
$\mathcal{F} = i_*\mathcal{O}_X$ 是有限型拟凝聚
$\mathcal{O}_{\mathbf{P}^1_R}$-模，其 $H^0$ 等于上面算得的非有限生成模
$M$。

\begin{lemma}
\label{examples-lemma-nonfinite-h0}
非有限的 $H^0$。
\begin{enumerate}
\item 存在环 $R$ 和 $R$ 上的射影概形 $X$，使得
$H^0(X, \mathcal{O}_X)$ 不是有限 $R$-模。
\item 存在环 $R$ 和 $\mathbf{P}^1_R$ 上的有限型拟凝聚模
$\mathcal{F}$，使得 $H^0(\mathbf{P}^1_R, \mathcal{F})$
不是有限 $R$-模。
\end{enumerate}
\end{lemma}

\begin{proof}
见上述讨论。
\end{proof}









\section{射影性不在基底上局部}
\label{examples-section-non-descending-property-projective}

\noindent
在关于下降的章中，我们已经看到，许多态射性质在基底上局部，甚至在
fpqc 拓扑中也是如此。见《下降》的章节
\ref{descent-section-descending-properties-morphisms}、
\ref{descent-section-descending-properties-morphisms-fpqc} 和
\ref{descent-section-descending-properties-morphisms-fppf}。
对态射的射影性而言，情形并非如此。

\begin{lemma}
\label{examples-lemma-non-descending-property-projective}
性质
\begin{enumerate}
\item[] $\mathcal{P}(f) =$``$f$ 是射影的''，以及
\item[] $\mathcal{P}(f) =$``$f$ 是拟射影的''
\end{enumerate}
在基底上都不是 Zariski 局部的。因此，它们更不是在基底上 fpqc 局部的。
\end{lemma}

\begin{proof}
依照 Hironaka \cite[例 B.3.4.1]{H}，我们构造光滑复三维簇之间的真态射
$f:V_Y\to Y$；它在 Zariski 局部是射影的，但并非射影的。由于 $f$
是真而非射影的，它也不是拟射影的。

\medskip\noindent
令 $Y$ 为复数域 ${\mathbf C}$ 上的射影三维空间。取 $Y$ 中的光滑二次曲线
$C$ 和 $D$，使得闭子概形 $C\cap D$ 约化，且由两个复点 $P$ 和 $Q$
组成。（例如，取
$C=\{ [x,y,z,w]: xy=z^2, w=0\}$，$D=\{ [x,y,z,w]:
xy=w^2, z=0\}$，$P=[1,0,0,0]$，且
$Q=[0,1,0,0]$。）
在 $Y-Q$ 上，先沿曲线 $C$ 吹起，再沿曲线 $D$ 的严格变换吹起
（《除子》定义 \ref{divisors-definition-strict-transform}）。在 $Y-P$ 上，
先沿曲线 $D$ 吹起，再沿曲线 $C$ 的严格变换吹起。在 $Y-P-Q$ 上，
如此构造的两个簇典范同构，故可把它们沿 $Y-P-Q$ 粘合。所得
$V_Y$ 是 ${\mathbf C}$ 上的光滑真三维簇。态射 $f:V_Y\to Y$ 是真的，
并且在 Zariski 局部是射影的（因为它在 $Y-P$ 和 $Y-Q$ 上都是吹起）；
这由《除子》引理 \ref{divisors-lemma-blowing-up-projective} 得到。
我们将证明 $V_Y$ 在 ${\mathbf C}$ 上不是射影的。这将推出 $f$ 不是射影的。

\medskip\noindent
为此，令 $L$ 为 $C-P-Q$ 的一个复点在 $V_Y$ 中的逆像，令 $M$ 为
$D-P-Q$ 的一个复点的逆像。则 $L$ 和 $M$ 都同构于射影直线
${\mathbf P}^1_{{\mathbf C}}$。接着，在 $V_Y$ 中令 $E$ 为
$C\cup D\subset Y$ 在 $V_Y$ 中的逆像；于是
$E\rightarrow C\cup D$ 是真态射，在
$(C\cup D)-\{P,Q\}$ 上的纤维同构于 ${\mathbf P}^1$。点 $P$ 在 $E$
中的逆像是两条直线 $L_0$ 和 $M_0$ 的并，并且在 $E$ 上有循环的有理等价
$L\sim L_0+M_0$ 和 $M\sim M_0$（这里利用 $C$ 和 $D$ 同构于
${\mathbf P}^1$）。注意，这种不对称性来自我们吹起两条曲线的次序。
在 $Q$ 附近，发生相反的情形。因此，$Q$ 的逆像是两条直线
$L_0'$ 和 $M_0'$ 的并，并且在 $E$ 上有有理等价
$L\sim L_0'$ 和 $M\sim L_0'+M_0'$。结合这些等价关系，得到在 $E$
上 $L_0+M_0'\sim 0$，因而在 $V_Y$ 上也如此。若 $V_Y$ 在
${\mathbf C}$ 上射影，则它具有丰富线丛 $H$，而该线丛在 $V_Y$ 的所有曲线
上的次数都 $> 0$。特别地，$H$ 在 $L_0+M_0'$ 上的次数为正；但在线性丛的
次数在域上真概形的模有理等价的一维循环上良定义，这给出矛盾
（《Chow 同调》引理
\ref{chow-lemma-proper-pushforward-rational-equivalence}
和引理 \ref{chow-lemma-factors}）。
所以 $V_Y$ 在 ${\mathbf C}$ 上不是射影的。
\end{proof}

\noindent
换一种术语说，Hironaka 的三维簇 $V_Y$ 是 $Y$ 沿约化子概形
$C\cup D$ 的吹起 $Y'$ 的一个小解消；这里 $Y'$ 有两个结点奇点。
若把 $Z$ 定义为先沿 $C$ 吹起 $Y$，再沿 $D$ 的严格变换吹起所得的概形，
则 $Z$ 是光滑射影三维簇，而非射影三维簇 $V_Y$ 与 $Z$ 的差别是在
$Y-P$ 上做一次 ``flop''。


\section{射影概形的非有效下降数据}
\label{examples-section-non-effective-descent-projective}

\noindent
在关于下降的章中，我们已看到，相对于 fpqc 态射的概形下降数据对于若干类
态射是有效的。特别地，仿射态射以及更一般的拟仿射态射满足 fpqc 覆盖的下降
（《下降》引理 \ref{descent-lemma-quasi-affine}）。对射影态射而言并非如此。

\begin{lemma}
\label{examples-lemma-non-effective-descent-projective}
存在概形的 étale 覆盖 $X\to S$，以及相对于 $X\to S$ 的下降数据
$(V/X,\varphi)$，使得 $V\to X$ 是射影的，但该下降数据在概形范畴中
不是有效的。
\end{lemma}

\begin{proof}
我们仿照 Hironaka 的一个例子：一个不是概形的三维光滑分离复代数空间
\cite[例 B.3.4.2]{H}。

\medskip\noindent
考虑群 $G = \mathbf{Z}/2 = \{1, g\}$ 在复数域上的射影三维空间
$\mathbf{P}^3$ 上的作用：
$$
g[x,y,z,w] = [y,x,w,z].
$$
除 ${\mathbf P}^3$ 中两条不交直线
$L_1=\{ [x,x,z,z]\}$ 和 $L_2=\{ [x,-x,z,-z]\}$ 外，该作用是自由的。
令 $Y={\mathbf P}^3-(L_1\cup L_2)$。存在 ${\mathbf C}$ 上的光滑拟射影概形
$S=Y/G$，使得 $Y\to S$ 是 $G$-挠子（《群胚》定义
\ref{groupoids-definition-principal-homogeneous-space}）。
具体地，可以把 $S$ 定义为开子集 $Y$ 在 ${\mathbf P}^3$ 的态射
\begin{align*}
{\mathbf P}^3 & \to \text{Proj } {\mathbf C}[x,y,z,w]^G\\
   & = \text{Proj } {\mathbf C}[u_0,u_1,v_0,v_1,v_2]/(v_0v_1=v_2^2),
\end{align*}
下的像，其中 $u_0=x+y$、$u_1=z+w$、$v_0=(x-y)^2$、
$v_1=(z-w)^2$ 且 $v_2=(x-y)(z-w)$；环的分次规定为
$u_0,u_1$ 的次数为 1，而 $v_0,v_1,v_2$ 的次数为 2。

\medskip\noindent
令 $C=\{ [x,y,z,w]: xy=z^2, w=0\}$，并令
$D=\{ [x,y,z,w]: xy=w^2, z=0\}$。它们是 ${\mathbf P}^3$ 中的光滑二次
曲线，包含于 $G$-不变开子集 $Y$，且 $g(C)=D$。此外，$C\cap D$
由两点 $P:=[1,0,0,0]$ 和 $Q:=[0,1,0,0]$ 组成，这两个点在 $G$
的作用下彼此交换。

\medskip\noindent
令 $V_Y\to Y$ 为如下概形：在 $Y-P$ 上，先沿 $D$ 吹起，再沿 $C$
的严格变换吹起；在 $Y-Q$ 上，先沿 $C$ 吹起，再沿 $D$ 的严格变换吹起。
（这与引理 \ref{examples-lemma-non-descending-property-projective} 的证明中的构造相同，
不同之处是，这里的 $Y$ 表示 ${\mathbf P}^3$ 的开子集，而不是整个
${\mathbf P}^3$。）于是 $G$ 在 $Y$ 上的作用提升为 $G$ 在 $V_Y$ 上的作用，
该作用交换 $Y-P$ 和 $Y-Q$ 的逆像。$G$ 在 $V_Y$ 上的这个作用给出
$V_Y$ 上相对于 $G$-挠子 $Y\to S$ 的下降数据
$(V_Y/Y,\varphi_Y)$。如引理
\ref{examples-lemma-non-descending-property-projective} 的证明所示，态射
$V_Y\to Y$ 是真的，但不是射影的。

\medskip\noindent
令 $X$ 为开子集 $Y-P$ 和 $Y-Q$ 的不交并；于是有满 étale 态射
$X\to Y\to S$。令 $V$ 为 $V_Y\to Y$ 到 $X$ 的拉回；则态射
$V\to X$ 是射影的，因为 $V_Y\to Y$ 在开子集 $Y-P$ 和 $Y-Q$
的每一个上都是吹起。此外，下降数据 $(V_Y/Y,\varphi_Y)$ 拉回为相对于
étale 覆盖 $X\to S$ 的下降数据 $(V/X,\varphi)$。

\medskip\noindent
假设该下降数据在概形范畴中有效。也就是说，存在概形 $U\to S$，其拉回连同
下降数据就是态射 $V\to X$。于是 $U$ 将是 $V_Y$ 对其 $G$-作用的商。
$$
\xymatrix{
V \ar[r]\ar[d]&  X\ar[d] \\
V_Y \ar[r]\ar[d]& Y\ar[d] \\
U \ar[r]& S
}
$$

\medskip\noindent
令 $E$ 为 $C\cup D\subset Y$ 在 $V_Y$ 中的逆像；于是
$E\rightarrow C\cup D$ 是真态射，并且在 $(C\cup D)-\{P,Q\}$
上的纤维同构于 ${\mathbf P}^1$。$P$ 在 $E$ 中的逆像是两条直线
$L_0$ 和 $M_0$ 的并。因此，$Q=g(P)$ 在 $E$ 中的逆像是两条直线
$L_0'=g(M_0)$ 和 $M_0'=g(L_0)$ 的并。如引理
\ref{examples-lemma-non-descending-property-projective} 的证明所示，在 $E$ 上有
有理等价 $L_0+M_0'=L_0+g(L_0)\sim 0$。

\medskip\noindent
由闭子概形的下降，存在曲线 $L_1\subset U$（同构于 ${\mathbf P}^1$），
其在 $V_Y$ 中的逆像是 $L_0\cup g(L_0)$。（使用《下降》引理
\ref{descent-lemma-affine}，并注意闭浸入是仿射态射。）
令 $R$ 为 $L_1$ 的一个复点。由于假设 $U$ 是概形，可在局部环
% P12-ERR-0375
$O_{U,R}$ 中选择函数 $f$，它在 $R$ 处消失，但不在整条曲线 $L_1$
上消失。令 $D_{\text{loc}}$ 为 $\Spec O_{U,R}$ 中闭子集
$\{f = 0\}$ 的一个不可约分支；则 $D_{\text{loc}}$ 的余维为 1。
$D_{\text{loc}}$ 在 $U$ 中的闭包是 $U$ 中的不可约除子 $D_U$，它包含点
$R$，但不包含整条曲线 $L_1$。$D_U$ 在 $V_Y$ 中的逆像是有效除子 $D$，
它与 $L_0\cup g(L_0)$ 相交，但不包含曲线 $L_0$ 或 $g(L_0)$ 中的任一条。

\medskip\noindent
由于复三维簇 $V_Y$ 光滑，
% P12-ERR-0376
$O(D)$ 是 $V_Y$ 上的线丛。这里我们用到了正则局部环是因子环，换言之是
UFD；见《代数进阶》引理 \ref{more-algebra-lemma-regular-local-UFD}。
$O(D)$ 到真曲面 $E\subset V_Y$ 的限制是线丛；根据我们关于 $D$ 的信息，
它在一维循环 $L_0+g(L_0)$ 上的次数为正。由于在 $E$ 上
$L_0+g(L_0)\sim 0$，这与线丛的次数在域上真概形的模有理等价的一维循环上
良定义相矛盾（《Chow 同调》引理
\ref{chow-lemma-proper-pushforward-rational-equivalence}
和引理 \ref{chow-lemma-factors}）。因此，下降数据 $(V/X,\varphi)$
事实上不是有效的；也就是说，$U$ 不作为概形存在。
\end{proof}

\noindent
在这个例中，该下降数据在代数空间范畴中{\it 是}有效的。更确切地说，
$U$ 作为 ${\mathbf C}$ 上的三维光滑分离代数空间存在；例如，这由
《代数空间》引理 \ref{spaces-lemma-quotient} 得到。Hironaka 的三维空间
$U$ 是光滑拟射影三维簇 $S$ 沿不可约结点曲线 $(C\cup D)/G$ 的吹起
$S'$ 的一个小解消；三维簇 $S'$ 有一个结点奇点。$S'$ 的另一个小解消
（与 $U$ 相差一次 ``flop''）仍是一个不是概形的代数空间。






\section{全空间不是概形的曲线族}
\label{examples-section-family-of-curves}

\noindent
在《Quot》章节 \ref{quot-section-curves} 中，我们把概形 $S$ 上的曲线族
定义为从代数空间 $X$ 到 $S$ 的相对维数 $\leq 1$ 的真、平坦、有限表现态射。
若 $S$ 是完备 Noether 局部环的谱，则 $X$ 是概形；见《空间态射进阶》引理
\ref{spaces-more-morphisms-lemma-projective-over-complete}。
本节说明，这在一般情形下并不成立。

\medskip\noindent
设 $k$ 为域。我们从真平坦态射
$$
Y \longrightarrow \mathbf{A}^1_k
$$
以及位于 $0 \in \mathbf{A}^1_k(k)$ 上方的点 $y \in Y(k)$ 出发，
并要求它们满足如下性质：
\begin{enumerate}
\item 纤维 $Y_0$ 是 $k$ 上的光滑几何不可约曲线；
\item 对任意支配 $\mathbf{A}^1_k$ 的真闭子概形 $T \subset Y$，交
$T \cap Y_0$ 都包含至少一个不同于 $y$ 的点。
\end{enumerate}
给定这样的曲面，我们如下构造所需例子：
$$
\xymatrix{
Y \ar[rd] & Z \ar[d] \ar[l] \ar[r] & X \ar[ld] \\
& \mathbf{A}^1_k
}
$$
这里 $Z \to Y$ 是 $Y$ 在 $y$ 处的吹起。令 $E \subset Z$ 为例外除子，
并令 $C \subset Z$ 为 $Y_0$ 的严格变换。概形论意义下有
$Z_0 = E \cup C$（为看出这一点，可用 $Y$ 在 $y$ 处光滑以及
$Y \to \mathbf{A}^1_k$ 在 $y$ 处光滑）。由 Artin 的结果
（\cite{ArtinII}；使用《半稳定约化》引理
\ref{models-lemma-properties-form} 说明 $C$ 的法丛是负的），
可在 $Z$ 中吹下曲线 $C$，得到图中的代数空间 $X$。令
$x \in X(k)$ 为 $C$ 的像。

\medskip\noindent
我们断言 $X$ 不是概形。事实上，若它是概形，则 $x$ 有仿射开邻域
$U \subset X$。置 $T = X \setminus U$。则 $T$ 支配
$\mathbf{A}^1_k$（因为 $X \to \mathbf{A}^1_k$ 的纤维是真且维数为 $1$ 的，而
$U \to \mathbf{A}^1_k$ 的纤维是仿射的，因而二者不同）。令
$T' \subset Z$ 为同构地映到 $T$ 的闭子概形（因为 $x \not \in T$）。
于是
% P12-ERR-0377
$T'$ 在 $X$ 中的像与上面的条件 (2) 矛盾（因为 $T' \cap Z_0$
包含于吹起 $Z \to Y$ 的例外除子 $E$）。

\medskip\noindent
为结束讨论，还需构造这样的 $Y$。我们假设 $k$ 的特征不是 $3$。
写成 $\mathbf{A}^1_k = \Spec(k[t])$，并在 $\mathbf{P}^2_{k[t]}$ 中取
$$
Y \quad : \quad T_0^3 + T_1^3 + T_2^3 - tT_0T_1T_2 = 0
$$
它在 $t = 0$ 处的纤维是光滑射影亏格为 $1$ 的曲线。在仿射块
% P12-ERR-0378
$V_+(T_0)$ 上，得到仿射方程
$$
1 + x^3 + y^3 - txy = 0
$$
它定义 $k$ 上的光滑曲面。由对称性，其他仿射块上也成立，故 $Y$ 是光滑
曲面。最后，从该仿射方程还可看出分式域是 $k(x, y)$，因而 $Y$ 是有理曲面。
有理曲面的 Picard 群有限生成
% P12-ERR-0379
（在此插入将来的参考文献）。因此，为选出具有性质 (2) 的
$y \in Y_0(k)$，只需选取 $y$ 使得
\begin{equation}
\label{examples-equation-three}
\mathcal{O}_{Y_0}(ny) \not \in \Im(\Pic(Y) \to \Pic(Y_0))
\text{ 对所有 }n > 0
\end{equation}
事实上，与 (2) 矛盾的 $T$ 的所有 $1$-维不可约分支之和将给出一个有效
Cartier 除子，
% P12-ERR-0380
其交 $Y_0$ 为某个 $n \geq 1$ 的除子 $ny$，由此可知
$\mathcal{O}_{Y_0}(ny)$ 位于限制映射的像中。注意，因为 $Y_0$ 的亏格
$\geq 1$，映射
$$
Y_0(k) \to \Pic(Y_0),\quad y \mapsto \mathcal{O}_{Y_0}(y)
$$
是单射。现在，若 $k$ 是不可数代数闭域，则由 $\Pic(Y)$ 的可数性和刚才的
观察，可以找到满足 (\ref{examples-equation-three})、从而满足 (2) 的点
$y \in Y_0(k)$。

\begin{lemma}
\label{examples-lemma-family-of-curves-not-scheme}
存在域 $k$ 和曲线族 $X \to \mathbf{A}^1_k$，使得 $X$ 不是概形。
\end{lemma}

\begin{proof}
见上述讨论。
\end{proof}






\section{导出基变换}
\label{examples-section-derived-base-change}

\noindent
设 $R \to R'$ 为环映射。在《代数进阶》章节
\ref{more-algebra-section-derived-base-change} 中，我们构造了导出基变换函子
$- \otimes_R^\mathbf{L} R' : D(R) \to D(R')$。接着，设
$R \to A$ 为第二个环映射。图示为
$$
\xymatrix{
	A \ar[r] & A \otimes_R R' \ar@{=}[r] & A' \\
R \ar[u] \ar[r] & R' \ar[u] \ar[ur]
}
$$
给定 $A$-模 $M$，张量积 $M \otimes_R R'$ 是
$A \otimes_R R'$-模，即 $A'$-模。对环映射 $A \to A'$，有导出函子
$$
- \otimes_A^\mathbf{L} A' : D(A) \longrightarrow D(A')
$$
但一般而言，该函子与 $- \otimes_R^\mathbf{L} R'$ 不一致。更确切地说，
对 $K \in D(A)$，在《代数进阶》方程
(\ref{more-algebra-equation-comparison-map}) 中构造的 $D(R')$ 中的典范映射
$$
K \otimes_R^{\mathbf{L}} R' \longrightarrow
K \otimes_A^{\mathbf{L}} A'
$$
一般不是同构。因此可以问，是否存在一个“导出基变换函子”
$T : D(A) \to D(A')$，即一个使得 $T(K)$ 在 $D(R')$ 中映到
$K \otimes_R^\mathbf{L} R'$ 的函子。本节将说明，一般而言它不存在。

\medskip\noindent
设 $k$ 为域。置 $R = k[x, y]$、$R' = R/(xy)$ 且 $A = R/(x^2)$。
$D(R')$ 中的对象 $A \otimes_R^\mathbf{L} R'$ 由
$$
x^2 : R' \longrightarrow R'
$$
表示，并且 $H^0(A \otimes_R^\mathbf{L} R') = A \otimes_R R'$。
我们断言，不存在 $D(A \otimes_R R')$ 的对象 $E$，使它在 $D(R')$
中映到 $A \otimes_R^\mathbf{L} R'$。事实上，对这样的 $E$，模
$H^0(E)$ 将是自由的，因而 $E$ 将分解成
$H^0(E)[0] \oplus H^{-1}(E)[1]$。但容易看出，在 $D(R')$ 中，
$A \otimes_R^\mathbf{L} R'$ 不同构于其各上同调群的直和。

\begin{lemma}
\label{examples-lemma-no-derived-base-change}
设 $R \to R'$ 和 $R \to A$ 为环映射。一般而言，不存在三角范畴的函子
$T : D(A) \to D(A \otimes_R R')$，使得 $A$-模 $M$ 给出
$D(A \otimes_R R')$ 中的对象 $T(M)$，且该对象在映射
$D(A \otimes_R R') \to D(R')$ 下映到
$M \otimes_R^\mathbf{L} R'$。
\end{lemma}

\begin{proof}
见上述讨论。
\end{proof}




\section{一个有趣的紧对象}
\label{examples-section-interesting-compact}


\noindent
设 $R$ 为环，并设 $(A, \text{d})$ 为微分分次 $R$-代数。若 $A = R$，
则我们知道 $D(A, \text{d}) = D(R)$ 的每个紧对象都由有限投射模的有限复形
表示。换言之，紧对象是完美的；见《代数进阶》命题
\ref{more-algebra-proposition-perfect-is-compact}。用微分分次模的语言，
相应问题是：“$D(A, \text{d})$ 的每个紧对象是否都可由一个有限且分次投射的
微分分次 $A$-模 $P$ 表示？”

\medskip\noindent
对一般的微分分次代数，这并不成立。事实上，设 $k$ 为特征 $2$ 的域
（这样便不必顾虑符号）。令 $A = k[x, y]/(y^2)$，其中
\begin{enumerate}
\item $x$ 的次数为 $0$；
\item $y$ 的次数为 $-1$；
\item $\text{d}(x) = 0$；并且
\item $\text{d}(y) = x^2 + x$。
\end{enumerate}
于是 $x : A \to A$ 是 $K(A, \text{d})$ 中的投影算子。因此
$$
A = \Ker(x) \oplus \Im(1 - x)
$$
在 $K(A, \text{d})$ 中成立；见《微分分次代数》引理
\ref{dga-lemma-homotopy-direct-sums} 和《导出范畴》引理
\ref{derived-lemma-projectors-have-images-triangulated}。
显然，$A$ 是 $D(A, \text{d})$ 的紧对象。由《导出范畴》引理
\ref{derived-lemma-compact-objects-subcategory}，
$\Ker(x)$ 也是 $D(A, \text{d})$ 的紧对象。

\medskip\noindent
接着，假设 $M$ 是表示 $\Ker(x)$ 的微分分次（右）$A$-模，并假设
$M$ 作为分次 $A$-模是有限且投射的。由于
$k[x, y]/(y^2)$ 上的每个有限分次投射模都是分次自由的，可知 $M$ 作为分次
$k[x, y]/(y^2)$-模是有限自由的（即忘掉微分以后）。置
$N = M/M(x^2 + x)$。考虑正合列
$$
0 \to M \xrightarrow{x^2 + x} M \to N \to 0
$$
由于 $x^2 + x$ 的次数为 $0$，位于 $A$ 的中心，并且
$\text{d}(x^2 + x) = 0$，可知这是微分分次 $A$-模的短正合列。此外，
由 $\text{d}(y) = x^2 + x$ 可知 $N$ 上的微分是线性的。映射
$$
H^{-1}(N) \to H^0(M)
\quad\text{且}\quad
H^0(M) \to H^0(N)
$$
是同构，因为 $M \cong \Ker(x)$ 在 $D(A, \text{d})$ 中成立，所以
$H^*(M) = H^0(M) = k$。计算边界映射可知，作为分次模有
$H^*(N) = k[x, y]/(x, y^2)$；我们略去细节。因为 $N$ 是自由
$k[x, y]/(y^2, x^2 + x)$-模，存在与微分相容的解消
$$
\ldots \to N[2] \xrightarrow{y} N[1] \xrightarrow{y} N \to N/Ny \to 0
$$
由于 $N$ 有界且 $H^0(N) = k[x,y]/(x, y^2)$，由《同调》引理
\ref{homology-lemma-first-quadrant-ss} 可知 $H^0(N/Ny) = k[x]/(x)$。
但 $N/Ny$ 是自由 $k[x]/(x^2 + x) = k \times k$-模的有限复形，
故其上同调的维数必为偶数，矛盾。

\begin{lemma}
\label{examples-lemma-no-good-representatif-compact-object}
存在微分分次代数 $(A, \text{d})$ 和 $D(A, \text{d})$ 的紧对象 $E$，
使得 $E$ 不能由有限且分次投射的微分分次 $A$-模表示。
\end{lemma}

\begin{proof}
见上述讨论。
\end{proof}








\section{两个微分分次范畴}
\label{examples-section-nongraded-differential-graded}

\noindent
本节构造两个满足《微分分次代数》情形 \ref{dga-situation-ABC} 中公理
(A)、(B) 和 (C) 的微分分次范畴，而它们的对象并不带
$\mathbf{Z}$-分次。

\medskip\noindent
{\bf 例 I。}设 $X$ 为拓扑空间。以
% P12-ERR-0381
$\underline{\mathbf{Z}}$ 表示取值为 $\mathbf{Z}$ 的常值层。设 $A$ 为
$\underline{\mathbf{Z}}$-挠子。在此情形下，我们称 Abel 群层
$\mathcal{F}$ 是{\it $A$-分次的}，若对给定的局部截面
$a \in A(U)$，存在投影算子
$p_a : \mathcal{F}|_U \to \mathcal{F}|_U$，使得每当存在局部同构
$\underline{\mathbf{Z}}|_U \to A|_U$ 时，都有
$\mathcal{F}|_U = \bigoplus_{n \in \mathbf{Z}} p_n(\mathcal{F})$。
换一种说法，Abel 层 $\mathcal{F}$ 在 $X$ 上局部具有
$\mathbf{Z}$-分次，但在交叠处，不同分次选择相差由挠子 $A$ 的转移函数
给出的次数平移。
如果 $\mathcal{F}$ 是 $A$-分次 Abel 层，而
$\text{d} : \mathcal{F} \to \mathcal{F}$ 是满足
$\text{d}^2 = 0$ 的微分，并且对 $A$ 的每个局部截面 $a$ 都有
$p_{a + 1} \circ \text{d} = \text{d} \circ p_a$，则称偶对
$(\mathcal{F}, \text{d})$ 是{\it Abel 层的 $A$-分次复形}。
换言之，$\text{d}(p_a(\mathcal{F}))$ 包含于
$p_{a + 1}(\mathcal{F})$。

\medskip\noindent
接着，考虑范畴 $\mathcal{A}$，其中
\begin{enumerate}
\item 对象是 Abel 层的 $A$-分次复形；并且
\item 对对象 $(\mathcal{F}, \text{d})$、$(\mathcal{G}, \text{d})$，置
$$
\Hom_\mathcal{A}((\mathcal{F}, \text{d}), (\mathcal{G}, \text{d})) =
\bigoplus \Hom^n(\mathcal{F}, \mathcal{G})
$$
其中 $\Hom^n(\mathcal{F}, \mathcal{G})$ 是满足下述条件的 Abel 层映射
$f$ 所成的群：对 $A$ 的所有局部截面 $a$，都有
$f(p_a(\mathcal{F})) \subset p_{a + n}(\mathcal{G})$。微分取为
$\text{d}(f) = \text{d} \circ f - (-1)^n f \circ \text{d}$；见
《微分分次代数》例 \ref{dga-example-category-complexes}。
\end{enumerate}
我们略去验证：这确实是满足 (A)、(B) 和 (C) 的微分分次范畴。所有性质都可
在 $X$ 上局部验证；在那里，我们恰好恢复 Abel 层复形所成的微分分次范畴。
因而得到三角范畴 $K(\mathcal{A})$。

\medskip\noindent
$X$ 的扭曲导出范畴。注意，给定 $\mathcal{A}$ 的对象
$(\mathcal{F}, \text{d})$，存在良定义的 $A$-分次上同调层
$H(\mathcal{F}, \text{d})$。因此，$K(\mathcal{A})$ 中的拟同构之含义是
清楚的。将拟同构取逆，就得到{\it $A$-分次层复形的导出范畴
$D(\mathcal{A})$}。若 $A$ 是平凡挠子，则 $D(\mathcal{A})$ 等于
$D(X)$；但对
% P12-ERR-0382
非零挠子，会得到 $X$ 的一种{\it 扭曲}导出范畴。

\medskip\noindent
{\bf 例 II。}设 $C$ 为完美域 $k$ 上的光滑曲线，且其特征为 $2$。
则 $\Omega_{C/k}$ 带有典范平方根。也就是说，可以写成
$\Omega_{C/k} = \mathcal{L}^{\otimes 2}$，使得对 $C$ 上每个局部函数
$f$，截面 $\text{d}(f)$ 都等于某个 $\mathcal{L}$ 的局部截面 $s$
的 $s^{\otimes 2}$。其“理由”是
$$
\text{d}(a_0 + a_1t + \ldots +a_dt^d) =
(\sum\nolimits_{i\text{ 为奇数}} a_i^{1/2} t^{(i - 1)/2})^2\text{d}t
$$
% P12-ERR-0383
（在此插入将来的参考文献）。特别地，这确定了典范联络
$$
\nabla_{can} :
\Omega_{C/k}
\longrightarrow
\Omega_{C/k} \otimes_{\mathcal{O}_C} \Omega_{C/k}
$$
其 $2$-曲率为零（即唯一使平方的导数为零的联络）。注意，带联络的向量丛范畴
是张量范畴，故对所有 $n \in \mathbf{Z}$，我们还在可逆层
$\Omega_{C/k}^{\otimes n}$ 上得到典范联络 $\nabla_{can}$。

\medskip\noindent
令 $\mathcal{A}$ 为如下范畴：
\begin{enumerate}
\item 对象是偶对 $(\mathcal{F}, \nabla)$，其中 $\mathcal{F}$ 是带有联络
$$
\nabla :
\mathcal{F}
\longrightarrow
\mathcal{F} \otimes_{\mathcal{O}_C} \Omega_{C/k}
$$
的有限局部自由层，并且该联络的 $2$-曲率为零；并且
\item $(\mathcal{F}, \nabla_\mathcal{F})$ 和
$(\mathcal{G}, \nabla_\mathcal{G})$ 之间的态射由
$$
\Hom_\mathcal{A}((\mathcal{F}, \nabla_\mathcal{F}),
(\mathcal{G}, \nabla_\mathcal{G})) =
\bigoplus \Hom_{\mathcal{O}_C}(\mathcal{F},
\mathcal{G} \otimes_{\mathcal{O}_C} \Omega_{C/k}^{\otimes n})
$$
给出。对次数为 $n$ 的元素
$f : \mathcal{F} \to \mathcal{G} \otimes \Omega_{C/k}^{\otimes n}$，置
$$
\text{d}(f) =
\nabla_{\mathcal{G} \otimes \Omega_{C/k}^{\otimes n}} \circ f +
f \circ \nabla_\mathcal{F}
$$
其中采用适当的等同。
\end{enumerate}
我们略去验证：这构成具有性质 (A)、(B)、(C) 的微分分次范畴。因此得到三角
同伦范畴 $K(\mathcal{A})$。

\medskip\noindent
若 $C = \mathbf{P}^1_k$，则 $K(\mathcal{A})$ 是零范畴。然而，若 $C$
是亏格 $> 1$ 的光滑真曲线，则 $K(\mathcal{A})$ 非零。事实上，假设
$\mathcal{N}$ 是次数满足 $0 \leq d < g - 1$ 且具有非零截面
$\sigma$ 的可逆层。置
$(\mathcal{F}, \nabla_\mathcal{F}) = (\mathcal{O}_C, \text{d})$
以及
$(\mathcal{G}, \nabla_\mathcal{G}) = (\mathcal{N}^{\otimes 2}, \nabla_{can})$。
可见
% P12-ERR-0384
$$
\Hom_\mathcal{A}^n((\mathcal{F}, \nabla_\mathcal{F}),
(\mathcal{G}, \nabla_\mathcal{G})) =
\left\{
\begin{matrix}
0 & \text{若} & n < 0 \\
\Gamma(C, \mathcal{N}^{\otimes 2}) & \text{若} & n = 0 \\
\Gamma(C, \mathcal{N}^{\otimes 2} \otimes \Omega_{C/k}) & \text{若} & n = 1
\end{matrix}
\right.
$$
% P12-ERR-0385
第一个 $0$ 是因为，由条件 $d < g - 1$，
$\mathcal{N}^{\otimes 2} \otimes \Omega_{C/k}^{\otimes -1}$ 的次数为负。
现在，截面 $\sigma^{\otimes 2}$ 的导数等于零，因而同态群
$$
\Hom_{K(\mathcal{A})}((\mathcal{F}, \nabla_\mathcal{F}),
(\mathcal{G}, \nabla_\mathcal{G}))
$$
非零。







\section{真代数空间的叠不是代数叠}
\label{examples-section-proper-spaces-not-algebraic}

\noindent
在《Quot》章节 \ref{quot-section-stack-of-spaces} 中，我们引入并研究了群胚叠
$$
p'_{fp, flat, proper} :
\Spacesstack'_{fp, flat, proper}
\longrightarrow
\Sch_{fppf}
$$
它在概形 $S$ 上的截面范畴是 $S$ 上平坦、真、有限表现的代数空间所成的
范畴。我们证明了它满足许多 Artin 公理。本节将
% P12-ERR-0386
说明这个叠为什么不是代数叠：一般而言，形式有效性不成立。

\medskip\noindent
典范例利用如下事实：$g$ 维阿贝尔簇的泛形变空间有 $g^2$ 个形式参数，
而任何有效形式形变都可定义在维数 $\leq g(g + 1)/2$ 的完备局部环上。
我们的例子将通过写下 K3 曲面的适当非有效形变来构造。我们只勾勒论证，
不提供全部细节。

\medskip\noindent
令 $k = \mathbf{C}$ 为复数域。设 $X \subset \mathbf{P}^3_k$ 为 $k$ 上的
光滑 $4$ 次曲面。我们有
$\omega_X \cong \Omega^2_{X/k} \cong \mathcal{O}_X$。最后，还有
$\dim_k H^0(X, T_{X/k}) = 0$、
$\dim_k H^1(X, T_{X/k}) = 20$，以及
$\dim_k H^2(X, T_{X/k}) = 0$。
由于 $X$ 在 $k$ 上光滑，$L_{X/k} = \Omega_{X/k}$；又因为
$\Ext^i_{\mathcal{O}_X}(\Omega_{X/k}, \mathcal{O}_X) =
H^i(X, T_{X/k})$，并利用《余切复形》引理
\ref{cotangent-lemma-find-obstruction-ringed-topoi}，可知 $X$ 在
$$
k[[x_1, \ldots, x_{20}]]
$$
上存在泛形变。假设该泛形变是有效的（定义见《Artin 公理》章节
\ref{artin-section-formal-objects}）。那么将得到平坦真态射
$$
f : Y \longrightarrow \Spec(k[[x_1, \ldots, x_{20}]])
$$
其中 $Y$ 是恢复该泛形变的代数空间。下面说明这是不可能的。因为 $Y$ 分离，
可找到仿射开子概形 $V \subset Y$。由于 $Y$ 的特殊纤维 $X$ 光滑，可知
$f$ 光滑。因此，$Y$ 在正则空间上光滑，故为正则的；从而 $V$ 在 $Y$
中的补集 $D$ 是有效 Cartier 除子。于是 $\mathcal{O}_Y(D)$ 是
$\Pic(Y)$ 的非平凡元素（为证明这一点，可说明域上真光滑代数空间中非空仿射
开集的补集总在 Picard 群中给出
% P12-ERR-0387
一个非平凡元素，再把它应用于 $f$ 的泛纤维）。最后，为得到矛盾，我们证明
$\Pic(Y) = 0$。事实上，映射 $\Pic(Y) \to \Pic(X)$ 是单射，因为
$H^1(X, \mathcal{O}_X) = 0$（所以 $\mathcal{O}_X$ 到
% P12-ERR-0388
$Y \times \Spec(k[[x_i]]/\mathfrak m^n)$ 的所有形变都是平凡的），
并且由 Grothendieck 存在定理（它断言：在各加厚上给出相同层的凝聚模彼此
同构）。若 $X$ 足够一般，则 $\Pic(X) = \mathbf{Z}$，由
$\mathcal{O}_X(1)$ 生成。因此只须说明，对 $n > 0$，
$\mathcal{O}_X(n)$ 不形变到一阶邻域\footnote{只要映射
$c_1 : \Pic(X) \to H^1(X, \Omega_{X/k})$ 是单射，该论证就成立；
对任意特征为零的域 $k$ 以及 $\mathbf{P}^3_k$ 中任意光滑 $4$ 次超曲面
$X$，这一点都成立。}。考虑杯积
$$
H^1(X, \Omega_{X/k}) \times H^1(X, T_{X/k})
\longrightarrow H^2(X, \mathcal{O}_X)
$$
由凝聚对偶，这是非退化配对。计算表明，Hodge 上同调中的 Chern 类
$c_1(\mathcal{O}_X(n)) \in H^1(X, \Omega_{X/k})$ 非零。
因此存在一个一阶形变，它与 $c_1(\mathcal{O}_X(n))$ 的杯积非零。
最后可说明，该杯积就是提升的障碍类。

\begin{lemma}
\label{examples-lemma-proper-spaces-not-algebraic}
群胚叠
$$
p'_{fp, flat, proper} :
\Spacesstack'_{fp, flat, proper}
\longrightarrow
\Sch_{fppf}
$$
在概形 $S$ 上的截面范畴是 $S$ 上平坦、真、有限表现的代数空间所成的范畴
（见《Quot》章节 \ref{quot-section-stack-of-spaces}）；它不是代数叠。
\end{lemma}

\begin{proof}
若它是代数叠，则每个形式对象都将是有效的；见《Artin 公理》引理
\ref{artin-lemma-effective}。上述讨论
% P12-ERR-0389
表明，在基变换到 $\Spec(\mathbf{C})$ 后并非如此。结论得证。
\end{proof}






\section{非代数 Hom-叠的一个例}
\label{examples-section-non-algebraic-hom-stack}

\noindent
设 $\mathcal{Y}, \mathcal{Z}$ 为概形 $S$ 上的代数叠。{\it Hom-叠}
$\underline{\Mor}_S(\mathcal{Y}, \mathcal{Z})$ 是 $S$ 上的群胚叠，
其在概形 $T$ 上的截面范畴为
$$
\Mor_T(\mathcal{Y} \times_S T, \mathcal{Z} \times_S T)
$$
这里的对象是 $1$-态射，而态射是 $2$-态射。我们略去证明它确实是
$(\Sch/S)_{fppf}$ 上的群胚叠
% P12-ERR-0390
（在此插入将来的参考文献）。当然，一般而言，Hom-叠并非代数叠。
本节给出一个反例，其中 $\mathcal{Y}$ 可由 $S$ 上的真平坦概形表示，
而 $\mathcal{Z}$ 在 $S$ 上光滑且真。

\medskip\noindent
设 $k$ 为代数闭域，但不是有限域的代数闭包。令
$S = \Spec(k[[t]])$ 且 $S_n = \Spec(k[t]/(t^n)) \subset S$。
设 $f : X \to S$ 为满足下列条件的映射：
\begin{enumerate}
\item $f$ 射影且平坦，并且它的纤维是几何连通曲线；
\item
% P12-ERR-0391
纤维 $X_0 = X \times_S S_0$ 是结点曲线，其光滑不可约分支的对偶图含有
一个由有理曲线组成的圈；
\item $X$ 是正则概形。
\end{enumerate}
例如，为构造这样的曲面 $X$，可在 $\mathbf{P}^2_{k[[t]]}$ 中取
$$
X\quad :\quad T_0T_1T_2 - t(T_0^3 + T_1^3 + T_2^3) = 0
$$
令 $A_0$ 为 $k$ 上的非零阿贝尔簇，例如一条椭圆曲线。令
$A = A_0 \times_{\Spec(k)} S$ 为与 $A_0$ 对应的 $S$ 上常值 Abel
概形\footnote{基底 $S$ 上的 Abel 概形是 $S$ 上平坦、真、有限表现的群概形，
且其所有纤维都是阿贝尔簇。}。我们将证明叠
% P12-ERR-0392
$\mathcal{X} = \underline{\Mor}_S(X, [S/A]))$ 不是代数叠。

\medskip\noindent
回忆，$[S/A]$ 一方面是 $A$ 在 $S$ 上平凡作用的商叠，另一方面等于分类
fppf $A$-挠子的叠；见《叠的例》命题
\ref{examples-stacks-proposition-equal-quotient-stacks}。
注意，$[S/A] = [\Spec(k)/A_0] \times_{\Spec(k)} S$。这使我们可以如下描述
概形 $T$ 上的纤维范畴：
\begin{align*}
\mathcal{X}_T
& =
\underline{\Mor}_S(X, [S/A])_T \\
& =
\Mor_T(X \times_S T, [S/A] \times_S T) \\
& =
\Mor_S(X \times_S T, [S/A]) \\
& =
\Mor_{\Spec(k)}(X \times_S T, [\Spec(k)/A_0])
\end{align*}
这里 $T$ 是任意 $S$-概形。换言之，群胚 $\mathcal{X}_T$ 是
$X \times_S T$ 上 fppf $A_0$-挠子所成的群胚。在讨论
$\mathcal{X}$ 为什么不是代数叠之前，需要几个引理。

\begin{lemma}
\label{examples-lemma-torsors-over-two-dimensional-regular}
设 $W$ 为二维正则整 Noether 概形，分式域为 $K$。设
$G \to W$ 为 Abel 概形。则映射
$H^1_{fppf}(W, G) \to H^1_{fppf}(\Spec(K), G)$ 是单射。
\end{lemma}

\begin{proof}[证明概要]
设 $P \to W$ 为在泛点处平凡的 fppf $G$-挠子。于是有 $W$ 上的态射
$\Spec(K) \to P$，并可取其概形论像 $Z \subset P$。由于
$P \to W$ 是真的（因为它是 $W$ 上真群代数空间的挠子），可知
$Z \to W$ 是真双有理态射。由《域上的空间》引理
\ref{spaces-over-fields-lemma-finite-in-codim-1}，态射 $Z \to W$
在除 $W$ 的有限多个闭点以外之处都是有限的。由
% P12-ERR-0393
（在此插入关于通过曲面二次变换的吹起消解态射不定性的将来参考文献），
$Z \to W$ 的几何纤维的不可约分支都是有理曲线。由《空间中的群胚进阶》引理
\ref{spaces-more-groupoids-lemma-no-nonconstant-morphism-from-P1-to-group}，
不存在从有理曲线到群概形或其挠子的非常值态射。因此 $Z \to W$ 是有限的，
从而 $Z$ 是概形，并且由《态射》引理
\ref{morphisms-lemma-finite-birational-over-normal}，$Z \to W$ 是同构。
换言之，挠子 $P$ 是平凡的。
\end{proof}

\begin{remark}
\label{examples-remark-torsors-over-regular}
可以去掉引理 \ref{examples-lemma-torsors-over-two-dimensional-regular} 中
$W$ 的 $2$-维条件。也就是说，若 $W$ 是 Noether、正则、整的，则：
(1) 对每个 $G$-挠子 $P$，$P(W)\to P(K)$ 是双射；并且
(2) $H^1(W,G)\to H^1(K,G)$ 是单射。证明概要如下。
首先，(2) 显然由 (1) 得到。对 (1)，单射性显然；满射性则像当前证明一样，
归结为证明有理截面的概形论闭包是截面。这个问题在 $W$ 上是 fppf 局部的，
所以经过 \'etale 基变换后，可假设 $P=G$（注意 \'etale 覆盖保持正则性）。
然后注意，$G=\underline{\mathrm{Pic}}^0_{G'/W}$，其中 $G'$ 是对偶 Abel
概形。结论来自如下事实：$G'$ 是正则概形，故泛定义的可逆层可以延拓。
事实上，可把“$W$ 正则”弱化为“光滑 $W$-概形局部为因子概形”。
Laurent Moret-Bailly 从 Raynaud 那里学到了这个 Picard 技巧；
Raynaud 的著作中某处应当有相应参考文献。
\end{remark}

\begin{lemma}
\label{examples-lemma-torsors-over-field-torsion}
设 $G$ 为域 $K$ 上的光滑交换群代数空间。则
$H^1_{fppf}(\Spec(K), G)$ 是挠群。
\end{lemma}

\begin{proof}
作为 $G$ 的形式，$\Spec(K)$ 上的每个 $G$-挠子 $P$ 都在 $K$ 上光滑。
因此，$P$ 在某个有限可分扩张 $L/K$ 上有点。设 $[L : K] = n$。
以 $[n](P)$ 表示如下 $G$-挠子：其类是 $P$ 的类在
$H^1_{fppf}(\Spec(K), G)$ 中的 $n$ 倍。存在 $K$ 上代数空间的典范态射
$$
P \times_{\Spec(K)} \ldots \times_{\Spec(K)} P \to [n](P)
$$
由于 $G$ 交换，该态射是对称的，故它通过商
$$
(P \times_{\Spec(K)} \ldots \times_{\Spec(K)} P)/S_n
$$
分解。另一方面，态射 $\Spec(L) \to P$ 定义了态射
$$
(\Spec(L) \times_{\Spec(K)} \ldots \times_{\Spec(K)} \Spec(L))/S_n
\longrightarrow (P \times_{\Spec(K)} \ldots \times_{\Spec(K)} P)/S_n
$$
读者可以验证，左边的概形有一个 $K$-有理点。因此，$[n](P)$ 是平凡挠子。
\end{proof}

\noindent
为证明 $\mathcal{X} = \underline{\Mor}_S(X, [S/A])$ 不是代数叠，
由《Artin 公理》引理 \ref{artin-lemma-effective}，只须证明如下结论。

\begin{lemma}
\label{examples-lemma-not-essentially-surjective}
典范映射 $\mathcal{X}(S) \to \lim \mathcal{X}(S_n)$ 不是本质满的。
\end{lemma}

\begin{proof}[证明概要]
展开定义，只须验证 $H^1(X, A_0) \to \lim H^1(X_n, A_0)$ 不是满射。
由于 $X$ 正则且射影，由引理 \ref{examples-lemma-torsors-over-field-torsion} 和
\ref{examples-lemma-torsors-over-two-dimensional-regular}，$X$ 上的每个
$A_0$-挠子都是挠的。特别地，群 $H^1(X, A_0)$ 是挠群。因此只须说明：
(a) 群 $H^1(X_0, A_0)$ 不是挠群；以及
(b) 对所有 $n$，映射
$H^1(X_{n + 1}, A_0) \to H^1(X_n, A_0)$ 都是满射。

\medskip\noindent
关于 (a)。在结点曲线 $X_0$ 上，通过在 $X_0$ 的每个分支上取平凡
$A_0$-挠子，并以 $A_0$ 的非挠点作为各结点上方的同构来粘合，可以构造
非挠 $A_0$-挠子 $P_0$。更确切地说，令 $x \in X_0$ 为由有理曲线组成的
一个圈中出现的结点。令 $X'_0 \to X_0$ 为 $X_0$ 在
$X_0 \setminus \{x\}$ 中的正规化。令 $x', x'' \in X'_0$ 为映到
% P12-ERR-0394
$x_0$ 的两个点。然后取 $A_0 \times_{\Spec(k)} X'_0$，并利用非挠点
$a_0 \in A_0(k)$ 给出的平移 $A_0 \to A_0$，把
$A_0 \times {x'}$ 与 $A_0 \times \{x''\}$ 等同（这样的非挠点存在，
因为 $k$ 代数闭但不是有限域的代数闭包；这实际上并非容易证明）。
可以说明该粘合是代数空间（事实上，可以说明它是概形），并且它是 $X_0$
上的
% P12-ERR-0395
非挠 $A_0$-挠子。它之所以非挠，是因为若 $[n](P_0)$ 有截面，则该截面
产生态射 $s : X'_0 \to A_0$，使得在 $A_0(k)$ 的群律下有
$[n](a_0) = s(x') - s(x'')$。然而，由于圈的不可约分支都是有理曲线，
% P12-ERR-0396
截面 $s$ 在其上为常值（《空间中的群胚进阶》引理
\ref{spaces-more-groupoids-lemma-no-nonconstant-morphism-from-P1-to-group}）。
故 $s(x') = s(x'')$，得到矛盾。

\medskip\noindent
关于 (b)。形变理论表明，把 $A_0$-挠子 $P_n \to X_n$ 形变为
$A_0$-挠子 $P_{n + 1} \to X_{n + 1}$ 的障碍位于
$H^2(X_0, \omega)$ 中，其中 $\omega$ 是 $X_0$ 上的某个向量丛。
后一个群因为 $X_0$ 是曲线而消失，这就证明了断言。
\end{proof}

\begin{proposition}
\label{examples-proposition-nonalghomstack}
叠 $\mathcal{X} = \underline{\Mor}_S(X, [S/A])$ 不是代数叠。
\end{proposition}

\begin{proof}
见上述讨论。
\end{proof}

\begin{remark}
\label{examples-remark-contradict-aoki}
命题 \ref{examples-proposition-nonalghomstack} 与 \cite[定理 1.1]{AokiHomStacks}
矛盾。问题在于 $\underline{\Mor}_S(X, [S/A])$ 的形式对象不是有效的。
勘误 \cite{AokiHomStacksErr} 也提到了同一问题；这是对
\cite{AokiHomStacks} 的勘误。
遗憾的是，该勘误接着断言：若 $\mathcal{Z}$ 分离，则
$\underline{\Mor}_S(\mathcal{Y}, \mathcal{Z})$ 是代数叠；
这也与命题 \ref{examples-proposition-nonalghomstack} 矛盾，因为 $[S/A]$ 是分离的。
\end{remark}









\section{不满足强形式有效性的代数叠}
\label{examples-section-non-formal-effectiveness}

\noindent
这是 \cite[例 4.12]{Bhatt-Algebraize}。设 $k$ 为代数闭域，并设 $A$ 为
$k$ 上的阿贝尔簇。假设 $A(k)$ 不是挠群（若 $k$ 不是有限域的代数闭包，
这总成立）。令 $\mathcal{X} = [\Spec(k)/A]$。我们断言，存在理想
$I \subset k[x, y]$，使得
$$
\mathcal{X}_{\Spec(k[x, y]^\wedge)}
\longrightarrow
\lim \mathcal{X}_{\Spec(k[x, y]/I^n)}
$$
不是本质满的。事实上，令 $I$ 为由 $xy(x + y - 1)$ 生成的理想。
于是 $X_0 = V(I)$ 由三份 $\mathbf{A}^1_k$ 沿三个点粘成三角形而得。
因此，可以在 $X_0$ 的各不可约分支上取平凡挠子，并使用非挠点给出的平移
进行粘合，从而构造 $X_0$ 上的无限阶 $A$-挠子 $P_0$。与引理
\ref{examples-lemma-not-essentially-surjective} 的证明完全相同，可以把 $P_0$
提升为 $X_n = \Spec(k[x, y]/I^n)$ 上的挠子 $P_n$。
由于 $k[x, y]^\wedge$ 是二维正则整环，由引理
\ref{examples-lemma-torsors-over-two-dimensional-regular} 和
\ref{examples-lemma-torsors-over-field-torsion} 可知，
$\Spec(k[x, y]^\wedge)$ 上任意 $A$-挠子 $P$ 都是挠的。
所以这族挠子不在所显示函子的像中。

\begin{lemma}
\label{examples-lemma-non-formal-effectiveness}
设 $k$ 为不是有限域
% P12-ERR-0397
闭包的代数闭域。设 $A$ 为 $k$ 上的阿贝尔簇，并令
$\mathcal{X} = [\Spec(k)/A]$。存在具有满转移映射的 $k$-代数逆系
$R_n$，其转移映射的核局部幂零；并存在位于系统
% P12-ERR-0398
$(\Spec(R_n))$ 上方的 $\mathcal{X}$ 的系统 $(\xi_n)$，使得该系统在
《Artin 公理》注 \ref{artin-remark-strong-effectiveness} 的意义下不是
有效的。
\end{lemma}

\begin{proof}
见上述讨论。
\end{proof}








\section{Grothendieck 存在定理的一个反例}
\label{examples-section-Grothendieck-existence}

\noindent
设 $k$ 为域，并令 $A = k[[t]]$。令 $X$ 为如下粘合：把
$U = \Spec(A[x])$ 和 $V = \Spec(A[y])$ 通过等同
$$
U \setminus \{0_U\} \longrightarrow V \setminus \{0_V\}
$$
粘合，其中该等同把 $x$ 送到 $y$，而
% P12-ERR-0399
$0_U \in U$ 和 $O_V \in V$ 是分别对应于极大理想 $(x, t)$ 和
$(y, t)$ 的点。置 $A_n = A/(t^n)$，并置
$X_n = X \times_{\Spec(A)} \Spec(A_n)$。令 $\mathcal{F}_n$ 为
$X_n$ 上的凝聚层；它对应于 $A_n[x]$-模
$A_n[x]/(x) \cong A_n$ 和 $A_n[y]$-模 $0$，二者按显然方式粘合。
令 $\mathcal{I} \subset \mathcal{O}_X$ 为
% P12-ERR-0400
由 $t$ 生成的理想层。则 $(\mathcal{F}_n)$ 是《概形的上同调》章节
\ref{coherent-section-existence-proper-support} 中定义的范畴
$\textit{Coh}_{\text{支撑在 }A\text{ 上为固有}}(X, \mathcal{I})$ 的对象。
另一方面，该对象不在《概形的上同调》方程
(\ref{coherent-equation-completion-functor-proper-over-A}) 中函子的像内。
事实上，若它
% P12-ERR-0401
在那里，则将存在有限 $A[x]$-模 $M$、有限 $A[y]$-模 $N$ 和同构
$M[1/t] \cong N[1/t]$，使得对所有 $n$ 都有
$M/t^nM \cong A_n[x]/(x)$ 且 $N/t^nN = 0$。容易看出这是不可能的。

\begin{lemma}
\label{examples-lemma-counter-Grothendieck-existence}
凝聚层代数化的反例。
\begin{enumerate}
\item 若删去 $X \to \Spec(A)$ 分离的假设，则《概形的上同调》定理
\ref{coherent-theorem-grothendieck-existence} 中陈述的 Grothendieck
存在定理为假。
\item 若删去 $X \to S$ 分离的假设，则《Quot》定理
\ref{quot-theorem-coherent-algebraic-general} 和
\ref{quot-theorem-coherent-algebraic} 中的凝聚层叠
$\Cohstack_{X/B}$ 一般不是代数叠。
\item 若删去 $X \to B$ 分离的假设，则《Quot》命题
\ref{quot-proposition-quot} 中的函子
$\Quotfunctor_{\mathcal{F}/X/B}$ 一般不是代数空间。
\end{enumerate}
\end{lemma}

\begin{proof}
第 (1) 部分已在上面说明。这表明 $\textit{Coh}_{X/A}$ 不满足
《Artin 公理》章节 \ref{artin-section-axioms} 的公理 [4]。
因此，由《Artin 公理》引理 \ref{artin-lemma-effective}，它不可能是代数叠。
这就说明了 (2)。为看出 (3)，注意对所有 $n$ 都有相容的满射
$\mathcal{O}_{X_n} \to \mathcal{F}_n$。因此，
$\Quotfunctor_{\mathcal{O}_X/X/A}$ 不满足公理 [4]，并可像前面一样看出
(3) 成立。
\end{proof}






\section{仿射形式代数空间}
\label{examples-section-affine-formal-algebraic-space}

\noindent
设 $K$ 为域，并设 $(V_i)_{i \in I}$ 为 $K$ 上非零向量空间的有向逆系，
其转移映射为满射且 $\lim V_i = 0$；见章节 \ref{examples-section-zero-limit}。
令 $R_i = K \oplus V_i$ 为 $K$-代数，其中 $V_i$ 是平方为零的理想。
于是 $R_i$ 构成 $K$-代数的逆系，其转移映射为满射，核为幂零理想，
并且 $\lim R_i = K$。仿射形式代数空间
$X = \colim \Spec(R_i)$ 是一个并非 McQuillan 的仿射形式代数空间之例。

\begin{lemma}
\label{examples-lemma-affine-not-mcquillan}
存在一个并非 McQuillan 的仿射形式代数空间。
\end{lemma}

\begin{proof}
见上述讨论。
\end{proof}

\noindent
令 $0 \to W_i \to V_i \to K \to 0$ 为章节 \ref{examples-section-zero-limit}
中的正合列系统。令
$A_i = K[V_i]/(ww'; w, w' \in W_i)$。
于是存在相容的满射系统 $A_i \to K[t]$，其核幂零；而转移映射
$A_i \to A_j$ 也是核幂零的满射。回忆，$V_i$ 在 $K$ 上自由，其基由
$s \in S_i$ 给出。于是，若 $K$ 的特征为零，$A_i$ 的次数 $d$ 部分在
$K$ 上自由，其基由 $s^d$（$s \in S_i$）给出，而每个基元都映到 $t^d$。
故 $A_i$ 的次数 $d$ 部分所成的逆系同构于向量空间 $V_i$ 所成的逆系。
由于 $\lim V_i = 0$，可知 $\lim A_i = K$，至少当 $K$ 的特征为零时如此。
这给出一个仿射形式代数空间的例，其“正则函数”不能分离点。

\begin{lemma}
\label{examples-lemma-affine-formal-functions-do-not-separate-points}
存在仿射形式代数空间 $X$，其正则函数在如下意义下不能分离点：
若按《形式空间》定义
\ref{formal-spaces-definition-affine-formal-algebraic-space}
写成 $X = \colim X_\lambda$，则
$\lim \Gamma(X_\lambda, \mathcal{O}_{X_\lambda})$ 是域，但
$X_{red}$ 有无穷多个点。
\end{lemma}

\begin{proof}
见上述讨论。
\end{proof}

\noindent
设 $K$、$I$ 和 $(V_i)$ 如上。考虑系统
$$
\Phi = (\Lambda, J_i \subset \Lambda, (M_i) \to (V_i))
$$
其中 $\Lambda$ 是增广 $K$-代数；对 $i \in I$，
$J_i \subset \Lambda$ 是平方为零的理想；$(M_i) \to (V_i)$ 是
$K$-向量空间逆系的映射，使得对每个 $i$，$M_i \to V_i$ 都是满射；
$M_i$ 带有 $\Lambda$-模结构；转移映射 $M_i \to M_j$（$i > j$）
为 $\Lambda$-线性的；并且当 $i > j$ 时，
$J_j M_i \subset \Ker(M_i \to M_j)$。
断言：存在上述系统，使得对所有 $i > j$ 都有 $M_j = M_i/J_j M_i$。

\medskip\noindent
若断言成立，就得到仿射形式代数空间的可表态射
$$
\colim_{i \in I} \Spec(\Lambda/J_i \oplus M_i)
\longrightarrow
\text{Spf}(\lim \Lambda/J_i)
$$
其源不是 McQuillan 的，而靶是 McQuillan 的。这里
$\Lambda/J_i \oplus M_i$ 带有通常的 $\Lambda/J_i$-代数结构，
其中 $M_i$ 是平方为零的理想。可表性恰好转化为条件：
当 $i > j$ 时，$M_i/J_jM_i = M_j$。该态射的源不是 McQuillan 的，
因为投影
% P12-ERR-0402
$\lim_{i \in I} M_i \to M_i$ 不是满射。这是因为映射
$\lim V_i \to V_i$ 不是满射，并且有满射 $M_i \to V_i$。
略去一些细节。

\medskip\noindent
断言的证明。首先注意，至少存在一个系统，即
$$
\Phi_0 = (K, J_i = (0), (V_i) \xrightarrow{\text{id}} (V_i))
$$
给定系统 $\Phi$，我们将证明存在系统态射 $\Phi \to \Phi'$（系统态射按
显然方式定义），使得 $\Ker(M_i/J_j M_i \to M_j)$ 在
$M'_i/J'_j M'_i$ 中的像为零。做到这一点以后，可以使用通常技巧：
对 $n \geq 1$ 归纳地置 $\Phi_n = (\Phi_{n - 1})'$，再取
$\Phi = \colim \Phi_n$，得到具有所需性质的系统。略去细节。

\medskip\noindent
给定 $\Phi$，构造 $\Phi'$。考虑三元组
$u = (i, j, \xi)$ 所成的集合 $U$，其中 $i > j$ 且
$\xi \in \Ker(M_i \to M_j)$。以 $s, t : U \to I$ 表示映射
$s(i, j, \xi) = i$ 和 $t(i, j, \xi) = j$。令
$\xi_u \in M_{s(u)}$ 为 $u$ 的第三个分量。取
$$
\Lambda' = \Lambda[x_u; u \in U]/(x_u x_{u'}; u, u' \in U)
$$
其中增广 $\Lambda' \to K$ 由 $\Lambda$ 的增广给出，并把 $x_u$
送到零。取
% P12-ERR-0403
$J'_k = J_k \Lambda' + (x_{u,\ t(u) \geq k})$。置
$$
M'_i = M_i \oplus \bigoplus\nolimits_{s(u) \geq i} K\epsilon_{i, u}
$$
作为 $i > j$ 时的转移映射 $M'_i \to M'_j$，我们使用给定映射
$M_i \to M_j$，并把 $\epsilon_{i, u}$ 送到 $\epsilon_{j, u}$。
映射 $M'_i \to V_i$ 诱导给定映射 $M_i \to V_i$，并把
$\epsilon_{i, u}$ 送到零。最后，令 $\Lambda'$ 如下作用在 $M'_i$ 上：
对 $\lambda \in \Lambda$，在 $M_i$ 上按 $\Lambda$-模结构作用，并通过增广
$\Lambda \to K$ 作用在 $\epsilon_{i, u}$ 上。对所有 $i$，元素 $x_u$
在 $M_i$ 上作用为 $0$。最后定义
$$
x_u \epsilon_{i, u} = \text{元素 }\xi_u\text{ 映至 }M_i
$$
并令所有其他乘积 $x_{u'} \epsilon_{i, u}$ 等于零。所显示公式有意义，
因为 $s(u) \geq i$ 且 $\xi_u \in M_{s(u)}$。主要需要
% P12-ERR-0404
检验的是：当 $i > j$ 时，$J'_j M'_i \subset M'_i$ 在 $M'_j$ 中映为零；
以及
% P12-ERR-0405
$\Ker(M_i \to M_j)$ 在 $M'_i/J_j M'_i$ 中映为零。后一个事实的理由是，
对任意 $\xi \in \Ker(M_i \to M_j)$，都有
$\xi = x_{(i, j, \xi)} \epsilon_{i, (i, j, \xi)} \in J'_j M'_i$。
我们略去细节。

\begin{lemma}
\label{examples-lemma-representable-morphism-affine-formal-not-mcquillan-top}
存在仿射形式代数空间的可表态射 $f : X \to Y$，使得 $Y$ 是
McQuillan 的，而 $X$ 不是 McQuillan 的。
\end{lemma}

\begin{proof}
见上述讨论。
\end{proof}









\section{平坦映射不一定是有限表现平坦映射的滤过余极限}
\label{examples-section-flat-not-colimit-flat-finitely-presented}

\noindent
本节的目标是给出一个平坦环映射的例，它不是平坦有限表现环映射的滤过
余极限。\cite{gabber-nonexcellent} 证明了：若 $A$ 是维数为 $1$、剩余特征
为零的非优局部环，则到其完备化的（平坦）环映射
$A \to A^\wedge$ 不是有限型平坦环映射的滤过余极限。本节例子的源将是优环。
我们鼓励读者提交其他例子；若你有例子，请发送电子邮件到
\href{mailto:stacks.project@gmail.com}{stacks.project@gmail.com}。

\medskip\noindent
为进行构造，固定素数 $p$，并令
$A = \mathbf{F}_p[x_1, \ldots, x_n]$。选取 $A$ 的绝对整闭包 $A^+$；
也就是说，$A^+$ 是 $A$ 在其分式域的某个代数闭包中的整闭包。
\cite[\S 6.7]{HHBigCM} 证明了 $A \to A^+$ 平坦。

\medskip\noindent
我们断言，当 $n \geq 3$ 时，$A$-代数 $A^+$ 不是有限表现平坦
$A$-代数的滤过余极限。

\medskip\noindent
我们勾勒 $n = 3$ 情形的论证，把推广到更高 $n$ 的工作留给读者。
只须对映射 $R \to R^+$ 证明类似陈述，其中 $R$ 是 $A$ 在原点处的严格
Hensel 化，而 $R^+$ 是它的绝对整闭包。注意，$R$ 是 Hensel 正则局部环，
其剩余域 $k$ 是 $\mathbf{F}_p$ 的代数闭包。

\medskip\noindent
选取 $k$ 上的普通 Abel 曲面 $X$ 和 $X$ 上的甚丰富线丛 $L$。
截面环 $\Gamma_*(X, L) = \bigoplus_n H^0(X,L^n)$ 是 $X$ 关于 $L$
的仿射锥的坐标环。当 $L$ 足够正时，它是正规环。令 $S$ 表示
$\Gamma_*(X, L)$ 在
% P12-ERR-0406
锥顶点处的 Hensel 化。则 $S$ 是维数为 $3$ 的 Hensel Noether 正规整环。
取有限型 $k$-代数 $\Gamma_*(X, L)$ 的一个 Noether 正规化并作 Hensel 化，
得到有限单射 $R \to S$。由于 $R^+$ 是 $R$ 的绝对整闭包，还可固定嵌入
$S \to R^+$。因此，$R^+$ 也是 $S$ 的绝对整闭包。为说明 $R^+$ 不是
平坦 $R$-代数的滤过余极限，只须说明：
\begin{enumerate}
\item 若 $S \to P$ 为环映射，且 $P$ 在 $R$ 上平坦、有限型，则存在环映射
$S \to T$，其中 $T$ 在 $R$ 上有限平坦。
\item 对任意有限环映射 $S \to T$，环 $T$ 都不是 $R$-平坦的。
\end{enumerate}
事实上，因为 $S$ 在 $R$ 上有限表现，若可以把
$R^+ = \colim_i P_i$ 写成有限表现平坦 $R$-代数 $P_i$ 的滤过余极限，
则对 $i \gg 0$，$S \to R^+$ 将分解为 $S \to P_i \to R^+$，
这与上述两个断言矛盾。断言 (1) 来自 $R$ 是 Hensel 环这一事实和一个切片
论证；见《态射进阶》引理
\ref{more-morphisms-lemma-qf-fp-flat-neighbourhood-dominates-fppf}。
第 (2) 部分已在 \cite{BhattSmallCMMod} 中证明；为方便读者，我们回忆该论证。

\medskip\noindent
令 $U \subset \Spec(S)$ 为穿孔谱，于是有自然映射
$X \leftarrow U \subset \Spec(S)$。第一个映射给出等同
$H^1(U, \mathcal{O}_U) \simeq H^1(X, \mathcal{O}_X)$。
通过完美化的 Witt 向量并使用 Artin--Schreier 序列\footnote{这里用到：
$S$ 是特征 $p$ 的严格 Hensel 局部环，因而
$S \to S$，$f \mapsto f^p - f$ 是满射。又因为 $S$ 是正规整环，
$\Gamma(U, \mathcal{O}_U) = S$。所以
$H^1_\etale(U, \mathbf{Z}/p)$ 是由 $f \mapsto f^p - f$ 诱导的映射
$H^1(U, \mathcal{O}_U) \to H^1(U, \mathcal{O}_U)$ 的核。}，
得到等同
$H^1_\etale(U, \mathbf{Z}_p) \simeq H^1_\etale(X, \mathbf{Z}_p)$。
特别地，这个群是秩为 $2$ 的有限自由 $\mathbf{Z}_p$-模
（因为 $X$ 是普通的）。为得到矛盾，假设存在上述 (2) 中的
$R$-平坦 $T$。令 $V \subset \Spec(T)$ 表示 $U$ 的逆像，并以
$f : V \to U$ 表示诱导的有限满射。由于 $U$ 正规，$U_\etale$ 上存在迹映射
$f_*\mathbf{Z}_p \to \mathbf{Z}_p$，它与拉回
$\mathbf{Z}_p \to f_*\mathbf{Z}_p$ 的复合是乘以
$d = \deg(f)$。传到上同调，并利用
$H^1_\etale(U, \mathbf{Z}_p)$ 无挠，可知
$H^1_\etale(V, \mathbf{Z}_p)$ 非零。由于不存在 $R^1\lim$ 干扰，有
$H^1_\etale(V, \mathbf{Z}_p) \simeq \lim H^1_\etale(V, \mathbf{Z}/p^n)$，
故群 $H^1(V_\etale,\mathbf{Z}/p)$ 必非零。由于 $T$ 是 $R$-平坦的，
有 $\Gamma(V, \mathcal{O}_V) = T$，这是严格 Hensel 的；Artin--Schreier
序列表明 $H^1(V, \mathcal{O}_V) \neq 0$。这等价于
$H^2_\mathfrak m(T) \neq 0$，其中 $\mathfrak m \subset R$ 是极大理想。
于是得到矛盾，因为 $T$ 作为 $R$-模有限平坦（即有限自由），而
$H^2_\mathfrak m(R) = 0$。这个矛盾证明了 (2)。

\begin{lemma}
\label{examples-lemma-weird-flat-map}
存在交换环 $A$ 和平坦 $A$-代数 $B$，使得该代数不能写成有限表现平坦
$A$-代数的滤过余极限。事实上，可以选取 $A$ 为有限型
$\mathbf{F}_p$-代数，也可以选为剩余域特征为 $0$ 的 $1$-维 Noether
局部环。
\end{lemma}

\begin{proof}
见上述讨论。
\end{proof}

\section{模范畴对挠模取商}
\label{examples-section-serre-quotient-modulo-torsion-modules}

\noindent
挠群范畴是全体 Abel 群范畴的 Serre 子范畴
（《同调》定义 \ref{homology-definition-serre-subcategory}）。更一般地，对任意环
$A$，挠 $A$-模范畴是全体 $A$-模范畴的 Serre 子范畴；见《代数进阶》第
\ref{more-algebra-section-abelian-categories-modules} 节。若 $A$ 是整环，则其商范畴
（《同调》引理 \ref{homology-lemma-serre-subcategory-is-kernel}）等价于分式域上的
向量空间范畴。这可由下面更一般的命题推出。

\begin{proposition}
\label{examples-proposition-localization-and-serre-quotients}
设 $A$ 为环，$S$ 为 $A$ 的乘法子集。以 $\text{Mod}_A$ 表示 $A$-模范畴，
以 $\mathcal{T}$ 表示其如下 Serre 子范畴：其中每个模的任一元素都被 $S$ 的某个
元素零化。则存在典范等价
$\text{Mod}_A/\mathcal{T} \rightarrow \text{Mod}_{S^{-1}A}$。
\end{proposition}

\begin{proof}
由 $M \mapsto M \otimes_A S^{-1}A$ 给出的函子
$\text{Mod}_A \to \text{Mod}_{S^{-1}A}$ 是正合的（由《代数》命题
\ref{algebra-proposition-localization-exact}），并把 $\mathcal{T}$ 中的模映为零。
因此，由《同调》引理 \ref{homology-lemma-serre-subcategory-is-kernel} 给出的泛性质，
该函子下降为函子
$\text{Mod}_A/\mathcal{T} \to \text{Mod}_{S^{-1}A}$。

\medskip\noindent
反过来，任意满足 $M \otimes_A S^{-1}A = 0$ 的 $A$-模 $M$ 都是
$\mathcal{T}$ 的对象，因为
$M \otimes_A S^{-1}A \cong S^{-1} M$
（《代数》引理 \ref{algebra-lemma-tensor-localization}）。因而《同调》引理
\ref{homology-lemma-quotient-by-kernel-exact-functor} 表明函子
$\text{Mod}_A/\mathcal{T} \to \text{Mod}_{S^{-1}A}$ 是忠实的。

\medskip\noindent
此外，这个嵌入还是本质满的：$S^{-1}A$-模 $N$ 的一个原像是 $N_A$，即把
$N$ 看作 $A$-模；这是因为把 $x \otimes a/s$ 映为 $(a/s) \cdot x$ 的典范映射
$N_A \otimes_A S^{-1}A \to N$ 是 $S^{-1}A$-模的同构。
\end{proof}

\begin{proposition}
\label{examples-proposition-quotient-by-torsion-modules}
设 $A$ 为环。以 $Q(A)$ 表示其全商环（如《代数》例
\ref{algebra-example-localize-at-prime}）。以 $\text{Mod}_A$ 表示 $A$-模范畴，
以 $\mathcal{T}$ 表示其挠模 Serre 子范畴，并以 $\text{Mod}_{Q(A)}$ 表示
$Q(A)$-模范畴。则存在典范等价
$\text{Mod}_A/\mathcal{T} \rightarrow \text{Mod}_{Q(A)}$。
\end{proposition}

\begin{proof}
将命题 \ref{examples-proposition-localization-and-serre-quotients} 应用于乘法子集
$S = \{f \in A \mid f \text{ 为非零因子于 }A\}$ 即立即得到结论，
因为一个模是挠模，当且仅当它的每个元素分别被 $S$ 的某个元素零化。
\end{proof}

\begin{proposition}
\label{examples-proposition-quotient-by-finitely-generated-torsion-modules}
设 $A$ 为 Noether 整环，$K$ 为其分式域。以 $\text{Mod}_A^{fg}$ 表示有限生成
$A$-模范畴，以 $\mathcal{T}^{fg}$ 表示其有限生成挠模 Serre 子范畴。则
$\text{Mod}_A^{fg}/\mathcal{T}^{fg}$ 典范等价于有限维 $K$-向量空间范畴。
\end{proposition}

\begin{proof}
命题 \ref{examples-proposition-quotient-by-torsion-modules} 给出的等价沿嵌入
$\text{Mod}_A^{fg}/\mathcal{T}^{fg} \to \text{Mod}_A/\mathcal{T}$ 限制为等价
$\text{Mod}_A^{fg}/\mathcal{T}^{fg} \to \text{Vect}_K^{fd}$。
Noether 假设保证 $\text{Mod}_A^{fg}$ 是 Abel 范畴（见《代数进阶》第
\ref{more-algebra-section-abelian-categories-modules} 节），并保证典范函子
$\text{Mod}_A^{fg}/\mathcal{T}^{fg} \to \text{Mod}_A/\mathcal{T}$ 是满的
（否则，有限生成模的挠子模可能不是 $\mathcal{T}^{fg}$ 的对象）。
\end{proof}

\begin{proposition}
\label{examples-proposition-quotient-abelian-groups-by-torsion-groups}
Abel 群范畴对其挠群 Serre 子范畴的商范畴是 $\mathbf{Q}$-向量空间范畴。
\end{proposition}

\begin{proof}
该断言直接由命题 \ref{examples-proposition-quotient-by-torsion-modules} 得出。
\end{proof}




\section{不同的余极限拓扑}
\label{examples-section-colimit-topology}

\noindent
本例见 \cite[例 1.2，第 553 页]{TSH}。令
$G_n = \mathbf{Q} \times \mathbf{R}^n$，$n \geq 1$；把它视为加法拓扑群，并赋予
通常的（欧氏）拓扑。考虑闭嵌入 $G_n \to G_{n + 1}$，它把
$(x_0, \ldots, x_n)$ 映为 $(x_0, \ldots, x_n, 0)$。我们断言，给
$G = \colim G_n$ 赋予由下式定义的拓扑：
$$
U \subset G\text{ 为开集} \Leftrightarrow G_n \cap U\text{ 为开集 }\forall n
$$
则所得空间不是拓扑群。

\medskip\noindent
为说明这一点，考虑集合
$$
U = \{(x_0, x_1, x_2, \ldots)\text{，满足 }
|x_j| < |\cos(jx_0)| \text{ 对 } j > 0\}
$$
% P12-ERR-0407
利用 $\pi$ 不是有理数，容易从 $jx_0$ 绝不会是 $\pi/2$ 的整数倍推出
$U \cap G_n$ 是开集。由于 $0 \in U$，若上述拓扑使 $G$ 成为拓扑群，则存在
$0$ 的开邻域 $V \subset G$，使得 $V + V \subset U$。于是，对每个
$j \geq 0$，都存在 $\epsilon_j > 0$，使得当 $|x_j| < \epsilon_j$ 时有
$(0, \ldots, 0, x_j, 0, \ldots) \in V$。由于 $V + V \subset U$，当
$|x_0| < \epsilon_0$ 且 $|x_j| < \epsilon_j$ 时，必有
$$
(x_0, 0, \ldots, 0, x_j, 0, \ldots) \in U
$$
然而，取充分大的 $j$ 使 $j \epsilon_0 > \pi/2$，就可以选取
$x_0 \in \mathbf{Q}$，使得 $|\cos(jx_0)|$ 小于 $\epsilon_j$，进而存在
$x_j$ 满足 $|\cos(jx_0)| < |x_j| < \epsilon_j$。这个矛盾证明了该断言。

\begin{lemma}
\label{examples-lemma-colimit-topology}
存在（Abel）拓扑群组成的系统 $G_1 \to G_2 \to G_3 \to \ldots$，使得在
拓扑空间范畴中取的 $\colim G_n$ 不同于在拓扑群范畴中取的 $\colim G_n$。
\end{lemma}

\begin{proof}
见上述讨论。
\end{proof}







\section{普遍下沉但不是 V 覆盖}
\label{examples-section-universally-submersive-not-V}

\noindent
设 $A$ 为赋值环。设 $\mathfrak p \subset A$ 为既非极小素理想也非极大理想的
素理想。（一个便于记忆的情形是 $A$ 恰有三个素理想，而 $\mathfrak p$ 是
``中间''那个。）考虑仿射概形的态射
$$
\Spec(A_\mathfrak p) \amalg \Spec(A/\mathfrak p) \longrightarrow \Spec(A)
$$
我们断言它是普遍下沉的。为证明这一点，设
$\Spec(B) \to \Spec(A)$ 是由环映射 $A \to B$ 给出的仿射概形态射。于是必须证明
$$
\Spec(B_\mathfrak p) \amalg \Spec(B/\mathfrak pB) \to \Spec(B)
$$
是下沉的。首先，它是满射。其次，设 $T \subset \Spec(B)$ 为子集，并且
$T_1 = \Spec(B_\mathfrak p) \cap T$ 与
$T_2 = \Spec(B/\mathfrak p B) \cap T$ 都是闭集。由于 $T_1$ 和 $T_2$
都是 $B$-代数的谱，可知 $T$ 是某个 $B$-代数的谱的像。因此，要证明 $T$
为闭集，只须证明 $T$ 在专化下稳定；见《代数》引理
\ref{algebra-lemma-image-stable-specialization-closed}。为此，设
$p \leadsto q$ 是 $\Spec(B)$ 中点的一个专化，且 $p \in T$。设 $A'$ 为赋值环，
并设 $\Spec(A') \to \Spec(B)$ 为一个态射，使一般点 $\eta$（属于
$\Spec(A')$）映到 $p$、闭点 $s$（属于 $\Spec(A')$）映到 $q$；见《概形》引理
\ref{schemes-lemma-points-specialize}。注意，复合
$\gamma : \Spec(A') \to \Spec(A)$ 的像恰为点 $\xi \in \Spec(A)$ 之集，
这些点满足 $\gamma(\eta) \leadsto \xi \leadsto \gamma(s)$（略去细节）。
若 $\mathfrak p \not \in \Im(\gamma)$，则或者
$p, q \in \Spec(B_\mathfrak p)$，或者
$p, q \in \Spec(B/\mathfrak pB)$。此时，由 $T_1$（相应地，$T_2$）为闭集，
可得 $q \in T_1$（相应地，$q \in T_2$），因而 $q \in T$。
最后，设 $\mathfrak p \in \Im(\gamma)$，比如
$\mathfrak p = \gamma(r)$。则有专化 $p \leadsto r$ 和 $r \leadsto q$。
在这种情形下，$p, r \in \Spec(B_\mathfrak p)$，并且
$r, q \in \Spec(B/\mathfrak pB)$。于是，
% P12-ERR-0408
先得 $r \in T_1 \subset T$；再由 $r$ 映到 $\mathfrak p$ 得
$r \in T_2$；最后得到所需的 $q \in T_2 \subset T$。

\medskip\noindent
另一方面，我们断言单元素族
$$
\{\Spec(A_\mathfrak p) \amalg \Spec(A/\mathfrak p) \longrightarrow \Spec(A)\}
$$
不是 V 覆盖。见《拓扑》定义 \ref{topologies-definition-V-covering}。事实上，
% P12-ERR-0409
若它是 V 覆盖，就会有赋值环扩张 $A \subset B$，使得
$\Spec(B) \to \Spec(A)$ 分解经过
$\Spec(A_\mathfrak p) \amalg \Spec(A/\mathfrak p)$。
% P12-ERR-0410
这将推出 $\Spec(A')$ 不连通，荒谬。

\begin{lemma}
\label{examples-lemma-universally-submersive-not-V}
存在仿射概形的态射 $X \to Y$，它是普遍下沉的，但 $\{X \to Y\}$
不是 V 覆盖。
\end{lemma}

\begin{proof}
见上述讨论。
\end{proof}









\section{整数谱不是拟紧的}
\label{examples-section-canonical}

\noindent
当然，本节标题并不是把整数谱看作拓扑空间而言的，因为任意谱作为拓扑空间都是
拟紧的（《代数》引理 \ref{algebra-lemma-quasi-compact}）。不，这里所指的是概形范畴
上的典范拓扑中的整数谱，以及一个位点中拟紧对象的定义
（《位点》定义 \ref{sites-definition-quasi-compact}）。

\medskip\noindent
设 $U$ 为素数集 $P$ 上的一个非主超滤子。对一个子集 $T \subset P$，以
$T^c = P \setminus T$ 表示其补集。对 $A \in U$，令
$S_A \subset \mathbf{Z}$ 为由 $p \in A$ 生成的乘法子集。置
$$
\mathbf{Z}_A = S_A^{-1}\mathbf{Z}
$$
若把 $P$ 看作 $\Spec(\mathbf{Z})$ 的闭点集，则注意到
$\Spec(\mathbf{Z}_A) = \{(0)\} \cup A^c \subset \Spec(\mathbf{Z})$。
若 $A, B \in U$，则 $A \cap B \in U$ 且 $A \cup B \in U$，并且有
$$
\mathbf{Z}_{A \cap B} =
\mathbf{Z}_A \times_{\mathbf{Z}_{A \cup B}} \mathbf{Z}_B
$$
（环的纤维积）。特别地，对任意整数 $n$ 和元素
$A_1, \ldots, A_n \in U$，态射
$$
\Spec(\mathbf{Z}_{A_1}) \amalg \ldots \amalg \Spec(\mathbf{Z}_{A_n})
\longrightarrow \Spec(\mathbf{Z})
$$
% P12-ERR-0411
都分解经过某个 $p$ 的 $\Spec(\mathbf{Z}[1/p])$
（事实上，任取 $p \in A_1 \cap \ldots \cap A_n$ 即可）。因此，平坦态射族
$\{\Spec(\mathbf{Z}_A) \to \Spec(\mathbf{Z})\}_{A \in U}$ 联合满射，
但其任何有限子族都不是联合满射。

\medskip\noindent
对一个 $\mathbf{Z}$-模 $M$，置
$$
M_A = S_A^{-1}M = M \otimes_{\mathbf{Z}} \mathbf{Z}_A
$$
断言 I：对每个 $\mathbf{Z}$-模 $M$，都有
$$
M =
\text{等化子}\left(
\xymatrix{
\prod\nolimits_{A \in U} M_A \ar@<1ex>[r] \ar@<-1ex>[r] &
\prod\nolimits_{A, B \in U} M_{A \cup B}
}
\right)
$$
首先，设 $M$ 无挠。则对所有 $A \in U$，都有 $M_A \subset M_P$。
于是只须证明
$$
M = \bigcap\nolimits_{A \in U} M_A\text{，交取于 }M_P = M \otimes \mathbf{Q}
$$
事实上，因为 $U$ 非主，对任意素数 $p$ 都有 $\{p\}^c \in U$。此外，
$M_{\{p\}^c} = M_{(p)}$ 等于在素理想 $(p)$ 处的局部化。
% P12-ERR-0412
所以，上式是显然的，因为已经有 $M_{(2)} \cap M_{(3)} = M$。
其次，设 $M$ 为挠模。则有
$$
M = \bigoplus\nolimits_{p \in P} M[p^\infty]
$$
相应地有
$$
M_A = \bigoplus\nolimits_{p \not \in A} M[p^\infty]
$$
因为我们正对 $A$ 中的素数作局部化。设
$(x_A) \in \prod M_A$ 属于该等化子。记
$x_p = x_{\{p\}^c} \in M[p^{\infty}]$。等化子性质给出
$$
x_A = (x_p)_{p \not \in A}
$$
特别地，它说明除有限多个 $p \not \in A$ 外，$x_p$ 都为零。为完成挠模情形的
证明，只须说明除有限多个素数 $p$ 外，$x_p$ 都为零。若非如此，把
$\{p \in P \mid x_p \not = 0\} = T \amalg T'$ 写成两个无限集合的不交并。
由于 $U$ 是超滤子，或者 $T \not \in U$，或者 $T' \not \in U$
（事实上，若 $T, T'$ 都属于 $U$，则 $U$ 包含
$T \cap T' = \emptyset$，这是不允许的）。不妨设 $T \not \in U$。于是
$T = A^c$，这与上面提到的有限性矛盾。
最后，设 $M$ 为一般的模。考察短正合列
$$
0 \to M_{tors} \to M \to M/M_{tors} \to 0
$$
以及下面的大交换图
$$
\xymatrix{
M_{tors} \ar[r] \ar[d] &
\prod\nolimits_{A \in U} M_{tors, A} \ar@<1ex>[r] \ar@<-1ex>[r] \ar[d] &
\prod\nolimits_{A, B \in U} M_{tors, A \cup B} \ar[d] \\
M \ar[r] \ar[d] &
\prod\nolimits_{A \in U} M_A \ar@<1ex>[r] \ar@<-1ex>[r] \ar[d] &
\prod\nolimits_{A, B \in U} M_{A \cup B} \ar[d] \\
M/M_{tors} \ar[r] &
\prod\nolimits_{A \in U} (M/M_{tors})_A \ar@<1ex>[r] \ar@<-1ex>[r] &
\prod\nolimits_{A, B \in U} (M/M_{tors})_{A \cup B} \\
}
$$
% P12-ERR-0413
利用各列的正合性，以及对挠模 $M_{tors}$ 和无挠模 $M/M_{tors}$ 已得的结论
作图追逐，就为 $M$ 证明了断言 I。这给出了下述引理中现象的一个例子。

\begin{lemma}
\label{examples-lemma-non-fpqc-descent}
存在环 $A$ 和平坦环映射的无限族 $\{A \to A_i\}_{i \in I}$，使得对每个
$A$-模 $M$，都有
$$
M =
\text{等化子}\left(
\xymatrix{
\prod\nolimits_{i \in I} M \otimes_A A_i \ar@<1ex>[r] \ar@<-1ex>[r] &
\prod\nolimits_{i, j \in I} M \otimes_A A_i \otimes_A A_j
}
\right)
$$
但不存在具有同一性质的有限子族。
\end{lemma}

\begin{proof}
见上述讨论。
\end{proof}

\noindent
继续使用素数集 $P$ 上的非主超滤子 $U$。设 $R$ 为环。对 $A \in U$，记
$R_A = S_A^{-1}R = R \otimes \mathbf{Z}_A$。
断言 II：给定闭子集 $T_A \subset \Spec(R_A)$（$A \in U$），并满足
$$
(\Spec(R_{A \cup B}) \to \Spec(R_A))^{-1}T_A =
(\Spec(R_{A \cup B}) \to \Spec(R_B))^{-1}T_B
$$
对所有 $A, B \in U$ 都成立，则存在闭子集 $T \subset \Spec(R)$，使得对所有
$A \in U$，都有 $T_A = (\Spec(R_A) \to \Spec(R))^{-1}(T)$。
对 $A \in U$，设 $I_A \subset R_A$ 为定义 $T_A$ 的根理想。则粘合条件蕴含
$S_{A \cup B}^{-1}I_A = S_{A \cup B}^{-1}I_B$
在 $R_{A \cup B}$ 中对所有 $A, B \in U$ 成立（因为局部化保持根理想这一性质）。
设 $I' \subset R$ 为映入 $I_P \subset R_P = R \otimes \mathbf{Q}$ 的元素之集。
于是对 $A \in U$ 可见：
\begin{enumerate}
\item $I_A \subset I'_A = S_A^{-1}I'$；并且
\item $M_A = I'_A/I_A$ 是挠模。
\end{enumerate}
当然，对 $A, B \in U$，得到典范等同
$S_{A \cup B}^{-1}M_A = S_{A \cup B}^{-1}M_B$。把挠模 $M_A$ 分解为它们的
$p$-准素分量后，读者容易证明：存在 $p$-幂挠 $R$-模 $M_p$，使得
$$
M_A = \bigoplus\nolimits_{p \not \in A} M_p
$$
并与上述典范等同相容。置 $M = \bigoplus_{p \in P} M_p$，便得到典范同构
$M_A = S_A^{-1}M$，且它们与上述典范等同相容。于是得到 $R$-模的典范映射
$$
I' \longrightarrow M
$$
% P12-ERR-0414
它对所有 $A \in U$ 都恢复映射 $I_A \to M_A$。这是上述所有相容性以及先前所证
断言的结果；该断言说 $M$ 是从 $\prod_{A \in U} M_A$ 到
$\prod_{A, B \in U} M_{A \cup B}$ 的两个映射的等化子。
令 $I = \Ker(I' \to M)$。则 $I$ 是理想，并且 $T = V(I)$ 是一个闭子集；
它对所有 $A \in U$ 都恢复闭子集 $T_A$。这证明了断言 II。

\begin{lemma}
\label{examples-lemma-Z-not-quasi-compact}
概形 $\Spec(\mathbf{Z})$ 在概形范畴的典范拓扑中不是拟紧的。
\end{lemma}

\begin{proof}
采用上述记号，考虑态射族
$$
\mathcal{W} = \{\Spec(\mathbf{Z}_A) \to \Spec(\mathbf{Z})\}_{A \in U}
$$
由《下降》引理 \ref{descent-lemma-universal-effective-epimorphism} 和上面证明的两个
断言，这是一个普遍有效满态射。在任意有纤维积的范畴中，普遍有效满态射给
$\mathcal{C}$ 赋予位点结构（模去一些容易解决的集合论问题），从而定义典范拓扑。
因此，$\mathcal{W}$ 是典范拓扑的一个覆盖。另一方面，我们已在上面看到，
任何有限子族
$$
\{\Spec(\mathbf{Z}_{A_i}) \to \Spec(\mathbf{Z})\}_{i = 1, \ldots, n},\quad
n \in \mathbf{N}, A_1, \ldots, A_n \in U
$$
都分解经过某个 $p$ 的 $\Spec(\mathbf{Z}[1/p])$。因此，这个有限族不可能是
普遍有效满态射；更一般地，任何普遍有效满态射
$\{g_j : T_j \to \Spec(\mathbf{Z})\}$ 都不能加细
$\{\Spec(\mathbf{Z}_{A_i}) \to \Spec(\mathbf{Z})\}_{i = 1, \ldots, n}$。
由《位点》定义 \ref{sites-definition-quasi-compact}，这说明
$\Spec(\mathbf{Z})$ 在典范拓扑中不是拟紧的。要看这里的拟紧概念与通常的
拓扑斯论定义一致，见《位点》引理 \ref{sites-lemma-quasi-compact}。
\end{proof}
% Generated mechanically from the frozen exercises.tex target.
% Do not edit this generated file; edit the bound locale source instead.

% OK, start here.
%

\setcounter{chapter}{110}
\chapter{练习}



\phantomsection
\label{exercises-section-phantom}



\section{代数}
\label{exercises-section-algebra}

\noindent
本节首先收录若干杂题。

\begin{exercise}
\label{exercises-exercise-isomorphism-localizations}
设 $A$ 为环，${\mathfrak m}$ 为极大理想。在 $A[X]$ 中令
$\tilde {\mathfrak m}_1 = ({\mathfrak m}, X)$ 且
$\tilde {\mathfrak m}_2 = ({\mathfrak m}, X-1)$。证明
$$
A[X]_{\tilde {\mathfrak m}_1} \cong A[X]_{\tilde {\mathfrak m}_2}.
$$
\end{exercise}

\begin{exercise}
\label{exercises-exercise-coherent}
% P12-ERR-0415
求一个非 Noether 环 $R$ 的例子，使得 $R$ 的每个有限生成理想作为
$R$-模都是有限表示的。（若一个环具有后一性质，则称它是{\it 凝聚环}。）
\end{exercise}

\begin{exercise}
\label{exercises-exercise-flat-ideals-pid}
设 $(A, {\mathfrak m}, k)$ 为 Noether 局部环。对任意有限 $A$-模 $M$，
定义 $r(M)$ 为 $M$ 作为 $A$-模的最少生成元个数。由 NAK，此数等于
$\dim_k M/{\mathfrak m} M = \dim_k M \otimes_A k$。
\begin{enumerate}
\item 证明 $r(M \otimes_A N) = r(M)r(N)$。
\item 设 $I\subset A $ 为满足 $r(I) > 1$ 的理想。证明
$r(I^2) < r(I)^2$。
\item 推出：若 $A$ 中每个理想都是平坦模，则 $A$ 是主理想整环
（或域）。
\end{enumerate}
\end{exercise}

\begin{exercise}
\label{exercises-exercise-not-isomorphic}
设 $k$ 为域。证明下列每一对 $k$-代数不同构：
\begin{enumerate}
\item $k[x_1, \ldots, x_n]$ 与 $k[x_1, \ldots, x_{n + 1}]$，
其中 $n\geq 1$ 任意；
\item $k[a, b, c, d, e, f]/(ab + cd + ef)$ 与
$k[x_1, \ldots, x_n]$，其中 $n = 5$；
\item $k[a, b, c, d, e, f]/(ab + cd + ef)$ 与
$k[x_1, \ldots, x_n]$，其中 $n = 6$。
\end{enumerate}
\end{exercise}

\begin{remark}
\label{exercises-remark-simple-geometric}
当然，本练习的意图是在每种情形中找到简单论证，而不是应用某个“大”定理。
不过，以一般原理为指引仍然很有益。
\end{remark}

\begin{exercise}
\label{exercises-exercise-silly}
代数。（这是个很初等且应当容易的问题。）
\begin{enumerate}
\item 给出一个环 $A$ 及一个不分裂的 $A$-模短正合列的例子
$$
0 \to M_1 \to M_2 \to M_3 \to 0.
$$
\item 给出一个如上的不分裂 $A$-模正合列，以及一个忠实平坦映射
$A \to B$，使得
$$
0 \to M_1\otimes_A B \to M_2\otimes_A B \to M_3\otimes_A B \to 0.
$$
作为 $B$-模正合列是分裂的。
\end{enumerate}
\end{exercise}

\begin{exercise}
\label{exercises-exercise-field-kummer}
设域 $k$ 含有一个本原 $n$ 次单位根 $\zeta$。这意味着
$\zeta^n = 1$，但 $\zeta^m\not = 1$ 对 $0 < m < n$ 成立。
\begin{enumerate}
\item 证明 $k$ 的特征与 $n$ 互素。
\item 设 $a \in k$ 是 $k$ 中一个元素。
% P12-ERR-0416
若 $d$ 是 $n$ 的任意满足 $n \geq d > 1$ 的因子，则该元素在 $k$ 中
不是 $d$ 次幂。证明 $k[x]/(x^n-a)$ 是域。（提示：考虑
$x^n-a$ 的一个分裂域，并使用 Galois 理论。）
\end{enumerate}
\end{exercise}

\begin{exercise}
\label{exercises-exercise-valuation}
设 $\nu : k[x]\setminus \{0\}  \to {\mathbf Z}$ 为满足下列性质的映射：
只要 $f$、$g$ 非零，就有 $\nu(fg) = \nu(f) + \nu(g)$；只要
$f$、$g$、$f + g$ 非零，就有
% P12-ERR-0417
$\nu(f + g) \geq min(\nu(f), \nu(g))$；并且
$\nu(c) = 0$ 对所有 $c\in k^*$ 都成立。
\begin{enumerate}
\item 证明：若 $f$、$g$ 和 $f + g$ 均非零，且
$\nu(f) \not = \nu(g)$，则
$\nu(f + g) = min(\nu(f), \nu(g))$。
\item 证明：若 $f = \sum a_i x^i$ 且 $f\not = 0$，则
$\nu(f) \geq min(\{i\nu(x)\}_{a_i\not = 0})$。何时等号成立？
\item 证明：若 $\nu$ 取到负值，则
$\nu(f) = -n \deg(f)$ 对某个 $n\in {\mathbf N}$ 成立。
\item 设 $\nu(x) \geq 0$。证明
$\{f \mid f = 0, \ or\ \nu(f) > 0\}$ 是 $k[x]$ 的素理想。
\item 描述所有可能的 $\nu$。
\end{enumerate}
\end{exercise}


\noindent
设 $A$ 为环。若元素 $e \in A$ 满足 $e^2 = e$，则称它为{\it 幂等元}。
元素 $1$ 和 $0$ 总是幂等元。
% P12-ERR-0418
{\it 非平凡幂等元}是一个不等于零的幂等元。若两个幂等元
$e, e' \in A$ 满足 $ee' = 0$，则称它们{\it 正交}。

\begin{exercise}
\label{exercises-exercise-product}
设 $A$ 为环。证明：$A$ 是两个非零环的直积，当且仅当 $A$
具有非平凡幂等元。
\end{exercise}

\begin{exercise}
\label{exercises-exercise-lift-idempotents}
设 $A$ 为环，$I \subset A$ 为局部幂零理想。证明映射
$A \to A/I$ 在幂等元上诱导双射。（提示：先证明 $I$ 为幂零理想时的结论，
这可能更容易。然后用“绝对 Noether 约化”化归到幂零情形。）
\end{exercise}







\section{余极限}
\label{exercises-section-colimits}


\begin{definition}
\label{exercises-definition-directed-poset}
{\it 有向集}是一个非空集合 $I$，其上赋有预序 $\leq$，并且对任意一对
$i, j \in I$，存在 $k \in I$ 使得 $i \leq k$ 且 $j \leq k$。
$I$ 上的一个{\it 环系}由每个 $i \in I$ 所对应的环 $A_i$，以及每当
$i \leq j$ 时的环映射 $\varphi_{ij} : A_i \to A_j$ 组成，并要求：
只要 $i \leq j \leq k$，复合 $A_i \to A_j \to A_k$ 就等于
$A_i \to A_k$。
\end{definition}

\noindent
群系、固定环上的模系、域上的向量空间系等可类似定义。

\begin{exercise}
\label{exercises-exercise-directed-colimit}
设 $I$ 为有向集，$(A_i, \varphi_{ij})$ 为 $I$ 上的环系。
证明存在环 $A$ 及映射 $\varphi_i : A_i \to A$，使得对所有
$i \leq j$ 都有 $\varphi_j \circ \varphi_{ij} = \varphi_i$，并具有
下列泛性质：给定任意环 $B$ 及映射 $\psi_i : A_i \to B$，若对所有
$i \leq j$ 都有 $\psi_j \circ \varphi_{ij} = \psi_i$，则存在唯一的
环映射 $\psi : A \to B$ 使得 $\psi_i = \psi \circ \varphi_i$。
\end{exercise}

\begin{definition}
\label{exercises-definition-colimit}
练习 \ref{exercises-exercise-directed-colimit} 中构造的环 $A$ 称为该系统的
{\it 余极限}，记作 $\colim A_i$。
\end{definition}

\begin{exercise}
\label{exercises-exercise-prime-in-colimit}
% P12-ERR-0419
设 $(I, \geq)$ 为有向集，$(A_i, \varphi_{ij})$ 为 $I$ 上的环系，
其余极限为 $A$。证明存在双射
$$
\Spec(A) = \{(\mathfrak p_i)_{i \in I} \mid
\mathfrak p_i \subset A_i \text{ 且 }
\mathfrak p_i = \varphi_{ij}^{-1}(\mathfrak p_j)\ \forall i \leq j\}
\subset \prod\nolimits_{i \in I} \Spec(A_i)
$$
等式右侧的集合是各集合 $\Spec(A_i)$ 的极限，记作
$\lim \Spec(A_i)$。
\end{exercise}

\begin{exercise}
\label{exercises-exercise-colimit-surjective}
设 $(I, \geq)$ 为有向集，$(A_i, \varphi_{ij})$ 为 $I$ 上的环系，
其余极限为 $A$。设对所有 $i \leq j$，映射
$\Spec(A_j) \to \Spec(A_i)$ 都是满射。证明对所有 $i$，
$\Spec(A) \to \Spec(A_i)$ 都是满射。（提示：可以尝试使用
Tychonoff 定理；不过，由《代数》引理
\ref{algebra-lemma-in-image} 也可给出一个基本上平凡的直接代数证明。）
\end{exercise}

\begin{exercise}
\label{exercises-exercise-integral-colimit-finite}
设 $A \subset B$ 为整环扩张。证明
$\Spec(B) \to \Spec(A)$ 是满射。
% P12-ERR-0420
使用上面的练习、课堂上已证明的有限环扩张情形，并证明
$B = \colim B_i$ 是有限扩张 $A \subset B_i$ 的有向余极限。
\end{exercise}

\begin{exercise}
\label{exercises-exercise-colimit-tensor}
设 $(I, \geq)$ 为有向集。设 $A$ 为环，
$(N_i, \varphi_{i, i'})$ 为由 $I$ 标号的 $A$-模有向系，并设
$M$ 为另一个 $A$-模。证明
$$
\colim_{i\in I} M \otimes_A N_i\cong
M \otimes_A \Big( \colim_{i\in I} N_i\Big).
$$
\end{exercise}

\begin{definition}
\label{exercises-definition-finite-presentation}
若 $R$-模 $M$ 同构于有限秩自由模之间某个映射
$ R^{\oplus n} \to R^{\oplus m}$ 的余核，则称它在 $R$ 上
{\it 有限表示}。
\end{definition}

\begin{exercise}
\label{exercises-exercise-colimit-modules}
证明任意环上的任意模：
\begin{enumerate}
\item 是其有限生成子模的余极限；
\item 可按某种方式写成有限表示模的余极限。
\end{enumerate}
\end{exercise}





\section{加性范畴与 Abel 范畴}
\label{exercises-section-additive}

\begin{exercise}
\label{exercises-exercise-filtered-vector-spaces}
设 $k$ 为域，$\mathcal{C}$ 为 $k$ 上滤过向量空间所成的范畴。
任意范畴中滤过对象的定义见《同调代数》定义
\ref{homology-definition-filtered}。
\begin{enumerate}
\item 证明这是一个加性范畴（并仔细说明两个对象的直和是什么）。
\item 设 $f : (V, F) \to (W, F)$ 为 $\mathcal{C}$ 的一个态射。
证明 $f$ 有核与余核（并准确说明 $f$ 的核与余核是什么）。
\item 给出 $\mathcal{C}$ 中一个映射的例子，使得典范映射
$\Coim(f) \to \Im(f)$ 不是同构。
\end{enumerate}
\end{exercise}

\begin{exercise}
\label{exercises-exercise-torsion-free}
设 $R$ 为 Noether 整环，$\mathcal{C}$ 为有限生成无挠 $R$-模
所成的范畴。
\begin{enumerate}
\item 证明这是一个加性范畴。
\item 设 $f : N \to M$ 为 $\mathcal{C}$ 的一个态射。证明 $f$
有核与余核（务必准确地定义 $f$ 的核与余核）。
\item 给出一个 Noether 整环 $R$ 以及 $\mathcal{C}$ 中的一个映射，
使得典范映射 $\Coim(f) \to \Im(f)$ 不是同构。
\end{enumerate}
\end{exercise}

\begin{exercise}
\label{exercises-exercise-other}
给出一个加性且具有核与余核的范畴的例子，但它不同于练习
\ref{exercises-exercise-filtered-vector-spaces} 和
\ref{exercises-exercise-torsion-free} 中的范畴。
\end{exercise}





\section{张量积}
\label{exercises-section-tensor-product}

\noindent
张量积在《代数》第 \ref{algebra-section-tensor-product} 节中引入。
设 $R$ 为环，$\text{Mod}_R$ 为 $R$-模范畴。我们称函子
$F : \text{Mod}_R \to \text{Mod}_R$
\begin{enumerate}
\item 是加性的，如果对任意 $R$-模 $M, N$，映射
$F : \Hom_R(M, N) \to \Hom_R(F(M), F(N))$
都是 Abel 群同态；见《同调代数》定义
\ref{homology-definition-preadditive}；
% P12-ERR-0421
\item 是 $R$-线性的，如果对任意 $R$-模 $M, N$，映射
$F : \Hom_R(M, N) \to \Hom_R(F(M), F(N))$ 都是 $R$-线性的；
\item 是右正合的，如果对任意短正合列
$0 \to M_1 \to M_2 \to M_3 \to 0$，序列
$F(M_1) \to F(M_2) \to F(M_3) \to 0$ 都正合；
\item 是左正合的，如果对任意短正合列
$0 \to M_1 \to M_2 \to M_3 \to 0$，序列
$0 \to F(M_1) \to F(M_2) \to F(M_3)$ 都正合；
\item 与直和交换，如果给定集合 $I$ 及 $R$-模 $M_i$，各映射
$F(M_i) \to F(\bigoplus M_i)$ 诱导同构
$\bigoplus F(M_i) = F(\bigoplus M_i)$。
\end{enumerate}

\begin{exercise}
\label{exercises-exercise-characterize-tensor-functor}
设 $R$ 为环，沿用上述记号。
\begin{enumerate}
\item 给出环 $R$ 及加性函子
$F : \text{Mod}_R \to \text{Mod}_R$ 的例子，使得该函子不是
$R$-线性的。
\item 设 $N$ 为 $R$-模。证明函子 $F(M) = M \otimes_R N$
是 $R$-线性的、右正合的，并且与直和交换。
\item 反过来，证明任意函子 $F : \text{Mod}_R \to \text{Mod}_R$，
若它是 $R$-线性的、右正合的且与直和交换，则它具有形式
$F(M) = M \otimes_R N$，其中 $N$ 是某个 $R$-模。
\item 证明：若在 (3) 中去掉 $F$ 与直和交换的假设，则结论不再成立。
\end{enumerate}
\end{exercise}






\section{平坦环映射}
\label{exercises-section-flat}

\begin{exercise}
\label{exercises-exercise-localization-flat}
设 $S$ 为环 $A$ 的乘法子集。
\begin{enumerate}
\item 对 $A$-模 $M$，证明 $S^{-1}M = S^{-1}A \otimes_A M$。
\item 证明 $S^{-1}A$ 在 $A$ 上平坦。
\end{enumerate}
\end{exercise}

\begin{exercise}
\label{exercises-exercise-examples-not-flat}
在下列各情形中，求一个 $A$-模单射 $M_1 \to M_2$，使得
$M_1\otimes N \to M_2 \otimes N$ 不是单射：
\begin{enumerate}
\item $A = k[x, y]$ 且 $N = (x, y) \subset A$。（这里及以下，
$k$ 都是域。）
\item $A = k[x, y]$ 且 $N = A/(x, y)$。
\end{enumerate}
\end{exercise}

\begin{exercise}
\label{exercises-exercise-flat-not-projective}
% P12-ERR-0422
给出环 $A$ 及有限 $A$-模 $M$ 的例子，使得它是平坦的但不是
投射 $A$-模。
\end{exercise}

\begin{remark}
\label{exercises-remark-flat-not-projective}
若 $M$ 在 $A$ 上有限表示且平坦，则 $M$ 在 $A$ 上投射。
因此，你的例子必须涉及一个非 Noether 环 $A$。我知道一个例子，
其中 $A$ 是 ${\mathcal C}^\infty$-函数在 ${\mathbf R}$ 上的环。
\end{remark}

\begin{exercise}
\label{exercises-exercise-flat-not-free-dvr}
求 ${\mathbf Z}_{(2)}$ 上一个平坦但不自由的模。
\end{exercise}

\begin{exercise}
\label{exercises-exercise-flat-deformations}
平坦形变。
\begin{enumerate}
\item 设 $k$ 为域，$k[\epsilon]$ 为对偶数环
$k[\epsilon] = k[x]/(x^2)$，且 $\epsilon = \bar x$。证明：对任意
$k$-代数 $A$，存在平坦 $k[\epsilon]$-代数 $B$，使得
$A$ 同构于 $B/\epsilon B$。
\item 设 $k = {\mathbf F}_p = {\mathbf Z}/p{\mathbf Z}$ 且
$$
A = k[x_1, x_2, x_3, x_4, x_5, x_6]/
(x_1^p, x_2^p, x_3^p, x_4^p, x_5^p, x_6^p).
$$
证明存在平坦 ${\mathbf Z}/p^2{\mathbf Z}$-代数 $B$，使得
$B/pB$ 同构于 $A$。（因此，这里 $p$ 扮演 $\epsilon$ 的角色。）
\item 现在令 $p = 2$，并对
$k = {\mathbf F}_2 = {\mathbf Z}/2{\mathbf Z}$ 及
$$
A = k[x_1, x_2, x_3, x_4, x_5, x_6]/
(x_1^2, x_2^2, x_3^2, x_4^2, x_5^2, x_6^2, x_1x_2 + x_3x_4 + x_5x_6).
$$
考虑同一问题。不过，在这种情形下证明：{\it 不}存在平坦
${\mathbf Z}/4{\mathbf Z}$-代数 $B$ 使得 $B/2B$ 同构于 $A$。
（找出其中的窍门！同一例子对任意特征 $p > 0$ 都成立，只是计算更困难。）
\end{enumerate}
\end{exercise}

\begin{exercise}
\label{exercises-exercise-flat-given-residue-field-extension}
设 $(A, {\mathfrak m}, k)$ 为局部环，$k'/k$ 为有限域扩张。
证明存在局部环之间的平坦局部映射 $A \to B$，使得
${\mathfrak m}_B = {\mathfrak m} B$，并且
$B/{\mathfrak m} B$ 作为 $k$-代数同构于 $k'$。（提示：先处理
$k \subset k'$ 由单个元素生成的情形。）
\end{exercise}

\begin{remark}
\label{exercises-remark-flat-given-residue-field-extension-general}
同一结果对任意域扩张 $K/k$ 都成立。
\end{remark}

\section{环的谱}
\label{exercises-section-spectrum-ring}

\begin{exercise}
\label{exercises-exercise-spec-Z}
计算集合 $\Spec(\mathbf{Z})$ 并描述其拓扑。
\end{exercise}

\begin{exercise}
\label{exercises-exercise-basis-opens-standard}
设 $A$ 为任意环。对 $f\in A$，定义
$D(f):= \{\mathfrak p \subset A \mid f \not \in \mathfrak p\}$。
证明开子集 $D(f)$ 构成 $\Spec(A)$ 拓扑的一组基。
\end{exercise}

\begin{exercise}
\label{exercises-exercise-radical-ideals-closed}
证明映射 $I\mapsto V(I)$ 定义一个自然双射
$$
\{I\subset A\text{ 满足 }I = \sqrt{I}\}
\longrightarrow
\{T\subset \Spec(A)\text{ 为闭集}\}
$$
\end{exercise}

\begin{definition}
\label{exercises-definition-quasi-compact}
若对任意开覆盖 $X = \bigcup_{i\in I} U_i$，都存在有限子集
$\{i_1, \ldots, i_n\}\subset I$ 使得
$X = U_{i_1}\cup\ldots U_{i_n}$，则称拓扑空间 $X$ {\it 拟紧}。
\end{definition}

\begin{exercise}
\label{exercises-exercise-spec-quasi-compact}
证明对任意环 $A$，$\Spec(A)$ 都拟紧。
\end{exercise}

\begin{definition}
\label{exercises-definition-Hausdorff}
若拓扑空间 $X$ 中任意两个不同的点 $x, y\in X$，$x\not = y$，
都有一个开子集仅含其中一点而不含另一点，则称 $X$ 满足分离公理
$T_0$。若对任意一对 $x, y\in X$，$x\not = y$，都存在不交的
开子集 $U, V$，使得 $x\in U$ 且 $y\in V$，则称 $X$
{\it 豪斯多夫}。
\end{definition}

\begin{exercise}
\label{exercises-exercise-not-hausdorff}
证明一般而言 $\Spec(A)$ {\bf 不}是豪斯多夫的。
证明 $\Spec(A)$ 满足 $T_0$。给出一个不满足 $T_0$ 的拓扑空间
$X$ 的例子。
\end{exercise}

\begin{remark}
\label{exercises-remark-not-hausdorff}
“紧”一词通常保留给既拟紧又豪斯多夫的空间。
\end{remark}

\begin{definition}
\label{exercises-definition-irreducible}
若拓扑空间 $X$ 非空，并且每当 $X = Z_1\cup Z_2$、其中
$Z_1, Z_2\subset X$ 均为闭集时，都有 $Z_1 = X$ 或
$Z_2 = X$，则称 $X$ {\it 不可约}。拓扑空间的子集
$T\subset X$ 若在 $X$ 诱导的拓扑下是不可约拓扑空间，则称
$T$ {\it 不可约}。由此定义可知，该子集不可约，当且仅当
$T$ 在 $X$ 中的闭包 $\bar T$ 不可约。
\end{definition}

\begin{exercise}
\label{exercises-exercise-irreducible-spec}
证明 $\Spec(A)$ 不可约，当且仅当 $Nil(A)$ 是素理想；并证明在
这种情形下，它是 $A$ 的唯一极小素理想。
\end{exercise}

\begin{exercise}
\label{exercises-exercise-irreducible-prime}
证明闭子集 $T\subset \Spec(A)$ 不可约，当且仅当存在素理想
${\mathfrak p}\subset A$ 使它具有形式 $T = V({\mathfrak p})$。
\end{exercise}

\begin{definition}
\label{exercises-definition-generic-point}
设 $X$ 为不可约拓扑空间。若点 $x$ 满足 $X$ 等于单点子集
$\{x\}$ 的闭包，则称其为 $X$ 的{\it 泛点}。
\end{definition}

\begin{exercise}
\label{exercises-exercise-irreducible-T0-at-most-one-generic}
证明在 $T_0$ 空间 $X$ 中，每个不可约闭子集至多有一个泛点。
\end{exercise}

\begin{exercise}
\label{exercises-exercise-spec-sober}
证明在 $\Spec(A)$ 中，每个不可约闭子集{\it 确实}有泛点。
事实上，证明映射 ${\mathfrak p} \mapsto
\overline{\{{\mathfrak p}\}}$ 是从 $\Spec(A)$ 到 $X$ 的不可约
闭子集所成集合的双射。
% P12-ERR-0423
\end{exercise}

\begin{exercise}
\label{exercises-exercise-irreducible-subset-not-generic}
给出一个例子，说明 $\Spec(\mathbf{Z})$ 的不可约子集不一定有泛点。
\end{exercise}

\begin{definition}
\label{exercises-definition-Noetherian-space}
若 $X$ 的任意闭子集降链 $Z_1\supset Z_2 \supset Z_3\supset \ldots$
都稳定，则称拓扑空间 $X$ 是{\it Noether} 的。（若任意闭子集升链
都稳定，则称它是{\it Artin} 的。）
\end{definition}

\begin{exercise}
\label{exercises-exercise-Noetherian-spec}
证明：若环 $A$ 是 Noether 环，则拓扑空间 $\Spec(A)$ 是
Noether 空间。给出一个例子，说明逆命题不成立。（如果愿意，也可对
Artin 情形做同样的事。）
\end{exercise}

\begin{definition}
\label{exercises-definition-irreducible-component}
极大不可约子集 $T\subset X$ 称为 $X$ 的一个
{\it 不可约分支}。$X$ 的任何不可约分支自动是 $X$ 的闭子集。
\end{definition}

\begin{exercise}
\label{exercises-exercise-irreducible-in-irreducible}
证明 $X$ 的任意不可约子集都包含于 $X$ 的某个不可约分支中。
\end{exercise}

\begin{exercise}
\label{exercises-exercise-Noetherian-finite-nr-irreducible}
证明 Noether 拓扑空间 $X$ 只有有限多个不可约分支，记为
$X_1, \ldots, X_n$，且 $X = X_1\cup X_2\cup\ldots\cup X_n$。
（注意，任意 $X$ 总是其不可约分支的并；不过，例如若
$X = {\mathbf R}$ 取通常拓扑，则 $X$ 的不可约分支都是单点子集，
这并没有多大意思。）
\end{exercise}

\begin{exercise}
\label{exercises-exercise-irreducible-components-minimal-primes}
证明 $\Spec(A)$ 的不可约分支对应于 $A$ 的极小素理想。
\end{exercise}

\begin{definition}
\label{exercises-definition-closed}
若 $x\in X$ 满足 $\overline{\{x\}} = \{ x\}$，则称该点
为{\it 闭点}。设 $x, y$ 是 $X$ 的点。若
$x\in \overline{\{y\}}$，则称 $x$ 是 $y$ 的一个{\it 专化}，
或称 $y$ 是 $x$ 的一个{\it 一般化}。
\end{definition}

\begin{exercise}
\label{exercises-exercise-closed-maximal}
证明 $\Spec(A)$ 的闭点对应于 $A$ 的极大理想。
\end{exercise}

\begin{exercise}
\label{exercises-exercise-generalization}
证明 ${\mathfrak p}$ 是 $\Spec(A)$ 中 ${\mathfrak q}$ 的一般化，
当且仅当 ${\mathfrak p}\subset {\mathfrak q}$。用一般化与专化
刻画闭点、极大理想、泛点和极小素理想。（这里采用如下术语：
可约拓扑空间 $X$ 的一个点若是 $X$ 某个不可约分支的泛点，就称它为
% P12-ERR-0424
一个泛点。）
\end{exercise}

\begin{exercise}
\label{exercises-exercise-disjoint-closed-spec}
设 $I$ 和 $J$ 为 $A$ 的理想。$V(I)$ 与 $V(J)$ 不交的条件是什么？
\end{exercise}

\begin{definition}
\label{exercises-definition-connected-component}
若拓扑空间 $X$ 非空，且不能写成两个非空不交开子集的并，则称
它{\it 连通}。它的极大连通子集称为 $X$ 的一个
{\it 连通分支}。$X$ 的每个点都包含于 $X$ 的某个连通分支中，
而 $X$ 的每个连通分支在 $X$ 中都是闭的。（但一般而言，连通分支
在 $X$ 中未必是开的。）
\end{definition}

\begin{exercise}
\label{exercises-exercise-disconnected-spec}
设 $A$ 为非零环。证明 $\Spec(A)$ 不连通，当且仅当对某些非零环
$B, C$ 有 $A\cong B \times C$。
\end{exercise}

\begin{exercise}
\label{exercises-exercise-connected-component-stable-generalization}
设 $T$ 为 $\Spec(A)$ 的一个连通分支。证明 $T$ 在一般化下稳定。
若 $A$ 是 Noether 环，证明 $T$ 是 $\Spec(A)$ 的开子集。
（注：例如，当 $A$ 是无穷多个 ${\mathbf F}_2$ 的直积时，该结论错误。
此环的谱由无穷多个闭点组成。）
\end{exercise}

\begin{exercise}
\label{exercises-exercise-primes-kx}
计算 $\Spec(k[x])$，即描述该环的素理想、可能的专化以及拓扑。
（在 $k$ 为代数闭域时完成计算，也在 $k$ 不是代数闭域时完成计算。）
\end{exercise}

\begin{exercise}
\label{exercises-exercise-primes-kxy}
计算 $\Spec(k[x, y])$，其中 $k$ 是代数闭域。
[提示：使用态射 $\varphi : \Spec(k[x, y]) \to \Spec(k[x])$；若
$\varphi({\mathfrak p}) = (0)$，则关于
$S = \{f\in k[x] \mid f \not = 0\}$ 局部化，并使用课堂上关于
局部化与 $\Spec$ 的结果。]
（你认为代数几何学家为什么称它为仿射 2-空间？）
\end{exercise}

\begin{exercise}
\label{exercises-exercise-primes-Zy}
计算 $\Spec(\mathbf{Z}[y])$。
[提示：同上。]（$\mathbf{Z}$ 上的仿射 1-空间。）
\end{exercise}






\section{局部化}
\label{exercises-section-localization}


\begin{exercise}
\label{exercises-exercise-submodule-localization}
设 $A$ 为环，$S \subset A$ 为乘法子集，$M$ 为 $A$-模，
$N \subset S^{-1}M$ 为 $S^{-1}A$-子模。证明存在 $A$-子模
$N' \subset M$ 使得 $N = S^{-1}N'$。（这一有用结果特别适用于
$S^{-1}A$ 的理想。）
\end{exercise}

\begin{exercise}
\label{exercises-exercise-localize-zero}
设 $A$ 为环，$M$ 为 $A$-模，$m \in M$。
\begin{enumerate}
\item 证明 $I = \{a \in A \mid am = 0\}$ 是 $A$ 的理想。
\item 对 $A$ 的素理想 $\mathfrak p$，证明 $m$ 在
$M_\mathfrak p$ 中的像为零，当且仅当 $I \not \subset \mathfrak p$。
\item 证明 $m$ 为零，当且仅当 $m$ 在所有 $A$ 的素理想
$\mathfrak p$ 所对应的 $M_\mathfrak p$ 中的像都为零。
\item 证明 $m$ 为零，当且仅当 $m$ 在所有 $A$ 的极大理想
$\mathfrak m$ 所对应的 $M_\mathfrak m$ 中的像都为零。
\item 证明 $M = 0$，当且仅当对所有极大理想 $\mathfrak m$，
$M_{\mathfrak m}$ 都为零。
\end{enumerate}
\end{exercise}

\begin{exercise}
\label{exercises-exercise-localization-is-field}
求一对 $(A, f)$，其中 $A$ 是具有三个或更多两两不同素理想的整环，
$f \in A$ 是一个元素，并且主局部化
$A_f = \{1, f, f^2, \ldots \}^{-1}A$ 是域。
\end{exercise}

\begin{exercise}
\label{exercises-exercise-localize-finite-module-zero}
设 $A$ 为环，$M$ 为有限 $A$-模，$S \subset A$ 为乘法集。
假设 $S^{-1}M = 0$。证明存在 $f \in S$，使得主局部化
$M_f = \{1, f, f^2, \ldots \}^{-1}M$ 为零。
\end{exercise}

\begin{exercise}
\label{exercises-exercise-localization-is-quotient}
给出三元组 $(A, I, S)$ 的例子，其中 $A$ 是环，
$0 \not = I \not = A$ 是真非零理想，$S \subset A$ 是乘法子集，
并且作为 $A$-代数有 $A/I \cong S^{-1}A$。
\end{exercise}




\section{Nakayama 引理}
\label{exercises-section-nakayama}

\begin{exercise}
\label{exercises-exercise-nakayama}
设 $A$ 为环，$I$ 为 $A$ 的理想，$M$ 为 $A$-模，并设
$x_1, \ldots, x_n \in M$。假设：
\begin{enumerate}
\item $M/IM$ 由 $x_1, \ldots, x_n$ 生成；
\item $M$ 是有限 $A$-模；
\item $I$ 包含于 $A$ 的每个极大理想中。
\end{enumerate}
证明 $x_1, \ldots, x_n$ 生成 $M$。（建议解法：使用练习
\ref{exercises-exercise-localize-zero} 以及局部化的正合性，化归到 $A$
的一个极大理想处的局部化。然后考察 $M$ 模去由
$x_1, \ldots, x_n$ 生成的子模所得的商，化归到课堂上的
Nakayama 引理。）
\end{exercise}







\section{长度}
\label{exercises-section-length}

\begin{definition}
\label{exercises-definition-length}
设 $A$ 为环，$M$ 为 $A$-模。
% P12-ERR-0425
$M$ 作为 $R$-模的{\it 长度}为
$$
\text{length}_A(M)
=
\sup
\{
n
\mid
\exists\ 0 = M_0 \subset M_1 \subset \ldots \subset M_n = M,
\text{ }M_i \not = M_{i + 1}
\}.
$$
换言之，它是子模链长度的上确界。
\end{definition}

\begin{exercise}
\label{exercises-exercise-length-is-one}
证明环 $A$ 上的模 $M$ 长度为 $1$，当且仅当它同构于
$A/\mathfrak m$，其中 $\mathfrak m$ 是 $A$ 的某个极大理想。
\end{exercise}

\begin{exercise}
\label{exercises-exercise-length-easy}
计算下列各模在相应环上的长度。简要（！）说明你的答案。
（可自由使用短正合列中长度函数的可加性；见《代数》引理
\ref{algebra-lemma-length-additive}。）
\begin{enumerate}
\item $\mathbf{Z}/120\mathbf{Z}$ 在 $\mathbf{Z}$ 上的长度。
\item $\mathbf{C}[x]/(x^{100} + x + 1)$ 在 $\mathbf{C}[x]$ 上的长度。
\item $\mathbf{R}[x]/(x^4 + 2x^2 + 1)$ 在 $\mathbf{R}[x]$ 上的长度。
\end{enumerate}
\end{exercise}

\begin{exercise}
\label{exercises-exercise-compute-length}
设 $A = k[x, y]_{(x, y)}$ 为仿射平面在原点处的局部环。
可对域 $k$ 作任意你喜欢的假设。设
$f = x^3 + x^2y^2 + y^{100}$ 且 $g = y^3 - x^{999}$。
$A/(f, g)$ 作为 $A$-模的长度是多少？（一种可能的方法是：
考察 $f$ 和 $g$ 在形如
% P12-ERR-0426
$A/{\mathfrak m}_A^n= k[x, y]/(x, y)^n$ 的商中生成的理想，其中
$n$ 变化。尝试寻找 $n$，使得
$A/(f, g)+{\mathfrak m}_A^n \cong A/(f, g)+{\mathfrak m}_A^{n + 1}$，
并使用 NAK。）
\end{exercise}






\section{伴随素理想}
\label{exercises-section-ass}

\noindent
伴随素理想在《代数》第 \ref{algebra-section-ass} 节中讨论。

\begin{exercise}
\label{exercises-exercise-compute-ass}
计算下列各模的伴随素理想集合。
\begin{enumerate}
\item $R = k[x, y]$ 且 $M = R/(xy(x + y))$；
\item $R = \mathbf{Z}[x]$ 且 $M = R/(300x + 75)$；
\item $R = k[x, y, z]$ 且 $M = R/(x^3, x^2y, xz)$。
\end{enumerate}
这里照例 $k$ 是域。
\end{exercise}

\begin{exercise}
\label{exercises-exercise-prime-power-not-primary}
给出 Noether 环 $R$ 及素理想 $\mathfrak p$ 的例子，使得
$\mathfrak p$ 不是 $R/\mathfrak p^2$ 的唯一伴随素理想。
\end{exercise}

\begin{exercise}
\label{exercises-exercise-product-primes-not-only-associated}
设 $R$ 为 Noether 环，具有不可比的素理想
$\mathfrak p$、$\mathfrak q$；即
$\mathfrak p \not \subset \mathfrak q$ 且
$\mathfrak q \not \subset \mathfrak p$。
\begin{enumerate}
\item 证明对 $N = R/(\mathfrak p \cap \mathfrak q)$ 有
$\text{Ass}(N) = \{\mathfrak p, \mathfrak q\}$。
\item 用例子说明，模 $M = R/\mathfrak p \mathfrak q$ 可以有一个
既不等于 $\mathfrak p$ 也不等于 $\mathfrak q$ 的伴随素理想。
\end{enumerate}
\end{exercise}





\section{Ext 群}
\label{exercises-section-ext}

\noindent
Ext 群在《代数》第 \ref{algebra-section-ext} 节中定义。

\begin{exercise}
\label{exercises-exercise-compute-ext-abelian-groups}
计算所给各模的全部 Ext 群 $\Ext^i(M, N)$；这里取
$\mathbf{Z}$-模范畴（亦称 Abel 群范畴）。
\begin{enumerate}
\item $M = \mathbf{Z}$ 且 $N = \mathbf{Z}$；
\item $M = \mathbf{Z}/4\mathbf{Z}$ 且 $N = \mathbf{Z}/8\mathbf{Z}$；
\item $M = \mathbf{Q}$ 且 $N = \mathbf{Z}/2\mathbf{Z}$；
\item $M = \mathbf{Z}/2\mathbf{Z}$ 且 $N = \mathbf{Q}/\mathbf{Z}$。
\end{enumerate}
\end{exercise}

\begin{exercise}
\label{exercises-exercise-compute-in-regular}
设 $R = k[x, y]$，其中 $k$ 是域。
\begin{enumerate}
\item 直接证明 Koszul 复形
$$
0 \to R
\xrightarrow{
\left(
\begin{matrix}
y \\
-x
\end{matrix}
\right)
}
R^{\oplus 2} \xrightarrow{(x, y)} R \xrightarrow{f \mapsto f(0, 0)} k \to 0
$$
是正合的。
\item 计算 $\Ext^i_R(k, k)$，其中把 $k = R/(x, y)$ 视为 $R$-模。
\end{enumerate}
\end{exercise}

\begin{exercise}
\label{exercises-exercise-infinitely-many-nonzero-ext}
给出 Noether 环 $R$ 以及有限模 $M$、$N$ 的例子，使得对所有
$i \geq 0$，$\Ext^i_R(M, N)$ 都非零。
\end{exercise}

\begin{exercise}
\label{exercises-exercise-infinite-ext}
给出环 $R$ 和理想 $I$ 的例子，使得
$\Ext^1_R(R/I, R/I)$ 不是有限 $R$-模。
（由《代数》引理 \ref{algebra-lemma-ext-noetherian} 可知，当
$R$ 为 Noether 环时不可能发生这种情形。）
\end{exercise}







\section{深度}
\label{exercises-section-depth}

\noindent
深度在《代数》第 \ref{algebra-section-depth} 节中定义，并在
《对偶化复形》第 \ref{dualizing-section-depth} 节中进一步研究。

\begin{exercise}
\label{exercises-exercise-compute-depth}
设 $R$ 为环，$I \subset R$ 为理想，$M$ 为 $R$-模。
在下列各情形中计算 $\text{depth}_I(M)$。
\begin{enumerate}
\item $R = \mathbf{Z}$，$I = (30)$，$M = \mathbf{Z}$；
\item $R = \mathbf{Z}$，$I = (30)$，$M = \mathbf{Z}/(300)$；
\item $R = \mathbf{Z}$，$I = (30)$，$M = \mathbf{Z}/(7)$；
\item $R = k[x, y, z]/(x^2 + y^2 + z^2)$，$I = (x, y, z)$，$M = R$；
\item $R = k[x, y, z, w]/(xz, xw, yz, yw)$，$I = (x, y, z, w)$，
$M = R$。
\end{enumerate}
这里 $k$ 是域。在最后两种情形中，可自由地在极大理想 $I$ 处局部化。
\end{exercise}

\begin{exercise}
\label{exercises-exercise-depth-not-inherited-localization}
给出深度 $\geq 1$ 的 Noether 局部环 $(R, \mathfrak m, \kappa)$
以及素理想 $\mathfrak p$ 的例子，使得
\begin{enumerate}
\item $\text{depth}_\mathfrak m(R) \geq 1$；
\item $\text{depth}_\mathfrak p(R_\mathfrak p) = 0$；
\item $\dim(R_\mathfrak p) \geq 1$。
\end{enumerate}
若不要求（3），这个练习就太容易了。为什么？
\end{exercise}

\begin{exercise}
\label{exercises-exercise-depth-torsion-free}
设 $(R, \mathfrak m)$ 为 Noether 局部整环，$M$ 为有限 $R$-模。
\begin{enumerate}
\item 若 $M$ 无挠，证明 $M$ 在 $R$ 上的深度至少为 $1$。
\item 给出深度恰为 $1$ 的例子。
\end{enumerate}
\end{exercise}

\begin{exercise}
\label{exercises-exercise-depth-examples}
对每一对 $m \geq n \geq 0$，给出 Noether 局部环 $R$ 的例子，使得
$\dim(R) = m$ 且 $\text{depth}(R) = n$。
\end{exercise}

\begin{exercise}
\label{exercises-exercise-make-depth-1}
设 $(R, \mathfrak m)$ 为 Noether 局部环，$M$ 为有限 $R$-模。
证明存在典范短正合列
$$
0 \to K \to M \to Q \to 0
$$
满足下列条件：
\begin{enumerate}
\item $\text{depth}(Q) \geq 1$；
\item $K$ 为零，或 $\text{Supp}(K) = \{\mathfrak m\}$；
\item $\text{length}_R(K) < \infty$。
\end{enumerate}
提示：利用 Noether 性证明存在满足（2）的极大子模 $K$，再证明
$Q = M/K$ 满足（1），且 $K$ 满足（3）。
\end{exercise}

\begin{exercise}
\label{exercises-exercise-make-depth-2}
设 $(R, \mathfrak m)$ 为 Noether 局部环，$M$ 为深度
$\geq 2$ 的有限 $R$-模。设 $N \subset M$ 为非零子模。
\begin{enumerate}
\item 证明 $\text{depth}(N) \geq 1$。
\item 证明 $\text{depth}(N) = 1$，当且仅当商模 $M/N$ 满足
$\text{depth}(M/N) = 0$。
\item 证明存在子模 $N' \subset M$，其中 $N \subset N'$ 且余长度有限，
即 $\text{length}_R(N'/N) < \infty$，使得 $N'$ 的深度
$\geq 2$。提示：把练习 \ref{exercises-exercise-make-depth-1} 应用于 $M/N$，
并取 $N'$ 为 $K$ 的逆像。
\end{enumerate}
\end{exercise}

\begin{exercise}
\label{exercises-exercise-Hartshorne-reduced}
设 $(R, \mathfrak m)$ 为 Noether 局部环。假设 $R$ 既约，即
$R$ 没有非零幂零元。再假设 $R$ 有两个不同的极小素理想
$\mathfrak p$ 和 $\mathfrak q$。
\begin{enumerate}
\item 证明 $R$-模序列
% P12-ERR-0427
$$
0 \to R \to R/\mathfrak p \oplus R/\mathfrak q \to
R/\mathfrak p + \mathfrak q \to 0
$$
是正合的（逐处检查）。这些映射分别为
$x \mapsto (x \bmod \mathfrak p, x \bmod \mathfrak q)$
以及 $(y \bmod \mathfrak p, z \bmod \mathfrak q) \mapsto
(y - z \bmod \mathfrak p + \mathfrak q)$。
\item 证明若 $\text{depth}(R) \geq 2$，则
$\dim(R/\mathfrak p + \mathfrak q) \geq 1$。
\item 证明若 $\text{depth}(R) \geq 2$，则
$U = \Spec(R) \setminus \{\mathfrak m\}$ 是连通拓扑空间。
\end{enumerate}
这证明了 Hartshorne 连通性定理的一个非常特殊的情形。该定理断言：
深度 $\text{depth} \geq 2$ 的 Noether 局部环之去心谱 $U$ 是连通的。
\end{exercise}

\begin{exercise}
\label{exercises-exercise-depth-2}
设 $(R, \mathfrak m)$ 为 Noether 局部环。设
$x, y \in \mathfrak m$ 为长度 $2$ 的正则序列。证明对任意
$n \geq 2$，不存在 $a, b \in R$ 使得
$$
x^{n - 1}y^{n - 1} = a x^n + b y^n
$$
建议：先尝试 $n = 2$，以了解论证方法。
注：此结果有一个极为广泛的推广，称为单项式猜想。
\end{exercise}







\section{Cohen--Macaulay 模与环}
\label{exercises-section-CM}

\noindent
Cohen--Macaulay 模在《代数》第 \ref{algebra-section-CM} 节中研究，
Cohen--Macaulay 环则在《代数》第
\ref{algebra-section-CM-ring} 节中研究。

\begin{exercise}
\label{exercises-exercise-examples-CM}
请在下列各情形中回答“是”或“否”，无需解释或证明。
\begin{enumerate}
\item 设 $p$ 为素数。局部环 $\mathbf{Z}_{(p)}$ 是
Cohen--Macaulay 局部环吗？
\item 设 $p$ 为素数。局部环 $\mathbf{Z}_{(p)}$ 是正则局部环吗？
\item 设 $k$ 为域。局部环 $k[x]_{(x)}$ 是
Cohen--Macaulay 局部环吗？
\item 设 $k$ 为域。局部环 $k[x]_{(x)}$ 是正则局部环吗？
\item 设 $k$ 为域。局部环 $(k[x, y]/(y^2 - x^3))_{(x, y)} =
k[x, y]_{(x, y)}/(y^2 - x^3)$ 是 Cohen--Macaulay 局部环吗？
\item 设 $k$ 为域。局部环 $(k[x, y]/(y^2, xy))_{(x, y)} =
k[x, y]_{(x, y)}/(y^2, xy)$ 是 Cohen--Macaulay 局部环吗？
\end{enumerate}
\end{exercise}













\section{奇点}
\label{exercises-section-singularities}

\begin{exercise}
\label{exercises-exercise-singularities}
任取域 $k$。设 $A = k[[x, y]]/(f)$ 且
$B = k[[u, v]]/(g)$，其中 $f = xy$，$g = uv + \delta$，并且
$\delta \in (u, v)^3$。证明 $A$ 与 $B$ 是同构的环。
\end{exercise}

\begin{remark}
\label{exercises-remark-singularities}
设曲线定义在域 $k$ 上。若曲线在某点的完备局部环形如
$k'[[x, y]]/(f)$，其中（a）$k'$ 是 $k$ 的有限可分扩张；
（b）$f$ 的初始项次数为二，即形如
$q = ax^2 + bxy + cy^2$，其中 $a, b, c\in k'$ 不全为零；
（c）$q$ 是 $k'$ 上的非退化二次型（在特征 2 时，这意味着
$b$ 非零），则称该点为普通二重点。一般而言，每个同构类的序对
$(k', q)$ 都对应这类环的一个同构类。
\end{remark}

\begin{exercise}
\label{exercises-exercise-periodic-resolution}
设 $R$ 为环，$n \geq 1$。设 $A$、$B$ 为系数在 $R$ 中的
$n \times n$ 矩阵，并且对 $R$ 中某个非零因子 $f$ 有
$AB = f 1_{n \times n}$。令 $S = R/(f)$。证明
$$
\ldots \to
S^{\oplus n} \xrightarrow{B}
S^{\oplus n} \xrightarrow{A}
S^{\oplus n} \xrightarrow{B}
S^{\oplus n} \to \ldots
$$
是正合的。
\end{exercise}







\section{可构造集}
\label{exercises-section-constructible}

\noindent
设 $k$ 为代数闭域，例如复数域 $\mathbf{C}$。设 $n \geq 0$。
多项式 $f \in k[x_1, \ldots, x_n]$ 通过取值给出函数
$f : k^n \to k$。若子集 $Z \subset k^n$ 是一族多项式的公共零点集，
则称该子集为{\it 代数集}。

\begin{exercise}
\label{exercises-exercise-finite-nr-equations}
证明代数集总能写成有限多个多项式的零点集。
\end{exercise}

\noindent
沿用上述记号。若子集 $E \subset k^n$ 是有限多个形如
$Z \cap \{f \not = 0\}$ 的集合之并，其中 $f$ 是多项式，
则称该子集{\it 可构造}。

\begin{exercise}
\label{exercises-exercise-constructible-classical}
证明下列各点：
\begin{enumerate}
\item 可构造集的补集是可构造集；
\item 有限多个可构造集的并是可构造集；
\item 有限多个可构造集的交是可构造集；
\item 任意可构造集 $E$ 都可写成有限不交并
$E = \coprod E_i$，其中每个 $E_i$ 都形如
$Z \cap \{f \not = 0\}$，$Z$ 为代数集，$f$ 为多项式。
\end{enumerate}
\end{exercise}

\begin{exercise}
\label{exercises-exercise-division-with-remainder}
设 $R$ 为环，并设
$f = a_d x^d + a_{d - 1} x^{d - 1} + \ldots + a_0 \in R[x]$。
（照例，这一记号意味着 $a_0, \ldots, a_d \in R$。）设 $g \in R[x]$。
证明可以找到 $N \geq 0$ 以及 $r, q \in R[x]$，使得
$$
a_d^N g = q f + r
$$
且 $\deg(r) < d$；即对某些 $c_i \in R$ 有
$r = c_0 + c_1 x + \ldots + c_{d - 1}x^{d - 1}$。
\end{exercise}







\section{Hilbert 零点定理}
\label{exercises-section-Hilbert-Nullstellensatz}


\begin{exercise}
\label{exercises-exercise-uncountable}
{\it 一个仅用复数的有趣论证！}
设 ${\mathbf C}$ 为复数域，$V$ 为 ${\mathbf C}$ 上的向量空间。
线性算子 $T : V \to V$ 的谱，是满足算子
$T - \lambda \text{id}_V$ 不可逆的复数 $\lambda \in {\mathbf C}$
所成的集合。
\begin{enumerate}
\item 证明 $\mathbf{C}(X)$ 在 ${\mathbf C}$ 上具有不可数维数。
\item 证明：若 $V$ 的维数有限或可数，则 $V$ 上的任意线性算子
都有非空谱。
\item 证明：若有限生成 ${\mathbf C}$-代数 $R$ 是域，则映射
${\mathbf C}\to R$ 是同构。
\item 证明 ${\mathbf C}[x_1, \ldots, x_n]$ 的任意极大理想
${\mathfrak m}$ 都形如
$(x_1-\alpha_1, \ldots, x_n-\alpha_n)$，其中某些
$\alpha_i \in {\mathbf C}$。
\end{enumerate}
\end{exercise}

\begin{remark}
\label{exercises-remark-HNSS}
设 $k$ 为域。那么，对每个整数 $n\in {\mathbf N}$ 以及每个极大理想
${\mathfrak m} \subset k[x_1, \ldots, x_n]$，商
$k[x_1, \ldots, x_n]/{\mathfrak m}$ 都是 $k$ 的有限域扩张。
这将在课程后面证明。当然（请检查），它蕴含有限生成
$k$-代数之极大理想的类似结论。上面的练习在
$k = {\mathbf C}$ 的情形证明了这一点。
\end{remark}

\begin{exercise}
\label{exercises-exercise-Hilbert-Nullstellensatz}
设 $k$ 为域。请使用注 \ref{exercises-remark-HNSS}。
\begin{enumerate}
\item 设 $R$ 为 $k$-代数。假设 $\dim_k R < \infty$ 且
$R$ 是整环。证明 $R$ 是域。
\item 假设 $R$ 是有限生成 $k$-代数，且 $f\in R$ 不是幂零元。
证明存在极大理想 ${\mathfrak m} \subset R$，使得
$f\not\in {\mathfrak m}$。
\item 用例子说明，当 $R$ 不是域上的有限型代数时，该命题不成立。
\item 证明任意根理想 $I \subset {\mathbf C}[x_1, \ldots, x_n]$
都是所有包含它的极大理想之交。
\end{enumerate}
\end{exercise}

\begin{remark}
\label{exercises-remark-Hilbert-Nullstellensatz}
这就是 Hilbert 零点定理。它断言
$\Spec(k[x_1, \ldots, x_n])$ 的闭子集
（由前面的练习，它们对应于根理想）由其中包含的闭点决定。
\end{remark}

\begin{exercise}
\label{exercises-exercise-product-matrices-ring}
设
$A =
{\mathbf C}[x_{11}, x_{12}, x_{21}, x_{22}, y_{11}, y_{12}, y_{21}, y_{22}]$。
设 $I$ 为 $A$ 中由矩阵 $XY$ 的各个元素生成的理想，其中
$$
X = \left(
\begin{matrix}
x_{11} & x_{12}\\
x_{21} & x_{22}
\end{matrix}
\right)
\quad\text{且}\quad
Y = \left(
\begin{matrix}
y_{11} & y_{12}\\
y_{21} & y_{22}
\end{matrix}
\right).
$$
求 $\Spec(A)$ 的闭子集 $V(I)$ 的不可约分支。
（这里是要描述这些分支并给出每个分支的方程；不必证明所写方程定义
素理想。）提示：
\begin{enumerate}
\item 可以使用 Hilbert 零点定理；只需找出覆盖 $V(I)$
之闭点集的不可约局部闭子集。
\item 有两个容易找到的分支。
\item 不可约集在连续映射下的像不可约。
\end{enumerate}
\end{exercise}




\section{维数}
\label{exercises-section-dimension}

\begin{exercise}
\label{exercises-exercise-dimension-bigger-one-finite-nr-primes}
构造仅有有限多个素理想、但维数 $> 1$ 的环 $A$。
\end{exercise}

\begin{exercise}
\label{exercises-exercise-hypersurface-in-A2-dimension-one}
设 $f \in \mathbf{C}[x, y]$ 为非常数多项式。
证明 $\mathbf{C}[x, y]/(f)$ 的维数为 $1$。
\end{exercise}

\begin{exercise}
\label{exercises-exercise-dimension-polynomial-ring}
设 $(R, \mathfrak m)$ 为 Noether 局部环。设 $n \geq 1$。
在多项式环 $R[x_1, \ldots, x_n]$ 中令
$\mathfrak m' = (\mathfrak m, x_1, \ldots, x_n)$。证明
$$
\dim(R[x_1, \ldots, x_n]_{\mathfrak m'}) = \dim(R) + n.
$$
\end{exercise}





\section{链环}
\label{exercises-section-catenary}

\begin{definition}
\label{exercises-definition-catenary}
若对任意三个素理想
${\mathfrak p}_1 \subset {\mathfrak p}_2 \subset {\mathfrak p}_3$ 均有
$$
ht({\mathfrak p}_3 / {\mathfrak p}_1) = ht({\mathfrak p}_3/{\mathfrak p}_2) +
ht({\mathfrak p}_2/{\mathfrak p}_1).
$$
则称 Noether 环 $A$ 为{\it 链环}。这里
$ht(\mathfrak p/\mathfrak q)$ 表示 $\mathfrak p/\mathfrak q$
在环 $A/\mathfrak q$ 中的高度。写成公式即
$$
ht(\mathfrak p/\mathfrak q) =
\dim(A_\mathfrak p/\mathfrak qA_\mathfrak p) =
\dim((A/\mathfrak q)_\mathfrak p) =
\dim((A/\mathfrak q)_{\mathfrak p/\mathfrak q})
$$
若 $T \subset T' \subset X$，其中 $T$ 与 $T'$ 为闭不可约子集，
则存在一条不可约闭子集的极大链
$$
T = T_0 \subset T_1 \subset \ldots \subset T_n = T'
$$
并且每条这样的链都有相同的有限长度，则称拓扑空间 $X$ {\it 链状}。
\end{definition}

\begin{exercise}
\label{exercises-exercise-catenary-the-same}
证明：对 Noether 环而言，《代数》定义
\ref{algebra-definition-catenary} 中定义的链状性概念，与定义
\ref{exercises-definition-catenary} 中的概念一致。
\end{exercise}

\begin{exercise}
\label{exercises-exercise-Noetherian-local-domain-dim-2-catenary}
证明维数为 $2$ 的 Noether 局部整环是链环。
\end{exercise}

\begin{exercise}
\label{exercises-exercise-finite-type-over-field-catenary}
设 $k$ 为域。证明有限型 $k$-代数是链环。
\end{exercise}

\begin{exercise}
\label{exercises-exercise-example-no-dim-function}
给出一个有限、清醒且链状的拓扑空间 $X$，使它不存在维数函数
$\delta : X \to \mathbf{Z}$。这里，映射
$\delta : X \to \mathbf{Z}$ 称为维数函数，是指对
$x, y \in X$ 有：
\begin{enumerate}
\item 若 $x \leadsto y$ 且 $x \not = y$，则
$\delta(x) > \delta(y)$；
\item 若 $x \leadsto y$ 且 $\delta(x) \geq \delta(y) + 2$，
则存在 $z \in X$，使得 $x \leadsto z \leadsto y$ 且
$\delta(x) > \delta(z) > \delta(y)$。
\end{enumerate}
清楚地描述你的空间，并简洁说明为什么它不可能有维数函数。
\end{exercise}




\section{分式域}
\label{exercises-section-fraction-fields}

\begin{exercise}
\label{exercises-exercise-find-fraction-field}
考察整环
$$
{\mathbf Q}[r, s, t]/(s^2-(r-1)(r-2)(r-3), t^2-(r + 1)(r + 2)(r + 3)).
$$
找出一个形如 ${\mathbf Q}[x, y]/(f)$ 的整环，使其分式域与之同构。
\end{exercise}



\section{超越次数}
\label{exercises-section-transcendence}

\begin{exercise}
\label{exercises-exercise-algebraic-extension}
设 $K'/K/k$ 为域扩张，并且 $K'$ 在 $K$ 上代数。
证明 $\text{trdeg}_k(K) = \text{trdeg}_k(K')$。
（提示：证明若 $x_1, \ldots, x_d \in K$ 在 $k$ 上代数无关，且
$d < \text{trdeg}_k(K')$，则 $k(x_1, \ldots, x_d) \subset K$
不可能是代数扩张。）
\end{exercise}

\begin{exercise}
\label{exercises-exercise-growth-powers-subvector-space}
设 $k$ 为域，$K/k$ 为超越次数 $d$ 的有限生成扩张。
若 $V, W \subset K$ 是有限维 $k$-向量子空间，记
$$
VW = \{f \in K \mid f = \sum\nolimits_{i = 1, \ldots, n} v_i w_i
\text{，其中某个 }n\text{ 且 }v_i \in V, w_i \in W\}
$$
这是有限维 $k$-向量子空间。令 $V^2 = VV$、$V^3 = V V^2$，依此类推。
\begin{enumerate}
\item 证明可找到 $V \subset K$ 和 $\epsilon > 0$，使得对所有
$n \geq 1$ 都有 $\dim V^n \geq \epsilon n^d$。
\item 反过来，证明对每个有限维 $V \subset K$，存在 $C > 0$，
使得对所有 $n \geq 1$ 都有 $\dim V^n \leq C n^d$。
（一种可能的做法是：先对 $k[x_1, \ldots, x_d]$ 的向量子空间证明；
再对 $k(x_1, \ldots, x_d)$ 的向量子空间证明；最后，若
$K/k(x_1, \ldots, x_d)$ 是有限扩张，就选取 $K$ 在
$k(x_1, \ldots, x_d)$ 上的一组基，并按此基展开论证。）
\item 推出：可以用 $K$ 的有限维向量子空间各次幂的增长来重新定义
超越次数。
\end{enumerate}
这与 $k$ 上（非交换）代数的 Gelfand--Kirillov 维数有关。
\end{exercise}








\section{纤维维数}
\label{exercises-section-dimension-fibres}

\noindent
下面是一些与维数公式有关的问题；见《代数》第
\ref{algebra-section-dimension-formula} 节。

\begin{exercise}
\label{exercises-exercise-nr-components-fibre}
设 $k$ 为你偏爱的代数闭域。下文中 $k[x]$ 和 $k[x, y]$
表示多项式环。
\begin{enumerate}
\item 对每个整数 $n \geq 0$，找出整环的有限型扩张
$k[x] \subset A$，使得 $A/xA$ 的谱恰有 $n$ 个不可约分支。
\item 构造整环的有限型扩张 $k[x] \subset A$，使得对每个
$\alpha \in k$，$A/(x - \alpha)A$ 的谱都非空且可约。
\item 构造整环的有限型扩张 $k[x, y] \subset A$，使得对所有
$(\alpha, \beta) \in k^2 \setminus \{(0, 0)\}$，
$A/(x - \alpha, y - \beta)A$ 的谱都不可约
\footnote{回忆：不可约蕴含非空。}，而 $A/(x, y)A$ 的谱非空且可约。
\end{enumerate}
\end{exercise}

\begin{exercise}
\label{exercises-exercise-codim-1}
设 $k$ 为你偏爱的代数闭域。设 $n \geq 1$。
令 $k[x_1, \ldots, x_n]$ 为多项式环，并令
$\mathfrak m = (x_1, \ldots, x_n)$。设
$k[x_1, \ldots, x_n] \subset A$ 为整环的有限型扩张。
令 $d = \dim(A)$。
\begin{enumerate}
\item 若 $A/\mathfrak mA \not = 0$，证明
$d - 1 \geq \dim(A/\mathfrak m A) \geq d - n$。
\item 用例子说明每个数值都能出现。
\item 用例子说明 $\Spec(A/\mathfrak m A)$ 可以有维数不同的
不可约分支。
\end{enumerate}
\end{exercise}




\section{有限局部自由模}
\label{exercises-section-finite-locally-free}

\begin{definition}
\label{exercises-definition-finite-locally-free}
设 $A$ 为环。回忆：{\it 有限局部自由} $A$-模 $M$，是指满足如下条件的模：
对每个 ${\mathfrak p} \in \Spec(A)$，都存在
$f\in A$、$f \not \in {\mathfrak p}$，使 $M_f$ 是有限自由
$A_f$-模。若 $M$ 是秩为 $1$ 的有限局部自由模，即对每个
${\mathfrak p} \in \Spec(A)$，都存在
$f\in A$、$f \not \in \mathfrak p$，使得作为 $A_f$-模有
$M_f \cong A_f$，则称 $M$ 为{\it 可逆模}。
\end{definition}

\begin{exercise}
\label{exercises-exercise-tensor-finite-locally-free}
证明有限局部自由模的张量积仍有限局部自由。证明两个可逆模的张量积
仍可逆。
\end{exercise}

\begin{definition}
\label{exercises-definition-class-group}
设 $A$ 为环。$A$ 的{\it 类群}，有时也称 $A$ 的{\it Picard 群}，
是可逆 $A$-模的同构类集合 $\Pic(A)$，其群运算由张量积定义
（见练习 \ref{exercises-exercise-tensor-finite-locally-free}）。
\end{definition}

\noindent
注意，$A$ 的类群平凡，当且仅当每个可逆模都同构于秩为 1 的自由模。

\begin{exercise}
\label{exercises-exercise-class-group-trivial}
证明下列各环的类群平凡：
\begin{enumerate}
\item 多项式环 $A = k[x]$，其中 $k$ 为域；
\item 整数环 $A = \mathbf{Z}$；
\item 多项式环 $A = k[x, y]$，其中 $k$ 为域；
\item 商环 $k[x, y]/(xy)$，其中 $k$ 为域。
\end{enumerate}
\end{exercise}

\begin{exercise}
\label{exercises-exercise-class-group-not-trivial}
设 $k$ 为特征不为 $2$ 的域，并设
$f(x) = (x - t_1) \ldots (x - t_n)$，其中
$t_1, \ldots, t_n \in k$ 两两不同，$n \geq 3$ 为奇数。
证明环 $A = k[x, y]/(y^2 - f(x))$ 的类群非平凡。
（提示：证明理想 $(y, x - t_1)$ 定义 $\Pic(A)$ 的一个非平凡元素。）
\end{exercise}

\begin{exercise}
\label{exercises-exercise-trace-det}
设 $A$ 为环。
\begin{enumerate}
\item 假设 $M$ 是有限局部自由 $A$-模，且
$\varphi : M \to M$ 是自同态。定义或构造 $\varphi$ 的{\it 迹}
与{\it 行列式}，并证明你的构造“关于三元组 $(A, M, \varphi)$
具有函子性”。
\item 证明：若 $M, N$ 是有限局部自由 $A$-模，并且
$\varphi : M \to N$、$\psi : N \to M$，则
% P12-ERR-0428
$\text{Trace}(\varphi \circ \psi) = \text{Trace}(\psi \circ \varphi)$ 且
$\det(\varphi \circ \psi) = \det(\psi \circ \varphi)$。
\item 当 $M$ 有限局部自由时，证明 $\text{Trace}$ 定义
$A$-线性映射 $\text{End}_A(M) \to A$，而 $\det$ 定义乘法映射
$\text{End}_A(M) \to A$。
\end{enumerate}
\end{exercise}

\begin{exercise}
\label{exercises-exercise-trace-det-rings}
现在假设 $B$ 是一个 $A$-代数，且作为 $A$-模有限局部自由；
换言之，$B$ 是有限局部自由 $A$-代数。
\begin{enumerate}
\item 使用练习 \ref{exercises-exercise-trace-det} 中的 $\text{Trace}$ 和
$\det$ 定义 $\text{Trace}_{B/A}$ 与 $\text{Norm}_{B/A}$。
\item 设 $b\in B$，并设 $\pi : \Spec(B) \to \Spec(A)$ 为诱导态射。
证明 $\pi(V(b)) = V(\text{Norm}_{B/A}(b))$。
（回忆 $V(f) = \{ {\mathfrak p} \mid f \in {\mathfrak p}\}$。）
\item （基变换。）假设 $i : A \to A'$ 是环映射。令
$B' = B \otimes_A A'$。说明为什么 $i(\text{Norm}_{B/A}(b))$ 等于
$\text{Norm}_{B'/A'}(b \otimes 1)$。
\item 当 $B = A \times A \times A \times \ldots \times A$ 且
$b = (a_1, \ldots, a_n)$ 时，计算 $\text{Norm}_{B/A}(b)$。
\item 对有限平坦映射
${\mathbf Q}[x] \to {\mathbf Q}[y]$、$x \to y^n$，计算
$y-y^3$ 的范数。（提示：使用“基变换”
$A = {\mathbf Q}[x] \subset A' = {\mathbf Q}(\zeta_n)(x^{1/n})$。）
\end{enumerate}
\end{exercise}



\section{粘合}
\label{exercises-section-glueing}

\begin{exercise}
\label{exercises-exercise-cover}
假设 $A$ 是环，$M$ 是 $A$-模。设 $f_i$、$i \in I$
为 $A$ 的一族元素，并且
$$
\Spec(A) = \bigcup D(f_i).
$$
\begin{enumerate}
\item 证明若 $M_{f_i}$ 是有限 $A_{f_i}$-模，则 $M$ 是有限
$A$-模。
\item 证明若 $M_{f_i}$ 是平坦 $A_{f_i}$-模，则 $M$ 是平坦
$A$-模。（仔细想想，这一点其实有些显然。）
\end{enumerate}
\end{exercise}

\begin{remark}
\label{exercises-remark-cover}
用代数几何语言来说，这意味着“有限生成”或“平坦”这一性质
（在适当意义下）对 Zariski 拓扑是局部的。还可以对“有限表示”这一
性质证明相同结论。
\end{remark}

\begin{exercise}
\label{exercises-exercise-cover-ring-map}
假设 $A \to B$ 是环映射。设 $f_i \in A$、$i \in I$ 以及
$g_j \in B$、$j \in J$ 为两族元素，并且
$$
\Spec(A) = \bigcup D(f_i)
\quad\text{且}\quad
\Spec(B) = \bigcup D(g_j).
$$
证明：若对所有 $i, j$，$A_{f_i} \to B_{f_ig_j}$ 都是有限型的，
则 $A \to B$ 是有限型的。
\end{exercise}




\section{上升与下降}
\label{exercises-section-going-up}


\begin{definition}
\label{exercises-definition-GU-GD}
设 $\phi : A \to B$ 为环同态。若满足下列条件，则称 $\phi$
满足{\it 上升定理}：
\begin{itemize}
\item[(GU)] 对任意
${\mathfrak p}, {\mathfrak p}' \in \Spec(A)$，若
${\mathfrak p} \subset {\mathfrak p}'$，并且任取位于
${\mathfrak p}$ 上方的 $P \in \Spec(B)$，都存在位于
${\mathfrak p}'$ 上方的 $P'\in \Spec(B)$，使得 $P \subset P'$。
\end{itemize}
类似地，若满足下列条件，则称 $\phi$ 满足{\it 下降定理}：
\begin{itemize}
\item[(GD)] 对任意
${\mathfrak p}, {\mathfrak p}' \in \Spec(A)$，若
${\mathfrak p} \subset {\mathfrak p}'$，并且任取位于
${\mathfrak p}'$ 上方的 $P' \in \Spec(B)$，都存在位于
${\mathfrak p}$ 上方的 $P\in \Spec(B)$，使得 $P \subset P'$。
\end{itemize}
\end{definition}

\begin{exercise}
\label{exercises-exercise-GU-GD}
在下列每一种情形中，判断（GU）与（GD）是否成立，并说明理由。
（可以使用能找到的任何命题、定理或引理，但要逐一检查假设。）
\begin{enumerate}
\item $k$ 是域，$A = k$，$B = k[x]$。
\item $k$ 是域，$A = k[x]$，$B = k[x, y]$。
\item $A = {\mathbf Z}$，$B = {\mathbf Z}[1/11]$。
\item $k$ 是代数闭域，$A = k[x, y]$，
$B = k[x, y, z]/(x^2-y, z^2-x)$。
\item $A = {\mathbf Z}$，$B = {\mathbf Z}[i, 1/(2 + i)]$。
\item $A = {\mathbf Z}$，$B = {\mathbf Z}[i, 1/(14 + 7i)]$。
\item $k$ 是代数闭域，$A = k[x]$，
$B = k[x, y, 1/(xy-1)]/(y^2-y)$。
\end{enumerate}
\end{exercise}

\begin{exercise}
\label{exercises-exercise-image}
设 $A$ 为环，$B = A[x]$ 为 $A$ 上的单变量多项式代数。设
$f = a_0 + a_1 x + \ldots + a_r x^r \in B = A[x]$。
仔细证明 $D(f)$ 在 $\Spec(A)$ 中的像等于
$D(a_0) \cup \ldots \cup D(a_r)$。
\end{exercise}

\begin{exercise}
\label{exercises-exercise-images}
设 $k$ 为代数闭域。计算下列映射在 $\Spec(k[x, y])$ 中的像：
\begin{enumerate}
\item $\Spec(k[x, yx^{-1}]) \to \Spec(k[x, y])$，其中
$k[x, y] \subset k[x, yx^{-1}] \subset k[x, y, x^{-1}]$。
（提示：为避免混淆，给元素 $yx^{-1}$ 另取一个名字。）
\item $\Spec(k[x, y, a, b]/(ax-by-1))\to \Spec(k[x, y])$。
\item $\Spec(k[t, 1/(t-1)]) \to \Spec(k[x, y])$，它由
$x \mapsto t^2$ 和 $y \mapsto t^3$ 诱导。
\item $k = {\mathbf C}$（复数），
$\Spec(k[s, t]/(s^3 + t^3-1)) \to \Spec(k[x, y])$，其中
$x\mapsto s^2$，$y \mapsto t^2$。
\end{enumerate}
\end{exercise}

\begin{remark}
\label{exercises-remark-elimination-theory}
求上述像通常要使用消元理论。
\end{remark}




\section{Fitting 理想}
\label{exercises-section-fitting-ideals}

\begin{exercise}
\label{exercises-exercise-fitting}
设 $R$ 为环，$M$ 为有限 $R$-模。选取 $M$ 的一个表示
$$
\bigoplus\nolimits_{j \in J} R \longrightarrow R^{\oplus n}
\longrightarrow M \longrightarrow 0.
$$
注意，映射 $R^{\oplus n} \to M$ 由 $M$ 的一列元素
$x_1, \ldots, x_n$ 给出；这些元素 $x_i$ 是 $M$ 的{\it 生成元}。
映射 $\bigoplus_{j \in J} R \to R^{\oplus n}$ 由系数在 $R$ 中的
$n \times J$ 矩阵 $A$ 给出。换言之，
$A = (a_{ij})_{i = 1, \ldots, n, j \in J}$。$A$ 的各列
$(a_{1j}, \ldots, a_{nj})$（$j \in J$）称为{\it 关系}。
任何满足 $\sum r_i x_i = 0$ 的向量 $(r_i) \in R^{\oplus n}$
都是 $A$ 的各列的线性组合。当然，每个有限 $R$-模都有许多不同的表示。
\begin{enumerate}
\item 证明由 $A$ 的 $(n - k) \times (n - k)$ 子式生成的理想
与表示的选择无关。该理想称为 $M$ 的{\it 第 $k$ 个 Fitting 理想}，
记作 $Fit_k(M)$。
\item 证明
$Fit_0(M) \subset Fit_1(M) \subset Fit_2(M) \subset \ldots$。
（提示：利用按一列展开来计算行列式。）
\item 证明下列条件等价：
\begin{enumerate}
\item $Fit_{r - 1}(M) = (0)$ 且 $Fit_r(M) = R$；
\item $M$ 是秩为 $r$ 的局部自由模。
\end{enumerate}
\end{enumerate}
\end{exercise}




\section{Hilbert 函数}
\label{exercises-section-hilbert}

\begin{definition}
\label{exercises-definition-numerical-polynomial}
{\it 数值多项式}是满足如下条件的多项式
$f(x) \in {\mathbf Q}[x]$：对每个整数 $n$，都有
$f(n) \in {\mathbf Z}$。
\end{definition}

\begin{definition}
\label{exercises-definition-graded-module}
环 $A$ 上的{\it 分次模} $M$，是带有 $A$-子模直和分解
$
\bigoplus\nolimits_{n \in {\mathbf Z}} M_n
$
的 $A$-模 $M$。若所有 $M_n$ 都是有限 $A$-模，则称 $M$
{\it 局部有限}。假设 $A$ 是 Noether 环，并且 $\varphi$
是有限 $A$-模上的一个{\it Euler--Poincar\'e 函数}。
这意味着：对每个有限生成 $A$-模 $M$，给定一个整数
$\varphi(M) \in {\mathbf Z}$；并且对每个短正合列
$$
0
\longrightarrow
M'
\longrightarrow
M
\longrightarrow
M''
\longrightarrow
0
$$
都有 $\varphi(M) = \varphi(M') + \varphi(M'')$。
局部有限分次模 $M$（关于 $\varphi$）的{\it Hilbert 函数}是函数
$\chi_\varphi(M, n) = \varphi(M_n)$。若存在某个数值多项式
$P_\varphi$，使得对所有充分大的整数 $n$ 都有
$\chi_\varphi(M, n) = P_\varphi(n)$，则称 $M$ 有
{\it Hilbert 多项式}。
\end{definition}

\begin{definition}
\label{exercises-definition-graded-algebra}
{\it 分次 $A$-代数}是分次 $A$-模
$B = \bigoplus_{n \geq 0} B_n$，连同一个 $A$-双线性映射
$$
B \times B \longrightarrow B, \ (b, b') \longmapsto bb'
$$
使 $B$ 成为 $A$-代数，且 $B_n \cdot B_m \subset B_{n + m}$。
最后，{\it 分次 $A$-代数 $B$ 上的分次模 $M$}，由一个分次
$A$-模 $M$ 连同一个相容的 $B$-模结构给出，并满足
$B_n \cdot M_d \subset M_{n + d}$。现在可以定义
{\it 分次模或分次环的同态}、{\it 分次子模}、{\it 分次理想}、
{\it 分次模的正合列}等等。
\end{definition}

\begin{exercise}
\label{exercises-exercise-Euler-Poincare-field}
设 $A = k$ 为域。在此情形中，有限 $A$-模上所有可能的
Euler--Poincar\'e 函数是什么？
\end{exercise}

\begin{exercise}
\label{exercises-exercise-Euler-Poincare-Z}
设 $A ={\mathbf Z}$。在此情形中，有限 $A$-模上所有可能的
Euler--Poincar\'e 函数是什么？
\end{exercise}

\begin{exercise}
\label{exercises-exercise-Euler-Poincare-node}
设 $A = k[x, y]/(xy)$，其中 $k$ 代数闭。在此情形中，有限
$A$-模上所有可能的 Euler--Poincar\'e 函数是什么？
\end{exercise}

\begin{exercise}
\label{exercises-exercise-kernel-locally-finite}
假设 $A$ 是 Noether 环。证明局部有限分次 $A$-模之间的映射之核
仍局部有限。
\end{exercise}

\begin{exercise}
\label{exercises-exercise-no-hilbert}
设 $k$ 为域，$A = k$，并设 $B = k[x, y]$ 的分次由
$\deg(x) = 2$ 与 $\deg(y) = 3$ 决定。令
$\varphi(M) = \dim_k(M)$。计算 $B$ 作为分次 $k$-模的
Hilbert 函数。在此情形中有 Hilbert 多项式吗？
\end{exercise}

\begin{exercise}
\label{exercises-exercise-no-hilbert-or-is-there}
设 $k$ 为域，$A = k$，并设 $B = k[x, y]/(x^2, xy)$ 的分次由
$\deg(x) = 2$ 与 $\deg(y) = 3$ 决定。令
$\varphi(M) = \dim_k(M)$。计算 $B$ 作为分次 $k$-模的
Hilbert 函数。在此情形中有 Hilbert 多项式吗？
\end{exercise}

\begin{exercise}
\label{exercises-exercise-hilbert-to-compute}
设 $k$ 为域，$A = k$。令 $\varphi(M) = \dim_k(M)$。
固定 $d\in {\mathbf N}$。考察分次 $A$-代数
$B = k[x, y, z]/(x^d + y^d + z^d)$，其中 $x, y, z$
各自的次数都是 $1$。计算 $B$ 的 Hilbert 函数。在此情形中
有 Hilbert 多项式吗？
\end{exercise}



\section{环的 Proj}
\label{exercises-section-proj-ring}

\begin{definition}
\label{exercises-definition-homogeneous-ideal}
设 $R$ 为分次环。{\it 齐次}理想就是同时也是 $R$ 的分次子模的理想
$I \subset R$。等价地，它是由齐次元素生成的理想。再等价地，若
$f \in I$，并且
$$
f = f_0 + f_1 + \ldots + f_n
$$
是 $f$ 在 $R$ 中分解为齐次部分的分解，则对每个 $i$ 都有
$f_i \in I$。
\end{definition}

\begin{definition}
\label{exercises-definition-Proj-R}
定义分次环 $R$ 的{\it 齐次谱 $\text{Proj}(R)$}为 $R$ 的
齐次素理想 ${\mathfrak p}$ 所成的集合，其中要求
$R_{+} \not \subset {\mathfrak p}$。注意 $\text{Proj}(R)$
是 $\Spec(R)$ 的子集，因而带有自然的诱导拓扑。
\end{definition}

\begin{definition}
\label{exercises-definition-Dplus-Vplus}
设 $R = \oplus_{d \geq 0} R_d$ 为分次环，设 $f\in R_d$，并假设
$d \geq 1$。定义 {\it $R_{(f)}$} 为 $R_f$ 的子环，由所有形如
$r/f^n$ 的元素组成，其中 $r$ 齐次且 $\deg(r) = nd$。此外，定义
$$
D_{+}(f) = \{ {\mathfrak p} \in \text{Proj}(R) | f \not\in {\mathfrak p} \}.
$$
最后，对齐次理想 $I \subset R$，定义
$V_{+}(I) = V(I) \cap \text{Proj}(R)$。
\end{definition}

\begin{exercise}
\label{exercises-exercise-topology-proj}
关于 $\text{Proj}(R)$ 上的拓扑。使用上述定义与记号，证明下列命题。
\begin{enumerate}
\item 证明 $D_{+}(f)$ 在 $\text{Proj}(R)$ 中是开的。
\item 证明 $D_{+}(ff') = D_{+}(f) \cap D_{+}(f')$。
\item 设 $g = g_0 + \ldots + g_m$ 是 $R$ 的元素，其中
$g_i \in R_i$。用 $D_{+}(g_i)$（$i \geq 1$）以及
$D(g_0) \cap \text{Proj}(R)$ 表示
$D(g) \cap \text{Proj}(R)$。无需证明。
\item 设 $g\in R_0$ 为次数 $0$ 的齐次元素。用一族适当的
正次数齐次元素 $f_\alpha \in R$ 的 $D_{+}(f_\alpha)$
表示 $D(g) \cap \text{Proj}(R)$。
\item 证明开集族 $\{D_{+}(f)\}$ 构成 $\text{Proj}(R)$
拓扑的一组基。
\item
\label{exercises-item-bijection}
证明存在典范双射 $D_{+}(f) \to \Spec(R_{(f)})$。
% P12-ERR-0429
（提示：模仿 $\Spec$ 情形的证明，但在某一步加入一个理想的根。）
\item 证明（\ref{exercises-item-bijection}）中的映射是同胚。
\item 给出 $R$ 的例子，使得 $\text{Proj}(R)$ 不拟紧。无需证明。
\item 证明任意闭子集 $T \subset \text{Proj}(R)$ 都形如
$V_{+}(I)$，其中 $I \subset R$ 是某个齐次理想。
\end{enumerate}
\end{exercise}

\begin{remark}
\label{exercises-remark-continuous-proj-spec}
存在连续映射 $ \text{Proj}(R) \longrightarrow \Spec(R_0) $。
\end{remark}

\begin{exercise}
\label{exercises-exercise-iso-polynomial-ring-one-variable}
若 $R = A[X]$ 且 $\deg(X) = 1$，证明自然映射
$\text{Proj}(R) \to \Spec(A)$ 是双射，事实上是同胚。
\end{exercise}

\begin{exercise}
\label{exercises-exercise-blowing-up-I}
爆破：第一部分。
在本练习中，$R = Bl_I(A) = A \oplus I \oplus I^2 \oplus \ldots$。
考察自然映射 $b : \text{Proj}(R) \to \Spec(A)$。
令 $U = \Spec(A) - V(I)$。证明
$$
b : b^{-1}(U) \longrightarrow U
$$
是同胚。因此可以把 $U$ 看作 $\text{Proj}(R)$ 的开子集。
设 $Z \subset \Spec(A)$ 为不可约闭子概形，泛点为
$\xi \in Z$。假设 $\xi \not\in V(I)$；换言之，
$Z \not\subset V(I)$；换言之，$\xi \in U$；换言之，
$Z\cap U \not = \emptyset$。定义 $Z$ 的{\it 严格变换}
$Z'$ 为位于 $\xi$ 上方的唯一一点 $\xi'$ 的闭包。另一种说法是，
$Z'$ 是局部闭子集 $Z\cap U \subset U \subset \text{Proj}(R)$
在 $\text{Proj}(R)$ 中的闭包。
\end{exercise}

\begin{exercise}
\label{exercises-exercise-blowing-up-II}
爆破：第二部分。
设 $A = k[x, y]$，其中 $k$ 是域，并设 $I = (x, y)$。
设 $R$ 为 $A$ 关于 $I$ 的爆破代数。
\begin{enumerate}
\item 证明 $Z_1 = V(\{x\})$ 与 $Z_2 = V(\{y\})$ 的严格变换不交。
\item 证明 $Z_1 = V(\{x\})$ 与 $Z_2 = V(\{x-y^2\})$
的严格变换相交。
\item 找出理想 $J \subset A$，使得 $V(J) = V(I)$，并且
$Z_1 = V(\{x\})$ 与 $Z_2 = V(\{x-y^2\})$ 在沿 $J$
的爆破中的严格变换不交。
\end{enumerate}
\end{exercise}

\begin{exercise}
\label{exercises-exercise-proj-when-empty}
设 $R$ 为分次环。
\begin{enumerate}
% P12-ERR-0430
\item 证明若对所有 $n >> 0$ 都有 $R_n = (0)$，则
$\text{Proj}(R)$ 为空。
\item 证明若 $R$ 是整环且 $R_{+}$ 非零，则
$\text{Proj}(R)$ 是不可约拓扑空间。（回忆：空拓扑空间不可约。）
\end{enumerate}
\end{exercise}

\begin{exercise}
\label{exercises-exercise-blowing-up-III}
爆破：第三部分。
考察严格变换定义中的 $A$、$I$、$U$ 与 $Z$。设
$Z = V({\mathfrak p})$，其中 ${\mathfrak p}$ 为素理想。令
$\bar A = A/{\mathfrak p}$，并令 $\bar I$ 为 $I$ 在 $\bar A$
中的像。
\begin{enumerate}
\item 证明存在满环映射
% P12-ERR-0431
$R: = Bl_I(A) \to \bar R: = Bl_{\bar I}(\bar A)$。
\item 证明上述环映射诱导从 $\text{Proj}(\bar R)$ 到 $Z$
的严格变换 $Z'$ 上的双射。（这并不容易。
% P12-ERR-0432
提示：使用上面的 5(b)。）
\item 推出严格变换 $Z' = V_{+}(P)$，其中 $P \subset R$
是由 $P_d = I^d \cap {\mathfrak p}$ 定义的齐次理想。
\item 假设 $Z_1 = V({\mathfrak p})$ 与
$Z_2 = V({\mathfrak q})$ 是由素理想定义的不可约闭子集，并且
$Z_1 \not \subset Z_2$、$Z_2 \not \subset Z_1$。
证明爆破理想 $I = {\mathfrak p} + {\mathfrak q}$ 会分离
$Z_1$ 与 $Z_2$ 的严格变换；即
$Z_1' \cap Z_2' = \emptyset$。（提示：考察（c）中的齐次理想
$P$ 与 $Q$，并考察 $V(P + Q)$。）
\end{enumerate}
\end{exercise}


\section{一维 Cohen--Macaulay 环}
\label{exercises-section-CM-dim-1}

\begin{definition}
\label{exercises-definition-CM}
若 Noether 局部环 $A$ 的维数为 $d$，并且 $A$ 存在参数系
$x_1, \ldots, x_d$，使得对 $i = 1, \ldots, d$，$x_i$
在 $A/(x_1, \ldots, x_{i-1})$ 中都是非零因子，则称该环
为维数 $d$ 的{\it Cohen--Macaulay} 环。
\end{definition}

\begin{exercise}
\label{exercises-exercise-CM-dim-1-I}
一维 Cohen--Macaulay 环。第一部分：理论。
\begin{enumerate}
% P12-ERR-0433
\item 设 $(A, {\mathfrak m})$ 为局部 Noether 环，且
$\dim A = 1$。证明若 $x\in {\mathfrak m}$ 不是零因子，则
\begin{enumerate}
\item $\dim A/xA = 0$；换言之，$A/xA$ 是 Artin 环；再换言之，
$\{x\}$ 是 $A$ 的参数系。
% P12-ERR-0434
\item $A$ 没有嵌入素理想。
\end{enumerate}
\item 反过来，设 $(A, {\mathfrak m})$ 为维数为 $1$ 的局部 Noether 环。
证明若 $A$ 没有嵌入素理想，则 ${\mathfrak m}$ 中存在非零因子。
\end{enumerate}
\end{exercise}

\begin{exercise}
\label{exercises-exercise-CM-dim-1-II}
一维 Cohen--Macaulay 环。第二部分：例子。
\begin{enumerate}
\item 设 $A$ 为 $k[x, y]/(x^2, xy)$ 在 $(x, y)$ 处的局部环。
\begin{enumerate}
\item 证明 $A$ 的维数为 1。
\item 证明 ${\mathfrak m}\subset A$ 的每个元素都是零因子。
\item 找出 $z\in {\mathfrak m}$，使得 $\dim A/zA = 0$
（无需证明）。
\end{enumerate}
\item 设 $A$ 为 $k[x, y]/(x^2)$ 在 $(x, y)$ 处的局部环。
在 ${\mathfrak m}$ 中找出一个非零因子（无需证明）。
\end{enumerate}
\end{exercise}

\begin{exercise}
\label{exercises-exercise-embedding-dim-1}
嵌入维数为 $1$ 的局部环。
假设 $(A, {\mathfrak m}, k)$ 是嵌入维数为 $1$ 的 Noether 局部环，
即
$$
\dim_k {\mathfrak m}/{\mathfrak m}^2 = 1.
$$
证明函数 $f(n) = \dim_k {\mathfrak m}^n/{\mathfrak m}^{n + 1}$
或者恒等于 $1$，或者其值依次为
$$
1, 1, \ldots, 1, 0, 0, 0, 0, 0, \ldots
$$
\end{exercise}

\begin{exercise}
\label{exercises-exercise-regular-local-dim-1}
维数为 $1$ 的正则局部环。
假设 $(A, {\mathfrak m}, k)$ 是维数为 $1$ 的正则 Noether 局部环。
回忆：这意味着 $A$ 的维数为 $1$，嵌入维数也为 $1$，即
$$
\dim_k {\mathfrak m}/{\mathfrak m}^2 = 1.
$$
任取 $x\in{\mathfrak m}$，使其在
${\mathfrak m}/{\mathfrak m}^2$ 中的类非零。
\begin{enumerate}
\item 证明对 ${\mathfrak m}$ 的每个元素 $y$，都存在整数 $n$，
使得 $y$ 可写成 $y = ux^n$，其中 $u\in A^\ast$ 是单位。
\item 证明 $x$ 在 $A$ 中是非零因子。
\item 推出 $A$ 是整环。
\end{enumerate}
\end{exercise}

\begin{exercise}
\label{exercises-exercise-nonzerodivisor-graded}
设 $(A, {\mathfrak m}, k)$ 为 Noether 局部环，其伴随分次环为
$Gr_{\mathfrak m}(A)$。
\begin{enumerate}
\item 假设 $x\in {\mathfrak m}^d$ 映到
$Gr_{\mathfrak m}(A)$ 的次数 $d$ 部分中的非零因子
$\bar x \in {\mathfrak m}^d/{\mathfrak m}^{d + 1}$。
证明 $x$ 是非零因子。
\item 假设 $A$ 的深度至少为 $1$；也就是说，假设存在非零因子
$y \in {\mathfrak m}$。在此情形中可以加强结论：只假设
$x\in {\mathfrak m}^d$ 映到 $Gr_{\mathfrak m}(A)$ 的次数 $d$
部分中的元素 $\bar x \in {\mathfrak m}^d/{\mathfrak m}^{d + 1}$，
且它在充分高的次数上是非零因子：即 $\exists N$，使得对所有
$n \geq N$，乘以 $\bar x$ 的映射
$$
{\mathfrak m}^n/{\mathfrak m}^{n + 1} \longrightarrow
{\mathfrak m}^{n + d}/{\mathfrak m}^{n + d + 1}
$$
都是单射。证明此时 $x$ 是非零因子。
\end{enumerate}
\end{exercise}

\begin{exercise}
\label{exercises-exercise-embedding-2-dim-1}
假设 $(A, {\mathfrak m}, k)$ 是维数为 $1$ 的 Noether 局部环。
再假设 $A$ 的嵌入维数为 $2$，即
$$
\dim_k {\mathfrak m}/{\mathfrak m}^2 = 2.
$$
记号：$f(n) = \dim_k {\mathfrak m}^n/{\mathfrak m}^{n + 1}$。
选取生成元 $x, y \in {\mathfrak m}$，并对某个齐次理想 $I$ 写成
$Gr_{\mathfrak m}(A) = k[\bar x, \bar y]/I$。
\begin{enumerate}
\item 证明存在齐次元素 $F\in k[\bar x, \bar y]$，使得
$I \subset (F)$，并且在所有充分高的次数上取等号。
\item 证明 $f(n) \leq n + 1$。
\item 证明若 $f(n) < n + 1$，则 $n \geq \deg(F)$。
\item 证明若 $f(n) < n + 1$，则 $f(n + 1) \leq f(n)$。
% P12-ERR-0443
\item 证明对所有 $n >> 0$，都有 $f(n) = \deg(F)$。
\end{enumerate}
\end{exercise}

\begin{exercise}
\label{exercises-exercise-CM-dim-1-embedding-dim-2}
维数为 1、嵌入维数为 2 的 Cohen--Macaulay 环。
假设 $(A, {\mathfrak m}, k)$ 是维数为 $1$ 的 Cohen--Macaulay
Noether 局部环。再假设 $A$ 的嵌入维数为 $2$，即
$$
\dim_k {\mathfrak m}/{\mathfrak m}^2 = 2.
$$
记号 $f$、$F$、$x, y\in {\mathfrak m}$、$I$ 同练习
\ref{exercises-exercise-embedding-2-dim-1}。请自由使用上面各题的结果。
\begin{enumerate}
\item 假设 $z\in {\mathfrak m}$ 的类在
${\mathfrak m}/{\mathfrak m}^2$ 中是线性形式
$\alpha \bar x + \beta \bar y \in k[\bar x, \bar y]$，并且它与
$F$ 互素。
\begin{enumerate}
\item 证明 $z$ 在 $A$ 上是非零因子。
\item 令 $d = \deg(F)$。证明对所有充分大的 $n$ 都有
${\mathfrak m}^n = z^{n + 1-d}{\mathfrak m}^{d-1}$。
（提示：先利用关于 $Gr_{\mathfrak m}(A)$ 的已知结果，证明
$z^{n + 1-d}{\mathfrak m}^{d-1} \to
{\mathfrak m}^n/{\mathfrak m}^{n + 1}$ 是满射，再使用 NAK。）
\end{enumerate}
\item $k$ 满足什么条件才能保证存在这样的 $z$？
（无需证明；这太容易了。）

\noindent
现在假设存在上述 $z$。事实证明这是无害的假设（其含义是：
为了得到下文（d）和（e）的结果，可以化归到它成立的情形）。
\item 证明对所有 $\ell \geq d$ 都有
${\mathfrak m}^\ell = z^{\ell - d + 1} {\mathfrak m}^{d-1}$。
\item 推出 $I = (F)$。
% P12-ERR-0435
\item 推出函数 $f$ 的值依次为
$$
2, 3, 4, \ldots, d-1, d, d, d, d, d, d, d, \ldots
$$
\end{enumerate}
\end{exercise}

\begin{remark}
\label{exercises-remark-CM-dim-1-embedding-dim-2}
这表明，维数为 1、嵌入维数为 2 的 Cohen--Macaulay 局部 Noether 环
应当形如 $B/FB$，其中 $B$ 是二维正则局部环。在环具有适当“良好性”
的条件下，这基本上是正确的。
\end{remark}






\section{无穷多个素数}
\label{exercises-section-many-primes}

\noindent
本节收集了一些关于无穷多个素数不可逆之环的奇特问题。

\begin{exercise}
\label{exercises-exercise-not-in-Q}
给出具有下列两个性质的有限型 ${\mathbf Z}$-代数 $R$：
\begin{enumerate}
\item 不存在环映射 $R \to {\mathbf Q}$。
\item 对每个素数 $p$，都存在极大理想 ${\mathfrak m} \subset R$，
使得 $R/{\mathfrak m} \cong {\mathbf F}_p$。
\end{enumerate}
\end{exercise}

\begin{exercise}
\label{exercises-exercise-strange-fp-1}
对 $f \in {\mathbf Z}[x, u]$，定义
$f_p(x) = f(x, x^p) \bmod p \in {\mathbf F}_p[x]$。
给出满足下列两个性质的 $f \in {\mathbf Z}[x, u]$：
\begin{enumerate}
\item 有无穷多个 $p$，使 $f_p$ 在 ${\mathbf F}_p$ 中没有零点。
% P12-ERR-0443
\item 对所有 $p >> 0$，多项式 $f_p$ 都有一次因子或二次因子。
\end{enumerate}
\end{exercise}

\begin{exercise}
\label{exercises-exercise-strange-fp-2}
对 $f \in {\mathbf Z}[x, y, u, v]$，定义
$f_p(x, y) = f(x, y, x^p, y^p) \bmod p \in {\mathbf F}_p[x, y]$。
给出一个“有趣”的 $f$，使得对所有
% P12-ERR-0443
$p >> 0$，$f_p$ 都可约。例如，$f = xv-yu$ 对应
$f_p = xy^p-x^py = xy(x^{p-1}-y^{p-1})$，这是“无趣”的；
任何仅依赖 $x, u$ 的 $f$ 也是“无趣”的；等等。
\end{exercise}

\begin{remark}
\label{exercises-remark-strange-fp}
设 $h \in {\mathbf Z}[y]$ 为次数 $d$ 的首一多项式。那么：
\begin{enumerate}
\item 映射 $A = {\mathbf Z}[x] \to B ={\mathbf Z}[y]$、
$x \mapsto h$ 是秩为 $d$ 的有限局部自由映射。
% P12-ERR-0436
\item 对所有素数 $p$，映射
$A_p = {\mathbf F}_p[x]\to B_p = {\mathbf F}_p[y]$、
$y \mapsto h(y) \bmod p$ 都是秩为 $d$ 的有限局部自由映射。
\end{enumerate}
\end{remark}

\begin{exercise}
\label{exercises-exercise-strange-fp-3}
设 $h, A, B, A_p, B_p$ 同上注。对 $f \in {\mathbf Z}[x, u]$，
定义 $f_p(x) = f(x, x^p) \bmod p \in {\mathbf F}_p[x]$。
对 $g \in {\mathbf Z}[y, v]$，定义
$g_p(y) = g(y, y^p) \bmod p \in {\mathbf F}_p[y]$。
\begin{enumerate}
\item 给出 $h$ 和 $g$ 的例子，使得不存在满足下述性质的 $f$：
% P12-ERR-0437
$$
f_p  =  Norm_{B_p/A_p}(g_p).
$$
\item 证明对上述任意 $h$ 和 $g$，都存在非零 $f$，使得对所有
$p$ 都有
$$
Norm_{B_p/A_p}(g_p)\quad\text{整除}\quad f_p .
$$
如果愿意，可以只处理 $h = y^n$ 甚至 $n = 2$ 的情形，
不过一般情形也成立。
% P12-ERR-0438
\item 讨论这与前一题组的练习 6 和 7 有何关系。
\end{enumerate}
\end{exercise}

\begin{exercise}
\label{exercises-exercise-strange-fp-unsolved}
未解决的问题。它们可能非常困难，也可能很容易。我不知道。
\begin{enumerate}
\item 是否存在 $f \in {\mathbf Z}[x, u]$，使得 $f_p$ 对无穷多个
$p$ 都不可约？（提示：是的，$f(x, u) = u - x - 1$ 会发生这种情形，
$f(x, u) = u^2 - x^2 + 1$ 也会。）
\item 设 $f \in {\mathbf Z}[x, u]$ 非零，并假设对所有充分大的
$p$，都有 $\deg_x(f_p) = dp + d'$。（换言之，
$\deg_u(f) = d$，且 $f$ 中 $u^d$ 的系数 $c$ 满足
$\deg_x(c) = d'$。）假设可以写成 $d = d_1 + d_2$ 和
$d' = d'_1 + d'_2$，其中 $d_1, d_2 > 0$ 且
$d'_1, d'_2 \geq 0$；并且对所有充分大的 $p$，存在分解
$$
f_p = f_{1, p} f_{2, p}
$$
满足 $\deg_x(f_{1, p}) = d_1p + d'_1$。% P12-ERR-0439
$f$ 是否必定来自练习 4
中的范数构造？（更准确地说，是否存在 $h$ 和 $g$，使得
% P12-ERR-0437
$Norm_{B_p/A_p}(g_p)$ 对所有
% P12-ERR-0443
$p >> 0$ 都整除 $f_p$？）
\item 与（b）类似的问题，但现在假设
$f \in {\mathbf Z}[x_1, x_2, u_1, u_2]$ 不可约，并且只假设
$f_p(x_1, x_2) = f(x_1, x_2, x_1^p, x_2^p) \bmod p$ 对所有
% P12-ERR-0443
$p >> 0$ 都可分解。
\end{enumerate}
\end{exercise}


















\section{滤过导出范畴}
\label{exercises-section-filtered-derived}

\noindent
为完成本节练习，请阅读《同调代数》第
\ref{homology-section-filtrations} 节的材料。我们说 $A$ 是
$\mathcal{A}$ 的滤过对象，是指 $A$ 带有一个滤过 $F$，
但在记号中略去它。

\begin{exercise}
\label{exercises-exercise-split-injective}
设 $\mathcal{A}$ 为 Abel 范畴，$I$ 为 $\mathcal{A}$ 的滤过对象。
假设 $I$ 上的滤过有限，并且每个 $\text{gr}^p(I)$ 都是
$\mathcal{A}$ 的内射对象。证明存在同构
$I \cong \bigoplus \text{gr}^p(I)$，使得滤过 $F^p(I)$ 对应于
% P12-ERR-0440
$\bigoplus_{p' \geq p} \text{gr}^p(I)$。
\end{exercise}

\begin{exercise}
\label{exercises-exercise-filtered-injective}
设 $\mathcal{A}$ 为 Abel 范畴，$I$ 为 $\mathcal{A}$ 的滤过对象。
假设 $I$ 上的滤过有限。证明下列条件等价：
\begin{enumerate}
\item 对任意滤过对象的实线图
$$
\xymatrix{
A \ar[r]_\alpha \ar[d] & B \ar@{-->}[ld] \\
I &
}
$$
若（\romannumeral1）$A$ 与 $B$ 上的滤过有限，并且
% P12-ERR-0441
（\romannumeral2）$\text{gr}(\alpha)$ 为单射，则存在使图交换的虚线箭头。
\item 每个 $\text{gr}^p I$ 都是内射对象。
\end{enumerate}
\end{exercise}

\noindent
注意，给定有限滤过对象间的态射 $\alpha : A \to B$，
% P12-ERR-0441
$\text{gr}(\alpha)$ 为单射，等价于 $\alpha$ 在范畴
$\text{Fil}(\mathcal{A})$ 中是{\it 严格单态射}。事实上，成为单态射
意味着 $\Ker(\alpha) = 0$，而严格意味着这还蕴含
$\Ker(\text{gr}(\alpha)) = 0$。见《同调代数》引理
\ref{homology-lemma-characterize-strict}。（我们只把“单射”一词用于
Abel 范畴中的态射，尽管它在任何具有核的加性范畴中都有意义。）
% P12-ERR-0442
上面的练习说明了下述定义的合理性。

\begin{definition}
\label{exercises-definition-injective-filtered}
设 $\mathcal{A}$ 为 Abel 范畴，$I$ 为 $\mathcal{A}$ 的滤过对象。
假设 $I$ 上的滤过有限。若每个 $\text{gr}^p(I)$ 都是
$\mathcal{A}$ 的内射对象，则称 $I$ {\it 滤过内射}。
\end{definition}

\noindent
作如下定义，以免处处重复“具有有限滤过”。

\begin{definition}
\label{exercises-definition-finite-filtration-category}
设 $\mathcal{A}$ 为 Abel 范畴。记 $\text{Fil}(\mathcal{A})$
的全子范畴 {\it $\text{Fil}^f(\mathcal{A})$} 为这样的对象
$A \in \Ob(\text{Fil}(\mathcal{A}))$ 所成的范畴：其滤过有限。
\end{definition}

\begin{exercise}
\label{exercises-exercise-inject-into-injective}
设 $\mathcal{A}$ 为 Abel 范畴，并假设 $\mathcal{A}$ 有足够多内射对象。
设 $A$ 为 $\text{Fil}^f(\mathcal{A})$ 的对象。证明存在严格单态射
$\alpha : A \to I$，把 $A$ 嵌入
$\text{Fil}^f(\mathcal{A})$ 的滤过内射对象 $I$。
\end{exercise}

\begin{definition}
\label{exercises-definition-filtered-quasi-isomorphism}
设 $\mathcal{A}$ 为 Abel 范畴，并设
$\alpha : K^\bullet \to L^\bullet$ 为
$\text{Fil}(\mathcal{A})$ 中复形的态射。若对每个
$p \in \mathbf{Z}$，态射
$\text{gr}^p(K^\bullet) \to \text{gr}^p(L^\bullet)$
都是拟同构，则称 $\alpha$ 为{\it 滤过拟同构}。
\end{definition}

\begin{definition}
\label{exercises-definition-filtered-acyclic}
设 $\mathcal{A}$ 为 Abel 范畴，并设 $K^\bullet$ 为
$\text{Fil}^f(\mathcal{A})$ 中的复形。若对每个
$p \in \mathbf{Z}$，复形 $\text{gr}^p(K^\bullet)$ 都无环，
则称 $K^\bullet$ {\it 滤过无环}。
\end{definition}

\begin{exercise}
\label{exercises-exercise-filtered-quasi-isomorphism}
设 $\mathcal{A}$ 为 Abel 范畴，并设
$\alpha : K^\bullet \to L^\bullet$ 为
$\text{Fil}^f(\mathcal{A})$ 中下有界复形的态射。
（注意上标 $f$。）证明下列条件等价：
\begin{enumerate}
\item $\alpha$ 是滤过拟同构；
\item 对每个 $p \in \mathbf{Z}$，映射
$\alpha : F^pK^\bullet \to F^pL^\bullet$ 是拟同构；
\item 对每个 $p \in \mathbf{Z}$，映射
$\alpha : K^\bullet/F^pK^\bullet \to L^\bullet/F^pL^\bullet$
是拟同构；
\item $\alpha$ 的锥（见《导出范畴》定义
\ref{derived-definition-cone}）是滤过无环复形。
\end{enumerate}
此外，证明若 $\alpha$ 是滤过拟同构，则 $\alpha$ 也是通常的拟同构。
\end{exercise}

\begin{exercise}
\label{exercises-exercise-injective-resolution}
设 $\mathcal{A}$ 为 Abel 范畴，并假设 $\mathcal{A}$ 有足够多内射对象。
设 $A$ 为 $\text{Fil}^f(\mathcal{A})$ 的对象。证明存在
$\text{Fil}^f(\mathcal{A})$ 的复形 $I^\bullet$，以及态射
$A[0] \to I^\bullet$，使得
\begin{enumerate}
\item 每个 $I^p$ 都滤过内射；
\item 当 $p < 0$ 时 $I^p = 0$；
\item $A[0] \to I^\bullet$ 是滤过拟同构。
\end{enumerate}
\end{exercise}

\begin{exercise}
\label{exercises-exercise-injective-resolution-complex}
设 $\mathcal{A}$ 为 Abel 范畴，并假设 $\mathcal{A}$ 有足够多内射对象。
设 $K^\bullet$ 为 $\text{Fil}^f(\mathcal{A})$ 的对象所成的
下有界复形。证明存在滤过拟同构
$\alpha : K^\bullet \to I^\bullet$，其中 $I^\bullet$ 是
$\text{Fil}^f(\mathcal{A})$ 的复形，其各项 $I^n$ 滤过内射，
并且它下有界。事实上，可以选取 $\alpha$，使每个 $\alpha^n$
都是严格单态射。
\end{exercise}

\begin{exercise}
\label{exercises-exercise-morphisms-lift}
设 $\mathcal{A}$ 为 Abel 范畴。考察实线图
$$
\xymatrix{
K^\bullet \ar[r]_\alpha \ar[d]_\gamma & L^\bullet \ar@{-->}[dl]^\beta \\
I^\bullet
}
$$
其中各对象是 $\text{Fil}^f(\mathcal{A})$ 中的复形。假设
$K^\bullet$、$L^\bullet$、$I^\bullet$ 下有界，并且每个
$I^n$ 都是滤过内射对象。再假设 $\alpha$ 是滤过拟同构。
\begin{enumerate}
\item 存在复形的映射 $\beta$，使图在同伦意义下交换。
\item 若 $\alpha$ 在每个次数上都是严格单态射，则可以找到
使图严格交换的 $\beta$。
\end{enumerate}
\end{exercise}

\begin{exercise}
\label{exercises-exercise-acyclic-is-zero}
设 $\mathcal{A}$ 为 Abel 范畴。
% P12-ERR-0444
设 $K^\bullet$、$K^\bullet$ 为 $\text{Fil}^f(\mathcal{A})$ 中的复形。
假设
\begin{enumerate}
\item $K^\bullet$ 下有界且滤过无环；
\item $I^\bullet$ 下有界，并且由滤过内射对象组成。
\end{enumerate}
则任意态射 $K^\bullet \to I^\bullet$ 都同伦于零。
\end{exercise}

\begin{exercise}
\label{exercises-exercise-morphisms-equal-up-to-homotopy}
设 $\mathcal{A}$ 为 Abel 范畴。考察实线图
$$
\xymatrix{
K^\bullet \ar[r]_\alpha \ar[d]_\gamma & L^\bullet \ar@{-->}[dl]^{\beta_i} \\
I^\bullet
}
$$
其中各对象是 $\text{Fil}^f(\mathcal{A})$ 中的复形。
假设 $K^\bullet$、$L^\bullet$、$I^\bullet$ 下有界，每个 $I^n$
都是滤过内射对象，并且 $\alpha$ 是滤过拟同构。任何两个使图在
同伦意义下交换的态射 $\beta_1, \beta_2$ 彼此同伦。
\end{exercise}







\section{正则函数}
\label{exercises-section-regular-functions}

\begin{exercise}
\label{exercises-exercise-extra-function}
考察由 $\mathbf{C}^2$ 中方程 $t^2 = s^5 + 8$ 给出的仿射曲线
$X$，坐标为 $s, t$。设 $x \in X$ 为坐标 $(1, 3)$ 的点。
令 $U = X \setminus \{x\}$。证明 $U$ 上存在一个正则函数，
它不是 $\mathbf{C}^2$ 上正则函数的限制；也就是说，它不是
$s$ 与 $t$ 的多项式在 $U$ 上的限制。
\end{exercise}

\begin{exercise}
\label{exercises-exercise-no-extra-function}
设 $n \geq 2$，并设 $E \subset \mathbf{C}^n$ 为有限子集。
证明 $\mathbf{C}^n \setminus E$ 上的任意正则函数都是多项式。
\end{exercise}

\begin{exercise}
\label{exercises-exercise-cone}
设 $X \subset \mathbf{C}^n$ 为仿射簇。若
$x = (a_1, \ldots, a_n) \in X$ 且 $\lambda \in \mathbf{C}$
蕴含 $(\lambda a_1, \ldots, \lambda a_n) \in X$，则称 $X$
为{\it 锥}。当然，若 $\mathfrak p \subset \mathbf{C}[x_1, \ldots, x_n]$
是由 $x_1, \ldots, x_n$ 的齐次多项式生成的素理想，则仿射簇
$X = V(\mathfrak p) \subset \mathbf{C}^n$ 是锥。证明反过来，
与锥对应的素理想 $\mathfrak p \subset \mathbf{C}[x_1, \ldots, x_n]$
也可由 $x_1, \ldots, x_n$ 的齐次多项式生成。
\end{exercise}

\begin{exercise}
\label{exercises-exercise-extra-function-cone}
给出一个为锥的仿射簇 $X \subset \mathbf{C}^n$（见练习
\ref{exercises-exercise-cone}），以及 $U = X \setminus \{(0, \ldots, 0)\}$
上的正则函数 $f$，使它不是 $\mathbf{C}^n$ 上多项式函数的限制。
\end{exercise}

\begin{exercise}
\label{exercises-exercise-regular-functions}
本练习考察非代数闭域上的正则函数。设 $k$ 为域。
设 $Z \subset k^n$ 为 Zariski 局部闭子集；即存在理想
$I \subset J \subset k[x_1, \ldots, x_n]$，使得
$$
Z = \{a \in k^n \mid
f(a) = 0\ \forall\ f \in I,\ \exists\ g \in J,\ g(a) \not = 0\}.
$$
若对每个 $z \in Z$，都存在 Zariski 开邻域
$z \in U \subset Z$ 和多项式 $f, g \in k[x_1, \ldots, x_n]$，
使得对所有 $u \in U$ 都有 $g(u) \not = 0$，并且对所有
$u \in U$ 都有 $\varphi(u) = f(u)/g(u)$，则称函数
$\varphi : Z \to k$ {\it 正则}。
\begin{enumerate}
\item 若 $k = \bar k$ 且 $Z = k^n$，证明正则函数由多项式给出。
（只在以前未见过这一论证时完成此题。）
\item 若 $k$ 有限，证明（a）每个函数 $\varphi$ 都正则；
（b）正则函数环在 $k$ 上有限维。（如果愿意，可以取
$Z = k^n$，甚至取 $n = 1$。）
\item 若 $k = \mathbf{R}$，给出 $Z = \mathbf{R}$ 上不由多项式
给出的正则函数。
\item 若 $k = \mathbf{Q}_p$，给出 $Z = \mathbf{Q}_p$ 上不由
多项式给出的正则函数。
\end{enumerate}
\end{exercise}






\section{层}
\label{exercises-section-sheaves}

\noindent
范畴 $\mathcal{C}$ 的态射 $f : X \to Y$ 称为{\it 单态射}，
是指对每一对态射 $a, b : W \to X$，都有
$f \circ a = f \circ b \Rightarrow a = b$。集合范畴中的单态射
就是集合的单射。

\begin{exercise}
\label{exercises-exercise-mono-sheaves-sets}
% P12-ERR-0445
仔细证明：集合层的映射是单态射（在集合层范畴中），当且仅当它在
所有茎上诱导的映射都是单射。
\end{exercise}

\noindent
范畴 $\mathcal{C}$ 的态射 $f : X \to Y$ 称为{\it 同构}，
是指存在态射 $g : Y \to X$，使得
$f \circ g = \text{id}_Y$ 且 $g \circ f = \text{id}_X$。
集合范畴中的同构就是集合的双射。

\begin{exercise}
\label{exercises-exercise-isomorphism-sheaves-sets}
仔细证明：集合层的映射是同构（在集合层范畴中），当且仅当它在
所有茎上诱导的映射都是双射。
\end{exercise}

\noindent
范畴 $\mathcal{C}$ 的态射 $f : X \to Y$ 称为{\it 满态射}，
是指对每一对态射 $a, b : Y \to Z$，都有
$a \circ f = b \circ f \Rightarrow a = b$。集合范畴中的满态射
就是集合的满射。

\begin{exercise}
\label{exercises-exercise-epi-sheaves-sets}
仔细证明：集合层的映射是满态射（在集合层范畴中），当且仅当它在
所有茎上诱导的映射都是满射。
\end{exercise}

\begin{exercise}
\label{exercises-exercise-adjoint-push-pull}
设 $f : X \to Y$ 为拓扑空间的映射。证明{\bf 集合}层的推前 $f_\ast$
与拉回 $f^{-1}$ 构成一对伴随函子。
\end{exercise}

\begin{exercise}
\label{exercises-exercise-j-shriek}
设 $j : U \to X$ 为开浸入。证明：
\begin{enumerate}
\item 拉回 $j^{-1} : \Sh(X) \to \Sh(U)$ 有左伴随
$j_{!} : \Sh(U) \to \Sh(X)$，称为{\it 以空集延拓}。
\item 刻画 $j_{!}({\mathcal G})$ 的茎，其中
$\mathcal{G} \in \Sh(U)$。
\item 拉回 $j^{-1} : \textit{Ab}(X) \to \textit{Ab}(U)$ 有左伴随
$j_{!} : \textit{Ab}(U) \to \textit{Ab}(X)$，称为{\it 零延拓}。
\item 刻画 $j_{!}({\mathcal G})$ 的茎，其中
$\mathcal{G} \in \textit{Ab}(U)$。
\end{enumerate}
注意，零延拓与以空集延拓不同！
\end{exercise}

\begin{exercise}
\label{exercises-exercise-not-locally-generated-by-sections}
令 $X = \mathbf{R}$，取通常拓扑。令
$\mathcal{O}_X = \underline{\mathbf{Z}/2\mathbf{Z}}_X$。
令 $i : Z = \{0\} \to X$ 为包含映射，并令
$\mathcal{O}_Z = \underline{\mathbf{Z}/2\mathbf{Z}}_Z$。
证明下列各点（前三点直接来自定义；若定义尚不清楚，应借此阐明）：
\begin{enumerate}
\item $i_*\mathcal{O}_Z$ 是摩天层。
\item 存在典范满射
$\underline{\mathbf{Z}/2\mathbf{Z}}_X \to
i_*\underline{\mathbf{Z}/2\mathbf{Z}}_Z$。记其核为
$\mathcal{I} \subset \mathcal{O}_X$。
\item $\mathcal{I}$ 是 $\mathcal{O}_X$ 的理想层。
\item $X$ 上的层 $\mathcal{I}$ 不可能由截面局部生成
（含义见《模》定义 \ref{modules-definition-locally-generated}）。
\end{enumerate}
\end{exercise}

\begin{exercise}
\label{exercises-exercise-quotient-j-shriek-Z}
设 $X$ 为拓扑空间，${\mathcal F}$ 为 $X$ 上的 Abel 群层。
证明 ${\mathcal F}$ 是某个（可能非常大的）层直和的商，其中每一项
都形如
$$
j_{!}(\underline{{\mathbf Z}}_U)
$$
这里 $U \subset X$ 开，$\underline{{\mathbf Z}}_U$ 表示 $U$
上取值为 ${\mathbf Z}$ 的常值层。
\end{exercise}

\begin{remark}
\label{exercises-remark-direct-sum-stalk-abelian}
设 $X$ 为拓扑空间。在 Abel 群层范畴中，一族层
$\{{\mathcal F}_i\}_{i\in I}$ 的直和，是预层
% P12-ERR-0446
$U \mapsto \oplus {\mathcal F}_i(U)$ 的层化。因此直和在点
$x$ 处的茎，是各 ${\mathcal F}_i$ 在 $x$ 处之茎的直和。
\end{remark}

\begin{exercise}
\label{exercises-exercise-product-over-points}
设 $X$ 为拓扑空间，并给定由 $x \in X$ 标号的一族 Abel 群
$A_x$。证明规则 $U \mapsto \prod_{x \in U} A_x$ 配以显然的限制映射，
定义 Abel 群层 $\mathcal{G}$。用例子说明，通常对
$x \in X$ 并没有 $\mathcal{G}_x = A_x$。
\end{exercise}

\begin{exercise}
\label{exercises-exercise-modified-product-over-points}
设 $X$、$A_x$、$\mathcal{G}$ 同练习
\ref{exercises-exercise-product-over-points}。设 $\mathcal{B}$ 为 $X$
拓扑的一组基；见《拓扑》定义 \ref{topology-definition-base}。
对 $U \in \mathcal{B}$，设 $A_U$ 为子群
$A_U \subset \mathcal{G}(U) = \prod_{x \in U} A_x$。假设当
$U \subset V$ 且 $U, V \in \mathcal{B}$ 时，限制映射把
$A_V$ 映入 $A_U$。对开集 $U \subset X$，令
$$
\mathcal{F}(U) =
\left\{
(s_x)_{x \in U}
\middle|
\begin{matrix}
\text{ 对每个 }x\text{ 属于 }U\text{，存在 } V \in \mathcal{B} \\
x \in V \subset U\text{，使得 } (s_y)_{y \in V} \in A_V
\end{matrix}
\right\}
$$
证明 $\mathcal{F}$ 定义 $X$ 上的 Abel 群层。用例子说明，
通常对 $U \in \mathcal{B}$ 并没有 $\mathcal{F}(U) = A_U$。
\end{exercise}

\begin{exercise}
\label{exercises-exercise-exact-but-not-a-stalk-functor}
给出拓扑空间 $X$ 和函子
$$
F : \Sh(X) \longrightarrow \textit{Sets}
$$
的例子，使该函子正合（与有限积、等化子、有限余积及余等化子交换；
见《范畴》 \ref{categories-section-exact-functor} 节），但不存在
$x \in X$，使 $F$ 同构于茎函子
$\mathcal{F} \mapsto \mathcal{F}_x$。
\end{exercise}




\section{概形}
\label{exercises-section-schemes}

\noindent
令 $LRS$ 为局部环空间范畴。仿射概形是 $LRS$ 中与某个形如
$(\Spec(A), {\mathcal O}_{\Spec(A)})$ 的序对在 $LRS$ 中同构的对象。
概形是 $LRS$ 的对象 $(X, {\mathcal O}_X)$，并满足：每个点
$x\in X$ 都有开邻域 $U \subset X$，使得序对
$(U, {\mathcal O}_X|_U)$ 是仿射概形。

\begin{exercise}
\label{exercises-exercise-one-point}
找出一个不是概形的 $1$-点局部环空间。
\end{exercise}

\begin{exercise}
\label{exercises-exercise-two-points}
假设概形 $X$ 的底拓扑空间有 2 个点。证明 $X$ 是仿射概形。
\end{exercise}

\begin{exercise}
\label{exercises-exercise-discrete-finite-set-points}
假设概形 $X$ 的底拓扑空间是有限离散集。证明 $X$ 是仿射概形。
\end{exercise}

\begin{exercise}
\label{exercises-exercise-three-points}
证明存在具有三个点的非仿射概形。
\end{exercise}

\begin{exercise}
\label{exercises-exercise-quasi-compact-closed-point}
假设 $X$ 是非空拟紧概形。证明 $X$ 有闭点。
\end{exercise}

\begin{remark}
\label{exercises-remark-open-immersion}
若 $(X, {\mathcal O}_X)$ 是环空间，且 $U \subset X$ 是开子集，
则 $(U, {\mathcal O}_X|_U)$ 是环空间。记
${\mathcal O}_U = {\mathcal O}_X|_U$。
% P12-ERR-0447
存在典范的环空间态射
$$
j : (U, {\mathcal O}_U) \longrightarrow (X, {\mathcal O}_X).
$$
若 $(X, {\mathcal O}_X)$ 是局部环空间，则
$(U, {\mathcal O}_U)$ 也是局部环空间，且 $j$ 是局部环空间态射。
若 $(X, {\mathcal O}_X)$ 是概形，则 $(U, {\mathcal O}_U)$
也是概形，且 $j$ 是概形态射。称 $(U, {\mathcal O}_U)$ 为
$(X, {\mathcal O}_X)$ 的{\it 开子概形}，并称 $j$ 为{\it 开浸入}。
更一般地，任何态射
$j' : (V, {\mathcal O}_V) \to (X, {\mathcal O}_X)$，若与上述某个
态射 $j : (U, {\mathcal O}_U) \to (X, {\mathcal O}_X)$
{\it 同构}，也称为开浸入。
\end{remark}

\begin{exercise}
\label{exercises-exercise-open-affine-not-affine}
给出仿射概形 $(X, {\mathcal O}_X)$ 以及开集 $U \subset X$ 的例子，
使得
% P12-ERR-0448
$(U, {\mathcal O}_X|U)$ 不是仿射概形。
\end{exercise}

\begin{exercise}
\label{exercises-exercise-morphism-does-not-extend}
% P12-ERR-0449
给出一对仿射概形 $(X, {\mathcal O}_X)$、$(Y, {\mathcal O}_Y)$，
$X$ 的开子概形 $(U, {\mathcal O}_X|_U)$，以及概形态射
$(U, {\mathcal O}_X|_U) \to (Y, {\mathcal O}_Y)$，使它不能延拓为
概形态射 $(X, {\mathcal O}_X) \to (Y, {\mathcal O}_Y)$。
\end{exercise}

\begin{exercise}
\label{exercises-exercise-closed-subscheme-does-not-extend}
（这相当困难。）
% P12-ERR-0450
给出概形 $X$、开子概形 $U \subset X$ 和闭子概形 $Z \subset U$
的例子，使 $Z$ 不能延拓为 $X$ 的闭子概形。
\end{exercise}

\begin{exercise}
\label{exercises-exercise-not-morphism-schemes}
给出概形 $X$、域 $K$ 以及环空间态射
$\Spec(K) \to X$ 的例子，使该态射不是概形态射。
\end{exercise}

\begin{exercise}
\label{exercises-exercise-just-kidding}
完成 \cite[Chapter II]{H} 第 1 节与第 2 节中的所有练习……
\ \ 开玩笑！
\end{exercise}

\begin{definition}
\label{exercises-definition-integral}
若概形 $X$ 非空，并且对每个非空仿射开集 $U \subset X$，环
$\Gamma(U, \mathcal{O}_X) = \mathcal{O}_X(U)$ 都是整环，
则称 $X$ {\it 整}。
\end{definition}

\begin{exercise}
\label{exercises-exercise-morphism-integral-schemes-surjective-stalks-not-closed}
给出{\it 整}概形的态射 $f : X \to Y$ 的例子，使得对所有
$x\in X$，诱导映射
${\mathcal O}_{Y, f(x)} \to {\mathcal O}_{X, x}$ 都是满射，
但 $f$ 不是闭浸入。
\end{exercise}

\begin{exercise}
\label{exercises-exercise-fibre-product-affines-not-affine}
给出纤维积 $X \times_S Y$ 的例子，使 $X$ 与 $Y$ 都仿射，
但 $X \times_S Y$ 不仿射。
\end{exercise}

\begin{remark}
\label{exercises-remark-separated-base-fibre-product-affines-affine}
事实证明，当 $S$ 分离时，这不可能发生。你知道为什么吗？
\end{remark}

\begin{exercise}
\label{exercises-exercise-not-geometrically-integral}
% P12-ERR-0451
给出概形 $V$ 的例子，使它是 ${\mathbf Q}$ 上有限型的一维整概形，
但 $\Spec({\mathbf C}) \times_{\Spec({\mathbf Q})} V$ 不整。
\end{exercise}

\begin{exercise}
\label{exercises-exercise-not-geometrically-reduced}
% P12-ERR-0452
给出概形 $V$ 的例子，使它是域 $k$ 上有限型的一维整概形，
但对某个有限域扩张 $k'/k$，
$\Spec(k') \times_{\Spec(k)} V$ 不是既约的。
\end{exercise}

\begin{remark}
\label{exercises-remark-affine-dimension}
% P12-ERR-0453
若你的概形是仿射的，则其维数等于底层环的 Krull 维数。
因此可以利用上学期的结果来计算维数。
\end{remark}






\section{态射}
\label{exercises-section-morphisms}

\noindent
一个重要问题是：给定态射 $\pi : X \to S$，
该态射是否有截面或有理截面。下面是一些例题。

\begin{exercise}
\label{exercises-exercise-no-section}
考虑概形态射
$$
\pi :
X = \Spec(\mathbf{C}[x, t, 1/xt])
\longrightarrow
S = \Spec(\mathbf{C}[t]).
$$
\begin{enumerate}
\item 证明不存在态射 $\sigma : S \to X$
使得 $\pi \circ \sigma = \text{id}_S$。
\item 证明存在非空开集 $U \subset S$ 以及态射
$\sigma : U \to X$，使得 $\pi \circ \sigma = \text{id}_U$。
\end{enumerate}
\end{exercise}

\begin{exercise}
\label{exercises-exercise-no-rational-section}
考虑概形态射
$$
\pi :
X = \Spec(\mathbf{C}[x, t]/(x^2 + t))
\longrightarrow
S = \Spec(\mathbf{C}[t]).
$$
证明不存在非空开集 $U \subset S$ 以及态射
$\sigma : U \to X$，使得 $\pi \circ \sigma = \text{id}_U$。
\end{exercise}

\begin{exercise}
\label{exercises-exercise-has-rational-section}
设 $A, B, C \in \mathbf{C}[t]$ 为非零多项式。考虑概形态射
$$
\pi :
X = \Spec(\mathbf{C}[x, y, t]/(A + Bx^2 + Cy^2))
\longrightarrow
S = \Spec(\mathbf{C}[t]).
$$
证明存在非空开集 $U \subset S$ 以及态射
$\sigma : U \to X$，使得 $\pi \circ \sigma = \text{id}_U$。
（提示：在形式上，对某个 $d$，把
$x = X/Z$、$y = Y/Z$ 写成次数 $\leq d$ 的某些
$X, Y, Z \in \mathbf{C}[t]$，并求出它们解此方程的条件。
% P12-ERR-0454
然后用维数理论证明：若 $d >> 0$，便可找到解此方程的非零
$X, Y, Z$。）
\end{exercise}

\begin{remark}
\label{exercises-remark-tsen}
练习 \ref{exercises-exercise-has-rational-section}
是“Tsen 定理”的一个特例。
练习 \ref{exercises-exercise-no-section-curve} 表明，此方法仅限于低次方程
（当底与纤维的维数均为 1 时，即二次曲线）。
\end{remark}

\begin{exercise}
\label{exercises-exercise-no-section-curve}
考虑概形态射
$$
\pi :
X = \Spec(\mathbf{C}[x, y, t]
/(1 + t x^3 + t^2 y^3))
\longrightarrow
S = \Spec(\mathbf{C}[t])
$$
证明不存在非空开集 $U \subset S$ 以及态射
$\sigma : U \to X$，使得 $\pi \circ \sigma = \text{id}_U$。
\end{exercise}

\begin{exercise}
\label{exercises-exercise-no-section-surface}
考虑概形
$$
X = \Spec(\mathbf{C}[\{x_i\}_{i = 1}^{8}, s, t]
/(1 + s x_1^3 + s^2 x_2^3 +
t x_3^3 + st x_4^3 + s^2t x_5^3 +
t^2 x_6^3 + st^2 x_7^3 + s^2t^2 x_8^3))
$$
以及
$$
S = \Spec(\mathbf{C}[s, t])
$$
和概形态射
$$
\pi : X \longrightarrow S
$$
证明不存在非空开集 $U \subset S$ 以及态射
$\sigma : U \to X$，使得 $\pi \circ \sigma = \text{id}_U$。
\end{exercise}

\begin{exercise}
\label{exercises-exercise-for-number-theorists}
（给数论学家。）给出闭子概形
$$
Z \subset \Spec\left({\mathbf Z}[x, \frac{1 }{ x(x-1)(2x-1)}]\right)
$$
使得态射 $Z \to \Spec({\mathbf Z})$ 有限且满。
\end{exercise}

\begin{exercise}
\label{exercises-exercise-quasi-section}
若你不喜欢数论，可以尝试如下变体。考察
$$
\Spec\left({\mathbf F}_p[t, x, \frac{1 }{ x(x-t)(tx-1)}]\right)
\longrightarrow
\Spec({\mathbf F}_p[t])
$$
并试着找出上方概形的一个闭子概形，使其有限满射到下方概形。
% P12-ERR-0455
（这里取有限基域有理论上的原因；不过在这个特例中可能并非必要。）
\end{exercise}

\begin{remark}
\label{exercises-remark-interpretation-skolem-noether}
在练习 \ref{exercises-exercise-for-number-theorists} 与
\ref{exercises-exercise-quasi-section} 中，给定态射 $f : X \to S$，
其中 $S$ 是 Dedekind 环 $A$ 的谱，而 $f$ 平坦且满，并具有
几何不可约的泛纤维。在两种情形中，结论都是：存在有限满态射
$S' \to S$，使得基变换 $X_{S'} \to S'$ 确有截面。
事实证明，若 $A$ 优良，其极大理想处的剩余域是有限域的代数扩张，
且当 $A'$ 是 $A$ 在其分式域的有限扩张中的整闭包时
$\Pic(A')$ 为挠群，则上述结论成立；见 \cite{MB1}。
例如，$A = \mathbf{Z}$ 或 $A = \mathbf{F}_p[x]$，或它们的有限扩张，
都满足这一条件。然而，存在一个 Dedekind 环 $A$，其极大素理想处的
剩余域有限，而其 Picard 群含有非挠元；见 \cite{Goldman}。
对这种环的谱，上述结论为假。
练习 \ref{exercises-exercise-no-quasi-section} 给出了该结论不成立的几何例子。
\end{remark}

\begin{exercise}
\label{exercises-exercise-no-quasi-section}
% P12-ERR-0456
证明存在 $f \in \mathbf{C}[x, t]$，它不被任意
$\alpha \in \mathbf{C}$ 所对应的 $t - \alpha$ 整除，并且不存在
任何 $Z \subset \Spec(\mathbf{C}[x, t, 1/f])$
有限满射到 $\Spec(\mathbf{C}[t])$。
（我认为 $f(x, t) = (xt - 2)(x - t + 3)$ 可行。
不过我觉得，要证明任一候选式具有所需性质并不容易。）
\end{exercise}

\begin{exercise}
\label{exercises-exercise-finite}
设 $A \to B$ 为有限型环映射。假设
$\Spec(B) \to \Spec(A)$ 可分解为某个 $n$ 下的闭浸入
$\Spec(B) \to \mathbf{P}^n_A$。
证明 $A \to B$ 是有限环映射，即 $B$ 作为 $A$-模是有限的。
提示：若 $A$ Noether（请直接假设这一点），则可利用如下事实论证：
对每个闭子概形 $Z \subset \mathbf{P}^n_A$ 和每个
$i \in \mathbf{Z}$，$H^i(Z, \mathcal{O}_Z)$ 都是有限 $A$-模。
\end{exercise}

\begin{exercise}
\label{exercises-exercise-projective-finite}
设 $k$ 为代数闭域。设 $f : X \to Y$ 为射影簇的态射，
并且对每个闭点 $y \in Y$，$f^{-1}(\{y\})$ 都是有限集。
证明 $f$ 作为概形态射是有限的。
提示：（a）有限性是局部性质；（b）试着归约到练习
\ref{exercises-exercise-finite}；（c）利用闭浸入
$X \to \mathbf{P}^n_k$ 得到 $Y$ 上的闭浸入
$X \to \mathbf{P}^n_Y$。
\end{exercise}






\section{切空间}
\label{exercises-section-tangent-space}

\begin{definition}
\label{exercises-definition-dual-numbers}
对任意环 $R$，用 $R[\epsilon]$ 表示{\it 对偶数}环。
作为 $R$-模，它以 $1$、$\epsilon$ 为基自由。
环结构由规定 $\epsilon^2 = 0$ 得到。
\end{definition}

\begin{exercise}
\label{exercises-exercise-tangent-space-Zariski}
设 $f : X \to S$ 为概形态射。设 $x \in X$ 为一点，并令
$s = f(x)$。考虑实线所成的交换图
$$
\xymatrix{
\Spec(\kappa(x)) \ar[r] \ar[dr] \ar@/^1pc/[rr] &
\Spec(\kappa(x)[\epsilon]) \ar@{.>}[r] \ar[d]&
X \ar[d] \\
&
\Spec(\kappa(s)) \ar[r] &
S
}
$$
其中弯曲箭头是 $\Spec(\kappa(x))$ 到 $X$ 的典范态射。
% P12-ERR-0457
若 $\kappa(x) = \kappa(s)$，证明使该图交换的虚线箭头集合
与线性映射集合
$$
\Hom_{\kappa(x)}(
\frac{\mathfrak m_x}{\mathfrak m_x^2 + \mathfrak m_s\mathcal{O}_{X, x}},
\kappa(x))
$$
一一对应。换言之，请描述这样的双射。
（当 $\kappa(x) \supset \kappa(s)$ 为可分代数扩张时，
这一结论更一般地也成立。）
\end{exercise}

\begin{definition}
\label{exercises-definition-tangent-space}
设 $f : X \to S$ 为概形态射，$x \in X$。
我们把练习 \ref{exercises-exercise-tangent-space-Zariski} 中的虚线箭头集合
称为{\it $X$ 在 $S$ 上的切空间}，并记作 $T_{X/S, x}$。
此空间的一个元素称为 $X/S$ 在 $x$ 处的{\it 切向量}。
\end{definition}

\begin{exercise}
\label{exercises-exercise-simple-push-out}
对任意域 $K$，证明图
$$
\xymatrix{
\Spec(K) \ar[r] \ar[d] & \Spec(K[\epsilon_1]) \ar[d] \\
\Spec(K[\epsilon_2]) \ar[r] &
\Spec(K[\epsilon_1, \epsilon_2]/(\epsilon_1\epsilon_2))
}
$$
是概形范畴中的推出图。
（这里仍有 $\epsilon_i^2 = 0$。）
\end{exercise}

\begin{exercise}
\label{exercises-exercise-tangent-space-vectors-space}
设 $f : X \to S$ 为概形态射，$x \in X$。
利用练习 \ref{exercises-exercise-simple-push-out} 以及适当的态射
$$
\Spec(K[\epsilon])
\longrightarrow
\Spec(K[\epsilon_1, \epsilon_2]/(\epsilon_1\epsilon_2)).
$$
定义切向量的加法。类似地，定义切向量的标量乘法（这更容易）。
证明你的构造使 $T_{X/S, x}$ 成为 $\kappa(x)$-向量空间。
\end{exercise}

\begin{exercise}
\label{exercises-exercise-compute-TS}
设 $k$ 为域。考虑结构态射
$f : X = \mathbf{A}^1_k \to \Spec(k) = S$。
\begin{enumerate}
\item 设 $x \in X$ 为闭点。$T_{X/S, x}$ 的维数是多少？
\item 设 $\eta \in X$ 为泛点。$T_{X/S, \eta}$ 的维数是多少？
\item 现在把 $X$ 看作 $\Spec(\mathbf{Z})$ 上的概形。
$T_{X/\mathbf{Z}, x}$ 与 $T_{X/\mathbf{Z}, \eta}$ 的维数分别是多少？
\end{enumerate}
\end{exercise}

\begin{remark}
\label{exercises-remark-tangent-space-relative}
练习 \ref{exercises-exercise-compute-TS} 说明了为何必须考虑 $X$ 在 $S$
上的切空间，才能得到良好的概念。
\end{remark}

\begin{exercise}
\label{exercises-exercise-compute-TS-field}
考虑概形态射
$$
f : X = \Spec(\mathbf{F}_p(t))
\longrightarrow
\Spec(\mathbf{F}_p(t^p)) = S
$$
计算 $X/S$ 在 $X$ 的唯一一点处的切空间。
这不是很奇怪吗？若取与
$\mathbf{F}_p[t^p] \to \mathbf{F}_p[t]$ 对应的概形态射，
你认为会发生什么？
\end{exercise}

\begin{exercise}
\label{exercises-exercise-compute-TS-cusp}
设 $k$ 为域。对
$X = \Spec(k[x, y]/(x^2 - y^3))$，计算 $X/k$ 在点
$x = (0, 0)$ 处的切空间。
\end{exercise}

\begin{exercise}
\label{exercises-exercise-map-tangent-spaces}
设 $f : X \to Y$ 为 $S$ 上的概形态射。设 $x \in X$，
并令 $y = f(x)$。假设自然映射
$\kappa(y) \to \kappa(x)$ 为双射。利用定义证明 $f$
诱导自然线性映射
$$
\text{d}f : T_{X/S, x} \longrightarrow T_{Y/S, y}.
$$
在 $\kappa(x) = \kappa(s)$ 的情形下，通过练习
\ref{exercises-exercise-tangent-space-Zariski} 将它与局部环上的情形对应起来。
\end{exercise}

\begin{exercise}
\label{exercises-exercise-Jacobian}
设 $k$ 为代数闭域。设
% P12-ERR-0458
\begin{eqnarray*}
f : \mathbf{A}_k^n & \longrightarrow & \mathbf{A}^m_k \\
(x_1, \ldots, x_n) & \longmapsto & (f_1(x_i), \ldots, f_m(x_i))
\end{eqnarray*}
为 $k$ 上的概形态射。它由 $n$ 个变量中的 $m$ 个多项式
$f_1, \ldots, f_m$ 给出。考虑矩阵
% P12-ERR-0459
$$
A = \left( \frac{\partial f_j}{\partial x_i} \right)
$$
设 $x \in \mathbf{A}^n_k$ 为闭点，并令 $y =  f(x)$。
证明切空间上的映射
$T_{\mathbf{A}^n_k/k, x} \to T_{\mathbf{A}^m_k/k, y}$
由矩阵 $A$ 在点 $x$ 处的取值给出。
\end{exercise}













\section{拟凝聚层}
\label{exercises-section-quasi-coherent}

\begin{definition}
\label{exercises-definition-quasi-coherent}
设 $X$ 为概形。若对每个仿射开集
$\Spec(R) = U \subset X$，$\mathcal{O}_X$-模层
$\mathcal{F}$ 的限制 $\mathcal{F}|_U$ 都可写成某个
$R$-模 $M$ 所对应的 $\widetilde M$，则称它
是{\it 拟凝聚的}。
\end{definition}

\noindent
% P12-ERR-0460
只需在 $X$ 的一个仿射开覆盖的成员上检验此条件即可。
更多结果见概形篇第 \ref{schemes-section-quasi-coherent} 节。

\begin{definition}
\label{exercises-definition-specialization}
设 $X$ 为拓扑空间，$x, x' \in X$。
当且仅当 $x \in \overline{\{x'\}}$ 时，称 $x$ 是 $x'$ 的
{\it 特化}。
\end{definition}

\begin{exercise}
\label{exercises-exercise-quasi-coherent-specialization-points}
设 $X$ 为概形，$x, x' \in X$，并设 $\mathcal{F}$ 为
$\mathcal{O}_X$-模的拟凝聚层。假设（a）$x$ 是 $x'$ 的特化，
并且（b）$\mathcal{F}_{x'} \not = 0$。
证明 $\mathcal{F}_x \not = 0$。
\end{exercise}

\begin{exercise}
\label{exercises-exercise-O-module-specialization-points}
找出概形 $X$、点 $x, x' \in X$ 以及
$\mathcal{O}_X$-模层 $\mathcal{F}$ 的例子，使得
（a）$x$ 是 $x'$ 的特化，并且（b）
$\mathcal{F}_{x'} \not = 0$ 而 $\mathcal{F}_x = 0$。
\end{exercise}

\begin{definition}
\label{exercises-definition-Noetherian-scheme}
当且仅当对每个点 $x \in X$，都存在仿射开集
$\Spec(R) = U \subset X$，其中 $R$ 为 Noether 环时，
称概形 $X$ {\it 局部 Noether}。若概形局部 Noether 且拟紧，
则称其为 {\it Noether 概形}。
\end{definition}

\noindent
若 $X$ 局部 Noether，则 $X$ 的任意仿射开集都是 Noether 环的谱；
见性质篇引理 \ref{properties-lemma-locally-Noetherian}。

\begin{definition}
\label{exercises-definition-coherent}
设 $X$ 为局部 Noether 概形，并设 $\mathcal{F}$ 为
$\mathcal{O}_X$-模的拟凝聚层。若对每个点 $x \in X$，
都存在仿射开集 $\Spec(R) = U \subset X$，使得
$\mathcal{F}|_U$ 同构于某个有限 $R$-模 $M$ 所对应的
$\widetilde M$，则称 $\mathcal{F}$ {\it 凝聚}。
\end{definition}

\begin{exercise}
\label{exercises-exercise-extend-quasi-coherent}
设 $X = \Spec(R)$ 为仿射概形。
\begin{enumerate}
\item 设 $f \in R$，并设 $\mathcal{G}$ 为开子概形
$D(f)$ 上的拟凝聚 $\mathcal{O}_{D(f)}$-模层。
证明存在拟凝聚 $\mathcal{O}_X$-模层 $\mathcal{F}$，
使 $\mathcal{G} = \mathcal{F}|_{D(f)}$。
\item 设 $I \subset R$ 为理想。令 $i : Z \to X$ 为与 $I$
对应的 $X$ 的闭子概形。设 $\mathcal{G}$ 为闭子概形 $Z$ 上的
拟凝聚 $\mathcal{O}_Z$-模层。证明存在拟凝聚
$\mathcal{O}_X$-模层 $\mathcal{F}$，使得
$\mathcal{G} = i^*\mathcal{F}$。（为什么这很显然？）
\item 假设 $R$ Noether。设 $f \in R$，并设 $\mathcal{G}$
为开子概形 $D(f)$ 上的凝聚 $\mathcal{O}_{D(f)}$-模层。
证明存在凝聚 $\mathcal{O}_X$-模层 $\mathcal{F}$，
使 $\mathcal{G} = \mathcal{F}|_{D(f)}$。
\end{enumerate}
\end{exercise}

\begin{remark}
\label{exercises-remark-extend-off-open}
若 $U \to X$ 是拟紧浸入，则 $U$ 上任意拟凝聚层都是 $X$ 上
某个拟凝聚层的限制。若 $X$ 为 Noether 概形且 $U \subset X$
为开集，则 $U$ 上任意凝聚层都是 $X$ 上某个凝聚层的限制。
当然，上面的练习更容易，不应使用这些一般事实。
\end{remark}

















\section{Proj 与射影概形}
\label{exercises-section-proj}

\begin{exercise}
\label{exercises-exercise-graded-ring-specified-result}
给出分次环 $S$ 的例子，使得
\begin{enumerate}
\item $\text{Proj}(S)$ 仿射且非空；
\item $\text{Proj}(S)$ 整且非空，但对任意 $n\geq 0$、
任意环 $A$，均不同构于 ${\mathbf P}^n_A$。
\end{enumerate}
\end{exercise}

\begin{exercise}
\label{exercises-exercise-nonconstant-morphism-proj}
给出 $\Spec({\mathbf C})$ 上概形的非常值态射
${\mathbf P}^1_{\mathbf C} \to {\mathbf P}^5_{\mathbf C}$ 的例子。
\end{exercise}

\begin{exercise}
\label{exercises-exercise-isomorphism-P1-conic}
给出概形同构的例子
$$
{\mathbf P}^1_{\mathbf C} \to
\text{Proj}({\mathbf C}[X_0, X_1, X_2]/(X_0^2 + X_1^2 + X_2^2))
$$
\end{exercise}

\begin{exercise}
\label{exercises-exercise-family-special-fibre-different}
给出概形态射
$f : X \to {\mathbf A}^1_{\mathbf C} = \Spec({\mathbf C}[T])$
的例子，使 $f$ 在 $t \in {\mathbf A}^1_{\mathbf C}$ 上的
（概形论）纤维 $X_t$ 满足：
（a）当 $t$ 是不等于 $0$ 的闭点时，它同构于
${\mathbf P}^1_{\mathbf C}$；
（b）当 $t = 0$ 时，它不同构于 ${\mathbf P}^1_{\mathbf C}$。
我们把 $X_0$ 称为该态射的{\it 特殊纤维}。
这可以用许多许多种方式做到。试着给出还满足下列每一项附加限制的例子
（除非不可能）：
\begin{enumerate}
\item 能否使特殊纤维射影？
\item 能否使特殊纤维不可约且射影？
\item 能否使特殊纤维整且射影？
\item 能否使特殊纤维光滑且射影？
\item 能否使 $f$ 为平坦态射？
这仅表示：对每个仿射开集 $\Spec(A) \subset X$，诱导环映射
$\mathbf{C}[t] \to A$ 都平坦；在本例中，这表示 $t$ 的任意
非零多项式在 $A$ 上都是非零因子。
\item 能否使 $f$ 为平坦射影态射？
\item 能否使 $f$ 平坦、射影且特殊纤维既约？
\item 能否使 $f$ 平坦、射影且特殊纤维不可约？
\item 能否使 $f$ 平坦、射影且特殊纤维整？
\end{enumerate}
若把 ${\mathbf P}^1_{\mathbf C}$ 换成 ${\mathbf C}$ 上的另一个簇，
你认为会发生什么？
（取决于你要求上述哪个变体，这可能会变得非常困难。）
\end{exercise}

\begin{exercise}
\label{exercises-exercise-affine-onto-projective-space}
设 $n \geq 1$ 为任意正整数。给出以仿射概形 $X$ 为源的满态射
$X \to {\mathbf P}^n_{\mathbf C}$ 的例子。
\end{exercise}

\begin{exercise}
\label{exercises-exercise-morphism-proj}
$\text{Proj}$ 之间的映射。设 $R$ 与 $S$ 为分次环。
假设有环映射
$$
\psi : R \to S
$$
以及整数 $e \geq 1$，使得对所有 $d \geq 0$，
$\psi(R_d) \subset S_{de}$。
（按我们的约定，除非 $e = 1$，这不是分次环同态。）
\begin{enumerate}
\item 对哪些元素 $\mathfrak p \in \text{Proj}(S)$，
$\text{Proj}(R)$ 中有定义良好的对应点？换言之，找出适当的开集
$U \subset \text{Proj}(S)$，使 $\psi$ 定义连续映射
$r_\psi : U \to \text{Proj}(R)$。
\item 给出 $U \not = \text{Proj}(S)$ 的例子。
\item 给出 $U = \text{Proj}(S)$ 的例子。
\item （不要把它写下来。）说服自己：连续映射
$U \to \text{Proj}(R)$ 典范地带有层上的映射，从而 $r_\psi$
是概形态射：
$$
\text{Proj}(S) \supset U \longrightarrow \text{Proj}(R).
$$
\item 若
$R = \bigoplus_{d \geq 0} S_{de}$（作为分次环，把
$S_e$、$S_{2e}$ 等分别放在次数 $1$、$2$ 等处），且 $\psi$
是包含映射，你能对该映射说些什么？
\end{enumerate}
\end{exercise}

\noindent
{\bf 记号。}设 $R$ 为如上的分次环，$n \geq 0$ 为整数，
并令 $X = \text{Proj}(R)$。则 $X$ 上存在唯一的拟凝聚
${\mathcal O}_X$-模 ${\mathcal O}_X(n)$，使得对每个正次数齐次元
$f \in R$，${\mathcal O}_X |_{D_{+}(f)}$ 是与
$R_{(f)} = (R_f)_0$-模 $(R_f)_n$ 相伴的拟凝聚层
% P12-ERR-0461
（$ = $在 $R_f = R[1/f]$ 中次数为 $n$ 的齐次元）。
见 \cite[page 116+]{H}。注意，存在自然映射
$$
{\mathcal O}_X(n_1) \otimes_{{\mathcal O}_X} {\mathcal O}_X(n_2)
\longrightarrow
{\mathcal O}_X(n_1 + n_2)
$$

\begin{exercise}
\label{exercises-exercise-pathologies-proj}
$\text{Proj}$ 中的病态现象。
给出如上的 $R$ 的例子，使得
\begin{enumerate}
\item ${\mathcal O}_X(1)$ 不是可逆 ${\mathcal O}_X$-模。
\item ${\mathcal O}_X(1)$ 可逆，但自然映射
${\mathcal O}_X(1) \otimes_{{\mathcal O}_X} {\mathcal O}_X(1) \to
{\mathcal O}_X(2)$ 不是同构。
\end{enumerate}
\end{exercise}

\begin{exercise}
\label{exercises-exercise-finitely-many-points-in-affine}
设 $S$ 为分次环，并令 $X = \text{Proj}(S)$。
证明 $X$ 的任意有限点集都包含在某个标准仿射开集中。
\end{exercise}

\begin{exercise}
\label{exercises-exercise-prepare-glueing}
设 $S$ 为分次环，$X = \text{Proj}(S)$。
设 $Z, Z' \subset X$ 为两个闭子概形，并设
$\varphi : Z \to Z'$ 为同构。假设 $Z \cap Z' = \emptyset$。
证明：对任意 $z \in Z$，都存在仿射开集 $U \subset X$，使得
$z \in U$、$\varphi(z) \in U$ 且
$\varphi(Z \cap U) = Z' \cap U$。
（提示：使用练习 \ref{exercises-exercise-finitely-many-points-in-affine}
以及类似于概形篇引理
\ref{schemes-lemma-standard-open-two-affines} 的方法。）
\end{exercise}







\section{从射影直线出发的态射}
\label{exercises-section-from-P1}

\noindent
本节研究从 $\mathbf{P}^1$ 到射影概形的态射。

\begin{exercise}
\label{exercises-exercise-from-generic-point}
设 $k$ 为域，并设 $k[t] \subset k(t)$ 为多项式环到其分式域的包含。
设 $X$ 为 $k$ 上有限型概形。证明：对任意 $k$ 上的态射
$$
\varphi : \Spec(k(t)) \longrightarrow X
$$
都存在非零 $f \in k[t]$ 以及 $k$ 上的态射
$\psi : \Spec(k[t, 1/f]) \to X$，使 $\varphi$ 为复合
$$
\Spec(k(t)) \longrightarrow \Spec(k[t, 1/f]) \longrightarrow X
$$
\end{exercise}

\begin{exercise}
\label{exercises-exercise-from-generic-point-Pn}
设 $k$ 为域，并设 $k[t] \subset k(t)$ 为多项式环到其分式域的包含。
证明：对任意 $k$ 上的态射
$$
\varphi : \Spec(k(t)) \longrightarrow \mathbf{P}^n_k
$$
都存在 $k$ 上的态射
$\psi : \Spec(k[t]) \to \mathbf{P}^n_k$，使 $\varphi$ 为复合
$$
\Spec(k(t)) \longrightarrow \Spec(k[t]) \longrightarrow \mathbf{P}^n_k
$$
提示：$\varphi$ 的像包含在某个 $i$ 所对应的标准开集
$D_+(T_i)$ 中；然后证明可以“清除分母”。
\end{exercise}

\begin{exercise}
\label{exercises-exercise-from-generic-point-projective}
设 $k$ 为域，并设 $k[t] \subset k(t)$ 为多项式环到其分式域的包含。
设 $X$ 为 $k$ 上的射影概形。证明：对任意 $k$ 上的态射
$$
\varphi : \Spec(k(t)) \longrightarrow X
$$
都存在 $k$ 上的态射 $\psi : \Spec(k[t]) \to X$，
使 $\varphi$ 为复合
$$
\Spec(k(t)) \longrightarrow \Spec(k[t]) \longrightarrow X
$$
提示：使用练习 \ref{exercises-exercise-from-generic-point-Pn}。
\end{exercise}

\begin{exercise}
\label{exercises-exercise-from-generic-point-P1-to-projective}
设 $k$ 为域，$X$ 为 $k$ 上的射影概形。
令 $K$ 为 $\mathbf{P}^1_k$ 的函数域（见下方提示）。
证明：对任意 $k$ 上的态射
$$
\varphi : \Spec(K) \longrightarrow X
$$
都存在 $k$ 上的态射 $\psi : \mathbf{P}^1_k \to X$，
使 $\varphi$ 为复合
$$
\Spec(k(t)) \longrightarrow \mathbf{P}^1_k \longrightarrow X
$$
提示：对仿射开覆盖
$\mathbf{P}^1_k = D_+(T_0) \cup D_+(T_1)$ 的两个部分分别使用
练习 \ref{exercises-exercise-from-generic-point-projective}；利用
$D_+(T_0)$ 是多项式环的谱，并且 $K$ 是该多项式环的分式域。
\end{exercise}







\section{从曲面到曲线的态射}
\label{exercises-section-from-surfaces-to-curves}

\begin{exercise}
\label{exercises-exercise-points-projective-space}
设 $R$ 为环。设 $R \to k$ 为从 $R$ 到某个域的映射。
设 $n \geq 0$。证明
$$
\Mor_{\Spec(R)}(\Spec(k), \mathbf{P}^n_R)
=
(k^{n + 1} \setminus \{0\})/k^*
$$
其中 $k^*$ 通过标量乘法作用在 $k^{n + 1}$ 上。
此后，用 $(x_0 : \ldots : x_n)$ 表示与元素
$(x_0, \ldots, x_n) \in k^{n + 1} \setminus \{0\}$ 的等价类
对应的态射 $\Spec(k) \to \mathbf{P}^n_k$。
\end{exercise}

\begin{exercise}
\label{exercises-exercise-curve-projective-plane}
设 $k$ 为域，并设 $Z \subset \mathbf{P}^2_k$ 为不可约既约闭子概形。
证明下列三种情形之一成立：（a）$Z$ 是闭点；或
% P12-ERR-0462
（b）存在次数 $> 0$ 的齐次不可约多项式
$F \in k[X_0, X_1, X_2]$，使 $Z = V_{+}(F)$；或
（c）$Z = \mathbf{P}^2_k$。（提示：考察一个标准仿射开集。）
\end{exercise}

\begin{exercise}
\label{exercises-exercise-bezout}
% P12-ERR-0463
设 $k$ 为域。设 $Z_1, Z_2 \subset \mathbf{P}^2_k$
为形如 $V_{+}(F)$ 的不可约闭子概形，其中
$F_i \in k[X_0, X_1, X_2]$ 是次数 $> 0$ 的齐次不可约多项式。
证明 $Z_1 \cap Z_2$ 非空。
（提示：用维数理论估计
$k[X_0, X_1, X_2]/(F_1, F_2)$ 在 $0$ 处局部环的维数。）
\end{exercise}

\begin{exercise}
\label{exercises-exercise-no-nonconstant-morphism-proj}
证明不存在 $\Spec(\mathbf{C})$ 上概形的非常值态射
$\mathbf{P}^2_{\mathbf{C}} \to \mathbf{P}^1_{\mathbf{C}}$。
这里，{\it 常值态射}是指其像仅为一个点的态射。
（提示：若态射不是常值的，考察 $0$ 与 $\infty$ 上的纤维，
并论证它们必然相交，从而得到矛盾。）
\end{exercise}

\begin{exercise}
\label{exercises-exercise-nonconstant-morphism}
设 $k$ 为域。假设 $X \subset \mathbf{P}^3_k$ 是由单个齐次方程
$F \in k[X_0, X_1, X_2, X_3]$ 给出的闭子概形。换言之，
$$
X = \text{Proj}(k[X_0, X_1, X_2, X_3]/(F)) \subset \mathbf{P}^3_k
$$
如课程中所述。假设
$$
F = X_0 G + X_1 H
$$
其中 $G, H \in k[X_0, X_1, X_2, X_3]$ 是正次数齐次多项式。
证明：若 $X_0, X_1, G, H$ 没有公共零点，则存在
$\Spec(k)$ 上概形的非常值态射
$$
X \longrightarrow \mathbf{P}^1_k
$$
它在域值点上（见练习 \ref{exercises-exercise-points-projective-space}），
只要 $x_0$ 或 $x_1$ 非零，就表现为
$(x_0 : x_1 : x_2 : x_3) \mapsto (x_0 : x_1)$。
\end{exercise}





\section{可逆层}
\label{exercises-section-invertible-sheaves}

\begin{definition}
\label{exercises-definition-invertible-sheaf}
设 $X$ 为局部环空间。$X$ 上的一个
{\it 可逆 ${\mathcal O}_X$-模}，是 ${\mathcal O}_X$-模层
${\mathcal L}$，使得每一点都有开邻域 $U \subset X$，
且 ${\mathcal L}|_U$ 作为 ${\mathcal O}_U$-模同构于
${\mathcal O}_U$。若 ${\mathcal L}$ 作为
${\mathcal O}_X$-模同构于 ${\mathcal O}_X$，则称它平凡。
\end{definition}

\begin{exercise}
\label{exercises-exercise-general-facts-invertible}
一般事实。
\begin{enumerate}
\item 证明概形 $X$ 上的可逆 ${\mathcal O}_X$-模是拟凝聚的。
\item 假设 $X\to Y$ 是局部环空间的态射，并且
% P12-ERR-0464
${\mathcal L}$ 是可逆 ${\mathcal O}_Y$-模。
% P12-ERR-0465
证明 $f^\ast {\mathcal L}$ 是可逆 ${\mathcal O}_X$ 模。
\end{enumerate}
\end{exercise}

\begin{exercise}
\label{exercises-exercise-invertible-algebra}
代数。
\begin{enumerate}
\item 证明仿射概形 $\Spec(A)$ 上的可逆
${\mathcal O}_X$-模对应于 $A$-模 $M$，后者
（i）有限，（ii）射影，（iii）局部自由且秩为 1，
从而（iv）平坦，并且（v）有限表现。
% P12-ERR-0466
（可以随意引用上学期课程中的结论，或代数教材中的结论。）
\item 假设 $A$ 为整环，$M$ 是（a）中那样的模。
% P12-ERR-0467
证明 $M$ 作为 $A$-模同构于某个理想 $I \subset A$，
且对每个素理想 ${\mathfrak p}$，$IA_{\mathfrak p}$ 都是主理想。
\end{enumerate}
\end{exercise}

\begin{definition}
\label{exercises-definition-invertible-module}
设 $R$ 为环。所谓{\it 可逆模 $M$}，是指 $R$-模 $M$
满足：$\widetilde M$ 是 $R$ 的谱上的可逆层。
若 $M$ 作为 $R$-模满足 $M \cong R$，则称该模{\it 平凡}。
\end{definition}

\noindent
换言之，当且仅当 $M$ 满足以下全部条件时，它才可逆：
它平坦、有限表现、射影，并且局部自由且秩为 1。
（当然，仅要求它局部自由且秩为 1 就足够了。）

\begin{exercise}
\label{exercises-exercise-simple-examples-invertible}
简单例子。
\begin{enumerate}
\item
\label{exercises-item-affine-line}
设 $k$ 为域，$A = k[x]$。
证明 $X = \Spec(A)$ 仅有平凡的可逆
${\mathcal O}_X$-模。换言之，证明每个可逆 $A$-模都自由且秩为 1。
\item
\label{exercises-item-affine-line-with-0-and-1-identified}
设 $A$ 为环
$$
A = \{ f\in k[x] \mid f(0) = f(1) \}.
$$
证明除非 $k = {\mathbf F}_2$，否则存在非平凡可逆 $A$-模。
（提示：把 $\Spec(A)$ 视为在
${\mathbf A}^1_k = \Spec(k[x])$ 中将 $0$ 与 $1$ 粘合。）
\item
\label{exercises-item-affine-line-with-cusp}
对环 $A = k[x^2, x^3] \subset k[x]$ 回答
（\ref{exercises-item-affine-line-with-0-and-1-identified}）中的同一问题
（但现在 $k = {\mathbf F}_2$ 也可行）。
\end{enumerate}
\end{exercise}

\begin{exercise}
\label{exercises-exercise-higher-dimension-invertible}
更高维情形。
\begin{enumerate}
\item 证明二维仿射空间上的每个可逆层都平凡。更精确地，设
${\mathbf A}^2_k = \Spec(k[x, y])$，其中 $k$ 为域。
证明 ${\mathbf A}^2_k$ 上的每个可逆层都平凡。
（提示：一种做法是考察对应模 $M$，再考察
$M \otimes_{k[x, y]} k(x)[y]$，然后用练习
\ref{exercises-exercise-simple-examples-invertible}
（\ref{exercises-item-affine-line}）找出它的一个生成元；之后你仍需思考。
% P12-ERR-0468
另一种做法是使用练习 \ref{exercises-exercise-invertible-algebra}，
并利用我们关于多项式环中理想的知识：高度一的素理想由一个
不可约多项式生成；之后你仍需思考。）
\item 证明二维仿射空间的任意开子概形上的每个可逆层都平凡。
更精确地，设 $U \subset {\mathbf A}^2_k$ 为开子概形，其中 $k$
为域。证明 $U$ 上的每个可逆层都平凡。提示：证明 $U$ 上每个
可逆层都可延拓为 ${\mathbf A}^2_k$ 上的可逆层。这并不容易，
但可在 \cite{H} 中找到。
\item 找出域上去掉顶点的锥面上的一个非平凡可逆层。
更精确地，设 $k$ 为域，$C = \Spec(k[x, y, z]/(xy-z^2))$。
令 $U = C \setminus \{ (x, y, z) \}$。找出 $U$ 上的一个
非平凡可逆层。提示：计算 $U$ 上可逆层的同构类群，可能比
仅找出一个更容易。注意，$U$ 被开集
$\Spec(k[x, y, z, 1/x]/(xy-z^2))$ 与
$\Spec(k[x, y, z, 1/y]/(xy-z^2))$ 覆盖，
而它们是“容易”处理的。
\end{enumerate}
\end{exercise}

\begin{definition}
\label{exercises-definition-picard-group}
设 $X$ 为局部环空间。$X$ 的{\it Picard 群}是可逆
$\mathcal{O}_X$-模的同构类集合 $\Pic(X)$，其加法由张量积给出。
见模篇定义 \ref{modules-definition-pic}。
对环 $R$，令 $\Pic(R) = \Pic(\Spec(R))$。
\end{definition}

\begin{exercise}
\label{exercises-exercise-traverso}
设 $R$ 为环。
\begin{enumerate}
\item 证明：若 $R$ 是 Noether 正规整环，则
$\Pic(R) = \Pic(R[t])$。[提示：存在映射 $R[t] \to R$，
$t \mapsto 0$，它是映射 $R \to R[t]$ 的左逆。
因此只需证明：若可逆 $R[t]$-模 $M$ 满足
$M/tM \cong R$，则它自由且秩为 $1$。令 $K$ 为 $R$ 的分式域。
选取平凡化 $K[t] \to M \otimes_{R[t]} K[t]$；由练习
\ref{exercises-exercise-simple-examples-invertible}（\ref{exercises-item-affine-line}），
这是可能的。调整它，使其与上面 $M/tM$ 的平凡化相容。
证明它事实上是 $M$ 在 $R[t]$ 上的一个平凡化
（正规性在此处起作用）。]
\item 设 $k$ 为域。证明
$\Pic(k[x^2, x^3, t]) \not = \Pic(k[x^2, x^3])$。
\end{enumerate}
\end{exercise}














\section{{\v C}ech 上同调}
\label{exercises-section-cech-cohomology}

\begin{exercise}
\label{exercises-exercise-cech-cohomology}
{\v C}ech 上同调。这里 $k$ 为域。
\begin{enumerate}
\item 设 $X$ 为带有开覆盖
${\mathcal U} : X = U_1 \cup U_2$ 的概形，其中
$U_1 = \Spec(k[x])$、$U_2 =  \Spec(k[y])$，
$U_1 \cap U_2 = \Spec(k[z, 1/z])$，而开浸入
$U_1 \cap U_2 \to U_1$ 与 $U_1 \cap U_2 \to U_2$
分别由 $x \mapsto z$ 与 $y \mapsto z$ 决定
（我确实就是这个意思）。
（我们已在讲课中见过这样的 $X$ 存在；它是原点加倍的仿射直线。）
计算 ${\check H}^1({\mathcal U}, {\mathcal O})$；
例如，给出它作为 $k$-向量空间的一组基。
\item 对
${\check H}^1({\mathcal U}, {\mathcal O})$ 中的每个元素，
构造 ${\mathcal O}_X$-模层的正合列
$$
0 \to {\mathcal O}_X \to E \to {\mathcal O}_X \to 0
$$
使边界 $\delta(1) \in {\check H}^1({\mathcal U},
{\mathcal O})$ 等于给定元素。
（问题的一部分是赋予这句话意义；另见下文。
抽象地证明这样的对象必然存在也可以。）
\end{enumerate}
\end{exercise}

\begin{definition}
\label{exercises-definition-delta}
（delta 的定义。）假设
$$
0 \to {\mathcal F}_1 \to {\mathcal F}_2 \to {\mathcal F}_3 \to 0
$$
是任意拓扑空间 $X$ 上 Abel 层的短正合列。边界映射
$\delta : H^0(X, {\mathcal F}_3) \to {\check H}^1(X, {\mathcal F}_1)$
定义如下。取元素 $\tau \in H^0(X, {\mathcal F}_3)$。
选取开覆盖 ${\mathcal U} : X = \bigcup_{i\in I} U_i$，使对每个
$i$，都存在截面 $\tilde \tau_i \in {\mathcal F}_2$，
它提升 $\tau$ 在 $U_i$ 上的限制。然后考虑赋值
$$
(i_0, i_1) \longmapsto
\tilde \tau_{i_0}|_{U_{i_0i_1}} - \tilde \tau_{i_1}|_{U_{i_0i_1}}.
$$
这是 {\v C}ech 复形
${\check C}^\ast({\mathcal U}, {\mathcal F}_2)$ 中显然的 1-余边界。
但注意到（把 ${\mathcal F}_1$ 看作 ${\mathcal F}_2$ 的子层），
右端总是 $U_{i_0i_1}$ 上 ${\mathcal F}_1$ 的截面。
% P12-ERR-0469
因此可见，该赋值在复形
${\check C}^\ast({\mathcal U}, {\mathcal F}_2)$ 中定义了一个
1-上链。按照定义，这个 1-上链的上同调类就是
{\it $\delta(\tau)$}。
\end{definition}




\section{上同调}
\label{exercises-section-cohomology}

\begin{exercise}
\label{exercises-exercise-cohomology-not-zero}
令 $X = \mathbf{R}$，赋予通常的 Euclid 拓扑。
仅使用上同调的形式性质（函子性与长正合上同调列），证明存在
$X$ 上的层 $\mathcal{F}$，使 $H^1(X, \mathcal{F})$ 非零。
\end{exercise}

\begin{exercise}
\label{exercises-exercise-mayer-vietoris}
设拓扑空间 $X = U \cup V$ 写成两个开集的并。
则有 Mayer--Vietoris 长正合列
$$
0 \to
H^0(X, \mathcal{F}) \to
H^0(U, \mathcal{F}) \oplus H^0(V, \mathcal{F}) \to
H^0(U \cap V, \mathcal{F}) \to
H^1(X, \mathcal{F}) \to \ldots
$$
内射层的什么性质对于构造 Mayer--Vietoris 长正合列至关重要？
为什么该性质成立？
\end{exercise}

\begin{exercise}
\label{exercises-exercise-cohomology-two-point-space}
设 $X$ 为拓扑空间。
\begin{enumerate}
\item 若 $X$ 至多有 $2$ 个点，证明对 $i > 0$，
$H^i(X, \mathcal{F})$ 为零。
\item 若 $X$ 有 $3$ 个点，又如何？
\end{enumerate}
\end{exercise}

\begin{exercise}
\label{exercises-exercise-cohomology-spec-local-ring}
设 $X$ 为局部环的谱。证明：对任意 Abel 群层
$\mathcal{F}$ 和任意 $i > 0$，$H^i(X, \mathcal{F})$ 为零。
\end{exercise}

\begin{exercise}
\label{exercises-exercise-affine-morphism}
设 $f : X \to Y$ 为概形的仿射态射。证明：对任意拟凝聚
$\mathcal{O}_X$-模 $\mathcal{F}$，
$H^i(X, \mathcal{F}) = H^i(Y, f_*\mathcal{F})$。
你可以随意对 $X$ 与 $Y$ 加上一些条件，并使用如下相合性：
对分离概形的仿射开覆盖与拟凝聚层，{\v C}ech 上同调与上同调相合。
\end{exercise}

\begin{exercise}
\label{exercises-exercise-affine-morphism-to-Pn}
设 $A$ 为环，并设
$\mathbf{P}^n_A = \text{Proj}(A[T_0, \ldots, T_n])$
为 $A$ 上的射影空间。

令
$\mathbf{A}^{n + 1}_A = \Spec(A[T_0, \ldots, T_n])$，并令
$$
U = \bigcup\nolimits_{i = 0, \ldots, n} D(T_i) \subset \mathbf{A}^{n + 1}_A
$$
为闭浸入 $0 : \Spec(A) \to \mathbf{A}^{n + 1}_A$ 的像的补集。
构造仿射满态射
$$
f : U \longrightarrow \mathbf{P}^n_A
$$
并证明
$f_*\mathcal{O}_U =
\bigoplus_{d \in \mathbf{Z}} \mathcal{O}_{\mathbf{P}^n_A}(d)$。
更一般地，证明对分次 $A[T_0, \ldots, T_n]$-模 $M$ 有
$$
f_*(\widetilde{M}|_U) =
\bigoplus\nolimits_{d \in \mathbf{Z}} \widetilde{M(d)}
$$
其中左端的 $\widetilde{M}$ 是 $\mathbf{A}^{n + 1}_A$ 上
与 $M$ 相伴的拟凝聚层，而右端的 $\widetilde{M(d)}$
是与分次模 $M(d)$ 相伴的拟凝聚层。
\end{exercise}

\begin{exercise}
\label{exercises-exercise-compute-Pn}
设 $A$ 为环，并设
$\mathbf{P}^n_A = \text{Proj}(A[T_0, \ldots, T_n])$
为 $A$ 上的射影空间。仔细计算 $\mathbf{P}^n_A$ 上结构层的
Serre 扭曲 $\mathcal{O}_{\mathbf{P}^n_A}(d)$ 的上同调。
你可以使用 {\v C}ech 上同调，以及分离概形的仿射开覆盖与
拟凝聚层所具有的 {\v C}ech 上同调和上同调相合性。
\end{exercise}

\begin{exercise}
\label{exercises-exercise-cohomology-hypersurface}
设 $A$ 为环，并设
$\mathbf{P}^n_A = \text{Proj}(A[T_0, \ldots, T_n])$
为 $A$ 上的射影空间。设 $F \in A[T_0, \ldots, T_n]$
为次数 $d$ 的齐次多项式。令 $X \subset \mathbf{P}^n_A$
为与 $A[T_0, \ldots, T_n]$ 的分次理想 $(F)$ 对应的闭子概形。
关于 $H^i(X, \mathcal{O}_X)$，你能说些什么？
\end{exercise}

\begin{exercise}
\label{exercises-exercise-characterize-finite-product-fields}
设环 $R$ 满足：对任意左正合函子
$F : \text{Mod}_R \to \textit{Ab}$，对 $i > 0$ 都有
$R^iF = 0$。证明 $R$ 是有限多个域的乘积。
\end{exercise}







\section{更多上同调}
\label{exercises-section-more-cohomology}

\begin{exercise}
\label{exercises-exercise-cohomology-coordinate-axes}
设 $k$ 为域。设 $X \subset \mathbf{P}^n_k$ 为“坐标十字”。
具体地，令 $X$ 由齐次方程
$$
T_i T_j = 0\text{ 对 }i > j > 0
$$
定义；照例写
$\mathbf{P}^n_k = \text{Proj}(k[T_0, \ldots, T_n])$。
换言之，$X$ 是与多项式环的商
$k[T_0, \ldots, T_n]/(T_iT_j; i > j > 0)$ 对应的闭子概形。
对所有 $i$ 计算 $H^i(X, \mathcal{O}_X)$。
提示：使用 {\v C}ech 上同调。
\end{exercise}

\begin{exercise}
\label{exercises-exercise-compute-cohomology-punctured}
设 $A$ 为环，$I = (f_1, \ldots, f_t)$ 为 $A$ 的有限生成理想。
令 $U \subset \Spec(A)$ 为 $V(I)$ 的补集。
对任意 $A$-模 $M$，写出一个 $A$-模复形
（用 $A$、$f_1, \ldots, f_t$、$M$ 表示），
使其上同调群给出 $H^n(U, \widetilde{M})$。
\end{exercise}

\begin{exercise}
\label{exercises-exercise-compute-cohomology-affine-space-punctured}
设 $k$ 为域。令 $U \subset \mathbf{A}^d_k$ 为
$\mathbf{A}^d_k$ 中闭点 $0$ 的补集。
对所有 $n$ 计算 $H^n(U, \mathcal{O}_U)$。
\end{exercise}

\begin{exercise}
\label{exercises-exercise-find-curve-genus-one-hundred}
设 $k$ 为域。显式找出 $k$ 上射影、维数为 $1$ 的概形 $X$，
使 $H^0(X, \mathcal{O}_X) = k$ 且
$\dim_k H^1(X, \mathcal{O}_X) = 100$。
\end{exercise}

\begin{exercise}
\label{exercises-exercise-degree-2-cover}
设 $f : X \to Y$ 为次数 $2$ 的有限局部自由态射。
假设 $X$ 与 $Y$ 都是整概形，并且 $2$ 在 $Y$ 的结构层中可逆，
即 $2 \in \Gamma(Y, \mathcal{O}_Y)$ 可逆。
证明 $\mathcal{O}_Y$-模映射
$$
f^\sharp : \mathcal{O}_Y \longrightarrow f_*\mathcal{O}_X
$$
有左逆，即存在 $\mathcal{O}_Y$-模映射
$\tau : f_*\mathcal{O}_X \to \mathcal{O}_Y$，使得
$\tau \circ f^\sharp = \text{id}$。
由此推出 $H^n(Y, \mathcal{O}_Y) \to H^n(X, \mathcal{O}_X)$
为单射\footnote{确实存在整概形之间次数为 $2$ 的有限局部自由态射
$X \to Y$，使得映射
$H^1(Y, \mathcal{O}_Y) \to H^1(X, \mathcal{O}_X)$ 不是单射。}。
\end{exercise}

\begin{exercise}
\label{exercises-exercise-pic}
设 $X$ 为概形（或局部环空间）。规则
$U \mapsto \mathcal{O}_X(U)^*$ 定义一个群层，记为
$\mathcal{O}_X^*$。简要说明为何 $X$ 的 Picard 群
（定义 \ref{exercises-definition-picard-group}）等于
$H^1(X, \mathcal{O}_X^*)$。
\end{exercise}

\begin{exercise}
\label{exercises-exercise-pic-nontrivial}
给出具有非平凡 $\Pic(X)$ 的仿射概形 $X$ 的例子。
利用练习 \ref{exercises-exercise-pic} 推出：对任意这样的 $X$，
$H^1(X, -)$ 都不是零函子。
\end{exercise}

\begin{exercise}
\label{exercises-exercise-kill-cohomology-complement}
设 $A$ 为环，$I = (f_1, \ldots, f_t)$ 为 $A$ 的有限生成理想，
并令 $U \subset \Spec(A)$ 为 $V(I)$ 的补集。
给定拟凝聚 $\mathcal{O}_{\Spec(A)}$-模 $\mathcal{F}$
以及满足 $p > 0$ 的 $\xi \in H^p(U, \mathcal{F})$，
证明存在 $n > 0$，使得对 $i = 1, \ldots, t$ 有
$f_i^n \xi = 0$。提示：一种可能的做法是使用你在练习
\ref{exercises-exercise-compute-cohomology-punctured} 中找到的复形。
\end{exercise}

\begin{exercise}
\label{exercises-exercise-h0-complement}
设 $A$ 为环，$I = (f_1, \ldots, f_t)$ 为 $A$ 的有限生成理想，
并令 $U \subset \Spec(A)$ 为 $V(I)$ 的补集。
设 $M$ 为 $A$-模，其 $I$-挠子模为零，即
$0 = \Ker((f_1, \ldots, f_t) : M \to M^{\oplus t})$。
证明存在典范同构
$$
H^0(U, \widetilde{M}) = \colim \Hom_A(I^n, M).
$$
警告：这并非显然。
\end{exercise}

\begin{exercise}
\label{exercises-exercise-Noetherian}
设 $A$ 为 Noether 环，$I$ 为 $A$ 的理想，$M$ 为 $A$-模。
令 $M[I^\infty]$ 为由下式定义的 $I$-幂挠元集合：
$$
M[I^\infty] = \{x \in M \mid
\text{ there exists an }n \geq 1\text{ 使得 }I^nx = 0\}
$$
令 $M' = M/M[I^\infty]$。则 $M'$ 的 $I$-幂挠子模为零。
证明
$$
\colim \Hom_A(I^n, M) = \colim \Hom_A(I^n, M').
$$
警告：这并非显然。提示：（1）试着归约到 $M$ 有限；
（2）若 $N = M[I^\infty]$ 且 $M$ 有限，证明
$\Ext^1_A(I^n, N)$ 的任意元素对某个 $m > 0$ 映到
$\Ext^1_A(I^{n + m}, N)$ 中的零；
（3）对 $\Hom_A(I^n, N)$ 证明与（2）相同的结论；
% P12-ERR-0470
（3）考察长正合列
$$
0 \to \Hom_A(I^n, M[I^\infty]) \to \Hom_A(I^n, M) \to \Hom_A(I^n, M')
\to \Ext^1_A(I^n, M[I^\infty])
$$
其中 $M$ 有限，并与 $I^{n + m}$ 所对应的正合列比较以得出结论。
\end{exercise}








\section{再论上同调}
\label{exercises-section-more-more-cohomology}


\begin{exercise}
\label{exercises-exercise-nonkillable}
构造如下例子：一个域 $k$、$k$ 上的一条曲线 $X$、
一个可逆 $\mathcal{O}_X$-模 $\mathcal{L}$，以及上同调类
$\xi \in H^1(X, \mathcal{L})$，使其具有如下性质：
对概形的每个有限满态射 $\pi : Y \to X$，元素 $\xi$
都拉回为 $H^1(Y, \pi^*\mathcal{L})$ 中的非零元素。
提示：构造 $X$、$k$、$\mathcal{L}$，使得存在短正合列
$0 \to \mathcal{L} \to \mathcal{O}_X \to i_*\mathcal{O}_Z \to 0$，
其中 $Z \subset X$ 是由多于 $1$ 个闭点组成的闭子概形。
% P12-ERR-0471
然后考察拉回时会发生什么。
\end{exercise}

\begin{exercise}
\label{exercises-exercise-cohomology-cycle-curves}
设 $k$ 为代数闭域，$X$ 为 $1$-维射影概形。
假设 $X$ 含有一个曲线圈，即假设存在 $n \geq 2$、
两两不同的 $1$-维整闭子概形 $C_1, \ldots, C_n$ 以及两两不同的闭点
$x_1, \ldots, x_n \in X$，使得 $x_n \in C_n \cap C_1$，
且对 $i = 1, \ldots, n - 1$ 有 $x_i \in C_i \cap C_{i + 1}$。
证明 $H^1(X, \mathcal{O}_X)$ 非零。
提示：令 $\mathcal{F}$ 为映射
% P12-ERR-0472
$\mathcal{O}_X \to \bigoplus \mathcal{O}_{C_i}$ 的像，
并利用 $\kappa(x_i) = k$ 与
$H^0(C_i, \mathcal{O}_{C_i}) = k$，证明
$H^1(X, \mathcal{F})$ 非零。
还要利用 Grothendieck 定理所给出的 $H^2(X, -) = 0$。
\end{exercise}

\begin{exercise}
\label{exercises-exercise-surface}
设 $X$ 为代数闭域 $k$ 上的射影曲面。
证明存在真闭子概形 $Z \subset X$，使 $H^1(Z, \mathcal{O}_Z)$
非零。提示：使用练习 \ref{exercises-exercise-cohomology-cycle-curves}。
\end{exercise}

\begin{exercise}
\label{exercises-exercise-surface-better}
设 $X$ 为代数闭域 $k$ 上的射影曲面。
证明对每个 $n \geq 0$，都存在真闭子概形 $Z \subset X$，
使得 $\dim_k H^1(Z, \mathcal{O}_Z) > n$。
只需说明如何修改练习 \ref{exercises-exercise-surface} 与
\ref{exercises-exercise-cohomology-cycle-curves} 中的论证；不必给出全部细节。
\end{exercise}

\begin{exercise}
\label{exercises-exercise-surface-h2-nonzero}
设 $X$ 为代数闭域 $k$ 上的射影曲面。
证明存在凝聚 $\mathcal{O}_X$-模 $\mathcal{F}$，使得
$H^2(X, \mathcal{F})$ 非零。
提示：使用练习 \ref{exercises-exercise-surface-better} 的结论以及一个巧妙选取的
正合列。
\end{exercise}

\begin{exercise}
\label{exercises-exercise-kunneth}
设 $X$ 与 $Y$ 为域 $k$ 上的概形
（可以随意假设 $X$ 与 $Y$ 性质良好，例如在 $k$ 上 qcqs 或射影）。
令 $X \times Y = X \times_{\Spec(k)} Y$，其投影为
$p : X \times Y \to X$ 与 $q : X \times Y \to Y$。
对拟凝聚 $\mathcal{O}_X$-模 $\mathcal{F}$ 和拟凝聚
$\mathcal{O}_Y$-模 $\mathcal{G}$，证明
$$
H^n(X \times Y,
p^*\mathcal{F} \otimes_{\mathcal{O}_{X \times Y}} q^*\mathcal{G}) =
\bigoplus\nolimits_{a + b = n}
H^a(X, \mathcal{F}) \otimes_k H^b(Y, \mathcal{G})
$$
或者只证明取维数后该等式成立。
“干净”的解答可获额外加分。
\end{exercise}

\begin{exercise}
\label{exercises-exercise-Oab}
设 $k$ 为域。令
% P12-ERR-0473
$X = \mathbf{P}|^1 \times \mathbf{P}^1$
为 $k$ 上射影直线与自身的乘积，投影为
$p : X \to \mathbf{P}^1_k$ 与 $q : X \to \mathbf{P}^1_k$。
令
$$
\mathcal{O}(a, b) = p^*\mathcal{O}_{\mathbf{P}^1_k}(a)
\otimes_{\mathcal{O}_X} q^*\mathcal{O}_{\mathbf{P}^1_k}(b)
$$
对所有 $i, a, b$，计算 $H^i(X, \mathcal{O}(a, b))$ 的维数。
提示：使用练习 \ref{exercises-exercise-kunneth}。
\end{exercise}











\section{上同调与 Hilbert 多项式}
\label{exercises-section-cohomology-hilbert-polynomials}

\begin{situation}
\label{exercises-situation-hilbert-polynomial}
设 $k$ 为域。设 $X = \mathbf{P}^n_k$ 为 $n$-维射影空间，
并设 $\mathcal{F}$ 为凝聚 $\mathcal{O}_X$-模。回顾
$$
\chi(X, \mathcal{F}) =
\sum\nolimits_{i = 0}^n (-1)^i \dim_k H^i(X, \mathcal{F})
$$
回顾，$\mathcal{F}$ 的{\it Hilbert 多项式}是函数
$$
t \longmapsto \chi(X, \mathcal{F}(t))
$$
还回顾
$\mathcal{F}(t) = \mathcal{F} \otimes_{\mathcal{O}_X} \mathcal{O}_X(t)$，
其中 $\mathcal{O}_X(t)$ 是结构层的第 $t$ 次扭曲，如构造篇定义
\ref{constructions-definition-twist}。
我们已在簇篇第 \ref{varieties-subsection-hilbert} 小节证明，
Hilbert 多项式是 $t$ 的多项式。
\end{situation}

\begin{exercise}
\label{exercises-exercise-hilbert-pol-easy}
在情形 \ref{exercises-situation-hilbert-polynomial} 中：
\begin{enumerate}
\item 若 $P(t)$ 是 $\mathcal{F}$ 的 Hilbert 多项式，
$\mathcal{F}(-13)$ 的 Hilbert 多项式是什么？
\item 若 $P_i$ 是 $\mathcal{F}_i$ 的 Hilbert 多项式，
$\mathcal{F}_1 \oplus \mathcal{F}_2$ 的 Hilbert 多项式是什么？
\item 若 $P_i$ 是 $\mathcal{F}_i$ 的 Hilbert 多项式，
且 $\mathcal{F}$ 是满射 $\mathcal{F}_1 \to \mathcal{F}_2$ 的核，
$\mathcal{F}$ 的 Hilbert 多项式是什么？
\end{enumerate}
\end{exercise}

\begin{exercise}
\label{exercises-exercise-find-given-hilbert-pol-dim-1}
在情形 \ref{exercises-situation-hilbert-polynomial} 中，假设 $n \geq 1$。
找出 Hilbert 多项式为 $t - 101$ 的凝聚层。
\end{exercise}

\begin{exercise}
\label{exercises-exercise-find-given-hilbert-pol-dim-2}
在情形 \ref{exercises-situation-hilbert-polynomial} 中，假设 $n \geq 2$。
找出 Hilbert 多项式为 $t^2/2 + t/2 - 1$ 的凝聚层。
（这有点棘手；只证明存在这样的凝聚层即可。）
\end{exercise}

\begin{exercise}
\label{exercises-exercise-bound-genus-in-degree}
在情形 \ref{exercises-situation-hilbert-polynomial} 中，假设 $n \geq 2$
且 $k$ 代数闭。设 $C \subset X$ 为维数 $1$ 的整闭子概形。
换言之，$C$ 是射影曲线。
设 $d t + e$ 为 $\mathcal{O}_C$ 作为 $X$ 上凝聚层的
Hilbert 多项式。
\begin{enumerate}
\item 给出 $e$ 的一个上界。（提示：利用
$\mathcal{O}_C(t)$ 仅在次数 $0$ 与 $1$ 有上同调，
并研究 $H^0(C, \mathcal{O}_C)$。）
\item 选取 $\mathcal{O}_X(1)$ 的整体截面 $s$，使它与 $C$
横截相交，即存在两两不同的闭点 $c_1, \ldots, c_r \in C$
以及短正合列
$$
0 \to \mathcal{O}_C \xrightarrow{s} \mathcal{O}_C(1) \to
\bigoplus\nolimits_{i = 1, \ldots, r} k_{c_i} \to 0
$$
其中 $k_{c_i}$ 是在 $c_i$ 处取值为 $k$ 的摩天层。
（这样的 $s$ 存在；请直接使用这一点。）
证明 $r = d$。（提示：扭曲该正合列并观察所得结果。）
\item 扭曲该短正合列后，对任意 $t$ 得到 $k$-线性映射
$\varphi_t : \Gamma(C, \mathcal{O}_C(t)) \to \bigoplus_{i = 1, \ldots, d} k$。
% P12-ERR-0474
证明若 $t \geq d - 1$，该映射为满射。
\item 用 $d$ 给出 $e$ 的一个下界。（提示：使用（3）的结果和消失性，
证明 $H^1(C, \mathcal{O}_C(d - 2)) = 0$。）
\end{enumerate}
\end{exercise}

\begin{exercise}
\label{exercises-exercise-three-quadrics-in-plane}
在情形 \ref{exercises-situation-hilbert-polynomial} 中，假设 $n = 2$。
设 $s_1, s_2, s_3 \in \Gamma(X, \mathcal{O}_X(2))$
为三个二次方程。考虑凝聚层
$$
\mathcal{F} = \Coker\left(\mathcal{O}_X(-2)^{\oplus 3}
\xrightarrow{s_1, s_2, s_3} \mathcal{O}_X\right)
$$
列出这样的 $\mathcal{F}$ 所有可能的 Hilbert 多项式。
（试着把射影平面中二次曲线的交可视化。）
\end{exercise}





\section{曲线}
\label{exercises-section-curves}

\begin{exercise}
\label{exercises-exercise-closed-subset-vanshing}
设 $k$ 为代数闭域，$X$ 为 $k$ 上的射影曲线。
设 $\mathcal{L}$ 为可逆 $\mathcal{O}_X$-模，并设
$s_0, \ldots, s_n \in H^0(X, \mathcal{L})$ 为 $\mathcal{L}$
的整体截面。证明存在自然闭子概形
$$
Z \subset
\mathbf{P}^n \times X
$$
使闭点 $((\lambda_0 : \ldots : \lambda_n), x)$ 属于 $Z$
当且仅当截面 $\lambda_0 s_0 + \ldots + \lambda_n s_n$
在 $x$ 处消失。（提示：在仿射局部描述 $Z$。）
\end{exercise}

\begin{exercise}
\label{exercises-exercise-diagonal-cartier}
设 $k$ 为代数闭域，$X$ 为 $k$ 上的光滑曲线。
设 $r \geq 1$。证明闭子集
$$
D \subset X \times X^r = X^{r + 1}
$$
具有可逆理想层；其闭点是满足某个 $i$ 下 $x = x_i$ 的元组
$(x, x_1, \ldots, x_r)$。换言之，证明 $D$ 是有效 Cartier 除子。
% P12-ERR-0475
提示：归约到 $r = 1$，并利用 $X$ 是光滑曲线这一点来研究对角线
（可查阅书籍）。
\end{exercise}

\begin{exercise}
\label{exercises-exercise-one-divisor-in-another}
设 $k$ 为代数闭域，$X$ 为 $k$ 上的光滑射影曲线。
设 $T$ 为 $k$ 上有限型概形，并设
$$
D_1 \subset X \times T
\quad\text{且}\quad
D_2 \subset X \times T
$$
为两个有效 Cartier 除子，使对 $t \in T$，纤维
$D_{i, t} \subset X_t$ 不稠密（即不含曲线 $X_t$ 的泛点）。
证明存在典范闭子概形 $Z \subset T$，使闭点 $t \in T$ 属于
$Z$，当且仅当 $D_1$、$D_2$ 的概形论纤维
$D_{1, t}$、$D_{2, t}$ 满足
$$
D_{1, t} \subset D_{2, t}
$$
（作为 $X_t$ 的闭子概形）。
提示：证明可以在缩小 $T$ 后假设 $T = \Spec(A)$ 仿射，
并存在仿射开集 $U \subset X$ 使 $D_i \subset U \times T$。
然后证明 $M_1 = \Gamma(D_1, \mathcal{O}_{D_1})$ 是有限局部自由
$A$-模（这里需要一些非平凡的代数——问问你的朋友）。
缩小 $T$ 后，可假设 $M_1$ 是自由 $A$-模。然后考察
$$
\Gamma(U \times T, \mathcal{I}_{D_2}) \to M_1 = A^{\oplus N}
$$
并把 $Z$ 定义为由该映射像中向量的系数所生成的理想切出的闭子概形。
说明这为何成立（这里可能还需要一点交换代数）。
\end{exercise}

\begin{exercise}
\label{exercises-exercise-closed-subset-vanshing-r}
设 $k$ 为代数闭域，$X$ 为 $k$ 上的光滑射影曲线。
设 $\mathcal{L}$ 为可逆 $\mathcal{O}_X$-模，并设
$s_0, \ldots, s_n \in H^0(X, \mathcal{L})$ 为 $\mathcal{L}$
的整体截面。设 $r \geq 1$。证明存在自然闭子概形
$$
Z \subset
\mathbf{P}^n \times X \times \ldots \times X = \mathbf{P}^n \times X^r
$$
使闭点 $((\lambda_0 : \ldots : \lambda_n), x_1, \ldots, x_r)$
属于 $Z$，当且仅当截面
$s_\lambda = \lambda_0 s_0 + \ldots + \lambda_n s_n$
在除子 $D = x_1 + \ldots + x_r$ 上消失，即截面 $s_\lambda$
属于 $\mathcal{L}(-D)$。
% P12-ERR-0476
提示：说明如何由练习 \ref{exercises-exercise-diagonal-cartier} 与
\ref{exercises-exercise-one-divisor-in-another} 的结论结合得出这一点。
\end{exercise}

\begin{exercise}
\label{exercises-exercise-effective}
设 $k$ 为代数闭域，$X$ 为 $k$ 上的光滑射影曲线。
设 $\mathcal{L}$ 为可逆 $\mathcal{O}_X$-模。
证明存在自然闭子集
$$
Z \subset X^r
$$
使 $X^r$ 的闭点 $(x_1, \ldots, x_r)$ 属于 $Z$，当且仅当
$\mathcal{L}(-x_1 - \ldots -x_r)$ 有非零整体截面。
提示：使用练习 \ref{exercises-exercise-closed-subset-vanshing-r}。
\end{exercise}

\begin{exercise}
\label{exercises-exercise-difference-effective}
设 $k$ 为代数闭域，$X$ 为 $k$ 上的光滑射影曲线。
设 $r \geq s$ 为整数。证明存在自然闭子集
$$
Z \subset X^r \times X^s
$$
使 $X^r \times X^s$ 的闭点
$(x_1, \ldots, x_r, y_1, \ldots, y_s)$ 属于 $Z$，当且仅当
$x_1 + \ldots + x_r - y_1 - \ldots - y_s$
线性等价于某个有效除子。
% P12-ERR-0477
提示：选取次数非常高的辅助可逆模 $\mathcal{L}$，使对任意次数
为 $r$ 的有效除子 $D$，$\mathcal{L}(-D)$ 都有非消失截面。
然后两次使用练习 \ref{exercises-exercise-effective} 的结论。
\end{exercise}

\begin{exercise}
\label{exercises-exercise-genus-7}
选择你最喜欢的代数闭域 $k$。尽你所能，确定某条亏格为 $7$ 的
曲线上可能存在的所有 $\mathfrak g^r_d$。同时还要试着：
\begin{enumerate}
\item 确定在哪些情形下 $\mathfrak g^r_d$ 无基点；
\item 确定在哪些情形下 $\mathfrak g^r_d$ 给出到
$\mathbf{P}^r$ 的闭嵌入。
\end{enumerate}
若假设曲线是“一般的”，完成同样的工作
（自行给出“一般”的概念——这可能比上面的问题容易）。
若假设曲线是超椭圆曲线，完成同样的工作。
若假设曲线是三角曲线（且不是超椭圆曲线），完成同样的工作。等等。
\end{exercise}





\section{模空间}
\label{exercises-section-moduli}

\noindent
本节考察代数几何对象的模空间的一些朴素处理方式。

\medskip\noindent
设 $k$ 为代数闭域。假设 $M$ 是 $k$ 上的一个模问题。
这里不精确定义其含义，但在每个练习中，我们的意思应当是清楚的。
要理解下面的内容，只需知道 $M$ 在 $k$ 上有哪些对象、$M$
在 $k$ 上的对象之间有哪些同构，以及簇上的 $M$-对象族是什么。
% P12-ERR-0478
若下列命题成立，就称 $M$ 的{\it 模数}为 $d \geq 0$：
\begin{enumerate}
\item 存在有限多个族 $X_i \to V_i$（$i = 1, \ldots, n$），
使 $M$ 在 $k$ 上的每个对象都同构于其中某个族的一条纤维，
且 $\max \dim(V_i) = d$；
\item 不可能用更小的 $d$ 做到这一点。
\end{enumerate}
这实际上只是文献中更好概念的一种非常粗略的近似。

\begin{exercise}
\label{exercises-exercise-number-moduli-lines-conics}
设 $k$ 为代数闭域，并设 $d \geq 1$、$n \geq 1$。
规定 $P^n$ 中 $d$ 次超曲面的模空间如下：
\begin{enumerate}
\item 一个对象是次数为 $d$ 的超曲面
$X \subset \mathbf{P}^n_k$；
\item 两个对象 $X \subset \mathbf{P}^n_k$ 与
$Y \subset \mathbf{P}^n_k$ 之间的一个同构，是满足
% P12-ERR-0479
$g(X) = Y$ 的元素 $g \in \text{PGL}_n(k)$；
\item 簇 $V$ 上的一个超曲面族，是闭子概形
$X \subset \mathbf{P}^n_V$，使对所有 $v \in V$，
$X \to V$ 的概形论纤维 $X_v$ 是 $\mathbf{P}^n_v$ 中的一条
% P12-ERR-0480
超曲面。
\end{enumerate}
计算下列模数（如果可以的话——它们会逐渐变难）：
\begin{enumerate}
\item $n = 1$ 而 $d$ 任意时的模数；
\item $n = 2$ 且 $d = 1$ 时的模数；
\item $n = 2$ 且 $d = 2$ 时的模数；
\item $n \geq 1$ 且 $d = 2$ 时的模数；
\item $n = 2$ 且 $d = 3$ 时的模数；
\item $n = 3$ 且 $d = 3$ 时的模数；
\item $n = 2$ 且 $d = 4$ 时的模数。
\end{enumerate}
\end{exercise}

\begin{exercise}
\label{exercises-exercise-number-moduli-hyperelliptic}
设 $k$ 为代数闭域，并设 $g \geq 2$。
规定亏格 $g$ 的超椭圆曲线的模空间如下：
\begin{enumerate}
\item 一个对象是亏格 $g$ 的光滑射影超椭圆曲线 $X$；
\item 两个对象 $X$ 与 $Y$ 之间的一个同构，是 $k$ 上概形的同构
$X \to Y$；
\item 簇 $V$ 上亏格 $g$ 的一个超椭圆曲线族，是一个
真平坦\footnote{去掉此假设也不会改变问题的答案。}态射
% P12-ERR-0481
$X \to Y$，使 $X \to V$ 的所有概形论纤维都是亏格 $g$ 的
光滑射影超椭圆曲线。
\end{enumerate}
证明其模数为 $2g - 1$。
\end{exercise}





\section{整体 Ext}
\label{exercises-section-global-exts}

\begin{exercise}
\label{exercises-exercise-ext-line-plane-p3}
设 $k$ 为域，$X = \mathbf{P}^3_k$。设 $L \subset X$ 与
$P \subset X$ 分别为直线和平面，把它们看作分别由 $1$ 个和
$2$ 个线性方程切出的闭子概形。对所有 $i$，计算
$$
\Ext^i_X(\mathcal{O}_L, \mathcal{O}_P)
$$
的维数。务必同时处理 $L$ 包含在 $P$ 中以及 $L$ 不包含在
$P$ 中的情形。
\end{exercise}

\begin{exercise}
\label{exercises-exercise-ext-point}
设 $k$ 为域，$X = \mathbf{P}^n_k$。
设 $Z \subset X$ 为一个闭的 $k$-有理点，并把它看作闭子概形；
例如，齐次坐标为 $(1 : 0 : \ldots : 0)$ 的点。
对所有 $i$，计算
$$
\Ext^i_X(\mathcal{O}_Z, \mathcal{O}_Z)
$$
的维数。
\end{exercise}

\begin{exercise}
\label{exercises-exercise-ext-pairings}
设 $X$ 为环空间。对 $\mathcal{O}_X$-模
$\mathcal{F}, \mathcal{G}, \mathcal{H}$ 定义杯积映射
$$
\Ext^i_X(\mathcal{G}, \mathcal{H})
\times
\Ext^j_X(\mathcal{F}, \mathcal{G})
\longrightarrow
\Ext^{i + j}_X(\mathcal{F}, \mathcal{H})
$$
（提示：这是一个极其一般的构造。）
\end{exercise}

\begin{exercise}
\label{exercises-exercise-move-out-locally-free}
设 $X$ 为环空间。设 $\mathcal{E}$ 为有限局部自由
$\mathcal{O}_X$-模，其对偶为
$\mathcal{E}^\vee = \SheafHom_{\mathcal{O}_X}(\mathcal{E}, \mathcal{O}_X)$。
证明下列命题：
\begin{enumerate}
\item $\SheafExt^i_{\mathcal{O}_X}(
\mathcal{F} \otimes_{\mathcal{O}_X} \mathcal{E}, \mathcal{G}) =
\SheafExt^i_{\mathcal{O}_X}(
\mathcal{F},
\mathcal{E}^\vee \otimes_{\mathcal{O}_X} \mathcal{G}) =
\SheafExt^i_{\mathcal{O}_X}(
\mathcal{F}, \mathcal{G}) \otimes_{\mathcal{O}_X} \mathcal{E}^\vee$；
\item $\Ext^i_X(
\mathcal{F} \otimes_{\mathcal{O}_X} \mathcal{E}, \mathcal{G}) =
\Ext^i_X(\mathcal{F},
\mathcal{E}^\vee \otimes_{\mathcal{O}_X} \mathcal{G})$。
\end{enumerate}
这里 $\mathcal{F}$ 与 $\mathcal{G}$ 是 $\mathcal{O}_X$-模。
推出
$$
\Ext^i_X(\mathcal{E}, \mathcal{G}) =
H^i(X, \mathcal{E}^\vee \otimes_{\mathcal{O}_X} \mathcal{G})
$$
\end{exercise}

\begin{exercise}
\label{exercises-exercise-trace-on-ext}
设 $X$ 为环空间，$\mathcal{E}$ 为有限局部自由
$\mathcal{O}_X$-模。对所有 $i$ 构造迹映射
$$
\Ext^i_X(\mathcal{E}, \mathcal{E}) \to H^i(X, \mathcal{O}_X)
$$
把它推广为对任意 $\mathcal{O}_X$-模 $\mathcal{F}$ 的迹映射
$$
\Ext^i_X(\mathcal{E}, \mathcal{E} \otimes_{\mathcal{O}_X} \mathcal{F})
\to H^i(X, \mathcal{F})
$$
\end{exercise}

\begin{exercise}
\label{exercises-exercise-duality}
设 $k$ 为域，$X = \mathbf{P}^d_k$。令
$\omega_{X/k} = \mathcal{O}_X(-d - 1)$。
证明：对有限局部自由模 $\mathcal{E}$、$\mathcal{F}$，
Ext 上的杯积与 Ext 上的迹映射结合得到
$$
\Ext^i_X(\mathcal{E}, \mathcal{F} \otimes_{\mathcal{O}_X} \omega_{X/k})
\times
\Ext^{d - i}_X(\mathcal{F}, \mathcal{E})
\to
\Ext_X^d(\mathcal{F}, \mathcal{F} \otimes_{\mathcal{O}_X} \omega_{X/k})
\to
H^d(X, \omega_{X/k}) = k
$$
是非退化配对。提示：你可以在此设定中重新证明对偶性，
也可以归约到层上同调，再应用讲课中证明的 Serre 对偶定理。
\end{exercise}








\section{除子}
\label{exercises-section-divisors}

\noindent
为方便起见，我们把所有相关定义汇集于此。

\begin{definition}
\label{exercises-definition-divisor}
以下始终令 $S$ 为任意概形，并令 $X$ 为 Noether 整概形。
\begin{enumerate}
\item $X$ 上的一个{\it Weil 除子}是素除子 $Z_i$ 的整数系数形式线性组合
$\Sigma n_i[Z_i]$。
\item 所谓{\it 素除子}，是整闭子概形 $Z \subset X$，其泛点
$\xi \in Z$ 满足 ${\mathcal O}_{X, \xi}$ 的维数为 $1$。
当 $\xi \in Z \subset X$ 如上时，我们使用记号
${\mathcal O}_{X, Z} = {\mathcal O}_{X, \xi}$。
注意，${\mathcal O}_{X, Z} \subset K(X)$ 是 $X$ 的函数域的子环。
\item 与有理函数 $f \in K(X)^\ast$ {\it 相伴的 Weil 除子}
是和 $\Sigma v_Z(f)[Z]$。这里 $v_Z(f)$ 定义如下：
\begin{enumerate}
\item 若 $f \in {\mathcal O}_{X, Z}^\ast$，则 $v_Z(f) = 0$。
\item 若 $f \in {\mathcal O}_{X, Z}$，则
$$
v_Z(f) = \text{length}_{{\mathcal O}_{X, Z}}({\mathcal O}_{X, Z}/(f)).
$$
\item 若 $f = \frac{a}{b}$，其中
$a, b \in {\mathcal O}_{X, Z}$，则
$$
v_Z(f) = \text{length}_{{\mathcal O}_{X, Z}}({\mathcal O}_{X, Z}/(a)) -
\text{length}_{{\mathcal O}_{X, Z}}({\mathcal O}_{X, Z}/(b)).
$$
\end{enumerate}
\item 概形 $S$ 上的一个{\it 有效 Cartier 除子}是闭子概形
$D \subset S$，使得每个点 $d\in D$ 都有 $S$ 中的仿射开邻域
% P12-ERR-0482
$\Spec(A) = U \subset S$，并且 $D \cap U = \Spec(A/(f))$，
其中 $f \in A$ 为非零因子。
\item 与 Noether 整概形 $X$ 的有效 Cartier 除子 $D \subset X$
{\it 相伴的 Weil 除子 $[D]$}定义为和
$\Sigma v_Z(D)[Z]$，其中 $v_Z(D)$ 定义如下：
\begin{enumerate}
\item 若 $Z$ 的泛点 $\xi$ 不属于 $D$，则 $v_Z(D) = 0$。
\item 若 $Z$ 的泛点 $\xi$ 属于 $D$，则
$$
v_Z(D) = \text{length}_{{\mathcal O}_{X, Z}}({\mathcal O}_{X, Z}/(f))
$$
其中 $f \in {\mathcal O}_{X, Z} = {\mathcal O}_{X, \xi}$ 是在
$\xi$ 的某个仿射邻域中定义 $D$ 的非零因子
（如上面的（4））。
\end{enumerate}
\item 设 $S$ 为概形。{\it 全商环层 ${\mathcal K}_S$}是
${\mathcal O}_S$-代数层，它是如下预层 ${\mathcal K}'$ 的层化。
对开集 $U \subset S$，令
${\mathcal K}'(U) = S_U^{-1}{\mathcal O}_S(U)$，其中
$S_U \subset {\mathcal O}_S(U)$ 是由如下截面组成的乘法子集：
$f \in {\mathcal O}_S(U)$，且对每个 $u\in U$，$f$ 在
${\mathcal O}_{S, u}$ 中的芽都是非零因子。
特别地，$S_U$ 的元素全都是非零因子。因此
${\mathcal O}_S$ 是 ${\mathcal K}_S$ 的子层，并得到短正合列
$$
0 \to {\mathcal O}_S^\ast \to {\mathcal K}_S^\ast \to
{\mathcal K}_S^\ast/{\mathcal O}_S^\ast \to 0.
$$
\item 概形 $S$ 上的一个{\it Cartier 除子}是商层
${\mathcal K}_S^\ast/{\mathcal O}_S^\ast$ 的整体截面。
\item 在 Noether 整概形 $X$ 上，与 Cartier 除子
$\tau \in \Gamma(X, {\mathcal K}_X^\ast/{\mathcal O}_X^\ast)$
{\it 相伴的 Weil 除子}是和 $\Sigma v_Z(\tau)[Z]$，其中
% P12-ERR-0483
$v_Z(\tau)$ 按下列步骤定义：
\begin{enumerate}
\item 若 $\tau$ 在 $Z$ 的泛点 $\xi$ 处的芽为零——换言之，
$\tau$ 在茎 $({\mathcal K}^\ast/{\mathcal O}^\ast)_\xi$ 中的像
为“零”——则 $v_Z(\tau) = 0$。
\item 找出仿射开邻域 $\Spec(A) = U \subset X$，使
$\tau|_U$ 是某个截面 $f \in {\mathcal K}(U)$ 的像，
而且 $f = a/b$，其中 $a, b \in A$。然后令
$$
v_Z(f) = \text{length}_{{\mathcal O}_{X, Z}}({\mathcal O}_{X, Z}/(a)) -
\text{length}_{{\mathcal O}_{X, Z}}({\mathcal O}_{X, Z}/(b)).
$$
\end{enumerate}
\end{enumerate}
\end{definition}

\begin{remarks}
\label{exercises-remarks-divisors}
下面是一些简单的说明。
\begin{enumerate}
\item 在 Noether 整概形 $X$ 上，层 ${\mathcal K}_X$ 是取值为
函数域 $K(X)$ 的常层。
\item 为使上述定义有意义，需要证明：对一维 Noether 局部整环
${\mathcal O}$ 的任意一对非零元素 $(a, b)$，有
$$
\text{length}_{\mathcal O}({\mathcal O}/(ab)) =
\text{length}_{\mathcal O}({\mathcal O}/(a)) +
\text{length}_{\mathcal O}({\mathcal O}/(b))
$$
这将在讲课中完成。
\end{enumerate}
\end{remarks}

\begin{exercise}
\label{exercises-exercise-effective-cartier-cartier}
（在任意概形上。）描述如何给一个有效 Cartier 除子指派
一个 Cartier 除子。
\end{exercise}

\begin{exercise}
\label{exercises-exercise-rational-function-cartier}
（在整概形上。）描述如何给有理函数 $f$ 指派 Cartier 除子
$D$，使与 $D$ 相伴的 Weil 除子和与 $f$ 相伴的 Weil 除子相同。
（这很显然。）
\end{exercise}

\begin{exercise}
\label{exercises-exercise-weil-not-cartier}
给出簇上的 Weil 除子的例子，使它不是与任何 Cartier 除子相伴的
Weil 除子。
\end{exercise}

\begin{exercise}
\label{exercises-exercise-weil-Q-cartier}
给出簇上的 Weil 除子 $D$ 的例子，使它不是与任何 Cartier 除子
相伴的 Weil 除子，但对某个 $n > 1$，$nD$ 是与一个 Cartier
除子相伴的 Weil 除子。
\end{exercise}

\begin{exercise}
\label{exercises-exercise-weil-not-Q-cartier}
给出簇上的 Weil 除子 $D$ 的例子，使它不是与任何 Cartier 除子
相伴的 Weil 除子，并且对任意 $n > 1$，$nD$ 都不是与任何
Cartier 除子相伴的 Weil 除子。
（提示：考虑一个锥，例如 $\mathbf{A}^4_k$ 中的
$X : xy - zw = 0$。试着证明 $D = [x = 0, z = 0]$ 可行。）
\end{exercise}

\begin{exercise}
\label{exercises-exercise-cartier-not-difference-effective-cartier}
在域上有限型分离概形 $X$ 上：
给出一个 Cartier 除子，使它不能写成两个有效 Cartier 除子的差。
提示：找一个没有任何非空有效 Cartier 除子的 $X$，例如
\cite[III Exercise 5.9]{H} 中构造的概形。
甚至有 $X$ 为簇的例子——即练习 \ref{exercises-exercise-nonprojective}
中的簇。
\end{exercise}

\begin{exercise}
\label{exercises-exercise-nonprojective}
非射影真簇的例子。设 $k$ 为域。设 $L \subset \mathbf{P}^3_k$
为直线，$C \subset \mathbf{P}^3_k$ 为非奇异二次曲线。
假设 $C \cap L = \emptyset$。选取同构 $\varphi : L \to C$。
令 $X$ 为通过 $\varphi$ 将 $C$ 与 $L$ 粘合所得的 $k$-簇。
换言之，存在满的真双有理态射
$$
\pi : \mathbf{P}^3_k \longrightarrow X
$$
以及开集 $U \subset X$，使
$\pi : \pi^{-1}(U) \to U$ 是同构，
$\pi^{-1}(U) = \mathbf{P}^3_k \setminus (L \cup C)$，
且 $\pi|_L = \pi|_C \circ \varphi$。
（这些条件尚不能唯一确定 $X$。为此还需指定 $X$ 在
$Z = X \setminus U$ 的点处的结构层。）
证明 $X$ 存在、是真簇，但不是射影的。
（提示：存在性使用练习 \ref{exercises-exercise-prepare-glueing} 的结论；
非射影性使用 $\Pic(\mathbf{P}^3_k) = \mathbf{Z}$，证明 $X$
不可能有充足可逆层。）
\end{exercise}




\section{微分}
\label{exercises-section-differentials}

\noindent
{\bf 定义与结果。}K\"ahler 微分。
\begin{enumerate}
\item 设 $R \to A$ 为环映射。$A$ 在 $R$ 上的
{\it K\"ahler 微分模}记作 $\Omega_{A/R}$。
它由元素 $\text{d}a$（$a \in A$）生成，并满足关系：
$$
\text{d}(a_1 + a_2) = \text{d}a_1 + \text{d}a_2,\quad
\text{d}(a_1a_2) = a_1\text{d}a_2 + a_2\text{d}a_1,\quad
\text{d}r = 0
$$
典范泛 $R$-导子 $\text{d} : A \to \Omega_{A/R}$
把 $a\mapsto \text{d}a$。
\item 考虑定义理想 $I$ 的短正合列
$$
0 \to I \to A \otimes_R A \to A \to 0
$$
存在典范导子 $\text{d} : A \to I/I^2$，把 $a$ 映到
$a \otimes 1 - 1 \otimes a$ 的类。这给出了 $A$ 在 $R$ 上的
微分模的另一种表示，换言之，
$$
(I/I^2, \text{d}) \cong (\Omega_{A/R}, \text{d}).
$$
\item 对满足 $S_R \subset R$ 的像包含在 $S_A \subset A$ 中的
乘法子集 $S_R$、$S_A$，有
$$
\Omega_{S_A^{-1}A / S_R^{-1}R} =
S_A^{-1}\Omega_{A/R}.
$$
\item 若 $A$ 是有限表现 $R$-代数，则 $\Omega_{A/R}$ 是有限表现
$A$-模。因此，在这种情形下，$\Omega_{A/R}$ 的
{\it Fitting} 理想有定义。
\item 设 $f : X \to S$ 为概形态射。存在拟凝聚
${\mathcal O}_X$-模层 $\Omega_{X/S}$ 以及
% P12-ERR-0484
${\mathcal O}_S$-线性导子
$$
\text{d} : {\mathcal O}_X \longrightarrow \Omega_{X/S}
$$
使得对满足 $f(U) \subset V$ 的任意仿射开集
$\Spec(A) = U \subset X$、$\Spec(R) = V \subset S$，
都有
$$
\Gamma(\Spec(A), \Omega_{X/S}) = \Omega_{A/R}
$$
并且与 $\text{d}$ 相容。
\end{enumerate}

\begin{exercise}
\label{exercises-exercise-dual-numbers}
设 $k[\epsilon]$ 为域 $k$ 上的对偶数环，即
$\epsilon^2 = 0$。
\begin{enumerate}
\item 考虑环映射
$$
R = k[\epsilon] \to A = k[x, \epsilon]/(\epsilon x)
$$
证明 $\Omega_{A/R}$ 的 Fitting 理想为
（从第零 Fitting 理想开始）
$$
(\epsilon), A, A, \ldots
$$
\item 考虑映射
$R = k[t] \to A = k[x, y, t]/(x(y-t)(y-1), x(x-t))$。
证明 $\Omega_{A/R}$ 在 $A$ 中的 Fitting 理想为
（为简单起见，假设 $k$ 的特征为零）
$$
x(2x-t)(2y-t-1)A, \ (x, y, t)\cap (x, y-1, t), \ A, \ A, \ldots
$$
% P12-ERR-0485
因此，第 $0$-个 Fitting 理想由 $A$ 的单个元素切出，
第 $1$ 个 Fitting 理想定义 $\Spec(A)$ 的两个闭点，
其余理想全都平凡。
\item 考虑映射 $R = k[t] \to A = k[x, y, t]/(xy-t^n)$。
计算 $\Omega_{A/R}$ 的 Fitting 理想。
\end{enumerate}
\end{exercise}

\begin{remark}
\label{exercises-remark-fitting-omega-not-sings}
$\Omega_{X/S}$ 的第 $k$ 个 Fitting 理想常用于定义态射
$X \to S$ 的奇异概形，其中 $X$ 在 $S$ 上的相对维数为 $k$。
% P12-ERR-0486
但正如（a）表明的，当族的纤维维数不是“常数”（例如它不平坦）时，
这样做必须小心。正如（b）表明的，平坦性也不保证这种做法有效
（而且，这确实是一个平坦族）。在“良好情形”中——例如（c）——
对于曲线族，我们期望第 $0$ 个 Fitting 理想为零，而第 $1$ 个
Fitting 理想在概形论意义下定义奇异轨迹。
\end{remark}

\begin{exercise}
\label{exercises-exercise-formally-smooth}
假设 $R$ 是环且
$$
A = R[x_1, \ldots, x_n]/(f_1, \ldots, f_n).
$$
注意，我们假设 $A$ 的表示中方程数与变量数相同。因此偏导矩阵
$$
( \partial f_i / \partial x_j )
$$
是 $n \times n$ 矩阵，即方阵。假设其行列式作为 $A$ 的元素可逆。
注意，在有 $n$ 个生成元与 $n$ 个关系的本情形中，这恰好是条件
$\Omega_{A/R} = (0)$。设 $\pi : B' \to B$ 为 $R$-代数的满射，
其核 $J$ 的平方为零（作为 $B'$ 的理想）。
设 $\varphi : A \to B$ 为 $R$-代数同态。
证明存在唯一的 $R$-代数同态 $\varphi' : A \to B'$，
使得 $\varphi = \pi \circ \varphi'$。
\end{exercise}

\begin{exercise}
\label{exercises-exercise-formally-smooth-one-equation}
把练习 \ref{exercises-exercise-formally-smooth} 的结论推广到
$A = R[x, y]/(f)$ 的情形。
\end{exercise}

\begin{exercise}
\label{exercises-exercise-Jacobian-criterion}
设 $k$ 为域，
$f_1, \ldots, f_c \in k[x_1, \ldots, x_n]$，并令
$A = k[x_1, \ldots, x_n]/(f_1, \ldots, f_c)$。
假设 $f_j(0, \ldots, 0) = 0$。这意味着
$\mathfrak m = (x_1, \ldots, x_n)A$ 是极大理想。
证明：若矩阵
$$
(\partial f_j/ \partial x_i)|_{(x_1, \ldots, x_n) = (0, \ldots, 0)}
$$
的秩为 $c$，则局部环 $A_\mathfrak m$ 正则。
在这种情形下，$A_\mathfrak m$ 的维数是多少？
给出 $A_\mathfrak m$ 正则但上述秩小于 $c$ 的例子，
以证明逆命题为假；在你的例子中，$A_\mathfrak m$ 的维数是多少？
\end{exercise}





\section{概形，2007 年秋季期末考试}
\label{exercises-section-final-exam-fall-2007}

\noindent
% P12-ERR-0487
这些是哥伦比亚大学 2007 年春季一门概形课程的期末考试题。

\begin{exercise}[定义]
\label{exercises-exercise-definitions}
给出下列概念的定义。
\begin{enumerate}
\item $X$ 是{\it 概形}
\item 概形态射 $f : X \to Y$ 是{\it 有限的}
% P12-ERR-0488
\item 概形态射 $f : X \to Y$ 是{\it 有限型的}
\item 概形 $X$ 是{\it Noether 的}
\item 概形 $X$ 上的 ${\mathcal O}_X$-模 ${\mathcal L}$ 是
{\it 可逆的}
\item 代数闭域上非奇异射影曲线的{\it 亏格}
\end{enumerate}
\end{exercise}

\begin{exercise}
\label{exercises-exercise-kill-global-sections}
设 $X = \Spec({\mathbf Z}[x, y])$，并设 ${\mathcal F}$ 为拟凝聚
${\mathcal O}_X$-模。假设 ${\mathcal F}$ 限制到标准仿射开集
$D(x)$ 后为零。
\begin{enumerate}
\item 证明 ${\mathcal F}$ 的每个整体截面 $s$ 都被 $x$ 的某个幂
消去，即对某个 $n\in {\mathbf N}$ 有 $x^ns = 0$。
\item 若不假设 ${\mathcal F}$ 拟凝聚，你认为同一结论仍成立吗？
\end{enumerate}
\end{exercise}

\begin{exercise}
\label{exercises-exercise-empty-fibre-empty}
假设 $X \to \Spec(R)$ 是真态射，$R$ 是剩余域为 $k$ 的离散赋值环。
假设 $X \times_{\Spec(R)} \Spec(k)$ 是空概形。
证明 $X$ 是空概形。
\end{exercise}

\begin{exercise}
\label{exercises-exercise-curve-p1-p1}
考虑复数域 ${\mathbf C}$ 上的射影\footnote{射影嵌入为
$((X_0, X_1), (Y_0, Y_1))\mapsto (X_0Y_0, X_0Y_1, X_1Y_0, X_1Y_1)$，
换言之，$(x, y)\mapsto (1, y, x, xy)$。}簇
$$
{\mathbf P}^1 \times {\mathbf P}^1 = {\mathbf P}^1_{{\mathbf C}}
\times_{\Spec({\mathbf C})} {\mathbf P}^1_{\mathbf C}
$$
它被四个仿射块覆盖，分别对应于 ${\mathbf P}^1_{\mathbf C}$
两个因子的标准仿射块对。例如，在第一个因子上使用齐次坐标
$X_0, X_1$，在第二个因子上使用 $Y_0, Y_1$。
令 $x = X_1/X_0$、$y = Y_1/Y_0$。则四个仿射开块是下列环的谱：
$$
{\mathbf C}[x, y], \quad
{\mathbf C}[x^{-1}, y], \quad
{\mathbf C}[x, y^{-1}], \quad
{\mathbf C}[x^{-1}, y^{-1}].
$$
设 $X \subset {\mathbf P}^1 \times {\mathbf P}^1$ 为闭子概形，
它是第一个仿射块中由方程
$$
y^3(x^4 + 1) = x^4 -1.
$$
给出的闭子集的闭包。
\begin{enumerate}
\item 证明 $X$ 包含在四个仿射开块中第一个与最后一个的并中。
\item 证明 $X$ 是非奇异射影曲线。
% P12-ERR-0489
\item 考虑态射 $pr_2 : X \to {\mathbf P}^1$（到第一个因子的投影）。
在第一个仿射块上，它是映射 $(x, y) \mapsto x$。
简要说明为何它的次数为 $3$。
\item 计算映射 $pr_2 : X \to {\mathbf P}^1$ 的分歧点与分歧指数。
\item 计算 $X$ 的亏格。
\end{enumerate}
\end{exercise}

\begin{exercise}
\label{exercises-exercise-finite-type-over-Z}
设 $X \to \Spec({\mathbf Z})$ 为有限型态射。
假设有无穷多个素数 $p$，使得
$X \times_{\Spec({\mathbf Z})} \Spec({\mathbf F}_p)$ 非空。
\begin{enumerate}
\item 证明 $X \times_{\Spec({\mathbf Z})}\Spec(\mathbf{Q})$ 非空。
\item 若把“有限型”条件换成“局部有限型”条件，你认为同一结论仍成立吗？
\end{enumerate}
\end{exercise}





\section{概形，2009 年春季期末考试}
\label{exercises-section-final-exam-spring-2009}

\noindent
这些是哥伦比亚大学 2009 年春季一门概形课程的期末考试题。

\begin{exercise}
\label{exercises-exercise-Noetherian-coherent}
设 $X$ 为 Noether 概形，$\mathcal{F}$ 为 $X$ 上的凝聚层，
$x \in X$ 为一点。假设
$\text{Supp}(\mathcal{F}) = \{ x \}$。
\begin{enumerate}
\item 证明 $x$ 是 $X$ 的闭点。
\item 证明 $H^0(X, \mathcal{F})$ 非零。
\item 证明 $\mathcal{F}$ 由整体截面生成。
\item 证明对 $p > 0$，$H^p(X, \mathcal{F}) = 0$。
\end{enumerate}
\end{exercise}

\begin{remark}
\label{exercises-remark-invertible-projective-space}
设 $k$ 为域，并令
$\mathbf{P}^2_k = \text{Proj}(k[X_0, X_1, X_2])$。
$\mathbf{P}^2_k$ 上的任意可逆层都同构于某个
$n \in \mathbf{Z}$ 所对应的
$\mathcal{O}_{\mathbf{P}^2_k}(n)$。回顾
$$
\Gamma(\mathbf{P}^2_k, \mathcal{O}_{\mathbf{P}^2_k}(n)) =
k[X_0, X_1, X_2]_n
$$
是多项式环的次数 $n$ 部分。对 $\mathbf{P}^2_k$ 上的拟凝聚层
$\mathcal{F}$，照例令
$\mathcal{F}(n) =
\mathcal{F}
\otimes_{\mathcal{O}_{\mathbf{P}^2_k}}
\mathcal{O}_{\mathbf{P}^2_k}(n)$。
\end{remark}

\begin{exercise}
\label{exercises-exercise-nonsplit-vectorbundle}
设 $k$ 为域。设 $\mathcal{E}$ 为 $\mathbf{P}^2_k$ 上的向量丛，
即有限局部自由 $\mathcal{O}_{\mathbf{P}^2_k}$-模。
% P12-ERR-0490
若 $\mathcal{E}$ 同构于可逆
$\mathcal{O}_{\mathbf{P}^2_k}$-模的直和，则称 $\mathcal{E}$
{\it 分裂}。
\begin{enumerate}
\item 证明 $\mathcal{E}$ 分裂当且仅当 $\mathcal{E}(n)$ 分裂。
\item 证明若 $\mathcal{E}$ 分裂，则对所有 $n \in \mathbf{Z}$，
$H^1({\mathbf{P}^2_k}, \mathcal{E}(n)) = 0$。
\item 设
$$
\varphi :
\mathcal{O}_{\mathbf{P}^2_k}
\longrightarrow
\mathcal{O}_{\mathbf{P}^2_k}(1) \oplus
\mathcal{O}_{\mathbf{P}^2_k}(1) \oplus
\mathcal{O}_{\mathbf{P}^2_k}(1)
$$
由线性形式
$L_0, L_1, L_2 \in \Gamma(\mathbf{P}^2_k, \mathcal{O}_{\mathbf{P}^2_k}(1))$
给出。假设对某个 $i$ 有 $L_i \not = 0$。
$L_0, L_1, L_2$ 应满足什么条件，才能使 $\varphi$ 的余核为向量丛？
为什么？
% P12-ERR-0491
\item 给出这样的 $\varphi$ 的例子。
\item 证明 $\Coker(\varphi)$ 不分裂（若它是向量丛）。
\end{enumerate}
\end{exercise}

\begin{remark}
\label{exercises-remark-recall-dimension-theory}
可以自由使用下列维数理论事实（若需要，还可自行添加）。
\begin{enumerate}
\item 概形的维数是不可约闭子集链长度的上确界。
\item 域上有限型概形的维数是其仿射开集维数的最大值。
\item Noether 概形的维数是其不可约分支维数的最大值。
\item 仿射概形的维数等于对应环的维数。
\item 设 $k$ 为域，$A$ 为有限型 $k$-代数。
% P12-ERR-0492
若 $A$ 是整环且 $x \not = 0$，则
$\dim(A) = \dim(A/xA) + 1$。
\end{enumerate}
\end{remark}

\begin{exercise}
\label{exercises-exercise-irreducible-fibres-same-dimension-irreducible}
设 $k$ 为域，$X$ 为 $k$ 上射影既约概形。
设 $f : X \to \mathbf{P}^1_k$ 为 $k$ 上的概形态射。
假设存在整数 $d \geq 0$，使对每个点
$t \in \mathbf{P}^1_k$，纤维 $X_t = f^{-1}(t)$
都不可约且维数为 $d$。（回顾：不可约空间非空。）
\begin{enumerate}
\item 证明 $\dim(X) = d + 1$。
\item 设 $X_0 \subset X$ 为 $X$ 的一个维数为 $d + 1$
的不可约分支。证明对每个 $t \in \mathbf{P}^1_k$，
纤维 $X_{0, t}$ 的维数为 $d$。
\item 由上面各项，关于 $X_t$ 与 $X_{0, t}$ 能得出什么结论？
\item 证明 $X$ 不可约。
\end{enumerate}
\end{exercise}

\begin{remark}
\label{exercises-remark-chi}
给定域 $k$ 上的射影概形 $X$ 以及 $X$ 上的凝聚层
$\mathcal{F}$，令
$$
\chi(X, \mathcal{F}) =
\sum\nolimits_{i \geq 0} (-1)^i\dim_k H^i(X, \mathcal{F}).
$$
\end{remark}

\begin{exercise}
\label{exercises-exercise-complete-intersection}
设 $k$ 为域。写
$\mathbf{P}^3_k = \text{Proj}(k[X_0, X_1, X_2, X_3])$。
设 $C \subset \mathbf{P}^3_k$ 为
{\it $(5, 6)$ 型完全交曲线}。这意味着存在
$F \in k[X_0, X_1, X_2, X_3]_5$ 与
$G \in k[X_0, X_1, X_2, X_3]_6$，使
$$
C = \text{Proj}(k[X_0, X_1, X_2, X_3]/(F, G))
$$
是维数为 $1$ 的簇。（簇意味着既约且不可约，但若愿意，
可以直接假设 $C$ 非奇异。）
令 $i : C \to \mathbf{P}^3_k$ 为相应的闭浸入。
完全交还意味着
$$
\xymatrix{
0 \ar[r] &
\mathcal{O}_{\mathbf{P}^3_k}(-11)
\ar[r]^-{
\left(
\begin{matrix}
-G \\
F
\end{matrix}
\right)
} &
\mathcal{O}_{\mathbf{P}^3_k}(-5) \oplus \mathcal{O}_{\mathbf{P}^3_k}(-6)
\ar[r]^-{(F, G)} &
\mathcal{O}_{\mathbf{P}^3_k} \ar[r] &
i_*\mathcal{O}_C \ar[r] &
0
}
$$
是层的正合列。请使用这些事实：
\begin{enumerate}
\item 对任意 $n \in \mathbf{Z}$ 计算
$\chi(C, i^*\mathcal{O}_{\mathbf{P}^3_k}(n))$；
\item 计算 $H^1(C, \mathcal{O}_C)$ 的维数。
\end{enumerate}
\end{exercise}

\begin{exercise}
\label{exercises-exercise-glueing}
设 $k$ 为域。考虑环
\begin{align*}
A & = k[x, y]/(xy) \\
B & = k[u, v]/(uv) \\
C & = k[t, t^{-1}] \times k[s, s^{-1}]
\end{align*}
以及 $k$-代数映射
$$
\begin{matrix}
A \longrightarrow C, &
x \mapsto (t, 0), &
y \mapsto (0, s) \\
B \longrightarrow C, &
u \mapsto (t^{-1}, 0), &
v \mapsto (0, s^{-1})
\end{matrix}
$$
事实是，这些映射诱导同构 $A_{x + y} \to C$ 与
$B_{u + v} \to C$。因此映射 $A \to C$ 与 $B \to C$
把 $\Spec(C)$ 分别等同于 $\Spec(A)$ 与 $\Spec(B)$ 的开子集。
令 $X$ 为沿 $\Spec(C)$ 粘合 $\Spec(A)$ 与 $\Spec(B)$
所得的概形：
$$
X = \Spec(A) \amalg_{\Spec(C)} \Spec(B).
$$
正如课程中所见，这样的概形存在，并有仿射开集
$\Spec(A) \subset X$ 与 $\Spec(B) \subset X$；
它们的交恰好是 $\Spec(C)$，通过上述映射与二者各自的一个开集等同。
\begin{enumerate}
\item 为什么 $X$ 分离？
\item 为什么 $X$ 在 $k$ 上有限型？
\item 计算 $H^1(X, \mathcal{O}_X)$，或者说，它的维数是多少？
\item 如何以更几何的方式描述 $X$？
\end{enumerate}
\end{exercise}





\section{概形，2010 年秋季期末考试}
\label{exercises-section-final-exam-fall-2010}

\noindent
这些是哥伦比亚大学 2010 年秋季一门概形课程的期末考试题。

\begin{exercise}[定义]
\label{exercises-exercise-definitions-fall-2010}
给出下列概念的定义。
\begin{enumerate}
\item 分离概形；
\item 概形的拟紧态射；
\item 概形的仿射态射；
\item 环的乘法子集；
\item Noether 概形；
\item 簇。
\end{enumerate}
\end{exercise}

\begin{exercise}
\label{exercises-exercise-prime-avoidance}
素理想回避。
\begin{enumerate}
\item 设 $A$ 为环，$I \subset A$ 为理想，并设
$\mathfrak q_1$、$\mathfrak q_2$ 为满足
$I \not \subset \mathfrak q_i$ 的素理想。
证明 $I \not \subset \mathfrak q_1 \cup \mathfrak q_2$。
\item （1）的几何解释是什么？
\item 设 $X = \text{Proj}(S)$，其中 $S$ 为某个分次环。
设 $x_1, x_2 \in X$。证明存在同时包含 $x_1$ 与 $x_2$
的标准开集 $D_{+}(F)$。
\end{enumerate}
\end{exercise}

\begin{exercise}
\label{exercises-exercise-compose-affine}
为什么仿射态射的复合仍是仿射态射？
\end{exercise}

\begin{exercise}[例子]
\label{exercises-exercise-examples}
给出下列对象的例子：
\begin{enumerate}
\item 域 $k$ 上的可约射影概形；
\item 具有 100 个点的概形；
\item 概形的非仿射态射。
\end{enumerate}
\end{exercise}

\begin{exercise}
\label{exercises-exercise-chevalley-hilbert-nullstellensatz}
Chevalley 定理与 Hilbert 零点定理。
\begin{enumerate}
\item 设 $\mathfrak p \subset \mathbf{Z}[x_1, \ldots, x_n]$
为极大理想。Chevalley 定理对
$\mathfrak p \cap \mathbf{Z}$ 蕴含什么？
\item 进而，Hilbert 零点定理对 $\kappa(\mathfrak p)$ 蕴含什么？
\end{enumerate}
\end{exercise}

\begin{exercise}
\label{exercises-exercise-P0}
设 $A$ 为环。令 $S = A[X]$，把它看成分次 $A$-代数，
其中 $X$ 的次数为 $1$。证明
$\text{Proj}(S) \cong \Spec(A)$（作为 $A$ 上的概形）。
\end{exercise}

\begin{exercise}
\label{exercises-exercise-finite-is-projective}
设 $A \to B$ 为有限环映射。
证明 $\Spec(B)$ 是 $\Spec(A)$ 上的 H-射影概形。
\end{exercise}

\begin{exercise}
\label{exercises-exercise-not-geometrically-irreducible}
给出域 $k$ 上概形 $X$ 的例子，使 $X$ 不可约，
但对某个有限扩张 $k'/k$，基变换
$X_{k'} = X \times_{\Spec(k)} \Spec(k')$ 连通但可约。
\end{exercise}





\section{概形，2011 年春季期末考试}
\label{exercises-section-final-exam-spring-2011}

\noindent
这些是哥伦比亚大学 2011 年春季一门概形课程的期末考试题。

\begin{exercise}[定义]
\label{exercises-exercise-definitions-spring-2011}
给出斜体概念的定义。
\begin{enumerate}
\item {\it 分离}概形；
\item 概形的{\it 普遍闭}态射；
\item 对包含在同一个域中的局部环 $A, B$，所谓
{\it $A$ 控制 $B$}；
\item 概形 $X$ 的{\it 维数}；
\item 概形 $X$ 的不可约闭子概形 $Y$ 的{\it 余维数}。
\end{enumerate}
\end{exercise}

\begin{exercise}[结果]
\label{exercises-exercise-results-spring-2011}
陈述某个与课程中讨论的事实形式等价的命题。
\begin{enumerate}
\item 例如，簇的态射 $X \to Y$ 的真性赋值判别法。
\item 对域 $k$ 上的簇 $X$，$\dim(X)$ 与 $X$ 的函数域
$k(X)$ 之间的关系。
\item 填空：$k$ 上非奇异射影曲线与非常值态射的范畴，
反等价于 $\ldots\ldots\ldots$。
\item Noether 正规化。
\item Jacobian 判别法。
\end{enumerate}
\end{exercise}

\begin{exercise}
\label{exercises-exercise-genus-plane-curve}
设 $k$ 为域。设 $F \in k[X_0, X_1, X_2]$ 为次数 $d$ 的齐次形式。
假设 $C = V_{+}(F) \subset \mathbf{P}^2_k$ 是 $k$ 上的光滑曲线。
% P12-ERR-0493
记 $i : C \to \mathbf{P}^2_k$ 为相应的闭浸入。
\begin{enumerate}
\item 证明 $\mathbf{P}^2_k$ 上存在凝聚层的短正合列
$$
0 \to \mathcal{O}_{\mathbf{P}^2_k}(-d)
\to \mathcal{O}_{\mathbf{P}^2_k} \to i_*\mathcal{O}_C \to 0
$$
请说明映射是什么，并简要说明为何该列正合。
\item 推出 $H^0(C, \mathcal{O}_C) = k$。
\item 计算 $C$ 的亏格。
\item 现在假设 $P = (0 : 0 : 1)$ 不在 $C$ 上。
证明由 $(a_0 : a_1 : a_2) \mapsto (a_0 : a_1)$ 给出的
$\pi : C \to \mathbf{P}^1_k$ 的次数为 $d$。
\item 假设 $k$ 代数闭，所有分歧指数（“$e_i$”）均为
$1$ 或 $2$，且 $k$ 的特征不等于 $2$。
$\pi : C \to \mathbf{P}^1_k$ 有多少个分歧点？
\item 用 $F$ 表示，你认为 $\pi$ 的分歧点集合的一组方程是什么？
\item 你能猜出 $K_C$ 吗？
\end{enumerate}
\end{exercise}

\begin{exercise}
\label{exercises-exercise-Pic-triangle}
设 $k$ 为域。设 $X$ 为 $k$ 上的一个“三角形”，即取三份
$\mathbf{A}^1_k$，把第一份的 $0$ 与第二份的 $1$ 粘合，
第二份的 $0$ 与第三份的 $1$ 粘合，再把第三份的 $0$ 与第一份的
$1$ 粘合，从而得到 $X$。事实证明，
$X$ 同构于 $\Spec(k[x, y]/(xy(x + y + 1)))$；可以直接使用这一点。
计算 $X$ 的 Picard 群。
\end{exercise}

\begin{exercise}
\label{exercises-exercise-birational-morphism-curves-ample}
设 $k$ 为域。设 $\pi : X \to Y$ 为曲线的有限双有理态射，
其中 $X$ 是 $k$ 上的射影非奇异曲线。
由课程内容可知，$Y$ 是真曲线，$\pi$ 是 $Y$ 的正规化态射。
课程中还已看到，存在稠密开集 $V \subset Y$，使
$U = \pi^{-1}(V)$ 是 $X$ 中的稠密开集，且
$\pi : U \to V$ 是同构。
\begin{enumerate}
\item 证明存在有效 Cartier 除子 $D \subset X$，使
$D \subset U$，并且 $\mathcal{O}_X(D)$ 在 $X$ 上充足。
\item 设 $D$ 如（1）。证明 $E = \pi(D)$ 是 $Y$ 上的有效
Cartier 除子。
\item 简要说明为何：
\begin{enumerate}
\item 映射 $\mathcal{O}_Y \to \pi_*\mathcal{O}_X$ 有凝聚余核
$Q$，且它支撑在 $Y \setminus V$ 上；
\item 对每个 $n$，都有相应映射
$\mathcal{O}_Y(nE) \to \pi_*\mathcal{O}_X(nD)$，
其余核同构于 $Q$。
\end{enumerate}
\item 证明
$\dim_k H^0(X, \mathcal{O}_X(nD)) - \dim_k H^0(Y, \mathcal{O}_Y(nE))$
有界（被什么量所界？），并推出当 $n$ 很大时，可逆层
$\mathcal{O}_Y(nE)$ 有许多截面（为什么？）。
\end{enumerate}
\end{exercise}






\section{概形，2011 年秋季期末考试}
\label{exercises-section-final-exam-fall-2011}

\noindent
这些是哥伦比亚大学 2011 年秋季一门交换代数课程的期末考试题。


\begin{exercise}[定义]
\label{exercises-exercise-definitions-fall-2011}
给出斜体概念的定义。
\begin{enumerate}
\item {\it Noether} 环；
\item {\it Noether} 概形；
\item {\it 有限}环同态；
\item 概形的{\it 有限}态射；
\item 环的{\it 维数}。
\end{enumerate}
\end{exercise}

\begin{exercise}[结果]
\label{exercises-exercise-results-fall-2011}
陈述某个与课程中讨论的事实形式等价的命题。
\begin{enumerate}
\item Zariski 主定理。
\item Noether 正规化。
\item 中国剩余定理。
\item 有限环映射的上升定理。
\end{enumerate}
\end{exercise}

\begin{exercise}
\label{exercises-exercise-dimension-of-ring}
设 $(A, \mathfrak m, \kappa)$ 为 Noether 局部环，
其剩余域的特征不为 $2$。
假设 $\mathfrak m$ 由三个元素 $x, y, z$ 生成，
并且在 $A$ 中 $x^2 + y^2 + z^2 = 0$。
\begin{enumerate}
\item $\dim(A)$ 的可能取值是什么？
\item 给出例子，说明每个取值都可能出现。
\item 证明若 $\dim(A) = 2$，则 $A$ 是整环。
（提示：考察 $\bigoplus_{n \geq 0} \mathfrak m^n/\mathfrak m^{n + 1}$。）
\end{enumerate}
\end{exercise}

\begin{exercise}
\label{exercises-exercise-localization}
设 $A$ 为环，$S \subset T \subset A$ 为乘法子集。假设
$$
\{\mathfrak q \mid \mathfrak q \cap S = \emptyset\} =
\{\mathfrak q \mid \mathfrak q \cap T = \emptyset\}.
$$
证明 $S^{-1}A \to T^{-1}A$ 是同构。
\end{exercise}

\begin{exercise}
\label{exercises-exercise-locus-of-rank-1}
设 $k$ 为代数闭域。令
$$
V_0 = \{ A \in \text{Mat}(3 \times 3, k) \mid \text{rank}(A) = 1\}
\subset \text{Mat}(3 \times 3, k) = k^9.
$$
\begin{enumerate}
\item 证明 $V_0$ 是某个（Zariski）局部闭子集
$V \subset \mathbf{A}^9_k$ 的闭点集合。
\item $V$ 不可约吗？
\item $\dim(V)$ 是多少？
\end{enumerate}
\end{exercise}

\begin{exercise}
\label{exercises-exercise-not-complete-intersection}
证明 $\mathbf{C}[x, y]$ 中的理想 $(x^2, xy, y^2)$
不能由 $2$ 个元素生成。
\end{exercise}

\begin{exercise}
\label{exercises-exercise-finite-projection}
设 $f \in \mathbf{C}[x, y]$ 为非常值多项式。
证明对某些 $\alpha, \beta \in \mathbf{C}$，$\mathbf{C}$-代数映射
$$
\mathbf{C}[t] \longrightarrow \mathbf{C}[x, y]/(f),\quad
t \longmapsto \alpha x + \beta y
$$
是有限的。
\end{exercise}

\begin{exercise}
\label{exercises-exercise-union-of-two-affines}
证明：给定有限多个点
$p_1, \ldots, p_n \in \mathbf{C}^2$，概形
$\mathbf{A}^2_\mathbf{C} \setminus \{p_1, \ldots, p_n\}$
是两个仿射开集的并。
\end{exercise}

\begin{exercise}
\label{exercises-exercise-surjection-A1-P1}
证明存在概形的满态射
$\mathbf{A}^1_\mathbf{C} \to \mathbf{P}^1_\mathbf{C}$。
（“满”仅表示在底层点集上为满射。）
\end{exercise}

\begin{exercise}
\label{exercises-exercise-almost-surjective}
设 $k$ 为代数闭域，$A \subset B$ 为整环扩张，并且二者都是
有限型 $k$-代数。用下列步骤证明
$\Spec(B) \to \Spec(A)$ 的像包含 $\Spec(A)$ 的一个非空开子集：
\begin{enumerate}
\item 当 $A \to B$ 还是有限映射时证明该结论。
\item 当 $B$ 的分式域是 $A$ 的分式域的有限扩张时证明该结论。
\item 把命题归约到前一种情形。
\end{enumerate}
\end{exercise}






\section{概形，2013 年秋季期末考试}
\label{exercises-section-final-exam-fall-2013}

\noindent
这些是哥伦比亚大学 2013 年秋季一门交换代数课程的期末考试题。

\begin{exercise}[定义]
\label{exercises-exercise-definitions-fall-2013}
给出斜体概念的定义。
\begin{enumerate}
\item 环的{\it 根理想}；
\item {\it 有限型}环同态；
\item {\it Weil 式微分}；
\item {\it 概形}。
\end{enumerate}
\end{exercise}

\begin{exercise}[结果]
\label{exercises-exercise-results-fall-2013}
陈述某个与课程中讨论的事实形式等价的命题。
\begin{enumerate}
\item 关于分次模的 Hilbert 多项式的结果。
\item Noether 局部环 $(R, \mathfrak m)$ 的维数与
$\bigoplus_{n \geq 0} \mathfrak m^n/\mathfrak m^{n + 1}$。
\item Riemann--Roch 定理。
\item Clifford 定理。
\item Chevalley 定理。
\end{enumerate}
\end{exercise}

\begin{exercise}
\label{exercises-exercise-surjective-after-localization}
设 $A \to B$ 为环映射，$S \subset A$ 为乘法子集。
假设 $A \to B$ 有限型，并且
$S^{-1}A \to S^{-1}B$ 为满射。证明存在 $f \in S$，
使 $A_f \to B_f$ 为满射。
\end{exercise}

\begin{exercise}
\label{exercises-exercise-injective-local-ring-map-not-surjective}
给出局部环的单射局部同态 $A \to B$ 的例子，使得
$\Spec(B) \to \Spec(A)$ 不满。
\end{exercise}

\begin{situation}[平面曲线的记号]
\label{exercises-situation-curve-in-the-plane}
设 $k$ 为代数闭域。设
$F(X_0, X_1, X_2) \in k[X_0, X_1, X_2]$ 为次数 $d$ 的
齐次不可约多项式。令
$$
D = V(F) \subset \mathbf{P}^2
$$
为 $F$ 的零点所给出的射影平面曲线。

令 $x = X_1/X_0$、$y = X_2/X_0$，并令
$f(x, y) = X_0^{-d}F(X_0, X_1, X_2) = F(1, x, y)$。

% P12-ERR-0494
用 $K$ 表示整环 $k[x, y]/(f)$ 的分式域。
令 $C$ 为与 $K$ 对应的抽象曲线。
回顾（来自讲课）：存在满映射 $C \to D$，它在 $D$ 的非奇异轨迹上
为双射，并且当 $D$ 非奇异时为同构。
令 $f_x = \partial f/\partial x$、$f_y = \partial f/\partial y$。
最后，用 $\omega = \text{d}x/f_y = - \text{d}y/f_x$ 表示讲课中讨论的
$\Omega_{K/k}$ 的元素。用 $K_C$ 表示 $\omega$ 的零点与极点除子。
\end{situation}

\begin{exercise}
\label{exercises-exercise-node-in-the-plane}
在情形 \ref{exercises-situation-curve-in-the-plane} 中，假设 $d \geq 3$，
且曲线 $D$ 恰有一个奇点，即 $P = (1 : 0 : 0)$。
进一步假设在 $P = (0, 0)$ 附近有展开式
$$
f(x, y) = xy + h.o.t
$$
则 $C$ 在 $P$ 上方有两个点 $v$ 与 $w$，由下式刻画：
$$
v(x) = 1, v(y) > 1
\quad\text{且}\quad
w(x) > 1, w(y) = 1
$$
\begin{enumerate}
\item 证明 $\Omega_{K/k}$ 的元素
$\omega = \text{d}x/f_y = - \text{d}y/f_x$
在 $v$ 与 $w$ 处都有一阶极点。
（$\omega$ 在非奇异点处的行为如讲课中所述。）
\item 讲课中已证明，$\omega$ 在通过映射 $C \to D$ 拉回到
$C$ 的除子 $X_0 = 0$ 处消失至 $d - 3$ 阶。
结合（1）的信息，$C$ 上 $\omega$ 的零点与极点除子的次数是多少？
\item 曲线 $C$ 的亏格是多少？
\end{enumerate}
\end{exercise}

\begin{exercise}
\label{exercises-exercise-smooth-plane-curve-linear-system}
在情形 \ref{exercises-situation-curve-in-the-plane} 中，假设 $d = 5$，
且曲线 $C = D$ 非奇异。讲课中已证明 $C$ 的亏格为 $6$，
而线性系 $K_C$ 由
$$
L(K_C) = \{h\omega \mid h \in k[x, y],\ \deg(h) \leq 2\}
$$
给出，其中 $\deg$ 表示总次数\footnote{这里得到 $\leq 2$，
因为 $d - 3 = 5 - 3 = 2$。}。
设 $P_1, P_2, P_3, P_4, P_5 \in D$ 为位于仿射开集
$X_0 \not = 0$ 中的两两不同的点。
用 $\sum P_i = P_1 + P_2 + P_3 + P_4 + P_5$ 表示 $C$ 的相应除子。
\begin{enumerate}
\item 用多项式描述 $L(K_C - \sum P_i)$。
\item $l(\sum P_i)$ 有哪些可能值？
\end{enumerate}
\end{exercise}

\begin{exercise}
\label{exercises-exercise-rational-curve-high-degree}
写出情形 \ref{exercises-situation-curve-in-the-plane} 中的一个 $F$，
使 $d = 100$ 且 $C$ 的亏格为 $0$。
\end{exercise}

\begin{exercise}
\label{exercises-exercise-high-degree-curve-quadratic-equation}
设 $k$ 为代数闭域。设 $K/k$ 为超越次数 $1$ 的有限生成域扩张。
% P12-ERR-0495
令 $C$ 为与 $K$ 对应的抽象曲线。设 $V \subset K$ 为一个
$g^r_d$，$\Phi : C \to \mathbf{P}^r$ 为相应态射。
证明：若 $V$ 是完备线性系，且 $d$ 相对于 $C$ 的亏格足够大，
则 $C$ 的像包含在某个二次超曲面中\footnote{二次超曲面是次数为
$2$ 的超曲面，即 $\mathbf{P}^r$ 中某个次数为 $2$ 的齐次多项式的
零点集。}。（加分题：给出所需次数的良好界。）
\end{exercise}

\begin{exercise}
\label{exercises-exercise-surjective-map-a2-p2}
记号如情形 \ref{exercises-situation-curve-in-the-plane}。
设 $U \subset \mathbf{P}^2_k$ 为补集是 $D$ 的开子概形。
描述 $k$-代数 $A = \mathcal{O}_{\mathbf{P}^2_k}(U)$。
给出 $A$ 作为 $k$-代数的生成元个数的上界。
\end{exercise}




\section{概形，2014 年春季期末考试}
\label{exercises-section-final-exam-spring-2014}

\noindent
这些是哥伦比亚大学 2014 年春季一门概形课程的期末考试题。

\begin{exercise}[定义]
\label{exercises-exercise-definitions-spring-2014}
设 $(X, \mathcal{O}_X)$ 为概形。给出斜体概念的定义。
\begin{enumerate}
\item $X$ 在点 $x$ 处的{\it 局部环}；
\item $\mathcal{O}_X$-模的{\it 拟凝聚}层；
% P12-ERR-0496
\item $\mathcal{O}_X$-模的{\it 凝聚}层（请假设 $X$ 局部 Noether；
\item $X$ 的{\it 仿射开集}；
\item {\it 概形的有限态射} $X \to Y$。
\end{enumerate}
\end{exercise}

\begin{exercise}[定理]
\label{exercises-exercise-results-spring-2014}
对每一项，精确陈述讲课中讨论的一个相关非平凡事实。
\begin{enumerate}
\item 关于簇的多重亏格的双有理不变性；
\item 仿射态射是局部性质；
\item 态射的概形论纤维的拓扑；
\item 真性赋值判别法。
\end{enumerate}
\end{exercise}

\begin{exercise}
\label{exercises-exercise-miss-curve}
设 $X = \mathbf{A}^2_\mathbf{C}$，其中 $\mathbf{C}$ 为复数域。
所谓{\it 直线}，是指由一个线性方程 $ax + by + c = 0$
定义的 $X$ 的闭子概形，其中
$a, b, c \in \mathbf{C}$ 且 $(a, b) \not = (0, 0)$。
所谓{\it 曲线}，是指维数为 $1$ 的不可约（因而非空）闭子概形
$C \subset X$。所谓{\it 二次曲线}，是指由一个二次方程
$ax^2 + bxy + cy^2 + dx + ey + f = 0$ 定义的曲线，其中
$a, b, c, d, e, f \in \mathbf{C}$ 且
$(a, b, c) \not = (0, 0, 0)$。
\begin{enumerate}
\item 找一条曲线 $C$，使每条直线都与 $C$ 有非空交。
\item 找一条曲线 $C$，使每条直线与每条二次曲线都与 $C$ 有非空交。
\item 证明对每条曲线 $C$，都存在另一条曲线，使
$C \cap C' = \emptyset$。
\end{enumerate}
\end{exercise}

\begin{exercise}
\label{exercises-exercise-normal-bundle-exceptional-curve}
设 $k$ 为域。设 $b : X \to \mathbf{A}^2_k$ 为仿射平面在原点处的爆破。
换言之，若 $\mathbf{A}^2_k = \Spec(k[x, y])$，则
$X = \text{Proj}(\bigoplus_{n \geq 0} \mathfrak m^n)$，
其中 $\mathfrak m = (x, y) \subset k[x, y]$。
证明下列命题：
\begin{enumerate}
\item $b$ 在原点上的概形论纤维 $E$ 同构于 $\mathbf{P}^1_k$；
\item $E$ 是 $X$ 上的有效 Cartier 除子；
\item $\mathcal{O}_X(-E)$ 限制到 $E$ 后是次数为 $1$ 的线丛。
\end{enumerate}
（回顾：$\mathcal{O}_X(-E)$ 是 $E$ 在 $X$ 中的理想层。）
\end{exercise}

\begin{exercise}
\label{exercises-exercise-surjective-map-affine-variety-projective-variety}
设 $k$ 为域，$X$ 为 $k$ 上的射影簇。
证明存在 $k$ 上的仿射簇 $U$ 以及簇的满态射 $U \to X$。
\end{exercise}

\begin{exercise}
\label{exercises-exercise-vandermonde}
设 $k$ 为特征 $p > 0$ 且不同于 $2,3$ 的域。
考虑 $\mathbf{P}^n_k$ 中由下列方程定义的闭子概形 $X$：
$$
\sum\nolimits_{i = 0, \ldots, n} X_i = 0,\quad
\sum\nolimits_{i = 0, \ldots, n} X_i^2 = 0,\quad
\sum\nolimits_{i = 0, \ldots, n} X_i^3 = 0
$$
% P12-ERR-0497
对哪些数对 $(n, p)$，这个簇是奇异的？
\end{exercise}







\section{交换代数，2016 年秋季期末考试}
\label{exercises-section-final-exam-fall-2016}

\noindent
这些是哥伦比亚大学 2016 年秋季一门交换代数课程的期末考试题。

\begin{exercise}[定义]
\label{exercises-exercise-definitions-fall-2016}
设 $R$ 为环。给出斜体概念的定义。
\begin{enumerate}
\item $R$ 在素理想 $\mathfrak p$ 处的{\it 局部环}；
\item {\it 有限} $R$-模；
\item {\it 有限表现} $R$-模；
\item $R$ 是{\it 正则的}；
\item $R$ 是{\it 链状的}；
\item $R$ 是{\it Cohen--Macaulay 的}。
\end{enumerate}
\end{exercise}

\begin{exercise}[定理]
\label{exercises-exercise-results-fall-2016}
对每一项，精确陈述讲课中讨论的一个相关非平凡事实。
\begin{enumerate}
\item 正则环；
\item Cohen--Macaulay 模的相伴素理想；
\item 域上有限型整环的维数；
\item Chevalley 定理。
\end{enumerate}
\end{exercise}

\begin{exercise}
\label{exercises-exercise-finite-injective}
设 $A \to B$ 为满足下列条件的环映射：
\begin{enumerate}
\item $A$ 是极大理想为 $\mathfrak m$ 的局部环；
\item $A \to B$ 是有限环映射\footnote{回顾：这表示 $B$
作为 $A$-模是有限的。}；
\item $A \to B$ 是单射（把 $A$ 看作 $B$ 的子环）。
\end{enumerate}
证明存在素理想 $\mathfrak q \subset B$，使
$\mathfrak m = A \cap \mathfrak q$。
\end{exercise}

\begin{exercise}
\label{exercises-exercise-find-components}
设 $k$ 为域，$R = k[x, y, z, w]$。考虑理想
$I = (xy, xz, xw)$。$V(I) \subset \Spec(R)$ 的不可约分支是什么，
其维数分别是多少？
\end{exercise}

\begin{exercise}
\label{exercises-exercise-no-nonconstant-morphism}
设 $k$ 为域，$A = k[x, x^{-1}]$、$B = k[y]$。
证明任意 $k$-代数映射 $\varphi : A \to B$ 都把 $x$ 映到常数。
\end{exercise}

\begin{exercise}
\label{exercises-exercise-regular-over-Z}
考虑环 $R = \mathbf{Z}[x, y]/(xy - 7)$。
证明 $R$ 正则。
\end{exercise}

\noindent
给定 Noether 局部环 $(R, \mathfrak m, \kappa)$，对
$n \geq 0$ 令
$\varphi_R(n) = \dim_\kappa(\mathfrak m^n/\mathfrak m^{n + 1})$。

\begin{exercise}
\label{exercises-exercise-hilbert-function-allowed}
是否存在 Noether 局部环 $R$，使对所有 $n \geq 0$ 都有
$\varphi_R(n) = n + 1$？
\end{exercise}

\begin{exercise}
\label{exercises-exercise-hilbert-function-prohibited}
设 $R$ 为 Noether 局部环。假设
$\varphi_R(0) = 1$、$\varphi_R(1) = 3$、
$\varphi_R(2) = 5$。证明 $\varphi_R(3) \leq 7$。
\end{exercise}





\section{概形，2017 年春季期末考试}
\label{exercises-section-final-exam-spring-2017}

\noindent
这些是哥伦比亚大学 2017 年春季一门概形课程的期末考试题。

\begin{exercise}[定义]
\label{exercises-exercise-definitions-spring-2017}
设 $f : X \to Y$ 为概形态射。简要给出斜体概念的定义。
\begin{enumerate}
\item $f$ 在 $y \in Y$ 处的{\it 概形论纤维}；
\item $f$ 是{\it 有限态射}；
\item {\it 拟凝聚} $\mathcal{O}_X$-模；
% P12-ERR-0498
\item $X$ 是{\it 簇}；
\item $f$ 是{\it 光滑态射}；
\item $f$ 是{\it 真态射}。
\end{enumerate}
\end{exercise}

\begin{exercise}[定理]
\label{exercises-exercise-results-spring-2017}
对每一项，精确而简要地陈述讲课中讨论的一个相关非平凡事实。
\begin{enumerate}
\item 拟凝聚层的前推；
\item 射影簇上凝聚层的上同调；
\item 域上射影概形的 Serre 对偶；
\item Riemann--Hurwitz 公式。
\end{enumerate}
\end{exercise}

\begin{exercise}
\label{exercises-exercise-compute-degree}
设 $k$ 为代数闭域。设 $\ell > 100$ 为不同于 $k$ 的特征的素数。
设 $X$ 为仿射曲线
$$
y^\ell = x(x - 1)^3
$$
在 $\mathbf{A}^2_k$ 中所给出的非奇异射影模型。回答下列问题：
\begin{enumerate}
\item $X$ 的亏格是多少？
\item 给出 $X$ 的脉性\footnote{脉性是从 $X$ 到
$\mathbf{P}^1_k$ 的非常值态射的最小次数。}的一个上界。
\end{enumerate}
\end{exercise}

\begin{exercise}
\label{exercises-exercise-count-components}
设 $k$ 为代数闭域。设 $X$ 为 $k$ 上既约射影概形，
其所有不可约分支的维数都等于 $1$。
设 $\omega_{X/k}$ 为相对对偶模。
证明若 $\dim_k H^1(X, \omega_{X/k}) > 1$，则 $X$ 不连通。
\end{exercise}

\begin{exercise}
\label{exercises-exercise-sections-L-L-inverse}
给出概形 $X$ 与非平凡可逆 $\mathcal{O}_X$-模
$\mathcal{L}$ 的例子，使 $H^0(X, \mathcal{L})$ 与
$H^0(X, \mathcal{L}^{\otimes -1})$ 均非零。
\end{exercise}

\begin{exercise}
\label{exercises-exercise-in-product}
设 $k$ 为代数闭域，$g \geq 3$。
设 $X$ 与 $X'$ 为 $k$ 上亏格分别为 $g$ 与 $g + 1$ 的
光滑射影曲线。设 $Y \subset X \times X'$ 为曲线，并且投影
$Y \to X$ 与 $Y \to X'$ 都不是常值的。
证明 $Y$ 的非奇异射影模型的亏格 $\geq 2g + 1$。
\end{exercise}

\begin{exercise}
\label{exercises-exercise-finitely-many}
设 $k$ 为有限域，$g > 1$。勾勒下列命题的证明：
$k$ 上亏格为 $g$ 的光滑射影曲线仅有有限多个同构类。
（即使只是尝试作答也可得分。）
\end{exercise}








\section{交换代数，2017 年秋季期末考试}
\label{exercises-section-final-exam-fall-2017}

\noindent
这些是哥伦比亚大学 2017 年秋季一门交换代数课程的期末考试题。

\begin{exercise}[定义]
\label{exercises-exercise-definitions-fall-2017}
简要给出斜体概念的定义。
\begin{enumerate}
\item 函子 $F : \mathcal{A} \to \mathcal{B}$ 的{\it 左伴随}；
\item 域扩张 $L/K$ 的{\it 超越次数}；
\item 经典仿射簇 $X \subset k^n$ 上的{\it 正则函数}；
\item 拓扑空间上的{\it 层}；
\item {\it 局部环}；
\item 概形态射 $f : X \to Y$ 是{\it 仿射的}。
\end{enumerate}
\end{exercise}

\begin{exercise}[定理]
\label{exercises-exercise-results-fall-2017}
对每一项，精确而简要地陈述讲课中讨论的一个相关非平凡事实
（若不止一个，只需任选其一）。
\begin{enumerate}
\item Yoneda 引理；
\item Mayer--Vietoris；
\item 维数与上同调；
\item Hilbert 多项式；
\item 射影空间的对偶性。
\end{enumerate}
\end{exercise}

\begin{exercise}
\label{exercises-exercise-compute-dimension}
设 $k$ 为代数闭域。考虑带 Zariski 拓扑与坐标
$x_1, x_2, x_3, x_4, x_5$ 的 $k^5$ 中由方程
$$
x_1^2 - x_4 = 0,\quad
x_2^5 - x_5 = 0,\quad
x_3^2 + x_3 + x_4 + x_5 = 0
$$
给出的闭子集 $X$。$X$ 的维数是多少？为什么？
\end{exercise}

\begin{exercise}
\label{exercises-exercise-can-there-be}
设 $k$ 为域，$X = \mathbf{P}^1_k$ 为 $k$ 上的 $1$-维射影空间。
设 $\mathcal{E}$ 为有限局部自由 $\mathcal{O}_X$-模。
对 $d \in \mathbf{Z}$，用
$\mathcal{E}(d) = \mathcal{E} \otimes_{\mathcal{O}_X} \mathcal{O}_X(d)$
表示 $\mathcal{E}$ 的第 $d$ 次 Serre 扭曲，并令
$h^i(X, \mathcal{E}(d)) = \dim_k H^i(X, \mathcal{E}(d))$。
\begin{enumerate}
\item 为什么不存在满足
$h^0(X, \mathcal{E}) = 5$ 且 $h^0(X, \mathcal{E}(1)) = 4$
的 $\mathcal{E}$？
\item 为什么不存在满足
$h^1(X, \mathcal{E}(1)) = 5$ 且 $h^1(X, \mathcal{E}) = 4$
的 $\mathcal{E}$？
\item 对哪些 $a \in \mathbf{Z}$，可能存在 $X$ 上的向量丛
$\mathcal{E}$，使
$$
\begin{matrix}
h^0(X, \mathcal{E})\phantom{(1)} = 1 &
h^1(X, \mathcal{E})\phantom{(1)} = 1 \\
h^0(X, \mathcal{E}(1)) = 2 &
h^1(X, \mathcal{E}(1)) = 0 \\
h^0(X, \mathcal{E}(2)) = 4 &
h^1(X, \mathcal{E}(2)) = a
\end{matrix}
$$
\end{enumerate}
欢迎并鼓励不完整的答案。
\end{exercise}

\begin{exercise}
\label{exercises-exercise-banana}
设拓扑空间 $X$ 是两个闭子集 $Y$ 与 $Z$ 的并
$X = Y \cup Z$，其交记为 $W = Y \cap Z$。
用 $i : Y \to X$、$j : Z \to X$ 与 $k : W \to X$
表示包含映射。
\begin{enumerate}
\item 证明存在层的短正合列
$$
0 \to \underline{\mathbf{Z}}_X \to
i_*(\underline{\mathbf{Z}}_Y) \oplus
j_*(\underline{\mathbf{Z}}_Z) \to
k_*(\underline{\mathbf{Z}}_W) \to 0
$$
其中 $\underline{\mathbf{Z}}_X$ 表示 $X$ 上取值为
$\mathbf{Z}$ 的常层，等等。
\item 关于 $\mathbf{Z}$-系数的 $X$、$Y$、$Z$、$W$
上同调群之间的关系，能得出什么结论？
\end{enumerate}
\end{exercise}

\begin{exercise}
\label{exercises-exercise-cohomology-infinite-punctured}
设 $k$ 为域。令 $A = k[x_1, x_2, x_3, \ldots]$ 为无穷多个变量的
多项式环。用 $\mathfrak m$ 表示由所有变量生成的 $A$ 的极大理想。
令 $X = \Spec(A)$，$U = X \setminus \{\mathfrak m\}$。
\begin{enumerate}
\item 证明 $H^1(U, \mathcal{O}_U) = 0$。
提示：作 {\v C}ech 上同调计算。
\item 对 $i \geq 1$，你猜测 $H^i(U, \mathcal{O}_U)$ 是什么？
\end{enumerate}
\end{exercise}

\begin{exercise}
\label{exercises-exercise-principal}
设 $A$ 为局部环，$a \in A$ 为非零因子。
设 $I, J \subset A$ 为满足 $IJ = (a)$ 的理想。
证明理想 $I$ 是主理想，即由一个元素生成
（该元素也将是非零因子）。
\end{exercise}




\section{概形，2018 年春季期末考试}
\label{exercises-section-final-exam-spring-2018}

\noindent
这些是哥伦比亚大学 2018 年春季一门概形课程的期末考试题。

\begin{exercise}[定义]
\label{exercises-exercise-definitions-spring-2018}
简要给出斜体概念的定义。
设 $k$ 为代数闭域，$X$ 为 $k$ 上的射影曲线。
\begin{enumerate}
\item $k$ 上的{\it 光滑}代数；
\item $X$ 上可逆 $\mathcal{O}_X$-模的{\it 次数}；
\item $X$ 的{\it 亏格}；
\item $X$ 的{\it Weil 除子类群}；
\item $X$ 是{\it 超椭圆的}；
\item $k$ 上光滑射影曲面中两条曲线的{\it 交数}。
\end{enumerate}
\end{exercise}

\begin{exercise}[定理]
\label{exercises-exercise-results-spring-2018}
对每一项，精确而简要地陈述讲课中讨论的一个相关非平凡事实
（若不止一个，只需任选其一）。
\begin{enumerate}
\item Riemann--Hurwitz 定理；
\item Clifford 定理；
\item 光滑射影曲面之间映射的分解；
\item Hodge 指标定理；
\item 有限域上曲线的 Riemann 猜想。
\end{enumerate}
\end{exercise}

\begin{exercise}
\label{exercises-exercise-hyperelliptic}
设 $k$ 为代数闭域。设 $X \subset \mathbf{P}^3_k$ 为次数 $d$、
亏格 $\geq 2$ 的光滑曲线。假设 $X$ 不包含在任何平面中，
并且 $\mathbf{P}^3_k$ 中有一条直线 $\ell$ 与 $X$ 交于
$d - 2$ 个点。证明 $X$ 是超椭圆曲线。
\end{exercise}

\begin{exercise}
\label{exercises-exercise-singular}
设 $k$ 为代数闭域。设 $X$ 为射影曲线，其两两不同的奇点为
$p_1, \ldots, p_n$。说明为何 $X$ 的正规化的亏格至多为
$-n + \dim_k H^1(X, \mathcal{O}_X)$。
\end{exercise}

\begin{exercise}
\label{exercises-exercise-how-many-blowups}
设 $k$ 为域，$X = \Spec(k[x, y])$ 为仿射 $2$-空间。令
$$
I = (x^3, x^2y, xy^2, y^3) \subset k[x, y].
$$
令 $Y \subset X$ 为与 $I$ 对应的闭子概形。
设 $b : X' \to X$ 为理想 $(x, y)$ 的爆破，即仿射空间在原点处的
爆破。
\begin{enumerate}
\item 证明概形论逆像 $b^{-1}Y \subset X'$ 是有效 Cartier 除子。
% P12-ERR-0499
\item 给出理想 $J \subset k[x, y]$ 的例子，使
$I \subset J \subset (x, y)$，并且若 $Z \subset X$ 是与 $J$
对应的闭子概形，则概形论逆像 $b^{-1}Z$ 不是有效 Cartier 除子。
\end{enumerate}
\end{exercise}

\begin{exercise}
\label{exercises-exercise-rational-surface}
设 $k$ 为代数闭域。考虑下列类型的曲面：
\begin{enumerate}
\item $S = C_1 \times C_2$，其中 $C_1$ 与 $C_2$ 是光滑射影曲线；
\item $S = C_1 \times C_2$，其中 $C_1$ 与 $C_2$ 是光滑射影曲线，
且 $C_1$ 的亏格 $> 0$；
\item $S \subset \mathbf{P}^3_k$ 是次数 $4$ 的超曲面；
\item $S \subset \mathbf{P}^3_k$ 是次数 $4$ 的光滑超曲面。
\end{enumerate}
对每种类型，简要说明为何该类型的曲面类包含或不包含有理曲面。
\end{exercise}

\begin{exercise}
\label{exercises-exercise-self-square-line}
设 $k$ 为代数闭域。设 $S \subset \mathbf{P}^3_k$ 为次数 $d$
的光滑超曲面。假设 $S$ 包含直线 $\ell$。
% P12-ERR-0500
把 $\ell$ 看作 $S$ 上的除子时，它的自交数是多少？
\end{exercise}

\section{交换代数，2019 年秋季期末考试}
\label{exercises-section-final-exam-fall-2019}

\noindent
以下是哥伦比亚大学 2019 年秋季交换代数课程期末考试的试题。

\begin{exercise}[定义]
\label{exercises-exercise-definitions-fall-2019}
简要定义下列斜体概念。
\begin{enumerate}
\item Noether 拓扑空间的一个 {\it 可构造子集}，
\item $R$-模 $M$ 在素理想 $\mathfrak p$ 处的 {\it 局部化}，
\item Noether 局部环 $(A, \mathfrak m, \kappa)$ 上模的 {\it 长度}，
\item 环 $R$ 上的 {\it 射影模}，以及
\item Noether 局部环 $(A, \mathfrak m, \kappa)$ 上的
{\it Cohen--Macaulay 模}。
\end{enumerate}
\end{exercise}

\begin{exercise}[定理]
\label{exercises-exercise-results-fall-2019}
对下列每一项，准确而简要地陈述一个课堂上讨论过的相关非平凡事实
（若有多个，只需任选一个）。
\begin{enumerate}
\item 可构造集的像，
\item Hilbert 零点定理，
\item 域上有限型代数的维数，
\item Noether 正规化，以及
\item 正则局部环。
\end{enumerate}
\end{exercise}

\noindent
回顾：对于环 $R$ 与理想 $I \subset R$，$V(I)$ 表示点
$\mathfrak p \in \Spec(R)$ 所成的集合，其中 $I \subset \mathfrak p$。

\begin{exercise}[构造素理想]
\label{exercises-exercise-infinitely-many-primes}
构造无穷多个互不相同的素理想
$\mathfrak p \subset \mathbf{C}[x, y]$，
使 $V(\mathfrak p)$ 包含 $(x, y)$ 与 $(x - 1, y - 1)$。
\end{exercise}

\begin{exercise}[不存在这样的素理想]
\label{exercises-exercise-no-prime}
令 $R = \mathbf{C}[x, y, z]/(xy)$。简要论证：不存在素理想
$\mathfrak p \subset R$，使得
$V(\mathfrak p)$ 同时包含 $(x, y - 1, z - 5)$
与 $(x - 1, y, z - 7)$。
\end{exercise}

\begin{exercise}[Frobenius]
\label{exercises-exercise-frobenius}
令 $p$ 为素数（为简化公式，你可以假设 $p = 2$）。
令 $R$ 为满足 $R$ 中 $p = 0$ 的环。
\begin{enumerate}
\item 证明映射 $F : R \to R$，$x \mapsto x^p$，是环同态。
\item 证明 $\Spec(F) : \Spec(R) \to \Spec(R)$ 是恒等映射。
\end{enumerate}
\end{exercise}

\noindent
回顾：拓扑空间中点的一个 {\it 特化} $x \leadsto y$，仅仅意指
$y$ 位于 $x$ 的闭包中。若不存在异于 $x$ 与 $y$ 的点 $z$，
使 $x \leadsto z$ 且 $z \leadsto y$，则称 $x \leadsto y$
为 {\it 直接特化}。

\begin{exercise}[维数]
\label{exercises-exercise-dimension}
设索伯拓扑空间 $X$ 含有 $5$ 个互不相同的点
$x, y, z, u, v$，它们之间有如下特化关系：
$$
\xymatrix{
x \ar[d] \ar[r] & u & v \ar[l] \ar[dl] \\
y \ar[r] & z
}
$$
这样的 $X$ 至少可以是多少维？若 $X$ 是域上有限型代数的谱，
且 $x \leadsto u$ 是直接特化，那么关于特化
$v \leadsto z$ 你能得出什么结论？
\end{exercise}

\begin{exercise}[Tor 计算]
\label{exercises-exercise-tor-computation}
令 $R = \mathbf{C}[x, y, z]$。令 $M = R/(x, z)$ 且
$N = R/(y, z)$。对于哪些 $i \in \mathbf{Z}$，
$\text{Tor}_i^R(M, N)$ 非零？
\end{exercise}

\begin{exercise}
\label{exercises-exercise-depth-goes-up}
令 $A \to B$ 为 Noether 局部环之间的平坦局部同态。
证明：若 $A$ 的深度为 $k$，则 $B$ 的深度至少为 $k$。
\end{exercise}

\section{代数几何，2020 年春季期末考试}
\label{exercises-section-final-exam-spring-2020}

\noindent
以下是哥伦比亚大学 2020 年春季代数几何课程期末考试的试题。

\begin{exercise}[定义]
\label{exercises-exercise-definitions-spring-2020}
简要定义下列斜体概念。
\begin{enumerate}
\item 一个 {\it 概形}，
\item 概形的一个 {\it 态射}，
\item 概形上的一个 {\it 拟凝聚模}，
\item 域 $k$ 上的一个 {\it 代数簇}，
\item 域 $k$ 上的一条 {\it 曲线}，
\item 概形的一个 {\it 有限态射}，
\item 拓扑空间上阿贝尔群层 $\mathcal{F}$ 的 {\it 上同调}，其所在空间为 $X$，
\item 概形 $X$（维数为 $d$，在域 $k$ 上适当）上的 {\it 对偶层}，以及
\item 从代数簇 $X$ 到代数簇 $Y$ 的一个 {\it 有理映射}。
\end{enumerate}
\end{exercise}

\begin{exercise}[定理]
\label{exercises-exercise-results-spring-2020}
对下列每一项，准确而简要地陈述一个课堂上讨论过的相关非平凡事实
（若有多个，只需任选一个）。
\begin{enumerate}
\item Noether 拓扑空间 $X$（维数为 $d$）上阿贝尔层的上同调，
\item 域 $k$ 上光滑代数簇的微分层 $\Omega^1_{X/k}$，
\item 光滑射影代数簇的对偶层 $\omega_X$（代数簇为 $X$，定义在域 $k$ 上），
\item 代数闭域 $k$ 上一条光滑适当的亏格 $0$ 曲线，以及
\item 次数为 $d$ 的平面曲线的亏格。
\end{enumerate}
\end{exercise}

\begin{exercise}
\label{exercises-exercise-coh-P1-infty-doubled}
设 $k$ 为域，$X$ 为 $k$ 上的概形。假设 $X = X_1 \cup X_2$
是一个开覆盖，其中 $X_1$、$X_2$ 都同构于 $\mathbf{P}^1_k$，
而 $X_1 \cap X_2$ 同构于 $\mathbf{A}^1_k$。（例如，将
$\mathbf{P}^1_k$ 的 $\infty$ 点复制一份就可得到这样的概形。）
证明 $\dim_k H^1(X, \mathcal{O}_X)$ 是无穷的。
\end{exercise}

\begin{exercise}
\label{exercises-exercise-no-morphism}
设 $k$ 为代数闭域。令 $Y$ 为亏格 $10$ 的光滑射影曲线。
为一条光滑射影曲线 $X$ 的亏格给出一个好的下界，使得存在一个
非恒定态射 $f : X \to Y$，且它不是同构。
\end{exercise}

\begin{exercise}
\label{exercises-exercise-element-order-n-in-pic}
设 $k$ 为特征 $0$ 的代数闭域。令
$$
X : T_0^d + T_1^d - T_2^d = 0 \subset \mathbf{P}^2_k
$$
为次数 $d \geq 3$ 的 Fermat 曲线。考虑 $X$ 上的闭点
$p = [1 : 0 : 1]$ 与 $q = [0 : 1 : 1]$。置 $D = [p] - [q]$。
\begin{enumerate}
\item 证明 $D$ 在 Weil 除子类群中非平凡。
\item 证明 $d D$ 在 Weil 除子类群中平凡。
（提示：尝试证明 $d[p]$ 与 $d[q]$ 都是平面中一条直线
与 $X$ 的交。）
\end{enumerate}
\end{exercise}

\begin{exercise}
\label{exercises-exercise-number-equations}
设 $k$ 为代数闭域。考虑 $2$-重嵌入
$$
\varphi : \mathbf{P}^2 \longrightarrow \mathbf{P}^5
$$
按照课堂材料/记号，这是态射
$$
\varphi = \varphi_{\mathcal{O}_{\mathbf{P}^2}(2)} :
\mathbf{P}^2
\longrightarrow
\mathbf{P}(\Gamma(\mathbf{P}^2, \mathcal{O}_{\mathbf{P}^2}(2)))
$$
用齐次坐标表示，它为
$$
[a_0 : a_1 : a_2] \longmapsto
[a_0^2 : a_0a_1 : a_0a_2 : a_1^2 : a_1a_2 : a_2^2]
$$
这是一个闭浸入（请直接使用这一点）。令
$I \subset k[T_0, \ldots, T_5]$ 为 $\varphi(\mathbf{P}^2)$ 的齐次理想；
也就是说，齐次部分 $I_d$ 的元素是次数为 $d$ 的齐次多项式
$F(T_0, \ldots, T_5)$，它们在闭子概形 $\varphi(\mathbf{P}^2)$
上的限制为零。计算 $\dim_k I_d$ 关于 $d$ 的函数。
\end{exercise}

\begin{exercise}
\label{exercises-exercise-dualizing-sheaf-rank-1}
设 $k$ 为代数闭域。令 $X$ 为维数为 $d$ 的 $k$-上适当概形，
其对偶模为 $\omega_X$。已知：
\begin{enumerate}
\item
$\text{Ext}^i_X(\mathcal{F}, \omega_X) \times
H^{d - i}(X, \mathcal{F}) \to H^d(X, \omega_X) \xrightarrow{t} k$
对所有 $i$ 以及所有相干 $\mathcal{O}_X$-模 $\mathcal{F}$，
均为非退化的，以及
\item $\omega_X$ 是某个秩 $r$ 的有限局部自由模。
\end{enumerate}
证明 $r = 1$。（提示：考察取一个支集在闭点上的适当模时会发生什么，
例如取 $\mathcal{F}$。）
\end{exercise}

\section{交换代数，2021 年秋季期末考试}
\label{exercises-section-final-exam-fall-2021}

\noindent
以下是哥伦比亚大学 2021 年秋季交换代数课程期末考试的试题。

\begin{exercise}[定义]
\label{exercises-exercise-definitions-fall-2021}
简要定义下列斜体概念。
\begin{enumerate}
\item 环 $A$ 的一个 {\it 乘法子集}，
\item 一个 {\it Artinian 环} $A$，
\item 环 $A$ 的 {\it 谱}作为拓扑空间，
\item 一个 {\it 平坦环映射} $A \to B$，
\item $A$ 中素理想 $\mathfrak p$ 的 {\it 高度}，以及
\item 环 $A$ 上的函子 {\it $\text{Tor}^A_i(-, -)$}。
\end{enumerate}
\end{exercise}

\begin{exercise}[定理]
\label{exercises-exercise-results-fall-2021}
对下列每一项，准确而简要地陈述一个课堂上讨论过的相关非平凡事实
（若有多个，只需任选一个）。
\begin{enumerate}
\item Artinian 环，
\item 平坦性与素理想，
\item 对于 Noether 局部环 $(A, \mathfrak m)$，$A/\mathfrak m^n$ 的长度，
\item 万有链式 Noether 环的维数公式，
\item Noether 局部环的完备化，以及
\item Artinian 局部环的 Matlis 对偶。
\end{enumerate}
\end{exercise}

\begin{exercise}[单位]
\label{exercises-exercise-simple-units}
$\mathbf{Z}[x, 1/x]$ 的单位群作为阿贝尔群具有什么结构？
无需解释。
\end{exercise}

\begin{exercise}[理想]
\label{exercises-exercise-ideals}
令 $A = \mathbf{F}_2[x, y]/(x^2, xy, y^2)$，并将
$\overline{x}$ 与 $\overline{y}$ 记为 $x$ 与 $y$ 在 $A$ 中的像。
列出 $A$ 的理想。无需解释。
\end{exercise}

\begin{exercise}[Tor 与 Ext]
\label{exercises-exercise-compute-tor-ext}
设 $(A, \mathfrak m, \kappa)$ 为 Noether 局部环。
置 $\varphi(n) = \dim_\kappa \mathfrak m^n/\mathfrak m^{n + 1}$。
\begin{enumerate}
\item 证明 $\text{Tor}_1^A(A/\mathfrak m^n, \kappa)$
作为 $\kappa$-向量空间的维数为 $\varphi(n)$。
\item 证明 $\text{Ext}^1_A(A/\mathfrak m^n, \kappa)$
作为 $\kappa$-向量空间的维数为 $\varphi(n)$。
\end{enumerate}
\end{exercise}

\begin{exercise}[两个向量]
\label{exercises-exercise-two-vectors}
% P12-ERR-0501
令 $A = \mathbf{Z}[a_1, a_2, a_3, b_1, b_2, b3]$。
令 $a = (a_1, a_2, a_3)$ 且 $b = (b_1, b_2, b_3)$ 为 $A^{\oplus 3}$ 中的向量。
考虑集合
$$
Z = \{\mathfrak p \in \Spec(A) \mid a, b
\text{ map to linearly dependent vectors of }
\kappa(\mathfrak p)^{\oplus 3}\}
$$
\begin{enumerate}
% P12-ERR-0502
\item 证明 $Z$ 是 $\Spec(A)$ 的闭子集。
\item $Z$ 的维数 $\dim(Z)$ 是多少？
\item 如果将 $\mathbf{Z}$ 替换为一个域，$\dim(Z)$ 会怎样变化？
\end{enumerate}
\end{exercise}

\begin{exercise}[内射模]
\label{exercises-exercise-injective-artinian-local}
设 $(A, \mathfrak m, \kappa)$ 为 Artinian 局部环。
% P12-ERR-0503
假设 $A$ 作为 $A$-模是内射的。证明
$\Hom_A(\kappa, A)$ 的维数为 $1$ 是一个
$\kappa$-向量空间。
\end{exercise}

\section{代数几何，2022 年春季期末考试}
\label{exercises-section-final-exam-spring-2022}

\noindent
以下是哥伦比亚大学 2022 年春季代数几何课程期末考试的试题。

\begin{exercise}[定义]
\label{exercises-exercise-definitions-spring-2022}
简要定义下列斜体概念。
\begin{enumerate}
\item 一个 {\it 概形}，
\item 概形 $X$ 上的一个 {\it 拟凝聚模}，
\item 概形的一个 {\it 平坦}态射 $X \to Y$，
\item 概形的一个 {\it 有限}态射 $X \to Y$，
\item 基概形 $S$ 上的一个 {\it 群概形} $G$，
\item 基概形 $S$ 上的一个 {\it 代数簇族}，
\item 域 $k$ 上代数簇 $X$ 的闭点 $x$ 的 {\it 次数}，
\item 点 $p = (a_0 : \ldots : a_n)$ 在 $\mathbf{P}^n(\mathbf{Q})$
中的通常 {\it 对数高度}，以及
\item 一个 {\it $C_i$ 域}。
\end{enumerate}
\end{exercise}

\begin{exercise}[定理]
\label{exercises-exercise-results-spring-2022}
对下列每一项，准确而简要地陈述一个课堂上讨论过的相关非平凡事实
（若有多个，只需任选一个）。
\begin{enumerate}
\item 从概形 $X$ 到仿射概形 $\Spec(A)$ 的态射，
\item 仿射概形 $X$ 上拟凝聚模 $\mathcal{F}$ 的上同调，
\item 域 $k$ 上 $\mathbf{P}^1_k$ 的 Picard 群，
\item 当 $S$ 为 Noether 概形时，平坦适当态射
$X \to S$ 的纤维维数，
\item 概形 $S$ 上的 $\mathbf{G}_m$-等变模，以及
\item 交的 Bézout 定理（若愿意，可限制在某个特殊情形）。
\end{enumerate}
\end{exercise}

\begin{exercise}[三次超曲面]
\label{exercises-exercise-cubics}
令 $F \in \mathbf{C}[T_0, \ldots, T_n]$ 为次数 $3$ 的齐次多项式。
给定 $3$ 个向量 $x, y, z \in \mathbf{C}^{n + 1}$，考虑条件
$$
(*)\quad
F(\lambda x + \mu y + \nu z) = 0 \text{ 于 } \mathbf{C}[\lambda, \mu, \nu]
$$
\begin{enumerate}
\item 所有 $x, y, z$ 的选择所成空间的维数是多少？
\item 条件 (*) 对 $x$、$y$、$z$ 的坐标给出了多少个方程？
\item 满足 (*) 的所有三元组 $x, y, z$ 所成空间的预期维数是多少？
\item $x, y, z$ 线性相关的所有三元组所成空间的维数是多少？
\item 推出：在 $\mathbf{P}^n$ 中次数为 $3$ 的超曲面上，
当 $n \geq a$（其中 $a$ 由你求出）时，我们预期能找到维数为 $2$
的线性子空间。
\end{enumerate}
\end{exercise}

\begin{exercise}[高度]
\label{exercises-exercise-heights}
设 $K$ 为域。令 $h_n : \mathbf{P}^n(K) \to \mathbf{R}$（$n \geq 0$）
为满足课堂中讨论的 $2$ 条公理的一族函数。令 $X$ 为
$K$ 上的射影代数簇。令 $\mathcal{L}$ 为可逆的
$\mathcal{O}_X$-模，并回顾课堂中构造的关联高度函数
$h_\mathcal{L} : X(K) \to \mathbf{R}$。令 $\alpha : X \to X$
为 $X$ 上定义在 $K$ 上的自同构。
\begin{enumerate}
\item 证明 $P \mapsto h_\mathcal{L}(\alpha(P))$
与函数 $h_{\alpha^*\mathcal{L}}$ 相差有界量。（提示：回顾若存在态射
$\varphi : X \to \mathbf{P}^n$，使得
$\mathcal{L} = \varphi^*\mathcal{O}_{\mathbf{P}^n}(1)$，则按构造
$h_\mathcal{L}(P) = h_n(\varphi(P))$，并利用这一点。
一般地，将 $\mathcal{L}$ 写成两个此类模之差。）
\item 假设 $h_\mathcal{L}(P) - h_\mathcal{L}(\alpha(P))$
在 $X(K)$ 上无界。证明对于
$\mathcal{N} = \mathcal{L} \otimes \alpha^*\mathcal{L}^{\otimes -1}$，
$h_\mathcal{N}$ 在 $X(K)$ 上无界。
\item 假设 $X$ 是椭圆曲线，且 $\mathcal{L}$ 是 $X$ 上对称的充足可逆模，
使得 $h_\mathcal{L}$ 在 $X(K)$ 上无界。证明存在一个次数为 $0$ 的可逆模
$\mathcal{N}$，使 $h_\mathcal{N}$ 无界。
（提示：回顾 $X$ 是维数为 $1$ 的阿贝尔簇。因此由课堂结果，
$h_\mathcal{L}$ 除一个常数外是二次的。选取适当的点
$P_0 \in X(K)$。令 $\alpha : X \to X$ 为按 $P_0$ 平移。
考虑 $P \mapsto h_\mathcal{L}(P) - h_\mathcal{L}(P + P_0)$。
应用你在上面证明的结果。）
\end{enumerate}
\end{exercise}

\begin{exercise}[单态射]
\label{exercises-exercise-monomorphisms}
设 $f : X \to Y$ 是概形范畴中的单态射：对于任意一对概形态射
$a, b : T \to X$，若 $f \circ a = f \circ b$，则 $a = b$。
证明 $f$ 在点上是单射。% P12-ERR-0504
你的论证还说明了其他什么吗？
\end{exercise}

\begin{exercise}[不动点]
\label{exercises-exercise-fix-points}
设 $k$ 为代数闭域。
\begin{enumerate}
\item 若 $G = \mathbf{G}_{m, k}$，证明当 $G$ 作用于 $k$ 上的射影代数簇
$X$ 时，该作用有不动点；即证明存在 $x \in X(k)$，使得对所有
$g \in G(k)$ 都有 $a(g, x) = x$。
\item 对 $G = (\mathbf{G}_{m, k})^n$ 作同样证明，其中它等于
$n \geq 1$ 个乘法群的直积。
\item 给出一个连通群概形 $G$ 作用于光滑射影代数簇 $X$ 的例子，
使该作用没有不动点。
\end{enumerate}
\end{exercise}

\section{代数几何，2025 年春季期末考试}
\label{exercises-section-final-exam-spring-2025}

\noindent
以下是哥伦比亚大学 2025 年春季代数几何课程期末考试的试题。

\begin{exercise}[定义]
\label{exercises-exercise-definitions-spring-2025}
简要定义下列斜体概念。
\begin{enumerate}
\item 一个 {\it 概形} $X$，
\item 概形 $X$ 上的一个 {\it 拟凝聚模}，
\item 概形 $X$ 的 {\it Picard 群}，
% P12-ERR-0505
\item 给定概形态射 $f : X \to Y$ 与 $\mathcal{O}_X$-模
$\mathcal{F}$，其 {\it 推前} $f_*\mathcal{F}$，
\item 概形的一个 {\it 闭浸入}，以及
\item 给定域 $k$ 上的一个 {\it 代数簇}。
\end{enumerate}
\end{exercise}

\begin{exercise}[定理]
\label{exercises-exercise-results-spring-2025}
% P12-ERR-0506
准确而简洁地陈述一个课堂上讨论过的、与下列每项相关的非平凡事实
（若有多个，只需任选一个）。
\begin{enumerate}
\item 概形 $X$ 的拓扑空间，
\item 概形范畴上的函子的一个可表性判据，
\item 代数闭域 $k$ 上光滑射影曲线 $X$ 的 Picard 函子，
\item 代数闭域 $k$ 上光滑射影曲线 $X$ 的 Picard 群中的挠元，
\item 一个关于光滑射影代数簇 $X$、$Y$、$Z$ 的乘积
$X \times Y \times Z$ 的 Picard 群的定理，其中这些代数簇定义在
代数闭域 $k$ 上。
\end{enumerate}
\end{exercise}

\begin{exercise}
\label{exercises-exercise-affine-line-infinite-nr-points}
设 $k$ 为域。解释为什么 $\Spec(k[x])$ 有无穷多个点。
\end{exercise}

\begin{exercise}
\label{exercises-exercise-hypersurface-variety}
设 $k$ 为域，$x_1, \ldots, x_n$ 为变量。
令 $f \in k[x_1, \ldots, x_n]$ 为非常数多项式。
令 $X = \Spec(k[x_1, \ldots, x_n]/(f))$。为了使 $X$ 成为
$k$ 上的代数簇，$f$ 应具有什么性质？
\end{exercise}

\begin{exercise}
\label{exercises-exercise-find-singular-point}
找出 $\mathbf{C}[x, y]/(x^5 - 5xy^4 + 20y - 16)$ 的谱上
一个奇点（无需找出全部奇点），并将其视为 $\mathbf{C}$ 上的代数簇。
\end{exercise}

\begin{exercise}
\label{exercises-exercise-kernel-restriction-pic-curve}
设 $X$ 为代数闭域 $k$ 上的光滑射影曲线。令 $x \in X(k)$
为 $X$ 上的闭点。令 $U = X \setminus \{x\}$。
关于限制映射 $\Pic(X) \to \Pic(U)$ 的核，你能说些什么？
\end{exercise}

\begin{exercise}
\label{exercises-exercise-sections-degree-2g-minus-4}
令 $X$ 为代数闭域 $k$ 上亏格 $g \geq 100$ 的光滑射影曲线。
% P12-ERR-0507
对于次数为 $2g - 4$ 的 $X$ 上可逆 $\mathcal{O}_X$-模
$\mathcal{L}$，$h^0(\mathcal{L}) = \dim_k H^0(X, \mathcal{L})$
可以取哪些值？请尽可能完整地回答。
\end{exercise}
% Generated mechanically from the frozen guide.tex target.
% Do not edit this generated file; edit the bound locale source instead.

% OK, start here.
%

\setcounter{chapter}{111}
\chapter{文献指南}


\phantomsection
\label{guide-section-phantom}


\section{简短导论文章}
\label{guide-section-short-introductions}


\begin{itemize}
\item Barbara Fantechi: \emph{Stacks for Everybody} \cite{fantechi_stacks}
\item Dan Edidin: \emph{What is a stack?} \cite{edidin_whatis}
\item Dan Edidin: \emph{Notes on the construction of the moduli space of
curves} \cite{edidin_notes}
\item Angelo Vistoli: \emph{Intersection theory on algebraic
stacks and on their moduli spaces}，尤其是附录。
\cite{vistoli_intersection}
\end{itemize}


\section{经典参考文献}
\label{guide-section-classics}

\begin{itemize}
\item Mumford: \emph{Picard groups of moduli problems}
\cite{mumford_picard}
\begin{quote}
本文从未使用 ''stack'' 一词，但这一概念隐含其中；作者计算了椭圆曲线模叠的 Picard 群。
\end{quote}
\item Deligne, Mumford: \emph{The irreducibility of the space of curves of
given genus} \cite{DM}
\begin{quote}
这篇有影响力的论文引入了 ''algebraic stacks''；如今通常称为 Deligne--Mumford 叠（具有可表示对角且允许由概形给出的 \'etale 展示的叠）。其中有许多 \emph{未给出证明} 的基础性结果。论文利用叠给出了亏格为 $g$ 的曲线模空间不可约性的两个证明。
\end{quote}
\item Artin: \emph{Versal deformations and algebraic stacks}
\cite{ArtinVersal}
\begin{quote}
本文引入了推广 Deligne--Mumford 叠的 ''algebraic stacks''，现在通常称为 \emph{Artin stacks}；其对角可表示并允许由概形给出的光滑展示。本文给出称为 Artin 判据的形变论判据，可在不显式构造展示的情况下证明给定模叠是 Artin 叠。
\end{quote}
\end{itemize}




\section{书籍与在线讲义}
\label{guide-section-books}


\begin{itemize}
\item Laumon, Moret-Bailly: \emph{Champs Alg\'ebriques} \cite{LM-B}
\begin{quote}
这是目前关于叠且包含许多基础结果的最详尽参考书。它假定读者熟悉代数空间，并频繁引用 Knutson 的书 \cite{Kn}。第 12 章关于代数叠 lisse-\'etale site 的函子性存在一个错误。无需担心，因为 Martin Olsson 已修补该错误（见 \cite{olsson_sheaves}），其余各章的结果（可能需要略作修改）都是正确的。
\end{quote}
\item The Stacks Project Authors: \emph{Stacks Project}
\cite{stacks-project}.
\begin{quote}
你正在阅读它！
\end{quote}
\item Anton Geraschenko:
\emph{Lecture notes for Martin Olsson's class on stacks} \cite{olsson_stacks}
\begin{quote}
本课程先系统发展代数空间理论，再引入代数叠（第 27 讲才首次定义！）。除基本性质外，课程涵盖 Deligne--Mumford 条件与具有非分歧对角之间的等价性、Artin 叠上的 lisse-\'etale site、拟凝聚层理论、Keel--Mori 定理、上同调下降以及 gerbe（及其与 Brauer 群的关系）。另有一些练习。
\end{quote}
\item
Behrend, Conrad, Edidin, Fantechi, Fulton, G\"ottsche, and Kresch:
\emph{Algebraic stacks}, online notes for a book being currently written
\cite{stacks_book}
\begin{quote}
本书旨在友好地介绍叠，重点放在例子与应用上，不要求读者具有高深背景。不同于 \cite{LM-B}，它不假定读者已经掌握代数空间理论；相反，先引入 Deligne--Mumford 叠，并将代数空间作为特殊情形，目标之一是发展足够理论以证明 \cite{DM} 中的断言。第二部分将发展 Artin 叠的一般理论。目前 Kresch 网站上只有书的一小部分可用。
\end{quote}
\item Olsson, Martin: \emph{Algebraic spaces and stacks}, \cite{olsson_book}
\begin{quote}
这是代数空间与代数叠的高度推荐入门读物，起点是已经掌握 Hartshorne 代数几何的读者。
\end{quote}
\end{itemize}


\section{关于叠基础的相关参考文献}
\label{guide-section-related}


\begin{itemize}
\item Vistoli:
\emph{Notes on Grothendieck topologies, fibered categories and descent theory}
\cite{vistoli_fga}
\begin{quote}
包含关于纤维范畴、叠以及 fpqc 拓扑中的下降理论的有用事实，并给出严谨证明。
\end{quote}
\item Knutson: \emph{Algebraic Spaces} \cite{Kn}
\begin{quote}
本书源自作者在 Michael Artin 指导下的博士论文，包含代数空间理论的基础。书 \cite{LM-B} 经常引用它。另见 Artin 关于代数空间的论文：\cite{Artin-Algebraic-Approximation}、
\cite{ArtinI}、\cite{Artin-Implicit-Function}、
\cite{ArtinII}、\cite{Artin-Construction-Techniques}、
\cite{Artin-Algebraic-Spaces}、\cite{Artin-Theorem-Representability} 以及
\cite{ArtinVersal}
\end{quote}
\item Grothendieck et al, \emph{Th\'eorie des Topos et Cohomologie \'Etale des
Sch\'emas I, II, III}，亦称 SGA4 \cite{SGA4}
\begin{quote}
第 1 卷包含关于宇宙、site 与纤维范畴的许多一般事实。''champ'' 一词（法语中表示 ''stack''）出现在 Deligne 的 Expos\'e XVIII。
\end{quote}
\item Jean Giraud: \emph{Cohomologie non ab\'elienne} \cite{giraud}
\begin{quote}
本书讨论一般 site 上的纤维范畴、叠、torsor 与 gerbe，但不讨论代数叠。例如，若 $G$ 是 $X$ 上的阿贝尔群层，则如同 $H^1(X, G)$ 可与 $G$-torsor 识别一样，$H^2(X, G)$ 可与适当定义的 $G$-gerbe 集合识别。当 $G$ 非阿贝尔时，$H^2(X, G)$ 定义为 $G$-gerbe 的集合。
\end{quote}
\item Kelly and Street: \emph{Review of the elements of 2-categories}
\cite{kelly-street}
\begin{quote}
叠的范畴形成一个 2-范畴，虽然后者是 2-范畴的一种简单类型，其中 2-态射都是可逆的。这是关于一般 2-范畴的参考文献。作者没有使用过它，因而无法评价其有用程度；另请注意 \cite{stacks-project} 包含一些关于 2-范畴的基础内容。
\end{quote}
\end{itemize}




\section{文献中的论文}
\label{guide-section-papers}

\noindent
下面列出包含叠与代数空间基本结果的研究论文。摘要的目的仅是指出对叠理论有贡献的结果；在许多情况下，这些结果只是论文主要目标的附带结果。我们将论文分为若干类别，有些论文会属于多个类别。


\subsection{形变理论与代数叠}
\label{guide-subsection-deformation-theory}

\noindent
Artin 的前三篇论文虽不涉及叠，却包含强有力的结果；前两篇对
\cite{ArtinVersal} 至关重要。
\begin{itemize}
\item Artin: \emph{Algebraic approximation of structures over
complete local rings} \cite{Artin-Algebraic-Approximation}
\begin{quote}
论文证明，在适当的温和假设下，任意有效形式形变都可以近似：若
$F: (\Sch/S) \to (\textit{Sets})$ 是一个逆变函子，
在有限型域或优良 DVR 上局部有限表示（基概形为 $S$），且 $s \in S$，
$\hat{\xi} \in F(\hat{\mathcal{O}}_{S, s})$ 是有效形式形变，
则对任意 $n > 0$，存在残余平凡的 \'etale 邻域
$(S', s') \to (S, s)$ 以及 $\xi' \in F(S')$，使得 $\xi'$ 与
$\hat{\xi}$ 一致到 $n$ 阶（即在
$F(\mathcal{O}_{S, s} / \mathfrak m^n)$ 中具有相同限制）。
\end{quote}
\item
Artin: \emph{Algebraization of formal moduli I} \cite{ArtinI}
\begin{quote}
论文证明，在温和假设下任意有效形式 versal 形变都可代数化。设
$F: (\Sch/S) \to (\textit{Sets})$ 为逆变函子，在有限型域或优良 DVR
上局部有限表示（基概形为 $S$），$s \in S$ 是局部闭点，$\hat A$ 是完备 Noether 局部
$\mathcal{O}_S$-代数，剩余域 $k'$ 是 $k(s)$ 的有限扩张，
并且 $\hat{\xi} \in F(\hat A)$ 是元素 $\xi_0 \in F(k')$ 的有效形式
versal 形变。则存在有限型 $S$-概形 $X$、闭点 $x \in X$（剩余域
$k(x) = k'$）以及 $\xi \in F(X)$，存在同构
$\hat{\mathcal{O}}_{X, x} \cong \hat{A}$，并且在每个
$F(\hat A / \mathfrak m^n)$ 中识别 $\xi$ 与 $\hat{\xi}$ 的限制。
若 $\hat{\xi}$ 是普遍形变，则代数化是唯一的。文中还给出了 Hilbert
概形与 Picard 概形可表示性的应用。
\end{quote}
\item Artin: \emph{Algebraization of formal moduli. II}
\cite{ArtinII}
\begin{quote}
粗略地说，论文证明：若能在 $Y' \subset X'$ 附近形式局部地收缩闭子集
$Y'$，则存在全局态射 $X' \to X$ 收缩 $Y$，其中 $X$ 是代数空间。
\end{quote}
\item
Artin: \emph{Versal deformations and algebraic stacks} \cite{ArtinVersal}
\begin{quote}
这篇重要论文建立在
\cite{Artin-Algebraic-Approximation} 与 \cite{ArtinI} 的工作之上。
论文引入 Artin 判据：通过检验形变论性质即可证明叠的代数性。更准确地说
（但仍非完全严格地说），Artin 在点 $x \in \mathcal{X}(k)$ 附近为极限保持叠
$\mathcal{X}$ 构造如下展示：假设叠 $\mathcal{X}$ 满足 Schlessinger
判据（\cite{Sch}），则存在 $x$ 的形式 versal 形变
$\hat{\xi} \in \lim \mathcal{X}(\hat A / \mathfrak m^n)$。若形式形变有效
（即 $\mathcal{X}(\hat{A}) \to \lim \mathcal{X}(\hat A / \mathfrak m^n)$
为双射），则得到有效形式 versal 形变 $\xi \in \mathcal{X}(\hat A)$。
利用 \cite{ArtinI} 的结果，可取有限型概形 $U$ 及元素
$\xi_U: U \to \mathcal{X}$，它在 $x$ 上方的点 $u \in U$ 处形式 versal。
若进一步假设 $\mathcal{X}$ 具有满足某些条件的形变与障碍理论（即与
\'etale 局部化、完备化兼容且满足可构造性条件），则第 4 节证明形式 versal
性是开条件；因此缩小 $U$ 后，$U \to \mathcal{X}$ 光滑。
Artin 还证明，任何具有概形 fppf 展示的叠都有概形光滑展示；特别地，
可用平坦、分离、有限表示的群概形构造商叠。
\end{quote}
\item Conrad, de Jong: \emph{Approximation of Versal Deformations}
\cite{conrad-dejong}
\begin{quote}
本文通过应用 Popescu 的强结果研究 Artin 的代数化结果：若 $A$ 是
Noether 环、$B$ 是 Noether $A$-代数，则映射 $A \to B$ 为正规态射，当且仅当
$B$ 是光滑 $A$-代数的直接极限。容易看出，Popescu 结果蕴含任意优良概形
上的 Artin 近似（优良性意味着对局部环 $A$，从 Hensel 化到完备化的映射
$A^{\text{h}} \to \hat A$ 为正规态射）。
论文给出 \cite{ArtinI} 主定理的一个适用于任意优良基概形、任意
$s \in S$ 的 ''groupoid'' 推广。因此 \cite{ArtinVersal} 中的结果在任意
优良基上成立。论文还讨论代数化的 \'etale-局部唯一性，以及对象的自同构群
是否自然作用在代数化的 Hensel 化上。
\end{quote}
\item Jason Starr: \emph{Artin's axioms, composition, and moduli spaces}
\cite{starr_artin}
\begin{quote}
论文证明 Artin 的代数化公理与 1-态射的复合相容。
\end{quote}
\item Martin Olsson: \emph{Deformation theory of representable
morphism of algebraic stacks} \cite{olsson_deformation}
\begin{quote}
本文用余切复形将概形态射的标准形变理论结果推广到代数叠的可表示态射。
这些结果不能视为 Illusie 一般理论的推论，因为可表示态射
$X \to \mathcal{X}$ 的余切复形并非由环化拓扑的态射定义（lisse-\'etale
site 不具有函子性）。
\end{quote}
\end{itemize}

\subsection{粗模空间}
\label{guide-subsection-coarse-moduli-spaces}

\noindent
讨论粗模空间的论文。
\begin{itemize}
\item Keel, Mori: \emph{Quotients in Groupoids} \cite{K-M}
\begin{quote}
长期以来，人们似乎以 ''folklore'' 形式知道分离 Deligne--Mumford 叠具有粗模空间。本文给出下述定理的严格（虽简略）证明：若 $\mathcal{X}$ 是 Noether 基概形上局部有限型的 Artin 叠，且惯性叠 $I_\mathcal{X} \to \mathcal{X}$ 有限，则存在粗模空间 $\phi : \mathcal{X} \to Y$，其中 $\phi$ 分离，$Y$ 是在 $S$ 上局部有限型的代数空间。惯性有限恰好是正确条件：存在分离的粗模空间 $\phi : \mathcal{X} \to Y$，其中 $\phi$ 分离，当且仅当惯性有限。
\end{quote}
\item Conrad: \emph{The Keel-Mori Theorem via Stacks} \cite{conrad}
\begin{quote}
Keel 与 Mori 的论文 \cite{K-M} 使用 groupoid 语言，读者往往觉得难以把握。Brian Conrad 给出了堆论版本的证明；虽然使用了精细的叠语言，却相当透明。Conrad 还去除了 Noether 假设。
\end{quote}
\item Rydh: \emph{Existence of quotients by finite groups and coarse moduli
spaces} \cite{rydh_quotients}
\begin{quote}
Rydh 去除了 \cite{K-M} 与 \cite{conrad} 中要求 $\mathcal{X}$ 在某个基上有限表示的假设。
\end{quote}
\item
Abramovich, Olsson, Vistoli: \emph{Tame stacks in positive characteristic}
\cite{tame}
\begin{quote}
作者将 \emph{tame Artin stack} 定义为具有有限惯性的 Artin 叠，并要求当
$\phi : \mathcal{X} \to Y$ 是粗模空间时，$\phi_*$ 在拟凝聚层上正合。他们证明，对有限惯性的 Artin 叠，下列条件等价：$\mathcal{X}$ tame；$\mathcal{X}$ 的稳定子群线性约化；以及在粗模空间上 \'etale 局部地，$\mathcal{X}$ 是仿射概形被线性约化群概形作用所得的商。对 tame Artin 叠，粗模空间具有特别良好的性质：例如粗模空间与任意基变换交换，而一般有限惯性 Artin 叠的粗模空间只与平坦基变换交换。
\end{quote}
\item Alper: \emph{Good moduli spaces for Artin stacks} \cite{alper_good}
\begin{quote}
对于具有无限仿射稳定子群（必然不分离）的一般 Artin 叠，粗模空间常常不存在。最简单的例子是
$[\mathbf{A}^1 / \mathbf{G}_m]$。本文将拟紧态射 $\phi : \mathcal{X} \to Y$ 称为
\emph{good moduli space}，若 $\mathcal{O}_Y \to \phi_* \mathcal{O}_\mathcal{X}$ 为同构，
且 $\phi_*$ 在拟凝聚层上正合。
这一概念既推广了 \cite{tame} 中的 tame Artin 叠，也概括了 Mumford 的几何不变理论：若 $G$ 是作用在线性空间 $X \subset \mathbf{P}^n$ 上的约化群，则半稳定轨迹的商叠到 GIT 商的态射 $[X^{ss}/G] \to X//G$ 是 good moduli space。该概念具有许多良好几何性质：
(1) $\phi$ 满射、普遍闭且普遍下沉；(2) $\phi$ 按闭包等价关系识别
$Y$ 与 $\mathcal{X}$ 中的点；(3) $\phi$ 对映到代数空间具有泛性质；
(4) good moduli space 在任意基变换下稳定；(5) Artin 叠上的向量丛下降到 good moduli space，当且仅当闭点处表示是平凡的。
\end{quote}
\end{itemize}


\subsection{交理论}
\label{guide-subsection-intersection-theory}

\noindent
讨论代数叠上交理论的论文。
\begin{itemize}
\item
Vistoli: \emph{Intersection theory on algebraic stacks and on their moduli
spaces} \cite{vistoli_intersection}
\begin{quote}
本文为 Deligne--Mumford 叠建立带有有理系数的交理论基础。若
$\mathcal{X}$ 是分离 Deligne--Mumford 叠，则带有有理系数的 Chow 群
$\CH_*(\mathcal{X})$ 定义为维数为 $k$ 的整闭子叠所生成的自由阿贝尔群模去有理等价。它具有平坦拉回、适当推前以及正则局部嵌入的广义 Gysin 同态。若 $\phi : \mathcal{X} \to Y$ 是模空间（即在几何点上为双射的适当态射），则存在诱导推前 $\CH_*(\mathcal{X}) \to \CH_k(Y)$，且为同构。
\end{quote}
\item Edidin, Graham: \emph{Equivariant Intersection Theory}
\cite{edidin-graham}
\begin{quote}
本文旨在为代数群 $G$ 作用于代数空间 $X$ 的商叠 $[X/G]$ 建立带整数系数的交理论，即建立 $X$ 的 $G$-等变交理论。仅用不变循环定义的等变 Chow 群不具备良好性质。作者推广 Totaro 对 $BG$ 的定义，并利用向量丛 $V \to X$ 给出自然同构 $\CH_i(X) \cong \CH_i(V)$，定义如下。记住定义中的 $\CH_i^G(X)$：
令 $\dim(X) = n$、$\dim(G) = g$。对每个 $i$，取一个 $l$ 维 $G$-表示 $V$，使 $G$ 在开集 $U \subset V$ 上自由作用，且补集的余维 $d > n - i$。于是 $X_G = [X \times U / G]$ 是代数空间（甚至可取为概形）。定义
$\CH_i^G(X) = \CH_{i + l - g}(X_G)$。对商叠定义
$\CH_i( [X/G]) = \CH_{i + g}^G(X) = \CH_{i + l}(X_G)$。
特别地，若 $i > \dim [X/G] = n - g$，则 $\CH_i([X/G]) = 0$，但当 $i \ll 0$ 时可能非零。例如
$\CH_i(B \mathbf{G}_m) = \mathbf{Z}$（$i \le 0$）。作者证明这些等变 Chow 群具有与普通 Chow 群相同的函子性；此外还证明若
$[X / G] \cong [Y / H]$，则 $\CH_i([X/G]) = \CH_i([Y/H])$，所以定义与叠如何表示为商叠无关。
\end{quote}
\item
Kresch: \emph{Cycle Groups for Artin Stacks} \cite{kresch_cycle}
\begin{quote}
Kresch 为任意 Artin 叠定义 Chow 群；在商叠情形，这一定义与 \cite{edidin-graham} 中 Edidin 和 Graham 的定义一致。对于具有仿射稳定子群的代数叠，该理论满足通常性质。
\end{quote}
\item Behrend and Fantechi: \emph{The intrinsic normal cone}
\cite{behrend-fantechi}
\begin{quote}
Behrend 与 Fantechi 将 Li 和 Tian 的构造推广到 Deligne--Mumford 叠，构造了虚基本类。
\end{quote}
\end{itemize}


\subsection{商叠}
\label{guide-subsection-quotient-stacks}

\noindent
商叠\footnote{在文献中，\emph{quotient stack} 常指形如 $[X/G]$ 的叠，其中 $X$ 为代数空间而 $G$ 为 $\text{GL}_n$ 的子群概形，而非任意平坦群概形。}
是 Artin 叠中非常重要的子类，涵盖代数几何学家研究的几乎所有模叠。商叠
$[X/G]$ 的几何就是 $X$ 的 $G$-等变几何。证明某性质对商叠成立通常更容易，
有些结果也只对商叠已知。下列论文讨论：代数叠何时是全局商叠？代数叠是否
在局部上是商叠？
\begin{itemize}
\item Laumon, Moret-Bailly: \cite[Chapter 6]{LM-B}
\begin{quote}
第 6 章包含关于代数叠局部与整体结构的若干事实。证明了：$S$ 上的代数叠
$\mathcal{X}$ 是商叠 $[Y/G]$（其中 $Y$ 可取代数空间、概形或仿射概形）
且 $G$ 为有限群，当且仅当存在代数空间、概形或仿射概形 $Y'$ 以及有限
\'etale 态射 $Y' \to \mathcal{X}$。还证明任意 $S$ 上的 Deligne--Mumford
叠及任意 $x : \Spec(K) \to \mathcal{X}$ 都有可表示、\'etale 且分离的态射
$\phi : [X/G] \to \mathcal{X}$，其中 $G$ 是作用于 $S$ 上仿射概形的有限群，
且满足 $\Spec(K) = [X/G] \times_\mathcal{X} \Spec(K)$。还详细讨论了具有几何连通纤维的展示的存在性。
\end{quote}
\item
Edidin, Hassett, Kresch, Vistoli:
\emph{Brauer Groups and Quotient stacks} \cite{ehkv}
\begin{quote}
首先，他们建立了关于代数叠何时为商叠的一些基础（虽不困难）事实（本文中
代数叠总假定在 Noether 概形上有限型）。对代数叠 $\mathcal{X}$，以下条件等价：
$\mathcal{X}$ 是商叠；存在向量丛 $V \to \mathcal{X}$，使每个几何点的稳定子群忠实作用于纤维；存在向量丛 $V \to \mathcal{X}$ 及局部闭子叠
$V^0 \subset V$，使 $V^0$ 可表示且满射到 $\mathcal{X}$。他们证明若存在来自代数空间的有限平坦覆盖，则代数叠是商叠；任意具有一般点平凡稳定子的光滑 Deligne--Mumford 叠都是商叠。
他们证明，Noether 概形 $X$ 上对应 $\beta \in H^2(X, \mathbf{G}_m)$ 的
$\mathbf{G}_m$-gerbe 是商叠，当且仅当 $\beta$ 在 Brauer 映射
$\text{Br}(X) \to \text{Br}'(X)$ 的像中。由此构造了一个不是商叠的非分离
Deligne--Mumford 叠。
\end{quote}
\item Totaro: \emph{The resolution property for schemes and stacks}
\cite{totaro_resolution}
\begin{quote}
若每个拟凝聚层都是某个向量丛的商，则叠具有 resolution property。第一主定理：
若 $\mathcal{X}$ 是闭点处稳定子群仿射的正规 Noether 代数叠，则下列等价：
(1) $\mathcal{X}$ 具有 resolution property；(2) $\mathcal{X} = [Y/\text{GL}_n]$，其中 $Y$ 拟仿射。
若 $\mathcal{X}$ 在域上有限型，则还等价于：(3)
$\mathcal{X} = [\Spec(A)/G]$，其中 $G$ 是在 $k$ 上有限型的仿射群概形。
商叠具有 resolution property 的蕴含由 Thomason 证明。第二主定理：
若 $\mathcal{X}$ 是域上的光滑 Deligne--Mumford 叠，惯性群
$I_\mathcal{X} \to \mathcal{X}$ 有限且一般点平凡，粗模空间是具有仿射对角的概形，则 $\mathcal{X}$ 具有 resolution property。另一个有趣结果是：若 $\mathcal{X}$ 是具有 resolution property 的 Noether 代数叠，则其对角仿射当且仅当闭点的稳定子群仿射。
\end{quote}
\item Kresch: \emph{On the Geometry of Deligne-Mumford Stacks}
\cite{kresch_geometry}
\begin{quote}
本文总结域上有限型 Deligne--Mumford 叠的一般结构结果，并包含关于商叠的有趣结果。证明任意光滑、分离、一般点 tame 且粗模空间拟射影的 Deligne--Mumford 叠都是商叠 $[Y/G]$，其中 $Y$ 拟射影、$G$ 为代数群。若 $\mathcal{X}$ 为粗模空间是概形的 Deligne--Mumford 叠，则 $\mathcal{X}$ 在 Zariski 局部为商叠，当且仅当它具有由概形被有限群商得到的叠的 Zariski 开覆盖。
若 $\mathcal{X}$ 是特征 0 域上 proper 的 Deligne--Mumford 叠、粗模空间为 $Y$，则：
$Y$ 射影且 $\mathcal{X}$ 为商叠，当且仅当 $Y$ 射影且 $\mathcal{X}$ 具有生成层，当且仅当 $\mathcal{X}$ 可闭嵌入一个具有射影粗模空间的光滑 proper DM 叠。由此启发定义：若 Deligne--Mumford 叠可闭嵌入具有射影粗模空间的光滑 proper Deligne--Mumford 叠，则称其为 \emph{projective}。
\end{quote}
\item Kresch, Vistoli \emph{On coverings of Deligne-Mumford stacks and
surjectivity of the Brauer map} \cite{kresch-vistoli}
\begin{quote}
在特征 0 且固定 $n$ 时，以下两条等价：(1) 每个 $n$ 维光滑 Deligne--Mumford
叠都是商叠；(2) $n$ 维光滑概形的 Azumaya Brauer 群等于上同调 Brauer 群。
\end{quote}
\item Kresch: \emph{Cycle Groups for Artin Stacks} \cite{kresch_cycle}
\begin{quote}
证明域上有限型且约化、稳定子群仿射的 Artin 叠具有按商叠给出的分层。
\end{quote}
\item Abramovich-Vistoli:
\emph{Compactifying the space of stable maps} \cite{abramovich-vistoli}
\begin{quote}
引理 2.2.3 证明：任意分离 Deligne--Mumford 叠在粗模空间上 \'etale 局部为
商叠 $[U/G]$，其中 $U$ 仿射、$G$ 为有限群。\cite[Theorem 2.12]{olsson_homstacks} 说明在此论证中 $G$ 甚至就是稳定子群。
\end{quote}
\item Abramovich, Olsson, Vistoli:
\emph{Tame stacks in positive characteristic} \cite{tame}
\begin{quote}
本文证明 tame Artin 叠在粗模空间上 \'etale 局部为仿射概形被稳定子群作用所得的商叠。
\end{quote}
\item Alper: \emph{On the local quotient structure of Artin stacks}
\cite{alper_quotient}
\begin{quote}
猜想：对 Artin 叠 $\mathcal{X}$ 与闭点 $x \in \mathcal{X}$，若稳定子群线性约化，
则存在带 $V$ 为代数空间的 \'etale 态射 $[V/G_x] \to \mathcal{X}$。文中给出
这一猜想的若干证据。基于 \cite{tame} 思想的简单形变理论论证表明，该结论在
形式局部成立。文章给出 Luna \'etale slice 定理的叠论证明：若
$\mathcal{X} = [\Spec(A)/G]$ 且 $G$ 线性约化，则在 GIT 商
$\Spec(A^G)$ 上 \'etale 局部地，$\mathcal{X}$ 是由稳定子群作商得到的商叠。
\end{quote}
\end{itemize}

\subsection{上同调}
\label{guide-subsection-cohomology}

\noindent
讨论代数叠上层的上同调的论文。
\begin{itemize}
\item Olsson: \emph{Sheaves on Artin stacks} \cite{olsson_sheaves}
\begin{quote}
本文发展拟凝聚层与可构造层理论，并证明基本上同调性质。它纠正了
\cite{LM-B} 关于 lisse-\'etale site 函子性的错误，并构造余切复形。此外证明：
proper 态射的 Grothendieck 基本定理、Grothendieck 存在定理、Zariski 连通性定理，
以及 proper 推前的凝聚层与可构造层有限性定理。
\end{quote}
\item Behrend: \emph{Derived $l$-adic categories for algebraic stacks}
\cite{behrend_derived}
\begin{quote}
证明代数叠的 Lefschetz 迹公式。
\end{quote}
\item Behrend: \emph{Cohomology of stacks} \cite{behrend_cohomology}
\begin{quote}
定义微分叠的 de Rham 上同调与拓扑叠的奇异上同调。
\end{quote}
\item Faltings:
\emph{Finiteness of coherent cohomology for proper fppf stacks}
\cite{faltings_finiteness}
\begin{quote}
证明 proper 态射下凝聚层直接像的凝聚性。
\end{quote}
\item Abramovich, Corti, Vistoli:
\emph{Twisted bundles and admissible covers} \cite{acv}
\begin{quote}
附录包含 tame Deligne--Mumford 叠的 \'etale 上同调 proper 基变换定理。
\end{quote}
\end{itemize}


\subsection{由概形给出的有限覆盖的存在性}
\label{guide-subsection-finite-covers}

\noindent
Deligne--Mumford 叠具有由概形给出的有限覆盖，是一项重要结果。在
Deligne--Mumford 叠的交理论中，它是为不可表示态射定义 proper 推前的关键。
关于 $\overline{\mathcal{M}}_g$ 的若干结果依赖于存在由\emph{光滑}概形给出的有限覆盖，
Looijenga 已证明了这一点。这个方向上较早的结果或许是
\cite[Theorem 6.1]{seshadri_quotients}，它处理了等变情形。
\begin{itemize}
\item Vistoli: \emph{Intersection theory on algebraic stacks and on their
moduli spaces} \cite{vistoli_intersection}
\begin{quote}
若 $\mathcal{X}$ 是具有模空间的 Deligne--Mumford 叠（即在几何点上为双射的
proper 态射），则存在来自概形 $X$ 的有限态射 $X \to \mathcal{X}$。
\end{quote}
\item Laumon, Moret-Bailly: \cite[Chapter 16]{LM-B}
\begin{quote}
Zariski 主定理的应用是：定理 16.6 证明，若 $\mathcal{X}$ 是 Noether 概形上
有限型的 Deligne--Mumford 叠，则存在有限、满射且一般为 \'etale 的态射
$Z \to \mathcal{X}$，其中 $Z$ 为概形。推论 16.6.2 还证明，任意 Noether
正规代数空间同构于有限群 $G$ 作用在正规概形 $X$ 上所得的代数空间商 $X'/G$。
\end{quote}
\item Edidin, Hassett, Kresch, Vistoli:
\emph{Brauer Groups and Quotient stacks} \cite{ehkv}
\begin{quote}
定理 2.7 表明：若 $\mathcal{X}$ 是 Noether 基概形 $S$ 上有限型的代数叠，
则对角 $\mathcal{X} \to \mathcal{X} \times_S \mathcal{X}$ 拟有限，当且仅当
存在来自概形 $X$ 的有限满射 $X \to F$。
\end{quote}
\item Kresch, Vistoli: \emph{On coverings of Deligne-Mumford stacks and
surjectivity of the Brauer map} \cite{kresch-vistoli}
\begin{quote}
证明：域上有限型、光滑且分离、粗模空间拟射影的 Deligne--Mumford 叠具有来自
光滑拟射影概形的有限平坦覆盖。
\end{quote}
\item Olsson: \emph{On proper coverings of Artin stacks} \cite{olsson_proper}
\begin{quote}
证明若 $\mathcal{X}$ 是 $S$ 上分离且有限型的 Artin 叠，则存在来自在 $S$ 上
拟射影的概形 $X$ 的 proper 满射 $X \to \mathcal{X}$。应用包括 proper 态射下
直接像层的凝聚性与可构造性，以及 proper Artin 叠的 Grothendieck 存在定理。
\end{quote}
\item Rydh: \emph{Noetherian approximation of algebraic spaces and stacks}
\cite{rydh_approx}
\begin{quote}
本文定理 B 如下：设 $X$ 是拟紧代数叠，具有拟有限分离对角（resp.\ 拟紧
Deligne--Mumford 叠具有拟紧分离对角）。则存在概形 $Z$ 及有限、有限表示且
满射的态射 $Z \to X$，它在稠密拟紧开子叠 $U \subset X$ 上是平坦的（resp.\ \'etale 的）。
\end{quote}
\end{itemize}


\subsection{刚化}
\label{guide-subsection-rigidification}

\noindent
刚化是从惯性中去除一个平坦子群的过程。例如，若 $X$ 是射影簇，
从 Picard 叠到 Picard 概形的态射就是对自同构群
$\mathbf{G}_m$ 的刚化。
\begin{itemize}
\item Abramovich, Corti, Vistoli:
\emph{Twisted bundles and admissible covers} \cite{acv}
\begin{quote}
设 $\mathcal{X}$ 是 $S$ 上的代数叠，$H$ 是 $S$ 上平坦、有限表示且分离的群概形。
假设对每个 $\xi \in \mathcal{X}(T)$ 都有嵌入
$H(T) \hookrightarrow \text{Aut}_{\mathcal{X}(T)}(\xi)$，并且在拉回下相容：
对每个覆盖 $f: T \rightarrow T'$ 上的箭头 $\phi : \xi \rightarrow \xi'$ 与
$g \in H(T')$，都有 $g \circ \phi = \phi \circ f^*g$。
则存在代数叠 $\mathcal{X}/H$ 及态射
$\rho : \mathcal{X} \rightarrow \mathcal{X}/H$，它是 fppf gerbe，并且对每个
$\xi \in \mathcal{X}(T)$，态射
$\text{Aut}_{\mathcal{X}(T)} (\xi) \rightarrow
\text{Aut}_{\mathcal{X}/H (T)} (\xi)$
为满射，核为 $H(T)$。
\end{quote}
\item Romagny: \emph{Group actions on stacks and applications}
\cite{romagny_actions}
\begin{quote}
讨论群作用相对于刚化的行为。
\end{quote}
\item Abramovich, Graber, Vistoli:
\emph{Gromov-Witten theory for Deligne-Mumford stacks} \cite{agv}
\begin{quote}
附录按照 \cite{acv} 总结刚化，并给出两种替代解释。本文还包含沿闭子叠粘合代数叠以及对线丛取根的构造。
\end{quote}
\item
Abramovich, Olsson, Vistoli: \emph{Tame stacks in positive characteristic}
(\cite{tame})
\begin{quote}
附录处理更复杂的情形：惯性中的平坦子叠
$H \subset I_\mathcal{X}$ 是正规但不一定中心的。
\end{quote}
\end{itemize}

\subsection{叠曲线}
\label{guide-subsection-stacky-curves}

\noindent
讨论叠曲线的论文。
\begin{itemize}
\item Abramovich, Vistoli: \emph{Compactifying the space of stable maps}
\cite{abramovich-vistoli}
\begin{quote}
本文引入 \emph{twisted curves}。从稳定曲线到代数叠的稳定映射的模空间通常并不紧。
作者利用扭曲曲线构造模叠；当目标是 tame Deligne--Mumford 叠且粗模空间射影时，
该模叠是 proper 的。
\end{quote}
\item Behrend, Noohi: \emph{Uniformization of Deligne-Mumford curves}
\cite{behrend-noohi}
\begin{quote}
证明 Deligne--Mumford 解析曲线的统一化定理。
\end{quote}
\end{itemize}

\subsection{Hilbert、Quot、Hom 与分支簇叠}
\label{guide-subsection-hilbert-quot-hom}

\noindent
讨论 Hilbert 概形及类似对象的论文。
\begin{itemize}
\item Vistoli: \emph{The Hilbert stack and the theory of moduli of families}
\cite{vistoli_hilbert}
\begin{quote}
若 $\mathcal{X}$ 是在局部 Noether 且局部分离代数空间 $S$ 上分离、局部有限型的代数叠，
Vistoli 定义 Hilbert 叠 $\mathcal{H}\text{ilb}(\mathcal{F} / S)$，参数化来自 proper 概形的有限且非分歧态射。论文未证明地声称 $\mathcal{H}\text{ilb}(\mathcal{F} / S)$ 是代数叠。由此证明：若 $\mathcal{X}$ 如上，则当 $T$ 在 $S$ 上 proper 且平坦时，Hom 叠 $\mathcal{H} \text{om}_S(T, \mathcal{X})$ 是代数叠。
\end{quote}
\item Olsson, Starr: \emph{Quot functors for Deligne-Mumford stacks}
\cite{olsson-starr}
\begin{quote}
若 $\mathcal{X}$ 是代数空间 $S$ 上分离且局部有限表示的 Deligne--Mumford 叠，
$\mathcal{F}$ 是局部有限表示的 $\mathcal{O}_\mathcal{X}$-模，则 Quot 函子
$\text{Quot}(\mathcal{F} / \mathcal{X} / S)$ 由一个在 $S$ 上分离且局部有限表示的代数空间表示。
本文还定义生成层，并证明对由有限群作用于概形得到的全局商 tame、分离
Deligne--Mumford 叠存在生成层。
\end{quote}
\item Olsson: \emph{Hom-stacks and Restrictions of Scalars}
\cite{olsson_homstacks}
\begin{quote}
设 $\mathcal{X}$ 与 $\mathcal{Y}$ 是代数空间 $S$ 上局部有限表示、对角有限的 Artin 叠，
其中 $\mathcal{X}$ 在 $S$ 上 proper 且平坦，并且在 $S$ 上 fppf 局部地，
$\mathcal{X}$ 有来自代数空间的有限有限表示平坦覆盖（例如 $\mathcal{X}$ 是
Deligne--Mumford 叠或 tame Artin 叠）。则 $\Hom_S(\mathcal{X}, \mathcal{Y})$
是 $S$ 上局部有限表示的 Artin 叠。
\end{quote}
\item Alexeev and Knutson: \emph{Complete moduli spaces of branchvarieties}
(\cite{alexeev-knutson})
\begin{quote}
他们将 $\mathbf{P}^n$ 的 branchvariety 定义为来自\emph{约化}概形 $X$ 的有限态射
$X \rightarrow \mathbf{P}^n$。证明固定 Hilbert 多项式及 $i$ 维分支的总次数后，
branchvariety 的模叠是具有有限稳定子的 proper Artin 叠。并比较 branchvariety 叠
与 Hilbert 概形、Chow 概形和稳定映射模空间。
\end{quote}
\item Lieblich: \emph{Remarks on the stack of coherent algebras}
\cite{lieblich_remarks}
\begin{quote}
本文通过把叠构造成概形 $Y$ 的结构层上的代数叠，构造 Alexeev--Knutson
branch-variety 叠在概形 $Y$ 上的推广，并给出 $\text{Quot}$ 与 $\Hom$ 空间的存在性证明。
\end{quote}
\item Starr: \emph{Artin's axioms, composition, and moduli spaces}
\cite{starr_artin}
\begin{quote}
作为主结果的应用，给出 Vistoli Hilbert 叠 \cite{vistoli_hilbert} 与 Alexeev--Knutson
branchvariety 叠 \cite{alexeev-knutson} 的共同推广。若 $\mathcal{X}$ 是优良概形
$S$ 上局部有限型、对角有限的代数叠，则参数化从带有 $G$-ample 线丛 $L$ 的
proper 代数空间 $T$ 到 $\mathcal{X}$ 的态射 $g: T \rightarrow \mathcal{X}$ 的叠
$\mathcal{H}$ 是 $S$ 上局部有限型 Artin 叠。
\end{quote}
\item Lundkvist and Skjelnes:
\emph{Non-effective deformations of Grothendieck's Hilbert functor}
\cite{lundkvist-skjelnes}
\begin{quote}
证明非分离概形的 Hilbert 函子不可表示，因为存在非有效形变。
\end{quote}
\item Halpern-Leistner and Preygel:
\emph{Mapping stacks and categorical notions of properness}
\cite{HL-P}
\begin{quote}
本文证明，在比 Olsson 论文更一般或至少不同的源与目标假设下，Hom 叠仍是代数的。
\end{quote}
\end{itemize}



\subsection{环面叠}
\label{guide-subsection-toric}

\noindent
环面叠提供了大量例子，也是检验猜想的自然场所：环面叠的几何与其叠扇的组合学
之间存在字典，正如环面簇在概形理论中提供例子与反例一样。
\begin{itemize}
\item Borisov, Chen and Smith: \emph{The orbifold Chow ring of toric
Deligne-Mumford stacks} \cite{bcs}
\begin{quote}
受 Cox 对环面簇构造的启发，本文将光滑环面 DM 叠定义为与称为
\emph{stacky fan} 的组合对象相联系的显式商叠。
\end{quote}
\item Iwanari: \emph{The category of toric stacks} \cite{iwanari_toric}
\begin{quote}
本文将 \emph{toric triple} 定义为光滑 Deligne--Mumford 叠
$\mathcal{X}$，带有稠密像的开浸入 $\mathbf{G}_m \hookrightarrow \mathcal{X}$
（因此 $\mathcal{X}$ 是 orbifold）以及作用
$\mathcal{X} \times \mathbf{G}_m \rightarrow \mathcal{X}$。证明 toric triple
的 2-范畴与 stacky fan 的 1-范畴等价，并讨论 toric triple 与
\cite{bcs} 中光滑环面 DM 叠定义的关系。
\end{quote}
\item Iwanari: \emph{Integral Chow rings for toric stacks}
\cite{iwanari_chow}
\begin{quote}
将 Cox 对环面簇的 $\Delta$-collections 推广到环面 orbifold。
\end{quote}
\item Perroni: \emph{A note on toric Deligne-Mumford stacks}
\cite{perroni}
\begin{quote}
将 Cox 的 $\Delta$-collections 及 Iwanari 的论文
\cite{iwanari_chow} 推广到一般光滑环面 DM 叠。
\end{quote}
\item Fantechi, Mann, and Nironi: \emph{Smooth toric DM stacks}
\cite{fmn}
\begin{quote}
本文将光滑环面 DM 叠定义为带有 DM 环面 $\mathcal{T}$ 作用的光滑 DM 叠
$\mathcal{X}$（即 Picard 叠同构于 $T \times BG$，其中 $G$ 有限），且具有同构于
$\mathcal{T}$ 的开稠密轨道。作者给出“自底向上”的描述，并证明光滑环面 DM 叠与
stacky fan 等价。
\end{quote}
\item Geraschenko and Satriano: \emph{Toric Stacks I and II}
\cite{gs_toric1} and \cite{gs_toric2}
\begin{quote}
这两篇论文将环面叠定义为环面簇被其环面的某个子群作商得到的叠。generically
stacky toric stack 定义为环面叠中的环面不变子叠；该定义包含并推广了早期定义。
第一篇建立叠扇组合学与相应叠几何之间的字典，并给出光滑环面叠的模解释，
推广 \cite{perroni} 中的解释。第二篇证明环面叠的内在刻画。
\end{quote}
\end{itemize}



\subsection{形式函数定理与 Grothendieck 存在定理}
\label{guide-subsection-theorem-formal-functions}

\noindent
这些论文推广形式函数定理
\cite[III.4.1.5]{EGA}（有时称为 proper 态射的 Grothendieck 基本定理）
以及 Grothendieck 存在定理 \cite[III.5.1.4]{EGA}。
\begin{itemize}
\item Knutson: \emph{Algebraic spaces} \cite[Chapter V]{Kn}
\begin{quote}
将这些定理推广到代数空间。
\end{quote}
\item Abramovich-Vistoli: \emph{Compactifying the space of stable maps}
\cite[A.1.1]{abramovich-vistoli}
\begin{quote}
将这些定理推广到 tame Deligne--Mumford 叠。
\end{quote}
\item Olsson and Starr: \emph{Quot functors for Deligne-Mumford stacks}
\cite{olsson-starr}
\begin{quote}
将这些定理推广到分离 Deligne--Mumford 叠。
\end{quote}
\item Olsson: \emph{On proper coverings of Artin stacks}
\cite{olsson_proper}
\begin{quote}
给出到 proper Artin 叠的推广。
\end{quote}
\item Conrad: \emph{Formal GAGA on Artin stacks} \cite{conrad_gaga}
\begin{quote}
给出到 proper Artin 叠的推广，并证明形式 GAGA 定理。
\end{quote}
\item Olsson: \emph{Sheaves on Artin stacks} \cite{olsson_sheaves}
\begin{quote}
给出到 proper Artin 叠的推广的另一种证明。
\end{quote}
\end{itemize}


\subsection{叠上的群作用}
\label{guide-subsection-group-actions}

\noindent
代数叠上的群作用自然出现。例如，对称群 $S_n$ 作用于
$\overline{\mathcal{M}}_{g, n}$；若群 $G$ 作用于概形 $X$，则
$\text{Aut}(X)$ 中 $G$ 的正规化子作用于 $[X/G]$。此外，环面作用经常出现在
Gromov--Witten 理论中。
\begin{itemize}
\item Romagny: \emph{Group actions on stacks and applications}
\cite{romagny_actions}
\begin{quote}
本文精确说明群如何作用于代数叠，并证明固定点的存在性以及群概形作用在代数叠上时商的存在性。另见 Romagny 较早的笔记 \cite{romagny_notes}。
\end{quote}
\end{itemize}

\subsection{对线丛取根}
\label{guide-subsection-root-stacks}

\noindent
这一有用构造由 Cadman 以及 Abramovich、Graber、Vistoli 独立发现。给定带有效 Cartier 除子
$D$ 的概形 $X$，$r$ 次根叠是沿 $D$ 在 $X$ 上分歧的 Artin 叠，在 $D$ 上有
$\mu_r$ 稳定子群，而远离 $D$ 时类似概形。
\begin{itemize}
\item Charles Cadman
\emph{Using Stacks to Impose Tangency Conditions on Curves}
\cite{cadman}
\item Abramovich, Graber, Vistoli: \emph{Gromov-Witten theory for
Deligne-Mumford stacks} \cite{agv}
\end{itemize}

\subsection{其他论文}
\label{guide-subsection-other}

\noindent
其他论文的杂集。
\begin{itemize}
\item Lieblich: \emph{Moduli of twisted sheaves} \cite{lieblich_twisted}
\begin{quote}
本文概述 gerbe 与扭曲层。若 $\mathcal{X} \rightarrow X$ 是 $\mu_n$-gerbe，
$X$ 是具有光滑连通几何纤维的射影相对曲面，则证明 $\mathcal{X}$-扭曲半稳定层的
叠是在 $S$ 上局部有限表示的 Artin 叠。本文还发展 Artin 叠上伴随点与层纯性的理论。
\end{quote}
\item Lieblich, Osserman:
\emph{Functorial reconstruction theorem for stacks}
\cite{lieblich-osserman}
\begin{quote}
证明代数叠何时可以由其伴随函子重构的一些出人意料且有趣的结果。
\end{quote}
\item David Rydh:
\emph{Noetherian approximation of algebraic spaces and stacks}
\cite{rydh_approx}
\begin{quote}
证明每个具有拟有限对角的拟紧代数叠都可由 Noether 叠近似。应用包括去除
Chevalley、Serre、Zariski 与 Chow 结果中的 Noether 假设。
\end{quote}
\end{itemize}



\section{其他领域中的叠}
\label{guide-section-stacks-areas}

\begin{itemize}
\item Behrend and Noohi: \emph{Uniformization of Deligne-Mumford curves}
\cite{behrend-noohi}
\begin{quote}
概述拓扑叠、解析叠与代数叠，并进行比较。
\end{quote}
\item Behrang Noohi: \emph{Foundations of topological stacks I} \cite{noohi}
\item David Metzler: \emph{Topological and smooth stacks} \cite{metzler}
\end{itemize}


\section{高阶叠}
\label{guide-section-higher-stacks}

\begin{itemize}
\item Lurie: \emph{Higher topos theory} \cite{lurie_topos}
\item Lurie: \emph{Derived Algebraic Geometry I - V}
\cite{dag1}, \cite{dag2}, \cite{dag3}, \cite{dag4}, \cite{dag5}
\item To\"en: \emph{Higher and derived stacks: a global overview}
\cite{toen_higher}
\item To\"en and Vezzosi: \emph{Homotopical algebraic geometry I, II}
\cite{hag1}, \cite{hag2}
\end{itemize}
% Generated mechanically from the frozen desirables.tex target.
% Do not edit this generated file; edit the bound locale source instead.

% OK, start here.
%

\setcounter{chapter}{112}
\chapter{待增内容}


\phantomsection
\label{desirables-section-phantom}



\section{引言}
\label{desirables-section-introduction}

\noindent
这里基本上只是列出我们希望加入 Stacks 项目的内容。由于我们持续向
Stacks 项目添加材料，这份清单总是会稍微落后于项目的当前状态。事实上，
试图列出应当添加的内容也许是个错误，因为这份清单似乎不可能保持最新。

\medskip\noindent
最后更新：2017 年 8 月 31 日，星期四。


\section{约定}
\label{desirables-section-conventions}

\noindent
我们应当增加一章，简要列出文档中使用的约定。
该章节已经存在，见
Conventions，第 \ref{conventions-section-comments} 节，
但那里还可以补充许多内容。尤其有用的是找出“隐藏的”约定和未明说的假设，
并将它们放在该章节中。


\section{位点与拓扑斯}
\label{desirables-section-sites}

\noindent
我们有一章讨论位点与层，见
Sites，第 \ref{sites-section-introduction} 节。
我们有一章讨论环化位点（以及拓扑斯）及其上的模，见
Modules on Sites，第 \ref{sites-modules-section-introduction} 节。
我们还有一章讨论这一背景下的上同调，见
Cohomology on Sites，第 \ref{sites-cohomology-section-introduction} 节。
不过仍可加入更多内容，尤其是上同调章节。


\section{叠}
\label{desirables-section-stacks}

\noindent
我们有一章讨论（抽象的）叠，见
Stacks，第 \ref{stacks-section-introduction} 节。
如果能够做到以下几点就很好：
\begin{enumerate}
\item 改进关于“叠化”的讨论，
\item 给出叠化的例子，
\item 总体上增加更多例子，
\item 改进关于 gerbe 的讨论。
\end{enumerate}
尚未加入的一个结果例子是：给定阿贝尔群层 $\mathcal{F}$
在范畴 $\mathcal{C}$ 上，由 $\mathcal{F}$ 带化的 gerbe 的等价类集合与
$H^2(\mathcal{C}, \mathcal{F})$ 双射。


\section{单纯方法}
\label{desirables-section-simplicial}

\noindent
我们有一章讨论单纯方法，见
Simplicial，第 \ref{simplicial-section-introduction} 节。
这一章需要审阅和改进。关于单纯同伦（也称为组合同伦）与 Kan 复形之间
关系的讨论应当加强。
我们还有一章讨论单纯空间，见
Simplicial Spaces，第 \ref{spaces-simplicial-section-introduction} 节。
该章简要讨论单纯拓扑空间、单纯位点和单纯拓扑斯。
我们可以进一步发展“单纯代数几何”，讨论单纯概形（或单纯代数空间、或
单纯代数叠），并处理几何问题、它们的上同调等。


\section{概形的上同调}
\label{desirables-section-schemes-cohomology}

\noindent
已经有一章讨论拟凝聚层的上同调，见
Cohomology of Schemes，第 \ref{coherent-section-introduction} 节。
我们有一章讨论概形上的拟凝聚层的导出范畴，见
Derived Categories of Schemes，第 \ref{perfect-section-introduction} 节。
我们有一章讨论 Noether 概形的对偶性以及概形态射的相对对偶性，见
Duality for Schemes，第 \ref{duality-section-introduction} 节。
我们还有关于概形的 \'etale 上同调以及概形的结晶上同调的章节。不过这些
章节中的大部分材料都非常基础，其中还有许多内容可以、也应当补充。


\section{Schlessinger 风格的形变理论（\`a la Schlessinger）}
\label{desirables-section-deformation-schlessinger}

\noindent
我们有一章讨论这方面的材料，见
Formal Deformation Theory，第 \ref{formal-defos-section-introduction} 节。
我们有一章讨论一般理论的例子，见
Deformation Problems，第 \ref{examples-defos-section-introduction} 节。
我们还有一章，见
Deformation Theory，第 \ref{defos-section-introduction} 节，
其中讨论环（以及模）的形变、环化空间（以及模层）的形变、以及环化拓扑斯
（以及模层）的形变。
本章使用朴素余切复形描述障碍、一阶形变和无穷小自同构。这些材料后来已
在模叠的代数性问题中得到一些应用。
另有一章讨论完整余切复形，见
Cotangent，第 \ref{cotangent-section-introduction} 节。


\section{代数叠的定义}
\label{desirables-section-definition-algebraic-stacks}

\noindent
代数叠是在带有 fppf 拓扑的概形范畴上的群胚中的叠，其对角态射可由代数
空间表示，并且是某个概形到它的满射光滑态射的目标。
见 Algebraic Stacks，第 \ref{algebraic-section-algebraic-stacks} 节。
“Deligne--Mumford 叠”是满足下述条件的代数叠：存在一个概形以及从该概形
到它的满射 \'etale 态射，正如 Deligne 和 Mumford 的论文 \cite{DM} 中所述；
见 Algebraic Stacks，定义 \ref{algebraic-definition-deligne-mumford}。
我们将“Artin 叠”一词保留给 Artin 的论文中所讨论的叠，见
\cite{ArtinI}、\cite{ArtinII} 和 \cite{ArtinVersal}。
一种可能的定义是：Artin 叠是局部 Noether 概形 $S$ 上的代数叠
$\mathcal{X}$，使得 $\mathcal{X} \to S$ 局部有限型
\footnote{也就是说，这些恰好是满足 Artin 公理 [-1]、[0]、[1]、[2]、[3]、[4]、[5]
的 $S$ 上代数叠，见 Artin's Axioms，第 \ref{artin-section-axioms} 节。}。


\section{概形、代数空间与代数叠的例子}
\label{desirables-section-examples-stacks}

\noindent
Stacks 项目目前有两章讨论模叠及其性质，见
Moduli Stacks，第 \ref{moduli-section-introduction} 节，以及
Moduli of Curves，第 \ref{moduli-curves-section-introduction} 节。
随着时间推移，我们打算加入更多例子，例如：
\begin{enumerate}
\item $\mathcal{A}_g$，即亏格为 $g$ 的主极化阿贝尔概形，
\item $\mathcal{A}_1 = \mathcal{M}_{1, 1}$，即带 $1$ 个标记点的光滑射影亏格 $1$ 曲线，
\item $\mathcal{M}_{g, n}$，即带有 $n$ 个两两不同且带标号点的光滑射影亏格 $g$ 曲线，
\item $\overline{\mathcal{M}}_{g, n}$，即带 $n$ 个标记点的稳定结点射影亏格 $g$ 曲线，
\item $\SheafHom_S(\mathcal{X}, \mathcal{Y})$，态射的模空间
（对叠 $\mathcal{X}$、$\mathcal{Y}$ 以及基概形 $S$ 施加适当条件），
\item $\textit{Bun}_G(X) = \SheafHom_S(X, BG)$，几何 Langlands 纲领中的 $G$-丛叠
（对概形 $X$、群概形 $G$ 以及基概形 $S$ 施加适当条件），
\item $\Picardstack_{\mathcal{X}/S}$，即与基概形（或基空间）上的代数叠关联的 Picard 叠。
\end{enumerate}
更一般地说，Stacks 项目在具有几何意义的例子方面仍然有所欠缺。


\section{代数叠的性质}
\label{desirables-section-stacks-properties}

\noindent
这也许是比较容易开展的项目之一，因为大部分基础理论现在已经具备。
当然，这些实际上是叠的态射的性质。我们可以定义奇点（模去光滑因子）等；
还可以证明连通正规叠是不可约的，等等。


\section{代数叠的光滑 \'etale 位点}
\label{desirables-section-lisse-etale}

\noindent
这一内容已在 Cohomology of Stacks，第 \ref{stacks-cohomology-section-lisse-etale} 节中介绍。
为了说明它对于代数叠的 $1$-态射并非函子性，Examples，第
\ref{examples-section-lisse-etale-not-functorial} 节中讨论了一个例子。
当然还可以说得更多，但事实证明，尽可能使用“大” \'etale 位点来证明结果
非常有用。


\section{你一直想知道却不敢问的事情}
\label{desirables-section-stacks-fun-lemmas}

\noindent
会有许多引理被反复使用：它们很有用，却没有在文献中得到明确提及，或者
很难找到参考文献。这里就是一个技巧袋。

\medskip\noindent
例：给定概形中的两个群胚 $R\Rightarrow U$ 和
$R' \Rightarrow U'$，仅用概形中的群胚来描述一个
$1$-态射 $[U/R] \to [U'/R']$ 到底意味着什么。


\section{叠上的拟凝聚层}
\label{desirables-section-quasi-coherent}

\noindent
这些内容在 Cohomology of Stacks 章节中定义和讨论，见
第 \ref{stacks-cohomology-section-introduction} 节。
模的导出范畴在 Derived Categories of Stacks 章节中讨论，见
第 \ref{stacks-perfect-section-introduction} 节。
这些章节还可以补充许多内容。


\section{平坦与光滑}
\label{desirables-section-flat-smooth}

\noindent
Artin 定理说明，从概形出发的平坦满射可以替代满射光滑条件。
这一结果现在已作为 Criteria for Representability，第
\ref{criteria-theorem-bootstrap} 条定理提供。


\section{Artin 表示定理}
\label{desirables-section-representability}

\noindent
这一内容在 Artin's Axioms 章节中讨论，见
第 \ref{artin-section-introduction} 节。
我们还有一个应用，见
Quot，第 \ref{quot-theorem-coherent-algebraic} 条定理。
还应当有更多应用，并且该章节本身也需要清理。


\section{DM 叠由概形有限覆盖}
\label{desirables-section-dm-finite-cover}

\noindent
对于代数空间，我们已经有相应结果，见 Limits of Spaces，第
\ref{spaces-limits-section-finite-cover} 节。
缺少的是 DM 叠和拟 DM 叠的结果。


\section{Martin Olsson 关于真性的论文}
\label{desirables-section-proper-parametrization}

\noindent
该论文证明两种真性的概念相同。第一部分现在以代数叠的 Chow 引理形式
提供，见 More on Morphisms of Stacks，第
\ref{stacks-more-morphisms-theorem-chow-finite-type} 条定理。
作为推论，在某些情形下检验代数叠的赋值判据时，只需使用 DVR，见
More on Morphisms of Stacks，第
\ref{stacks-more-morphisms-section-Noetherian-valuative-criterion} 节。


\section{相干层的真推前}
\label{desirables-section-proper-pushforward}

\noindent
现在我们已经有了代数叠的 Chow 引理，可以开始研究这一内容，见上一节。


\section{Keel 与 Mori}
\label{desirables-section-keel-mori}

\noindent
见 \cite{K-M}。他们的结果已经加入 More on Morphisms of Stacks，第
\ref{stacks-more-morphisms-section-Keel-Mori} 节。


\section{在此处添加更多内容}
\label{desirables-section-add-more}

\noindent
其实，我们根本不应当把这份清单作为 Stacks 项目本身的一部分开始！
其他地方有一份待办清单，更容易更新。
% Generated mechanically from the frozen coding.tex target.
% Do not edit this generated file; edit the bound locale source instead.

% OK, start here.
%

\setcounter{chapter}{113}
\chapter{编码风格}



\phantomsection
\label{coding-section-phantom}


\section{风格说明清单}
\label{coding-section-style}

\noindent
这些说明会随时间改变，但现在先列出一些，希望能够促进 LaTeX 风格的一致性。
我们把源文件的内容称为“代码\footnote{这全是 Knuth 的责任。见 \cite{Knuth}。}”。

\begin{enumerate}
\item 所有 tex 文件中的每一行最多 80 个字符。
\item tex 文件中不要使用缩进。使用编辑器的语法高亮，而不是缩进，来直观显示
环境等结构。
\item 使用
\begin{verbatim}
\medskip\noindent
\end{verbatim}
开始一个新段落；环境之后紧接着开始新段落时，使用
\begin{verbatim}
\noindent
\end{verbatim}
。
\item 如果可能，不要把数学公式的代码跨行拆开。如果包含美元符号的完整代码
无法放在一行中，那么就在下一行的第一个字符处开始第一个美元符号。如果这样
仍然放不下，就在数学上合理的位置拆分代码。
\item 陈列的数学公式应按如下方式编码：
\begin{verbatim}
$$
...
...
$$
\end{verbatim}
换句话说，在单独一行上以双美元符号开始，并以同样方式结束。
\item {\it 不要使用任何宏}。理由是：这样更容易阅读 tex 文件，也更容易开始编辑
任意部分，而不必先学会无数宏；并且这不会使书写变得更困难或更耗时。
当然，其缺点是同一个数学对象在文本的不同位置可能采用不同的 TeX 写法，
但这应当很容易发现。
\item 我们使用的定理环境包括：
“theorem”、“proposition”、“lemma”（plain），
“definition”、“example”、“exercise”、“situation”（definition），
“remark”、“remarks”（remark）。当然也有“proof”环境。
\item 环境 “foo” 应按如下方式编码：
\begin{verbatim}
\begin{foo}
...
...
\end{foo}
\end{verbatim}
其方式与陈列公式的编码方式类似。
\item 不使用 “corollary”，而只使用 “lemma” 环境，因为该结果很可能会被用来
证明下一个更大的定理。
\item 每个引理、命题或定理之后都直接给出该引理、命题或定理的证明。请不要
使用嵌套证明。
\item 文件 preamble.tex、chapters.tex 和 fdl.tex 是特殊的 tex 文件。除此之外，
每个 tex 文件都具有如下结构：
\begin{verbatim}
\title{Title}
...
...
\end{verbatim}
\item 尽量给引理、命题、定理，甚至注记、练习和其他环境添加标签。
如果给引理加标签，可以使用如下形式：
\begin{verbatim}
\begin{lemma}
\label{coding-lemma-bar}
...
\end{lemma}
\end{verbatim}
其他环境同理。换句话说，名为 “foo” 的环境的标签以 “foo-” 开头。
此外，请使所有标签只由小写字母、数字和符号 “-” 构成。
\item 永远不要称某个结果为“上面的引理”（或命题等）。应改为使用：
\begin{verbatim}
Lemma \ref{coding-lemma-bar} above
\end{verbatim}
这样以后移动引理基本上不会造成问题。
\item 跨文件引用。要引用文件 foo.tex 中标签为 “lemma-bar” 的引理，而该文件
标题为 “Foo”，请使用如下代码：
\begin{verbatim}
Foo, Lemma \ref{foo-lemma-bar}
\end{verbatim}
如果这不起作用，请查看文件 preamble.tex，找出应使用的正确表达式。
这样输出文件中会出现 “Foo, Lemma $<$link$>$”，从而清楚表明该链接指向文件外部。
\item 如果可能，避免在证明环境中使用前向引用。（应当可以为此编写自动化测试。）
\item 不要以数学符号开始任何句子。
\item 不要在引理、命题或定理之前紧接着写“这由下面内容推出”之类的句子。
每个句子都以句号结束。
\item 在每个引理、命题和定理中陈述全部假设。这样读者更容易判断某个引理、
命题或定理是否适用于自己的具体问题。
\item 保持证明简短；在 pdf 或 dvi 中少于 1 页。总可以通过把证明拆分到引理等
环境中来做到这一点。
\item 在定义性质 foobar 时，在 definition 环境内部的代码中使用
\begin{verbatim}
{\it foobar}
\end{verbatim}
。定义出现在文档正文中时也同样处理。这样读者就能更容易看出正在定义的对象。
\item 把将在所在章节之外使用的定义放入它自己的 definition 环境中。临时定义
可以写在正文中。数学构造是一个棘手的情形（它们往往是涉及相互关联引理的定义）。
一个可能的好办法是把它们放在自己的短章节中，使用户可以引用该章节，而不是
引用某个定义。
\item 除非公式确实会在正文中被引用，否则不要给公式编号。以后总可以添加标签。
\item 在引理、命题和定理的陈述及证明中保持句子简短。例如，与其写
“设 $R$ 是一个环，且设 $M$ 是一个 $R$-模。”，不如写
“设 $R$ 是一个环。设 $M$ 是一个 $R$-模。”理由是：这会使证明和陈述中较难
解析的部分更容易理解。
\item 使用
\begin{verbatim}
\section
\end{verbatim}
命令创建章节，但尽量避免使用小节和小小节。
\item 避免使用复杂的 latex 构造。
\end{enumerate}
% Generated mechanically from the frozen obsolete.tex target.
% Do not edit this generated file; edit the bound locale source instead.

% OK, start here.
%

\setcounter{chapter}{114}
\chapter{已废弃内容}


\phantomsection
\label{obsolete-section-phantom}





\section{引言}
\label{obsolete-section-introduction}

\noindent
本章收录一些已经变得“已废弃”的引理（见 \cite{Miller}）。


\section{预备知识}
\label{obsolete-section-preliminaries}

\begin{remark}
\label{obsolete-remark-composition-of-adjoints-isomorphic-to-identity}
本注记原先包含的信息，现在已由下列两个引理的结合涵盖：
Categories，引理 \ref{categories-lemma-adjoint-fully-faithful} 与
\ref{categories-lemma-left-adjoint-composed-fully-faithful}。
\end{remark}


\section{同调代数}
\label{obsolete-section-homological-algebra}

\begin{remark}
\label{obsolete-remark-weak-serre-subcategory}
下列注记已经废弃，因为它们已由
Homology，引理 \ref{homology-lemma-biregular-ss-converges} 与
\ref{homology-lemma-first-quadrant-ss} 涵盖。
设 $\mathcal{A}$ 是阿贝尔范畴，且
$\mathcal{C} \subset \mathcal{A}$ 是弱 Serre 子范畴（见
Homology，定义 \ref{homology-definition-serre-subcategory}）。
设 $K^{\bullet, \bullet}$ 是适用
Homology，引理 \ref{homology-lemma-first-quadrant-ss} 的双复形，并且对某个
$r \geq 0$，所有对象 ${}'E_r^{p, q}$ 都属于 $\mathcal{C}$。
则所有上同调群 $H^n(sK^\bullet)$ 都属于 $\mathcal{C}$。事实上，这些假设意味着
${}'d_r^{p, q}$ 的核与像都在 $\mathcal{C}$ 中，于是每个
${}'E_{r + 1}^{p, q}$ 也在 $\mathcal{C}$ 中。归纳可知每个
${}'E_\infty^{p, q}$ 都在 $\mathcal{C}$ 中。因此每个
$H^n(sK^\bullet)$ 都有一个有限过滤，其子商属于 $\mathcal{C}$。
利用 $\mathcal{C}$ 对扩张封闭，可得 $H^n(sK^\bullet)$ 属于
$\mathcal{C}$，如所述。
对 $K^{\bullet, \bullet}$ 关联的第二个谱序列也有同样结论。类似地，若
$(K^\bullet, F)$ 是适用
Homology，引理 \ref{homology-lemma-biregular-ss-converges} 的过滤复形，并且对某个
$r \geq 0$，所有对象 $E_r^{p, q}$ 都属于 $\mathcal{C}$，则每个
$H^n(K^\bullet)$ 都是 $\mathcal{C}$ 中的对象。
\end{remark}




\section{已废弃的代数引理}
\label{obsolete-section-algebra}

\begin{lemma}
\label{obsolete-lemma-finite-presentation-module-independent}
设 $M$ 是有限表示的 $R$-模。对任意满射
$\alpha : R^{\oplus n} \to M$，$\alpha$ 的核是有限生成的 $R$-模。
\end{lemma}

\begin{proof}
这是 Algebra，引理 \ref{algebra-lemma-extension} 的特殊情形。
\end{proof}

\begin{lemma}
\label{obsolete-lemma-p-ring-map}
设 $\varphi : R \to S$ 是环映射。若
\begin{enumerate}
\item 对任意 $x \in S$，存在 $n > 0$ 使得
$x^n$ 属于 $\varphi$ 的像；并且
\item 对任意 $x \in \Ker(\varphi)$，存在 $n > 0$
使得 $x^n = 0$；
\end{enumerate}
则 $\varphi$ 在谱空间上诱导同胚。给定素数 $p$，若
\begin{enumerate}
\item[(a)] $S$ 作为 $R$-代数由元素 $x$ 生成，并且对每个这样的 x，
存在 $n > 0$ 使 $x^{p^n} \in \varphi(R)$ 且
$p^nx \in \varphi(R)$；并且
\item[(b)] $\varphi$ 的核由幂零元素生成，
\end{enumerate}
则 (1) 与 (2) 成立；对任意环映射 $R \to R'$，
环映射 $R' \to R' \otimes_R S$ 也满足 (a)、(b)、(1) 和 (2)，
尤其在谱空间上诱导同胚。
\end{lemma}

\begin{proof}
这是 Algebra，引理 \ref{algebra-lemma-powers} 与
\ref{algebra-lemma-p-ring-map} 的组合。
\end{proof}

\noindent
下面的技术性引理说明：对于在适当意义下本质有限型的环映射，
可以把纤维中的任意关系序列提升到整个空间。

\begin{lemma}
\label{obsolete-lemma-lift-elements-ideal}
设 $R \to S$ 是环映射。
设 $\mathfrak p \subset R$ 是素理想，
$\mathfrak q \subset S$ 是位于 $\mathfrak p$ 上方的素理想。
假设 $S_{\mathfrak q}$ 在 $R_\mathfrak p$ 上本质有限型。
给定
\begin{enumerate}
\item 整数 $n \geq 0$；
\item 素理想 $\mathfrak a \subset \kappa(\mathfrak p)[x_1, \ldots, x_n]$；
\item 满的 $\kappa(\mathfrak p)$-同态
$$
\psi : (\kappa(\mathfrak p)[x_1, \ldots, x_n])_{\mathfrak a}
\longrightarrow
S_{\mathfrak q}/\mathfrak p S_{\mathfrak q},
$$
以及
\item $\Ker(\psi)$ 中的元素 $\overline{f}_1, \ldots, \overline{f}_e$。
\end{enumerate}
则存在
\begin{enumerate}
\item 整数 $m \geq 0$；
\item 元素 $g \in S$，$g \not\in \mathfrak q$；
\item 映射
$$
\Psi :
R[x_1, \ldots, x_n, x_{n + 1}, \ldots, x_{n + m}]
\longrightarrow
S_g,
$$
以及
\item $\Ker(\Psi)$ 中的元素
$f_1, \ldots, f_e, f_{e + 1}, \ldots, f_{e + m}$
\end{enumerate}
使得
\begin{enumerate}
\item 下图交换：
$$
\xymatrix{
R[x_1, \ldots, x_{n + m}] \ar[d]_\Psi
\ar[rr]_-{x_{n + j} \mapsto 0} & &
(\kappa(\mathfrak p)[x_1, \ldots, x_n])_{\mathfrak a} \ar[d]^\psi \\
S_g \ar[rr] & &
S_{\mathfrak q}/\mathfrak p S_{\mathfrak q}
},
$$
\item 元素 $f_i$（$i \leq n$）在局部环中映到
$\overline{f}_i$ 的某个单位倍；
$$
(\kappa(\mathfrak p)[x_1, \ldots, x_{n + m}])_{
(\mathfrak a, x_{n + 1}, \ldots, x_{n + m})},
$$
\item 元素 $f_{e + j}$ 在同一局部环中映到
$x_{n + j}$ 的某个单位倍；并且
\item 诱导映射 $R[x_1, \ldots, x_{n + m}]_{\mathfrak b}
\to S_{\mathfrak q}$ 是满射，其中
$\mathfrak b = \Psi^{-1}(\mathfrak qS_g)$。
\end{enumerate}
\end{lemma}

\begin{proof}
我们断言，只需在 $R$ 与 $S$ 都是局部环，
且极大理想分别为 $\mathfrak p$ 与 $\mathfrak q$ 的情形证明引理。
事实上，假设已经构造出
$$
\Psi' : R_{\mathfrak p}[x_1, \ldots, x_{n + m}]
\longrightarrow
S_{\mathfrak q}
$$
以及 $f_1', \ldots, f_{e + m}' \in R_{\mathfrak p}[x_1, \ldots, x_{n + m}]$，
并满足所有要求。则存在元素
$f \in R$，$f \not \in \mathfrak p$，使每个
$ff_k'$ 都来自某个 $f_k \in R[x_1, \ldots, x_{n + m}]$。
此外，对某个适当的 $g \in S$，$g \not \in \mathfrak q$，
元素 $\Psi'(x_i)$ 是某些元素 $y_i \in S_g$ 的像。
令 $\Psi$ 为由规则 $\Psi(x_i) = y_i$ 定义的 $R$-代数映射。
由于 $\Psi(f_i)$ 在局部化 $S_{\mathfrak q}$ 中为零，
必要时替换 $g$ 后可假设 $\Psi(f_i) = 0$。断言得证。

\medskip\noindent
因此可假设 $R$ 与 $S$ 都是局部环，
其极大理想分别为 $\mathfrak p$ 与 $\mathfrak q$。
取 $y_1, \ldots, y_n \in S$，使得
$y_i \bmod \mathfrak pS = \psi(x_i)$。
取元素 $y_{n + 1}, \ldots, y_{n + m} \in S$，使它们生成某个
$R$-子代数，而 $S$ 是该子代数的局部化。
这样的元素存在，因为假设 $S$ 在 $R$ 上本质有限型。
由于 $\psi$ 是满射，可对某个
$h_j \in \kappa(\mathfrak p)[x_1, \ldots, x_n]_{\mathfrak a}$ 写成
$y_{n + j} \bmod \mathfrak pS = \psi(h_j)$。
写 $h_j = g_j/d$，其中 $g_j \in \kappa(\mathfrak p)[x_1, \ldots, x_n]$，
公共分母 $d \in \kappa(\mathfrak p)[x_1, \ldots, x_n]$ 满足
$d \not \in \mathfrak a$。取 $g_j$ 与 $d$ 的提升
$G_j, D \in R[x_1, \ldots, x_n]$。置
$y_{n + j}' = D(y_1, \ldots, y_n) y_{n + j} - G_j(y_1, \ldots, y_n)$。
由构造，$y_{n + j}' \in \mathfrak p S$。
显然，$y_1, \ldots, y_n, y_n', \ldots, y_{n + m}'$
生成 $S$ 的一个 $R$-子代数，其局部化为 $S$。
定义
$$
\Psi : R[x_1, \ldots, x_{n + m}] \to S
$$
为把 $x_i$（$i = 1, \ldots, n$）映到 $y_i$，
并把 $x_{n + j}$（$j = 1, \ldots, m$）映到 $y'_{n + j}$ 的映射。
由构造，性质 (1) 与 (4) 显然成立。此外，理想
$\mathfrak b$ 满射到多项式环 $\kappa(\mathfrak p)[x_1, \ldots, x_{n + m}]$ 中的理想
$(\mathfrak a, x_{n + 1}, \ldots, x_{n + m})$。

\medskip\noindent
记 $J = \Ker(\Psi)$。有短正合列
$$
0 \to J_{\mathfrak b}
\to R[x_1, \ldots, x_{n + m}]_{\mathfrak b}
\to S_{\mathfrak q}
\to 0.
$$
满射性来自上面对
$y_1, \ldots, y_n, y_n', \ldots, y_{n + m}'$ 的选择。
这意味着
$$
J_{\mathfrak b}/ \mathfrak pJ_{\mathfrak b}
\to \kappa(\mathfrak p)[x_1, \ldots, x_{n + m}]_{
(\mathfrak a, x_{n + 1}, \ldots, x_{n + m})}
\to S_{\mathfrak q}/\mathfrak pS_{\mathfrak q}
\to 0
$$
正合。由构造，在最后一个映射下 $x_i$ 映到 $\psi(x_i)$，
$x_{n + j}$ 映到零。
因此容易按引理的 (2) 与 (3) 选择 $f_i$。
\end{proof}

\begin{lemma}
\label{obsolete-lemma-smooth-independent-presentation}
设 $R \to S$ 是有限表示的环映射。
若对 $S$ 在 $R$ 上的某个表示 $\alpha$，其
朴素余切复形 $\NL(\alpha)$ 拟同构于置于次数 $0$ 的
有限投射 $S$-模，则
对任意表示都成立。
\end{lemma}

\begin{proof}
立即由 Algebra，引理 \ref{algebra-lemma-NL-homotopy} 得出。
\end{proof}

\begin{remark}[投射分解]
\label{obsolete-remark-projective-resolution}
设 $R$ 是环。
对任意集合 $S$，记 $F(S)$ 为以 $S$ 为基的自由 $R$-模。
则任意左 $R$-模都有如下两步分解
$$
F(M \times M) \oplus F(R \times M) \to F(M) \to M \to 0.
$$
第一个映射由规则
$$
[m_1, m_2] \oplus [r, m] \mapsto [m_1 + m_2] - [m_1] - [m_2] + [rm] - r[m].
$$
\end{remark}

\begin{lemma}
\label{obsolete-lemma-spec-localization-first}
设 $S$ 是 $A$ 的乘法闭集。则映射
$$
f: \Spec(S^{-1}A)\longrightarrow \Spec(A)
$$
由典范环映射
$A \to S^{-1}A$ 诱导的映射到其像上是同胚，且
$\Im(f) = \{ \mathfrak p \in \Spec(A) : \mathfrak p\cap S = \emptyset \}$。
\end{lemma}

\begin{proof}
这是 Algebra，引理 \ref{algebra-lemma-spec-localization} 的重复。
\end{proof}

\begin{lemma}
\label{obsolete-lemma-finite-type-flat-over-integral-algebra}
设 $A \to B$ 是有限型平坦环映射，且 $A$ 是整
环。则 $B$ 是有限表示的 $A$-代数。
\end{lemma}

\begin{proof}
这是 More on Flatness，命题
\ref{flat-proposition-flat-finite-type-finite-presentation-domain} 的特殊情形。
\end{proof}

\begin{lemma}
\label{obsolete-lemma-helper-finite-type-flat-finite-presentation}
设 $R$ 是分式域为 $K$ 的整环。
设 $S = R[x_1, \ldots, x_n]$ 是 $R$ 上的多项式环。
设 $M$ 是有限 $S$-模，并假设 $M$ 平坦于 $R$。
若对每个子环 $R \subset R' \subset K$ 且 $R \not = R'$，
$M \otimes_R R'$ 在
$S \otimes_R R'$ 上有限表示，则 $M$ 在 $S$ 上有限表示。
\end{lemma}

\begin{proof}
该引理成立，是因为即使不作引理陈述中对每个 $R'$ 的
$M \otimes_R R'$ 有限表示这一假设，$M$ 仍是有限表示的。
这由 More on Flatness，命题
\ref{flat-proposition-flat-finite-type-finite-presentation-domain} 得出。
这个引理原先有一个错误的证明（感谢 Ofer Gabber
发现其中的缺口），并曾用于所引命题的另一种证明。要恢复这个引理，
在 $R$ 为局部正规整环时，我们需要只使用交换代数各章结果的
正确论证；若有论证，请发送邮件至
\href{mailto:stacks.project@gmail.com}{stacks.project@gmail.com}。
\end{proof}

\begin{lemma}
\label{obsolete-lemma-relative-effective-cartier-algebra}
设 $A \to B$ 是环映射，且 $f \in B$。假设
\begin{enumerate}
\item $A \to B$ 平坦；
\item $f$ 是非零因子；并且
\item $A \to B/fB$ 平坦。
\end{enumerate}
于是对任意理想 $I \subset A$，映射
$f : B/IB \to B/IB$ 是单射。
\end{lemma}

\begin{proof}
注意，$IB = I \otimes_A B$ 以及 $I(B/fB) = I \otimes_A B/fB$。
这是因为 $B$ 与 $B/fB$ 都平坦于 $A$。
特别地，$IB/fIB \cong I \otimes_A B/fB$ 单射地映入
$B/fB$。因此将蛇形引理应用于下列图表即可得到结论。
$$
\xymatrix{
0 \ar[r] &
I \otimes_A B \ar[r] \ar[d]^f &
B \ar[r] \ar[d]^f &
B/IB \ar[r] \ar[d]^f &
0 \\
0 \ar[r] &
I \otimes_A B \ar[r] &
B \ar[r] &
B/IB \ar[r] &
0
}
$$
其中各行都是正合的。
\end{proof}

\begin{lemma}
\label{obsolete-lemma-faithfully-flat-injective}
若 $R \to S$ 是忠实平坦的环映射，则对每个 $R$-模
$M$，映射 $M \to S \otimes_R M$，$x \mapsto 1 \otimes x$ 是单射。
\end{lemma}

\begin{proof}
这是 Algebra，引理
\ref{algebra-lemma-faithfully-flat-universally-injective} 的重复。
\end{proof}

\begin{remark}
\label{obsolete-remark-section-colimits}
这个参考/标签原先指向 Smoothing Ring Maps 一章中的某个小节，
但相关材料后来已并入 Algebra 的小节
\ref{algebra-section-colimits-flat}。
\end{remark}

\begin{lemma}
\label{obsolete-lemma-not-domain}
设 $(R, \mathfrak m)$ 是维数为 $1$ 的约化诺特局部环，
且 $x \in \mathfrak m$ 是非零因子。设
$\mathfrak q_1, \ldots, \mathfrak q_r$ 是 $R$ 的极小素理想。则
$$
\text{length}_R(R/(x)) = \sum\nolimits_i \text{ord}_{R/\mathfrak q_i}(x)
$$
\end{lemma}

\begin{proof}
这是 Chow Homology，引理
\ref{chow-lemma-additivity-divisors-restricted} 的特殊（非常容易）情形。
\end{proof}

\begin{lemma}
\label{obsolete-lemma-bound-primes}
设 $A$ 是维数为 $2$ 的诺特局部正规整环。
对非零的 $f \in \mathfrak m$，记
$\text{div}(f) = \sum n_i (\mathfrak p_i)$
为 $f$ 在 $A$ 的去心谱上的因子。
置 $|f| = \sum n_i$。存在整数 $N$ 与 $M$，
使得对所有 $g \in \mathfrak m^N$ 都有 $|f + g| \leq M$。
\end{lemma}

\begin{proof}
取 $h \in \mathfrak m$ 使得 $f, h$ 是 $A$ 中的正则序列
（这由 Algebra，引理 \ref{algebra-lemma-criterion-normal} 与
\ref{algebra-lemma-depth-drops-by-one} 推出）。
取 $M = \text{length}_A(A/(f, h))$，并取满足
$\mathfrak m^N \subset (f, h)$ 的任意整数 $N$ 来证明引理。这样的
$N$ 存在，因为 $\sqrt{(f, h)} = \mathfrak m$。注意，对所有 $g \in \mathfrak m^N$，
都有 $M = \text{length}_A(A/(f + g, h))$，因为
$(f, h) = (f + g, h)$。这也说明 $f + g, h$
同样是 $A$ 中的正则序列，见
Algebra，引理 \ref{algebra-lemma-reformulate-CM}。
现在设 $\text{div}(f + g ) = \sum m_j (\mathfrak q_j)$。
考虑映射
$$
c : A/(f + g) \longrightarrow \prod A/\mathfrak q_j^{(m_j)}
$$
其中 $\mathfrak q_j^{(m_j)}$ 是符号幂，见
Algebra，小节 \ref{algebra-section-symbolic-power}。
由于 $A$ 正规，可知 $A_{\mathfrak q_i}$ 是
离散赋值环，因此
$$
A_{\mathfrak q_i}/(f + g) =
A_{\mathfrak q_i}/\mathfrak q_i^{m_i} A_{\mathfrak q_i} =
(A/\mathfrak q_i^{(m_i)})_{\mathfrak q_i}
$$
由于 $V(f + g, h) = \{\mathfrak m\}$，这意味着 $c$ 在令 $h$ 可逆后
成为同构（略去一个细节）。由于 $h$ 是 $A/(f + g)$ 上的非零因子，
可知 $A/(f + g, h)$ 的长度等于 Chow Homology 小节
\ref{chow-section-periodic-complexes} 中定义的 Herbrand 商
$e_A(A/(f + g), 0, h)$。
类似地，$A/(h, \mathfrak q_j^{(m_j)})$ 的长度等于
$e_A(A/\mathfrak q_j^{(m_j)}, 0, h)$。于是
\begin{align*}
M & = \text{length}_A(A/(f + g, h) \\
& =
e_A(A/(f + g), 0, h) \\
& =
\sum\nolimits_i e_A(A/\mathfrak q_j^{(m_j)}, 0, h) \\
& =
\sum\nolimits_i \sum\nolimits_{m = 0, \ldots, m_j - 1}
e_A(\mathfrak q_j^{(m)}/\mathfrak q_j^{(m + 1)}, 0, h)
\end{align*}
这些等式由 Chow Homology，引理
\ref{chow-lemma-additivity-periodic-length} 与
\ref{chow-lemma-finite-periodic-length} 推出；特别使用了
上面所述 $c$ 的余核具有有限长度这一事实。
容易证明，
$e_A(\mathfrak q^{(m)}/\mathfrak q^{(m + 1)}, 0, h)$
由 Nakayama 引理可知至少为 $1$。引理得证。
\end{proof}

\begin{lemma}
\label{obsolete-lemma-flat-over-gorenstein-gorenstein-fibre}
设 $A \to B$ 是诺特局部环之间的平坦局部同态。
若 $A$ 与 $B/\mathfrak m_A B$ 都是 Gorenstein 环，则 $B$ 是 Gorenstein 环。
\end{lemma}

\begin{proof}
立即由 Dualizing Complexes，引理
\ref{dualizing-lemma-flat-under-gorenstein} 得出。
\end{proof}

\begin{lemma}
\label{obsolete-lemma-kill-local}
设 $(A, \mathfrak m)$ 是诺特局部环。
设 $I \subset A$ 是理想，$M$ 是有限 $A$-模。
设 $s$ 是整数。假设
\begin{enumerate}
\item $A$ 有对偶化复形；
\item 若 $\mathfrak p \not \in V(I)$ 且
$V(\mathfrak p) \cap V(I) \not = \{\mathfrak m\}$，则
$\text{depth}_{A_\mathfrak p}(M_\mathfrak p) + \dim(A/\mathfrak p) > s$。
\end{enumerate}
则存在 $n > 0$ 与理想 $J \subset A$，满足
$V(J) \cap V(I) = \{\mathfrak m\}$，且 $JI^n$ 湮灭
$i \leq s$ 时的 $H^i_\mathfrak m(M)$。
\end{lemma}

\begin{proof}
根据
Local Cohomology，引理 \ref{local-cohomology-lemma-sitting-in-degrees}，
只需对有限 $A$-模
$E^i = \text{Ext}^{-i}_A(M, \omega_A^\bullet)$
（$i \leq s$）证明这一点。$E^0 \oplus \ldots \oplus E^s$ 的支集 $Z$
在 $\Spec(A)$ 中闭合，且不包含满足 (2) 的任何素理想。
因此对某个符合引理陈述的 $J$，它包含于
$V(JI^n)$ 中。
\end{proof}

\begin{lemma}
\label{obsolete-lemma-algebraize-local-cohomology-bis-bis}
设 $(A, \mathfrak m)$ 是诺特局部环。
设 $I \subset A$ 是理想，$M$ 是有限 $A$-模。
设 $s$ 与 $d$ 是整数。假设
\begin{enumerate}
\item[(a)] $A$ 有对偶化复形；
\item[(b)] $\text{cd}(A, I) \leq d$；
\item[(c)] 若 $\mathfrak p \not \in V(I)$，则
$\text{depth}_{A_\mathfrak p}(M_\mathfrak p) > s$，或
$\text{depth}_{A_\mathfrak p}(M_\mathfrak p) + \dim(A/\mathfrak p) > d + s$。
\end{enumerate}
则 Algebraic and Formal Geometry，引理
\ref{algebraization-lemma-algebraize-local-cohomology-bis} 对
$A, I, \mathfrak m, M$ 成立，并且
$H^i_\mathfrak m(M) \to \lim H^i_\mathfrak m(M/I^nM)$
在 $i \leq s$ 时是同构；这些模
被 $I$ 的某个幂湮灭。
\end{lemma}

\begin{proof}
假设由更一般的 Algebraic and Formal Geometry，引理
\ref{algebraization-lemma-algebraize-local-cohomology-bis}，以及
Algebraic and Formal Geometry，引理
\ref{algebraization-lemma-bootstrap-bis-bis} 给出。
然后 Algebraic and Formal Geometry，引理
\ref{algebraization-lemma-algebraize-local-cohomology-bis} 的结论
给出第二个断言。
\end{proof}

\begin{lemma}
\label{obsolete-lemma-combine-one}
在 Algebraic and Formal Geometry，情形
\ref{algebraization-situation-bootstrap} 中，
有 $H^s_\mathfrak a(M) = \lim H^s_\mathfrak a(M/I^nM)$。
\end{lemma}

\begin{proof}
这立即由 Algebraic and Formal Geometry，定理
\ref{algebraization-theorem-final-bootstrap} 得出。
这个引理的原版本带有额外假设，
已被该定理取代。
\end{proof}

\begin{lemma}
\label{obsolete-lemma-fully-faithful-simple-one}
设 $A$ 是诺特环，且 $f \in \mathfrak a$ 属于 $A$ 的某个理想。
设 $U = \Spec(A) \setminus V(\mathfrak a)$。假设
\begin{enumerate}
\item $A$ 有对偶化复形，且关于 $f$ 完备；
\item $A_f$ 满足 $(S_2)$，并且对每个极小素理想 $\mathfrak p \subset A$，
$f \not \in \mathfrak p$，且对
$\mathfrak q \in V(\mathfrak p) \cap V(\mathfrak a)$ 有
$\dim((A/\mathfrak p)_\mathfrak q) \geq 3$。
\end{enumerate}
则完备化函子
$$
\textit{Coh}(\mathcal{O}_U)
\longrightarrow
\textit{Coh}(U, I\mathcal{O}_U),
\quad
\mathcal{F} \longmapsto \mathcal{F}^\wedge
$$
在有限局部自由对象构成的满子范畴上是全忠实的。
\end{lemma}

\begin{proof}
这是 Algebraic and Formal Geometry，引理
\ref{algebraization-lemma-fully-faithful-simple-two} 的特殊情形。
\end{proof}





\section{与 ZMT 相关的引理}
\label{obsolete-section-ZMT}

\noindent
本节的引理原先用于证明（代数版本的）Zariski 主定理，
即 Algebra，定理 \ref{algebra-theorem-main-theorem}。

\begin{lemma}
\label{obsolete-lemma-change-equation-multiply}
设 $R$ 是环，且 $\varphi : R[x] \to S$ 是
环映射。设 $t \in S$。若 $t$ 在
$R[x]$ 上整，则存在 $\ell \geq 0$，使得对每个 $a \in R$，元素
$\varphi(a)^\ell t$ 关于由下式定义的
$\varphi_a : R[y] \to S$ 是整的：
$y \mapsto \varphi(ax)$ 且 $r \mapsto \varphi(r)$，
其中 $r\in R$。
\end{lemma}

\begin{proof}
写成 $t^d + \sum_{i < d} \varphi(f_i)t^i = 0$，其中
$f_i \in R[x]$。令 $\ell$ 为所有
$f_i$ 关于 $x$ 的最大次数。将等式乘以 $\varphi(a)^\ell$，得到
$\varphi(a)^\ell t^d + \sum_{i < d} \varphi(a^\ell f_i)t^i = 0$。
注意，每个 $\varphi(a^\ell f_i)$ 都在
$\varphi_a$ 的像中。结论由
Algebra，引理 \ref{algebra-lemma-make-integral-trivial} 得出。
\end{proof}

\begin{lemma}
\label{obsolete-lemma-make-integral-less-trivial}
设 $\varphi : R \to S$ 是环映射。
设 $t \in S$ 满足关系
$\varphi(a_0) + \varphi(a_1)t + \ldots + \varphi(a_n) t^n = 0$。
置 $u_n = \varphi(a_n)$，$u_{n-1} = u_n t + \varphi(a_{n-1})$，
依此递推至 $u_1 = u_2 t + \varphi(a_1)$。
则 $u_n, u_{n-1}, \ldots, u_1$ 以及
$u_nt, u_{n-1}t, \ldots, u_1t$ 都在 $R$ 上整，
并且 $S$ 中的理想 $(\varphi(a_0), \ldots, \varphi(a_n))$ 与
$(u_n, \ldots, u_1)$ 相等。
\end{lemma}

\begin{proof}
我们对 $n$ 作归纳证明。由于 $u_n = \varphi(a_n)$，由
Algebra，引理 \ref{algebra-lemma-make-integral-trivial}
可得 $u_nt$ 在 $R$ 上整。当然
$u_n = \varphi(a_n)$ 在 $R$ 上整。于是
$u_{n - 1} = u_n t  + \varphi(a_{n - 1})$ 在 $R$ 上整（见
Algebra，引理 \ref{algebra-lemma-integral-closure-is-ring}），
并且有
$$
\varphi(a_0) + \varphi(a_1)t + \ldots + \varphi(a_{n - 1})t^{n - 1} +
u_{n - 1}t^{n - 1} = 0.
$$
因此将归纳假设应用于映射
$S' \to S$（其中 $S'$ 是 $S$ 中 $R$ 的整闭包）以及上面的等式，
可知 $u_{n-1}, \ldots, u_1$ 与 $u_{n-1}t, \ldots, u_1t$
也都属于 $S'$。关于理想的断言直接由
这些元素的形式以及 $u_1t + \varphi(a_0) = 0$ 得出。
\end{proof}

\begin{lemma}
\label{obsolete-lemma-make-integral-not-in-ideal}
设 $\varphi : R \to S$ 是环映射。
设 $t \in S$ 满足
关系 $\varphi(a_0) + \varphi(a_1)t + \ldots + \varphi(a_n) t^n = 0$。
设 $J \subset S$ 是理想，并且至少存在一个 $i$ 使
$\varphi(a_i) \not \in J$。
则存在 $u \in S$，$u \not\in J$，使得
$u$ 与 $ut$ 都在 $R$ 上整。
\end{lemma}

\begin{proof}
这立即由引理 \ref{obsolete-lemma-make-integral-less-trivial} 得出，
因为某个元素 $u_i$ 不属于 $J$。
\end{proof}

\noindent
下面两个引理给出一种描述由一个（非退化）等式
切出的 $\mathbf{P}^1_R$ 闭子概形的方法。

\begin{lemma}
\label{obsolete-lemma-P1}
设 $R$ 是环。
设 $F(X, Y) \in R[X, Y]$ 是次数为
$d$ 的齐次多项式。假设对 $R$ 的每个素理想 $\mathfrak p$，
$F$ 至少有一个系数不属于 $\mathfrak p$。
设 $S = R[X, Y]/(F)$ 为分次环。
则对所有 $n \geq d$，$R$-模 $S_n$
都是秩为 $d$ 的有限局部自由模。
\end{lemma}

\begin{proof}
$R$-模 $S_n$ 有如下表示：
$$
R[X, Y]_{n-d} \to R[X, Y]_n \to S_n \to 0.
$$
因此由 Algebra，引理 \ref{algebra-lemma-cokernel-flat}，
只需证明乘以 $F$ 诱导单射
$\kappa(\mathfrak p)[X, Y]
\to \kappa(\mathfrak p)[X, Y]$
对所有素理想 $\mathfrak p$ 都成立。
这由 $F$ 模 $\mathfrak p$ 不映为零多项式的假设立即可见。
秩的断言也同样显然。
\end{proof}

\begin{lemma}
\label{obsolete-lemma-rel-prime-pols}
设 $k$ 是域，且 $F, G \in k[X, Y]$ 齐次，
次数分别为 $d, e$。假设 $F, G$ 互素。
则乘以 $G$ 在 $S = k[X, Y]/(F)$ 上是单射。
\end{lemma}

\begin{proof}
这是一种定义“互素”的方式。若采用另一种
定义，可以证明它与此定义等价。
\end{proof}

\begin{lemma}
\label{obsolete-lemma-P1-localize}
设 $R$ 是环。设 $F(X, Y) \in R[X, Y]$ 是次数为
$d$ 的齐次多项式，且分次环 $S = R[X, Y]/(F)$。
设 $\mathfrak p \subset R$ 是素理想，且
$F$ 的某个系数不属于 $\mathfrak p$。
则存在 $f \in R$，$f \not\in \mathfrak p$，
整数 $e$ 及 $G \in R[X, Y]_e$，
使得乘以 $G$ 对所有 $n \geq d$ 诱导同构
$(S_n)_f \to (S_{n + e})_f$。
\end{lemma}

\begin{proof}
在证明过程中，我们可以用 $R_f$ 替换 $R$，其中
$f\in R$ 且 $f\not\in \mathfrak p$（有限次进行）。
首先作这样的替换，使得
$F$ 的某个系数在 $R$ 中可逆。
特别地，由引理 \ref{obsolete-lemma-P1}，现在对 $n \geq d$，模 $S_n$
是秩为 $d$ 的局部自由模。取任意 $G \in R[X, Y]_e$，使其在
$G$ 在 $\kappa(\mathfrak p)[X, Y]$ 中的像与 $F(X, Y)$ 的像
互素（对某个 $e$ 这是可能的）。
将 Algebra，引理 \ref{algebra-lemma-cokernel-flat} 应用于由乘以 $G$
从 $S_d \to S_{d + e}$ 诱导的映射。由 $G$ 的选择及
引理 \ref{obsolete-lemma-rel-prime-pols}，可知
$S_d \otimes \kappa(\mathfrak p) \to S_{d + e} \otimes \kappa(\mathfrak p)$
是双射。因此，对适当的 $f$ 用 $R_f$ 替换 $R$ 后，可设
$G : S_d \to S_{d + e}$ 是双射。这又意味着
对 $R$ 的所有素理想 $\mathfrak p'$，$G$ 在
$\kappa(\mathfrak p')[X, Y]$ 中的像与 $F$ 的像
互素。于是再次由 Algebra，引理 \ref{algebra-lemma-cokernel-flat}
可知所有映射
$G : S_d \to S_{d + e}$（$n \geq d$）都是同构。
\end{proof}

\begin{remark}
\label{obsolete-remark-algebra}
设 $R$ 是环。假设有 $F \in R[X, Y]_d$ 与
$G \in R[X, Y]_e$，置 $S = R[X, Y]/(F)$，并满足：(1) 对
所有 $n \geq d$，$S_n$ 是秩为 $d$ 的有限局部自由模；(2) 乘以 $G$
对所有 $n \geq d$ 诱导同构 $S_n \to S_{n + e}$。此时
可以如下定义有限局部自由的 $R$-代数
$A$：
\begin{enumerate}
\item 作为 $R$-模，$A = S_{ed}$；
\item 乘法 $A \times A \to A$ 由下列规则给出：
在 $S_{2ed}$ 中，$H_1 H_2 = H_3$ 当且仅当 $G^d H_3 = H_1 H_2$。
\end{enumerate}
这是合理的，因为乘以 $G^d$
诱导双射 $S_{de} \to S_{2de}$。
容易看出这定义了一个环结构。
需要注意一个容易混淆的事实：元素 $G^d$
定义了环 $A$ 的单位元。
\end{remark}

\begin{lemma}
\label{obsolete-lemma-finite-after-localization}
设 $R$ 是环，且 $f \in R$。
设有 $S$、$S'$ 及实线箭头，
组成下列环的交换图：
$$
\xymatrix{
& S'' \ar@{-->}[rd] \ar@{-->}[dd] &
\\
R \ar[rr] \ar@{-->}[ru] \ar[d] &  & S \ar[d]
\\
R_f \ar[r] & S' \ar[r] & S_f
}
$$
假设 $R_f \to S'$ 是有限映射。则可找到
有限环映射 $R \to S''$ 及图中所示的虚线箭头，
使得 $S' = (S'')_f$。
\end{lemma}

\begin{proof}
具体地，设 $S'$ 由
$x_i$ 在 $R_f$ 上生成，其中 $i = 1, \ldots, w$。令 $P_i(t) \in R_f[t]$
为满足 $P_i(x_i) = 0$ 的首一多项式。设
$P_i$ 的次数为 $d_i > 0$。写成
$P_i(t) = t^{d_i} + \sum_{j < d_i} (a_{ij}/f^n) t^j$，
其中 $n$ 统一。再将
$x_i$ 在 $S_f$ 中的像写成 $g_i / f^n$，
其中 $g_i \in S$ 适当。于是元素
$\xi_i = f^{nd_i} g_i^{d_i} + \sum_{j < d_i} f^{n(d_i - j)} a_{ij} g_i^j$
$S$ 被 $f$ 的某个幂湮灭。
因此将 $n$ 增大为 $n'$（这会把
$g_i$ 替换为 $f^{n' - n}g_i$）后，可设 $\xi_i = 0$。
于是 $S'$ 由元素
$f^n x_i$ 生成；每个元素都是下列首一多项式的根：
$Q_i(t) = t^{d_i} + \sum_{j < d_i} f^{n(d_i - j)} a_{ij} t^j$，
其系数在 $R$ 中。按构造，
$Q_i(f^ng_i) = 0$ 在 $S$ 中。因此得到有限 $R$-代数
$S'' = R[z_1, \ldots, z_w]/(Q_1(z_1), \ldots, Q_w(z_w))$，
它适合于上面的交换图。
映射 $\alpha : S'' \to S$ 将 $z_i$ 映到 $f^ng_i$，
映射 $\beta : S'' \to S'$ 将 $z_i$ 映到 $f^nx_i$。
此时 $\beta$ 可能尚未诱导
同构 $(S'')_f \cong S'$。
目前只知道该映射是满射。问题在于，可能存在
元素 $h/f^n \in (S'')_f$ 映到 $S'$ 中的零但自身不为零。此时 $\beta(h)$
是 $S$ 中的元素，且对某个 $N$ 有 $f^N \beta(h) = 0$。因此 $f^N h$ 属于理想
$J = \{h \in S'' \mid \alpha(h) = 0 \text{ 且 }
\beta(h) = 0\}$ 的 $S''$。容易看出 $S''/J$ 即可满足要求。
\end{proof}




\section{形式光滑的环映射}
\label{obsolete-section-formally-smooth}

\begin{lemma}
\label{obsolete-lemma-formally-smooth-smooth}
设 $R$ 是环，$S$ 是 $R$-代数。
若 $S$ 有限表示且形式光滑于 $R$，
则 $S$ 光滑于 $R$。
\end{lemma}

\begin{proof}
见 Algebra，命题 \ref{algebra-proposition-smooth-formally-smooth}。
\end{proof}

\begin{remark}
\label{obsolete-remark-equation-derivatives}
该标签原先指向 Algebraization of Formal Spaces，命题
\ref{restricted-proposition-approximate} 的证明中的一个等式，
由于材料重新编排而不再使用。
\end{remark}

\begin{remark}
\label{obsolete-remark-equation-ci}
该标签原先指向 Algebraization of Formal Spaces，命题
\ref{restricted-proposition-approximate} 的证明中的一个等式，
由于材料重新编排而不再使用。
\end{remark}

\begin{remark}
\label{obsolete-remark-equation-in-ideal}
该标签原先指向 Algebraization of Formal Spaces，命题
\ref{restricted-proposition-approximate} 的证明中的一个等式，
由于材料重新编排而不再使用。
\end{remark}

\begin{remark}
\label{obsolete-remark-equation-derivatives-analogue}
该标签原先指向 Algebraization of Formal Spaces，命题
\ref{restricted-proposition-approximate} 的证明中的一个等式，
由于材料重新编排而不再使用。
\end{remark}

\begin{remark}
\label{obsolete-remark-equation-go-down}
该标签原先指向 Algebraization of Formal Spaces，引理
\ref{restricted-lemma-lift-approximation} 的证明中的一个等式，
由于材料重新编排而不再使用。
\end{remark}

\begin{lemma}
\label{obsolete-lemma-get-morphism-general}
设 $A$ 是诺特环，且 $I \subset A$ 是理想。
设 $t$ 是 $I$ 的最小生成元个数。
设 $C$ 是关于 $I$-进拓扑完备的诺特 $A$-代数。
存在只依赖于 $I \subset A \to C$ 的整数 $d \geq 0$，
具有如下性质：给定
\begin{enumerate}
\item $c \geq 0$ 及 Algebraization of Formal Spaces，
等式 (\ref{restricted-equation-C-prime}) 中的 $B$，满足对 $a \in I^c$，
乘以 $a$ 在 $\NL_{B/A}^\wedge$ 上于 $D(B)$ 中为零；
\item 整数 $n > 2t\max(c, d)$；
\item $A/I^n$-代数映射 $\psi_n : B/I^nB \to C/I^nC$；
\end{enumerate}
存在 $A$-代数映射 $\varphi : B \to C$，使得
$\psi_n \bmod I^{m - c} = \varphi \bmod I^{m - c}$，
其中 $m = \lfloor \frac{n}{t} \rfloor$。
\end{lemma}

\begin{proof}
这个引理已被更强的
Algebraization of Formal Spaces，引理
\ref{restricted-lemma-get-morphism-general-better} 所废止。
事实上，我们将由该引理推出本引理。

\medskip\noindent
设 $I \subset A \to C$ 如上面的陈述。
记 $d(\text{Gr}_I(C))$ 与 $q(\text{Gr}_I(C))$ 为
Local Cohomology，小节 \ref{local-cohomology-section-uniform} 中得到的整数。
注意，$t$ 是 $IC$ 的最小生成元个数的上界，因此
$d(\text{Gr}_I(C)) + 1 \leq t$；见
Local Cohomology，小节 \ref{local-cohomology-section-uniform} 中的讨论。
可假设 $t \geq 1$，因为否则引理没有内容。
我们断言取
$$
d = q(\text{Gr}_I(C))
$$
具体地，设 $c$、$B$、$n$、$\psi_n$ 如上面的陈述。
则有
$$
n > 2t\max(c, d) \Rightarrow n \geq 2tc + 1 \Rightarrow
n \geq 2(d(\text{Gr}_I(C)) + 1)c + 1
$$
另一方面，有
$$
n > 2t\max(c, d) \Rightarrow n > t(c + d) \Rightarrow
n \geq q(C) + tc \geq q(\text{Gr}_I(C)) + (d(\text{Gr}_I(C)) + 1)c
$$
因此 Algebraization of Formal Spaces，引理
\ref{restricted-lemma-get-morphism-general-better} 的假设
得到满足，从而得到 $A$-代数同态
$\varphi : B \to C$，它与 $\psi_n$ 在
模 $I^{n - (d(\text{Gr}_I(C)) + 1)c}C$ 上同余。
由于
\begin{align*}
n - (d(\text{Gr}_I(C)) + 1)c
& = \frac{n}{t} + \frac{(t - 1)n}{t} - (d(\text{Gr}_I(C)) + 1)c \\
& \geq \frac{n}{t} + \frac{(d(\text{Gr}_I(C))n}{t} - (d(\text{Gr}_I(C)) + 1)c \\
& > \frac{n}{t} + \frac{d(\text{Gr}_I(C))2tc}{t} - (d(\text{Gr}_I(C)) + 1)c \\
& = \frac{n}{t} + 2d(\text{Gr}_I(C))c - (d(\text{Gr}_I(C)) + 1)c \\
& = \frac{n}{t} + d(\text{Gr}_I(C))c - c \\
& \geq m - c
\end{align*}
可知如所需，$\varphi$ 与 $\psi_n$ 在模 $I^{m - c}C$ 上同余。
\end{proof}




\section{网站与层}
\label{obsolete-section-sites}

\begin{remark}[不存在从零延拓到推前的映射]
\label{obsolete-remark-from-shriek-to-star}
设 $U$ 是网站 $\mathcal{C}$ 的对象。对 $\mathcal{C}/U$ 上的任意阿贝尔层
$\mathcal{G}$，可以问是否存在典范映射
$$
c : j_{U!}\mathcal{G} \longrightarrow j_{U*}\mathcal{G}
$$
构造这样的映射等价于构造映射
$j_U^{-1}j_{U!}\mathcal{G} \to \mathcal{G}$。
注意限制与层化交换。
因此可以使用
Modules on Sites，引理 \ref{sites-modules-lemma-extension-by-zero} 中的预层。
所以只需对 $V/U$ 定义映射
$$
\bigoplus\nolimits_{\varphi \in \Mor_\mathcal{C}(V, U)}
\mathcal{G}(V \xrightarrow{\varphi} U)
\longrightarrow
\mathcal{G}(V/U)
$$
使其与限制相容。看起来可以令该映射在除以下求和项外的
所有求和项上为零：当 $\varphi$ 是结构态射 $\varphi_0 : V \to U$ 时，
在该项上取值 $1$。
然而这与限制映射不相容：具体地，若
$\alpha : V' \to V$ 是 $\mathcal{C}$ 中的态射，记
$V'/U$ 为 $\mathcal{C}/U$ 中结构
态射为 $\varphi'_0 = \varphi_0 \circ \alpha$ 的对象。
需要检查下图
$$
\xymatrix{
\bigoplus\nolimits_{\varphi \in \Mor_\mathcal{C}(V, U)}
\mathcal{G}(V \xrightarrow{\varphi} U)
\ar[d] \ar[r] &
\mathcal{G}(V/U) \ar[d] \\
\bigoplus\nolimits_{\varphi' \in \Mor_\mathcal{C}(V', U)}
\mathcal{G}(V' \xrightarrow{\varphi'} U)
\ar[r] &
\mathcal{G}(V'/U)
}
$$
交换。问题在于，可能存在不同于 $\varphi_0$ 的态射
$\varphi : V \to U$，满足
$\varphi \circ \alpha = \varphi'_0$。
因此左侧竖直箭头会把对应于 $\varphi$ 的求和项送到下方
水平箭头取值为 $1$ 的求和项，图表几乎肯定不交换。
\end{remark}




\section{上同调}
\label{obsolete-section-cohomology}

\noindent
关于上同调的已废弃引理。

\begin{lemma}
\label{obsolete-lemma-ML-general}
设 $I$ 是环 $A$ 的理想，$X$ 是位于 $\Spec(A)$ 上的概形。设
$$
\ldots \to \mathcal{F}_3 \to \mathcal{F}_2 \to \mathcal{F}_1
$$
是 $\mathcal{O}_X$-模的逆系统，
满足 $\mathcal{F}_n = \mathcal{F}_{n + 1}/I^n\mathcal{F}_{n + 1}$。
假设
$$
\bigoplus\nolimits_{n \geq 0} H^1(X, I^n\mathcal{F}_{n + 1})
$$
作为分次
$\bigoplus_{n \geq 0} I^n/I^{n + 1}$-模满足升链条件。
则逆系统 $M_n = \Gamma(X, \mathcal{F}_n)$ 满足
Mittag--Leffler 条件。
\end{lemma}

\begin{proof}
这是更一般的
Cohomology，引理 \ref{cohomology-lemma-ML-general} 的特殊情形。
\end{proof}

\begin{lemma}
\label{obsolete-lemma-ML-general-better}
设 $I$ 是环 $A$ 的理想，$X$ 是 $\Spec(A)$ 上的概形。设
$$
\ldots \to \mathcal{F}_3 \to \mathcal{F}_2 \to \mathcal{F}_1
$$
为 $\mathcal{O}_X$-模的逆系统，
满足 $\mathcal{F}_n = \mathcal{F}_{n + 1}/I^n\mathcal{F}_{n + 1}$。
给定 $n$，定义
$$
H^1_n =
\bigcap\nolimits_{m \geq n}
\Im\left(
H^1(X, I^n\mathcal{F}_{m + 1}) \to H^1(X, I^n\mathcal{F}_{n + 1})
\right)
$$
若 $\bigoplus H^1_n$ 作为分次
$\bigoplus_{n \geq 0} I^n/I^{n + 1}$-模满足升链条件，则逆系统
$M_n = \Gamma(X, \mathcal{F}_n)$ 满足 Mittag--Leffler 条件。
\end{lemma}

\begin{proof}
这是更一般的
Cohomology，引理 \ref{cohomology-lemma-ML-general-better} 的特殊情形。
\end{proof}

\begin{lemma}
\label{obsolete-lemma-topology-I-adic-general}
设 $I$ 是环 $A$ 的有限生成理想。
设 $X$ 是位于 $\Spec(A)$ 上的概形。设
$$
\ldots \to \mathcal{F}_3 \to \mathcal{F}_2 \to \mathcal{F}_1
$$
是 $\mathcal{O}_X$-模的逆系统，满足
$\mathcal{F}_n = \mathcal{F}_{n + 1}/I^n\mathcal{F}_{n + 1}$。假设
$$
\bigoplus\nolimits_{n \geq 0} H^0(X, I^n\mathcal{F}_{n + 1})
$$
作为分次
$\bigoplus_{n \geq 0} I^n/I^{n + 1}$-模满足升链条件。
则 $M = \lim \Gamma(X, \mathcal{F}_n)$ 上的极限拓扑
就是 $I$-进拓扑。
\end{lemma}

\begin{proof}
这是更一般的
Cohomology，引理 \ref{cohomology-lemma-topology-I-adic-general} 的特殊情形。
\end{proof}

\begin{lemma}
\label{obsolete-lemma-pullback-K-flat}
设 $(\Sh(\mathcal{C}), \mathcal{O}_\mathcal{C})$ 是环化拓扑斯。
对任意 $\mathcal{O}_\mathcal{C}$-模复形 $\mathcal{G}^\bullet$，
存在拟同构 $\mathcal{K}^\bullet \to \mathcal{G}^\bullet$，使得
对任意环化拓扑斯态射
$f : (\Sh(\mathcal{D}), \mathcal{O}_\mathcal{D}) \to
(\Sh(\mathcal{C}), \mathcal{O}_\mathcal{C})$，$f^*\mathcal{K}^\bullet$ 是
$\mathcal{O}_\mathcal{D}$-模的 K-平坦复形。
\end{lemma}

\begin{proof}
这由 Cohomology on Sites，引理
\ref{sites-cohomology-lemma-K-flat-resolution} 与
\ref{sites-cohomology-lemma-pullback-K-flat} 得出。
\end{proof}

\begin{remark}
\label{obsolete-remark-pullback-K-flat}
本注记原先讨论 K-平坦复形的拉回是否仍为 K-平坦，
但现在已被
Cohomology on Sites，引理 \ref{sites-cohomology-lemma-pullback-K-flat} 废止。
\end{remark}

\noindent
下面的引理在特殊情形下计算
导出极限的上同调层。

\begin{lemma}
\label{obsolete-lemma-Rlim-of-system}
设 $(\mathcal{C}, \mathcal{O})$ 是环化网站。设 $(K_n)$
是 $D(\mathcal{O})$ 中对象的逆系统。
设 $\mathcal{B} \subset \Ob(\mathcal{C})$ 是子集。
设 $d \in \mathbf{N}$。假设
\begin{enumerate}
\item 对所有 $n$，$K_n$ 是 $D^+(\mathcal{O})$ 中的对象；
\item 对 $q \in \mathbf{Z}$，存在
$n(q)$，使得当 $n \geq n(q)$ 时，$H^q(K_{n + 1}) \to H^q(K_n)$ 是同构；
\item $\mathcal{C}$ 的每个对象都有一个覆盖，其成员是
$\mathcal{B}$ 的元素；
\item 对每个 $U \in \mathcal{B}$，有 $H^p(U, H^q(K_n)) = 0$，
对所有 $q$ 及 $p > d$ 均成立。
\end{enumerate}
则对所有 $m \in \mathbf{Z}$，有 $H^m(R\lim K_n) = \lim H^m(K_n)$。
\end{lemma}

\begin{proof}
置 $K = R\lim K_n$。设 $U \in \mathcal{B}$。对每个 $n$ 有谱
序列
$$
H^p(U, H^q(K_n)) \Rightarrow H^{p + q}(U, K_n)
$$
它收敛，因为 $K_n$ 有下界；见
Derived Categories，引理 \ref{derived-lemma-two-ss-complex-functor}。
固定 $m \in \mathbf{Z}$，由假设 (4) 可知，只有
$H^p(U, H^q(K_n))$ 对 $H^m(U, K_n)$ 有贡献，
其中 $0 \leq p \leq d$ 且 $m - d \leq q \leq m$。由假设 (2)，
当 $n \geq \max{n(m), \ldots, n(m - d)}$ 时，
$H^m(U, K_{n + 1}) \to H^m(U, K_n)$ 是同构。函子 $R\Gamma(U, -)$
与导出极限交换，依据
Injectives，引理 \ref{injectives-lemma-RF-commutes-with-Rlim}。
因此
$$
H^m(U, K) = H^m(R\lim R\Gamma(U, K_n))
$$
另一方面，我们刚刚看到复形 $R\Gamma(U, K_n)$ 的上同调群最终恒定。因此由
More on Algebra，注记 \ref{more-algebra-remark-compare-derived-limit}，
可知对某个与 $U \in \mathcal{B}$ 无关的界，所有 $n \gg 0$ 都有
$H^m(U, K)$ 等于 $H^m(U, K_n)$。取这样的 $n$。最后回忆，$H^m(K)$ 是预层
$U \mapsto H^m(U, K)$ 的层化，而 $H^m(K_n)$ 是预层
$U \mapsto H^m(U, K_n)$ 的层化。在 $\mathcal{B}$ 的元素上，这些预层取值相同。
因此假设 (3) 保证层化也相同。
引理得证。
\end{proof}

\begin{lemma}
\label{obsolete-lemma-trivialities-cohomological-descent-abelian}
在 Simplicial Spaces，情形
\ref{spaces-simplicial-situation-simplicial-site} 中，
设 $a_0$ 是指向网站 $\mathcal{D}$ 的增广，如
Simplicial Spaces，注记 \ref{spaces-simplicial-remark-augmentation-site} 所述。
给定严格满弱 Serre 子范畴
$$
\mathcal{A} \subset \textit{Ab}(\mathcal{D}),\quad
\mathcal{A}_n \subset \textit{Ab}(\mathcal{C}_n)
$$
则
\begin{enumerate}
\item[(1)]
$\mathcal{C}_{total}$ 上所有满足其到每个 $\mathcal{C}_n$ 的限制属于 $\mathcal{A}_n$（对所有 $n$）的
阿贝尔层 $\mathcal{F}$ 构成严格满弱 Serre 子范畴
$\mathcal{A}_{total} \subset \textit{Ab}(\mathcal{C}_{total})$。
\end{enumerate}
若对所有 $n$，$a_n^{-1}$ 将 $\mathcal{A}$ 映入 $\mathcal{A}_n$，则
\begin{enumerate}
\item[(2)] $a^{-1}$ 将 $\mathcal{A}$ 映入 $\mathcal{A}_{total}$；
\item[(3)] $a^{-1}$ 将 $D_\mathcal{A}(\mathcal{D})$ 映入
$D_{\mathcal{A}_{total}}(\mathcal{C}_{total})$。
\end{enumerate}
若对所有 $n, q$，$R^qa_{n, *}$ 将 $\mathcal{A}_n$ 映入 $\mathcal{A}$，则
\begin{enumerate}
\item[(4)] 对所有 $q$，$R^qa_*$ 将 $\mathcal{A}_{total}$ 映入 $\mathcal{A}$；
并且
\item[(5)] $Ra_*$ 将 $D_{\mathcal{A}_{total}}^+(\mathcal{C}_{total})$
映入 $D_\mathcal{A}^+(\mathcal{D})$。
\end{enumerate}
\end{lemma}

\begin{proof}
唯一有实质内容的断言是 (4) 与 (5)。
(4) 由下列谱序列推出：
Simplicial Spaces，引理
\ref{spaces-simplicial-lemma-augmentation-spectral-sequence}
以及 Homology，引理 \ref{homology-lemma-biregular-ss-converges}。
然后考虑由截断给出的
$D_{\mathcal{A}_{total}}^+(\mathcal{C}_{total})$ 中对象
$K$ 的典范过滤所关联的谱序列，即得 (5)。
细节略去。
\end{proof}

\begin{remark}
\label{obsolete-remark-cohomology-topics}
该标签原先指向上调章节中列出待处理主题的一个小节。
\end{remark}

\begin{remark}
\label{obsolete-remark-sites-cohomology-topics}
该标签原先指向上调章节中列出待处理主题的一个小节。
\end{remark}

\begin{remark}
\label{obsolete-remark-V-implies-C}
该标签原先指向 Cohomology on Sites，引理
\ref{sites-cohomology-lemma-V-implies-C-general} 在
Cohomology on Sites，引理 \ref{sites-cohomology-lemma-compare-qc-zar}
所述情形下的特殊情形。
\end{remark}

\begin{remark}
\label{obsolete-remark-V-implies-cohomology}
该标签原先指向 Cohomology on Sites，引理
\ref{sites-cohomology-lemma-V-implies-cohomology-general} 在
Cohomology on Sites，引理 \ref{sites-cohomology-lemma-compare-qc-zar}
所述情形下的特殊情形。
\end{remark}

\begin{remark}
\label{obsolete-remark-induction-step-V-C}
该标签原先指向 Cohomology on Sites，引理
\ref{sites-cohomology-lemma-induction-step-V-C-general} 在
Cohomology on Sites，引理 \ref{sites-cohomology-lemma-compare-qc-zar}
所述情形下的特殊情形。
\end{remark}

\begin{remark}
\label{obsolete-remark-V-implies-C-etale-fppf}
该标签原先指向 Cohomology on Sites，引理
\ref{sites-cohomology-lemma-V-implies-C-general} 在
\'Etale Cohomology，引理 \ref{etale-cohomology-lemma-compare-fppf-etale}
所述情形下的特殊情形。
\end{remark}

\begin{remark}
\label{obsolete-remark-V-implies-cohomology-etale-fppf}
该标签原先指向 Cohomology on Sites，引理
\ref{sites-cohomology-lemma-V-implies-cohomology-general} 在
\'Etale Cohomology，引理 \ref{etale-cohomology-lemma-compare-fppf-etale}
所述情形下的特殊情形。
\end{remark}

\begin{remark}
\label{obsolete-remark-induction-step-V-C-etale-fppf}
该标签原先指向 Cohomology on Sites，引理
\ref{sites-cohomology-lemma-induction-step-V-C-general} 在
\'Etale Cohomology，引理 \ref{etale-cohomology-lemma-compare-fppf-etale}
所述情形下的特殊情形。
\end{remark}

\begin{remark}
\label{obsolete-remark-V-implies-C-etale-ph}
该标签原先指向 Cohomology on Sites，引理
\ref{sites-cohomology-lemma-V-implies-C-general} 在
\'Etale Cohomology，引理 \ref{etale-cohomology-lemma-compare-ph-etale}
所述情形下的特殊情形。
\end{remark}

\begin{remark}
\label{obsolete-remark-V-implies-cohomology-etale-ph}
该标签原先指向 Cohomology on Sites，引理
\ref{sites-cohomology-lemma-V-implies-cohomology-general} 在
\'Etale Cohomology，引理 \ref{etale-cohomology-lemma-compare-ph-etale}
所述情形下的特殊情形。
\end{remark}

\begin{remark}
\label{obsolete-remark-V-implies-cohomology-etale-ph-extra}
该标签原先指向 Cohomology on Sites，引理
\ref{sites-cohomology-lemma-V-implies-cohomology-extra-general} 在
\'Etale Cohomology，引理 \ref{etale-cohomology-lemma-compare-ph-etale}
所述情形下的特殊情形。
\end{remark}

\begin{remark}
\label{obsolete-remark-make-class-zero}
该标签原先指向 Cohomology on Sites，引理
\ref{sites-cohomology-lemma-make-class-zero-general} 在
\'Etale Cohomology，引理 \ref{etale-cohomology-lemma-compare-ph-etale}
所述情形下的特殊情形。
\end{remark}

\begin{remark}
\label{obsolete-remark-induction-step-V-C-etale-ph}
该标签原先指向 Cohomology on Sites，引理
\ref{sites-cohomology-lemma-induction-step-V-C-general} 在
\'Etale Cohomology，引理 \ref{etale-cohomology-lemma-compare-ph-etale}
所述情形下的特殊情形。
\end{remark}

\begin{remark}
\label{obsolete-remark-V-implies-C-etale-h}
该标签原先指向 Cohomology on Sites，引理
\ref{sites-cohomology-lemma-V-implies-C-general} 在
\'Etale Cohomology，引理 \ref{etale-cohomology-lemma-compare-h-etale}
所述情形下的特殊情形。
\end{remark}

\begin{remark}
\label{obsolete-remark-V-implies-cohomology-etale-h}
该标签原先指向 Cohomology on Sites，引理
\ref{sites-cohomology-lemma-V-implies-cohomology-general} 在
\'Etale Cohomology，引理 \ref{etale-cohomology-lemma-compare-h-etale}
所述情形下的特殊情形。
\end{remark}

\begin{remark}
\label{obsolete-remark-V-implies-cohomology-etale-h-extra}
该标签原先指向 Cohomology on Sites，引理
\ref{sites-cohomology-lemma-V-implies-cohomology-extra-general} 在
\'Etale Cohomology，引理 \ref{etale-cohomology-lemma-compare-h-etale}
所述情形下的特殊情形。
\end{remark}

\begin{remark}
\label{obsolete-remark-induction-step-V-C-etale-h}
该标签原先指向 Cohomology on Sites，引理
\ref{sites-cohomology-lemma-induction-step-V-C-general} 在
\'Etale Cohomology，引理 \ref{etale-cohomology-lemma-compare-h-etale}
所述情形下的特殊情形。
\end{remark}

\begin{remark}
\label{obsolete-remark-how-used}
该标签原先位于 \'etale 上同调章节，但因重新组织已不适合留在那里。
标签内容如下：
\'Etale Cohomology，引理 \ref{etale-cohomology-lemma-when-ctf}
可用于证明：若 $f : X \to Y$ 是概形的分离有限型态射且 $Y$ 诺特，
则 $Rf_!$ 诱导函子
$D_{ctf}(X_\etale, \Lambda) \to D_{ctf}(Y_\etale, \Lambda)$。
当 $Y$ 是域的谱且 $X$ 是曲线时，这一论证的例子见 The Trace Formula，
命题 \ref{trace-proposition-projective-curve-constructible-cohomology}。
\end{remark}

\begin{lemma}
\label{obsolete-lemma-glue-f-upper-shriek}
设 $f : X \to Y$ 是概形的局部拟有限态射。
存在唯一函子
$f^! : \textit{Ab}(Y_\etale) \to \textit{Ab}(X_\etale)$，满足
\begin{enumerate}
\item 对任意开浸入 $j : U \to X$，若 $f \circ j$ 分离，则存在典范同构
$j^! \circ f^! = (f \circ j)^!$；并且
\item 当 $U \subset U' \subset X$ 时，这些同构与 More \'Etale Cohomology，引理
\ref{more-etale-lemma-upper-shriek-restriction} 中的同构相容。
\end{enumerate}
\end{lemma}

\begin{proof}
立即由 More \'Etale Cohomology，引理
\ref{more-etale-lemma-lqf-f-upper-shriek} 与
\ref{more-etale-lemma-upper-shriek-restriction} 得出。
\end{proof}

\begin{proposition}
\label{obsolete-proposition-lqf-shriek}
设 $f : X \to Y$ 是局部拟有限态射。存在伴随函子
$f_! : \textit{Ab}(X_\etale) \to \textit{Ab}(Y_\etale)$
与 $f^! : \textit{Ab}(Y_\etale) \to \textit{Ab}(X_\etale)$，
具有以下性质
\begin{enumerate}
\item 函子 $f^!$ 是 More \'Etale Cohomology 中由
引理 \ref{more-etale-lemma-lqf-f-upper-shriek} 构造的函子；
\item 对任意开浸入 $j : U \to X$，若 $f \circ j$ 分离，存在典范同构
$f_! \circ j_! = (f \circ j)_!$；并且
\item 当 $U \subset U' \subset X$ 时，这些同构与 More \'Etale Cohomology，
引理 \ref{more-etale-lemma-f-shriek-composition} 中的同构相容。
\end{enumerate}
\end{proposition}

\begin{proof}
见 More \'Etale Cohomology，小节
\ref{more-etale-section-finite-support} 与
\ref{more-etale-section-duality-locally-quasi-finite}。
\end{proof}

\begin{lemma}
\label{obsolete-lemma-lqf-f-upper-shriek-stalk}
设 $f : X \to Y$ 是局部拟有限的概形态射。
对阿贝尔群 $A$ 与几何点
$\overline{y} : \Spec(k) \to Y$，有
$f^!(\overline{y}_*A) = \prod\nolimits_{f(\overline{x}) = \overline{y}}
\overline{x}_*A$.
\end{lemma}

\begin{proof}
由 More \'Etale Cohomology，引理
\ref{more-etale-lemma-lqf-f-upper-shriek} 中对应的陈述推出。
\end{proof}

\begin{lemma}
\label{obsolete-lemma-lqf-f-shriek-composition}
设 $f : X \to Y$ 与 $g : Y \to Z$ 是可复合的局部拟有限概形态射。
则 $g_! \circ f_! = (g \circ f)_!$，并且
$f^! \circ g^! = (g \circ f)^!$。
\end{lemma}

\begin{proof}
这是 More \'Etale Cohomology 中引理
\ref{more-etale-lemma-lqf-shriek-composition} 与
\ref{more-etale-lemma-upper-shriek-restriction} 的组合。
\end{proof}









\section{微分分次代数}
\label{obsolete-section-dga}


\begin{lemma}
\label{obsolete-lemma-P-not-preserved-base-change}
设 $(A, \text{d})$ 与 $(B, \text{d})$ 是微分分次代数。
设 $N$ 是具有性质 (P) 的微分分次 $(A, B)$-双模，设 $M$ 是具有性质 (P)
的微分分次 $A$-模。则 $Q = M \otimes_A N$ 是微分分次 $B$-模，
在 $D(B)$ 中表示 $M \otimes_A^\mathbf{L} N$，并带有滤链
$$
0 = F_{-1}Q \subset F_0Q \subset F_1Q \subset \ldots \subset Q
$$
其微分分次子模满足 $Q = \bigcup F_pQ$；包含映射
$F_iQ \to F_{i + 1}Q$ 是可容许单射；商
$F_{i + 1}Q/F_iQ$ 作为微分分次 $B$-模同构于
$(A \otimes_R B)[k]$ 的直和。
\end{lemma}

\begin{proof}
在 $M$ 与 $N$ 上选取滤链 $F_\bullet$。考虑 $Q = M \otimes_A N$ 上由下式给出的滤链
$$
F_n(Q) = \sum\nolimits_{i + j = n}
F_i(M) \otimes_A F_j(N)
$$
这显然是微分分次 $B$-子模。我们有
$$
F_n(Q)/F_{n - 1}(Q) =
\bigoplus\nolimits_{i + j = n}
F_i(M)/F_{i - 1}(M) \otimes_A F_j(N)/F_{j - 1}(N)
$$
例如，这是因为 $M$ 的滤链在分次 $A$-模范畴中是可分裂的。由于假设右侧的商同构于
$A$ 与 $A \otimes_R B$ 的移位直和，并且
$A \otimes_A (A \otimes_R B) = A \otimes_R B$，
所以左侧作为微分分次 $B$-模是 $A \otimes_R B$ 的移位直和。
（警告：$Q$ 没有 $(A, B)$-双模结构。）
这证明了引理的第一个断言。第二个断言直接由
Differential Graded Algebra，引理 \ref{dga-lemma-derived-bc} 中函子的定义得到。
\end{proof}












\section{单纯方法}
\label{obsolete-section-simplicial}


\begin{lemma}
\label{obsolete-lemma-equiv}
假设与记号同 Simplicial，引理 \ref{simplicial-lemma-section}。
存在态射 $f$ 的截面 $g : U \to V$，且复合 $g \circ f$ 在 $V$ 上与恒等态射同伦等价。
特别地，态射 $f$ 是同伦等价。
\end{lemma}

\begin{proof}
立即由 Simplicial，引理 \ref{simplicial-lemma-section} 与
\ref{simplicial-lemma-trivial-kan-homotopy} 得出。
\end{proof}

\begin{lemma}
\label{obsolete-lemma-cosk-hom-deltak}
设 $\mathcal{C}$ 是具有有限余积与有限极限的范畴。设 $X$ 是 $\mathcal{C}$ 中的对象。
设 $k \geq 0$。规范映射
$$
\Hom(\Delta[k], X)
\longrightarrow
\text{cosk}_1 \text{sk}_1 \Hom(\Delta[k], X)
$$
是同构。
\end{lemma}

\begin{proof}
对任意单纯对象 $V$，有
\begin{eqnarray*}
\Mor(V, \text{cosk}_1 \text{sk}_1 \Hom(\Delta[k], X))
& = &
\Mor(\text{sk}_1 V, \text{sk}_1 \Hom(\Delta[k], X)) \\
& = &
\Mor(i_{1!} \text{sk}_1 V, \Hom(\Delta[k], X)) \\
& = &
\Mor(i_{1!} \text{sk}_1 V \times \Delta[k], X)
\end{eqnarray*}
第一个等式来自 $\text{sk}$ 与 $\text{cosk}$ 的伴随性，第二个等式来自
$i_{1!}$ 与 $\text{sk}_1$ 的伴随性，第三个等式来自
Simplicial，定义 \ref{simplicial-definition-hom-from-simplicial-set}，
其中最后的 $X$ 表示取值为 $X$ 的常值单纯对象。由 Simplicial，引理
\ref{simplicial-lemma-augmentation-howto}，该集合中的元素只依赖于
$i_{1!} \text{sk}_1 V \times \Delta[k]$ 的 $0$ 次与 $1$ 次项。
这些项与 $V \times \Delta[k]$ 的 $0$ 次与 $1$ 次项一致；见
Simplicial，引理 \ref{simplicial-lemma-recovering-U-for-real}。
因此上面的集合等于
$\Mor(V \times \Delta[k], X) = \Mor(V, \Hom(\Delta[k], X))$.
\end{proof}

\begin{lemma}
\label{obsolete-lemma-cosk0-hom-deltak}
设 $\mathcal{C}$ 是范畴。设 $X$ 是 $\mathcal{C}$ 中的对象，并且自积
$X \times \ldots \times X$ 存在。设 $k \geq 0$，令 $C[k]$ 如
Simplicial，例 \ref{simplicial-example-simplex-cosimplicial-set} 中所述。
采用 Simplicial，引理 \ref{simplicial-lemma-morphism-into-product} 的记号，
规范映射
$$
\Hom(C[k], X)_1
\longrightarrow
(\text{cosk}_0 \text{sk}_0 \Hom(C[k], X))_1
$$
可识别为映射
$$
\prod\nolimits_{\alpha : [k] \to [1]} X
\longrightarrow
X \times X
$$
它是投影到 $\alpha$ 为常值映射的那些因子的投影。
\end{lemma}

\begin{proof}
这在 Hypercoverings，引理 \ref{hypercovering-lemma-covering} 的证明中给出。
\end{proof}




\section{关于概形的结果}
\label{obsolete-section-devissage}

\noindent
看似多余的引理。

\begin{lemma}
\label{obsolete-lemma-stein-projective}
设 $(R, \mathfrak m, \kappa)$ 是局部环。设 $X \subset \mathbf{P}^n_R$ 是闭子概形。
假设 $R = \Gamma(X, \mathcal{O}_X)$。则特殊纤维 $X_k$ 几何连通。
\end{lemma}

\begin{proof}
这是 More on Morphisms，定理
\ref{more-morphisms-theorem-stein-factorization-general}.
\end{proof}

\begin{lemma}
\label{obsolete-lemma-property-irreducible-higher-rank}
设 $X$ 是诺特概形。设 $Z_0 \subset X$ 是不可约闭子集，其一般点为 $\xi$。
设 $\mathcal{P}$ 是 $X$ 上凝聚层的性质，满足
\begin{enumerate}
\item 对任意凝聚层的短正合列，若其中三个对象中的两个具有性质 $\mathcal{P}$，
则第三个也具有该性质。
\item 若凝聚层的直和具有 $\mathcal{P}$，则该性质对两个直项均成立。
\item 对每个积分闭子概形 $Z \subset Z_0 \subset X$，若 $Z \not = Z_0$，且每个拟凝聚理想层
$\mathcal{I} \subset \mathcal{O}_Z$，都有 $(Z \to X)_*\mathcal{I}$ 具有 $\mathcal{P}$。
\item 存在 $X$ 上的凝聚层 $\mathcal{G}$，使得
\begin{enumerate}
\item $\text{Supp}(\mathcal{G}) = Z_0$；
\item $\mathcal{G}_\xi$ 被 $\mathfrak m_\xi$ 湮灭；
\item $\mathcal{G}$ 具有性质 $\mathcal{P}$。
\end{enumerate}
\end{enumerate}
则 $X$ 上支集包含于 $Z_0$ 的每个凝聚层 $\mathcal{F}$ 都具有性质 $\mathcal{P}$。
\end{lemma}

\begin{proof}
证明是 Cohomology of Schemes，引理
\ref{coherent-lemma-property-irreducible} 证明的变体。
以完全相同的方式可见，支集严格包含于 $Z_0$ 的任何凝聚层都具有性质 $\mathcal{P}$。

\medskip\noindent
考虑如 (3) 中的凝聚层 $\mathcal{G}$。由 Cohomology of Schemes，引理
\ref{coherent-lemma-prepare-filter-irreducible}，在 $Z_0$ 上存在理想层 $\mathcal{I}$，以及短正合列
$$
0 \to
\left((Z_0 \to X)_*\mathcal{I}\right)^{\oplus r} \to
\mathcal{G} \to
\mathcal{Q} \to 0
$$
其中 $\mathcal{Q}$ 的支集严格包含于 $Z_0$。特别地，由于 $\mathcal{G}$ 的支集等于 $Z$，
有 $r > 0$ 且 $\mathcal{I}$ 非零。由于 $\mathcal{Q}$ 具有性质 $\mathcal{P}$，我们得到
$\left((Z_0 \to X)_*\mathcal{I}\right)^{\oplus r}$ 也具有性质 $\mathcal{P}$。
由 (2) 推出 $(Z_0 \to X)_*\mathcal{I}$ 具有性质 $\mathcal{P}$。
将此放入 Cohomology of Schemes，引理
\ref{coherent-lemma-property-irreducible} 证明中的适当位置，即得引理。
细节略。
\end{proof}

\begin{lemma}
\label{obsolete-lemma-property-higher-rank}
设 $X$ 是诺特概形。设 $\mathcal{P}$ 是 $X$ 上凝聚层的性质，满足
\begin{enumerate}
\item 对任意凝聚层的短正合列，若其中三个对象中的两个具有性质 $\mathcal{P}$，
则第三个也具有该性质。
\item 若凝聚层的直和具有 $\mathcal{P}$，则该性质对两个直项均成立。
\item 对每个一般点为 $\xi$ 的积分闭子概形 $Z \subset X$，存在凝聚层 $\mathcal{G}$，使得
\begin{enumerate}
\item $\text{Supp}(\mathcal{G}) = Z$；
\item $\mathcal{G}_\xi$ 被 $\mathfrak m_\xi$ 湮灭；
\item $\mathcal{G}$ 具有性质 $\mathcal{P}$。
\end{enumerate}
\end{enumerate}
则 $X$ 上每个凝聚层都具有性质 $\mathcal{P}$。
\end{lemma}

\begin{proof}
这由引理 \ref{obsolete-lemma-property-irreducible-higher-rank} 得出，方式与
Cohomology of Schemes，引理 \ref{coherent-lemma-property} 由
Cohomology of Schemes，引理 \ref{coherent-lemma-property-irreducible} 得出的方式完全相同。
\end{proof}

\begin{lemma}
\label{obsolete-lemma-section-maps-back-into}
设 $X$ 是概形。设 $\mathcal{L}$ 是可逆 $\mathcal{O}_X$-模。
设 $s \in \Gamma(X, \mathcal{L})$ 是截面。设 $\mathcal{F}' \subset \mathcal{F}$ 是拟凝聚
$\mathcal{O}_X$-模。假设
\begin{enumerate}
\item $X$ 拟紧；
\item $\mathcal{F}$ 是有限型；
\item $\mathcal{F}'|_{X_s} = \mathcal{F}|_{X_s}$。
\end{enumerate}
则存在 $n \geq 0$，使得 $\mathcal{F}$ 上乘以 $s^n$ 的映射经由 $\mathcal{F}'$ 分解。
\end{lemma}

\begin{proof}
换言之，我们要证明
$s^n\mathcal{F} \subset
\mathcal{F}' \otimes_{\mathcal{O}_X} \mathcal{L}^{\otimes n}$
对某个 $n \geq 0$ 成立。换言之，商映射
$\mathcal{F} \to \mathcal{F}/\mathcal{F}'$ 在乘以 $s$ 的某个幂后变为零。
这由 Properties，引理
\ref{properties-lemma-section-maps-backwards}.
\end{proof}

\begin{lemma}
\label{obsolete-lemma-bound-degree-in-nbhd-generic-point}
设 $f : X \to Y$ 是概形态射。假设
\begin{enumerate}
\item $X$ 与 $Y$ 是积分概形；
\item $f$ 局部有限型且占优；
\item $f$ 拟紧或分离；
\item $f$ 泛有限，即 Morphisms，引理 \ref{morphisms-lemma-finite-degree} 的
(1)--(5) 中至少有一个成立。
\end{enumerate}
则存在非空开集 $V \subset Y$，使得 $f^{-1}(V) \to V$ 是次数为 $\deg(X/Y)$ 的有限局部自由态射。
特别地，$f^{-1}(V) \to V$ 的纤维次数由 $\deg(X/Y)$ 控制。
\end{lemma}

\begin{proof}
可选取 $V$ 使 $f^{-1}(V) \to V$ 有限。再缩小 $V$，由一般平坦性可设
$f^{-1}(V) \to V$ 平坦且有限表示（Morphisms，命题
\ref{morphisms-proposition-generic-flatness}）。由 Morphisms，引理
\ref{morphisms-lemma-finite-flat}，该态射是有限局部自由的。
由于 $V$ 不可约，该态射具有固定次数。最后一个断言由此以及 Morphisms，引理
\ref{morphisms-lemma-finite-locally-free-universally-bounded} 得出。
\end{proof}







\section{簇的导出范畴}
\label{obsolete-section-equiv}

\noindent
以下若干引理原属簇的导出范畴一章，
但现已不再需要。

\begin{lemma}
\label{obsolete-lemma-preserves-Coh}
设 $k$ 是域。设 $X$ 是 $k$ 上有限型分离且正则的概形。设
$F : D_{perf}(\mathcal{O}_X) \to D_{perf}(\mathcal{O}_X)$ 是 $k$-线性正合函子。
假设对每个满足 $\dim(\text{Supp}(\mathcal{F})) = 0$ 的凝聚
$\mathcal{O}_X$-模 $\mathcal{F}$，存在 $k$-向量空间同构
$$
\Hom_X(\mathcal{F}, M) = \Hom_X(\mathcal{F}, F(M))
$$
且该同构在 $D_{perf}(\mathcal{O}_X)$ 中关于 $M$ 函子性。
则存在定义在 $k$ 上的自同构 $f : X \to X$，它在底层拓扑空间上诱导恒等映射\footnote{这通常迫使 $f$ 为恒等映射；见 Varieties，引理 \ref{varieties-lemma-automorphism}。}
，以及可逆 $\mathcal{O}_X$-模 $\mathcal{L}$，使得 $F$ 与
$F'(M) = f^*M \otimes_{\mathcal{O}_X}^\mathbf{L} \mathcal{L}$ 是兄弟函子。
\end{lemma}

\begin{proof}
由 Derived Categories of Varieties，引理 \ref{equiv-lemma-duality-at-point}，
对每个支集为闭点的凝聚 $\mathcal{O}_X$-模 $\mathcal{F}$，都有同构
$$
H^0(X, M \otimes^\mathbf{L}_{\mathcal{O}_X} \mathcal{F}) =
H^0(X, F(M) \otimes^\mathbf{L}_{\mathcal{O}_X} \mathcal{F})
$$
，且关于 $M$ 函子性。

\medskip\noindent
令 $x \in X$ 为闭点，并将上述结论应用于在 $x$ 处取值为 $\kappa(x)$ 的叠置层
$\mathcal{F} = \mathcal{O}_x$。得到
$$
\dim_{\kappa(x)} 
\text{Tor}^{\mathcal{O}_{X, x}}_p(M_x, \kappa(x)) =
\dim_{\kappa(x)} 
\text{Tor}^{\mathcal{O}_{X, x}}_p(F(M)_x, \kappa(x))
$$
对所有 $p \in \mathbf{Z}$ 成立。特别地，若 $H^i(M) = 0$ 对 $i > 0$ 成立，
则 $H^i(F(M)) = 0$ 对 $i > 0$，这由 Derived Categories of Varieties，引理
\ref{equiv-lemma-orthogonal-point-sheaf}.

\medskip\noindent
若 $\mathcal{E}$ 是秩为 $r$ 的局部自由模，则 $F(\mathcal{E})$ 也是秩为 $r$ 的局部自由模。
这是因为在 $\mathcal{O}_{X, x}$ 上的完美复形 $K$ 若满足
$$
\dim_{\kappa(x)} \text{Tor}^{\mathcal{O}_{X, x}}_i(K, \kappa(x)) =
\left\{
\begin{matrix}
r & \text{若} & i = 0 \\
0 & \text{若} & i \not = 0
\end{matrix}
\right.
$$
则等于置于次数 $0$ 的秩为 $r$ 的自由模。例见 More on Algebra，引理
\ref{more-algebra-lemma-lift-perfect-from-residue-field}.

\medskip\noindent
若 $M$ 的支集包含于闭子概形 $Z \subset X$，则 $F(M)$ 的支集也包含于 $Z$。
这是显然的，因为对 $x \not \in Z$ 有
$M \otimes_{\mathcal{O}_X}^\mathbf{L} \mathcal{O}_x = 0$，于是 $F(M)$ 也满足同样结论，
再由 Derived Categories of Varieties，引理
\ref{equiv-lemma-orthogonal-point-sheaf}.

\medskip\noindent
特别地，$F(\mathcal{O}_x)$ 的支集为 $\{x\}$。令 $i \in \mathbf{Z}$ 为满足
$H^i(\mathcal{O}_x) \not = 0$ 的最小整数。我们知道 $i \leq 0$。
若 $i < 0$，则存在态射
$\mathcal{O}_x[-i] \to F(\mathcal{O}_x)$
，这与所有态射 $\mathcal{O}_x[-i] \to \mathcal{O}_x$ 均为零相矛盾。
因此 $F(\mathcal{O}_x) = \mathcal{H}[0]$，其中 $\mathcal{H}$ 是集中在 $x$ 的叠置层。

\medskip\noindent
设 $\mathcal{G}$ 是满足 $\dim(\text{Supp}(\mathcal{G})) = 0$ 的凝聚 $\mathcal{O}_X$-模。
则存在滤链
$$
0 = \mathcal{G}_0 \subset \mathcal{G}_1 \subset \ldots \subset
\mathcal{G}_n = \mathcal{G}
$$
使得对 $n \geq i \geq 1$，商 $\mathcal{G}_i/\mathcal{G}_{i - 1}$ 同构于某个闭点
$x_i \in X$ 处的 $\mathcal{O}_{x_i}$。于是得到三角形
$$
F(\mathcal{G}_{i - 1}) \to F(\mathcal{G}_i) \to F(\mathcal{O}_{x_i})
$$
利用归纳法可知 $F(\mathcal{G}_i)$ 是置于次数 $0$ 的凝聚层。

\medskip\noindent
设 $\mathcal{G}$ 是凝聚 $\mathcal{O}_X$-模。我们知道对 $i > 0$ 有
$H^i(F(\mathcal{G})) = 0$。为导出矛盾，假设对某个 $i < 0$，
$H^i(F(\mathcal{G}))$ 非零。取具有此性质的最小 $i$，于是存在
$H^i(F(\mathcal{G}))[-i] \to F(\mathcal{G})$ 的 $D_{perf}(\mathcal{O}_X)$ 中态射。
取 $H^i(F(\mathcal{G}))$ 支集中的闭点 $x \in X$。由 More on Algebra，引理
\ref{more-algebra-lemma-kollar-kovacs-pseudo-coherent}
存在 $n > 0$，使得
$$
H^i(F(\mathcal{G}))_x \otimes_{\mathcal{O}_{X, x}}
\mathcal{O}_{X, x}/\mathfrak m_x^n
\longrightarrow
\text{Tor}^{\mathcal{O}_{X, x}}_{-i}(F(\mathcal{G})_x,
\mathcal{O}_{X, x}/\mathfrak m_x^n)
$$
非零。接着取 $m \geq 1$，考虑短正合列
$$
0 \to \mathfrak m_x^m \mathcal{G} \to \mathcal{G} \to
\mathcal{G}/\mathfrak m_x^m\mathcal{G} \to 0
$$
由上面的结论，$F(\mathcal{G}/\mathfrak m_x^m\mathcal{G})$ 是置于次数 $0$ 的层。
因此 $H^i(F(\mathfrak m_x^m \mathcal{G})) \to H^i(F(\mathcal{G}))$ 是同构。
考虑交换图表
$$
\xymatrix{
H^i(F(\mathfrak m_x^m\mathcal{G}))_x \otimes_{\mathcal{O}_{X, x}}
\mathcal{O}_{X, x}/\mathfrak m_x^n \ar[r] \ar[d] &
\text{Tor}^{\mathcal{O}_{X, x}}_{-i}(F(\mathfrak m_x^m\mathcal{G})_x,
\mathcal{O}_{X, x}/\mathfrak m_x^n) \ar[d] \\
H^i(F(\mathcal{G}))_x \otimes_{\mathcal{O}_{X, x}}
\mathcal{O}_{X, x}/\mathfrak m_x^n \ar[r] &
\text{Tor}^{\mathcal{O}_{X, x}}_{-i}(F(\mathcal{G})_x,
\mathcal{O}_{X, x}/\mathfrak m_x^n)
}
$$
由于左侧竖直箭头是同构且底部箭头非零，得到右侧竖直箭头对所有 $m \geq 1$ 均非零。
另一方面，由证明的第一段，该箭头同构于下列箭头
$$
\text{Tor}^{\mathcal{O}_{X, x}}_{-i}(\mathfrak m_x^m\mathcal{G}_x,
\mathcal{O}_{X, x}/\mathfrak m_x^n)
\longrightarrow
\text{Tor}^{\mathcal{O}_{X, x}}_{-i}(\mathcal{G}_x,
\mathcal{O}_{X, x}/\mathfrak m_x^n)
$$
然而由 More on Algebra，引理 \ref{more-algebra-lemma-tor-annihilated}，当 $m \gg n$ 时该箭头为零，
这就是所需的矛盾。

\medskip\noindent
因此 $F$ 保持凝聚模。由 Derived Categories of Varieties，引理
\ref{equiv-lemma-exact-functor-preserving-Coh}
，可知 $F$ 是 Fourier--Mukai 函子 $F'$ 的兄弟函子，其中该函子由凝聚
$\mathcal{O}_{X \times X}$-模 $\mathcal{K}$ 给出，该模经 $\text{pr}_1$ 在 $X$ 上平坦，
经 $\text{pr}_2$ 在 $X$ 上有限。由于 $F(\mathcal{O}_X)$ 是置于次数 $0$ 的可逆
$\mathcal{O}_X$-模 $\mathcal{L}$，我们得到
$$
\mathcal{L} \cong F(\mathcal{O}_X) \cong F'(\mathcal{O}_X) \cong
\text{pr}_{2, *}\mathcal{K}
$$
于是由 Functors and Morphisms，引理 \ref{functors-lemma-pushforward-invertible-pre}，
存在态射 $s : X \to X \times X$，满足 $\text{pr}_2 \circ s = \text{id}_X$ 且
$\mathcal{K} = s_*\mathcal{L}$。令 $f = \text{pr}_1 \circ s$。则有
\begin{align*}
F'(M)
& =
R\text{pr}_{2, *}(L\text{pr}_1^*K \otimes \mathcal{K}) \\
& =
R\text{pr}_{2, *}(L\text{pr}_1^*M \otimes s_*\mathcal{L}) \\
& =
R\text{pr}_{2, *}(Rs_*(Lf^*M \otimes \mathcal{L})) \\
& =
Lf^*M \otimes \mathcal{L}
\end{align*}
第三步使用了 Derived Categories of Schemes，引理
\ref{perfect-lemma-cohomology-base-change}。由于对所有闭点 $x \in X$，模
$F(\mathcal{O}_x)$ 的支集在 $x$，可知 $f$ 在 $X$ 的底层拓扑空间上诱导恒等映射。
仍需证明 $f$ 是同构，这将在下一段完成。

\medskip\noindent
令 $x \in X$ 为闭点。对 $n \geq 1$，记 $\mathcal{O}_{x, n}$ 为在 $x$ 处取值
$\mathcal{O}_{X, x}/\mathfrak m_x^n$ 的叠置层。我们有
$$
\Hom_X(\mathcal{O}_{x, m}, \mathcal{O}_{x, n}) \cong
\Hom_X(\mathcal{O}_{x, m}, F(\mathcal{O}_{x, n})) \cong
\Hom_X(\mathcal{O}_{x, m}, f^*\mathcal{O}_{x, n} \otimes \mathcal{L})
$$
这些同构关于各 $\mathcal{O}_{x, n}$ 之间的 $\mathcal{O}_X$-模同态函子性。
（第一个同构由假设给出，第二个同构因为 $F$ 与 $F'$ 是兄弟函子。）
当 $m \geq n$ 时，通过对 $\mathcal{O}_{x, n}$ 的作用，有
$\mathcal{O}_{X, x}/\mathfrak m^n = \Hom_X(\mathcal{O}_{x, m}, \mathcal{O}_{x, n})$。
因此 $f^\sharp : \mathcal{O}_{X, x}/\mathfrak m_x^n \to
\mathcal{O}_{X, x}/\mathfrak m_x^n$ 对所有 $n$ 均为双射。
所以 $f$ 在闭点处的完备局部环上诱导同构，因而是 \'etale
（\'Etale Morphisms，引理 \ref{etale-lemma-characterize-etale-completions}）。
考察闭点可知，$\Delta_f : X \to X \times_{f, X, f} X$（由于 $f$ 是 \'etale，这是开浸入）
是双射，因而是同构。因此 $f$ 是单态射。最后由 Descent，引理
\ref{descent-lemma-flat-surjective-quasi-compact-monomorphism-isomorphism}
知它是开浸入，从而 $f$ 是同构。
\end{proof}










\section{正则真情形下的可表性}
\label{obsolete-section-regular-proper}

\noindent
本节已废弃，因为我们改进了 Derived Categories of Varieties，
定理 \ref{equiv-theorem-bondal-van-den-bergh}，使其适用于域上的所有真概形
（此前我们只对域上的射影概形证明了该结论）。

\begin{lemma}
\label{obsolete-lemma-trace-map}
\begin{reference}
此处的证明沿用
\cite[Remark 3.4]{MS}
\end{reference}
设 $f : X' \to X$ 是积分诺特概形之间的真双有理态射，且 $X$ 正则。
映射 $\mathcal{O}_X \to Rf_*\mathcal{O}_{X'}$ 在 $D(\mathcal{O}_X)$ 中典范分裂。
\end{lemma}

\begin{proof}
在 $D(\mathcal{O}_X)$ 中置 $E = Rf_*\mathcal{O}_{X'}$。由 Derived Categories of Schemes，引理
\ref{perfect-lemma-direct-image-coherent}，有 $E$ 属于 $D^b_{\textit{Coh}}(\mathcal{O}_X)$。
由 Derived Categories of Schemes，引理 \ref{perfect-lemma-perfect-on-regular}，可知 $E$ 是
$D(\mathcal{O}_X)$ 中的完美对象。
由于 $\mathcal{O}_{X'}$ 是代数层，由 Cohomology，注记
\ref{cohomology-remark-cup-product} 有相对杯积
$\mu : E \otimes_{\mathcal{O}_X}^\mathbf{L} E \to E$。
令 $\sigma : E \otimes E^\vee \to E^\vee \otimes E$ 为对称幺半范畴
$D(\mathcal{O}_X)$ 中的交换约束（Cohomology，引理
\ref{cohomology-lemma-symmetric-monoidal-derived}）。记
$\eta : \mathcal{O}_X \to E \otimes E^\vee$ 与
$\epsilon : E^\vee \otimes E \to \mathcal{O}_X$ 为 Cohomology，例
\ref{cohomology-example-dual-derived} 中构造的映射。于是可考虑映射
$$
E \xrightarrow{\eta \otimes 1} E \otimes E^\vee \otimes E
\xrightarrow{\sigma \otimes 1} E^\vee \otimes E \otimes E
\xrightarrow{1 \otimes \mu} E^\vee \otimes E
\xrightarrow{\epsilon} \mathcal{O}_X
$$
我们断言该映射是引理陈述中映射的单侧逆。只需证明复合
$\mathcal{O}_X \to \mathcal{O}_X$ 是恒等映射即可。可以在 $X$ 的一般点处
（或在 $f$ 为同构的 $X$ 的某个开子概形上）进行验证。此时
$E = \mathcal{O}_X$，而 $\mu$ 是通常的乘法映射，结论显然。
\end{proof}

\begin{lemma}
\label{obsolete-lemma-characterize-dbcoh-proper-regular}
设 $X$ 是域 $k$ 上的正则真概形。设
$K \in \Ob(D_\QCoh(\mathcal{O}_X))$。以下条件等价：
\begin{enumerate}
\item $K \in D^b_{\textit{Coh}}(\mathcal{O}_X) = D_{perf}(\mathcal{O}_X)$；
\item 对 $D(\mathcal{O}_X)$ 中所有完美对象 $E$，有
$\sum_{i \in \mathbf{Z}} \dim_k \Ext^i_X(E, K) < \infty$。
\end{enumerate}
\end{lemma}

\begin{proof}
条件 (1) 中的等式由 Derived Categories of Schemes，引理
\ref{perfect-lemma-perfect-on-regular} 给出。(1) $\Rightarrow$ (2) 由
Derived Categories of Varieties，引理 \ref{equiv-lemma-finiteness} 得出。
(2) $\Rightarrow$ (1) 由 More on Morphisms，引理
\ref{more-morphisms-lemma-characterize-relatively-perfect}.
\end{proof}

\begin{lemma}
\label{obsolete-lemma-bondal-van-den-bergh}
设 $X$ 是域 $k$ 上的正则真概形。
\begin{enumerate}
\item 设 $F : D_{perf}(\mathcal{O}_X)^{opp} \to \text{Vect}_k$ 是 $k$-线性上同调函子，满足
$$
\sum\nolimits_{n \in \mathbf{Z}} \dim_k F(E[n]) < \infty
$$
对所有 $E \in D_{perf}(\mathcal{O}_X)$ 成立。则 $F$ 同构于形如
$E \mapsto \Hom_X(E, K)$ 的函子，其中 $K \in D_{perf}(\mathcal{O}_X)$。
\item 设 $G : D_{perf}(\mathcal{O}_X) \to \text{Vect}_k$ 是 $k$-线性同调函子，满足
$$
\sum\nolimits_{n \in \mathbf{Z}} \dim_k G(E[n]) < \infty
$$
对所有 $E \in D_{perf}(\mathcal{O}_X)$ 成立。则 $G$ 同构于形如
$E \mapsto \Hom_X(K, E)$ 的函子，其中 $K \in D_{perf}(\mathcal{O}_X)$。
\end{enumerate}
\end{lemma}

\begin{proof}
这由 Derived Categories of Varieties，定理
\ref{equiv-theorem-bondal-van-den-bergh} 与
引理 \ref{equiv-lemma-homological-representable} 得出。
下面也给出另一种证明。

\medskip\noindent
证明 (1)。导出范畴 $D_\QCoh(\mathcal{O}_X)$ 有直和，是紧生成的，且
$D_{perf}(\mathcal{O}_X)$ 是紧对象的全子范畴；见
Derived Categories of Schemes，引理
\ref{perfect-lemma-quasi-coherence-direct-sums},
定理 \ref{perfect-theorem-bondal-van-den-Bergh} 与
命题 \ref{perfect-proposition-compact-is-perfect}。
由 Derived Categories of Varieties，引理 \ref{equiv-lemma-van-den-bergh}，可假设
$F(E) = \Hom_X(E, K)$，其中 $K \in \Ob(D_\QCoh(\mathcal{O}_X))$。
于是由引理 \ref{obsolete-lemma-characterize-dbcoh-proper-regular}，$K$ 属于
$D^b_{\textit{Coh}}(\mathcal{O}_X)$。

\medskip\noindent
证明 (2)。考虑 $D_{perf}(\mathcal{O}_X)$ 上的逆变函子 $E \mapsto E^\vee$；见
Cohomology，引理 \ref{cohomology-lemma-dual-perfect-complex}。
该函子是 $D_{perf}(\mathcal{O}_X)$ 的正合反自等价。
因此可将 (1) 应用于函子 $F(E) = G(E^\vee)$，得到
$K \in D_{perf}(\mathcal{O}_X)$，满足 $G(E^\vee) = \Hom_X(E, K)$。
于是 $G(E) = \Hom_X(E^\vee, K) = \Hom_X(K^\vee, E)$，
故取 $K^\vee$ 即可。
\end{proof}









\section{商函子}
\label{obsolete-section-quotients}

\begin{lemma}
\label{obsolete-lemma-factors-through-quotient}
设 $S = \Spec(R)$ 是仿射概形。设 $X$ 是 $S$ 上的代数空间。
设 $q_i : \mathcal{F} \to \mathcal{Q}_i$（$i = 1, 2$）是拟凝聚
$\mathcal{O}_X$-模的满射。假设 $\mathcal{Q}_1$ 在 $S$ 上平坦。
设 $T \to S$ 是概形的拟紧态射，并且存在如下分解
$$
\xymatrix{
& \mathcal{F}_T \ar[rd]^{q_{2, T}} \ar[ld]_{q_{1, T}} \\
\mathcal{Q}_{1, T} & & \mathcal{Q}_{2, T} \ar@{..>}[ll]
}
$$
则存在闭子概形 $Z \subset S$，使得 (a) $T \to S$ 经由 $Z$ 分解，且 (b)
$q_{1, Z}$ 经由 $q_{2, Z}$ 分解。若 $\Ker(q_2)$ 是有限型
$\mathcal{O}_X$-模且 $X$ 拟紧，则可取 $Z \to S$ 为有限表示。
\end{lemma}

\begin{proof}
将 Flatness on Spaces，引理 \ref{spaces-flat-lemma-F-zero-somewhat-closed}
应用于映射 $\Ker(q_2) \to \mathcal{Q}_1$。
\end{proof}







\section{空间与 fpqc 覆盖}
\label{obsolete-section-fpqc}

\noindent
本节材料因 Gabber 的论证而废弃；该论证表明代数空间相对于 fpqc 覆盖满足层条件。
请参阅
Properties of Spaces，小节 \ref{spaces-properties-section-fpqc}。

\begin{lemma}
\label{obsolete-lemma-separated-fpqc}
设 $S$ 是概形。设 $X$ 是 $S$ 上的代数空间。设
$\{f_i : T_i \to T\}_{i \in I}$ 是 $S$ 上概形的 fpqc 覆盖。
则映射
$$
\Mor_S(T, X)
\longrightarrow
\prod\nolimits_{i \in I} \Mor_S(T_i, X)
$$
是单射。
\end{lemma}

\begin{proof}
立即由 Properties of Spaces，命题
\ref{spaces-properties-proposition-sheaf-fpqc}.
\end{proof}

\begin{lemma}
\label{obsolete-lemma-sheaf-fpqc-open-covering}
设 $S$ 是概形。设 $X$ 是 $S$ 上的代数空间。设
$X = \bigcup_{j \in J} X_j$ 是 Zariski 覆盖，见
Spaces，定义 \ref{spaces-definition-Zariski-open-covering}。
若每个 $X_j$ 对 fpqc 拓扑满足层性质，则 $X$ 对 fpqc 拓扑也满足层性质。
\end{lemma}

\begin{proof}
这是因为所有代数空间对 fpqc 拓扑都满足层性质；见
Properties of Spaces，命题
\ref{spaces-properties-proposition-sheaf-fpqc}.
\end{proof}

\begin{lemma}
\label{obsolete-lemma-sheaf-fpqc-quasi-separated}
设 $S$ 是概形。设 $X$ 是 $S$ 上的代数空间。若 $X$ 在 Zariski 局部上相对于 $S$
拟分离，则 $X$ 对 fpqc 拓扑满足层条件。
\end{lemma}

\begin{proof}
立即由一般的 Properties of Spaces，命题
\ref{spaces-properties-proposition-sheaf-fpqc}.
\end{proof}

\begin{remark}
\label{obsolete-remark-proof-works-when}
本注记原用于讨论引理 \ref{obsolete-lemma-sheaf-fpqc-quasi-separated}（2009 年 12 月 18 日）
的原证明在多大程度上可以推广。
\end{remark}






\section{非常良好的代数空间}
\label{obsolete-section-very-reasonable}

\noindent
略已废弃的材料。

\begin{lemma}
\label{obsolete-lemma-reasonable-kolmogorov}
设 $S$ 是概形。设 $X$ 是 $S$ 上的良好代数空间。
则 $|X|$ 是 Kolmogorov 空间（见 Topology，定义
\ref{topology-definition-generic-point}）。
\end{lemma}

\begin{proof}
由定义与 Decent Spaces，引理 \ref{decent-spaces-lemma-kolmogorov} 得出。
\end{proof}

\noindent
本节其余部分讨论非常良好的代数空间。若存在概形 $U$ 及满射、\'etale、拟紧
态射 $U \to X$，则 $X$ 非常良好；见 Decent Spaces，引理
\ref{decent-spaces-lemma-characterize-very-reasonable}。

\begin{lemma}
\label{obsolete-lemma-scheme-very-reasonable}
概形是非常良好的。
\end{lemma}

\begin{proof}
这是因为恒等映射是拟紧满射 \'etale 态射。
\end{proof}

\begin{lemma}
\label{obsolete-lemma-very-reasonable-Zariski-local}
设 $S$ 是概形。设 $X$ 是 $S$ 上的代数空间。
若存在 Zariski 开覆盖 $X = \bigcup X_i$，使每个 $X_i$ 都非常良好，
则 $X$ 非常良好。
\end{lemma}

\begin{proof}
这是 Decent Spaces，引理 \ref{decent-spaces-lemma-properties-local} 的情形 $(\epsilon)$。
\end{proof}

\begin{lemma}
\label{obsolete-lemma-quasi-separated-very-reasonable}
Zariski 局部拟分离的代数空间非常良好。特别地，任何拟分离代数空间都非常良好。
\end{lemma}

\begin{proof}
这是 Decent Spaces，引理 \ref{decent-spaces-lemma-bounded-fibres} 的一个推论。
\end{proof}

\begin{lemma}
\label{obsolete-lemma-representable-very-reasonable}
设 $S$ 是概形。设 $X$、$Y$ 是 $S$ 上的代数空间。
设 $Y \to X$ 是可表态射。若 $X$ 非常良好，则 $Y$ 也非常良好。
\end{lemma}

\begin{proof}
这是 Decent Spaces，引理 \ref{decent-spaces-lemma-representable-properties} 的情形 $(\epsilon)$。
\end{proof}

\begin{remark}
\label{obsolete-remark-very-reasonable-Zariski-locally-quasi-separated}
非常良好的代数空间构成一个严格大于 Zariski 局部拟分离代数空间的集合。
考虑形如 $X = [U/G]$ 的代数空间（见 Spaces，定义
\ref{spaces-definition-quotient}），其中 $G$ 是有限群，在非拟分离概形 $U$ 上无不动点地作用。
在此情形，$U \times_X U = U \times G$，且显然到 $U$ 的两个投影都是拟紧的，
故 $X$ 非常良好。另一方面，对角映射
$U \times_X U \to U \times U$ 不是拟紧的，因此该代数空间不是拟分离的。
现在取特征为 $\not = 2$ 的域 $k$ 上的无限仿射空间 $U$，并将零点加倍；见
Schemes，例 \ref{schemes-example-not-quasi-separated}。
令 $0_1, 0_2$ 为 $U$ 的两个零点。令 $G = \{+1, -1\}$，令 $-1$ 在所有坐标上按
$-1$ 作用，并交换 $0_1$ 与 $0_2$。则 $[U/G]$ 非常良好，但不是 Zariski 局部拟分离的
（细节略）。
\end{remark}

\noindent
警告：使用以下引理时须谨慎，因为其中的概形 $U_i$ 不一定分离，甚至不一定拟分离。

\begin{lemma}
\label{obsolete-lemma-very-reasonable-quasi-compact-pieces}
设 $S$ 是概形。设 $X$ 是 $S$ 上的非常良好代数空间。
存在一族概形 $U_i$ 及态射 $U_i \to X$，使得
\begin{enumerate}
\item 每个 $U_i$ 是拟紧概形；
\item 每个 $U_i \to X$ 是 \'etale；
\item 两个投影 $U_i \times_X U_i \to U_i$ 都是拟紧的；
\item 态射 $\coprod U_i \to X$ 是满射（且为 \'etale）。
\end{enumerate}
\end{lemma}

\begin{proof}
Decent Spaces，定义 \ref{decent-spaces-definition-very-reasonable} 表明存在满足 (2)、(3)、(4) 的
$U_i \to X$。固定 $i$，置 $R_i = U_i \times_X U_i$，记投影为
$s, t : R_i \to U_i$。对任意仿射开集 $W \subset U_i$，开集
$W' = t(s^{-1}(W)) \subset U_i$ 是拟紧的 $R_i$-不变开集（见
Groupoids，引理 \ref{groupoids-lemma-constructing-invariant-opens}）。
因此 $W'$ 是拟紧概形，$W' \to X$ 是 \'etale，且
$W' \times_X W' = s^{-1}(W') = t^{-1}(W')$，所以两个投影
$W' \times_X W' \to W'$ 都是拟紧的。令 $W \subset U_i$ 遍历 $U_i$ 的仿射开覆盖中的成员，
所得族 $W' \to X$ 即满足要求。
\end{proof}




\section{关于代数空间的已废弃引理}
\label{obsolete-section-obsolete-on-spaces}

\noindent
看似多余或不再用于正文的引理。

\begin{lemma}
\label{obsolete-lemma-vanishing-surjective}
在 Cohomology of Spaces，情形 \ref{spaces-cohomology-situation-vanishing} 中，
态射 $p : X \to \Spec(A)$ 是满射。
\end{lemma}

\begin{proof}
本引理原用于 Cohomology of Spaces，命题
\ref{spaces-cohomology-proposition-vanishing-affine} 的证明，
但现在它是该命题的推论。
\end{proof}

\begin{lemma}
\label{obsolete-lemma-vanishing-universally-closed}
在 Cohomology of Spaces，情形 \ref{spaces-cohomology-situation-vanishing} 中，
态射 $p : X \to \Spec(A)$ 是泛闭的。
\end{lemma}

\begin{proof}
本引理原用于 Cohomology of Spaces，命题
\ref{spaces-cohomology-proposition-vanishing-affine} 的证明，
但现在它是该命题的推论。
\end{proof}

\begin{remark}
\label{obsolete-remark-equation-first}
此标签原指向 Formal Spaces，引理 \ref{formal-spaces-lemma-iff-adic-star} 证明中的一个方程。
\end{remark}

\begin{remark}
\label{obsolete-remark-equation-second}
此标签原指向 Formal Spaces，引理 \ref{formal-spaces-lemma-iff-adic-star} 证明中的一个方程。
\end{remark}





\section{关于代数叠的已废弃引理}
\label{obsolete-section-obsolete-on-stacks}

\noindent
看似多余或不再用于正文的引理。

\begin{lemma}
\label{obsolete-lemma-infinite-sequence}
设 $S$ 是局部诺特概形。设 $\mathcal{X}$ 是在 $(\Sch/S)_{fppf}$ 上纤维化于群胚的范畴，
具有 (RS*)。设 $x$ 是 $\mathcal{X}$ 中位于 $S$ 上有限型仿射概形 $U$ 上的对象。
设 $u_n \in U$（$n \geq 1$）是两两不同的有限型点，并且对所有 $n$，$x$ 在 $u_n$ 处都不是 versal 的。
将 $u_n$ 换为子列后，存在态射
$$
x \to x_1 \to x_2 \to \ldots
\quad\text{于 }\mathcal{X}\text{ 位于下列对象之上： }\quad
U \to U_1 \to U_2 \to \ldots
$$
它们定义在 $S$ 上，并满足
\begin{enumerate}
\item 对每个 $n$，态射 $U \to U_n$ 是一阶加厚；
\item 对每个 $n$，有短正合列
$$
0 \to \kappa(u_n) \to \mathcal{O}_{U_n} \to \mathcal{O}_{U_{n - 1}} \to 0
$$
其中 $n = 1$ 时 $U_0 = U$；
\item 对每个 $n$，{\bf 不}存在由 $u_n$ 在其中的开邻域 $W \subset U_n$ 与态射
$\alpha : x_n|_W \to x$ 构成的对 $(W, \alpha)$，使得复合
$$
x|_{U \cap W} \xrightarrow{\text{下列对象的限制： }x \to x_n}
x_n|_W \xrightarrow{\alpha} x
$$
是典范态射 $x|_{U \cap W} \to x$。
\end{enumerate}
\end{lemma}

\begin{proof}
本引理原用于 versal 性开放性判据的证明
（Artin's Axioms，引理 \ref{artin-lemma-SGE-implies-openness-versality}），但已被
Artin's Axioms，引理 \ref{artin-lemma-infinite-sequence-pre} 取代，后者可直接推出本引理。
具体地，将 $u_n$（$n \geq 1$）换为子列后，可假设这些点之间没有专化关系；见
Properties，引理 \ref{properties-lemma-thin-infinite-sequence}。然后应用
Artin's Axioms，引理 \ref{artin-lemma-infinite-sequence-pre}，即可完成证明。
\end{proof}







\section{概形余切复形的变体}
\label{obsolete-section-cotangent-schemes-variant}

\noindent
本节给出概形态射余切复形的另一种构造。目前本节位于已废弃章节中，
因为在应用中可以使用 Cotangent，小节
\ref{cotangent-section-cotangent-schemes-variant} 所讨论的较容易版本。

\medskip\noindent
设 $f : X \to Y$ 是概形态射。令 $\mathcal{C}_{X/Y}$ 为如下交换图表为对象的范畴
\begin{equation}
\label{obsolete-equation-object}
\vcenter{
\xymatrix{
X \ar[d] & U \ar[l] \ar[d] \ar[r]_i & A \ar[ld] \\
Y & V \ar[l]
}
}
\end{equation}
，其中各项为概形且
\begin{enumerate}
\item $U$ 是 $X$ 的开子概形；
\item $V$ 是 $Y$ 的开子概形；
\item 在 $V$ 上存在同构 $A = V \times \Spec(P)$，其中 $P$ 是
  $\mathbf{Z}$ 上的多项式代数（变量取自某个集合）。
\end{enumerate}
换言之，$A$ 是 $V$ 上的（无限维）仿射空间。态射由交换图表给出。

\medskip\noindent
{\bf 记号。} $\mathcal{C}_{X/Y}$ 的对象，即图表
(\ref{obsolete-equation-object})，常记为 $U \to A$，其中约定：(a) $U$ 是 $X$ 的开子概形；(b)
$U \to A$ 是在 $Y$ 上的态射；(c) 结构态射 $A \to Y$ 的像是开集 $V \subset Y$；(d)
$A \to V$ 是仿射空间。我们写 $U \to A/V$ 表示 $V \subset Y$ 是 $A \to Y$ 的像。
记 $X_{Zar}$ 为 $X$ 的小 Zariski site。存在遗忘函子
$$
\mathcal{C}_{X/Y} \to X_{Zar},\ (U \to A) \mapsto U
\quad\text{且}\quad
\mathcal{C}_{X/Y} \mapsto Y_{Zar},\ (U \to A/V) \mapsto V.
$$

\begin{lemma}
\label{obsolete-lemma-category-fibred}
设 $X \to Y$ 是概形态射。
\begin{enumerate}
\item 范畴 $\mathcal{C}_{X/Y}$ 在 $X_{Zar}$ 上纤维化；
\item 范畴 $\mathcal{C}_{X/Y}$ 在 $Y_{Zar}$ 上纤维化；
\item 范畴 $\mathcal{C}_{X/Y}$ 在如下二元组 $(U, V)$ 的范畴上纤维化：
  $U \subset X$、$V \subset Y$ 为开集且 $f(U) \subset V$。
\end{enumerate}
\end{lemma}

\begin{proof}
对 (1)。给定 $\mathcal{C}_{X/Y}$ 的对象 $U \to A$ 及 $X_{Zar}$ 中的态射
$U' \to U$，考虑 $\mathcal{C}_{X/Y}$ 的对象 $i' : U' \to A$，其中 $i'$ 是
$i$ 与 $U' \to U$ 的复合。$\mathcal{C}_{X/Y}$ 中的态射
$(U' \to A) \to (U \to A)$ 在 $X_{Zar}$ 上强笛卡尔。

\medskip\noindent
对 (2)。给定对象 $U \to A/V$ 与 $V' \to V$，置
$U' = U \cap f^{-1}(V')$、$A' = V' \times_V A$，即可得到在 $V' \to V$ 上强笛卡尔的
态射 $(U' \to A') \to (U \to A)$。

\medskip\noindent
对 (3)。记 (3) 中的范畴为 $(X/Y)_{Zar}$。给定 $U \to A/V$ 及
$(X/Y)_{Zar}$ 中的态射 $(U', V') \to (U, V)$，考虑
$A' = V' \times_V A$。则态射 $(U' \to A'/V') \to (U \to A/V)$
在 $\mathcal{C}_{X/Y}$ 中关于 $(X/Y)_{Zar}$ 强笛卡尔。
\end{proof}

\noindent
使用从 $X_{Zar}$ 继承的拓扑（见 Stacks，小节 \ref{stacks-section-topology}），
在 $\mathcal{C}_{X/Y}$ 上得到拓扑 $\tau_X$。除非另有说明，我们考虑的都是
$\mathcal{C}_{X/Y}$ 上的这个拓扑。准确地说，一族态射
$\{(U_i \to A_i) \to (U \to A)\}$ 是 $\mathcal{C}_{X/Y}$ 的覆盖，当且仅当
\begin{enumerate}
\item $U = \bigcup U_i$；
\item 对所有 $i$，$A_i = A$。
\end{enumerate}
若在 (2) 中允许 $A_i \cong A$，得到相同的层集合。函子 $u$ 定义拓扑斯态射
$\pi : \Sh(\mathcal{C}_{X/Y}) \to \Sh(X_{Zar})$.

\medskip\noindent
site $\mathcal{C}_{X/Y}$ 带有若干环层。
\begin{enumerate}
\item 由规则
$(U \to A) \mapsto \mathcal{O}(A)$.
给出的层 $\mathcal{O}$；
\item 由规则 $(U \to A) \mapsto \mathcal{O}(U)$ 给出的层
$\underline{\mathcal{O}}_X = \pi^{-1}\mathcal{O}_X$；
\item 由规则
$(U \to A/V) \mapsto \mathcal{O}(V)$.
给出的层 $\underline{\mathcal{O}}_Y$。
\end{enumerate}
得到环化拓扑斯的态射
\begin{equation}
\label{obsolete-equation-pi-schemes}
\vcenter{
\xymatrix{
(\Sh(\mathcal{C}_{X/Y}), \underline{\mathcal{O}}_X) \ar[r]_i \ar[d]_\pi &
(\Sh(\mathcal{C}_{X/Y}), \mathcal{O}) \\
(\Sh(X_{Zar}), \mathcal{O}_X)
}
}
\end{equation}
态射 $i$ 在底层拓扑斯上是恒等态射，且
$i^\sharp : \mathcal{O} \to \underline{\mathcal{O}}_X$ 是显然的映射。
映射 $\pi$ 是 Cohomology on Sites，情形
\ref{sites-cohomology-situation-fibred-category} 的特殊情形。
下列导出函子将在后文发挥重要作用：
$
Li^* : D(\mathcal{O}) \longrightarrow D(\underline{\mathcal{O}}_X)
$
是 $Ri_* = i_* : D(\underline{\mathcal{O}}_X) \to D(\mathcal{O})$ 的左伴随，以及
$
L\pi_! : D(\underline{\mathcal{O}}_X) \longrightarrow D(\mathcal{O}_X)
$
是下列函子的左伴随：
$\pi^* = \pi^{-1} : D(\mathcal{O}_X) \to D(\underline{\mathcal{O}}_X)$.

\begin{remark}
\label{obsolete-remark-different-topologies}
取从 $Y_{Zar}$ 继承的拓扑，可在 $\mathcal{C}_{X/Y}$ 上得到第二个拓扑 $\tau_Y$。
还有第三个拓扑 $\tau_{X \to Y}$：一族态射
$\{(U_i \to A_i) \to (U \to A)\}$ 是覆盖，当且仅当
$U = \bigcup U_i$、$V = \bigcup V_i$ 且 $A_i \cong V_i \times_V A$。
这是从 site $(X/Y)_{Zar}$ 上的拓扑继承而来的；其底层范畴是引理
\ref{obsolete-lemma-category-fibred} 第 (3) 部分中的二元组 $(U, V)$ 的范畴。
$(X/Y)_{Zar}$ 的覆盖是满足 $U = \bigcup U_i$ 与 $V = \bigcup V_i$ 的族
$\{(U_i, V_i) \to (U, V)\}$。存在拓扑斯态射
$$
\xymatrix{
\Sh(\mathcal{C}_{X/Y})
= \Sh(\mathcal{C}_{X/Y}, \tau_X) &
\Sh(\mathcal{C}_{X/Y}, \tau_{X \to Y}) \ar[l] \ar[r] &
\Sh(\mathcal{C}_{X/Y}, \tau_Y)
}
$$
（回忆 $\tau_X$ 是我们的“默认”拓扑）。这些箭头的拉回函子是层化，而推出在底层预层上是恒等函子。
拓扑斯图表
$$
\xymatrix{
\Sh(X_{Zar}) \ar[d]^f & \Sh(\mathcal{C}_{X/Y}) \ar[l]^\pi &
\Sh(\mathcal{C}_{X/Y}, \tau_{X \to Y}) \ar[l] \ar[d] \\
\Sh(Y_{Zar}) & & \Sh(\mathcal{C}_{X/Y}, \tau_Y) \ar[ll]
}
$$
{\bf 不}交换。事实上，$Y$ 上非零阿贝尔层的拉回是
$(\mathcal{C}_{X/Y}, \tau_{X \to Y})$ 上的非零阿贝尔层，
但确实可以找到该层在 $X$ 上拉回为零的例子。注意，$Y_{Zar}$ 上任意预层
$\mathcal{F}$ 都通过把集合 $\mathcal{F}(V)$ 赋给 $(U \to A/V)$ 的规则，
给出 $\mathcal{C}_{Y/X}$ 上的层 $\underline{\mathcal{F}}$。
即使 $\mathcal{F}$ 恰为层，一般也不成立
$\underline{\mathcal{F}} = \pi^{-1}f^{-1}\mathcal{F}$。这与上图不交换有关：
经下方水平箭头与右侧竖直箭头，可把 $\underline{\mathcal{F}}$ 描述为
$\mathcal{F}$ 到 $\Sh(\mathcal{C}_{X/Y}, \tau_{X \to Y})$ 的拉回之推出。
层 $\underline{\mathcal{O}}_Y$ 即为一例。
但确实存在映射 $\underline{\mathcal{F}} \to \pi^{-1}f^{-1}\mathcal{F}$，
它在应用 $\pi_!$ 后（如后文所见）变成同构。
\end{remark}










\section{平坦模的形变与障碍}
\label{obsolete-section-modules}

\noindent
本节概述环 $\Lambda$ 上任意代数空间 $X$ 的凝聚层叠之形变理论构造。
由于 Quot，小节 \ref{quot-section-not-flat} 中已有改进的讨论，本材料已废弃。

\medskip\noindent
设定如下。假设给定
\begin{enumerate}
\item 环 $\Lambda$；
\item $\Lambda$ 上的代数空间 $X$；
\item $\Lambda$-代数 $A$，置 $X_A = X \times_{\Spec(\Lambda)} \Spec(A)$；
\item 在 $A$ 上平坦的有限表示 $\mathcal{O}_{X_A}$-模 $\mathcal{F}$。
\end{enumerate}
在这种情形下，考虑所有可能的满射
$$
0 \to I \to A' \to A \to 0
$$
其中 $A'$ 是 $\Lambda$-代数，其核 $I$ 是 $A'$ 中平方为零的理想。
给定 $A'$，得到 $\Spec(\Lambda)$ 上代数空间的一阶加厚 $X_A \to X_{A'}$。
对每个这样的加厚，考虑把 $\mathcal{F}$ 提升为 $X_{A'}$ 上在 $A'$ 上平坦的有限表示模
$\mathcal{F}'$ 的问题。我们希望在此设定下复现 Deformation Theory，引理
\ref{defos-lemma-flat-ringed-topoi} 的结果。

\medskip\noindent
更准确地，令 $\textit{Lift}(\mathcal{F}, A')$ 表示如下二元组
$(\mathcal{F}', \alpha)$ 的范畴：$\mathcal{F}'$ 是 $X_{A'}$ 上在 $A'$ 上平坦的有限表示模，
而 $\alpha : \mathcal{F}'|_{X_A} \to \mathcal{F}$ 是同构。
态射 $(\mathcal{F}'_1, \alpha_1) \to (\mathcal{F}'_2, \alpha_2)$ 是与
$\alpha_1$、$\alpha_2$ 相容的同构 $\mathcal{F}'_1 \to \mathcal{F}'_2$。
$\textit{Lift}(\mathcal{F}, A')$ 的同构类集合记为 $\text{Lift}(\mathcal{F}, A')$。

\medskip\noindent
设 $\mathcal{G}$ 是 $X_\etale$ 上在 $A$ 上平坦的
$\mathcal{O}_X \otimes_\Lambda A$-模层。定义二元组
$(\mathcal{G}', \beta)$ 的范畴 $\textit{Lift}(\mathcal{G}, A')$，其中 $\mathcal{G}'$ 是
在 $A'$ 上平坦的 $\mathcal{O}_X \otimes_\Lambda A'$-模层，而 $\beta$ 是同构
$\mathcal{G}' \otimes_{A'} A \to \mathcal{G}$。

\begin{lemma}
\label{obsolete-lemma-equivalence}
记号与假设如上。记投影为 $p : X_A \to X$。给定 $A'$，记投影为
$p' : X_{A'} \to X$。函子 $p'_*$ 在下列范畴之间诱导等价：
\begin{enumerate}
\item 范畴 $\textit{Lift}(\mathcal{F}, A')$；
\item 范畴 $\textit{Lift}(p_*\mathcal{F}, A')$。
\end{enumerate}
\end{lemma}

\begin{proof}
FIXME.
\end{proof}

\noindent
设 $\mathcal{H}$ 是 $\mathcal{C}_{X/\Lambda}$ 上在 $A$ 上平坦的
$\mathcal{O} \otimes_\Lambda A$-模层。定义范畴
$\textit{Lift}_\mathcal{O}(\mathcal{H}, A')$，其对象是二元组
$(\mathcal{H}', \gamma)$，其中 $\mathcal{H}'$ 是在 $A'$ 上平坦的
$\mathcal{O} \otimes_\Lambda A'$-模层，而
$\gamma : \mathcal{H}' \otimes_A A' \to \mathcal{H}$ 是
$\mathcal{O} \otimes_\Lambda A$-模的同构。

\medskip\noindent
设 $\mathcal{G}$ 是 $X_\etale$ 上在 $A$ 上平坦的
$\mathcal{O}_X \otimes_\Lambda A$-模层。考虑 Cotangent，方程
(\ref{cotangent-equation-pi-spaces}) 中的态射 $i$ 与 $\pi$。
记 $\underline{\mathcal{G}} = \pi^{-1}(\mathcal{G})$。它由规则
$(U \to \mathbf{A}) \mapsto \mathcal{G}(U)$ 直接给出，因而是
$\underline{\mathcal{O}}_X \otimes_\Lambda A$-模层。
以 $i_*\underline{\mathcal{G}}$ 表示同一个层，但将其视为
$\mathcal{O} \otimes_\Lambda A$-模层。

\begin{lemma}
\label{obsolete-lemma-second-equivalence}
记号与假设同上。函子 $\pi_!$ 在下列两个范畴之间诱导范畴等价：
\begin{enumerate}
\item 范畴 $\textit{Lift}_\mathcal{O}(i_*\underline{\mathcal{G}}, A')$；
\item 范畴 $\textit{Lift}(\mathcal{G}, A')$。
\end{enumerate}
\end{lemma}

\begin{proof}
FIXME.
\end{proof}

\begin{lemma}
\label{obsolete-lemma-second-equivalence-obs}
记号与假设同引理 \ref{obsolete-lemma-second-equivalence}。考虑
$$
L = L(\Lambda, X, A, \mathcal{G}) = L\pi_!(Li^*(i_*(\underline{\mathcal{G}})))
$$
这一 $D(\mathcal{O}_X \otimes_\Lambda A)$ 中的对象。给定核为平方零理想
$I$ 的 $\Lambda$-代数满射 $A' \to A$，则有：
\begin{enumerate}
\item 范畴 $\textit{Lift}(\mathcal{G}, A')$ 非空，当且仅当某个类
$\xi \in \Ext^2_{\mathcal{O}_X \otimes A}(L, \mathcal{G} \otimes_A I)$
为零。
\item 若 $\textit{Lift}(\mathcal{G}, A')$ 非空，则
$\text{Lift}(\mathcal{G}, A')$ 是
$\Ext^1_{\mathcal{O}_X \otimes A}(L, \mathcal{G} \otimes_A I)$
作用下的主齐性集。
\item 给定一个提升 $\mathcal{G}'$，其拉回为
$\text{id}_\mathcal{G}$ 的 $\mathcal{G}'$ 自同构所成集合典范同构于
$\Ext^0_{\mathcal{O}_X \otimes A}(L, \mathcal{G} \otimes_A I)$.
\end{enumerate}
\end{lemma}

\begin{proof}
FIXME.
\end{proof}

\noindent
最后，将以上各项汇总如下。

\begin{proposition}
\label{obsolete-proposition-conclusion}
设 $\Lambda$、$X$、$A$、$\mathcal{F}$ 如上。存在 $D(X_A)$ 中的典范对象
$L = L(\Lambda, X, A, \mathcal{F})$，使得对任意核为平方零理想 $I$ 的
$\Lambda$-代数满射 $A' \to A$，均有：
\begin{enumerate}
\item 范畴 $\textit{Lift}(\mathcal{F}, A')$ 非空，当且仅当某个类
$\xi \in \Ext^2_{X_A}(L, \mathcal{F} \otimes_A I)$ 为零。
\item 若 $\textit{Lift}(\mathcal{F}, A')$ 非空，则
$\text{Lift}(\mathcal{F}, A')$ 是
$\Ext^1_{X_A}(L, \mathcal{F} \otimes_A I)$ 作用下的主齐性集。
\item 给定一个提升 $\mathcal{F}'$，其拉回为
$\text{id}_\mathcal{F}$ 的 $\mathcal{F}'$ 自同构所成集合典范同构于
$\Ext^0_{X_A}(L, \mathcal{F} \otimes_A I)$.
\end{enumerate}
\end{proposition}

\begin{proof}
FIXME.
\end{proof}

\begin{lemma}
\label{obsolete-lemma-pseudo-coherent}
在命题 \ref{obsolete-proposition-conclusion} 的情形中，若
$X \to \Spec(\Lambda)$ 局部有限型且 $\Lambda$ 为诺特环，则 $L$ 伪凝聚。
\end{lemma}

\begin{proof}
FIXME.
\end{proof}






\section{非平坦情形下的凝聚层叠}
\label{obsolete-section-not-flat}

\noindent
Quot 的定理 \ref{quot-theorem-coherent-algebraic} 实际上不需要
$f : X \to B$ 平坦这一假设。本节依据导出代数几何中的思想修改证明方法，
以绕开平坦性假设。Quot 的小节 \ref{quot-section-not-flat}
采用一种完全不同的方法得到完全相同的结果；因此本节的方法已经废弃。

\medskip\noindent
Quot 的定理 \ref{quot-theorem-coherent-algebraic} 的证明中，
唯一使用平坦性的步骤是应用 Quot 的引理
\ref{quot-lemma-coherent-defo-thy}。该引理用于构造 Artin's Axioms
小节 \ref{artin-section-dual} 中的那种障碍理论。其证明依赖
Deformation Theory 的引理 \ref{defos-lemma-flat-ringed-topoi} 与
\ref{defos-lemma-verify-iv-ringed-topoi}；这两个引理见 Deformation Theory
小节 \ref{defos-section-flat-ringed-topoi}。$f$ 的平坦性假设正由此产生。
继续之前请注意，即使不假设 $f$ 平坦，Deformation Theory 的引理
\ref{defos-lemma-flat-ringed-topoi} 中结论 (2) 与 (3) 仍然成立，
因为它们依赖的 Deformation Theory 引理
\ref{defos-lemma-inf-ext-rel-ringed-topoi} 与
\ref{defos-lemma-inf-map-rel-ringed-topoi} 均不含任何平坦性假设。

\medskip\noindent
在给出细节之前，我们先从导出代数几何的角度说明这一构造的动机，
因为这有助于澄清下文。设 $A$ 是局部诺特基底 $S$ 上的有限型代数。
以 $X \otimes^\mathbf{L} A$ 表示 $X$ 到 $A$ 的“导出基变换”，并以
$i : X_A \to X \otimes^\mathbf{L} A$ 表示典范嵌入态射。
对象 $X \otimes^\mathbf{L} A$ 在 Stacks 项目中（尚）无定义；
可以把它想成配备单纯环层
$\mathcal{O}_{X \otimes^\mathbf{L} A}$ 的代数空间 $X_A$，其同调层为
$$
H_i(\mathcal{O}_{X \otimes^\mathbf{L} A}) =
\text{Tor}^{\mathcal{O}_S}_i(\mathcal{O}_X, A).
$$
态射 $X \otimes^\mathbf{L} A \to \Spec(A)$ 是平坦的
（该单纯环层的各项均在 $A$ 上平坦），故平坦模形变的通常理论适用于它。
由此可见，我们用下列群得到一个障碍理论：
$$
\Ext^i_{X \otimes^\mathbf{L} A}(i_*\mathcal{F},
i_*\mathcal{F} \otimes_A M)
$$
其中 $i = 0, 1, 2$ 分别对应无穷小自同构、无穷小形变和障碍。
注意，$i_*\mathcal{F}$ 到 $X \otimes^\mathbf{L} A'$ 的平坦形变自动具有
$i'_*\mathcal{F}'$ 的形式，其中 $\mathcal{F}'$ 是 $\mathcal{F}$ 的平坦形变。
由函子 $Li^*$ 与 $i_* = Ri_*$ 的伴随性，这些 Ext 群等于
$$
\Ext^i_{X_A}(Li^*(i_*\mathcal{F}), \mathcal{F} \otimes_A M)
$$
因此所得障碍群与 Quot 的引理 \ref{quot-lemma-coherent-defo-thy}
证明中的形式完全相同；唯一的变化，是将第一次出现的 $\mathcal{F}$
替换为复形 $Li^*(i_*\mathcal{F})$。

\medskip\noindent
下文将“直接”构造 $E(\mathcal{F}) = Li^*(i_*\mathcal{F})$，并直接证明它与
$\mathcal{F}$ 的形变理论之间的关系，由此证明该引理的非平坦版本。
事实上，构造 $\tau_{\geq -2}E(\mathcal{F})$ 已经足够，因为我们只关心
$\Ext^i_{X_A}(Li^*(i_*\mathcal{F}), \mathcal{F} \otimes_A M)$
在 $i = 0, 1, 2$ 时的 Ext 群。我们甚至可以确定其上同调层：
$$
H^i(E(\mathcal{F})) =
\left\{
\begin{matrix}
0 & \text{若 }i > 0 \\
\mathcal{F} & \text{若 } i = 0 \\
0 & \text{若 } i = -1 \\
\text{Tor}_1^{\mathcal{O}_S}(\mathcal{O}_X, A)
\otimes_{\mathcal{O}_X} \mathcal{F} &
\text{若 } i = -2
\end{matrix}
\right.
$$
这一观察将指导下面各条注记中对 $E(\mathcal{F})$ 的构造。

\begin{remark}[直接构造]
\label{obsolete-remark-construction-E}
设 $S$ 为概形，$f : X \to B$ 为 $S$ 上代数空间的态射，
$U$ 为另一个 $B$ 上代数空间。以 $q : X \times_B U \to U$
表示第二投影。考虑优三角
$$
Lq^*L_{U/B} \to L_{X \times_B U/B} \to E \to Lq^*L_{U/B}[1]
$$
，见 Cotangent 小节 \ref{cotangent-section-fibre-product}。
对任意 $\mathcal{O}_{X \times_B U}$-模层 $\mathcal{F}$，都有阿蒂亚类
$$
\mathcal{F} \to
L_{X \times_B U/B}
\otimes_{\mathcal{O}_{X \times_B U}}^\mathbf{L} \mathcal{F}[1]
$$
（见 Cotangent 小节 \ref{cotangent-section-atiyah-general}）。
将其与到 $E$ 的态射复合，并在 $D(\mathcal{O}_{X \times_B U})$ 中选取优三角
$$
E(\mathcal{F}) \to \mathcal{F} \to
\mathcal{F} \otimes_{\mathcal{O}_{X \times_B U}}^\mathbf{L} E[1] \to
E(\mathcal{F})[1]
$$
。由构造，阿蒂亚类可提升为态射
$$
e_\mathcal{F} :
E(\mathcal{F})
\longrightarrow
Lq^*L_{U/B} \otimes_{\mathcal{O}_{X \times_B U}}^\mathbf{L} \mathcal{F}[1]
$$
，并嵌入下列优三角态射：
$$
\xymatrix{
\mathcal{F} \otimes^\mathbf{L} Lq^*L_{U/B}[1] \ar[r] &
\mathcal{F} \otimes^\mathbf{L} L_{X \times_B U/B}[1] \ar[r] &
\mathcal{F} \otimes^\mathbf{L} E[1] \\
E(\mathcal{F}) \ar[r] \ar[u]^{e_\mathcal{F}} &
\mathcal{F} \ar[r] \ar[u]^{Atiyah} &
\mathcal{F} \otimes^\mathbf{L} E[1] \ar[u]^{=}
}
$$
给定 $S, B, X, f, U, \mathcal{F}$ 后，我们固定 $E(\mathcal{F})$
与 $e_\mathcal{F}$ 的一种选取。
\end{remark}

\begin{remark}[障碍类的构造]
\label{obsolete-remark-construction-ob}
采用注记 \ref{obsolete-remark-construction-E} 的记号，设 $i : U \to U'$ 是
$U$ 在 $B$ 上的一阶增厚。设
% P12-ERR-0508
$\mathcal{I} \subset \mathcal{O}_{U'}$ 是定义 $B$ 在 $B'$ 中嵌入的
拟凝聚理想层。基本三角
$$
Li^*L_{U'/B} \to L_{U/B} \to L_{U/U'} \to Li^*L_{U'/B}[1]
$$
与态射 $L_{U/U'} \to \mathcal{I}[1]$ 共同确定态射
$e_{U'} : L_{U/B} \to \mathcal{I}[1]$。将它同上一注记中的态射
$e_\mathcal{F}$ 结合，得到
$$
(\text{id}_\mathcal{F} \otimes Lq^*e_{U'}) \cup e_\mathcal{F} :
E(\mathcal{F})
\longrightarrow
\mathcal{F} \otimes_{\mathcal{O}_{X \times_B U}} q^*\mathcal{I}[2]
$$
（这里还同从导出张量积到通常张量积的态射作了复合）。换言之，我们得到元素
$$
\xi_{U'} \in
\Ext^2_{\mathcal{O}_{X \times_B U}}(
E(\mathcal{F}),
\mathcal{F} \otimes_{\mathcal{O}_{X \times_B U}} q^*\mathcal{I})
$$
\end{remark}

\begin{lemma}
\label{obsolete-lemma-ob-is-obstruction}
在注记 \ref{obsolete-remark-construction-ob} 的情形中，假设 $\mathcal{F}$ 在 $U$ 上平坦。
则类 $\xi_{U'}$ 为零，是存在满足 $i^*\mathcal{F}' \cong \mathcal{F}$
且在 $U'$ 上平坦的 $\mathcal{O}_{X \times_B U'}$-模 $\mathcal{F}'$
的充分必要条件。
\end{lemma}

\begin{proof}[证明（概要）]
我们将使用 Deformation Theory 的引理
\ref{defos-lemma-inf-obs-ext-rel-ringed-topoi} 中的判据。
简记 $\mathcal{O} = \mathcal{O}_{X \times_B U}$ 以及
$\mathcal{O}' = \mathcal{O}_{X \times_B U'}$。考虑短正合列
$$
0 \to \mathcal{I} \to \mathcal{O}_{U'} \to \mathcal{O}_U \to 0.
$$
设 $\mathcal{J} \subset \mathcal{O}'$ 是定义 $X \times_B U$ 的拟凝聚理想层。
由上可得正合列
$$
\text{Tor}_1^{\mathcal{O}_B}(\mathcal{O}_X, \mathcal{O}_U) \to
q^*\mathcal{I} \to \mathcal{J} \to 0
$$
其中 $\text{Tor}_1^{\mathcal{O}_B}(\mathcal{O}_X, \mathcal{O}_U)$ 是下式的简写：
$$
\text{Tor}_1^{h^{-1}\mathcal{O}_B}(p^{-1}\mathcal{O}_X, q^{-1}\mathcal{O}_U)
\otimes_{(p^{-1}\mathcal{O}_X\otimes_{h^{-1}\mathcal{O}_B}q^{-1}\mathcal{O}_U)}
\mathcal{O}.
$$
与 $\mathcal{F}$ 作张量积，得到正合列
$$
\mathcal{F} \otimes_\mathcal{O}
\text{Tor}_1^{\mathcal{O}_B}(\mathcal{O}_X, \mathcal{O}_U) \to
\mathcal{F} \otimes_\mathcal{O}
q^*\mathcal{I} \to
\mathcal{F} \otimes_\mathcal{O} \mathcal{J} \to 0
$$
（注意，相对于 Deformation Theory 的引理
\ref{defos-lemma-inf-obs-ext-rel-ringed-topoi} 中的记号，字母
$\mathcal{I}$ 与 $\mathcal{J}$ 的作用恰好对调。）该引理的条件 (1)
要求上面的最后一个态射为同构，亦即第一个态射为零。
这一态射是否为零，可在几何点
$\overline{z} = (\overline{x}, \overline{u}) : \Spec(k) \to X \times_B U$
处的茎上检验。置 $R = \mathcal{O}_{B, \overline{b}}$、
$A = \mathcal{O}_{X, \overline{x}}$、$B = \mathcal{O}_{U, \overline{u}}$，以及
$C = \mathcal{O}_{\overline{z}}$。由 Cotangent 的引理
\ref{cotangent-lemma-fibre-product} 和 $E(\mathcal{F})$ 的定义三角可见
$$
H^{-2}(E(\mathcal{F}))_{\overline{z}} =
\mathcal{F}_{\overline{z}} \otimes \text{Tor}_1^R(A, B)
$$
因此态射 $\xi_{U'}$ 诱导态射
$$
\mathcal{F}_{\overline{z}} \otimes \text{Tor}_1^R(A, B)
\longrightarrow
\mathcal{F}_{\overline{z}} \otimes_B \mathcal{I}_{\overline{u}}
$$
我们断言，它与上述态射的茎相同（证明从略；这是一个纯粹的环论陈述）。
由此可见，Deformation Theory 的引理
\ref{defos-lemma-inf-obs-ext-rel-ringed-topoi} 的条件 (1) 等价于下列态射为零：
$H^{-2}(\xi_{U'}) :
H^{-2}(E(\mathcal{F})) \to \mathcal{F} \otimes \mathcal{I}$.

\medskip\noindent
为完成证明，我们说明：在条件 (1) 成立的假设下，条件 (2) 等价于
$\xi_{U'}$ 为零。在证明的余下部分，以 $\mathcal{F} \otimes \mathcal{I}$ 表示
$\mathcal{F} \otimes_\mathcal{O} q^*\mathcal{I} =
\mathcal{F} \otimes_\mathcal{O} \mathcal{J}$。考察谱序列
$$
\Ext^i(H^{-j}(E(\mathcal{F})), \mathcal{F} \otimes \mathcal{I})
\Rightarrow
\Ext^{i + j}(E(\mathcal{F}), \mathcal{F} \otimes \mathcal{I})
$$
并利用 $H^0(E(\mathcal{F})) = \mathcal{F}$ 与
$H^{-1}(E(\mathcal{F})) = 0$，可知存在正合列
$$
0 \to
\Ext^2(\mathcal{F}, \mathcal{F} \otimes \mathcal{I}) \to
\Ext^2(E(\mathcal{F}), \mathcal{F} \otimes \mathcal{I}) \to
\Hom(H^{-2}(E(\mathcal{F})), \mathcal{F} \otimes \mathcal{I})
$$
因此元素 $\xi_{U'}$ 属于
$\Ext^2(\mathcal{F}, \mathcal{F} \otimes \mathcal{I})$。
只需说明它与 Deformation Theory 的引理
\ref{defos-lemma-inf-obs-ext-rel-ringed-topoi} 中的元素相同即可完成证明；
这一验证从略。
\end{proof}

\begin{lemma}
\label{obsolete-lemma-coherent-defo-thy-general}
在 Quot 的情形 \ref{quot-situation-coherent} 中，假设 $S$ 是局部诺特概形且
$S = B$。令 $\mathcal{X} = \textit{Coh}_{X/B}$。则 $\mathcal{X}$
具有 versal 性的开性（见 Artin's Axioms 的定义
\ref{artin-definition-openness-versality}）。
\end{lemma}

\begin{proof}[证明（概要）]
设 $U \to S$ 为有限型概形态射，$x$ 为 $U$ 上的 $\mathcal{X}$-对象，
$u_0 \in U$ 为有限型点，且 $x$ 在 $u_0$ 处为 versal。缩小 $U$ 后，
可假设 $u_0$ 是闭点（Morphisms 的引理
\ref{morphisms-lemma-point-finite-type}），并且 $U = \Spec(A)$，
而 $U \to S$ 的像落在 $S$ 的仿射开集 $\Spec(\Lambda)$ 中。
我们将使用 Artin's Axioms 的引理 \ref{artin-lemma-dual-openness}
证明本引理。设 $\mathcal{F}$ 是 $X_A = \Spec(A) \times_S X$ 上
对应于给定对象 $x$、且在 $A$ 上平坦的凝聚模。

\medskip\noindent
按注记 \ref{obsolete-remark-construction-E} 选取 $E(\mathcal{F})$ 与
$e_\mathcal{F}$。由 $E(\mathcal{F})$ 的上同调层描述可知，对任意
$A$-模 $M$ 都有
$$
\Ext^1(E(\mathcal{F}), \mathcal{F} \otimes_A M) =
\Ext^1(\mathcal{F}, \mathcal{F} \otimes_A M)
$$
。结合这一点与 Deformation Theory 的引理
\ref{defos-lemma-inf-ext-rel-ringed-topoi}，得到函子同构
$$
T_x(M) = \Ext^1_{X_A}(E(\mathcal{F}), \mathcal{F} \otimes_A M)
$$
。由引理 \ref{obsolete-lemma-ob-is-obstruction}，对任意核为平方零理想 $I$ 的
$\Lambda$-代数满射 $A' \to A$，都有障碍类
$$
\xi_{A'} \in \Ext^2_{X_A}(E(\mathcal{F}), \mathcal{F} \otimes_A I)
$$
。将 Derived Categories of Spaces 的引理
\ref{spaces-perfect-lemma-compute-ext} 应用于 Ext 群
$\Ext^i_{X_A}(E(\mathcal{F}), \mathcal{F} \otimes_A M)$
在 $i \leq m$、$m = 2$ 时的计算。我们略去
$E(\mathcal{F})$ 属于 $D^-_{\textit{Coh}}$ 的验证；提示：使用
Cotangent 的引理 \ref{cotangent-lemma-cotangent-finite}。
于是得到完美对象 $K \in D(A)$ 以及函子性同构
$$
H^i(K \otimes_A^\mathbf{L} M)
\longrightarrow
\Ext^i_{X_A}(E(\mathcal{F}), \mathcal{F} \otimes_A M)
$$
，其中 $i \leq m$，且这些同构与边界态射相容。对象 $K$ 连同上述同一化，
给出了 Artin's Axioms 的情形 \ref{artin-situation-dual} 中的一组数据。
最后，Artin's Axioms 的引理 \ref{artin-lemma-dual-obstruction}
之条件 (iv)，可由 Deformation Theory 的引理
\ref{defos-lemma-verify-iv-ringed-topoi} 的一个变体推出；
我们略去该变体的陈述与证明。因此 Artin's Axioms 的引理
\ref{artin-lemma-dual-openness} 适用，引理得证。
\end{proof}

\begin{theorem}
\label{obsolete-theorem-coherent-algebraic-general}
设 $S$ 为概形，$f : X \to B$ 为 $S$ 上代数空间的态射。
假设 $f$ 有限表示且分离。则 $\textit{Coh}_{X/B}$ 是 $S$ 上的代数叠。
\end{theorem}

\begin{proof}
本定理是 Quot 的定理 \ref{quot-theorem-coherent-algebraic-general}
的一份副本。之所以在此保留它，是因为本节材料给出了第二种证明
（如本节开头所述）。具体而言，论证与 Quot 的定理
\ref{quot-theorem-coherent-algebraic} 的证明完全相同，
但以引理 \ref{obsolete-lemma-coherent-defo-thy-general} 取代 Quot 的引理
\ref{quot-lemma-coherent-defo-thy}。
\end{proof}







\section{修改}
\label{obsolete-section-modifications}

\noindent
下面是关于 Algebraization of Formal Spaces 中式
(\ref{restricted-equation-modification}) 所给范畴的一项已废弃结果。
现行材料请参见 Algebraization of Formal Spaces 的小节
\ref{restricted-section-modifications}。

\begin{lemma}
\label{obsolete-lemma-henselian}
设 $(A, \mathfrak m, \kappa)$ 为诺特局部环。Algebraization of Formal Spaces
的式 (\ref{restricted-equation-modification}) 给出的 $A$ 的范畴，
等价于同一式 (\ref{restricted-equation-modification}) 给出的
$A$ 的亨泽尔化 $A^h$ 的范畴。
\end{lemma}

\begin{proof}
这是 Algebraization of Formal Spaces 的引理
\ref{restricted-lemma-equivalence-to-completion} 的特殊情形。
\end{proof}

\noindent
下面关于有理奇点的引理，在曲面奇点消解一章中已不再需要。

\begin{lemma}
\label{obsolete-lemma-double-dual-rational}
在 Resolution of Surfaces 的情形 \ref{resolve-situation-rational} 中，
设 $M$ 为有限反身 $A$-模。以 $M \otimes_A \mathcal{O}_X$
表示相应 $\mathcal{O}_S$-模的拉回。则
$M \otimes_A \mathcal{O}_X$ 满射到其双对偶。
\end{lemma}

\begin{proof}
令 $\mathcal{F} = (M \otimes_A \mathcal{O}_X)^{**}$ 为双对偶，
并令 $\mathcal{F}' \subset \mathcal{F}$ 为求值态射
$M \otimes_A \mathcal{O}_X \to \mathcal{F}$ 的像。于是有短正合列
$$
0 \to \mathcal{F}' \to \mathcal{F} \to \mathcal{Q} \to 0
$$
。由于 $X$ 正规，对于余维 $1$ 的点，局部环 $\mathcal{O}_{X, x}$
均为离散赋值环（见 Properties 的引理
\ref{properties-lemma-criterion-normal}）。因此由 More on Algebra 的引理
\ref{more-algebra-lemma-cokernel-map-double-dual-dvr}，
在这类点上有 $\mathcal{Q}_x = 0$。故 $\mathcal{Q}$ 支撑在有限多个闭点上，
并由 Cohomology of Schemes 的引理
\ref{coherent-lemma-coherent-support-dimension-0} 得到它由整体截面生成。
我们得到正合列
$$
0 \to H^0(X, \mathcal{F}') \to H^0(X, \mathcal{F}) \to H^0(X, \mathcal{Q}) \to 0
$$
，因为 $\mathcal{F}'$ 由整体截面生成（Resolution of Surfaces 的引理
\ref{resolve-lemma-globally-generated}）。由于 $X \to \Spec(A)$
在闭点的补集上为同构，且 $M$ 反身，可见态射
$$
M \to H^0(X, \mathcal{F}') \to H^0(X, \mathcal{F})
$$
在 $A$ 的任意非极大素理想处局部化后诱导同构。由 More on Algebra 的引理
\ref{more-algebra-lemma-check-isomorphism-via-depth-and-ass}，以及正规环上的
反身模具有性质 $(S_2)$ 这一事实（More on Algebra 的引理
\ref{more-algebra-lemma-reflexive-over-normal}），这些态射因而都是同构。
故 $\mathcal{Q} = 0$，正如所需。
\end{proof}



\section{相交理论}
\label{obsolete-section-intersection-theory}

\begin{lemma}
\label{obsolete-lemma-good-blowing-up}
设 $b : X' \to X$ 是域 $k$ 上光滑射影概形沿其光滑闭子概形
$Z \subset X$ 的吹起。图示为
$$
\xymatrix{
E \ar[r]_j \ar[d]_\pi & X' \ar[d]^b \\
Z \ar[r]^i & X
}
$$
假设存在 $K_0(X)$ 中的元素，其在 $Z$ 上的限制等于
$K_0(Z)$ 中 $\mathcal{C}_{Z/X}$ 的类。则对某个
$\alpha'' \in K_0(X')$，在 $K_0(X')$ 中有
$[Lb^*\mathcal{O}_Z] = [\mathcal{O}_E] \cdot \alpha''$。
\end{lemma}

\begin{proof}
概形 $X$、$X'$、$E$、$Z$ 均在 $k$ 上光滑且射影，因而
$K'_0(X) = K_0(X) = K_0(\textit{Vect}(X)) =
K_0(D^b_{\textit{Coh}}(X)))$，其余 $3$ 个对象亦然。参见
Derived Categories of Schemes 的引理
\ref{perfect-lemma-Noetherian-Kprime}、\ref{perfect-lemma-Kprime-K}
与 \ref{perfect-lemma-K-is-old-K}。本证明中将任意换用这些版本。
考虑定义 $\mathcal{F}$ 的有限局部自由 $\mathcal{O}_E$-模短正合列
$$
0 \to \mathcal{F} \to \pi^*\mathcal{C}_{Z/X} \to \mathcal{C}_{E/X'} \to 0
$$
% P12-ERR-0509
。注意，$\mathcal{C}_{E/X'} = \mathcal{O}_{X'}(-E)|_E$
是可逆 $\mathcal{O}_X$-模 $\mathcal{O}_{X'}(-E)$ 的限制。
取 $\alpha \in K_0(X)$，使得在 $K_0(Z)$ 中
$i^*\alpha = [\mathcal{C}_{Z/X}]$。令
$\alpha' = b^*\alpha - [\mathcal{O}_{X'}(-E)]$。
于是 $j^*\alpha' = [\mathcal{F}]$。由 Weil Cohomology Theories 的引理
\ref{weil-lemma-lambda-operations}，推出
$j^*\lambda^i(\alpha') = [\wedge^i(\mathcal{F})]$。
% P12-ERR-0510 P12-ERR-0511
这意味着在 $K_0(X)$ 中
$[\mathcal{O}_E] \cdot \alpha' = [\wedge^i\mathcal{F}]$；参见
Derived Categories of Schemes 的引理
\ref{perfect-lemma-projection-formula}。令 $r$ 为 $Z$ 的不可约分支在 $X$
中的最大余维。一个从略的计算表明
% P12-ERR-0512
$H^{-i}(Lb^*\mathcal{O}_Z) = \wedge^i\mathcal{F}$
当 $i \geq 0, 1, \ldots, r - 1$ 时成立，而在其余次数为零。
因此在 $K_0(X)$ 中有
\begin{align*}
[Lb^*\mathcal{O}_Z] & =
\sum\nolimits_{i = 0, \ldots, r - 1} (-1)^i[\wedge^i\mathcal{F}] \\
& =
\sum\nolimits_{i = 0, \ldots, r - 1} (-1)^i[\mathcal{O}_E] \lambda^i(\alpha') \\
& =
[\mathcal{O}_E] \left(\sum\nolimits_{i = 0, \ldots, r - 1}
(-1)^i \lambda^i(\alpha')\right)
\end{align*}
取 $\alpha'' = \sum_{i = 0, \ldots, r - 1} (-1)^i \lambda^i(\alpha')$
即得本引理。
\end{proof}



\begin{lemma}
\label{obsolete-lemma-gysin-factors-principal}
设 $(S, \delta)$ 如 Chow Homology 的情形 \ref{chow-situation-setup}。
设 $X$ 在 $S$ 上局部有限型。假设 $X$ 整，且 $n = \dim_\delta(X)$。
设 $a \in \Gamma(X, \mathcal{O}_X)$ 为非零函数，并设
$i : D = Z(a) \to X$ 为 $a$ 的零概形之闭浸入。令 $f \in R(X)^*$。
此时在 $A_{n - 2}(D)$ 中有 $i^*\text{div}_X(f) = 0$。
\end{lemma}

\begin{proof}
这是 Chow Homology 的引理 \ref{chow-lemma-gysin-factors-general} 的特殊情形。
\end{proof}

\begin{remark}
\label{obsolete-remark-not-true-not-quasi-compact}
本注记过去曾说，尚不清楚 Chow Homology 的引理
\ref{chow-lemma-cycles-k-group} 中的箭头在一般情况下是否为同构。
不过，现在我们已经找到了这一事实的证明。
\end{remark}







\subsection{吹起引理}
\label{obsolete-section-blowing-up-lemmas}

\noindent
本节证明若干引理：在吹起上用适当的有效 Cartier 因子表示 Cartier 因子。
这些引理见 \cite[Section 2.4]{F}。我们调整了表述，使其在非有限型情形下
同样适用。引理 \ref{obsolete-lemma-blowing-up-intersections} 中的态射 $b$
可能是无穷多个吹起的复合，但在任意给定的拟紧开集 $W \subset X$ 上，
只需有限多个吹起（这正是同上引文中的结果；loc.\ cit.）。

\begin{lemma}
\label{obsolete-lemma-push-pull-effective-Cartier}
设 $(S, \delta)$ 如 Chow Homology 的情形 \ref{chow-situation-setup}。
设 $X$、$Y$ 在 $S$ 上局部有限型，$f : X \to Y$ 为固有态射，
$D \subset Y$ 为有效 Cartier 因子。假设 $X$、$Y$ 整，
$n = \dim_\delta(X) = \dim_\delta(Y)$，且 $f$ 支配。则
$$
f_*[f^{-1}(D)]_{n - 1} = [R(X) : R(Y)] [D]_{n - 1}.
$$
特别地，若 $f$ 双有理，则 $f_*[f^{-1}(D)]_{n - 1} = [D]_{n - 1}$。
\end{lemma}

\begin{proof}
这立即来自 Chow Homology 的引理 \ref{chow-lemma-equal-c1-as-cycles}，
以及 $D$ 是 $\mathcal{O}_X(D)$ 的典范截面 $1_D$ 的零概形这一事实。
\end{proof}

\begin{lemma}
\label{obsolete-lemma-blowing-up-denominators}
设 $(S, \delta)$ 如 Chow Homology 的情形 \ref{chow-situation-setup}。
设 $X$ 在 $S$ 上局部有限型，并假设 $X$ 整且 $\dim_\delta(X) = n$。
设 $\mathcal{L}$ 为可逆 $\mathcal{O}_X$-模，$s$ 为 $\mathcal{L}$ 的非零亚纯截面。
令 $U \subset X$ 为使 $s$ 对应于 $\mathcal{L}$ 在 $U$ 上一个截面的最大开子概形。
则存在射影态射
$$
\pi : X' \longrightarrow X
$$
，满足
\begin{enumerate}
\item $X'$ 整；
\item $\pi|_{\pi^{-1}(U)} : \pi^{-1}(U) \to U$ 是同构；
\item 存在有效 Cartier 因子 $D, E \subset X'$，使得
$$
\pi^*\mathcal{L} = \mathcal{O}_{X'}(D - E),
$$
\item 通过上述同构，亚纯截面 $s$ 对应于亚纯截面
$1_D \otimes (1_E)^{-1}$（见 Divisors 的定义
\ref{divisors-definition-invertible-sheaf-effective-Cartier-divisor}）；
\item 有
$$
\pi_*([D]_{n - 1} - [E]_{n - 1}) = \text{div}_\mathcal{L}(s)
$$
在 $Z_{n - 1}(X)$ 中成立。
\end{enumerate}
\end{lemma}

\begin{proof}
设 $\mathcal{I} \subset \mathcal{O}_X$ 为 $s$ 的拟凝聚分母理想层，
见 Divisors 的定义
\ref{divisors-definition-regular-meromorphic-ideal-denominators}。
由 Divisors 的引理 \ref{divisors-lemma-blowing-up-denominators}
得到 (2)、(3)、(4)。由 Divisors 的引理
\ref{divisors-lemma-blow-up-integral-scheme} 得到 (1)。
由 Divisors 的引理 \ref{divisors-lemma-blowing-up-projective}，
态射 $\pi$ 为射影态射。尚需证明 (5)。由 Chow Homology 的引理
\ref{chow-lemma-equal-c1-as-cycles}，有
$$
\pi_*(\text{div}_{\mathcal{L}'}(s')) = \text{div}_\mathcal{L}(s).
$$
因此只需证明
$\text{div}_{\mathcal{L}'}(s') = [D]_{n - 1} - [E]_{n - 1}$。
这由等式 $s' = 1_D \otimes 1_E^{-1}$ 和可加性推出；
见 Divisors 的引理 \ref{divisors-lemma-c1-additive}。
\end{proof}

\begin{definition}
\label{obsolete-definition-epsilon}
设 $(S, \delta)$ 如 Chow Homology 的情形 \ref{chow-situation-setup}。
设 $X$ 在 $S$ 上局部有限型，并假设 $X$ 整且 $\dim_\delta(X) = n$。
设 $D_1, D_2$ 为 $X$ 中的两个有效 Cartier 因子，
$Z \subset X$ 为满足 $\dim_\delta(Z) = n - 1$ 的整闭子概形。
这一情形的{\it $\epsilon$-不变量}为
$$
\epsilon_Z(D_1, D_2) = n_Z \cdot m_Z
$$
，其中 $n_Z$（resp.\ $m_Z$）是 $Z$ 在 $(n - 1)$-循环
$[D_1]_{n - 1}$（resp.\ $[D_2]_{n - 1}$）中的系数。
\end{definition}

\begin{lemma}
\label{obsolete-lemma-two-divisors}
设 $(S, \delta)$ 如 Chow Homology 的情形 \ref{chow-situation-setup}。
设 $X$ 在 $S$ 上局部有限型，并假设 $X$ 整且 $\dim_\delta(X) = n$。
设 $D_1, D_2$ 为 $X$ 中的两个有效 Cartier 因子，
$Z$ 为概形 $D_1 \cap D_2$ 的开闭子概形。假设
$\dim_\delta(D_1 \cap D_2 \setminus Z) \leq n - 2$。
则存在态射 $b : X' \to X$ 以及 $X'$ 上的 Cartier 因子
$D_1', D_2', E$，满足下列性质：
\begin{enumerate}
\item $X'$ 整；
\item $b$ 为射影态射；
\item $b$ 是 $X$ 沿闭子概形 $Z$ 的吹起；
\item $E = b^{-1}(Z)$；
\item $b^{-1}(D_1) = D'_1 + E$，且 $b^{-1}D_2 = D_2' + E$；
\item $\dim_\delta(D'_1 \cap D'_2) \leq n - 2$；并且若
$Z = D_1 \cap D_2$，则 $D'_1 \cap D'_2 = \emptyset$；
\item 对每个满足 $\dim_\delta(W') = n - 1$ 的整闭子概形 $W'$，均有：
\begin{enumerate}
\item 若 $\epsilon_{W'}(D'_1, E) > 0$，则置 $W = b(W')$ 后有
$\dim_\delta(W) = n - 1$，且
$$
\epsilon_{W'}(D'_1, E) < \epsilon_W(D_1, D_2),
$$
\item 若 $\epsilon_{W'}(D'_2, E) > 0$，则置 $W = b(W')$ 后有
$\dim_\delta(W) = n - 1$，且
$$
\epsilon_{W'}(D'_2, E) < \epsilon_W(D_1, D_2),
$$
\end{enumerate}
\end{enumerate}
\end{lemma}

\begin{proof}
注意，拟凝聚理想层
$\mathcal{I} = \mathcal{I}_{D_1} + \mathcal{I}_{D_2}$
定义了概形论交 $D_1 \cap D_2 \subset X$。由于 $Z$ 是
$D_1 \cap D_2$ 的若干连通分支之并，对每个 $z \in Z$，态射
$\mathcal{O}_{X, z} \to \mathcal{O}_{Z, z}$ 的核等于 $\mathcal{I}_z$。
令 $b : X' \to X$ 为 $X$ 沿 $Z$ 的吹起（所以在 $Z$ 附近 Zariski 局部，
它是 $X$ 沿 $\mathcal{I}$ 的吹起）。以 $E = b^{-1}(Z)$ 表示相应的
有效 Cartier 因子；见 Divisors 的引理
\ref{divisors-lemma-blowing-up-gives-effective-Cartier-divisor}。
% P12-ERR-0513
由于 $Z \subset D_1$，有 $E \subset f^{-1}(D_1)$，故对某个有效
Cartier 因子 $D'_1 \subset X'$，有 $D_1 = D_1' + E$；见 Divisors 的引理
\ref{divisors-lemma-difference-effective-Cartier-divisors}。
同理，$D_2 = D_2' + E$。这就处理了断言 (1) -- (5)。

\medskip\noindent
注意，若 $W'$ 如 (7) (a) 或 (7) (b) 所述，则 $W'$ 的像 $W$
包含于 $D_1 \cap D_2$。若 $W$ 不包含于 $Z$，则 $b$ 在 $W$ 的一般点处
为同构，从而 $\dim_\delta(W) = \dim_\delta(W') = n - 1$，
这与假设 $\dim_\delta(D_1 \cap D_2 \setminus Z) \leq n - 2$ 矛盾。
故 $W \subset Z$。这意味着为证明 (6) 与 (7)，可以在 $X$ 上 $Z$ 的附近局部工作。

\medskip\noindent
因此可假设 $X = \Spec(A)$，其中 $A$ 为诺特整环，并且
$D_1 = \Spec(A/a)$、$D_2 = \Spec(A/b)$、$Z = D_1 \cap D_2$。
置 $I = (a, b)$。由于 $A$ 为整环且 $a, b \not = 0$，可用两个片覆盖吹起，即
$U = \Spec(A[s]/(as - b))$ 与 $V = \Spec(A[t]/(bt -a))$。
这两个片通过同构
$A[s, s^{-1}]/(as - b) \cong A[t, t^{-1}]/(bt - a)$
黏合，该同构将 $s$ 映到 $t^{-1}$。有效 Cartier 因子 $E$ 由
$\Spec(A[s]/(as - b, a)) \subset U$ 与
$\Spec(A[t]/(bt - a, b)) \subset V$ 描述。闭子概形 $D'_1$ 对应于
% P12-ERR-0514
$\Spec(A[t]/(bt - a, t)) \subset U$。闭子概形 $D'_2$ 对应于
$\Spec(A[s]/(as -b, s)) \subset V$。由于“$ts = 1$”，可见
$D'_1 \cap D'_2 = \emptyset$。

\medskip\noindent
假设有高度为一的素理想 $\mathfrak q \subset A[s]/(as - b)$，且
$s, a \in \mathfrak q$。令 $\mathfrak p \subset A$ 为 $A$ 中相应的素理想。
注意 $a, b \in \mathfrak p$。由维数公式亦有
$\dim(A_{\mathfrak p}) = 1$。最后需要证明的断言是
$$
\text{ord}_{A_{\mathfrak p}}(a)
\text{ord}_{A_{\mathfrak p}}(b)
>
\text{ord}_{B_{\mathfrak q}}(a)
\text{ord}_{B_{\mathfrak q}}(s)
$$
，其中 $B = A[s]/(as - b)$。由 Algebra 的引理
\ref{algebra-lemma-quasi-finite-extension-dim-1}，对 $x = a, b$ 有
$\text{ord}_{A_{\mathfrak p}}(x) \geq \text{ord}_{B_{\mathfrak q}}(x)$。
又因 $\text{ord}_{B_{\mathfrak q}}(s) > 0$，由 $\text{ord}$ 函数的可加性
以及 $as = b$，所需不等式成立。
\end{proof}

\begin{definition}
\label{obsolete-definition-locally-finite-sum-effective-Cartier-divisors}
设 $X$ 为概形，$\{D_i\}_{i \in I}$ 为 $X$ 上一族局部有限的
有效 Cartier 因子。给定函数
$I \to \mathbf{Z}_{\geq 0}$，$i \mapsto n_i$。这些有效 Cartier 因子的
{\it 和} $D = \sum n_i D_i$，是唯一满足下列性质的有效 Cartier 因子
$D \subset X$：对任意拟紧开集 $U \subset X$，
$D|_U = \sum_{D_i \cap U \not = \emptyset} n_iD_i|_U$ 是 Divisors 的定义
\ref{divisors-definition-sum-effective-Cartier-divisors} 中所定义的和。
\end{definition}

\begin{lemma}
\label{obsolete-lemma-sum-divisors-associated-Weil}
设 $(S, \delta)$ 如 Chow Homology 的情形 \ref{chow-situation-setup}。
设 $X$ 在 $S$ 上局部有限型，并假设 $X$ 整且 $\dim_\delta(X) = n$。
设 $\{D_i\}_{i \in I}$ 为 $X$ 上一族局部有限的有效 Cartier 因子，
并对 $i \in I$ 给定 $n_i \geq 0$。则
$$
[D]_{n - 1} = \sum\nolimits_i n_i[D_i]_{n - 1}
$$
在 $Z_{n - 1}(X)$ 中成立。
\end{lemma}

\begin{proof}
由于要证明的是循环的等式，可以在 $X$ 上局部工作。因此问题化为有限和，
再由归纳化为两个有效 Cartier 因子的和 $D = D_1 + D_2$。
由 Chow Homology 的引理 \ref{chow-lemma-compute-c1} 可见
$D_1 = \text{div}_{\mathcal{O}_X(D_1)}(1_{D_1})$，其中
$1_{D_1}$ 表示 $\mathcal{O}_X(D_1)$ 的典范截面。$D_2$ 与 $D$ 当然也有同样的陈述。
在同一化 $\mathcal{O}_X(D) = \mathcal{O}_X(D_1) \otimes \mathcal{O}_X(D_2)$
下，$1_D = 1_{D_1} \otimes 1_{D_2}$，故结论由 Divisors 的引理
\ref{divisors-lemma-c1-additive} 得出。
\end{proof}

\begin{lemma}
\label{obsolete-lemma-blowing-up-intersections}
设 $(S, \delta)$ 如 Chow Homology 的情形 \ref{chow-situation-setup}。
设 $X$ 在 $S$ 上局部有限型，并假设 $X$ 整且 $\dim_\delta(X) = d$。
设 $\{D_i\}_{i \in I}$ 为 $X$ 上一族局部有限的有效 Cartier 因子。
假设对所有满足 $\#\{i, j, k\} = 3$ 的 $\{i, j, k\} \subset I$，都有
$D_i \cap D_j \cap D_k = \emptyset$。则存在
\begin{enumerate}
\item 开子概形 $U \subset X$，满足
$\dim_\delta(X \setminus U) \leq d - 3$；
\item 态射 $b : U' \to U$；
\item $U'$ 上的有效 Cartier 因子 $\{D'_j\}_{j \in J}$；
\end{enumerate}
它们具有下列性质：
\begin{enumerate}
\item $b$ 是固有态射 $b : U' \to U$；
\item $U'$ 整；
\item 在各 $D_i|_U$ 两两交之并的补集上，$b$ 为同构；
\item $\{D'_j\}_{j \in J}$ 是 $U'$ 上一族局部有限的有效 Cartier 因子；
\item 若 $j \not = j'$，则 $\dim_\delta(D'_j \cap D'_{j'}) \leq d - 2$；
\item 对某些 $n_{ij} \geq 0$，有
$b^{-1}(D_i|_U) = \sum n_{ij} D'_j$。
\end{enumerate}
此外，若 $X$ 拟紧，则上面可以假设 $U = X$。
\end{lemma}

\begin{proof}
先证明拟紧情形，因为这也许是最有意思的情形。在此情形中，归纳构造一列吹起
$$
X = X_0 \xleftarrow{b_0} X_1 \xleftarrow{b_1} X_2 \leftarrow \ldots
$$
以及有效 Cartier 因子的有限集 $\{D_{n, i}\}_{i \in I_n}$。
在每一步，任意三重交 $D_{n, i} \cap D_{n, j} \cap D_{n, k}$ 都为空。
此外，对每个 $n \geq 0$，都有
$I_{n + 1} = I_n \amalg P(I_n)$，其中 $P(I_n)$ 表示 $I_n$ 中元素对的集合。
最后，还有
$$
b_n^{-1}(D_{n, i}) = D_{n + 1, i} +
\sum\nolimits_{i' \in I_n, i' \not = i} D_{n + 1, \{i, i'\}}
$$
从而对每个 $n \geq 0$，$(b_0 \circ \ldots \circ b_n)^{-1}(D_i)$
都是各因子 $D_{n + 1, j}$（$j \in I_{n + 1}$）的非负整系数组合。

\medskip\noindent
为开始归纳，置 $X_0 = X$、$I_0 = I$、$D_{0, i} = D_i$。

\medskip\noindent
给定 $(X_n, \{D_{n, i}\}_{i \in I_n})$，令 $X_{n + 1}$ 为 $X_n$
沿闭子概形
$Z_n = \bigcup_{\{i, i'\} \in P(I_n)} D_{n, i} \cap D_{n, i'}$
的吹起。由关于三重交的假设，各闭子概形
$D_{n, i} \cap D_{n, i'}$ 两两不交。换言之，可以写成
$Z_n = \coprod_{\{i, i'\} \in P(I_n)} D_{n, i} \cap D_{n, i'}$.
此外，在 $D_{n, i} \cap D_{n, i'}$ 的 Zariski 邻域中，态射 $b_n$
等于概形 $X_n$ 沿闭子概形 $D_{n, i} \cap D_{n, i'}$ 的吹起，
故引理 \ref{obsolete-lemma-two-divisors} 的结论适用。因此置
% P12-ERR-0515
$D_{n + 1, \{i, i'\}} = b_n^{-1}(D_i \cap D_{i'})$，
便得到有效 Cartier 因子。各 Cartier 因子
$D_{n + 1, \{i, i'\}}$ 两两不交。显然，对每个
$i' \in I_n$、$i' \not = i$，都有
$b_n^{-1}(D_{n, i}) \supset D_{n + 1, \{i, i'\}}$。
因此应用 Divisors 的引理
\ref{divisors-lemma-difference-effective-Cartier-divisors}，可见确有
% P12-ERR-0516
$b^{-1}(D_{n, i}) = D_{n + 1, i} +
\sum\nolimits_{i' \in I_n, i' \not = i} D_{n + 1, \{i, i'\}}$
，其中 $D_{n + 1, i}$ 是 $X_{n + 1}$ 上的某个有效 Cartier 因子。
在 $D_{n + 1, \{i, i'\}}$ 的邻域中，这些因子 $D_{n + 1, i}$
扮演引理 \ref{obsolete-lemma-two-divisors} 中带撇因子的角色。特别地，
由引理 \ref{obsolete-lemma-two-divisors} 第 (6) 部分可知，若
$i \not = i'$ 且 $i, i' \in I_n$，则
$D_{n + 1, i} \cap D_{n + 1, i'} = \emptyset$。
这已经蕴含各因子 $D_{n + 1, i}$ 的三重交为零。

\medskip\noindent
至此，可以利用 $X$ 的拟紧性断定不变量
\begin{equation}
\label{obsolete-equation-invariant}
\epsilon(X, \{D_i\}_{i \in I}) =
\max\{\epsilon_Z(D_i, D_{i'}) \mid
Z \subset X,
\dim_\delta(Z) = d - 1,
\{i, i'\} \in P(I)\}
\end{equation}
是有限的；毕竟每个 $D_i$ 至多只有有限多个不可约分支。我们断言，对某个
$n$，不变量 $\epsilon(X_n, \{D_{n, i}\}_{i \in I_n})$ 为零。
否则，由引理 \ref{obsolete-lemma-two-divisors} 将得到严格递减序列
$$
\epsilon(X, \{D_i\}_{i \in I})
=
\epsilon(X_0, \{D_{0, i}\}_{i \in I_0})
>
\epsilon(X_1, \{D_{1, i}\}_{i \in I_1})
>
\ldots
$$
，其各项均为正整数，这就产生矛盾。取 $n$ 使得不变量
$\epsilon(X_n, \{D_{n, i}\}_{i \in I_n})$ 为零。这意味着不存在整闭子概形
$Z \subset X_n$ 和指标对 $i, i' \in I_n$，使得
$\epsilon_Z(D_{n, i}, D_{n, i'}) > 0$。换言之，对所有
$\{i, i'\} \in P(I_n)$，都有
% P12-ERR-0517
$\dim_\delta(D_{n, i}, D_{n, i'}) \leq d - 2$，正如所需。

\medskip\noindent
下面转向不再假设概形 $X$ 拟紧的一般情形。证明第一部分的思路存在一个问题：
可能得到一列无穷吹起，其中心支配 $X$ 的某个固定点。为避免这一点，
我们在每一步删去适当的余维 $\geq 3$ 的闭子集。具体而言，
将归纳构造如下形状的一列态射：
$$
\xymatrix{
X = X_0 \\
U_0 \ar[u]^{j_0} & X_1 \ar[l]_{b_0} \\
 & U_1 \ar[u]^{j_1} & X_2 \ar[l]_{b_1} \\
 & & U_2 \ar[u]^{j_2} & X_3 \ar[l]_{b_2}
}
$$
每个态射 $j_n : U_n \to X_n$ 都是开浸入。每个态射
$b_n : X_{n + 1} \to U_n$ 都是整概形之间的固有双有理态射。
同拟紧情形一样，在 $X_n$ 上有有效 Cartier 因子
$\{D_{n, i}\}_{i \in I_n}$。在每一步，任意三重交
$D_{n, i} \cap D_{n, j} \cap D_{n, k}$ 都为空。此外，对每个
$n \geq 0$，有 $I_{n + 1} = I_n \amalg P(I_n)$，其中 $P(I_n)$
表示 $I_n$ 中元素对的集合。最后，将使得
$$
b_n^{-1}(D_{n, i}|_{U_n}) = D_{n + 1, i} +
\sum\nolimits_{i' \in I_n, i' \not = i} D_{n + 1, \{i, i'\}}
$$

\medskip\noindent
归纳从置 $X_0 = X$、$I_0 = I$、$D_{0, i} = D_i$ 开始。

\medskip\noindent
给定 $(X_n, \{D_{n, i}\})$ 后，按如下方式构造开子概形 $U_n$。
对每个指标对 $\{i, i'\} \in P(I_n)$，考虑闭子概形
$D_{n, i} \cap D_{n, i'}$。它有 $\delta$-维数为 $d - 2$ 的“良”不可约分支，
以及 $\delta$-维数为 $d - 1$ 的“坏”不可约分支。置
$$
\text{Bad}(i, i')
=
\bigcup\nolimits_{W \subset D_{n, i} \cap D_{n, i'}
\text{ irred.\ comp. with }\dim_\delta(W) = d - 1} W
$$
，类似地置
$$
\text{Good}(i, i')
=
\bigcup\nolimits_{W \subset D_{n, i} \cap D_{n, i'}
\text{ irred.\ comp. with }\dim_\delta(W) = d - 2} W.
$$
于是 $D_{n, i} \cap D_{n, i'} = \text{Bad}(i, i') \cup \text{Good}(i, i')$，
并且
$\dim_\delta(\text{Bad}(i, i') \cap \text{Good}(i, i')) \leq d - 3$。
我们选取 $U_n$ 如下：
$$
U_n
=
X_n
\setminus
\bigcup\nolimits_{\{i, i'\} \in P(I_n)}
\text{Bad}(i, i') \cap \text{Good}(i, i').
$$
由各因子 $D_{n, i}$ 的三重交条件可见，该并实际上是不交并。
此外，作为概形有
$$
D_{n, i}|_{U_n} \cap D_{n, i'}|_{U_n}
=
Z_{n, i, i'} \amalg G_{n, i, i'}
$$
，其中 $Z_{n, i, i'}$ 为 $\delta$-等维的，维数为 $d - 1$；
$G_{n, i, i'}$ 为 $\delta$-等维的，维数为 $d - 2$。
（因此从拓扑上看，$Z_{n, i, i'}$ 是坏分支之并去掉它们同良分支的交。）
最后置
$$
Z_n =
\bigcup\nolimits_{\{i, i'\} \in P(I_n)} Z_{n, i, i'} =
\coprod\nolimits_{\{i, i'\} \in P(I_n)} Z_{n, i, i'},
$$
% P12-ERR-0518
，并令 $b_n : X_{n + 1} \to X_n$ 为沿 $Z_n$ 的吹起。注意，在每个集合
$D_{n, i}|_{U_n} \cap D_{n, i'}|_{U_n}$ 的邻域中，引理
\ref{obsolete-lemma-two-divisors} 都适用于态射 $b_n : X_{n + 1} \to X_n$。
因此，完全如证明第一部分那样，对 $\{i, i'\} \in P(I_n)$ 得到有效
Cartier 因子 $D_{n + 1, \{i, i'\}}$，并对 $i \in I_n$
得到有效 Cartier 因子 $D_{n + 1, i}$，使得
$b_n^{-1}(D_{n, i}|_{U_n}) = D_{n + 1, i} +
\sum\nolimits_{i' \in I_n, i' \not = i} D_{n + 1, \{i, i'\}}$.
对每个 $n$，以 $\pi_n : X_n \to X$ 表示复合
$j_0 \circ \ldots \circ j_{n - 1} \circ b_{n - 1}$ 所得的态射。

\medskip\noindent
{\bf 断言：}给定任意拟紧开集 $V \subset X$，则对所有充分大的 $n$，态射
$$
\pi_n^{-1}(V) \leftarrow \pi_{n + 1}^{-1}(V) \leftarrow \ldots
$$
全为同构。事实上，若态射
$\pi_n^{-1}(V) \leftarrow \pi_{n + 1}^{-1}(V)$ 不是同构，则对某个
$\{i, i'\} \in P(I_n)$，有
$Z_{n, i, i'} \cap \pi_n^{-1}(V) \not = \emptyset$。因此存在不可约分支
$W \subset D_{n, i} \cap D_{n, i'}$，满足 $\dim_\delta(W) = d - 1$。
特别地，$\epsilon_W(D_{n, i}, D_{n, i'}) > 0$。反复应用引理
\ref{obsolete-lemma-two-divisors} 可见
$$
\epsilon_W(D_{n, i}, D_{n, i'})
<
\epsilon(V, \{D_i|_V\}) - n
$$
，其中 $\epsilon(V, \{D_i|_V\})$ 如 (\ref{obsolete-equation-invariant}) 所定义。
由于 $V$ 拟紧，有 $\epsilon(V, \{D_i|_V\}) < \infty$；取
$n > \epsilon(V, \{D_i|_V\})$ 即得结论。

\medskip\noindent
注意，差集 $X_n \setminus U_n$ 由构造满足
$\dim_\delta(X_n \setminus U_n) \leq d - 3$。
令 $T_n = \pi_n(X_n \setminus U_n)$ 为其在 $X$ 中的像。
沿上述态射图逐步考察，并利用每个 $b_n$ 都是闭态射，可知对每个 $n$，
$T_0 \cup \ldots \cup T_n$ 都是 $X$ 的闭子集。由构造，任意
$t \in T_n$ 均满足 $\delta(t) \leq d - 3$。因此
$\overline{T_n} \subset X$ 是满足 $\dim_\delta(T_n) \leq d - 3$ 的闭子集。
由上述断言，对任意拟紧开集 $V \subset X$，仅有有限多个 $n$ 满足
$T_n \cap V \not = \emptyset$。故 $\{\overline{T_n}\}_{n \geq 0}$
是一族局部有限的闭子集，可以置
$U = X \setminus \bigcup \overline{T_n}$。这就是引理中的 $U$。

\medskip\noindent
由 $U$ 的构造，$U_n \cap \pi_n^{-1}(U) = \pi_n^{-1}(U)$。因此所有态射
$$
b_n : \pi_{n + 1}^{-1}(U) \longrightarrow \pi_n^{-1}(U)
$$
均为固有态射。此外，由上述断言，它们在 $X$ 的每个拟紧开集上最终成为同构。
因此可以定义
$$
U' = \lim_n \pi_n^{-1}(U).
$$
诱导态射 $b : U' \to U$ 是固有的，因为此性质在 $U$ 上是局部的，
而在每个拟紧开集上该极限都稳定。类似地，利用构造中的嵌入
$I_n \to I_{n + 1}$，置 $J = \bigcup_{n \geq 0} I_n$。
对 $j \in J$，选取 $n_0$ 使 $j$ 对应于某个 $i \in I_{n_0}$，并定义
$D'_j = \lim_{n \geq n_0} D_{n, i}$。这同样有意义，因为态射在 $X$ 上
局部稳定。引理的其余断言与 $X$ 拟紧的情形同样验证。
\end{proof}









\section{相交因子的交换性}
\label{obsolete-section-commutativity}

\noindent
本节结果原本用于给出 Chow Homology 小节
\ref{chow-section-commutativity} 中各引理的另一种证明，
以及 Chow Homology 的引理 \ref{chow-lemma-gysin-commutes-gysin}
的一个弱版本。

\begin{lemma}
\label{obsolete-lemma-improved-additivity}
设 $(S, \delta)$ 如 Chow Homology 的情形 \ref{chow-situation-setup}。
设 $X$ 在 $S$ 上局部有限型，且
$\{i_j : D_j \to X \}_{j \in J}$ 是 $X$ 上一族局部有限的有效
Cartier 因子。对 $j\in J$，设 $n_j > 0$。置
$D = \sum_{j \in J} n_j D_j$，并以 $i : D \to X$ 表示包含态射。
设 $\alpha \in Z_{k + 1}(X)$。则
$$
p : \coprod\nolimits_{j \in J} D_j \longrightarrow D
$$
为固有态射，且
$$
i^*\alpha = p_*\left(\sum n_j i_j^*\alpha\right)
$$
在 $\CH_k(D)$ 中成立。
\end{lemma}

\begin{proof}
由于无穷个有理等价之和涉及一个细微问题，本引理的证明比预期稍长。
在拟紧情形中，族 $D_j$ 是有限的，结论十分容易，直接由 Chow Homology
的引理 \ref{chow-lemma-compute-c1}、Divisors 的引理
\ref{divisors-lemma-c1-additive} 以及各项定义推出。

\medskip\noindent
由于族 $\{D_j\}_{j \in J}$ 局部有限，态射 $p$ 为固有态射。写成
$\alpha = \sum_{a \in A} m_a [W_a]$，其中 $W_a \subset X$
为 $\delta$-维数 $k + 1$ 的整闭子概形。以 $i_a : W_a \to X$
表示闭浸入。我们假设对所有 $a \in A$，有 $m_a \not = 0$，
且 $\{W_a\}_{a \in A}$ 在 $X$ 上局部有限。

\medskip\noindent
注意，由 Chow Homology 的定义 \ref{chow-definition-gysin-homomorphism}，
类 $i^*\alpha$ 是某个循环 $\sum m_a\beta_a$ 的类，
其中 $\beta_a \in Z_k(W_a \cap D)$。具体而言，若
$W_a \not \subset D$，则 $\beta_a = [D \cap W_a]_k$；若
$W_a \subset D$，则 $\beta_a$ 是表示
$c_1(\mathcal{O}_X(D)) \cap [W_a]$ 的循环。

\medskip\noindent
对每个 $a \in A$，写成 $J = J_{a, 1} \amalg J_{a, 2} \amalg J_{a, 3}$，其中
\begin{enumerate}
\item $j \in J_{a, 1}$ 当且仅当 $W_a \cap D_j = \emptyset$；
% P12-ERR-0519
\item $j \in J_{a, 2}$ 当且仅当
$W_a \not = W_a \cap D_1 \not = \emptyset$；
\item $j \in J_{a, 3}$ 当且仅当 $W_a \subset D_j$。
\end{enumerate}
由于族 $\{D_j\}$ 局部有限，$J_{a, 3}$ 是有限集。
对每个 $a \in A$ 与 $j \in J$，按如下方式选取循环
$\beta_{a, j} \in Z_k(W_a \cap D_j)$：
\begin{enumerate}
\item 若 $j \in J_{a, 1}$，置 $\beta_{a, j} = 0$；
\item 若 $j \in J_{a, 2}$，置 $\beta_{a, j} = [D_j \cap W_a]_k$；
\item 若 $j \in J_{a, 3}$，选取 $\beta_{a, j} \in Z_k(W_a)$，
% P12-ERR-0520
使其表示 $c_1(i_a^*\mathcal{O}_X(D_j)) \cap [W_j]$。
\end{enumerate}
我们断言
$$
\beta_a \sim_{rat}
\sum\nolimits_{j \in J} n_j \beta_{a, j}
$$
在 $\CH_k(W_a \cap D)$ 中成立。

\medskip\noindent
情形 I：$W_a \not \subset D$。此时 $J_{a, 3} = \emptyset$。
所以只需证明循环等式
$[D \cap W_a]_k = \sum n_j [D_j \cap W_a]_k$。
这正是引理 \ref{obsolete-lemma-sum-divisors-associated-Weil}。

\medskip\noindent
情形 II：$W_a \subset D$。此时 $\beta_a$ 是表示
$c_1(i_a^*\mathcal{O}_X(D)) \cap [W_a]$ 的循环。写成
$D = D_{a, 1} + D_{a, 2} + D_{a, 3}$，其中
$D_{a, s} = \sum_{j \in J_{a, s}} n_jD_j$。
由 Divisors 的引理 \ref{divisors-lemma-c1-additive}，有
\begin{eqnarray*}
c_1(i_a^*\mathcal{O}_X(D)) \cap [W_a] & = &
c_1(i_a^*\mathcal{O}_X(D_{a, 1})) \cap [W_a] +
c_1(i_a^*\mathcal{O}_X(D_{a, 2})) \cap [W_a] \\
& &
 + c_1(i_a^*\mathcal{O}_X(D_{a, 3})) \cap [W_a].
\end{eqnarray*}
显然，和式第一项为零。由于 $J_{a, 3}$ 有限，最后一项等于
% P12-ERR-0521
$\sum\nolimits_{j \in J_{a, 3}} n_jc_1(i_a^*\mathcal{L}_j) \cap [W_a]$；
见 Divisors 的引理 \ref{divisors-lemma-c1-additive}。
它由 $\sum_{j \in J_{a, 3}} n_j \beta_{a, j}$ 表示。
最后，由情形 I，中间一项由循环
$\sum\nolimits_{j \in J_{a, 2}} n_j[D_j \cap W_a]_k =
\sum_{j \in J_{a, 2}} n_j\beta_{a, j}$ 表示。
因此此情形下断言成立。

\medskip\noindent
现在可以完成引理的证明。由 $\beta_a$ 的选取，有
% P12-ERR-0522
$i^*D \sim_{rat} \sum m_a\beta_a$。对每个 $a$，有
% P12-ERR-0523
$\beta_a \sim_{rat} \sum_j \beta_{a, j}$，该有理等价发生在
$D \cap W_a$ 上。由于各闭子概形 $D \cap W_a$ 在 $D$ 上局部有限，
还可在 $D$ 上得到
$\sum m_a \beta_a \sim_{rat} \sum_{a, j} m_a\beta_{a, j}$！
（见 Chow Homology 的注记
\ref{chow-remark-infinite-sums-rational-equivalences}。）
最后，$\sum_a m_a\beta_{a, j}$（视为 $D_j$ 上的循环）显然表示
$i_j^*\alpha$，故 $\sum_{a, j} m_a\beta_{a, j}$ 表示
$p_* \sum_j i_j^*\alpha$，结论得证。
\end{proof}

\begin{lemma}
\label{obsolete-lemma-commutativity-effective-Cartier-proper-intersection}
设 $(S, \delta)$ 如 Chow Homology 的情形 \ref{chow-situation-setup}。
设 $X$ 在 $S$ 上局部有限型，并假设 $X$ 整且 $\dim_\delta(X) = n$。
设 $D$、$D'$ 为 $X$ 上的有效 Cartier 因子，并假设
$\dim_\delta(D \cap D') = n - 2$。以 $i : D \to X$
（resp.\ $i' : D' \to X$）表示相应的闭浸入。则
\begin{enumerate}
\item 存在循环 $\alpha \in Z_{n - 2}(D \cap D')$，其到 $D$ 的前推表示
$i^*[D']_{n - 1} \in \CH_{n - 2}(D)$，而其到 $D'$ 的前推表示
$(i')^*[D]_{n - 1} \in \CH_{n - 2}(D')$；
\item 有
$$
D \cdot [D']_{n - 1}
=
D' \cdot [D]_{n - 1}
$$
在 $\CH_{n - 2}(X)$ 中成立。
\end{enumerate}
\end{lemma}

\begin{proof}
第 (2) 部分是第 (1) 部分的直接推论。写成
$[D]_{n - 1} = \sum n_a[Z_a]$ 与
$[D']_{n - 1} = \sum m_b[Z_b]$，其中 $Z_a$ 是 $D$ 的不可约分支，
$[Z_b]$ 是 $D'$ 的不可约分支。由 Chow Homology 的定义
\ref{chow-definition-gysin-homomorphism}，有
$i^*D' = \sum m_b i^*[Z_b]$ 以及
$(i')^*D = \sum n_a(i')^*[Z_a]$。由假设，没有不可约分支 $Z_b$
包含于 $D$，故由定义 $i^*[Z_b] = [Z_b\cap D]_{n - 2}$。
同理，$(i')^*[Z_a] = [Z_a \cap D']_{n - 2}$。因此要证明循环等式
$$
\sum n_a[Z_a \cap D']_{n - 2} = \sum m_b[Z_b \cap D]_{n - 2}
$$
；二者确实都支撑在 $D \cap D'$ 上。设 $W \subset X$ 为满足
$\dim_\delta(W) = n - 2$ 的整闭子概形，并设 $\xi \in W$ 为其一般点。
置 $R = \mathcal{O}_{X, \xi}$；这是诺特局部整环。注意
$\dim(R) = 2$。设 $f \in R$（resp.\ $f' \in R$）为定义 $D$
（resp.\ $D'$）之理想的元素。由假设，$\dim(R/(f, f')) = 0$。
设 $\mathfrak q'_1, \ldots, \mathfrak q'_t \subset R$ 为 $(f')$ 上方的
极小素理想，$\mathfrak q_1, \ldots, \mathfrak q_s \subset R$ 为
$(f)$ 上方的极小素理想。上述等式归结为
$$
\sum_{i = 1, \ldots, s}
\text{length}_{R_{\mathfrak q_i}}(R_{\mathfrak q_i}/(f))
\text{ord}_{R/\mathfrak q_i}(f')
=
\sum_{j = 1, \ldots, t}
\text{length}_{R_{\mathfrak q'_j}}(R_{\mathfrak q'_j}/(f'))
\text{ord}_{R/\mathfrak q'_j}(f).
$$
将 Chow Homology 的引理 \ref{chow-lemma-length-multiplication}
应用于 $M = R/(f)$，该等式左边等于
$$
\text{length}_R(R/(f, f'))
-
\text{length}_R(\Ker(f' : R/(f) \to R/(f)))
$$
。注意，$\Ker(f' : R/(f) \to R/(f))$ 通过态射
$x \bmod (f) \mapsto f'x \bmod (ff')$ 典范同构于
$((f) \cap (f'))/(ff')$。所以左边为
$$
\text{length}_R(R/(f, f'))
-
\text{length}_R((f) \cap (f')/(ff'))
$$
。由于该表达式关于 $f$ 与 $f'$ 对称，结论得证。
\end{proof}

\begin{lemma}
\label{obsolete-lemma-commutativity-effective-Cartier-proper-intersection-infinite}
设 $(S, \delta)$ 如 Chow Homology 的情形 \ref{chow-situation-setup}。
设 $X$ 在 $S$ 上局部有限型，并假设 $X$ 整且 $\dim_\delta(X) = n$。
设 $\{D_j\}_{j \in J}$ 为 $X$ 上一族局部有限的有效 Cartier 因子，
$n_j, m_j \geq 0$ 为两族非负整数。置
$D = \sum n_j D_j$ 与 $D' = \sum m_j D_j$。假设对每个
$j \not = j'$ 都有 $\dim_\delta(D_j \cap D_{j'}) = n - 2$。
则在 $\CH_{n - 2}(X)$ 中，
$D \cdot [D']_{n - 1} = D' \cdot [D]_{n - 1}$。
\end{lemma}

\begin{proof}
当各和式有限时，例如 $X$ 拟紧时，本引理是引理
\ref{obsolete-lemma-sum-divisors-associated-Weil} 与
\ref{obsolete-lemma-commutativity-effective-Cartier-proper-intersection}
的直接推论。因此建议读者略过本证明。

\medskip\noindent
下面给出一般情形的证明。设 $i_j : D_j \to X$ 为闭浸入，
以 $p : \coprod D_j \to X$ 表示各态射 $i_j$ 的余积。
设 $\{Z_a\}_{a \in A}$ 为 $\bigcup D_j$ 的不可约分支所成的族。
对每个 $j$，写成
$$
[D_j]_{n - 1} = \sum d_{j, a}[Z_a].
$$
由引理 \ref{obsolete-lemma-sum-divisors-associated-Weil}，有
$$
[D]_{n - 1} = \sum n_j d_{j, a} [Z_a],
\quad
[D']_{n - 1} = \sum m_j d_{j, a} [Z_a].
$$
由引理 \ref{obsolete-lemma-improved-additivity}，有
$$
D \cdot [D']_{n - 1} = p_*\left(\sum n_j i_j^*[D']_{n - 1} \right),
\quad
D' \cdot [D]_{n - 1} = p_*\left(\sum m_{j'} i_{j'}^*[D]_{n - 1} \right).
$$
如 Gysin 同态的定义（见 Chow Homology 的定义
\ref{chow-definition-gysin-homomorphism}）那样，选取
$D_j \cap Z_a$ 上表示 $i_j^*[Z_a]$ 的循环 $\beta_{a, j}$。
（注意，若 $Z_a$ 不包含于 $D_j$，则事实上
$\beta_{a, j} = [D_j \cap Z_a]_{n - 2}$，即此时并无选择。）
由于 $p$ 限制在每个 $D_j$ 上都是闭浸入，可以而且将把
$\beta_{a, j}$ 视为 $X$ 上的循环。代入上面所得
$[D]_{n - 1}$ 与 $[D']_{n - 1}$ 的公式，得到
$$
D \cdot [D']_{n - 1} =
\sum\nolimits_{j, j', a} n_j m_{j'} d_{j', a} \beta_{a, j},
\quad
D' \cdot [D]_{n - 1} =
\sum\nolimits_{j, j', a} m_{j'} n_j d_{j, a} \beta_{a, j'}.
$$
此外，在同样约定下，还有
$$
D_j \cdot [D_{j'}]_{n - 1} = \sum d_{j', a} \beta_{a, j}.
$$
用这些记号，引理
\ref{obsolete-lemma-commutativity-effective-Cartier-proper-intersection}
（亦见其证明）表明：当 $j \not = j'$ 时，循环
$\sum d_{j', a} \beta_{a, j}$ 与 $\sum d_{j, a} \beta_{a, j'}$
作为循环相等！因此
\begin{eqnarray*}
D \cdot [D']_{n - 1}
& = &
\sum\nolimits_{j, j', a} n_j m_{j'} d_{j', a} \beta_{a, j} \\
& = &
\sum\nolimits_{j \not = j'} n_j m_{j'}
\left(\sum\nolimits_a d_{j', a} \beta_{a, j}\right) +
\sum\nolimits_{j, a} n_j m_j d_{j, a} \beta_{a, j} \\
& = &
\sum\nolimits_{j \not = j'} n_j m_{j'}
\left(\sum\nolimits_a d_{j, a} \beta_{a, j'}\right) +
\sum\nolimits_{j, a} n_j m_j d_{j, a} \beta_{a, j} \\
& = &
\sum\nolimits_{j, j', a} m_{j'} n_j d_{j, a} \beta_{a, j'} \\
& = &
D' \cdot [D]_{n - 1}
\end{eqnarray*}
，结论得证。
\end{proof}

\begin{lemma}
\label{obsolete-lemma-commutativity-effective-Cartier}
设 $(S, \delta)$ 如 Chow Homology 的情形 \ref{chow-situation-setup}。
设 $X$ 在 $S$ 上局部有限型，并假设 $X$ 整且 $\dim_\delta(X) = n$。
设 $D$、$D'$ 为 $X$ 上的有效 Cartier 因子。则
$$
D \cdot [D']_{n - 1} = D' \cdot [D]_{n - 1}
$$
在 $\CH_{n - 2}(X)$ 中成立。
\end{lemma}

\begin{proof}[引理 \ref{obsolete-lemma-commutativity-effective-Cartier} 的第一种证明]
先证明 $X$ 拟紧的情形。此时，对 $X$ 和由有效 Cartier 因子构成的
二元集合 $\{D, D'\}$ 应用引理 \ref{obsolete-lemma-blowing-up-intersections}。
于是得到固有态射 $b : X' \to X$，以及一族有限的有效 Cartier 因子
$D'_j \subset X'$；它们两两以余维 $\geq 2$ 相交，并且
$b^{-1}(D) = \sum n_j D'_j$、$b^{-1}(D') = \sum m_j D'_j$。
注意，在 $Z_{n - 1}(X)$ 中
$b_*[b^{-1}(D)]_{n - 1} = [D]_{n - 1}$，对 $D'$ 亦然；
见引理 \ref{obsolete-lemma-push-pull-effective-Cartier}。因此由 Chow Homology
的引理 \ref{chow-lemma-pushforward-cap-c1}，有
$$
D \cdot [D']_{n - 1} = b_*\left(b^{-1}(D) \cdot [b^{-1}(D')]_{n - 1}\right)
$$
在 $\CH_{n - 2}(X)$ 中成立；另一项亦然。因此本引理由等式
$b^{-1}(D) \cdot [b^{-1}(D')]_{n - 1} = b^{-1}(D') \cdot [b^{-1}(D)]_{n - 1}$
推出；该等式由引理
\ref{obsolete-lemma-commutativity-effective-Cartier-proper-intersection-infinite}
在 $\CH_{n - 2}(X')$ 中给出。

\medskip\noindent
注意，上述证明中引用的每个引理在一般情形下也成立（此时不假设
$X$ 拟紧）。一般情形唯一的细小变化是：应用引理
\ref{obsolete-lemma-blowing-up-intersections} 得到的态射
$b : U' \to U$，其目标是开集 $U \subset X$，且该开集的补集余维
$\geq 3$。因此由 Chow Homology 的引理
\ref{chow-lemma-restrict-to-open}，有
$\CH_{n - 2}(U) = \CH_{n - 2}(X)$；以 $U$ 替换 $X$ 后，
证明余下部分原封不动地成立。
\end{proof}

\begin{proof}[引理 \ref{obsolete-lemma-commutativity-effective-Cartier} 的第二种证明]
设 $\mathcal{I} = \mathcal{O}_X(-D)$ 与
$\mathcal{I}' = \mathcal{O}_X(-D')$ 分别为 $D$ 与 $D'$ 的可逆理想层。
记
$\mathcal{I}_{D'} = \mathcal{I} \otimes_{\mathcal{O}_X} \mathcal{O}_{D'}$
以及
$\mathcal{I}'_D = \mathcal{I}' \otimes_{\mathcal{O}_X} \mathcal{O}_D$。
将包含态射 $\mathcal{I} \to \mathcal{O}_X$ 限制到 $D'$，得到态射
$$
\varphi : \mathcal{I}_{D'} \longrightarrow \mathcal{O}_{D'}
$$
；类似地得到
$$
\psi : \mathcal{I}'_D \longrightarrow \mathcal{O}_D
$$
。显然
$$
\Coker(\varphi)
\cong
\mathcal{O}_{D \cap D'}
\cong
\Coker(\psi)
$$
，并且
$$
\Ker(\varphi)
\cong
\frac{\mathcal{I} \cap \mathcal{I}'}{\mathcal{I}\mathcal{I}'}
\cong
\Ker(\psi).
$$
因此可见
$$
\gamma =
[\mathcal{I}_{D'}] - [\mathcal{O}_{D'}]
=
[\mathcal{I}'_D] - [\mathcal{O}_D]
$$
在 $K_0(\textit{Coh}_{\leq n - 1}(X))$ 中成立。另一方面，显然
$$
[\mathcal{I}'_D]_{n - 1} = [D]_{n - 1}, \quad
[\mathcal{I}_{D'}]_{n - 1} = [D']_{n - 1}.
$$
并且
$$
\mathcal{O}_X(D') \otimes \mathcal{I}'_D = \mathcal{O}_D, \quad
\mathcal{O}_X(D) \otimes \mathcal{I}_{D'} = \mathcal{O}_{D'}.
$$
将 Chow Homology 的引理 \ref{chow-lemma-coherent-sheaf-cap-c1}
应用两次，这意味着元素 $\gamma$ 属于 $B_{n - 2}(X)$，并同时映到
$c_1(\mathcal{O}_X(D')) \cap [D]_{n - 1}$ 与
$c_1(\mathcal{O}_X(D)) \cap [D']_{n - 1}$，结论得证
（因为态射 $B_{n - 2}(X) \to \CH_{n - 2}(X)$ 是良定义的——
这正是本证明的关键）。
\end{proof}











\section{正则固有模型上的对偶化模}
\label{obsolete-section-dualizing}

\noindent
在 Semistable Reduction 的情形 \ref{models-situation-regular-model} 中，
令 $\omega_{X/R}^\bullet = f^!\mathcal{O}_{\Spec(R)}$ 为态射
$f : X \to \Spec(R)$ 的相对对偶化复形；见 Duality for Schemes 的注记
\ref{duality-remark-relative-dualizing-complex}。由 Semistable Reduction
的引理 \ref{models-lemma-gorenstein}，$f$ 是相对维数 $1$ 的戈伦斯坦态射，
故可以使用 Duality for Schemes 的引理
\ref{duality-lemma-affine-flat-Noetherian-gorenstein},
\ref{duality-lemma-CM-shriek} 与
\ref{duality-lemma-gorenstein-CM-morphism}
得到
$$
\omega_{X/R}^\bullet = \omega_X[1]
$$
，其中 $\omega_X$ 是某个可逆 $\mathcal{O}_X$-模。这个可逆模通常称为
$X$ 在 $R$ 上的{\it 相对对偶化模}。
由于 $R$ 正则（因而戈伦斯坦）且维数为 $1$，
$\omega_R^\bullet = R[1]$ 是 $R$ 的归一化对偶化复形。所以
$\omega_X = H^{-2}(f^!\omega_R^\bullet)$；由此可见，$\omega_X$
不仅是相对对偶化模，也是对偶化模，见 Duality for Schemes 的例
\ref{duality-example-proper-over-local}。因此 $\omega_X$ 表示函子
$$
\textit{Coh}(\mathcal{O}_X) \to \textit{Sets},\quad
\mathcal{F} \mapsto \Hom_R(H^1(X, \mathcal{F}), R)
$$
，这是 Duality for Schemes 的引理
\ref{duality-lemma-dualizing-module-proper-over-A}。这给出了
Semistable Reduction 的情形 \ref{models-situation-regular-model}
中相对对偶化模的另一种定义。$\omega_X$ 的形成与任意基变换可换
（对任意给定相对维数的固有戈伦斯坦态射均如此）；这由相对对偶化复形
的相应事实推出，该事实见 Duality for Schemes 的注记
\ref{duality-remark-relative-dualizing-complex}，其源自 Duality for Schemes
的引理 \ref{duality-lemma-proper-flat-base-change}。因此 $\omega_X$
拉回为 Algebraic Curves 的引理 \ref{curves-lemma-duality-dim-1-CM}
中讨论的 $C$ 在 $K$ 上的对偶化模 $\omega_C$。注意，由 Algebraic Curves
的引理 \ref{curves-lemma-duality-dim-1}，$\omega_C$ 同构于
$\Omega_{C/K}$。类似地，$\omega_X|_{X_k}$ 是 $X_k$ 在 $k$ 上的
对偶化模 $\omega_{X_k}$。

\begin{lemma}
\label{obsolete-lemma-dualizing-components}
在 Semistable Reduction 的情形 \ref{models-situation-regular-model} 中，
$C_i$ 在 $k$ 上的对偶化模为
$$
\omega_{C_i} = \omega_X(C_i)|_{C_i}
$$
，其中 $\omega_X$ 如上。
\end{lemma}

\begin{proof}
设 $t : C_i \to X$ 为闭浸入。由于 $t$ 是有效 Cartier 因子的包含，
由 Duality for Schemes 的引理
\ref{duality-lemma-twisted-inverse-image-closed} 与
\ref{duality-lemma-sheaf-with-exact-support-effective-Cartier} 可知，
对每个可逆 $\mathcal{O}_X$-模 $\mathcal{L}$，都有
$t^!(\mathcal{L}) = \mathcal{L}(C_i)|_{C_i}$。考虑交换图
$$
\xymatrix{
C_i \ar[r]_t \ar[d]_g & X \ar[d]^f \\
\Spec(k) \ar[r]^s & \Spec(R)
}
$$
。注意，$C_i$ 是戈伦斯坦曲线（Semistable Reduction 的引理
\ref{models-lemma-gorenstein}），其可逆对偶化模 $\omega_{C_i}$
由性质 $\omega_{C_i}[0] = g^!\mathcal{O}_{\Spec(k)}$ 刻画。
参见 Algebraic Curves 的引理 \ref{curves-lemma-duality-dim-1} 及其证明，
以及 Algebraic Curves 的引理 \ref{curves-lemma-duality-dim-1-CM}
与 \ref{curves-lemma-rr}。另一方面，$s^!(R[1]) = k$，故
$$
\omega_{C_i}[0] =
g^! s^!(R[1]) = t^!f^!(R[1]) = t^!\omega_X
$$
。结合以上各式即得引理结论。
\end{proof}






\section{重复项与拆分出的引用}
\label{obsolete-section-duplicates}

\noindent
本节收集重复项，以及那些过去同时陈述若干件事、如今已拆分为各组成部分的引用。

\begin{lemma}
\label{obsolete-lemma-characterize-epi-mono}
设 $X$ 为拓扑空间。
\begin{enumerate}
\item 预层范畴中的满态射（resp.\ 单态射）恰好是预层的满射
（resp.\ 单射）。
\item 层范畴中的满态射（resp.\ 单态射）恰好是层的满射
（resp.\ 单射），也恰好是在所有茎上均为满射（resp.\ 单射）的态射。
\item 集合预层的满射（resp.\ 单射）态射的层化仍为满射（resp.\ 单射）。
\end{enumerate}
\end{lemma}

\begin{proof}
合并 Sheaves 的引理 \ref{sheaves-lemma-characterize-epi-mono}
与 \ref{sheaves-lemma-sheafify-epi-mono} 即得。
\end{proof}

\begin{lemma}
\label{obsolete-lemma-directed-colimit-finite-type}
设 $X$ 为概形，并假设 $X$ 拟紧且拟分离。
设 $\mathcal{F}$ 为拟凝聚 $\mathcal{O}_X$-模。
则 $\mathcal{F}$ 是其有限型拟凝聚子模的滤过余极限。
\end{lemma}

\begin{proof}
这是 Properties 的引理
\ref{properties-lemma-quasi-coherent-colimit-finite-type} 的重复项。
\end{proof}

\begin{lemma}
\label{obsolete-lemma-points-monomorphism}
设 $S$ 为概形，$X$ 为 $S$ 上的代数空间。态射
$\{\Spec(k) \to X \text{ monomorphism}\} \to |X|$ 为单射。
\end{lemma}

\begin{proof}
这是 Properties of Spaces 的引理
\ref{spaces-properties-lemma-points-monomorphism} 的重复项。
\end{proof}

\begin{theorem}
\label{obsolete-theorem-equivalence-sheaves-point}
设 $S = \Spec(K)$，其中 $K$ 为域。设 $\overline{s}$ 为 $S$ 的几何点，
$G = \text{Gal}_{\kappa(s)}$ 表示绝对伽罗瓦群。则存在范畴等价
$\Sh(S_\etale) \to G\textit{-Sets}$,
$\mathcal{F} \mapsto \mathcal{F}_{\overline{s}}$.
\end{theorem}

\begin{proof}
这是 \'Etale Cohomology 的定理
\ref{etale-cohomology-theorem-equivalence-sheaves-point} 的重复项。
\end{proof}

\begin{remark}
\label{obsolete-remark-tangent-spaces}
你来到这里是因为标签重复。实际内容请见 Formal Deformation Theory
的小节 \ref{formal-defos-section-tangent-spaces}。
\end{remark}

\begin{lemma}
\label{obsolete-lemma-locally-ringed-space-direct-summand-free}
设 $X$ 为局部环化空间。有限自由 $\mathcal{O}_X$-模的直和因子有限局部自由。
\end{lemma}

\begin{proof}
这是 Modules 的引理
\ref{modules-lemma-direct-summand-of-locally-free-is-locally-free} 的重复项。
\end{proof}

\begin{lemma}
\label{obsolete-lemma-characterize-injective}
设 $R$ 为环，$E$ 为 $R$-模。下列条件等价：
\begin{enumerate}
\item $E$ 是内射 $R$-模；
\item 给定理想 $I \subset R$ 和模态射 $\varphi : I \to E$，
存在 $\varphi$ 到 $R$-模态射 $R \to E$ 的延拓。
\end{enumerate}
\end{lemma}

\begin{proof}
这是 Baer 判据；见 Injectives 的引理 \ref{injectives-lemma-criterion-baer}。
\end{proof}

\begin{lemma}
\label{obsolete-lemma-periodic-length}
设 $R$ 为局部环。
\begin{enumerate}
\item 若 $(M, N, \varphi, \psi)$ 是 $2$-周期复形，且 $M$、$N$ 长度有限，
则 $e_R(M, N, \varphi, \psi) = \text{length}_R(M) - \text{length}_R(N)$。
\item 若 $(M, \varphi, \psi)$ 是 $(2, 1)$-周期复形，且 $M$ 长度有限，
则 $e_R(M, \varphi, \psi) = 0$。
\item 假设有 $2$-周期复形的短正合列
$$
0 \to (M_1, N_1, \varphi_1, \psi_1)
\to (M_2, N_2, \varphi_2, \psi_2)
\to (M_3, N_3, \varphi_3, \psi_3)
\to 0
$$
。若三者中有两个的上同调模长度有限，则第三个也如此，并且有
$$
e_R(M_2, N_2, \varphi_2, \psi_2) =
e_R(M_1, N_1, \varphi_1, \psi_1) +
e_R(M_3, N_3, \varphi_3, \psi_3).
$$
\end{enumerate}
\end{lemma}

\begin{proof}
这由 Chow Homology 的引理 \ref{chow-lemma-additivity-periodic-length}
与 \ref{chow-lemma-finite-periodic-length} 推出。
\end{proof}


\begin{lemma}
\label{obsolete-lemma-extensions-of-rings}
设 $A$ 为环，$I$ 为 $A$-模。
\begin{enumerate}
\item 以 $I$ 为平方零理想的环扩张
$0 \to I \to A' \to A \to 0$ 所成集合，典范双射于
$\Ext^1_A(\NL_{A/\mathbf{Z}}, I)$。
\item 给定环同态 $A \to B$、$B$-模 $N$、$A$-模态射
$c : I \to N$，并给定下列核平方为零的环扩张：
\begin{enumerate}
\item[(a)] $0 \to I \to A' \to A \to 0$，对应于
$\alpha \in \Ext^1_A(\NL_{A/\mathbf{Z}}, I)$；
\item[(b)] $0 \to N \to B' \to B \to 0$，对应于
$\beta \in \Ext^1_B(\NL_{B/\mathbf{Z}}, N)$；
\end{enumerate}
则存在符合 Deformation Theory 的式 (\ref{defos-equation-to-solve})
的态射 $A' \to B'$，当且仅当 $\beta$ 与 $\alpha$ 映到
$\Ext^1_A(\NL_{A/\mathbf{Z}}, N)$ 中的同一元素。
\end{enumerate}
\end{lemma}

\begin{proof}
这由 Deformation Theory 的引理
\ref{defos-lemma-extensions-of-algebras} 与
\ref{defos-lemma-extensions-of-algebras-functorial} 推出。
\end{proof}

\begin{lemma}
\label{obsolete-lemma-extensions-of-ringed-spaces}
设 $(S, \mathcal{O}_S)$ 为环化空间，$\mathcal{J}$ 为 $\mathcal{O}_S$-模。
\begin{enumerate}
\item 以 $\mathcal{J}$ 为平方零理想的环层扩张
$0 \to \mathcal{J} \to \mathcal{O}_{S'} \to \mathcal{O}_S \to 0$
所成集合，典范双射于
$\Ext^1_{\mathcal{O}_S}(\NL_{S/\mathbf{Z}}, \mathcal{J})$。
\item 给定环化空间态射
$f : (X, \mathcal{O}_X) \to (S, \mathcal{O}_S)$、$\mathcal{O}_X$-模
$\mathcal{G}$、$f$-态射 $c : \mathcal{J} \to \mathcal{G}$，
并给定下列核平方为零的环层扩张：
\begin{enumerate}
\item[(a)] $0 \to \mathcal{J} \to \mathcal{O}_{S'} \to \mathcal{O}_S \to 0$
，对应于
$\alpha \in \Ext^1_{\mathcal{O}_S}(\NL_{S/\mathbf{Z}}, \mathcal{J})$；
\item[(b)] $0 \to \mathcal{G} \to \mathcal{O}_{X'} \to \mathcal{O}_X \to 0$
，对应于
$\beta \in \Ext^1_{\mathcal{O}_X}(\NL_{X/\mathbf{Z}}, \mathcal{G})$；
\end{enumerate}
则存在符合 Deformation Theory 的式
(\ref{defos-equation-to-solve-ringed-spaces}) 的态射 $X' \to S'$，
当且仅当 $\beta$ 与 $\alpha$ 映到
$\Ext^1_{\mathcal{O}_X}(Lf^*\NL_{S/\mathbf{Z}}, \mathcal{G})$
中的同一元素。
\end{enumerate}
\end{lemma}

\begin{proof}
这由 Deformation Theory 的引理
\ref{defos-lemma-extensions-of-relative-ringed-spaces} 与
\ref{defos-lemma-extensions-of-relative-ringed-spaces-functorial} 推出。
\end{proof}

\begin{lemma}
\label{obsolete-lemma-extensions-of-ringed-topoi}
设 $(\Sh(\mathcal{B}), \mathcal{O}_\mathcal{B})$ 为环化拓扑斯，
$\mathcal{J}$ 为 $\mathcal{O}_\mathcal{B}$-模。
\begin{enumerate}
\item 以 $\mathcal{J}$ 为平方零理想的环层扩张
$0 \to \mathcal{J} \to \mathcal{O}_{\mathcal{B}'} \to
\mathcal{O}_\mathcal{B} \to 0$
所成集合，典范双射于
$\Ext^1_{\mathcal{O}_\mathcal{B}}(
\NL_{\mathcal{O}_\mathcal{B}/\mathbf{Z}}, \mathcal{J})$.
\item 给定环化拓扑斯态射
$f : (\Sh(\mathcal{C}), \mathcal{O}) \to
(\Sh(\mathcal{B}), \mathcal{O}_\mathcal{B})$、$\mathcal{O}$-模
$\mathcal{G}$、$f^{-1}\mathcal{O}_\mathcal{B}$-模态射
$c : f^{-1}\mathcal{J} \to \mathcal{G}$，并给定下列核平方为零的环层扩张：
\begin{enumerate}
\item[(a)] $0 \to \mathcal{J} \to \mathcal{O}_{\mathcal{B}'} \to
\mathcal{O}_\mathcal{B} \to 0$，对应于
$\alpha \in \Ext^1_{\mathcal{O}_\mathcal{B}}(
\NL_{\mathcal{O}_\mathcal{B}/\mathbf{Z}}, \mathcal{J})$；
\item[(b)] $0 \to \mathcal{G} \to \mathcal{O}' \to \mathcal{O} \to 0$
，对应于
$\beta \in \Ext^1_\mathcal{O}(\NL_{\mathcal{O}/\mathbf{Z}}, \mathcal{G})$
\end{enumerate}
则存在符合 Deformation Theory 的式
(\ref{defos-equation-to-solve-ringed-topoi}) 的态射
$(\Sh(\mathcal{C}), \mathcal{O}') \to
(\Sh(\mathcal{B}, \mathcal{O}_{\mathcal{B}'})$，当且仅当
$\beta$ 与 $\alpha$ 映到下列 Ext 群中的同一元素：
$\Ext^1_\mathcal{O}(
Lf^*\NL_{\mathcal{O}_\mathcal{B}/\mathbf{Z}}, \mathcal{G})$.
\end{enumerate}
\end{lemma}

\begin{proof}
这由 Deformation Theory 的引理
\ref{defos-lemma-extensions-of-relative-ringed-topoi} 与
\ref{defos-lemma-extensions-of-relative-ringed-topoi-functorial} 推出。
\end{proof}

\begin{remark}
\label{obsolete-remark-examples-formal-defos}
这个标签过去指向一个描述若干形变问题例子的章节。现在每个例子都有自己的章节。
见 Deformation Problems 的小节
\ref{examples-defos-section-finite-projective-modules},
\ref{examples-defos-section-representations},
\ref{examples-defos-section-continuous-representations} 与
\ref{examples-defos-section-graded-algebras}.
\end{remark}

\begin{lemma}
\label{obsolete-lemma-examples-have-RS}
Deformation Problems 的例
\ref{examples-defos-example-finite-projective-modules},
\ref{examples-defos-example-representations},
\ref{examples-defos-example-continuous-representations} 与
\ref{examples-defos-example-graded-algebras}
满足 Rim--Schlessinger 条件 (RS)。
\end{lemma}

\begin{proof}
这由 Deformation Problems 的引理
\ref{examples-defos-lemma-finite-projective-modules-RS},
\ref{examples-defos-lemma-representations-RS},
\ref{examples-defos-lemma-continuous-representations-RS} 与
\ref{examples-defos-lemma-graded-algebras-RS} 推出。
\end{proof}

\begin{lemma}
\label{obsolete-lemma-tangent-and-inf}
有下列典范 $k$-向量空间同一化：
\begin{enumerate}
\item 在 Deformation Problems 的例
\ref{examples-defos-example-finite-projective-modules}
中，若 $x_0 = (k, V)$，则 $T_{x_0}\mathcal{F} = (0)$ 与
$\text{Inf}_{x_0}(\mathcal{F}) = \text{End}_k(V)$ 均为有限维。
\item 在 Deformation Problems 的例
\ref{examples-defos-example-representations}
中，若 $x_0 = (k, V, \rho_0)$，则
$T_{x_0}\mathcal{F} = \Ext^1_{k[\Gamma]}(V, V) = H^1(\Gamma, \text{End}_k(V))$
与 $\text{Inf}_{x_0}(\mathcal{F}) = H^0(\Gamma, \text{End}_k(V))$
在 $\Gamma$ 有限生成时均为有限维。
\item 在 Deformation Problems 的例
\ref{examples-defos-example-continuous-representations}
中，若 $x_0 = (k, V, \rho_0)$，则
$T_{x_0}\mathcal{F} = H^1_{cont}(\Gamma, \text{End}_k(V))$
且
$\text{Inf}_{x_0}(\mathcal{F}) = H^0_{cont}(\Gamma, \text{End}_k(V))$
在 $\Gamma$ 拓扑有限生成时均为有限维。
\item 在 Deformation Problems 的例
\ref{examples-defos-example-graded-algebras}
中，若 $x_0 = (k, P)$，则
$T_{x_0}\mathcal{F}$ 与 $\text{Inf}_{x_0}(\mathcal{F}) = \text{Der}_k(P, P)$
在 $P$ 为 $k$ 上有限生成时均为有限维。
\end{enumerate}
\end{lemma}

\begin{proof}
这由 Deformation Problems 的引理
\ref{examples-defos-lemma-finite-projective-modules-TI},
\ref{examples-defos-lemma-representations-TI},
\ref{examples-defos-lemma-continuous-representations-TI} 与
\ref{examples-defos-lemma-graded-algebras-TI} 推出。
\end{proof}
% Generated mechanically from the frozen fdl-zh.tex target.
% Do not edit this generated file; edit the bound locale source instead.

% OK, start here.
%
%---------------------------------------------------------------------
%\tableofcontents

\setcounter{chapter}{115}
\chapter{GNU 自由文档许可证（GNU Free Documentation License）}

\phantomsection
\label{fdl-zh-section-phantom}

\begin{center}

      1.2 版，2002 年 11 月


版权所有 \copyright 2000, 2001, 2002 Free Software Foundation, Inc.

\bigskip

    51 Franklin St, Fifth Floor, Boston, MA  02110-1301  USA

\bigskip

人人均可复制和分发本许可证文档的逐字副本，
但不得对其作任何更改。
\end{center}

This is an unofficial translation of the GNU Free Documentation License into
Simplified Chinese. It was not published by the Free Software Foundation, and
does not legally state the distribution terms for documentation that uses the
GNU FDL---only the original English text of the GNU FDL does that. However, we
hope that this translation will help Simplified Chinese speakers understand
the GNU FDL better.

这是 GNU 自由文档许可证的非官方简体中文译文。它并非由自由软件基金会发布，也不构成采用
GNU FDL 的文档所适用的法律分发条款；只有 GNU FDL 的英文原文具有该等法律效力。不过，
我们希望本译文能帮助简体中文读者理解 GNU FDL。本译文必须与英文原文一并提供；如译文与
英文原文有任何分歧，以英文原文为准。

You may publish this translation, modified or unmodified, only under the terms
at https://www.gnu.org/licenses/translations.html.

只有遵守 https://www.gnu.org/licenses/translations.html 所列条款，方可发布本译文，
无论是否经过修改。


\begin{center}
{\bf\large 前言（Preamble）}
\end{center}

本许可证旨在使手册、教科书或其他具有实际功能和用途的文档获得自由意义上的“自由”：
确保每个人都切实享有复制和再分发该文档的自由，无论是否修改，也无论出于商业还是非商业
目的。其次，本许可证为作者和出版者保留一种使其工作获得署名认可的方式，同时不使他们因
他人所作的修改而承担责任。

本许可证是一种“copyleft”（著佐权）许可证；这意味着文档的衍生作品本身也必须在同一
意义上保持自由。它与 GNU 通用公共许可证相配合；后者是为自由软件设计的著佐权许可证。

我们设计本许可证，是为了将其用于自由软件的手册，因为自由软件需要自由文档：自由程序
应当附有向读者提供与该软件相同自由的手册。不过，本许可证并不限于软件手册；任何文字作品
都可以采用它，而不论主题为何，也不论是否以印刷书籍出版。我们尤其建议以教学或参考为目的
的作品采用本许可证。


\section{适用范围与定义（APPLICABILITY AND DEFINITIONS）}
\label{fdl-zh-section-applicability-and-definitions}

本许可证适用于以任何媒介承载的任何手册或其他作品，只要其中载有版权持有人作出的声明，
说明该作品可依本许可证条款分发。该声明授予一项全球范围、免版税、期限不受限制的许可，
允许依本文所述条件使用该作品。下文所称的\textbf{“文档”}，指任何此类手册或作品。
任何公众成员都是被许可人，本文以\textbf{“您”}称呼被许可人。若您以版权法要求获得许可
的方式复制、修改或分发该作品，即表示您接受本许可证。

文档的\textbf{“修改版”}，指任何包含该文档或其一部分的作品，无论其内容是逐字复制，
还是经过修改和／或翻译成另一种语言。

\textbf{“次要章节”}，指文档中有标题的附录或前置章节；该章节只讨论文档的出版者或作者
与文档总体主题（或相关事项）之间的关系，不包含任何可直接属于该总体主题的内容。（因此，
若文档的一部分是数学教科书，次要章节不得讲解任何数学。）这种关系可以是与该主题或相关
事项的历史联系，也可以是对这些事项所持的法律、商业、哲学、伦理或政治立场。

\textbf{“不变章节”}，指某些次要章节，其标题在说明文档依本许可证发布的声明中被指定为
不变章节的标题。若某一章节不符合上述次要章节的定义，就不得将其指定为不变章节。文档可以
没有不变章节。若文档未指明任何不变章节，则不存在不变章节。

\textbf{“封面文字”}，指说明文档依本许可证发布的声明中列为前封面文字或后封面文字的
某些简短文字。前封面文字最多为 5 个词，后封面文字最多为 25 个词。

文档的\textbf{“透明”}副本，指采用一种机器可读格式表示的副本；该格式的规范可供公众
取得，并适合使用通用文本编辑器直接修订文档，或者（对于由像素组成的图像）使用通用绘图
程序，或者（对于线条图）使用某种广泛可得的绘图编辑器修订；该格式还适合输入排版程序，
或适合自动转换为多种可输入排版程序的格式。若某一原本属于透明格式的文件，其标记方式或
无标记方式被特意安排成阻碍或劝阻读者日后修改，则该副本并不透明。若图像格式被用于承载
大量文字，该格式并不透明。不是“透明”的副本称为\textbf{“不透明”}副本。

适合作为透明副本的格式包括：无标记的纯 ASCII、Texinfo 输入格式、LaTeX 输入格式、采用
公开可得 DTD 的 SGML 或 XML，以及符合标准、为便于人工修改而设计的简单 HTML、PostScript
或 PDF。透明图像格式包括 PNG、XCF 和 JPG。不透明格式包括：只能由专有文字处理器读取和
编辑的专有格式；其 DTD 和／或处理工具并非普遍可得的 SGML 或 XML；以及某些文字处理器
仅为输出目的而生成的 HTML、PostScript 或 PDF。

对于印刷书籍，\textbf{“标题页”}指标题页本身，以及为清晰容纳本许可证要求在标题页出现
的材料所必需的后续页面。对于本身没有标题页的格式，“标题页”指正文开始之前、接近作品标题
最醒目处的文字。

一个\textbf{“标题为 XYZ”}的章节，是指文档中有标题的子单元，其标题恰为 XYZ，或者在
另一种语言对 XYZ 的译文之后以括号包含 XYZ。（这里的 XYZ 代表下文提到的某个特定章节名，
例如\textbf{“Acknowledgements”}（致谢）、\textbf{“Dedications”}（献辞）、
\textbf{“Endorsements”}（背书）或\textbf{“History”}（历史）。）当您修改文档时，
\textbf{“保留该标题”}，是指该章节仍依本定义构成“标题为 XYZ”的章节。

文档可以在声明本许可证适用于该文档的通知旁列出担保免责声明。这些担保免责声明视为通过
引用纳入本许可证，但仅限于免除担保责任：这些担保免责声明可能产生的任何其他含义均属无效，
且不影响本许可证的含义。


\section{逐字复制（VERBATIM COPYING）}
\label{fdl-zh-section-verbatim-copying}

您可以在任何媒介中复制和分发文档，无论出于商业还是非商业目的，条件是：所有副本均重现
本许可证、版权声明以及说明本许可证适用于该文档的许可声明；并且您不得在本许可证的条件
之外附加任何其他条件。您不得使用技术措施妨碍或控制对您制作或分发之副本的阅读或进一步
复制。不过，您可以因提供副本而收取报酬。若您分发的副本数量足够大，还必须遵守第 3 节的
条件。

您也可以按上述相同条件出借副本，并可以公开展示副本。


\section{大量复制（COPYING IN QUANTITY）}
\label{fdl-zh-section-copying-in-quantity}


若您出版的文档印刷副本（或采用通常带有印刷封面的媒介制作的副本）超过 100 份，并且文档
的许可声明要求封面文字，则必须为副本加上清晰易读地载有全部封面文字的封面：前封面载有
前封面文字，后封面载有后封面文字。两个封面还必须清晰易读地表明您是这些副本的出版者。
前封面必须列出完整标题，标题中的所有词都应同样醒目可见。除上述内容外，您可以在封面上
增添其他材料。若改动仅限于封面，同时保留文档标题并满足这些条件，则在其他方面可视为逐字
复制。

若任一封面所要求的文字太长，无法清晰排入，您应将所列文字中的前若干项（在合理范围内尽量
多排）置于实际封面，并将其余内容续排在相邻页面。

若您出版或分发的文档不透明副本超过 100 份，则必须满足以下二者之一：随每一份不透明副本
附上一份机器可读的透明副本；或者在每一份不透明副本内或随附材料中说明一个计算机网络位置，
使一般网络公众能够通过公共标准网络协议下载文档的完整透明副本，且该透明副本不含附加材料。
若采用后一方案，则在开始大量分发不透明副本时，您必须采取合理审慎的措施，确保该透明副本
在所述位置持续可访问，直至您最后一次向公众分发该版本的不透明副本（无论直接分发还是通过
代理商或零售商分发）之后至少一年。

我们请求但不强制要求您在大量再分发副本之前及早联系文档作者，使其有机会向您提供文档的
更新版本。


\section{修改（MODIFICATIONS）}
\label{fdl-zh-section-modifications}

您可以按照上述第 2 节和第 3 节的条件复制和分发文档的修改版，但须将修改版完全依本许可证
发布，并由修改版承担“文档”的角色，从而允许任何持有修改版副本的人分发和修改该修改版。
此外，您必须在修改版中完成以下事项：

\begin{itemize}
\item[A.]
   在标题页（以及封面，如有）使用一个不同于文档标题和任何先前版本标题的标题；先前版本的
   标题（如有）应列在文档的“History”（历史）章节中。若某一先前版本的原出版者许可，您
   可以使用与该版本相同的标题。

\item[B.]
   在标题页上以作者身份列出一个或多个对修改版中的修改负有作者责任的个人或实体，同时列出
   文档至少五位主要作者（若主要作者不足五位，则列出全部主要作者），除非这些作者免除您的
   此项义务。

\item[C.]
   在标题页上以出版者身份写明修改版出版者的名称。

\item[D.]
   保留文档的全部版权声明。

\item[E.]
   在其他版权声明旁，为您所作的修改添加适当的版权声明。

\item[F.]
   紧接版权声明之后，按照下文附录所示格式加入许可声明，准许公众依本许可证条款使用修改版。

\item[G.]
   在该许可声明中保留文档许可声明所列不变章节和所要求封面文字的完整清单。

\item[H.]
   包含一份未经改动的本许可证副本。

\item[I.]
   保留标题为“History”（历史）的章节，保留其标题，并在其中加入一项，至少写明标题页所列
   修改版的标题、年份、新作者和出版者。若文档没有标题为“History”的章节，则创建该章节，
   写明文档标题页所列的标题、年份、作者和出版者，然后再按前一句所述加入一项描述修改版。

\item[J.]
   保留文档中为公众访问其透明副本而提供的网络位置（如有），并同样保留文档所依据的先前版本
   中所给的网络位置。这些信息可以放在“History”章节中。若某一作品在文档本身之前至少四年
   已经出版，或者该网络位置所指版本的原出版者允许，您可以省略该作品的网络位置。

\item[K.]
   对于任何标题为“Acknowledgements”（致谢）或“Dedications”（献辞）的章节，保留该章节
   的标题，并在该章节中保留其中每一项贡献者致谢和／或献辞的全部实质内容与语气。

\item[L.]
   保留文档的全部不变章节，不得改动其文字和标题。章节编号或同等标识不视为章节标题的一部分。

\item[M.]
   删除任何标题为“Endorsements”（背书）的章节。修改版不得包含此类章节。

\item[N.]
   不得将任何既有章节改名为“Endorsements”，也不得将其改成与任何不变章节标题冲突的标题。

\item[O.]
   保留所有担保免责声明。
\end{itemize}

若修改版包含新的前置章节或附录，且这些部分符合次要章节的定义并且不含从文档复制的材料，
您可以自行选择将其中一些或全部指定为不变章节。为此，请把它们的标题加入修改版许可声明中的
不变章节清单。这些标题必须与其他所有章节标题不同。

您可以加入标题为“Endorsements”（背书）的章节，但其中只能包含各方对您的修改版所作的
背书，例如同行评审声明，或说明某一组织已认可该文本为某项标准的权威定义。

您可以在修改版封面文字清单的末尾加入一段最多五个词的前封面文字，以及一段最多二十五个词
的后封面文字。任何一个实体（或由其安排）只能加入一段前封面文字和一段后封面文字。若文档
已经在同一封面包含先前由您加入、或由您所代表的同一实体安排加入的封面文字，您不得再加入
另一段；但在取得添加旧文字之先前出版者的明确许可后，可以替换旧文字。

文档的作者和出版者并未通过本许可证准许他人使用其姓名作宣传，也未准许他人声称或暗示其为
任何修改版背书。


\section{合并文档（COMBINING DOCUMENTS）}
\label{fdl-zh-section-combining-documents}


您可以依上述第 4 节为修改版规定的条款，将文档与其他依本许可证发布的文档合并，但条件是：
合并作品中必须包含所有原文档的全部不变章节，且不得改动；必须在合并作品的许可声明中将它们
全部列为不变章节；并且必须保留其全部担保免责声明。

合并作品只需包含本许可证的一份副本，多个完全相同的不变章节可以合并为一份。若存在标题相同
但内容不同的多个不变章节，应在每一章节标题末尾的括号内加入该章节原作者或原出版者的姓名
（如已知），否则加入一个唯一编号，使各标题彼此不同。合并作品许可声明中的不变章节清单也须
对相应章节标题作同样调整。

合并时，必须把各原文档中标题为“History”（历史）的章节合并成一个标题为“History”的章节；
同样合并所有标题为“Acknowledgements”（致谢）的章节和所有标题为“Dedications”（献辞）的
章节。必须删除所有标题为“Endorsements”（背书）的章节。

\section{文档合集（COLLECTIONS OF DOCUMENTS）}
\label{fdl-zh-section-collections-of-documents}

您可以制作由本文档和其他依本许可证发布的文档组成的合集，并以合集内的一份许可证副本替代
各文档分别包含的许可证副本，但条件是在所有其他方面，对每一份文档都遵守本许可证关于逐字
复制的规则。

您可以从此类合集中取出一份文档，并依本许可证单独分发，但必须在取出的文档中加入一份本许可
证副本，并在有关该文档逐字复制的所有其他方面遵守本许可证。


\section{与独立作品聚合（AGGREGATION WITH INDEPENDENT WORKS）}
\label{fdl-zh-section-aggregation-with-independent-works}


若将文档或其衍生作品与其他相互独立的文档或作品一同汇编在某一存储或分发媒介卷册中或其上，
且汇编产生的版权并未被用于在各独立作品所允许的范围之外限制汇编物使用者的法律权利，则该汇编
称为“聚合物”。当文档被纳入聚合物时，本许可证不适用于聚合物中本身并非文档衍生作品的其他作品。

若第 3 节的封面文字要求适用于这些文档副本，并且文档少于整个聚合物的一半，则文档的封面文字
可以放在聚合物内部包围该文档的封面上；若文档采用电子形式，则可放在相当于封面的电子界面上。
否则，封面文字必须出现在包围整个聚合物的印刷封面上。


\section{翻译（TRANSLATION）}
\label{fdl-zh-section-translation}


翻译视为一种修改，因此您可以依第 4 节的条款分发文档的译本。以译文替换不变章节，须取得
相应版权持有人的特别许可；但您可以在这些不变章节的原文之外，附加其中一部分或全部的译文。
您可以附上本许可证、文档中全部许可声明以及任何担保免责声明的译文，但同时必须包含本许可证
的英文原版，以及上述声明和免责声明的原版。若译文与本许可证、声明或免责声明的原版有任何
分歧，以原版为准。

若文档中有标题为“Acknowledgements”（致谢）、“Dedications”（献辞）或“History”（历史）
的章节，则第 4 节关于保留其标题（定义见第 1 节）的要求，通常意味着必须更改实际标题。


\section{终止（TERMINATION）}
\label{fdl-zh-section-termination}


除本许可证明确允许的情形外，您不得复制、修改、再许可或分发文档。任何其他复制、修改、
再许可或分发文档的企图均属无效，并将自动终止您依本许可证享有的权利。不过，只要从您处
获得副本或权利的各方继续完全遵守本许可证，其许可就不会因此终止。


\section{本许可证的未来修订版（FUTURE REVISIONS OF THIS LICENSE）}
\label{fdl-zh-section-future-revisions-of-this-license}


自由软件基金会可能不时发布 GNU 自由文档许可证的新修订版。新版本在精神上将与本版本相似，
但可能在细节上有所不同，以处理新的问题或关切。参见 https://www.gnu.org/copyleft/。

本许可证的每一版本均有不同的版本号。若文档说明本许可证的某个指定编号版本“或任何后续
版本”适用于该文档，您可以选择遵守该指定版本的条款和条件，也可以选择遵守自由软件基金会
此后正式发布（而非作为草案发布）的任何后续版本。若文档未指定本许可证的版本号，您可以
选择自由软件基金会曾经正式发布（而非作为草案发布）的任何版本。


\section{附录：如何将本许可证用于您的文档（ADDENDUM: How to use this License for your documents）}
\label{fdl-zh-section-addendum-how-to-use-this-license-for-your-documents}

若要将本许可证用于您撰写的文档，请在文档中包含一份本许可证副本，并在标题页之后紧接着
加入以下版权和许可声明：

\bigskip
\begin{quote}
    Copyright \copyright  YEAR  YOUR NAME.
    Permission is granted to copy, distribute and/or modify this document
    under the terms of the GNU Free Documentation License, Version 1.2
    or any later version published by the Free Software Foundation;
    with no Invariant Sections, no Front-Cover Texts, and no Back-Cover Texts.
    A copy of the license is included in the section entitled "GNU
    Free Documentation License".
\end{quote}
\bigskip

若文档包含不变章节、前封面文字和后封面文字，请用以下文字替换“with...Texts.”一行：

\bigskip
\begin{quote}
    with the Invariant Sections being LIST THEIR TITLES, with the
    Front-Cover Texts being LIST, and with the Back-Cover Texts being LIST.
\end{quote}
\bigskip

若文档有不变章节而没有封面文字，或者三者存在其他组合，请合并上述两种写法以符合实际情况。

若您的文档含有实质性的程序代码示例，我们建议同时依您选择的自由软件许可证（例如 GNU 通用
公共许可证）发布这些示例，使其可以用于自由软件。

%---------------------------------------------------------------------

\appendix
% Generated mechanically from the frozen index.tex target.
% Do not edit this generated file; edit the bound locale source instead.
\chapter{自动生成的索引}

\phantomsection
\label{index-section-phantom}


\frenchspacing


\begin{multicols}{2}[\section{按英文源词条字母顺序排列的定义}\label{index-section-alphabetized}]
\noindent
{\it $(-1)$-指标}
见 \ref{models-definition-type-minus-one}

\noindent
{\it $(-2)$-指标}
见 \ref{models-definition-type-minus-two}

\noindent
{\it $(2, 1)$-范畴}
见 \ref{categories-definition-2-1-category}

\noindent
{\it $(2, 1)$-周期复形}
见 \ref{chow-definition-periodic-complex}

\noindent
{\it $(\mathcal{F}_n)$ 可典范地延拓到 $X$}
见 \ref{algebraization-definition-canonically-algebraizable}

\noindent
{\it $(\mathcal{F}_n)$ 可延拓到 $X$}
见 \ref{algebraization-definition-algebraizable}

\noindent
{\it $(\mathcal{F}_n)$ 满足 $(a, b)$-不等式组}
见 \ref{algebraization-definition-s-d-inequalities}

\noindent
{\it $(\mathcal{F}_n)$ 满足严格 $(a, b)$-不等式组}
见 \ref{algebraization-definition-s-d-inequalities}

\noindent
{\it $(\textit{Spaces}/S)_\etale$}
见 \ref{spaces-topologies-definition-big-etale-site}

\noindent
{\it $(\textit{Spaces}/S)_{fppf}$}
见 \ref{spaces-topologies-definition-big-fppf-site}

\noindent
{\it $(\textit{Spaces}/S)_{ph}$}
见 \ref{spaces-topologies-definition-big-ph-site}

\noindent
{\it $(\textit{Spaces}/X)_\etale$}
见 \ref{spaces-topologies-definition-big-small-etale}

\noindent
{\it $(\textit{Spaces}/X)_{fppf}$}
见 \ref{spaces-topologies-definition-big-small-fppf}

\noindent
{\it $(\textit{Spaces}/X)_{ph}$}
见 \ref{spaces-topologies-definition-big-small-ph}

\noindent
{\it $(A, B)$-双模}
见 \ref{algebra-definition-bimodule}

\noindent
{\it $(A, B)$-双模}
见 \ref{dga-definition-bimodule}

\noindent
{\it $(R_k)$}
见 \ref{algebra-definition-conditions}

\noindent
{\it $(R_k)$}
见 \ref{properties-definition-Rk}

\noindent
{\it $(S_k)$}
见 \ref{algebra-definition-conditions}

\noindent
{\it $(S_k)$}
见 \ref{properties-definition-Rk}

\noindent
{\it $(S_k)$}
见 \ref{coherent-definition-depth}

\noindent
{\it $(S_k)$}
见 \ref{coherent-definition-depth}

\noindent
{\it $(U', R', s', t', c')$ 在 $(U, R, s, t, c)$ 上是笛卡尔的}
见 \ref{groupoids-definition-cartesian-morphism}

\noindent
{\it $1$-态射}
见 \ref{categories-definition-2-category}

\noindent
{\it $S$ 上代数叠的 $2$-范畴}
见 \ref{algebraic-definition-morphism-algebraic-stacks}

\noindent
{\it $\mathcal{C}$ 上群胚纤维化范畴 所成的 $2$-范畴}
见 \ref{categories-definition-categories-fibred-in-groupoids-over-C}

\noindent
{\it $\mathcal{C}$ 上集合胚纤维化范畴 所成的 $2$-范畴}
见 \ref{categories-definition-categories-fibred-in-setoids-over-C}

\noindent
{\it $\mathcal{C}$ 上集合纤维化范畴 所成的 $2$-范畴}
见 \ref{categories-definition-categories-fibred-in-sets-over-C}

\noindent
{\it $\mathcal{C}$ 上的范畴所成的 $2$-范畴}
见 \ref{categories-definition-categories-over-C}

\noindent
{\it $\mathcal{C}$ 上纤维化范畴 所成的 $2$-范畴}
见 \ref{categories-definition-fibred-categories-over-C}

\noindent
{\it $\mathcal{C}$ 上群胚中的栈的 $2$-范畴}
见 \ref{stacks-definition-stacks-in-groupoids-over-C}

\noindent
{\it $\mathcal{C}$ 上类集合中的栈的 $2$-范畴}
见 \ref{stacks-definition-stacks-in-setoids-over-C}

\noindent
{\it $\mathcal{C}$ 上栈的 $2$-范畴}
见 \ref{stacks-definition-stacks-over-C}

\noindent
{\it $2$-范畴}
见 \ref{categories-definition-2-category}

\noindent
{\it $2$-态射}
见 \ref{categories-definition-2-category}

\noindent
{\it $2$-周期复形}
见 \ref{chow-definition-periodic-complex}

\noindent
{\it 相对于 $I$ 为 $\alpha$-小}
见 \ref{injectives-definition-small}

\noindent
{\it $\delta$ 与 $\gamma$ 兼容}
见 \ref{crystalline-definition-compatible}

\noindent
{\it $T$ 的 $\delta$-维数}
见 \ref{spaces-chow-definition-delta-dimension}

\noindent
{\it $Z$ 的 $\delta$-维数}
见 \ref{chow-definition-delta-dimension}

\noindent
{\it 从 $\mathcal{A}$ 到 $\mathcal{D}$ 的 $\delta$-函子}
见 \ref{derived-definition-delta-functor}

\noindent
{\it $\delta$-函子}
见 \ref{homology-definition-cohomological-delta-functor}

\noindent
{\it $A$ 的 $\delta$-不变量}
见 \ref{varieties-definition-delta-invariant-algebra}

\noindent
{\it $X$ 在 $x$ 处的 $\delta$-不变量}
见 \ref{varieties-definition-delta-invariant}

\noindent
{\it $\delta(\tau)$}
见 \ref{exercises-definition-delta}

\noindent
{\it $\delta^n_j : [n-1] \to [n]$}
见 \ref{simplicial-definition-face-degeneracy}

\noindent
{\it $\ell$-进上同调}
见 \ref{trace-definition-cohomology-l-adic}

\noindent
{\it $\ell$-进层}
见 \ref{trace-definition-l-adic-sheaf}

\noindent
{\it $\epsilon$-不变量}
见 \ref{obsolete-definition-epsilon}

\noindent
{\it $\Ext$-群}
见 \ref{homology-definition-ext-group}

\noindent
{\it $\Hom(U, V)$}
见 \ref{simplicial-definition-hom-deltak-cosimplicial}

\noindent
{\it $\Hom(U, V)$}
见 \ref{simplicial-definition-hom-deltak-simplicial}

\noindent
{\it $\Hom(U, V)$}
见 \ref{simplicial-definition-hom-from-simplicial-set}

\noindent
{\it $\kappa$-生成的}
见 \ref{properties-definition-kappa-generated}

\noindent
{\it $\mathbf{Z}_\ell$-层}
见 \ref{trace-definition-l-adic-sheaf}

\noindent
{\it $\mathcal{A}^0$}
见 \ref{dga-definition-H0-of-graded-category}

\noindent
{\it $\mathcal{C}_\Lambda$}
见 \ref{formal-defos-definition-CLambda}

\noindent
{\it $\mathcal{F}$ 在 $x$ 处的长度为 $d$}
见 \ref{spaces-chow-definition-length-at-x}

\noindent
{\it $\mathcal{F}$ 在维数 $\geq n$ 上关于 $S$ 平坦}
见 \ref{flat-definition-flat-dimension-n}

\noindent
{\it $\mathcal{F}$ 在维数 $\geq n$ 上相对于 $Y$ 平坦}
见 \ref{spaces-flat-definition-flat-dimension-n}

\noindent
{\it $\mathcal{G}$-扭子}
见 \ref{cohomology-definition-torsor}

\noindent
{\it $\mathcal{G}$-扭子}
见 \ref{sites-cohomology-definition-torsor}

\noindent
{\it $\mathcal{I}$ 在 $\mathcal{J}$ 中共尾}
见 \ref{categories-definition-cofinal}

\noindent
{\it $\mathcal{I}$ 在 $\mathcal{J}$ 中是始的}
见 \ref{categories-definition-initial}

\noindent
{\it $\mathcal{K}_X$}
见 \ref{divisors-definition-sheaf-meromorphic-functions}

\noindent
{\it $\mathcal{K}_X$}
见 \ref{spaces-divisors-definition-sheaf-meromorphic-functions}

\noindent
{\it $\mathcal{O}^*$}
见 \ref{sites-modules-definition-invertible-sheaf}

\noindent
{\it $\mathcal{O}_1$-导子}
见 \ref{modules-definition-derivation}

\noindent
{\it $\mathcal{O}_1$-导子}
见 \ref{sites-modules-definition-derivation}

\noindent
{\it $\mathcal{O}_\mathcal{X}$-模}
见 \ref{stacks-sheaves-definition-modules}

\noindent
{\it $\mathcal{S}$ 配备从 $\mathcal{C}$ 继承的拓扑}
见 \ref{stacks-definition-topology-inherited}

\noindent
{\it $\mathcal{S}_F$}
见 \ref{categories-definition-split-fibred-category}

\noindent
{\it $\mathcal{S}_F$}
见 \ref{categories-definition-split-category-fibred-in-groupoids}

\noindent
{\it $\mathcal{X}$ 在 $\mathcal{Y}$ 上相对可表}
见 \ref{categories-definition-representable-map-categories-fibred-in-groupoids}

\noindent
{\it $\mathfrak g^r_d$}
见 \ref{curves-definition-linear-series}

\noindent
{\it $\mathit{QC}(\mathcal{A}, \text{d})$}
见 \ref{sdga-definition-cartesian}

\noindent
{\it $\mathit{QC}(\mathcal{O})$}
见 \ref{sites-cohomology-definition-cartesian}

\noindent
{\it $\phi$ 位于 $f$ 之上}
见 \ref{categories-definition-fibre-category}

\noindent
{\it $\Sh(\mathcal{C})$}
见 \ref{sites-definition-category-sheaves-sets}

\noindent
{\it $\sigma^n_j : [n + 1]  \to [n]$}
见 \ref{simplicial-definition-face-degeneracy}

\noindent
{\it $\tau$ $G$-扭子}
见 \ref{groupoids-definition-principal-homogeneous-space}

\noindent
{\it $\tau$ $G$-扭子}
见 \ref{spaces-groupoids-definition-principal-homogeneous-space}

\noindent
{\it 在基底上为 $\tau$-局部的}
见 \ref{descent-definition-property-morphisms-local}

\noindent
{\it 在基底上为 $\tau$-局部的}
见 \ref{spaces-descent-definition-property-morphisms-local}

\noindent
{\it 在源上 $\tau$ 局部的}
见 \ref{descent-definition-property-morphisms-local-source}

\noindent
{\it 在源上 $\tau$ 局部的}
见 \ref{spaces-descent-definition-property-morphisms-local-source}

\noindent
{\it 在目标上为 $\tau$-局部的}
见 \ref{descent-definition-property-morphisms-local}

\noindent
{\it 在目标上为 $\tau$-局部的}
见 \ref{spaces-descent-definition-property-morphisms-local}

\noindent
{\it $\tau$ 扭子}
见 \ref{groupoids-definition-principal-homogeneous-space}

\noindent
{\it $\tau$ 扭子}
见 \ref{spaces-groupoids-definition-principal-homogeneous-space}

\noindent
{\it $\tau$-覆盖}
见 \ref{etale-cohomology-definition-tau-covering}

\noindent
{\it $\textit{Adeq}((\Sch/S)_\tau, \mathcal{O})$}
见 \ref{adequate-definition-category-adequate-modules}

\noindent
{\it $\textit{Adeq}(\mathcal{O})$}
见 \ref{adequate-definition-category-adequate-modules}

\noindent
{\it $\textit{Adeq}(S)$}
见 \ref{adequate-definition-category-adequate-modules}

\noindent
{\it $\text{Fil}^f(\mathcal{A})$}
见 \ref{exercises-definition-finite-filtration-category}

\noindent
{\it $\underline{U}$}
见 \ref{sites-definition-representable-sheaf}

\noindent
{\it $\varphi$-导子}
见 \ref{modules-definition-derivation}

\noindent
{\it $\varphi$-导子}
见 \ref{sites-modules-definition-derivation}

\noindent
{\it $\widehat{\mathcal{C}}_\Lambda$}
见 \ref{formal-defos-definition-completion-CLambda}

\noindent
{\it $A$ 是 $A_1, \ldots, A_n$ 的楔}
见 \ref{varieties-definition-wedge}

\noindent
{\it $A$-双导子}
见 \ref{cotangent-definition-biderivation}

\noindent
{\it 从 $X$ 到 $Y$ 的 $B$-有理映射}
见 \ref{spaces-morphisms-definition-rational-map}

\noindent
{\it $c$-adic}
见 \ref{formal-spaces-definition-weakly-adic}

\noindent
{\it $C$}
见 \ref{descent-definition-C}

\noindent
{\it $C_r$}
见 \ref{etale-cohomology-definition-Cr}

\noindent
{\it $C_{S/R}$}
见 \ref{algebra-definition-universal-thickening}

\noindent
{\it $d(M)$}
见 \ref{algebra-definition-d}

\noindent
{\it $D_c(X_\etale, \Lambda)$}
见 \ref{etale-cohomology-definition-c}

\noindent
{\it $D_{ctf}(X_\etale, \Lambda)$}
见 \ref{etale-cohomology-definition-ctf}

\noindent
{\it $f s^{-1}$}
见 \ref{categories-definition-right-localization-as-fraction}

\noindent
{\it $f$ 在 $x$ 处具有相对维数 $d$}
见 \ref{spaces-morphisms-definition-dimension-fibre}

\noindent
{\it $F$ 在 $G$ 上相对可表}
见 \ref{categories-definition-representable-map-presheaves}

\noindent
{\it $f$-丰沛}
见 \ref{morphisms-definition-relatively-ample}

\noindent
{\it $f$-丰沛}
见 \ref{spaces-divisors-definition-relatively-ample}

\noindent
{\it $X/S$ 上（相对于 $\sigma$）的 $F$-晶体}
见 \ref{crystalline-definition-F-crystal}

\noindent
{\it $f$-映射 $\varphi : \mathcal{G} \to \mathcal{F}$}
见 \ref{spaces-properties-definition-f-map}

\noindent
{\it $f$-映射 $\xi : \mathcal{G} \to \mathcal{F}$}
见 \ref{sheaves-definition-f-map}

\noindent
{\it $f$-相对丰沛}
见 \ref{morphisms-definition-relatively-ample}

\noindent
{\it $f$-相对丰沛}
见 \ref{spaces-divisors-definition-relatively-ample}

\noindent
{\it $f$-相对甚丰沛}
见 \ref{morphisms-definition-very-ample}

\noindent
{\it $f$-甚丰沛}
见 \ref{morphisms-definition-very-ample}

\noindent
{\it $f^{-1}\mathcal{S}$}
见 \ref{stacks-definition-pullback-stack}

\noindent
{\it $f_*$}
见 \ref{descent-definition-pushforward}

\noindent
{\it $f_*\mathcal{S}$}
见 \ref{stacks-definition-pushforward-stack}

\noindent
{\it $G$-等变拟凝聚 $\mathcal{O}_X$-模}
见 \ref{groupoids-definition-equivariant-module}

\noindent
{\it $G$-等变拟凝聚 $\mathcal{O}_X$-模}
见 \ref{spaces-groupoids-definition-equivariant-module}

\noindent
{\it $G$-等变态射}
见 \ref{groupoids-definition-action-group-scheme}

\noindent
{\it $G$-等变态射}
见 \ref{spaces-groupoids-definition-action-group-space}

\noindent
{\it $G$-不变的}
见 \ref{groupoids-quotients-definition-invariant}

\noindent
{\it $G$-模}
见 \ref{etale-cohomology-definition-G-module-continuous}

\noindent
{\it $G$-集}
见 \ref{pione-definition-G-set-continuous}

\noindent
{\it $\tau$ 拓扑中的 $G$-扭子}
见 \ref{groupoids-definition-principal-homogeneous-space}

\noindent
{\it $\tau$ 拓扑中的 $G$-扭子}
见 \ref{spaces-groupoids-definition-principal-homogeneous-space}

\noindent
{\it $G$-扭子}
见 \ref{groupoids-definition-principal-homogeneous-space}

\noindent
{\it $f$ 在 $P$ 上的 $G$-迹}
见 \ref{trace-definition-trace-G}

\noindent
{\it $G\textit{-Sets}$}
见 \ref{pione-definition-G-set-continuous}

\noindent
{\it $g_!\mathcal{F} = (g_{p!}\mathcal{F})^\#$}
见 \ref{sites-modules-definition-g-shriek}

\noindent
{\it $g_{p!}\mathcal{F}$}
见 \ref{sites-modules-definition-g-shriek}

\noindent
{\it $H^{i + k}(A^\bullet) \longrightarrow H^i(A[k]^\bullet)$}
见 \ref{homology-definition-cohomology-shift}

\noindent
{\it $H_1$-正则理想}
见 \ref{more-algebra-definition-regular-ideal}

\noindent
{\it $H_1$-正则浸入}
见 \ref{divisors-definition-regular-immersion}

\noindent
{\it $H_1$-正则浸入}
见 \ref{spaces-more-morphisms-definition-regular-immersion}

\noindent
{\it $H_1$-正则}
见 \ref{more-algebra-definition-koszul-regular-sequence}

\noindent
{\it $H_1$-正则}
见 \ref{divisors-definition-regular-ideal-sheaf}

\noindent
{\it $H_{i + k}(A_\bullet) \rightarrow H_i(A[k]_\bullet)$}
见 \ref{homology-definition-homology-shift}

\noindent
{\it $I$-进完备}
见 \ref{algebra-definition-complete}

\noindent
{\it $I$-进完备}
见 \ref{algebra-definition-complete}

\noindent
{\it $I$-深度}
见 \ref{algebra-definition-depth}

\noindent
{\it $I$-深度}
见 \ref{local-cohomology-definition-depth-complex}

\noindent
{\it $I$-幂挠模}
见 \ref{more-algebra-definition-f-power-torsion}

\noindent
{\it $I$-投射}
见 \ref{more-algebra-definition-near-projective}

\noindent
{\it $\mathcal{E}$ 的第 $i$ 个 Chern 类}
见 \ref{spaces-chow-definition-chern-classes}

\noindent
{\it 第 $i$ 陈类}
见 \ref{chow-definition-chern-classes-final}

\noindent
{\it $M$ 的第 $i$ 个 Chow 群}
见 \ref{weil-definition-chow-group-motives}

\noindent
{\it 第 $i$ 扩张群}
见 \ref{derived-definition-ext}

\noindent
{\it $F$ 的第 $i$ 个右导出函子 $R^iF$}
见 \ref{derived-definition-higher-derived-functors}

\noindent
{\it 与 $\mathcal{F}$ 关联的 $k$ 维循环}
见 \ref{chow-definition-cycle-associated-to-coherent-sheaf}

\noindent
{\it 与 $\mathcal{F}$ 关联的 $k$ 维循环}
见 \ref{spaces-chow-definition-cycle-associated-to-coherent-sheaf}

\noindent
{\it $Y$ 的伴随 $k$ 维闭链}
见 \ref{spaces-chow-definition-cycle-associated-to-closed-subscheme}

\noindent
{\it 与 $Z$ 相伴的 $k$-循环}
见 \ref{chow-definition-cycle-associated-to-closed-subscheme}

\noindent
{\it $k$ 维循环}
见 \ref{chow-definition-cycles}

\noindent
{\it $k$ 维循环}
见 \ref{spaces-chow-definition-cycles}

\noindent
{\it $k$-移位链复形 $A[k]_\bullet$}
见 \ref{homology-definition-shift}

\noindent
{\it $k$-移位上链复形 $A[k]^\bullet$}
见 \ref{homology-definition-shift-cochain}

\noindent
{\it $k$-移位模}
见 \ref{dga-definition-shift}

\noindent
{\it 第 $k$ 个 Fitting 理想}
见 \ref{more-algebra-definition-fitting-ideal}

\noindent
{\it 第 $k$ 个移位 $A$-模}
见 \ref{dga-definition-shift-graded-module}

\noindent
{\it 第 $k$ 个移位 $A$-模}
见 \ref{dga-definition-shift-graded-module}

\noindent
{\it $\mathcal{F}$ 的 $L$-函数}
见 \ref{trace-definition-L-function-finite-ring}

\noindent
{\it $\mathcal{F}$ 的 $L$-函数}
见 \ref{trace-definition-L-function-l-adic}

\noindent
{\it $M \mapsto M^\vee$}
见 \ref{more-algebra-definition-simple-functors}

\noindent
{\it $M$-$H_1$-正则}
见 \ref{more-algebra-definition-koszul-regular-sequence}

\noindent
{\it $M$-Koszul-正则}
见 \ref{more-algebra-definition-koszul-regular-sequence}

\noindent
{\it 相对于 $R$ 为 $m$-伪凝聚}
见 \ref{more-algebra-definition-relatively-pseudo-coherent}

\noindent
{\it 相对于 $R$ 为 $m$-伪凝聚}
见 \ref{more-algebra-definition-relatively-pseudo-coherent}

\noindent
{\it 相对于 $S$ 是 $m$-伪凝聚的}
见 \ref{more-morphisms-definition-relative-pseudo-coherence}

\noindent
{\it 相对于 $S$ 是 $m$-伪凝聚的}
见 \ref{more-morphisms-definition-relative-pseudo-coherence}

\noindent
{\it 相对于 $Y$ 为 $m$-伪凝聚}
见 \ref{spaces-more-morphisms-definition-relative-pseudo-coherence}

\noindent
{\it 相对于 $Y$ 为 $m$-伪凝聚}
见 \ref{spaces-more-morphisms-definition-relative-pseudo-coherence}

\noindent
{\it $m$-伪凝聚}
见 \ref{more-algebra-definition-pseudo-coherent}

\noindent
{\it $m$-伪凝聚}
见 \ref{more-algebra-definition-pseudo-coherent}

\noindent
{\it $m$-伪凝聚}
见 \ref{cohomology-definition-pseudo-coherent}

\noindent
{\it $m$-伪凝聚}
见 \ref{cohomology-definition-pseudo-coherent}

\noindent
{\it $m$-伪凝聚}
见 \ref{sites-cohomology-definition-pseudo-coherent}

\noindent
{\it $m$-伪凝聚}
见 \ref{sites-cohomology-definition-pseudo-coherent}

\noindent
{\it $M$-拟正则}
见 \ref{algebra-definition-quasi-regular-sequence}

\noindent
{\it $I$ 中的 $M$-正则序列}
见 \ref{algebra-definition-regular-sequence}

\noindent
{\it $M$-正则序列}
见 \ref{algebra-definition-regular-sequence}

\noindent
{\it $m$-正则}
见 \ref{varieties-definition-regularity}

\noindent
{\it $U$ 的 $n$-单形}
见 \ref{simplicial-definition-terminology-simplicial-sets}

\noindent
{\it $\mathcal{C}$ 的 $n$-截断单纯对象}
见 \ref{simplicial-definition-truncated-simplicial-object}

\noindent
{\it $n$ 阶无穷小邻域}
见 \ref{more-morphisms-definition-first-order-infinitesimal-neighbourhood}

\noindent
{\it $n$ 阶无穷小邻域}
见 \ref{spaces-more-morphisms-definition-first-order-infinitesimal-neighbourhood}

\noindent
{\it $n$ 阶加厚}
见 \ref{more-morphisms-definition-thickening}

\noindent
{\it $n$ 阶加厚}
见 \ref{spaces-more-morphisms-definition-thickening}

\noindent
{\it $R$-双线性的}
见 \ref{algebra-definition-bilinear}

\noindent
{\it $R$-导子}
见 \ref{algebra-definition-derivation}

\noindent
{\it $R$-等价}
见 \ref{groupoids-quotients-definition-geometric-orbits}

\noindent
{\it $R$-不变的}
见 \ref{groupoids-definition-invariant-open}

\noindent
{\it $R$-不变的}
见 \ref{groupoids-definition-invariant-open}

\noindent
{\it $R$-不变的}
见 \ref{groupoids-definition-invariant-open}

\noindent
{\it $R$-不变的}
见 \ref{spaces-groupoids-definition-invariant-open}

\noindent
{\it $R$-不变的}
见 \ref{spaces-groupoids-definition-invariant-open}

\noindent
{\it $R$-不变的}
见 \ref{spaces-groupoids-definition-invariant-open}

\noindent
{\it $R$-不变的}
见 \ref{groupoids-quotients-definition-invariant}

\noindent
{\it $R$-线性范畴 $\mathcal{A}$}
见 \ref{dga-definition-linear-category}

\noindent
{\it $R$-线性函子}
见 \ref{dga-definition-functor-linear-categories}

\noindent
{\it $R$-线性}
见 \ref{formal-defos-definition-linear}

\noindent
{\it 有限表示的 $R$-模}
见 \ref{algebra-definition-module-finite-type}

\noindent
{\it $R$-轨道}
见 \ref{groupoids-quotients-definition-orbit}

\noindent
{\it $R$-轨道}
见 \ref{groupoids-quotients-definition-geometric-orbits}

\noindent
{\it $R$-完美}
见 \ref{more-algebra-definition-relatively-perfect}

\noindent
{\it $R\text{-}G$-模}
见 \ref{etale-cohomology-definition-G-module-continuous}

\noindent
{\it $R_{(f)}$}
见 \ref{exercises-definition-Dplus-Vplus}

\noindent
{\it $S$ 是有限型 $R$-代数}
见 \ref{algebra-definition-finite-type}

\noindent
{\it $S$-双有理的}
见 \ref{morphisms-definition-birational-schemes}

\noindent
{\it $S$-导子 $D : \mathcal{O}_{X/S} \to \mathcal{F}$}
见 \ref{crystalline-definition-global-derivation}

\noindent
{\it $S$-导子}
见 \ref{modules-definition-differentials}

\noindent
{\it $S$-完美的}
见 \ref{perfect-definition-relatively-perfect}

\noindent
{\it $S$-纯}
见 \ref{flat-definition-pure}

\noindent
{\it $S$-纯}
见 \ref{flat-definition-pure}

\noindent
{\it 从 $X$ 到 $Y$ 的 $S$-有理映射}
见 \ref{morphisms-definition-rational-map}

\noindent
{\it $s^{-1} f$}
见 \ref{categories-definition-left-localization-as-fraction}

\noindent
{\it $T$ 在 $Y$ 上固有}
见 \ref{spaces-perfect-definition-proper-over-base}

\noindent
{\it $U$-可容许吹起}
见 \ref{divisors-definition-admissible-blowup}

\noindent
{\it $U$-可容许吹起}
见 \ref{spaces-divisors-definition-admissible-blowup}

\noindent
{\it $x$ 是 $X$ 上余维为 $d$ 的点}
见 \ref{spaces-properties-definition-dimension-local-ring}

\noindent
{\it $x$ 是 $X$ 的伴随点}
见 \ref{spaces-divisors-definition-weakly-associated}

\noindent
{\it $x$ 与 $\mathcal{F}$ 伴随}
见 \ref{spaces-divisors-definition-weakly-associated}

\noindent
{\it $X$ 在 $x$ 处正则}
见 \ref{spaces-properties-definition-regular-at-point}

\noindent
{\it $x$ 位于 $U$ 之上}
见 \ref{categories-definition-fibre-category}

\noindent
{\it $X_{affine, \etale}$}
见 \ref{spaces-properties-definition-affine-etale-site}

\noindent
{\it $X_{spaces, \etale}$}
见 \ref{spaces-properties-definition-spaces-etale-site}

\noindent
{\it $Y$ 在 $X$ 上是笛卡尔的}
见 \ref{spaces-simplicial-definition-cartesian-morphism}

\noindent
{\it $Y$-导子}
见 \ref{sites-modules-definition-sheaf-differentials}

\noindent
{\it $Y$-完美}
见 \ref{spaces-more-morphisms-definition-relatively-perfect}

\noindent
{\it $Y$-纯}
见 \ref{spaces-flat-definition-pure}

\noindent
{\it $Y$-纯}
见 \ref{spaces-flat-definition-pure}

\noindent
{\it $Z$ 在 $S$ 上固有}
见 \ref{coherent-definition-proper-over-base}

\noindent
{\it （加性）Herbrand 商}
见 \ref{chow-definition-periodic-length}

\noindent
{\it $f$ 与 $g$ 的 2-纤维积}
见 \ref{categories-definition-2-fibre-products}

\noindent
{\it 从 $f$ 到 $g$ 的 2-态射}
见 \ref{sites-definition-2-morphism-topoi}

\noindent
{\it 从 $f$ 到 $g$ 的 2-态射}
见 \ref{sites-modules-definition-2-morphism-ringed-topoi}

\noindent
{\it 在 $\mathfrak q$ 处平展}
见 \ref{algebra-definition-etale}

\noindent
{\it 在 $x \in X$ 处平展}
见 \ref{morphisms-definition-etale}

\noindent
{\it 在 $x \in X$ 处平展}
见 \ref{etale-definition-etale-schemes-1}

\noindent
{\it 在 $x$ 处 \'etale}
见 \ref{spaces-morphisms-definition-etale}

\noindent
{\it $T$ 的平展覆盖}
见 \ref{topologies-definition-etale-covering}

\noindent
{\it $X$ 的平展覆盖}
见 \ref{spaces-topologies-definition-etale-covering}

\noindent
{\it 平展覆盖}
见 \ref{etale-cohomology-definition-etale-covering-initial}

\noindent
{\it 平展覆盖}
见 \ref{etale-cohomology-definition-etale-covering}

\noindent
{\it \'etale 等价关系}
见 \ref{spaces-definition-etale-equivalence-relation}

\noindent
{\it 局部环的平展（\'etale）同态}
见 \ref{etale-definition-etale-ring}

\noindent
{\it 在源与目标上为 \'etale 局部的}
见 \ref{descent-definition-local-source-target}

\noindent
{\it 在源与目标上为 \'etale 局部的}
见 \ref{descent-definition-local-source-target-at-point}

\noindent
{\it $S$ 在 $\overline{s}$ 处的平展局部环}
见 \ref{etale-cohomology-definition-etale-local-rings}

\noindent
{\it $X$ 在 $\overline{x}$ 处的平展局部环}
见 \ref{spaces-properties-definition-etale-local-rings}

\noindent
{\it \'etale 局部可构造}
见 \ref{spaces-properties-definition-locally-constructible}

\noindent
{\it \'etale 局部}
见 \ref{descent-definition-local-at-point}

\noindent
{\it \'etale 邻域}
见 \ref{etale-cohomology-definition-geometric-point}

\noindent
{\it \'etale 邻域}
见 \ref{spaces-properties-definition-etale-neighbourhood}

\noindent
{\it $(S, s)$ 的平展邻域}
见 \ref{more-morphisms-definition-etale-neighbourhood}

\noindent
{\it \'etale 层}
见 \ref{stacks-sheaves-definition-sheaves}

\noindent
{\it \'etale 拓扑斯}
见 \ref{etale-cohomology-definition-etale-topos}

\noindent
{\it \'etale 拓扑斯}
见 \ref{spaces-properties-definition-etale-topos}

\noindent
{\it 在源与靶上相对于平展—光滑拓扑为局部的}
见 \ref{spaces-descent-definition-etale-smooth-local-source-target}

\noindent
{\it 平展}
见 \ref{algebra-definition-etale}

\noindent
{\it 平展}
见 \ref{morphisms-definition-etale}

\noindent
{\it 平展}
见 \ref{descent-definition-etale-morphism-germs}

\noindent
{\it 平展}
见 \ref{etale-definition-etale-schemes-1}

\noindent
{\it 平展}
见 \ref{etale-cohomology-definition-etale-morphism}

\noindent
{\it 平展}
见 \ref{spaces-properties-definition-etale}

\noindent
{\it 平展}
见 \ref{stacks-morphisms-definition-etale}

\noindent
{\it 存在 Serre 函子}
见 \ref{equiv-definition-Serre-functor}

\noindent
{\it $X$ 上的阿贝尔预层}
见 \ref{sheaves-definition-abelian-presheaves}

\noindent
{\it 阿贝尔预层}
见 \ref{etale-cohomology-definition-presheaf}

\noindent
{\it $X$ 上的阿贝尔层}
见 \ref{sheaves-definition-abelian-sheaf}

\noindent
{\it 阿贝尔层}
见 \ref{etale-cohomology-definition-category-sheaves}

\noindent
{\it 阿贝尔簇}
见 \ref{groupoids-definition-abelian-variety}

\noindent
{\it 阿贝尔}
见 \ref{homology-definition-abelian-category}

\noindent
{\it $X$ 的绝对 Frobenius}
见 \ref{varieties-definition-absolute-frobenius}

\noindent
{\it 绝对 Galois 群}
见 \ref{etale-cohomology-definition-algebraic-geometric-point}

\noindent
{\it 绝对分歧指数}
见 \ref{more-algebra-definition-mixed}

\noindent
{\it 绝对弱正规化}
见 \ref{morphisms-definition-seminormalization}

\noindent
{\it 绝对平坦的}
见 \ref{more-algebra-definition-weakly-etale}

\noindent
{\it 绝对平坦的}
见 \ref{more-algebra-definition-weakly-etale}

\noindent
{\it 绝对平坦的}
见 \ref{more-morphisms-definition-weakly-etale}

\noindent
{\it 绝对整闭}
见 \ref{more-algebra-definition-absolutely-integrally-closed}

\noindent
{\it 绝对弱正规}
见 \ref{morphisms-definition-seminormal-ring}

\noindent
{\it 绝对弱正规}
见 \ref{morphisms-definition-seminormal}

\noindent
{\it 收敛于 $H(K)$}
见 \ref{homology-definition-filtered-differential-ss-converges}

\noindent
{\it 收敛于 $H^*(K^\bullet)$}
见 \ref{homology-definition-filtered-complex-ss-converges}

\noindent
{\it 收敛于 $H^n(\text{Tot}(K^{\bullet, \bullet}))$}
见 \ref{homology-definition-ss-double-complex-converge}

\noindent
{\it 收敛于 $H^n(\text{Tot}(K^{\bullet, \bullet}))$}
见 \ref{homology-definition-ss-double-complex-converge}

\noindent
{\it $G$ 在代数空间 $X/B$ 上的作用}
见 \ref{spaces-groupoids-definition-action-group-space}

\noindent
{\it $G$ 在概形 $X/S$ 上的作用}
见 \ref{groupoids-definition-action-group-scheme}

\noindent
{\it 自由作用于}
见 \ref{spaces-definition-quotient}

\noindent
{\it 对 $LF$ 无环}
见 \ref{derived-definition-derived-functor}

\noindent
{\it 对 $RF$ 无环}
见 \ref{derived-definition-derived-functor}

\noindent
{\it 无环}
见 \ref{homology-definition-quasi-isomorphism}

\noindent
{\it 无环}
见 \ref{homology-definition-quasi-isomorphism-cochain}

\noindent
{\it 加性幺半范畴}
见 \ref{homology-definition-additive-monoidal}

\noindent
{\it 加性}
见 \ref{homology-definition-preadditive}

\noindent
{\it 加性}
见 \ref{homology-definition-additive-category}

\noindent
{\it 适足}
见 \ref{adequate-definition-adequate-functor}

\noindent
{\it 适足}
见 \ref{adequate-definition-adequate}

\noindent
{\it 进可构造}
见 \ref{proetale-definition-adic}

\noindent
{\it 进可构造}
见 \ref{proetale-definition-adic-constructible}

\noindent
{\it 进 lisse}
见 \ref{proetale-definition-adic}

\noindent
{\it 进 lisse}
见 \ref{proetale-definition-adic-constructible}

\noindent
{\it 进态射}
见 \ref{formal-spaces-definition-adic-morphism}

\noindent
{\it adic*}
见 \ref{formal-spaces-definition-types-affine-formal-algebraic-space}

\noindent
{\it adic}
见 \ref{more-algebra-definition-topological-ring}

\noindent
{\it adic}
见 \ref{formal-spaces-definition-adic-homomorphism}

\noindent
{\it adic}
见 \ref{formal-spaces-definition-types-affine-formal-algebraic-space}

\noindent
{\it 可容许满射}
见 \ref{dga-definition-admissible-ses}

\noindent
{\it 可容许单射}
见 \ref{dga-definition-admissible-ses}

\noindent
{\it 容许关系}
见 \ref{chow-definition-determinant}

\noindent
{\it 可容许短正合列}
见 \ref{dga-definition-admissible-ses}

\noindent
{\it 容许的}
见 \ref{more-algebra-definition-topological-ring}

\noindent
{\it 容许的}
见 \ref{chow-definition-determinant}

\noindent
{\it 仿射 $n$ 维空间（在 $R$ 上）}
见 \ref{constructions-definition-affine-n-space}

\noindent
{\it 仿射 $n$ 维空间（在 $S$ 上）}
见 \ref{constructions-definition-affine-n-space}

\noindent
{\it 仿射吹起代数}
见 \ref{algebra-definition-blow-up}

\noindent
{\it 与 $\mathcal{A}$ 关联的仿射锥}
见 \ref{constructions-definition-cone}

\noindent
{\it 仿射形式代数空间}
见 \ref{formal-spaces-definition-affine-formal-algebraic-space}

\noindent
{\it 仿射概形}
见 \ref{schemes-definition-affine-scheme}

\noindent
{\it 仿射分层数}
见 \ref{more-morphisms-definition-affine-stratification-number}

\noindent
{\it 仿射分层}
见 \ref{more-morphisms-definition-affine-stratification}

\noindent
{\it 仿射簇}
见 \ref{varieties-definition-variety-type}

\noindent
{\it 仿射}
见 \ref{morphisms-definition-affine}

\noindent
{\it 仿射}
见 \ref{spaces-morphisms-definition-affine}

\noindent
{\it 仿射}
见 \ref{stacks-morphisms-definition-affine}

\noindent
{\it 代数 $k$-概形}
见 \ref{varieties-definition-algebraic-scheme}

\noindent
{\it $k$ 在 $K$ 中的代数闭包}
见 \ref{fields-definition-algebraically-closed-in}

\noindent
{\it 代数闭包}
见 \ref{fields-definition-algebraic-closure}

\noindent
{\it 代数扩张}
见 \ref{fields-definition-algebraic}

\noindent
{\it $S$ 上的代数空间}
见 \ref{spaces-definition-algebraic-space}

\noindent
{\it $Z$ 上的代数空间结构}
见 \ref{spaces-properties-definition-reduced-induced-space}

\noindent
{\it $S$ 上的代数叠}
见 \ref{algebraic-definition-algebraic-stack}

\noindent
{\it $Z$ 上的代数叠结构}
见 \ref{stacks-properties-definition-reduced-induced-stack}

\noindent
{\it 代数叠}
见 \ref{stacks-introduction-definition-algebraic-stack}

\noindent
{\it $K$ 中代数闭}
见 \ref{fields-definition-algebraically-closed-in}

\noindent
{\it 代数闭域}
见 \ref{fields-definition-algebraically-closed}

\noindent
{\it 代数无关}
见 \ref{fields-definition-transcendence}

\noindent
{\it 代数的}
见 \ref{fields-definition-algebraic}

\noindent
{\it 代数的}
见 \ref{fields-definition-separable-algebraic}

\noindent
{\it 代数的}
见 \ref{etale-cohomology-definition-algebraic-geometric-point}

\noindent
{\it 代数的}
见 \ref{criteria-definition-algebraic}

\noindent
{\it 几乎余连续的}
见 \ref{sites-definition-almost-cocontinuous}

\noindent
{\it 几乎整于 $R$}
见 \ref{algebra-definition-almost-integral}

\noindent
{\it $X$ 的改造}
见 \ref{morphisms-definition-alteration}

\noindent
{\it $X$ 的改造}
见 \ref{spaces-over-fields-definition-alteration}

\noindent
{\it 交错 {\v C}ech 复形}
见 \ref{cohomology-definition-alternating-cech-complex}

\noindent
{\it 交错 {\v C}ech 复形}
见 \ref{spaces-cohomology-definition-alternating-cech-complex}

\noindent
{\it 并合和}
见 \ref{categories-definition-coproducts}

\noindent
{\it $X$ 上的丰沛可逆模族}
见 \ref{morphisms-definition-family-ample-invertible-modules}

\noindent
{\it 在 $X/S$ 上丰沛}
见 \ref{morphisms-definition-relatively-ample}

\noindent
{\it 在 $X/Y$ 上充足}
见 \ref{spaces-divisors-definition-relatively-ample}

\noindent
{\it 丰沛}
见 \ref{properties-definition-ample}

\noindent
{\it $f$-幂挠模}
见 \ref{more-algebra-definition-f-power-torsion}

\noindent
{\it $R$ 的定义理想}
见 \ref{algebra-definition-ideal-definition}

\noindent
{\it 解析无分歧}
见 \ref{algebra-definition-analytically-unramified}

\noindent
{\it 解析无分歧}
见 \ref{algebra-definition-analytically-unramified}

\noindent
{\it $m$ 的零化子}
见 \ref{algebra-definition-annihilator}

\noindent
{\it $M$ 的零化子}
见 \ref{algebra-definition-annihilator}

\noindent
{\it 零化子}
见 \ref{modules-definition-annihilator-sheaf}

\noindent
{\it 用完美复形逼近的性质成立}
见 \ref{perfect-definition-approximation}

\noindent
{\it 用完美复形逼近的性质成立}
见 \ref{spaces-perfect-definition-approximation}

\noindent
{\it 三元组满足逼近性质}
见 \ref{perfect-definition-approximation-holds}

\noindent
{\it 三元组满足逼近性质}
见 \ref{spaces-perfect-definition-approximation-holds}

\noindent
{\it 算术 Frobenius}
见 \ref{trace-definition-arithmetic-frobenius}

\noindent
{\it Artin}
见 \ref{algebra-definition-artinian}

\noindent
{\it Artin}
见 \ref{homology-definition-Artinian}

\noindent
{\it Artin}
见 \ref{homology-definition-Artinian}

\noindent
{\it Artin}
见 \ref{exercises-definition-Noetherian-space}

\noindent
{\it 相联系的 \'etale 位点}
见 \ref{stacks-sheaves-definition-inherited-topologies}

\noindent
{\it 相联系的仿射 \'etale 位点}
见 \ref{stacks-sheaves-definition-alternative-inherited-topologies}

\noindent
{\it 相联系的仿射 fppf 位点}
见 \ref{stacks-sheaves-definition-alternative-inherited-topologies}

\noindent
{\it 相联系的仿射位点}
见 \ref{stacks-sheaves-definition-alternative}

\noindent
{\it 相联系的仿射光滑位点}
见 \ref{stacks-sheaves-definition-alternative-inherited-topologies}

\noindent
{\it 相联系的仿射 syntomic 位点}
见 \ref{stacks-sheaves-definition-alternative-inherited-topologies}

\noindent
{\it 相联系的仿射 Zariski 位点}
见 \ref{stacks-sheaves-definition-alternative-inherited-topologies}

\noindent
{\it 相联系的 fppf 位点}
见 \ref{stacks-sheaves-definition-inherited-topologies}

\noindent
{\it 伴随分次环}
见 \ref{modules-definition-gamma-star}

\noindent
{\it 相联系的 fppf 拓扑斯态射}
见 \ref{stacks-sheaves-definition-morphism}

\noindent
{\it $X$ 的伴随点}
见 \ref{divisors-definition-associated}

\noindent
{\it 相伴单复形}
见 \ref{homology-definition-associated-simple-complex}

\noindent
{\it 相联系的光滑位点}
见 \ref{stacks-sheaves-definition-inherited-topologies}

\noindent
{\it 相联系的 syntomic 位点}
见 \ref{stacks-sheaves-definition-inherited-topologies}

\noindent
{\it 相伴总复形}
见 \ref{homology-definition-associated-simple-complex}

\noindent
{\it 相联系的 Zariski 位点}
见 \ref{stacks-sheaves-definition-inherited-topologies}

\noindent
{\it 相伴}
见 \ref{algebra-definition-associated}

\noindent
{\it 相伴}
见 \ref{divisors-definition-associated}

\noindent
{\it 相伴元}
见 \ref{algebra-definition-irreducible-prime-element}

\noindent
{\it 相对维数为 $1$ 且至多为结点}
见 \ref{curves-definition-nodal-family}

\noindent
{\it 相对维数为 $1$ 且至多为结点}
见 \ref{spaces-more-morphisms-definition-nodal-family}

\noindent
{\it Atiyah 类}
见 \ref{cotangent-definition-atiyah-class}

\noindent
{\it Atiyah 类}
见 \ref{cotangent-definition-atiyah-class-general}

\noindent
{\it $U$ 朝向 $\mathcal{C}$ 中对象 $X$ 的增广 $\epsilon : U \to X$}
见 \ref{simplicial-definition-augmentation}

\noindent
{\it 自伴随}
见 \ref{more-algebra-definition-auto-ass}

\noindent
{\it $x$ 的自同构函子}
见 \ref{formal-defos-definition-automorphism-functor}

\noindent
{\it $E$ 在 $F$ 上的自同构}
见 \ref{fields-definition-automorphisms}

\noindent
{\it $E/F$ 的自同构}
见 \ref{fields-definition-automorphisms}

\noindent
{\it B\'ezout 整环}
见 \ref{more-algebra-definition-bezout}

\noindent
{\it $F'$ 到 $S$ 的基变换}
见 \ref{spaces-definition-base-change}

\noindent
{\it 基变换}
见 \ref{algebra-definition-base-change}

\noindent
{\it 基变换}
见 \ref{algebra-definition-base-change}

\noindent
{\it 基变换}
见 \ref{schemes-definition-base-change}

\noindent
{\it 基变换}
见 \ref{schemes-definition-base-change}

\noindent
{\it 基变换}
见 \ref{schemes-definition-base-change}

\noindent
{\it 基变换}
见 \ref{groupoids-quotients-definition-base-change}

\noindent
{\it 沿 $f$ 的底扩张}
见 \ref{descent-definition-effective-descent}

\noindent
{\it $X$ 上拓扑的一个基}
见 \ref{topology-definition-base}

\noindent
{\it 基点}
见 \ref{pione-definition-fundamental-group}

\noindent
{\it $X$ 上拓扑的一个基底}
见 \ref{topology-definition-base}

\noindent
{\it $S$ 的大 $\tau$-位点}
见 \ref{etale-cohomology-definition-tau-site}

\noindent
{\it 大 $\tau$-拓扑斯}
见 \ref{etale-cohomology-definition-etale-topos}

\noindent
{\it $S$ 的大平展位点}
见 \ref{topologies-definition-big-small-etale}

\noindent
{\it $S$ 上的大平展位点}
见 \ref{etale-cohomology-definition-big-etale-site}

\noindent
{\it 大平展位点}
见 \ref{topologies-definition-big-etale-site}

\noindent
{\it $S$ 的大仿射平展位点}
见 \ref{topologies-definition-big-small-etale}

\noindent
{\it $S$ 的大仿射 fppf 位点}
见 \ref{topologies-definition-big-small-fppf}

\noindent
{\it $S$ 的大仿射 h 位点}
见 \ref{flat-definition-big-small-h}

\noindent
{\it $S$ 的大仿射 ph 位点}
见 \ref{topologies-definition-big-small-ph}

\noindent
{\it $S$ 的大仿射 pro-平展位点}
见 \ref{proetale-definition-big-small-proetale}

\noindent
{\it $S$ 的大仿射光滑位点}
见 \ref{topologies-definition-big-small-smooth}

\noindent
{\it $S$ 的大仿射 syntomic 位点}
见 \ref{topologies-definition-big-small-syntomic}

\noindent
{\it $S$ 的大仿射 Zariski 位点}
见 \ref{topologies-definition-big-small-Zariski}

\noindent
{\it 大晶体位点}
见 \ref{crystalline-definition-big-crystalline-site}

\noindent
{\it $S$ 的大 fppf 位点}
见 \ref{topologies-definition-big-small-fppf}

\noindent
{\it 大 fppf 位点}
见 \ref{topologies-definition-big-fppf-site}

\noindent
{\it $S$ 的大 h 位点}
见 \ref{flat-definition-big-small-h}

\noindent
{\it 大 h 站点}
见 \ref{flat-definition-big-h-site}

\noindent
{\it $S$ 的大 ph 位点}
见 \ref{topologies-definition-big-small-ph}

\noindent
{\it 大 ph 位点}
见 \ref{topologies-definition-big-ph-site}

\noindent
{\it $S$ 的大 pro-平展位点}
见 \ref{proetale-definition-big-small-proetale}

\noindent
{\it 大 pro-\'etale 位点}
见 \ref{proetale-definition-big-proetale-site}

\noindent
{\it $S$ 的大光滑位点}
见 \ref{topologies-definition-big-small-smooth}

\noindent
{\it 大光滑位点}
见 \ref{topologies-definition-big-smooth-site}

\noindent
{\it $S$ 的大 syntomic 位点}
见 \ref{topologies-definition-big-small-syntomic}

\noindent
{\it 大 syntomic 位点}
见 \ref{topologies-definition-big-syntomic-site}

\noindent
{\it $S$ 的大 Zariski 位点}
见 \ref{topologies-definition-big-small-Zariski}

\noindent
{\it 大 Zariski 位点}
见 \ref{topologies-definition-big-zariski-site}

\noindent
{\it 大的}
见 \ref{etale-cohomology-definition-big-etale-site}

\noindent
{\it 双有理}
见 \ref{morphisms-definition-birational-schemes}

\noindent
{\it 双有理}
见 \ref{morphisms-definition-birational}

\noindent
{\it 双有理}
见 \ref{spaces-morphisms-definition-birational-spaces}

\noindent
{\it 双有理}
见 \ref{decent-spaces-definition-birational}

\noindent
{\it $f$ 的次数为 $p$ 的双变类 $c$}
见 \ref{chow-definition-bivariant-class}

\noindent
{\it $f$ 的次数为 $p$ 的双变类 $c$}
见 \ref{spaces-chow-definition-bivariant-class}

\noindent
{\it $X$ 在 $x$ 处的吹起 $X' \to X$}
见 \ref{spaces-resolve-definition-blowup-at-point}

\noindent
{\it $X$ 沿 $Z$ 的吹起}
见 \ref{divisors-definition-blow-up}

\noindent
{\it $X$ 沿 $Z$ 的吹起}
见 \ref{spaces-divisors-definition-blow-up}

\noindent
{\it $X$ 关于理想层 $\mathcal{I}$ 的吹起}
见 \ref{divisors-definition-blow-up}

\noindent
{\it $X$ 关于理想层 $\mathcal{I}$ 的吹起}
见 \ref{spaces-divisors-definition-blow-up}

\noindent
{\it 吹起代数}
见 \ref{algebra-definition-blow-up}

\noindent
{\it 上有界}
见 \ref{homology-definition-bounded-ss}

\noindent
{\it 上有界}
见 \ref{derived-definition-complexes-notation}

\noindent
{\it 下有界}
见 \ref{homology-definition-bounded-ss}

\noindent
{\it 下有界}
见 \ref{derived-definition-complexes-notation}

\noindent
{\it 有界导出范畴}
见 \ref{derived-definition-unbounded-derived-category}

\noindent
{\it 有界滤过导出范畴}
见 \ref{derived-definition-filtered-derived-bounded}

\noindent
{\it 有界}
见 \ref{homology-definition-bounded-ss}

\noindent
{\it 有界}
见 \ref{derived-definition-complexes-notation}

\noindent
{\it 界定 $f$ 的纤维次数}
见 \ref{morphisms-definition-universally-bounded}

\noindent
{\it Bourbaki-固有的}
见 \ref{topology-definition-proper-map}

\noindent
{\it Brauer 群}
见 \ref{brauer-definition-brauer-group}

\noindent
{\it Brauer 群}
见 \ref{etale-cohomology-definition-brauer-group}

\noindent
{\it 典范下降数据}
见 \ref{stacks-definition-effective-descent-datum}

\noindent
{\it 典范下降数据}
见 \ref{descent-definition-descent-datum-effective-quasi-coherent}

\noindent
{\it 典范下降数据}
见 \ref{descent-definition-effective}

\noindent
{\it 典范下降数据}
见 \ref{descent-definition-effective-family}

\noindent
{\it 典范下降数据}
见 \ref{spaces-descent-definition-descent-datum-effective-quasi-coherent}

\noindent
{\it 典范下降数据}
见 \ref{spaces-descent-definition-effective}

\noindent
{\it 典范下降数据}
见 \ref{spaces-descent-definition-effective-family}

\noindent
{\it $T$ 上的典范概形结构}
见 \ref{morphisms-definition-scheme-structure-connected-component}

\noindent
{\it 典范截面}
见 \ref{divisors-definition-invertible-sheaf-effective-Cartier-divisor}

\noindent
{\it 典范拓扑}
见 \ref{sites-definition-canonical-topology}

\noindent
{\it Cartan--Eilenberg 分解}
见 \ref{derived-definition-cartan-eilenberg}

\noindent
{\it 笛卡尔的}
见 \ref{categories-definition-cartesian}

\noindent
{\it 笛卡尔的}
见 \ref{groupoids-definition-cartesian-morphism}

\noindent
{\it 笛卡尔的}
见 \ref{spaces-simplicial-definition-cartesian-sheaf}

\noindent
{\it 笛卡尔的}
见 \ref{spaces-simplicial-definition-cartesian-sheaf}

\noindent
{\it 笛卡尔的}
见 \ref{spaces-simplicial-definition-cartesian-sheaf}

\noindent
{\it 笛卡尔的}
见 \ref{spaces-simplicial-definition-cartesian-sheaf}

\noindent
{\it 笛卡尔的}
见 \ref{spaces-simplicial-definition-cartesian-derived}

\noindent
{\it 笛卡尔的}
见 \ref{spaces-simplicial-definition-cartesian-derived-modules}

\noindent
{\it 笛卡尔的}
见 \ref{spaces-simplicial-definition-cartesian-morphism}

\noindent
{\it Cartier 除子}
见 \ref{exercises-definition-divisor}

\noindent
{\it $\mathcal{C}$ 中的范畴模空间}
见 \ref{stacks-more-morphisms-definition-categorical-quotient}

\noindent
{\it 范畴模空间}
见 \ref{stacks-more-morphisms-definition-categorical-quotient}

\noindent
{\it $\mathcal{C}$ 中的范畴商}
见 \ref{groupoids-quotients-definition-categorical}

\noindent
{\it 概形中的范畴商}
见 \ref{groupoids-quotients-definition-categorical}

\noindent
{\it 概形范畴中的范畴商}
见 \ref{groupoids-quotients-definition-categorical}

\noindent
{\it 范畴商}
见 \ref{groupoids-quotients-definition-categorical}

\noindent
{\it 范畴紧致}
见 \ref{categories-definition-compact-object}

\noindent
{\it $\mathcal{F}$ 的形式对象范畴 $\widehat{\mathcal{F}}$}
见 \ref{formal-defos-definition-formal-objects}

\noindent
{\it 在 $\mathcal{C}$ 上余纤维化为群胚的范畴}
见 \ref{formal-defos-definition-category-cofibred-groupoids}

\noindent
{\it 离散范畴纤维化范畴}
见 \ref{categories-definition-category-fibred-sets}

\noindent
{\it 集合胚纤维化范畴}
见 \ref{categories-definition-category-fibred-setoids}

\noindent
{\it 集合纤维化范畴}
见 \ref{categories-definition-category-fibred-sets}

\noindent
{\it （上链）复形范畴}
见 \ref{derived-definition-complexes-notation}

\noindent
{\it $\mathcal{A}$ 的复形范畴}
见 \ref{dga-definition-homotopy-category-of-dga-category}

\noindent
{\it $\mathcal{C}$ 的余单纯对象范畴}
见 \ref{simplicial-definition-cosimplicial-object}

\noindent
{\it $\mathcal{A}$ 的有限滤过对象范畴}
见 \ref{derived-definition-finite-filtered}

\noindent
{\it $\mathcal{A}$ 的分次对象范畴}
见 \ref{homology-definition-graded}

\noindent
{\it $\mathcal{C}$ 上函子中的群胚范畴}
见 \ref{formal-defos-definition-groupoid-in-functors}

\noindent
{\it 集合层范畴}
见 \ref{etale-cohomology-definition-category-sheaves}

\noindent
{\it $\mathcal{C}$ 的单纯对象范畴}
见 \ref{simplicial-definition-simplicial-object}

\noindent
{\it 范畴}
见 \ref{categories-definition-category}

\noindent
{\it 链状}
见 \ref{topology-definition-catenary}

\noindent
{\it 链状}
见 \ref{algebra-definition-catenary}

\noindent
{\it 链状}
见 \ref{properties-definition-catenary}

\noindent
{\it 链状}
见 \ref{decent-spaces-definition-catenary}

\noindent
{\it 链状}
见 \ref{exercises-definition-catenary}

\noindent
{\it 链状}
见 \ref{exercises-definition-catenary}

\noindent
{\it 有中心}
见 \ref{algebra-definition-valuation-ring}

\noindent
{\it 中心}
见 \ref{divisors-definition-blow-up}

\noindent
{\it 中心}
见 \ref{spaces-divisors-definition-blow-up}

\noindent
{\it 中心的}
见 \ref{brauer-definition-central}

\noindent
{\it 不可约闭子集链}
见 \ref{topology-definition-Krull}

\noindent
{\it 素理想链}
见 \ref{algebra-definition-chain-primes}

\noindent
{\it $\mathcal{X}'$ 的基变换}
见 \ref{algebraic-definition-change-of-base}

\noindent
{\it 特征}
见 \ref{fields-definition-characteristic}

\noindent
{\it $X$ 上 $\mathcal{E}$ 的陈类}
见 \ref{chow-definition-chern-classes}

\noindent
{\it $E$ 的陈类已定义}
见 \ref{chow-definition-defined-on-perfect}

\noindent
{\it 拉回选择}
见 \ref{categories-definition-pullback-functor-fibred-category}

\noindent
{\it Chow 上同调}
见 \ref{chow-definition-chow-cohomology}

\noindent
{\it Chow 上同调}
见 \ref{spaces-chow-definition-chow-cohomology}

\noindent
{\it $X$ 上模有理等价的 $k$ 维循环 Chow 群}
见 \ref{chow-definition-rational-equivalence}

\noindent
{\it $X$ 上模有理等价的 $k$ 维循环 Chow 群}
见 \ref{spaces-chow-definition-rational-equivalence}

\noindent
{\it $X$ 上 $k$ 维循环的 Chow 群}
见 \ref{chow-definition-rational-equivalence}

\noindent
{\it $X$ 上 $k$ 维循环的 Chow 群}
见 \ref{spaces-chow-definition-rational-equivalence}

\noindent
{\it $A$ 的类群}
见 \ref{exercises-definition-class-group}

\noindent
{\it 经典情形}
见 \ref{formal-defos-definition-CLambda}

\noindent
{\it 经典生成元}
见 \ref{derived-definition-generators}

\noindent
{\it 经典 Weil 上同调理论}
见 \ref{weil-definition-weil-cohomology-theory-classical}

\noindent
{\it 经典的}
见 \ref{formal-spaces-definition-types-affine-formal-algebraic-space}

\noindent
{\it 环化空间的闭浸入}
见 \ref{modules-definition-closed-immersion}

\noindent
{\it 闭浸入}
见 \ref{sites-definition-immersion-topoi}

\noindent
{\it 闭浸入}
见 \ref{schemes-definition-closed-immersion-locally-ringed-spaces}

\noindent
{\it 闭浸入}
见 \ref{schemes-definition-immersion}

\noindent
{\it 闭浸入}
见 \ref{spaces-definition-immersion}

\noindent
{\it 闭浸入}
见 \ref{formal-spaces-definition-closed-immersion}

\noindent
{\it 闭浸入}
见 \ref{stacks-properties-definition-immersion}

\noindent
{\it 闭子群概形}
见 \ref{groupoids-definition-closed-subgroup-scheme}

\noindent
{\it 闭子概形}
见 \ref{schemes-definition-immersion}

\noindent
{\it $X$ 的与理想层 $\mathcal{I}$ 相伴的闭子空间}
见 \ref{schemes-definition-closed-subspace}

\noindent
{\it 闭子空间}
见 \ref{spaces-definition-immersion}

\noindent
{\it 闭子叠}
见 \ref{stacks-properties-definition-substacks}

\noindent
{\it 闭子拓扑斯}
见 \ref{sites-definition-closed-subtopos}

\noindent
{\it 闭的}
见 \ref{topology-definition-proper-map}

\noindent
{\it 闭的}
见 \ref{spaces-morphisms-definition-closed}

\noindent
{\it 闭的}
见 \ref{stacks-morphisms-definition-closed}

\noindent
{\it 闭的}
见 \ref{exercises-definition-closed}

\noindent
{\it 概形中的粗商}
见 \ref{groupoids-quotients-definition-coarse}

\noindent
{\it 粗商}
见 \ref{groupoids-quotients-definition-coarse}

\noindent
{\it 更粗}
见 \ref{sites-definition-finer}

\noindent
{\it 余笛卡尔}
见 \ref{categories-definition-cocartesian}

\noindent
{\it 余连续的}
见 \ref{sites-definition-cocontinuous}

\noindent
{\it 余循环条件}
见 \ref{stacks-definition-descent-data}

\noindent
{\it 余循环条件}
见 \ref{descent-definition-descent-datum-quasi-coherent}

\noindent
{\it 余循环条件}
见 \ref{descent-definition-descent-datum-modules}

\noindent
{\it 余循环条件}
见 \ref{descent-definition-descent-datum}

\noindent
{\it 余循环条件}
见 \ref{etale-cohomology-definition-descent-datum}

\noindent
{\it 余循环条件}
见 \ref{spaces-descent-definition-descent-datum-quasi-coherent}

\noindent
{\it 余循环条件}
见 \ref{spaces-descent-definition-descent-datum}

\noindent
{\it 余维数}
见 \ref{topology-definition-codimension}

\noindent
{\it 反向}
见 \ref{categories-definition-codirected}

\noindent
{\it 反向}
见 \ref{categories-definition-codirected}

\noindent
{\it 系数环}
见 \ref{algebra-definition-coefficient-ring}

\noindent
{\it 余等化子}
见 \ref{categories-definition-coequalizers}

\noindent
{\it 余滤过}
见 \ref{categories-definition-codirected}

\noindent
{\it 余滤过}
见 \ref{categories-definition-codirected}

\noindent
{\it 共尾}
见 \ref{categories-definition-cofinal}

\noindent
{\it Cohen 环}
见 \ref{algebra-definition-cohen-ring}

\noindent
{\it 在 $x$ 处 Cohen--Macaulay}
见 \ref{more-morphisms-definition-CM}

\noindent
{\it 在 $x$ 处 Cohen--Macaulay}
见 \ref{spaces-more-morphisms-definition-CM}

\noindent
{\it Cohen--Macaulay 态射}
见 \ref{more-morphisms-definition-CM}

\noindent
{\it Cohen--Macaulay 态射}
见 \ref{spaces-more-morphisms-definition-CM}

\noindent
{\it Cohen--Macaulay}
见 \ref{algebra-definition-CM}

\noindent
{\it Cohen--Macaulay}
见 \ref{algebra-definition-module-CM}

\noindent
{\it Cohen--Macaulay}
见 \ref{algebra-definition-local-ring-CM}

\noindent
{\it Cohen--Macaulay}
见 \ref{algebra-definition-ring-CM}

\noindent
{\it Cohen--Macaulay}
见 \ref{properties-definition-Cohen-Macaulay}

\noindent
{\it Cohen--Macaulay}
见 \ref{coherent-definition-Cohen-Macaulay}

\noindent
{\it Cohen--Macaulay}
见 \ref{exercises-definition-CM}

\noindent
{\it 凝聚 $\mathcal{O}_X$-模}
见 \ref{modules-definition-coherent}

\noindent
{\it 凝聚模}
见 \ref{algebra-definition-coherent}

\noindent
{\it 凝聚环}
见 \ref{algebra-definition-coherent}

\noindent
{\it 凝聚}
见 \ref{sites-modules-definition-site-local}

\noindent
{\it 凝聚}
见 \ref{spaces-cohomology-definition-coherent}

\noindent
{\it 凝聚}
见 \ref{stacks-cohomology-definition-coherent}

\noindent
{\it 凝聚}
见 \ref{exercises-definition-coherent}

\noindent
{\it 上同调 $\delta$-函子}
见 \ref{homology-definition-cohomological-delta-functor}

\noindent
{\it $f$ 的上同调维数}
见 \ref{etale-cohomology-definition-cd-f}

\noindent
{\it $I$ 在 $A$ 中的上同调维数}
见 \ref{local-cohomology-definition-cd}

\noindent
{\it $X$ 的上同调维数}
见 \ref{etale-cohomology-definition-cd}

\noindent
{\it 上同调的}
见 \ref{derived-definition-homological}

\noindent
{\it 上同调模}
见 \ref{chow-definition-periodic-complex}

\noindent
{\it 上同调模}
见 \ref{chow-definition-periodic-complex}

\noindent
{\it $K$ 的带紧支撑上同调}
见 \ref{more-etale-definition-cohomology-compact-support}

\noindent
{\it $f$ 的余像}
见 \ref{homology-definition-kernel}

\noindent
{\it 余核}
见 \ref{homology-definition-kernel}

\noindent
{\it 余极限}
见 \ref{categories-definition-colimit}

\noindent
{\it 余极限}
见 \ref{exercises-definition-colimit}

\noindent
{\it 组合等价}
见 \ref{sites-definition-combinatorial-tautological}

\noindent
{\it 可交换}
见 \ref{dga-definition-cdga}

\noindent
{\it 紧对象}
见 \ref{derived-definition-compact-object}

\noindent
{\it 紧生成}
见 \ref{derived-definition-compactly-generated}

\noindent
{\it $K$ 的紧支撑上同调}
见 \ref{more-etale-definition-cohomology-compact-support}

\noindent
{\it 与微分分次结构相容}
见 \ref{dpa-definition-divided-powers-dga}

\noindent
{\it 与三角结构相容}
见 \ref{derived-definition-localization}

\noindent
{\it $x$ 处 $\mathcal{F}/X/S$ 的完全 d\'evissage}
见 \ref{flat-definition-complete-devissage-at-x}

\noindent
{\it $s$ 上 $\mathcal{F}/X/S$ 的完全 d\'evissage}
见 \ref{flat-definition-complete-devissage}

\noindent
{\it $N/S/R$ 在 $\mathfrak q$ 处的完全 d\'evissage}
见 \ref{flat-definition-complete-devissage-at-x-algebra}

\noindent
{\it $\mathfrak r$ 上 $N/S/R$ 的完全 d\'evissage}
见 \ref{flat-definition-complete-devissage-algebra}

\noindent
{\it （在 $k$ 上的）完全交}
见 \ref{algebra-definition-lci-local-ring}

\noindent
{\it 完全交}
见 \ref{dpa-definition-lci}

\noindent
{\it 完备局部环}
见 \ref{algebra-definition-complete-local-ring}

\noindent
{\it 完备化主局部化}
见 \ref{restricted-definition-completed-principal-localization}

\noindent
{\it 完备张量积}
见 \ref{formal-spaces-definition-toplogy-tensor-product}

\noindent
{\it 完全分解}
见 \ref{more-morphisms-definition-cd-morphism}

\noindent
{\it 完全分解}
见 \ref{more-morphisms-definition-cd-morphism}

\noindent
{\it 完全正规}
见 \ref{algebra-definition-almost-integral}

\noindent
{\it $(U, R, s, t, c)$ 的完备化 $(U, R, s, t, c)^{\wedge}$}
见 \ref{formal-defos-definition-completion-groupoid-in-functors}

\noindent
{\it $\mathcal{F}$ 的完备化}
见 \ref{formal-defos-definition-completion}

\noindent
{\it $X$ 沿 $T$ 的完备化}
见 \ref{formal-spaces-definition-completion}

\noindent
{\it $X$ 沿 $T$ 的完备化}
见 \ref{formal-spaces-definition-completion-formal-algebraic-space}

\noindent
{\it $X$ 沿 $Z$ 的完备化}
见 \ref{formal-spaces-definition-completion-subspace}

\noindent
{\it 复形}
见 \ref{homology-definition-exact}

\noindent
{\it 复合 $f\circ g$}
见 \ref{sites-definition-topos}

\noindent
{\it $\varphi$ 与 $\psi$ 的复合}
见 \ref{sheaves-definition-composition-f-maps}

\noindent
{\it 芽态射的复合}
见 \ref{descent-definition-germs}

\noindent
{\it 环化位点态射的复合}
见 \ref{sites-modules-definition-ringed-site}

\noindent
{\it 赋环空间态射的复合}
见 \ref{sheaves-definition-composition-maps-ringed-spaces}

\noindent
{\it 环化拓扑斯态射的复合}
见 \ref{sites-modules-definition-ringed-topos}

\noindent
{\it 复合}
见 \ref{categories-definition-2-category}

\noindent
{\it 复合}
见 \ref{sites-definition-composition-morphisms-sites}

\noindent
{\it $K$ 与 $L$ 在 $\Omega$ 中的合成域}
见 \ref{fields-definition-compositum}

\noindent
{\it 计算出}
见 \ref{derived-definition-computes}

\noindent
{\it 计算出}
见 \ref{derived-definition-computes}

\noindent
{\it 条件 (RS)}
见 \ref{formal-defos-definition-RS}

\noindent
{\it 条件 (RS)}
见 \ref{artin-definition-RS}

\noindent
{\it 条件 (RS*)}
见 \ref{artin-definition-RS-star}

\noindent
{\it 条件 (S1) 与 (S2)}
见 \ref{formal-defos-definition-S1-S2}

\noindent
{\it 锥 $\pi : C \to S$（在 $S$ 上）}
见 \ref{constructions-definition-abstract-cone}

\noindent
{\it 与 $\mathcal{A}$ 关联的锥}
见 \ref{constructions-definition-cone}

\noindent
{\it 锥}
见 \ref{derived-definition-cone}

\noindent
{\it 锥}
见 \ref{dga-definition-cone}

\noindent
{\it 锥}
见 \ref{sdga-definition-cone}

\noindent
{\it 连通分支}
见 \ref{topology-definition-connected-components}

\noindent
{\it 连通分支}
见 \ref{exercises-definition-connected-component}

\noindent
{\it 连通}
见 \ref{categories-definition-category-connected}

\noindent
{\it 连通}
见 \ref{topology-definition-connected-components}

\noindent
{\it 连通}
见 \ref{exercises-definition-connected-component}

\noindent
{\it $Z$ 在 $X$ 中的余法代数 $\mathcal{C}_{Z/X, *}$}
见 \ref{divisors-definition-conormal-sheaf}

\noindent
{\it $Z$ 在 $X$ 中的余法代数 $\mathcal{C}_{Z/X, *}$}
见 \ref{spaces-more-morphisms-definition-conormal-algebra}

\noindent
{\it $f$ 的余法代数}
见 \ref{divisors-definition-conormal-sheaf}

\noindent
{\it $i$ 的余法代数}
见 \ref{spaces-more-morphisms-definition-conormal-algebra}

\noindent
{\it 余法模}
见 \ref{algebra-definition-universal-thickening}

\noindent
{\it $Z$ 在 $X$ 中的余法层 $\mathcal{C}_{Z/X}$}
见 \ref{morphisms-definition-conormal-sheaf}

\noindent
{\it $Z$ 在 $X$ 中的余法层 $\mathcal{C}_{Z/X}$}
见 \ref{spaces-more-morphisms-definition-conormal-sheaf}

\noindent
{\it $i$ 的余法层}
见 \ref{morphisms-definition-conormal-sheaf}

\noindent
{\it $i$ 的余法层}
见 \ref{spaces-more-morphisms-definition-conormal-sheaf}

\noindent
{\it $Z$ 在 $X$ 上的余法层}
见 \ref{more-morphisms-definition-universal-thickening}

\noindent
{\it $Z$ 在 $X$ 上的余法层}
见 \ref{spaces-more-morphisms-definition-universal-thickening}

\noindent
{\it 保守}
见 \ref{sites-definition-enough-points}

\noindent
{\it 取值为 $A$ 的常值预层}
见 \ref{sheaves-definition-constant-presheaf}

\noindent
{\it 值为 $A$ 的常值层}
见 \ref{sheaves-definition-constant-sheaf}

\noindent
{\it 值为 $A$ 的常值层}
见 \ref{etale-cohomology-definition-finite-locally-constant}

\noindent
{\it 值为 $E$ 的常值层}
见 \ref{etale-cohomology-definition-finite-locally-constant}

\noindent
{\it 值为 $M$ 的常值层}
见 \ref{etale-cohomology-definition-finite-locally-constant}

\noindent
{\it 常值层}
见 \ref{sites-modules-definition-locally-constant}

\noindent
{\it 常值层}
见 \ref{etale-cohomology-definition-additive-sheaf}

\noindent
{\it 常值层}
见 \ref{etale-cohomology-definition-finite-locally-constant}

\noindent
{\it 常值层}
见 \ref{etale-cohomology-definition-finite-locally-constant}

\noindent
{\it 常值层}
见 \ref{etale-cohomology-definition-finite-locally-constant}

\noindent
{\it 可构造 $\Lambda$-层}
见 \ref{proetale-definition-adic}

\noindent
{\it 可构造}
见 \ref{topology-definition-constructible}

\noindent
{\it 可构造}
见 \ref{etale-cohomology-definition-constructible}

\noindent
{\it 可构造}
见 \ref{etale-cohomology-definition-constructible}

\noindent
{\it 可构造}
见 \ref{etale-cohomology-definition-constructible}

\noindent
{\it 可构造}
见 \ref{proetale-definition-constructible}

\noindent
{\it 可构造}
见 \ref{proetale-definition-Dbc}

\noindent
{\it $x$ 的内容理想}
见 \ref{more-algebra-definition-content-ideal}

\noindent
{\it 连续群上同调群}
见 \ref{etale-cohomology-definition-galois-cohomology}

\noindent
{\it 连续的}
见 \ref{sites-definition-continuous}

\noindent
{\it 反变}
见 \ref{categories-definition-contravariant}

\noindent
{\it 收敛于 $H^*(K^\bullet)$}
见 \ref{homology-definition-filtered-complex-ss-converges}

\noindent
{\it 收敛于 $H^n(\text{Tot}(K^{\bullet, \bullet}))$}
见 \ref{homology-definition-ss-double-complex-converge}

\noindent
{\it 收敛于 $H^n(\text{Tot}(K^{\bullet, \bullet}))$}
见 \ref{homology-definition-ss-double-complex-converge}

\noindent
{\it 余积}
见 \ref{categories-definition-coproducts}

\noindent
{\it 余积}
见 \ref{categories-definition-coproduct}

\noindent
{\it 余正则}
见 \ref{homology-definition-bounded-ss}

\noindent
{\it 余单纯阿贝尔群}
见 \ref{simplicial-definition-cosimplicial-object}

\noindent
{\it $\mathcal{C}$ 的余单纯对象 $U$}
见 \ref{simplicial-definition-cosimplicial-object}

\noindent
{\it 余单纯集}
见 \ref{simplicial-definition-cosimplicial-object}

\noindent
{\it $X$ 在 $Y$ 上的余切复形 $L_{X/Y}$}
见 \ref{cotangent-definition-cotangent-morphism-schemes}

\noindent
{\it $X$ 在 $Y$ 上的余切复形 $L_{X/Y}$}
见 \ref{cotangent-definition-cotangent-morphism-spaces}

\noindent
{\it 余切复形}
见 \ref{cotangent-definition-cotangent-complex-ring-map}

\noindent
{\it 余切复形}
见 \ref{cotangent-definition-cotangent-complex-morphism-sheaves-rings}

\noindent
{\it 余切复形}
见 \ref{cotangent-definition-cotangent-complex-morphism-ringed-spaces}

\noindent
{\it 余切复形}
见 \ref{cotangent-definition-cotangent-complex-morphism-ringed-topoi}

\noindent
{\it 可数指标的}
见 \ref{formal-spaces-definition-countable}

\noindent
{\it $\mathcal{C}$ 的覆盖}
见 \ref{sites-definition-site}

\noindent
{\it 覆盖}
见 \ref{etale-cohomology-definition-site}

\noindent
{\it 覆盖}
见 \ref{hypercovering-definition-covering-SR}

\noindent
{\it 覆盖}
见 \ref{hypercovering-definition-covering-SR}

\noindent
{\it 覆盖 $F$}
见 \ref{schemes-definition-representable-by-open-immersions}

\noindent
{\it $\mathcal{O}_{X/S}$-模晶体}
见 \ref{crystalline-definition-modules}

\noindent
{\it 有限局部自由模中的晶体}
见 \ref{crystalline-definition-crystal-quasi-coherent-modules}

\noindent
{\it 拟凝聚模中的晶体}
见 \ref{crystalline-definition-crystal-quasi-coherent-modules}

\noindent
{\it 晶体位点}
见 \ref{crystalline-definition-crystalline-site}

\noindent
{\it 曲线}
见 \ref{varieties-definition-curve}

\noindent
{\it 曲线}
见 \ref{etale-cohomology-definition-variety}

\noindent
{\it $X$ 上的循环}
见 \ref{chow-definition-cycles}

\noindent
{\it $X$ 上的循环}
见 \ref{spaces-chow-definition-cycles}

\noindent
{\it $\mathcal{B}$ 相对于 $\mathcal{A}$ 的 de Rham 复形}
见 \ref{modules-definition-de-rham-complex}

\noindent
{\it $Y \subset X$（在 $S$ 上）的对数极点德拉姆复形}
见 \ref{derham-definition-log-complex}

\noindent
{\it 对 $Y \subset X$（在 $S$ 上）定义了对数极点德拉姆复形}
见 \ref{derham-definition-local-product}

\noindent
{\it de Rham 复形}
见 \ref{modules-definition-de-rham-complex-morphism-ringed-spaces}

\noindent
{\it decent}
见 \ref{decent-spaces-definition-very-reasonable}

\noindent
{\it decent}
见 \ref{decent-spaces-definition-relative-conditions}

\noindent
{\it decent}
见 \ref{stacks-morphisms-definition-decent}

\noindent
{\it $\mathfrak m$ 的分解群}
见 \ref{more-algebra-definition-decomposition-inertia}

\noindent
{\it 递减滤过}
见 \ref{homology-definition-filtered}

\noindent
{\it Dedekind 整环}
见 \ref{algebra-definition-dedekind-domain}

\noindent
{\it 在点 $x \in X$ 处有定义}
见 \ref{morphisms-definition-domain-of-definition}

\noindent
{\it 在点 $x \in |X|$ 处有定义}
见 \ref{spaces-morphisms-definition-domain-of-definition}

\noindent
{\it 定义一个结点奇点}
见 \ref{curves-definition-multicross}

\noindent
{\it 定义一个结点奇点}
见 \ref{curves-definition-nodal}

\noindent
{\it 定义一个有理奇点}
见 \ref{resolve-definition-reduce-to-rational}

\noindent
{\it 形变范畴}
见 \ref{formal-defos-definition-deformation-category}

\noindent
{\it $x$ 的退化度}
见 \ref{simplicial-definition-terminology-simplicial-sets}

\noindent
{\it 在 $E_r$ 处退化}
见 \ref{homology-definition-limit-spectral-sequence}

\noindent
{\it 退化}
见 \ref{simplicial-definition-terminology-simplicial-sets}

\noindent
{\it $\mathcal{X}$ 在 $\mathcal{Y}$ 上次数 $d$ 的有限 Hilbert 叠}
见 \ref{examples-stacks-definition-hilbert-d-stack}

\noindent
{\it $X$ 在 $Y$ 上的次数}
见 \ref{morphisms-definition-degree}

\noindent
{\it $X$ 在 $Y$ 上的次数}
见 \ref{spaces-over-fields-definition-degree}

\noindent
{\it $Z$ 相对于 $\mathcal{L}$ 的次数}
见 \ref{varieties-definition-degree}

\noindent
{\it 零维循环的次数}
见 \ref{chow-definition-degree-zero-cycle}

\noindent
{\it 零维循环的次数}
见 \ref{spaces-chow-definition-degree-zero-cycle}

\noindent
{\it 不可分度}
见 \ref{fields-definition-insep-degree}

\noindent
{\it 次数}
见 \ref{fields-definition-degree}

\noindent
{\it 次数}
见 \ref{morphisms-definition-finite-locally-free}

\noindent
{\it 次数}
见 \ref{varieties-definition-degree-invertible-sheaf}

\noindent
{\it 次数}
见 \ref{varieties-definition-degree-invertible-sheaf}

\noindent
{\it 次数}
见 \ref{spaces-morphisms-definition-finite-locally-free}

\noindent
{\it Deligne--Mumford 叠}
见 \ref{algebraic-definition-deligne-mumford}

\noindent
{\it 某点处的深度 $k$}
见 \ref{coherent-definition-depth}

\noindent
{\it 某点处的深度 $k$}
见 \ref{coherent-definition-depth}

\noindent
{\it 深度}
见 \ref{algebra-definition-depth}

\noindent
{\it 深度}
见 \ref{local-cohomology-definition-depth-complex}

\noindent
{\it 导子}
见 \ref{algebra-definition-derivation}

\noindent
{\it $(\mathcal{A}, \text{d})$ 的导出范畴}
见 \ref{sdga-definition-derived-category}

\noindent
{\it $(A, \text{d})$ 的导出范畴}
见 \ref{dga-definition-unbounded-derived-category}

\noindent
{\it $\mathcal{A}$ 的导出范畴}
见 \ref{derived-definition-unbounded-derived-category}

\noindent
{\it 上同调层拟凝聚的 $\mathcal{O}_\mathcal{X}$-模导出范畴}
见 \ref{stacks-perfect-definition-derived}

\noindent
{\it 具有拟凝聚上同调层的 $\mathcal{O}_X$-模导出范畴}
见 \ref{spaces-perfect-definition-derived-quasi-coherent}

\noindent
{\it 导出余极限}
见 \ref{derived-definition-derived-colimit}

\noindent
{\it 关于 $\mathcal{I}$ 导出完备}
见 \ref{algebraization-definition-derived-complete}

\noindent
{\it 关于 $I$ 导出完备}
见 \ref{more-algebra-definition-derived-complete}

\noindent
{\it 关于 $I$ 导出完备}
见 \ref{more-algebra-definition-derived-complete}

\noindent
{\it 导出等价}
见 \ref{equiv-definition-derived-equivalent}

\noindent
{\it 导出内 Hom}
见 \ref{sdga-definition-pushforward}

\noindent
{\it 导出极限}
见 \ref{derived-definition-derived-limit}

\noindent
{\it 导出拉回}
见 \ref{sdga-definition-pullback}

\noindent
{\it 导出推前}
见 \ref{sdga-definition-pushforward}

\noindent
{\it 导出张量积}
见 \ref{more-algebra-definition-derived-tor}

\noindent
{\it 导出张量积}
见 \ref{cohomology-definition-derived-tor}

\noindent
{\it 导出张量积}
见 \ref{sites-cohomology-definition-derived-tor}

\noindent
{\it 导出张量积}
见 \ref{sdga-definition-pullback}

\noindent
{\it 拟凝聚层下降数据 $(\mathcal{F}_i, \varphi_{ij})$}
见 \ref{descent-definition-descent-datum-quasi-coherent}

\noindent
{\it 拟凝聚层下降数据 $(\mathcal{F}_i, \varphi_{ij})$}
见 \ref{spaces-descent-definition-descent-datum-quasi-coherent}

\noindent
{\it 模相对于 $R \to A$ 的下降数据 $(N, \varphi)$}
见 \ref{descent-definition-descent-datum-modules}

\noindent
{\it 相对于态射族 $\{X_i \to S\}$ 的下降数据 $(V_i, \varphi_{ij})$}
见 \ref{descent-definition-descent-datum-for-family-of-morphisms}

\noindent
{\it 相对于族 $\{X_i \to X\}$ 的下降数据 $(V_i, \varphi_{ij})$}
见 \ref{spaces-descent-definition-descent-datum-for-family-of-morphisms}

\noindent
{\it $\mathcal{S}$ 中相对于族 $\{f_i : U_i \to U\}$ 的下降数据 $(X_i, \varphi_{ij})$}
见 \ref{stacks-definition-descent-data}

\noindent
{\it $V/X/S$ 的下降数据}
见 \ref{descent-definition-descent-datum}

\noindent
{\it $V/Y/X$ 的下降数据}
见 \ref{spaces-descent-definition-descent-datum}

\noindent
{\it 相对于 $X \to S$ 的下降数据}
见 \ref{descent-definition-descent-datum}

\noindent
{\it 相对于 $Y \to X$ 的下降数据}
见 \ref{spaces-descent-definition-descent-datum}

\noindent
{\it 下降数据}
见 \ref{etale-cohomology-definition-descent-datum}

\noindent
{\it 下降数据}
见 \ref{etale-cohomology-definition-descent-datum-modules}

\noindent
{\it 模的下降态射}
见 \ref{descent-definition-effective-descent}

\noindent
{\it $(M, \varphi, \psi)$ 的行列式}
见 \ref{chow-definition-periodic-determinant}

\noindent
{\it 有限长度 $R$-模 $M$ 的行列式}
见 \ref{chow-definition-determinant}

\noindent
{\it $\varphi$ 的微分 $d \varphi : T \mathcal{F} \to T \mathcal{G}$}
见 \ref{formal-defos-definition-differential}

\noindent
{\it 微分分次 $(\mathcal{A}, \mathcal{B})$-双模}
见 \ref{sdga-definition-dg-bimodule}

\noindent
{\it 微分分次 $(A, B)$-双模}
见 \ref{dga-definition-bimodule}

\noindent
{\it 微分分次 $\mathcal{A}$-模}
见 \ref{sdga-definition-dgm}

\noindent
{\it $R$ 上的微分分次代数}
见 \ref{dga-definition-dga}

\noindent
{\it $R$ 上的微分分次范畴 $\mathcal{A}$}
见 \ref{dga-definition-dga-category}

\noindent
{\it 微分分次直和}
见 \ref{dga-definition-dg-direct-sum}

\noindent
{\it 微分分次模}
见 \ref{dga-definition-dgm}

\noindent
{\it 微分分次模}
见 \ref{sdga-definition-dgm}

\noindent
{\it 微分对象}
见 \ref{homology-definition-differential-object}

\noindent
{\it $k$ 阶微分算子 $D : \mathcal{F} \to \mathcal{G}$}
见 \ref{modules-definition-differential-operators}

\noindent
{\it $k$ 阶微分算子 $D : \mathcal{F} \to \mathcal{G}$}
见 \ref{sites-modules-definition-differential-operators}

\noindent
{\it $k$ 阶微分算子 $D : M \to N$}
见 \ref{algebra-definition-differential-operators}

\noindent
{\it $X/S$ 上的 $k$ 阶微分算子}
见 \ref{modules-definition-relative-differential-operators}

\noindent
{\it 不同理想}
见 \ref{discriminant-definition-different}

\noindent
{\it 维数函数}
见 \ref{topology-definition-dimension-function}

\noindent
{\it $\mathcal{X}$ 在 $x$ 处的维数}
见 \ref{stacks-properties-definition-dimension-at-point}

\noindent
{\it $X$ 在 $x$ 处的维数}
见 \ref{properties-definition-dimension}

\noindent
{\it $X$ 在 $x$ 处的维数}
见 \ref{spaces-properties-definition-dimension-at-point}

\noindent
{\it $\mathcal{X}$ 在 $x$ 处的局部环的维数}
见 \ref{stacks-geometry-definition-dimension-local-ring}

\noindent
{\it $X$ 在 $x$ 处之局部环的维数}
见 \ref{spaces-properties-definition-dimension-local-ring}

\noindent
{\it $f$ 在 $x$ 处之纤维的局部环维数}
见 \ref{spaces-morphisms-definition-dimension-fibre}

\noindent
{\it 维数}
见 \ref{topology-definition-Krull}

\noindent
{\it 维数}
见 \ref{properties-definition-dimension}

\noindent
{\it 维数}
见 \ref{spaces-properties-definition-dimension}

\noindent
{\it 维数}
见 \ref{stacks-properties-definition-dimension}

\noindent
{\it 直像函子}
见 \ref{sites-definition-localize}

\noindent
{\it 直像函子}
见 \ref{sites-modules-definition-localize-ringed-site}

\noindent
{\it 紧支撑直像}
见 \ref{more-etale-definition-f-shriek-separated}

\noindent
{\it 紧支撑直像}
见 \ref{more-etale-definition-f-shriek-lqf}

\noindent
{\it 直接像}
见 \ref{etale-cohomology-definition-direct-image-presheaf}

\noindent
{\it 直接像}
见 \ref{etale-cohomology-definition-direct-image-sheaf}

\noindent
{\it 直和 dévissage}
见 \ref{algebra-definition-devissage}

\noindent
{\it 直和}
见 \ref{homology-definition-direct-sum}

\noindent
{\it 有向逆系}
见 \ref{categories-definition-directed-system}

\noindent
{\it 有向偏序集}
见 \ref{categories-definition-directed-set}

\noindent
{\it 有向集}
见 \ref{categories-definition-directed-set}

\noindent
{\it 有向集}
见 \ref{exercises-definition-directed-poset}

\noindent
{\it 有向系统}
见 \ref{categories-definition-directed-system}

\noindent
{\it 有向系统}
见 \ref{algebra-definition-directed-system}

\noindent
{\it 有向}
见 \ref{categories-definition-directed}

\noindent
{\it 有向}
见 \ref{categories-definition-directed}

\noindent
{\it 离散 $G$-模}
见 \ref{etale-cohomology-definition-G-module-continuous}

\noindent
{\it 离散 $G$-集}
见 \ref{pione-definition-G-set-continuous}

\noindent
{\it 离散赋值环}
见 \ref{algebra-definition-value-group}

\noindent
{\it 离散的}
见 \ref{categories-definition-discrete}

\noindent
{\it $L/K$ 的判别式}
见 \ref{fields-definition-discriminant}

\noindent
{\it $M$ 与 $M'$ 之间的距离}
见 \ref{algebra-definition-distance}

\noindent
{\it $K(\mathcal{A})$ 的可区别三角}
见 \ref{derived-definition-distinguished-triangle}

\noindent
{\it 可区别三角}
见 \ref{derived-definition-triangulated-category}

\noindent
{\it 可区别三角}
见 \ref{dga-definition-distinguished-triangle}

\noindent
{\it 除幂 $A$-导子}
见 \ref{crystalline-definition-derivation}

\noindent
{\it $B$ 中 $J$ 相对于 $(A, I, \gamma)$ 的除幂包络}
见 \ref{crystalline-definition-divided-power-envelope}

\noindent
{\it 除幂环}
见 \ref{dpa-definition-divided-power-ring}

\noindent
{\it 除幂概形}
见 \ref{crystalline-definition-divided-power-scheme}

\noindent
{\it 除幂结构 $\gamma$}
见 \ref{crystalline-definition-divided-power-structure}

\noindent
{\it 除幂结构}
见 \ref{dpa-definition-divided-powers}

\noindent
{\it 除幂结构}
见 \ref{dpa-definition-divided-powers-graded}

\noindent
{\it 相对于 $(S, \mathcal{I}, \gamma)$ 的 $X$ 的除幂加厚}
见 \ref{crystalline-definition-divided-power-thickening-X}

\noindent
{\it 除幂加厚}
见 \ref{crystalline-definition-affine-thickening}

\noindent
{\it 除幂加厚}
见 \ref{crystalline-definition-divided-power-thickening}

\noindent
{\it 在 $S$ 上为 DM}
见 \ref{stacks-morphisms-definition-absolute-separated}

\noindent
{\it DM}
见 \ref{stacks-morphisms-definition-separated}

\noindent
{\it DM}
见 \ref{stacks-morphisms-definition-absolute-separated}

\noindent
{\it 定义域}
见 \ref{morphisms-definition-domain-of-definition}

\noindent
{\it 定义域}
见 \ref{spaces-morphisms-definition-domain-of-definition}

\noindent
{\it 整环}
见 \ref{fields-definition-domain}

\noindent
{\it 支配}
见 \ref{morphisms-definition-dominant}

\noindent
{\it 支配}
见 \ref{morphisms-definition-dominant-rational}

\noindent
{\it 支配}
见 \ref{spaces-morphisms-definition-dominant}

\noindent
{\it 支配}
见 \ref{spaces-morphisms-definition-dominant-rational}

\noindent
{\it 支配}
见 \ref{algebra-definition-valuation-ring}

\noindent
{\it 支配}
见 \ref{algebra-definition-domination}

\noindent
{\it dotted arrow}
见 \ref{stacks-morphisms-definition-fill-in-diagram}

\noindent
{\it 双复形}
见 \ref{homology-definition-double-complex}

\noindent
{\it 对偶数}
见 \ref{varieties-definition-dual-numbers}

\noindent
{\it 对偶数}
见 \ref{exercises-definition-dual-numbers}

\noindent
{\it dualizing complex normalized relative to $\omega_S^\bullet$}
见 \ref{duality-definition-good-dualizing}

\noindent
{\it 对偶复形}
见 \ref{dualizing-definition-dualizing}

\noindent
{\it 对偶复形}
见 \ref{duality-definition-dualizing-scheme}

\noindent
{\it 对偶复形}
见 \ref{spaces-duality-definition-dualizing-scheme}

\noindent
{\it 有效 Cartier 除子}
见 \ref{divisors-definition-effective-Cartier-divisor}

\noindent
{\it 有效 Cartier 除子}
见 \ref{spaces-divisors-definition-effective-Cartier-divisor}

\noindent
{\it 有效 Cartier 除子}
见 \ref{exercises-definition-divisor}

\noindent
{\it 模的有效下降态射}
见 \ref{descent-definition-effective-descent}

\noindent
{\it 有效满态射}
见 \ref{sites-definition-universal-effective-epimorphisms}

\noindent
{\it 有效}
见 \ref{stacks-definition-effective-descent-datum}

\noindent
{\it 有效}
见 \ref{descent-definition-descent-datum-effective-quasi-coherent}

\noindent
{\it 有效}
见 \ref{descent-definition-descent-datum-effective-module}

\noindent
{\it 有效}
见 \ref{descent-definition-effective}

\noindent
{\it 有效}
见 \ref{descent-definition-effective-family}

\noindent
{\it 有效}
见 \ref{chow-definition-effective-cycle}

\noindent
{\it 有效}
见 \ref{etale-cohomology-definition-descent-datum}

\noindent
{\it 有效}
见 \ref{etale-cohomology-definition-effective-modules}

\noindent
{\it 有效}
见 \ref{relative-cycles-definition-effective}

\noindent
{\it 有效}
见 \ref{spaces-descent-definition-descent-datum-effective-quasi-coherent}

\noindent
{\it 有效}
见 \ref{spaces-descent-definition-effective}

\noindent
{\it 有效}
见 \ref{spaces-descent-definition-effective-family}

\noindent
{\it 有效}
见 \ref{artin-definition-effective}

\noindent
{\it Eilenberg--Mac Lane 对象 $K(A, k)$}
见 \ref{simplicial-definition-eilenberg-maclane}

\noindent
{\it 环同态 $R \to S$ 在 $\mathfrak q$ 处的初等 \'etale 局部化}
见 \ref{flat-definition-elementary-etale-neighbourhood}

\noindent
{\it 初等平展邻域}
见 \ref{more-morphisms-definition-etale-neighbourhood}

\noindent
{\it 初等平展邻域}
见 \ref{decent-spaces-definition-elemenary-etale-neighbourhood}

\noindent
{\it 初等可区别方块}
见 \ref{spaces-perfect-definition-elementary-distinguished-square}

\noindent
{\it 初等因子整环}
见 \ref{more-algebra-definition-bezout}

\noindent
{\it 在 $A$ 中相对于 $R$ 为初等标准}
见 \ref{smoothing-definition-strictly-standard}

\noindent
{\it 嵌入伴随点}
见 \ref{divisors-definition-embedded}

\noindent
{\it 嵌入伴随素理想}
见 \ref{algebra-definition-embedded-primes}

\noindent
{\it 嵌入分量}
见 \ref{divisors-definition-embedded}

\noindent
{\it 嵌入点}
见 \ref{divisors-definition-embedded}

\noindent
{\it $R$ 的嵌入素理想}
见 \ref{algebra-definition-embedded-primes}

\noindent
{\it $X$ 在 $x$ 处的嵌入维数}
见 \ref{varieties-definition-embed-dim}

\noindent
{\it $X/k$ 在 $x$ 处的嵌入维数}
见 \ref{varieties-definition-embedding-dimension}

\noindent
{\it 嵌入}
见 \ref{sites-definition-embedding}

\noindent
{\it 足够多的 $P$ 对象}
见 \ref{sites-definition-w-contractible}

\noindent
{\it 有足够多的内射对象}
见 \ref{homology-definition-enough-injectives}

\noindent
{\it 有足够多的投射对象}
见 \ref{homology-definition-enough-projectives}

\noindent
{\it 足够多的弱可缩对象}
见 \ref{sites-definition-w-contractible}

\noindent
{\it 包络}
见 \ref{chow-definition-envelope}

\noindent
{\it 满态射}
见 \ref{categories-definition-mono-epi}

\noindent
{\it 等化子}
见 \ref{categories-definition-equalizers}

\noindent
{\it 等维}
见 \ref{topology-definition-equidimensional}

\noindent
{\it 等维}
见 \ref{relative-cycles-definition-equidimensional}

\noindent
{\it 范畴等价}
见 \ref{categories-definition-equivalence-categories}

\noindent
{\it $B$ 上 $U$ 的等价关系}
见 \ref{spaces-groupoids-definition-equivalence-relation}

\noindent
{\it $U$ 在 $S$ 上的一个等价关系}
见 \ref{groupoids-definition-equivalence-relation}

\noindent
{\it 等价型}
见 \ref{models-definition-type-equivalent}

\noindent
{\it 等价}
见 \ref{categories-definition-equivalence}

\noindent
{\it 等价}
见 \ref{derived-definition-yoneda-extension}

\noindent
{\it 等价}
见 \ref{morphisms-definition-rational-map}

\noindent
{\it 等价}
见 \ref{etale-cohomology-definition-brauer-equivalent}

\noindent
{\it 等价}
见 \ref{spaces-morphisms-definition-rational-map}

\noindent
{\it 等变拟凝聚 $\mathcal{O}_X$-模}
见 \ref{groupoids-definition-equivariant-module}

\noindent
{\it 等变拟凝聚 $\mathcal{O}_X$-模}
见 \ref{spaces-groupoids-definition-equivariant-module}

\noindent
{\it 等变}
见 \ref{groupoids-definition-action-group-scheme}

\noindent
{\it 等变}
见 \ref{spaces-groupoids-definition-action-group-space}

\noindent
{\it 本质扩张}
见 \ref{dualizing-definition-essential}

\noindent
{\it 本质满射}
见 \ref{formal-defos-definition-essential-surjection}

\noindent
{\it 本质常值逆系}
见 \ref{categories-definition-essentially-constant-system}

\noindent
{\it 本质常值系}
见 \ref{categories-definition-essentially-constant-system}

\noindent
{\it 值}
见 \ref{categories-definition-essentially-constant-diagram}

\noindent
{\it 值}
见 \ref{categories-definition-essentially-constant-diagram}

\noindent
{\it 本质有限表示的}
见 \ref{algebra-definition-essentially-finite-p-t}

\noindent
{\it 本质有限型的}
见 \ref{algebra-definition-essentially-finite-p-t}

\noindent
{\it 本质满射}
见 \ref{categories-definition-faithful}

\noindent
{\it 本质的}
见 \ref{dualizing-definition-essential}

\noindent
{\it 本质的}
见 \ref{dualizing-definition-essential}

\noindent
{\it $\mathcal{F}$ 的 Euler 示性数}
见 \ref{varieties-definition-euler-characteristic}

\noindent
{\it $\mathcal{F}$ 的 Euler 示性数}
见 \ref{spaces-over-fields-definition-euler-characteristic}

\noindent
{\it Euler--Poincar\'e 函数}
见 \ref{exercises-definition-graded-module}

\noindent
{\it 处处有定义}
见 \ref{derived-definition-everywhere-defined}

\noindent
{\it 处处有定义}
见 \ref{derived-definition-everywhere-defined}

\noindent
{\it 在 $x_i$ 处正合}
见 \ref{homology-definition-exact}

\noindent
{\it 在 $y$ 处正合}
见 \ref{homology-definition-exact}

\noindent
{\it 正合复形}
见 \ref{homology-definition-exact}

\noindent
{\it 正合偶}
见 \ref{homology-definition-exact-couple}

\noindent
{\it 正合函子}
见 \ref{derived-definition-exact-functor-triangulated-categories}

\noindent
{\it 分次模的正合列}
见 \ref{exercises-definition-graded-algebra}

\noindent
{\it 正合列}
见 \ref{homology-definition-exact}

\noindent
{\it 正合}
见 \ref{categories-definition-exact}

\noindent
{\it 正合}
见 \ref{homology-definition-exact}

\noindent
{\it 正合}
见 \ref{chow-definition-periodic-complex}

\noindent
{\it 优秀概形}
见 \ref{more-algebra-definition-excellent}

\noindent
{\it 优秀概形}
见 \ref{morphisms-definition-excellent}

\noindent
{\it 例外除子}
见 \ref{divisors-definition-blow-up}

\noindent
{\it 例外除子}
见 \ref{spaces-divisors-definition-blow-up}

\noindent
{\it 穷尽}
见 \ref{homology-definition-filtered}

\noindent
{\it existence part of the valuative criterion}
见 \ref{stacks-morphisms-definition-existence}

\noindent
{\it 延拓}
见 \ref{dpa-definition-extends}

\noindent
{\it 以 $A$ 为核的 $B$ 的扩张 $E$}
见 \ref{homology-definition-extension}

\noindent
{\it $\mathcal{F}$ 的以 $0$ 延拓 $j_!\mathcal{F}$}
见 \ref{sheaves-definition-j-shriek-structures}

\noindent
{\it $\mathcal{F}$ 的以 $e$ 延拓 $j_!\mathcal{F}$}
见 \ref{sheaves-definition-j-shriek-structures}

\noindent
{\it $\mathcal{F}$ 的以 $0$ 延拓 $j_{p!}\mathcal{F}$}
见 \ref{sheaves-definition-j-shriek-structures}

\noindent
{\it $\mathcal{F}$ 的以 $e$ 延拓 $j_{p!}\mathcal{F}$}
见 \ref{sheaves-definition-j-shriek-structures}

\noindent
{\it 以 $0$ 延拓}
见 \ref{sheaves-definition-j-shriek-structures}

\noindent
{\it 以 $0$ 延拓}
见 \ref{sheaves-definition-j-shriek-structures}

\noindent
{\it 零延拓}
见 \ref{sites-modules-definition-localize-ringed-site}

\noindent
{\it 零延拓}
见 \ref{etale-cohomology-definition-extension-zero}

\noindent
{\it 零延拓}
见 \ref{etale-cohomology-definition-extension-zero}

\noindent
{\it 零延拓}
见 \ref{proetale-definition-extension-zero}

\noindent
{\it 零延拓}
见 \ref{proetale-definition-extension-zero}

\noindent
{\it $\mathcal{F}$ 的以空集延拓 $j_!\mathcal{F}$}
见 \ref{sheaves-definition-j-shriek}

\noindent
{\it $\mathcal{F}$ 的以空集延拓 $j_{p!}\mathcal{F}$}
见 \ref{sheaves-definition-j-shriek}

\noindent
{\it $\mathcal{G}$ 以空集作的延拓}
见 \ref{sites-definition-localize}

\noindent
{\it 离散赋值环的扩张}
见 \ref{more-algebra-definition-extension-discrete-valuation-rings}

\noindent
{\it 赋值环的扩张}
见 \ref{more-algebra-definition-extension-valuation-rings}

\noindent
{\it 极端不连通的}
见 \ref{topology-definition-extremally-disconnected}

\noindent
{\it $x$ 的面}
见 \ref{simplicial-definition-terminology-simplicial-sets}

\noindent
{\it 忠实平坦}
见 \ref{algebra-definition-flat}

\noindent
{\it 忠实平坦}
见 \ref{algebra-definition-flat}

\noindent
{\it 忠实平坦}
见 \ref{etale-definition-flat-rings}

\noindent
{\it 忠实平坦}
见 \ref{etale-definition-flat-schemes}

\noindent
{\it 忠实}
见 \ref{categories-definition-faithful}

\noindent
{\it 具有固定靶的态射族}
见 \ref{sites-definition-family-morphisms-fixed-target}

\noindent
{\it 具有固定靶的态射族}
见 \ref{etale-cohomology-definition-family-morphisms-fixed-target}

\noindent
{\it 纤维范畴}
见 \ref{categories-definition-fibre-category}

\noindent
{\it $f$ 在 $s$ 上方的纤维}
见 \ref{schemes-definition-fibre}

\noindent
{\it $V$ 与 $W$ 在 $U$ 上的纤维积}
见 \ref{simplicial-definition-fibre-product}

\noindent
{\it $V$ 与 $W$ 在 $U$ 上的纤维积}
见 \ref{simplicial-definition-fibre-product-cosimplicial-objects}

\noindent
{\it 纤维积}
见 \ref{categories-definition-fibre-products}

\noindent
{\it 纤维积}
见 \ref{schemes-definition-fibre-product}

\noindent
{\it $\mathcal{C}$ 上的纤维化范畴}
见 \ref{categories-definition-fibred-category}

\noindent
{\it 以群胚为纤维}
见 \ref{categories-definition-fibred-groupoids}

\noindent
{\it fibres of $f$ are universally bounded}
见 \ref{morphisms-definition-universally-bounded}

\noindent
{\it fibres of $f$ are universally bounded}
见 \ref{decent-spaces-definition-universally-bounded}

\noindent
{\it 域扩张}
见 \ref{fields-definition-extension}

\noindent
{\it 有理函数域}
见 \ref{morphisms-definition-function-field}

\noindent
{\it 有理函数域}
见 \ref{spaces-over-fields-definition-function-field}

\noindent
{\it 域}
见 \ref{fields-definition-field}

\noindent
{\it 滤过无环}
见 \ref{derived-definition-filtered-acyclic}

\noindent
{\it 滤过无环}
见 \ref{exercises-definition-filtered-acyclic}

\noindent
{\it $\mathcal{A}$ 的滤过复形 $K^\bullet$}
见 \ref{homology-definition-filtered-complex}

\noindent
{\it $\mathcal{A}$ 的滤过导出范畴}
见 \ref{derived-definition-filtered-derived}

\noindent
{\it 滤过导出函子}
见 \ref{trace-definition-filtered-derived-functors}

\noindent
{\it 滤过微分对象}
见 \ref{homology-definition-filtered-differential}

\noindent
{\it 滤过内射}
见 \ref{derived-definition-filtered-complexes-notation}

\noindent
{\it 滤过内射}
见 \ref{trace-definition-filtered}

\noindent
{\it 滤过内射}
见 \ref{exercises-definition-injective-filtered}

\noindent
{\it $\mathcal{A}$ 的滤过对象}
见 \ref{homology-definition-filtered}

\noindent
{\it 滤过拟同构}
见 \ref{derived-definition-filtered-acyclic}

\noindent
{\it 滤过拟同构}
见 \ref{trace-definition-filtered}

\noindent
{\it 滤过拟同构}
见 \ref{exercises-definition-filtered-quasi-isomorphism}

\noindent
{\it 滤过}
见 \ref{categories-definition-directed}

\noindent
{\it 滤过}
见 \ref{categories-definition-directed}

\noindent
{\it 终对象}
见 \ref{categories-definition-final-object-2-category}

\noindent
{\it 终对象}
见 \ref{categories-definition-initial-final}

\noindent
{\it 更细}
见 \ref{sites-definition-finer}

\noindent
{\it 有限 $\text{Tor}$-维数}
见 \ref{trace-definition-finite-tor-dimension}

\noindent
{\it 有限 $R$-模}
见 \ref{algebra-definition-module-finite-type}

\noindent
{\it 有限自由}
见 \ref{sites-modules-definition-global}

\noindent
{\it 有限整体维数}
见 \ref{algebra-definition-finite-gl-dim}

\noindent
{\it 有限内射维数}
见 \ref{more-algebra-definition-injective-dimension}

\noindent
{\it 有限局部常值}
见 \ref{sites-modules-definition-locally-constant}

\noindent
{\it 有限局部常值}
见 \ref{etale-cohomology-definition-finite-locally-constant}

\noindent
{\it 有限局部常值}
见 \ref{etale-cohomology-definition-finite-locally-constant}

\noindent
{\it 秩为 $r$ 的有限局部自由}
见 \ref{algebra-definition-locally-free}

\noindent
{\it 秩为 $r$ 的有限局部自由}
见 \ref{modules-definition-locally-free}

\noindent
{\it 有限局部自由}
见 \ref{algebra-definition-locally-free}

\noindent
{\it 有限局部自由}
见 \ref{modules-definition-locally-free}

\noindent
{\it 有限局部自由}
见 \ref{sites-modules-definition-site-local}

\noindent
{\it 有限局部自由}
见 \ref{morphisms-definition-finite-locally-free}

\noindent
{\it 有限局部自由}
见 \ref{spaces-morphisms-definition-finite-locally-free}

\noindent
{\it 有限局部自由}
见 \ref{exercises-definition-finite-locally-free}

\noindent
{\it 有限阶加厚}
见 \ref{more-morphisms-definition-thickening}

\noindent
{\it 有限阶加厚}
见 \ref{spaces-more-morphisms-definition-thickening}

\noindent
{\it 在 $x \in X$ 处有限表示}
见 \ref{morphisms-definition-finite-presentation}

\noindent
{\it 在 $x$ 处有限表示}
见 \ref{spaces-morphisms-definition-locally-finite-presentation}

\noindent
{\it 有限表示}
见 \ref{algebra-definition-finite-type}

\noindent
{\it 有限表示}
见 \ref{modules-definition-finite-presentation}

\noindent
{\it 有限表示}
见 \ref{morphisms-definition-finite-presentation}

\noindent
{\it 有限表示}
见 \ref{exercises-definition-finite-presentation}

\noindent
{\it 有限投射维数}
见 \ref{algebra-definition-finite-proj-dim}

\noindent
{\it 有限投射维数}
见 \ref{more-algebra-definition-projective-dimension}

\noindent
{\it 有限 Tor 维数}
见 \ref{more-algebra-definition-tor-amplitude}

\noindent
{\it 有限 Tor 维数}
见 \ref{more-algebra-definition-tor-amplitude}

\noindent
{\it 有限 Tor 维数}
见 \ref{cohomology-definition-tor-amplitude}

\noindent
{\it 有限 Tor 维数}
见 \ref{sites-cohomology-definition-tor-amplitude}

\noindent
{\it 在 $x \in X$ 处为有限型}
见 \ref{morphisms-definition-finite-type}

\noindent
{\it 在 $x$ 处有限型}
见 \ref{spaces-morphisms-definition-locally-finite-type}

\noindent
{\it 有限型点}
见 \ref{morphisms-definition-finite-type-point}

\noindent
{\it 有限型点}
见 \ref{spaces-morphisms-definition-finite-type-point}

\noindent
{\it 有限型点}
见 \ref{stacks-morphisms-definition-finite-type-point}

\noindent
{\it 有限型}
见 \ref{algebra-definition-finite-type}

\noindent
{\it 有限型}
见 \ref{modules-definition-finite-type}

\noindent
{\it 有限型}
见 \ref{morphisms-definition-finite-type}

\noindent
{\it 有限型}
见 \ref{formal-spaces-definition-finite-type}

\noindent
{\it 有限生成 $R$-模}
见 \ref{algebra-definition-module-finite-type}

\noindent
{\it 有限生成域扩张}
见 \ref{fields-definition-generated-by}

\noindent
{\it 有限表示 $R$-模}
见 \ref{algebra-definition-module-finite-type}

\noindent
{\it 相对于 $R$ 有限表示}
见 \ref{more-algebra-definition-relatively-finitely-presented}

\noindent
{\it 相对于 $S$ 有限表现}
见 \ref{more-morphisms-definition-relatively-finitely-presented-sheaf}

\noindent
{\it 有限}
见 \ref{fields-definition-degree}

\noindent
{\it 有限}
见 \ref{algebra-definition-finite-ring-map}

\noindent
{\it 有限}
见 \ref{brauer-definition-finite}

\noindent
{\it 有限}
见 \ref{homology-definition-filtered}

\noindent
{\it 有限}
见 \ref{morphisms-definition-integral}

\noindent
{\it 有限}
见 \ref{spaces-morphisms-definition-integral}

\noindent
{\it 有限}
见 \ref{stacks-morphisms-definition-integral}

\noindent
{\it 第一陈类}
见 \ref{chow-definition-first-chern-class}

\noindent
{\it 一阶无穷小邻域}
见 \ref{more-morphisms-definition-first-order-infinitesimal-neighbourhood}

\noindent
{\it 一阶无穷小邻域}
见 \ref{spaces-more-morphisms-definition-first-order-infinitesimal-neighbourhood}

\noindent
{\it 一阶加厚}
见 \ref{more-morphisms-definition-thickening}

\noindent
{\it 一阶加厚}
见 \ref{spaces-more-morphisms-definition-thickening}

\noindent
{\it 一阶加厚}
见 \ref{stacks-more-morphisms-definition-first-order-thickening}

\noindent
{\it 松弛}
见 \ref{cohomology-definition-flasque}

\noindent
{\it 松弛}
见 \ref{cohomology-definition-flasque}

\noindent
{\it 平坦（相应地，忠实平坦）}
见 \ref{etale-definition-flat-rings}

\noindent
{\it 在 $x \in X$ 处平坦}
见 \ref{etale-definition-flat-schemes}

\noindent
{\it 在 $x$ 处相对于 $Y$ 平坦}
见 \ref{spaces-morphisms-definition-flat-module}

\noindent
{\it 在 $x$ 处平坦}
见 \ref{modules-definition-flat-at-point}

\noindent
{\it 在 $x$ 处平坦}
见 \ref{modules-definition-flat-morphism}

\noindent
{\it 在 $x$ 处平坦}
见 \ref{spaces-morphisms-definition-flat}

\noindent
{\it 在 $x$ 处平坦}
见 \ref{stacks-morphisms-definition-flat-at-point}

\noindent
{\it 在点 $x \in X$ 处平坦}
见 \ref{morphisms-definition-flat}

\noindent
{\it 平坦基变换性质}
见 \ref{stacks-cohomology-definition-flat-base-change}

\noindent
{\it 平坦基变换}
见 \ref{groupoids-quotients-definition-base-change}

\noindent
{\it 平坦群概形}
见 \ref{groupoids-definition-smooth-group-scheme}

\noindent
{\it $R$ 上平坦的局部完全交}
见 \ref{algebra-definition-lci}

\noindent
{\it 在 $(\Sh(\mathcal{D}), \mathcal{O}')$ 上平坦}
见 \ref{sites-modules-definition-flat-module}

\noindent
{\it 在 $S$ 上于点 $x \in X$ 处平坦}
见 \ref{morphisms-definition-flat}

\noindent
{\it 在 $S$ 上平坦}
见 \ref{morphisms-definition-flat}

\noindent
{\it 在 $x \in X$ 处关于 $Y$ 平坦}
见 \ref{etale-definition-flat-schemes}

\noindent
{\it 在点 $x \in X$ 处相对于 $Y$ 平坦}
见 \ref{modules-definition-flat-module}

\noindent
{\it 在 $Y$ 上平坦}
见 \ref{modules-definition-flat-module}

\noindent
{\it 在 $Y$ 上平坦}
见 \ref{spaces-morphisms-definition-flat-module}

\noindent
{\it $\alpha$ 沿 $f$ 的平坦拉回}
见 \ref{chow-definition-flat-pullback}

\noindent
{\it $\alpha$ 沿 $f$ 的平坦拉回}
见 \ref{spaces-chow-definition-flat-pullback}

\noindent
{\it 平坦-fppf site}
见 \ref{stacks-cohomology-definition-lisse-etale}

\noindent
{\it 扁平化分层}
见 \ref{flat-definition-flattening-stratification}

\noindent
{\it 扁平化分层}
见 \ref{flat-definition-flattening-stratification}

\noindent
{\it 平坦}
见 \ref{algebra-definition-flat}

\noindent
{\it 平坦}
见 \ref{algebra-definition-flat}

\noindent
{\it 平坦}
见 \ref{modules-definition-flat}

\noindent
{\it 平坦}
见 \ref{modules-definition-flat-morphism}

\noindent
{\it 平坦}
见 \ref{sites-modules-definition-flat}

\noindent
{\it 平坦}
见 \ref{sites-modules-definition-flat}

\noindent
{\it 平坦}
见 \ref{sites-modules-definition-flat}

\noindent
{\it 平坦}
见 \ref{sites-modules-definition-flat}

\noindent
{\it 平坦}
见 \ref{sites-modules-definition-flat-morphism}

\noindent
{\it 平坦}
见 \ref{sites-modules-definition-flat-morphism}

\noindent
{\it 平坦}
见 \ref{morphisms-definition-flat}

\noindent
{\it 平坦}
见 \ref{etale-definition-flat-rings}

\noindent
{\it 平坦}
见 \ref{etale-definition-flat-schemes}

\noindent
{\it 平坦}
见 \ref{spaces-morphisms-definition-flat}

\noindent
{\it 平坦}
见 \ref{restricted-definition-flat}

\noindent
{\it 平坦}
见 \ref{stacks-morphisms-definition-flat}

\noindent
{\it 形式代数空间}
见 \ref{formal-spaces-definition-formal-algebraic-space}

\noindent
{\it $\mathcal{X}$ 穿过 $x_0$ 的形式分支}
见 \ref{stacks-geometry-definition-formal-branches}

\noindent
{\it 形式修改}
见 \ref{restricted-definition-formal-modification}

\noindent
{\it $\mathcal{F}$ 的形式对象 $\xi = (R, \xi_n, f_n)$}
见 \ref{formal-defos-definition-formal-objects}

\noindent
{\it 形式对象}
见 \ref{artin-definition-formal-objects}

\noindent
{\it 形式谱}
见 \ref{formal-spaces-definition-affine-formal-spectrum}

\noindent
{\it 在 $R$ 上形式平展}
见 \ref{algebra-definition-formally-etale}

\noindent
{\it 形式平展的}
见 \ref{more-morphisms-definition-formally-etale}

\noindent
{\it 形式平展的}
见 \ref{spaces-more-morphisms-definition-formally-smooth-etale-unramified}

\noindent
{\it 形式平展的}
见 \ref{spaces-more-morphisms-definition-formally-etale}

\noindent
{\it 形式链状}
见 \ref{more-algebra-definition-formally-catenary}

\noindent
{\it 在 $G$ 下形式主齐}
见 \ref{groupoids-definition-pseudo-torsor}

\noindent
{\it 在 $G$ 下形式主齐}
见 \ref{spaces-groupoids-definition-pseudo-torsor}

\noindent
{\it 对 $\mathfrak n$-进拓扑形式光滑}
见 \ref{more-algebra-definition-formally-smooth-adic}

\noindent
{\it 在 $R$ 上形式光滑}
见 \ref{algebra-definition-formally-smooth}

\noindent
{\it 在 $R$ 上形式光滑}
见 \ref{more-algebra-definition-formally-smooth}

\noindent
{\it 形式光滑的}
见 \ref{more-morphisms-definition-formally-smooth}

\noindent
{\it 形式光滑的}
见 \ref{spaces-more-morphisms-definition-formally-smooth-etale-unramified}

\noindent
{\it 形式光滑的}
见 \ref{spaces-more-morphisms-definition-formally-smooth}

\noindent
{\it 形式光滑的}
见 \ref{stacks-more-morphisms-definition-formally-smooth}

\noindent
{\it 在 $R$ 上形式无分歧}
见 \ref{algebra-definition-formally-unramified}

\noindent
{\it 形式非分歧的}
见 \ref{more-morphisms-definition-formally-unramified}

\noindent
{\it 形式非分歧的}
见 \ref{spaces-more-morphisms-definition-formally-smooth-etale-unramified}

\noindent
{\it 形式非分歧的}
见 \ref{spaces-more-morphisms-definition-formally-unramified}

\noindent
{\it Fourier--Mukai 函子}
见 \ref{equiv-definition-fourier-mukai-functor}

\noindent
{\it Fourier--Mukai 核}
见 \ref{equiv-definition-fourier-mukai-functor}

\noindent
{\it $T$ 的 fppf 覆盖}
见 \ref{topologies-definition-fppf-covering}

\noindent
{\it $X$ 的 fppf 覆盖}
见 \ref{spaces-topologies-definition-fppf-covering}

\noindent
{\it fppf 层}
见 \ref{stacks-sheaves-definition-sheaves}

\noindent
{\it $T$ 的 fpqc 覆盖}
见 \ref{topologies-definition-fpqc-covering}

\noindent
{\it $X$ 的 fpqc 覆盖}
见 \ref{spaces-topologies-definition-fpqc-covering}

\noindent
{\it fpqc 覆盖}
见 \ref{etale-cohomology-definition-fpqc-covering}

\noindent
{\it 自由 $\mathcal{O}$-模}
见 \ref{sites-modules-definition-global}

\noindent
{\it $\mathcal{G}$ 上的自由阿贝尔预层}
见 \ref{etale-cohomology-definition-free-abelian-presheaf}

\noindent
{\it 自由阿贝尔预层}
见 \ref{sites-modules-definition-free-abelian-presheaf-on}

\noindent
{\it 自由阿贝尔层}
见 \ref{sites-modules-definition-free-abelian-sheaf-on}

\noindent
{\it 自由模}
见 \ref{more-algebra-definition-simple-functors}

\noindent
{\it 自由的}
见 \ref{groupoids-definition-free-action}

\noindent
{\it 自由的}
见 \ref{spaces-groupoids-definition-free-action}

\noindent
{\it 满子范畴}
见 \ref{categories-definition-subcategory}

\noindent
{\it 全忠实}
见 \ref{categories-definition-faithful}

\noindent
{\it 函数域}
见 \ref{morphisms-definition-function-field}

\noindent
{\it 函数域}
见 \ref{spaces-over-fields-definition-function-field}

\noindent
{\it $R$-线性范畴的函子}
见 \ref{dga-definition-functor-linear-categories}

\noindent
{\it $R$ 上微分分次范畴的函子}
见 \ref{dga-definition-functor-dga-categories}

\noindent
{\it $R$ 上分次范畴的函子}
见 \ref{dga-definition-functor-graded-categories}

\noindent
{\it 函子}
见 \ref{categories-definition-functor-monoidal-categories}

\noindent
{\it 函子}
见 \ref{categories-definition-functor-symmetric-monoidal-categories}

\noindent
{\it 函子性内射嵌入}
见 \ref{homology-definition-functorial-injective-embedding}

\noindent
{\it 函子性投射满射}
见 \ref{homology-definition-functorial-projective-surjections}

\noindent
{\it 函子}
见 \ref{categories-definition-functor}

\noindent
{\it 函子}
见 \ref{categories-definition-functor-into-2-category}

\noindent
{\it 基本群}
见 \ref{pione-definition-fundamental-group}

\noindent
{\it G-环}
见 \ref{more-algebra-definition-G-ring}

\noindent
{\it G-概形}
见 \ref{properties-definition-G-scheme}

\noindent
{\it 在 $\mathfrak q$ 处 G-无分歧}
见 \ref{algebra-definition-unramified}

\noindent
{\it 在 $x \in X$ 处 G-非分歧}
见 \ref{morphisms-definition-unramified}

\noindent
{\it 在 $x$ 处 G-非分歧}
见 \ref{spaces-morphisms-definition-unramified}

\noindent
{\it G-非分歧}
见 \ref{algebra-definition-unramified}

\noindent
{\it G-非分歧}
见 \ref{morphisms-definition-unramified}

\noindent
{\it G-非分歧}
见 \ref{spaces-morphisms-definition-unramified}

\noindent
{\it Galois 范畴}
见 \ref{pione-definition-galois-category}

\noindent
{\it $K$ 以 $M$ 为系数的 Galois 上同调群}
见 \ref{etale-cohomology-definition-galois-cohomology}

\noindent
{\it Galois 上同调群}
见 \ref{etale-cohomology-definition-galois-cohomology}

\noindent
{\it Galois 群}
见 \ref{fields-definition-galois-group}

\noindent
{\it Galois 扩张}
见 \ref{fields-definition-galois}

\noindent
{\it Galois 扩张}
见 \ref{fields-definition-separable-algebraic}

\noindent
{\it 一般化可沿 $f$ 提升}
见 \ref{topology-definition-lift-specializations}

\noindent
{\it 一般化}
见 \ref{topology-definition-specialization}

\noindent
{\it 一般化}
见 \ref{exercises-definition-closed}

\noindent
{\it 一般化型的}
见 \ref{topology-definition-lift-specializations}

\noindent
{\it 由 $r$ 个整体截面生成}
见 \ref{sites-modules-definition-global}

\noindent
{\it 由有限多个整体截面生成}
见 \ref{sites-modules-definition-global}

\noindent
{\it 由整体截面生成}
见 \ref{modules-definition-globally-generated}

\noindent
{\it 由整体截面生成}
见 \ref{sites-modules-definition-global}

\noindent
{\it 生成域扩张}
见 \ref{fields-definition-generated-by}

\noindent
{\it 生成}
见 \ref{modules-definition-globally-generated}

\noindent
{\it 生成元}
见 \ref{derived-definition-generators}

\noindent
{\it 生成元}
见 \ref{injectives-definition-grothendieck-conditions}

\noindent
{\it 泛点}
见 \ref{topology-definition-generic-point}

\noindent
{\it 泛点}
见 \ref{exercises-definition-generic-point}

\noindent
{\it 亏格}
见 \ref{pic-definition-genus}

\noindent
{\it 亏格}
见 \ref{curves-definition-genus}

\noindent
{\it 几何 Frobenius}
见 \ref{trace-definition-geometric-frobenius}

\noindent
{\it 几何 Frobenius}
见 \ref{trace-definition-geometric-frobenius-on-stalk}

\noindent
{\it 几何亏格}
见 \ref{curves-definition-geometric-genus}

\noindent
{\it 位于 $x$ 上方的几何点}
见 \ref{spaces-properties-definition-geometric-point}

\noindent
{\it 几何点}
见 \ref{etale-cohomology-definition-geometric-point}

\noindent
{\it 几何点}
见 \ref{spaces-properties-definition-geometric-point}

\noindent
{\it 几何商}
见 \ref{groupoids-quotients-definition-geometric}

\noindent
{\it 在 $k$ 上几何连通}
见 \ref{algebra-definition-geometrically-connected}

\noindent
{\it 几何连通}
见 \ref{varieties-definition-geometrically-connected}

\noindent
{\it 几何连通}
见 \ref{spaces-over-fields-definition-geometrically-connected}

\noindent
{\it 在 $k$ 上几何整}
见 \ref{algebra-definition-geometrically-integral}

\noindent
{\it 几何整}
见 \ref{varieties-definition-geometrically-integral}

\noindent
{\it 几何整}
见 \ref{spaces-over-fields-definition-geometrically-integral}

\noindent
{\it 在 $k$ 上几何不可约}
见 \ref{algebra-definition-geometrically-irreducible}

\noindent
{\it 几何不可约}
见 \ref{varieties-definition-geometrically-irreducible}

\noindent
{\it 几何不可约}
见 \ref{spaces-over-fields-definition-geometrically-irreducible}

\noindent
{\it 在 $x$ 处几何正规}
见 \ref{varieties-definition-geometrically-normal}

\noindent
{\it 几何正规}
见 \ref{algebra-definition-geometrically-normal}

\noindent
{\it 几何正规}
见 \ref{varieties-definition-geometrically-normal}

\noindent
{\it 在 $x$ 处几何逐点整}
见 \ref{varieties-definition-geometrically-integral}

\noindent
{\it 几何逐点整}
见 \ref{varieties-definition-geometrically-integral}

\noindent
{\it 在 $x$ 处几何既约}
见 \ref{varieties-definition-geometrically-reduced}

\noindent
{\it 在 $x$ 处几何既约}
见 \ref{spaces-over-fields-definition-geometrically-reduced}

\noindent
{\it 在 $k$ 上几何既约}
见 \ref{algebra-definition-geometrically-reduced}

\noindent
{\it 几何既约}
见 \ref{varieties-definition-geometrically-reduced}

\noindent
{\it 几何既约}
见 \ref{spaces-over-fields-definition-geometrically-reduced}

\noindent
{\it 在 $x$ 处几何正则}
见 \ref{varieties-definition-geometrically-regular}

\noindent
{\it 在 $k$ 上几何正则}
见 \ref{varieties-definition-geometrically-regular}

\noindent
{\it 几何正则}
见 \ref{algebra-definition-geometrically-regular}

\noindent
{\it 在 $x$ 处几何单分支}
见 \ref{properties-definition-unibranch}

\noindent
{\it 在 $x$ 处几何单分支}
见 \ref{spaces-properties-definition-unibranch}

\noindent
{\it 在 $x$ 处几何单分支}
见 \ref{stacks-properties-definition-number-of-geometric-branches}

\noindent
{\it 几何单枝}
见 \ref{more-algebra-definition-unibranch}

\noindent
{\it 几何单枝}
见 \ref{properties-definition-unibranch}

\noindent
{\it 几何单枝}
见 \ref{spaces-properties-definition-unibranch}

\noindent
{\it ……上的 gerbe（胚）}
见 \ref{stacks-definition-gerbe-over-stack-in-groupoids}

\noindent
{\it ……上的 gerbe（胚）}
见 \ref{stacks-morphisms-definition-gerbe}

\noindent
{\it gerbe（胚）}
见 \ref{stacks-definition-gerbe}

\noindent
{\it gerbe（胚）}
见 \ref{stacks-morphisms-definition-gerbe}

\noindent
{\it $X$ 在 $x$ 处的芽}
见 \ref{descent-definition-germs}

\noindent
{\it $k$ 上的整体完全交}
见 \ref{algebra-definition-lci-field}

\noindent
{\it 整体维数}
见 \ref{algebra-definition-finite-gl-dim}

\noindent
{\it 全局有限表示}
见 \ref{sites-modules-definition-global}

\noindent
{\it 全局 Lefschetz 数}
见 \ref{trace-definition-global-lefschetz-number}

\noindent
{\it 全局表示}
见 \ref{sites-modules-definition-global}

\noindent
{\it 整体截面}
见 \ref{sites-definition-global-sections}

\noindent
{\it 下降}
见 \ref{algebra-definition-going-up-down}

\noindent
{\it 上升}
见 \ref{algebra-definition-going-up-down}

\noindent
{\it 下降定理}
见 \ref{exercises-definition-GU-GD}

\noindent
{\it 上升定理}
见 \ref{exercises-definition-GU-GD}

\noindent
{\it 良商}
见 \ref{groupoids-quotients-definition-good}

\noindent
{\it 良好约化}
见 \ref{models-definition-good}

\noindent
{\it 良好分层}
见 \ref{topology-definition-good-stratification}

\noindent
{\it Gorenstein at $x$}
见 \ref{duality-definition-gorenstein-morphism}

\noindent
{\it Gorenstein at $x$}
见 \ref{spaces-more-morphisms-definition-gorenstein}

\noindent
{\it Gorenstein morphism}
见 \ref{duality-definition-gorenstein-morphism}

\noindent
{\it Gorenstein morphism}
见 \ref{spaces-more-morphisms-definition-gorenstein}

\noindent
{\it Gorenstein}
见 \ref{dualizing-definition-gorenstein}

\noindent
{\it Gorenstein}
见 \ref{dualizing-definition-gorenstein}

\noindent
{\it Gorenstein}
见 \ref{duality-definition-gorenstein}

\noindent
{\it 分次 $(\mathcal{A}, \mathcal{B})$-双模}
见 \ref{sdga-definition-bimodule}

\noindent
{\it 分次 $(A, B)$-双模}
见 \ref{dga-definition-bimodule}

\noindent
{\it 分次 $\mathcal{A}$-模}
见 \ref{sdga-definition-gm}

\noindent
{\it 分次 $A$-代数}
见 \ref{exercises-definition-graded-algebra}

\noindent
{\it $R$ 上的分次范畴 $\mathcal{A}$}
见 \ref{dga-definition-graded-category}

\noindent
{\it 分次直和}
见 \ref{dga-definition-graded-direct-sum}

\noindent
{\it 分次函子}
见 \ref{dga-definition-functor-graded-categories}

\noindent
{\it 分次理想}
见 \ref{exercises-definition-graded-algebra}

\noindent
{\it 分次内射}
见 \ref{sdga-definition-graded-injective}

\noindent
{\it 分次 $A$-代数 $B$ 上的分次模 $M$}
见 \ref{exercises-definition-graded-algebra}

\noindent
{\it 分次模}
见 \ref{sdga-definition-gm}

\noindent
{\it 分次模}
见 \ref{exercises-definition-graded-module}

\noindent
{\it 分次子模}
见 \ref{exercises-definition-graded-algebra}

\noindent
{\it $\mathbf{Z}$ 上的格拉斯曼概形}
见 \ref{constructions-definition-grassmannian}

\noindent
{\it $R$ 上的格拉斯曼概形}
见 \ref{constructions-definition-grassmannian}

\noindent
{\it $S$ 上的格拉斯曼概形}
见 \ref{constructions-definition-grassmannian}

\noindent
{\it Grothendieck 阿贝尔范畴}
见 \ref{injectives-definition-grothendieck-conditions}

\noindent
{\it $X$ 的 Grothendieck 群}
见 \ref{perfect-definition-K-group}

\noindent
{\it $X$ 上凝聚层的 Grothendieck 群}
见 \ref{perfect-definition-K-group}

\noindent
{\it $B$ 上的群代数空间}
见 \ref{spaces-groupoids-definition-group-space}

\noindent
{\it 群上同调群}
见 \ref{etale-cohomology-definition-galois-cohomology}

\noindent
{\it $x'$ 在 $x$ 上的无穷小自同构群}
见 \ref{formal-defos-definition-relative-infinitesimal-auts}

\noindent
{\it $x_0$ 的无穷小自同构群}
见 \ref{formal-defos-definition-infinitesimal-auts}

\noindent
{\it $S$ 上的群概形}
见 \ref{groupoids-definition-group-scheme}

\noindent
{\it $B$ 上代数空间中的群胚}
见 \ref{spaces-groupoids-definition-groupoid}

\noindent
{\it $\mathcal{C}$ 上函子中的群胚}
见 \ref{formal-defos-definition-groupoid-in-functors}

\noindent
{\it $S$ 上的群胚}
见 \ref{groupoids-definition-groupoid}

\noindent
{\it $S$ 上的群胚概形}
见 \ref{groupoids-definition-groupoid}

\noindent
{\it 群胚}
见 \ref{categories-definition-groupoid}

\noindent
{\it Gysin 同态}
见 \ref{chow-definition-gysin-homomorphism}

\noindent
{\it Gysin 同态}
见 \ref{spaces-chow-definition-gysin-homomorphism}

\noindent
{\it Gysin 映射}
见 \ref{chow-definition-lci-gysin}

\noindent
{\it $T$ 的 h 覆盖}
见 \ref{flat-definition-h-covering}

\noindent
{\it H-射影}
见 \ref{morphisms-definition-projective}

\noindent
{\it H-拟射影}
见 \ref{morphisms-definition-quasi-projective}

\noindent
{\it 具有对象对的余积}
见 \ref{categories-definition-has-coproducts-of-pairs}

\noindent
{\it 有足够多的点}
见 \ref{sites-definition-enough-points}

\noindent
{\it 具有纤维积}
见 \ref{categories-definition-has-fibre-products}

\noindent
{\it 具有对象对的乘积}
见 \ref{categories-definition-has-products-of-pairs}

\noindent
{\it has property $(\beta)$}
见 \ref{decent-spaces-definition-relative-conditions}

\noindent
{\it has property $(\beta)$}
见 \ref{decent-spaces-definition-relative-conditions}

\noindent
{\it 在 $x$ 处具有性质 $\mathcal{P}$}
见 \ref{spaces-properties-definition-property-at-point}

\noindent
{\it 在 $x$ 处具有性质 $\mathcal{P}$}
见 \ref{stacks-properties-definition-property-at-point}

\noindent
{\it 具有性质 $\mathcal{P}$}
见 \ref{spaces-properties-definition-type-property}

\noindent
{\it 具有性质 $\mathcal{P}$}
见 \ref{spaces-morphisms-definition-P}

\noindent
{\it 具有性质 $\mathcal{P}$}
见 \ref{stacks-properties-definition-type-property}

\noindent
{\it 具有性质 $\mathcal{P}$}
见 \ref{stacks-morphisms-definition-P}

\noindent
{\it 具有性质 $\mathcal{P}$}
见 \ref{stacks-morphisms-definition-etale-smooth-P}

\noindent
{\it 在 $x$ 处具有性质 $\mathcal{Q}$}
见 \ref{spaces-morphisms-definition-P-at-point}

\noindent
{\it 豪斯多夫}
见 \ref{exercises-definition-Hausdorff}

\noindent
{\it 高度}
见 \ref{algebra-definition-height}

\noindent
{\it henselian local ring of $X$ at $x$}
见 \ref{decent-spaces-definition-henselian-local-ring}

\noindent
{\it Hensel 对}
见 \ref{more-algebra-definition-henselian-pair}

\noindent
{\it Hensel}
见 \ref{algebra-definition-henselian}

\noindent
{\it Hensel}
见 \ref{etale-cohomology-definition-henselian}

\noindent
{\it $\mathcal{O}_{S, s}$ 的 Hensel 化}
见 \ref{etale-cohomology-definition-etale-local-rings}

\noindent
{\it $S$ 在 $s$ 处的 Hensel 化}
见 \ref{etale-cohomology-definition-etale-local-rings}

\noindent
{\it Hensel 化}
见 \ref{algebra-definition-henselization}

\noindent
{\it 高阶直接像}
见 \ref{etale-cohomology-definition-higher-direct-images}

\noindent
{\it Hilbert 函数}
见 \ref{exercises-definition-graded-module}

\noindent
{\it Hilbert 多项式}
见 \ref{algebra-definition-hilbert-polynomial}

\noindent
{\it Hilbert 多项式}
见 \ref{varieties-definition-hilbert-polynomial}

\noindent
{\it Hilbert 多项式}
见 \ref{exercises-definition-graded-module}

\noindent
{\it 霍奇滤过}
见 \ref{derham-definition-hodge-filtration}

\noindent
{\it 齐次谱 $\text{Proj}(R)$}
见 \ref{exercises-definition-Proj-R}

\noindent
{\it $\mathcal{A}$ 在 $S$ 上的齐次谱}
见 \ref{constructions-definition-relative-proj}

\noindent
{\it $\mathcal{A}$ 在 $X$ 上的齐次谱}
见 \ref{spaces-divisors-definition-relative-proj}

\noindent
{\it 齐次谱}
见 \ref{algebra-definition-proj}

\noindent
{\it 齐次谱}
见 \ref{constructions-definition-structure-sheaf}

\noindent
{\it 齐次}
见 \ref{exercises-definition-homogeneous-ideal}

\noindent
{\it 同调的}
见 \ref{derived-definition-homological}

\noindent
{\it $K$ 的同调}
见 \ref{hypercovering-definition-homology}

\noindent
{\it 同调}
见 \ref{homology-definition-differential-object-homology}

\noindent
{\it 微分分次 $(\mathcal{A}, \mathcal{B})$-双模同态}
见 \ref{sdga-definition-dg-bimodule}

\noindent
{\it 微分分次 $\mathcal{A}$-模同态}
见 \ref{sdga-definition-dgm}

\noindent
{\it 微分分次 $\mathcal{O}$-代数同态}
见 \ref{sdga-definition-dga}

\noindent
{\it 微分分次代数同态}
见 \ref{dga-definition-homomorphism-dga}

\noindent
{\it 微分分次模同态}
见 \ref{dga-definition-dgm}

\noindent
{\it 除幂环同态}
见 \ref{dpa-definition-divided-power-ring}

\noindent
{\it 除幂加厚同态}
见 \ref{crystalline-definition-affine-thickening}

\noindent
{\it 分次 $(\mathcal{A}, \mathcal{B})$-双模同态}
见 \ref{sdga-definition-bimodule}

\noindent
{\it 分次 $\mathcal{A}$-模同态}
见 \ref{sdga-definition-gm}

\noindent
{\it 分次 $\mathcal{O}$-代数同态}
见 \ref{sdga-definition-ga}

\noindent
{\it 系统同态}
见 \ref{algebra-definition-homomorphism-directed-systems}

\noindent
{\it 拓扑群同态}
见 \ref{topology-definition-topological-group}

\noindent
{\it 拓扑模同态}
见 \ref{topology-definition-topological-module}

\noindent
{\it 拓扑模同态}
见 \ref{more-algebra-definition-topological-ring}

\noindent
{\it 拓扑环同态}
见 \ref{topology-definition-topological-ring}

\noindent
{\it 拓扑环同态}
见 \ref{more-algebra-definition-topological-ring}

\noindent
{\it 分次模或分次环的同态}
见 \ref{exercises-definition-graded-algebra}

\noindent
{\it 同伦}
见 \ref{simplicial-definition-homotopy}

\noindent
{\it 同伦}
见 \ref{simplicial-definition-homotopy-cosimplicial}

\noindent
{\it 同伦}
见 \ref{dga-definition-homotopy}

\noindent
{\it 同伦}
见 \ref{sdga-definition-homotopy}

\noindent
{\it $f$ 与 $g$ 之间的同伦}
见 \ref{dga-definition-homotopy}

\noindent
{\it $f$ 与 $g$ 之间的同伦}
见 \ref{sdga-definition-homotopy}

\noindent
{\it $\mathcal{A}$ 的同伦范畴}
见 \ref{dga-definition-homotopy-category-of-dga-category}

\noindent
{\it 同伦范畴}
见 \ref{dga-definition-complexes-notation}

\noindent
{\it 同伦范畴}
见 \ref{sdga-definition-complexes-notation}

\noindent
{\it 同伦余极限}
见 \ref{derived-definition-derived-colimit}

\noindent
{\it 同伦等价}
见 \ref{homology-definition-homotopy-equivalent}

\noindent
{\it 同伦等价}
见 \ref{homology-definition-homotopy-equivalent-cochain}

\noindent
{\it 同伦等价}
见 \ref{simplicial-definition-homotopy-equivalent}

\noindent
{\it 同伦等价}
见 \ref{homology-definition-homotopy-equivalent}

\noindent
{\it 同伦等价}
见 \ref{homology-definition-homotopy-equivalent-cochain}

\noindent
{\it 同伦等价}
见 \ref{simplicial-definition-homotopy-equivalent}

\noindent
{\it 从 $a$ 到 $b$ 的同伦}
见 \ref{simplicial-definition-homotopy}

\noindent
{\it 从 $a$ 到 $b$ 的同伦}
见 \ref{simplicial-definition-homotopy-cosimplicial}

\noindent
{\it 同伦极限}
见 \ref{derived-definition-derived-limit}

\noindent
{\it 水平}
见 \ref{categories-definition-horizontal-composition}

\noindent
{\it 水平}
见 \ref{categories-definition-2-category}

\noindent
{\it $\mathcal{G}$ 的超覆盖}
见 \ref{hypercovering-definition-hypercovering-variant}

\noindent
{\it $X$ 的超覆盖}
见 \ref{hypercovering-definition-hypercovering}

\noindent
{\it 超覆盖}
见 \ref{hypercovering-definition-hypercovering-variant}

\noindent
{\it 定义理想}
见 \ref{more-algebra-definition-topological-ring}

\noindent
{\it $s$ 的分母理想层}
见 \ref{divisors-definition-regular-meromorphic-ideal-denominators}

\noindent
{\it 诱导局部环同构}
见 \ref{proetale-definition-local-isomorphism}

\noindent
{\it $\varphi$ 的像}
见 \ref{sites-definition-image}

\noindent
{\it $f$ 的像}
见 \ref{homology-definition-kernel}

\noindent
{\it 给定 $\delta$-函子下该短正合列的像}
见 \ref{derived-definition-delta-functor}

\noindent
{\it 直接特化}
见 \ref{topology-definition-dimension-function}

\noindent
{\it 浸入}
见 \ref{schemes-definition-immersion}

\noindent
{\it 浸入}
见 \ref{spaces-definition-immersion}

\noindent
{\it 浸入}
见 \ref{stacks-properties-definition-immersion}

\noindent
{\it $\mathcal{F}$ 在 $s$ 上方的不纯度}
见 \ref{flat-definition-impurity}

\noindent
{\it $\mathcal{F}$ 在 $y$ 上方的不纯度}
见 \ref{spaces-flat-definition-impurity}

\noindent
{\it 位于同一同伦类}
见 \ref{simplicial-definition-homotopy}

\noindent
{\it 位于同一同伦类}
见 \ref{simplicial-definition-homotopy-cosimplicial}

\noindent
{\it ind-\'etale}
见 \ref{proetale-definition-ind-etale}

\noindent
{\it ind-拟仿射的}
见 \ref{more-morphisms-definition-ind-quasi-affine}

\noindent
{\it ind-拟仿射的}
见 \ref{more-morphisms-definition-ind-quasi-affine}

\noindent
{\it ind-Zariski}
见 \ref{proetale-definition-ind-zariski}

\noindent
{\it 不可分解的}
见 \ref{dualizing-definition-indecomposable}

\noindent
{\it 诱导滤过}
见 \ref{homology-definition-filtered}

\noindent
{\it 诱导滤过}
见 \ref{homology-definition-filtration-cohomology-filtered-differential}

\noindent
{\it 诱导滤过}
见 \ref{homology-definition-filtration-cohomology-filtered-complex}

\noindent
{\it $\mathcal{C}$ 中以 $I$ 为指标的归纳系统}
见 \ref{categories-definition-system-over-poset}

\noindent
{\it $\mathcal{S}$ 的惯性纤维化范畴 $\mathcal{I}_\mathcal{S}$}
见 \ref{categories-definition-inertia-fibred-category}

\noindent
{\it $\mathfrak m$ 的惯性群}
见 \ref{more-algebra-definition-decomposition-inertia}

\noindent
{\it 始的}
见 \ref{categories-definition-initial-final}

\noindent
{\it 始的}
见 \ref{categories-definition-initial}

\noindent
{\it 内射包}
见 \ref{dualizing-definition-injective-hull}

\noindent
{\it $A$ 的内射消解}
见 \ref{derived-definition-injective-resolution}

\noindent
{\it $K^\bullet$ 的内射消解}
见 \ref{derived-definition-injective-resolution}

\noindent
{\it 内射振幅位于 $[a, b]$}
见 \ref{more-algebra-definition-injective-dimension}

\noindent
{\it 单射的}
见 \ref{sheaves-definition-injective-surjective}

\noindent
{\it 单射的}
见 \ref{sheaves-definition-injective-surjective}

\noindent
{\it 单射的}
见 \ref{sites-definition-presheaves-injective-surjective}

\noindent
{\it 单射的}
见 \ref{sites-definition-sheaves-injective-surjective}

\noindent
{\it 单射的}
见 \ref{homology-definition-injective-surjective}

\noindent
{\it 单射的}
见 \ref{homology-definition-injective}

\noindent
{\it 单射的}
见 \ref{more-algebra-definition-projective}

\noindent
{\it 不可分次数}
见 \ref{fields-definition-insep-degree}

\noindent
{\it $\mathcal{O}_X$ 在 $\mathcal{A}$ 中的整闭包}
见 \ref{morphisms-definition-integral-closure}

\noindent
{\it $\mathcal{O}_X$ 在 $\mathcal{A}$ 中的整闭包}
见 \ref{spaces-morphisms-definition-integral-closure}

\noindent
{\it 整闭包}
见 \ref{algebra-definition-integral-closure}

\noindent
{\it 整环}
见 \ref{fields-definition-domain}

\noindent
{\it 整于 $I$}
见 \ref{algebra-definition-integral-over-ideal}

\noindent
{\it 整于 $R$}
见 \ref{algebra-definition-integral-ring-map}

\noindent
{\it 整闭}
见 \ref{algebra-definition-integral-closure}

\noindent
{\it 整}
见 \ref{algebra-definition-integral-ring-map}

\noindent
{\it 整}
见 \ref{properties-definition-integral}

\noindent
{\it 整}
见 \ref{morphisms-definition-integral}

\noindent
{\it 整}
见 \ref{spaces-morphisms-definition-integral}

\noindent
{\it 整}
见 \ref{spaces-over-fields-definition-integral-algebraic-space}

\noindent
{\it 整}
见 \ref{stacks-morphisms-definition-integral}

\noindent
{\it 整}
见 \ref{stacks-morphisms-definition-integral-algebraic-stack}

\noindent
{\it 整}
见 \ref{exercises-definition-integral}

\noindent
{\it 内部}
见 \ref{topology-definition-nowhere-dense}

\noindent
{\it 正常相交}
见 \ref{intersection-definition-proper-intersection}

\noindent
{\it 正常相交}
见 \ref{intersection-definition-proper-intersection}

\noindent
{\it 交数}
见 \ref{varieties-definition-intersection-number}

\noindent
{\it 交数}
见 \ref{spaces-over-fields-definition-intersection-number}

\noindent
{\it 与 $\mathcal{E}$ 的第 $j$ 陈类相交}
见 \ref{chow-definition-cap-chern-classes}

\noindent
{\it 与 $\mathcal{L}$ 的第一陈类相交}
见 \ref{chow-definition-cap-c1}

\noindent
{\it 与 $\mathcal{L}$ 的第一陈类相交}
见 \ref{spaces-chow-definition-cap-c1}

\noindent
{\it 逆像 $f^{-1}(Z)$，即闭子概形 $Z$ 的逆像}
见 \ref{schemes-definition-inverse-image-closed-subscheme}

\noindent
{\it 闭子空间 $Z$ 的逆像 $f^{-1}(Z)$}
见 \ref{spaces-morphisms-definition-inverse-image-closed-subspace}

\noindent
{\it 逆像}
见 \ref{etale-cohomology-definition-inverse-image}

\noindent
{\it $\mathcal{C}$ 中以 $I$ 为指标的逆系统}
见 \ref{categories-definition-system-over-poset}

\noindent
{\it 可逆 $\mathcal{O}_X$-模}
见 \ref{modules-definition-invertible}

\noindent
{\it 可逆 ${\mathcal O}_X$-模}
见 \ref{exercises-definition-invertible-sheaf}

\noindent
{\it 可逆模 $M$}
见 \ref{exercises-definition-invertible-module}

\noindent
{\it 可逆模}
见 \ref{exercises-definition-finite-locally-free}

\noindent
{\it 与 $D$ 关联的可逆层 $\mathcal{O}_S(D)$}
见 \ref{divisors-definition-invertible-sheaf-effective-Cartier-divisor}

\noindent
{\it 与 $D$ 关联的可逆层 $\mathcal{O}_X(D)$}
见 \ref{spaces-divisors-definition-invertible-sheaf-effective-Cartier-divisor}

\noindent
{\it 可逆的}
见 \ref{categories-definition-invertible}

\noindent
{\it 可逆的}
见 \ref{more-algebra-definition-invertible}

\noindent
{\it 可逆的}
见 \ref{sites-modules-definition-invertible-sheaf}

\noindent
{\it 不可约分支}
见 \ref{topology-definition-irreducible-components}

\noindent
{\it 不可约分支}
见 \ref{exercises-definition-irreducible-component}

\noindent
{\it 不可约}
见 \ref{topology-definition-irreducible-components}

\noindent
{\it 不可约}
见 \ref{algebra-definition-irreducible-prime-element}

\noindent
{\it 不可约}
见 \ref{exercises-definition-irreducible}

\noindent
{\it 不可约}
见 \ref{exercises-definition-irreducible}

\noindent
{\it 孤立点}
见 \ref{topology-definition-isolated-point}

\noindent
{\it 同构}
见 \ref{categories-definition-isomorphism}

\noindent
{\it J-0}
见 \ref{more-algebra-definition-J}

\noindent
{\it J-1}
见 \ref{more-algebra-definition-J}

\noindent
{\it J-2 的}
见 \ref{more-algebra-definition-J}

\noindent
{\it J-2 的}
见 \ref{morphisms-definition-J}

\noindent
{\it Jacobson 环}
见 \ref{algebra-definition-ring-jacobson}

\noindent
{\it Jacobson 空间}
见 \ref{topology-definition-space-jacobson}

\noindent
{\it Jacobson 空间}
见 \ref{properties-definition-jacobson}

\noindent
{\it Japanese 概形}
见 \ref{algebra-definition-N}

\noindent
{\it Japanese 概形}
见 \ref{properties-definition-nagata}

\noindent
{\it K-平坦}
见 \ref{more-algebra-definition-K-flat}

\noindent
{\it K-平坦}
见 \ref{cohomology-definition-K-flat}

\noindent
{\it K-平坦}
见 \ref{sites-cohomology-definition-K-flat}

\noindent
{\it K-内射}
见 \ref{derived-definition-K-injective}

\noindent
{\it K-内射}
见 \ref{sdga-definition-K-injective}

\noindent
{\it K\"ahler 不同理想}
见 \ref{discriminant-definition-kahler-different}

\noindent
{\it Kan 复形}
见 \ref{simplicial-definition-kan}

\noindent
{\it Kan 纤维化}
见 \ref{simplicial-definition-kan}

\noindent
{\it Kaplansky dévissage}
见 \ref{algebra-definition-devissage}

\noindent
{\it Karoubi 完备}
见 \ref{homology-definition-karoubian}

\noindent
{\it $F$ 的核}
见 \ref{derived-definition-kernel-category}

\noindent
{\it $H$ 的核}
见 \ref{derived-definition-kernel-category}

\noindent
{\it 函子 $F$ 的核}
见 \ref{homology-definition-kernel-category}

\noindent
{\it 核}
见 \ref{homology-definition-kernel}

\noindent
{\it Kolmogorov 空间}
见 \ref{topology-definition-generic-point}

\noindent
{\it 在 $x$ 处为 Koszul}
见 \ref{more-morphisms-definition-lci}

\noindent
{\it 在 $x$ 处为 Koszul}
见 \ref{spaces-more-morphisms-definition-lci}

\noindent
{\it $f_1, \ldots, f_n$ 上的 Koszul 复形}
见 \ref{modules-definition-koszul-complex}

\noindent
{\it $f_1, \ldots, f_r$ 上的 Koszul 复形}
见 \ref{more-algebra-definition-koszul-complex}

\noindent
{\it Koszul 复形}
见 \ref{more-algebra-definition-koszul}

\noindent
{\it Koszul 复形}
见 \ref{modules-definition-koszul}

\noindent
{\it Koszul 态射}
见 \ref{more-morphisms-definition-lci}

\noindent
{\it Koszul 态射}
见 \ref{spaces-more-morphisms-definition-lci}

\noindent
{\it Koszul-正则理想}
见 \ref{more-algebra-definition-regular-ideal}

\noindent
{\it Koszul-正则浸入}
见 \ref{divisors-definition-regular-immersion}

\noindent
{\it Koszul-正则浸入}
见 \ref{spaces-more-morphisms-definition-regular-immersion}

\noindent
{\it Koszul-正则}
见 \ref{more-algebra-definition-koszul-regular-sequence}

\noindent
{\it Koszul-正则}
见 \ref{divisors-definition-regular-ideal-sheaf}

\noindent
{\it Koszul}
见 \ref{stacks-morphisms-definition-lci}

\noindent
{\it $X$ 在 $x$ 处的 Krull 维数}
见 \ref{topology-definition-Krull}

\noindent
{\it Krull 维数}
见 \ref{topology-definition-Krull}

\noindent
{\it Krull 维数}
见 \ref{algebra-definition-Krull}

\noindent
{\it $V$ 中的格}
见 \ref{algebra-definition-lattice}

\noindent
{\it 对 $F$ 左无环}
见 \ref{derived-definition-derived-functor}

\noindent
{\it 左伴随}
见 \ref{categories-definition-adjoint}

\noindent
{\it 左可容许}
见 \ref{derived-definition-admissible}

\noindent
{\it 左可导}
见 \ref{derived-definition-everywhere-defined}

\noindent
{\it 左导出函子 $LF$ 在此处有定义}
见 \ref{derived-definition-right-derived-functor-defined}

\noindent
{\it $F$ 的左导出函子}
见 \ref{derived-definition-derived-functor}

\noindent
{\it 左对偶}
见 \ref{categories-definition-dual}

\noindent
{\it 左正合}
见 \ref{categories-definition-exact}

\noindent
{\it 左乘法系}
见 \ref{categories-definition-multiplicative-system}

\noindent
{\it 左正交}
见 \ref{derived-definition-orthogonal}

\noindent
{\it Leibniz 法则}
见 \ref{algebra-definition-derivation}

\noindent
{\it Leibniz 法则}
见 \ref{modules-definition-derivation}

\noindent
{\it Leibniz 法则}
见 \ref{sites-modules-definition-derivation}

\noindent
{\it 长度}
见 \ref{topology-definition-Krull}

\noindent
{\it 长度}
见 \ref{algebra-definition-length}

\noindent
{\it 长度}
见 \ref{algebra-definition-chain-primes}

\noindent
{\it 长度}
见 \ref{exercises-definition-length}

\noindent
{\it 位于}
见 \ref{etale-cohomology-definition-geometric-point}

\noindent
{\it 位于}
见 \ref{artin-definition-formal-objects}

\noindent
{\it $x$ 沿 $f$ 的提升}
见 \ref{formal-defos-definition-lifts}

\noindent
{\it 提升}
见 \ref{categories-definition-fibre-category}

\noindent
{\it 提升}
见 \ref{categories-definition-fibre-category}

\noindent
{\it 保持极限}
见 \ref{spaces-limits-definition-locally-finite-presentation}

\noindent
{\it 保持极限}
见 \ref{spaces-limits-definition-locally-finite-presentation}

\noindent
{\it 保持极限}
见 \ref{artin-definition-limit-preserving}

\noindent
{\it 保持极限}
见 \ref{stacks-limits-definition-limit-preserving}

\noindent
{\it 极限}
见 \ref{categories-definition-limit}

\noindent
{\it 极限}
见 \ref{homology-definition-limit-spectral-sequence}

\noindent
{\it 次数为 $d$、维数为 $r$ 的线性系}
见 \ref{curves-definition-linear-series}

\noindent
{\it 线性适足的}
见 \ref{adequate-definition-adequate-functor}

\noindent
{\it 在 $\Omega$ 中于 $k$ 上线性无交}
见 \ref{fields-definition-linearly-disjoint}

\noindent
{\it 线性拓扑化}
见 \ref{more-algebra-definition-topological-ring}

\noindent
{\it 线性拓扑化}
见 \ref{more-algebra-definition-topological-ring}

\noindent
{\it 光滑-\'etale site}
见 \ref{stacks-cohomology-definition-lisse-etale}

\noindent
{\it lisse}
见 \ref{proetale-definition-adic}

\noindent
{\it lisse}
见 \ref{trace-definition-l-adic-sheaf}

\noindent
{\it 局部完全交态射}
见 \ref{more-morphisms-definition-lci}

\noindent
{\it 局部完全交态射}
见 \ref{spaces-more-morphisms-definition-lci}

\noindent
{\it 局部完全交态射}
见 \ref{stacks-morphisms-definition-lci}

\noindent
{\it $k$ 上的局部完全交}
见 \ref{algebra-definition-lci-field}

\noindent
{\it $k$ 上的局部完全交}
见 \ref{morphisms-definition-syntomic}

\noindent
{\it 局部完全交}
见 \ref{more-algebra-definition-local-complete-intersection}

\noindent
{\it 局部完全交}
见 \ref{dpa-definition-lci}

\noindent
{\it 局部同态}
见 \ref{algebra-definition-local-ring}

\noindent
{\it 在 $\tau$-拓扑中为局部的}
见 \ref{descent-definition-property-local}

\noindent
{\it 局部同构}
见 \ref{proetale-definition-local-isomorphism}

\noindent
{\it 局部 Lefschetz 数}
见 \ref{trace-definition-local-lefschetz-number}

\noindent
{\it 关于 $\tau$-拓扑在基底上为局部的}
见 \ref{descent-definition-property-morphisms-local}

\noindent
{\it 关于 $\tau$-拓扑在基底上为局部的}
见 \ref{spaces-descent-definition-property-morphisms-local}

\noindent
{\it 对 $\tau$-拓扑在源上局部}
见 \ref{descent-definition-property-morphisms-local-source}

\noindent
{\it 对 $\tau$-拓扑在源上局部}
见 \ref{spaces-descent-definition-property-morphisms-local-source}

\noindent
{\it 局部环映射 $\varphi : R \to S$}
见 \ref{algebra-definition-local-ring}

\noindent
{\it $X$ 在 $x$ 处的局部环}
见 \ref{schemes-definition-locally-ringed-space}

\noindent
{\it 纤维在 $\mathfrak q$ 处的局部环}
见 \ref{algebra-definition-fibre}

\noindent
{\it 局部环}
见 \ref{algebra-definition-local-ring}

\noindent
{\it 局部化态射}
见 \ref{sites-definition-localize}

\noindent
{\it 局部化态射}
见 \ref{sites-definition-localize-topos}

\noindent
{\it 局部化态射}
见 \ref{sites-modules-definition-localize-ringed-site}

\noindent
{\it 局部化态射}
见 \ref{sites-modules-definition-localize-ringed-topos}

\noindent
{\it $A$ 关于 $S$ 的局部化}
见 \ref{algebra-definition-localization}

\noindent
{\it 环化位点 $(\mathcal{C}, \mathcal{O})$ 在对象 $U$ 处的局部化}
见 \ref{sites-modules-definition-localize-ringed-site}

\noindent
{\it 环化拓扑斯 $(\Sh(\mathcal{C}), \mathcal{O})$ 在 $\mathcal{F}$ 处的局部化}
见 \ref{sites-modules-definition-localize-ringed-topos}

\noindent
{\it 位点 $\mathcal{C}$ 在对象 $U$ 处的 局部化}
见 \ref{sites-definition-localize}

\noindent
{\it 拓扑斯 $\Sh(\mathcal{C})$ 在 $\mathcal{F}$ 处的局部化}
见 \ref{sites-definition-localize-topos}

\noindent
{\it 局部化}
见 \ref{algebra-definition-localization-module}

\noindent
{\it 局部化第 $p$ 陈类}
见 \ref{chow-definition-localized-chern}

\noindent
{\it 局部化陈特征}
见 \ref{chow-definition-localized-chern}

\noindent
{\it 局部具有 $P$ 性质}
见 \ref{properties-definition-locally-P}

\noindent
{\it 相对于 $K$ 在 $\overline{x}$ 处局部无环}
见 \ref{etale-cohomology-definition-locally-acyclic}

\noindent
{\it 相对于 $K$ 局部无环}
见 \ref{etale-cohomology-definition-locally-acyclic}

\noindent
{\it 局部无环}
见 \ref{etale-cohomology-definition-locally-acyclic}

\noindent
{\it 局部进*}
见 \ref{formal-spaces-definition-types-formal-algebraic-spaces}

\noindent
{\it 局部代数 $k$-概形}
见 \ref{varieties-definition-algebraic-scheme}

\noindent
{\it 局部闭浸入}
见 \ref{schemes-definition-immersion}

\noindent
{\it 局部闭子空间}
见 \ref{spaces-definition-immersion}

\noindent
{\it 局部闭子叠}
见 \ref{stacks-properties-definition-substacks}

\noindent
{\it 局部连通}
见 \ref{topology-definition-locally-connected}

\noindent
{\it 局部常值}
见 \ref{sites-modules-definition-locally-constant}

\noindent
{\it 局部常值}
见 \ref{etale-cohomology-definition-finite-locally-constant}

\noindent
{\it 局部常值}
见 \ref{etale-cohomology-definition-finite-locally-constant}

\noindent
{\it 局部常值}
见 \ref{etale-cohomology-definition-finite-locally-constant}

\noindent
{\it 局部可构造}
见 \ref{topology-definition-constructible}

\noindent
{\it 局部可数指标且经典}
见 \ref{formal-spaces-definition-types-formal-algebraic-spaces}

\noindent
{\it 局部可数指标}
见 \ref{formal-spaces-definition-types-formal-algebraic-spaces}

\noindent
{\it 局部有限}
见 \ref{topology-definition-locally-finite}

\noindent
{\it 局部有限}
见 \ref{cohomology-definition-covering-locally-finite}

\noindent
{\it 局部有限}
见 \ref{exercises-definition-graded-module}

\noindent
{\it 局部自由}
见 \ref{algebra-definition-locally-free}

\noindent
{\it 局部自由}
见 \ref{modules-definition-locally-free}

\noindent
{\it 局部自由}
见 \ref{sites-modules-definition-site-local}

\noindent
{\it 局部由 $r$ 个截面生成}
见 \ref{sites-modules-definition-site-local}

\noindent
{\it 局部由截面生成}
见 \ref{modules-definition-locally-generated}

\noindent
{\it 局部由截面生成}
见 \ref{sites-modules-definition-site-local}

\noindent
{\it 局部具有有限 Tor 维数}
见 \ref{cohomology-definition-tor-amplitude}

\noindent
{\it 局部具有有限 Tor 维数}
见 \ref{sites-cohomology-definition-tor-amplitude}

\noindent
{\it 局部幂零的}
见 \ref{algebra-definition-locally-nilpotent-ideal}

\noindent
{\it 局部 Noether}
见 \ref{topology-definition-noetherian}

\noindent
{\it 局部 Noether}
见 \ref{properties-definition-noetherian}

\noindent
{\it 局部 Noether}
见 \ref{formal-spaces-definition-types-formal-algebraic-spaces}

\noindent
{\it 局部 Noether}
见 \ref{exercises-definition-Noetherian-scheme}

\noindent
{\it 在 $S$ 上局部有限表示}
见 \ref{spaces-limits-definition-locally-finite-presentation}

\noindent
{\it 局部有限表示}
见 \ref{morphisms-definition-finite-presentation}

\noindent
{\it 局部有限表示}
见 \ref{spaces-morphisms-definition-locally-finite-presentation}

\noindent
{\it 局部有限表示}
见 \ref{spaces-limits-definition-locally-finite-presentation}

\noindent
{\it 局部有限表示}
见 \ref{spaces-limits-definition-locally-finite-presentation}

\noindent
{\it 局部有限表示}
见 \ref{stacks-morphisms-definition-locally-finite-presentation}

\noindent
{\it 局部有限型}
见 \ref{morphisms-definition-finite-type}

\noindent
{\it 局部有限型}
见 \ref{spaces-morphisms-definition-locally-finite-type}

\noindent
{\it 局部有限型}
见 \ref{formal-spaces-definition-finite-type}

\noindent
{\it 局部有限型}
见 \ref{stacks-morphisms-definition-locally-finite-type}

\noindent
{\it 局部为 $P$ 型}
见 \ref{morphisms-definition-locally-P}

\noindent
{\it 局部主闭子概形}
见 \ref{divisors-definition-effective-Cartier-divisor}

\noindent
{\it 局部主闭子空间}
见 \ref{spaces-divisors-definition-effective-Cartier-divisor}

\noindent
{\it 局部投射}
见 \ref{properties-definition-locally-projective}

\noindent
{\it 局部投射}
见 \ref{morphisms-definition-projective}

\noindent
{\it 局部投射}
见 \ref{spaces-properties-definition-locally-projective}

\noindent
{\it 局部拟凝聚}
见 \ref{crystalline-definition-modules}

\noindent
{\it 局部拟凝聚}
见 \ref{stacks-sheaves-definition-locally-quasi-coherent}

\noindent
{\it 局部拟紧}
见 \ref{topology-definition-locally-quasi-compact}

\noindent
{\it 局部拟有限}
见 \ref{morphisms-definition-quasi-finite}

\noindent
{\it 局部拟有限}
见 \ref{spaces-morphisms-definition-locally-quasi-finite}

\noindent
{\it 局部拟有限}
见 \ref{stacks-morphisms-definition-locally-quasi-finite}

\noindent
{\it 局部拟射影}
见 \ref{morphisms-definition-quasi-projective}

\noindent
{\it 局部环化位点}
见 \ref{sites-modules-definition-locally-ringed}

\noindent
{\it 局部环层空间 $(X, \mathcal{O}_X)$}
见 \ref{schemes-definition-locally-ringed-space}

\noindent
{\it 局部环化}
见 \ref{sites-modules-definition-locally-ringed-topos}

\noindent
{\it 在 $S$ 上局部分离}
见 \ref{spaces-definition-separated}

\noindent
{\it 局部分离}
见 \ref{spaces-properties-definition-separated}

\noindent
{\it 局部分离}
见 \ref{spaces-properties-definition-separated}

\noindent
{\it 局部分离}
见 \ref{spaces-morphisms-definition-separated}

\noindent
{\it 局部平凡}
见 \ref{groupoids-definition-principal-homogeneous-space}

\noindent
{\it 局部平凡}
见 \ref{spaces-groupoids-definition-principal-homogeneous-space}

\noindent
{\it 局部弱进}
见 \ref{formal-spaces-definition-types-formal-algebraic-spaces}

\noindent
{\it 局部的}
见 \ref{properties-definition-property-local}

\noindent
{\it 局部的}
见 \ref{morphisms-definition-property-local}

\noindent
{\it 极大 Cohen--Macaulay}
见 \ref{algebra-definition-maximal-CM}

\noindent
{\it McQuillan}
见 \ref{formal-spaces-definition-types-affine-formal-algebraic-space}

\noindent
{\it 亚纯函数}
见 \ref{divisors-definition-sheaf-meromorphic-functions}

\noindent
{\it 亚纯函数}
见 \ref{spaces-divisors-definition-sheaf-meromorphic-functions}

\noindent
{\it $\mathcal{F}$ 的亚纯截面}
见 \ref{divisors-definition-meromorphic-section}

\noindent
{\it $\mathcal{F}$ 的亚纯截面}
见 \ref{spaces-divisors-definition-meromorphic-section}

\noindent
{\it 极小模型}
见 \ref{models-definition-minimal-model}

\noindent
{\it 最小多项式}
见 \ref{fields-definition-minimal-polynomial}

\noindent
{\it 极小}
见 \ref{models-definition-type-minimal}

\noindent
{\it 极小}
见 \ref{formal-defos-definition-minimal-versal}

\noindent
{\it 极小}
见 \ref{formal-defos-definition-minimal-groupoid-in-functors}

\noindent
{\it 极小 versal}
见 \ref{formal-defos-definition-minimal-versal}

\noindent
{\it Mittag--Leffler 条件}
见 \ref{homology-definition-Mittag-Leffler}

\noindent
{\it 模的 Mittag--Leffler 定向系统}
见 \ref{algebra-definition-ML-inductive-system}

\noindent
{\it Mittag--Leffler}
见 \ref{algebra-definition-ML-system}

\noindent
{\it Mittag--Leffler}
见 \ref{algebra-definition-mittag-leffler-module}

\noindent
{\it 混合特征}
见 \ref{more-algebra-definition-mixed}

\noindent
{\it ML}
见 \ref{homology-definition-Mittag-Leffler}

\noindent
{\it $X$ 的改造}
见 \ref{morphisms-definition-modification}

\noindent
{\it $X$ 的改造}
见 \ref{spaces-over-fields-definition-modification}

\noindent
{\it 微分模}
见 \ref{algebra-definition-differentials}

\noindent
{\it 微分模}
见 \ref{modules-definition-module-differentials}

\noindent
{\it 微分模}
见 \ref{sites-modules-definition-module-differentials}

\noindent
{\it K\"ahler 微分模}
见 \ref{algebra-definition-differentials}

\noindent
{\it $k$ 阶主部件模}
见 \ref{algebra-definition-module-principal-parts}

\noindent
{\it $k$ 阶主部件模}
见 \ref{modules-definition-module-principal-parts}

\noindent
{\it $k$ 阶主部件模}
见 \ref{sites-modules-definition-module-principal-parts}

\noindent
{\it 模值函子}
见 \ref{adequate-definition-module-valued-functor}

\noindent
{\it 亏格为 $g$ 的光滑真曲线的模叠}
见 \ref{moduli-curves-definition-deligne-mumford-smooth}

\noindent
{\it 光滑真曲线的模叠}
见 \ref{moduli-curves-definition-deligne-mumford-smooth}

\noindent
{\it 亏格为 $g$ 的稳定曲线模叠}
见 \ref{moduli-curves-definition-deligne-mumford}

\noindent
{\it 稳定曲线模叠}
见 \ref{moduli-curves-definition-deligne-mumford}

\noindent
{\it 幺半范畴}
见 \ref{categories-definition-monoidal-category}

\noindent
{\it 单态射}
见 \ref{categories-definition-mono-epi}

\noindent
{\it 单态射}
见 \ref{schemes-definition-monomorphism}

\noindent
{\it 单态射}
见 \ref{spaces-morphisms-definition-monomorphism}

\noindent
{\it 单态射}
见 \ref{formal-spaces-definition-monomorphism}

\noindent
{\it 单态射}
见 \ref{stacks-properties-definition-monomorphism}

\noindent
{\it 滤过对象的态射 $(A, F) \to (B, F)$}
见 \ref{homology-definition-filtered}

\noindent
{\it 下降数据态射 $(N, \varphi) \to (N', \varphi')$}
见 \ref{descent-definition-descent-datum-modules}

\noindent
{\it $\mathcal{C}$ 上函子中的群胚态射 $(U, R, s, t, c) \to (U', R', s', t', c')$}
见 \ref{formal-defos-definition-groupoid-in-functors}

\noindent
{\it 下降数据态射 $\psi : (\mathcal{F}_i, \varphi_{ij}) \to (\mathcal{F}'_i, \varphi'_{ij})$}
见 \ref{descent-definition-descent-datum-quasi-coherent}

\noindent
{\it 下降数据态射 $\psi : (\mathcal{F}_i, \varphi_{ij}) \to (\mathcal{F}'_i, \varphi'_{ij})$}
见 \ref{spaces-descent-definition-descent-datum-quasi-coherent}

\noindent
{\it $B$ 上群代数空间的态射 $\psi : (G, m) \to (G', m')$}
见 \ref{spaces-groupoids-definition-group-space}

\noindent
{\it $S$ 上群概形的态射 $\psi : (G, m) \to (G', m')$}
见 \ref{groupoids-definition-group-scheme}

\noindent
{\it 态射 $\psi : (V_i, \varphi_{ij}) \to (V'_i, \varphi'_{ij})$}
见 \ref{descent-definition-descent-datum-for-family-of-morphisms}

\noindent
{\it 态射 $\psi : (V_i, \varphi_{ij}) \to (V'_i, \varphi'_{ij})$}
见 \ref{spaces-descent-definition-descent-datum-for-family-of-morphisms}

\noindent
{\it 下降数据的态射 $\psi : (X_i, \varphi_{ij}) \to (X'_i, \varphi'_{ij})$}
见 \ref{stacks-definition-descent-data}

\noindent
{\it $\mathcal{B}$ 上 $\mathcal{O}$-模预层的态射 $\varphi : \mathcal{F} \to \mathcal{G}$}
见 \ref{sheaves-definition-sheaf-modules-basis}

\noindent
{\it $\mathcal{O}$-模预层的态射 $\varphi : \mathcal{F} \to \mathcal{G}$}
见 \ref{sheaves-definition-presheaf-modules}

\noindent
{\it $\mathcal{O}$-模预层的态射 $\varphi : \mathcal{F} \to \mathcal{G}$}
见 \ref{sites-modules-definition-presheaf-modules}

\noindent
{\it $\mathcal{B}$ 上集合预层的态射 $\varphi : \mathcal{F} \to \mathcal{G}$}
见 \ref{sheaves-definition-presheaf-basis}

\noindent
{\it $X$ 上集合预层的态射 $\varphi : \mathcal{F} \to \mathcal{G}$}
见 \ref{sheaves-definition-presheaf}

\noindent
{\it 取值于 $\mathcal{C}$ 的预层态射 $\varphi : \mathcal{F} \to \mathcal{G}$}
见 \ref{sheaves-definition-presheaf-values-in-category}

\noindent
{\it $\mathcal{B}$ 上取值于 $\mathcal{C}$ 的预层态射 $\varphi : \mathcal{F} \to \mathcal{G}$}
见 \ref{sheaves-definition-sheaf-structures-basis}

\noindent
{\it 形式对象的态射 $a : \xi \to \eta$}
见 \ref{formal-defos-definition-formal-objects}

\noindent
{\it $S$ 上群胚概形的态射 $f : (U, R, s, t, c) \to (U', R', s', t', c')$}
见 \ref{groupoids-definition-groupoid}

\noindent
{\it $B$ 上代数空间中的群胚态射 $f : (U, R, s, t, c) \to (U', R', s', t', c')$}
见 \ref{spaces-groupoids-definition-groupoid}

\noindent
{\it 相对于 $Y \to X$ 的下降数据态射 $f : (V/Y, \varphi) \to (V'/Y, \varphi')$}
见 \ref{spaces-descent-definition-descent-datum}

\noindent
{\it $S$ 上代数空间的态射 $f : F \to F'$}
见 \ref{spaces-definition-morphism-algebraic-spaces}

\noindent
{\it 态射 $f : p \to p'$}
见 \ref{sites-definition-morphism-points}

\noindent
{\it 态射 $f : X \to Y$，即 $S$ 上概形的态射}
见 \ref{schemes-definition-base-change}

\noindent
{\it 相对于 $X \to S$ 的下降数据态射 $g : (V/X, \varphi) \to (V'/X, \varphi')$}
见 \ref{descent-definition-descent-datum}

\noindent
{\it 从 $\mathcal{U}$ 到 $\mathcal{V}$ 的态射}
见 \ref{sites-definition-morphism-coverings}

\noindent
{\it 从 $F$ 到 $G$ 的 $\delta$-函子态射}
见 \ref{homology-definition-morphism-delta-functors}

\noindent
{\it $\mathcal{G}$-扭子的态射}
见 \ref{cohomology-definition-torsor}

\noindent
{\it $\mathcal{G}$-扭子的态射}
见 \ref{sites-cohomology-definition-torsor}

\noindent
{\it $G$-模的态射}
见 \ref{etale-cohomology-definition-G-module-continuous}

\noindent
{\it $G$-集的态射}
见 \ref{pione-definition-G-set-continuous}

\noindent
{\it $n$-截断单纯对象的态射}
见 \ref{simplicial-definition-truncated-simplicial-object}

\noindent
{\it $R\text{-}G$-模的态射}
见 \ref{etale-cohomology-definition-G-module-continuous}

\noindent
{\it \'etale 邻域态射}
见 \ref{etale-cohomology-definition-geometric-point}

\noindent
{\it \'etale 邻域态射}
见 \ref{spaces-properties-definition-etale-neighbourhood}

\noindent
{\it 态射}
见 \ref{more-morphisms-definition-etale-neighbourhood}

\noindent
{\it $X$ 上阿贝尔预层的态射}
见 \ref{sheaves-definition-abelian-presheaves}

\noindent
{\it 态射}
见 \ref{formal-spaces-definition-affine-formal-algebraic-space}

\noindent
{\it 仿射概形之间的态射}
见 \ref{schemes-definition-affine-scheme}

\noindent
{\it 锥态射}
见 \ref{constructions-definition-abstract-cone}

\noindent
{\it 余单纯对象的态射 $U \to U'$}
见 \ref{simplicial-definition-cosimplicial-object}

\noindent
{\it 微分对象的态射}
见 \ref{homology-definition-differential-object}

\noindent
{\it 除幂概形态射}
见 \ref{crystalline-definition-divided-power-scheme}

\noindent
{\it 相对于 $(S, \mathcal{I}, \gamma)$ 的 $X$ 的除幂加厚态射}
见 \ref{crystalline-definition-divided-power-thickening-X}

\noindent
{\it 虚线箭头的态射}
见 \ref{categories-definition-dotted-arrows}

\noindent
{\it 虚线箭头的态射}
见 \ref{stacks-morphisms-definition-fill-in-diagram}

\noindent
{\it morphism of elementary \'etale neighbourhoods}
见 \ref{decent-spaces-definition-elemenary-etale-neighbourhood}

\noindent
{\it 正合偶的态射}
见 \ref{homology-definition-exact-couple}

\noindent
{\it 扩张的态射}
见 \ref{homology-definition-extension}

\noindent
{\it $\mathcal{C}$ 中从 $\mathcal{U}$ 到 $\mathcal{V}$ 的固定靶映射族态射}
见 \ref{sites-definition-morphism-coverings}

\noindent
{\it 态射}
见 \ref{formal-spaces-definition-formal-algebraic-space}

\noindent
{\it 形式对象的态射}
见 \ref{artin-definition-formal-objects}

\noindent
{\it 函子态射}
见 \ref{categories-definition-transformation-functors}

\noindent
{\it 芽态射}
见 \ref{descent-definition-germs}

\noindent
{\it 在 $(U, R, s, t, c)$ 上笛卡尔的群胚概形态射}
见 \ref{groupoids-definition-cartesian-morphism}

\noindent
{\it 提升的态射}
见 \ref{formal-defos-definition-lifts}

\noindent
{\it 局部环化位点的态射}
见 \ref{sites-modules-definition-morphism-locally-ringed-topoi}

\noindent
{\it 局部环层空间之间的态射}
见 \ref{schemes-definition-locally-ringed-space}

\noindent
{\it 局部环化拓扑斯的态射}
见 \ref{sites-modules-definition-morphism-locally-ringed-topoi}

\noindent
{\it 模值函子之间的态射}
见 \ref{adequate-definition-module-valued-functor}

\noindent
{\it Postnikov 系统的态射}
见 \ref{derived-definition-postnikov-system}

\noindent
{\it 预形变范畴的态射}
见 \ref{formal-defos-definition-predeformation-category}

\noindent
{\it $\mathcal{X}$ 上预层的态射}
见 \ref{stacks-sheaves-definition-presheaves}

\noindent
{\it 伪 $\mathcal{G}$-扭子的态射}
见 \ref{cohomology-definition-torsor}

\noindent
{\it 伪 $\mathcal{G}$-扭子的态射}
见 \ref{sites-cohomology-definition-torsor}

\noindent
{\it 环化位点的态射}
见 \ref{sites-modules-definition-ringed-site}

\noindent
{\it 赋环空间态射}
见 \ref{sheaves-definition-ringed-space}

\noindent
{\it 环化拓扑斯的态射}
见 \ref{sites-modules-definition-ringed-topos}

\noindent
{\it 概形态射}
见 \ref{schemes-definition-scheme}

\noindent
{\it $\mathcal{O}$-模层态射}
见 \ref{sheaves-definition-sheaf-modules}

\noindent
{\it $\mathcal{O}$-模层态射}
见 \ref{sites-modules-definition-sheaf-modules}

\noindent
{\it $\mathcal{B}$ 上集合层的态射}
见 \ref{sheaves-definition-sheaf-basis}

\noindent
{\it 集合层态射}
见 \ref{sheaves-definition-sheaf}

\noindent
{\it 单纯对象的态射 $U \to U'$}
见 \ref{simplicial-definition-simplicial-object}

\noindent
{\it 模导出范畴的单纯系统的态射}
见 \ref{spaces-simplicial-definition-cartesian-derived-modules}

\noindent
{\it 导出范畴的单纯系统的态射}
见 \ref{spaces-simplicial-definition-cartesian-derived}

\noindent
{\it 位点态射}
见 \ref{sites-definition-morphism-sites}

\noindent
{\it 谱序列的态射}
见 \ref{homology-definition-spectral-sequence}

\noindent
{\it 加厚的态射}
见 \ref{more-morphisms-definition-thickening}

\noindent
{\it 加厚的态射}
见 \ref{spaces-more-morphisms-definition-thickening}

\noindent
{\it 加厚的态射}
见 \ref{stacks-more-morphisms-definition-thickening}

\noindent
{\it 托扑斯态射}
见 \ref{sites-definition-topos}

\noindent
{\it 三角的态射}
见 \ref{derived-definition-triangle}

\noindent
{\it $S$ 上向量丛的态射}
见 \ref{constructions-definition-abstract-vector-bundle}

\noindent
{\it 预层态射}
见 \ref{sites-definition-presheaves-sets}

\noindent
{\it $\mathcal{Z}$ 上增厚的态射}
见 \ref{stacks-more-morphisms-definition-thickening}

\noindent
{\it $B$ 上加厚的态射}
见 \ref{spaces-more-morphisms-definition-thickening}

\noindent
{\it $S$ 上加厚的态射}
见 \ref{more-morphisms-definition-thickening}

\noindent
{\it $\mathcal{P}$ 型态射对 $\tau$-覆盖满足下降}
见 \ref{descent-definition-descending-types-morphisms}

\noindent
{\it 态射}
见 \ref{sites-definition-presheaf}

\noindent
{\it 态射}
见 \ref{trace-definition-l-adic-sheaf}

\noindent
{\it 多重交叉奇点}
见 \ref{curves-definition-multicross}

\noindent
{\it $R$ 的乘法子集}
见 \ref{algebra-definition-multiplicative-subset}

\noindent
{\it 乘法系}
见 \ref{categories-definition-multiplicative-system}

\noindent
{\it $M$ 关于定义理想 $I$ 的重数}
见 \ref{intersection-definition-multiplicity}

\noindent
{\it $Z$ 在 $\mathcal{F}$ 中的重数}
见 \ref{spaces-chow-definition-cycle-associated-to-coherent-sheaf}

\noindent
{\it $Z$ 在 $Y$ 中的重数}
见 \ref{spaces-chow-definition-cycle-associated-to-closed-subscheme}

\noindent
{\it $Z'$ 在 $\mathcal{F}$ 中的重数}
见 \ref{chow-definition-cycle-associated-to-coherent-sheaf}

\noindent
{\it $Z'$ 在 $Z$ 中的重数}
见 \ref{chow-definition-cycle-associated-to-closed-subscheme}

\noindent
{\it $\mathcal{X}$ 穿过 $x_0$ 的一个形式分支的重数}
见 \ref{stacks-geometry-definition-multiplicity-formal-branches}

\noindent
{\it 重数}
见 \ref{chow-definition-periodic-length}

\noindent
{\it 重数}
见 \ref{stacks-geometry-definition-multiplicity}

\noindent
{\it N-1}
见 \ref{algebra-definition-N}

\noindent
{\it N-2}
见 \ref{algebra-definition-N}

\noindent
{\it Nagata 环}
见 \ref{algebra-definition-nagata}

\noindent
{\it Nagata 概形}
见 \ref{properties-definition-nagata}

\noindent
{\it $f$ 的朴素余切复形}
见 \ref{more-morphisms-definition-netherlander}

\noindent
{\it $f$ 的朴素余切复形}
见 \ref{spaces-more-morphisms-definition-netherlander}

\noindent
{\it 朴素余切复形}
见 \ref{algebra-definition-naive-cotangent-complex}

\noindent
{\it 朴素余切复形}
见 \ref{modules-definition-naive-cotangent-complex}

\noindent
{\it 朴素余切复形}
见 \ref{modules-definition-cotangent-complex-morphism-ringed-topoi}

\noindent
{\it 朴素余切复形}
见 \ref{sites-modules-definition-naive-cotangent-complex}

\noindent
{\it 朴素余切复形}
见 \ref{sites-modules-definition-cotangent-complex-morphism-ringed-topoi}

\noindent
{\it 朴素障碍理论}
见 \ref{artin-definition-naive-obstruction-theory}

\noindent
{\it 朴素 rig-平坦的}
见 \ref{restricted-definition-naively-rig-flat}

\noindent
{\it 自然变换}
见 \ref{categories-definition-transformation-functors}

\noindent
{\it 幂零的}
见 \ref{algebra-definition-locally-nilpotent-ideal}

\noindent
{\it 结点}
见 \ref{curves-definition-multicross}

\noindent
{\it 结点}
见 \ref{curves-definition-nodal}

\noindent
{\it Noether}
见 \ref{topology-definition-noetherian}

\noindent
{\it Noether}
见 \ref{homology-definition-Noetherian}

\noindent
{\it Noether}
见 \ref{homology-definition-Noetherian}

\noindent
{\it Noether}
见 \ref{properties-definition-noetherian}

\noindent
{\it Noether}
见 \ref{spaces-properties-definition-noetherian}

\noindent
{\it Noether}
见 \ref{formal-spaces-definition-types-affine-formal-algebraic-space}

\noindent
{\it Noether}
见 \ref{stacks-morphisms-definition-noetherian}

\noindent
{\it Noether}
见 \ref{exercises-definition-Noetherian-space}

\noindent
{\it Noether}
见 \ref{exercises-definition-Noetherian-scheme}

\noindent
{\it 非退化}
见 \ref{crystalline-definition-F-crystal}

\noindent
{\it $X$ 的非奇异投射模型}
见 \ref{curves-definition-normal-projective-model}

\noindent
{\it 非奇异概形}
见 \ref{properties-definition-regular}

\noindent
{\it 非平凡解}
见 \ref{etale-cohomology-definition-Cr}

\noindent
{\it 在 $x$ 处正规}
见 \ref{more-morphisms-definition-normal}

\noindent
{\it 法丛}
见 \ref{divisors-definition-normal-cone}

\noindent
{\it 法丛}
见 \ref{spaces-more-morphisms-definition-normal-cone}

\noindent
{\it $E$ 在 $F$ 上的正规闭包}
见 \ref{fields-definition-normal-closure}

\noindent
{\it 法锥 $C_ZX$}
见 \ref{divisors-definition-normal-cone}

\noindent
{\it 法锥 $C_ZX$}
见 \ref{spaces-more-morphisms-definition-normal-cone}

\noindent
{\it 正规交叉除子}
见 \ref{etale-definition-normal-crossings}

\noindent
{\it 正规态射}
见 \ref{more-morphisms-definition-normal}

\noindent
{\it $X$ 在 $Y$ 中的正规化}
见 \ref{morphisms-definition-normalization-X-in-Y}

\noindent
{\it $X$ 在 $Y$ 中的正规化}
见 \ref{spaces-morphisms-definition-normalization-X-in-Y}

\noindent
{\it normalization}
见 \ref{morphisms-definition-normalization}

\noindent
{\it normalization}
见 \ref{spaces-morphisms-definition-normalization}

\noindent
{\it normalization}
见 \ref{stacks-morphisms-definition-normalization}

\noindent
{\it $X$ 在 $x$ 处的正规化吹起}
见 \ref{resolve-definition-normalized-blowup}

\noindent
{\it $X$ 在 $x$ 处的正规化吹起}
见 \ref{spaces-resolve-definition-normalized-blowup}

\noindent
{\it 归一化}
见 \ref{formal-defos-definition-minimal-groupoid-in-functors}

\noindent
{\it 正规}
见 \ref{fields-definition-normal}

\noindent
{\it 正规}
见 \ref{fields-definition-separable-algebraic}

\noindent
{\it 正规}
见 \ref{algebra-definition-domain-normal}

\noindent
{\it 正规}
见 \ref{algebra-definition-ring-normal}

\noindent
{\it 正规}
见 \ref{properties-definition-normal}

\noindent
{\it 范数}
见 \ref{fields-definition-trace-norm}

\noindent
{\it 无处稠密}
见 \ref{topology-definition-nowhere-dense}

\noindent
{\it 数域}
见 \ref{fields-definition-number-field}

\noindent
{\it $A$ 的分支数}
见 \ref{more-algebra-definition-number-of-branches}

\noindent
{\it $X$ 在 $x$ 处的分支数}
见 \ref{properties-definition-number-of-branches}

\noindent
{\it $X$ 在 $x$ 处的分支数}
见 \ref{decent-spaces-definition-number-of-branches}

\noindent
{\it $\mathcal{X}$ 在 $x$ 处的几何分支数}
见 \ref{stacks-properties-definition-number-of-geometric-branches}

\noindent
{\it $A$ 的几何分支数}
见 \ref{more-algebra-definition-number-of-branches}

\noindent
{\it $X$ 在 $x$ 处的几何分支数}
见 \ref{properties-definition-number-of-branches}

\noindent
{\it $X$ 在 $x$ 处的几何分支数}
见 \ref{spaces-properties-definition-number-of-geometric-branches}

\noindent
{\it 数值多项式}
见 \ref{algebra-definition-numerical-polynomial}

\noindent
{\it 数值多项式}
见 \ref{exercises-definition-numerical-polynomial}

\noindent
{\it 与 $X$ 关联的数值型}
见 \ref{models-definition-numerical-type-model}

\noindent
{\it 亏格为 $g$ 的数值型}
见 \ref{models-definition-genus}

\noindent
{\it 数值型}
见 \ref{models-definition-type}

\noindent
{\it 障碍模}
见 \ref{artin-definition-obstruction-theory}

\noindent
{\it 障碍理论}
见 \ref{artin-definition-obstruction-theory}

\noindent
{\it 障碍}
见 \ref{artin-definition-obstruction-theory}

\noindent
{\it 在 $S$ 上相对有限表现}
见 \ref{more-morphisms-definition-relatively-finitely-presented-sheaf}

\noindent
{\it 有限表示的}
见 \ref{sites-modules-definition-site-local}

\noindent
{\it 有限表示的}
见 \ref{spaces-morphisms-definition-locally-finite-presentation}

\noindent
{\it 有限表示的}
见 \ref{stacks-morphisms-definition-locally-finite-presentation}

\noindent
{\it 有限型的}
见 \ref{sites-modules-definition-site-local}

\noindent
{\it 有限型的}
见 \ref{spaces-morphisms-definition-locally-finite-type}

\noindent
{\it 有限型的}
见 \ref{stacks-morphisms-definition-locally-finite-type}

\noindent
{\it Oka 族}
见 \ref{algebra-definition-oka-family}

\noindent
{\it $x$ 处 $\mathcal{F}/X/S$ 的一步 d\'evissage}
见 \ref{flat-definition-one-step-devissage-at-x}

\noindent
{\it $s$ 上 $\mathcal{F}/X/S$ 的一步 d\'evissage}
见 \ref{flat-definition-one-step-devissage}

\noindent
{\it 开浸入}
见 \ref{sites-definition-immersion-topoi}

\noindent
{\it 开浸入}
见 \ref{schemes-definition-immersion-locally-ringed-spaces}

\noindent
{\it 开浸入}
见 \ref{schemes-definition-immersion}

\noindent
{\it 开浸入}
见 \ref{spaces-definition-immersion}

\noindent
{\it 开浸入}
见 \ref{stacks-properties-definition-immersion}

\noindent
{\it 开子群概形}
见 \ref{groupoids-definition-closed-subgroup-scheme}

\noindent
{\it 开子概形}
见 \ref{schemes-definition-immersion}

\noindent
{\it 与 $U$ 关联的 $(X, \mathcal{O})$ 的开子空间}
见 \ref{sheaves-definition-restriction}

\noindent
{\it $X$ 的与 $U$ 相伴的开子空间}
见 \ref{schemes-definition-open-subspace}

\noindent
{\it 开子空间}
见 \ref{spaces-definition-immersion}

\noindent
{\it 开子叠}
见 \ref{stacks-properties-definition-substacks}

\noindent
{\it 开子拓扑斯}
见 \ref{sites-definition-open-subtopos}

\noindent
{\it 通用性的开性}
见 \ref{artin-definition-openness-versality}

\noindent
{\it 通用性的开性}
见 \ref{artin-definition-openness-versality}

\noindent
{\it 开的}
见 \ref{morphisms-definition-open}

\noindent
{\it 开的}
见 \ref{trace-definition-open}

\noindent
{\it 开的}
见 \ref{spaces-morphisms-definition-open}

\noindent
{\it 开的}
见 \ref{stacks-morphisms-definition-open}

\noindent
{\it 反代数}
见 \ref{brauer-definition-opposite}

\noindent
{\it 对偶范畴}
见 \ref{categories-definition-opposite}

\noindent
{\it 反微分分次代数}
见 \ref{dga-definition-opposite-dga}

\noindent
{\it $R$ 的轨道空间}
见 \ref{groupoids-quotients-definition-orbit-space}

\noindent
{\it 轨道}
见 \ref{groupoids-quotients-definition-orbit}

\noindent
{\it 轨道}
见 \ref{groupoids-quotients-definition-geometric-orbits}

\noindent
{\it 沿 $R$ 的消失阶}
见 \ref{algebra-definition-ord}

\noindent
{\it $f$ 沿 $Z$ 的消失阶}
见 \ref{divisors-definition-order-vanishing}

\noindent
{\it $f$ 沿 $Z$ 的消失阶}
见 \ref{spaces-over-fields-definition-order-vanishing}

\noindent
{\it $s$ 沿 $Z$ 的消失阶}
见 \ref{divisors-definition-order-vanishing-meromorphic}

\noindent
{\it $s$ 沿 $Z$ 的消失阶}
见 \ref{spaces-over-fields-definition-order-vanishing-meromorphic}

\noindent
{\it 有序 {\v C}ech 复形}
见 \ref{cohomology-definition-ordered-cech-complex}

\noindent
{\it 普通二重点}
见 \ref{curves-definition-multicross}

\noindent
{\it 普通二重点}
见 \ref{curves-definition-nodal}

\noindent
{\it $K$ 在 $k$ 上的 p-基}
见 \ref{more-algebra-definition-p-basis}

\noindent
{\it 在 $k$ 上 p-无关}
见 \ref{more-algebra-definition-p-basis}

\noindent
{\it 关于 $\tau$-拓扑寄生的}
见 \ref{descent-definition-parasitic}

\noindent
{\it 寄生的}
见 \ref{descent-definition-parasitic}

\noindent
{\it 寄生的}
见 \ref{stacks-cohomology-definition-parasitic}

\noindent
{\it 偏序}
见 \ref{categories-definition-directed-set}

\noindent
{\it 偏序集}
见 \ref{categories-definition-directed-set}

\noindent
{\it 分划}
见 \ref{topology-definition-paritition}

\noindent
{\it 分块}
见 \ref{topology-definition-paritition}

\noindent
{\it 在 $x$ 处完美}
见 \ref{spaces-more-morphisms-definition-perfect}

\noindent
{\it 完美闭包}
见 \ref{algebra-definition-perfection}

\noindent
{\it 相对于 $R$ 完美}
见 \ref{more-algebra-definition-relatively-perfect}

\noindent
{\it 相对于 $S$ 完美}
见 \ref{perfect-definition-relatively-perfect}

\noindent
{\it 相对于 $Y$ 完美}
见 \ref{spaces-more-morphisms-definition-relatively-perfect}

\noindent
{\it 完美环映射}
见 \ref{more-algebra-definition-pseudo-coherent-perfect}

\noindent
{\it 完美的}
见 \ref{algebra-definition-perfect}

\noindent
{\it 完美的}
见 \ref{more-algebra-definition-perfect}

\noindent
{\it 完美的}
见 \ref{more-algebra-definition-perfect}

\noindent
{\it 完美的}
见 \ref{cohomology-definition-perfect}

\noindent
{\it 完美的}
见 \ref{cohomology-definition-perfect}

\noindent
{\it 完美的}
见 \ref{sites-cohomology-definition-perfect}

\noindent
{\it 完美的}
见 \ref{sites-cohomology-definition-perfect}

\noindent
{\it 完美的}
见 \ref{more-morphisms-definition-perfect}

\noindent
{\it 完美的}
见 \ref{trace-definition-perfect}

\noindent
{\it 完美的}
见 \ref{spaces-more-morphisms-definition-perfect}

\noindent
{\it $T$ 的 ph 覆盖}
见 \ref{topologies-definition-ph-covering}

\noindent
{\it $X$ 的 ph 覆盖}
见 \ref{spaces-topologies-definition-ph-covering}

\noindent
{\it Picard 函子}
见 \ref{pic-definition-picard-functor}

\noindent
{\it $A$ 的 Picard 群}
见 \ref{exercises-definition-class-group}

\noindent
{\it $T$ 的 Picard 群}
见 \ref{models-definition-picard-group}

\noindent
{\it $X$ 的 Picard 群}
见 \ref{exercises-definition-picard-group}

\noindent
{\it Picard 群}
见 \ref{modules-definition-pic}

\noindent
{\it Picard 群}
见 \ref{sites-modules-definition-pic}

\noindent
{\it PID}
见 \ref{algebra-definition-PID}

\noindent
{\it 位点 $\mathcal{C}$ 的点 $p$}
见 \ref{sites-definition-point}

\noindent
{\it 点 $p$}
见 \ref{sites-definition-point-topology}

\noindent
{\it 拓扑斯 $\Sh(\mathcal{C})$ 的点}
见 \ref{sites-definition-point-topos}

\noindent
{\it 点}
见 \ref{spaces-properties-definition-points}

\noindent
{\it 点}
见 \ref{stacks-properties-definition-points}

\noindent
{\it pond\'eration}
见 \ref{more-morphisms-definition-weighting}

\noindent
{\it Postnikov 系统}
见 \ref{derived-definition-postnikov-system}

\noindent
{\it 预进}
见 \ref{more-algebra-definition-topological-ring}

\noindent
{\it 预可容许}
见 \ref{more-algebra-definition-topological-ring}

\noindent
{\it 预等价关系}
见 \ref{groupoids-definition-equivalence-relation}

\noindent
{\it 预等价关系}
见 \ref{spaces-groupoids-definition-equivalence-relation}

\noindent
{\it 预关系}
见 \ref{groupoids-definition-equivalence-relation}

\noindent
{\it 预关系}
见 \ref{spaces-groupoids-definition-equivalence-relation}

\noindent
{\it 预三角范畴}
见 \ref{derived-definition-triangulated-category}

\noindent
{\it 预三角子范畴}
见 \ref{derived-definition-triangulated-subcategory}

\noindent
{\it 预加性}
见 \ref{homology-definition-preadditive}

\noindent
{\it 预形变范畴}
见 \ref{formal-defos-definition-predeformation-category}

\noindent
{\it 预序集}
见 \ref{categories-definition-directed-set}

\noindent
{\it 预序}
见 \ref{categories-definition-directed-set}

\noindent
{\it 由 $(U, R, s, t, c)$ 给出的 $\mathcal{F}$ 的表示}
见 \ref{formal-defos-definition-presentation}

\noindent
{\it 表示}
见 \ref{spaces-definition-presentation}

\noindent
{\it 表示}
见 \ref{algebraic-definition-presentation}

\noindent
{\it 任意基变换下保持}
见 \ref{schemes-definition-preserved-by-base-change}

\noindent
{\it 任意基变换下保持}
见 \ref{schemes-definition-preserved-by-base-change}

\noindent
{\it 基变换下保持}
见 \ref{schemes-definition-preserved-by-base-change}

\noindent
{\it 基变换下保持}
见 \ref{schemes-definition-preserved-by-base-change}

\noindent
{\it $\mathcal{B}$ 上的集合预层 $\mathcal{F}$}
见 \ref{sheaves-definition-presheaf-basis}

\noindent
{\it $X$ 上的集合预层 $\mathcal{F}$}
见 \ref{sheaves-definition-presheaf}

\noindent
{\it $X$ 上取值于 $\mathcal{C}$ 的预层 $\mathcal{F}$}
见 \ref{sheaves-definition-presheaf-values-in-category}

\noindent
{\it $\mathcal{B}$ 上取值于 $\mathcal{C}$ 的预层 $\mathcal{F}$}
见 \ref{sheaves-definition-sheaf-structures-basis}

\noindent
{\it $\mathcal{B}$ 上的 $\mathcal{O}$-模预层 $\mathcal{F}$}
见 \ref{sheaves-definition-sheaf-modules-basis}

\noindent
{\it $\mathcal{O}$-模预层}
见 \ref{sheaves-definition-presheaf-modules}

\noindent
{\it $\mathcal{O}$-模预层}
见 \ref{sites-modules-definition-presheaf-modules}

\noindent
{\it $X$ 上的阿贝尔群预层}
见 \ref{sheaves-definition-abelian-presheaves}

\noindent
{\it 从 $x$ 到 $y$ 的同构预层}
见 \ref{stacks-definition-mor-presheaf}

\noindent
{\it $\mathcal{X}$ 上的模预层}
见 \ref{stacks-sheaves-definition-modules}

\noindent
{\it 从 $x$ 到 $y$ 的态射预层}
见 \ref{stacks-definition-mor-presheaf}

\noindent
{\it $\mathcal{C}$ 上的集合预层}
见 \ref{categories-definition-presheaf}

\noindent
{\it 集合预层}
见 \ref{sites-definition-presheaves-sets}

\noindent
{\it 集合预层}
见 \ref{etale-cohomology-definition-presheaf}

\noindent
{\it $\mathcal{X}$ 上的预层}
见 \ref{stacks-sheaves-definition-presheaves}

\noindent
{\it 预层}
见 \ref{categories-definition-presheaf}

\noindent
{\it 预层}
见 \ref{sites-definition-presheaf}

\noindent
{\it 预稳定曲线族}
见 \ref{moduli-curves-definition-prestable}

\noindent
{\it 素除子}
见 \ref{divisors-definition-Weil-divisor}

\noindent
{\it 素除子}
见 \ref{spaces-over-fields-definition-Weil-divisor}

\noindent
{\it 素除子}
见 \ref{exercises-definition-divisor}

\noindent
{\it $F$ 的素子域}
见 \ref{fields-definition-characteristic}

\noindent
{\it 素的}
见 \ref{algebra-definition-irreducible-prime-element}

\noindent
{\it 与 $f$ 关联的主除子}
见 \ref{chow-definition-principal-divisor}

\noindent
{\it 与 $f$ 关联的主除子}
见 \ref{spaces-chow-definition-principal-divisor}

\noindent
{\it $B$ 上的主齐性 $G$-空间}
见 \ref{spaces-groupoids-definition-principal-homogeneous-space}

\noindent
{\it 主齐性空间}
见 \ref{groupoids-definition-principal-homogeneous-space}

\noindent
{\it 主齐性空间}
见 \ref{spaces-groupoids-definition-principal-homogeneous-space}

\noindent
{\it 主理想整环}
见 \ref{algebra-definition-PID}

\noindent
{\it 与 $f$ 关联的主 Weil 除子}
见 \ref{divisors-definition-principal-divisor}

\noindent
{\it 与 $f$ 关联的主 Weil 除子}
见 \ref{spaces-over-fields-definition-principal-divisor}

\noindent
{\it $T$ 的 pro-平展覆盖}
见 \ref{proetale-definition-fpqc-covering}

\noindent
{\it 积 $U \times V$ 存在}
见 \ref{simplicial-definition-product-with-simplicial-set}

\noindent
{\it $U$ 与 $V$ 的积 $U \times V$}
见 \ref{simplicial-definition-product-with-simplicial-set}

\noindent
{\it 积范畴}
见 \ref{categories-definition-product-category}

\noindent
{\it $U$ 与 $V$ 的积}
见 \ref{simplicial-definition-product}

\noindent
{\it $U$ 与 $V$ 的积}
见 \ref{simplicial-definition-product-cosimplicial-objects}

\noindent
{\it 乘积}
见 \ref{categories-definition-products}

\noindent
{\it 乘积}
见 \ref{categories-definition-product}

\noindent
{\it 仿有限群}
见 \ref{topology-definition-profinite-group}

\noindent
{\it 仿有限的}
见 \ref{topology-definition-profinite}

\noindent
{\it 射影 $n$ 维空间（在 $\mathbf{Z}$ 上）}
见 \ref{constructions-definition-projective-space}

\noindent
{\it 射影 $n$ 维空间（在 $R$ 上）}
见 \ref{constructions-definition-projective-space}

\noindent
{\it 射影 $n$ 维空间（在 $S$ 上）}
见 \ref{constructions-definition-projective-space}

\noindent
{\it 与 $\mathcal{E}$ 相伴的射影丛}
见 \ref{constructions-definition-projective-bundle}

\noindent
{\it 射影覆盖}
见 \ref{dualizing-definition-projective-cover}

\noindent
{\it 投射维数}
见 \ref{algebra-definition-finite-proj-dim}

\noindent
{\it 射影包络}
见 \ref{dualizing-definition-projective-cover}

\noindent
{\it $A$ 的投射消解}
见 \ref{derived-definition-projective-resolution}

\noindent
{\it $K^\bullet$ 的投射消解}
见 \ref{derived-definition-projective-resolution}

\noindent
{\it $\mathcal{C}$ 中以 $I$ 为指标的投射系统}
见 \ref{categories-definition-system-over-poset}

\noindent
{\it 投射簇}
见 \ref{varieties-definition-variety-type}

\noindent
{\it 投射振幅位于 $[a, b]$}
见 \ref{more-algebra-definition-projective-dimension}

\noindent
{\it 投射}
见 \ref{algebra-definition-projective}

\noindent
{\it 投射}
见 \ref{homology-definition-projective}

\noindent
{\it 投射}
见 \ref{morphisms-definition-projective}

\noindent
{\it 投射}
见 \ref{trace-definition-filtered}

\noindent
{\it 固有相对循环}
见 \ref{relative-cycles-definition-proper}

\noindent
{\it 固有簇}
见 \ref{varieties-definition-variety-type}

\noindent
{\it 具有性质 $\mathcal{P}$}
见 \ref{spaces-definition-relative-representable-property}

\noindent
{\it 具有性质 $\mathcal{P}$}
见 \ref{bootstrap-definition-property-transformation}

\noindent
{\it 具有性质 $\mathcal{P}$}
见 \ref{algebraic-definition-relative-representable-property}

\noindent
{\it 固有}
见 \ref{topology-definition-proper-map}

\noindent
{\it 固有}
见 \ref{morphisms-definition-proper}

\noindent
{\it 固有}
见 \ref{spaces-morphisms-definition-proper}

\noindent
{\it 固有}
见 \ref{formal-spaces-definition-proper}

\noindent
{\it 固有}
见 \ref{stacks-morphisms-definition-proper}

\noindent
{\it 可前表示}
见 \ref{formal-defos-definition-prorepresentable}

\noindent
{\it 可前表示}
见 \ref{formal-defos-definition-prorepresentable-groupoid-in-functors}

\noindent
{\it 伪 $\mathcal{G}$-扭子}
见 \ref{cohomology-definition-torsor}

\noindent
{\it 伪 $\mathcal{G}$-扭子}
见 \ref{sites-cohomology-definition-torsor}

\noindent
{\it 伪 $G$-扭子}
见 \ref{groupoids-definition-pseudo-torsor}

\noindent
{\it 伪 $G$-扭子}
见 \ref{spaces-groupoids-definition-pseudo-torsor}

\noindent
{\it 伪函子}
见 \ref{categories-definition-functor-into-2-category}

\noindent
{\it 伪扭子}
见 \ref{cohomology-definition-torsor}

\noindent
{\it 伪扭子}
见 \ref{sites-cohomology-definition-torsor}

\noindent
{\it 拟链状的}
见 \ref{stacks-geometry-definition-pseudo-catenary}

\noindent
{\it 在 $x$ 处伪凝聚}
见 \ref{spaces-more-morphisms-definition-pseudo-coherent}

\noindent
{\it 相对于 $R$ 伪凝聚}
见 \ref{more-algebra-definition-relatively-pseudo-coherent}

\noindent
{\it 相对于 $R$ 伪凝聚}
见 \ref{more-algebra-definition-relatively-pseudo-coherent}

\noindent
{\it 相对于 $S$ 是伪凝聚的}
见 \ref{more-morphisms-definition-relative-pseudo-coherence}

\noindent
{\it 相对于 $S$ 是伪凝聚的}
见 \ref{more-morphisms-definition-relative-pseudo-coherence}

\noindent
{\it 相对于 $Y$ 伪凝聚}
见 \ref{spaces-more-morphisms-definition-relative-pseudo-coherence}

\noindent
{\it 相对于 $Y$ 伪凝聚}
见 \ref{spaces-more-morphisms-definition-relative-pseudo-coherence}

\noindent
{\it 伪凝聚环映射}
见 \ref{more-algebra-definition-pseudo-coherent-perfect}

\noindent
{\it 伪凝聚的}
见 \ref{more-algebra-definition-pseudo-coherent}

\noindent
{\it 伪凝聚的}
见 \ref{more-algebra-definition-pseudo-coherent}

\noindent
{\it 伪凝聚的}
见 \ref{cohomology-definition-pseudo-coherent}

\noindent
{\it 伪凝聚的}
见 \ref{cohomology-definition-pseudo-coherent}

\noindent
{\it 伪凝聚的}
见 \ref{sites-cohomology-definition-pseudo-coherent}

\noindent
{\it 伪凝聚的}
见 \ref{sites-cohomology-definition-pseudo-coherent}

\noindent
{\it 伪凝聚的}
见 \ref{more-morphisms-definition-pseudo-coherent}

\noindent
{\it 伪凝聚的}
见 \ref{spaces-more-morphisms-definition-pseudo-coherent}

\noindent
{\it $\mathcal{F}$ 的拉回 $x^{-1}\mathcal{F}$}
见 \ref{stacks-sheaves-definition-pullback}

\noindent
{\it 拉回函子}
见 \ref{categories-definition-pullback-functor-fibred-category}

\noindent
{\it 拉回函子}
见 \ref{stacks-definition-pullback-functor}

\noindent
{\it 拉回函子}
见 \ref{descent-definition-pullback-functor}

\noindent
{\it 拉回函子}
见 \ref{descent-definition-pullback-functor-family}

\noindent
{\it 拉回函子}
见 \ref{spaces-descent-definition-pullback-functor}

\noindent
{\it 拉回函子}
见 \ref{spaces-descent-definition-pullback-functor-family}

\noindent
{\it $\mathcal{S}$ 沿 $f$ 的拉回}
见 \ref{stacks-definition-pullback-stack}

\noindent
{\it $D$ 沿 $f$ 的拉回有定义}
见 \ref{divisors-definition-pullback-effective-Cartier-divisor}

\noindent
{\it $D$ 沿 $f$ 的拉回有定义}
见 \ref{spaces-divisors-definition-pullback-effective-Cartier-divisor}

\noindent
{\it $S$ 沿 $f$ 的拉回}
见 \ref{sites-definition-pullback-sieve}

\noindent
{\it 有效 Cartier 除子的拉回}
见 \ref{divisors-definition-pullback-effective-Cartier-divisor}

\noindent
{\it 有效 Cartier 除子的拉回}
见 \ref{spaces-divisors-definition-pullback-effective-Cartier-divisor}

\noindent
{\it 对 $f$ 而言亚纯函数的拉回有定义}
见 \ref{divisors-definition-pullback-meromorphic-sections}

\noindent
{\it 对 $f$ 而言亚纯函数的拉回有定义}
见 \ref{spaces-divisors-definition-pullback-meromorphic-sections}

\noindent
{\it 拉回}
见 \ref{sheaves-definition-pushforward}

\noindent
{\it 拉回}
见 \ref{sites-modules-definition-pushforward}

\noindent
{\it 拉回}
见 \ref{groupoids-definition-restrict-relation}

\noindent
{\it 拉回}
见 \ref{etale-cohomology-definition-inverse-image}

\noindent
{\it 拉回}
见 \ref{spaces-groupoids-definition-restrict-relation}

\noindent
{\it 在 $y$ 上方纯}
见 \ref{spaces-flat-definition-pure}

\noindent
{\it 在 $y$ 上方纯}
见 \ref{spaces-flat-definition-pure}

\noindent
{\it 沿 $X_s$ 纯}
见 \ref{flat-definition-pure}

\noindent
{\it 沿 $X_s$ 纯}
见 \ref{flat-definition-pure}

\noindent
{\it 纯扩张模}
见 \ref{adequate-definition-pure-ext}

\noindent
{\it 纯内射分解}
见 \ref{adequate-definition-pure-resolution}

\noindent
{\it 纯内射}
见 \ref{adequate-definition-pure}

\noindent
{\it 纯射影分解}
见 \ref{adequate-definition-pure-resolution}

\noindent
{\it 纯射影}
见 \ref{adequate-definition-pure}

\noindent
{\it 相对于 $S$ 纯}
见 \ref{flat-definition-pure}

\noindent
{\it 相对于 $S$ 纯}
见 \ref{flat-definition-pure}

\noindent
{\it 相对于 $Y$ 纯}
见 \ref{spaces-flat-definition-pure}

\noindent
{\it 相对于 $Y$ 纯}
见 \ref{spaces-flat-definition-pure}

\noindent
{\it 纯不可分}
见 \ref{fields-definition-purely-inseparable}

\noindent
{\it 纯不可分}
见 \ref{fields-definition-purely-inseparable}

\noindent
{\it 纯不可分}
见 \ref{fields-definition-purely-inseparable}

\noindent
{\it 纯不可分}
见 \ref{fields-definition-separable-algebraic}

\noindent
{\it 纯超越扩张}
见 \ref{fields-definition-transcendence}

\noindent
{\it 纯}
见 \ref{algebra-definition-pure-ideal}

\noindent
{\it $\mathcal{S}$ 沿 $f$ 的推前}
见 \ref{stacks-definition-pushforward-stack}

\noindent
{\it 推前}
见 \ref{sheaves-definition-pushforward}

\noindent
{\it 推前}
见 \ref{sites-definition-pushforward-algebraic-structures}

\noindent
{\it 推前}
见 \ref{sites-modules-definition-pushforward}

\noindent
{\it 推前}
见 \ref{chow-definition-proper-pushforward}

\noindent
{\it 推前}
见 \ref{etale-cohomology-definition-direct-image-presheaf}

\noindent
{\it 推前}
见 \ref{etale-cohomology-definition-direct-image-sheaf}

\noindent
{\it 推前}
见 \ref{spaces-chow-definition-proper-pushforward}

\noindent
{\it $V$ 与 $W$ 在 $U$ 上的推出}
见 \ref{simplicial-definition-push-out}

\noindent
{\it 推出}
见 \ref{categories-definition-pushouts}

\noindent
{\it qc 覆盖}
见 \ref{sites-cohomology-definition-covering-LC}

\noindent
{\it 拟仿射}
见 \ref{properties-definition-quasi-affine}

\noindent
{\it 拟仿射}
见 \ref{morphisms-definition-quasi-affine}

\noindent
{\it 拟仿射}
见 \ref{spaces-morphisms-definition-quasi-affine}

\noindent
{\it 拟凝聚 $\mathcal{O}_\mathcal{X}$-模}
见 \ref{stacks-sheaves-definition-quasi-coherent}

\noindent
{\it $(U, R, s, t, c)$ 上的拟凝聚模}
见 \ref{groupoids-definition-groupoid-module}

\noindent
{\it $(U, R, s, t, c)$ 上的拟凝聚模}
见 \ref{spaces-groupoids-definition-groupoid-module}

\noindent
{\it $\mathcal{X}$ 上的拟凝聚模}
见 \ref{stacks-sheaves-definition-quasi-coherent}

\noindent
{\it 拟凝聚 $\mathcal{O}_X$-模层}
见 \ref{modules-definition-quasi-coherent}

\noindent
{\it 拟凝聚}
见 \ref{sites-modules-definition-site-local}

\noindent
{\it 拟凝聚}
见 \ref{etale-cohomology-definition-ringed-site}

\noindent
{\it 拟凝聚}
见 \ref{crystalline-definition-modules}

\noindent
{\it 拟凝聚}
见 \ref{spaces-properties-definition-quasi-coherent}

\noindent
{\it 拟凝聚}
见 \ref{exercises-definition-quasi-coherent}

\noindent
{\it 拟紧}
见 \ref{topology-definition-quasi-compact}

\noindent
{\it 拟紧}
见 \ref{topology-definition-quasi-compact}

\noindent
{\it 拟紧}
见 \ref{sites-definition-quasi-compact}

\noindent
{\it 拟紧}
见 \ref{sites-definition-quasi-compact-topos}

\noindent
{\it 拟紧}
见 \ref{sites-definition-quasi-compact-topos}

\noindent
{\it 拟紧}
见 \ref{schemes-definition-quasi-compact}

\noindent
{\it 拟紧}
见 \ref{spaces-properties-definition-quasi-compact}

\noindent
{\it 拟紧}
见 \ref{spaces-morphisms-definition-quasi-compact}

\noindent
{\it 拟紧}
见 \ref{formal-spaces-definition-quasi-compact}

\noindent
{\it 拟紧}
见 \ref{formal-spaces-definition-quasi-compact-morphism}

\noindent
{\it 拟紧}
见 \ref{stacks-properties-definition-quasi-compact}

\noindent
{\it 拟紧}
见 \ref{stacks-morphisms-definition-quasi-compact}

\noindent
{\it 拟紧}
见 \ref{exercises-definition-quasi-compact}

\noindent
{\it 在 $S$ 上拟 DM}
见 \ref{stacks-morphisms-definition-absolute-separated}

\noindent
{\it 拟 DM}
见 \ref{stacks-morphisms-definition-separated}

\noindent
{\it 拟 DM}
见 \ref{stacks-morphisms-definition-absolute-separated}

\noindent
{\it 拟优秀概形}
见 \ref{more-algebra-definition-excellent}

\noindent
{\it 拟优秀概形}
见 \ref{morphisms-definition-excellent}

\noindent
{\it 在 $\mathfrak q$ 处拟有限}
见 \ref{algebra-definition-quasi-finite}

\noindent
{\it 在 $x$ 处准有限}
见 \ref{spaces-morphisms-definition-locally-quasi-finite}

\noindent
{\it 在点 $x \in X$ 处拟有限}
见 \ref{morphisms-definition-quasi-finite}

\noindent
{\it 拟有限}
见 \ref{algebra-definition-quasi-finite}

\noindent
{\it 拟有限}
见 \ref{morphisms-definition-quasi-finite}

\noindent
{\it 拟有限}
见 \ref{spaces-morphisms-definition-locally-quasi-finite}

\noindent
{\it 拟有限}
见 \ref{stacks-morphisms-definition-quasi-finite}

\noindent
{\it 拟逆}
见 \ref{categories-definition-equivalence-categories}

\noindent
{\it 拟同构}
见 \ref{homology-definition-quasi-isomorphism}

\noindent
{\it 拟同构}
见 \ref{homology-definition-quasi-isomorphism-cochain}

\noindent
{\it 拟等平凡}
见 \ref{groupoids-definition-principal-homogeneous-space}

\noindent
{\it 拟等平凡}
见 \ref{spaces-groupoids-definition-principal-homogeneous-space}

\noindent
{\it 拟投射簇}
见 \ref{varieties-definition-variety-type}

\noindent
{\it 拟射影}
见 \ref{morphisms-definition-quasi-projective}

\noindent
{\it 拟固有的}
见 \ref{topology-definition-proper-map}

\noindent
{\it 拟正则理想}
见 \ref{more-algebra-definition-regular-ideal}

\noindent
{\it 拟正则浸入}
见 \ref{divisors-definition-regular-immersion}

\noindent
{\it 拟正则浸入}
见 \ref{spaces-more-morphisms-definition-regular-immersion}

\noindent
{\it 拟正则序列}
见 \ref{algebra-definition-quasi-regular-sequence}

\noindent
{\it 拟正则}
见 \ref{divisors-definition-regular-ideal-sheaf}

\noindent
{\it 在 $S$ 上拟分离}
见 \ref{spaces-definition-separated}

\noindent
{\it 在 $S$ 上拟分离}
见 \ref{stacks-morphisms-definition-absolute-separated}

\noindent
{\it 拟分离}
见 \ref{schemes-definition-separated}

\noindent
{\it 拟分离}
见 \ref{schemes-definition-separated}

\noindent
{\it 拟分离}
见 \ref{spaces-properties-definition-separated}

\noindent
{\it 拟分离}
见 \ref{spaces-properties-definition-separated}

\noindent
{\it 拟分离}
见 \ref{spaces-morphisms-definition-separated}

\noindent
{\it 拟分离}
见 \ref{formal-spaces-definition-separated}

\noindent
{\it 拟分离}
见 \ref{formal-spaces-definition-separated-morphism}

\noindent
{\it 拟分离}
见 \ref{stacks-morphisms-definition-separated}

\noindent
{\it 拟分离}
见 \ref{stacks-morphisms-definition-absolute-separated}

\noindent
{\it 拟清醒}
见 \ref{topology-definition-generic-point}

\noindent
{\it 在 $u$ 上拟分裂}
见 \ref{spaces-more-groupoids-definition-split-at-point}

\noindent
{\it $R$ 在 $u$ 上的拟分裂}
见 \ref{spaces-more-groupoids-definition-split-at-point}

\noindent
{\it 商范畴 $\mathcal{D}/\mathcal{B}$}
见 \ref{derived-definition-quotient-category}

\noindent
{\it 余纤维化为群胚的商范畴 $[U/R] \to \mathcal{C}$}
见 \ref{formal-defos-definition-quotient}

\noindent
{\it 商滤过}
见 \ref{homology-definition-filtered}

\noindent
{\it 商函子}
见 \ref{derived-definition-quotient-category}

\noindent
{\it 商态射 $U \to [U/R]$}
见 \ref{formal-defos-definition-quotient}

\noindent
{\it $U$ 对 $G$ 的商}
见 \ref{spaces-definition-quotient}

\noindent
{\it 商可由代数空间表示}
见 \ref{spaces-groupoids-definition-representable-quotient}

\noindent
{\it 商可由代数空间表示}
见 \ref{spaces-groupoids-definition-representable-quotient}

\noindent
{\it 商层 $U/R$}
见 \ref{groupoids-definition-quotient-sheaf}

\noindent
{\it 商层 $U/R$}
见 \ref{spaces-groupoids-definition-quotient-sheaf}

\noindent
{\it 商叠}
见 \ref{spaces-groupoids-definition-quotient-stack}

\noindent
{\it 商叠}
见 \ref{spaces-groupoids-definition-quotient-stack}

\noindent
{\it 商}
见 \ref{homology-definition-injective-surjective}

\noindent
{\it 根态射}
见 \ref{morphisms-definition-universally-injective}

\noindent
{\it 根态射}
见 \ref{spaces-more-morphisms-definition-radicial}

\noindent
{\it 分歧指数}
见 \ref{more-algebra-definition-extension-discrete-valuation-rings}

\noindent
{\it 秩 $r$}
见 \ref{sites-modules-definition-invertible-sheaf}

\noindent
{\it 秩}
见 \ref{algebra-definition-rank}

\noindent
{\it 秩}
见 \ref{more-algebra-definition-rank}

\noindent
{\it 秩}
见 \ref{morphisms-definition-finite-locally-free}

\noindent
{\it 秩}
见 \ref{divisors-definition-rank}

\noindent
{\it 秩}
见 \ref{spaces-morphisms-definition-finite-locally-free}

\noindent
{\it $X$ 上的有理函数}
见 \ref{morphisms-definition-rational-function}

\noindent
{\it $X$ 上的有理函数}
见 \ref{spaces-morphisms-definition-rational-function}

\noindent
{\it 从 $X$ 到 $Y$ 的有理映射}
见 \ref{morphisms-definition-rational-map}

\noindent
{\it 从 $X$ 到 $Y$ 的有理映射}
见 \ref{spaces-morphisms-definition-rational-map}

\noindent
{\it 有理等价于零}
见 \ref{chow-definition-rational-equivalence}

\noindent
{\it 有理等价于零}
见 \ref{spaces-chow-definition-rational-equivalence}

\noindent
{\it 有理等价}
见 \ref{chow-definition-rational-equivalence}

\noindent
{\it 有理等价}
见 \ref{spaces-chow-definition-rational-equivalence}

\noindent
{\it reasonable}
见 \ref{decent-spaces-definition-very-reasonable}

\noindent
{\it reasonable}
见 \ref{decent-spaces-definition-relative-conditions}

\noindent
{\it 诱导既约代数空间结构}
见 \ref{spaces-properties-definition-reduced-induced-space}

\noindent
{\it 既约诱导代数叠结构}
见 \ref{stacks-properties-definition-reduced-induced-stack}

\noindent
{\it 约化诱导概形结构}
见 \ref{schemes-definition-reduced-induced-scheme}

\noindent
{\it 约化的}
见 \ref{schemes-definition-reduced}

\noindent
{\it $\mathcal{X}$ 的既约化 $\mathcal{X}_{red}$}
见 \ref{stacks-properties-definition-reduced-induced-stack}

\noindent
{\it $X$ 的既约化 $X_{red}$}
见 \ref{schemes-definition-reduced-induced-scheme}

\noindent
{\it $X$ 的既约化 $X_{red}$}
见 \ref{spaces-properties-definition-reduced-induced-space}

\noindent
{\it 对 $A$ 可以归约到有理奇点}
见 \ref{resolve-definition-reduce-to-rational}

\noindent
{\it Rees 代数}
见 \ref{algebra-definition-blow-up}

\noindent
{\it 加细}
见 \ref{sites-definition-morphism-coverings}

\noindent
{\it 加细}
见 \ref{topology-definition-paritition}

\noindent
{\it 反身包}
见 \ref{more-algebra-definition-reflexive-hull}

\noindent
{\it 反身包}
见 \ref{divisors-definition-reflexive}

\noindent
{\it 反身}
见 \ref{more-algebra-definition-reflexive}

\noindent
{\it 反身}
见 \ref{divisors-definition-reflexive}

\noindent
{\it 在 $x$ 处正则}
见 \ref{more-morphisms-definition-regular}

\noindent
{\it 正则理想}
见 \ref{more-algebra-definition-regular-ideal}

\noindent
{\it 正则浸入}
见 \ref{divisors-definition-regular-immersion}

\noindent
{\it 在余维 $\leq k$ 处正则}
见 \ref{algebra-definition-conditions}

\noindent
{\it 在余维 $k$ 内正则}
见 \ref{properties-definition-Rk}

\noindent
{\it 正则局部环}
见 \ref{algebra-definition-regular-local}

\noindent
{\it 正则轨迹}
见 \ref{properties-definition-singular-locus}

\noindent
{\it 正则态射}
见 \ref{more-morphisms-definition-regular}

\noindent
{\it 正则截面}
见 \ref{divisors-definition-regular-section}

\noindent
{\it 正则截面}
见 \ref{spaces-divisors-definition-regular-section}

\noindent
{\it 正则序列}
见 \ref{algebra-definition-regular-sequence}

\noindent
{\it 正则参数系}
见 \ref{algebra-definition-regular-local}

\noindent
{\it 正则概形}
见 \ref{algebra-definition-regular}

\noindent
{\it 正则概形}
见 \ref{homology-definition-bounded-ss}

\noindent
{\it 正则概形}
见 \ref{more-algebra-definition-regular}

\noindent
{\it 正则概形}
见 \ref{properties-definition-regular}

\noindent
{\it 正则概形}
见 \ref{divisors-definition-regular-ideal-sheaf}

\noindent
{\it 正则概形}
见 \ref{divisors-definition-regular-meromorphic-section}

\noindent
{\it 正则概形}
见 \ref{spaces-divisors-definition-regular-meromorphic-section}

\noindent
{\it 关系}
见 \ref{algebra-definition-relation}

\noindent
{\it 关系}
见 \ref{groupoids-definition-equivalence-relation}

\noindent
{\it 关系}
见 \ref{spaces-groupoids-definition-equivalence-relation}

\noindent
{\it 相对 $H_1$-正则浸入}
见 \ref{divisors-definition-relative-H1-regular-immersion}

\noindent
{\it $X/S$ 上的相对 $r$-循环}
见 \ref{relative-cycles-definition-relative-cycles}

\noindent
{\it $\mathcal{F}$ 在 $X$ 中相对于 $S$ 的相对伴随点集}
见 \ref{divisors-definition-relative-assassin}

\noindent
{\it $N$ 在 $S/R$ 上的相对相伴素理想集}
见 \ref{algebra-definition-relative-assassin}

\noindent
{\it 相对余切空间}
见 \ref{formal-defos-definition-tangent-space-ring}

\noindent
{\it 在 $x$ 处相对维数 $\leq d$}
见 \ref{morphisms-definition-relative-dimension-d}

\noindent
{\it 相对维数 $\leq d$}
见 \ref{morphisms-definition-relative-dimension-d}

\noindent
{\it 相对维数 $\leq d$}
见 \ref{spaces-morphisms-definition-relative-dimension}

\noindent
{\it 相对维数 $d$}
见 \ref{morphisms-definition-relative-dimension-d}

\noindent
{\it 相对维数 $d$}
见 \ref{spaces-morphisms-definition-relative-dimension}

\noindent
{\it $S/R$ 在 $\mathfrak q$ 处的相对维数}
见 \ref{algebra-definition-relative-dimension}

\noindent
{\it ……的相对维数}
见 \ref{algebra-definition-relative-dimension}

\noindent
{\it 相对维数}
见 \ref{stacks-geometry-definition-relative-dimension-for-stacks}

\noindent
{\it 相对对偶化复形}
见 \ref{dualizing-definition-relative-dualizing-complex}

\noindent
{\it 相对对偶化复形}
见 \ref{duality-definition-relative-dualizing-complex}

\noindent
{\it 相对对偶化复形}
见 \ref{spaces-duality-definition-relative-dualizing-proper-flat}

\noindent
{\it 相对对偶层}
见 \ref{moduli-curves-definition-relative-dualizing-sheaf}

\noindent
{\it 相对有效 Cartier 除子}
见 \ref{divisors-definition-relative-effective-Cartier-divisor}

\noindent
{\it 相对有效 Cartier 除子}
见 \ref{spaces-divisors-definition-relative-effective-Cartier-divisor}

\noindent
{\it $X/S$ 的相对 Frobenius 态射}
见 \ref{varieties-definition-relative-frobenius}

\noindent
{\it 相对整体完全交}
见 \ref{algebra-definition-relative-global-complete-intersection}

\noindent
{\it $\mathcal{A}$ 在 $S$ 上的相对齐次谱}
见 \ref{constructions-definition-relative-proj}

\noindent
{\it $\mathcal{A}$ 在 $X$ 上的相对齐次谱}
见 \ref{spaces-divisors-definition-relative-proj}

\noindent
{\it $\mathcal{S}$ 在 $\mathcal{S}'$ 上的相对惯性}
见 \ref{categories-definition-inertia-fibred-category}

\noindent
{\it $\mathcal{A}$ 在 $S$ 上的相对 Proj}
见 \ref{constructions-definition-relative-proj}

\noindent
{\it $\mathcal{A}$ 在 $X$ 上的相对 Proj}
见 \ref{spaces-divisors-definition-relative-proj}

\noindent
{\it 相对拟正则浸入}
见 \ref{divisors-definition-relative-H1-regular-immersion}

\noindent
{\it $x$ 的相对自同构层}
见 \ref{stacks-morphisms-definition-isom}

\noindent
{\it 从 $x_1$ 到 $x_2$ 的相对同构层}
见 \ref{stacks-morphisms-definition-isom}

\noindent
{\it $\mathcal{A}$ 在 $S$ 上的相对谱}
见 \ref{constructions-definition-relative-spec}

\noindent
{\it $\mathcal{A}$ 在 $X$ 上的相对谱}
见 \ref{spaces-morphisms-definition-relative-spec}

\noindent
{\it $\mathcal{F}$ 在 $X$ 中相对于 $S$ 的相对弱伴随点集}
见 \ref{divisors-definition-relative-weak-assassin}

\noindent
{\it $\mathcal{F}$ 在 $X$ 中相对于 $Y$ 的相对弱伴随点集}
见 \ref{spaces-divisors-definition-relative-weak-assassin}

\noindent
{\it 相对丰沛}
见 \ref{morphisms-definition-relatively-ample}

\noindent
{\it 相对丰沛}
见 \ref{spaces-divisors-definition-relatively-ample}

\noindent
{\it 相对保持极限}
见 \ref{spaces-limits-definition-locally-finite-presentation}

\noindent
{\it 互素}
见 \ref{fields-definition-relatively-prime}

\noindent
{\it 相对甚丰沛}
见 \ref{morphisms-definition-very-ample}

\noindent
{\it 由概形表示}
见 \ref{schemes-definition-representable-functor}

\noindent
{\it 可由代数空间表示}
见 \ref{bootstrap-definition-morphism-representable-by-spaces}

\noindent
{\it 可由代数空间表示}
见 \ref{algebraic-definition-representable-by-algebraic-spaces}

\noindent
{\it 可由 $S$ 上的代数空间表示}
见 \ref{algebraic-definition-representable-by-algebraic-space}

\noindent
{\it 由开浸入表示}
见 \ref{schemes-definition-representable-by-open-immersions}

\noindent
{\it 可表商}
见 \ref{groupoids-definition-representable-quotient}

\noindent
{\it 可表商}
见 \ref{groupoids-definition-representable-quotient}

\noindent
{\it 可表商}
见 \ref{spaces-groupoids-definition-representable-quotient}

\noindent
{\it 可表商}
见 \ref{spaces-groupoids-definition-representable-quotient}

\noindent
{\it 可表层}
见 \ref{sites-definition-representable-sheaf}

\noindent
{\it 可表示}
见 \ref{categories-definition-representable-functor}

\noindent
{\it 可表示}
见 \ref{categories-definition-representable-morphism}

\noindent
{\it 可表示}
见 \ref{categories-definition-representable-map-presheaves}

\noindent
{\it 可表示}
见 \ref{categories-definition-representable-fibred-category}

\noindent
{\it 可表示}
见 \ref{categories-definition-representable-map-categories-fibred-in-groupoids}

\noindent
{\it 可表示}
见 \ref{schemes-definition-representable-functor}

\noindent
{\it 可表示}
见 \ref{formal-defos-definition-representable}

\noindent
{\it 残余次数}
见 \ref{more-algebra-definition-extension-discrete-valuation-rings}

\noindent
{\it 残余次数}
见 \ref{more-algebra-definition-extension-valuation-rings}

\noindent
{\it $\mathcal{X}$ 在 $x$ 处的剩余胚存在}
见 \ref{stacks-properties-definition-residual-gerbe}

\noindent
{\it $\mathcal{X}$ 在 $x$ 处的剩余胚}
见 \ref{stacks-properties-definition-residual-gerbe}

\noindent
{\it residual space of $X$ at $x$}
见 \ref{decent-spaces-definition-residual-space}

\noindent
{\it 剩余次数}
见 \ref{more-algebra-definition-extension-discrete-valuation-rings}

\noindent
{\it 剩余次数}
见 \ref{more-algebra-definition-extension-valuation-rings}

\noindent
{\it $X$ 在 $x$ 处的剩余域}
见 \ref{schemes-definition-locally-ringed-space}

\noindent
{\it $X$ 在 $x$ 处的剩余域}
见 \ref{decent-spaces-definition-residue-field}

\noindent
{\it 剩余域}
见 \ref{algebra-definition-local-ring}

\noindent
{\it 消解函子}
见 \ref{derived-definition-localization-functor}

\noindent
{\it 由有限自由 $R$-模给出的 $M$ 的消解}
见 \ref{algebra-definition-finite-free-resolution}

\noindent
{\it 由自由 $R$-模给出的 $M$ 的消解}
见 \ref{algebra-definition-finite-free-resolution}

\noindent
{\it 由正规化吹起给出的奇点消解}
见 \ref{resolve-definition-resolution-surface}

\noindent
{\it 由正规化吹起给出的奇点消解}
见 \ref{spaces-resolve-definition-resolution-surface}

\noindent
{\it 奇点消解}
见 \ref{resolve-definition-resolution}

\noindent
{\it 奇点消解}
见 \ref{spaces-resolve-definition-resolution}

\noindent
{\it 消解性质}
见 \ref{perfect-definition-resolution-property}

\noindent
{\it 消解性质}
见 \ref{spaces-perfect-definition-resolution-property}

\noindent
{\it 消解}
见 \ref{algebra-definition-finite-free-resolution}

\noindent
{\it $(U, R, s, t, c)$ 到 $\mathcal{C}'$ 的限制 $(U, R, s, t, c)|_{\mathcal{C}'}$}
见 \ref{formal-defos-definition-restricting-groupoids-in-functors}

\noindent
{\it $(U, R, s, t, c)$ 到 $U'$ 的限制}
见 \ref{groupoids-definition-restrict-groupoid}

\noindent
{\it $(U, R, s, t, c)$ 到 $U'$ 的限制}
见 \ref{spaces-groupoids-definition-restrict-groupoid}

\noindent
{\it $\mathcal{F}$ 到 $\mathcal{C}/U$ 的限制}
见 \ref{sites-definition-localize}

\noindent
{\it $\mathcal{F}$ 到 $\mathcal{C}/U$ 的限制}
见 \ref{sites-modules-definition-localize-ringed-site}

\noindent
{\it $\mathcal{F}$ 到 $U_\etale$ 的限制}
见 \ref{stacks-sheaves-definition-pullback}

\noindent
{\it $\mathcal{G}$ 到 $U$ 的限制}
见 \ref{sheaves-definition-restriction}

\noindent
{\it $\mathcal{G}$ 到 $U$ 的限制}
见 \ref{sheaves-definition-restriction}

\noindent
{\it $\mathcal{G}$ 到 $U$ 的限制}
见 \ref{sheaves-definition-restriction}

\noindent
{\it 到小平展位点的限制}
见 \ref{topologies-definition-restriction-small-etale}

\noindent
{\it 到小平展位点的限制}
见 \ref{spaces-topologies-definition-restriction-small-etale}

\noindent
{\it 到小 pro-\'etale 位点的限制}
见 \ref{proetale-definition-restriction-small-proetale}

\noindent
{\it 到小 Zariski 位点的限制}
见 \ref{topologies-definition-restriction-small-zariski}

\noindent
{\it 限制}
见 \ref{groupoids-definition-restrict-relation}

\noindent
{\it 限制}
见 \ref{spaces-groupoids-definition-restrict-relation}

\noindent
{\it 逆紧}
见 \ref{topology-definition-quasi-compact}

\noindent
{\it 在 $(A, I)$ 上 rig-平展}
见 \ref{restricted-definition-rig-etale-homomorphism}

\noindent
{\it rig-\,'etale 的}
见 \ref{restricted-definition-rig-etale}

\noindent
{\it rig-闭}
见 \ref{restricted-definition-rig-closed}

\noindent
{\it rig-etale 的}
见 \ref{restricted-definition-rig-etale-continuous-homomorphism}

\noindent
{\it rig-平坦的}
见 \ref{restricted-definition-rig-flat-continuous-homomorphism}

\noindent
{\it rig-平坦的}
见 \ref{restricted-definition-rig-flat}

\noindent
{\it 在 $(A, I)$ 上 rig-光滑}
见 \ref{restricted-definition-rig-smooth-homomorphism}

\noindent
{\it rig-光滑的}
见 \ref{restricted-definition-rig-smooth-continuous-homomorphism}

\noindent
{\it rig-光滑的}
见 \ref{restricted-definition-rig-smooth}

\noindent
{\it rig-满射的}
见 \ref{restricted-definition-rig-surjective}

\noindent
{\it 对 $F$ 右无环}
见 \ref{derived-definition-derived-functor}

\noindent
{\it 右伴随}
见 \ref{categories-definition-adjoint}

\noindent
{\it 右可容许}
见 \ref{derived-definition-admissible}

\noindent
{\it 右可导}
见 \ref{derived-definition-everywhere-defined}

\noindent
{\it 右导出函子 $RF$ 在此处有定义}
见 \ref{derived-definition-right-derived-functor-defined}

\noindent
{\it $F$ 的右导出函子}
见 \ref{derived-definition-derived-functor}

\noindent
{\it 右对偶}
见 \ref{categories-definition-dual}

\noindent
{\it 右正合}
见 \ref{categories-definition-exact}

\noindent
{\it 右乘法系}
见 \ref{categories-definition-multiplicative-system}

\noindent
{\it 右正交}
见 \ref{derived-definition-orthogonal}

\noindent
{\it $X$ 上的有理函数环}
见 \ref{morphisms-definition-ring-of-rational-functions}

\noindent
{\it $X$ 上的有理函数环}
见 \ref{spaces-morphisms-definition-ring-of-rational-functions}

\noindent
{\it 环化位点}
见 \ref{sites-modules-definition-ringed-site}

\noindent
{\it 环化位点}
见 \ref{etale-cohomology-definition-ringed-site}

\noindent
{\it 赋环空间}
见 \ref{sheaves-definition-ringed-space}

\noindent
{\it 环化拓扑斯}
见 \ref{sites-modules-definition-ringed-topos}

\noindent
{\it 满足赋值判据的存在性部分}
见 \ref{schemes-definition-valuative-criterion}

\noindent
{\it 满足赋值判据的存在性部分}
见 \ref{spaces-morphisms-definition-valuative-criterion}

\noindent
{\it 满足 fpqc 拓扑的层性质}
见 \ref{topologies-definition-sheaf-property-fpqc}

\noindent
{\it 满足 fpqc 拓扑的层性质}
见 \ref{etale-cohomology-definition-sheaf-property-fpqc}

\noindent
{\it 对给定态射族满足层性质}
见 \ref{topologies-definition-sheaf-property-fpqc}

\noindent
{\it 满足 V 拓扑的层性质}
见 \ref{topologies-definition-sheaf-property-V}

\noindent
{\it 满足 Zariski 拓扑的层性质}
见 \ref{schemes-definition-representable-by-open-immersions}

\noindent
{\it 满足赋值判据的唯一性部分}
见 \ref{schemes-definition-valuative-criterion}

\noindent
{\it 满足赋值判据的唯一性部分}
见 \ref{spaces-morphisms-definition-valuative-criterion}

\noindent
{\it 满足赋值判据}
见 \ref{spaces-morphisms-definition-valuative-criterion}

\noindent
{\it 饱和}
见 \ref{categories-definition-saturated-multiplicative-system}

\noindent
{\it 饱和}
见 \ref{derived-definition-saturated}

\noindent
{\it $R$ 上的概形}
见 \ref{schemes-definition-base-change}

\noindent
{\it $S$ 上的概形}
见 \ref{schemes-definition-base-change}

\noindent
{\it $Z$ 上的概形结构}
见 \ref{schemes-definition-reduced-induced-scheme}

\noindent
{\it $U$ 在 $X$ 中的概形论闭包}
见 \ref{morphisms-definition-scheme-theoretically-dense}

\noindent
{\it $U$ 在 $X$ 中的概形论闭包}
见 \ref{spaces-morphisms-definition-scheme-theoretically-dense}

\noindent
{\it 概形论纤维 $X_s$，即 $f$ 在 $s$ 上方的纤维}
见 \ref{schemes-definition-fibre}

\noindent
{\it 概形论像}
见 \ref{morphisms-definition-scheme-theoretic-image}

\noindent
{\it 概形论像}
见 \ref{spaces-morphisms-definition-scheme-theoretic-image}

\noindent
{\it 概形论像}
见 \ref{stacks-morphisms-definition-scheme-theoretic-image}

\noindent
{\it 概形论交}
见 \ref{morphisms-definition-scheme-theoretic-intersection-union}

\noindent
{\it 概形论交}
见 \ref{spaces-morphisms-definition-scheme-theoretic-intersection-union}

\noindent
{\it $\mathcal{F}$ 的概形论支集}
见 \ref{morphisms-definition-scheme-theoretic-support}

\noindent
{\it $\mathcal{F}$ 的概形论支集}
见 \ref{spaces-morphisms-definition-scheme-theoretic-support}

\noindent
{\it 概形论并}
见 \ref{morphisms-definition-scheme-theoretic-intersection-union}

\noindent
{\it 概形论并}
见 \ref{spaces-morphisms-definition-scheme-theoretic-intersection-union}

\noindent
{\it 在 $X$ 中概形论稠密}
见 \ref{morphisms-definition-scheme-theoretically-dense}

\noindent
{\it 在 $X$ 中概形论稠密}
见 \ref{spaces-morphisms-definition-scheme-theoretically-dense}

\noindent
{\it 概形}
见 \ref{schemes-definition-scheme}

\noindent
{\it 紧支撑截面}
见 \ref{more-etale-definition-compact-support}

\noindent
{\it $X$ 上的半可表示对象}
见 \ref{hypercovering-definition-SR}

\noindent
{\it 半可表示对象}
见 \ref{hypercovering-definition-SR}

\noindent
{\it $X$ 在 $Y$ 中的半正规化}
见 \ref{morphisms-definition-relative-seminormalization}

\noindent
{\it 半正规化}
见 \ref{morphisms-definition-seminormalization}

\noindent
{\it 半正规}
见 \ref{morphisms-definition-seminormal-ring}

\noindent
{\it 半正规}
见 \ref{morphisms-definition-seminormal}

\noindent
{\it 半稳定曲线族}
见 \ref{moduli-curves-definition-semistable}

\noindent
{\it 半稳定约化}
见 \ref{models-definition-semistable}

\noindent
{\it 可分次数}
见 \ref{fields-definition-separable-degree}

\noindent
{\it 可分次数}
见 \ref{fields-definition-insep-degree}

\noindent
{\it 在 $k$ 上可分}
见 \ref{algebra-definition-separable-field-extension}

\noindent
{\it 可分解}
见 \ref{more-algebra-definition-solution}

\noindent
{\it 可分}
见 \ref{fields-definition-separable}

\noindent
{\it 可分}
见 \ref{fields-definition-separable}

\noindent
{\it 可分}
见 \ref{fields-definition-separable}

\noindent
{\it 可分}
见 \ref{fields-definition-separable-algebraic}

\noindent
{\it 在 $k$ 上可分生成}
见 \ref{algebra-definition-separable-field-extension}

\noindent
{\it 分离群概形}
见 \ref{groupoids-definition-smooth-group-scheme}

\noindent
{\it 在 $S$ 上分离}
见 \ref{spaces-definition-separated}

\noindent
{\it 在 $S$ 上分离}
见 \ref{stacks-morphisms-definition-absolute-separated}

\noindent
{\it 分离预层}
见 \ref{etale-cohomology-definition-sheaf}

\noindent
{\it 分离}
见 \ref{topology-definition-separated}

\noindent
{\it 分离}
见 \ref{sheaves-definition-separated}

\noindent
{\it 分离}
见 \ref{sites-definition-separated}

\noindent
{\it 分离}
见 \ref{sites-definition-presheaf-separated-topology}

\noindent
{\it 分离}
见 \ref{homology-definition-filtered}

\noindent
{\it 分离}
见 \ref{schemes-definition-separated}

\noindent
{\it 分离}
见 \ref{schemes-definition-separated}

\noindent
{\it 分离}
见 \ref{spaces-properties-definition-separated}

\noindent
{\it 分离}
见 \ref{spaces-properties-definition-separated}

\noindent
{\it 分离}
见 \ref{spaces-morphisms-definition-separated}

\noindent
{\it 分离}
见 \ref{formal-spaces-definition-separated}

\noindent
{\it 分离}
见 \ref{formal-spaces-definition-separated-morphism}

\noindent
{\it 分离}
见 \ref{stacks-morphisms-definition-separated}

\noindent
{\it 分离}
见 \ref{stacks-morphisms-definition-absolute-separated}

\noindent
{\it 分离 $R$-轨道}
见 \ref{groupoids-quotients-definition-set-theoretically-invariant}

\noindent
{\it 分离轨道}
见 \ref{groupoids-quotients-definition-set-theoretically-invariant}

\noindent
{\it Serre 函子}
见 \ref{equiv-definition-Serre-functor}

\noindent
{\it Serre 子范畴}
见 \ref{homology-definition-serre-subcategory}

\noindent
{\it 集合论等价关系}
见 \ref{groupoids-quotients-definition-set-theoretic-equivalence}

\noindent
{\it 集合论预等价关系}
见 \ref{groupoids-quotients-definition-set-theoretic-equivalence}

\noindent
{\it 在集合论意义下 $R$-不变}
见 \ref{groupoids-definition-invariant-open}

\noindent
{\it 在集合论意义下 $R$-不变}
见 \ref{groupoids-quotients-definition-set-theoretically-invariant}

\noindent
{\it 集合胚（setoid）}
见 \ref{categories-definition-setoid}

\noindent
{\it $\mathcal{B}$ 上的 $\mathcal{O}$-模层 $\mathcal{F}$}
见 \ref{sheaves-definition-sheaf-modules-basis}

\noindent
{\it $\mathcal{B}$ 上的集合层 $\mathcal{F}$}
见 \ref{sheaves-definition-sheaf-basis}

\noindent
{\it $X$ 上的集合层 $\mathcal{F}$}
见 \ref{sheaves-definition-sheaf}

\noindent
{\it $\mathcal{B}$ 上取值于 $\mathcal{C}$ 的层 $\mathcal{F}$}
见 \ref{sheaves-definition-sheaf-structures-basis}

\noindent
{\it 与 $\mathcal{F}$ 关联的层}
见 \ref{sites-definition-associated-sheaf}

\noindent
{\it 与 $\mathcal{F}$ 关联的层}
见 \ref{sites-definition-associated-sheaf-topology}

\noindent
{\it 与模 $M$ 及环映射 $\alpha$ 关联的层}
见 \ref{modules-definition-sheaf-associated}

\noindent
{\it 与模 $M$ 关联的层}
见 \ref{modules-definition-sheaf-associated}

\noindent
{\it 关于 \'etale 拓扑的层}
见 \ref{stacks-sheaves-definition-sheaves}

\noindent
{\it 关于 fppf 拓扑的层}
见 \ref{stacks-sheaves-definition-sheaves}

\noindent
{\it 关于光滑拓扑的层}
见 \ref{stacks-sheaves-definition-sheaves}

\noindent
{\it 关于 syntomic 拓扑的层}
见 \ref{stacks-sheaves-definition-sheaves}

\noindent
{\it 关于 Zariski 拓扑的层}
见 \ref{stacks-sheaves-definition-sheaves}

\noindent
{\it 与 $\mathcal{F}$ 相联系的 $\mathcal{O}$-模层}
见 \ref{descent-definition-structure-sheaf}

\noindent
{\it 与 $\mathcal{F}$ 相联系的 $\mathcal{O}$-模层}
见 \ref{descent-definition-structure-sheaf}

\noindent
{\it $\mathcal{O}$-模层}
见 \ref{sheaves-definition-sheaf-modules}

\noindent
{\it $\mathcal{O}$-模层}
见 \ref{sites-modules-definition-sheaf-modules}

\noindent
{\it $\mathcal{O}_\mathcal{X}$-模层}
见 \ref{stacks-sheaves-definition-modules}

\noindent
{\it $X$ 上的 $R$-不变函数层}
见 \ref{groupoids-quotients-definition-functions}

\noindent
{\it $X$ 上的阿贝尔群层}
见 \ref{sheaves-definition-abelian-sheaf}

\noindent
{\it $x$ 的自同构层}
见 \ref{stacks-morphisms-definition-isom}

\noindent
{\it 微分分次 $\mathcal{O}$-代数层}
见 \ref{sdga-definition-dga}

\noindent
{\it 微分分次代数层}
见 \ref{sdga-definition-dga}

\noindent
{\it $X$ 在 $S$ 上的微分层 $\Omega_{X/S}$}
见 \ref{modules-definition-differentials}

\noindent
{\it $X$ 在 $S$ 上的微分层 $\Omega_{X/S}$}
见 \ref{morphisms-definition-sheaf-differentials}

\noindent
{\it $X$ 在 $Y$ 上的微分层 $\Omega_{X/Y}$}
见 \ref{sites-modules-definition-sheaf-differentials}

\noindent
{\it $X$ 在 $Y$ 上的微分层 $\Omega_{X/Y}$}
见 \ref{spaces-more-morphisms-definition-sheaf-differentials}

\noindent
{\it 分次 $\mathcal{O}$-代数层}
见 \ref{sdga-definition-ga}

\noindent
{\it 分次代数层}
见 \ref{sdga-definition-ga}

\noindent
{\it 从 $x_1$ 到 $x_2$ 的同构层}
见 \ref{stacks-morphisms-definition-isom}

\noindent
{\it $X$ 上的亚纯函数层}
见 \ref{divisors-definition-sheaf-meromorphic-functions}

\noindent
{\it $X$ 上的亚纯函数层}
见 \ref{spaces-divisors-definition-sheaf-meromorphic-functions}

\noindent
{\it 全商环层 ${\mathcal K}_S$}
见 \ref{exercises-definition-divisor}

\noindent
{\it 层论意义下为空}
见 \ref{sites-definition-empty}

\noindent
{\it 层}
见 \ref{sheaves-definition-sheaf-values-in-category}

\noindent
{\it 层}
见 \ref{sites-definition-sheaf-sets}

\noindent
{\it 层}
见 \ref{sites-definition-sheaf}

\noindent
{\it 层}
见 \ref{sites-definition-sheaf-sets-topology}

\noindent
{\it 层}
见 \ref{etale-cohomology-definition-sheaf}

\noindent
{\it 层}
见 \ref{stacks-sheaves-definition-sheaves}

\noindent
{\it 移位}
见 \ref{homology-definition-graded-shift}

\noindent
{\it 短正合列}
见 \ref{homology-definition-exact}

\noindent
{\it 同胞对象}
见 \ref{equiv-definition-siblings}

\noindent
{\it 同胞对象}
见 \ref{equiv-definition-siblings-geometric}

\noindent
{\it 同胞}
见 \ref{equiv-definition-siblings}

\noindent
{\it 同胞}
见 \ref{equiv-definition-siblings-geometric}

\noindent
{\it 由态射 $f_i$ 生成的 $U$ 上的筛}
见 \ref{sites-definition-sieve-generated}

\noindent
{\it $U$ 上的筛 $S$}
见 \ref{sites-definition-sieve}

\noindent
{\it 相似}
见 \ref{etale-cohomology-definition-brauer-equivalent}

\noindent
{\it 单}
见 \ref{algebra-definition-simple-module}

\noindent
{\it 单}
见 \ref{brauer-definition-simple}

\noindent
{\it 单}
见 \ref{brauer-definition-simple}

\noindent
{\it 单}
见 \ref{homology-definition-simple}

\noindent
{\it 单纯 $\mathcal{A}_\bullet$-模}
见 \ref{sites-cohomology-definition-simplicial-module}

\noindent
{\it 单纯阿贝尔群}
见 \ref{simplicial-definition-simplicial-object}

\noindent
{\it $\mathcal{C}$ 的单纯对象 $U$}
见 \ref{simplicial-definition-simplicial-object}

\noindent
{\it 与 $f$ 关联的单纯概形}
见 \ref{spaces-simplicial-definition-fibre-products-simplicial-scheme}

\noindent
{\it 单纯集}
见 \ref{simplicial-definition-simplicial-object}

\noindent
{\it 单纯 $\mathcal{A}_\bullet$-模层}
见 \ref{sites-cohomology-definition-simplicial-module}

\noindent
{\it 模导出范畴的单纯系统}
见 \ref{spaces-simplicial-definition-cartesian-derived-modules}

\noindent
{\it 导出范畴的单纯系统}
见 \ref{spaces-simplicial-definition-cartesian-derived}

\noindent
{\it $A$ 在 $R$ 上的奇异理想}
见 \ref{smoothing-definition-singular-ideal}

\noindent
{\it 奇异轨迹}
见 \ref{properties-definition-singular-locus}

\noindent
{\it $X$ 的奇点至多为结点型}
见 \ref{curves-definition-nodal}

\noindent
{\it 位点}
见 \ref{sites-definition-site}

\noindent
{\it 位点}
见 \ref{etale-cohomology-definition-site}

\noindent
{\it 大小}
见 \ref{injectives-definition-size}

\noindent
{\it 除环}
见 \ref{brauer-definition-skew-field}

\noindent
{\it 在 $x$ 处取值为 $A$ 的摩天楼层}
见 \ref{sheaves-definition-skyscraper-sheaf}

\noindent
{\it 摩天楼层}
见 \ref{sheaves-definition-skyscraper-sheaf}

\noindent
{\it 摩天楼层}
见 \ref{sheaves-definition-skyscraper-sheaf}

\noindent
{\it 摩天楼层}
见 \ref{sheaves-definition-skyscraper-sheaf}

\noindent
{\it 摩天楼层}
见 \ref{sheaves-definition-skyscraper-sheaf}

\noindent
{\it 摩天楼层}
见 \ref{sites-definition-pushforward-point}

\noindent
{\it $S$ 的小 $\tau$-位点}
见 \ref{etale-cohomology-definition-tau-site}

\noindent
{\it 小 \'etale 站点 $X_\etale$}
见 \ref{spaces-properties-definition-etale-site}

\noindent
{\it $S$ 的小平展位点}
见 \ref{topologies-definition-big-small-etale}

\noindent
{\it $S$ 上的小平展位点}
见 \ref{etale-cohomology-definition-big-etale-site}

\noindent
{\it 小平展位点}
见 \ref{formal-spaces-definition-etale-sites}

\noindent
{\it 小 \'etale 拓扑斯}
见 \ref{etale-cohomology-definition-etale-topos}

\noindent
{\it 小 \'etale 拓扑斯}
见 \ref{spaces-properties-definition-etale-topos}

\noindent
{\it $S$ 的小仿射平展位点}
见 \ref{topologies-definition-big-small-etale}

\noindent
{\it $S$ 的小仿射 Zariski 位点}
见 \ref{topologies-definition-big-small-Zariski}

\noindent
{\it 小扩张}
见 \ref{algebra-definition-small-extension}

\noindent
{\it 小扩张}
见 \ref{formal-defos-definition-small-extension}

\noindent
{\it $S$ 的小 pro-平展位点}
见 \ref{proetale-definition-big-small-proetale}

\noindent
{\it 小 Zariski 位点 $F_{Zar}$}
见 \ref{spaces-definition-small-Zariski-site}

\noindent
{\it $S$ 的小 Zariski 位点}
见 \ref{topologies-definition-big-small-Zariski}

\noindent
{\it 小 Zariski 位点}
见 \ref{etale-cohomology-definition-big-etale-site}

\noindent
{\it 小 Zariski 拓扑斯}
见 \ref{etale-cohomology-definition-etale-topos}

\noindent
{\it 在 $\mathfrak q$ 处光滑}
见 \ref{algebra-definition-smooth-at-prime}

\noindent
{\it 在 $x \in X$ 处光滑}
见 \ref{morphisms-definition-smooth}

\noindent
{\it 在 $x$ 处光滑}
见 \ref{spaces-morphisms-definition-smooth}

\noindent
{\it $T$ 的光滑覆盖}
见 \ref{topologies-definition-smooth-covering}

\noindent
{\it $X$ 的光滑覆盖}
见 \ref{spaces-topologies-definition-smooth-covering}

\noindent
{\it 光滑群概形}
见 \ref{groupoids-definition-smooth-group-scheme}

\noindent
{\it 光滑群胚}
见 \ref{algebraic-definition-smooth-groupoid}

\noindent
{\it 在源与靶上关于光滑拓扑为局部的}
见 \ref{spaces-descent-definition-local-source-target}

\noindent
{\it smooth 局部}
见 \ref{descent-definition-local-at-point}

\noindent
{\it 相对维数 $d$ 的光滑态射}
见 \ref{morphisms-definition-smooth-relative-dimension}

\noindent
{\it 光滑层}
见 \ref{stacks-sheaves-definition-sheaves}

\noindent
{\it 光滑簇}
见 \ref{varieties-definition-variety-type}

\noindent
{\it 光滑}
见 \ref{algebra-definition-smooth}

\noindent
{\it 光滑}
见 \ref{morphisms-definition-smooth}

\noindent
{\it 光滑}
见 \ref{descent-definition-etale-morphism-germs}

\noindent
{\it 光滑}
见 \ref{spaces-morphisms-definition-smooth}

\noindent
{\it 光滑}
见 \ref{formal-defos-definition-smooth-morphism}

\noindent
{\it 光滑}
见 \ref{formal-defos-definition-cofibered-groupoid-projection-smooth}

\noindent
{\it 光滑}
见 \ref{formal-defos-definition-smooth-groupoid-in-functors}

\noindent
{\it 光滑}
见 \ref{stacks-morphisms-definition-smooth}

\noindent
{\it 光滑}
见 \ref{stacks-introduction-definition-smooth}

\noindent
{\it 清醒}
见 \ref{topology-definition-generic-point}

\noindent
{\it $A \subset B$ 的解}
见 \ref{more-algebra-definition-solution}

\noindent
{\it 从 $\mathcal{C}$ 到 $\mathcal{D}$ 的特殊余连续函子 $u$}
见 \ref{sites-definition-special-cocontinuous-functor}

\noindent
{\it 特化可沿 $f$ 提升}
见 \ref{topology-definition-lift-specializations}

\noindent
{\it 特化}
见 \ref{topology-definition-specialization}

\noindent
{\it 特化}
见 \ref{exercises-definition-closed}

\noindent
{\it 特化}
见 \ref{exercises-definition-specialization}

\noindent
{\it 特化型的}
见 \ref{topology-definition-lift-specializations}

\noindent
{\it 与 $(A, d, \alpha)$ 关联的谱序列}
见 \ref{homology-definition-differential-object-selfmap}

\noindent
{\it 与正合偶关联的谱序列}
见 \ref{homology-definition-spectral-sequence-associated-exact-couple}

\noindent
{\it $\mathcal{A}$ 中的谱序列}
见 \ref{homology-definition-spectral-sequence}

\noindent
{\it 谱的}
见 \ref{topology-definition-spectral-space}

\noindent
{\it 谱的}
见 \ref{topology-definition-spectral-space}

\noindent
{\it $\mathcal{A}$ 在 $S$ 上的谱}
见 \ref{constructions-definition-relative-spec}

\noindent
{\it $\mathcal{A}$ 在 $X$ 上的谱}
见 \ref{spaces-morphisms-definition-relative-spec}

\noindent
{\it 谱}
见 \ref{algebra-definition-spectrum-ring}

\noindent
{\it 谱}
见 \ref{schemes-definition-structure-sheaf}

\noindent
{\it 分裂群胚纤维化范畴}
见 \ref{categories-definition-split-category-fibred-in-groupoids}

\noindent
{\it 可裂等化子}
见 \ref{descent-definition-split-equalizer}

\noindent
{\it 分裂纤维化范畴}
见 \ref{categories-definition-split-fibred-category}

\noindent
{\it 分裂结点}
见 \ref{curves-definition-split-node}

\noindent
{\it 在 $u$ 上分裂}
见 \ref{spaces-more-groupoids-definition-split-at-point}

\noindent
{\it 分裂}
见 \ref{brauer-definition-splitting}

\noindent
{\it $P$ 在 $F$ 上的分裂域}
见 \ref{fields-definition-splitting-field}

\noindent
{\it 分裂域}
见 \ref{brauer-definition-splitting}

\noindent
{\it $R$ 在 $u$ 上的分裂}
见 \ref{spaces-more-groupoids-definition-split-at-point}

\noindent
{\it 分裂}
见 \ref{homology-definition-ses-split}

\noindent
{\it 分裂}
见 \ref{simplicial-definition-split}

\noindent
{\it 分裂}
见 \ref{simplicial-definition-split}

\noindent
{\it 代数空间中的群胚 $(U, R, s, t, c)$ 的稳定子}
见 \ref{spaces-groupoids-definition-stabilizer-groupoid}

\noindent
{\it 群胚概形 $(U, R, s, t, c)$ 的稳定子}
见 \ref{groupoids-definition-stabilizer-groupoid}

\noindent
{\it 稳定曲线族}
见 \ref{moduli-curves-definition-stable}

\noindent
{\it 基变换下稳定}
见 \ref{morphisms-definition-property-local}

\noindent
{\it 复合下稳定}
见 \ref{morphisms-definition-property-local}

\noindent
{\it 对一般化稳定}
见 \ref{topology-definition-specialization}

\noindent
{\it 对特化稳定}
见 \ref{topology-definition-specialization}

\noindent
{\it 稳定自由}
见 \ref{more-algebra-definition-stably-free}

\noindent
{\it 稳定同构}
见 \ref{more-algebra-definition-stably-free}

\noindent
{\it 离散范畴中的栈}
见 \ref{stacks-definition-stack-in-sets}

\noindent
{\it 群胚中的栈}
见 \ref{stacks-definition-stack-in-groupoids}

\noindent
{\it 类集合（setoid）中的栈}
见 \ref{stacks-definition-stack-in-sets}

\noindent
{\it 集合中的栈}
见 \ref{stacks-definition-stack-in-sets}

\noindent
{\it 栈}
见 \ref{stacks-definition-stack}

\noindent
{\it 茎}
见 \ref{etale-cohomology-definition-stalk}

\noindent
{\it 茎}
见 \ref{trace-definition-torsion-l-adic-sheaf}

\noindent
{\it 茎}
见 \ref{spaces-properties-definition-stalk}

\noindent
{\it 标准 $\tau$-覆盖}
见 \ref{etale-cohomology-definition-standard-tau}

\noindent
{\it 标准平展覆盖}
见 \ref{topologies-definition-standard-etale}

\noindent
{\it 标准平展的}
见 \ref{algebra-definition-standard-etale}

\noindent
{\it 标准平展的}
见 \ref{morphisms-definition-etale}

\noindent
{\it 标准平展的}
见 \ref{etale-cohomology-definition-standard-etale}

\noindent
{\it 标准 fppf 覆盖}
见 \ref{topologies-definition-standard-fppf}

\noindent
{\it 标准 fpqc 覆盖}
见 \ref{topologies-definition-standard-fpqc}

\noindent
{\it 标准 h 覆盖}
见 \ref{flat-definition-standard-h}

\noindent
{\it 标准开覆盖}
见 \ref{schemes-definition-standard-covering}

\noindent
{\it 标准开覆盖}
见 \ref{schemes-definition-standard-covering}

\noindent
{\it 标准开覆盖}
见 \ref{constructions-definition-standard-covering}

\noindent
{\it 标准开集}
见 \ref{algebra-definition-Zariski-topology}

\noindent
{\it 标准 ph 覆盖}
见 \ref{topologies-definition-standard-ph-covering}

\noindent
{\it 标准 pro-\'etale 覆盖}
见 \ref{proetale-definition-standard-proetale}

\noindent
{\it $\mathcal{B}$ 在 $\mathcal{A}$ 上的标准消解}
见 \ref{cotangent-definition-standard-resolution-sheaves-rings}

\noindent
{\it $B$ 在 $A$ 上的标准消解}
见 \ref{cotangent-definition-standard-resolution}

\noindent
{\it 标准缩小}
见 \ref{flat-definition-shrink}

\noindent
{\it 标准缩小}
见 \ref{flat-definition-shrink-complete}

\noindent
{\it $R$ 上的标准光滑代数}
见 \ref{algebra-definition-standard-smooth}

\noindent
{\it 标准光滑覆盖}
见 \ref{topologies-definition-standard-smooth}

\noindent
{\it 标准光滑}
见 \ref{algebra-definition-standard-smooth}

\noindent
{\it 标准光滑}
见 \ref{algebra-definition-standard-smooth}

\noindent
{\it 标准光滑}
见 \ref{morphisms-definition-smooth}

\noindent
{\it 标准 syntomic 覆盖}
见 \ref{topologies-definition-standard-syntomic}

\noindent
{\it 标准 syntomic}
见 \ref{morphisms-definition-syntomic}

\noindent
{\it 标准 V 覆盖}
见 \ref{topologies-definition-standard-V-covering}

\noindent
{\it 标准 Zariski 覆盖}
见 \ref{topologies-definition-standard-Zariski}

\noindent
{\it 层}
见 \ref{topology-definition-stratification}

\noindent
{\it 分层}
见 \ref{topology-definition-stratification}

\noindent
{\it $\mathcal{O}_{S, s}$ 的严格 Hensel 化}
见 \ref{etale-cohomology-definition-etale-local-rings}

\noindent
{\it $R$ 关于 $\kappa \subset \kappa^{sep}$ 的严格 Hensel 化}
见 \ref{algebra-definition-henselization}

\noindent
{\it $S$ 在 $\overline{s}$ 处的严格 Hensel 化}
见 \ref{etale-cohomology-definition-etale-local-rings}

\noindent
{\it $X$ 在 $\overline{x}$ 处的严格 Hensel 化}
见 \ref{spaces-properties-definition-etale-local-rings}

\noindent
{\it 严格 Hensel 化}
见 \ref{algebra-definition-henselization}

\noindent
{\it 拓扑空间的严格映射}
见 \ref{topology-definition-submersive}

\noindent
{\it 增厚的严格态射}
见 \ref{defos-definition-strict-morphism-thickenings}

\noindent
{\it 增厚的严格态射}
见 \ref{defos-definition-strict-morphism-thickenings-ringed-topoi}

\noindent
{\it 严格正规交叉除子}
见 \ref{etale-definition-strict-normal-crossings}

\noindent
{\it $M$ 沿 $R \to R[\frac{I}{a}]$ 的严格变换}
见 \ref{more-algebra-definition-strict-transform}

\noindent
{\it 严格变换}
见 \ref{divisors-definition-strict-transform}

\noindent
{\it 严格变换}
见 \ref{divisors-definition-strict-transform}

\noindent
{\it 严格变换}
见 \ref{spaces-divisors-definition-strict-transform}

\noindent
{\it 严格变换}
见 \ref{spaces-divisors-definition-strict-transform}

\noindent
{\it 严格可交换}
见 \ref{dga-definition-cdga}

\noindent
{\it 严格满}
见 \ref{categories-definition-subcategory}

\noindent
{\it 严格 Hensel}
见 \ref{algebra-definition-henselian}

\noindent
{\it 严格 Hensel}
见 \ref{etale-cohomology-definition-strictly-henselian}

\noindent
{\it 严格完美}
见 \ref{cohomology-definition-strictly-perfect}

\noindent
{\it 严格完美}
见 \ref{sites-cohomology-definition-strictly-perfect}

\noindent
{\it 在 $A$ 中相对于 $R$ 为严格标准}
见 \ref{smoothing-definition-strictly-standard}

\noindent
{\it 严格}
见 \ref{homology-definition-strict}

\noindent
{\it 强生成元}
见 \ref{derived-definition-generators}

\noindent
{\it $R$ 在 $u$ 上的强分裂}
见 \ref{spaces-more-groupoids-definition-split-at-point}

\noindent
{\it 更强}
见 \ref{sites-definition-finer}

\noindent
{\it 强 $\mathcal{C}$-笛卡尔态射}
见 \ref{categories-definition-cartesian-over-C}

\noindent
{\it 强笛卡尔态射}
见 \ref{categories-definition-cartesian-over-C}

\noindent
{\it 在 $u$ 上强分裂}
见 \ref{spaces-more-groupoids-definition-split-at-point}

\noindent
{\it 在 $R$ 上强超越}
见 \ref{algebra-definition-strongly-transcendental}

\noindent
{\it 结构态射}
见 \ref{schemes-definition-base-change}

\noindent
{\it $\mathcal{S}$ 上从 $\mathcal{C}$ 继承的位点结构}
见 \ref{stacks-definition-topology-inherited}

\noindent
{\it 结构层 $\mathcal{O}_{\Spec(R)}$，即 $R$ 的谱的结构层}
见 \ref{schemes-definition-structure-sheaf}

\noindent
{\it 结构层 $\mathcal{O}_{\text{Proj}(S)}$（关于 $S$）}
见 \ref{constructions-definition-structure-sheaf}

\noindent
{\it $\mathcal{X}$ 的结构层}
见 \ref{stacks-sheaves-definition-structure-sheaf}

\noindent
{\it 大位点 $(\Sch/S)_\tau$ 的结构层}
见 \ref{descent-definition-structure-sheaf}

\noindent
{\it 结构层}
见 \ref{descent-definition-structure-sheaf}

\noindent
{\it 结构层}
见 \ref{sites-modules-definition-ringed-site}

\noindent
{\it 结构层}
见 \ref{sites-modules-definition-ringed-topos}

\noindent
{\it 结构层}
见 \ref{etale-cohomology-definition-structure-sheaf}

\noindent
{\it 结构层}
见 \ref{spaces-properties-definition-structure-sheaf}

\noindent
{\it 子 $2$-范畴}
见 \ref{categories-definition-sub-2-category}

\noindent
{\it $X$ 上拓扑的一个子基}
见 \ref{topology-definition-subbase}

\noindent
{\it $X$ 上拓扑的一个子基底}
见 \ref{topology-definition-subbase}

\noindent
{\it 次典范的}
见 \ref{sites-definition-weaker-than-canonical}

\noindent
{\it 子范畴}
见 \ref{categories-definition-subcategory}

\noindent
{\it 子域}
见 \ref{fields-definition-field}

\noindent
{\it 子函子 $H \subset F$}
见 \ref{schemes-definition-representable-by-open-immersions}

\noindent
{\it 商映射}
见 \ref{topology-definition-submersive}

\noindent
{\it 商映射}
见 \ref{morphisms-definition-submersive}

\noindent
{\it 商映射}
见 \ref{spaces-morphisms-definition-submersive}

\noindent
{\it 商映射}
见 \ref{stacks-morphisms-definition-submersive}

\noindent
{\it 子对象}
见 \ref{homology-definition-injective-surjective}

\noindent
{\it 子预层}
见 \ref{sheaves-definition-injective-surjective}

\noindent
{\it 子预层}
见 \ref{sites-definition-sub-presheaf}

\noindent
{\it 由 $s_i$ 生成的子层}
见 \ref{modules-definition-generated-by-local-sections}

\noindent
{\it 被 $\mathcal{I}$ 零化的截面子层}
见 \ref{properties-definition-subsheaf-sections-annihilated-by-ideal}

\noindent
{\it 被 $\mathcal{I}$ 零化的截面子层}
见 \ref{spaces-limits-definition-subsheaf-sections-annihilated-by-ideal}

\noindent
{\it 支集位于 $T$ 上的截面子层}
见 \ref{properties-definition-subsheaf-sections-supported-on-closed}

\noindent
{\it 支集位于 $T$ 上的截面子层}
见 \ref{spaces-limits-definition-subsheaf-sections-supported-on-closed}

\noindent
{\it 子层}
见 \ref{sheaves-definition-injective-surjective}

\noindent
{\it 子拓扑斯}
见 \ref{sites-definition-subtopos}

\noindent
{\it 有效 Cartier 除子 $D_1$ 与 $D_2$ 的和}
见 \ref{divisors-definition-sum-effective-Cartier-divisors}

\noindent
{\it 有效 Cartier 除子 $D_1$ 与 $D_2$ 的和}
见 \ref{spaces-divisors-definition-sum-effective-Cartier-divisors}

\noindent
{\it 和}
见 \ref{obsolete-definition-locally-finite-sum-effective-Cartier-divisors}

\noindent
{\it $\mathcal{F}$ 的支集}
见 \ref{modules-definition-support}

\noindent
{\it $\mathcal{F}$ 的支集}
见 \ref{etale-cohomology-definition-support}

\noindent
{\it $\mathcal{F}$ 的支集}
见 \ref{spaces-properties-definition-support}

\noindent
{\it $\sigma$ 的支集}
见 \ref{etale-cohomology-definition-support}

\noindent
{\it $\sigma$ 的支集}
见 \ref{spaces-properties-definition-support}

\noindent
{\it $M$ 的支集}
见 \ref{algebra-definition-support-module}

\noindent
{\it $s$ 的支集}
见 \ref{modules-definition-support}

\noindent
{\it 支集在 $T$ 上}
见 \ref{perfect-definition-supported-on}

\noindent
{\it 支集在 $T$ 上}
见 \ref{spaces-perfect-definition-supported-on}

\noindent
{\it 支撑}
见 \ref{chow-definition-support-cycle}

\noindent
{\it 满}
见 \ref{sheaves-definition-injective-surjective}

\noindent
{\it 满}
见 \ref{sheaves-definition-injective-surjective}

\noindent
{\it 满}
见 \ref{sites-definition-presheaves-injective-surjective}

\noindent
{\it 满}
见 \ref{sites-definition-sheaves-injective-surjective}

\noindent
{\it 满}
见 \ref{homology-definition-injective-surjective}

\noindent
{\it 满}
见 \ref{morphisms-definition-surjective}

\noindent
{\it 满}
见 \ref{spaces-morphisms-definition-surjective}

\noindent
{\it 满}
见 \ref{formal-spaces-definition-surjective}

\noindent
{\it 满}
见 \ref{stacks-properties-definition-surjective}

\noindent
{\it 与 $M, a, b$ 关联的符号}
见 \ref{chow-definition-symbol-M}

\noindent
{\it 符号幂}
见 \ref{algebra-definition-symbolic-power}

\noindent
{\it 符号}
见 \ref{chow-definition-determinant}

\noindent
{\it 对称幺半范畴}
见 \ref{categories-definition-symmetric-monoidal-category}

\noindent
{\it 在 $x \in X$ 处 syntomic}
见 \ref{morphisms-definition-syntomic}

\noindent
{\it 在 $x$ 处 syntomic}
见 \ref{spaces-morphisms-definition-syntomic}

\noindent
{\it $T$ 的 syntomic 覆盖}
见 \ref{topologies-definition-syntomic-covering}

\noindent
{\it $X$ 的 syntomic 覆盖}
见 \ref{spaces-topologies-definition-syntomic-covering}

\noindent
{\it 相对维数 $d$ 的 syntomic 态射}
见 \ref{morphisms-definition-syntomic-relative-dimension}

\noindent
{\it syntomic 层}
见 \ref{stacks-sheaves-definition-sheaves}

\noindent
{\it syntomic}
见 \ref{algebra-definition-lci}

\noindent
{\it syntomic}
见 \ref{morphisms-definition-syntomic}

\noindent
{\it syntomic}
见 \ref{spaces-morphisms-definition-syntomic}

\noindent
{\it $(X_i, f_{i'i})$ 上的层系统 $(\mathcal{F}_i, \varphi_{i'i})$}
见 \ref{etale-cohomology-definition-inverse-system-sheaves}

\noindent
{\it 以 $I$ 为指标的 $R$-模系统 $(M_i, \mu_{ij})$}
见 \ref{algebra-definition-directed-system}

\noindent
{\it $R$ 的一组参数系}
见 \ref{algebra-definition-regular-local}

\noindent
{\it 环系}
见 \ref{exercises-definition-directed-poset}

\noindent
{\it $\mathcal{C}$ 中以 $I$ 为指标的系统}
见 \ref{categories-definition-system-over-poset}

\noindent
{\it $\mathfrak m$ 的驯惯性群}
见 \ref{more-algebra-definition-wild-inertia}

\noindent
{\it 驯符号}
见 \ref{chow-definition-tame-symbol}

\noindent
{\it 相对于 $A$ 驯分歧}
见 \ref{more-algebra-definition-types-of-extensions}

\noindent
{\it $\mathcal{F}$ 的切空间 $T \mathcal{F}$}
见 \ref{formal-defos-definition-tangent-space}

\noindent
{\it $F$ 的切空间 $TF$}
见 \ref{formal-defos-definition-tangent-space-over-R}

\noindent
{\it $X$ 在 $x$ 处相对于 $S$ 的切空间}
见 \ref{varieties-definition-tangent-space}

\noindent
{\it $X$ 在 $S$ 上的切空间}
见 \ref{exercises-definition-tangent-space}

\noindent
{\it 切向量}
见 \ref{varieties-definition-tangent-space}

\noindent
{\it 切向量}
见 \ref{exercises-definition-tangent-space}

\noindent
{\it 自明等价}
见 \ref{sites-definition-combinatorial-tautological}

\noindent
{\it taut}
见 \ref{formal-spaces-definition-taut}

\noindent
{\it 张量幂}
见 \ref{modules-definition-powers}

\noindent
{\it 张量积微分分次代数}
见 \ref{dga-definition-tensor-product}

\noindent
{\it 张量积}
见 \ref{formal-spaces-definition-toplogy-tensor-product}

\noindent
{\it $\mathcal{A}$ 的复形的逐项分裂正合列}
见 \ref{derived-definition-split-ses}

\noindent
{\it 逐项分裂单射 $\alpha : A^\bullet \to B^\bullet$}
见 \ref{derived-definition-termwise-split-map}

\noindent
{\it 逐项分裂满射 $\beta : B^\bullet \to C^\bullet$}
见 \ref{derived-definition-termwise-split-map}

\noindent
{\it $f : X \to Y$ 在 $y$ 处的纤维几何既约}
见 \ref{spaces-more-morphisms-definition-geometrically-reduced-fibre}

\noindent
{\it $f$ 在 $y$ 上的纤维局部 Noether}
见 \ref{spaces-divisors-definition-locally-Noetherian-fibre}

\noindent
{\it $X$ 在 $z$ 上的纤维在 $x$ 处相对于 $Y$ 在 $z$ 上的纤维平坦}
见 \ref{spaces-more-morphisms-definition-module-flat-on-fibre}

\noindent
{\it $X$ 在 $z$ 上的纤维相对于 $Y$ 在 $z$ 上的纤维平坦}
见 \ref{spaces-more-morphisms-definition-module-flat-on-fibre}

\noindent
{\it $f$ 的纤维局部 Noether}
见 \ref{spaces-divisors-definition-locally-Noetherian-fibre}

\noindent
{\it $S$ 上从 $X$ 到 $Y$ 的相对等价的 Fourier--Mukai 核}
见 \ref{equiv-definition-relative-equivalence-kernel}

\noindent
{\it $X$ 上的函数就是 $U$ 上的 $R$-不变函数}
见 \ref{groupoids-quotients-definition-functions}

\noindent
{\it $f$ 的 Gysin 映射存在}
见 \ref{chow-definition-lci-gysin}

\noindent
{\it 相对维数}
见 \ref{stacks-geometry-definition-relative-dimension}

\noindent
{\it $\mathcal{F}$ 到其在 $z$ 上纤维的限制在 $x$ 处相对于 $Y$ 在 $z$ 上的纤维平坦}
见 \ref{spaces-more-morphisms-definition-module-flat-on-fibre}

\noindent
{\it $\mathcal{Z}$ 上的增厚}
见 \ref{stacks-more-morphisms-definition-thickening}

\noindent
{\it $B$ 上的加厚}
见 \ref{spaces-more-morphisms-definition-thickening}

\noindent
{\it $S$ 上的加厚}
见 \ref{more-morphisms-definition-thickening}

\noindent
{\it 加厚}
见 \ref{more-morphisms-definition-thickening}

\noindent
{\it 加厚}
见 \ref{spaces-more-morphisms-definition-thickening}

\noindent
{\it 加厚}
见 \ref{stacks-more-morphisms-definition-thickening}

\noindent
{\it $T$ 的拓扑亏格}
见 \ref{models-definition-top-genus}

\noindent
{\it 拓扑群}
见 \ref{topology-definition-topological-group}

\noindent
{\it 拓扑模}
见 \ref{topology-definition-topological-module}

\noindent
{\it 拓扑模}
见 \ref{more-algebra-definition-topological-ring}

\noindent
{\it 拓扑环}
见 \ref{topology-definition-topological-ring}

\noindent
{\it 拓扑环}
见 \ref{more-algebra-definition-topological-ring}

\noindent
{\it 拓扑空间}
见 \ref{spaces-properties-definition-topological-space}

\noindent
{\it 拓扑空间}
见 \ref{stacks-properties-definition-topological-space}

\noindent
{\it 拓扑幂零的}
见 \ref{formal-spaces-definition-weakly-admissible}

\noindent
{\it 在……上拓扑有限型}
见 \ref{formal-spaces-definition-topologically-finite-type}

\noindent
{\it 与 $\mathcal{C}$ 关联的拓扑}
见 \ref{sites-definition-topology-associated-site}

\noindent
{\it $\mathcal{C}$ 上的拓扑}
见 \ref{sites-definition-topology}

\noindent
{\it 托扑斯}
见 \ref{sites-definition-topos}

\noindent
{\it Tor 维数 $\leq d$}
见 \ref{more-algebra-definition-tor-amplitude}

\noindent
{\it Tor 维数 $\leq d$}
见 \ref{cohomology-definition-tor-amplitude}

\noindent
{\it Tor 维数 $\leq d$}
见 \ref{sites-cohomology-definition-tor-amplitude}

\noindent
{\it 在 $B$ 上 Tor 无关}
见 \ref{spaces-perfect-definition-tor-independent}

\noindent
{\it 在 $R$ 上 Tor 无关}
见 \ref{more-algebra-definition-tor-independent}

\noindent
{\it 在 $S$ 上 Tor 独立}
见 \ref{perfect-definition-tor-independent}

\noindent
{\it Tor 振幅位于 $[a, b]$}
见 \ref{more-algebra-definition-tor-amplitude}

\noindent
{\it Tor 振幅位于 $[a, b]$}
见 \ref{cohomology-definition-tor-amplitude}

\noindent
{\it Tor 振幅位于 $[a, b]$}
见 \ref{sites-cohomology-definition-tor-amplitude}

\noindent
{\it 无挠}
见 \ref{more-algebra-definition-torsion}

\noindent
{\it 无挠}
见 \ref{divisors-definition-torsion}

\noindent
{\it 挠模}
见 \ref{more-algebra-definition-torsion}

\noindent
{\it 挠的}
见 \ref{more-algebra-definition-torsion}

\noindent
{\it 挠的}
见 \ref{divisors-definition-torsion}

\noindent
{\it 挠的}
见 \ref{trace-definition-torsion-l-adic-sheaf}

\noindent
{\it 扭子}
见 \ref{cohomology-definition-torsor}

\noindent
{\it 扭子}
见 \ref{sites-cohomology-definition-torsor}

\noindent
{\it Tor}
见 \ref{cohomology-definition-tor}

\noindent
{\it Tor}
见 \ref{sites-cohomology-definition-tor}

\noindent
{\it $X$ 上 $\mathcal{E}$ 的全陈类}
见 \ref{chow-definition-chern-classes}

\noindent
{\it $\mathcal{E}$ 的总 Chern 类}
见 \ref{spaces-chow-definition-chern-classes}

\noindent
{\it 全陈类}
见 \ref{chow-definition-chern-classes-final}

\noindent
{\it $G$ 的全左导出函子}
见 \ref{trace-definition-derived-functor}

\noindent
{\it $F$ 的全右导出函子}
见 \ref{trace-definition-derived-functor}

\noindent
{\it 全无环}
见 \ref{sites-cohomology-definition-limp}

\noindent
{\it 全不连通}
见 \ref{topology-definition-totally-disconnected}

\noindent
{\it 相对于 $A$ 全分歧}
见 \ref{more-algebra-definition-types-of-extensions}

\noindent
{\it 塔}
见 \ref{fields-definition-tower}

\noindent
{\it 迹元素}
见 \ref{discriminant-definition-trace-element}

\noindent
{\it 迹配对}
见 \ref{fields-definition-trace-pairing}

\noindent
{\it 迹}
见 \ref{fields-definition-trace-norm}

\noindent
{\it 迹}
见 \ref{etale-cohomology-definition-trace-map}

\noindent
{\it 迹}
见 \ref{trace-definition-trace}

\noindent
{\it 超越基}
见 \ref{fields-definition-transcendence}

\noindent
{\it $x/f(x)$ 的超越次数}
见 \ref{spaces-morphisms-definition-dimension-fibre}

\noindent
{\it 超越次数}
见 \ref{fields-definition-transcendence-degree}

\noindent
{\it 过渡映射}
见 \ref{categories-definition-system-over-poset}

\noindent
{\it 与 $0 \to K \to L \to M \to 0$ 关联的三角}
见 \ref{dga-definition-distinguished-triangle}

\noindent
{\it 与逐项分裂复形列关联的三角}
见 \ref{derived-definition-split-ses}

\noindent
{\it 三角}
见 \ref{derived-definition-triangle}

\noindent
{\it 导出范畴中拟凝聚对象所成的三角范畴}
见 \ref{stacks-sheaves-definition-QC}

\noindent
{\it 三角范畴}
见 \ref{derived-definition-triangulated-category}

\noindent
{\it 三角函子}
见 \ref{derived-definition-exact-functor-triangulated-categories}

\noindent
{\it 三角子范畴}
见 \ref{derived-definition-triangulated-subcategory}

\noindent
{\it 平凡 $\mathcal{G}$-扭子}
见 \ref{cohomology-definition-torsor}

\noindent
{\it 平凡 $\mathcal{G}$-扭子}
见 \ref{sites-cohomology-definition-torsor}

\noindent
{\it 平凡下降数据}
见 \ref{stacks-definition-effective-descent-datum}

\noindent
{\it 平凡下降数据}
见 \ref{descent-definition-descent-datum-effective-quasi-coherent}

\noindent
{\it 平凡下降数据}
见 \ref{descent-definition-effective}

\noindent
{\it 平凡下降数据}
见 \ref{spaces-descent-definition-descent-datum-effective-quasi-coherent}

\noindent
{\it 平凡下降数据}
见 \ref{spaces-descent-definition-effective}

\noindent
{\it 平凡 Kan 纤维化}
见 \ref{simplicial-definition-trivial-kan}

\noindent
{\it 平凡}
见 \ref{more-algebra-definition-invertible}

\noindent
{\it 平凡}
见 \ref{modules-definition-invertible}

\noindent
{\it 平凡}
见 \ref{groupoids-definition-pseudo-torsor}

\noindent
{\it 平凡}
见 \ref{spaces-groupoids-definition-pseudo-torsor}

\noindent
{\it 平凡}
见 \ref{exercises-definition-invertible-module}

\noindent
{\it $\text{Proj}(S)$ 的结构层扭转}
见 \ref{constructions-definition-twist}

\noindent
{\it 结构层的扭曲}
见 \ref{constructions-definition-projective-bundle}

\noindent
{\it 双侧可容许}
见 \ref{derived-definition-admissible}

\noindent
{\it 代数结构类型}
见 \ref{sheaves-definition-algebraic-structure}

\noindent
{\it UFD}
见 \ref{algebra-definition-UFD}

\noindent
{\it $\mathcal{F}$ 的底层集合预层}
见 \ref{sheaves-definition-underlying-presheaf-sets}

\noindent
{\it 在 $x$ 处单分支}
见 \ref{properties-definition-unibranch}

\noindent
{\it 在 $x$ 处单分支}
见 \ref{decent-spaces-definition-unibranch}

\noindent
{\it 单分支}
见 \ref{more-algebra-definition-unibranch}

\noindent
{\it 单分支}
见 \ref{properties-definition-unibranch}

\noindent
{\it 单分支}
见 \ref{decent-spaces-definition-unibranch}

\noindent
{\it $\mathcal{C}$ 中的一致范畴模空间}
见 \ref{stacks-more-morphisms-definition-categorical-quotient}

\noindent
{\it 一致范畴模空间}
见 \ref{stacks-more-morphisms-definition-categorical-quotient}

\noindent
{\it 一致范畴商}
见 \ref{groupoids-quotients-definition-universal-categorical}

\noindent
{\it 一致化元}
见 \ref{algebra-definition-uniformizer}

\noindent
{\it 一致地}
见 \ref{groupoids-quotients-definition-topological}

\noindent
{\it 唯一分解整环}
见 \ref{algebra-definition-UFD}

\noindent
{\it uniqueness part of the valuative criterion}
见 \ref{stacks-morphisms-definition-uniqueness}

\noindent
{\it 泛 $\delta$-函子}
见 \ref{homology-definition-universal-delta-functor}

\noindent
{\it 泛 $\varphi$-导子}
见 \ref{modules-definition-module-differentials}

\noindent
{\it 泛 $\varphi$-导子}
见 \ref{sites-modules-definition-module-differentials}

\noindent
{\it 泛 $S$-导子}
见 \ref{morphisms-definition-sheaf-differentials}

\noindent
{\it 泛 $Y$-导子}
见 \ref{sites-modules-definition-sheaf-differentials}

\noindent
{\it 泛 $Y$-导子}
见 \ref{spaces-more-morphisms-definition-sheaf-differentials}

\noindent
{\it 泛范畴商}
见 \ref{groupoids-quotients-definition-universal-categorical}

\noindent
{\it 泛有效满态射}
见 \ref{sites-definition-universal-effective-epimorphisms}

\noindent
{\it 泛一阶加厚}
见 \ref{algebra-definition-universal-thickening}

\noindent
{\it 泛一阶加厚}
见 \ref{more-morphisms-definition-universal-thickening}

\noindent
{\it 泛一阶加厚}
见 \ref{spaces-more-morphisms-definition-universal-thickening}

\noindent
{\it $\mathcal{F}$ 的泛扁平化存在}
见 \ref{flat-definition-flattening}

\noindent
{\it $\mathcal{F}$ 的泛扁平化存在}
见 \ref{spaces-flat-definition-flattening}

\noindent
{\it $X$ 的泛扁平化存在}
见 \ref{flat-definition-flattening}

\noindent
{\it $X$ 的泛扁平化存在}
见 \ref{spaces-flat-definition-flattening}

\noindent
{\it 普遍同胚}
见 \ref{morphisms-definition-universal-homeomorphism}

\noindent
{\it 普遍同胚}
见 \ref{spaces-morphisms-definition-universal-homeomorphism}

\noindent
{\it 普遍同胚}
见 \ref{stacks-morphisms-definition-universal-homeomorphism}

\noindent
{\it 普遍 $S$-纯}
见 \ref{flat-definition-pure}

\noindent
{\it 普遍 $Y$-纯}
见 \ref{spaces-flat-definition-pure}

\noindent
{\it 普遍链状}
见 \ref{algebra-definition-universally-catenary}

\noindent
{\it 普遍链状}
见 \ref{morphisms-definition-universally-catenary}

\noindent
{\it 普遍链状}
见 \ref{decent-spaces-definition-universally-catenary}

\noindent
{\it 普遍闭}
见 \ref{topology-definition-proper-map}

\noindent
{\it 普遍闭}
见 \ref{schemes-definition-universally-closed}

\noindent
{\it 普遍闭}
见 \ref{spaces-morphisms-definition-closed}

\noindent
{\it 普遍闭}
见 \ref{stacks-morphisms-definition-closed}

\noindent
{\it 普遍正合}
见 \ref{algebra-definition-universally-injective}

\noindent
{\it 普遍单射}
见 \ref{algebra-definition-universally-injective}

\noindent
{\it 普遍单射}
见 \ref{morphisms-definition-universally-injective}

\noindent
{\it 普遍单射}
见 \ref{descent-definition-universally-injective}

\noindent
{\it 普遍单射}
见 \ref{spaces-morphisms-definition-universally-injective}

\noindent
{\it 普遍单射}
见 \ref{stacks-morphisms-definition-universally-injective}

\noindent
{\it 泛 Japanese 概形}
见 \ref{algebra-definition-nagata}

\noindent
{\it 泛 Japanese 概形}
见 \ref{properties-definition-nagata}

\noindent
{\it 相对于 $K$ 普遍局部无环}
见 \ref{etale-cohomology-definition-locally-acyclic}

\noindent
{\it 普遍局部无环}
见 \ref{etale-cohomology-definition-locally-acyclic}

\noindent
{\it 普遍开}
见 \ref{morphisms-definition-open}

\noindent
{\it 普遍开}
见 \ref{spaces-morphisms-definition-open}

\noindent
{\it 普遍开}
见 \ref{stacks-morphisms-definition-open}

\noindent
{\it 在 $y$ 上方普遍纯}
见 \ref{spaces-flat-definition-pure}

\noindent
{\it 沿 $X_s$ 普遍纯}
见 \ref{flat-definition-pure}

\noindent
{\it 相对于 $S$ 普遍纯}
见 \ref{flat-definition-pure}

\noindent
{\it 相对于 $Y$ 普遍纯}
见 \ref{spaces-flat-definition-pure}

\noindent
{\it 普遍商映射}
见 \ref{morphisms-definition-submersive}

\noindent
{\it 普遍商映射}
见 \ref{spaces-morphisms-definition-submersive}

\noindent
{\it 普遍商映射}
见 \ref{stacks-morphisms-definition-submersive}

\noindent
{\it 泛地}
见 \ref{groupoids-quotients-definition-topological}

\noindent
{\it 无阻碍}
见 \ref{formal-defos-definition-cofibered-groupoid-projection-smooth}

\noindent
{\it 在 $\mathfrak q$ 处无分歧}
见 \ref{algebra-definition-unramified}

\noindent
{\it 在 $x \in X$ 处非分歧}
见 \ref{morphisms-definition-unramified}

\noindent
{\it 在 $x$ 处非分歧}
见 \ref{etale-definition-unramified-schemes}

\noindent
{\it 在 $x$ 处非分歧}
见 \ref{spaces-morphisms-definition-unramified}

\noindent
{\it 取值于 $\Lambda$ 的 $\text{GL}_2(\mathbf{A})$ 上非分歧尖点形式}
见 \ref{trace-definition-unramified}

\noindent
{\it 局部环的非分歧同态}
见 \ref{etale-definition-unramified-rings}

\noindent
{\it 相对于 $A$ 无分歧}
见 \ref{more-algebra-definition-types-of-extensions}

\noindent
{\it 无分歧}
见 \ref{algebra-definition-unramified}

\noindent
{\it 无分歧}
见 \ref{morphisms-definition-unramified}

\noindent
{\it 无分歧}
见 \ref{etale-definition-unramified-schemes}

\noindent
{\it 无分歧}
见 \ref{spaces-morphisms-definition-unramified}

\noindent
{\it 无分歧}
见 \ref{stacks-morphisms-definition-unramified}

\noindent
{\it $T$ 的 V 覆盖}
见 \ref{topologies-definition-V-covering}

\noindent
{\it 赋值环}
见 \ref{algebra-definition-valuation-ring}

\noindent
{\it 赋值}
见 \ref{algebra-definition-value-group}

\noindent
{\it 值群}
见 \ref{algebra-definition-value-group}

\noindent
{\it $LF$ 在 $X$ 处的值}
见 \ref{derived-definition-right-derived-functor-defined}

\noindent
{\it $RF$ 在 $X$ 处的值}
见 \ref{derived-definition-right-derived-functor-defined}

\noindent
{\it 本质常值}
见 \ref{categories-definition-essentially-constant-diagram}

\noindent
{\it 本质常值}
见 \ref{categories-definition-essentially-constant-diagram}

\noindent
{\it 簇}
见 \ref{varieties-definition-variety}

\noindent
{\it 簇}
见 \ref{etale-cohomology-definition-variety}

\noindent
{\it 向量丛 $\pi : V \to S$（在 $S$ 上）}
见 \ref{constructions-definition-abstract-vector-bundle}

\noindent
{\it 与 $\mathcal{E}$ 关联的向量丛}
见 \ref{constructions-definition-vector-bundle}

\noindent
{\it $\mathcal{X}$ 在 $x_0$ 处的 versal 环}
见 \ref{stacks-geometry-definition-versal-ring-at-x}

\noindent
{\it versal}
见 \ref{formal-defos-definition-versal}

\noindent
{\it versal}
见 \ref{artin-definition-versal-formal-object}

\noindent
{\it versal}
见 \ref{artin-definition-versal}

\noindent
{\it 竖直}
见 \ref{categories-definition-2-category}

\noindent
{\it 在 $X/S$ 上甚丰沛}
见 \ref{morphisms-definition-very-ample}

\noindent
{\it very reasonable}
见 \ref{decent-spaces-definition-very-reasonable}

\noindent
{\it very reasonable}
见 \ref{decent-spaces-definition-relative-conditions}

\noindent
{\it 视作 $S'$ 上的代数空间}
见 \ref{spaces-definition-base-change}

\noindent
{\it 视为 $S'$ 上的代数叠}
见 \ref{algebraic-definition-viewed-as}

\noindent
{\it w-可缩的}
见 \ref{proetale-definition-w-contractible}

\noindent
{\it w-局部的}
见 \ref{proetale-definition-w-local}

\noindent
{\it w-局部的}
见 \ref{proetale-definition-w-local}

\noindent
{\it 弱 $R$-轨道}
见 \ref{groupoids-quotients-definition-geometric-orbits}

\noindent
{\it 弱维数 $\leq d$}
见 \ref{more-algebra-definition-weak-dimension}

\noindent
{\it 弱函子}
见 \ref{categories-definition-functor-into-2-category}

\noindent
{\it 弱生成元}
见 \ref{derived-definition-generators}

\noindent
{\it 弱定义理想}
见 \ref{formal-spaces-definition-weakly-admissible}

\noindent
{\it $X$ 在 $Y$ 中的弱正规化}
见 \ref{morphisms-definition-relative-seminormalization}

\noindent
{\it 弱正规化}
见 \ref{morphisms-definition-weak-normalization}

\noindent
{\it 弱轨道}
见 \ref{groupoids-quotients-definition-geometric-orbits}

\noindent
{\it 弱 Serre 子范畴}
见 \ref{homology-definition-serre-subcategory}

\noindent
{\it $A \subset B$ 的弱解}
见 \ref{more-algebra-definition-solution}

\noindent
{\it 弱于典范拓扑}
见 \ref{sites-definition-weaker-than-canonical}

\noindent
{\it 更弱}
见 \ref{sites-definition-finer}

\noindent
{\it 弱 $R$-等价}
见 \ref{groupoids-quotients-definition-geometric-orbits}

\noindent
{\it 弱 \'etale 的}
见 \ref{more-algebra-definition-weakly-etale}

\noindent
{\it 弱 \'etale 的}
见 \ref{more-morphisms-definition-weakly-etale}

\noindent
{\it 弱 adic}
见 \ref{formal-spaces-definition-weakly-adic}

\noindent
{\it 弱 adic}
见 \ref{formal-spaces-definition-types-affine-formal-algebraic-space}

\noindent
{\it 弱可容许}
见 \ref{formal-spaces-definition-weakly-admissible}

\noindent
{\it $X$ 的弱伴随点}
见 \ref{divisors-definition-weakly-associated}

\noindent
{\it $X$ 的弱伴随点}
见 \ref{spaces-divisors-definition-weakly-associated}

\noindent
{\it 弱相伴}
见 \ref{algebra-definition-weakly-associated}

\noindent
{\it 弱相伴}
见 \ref{divisors-definition-weakly-associated}

\noindent
{\it 弱相伴}
见 \ref{spaces-divisors-definition-weakly-associated}

\noindent
{\it 弱可缩的}
见 \ref{sites-definition-w-contractible}

\noindent
{\it 弱收敛于 $H(K)$}
见 \ref{homology-definition-filtered-differential-ss-converges}

\noindent
{\it 弱收敛于 $H^*(K^\bullet)$}
见 \ref{homology-definition-filtered-complex-ss-converges}

\noindent
{\it 弱收敛于 $H^n(\text{Tot}(K^{\bullet, \bullet}))$}
见 \ref{homology-definition-ss-double-complex-converge}

\noindent
{\it 弱收敛于 $H^n(\text{Tot}(K^{\bullet, \bullet}))$}
见 \ref{homology-definition-ss-double-complex-converge}

\noindent
{\it 弱正规}
见 \ref{morphisms-definition-weakly-normal}

\noindent
{\it 弱预-adic}
见 \ref{formal-spaces-definition-weakly-adic}

\noindent
{\it 弱预可容许}
见 \ref{formal-spaces-definition-weakly-admissible}

\noindent
{\it 弱无分歧}
见 \ref{more-algebra-definition-extension-discrete-valuation-rings}

\noindent
{\it 弱无分歧}
见 \ref{more-algebra-definition-extension-valuation-rings}

\noindent
{\it 赋权}
见 \ref{more-morphisms-definition-weighting}

\noindent
{\it Weil 上同调理论}
见 \ref{weil-definition-weil-cohomology-theory}

\noindent
{\it 与有效 Cartier 除子 $D \subset X$ 关联的 Weil 除子 $[D]$}
见 \ref{exercises-definition-divisor}

\noindent
{\it 与 $\mathcal{L}$ 关联的 Weil 除子}
见 \ref{chow-definition-divisor-invertible-sheaf}

\noindent
{\it 与 $\mathcal{L}$ 关联的 Weil 除子}
见 \ref{spaces-chow-definition-divisor-invertible-sheaf}

\noindent
{\it 与 $s$ 关联的 Weil 除子}
见 \ref{divisors-definition-divisor-invertible-sheaf}

\noindent
{\it 与 $s$ 关联的 Weil 除子}
见 \ref{chow-definition-divisor-invertible-sheaf}

\noindent
{\it 与 $s$ 关联的 Weil 除子}
见 \ref{spaces-over-fields-definition-divisor-invertible-sheaf}

\noindent
{\it 与 $s$ 关联的 Weil 除子}
见 \ref{spaces-chow-definition-divisor-invertible-sheaf}

\noindent
{\it 相伴的 Weil 除子}
见 \ref{exercises-definition-divisor}

\noindent
{\it 与有理函数 $f \in K(X)^\ast$ 关联的 Weil 除子}
见 \ref{exercises-definition-divisor}

\noindent
{\it 与 $\mathcal{L}$ 关联的 Weil 除子类}
见 \ref{divisors-definition-divisor-invertible-sheaf}

\noindent
{\it 与 $\mathcal{L}$ 关联的 Weil 除子类}
见 \ref{spaces-over-fields-definition-divisor-invertible-sheaf}

\noindent
{\it Weil 除子类群}
见 \ref{divisors-definition-class-group}

\noindent
{\it Weil 除子类群}
见 \ref{spaces-over-fields-definition-class-group}

\noindent
{\it Weil 除子}
见 \ref{divisors-definition-Weil-divisor}

\noindent
{\it Weil 除子}
见 \ref{spaces-over-fields-definition-Weil-divisor}

\noindent
{\it Weil 除子}
见 \ref{exercises-definition-divisor}

\noindent
{\it 近乎仿射的}
见 \ref{stacks-more-morphisms-definition-well-nigh-affine}

\noindent
{\it 将半可表示对象对应到预层的函子}
见 \ref{hypercovering-definition-SR-F}

\noindent
{\it $\mathfrak m$ 的野惯性群}
见 \ref{more-algebra-definition-wild-inertia}

\noindent
{\it Yoneda 扩张}
见 \ref{derived-definition-yoneda-extension}

\noindent
{\it $T$ 的 Zariski 覆盖}
见 \ref{topologies-definition-zariski-covering}

\noindent
{\it $X$ 的 Zariski 覆盖}
见 \ref{spaces-topologies-definition-zariski-covering}

\noindent
{\it Zariski 覆盖}
见 \ref{spaces-definition-Zariski-open-covering}

\noindent
{\it 在 $S$ 上 Zariski 局部拟分离}
见 \ref{spaces-definition-separated}

\noindent
{\it Zariski 局部拟分离}
见 \ref{spaces-properties-definition-separated}

\noindent
{\it Zariski 局部拟分离}
见 \ref{spaces-properties-definition-separated}

\noindent
{\it Zariski 对}
见 \ref{more-algebra-definition-zariski-pair}

\noindent
{\it Zariski 层}
见 \ref{stacks-sheaves-definition-sheaves}

\noindent
{\it Zariski 拓扑斯}
见 \ref{etale-cohomology-definition-etale-topos}

\noindent
{\it Zariski、\'etale、光滑、syntomic 或 fppf 覆盖}
见 \ref{crystalline-definition-big-crystalline-site}

\noindent
{\it 扎里斯基}
见 \ref{algebra-definition-Zariski-topology}

\noindent
{\it 零对象}
见 \ref{homology-definition-zero-object}

\noindent
{\it 零概形}
见 \ref{divisors-definition-zero-scheme-s}

\noindent
{\it 零概形}
见 \ref{spaces-divisors-definition-zero-scheme-s}

\noindent
{\it $\mathcal{A}$ 的第零 $K$-群}
见 \ref{homology-definition-K-zero}

\noindent
{\it $\mathcal{D}$ 的第零 $K$-群}
见 \ref{derived-definition-K-zero}

\noindent
{\it 第零 {\v C}ech 上同调群}
见 \ref{etale-cohomology-definition-0-cech}

\noindent
{\it {\v C}ech 上同调群}
见 \ref{cohomology-definition-cech-complex}

\noindent
{\it {\v C}ech 上同调群}
见 \ref{sites-cohomology-definition-cech-complex}

\noindent
{\it {\v C}ech 上同调群}
见 \ref{etale-cohomology-definition-cech-complex}

\noindent
{\it {\v C}ech 复形}
见 \ref{cohomology-definition-cech-complex}

\noindent
{\it {\v C}ech 复形}
见 \ref{sites-cohomology-definition-cech-complex}

\noindent
{\it {\v C}ech 复形}
见 \ref{etale-cohomology-definition-cech-complex}

\end{multicols}

\begin{multicols}{2}[\section{按章节列出的定义}\label{index-section-per-chapter}]

\medskip\noindent
{\bf 引言}

\medskip

\medskip\noindent
{\bf 约定}

\medskip

\medskip\noindent
{\bf 集合论}

\medskip

\medskip\noindent
{\bf 范畴}

\medskip

\noindent
见 \ref{categories-definition-category}：
{\it 范畴}


\noindent
见 \ref{categories-definition-isomorphism}：
{\it 同构}


\noindent
见 \ref{categories-definition-groupoid}：
{\it 群胚}


\noindent
见 \ref{categories-definition-functor}：
{\it 函子}


\noindent
见 \ref{categories-definition-faithful}：
{\it 忠实},
{\it 全忠实},
{\it 本质满射}


\noindent
见 \ref{categories-definition-subcategory}：
{\it 子范畴},
{\it 满子范畴},
{\it 严格满}


\noindent
见 \ref{categories-definition-transformation-functors}：
{\it 自然变换},
{\it 函子态射}


\noindent
见 \ref{categories-definition-equivalence-categories}：
{\it 范畴等价},
{\it 拟逆}


\noindent
见 \ref{categories-definition-product-category}：
{\it 积范畴}


\noindent
见 \ref{categories-definition-opposite}：
{\it 对偶范畴}


\noindent
见 \ref{categories-definition-contravariant}：
{\it 反变}


\noindent
见 \ref{categories-definition-presheaf}：
{\it $\mathcal{C}$ 上的集合预层},
{\it 预层}


\noindent
见 \ref{categories-definition-representable-functor}：
{\it 可表示}


\noindent
见 \ref{categories-definition-products}：
{\it 乘积}


\noindent
见 \ref{categories-definition-has-products-of-pairs}：
{\it 具有对象对的乘积}


\noindent
见 \ref{categories-definition-coproducts}：
{\it 余积},
{\it 并合和}


\noindent
见 \ref{categories-definition-has-coproducts-of-pairs}：
{\it 具有对象对的余积}


\noindent
见 \ref{categories-definition-fibre-products}：
{\it 纤维积}


\noindent
见 \ref{categories-definition-cartesian}：
{\it 笛卡尔的}


\noindent
见 \ref{categories-definition-has-fibre-products}：
{\it 具有纤维积}


\noindent
见 \ref{categories-definition-representable-morphism}：
{\it 可表示}


\noindent
见 \ref{categories-definition-representable-map-presheaves}：
{\it 可表示},
{\it $F$ 在 $G$ 上相对可表}


\noindent
见 \ref{categories-definition-pushouts}：
{\it 推出}


\noindent
见 \ref{categories-definition-cocartesian}：
{\it 余笛卡尔}


\noindent
见 \ref{categories-definition-equalizers}：
{\it 等化子}


\noindent
见 \ref{categories-definition-coequalizers}：
{\it 余等化子}


\noindent
见 \ref{categories-definition-initial-final}：
{\it 始的},
{\it 终对象}


\noindent
见 \ref{categories-definition-mono-epi}：
{\it 单态射},
{\it 满态射}


\noindent
见 \ref{categories-definition-limit}：
{\it 极限}


\noindent
见 \ref{categories-definition-colimit}：
{\it 余极限}


\noindent
见 \ref{categories-definition-product}：
{\it 乘积}


\noindent
见 \ref{categories-definition-coproduct}：
{\it 余积}


\noindent
见 \ref{categories-definition-category-connected}：
{\it 连通}


\noindent
见 \ref{categories-definition-cofinal}：
{\it $\mathcal{I}$ 在 $\mathcal{J}$ 中共尾},
{\it 共尾}


\noindent
见 \ref{categories-definition-initial}：
{\it $\mathcal{I}$ 在 $\mathcal{J}$ 中是始的},
{\it 始的}


\noindent
见 \ref{categories-definition-directed}：
{\it 有向},
{\it 滤过},
{\it 有向},
{\it 滤过}


\noindent
见 \ref{categories-definition-codirected}：
{\it 反向},
{\it 余滤过},
{\it 反向},
{\it 余滤过}


\noindent
见 \ref{categories-definition-directed-set}：
{\it 预序},
{\it 预序集},
{\it 有向集},
{\it 偏序},
{\it 偏序集},
{\it 有向偏序集}


\noindent
见 \ref{categories-definition-system-over-poset}：
{\it $\mathcal{C}$ 中以 $I$ 为指标的系统},
{\it $\mathcal{C}$ 中以 $I$ 为指标的归纳系统},
{\it $\mathcal{C}$ 中以 $I$ 为指标的逆系统},
{\it $\mathcal{C}$ 中以 $I$ 为指标的投射系统},
{\it 过渡映射}


\noindent
见 \ref{categories-definition-directed-system}：
{\it 有向系统},
{\it 有向逆系}


\noindent
见 \ref{categories-definition-essentially-constant-diagram}：
{\it 值},
{\it 本质常值},
{\it 值},
{\it 本质常值}


\noindent
见 \ref{categories-definition-essentially-constant-system}：
{\it 本质常值系},
{\it 本质常值逆系}


\noindent
见 \ref{categories-definition-exact}：
{\it 左正合},
{\it 右正合},
{\it 正合}


\noindent
见 \ref{categories-definition-adjoint}：
{\it 左伴随},
{\it 右伴随}


\noindent
见 \ref{categories-definition-compact-object}：
{\it 范畴紧致}


\noindent
见 \ref{categories-definition-multiplicative-system}：
{\it 左乘法系},
{\it 右乘法系},
{\it 乘法系}


\noindent
见 \ref{categories-definition-left-localization-as-fraction}：
{\it $s^{-1} f$}


\noindent
见 \ref{categories-definition-right-localization-as-fraction}：
{\it $f s^{-1}$}


\noindent
见 \ref{categories-definition-saturated-multiplicative-system}：
{\it 饱和}


\noindent
见 \ref{categories-definition-horizontal-composition}：
{\it 水平}


\noindent
见 \ref{categories-definition-2-category}：
{\it $2$-范畴},
{\it $1$-态射},
{\it $2$-态射},
{\it 竖直},
{\it 复合},
{\it 水平}


\noindent
见 \ref{categories-definition-sub-2-category}：
{\it 子 $2$-范畴}


\noindent
见 \ref{categories-definition-equivalence}：
{\it 等价}


\noindent
见 \ref{categories-definition-functor-into-2-category}：
{\it 函子},
{\it 弱函子},
{\it 伪函子}


\noindent
见 \ref{categories-definition-2-1-category}：
{\it $(2, 1)$-范畴}


\noindent
见 \ref{categories-definition-final-object-2-category}：
{\it 终对象}


\noindent
见 \ref{categories-definition-2-fibre-products}：
{\it $f$ 与 $g$ 的 2-纤维积}


\noindent
见 \ref{categories-definition-categories-over-C}：
{\it $\mathcal{C}$ 上的范畴所成的 $2$-范畴}


\noindent
见 \ref{categories-definition-fibre-category}：
{\it 纤维范畴},
{\it 提升},
{\it $x$ 位于 $U$ 之上},
{\it 提升},
{\it $\phi$ 位于 $f$ 之上}


\noindent
见 \ref{categories-definition-cartesian-over-C}：
{\it 强笛卡尔态射},
{\it 强 $\mathcal{C}$-笛卡尔态射}


\noindent
见 \ref{categories-definition-fibred-category}：
{\it $\mathcal{C}$ 上的纤维化范畴}


\noindent
见 \ref{categories-definition-pullback-functor-fibred-category}：
{\it 拉回选择},
{\it 拉回函子}


\noindent
见 \ref{categories-definition-fibred-categories-over-C}：
{\it $\mathcal{C}$ 上纤维化范畴 所成的 $2$-范畴}


\noindent
见 \ref{categories-definition-inertia-fibred-category}：
{\it $\mathcal{S}$ 在 $\mathcal{S}'$ 上的相对惯性},
{\it $\mathcal{S}$ 的惯性纤维化范畴 $\mathcal{I}_\mathcal{S}$}


\noindent
见 \ref{categories-definition-fibred-groupoids}：
{\it 以群胚为纤维}


\noindent
见 \ref{categories-definition-categories-fibred-in-groupoids-over-C}：
{\it $\mathcal{C}$ 上群胚纤维化范畴 所成的 $2$-范畴}


\noindent
见 \ref{categories-definition-split-fibred-category}：
{\it 分裂纤维化范畴},
{\it $\mathcal{S}_F$}


\noindent
见 \ref{categories-definition-split-category-fibred-in-groupoids}：
{\it 分裂群胚纤维化范畴},
{\it $\mathcal{S}_F$}


\noindent
见 \ref{categories-definition-discrete}：
{\it 离散的}


\noindent
见 \ref{categories-definition-category-fibred-sets}：
{\it 集合纤维化范畴},
{\it 离散范畴纤维化范畴}


\noindent
见 \ref{categories-definition-categories-fibred-in-sets-over-C}：
{\it $\mathcal{C}$ 上集合纤维化范畴 所成的 $2$-范畴}


\noindent
见 \ref{categories-definition-setoid}：
{\it 集合胚（setoid）}


\noindent
见 \ref{categories-definition-category-fibred-setoids}：
{\it 集合胚纤维化范畴}


\noindent
见 \ref{categories-definition-categories-fibred-in-setoids-over-C}：
{\it $\mathcal{C}$ 上集合胚纤维化范畴 所成的 $2$-范畴}


\noindent
见 \ref{categories-definition-representable-fibred-category}：
{\it 可表示}


\noindent
见 \ref{categories-definition-representable-map-categories-fibred-in-groupoids}：
{\it 可表示},
{\it $\mathcal{X}$ 在 $\mathcal{Y}$ 上相对可表}


\noindent
见 \ref{categories-definition-monoidal-category}：
{\it 幺半范畴}


\noindent
见 \ref{categories-definition-functor-monoidal-categories}：
{\it 函子}


\noindent
见 \ref{categories-definition-invertible}：
{\it 可逆的}


\noindent
见 \ref{categories-definition-dual}：
{\it 左对偶},
{\it 右对偶}


\noindent
见 \ref{categories-definition-symmetric-monoidal-category}：
{\it 对称幺半范畴}


\noindent
见 \ref{categories-definition-functor-symmetric-monoidal-categories}：
{\it 函子}


\noindent
见 \ref{categories-definition-dotted-arrows}：
{\it 虚线箭头的态射}


\medskip\noindent
{\bf 拓扑}

\medskip

\noindent
见 \ref{topology-definition-separated}：
{\it 分离}


\noindent
见 \ref{topology-definition-base}：
{\it $X$ 上拓扑的一个基},
{\it $X$ 上拓扑的一个基底}


\noindent
见 \ref{topology-definition-subbase}：
{\it $X$ 上拓扑的一个子基},
{\it $X$ 上拓扑的一个子基底}


\noindent
见 \ref{topology-definition-submersive}：
{\it 拓扑空间的严格映射},
{\it 商映射}


\noindent
见 \ref{topology-definition-connected-components}：
{\it 连通},
{\it 连通分支}


\noindent
见 \ref{topology-definition-totally-disconnected}：
{\it 全不连通}


\noindent
见 \ref{topology-definition-locally-connected}：
{\it 局部连通}


\noindent
见 \ref{topology-definition-irreducible-components}：
{\it 不可约},
{\it 不可约分支}


\noindent
见 \ref{topology-definition-generic-point}：
{\it 泛点},
{\it Kolmogorov 空间},
{\it 拟清醒},
{\it 清醒}


\noindent
见 \ref{topology-definition-noetherian}：
{\it Noether},
{\it 局部 Noether}


\noindent
见 \ref{topology-definition-Krull}：
{\it 不可约闭子集链},
{\it 长度},
{\it 维数},
{\it Krull 维数},
{\it $X$ 在 $x$ 处的 Krull 维数}


\noindent
见 \ref{topology-definition-equidimensional}：
{\it 等维}


\noindent
见 \ref{topology-definition-codimension}：
{\it 余维数}


\noindent
见 \ref{topology-definition-catenary}：
{\it 链状}


\noindent
见 \ref{topology-definition-quasi-compact}：
{\it 拟紧},
{\it 拟紧},
{\it 逆紧}


\noindent
见 \ref{topology-definition-locally-quasi-compact}：
{\it 局部拟紧}


\noindent
见 \ref{topology-definition-constructible}：
{\it 可构造},
{\it 局部可构造}


\noindent
见 \ref{topology-definition-proper-map}：
{\it 闭的},
{\it Bourbaki-固有的},
{\it 拟固有的},
{\it 普遍闭},
{\it 固有}


\noindent
见 \ref{topology-definition-space-jacobson}：
{\it Jacobson 空间}


\noindent
见 \ref{topology-definition-specialization}：
{\it 特化},
{\it 一般化},
{\it 对特化稳定},
{\it 对一般化稳定}


\noindent
见 \ref{topology-definition-lift-specializations}：
{\it 特化可沿 $f$ 提升},
{\it 特化型的},
{\it 一般化可沿 $f$ 提升},
{\it 一般化型的}


\noindent
见 \ref{topology-definition-dimension-function}：
{\it 直接特化},
{\it 维数函数}


\noindent
见 \ref{topology-definition-nowhere-dense}：
{\it 内部},
{\it 无处稠密}


\noindent
见 \ref{topology-definition-profinite}：
{\it 仿有限的}


\noindent
见 \ref{topology-definition-spectral-space}：
{\it 谱的},
{\it 谱的}


\noindent
见 \ref{topology-definition-extremally-disconnected}：
{\it 极端不连通的}


\noindent
见 \ref{topology-definition-isolated-point}：
{\it 孤立点}


\noindent
见 \ref{topology-definition-paritition}：
{\it 分划},
{\it 分块},
{\it 加细}


\noindent
见 \ref{topology-definition-good-stratification}：
{\it 良好分层}


\noindent
见 \ref{topology-definition-stratification}：
{\it 分层},
{\it 层}


\noindent
见 \ref{topology-definition-locally-finite}：
{\it 局部有限}


\noindent
见 \ref{topology-definition-topological-group}：
{\it 拓扑群},
{\it 拓扑群同态}


\noindent
见 \ref{topology-definition-profinite-group}：
{\it 仿有限群}


\noindent
见 \ref{topology-definition-topological-ring}：
{\it 拓扑环},
{\it 拓扑环同态}


\noindent
见 \ref{topology-definition-topological-module}：
{\it 拓扑模},
{\it 拓扑模同态}


\medskip\noindent
{\bf 空间上的层}

\medskip

\noindent
见 \ref{sheaves-definition-presheaf}：
{\it $X$ 上的集合预层 $\mathcal{F}$},
{\it $X$ 上集合预层的态射 $\varphi : \mathcal{F} \to \mathcal{G}$}


\noindent
见 \ref{sheaves-definition-constant-presheaf}：
{\it 取值为 $A$ 的常值预层}


\noindent
见 \ref{sheaves-definition-abelian-presheaves}：
{\it $X$ 上的阿贝尔群预层},
{\it $X$ 上的阿贝尔预层},
{\it $X$ 上阿贝尔预层的态射}


\noindent
见 \ref{sheaves-definition-presheaf-values-in-category}：
{\it $X$ 上取值于 $\mathcal{C}$ 的预层 $\mathcal{F}$},
{\it 取值于 $\mathcal{C}$ 的预层态射 $\varphi : \mathcal{F} \to \mathcal{G}$}


\noindent
见 \ref{sheaves-definition-underlying-presheaf-sets}：
{\it $\mathcal{F}$ 的底层集合预层}


\noindent
见 \ref{sheaves-definition-presheaf-modules}：
{\it $\mathcal{O}$-模预层},
{\it $\mathcal{O}$-模预层的态射 $\varphi : \mathcal{F} \to \mathcal{G}$}


\noindent
见 \ref{sheaves-definition-sheaf}：
{\it $X$ 上的集合层 $\mathcal{F}$},
{\it 集合层态射}


\noindent
见 \ref{sheaves-definition-constant-sheaf}：
{\it 值为 $A$ 的常值层}


\noindent
见 \ref{sheaves-definition-abelian-sheaf}：
{\it $X$ 上的阿贝尔层},
{\it $X$ 上的阿贝尔群层}


\noindent
见 \ref{sheaves-definition-sheaf-values-in-category}：
{\it 层}


\noindent
见 \ref{sheaves-definition-sheaf-modules}：
{\it $\mathcal{O}$-模层},
{\it $\mathcal{O}$-模层态射}


\noindent
见 \ref{sheaves-definition-separated}：
{\it 分离}


\noindent
见 \ref{sheaves-definition-algebraic-structure}：
{\it 代数结构类型}


\noindent
见 \ref{sheaves-definition-injective-surjective}：
{\it 子预层},
{\it 子层},
{\it 单射的},
{\it 满},
{\it 单射的},
{\it 满}


\noindent
见 \ref{sheaves-definition-f-map}：
{\it $f$-映射 $\xi : \mathcal{G} \to \mathcal{F}$}


\noindent
见 \ref{sheaves-definition-composition-f-maps}：
{\it $\varphi$ 与 $\psi$ 的复合}


\noindent
见 \ref{sheaves-definition-ringed-space}：
{\it 赋环空间},
{\it 赋环空间态射}


\noindent
见 \ref{sheaves-definition-composition-maps-ringed-spaces}：
{\it 赋环空间态射的复合}


\noindent
见 \ref{sheaves-definition-pushforward}：
{\it 推前},
{\it 拉回}


\noindent
见 \ref{sheaves-definition-skyscraper-sheaf}：
{\it 在 $x$ 处取值为 $A$ 的摩天楼层},
{\it 摩天楼层},
{\it 摩天楼层},
{\it 摩天楼层},
{\it 摩天楼层}


\noindent
见 \ref{sheaves-definition-presheaf-basis}：
{\it $\mathcal{B}$ 上的集合预层 $\mathcal{F}$},
{\it $\mathcal{B}$ 上集合预层的态射 $\varphi : \mathcal{F} \to \mathcal{G}$}


\noindent
见 \ref{sheaves-definition-sheaf-basis}：
{\it $\mathcal{B}$ 上的集合层 $\mathcal{F}$},
{\it $\mathcal{B}$ 上集合层的态射}


\noindent
见 \ref{sheaves-definition-sheaf-structures-basis}：
{\it $\mathcal{B}$ 上取值于 $\mathcal{C}$ 的预层 $\mathcal{F}$},
{\it $\mathcal{B}$ 上取值于 $\mathcal{C}$ 的预层态射 $\varphi : \mathcal{F} \to \mathcal{G}$},
{\it $\mathcal{B}$ 上取值于 $\mathcal{C}$ 的层 $\mathcal{F}$}


\noindent
见 \ref{sheaves-definition-sheaf-modules-basis}：
{\it $\mathcal{B}$ 上的 $\mathcal{O}$-模预层 $\mathcal{F}$},
{\it $\mathcal{B}$ 上 $\mathcal{O}$-模预层的态射 $\varphi : \mathcal{F} \to \mathcal{G}$},
{\it $\mathcal{B}$ 上的 $\mathcal{O}$-模层 $\mathcal{F}$}


\noindent
见 \ref{sheaves-definition-restriction}：
{\it $\mathcal{G}$ 到 $U$ 的限制},
{\it $\mathcal{G}$ 到 $U$ 的限制},
{\it 与 $U$ 关联的 $(X, \mathcal{O})$ 的开子空间},
{\it $\mathcal{G}$ 到 $U$ 的限制}


\noindent
见 \ref{sheaves-definition-j-shriek}：
{\it $\mathcal{F}$ 的以空集延拓 $j_{p!}\mathcal{F}$},
{\it $\mathcal{F}$ 的以空集延拓 $j_!\mathcal{F}$}


\noindent
见 \ref{sheaves-definition-j-shriek-structures}：
{\it $\mathcal{F}$ 的以 $0$ 延拓 $j_{p!}\mathcal{F}$},
{\it $\mathcal{F}$ 的以 $0$ 延拓 $j_!\mathcal{F}$},
{\it $\mathcal{F}$ 的以 $e$ 延拓 $j_{p!}\mathcal{F}$},
{\it $\mathcal{F}$ 的以 $e$ 延拓 $j_!\mathcal{F}$},
{\it 以 $0$ 延拓},
{\it 以 $0$ 延拓}


\medskip\noindent
{\bf 位点与层}

\medskip

\noindent
见 \ref{sites-definition-presheaves-sets}：
{\it 集合预层},
{\it 预层态射}


\noindent
见 \ref{sites-definition-presheaf}：
{\it 预层},
{\it 态射}


\noindent
见 \ref{sites-definition-presheaves-injective-surjective}：
{\it 单射的},
{\it 满}


\noindent
见 \ref{sites-definition-sub-presheaf}：
{\it 子预层}


\noindent
见 \ref{sites-definition-image}：
{\it $\varphi$ 的像}


\noindent
见 \ref{sites-definition-family-morphisms-fixed-target}：
{\it 具有固定靶的态射族}


\noindent
见 \ref{sites-definition-site}：
{\it 位点},
{\it $\mathcal{C}$ 的覆盖}


\noindent
见 \ref{sites-definition-sheaf-sets}：
{\it 层}


\noindent
见 \ref{sites-definition-category-sheaves-sets}：
{\it $\Sh(\mathcal{C})$}


\noindent
见 \ref{sites-definition-sheaf}：
{\it 层}


\noindent
见 \ref{sites-definition-morphism-coverings}：
{\it $\mathcal{C}$ 中从 $\mathcal{U}$ 到 $\mathcal{V}$ 的固定靶映射族态射},
{\it 从 $\mathcal{U}$ 到 $\mathcal{V}$ 的态射},
{\it 加细}


\noindent
见 \ref{sites-definition-combinatorial-tautological}：
{\it 组合等价},
{\it 自明等价}


\noindent
见 \ref{sites-definition-separated}：
{\it 分离}


\noindent
见 \ref{sites-definition-associated-sheaf}：
{\it 与 $\mathcal{F}$ 关联的层}


\noindent
见 \ref{sites-definition-sheaves-injective-surjective}：
{\it 单射的},
{\it 满}


\noindent
见 \ref{sites-definition-universal-effective-epimorphisms}：
{\it 有效满态射},
{\it 泛有效满态射}


\noindent
见 \ref{sites-definition-weaker-than-canonical}：
{\it 弱于典范拓扑},
{\it 次典范的}


\noindent
见 \ref{sites-definition-representable-sheaf}：
{\it 可表层},
{\it $\underline{U}$}


\noindent
见 \ref{sites-definition-continuous}：
{\it 连续的}


\noindent
见 \ref{sites-definition-morphism-sites}：
{\it 位点态射}


\noindent
见 \ref{sites-definition-composition-morphisms-sites}：
{\it 复合}


\noindent
见 \ref{sites-definition-topos}：
{\it 托扑斯},
{\it 托扑斯态射},
{\it 复合 $f\circ g$}


\noindent
见 \ref{sites-definition-quasi-compact}：
{\it 拟紧}


\noindent
见 \ref{sites-definition-quasi-compact-topos}：
{\it 拟紧},
{\it 拟紧}


\noindent
见 \ref{sites-definition-cocontinuous}：
{\it 余连续的}


\noindent
见 \ref{sites-definition-localize}：
{\it 位点 $\mathcal{C}$ 在对象 $U$ 处的 局部化},
{\it 局部化态射},
{\it 直像函子},
{\it $\mathcal{F}$ 到 $\mathcal{C}/U$ 的限制},
{\it $\mathcal{G}$ 以空集作的延拓}


\noindent
见 \ref{sites-definition-special-cocontinuous-functor}：
{\it 从 $\mathcal{C}$ 到 $\mathcal{D}$ 的特殊余连续函子 $u$}


\noindent
见 \ref{sites-definition-localize-topos}：
{\it 拓扑斯 $\Sh(\mathcal{C})$ 在 $\mathcal{F}$ 处的局部化},
{\it 局部化态射}


\noindent
见 \ref{sites-definition-point-topos}：
{\it 拓扑斯 $\Sh(\mathcal{C})$ 的点}


\noindent
见 \ref{sites-definition-point}：
{\it 位点 $\mathcal{C}$ 的点 $p$}


\noindent
见 \ref{sites-definition-pushforward-point}：
{\it 摩天楼层}


\noindent
见 \ref{sites-definition-2-morphism-topoi}：
{\it 从 $f$ 到 $g$ 的 2-态射}


\noindent
见 \ref{sites-definition-morphism-points}：
{\it 态射 $f : p \to p'$}


\noindent
见 \ref{sites-definition-enough-points}：
{\it 保守},
{\it 有足够多的点}


\noindent
见 \ref{sites-definition-w-contractible}：
{\it 弱可缩的},
{\it 足够多的弱可缩对象},
{\it 足够多的 $P$ 对象}


\noindent
见 \ref{sites-definition-empty}：
{\it 层论意义下为空}


\noindent
见 \ref{sites-definition-almost-cocontinuous}：
{\it 几乎余连续的}


\noindent
见 \ref{sites-definition-embedding}：
{\it 嵌入}


\noindent
见 \ref{sites-definition-subtopos}：
{\it 子拓扑斯}


\noindent
见 \ref{sites-definition-open-subtopos}：
{\it 开子拓扑斯}


\noindent
见 \ref{sites-definition-closed-subtopos}：
{\it 闭子拓扑斯}


\noindent
见 \ref{sites-definition-immersion-topoi}：
{\it 开浸入},
{\it 闭浸入}


\noindent
见 \ref{sites-definition-pushforward-algebraic-structures}：
{\it 推前}


\noindent
见 \ref{sites-definition-global-sections}：
{\it 整体截面}


\noindent
见 \ref{sites-definition-sieve}：
{\it $U$ 上的筛 $S$}


\noindent
见 \ref{sites-definition-sieve-generated}：
{\it 由态射 $f_i$ 生成的 $U$ 上的筛}


\noindent
见 \ref{sites-definition-pullback-sieve}：
{\it $S$ 沿 $f$ 的拉回}


\noindent
见 \ref{sites-definition-topology}：
{\it $\mathcal{C}$ 上的拓扑}


\noindent
见 \ref{sites-definition-finer}：
{\it 更细},
{\it 更强},
{\it 更粗},
{\it 更弱}


\noindent
见 \ref{sites-definition-sheaf-sets-topology}：
{\it 层}


\noindent
见 \ref{sites-definition-canonical-topology}：
{\it 典范拓扑}


\noindent
见 \ref{sites-definition-topology-associated-site}：
{\it 与 $\mathcal{C}$ 关联的拓扑}


\noindent
见 \ref{sites-definition-presheaf-separated-topology}：
{\it 分离}


\noindent
见 \ref{sites-definition-associated-sheaf-topology}：
{\it 与 $\mathcal{F}$ 关联的层}


\noindent
见 \ref{sites-definition-point-topology}：
{\it 点 $p$}


\medskip\noindent
{\bf 叠}

\medskip

\noindent
见 \ref{stacks-definition-mor-presheaf}：
{\it 从 $x$ 到 $y$ 的态射预层},
{\it 从 $x$ 到 $y$ 的同构预层}


\noindent
见 \ref{stacks-definition-descent-data}：
{\it $\mathcal{S}$ 中相对于族 $\{f_i : U_i \to U\}$ 的下降数据 $(X_i, \varphi_{ij})$},
{\it 余循环条件},
{\it 下降数据的态射 $\psi : (X_i, \varphi_{ij}) \to (X'_i, \varphi'_{ij})$}


\noindent
见 \ref{stacks-definition-pullback-functor}：
{\it 拉回函子}


\noindent
见 \ref{stacks-definition-effective-descent-datum}：
{\it 平凡下降数据},
{\it 典范下降数据},
{\it 有效}


\noindent
见 \ref{stacks-definition-stack}：
{\it 栈}


\noindent
见 \ref{stacks-definition-stacks-over-C}：
{\it $\mathcal{C}$ 上栈的 $2$-范畴}


\noindent
见 \ref{stacks-definition-stack-in-groupoids}：
{\it 群胚中的栈}


\noindent
见 \ref{stacks-definition-stacks-in-groupoids-over-C}：
{\it $\mathcal{C}$ 上群胚中的栈的 $2$-范畴}


\noindent
见 \ref{stacks-definition-stack-in-sets}：
{\it 类集合（setoid）中的栈},
{\it 集合中的栈},
{\it 离散范畴中的栈}


\noindent
见 \ref{stacks-definition-stacks-in-setoids-over-C}：
{\it $\mathcal{C}$ 上类集合中的栈的 $2$-范畴}


\noindent
见 \ref{stacks-definition-topology-inherited}：
{\it $\mathcal{S}$ 上从 $\mathcal{C}$ 继承的位点结构},
{\it $\mathcal{S}$ 配备从 $\mathcal{C}$ 继承的拓扑}


\noindent
见 \ref{stacks-definition-gerbe}：
{\it gerbe（胚）}


\noindent
见 \ref{stacks-definition-gerbe-over-stack-in-groupoids}：
{\it ……上的 gerbe（胚）}


\noindent
见 \ref{stacks-definition-pushforward-stack}：
{\it $f_*\mathcal{S}$},
{\it $\mathcal{S}$ 沿 $f$ 的推前}


\noindent
见 \ref{stacks-definition-pullback-stack}：
{\it $f^{-1}\mathcal{S}$},
{\it $\mathcal{S}$ 沿 $f$ 的拉回}


\medskip\noindent
{\bf 域}

\medskip

\noindent
见 \ref{fields-definition-field}：
{\it 域},
{\it 子域}


\noindent
见 \ref{fields-definition-domain}：
{\it 整环},
{\it 整环}


\noindent
见 \ref{fields-definition-characteristic}：
{\it 特征},
{\it $F$ 的素子域}


\noindent
见 \ref{fields-definition-extension}：
{\it 域扩张}


\noindent
见 \ref{fields-definition-tower}：
{\it 塔}


\noindent
见 \ref{fields-definition-generated-by}：
{\it 生成域扩张},
{\it 有限生成域扩张}


\noindent
见 \ref{fields-definition-degree}：
{\it 次数},
{\it 有限}


\noindent
见 \ref{fields-definition-number-field}：
{\it 数域}


\noindent
见 \ref{fields-definition-algebraic}：
{\it 代数的},
{\it 代数扩张}


\noindent
见 \ref{fields-definition-minimal-polynomial}：
{\it 最小多项式}


\noindent
见 \ref{fields-definition-algebraically-closed}：
{\it 代数闭域}


\noindent
见 \ref{fields-definition-algebraic-closure}：
{\it 代数闭包}


\noindent
见 \ref{fields-definition-relatively-prime}：
{\it 互素}


\noindent
见 \ref{fields-definition-separable}：
{\it 可分},
{\it 可分},
{\it 可分}


\noindent
见 \ref{fields-definition-separable-degree}：
{\it 可分次数}


\noindent
见 \ref{fields-definition-purely-inseparable}：
{\it 纯不可分},
{\it 纯不可分},
{\it 纯不可分}


\noindent
见 \ref{fields-definition-insep-degree}：
{\it 可分次数},
{\it 不可分次数},
{\it 不可分度}


\noindent
见 \ref{fields-definition-normal}：
{\it 正规}


\noindent
见 \ref{fields-definition-automorphisms}：
{\it $E$ 在 $F$ 上的自同构},
{\it $E/F$ 的自同构}


\noindent
见 \ref{fields-definition-splitting-field}：
{\it $P$ 在 $F$ 上的分裂域}


\noindent
见 \ref{fields-definition-normal-closure}：
{\it $E$ 在 $F$ 上的正规闭包}


\noindent
见 \ref{fields-definition-trace-norm}：
{\it 迹},
{\it 范数}


\noindent
见 \ref{fields-definition-trace-pairing}：
{\it 迹配对}


\noindent
见 \ref{fields-definition-discriminant}：
{\it $L/K$ 的判别式}


\noindent
见 \ref{fields-definition-galois}：
{\it Galois 扩张}


\noindent
见 \ref{fields-definition-galois-group}：
{\it Galois 群}


\noindent
见 \ref{fields-definition-transcendence}：
{\it 代数无关},
{\it 纯超越扩张},
{\it 超越基}


\noindent
见 \ref{fields-definition-transcendence-degree}：
{\it 超越次数}


\noindent
见 \ref{fields-definition-algebraically-closed-in}：
{\it $k$ 在 $K$ 中的代数闭包},
{\it $K$ 中代数闭}


\noindent
见 \ref{fields-definition-compositum}：
{\it $K$ 与 $L$ 在 $\Omega$ 中的合成域}


\noindent
见 \ref{fields-definition-linearly-disjoint}：
{\it 在 $\Omega$ 中于 $k$ 上线性无交}


\noindent
见 \ref{fields-definition-separable-algebraic}：
{\it 代数的},
{\it 可分},
{\it 纯不可分},
{\it 正规},
{\it Galois 扩张}


\medskip\noindent
{\bf 交换代数}

\medskip

\noindent
见 \ref{algebra-definition-module-finite-type}：
{\it 有限 $R$-模},
{\it 有限生成 $R$-模},
{\it 有限表示 $R$-模},
{\it 有限表示的 $R$-模}


\noindent
见 \ref{algebra-definition-finite-type}：
{\it 有限型},
{\it $S$ 是有限型 $R$-代数},
{\it 有限表示}


\noindent
见 \ref{algebra-definition-finite-ring-map}：
{\it 有限}


\noindent
见 \ref{algebra-definition-directed-system}：
{\it 以 $I$ 为指标的 $R$-模系统 $(M_i, \mu_{ij})$},
{\it 有向系统}


\noindent
见 \ref{algebra-definition-homomorphism-directed-systems}：
{\it 系统同态}


\noindent
见 \ref{algebra-definition-multiplicative-subset}：
{\it $R$ 的乘法子集}


\noindent
见 \ref{algebra-definition-localization}：
{\it $A$ 关于 $S$ 的局部化}


\noindent
见 \ref{algebra-definition-localization-module}：
{\it 局部化}


\noindent
见 \ref{algebra-definition-relation}：
{\it 关系}


\noindent
见 \ref{algebra-definition-bilinear}：
{\it $R$-双线性的}


\noindent
见 \ref{algebra-definition-bimodule}：
{\it $(A, B)$-双模}


\noindent
见 \ref{algebra-definition-base-change}：
{\it 基变换},
{\it 基变换}


\noindent
见 \ref{algebra-definition-spectrum-ring}：
{\it 谱}


\noindent
见 \ref{algebra-definition-Zariski-topology}：
{\it 扎里斯基},
{\it 标准开集}


\noindent
见 \ref{algebra-definition-local-ring}：
{\it 局部环},
{\it 剩余域},
{\it 局部同态},
{\it 局部环映射 $\varphi : R \to S$}


\noindent
见 \ref{algebra-definition-oka-family}：
{\it Oka 族}


\noindent
见 \ref{algebra-definition-locally-nilpotent-ideal}：
{\it 局部幂零的},
{\it 幂零的}


\noindent
见 \ref{algebra-definition-ring-jacobson}：
{\it Jacobson 环}


\noindent
见 \ref{algebra-definition-integral-ring-map}：
{\it 整于 $R$},
{\it 整}


\noindent
见 \ref{algebra-definition-integral-closure}：
{\it 整闭包},
{\it 整闭}


\noindent
见 \ref{algebra-definition-domain-normal}：
{\it 正规}


\noindent
见 \ref{algebra-definition-almost-integral}：
{\it 几乎整于 $R$},
{\it 完全正规}


\noindent
见 \ref{algebra-definition-ring-normal}：
{\it 正规}


\noindent
见 \ref{algebra-definition-integral-over-ideal}：
{\it 整于 $I$}


\noindent
见 \ref{algebra-definition-flat}：
{\it 平坦},
{\it 忠实平坦},
{\it 平坦},
{\it 忠实平坦}


\noindent
见 \ref{algebra-definition-support-module}：
{\it $M$ 的支集}


\noindent
见 \ref{algebra-definition-annihilator}：
{\it $m$ 的零化子},
{\it $M$ 的零化子}


\noindent
见 \ref{algebra-definition-going-up-down}：
{\it 上升},
{\it 下降}


\noindent
见 \ref{algebra-definition-separable-field-extension}：
{\it 在 $k$ 上可分生成},
{\it 在 $k$ 上可分}


\noindent
见 \ref{algebra-definition-geometrically-reduced}：
{\it 在 $k$ 上几何既约}


\noindent
见 \ref{algebra-definition-perfect}：
{\it 完美的}


\noindent
见 \ref{algebra-definition-perfection}：
{\it 完美闭包}


\noindent
见 \ref{algebra-definition-geometrically-irreducible}：
{\it 在 $k$ 上几何不可约}


\noindent
见 \ref{algebra-definition-geometrically-connected}：
{\it 在 $k$ 上几何连通}


\noindent
见 \ref{algebra-definition-geometrically-integral}：
{\it 在 $k$ 上几何整}


\noindent
见 \ref{algebra-definition-valuation-ring}：
{\it 支配},
{\it 赋值环},
{\it 有中心}


\noindent
见 \ref{algebra-definition-value-group}：
{\it 值群},
{\it 赋值},
{\it 离散赋值环}


\noindent
见 \ref{algebra-definition-length}：
{\it 长度}


\noindent
见 \ref{algebra-definition-simple-module}：
{\it 单}


\noindent
见 \ref{algebra-definition-artinian}：
{\it Artin}


\noindent
见 \ref{algebra-definition-essentially-finite-p-t}：
{\it 本质有限型的},
{\it 本质有限表示的}


\noindent
见 \ref{algebra-definition-proj}：
{\it 齐次谱}


\noindent
见 \ref{algebra-definition-numerical-polynomial}：
{\it 数值多项式}


\noindent
见 \ref{algebra-definition-ideal-definition}：
{\it $R$ 的定义理想}


\noindent
见 \ref{algebra-definition-hilbert-polynomial}：
{\it Hilbert 多项式}


\noindent
见 \ref{algebra-definition-d}：
{\it $d(M)$}


\noindent
见 \ref{algebra-definition-chain-primes}：
{\it 素理想链},
{\it 长度}


\noindent
见 \ref{algebra-definition-Krull}：
{\it Krull 维数}


\noindent
见 \ref{algebra-definition-height}：
{\it 高度}


\noindent
见 \ref{algebra-definition-regular-local}：
{\it $R$ 的一组参数系},
{\it 正则局部环},
{\it 正则参数系}


\noindent
见 \ref{algebra-definition-associated}：
{\it 相伴}


\noindent
见 \ref{algebra-definition-symbolic-power}：
{\it 符号幂}


\noindent
见 \ref{algebra-definition-relative-assassin}：
{\it $N$ 在 $S/R$ 上的相对相伴素理想集}


\noindent
见 \ref{algebra-definition-weakly-associated}：
{\it 弱相伴}


\noindent
见 \ref{algebra-definition-embedded-primes}：
{\it 嵌入伴随素理想},
{\it $R$ 的嵌入素理想}


\noindent
见 \ref{algebra-definition-regular-sequence}：
{\it $M$-正则序列},
{\it $I$ 中的 $M$-正则序列},
{\it 正则序列}


\noindent
见 \ref{algebra-definition-quasi-regular-sequence}：
{\it $M$-拟正则},
{\it 拟正则序列}


\noindent
见 \ref{algebra-definition-blow-up}：
{\it 吹起代数},
{\it Rees 代数},
{\it 仿射吹起代数}


\noindent
见 \ref{algebra-definition-finite-free-resolution}：
{\it 消解},
{\it 由自由 $R$-模给出的 $M$ 的消解},
{\it 由有限自由 $R$-模给出的 $M$ 的消解}


\noindent
见 \ref{algebra-definition-depth}：
{\it $I$-深度},
{\it 深度}


\noindent
见 \ref{algebra-definition-projective}：
{\it 投射}


\noindent
见 \ref{algebra-definition-locally-free}：
{\it 局部自由},
{\it 有限局部自由},
{\it 秩为 $r$ 的有限局部自由}


\noindent
见 \ref{algebra-definition-universally-injective}：
{\it 普遍单射},
{\it 普遍正合}


\noindent
见 \ref{algebra-definition-devissage}：
{\it 直和 dévissage},
{\it Kaplansky dévissage}


\noindent
见 \ref{algebra-definition-ML-system}：
{\it Mittag--Leffler}


\noindent
见 \ref{algebra-definition-ML-inductive-system}：
{\it 模的 Mittag--Leffler 定向系统}


\noindent
见 \ref{algebra-definition-domination}：
{\it 支配}


\noindent
见 \ref{algebra-definition-mittag-leffler-module}：
{\it Mittag--Leffler}


\noindent
见 \ref{algebra-definition-coherent}：
{\it 凝聚模},
{\it 凝聚环}


\noindent
见 \ref{algebra-definition-complete}：
{\it $I$-进完备},
{\it $I$-进完备}


\noindent
见 \ref{algebra-definition-rank}：
{\it 秩}


\noindent
见 \ref{algebra-definition-CM}：
{\it Cohen--Macaulay}


\noindent
见 \ref{algebra-definition-maximal-CM}：
{\it 极大 Cohen--Macaulay}


\noindent
见 \ref{algebra-definition-module-CM}：
{\it Cohen--Macaulay}


\noindent
见 \ref{algebra-definition-local-ring-CM}：
{\it Cohen--Macaulay}


\noindent
见 \ref{algebra-definition-ring-CM}：
{\it Cohen--Macaulay}


\noindent
见 \ref{algebra-definition-catenary}：
{\it 链状}


\noindent
见 \ref{algebra-definition-universally-catenary}：
{\it 普遍链状}


\noindent
见 \ref{algebra-definition-pure-ideal}：
{\it 纯}


\noindent
见 \ref{algebra-definition-finite-proj-dim}：
{\it 有限投射维数},
{\it 投射维数}


\noindent
见 \ref{algebra-definition-finite-gl-dim}：
{\it 有限整体维数},
{\it 整体维数}


\noindent
见 \ref{algebra-definition-regular}：
{\it 正则概形}


\noindent
见 \ref{algebra-definition-fibre}：
{\it 纤维在 $\mathfrak q$ 处的局部环}


\noindent
见 \ref{algebra-definition-uniformizer}：
{\it 一致化元}


\noindent
见 \ref{algebra-definition-irreducible-prime-element}：
{\it 相伴元},
{\it 不可约},
{\it 素的}


\noindent
见 \ref{algebra-definition-UFD}：
{\it 唯一分解整环},
{\it UFD}


\noindent
见 \ref{algebra-definition-PID}：
{\it 主理想整环},
{\it PID}


\noindent
见 \ref{algebra-definition-dedekind-domain}：
{\it Dedekind 整环}


\noindent
见 \ref{algebra-definition-ord}：
{\it 沿 $R$ 的消失阶}


\noindent
见 \ref{algebra-definition-lattice}：
{\it $V$ 中的格}


\noindent
见 \ref{algebra-definition-distance}：
{\it $M$ 与 $M'$ 之间的距离}


\noindent
见 \ref{algebra-definition-quasi-finite}：
{\it 在 $\mathfrak q$ 处拟有限},
{\it 拟有限}


\noindent
见 \ref{algebra-definition-strongly-transcendental}：
{\it 在 $R$ 上强超越}


\noindent
见 \ref{algebra-definition-relative-dimension}：
{\it $S/R$ 在 $\mathfrak q$ 处的相对维数},
{\it ……的相对维数}


\noindent
见 \ref{algebra-definition-derivation}：
{\it 导子},
{\it $R$-导子},
{\it Leibniz 法则}


\noindent
见 \ref{algebra-definition-differentials}：
{\it K\"ahler 微分模},
{\it 微分模}


\noindent
见 \ref{algebra-definition-differential-operators}：
{\it $k$ 阶微分算子 $D : M \to N$}


\noindent
见 \ref{algebra-definition-module-principal-parts}：
{\it $k$ 阶主部件模}


\noindent
见 \ref{algebra-definition-naive-cotangent-complex}：
{\it 朴素余切复形}


\noindent
见 \ref{algebra-definition-lci-field}：
{\it $k$ 上的整体完全交},
{\it $k$ 上的局部完全交}


\noindent
见 \ref{algebra-definition-lci-local-ring}：
{\it （在 $k$ 上的）完全交}


\noindent
见 \ref{algebra-definition-lci}：
{\it syntomic},
{\it $R$ 上平坦的局部完全交}


\noindent
见 \ref{algebra-definition-relative-global-complete-intersection}：
{\it 相对整体完全交}


\noindent
见 \ref{algebra-definition-smooth}：
{\it 光滑}


\noindent
见 \ref{algebra-definition-standard-smooth}：
{\it $R$ 上的标准光滑代数},
{\it 标准光滑},
{\it 标准光滑}


\noindent
见 \ref{algebra-definition-smooth-at-prime}：
{\it 在 $\mathfrak q$ 处光滑}


\noindent
见 \ref{algebra-definition-formally-smooth}：
{\it 在 $R$ 上形式光滑}


\noindent
见 \ref{algebra-definition-small-extension}：
{\it 小扩张}


\noindent
见 \ref{algebra-definition-etale}：
{\it 平展},
{\it 在 $\mathfrak q$ 处平展}


\noindent
见 \ref{algebra-definition-standard-etale}：
{\it 标准平展的}


\noindent
见 \ref{algebra-definition-formally-unramified}：
{\it 在 $R$ 上形式无分歧}


\noindent
见 \ref{algebra-definition-universal-thickening}：
{\it 泛一阶加厚},
{\it 余法模},
{\it $C_{S/R}$}


\noindent
见 \ref{algebra-definition-formally-etale}：
{\it 在 $R$ 上形式平展}


\noindent
见 \ref{algebra-definition-unramified}：
{\it 无分歧},
{\it G-非分歧},
{\it 在 $\mathfrak q$ 处无分歧},
{\it 在 $\mathfrak q$ 处 G-无分歧}


\noindent
见 \ref{algebra-definition-henselian}：
{\it Hensel},
{\it 严格 Hensel}


\noindent
见 \ref{algebra-definition-henselization}：
{\it Hensel 化},
{\it $R$ 关于 $\kappa \subset \kappa^{sep}$ 的严格 Hensel 化},
{\it 严格 Hensel 化}


\noindent
见 \ref{algebra-definition-conditions}：
{\it $(R_k)$},
{\it 在余维 $\leq k$ 处正则},
{\it $(S_k)$}


\noindent
见 \ref{algebra-definition-complete-local-ring}：
{\it 完备局部环}


\noindent
见 \ref{algebra-definition-coefficient-ring}：
{\it 系数环}


\noindent
见 \ref{algebra-definition-cohen-ring}：
{\it Cohen 环}


\noindent
见 \ref{algebra-definition-N}：
{\it N-1},
{\it N-2},
{\it Japanese 概形}


\noindent
见 \ref{algebra-definition-nagata}：
{\it 泛 Japanese 概形},
{\it Nagata 环}


\noindent
见 \ref{algebra-definition-analytically-unramified}：
{\it 解析无分歧},
{\it 解析无分歧}


\noindent
见 \ref{algebra-definition-geometrically-normal}：
{\it 几何正规}


\noindent
见 \ref{algebra-definition-geometrically-regular}：
{\it 几何正则}


\medskip\noindent
{\bf Brauer 群}

\medskip

\noindent
见 \ref{brauer-definition-finite}：
{\it 有限}


\noindent
见 \ref{brauer-definition-skew-field}：
{\it 除环}


\noindent
见 \ref{brauer-definition-simple}：
{\it 单},
{\it 单}


\noindent
见 \ref{brauer-definition-central}：
{\it 中心的}


\noindent
见 \ref{brauer-definition-opposite}：
{\it 反代数}


\noindent
见 \ref{brauer-definition-brauer-group}：
{\it Brauer 群}


\noindent
见 \ref{brauer-definition-splitting}：
{\it 分裂},
{\it 分裂域}


\medskip\noindent
{\bf 同调代数}

\medskip

\noindent
见 \ref{homology-definition-preadditive}：
{\it 预加性},
{\it 加性}


\noindent
见 \ref{homology-definition-zero-object}：
{\it 零对象}


\noindent
见 \ref{homology-definition-direct-sum}：
{\it 直和}


\noindent
见 \ref{homology-definition-additive-category}：
{\it 加性}


\noindent
见 \ref{homology-definition-kernel}：
{\it 核},
{\it 余核},
{\it $f$ 的余像},
{\it $f$ 的像}


\noindent
见 \ref{homology-definition-karoubian}：
{\it Karoubi 完备}


\noindent
见 \ref{homology-definition-abelian-category}：
{\it 阿贝尔}


\noindent
见 \ref{homology-definition-injective-surjective}：
{\it 单射的},
{\it 满},
{\it 子对象},
{\it 商}


\noindent
见 \ref{homology-definition-exact}：
{\it 复形},
{\it 在 $y$ 处正合},
{\it 在 $x_i$ 处正合},
{\it 正合},
{\it 正合列},
{\it 正合复形},
{\it 短正合列}


\noindent
见 \ref{homology-definition-ses-split}：
{\it 分裂}


\noindent
见 \ref{homology-definition-extension}：
{\it 以 $A$ 为核的 $B$ 的扩张 $E$},
{\it 扩张的态射}


\noindent
见 \ref{homology-definition-ext-group}：
{\it $\Ext$-群}


\noindent
见 \ref{homology-definition-simple}：
{\it 单}


\noindent
见 \ref{homology-definition-Artinian}：
{\it Artin},
{\it Artin}


\noindent
见 \ref{homology-definition-Noetherian}：
{\it Noether},
{\it Noether}


\noindent
见 \ref{homology-definition-serre-subcategory}：
{\it Serre 子范畴},
{\it 弱 Serre 子范畴}


\noindent
见 \ref{homology-definition-kernel-category}：
{\it 函子 $F$ 的核}


\noindent
见 \ref{homology-definition-K-zero}：
{\it $\mathcal{A}$ 的第零 $K$-群}


\noindent
见 \ref{homology-definition-cohomological-delta-functor}：
{\it 上同调 $\delta$-函子},
{\it $\delta$-函子}


\noindent
见 \ref{homology-definition-morphism-delta-functors}：
{\it 从 $F$ 到 $G$ 的 $\delta$-函子态射}


\noindent
见 \ref{homology-definition-universal-delta-functor}：
{\it 泛 $\delta$-函子}


\noindent
见 \ref{homology-definition-homotopy-equivalent}：
{\it 同伦等价},
{\it 同伦等价}


\noindent
见 \ref{homology-definition-quasi-isomorphism}：
{\it 拟同构},
{\it 无环}


\noindent
见 \ref{homology-definition-homotopy-equivalent-cochain}：
{\it 同伦等价},
{\it 同伦等价}


\noindent
见 \ref{homology-definition-quasi-isomorphism-cochain}：
{\it 拟同构},
{\it 无环}


\noindent
见 \ref{homology-definition-shift}：
{\it $k$-移位链复形 $A[k]_\bullet$}


\noindent
见 \ref{homology-definition-homology-shift}：
{\it $H_{i + k}(A_\bullet) \rightarrow H_i(A[k]_\bullet)$}


\noindent
见 \ref{homology-definition-shift-cochain}：
{\it $k$-移位上链复形 $A[k]^\bullet$}


\noindent
见 \ref{homology-definition-cohomology-shift}：
{\it $H^{i + k}(A^\bullet) \longrightarrow H^i(A[k]^\bullet)$}


\noindent
见 \ref{homology-definition-graded}：
{\it $\mathcal{A}$ 的分次对象范畴}


\noindent
见 \ref{homology-definition-graded-shift}：
{\it 移位}


\noindent
见 \ref{homology-definition-additive-monoidal}：
{\it 加性幺半范畴}


\noindent
见 \ref{homology-definition-double-complex}：
{\it 双复形}


\noindent
见 \ref{homology-definition-associated-simple-complex}：
{\it 相伴单复形},
{\it 相伴总复形}


\noindent
见 \ref{homology-definition-filtered}：
{\it 递减滤过},
{\it $\mathcal{A}$ 的滤过对象},
{\it 滤过对象的态射 $(A, F) \to (B, F)$},
{\it 诱导滤过},
{\it 商滤过},
{\it 有限},
{\it 分离},
{\it 穷尽}


\noindent
见 \ref{homology-definition-strict}：
{\it 严格}


\noindent
见 \ref{homology-definition-spectral-sequence}：
{\it $\mathcal{A}$ 中的谱序列},
{\it 谱序列的态射}


\noindent
见 \ref{homology-definition-limit-spectral-sequence}：
{\it 极限},
{\it 在 $E_r$ 处退化}


\noindent
见 \ref{homology-definition-exact-couple}：
{\it 正合偶},
{\it 正合偶的态射}


\noindent
见 \ref{homology-definition-spectral-sequence-associated-exact-couple}：
{\it 与正合偶关联的谱序列}


\noindent
见 \ref{homology-definition-differential-object}：
{\it 微分对象},
{\it 微分对象的态射}


\noindent
见 \ref{homology-definition-differential-object-homology}：
{\it 同调}


\noindent
见 \ref{homology-definition-differential-object-selfmap}：
{\it 与 $(A, d, \alpha)$ 关联的谱序列}


\noindent
见 \ref{homology-definition-filtered-differential}：
{\it 滤过微分对象}


\noindent
见 \ref{homology-definition-filtration-cohomology-filtered-differential}：
{\it 诱导滤过}


\noindent
见 \ref{homology-definition-filtered-differential-ss-converges}：
{\it 弱收敛于 $H(K)$},
{\it 收敛于 $H(K)$}


\noindent
见 \ref{homology-definition-filtered-complex}：
{\it $\mathcal{A}$ 的滤过复形 $K^\bullet$}


\noindent
见 \ref{homology-definition-filtration-cohomology-filtered-complex}：
{\it 诱导滤过}


\noindent
见 \ref{homology-definition-bounded-ss}：
{\it 正则概形},
{\it 余正则},
{\it 有界},
{\it 下有界},
{\it 上有界}


\noindent
见 \ref{homology-definition-filtered-complex-ss-converges}：
{\it 弱收敛于 $H^*(K^\bullet)$},
{\it 收敛于 $H^*(K^\bullet)$},
{\it 收敛于 $H^*(K^\bullet)$}


\noindent
见 \ref{homology-definition-ss-double-complex-converge}：
{\it 弱收敛于 $H^n(\text{Tot}(K^{\bullet, \bullet}))$},
{\it 收敛于 $H^n(\text{Tot}(K^{\bullet, \bullet}))$},
{\it 收敛于 $H^n(\text{Tot}(K^{\bullet, \bullet}))$},
{\it 弱收敛于 $H^n(\text{Tot}(K^{\bullet, \bullet}))$},
{\it 收敛于 $H^n(\text{Tot}(K^{\bullet, \bullet}))$},
{\it 收敛于 $H^n(\text{Tot}(K^{\bullet, \bullet}))$}


\noindent
见 \ref{homology-definition-injective}：
{\it 单射的}


\noindent
见 \ref{homology-definition-enough-injectives}：
{\it 有足够多的内射对象}


\noindent
见 \ref{homology-definition-functorial-injective-embedding}：
{\it 函子性内射嵌入}


\noindent
见 \ref{homology-definition-projective}：
{\it 投射}


\noindent
见 \ref{homology-definition-enough-projectives}：
{\it 有足够多的投射对象}


\noindent
见 \ref{homology-definition-functorial-projective-surjections}：
{\it 函子性投射满射}


\noindent
见 \ref{homology-definition-Mittag-Leffler}：
{\it Mittag--Leffler 条件},
{\it ML}


\medskip\noindent
{\bf 导出范畴}

\medskip

\noindent
见 \ref{derived-definition-triangle}：
{\it 三角},
{\it 三角的态射}


\noindent
见 \ref{derived-definition-triangulated-category}：
{\it 三角范畴},
{\it 可区别三角},
{\it 预三角范畴}


\noindent
见 \ref{derived-definition-exact-functor-triangulated-categories}：
{\it 正合函子},
{\it 三角函子}


\noindent
见 \ref{derived-definition-triangulated-subcategory}：
{\it 预三角子范畴},
{\it 三角子范畴}


\noindent
见 \ref{derived-definition-homological}：
{\it 同调的},
{\it 上同调的}


\noindent
见 \ref{derived-definition-delta-functor}：
{\it 从 $\mathcal{A}$ 到 $\mathcal{D}$ 的 $\delta$-函子},
{\it 给定 $\delta$-函子下该短正合列的像}


\noindent
见 \ref{derived-definition-localization}：
{\it 与三角结构相容}


\noindent
见 \ref{derived-definition-saturated}：
{\it 饱和}


\noindent
见 \ref{derived-definition-kernel-category}：
{\it $F$ 的核},
{\it $H$ 的核}


\noindent
见 \ref{derived-definition-quotient-category}：
{\it 商范畴 $\mathcal{D}/\mathcal{B}$},
{\it 商函子}


\noindent
见 \ref{derived-definition-complexes-notation}：
{\it （上链）复形范畴},
{\it 下有界},
{\it 上有界},
{\it 有界}


\noindent
见 \ref{derived-definition-cone}：
{\it 锥}


\noindent
见 \ref{derived-definition-termwise-split-map}：
{\it 逐项分裂单射 $\alpha : A^\bullet \to B^\bullet$},
{\it 逐项分裂满射 $\beta : B^\bullet \to C^\bullet$}


\noindent
见 \ref{derived-definition-split-ses}：
{\it $\mathcal{A}$ 的复形的逐项分裂正合列},
{\it 与逐项分裂复形列关联的三角}


\noindent
见 \ref{derived-definition-distinguished-triangle}：
{\it $K(\mathcal{A})$ 的可区别三角}


\noindent
见 \ref{derived-definition-unbounded-derived-category}：
{\it $\mathcal{A}$ 的导出范畴},
{\it 有界导出范畴}


\noindent
见 \ref{derived-definition-finite-filtered}：
{\it $\mathcal{A}$ 的有限滤过对象范畴}


\noindent
见 \ref{derived-definition-filtered-acyclic}：
{\it 滤过拟同构},
{\it 滤过无环}


\noindent
见 \ref{derived-definition-filtered-derived}：
{\it $\mathcal{A}$ 的滤过导出范畴}


\noindent
见 \ref{derived-definition-filtered-derived-bounded}：
{\it 有界滤过导出范畴}


\noindent
见 \ref{derived-definition-right-derived-functor-defined}：
{\it 右导出函子 $RF$ 在此处有定义},
{\it $RF$ 在 $X$ 处的值},
{\it 左导出函子 $LF$ 在此处有定义},
{\it $LF$ 在 $X$ 处的值}


\noindent
见 \ref{derived-definition-everywhere-defined}：
{\it 右可导},
{\it 处处有定义},
{\it 左可导},
{\it 处处有定义}


\noindent
见 \ref{derived-definition-computes}：
{\it 计算出},
{\it 计算出}


\noindent
见 \ref{derived-definition-derived-functor}：
{\it $F$ 的右导出函子},
{\it $F$ 的左导出函子},
{\it 对 $F$ 右无环},
{\it 对 $RF$ 无环},
{\it 对 $F$ 左无环},
{\it 对 $LF$ 无环}


\noindent
见 \ref{derived-definition-higher-derived-functors}：
{\it $F$ 的第 $i$ 个右导出函子 $R^iF$}


\noindent
见 \ref{derived-definition-injective-resolution}：
{\it $A$ 的内射消解},
{\it $K^\bullet$ 的内射消解}


\noindent
见 \ref{derived-definition-projective-resolution}：
{\it $A$ 的投射消解},
{\it $K^\bullet$ 的投射消解}


\noindent
见 \ref{derived-definition-cartan-eilenberg}：
{\it Cartan--Eilenberg 分解}


\noindent
见 \ref{derived-definition-localization-functor}：
{\it 消解函子}


\noindent
见 \ref{derived-definition-filtered-complexes-notation}：
{\it 滤过内射}


\noindent
见 \ref{derived-definition-ext}：
{\it 第 $i$ 扩张群}


\noindent
见 \ref{derived-definition-yoneda-extension}：
{\it Yoneda 扩张},
{\it 等价}


\noindent
见 \ref{derived-definition-K-zero}：
{\it $\mathcal{D}$ 的第零 $K$-群}


\noindent
见 \ref{derived-definition-K-injective}：
{\it K-内射}


\noindent
见 \ref{derived-definition-derived-colimit}：
{\it 导出余极限},
{\it 同伦余极限}


\noindent
见 \ref{derived-definition-derived-limit}：
{\it 导出极限},
{\it 同伦极限}


\noindent
见 \ref{derived-definition-generators}：
{\it 经典生成元},
{\it 强生成元},
{\it 弱生成元},
{\it 生成元}


\noindent
见 \ref{derived-definition-compact-object}：
{\it 紧对象}


\noindent
见 \ref{derived-definition-compactly-generated}：
{\it 紧生成}


\noindent
见 \ref{derived-definition-orthogonal}：
{\it 右正交},
{\it 左正交}


\noindent
见 \ref{derived-definition-admissible}：
{\it 右可容许},
{\it 左可容许},
{\it 双侧可容许}


\noindent
见 \ref{derived-definition-postnikov-system}：
{\it Postnikov 系统},
{\it Postnikov 系统的态射}


\medskip\noindent
{\bf 单纯方法}

\medskip

\noindent
见 \ref{simplicial-definition-face-degeneracy}：
{\it $\delta^n_j : [n-1] \to [n]$},
{\it $\sigma^n_j : [n + 1]  \to [n]$}


\noindent
见 \ref{simplicial-definition-simplicial-object}：
{\it $\mathcal{C}$ 的单纯对象 $U$},
{\it 单纯集},
{\it 单纯阿贝尔群},
{\it 单纯对象的态射 $U \to U'$},
{\it $\mathcal{C}$ 的单纯对象范畴}


\noindent
见 \ref{simplicial-definition-cosimplicial-object}：
{\it $\mathcal{C}$ 的余单纯对象 $U$},
{\it 余单纯集},
{\it 余单纯阿贝尔群},
{\it 余单纯对象的态射 $U \to U'$},
{\it $\mathcal{C}$ 的余单纯对象范畴}


\noindent
见 \ref{simplicial-definition-product}：
{\it $U$ 与 $V$ 的积}


\noindent
见 \ref{simplicial-definition-fibre-product}：
{\it $V$ 与 $W$ 在 $U$ 上的纤维积}


\noindent
见 \ref{simplicial-definition-push-out}：
{\it $V$ 与 $W$ 在 $U$ 上的推出}


\noindent
见 \ref{simplicial-definition-product-cosimplicial-objects}：
{\it $U$ 与 $V$ 的积}


\noindent
见 \ref{simplicial-definition-fibre-product-cosimplicial-objects}：
{\it $V$ 与 $W$ 在 $U$ 上的纤维积}


\noindent
见 \ref{simplicial-definition-terminology-simplicial-sets}：
{\it $U$ 的 $n$-单形},
{\it $x$ 的面},
{\it $x$ 的退化度},
{\it 退化}


\noindent
见 \ref{simplicial-definition-truncated-simplicial-object}：
{\it $\mathcal{C}$ 的 $n$-截断单纯对象},
{\it $n$-截断单纯对象的态射}


\noindent
见 \ref{simplicial-definition-product-with-simplicial-set}：
{\it $U$ 与 $V$ 的积 $U \times V$},
{\it 积 $U \times V$ 存在}


\noindent
见 \ref{simplicial-definition-hom-deltak-cosimplicial}：
{\it $\Hom(U, V)$}


\noindent
见 \ref{simplicial-definition-hom-deltak-simplicial}：
{\it $\Hom(U, V)$}


\noindent
见 \ref{simplicial-definition-hom-from-simplicial-set}：
{\it $\Hom(U, V)$}


\noindent
见 \ref{simplicial-definition-split}：
{\it 分裂},
{\it 分裂}


\noindent
见 \ref{simplicial-definition-augmentation}：
{\it $U$ 朝向 $\mathcal{C}$ 中对象 $X$ 的增广 $\epsilon : U \to X$}


\noindent
见 \ref{simplicial-definition-eilenberg-maclane}：
{\it Eilenberg--Mac Lane 对象 $K(A, k)$}


\noindent
见 \ref{simplicial-definition-homotopy}：
{\it 从 $a$ 到 $b$ 的同伦},
{\it 同伦},
{\it 位于同一同伦类}


\noindent
见 \ref{simplicial-definition-homotopy-equivalent}：
{\it 同伦等价},
{\it 同伦等价}


\noindent
见 \ref{simplicial-definition-homotopy-cosimplicial}：
{\it 从 $a$ 到 $b$ 的同伦},
{\it 同伦},
{\it 位于同一同伦类}


\noindent
见 \ref{simplicial-definition-trivial-kan}：
{\it 平凡 Kan 纤维化}


\noindent
见 \ref{simplicial-definition-kan}：
{\it Kan 纤维化},
{\it Kan 复形}


\medskip\noindent
{\bf 代数进阶}

\medskip

\noindent
见 \ref{more-algebra-definition-stably-free}：
{\it 稳定同构},
{\it 稳定自由}


\noindent
见 \ref{more-algebra-definition-fitting-ideal}：
{\it 第 $k$ 个 Fitting 理想}


\noindent
见 \ref{more-algebra-definition-zariski-pair}：
{\it Zariski 对}


\noindent
见 \ref{more-algebra-definition-henselian-pair}：
{\it Hensel 对}


\noindent
见 \ref{more-algebra-definition-absolutely-integrally-closed}：
{\it 绝对整闭}


\noindent
见 \ref{more-algebra-definition-auto-ass}：
{\it 自伴随}


\noindent
见 \ref{more-algebra-definition-torsion}：
{\it 挠的},
{\it 无挠},
{\it 挠模}


\noindent
见 \ref{more-algebra-definition-rank}：
{\it 秩}


\noindent
见 \ref{more-algebra-definition-reflexive}：
{\it 反身}


\noindent
见 \ref{more-algebra-definition-reflexive-hull}：
{\it 反身包}


\noindent
见 \ref{more-algebra-definition-content-ideal}：
{\it $x$ 的内容理想}


\noindent
见 \ref{more-algebra-definition-strict-transform}：
{\it $M$ 沿 $R \to R[\frac{I}{a}]$ 的严格变换}


\noindent
见 \ref{more-algebra-definition-koszul}：
{\it Koszul 复形}


\noindent
见 \ref{more-algebra-definition-koszul-complex}：
{\it $f_1, \ldots, f_r$ 上的 Koszul 复形}


\noindent
见 \ref{more-algebra-definition-koszul-regular-sequence}：
{\it $M$-Koszul-正则},
{\it $M$-$H_1$-正则},
{\it Koszul-正则},
{\it $H_1$-正则}


\noindent
见 \ref{more-algebra-definition-regular-ideal}：
{\it 正则理想},
{\it Koszul-正则理想},
{\it $H_1$-正则理想},
{\it 拟正则理想}


\noindent
见 \ref{more-algebra-definition-local-complete-intersection}：
{\it 局部完全交}


\noindent
见 \ref{more-algebra-definition-topological-ring}：
{\it 拓扑环},
{\it 拓扑模},
{\it 拓扑模同态},
{\it 拓扑环同态},
{\it 线性拓扑化},
{\it 线性拓扑化},
{\it 定义理想},
{\it 预可容许},
{\it 容许的},
{\it 预进},
{\it adic}


\noindent
见 \ref{more-algebra-definition-formally-smooth}：
{\it 在 $R$ 上形式光滑}


\noindent
见 \ref{more-algebra-definition-formally-smooth-adic}：
{\it 对 $\mathfrak n$-进拓扑形式光滑}


\noindent
见 \ref{more-algebra-definition-regular}：
{\it 正则概形}


\noindent
见 \ref{more-algebra-definition-p-basis}：
{\it 在 $k$ 上 p-无关},
{\it $K$ 在 $k$ 上的 p-基}


\noindent
见 \ref{more-algebra-definition-J}：
{\it J-0},
{\it J-1},
{\it J-2 的}


\noindent
见 \ref{more-algebra-definition-G-ring}：
{\it G-环}


\noindent
见 \ref{more-algebra-definition-excellent}：
{\it 拟优秀概形},
{\it 优秀概形}


\noindent
见 \ref{more-algebra-definition-projective}：
{\it 单射的}


\noindent
见 \ref{more-algebra-definition-simple-functors}：
{\it $M \mapsto M^\vee$},
{\it 自由模}


\noindent
见 \ref{more-algebra-definition-K-flat}：
{\it K-平坦}


\noindent
见 \ref{more-algebra-definition-derived-tor}：
{\it 导出张量积}


\noindent
见 \ref{more-algebra-definition-tor-independent}：
{\it 在 $R$ 上 Tor 无关}


\noindent
见 \ref{more-algebra-definition-pseudo-coherent}：
{\it $m$-伪凝聚},
{\it 伪凝聚的},
{\it $m$-伪凝聚},
{\it 伪凝聚的}


\noindent
见 \ref{more-algebra-definition-tor-amplitude}：
{\it Tor 振幅位于 $[a, b]$},
{\it 有限 Tor 维数},
{\it Tor 维数 $\leq d$},
{\it 有限 Tor 维数}


\noindent
见 \ref{more-algebra-definition-projective-dimension}：
{\it 有限投射维数},
{\it 投射振幅位于 $[a, b]$}


\noindent
见 \ref{more-algebra-definition-injective-dimension}：
{\it 有限内射维数},
{\it 内射振幅位于 $[a, b]$}


\noindent
见 \ref{more-algebra-definition-near-projective}：
{\it $I$-投射}


\noindent
见 \ref{more-algebra-definition-perfect}：
{\it 完美的},
{\it 完美的}


\noindent
见 \ref{more-algebra-definition-relatively-finitely-presented}：
{\it 相对于 $R$ 有限表示}


\noindent
见 \ref{more-algebra-definition-relatively-pseudo-coherent}：
{\it 相对于 $R$ 为 $m$-伪凝聚},
{\it 相对于 $R$ 伪凝聚},
{\it 相对于 $R$ 为 $m$-伪凝聚},
{\it 相对于 $R$ 伪凝聚}


\noindent
见 \ref{more-algebra-definition-pseudo-coherent-perfect}：
{\it 伪凝聚环映射},
{\it 完美环映射}


\noindent
见 \ref{more-algebra-definition-relatively-perfect}：
{\it $R$-完美},
{\it 相对于 $R$ 完美}


\noindent
见 \ref{more-algebra-definition-f-power-torsion}：
{\it $I$-幂挠模},
{\it $f$-幂挠模}


\noindent
见 \ref{more-algebra-definition-derived-complete}：
{\it 关于 $I$ 导出完备},
{\it 关于 $I$ 导出完备}


\noindent
见 \ref{more-algebra-definition-weakly-etale}：
{\it 绝对平坦的},
{\it 弱 \'etale 的},
{\it 绝对平坦的}


\noindent
见 \ref{more-algebra-definition-weak-dimension}：
{\it 弱维数 $\leq d$}


\noindent
见 \ref{more-algebra-definition-unibranch}：
{\it 单分支},
{\it 几何单枝}


\noindent
见 \ref{more-algebra-definition-number-of-branches}：
{\it $A$ 的分支数},
{\it $A$ 的几何分支数}


\noindent
见 \ref{more-algebra-definition-formally-catenary}：
{\it 形式链状}


\noindent
见 \ref{more-algebra-definition-extension-discrete-valuation-rings}：
{\it 离散赋值环的扩张},
{\it 分歧指数},
{\it 弱无分歧},
{\it 残余次数},
{\it 剩余次数}


\noindent
见 \ref{more-algebra-definition-types-of-extensions}：
{\it 相对于 $A$ 无分歧},
{\it 相对于 $A$ 驯分歧},
{\it 相对于 $A$ 全分歧}


\noindent
见 \ref{more-algebra-definition-decomposition-inertia}：
{\it $\mathfrak m$ 的分解群},
{\it $\mathfrak m$ 的惯性群}


\noindent
见 \ref{more-algebra-definition-wild-inertia}：
{\it $\mathfrak m$ 的野惯性群},
{\it $\mathfrak m$ 的驯惯性群}


\noindent
见 \ref{more-algebra-definition-mixed}：
{\it 混合特征},
{\it 绝对分歧指数}


\noindent
见 \ref{more-algebra-definition-solution}：
{\it $A \subset B$ 的弱解},
{\it $A \subset B$ 的解},
{\it 可分解}


\noindent
见 \ref{more-algebra-definition-invertible}：
{\it 可逆的},
{\it 平凡}


\noindent
见 \ref{more-algebra-definition-extension-valuation-rings}：
{\it 赋值环的扩张},
{\it 弱无分歧},
{\it 残余次数},
{\it 剩余次数}


\noindent
见 \ref{more-algebra-definition-bezout}：
{\it B\'ezout 整环},
{\it 初等因子整环}


\medskip\noindent
{\bf 环同态的光滑化}

\medskip

\noindent
见 \ref{smoothing-definition-singular-ideal}：
{\it $A$ 在 $R$ 上的奇异理想}


\noindent
见 \ref{smoothing-definition-strictly-standard}：
{\it 在 $A$ 中相对于 $R$ 为初等标准},
{\it 在 $A$ 中相对于 $R$ 为严格标准}


\medskip\noindent
{\bf 模层}

\medskip

\noindent
见 \ref{modules-definition-globally-generated}：
{\it 由整体截面生成},
{\it 生成}


\noindent
见 \ref{modules-definition-generated-by-local-sections}：
{\it 由 $s_i$ 生成的子层}


\noindent
见 \ref{modules-definition-support}：
{\it $\mathcal{F}$ 的支集},
{\it $s$ 的支集}


\noindent
见 \ref{modules-definition-locally-generated}：
{\it 局部由截面生成}


\noindent
见 \ref{modules-definition-finite-type}：
{\it 有限型}


\noindent
见 \ref{modules-definition-quasi-coherent}：
{\it 拟凝聚 $\mathcal{O}_X$-模层}


\noindent
见 \ref{modules-definition-sheaf-associated}：
{\it 与模 $M$ 及环映射 $\alpha$ 关联的层},
{\it 与模 $M$ 关联的层}


\noindent
见 \ref{modules-definition-finite-presentation}：
{\it 有限表示}


\noindent
见 \ref{modules-definition-coherent}：
{\it 凝聚 $\mathcal{O}_X$-模}


\noindent
见 \ref{modules-definition-closed-immersion}：
{\it 环化空间的闭浸入}


\noindent
见 \ref{modules-definition-locally-free}：
{\it 局部自由},
{\it 有限局部自由},
{\it 秩为 $r$ 的有限局部自由}


\noindent
见 \ref{modules-definition-flat}：
{\it 平坦}


\noindent
见 \ref{modules-definition-flat-at-point}：
{\it 在 $x$ 处平坦}


\noindent
见 \ref{modules-definition-flat-morphism}：
{\it 在 $x$ 处平坦},
{\it 平坦}


\noindent
见 \ref{modules-definition-flat-module}：
{\it 在点 $x \in X$ 处相对于 $Y$ 平坦},
{\it 在 $Y$ 上平坦}


\noindent
见 \ref{modules-definition-annihilator-sheaf}：
{\it 零化子}


\noindent
见 \ref{modules-definition-koszul}：
{\it Koszul 复形}


\noindent
见 \ref{modules-definition-koszul-complex}：
{\it $f_1, \ldots, f_n$ 上的 Koszul 复形}


\noindent
见 \ref{modules-definition-invertible}：
{\it 可逆 $\mathcal{O}_X$-模},
{\it 平凡}


\noindent
见 \ref{modules-definition-powers}：
{\it 张量幂}


\noindent
见 \ref{modules-definition-gamma-star}：
{\it 伴随分次环}


\noindent
见 \ref{modules-definition-pic}：
{\it Picard 群}


\noindent
见 \ref{modules-definition-derivation}：
{\it $\mathcal{O}_1$-导子},
{\it $\varphi$-导子},
{\it Leibniz 法则}


\noindent
见 \ref{modules-definition-module-differentials}：
{\it 微分模},
{\it 泛 $\varphi$-导子}


\noindent
见 \ref{modules-definition-differentials}：
{\it $S$-导子},
{\it $X$ 在 $S$ 上的微分层 $\Omega_{X/S}$}


\noindent
见 \ref{modules-definition-differential-operators}：
{\it $k$ 阶微分算子 $D : \mathcal{F} \to \mathcal{G}$}


\noindent
见 \ref{modules-definition-module-principal-parts}：
{\it $k$ 阶主部件模}


\noindent
见 \ref{modules-definition-relative-differential-operators}：
{\it $X/S$ 上的 $k$ 阶微分算子}


\noindent
见 \ref{modules-definition-de-rham-complex}：
{\it $\mathcal{B}$ 相对于 $\mathcal{A}$ 的 de Rham 复形}


\noindent
见 \ref{modules-definition-de-rham-complex-morphism-ringed-spaces}：
{\it de Rham 复形}


\noindent
见 \ref{modules-definition-naive-cotangent-complex}：
{\it 朴素余切复形}


\noindent
见 \ref{modules-definition-cotangent-complex-morphism-ringed-topoi}：
{\it 朴素余切复形}


\medskip\noindent
{\bf 位点上的模}

\medskip

\noindent
见 \ref{sites-modules-definition-free-abelian-presheaf-on}：
{\it 自由阿贝尔预层}


\noindent
见 \ref{sites-modules-definition-free-abelian-sheaf-on}：
{\it 自由阿贝尔层}


\noindent
见 \ref{sites-modules-definition-ringed-site}：
{\it 环化位点},
{\it 结构层},
{\it 环化位点的态射},
{\it 环化位点态射的复合}


\noindent
见 \ref{sites-modules-definition-ringed-topos}：
{\it 环化拓扑斯},
{\it 结构层},
{\it 环化拓扑斯的态射},
{\it 环化拓扑斯态射的复合}


\noindent
见 \ref{sites-modules-definition-2-morphism-ringed-topoi}：
{\it 从 $f$ 到 $g$ 的 2-态射}


\noindent
见 \ref{sites-modules-definition-presheaf-modules}：
{\it $\mathcal{O}$-模预层},
{\it $\mathcal{O}$-模预层的态射 $\varphi : \mathcal{F} \to \mathcal{G}$}


\noindent
见 \ref{sites-modules-definition-sheaf-modules}：
{\it $\mathcal{O}$-模层},
{\it $\mathcal{O}$-模层态射}


\noindent
见 \ref{sites-modules-definition-pushforward}：
{\it 推前},
{\it 拉回}


\noindent
见 \ref{sites-modules-definition-g-shriek}：
{\it $g_{p!}\mathcal{F}$},
{\it $g_!\mathcal{F} = (g_{p!}\mathcal{F})^\#$}


\noindent
见 \ref{sites-modules-definition-global}：
{\it 自由 $\mathcal{O}$-模},
{\it 有限自由},
{\it 由整体截面生成},
{\it 由 $r$ 个整体截面生成},
{\it 由有限多个整体截面生成},
{\it 全局表示},
{\it 全局有限表示}


\noindent
见 \ref{sites-modules-definition-localize-ringed-site}：
{\it 环化位点 $(\mathcal{C}, \mathcal{O})$ 在对象 $U$ 处的局部化},
{\it 局部化态射},
{\it 直像函子},
{\it $\mathcal{F}$ 到 $\mathcal{C}/U$ 的限制},
{\it 零延拓}


\noindent
见 \ref{sites-modules-definition-localize-ringed-topos}：
{\it 环化拓扑斯 $(\Sh(\mathcal{C}), \mathcal{O})$ 在 $\mathcal{F}$ 处的局部化},
{\it 局部化态射}


\noindent
见 \ref{sites-modules-definition-site-local}：
{\it 局部自由},
{\it 有限局部自由},
{\it 局部由截面生成},
{\it 局部由 $r$ 个截面生成},
{\it 有限型的},
{\it 拟凝聚},
{\it 有限表示的},
{\it 凝聚}


\noindent
见 \ref{sites-modules-definition-flat}：
{\it 平坦},
{\it 平坦},
{\it 平坦},
{\it 平坦}


\noindent
见 \ref{sites-modules-definition-flat-morphism}：
{\it 平坦},
{\it 平坦}


\noindent
见 \ref{sites-modules-definition-flat-module}：
{\it 在 $(\Sh(\mathcal{D}), \mathcal{O}')$ 上平坦}


\noindent
见 \ref{sites-modules-definition-invertible-sheaf}：
{\it 秩 $r$},
{\it 可逆的},
{\it $\mathcal{O}^*$}


\noindent
见 \ref{sites-modules-definition-pic}：
{\it Picard 群}


\noindent
见 \ref{sites-modules-definition-derivation}：
{\it $\mathcal{O}_1$-导子},
{\it $\varphi$-导子},
{\it Leibniz 法则}


\noindent
见 \ref{sites-modules-definition-module-differentials}：
{\it 微分模},
{\it 泛 $\varphi$-导子}


\noindent
见 \ref{sites-modules-definition-sheaf-differentials}：
{\it $Y$-导子},
{\it $X$ 在 $Y$ 上的微分层 $\Omega_{X/Y}$},
{\it 泛 $Y$-导子}


\noindent
见 \ref{sites-modules-definition-differential-operators}：
{\it $k$ 阶微分算子 $D : \mathcal{F} \to \mathcal{G}$}


\noindent
见 \ref{sites-modules-definition-module-principal-parts}：
{\it $k$ 阶主部件模}


\noindent
见 \ref{sites-modules-definition-naive-cotangent-complex}：
{\it 朴素余切复形}


\noindent
见 \ref{sites-modules-definition-cotangent-complex-morphism-ringed-topoi}：
{\it 朴素余切复形}


\noindent
见 \ref{sites-modules-definition-locally-ringed}：
{\it 局部环化位点}


\noindent
见 \ref{sites-modules-definition-locally-ringed-topos}：
{\it 局部环化}


\noindent
见 \ref{sites-modules-definition-morphism-locally-ringed-topoi}：
{\it 局部环化拓扑斯的态射},
{\it 局部环化位点的态射}


\noindent
见 \ref{sites-modules-definition-locally-constant}：
{\it 常值层},
{\it 局部常值},
{\it 有限局部常值}


\medskip\noindent
{\bf 内射对象}

\medskip

\noindent
见 \ref{injectives-definition-small}：
{\it 相对于 $I$ 为 $\alpha$-小}


\noindent
见 \ref{injectives-definition-grothendieck-conditions}：
{\it 生成元},
{\it Grothendieck 阿贝尔范畴}


\noindent
见 \ref{injectives-definition-size}：
{\it 大小}


\medskip\noindent
{\bf 层上同调}

\medskip

\noindent
见 \ref{cohomology-definition-torsor}：
{\it 伪扭子},
{\it 伪 $\mathcal{G}$-扭子},
{\it 伪 $\mathcal{G}$-扭子的态射},
{\it 扭子},
{\it $\mathcal{G}$-扭子},
{\it $\mathcal{G}$-扭子的态射},
{\it 平凡 $\mathcal{G}$-扭子}


\noindent
见 \ref{cohomology-definition-cech-complex}：
{\it {\v C}ech 复形},
{\it {\v C}ech 上同调群}


\noindent
见 \ref{cohomology-definition-flasque}：
{\it 松弛},
{\it 松弛}


\noindent
见 \ref{cohomology-definition-alternating-cech-complex}：
{\it 交错 {\v C}ech 复形}


\noindent
见 \ref{cohomology-definition-ordered-cech-complex}：
{\it 有序 {\v C}ech 复形}


\noindent
见 \ref{cohomology-definition-covering-locally-finite}：
{\it 局部有限}


\noindent
见 \ref{cohomology-definition-K-flat}：
{\it K-平坦}


\noindent
见 \ref{cohomology-definition-derived-tor}：
{\it 导出张量积}


\noindent
见 \ref{cohomology-definition-tor}：
{\it Tor}


\noindent
见 \ref{cohomology-definition-strictly-perfect}：
{\it 严格完美}


\noindent
见 \ref{cohomology-definition-pseudo-coherent}：
{\it $m$-伪凝聚},
{\it 伪凝聚的},
{\it $m$-伪凝聚},
{\it 伪凝聚的}


\noindent
见 \ref{cohomology-definition-tor-amplitude}：
{\it Tor 振幅位于 $[a, b]$},
{\it 有限 Tor 维数},
{\it 局部具有有限 Tor 维数},
{\it Tor 维数 $\leq d$}


\noindent
见 \ref{cohomology-definition-perfect}：
{\it 完美的},
{\it 完美的}


\medskip\noindent
{\bf 位点上同调}

\medskip

\noindent
见 \ref{sites-cohomology-definition-torsor}：
{\it 伪扭子},
{\it 伪 $\mathcal{G}$-扭子},
{\it 伪 $\mathcal{G}$-扭子的态射},
{\it 扭子},
{\it $\mathcal{G}$-扭子},
{\it $\mathcal{G}$-扭子的态射},
{\it 平凡 $\mathcal{G}$-扭子}


\noindent
见 \ref{sites-cohomology-definition-cech-complex}：
{\it {\v C}ech 复形},
{\it {\v C}ech 上同调群}


\noindent
见 \ref{sites-cohomology-definition-limp}：
{\it 全无环}


\noindent
见 \ref{sites-cohomology-definition-K-flat}：
{\it K-平坦}


\noindent
见 \ref{sites-cohomology-definition-derived-tor}：
{\it 导出张量积}


\noindent
见 \ref{sites-cohomology-definition-tor}：
{\it Tor}


\noindent
见 \ref{sites-cohomology-definition-covering-LC}：
{\it qc 覆盖}


\noindent
见 \ref{sites-cohomology-definition-simplicial-module}：
{\it 单纯 $\mathcal{A}_\bullet$-模},
{\it 单纯 $\mathcal{A}_\bullet$-模层}


\noindent
见 \ref{sites-cohomology-definition-cartesian}：
{\it $\mathit{QC}(\mathcal{O})$}


\noindent
见 \ref{sites-cohomology-definition-strictly-perfect}：
{\it 严格完美}


\noindent
见 \ref{sites-cohomology-definition-pseudo-coherent}：
{\it $m$-伪凝聚},
{\it 伪凝聚的},
{\it $m$-伪凝聚},
{\it 伪凝聚的}


\noindent
见 \ref{sites-cohomology-definition-tor-amplitude}：
{\it Tor 振幅位于 $[a, b]$},
{\it 有限 Tor 维数},
{\it 局部具有有限 Tor 维数},
{\it Tor 维数 $\leq d$}


\noindent
见 \ref{sites-cohomology-definition-perfect}：
{\it 完美的},
{\it 完美的}


\medskip\noindent
{\bf 微分分次代数}

\medskip

\noindent
见 \ref{dga-definition-dga}：
{\it $R$ 上的微分分次代数}


\noindent
见 \ref{dga-definition-homomorphism-dga}：
{\it 微分分次代数同态}


\noindent
见 \ref{dga-definition-cdga}：
{\it 可交换},
{\it 严格可交换}


\noindent
见 \ref{dga-definition-tensor-product}：
{\it 张量积微分分次代数}


\noindent
见 \ref{dga-definition-dgm}：
{\it 微分分次模},
{\it 微分分次模同态}


\noindent
见 \ref{dga-definition-shift}：
{\it $k$-移位模}


\noindent
见 \ref{dga-definition-homotopy}：
{\it $f$ 与 $g$ 之间的同伦},
{\it 同伦}


\noindent
见 \ref{dga-definition-complexes-notation}：
{\it 同伦范畴}


\noindent
见 \ref{dga-definition-cone}：
{\it 锥}


\noindent
见 \ref{dga-definition-admissible-ses}：
{\it 可容许单射},
{\it 可容许满射},
{\it 可容许短正合列}


\noindent
见 \ref{dga-definition-distinguished-triangle}：
{\it 与 $0 \to K \to L \to M \to 0$ 关联的三角},
{\it 可区别三角}


\noindent
见 \ref{dga-definition-opposite-dga}：
{\it 反微分分次代数}


\noindent
见 \ref{dga-definition-shift-graded-module}：
{\it 第 $k$ 个移位 $A$-模},
{\it 第 $k$ 个移位 $A$-模}


\noindent
见 \ref{dga-definition-unbounded-derived-category}：
{\it $(A, \text{d})$ 的导出范畴}


\noindent
见 \ref{dga-definition-linear-category}：
{\it $R$-线性范畴 $\mathcal{A}$}


\noindent
见 \ref{dga-definition-functor-linear-categories}：
{\it $R$-线性范畴的函子},
{\it $R$-线性函子}


\noindent
见 \ref{dga-definition-graded-category}：
{\it $R$ 上的分次范畴 $\mathcal{A}$}


\noindent
见 \ref{dga-definition-functor-graded-categories}：
{\it $R$ 上分次范畴的函子},
{\it 分次函子}


\noindent
见 \ref{dga-definition-H0-of-graded-category}：
{\it $\mathcal{A}^0$}


\noindent
见 \ref{dga-definition-graded-direct-sum}：
{\it 分次直和}


\noindent
见 \ref{dga-definition-dga-category}：
{\it $R$ 上的微分分次范畴 $\mathcal{A}$}


\noindent
见 \ref{dga-definition-functor-dga-categories}：
{\it $R$ 上微分分次范畴的函子}


\noindent
见 \ref{dga-definition-homotopy-category-of-dga-category}：
{\it $\mathcal{A}$ 的复形范畴},
{\it $\mathcal{A}$ 的同伦范畴}


\noindent
见 \ref{dga-definition-dg-direct-sum}：
{\it 微分分次直和}


\noindent
见 \ref{dga-definition-bimodule}：
{\it $(A, B)$-双模},
{\it 分次 $(A, B)$-双模},
{\it 微分分次 $(A, B)$-双模}


\medskip\noindent
{\bf 除幂代数}

\medskip

\noindent
见 \ref{dpa-definition-divided-powers}：
{\it 除幂结构}


\noindent
见 \ref{dpa-definition-divided-power-ring}：
{\it 除幂环},
{\it 除幂环同态}


\noindent
见 \ref{dpa-definition-extends}：
{\it 延拓}


\noindent
见 \ref{dpa-definition-divided-powers-graded}：
{\it 除幂结构}


\noindent
见 \ref{dpa-definition-divided-powers-dga}：
{\it 与微分分次结构相容}


\noindent
见 \ref{dpa-definition-lci}：
{\it 完全交},
{\it 局部完全交}


\medskip\noindent
{\bf 微分分次层}

\medskip

\noindent
见 \ref{sdga-definition-ga}：
{\it 分次 $\mathcal{O}$-代数层},
{\it 分次代数层},
{\it 分次 $\mathcal{O}$-代数同态}


\noindent
见 \ref{sdga-definition-gm}：
{\it 分次 $\mathcal{A}$-模},
{\it 分次模},
{\it 分次 $\mathcal{A}$-模同态}


\noindent
见 \ref{sdga-definition-bimodule}：
{\it 分次 $(\mathcal{A}, \mathcal{B})$-双模},
{\it 分次 $(\mathcal{A}, \mathcal{B})$-双模同态}


\noindent
见 \ref{sdga-definition-dga}：
{\it 微分分次 $\mathcal{O}$-代数层},
{\it 微分分次代数层},
{\it 微分分次 $\mathcal{O}$-代数同态}


\noindent
见 \ref{sdga-definition-dgm}：
{\it 微分分次 $\mathcal{A}$-模},
{\it 微分分次模},
{\it 微分分次 $\mathcal{A}$-模同态}


\noindent
见 \ref{sdga-definition-dg-bimodule}：
{\it 微分分次 $(\mathcal{A}, \mathcal{B})$-双模},
{\it 微分分次 $(\mathcal{A}, \mathcal{B})$-双模同态}


\noindent
见 \ref{sdga-definition-homotopy}：
{\it $f$ 与 $g$ 之间的同伦},
{\it 同伦}


\noindent
见 \ref{sdga-definition-complexes-notation}：
{\it 同伦范畴}


\noindent
见 \ref{sdga-definition-cone}：
{\it 锥}


\noindent
见 \ref{sdga-definition-graded-injective}：
{\it 分次内射}


\noindent
见 \ref{sdga-definition-K-injective}：
{\it K-内射}


\noindent
见 \ref{sdga-definition-derived-category}：
{\it $(\mathcal{A}, \text{d})$ 的导出范畴}


\noindent
见 \ref{sdga-definition-pullback}：
{\it 导出张量积},
{\it 导出拉回}


\noindent
见 \ref{sdga-definition-pushforward}：
{\it 导出内 Hom},
{\it 导出推前}


\noindent
见 \ref{sdga-definition-cartesian}：
{\it $\mathit{QC}(\mathcal{A}, \text{d})$}


\medskip\noindent
{\bf 超覆盖}

\medskip

\noindent
见 \ref{hypercovering-definition-SR}：
{\it 半可表示对象},
{\it $X$ 上的半可表示对象}


\noindent
见 \ref{hypercovering-definition-SR-F}：
{\it 将半可表示对象对应到预层的函子}


\noindent
见 \ref{hypercovering-definition-covering-SR}：
{\it 覆盖},
{\it 覆盖}


\noindent
见 \ref{hypercovering-definition-hypercovering}：
{\it $X$ 的超覆盖}


\noindent
见 \ref{hypercovering-definition-homology}：
{\it $K$ 的同调}


\noindent
见 \ref{hypercovering-definition-hypercovering-variant}：
{\it $\mathcal{G}$ 的超覆盖},
{\it 超覆盖}


\medskip\noindent
{\bf 概形}

\medskip

\noindent
见 \ref{schemes-definition-locally-ringed-space}：
{\it 局部环层空间 $(X, \mathcal{O}_X)$},
{\it $X$ 在 $x$ 处的局部环},
{\it $X$ 在 $x$ 处的剩余域},
{\it 局部环层空间之间的态射}


\noindent
见 \ref{schemes-definition-immersion-locally-ringed-spaces}：
{\it 开浸入}


\noindent
见 \ref{schemes-definition-open-subspace}：
{\it $X$ 的与 $U$ 相伴的开子空间}


\noindent
见 \ref{schemes-definition-closed-immersion-locally-ringed-spaces}：
{\it 闭浸入}


\noindent
见 \ref{schemes-definition-closed-subspace}：
{\it $X$ 的与理想层 $\mathcal{I}$ 相伴的闭子空间}


\noindent
见 \ref{schemes-definition-standard-covering}：
{\it 标准开覆盖},
{\it 标准开覆盖}


\noindent
见 \ref{schemes-definition-structure-sheaf}：
{\it 结构层 $\mathcal{O}_{\Spec(R)}$，即 $R$ 的谱的结构层},
{\it 谱}


\noindent
见 \ref{schemes-definition-affine-scheme}：
{\it 仿射概形},
{\it 仿射概形之间的态射}


\noindent
见 \ref{schemes-definition-scheme}：
{\it 概形},
{\it 概形态射}


\noindent
见 \ref{schemes-definition-immersion}：
{\it 开浸入},
{\it 开子概形},
{\it 闭浸入},
{\it 闭子概形},
{\it 浸入},
{\it 局部闭浸入}


\noindent
见 \ref{schemes-definition-reduced}：
{\it 约化的}


\noindent
见 \ref{schemes-definition-reduced-induced-scheme}：
{\it $Z$ 上的概形结构},
{\it 约化诱导概形结构},
{\it $X$ 的既约化 $X_{red}$}


\noindent
见 \ref{schemes-definition-representable-functor}：
{\it 由概形表示},
{\it 可表示}


\noindent
见 \ref{schemes-definition-representable-by-open-immersions}：
{\it 满足 Zariski 拓扑的层性质},
{\it 子函子 $H \subset F$},
{\it 由开浸入表示},
{\it 覆盖 $F$}


\noindent
见 \ref{schemes-definition-fibre-product}：
{\it 纤维积}


\noindent
见 \ref{schemes-definition-inverse-image-closed-subscheme}：
{\it 逆像 $f^{-1}(Z)$，即闭子概形 $Z$ 的逆像}


\noindent
见 \ref{schemes-definition-base-change}：
{\it $S$ 上的概形},
{\it 结构态射},
{\it $R$ 上的概形},
{\it 态射 $f : X \to Y$，即 $S$ 上概形的态射},
{\it 基变换},
{\it 基变换},
{\it 基变换}


\noindent
见 \ref{schemes-definition-preserved-by-base-change}：
{\it 任意基变换下保持},
{\it 基变换下保持},
{\it 任意基变换下保持},
{\it 基变换下保持}


\noindent
见 \ref{schemes-definition-fibre}：
{\it 概形论纤维 $X_s$，即 $f$ 在 $s$ 上方的纤维},
{\it $f$ 在 $s$ 上方的纤维}


\noindent
见 \ref{schemes-definition-quasi-compact}：
{\it 拟紧}


\noindent
见 \ref{schemes-definition-universally-closed}：
{\it 普遍闭}


\noindent
见 \ref{schemes-definition-valuative-criterion}：
{\it 满足赋值判据的存在性部分},
{\it 满足赋值判据的唯一性部分}


\noindent
见 \ref{schemes-definition-separated}：
{\it 分离},
{\it 拟分离},
{\it 分离},
{\it 拟分离}


\noindent
见 \ref{schemes-definition-monomorphism}：
{\it 单态射}


\medskip\noindent
{\bf 概形的构造}

\medskip

\noindent
见 \ref{constructions-definition-relative-spec}：
{\it $\mathcal{A}$ 在 $S$ 上的相对谱},
{\it $\mathcal{A}$ 在 $S$ 上的谱}


\noindent
见 \ref{constructions-definition-affine-n-space}：
{\it 仿射 $n$ 维空间（在 $S$ 上）},
{\it 仿射 $n$ 维空间（在 $R$ 上）}


\noindent
见 \ref{constructions-definition-vector-bundle}：
{\it 与 $\mathcal{E}$ 关联的向量丛}


\noindent
见 \ref{constructions-definition-abstract-vector-bundle}：
{\it 向量丛 $\pi : V \to S$（在 $S$ 上）},
{\it $S$ 上向量丛的态射}


\noindent
见 \ref{constructions-definition-cone}：
{\it 与 $\mathcal{A}$ 关联的锥},
{\it 与 $\mathcal{A}$ 关联的仿射锥}


\noindent
见 \ref{constructions-definition-abstract-cone}：
{\it 锥 $\pi : C \to S$（在 $S$ 上）},
{\it 锥态射}


\noindent
见 \ref{constructions-definition-standard-covering}：
{\it 标准开覆盖}


\noindent
见 \ref{constructions-definition-structure-sheaf}：
{\it 结构层 $\mathcal{O}_{\text{Proj}(S)}$（关于 $S$）},
{\it 齐次谱}


\noindent
见 \ref{constructions-definition-twist}：
{\it $\text{Proj}(S)$ 的结构层扭转}


\noindent
见 \ref{constructions-definition-projective-space}：
{\it 射影 $n$ 维空间（在 $\mathbf{Z}$ 上）},
{\it 射影 $n$ 维空间（在 $S$ 上）},
{\it 射影 $n$ 维空间（在 $R$ 上）}


\noindent
见 \ref{constructions-definition-relative-proj}：
{\it $\mathcal{A}$ 在 $S$ 上的相对齐次谱},
{\it $\mathcal{A}$ 在 $S$ 上的齐次谱},
{\it $\mathcal{A}$ 在 $S$ 上的相对 Proj}


\noindent
见 \ref{constructions-definition-projective-bundle}：
{\it 与 $\mathcal{E}$ 相伴的射影丛},
{\it 结构层的扭曲}


\noindent
见 \ref{constructions-definition-grassmannian}：
{\it $\mathbf{Z}$ 上的格拉斯曼概形},
{\it $S$ 上的格拉斯曼概形},
{\it $R$ 上的格拉斯曼概形}


\medskip\noindent
{\bf 概形的性质}

\medskip

\noindent
见 \ref{properties-definition-integral}：
{\it 整}


\noindent
见 \ref{properties-definition-property-local}：
{\it 局部的}


\noindent
见 \ref{properties-definition-locally-P}：
{\it 局部具有 $P$ 性质}


\noindent
见 \ref{properties-definition-noetherian}：
{\it 局部 Noether},
{\it Noether}


\noindent
见 \ref{properties-definition-jacobson}：
{\it Jacobson 空间}


\noindent
见 \ref{properties-definition-normal}：
{\it 正规}


\noindent
见 \ref{properties-definition-Cohen-Macaulay}：
{\it Cohen--Macaulay}


\noindent
见 \ref{properties-definition-regular}：
{\it 正则概形},
{\it 非奇异概形}


\noindent
见 \ref{properties-definition-dimension}：
{\it 维数},
{\it $X$ 在 $x$ 处的维数}


\noindent
见 \ref{properties-definition-catenary}：
{\it 链状}


\noindent
见 \ref{properties-definition-Rk}：
{\it 在余维 $k$ 内正则},
{\it $(R_k)$},
{\it $(S_k)$}


\noindent
见 \ref{properties-definition-nagata}：
{\it Japanese 概形},
{\it 泛 Japanese 概形},
{\it Nagata 概形}


\noindent
见 \ref{properties-definition-G-scheme}：
{\it G-概形}


\noindent
见 \ref{properties-definition-singular-locus}：
{\it 正则轨迹},
{\it 奇异轨迹}


\noindent
见 \ref{properties-definition-unibranch}：
{\it 在 $x$ 处单分支},
{\it 在 $x$ 处几何单分支},
{\it 单分支},
{\it 几何单枝}


\noindent
见 \ref{properties-definition-number-of-branches}：
{\it $X$ 在 $x$ 处的分支数},
{\it $X$ 在 $x$ 处的几何分支数}


\noindent
见 \ref{properties-definition-quasi-affine}：
{\it 拟仿射}


\noindent
见 \ref{properties-definition-locally-projective}：
{\it 局部投射}


\noindent
见 \ref{properties-definition-kappa-generated}：
{\it $\kappa$-生成的}


\noindent
见 \ref{properties-definition-subsheaf-sections-annihilated-by-ideal}：
{\it 被 $\mathcal{I}$ 零化的截面子层}


\noindent
见 \ref{properties-definition-subsheaf-sections-supported-on-closed}：
{\it 支集位于 $T$ 上的截面子层}


\noindent
见 \ref{properties-definition-ample}：
{\it 丰沛}


\medskip\noindent
{\bf 概形的态射}

\medskip

\noindent
见 \ref{morphisms-definition-scheme-theoretic-intersection-union}：
{\it 概形论交},
{\it 概形论并}


\noindent
见 \ref{morphisms-definition-scheme-theoretic-support}：
{\it $\mathcal{F}$ 的概形论支集}


\noindent
见 \ref{morphisms-definition-scheme-theoretic-image}：
{\it 概形论像}


\noindent
见 \ref{morphisms-definition-scheme-theoretically-dense}：
{\it $U$ 在 $X$ 中的概形论闭包},
{\it 在 $X$ 中概形论稠密}


\noindent
见 \ref{morphisms-definition-dominant}：
{\it 支配}


\noindent
见 \ref{morphisms-definition-surjective}：
{\it 满}


\noindent
见 \ref{morphisms-definition-universally-injective}：
{\it 普遍单射},
{\it 根态射}


\noindent
见 \ref{morphisms-definition-affine}：
{\it 仿射}


\noindent
见 \ref{morphisms-definition-family-ample-invertible-modules}：
{\it $X$ 上的丰沛可逆模族}


\noindent
见 \ref{morphisms-definition-quasi-affine}：
{\it 拟仿射}


\noindent
见 \ref{morphisms-definition-property-local}：
{\it 局部的},
{\it 基变换下稳定},
{\it 复合下稳定}


\noindent
见 \ref{morphisms-definition-locally-P}：
{\it 局部为 $P$ 型}


\noindent
见 \ref{morphisms-definition-finite-type}：
{\it 在 $x \in X$ 处为有限型},
{\it 局部有限型},
{\it 有限型}


\noindent
见 \ref{morphisms-definition-finite-type-point}：
{\it 有限型点}


\noindent
见 \ref{morphisms-definition-universally-catenary}：
{\it 普遍链状}


\noindent
见 \ref{morphisms-definition-J}：
{\it J-2 的}


\noindent
见 \ref{morphisms-definition-excellent}：
{\it 拟优秀概形},
{\it 优秀概形}


\noindent
见 \ref{morphisms-definition-quasi-finite}：
{\it 在点 $x \in X$ 处拟有限},
{\it 局部拟有限},
{\it 拟有限}


\noindent
见 \ref{morphisms-definition-finite-presentation}：
{\it 在 $x \in X$ 处有限表示},
{\it 局部有限表示},
{\it 有限表示}


\noindent
见 \ref{morphisms-definition-open}：
{\it 开的},
{\it 普遍开}


\noindent
见 \ref{morphisms-definition-submersive}：
{\it 商映射},
{\it 普遍商映射}


\noindent
见 \ref{morphisms-definition-flat}：
{\it 在点 $x \in X$ 处平坦},
{\it 在 $S$ 上于点 $x \in X$ 处平坦},
{\it 平坦},
{\it 在 $S$ 上平坦}


\noindent
见 \ref{morphisms-definition-scheme-structure-connected-component}：
{\it $T$ 上的典范概形结构}


\noindent
见 \ref{morphisms-definition-relative-dimension-d}：
{\it 在 $x$ 处相对维数 $\leq d$},
{\it 相对维数 $\leq d$},
{\it 相对维数 $d$}


\noindent
见 \ref{morphisms-definition-syntomic}：
{\it 在 $x \in X$ 处 syntomic},
{\it syntomic},
{\it $k$ 上的局部完全交},
{\it 标准 syntomic}


\noindent
见 \ref{morphisms-definition-syntomic-relative-dimension}：
{\it 相对维数 $d$ 的 syntomic 态射}


\noindent
见 \ref{morphisms-definition-conormal-sheaf}：
{\it $Z$ 在 $X$ 中的余法层 $\mathcal{C}_{Z/X}$},
{\it $i$ 的余法层}


\noindent
见 \ref{morphisms-definition-sheaf-differentials}：
{\it $X$ 在 $S$ 上的微分层 $\Omega_{X/S}$},
{\it 泛 $S$-导子}


\noindent
见 \ref{morphisms-definition-smooth}：
{\it 在 $x \in X$ 处光滑},
{\it 光滑},
{\it 标准光滑}


\noindent
见 \ref{morphisms-definition-smooth-relative-dimension}：
{\it 相对维数 $d$ 的光滑态射}


\noindent
见 \ref{morphisms-definition-unramified}：
{\it 在 $x \in X$ 处非分歧},
{\it 在 $x \in X$ 处 G-非分歧},
{\it 无分歧},
{\it G-非分歧}


\noindent
见 \ref{morphisms-definition-etale}：
{\it 在 $x \in X$ 处平展},
{\it 平展},
{\it 标准平展的}


\noindent
见 \ref{morphisms-definition-relatively-ample}：
{\it 相对丰沛},
{\it $f$-相对丰沛},
{\it 在 $X/S$ 上丰沛},
{\it $f$-丰沛}


\noindent
见 \ref{morphisms-definition-very-ample}：
{\it 相对甚丰沛},
{\it $f$-相对甚丰沛},
{\it 在 $X/S$ 上甚丰沛},
{\it $f$-甚丰沛}


\noindent
见 \ref{morphisms-definition-quasi-projective}：
{\it 拟射影},
{\it H-拟射影},
{\it 局部拟射影}


\noindent
见 \ref{morphisms-definition-proper}：
{\it 固有}


\noindent
见 \ref{morphisms-definition-projective}：
{\it 投射},
{\it H-射影},
{\it 局部投射}


\noindent
见 \ref{morphisms-definition-integral}：
{\it 整},
{\it 有限}


\noindent
见 \ref{morphisms-definition-universal-homeomorphism}：
{\it 普遍同胚}


\noindent
见 \ref{morphisms-definition-seminormal-ring}：
{\it 半正规},
{\it 绝对弱正规}


\noindent
见 \ref{morphisms-definition-seminormal}：
{\it 半正规},
{\it 绝对弱正规}


\noindent
见 \ref{morphisms-definition-seminormalization}：
{\it 半正规化},
{\it 绝对弱正规化}


\noindent
见 \ref{morphisms-definition-finite-locally-free}：
{\it 有限局部自由},
{\it 秩},
{\it 次数}


\noindent
见 \ref{morphisms-definition-rational-map}：
{\it 等价},
{\it 从 $X$ 到 $Y$ 的有理映射},
{\it 从 $X$ 到 $Y$ 的 $S$-有理映射}


\noindent
见 \ref{morphisms-definition-rational-function}：
{\it $X$ 上的有理函数}


\noindent
见 \ref{morphisms-definition-ring-of-rational-functions}：
{\it $X$ 上的有理函数环}


\noindent
见 \ref{morphisms-definition-function-field}：
{\it 函数域},
{\it 有理函数域}


\noindent
见 \ref{morphisms-definition-domain-of-definition}：
{\it 在点 $x \in X$ 处有定义},
{\it 定义域}


\noindent
见 \ref{morphisms-definition-dominant-rational}：
{\it 支配}


\noindent
见 \ref{morphisms-definition-birational-schemes}：
{\it 双有理},
{\it $S$-双有理的}


\noindent
见 \ref{morphisms-definition-birational}：
{\it 双有理}


\noindent
见 \ref{morphisms-definition-degree}：
{\it $X$ 在 $Y$ 上的次数}


\noindent
见 \ref{morphisms-definition-modification}：
{\it $X$ 的改造}


\noindent
见 \ref{morphisms-definition-alteration}：
{\it $X$ 的改造}


\noindent
见 \ref{morphisms-definition-integral-closure}：
{\it $\mathcal{O}_X$ 在 $\mathcal{A}$ 中的整闭包}


\noindent
见 \ref{morphisms-definition-normalization-X-in-Y}：
{\it $X$ 在 $Y$ 中的正规化}


\noindent
见 \ref{morphisms-definition-normalization}：
{\it normalization}


\noindent
见 \ref{morphisms-definition-relative-seminormalization}：
{\it $X$ 在 $Y$ 中的半正规化},
{\it $X$ 在 $Y$ 中的弱正规化}


\noindent
见 \ref{morphisms-definition-weak-normalization}：
{\it 弱正规化}


\noindent
见 \ref{morphisms-definition-weakly-normal}：
{\it 弱正规}


\noindent
见 \ref{morphisms-definition-universally-bounded}：
{\it 界定 $f$ 的纤维次数},
{\it fibres of $f$ are universally bounded}


\medskip\noindent
{\bf 概形上同调}

\medskip

\noindent
见 \ref{coherent-definition-depth}：
{\it 某点处的深度 $k$},
{\it 某点处的深度 $k$},
{\it $(S_k)$},
{\it $(S_k)$}


\noindent
见 \ref{coherent-definition-Cohen-Macaulay}：
{\it Cohen--Macaulay}


\noindent
见 \ref{coherent-definition-proper-over-base}：
{\it $Z$ 在 $S$ 上固有}


\medskip\noindent
{\bf 除子}

\medskip

\noindent
见 \ref{divisors-definition-associated}：
{\it 相伴},
{\it $X$ 的伴随点}


\noindent
见 \ref{divisors-definition-embedded}：
{\it 嵌入伴随点},
{\it 嵌入点},
{\it 嵌入分量}


\noindent
见 \ref{divisors-definition-weakly-associated}：
{\it 弱相伴},
{\it $X$ 的弱伴随点}


\noindent
见 \ref{divisors-definition-relative-assassin}：
{\it $\mathcal{F}$ 在 $X$ 中相对于 $S$ 的相对伴随点集}


\noindent
见 \ref{divisors-definition-relative-weak-assassin}：
{\it $\mathcal{F}$ 在 $X$ 中相对于 $S$ 的相对弱伴随点集}


\noindent
见 \ref{divisors-definition-torsion}：
{\it 挠的},
{\it 无挠}


\noindent
见 \ref{divisors-definition-rank}：
{\it 秩}


\noindent
见 \ref{divisors-definition-reflexive}：
{\it 反身包},
{\it 反身}


\noindent
见 \ref{divisors-definition-effective-Cartier-divisor}：
{\it 局部主闭子概形},
{\it 有效 Cartier 除子}


\noindent
见 \ref{divisors-definition-sum-effective-Cartier-divisors}：
{\it 有效 Cartier 除子 $D_1$ 与 $D_2$ 的和}


\noindent
见 \ref{divisors-definition-pullback-effective-Cartier-divisor}：
{\it $D$ 沿 $f$ 的拉回有定义},
{\it 有效 Cartier 除子的拉回}


\noindent
见 \ref{divisors-definition-invertible-sheaf-effective-Cartier-divisor}：
{\it 与 $D$ 关联的可逆层 $\mathcal{O}_S(D)$},
{\it 典范截面}


\noindent
见 \ref{divisors-definition-regular-section}：
{\it 正则截面}


\noindent
见 \ref{divisors-definition-zero-scheme-s}：
{\it 零概形}


\noindent
见 \ref{divisors-definition-relative-effective-Cartier-divisor}：
{\it 相对有效 Cartier 除子}


\noindent
见 \ref{divisors-definition-conormal-sheaf}：
{\it $Z$ 在 $X$ 中的余法代数 $\mathcal{C}_{Z/X, *}$},
{\it $f$ 的余法代数}


\noindent
见 \ref{divisors-definition-normal-cone}：
{\it 法锥 $C_ZX$},
{\it 法丛}


\noindent
见 \ref{divisors-definition-regular-ideal-sheaf}：
{\it 正则概形},
{\it Koszul-正则},
{\it $H_1$-正则},
{\it 拟正则}


\noindent
见 \ref{divisors-definition-regular-immersion}：
{\it 正则浸入},
{\it Koszul-正则浸入},
{\it $H_1$-正则浸入},
{\it 拟正则浸入}


\noindent
见 \ref{divisors-definition-relative-H1-regular-immersion}：
{\it 相对拟正则浸入},
{\it 相对 $H_1$-正则浸入}


\noindent
见 \ref{divisors-definition-sheaf-meromorphic-functions}：
{\it $X$ 上的亚纯函数层},
{\it $\mathcal{K}_X$},
{\it 亚纯函数}


\noindent
见 \ref{divisors-definition-meromorphic-section}：
{\it $\mathcal{F}$ 的亚纯截面}


\noindent
见 \ref{divisors-definition-pullback-meromorphic-sections}：
{\it 对 $f$ 而言亚纯函数的拉回有定义}


\noindent
见 \ref{divisors-definition-regular-meromorphic-section}：
{\it 正则概形}


\noindent
见 \ref{divisors-definition-regular-meromorphic-ideal-denominators}：
{\it $s$ 的分母理想层}


\noindent
见 \ref{divisors-definition-Weil-divisor}：
{\it 素除子},
{\it Weil 除子}


\noindent
见 \ref{divisors-definition-order-vanishing}：
{\it $f$ 沿 $Z$ 的消失阶}


\noindent
见 \ref{divisors-definition-principal-divisor}：
{\it 与 $f$ 关联的主 Weil 除子}


\noindent
见 \ref{divisors-definition-class-group}：
{\it Weil 除子类群}


\noindent
见 \ref{divisors-definition-order-vanishing-meromorphic}：
{\it $s$ 沿 $Z$ 的消失阶}


\noindent
见 \ref{divisors-definition-divisor-invertible-sheaf}：
{\it 与 $s$ 关联的 Weil 除子},
{\it 与 $\mathcal{L}$ 关联的 Weil 除子类}


\noindent
见 \ref{divisors-definition-blow-up}：
{\it $X$ 沿 $Z$ 的吹起},
{\it $X$ 关于理想层 $\mathcal{I}$ 的吹起},
{\it 例外除子},
{\it 中心}


\noindent
见 \ref{divisors-definition-strict-transform}：
{\it 严格变换},
{\it 严格变换}


\noindent
见 \ref{divisors-definition-admissible-blowup}：
{\it $U$-可容许吹起}


\medskip\noindent
{\bf 概形的极限}

\medskip

\medskip\noindent
{\bf 代数簇}

\medskip

\noindent
见 \ref{varieties-definition-variety}：
{\it 簇}


\noindent
见 \ref{varieties-definition-geometrically-reduced}：
{\it 在 $x$ 处几何既约},
{\it 几何既约}


\noindent
见 \ref{varieties-definition-geometrically-connected}：
{\it 几何连通}


\noindent
见 \ref{varieties-definition-geometrically-irreducible}：
{\it 几何不可约}


\noindent
见 \ref{varieties-definition-geometrically-integral}：
{\it 在 $x$ 处几何逐点整},
{\it 几何逐点整},
{\it 几何整}


\noindent
见 \ref{varieties-definition-geometrically-normal}：
{\it 在 $x$ 处几何正规},
{\it 几何正规}


\noindent
见 \ref{varieties-definition-geometrically-regular}：
{\it 在 $x$ 处几何正则},
{\it 在 $k$ 上几何正则}


\noindent
见 \ref{varieties-definition-dual-numbers}：
{\it 对偶数}


\noindent
见 \ref{varieties-definition-tangent-space}：
{\it $X$ 在 $x$ 处相对于 $S$ 的切空间},
{\it 切向量}


\noindent
见 \ref{varieties-definition-algebraic-scheme}：
{\it 代数 $k$-概形},
{\it 局部代数 $k$-概形}


\noindent
见 \ref{varieties-definition-variety-type}：
{\it 仿射簇},
{\it 投射簇},
{\it 拟投射簇},
{\it 固有簇},
{\it 光滑簇}


\noindent
见 \ref{varieties-definition-euler-characteristic}：
{\it $\mathcal{F}$ 的 Euler 示性数}


\noindent
见 \ref{varieties-definition-regularity}：
{\it $m$-正则}


\noindent
见 \ref{varieties-definition-hilbert-polynomial}：
{\it Hilbert 多项式}


\noindent
见 \ref{varieties-definition-absolute-frobenius}：
{\it $X$ 的绝对 Frobenius}


\noindent
见 \ref{varieties-definition-relative-frobenius}：
{\it $X/S$ 的相对 Frobenius 态射}


\noindent
见 \ref{varieties-definition-delta-invariant-algebra}：
{\it $A$ 的 $\delta$-不变量}


\noindent
见 \ref{varieties-definition-delta-invariant}：
{\it $X$ 在 $x$ 处的 $\delta$-不变量}


\noindent
见 \ref{varieties-definition-wedge}：
{\it $A$ 是 $A_1, \ldots, A_n$ 的楔}


\noindent
见 \ref{varieties-definition-curve}：
{\it 曲线}


\noindent
见 \ref{varieties-definition-degree-invertible-sheaf}：
{\it 次数},
{\it 次数}


\noindent
见 \ref{varieties-definition-intersection-number}：
{\it 交数}


\noindent
见 \ref{varieties-definition-degree}：
{\it $Z$ 相对于 $\mathcal{L}$ 的次数}


\noindent
见 \ref{varieties-definition-embed-dim}：
{\it $X$ 在 $x$ 处的嵌入维数}


\noindent
见 \ref{varieties-definition-embedding-dimension}：
{\it $X/k$ 在 $x$ 处的嵌入维数}


\medskip\noindent
{\bf 概形上的拓扑}

\medskip

\noindent
见 \ref{topologies-definition-zariski-covering}：
{\it $T$ 的 Zariski 覆盖}


\noindent
见 \ref{topologies-definition-standard-Zariski}：
{\it 标准 Zariski 覆盖}


\noindent
见 \ref{topologies-definition-big-zariski-site}：
{\it 大 Zariski 位点}


\noindent
见 \ref{topologies-definition-big-small-Zariski}：
{\it $S$ 的大 Zariski 位点},
{\it $S$ 的小 Zariski 位点},
{\it $S$ 的大仿射 Zariski 位点},
{\it $S$ 的小仿射 Zariski 位点}


\noindent
见 \ref{topologies-definition-restriction-small-zariski}：
{\it 到小 Zariski 位点的限制}


\noindent
见 \ref{topologies-definition-etale-covering}：
{\it $T$ 的平展覆盖}


\noindent
见 \ref{topologies-definition-standard-etale}：
{\it 标准平展覆盖}


\noindent
见 \ref{topologies-definition-big-etale-site}：
{\it 大平展位点}


\noindent
见 \ref{topologies-definition-big-small-etale}：
{\it $S$ 的大平展位点},
{\it $S$ 的小平展位点},
{\it $S$ 的大仿射平展位点},
{\it $S$ 的小仿射平展位点}


\noindent
见 \ref{topologies-definition-restriction-small-etale}：
{\it 到小平展位点的限制}


\noindent
见 \ref{topologies-definition-smooth-covering}：
{\it $T$ 的光滑覆盖}


\noindent
见 \ref{topologies-definition-standard-smooth}：
{\it 标准光滑覆盖}


\noindent
见 \ref{topologies-definition-big-smooth-site}：
{\it 大光滑位点}


\noindent
见 \ref{topologies-definition-big-small-smooth}：
{\it $S$ 的大光滑位点},
{\it $S$ 的大仿射光滑位点}


\noindent
见 \ref{topologies-definition-syntomic-covering}：
{\it $T$ 的 syntomic 覆盖}


\noindent
见 \ref{topologies-definition-standard-syntomic}：
{\it 标准 syntomic 覆盖}


\noindent
见 \ref{topologies-definition-big-syntomic-site}：
{\it 大 syntomic 位点}


\noindent
见 \ref{topologies-definition-big-small-syntomic}：
{\it $S$ 的大 syntomic 位点},
{\it $S$ 的大仿射 syntomic 位点}


\noindent
见 \ref{topologies-definition-fppf-covering}：
{\it $T$ 的 fppf 覆盖}


\noindent
见 \ref{topologies-definition-standard-fppf}：
{\it 标准 fppf 覆盖}


\noindent
见 \ref{topologies-definition-big-fppf-site}：
{\it 大 fppf 位点}


\noindent
见 \ref{topologies-definition-big-small-fppf}：
{\it $S$ 的大 fppf 位点},
{\it $S$ 的大仿射 fppf 位点}


\noindent
见 \ref{topologies-definition-standard-ph-covering}：
{\it 标准 ph 覆盖}


\noindent
见 \ref{topologies-definition-ph-covering}：
{\it $T$ 的 ph 覆盖}


\noindent
见 \ref{topologies-definition-big-ph-site}：
{\it 大 ph 位点}


\noindent
见 \ref{topologies-definition-big-small-ph}：
{\it $S$ 的大 ph 位点},
{\it $S$ 的大仿射 ph 位点}


\noindent
见 \ref{topologies-definition-fpqc-covering}：
{\it $T$ 的 fpqc 覆盖}


\noindent
见 \ref{topologies-definition-standard-fpqc}：
{\it 标准 fpqc 覆盖}


\noindent
见 \ref{topologies-definition-sheaf-property-fpqc}：
{\it 对给定态射族满足层性质},
{\it 满足 fpqc 拓扑的层性质}


\noindent
见 \ref{topologies-definition-standard-V-covering}：
{\it 标准 V 覆盖}


\noindent
见 \ref{topologies-definition-V-covering}：
{\it $T$ 的 V 覆盖}


\noindent
见 \ref{topologies-definition-sheaf-property-V}：
{\it 满足 V 拓扑的层性质}


\medskip\noindent
{\bf 下降}

\medskip

\noindent
见 \ref{descent-definition-descent-datum-quasi-coherent}：
{\it 拟凝聚层下降数据 $(\mathcal{F}_i, \varphi_{ij})$},
{\it 余循环条件},
{\it 下降数据态射 $\psi : (\mathcal{F}_i, \varphi_{ij}) \to (\mathcal{F}'_i, \varphi'_{ij})$}


\noindent
见 \ref{descent-definition-descent-datum-effective-quasi-coherent}：
{\it 平凡下降数据},
{\it 典范下降数据},
{\it 有效}


\noindent
见 \ref{descent-definition-descent-datum-modules}：
{\it 模相对于 $R \to A$ 的下降数据 $(N, \varphi)$},
{\it 余循环条件},
{\it 下降数据态射 $(N, \varphi) \to (N', \varphi')$}


\noindent
见 \ref{descent-definition-descent-datum-effective-module}：
{\it 有效}


\noindent
见 \ref{descent-definition-split-equalizer}：
{\it 可裂等化子}


\noindent
见 \ref{descent-definition-universally-injective}：
{\it 普遍单射}


\noindent
见 \ref{descent-definition-C}：
{\it $C$}


\noindent
见 \ref{descent-definition-effective-descent}：
{\it 沿 $f$ 的底扩张},
{\it 模的下降态射},
{\it 模的有效下降态射}


\noindent
见 \ref{descent-definition-pushforward}：
{\it $f_*$}


\noindent
见 \ref{descent-definition-structure-sheaf}：
{\it 大位点 $(\Sch/S)_\tau$ 的结构层},
{\it 结构层},
{\it 与 $\mathcal{F}$ 相联系的 $\mathcal{O}$-模层},
{\it 与 $\mathcal{F}$ 相联系的 $\mathcal{O}$-模层}


\noindent
见 \ref{descent-definition-parasitic}：
{\it 寄生的},
{\it 关于 $\tau$-拓扑寄生的}


\noindent
见 \ref{descent-definition-property-local}：
{\it 在 $\tau$-拓扑中为局部的}


\noindent
见 \ref{descent-definition-germs}：
{\it $X$ 在 $x$ 处的芽},
{\it 芽态射},
{\it 芽态射的复合}


\noindent
见 \ref{descent-definition-etale-morphism-germs}：
{\it 平展},
{\it 光滑}


\noindent
见 \ref{descent-definition-local-at-point}：
{\it \'etale 局部},
{\it smooth 局部}


\noindent
见 \ref{descent-definition-property-morphisms-local}：
{\it 在基底上为 $\tau$-局部的},
{\it 在目标上为 $\tau$-局部的},
{\it 关于 $\tau$-拓扑在基底上为局部的}


\noindent
见 \ref{descent-definition-property-morphisms-local-source}：
{\it 在源上 $\tau$ 局部的},
{\it 对 $\tau$-拓扑在源上局部}


\noindent
见 \ref{descent-definition-local-source-target}：
{\it 在源与目标上为 \'etale 局部的}


\noindent
见 \ref{descent-definition-local-source-target-at-point}：
{\it 在源与目标上为 \'etale 局部的}


\noindent
见 \ref{descent-definition-descent-datum}：
{\it $V/X/S$ 的下降数据},
{\it 余循环条件},
{\it 相对于 $X \to S$ 的下降数据},
{\it 相对于 $X \to S$ 的下降数据态射 $g : (V/X, \varphi) \to (V'/X, \varphi')$}


\noindent
见 \ref{descent-definition-descent-datum-for-family-of-morphisms}：
{\it 相对于态射族 $\{X_i \to S\}$ 的下降数据 $(V_i, \varphi_{ij})$},
{\it 态射 $\psi : (V_i, \varphi_{ij}) \to (V'_i, \varphi'_{ij})$}


\noindent
见 \ref{descent-definition-pullback-functor}：
{\it 拉回函子}


\noindent
见 \ref{descent-definition-pullback-functor-family}：
{\it 拉回函子}


\noindent
见 \ref{descent-definition-effective}：
{\it 平凡下降数据},
{\it 典范下降数据},
{\it 有效}


\noindent
见 \ref{descent-definition-effective-family}：
{\it 典范下降数据},
{\it 有效}


\noindent
见 \ref{descent-definition-descending-types-morphisms}：
{\it $\mathcal{P}$ 型态射对 $\tau$-覆盖满足下降}


\medskip\noindent
{\bf 概形的导出范畴}

\medskip

\noindent
见 \ref{perfect-definition-supported-on}：
{\it 支集在 $T$ 上}


\noindent
见 \ref{perfect-definition-approximation-holds}：
{\it 三元组满足逼近性质}


\noindent
见 \ref{perfect-definition-approximation}：
{\it 用完美复形逼近的性质成立}


\noindent
见 \ref{perfect-definition-tor-independent}：
{\it 在 $S$ 上 Tor 独立}


\noindent
见 \ref{perfect-definition-relatively-perfect}：
{\it 相对于 $S$ 完美},
{\it $S$-完美的}


\noindent
见 \ref{perfect-definition-resolution-property}：
{\it 消解性质}


\noindent
见 \ref{perfect-definition-K-group}：
{\it $X$ 的 Grothendieck 群},
{\it $X$ 上凝聚层的 Grothendieck 群}


\medskip\noindent
{\bf 态射进阶}

\medskip

\noindent
见 \ref{more-morphisms-definition-thickening}：
{\it 加厚},
{\it 一阶加厚},
{\it 有限阶加厚},
{\it $n$ 阶加厚},
{\it 加厚的态射},
{\it $S$ 上的加厚},
{\it $S$ 上加厚的态射}


\noindent
见 \ref{more-morphisms-definition-first-order-infinitesimal-neighbourhood}：
{\it 一阶无穷小邻域},
{\it $n$ 阶无穷小邻域}


\noindent
见 \ref{more-morphisms-definition-formally-unramified}：
{\it 形式非分歧的}


\noindent
见 \ref{more-morphisms-definition-universal-thickening}：
{\it 泛一阶加厚},
{\it $Z$ 在 $X$ 上的余法层}


\noindent
见 \ref{more-morphisms-definition-formally-etale}：
{\it 形式平展的}


\noindent
见 \ref{more-morphisms-definition-formally-smooth}：
{\it 形式光滑的}


\noindent
见 \ref{more-morphisms-definition-netherlander}：
{\it $f$ 的朴素余切复形}


\noindent
见 \ref{more-morphisms-definition-normal}：
{\it 在 $x$ 处正规},
{\it 正规态射}


\noindent
见 \ref{more-morphisms-definition-regular}：
{\it 在 $x$ 处正则},
{\it 正则态射}


\noindent
见 \ref{more-morphisms-definition-CM}：
{\it 在 $x$ 处 Cohen--Macaulay},
{\it Cohen--Macaulay 态射}


\noindent
见 \ref{more-morphisms-definition-etale-neighbourhood}：
{\it $(S, s)$ 的平展邻域},
{\it 态射},
{\it 初等平展邻域}


\noindent
见 \ref{more-morphisms-definition-relatively-finitely-presented-sheaf}：
{\it 相对于 $S$ 有限表现},
{\it 在 $S$ 上相对有限表现}


\noindent
见 \ref{more-morphisms-definition-relative-pseudo-coherence}：
{\it 相对于 $S$ 是 $m$-伪凝聚的},
{\it 相对于 $S$ 是伪凝聚的},
{\it 相对于 $S$ 是 $m$-伪凝聚的},
{\it 相对于 $S$ 是伪凝聚的}


\noindent
见 \ref{more-morphisms-definition-pseudo-coherent}：
{\it 伪凝聚的}


\noindent
见 \ref{more-morphisms-definition-perfect}：
{\it 完美的}


\noindent
见 \ref{more-morphisms-definition-lci}：
{\it 在 $x$ 处为 Koszul},
{\it Koszul 态射},
{\it 局部完全交态射}


\noindent
见 \ref{more-morphisms-definition-weakly-etale}：
{\it 弱 \'etale 的},
{\it 绝对平坦的}


\noindent
见 \ref{more-morphisms-definition-ind-quasi-affine}：
{\it ind-拟仿射的},
{\it ind-拟仿射的}


\noindent
见 \ref{more-morphisms-definition-affine-stratification}：
{\it 仿射分层}


\noindent
见 \ref{more-morphisms-definition-affine-stratification-number}：
{\it 仿射分层数}


\noindent
见 \ref{more-morphisms-definition-weighting}：
{\it 赋权},
{\it pond\'eration}


\noindent
见 \ref{more-morphisms-definition-cd-morphism}：
{\it 完全分解},
{\it 完全分解}


\medskip\noindent
{\bf 平坦性进阶}

\medskip

\noindent
见 \ref{flat-definition-one-step-devissage}：
{\it $s$ 上 $\mathcal{F}/X/S$ 的一步 d\'evissage}


\noindent
见 \ref{flat-definition-one-step-devissage-at-x}：
{\it $x$ 处 $\mathcal{F}/X/S$ 的一步 d\'evissage}


\noindent
见 \ref{flat-definition-shrink}：
{\it 标准缩小}


\noindent
见 \ref{flat-definition-complete-devissage}：
{\it $s$ 上 $\mathcal{F}/X/S$ 的完全 d\'evissage}


\noindent
见 \ref{flat-definition-complete-devissage-at-x}：
{\it $x$ 处 $\mathcal{F}/X/S$ 的完全 d\'evissage}


\noindent
见 \ref{flat-definition-shrink-complete}：
{\it 标准缩小}


\noindent
见 \ref{flat-definition-elementary-etale-neighbourhood}：
{\it 环同态 $R \to S$ 在 $\mathfrak q$ 处的初等 \'etale 局部化}


\noindent
见 \ref{flat-definition-complete-devissage-algebra}：
{\it $\mathfrak r$ 上 $N/S/R$ 的完全 d\'evissage}


\noindent
见 \ref{flat-definition-complete-devissage-at-x-algebra}：
{\it $N/S/R$ 在 $\mathfrak q$ 处的完全 d\'evissage}


\noindent
见 \ref{flat-definition-impurity}：
{\it $\mathcal{F}$ 在 $s$ 上方的不纯度}


\noindent
见 \ref{flat-definition-pure}：
{\it 沿 $X_s$ 纯},
{\it 沿 $X_s$ 普遍纯},
{\it 沿 $X_s$ 纯},
{\it 普遍 $S$-纯},
{\it 相对于 $S$ 普遍纯},
{\it $S$-纯},
{\it 相对于 $S$ 纯},
{\it $S$-纯},
{\it 相对于 $S$ 纯}


\noindent
见 \ref{flat-definition-flat-dimension-n}：
{\it $\mathcal{F}$ 在维数 $\geq n$ 上关于 $S$ 平坦}


\noindent
见 \ref{flat-definition-flattening}：
{\it $\mathcal{F}$ 的泛扁平化存在},
{\it $X$ 的泛扁平化存在}


\noindent
见 \ref{flat-definition-flattening-stratification}：
{\it 扁平化分层},
{\it 扁平化分层}


\noindent
见 \ref{flat-definition-h-covering}：
{\it $T$ 的 h 覆盖}


\noindent
见 \ref{flat-definition-big-h-site}：
{\it 大 h 站点}


\noindent
见 \ref{flat-definition-standard-h}：
{\it 标准 h 覆盖}


\noindent
见 \ref{flat-definition-big-small-h}：
{\it $S$ 的大 h 位点},
{\it $S$ 的大仿射 h 位点}


\medskip\noindent
{\bf 群胚概形}

\medskip

\noindent
见 \ref{groupoids-definition-equivalence-relation}：
{\it 预关系},
{\it 关系},
{\it 预等价关系},
{\it $U$ 在 $S$ 上的一个等价关系}


\noindent
见 \ref{groupoids-definition-restrict-relation}：
{\it 限制},
{\it 拉回}


\noindent
见 \ref{groupoids-definition-group-scheme}：
{\it $S$ 上的群概形},
{\it $S$ 上群概形的态射 $\psi : (G, m) \to (G', m')$}


\noindent
见 \ref{groupoids-definition-closed-subgroup-scheme}：
{\it 闭子群概形},
{\it 开子群概形}


\noindent
见 \ref{groupoids-definition-smooth-group-scheme}：
{\it 光滑群概形},
{\it 平坦群概形},
{\it 分离群概形}


\noindent
见 \ref{groupoids-definition-abelian-variety}：
{\it 阿贝尔簇}


\noindent
见 \ref{groupoids-definition-action-group-scheme}：
{\it $G$ 在概形 $X/S$ 上的作用},
{\it 等变},
{\it $G$-等变态射}


\noindent
见 \ref{groupoids-definition-free-action}：
{\it 自由的}


\noindent
见 \ref{groupoids-definition-pseudo-torsor}：
{\it 伪 $G$-扭子},
{\it 在 $G$ 下形式主齐},
{\it 平凡}


\noindent
见 \ref{groupoids-definition-principal-homogeneous-space}：
{\it 主齐性空间},
{\it $G$-扭子},
{\it $\tau$ 拓扑中的 $G$-扭子},
{\it $\tau$ $G$-扭子},
{\it $\tau$ 扭子},
{\it 拟等平凡},
{\it 局部平凡}


\noindent
见 \ref{groupoids-definition-equivariant-module}：
{\it $G$-等变拟凝聚 $\mathcal{O}_X$-模},
{\it 等变拟凝聚 $\mathcal{O}_X$-模}


\noindent
见 \ref{groupoids-definition-groupoid}：
{\it $S$ 上的群胚概形},
{\it $S$ 上的群胚},
{\it $S$ 上群胚概形的态射 $f : (U, R, s, t, c) \to (U', R', s', t', c')$}


\noindent
见 \ref{groupoids-definition-groupoid-module}：
{\it $(U, R, s, t, c)$ 上的拟凝聚模}


\noindent
见 \ref{groupoids-definition-stabilizer-groupoid}：
{\it 群胚概形 $(U, R, s, t, c)$ 的稳定子}


\noindent
见 \ref{groupoids-definition-restrict-groupoid}：
{\it $(U, R, s, t, c)$ 到 $U'$ 的限制}


\noindent
见 \ref{groupoids-definition-invariant-open}：
{\it 在集合论意义下 $R$-不变},
{\it $R$-不变的},
{\it $R$-不变的},
{\it $R$-不变的}


\noindent
见 \ref{groupoids-definition-quotient-sheaf}：
{\it 商层 $U/R$}


\noindent
见 \ref{groupoids-definition-representable-quotient}：
{\it 可表商},
{\it 可表商}


\noindent
见 \ref{groupoids-definition-cartesian-morphism}：
{\it 笛卡尔的},
{\it $(U', R', s', t', c')$ 在 $(U, R, s, t, c)$ 上是笛卡尔的},
{\it 在 $(U, R, s, t, c)$ 上笛卡尔的群胚概形态射}


\medskip\noindent
{\bf 群胚概形进阶}

\medskip

\medskip\noindent
{\bf 概形的平展态射}

\medskip

\noindent
见 \ref{etale-definition-unramified-rings}：
{\it 局部环的非分歧同态}


\noindent
见 \ref{etale-definition-unramified-schemes}：
{\it 在 $x$ 处非分歧},
{\it 无分歧}


\noindent
见 \ref{etale-definition-flat-rings}：
{\it 平坦},
{\it 忠实平坦},
{\it 平坦（相应地，忠实平坦）}


\noindent
见 \ref{etale-definition-flat-schemes}：
{\it 在 $x \in X$ 处关于 $Y$ 平坦},
{\it 在 $x \in X$ 处平坦},
{\it 平坦},
{\it 忠实平坦}


\noindent
见 \ref{etale-definition-etale-ring}：
{\it 局部环的平展（\'etale）同态}


\noindent
见 \ref{etale-definition-etale-schemes-1}：
{\it 在 $x \in X$ 处平展},
{\it 平展}


\noindent
见 \ref{etale-definition-strict-normal-crossings}：
{\it 严格正规交叉除子}


\noindent
见 \ref{etale-definition-normal-crossings}：
{\it 正规交叉除子}


\medskip\noindent
{\bf Chow 同调}

\medskip

\noindent
见 \ref{chow-definition-periodic-complex}：
{\it $2$-周期复形},
{\it 上同调模},
{\it 正合},
{\it $(2, 1)$-周期复形},
{\it 上同调模}


\noindent
见 \ref{chow-definition-periodic-length}：
{\it 重数},
{\it （加性）Herbrand 商}


\noindent
见 \ref{chow-definition-delta-dimension}：
{\it $Z$ 的 $\delta$-维数}


\noindent
见 \ref{chow-definition-cycles}：
{\it $X$ 上的循环},
{\it $k$ 维循环}


\noindent
见 \ref{chow-definition-support-cycle}：
{\it 支撑}


\noindent
见 \ref{chow-definition-effective-cycle}：
{\it 有效}


\noindent
见 \ref{chow-definition-cycle-associated-to-closed-subscheme}：
{\it $Z'$ 在 $Z$ 中的重数},
{\it 与 $Z$ 相伴的 $k$-循环}


\noindent
见 \ref{chow-definition-cycle-associated-to-coherent-sheaf}：
{\it $Z'$ 在 $\mathcal{F}$ 中的重数},
{\it 与 $\mathcal{F}$ 关联的 $k$ 维循环}


\noindent
见 \ref{chow-definition-proper-pushforward}：
{\it 推前}


\noindent
见 \ref{chow-definition-flat-pullback}：
{\it $\alpha$ 沿 $f$ 的平坦拉回}


\noindent
见 \ref{chow-definition-principal-divisor}：
{\it 与 $f$ 关联的主除子}


\noindent
见 \ref{chow-definition-rational-equivalence}：
{\it 有理等价于零},
{\it 有理等价},
{\it $X$ 上 $k$ 维循环的 Chow 群},
{\it $X$ 上模有理等价的 $k$ 维循环 Chow 群}


\noindent
见 \ref{chow-definition-envelope}：
{\it 包络}


\noindent
见 \ref{chow-definition-divisor-invertible-sheaf}：
{\it 与 $s$ 关联的 Weil 除子},
{\it 与 $\mathcal{L}$ 关联的 Weil 除子}


\noindent
见 \ref{chow-definition-cap-c1}：
{\it 与 $\mathcal{L}$ 的第一陈类相交}


\noindent
见 \ref{chow-definition-gysin-homomorphism}：
{\it Gysin 同态}


\noindent
见 \ref{chow-definition-bivariant-class}：
{\it $f$ 的次数为 $p$ 的双变类 $c$}


\noindent
见 \ref{chow-definition-chow-cohomology}：
{\it Chow 上同调}


\noindent
见 \ref{chow-definition-first-chern-class}：
{\it 第一陈类}


\noindent
见 \ref{chow-definition-chern-classes}：
{\it $X$ 上 $\mathcal{E}$ 的陈类},
{\it $X$ 上 $\mathcal{E}$ 的全陈类}


\noindent
见 \ref{chow-definition-cap-chern-classes}：
{\it 与 $\mathcal{E}$ 的第 $j$ 陈类相交}


\noindent
见 \ref{chow-definition-chern-classes-final}：
{\it 第 $i$ 陈类},
{\it 全陈类}


\noindent
见 \ref{chow-definition-degree-zero-cycle}：
{\it 零维循环的次数}


\noindent
见 \ref{chow-definition-defined-on-perfect}：
{\it $E$ 的陈类已定义}


\noindent
见 \ref{chow-definition-localized-chern}：
{\it 局部化陈特征},
{\it 局部化第 $p$ 陈类}


\noindent
见 \ref{chow-definition-lci-gysin}：
{\it $f$ 的 Gysin 映射存在},
{\it Gysin 映射}


\noindent
见 \ref{chow-definition-determinant}：
{\it 容许的},
{\it 符号},
{\it 容许关系},
{\it 有限长度 $R$-模 $M$ 的行列式}


\noindent
见 \ref{chow-definition-periodic-determinant}：
{\it $(M, \varphi, \psi)$ 的行列式}


\noindent
见 \ref{chow-definition-symbol-M}：
{\it 与 $M, a, b$ 关联的符号}


\noindent
见 \ref{chow-definition-tame-symbol}：
{\it 驯符号}


\medskip\noindent
{\bf 相交理论}

\medskip

\noindent
见 \ref{intersection-definition-proper-intersection}：
{\it 正常相交},
{\it 正常相交}


\noindent
见 \ref{intersection-definition-multiplicity}：
{\it $M$ 关于定义理想 $I$ 的重数}


\medskip\noindent
{\bf 曲线的 Picard 概形}

\medskip

\noindent
见 \ref{pic-definition-picard-functor}：
{\it Picard 函子}


\noindent
见 \ref{pic-definition-genus}：
{\it 亏格}


\medskip\noindent
{\bf Weil 上同调理论}

\medskip

\noindent
见 \ref{weil-definition-chow-group-motives}：
{\it $M$ 的第 $i$ 个 Chow 群}


\noindent
见 \ref{weil-definition-weil-cohomology-theory-classical}：
{\it 经典 Weil 上同调理论}


\noindent
见 \ref{weil-definition-weil-cohomology-theory}：
{\it Weil 上同调理论}


\medskip\noindent
{\bf 适足模}

\medskip

\noindent
见 \ref{adequate-definition-module-valued-functor}：
{\it 模值函子},
{\it 模值函子之间的态射}


\noindent
见 \ref{adequate-definition-adequate-functor}：
{\it 适足},
{\it 线性适足的}


\noindent
见 \ref{adequate-definition-adequate}：
{\it 适足}


\noindent
见 \ref{adequate-definition-category-adequate-modules}：
{\it $\textit{Adeq}(\mathcal{O})$},
{\it $\textit{Adeq}((\Sch/S)_\tau, \mathcal{O})$},
{\it $\textit{Adeq}(S)$}


\noindent
见 \ref{adequate-definition-pure}：
{\it 纯射影},
{\it 纯内射}


\noindent
见 \ref{adequate-definition-pure-resolution}：
{\it 纯射影分解},
{\it 纯内射分解}


\noindent
见 \ref{adequate-definition-pure-ext}：
{\it 纯扩张模}


\medskip\noindent
{\bf 对偶化复形}

\medskip

\noindent
见 \ref{dualizing-definition-essential}：
{\it 本质的},
{\it 本质扩张},
{\it 本质的}


\noindent
见 \ref{dualizing-definition-projective-cover}：
{\it 射影覆盖},
{\it 射影包络}


\noindent
见 \ref{dualizing-definition-injective-hull}：
{\it 内射包}


\noindent
见 \ref{dualizing-definition-indecomposable}：
{\it 不可分解的}


\noindent
见 \ref{dualizing-definition-dualizing}：
{\it 对偶复形}


\noindent
见 \ref{dualizing-definition-gorenstein}：
{\it Gorenstein},
{\it Gorenstein}


\noindent
见 \ref{dualizing-definition-relative-dualizing-complex}：
{\it 相对对偶化复形}


\medskip\noindent
{\bf 概形的对偶性}

\medskip

\noindent
见 \ref{duality-definition-dualizing-scheme}：
{\it 对偶复形}


\noindent
见 \ref{duality-definition-good-dualizing}：
{\it dualizing complex normalized relative to $\omega_S^\bullet$}


\noindent
见 \ref{duality-definition-gorenstein}：
{\it Gorenstein}


\noindent
见 \ref{duality-definition-gorenstein-morphism}：
{\it Gorenstein at $x$},
{\it Gorenstein morphism}


\noindent
见 \ref{duality-definition-relative-dualizing-complex}：
{\it 相对对偶化复形}


\medskip\noindent
{\bf 判别式与不同理想}

\medskip

\noindent
见 \ref{discriminant-definition-trace-element}：
{\it 迹元素}


\noindent
见 \ref{discriminant-definition-kahler-different}：
{\it K\"ahler 不同理想}


\noindent
见 \ref{discriminant-definition-different}：
{\it 不同理想}


\medskip\noindent
{\bf de Rham 上同调}

\medskip

\noindent
见 \ref{derham-definition-hodge-filtration}：
{\it 霍奇滤过}


\noindent
见 \ref{derham-definition-local-product}：
{\it 对 $Y \subset X$（在 $S$ 上）定义了对数极点德拉姆复形}


\noindent
见 \ref{derham-definition-log-complex}：
{\it $Y \subset X$（在 $S$ 上）的对数极点德拉姆复形}


\medskip\noindent
{\bf 局部上同调}

\medskip

\noindent
见 \ref{local-cohomology-definition-cd}：
{\it $I$ 在 $A$ 中的上同调维数}


\noindent
见 \ref{local-cohomology-definition-depth-complex}：
{\it $I$-深度},
{\it 深度}


\medskip\noindent
{\bf 代数几何与形式几何}

\medskip

\noindent
见 \ref{algebraization-definition-derived-complete}：
{\it 关于 $\mathcal{I}$ 导出完备}


\noindent
见 \ref{algebraization-definition-algebraizable}：
{\it $(\mathcal{F}_n)$ 可延拓到 $X$}


\noindent
见 \ref{algebraization-definition-canonically-algebraizable}：
{\it $(\mathcal{F}_n)$ 可典范地延拓到 $X$}


\noindent
见 \ref{algebraization-definition-s-d-inequalities}：
{\it $(\mathcal{F}_n)$ 满足 $(a, b)$-不等式组},
{\it $(\mathcal{F}_n)$ 满足严格 $(a, b)$-不等式组}


\medskip\noindent
{\bf 代数曲线}

\medskip

\noindent
见 \ref{curves-definition-normal-projective-model}：
{\it $X$ 的非奇异投射模型}


\noindent
见 \ref{curves-definition-linear-series}：
{\it 次数为 $d$、维数为 $r$ 的线性系},
{\it $\mathfrak g^r_d$}


\noindent
见 \ref{curves-definition-genus}：
{\it 亏格}


\noindent
见 \ref{curves-definition-geometric-genus}：
{\it 几何亏格}


\noindent
见 \ref{curves-definition-multicross}：
{\it 多重交叉奇点},
{\it 结点},
{\it 普通二重点},
{\it 定义一个结点奇点}


\noindent
见 \ref{curves-definition-nodal}：
{\it 结点},
{\it 普通二重点},
{\it 定义一个结点奇点},
{\it $X$ 的奇点至多为结点型}


\noindent
见 \ref{curves-definition-split-node}：
{\it 分裂结点}


\noindent
见 \ref{curves-definition-nodal-family}：
{\it 相对维数为 $1$ 且至多为结点}


\medskip\noindent
{\bf 曲面的消解}

\medskip

\noindent
见 \ref{resolve-definition-normalized-blowup}：
{\it $X$ 在 $x$ 处的正规化吹起}


\noindent
见 \ref{resolve-definition-reduce-to-rational}：
{\it 定义一个有理奇点},
{\it 对 $A$ 可以归约到有理奇点}


\noindent
见 \ref{resolve-definition-resolution}：
{\it 奇点消解}


\noindent
见 \ref{resolve-definition-resolution-surface}：
{\it 由正规化吹起给出的奇点消解}


\medskip\noindent
{\bf 半稳定约化}

\medskip

\noindent
见 \ref{models-definition-type}：
{\it 数值型}


\noindent
见 \ref{models-definition-type-equivalent}：
{\it 等价型}


\noindent
见 \ref{models-definition-genus}：
{\it 亏格为 $g$ 的数值型}


\noindent
见 \ref{models-definition-type-minus-one}：
{\it $(-1)$-指标}


\noindent
见 \ref{models-definition-top-genus}：
{\it $T$ 的拓扑亏格}


\noindent
见 \ref{models-definition-type-minimal}：
{\it 极小}


\noindent
见 \ref{models-definition-type-minus-two}：
{\it $(-2)$-指标}


\noindent
见 \ref{models-definition-picard-group}：
{\it $T$ 的 Picard 群}


\noindent
见 \ref{models-definition-minimal-model}：
{\it 极小模型}


\noindent
见 \ref{models-definition-numerical-type-model}：
{\it 与 $X$ 关联的数值型}


\noindent
见 \ref{models-definition-semistable}：
{\it 半稳定约化}


\noindent
见 \ref{models-definition-good}：
{\it 良好约化}


\medskip\noindent
{\bf 函子与态射}

\medskip

\medskip\noindent
{\bf 代数簇的导出范畴}

\medskip

\noindent
见 \ref{equiv-definition-Serre-functor}：
{\it 存在 Serre 函子},
{\it Serre 函子}


\noindent
见 \ref{equiv-definition-fourier-mukai-functor}：
{\it Fourier--Mukai 函子},
{\it Fourier--Mukai 核}


\noindent
见 \ref{equiv-definition-siblings}：
{\it 同胞对象},
{\it 同胞}


\noindent
见 \ref{equiv-definition-siblings-geometric}：
{\it 同胞对象},
{\it 同胞}


\noindent
见 \ref{equiv-definition-relative-equivalence-kernel}：
{\it $S$ 上从 $X$ 到 $Y$ 的相对等价的 Fourier--Mukai 核}


\noindent
见 \ref{equiv-definition-derived-equivalent}：
{\it 导出等价}


\medskip\noindent
{\bf 概形的基本群}

\medskip

\noindent
见 \ref{pione-definition-G-set-continuous}：
{\it $G$-集},
{\it 离散 $G$-集},
{\it $G$-集的态射},
{\it $G\textit{-Sets}$}


\noindent
见 \ref{pione-definition-galois-category}：
{\it Galois 范畴}


\noindent
见 \ref{pione-definition-fundamental-group}：
{\it 基本群},
{\it 基点}


\medskip\noindent
{\bf 平展上同调}

\medskip

\noindent
见 \ref{etale-cohomology-definition-etale-covering-initial}：
{\it 平展覆盖}


\noindent
见 \ref{etale-cohomology-definition-presheaf}：
{\it 集合预层},
{\it 阿贝尔预层}


\noindent
见 \ref{etale-cohomology-definition-family-morphisms-fixed-target}：
{\it 具有固定靶的态射族}


\noindent
见 \ref{etale-cohomology-definition-site}：
{\it 位点},
{\it 覆盖}


\noindent
见 \ref{etale-cohomology-definition-sheaf}：
{\it 分离预层},
{\it 层}


\noindent
见 \ref{etale-cohomology-definition-category-sheaves}：
{\it 集合层范畴},
{\it 阿贝尔层}


\noindent
见 \ref{etale-cohomology-definition-0-cech}：
{\it 第零 {\v C}ech 上同调群}


\noindent
见 \ref{etale-cohomology-definition-fpqc-covering}：
{\it fpqc 覆盖}


\noindent
见 \ref{etale-cohomology-definition-sheaf-property-fpqc}：
{\it 满足 fpqc 拓扑的层性质}


\noindent
见 \ref{etale-cohomology-definition-descent-datum}：
{\it 下降数据},
{\it 余循环条件},
{\it 有效}


\noindent
见 \ref{etale-cohomology-definition-descent-datum-modules}：
{\it 下降数据}


\noindent
见 \ref{etale-cohomology-definition-effective-modules}：
{\it 有效}


\noindent
见 \ref{etale-cohomology-definition-ringed-site}：
{\it 环化位点},
{\it 拟凝聚}


\noindent
见 \ref{etale-cohomology-definition-cech-complex}：
{\it {\v C}ech 复形},
{\it {\v C}ech 上同调群}


\noindent
见 \ref{etale-cohomology-definition-free-abelian-presheaf}：
{\it $\mathcal{G}$ 上的自由阿贝尔预层}


\noindent
见 \ref{etale-cohomology-definition-tau-covering}：
{\it $\tau$-覆盖}


\noindent
见 \ref{etale-cohomology-definition-tau-site}：
{\it $S$ 的大 $\tau$-位点},
{\it $S$ 的小 $\tau$-位点}


\noindent
见 \ref{etale-cohomology-definition-standard-tau}：
{\it 标准 $\tau$-覆盖}


\noindent
见 \ref{etale-cohomology-definition-etale-topos}：
{\it \'etale 拓扑斯},
{\it 小 \'etale 拓扑斯},
{\it Zariski 拓扑斯},
{\it 小 Zariski 拓扑斯},
{\it 大 $\tau$-拓扑斯}


\noindent
见 \ref{etale-cohomology-definition-additive-sheaf}：
{\it 常值层}


\noindent
见 \ref{etale-cohomology-definition-structure-sheaf}：
{\it 结构层}


\noindent
见 \ref{etale-cohomology-definition-etale-morphism}：
{\it 平展}


\noindent
见 \ref{etale-cohomology-definition-standard-etale}：
{\it 标准平展的}


\noindent
见 \ref{etale-cohomology-definition-etale-covering}：
{\it 平展覆盖}


\noindent
见 \ref{etale-cohomology-definition-big-etale-site}：
{\it $S$ 上的大平展位点},
{\it $S$ 上的小平展位点},
{\it 大的},
{\it 小 Zariski 位点}


\noindent
见 \ref{etale-cohomology-definition-geometric-point}：
{\it 几何点},
{\it 位于},
{\it \'etale 邻域},
{\it \'etale 邻域态射}


\noindent
见 \ref{etale-cohomology-definition-stalk}：
{\it 茎}


\noindent
见 \ref{etale-cohomology-definition-support}：
{\it $\mathcal{F}$ 的支集},
{\it $\sigma$ 的支集}


\noindent
见 \ref{etale-cohomology-definition-henselian}：
{\it Hensel}


\noindent
见 \ref{etale-cohomology-definition-strictly-henselian}：
{\it 严格 Hensel}


\noindent
见 \ref{etale-cohomology-definition-etale-local-rings}：
{\it $S$ 在 $\overline{s}$ 处的平展局部环},
{\it $\mathcal{O}_{S, s}$ 的严格 Hensel 化},
{\it $\mathcal{O}_{S, s}$ 的 Hensel 化},
{\it $S$ 在 $\overline{s}$ 处的严格 Hensel 化},
{\it $S$ 在 $s$ 处的 Hensel 化}


\noindent
见 \ref{etale-cohomology-definition-direct-image-presheaf}：
{\it 直接像},
{\it 推前}


\noindent
见 \ref{etale-cohomology-definition-direct-image-sheaf}：
{\it 直接像},
{\it 推前}


\noindent
见 \ref{etale-cohomology-definition-higher-direct-images}：
{\it 高阶直接像}


\noindent
见 \ref{etale-cohomology-definition-inverse-image}：
{\it 逆像},
{\it 拉回}


\noindent
见 \ref{etale-cohomology-definition-inverse-system-sheaves}：
{\it $(X_i, f_{i'i})$ 上的层系统 $(\mathcal{F}_i, \varphi_{i'i})$}


\noindent
见 \ref{etale-cohomology-definition-algebraic-geometric-point}：
{\it 绝对 Galois 群},
{\it 代数的}


\noindent
见 \ref{etale-cohomology-definition-G-module-continuous}：
{\it $G$-模},
{\it 离散 $G$-模},
{\it $G$-模的态射},
{\it $R\text{-}G$-模},
{\it $R\text{-}G$-模的态射}


\noindent
见 \ref{etale-cohomology-definition-galois-cohomology}：
{\it 连续群上同调群},
{\it 群上同调群},
{\it Galois 上同调群},
{\it $K$ 以 $M$ 为系数的 Galois 上同调群}


\noindent
见 \ref{etale-cohomology-definition-brauer-equivalent}：
{\it 相似},
{\it 等价}


\noindent
见 \ref{etale-cohomology-definition-brauer-group}：
{\it Brauer 群}


\noindent
见 \ref{etale-cohomology-definition-finite-locally-constant}：
{\it 值为 $E$ 的常值层},
{\it 常值层},
{\it 局部常值},
{\it 有限局部常值},
{\it 值为 $A$ 的常值层},
{\it 常值层},
{\it 局部常值},
{\it 有限局部常值},
{\it 值为 $M$ 的常值层},
{\it 常值层},
{\it 局部常值}


\noindent
见 \ref{etale-cohomology-definition-trace-map}：
{\it 迹}


\noindent
见 \ref{etale-cohomology-definition-Cr}：
{\it $C_r$},
{\it 非平凡解}


\noindent
见 \ref{etale-cohomology-definition-variety}：
{\it 簇},
{\it 曲线}


\noindent
见 \ref{etale-cohomology-definition-extension-zero}：
{\it 零延拓},
{\it 零延拓}


\noindent
见 \ref{etale-cohomology-definition-constructible}：
{\it 可构造},
{\it 可构造},
{\it 可构造}


\noindent
见 \ref{etale-cohomology-definition-c}：
{\it $D_c(X_\etale, \Lambda)$}


\noindent
见 \ref{etale-cohomology-definition-ctf}：
{\it $D_{ctf}(X_\etale, \Lambda)$}


\noindent
见 \ref{etale-cohomology-definition-locally-acyclic}：
{\it 相对于 $K$ 在 $\overline{x}$ 处局部无环},
{\it 相对于 $K$ 局部无环},
{\it 相对于 $K$ 普遍局部无环},
{\it 局部无环},
{\it 普遍局部无环}


\noindent
见 \ref{etale-cohomology-definition-cd}：
{\it $X$ 的上同调维数}


\noindent
见 \ref{etale-cohomology-definition-cd-f}：
{\it $f$ 的上同调维数}


\medskip\noindent
{\bf 晶体上同调}

\medskip

\noindent
见 \ref{crystalline-definition-divided-power-envelope}：
{\it $B$ 中 $J$ 相对于 $(A, I, \gamma)$ 的除幂包络}


\noindent
见 \ref{crystalline-definition-compatible}：
{\it $\delta$ 与 $\gamma$ 兼容}


\noindent
见 \ref{crystalline-definition-affine-thickening}：
{\it 除幂加厚},
{\it 除幂加厚同态}


\noindent
见 \ref{crystalline-definition-derivation}：
{\it 除幂 $A$-导子}


\noindent
见 \ref{crystalline-definition-divided-power-structure}：
{\it 除幂结构 $\gamma$}


\noindent
见 \ref{crystalline-definition-divided-power-scheme}：
{\it 除幂概形},
{\it 除幂概形态射}


\noindent
见 \ref{crystalline-definition-divided-power-thickening}：
{\it 除幂加厚}


\noindent
见 \ref{crystalline-definition-divided-power-thickening-X}：
{\it 相对于 $(S, \mathcal{I}, \gamma)$ 的 $X$ 的除幂加厚},
{\it 相对于 $(S, \mathcal{I}, \gamma)$ 的 $X$ 的除幂加厚态射}


\noindent
见 \ref{crystalline-definition-big-crystalline-site}：
{\it Zariski、\'etale、光滑、syntomic 或 fppf 覆盖},
{\it 大晶体位点}


\noindent
见 \ref{crystalline-definition-crystalline-site}：
{\it 晶体位点}


\noindent
见 \ref{crystalline-definition-modules}：
{\it 局部拟凝聚},
{\it 拟凝聚},
{\it $\mathcal{O}_{X/S}$-模晶体}


\noindent
见 \ref{crystalline-definition-crystal-quasi-coherent-modules}：
{\it 拟凝聚模中的晶体},
{\it 有限局部自由模中的晶体}


\noindent
见 \ref{crystalline-definition-global-derivation}：
{\it $S$-导子 $D : \mathcal{O}_{X/S} \to \mathcal{F}$}


\noindent
见 \ref{crystalline-definition-F-crystal}：
{\it $X/S$ 上（相对于 $\sigma$）的 $F$-晶体},
{\it 非退化}


\medskip\noindent
{\bf pro-平展上同调}

\medskip

\noindent
见 \ref{proetale-definition-w-local}：
{\it w-局部的},
{\it w-局部的}


\noindent
见 \ref{proetale-definition-local-isomorphism}：
{\it 局部同构},
{\it 诱导局部环同构}


\noindent
见 \ref{proetale-definition-ind-zariski}：
{\it ind-Zariski}


\noindent
见 \ref{proetale-definition-ind-etale}：
{\it ind-\'etale}


\noindent
见 \ref{proetale-definition-w-contractible}：
{\it w-可缩的}


\noindent
见 \ref{proetale-definition-fpqc-covering}：
{\it $T$ 的 pro-平展覆盖}


\noindent
见 \ref{proetale-definition-standard-proetale}：
{\it 标准 pro-\'etale 覆盖}


\noindent
见 \ref{proetale-definition-big-proetale-site}：
{\it 大 pro-\'etale 位点}


\noindent
见 \ref{proetale-definition-big-small-proetale}：
{\it $S$ 的大 pro-平展位点},
{\it $S$ 的小 pro-平展位点},
{\it $S$ 的大仿射 pro-平展位点}


\noindent
见 \ref{proetale-definition-restriction-small-proetale}：
{\it 到小 pro-\'etale 位点的限制}


\noindent
见 \ref{proetale-definition-extension-zero}：
{\it 零延拓},
{\it 零延拓}


\noindent
见 \ref{proetale-definition-constructible}：
{\it 可构造}


\noindent
见 \ref{proetale-definition-adic}：
{\it 可构造 $\Lambda$-层},
{\it lisse},
{\it 进 lisse},
{\it 进可构造}


\noindent
见 \ref{proetale-definition-Dbc}：
{\it 可构造}


\noindent
见 \ref{proetale-definition-adic-constructible}：
{\it 进 lisse},
{\it 进可构造}


\medskip\noindent
{\bf 相对循环}

\medskip

\noindent
见 \ref{relative-cycles-definition-relative-cycles}：
{\it $X/S$ 上的相对 $r$-循环}


\noindent
见 \ref{relative-cycles-definition-equidimensional}：
{\it 等维}


\noindent
见 \ref{relative-cycles-definition-effective}：
{\it 有效}


\noindent
见 \ref{relative-cycles-definition-proper}：
{\it 固有相对循环}


\medskip\noindent
{\bf 平展上同调进阶}

\medskip

\noindent
见 \ref{more-etale-definition-f-shriek-separated}：
{\it 紧支撑直像}


\noindent
见 \ref{more-etale-definition-compact-support}：
{\it 紧支撑截面}


\noindent
见 \ref{more-etale-definition-f-shriek-lqf}：
{\it 紧支撑直像}


\noindent
见 \ref{more-etale-definition-cohomology-compact-support}：
{\it $K$ 的带紧支撑上同调},
{\it $K$ 的紧支撑上同调}


\medskip\noindent
{\bf 迹公式}

\medskip

\noindent
见 \ref{trace-definition-geometric-frobenius}：
{\it 几何 Frobenius}


\noindent
见 \ref{trace-definition-arithmetic-frobenius}：
{\it 算术 Frobenius}


\noindent
见 \ref{trace-definition-geometric-frobenius-on-stalk}：
{\it 几何 Frobenius}


\noindent
见 \ref{trace-definition-trace}：
{\it 迹}


\noindent
见 \ref{trace-definition-derived-functor}：
{\it $F$ 的全右导出函子},
{\it $G$ 的全左导出函子}


\noindent
见 \ref{trace-definition-filtered}：
{\it 滤过内射},
{\it 投射},
{\it 滤过拟同构}


\noindent
见 \ref{trace-definition-filtered-derived-functors}：
{\it 滤过导出函子}


\noindent
见 \ref{trace-definition-perfect}：
{\it 完美的}


\noindent
见 \ref{trace-definition-finite-tor-dimension}：
{\it 有限 $\text{Tor}$-维数}


\noindent
见 \ref{trace-definition-global-lefschetz-number}：
{\it 全局 Lefschetz 数}


\noindent
见 \ref{trace-definition-local-lefschetz-number}：
{\it 局部 Lefschetz 数}


\noindent
见 \ref{trace-definition-trace-G}：
{\it $f$ 在 $P$ 上的 $G$-迹}


\noindent
见 \ref{trace-definition-l-adic-sheaf}：
{\it $\mathbf{Z}_\ell$-层},
{\it $\ell$-进层},
{\it lisse},
{\it 态射}


\noindent
见 \ref{trace-definition-torsion-l-adic-sheaf}：
{\it 挠的},
{\it 茎}


\noindent
见 \ref{trace-definition-cohomology-l-adic}：
{\it $\ell$-进上同调}


\noindent
见 \ref{trace-definition-L-function-finite-ring}：
{\it $\mathcal{F}$ 的 $L$-函数}


\noindent
见 \ref{trace-definition-L-function-l-adic}：
{\it $\mathcal{F}$ 的 $L$-函数}


\noindent
见 \ref{trace-definition-open}：
{\it 开的}


\noindent
见 \ref{trace-definition-unramified}：
{\it 取值于 $\Lambda$ 的 $\text{GL}_2(\mathbf{A})$ 上非分歧尖点形式}


\medskip\noindent
{\bf 代数空间}

\medskip

\noindent
见 \ref{spaces-definition-relative-representable-property}：
{\it 具有性质 $\mathcal{P}$}


\noindent
见 \ref{spaces-definition-algebraic-space}：
{\it $S$ 上的代数空间}


\noindent
见 \ref{spaces-definition-morphism-algebraic-spaces}：
{\it $S$ 上代数空间的态射 $f : F \to F'$}


\noindent
见 \ref{spaces-definition-etale-equivalence-relation}：
{\it \'etale 等价关系}


\noindent
见 \ref{spaces-definition-presentation}：
{\it 表示}


\noindent
见 \ref{spaces-definition-immersion}：
{\it 开浸入},
{\it 开子空间},
{\it 闭浸入},
{\it 闭子空间},
{\it 浸入},
{\it 局部闭子空间}


\noindent
见 \ref{spaces-definition-Zariski-open-covering}：
{\it Zariski 覆盖}


\noindent
见 \ref{spaces-definition-small-Zariski-site}：
{\it 小 Zariski 位点 $F_{Zar}$}


\noindent
见 \ref{spaces-definition-separated}：
{\it 在 $S$ 上分离},
{\it 在 $S$ 上局部分离},
{\it 在 $S$ 上拟分离},
{\it 在 $S$ 上 Zariski 局部拟分离}


\noindent
见 \ref{spaces-definition-quotient}：
{\it 自由作用于},
{\it $U$ 对 $G$ 的商}


\noindent
见 \ref{spaces-definition-base-change}：
{\it $F'$ 到 $S$ 的基变换},
{\it 视作 $S'$ 上的代数空间}


\medskip\noindent
{\bf 代数空间的性质}

\medskip

\noindent
见 \ref{spaces-properties-definition-separated}：
{\it 分离},
{\it 局部分离},
{\it 拟分离},
{\it Zariski 局部拟分离},
{\it 分离},
{\it 局部分离},
{\it 拟分离},
{\it Zariski 局部拟分离}


\noindent
见 \ref{spaces-properties-definition-points}：
{\it 点}


\noindent
见 \ref{spaces-properties-definition-topological-space}：
{\it 拓扑空间}


\noindent
见 \ref{spaces-properties-definition-quasi-compact}：
{\it 拟紧}


\noindent
见 \ref{spaces-properties-definition-type-property}：
{\it 具有性质 $\mathcal{P}$}


\noindent
见 \ref{spaces-properties-definition-property-at-point}：
{\it 在 $x$ 处具有性质 $\mathcal{P}$}


\noindent
见 \ref{spaces-properties-definition-locally-constructible}：
{\it \'etale 局部可构造}


\noindent
见 \ref{spaces-properties-definition-dimension-at-point}：
{\it $X$ 在 $x$ 处的维数}


\noindent
见 \ref{spaces-properties-definition-dimension}：
{\it 维数}


\noindent
见 \ref{spaces-properties-definition-dimension-local-ring}：
{\it $X$ 在 $x$ 处之局部环的维数},
{\it $x$ 是 $X$ 上余维为 $d$ 的点}


\noindent
见 \ref{spaces-properties-definition-reduced-induced-space}：
{\it $Z$ 上的代数空间结构},
{\it 诱导既约代数空间结构},
{\it $X$ 的既约化 $X_{red}$}


\noindent
见 \ref{spaces-properties-definition-etale}：
{\it 平展}


\noindent
见 \ref{spaces-properties-definition-etale-site}：
{\it 小 \'etale 站点 $X_\etale$}


\noindent
见 \ref{spaces-properties-definition-spaces-etale-site}：
{\it $X_{spaces, \etale}$}


\noindent
见 \ref{spaces-properties-definition-affine-etale-site}：
{\it $X_{affine, \etale}$}


\noindent
见 \ref{spaces-properties-definition-etale-topos}：
{\it \'etale 拓扑斯},
{\it 小 \'etale 拓扑斯}


\noindent
见 \ref{spaces-properties-definition-f-map}：
{\it $f$-映射 $\varphi : \mathcal{G} \to \mathcal{F}$}


\noindent
见 \ref{spaces-properties-definition-geometric-point}：
{\it 几何点},
{\it 位于 $x$ 上方的几何点}


\noindent
见 \ref{spaces-properties-definition-etale-neighbourhood}：
{\it \'etale 邻域},
{\it \'etale 邻域态射}


\noindent
见 \ref{spaces-properties-definition-stalk}：
{\it 茎}


\noindent
见 \ref{spaces-properties-definition-support}：
{\it $\mathcal{F}$ 的支集},
{\it $\sigma$ 的支集}


\noindent
见 \ref{spaces-properties-definition-structure-sheaf}：
{\it 结构层}


\noindent
见 \ref{spaces-properties-definition-etale-local-rings}：
{\it $X$ 在 $\overline{x}$ 处的平展局部环},
{\it $X$ 在 $\overline{x}$ 处的严格 Hensel 化}


\noindent
见 \ref{spaces-properties-definition-unibranch}：
{\it 在 $x$ 处几何单分支},
{\it 几何单枝}


\noindent
见 \ref{spaces-properties-definition-number-of-geometric-branches}：
{\it $X$ 在 $x$ 处的几何分支数}


\noindent
见 \ref{spaces-properties-definition-noetherian}：
{\it Noether}


\noindent
见 \ref{spaces-properties-definition-regular-at-point}：
{\it $X$ 在 $x$ 处正则}


\noindent
见 \ref{spaces-properties-definition-quasi-coherent}：
{\it 拟凝聚}


\noindent
见 \ref{spaces-properties-definition-locally-projective}：
{\it 局部投射}


\medskip\noindent
{\bf 代数空间的态射}

\medskip

\noindent
见 \ref{spaces-morphisms-definition-separated}：
{\it 分离},
{\it 局部分离},
{\it 拟分离}


\noindent
见 \ref{spaces-morphisms-definition-surjective}：
{\it 满}


\noindent
见 \ref{spaces-morphisms-definition-open}：
{\it 开的},
{\it 普遍开}


\noindent
见 \ref{spaces-morphisms-definition-submersive}：
{\it 商映射},
{\it 普遍商映射}


\noindent
见 \ref{spaces-morphisms-definition-quasi-compact}：
{\it 拟紧}


\noindent
见 \ref{spaces-morphisms-definition-closed}：
{\it 闭的},
{\it 普遍闭}


\noindent
见 \ref{spaces-morphisms-definition-monomorphism}：
{\it 单态射}


\noindent
见 \ref{spaces-morphisms-definition-inverse-image-closed-subspace}：
{\it 闭子空间 $Z$ 的逆像 $f^{-1}(Z)$}


\noindent
见 \ref{spaces-morphisms-definition-scheme-theoretic-intersection-union}：
{\it 概形论交},
{\it 概形论并}


\noindent
见 \ref{spaces-morphisms-definition-scheme-theoretic-support}：
{\it $\mathcal{F}$ 的概形论支集}


\noindent
见 \ref{spaces-morphisms-definition-scheme-theoretic-image}：
{\it 概形论像}


\noindent
见 \ref{spaces-morphisms-definition-scheme-theoretically-dense}：
{\it $U$ 在 $X$ 中的概形论闭包},
{\it 在 $X$ 中概形论稠密}


\noindent
见 \ref{spaces-morphisms-definition-dominant}：
{\it 支配}


\noindent
见 \ref{spaces-morphisms-definition-universally-injective}：
{\it 普遍单射}


\noindent
见 \ref{spaces-morphisms-definition-affine}：
{\it 仿射}


\noindent
见 \ref{spaces-morphisms-definition-relative-spec}：
{\it $\mathcal{A}$ 在 $X$ 上的相对谱},
{\it $\mathcal{A}$ 在 $X$ 上的谱}


\noindent
见 \ref{spaces-morphisms-definition-quasi-affine}：
{\it 拟仿射}


\noindent
见 \ref{spaces-morphisms-definition-P}：
{\it 具有性质 $\mathcal{P}$}


\noindent
见 \ref{spaces-morphisms-definition-P-at-point}：
{\it 在 $x$ 处具有性质 $\mathcal{Q}$}


\noindent
见 \ref{spaces-morphisms-definition-locally-finite-type}：
{\it 局部有限型},
{\it 在 $x$ 处有限型},
{\it 有限型的}


\noindent
见 \ref{spaces-morphisms-definition-finite-type-point}：
{\it 有限型点}


\noindent
见 \ref{spaces-morphisms-definition-locally-quasi-finite}：
{\it 局部拟有限},
{\it 在 $x$ 处准有限},
{\it 拟有限}


\noindent
见 \ref{spaces-morphisms-definition-locally-finite-presentation}：
{\it 局部有限表示},
{\it 在 $x$ 处有限表示},
{\it 有限表示的}


\noindent
见 \ref{spaces-morphisms-definition-flat}：
{\it 平坦},
{\it 在 $x$ 处平坦}


\noindent
见 \ref{spaces-morphisms-definition-flat-module}：
{\it 在 $x$ 处相对于 $Y$ 平坦},
{\it 在 $Y$ 上平坦}


\noindent
见 \ref{spaces-morphisms-definition-dimension-fibre}：
{\it $f$ 在 $x$ 处之纤维的局部环维数},
{\it $x/f(x)$ 的超越次数},
{\it $f$ 在 $x$ 处具有相对维数 $d$}


\noindent
见 \ref{spaces-morphisms-definition-relative-dimension}：
{\it 相对维数 $\leq d$},
{\it 相对维数 $d$}


\noindent
见 \ref{spaces-morphisms-definition-syntomic}：
{\it syntomic},
{\it 在 $x$ 处 syntomic}


\noindent
见 \ref{spaces-morphisms-definition-smooth}：
{\it 光滑},
{\it 在 $x$ 处光滑}


\noindent
见 \ref{spaces-morphisms-definition-unramified}：
{\it 无分歧},
{\it 在 $x$ 处非分歧},
{\it G-非分歧},
{\it 在 $x$ 处 G-非分歧}


\noindent
见 \ref{spaces-morphisms-definition-etale}：
{\it 在 $x$ 处 \'etale}


\noindent
见 \ref{spaces-morphisms-definition-proper}：
{\it 固有}


\noindent
见 \ref{spaces-morphisms-definition-valuative-criterion}：
{\it 满足赋值判据的唯一性部分},
{\it 满足赋值判据的存在性部分},
{\it 满足赋值判据}


\noindent
见 \ref{spaces-morphisms-definition-integral}：
{\it 整},
{\it 有限}


\noindent
见 \ref{spaces-morphisms-definition-finite-locally-free}：
{\it 有限局部自由},
{\it 秩},
{\it 次数}


\noindent
见 \ref{spaces-morphisms-definition-rational-map}：
{\it 等价},
{\it 从 $X$ 到 $Y$ 的有理映射},
{\it 从 $X$ 到 $Y$ 的 $B$-有理映射}


\noindent
见 \ref{spaces-morphisms-definition-rational-function}：
{\it $X$ 上的有理函数}


\noindent
见 \ref{spaces-morphisms-definition-ring-of-rational-functions}：
{\it $X$ 上的有理函数环}


\noindent
见 \ref{spaces-morphisms-definition-domain-of-definition}：
{\it 在点 $x \in |X|$ 处有定义},
{\it 定义域}


\noindent
见 \ref{spaces-morphisms-definition-dominant-rational}：
{\it 支配}


\noindent
见 \ref{spaces-morphisms-definition-birational-spaces}：
{\it 双有理}


\noindent
见 \ref{spaces-morphisms-definition-integral-closure}：
{\it $\mathcal{O}_X$ 在 $\mathcal{A}$ 中的整闭包}


\noindent
见 \ref{spaces-morphisms-definition-normalization-X-in-Y}：
{\it $X$ 在 $Y$ 中的正规化}


\noindent
见 \ref{spaces-morphisms-definition-normalization}：
{\it normalization}


\noindent
见 \ref{spaces-morphisms-definition-universal-homeomorphism}：
{\it 普遍同胚}


\medskip\noindent
{\bf 合宜代数空间}

\medskip

\noindent
见 \ref{decent-spaces-definition-universally-bounded}：
{\it fibres of $f$ are universally bounded}


\noindent
见 \ref{decent-spaces-definition-very-reasonable}：
{\it decent},
{\it reasonable},
{\it very reasonable}


\noindent
见 \ref{decent-spaces-definition-residue-field}：
{\it $X$ 在 $x$ 处的剩余域}


\noindent
见 \ref{decent-spaces-definition-elemenary-etale-neighbourhood}：
{\it 初等平展邻域},
{\it morphism of elementary \'etale neighbourhoods}


\noindent
见 \ref{decent-spaces-definition-henselian-local-ring}：
{\it henselian local ring of $X$ at $x$}


\noindent
见 \ref{decent-spaces-definition-residual-space}：
{\it residual space of $X$ at $x$}


\noindent
见 \ref{decent-spaces-definition-relative-conditions}：
{\it has property $(\beta)$},
{\it has property $(\beta)$},
{\it decent},
{\it reasonable},
{\it very reasonable}


\noindent
见 \ref{decent-spaces-definition-birational}：
{\it 双有理}


\noindent
见 \ref{decent-spaces-definition-unibranch}：
{\it 在 $x$ 处单分支},
{\it 单分支}


\noindent
见 \ref{decent-spaces-definition-number-of-branches}：
{\it $X$ 在 $x$ 处的分支数}


\noindent
见 \ref{decent-spaces-definition-catenary}：
{\it 链状}


\noindent
见 \ref{decent-spaces-definition-universally-catenary}：
{\it 普遍链状}


\medskip\noindent
{\bf 代数空间上同调}

\medskip

\noindent
见 \ref{spaces-cohomology-definition-alternating-cech-complex}：
{\it 交错 {\v C}ech 复形}


\noindent
见 \ref{spaces-cohomology-definition-coherent}：
{\it 凝聚}


\medskip\noindent
{\bf 代数空间的极限}

\medskip

\noindent
见 \ref{spaces-limits-definition-locally-finite-presentation}：
{\it 保持极限},
{\it 局部有限表示},
{\it 在 $S$ 上局部有限表示},
{\it 保持极限},
{\it 局部有限表示},
{\it 相对保持极限}


\noindent
见 \ref{spaces-limits-definition-subsheaf-sections-annihilated-by-ideal}：
{\it 被 $\mathcal{I}$ 零化的截面子层}


\noindent
见 \ref{spaces-limits-definition-subsheaf-sections-supported-on-closed}：
{\it 支集位于 $T$ 上的截面子层}


\medskip\noindent
{\bf 代数空间上的除子}

\medskip

\noindent
见 \ref{spaces-divisors-definition-weakly-associated}：
{\it 弱相伴},
{\it $X$ 的弱伴随点},
{\it $x$ 与 $\mathcal{F}$ 伴随},
{\it $x$ 是 $X$ 的伴随点}


\noindent
见 \ref{spaces-divisors-definition-locally-Noetherian-fibre}：
{\it $f$ 在 $y$ 上的纤维局部 Noether},
{\it $f$ 的纤维局部 Noether}


\noindent
见 \ref{spaces-divisors-definition-relative-weak-assassin}：
{\it $\mathcal{F}$ 在 $X$ 中相对于 $Y$ 的相对弱伴随点集}


\noindent
见 \ref{spaces-divisors-definition-effective-Cartier-divisor}：
{\it 局部主闭子空间},
{\it 有效 Cartier 除子}


\noindent
见 \ref{spaces-divisors-definition-sum-effective-Cartier-divisors}：
{\it 有效 Cartier 除子 $D_1$ 与 $D_2$ 的和}


\noindent
见 \ref{spaces-divisors-definition-pullback-effective-Cartier-divisor}：
{\it $D$ 沿 $f$ 的拉回有定义},
{\it 有效 Cartier 除子的拉回}


\noindent
见 \ref{spaces-divisors-definition-invertible-sheaf-effective-Cartier-divisor}：
{\it 与 $D$ 关联的可逆层 $\mathcal{O}_X(D)$}


\noindent
见 \ref{spaces-divisors-definition-regular-section}：
{\it 正则截面}


\noindent
见 \ref{spaces-divisors-definition-zero-scheme-s}：
{\it 零概形}


\noindent
见 \ref{spaces-divisors-definition-relative-effective-Cartier-divisor}：
{\it 相对有效 Cartier 除子}


\noindent
见 \ref{spaces-divisors-definition-sheaf-meromorphic-functions}：
{\it $X$ 上的亚纯函数层},
{\it $\mathcal{K}_X$},
{\it 亚纯函数}


\noindent
见 \ref{spaces-divisors-definition-meromorphic-section}：
{\it $\mathcal{F}$ 的亚纯截面}


\noindent
见 \ref{spaces-divisors-definition-pullback-meromorphic-sections}：
{\it 对 $f$ 而言亚纯函数的拉回有定义}


\noindent
见 \ref{spaces-divisors-definition-regular-meromorphic-section}：
{\it 正则概形}


\noindent
见 \ref{spaces-divisors-definition-relative-proj}：
{\it $\mathcal{A}$ 在 $X$ 上的相对齐次谱},
{\it $\mathcal{A}$ 在 $X$ 上的齐次谱},
{\it $\mathcal{A}$ 在 $X$ 上的相对 Proj}


\noindent
见 \ref{spaces-divisors-definition-relatively-ample}：
{\it 相对丰沛},
{\it $f$-相对丰沛},
{\it 在 $X/Y$ 上充足},
{\it $f$-丰沛}


\noindent
见 \ref{spaces-divisors-definition-blow-up}：
{\it $X$ 沿 $Z$ 的吹起},
{\it $X$ 关于理想层 $\mathcal{I}$ 的吹起},
{\it 例外除子},
{\it 中心}


\noindent
见 \ref{spaces-divisors-definition-strict-transform}：
{\it 严格变换},
{\it 严格变换}


\noindent
见 \ref{spaces-divisors-definition-admissible-blowup}：
{\it $U$-可容许吹起}


\medskip\noindent
{\bf 域上的代数空间}

\medskip

\noindent
见 \ref{spaces-over-fields-definition-integral-algebraic-space}：
{\it 整}


\noindent
见 \ref{spaces-over-fields-definition-function-field}：
{\it 函数域},
{\it 有理函数域}


\noindent
见 \ref{spaces-over-fields-definition-degree}：
{\it $X$ 在 $Y$ 上的次数}


\noindent
见 \ref{spaces-over-fields-definition-Weil-divisor}：
{\it 素除子},
{\it Weil 除子}


\noindent
见 \ref{spaces-over-fields-definition-order-vanishing}：
{\it $f$ 沿 $Z$ 的消失阶}


\noindent
见 \ref{spaces-over-fields-definition-principal-divisor}：
{\it 与 $f$ 关联的主 Weil 除子}


\noindent
见 \ref{spaces-over-fields-definition-class-group}：
{\it Weil 除子类群}


\noindent
见 \ref{spaces-over-fields-definition-order-vanishing-meromorphic}：
{\it $s$ 沿 $Z$ 的消失阶}


\noindent
见 \ref{spaces-over-fields-definition-divisor-invertible-sheaf}：
{\it 与 $s$ 关联的 Weil 除子},
{\it 与 $\mathcal{L}$ 关联的 Weil 除子类}


\noindent
见 \ref{spaces-over-fields-definition-modification}：
{\it $X$ 的改造}


\noindent
见 \ref{spaces-over-fields-definition-alteration}：
{\it $X$ 的改造}


\noindent
见 \ref{spaces-over-fields-definition-geometrically-reduced}：
{\it 在 $x$ 处几何既约},
{\it 几何既约}


\noindent
见 \ref{spaces-over-fields-definition-geometrically-connected}：
{\it 几何连通}


\noindent
见 \ref{spaces-over-fields-definition-geometrically-irreducible}：
{\it 几何不可约}


\noindent
见 \ref{spaces-over-fields-definition-geometrically-integral}：
{\it 几何整}


\noindent
见 \ref{spaces-over-fields-definition-euler-characteristic}：
{\it $\mathcal{F}$ 的 Euler 示性数}


\noindent
见 \ref{spaces-over-fields-definition-intersection-number}：
{\it 交数}


\medskip\noindent
{\bf 代数空间上的拓扑}

\medskip

\noindent
见 \ref{spaces-topologies-definition-zariski-covering}：
{\it $X$ 的 Zariski 覆盖}


\noindent
见 \ref{spaces-topologies-definition-etale-covering}：
{\it $X$ 的平展覆盖}


\noindent
见 \ref{spaces-topologies-definition-big-etale-site}：
{\it $(\textit{Spaces}/S)_\etale$}


\noindent
见 \ref{spaces-topologies-definition-big-small-etale}：
{\it $(\textit{Spaces}/X)_\etale$}


\noindent
见 \ref{spaces-topologies-definition-restriction-small-etale}：
{\it 到小平展位点的限制}


\noindent
见 \ref{spaces-topologies-definition-smooth-covering}：
{\it $X$ 的光滑覆盖}


\noindent
见 \ref{spaces-topologies-definition-syntomic-covering}：
{\it $X$ 的 syntomic 覆盖}


\noindent
见 \ref{spaces-topologies-definition-fppf-covering}：
{\it $X$ 的 fppf 覆盖}


\noindent
见 \ref{spaces-topologies-definition-big-fppf-site}：
{\it $(\textit{Spaces}/S)_{fppf}$}


\noindent
见 \ref{spaces-topologies-definition-big-small-fppf}：
{\it $(\textit{Spaces}/X)_{fppf}$}


\noindent
见 \ref{spaces-topologies-definition-ph-covering}：
{\it $X$ 的 ph 覆盖}


\noindent
见 \ref{spaces-topologies-definition-big-ph-site}：
{\it $(\textit{Spaces}/S)_{ph}$}


\noindent
见 \ref{spaces-topologies-definition-big-small-ph}：
{\it $(\textit{Spaces}/X)_{ph}$}


\noindent
见 \ref{spaces-topologies-definition-fpqc-covering}：
{\it $X$ 的 fpqc 覆盖}


\medskip\noindent
{\bf 下降与代数空间}

\medskip

\noindent
见 \ref{spaces-descent-definition-descent-datum-quasi-coherent}：
{\it 拟凝聚层下降数据 $(\mathcal{F}_i, \varphi_{ij})$},
{\it 余循环条件},
{\it 下降数据态射 $\psi : (\mathcal{F}_i, \varphi_{ij}) \to (\mathcal{F}'_i, \varphi'_{ij})$}


\noindent
见 \ref{spaces-descent-definition-descent-datum-effective-quasi-coherent}：
{\it 平凡下降数据},
{\it 典范下降数据},
{\it 有效}


\noindent
见 \ref{spaces-descent-definition-property-morphisms-local}：
{\it 在基底上为 $\tau$-局部的},
{\it 在目标上为 $\tau$-局部的},
{\it 关于 $\tau$-拓扑在基底上为局部的}


\noindent
见 \ref{spaces-descent-definition-property-morphisms-local-source}：
{\it 在源上 $\tau$ 局部的},
{\it 对 $\tau$-拓扑在源上局部}


\noindent
见 \ref{spaces-descent-definition-local-source-target}：
{\it 在源与靶上关于光滑拓扑为局部的}


\noindent
见 \ref{spaces-descent-definition-etale-smooth-local-source-target}：
{\it 在源与靶上相对于平展—光滑拓扑为局部的}


\noindent
见 \ref{spaces-descent-definition-descent-datum}：
{\it $V/Y/X$ 的下降数据},
{\it 余循环条件},
{\it 相对于 $Y \to X$ 的下降数据},
{\it 相对于 $Y \to X$ 的下降数据态射 $f : (V/Y, \varphi) \to (V'/Y, \varphi')$}


\noindent
见 \ref{spaces-descent-definition-descent-datum-for-family-of-morphisms}：
{\it 相对于族 $\{X_i \to X\}$ 的下降数据 $(V_i, \varphi_{ij})$},
{\it 态射 $\psi : (V_i, \varphi_{ij}) \to (V'_i, \varphi'_{ij})$}


\noindent
见 \ref{spaces-descent-definition-pullback-functor}：
{\it 拉回函子}


\noindent
见 \ref{spaces-descent-definition-pullback-functor-family}：
{\it 拉回函子}


\noindent
见 \ref{spaces-descent-definition-effective}：
{\it 平凡下降数据},
{\it 典范下降数据},
{\it 有效}


\noindent
见 \ref{spaces-descent-definition-effective-family}：
{\it 典范下降数据},
{\it 有效}


\medskip\noindent
{\bf 空间的导出范畴}

\medskip

\noindent
见 \ref{spaces-perfect-definition-supported-on}：
{\it 支集在 $T$ 上}


\noindent
见 \ref{spaces-perfect-definition-derived-quasi-coherent}：
{\it 具有拟凝聚上同调层的 $\mathcal{O}_X$-模导出范畴}


\noindent
见 \ref{spaces-perfect-definition-proper-over-base}：
{\it $T$ 在 $Y$ 上固有}


\noindent
见 \ref{spaces-perfect-definition-elementary-distinguished-square}：
{\it 初等可区别方块}


\noindent
见 \ref{spaces-perfect-definition-approximation-holds}：
{\it 三元组满足逼近性质}


\noindent
见 \ref{spaces-perfect-definition-approximation}：
{\it 用完美复形逼近的性质成立}


\noindent
见 \ref{spaces-perfect-definition-tor-independent}：
{\it 在 $B$ 上 Tor 无关}


\noindent
见 \ref{spaces-perfect-definition-resolution-property}：
{\it 消解性质}


\medskip\noindent
{\bf 空间态射进阶}

\medskip

\noindent
见 \ref{spaces-more-morphisms-definition-radicial}：
{\it 根态射}


\noindent
见 \ref{spaces-more-morphisms-definition-conormal-sheaf}：
{\it $Z$ 在 $X$ 中的余法层 $\mathcal{C}_{Z/X}$},
{\it $i$ 的余法层}


\noindent
见 \ref{spaces-more-morphisms-definition-conormal-algebra}：
{\it $Z$ 在 $X$ 中的余法代数 $\mathcal{C}_{Z/X, *}$},
{\it $i$ 的余法代数}


\noindent
见 \ref{spaces-more-morphisms-definition-normal-cone}：
{\it 法锥 $C_ZX$},
{\it 法丛}


\noindent
见 \ref{spaces-more-morphisms-definition-sheaf-differentials}：
{\it $X$ 在 $Y$ 上的微分层 $\Omega_{X/Y}$},
{\it 泛 $Y$-导子}


\noindent
见 \ref{spaces-more-morphisms-definition-thickening}：
{\it 加厚},
{\it 一阶加厚},
{\it 有限阶加厚},
{\it $n$ 阶加厚},
{\it 加厚的态射},
{\it $B$ 上的加厚},
{\it $B$ 上加厚的态射}


\noindent
见 \ref{spaces-more-morphisms-definition-first-order-infinitesimal-neighbourhood}：
{\it 一阶无穷小邻域},
{\it $n$ 阶无穷小邻域}


\noindent
见 \ref{spaces-more-morphisms-definition-formally-smooth-etale-unramified}：
{\it 形式光滑的},
{\it 形式平展的},
{\it 形式非分歧的}


\noindent
见 \ref{spaces-more-morphisms-definition-formally-unramified}：
{\it 形式非分歧的}


\noindent
见 \ref{spaces-more-morphisms-definition-universal-thickening}：
{\it 泛一阶加厚},
{\it $Z$ 在 $X$ 上的余法层}


\noindent
见 \ref{spaces-more-morphisms-definition-formally-etale}：
{\it 形式平展的}


\noindent
见 \ref{spaces-more-morphisms-definition-formally-smooth}：
{\it 形式光滑的}


\noindent
见 \ref{spaces-more-morphisms-definition-netherlander}：
{\it $f$ 的朴素余切复形}


\noindent
见 \ref{spaces-more-morphisms-definition-module-flat-on-fibre}：
{\it $\mathcal{F}$ 到其在 $z$ 上纤维的限制在 $x$ 处相对于 $Y$ 在 $z$ 上的纤维平坦},
{\it $X$ 在 $z$ 上的纤维在 $x$ 处相对于 $Y$ 在 $z$ 上的纤维平坦},
{\it $X$ 在 $z$ 上的纤维相对于 $Y$ 在 $z$ 上的纤维平坦}


\noindent
见 \ref{spaces-more-morphisms-definition-CM}：
{\it 在 $x$ 处 Cohen--Macaulay},
{\it Cohen--Macaulay 态射}


\noindent
见 \ref{spaces-more-morphisms-definition-gorenstein}：
{\it Gorenstein at $x$},
{\it Gorenstein morphism}


\noindent
见 \ref{spaces-more-morphisms-definition-geometrically-reduced-fibre}：
{\it $f : X \to Y$ 在 $y$ 处的纤维几何既约}


\noindent
见 \ref{spaces-more-morphisms-definition-regular-immersion}：
{\it Koszul-正则浸入},
{\it $H_1$-正则浸入},
{\it 拟正则浸入}


\noindent
见 \ref{spaces-more-morphisms-definition-relative-pseudo-coherence}：
{\it 相对于 $Y$ 为 $m$-伪凝聚},
{\it 相对于 $Y$ 伪凝聚},
{\it 相对于 $Y$ 为 $m$-伪凝聚},
{\it 相对于 $Y$ 伪凝聚}


\noindent
见 \ref{spaces-more-morphisms-definition-pseudo-coherent}：
{\it 伪凝聚的},
{\it 在 $x$ 处伪凝聚}


\noindent
见 \ref{spaces-more-morphisms-definition-perfect}：
{\it 完美的},
{\it 在 $x$ 处完美}


\noindent
见 \ref{spaces-more-morphisms-definition-lci}：
{\it Koszul 态射},
{\it 局部完全交态射},
{\it 在 $x$ 处为 Koszul}


\noindent
见 \ref{spaces-more-morphisms-definition-relatively-perfect}：
{\it 相对于 $Y$ 完美},
{\it $Y$-完美}


\noindent
见 \ref{spaces-more-morphisms-definition-nodal-family}：
{\it 相对维数为 $1$ 且至多为结点}


\medskip\noindent
{\bf 代数空间上的平坦性}

\medskip

\noindent
见 \ref{spaces-flat-definition-impurity}：
{\it $\mathcal{F}$ 在 $y$ 上方的不纯度}


\noindent
见 \ref{spaces-flat-definition-pure}：
{\it 在 $y$ 上方纯},
{\it 在 $y$ 上方普遍纯},
{\it 在 $y$ 上方纯},
{\it 普遍 $Y$-纯},
{\it 相对于 $Y$ 普遍纯},
{\it $Y$-纯},
{\it 相对于 $Y$ 纯},
{\it $Y$-纯},
{\it 相对于 $Y$ 纯}


\noindent
见 \ref{spaces-flat-definition-flattening}：
{\it $\mathcal{F}$ 的泛扁平化存在},
{\it $X$ 的泛扁平化存在}


\noindent
见 \ref{spaces-flat-definition-flat-dimension-n}：
{\it $\mathcal{F}$ 在维数 $\geq n$ 上相对于 $Y$ 平坦}


\medskip\noindent
{\bf 代数空间中的群胚}

\medskip

\noindent
见 \ref{spaces-groupoids-definition-equivalence-relation}：
{\it 预关系},
{\it 关系},
{\it 预等价关系},
{\it $B$ 上 $U$ 的等价关系}


\noindent
见 \ref{spaces-groupoids-definition-restrict-relation}：
{\it 限制},
{\it 拉回}


\noindent
见 \ref{spaces-groupoids-definition-group-space}：
{\it $B$ 上的群代数空间},
{\it $B$ 上群代数空间的态射 $\psi : (G, m) \to (G', m')$}


\noindent
见 \ref{spaces-groupoids-definition-action-group-space}：
{\it $G$ 在代数空间 $X/B$ 上的作用},
{\it 等变},
{\it $G$-等变态射}


\noindent
见 \ref{spaces-groupoids-definition-free-action}：
{\it 自由的}


\noindent
见 \ref{spaces-groupoids-definition-pseudo-torsor}：
{\it 伪 $G$-扭子},
{\it 在 $G$ 下形式主齐},
{\it 平凡}


\noindent
见 \ref{spaces-groupoids-definition-principal-homogeneous-space}：
{\it 主齐性空间},
{\it $B$ 上的主齐性 $G$-空间},
{\it $\tau$ 拓扑中的 $G$-扭子},
{\it $\tau$ $G$-扭子},
{\it $\tau$ 扭子},
{\it 拟等平凡},
{\it 局部平凡}


\noindent
见 \ref{spaces-groupoids-definition-equivariant-module}：
{\it $G$-等变拟凝聚 $\mathcal{O}_X$-模},
{\it 等变拟凝聚 $\mathcal{O}_X$-模}


\noindent
见 \ref{spaces-groupoids-definition-groupoid}：
{\it $B$ 上代数空间中的群胚},
{\it $B$ 上代数空间中的群胚态射 $f : (U, R, s, t, c) \to (U', R', s', t', c')$}


\noindent
见 \ref{spaces-groupoids-definition-groupoid-module}：
{\it $(U, R, s, t, c)$ 上的拟凝聚模}


\noindent
见 \ref{spaces-groupoids-definition-stabilizer-groupoid}：
{\it 代数空间中的群胚 $(U, R, s, t, c)$ 的稳定子}


\noindent
见 \ref{spaces-groupoids-definition-restrict-groupoid}：
{\it $(U, R, s, t, c)$ 到 $U'$ 的限制}


\noindent
见 \ref{spaces-groupoids-definition-invariant-open}：
{\it $R$-不变的},
{\it $R$-不变的},
{\it $R$-不变的}


\noindent
见 \ref{spaces-groupoids-definition-quotient-sheaf}：
{\it 商层 $U/R$}


\noindent
见 \ref{spaces-groupoids-definition-representable-quotient}：
{\it 商可由代数空间表示},
{\it 可表商},
{\it 可表商},
{\it 商可由代数空间表示}


\noindent
见 \ref{spaces-groupoids-definition-quotient-stack}：
{\it 商叠},
{\it 商叠}


\medskip\noindent
{\bf 空间中的群胚进阶}

\medskip

\noindent
见 \ref{spaces-more-groupoids-definition-split-at-point}：
{\it 在 $u$ 上强分裂},
{\it $R$ 在 $u$ 上的强分裂},
{\it 在 $u$ 上分裂},
{\it $R$ 在 $u$ 上的分裂},
{\it 在 $u$ 上拟分裂},
{\it $R$ 在 $u$ 上的拟分裂}


\medskip\noindent
{\bf 自举}

\medskip

\noindent
见 \ref{bootstrap-definition-morphism-representable-by-spaces}：
{\it 可由代数空间表示}


\noindent
见 \ref{bootstrap-definition-property-transformation}：
{\it 具有性质 $\mathcal{P}$}


\medskip\noindent
{\bf 代数空间的推出}

\medskip

\medskip\noindent
{\bf 空间的 Chow 群}

\medskip

\noindent
见 \ref{spaces-chow-definition-delta-dimension}：
{\it $T$ 的 $\delta$-维数}


\noindent
见 \ref{spaces-chow-definition-cycles}：
{\it $X$ 上的循环},
{\it $k$ 维循环}


\noindent
见 \ref{spaces-chow-definition-length-at-x}：
{\it $\mathcal{F}$ 在 $x$ 处的长度为 $d$}


\noindent
见 \ref{spaces-chow-definition-cycle-associated-to-closed-subscheme}：
{\it $Z$ 在 $Y$ 中的重数},
{\it $Y$ 的伴随 $k$ 维闭链}


\noindent
见 \ref{spaces-chow-definition-cycle-associated-to-coherent-sheaf}：
{\it $Z$ 在 $\mathcal{F}$ 中的重数},
{\it 与 $\mathcal{F}$ 关联的 $k$ 维循环}


\noindent
见 \ref{spaces-chow-definition-proper-pushforward}：
{\it 推前}


\noindent
见 \ref{spaces-chow-definition-flat-pullback}：
{\it $\alpha$ 沿 $f$ 的平坦拉回}


\noindent
见 \ref{spaces-chow-definition-principal-divisor}：
{\it 与 $f$ 关联的主除子}


\noindent
见 \ref{spaces-chow-definition-rational-equivalence}：
{\it 有理等价于零},
{\it 有理等价},
{\it $X$ 上 $k$ 维循环的 Chow 群},
{\it $X$ 上模有理等价的 $k$ 维循环 Chow 群}


\noindent
见 \ref{spaces-chow-definition-divisor-invertible-sheaf}：
{\it 与 $s$ 关联的 Weil 除子},
{\it 与 $\mathcal{L}$ 关联的 Weil 除子}


\noindent
见 \ref{spaces-chow-definition-cap-c1}：
{\it 与 $\mathcal{L}$ 的第一陈类相交}


\noindent
见 \ref{spaces-chow-definition-gysin-homomorphism}：
{\it Gysin 同态}


\noindent
见 \ref{spaces-chow-definition-bivariant-class}：
{\it $f$ 的次数为 $p$ 的双变类 $c$}


\noindent
见 \ref{spaces-chow-definition-chow-cohomology}：
{\it Chow 上同调}


\noindent
见 \ref{spaces-chow-definition-chern-classes}：
{\it $\mathcal{E}$ 的第 $i$ 个 Chern 类},
{\it $\mathcal{E}$ 的总 Chern 类}


\noindent
见 \ref{spaces-chow-definition-degree-zero-cycle}：
{\it 零维循环的次数}


\medskip\noindent
{\bf 群胚的商}

\medskip

\noindent
见 \ref{groupoids-quotients-definition-invariant}：
{\it $R$-不变的},
{\it $G$-不变的}


\noindent
见 \ref{groupoids-quotients-definition-base-change}：
{\it 基变换},
{\it 平坦基变换}


\noindent
见 \ref{groupoids-quotients-definition-categorical}：
{\it 范畴商},
{\it $\mathcal{C}$ 中的范畴商},
{\it 概形范畴中的范畴商},
{\it 概形中的范畴商}


\noindent
见 \ref{groupoids-quotients-definition-universal-categorical}：
{\it 泛范畴商},
{\it 一致范畴商}


\noindent
见 \ref{groupoids-quotients-definition-orbit}：
{\it 轨道},
{\it $R$-轨道}


\noindent
见 \ref{groupoids-quotients-definition-geometric-orbits}：
{\it 弱 $R$-等价},
{\it $R$-等价},
{\it 弱轨道},
{\it 弱 $R$-轨道},
{\it 轨道},
{\it $R$-轨道}


\noindent
见 \ref{groupoids-quotients-definition-set-theoretically-invariant}：
{\it 在集合论意义下 $R$-不变},
{\it 分离轨道},
{\it 分离 $R$-轨道}


\noindent
见 \ref{groupoids-quotients-definition-set-theoretic-equivalence}：
{\it 集合论预等价关系},
{\it 集合论等价关系}


\noindent
见 \ref{groupoids-quotients-definition-orbit-space}：
{\it $R$ 的轨道空间}


\noindent
见 \ref{groupoids-quotients-definition-coarse}：
{\it 粗商},
{\it 概形中的粗商}


\noindent
见 \ref{groupoids-quotients-definition-topological}：
{\it 一致地},
{\it 泛地}


\noindent
见 \ref{groupoids-quotients-definition-functions}：
{\it $X$ 上的 $R$-不变函数层},
{\it $X$ 上的函数就是 $U$ 上的 $R$-不变函数}


\noindent
见 \ref{groupoids-quotients-definition-good}：
{\it 良商}


\noindent
见 \ref{groupoids-quotients-definition-geometric}：
{\it 几何商}


\medskip\noindent
{\bf 空间上同调进阶}

\medskip

\medskip\noindent
{\bf 单纯空间}

\medskip

\noindent
见 \ref{spaces-simplicial-definition-cartesian-sheaf}：
{\it 笛卡尔的},
{\it 笛卡尔的},
{\it 笛卡尔的},
{\it 笛卡尔的}


\noindent
见 \ref{spaces-simplicial-definition-cartesian-derived}：
{\it 导出范畴的单纯系统},
{\it 笛卡尔的},
{\it 导出范畴的单纯系统的态射}


\noindent
见 \ref{spaces-simplicial-definition-cartesian-derived-modules}：
{\it 模导出范畴的单纯系统},
{\it 笛卡尔的},
{\it 模导出范畴的单纯系统的态射}


\noindent
见 \ref{spaces-simplicial-definition-cartesian-morphism}：
{\it 笛卡尔的},
{\it $Y$ 在 $X$ 上是笛卡尔的}


\noindent
见 \ref{spaces-simplicial-definition-fibre-products-simplicial-scheme}：
{\it 与 $f$ 关联的单纯概形}


\medskip\noindent
{\bf 空间的对偶性}

\medskip

\noindent
见 \ref{spaces-duality-definition-dualizing-scheme}：
{\it 对偶复形}


\noindent
见 \ref{spaces-duality-definition-relative-dualizing-proper-flat}：
{\it 相对对偶化复形}


\medskip\noindent
{\bf 形式代数空间}

\medskip

\noindent
见 \ref{formal-spaces-definition-toplogy-tensor-product}：
{\it 张量积},
{\it 完备张量积}


\noindent
见 \ref{formal-spaces-definition-weakly-admissible}：
{\it 拓扑幂零的},
{\it 弱定义理想},
{\it 弱预可容许},
{\it 弱可容许}


\noindent
见 \ref{formal-spaces-definition-taut}：
{\it taut}


\noindent
见 \ref{formal-spaces-definition-adic-homomorphism}：
{\it adic}


\noindent
见 \ref{formal-spaces-definition-weakly-adic}：
{\it 弱预-adic},
{\it $c$-adic},
{\it 弱 adic}


\noindent
见 \ref{formal-spaces-definition-affine-formal-algebraic-space}：
{\it 仿射形式代数空间},
{\it 态射}


\noindent
见 \ref{formal-spaces-definition-types-affine-formal-algebraic-space}：
{\it McQuillan},
{\it 经典的},
{\it 弱 adic},
{\it adic},
{\it adic*},
{\it Noether}


\noindent
见 \ref{formal-spaces-definition-affine-formal-spectrum}：
{\it 形式谱}


\noindent
见 \ref{formal-spaces-definition-countable}：
{\it 可数指标的}


\noindent
见 \ref{formal-spaces-definition-formal-algebraic-space}：
{\it 形式代数空间},
{\it 态射}


\noindent
见 \ref{formal-spaces-definition-completion}：
{\it $X$ 沿 $T$ 的完备化}


\noindent
见 \ref{formal-spaces-definition-separated}：
{\it 拟分离},
{\it 分离}


\noindent
见 \ref{formal-spaces-definition-quasi-compact}：
{\it 拟紧}


\noindent
见 \ref{formal-spaces-definition-quasi-compact-morphism}：
{\it 拟紧}


\noindent
见 \ref{formal-spaces-definition-types-formal-algebraic-spaces}：
{\it 局部可数指标},
{\it 局部可数指标且经典},
{\it 局部弱进},
{\it 局部进*},
{\it 局部 Noether}


\noindent
见 \ref{formal-spaces-definition-adic-morphism}：
{\it 进态射}


\noindent
见 \ref{formal-spaces-definition-finite-type}：
{\it 局部有限型},
{\it 有限型}


\noindent
见 \ref{formal-spaces-definition-surjective}：
{\it 满}


\noindent
见 \ref{formal-spaces-definition-monomorphism}：
{\it 单态射}


\noindent
见 \ref{formal-spaces-definition-closed-immersion}：
{\it 闭浸入}


\noindent
见 \ref{formal-spaces-definition-topologically-finite-type}：
{\it 在……上拓扑有限型}


\noindent
见 \ref{formal-spaces-definition-separated-morphism}：
{\it 分离},
{\it 拟分离}


\noindent
见 \ref{formal-spaces-definition-proper}：
{\it 固有}


\noindent
见 \ref{formal-spaces-definition-etale-sites}：
{\it 小平展位点}


\noindent
见 \ref{formal-spaces-definition-completion-formal-algebraic-space}：
{\it $X$ 沿 $T$ 的完备化}


\noindent
见 \ref{formal-spaces-definition-completion-subspace}：
{\it $X$ 沿 $Z$ 的完备化}


\medskip\noindent
{\bf 形式空间的代数化}

\medskip

\noindent
见 \ref{restricted-definition-rig-smooth-homomorphism}：
{\it 在 $(A, I)$ 上 rig-光滑}


\noindent
见 \ref{restricted-definition-rig-etale-homomorphism}：
{\it 在 $(A, I)$ 上 rig-平展}


\noindent
见 \ref{restricted-definition-flat}：
{\it 平坦}


\noindent
见 \ref{restricted-definition-rig-closed}：
{\it rig-闭}


\noindent
见 \ref{restricted-definition-completed-principal-localization}：
{\it 完备化主局部化}


\noindent
见 \ref{restricted-definition-naively-rig-flat}：
{\it 朴素 rig-平坦的}


\noindent
见 \ref{restricted-definition-rig-flat-continuous-homomorphism}：
{\it rig-平坦的}


\noindent
见 \ref{restricted-definition-rig-flat}：
{\it rig-平坦的}


\noindent
见 \ref{restricted-definition-rig-smooth-continuous-homomorphism}：
{\it rig-光滑的}


\noindent
见 \ref{restricted-definition-rig-smooth}：
{\it rig-光滑的}


\noindent
见 \ref{restricted-definition-rig-etale-continuous-homomorphism}：
{\it rig-etale 的}


\noindent
见 \ref{restricted-definition-rig-etale}：
{\it rig-\,'etale 的}


\noindent
见 \ref{restricted-definition-rig-surjective}：
{\it rig-满射的}


\noindent
见 \ref{restricted-definition-formal-modification}：
{\it 形式修改}


\medskip\noindent
{\bf 再论曲面消解}

\medskip

\noindent
见 \ref{spaces-resolve-definition-blowup-at-point}：
{\it $X$ 在 $x$ 处的吹起 $X' \to X$}


\noindent
见 \ref{spaces-resolve-definition-normalized-blowup}：
{\it $X$ 在 $x$ 处的正规化吹起}


\noindent
见 \ref{spaces-resolve-definition-resolution}：
{\it 奇点消解}


\noindent
见 \ref{spaces-resolve-definition-resolution-surface}：
{\it 由正规化吹起给出的奇点消解}


\medskip\noindent
{\bf 形式形变理论}

\medskip

\noindent
见 \ref{formal-defos-definition-CLambda}：
{\it $\mathcal{C}_\Lambda$},
{\it 经典情形}


\noindent
见 \ref{formal-defos-definition-small-extension}：
{\it 小扩张}


\noindent
见 \ref{formal-defos-definition-tangent-space-ring}：
{\it 相对余切空间}


\noindent
见 \ref{formal-defos-definition-essential-surjection}：
{\it 本质满射}


\noindent
见 \ref{formal-defos-definition-completion-CLambda}：
{\it $\widehat{\mathcal{C}}_\Lambda$}


\noindent
见 \ref{formal-defos-definition-category-cofibred-groupoids}：
{\it 在 $\mathcal{C}$ 上余纤维化为群胚的范畴}


\noindent
见 \ref{formal-defos-definition-prorepresentable}：
{\it 可前表示}


\noindent
见 \ref{formal-defos-definition-predeformation-category}：
{\it 预形变范畴},
{\it 预形变范畴的态射}


\noindent
见 \ref{formal-defos-definition-formal-objects}：
{\it $\mathcal{F}$ 的形式对象范畴 $\widehat{\mathcal{F}}$},
{\it $\mathcal{F}$ 的形式对象 $\xi = (R, \xi_n, f_n)$},
{\it 形式对象的态射 $a : \xi \to \eta$}


\noindent
见 \ref{formal-defos-definition-completion}：
{\it $\mathcal{F}$ 的完备化}


\noindent
见 \ref{formal-defos-definition-smooth-morphism}：
{\it 光滑}


\noindent
见 \ref{formal-defos-definition-versal}：
{\it versal}


\noindent
见 \ref{formal-defos-definition-cofibered-groupoid-projection-smooth}：
{\it 光滑},
{\it 无阻碍}


\noindent
见 \ref{formal-defos-definition-S1-S2}：
{\it 条件 (S1) 与 (S2)}


\noindent
见 \ref{formal-defos-definition-linear}：
{\it $R$-线性}


\noindent
见 \ref{formal-defos-definition-tangent-space-over-R}：
{\it $F$ 的切空间 $TF$}


\noindent
见 \ref{formal-defos-definition-tangent-space}：
{\it $\mathcal{F}$ 的切空间 $T \mathcal{F}$}


\noindent
见 \ref{formal-defos-definition-differential}：
{\it $\varphi$ 的微分 $d \varphi : T \mathcal{F} \to T \mathcal{G}$}


\noindent
见 \ref{formal-defos-definition-minimal-versal}：
{\it 极小},
{\it 极小 versal}


\noindent
见 \ref{formal-defos-definition-RS}：
{\it 条件 (RS)}


\noindent
见 \ref{formal-defos-definition-deformation-category}：
{\it 形变范畴}


\noindent
见 \ref{formal-defos-definition-lifts}：
{\it $x$ 沿 $f$ 的提升},
{\it 提升的态射}


\noindent
见 \ref{formal-defos-definition-relative-infinitesimal-auts}：
{\it $x'$ 在 $x$ 上的无穷小自同构群}


\noindent
见 \ref{formal-defos-definition-infinitesimal-auts}：
{\it $x_0$ 的无穷小自同构群}


\noindent
见 \ref{formal-defos-definition-automorphism-functor}：
{\it $x$ 的自同构函子}


\noindent
见 \ref{formal-defos-definition-groupoid-in-functors}：
{\it $\mathcal{C}$ 上函子中的群胚范畴},
{\it $\mathcal{C}$ 上函子中的群胚},
{\it $\mathcal{C}$ 上函子中的群胚态射 $(U, R, s, t, c) \to (U', R', s', t', c')$}


\noindent
见 \ref{formal-defos-definition-representable}：
{\it 可表示}


\noindent
见 \ref{formal-defos-definition-restricting-groupoids-in-functors}：
{\it $(U, R, s, t, c)$ 到 $\mathcal{C}'$ 的限制 $(U, R, s, t, c)|_{\mathcal{C}'}$}


\noindent
见 \ref{formal-defos-definition-quotient}：
{\it 余纤维化为群胚的商范畴 $[U/R] \to \mathcal{C}$},
{\it 商态射 $U \to [U/R]$}


\noindent
见 \ref{formal-defos-definition-prorepresentable-groupoid-in-functors}：
{\it 可前表示}


\noindent
见 \ref{formal-defos-definition-completion-groupoid-in-functors}：
{\it $(U, R, s, t, c)$ 的完备化 $(U, R, s, t, c)^{\wedge}$}


\noindent
见 \ref{formal-defos-definition-smooth-groupoid-in-functors}：
{\it 光滑}


\noindent
见 \ref{formal-defos-definition-presentation}：
{\it 由 $(U, R, s, t, c)$ 给出的 $\mathcal{F}$ 的表示}


\noindent
见 \ref{formal-defos-definition-minimal-groupoid-in-functors}：
{\it 归一化},
{\it 极小}


\medskip\noindent
{\bf 形变理论}

\medskip

\noindent
见 \ref{defos-definition-strict-morphism-thickenings}：
{\it 增厚的严格态射}


\noindent
见 \ref{defos-definition-strict-morphism-thickenings-ringed-topoi}：
{\it 增厚的严格态射}


\medskip\noindent
{\bf 余切复形}

\medskip

\noindent
见 \ref{cotangent-definition-standard-resolution}：
{\it $B$ 在 $A$ 上的标准消解}


\noindent
见 \ref{cotangent-definition-cotangent-complex-ring-map}：
{\it 余切复形}


\noindent
见 \ref{cotangent-definition-biderivation}：
{\it $A$-双导子}


\noindent
见 \ref{cotangent-definition-atiyah-class}：
{\it Atiyah 类}


\noindent
见 \ref{cotangent-definition-standard-resolution-sheaves-rings}：
{\it $\mathcal{B}$ 在 $\mathcal{A}$ 上的标准消解}


\noindent
见 \ref{cotangent-definition-cotangent-complex-morphism-sheaves-rings}：
{\it 余切复形}


\noindent
见 \ref{cotangent-definition-atiyah-class-general}：
{\it Atiyah 类}


\noindent
见 \ref{cotangent-definition-cotangent-complex-morphism-ringed-spaces}：
{\it 余切复形}


\noindent
见 \ref{cotangent-definition-cotangent-complex-morphism-ringed-topoi}：
{\it 余切复形}


\noindent
见 \ref{cotangent-definition-cotangent-morphism-schemes}：
{\it $X$ 在 $Y$ 上的余切复形 $L_{X/Y}$}


\noindent
见 \ref{cotangent-definition-cotangent-morphism-spaces}：
{\it $X$ 在 $Y$ 上的余切复形 $L_{X/Y}$}


\medskip\noindent
{\bf 形变问题}

\medskip

\medskip\noindent
{\bf 代数叠}

\medskip

\noindent
见 \ref{algebraic-definition-representable-by-algebraic-space}：
{\it 可由 $S$ 上的代数空间表示}


\noindent
见 \ref{algebraic-definition-representable-by-algebraic-spaces}：
{\it 可由代数空间表示}


\noindent
见 \ref{algebraic-definition-relative-representable-property}：
{\it 具有性质 $\mathcal{P}$}


\noindent
见 \ref{algebraic-definition-algebraic-stack}：
{\it $S$ 上的代数叠}


\noindent
见 \ref{algebraic-definition-deligne-mumford}：
{\it Deligne--Mumford 叠}


\noindent
见 \ref{algebraic-definition-morphism-algebraic-stacks}：
{\it $S$ 上代数叠的 $2$-范畴}


\noindent
见 \ref{algebraic-definition-smooth-groupoid}：
{\it 光滑群胚}


\noindent
见 \ref{algebraic-definition-presentation}：
{\it 表示}


\noindent
见 \ref{algebraic-definition-viewed-as}：
{\it 视为 $S'$ 上的代数叠}


\noindent
见 \ref{algebraic-definition-change-of-base}：
{\it $\mathcal{X}'$ 的基变换}


\medskip\noindent
{\bf 叠的例子}

\medskip

\noindent
见 \ref{examples-stacks-definition-hilbert-d-stack}：
{\it $\mathcal{X}$ 在 $\mathcal{Y}$ 上次数 $d$ 的有限 Hilbert 叠}


\medskip\noindent
{\bf 代数叠上的层}

\medskip

\noindent
见 \ref{stacks-sheaves-definition-presheaves}：
{\it $\mathcal{X}$ 上的预层},
{\it $\mathcal{X}$ 上预层的态射}


\noindent
见 \ref{stacks-sheaves-definition-inherited-topologies}：
{\it 相联系的 Zariski 位点},
{\it 相联系的 \'etale 位点},
{\it 相联系的光滑位点},
{\it 相联系的 syntomic 位点},
{\it 相联系的 fppf 位点}


\noindent
见 \ref{stacks-sheaves-definition-sheaves}：
{\it Zariski 层},
{\it 关于 Zariski 拓扑的层},
{\it \'etale 层},
{\it 关于 \'etale 拓扑的层},
{\it 光滑层},
{\it 关于光滑拓扑的层},
{\it syntomic 层},
{\it 关于 syntomic 拓扑的层},
{\it fppf 层},
{\it 层},
{\it 关于 fppf 拓扑的层}


\noindent
见 \ref{stacks-sheaves-definition-morphism}：
{\it 相联系的 fppf 拓扑斯态射}


\noindent
见 \ref{stacks-sheaves-definition-structure-sheaf}：
{\it $\mathcal{X}$ 的结构层}


\noindent
见 \ref{stacks-sheaves-definition-modules}：
{\it $\mathcal{X}$ 上的模预层},
{\it $\mathcal{O}_\mathcal{X}$-模},
{\it $\mathcal{O}_\mathcal{X}$-模层}


\noindent
见 \ref{stacks-sheaves-definition-pullback}：
{\it $\mathcal{F}$ 的拉回 $x^{-1}\mathcal{F}$},
{\it $\mathcal{F}$ 到 $U_\etale$ 的限制}


\noindent
见 \ref{stacks-sheaves-definition-quasi-coherent}：
{\it $\mathcal{X}$ 上的拟凝聚模},
{\it 拟凝聚 $\mathcal{O}_\mathcal{X}$-模}


\noindent
见 \ref{stacks-sheaves-definition-locally-quasi-coherent}：
{\it 局部拟凝聚}


\noindent
见 \ref{stacks-sheaves-definition-alternative}：
{\it 相联系的仿射位点}


\noindent
见 \ref{stacks-sheaves-definition-alternative-inherited-topologies}：
{\it 相联系的仿射 Zariski 位点},
{\it 相联系的仿射 \'etale 位点},
{\it 相联系的仿射光滑位点},
{\it 相联系的仿射 syntomic 位点},
{\it 相联系的仿射 fppf 位点}


\noindent
见 \ref{stacks-sheaves-definition-QC}：
{\it 导出范畴中拟凝聚对象所成的三角范畴}


\medskip\noindent
{\bf 可表示性判据}

\medskip

\noindent
见 \ref{criteria-definition-algebraic}：
{\it 代数的}


\medskip\noindent
{\bf Artin 公理}

\medskip

\noindent
见 \ref{artin-definition-RS}：
{\it 条件 (RS)}


\noindent
见 \ref{artin-definition-formal-objects}：
{\it 形式对象},
{\it 形式对象的态射},
{\it 位于}


\noindent
见 \ref{artin-definition-effective}：
{\it 有效}


\noindent
见 \ref{artin-definition-limit-preserving}：
{\it 保持极限}


\noindent
见 \ref{artin-definition-versal-formal-object}：
{\it versal}


\noindent
见 \ref{artin-definition-versal}：
{\it versal}


\noindent
见 \ref{artin-definition-openness-versality}：
{\it 通用性的开性},
{\it 通用性的开性}


\noindent
见 \ref{artin-definition-RS-star}：
{\it 条件 (RS*)}


\noindent
见 \ref{artin-definition-obstruction-theory}：
{\it 障碍理论},
{\it 障碍模},
{\it 障碍}


\noindent
见 \ref{artin-definition-naive-obstruction-theory}：
{\it 朴素障碍理论}


\medskip\noindent
{\bf Quot 与 Hilbert 空间}

\medskip

\medskip\noindent
{\bf 代数叠的性质}

\medskip

\noindent
见 \ref{stacks-properties-definition-points}：
{\it 点}


\noindent
见 \ref{stacks-properties-definition-topological-space}：
{\it 拓扑空间}


\noindent
见 \ref{stacks-properties-definition-surjective}：
{\it 满}


\noindent
见 \ref{stacks-properties-definition-quasi-compact}：
{\it 拟紧}


\noindent
见 \ref{stacks-properties-definition-type-property}：
{\it 具有性质 $\mathcal{P}$}


\noindent
见 \ref{stacks-properties-definition-property-at-point}：
{\it 在 $x$ 处具有性质 $\mathcal{P}$}


\noindent
见 \ref{stacks-properties-definition-monomorphism}：
{\it 单态射}


\noindent
见 \ref{stacks-properties-definition-immersion}：
{\it 开浸入},
{\it 闭浸入},
{\it 浸入}


\noindent
见 \ref{stacks-properties-definition-substacks}：
{\it 开子叠},
{\it 闭子叠},
{\it 局部闭子叠}


\noindent
见 \ref{stacks-properties-definition-reduced-induced-stack}：
{\it $Z$ 上的代数叠结构},
{\it 既约诱导代数叠结构},
{\it $\mathcal{X}$ 的既约化 $\mathcal{X}_{red}$}


\noindent
见 \ref{stacks-properties-definition-residual-gerbe}：
{\it $\mathcal{X}$ 在 $x$ 处的剩余胚存在},
{\it $\mathcal{X}$ 在 $x$ 处的剩余胚}


\noindent
见 \ref{stacks-properties-definition-dimension-at-point}：
{\it $\mathcal{X}$ 在 $x$ 处的维数}


\noindent
见 \ref{stacks-properties-definition-dimension}：
{\it 维数}


\noindent
见 \ref{stacks-properties-definition-number-of-geometric-branches}：
{\it $\mathcal{X}$ 在 $x$ 处的几何分支数},
{\it 在 $x$ 处几何单分支}


\medskip\noindent
{\bf 代数叠的态射}

\medskip

\noindent
见 \ref{stacks-morphisms-definition-separated}：
{\it DM},
{\it 拟 DM},
{\it 分离},
{\it 拟分离}


\noindent
见 \ref{stacks-morphisms-definition-absolute-separated}：
{\it 在 $S$ 上为 DM},
{\it 在 $S$ 上拟 DM},
{\it 在 $S$ 上分离},
{\it 在 $S$ 上拟分离},
{\it DM},
{\it 拟 DM},
{\it 分离},
{\it 拟分离}


\noindent
见 \ref{stacks-morphisms-definition-isom}：
{\it $x$ 的相对自同构层},
{\it 从 $x_1$ 到 $x_2$ 的相对同构层},
{\it $x$ 的自同构层},
{\it 从 $x_1$ 到 $x_2$ 的同构层}


\noindent
见 \ref{stacks-morphisms-definition-quasi-compact}：
{\it 拟紧}


\noindent
见 \ref{stacks-morphisms-definition-noetherian}：
{\it Noether}


\noindent
见 \ref{stacks-morphisms-definition-affine}：
{\it 仿射}


\noindent
见 \ref{stacks-morphisms-definition-integral}：
{\it 整},
{\it 有限}


\noindent
见 \ref{stacks-morphisms-definition-open}：
{\it 开的},
{\it 普遍开}


\noindent
见 \ref{stacks-morphisms-definition-submersive}：
{\it 商映射},
{\it 普遍商映射}


\noindent
见 \ref{stacks-morphisms-definition-closed}：
{\it 闭的},
{\it 普遍闭}


\noindent
见 \ref{stacks-morphisms-definition-universally-injective}：
{\it 普遍单射}


\noindent
见 \ref{stacks-morphisms-definition-universal-homeomorphism}：
{\it 普遍同胚}


\noindent
见 \ref{stacks-morphisms-definition-P}：
{\it 具有性质 $\mathcal{P}$}


\noindent
见 \ref{stacks-morphisms-definition-locally-finite-type}：
{\it 局部有限型},
{\it 有限型的}


\noindent
见 \ref{stacks-morphisms-definition-finite-type-point}：
{\it 有限型点}


\noindent
见 \ref{stacks-morphisms-definition-locally-quasi-finite}：
{\it 局部拟有限}


\noindent
见 \ref{stacks-morphisms-definition-quasi-finite}：
{\it 拟有限}


\noindent
见 \ref{stacks-morphisms-definition-flat}：
{\it 平坦}


\noindent
见 \ref{stacks-morphisms-definition-flat-at-point}：
{\it 在 $x$ 处平坦}


\noindent
见 \ref{stacks-morphisms-definition-locally-finite-presentation}：
{\it 局部有限表示},
{\it 有限表示的}


\noindent
见 \ref{stacks-morphisms-definition-gerbe}：
{\it ……上的 gerbe（胚）},
{\it gerbe（胚）}


\noindent
见 \ref{stacks-morphisms-definition-smooth}：
{\it 光滑}


\noindent
见 \ref{stacks-morphisms-definition-etale-smooth-P}：
{\it 具有性质 $\mathcal{P}$}


\noindent
见 \ref{stacks-morphisms-definition-etale}：
{\it 平展}


\noindent
见 \ref{stacks-morphisms-definition-unramified}：
{\it 无分歧}


\noindent
见 \ref{stacks-morphisms-definition-proper}：
{\it 固有}


\noindent
见 \ref{stacks-morphisms-definition-scheme-theoretic-image}：
{\it 概形论像}


\noindent
见 \ref{stacks-morphisms-definition-fill-in-diagram}：
{\it dotted arrow},
{\it 虚线箭头的态射}


\noindent
见 \ref{stacks-morphisms-definition-uniqueness}：
{\it uniqueness part of the valuative criterion}


\noindent
见 \ref{stacks-morphisms-definition-existence}：
{\it existence part of the valuative criterion}


\noindent
见 \ref{stacks-morphisms-definition-lci}：
{\it 局部完全交态射},
{\it Koszul}


\noindent
见 \ref{stacks-morphisms-definition-normalization}：
{\it normalization}


\noindent
见 \ref{stacks-morphisms-definition-decent}：
{\it decent}


\noindent
见 \ref{stacks-morphisms-definition-integral-algebraic-stack}：
{\it 整}


\medskip\noindent
{\bf 代数叠的极限}

\medskip

\noindent
见 \ref{stacks-limits-definition-limit-preserving}：
{\it 保持极限}


\medskip\noindent
{\bf 代数叠上同调}

\medskip

\noindent
见 \ref{stacks-cohomology-definition-flat-base-change}：
{\it 平坦基变换性质}


\noindent
见 \ref{stacks-cohomology-definition-parasitic}：
{\it 寄生的}


\noindent
见 \ref{stacks-cohomology-definition-lisse-etale}：
{\it 光滑-\'etale site},
{\it 平坦-fppf site}


\noindent
见 \ref{stacks-cohomology-definition-coherent}：
{\it 凝聚}


\medskip\noindent
{\bf 叠的导出范畴}

\medskip

\noindent
见 \ref{stacks-perfect-definition-derived}：
{\it 上同调层拟凝聚的 $\mathcal{O}_\mathcal{X}$-模导出范畴}


\medskip\noindent
{\bf 代数叠导论}

\medskip

\noindent
见 \ref{stacks-introduction-definition-smooth}：
{\it 光滑}


\noindent
见 \ref{stacks-introduction-definition-algebraic-stack}：
{\it 代数叠}


\medskip\noindent
{\bf 叠态射进阶}

\medskip

\noindent
见 \ref{stacks-more-morphisms-definition-thickening}：
{\it 加厚},
{\it 加厚的态射},
{\it $\mathcal{Z}$ 上的增厚},
{\it $\mathcal{Z}$ 上增厚的态射}


\noindent
见 \ref{stacks-more-morphisms-definition-first-order-thickening}：
{\it 一阶加厚}


\noindent
见 \ref{stacks-more-morphisms-definition-formally-smooth}：
{\it 形式光滑的}


\noindent
见 \ref{stacks-more-morphisms-definition-categorical-quotient}：
{\it 范畴模空间},
{\it 一致范畴模空间},
{\it $\mathcal{C}$ 中的范畴模空间},
{\it $\mathcal{C}$ 中的一致范畴模空间}


\noindent
见 \ref{stacks-more-morphisms-definition-well-nigh-affine}：
{\it 近乎仿射的}


\medskip\noindent
{\bf 叠的几何}

\medskip

\noindent
见 \ref{stacks-geometry-definition-versal-ring-at-x}：
{\it $\mathcal{X}$ 在 $x_0$ 处的 versal 环}


\noindent
见 \ref{stacks-geometry-definition-multiplicity}：
{\it 重数}


\noindent
见 \ref{stacks-geometry-definition-formal-branches}：
{\it $\mathcal{X}$ 穿过 $x_0$ 的形式分支}


\noindent
见 \ref{stacks-geometry-definition-multiplicity-formal-branches}：
{\it $\mathcal{X}$ 穿过 $x_0$ 的一个形式分支的重数}


\noindent
见 \ref{stacks-geometry-definition-relative-dimension}：
{\it 相对维数}


\noindent
见 \ref{stacks-geometry-definition-relative-dimension-for-stacks}：
{\it 相对维数}


\noindent
见 \ref{stacks-geometry-definition-pseudo-catenary}：
{\it 拟链状的}


\noindent
见 \ref{stacks-geometry-definition-dimension-local-ring}：
{\it $\mathcal{X}$ 在 $x$ 处的局部环的维数}


\medskip\noindent
{\bf 模叠}

\medskip

\medskip\noindent
{\bf 曲线模空间}

\medskip

\noindent
见 \ref{moduli-curves-definition-deligne-mumford-smooth}：
{\it 光滑真曲线的模叠},
{\it 亏格为 $g$ 的光滑真曲线的模叠}


\noindent
见 \ref{moduli-curves-definition-relative-dualizing-sheaf}：
{\it 相对对偶层}


\noindent
见 \ref{moduli-curves-definition-prestable}：
{\it 预稳定曲线族}


\noindent
见 \ref{moduli-curves-definition-semistable}：
{\it 半稳定曲线族}


\noindent
见 \ref{moduli-curves-definition-stable}：
{\it 稳定曲线族}


\noindent
见 \ref{moduli-curves-definition-deligne-mumford}：
{\it 稳定曲线模叠},
{\it 亏格为 $g$ 的稳定曲线模叠}


\medskip\noindent
{\bf 例}

\medskip

\medskip\noindent
{\bf 练习}

\medskip

\noindent
见 \ref{exercises-definition-directed-poset}：
{\it 有向集},
{\it 环系}


\noindent
见 \ref{exercises-definition-colimit}：
{\it 余极限}


\noindent
见 \ref{exercises-definition-finite-presentation}：
{\it 有限表示}


\noindent
见 \ref{exercises-definition-quasi-compact}：
{\it 拟紧}


\noindent
见 \ref{exercises-definition-Hausdorff}：
{\it 豪斯多夫}


\noindent
见 \ref{exercises-definition-irreducible}：
{\it 不可约},
{\it 不可约}


\noindent
见 \ref{exercises-definition-generic-point}：
{\it 泛点}


\noindent
见 \ref{exercises-definition-Noetherian-space}：
{\it Noether},
{\it Artin}


\noindent
见 \ref{exercises-definition-irreducible-component}：
{\it 不可约分支}


\noindent
见 \ref{exercises-definition-closed}：
{\it 闭的},
{\it 特化},
{\it 一般化}


\noindent
见 \ref{exercises-definition-connected-component}：
{\it 连通},
{\it 连通分支}


\noindent
见 \ref{exercises-definition-length}：
{\it 长度}


\noindent
见 \ref{exercises-definition-catenary}：
{\it 链状},
{\it 链状}


\noindent
见 \ref{exercises-definition-finite-locally-free}：
{\it 有限局部自由},
{\it 可逆模}


\noindent
见 \ref{exercises-definition-class-group}：
{\it $A$ 的类群},
{\it $A$ 的 Picard 群}


\noindent
见 \ref{exercises-definition-GU-GD}：
{\it 上升定理},
{\it 下降定理}


\noindent
见 \ref{exercises-definition-numerical-polynomial}：
{\it 数值多项式}


\noindent
见 \ref{exercises-definition-graded-module}：
{\it 分次模},
{\it 局部有限},
{\it Euler--Poincar\'e 函数},
{\it Hilbert 函数},
{\it Hilbert 多项式}


\noindent
见 \ref{exercises-definition-graded-algebra}：
{\it 分次 $A$-代数},
{\it 分次 $A$-代数 $B$ 上的分次模 $M$},
{\it 分次模或分次环的同态},
{\it 分次子模},
{\it 分次理想},
{\it 分次模的正合列}


\noindent
见 \ref{exercises-definition-homogeneous-ideal}：
{\it 齐次}


\noindent
见 \ref{exercises-definition-Proj-R}：
{\it 齐次谱 $\text{Proj}(R)$}


\noindent
见 \ref{exercises-definition-Dplus-Vplus}：
{\it $R_{(f)}$}


\noindent
见 \ref{exercises-definition-CM}：
{\it Cohen--Macaulay}


\noindent
见 \ref{exercises-definition-injective-filtered}：
{\it 滤过内射}


\noindent
见 \ref{exercises-definition-finite-filtration-category}：
{\it $\text{Fil}^f(\mathcal{A})$}


\noindent
见 \ref{exercises-definition-filtered-quasi-isomorphism}：
{\it 滤过拟同构}


\noindent
见 \ref{exercises-definition-filtered-acyclic}：
{\it 滤过无环}


\noindent
见 \ref{exercises-definition-integral}：
{\it 整}


\noindent
见 \ref{exercises-definition-dual-numbers}：
{\it 对偶数}


\noindent
见 \ref{exercises-definition-tangent-space}：
{\it $X$ 在 $S$ 上的切空间},
{\it 切向量}


\noindent
见 \ref{exercises-definition-quasi-coherent}：
{\it 拟凝聚}


\noindent
见 \ref{exercises-definition-specialization}：
{\it 特化}


\noindent
见 \ref{exercises-definition-Noetherian-scheme}：
{\it 局部 Noether},
{\it Noether}


\noindent
见 \ref{exercises-definition-coherent}：
{\it 凝聚}


\noindent
见 \ref{exercises-definition-invertible-sheaf}：
{\it 可逆 ${\mathcal O}_X$-模}


\noindent
见 \ref{exercises-definition-invertible-module}：
{\it 可逆模 $M$},
{\it 平凡}


\noindent
见 \ref{exercises-definition-picard-group}：
{\it $X$ 的 Picard 群}


\noindent
见 \ref{exercises-definition-delta}：
{\it $\delta(\tau)$}


\noindent
见 \ref{exercises-definition-divisor}：
{\it Weil 除子},
{\it 素除子},
{\it 与有理函数 $f \in K(X)^\ast$ 关联的 Weil 除子},
{\it 有效 Cartier 除子},
{\it 与有效 Cartier 除子 $D \subset X$ 关联的 Weil 除子 $[D]$},
{\it 全商环层 ${\mathcal K}_S$},
{\it Cartier 除子},
{\it 相伴的 Weil 除子}


\medskip\noindent
{\bf 文献指南}

\medskip

\medskip\noindent
{\bf 待完善事项}

\medskip

\medskip\noindent
{\bf 编码规范}

\medskip

\medskip\noindent
{\bf 已废弃内容}

\medskip

\noindent
见 \ref{obsolete-definition-epsilon}：
{\it $\epsilon$-不变量}


\noindent
见 \ref{obsolete-definition-locally-finite-sum-effective-Cartier-divisors}：
{\it 和}


\medskip\noindent
{\bf GNU 自由文档许可证}

\medskip
\end{multicols}

% Generated mechanically from the frozen fdl.tex target.
% Do not edit this generated file; edit the bound locale source instead.

% OK, start here.
%
%---------------------------------------------------------------------
%\tableofcontents

\chapter{GNU Free Documentation License}

\phantomsection
\label{fdl-section-phantom}

\begin{center}

      Version 1.2, November 2002


Copyright \copyright 2000, 2001, 2002  Free Software Foundation, Inc.

\bigskip

    51 Franklin St, Fifth Floor, Boston, MA  02110-1301  USA

\bigskip

Everyone is permitted to copy and distribute verbatim copies
of this license document, but changing it is not allowed.
\end{center}


\begin{center}
{\bf\large Preamble}
\end{center}

The purpose of this License is to make a manual, textbook, or other
functional and useful document "free" in the sense of freedom: to
assure everyone the effective freedom to copy and redistribute it,
with or without modifying it, either commercially or noncommercially.
Secondarily, this License preserves for the author and publisher a way
to get credit for their work, while not being considered responsible
for modifications made by others.

This License is a kind of "copyleft", which means that derivative
works of the document must themselves be free in the same sense.  It
complements the GNU General Public License, which is a copyleft
license designed for free software.

We have designed this License in order to use it for manuals for free
software, because free software needs free documentation: a free
program should come with manuals providing the same freedoms that the
software does.  But this License is not limited to software manuals;
it can be used for any textual work, regardless of subject matter or
whether it is published as a printed book.  We recommend this License
principally for works whose purpose is instruction or reference.


\section{APPLICABILITY AND DEFINITIONS}
\label{fdl-section-applicability-and-definitions}

This License applies to any manual or other work, in any medium, that
contains a notice placed by the copyright holder saying it can be
distributed under the terms of this License.  Such a notice grants a
world-wide, royalty-free license, unlimited in duration, to use that
work under the conditions stated herein.  The \textbf{"Document"}, below,
refers to any such manual or work.  Any member of the public is a
licensee, and is addressed as \textbf{"you"}.  You accept the license if you
copy, modify or distribute the work in a way requiring permission
under copyright law.

A \textbf{"Modified Version"} of the Document means any work containing the
Document or a portion of it, either copied verbatim, or with
modifications and/or translated into another language.

A \textbf{"Secondary Section"} is a named appendix or a front-matter section of
the Document that deals exclusively with the relationship of the
publishers or authors of the Document to the Document's overall subject
(or to related matters) and contains nothing that could fall directly
within that overall subject.  (Thus, if the Document is in part a
textbook of mathematics, a Secondary Section may not explain any
mathematics.)  The relationship could be a matter of historical
connection with the subject or with related matters, or of legal,
commercial, philosophical, ethical or political position regarding
them.

The \textbf{"Invariant Sections"} are certain Secondary Sections whose titles
are designated, as being those of Invariant Sections, in the notice
that says that the Document is released under this License.  If a
section does not fit the above definition of Secondary then it is not
allowed to be designated as Invariant.  The Document may contain zero
Invariant Sections.  If the Document does not identify any Invariant
Sections then there are none.

The \textbf{"Cover Texts"} are certain short passages of text that are listed,
as Front-Cover Texts or Back-Cover Texts, in the notice that says that
the Document is released under this License.  A Front-Cover Text may
be at most 5 words, and a Back-Cover Text may be at most 25 words.

A \textbf{"Transparent"} copy of the Document means a machine-readable copy,
represented in a format whose specification is available to the
general public, that is suitable for revising the document
straightforwardly with generic text editors or (for images composed of
pixels) generic paint programs or (for drawings) some widely available
drawing editor, and that is suitable for input to text formatters or
for automatic translation to a variety of formats suitable for input
to text formatters.  A copy made in an otherwise Transparent file
format whose markup, or absence of markup, has been arranged to thwart
or discourage subsequent modification by readers is not Transparent.
An image format is not Transparent if used for any substantial amount
of text.  A copy that is not "Transparent" is called \textbf{"Opaque"}.

Examples of suitable formats for Transparent copies include plain
ASCII without markup, Texinfo input format, LaTeX input format, SGML
or XML using a publicly available DTD, and standard-conforming simple
HTML, PostScript or PDF designed for human modification.  Examples of
transparent image formats include PNG, XCF and JPG.  Opaque formats
include proprietary formats that can be read and edited only by
proprietary word processors, SGML or XML for which the DTD and/or
processing tools are not generally available, and the
machine-generated HTML, PostScript or PDF produced by some word
processors for output purposes only.

The \textbf{"Title Page"} means, for a printed book, the title page itself,
plus such following pages as are needed to hold, legibly, the material
this License requires to appear in the title page.  For works in
formats which do not have any title page as such, "Title Page" means
the text near the most prominent appearance of the work's title,
preceding the beginning of the body of the text.

A section \textbf{"Entitled XYZ"} means a named subunit of the Document whose
title either is precisely XYZ or contains XYZ in parentheses following
text that translates XYZ in another language.  (Here XYZ stands for a
specific section name mentioned below, such as \textbf{"Acknowledgements"},
\textbf{"Dedications"}, \textbf{"Endorsements"}, or \textbf{"History"}.)
To \textbf{"Preserve the Title"}
of such a section when you modify the Document means that it remains a
section "Entitled XYZ" according to this definition.

The Document may include Warranty Disclaimers next to the notice which
states that this License applies to the Document.  These Warranty
Disclaimers are considered to be included by reference in this
License, but only as regards disclaiming warranties: any other
implication that these Warranty Disclaimers may have is void and has
no effect on the meaning of this License.


\section{VERBATIM COPYING}
\label{fdl-section-verbatim-copying}

You may copy and distribute the Document in any medium, either
commercially or noncommercially, provided that this License, the
copyright notices, and the license notice saying this License applies
to the Document are reproduced in all copies, and that you add no other
conditions whatsoever to those of this License.  You may not use
technical measures to obstruct or control the reading or further
copying of the copies you make or distribute.  However, you may accept
compensation in exchange for copies.  If you distribute a large enough
number of copies you must also follow the conditions in section 3.

You may also lend copies, under the same conditions stated above, and
you may publicly display copies.


\section{COPYING IN QUANTITY}
\label{fdl-section-copying-in-quantity}


If you publish printed copies (or copies in media that commonly have
printed covers) of the Document, numbering more than 100, and the
Document's license notice requires Cover Texts, you must enclose the
copies in covers that carry, clearly and legibly, all these Cover
Texts: Front-Cover Texts on the front cover, and Back-Cover Texts on
the back cover.  Both covers must also clearly and legibly identify
you as the publisher of these copies.  The front cover must present
the full title with all words of the title equally prominent and
visible.  You may add other material on the covers in addition.
Copying with changes limited to the covers, as long as they preserve
the title of the Document and satisfy these conditions, can be treated
as verbatim copying in other respects.

If the required texts for either cover are too voluminous to fit
legibly, you should put the first ones listed (as many as fit
reasonably) on the actual cover, and continue the rest onto adjacent
pages.

If you publish or distribute Opaque copies of the Document numbering
more than 100, you must either include a machine-readable Transparent
copy along with each Opaque copy, or state in or with each Opaque copy
a computer-network location from which the general network-using
public has access to download using public-standard network protocols
a complete Transparent copy of the Document, free of added material.
If you use the latter option, you must take reasonably prudent steps,
when you begin distribution of Opaque copies in quantity, to ensure
that this Transparent copy will remain thus accessible at the stated
location until at least one year after the last time you distribute an
Opaque copy (directly or through your agents or retailers) of that
edition to the public.

It is requested, but not required, that you contact the authors of the
Document well before redistributing any large number of copies, to give
them a chance to provide you with an updated version of the Document.


\section{MODIFICATIONS}
\label{fdl-section-modifications}

You may copy and distribute a Modified Version of the Document under
the conditions of sections 2 and 3 above, provided that you release
the Modified Version under precisely this License, with the Modified
Version filling the role of the Document, thus licensing distribution
and modification of the Modified Version to whoever possesses a copy
of it.  In addition, you must do these things in the Modified Version:

\begin{itemize}
\item[A.]
   Use in the Title Page (and on the covers, if any) a title distinct
   from that of the Document, and from those of previous versions
   (which should, if there were any, be listed in the History section
   of the Document).  You may use the same title as a previous version
   if the original publisher of that version gives permission.

\item[B.]
   List on the Title Page, as authors, one or more persons or entities
   responsible for authorship of the modifications in the Modified
   Version, together with at least five of the principal authors of the
   Document (all of its principal authors, if it has fewer than five),
   unless they release you from this requirement.

\item[C.]
   State on the Title page the name of the publisher of the
   Modified Version, as the publisher.

\item[D.]
   Preserve all the copyright notices of the Document.

\item[E.]
   Add an appropriate copyright notice for your modifications
   adjacent to the other copyright notices.

\item[F.]
   Include, immediately after the copyright notices, a license notice
   giving the public permission to use the Modified Version under the
   terms of this License, in the form shown in the Addendum below.

\item[G.]
   Preserve in that license notice the full lists of Invariant Sections
   and required Cover Texts given in the Document's license notice.

\item[H.]
   Include an unaltered copy of this License.

\item[I.]
   Preserve the section Entitled "History", Preserve its Title, and add
   to it an item stating at least the title, year, new authors, and
   publisher of the Modified Version as given on the Title Page.  If
   there is no section Entitled "History" in the Document, create one
   stating the title, year, authors, and publisher of the Document as
   given on its Title Page, then add an item describing the Modified
   Version as stated in the previous sentence.

\item[J.]
   Preserve the network location, if any, given in the Document for
   public access to a Transparent copy of the Document, and likewise
   the network locations given in the Document for previous versions
   it was based on.  These may be placed in the "History" section.
   You may omit a network location for a work that was published at
   least four years before the Document itself, or if the original
   publisher of the version it refers to gives permission.

\item[K.]
   For any section Entitled "Acknowledgements" or "Dedications",
   Preserve the Title of the section, and preserve in the section all
   the substance and tone of each of the contributor acknowledgements
   and/or dedications given therein.

\item[L.]
   Preserve all the Invariant Sections of the Document,
   unaltered in their text and in their titles.  Section numbers
   or the equivalent are not considered part of the section titles.

\item[M.]
   Delete any section Entitled "Endorsements".  Such a section
   may not be included in the Modified Version.

\item[N.]
   Do not retitle any existing section to be Entitled "Endorsements"
   or to conflict in title with any Invariant Section.

\item[O.]
   Preserve any Warranty Disclaimers.
\end{itemize}

If the Modified Version includes new front-matter sections or
appendices that qualify as Secondary Sections and contain no material
copied from the Document, you may at your option designate some or all
of these sections as invariant.  To do this, add their titles to the
list of Invariant Sections in the Modified Version's license notice.
These titles must be distinct from any other section titles.

You may add a section Entitled "Endorsements", provided it contains
nothing but endorsements of your Modified Version by various
parties--for example, statements of peer review or that the text has
been approved by an organization as the authoritative definition of a
standard.

You may add a passage of up to five words as a Front-Cover Text, and a
passage of up to 25 words as a Back-Cover Text, to the end of the list
of Cover Texts in the Modified Version.  Only one passage of
Front-Cover Text and one of Back-Cover Text may be added by (or
through arrangements made by) any one entity.  If the Document already
includes a cover text for the same cover, previously added by you or
by arrangement made by the same entity you are acting on behalf of,
you may not add another; but you may replace the old one, on explicit
permission from the previous publisher that added the old one.

The author(s) and publisher(s) of the Document do not by this License
give permission to use their names for publicity for or to assert or
imply endorsement of any Modified Version.


\section{COMBINING DOCUMENTS}
\label{fdl-section-combining-documents}


You may combine the Document with other documents released under this
License, under the terms defined in section 4 above for modified
versions, provided that you include in the combination all of the
Invariant Sections of all of the original documents, unmodified, and
list them all as Invariant Sections of your combined work in its
license notice, and that you preserve all their Warranty Disclaimers.

The combined work need only contain one copy of this License, and
multiple identical Invariant Sections may be replaced with a single
copy.  If there are multiple Invariant Sections with the same name but
different contents, make the title of each such section unique by
adding at the end of it, in parentheses, the name of the original
author or publisher of that section if known, or else a unique number.
Make the same adjustment to the section titles in the list of
Invariant Sections in the license notice of the combined work.

In the combination, you must combine any sections Entitled "History"
in the various original documents, forming one section Entitled
"History"; likewise combine any sections Entitled "Acknowledgements",
and any sections Entitled "Dedications".  You must delete all sections
Entitled "Endorsements".

\section{COLLECTIONS OF DOCUMENTS}
\label{fdl-section-collections-of-documents}

You may make a collection consisting of the Document and other documents
released under this License, and replace the individual copies of this
License in the various documents with a single copy that is included in
the collection, provided that you follow the rules of this License for
verbatim copying of each of the documents in all other respects.

You may extract a single document from such a collection, and distribute
it individually under this License, provided you insert a copy of this
License into the extracted document, and follow this License in all
other respects regarding verbatim copying of that document.


\section{AGGREGATION WITH INDEPENDENT WORKS}
\label{fdl-section-aggregation-with-independent-works}


A compilation of the Document or its derivatives with other separate
and independent documents or works, in or on a volume of a storage or
distribution medium, is called an "aggregate" if the copyright
resulting from the compilation is not used to limit the legal rights
of the compilation's users beyond what the individual works permit.
When the Document is included in an aggregate, this License does not
apply to the other works in the aggregate which are not themselves
derivative works of the Document.

If the Cover Text requirement of section 3 is applicable to these
copies of the Document, then if the Document is less than one half of
the entire aggregate, the Document's Cover Texts may be placed on
covers that bracket the Document within the aggregate, or the
electronic equivalent of covers if the Document is in electronic form.
Otherwise they must appear on printed covers that bracket the whole
aggregate.


\section{TRANSLATION}
\label{fdl-section-translation}


Translation is considered a kind of modification, so you may
distribute translations of the Document under the terms of section 4.
Replacing Invariant Sections with translations requires special
permission from their copyright holders, but you may include
translations of some or all Invariant Sections in addition to the
original versions of these Invariant Sections.  You may include a
translation of this License, and all the license notices in the
Document, and any Warranty Disclaimers, provided that you also include
the original English version of this License and the original versions
of those notices and disclaimers.  In case of a disagreement between
the translation and the original version of this License or a notice
or disclaimer, the original version will prevail.

If a section in the Document is Entitled "Acknowledgements",
"Dedications", or "History", the requirement (section 4) to Preserve
its Title (section 1) will typically require changing the actual
title.


\section{TERMINATION}
\label{fdl-section-termination}


You may not copy, modify, sublicense, or distribute the Document except
as expressly provided for under this License.  Any other attempt to
copy, modify, sublicense or distribute the Document is void, and will
automatically terminate your rights under this License.  However,
parties who have received copies, or rights, from you under this
License will not have their licenses terminated so long as such
parties remain in full compliance.


\section{FUTURE REVISIONS OF THIS LICENSE}
\label{fdl-section-future-revisions-of-this-license}


The Free Software Foundation may publish new, revised versions
of the GNU Free Documentation License from time to time.  Such new
versions will be similar in spirit to the present version, but may
differ in detail to address new problems or concerns.  See
https://www.gnu.org/copyleft/.

Each version of the License is given a distinguishing version number.
If the Document specifies that a particular numbered version of this
License "or any later version" applies to it, you have the option of
following the terms and conditions either of that specified version or
of any later version that has been published (not as a draft) by the
Free Software Foundation.  If the Document does not specify a version
number of this License, you may choose any version ever published (not
as a draft) by the Free Software Foundation.


\section{ADDENDUM: How to use this License for your documents}
\label{fdl-section-addendum-how-to-use-this-license-for-your-documents}

To use this License in a document you have written, include a copy of
the License in the document and put the following copyright and
license notices just after the title page:

\bigskip
\begin{quote}
    Copyright \copyright  YEAR  YOUR NAME.
    Permission is granted to copy, distribute and/or modify this document
    under the terms of the GNU Free Documentation License, Version 1.2
    or any later version published by the Free Software Foundation;
    with no Invariant Sections, no Front-Cover Texts, and no Back-Cover Texts.
    A copy of the license is included in the section entitled "GNU
    Free Documentation License".
\end{quote}
\bigskip

If you have Invariant Sections, Front-Cover Texts and Back-Cover Texts,
replace the "with...Texts." line with this:

\bigskip
\begin{quote}
    with the Invariant Sections being LIST THEIR TITLES, with the
    Front-Cover Texts being LIST, and with the Back-Cover Texts being LIST.
\end{quote}
\bigskip

If you have Invariant Sections without Cover Texts, or some other
combination of the three, merge those two alternatives to suit the
situation.

If your document contains nontrivial examples of program code, we
recommend releasing these examples in parallel under your choice of
free software license, such as the GNU General Public License,
to permit their use in free software.

%---------------------------------------------------------------------

\backmatter
\bibliography{my}
\bibliographystyle{amsalpha}
\end{document}
